% library-book.tex -- GENERATED by (write-library-book-tex! ...) at load.  Do not hand-edit.
% The theorems and definitions of the tree, no proofs, one part per library.
% \vnbstmt -- a displayed statement that shrinks to the text width only
% when it would overflow it, so short statements are set at normal size and
% a wide one never runs into the margin.  Needs graphicx.
%
% The leading \leavevmode is not decoration.  amsthm defers the theorem head
% (`Theorem 1.1 (Title).') into \everypar, so it is injected at the start of
% the first paragraph of the body; without \leavevmode the opening \par is a
% no-op on an empty paragraph, the head lands in the SAME paragraph as the
% box below, and head-width + \linewidth overfills every statement by 300 pt.
% \leavevmode starts a paragraph for the head to attach to, and the \par then
% closes it before the box.
\providecommand{\vnbstmt}[1]{%
  \leavevmode\par\nobreak\smallskip\noindent
  \resizebox{\ifdim\width>\linewidth\linewidth\else\width\fi}{!}{$\displaystyle #1$}%
  \par\smallskip}
\providecommand{\vnbmissing}[1]{\textit{No theorem \texttt{#1} is installed in the library.}}
% \vnbprf -- one paragraph of the proof block.  Set RAGGED RIGHT: these are
% comma-separated lists of \texttt identifiers, which neither hyphenate nor
% break at their hyphens, so justified setting either overfills the measure or
% (under \sloppy) buys the fit with very loose interword glue.
\providecommand{\vnbprf}[1]{{\raggedright #1\par}}
% \vnbpart -- one PART of a statement that the library gives under several names:
% its own paragraph, the name as a run-in label (the user, 2026-09-21).
\providecommand{\vnbpart}[1]{\par\medskip{\raggedright\noindent\textup{\textbf{Part} #1.}\par}\nopagebreak\smallskip}
\providecommand{\vnbproofpart}[1]{\par\medskip\noindent\textbf{Proof of} #1\textbf{.}\par\nopagebreak}
\providecommand{\vnbstructdef}[2]{\par\medskip\noindent\textbf{Structure definition} (\texttt{#1}).\par\nopagebreak#2\par}
\providecommand{\vnbdef}[2]{\par\medskip\noindent\textbf{Definition} (\texttt{#1}).\par\nopagebreak#2\par}
\providecommand{\vnbaxiom}[2]{\par\medskip\noindent\textbf{Axiom} (\texttt{#1}).\par\nopagebreak#2\par}
\providecommand{\vnbassume}[3]{\par\medskip\noindent\textbf{Assumption} (\texttt{#1}; warrant: #2).\par\nopagebreak#3\par}
\providecommand{\vnbprop}[2]{\par\medskip\noindent\textbf{Proposition} (\texttt{#1}).\par\nopagebreak#2\par}
\providecommand{\vnbbookfile}[2]{\section*{\href{file:../#2.scm}{\texttt{#1}}}\addcontentsline{toc}{section}{\texttt{#1}}}
\tableofcontents

\part{Foundations}
\noindent Logic and sets; the number systems NN, ZZ, QQ, RR and CC with their arithmetic and order; tuples and lists; functions, bijections and intervals; finite sums; cardinality and ordinals; choice (Zermelo, Zorn); the combinatorial lemmas. Every other book stands on this one.
\par\noindent{\small 157 files.}

\vnbbookfile{structure-library/operation-properties}{structure-library/operation-properties}
\noindent{\small operation-properties.scm -- named properties of operations.}\par
\par\medskip\noindent\textbf{Definition} (\texttt{is-associative}).\;For all $\mathit{op}$ and $\mathit{crr}$, 
\begin{equation*}
\begin{array}{@{}l@{}} \operatorname{is\text{-}associative}(\mathit{op}, \mathit{crr}) \\ \Leftrightarrow\; \forall u, v, w \in \mathit{crr}.\; (\operatorname{op}(\operatorname{op}(u, v), w) = \operatorname{op}(u, \operatorname{op}(v, w))) \end{array}
\end{equation*}

\par\medskip\noindent\textbf{Definition} (\texttt{is-commutative}).\;For all $\mathit{op}$ and $\mathit{crr}$, 
\begin{equation*}
\operatorname{is\text{-}commutative}(\mathit{op}, \mathit{crr}) \Leftrightarrow \forall u, v \in \mathit{crr}.\; (\operatorname{op}(u, v) = \operatorname{op}(v, u))
\end{equation*}

\par\medskip\noindent\textbf{Definition} (\texttt{is-identity}).\;For all $\mathit{op}$, $\mathit{unit}$ and $\mathit{crr}$, 
\begin{equation*}
\operatorname{is\text{-}identity}(\mathit{op}, \mathit{unit}, \mathit{crr}) \Leftrightarrow \forall u \in \mathit{crr}.\; ((\operatorname{op}(\mathit{unit}, u) = u) \wedge (\operatorname{op}(u, \mathit{unit}) = u))
\end{equation*}

\par\medskip\noindent\textbf{Definition} (\texttt{has-inverses}).\;For all $\mathit{op}$, $\mathit{unit}$, $\mathit{invop}$ and $\mathit{crr}$, 
\begin{equation*}
\begin{array}{@{}l@{}} \operatorname{has\text{-}inverses}(\mathit{op}, \mathit{unit}, \mathit{invop}, \mathit{crr}) \\ \Leftrightarrow\; \forall u \in \mathit{crr}.\; ((\operatorname{op}(\operatorname{invop}(u), u) = \mathit{unit}) \wedge (\operatorname{op}(u, \operatorname{invop}(u)) = \mathit{unit})) \end{array}
\end{equation*}

\par\medskip\noindent\textbf{Definition} (\texttt{is-distributive}).\;For all $\mathit{addop}$, $\mathit{mulop}$ and $\mathit{crr}$, 
\begin{equation*}
\begin{array}{@{}l@{}} \operatorname{is\text{-}distributive}(\mathit{addop}, \mathit{mulop}, \mathit{crr}) \\ \Leftrightarrow\; \forall u, v, w \in \mathit{crr}. \\ \quad \operatorname{mulop}(u, \operatorname{addop}(v, w)) = \operatorname{addop}(\operatorname{mulop}(u, v), \operatorname{mulop}(u, w)) \wedge \\ \quad \operatorname{mulop}(\operatorname{addop}(u, v), w) = \operatorname{addop}(\operatorname{mulop}(u, w), \operatorname{mulop}(v, w)) \end{array}
\end{equation*}

\par\medskip\noindent\textbf{Definition} (\texttt{is-norm}).\;For all $\mathit{nm}$, $\mathit{addop}$, $\mathit{mulop}$, $\mathit{zr}$ and $\mathit{crr}$, 
\begin{equation*}
\begin{array}{@{}l@{}} \operatorname{is\text{-}norm}(\mathit{nm}, \mathit{addop}, \mathit{mulop}, \mathit{zr}, \mathit{crr}) \\ \Leftrightarrow\; \mathit{nm} \in (\mathit{crr} \to \mathbb{R}) \wedge \\ \forall a \in \mathit{crr}. \\ \quad 0 \leq \operatorname{nm}(a) \wedge \\ \quad (\operatorname{nm}(a) = 0) \Leftrightarrow (a = \mathit{zr}) \wedge \\ \quad \forall b \in \mathit{crr}. \\ \quad \quad (\operatorname{nm}(\operatorname{mulop}(a, b)) = \cdot(\operatorname{nm}(a), \operatorname{nm}(b))) \wedge (\operatorname{nm}(\operatorname{addop}(a, b)) \leq +(\operatorname{nm}(a), \operatorname{nm}(b))) \end{array}
\end{equation*}

\par\medskip\noindent\textbf{Definition} (\texttt{is-group-norm}).\;For all $\mathit{nm}$, $\mathit{op}$, $\mathit{invop}$, $\mathit{unit}$ and $\mathit{crr}$, 
\begin{equation*}
\begin{array}{@{}l@{}} \operatorname{is\text{-}group\text{-}norm}(\mathit{nm}, \mathit{op}, \mathit{invop}, \mathit{unit}, \mathit{crr}) \\ \Leftrightarrow\; \mathit{nm} \in (\mathit{crr} \to \mathbb{R}) \wedge \\ \forall u \in \mathit{crr}. \\ \quad 0 \leq \operatorname{nm}(u) \wedge \\ \quad (\operatorname{nm}(u) = 0) \Leftrightarrow (u = \mathit{unit}) \wedge \\ \quad \operatorname{nm}(\operatorname{invop}(u)) = \operatorname{nm}(u) \wedge \\ \quad \forall v \in \mathit{crr}.\; (\operatorname{nm}(\operatorname{op}(u, v)) \leq +(\operatorname{nm}(u), \operatorname{nm}(v))) \end{array}
\end{equation*}

\par\medskip\noindent\textbf{Definition} (\texttt{is-metric}).\;For all $\mathit{dst}$ and $\mathit{crr}$, 
\begin{equation*}
\begin{array}{@{}l@{}} \operatorname{is\text{-}metric}(\mathit{dst}, \mathit{crr}) \\ \Leftrightarrow\; \forall u \in \mathit{crr}. \\ \quad \operatorname{dst}(u, u) = 0 \wedge \\ \quad \forall v \in \mathit{crr}. \\ \quad \quad 0 \leq \operatorname{dst}(u, v) \wedge \\ \quad \quad (\operatorname{dst}(u, v) = 0) \Rightarrow (u = v) \wedge \\ \quad \quad \operatorname{dst}(u, v) = \operatorname{dst}(v, u) \wedge \\ \quad \quad \forall w \in \mathit{crr}.\; (\operatorname{dst}(u, w) \leq +(\operatorname{dst}(u, v), \operatorname{dst}(v, w))) \end{array}
\end{equation*}

\par\medskip\noindent\textbf{Definition} (\texttt{is-pseudometric}).\;For all $\mathit{dst}$ and $\mathit{crr}$, 
\begin{equation*}
\begin{array}{@{}l@{}} \operatorname{is\text{-}pseudometric}(\mathit{dst}, \mathit{crr}) \\ \Leftrightarrow\; \forall u \in \mathit{crr}. \\ \quad \operatorname{dst}(u, u) = 0 \wedge \\ \quad \forall v \in \mathit{crr}. \\ \quad \quad ((0 \leq \operatorname{dst}(u, v)) \wedge ((\operatorname{dst}(u, v) = \operatorname{dst}(v, u)) \wedge \forall w \in \mathit{crr}.\; (\operatorname{dst}(u, w) \leq +(\operatorname{dst}(u, v), \operatorname{dst}(v, w))))) \end{array}
\end{equation*}


\vnbbookfile{theorem-library/axioms}{theorem-library/axioms}
\noindent{\small axioms.scm -- axioms not yet derivable from primitive inferences}\par
\vnbdef{apply-tupling-2}{\vnbstmt{\forall f, a, b.\; (f(a, b) \simeq f({[}a, b]))}}
\vnbdef{apply-tupling-3}{\vnbstmt{\forall f, a, b, c.\; (f(a, b, c) \simeq f({[}a, b, c]))}}
\vnbaxiom{equality-transitivity}{\vnbstmt{\forall a, b, c.\; (((a = b) \wedge (b = c)) \Rightarrow (a = c))}}
\vnbaxiom{eq-subst-membership}{\vnbstmt{\forall a, b, s.\; (((a = b) \wedge (b \in s)) \Rightarrow (a \in s))}}
\vnbaxiom{quasi-eq-reflexivity}{\vnbstmt{\forall a.\; (a \simeq a)}}
\vnbaxiom{fun-apply-type}{\vnbstmt{\forall f, a, b, x.\; (((f \in (a \to b)) \wedge (x \in a)) \Rightarrow (f(x) \in b))}}
\vnbaxiom{quasi-eq-symmetry}{\vnbstmt{\forall a, b.\; ((a \simeq b) \Rightarrow (b \simeq a))}}
\vnbaxiom{quasi-eq-subst-membership}{\vnbstmt{\forall a, b, s.\; (((a \simeq b) \wedge (b \in s)) \Rightarrow (a \in s))}}
\vnbaxiom{quasi-eq-transitivity}{\vnbstmt{\forall a, b, c.\; (((a \simeq b) \wedge (b \simeq c)) \Rightarrow (a \simeq c))}}
\vnbaxiom{equality-symmetry}{\vnbstmt{\forall a, b.\; ((a = b) \Rightarrow (b = a))}}
\vnbaxiom{list-sethood}{\vnbstmt{\forall a, l.\; (((a \in \mathbf{Set}) \wedge (l \in \operatorname{tuples}(a))) \Rightarrow (l \in \mathbf{Set}))}}
\vnbaxiom{quasi-eq-def}{\vnbstmt{\forall a, b.\; ((a = a) \Rightarrow ((a \simeq b) \Leftrightarrow (a = b)))}}

\vnbbookfile{structure-library/set-vocabulary}{structure-library/set-vocabulary}
\noindent{\small set-vocabulary.scm -- two set constructors that the rest of the library is STATED with, and that nothing but the kernel's own constructors is needed to define.}\par
\par\medskip\noindent\textbf{Definition} (\texttt{difference}).\;$\operatorname{difference}(\operatorname{u\_}, \operatorname{b\_}) :=$ $\operatorname{complement\text{-}in}(\operatorname{u\_}, \operatorname{b\_})$
\par\medskip\noindent\textbf{Definition} (\texttt{singleton}).\;$\operatorname{singleton}(\operatorname{y\_}) :=$ $\operatorname{make\text{-}set}({[}\operatorname{y\_}])$

\vnbbookfile{structure-library/intersection-of}{structure-library/intersection-of}
\noindent{\small intersection-of.scm -- INTERSECTION-OF(c), the intersection of the members of a family c, as a DEFINED constructor.}\par
\par\medskip\noindent\textbf{Definition} (\texttt{intersection-of}).\;$\operatorname{intersection\text{-}of}(c) :=$
\begin{equation*}
\{\operatorname{iov\_} \in \operatorname{big\text{-}union}(\operatorname{iow\_}, c, \operatorname{iow\_}) : \forall \operatorname{iou\_} \in c.\; (\operatorname{iov\_} \in \operatorname{iou\_})\}
\end{equation*}


\vnbbookfile{structure-library/compose}{structure-library/compose}
\noindent{\small compose.scm -- COMPOSE: composition of functions, f o g.}\par
\par\medskip\noindent\textbf{Definition} (\texttt{compose}).\;$\operatorname{compose}(f, g) :=$ $(\lambda \operatorname{z\_} \in \operatorname{dom}(g).\; f(g(\operatorname{z\_})))$

\vnbbookfile{structure-library/id-fun}{structure-library/id-fun}
\noindent{\small id-fun.scm -- ID-FUN: the identity map of a set.}\par
\par\medskip\noindent\textbf{Definition} (\texttt{id-fun}).\;$\operatorname{id\text{-}fun}(\operatorname{a\_}) :=$ $(\lambda \operatorname{idfx\_} \in \operatorname{a\_}.\; \operatorname{idfx\_})$

\vnbbookfile{number-systems}{number-systems}
\noindent{\small number-systems.scm -- number system constants and axioms}\par
\par\medskip\noindent\textbf{Definition} (\texttt{rr-upper-bound}).\;For all $s$ and $b$, 
\begin{equation*}
\operatorname{rr\text{-}upper\text{-}bound}(s, b) \Leftrightarrow ((b \in \mathbb{R}) \wedge \forall x \in s.\; (x \leq b))
\end{equation*}

\par\medskip\noindent\textbf{Definition} (\texttt{rr-bounded-above}).\;For every $s$,  $\operatorname{rr\text{-}bounded\text{-}above}(s) \Leftrightarrow \exists b.\; \operatorname{rr\text{-}upper\text{-}bound}(s, b)$
\vnbdef{binary-minus-def}{\vnbstmt{\forall a, b.\; ((a - b) \simeq (a + -b))}}
\vnbdef{binary-divide-def}{\vnbstmt{\forall a, b.\; (/(a, b) \simeq (a \cdot \operatorname{recip}(b)))}}
\vnbdef{rr-abs-def}{\vnbstmt{\begin{array}{@{}l@{}} \forall x \in \mathbb{R}. \\ \quad ((0 \leq x) \Rightarrow (\lvert x \rvert = x)) \wedge (\neg (0 \leq x) \Rightarrow (\lvert x \rvert = -x)) \end{array}}}
\vnbdef{rr-max-def}{\vnbstmt{\begin{array}{@{}l@{}} \forall x, y \in \mathbb{R}. \\ \quad ((y \leq x) \Rightarrow (\operatorname{max}(x, y) = x)) \wedge (\neg (y \leq x) \Rightarrow (\operatorname{max}(x, y) = y)) \end{array}}}
\vnbdef{rr-min-def}{\vnbstmt{\begin{array}{@{}l@{}} \forall x, y \in \mathbb{R}. \\ \quad ((x \leq y) \Rightarrow (\operatorname{min}(x, y) = x)) \wedge (\neg (x \leq y) \Rightarrow (\operatorname{min}(x, y) = y)) \end{array}}}
\vnbdef{real-part-def}{\vnbstmt{\forall z.\; (\operatorname{real\text{-}part}(z) \simeq ((z + \operatorname{conjugate}(z)) \cdot \operatorname{recip}(2)))}}
\vnbdef{imag-part-def}{\vnbstmt{\forall z.\; (\operatorname{imag\text{-}part}(z) \simeq ((z - \operatorname{conjugate}(z)) \cdot \operatorname{recip}(+2i)))}}
\vnbaxiom{zz-one-in}{\vnbstmt{1 \in \mathbb{Z}}}
\vnbaxiom{rr-subset-cc}{\vnbstmt{\forall n \in \mathbb{R}.\; (n \in \mathbb{C})}}
\vnbaxiom{qq-recip-closed}{\vnbstmt{\forall a \in \mathbb{Q}.\; (\neg (a = 0) \Rightarrow (\operatorname{recip}(a) \in \mathbb{Q}))}}
\vnbaxiom{qq-is-set}{\vnbstmt{\mathbb{Q} \in \mathbf{Set}}}
\vnbaxiom{zz-mul-comm}{\vnbstmt{\forall a, b.\; (((a \in \mathbb{Z}) \wedge (b \in \mathbb{Z})) \Rightarrow ((a \cdot b) = (b \cdot a)))}}
\vnbaxiom{qq-add-comm}{\vnbstmt{\forall a, b.\; (((a \in \mathbb{Q}) \wedge (b \in \mathbb{Q})) \Rightarrow ((a + b) = (b + a)))}}
\vnbaxiom{rr-is-set}{\vnbstmt{\mathbb{R} \in \mathbf{Set}}}
\vnbaxiom{cc-mul-comm}{\vnbstmt{\forall a, b.\; (((a \in \mathbb{C}) \wedge (b \in \mathbb{C})) \Rightarrow ((a \cdot b) = (b \cdot a)))}}
\vnbaxiom{cc-conjugate-closed}{\vnbstmt{\forall a \in \mathbb{C}.\; (\operatorname{conjugate}(a) \in \mathbb{C})}}
\vnbaxiom{rr-add-zero}{\vnbstmt{\forall a \in \mathbb{R}.\; (a = a)}}
\vnbaxiom{rr-recip-inverse}{\vnbstmt{\forall a \in \mathbb{R}.\; (\neg (a = 0) \Rightarrow ((a \cdot \operatorname{recip}(a)) = 1))}}
\vnbaxiom{cc-conjugate-fixes-rr}{\vnbstmt{\forall a \in \mathbb{R}.\; (\operatorname{conjugate}(a) = a)}}
\vnbaxiom{zz-is-set}{\vnbstmt{\mathbb{Z} \in \mathbf{Set}}}
\vnbaxiom{cc-conjugate-mul}{\vnbstmt{\begin{array}{@{}l@{}} \forall a, b. \\ \quad ((a \in \mathbb{C}) \wedge (b \in \mathbb{C})) \Rightarrow (\operatorname{conjugate}((a \cdot b)) = (\operatorname{conjugate}(a) \cdot \operatorname{conjugate}(b))) \end{array}}}
\vnbaxiom{nn-succ-closed}{\vnbstmt{\forall n \in \mathbb{N}.\; ((n + 1) \in \mathbb{N})}}
\vnbaxiom{cc-conjugate-involution}{\vnbstmt{\forall a \in \mathbb{C}.\; (\operatorname{conjugate}(\operatorname{conjugate}(a)) = a)}}
\vnbaxiom{cc-one-in}{\vnbstmt{1 \in \mathbb{C}}}
\vnbaxiom{rr-recip-closed}{\vnbstmt{\forall a \in \mathbb{R}.\; (\neg (a = 0) \Rightarrow (\operatorname{recip}(a) \in \mathbb{R}))}}
\vnbaxiom{qq-add-zero}{\vnbstmt{\forall a \in \mathbb{Q}.\; (a = a)}}
\vnbaxiom{qq-subset-rr}{\vnbstmt{\forall n \in \mathbb{Q}.\; (n \in \mathbb{R})}}
\vnbaxiom{rr-leq-transitive}{\vnbstmt{\begin{array}{@{}l@{}} \forall a, b, c. \\ \quad (((a \in \mathbb{R}) \wedge ((b \in \mathbb{R}) \wedge (c \in \mathbb{R}))) \wedge ((a \leq b) \wedge (b \leq c))) \Rightarrow (a \leq c) \end{array}}}
\vnbaxiom{rr-leq-reflexive}{\vnbstmt{\forall a \in \mathbb{R}.\; (a \leq a)}}
\vnbaxiom{nn-mul-comm}{\vnbstmt{\forall a, b.\; (((a \in \mathbb{N}) \wedge (b \in \mathbb{N})) \Rightarrow ((a \cdot b) = (b \cdot a)))}}
\vnbaxiom{power-succ}{\vnbstmt{\forall x, n.\; (((x \in \mathbb{C}) \wedge (n \in \mathbb{N})) \Rightarrow (\operatorname{power}(x, (n + 1)) = (x \cdot \operatorname{power}(x, n))))}}
\vnbaxiom{qq-add-closed}{\vnbstmt{\forall a, b.\; (((a \in \mathbb{Q}) \wedge (b \in \mathbb{Q})) \Rightarrow ((a + b) \in \mathbb{Q}))}}
\vnbaxiom{nn-mul-closed}{\vnbstmt{\forall a, b.\; (((a \in \mathbb{N}) \wedge (b \in \mathbb{N})) \Rightarrow ((a \cdot b) \in \mathbb{N}))}}
\vnbaxiom{zz-add-closed}{\vnbstmt{\forall a, b.\; (((a \in \mathbb{Z}) \wedge (b \in \mathbb{Z})) \Rightarrow ((a + b) \in \mathbb{Z}))}}
\vnbaxiom{rr-add-closed}{\vnbstmt{\forall a, b.\; (((a \in \mathbb{R}) \wedge (b \in \mathbb{R})) \Rightarrow ((a + b) \in \mathbb{R}))}}
\vnbaxiom{cc-is-set}{\vnbstmt{\mathbb{C} \in \mathbf{Set}}}
\vnbaxiom{cc-zero-in}{\vnbstmt{0 \in \mathbb{C}}}
\vnbaxiom{cc-add-comm}{\vnbstmt{\forall a, b.\; (((a \in \mathbb{C}) \wedge (b \in \mathbb{C})) \Rightarrow ((a + b) = (b + a)))}}
\vnbaxiom{rr-zero-in}{\vnbstmt{0 \in \mathbb{R}}}
\vnbaxiom{nn-add-zero}{\vnbstmt{\forall a \in \mathbb{N}.\; (a = a)}}
\vnbaxiom{cc-add-zero}{\vnbstmt{\forall a \in \mathbb{C}.\; (a = a)}}
\vnbaxiom{rr-leq-add-compat}{\vnbstmt{\begin{array}{@{}l@{}} \forall a, b, c. \\ \quad (((a \in \mathbb{R}) \wedge ((b \in \mathbb{R}) \wedge (c \in \mathbb{R}))) \wedge (a \leq b)) \Rightarrow ((a + c) \leq (b + c)) \end{array}}}
\vnbaxiom{cc-generated-by-rr}{\vnbstmt{\forall z \in \mathbb{C}.\; \exists x, y \in \mathbb{R}.\; (z = (x + (y \cdot +i)))}}
\vnbaxiom{rr-neg-inverse}{\vnbstmt{\forall a \in \mathbb{R}.\; ((a + -a) = 0)}}
\vnbaxiom{zz-distributive}{\vnbstmt{\begin{array}{@{}l@{}} \forall a, b, c. \\ \quad ((a \in \mathbb{Z}) \wedge ((b \in \mathbb{Z}) \wedge (c \in \mathbb{Z}))) \Rightarrow ((a \cdot (b + c)) = ((a \cdot b) + (a \cdot c))) \end{array}}}
\vnbaxiom{qq-is-fraction}{\vnbstmt{\forall q \in \mathbb{Q}.\; \exists a, b \in \mathbb{Z}.\; (\neg (b = 0) \wedge (q = (a \cdot \operatorname{recip}(b))))}}
\vnbaxiom{qq-recip-inverse}{\vnbstmt{\forall a \in \mathbb{Q}.\; (\neg (a = 0) \Rightarrow ((a \cdot \operatorname{recip}(a)) = 1))}}
\vnbaxiom{cc-mul-assoc}{\vnbstmt{\begin{array}{@{}l@{}} \forall a, b, c. \\ \quad ((a \in \mathbb{C}) \wedge ((b \in \mathbb{C}) \wedge (c \in \mathbb{C}))) \Rightarrow (((a \cdot b) \cdot c) = (a \cdot (b \cdot c))) \end{array}}}
\vnbaxiom{nn-add-assoc}{\vnbstmt{\begin{array}{@{}l@{}} \forall a, b, c. \\ \quad ((a \in \mathbb{N}) \wedge ((b \in \mathbb{N}) \wedge (c \in \mathbb{N}))) \Rightarrow (((a + b) + c) = (a + (b + c))) \end{array}}}
\vnbaxiom{zz-one-mul}{\vnbstmt{\forall a \in \mathbb{Z}.\; (a = a)}}
\vnbaxiom{power-neg}{\vnbstmt{\begin{array}{@{}l@{}} \forall x, n. \\ \quad ((x \in \mathbb{C}) \wedge ((n \in \mathbb{N}) \wedge \neg (x = 0))) \Rightarrow (\operatorname{power}(x, -n) = \operatorname{recip}(\operatorname{power}(x, n))) \end{array}}}
\vnbaxiom{zz-mul-assoc}{\vnbstmt{\begin{array}{@{}l@{}} \forall a, b, c. \\ \quad ((a \in \mathbb{Z}) \wedge ((b \in \mathbb{Z}) \wedge (c \in \mathbb{Z}))) \Rightarrow (((a \cdot b) \cdot c) = (a \cdot (b \cdot c))) \end{array}}}
\vnbaxiom{qq-mul-comm}{\vnbstmt{\forall a, b.\; (((a \in \mathbb{Q}) \wedge (b \in \mathbb{Q})) \Rightarrow ((a \cdot b) = (b \cdot a)))}}
\vnbaxiom{cc-recip-closed}{\vnbstmt{\forall a \in \mathbb{C}.\; (\neg (a = 0) \Rightarrow (\operatorname{recip}(a) \in \mathbb{C}))}}
\vnbaxiom{nn-induction}{\vnbstmt{\begin{array}{@{}l@{}} \forall c.; \forall n \in \mathbb{N}. \\ \quad ((0 \in c) \wedge \forall n \in \mathbb{N}.\; ((n \in c) \Rightarrow ((n + 1) \in c))) \Rightarrow (n \in c) \end{array}}}
\vnbaxiom{qq-distributive}{\vnbstmt{\begin{array}{@{}l@{}} \forall a, b, c. \\ \quad ((a \in \mathbb{Q}) \wedge ((b \in \mathbb{Q}) \wedge (c \in \mathbb{Q}))) \Rightarrow ((a \cdot (b + c)) = ((a \cdot b) + (a \cdot c))) \end{array}}}
\vnbaxiom{nn-is-set}{\vnbstmt{\mathbb{N} \in \mathbf{Set}}}
\vnbaxiom{rr-mul-comm}{\vnbstmt{\forall a, b.\; (((a \in \mathbb{R}) \wedge (b \in \mathbb{R})) \Rightarrow ((a \cdot b) = (b \cdot a)))}}
\vnbaxiom{cc-recip-inverse}{\vnbstmt{\forall a \in \mathbb{C}.\; (\neg (a = 0) \Rightarrow ((a \cdot \operatorname{recip}(a)) = 1))}}
\vnbaxiom{zz-add-comm}{\vnbstmt{\forall a, b.\; (((a \in \mathbb{Z}) \wedge (b \in \mathbb{Z})) \Rightarrow ((a + b) = (b + a)))}}
\vnbaxiom{qq-add-assoc}{\vnbstmt{\begin{array}{@{}l@{}} \forall a, b, c. \\ \quad ((a \in \mathbb{Q}) \wedge ((b \in \mathbb{Q}) \wedge (c \in \mathbb{Q}))) \Rightarrow (((a + b) + c) = (a + (b + c))) \end{array}}}
\vnbaxiom{zz-generated-by-nn}{\vnbstmt{\forall a \in \mathbb{Z}.\; \exists n \in \mathbb{N}.\; ((a = n) \vee (a = -n))}}
\vnbaxiom{nn-zero-in}{\vnbstmt{0 \in \mathbb{N}}}
\vnbaxiom{qq-one-in}{\vnbstmt{1 \in \mathbb{Q}}}
\vnbaxiom{qq-neg-inverse}{\vnbstmt{\forall a \in \mathbb{Q}.\; ((a + -a) = 0)}}
\vnbaxiom{zz-mul-closed}{\vnbstmt{\forall a, b.\; (((a \in \mathbb{Z}) \wedge (b \in \mathbb{Z})) \Rightarrow ((a \cdot b) \in \mathbb{Z}))}}
\vnbaxiom{cc-self-conj-nonneg}{\vnbstmt{\forall a \in \mathbb{C}.\; (((a \cdot \operatorname{conjugate}(a)) \in \mathbb{R}) \wedge (0 \leq (a \cdot \operatorname{conjugate}(a))))}}
\vnbaxiom{rr-neg-closed}{\vnbstmt{\forall a \in \mathbb{R}.\; (-a \in \mathbb{R})}}
\vnbaxiom{cc-add-assoc}{\vnbstmt{\begin{array}{@{}l@{}} \forall a, b, c. \\ \quad ((a \in \mathbb{C}) \wedge ((b \in \mathbb{C}) \wedge (c \in \mathbb{C}))) \Rightarrow (((a + b) + c) = (a + (b + c))) \end{array}}}
\vnbaxiom{qq-zero-in}{\vnbstmt{0 \in \mathbb{Q}}}
\vnbaxiom{cc-conjugate-i}{\vnbstmt{\operatorname{conjugate}(+i) = -+i}}
\vnbaxiom{nn-add-comm}{\vnbstmt{\forall a, b.\; (((a \in \mathbb{N}) \wedge (b \in \mathbb{N})) \Rightarrow ((a + b) = (b + a)))}}
\vnbaxiom{rr-sup-upper}{\vnbstmt{\begin{array}{@{}l@{}} \forall s. \\ \quad ((s \subseteq \mathbb{R}) \wedge (\exists x.\; (x \in s) \wedge \operatorname{rr\text{-}bounded\text{-}above}(s))) \Rightarrow  \\ \quad \quad \operatorname{rr\text{-}upper\text{-}bound}(s, \operatorname{sup}(s)) \end{array}}}
\vnbaxiom{rr-add-assoc}{\vnbstmt{\begin{array}{@{}l@{}} \forall a, b, c. \\ \quad ((a \in \mathbb{R}) \wedge ((b \in \mathbb{R}) \wedge (c \in \mathbb{R}))) \Rightarrow (((a + b) + c) = (a + (b + c))) \end{array}}}
\vnbaxiom{cc-conjugate-add}{\vnbstmt{\begin{array}{@{}l@{}} \forall a, b. \\ \quad ((a \in \mathbb{C}) \wedge (b \in \mathbb{C})) \Rightarrow (\operatorname{conjugate}((a + b)) = (\operatorname{conjugate}(a) + \operatorname{conjugate}(b))) \end{array}}}
\vnbaxiom{cc-neg-closed}{\vnbstmt{\forall a \in \mathbb{C}.\; (-a \in \mathbb{C})}}
\vnbaxiom{rr-add-comm}{\vnbstmt{\forall a, b.\; (((a \in \mathbb{R}) \wedge (b \in \mathbb{R})) \Rightarrow ((a + b) = (b + a)))}}
\vnbaxiom{rr-one-in}{\vnbstmt{1 \in \mathbb{R}}}
\vnbaxiom{rr-mul-assoc}{\vnbstmt{\begin{array}{@{}l@{}} \forall a, b, c. \\ \quad ((a \in \mathbb{R}) \wedge ((b \in \mathbb{R}) \wedge (c \in \mathbb{R}))) \Rightarrow (((a \cdot b) \cdot c) = (a \cdot (b \cdot c))) \end{array}}}
\vnbaxiom{power-zero}{\vnbstmt{\forall x \in \mathbb{C}.\; (\operatorname{power}(x, 0) = 1)}}
\vnbaxiom{rr-distributive}{\vnbstmt{\begin{array}{@{}l@{}} \forall a, b, c. \\ \quad ((a \in \mathbb{R}) \wedge ((b \in \mathbb{R}) \wedge (c \in \mathbb{R}))) \Rightarrow ((a \cdot (b + c)) = ((a \cdot b) + (a \cdot c))) \end{array}}}
\vnbaxiom{nn-subset-zz}{\vnbstmt{\forall n \in \mathbb{N}.\; (n \in \mathbb{Z})}}
\vnbaxiom{cc-neg-inverse}{\vnbstmt{\forall a \in \mathbb{C}.\; ((a + -a) = 0)}}
\vnbaxiom{cc-mul-closed}{\vnbstmt{\forall a, b.\; (((a \in \mathbb{C}) \wedge (b \in \mathbb{C})) \Rightarrow ((a \cdot b) \in \mathbb{C}))}}
\vnbaxiom{cc-add-closed}{\vnbstmt{\forall a, b.\; (((a \in \mathbb{C}) \wedge (b \in \mathbb{C})) \Rightarrow ((a + b) \in \mathbb{C}))}}
\vnbaxiom{zz-neg-inverse}{\vnbstmt{\forall a \in \mathbb{Z}.\; ((a + -a) = 0)}}
\vnbaxiom{nn-one-mul}{\vnbstmt{\forall a \in \mathbb{N}.\; (a = a)}}
\vnbaxiom{qq-one-mul}{\vnbstmt{\forall a \in \mathbb{Q}.\; (a = a)}}
\vnbaxiom{rr-sup-least}{\vnbstmt{\begin{array}{@{}l@{}} \forall s, b. \\ \quad ((s \subseteq \mathbb{R}) \wedge (\exists x.\; (x \in s) \wedge (\operatorname{rr\text{-}bounded\text{-}above}(s) \wedge \operatorname{rr\text{-}upper\text{-}bound}(s, b)))) \Rightarrow  \\ \quad \quad \operatorname{sup}(s) \leq b \end{array}}}
\vnbaxiom{zz-subset-qq}{\vnbstmt{\forall n \in \mathbb{Z}.\; (n \in \mathbb{Q})}}
\vnbaxiom{nn-distributive}{\vnbstmt{\begin{array}{@{}l@{}} \forall a, b, c. \\ \quad ((a \in \mathbb{N}) \wedge ((b \in \mathbb{N}) \wedge (c \in \mathbb{N}))) \Rightarrow ((a \cdot (b + c)) = ((a \cdot b) + (a \cdot c))) \end{array}}}
\vnbaxiom{cc-i-squared}{\vnbstmt{(+i \cdot +i) = -1}}
\vnbaxiom{rr-sup-in}{\vnbstmt{\begin{array}{@{}l@{}} \forall s. \\ \quad ((s \subseteq \mathbb{R}) \wedge (\exists x.\; (x \in s) \wedge \operatorname{rr\text{-}bounded\text{-}above}(s))) \Rightarrow  \\ \quad \quad \operatorname{sup}(s) \in \mathbb{R} \end{array}}}
\vnbaxiom{nn-add-closed}{\vnbstmt{\forall a, b.\; (((a \in \mathbb{N}) \wedge (b \in \mathbb{N})) \Rightarrow ((a + b) \in \mathbb{N}))}}
\vnbaxiom{qq-neg-closed}{\vnbstmt{\forall a \in \mathbb{Q}.\; (-a \in \mathbb{Q})}}
\vnbaxiom{rr-one-mul}{\vnbstmt{\forall a \in \mathbb{R}.\; (a = a)}}
\vnbaxiom{nn-mul-assoc}{\vnbstmt{\begin{array}{@{}l@{}} \forall a, b, c. \\ \quad ((a \in \mathbb{N}) \wedge ((b \in \mathbb{N}) \wedge (c \in \mathbb{N}))) \Rightarrow (((a \cdot b) \cdot c) = (a \cdot (b \cdot c))) \end{array}}}
\vnbaxiom{zz-zero-in}{\vnbstmt{0 \in \mathbb{Z}}}
\vnbaxiom{cc-i-in}{\vnbstmt{+i \in \mathbb{C}}}
\vnbaxiom{rr-leq-mul-nonneg}{\vnbstmt{\forall a, b.\; ((((a \in \mathbb{R}) \wedge (b \in \mathbb{R})) \wedge ((0 \leq a) \wedge (0 \leq b))) \Rightarrow (0 \leq (a \cdot b)))}}
\vnbaxiom{zz-neg-closed}{\vnbstmt{\forall a \in \mathbb{Z}.\; (-a \in \mathbb{Z})}}
\vnbaxiom{qq-mul-assoc}{\vnbstmt{\begin{array}{@{}l@{}} \forall a, b, c. \\ \quad ((a \in \mathbb{Q}) \wedge ((b \in \mathbb{Q}) \wedge (c \in \mathbb{Q}))) \Rightarrow (((a \cdot b) \cdot c) = (a \cdot (b \cdot c))) \end{array}}}
\vnbaxiom{rr-leq-antisymmetric}{\vnbstmt{\forall a, b.\; ((((a \in \mathbb{R}) \wedge (b \in \mathbb{R})) \wedge ((a \leq b) \wedge (b \leq a))) \Rightarrow (a = b))}}
\vnbaxiom{cc-self-conj-real}{\vnbstmt{\forall a \in \mathbb{C}.\; ((a \cdot \operatorname{conjugate}(a)) \in \mathbb{R})}}
\vnbaxiom{rr-leq-total}{\vnbstmt{\forall a, b.\; (((a \in \mathbb{R}) \wedge (b \in \mathbb{R})) \Rightarrow ((a \leq b) \vee (b \leq a)))}}
\vnbaxiom{zz-add-assoc}{\vnbstmt{\begin{array}{@{}l@{}} \forall a, b, c. \\ \quad ((a \in \mathbb{Z}) \wedge ((b \in \mathbb{Z}) \wedge (c \in \mathbb{Z}))) \Rightarrow (((a + b) + c) = (a + (b + c))) \end{array}}}
\vnbaxiom{qq-mul-closed}{\vnbstmt{\forall a, b.\; (((a \in \mathbb{Q}) \wedge (b \in \mathbb{Q})) \Rightarrow ((a \cdot b) \in \mathbb{Q}))}}
\vnbaxiom{cc-distributive}{\vnbstmt{\begin{array}{@{}l@{}} \forall a, b, c. \\ \quad ((a \in \mathbb{C}) \wedge ((b \in \mathbb{C}) \wedge (c \in \mathbb{C}))) \Rightarrow ((a \cdot (b + c)) = ((a \cdot b) + (a \cdot c))) \end{array}}}
\vnbaxiom{rr-mul-closed}{\vnbstmt{\forall a, b.\; (((a \in \mathbb{R}) \wedge (b \in \mathbb{R})) \Rightarrow ((a \cdot b) \in \mathbb{R}))}}
\vnbaxiom{power-typing-nonneg}{\vnbstmt{\forall x, n.\; (((x \in \mathbb{C}) \wedge (n \in \mathbb{N})) \Rightarrow (\operatorname{power}(x, n) \in \mathbb{C}))}}
\vnbaxiom{cc-one-mul}{\vnbstmt{\forall a \in \mathbb{C}.\; (a = a)}}
\vnbaxiom{zz-add-zero}{\vnbstmt{\forall a \in \mathbb{Z}.\; (a = a)}}

\vnbbookfile{structure-library/nn-arith}{structure-library/nn-arith}
\noindent{\small nn-arith.scm -- the recursion equations for + and * on NN.}\par
\vnbdef{nn-add-succ}{\vnbstmt{\forall a, b \in \mathbb{N}.\; ((a + (b + 1)) = ((a + b) + 1))}}

\vnbbookfile{structure-library/list-recursion}{structure-library/list-recursion}
\noindent{\small list-recursion.scm -- CONS, and the recursive characterization of LENGTH.}\par
\vnbdef{empty-in-tuples}{\vnbstmt{\forall a.\; ({[}] \in \operatorname{tuples}(a))}}
\vnbdef{cons-in-tuples}{\vnbstmt{\forall a, x, l.\; (((x \in a) \wedge (l \in \operatorname{tuples}(a))) \Rightarrow (\operatorname{cons}(x, l) \in \operatorname{tuples}(a)))}}
\vnbdef{nth-cons-1}{\vnbstmt{\forall x \in \mathbf{Set}, l.\; (\operatorname{nth}(1, \operatorname{cons}(x, l)) = x)}}
\vnbdef{nth-cons-succ}{\vnbstmt{\begin{array}{@{}l@{}} \forall i, x, l. \\ \quad ((x \in \mathbf{Set}) \wedge ((i \in \mathbb{N}) \wedge ((1 \leq i) \wedge (i \leq \operatorname{length}(l))))) \Rightarrow  \\ \quad \quad \operatorname{nth}((i + 1), \operatorname{cons}(x, l)) = \operatorname{nth}(i, l) \end{array}}}
\vnbdef{length-cons}{\vnbstmt{\forall x \in \mathbf{Set}, l.\; (\operatorname{length}(\operatorname{cons}(x, l)) = (\operatorname{length}(l) + 1))}}
\vnbdef{makeset-cons}{\vnbstmt{\begin{array}{@{}l@{}} \forall x \in \mathbf{Set}, l. \\ \quad \operatorname{make\text{-}set}(\operatorname{cons}(x, l)) = (\{x\} \cup \operatorname{make\text{-}set}(l)) \end{array}}}
\vnbaxiom{nn-not-a-list}{\vnbstmt{\forall n \in \mathbb{N}.\; \neg (n \in \operatorname{tuples}(\mathbf{Set}))}}
\vnbaxiom{tuple-length-zero}{\vnbstmt{\forall a.; \forall l \in \operatorname{tuples}(a).\; ((\operatorname{length}(l) = 0) \Rightarrow (l = {[}]))}}
\vnbaxiom{tuple-cons-decompose}{\vnbstmt{\begin{array}{@{}l@{}} \forall a, n, l. \\ \quad ((n \in \mathbb{N}) \wedge ((l \in \operatorname{tuples}(a)) \wedge (\operatorname{length}(l) = (n + 1)))) \Rightarrow  \\ \quad \quad \exists x \in a, m \in \operatorname{tuples}(a).\; ((\operatorname{length}(m) = n) \wedge (l = \operatorname{cons}(x, m))) \end{array}}}

\vnbbookfile{structure-library/sqn}{structure-library/sqn}
\noindent{\small sqn.scm -- SQN(a): the sequences in a, i.e. FUN(NN, a).}\par
\par\medskip\noindent\textbf{Definition} (\texttt{sqn}).\;$a^{\mathbb{N}} :=$ $(\mathbb{N} \to a)$

\vnbbookfile{structure-library/order-predicates}{structure-library/order-predicates}
\noindent{\small order-predicates.scm -- derived order predicates on the kernel $<$=.}\par
\par\medskip\noindent\textbf{Definition} (\texttt{<}).\;For all $x$ and $y$,  $(x < y) \Leftrightarrow ((x \leq y) \wedge \neg (x = y))$
\par\medskip\noindent\textbf{Definition} (\texttt{pos-rr}).\;For every $r$,  $\operatorname{pos\text{-}rr}(r) \Leftrightarrow ((r \in \mathbb{R}) \wedge ((0 \leq r) \wedge \neg (0 = r)))$

\vnbbookfile{structure-library/nn-minus}{structure-library/nn-minus}
\noindent{\small nn-minus.scm -- truncated natural subtraction (monus), as VOCABULARY.}\par
\par\medskip\noindent\textbf{Definition} (\texttt{nn-minus}).\;Suppose $k \in \mathbb{N}$ and $l \in \mathbb{N}$. Then $\operatorname{nn\text{-}minus}(k, l) = \left(\text{if }(l \leq k)\text{ then }-(k, l)\text{ else }0\right)$

\vnbbookfile{structure-library/order-lemmas}{structure-library/order-lemmas}
\noindent{\small RETIRED 2026-09-17 (proven): nn-minus-succ-1, one-in-interval, succ-not-one, nn-minus-1-inj, nn-succ-le-cancel -- theorem-library/rake-intervals.scm, modulo 0.  succ-in-interval and pred-in-interval were FALSE as stated (the interval's upper bound was not typed: q = 3/2 with succ(3/2) interpreted as 0 refutes the ...}\par
\vnbaxiom{recip-defined-nonzero}{\vnbstmt{\forall q.\; ((\operatorname{recip}(q) = \operatorname{recip}(q)) \Rightarrow \neg (q = 0))}}

\vnbbookfile{structure-library/real-powers}{structure-library/real-powers}
\noindent{\small real-powers.scm -- the real power RPOW and the square root SQRT.}\par
\par\medskip\noindent\textbf{Definition} (\texttt{sqrt}).\;$\sqrt{\operatorname{a\_}} :=$ $(\iota\, \operatorname{x\_}.\;((\operatorname{x\_} \in \mathbb{R}) \wedge ((0 \leq \operatorname{x\_}) \wedge ((\operatorname{x\_} \cdot \operatorname{x\_}) = \operatorname{a\_}))))$

\vnbbookfile{structure-library/complex}{structure-library/complex}
\noindent{\small complex.scm -- CC as a metric space; completeness axiom}\par
\par\medskip\noindent\textbf{Definition} (\texttt{cc-ms}).\;
\begin{equation*}
\operatorname{cc\text{-}ms} = {[}\mathbb{C}, (\lambda {[}x, y] \in (\mathbb{C} \times \mathbb{C}).\; \operatorname{magnitude}(-(x, y)))]
\end{equation*}

\vnbdef{magnitude-def}{\vnbstmt{\begin{array}{@{}l@{}} \forall z. \\ \quad (\operatorname{magnitude}(z) \simeq \sqrt{((\operatorname{real\text{-}part}(z) \cdot \operatorname{real\text{-}part}(z)) + (\operatorname{imag\text{-}part}(z) \cdot \operatorname{imag\text{-}part}(z)))}) \end{array}}\par\noindent{\small\textit{Comment.} The right-hand side is defined iff $z$ is a complex number, so $\operatorname{magnitude}(z)$ is defined iff $z$ is a complex number.}\par
}
\vnbdef{cc-ms@pts}{\vnbstmt{(\operatorname{pts}(\operatorname{cc\text{-}ms}) \simeq \mathbb{C})}}
\vnbdef{cc-ms@dist}{\vnbstmt{(\operatorname{dist}(\operatorname{cc\text{-}ms}) \simeq (\lambda {[}x, y] \in (\mathbb{C} \times \mathbb{C}).\; \operatorname{magnitude}((x - y))))}}

\vnbbookfile{structure-library/reduce}{structure-library/reduce}
\noindent{\small reduce.scm -- REDUCE: left-fold of a binary op over a finite indexed family.  Bridges the kiddie n-ary surface form (+ x y z ...) and the adult finite-sum operator SUM-AG.}\par
\par\medskip\noindent\textbf{Definition} (\texttt{reduce}).\;For all $\mathit{op}$ and $f$,  $(\operatorname{reduce}(\mathit{op}, f, 1) \simeq f(0))$\par
Let $\mathit{op}$ and $f$ be arbitrary. Let $n \text{ be s.t. } ((n \in \mathbb{N}) \wedge (1 \leq n))$. Then $(\operatorname{reduce}(\mathit{op}, f, \operatorname{succ}(n)) \simeq \operatorname{op}(\operatorname{reduce}(\mathit{op}, f, n), f(n)))$
\par\medskip\noindent\textbf{Definition} (\texttt{fam-of-list}).\;Let $l$ be arbitrary. Let $i \text{ be s.t. } ((i \in \mathbb{N}) \wedge (\operatorname{succ}(i) \leq \operatorname{length}(l)))$. Then $\operatorname{fam\text{-}of\text{-}list}(l)(i) = \operatorname{nth}(\operatorname{succ}(i), l)$

\vnbbookfile{structure-library/ordinals}{structure-library/ordinals}
\noindent{\small ordinals.scm -- ordinal numbers: axioms, ordering, induction, recursion}\par
\vnbdef{ord-lt-iff}{\vnbstmt{\forall a, b.\; (\operatorname{ord\text{-}lt}(a, b) \Leftrightarrow (\operatorname{ord\text{-}le}(a, b) \wedge \neg (a = b)))}}
\vnbdef{limit-ord-iff}{\vnbstmt{\begin{array}{@{}l@{}} \forall a. \\ \quad \operatorname{limit\text{-}ord}(a) \Leftrightarrow ((a \in \mathrm{Ord}) \wedge (\neg (a = 0) \wedge \neg \exists b \in \mathrm{Ord}.\; (a = (b + 1)))) \end{array}}}
\vnbaxiom{ord-segment-membership}{\vnbstmt{\forall a \in \mathrm{Ord}, x.\; ((x \in \mathrm{OS}(a)) \Leftrightarrow \operatorname{ord\text{-}lt}(x, a))}}
\vnbaxiom{sup-ord-least}{\vnbstmt{\begin{array}{@{}l@{}} \forall a \in \mathbf{Set}, b \in \mathrm{Ord}. \\ \quad (\forall x \in a.\; (x \in \mathrm{Ord}) \wedge \forall c \in a.\; \operatorname{ord\text{-}le}(c, b)) \Rightarrow  \\ \quad \quad \operatorname{ord\text{-}le}(\operatorname{sup\text{-}ord}(a), b) \end{array}}}
\vnbaxiom{ord-succ-injective}{\vnbstmt{\begin{array}{@{}l@{}} \forall a, b. \\ \quad ((a \in \mathrm{Ord}) \wedge ((b \in \mathrm{Ord}) \wedge ((a + 1) = (b + 1)))) \Rightarrow (a = b) \end{array}}}
\vnbaxiom{transfinite-induction}{\vnbstmt{\begin{array}{@{}l@{}} \forall c.; \forall a \in \mathrm{Ord}. \\ \quad \forall a \in \mathrm{Ord}.\; (\forall b.\; (\operatorname{ord\text{-}lt}(b, a) \Rightarrow (b \in c)) \Rightarrow (a \in c)) \Rightarrow  \\ \quad \quad a \in c \end{array}}}
\vnbaxiom{ord-succ-in}{\vnbstmt{\forall a \in \mathrm{Ord}.\; ((a + 1) \in \mathrm{Ord})}}
\vnbaxiom{ord-zero-least}{\vnbstmt{\forall a \in \mathrm{Ord}.\; \operatorname{ord\text{-}le}(0, a)}}
\vnbaxiom{ord-segment-zero}{\vnbstmt{\mathrm{OS}(0) = \emptyset}}
\vnbaxiom{ord-segment-is-set}{\vnbstmt{\forall a \in \mathrm{Ord}.\; (\mathrm{OS}(a) \in \mathbf{Set})}}
\vnbaxiom{sup-ord-in}{\vnbstmt{\forall a \in \mathbf{Set}.\; (\forall x \in a.\; (x \in \mathrm{Ord}) \Rightarrow (\operatorname{sup\text{-}ord}(a) \in \mathrm{Ord}))}}
\vnbaxiom{sup-ord-empty}{\vnbstmt{\operatorname{sup\text{-}ord}(\emptyset) = 0}}
\vnbaxiom{limit-ord-is-sup}{\vnbstmt{\forall a.\; (\operatorname{limit\text{-}ord}(a) \Rightarrow (\operatorname{sup\text{-}ord}(\mathrm{OS}(a)) = a))}}
\vnbaxiom{ord-le-total}{\vnbstmt{\forall a, b.\; (((a \in \mathrm{Ord}) \wedge (b \in \mathrm{Ord})) \Rightarrow (\operatorname{ord\text{-}le}(a, b) \vee \operatorname{ord\text{-}le}(b, a)))}}
\vnbaxiom{ord-le-antisymm}{\vnbstmt{\forall a, b.\; ((\operatorname{ord\text{-}le}(a, b) \wedge \operatorname{ord\text{-}le}(b, a)) \Rightarrow (a = b))}}
\vnbaxiom{ord-segment-limit}{\vnbstmt{\begin{array}{@{}l@{}} \forall a, x. \\ \quad (\operatorname{limit\text{-}ord}(a) \wedge \operatorname{ord\text{-}lt}(x, a)) \Rightarrow  \\ \quad \quad \exists b.\; (\operatorname{ord\text{-}lt}(x, b) \wedge \operatorname{ord\text{-}lt}(b, a)) \end{array}}}
\vnbaxiom{ord-succ-nn}{\vnbstmt{\forall n \in \mathbb{N}.\; ((n + 1) = (n + 1))}}
\vnbaxiom{ord-segment-succ}{\vnbstmt{\begin{array}{@{}l@{}} \forall a \in \mathrm{Ord}, x. \\ \quad (x \in \mathrm{OS}((a + 1))) \Leftrightarrow ((x \in \mathrm{OS}(a)) \vee (x = a)) \end{array}}}
\vnbaxiom{ord-le-trans}{\vnbstmt{\forall a, b, c.\; ((\operatorname{ord\text{-}le}(a, b) \wedge \operatorname{ord\text{-}le}(b, c)) \Rightarrow \operatorname{ord\text{-}le}(a, c))}}
\vnbaxiom{ord-le-refl}{\vnbstmt{\forall a \in \mathrm{Ord}.\; \operatorname{ord\text{-}le}(a, a)}}
\vnbaxiom{ord-le-closure}{\vnbstmt{\forall a, b.\; (\operatorname{ord\text{-}le}(a, b) \Rightarrow ((a \in \mathrm{Ord}) \wedge (b \in \mathrm{Ord})))}}
\vnbaxiom{sup-ord-upper}{\vnbstmt{\begin{array}{@{}l@{}} \forall a \in \mathbf{Set}, b \in a. \\ \quad \forall x \in a.\; (x \in \mathrm{Ord}) \Rightarrow \operatorname{ord\text{-}le}(b, \operatorname{sup\text{-}ord}(a)) \end{array}}}
\vnbaxiom{ord-succ-above}{\vnbstmt{\forall a \in \mathrm{Ord}.\; \operatorname{ord\text{-}lt}(a, (a + 1))}}
\vnbaxiom{ord-le-nn-compat}{\vnbstmt{\forall m, n.\; (((m \in \mathbb{N}) \wedge (n \in \mathbb{N})) \Rightarrow (\operatorname{ord\text{-}le}(m, n) \Leftrightarrow (m \leq n)))}}
\vnbaxiom{burali-forti}{\vnbstmt{\neg (\mathrm{Ord} \in \mathbf{Set})}}
\vnbaxiom{sup-ord-succ-segment}{\vnbstmt{\forall a \in \mathrm{Ord}.\; (\operatorname{sup\text{-}ord}(\mathrm{OS}((a + 1))) = a)}}
\vnbaxiom{nn-subset-ord}{\vnbstmt{\forall n \in \mathbb{N}.\; (n \in \mathrm{Ord})}}
\vnbaxiom{ord-succ-immediate}{\vnbstmt{\begin{array}{@{}l@{}} \forall a, b. \\ \quad ((a \in \mathrm{Ord}) \wedge ((b \in \mathrm{Ord}) \wedge \operatorname{ord\text{-}lt}(a, b))) \Rightarrow \operatorname{ord\text{-}le}((a + 1), b) \end{array}}}

\vnbbookfile{structure-library/bijection}{structure-library/bijection}
\noindent{\small RETIRED 2026-09-17 (proven): inverse-bij-in-fun, inverse-bij-left, inverse-bij-right, inverse-bij-is-bijection -- were UNWARRANTED axioms; theorem-library/rake-inverse-bij.scm, modulo 0. RETIRED 2026-09-17 (proven): bijection-compose (was an `informal' axiom) -- theorem-library/rake-analysis2.scm, via the new ...}\par
\par\medskip\noindent\textbf{Definition} (\texttt{inverse-bij}).\;$\varphi^{-1} :=$ $(\lambda \operatorname{y\_} \in y.\; \varepsilon(\{\operatorname{x\_} \in x : (\varphi(\operatorname{x\_}) = \operatorname{y\_})\}))$
\vnbdef{bijection-membership-iff}{\vnbstmt{\begin{array}{@{}l@{}} \forall x, y, i. \\ \quad i \in \mathrm{Bij}(x, y) \\ \quad \Leftrightarrow\; i \in (x \to y) \wedge \\ \quad \forall a, b \in x.\; ((i(a) = i(b)) \Rightarrow (a = b)) \wedge \\ \quad \forall w \in y.\; \exists z \in x.\; (i(z) = w) \end{array}}}

\vnbbookfile{structure-library/cardinality}{structure-library/cardinality}
\noindent{\small cardinality.scm -- CARD, the cardinality of a set, DEFINED.}\par
\par\medskip\noindent\textbf{Definition} (\texttt{card}).\;$|\operatorname{a\_}| :=$
\begin{equation*}
(\iota\, \alpha.\;((\alpha \in \mathrm{Ord}) \wedge (\exists \varphi.\; (\varphi \in \mathrm{Bij}(\operatorname{a\_}, \mathrm{OS}(\alpha))) \wedge \forall \beta.\; (\operatorname{ord\text{-}lt}(\beta, \alpha) \Rightarrow \neg \exists \psi.\; (\psi \in \mathrm{Bij}(\operatorname{a\_}, \mathrm{OS}(\beta)))))))
\end{equation*}


\vnbbookfile{structure-library/injection}{structure-library/injection}
\noindent{\small RETIRED 2026-09-17 (proven): choose-n-0, choose-0-succ -- theorem-library/rake-combinatorics.scm. choose-succ (Pascal) stays: needs the CHOOSE-SET split + card-image-injection (\textasciitilde{}150 lines). RETIRED 2026-09-17 (proven): falling-in-nn -- theorem-library/rake-finsum-typing-nn.scm. choose-in-nn stays: it bottoms out in the ...}\par
\par\medskip\noindent\textbf{Definition} (\texttt{choose-set}).\;$\operatorname{choose\text{-}set}(n, m) :=$ $\{a \in \operatorname{power}(\mathrm{OS}(n)) : (|a| = m)\}$
\par\medskip\noindent\textbf{Definition} (\texttt{injective-star}).\;For every $f$, 
\begin{equation*}
\operatorname{injective\text{-}star}(f) \Leftrightarrow \forall \operatorname{u\_}, \operatorname{v\_}.\; ((f(\operatorname{u\_}) = f(\operatorname{v\_})) \Rightarrow (\operatorname{u\_} = \operatorname{v\_}))
\end{equation*}

\par\medskip\noindent\textbf{Definition} (\texttt{permutations}).\;$\operatorname{permutations}(n) :=$ $\operatorname{injection}(\mathrm{OS}(n), \mathrm{OS}(n))$
\par\medskip\noindent\textbf{Definition} (\texttt{factorial}).\; $\operatorname{factorial}(0) = 1$\par
Suppose $k \in \mathbb{N}$. Then $\operatorname{factorial}(\operatorname{succ}(k)) = \cdot(\operatorname{succ}(k), \operatorname{factorial}(k))$
\par\medskip\noindent\textbf{Definition} (\texttt{choose}).\;$\operatorname{choose}(n, m) :=$ $|\operatorname{choose\text{-}set}(n, m)|$
\par\medskip\noindent\textbf{Definition} (\texttt{falling}).\;For every $n$,  $(\operatorname{falling}(n, 0) \simeq 1)$\par
Let $n$ be arbitrary. Let $k \in \mathbb{N}$. Then $(\operatorname{falling}(n, \operatorname{succ}(k)) \simeq \cdot(-(n, k), \operatorname{falling}(n, k)))$
\vnbdef{injection-membership-iff}{\vnbstmt{\begin{array}{@{}l@{}} \forall x, y, i. \\ \quad (i \in \operatorname{injection}(x, y)) \Leftrightarrow ((i \in (x \to y)) \wedge \forall a, b \in x.\; ((i(a) = i(b)) \Rightarrow (a = b))) \end{array}}}
\vnbdef{image-membership-iff}{\vnbstmt{\forall i, s, w.\; ((w \in \operatorname{image}(i, s)) \Leftrightarrow ((w \in \mathbf{Set}) \wedge \exists x \in s.\; (i(x) = w)))}}
\vnbaxiom{image-set}{\vnbstmt{\forall i.; \forall s \in \mathbf{Set}.\; (\operatorname{image}(i, s) \in \mathbf{Set})}}

\vnbbookfile{structure-library/compose-typing}{structure-library/compose-typing}
\noindent{\small RETIRED 2026-09-17 (proven): range-membership, fun-range-membership, ran-subset-codomain, compose-type-2/3/4/5 -- theorem-library/rake-compose-typing.scm, all modulo 0 from dom-membership / dom-fun-membership / fun-codomain-iff / image-membership-iff.  They had been `proof'-warranted supports here; the}\par
\par\medskip\noindent\textbf{Definition} (\texttt{ran}).\;$\operatorname{ran}(f) :=$ $\operatorname{image}(f, \operatorname{dom}(f))$

\vnbbookfile{structure-library/inf-subsets}{structure-library/inf-subsets}
\noindent{\small inf-subsets.scm -- the class of infinite subsets of a given set.}\par
\par\medskip\noindent\textbf{Definition} (\texttt{inf-subsets}).\;$\mathrm{Inf}(a) :=$ $\{s \in \operatorname{power}(a) : \neg (|s| \in \mathbb{N})\}$
\vnbdef{inf-subsets-membership}{\vnbstmt{\begin{array}{@{}l@{}} \forall a, s. \\ \quad (s \in \mathrm{Inf}(a)) \Leftrightarrow ((s \in \mathbf{Set}) \wedge ((s \subseteq a) \wedge \neg (|s| \in \mathbb{N}))) \end{array}}}

\vnbbookfile{structure-library/finsum}{structure-library/finsum}
\noindent{\small finsum.scm -- FINSUM: summation of an abelian-group-valued function over an arbitrary finite set.}\par
\par\medskip\noindent\textbf{Definition} (\texttt{finsum}).\;$\operatorname{finsum}(\mathit{ag}, f, s) :=$ $\operatorname{sum\text{-}ag}(\mathit{ag}, \operatorname{enum\text{-}fam}(\mathit{ag}, f, \operatorname{fin\text{-}enum}(s), |s|), |s|)$
\par\medskip\noindent\textbf{Definition} (\texttt{fin-enum}).\;$\operatorname{fin\text{-}enum}(s) :=$ $\varepsilon(\mathrm{Bij}(\mathrm{OS}(|s|), s))$
\par\medskip\noindent\textbf{Definition} (\texttt{enum-fam}).\;$\operatorname{enum\text{-}fam}(\mathit{ag}, f, \varphi, n) :=$
\begin{equation*}
(\lambda i \in \mathbb{N}.\; \left(\text{if }(i \in \mathrm{OS}(n))\text{ then }f(\varphi(i))\text{ else }\operatorname{iden}(\mathit{ag})\right))
\end{equation*}

\par\medskip\noindent\textbf{Definition} (\texttt{insert-last}).\;$\operatorname{insert\text{-}last}(\varphi, x, n) :=$ $(\lambda i \in \mathbb{N}.\; \left(\text{if }(i \in \mathrm{OS}(n))\text{ then }\varphi(i)\text{ else }x\right))$
\par\medskip\noindent\textbf{Definition} (\texttt{sum-set}).\;$\operatorname{sum\text{-}set}(r, s, f) :=$ $\operatorname{finsum}(\operatorname{ring\text{-}additive\text{-}ag}(r), f, s)$

\vnbbookfile{structure-library/matrix}{structure-library/matrix}
\noindent{\small RETIRED 2026-09-17 (proven): finsum-two-support -- theorem-library/rake-finsum-laws2.scm RETIRED 2026-09-17 (proven): identmat-left-identity, identmat-right-identity -- theorem-library/rake-identmat.scm (matmul-entry + finsum-single-support + extensionality; no guard needed). RETIRED 2026-09-17 (proven): ...}\par
\par\medskip\noindent\textbf{Definition} (\texttt{size}).\;$\operatorname{size}(m) :=$ ${[}\operatorname{length}(m), \left(\text{if }(\operatorname{length}(m) = 0)\text{ then }0\text{ else }\operatorname{length}(\operatorname{nth}(1, m))\right)]$
\par\medskip\noindent\textbf{Definition} (\texttt{entry}).\;$\operatorname{entry}(m, i, j) :=$ $\operatorname{nth}(j, \operatorname{nth}(i, m))$
\par\medskip\noindent\textbf{Definition} (\texttt{mat}).\;$\operatorname{mat}(m, n, x) :=$
\begin{equation*}
\{p \in \operatorname{matrix}(x) : ((n \in \mathbb{N}) \wedge (((\operatorname{length}(p) = 0) \Rightarrow (m = 0)) \wedge (\neg (\operatorname{length}(p) = 0) \Rightarrow (\operatorname{size}(p) = {[}m, n]))))\}
\end{equation*}

\par\medskip\noindent\textbf{Definition} (\texttt{identmat}).\;$\operatorname{identmat}(a, n) :=$
\begin{equation*}
\begin{array}{@{}l@{}} \operatorname{matof}(n, n, \\ \quad (\lambda {[}i, j] \in (\operatorname{interval}(1, n) \times \operatorname{interval}(1, n)).\; \left(\text{if }(i = j)\text{ then }\operatorname{one}(a)\text{ else }\operatorname{zero}(a)\right))) \end{array}
\end{equation*}

\par\medskip\noindent\textbf{Definition} (\texttt{interval}).\;$\operatorname{interval}(a, b) :=$ $\{i \in \mathbb{N} : ((a \leq i) \wedge (i \leq b))\}$
\par\medskip\noindent\textbf{Definition} (\texttt{matof}).\;$\operatorname{matof}(m, n, g) :=$
\begin{equation*}
(\iota\, p.\;((p \in \operatorname{mat}(m, n, \operatorname{image}(g, (\operatorname{interval}(1, m) \times \operatorname{interval}(1, n))))) \wedge \forall i \in \operatorname{interval}(1, m), j \in \operatorname{interval}(1, n).\; (\operatorname{entry}(p, i, j) = g(i, j))))
\end{equation*}

\par\medskip\noindent\textbf{Definition} (\texttt{matmul}).\;$\operatorname{matmul}(a, p, q) :=$
\begin{equation*}
\begin{array}{@{}l@{}} \operatorname{matof}(\operatorname{nth}(1, \operatorname{size}(p)), \operatorname{nth}(2, \operatorname{size}(q)), \\ \quad (\lambda {[}i, k] \in (\operatorname{interval}(1, \operatorname{nth}(1, \operatorname{size}(p))) \times \operatorname{interval}(1, \operatorname{nth}(2, \operatorname{size}(q)))).\; \\ \quad \quad \operatorname{finsum}(\operatorname{ring\text{-}additive\text{-}ag}(a), \\ \quad \quad \quad (\lambda j \in \operatorname{interval}(1, \operatorname{nth}(2, \operatorname{size}(p))).\; \operatorname{mul}(a)(\operatorname{entry}(p, i, j), \operatorname{entry}(q, j, k))), \\ \quad \quad \quad \operatorname{interval}(1, \operatorname{nth}(2, \operatorname{size}(p)))))) \end{array}
\end{equation*}

\par\medskip\noindent\textbf{Definition} (\texttt{mat-ring}).\;$\operatorname{mat\text{-}ring}(a, n) :=$
\begin{equation*}
{[}\operatorname{mat}(n, n, \operatorname{carr}(a)), (\lambda {[}p, q] \in (\operatorname{mat}(n, n, \operatorname{carr}(a)) \times \operatorname{mat}(n, n, \operatorname{carr}(a))).\; \operatorname{matadd}(a, p, q)), (\lambda {[}p, q] \in (\operatorname{mat}(n, n, \operatorname{carr}(a)) \times \operatorname{mat}(n, n, \operatorname{carr}(a))).\; \operatorname{matmul}(a, p, q)), (\lambda p \in \operatorname{mat}(n, n, \operatorname{carr}(a)).\; \operatorname{matneg}(a, p)), \operatorname{zeromat}(a, n, n), \operatorname{identmat}(a, n)]
\end{equation*}

\par\medskip\noindent\textbf{Definition} (\texttt{zeromat}).\;$\operatorname{zeromat}(a, m, n) :=$
\begin{equation*}
\operatorname{matof}(m, n, (\lambda {[}i, j] \in (\operatorname{interval}(1, m) \times \operatorname{interval}(1, n)).\; \operatorname{zero}(a)))
\end{equation*}

\par\medskip\noindent\textbf{Definition} (\texttt{matadd}).\;$\operatorname{matadd}(a, p, q) :=$
\begin{equation*}
\begin{array}{@{}l@{}} \operatorname{matof}(\operatorname{nth}(1, \operatorname{size}(p)), \operatorname{nth}(2, \operatorname{size}(p)), \\ \quad (\lambda {[}i, j] \in (\operatorname{interval}(1, \operatorname{nth}(1, \operatorname{size}(p))) \times \operatorname{interval}(1, \operatorname{nth}(2, \operatorname{size}(p)))).\; \\ \quad \quad \operatorname{add}(a)(\operatorname{entry}(p, i, j), \operatorname{entry}(q, i, j)))) \end{array}
\end{equation*}

\par\medskip\noindent\textbf{Definition} (\texttt{matneg}).\;$\operatorname{matneg}(a, p) :=$
\begin{equation*}
\begin{array}{@{}l@{}} \operatorname{matof}(\operatorname{nth}(1, \operatorname{size}(p)), \operatorname{nth}(2, \operatorname{size}(p)), \\ \quad (\lambda {[}i, j] \in (\operatorname{interval}(1, \operatorname{nth}(1, \operatorname{size}(p))) \times \operatorname{interval}(1, \operatorname{nth}(2, \operatorname{size}(p)))).\; \\ \quad \quad \operatorname{neg}(a)(\operatorname{entry}(p, i, j)))) \end{array}
\end{equation*}

\par\medskip\noindent\textbf{Definition} (\texttt{matscale}).\;$\operatorname{matscale}(a, r, p) :=$
\begin{equation*}
\begin{array}{@{}l@{}} \operatorname{matof}(\operatorname{nth}(1, \operatorname{size}(p)), \operatorname{nth}(2, \operatorname{size}(p)), \\ \quad (\lambda {[}i, j] \in (\operatorname{interval}(1, \operatorname{nth}(1, \operatorname{size}(p))) \times \operatorname{interval}(1, \operatorname{nth}(2, \operatorname{size}(p)))).\; \\ \quad \quad \operatorname{mul}(a)(r, \operatorname{entry}(p, i, j)))) \end{array}
\end{equation*}

\par\medskip\noindent\textbf{Definition} (\texttt{unitrow}).\;$\operatorname{unitrow}(a, n, i) :=$
\begin{equation*}
\begin{array}{@{}l@{}} \operatorname{matof}(1, n, \\ \quad (\lambda {[}\mathit{rw}, \mathit{cl}] \in (\operatorname{interval}(1, 1) \times \operatorname{interval}(1, n)).\; \left(\text{if }(\mathit{cl} = i)\text{ then }\operatorname{one}(a)\text{ else }\operatorname{zero}(a)\right))) \end{array}
\end{equation*}

\par\medskip\noindent\textbf{Definition} (\texttt{block}).\;$\operatorname{block}(p, k, l) :=$
\begin{equation*}
\operatorname{matof}(k, l, (\lambda {[}i, j] \in (\operatorname{interval}(1, k) \times \operatorname{interval}(1, l)).\; \operatorname{entry}(p, i, j)))
\end{equation*}

\par\medskip\noindent\textbf{Definition} (\texttt{snoc-col}).\;$\operatorname{snoc\text{-}col}(w, n, x) :=$
\begin{equation*}
\begin{array}{@{}l@{}} \operatorname{matof}((n + 1), 1, \\ \quad (\lambda {[}\operatorname{i\_}, \operatorname{j\_}] \in (\operatorname{interval}(1, (n + 1)) \times \operatorname{interval}(1, 1)).\; \\ \quad \quad \left(\text{if }(\operatorname{i\_} = (n + 1))\text{ then }x\text{ else }\operatorname{entry}(w, \operatorname{i\_}, 1)\right))) \end{array}
\end{equation*}

\par\medskip\noindent\textbf{Definition} (\texttt{snoc-row}).\;$\operatorname{snoc\text{-}row}(c, n, r) :=$
\begin{equation*}
\begin{array}{@{}l@{}} \operatorname{matof}(1, (n + 1), \\ \quad (\lambda {[}\operatorname{i\_}, \operatorname{j\_}] \in (\operatorname{interval}(1, 1) \times \operatorname{interval}(1, (n + 1))).\; \\ \quad \quad \left(\text{if }(\operatorname{j\_} = (n + 1))\text{ then }r\text{ else }\operatorname{entry}(c, 1, \operatorname{j\_})\right))) \end{array}
\end{equation*}

\vnbdef{matrix-membership}{\vnbstmt{\begin{array}{@{}l@{}} \forall s, m. \\ \quad m \in \operatorname{matrix}(s) \\ \quad \Leftrightarrow\; m \in \operatorname{tuples}(\operatorname{tuples}(s)) \wedge \\ \quad \forall i, j \in \mathbb{N}. \\ \quad \quad (((1 \leq i) \wedge (i \leq \operatorname{length}(m))) \wedge ((1 \leq j) \wedge (j \leq \operatorname{length}(m)))) \Rightarrow  \\ \quad \quad \quad \operatorname{length}(\operatorname{nth}(i, m)) = \operatorname{length}(\operatorname{nth}(j, m)) \end{array}}}

\vnbbookfile{theorem-library/sep-congruence}{theorem-library/sep-congruence}
\noindent{\small sep-congruence.scm -- wave of 2026-11-09/10 (agent CNVS2), wired 2026-11-10 from scratchpad/cnvs2-1109/cnvs2-sep.scm. Library: core.}\par
\vnbprop{sep-congruence-ptwise-eq}{\vnbstmt{\begin{array}{@{}l@{}} \forall a, p, q. \\ \quad \forall x \in a.\; ((x \in p) \Leftrightarrow (x \in q)) \Rightarrow (\{x \in a : (x \in p)\} = \{x \in a : (x \in q)\}) \end{array}}}
\vnbprop{sep-congruence-ptwise}{\vnbstmt{\begin{array}{@{}l@{}} \forall a, p, q. \\ \quad \forall x \in a.\; ((x \in p) \Leftrightarrow (x \in q)) \Rightarrow (\{x \in a : (x \in p)\} \simeq \{x \in a : (x \in q)\}) \end{array}}}

\vnbbookfile{theorem-library/unique-choice}{theorem-library/unique-choice}
\noindent{\small unique-choice.scm -- wave of 2026-11-09/10 (agent CCR-1109-D4), wired 2026-11-10 from scratchpad/ccr1109d4/unique-choice.scm. Library: core.}\par
\vnbprop{uniq-choice-fun}{\vnbstmt{\begin{array}{@{}l@{}} \forall a \in \mathbf{Set}, c, q. \\ \quad \forall x \in a.\; \exists y \in c.\; ((y \in q(x)) \wedge \forall z \in c.\; ((z \in q(x)) \Rightarrow (z = y))) \Rightarrow  \\ \quad \quad \exists f \in (a \to c).\; \forall x \in a.\; (f(x) \in q(x)) \end{array}}}
\vnbprop{uniq-choice-fun-2}{\vnbstmt{\begin{array}{@{}l@{}} \forall a, b \in \mathbf{Set}, c, q. \\ \quad \forall x \in a, w \in b. \exists y \in c. \\ \quad \quad (y \in q(x, w)) \wedge \forall z \in c.\; ((z \in q(x, w)) \Rightarrow (z = y)) \Rightarrow \\ \quad \quad \exists f \in ((a \times b) \to c). \forall x \in a, w \in b. \\ \quad \quad \quad f(x, w) \in q(x, w) \end{array}}}

\vnbbookfile{theorem-library/ord-well-ordered-proof}{theorem-library/ord-well-ordered-proof}
\noindent{\small ord-well-ordered-proof.scm -- the well-ordering of ORD, PROVEN.}\par
\vnbprop{ord-well-ordered}{\vnbstmt{\begin{array}{@{}l@{}} \forall l. \\ \quad (\forall x \in l.\; (x \in \mathrm{Ord}) \wedge \exists w.\; (w \in l)) \Rightarrow  \\ \quad \quad \exists m \in l.\; \forall k \in l.\; \operatorname{ord\text{-}le}(m, k) \end{array}}}

\vnbbookfile{theorem-library/equality-basics}{theorem-library/equality-basics}
\noindent{\small equality-basics.scm -- eq-sym, eq-trans, neq-sym, PROVEN.}\par
\vnbprop{eq-sym}{\vnbstmt{\forall a, b.\; ((a = b) \Rightarrow (b = a))}}
\vnbprop{eq-trans}{\vnbstmt{\forall a, b, c.\; (((a = b) \wedge (b = c)) \Rightarrow (a = c))}}
\vnbprop{neq-sym}{\vnbstmt{\forall a, b.\; (\neg (a = b) \Rightarrow \neg (b = a))}}

\vnbbookfile{theorem-library/factorial-in-nn}{theorem-library/factorial-in-nn}
\noindent{\small factorial-in-nn.scm -- n! is a natural number, PROVEN. Batch 27-A follow-on, 2026-09-24.}\par
\vnbprop{factorial-in-nn}{\vnbstmt{\forall n \in \mathbb{N}.\; (\operatorname{factorial}(n) \in \mathbb{N})}}

\vnbbookfile{theorem-library/rake-compose-typing}{theorem-library/rake-compose-typing}
\noindent{\small rake-compose-typing.scm -- the seven RAN / composition typing supports of structure-library/compose-typing.scm, PROVEN modulo 0.}\par
\vnbprop{range-membership}{\vnbstmt{\forall f.; \forall a \in \operatorname{dom}(f).\; (f(a) \in \operatorname{ran}(f))}}
\vnbprop{fun-range-membership}{\vnbstmt{\forall x, f, a.\; (((f \in \mathrm{Fun}(x)) \wedge (a \in x)) \Rightarrow (f(a) \in \operatorname{ran}(f)))}}
\vnbprop{ran-subset-codomain}{\vnbstmt{\forall x, y.; \forall f \in (x \to y).\; (\operatorname{ran}(f) \subseteq y)}}
\vnbprop{compose-type-2}{\vnbstmt{\begin{array}{@{}l@{}} \forall x, f, g, a. \\ \quad ((f \in \mathrm{Fun}(x)) \wedge ((g \in \mathrm{Fun}(\operatorname{ran}(f))) \wedge (a \in x))) \Rightarrow (g(f(a)) \in \operatorname{ran}(g)) \end{array}}}
\vnbprop{compose-type-3}{\vnbstmt{\begin{array}{@{}l@{}} \forall x, f, g, h, a. \\ \quad ((f \in \mathrm{Fun}(x)) \wedge ((g \in \mathrm{Fun}(\operatorname{ran}(f))) \wedge ((h \in \mathrm{Fun}(\operatorname{ran}(g))) \wedge (a \in x)))) \Rightarrow  \\ \quad \quad h(g(f(a))) \in \operatorname{ran}(h) \end{array}}}
\vnbprop{compose-type-4}{\vnbstmt{\begin{array}{@{}l@{}} \forall x, f, g, h, k, a. \\ \quad ((f \in \mathrm{Fun}(x)) \wedge ((g \in \mathrm{Fun}(\operatorname{ran}(f))) \wedge ((h \in \mathrm{Fun}(\operatorname{ran}(g))) \wedge ((k \in \mathrm{Fun}(\operatorname{ran}(h))) \wedge (a \in x))))) \Rightarrow  \\ \quad \quad k(h(g(f(a)))) \in \operatorname{ran}(k) \end{array}}}
\vnbprop{compose-type-5}{\vnbstmt{\begin{array}{@{}l@{}} \forall x, f, g, h, k, m, a. \\ \quad f \in \mathrm{Fun}(x) \wedge \\ \quad g \in \mathrm{Fun}(\operatorname{ran}(f)) \wedge \\ \quad h \in \mathrm{Fun}(\operatorname{ran}(g)) \wedge \\ \quad k \in \mathrm{Fun}(\operatorname{ran}(h)) \wedge \\ \quad m \in \mathrm{Fun}(\operatorname{ran}(k)) \wedge \\ \quad a \in x \Rightarrow \\ \quad \quad m(k(h(g(f(a))))) \in \operatorname{ran}(m) \end{array}}}

\vnbbookfile{theorem-library/co-eq-lt-trans}{theorem-library/co-eq-lt-trans}
\noindent{\small co-eq-lt-trans.scm -- co-eq-lt-trans and co-lt-eq-trans, PROVEN.}\par
\vnbprop{co-eq-lt-trans}{\vnbstmt{\forall a, b, c.\; (((a = b) \wedge (b < c)) \Rightarrow (a < c))}}
\vnbprop{co-lt-eq-trans}{\vnbstmt{\forall a, b, c.\; (((a < b) \wedge (b = c)) \Rightarrow (a < c))}}

\vnbbookfile{theorem-library/interval-basics}{theorem-library/interval-basics}
\noindent{\small interval-basics.scm -- INTERVAL's read-offs, PROVEN from its definition.}\par
\vnbprop{interval-unfold}{\vnbstmt{\forall a, b.\; (\{i \in \mathbb{N} : ((a \leq i) \wedge (i \leq b))\} \simeq \operatorname{interval}(a, b))}}
\vnbprop{interval-in-set}{\vnbstmt{\forall a, b.\; (\operatorname{interval}(a, b) \in \mathbf{Set})}}
\vnbprop{interval-elt-in-nn}{\vnbstmt{\forall a, b.; \forall i \in \operatorname{interval}(a, b).\; (i \in \mathbb{N})}}
\vnbprop{interval-lo}{\vnbstmt{\forall a, b.; \forall i \in \operatorname{interval}(a, b).\; (a \leq i)}}
\vnbprop{interval-hi}{\vnbstmt{\forall a, b.; \forall i \in \operatorname{interval}(a, b).\; (i \leq b)}}

\vnbbookfile{theorem-library/interval-mem-intro}{theorem-library/interval-mem-intro}
\noindent{\small interval-mem-intro.scm -- membership INTRODUCTION for INTERVAL, proven.}\par
\vnbprop{interval-mem-intro}{\vnbstmt{\forall a, b.; \forall i \in \mathbb{N}.\; (((a \leq i) \wedge (i \leq b)) \Rightarrow (i \in \operatorname{interval}(a, b)))}}

\vnbbookfile{theorem-library/ord-segment-nn-succ-proof}{theorem-library/ord-segment-nn-succ-proof}
\noindent{\small ord-segment-nn-succ-proof.scm -- the NN-flavoured ORD-SEGMENT successor characterisation, proven from three ordinal primitives.}\par
\vnbprop{ord-segment-nn-succ}{\vnbstmt{\begin{array}{@{}l@{}} \forall n \in \mathbb{N}, k. \\ \quad (k \in \mathrm{OS}((n + 1))) \Leftrightarrow ((k \in \mathrm{OS}(n)) \vee (k = n)) \end{array}}}

\vnbbookfile{theorem-library/ord-segment-nn-subset-proof}{theorem-library/ord-segment-nn-subset-proof}
\noindent{\small ord-segment-nn-subset-proof.scm -- a member of ORD-SEGMENT(m), m in NN, is a natural number.  PROVEN, retiring the `informal' support of the same name (theorem-library/ord-segment-nn-subset.scm:9).}\par
\vnbprop{ord-segment-zero-no-members}{\vnbstmt{\forall k.\; \neg (k \in \mathrm{OS}(0))}}
\vnbprop{ord-segment-nn-subset}{\vnbstmt{\forall m \in \mathbb{N}, k \in \mathrm{OS}(m).\; (k \in \mathbb{N})}}

\vnbbookfile{theorem-library/inf-subsets-is-set}{theorem-library/inf-subsets-is-set}
\noindent{\small inf-subsets-is-set.scm -- INF-SUBSETS(A) is a set when A is, proven.}\par
\vnbprop{inf-subsets-is-set}{\vnbstmt{\forall a \in \mathbf{Set}.\; (\mathrm{Inf}(a) \in \mathbf{Set})}}

\vnbbookfile{theorem-library/sqn-laws}{theorem-library/sqn-laws}
\noindent{\small sqn-laws.scm -- the two SQN facts that structure-library/sqn.scm used to STAMP `definitional', proven here from the definition (the user, 2026-11-13: "Why is this a definition and not an axiom?" -- neither: each is a consequence).}\par
\vnbprop{sqn-membership}{\vnbstmt{\forall a, f.\; ((f \in a^{\mathbb{N}}) \Leftrightarrow (f \in (\mathbb{N} \to a)))}}
\vnbprop{sqn-sethood}{\vnbstmt{\forall a \in \mathbf{Set}.\; (a^{\mathbb{N}} \in \mathbf{Set})}}

\vnbbookfile{theorem-library/binary-minus-laws}{theorem-library/binary-minus-laws}
\noindent{\small binary-minus-laws.scm -- what follows from `binary-minus-def'.}\par
\vnbprop{rr-sub-in-rr}{\vnbstmt{\forall u, v.\; (((u \in \mathbb{R}) \wedge (v \in \mathbb{R})) \Rightarrow ((u - v) \in \mathbb{R}))}}
\vnbprop{zz-sub-in-zz}{\vnbstmt{\forall u, v.\; (((u \in \mathbb{Z}) \wedge (v \in \mathbb{Z})) \Rightarrow ((u - v) \in \mathbb{Z}))}}
\vnbprop{rr-add-in-rr}{\vnbstmt{\forall a, b.\; (((a \in \mathbb{R}) \wedge (b \in \mathbb{R})) \Rightarrow ((a + b) \in \mathbb{R}))}}
\vnbprop{rr-mul-in-rr}{\vnbstmt{\forall a, b.\; (((a \in \mathbb{R}) \wedge (b \in \mathbb{R})) \Rightarrow ((a \cdot b) \in \mathbb{R}))}}

\vnbbookfile{theorem-library/fun-apply-type-proof}{theorem-library/fun-apply-type-proof}
\noindent{\small fun-apply-type-proof.scm -- f : A -$>$ B and x in A give f(x) in B, PROVEN.}\par
\vnbprop{fun-apply-type-c}{\vnbstmt{\forall f, a, b, x.\; (((f \in (a \to b)) \wedge (x \in a)) \Rightarrow (f(x) \in b))}}

\vnbbookfile{theorem-library/image-subset-codomain}{theorem-library/image-subset-codomain}
\noindent{\small image-subset-codomain.scm -- the image of phi : dm -$>$ cod lies in cod.}\par
\vnbprop{image-subset-codomain}{\vnbstmt{\forall m, d.; \forall i \in (m \to d), w \in \operatorname{image}(i, m).\; (w \in d)}}

\vnbbookfile{theorem-library/rake-dc-on-nn}{theorem-library/rake-dc-on-nn}
\noindent{\small rake-dc-on-nn.scm -- dependent choice on NN, PROVEN.}\par
\par\medskip\noindent\textbf{Definition} (\texttt{dc-iter}).\;For all $s$, $a$ and $\mathit{nxt}$,  $(\operatorname{dc\text{-}iter}(s, a, \mathit{nxt}, 0) \simeq a)$\par
Let $s$, $a$ and $\mathit{nxt}$ be arbitrary. Let $n \in \mathbb{N}$. Then
\begin{equation*}
(\operatorname{dc\text{-}iter}(s, a, \mathit{nxt}, \operatorname{succ}(n)) \simeq \varepsilon(\{\operatorname{y\_} \in s : (\operatorname{y\_} \in \operatorname{nxt}(n, \operatorname{dc\text{-}iter}(s, a, \mathit{nxt}, n)))\}))
\end{equation*}

\vnbprop{dc-iter-in}{\vnbstmt{\begin{array}{@{}l@{}} \forall n \in \mathbb{N}, s.; \forall a \in s, x. \\ \quad \forall k \in \mathbb{N}, u \in s.\; \exists y \in s.\; (y \in x(k, u)) \Rightarrow  \\ \quad \quad \operatorname{dc\text{-}iter}(s, a, x, n) \in s \end{array}}}
\vnbprop{dc-iter-step}{\vnbstmt{\begin{array}{@{}l@{}} \forall n \in \mathbb{N}, s.; \forall a \in s, x. \\ \quad \forall k \in \mathbb{N}, u \in s.\; \exists y \in s.\; (y \in x(k, u)) \Rightarrow  \\ \quad \quad \operatorname{dc\text{-}iter}(s, a, x, (n + 1)) \in x(n, \operatorname{dc\text{-}iter}(s, a, x, n)) \end{array}}}
\vnbprop{dc-on-nn-pred}{\vnbstmt{\begin{array}{@{}l@{}} \forall x \in \mathbf{Set}, a \in x, t. \\ \quad \forall k \in \mathbb{N}, u \in x.\; \exists y \in x.\; (y \in t(k, u)) \Rightarrow  \\ \quad \quad \exists f \in (\mathbb{N} \to x). \\ \quad \quad \quad (f(0) = a) \wedge \forall k \in \mathbb{N}.\; (f((k + 1)) \in t(k, f(k))) \end{array}}}
\vnbprop{dc-on-nn}{\vnbstmt{\begin{array}{@{}l@{}} \forall x \in \mathbf{Set}, a \in x, r \in \mathbf{Set}. \\ \quad \forall k \in \mathbb{N}, u \in x.\; \exists y \in x.\; ({[}k, u, y] \in r) \Rightarrow  \\ \quad \quad \exists f \in (\mathbb{N} \to x). \\ \quad \quad \quad (f(0) = a) \wedge \forall k \in \mathbb{N}.\; ({[}k, f(k), f((k + 1))] \in r) \end{array}}}

\vnbbookfile{theorem-library/compose-apply-proof}{theorem-library/compose-apply-proof}
\noindent{\small compose-apply-proof.scm -- (COMPOSE f g)(x) = f(g(x)), PROVEN modulo 0.}\par
\vnbprop{compose-apply}{\vnbstmt{\begin{array}{@{}l@{}} \forall a, b, c.; \forall f \in (b \to c), g \in (a \to b), x \in a. \\ \quad \operatorname{compose}(f, g)(x) = f(g(x)) \end{array}}}
\vnbprop{compose-type}{\vnbstmt{\begin{array}{@{}l@{}} \forall a, b, c, f, g. \\ \quad ((a \in \mathbf{Set}) \wedge ((g \in (a \to b)) \wedge (f \in (b \to c)))) \Rightarrow  \\ \quad \quad \operatorname{compose}(f, g) \in (a \to c) \end{array}}}

\vnbbookfile{theorem-library/id-fun-laws}{theorem-library/id-fun-laws}
\noindent{\small id-fun-laws.scm -- the identity map and the composition laws over FUN.}\par
\vnbprop{id-fun-type}{\vnbstmt{\forall a \in \mathbf{Set}.\; (\operatorname{id\text{-}fun}(a) \in (a \to a))}}
\vnbprop{id-fun-apply}{\vnbstmt{\forall a.; \forall x \in a.\; (\operatorname{id\text{-}fun}(a)(x) = x)}}
\vnbprop{compose-id-left}{\vnbstmt{\begin{array}{@{}l@{}} \forall a, b, f. \\ \quad ((a \in \mathbf{Set}) \wedge ((b \in \mathbf{Set}) \wedge (f \in (a \to b)))) \Rightarrow (\operatorname{compose}(\operatorname{id\text{-}fun}(b), f) = f) \end{array}}}
\vnbprop{compose-id-right}{\vnbstmt{\forall a, b, f.\; (((a \in \mathbf{Set}) \wedge (f \in (a \to b))) \Rightarrow (\operatorname{compose}(f, \operatorname{id\text{-}fun}(a)) = f))}}
\vnbprop{compose-assoc}{\vnbstmt{\begin{array}{@{}l@{}} \forall a, b, c, u, f, g, h. \\ \quad ((a \in \mathbf{Set}) \wedge ((b \in \mathbf{Set}) \wedge ((f \in (a \to b)) \wedge ((g \in (b \to c)) \wedge (h \in (c \to u)))))) \Rightarrow  \\ \quad \quad \operatorname{compose}(h, \operatorname{compose}(g, f)) = \operatorname{compose}(\operatorname{compose}(h, g), f) \end{array}}}

\vnbbookfile{theorem-library/pair-tuple-sethood}{theorem-library/pair-tuple-sethood}
\noindent{\small pair-tuple-sethood.scm -- a pair [u,v] is a set, PROVEN modulo 0.}\par
\vnbprop{pair-in-cartesian}{\vnbstmt{\forall x, y, u, v.\; (((u \in x) \wedge (v \in y)) \Rightarrow ({[}u, v] \in (x \times y)))}}
\vnbprop{pair-tuple-is-set}{\vnbstmt{\forall x, y, u, v.\; (((u \in x) \wedge (v \in y)) \Rightarrow ({[}u, v] \in \mathbf{Set}))}}

\vnbbookfile{theorem-library/pair-laws}{theorem-library/pair-laws}
\noindent{\small pair-laws.scm -- plain facts about pairs and cartesian products, in CORE.}\par
\vnbprop{cartesian-nth}{\vnbstmt{\begin{array}{@{}l@{}} \forall x, a, b. \\ \quad (x \in (a \times b)) \Rightarrow ((\operatorname{nth}(1, x) \in a) \wedge (\operatorname{nth}(2, x) \in b)) \end{array}}}
\vnbprop{cartesian-pair-eq}{\vnbstmt{\forall x, a, b.\; ((x \in (a \times b)) \Rightarrow (x = {[}\operatorname{nth}(1, x), \operatorname{nth}(2, x)]))}}
\vnbprop{pair-in-cartesian-parts}{\vnbstmt{\forall x, y, u, v.\; (({[}u, v] \in (x \times y)) \Rightarrow ((u \in x) \wedge (v \in y)))}}
\vnbprop{node-pair-injective}{\vnbstmt{\forall k, a, j, b.\; (({[}k, a] = {[}j, b]) \Rightarrow ((k = j) \wedge (a = b)))}}

\vnbbookfile{theorem-library/nn-order-ord}{theorem-library/nn-order-ord}
\noindent{\small nn-order-ord.scm -- the elementary NN order block, PROVEN.  Seven facts that were `well-known' supports in structure-library/order-lemmas.scm are theorems here, every one billing `modulo 0'.}\par
\vnbprop{nn-one-in}{\vnbstmt{1 \in \mathbb{N}}}
\vnbprop{nn-zero-le}{\vnbstmt{\forall n \in \mathbb{N}.\; (0 \leq n)}}
\vnbprop{nn-le-succ}{\vnbstmt{\forall k \in \mathbb{N}.\; (k \leq (k + 1))}}
\vnbprop{nn-le-imp-neq-succ}{\vnbstmt{\forall k, j \in \mathbb{N}.\; ((j \leq k) \Rightarrow \neg (j = (k + 1)))}}
\vnbprop{nn-succ-mono}{\vnbstmt{\forall a, b \in \mathbb{N}.\; ((a \leq b) \Rightarrow ((a + 1) \leq (b + 1)))}}
\vnbprop{nn-one-le-succ}{\vnbstmt{\forall n \in \mathbb{N}.\; (1 \leq (n + 1))}}
\vnbprop{nn-le-succ-cases}{\vnbstmt{\forall k, j \in \mathbb{N}.\; ((j \leq (k + 1)) \Rightarrow ((j \leq k) \vee (j = (k + 1))))}}

\vnbbookfile{theorem-library/nn-not-le-succ-le}{theorem-library/nn-not-le-succ-le}
\noindent{\small nn-not-le-succ-le.scm -- NN is totally ordered and discrete, PROVEN.}\par
\vnbprop{nn-not-le-succ-le}{\vnbstmt{\forall m, n \in \mathbb{N}.\; (\neg (m \leq n) \Rightarrow ((n + 1) \leq m))}}

\vnbbookfile{theorem-library/nn-order-basics}{theorem-library/nn-order-basics}
\noindent{\small nn-order-basics.scm -- the elementary NN order facts, PROVEN.}\par
\vnbprop{nn-in-rr}{\vnbstmt{\forall k \in \mathbb{N}.\; (k \in \mathbb{R})}}
\vnbprop{rr-le-trans}{\vnbstmt{\forall x, y, z \in \mathbb{R}.\; (((x \leq y) \wedge (y \leq z)) \Rightarrow (x \leq z))}}
\vnbprop{rr-le-trans-c}{\vnbstmt{\forall x, y, z \in \mathbb{R}.\; (((x \leq y) \wedge (y \leq z)) \Rightarrow (x \leq z))}}
\vnbprop{nn-le-trans-guarded}{\vnbstmt{\forall a, b, c \in \mathbb{N}.\; (((a \leq b) \wedge (b \leq c)) \Rightarrow (a \leq c))}}
\vnbprop{nn-le-refl}{\vnbstmt{\forall k \in \mathbb{N}.\; (k \leq k)}}
\vnbprop{nn-le-add-right}{\vnbstmt{\forall n, m \in \mathbb{N}.\; (m \leq (m + n))}}
\vnbprop{nn-pair-upper-bound}{\vnbstmt{\begin{array}{@{}l@{}} \forall a, b. \\ \quad ((a \in \mathbb{N}) \wedge (b \in \mathbb{N})) \Rightarrow \exists c \in \mathbb{N}.\; ((a \leq c) \wedge (b \leq c)) \end{array}}}
\vnbprop{zz-in-rr}{\vnbstmt{\forall k \in \mathbb{Z}.\; (k \in \mathbb{R})}}

\vnbbookfile{theorem-library/mat-basics}{theorem-library/mat-basics}
\noindent{\small mat-basics.scm -- MAT's read-offs, PROVEN from its definition.}\par
\vnbprop{mat-unfold}{\vnbstmt{\begin{array}{@{}l@{}} \forall m, n, x. \\ \quad (\{p \in \operatorname{matrix}(x) : ((n \in \mathbb{N}) \wedge (((\operatorname{length}(p) = 0) \Rightarrow (m = 0)) \wedge (\neg (\operatorname{length}(p) = 0) \Rightarrow (\operatorname{size}(p) = {[}m, n]))))\} \simeq \operatorname{mat}(m, n, x)) \end{array}}}
\vnbprop{size-unfold}{\vnbstmt{\forall m.\; ({[}\operatorname{length}(m), \left(\text{if }(\operatorname{length}(m) = 0)\text{ then }0\text{ else }\operatorname{length}(\operatorname{nth}(1, m))\right)] \simeq \operatorname{size}(m))}}
\vnbprop{size-nonempty}{\vnbstmt{\begin{array}{@{}l@{}} \forall m. \\ \quad \neg (\operatorname{length}(m) = 0) \Rightarrow ({[}\operatorname{length}(m), \operatorname{length}(\operatorname{nth}(1, m))] \simeq \operatorname{size}(m)) \end{array}}}
\vnbprop{nth1-pair}{\vnbstmt{\forall a, b.\; (((a \in \mathbf{Set}) \wedge (b \in \mathbf{Set})) \Rightarrow (\operatorname{nth}(1, {[}a, b]) \simeq a))}}
\vnbprop{mat-in-matrix}{\vnbstmt{\forall m, n, x.; \forall q \in \operatorname{mat}(m, n, x).\; (q \in \operatorname{matrix}(x))}}
\vnbprop{mat-cols-in-nn}{\vnbstmt{\forall m, n, x.; \forall q \in \operatorname{mat}(m, n, x).\; (n \in \mathbb{N})}}
\vnbprop{mat-length}{\vnbstmt{\forall m, n, x.; \forall q \in \operatorname{mat}(m, n, x).\; (\operatorname{length}(q) = m)}}
\vnbprop{mat-rows-in-nn}{\vnbstmt{\forall m, n, x.; \forall q \in \operatorname{mat}(m, n, x).\; (m \in \mathbb{N})}}
\vnbprop{mat-size}{\vnbstmt{\forall m, n, x.; \forall q \in \operatorname{mat}(m, n, x).\; ((1 \leq m) \Rightarrow (\operatorname{size}(q) = {[}m, n]))}}
\vnbprop{mat-cols}{\vnbstmt{\forall m, n, x.; \forall q \in \operatorname{mat}(m, n, x).\; ((1 \leq m) \Rightarrow (\operatorname{length}(\operatorname{nth}(1, q)) = n))}}
\vnbprop{mat-size-rows}{\vnbstmt{\forall m, n, x.; \forall q \in \operatorname{mat}(m, n, x).\; (\operatorname{nth}(1, \operatorname{size}(q)) = m)}}
\vnbprop{mat-size-cols}{\vnbstmt{\forall m, n, x.; \forall q \in \operatorname{mat}(m, n, x).\; ((1 \leq m) \Rightarrow (\operatorname{nth}(2, \operatorname{size}(q)) = n))}}
\vnbprop{mat-intro}{\vnbstmt{\begin{array}{@{}l@{}} \forall m, n, x.; \forall q \in \operatorname{matrix}(x). \\ \quad ((n \in \mathbb{N}) \wedge ((1 \leq m) \wedge ((\operatorname{length}(q) = m) \wedge (\operatorname{length}(\operatorname{nth}(1, q)) = n)))) \Rightarrow  \\ \quad \quad q \in \operatorname{mat}(m, n, x) \end{array}}}
\vnbprop{mat-intro-empty}{\vnbstmt{\begin{array}{@{}l@{}} \forall n, x.; \forall q \in \operatorname{matrix}(x). \\ \quad ((n \in \mathbb{N}) \wedge (\operatorname{length}(q) = 0)) \Rightarrow (q \in \operatorname{mat}(0, n, x)) \end{array}}}
\vnbprop{nil-in-mat}{\vnbstmt{\forall n \in \mathbb{N}, x.\; ({[}] \in \operatorname{mat}(0, n, x))}}
\vnbprop{mat-size-cols-in-nn}{\vnbstmt{\forall m, n, x.; \forall q \in \operatorname{mat}(m, n, x).\; (\operatorname{nth}(2, \operatorname{size}(q)) \in \mathbb{N})}}
\vnbprop{mat-size-cols-pos}{\vnbstmt{\begin{array}{@{}l@{}} \forall m, n, x.; \forall q \in \operatorname{mat}(m, n, x). \\ \quad (1 \leq \operatorname{nth}(2, \operatorname{size}(q))) \Rightarrow (\operatorname{nth}(2, \operatorname{size}(q)) = n) \end{array}}}
\vnbprop{mat-empty-any-cols}{\vnbstmt{\forall c, x, q, k.\; (((q \in \operatorname{mat}(0, c, x)) \wedge (k \in \mathbb{N})) \Rightarrow (q \in \operatorname{mat}(0, k, x)))}}
\vnbprop{interval-1-0-empty}{\vnbstmt{\operatorname{interval}(1, 0) = \emptyset}}

\vnbbookfile{theorem-library/nn-pos-of-nonzero}{theorem-library/nn-pos-of-nonzero}
\noindent{\small nn-pos-of-nonzero.scm -- a nonzero natural is strictly positive.}\par
\vnbprop{nn-pos-of-nonzero}{\vnbstmt{\forall k \in \mathbb{N}.\; (\neg (k = 0) \Rightarrow (0 < k))}}

\vnbbookfile{theorem-library/interval-membership}{theorem-library/interval-membership}
\noindent{\small interval-membership.scm -- INTERVAL's membership IFF and the [1,1] instance, both PROVEN, both by citation of the read-offs already in the tree.}\par
\vnbprop{interval-membership}{\vnbstmt{\forall a, b, i.\; ((i \in \operatorname{interval}(a, b)) \Leftrightarrow ((i \in \mathbb{N}) \wedge ((a \leq i) \wedge (i \leq b))))}}
\vnbprop{one-in-interval-1}{\vnbstmt{1 \in \operatorname{interval}(1, 1)}}

\vnbbookfile{theorem-library/nn-not-lt-le-proof}{theorem-library/nn-not-lt-le-proof}
\noindent{\small nn-not-lt-le-proof.scm -- NOT (k $<$ j) =$>$ j $<$= k on NN, `modulo 0'.}\par
\vnbprop{nn-not-lt-le}{\vnbstmt{\forall k, j \in \mathbb{N}.\; (\neg (k < j) \Rightarrow (j \leq k))}}

\vnbbookfile{theorem-library/rr-order-basics}{theorem-library/rr-order-basics}
\noindent{\small rr-order-basics.scm -- the elementary RR order facts, PROVEN.}\par
\vnbprop{rr-lt-trans}{\vnbstmt{\forall x, y, z \in \mathbb{R}.\; (((x < y) \wedge (y < z)) \Rightarrow (x < z))}}
\vnbprop{rr-lt-le-trans}{\vnbstmt{\forall x, y, z \in \mathbb{R}.\; (((x < y) \wedge (y \leq z)) \Rightarrow (x < z))}}
\vnbprop{rr-le-lt-trans}{\vnbstmt{\forall x, y, z \in \mathbb{R}.\; (((x \leq y) \wedge (y < z)) \Rightarrow (x < z))}}
\vnbprop{rr-lt-implies-le}{\vnbstmt{\forall x, y \in \mathbb{R}.\; ((x < y) \Rightarrow (x \leq y))}}
\vnbprop{rr-le-diff-nonpos}{\vnbstmt{\forall u, v \in \mathbb{R}.\; ((u \leq v) \Rightarrow ((u - v) \leq 0))}}
\vnbprop{rr-lt-diff-pos}{\vnbstmt{\forall u, v \in \mathbb{R}.\; ((u < v) \Rightarrow (0 < (v - u)))}}
\vnbprop{rr-lt-diff-neg}{\vnbstmt{\forall u, v \in \mathbb{R}.\; ((u < v) \Rightarrow ((u - v) < 0))}}
\vnbprop{rr-diff-zero-eq}{\vnbstmt{\forall u, v \in \mathbb{R}.\; ((0 = (u - v)) \Rightarrow (u = v))}}
\vnbprop{rr-le-neg}{\vnbstmt{\forall u, v \in \mathbb{R}.\; ((u \leq v) \Rightarrow (-v \leq -u))}}
\vnbprop{rr-neg-eq-zero}{\vnbstmt{\forall u \in \mathbb{R}.\; ((-u = 0) \Rightarrow (u = 0))}}
\vnbprop{rr-sub-ne-zero}{\vnbstmt{\forall u, v \in \mathbb{R}.\; (\neg (u = v) \Rightarrow \neg ((u - v) = 0))}}
\vnbprop{rr-pos-ne-zero}{\vnbstmt{\forall c \in \mathbb{R}.\; ((0 < c) \Rightarrow \neg (c = 0))}}
\vnbprop{rr-lt-irrefl}{\vnbstmt{\forall a \in \mathbb{R}.\; \neg (a < a)}}
\vnbprop{rr-le-cases}{\vnbstmt{\forall u, v \in \mathbb{R}.\; ((u \leq v) \Rightarrow ((u < v) \vee (u = v)))}}
\vnbprop{rr-le-ne-lt}{\vnbstmt{\forall u, v \in \mathbb{R}.\; (((u \leq v) \wedge \neg (u = v)) \Rightarrow (u < v))}}
\vnbprop{rr-le-add}{\vnbstmt{\begin{array}{@{}l@{}} \forall x, y, u, v. \\ \quad ((x \in \mathbb{R}) \wedge ((y \in \mathbb{R}) \wedge ((u \in \mathbb{R}) \wedge ((v \in \mathbb{R}) \wedge ((x \leq y) \wedge (u \leq v)))))) \Rightarrow  \\ \quad \quad (x + u) \leq (y + v) \end{array}}}
\vnbprop{rr-lt-add}{\vnbstmt{\begin{array}{@{}l@{}} \forall x, y, u, v. \\ \quad ((x \in \mathbb{R}) \wedge ((y \in \mathbb{R}) \wedge ((u \in \mathbb{R}) \wedge ((v \in \mathbb{R}) \wedge ((x < y) \wedge (u \leq v)))))) \Rightarrow  \\ \quad \quad (x + u) < (y + v) \end{array}}}
\vnbprop{rr-add-nonneg}{\vnbstmt{\forall x, y.\; (((x \in \mathbb{R}) \wedge ((y \in \mathbb{R}) \wedge ((0 \leq x) \wedge (0 \leq y)))) \Rightarrow (0 \leq (x + y)))}}
\vnbprop{rr-le-from-diff-nonneg}{\vnbstmt{\forall x, y.\; (((x \in \mathbb{R}) \wedge ((y \in \mathbb{R}) \wedge (0 \leq (y - x)))) \Rightarrow (x \leq y))}}
\vnbprop{rr-double-nonneg}{\vnbstmt{\forall x \in \mathbb{R}.\; ((0 \leq (x + x)) \Rightarrow (0 \leq x))}}
\vnbprop{rr-no-zero-divisors}{\vnbstmt{\forall a, b.\; (((a \in \mathbb{R}) \wedge ((b \in \mathbb{R}) \wedge ((a \cdot b) = 0))) \Rightarrow ((a = 0) \vee (b = 0)))}}
\vnbprop{rr-cancel-mul-right}{\vnbstmt{\begin{array}{@{}l@{}} \forall u, v, c. \\ \quad ((u \in \mathbb{R}) \wedge ((v \in \mathbb{R}) \wedge ((c \in \mathbb{R}) \wedge (\neg (c = 0) \wedge ((u \cdot c) = (v \cdot c)))))) \Rightarrow (u = v) \end{array}}}
\vnbprop{rr-lt-scale-pos}{\vnbstmt{\begin{array}{@{}l@{}} \forall c, x, y. \\ \quad ((c \in \mathbb{R}) \wedge ((x \in \mathbb{R}) \wedge ((y \in \mathbb{R}) \wedge ((0 < c) \wedge (x < y))))) \Rightarrow ((c \cdot x) < (c \cdot y)) \end{array}}}
\vnbprop{rr-le-scale-nonneg}{\vnbstmt{\begin{array}{@{}l@{}} \forall c, x, y. \\ \quad ((c \in \mathbb{R}) \wedge ((x \in \mathbb{R}) \wedge ((y \in \mathbb{R}) \wedge ((0 \leq c) \wedge (x \leq y))))) \Rightarrow ((c \cdot x) \leq (c \cdot y)) \end{array}}}
\vnbprop{rr-sq-nonneg}{\vnbstmt{\forall x \in \mathbb{R}.\; (0 \leq (x \cdot x))}}
\vnbprop{rr-lt-trichotomy}{\vnbstmt{\forall u, v \in \mathbb{R}.\; ((u < v) \vee ((u = v) \vee (v < u)))}}
\vnbprop{rr-upper-of-two}{\vnbstmt{\begin{array}{@{}l@{}} \forall u, v. \\ \quad ((u \in \mathbb{R}) \wedge (v \in \mathbb{R})) \Rightarrow \exists w \in \mathbb{R}.\; ((u \leq w) \wedge (v \leq w)) \end{array}}}
\vnbprop{rr-upper-of-two-below}{\vnbstmt{\begin{array}{@{}l@{}} \forall u, v, m. \\ \quad ((u \in \mathbb{R}) \wedge ((v \in \mathbb{R}) \wedge ((m \in \mathbb{R}) \wedge ((u < m) \wedge (v < m))))) \Rightarrow  \\ \quad \quad \exists w \in \mathbb{R}.\; ((u \leq w) \wedge ((v \leq w) \wedge (w < m))) \end{array}}}
\vnbprop{rr-min-pos}{\vnbstmt{\begin{array}{@{}l@{}} \forall u, v. \\ \quad ((u \in \mathbb{R}) \wedge ((v \in \mathbb{R}) \wedge ((0 < u) \wedge (0 < v)))) \Rightarrow  \\ \quad \quad \exists w \in \mathbb{R}.\; ((0 < w) \wedge ((w \leq u) \wedge (w \leq v))) \end{array}}}
\vnbprop{rr-prod-le-prod}{\vnbstmt{\begin{array}{@{}l@{}} \forall a, b, u, v. \\ \quad ((a \in \mathbb{R}) \wedge ((b \in \mathbb{R}) \wedge ((u \in \mathbb{R}) \wedge ((v \in \mathbb{R}) \wedge (((0 \leq a) \wedge (a \leq b)) \wedge ((0 \leq u) \wedge (u \leq v))))))) \Rightarrow  \\ \quad \quad (a \cdot u) \leq (b \cdot v) \end{array}}}
\vnbprop{rr-nonneg-cancel-pos}{\vnbstmt{\forall c, d.\; (((c \in \mathbb{R}) \wedge ((d \in \mathbb{R}) \wedge ((0 < c) \wedge (0 \leq (c \cdot d))))) \Rightarrow (0 \leq d))}}
\vnbprop{rr-le-total}{\vnbstmt{\forall u, v \in \mathbb{R}.\; ((u \leq v) \vee (v \leq u))}}
\vnbprop{rr-lt-not-le}{\vnbstmt{\forall a, b \in \mathbb{R}.\; ((a < b) \Rightarrow \neg (b \leq a))}}
\vnbprop{rr-not-le-lt}{\vnbstmt{\forall a, b \in \mathbb{R}.\; (\neg (a \leq b) \Rightarrow (b < a))}}
\vnbprop{rr-mul-le-right}{\vnbstmt{\forall a, b, c \in \mathbb{R}.\; (((a \leq b) \wedge (0 \leq c)) \Rightarrow ((a \cdot c) \leq (b \cdot c)))}}

\vnbbookfile{theorem-library/co-order-trans-guarded}{theorem-library/co-order-trans-guarded}
\noindent{\small co-order-trans-guarded.scm -- the six `co-' order composition lemmas, PROVEN.  Four of them GUARDED on RR; two of them exactly as stated.}\par
\vnbprop{co-lt-trans}{\vnbstmt{\forall a, b, c \in \mathbb{R}.\; (((a < b) \wedge (b < c)) \Rightarrow (a < c))}}
\vnbprop{co-le-lt-trans}{\vnbstmt{\forall a, b, c \in \mathbb{R}.\; (((a \leq b) \wedge (b < c)) \Rightarrow (a < c))}}
\vnbprop{co-lt-le-trans}{\vnbstmt{\forall a, b, c \in \mathbb{R}.\; (((a < b) \wedge (b \leq c)) \Rightarrow (a < c))}}
\vnbprop{co-le-eq-trans}{\vnbstmt{\forall a, b, c.\; (((a \leq b) \wedge (b = c)) \Rightarrow (a \leq c))}}
\vnbprop{co-eq-le-trans}{\vnbstmt{\forall a, b, c.\; (((a = b) \wedge (b \leq c)) \Rightarrow (a \leq c))}}

\vnbbookfile{theorem-library/pos-rr-bridges}{theorem-library/pos-rr-bridges}
\noindent{\small pos-rr-bridges.scm -- the two projections of POS-RR, and one more positivity fact, PROVEN modulo 0.}\par
\vnbprop{rr-pos-rr-in-rr}{\vnbstmt{\forall t.\; (\operatorname{pos\text{-}rr}(t) \Rightarrow (t \in \mathbb{R}))}}
\vnbprop{rr-lt-of-pos-rr}{\vnbstmt{\forall t.\; (\operatorname{pos\text{-}rr}(t) \Rightarrow (0 < t))}}
\vnbprop{rr-one-plus-nonneg-pos}{\vnbstmt{\forall d \in \mathbb{R}.\; ((0 \leq d) \Rightarrow (0 < (1 + d)))}}

\vnbbookfile{theorem-library/rr-abs-basics}{theorem-library/rr-abs-basics}
\noindent{\small rr-abs-basics.scm -- absolute value on RR, PROVEN from its DEFINITION.}\par
\vnbprop{rr-abs-of-nonneg}{\vnbstmt{\forall x \in \mathbb{R}.\; ((0 \leq x) \Rightarrow (\lvert x \rvert = x))}}
\vnbprop{rr-abs-of-nonpos}{\vnbstmt{\forall x \in \mathbb{R}.\; ((x \leq 0) \Rightarrow (\lvert x \rvert = -x))}}
\vnbprop{rr-abs-of-neg}{\vnbstmt{\forall x \in \mathbb{R}.\; ((x < 0) \Rightarrow (\lvert x \rvert = -x))}}
\vnbprop{rr-abs-cases}{\vnbstmt{\forall x \in \mathbb{R}.\; (((0 \leq x) \wedge (\lvert x \rvert = x)) \vee (\neg (0 \leq x) \wedge (\lvert x \rvert = -x)))}}
\vnbprop{rr-abs-closed}{\vnbstmt{\forall a \in \mathbb{R}.\; (\lvert a \rvert \in \mathbb{R})}}
\vnbprop{rr-abs-nonneg}{\vnbstmt{\forall a \in \mathbb{R}.\; (0 \leq \lvert a \rvert)}}
\vnbprop{rr-abs-zero}{\vnbstmt{\forall a \in \mathbb{R}.\; ((\lvert a \rvert = 0) \Leftrightarrow (a = 0))}}
\vnbprop{rr-le-abs}{\vnbstmt{\forall x \in \mathbb{R}.\; (x \leq \lvert x \rvert)}}
\vnbprop{rr-neg-abs-le}{\vnbstmt{\forall x \in \mathbb{R}.\; (-\lvert x \rvert \leq x)}}
\vnbprop{rr-abs-bound}{\vnbstmt{\forall x, c \in \mathbb{R}.\; ((\lvert x \rvert \leq c) \Leftrightarrow ((-c \leq x) \wedge (x \leq c)))}}
\vnbprop{rr-abs-triangle}{\vnbstmt{\forall a, b.\; (((a \in \mathbb{R}) \wedge (b \in \mathbb{R})) \Rightarrow (\lvert (a + b) \rvert \leq (\lvert a \rvert + \lvert b \rvert)))}}
\vnbprop{rr-abs-triangle-c}{\vnbstmt{\forall a, b.\; (((a \in \mathbb{R}) \wedge (b \in \mathbb{R})) \Rightarrow (\lvert (a + b) \rvert \leq (\lvert a \rvert + \lvert b \rvert)))}}
\vnbprop{rr-abs-neg}{\vnbstmt{\forall a \in \mathbb{R}.\; (\lvert -a \rvert = \lvert a \rvert)}}
\vnbprop{rr-abs-sub-sym}{\vnbstmt{\forall u, v.\; (((u \in \mathbb{R}) \wedge (v \in \mathbb{R})) \Rightarrow (\lvert (u - v) \rvert = \lvert (v - u) \rvert))}}
\vnbprop{rr-abs-mult}{\vnbstmt{\forall a, b.\; (((a \in \mathbb{R}) \wedge (b \in \mathbb{R})) \Rightarrow (\lvert (a \cdot b) \rvert = (\lvert a \rvert \cdot \lvert b \rvert)))}}
\vnbprop{rr-abs-reverse-triangle}{\vnbstmt{\forall x, y \in \mathbb{R}.\; (\lvert (\lvert x \rvert - \lvert y \rvert) \rvert \leq \lvert (x - y) \rvert)}}
\vnbprop{rr-abs-sum-bound}{\vnbstmt{\begin{array}{@{}l@{}} \forall u, v, p, q, h, e. \\ \quad u \in \mathbb{R} \wedge \\ \quad v \in \mathbb{R} \wedge \\ \quad p \in \mathbb{R} \wedge \\ \quad q \in \mathbb{R} \wedge \\ \quad h \in \mathbb{R} \wedge \\ \quad e \in \mathbb{R} \wedge \\ \quad \lvert (u - p) \rvert \leq h \wedge \\ \quad \lvert (v - q) \rvert \leq h \wedge \\ \quad (h + h) = e \Rightarrow \\ \quad \quad \lvert ((u + v) - (p + q)) \rvert \leq e \end{array}}}
\vnbprop{rr-abs-prod-bound}{\vnbstmt{\begin{array}{@{}l@{}} \forall u, v, p, q, m, k, t, e. \\ \quad u \in \mathbb{R} \wedge \\ \quad v \in \mathbb{R} \wedge \\ \quad p \in \mathbb{R} \wedge \\ \quad q \in \mathbb{R} \wedge \\ \quad m \in \mathbb{R} \wedge \\ \quad k \in \mathbb{R} \wedge \\ \quad t \in \mathbb{R} \wedge \\ \quad e \in \mathbb{R} \wedge \\ \quad \lvert u \rvert \leq m \wedge \\ \quad \lvert q \rvert \leq k \wedge \\ \quad \lvert (u - p) \rvert \leq t \wedge \\ \quad \lvert (v - q) \rvert \leq t \wedge \\ \quad ((m + k) \cdot t) \leq e \Rightarrow \\ \quad \quad \lvert ((u \cdot v) - (p \cdot q)) \rvert \leq e \end{array}}}
\vnbprop{rr-abs-zero-value}{\vnbstmt{\lvert 0 \rvert = 0}}

\vnbbookfile{theorem-library/rr-max-basics}{theorem-library/rr-max-basics}
\noindent{\small rr-max-basics.scm -- the usable facts about MAX on RR, PROVEN from its defining equation.}\par
\vnbprop{rr-max-closed}{\vnbstmt{\forall x, y \in \mathbb{R}.\; (\operatorname{max}(x, y) \in \mathbb{R})}}
\vnbprop{rr-max-cases}{\vnbstmt{\forall x, y \in \mathbb{R}.\; ((\operatorname{max}(x, y) = x) \vee (\operatorname{max}(x, y) = y))}}
\vnbprop{rr-le-max-left}{\vnbstmt{\forall x, y \in \mathbb{R}.\; (x \leq \operatorname{max}(x, y))}}
\vnbprop{rr-le-max-right}{\vnbstmt{\forall x, y \in \mathbb{R}.\; (y \leq \operatorname{max}(x, y))}}
\vnbprop{rr-max-le}{\vnbstmt{\forall x, y, z \in \mathbb{R}.\; (((x \leq z) \wedge (y \leq z)) \Rightarrow (\operatorname{max}(x, y) \leq z))}}

\vnbbookfile{theorem-library/rr-min-basics}{theorem-library/rr-min-basics}
\noindent{\small rr-min-basics.scm -- the usable facts about MIN on RR, PROVEN from its defining equation, and the NN closure of both MIN and MAX.}\par
\vnbprop{rr-min-closed}{\vnbstmt{\forall x, y \in \mathbb{R}.\; (\operatorname{min}(x, y) \in \mathbb{R})}}
\vnbprop{rr-min-cases}{\vnbstmt{\forall x, y \in \mathbb{R}.\; ((\operatorname{min}(x, y) = x) \vee (\operatorname{min}(x, y) = y))}}
\vnbprop{rr-min-le-left}{\vnbstmt{\forall x, y \in \mathbb{R}.\; (\operatorname{min}(x, y) \leq x)}}
\vnbprop{rr-min-le-right}{\vnbstmt{\forall x, y \in \mathbb{R}.\; (\operatorname{min}(x, y) \leq y)}}
\vnbprop{rr-le-min}{\vnbstmt{\forall x, y, z \in \mathbb{R}.\; (((z \leq x) \wedge (z \leq y)) \Rightarrow (z \leq \operatorname{min}(x, y)))}}
\vnbprop{nn-max-closed}{\vnbstmt{\forall x, y \in \mathbb{N}.\; (\operatorname{max}(x, y) \in \mathbb{N})}}
\vnbprop{nn-min-closed}{\vnbstmt{\forall x, y \in \mathbb{N}.\; (\operatorname{min}(x, y) \in \mathbb{N})}}
\vnbprop{rr-min-one-subadditive}{\vnbstmt{\begin{array}{@{}l@{}} \forall x, y \in \mathbb{R}. \\ \quad ((0 \leq x) \wedge (0 \leq y)) \Rightarrow (\operatorname{min}(1, (x + y)) \leq (\operatorname{min}(1, x) + \operatorname{min}(1, y))) \end{array}}}

\vnbbookfile{theorem-library/cc-real-imag}{theorem-library/cc-real-imag}
\noindent{\small cc-real-imag.scm -- the real and imaginary parts of a complex number.}\par
\vnbprop{cc-re-im-of}{\vnbstmt{\begin{array}{@{}l@{}} \forall z \in \mathbb{C}, a, b \in \mathbb{R}. \\ \quad (z = (a + (b \cdot +i))) \Rightarrow ((\operatorname{real\text{-}part}(z) = a) \wedge (\operatorname{imag\text{-}part}(z) = b)) \end{array}}}
\vnbprop{real-part-in-rr}{\vnbstmt{\forall z \in \mathbb{C}.\; (\operatorname{real\text{-}part}(z) \in \mathbb{R})}}
\vnbprop{imag-part-in-rr}{\vnbstmt{\forall z \in \mathbb{C}.\; (\operatorname{imag\text{-}part}(z) \in \mathbb{R})}}
\vnbprop{cc-re-im-decompose}{\vnbstmt{\forall z \in \mathbb{C}.\; (z = (\operatorname{real\text{-}part}(z) + (+i \cdot \operatorname{imag\text{-}part}(z))))}}
\vnbprop{cc-plus-conj-is-2re}{\vnbstmt{\forall z \in \mathbb{C}.\; ((z + \operatorname{conjugate}(z)) = (2 \cdot \operatorname{real\text{-}part}(z)))}}
\vnbprop{cc-re-sq-le-mod-sq}{\vnbstmt{\forall z \in \mathbb{C}.\; ((\operatorname{real\text{-}part}(z) \cdot \operatorname{real\text{-}part}(z)) \leq (z \cdot \operatorname{conjugate}(z)))}}

\vnbbookfile{theorem-library/rr-recip-order}{theorem-library/rr-recip-order}
\noindent{\small rr-recip-order.scm -- the L1 rung of the "attach the loose ends" program: the reciprocal's SIGN, and the two field facts it stands on.}\par
\vnbprop{rr-mul-zero}{\vnbstmt{\forall a \in \mathbb{R}.\; ((a \cdot 0) = 0)}}
\vnbprop{rr-zero-lt-one}{\vnbstmt{0 < 1}}
\vnbprop{rr-mul-pos}{\vnbstmt{\forall a, b \in \mathbb{R}.\; (((0 < a) \wedge (0 < b)) \Rightarrow (0 < (a \cdot b)))}}
\vnbprop{rr-recip-pos}{\vnbstmt{\forall a \in \mathbb{R}.\; ((0 < a) \Rightarrow (0 < \operatorname{recip}(a)))}}

\vnbbookfile{theorem-library/pos-rr-of-lt}{theorem-library/pos-rr-of-lt}
\noindent{\small pos-rr-of-lt.scm -- 0 $<$ x IMPLIES POS-RR(x), the bridge the tree did not have, and `nn-recip-succ-pos' retired as a four-line consequence of it.}\par
\vnbprop{rr-pos-rr-of-lt}{\vnbstmt{\forall x \in \mathbb{R}.\; ((0 < x) \Rightarrow \operatorname{pos\text{-}rr}(x))}}
\vnbprop{nn-recip-succ-pos}{\vnbstmt{\forall n \in \mathbb{N}.\; \operatorname{pos\text{-}rr}(\operatorname{recip}((n + 1)))}}

\vnbbookfile{theorem-library/rr-halving}{theorem-library/rr-halving}
\noindent{\small rr-halving.scm -- every positive real splits into two equal positive halves, PROVEN.}\par
\vnbprop{rr-pos-halvable}{\vnbstmt{\forall s.\; (\operatorname{pos\text{-}rr}(s) \Rightarrow \exists d.\; (\operatorname{pos\text{-}rr}(d) \wedge ((d + d) = s)))}}

\vnbbookfile{theorem-library/rr-order-bundle}{theorem-library/rr-order-bundle}
\noindent{\small rr-order-bundle.scm -- four elementary real-order facts, PROVEN modulo 0.}\par
\vnbprop{rr-prod-nonpos-pos}{\vnbstmt{\forall u, v \in \mathbb{R}.\; (((0 < v) \wedge ((u \cdot v) \leq 0)) \Rightarrow (u \leq 0))}}
\vnbprop{rr-prod-nonpos-neg}{\vnbstmt{\forall u, v \in \mathbb{R}.\; (((v < 0) \wedge ((u \cdot v) \leq 0)) \Rightarrow (0 \leq u))}}
\vnbprop{rr-midpoint-between}{\vnbstmt{\forall u, v \in \mathbb{R}.\; ((u < v) \Rightarrow \exists w \in \mathbb{R}.\; ((u < w) \wedge (w < v)))}}
\vnbprop{rr-le-scale-nonneg-right}{\vnbstmt{\forall x, y, c \in \mathbb{R}.\; (((0 \leq c) \wedge (x \leq y)) \Rightarrow ((x \cdot c) \leq (y \cdot c)))}}

\vnbbookfile{theorem-library/rr-le-all-pos}{theorem-library/rr-le-all-pos}
\noindent{\small rr-le-all-pos.scm -- a real below every positive real is non-positive, PROVEN.}\par
\vnbprop{rr-le-all-pos-nonpos}{\vnbstmt{\forall x \in \mathbb{R}.\; (\forall s.\; (\operatorname{pos\text{-}rr}(s) \Rightarrow (x \leq s)) \Rightarrow (x \leq 0))}}

\vnbbookfile{theorem-library/subset-lemmas}{theorem-library/subset-lemmas}
\noindent{\small subset-lemmas.scm -- the four inclusion facts, PROVEN.}\par
\vnbprop{subset-mem-fwd}{\vnbstmt{\forall a, b, x.\; (((a \subseteq b) \wedge (x \in a)) \Rightarrow (x \in b))}}
\vnbprop{subset-trans}{\vnbstmt{\forall a, b, c.\; (((a \subseteq b) \wedge (b \subseteq c)) \Rightarrow (a \subseteq c))}}
\vnbprop{subclass-of-set-is-set}{\vnbstmt{\forall a.; \forall b \in \mathbf{Set}.\; ((a \subseteq b) \Rightarrow (a \in \mathbf{Set}))}}
\vnbprop{subset-refl}{\vnbstmt{\forall a.\; (a \subseteq a)}}

\vnbbookfile{theorem-library/fun-image-set}{theorem-library/fun-image-set}
\noindent{\small fun-image-set.scm -- the image of ANY class under a function into a SET is a set, proven with no citation of `image-set' (replacement).}\par
\vnbprop{fun-image-set}{\vnbstmt{\forall a, b, f, s.\; (((f \in (a \to b)) \wedge (b \in \mathbf{Set})) \Rightarrow (\operatorname{image}(f, s) \in \mathbf{Set}))}}

\vnbbookfile{theorem-library/rake-analysis2}{theorem-library/rake-analysis2}
\noindent{\small A file shared across books; the entries that belong to this one.}\par
\noindent{\small rake-analysis2.scm -- rake batch L (2026-09-17): five asserted leaves of the analysis / sets area, plus the general lemma that unblocked one of them.}\par
\vnbprop{fun-domain-in-set}{\vnbstmt{\forall a, b.; \forall f \in (a \to b).\; (a \in \mathbf{Set})}}
\vnbprop{bijection-compose}{\vnbstmt{\begin{array}{@{}l@{}} \forall x, y, z, i, a. \\ \quad ((i \in \mathrm{Bij}(x, y)) \wedge (a \in \mathrm{Bij}(y, z))) \Rightarrow  \\ \quad \quad (\lambda b \in x.\; a(i(b))) \in \mathrm{Bij}(x, z) \end{array}}}
\vnbprop{le-bound-mono}{\vnbstmt{\begin{array}{@{}l@{}} \forall p, a, b, n. \\ \quad ((p \in \mathbb{R}) \wedge ((a \in \mathbb{R}) \wedge ((b \in \mathbb{R}) \wedge ((n \in \mathbb{R}) \wedge ((0 \leq n) \wedge ((p \leq (a \cdot n)) \wedge (a \leq b))))))) \Rightarrow  \\ \quad \quad p \leq (b \cdot n) \end{array}}}

\vnbbookfile{theorem-library/rake-inverse-bij}{theorem-library/rake-inverse-bij}
\noindent{\small theorem-library/rake-inverse-bij.scm -- rake batch N (2026-09-17), part 1 of 2. --------------------------------------------------------------------------- NOTE, 2026-09-20 (batch 9-B).  This file was written while cardinality was AXIOMATISED under the name CARD and the defined constant was its companion}\par
\vnbprop{inverse-bij-right}{\vnbstmt{\forall m, d.; \forall i \in \mathrm{Bij}(m, d), y \in d.\; (i(i^{-1}(y)) = y)}}
\vnbprop{inverse-bij-left}{\vnbstmt{\forall m, d.; \forall i \in \mathrm{Bij}(m, d), x \in m.\; (i^{-1}(i(x)) = x)}}
\vnbprop{inverse-bij-in-fun}{\vnbstmt{\forall x, y.; \forall i \in \mathrm{Bij}(x, y).\; (i^{-1} \in (y \to x))}}
\vnbprop{inverse-bij-is-bijection}{\vnbstmt{\forall x, y.; \forall i \in \mathrm{Bij}(x, y).\; (i^{-1} \in \mathrm{Bij}(y, x))}}
\vnbprop{ord-segment-self}{\vnbstmt{\forall n \in \mathbb{N}.\; \neg (n \in \mathrm{OS}(n))}}

\vnbbookfile{theorem-library/nn-parity-proof}{theorem-library/nn-parity-proof}
\noindent{\small nn-parity-proof.scm -- Peano arithmetic on NN: from the two recursion equations (structure-library/nn-arith.scm) to parity and even-square.}\par
\vnbprop{nn-one-is-succ-zero}{\vnbstmt{1 = (0 + 1)}}
\vnbprop{nn-succ-inj}{\vnbstmt{\forall a, b \in \mathbb{N}.\; (((a + 1) = (b + 1)) \Rightarrow (a = b))}}
\vnbprop{nn-succ-nonzero}{\vnbstmt{\forall n \in \mathbb{N}.\; \neg ((n + 1) = 0)}}
\vnbprop{nn-succ-plus-one}{\vnbstmt{\forall n \in \mathbb{N}.\; ((n + 1) = (n + 1))}}
\vnbprop{nn-mul-succ}{\vnbstmt{\forall a, b \in \mathbb{N}.\; ((a \cdot (b + 1)) = ((a \cdot b) + a))}}
\vnbprop{nn-two-is-succ-succ-zero}{\vnbstmt{2 = ((0 + 1) + 1)}}
\vnbprop{nn-plus-two}{\vnbstmt{\forall x \in \mathbb{N}.\; ((x + 2) = ((x + 1) + 1))}}
\vnbprop{nn-two-mul-succ}{\vnbstmt{\forall k \in \mathbb{N}.\; ((2 \cdot (k + 1)) = (((2 \cdot k) + 1) + 1))}}
\vnbprop{nn-nonzero-is-succ}{\vnbstmt{\forall n \in \mathbb{N}.\; (\neg (n = 0) \Rightarrow \exists q \in \mathbb{N}.\; (n = (q + 1)))}}
\vnbprop{nn-parity}{\vnbstmt{\forall n \in \mathbb{N}.\; \exists k \in \mathbb{N}.\; ((n = (2 \cdot k)) \vee (n = ((2 \cdot k) + 1)))}}
\vnbprop{nn-zero-or-succ}{\vnbstmt{\forall y \in \mathbb{N}.\; ((y = 0) \vee \exists q \in \mathbb{N}.\; (y = (q + 1)))}}
\vnbprop{nn-parity-exclusive}{\vnbstmt{\forall x, y \in \mathbb{N}.\; \neg ((2 \cdot x) = ((2 \cdot y) + 1))}}
\vnbprop{nn-even-square}{\vnbstmt{\begin{array}{@{}l@{}} \forall p \in \mathbb{N}. \\ \quad \exists m \in \mathbb{N}.\; ((p \cdot p) = (2 \cdot m)) \Rightarrow \exists k \in \mathbb{N}.\; (p = (2 \cdot k)) \end{array}}}
\vnbprop{bt-succ-minus-1}{\vnbstmt{\forall n \in \mathbb{N}.\; (((n + 1) - 1) = n)}}

\vnbbookfile{theorem-library/nn-pred}{theorem-library/nn-pred}
\noindent{\small nn-pred.scm -- the predecessor on NN, DEFINED by definite description.}\par
\par\medskip\noindent\textbf{Definition} (\texttt{pred}).\;$\operatorname{pred}(\operatorname{n\_}) :=$ $(\iota\, \operatorname{m\_}.\;((\operatorname{m\_} \in \mathbb{N}) \wedge ((\operatorname{m\_} + 1) = \operatorname{n\_})))$
\vnbprop{pred-property}{\vnbstmt{\forall n \in \mathbb{N}.\; (\neg (n = 0) \Rightarrow ((\operatorname{pred}(n) \in \mathbb{N}) \wedge ((\operatorname{pred}(n) + 1) = n)))}}
\vnbprop{pred-in-nn}{\vnbstmt{\forall n \in \mathbb{N}.\; (\neg (n = 0) \Rightarrow (\operatorname{pred}(n) \in \mathbb{N}))}}
\vnbprop{pred-succ}{\vnbstmt{\forall n \in \mathbb{N}.\; (\neg (n = 0) \Rightarrow ((\operatorname{pred}(n) + 1) = n))}}
\vnbprop{pred-of-succ}{\vnbstmt{\forall m \in \mathbb{N}.\; (\operatorname{pred}((m + 1)) = m)}}

\vnbbookfile{theorem-library/ord-segment-arith}{theorem-library/ord-segment-arith}
\noindent{\small ord-segment-arith.scm -- the two bridges between segment membership and arithmetic on NN.}\par
\vnbprop{seg-mem-lt}{\vnbstmt{\forall n, j \in \mathbb{N}.\; ((j \in \mathrm{OS}(n)) \Leftrightarrow (j < n))}}
\vnbprop{seg-mem-succ-le}{\vnbstmt{\forall n, j \in \mathbb{N}.\; ((j \in \mathrm{OS}((n + 1))) \Leftrightarrow (j \leq n))}}

\vnbbookfile{theorem-library/finite-surgery}{theorem-library/finite-surgery}
\noindent{\small finite-surgery.scm -- the finite-surgery kit, first member: COLLAPSE-AT.}\par
\par\medskip\noindent\textbf{Definition} (\texttt{collapse-at}).\;$\operatorname{collapse\text{-}at}(\operatorname{k\_}) :=$ $(\lambda \operatorname{z\_} \in \mathbb{N}.\; \left(\text{if }(\operatorname{k\_} < \operatorname{z\_})\text{ then }\operatorname{pred}(\operatorname{z\_})\text{ else }\operatorname{z\_}\right))$
\vnbprop{nn-lt-succ-le}{\vnbstmt{\forall k, j \in \mathbb{N}.\; ((k < j) \Rightarrow ((k + 1) \leq j))}}
\vnbprop{collapse-at-lo}{\vnbstmt{\forall k, j \in \mathbb{N}.\; (\neg (k < j) \Rightarrow (\operatorname{collapse\text{-}at}(k)(j) = j))}}
\vnbprop{collapse-at-hi}{\vnbstmt{\forall k, j \in \mathbb{N}.\; ((k < j) \Rightarrow (\operatorname{collapse\text{-}at}(k)(j) = \operatorname{pred}(j)))}}
\vnbprop{collapse-at-in-nn}{\vnbstmt{\forall k, j \in \mathbb{N}.\; (\operatorname{collapse\text{-}at}(k)(j) \in \mathbb{N})}}
\vnbprop{collapse-at-lt}{\vnbstmt{\begin{array}{@{}l@{}} \forall n, k, j. \\ \quad ((n \in \mathbb{N}) \wedge ((k \in \mathbb{N}) \wedge ((j \in \mathbb{N}) \wedge ((k \leq n) \wedge ((j \leq n) \wedge \neg (j = k)))))) \Rightarrow  \\ \quad \quad \operatorname{collapse\text{-}at}(k)(j) < n \end{array}}}
\vnbprop{collapse-at-inj}{\vnbstmt{\begin{array}{@{}l@{}} \forall k, a, b. \\ \quad k \in \mathbb{N} \wedge \\ \quad a \in \mathbb{N} \wedge \\ \quad b \in \mathbb{N} \wedge \\ \quad \neg (a = k) \wedge \\ \quad \neg (b = k) \wedge \\ \quad \operatorname{collapse\text{-}at}(k)(a) = \operatorname{collapse\text{-}at}(k)(b) \Rightarrow \\ \quad \quad a = b \end{array}}}
\vnbprop{collapse-at-in-seg}{\vnbstmt{\begin{array}{@{}l@{}} \forall n, k, j. \\ \quad n \in \mathbb{N} \wedge \\ \quad k \in \mathbb{N} \wedge \\ \quad j \in \mathbb{N} \wedge \\ \quad k \in \mathrm{OS}((n + 1)) \wedge \\ \quad j \in \mathrm{OS}((n + 1)) \wedge \\ \quad \neg (j = k) \Rightarrow \\ \quad \quad \operatorname{collapse\text{-}at}(k)(j) \in \mathrm{OS}(n) \end{array}}}
\vnbprop{nn-le-ne-succ}{\vnbstmt{\forall a, k \in \mathbb{N}.\; (((a \leq k) \wedge \neg (k = a)) \Rightarrow ((a + 1) \leq k))}}

\vnbbookfile{theorem-library/pigeonhole-segments}{theorem-library/pigeonhole-segments}
\noindent{\small pigeonhole-segments.scm -- Track A of the CARD plan: finite pigeonhole for ordinal segments, the one piece of real mathematics the plan needs.}\par
\vnbprop{fun-into-seg-zero-absurd}{\vnbstmt{\forall a, f, x.\; (((f \in (a \to \mathrm{OS}(0))) \wedge (x \in a)) \Rightarrow \mathit{falsity})}}
\vnbprop{no-injection-seg-1-into-seg-0}{\vnbstmt{\forall f.\; \neg (f \in \operatorname{injection}(\mathrm{OS}(1), \mathrm{OS}(0)))}}
\vnbprop{pigeonhole-step}{\vnbstmt{\begin{array}{@{}l@{}} \forall n \in \mathbb{N}, g. \\ \quad \forall f.\; \neg (f \in \operatorname{injection}(\mathrm{OS}((n + 1)), \mathrm{OS}(n))) \Rightarrow  \\ \quad \quad \neg (g \in \operatorname{injection}(\mathrm{OS}(((n + 1) + 1)), \mathrm{OS}((n + 1)))) \end{array}}}
\vnbprop{pigeonhole-segments}{\vnbstmt{\forall n \in \mathbb{N}, f.\; \neg (f \in \operatorname{injection}(\mathrm{OS}((n + 1)), \mathrm{OS}(n)))}}
\vnbprop{pigeonhole-segments-gen}{\vnbstmt{\begin{array}{@{}l@{}} \forall m, n, f. \\ \quad ((m \in \mathbb{N}) \wedge ((n \in \mathbb{N}) \wedge (n < m))) \Rightarrow  \\ \quad \quad \neg (f \in \operatorname{injection}(\mathrm{OS}(m), \mathrm{OS}(n))) \end{array}}}

\vnbbookfile{theorem-library/enum-append}{theorem-library/enum-append}
\noindent{\small theorem-library/enum-append.scm -- enum-append-is-bijection, moved out of theorem-library/finsum-insert.scm on 2026-09-20 (batch 9-B, CARD := CARD-STAR).}\par
\vnbprop{enum-append-is-bijection}{\vnbstmt{\begin{array}{@{}l@{}} \forall n \in \mathbb{N}, g, t \in \mathbf{Set}, m \in \mathrm{Bij}(\mathrm{OS}(n), g). \\ \quad \neg (t \in g) \Rightarrow  \\ \quad \quad (\lambda i \in \mathrm{OS}((n + 1)).\; \left(\text{if }(i \in \mathrm{OS}(n))\text{ then }m(i)\text{ else }t\right)) \\ \quad \quad \in\; \mathrm{Bij}(\mathrm{OS}((n + 1)), (g \cup \{t\})) \end{array}}}

\vnbbookfile{theorem-library/union-laws}{theorem-library/union-laws}
\noindent{\small theorem-library/union-laws.scm -- union-empty-right and union-assoc, moved out of theorem-library/fin-subsets.scm on 2026-09-20 (batch 9-B, CARD := CARD-STAR).}\par
\vnbprop{union-empty-right}{\vnbstmt{\forall a.\; ((a \cup \emptyset) = a)}}
\vnbprop{union-assoc}{\vnbstmt{\forall a, b, c.\; (((a \cup b) \cup c) = (a \cup (b \cup c)))}}

\vnbbookfile{structure-library/bijection-derived}{structure-library/bijection-derived}
\noindent{\small bijection-derived.scm -- the three BIJECTION projection lemmas, PROVEN.}\par
\vnbprop{bijection-in-fun}{\vnbstmt{\forall x, y.; \forall i \in \mathrm{Bij}(x, y).\; (i \in (x \to y))}}
\vnbprop{bijection-injective}{\vnbstmt{\forall x, y.; \forall i \in \mathrm{Bij}(x, y), a, b \in x.\; ((i(a) = i(b)) \Rightarrow (a = b))}}
\vnbprop{bijection-surjective}{\vnbstmt{\forall x, y.; \forall i \in \mathrm{Bij}(x, y), w \in y.\; \exists z \in x.\; (i(z) = w)}}

\vnbbookfile{theorem-library/bijection-identity-proof}{theorem-library/bijection-identity-proof}
\noindent{\small bijection-identity-proof.scm -- the identity lambda on a SET is a bijection.}\par
\vnbprop{bijection-identity}{\vnbstmt{\forall x \in \mathbf{Set}.\; ((\lambda a \in x.\; a) \in \mathrm{Bij}(x, x))}}

\vnbbookfile{theorem-library/card-defined}{theorem-library/card-defined}
\noindent{\small card-defined.scm -- CARD, cardinality DEFINED rather than axiomatised.}\par
\vnbprop{cd-body-unique}{\vnbstmt{\begin{array}{@{}l@{}} \forall a, l, e. \\ \quad l \in \mathrm{Ord} \wedge \\ \quad \exists i.\; (i \in \mathrm{Bij}(a, \mathrm{OS}(l))) \wedge \\ \quad \forall b. \\ \quad \quad \operatorname{ord\text{-}lt}(b, l) \Rightarrow \neg \exists c.\; (c \in \mathrm{Bij}(a, \mathrm{OS}(b))) \wedge \\ \quad e \in \mathrm{Ord} \wedge \\ \quad \exists i.\; (i \in \mathrm{Bij}(a, \mathrm{OS}(e))) \wedge \\ \quad \forall b. \\ \quad \quad \operatorname{ord\text{-}lt}(b, e) \Rightarrow \neg \exists c.\; (c \in \mathrm{Bij}(a, \mathrm{OS}(b))) \Rightarrow \\ \quad \quad l = e \end{array}}}
\vnbprop{cd-seg-body}{\vnbstmt{\begin{array}{@{}l@{}} \forall n \in \mathbb{N}. \\ \quad n \in \mathrm{Ord} \wedge \\ \quad \exists i.\; (i \in \mathrm{Bij}(\mathrm{OS}(n), \mathrm{OS}(n))) \wedge \\ \quad \forall a. \\ \quad \quad \operatorname{ord\text{-}lt}(a, n) \Rightarrow  \\ \quad \quad \quad \neg \exists b.\; (b \in \mathrm{Bij}(\mathrm{OS}(n), \mathrm{OS}(a))) \end{array}}}
\vnbprop{card-segment}{\vnbstmt{\forall n \in \mathbb{N}.\; (|\mathrm{OS}(n)| = n)}}

\vnbbookfile{theorem-library/card-finite}{theorem-library/card-finite}
\noindent{\small card-finite.scm -- the finite layer of the DEFINED cardinal, CARD. --------------------------------------------------------------------------- NOTE, 2026-09-20 (batch 9-B).  This file was written while cardinality was AXIOMATISED under the name CARD and the defined constant was its companion}\par
\vnbprop{card-from-body}{\vnbstmt{\begin{array}{@{}l@{}} \forall a.; \forall l \in \mathrm{Ord}. \\ \quad \exists i.\; (i \in \mathrm{Bij}(a, \mathrm{OS}(l))) \wedge \\ \quad \forall b. \\ \quad \quad \operatorname{ord\text{-}lt}(b, l) \Rightarrow \neg \exists c.\; (c \in \mathrm{Bij}(a, \mathrm{OS}(b))) \Rightarrow \\ \quad \quad |a| = l \end{array}}}
\vnbprop{card-bij}{\vnbstmt{\begin{array}{@{}l@{}} \forall n \in \mathbb{N}, a, h, t. \\ \quad ((h \in \mathrm{Bij}(a, \mathrm{OS}(n))) \wedge (t \in \mathrm{Bij}(\mathrm{OS}(n), a))) \Rightarrow  \\ \quad \quad |a| = n \end{array}}}
\vnbprop{seg0-empty}{\vnbstmt{\emptyset = \mathrm{OS}(0)}}
\vnbprop{card-empty}{\vnbstmt{|\emptyset| = 0}}

\vnbbookfile{theorem-library/ord-no-injection}{theorem-library/ord-no-injection}
\noindent{\small ord-no-injection.scm -- NO SET RECEIVES AN INJECTIVE CLASS FUNCTION FROM ORD.}\par
\vnbprop{ord-no-injection-into-set}{\vnbstmt{\begin{array}{@{}l@{}} \forall d, i. \\ \quad d \in \mathbf{Set} \wedge \\ \quad \forall a \in \mathrm{Ord}.\; (i(a) \in d) \wedge \\ \quad \forall a_{1}, a_{2} \in \mathrm{Ord}.\; ((i(a_{1}) = i(a_{2})) \Rightarrow (a_{1} = a_{2})) \Rightarrow \\ \quad \quad \mathit{falsity} \end{array}}}

\vnbbookfile{theorem-library/zen-step}{theorem-library/zen-step}
\noindent{\small zen-step.scm -- ZEN's three recursion equations, collapsed into ONE.}\par
\par\medskip\noindent\textbf{Definition} (\texttt{zen}).\;For every $x$,  $(\operatorname{zen}(x, 0) \simeq \varepsilon(x))$\par
Let $x$ be arbitrary. Let $\alpha \in \mathrm{Ord}$. Then
\begin{equation*}
(\operatorname{zen}(x, (\alpha + 1)) \simeq \varepsilon(\{\operatorname{y\_} \in x : \forall \operatorname{c\_} \in \mathrm{Ord}.\; (\operatorname{ord\text{-}le}(\operatorname{c\_}, \alpha) \Rightarrow \neg (\operatorname{zen}(x, \operatorname{c\_}) = \operatorname{y\_}))\}))
\end{equation*}
\par
Let $x$ be arbitrary. Let $\mathit{lam} \text{ be s.t. } \operatorname{limit\text{-}ord}(\mathit{lam})$. Then
\begin{equation*}
(\operatorname{zen}(x, \mathit{lam}) \simeq \varepsilon(\{\operatorname{y\_} \in x : \forall \operatorname{c\_} \in \mathrm{Ord}.\; (\operatorname{ord\text{-}lt}(\operatorname{c\_}, \mathit{lam}) \Rightarrow \neg (\operatorname{zen}(x, \operatorname{c\_}) = \operatorname{y\_}))\}))
\end{equation*}

\vnbprop{ord-lt-succ-iff-le}{\vnbstmt{\forall l, z \in \mathrm{Ord}.\; (\operatorname{ord\text{-}lt}(z, (l + 1)) \Leftrightarrow \operatorname{ord\text{-}le}(z, l))}}
\vnbprop{zen-step}{\vnbstmt{\begin{array}{@{}l@{}} \forall a \in \mathrm{Ord}, x. \\ \quad (\operatorname{zen}(x, a) \simeq \varepsilon(\{y \in x : \forall c \in \mathrm{Ord}.\; (\operatorname{ord\text{-}lt}(c, a) \Rightarrow \neg (\operatorname{zen}(x, c) = y))\})) \end{array}}}
\vnbprop{zen-hits}{\vnbstmt{\begin{array}{@{}l@{}} \forall a \in \mathrm{Ord}, s. \\ \quad \exists x.\; (x \in \{y \in s : \forall c \in \mathrm{Ord}.\; (\operatorname{ord\text{-}lt}(c, a) \Rightarrow \neg (\operatorname{zen}(s, c) = y))\}) \Rightarrow  \\ \quad \quad \operatorname{zen}(s, a) \in \{y \in s : \forall c \in \mathrm{Ord}.\; (\operatorname{ord\text{-}lt}(c, a) \Rightarrow \neg (\operatorname{zen}(s, c) = y))\} \end{array}}}
\vnbprop{zen-exhausts}{\vnbstmt{\begin{array}{@{}l@{}} \forall s \in \mathbf{Set}. \exists a \in \mathrm{Ord}. \\ \quad \neg \exists x.\; (x \in \{y \in s : \forall c \in \mathrm{Ord}.\; (\operatorname{ord\text{-}lt}(c, a) \Rightarrow \neg (\operatorname{zen}(s, c) = y))\}) \end{array}}}

\vnbbookfile{theorem-library/rake-zermelo}{theorem-library/rake-zermelo}
\noindent{\small rake-zermelo.scm -- CARD-free Zermelo: every set bijects with an ordinal --------------------------------------------------------------------------- NOTE, 2026-09-20 (batch 9-B).  This file was written while cardinality was AXIOMATISED under the name CARD and the defined constant was its companion}\par
\vnbprop{zermelo-least-ordinal}{\vnbstmt{\begin{array}{@{}l@{}} \forall s \in \mathbf{Set}. \exists l \in \mathrm{Ord}. \\ \quad \neg \exists x.\; (x \in \{y \in s : \forall c \in \mathrm{Ord}.\; (\operatorname{ord\text{-}lt}(c, l) \Rightarrow \neg (\operatorname{zen}(s, c) = y))\}) \wedge \\ \quad \forall z. \\ \quad \quad \operatorname{ord\text{-}lt}(z, l) \Rightarrow  \\ \quad \quad \quad \exists x. \\ \quad \quad \quad \quad x \in \{y \in s : \forall c \in \mathrm{Ord}.\; (\operatorname{ord\text{-}lt}(c, z) \Rightarrow \neg (\operatorname{zen}(s, c) = y))\} \end{array}}}
\vnbprop{zen-enum-in-fun}{\vnbstmt{\begin{array}{@{}l@{}} \forall l \in \mathrm{Ord}, s \in \mathbf{Set}. \\ \quad \forall z. \\ \quad \quad \operatorname{ord\text{-}lt}(z, l) \Rightarrow  \\ \quad \quad \quad \exists x. \\ \quad \quad \quad \quad x \in \{y \in s : \forall c \in \mathrm{Ord}.\; (\operatorname{ord\text{-}lt}(c, z) \Rightarrow \neg (\operatorname{zen}(s, c) = y))\} \Rightarrow \\ \quad \quad (\lambda a \in \mathrm{OS}(l).\; \operatorname{zen}(s, a)) \in (\mathrm{OS}(l) \to s) \end{array}}}
\vnbprop{zen-enum-injective}{\vnbstmt{\begin{array}{@{}l@{}} \forall l \in \mathrm{Ord}, s \in \mathbf{Set}, a_{1}, a_{2} \in \mathrm{OS}(l). \\ \quad \forall z. \\ \quad \quad \operatorname{ord\text{-}lt}(z, l) \Rightarrow  \\ \quad \quad \quad \exists x. \\ \quad \quad \quad \quad x \in \{y \in s : \forall c \in \mathrm{Ord}.\; (\operatorname{ord\text{-}lt}(c, z) \Rightarrow \neg (\operatorname{zen}(s, c) = y))\} \wedge \\ \quad (\lambda a \in \mathrm{OS}(l).\; \operatorname{zen}(s, a))(a_{1}) = (\lambda a \in \mathrm{OS}(l).\; \operatorname{zen}(s, a))(a_{2}) \Rightarrow \\ \quad \quad a_{1} = a_{2} \end{array}}}
\vnbprop{zen-enum-surjective}{\vnbstmt{\begin{array}{@{}l@{}} \forall l \in \mathrm{Ord}, s \in \mathbf{Set}, w \in s. \\ \quad \neg \exists x.\; (x \in \{y \in s : \forall c \in \mathrm{Ord}.\; (\operatorname{ord\text{-}lt}(c, l) \Rightarrow \neg (\operatorname{zen}(s, c) = y))\}) \Rightarrow  \\ \quad \quad \exists z \in \mathrm{OS}(l). \\ \quad \quad \quad (\lambda a \in \mathrm{OS}(l).\; \operatorname{zen}(s, a))(z) = w \end{array}}}
\vnbprop{zermelo-bijection}{\vnbstmt{\forall s \in \mathbf{Set}.\; \exists l \in \mathrm{Ord}, i.\; (i \in \mathrm{Bij}(\mathrm{OS}(l), s))}}

\vnbbookfile{theorem-library/rake-ord-pigeonhole}{theorem-library/rake-ord-pigeonhole}
\noindent{\small rake-ord-pigeonhole.scm -- rake batch 5c, assignment 5c-U (2026-09-18). --------------------------------------------------------------------------- NOTE, 2026-09-20 (batch 9-B).  This file was written while cardinality was AXIOMATISED under the name CARD and the defined constant was its companion}\par
\vnbprop{ord-segment-pigeonhole-guarded}{\vnbstmt{\begin{array}{@{}l@{}} \forall l, e \in \mathbb{N}, f. \\ \quad \operatorname{ord\text{-}lt}(e, l) \Rightarrow \neg (f \in \operatorname{injection}(\mathrm{OS}(l), \mathrm{OS}(e))) \end{array}}}
\vnbprop{card-body-exists}{\vnbstmt{\begin{array}{@{}l@{}} \forall s \in \mathbf{Set}. \exists l \in \mathrm{Ord}. \\ \quad \exists i.\; (i \in \mathrm{Bij}(s, \mathrm{OS}(l))) \wedge \\ \quad \forall a. \\ \quad \quad \operatorname{ord\text{-}lt}(a, l) \Rightarrow \neg \exists b.\; (b \in \mathrm{Bij}(s, \mathrm{OS}(a))) \end{array}}}
\vnbprop{card-zermelo}{\vnbstmt{\begin{array}{@{}l@{}} \forall s \in \mathbf{Set}. \\ \quad (|s| \in \mathrm{Ord}) \wedge \exists i.\; (i \in \mathrm{Bij}(s, \mathrm{OS}(|s|))) \end{array}}}
\vnbprop{well-ordering-principle}{\vnbstmt{\forall s \in \mathbf{Set}.\; \exists i.\; (i \in \mathrm{Bij}(\mathrm{OS}(|s|), s))}}

\vnbbookfile{theorem-library/rake-card-star-laws}{theorem-library/rake-card-star-laws}
\noindent{\small rake-card-star-laws.scm -- rake batch 5c, assignment 5c-X (2026-09-18). --------------------------------------------------------------------------- NOTE, 2026-09-20 (batch 9-B).  This file was written while cardinality was AXIOMATISED under the name CARD and the defined constant was its companion}\par
\vnbprop{card-insert-curried}{\vnbstmt{\begin{array}{@{}l@{}} \forall a, x \in \mathbf{Set}. \\ \quad ((|a| \in \mathbb{N}) \wedge \neg (x \in a)) \Rightarrow  \\ \quad \quad |(a \cup \{x\})| = (|a| + 1) \end{array}}}
\vnbprop{seg-restrict-is-bijection}{\vnbstmt{\begin{array}{@{}l@{}} \forall n \in \mathbb{N}, b.; \forall g \in \mathrm{Bij}(\mathrm{OS}((n + 1)), b). \\ \quad (\lambda i \in \mathrm{OS}(n).\; g(i)) \in \mathrm{Bij}(\mathrm{OS}(n), \operatorname{image}(g, \mathrm{OS}(n))) \end{array}}}
\vnbprop{card-seg-image}{\vnbstmt{\begin{array}{@{}l@{}} \forall n \in \mathbb{N}, b.; \forall g \in \mathrm{Bij}(\mathrm{OS}((n + 1)), b). \\ \quad |\operatorname{image}(g, \mathrm{OS}(n))| = n \end{array}}}
\vnbprop{seg-image-subset}{\vnbstmt{\begin{array}{@{}l@{}} \forall n \in \mathbb{N}, b.; \forall g \in \mathrm{Bij}(\mathrm{OS}((n + 1)), b), w \in \operatorname{image}(g, \mathrm{OS}(n)). \\ \quad w \in b \end{array}}}
\vnbprop{seg-image-omits-top}{\vnbstmt{\begin{array}{@{}l@{}} \forall n \in \mathbb{N}, b.; \forall g \in \mathrm{Bij}(\mathrm{OS}((n + 1)), b). \\ \quad \neg (g(n) \in \operatorname{image}(g, \mathrm{OS}(n))) \end{array}}}
\vnbprop{seg-image-insert}{\vnbstmt{\begin{array}{@{}l@{}} \forall n \in \mathbb{N}, b.; \forall g \in \mathrm{Bij}(\mathrm{OS}((n + 1)), b). \\ \quad b = (\operatorname{image}(g, \mathrm{OS}(n)) \cup \{g(n)\}) \end{array}}}
\vnbprop{card-zero-is-empty}{\vnbstmt{\forall b \in \mathbf{Set}.\; ((|b| = 0) \Rightarrow (b = \emptyset))}}
\vnbprop{intersection-empty-transfer}{\vnbstmt{\begin{array}{@{}l@{}} \forall a \in \mathbf{Set}, b, c. \\ \quad (((a \cap b) = \emptyset) \wedge \forall w \in c.\; (w \in b)) \Rightarrow  \\ \quad \quad (a \cap c) = \emptyset \end{array}}}
\vnbprop{card-union-disjoint-ind}{\vnbstmt{\begin{array}{@{}l@{}} \forall n \in \mathbb{N}, a, b \in \mathbf{Set}. \\ \quad ((|a| \in \mathbb{N}) \wedge ((|b| = n) \wedge ((a \cap b) = \emptyset))) \Rightarrow  \\ \quad \quad |(a \cup b)| = (|a| + n) \end{array}}}
\vnbprop{card-union-disjoint-curried}{\vnbstmt{\begin{array}{@{}l@{}} \forall a, b \in \mathbf{Set}. \\ \quad ((|a| \in \mathbb{N}) \wedge ((|b| \in \mathbb{N}) \wedge ((a \cap b) = \emptyset))) \Rightarrow  \\ \quad \quad |(a \cup b)| = (|a| + |b|) \end{array}}}
\vnbprop{card-finite-induction-aux}{\vnbstmt{\begin{array}{@{}l@{}} \forall n \in \mathbb{N}, c.; \forall s \in \mathbf{Set}. \\ \quad \emptyset \in c \wedge \\ \quad \forall s, x \in \mathbf{Set}. \\ \quad \quad (((|s| \in \mathbb{N}) \wedge (s \in c)) \wedge \neg (x \in s)) \Rightarrow ((s \cup \{x\}) \in c) \wedge \\ \quad |s| = n \Rightarrow \\ \quad \quad s \in c \end{array}}}
\vnbprop{finite-set-induction}{\vnbstmt{\begin{array}{@{}l@{}} \forall c.; \forall s \in \mathbf{Set}. \\ \quad \emptyset \in c \wedge \\ \quad \forall s, x \in \mathbf{Set}. \\ \quad \quad (((|s| \in \mathbb{N}) \wedge (s \in c)) \wedge \neg (x \in s)) \Rightarrow ((s \cup \{x\}) \in c) \wedge \\ \quad |s| \in \mathbb{N} \Rightarrow \\ \quad \quad s \in c \end{array}}}
\vnbprop{injection-image-is-bijection}{\vnbstmt{\begin{array}{@{}l@{}} \forall m \in \mathbf{Set}, d.; \forall h \in \operatorname{injection}(m, d). \\ \quad (\lambda z \in m.\; h(z)) \in \mathrm{Bij}(m, \operatorname{image}(h, m)) \end{array}}}
\vnbprop{card-image-injection-curried}{\vnbstmt{\begin{array}{@{}l@{}} \forall m \in \mathbf{Set}, d.; \forall h \in \operatorname{injection}(m, d). \\ \quad (|m| \in \mathbb{N}) \Rightarrow (|\operatorname{image}(h, m)| = |m|) \end{array}}}

\vnbbookfile{theorem-library/card-laws}{theorem-library/card-laws}
\noindent{\small theorem-library/card-laws.scm -- the former CARD axioms whose counterpart has a different SHAPE, restated character for character as the axioms stood and proven (2026-09-20, batch 9-B: CARD := CARD-STAR, the user's decision).}\par
\vnbprop{card-in-ord}{\vnbstmt{\forall a \in \mathbf{Set}.\; (|a| \in \mathrm{Ord})}}
\vnbprop{card-insert}{\vnbstmt{\begin{array}{@{}l@{}} \forall a, x \in \mathbf{Set}. \\ \quad ((|a| \in \mathbb{N}) \wedge \neg (x \in a)) \Rightarrow  \\ \quad \quad |(a \cup \{x\})| = (|a| + 1) \end{array}}}
\vnbprop{card-finite-bij}{\vnbstmt{\begin{array}{@{}l@{}} \forall a \in \mathbf{Set}. \\ \quad (|a| \in \mathbb{N}) \Rightarrow \exists i.\; (i \in \mathrm{Bij}(\mathrm{OS}(|a|), a)) \end{array}}}
\vnbprop{card-union-disjoint}{\vnbstmt{\begin{array}{@{}l@{}} \forall a, b \in \mathbf{Set}. \\ \quad ((|a| \in \mathbb{N}) \wedge ((|b| \in \mathbb{N}) \wedge ((a \cap b) = \emptyset))) \Rightarrow  \\ \quad \quad |(a \cup b)| = (|a| + |b|) \end{array}}}
\vnbprop{card-image-injection}{\vnbstmt{\begin{array}{@{}l@{}} \forall m, d.; \forall i \in \operatorname{injection}(m, d). \\ \quad ((m \in \mathbf{Set}) \wedge (|m| \in \mathbb{N})) \Rightarrow (|\operatorname{image}(i, m)| = |m|) \end{array}}}

\vnbbookfile{theorem-library/card-singleton-proof}{theorem-library/card-singleton-proof}
\noindent{\small card-singleton-proof.scm -- the cardinality of a singleton is 1, proven.}\par
\vnbprop{card-singleton}{\vnbstmt{\forall x \in \mathbf{Set}.\; (|\{x\}| = 1)}}

\vnbbookfile{theorem-library/rake-fin-enum}{theorem-library/rake-fin-enum}
\noindent{\small theorem-library/rake-fin-enum.scm -- fin-enum-is-bijection, moved out of theorem-library/rake-inverse-bij.scm on 2026-09-20 (batch 9-B, CARD := CARD-STAR).}\par
\vnbprop{fin-enum-is-bijection}{\vnbstmt{\begin{array}{@{}l@{}} \forall s \in \mathbf{Set}. \\ \quad (|s| \in \mathbb{N}) \Rightarrow (\operatorname{fin\text{-}enum}(s) \in \mathrm{Bij}(\mathrm{OS}(|s|), s)) \end{array}}}

\vnbbookfile{theorem-library/difference-laws}{theorem-library/difference-laws}
\noindent{\small difference-laws.scm -- the two DIFFERENCE laws, PROVEN.}\par
\vnbprop{difference-membership}{\vnbstmt{\forall u, b, x.\; ((x \in \operatorname{difference}(u, b)) \Leftrightarrow ((x \in u) \wedge \neg (x \in b)))}}
\vnbprop{difference-set}{\vnbstmt{\forall x \in \mathbf{Set}, b.\; (\operatorname{difference}(x, b) \in \mathbf{Set})}}

\vnbbookfile{theorem-library/difference-subset}{theorem-library/difference-subset}
\noindent{\small difference-subset.scm -- 2026-11-10 (integration): DIFFERENCE(g, b) is a subset of g, in library core. The same statement as meta-core-laws' ctx-diff-sub (library meta), proved the same way; ba-gelfand-formula (operator-algebras) cites this one, since operator-algebras does not import meta.}\par
\vnbprop{difference-subset}{\vnbstmt{\forall g, b.\; (\operatorname{difference}(g, b) \subseteq g)}}

\vnbbookfile{theorem-library/intersection-of-laws}{theorem-library/intersection-of-laws}
\noindent{\small intersection-of-laws.scm -- the laws of INTERSECTION-OF, PROVEN.}\par
\vnbprop{intersection-of-unfold}{\vnbstmt{\begin{array}{@{}l@{}} \forall c. \\ \quad (\operatorname{intersection\text{-}of}(c) \simeq \{v \in \operatorname{big\text{-}union}(w, c, w) : \forall u \in c.\; (v \in u)\}) \end{array}}}
\vnbprop{intersection-of-membership}{\vnbstmt{\begin{array}{@{}l@{}} \forall c, x. \\ \quad (x \in \operatorname{intersection\text{-}of}(c)) \Leftrightarrow (\exists u.\; (u \in c) \wedge \forall u \in c.\; (x \in u)) \end{array}}}
\vnbprop{intersection-of-subset-member}{\vnbstmt{\forall c.; \forall a \in c.\; (\operatorname{intersection\text{-}of}(c) \subseteq a)}}
\vnbprop{intersection-of-empty-family}{\vnbstmt{\forall c.\; (\neg \exists u.\; (u \in c) \Rightarrow (\operatorname{intersection\text{-}of}(c) = \emptyset))}}
\vnbprop{intersection-of-is-set}{\vnbstmt{\forall c.\; (\operatorname{intersection\text{-}of}(c) \in \mathbf{Set})}}
\vnbprop{intersection-of-antitone}{\vnbstmt{\begin{array}{@{}l@{}} \forall c, d. \\ \quad ((c \subseteq d) \wedge \exists u.\; (u \in c)) \Rightarrow  \\ \quad \quad \operatorname{intersection\text{-}of}(d) \subseteq \operatorname{intersection\text{-}of}(c) \end{array}}}
\vnbprop{complement-in-big-union}{\vnbstmt{\begin{array}{@{}l@{}} \forall y, c, x. \\ \quad (x \in \operatorname{complement\text{-}in}(y, \operatorname{big\text{-}union}(w, c, w))) \Leftrightarrow ((x \in y) \wedge \forall u \in c.\; \neg (x \in u)) \end{array}}}
\vnbprop{complement-in-intersection-of}{\vnbstmt{\begin{array}{@{}l@{}} \forall y, c. \\ \quad \exists u.\; (u \in c) \Rightarrow  \\ \quad \quad \operatorname{complement\text{-}in}(y, \operatorname{intersection\text{-}of}(c)) = \operatorname{big\text{-}union}(w, c, \operatorname{complement\text{-}in}(y, w)) \end{array}}}

\vnbbookfile{theorem-library/rake-complement-laws}{theorem-library/rake-complement-laws}
\noindent{\small rake-complement-laws.scm -- two set-theoretic bricks the compactness duality needs, PROVEN (2026-09-20, batch 11-E; helper prefix `rcl-').}\par
\vnbprop{complement-in-double}{\vnbstmt{\forall y, a.\; ((a \subseteq y) \Rightarrow (\operatorname{complement\text{-}in}(y, \operatorname{complement\text{-}in}(y, a)) = a))}}
\vnbprop{family-of-subsets-is-set}{\vnbstmt{\forall y, c.\; (((y \in \mathbf{Set}) \wedge \forall u \in c.\; (u \subseteq y)) \Rightarrow (c \in \mathbf{Set}))}}

\vnbbookfile{theorem-library/rake-fun-codomain}{theorem-library/rake-fun-codomain}
\noindent{\small theorem-library/rake-fun-codomain.scm -- codomain widening for FUN, PROVEN.}\par
\vnbprop{fun-codomain-superset}{\vnbstmt{\forall f, a, b, c.\; (((f \in (a \to b)) \wedge (b \subseteq c)) \Rightarrow (f \in (a \to c)))}}

\vnbbookfile{theorem-library/rake-rr-order-leaves}{theorem-library/rake-rr-order-leaves}
\noindent{\small theorem-library/rake-rr-order-leaves.scm -- the two asserted order leaves of `deriv-pos-strictly-increasing', PROVEN.}\par
\vnbprop{rr-lt-from-diff-pos}{\vnbstmt{\forall x, y \in \mathbb{R}.\; ((0 < (y - x)) \Rightarrow (x < y))}}

\vnbbookfile{theorem-library/rr-sup-approx}{theorem-library/rr-sup-approx}
\noindent{\small rr-sup-approx.scm -- the SUP APPROXIMATION lemma, PROVEN.}\par
\vnbprop{rr-sup-approx}{\vnbstmt{\begin{array}{@{}l@{}} \forall s, d. \\ \quad ((s \subseteq \mathbb{R}) \wedge (\exists x.\; (x \in s) \wedge (\operatorname{rr\text{-}bounded\text{-}above}(s) \wedge ((d \in \mathbb{R}) \wedge (0 < d))))) \Rightarrow  \\ \quad \quad \exists w \in s.\; ((\operatorname{sup}(s) - d) < w) \end{array}}}

\vnbbookfile{theorem-library/nn-unbounded-in-rr}{theorem-library/nn-unbounded-in-rr}
\noindent{\small nn-unbounded-in-rr.scm -- the ARCHIMEDEAN PROPERTY, PROVEN.}\par
\vnbprop{nn-unbounded-in-rr}{\vnbstmt{\forall x \in \mathbb{R}.\; \exists n \in \mathbb{N}.\; (x < n)}}

\vnbbookfile{theorem-library/nn-recip-succ-small}{theorem-library/nn-recip-succ-small}
\noindent{\small nn-recip-succ-small.scm -- 1/(n+1) gets below any positive eps, PROVEN.}\par
\vnbprop{nn-recip-succ-small}{\vnbstmt{\forall s.\; (\operatorname{pos\text{-}rr}(s) \Rightarrow \exists n \in \mathbb{N}.\; (\operatorname{recip}((n + 1)) < s))}}

\vnbbookfile{theorem-library/discrete-space}{theorem-library/discrete-space}
\noindent{\small A file shared across books; the entries that belong to this one.}\par
\noindent{\small discrete-space.scm -- THE FIRST WITNESS OF A MEASURABLE SPACE, and the first outright inhabitant of TOP-SPACE.}\par
\vnbprop{power-mem-sethood}{\vnbstmt{\forall a.; \forall b \in \operatorname{power}(a).\; (b \in \mathbf{Set})}}
\vnbprop{power-mem-in}{\vnbstmt{\forall a, b, z.\; (((b \in \operatorname{power}(a)) \wedge (z \in b)) \Rightarrow (z \in a))}}
\vnbprop{power-mem-intro}{\vnbstmt{\forall a.; \forall b \in \mathbf{Set}.\; (\forall z \in b.\; (z \in a) \Rightarrow (b \in \operatorname{power}(a)))}}
\vnbprop{power-whole-in}{\vnbstmt{\forall a \in \mathbf{Set}.\; (a \in \operatorname{power}(a))}}
\vnbprop{power-empty-in}{\vnbstmt{\forall a.\; (\emptyset \in \operatorname{power}(a))}}
\vnbprop{power-inter-closed}{\vnbstmt{\begin{array}{@{}l@{}} \forall a, u, v. \\ \quad ((u \in \operatorname{power}(a)) \wedge (v \in \operatorname{power}(a))) \Rightarrow ((u \cap v) \in \operatorname{power}(a)) \end{array}}}
\vnbprop{power-complement-closed}{\vnbstmt{\forall a \in \mathbf{Set}, b.\; (\operatorname{complement\text{-}in}(a, b) \in \operatorname{power}(a))}}
\vnbprop{power-big-union-closed}{\vnbstmt{\begin{array}{@{}l@{}} \forall a, m. \\ \quad ((a \in \mathbf{Set}) \wedge (m \subseteq \operatorname{power}(a))) \Rightarrow (\operatorname{big\text{-}union}(u, m, u) \in \operatorname{power}(a)) \end{array}}}
\vnbprop{power-fun-union-closed}{\vnbstmt{\forall a.; \forall f \in (\mathbb{N} \to \operatorname{power}(a)).\; (\operatorname{big\text{-}union}(n, \mathbb{N}, f(n)) \in \operatorname{power}(a))}}

\vnbbookfile{structure-library/subtype-laws}{structure-library/subtype-laws}
\noindent{\small A file shared across books; the entries that belong to this one.}\par
\noindent{\small subtype-laws.scm -- the trivial subtype-subsumption laws, PROVEN.}\par
\vnbprop{bijection-is-injection}{\vnbstmt{\forall x, y.; \forall i \in \mathrm{Bij}(x, y).\; (i \in \operatorname{injection}(x, y))}}

\vnbbookfile{theorem-library/lambda-slot-apply}{theorem-library/lambda-slot-apply}
\noindent{\small lambda-slot-apply.scm -- taking an operation SLOT down to the surface.}\par
\vnbprop{lam-slot-add-apply}{\vnbstmt{\begin{array}{@{}l@{}} \forall m, a, b. \\ \quad ((a \in m) \wedge (b \in m)) \Rightarrow  \\ \quad \quad ((\lambda {[}\operatorname{x\_}, \operatorname{y\_}] \in (m \times m).\; (\operatorname{x\_} + \operatorname{y\_}))(a, b) \simeq (a + b)) \end{array}}}
\vnbprop{lam-slot-mul-apply}{\vnbstmt{\begin{array}{@{}l@{}} \forall m, a, b. \\ \quad ((a \in m) \wedge (b \in m)) \Rightarrow  \\ \quad \quad ((\lambda {[}\operatorname{x\_}, \operatorname{y\_}] \in (m \times m).\; (\operatorname{x\_} \cdot \operatorname{y\_}))(a, b) \simeq (a \cdot b)) \end{array}}}
\vnbprop{lam-slot-neg-apply}{\vnbstmt{\forall m.; \forall a \in m.\; ((\lambda x \in m.\; -x)(a) \simeq -a)}}

\vnbbookfile{theorem-library/nn-cancel}{theorem-library/nn-cancel}
\noindent{\small nn-cancel.scm -- CANCELLATION ON NN, proven in the core (2026-10-07).}\par
\vnbprop{nn-add-cancel}{\vnbstmt{\forall a, b, c \in \mathbb{N}.\; (((a + c) = (b + c)) \Rightarrow (a = b))}}
\vnbprop{nn-mul-cancel}{\vnbstmt{\forall a, b, c \in \mathbb{N}.\; ((\neg (c = 0) \wedge ((a \cdot c) = (b \cdot c))) \Rightarrow (a = b))}}

\vnbbookfile{theorem-library/nn-mod3-proof}{theorem-library/nn-mod3-proof}
\noindent{\small nn-mod3-proof.scm -- mod-3 arithmetic on NN, the mirror of nn-parity-proof.scm.}\par
\vnbprop{nn-three-is-succ3-zero}{\vnbstmt{3 = (((0 + 1) + 1) + 1)}}
\vnbprop{nn-plus-three}{\vnbstmt{\forall x \in \mathbb{N}.\; ((x + 3) = (((x + 1) + 1) + 1))}}
\vnbprop{nn-three-mul-succ}{\vnbstmt{\forall k \in \mathbb{N}.\; ((3 \cdot (k + 1)) = ((((3 \cdot k) + 1) + 1) + 1))}}
\vnbprop{nn-trichotomy-3}{\vnbstmt{\begin{array}{@{}l@{}} \forall n \in \mathbb{N}. \exists k \in \mathbb{N}. \\ \quad (n = (3 \cdot k)) \vee ((n = ((3 \cdot k) + 1)) \vee (n = (((3 \cdot k) + 1) + 1))) \end{array}}}
\vnbprop{nn-3-not-succ}{\vnbstmt{\forall x, y \in \mathbb{N}.\; \neg ((3 \cdot x) = ((3 \cdot y) + 1))}}
\vnbprop{nn-3-div-square}{\vnbstmt{\begin{array}{@{}l@{}} \forall p \in \mathbb{N}. \\ \quad \exists m \in \mathbb{N}.\; ((p \cdot p) = (3 \cdot m)) \Rightarrow \exists k \in \mathbb{N}.\; (p = (3 \cdot k)) \end{array}}}
\vnbprop{nn-3-cancel}{\vnbstmt{\forall x, y \in \mathbb{N}.\; (((3 \cdot x) = (3 \cdot y)) \Rightarrow (x = y))}}
\vnbprop{nn-lt-triple}{\vnbstmt{\forall k \in \mathbb{N}.\; (\neg (k = 0) \Rightarrow (k < (3 \cdot k)))}}

\vnbbookfile{theorem-library/nn-order-via-rr}{theorem-library/nn-order-via-rr}
\noindent{\small nn-order-via-rr.scm -- elementary NN order facts, PROVEN by taking the contradiction in RR.}\par
\vnbprop{nn-not-le-zero-pos}{\vnbstmt{\forall j \in \mathbb{N}.\; ((1 \leq j) \Rightarrow \neg (j \leq 0))}}
\vnbprop{nn-succ-le-antisym}{\vnbstmt{\forall a, b \in \mathbb{N}.\; (((a + 1) \leq b) \Rightarrow \neg (b \leq a))}}

\vnbbookfile{theorem-library/rake-intervals}{theorem-library/rake-intervals}
\noindent{\small rake-intervals.scm -- the interval / monus-predecessor supports of structure-library/order-lemmas.scm, PROVEN (batch K, 2026-09-17).}\par
\vnbprop{nn-succ-le-cancel}{\vnbstmt{\forall a, b \in \mathbb{N}.\; (((a + 1) \leq (b + 1)) \Rightarrow (a \leq b))}}
\vnbprop{nn-minus-succ-1}{\vnbstmt{\forall z \in \mathbb{N}.\; (\operatorname{nn\text{-}minus}((z + 1), 1) = z)}}
\vnbprop{one-in-interval}{\vnbstmt{\forall n \in \mathbb{N}.\; (1 \in \operatorname{interval}(1, (n + 1)))}}
\vnbprop{succ-not-one}{\vnbstmt{\forall q.; \forall z \in \operatorname{interval}(1, q).\; \neg ((z + 1) = 1)}}
\vnbprop{nn-minus-1-inj}{\vnbstmt{\begin{array}{@{}l@{}} \forall i, j. \\ \quad ((i \in \mathbb{N}) \wedge ((j \in \mathbb{N}) \wedge ((1 \leq i) \wedge ((1 \leq j) \wedge \neg (i = j))))) \Rightarrow  \\ \quad \quad \neg (\operatorname{nn\text{-}minus}(i, 1) = \operatorname{nn\text{-}minus}(j, 1)) \end{array}}}
\vnbprop{succ-in-interval}{\vnbstmt{\forall q \in \mathbb{N}, z \in \operatorname{interval}(1, q).\; ((z + 1) \in \operatorname{interval}(1, (q + 1)))}}
\vnbprop{pred-in-interval}{\vnbstmt{\begin{array}{@{}l@{}} \forall p \in \mathbb{N}, i \in \operatorname{interval}(1, (p + 1)). \\ \quad \neg (i = 1) \Rightarrow (\operatorname{nn\text{-}minus}(i, 1) \in \operatorname{interval}(1, p)) \end{array}}}

\vnbbookfile{theorem-library/zz-order}{theorem-library/zz-order}
\noindent{\small zz-order.scm -- the ORDER THEORY OF ZZ, proven.}\par
\vnbprop{zz-abs-in-nn}{\vnbstmt{\forall a \in \mathbb{Z}.\; (\lvert a \rvert \in \mathbb{N})}}
\vnbprop{zz-abs-cases}{\vnbstmt{\forall a \in \mathbb{Z}.\; ((a = \lvert a \rvert) \vee (a = -\lvert a \rvert))}}
\vnbprop{zz-trichotomy}{\vnbstmt{\forall a \in \mathbb{Z}.\; ((a < 0) \vee ((a = 0) \vee (0 < a)))}}
\vnbprop{zz-discrete}{\vnbstmt{\forall a \in \mathbb{Z}.\; ((0 < a) \Rightarrow (1 \leq a))}}
\vnbprop{zz-lt-succ-le}{\vnbstmt{\forall a, b \in \mathbb{Z}.\; ((a < b) \Rightarrow ((a + 1) \leq b))}}
\vnbprop{zz-nonneg-in-nn}{\vnbstmt{\forall a \in \mathbb{Z}.\; ((0 \leq a) \Rightarrow (a \in \mathbb{N}))}}

\vnbbookfile{theorem-library/rabasics-order}{theorem-library/rabasics-order}
\noindent{\small rabasics-order.scm -- calculus.pdf 1.1, 1.3, 1.9 (complex-analysis.pdf 1.2): the order basics.}\par
\vnbprop{rr-lub-unique}{\vnbstmt{\begin{array}{@{}l@{}} \forall s, a, b. \\ \quad \operatorname{rr\text{-}upper\text{-}bound}(s, a) \wedge \\ \quad \forall c.\; (\operatorname{rr\text{-}upper\text{-}bound}(s, c) \Rightarrow (a \leq c)) \wedge \\ \quad \operatorname{rr\text{-}upper\text{-}bound}(s, b) \wedge \\ \quad \forall c.\; (\operatorname{rr\text{-}upper\text{-}bound}(s, c) \Rightarrow (b \leq c)) \Rightarrow \\ \quad \quad a = b \end{array}}}
\vnbprop{nn-one-le-nonzero}{\vnbstmt{\forall c \in \mathbb{N}.\; (\neg (c = 0) \Rightarrow (1 \leq c))}}
\vnbprop{zz-lt-add-right}{\vnbstmt{\forall a, b, c \in \mathbb{Z}.\; ((a < b) \Rightarrow ((a + c) < (b + c)))}}
\vnbprop{zz-lt-mul-right}{\vnbstmt{\forall a, b \in \mathbb{Z}, c \in \mathbb{N}.\; (((a < b) \wedge \neg (c = 0)) \Rightarrow ((a \cdot c) < (b \cdot c)))}}

\vnbbookfile{theorem-library/zz-parity-proof}{theorem-library/zz-parity-proof}
\noindent{\small zz-parity-proof.scm -- parity on the INTEGERS, lifted from NN.}\par
\par\medskip\noindent\textbf{Definition} (\texttt{even}).\;For every $x$,  $\operatorname{even}(x) \Leftrightarrow \exists k \in \mathbb{Z}.\; (x = \cdot(2, k))$
\par\medskip\noindent\textbf{Definition} (\texttt{odd}).\;For every $x$,  $\operatorname{odd}(x) \Leftrightarrow \exists k \in \mathbb{Z}.\; (x = +(\cdot(2, k), 1))$
\vnbdef{even-def}{\vnbstmt{\forall x.\; (\operatorname{even}(x) \Leftrightarrow \exists k \in \mathbb{Z}.\; (x = (2 \cdot k)))}}
\vnbdef{odd-def}{\vnbstmt{\forall x.\; (\operatorname{odd}(x) \Leftrightarrow \exists k \in \mathbb{Z}.\; (x = ((2 \cdot k) + 1)))}}
\vnbprop{zz-parity}{\vnbstmt{\forall a \in \mathbb{Z}.\; (\operatorname{even}(a) \vee \operatorname{odd}(a))}}
\vnbprop{nn-2p-neq-1}{\vnbstmt{\forall p \in \mathbb{N}.\; \neg ((2 \cdot p) = 1)}}
\vnbprop{nn-2p-neq-neg1}{\vnbstmt{\forall p \in \mathbb{N}.\; \neg ((2 \cdot p) = -1)}}
\vnbprop{zz-2m-neq-1}{\vnbstmt{\forall m \in \mathbb{Z}.\; \neg ((2 \cdot m) = 1)}}
\vnbprop{zz-not-even-and-odd}{\vnbstmt{\forall x \in \mathbb{Z}.\; \neg (\operatorname{even}(x) \wedge \operatorname{odd}(x))}}
\vnbprop{zz-even-iff-not-odd}{\vnbstmt{\forall x \in \mathbb{Z}.\; (\operatorname{even}(x) \Leftrightarrow \neg \operatorname{odd}(x))}}
\vnbprop{zz-even-mul}{\vnbstmt{\forall x, y \in \mathbb{Z}.\; (\operatorname{even}(x) \Rightarrow \operatorname{even}((x \cdot y)))}}
\vnbprop{zz-odd-mul-odd}{\vnbstmt{\forall x, y \in \mathbb{Z}.\; ((\operatorname{odd}(x) \wedge \operatorname{odd}(y)) \Rightarrow \operatorname{odd}((x \cdot y)))}}
\vnbprop{zz-odd-sq-odd}{\vnbstmt{\forall x \in \mathbb{Z}.\; (\operatorname{odd}(x) \Rightarrow \operatorname{odd}((x \cdot x)))}}

\vnbbookfile{theorem-library/nn-pos-is-succ}{theorem-library/nn-pos-is-succ}
\noindent{\small nn-pos-is-succ.scm -- a natural at least 1 is a successor.}\par
\vnbprop{nn-pos-is-succ}{\vnbstmt{\forall n \in \mathbb{N}.\; ((1 \leq n) \Rightarrow \exists q \in \mathbb{N}.\; (n = (q + 1)))}}

\vnbbookfile{theorem-library/nn-order-proof}{theorem-library/nn-order-proof}
\noindent{\small nn-order-proof.scm -- three elementary facts about $<$= and + on NN, PROVEN.}\par
\vnbprop{nn-le-zero-is-zero}{\vnbstmt{\forall a \in \mathbb{N}.\; ((a \leq 0) \Rightarrow (a = 0))}}
\vnbprop{nn-add-le-mono}{\vnbstmt{\forall c, a, b \in \mathbb{N}.\; ((b \leq c) \Rightarrow ((a + b) \leq (a + c)))}}
\vnbprop{nn-le-add-left}{\vnbstmt{\forall a, b.\; (((a \in \mathbb{N}) \wedge (b \in \mathbb{N})) \Rightarrow (b \leq (a + b)))}}
\vnbprop{nn-le-antisym}{\vnbstmt{\forall a, b \in \mathbb{N}.\; (((a \leq b) \wedge (b \leq a)) \Rightarrow (a = b))}}
\vnbprop{nn-succ-not-le}{\vnbstmt{\forall n \in \mathbb{N}.\; \neg ((n + 1) \leq n)}}
\vnbprop{nn-le-gap}{\vnbstmt{\forall k, m \in \mathbb{N}.\; ((m \leq k) \Rightarrow \exists d \in \mathbb{N}.\; (k = (m + d)))}}

\vnbbookfile{theorem-library/rake-finsum-typing}{theorem-library/rake-finsum-typing}
\noindent{\small A file shared across books; the entries that belong to this one.}\par
\noindent{\small rake-finsum-typing.scm -- the finsum / monoid-power typing leaves of the rake (batch G), PROVEN, together with the guarded form of the one statement in the batch that is NOT provable as stated.}\par
\vnbprop{funcomp-succ-type}{\vnbstmt{\begin{array}{@{}l@{}} \forall x.; \forall q \in \mathbb{N}, f \in (\operatorname{interval}(1, (q + 1)) \to x). \\ \quad (\lambda z \in \operatorname{interval}(1, q).\; f((z + 1))) \in (\operatorname{interval}(1, q) \to x) \end{array}}}

\vnbbookfile{theorem-library/singleton-membership}{theorem-library/singleton-membership}
\noindent{\small singleton-membership.scm -- lemmas MOVED 2026-10-27 (integration of the 10-26/27 wave) out of theorem-library/field-ring-view.scm, so that files of library core may cite them (the library gate). The proofs are the originals, verbatim; the helper defines are a copy of field-ring-view.scm's.  Loaded immediately before ...}\par
\vnbprop{singleton-unfold}{\vnbstmt{\forall y.\; (\operatorname{singleton}(y) \simeq \operatorname{make\text{-}set}({[}y]))}}
\vnbprop{singleton-membership}{\vnbstmt{\forall y, x.\; ((x \in \operatorname{singleton}(y)) \Leftrightarrow ((x \in \mathbf{Set}) \wedge (x = y)))}}

\vnbbookfile{theorem-library/hom-laws}{theorem-library/hom-laws}
\noindent{\small A file shared across books; the entries that belong to this one.}\par
\noindent{\small hom-laws.scm -- the models of a structure form a CATEGORY: identity arrows, composition, and the hom-set as a set.}\par
\vnbprop{hom-semigroup-compose}{\vnbstmt{\begin{array}{@{}l@{}} \forall a, b, c, f, g. \\ \quad (\operatorname{is\text{-}hom\text{-}semigroup}(a, b, f) \wedge \operatorname{is\text{-}hom\text{-}semigroup}(b, c, g)) \Rightarrow  \\ \quad \quad \operatorname{is\text{-}hom\text{-}semigroup}(a, c, \operatorname{compose}(g, f)) \end{array}}}
\vnbprop{hom-monoid-compose}{\vnbstmt{\begin{array}{@{}l@{}} \forall a, b, c, f, g. \\ \quad (\operatorname{is\text{-}hom\text{-}monoid}(a, b, f) \wedge \operatorname{is\text{-}hom\text{-}monoid}(b, c, g)) \Rightarrow  \\ \quad \quad \operatorname{is\text{-}hom\text{-}monoid}(a, c, \operatorname{compose}(g, f)) \end{array}}}
\vnbprop{hom-comm-monoid-compose}{\vnbstmt{\begin{array}{@{}l@{}} \forall a, b, c, f, g. \\ \quad (\operatorname{is\text{-}hom\text{-}comm\text{-}monoid}(a, b, f) \wedge \operatorname{is\text{-}hom\text{-}comm\text{-}monoid}(b, c, g)) \Rightarrow  \\ \quad \quad \operatorname{is\text{-}hom\text{-}comm\text{-}monoid}(a, c, \operatorname{compose}(g, f)) \end{array}}}
\vnbprop{hom-group-compose}{\vnbstmt{\begin{array}{@{}l@{}} \forall a, b, c, f, g. \\ \quad (\operatorname{is\text{-}hom\text{-}group}(a, b, f) \wedge \operatorname{is\text{-}hom\text{-}group}(b, c, g)) \Rightarrow  \\ \quad \quad \operatorname{is\text{-}hom\text{-}group}(a, c, \operatorname{compose}(g, f)) \end{array}}}
\vnbprop{hom-abelian-group-compose}{\vnbstmt{\begin{array}{@{}l@{}} \forall a, b, c, f, g. \\ \quad (\operatorname{is\text{-}hom\text{-}abelian\text{-}group}(a, b, f) \wedge \operatorname{is\text{-}hom\text{-}abelian\text{-}group}(b, c, g)) \Rightarrow  \\ \quad \quad \operatorname{is\text{-}hom\text{-}abelian\text{-}group}(a, c, \operatorname{compose}(g, f)) \end{array}}}
\vnbprop{hom-ring-compose}{\vnbstmt{\begin{array}{@{}l@{}} \forall a, b, c, f, g. \\ \quad (\operatorname{is\text{-}hom\text{-}ring}(a, b, f) \wedge \operatorname{is\text{-}hom\text{-}ring}(b, c, g)) \Rightarrow  \\ \quad \quad \operatorname{is\text{-}hom\text{-}ring}(a, c, \operatorname{compose}(g, f)) \end{array}}}
\vnbprop{hom-commutative-ring-compose}{\vnbstmt{\begin{array}{@{}l@{}} \forall a, b, c, f, g. \\ \quad (\operatorname{is\text{-}hom\text{-}commutative\text{-}ring}(a, b, f) \wedge \operatorname{is\text{-}hom\text{-}commutative\text{-}ring}(b, c, g)) \Rightarrow  \\ \quad \quad \operatorname{is\text{-}hom\text{-}commutative\text{-}ring}(a, c, \operatorname{compose}(g, f)) \end{array}}}
\vnbprop{hom-integral-domain-compose}{\vnbstmt{\begin{array}{@{}l@{}} \forall a, b, c, f, g. \\ \quad (\operatorname{is\text{-}hom\text{-}integral\text{-}domain}(a, b, f) \wedge \operatorname{is\text{-}hom\text{-}integral\text{-}domain}(b, c, g)) \Rightarrow  \\ \quad \quad \operatorname{is\text{-}hom\text{-}integral\text{-}domain}(a, c, \operatorname{compose}(g, f)) \end{array}}}
\vnbprop{hom-euclidean-ring-compose}{\vnbstmt{\begin{array}{@{}l@{}} \forall a, b, c, f, g. \\ \quad (\operatorname{is\text{-}hom\text{-}euclidean\text{-}ring}(a, b, f) \wedge \operatorname{is\text{-}hom\text{-}euclidean\text{-}ring}(b, c, g)) \Rightarrow  \\ \quad \quad \operatorname{is\text{-}hom\text{-}euclidean\text{-}ring}(a, c, \operatorname{compose}(g, f)) \end{array}}}
\vnbprop{hom-pid-compose}{\vnbstmt{\begin{array}{@{}l@{}} \forall a, b, c, f, g. \\ \quad (\operatorname{is\text{-}hom\text{-}pid}(a, b, f) \wedge \operatorname{is\text{-}hom\text{-}pid}(b, c, g)) \Rightarrow  \\ \quad \quad \operatorname{is\text{-}hom\text{-}pid}(a, c, \operatorname{compose}(g, f)) \end{array}}}
\vnbprop{hom-field-ring-compose}{\vnbstmt{\begin{array}{@{}l@{}} \forall a, b, c, f, g. \\ \quad (\operatorname{is\text{-}hom\text{-}field\text{-}ring}(a, b, f) \wedge \operatorname{is\text{-}hom\text{-}field\text{-}ring}(b, c, g)) \Rightarrow  \\ \quad \quad \operatorname{is\text{-}hom\text{-}field\text{-}ring}(a, c, \operatorname{compose}(g, f)) \end{array}}}
\vnbprop{hom-normed-ag-compose}{\vnbstmt{\begin{array}{@{}l@{}} \forall a, b, c, f, g. \\ \quad (\operatorname{is\text{-}hom\text{-}normed\text{-}ag}(a, b, f) \wedge \operatorname{is\text{-}hom\text{-}normed\text{-}ag}(b, c, g)) \Rightarrow  \\ \quad \quad \operatorname{is\text{-}hom\text{-}normed\text{-}ag}(a, c, \operatorname{compose}(g, f)) \end{array}}}
\vnbprop{hom-normed-field-compose}{\vnbstmt{\begin{array}{@{}l@{}} \forall a, b, c, f, g. \\ \quad (\operatorname{is\text{-}hom\text{-}normed\text{-}field}(a, b, f) \wedge \operatorname{is\text{-}hom\text{-}normed\text{-}field}(b, c, g)) \Rightarrow  \\ \quad \quad \operatorname{is\text{-}hom\text{-}normed\text{-}field}(a, c, \operatorname{compose}(g, f)) \end{array}}}
\vnbprop{hom-module-compose}{\vnbstmt{\begin{array}{@{}l@{}} \forall a, b, c, f, g. \\ \quad (\operatorname{is\text{-}hom\text{-}module}(a, b, f) \wedge \operatorname{is\text{-}hom\text{-}module}(b, c, g)) \Rightarrow  \\ \quad \quad \operatorname{is\text{-}hom\text{-}module}(a, c, \operatorname{compose}(g, f)) \end{array}}}
\vnbprop{hom-vector-space-compose}{\vnbstmt{\begin{array}{@{}l@{}} \forall a, b, c, f, g. \\ \quad (\operatorname{is\text{-}hom\text{-}vector\text{-}space}(a, b, f) \wedge \operatorname{is\text{-}hom\text{-}vector\text{-}space}(b, c, g)) \Rightarrow  \\ \quad \quad \operatorname{is\text{-}hom\text{-}vector\text{-}space}(a, c, \operatorname{compose}(g, f)) \end{array}}}
\vnbprop{hom-normed-vector-space-compose}{\vnbstmt{\begin{array}{@{}l@{}} \forall a, b, c, f, g. \\ \quad (\operatorname{is\text{-}hom\text{-}normed\text{-}vector\text{-}space}(a, b, f) \wedge \operatorname{is\text{-}hom\text{-}normed\text{-}vector\text{-}space}(b, c, g)) \Rightarrow  \\ \quad \quad \operatorname{is\text{-}hom\text{-}normed\text{-}vector\text{-}space}(a, c, \operatorname{compose}(g, f)) \end{array}}}
\vnbprop{hom-complex-inner-product-space-compose}{\vnbstmt{\begin{array}{@{}l@{}} \forall a, b, c, f, g. \\ \quad (\operatorname{is\text{-}hom\text{-}complex\text{-}inner\text{-}product\text{-}space}(a, b, f) \wedge \operatorname{is\text{-}hom\text{-}complex\text{-}inner\text{-}product\text{-}space}(b, c, g)) \Rightarrow  \\ \quad \quad \operatorname{is\text{-}hom\text{-}complex\text{-}inner\text{-}product\text{-}space}(a, c, \operatorname{compose}(g, f)) \end{array}}}
\vnbprop{hom-metric-space-compose}{\vnbstmt{\begin{array}{@{}l@{}} \forall a, b, c, f, g. \\ \quad (\operatorname{is\text{-}hom\text{-}metric\text{-}space}(a, b, f) \wedge \operatorname{is\text{-}hom\text{-}metric\text{-}space}(b, c, g)) \Rightarrow  \\ \quad \quad \operatorname{is\text{-}hom\text{-}metric\text{-}space}(a, c, \operatorname{compose}(g, f)) \end{array}}}
\vnbprop{hom-pseudometric-space-compose}{\vnbstmt{\begin{array}{@{}l@{}} \forall a, b, c, f, g. \\ \quad (\operatorname{is\text{-}hom\text{-}pseudometric\text{-}space}(a, b, f) \wedge \operatorname{is\text{-}hom\text{-}pseudometric\text{-}space}(b, c, g)) \Rightarrow  \\ \quad \quad \operatorname{is\text{-}hom\text{-}pseudometric\text{-}space}(a, c, \operatorname{compose}(g, f)) \end{array}}}
\vnbprop{hom-c-metric-space-compose}{\vnbstmt{\begin{array}{@{}l@{}} \forall a, b, c, f, g. \\ \quad (\operatorname{is\text{-}hom\text{-}c\text{-}metric\text{-}space}(a, b, f) \wedge \operatorname{is\text{-}hom\text{-}c\text{-}metric\text{-}space}(b, c, g)) \Rightarrow  \\ \quad \quad \operatorname{is\text{-}hom\text{-}c\text{-}metric\text{-}space}(a, c, \operatorname{compose}(g, f)) \end{array}}}
\vnbprop{hom-top-space-compose}{\vnbstmt{\begin{array}{@{}l@{}} \forall s, t, w, f, g. \\ \quad (\operatorname{is\text{-}hom\text{-}top\text{-}space}(s, t, f) \wedge \operatorname{is\text{-}hom\text{-}top\text{-}space}(t, w, g)) \Rightarrow  \\ \quad \quad \operatorname{is\text{-}hom\text{-}top\text{-}space}(s, w, \operatorname{compose}(g, f)) \end{array}}}
\vnbprop{hom-metrizable-top-space-compose}{\vnbstmt{\begin{array}{@{}l@{}} \forall s, t, w, f, g. \\ \quad (\operatorname{is\text{-}hom\text{-}metrizable\text{-}top\text{-}space}(s, t, f) \wedge \operatorname{is\text{-}hom\text{-}metrizable\text{-}top\text{-}space}(t, w, g)) \Rightarrow  \\ \quad \quad \operatorname{is\text{-}hom\text{-}metrizable\text{-}top\text{-}space}(s, w, \operatorname{compose}(g, f)) \end{array}}}
\vnbprop{hom-ringoid-compose}{\vnbstmt{\begin{array}{@{}l@{}} \forall a, b, c, f, g. \\ \quad (\operatorname{is\text{-}hom\text{-}ringoid}(a, b, f) \wedge \operatorname{is\text{-}hom\text{-}ringoid}(b, c, g)) \Rightarrow  \\ \quad \quad \operatorname{is\text{-}hom\text{-}ringoid}(a, c, \operatorname{compose}(g, f)) \end{array}}}

\vnbbookfile{theorem-library/hom-kinds}{theorem-library/hom-kinds}
\noindent{\small A file shared across books; the entries that belong to this one.}\par
\noindent{\small hom-kinds.scm -- the category laws for the arrows the generic driver of hom-laws.scm does not reach: the FAMILY slot kinds, and the RELATION kind (SETOID, since batch 40; two carriers in batch 39).}\par
\vnbprop{hom-measurable-space-compose}{\vnbstmt{\begin{array}{@{}l@{}} \forall a, b, c, f, g. \\ \quad (\operatorname{is\text{-}hom\text{-}measurable\text{-}space}(a, b, f) \wedge \operatorname{is\text{-}hom\text{-}measurable\text{-}space}(b, c, g)) \Rightarrow  \\ \quad \quad \operatorname{is\text{-}hom\text{-}measurable\text{-}space}(a, c, \operatorname{compose}(g, f)) \end{array}}}
\vnbprop{hom-measure-space-compose}{\vnbstmt{\begin{array}{@{}l@{}} \forall a, b, c, f, g. \\ \quad (\operatorname{is\text{-}hom\text{-}measure\text{-}space}(a, b, f) \wedge \operatorname{is\text{-}hom\text{-}measure\text{-}space}(b, c, g)) \Rightarrow  \\ \quad \quad \operatorname{is\text{-}hom\text{-}measure\text{-}space}(a, c, \operatorname{compose}(g, f)) \end{array}}}
\vnbprop{hom-setoid-compose}{\vnbstmt{\begin{array}{@{}l@{}} \forall a, b, c, f, g. \\ \quad (\operatorname{is\text{-}hom\text{-}setoid}(a, b, f) \wedge \operatorname{is\text{-}hom\text{-}setoid}(b, c, g)) \Rightarrow  \\ \quad \quad \operatorname{is\text{-}hom\text{-}setoid}(a, c, \operatorname{compose}(g, f)) \end{array}}}

\vnbbookfile{theorem-library/hom-functors}{theorem-library/hom-functors}
\noindent{\small A file shared across books; the entries that belong to this one.}\par
\noindent{\small hom-functors.scm -- the Hom functors: Hom(a, -) covariant, Hom(-, b) contravariant, as maps on hom-sets.}\par
\vnbprop{post-id}{\vnbstmt{\begin{array}{@{}l@{}} \forall a, b, s, f. \\ \quad ((a \in \mathbf{Set}) \wedge ((b \in \mathbf{Set}) \wedge ((f \in s) \wedge (f \in (a \to b))))) \Rightarrow  \\ \quad \quad (\lambda \mathit{hmpf} \in s.\; \operatorname{compose}(\operatorname{id\text{-}fun}(b), \mathit{hmpf}))(f) = f \end{array}}}
\vnbprop{pre-id}{\vnbstmt{\begin{array}{@{}l@{}} \forall a, b, s, f. \\ \quad ((a \in \mathbf{Set}) \wedge ((f \in s) \wedge (f \in (a \to b)))) \Rightarrow  \\ \quad \quad (\lambda \mathit{hmpf} \in s.\; \operatorname{compose}(\mathit{hmpf}, \operatorname{id\text{-}fun}(a)))(f) = f \end{array}}}
\vnbprop{post-compose}{\vnbstmt{\begin{array}{@{}l@{}} \forall a, b, c, u, s, d, f, h, k. \\ \quad a \in \mathbf{Set} \wedge \\ \quad b \in \mathbf{Set} \wedge \\ \quad f \in (a \to b) \wedge \\ \quad h \in (b \to c) \wedge \\ \quad k \in (c \to u) \wedge \\ \quad f \in s \wedge \\ \quad \operatorname{compose}(h, f) \in d \Rightarrow \\ \quad \quad (\lambda \mathit{hmpf} \in s.\; \operatorname{compose}(\operatorname{compose}(k, h), \mathit{hmpf}))(f) \\ \quad \quad =\; (\lambda \mathit{hmpf} \in d.\; \operatorname{compose}(k, \mathit{hmpf}))((\lambda \mathit{hmpf} \in s.\; \operatorname{compose}(h, \mathit{hmpf}))(f)) \end{array}}}
\vnbprop{pre-compose}{\vnbstmt{\begin{array}{@{}l@{}} \forall a, b, c, u, s, d, f, h, k. \\ \quad a \in \mathbf{Set} \wedge \\ \quad b \in \mathbf{Set} \wedge \\ \quad k \in (a \to b) \wedge \\ \quad h \in (b \to c) \wedge \\ \quad f \in (c \to u) \wedge \\ \quad f \in s \wedge \\ \quad \operatorname{compose}(f, h) \in d \Rightarrow \\ \quad \quad (\lambda \mathit{hmpf} \in s.\; \operatorname{compose}(\mathit{hmpf}, \operatorname{compose}(h, k)))(f) \\ \quad \quad =\; (\lambda \mathit{hmpf} \in d.\; \operatorname{compose}(\mathit{hmpf}, k))((\lambda \mathit{hmpf} \in s.\; \operatorname{compose}(\mathit{hmpf}, h))(f)) \end{array}}}

\vnbbookfile{theorem-library/hom-views}{theorem-library/hom-views}
\noindent{\small A file shared across books; the entries that belong to this one.}\par
\noindent{\small hom-views.scm -- a view's action on arrows.}\par
\vnbprop{ringoid-as-ring-functorial}{\vnbstmt{\begin{array}{@{}l@{}} \forall a, b, f. \\ \quad \operatorname{is\text{-}hom\text{-}ringoid}(a, b, f) \Rightarrow  \\ \quad \quad \operatorname{is\text{-}hom\text{-}ring}(\operatorname{ringoid\text{-}as\text{-}ring}(a), \operatorname{ringoid\text{-}as\text{-}ring}(b), f) \end{array}}}

\vnbbookfile{theorem-library/makeset-basics}{theorem-library/makeset-basics}
\noindent{\small makeset-basics.scm -- reading membership out of a literal brace set.}\par
\vnbprop{ms-eq-symm}{\vnbstmt{\forall x, y.\; ((x = y) \Leftrightarrow (y = x))}}
\vnbprop{ms-index-2}{\vnbstmt{\forall i \in \mathbb{N}.\; (((1 \leq i) \wedge (i \leq 2)) \Rightarrow ((i = 1) \vee (i = 2)))}}
\vnbprop{ms-nth-2}{\vnbstmt{\begin{array}{@{}l@{}} \forall a, b, i. \\ \quad ((a \in \mathbf{Set}) \wedge ((b \in \mathbf{Set}) \wedge ((i \in \mathbb{N}) \wedge ((1 \leq i) \wedge (i \leq 2))))) \Rightarrow  \\ \quad \quad (\operatorname{nth}(i, {[}a, b]) = a) \vee (\operatorname{nth}(i, {[}a, b]) = b) \end{array}}}
\vnbprop{makeset2-membership}{\vnbstmt{\begin{array}{@{}l@{}} \forall x, a, b. \\ \quad ((a \in \mathbf{Set}) \wedge (b \in \mathbf{Set})) \Rightarrow ((x \in \operatorname{make\text{-}set}({[}a, b])) \Leftrightarrow ((x \in \mathbf{Set}) \wedge ((x = a) \vee (x = b)))) \end{array}}}
\vnbprop{makeset2-split}{\vnbstmt{\forall a, b \in \mathbf{Set}.\; (\operatorname{make\text{-}set}({[}a, b]) = (\{a\} \cup \{b\}))}}
\vnbprop{card-pair}{\vnbstmt{\forall a, b \in \mathbf{Set}.\; (\neg (a = b) \Rightarrow (|\operatorname{make\text{-}set}({[}a, b])| = 2))}}

\vnbbookfile{theorem-library/empty-cardinality}{theorem-library/empty-cardinality}
\noindent{\small empty-cardinality.scm -- the two BACKWARD rungs for "this set is empty, so its cardinal is 0".}\par
\vnbprop{card-empty-le}{\vnbstmt{\forall x.\; ((x = \emptyset) \Rightarrow (|x| \leq 0))}}
\vnbprop{makeset-of-empty-tuple}{\vnbstmt{\forall l.\; ((l = {[}]) \Rightarrow (\operatorname{make\text{-}set}(l) = \emptyset))}}

\vnbbookfile{theorem-library/tuples-induction}{theorem-library/tuples-induction}
\noindent{\small tuples-induction.scm -- list induction, DERIVED.}\par
\vnbprop{tuples-induction}{\vnbstmt{\begin{array}{@{}l@{}} \forall a, c.; \forall l \in \operatorname{tuples}(a). \\ \quad (({[}] \in c) \wedge \forall x, m.\; (((x \in a) \wedge ((m \in \operatorname{tuples}(a)) \wedge (m \in c))) \Rightarrow (\operatorname{cons}(x, m) \in c))) \Rightarrow  \\ \quad \quad l \in c \end{array}}}

\vnbbookfile{theorem-library/tuple-extensionality}{theorem-library/tuple-extensionality}
\noindent{\small tuple-extensionality.scm -- two tuples with the same length and the same entry at every index are equal; and, on top of it, matrix-entry-extensionality.}\par
\vnbprop{nth2-pair}{\vnbstmt{\forall a, b.\; (((a \in \mathbf{Set}) \wedge (b \in \mathbf{Set})) \Rightarrow (\operatorname{nth}(2, {[}a, b]) \simeq b))}}
\vnbprop{entry-unfold}{\vnbstmt{\forall m, i, j.\; (\operatorname{entry}(m, i, j) \simeq \operatorname{nth}(j, \operatorname{nth}(i, m)))}}
\vnbprop{tuple-extensionality}{\vnbstmt{\begin{array}{@{}l@{}} \forall a, l, m. \\ \quad l \in \operatorname{tuples}(a) \wedge \\ \quad m \in \operatorname{tuples}(a) \wedge \\ \quad \operatorname{length}(l) = \operatorname{length}(m) \wedge \\ \quad \forall k \in \mathbb{N}.\; (((1 \leq k) \wedge (k \leq \operatorname{length}(l))) \Rightarrow (\operatorname{nth}(k, l) = \operatorname{nth}(k, m))) \Rightarrow \\ \quad \quad l = m \end{array}}}
\vnbprop{matrix-entry-extensionality}{\vnbstmt{\begin{array}{@{}l@{}} \forall m, n, x, p, q. \\ \quad p \in \operatorname{mat}(m, n, x) \wedge \\ \quad q \in \operatorname{mat}(m, n, x) \wedge \\ \quad \forall i \in \operatorname{interval}(1, m), j \in \operatorname{interval}(1, n). \\ \quad \quad \operatorname{entry}(p, i, j) = \operatorname{entry}(q, i, j) \Rightarrow \\ \quad \quad p = q \end{array}}}

\vnbbookfile{theorem-library/la-z5-card}{theorem-library/la-z5-card}
\noindent{\small la-z5-card.scm -- counting: the cardinality of a cartesian product of finite sets (LA-Z5, 2026-10-25).  Prefix lz5c-.}\par
\vnbprop{card-cartesian}{\vnbstmt{\begin{array}{@{}l@{}} \forall a, b \in \mathbf{Set}. \\ \quad ((|a| \in \mathbb{N}) \wedge (|b| \in \mathbb{N})) \Rightarrow  \\ \quad \quad |(a \times b)| = (|a| \cdot |b|) \end{array}}}

\vnbbookfile{theorem-library/tuple-tabulation}{theorem-library/tuple-tabulation}
\noindent{\small tuple-tabulation.scm -- TABULATION: the missing generation principle for tuples and matrices, and the honest replacement for `matof-exists'.}\par
\vnbprop{tab-succ-add}{\vnbstmt{\forall a, b \in \mathbb{N}.\; (((a + 1) + b) = ((a + b) + 1))}}
\vnbprop{tuple-tabulation-offset}{\vnbstmt{\begin{array}{@{}l@{}} \forall n \in \mathbb{N}, a, f.; \forall k \in \mathbb{N}. \\ \quad \forall i \in \mathbb{N}.\; (((1 \leq i) \wedge (i \leq n)) \Rightarrow (f((k + i)) \in a)) \Rightarrow  \\ \quad \quad \exists l \in \operatorname{tuples}(a). \\ \quad \quad \quad ((\operatorname{length}(l) = n) \wedge \forall i \in \mathbb{N}.\; (((1 \leq i) \wedge (i \leq n)) \Rightarrow (\operatorname{nth}(i, l) = f((k + i))))) \end{array}}}
\vnbprop{row-tabulation-offset}{\vnbstmt{\begin{array}{@{}l@{}} \forall n \in \mathbb{N}, a, g, c.; \forall k \in \mathbb{N}. \\ \quad \forall j \in \mathbb{N}.\; (((1 \leq j) \wedge (j \leq n)) \Rightarrow (g(c, (k + j)) \in a)) \Rightarrow  \\ \quad \quad \exists l \in \operatorname{tuples}(a). \\ \quad \quad \quad ((\operatorname{length}(l) = n) \wedge \forall j \in \mathbb{N}.\; (((1 \leq j) \wedge (j \leq n)) \Rightarrow (\operatorname{nth}(j, l) = g(c, (k + j))))) \end{array}}}
\vnbprop{tuple-tabulation}{\vnbstmt{\begin{array}{@{}l@{}} \forall n \in \mathbb{N}, a, f. \\ \quad \forall i \in \mathbb{N}.\; (((1 \leq i) \wedge (i \leq n)) \Rightarrow (f(i) \in a)) \Rightarrow  \\ \quad \quad \exists l \in \operatorname{tuples}(a). \\ \quad \quad \quad (\operatorname{length}(l) = n) \wedge \forall i \in \mathbb{N}.\; (((1 \leq i) \wedge (i \leq n)) \Rightarrow (\operatorname{nth}(i, l) = f(i))) \end{array}}}
\vnbprop{row-tabulation}{\vnbstmt{\begin{array}{@{}l@{}} \forall n \in \mathbb{N}, a, g, c. \\ \quad \forall j \in \mathbb{N}.\; (((1 \leq j) \wedge (j \leq n)) \Rightarrow (g(c, j) \in a)) \Rightarrow  \\ \quad \quad \exists l \in \operatorname{tuples}(a). \\ \quad \quad \quad ((\operatorname{length}(l) = n) \wedge \forall j \in \mathbb{N}.\; (((1 \leq j) \wedge (j \leq n)) \Rightarrow (\operatorname{nth}(j, l) = g(c, j)))) \end{array}}}
\vnbprop{matrix-tabulation}{\vnbstmt{\begin{array}{@{}l@{}} \forall m \in \mathbb{N}, a, g, n.; \forall k \in \mathbb{N}. \\ \quad n \in \mathbb{N} \wedge \\ \quad \forall i, j \in \mathbb{N}. \\ \quad \quad (((1 \leq i) \wedge (i \leq m)) \wedge ((1 \leq j) \wedge (j \leq n))) \Rightarrow (g((k + i), j) \in a) \Rightarrow \\ \quad \quad \exists p \in \operatorname{tuples}(\operatorname{tuples}(a)). \\ \quad \quad \quad \operatorname{length}(p) = m \wedge \\ \quad \quad \quad \forall i \in \mathbb{N}.\; (((1 \leq i) \wedge (i \leq m)) \Rightarrow (\operatorname{length}(\operatorname{nth}(i, p)) = n)) \wedge \\ \quad \quad \quad \forall i, j \in \mathbb{N}. \\ \quad \quad \quad \quad (((1 \leq i) \wedge (i \leq m)) \wedge ((1 \leq j) \wedge (j \leq n))) \Rightarrow  \\ \quad \quad \quad \quad \quad \operatorname{nth}(j, \operatorname{nth}(i, p)) = g((k + i), j) \end{array}}}
\vnbprop{mat-tabulation-exists}{\vnbstmt{\begin{array}{@{}l@{}} \forall m, n, x, g. \\ \quad ((m \in \mathbb{N}) \wedge ((n \in \mathbb{N}) \wedge ((1 \leq m) \wedge \forall i \in \operatorname{interval}(1, m), j \in \operatorname{interval}(1, n).\; (g(i, j) \in x)))) \Rightarrow  \\ \quad \quad \exists p \in \operatorname{mat}(m, n, x). \forall i \in \operatorname{interval}(1, m), j \in \operatorname{interval}(1, n). \\ \quad \quad \quad \operatorname{entry}(p, i, j) = g(i, j) \end{array}}}
\vnbprop{matof-exists-image}{\vnbstmt{\begin{array}{@{}l@{}} \forall m, n, g. \\ \quad m \in \mathbb{N} \wedge \\ \quad n \in \mathbb{N} \wedge \\ \quad 1 \leq m \wedge \\ \quad \forall i \in \operatorname{interval}(1, m), j \in \operatorname{interval}(1, n).\; (g(i, j) \in \mathbf{Set}) \Rightarrow \\ \quad \quad \exists p \in \operatorname{mat}(m, n, \operatorname{image}(g, (\operatorname{interval}(1, m) \times \operatorname{interval}(1, n)))). \forall i \in \operatorname{interval}(1, m), j \in \operatorname{interval}(1, n). \\ \quad \quad \quad \operatorname{entry}(p, i, j) = g(i, j) \end{array}}}
\vnbprop{entry-of-matof-guarded}{\vnbstmt{\begin{array}{@{}l@{}} \forall m, n, g.; \forall u \in \operatorname{interval}(1, m), v \in \operatorname{interval}(1, n). \\ \quad m \in \mathbb{N} \wedge \\ \quad n \in \mathbb{N} \wedge \\ \quad 1 \leq m \wedge \\ \quad \forall i \in \operatorname{interval}(1, m), j \in \operatorname{interval}(1, n).\; (g(i, j) \in \mathbf{Set}) \Rightarrow \\ \quad \quad \operatorname{entry}(\operatorname{matof}(m, n, g), u, v) = g(u, v) \end{array}}}
\vnbprop{matof-exists}{\vnbstmt{\begin{array}{@{}l@{}} \forall m, n, g. \\ \quad ((m \in \mathbb{N}) \wedge ((n \in \mathbb{N}) \wedge \forall a \in \operatorname{interval}(1, m), b \in \operatorname{interval}(1, n).\; (g(a, b) \in \mathbf{Set}))) \Rightarrow  \\ \quad \quad \exists p \in \operatorname{mat}(m, n, \operatorname{image}(g, (\operatorname{interval}(1, m) \times \operatorname{interval}(1, n)))). \forall i \in \operatorname{interval}(1, m), j \in \operatorname{interval}(1, n). \\ \quad \quad \quad \operatorname{entry}(p, i, j) = g(i, j) \end{array}}}
\vnbprop{entry-of-matof}{\vnbstmt{\begin{array}{@{}l@{}} \forall m, n, g, i, j. \\ \quad m \in \mathbb{N} \wedge \\ \quad n \in \mathbb{N} \wedge \\ \quad \forall a \in \operatorname{interval}(1, m), b \in \operatorname{interval}(1, n).\; (g(a, b) \in \mathbf{Set}) \wedge \\ \quad i \in \operatorname{interval}(1, m) \wedge \\ \quad j \in \operatorname{interval}(1, n) \Rightarrow \\ \quad \quad \operatorname{entry}(\operatorname{matof}(m, n, g), i, j) = g(i, j) \end{array}}}

\vnbbookfile{theorem-library/matof-in-mat}{theorem-library/matof-in-mat}
\noindent{\small matof-in-mat.scm -- the tabulated matrix MATOF(m,n,g) is an m-by-n matrix over X when every value g(i,j) on the index box [1,m] x [1,n] lies in X.}\par
\vnbprop{tuples-mono}{\vnbstmt{\forall a, b.; \forall l \in \operatorname{tuples}(a).\; (\forall w \in a.\; (w \in b) \Rightarrow (l \in \operatorname{tuples}(b)))}}
\vnbprop{matrix-mono}{\vnbstmt{\forall a, b.; \forall m \in \operatorname{matrix}(a).\; (\forall w \in a.\; (w \in b) \Rightarrow (m \in \operatorname{matrix}(b)))}}
\vnbprop{mat-mono}{\vnbstmt{\begin{array}{@{}l@{}} \forall m, n, a, b.; \forall q \in \operatorname{mat}(m, n, a). \\ \quad \forall w \in a.\; (w \in b) \Rightarrow (q \in \operatorname{mat}(m, n, b)) \end{array}}}
\vnbprop{matof-in-mat}{\vnbstmt{\begin{array}{@{}l@{}} \forall m, n, x, g. \\ \quad ((m \in \mathbb{N}) \wedge ((n \in \mathbb{N}) \wedge \forall i \in \operatorname{interval}(1, m), j \in \operatorname{interval}(1, n).\; (g(i, j) \in x))) \Rightarrow  \\ \quad \quad \operatorname{matof}(m, n, g) \in \operatorname{mat}(m, n, x) \end{array}}}

\vnbbookfile{theorem-library/makeset-card-bound}{theorem-library/makeset-card-bound}
\noindent{\small makeset-card-bound.scm -- a tuple of length n has at most n distinct entries.}\par
\vnbprop{union-comm}{\vnbstmt{\forall a, b.\; ((a \cup b) = (b \cup a))}}
\vnbprop{union-singleton-absorb}{\vnbstmt{\forall s.; \forall y \in s.\; ((\{y\} \cup s) = s)}}
\vnbprop{card-union-singleton-bound}{\vnbstmt{\begin{array}{@{}l@{}} \forall s, y \in \mathbf{Set}. \\ \quad (|s| \in \mathbb{N}) \Rightarrow  \\ \quad \quad (|(\{y\} \cup s)| \in \mathbb{N}) \wedge (|(\{y\} \cup s)| \leq (|s| + 1)) \end{array}}}
\vnbprop{makeset-card-bound-strong}{\vnbstmt{\begin{array}{@{}l@{}} \forall n \in \mathbb{N}, a.; \forall l \in \operatorname{tuples}(a). \\ \quad (\operatorname{length}(l) = n) \Rightarrow  \\ \quad \quad (\operatorname{make\text{-}set}(l) \in \mathbf{Set}) \wedge ((|\operatorname{make\text{-}set}(l)| \in \mathbb{N}) \wedge (|\operatorname{make\text{-}set}(l)| \leq n)) \end{array}}}
\vnbprop{makeset-card-bound}{\vnbstmt{\begin{array}{@{}l@{}} \forall n \in \mathbb{N}, a.; \forall l \in \operatorname{tuples}(a). \\ \quad (\operatorname{length}(l) = n) \Rightarrow (|\operatorname{make\text{-}set}(l)| \leq n) \end{array}}}

\vnbbookfile{theorem-library/card-subset-nn}{theorem-library/card-subset-nn}
\noindent{\small card-subset-nn.scm -- a subset of a finite set is finite.}\par
\vnbprop{card-subset-nn}{\vnbstmt{\begin{array}{@{}l@{}} \forall x, s \in \mathbf{Set}. \\ \quad ((|x| \in \mathbb{N}) \wedge \forall z \in s.\; (z \in x)) \Rightarrow (|s| \in \mathbb{N}) \end{array}}}

\vnbbookfile{theorem-library/card-image-finite}{theorem-library/card-image-finite}
\noindent{\small card-image-finite.scm -- the image of a finite set is finite.}\par
\vnbprop{card-image-finite}{\vnbstmt{\forall i.; \forall x \in \mathbf{Set}.\; ((|x| \in \mathbb{N}) \Rightarrow (|\operatorname{image}(i, x)| \in \mathbb{N}))}}

\vnbbookfile{theorem-library/interval-card-in-nn}{theorem-library/interval-card-in-nn}
\noindent{\small interval-card-in-nn.scm -- an interval with a natural upper bound is finite.}\par
\vnbprop{interval-card-in-nn}{\vnbstmt{\forall a.; \forall b \in \mathbb{N}.\; (|\operatorname{interval}(a, b)| \in \mathbb{N})}}

\vnbbookfile{theorem-library/rake-finsum-laws2}{theorem-library/rake-finsum-laws2}
\noindent{\small A file shared across books; the entries that belong to this one.}\par
\noindent{\small theorem-library/rake-finsum-laws2.scm -- rake batch J, part 2: the three DISTRIBUTION laws of FINSUM (a ring element on either side, a module scalar) and the interval back-peel.  Part 1 is theorem-library/rake-finsum-laws.scm, which must load FIRST: this file cites its `enum-fam-value' and}\par
\vnbprop{interval-succ-insert}{\vnbstmt{\begin{array}{@{}l@{}} \forall n \in \mathbb{N}. \\ \quad \operatorname{interval}(1, (n + 1)) = (\operatorname{interval}(1, n) \cup \{(n + 1)\}) \end{array}}}

\vnbbookfile{theorem-library/rake-finsum-core}{theorem-library/rake-finsum-core}
\noindent{\small A file shared across books; the entries that belong to this one.}\par
\noindent{\small theorem-library/rake-finsum-core.scm -- the FINSUM CORE, proven. Rake batch M (2026-09-17).  Nine of the batch's ten leaves; the tenth (finsum-embed) is BLOCKED and its obstacle is stated at the end of this header.}\par
\vnbprop{cra-is-rag}{\vnbstmt{\forall g.\; (\operatorname{commutative\text{-}ring\text{-}additive\text{-}ag}(g) \simeq \operatorname{ring\text{-}additive\text{-}ag}(g))}}
\vnbprop{ord-segment-insert}{\vnbstmt{\forall n \in \mathbb{N}.\; (\mathrm{OS}((n + 1)) = (\mathrm{OS}(n) \cup \{n\}))}}

\vnbbookfile{theorem-library/nn-finite-subset-bounded}{theorem-library/nn-finite-subset-bounded}
\noindent{\small nn-finite-subset-bounded.scm -- a finite subset of NN is bounded.}\par
\vnbprop{nn-finite-subset-bounded}{\vnbstmt{\begin{array}{@{}l@{}} \forall t. \\ \quad ((t \subseteq \mathbb{N}) \wedge (|t| \in \mathbb{N})) \Rightarrow  \\ \quad \quad \exists n \in \mathbb{N}.\; \forall y \in \mathbb{N}.\; ((n \leq y) \Rightarrow \neg (y \in t)) \end{array}}}

\vnbbookfile{theorem-library/card-inequalities}{theorem-library/card-inequalities}
\noindent{\small card-inequalities.scm -- the two missing CARD inequalities. --------------------------------------------------------------------------- NOTE, 2026-09-20 (batch 9-B).  This file was written while cardinality was AXIOMATISED under the name CARD and the defined constant was its companion}\par
\vnbprop{complement-in-subset}{\vnbstmt{\forall a, b.\; (\operatorname{complement\text{-}in}(b, a) \subseteq b)}}
\vnbprop{union-complement-in}{\vnbstmt{\forall a, b.\; ((a \subseteq b) \Rightarrow (b = (a \cup \operatorname{complement\text{-}in}(b, a))))}}
\vnbprop{union-as-disjoint}{\vnbstmt{\forall a, b.\; ((a \cup b) = (a \cup \operatorname{complement\text{-}in}(b, a)))}}
\vnbprop{intersection-complement-in-empty}{\vnbstmt{\forall a, b.\; ((a \cap \operatorname{complement\text{-}in}(b, a)) = \emptyset)}}
\vnbprop{card-subset-mono}{\vnbstmt{\forall b \in \mathbf{Set}, a.\; (((|b| \in \mathbb{N}) \wedge (a \subseteq b)) \Rightarrow (|a| \leq |b|))}}
\vnbprop{card-union-bound}{\vnbstmt{\begin{array}{@{}l@{}} \forall a, b \in \mathbf{Set}. \\ \quad ((|a| \in \mathbb{N}) \wedge (|b| \in \mathbb{N})) \Rightarrow (|(a \cup b)| \leq (|a| + |b|)) \end{array}}}
\vnbprop{card-union-nn}{\vnbstmt{\begin{array}{@{}l@{}} \forall a, b \in \mathbf{Set}. \\ \quad ((|a| \in \mathbb{N}) \wedge (|b| \in \mathbb{N})) \Rightarrow (|(a \cup b)| \in \mathbb{N}) \end{array}}}

\vnbbookfile{theorem-library/rake-combinatorics}{theorem-library/rake-combinatorics}
\noindent{\small A file shared across books; the entries that belong to this one.}\par
\noindent{\small rake-combinatorics.scm -- BATCH P of the 2026-09-17 rake: NN combinatorics, the ring powers, the powerset insert-split, and the two CHOOSE boundary values.  Seventeen asserted supports, PROVEN; every statement is copied LITERALLY from its site.}\par
\vnbprop{succ-nn-ord}{\vnbstmt{\forall n \in \mathbb{N}.\; ((n + 1) = (n + 1))}}
\vnbprop{union-empty-left}{\vnbstmt{\forall a \in \mathbf{Set}.\; ((\emptyset \cup a) = a)}}
\vnbprop{ord-segment-trans}{\vnbstmt{\begin{array}{@{}l@{}} \forall m \in \mathbb{N}, k \in \mathrm{OS}(m), i \in \mathrm{OS}(k). \\ \quad i \in \mathrm{OS}(m) \end{array}}}
\vnbprop{power-zero-base}{\vnbstmt{\forall n \in \mathbb{N}.\; (\operatorname{power}(0, (n + 1)) = 0)}}
\vnbprop{power-insert-disjoint}{\vnbstmt{\begin{array}{@{}l@{}} \forall x, k \in \mathbf{Set}. \\ \quad \neg (k \in x) \Rightarrow  \\ \quad \quad (\operatorname{power}(x) \cap \operatorname{image}((\lambda s \in \operatorname{power}(x).\; (s \cup \{k\})), \operatorname{power}(x))) = \emptyset \end{array}}}
\vnbprop{choose-set-unfold}{\vnbstmt{\forall n, m.\; (\operatorname{choose\text{-}set}(n, m) \simeq \{a \in \operatorname{power}(\mathrm{OS}(n)) : (|a| = m)\})}}
\vnbprop{choose-0-succ}{\vnbstmt{\forall k \in \mathbb{N}.\; (\operatorname{choose}(0, (k + 1)) = 0)}}
\vnbprop{choose-n-0}{\vnbstmt{\forall n \in \mathbb{N}.\; (\operatorname{choose}(n, 0) = 1)}}
\vnbprop{power-insert-cover}{\vnbstmt{\begin{array}{@{}l@{}} \forall x, k \in \mathbf{Set}. \\ \quad \neg (k \in x) \Rightarrow  \\ \quad \quad \operatorname{power}((x \cup \{k\})) \\ \quad \quad =\; (\operatorname{power}(x) \cup \operatorname{image}((\lambda s \in \operatorname{power}(x).\; (s \cup \{k\})), \operatorname{power}(x))) \end{array}}}
\vnbprop{subset-of-empty-is-empty}{\vnbstmt{\forall w \in \mathbf{Set}.\; (\forall z \in w.\; (z \in \emptyset) \Rightarrow (w = \emptyset))}}
\vnbprop{power-of-empty}{\vnbstmt{\operatorname{power}(\emptyset) = \{\emptyset\}}}
\vnbprop{card-power-nn}{\vnbstmt{\forall x \in \mathbf{Set}.\; ((|x| \in \mathbb{N}) \Rightarrow (|\operatorname{power}(x)| \in \mathbb{N}))}}

\vnbbookfile{theorem-library/rake-combinatorics2}{theorem-library/rake-combinatorics2}
\noindent{\small rake-combinatorics2.scm -- BATCH U of the 2026-09-18 rake (batch 5, the off-bill supports): finite combinatorics.  Every statement is copied LITERALLY from its support site; the sites are named beside each proof. Six results, all `proven modulo 0'.}\par
\vnbprop{interval-card}{\vnbstmt{\forall n \in \mathbb{N}.\; (|\operatorname{interval}(1, n)| = n)}}
\vnbprop{choose-in-nn}{\vnbstmt{\forall n, m \in \mathbb{N}.\; (\operatorname{choose}(n, m) \in \mathbb{N})}}
\vnbprop{injection-from-empty}{\vnbstmt{\forall c \in \mathbf{Set}.\; (|\operatorname{injection}(\emptyset, c)| = 1)}}
\vnbprop{permutations-zero}{\vnbstmt{|\operatorname{permutations}(0)| = 1}}
\vnbprop{fibrewise-finite}{\vnbstmt{\begin{array}{@{}l@{}} \forall p.; \forall g, t \in \mathbf{Set}. \\ \quad ((|g| \in \mathbb{N}) \wedge (\forall y \in t.\; (p(y) \in g) \wedge \forall c \in g.\; (|\{x \in t : (p(x) = c)\}| \in \mathbb{N}))) \Rightarrow  \\ \quad \quad |t| \in \mathbb{N} \end{array}}}
\vnbprop{pigeonhole-infinite}{\vnbstmt{\begin{array}{@{}l@{}} \forall s, f \in \mathbf{Set}, p \in (s \to f). \\ \quad (\neg (|s| \in \mathbb{N}) \wedge (|f| \in \mathbb{N})) \Rightarrow  \\ \quad \quad \exists c \in f.\; \neg (|\{x \in s : (p(x) = c)\}| \in \mathbb{N}) \end{array}}}

\vnbbookfile{theorem-library/rake-choose-succ}{theorem-library/rake-choose-succ}
\noindent{\small rake-choose-succ.scm -- PASCAL'S RULE for the binomial coefficient, and the CHOOSE-SET lemmas it is assembled from.  Rake batch 5b, assignment 5b-G.}\par
\vnbprop{choose-card-unfold}{\vnbstmt{\forall n, m.\; (|\operatorname{choose\text{-}set}(n, m)| \simeq \operatorname{choose}(n, m))}}
\vnbprop{choose-set-mem-iff}{\vnbstmt{\begin{array}{@{}l@{}} \forall n, m, a. \\ \quad (a \in \operatorname{choose\text{-}set}(n, m)) \Leftrightarrow ((a \in \operatorname{power}(\mathrm{OS}(n))) \wedge (|a| = m)) \end{array}}}
\vnbprop{choose-set-avoids-point}{\vnbstmt{\forall n \in \mathbb{N}, m.; \forall a \in \operatorname{choose\text{-}set}(n, m).\; \neg (n \in a)}}
\vnbprop{choose-set-widen}{\vnbstmt{\forall n \in \mathbb{N}, m.; \forall a \in \operatorname{choose\text{-}set}(n, m).\; (a \in \operatorname{choose\text{-}set}((n + 1), m))}}
\vnbprop{choose-set-insert-in}{\vnbstmt{\begin{array}{@{}l@{}} \forall n, k \in \mathbb{N}, s \in \operatorname{choose\text{-}set}(n, k). \\ \quad (s \cup \{n\}) \in \operatorname{choose\text{-}set}((n + 1), (k + 1)) \end{array}}}
\vnbprop{choose-set-remove-restores}{\vnbstmt{\begin{array}{@{}l@{}} \forall n, a \in \mathbf{Set}. \\ \quad (n \in a) \Rightarrow ((\operatorname{difference}(a, \{n\}) \cup \{n\}) = a) \end{array}}}
\vnbprop{choose-set-remove-in}{\vnbstmt{\begin{array}{@{}l@{}} \forall n, k \in \mathbb{N}, a \in \operatorname{choose\text{-}set}((n + 1), (k + 1)). \\ \quad (n \in a) \Rightarrow (\operatorname{difference}(a, \{n\}) \in \operatorname{choose\text{-}set}(n, k)) \end{array}}}
\vnbprop{choose-set-avoid-in}{\vnbstmt{\begin{array}{@{}l@{}} \forall n \in \mathbb{N}, m.; \forall a \in \operatorname{choose\text{-}set}((n + 1), m). \\ \quad \neg (n \in a) \Rightarrow (a \in \operatorname{choose\text{-}set}(n, m)) \end{array}}}
\vnbprop{choose-set-insert-injective}{\vnbstmt{\begin{array}{@{}l@{}} \forall n \in \mathbb{N}, m.; \forall s, t \in \operatorname{choose\text{-}set}(n, m). \\ \quad ((s \cup \{n\}) = (t \cup \{n\})) \Rightarrow (s = t) \end{array}}}
\vnbprop{choose-set-split}{\vnbstmt{\begin{array}{@{}l@{}} \forall n, k \in \mathbb{N}. \\ \quad \operatorname{choose\text{-}set}((n + 1), (k + 1)) \\ \quad =\; (\operatorname{choose\text{-}set}(n, (k + 1)) \cup \operatorname{image}((\lambda s \in \operatorname{choose\text{-}set}(n, k).\; (s \cup \{n\})), \operatorname{choose\text{-}set}(n, k))) \end{array}}}
\vnbprop{choose-set-split-disjoint}{\vnbstmt{\begin{array}{@{}l@{}} \forall n, k \in \mathbb{N}. \\ \quad (\operatorname{choose\text{-}set}(n, (k + 1)) \cap \operatorname{image}((\lambda s \in \operatorname{choose\text{-}set}(n, k).\; (s \cup \{n\})), \operatorname{choose\text{-}set}(n, k))) \\ \quad =\; \emptyset \end{array}}}
\vnbprop{choose-succ}{\vnbstmt{\begin{array}{@{}l@{}} \forall n, k \in \mathbb{N}. \\ \quad \operatorname{choose}((n + 1), (k + 1)) = (\operatorname{choose}(n, k) + \operatorname{choose}(n, (k + 1))) \end{array}}}

\vnbbookfile{theorem-library/nn-infinite}{theorem-library/nn-infinite}
\noindent{\small nn-infinite.scm -- NN is infinite, and the two leaves that hang on it.}\par
\vnbprop{inf-subsets-unfold}{\vnbstmt{\forall a.\; (\mathrm{Inf}(a) \simeq \{s \in \operatorname{power}(a) : \neg (|s| \in \mathbb{N})\})}}
\vnbprop{nn-not-finite}{\vnbstmt{\neg (|\mathbb{N}| \in \mathbb{N})}}
\vnbprop{nn-in-inf-subsets}{\vnbstmt{\mathbb{N} \in \mathrm{Inf}(\mathbb{N})}}
\vnbprop{inf-subset-nn-unbounded}{\vnbstmt{\forall t \in \mathrm{Inf}(\mathbb{N}), u \in \mathbb{N}.\; \exists y \in \mathbb{N}.\; ((y \in t) \wedge (u < y))}}

\vnbbookfile{theorem-library/fin-subsets}{theorem-library/fin-subsets}
\noindent{\small fin-subsets.scm -- FIN-SUBSETS(a), the finite subsets of a, and the --------------------------------------------------------------------------- NOTE, 2026-09-20 (batch 9-B).  This file was written while cardinality was AXIOMATISED under the name CARD and the defined constant was its companion}\par
\par\medskip\noindent\textbf{Definition} (\texttt{fin-subsets}).\;$\operatorname{fin\text{-}subsets}(\operatorname{a\_}) :=$ $\{\operatorname{t\_} \in \operatorname{power}(\operatorname{a\_}) : (|\operatorname{t\_}| \in \mathbb{N})\}$
\vnbprop{fin-subsets-membership}{\vnbstmt{\begin{array}{@{}l@{}} \forall a, s. \\ \quad (s \in \operatorname{fin\text{-}subsets}(a)) \Leftrightarrow ((s \in \mathbf{Set}) \wedge (\forall z \in s.\; (z \in a) \wedge (|s| \in \mathbb{N}))) \end{array}}}
\vnbprop{fin-subsets-is-set}{\vnbstmt{\forall a \in \mathbf{Set}.\; (\operatorname{fin\text{-}subsets}(a) \in \mathbf{Set})}}
\vnbprop{fin-subsets-has-empty}{\vnbstmt{\forall a.\; (\emptyset \in \operatorname{fin\text{-}subsets}(a))}}
\vnbprop{fin-subsets-union-closed}{\vnbstmt{\begin{array}{@{}l@{}} \forall a, s, u. \\ \quad ((s \in \operatorname{fin\text{-}subsets}(a)) \wedge (u \in \operatorname{fin\text{-}subsets}(a))) \Rightarrow  \\ \quad \quad (s \cup u) \in \operatorname{fin\text{-}subsets}(a) \end{array}}}

\vnbbookfile{theorem-library/nn-pairing}{theorem-library/nn-pairing}
\noindent{\small nn-pairing.scm -- the Cantor pairing NN x NN -$>$ NN, and the arithmetic it needs.  Step one of the missing re-indexing mechanism.}\par
\par\medskip\noindent\textbf{Definition} (\texttt{trinum}).\; $(\operatorname{trinum}(0) \simeq 0)$\par
Suppose $\operatorname{n\_} \in \mathbb{N}$. Then $(\operatorname{trinum}(\operatorname{succ}(\operatorname{n\_})) \simeq +(\operatorname{trinum}(\operatorname{n\_}), \operatorname{succ}(\operatorname{n\_})))$
\par\medskip\noindent\textbf{Definition} (\texttt{nnpair}).\;$\operatorname{nnpair}(\operatorname{i\_}, \operatorname{j\_}) :=$ $\operatorname{trinum}((\operatorname{i\_} + \operatorname{j\_})) + \operatorname{j\_}$
\par\medskip\noindent\textbf{Definition} (\texttt{nnfst}).\;$\operatorname{nnfst}(\operatorname{n\_}) :=$
\begin{equation*}
(\iota\, \operatorname{a\_}.\;((\operatorname{a\_} \in \mathbb{N}) \wedge \exists \operatorname{b\_} \in \mathbb{N}.\; (\operatorname{nnpair}(\operatorname{a\_}, \operatorname{b\_}) = \operatorname{n\_})))
\end{equation*}

\par\medskip\noindent\textbf{Definition} (\texttt{nnsnd}).\;$\operatorname{nnsnd}(\operatorname{n\_}) :=$
\begin{equation*}
(\iota\, \operatorname{b\_}.\;((\operatorname{b\_} \in \mathbb{N}) \wedge \exists \operatorname{a\_} \in \mathbb{N}.\; (\operatorname{nnpair}(\operatorname{a\_}, \operatorname{b\_}) = \operatorname{n\_})))
\end{equation*}

\vnbprop{trinum-type}{\vnbstmt{\forall n \in \mathbb{N}.\; (\operatorname{trinum}(n) \in \mathbb{N})}}
\vnbprop{trinum-mono}{\vnbstmt{\forall b, a \in \mathbb{N}.\; ((a \leq b) \Rightarrow (\operatorname{trinum}(a) \leq \operatorname{trinum}(b)))}}
\vnbprop{nnpair-onto}{\vnbstmt{\forall n \in \mathbb{N}.\; \exists i, j \in \mathbb{N}.\; (\operatorname{nnpair}(i, j) = n)}}
\vnbprop{nnpair-diag-bound}{\vnbstmt{\begin{array}{@{}l@{}} \forall s, j. \\ \quad ((s \in \mathbb{N}) \wedge ((j \in \mathbb{N}) \wedge (j \leq s))) \Rightarrow (((\operatorname{trinum}(s) + j) + 1) \leq \operatorname{trinum}((s + 1))) \end{array}}}
\vnbprop{nnpair-cross}{\vnbstmt{\begin{array}{@{}l@{}} \forall s, t, j, k. \\ \quad ((s \in \mathbb{N}) \wedge ((t \in \mathbb{N}) \wedge ((j \in \mathbb{N}) \wedge ((k \in \mathbb{N}) \wedge ((j \leq s) \wedge ((s + 1) \leq t)))))) \Rightarrow  \\ \quad \quad ((\operatorname{trinum}(s) + j) + 1) \leq (\operatorname{trinum}(t) + k) \end{array}}}
\vnbprop{nnpair-inj}{\vnbstmt{\begin{array}{@{}l@{}} \forall i, j, \mathit{ii}, \mathit{jj}. \\ \quad ((i \in \mathbb{N}) \wedge ((j \in \mathbb{N}) \wedge ((\mathit{ii} \in \mathbb{N}) \wedge ((\mathit{jj} \in \mathbb{N}) \wedge (\operatorname{nnpair}(i, j) = \operatorname{nnpair}(\mathit{ii}, \mathit{jj})))))) \Rightarrow  \\ \quad \quad (i = \mathit{ii}) \wedge (j = \mathit{jj}) \end{array}}}
\vnbprop{nnpair-type}{\vnbstmt{\forall i, j.\; (((i \in \mathbb{N}) \wedge (j \in \mathbb{N})) \Rightarrow (\operatorname{nnpair}(i, j) \in \mathbb{N}))}}
\vnbprop{nnfst-nnpair}{\vnbstmt{\forall i, j.\; (((i \in \mathbb{N}) \wedge (j \in \mathbb{N})) \Rightarrow (\operatorname{nnfst}(\operatorname{nnpair}(i, j)) = i))}}
\vnbprop{nnsnd-nnpair}{\vnbstmt{\forall i, j.\; (((i \in \mathbb{N}) \wedge (j \in \mathbb{N})) \Rightarrow (\operatorname{nnsnd}(\operatorname{nnpair}(i, j)) = j))}}
\vnbprop{nnfst-type}{\vnbstmt{\forall n \in \mathbb{N}.\; (\operatorname{nnfst}(n) \in \mathbb{N})}}
\vnbprop{nnsnd-type}{\vnbstmt{\forall n \in \mathbb{N}.\; (\operatorname{nnsnd}(n) \in \mathbb{N})}}
\vnbprop{nn-flatten}{\vnbstmt{\begin{array}{@{}l@{}} \forall a.; \forall h \in (\mathbb{N} \to (\mathbb{N} \to a)). \exists e \in (\mathbb{N} \to a). \forall i, j. \\ \quad ((i \in \mathbb{N}) \wedge (j \in \mathbb{N})) \Rightarrow \exists m \in \mathbb{N}.\; (e(m) = h(i)(j)) \end{array}}}

\vnbbookfile{theorem-library/nn-nested-subset-chain-proof}{theorem-library/nn-nested-subset-chain-proof}
\noindent{\small nn-nested-subset-chain-proof.scm -- a descending NN-indexed family is a chain.}\par
\vnbprop{nn-nested-subset-chain-j}{\vnbstmt{\begin{array}{@{}l@{}} \forall m.; \forall j, k \in \mathbb{N}. \\ \quad (\forall k \in \mathbb{N}.\; (m(k) \in \mathbf{Set}) \wedge (\forall k \in \mathbb{N}.\; (m((k + 1)) \subseteq m(k)) \wedge (k \leq j))) \Rightarrow  \\ \quad \quad m(j) \subseteq m(k) \end{array}}}
\vnbprop{nn-nested-subset-chain}{\vnbstmt{\begin{array}{@{}l@{}} \forall t.; \forall k, j \in \mathbb{N}. \\ \quad (\forall k \in \mathbb{N}.\; (t(k) \in \mathbf{Set}) \wedge (\forall k \in \mathbb{N}.\; (t((k + 1)) \subseteq t(k)) \wedge (k \leq j))) \Rightarrow  \\ \quad \quad t(j) \subseteq t(k) \end{array}}}

\vnbbookfile{structure-library/rabasics-vocab}{structure-library/rabasics-vocab}
\noindent{\small rabasics-vocab.scm -- vocabulary for the RA-BASICS wave (calculus.pdf 1.13, 3.7; CA 1.7). DEFINITIONS ONLY (destination: structure-library/).  Agent RA-BASICS, 2026-10-26 (NEW DEFINITIONS, the user's to confirm; added under the rule "missing vocabulary is not a stopping point").}\par
\par\medskip\noindent\textbf{Definition} (\texttt{inf}).\;$\operatorname{inf}(\operatorname{infs\_}) :=$
\begin{equation*}
-\operatorname{sup}(\{\operatorname{infv\_} \in \mathbb{R} : \exists \operatorname{infx\_} \in \operatorname{infs\_}.\; (\operatorname{infv\_} = -\operatorname{infx\_})\})
\end{equation*}

\par\medskip\noindent\textbf{Definition} (\texttt{is-order-iso}).\;For all $\operatorname{oisf\_}$, $\operatorname{oisa\_}$ and $\operatorname{oisb\_}$, 
\begin{equation*}
\begin{array}{@{}l@{}} \operatorname{is\text{-}order\text{-}iso}(\operatorname{oisf\_}, \operatorname{oisa\_}, \operatorname{oisb\_}) \\ \Leftrightarrow\; \operatorname{oisf\_} \in \mathrm{Bij}(\operatorname{oisa\_}, \operatorname{oisb\_}) \wedge \\ \forall \operatorname{oisx\_}, \operatorname{oisy\_} \in \operatorname{oisa\_}. \\ \quad (\operatorname{oisx\_} < \operatorname{oisy\_}) \Rightarrow (\operatorname{oisf\_}(\operatorname{oisx\_}) < \operatorname{oisf\_}(\operatorname{oisy\_})) \end{array}
\end{equation*}


\vnbbookfile{structure-library/order-zorn}{structure-library/order-zorn}
\noindent{\small MOVED 2026-09-20 (batch 12-A) from theorem-library/: this file is VOCABULARY -- definitions, notation and warranted supports, not one proof -- and every theorem stated with it had to load below it.  Its load.scm slot is unchanged. order-zorn.scm -- abstract partial orders and ZORN'S LEMMA (statement only).}\par
\par\medskip\noindent\textbf{Definition} (\texttt{is-partial-order}).\;For all $\mathit{grd}$ and $\mathit{porel}$, 
\begin{equation*}
\begin{array}{@{}l@{}} \operatorname{is\text{-}partial\text{-}order}(\mathit{grd}, \mathit{porel}) \\ \Leftrightarrow\; \mathit{grd} \in \mathbf{Set} \wedge \\ \mathit{porel} \in \mathbf{Set} \wedge \\ \forall x \in \mathit{grd}.\; ({[}x, x] \in \mathit{porel}) \wedge \\ \forall x, y \in \mathit{grd}.\; ((({[}x, y] \in \mathit{porel}) \wedge ({[}y, x] \in \mathit{porel})) \Rightarrow (x = y)) \wedge \\ \forall x, y, z \in \mathit{grd}. \\ \quad (({[}x, y] \in \mathit{porel}) \wedge ({[}y, z] \in \mathit{porel})) \Rightarrow ({[}x, z] \in \mathit{porel}) \end{array}
\end{equation*}

\par\medskip\noindent\textbf{Definition} (\texttt{is-chain}).\;For all $\mathit{grd}$, $\mathit{porel}$ and $\mathit{ch}$, 
\begin{equation*}
\begin{array}{@{}l@{}} \operatorname{is\text{-}chain}(\mathit{grd}, \mathit{porel}, \mathit{ch}) \\ \Leftrightarrow\; (\mathit{ch} \subseteq \mathit{grd}) \wedge \forall x, y \in \mathit{ch}.\; (({[}x, y] \in \mathit{porel}) \vee ({[}y, x] \in \mathit{porel})) \end{array}
\end{equation*}

\par\medskip\noindent\textbf{Definition} (\texttt{is-upper-bound}).\;For all $\mathit{grd}$, $\mathit{porel}$, $\mathit{ch}$ and $b$, 
\begin{equation*}
\operatorname{is\text{-}upper\text{-}bound}(\mathit{grd}, \mathit{porel}, \mathit{ch}, b) \Leftrightarrow ((b \in \mathit{grd}) \wedge \forall x \in \mathit{ch}.\; ({[}x, b] \in \mathit{porel}))
\end{equation*}

\par\medskip\noindent\textbf{Definition} (\texttt{is-maximal}).\;For all $\mathit{grd}$, $\mathit{porel}$ and $\mathit{mx}$, 
\begin{equation*}
\begin{array}{@{}l@{}} \operatorname{is\text{-}maximal}(\mathit{grd}, \mathit{porel}, \mathit{mx}) \\ \Leftrightarrow\; (\mathit{mx} \in \mathit{grd}) \wedge \forall y \in \mathit{grd}.\; (({[}\mathit{mx}, y] \in \mathit{porel}) \Rightarrow ({[}y, \mathit{mx}] \in \mathit{porel})) \end{array}
\end{equation*}


\vnbbookfile{theorem-library/injective-star}{theorem-library/injective-star}
\noindent{\small injective-star.scm -- the bridge from the set-function vocabulary to the class-function one.}\par
\vnbprop{injection-is-injective-star}{\vnbstmt{\forall x, y.; \forall f \in \operatorname{injection}(x, y).\; \operatorname{injective\text{-}star}(f)}}

\vnbbookfile{theorem-library/zorn-proof}{theorem-library/zorn-proof}
\noindent{\small zorn-proof.scm -- proving ZORN'S LEMMA from the well-ordering principle, via the GREEDY-KEPT-CHAIN construction (avoids the cardinality/Hartogs comparison lemma the base lacks).  Multi-session build; this file grows as lemmas land. Started 2026-07-22.}\par
\par\medskip\noindent\textbf{Definition} (\texttt{zkept}).\;For all $\varphi$, $\mathit{grd}$ and $\mathit{porel}$,  $(\operatorname{zkept}(\varphi, \mathit{grd}, \mathit{porel}, 0) \simeq \emptyset)$\par
Let $\varphi$, $\mathit{grd}$ and $\mathit{porel}$ be arbitrary. Let $\alpha \in \mathrm{Ord}$. Then
\begin{equation*}
(\operatorname{zkept}(\varphi, \mathit{grd}, \mathit{porel}, (\alpha + 1)) \simeq (\operatorname{zkept}(\varphi, \mathit{grd}, \mathit{porel}, \alpha) \cup \operatorname{keep\text{-}set}(\varphi, \mathit{grd}, \mathit{porel}, \operatorname{zkept}(\varphi, \mathit{grd}, \mathit{porel}, \alpha), \alpha)))
\end{equation*}
\par
Let $\varphi$, $\mathit{grd}$ and $\mathit{porel}$ be arbitrary. Let $\mathit{lam} \text{ be s.t. } \operatorname{limit\text{-}ord}(\mathit{lam})$. Then
\begin{equation*}
(\operatorname{zkept}(\varphi, \mathit{grd}, \mathit{porel}, \mathit{lam}) \simeq \operatorname{big\text{-}union}(\beta, \mathrm{OS}(\mathit{lam}), \operatorname{zkept}(\varphi, \mathit{grd}, \mathit{porel}, \beta)))
\end{equation*}

\par\medskip\noindent\textbf{Definition} (\texttt{greedy-chain}).\;$\operatorname{greedy\text{-}chain}(\varphi, \mathit{grd}, \mathit{porel}) :=$ $\operatorname{zkept}(\varphi, \mathit{grd}, \mathit{porel}, |\mathit{grd}|)$
\par\medskip\noindent\textbf{Definition} (\texttt{keep-set}).\;$\operatorname{keep\text{-}set}(\varphi, \mathit{grd}, \mathit{porel}, \mathit{kset}, \alpha) :=$
\begin{equation*}
\{y \in \mathit{grd} : ((y = \varphi(\alpha)) \wedge \forall z \in \mathit{kset}.\; ({[}z, y] \in \mathit{porel}))\}
\end{equation*}

\vnbdef{keep-set-membership}{\vnbstmt{\begin{array}{@{}l@{}} \forall i, d, l, t, a, y. \\ \quad (y \in \operatorname{keep\text{-}set}(i, d, l, t, a)) \Leftrightarrow ((y \in d) \wedge ((y = i(a)) \wedge \forall z \in t.\; ({[}z, y] \in l))) \end{array}}}
\vnbprop{zorn-empty-is-chain}{\vnbstmt{\forall d \in \mathbf{Set}, l.\; \operatorname{is\text{-}chain}(d, l, \emptyset)}}
\vnbprop{zorn-chain-plus-keepset}{\vnbstmt{\begin{array}{@{}l@{}} \forall i, d, l, t, a. \\ \quad (\operatorname{is\text{-}partial\text{-}order}(d, l) \wedge \operatorname{is\text{-}chain}(d, l, t)) \Rightarrow  \\ \quad \quad \operatorname{is\text{-}chain}(d, l, (t \cup \operatorname{keep\text{-}set}(i, d, l, t, a))) \end{array}}}
\vnbprop{zkept-defined}{\vnbstmt{\forall a \in \mathrm{Ord}, i, d, l.\; (\operatorname{zkept}(i, d, l, a) = \operatorname{zkept}(i, d, l, a))}}
\vnbprop{zorn-zkept-monotone}{\vnbstmt{\begin{array}{@{}l@{}} \forall b \in \mathrm{Ord}, i, d, l.; \forall a \in \mathrm{Ord}. \\ \quad \operatorname{ord\text{-}le}(a, b) \Rightarrow (\operatorname{zkept}(i, d, l, a) \subseteq \operatorname{zkept}(i, d, l, b)) \end{array}}}
\vnbprop{zorn-zkept-limit-is-chain}{\vnbstmt{\begin{array}{@{}l@{}} \forall i, d, l, m. \\ \quad \operatorname{is\text{-}partial\text{-}order}(d, l) \wedge \\ \quad \operatorname{limit\text{-}ord}(m) \wedge \\ \quad \forall a \in \mathrm{OS}(m).\; \operatorname{is\text{-}chain}(d, l, \operatorname{zkept}(i, d, l, a)) \Rightarrow \\ \quad \quad \operatorname{is\text{-}chain}(d, l, \operatorname{big\text{-}union}(a_{2}, \mathrm{OS}(m), \operatorname{zkept}(i, d, l, a_{2}))) \end{array}}}
\vnbprop{zorn-zkept-is-chain}{\vnbstmt{\begin{array}{@{}l@{}} \forall a \in \mathrm{Ord}, i, d, l. \\ \quad \operatorname{is\text{-}partial\text{-}order}(d, l) \Rightarrow \operatorname{is\text{-}chain}(d, l, \operatorname{zkept}(i, d, l, a)) \end{array}}}
\vnbprop{greedy-chain-is-chain}{\vnbstmt{\begin{array}{@{}l@{}} \forall i, d, l. \\ \quad \operatorname{is\text{-}partial\text{-}order}(d, l) \Rightarrow \operatorname{is\text{-}chain}(d, l, \operatorname{greedy\text{-}chain}(i, d, l)) \end{array}}}

\vnbbookfile{theorem-library/zorn-route-two}{theorem-library/zorn-route-two}
\noindent{\small zorn-route-two.scm -- ZORN'S LEMMA, PROVED.}\par
\par\medskip\noindent\textbf{Definition} (\texttt{zup}).\;For all $\mathit{grd}$ and $\mathit{porel}$,  $(\operatorname{zup}(\mathit{grd}, \mathit{porel}, 0) \simeq \varepsilon(\mathit{grd}))$\par
Let $\mathit{grd}$ and $\mathit{porel}$ be arbitrary. Let $\alpha \in \mathrm{Ord}$. Then
\begin{equation*}
(\operatorname{zup}(\mathit{grd}, \mathit{porel}, (\alpha + 1)) \simeq \varepsilon(\{\operatorname{y\_} \in \mathit{grd} : \operatorname{is\text{-}strictly\text{-}below}(\mathit{porel}, \operatorname{zup}(\mathit{grd}, \mathit{porel}, \alpha), \operatorname{y\_})\}))
\end{equation*}
\par
Let $\mathit{grd}$ and $\mathit{porel}$ be arbitrary. Let $\mathit{lam} \text{ be s.t. } \operatorname{limit\text{-}ord}(\mathit{lam})$. Then
\begin{equation*}
(\operatorname{zup}(\mathit{grd}, \mathit{porel}, \mathit{lam}) \simeq \varepsilon(\{\operatorname{y\_} \in \mathit{grd} : \forall \operatorname{c\_} \in \mathrm{Ord}.\; (\operatorname{ord\text{-}lt}(\operatorname{c\_}, \mathit{lam}) \Rightarrow ({[}\operatorname{zup}(\mathit{grd}, \mathit{porel}, \operatorname{c\_}), \operatorname{y\_}] \in \mathit{porel}))\}))
\end{equation*}

\par\medskip\noindent\textbf{Definition} (\texttt{is-strictly-below}).\;For all $\mathit{porel}$, $x$ and $y$, 
\begin{equation*}
\operatorname{is\text{-}strictly\text{-}below}(\mathit{porel}, x, y) \Leftrightarrow (({[}x, y] \in \mathit{porel}) \wedge \neg ({[}y, x] \in \mathit{porel}))
\end{equation*}

\vnbprop{strictly-below-irreflexive}{\vnbstmt{\forall l, x.\; \neg \operatorname{is\text{-}strictly\text{-}below}(l, x, x)}}
\vnbprop{strictly-below-transitive}{\vnbstmt{\begin{array}{@{}l@{}} \forall d, l, x, y, z. \\ \quad \operatorname{is\text{-}partial\text{-}order}(d, l) \wedge \\ \quad x \in d \wedge \\ \quad y \in d \wedge \\ \quad z \in d \wedge \\ \quad \operatorname{is\text{-}strictly\text{-}below}(l, x, y) \wedge \\ \quad \operatorname{is\text{-}strictly\text{-}below}(l, y, z) \Rightarrow \\ \quad \quad \operatorname{is\text{-}strictly\text{-}below}(l, x, z) \end{array}}}
\vnbprop{strictly-below-not-equal}{\vnbstmt{\forall l, x, y.\; ((x = y) \Rightarrow \neg \operatorname{is\text{-}strictly\text{-}below}(l, x, y))}}
\vnbprop{no-maximal-strictly-above}{\vnbstmt{\begin{array}{@{}l@{}} \forall d, l.; \forall x \in d. \\ \quad \neg \exists \mathit{mx}.\; \operatorname{is\text{-}maximal}(d, l, \mathit{mx}) \Rightarrow  \\ \quad \quad \exists y \in d.\; \operatorname{is\text{-}strictly\text{-}below}(l, x, y) \end{array}}}
\vnbprop{po-reflexive}{\vnbstmt{\forall d, l, x.\; ((\operatorname{is\text{-}partial\text{-}order}(d, l) \wedge (x \in d)) \Rightarrow ({[}x, x] \in l))}}
\vnbprop{po-transitive}{\vnbstmt{\begin{array}{@{}l@{}} \forall d, l, x, y, z. \\ \quad (\operatorname{is\text{-}partial\text{-}order}(d, l) \wedge ((x \in d) \wedge ((y \in d) \wedge ((z \in d) \wedge (({[}x, y] \in l) \wedge ({[}y, z] \in l)))))) \Rightarrow  \\ \quad \quad {[}x, z] \in l \end{array}}}
\vnbprop{ord-lt-succ-cases}{\vnbstmt{\begin{array}{@{}l@{}} \forall a, b. \\ \quad ((a \in \mathrm{Ord}) \wedge ((b \in \mathrm{Ord}) \wedge \operatorname{ord\text{-}lt}(b, (a + 1)))) \Rightarrow (\operatorname{ord\text{-}lt}(b, a) \vee (b = a)) \end{array}}}
\vnbprop{ord-succ-lt-limit}{\vnbstmt{\begin{array}{@{}l@{}} \forall m, b. \\ \quad (\operatorname{limit\text{-}ord}(m) \wedge ((b \in \mathrm{Ord}) \wedge \operatorname{ord\text{-}lt}(b, m))) \Rightarrow \operatorname{ord\text{-}lt}((b + 1), m) \end{array}}}
\vnbprop{zup-tower}{\vnbstmt{\begin{array}{@{}l@{}} \forall a \in \mathrm{Ord}, d, l. \\ \quad \operatorname{is\text{-}partial\text{-}order}(d, l) \wedge \\ \quad \exists w.\; (w \in d) \wedge \\ \quad \forall h.\; (\operatorname{is\text{-}chain}(d, l, h) \Rightarrow \exists u.\; \operatorname{is\text{-}upper\text{-}bound}(d, l, h, u)) \wedge \\ \quad \forall x \in d.\; \exists y \in d.\; \operatorname{is\text{-}strictly\text{-}below}(l, x, y) \Rightarrow \\ \quad \quad \operatorname{zup}(d, l, a) \in d \wedge \\ \quad \quad \forall b. \\ \quad \quad \quad \operatorname{ord\text{-}lt}(b, a) \Rightarrow  \\ \quad \quad \quad \quad (\operatorname{zup}(d, l, b) \in d) \wedge \operatorname{is\text{-}strictly\text{-}below}(l, \operatorname{zup}(d, l, b), \operatorname{zup}(d, l, a)) \end{array}}}
\vnbprop{zup-injective}{\vnbstmt{\begin{array}{@{}l@{}} \forall d, l, a_{1}, a_{2}. \\ \quad \operatorname{is\text{-}partial\text{-}order}(d, l) \wedge \\ \quad \exists w.\; (w \in d) \wedge \\ \quad \forall h.\; (\operatorname{is\text{-}chain}(d, l, h) \Rightarrow \exists u.\; \operatorname{is\text{-}upper\text{-}bound}(d, l, h, u)) \wedge \\ \quad \forall x \in d.\; \exists y \in d.\; \operatorname{is\text{-}strictly\text{-}below}(l, x, y) \wedge \\ \quad a_{1} \in \mathrm{Ord} \wedge \\ \quad a_{2} \in \mathrm{Ord} \wedge \\ \quad \operatorname{zup}(d, l, a_{1}) = \operatorname{zup}(d, l, a_{2}) \Rightarrow \\ \quad \quad a_{1} = a_{2} \end{array}}}
\vnbprop{zorn-lemma}{\vnbstmt{\begin{array}{@{}l@{}} \forall d, l. \\ \quad \operatorname{is\text{-}partial\text{-}order}(d, l) \wedge \\ \quad \exists w.\; (w \in d) \wedge \\ \quad \forall h.\; (\operatorname{is\text{-}chain}(d, l, h) \Rightarrow \exists u.\; \operatorname{is\text{-}upper\text{-}bound}(d, l, h, u)) \Rightarrow \\ \quad \quad \exists x.\; \operatorname{is\text{-}maximal}(d, l, x) \end{array}}}

\vnbbookfile{theorem-library/card-cartesian-nn}{theorem-library/card-cartesian-nn}
\noindent{\small card-cartesian-nn.scm -- lemmas MOVED 2026-10-27 (integration of the 10-26/27 wave) out of theorem-library/monalg-is-ring.scm, so that files of library core may cite them (the library gate). The proofs are the originals, verbatim; the helper defines are a copy of monalg-is-ring.scm's.  Loaded immediately before ...}\par
\vnbprop{card-cartesian-nn}{\vnbstmt{\begin{array}{@{}l@{}} \forall x, y. \\ \quad ((x \in \mathbf{Set}) \wedge ((|x| \in \mathbb{N}) \wedge ((y \in \mathbf{Set}) \wedge (|y| \in \mathbb{N})))) \Rightarrow  \\ \quad \quad |(x \times y)| \in \mathbb{N} \end{array}}}

\vnbbookfile{theorem-library/nn-least-element}{theorem-library/nn-least-element}
\noindent{\small theorem-library/nn-least-element.scm =================================================================== The well-ordering of NN, PROVEN from the well-ordering of ORD.}\par
\vnbprop{nn-least-element}{\vnbstmt{\begin{array}{@{}l@{}} \forall t. \\ \quad ((t \subseteq \mathbb{N}) \wedge \exists n.\; (n \in t)) \Rightarrow  \\ \quad \quad \exists m \in t.\; \forall k \in t.\; (m \leq k) \end{array}}}

\vnbbookfile{theorem-library/pair-self-cartesian3}{theorem-library/pair-self-cartesian3}
\noindent{\small pair-self-cartesian3.scm -- three generic set lemmas: a set is the one member of its own singleton PAIR(a, a) (pair-self-member, pair-self-eq), and the ternary product of three sets is a set (cartesian3-set).}\par
\vnbprop{pair-self-member}{\vnbstmt{\forall a \in \mathbf{Set}.\; (a \in \{a\})}}
\vnbprop{pair-self-eq}{\vnbstmt{\forall a, x.\; (((a \in \mathbf{Set}) \wedge (x \in \{a\})) \Rightarrow (x = a))}}
\vnbprop{cartesian3-set}{\vnbstmt{\begin{array}{@{}l@{}} \forall u, v, w. \\ \quad ((u \in \mathbf{Set}) \wedge ((v \in \mathbf{Set}) \wedge (w \in \mathbf{Set}))) \Rightarrow ((u \times v) \in \mathbf{Set}) \end{array}}}

\vnbbookfile{theorem-library/th3-sort}{theorem-library/th3-sort}
\noindent{\small th3-sort.scm -- wave of 2026-11-09/10 (agent TH3-SORT), wired 2026-11-10 from scratchpad/th3-1109/sort/th3-sort.scm. Library: core.}\par
\vnbprop{srt-finite-max}{\vnbstmt{\begin{array}{@{}l@{}} \forall s, f \in \mathbf{Set}. \\ \quad ((|s| \in \mathbb{N}) \wedge (\forall z \in s.\; (f(z) \in \mathbb{R}) \wedge \exists w.\; (w \in s))) \Rightarrow  \\ \quad \quad \exists m \in s.\; \forall z \in s.\; (f(z) \leq f(m)) \end{array}}}
\vnbprop{srt-maximizer}{\vnbstmt{\begin{array}{@{}l@{}} \forall p, f \in \mathbf{Set}, q.; \forall w \in q. \\ \quad \forall x \in p.\; (f(x) \in \mathbb{R}) \wedge \\ \quad \forall x \in p.\; (|\{y \in p : (f(x) \leq f(y))\}| \in \mathbb{N}) \wedge \\ \quad \forall z \in q.\; (z \in p) \Rightarrow \\ \quad \quad \exists m \in q.\; \forall z \in q.\; (f(z) \leq f(m)) \end{array}}}
\vnbprop{srt-chain}{\vnbstmt{\begin{array}{@{}l@{}} \forall p, f \in \mathbf{Set}. \\ \quad (\forall x \in p.\; (f(x) \in \mathbb{R}) \wedge \forall x \in p.\; (|\{y \in p : (f(x) \leq f(y))\}| \in \mathbb{N})) \Rightarrow  \\ \quad \quad \exists t \in (\mathbb{N} \to \{u \in \operatorname{power}(p) : (|u| \in \mathbb{N})\}). \\ \quad \quad \quad t(0) = \emptyset \wedge \\ \quad \quad \quad \forall k \in \mathbb{N}. \\ \quad \quad \quad \quad \forall z \in p.\; (z \in t(k)) \wedge (t((k + 1)) = t(k)) \vee \\ \quad \quad \quad \quad \exists \mathit{srx} \in p. \\ \quad \quad \quad \quad \quad \neg (\mathit{srx} \in t(k)) \wedge \\ \quad \quad \quad \quad \quad \forall z \in p.\; (\neg (z \in t(k)) \Rightarrow (f(z) \leq f(\mathit{srx}))) \wedge \\ \quad \quad \quad \quad \quad t((k + 1)) = (t(k) \cup \{\mathit{srx}\}) \end{array}}}
\vnbprop{srt-inf-picks}{\vnbstmt{\begin{array}{@{}l@{}} \forall p, f \in \mathbf{Set}, t \in (\mathbb{N} \to \{u \in \operatorname{power}(p) : (|u| \in \mathbb{N})\}). \\ \quad \forall x \in p.\; (f(x) \in \mathbb{R}) \wedge \\ \quad \forall x \in p.\; (|\{y \in p : (f(x) \leq f(y))\}| \in \mathbb{N}) \wedge \\ \quad \neg (|p| \in \mathbb{N}) \wedge \\ \quad \forall k \in \mathbb{N}. \\ \quad \quad \forall z \in p.\; (z \in t(k)) \wedge (t((k + 1)) = t(k)) \vee \\ \quad \quad \exists \mathit{srx} \in p. \\ \quad \quad \quad \neg (\mathit{srx} \in t(k)) \wedge \\ \quad \quad \quad \forall z \in p.\; (\neg (z \in t(k)) \Rightarrow (f(z) \leq f(\mathit{srx}))) \wedge \\ \quad \quad \quad t((k + 1)) = (t(k) \cup \{\mathit{srx}\}) \Rightarrow \\ \quad \quad \exists h \in (\mathbb{N} \to p). \forall k \in \mathbb{N}. \\ \quad \quad \quad h(k) \in p \wedge \\ \quad \quad \quad \neg (h(k) \in t(k)) \wedge \\ \quad \quad \quad \forall z \in p.\; (\neg (z \in t(k)) \Rightarrow (f(z) \leq f(h(k)))) \wedge \\ \quad \quad \quad t((k + 1)) = (t(k) \cup \{h(k)\}) \end{array}}}
\vnbprop{srt-picks-sorted}{\vnbstmt{\begin{array}{@{}l@{}} \forall p, f \in \mathbf{Set}, t \in (\mathbb{N} \to \{u \in \operatorname{power}(p) : (|u| \in \mathbb{N})\}), h \in (\mathbb{N} \to p). \\ \quad \forall x \in p.\; (f(x) \in \mathbb{R}) \wedge \\ \quad \forall x \in p.\; (|\{y \in p : (f(x) \leq f(y))\}| \in \mathbb{N}) \wedge \\ \quad t(0) = \emptyset \wedge \\ \quad \forall k \in \mathbb{N}. \\ \quad \quad h(k) \in p \wedge \\ \quad \quad \neg (h(k) \in t(k)) \wedge \\ \quad \quad \forall z \in p.\; (\neg (z \in t(k)) \Rightarrow (f(z) \leq f(h(k)))) \wedge \\ \quad \quad t((k + 1)) = (t(k) \cup \{h(k)\}) \Rightarrow \\ \quad \quad h \in \mathrm{Bij}(\mathbb{N}, p) \wedge \\ \quad \quad \forall n \in \mathbb{N}.\; (f(h((n + 1))) \leq f(h(n))) \wedge \\ \quad \quad \forall n, m \in \mathbb{N}.\; ((m \leq n) \Rightarrow (f(h(n)) \leq f(h(m)))) \wedge \\ \quad \quad \forall c \in \mathbb{R}. \\ \quad \quad \quad \forall n \in \mathbb{N}.\; (c \leq f(h(n))) \Rightarrow \neg (|\{y \in p : (c \leq f(y))\}| \in \mathbb{N}) \end{array}}}
\vnbprop{srt-sort-infinite}{\vnbstmt{\begin{array}{@{}l@{}} \forall p, f \in \mathbf{Set}. \\ \quad \forall x \in p.\; (f(x) \in \mathbb{R}) \wedge \\ \quad \forall x \in p.\; (0 < f(x)) \wedge \\ \quad \forall d.\; (\operatorname{pos\text{-}rr}(d) \Rightarrow (|\{y \in p : (d \leq f(y))\}| \in \mathbb{N})) \wedge \\ \quad \neg (|p| \in \mathbb{N}) \Rightarrow \\ \quad \quad \exists h \in \mathrm{Bij}(\mathbb{N}, p). \\ \quad \quad \quad \forall n \in \mathbb{N}.\; (f(h((n + 1))) \leq f(h(n))) \wedge \\ \quad \quad \quad \forall d. \\ \quad \quad \quad \quad \operatorname{pos\text{-}rr}(d) \Rightarrow  \\ \quad \quad \quad \quad \quad \exists j \in \mathbb{N}.\; \forall n \in \mathbb{N}.\; ((j \leq n) \Rightarrow (f(h(n)) < d)) \end{array}}}
\vnbprop{srt-sort-finite}{\vnbstmt{\begin{array}{@{}l@{}} \forall p, f \in \mathbf{Set}. \\ \quad ((|p| \in \mathbb{N}) \wedge \forall x \in p.\; (f(x) \in \mathbb{R})) \Rightarrow  \\ \quad \quad \exists h \in \mathrm{Bij}(\mathrm{OS}(|p|), p). \forall n \in \mathbb{N}. \\ \quad \quad \quad ((n + 1) < |p|) \Rightarrow (f(h((n + 1))) \leq f(h(n))) \end{array}}}
\vnbprop{srt-superlevel-sort}{\vnbstmt{\begin{array}{@{}l@{}} \forall s \in \mathbf{Set}, f \in (s \to \mathbb{R}). \\ \quad \forall d.\; (\operatorname{pos\text{-}rr}(d) \Rightarrow (|\{x \in s : (d \leq f(x))\}| \in \mathbb{N})) \Rightarrow  \\ \quad \quad |\{x \in s : (0 < f(x))\}| \in \mathbb{N} \wedge \\ \quad \quad \exists h \in \mathrm{Bij}(\mathrm{OS}(|\{x \in s : (0 < f(x))\}|), \{x \in s : (0 < f(x))\}). \forall n \in \mathbb{N}. \\ \quad \quad \quad ((n + 1) < |\{x \in s : (0 < f(x))\}|) \Rightarrow (f(h((n + 1))) \leq f(h(n))) \vee \\ \quad \quad \neg (|\{x \in s : (0 < f(x))\}| \in \mathbb{N}) \wedge \\ \quad \quad \exists h \in \mathrm{Bij}(\mathbb{N}, \{x \in s : (0 < f(x))\}). \\ \quad \quad \quad \forall n \in \mathbb{N}.\; (f(h((n + 1))) \leq f(h(n))) \wedge \\ \quad \quad \quad \forall d. \\ \quad \quad \quad \quad \operatorname{pos\text{-}rr}(d) \Rightarrow  \\ \quad \quad \quad \quad \quad \exists j \in \mathbb{N}.\; \forall n \in \mathbb{N}.\; ((j \leq n) \Rightarrow (f(h(n)) < d)) \end{array}}}

\vnbbookfile{theorem-library/interval-widen}{theorem-library/interval-widen}
\noindent{\small interval-widen.scm -- an interval membership survives widening the top end.}\par
\vnbprop{interval-widen}{\vnbstmt{\begin{array}{@{}l@{}} \forall b, c \in \mathbb{N}, a.; \forall i \in \operatorname{interval}(a, b). \\ \quad (b \leq c) \Rightarrow (i \in \operatorname{interval}(a, c)) \end{array}}}

\vnbbookfile{theorem-library/cartesian-nth}{theorem-library/cartesian-nth}
\noindent{\small cartesian-nth.scm -- lemmas MOVED 2026-10-27 (integration of the 10-26/27 wave) out of theorem-library/product-measure.scm, so that files of library core may cite them (the library gate). The proofs are the originals, verbatim; the helper defines are a copy of product-measure.scm's.  Loaded immediately before ...}\par
\vnbprop{cartesian-membership-nth}{\vnbstmt{\begin{array}{@{}l@{}} \forall a, b, w. \\ \quad (w \in (a \times b)) \Leftrightarrow ((w = {[}\operatorname{nth}(1, w), \operatorname{nth}(2, w)]) \wedge ((\operatorname{nth}(1, w) \in a) \wedge (\operatorname{nth}(2, w) \in b))) \end{array}}}

\vnbbookfile{theorem-library/rr-zero-mul}{theorem-library/rr-zero-mul}
\noindent{\small rr-zero-mul.scm -- lemmas MOVED 2026-10-27 (integration of the 10-26/27 wave) out of theorem-library/taylor-proof.scm, so that files of library core may cite them (the library gate). The proofs are the originals, verbatim; the helper defines are a copy of taylor-proof.scm's.  Loaded immediately before taylor-proof.scm.}\par
\vnbprop{rr-zero-mul}{\vnbstmt{\forall v \in \mathbb{R}.\; ((0 \cdot v) = 0)}}

\vnbbookfile{theorem-library/rr-mul-le-cancel}{theorem-library/rr-mul-le-cancel}
\noindent{\small rr-mul-le-cancel.scm -- lemmas MOVED 2026-10-27 (integration of the 10-26/27 wave) out of theorem-library/primitive-uniform-limit.scm, so that files of library core may cite them (the library gate). The proofs are the originals, verbatim; the helper defines are a copy of primitive-uniform-limit.scm's.  Loaded ...}\par
\vnbprop{rr-mul-le-cancel-pos}{\vnbstmt{\forall a, c, m \in \mathbb{R}.\; (((0 < c) \wedge ((a \cdot c) \leq (m \cdot c))) \Rightarrow (a \leq m))}}

\vnbbookfile{theorem-library/rake-finsum-typing-nn}{theorem-library/rake-finsum-typing-nn}
\noindent{\small rake-finsum-typing-nn.scm -- two NN typing leaves of the rake, PROVEN.}\par
\vnbprop{nn-minus-in-nn}{\vnbstmt{\forall k, l \in \mathbb{N}.\; (\operatorname{nn\text{-}minus}(k, l) \in \mathbb{N})}}
\vnbprop{falling-in-nn-ind}{\vnbstmt{\begin{array}{@{}l@{}} \forall m, n \in \mathbb{N}. \\ \quad (\operatorname{falling}(n, m) \in \mathbb{N}) \wedge (((n - m) \in \mathbb{N}) \vee (\operatorname{falling}(n, m) = 0)) \end{array}}}
\vnbprop{falling-in-nn}{\vnbstmt{\forall n, m \in \mathbb{N}.\; (\operatorname{falling}(n, m) \in \mathbb{N})}}

\vnbbookfile{theorem-library/rake-asm-2}{theorem-library/rake-asm-2}
\noindent{\small A file shared across books; the entries that belong to this one.}\par
\noindent{\small rake-asm-2.scm -- RAKE-ASM agent RA-2 (2026-10-12): elementary laws that stood as ASSUMPTIONS (add-axiom! / support), each proven here under its EXACT name and statement (copied character for character from its definition site).}\par
\vnbprop{falling-zero-succ}{\vnbstmt{\forall k \in \mathbb{N}.\; (\operatorname{falling}(0, (k + 1)) = 0)}}
\vnbprop{falling-succ-succ}{\vnbstmt{\forall k, n \in \mathbb{N}.\; (\operatorname{falling}((n + 1), (k + 1)) = ((n + 1) \cdot \operatorname{falling}(n, k)))}}
\vnbprop{choose-times-factorial-all}{\vnbstmt{\forall n, m \in \mathbb{N}.\; ((\operatorname{choose}(n, m) \cdot \operatorname{factorial}(m)) = \operatorname{falling}(n, m))}}
\vnbprop{choose-times-factorial}{\vnbstmt{\begin{array}{@{}l@{}} \forall n, m \in \mathbb{N}. \\ \quad (m \leq n) \Rightarrow ((\operatorname{choose}(n, m) \cdot \operatorname{factorial}(m)) = \operatorname{falling}(n, m)) \end{array}}}

\vnbbookfile{theorem-library/rake-setoid2}{theorem-library/rake-setoid2}
\noindent{\small A file shared across books; the entries that belong to this one.}\par
\noindent{\small rake-setoid2.scm -- BATCH O of the 2026-09-17 rake: the SETOID quotient and the METRIC COMPLETION.  Twelve of the thirteen asserted leaves of the batch are PROVEN here; the thirteenth is a STATEMENT DEFECT and is recorded below. Every statement is its support's statement UNCHANGED (copied from the}\par
\vnbprop{rko2-dist-seq-at}{\vnbstmt{\forall m, f, g.; \forall n \in \mathbb{N}.\; (\operatorname{dist\text{-}seq}(m, f, g)(n) \simeq \operatorname{dist}(m)(f(n), g(n)))}}

\vnbbookfile{theorem-library/singleton-is-set}{theorem-library/singleton-is-set}
\noindent{\small singleton-is-set.scm -- lemmas MOVED 2026-10-27 (integration of the 10-26/27 wave) out of theorem-library/ca-lft-plane.scm, so that files of library core may cite them (the library gate). The proofs are the originals, verbatim; the helper defines are a copy of ca-lft-plane.scm's.  Loaded immediately before ...}\par
\vnbprop{singleton-is-set}{\vnbstmt{\forall y \in \mathbf{Set}.\; (\operatorname{singleton}(y) \in \mathbf{Set})}}

\vnbbookfile{theorem-library/rake-asm-1}{theorem-library/rake-asm-1}
\noindent{\small A file shared across books; the entries that belong to this one.}\par
\noindent{\small rake-asm-1.scm -- RAKE-ASM agent RA-1 (2026-10-12): elementary assumptions PROVEN.}\par
\vnbprop{injection-in-fun}{\vnbstmt{\forall x, y.; \forall i \in \operatorname{injection}(x, y).\; (i \in (x \to y))}}
\vnbprop{injection-injective}{\vnbstmt{\forall x, y.; \forall i \in \operatorname{injection}(x, y), a, b \in x.\; ((i(a) = i(b)) \Rightarrow (a = b))}}
\vnbprop{injection-set-iff}{\vnbstmt{\forall x, y.\; (((x \in \mathbf{Set}) \wedge (y \in \mathbf{Set})) \Rightarrow (\operatorname{injection}(x, y) \in \mathbf{Set}))}}
\vnbprop{bijection-set-iff}{\vnbstmt{\forall x, y.\; (((x \in \mathbf{Set}) \wedge (y \in \mathbf{Set})) \Rightarrow (\mathrm{Bij}(x, y) \in \mathbf{Set}))}}

\vnbbookfile{theorem-library/rake-asm-b5}{theorem-library/rake-asm-b5}
\noindent{\small A file shared across books; the entries that belong to this one.}\par
\noindent{\small rake-asm-b5.scm -- RAKE-ASM-B, agent RB-5 (2026-10-12).  Helper prefix `rb5-'.}\par
\vnbprop{rb5-nary-plus-3-cc}{\vnbstmt{\forall x, y, z \in \mathbb{C}.\; ((x + y + z) = ((x + y) + z))}}
\vnbprop{rb5-nary-plus-4-cc}{\vnbstmt{\begin{array}{@{}l@{}} \forall w, x, y, z \in \mathbb{C}. \\ \quad (w + x + y + z) = (((w + x) + y) + z) \end{array}}}
\vnbprop{rb5-nary-plus-5-cc}{\vnbstmt{\begin{array}{@{}l@{}} \forall v, w, x, y, z \in \mathbb{C}. \\ \quad (v + w + x + y + z) = ((((v + w) + x) + y) + z) \end{array}}}
\vnbprop{rb5-nary-times-3-cc}{\vnbstmt{\forall x, y, z \in \mathbb{C}.\; ((x \cdot y \cdot z) = ((x \cdot y) \cdot z))}}
\vnbprop{rb5-nary-times-4-cc}{\vnbstmt{\begin{array}{@{}l@{}} \forall w, x, y, z \in \mathbb{C}. \\ \quad (w \cdot x \cdot y \cdot z) = (((w \cdot x) \cdot y) \cdot z) \end{array}}}
\vnbprop{rb5-nary-times-5-cc}{\vnbstmt{\begin{array}{@{}l@{}} \forall v, w, x, y, z \in \mathbb{C}. \\ \quad (v \cdot w \cdot x \cdot y \cdot z) = ((((v \cdot w) \cdot x) \cdot y) \cdot z) \end{array}}}
\vnbprop{rb5-nary-plus-2-list-cc}{\vnbstmt{\begin{array}{@{}l@{}} \forall x, y \in \mathbb{C}. \\ \quad (x + y) = \operatorname{reduce}((\lambda {[}\operatorname{rbp\_}, \operatorname{rbq\_}] \in (\mathbb{C} \times \mathbb{C}).\; (\operatorname{rbp\_} + \operatorname{rbq\_})), \operatorname{fam\text{-}of\text{-}list}({[}x, y]), 2) \end{array}}}
\vnbprop{rb5-nary-plus-3-list-cc}{\vnbstmt{\begin{array}{@{}l@{}} \forall x, y, z \in \mathbb{C}. \\ \quad x + y + z \\ \quad =\; \operatorname{reduce}((\lambda {[}\operatorname{rbp\_}, \operatorname{rbq\_}] \in (\mathbb{C} \times \mathbb{C}).\; (\operatorname{rbp\_} + \operatorname{rbq\_})), \operatorname{fam\text{-}of\text{-}list}({[}x, y, z]), 3) \end{array}}}
\vnbprop{rb5-nary-plus-4-list-cc}{\vnbstmt{\begin{array}{@{}l@{}} \forall w, x, y, z \in \mathbb{C}. \\ \quad w + x + y + z \\ \quad =\; \operatorname{reduce}((\lambda {[}\operatorname{rbp\_}, \operatorname{rbq\_}] \in (\mathbb{C} \times \mathbb{C}).\; (\operatorname{rbp\_} + \operatorname{rbq\_})), \operatorname{fam\text{-}of\text{-}list}({[}w, x, y, z]), 4) \end{array}}}
\vnbprop{rb5-nary-plus-5-list-cc}{\vnbstmt{\begin{array}{@{}l@{}} \forall v, w, x, y, z \in \mathbb{C}. \\ \quad v + w + x + y + z \\ \quad =\; \operatorname{reduce}((\lambda {[}\operatorname{rbp\_}, \operatorname{rbq\_}] \in (\mathbb{C} \times \mathbb{C}).\; (\operatorname{rbp\_} + \operatorname{rbq\_})), \operatorname{fam\text{-}of\text{-}list}({[}v, w, x, y, z]), 5) \end{array}}}
\vnbprop{rb5-nary-times-2-list-cc}{\vnbstmt{\begin{array}{@{}l@{}} \forall x, y \in \mathbb{C}. \\ \quad (x \cdot y) = \operatorname{reduce}((\lambda {[}\operatorname{rbp\_}, \operatorname{rbq\_}] \in (\mathbb{C} \times \mathbb{C}).\; (\operatorname{rbp\_} \cdot \operatorname{rbq\_})), \operatorname{fam\text{-}of\text{-}list}({[}x, y]), 2) \end{array}}}
\vnbprop{rb5-nary-times-3-list-cc}{\vnbstmt{\begin{array}{@{}l@{}} \forall x, y, z \in \mathbb{C}. \\ \quad x \cdot y \cdot z \\ \quad =\; \operatorname{reduce}((\lambda {[}\operatorname{rbp\_}, \operatorname{rbq\_}] \in (\mathbb{C} \times \mathbb{C}).\; (\operatorname{rbp\_} \cdot \operatorname{rbq\_})), \operatorname{fam\text{-}of\text{-}list}({[}x, y, z]), 3) \end{array}}}
\vnbprop{rb5-nary-times-4-list-cc}{\vnbstmt{\begin{array}{@{}l@{}} \forall w, x, y, z \in \mathbb{C}. \\ \quad w \cdot x \cdot y \cdot z \\ \quad =\; \operatorname{reduce}((\lambda {[}\operatorname{rbp\_}, \operatorname{rbq\_}] \in (\mathbb{C} \times \mathbb{C}).\; (\operatorname{rbp\_} \cdot \operatorname{rbq\_})), \operatorname{fam\text{-}of\text{-}list}({[}w, x, y, z]), 4) \end{array}}}
\vnbprop{rb5-nary-times-5-list-cc}{\vnbstmt{\begin{array}{@{}l@{}} \forall v, w, x, y, z \in \mathbb{C}. \\ \quad v \cdot w \cdot x \cdot y \cdot z \\ \quad =\; \operatorname{reduce}((\lambda {[}\operatorname{rbp\_}, \operatorname{rbq\_}] \in (\mathbb{C} \times \mathbb{C}).\; (\operatorname{rbp\_} \cdot \operatorname{rbq\_})), \operatorname{fam\text{-}of\text{-}list}({[}v, w, x, y, z]), 5) \end{array}}}
\vnbprop{rb5-fibre-count}{\vnbstmt{\begin{array}{@{}l@{}} \forall p.; \forall m \in \mathbb{N}, g, t \in \mathbf{Set}. \\ \quad ((|g| \in \mathbb{N}) \wedge (\forall y \in t.\; (p(y) \in g) \wedge \forall c \in g.\; (|\{x \in t : (p(x) = c)\}| = m))) \Rightarrow  \\ \quad \quad |t| = (m \cdot |g|) \end{array}}}
\vnbprop{rb5-restrict-injection}{\vnbstmt{\begin{array}{@{}l@{}} \forall a, d, c, f. \\ \quad ((a \in \mathbf{Set}) \wedge ((f \in \operatorname{injection}(d, c)) \wedge (a \subseteq d))) \Rightarrow  \\ \quad \quad (\lambda x \in a.\; f(x)) \in \operatorname{injection}(a, c) \end{array}}}
\vnbprop{rb5-ext-at-b}{\vnbstmt{\begin{array}{@{}l@{}} \forall a, c, b, g, v. \\ \quad ((b \in \mathbf{Set}) \wedge (v \in c)) \Rightarrow  \\ \quad \quad (\lambda z \in (a \cup \{b\}).\; \left(\text{if }(z = b)\text{ then }v\text{ else }g(z)\right))(b) = v \end{array}}}
\vnbprop{rb5-ext-off-b}{\vnbstmt{\begin{array}{@{}l@{}} \forall a, c, b, g, v.; \forall w \in a. \\ \quad ((b \in \mathbf{Set}) \wedge (\neg (b \in a) \wedge (g \in (a \to c)))) \Rightarrow  \\ \quad \quad (\lambda z \in (a \cup \{b\}).\; \left(\text{if }(z = b)\text{ then }v\text{ else }g(z)\right))(w) = g(w) \end{array}}}
\vnbprop{rb5-ext-injection}{\vnbstmt{\begin{array}{@{}l@{}} \forall a, c, b, g, v. \\ \quad ((a \in \mathbf{Set}) \wedge ((b \in \mathbf{Set}) \wedge (\neg (b \in a) \wedge ((g \in \operatorname{injection}(a, c)) \wedge ((v \in c) \wedge \neg (v \in \operatorname{image}(g, a))))))) \Rightarrow  \\ \quad \quad (\lambda z \in (a \cup \{b\}).\; \left(\text{if }(z = b)\text{ then }v\text{ else }g(z)\right)) \in \operatorname{injection}((a \cup \{b\}), c) \end{array}}}
\vnbprop{rb5-ext-restrict}{\vnbstmt{\begin{array}{@{}l@{}} \forall a, c, b, g, v. \\ \quad ((a \in \mathbf{Set}) \wedge ((b \in \mathbf{Set}) \wedge (\neg (b \in a) \wedge ((g \in \operatorname{injection}(a, c)) \wedge ((v \in c) \wedge \neg (v \in \operatorname{image}(g, a))))))) \Rightarrow  \\ \quad \quad (\lambda x \in a.\; (\lambda z \in (a \cup \{b\}).\; \left(\text{if }(z = b)\text{ then }v\text{ else }g(z)\right))(x)) = g \end{array}}}
\vnbprop{rb5-free-value}{\vnbstmt{\begin{array}{@{}l@{}} \forall a, c, b, g, f. \\ \quad a \in \mathbf{Set} \wedge \\ \quad b \in \mathbf{Set} \wedge \\ \quad \neg (b \in a) \wedge \\ \quad f \in \operatorname{injection}((a \cup \{b\}), c) \wedge \\ \quad (\lambda x \in a.\; f(x)) = g \Rightarrow \\ \quad \quad (f(b) \in c) \wedge \neg (f(b) \in \operatorname{image}(g, a)) \end{array}}}
\vnbprop{rb5-ext-recover}{\vnbstmt{\begin{array}{@{}l@{}} \forall a, c, b, g, f. \\ \quad a \in \mathbf{Set} \wedge \\ \quad b \in \mathbf{Set} \wedge \\ \quad \neg (b \in a) \wedge \\ \quad g \in \operatorname{injection}(a, c) \wedge \\ \quad f \in \operatorname{injection}((a \cup \{b\}), c) \wedge \\ \quad (\lambda x \in a.\; f(x)) = g \Rightarrow \\ \quad \quad f = (\lambda z \in (a \cup \{b\}).\; \left(\text{if }(z = b)\text{ then }f(b)\text{ else }g(z)\right)) \end{array}}}
\vnbprop{rb5-complement-card}{\vnbstmt{\begin{array}{@{}l@{}} \forall a, c, g, m. \\ \quad a \in \mathbf{Set} \wedge \\ \quad c \in \mathbf{Set} \wedge \\ \quad |a| \in \mathbb{N} \wedge \\ \quad m \in \mathbb{N} \wedge \\ \quad |c| = (|a| + m) \wedge \\ \quad g \in \operatorname{injection}(a, c) \Rightarrow \\ \quad \quad |\operatorname{difference}(c, \operatorname{image}(g, a))| = m \end{array}}}
\vnbprop{rb5-fibre-image}{\vnbstmt{\begin{array}{@{}l@{}} \forall a, c, b, g. \\ \quad ((a \in \mathbf{Set}) \wedge ((c \in \mathbf{Set}) \wedge ((b \in \mathbf{Set}) \wedge (\neg (b \in a) \wedge (g \in \operatorname{injection}(a, c)))))) \Rightarrow  \\ \quad \quad \operatorname{image}((\lambda \mathit{rbc} \in \operatorname{difference}(c, \operatorname{image}(g, a)).\; \\ \quad \quad \quad \quad (\lambda z \in (a \cup \{b\}).\; \left(\text{if }(z = b)\text{ then }\mathit{rbc}\text{ else }g(z)\right))), \\ \quad \quad \quad \operatorname{difference}(c, \operatorname{image}(g, a))) \\ \quad \quad =\; \{x \in \operatorname{injection}((a \cup \{b\}), c) : ((\lambda f \in \operatorname{injection}((a \cup \{b\}), c).\; (\lambda \mathit{rsx} \in a.\; f(\mathit{rsx})))(x) = g)\} \end{array}}}
\vnbprop{rb5-fibre-card}{\vnbstmt{\begin{array}{@{}l@{}} \forall a, c, b, m, g. \\ \quad a \in \mathbf{Set} \wedge \\ \quad c \in \mathbf{Set} \wedge \\ \quad b \in \mathbf{Set} \wedge \\ \quad \neg (b \in a) \wedge \\ \quad |a| \in \mathbb{N} \wedge \\ \quad m \in \mathbb{N} \wedge \\ \quad |c| = (|a| + m) \wedge \\ \quad g \in \operatorname{injection}(a, c) \Rightarrow \\ \quad \quad |\{x \in \operatorname{injection}((a \cup \{b\}), c) : ((\lambda f \in \operatorname{injection}((a \cup \{b\}), c).\; (\lambda \mathit{rsx} \in a.\; f(\mathit{rsx})))(x) = g)\}| \\ \quad \quad =\; m \end{array}}}
\vnbprop{rb5-injection-finite}{\vnbstmt{\begin{array}{@{}l@{}} \forall a, c. \\ \quad ((a \in \mathbf{Set}) \wedge ((|a| \in \mathbb{N}) \wedge ((c \in \mathbf{Set}) \wedge (|c| \in \mathbb{N})))) \Rightarrow  \\ \quad \quad |\operatorname{injection}(a, c)| \in \mathbb{N} \end{array}}}
\vnbprop{injection-extension-recurrence}{\vnbstmt{\begin{array}{@{}l@{}} \forall a, c, b, m. \\ \quad \neg (b \in a) \wedge \\ \quad a \in \mathbf{Set} \wedge \\ \quad c \in \mathbf{Set} \wedge \\ \quad |a| \in \mathbb{N} \wedge \\ \quad m \in \mathbb{N} \wedge \\ \quad |c| = (|a| + m) \wedge \\ \quad b \in \mathbf{Set} \Rightarrow \\ \quad \quad |\operatorname{injection}((a \cup \{b\}), c)| = (m \cdot |\operatorname{injection}(a, c)|) \end{array}}}
\vnbprop{rb5-count-falling-aux}{\vnbstmt{\begin{array}{@{}l@{}} \forall m, n \in \mathbb{N}. \\ \quad (m \leq n) \Rightarrow (|\operatorname{injection}(\mathrm{OS}(m), \mathrm{OS}(n))| = \operatorname{falling}(n, m)) \end{array}}}
\vnbprop{injection-count-falling}{\vnbstmt{\begin{array}{@{}l@{}} \forall n, m \in \mathbb{N}. \\ \quad (m \leq n) \Rightarrow (|\operatorname{injection}(\mathrm{OS}(m), \mathrm{OS}(n))| = \operatorname{falling}(n, m)) \end{array}}}
\vnbprop{permutation-recurrence}{\vnbstmt{\forall n \in \mathbb{N}.\; (|\operatorname{permutations}((n + 1))| = ((n + 1) \cdot |\operatorname{permutations}(n)|))}}

\vnbbookfile{theorem-library/rr-not-lt-ge}{theorem-library/rr-not-lt-ge}
\noindent{\small rr-not-lt-ge.scm -- lemmas MOVED 2026-11-10 (integration of the 11-09/10 waves) out of theorem-library/ccr2-nv.scm into library core, for the library gate: cited by ccr4c-finite-char (algebra); RR order. The proofs are the originals, verbatim; the helper defines are a copy of the source's. rr-not-lt-ge}\par
\vnbprop{rr-not-lt-ge}{\vnbstmt{\forall a, b \in \mathbb{R}.\; (\neg (a < b) \Rightarrow (b \leq a))}}

\vnbbookfile{theorem-library/union-subset-both}{theorem-library/union-subset-both}
\noindent{\small union-subset-both.scm -- lemmas MOVED 2026-11-10 (integration of the 11-09/10 waves) out of theorem-library/sc-primitive.scm into library core, for the library gate: cited by esum-tonelli (measure); sets. The proofs are the originals, verbatim; the helper defines are a copy of the source's. union-subset-both}\par
\vnbprop{union-subset-both}{\vnbstmt{\forall a, b, c.\; (((a \subseteq c) \wedge (b \subseteq c)) \Rightarrow ((a \cup b) \subseteq c))}}

\part{Algebra}
\noindent Semigroups, monoids, groups, rings, integral domains and fields; ideals, Euclidean rings and Bezout; polynomials and monoid algebras; setoids, quotients and ringoids; homomorphisms, categories and views.
\par\noindent\textit{Builds on:} Foundations.
\par\noindent{\small 122 files.}

\vnbbookfile{structure-library/semigroup}{structure-library/semigroup}
\noindent{\small semigroup.scm -- SEMIGROUP structure}\par
\vnbstructdef{semigroup}{A semigroup $s$ is a 2-tuple whose components are listed below.\par
\begin{enumerate}[nosep]
\item $\operatorname{carr}(s) = \operatorname{nth}(1, s)$, a carrier (\texttt{carriers}): $\operatorname{carr}(s) \in \mathbf{Set}$.
\item $\operatorname{opr}(s) = \operatorname{nth}(2, s)$, an operation (\texttt{op}): $\operatorname{opr}(s) \in ((\operatorname{carr}(s) \times \operatorname{carr}(s)) \to \operatorname{carr}(s))$.
\end{enumerate}
The laws (the \texttt{property} and \texttt{law} clauses):
\begin{itemize}[nosep]
\item $\operatorname{is\text{-}associative}(\operatorname{opr}(s), \operatorname{carr}(s))$
\end{itemize}
The characteristic law (the axiom \texttt{is-semigroup}): $\operatorname{is\text{-}semigroup}(s)$ holds if and only if $\operatorname{length}(s) = 2$, every component has the typing listed with it, every law holds, and $s \in \mathbf{Set}$ (the structure is a set).
}
\par\medskip\noindent\textbf{Definition} (\texttt{hom-semigroup}).\;$\operatorname{hom\text{-}semigroup}(a, b) :=$
\begin{equation*}
\{\operatorname{homf\_} \in (\operatorname{carr}(a) \to \operatorname{carr}(b)) : \operatorname{is\text{-}hom\text{-}semigroup}(a, b, \operatorname{homf\_})\}
\end{equation*}

\vnbdef{is-hom-semigroup-def}{\vnbstmt{\begin{array}{@{}l@{}} \forall a, b, f. \\ \quad \operatorname{is\text{-}hom\text{-}semigroup}(a, b, f) \\ \quad \Leftrightarrow\; \operatorname{is\text{-}semigroup}(a) \wedge \\ \quad \operatorname{is\text{-}semigroup}(b) \wedge \\ \quad f \in (\operatorname{carr}(a) \to \operatorname{carr}(b)) \wedge \\ \quad \forall x_{1}, x_{2} \in \operatorname{carr}(a). \\ \quad \quad f(\operatorname{opr}(a)(x_{1}, x_{2})) = \operatorname{opr}(b)(f(x_{1}), f(x_{2})) \end{array}}}
\vnbdef{semigroup-class}{\vnbstmt{\forall s.\; ((s \in \mathit{semigroup}) \Leftrightarrow \operatorname{is\text{-}semigroup}(s))}}

\vnbbookfile{structure-library/monoid}{structure-library/monoid}
\noindent{\small RETIRED 2026-09-17 (proven): comm-monoid-opr-comm (was an unwarranted axiom, no citers) -- theorem-library/rake-finsum-core.scm RETIRED 2026-09-17 (proven): monoid-identity-in, monoid-carrier-closed-opr -- were UNWARRANTED axioms here; theorem-library/rake-algebra2.scm, projections of is-monoid-def, modulo 0.}\par
\vnbstructdef{monoid}{A monoid $s$ is a 3-tuple whose components are listed below.\par
\begin{enumerate}[nosep]
\item $\operatorname{carr}(s) = \operatorname{nth}(1, s)$, a carrier (\texttt{carriers}): $\operatorname{carr}(s) \in \mathbf{Set}$.
\item $\operatorname{opr}(s) = \operatorname{nth}(2, s)$, an operation (\texttt{op}): $\operatorname{opr}(s) \in ((\operatorname{carr}(s) \times \operatorname{carr}(s)) \to \operatorname{carr}(s))$.
\item $\operatorname{iden}(s) = \operatorname{nth}(3, s)$, a constant (\texttt{constant}): $\operatorname{iden}(s) \in \operatorname{carr}(s)$.
\end{enumerate}
The laws (the \texttt{property} and \texttt{law} clauses):
\begin{itemize}[nosep]
\item $\operatorname{is\text{-}associative}(\operatorname{opr}(s), \operatorname{carr}(s))$
\item $\operatorname{is\text{-}identity}(\operatorname{opr}(s), \operatorname{iden}(s), \operatorname{carr}(s))$
\end{itemize}
The characteristic law (the axiom \texttt{is-monoid}): $\operatorname{is\text{-}monoid}(s)$ holds if and only if $\operatorname{length}(s) = 3$, every component has the typing listed with it, every law holds, and $s \in \mathbf{Set}$ (the structure is a set).
}
\vnbstructdef{comm-monoid}{A comm monoid $s$ is a 3-tuple whose components are listed below.\par
\begin{enumerate}[nosep]
\item $\operatorname{carr}(s) = \operatorname{nth}(1, s)$, a carrier (\texttt{carriers}): $\operatorname{carr}(s) \in \mathbf{Set}$.
\item $\operatorname{opr}(s) = \operatorname{nth}(2, s)$, an operation (\texttt{op}): $\operatorname{opr}(s) \in ((\operatorname{carr}(s) \times \operatorname{carr}(s)) \to \operatorname{carr}(s))$.
\item $\operatorname{iden}(s) = \operatorname{nth}(3, s)$, a constant (\texttt{constant}): $\operatorname{iden}(s) \in \operatorname{carr}(s)$.
\end{enumerate}
The laws (the \texttt{property} and \texttt{law} clauses):
\begin{itemize}[nosep]
\item $\operatorname{is\text{-}associative}(\operatorname{opr}(s), \operatorname{carr}(s))$
\item $\operatorname{is\text{-}identity}(\operatorname{opr}(s), \operatorname{iden}(s), \operatorname{carr}(s))$
\item $\operatorname{is\text{-}commutative}(\operatorname{opr}(s), \operatorname{carr}(s))$
\end{itemize}
The characteristic law (the axiom \texttt{is-comm-monoid}): $\operatorname{is\text{-}comm\text{-}monoid}(s)$ holds if and only if $\operatorname{length}(s) = 3$, every component has the typing listed with it, every law holds, and $s \in \mathbf{Set}$ (the structure is a set).
}
\par\medskip\noindent\textbf{Definition} (\texttt{hom-comm-monoid}).\;$\operatorname{hom\text{-}comm\text{-}monoid}(a, b) :=$
\begin{equation*}
\{\operatorname{homf\_} \in (\operatorname{carr}(a) \to \operatorname{carr}(b)) : \operatorname{is\text{-}hom\text{-}comm\text{-}monoid}(a, b, \operatorname{homf\_})\}
\end{equation*}

\par\medskip\noindent\textbf{Definition} (\texttt{hom-monoid}).\;$\operatorname{hom\text{-}monoid}(a, b) :=$ $\{\operatorname{homf\_} \in (\operatorname{carr}(a) \to \operatorname{carr}(b)) : \operatorname{is\text{-}hom\text{-}monoid}(a, b, \operatorname{homf\_})\}$
\vnbdef{is-hom-monoid-def}{\vnbstmt{\begin{array}{@{}l@{}} \forall a, b, f. \\ \quad \operatorname{is\text{-}hom\text{-}monoid}(a, b, f) \\ \quad \Leftrightarrow\; \operatorname{is\text{-}monoid}(a) \wedge \\ \quad \operatorname{is\text{-}monoid}(b) \wedge \\ \quad f \in (\operatorname{carr}(a) \to \operatorname{carr}(b)) \wedge \\ \quad \forall x_{1}, x_{2} \in \operatorname{carr}(a). \\ \quad \quad f(\operatorname{opr}(a)(x_{1}, x_{2})) = \operatorname{opr}(b)(f(x_{1}), f(x_{2})) \wedge \\ \quad f(\operatorname{iden}(a)) = \operatorname{iden}(b) \end{array}}}
\vnbdef{monoid-class}{\vnbstmt{\forall s.\; ((s \in \mathit{monoid}) \Leftrightarrow \operatorname{is\text{-}monoid}(s))}}
\vnbdef{comm-monoid-class}{\vnbstmt{\forall s.\; ((s \in \operatorname{comm\text{-}monoid}) \Leftrightarrow \operatorname{is\text{-}comm\text{-}monoid}(s))}}
\vnbdef{is-hom-comm-monoid-def}{\vnbstmt{\begin{array}{@{}l@{}} \forall a, b, f. \\ \quad \operatorname{is\text{-}hom\text{-}comm\text{-}monoid}(a, b, f) \\ \quad \Leftrightarrow\; \operatorname{is\text{-}comm\text{-}monoid}(a) \wedge \\ \quad \operatorname{is\text{-}comm\text{-}monoid}(b) \wedge \\ \quad f \in (\operatorname{carr}(a) \to \operatorname{carr}(b)) \wedge \\ \quad \forall x_{1}, x_{2} \in \operatorname{carr}(a). \\ \quad \quad f(\operatorname{opr}(a)(x_{1}, x_{2})) = \operatorname{opr}(b)(f(x_{1}), f(x_{2})) \wedge \\ \quad f(\operatorname{iden}(a)) = \operatorname{iden}(b) \end{array}}}

\vnbbookfile{structure-library/group}{structure-library/group}
\noindent{\small group.scm -- GROUP structure}\par
\vnbstructdef{group}{A group $s$ is a 4-tuple whose components are listed below.\par
\begin{enumerate}[nosep]
\item $\operatorname{carr}(s) = \operatorname{nth}(1, s)$, a carrier (\texttt{carriers}): $\operatorname{carr}(s) \in \mathbf{Set}$.
\item $\operatorname{opr}(s) = \operatorname{nth}(2, s)$, an operation (\texttt{op}): $\operatorname{opr}(s) \in ((\operatorname{carr}(s) \times \operatorname{carr}(s)) \to \operatorname{carr}(s))$.
\item $\operatorname{iden}(s) = \operatorname{nth}(3, s)$, a constant (\texttt{constant}): $\operatorname{iden}(s) \in \operatorname{carr}(s)$.
\item $\operatorname{inv}(s) = \operatorname{nth}(4, s)$, an operation (\texttt{op}): $\operatorname{inv}(s) \in (\operatorname{carr}(s) \to \operatorname{carr}(s))$.
\end{enumerate}
The laws (the \texttt{property} and \texttt{law} clauses):
\begin{itemize}[nosep]
\item $\operatorname{is\text{-}associative}(\operatorname{opr}(s), \operatorname{carr}(s))$
\item $\operatorname{is\text{-}identity}(\operatorname{opr}(s), \operatorname{iden}(s), \operatorname{carr}(s))$
\item $\operatorname{has\text{-}inverses}(\operatorname{opr}(s), \operatorname{iden}(s), \operatorname{inv}(s), \operatorname{carr}(s))$
\end{itemize}
The characteristic law (the axiom \texttt{is-group}): $\operatorname{is\text{-}group}(s)$ holds if and only if $\operatorname{length}(s) = 4$, every component has the typing listed with it, every law holds, and $s \in \mathbf{Set}$ (the structure is a set).
}
\par\medskip\noindent\textbf{Definition} (\texttt{hom-group}).\;$\operatorname{hom\text{-}group}(a, b) :=$ $\{\operatorname{homf\_} \in (\operatorname{carr}(a) \to \operatorname{carr}(b)) : \operatorname{is\text{-}hom\text{-}group}(a, b, \operatorname{homf\_})\}$
\vnbdef{is-hom-group-def}{\vnbstmt{\begin{array}{@{}l@{}} \forall a, b, f. \\ \quad \operatorname{is\text{-}hom\text{-}group}(a, b, f) \\ \quad \Leftrightarrow\; \operatorname{is\text{-}group}(a) \wedge \\ \quad \operatorname{is\text{-}group}(b) \wedge \\ \quad f \in (\operatorname{carr}(a) \to \operatorname{carr}(b)) \wedge \\ \quad \forall x_{1}, x_{2} \in \operatorname{carr}(a). \\ \quad \quad f(\operatorname{opr}(a)(x_{1}, x_{2})) = \operatorname{opr}(b)(f(x_{1}), f(x_{2})) \wedge \\ \quad f(\operatorname{iden}(a)) = \operatorname{iden}(b) \wedge \\ \quad \forall x_{1} \in \operatorname{carr}(a).\; (f(\operatorname{inv}(a)(x_{1})) = \operatorname{inv}(b)(f(x_{1}))) \end{array}}}
\vnbdef{group-class}{\vnbstmt{\forall s.\; ((s \in \mathit{group}) \Leftrightarrow \operatorname{is\text{-}group}(s))}}

\vnbbookfile{structure-library/abelian-group}{structure-library/abelian-group}
\noindent{\small RETIRED 2026-09-17 (proven): abelian-group-assoc, abelian-group-right-id, abelian-group-inverse-unique -- theorem-library/rake-algebra2.scm (which also rebuilds the three MODULE-VECTOR-AG companions that proofs cite by name). abelian-group.scm -- ABELIAN-GROUP structure}\par
\vnbstructdef{abelian-group}{An abelian group $s$ is a 4-tuple whose components are listed below.\par
\begin{enumerate}[nosep]
\item $\operatorname{carr}(s) = \operatorname{nth}(1, s)$, a carrier (\texttt{carriers}): $\operatorname{carr}(s) \in \mathbf{Set}$.
\item $\operatorname{opr}(s) = \operatorname{nth}(2, s)$, an operation (\texttt{op}): $\operatorname{opr}(s) \in ((\operatorname{carr}(s) \times \operatorname{carr}(s)) \to \operatorname{carr}(s))$.
\item $\operatorname{iden}(s) = \operatorname{nth}(3, s)$, a constant (\texttt{constant}): $\operatorname{iden}(s) \in \operatorname{carr}(s)$.
\item $\operatorname{inv}(s) = \operatorname{nth}(4, s)$, an operation (\texttt{op}): $\operatorname{inv}(s) \in (\operatorname{carr}(s) \to \operatorname{carr}(s))$.
\end{enumerate}
The laws (the \texttt{property} and \texttt{law} clauses):
\begin{itemize}[nosep]
\item $\operatorname{is\text{-}associative}(\operatorname{opr}(s), \operatorname{carr}(s))$
\item $\operatorname{is\text{-}identity}(\operatorname{opr}(s), \operatorname{iden}(s), \operatorname{carr}(s))$
\item $\operatorname{has\text{-}inverses}(\operatorname{opr}(s), \operatorname{iden}(s), \operatorname{inv}(s), \operatorname{carr}(s))$
\item $\operatorname{is\text{-}commutative}(\operatorname{opr}(s), \operatorname{carr}(s))$
\end{itemize}
The characteristic law (the axiom \texttt{is-abelian-group}): $\operatorname{is\text{-}abelian\text{-}group}(s)$ holds if and only if $\operatorname{length}(s) = 4$, every component has the typing listed with it, every law holds, and $s \in \mathbf{Set}$ (the structure is a set).
}
\par\medskip\noindent\textbf{Definition} (\texttt{hom-abelian-group}).\;$\operatorname{hom\text{-}abelian\text{-}group}(a, b) :=$
\begin{equation*}
\{\operatorname{homf\_} \in (\operatorname{carr}(a) \to \operatorname{carr}(b)) : \operatorname{is\text{-}hom\text{-}abelian\text{-}group}(a, b, \operatorname{homf\_})\}
\end{equation*}

\vnbdef{is-hom-abelian-group-def}{\vnbstmt{\begin{array}{@{}l@{}} \forall a, b, f. \\ \quad \operatorname{is\text{-}hom\text{-}abelian\text{-}group}(a, b, f) \\ \quad \Leftrightarrow\; \operatorname{is\text{-}abelian\text{-}group}(a) \wedge \\ \quad \operatorname{is\text{-}abelian\text{-}group}(b) \wedge \\ \quad f \in (\operatorname{carr}(a) \to \operatorname{carr}(b)) \wedge \\ \quad \forall x_{1}, x_{2} \in \operatorname{carr}(a). \\ \quad \quad f(\operatorname{opr}(a)(x_{1}, x_{2})) = \operatorname{opr}(b)(f(x_{1}), f(x_{2})) \wedge \\ \quad f(\operatorname{iden}(a)) = \operatorname{iden}(b) \wedge \\ \quad \forall x_{1} \in \operatorname{carr}(a).\; (f(\operatorname{inv}(a)(x_{1})) = \operatorname{inv}(b)(f(x_{1}))) \end{array}}}
\vnbdef{abelian-group-class}{\vnbstmt{\forall s.\; ((s \in \operatorname{abelian\text{-}group}) \Leftrightarrow \operatorname{is\text{-}abelian\text{-}group}(s))}}

\vnbbookfile{structure-library/ring}{structure-library/ring}
\noindent{\small RETIRED 2026-09-17 (proven): ring-neg-mul-left -- theorem-library/rake-algebra2.scm RETIRED 2026-09-17 (proven): ring-carr-in-set -- theorem-library/rake-algebra.scm ring.scm -- RING, RING-PROD, ZERO-RING}\par
\vnbstructdef{ring}{A ring $s$ is a 6-tuple whose components are listed below.\par
\begin{enumerate}[nosep]
\item $\operatorname{carr}(s) = \operatorname{nth}(1, s)$, a carrier (\texttt{carriers}): $\operatorname{carr}(s) \in \mathbf{Set}$.
\item $\operatorname{add}(s) = \operatorname{nth}(2, s)$, an operation (\texttt{op}): $\operatorname{add}(s) \in ((\operatorname{carr}(s) \times \operatorname{carr}(s)) \to \operatorname{carr}(s))$.
\item $\operatorname{mul}(s) = \operatorname{nth}(3, s)$, an operation (\texttt{op}): $\operatorname{mul}(s) \in ((\operatorname{carr}(s) \times \operatorname{carr}(s)) \to \operatorname{carr}(s))$.
\item $\operatorname{neg}(s) = \operatorname{nth}(4, s)$, an operation (\texttt{op}): $\operatorname{neg}(s) \in (\operatorname{carr}(s) \to \operatorname{carr}(s))$.
\item $\operatorname{zero}(s) = \operatorname{nth}(5, s)$, a constant (\texttt{constant}): $\operatorname{zero}(s) \in \operatorname{carr}(s)$.
\item $\operatorname{one}(s) = \operatorname{nth}(6, s)$, a constant (\texttt{constant}): $\operatorname{one}(s) \in \operatorname{carr}(s)$.
\end{enumerate}
The laws (the \texttt{property} and \texttt{law} clauses):
\begin{itemize}[nosep]
\item $\operatorname{is\text{-}associative}(\operatorname{add}(s), \operatorname{carr}(s))$
\item $\operatorname{is\text{-}commutative}(\operatorname{add}(s), \operatorname{carr}(s))$
\item $\operatorname{is\text{-}identity}(\operatorname{add}(s), \operatorname{zero}(s), \operatorname{carr}(s))$
\item $\operatorname{has\text{-}inverses}(\operatorname{add}(s), \operatorname{zero}(s), \operatorname{neg}(s), \operatorname{carr}(s))$
\item $\operatorname{is\text{-}associative}(\operatorname{mul}(s), \operatorname{carr}(s))$
\item $\operatorname{is\text{-}identity}(\operatorname{mul}(s), \operatorname{one}(s), \operatorname{carr}(s))$
\item $\operatorname{is\text{-}distributive}(\operatorname{add}(s), \operatorname{mul}(s), \operatorname{carr}(s))$
\end{itemize}
The characteristic law (the axiom \texttt{is-ring}): $\operatorname{is\text{-}ring}(s)$ holds if and only if $\operatorname{length}(s) = 6$, every component has the typing listed with it, every law holds, and $s \in \mathbf{Set}$ (the structure is a set).
}
\par\medskip\noindent\textbf{Definition} (\texttt{ring-prod}).\;$\operatorname{ring\text{-}prod}(x, y) :=$
\begin{equation*}
{[}(\operatorname{carr}(x) \times \operatorname{carr}(y)), (\lambda {[}p, q] \in ((\operatorname{carr}(x) \times \operatorname{carr}(y)) \times (\operatorname{carr}(x) \times \operatorname{carr}(y))).\; {[}\operatorname{add}(x)(\operatorname{nth}(1, p), \operatorname{nth}(1, q)), \operatorname{add}(y)(\operatorname{nth}(2, p), \operatorname{nth}(2, q))]), (\lambda {[}p, q] \in ((\operatorname{carr}(x) \times \operatorname{carr}(y)) \times (\operatorname{carr}(x) \times \operatorname{carr}(y))).\; {[}\operatorname{mul}(x)(\operatorname{nth}(1, p), \operatorname{nth}(1, q)), \operatorname{mul}(y)(\operatorname{nth}(2, p), \operatorname{nth}(2, q))]), (\lambda p \in (\operatorname{carr}(x) \times \operatorname{carr}(y)).\; {[}\operatorname{neg}(x)(\operatorname{nth}(1, p)), \operatorname{neg}(y)(\operatorname{nth}(2, p))]), {[}\operatorname{zero}(x), \operatorname{zero}(y)], {[}\operatorname{one}(x), \operatorname{one}(y)]]
\end{equation*}

\par\medskip\noindent\textbf{Definition} (\texttt{zero-ring}).\;
\begin{equation*}
\begin{array}{@{}l@{}} \operatorname{zero\text{-}ring} \\ =\; {[}\operatorname{make\text{-}set}({[}0]), (\lambda {[}p, q] \in (\operatorname{make\text{-}set}({[}0]) \times \operatorname{make\text{-}set}({[}0])).\; 0), (\lambda {[}p, q] \in (\operatorname{make\text{-}set}({[}0]) \times \operatorname{make\text{-}set}({[}0])).\; 0), (\lambda p \in \operatorname{make\text{-}set}({[}0]).\; 0), 0, 0] \end{array}
\end{equation*}

\par\medskip\noindent\textbf{Definition} (\texttt{hom-ring}).\;$\operatorname{hom\text{-}ring}(a, b) :=$ $\{\operatorname{homf\_} \in (\operatorname{carr}(a) \to \operatorname{carr}(b)) : \operatorname{is\text{-}hom\text{-}ring}(a, b, \operatorname{homf\_})\}$
\vnbdef{ring-add-assoc}{\vnbstmt{\begin{array}{@{}l@{}} \forall s.; \forall a, b, c \in \operatorname{carr}(s). \\ \quad \operatorname{is\text{-}ring}(s) \Rightarrow (\operatorname{add}(s)(\operatorname{add}(s)(a, b), c) = \operatorname{add}(s)(a, \operatorname{add}(s)(b, c))) \end{array}}}
\vnbdef{ring-add-comm}{\vnbstmt{\begin{array}{@{}l@{}} \forall s.; \forall a, b \in \operatorname{carr}(s). \\ \quad \operatorname{is\text{-}ring}(s) \Rightarrow (\operatorname{add}(s)(a, b) = \operatorname{add}(s)(b, a)) \end{array}}}
\vnbdef{ring-add-left-id}{\vnbstmt{\forall s.; \forall a \in \operatorname{carr}(s).\; (\operatorname{is\text{-}ring}(s) \Rightarrow (\operatorname{add}(s)(\operatorname{zero}(s), a) = a))}}
\vnbdef{ring-add-left-inv}{\vnbstmt{\forall s.; \forall a \in \operatorname{carr}(s).\; (\operatorname{is\text{-}ring}(s) \Rightarrow (\operatorname{add}(s)(\operatorname{neg}(s)(a), a) = \operatorname{zero}(s)))}}
\vnbdef{ring-mul-assoc}{\vnbstmt{\begin{array}{@{}l@{}} \forall s.; \forall a, b, c \in \operatorname{carr}(s). \\ \quad \operatorname{is\text{-}ring}(s) \Rightarrow (\operatorname{mul}(s)(\operatorname{mul}(s)(a, b), c) = \operatorname{mul}(s)(a, \operatorname{mul}(s)(b, c))) \end{array}}}
\vnbdef{ring-mul-left-id}{\vnbstmt{\forall s.; \forall a \in \operatorname{carr}(s).\; (\operatorname{is\text{-}ring}(s) \Rightarrow (\operatorname{mul}(s)(\operatorname{one}(s), a) = a))}}
\vnbdef{ring-mul-right-id}{\vnbstmt{\forall s.; \forall a \in \operatorname{carr}(s).\; (\operatorname{is\text{-}ring}(s) \Rightarrow (\operatorname{mul}(s)(a, \operatorname{one}(s)) = a))}}
\vnbdef{ring-left-dist}{\vnbstmt{\begin{array}{@{}l@{}} \forall s.; \forall a, b, c \in \operatorname{carr}(s). \\ \quad \operatorname{is\text{-}ring}(s) \Rightarrow  \\ \quad \quad \operatorname{mul}(s)(a, \operatorname{add}(s)(b, c)) = \operatorname{add}(s)(\operatorname{mul}(s)(a, b), \operatorname{mul}(s)(a, c)) \end{array}}}
\vnbdef{ring-right-dist}{\vnbstmt{\begin{array}{@{}l@{}} \forall s.; \forall a, b, c \in \operatorname{carr}(s). \\ \quad \operatorname{is\text{-}ring}(s) \Rightarrow  \\ \quad \quad \operatorname{mul}(s)(\operatorname{add}(s)(a, b), c) = \operatorname{add}(s)(\operatorname{mul}(s)(a, c), \operatorname{mul}(s)(b, c)) \end{array}}}
\vnbdef{ring-zero-in}{\vnbstmt{\forall r.\; (\operatorname{is\text{-}ring}(r) \Rightarrow (\operatorname{zero}(r) \in \operatorname{carr}(r)))}}
\vnbdef{ring-carrier-closed-add}{\vnbstmt{\begin{array}{@{}l@{}} \forall r, a, b. \\ \quad (\operatorname{is\text{-}ring}(r) \wedge ((a \in \operatorname{carr}(r)) \wedge (b \in \operatorname{carr}(r)))) \Rightarrow (\operatorname{add}(r)(a, b) \in \operatorname{carr}(r)) \end{array}}}
\vnbdef{is-hom-ring-def}{\vnbstmt{\begin{array}{@{}l@{}} \forall a, b, f. \\ \quad \operatorname{is\text{-}hom\text{-}ring}(a, b, f) \\ \quad \Leftrightarrow\; \operatorname{is\text{-}ring}(a) \wedge \\ \quad \operatorname{is\text{-}ring}(b) \wedge \\ \quad f \in (\operatorname{carr}(a) \to \operatorname{carr}(b)) \wedge \\ \quad \forall x_{1}, x_{2} \in \operatorname{carr}(a). \\ \quad \quad f(\operatorname{add}(a)(x_{1}, x_{2})) = \operatorname{add}(b)(f(x_{1}), f(x_{2})) \wedge \\ \quad \forall x_{1}, x_{2} \in \operatorname{carr}(a). \\ \quad \quad f(\operatorname{mul}(a)(x_{1}, x_{2})) = \operatorname{mul}(b)(f(x_{1}), f(x_{2})) \wedge \\ \quad \forall x_{1} \in \operatorname{carr}(a).\; (f(\operatorname{neg}(a)(x_{1})) = \operatorname{neg}(b)(f(x_{1}))) \wedge \\ \quad f(\operatorname{zero}(a)) = \operatorname{zero}(b) \wedge \\ \quad f(\operatorname{one}(a)) = \operatorname{one}(b) \end{array}}}
\vnbdef{ring-class}{\vnbstmt{\forall s.\; ((s \in \mathit{ring}) \Leftrightarrow \operatorname{is\text{-}ring}(s))}}

\vnbbookfile{structure-library/setoid}{structure-library/setoid}
\noindent{\small RETIRED 2026-09-17 (proven, rake batch O): class-self, class-subset-carrier, class-eq-iff, class-disjoint, class-in-quotient, quotient-rep, descend-computes, quotient-universal, descend2-computes -- theorem-library/rake-setoid2.scm, all modulo 0.  This file now holds no support. RETIRED 2026-09-17 (proven): ...}\par
\vnbstructdef{setoid}{A setoid $s$ is a 2-tuple whose components are listed below.\par
\begin{enumerate}[nosep]
\item $\operatorname{pts}(s) = \operatorname{nth}(1, s)$, a carrier (\texttt{carriers}): $\operatorname{pts}(s) \in \mathbf{Set}$.
\item $\operatorname{rel}(s) = \operatorname{nth}(2, s)$, a relation (\texttt{relation}): $\operatorname{rel}(s) \in \operatorname{power}((\operatorname{pts}(s) \times \operatorname{pts}(s)))$.
\end{enumerate}
The laws (the \texttt{property} and \texttt{law} clauses):
\begin{itemize}[nosep]
\item $\operatorname{is\text{-}equivalence}(\operatorname{rel}(s), \operatorname{pts}(s))$
\end{itemize}
The characteristic law (the axiom \texttt{is-setoid}): $\operatorname{is\text{-}setoid}(s)$ holds if and only if $\operatorname{length}(s) = 2$, every component has the typing listed with it, every law holds, and $s \in \mathbf{Set}$ (the structure is a set).
}
\par\medskip\noindent\textbf{Definition} (\texttt{quotient}).\;$\operatorname{quotient}(s) :=$ $\operatorname{image}(\operatorname{proj}(s), \operatorname{pts}(s))$
\par\medskip\noindent\textbf{Definition} (\texttt{class}).\;$\operatorname{class}(s, a) :=$ $\{b \in \operatorname{pts}(s) : \operatorname{related}(s, a, b)\}$
\par\medskip\noindent\textbf{Definition} (\texttt{related}).\;$\operatorname{related}(s, a, b) :=$ ${[}a, b] \in \operatorname{rel}(s)$
\par\medskip\noindent\textbf{Definition} (\texttt{is-equivalence}).\;For all $\rho$ and $\mathit{crr}$, 
\begin{equation*}
\begin{array}{@{}l@{}} \operatorname{is\text{-}equivalence}(\rho, \mathit{crr}) \\ \Leftrightarrow\; \rho \in \operatorname{power}((\mathit{crr} \times \mathit{crr})) \wedge \\ \forall u \in \mathit{crr}.\; ({[}u, u] \in \rho) \wedge \\ \forall u, v \in \mathit{crr}.\; (({[}u, v] \in \rho) \Rightarrow ({[}v, u] \in \rho)) \wedge \\ \forall u, v, w \in \mathit{crr}. \\ \quad (({[}u, v] \in \rho) \wedge ({[}v, w] \in \rho)) \Rightarrow ({[}u, w] \in \rho) \end{array}
\end{equation*}

\par\medskip\noindent\textbf{Definition} (\texttt{proj}).\;$\operatorname{proj}(s) :=$ $(\lambda a \in \operatorname{pts}(s).\; \operatorname{class}(s, a))$
\par\medskip\noindent\textbf{Definition} (\texttt{respects}).\;$\operatorname{respects}(s, f) :=$
\begin{equation*}
\forall a, b \in \operatorname{pts}(s).\; (\operatorname{related}(s, a, b) \Rightarrow (f(a) = f(b)))
\end{equation*}

\par\medskip\noindent\textbf{Definition} (\texttt{descend}).\;$\operatorname{descend}(s, f) :=$
\begin{equation*}
(\lambda c \in \operatorname{quotient}(s).\; (\iota\, z.\;((z \in \mathbf{Set}) \wedge \exists a \in c.\; (z = f(a)))))
\end{equation*}

\par\medskip\noindent\textbf{Definition} (\texttt{respects2}).\;$\operatorname{respects2}(s, f) :=$
\begin{equation*}
\begin{array}{@{}l@{}} \forall a, b, \operatorname{a\_}, \operatorname{b\_} \in \operatorname{pts}(s). \\ \quad (\operatorname{related}(s, a, \operatorname{a\_}) \wedge \operatorname{related}(s, b, \operatorname{b\_})) \Rightarrow (f(a, b) = f(\operatorname{a\_}, \operatorname{b\_})) \end{array}
\end{equation*}

\par\medskip\noindent\textbf{Definition} (\texttt{descend2}).\;$\operatorname{descend2}(s, f) :=$
\begin{equation*}
\begin{array}{@{}l@{}} (\lambda {[}c, d] \in (\operatorname{quotient}(s) \times \operatorname{quotient}(s)).\; \\ \quad (\iota\, z.\;((z \in \mathbf{Set}) \wedge \exists a \in c, b \in d.\; (z = f(a, b))))) \end{array}
\end{equation*}

\par\medskip\noindent\textbf{Definition} (\texttt{hom-setoid}).\;$\operatorname{hom\text{-}setoid}(a, b) :=$ $\{\operatorname{homf\_} \in (\operatorname{pts}(a) \to \operatorname{pts}(b)) : \operatorname{is\text{-}hom\text{-}setoid}(a, b, \operatorname{homf\_})\}$
\vnbdef{is-hom-setoid-def}{\vnbstmt{\begin{array}{@{}l@{}} \forall a, b, f. \\ \quad \operatorname{is\text{-}hom\text{-}setoid}(a, b, f) \\ \quad \Leftrightarrow\; \operatorname{is\text{-}setoid}(a) \wedge \\ \quad \operatorname{is\text{-}setoid}(b) \wedge \\ \quad f \in (\operatorname{pts}(a) \to \operatorname{pts}(b)) \wedge \\ \quad \forall x_{1}, x_{2} \in \operatorname{pts}(a). \\ \quad \quad ({[}x_{1}, x_{2}] \in \operatorname{rel}(a)) \Rightarrow ({[}f(x_{1}), f(x_{2})] \in \operatorname{rel}(b)) \end{array}}}
\vnbdef{setoid-class}{\vnbstmt{\forall s.\; ((s \in \mathit{setoid}) \Leftrightarrow \operatorname{is\text{-}setoid}(s))}}

\vnbbookfile{structure-library/ringoid}{structure-library/ringoid}
\noindent{\small ringoid.scm -- a ring with a distinguished two-sided ideal.}\par
\vnbstructdef{ringoid}{A ringoid $s$ is a 7-tuple whose components are listed below.\par
\begin{enumerate}[nosep]
\item $\operatorname{carr}(s) = \operatorname{nth}(1, s)$, a carrier (\texttt{carriers}): $\operatorname{carr}(s) \in \mathbf{Set}$.
\item $\operatorname{add}(s) = \operatorname{nth}(2, s)$, an operation (\texttt{op}): $\operatorname{add}(s) \in ((\operatorname{carr}(s) \times \operatorname{carr}(s)) \to \operatorname{carr}(s))$.
\item $\operatorname{mul}(s) = \operatorname{nth}(3, s)$, an operation (\texttt{op}): $\operatorname{mul}(s) \in ((\operatorname{carr}(s) \times \operatorname{carr}(s)) \to \operatorname{carr}(s))$.
\item $\operatorname{neg}(s) = \operatorname{nth}(4, s)$, an operation (\texttt{op}): $\operatorname{neg}(s) \in (\operatorname{carr}(s) \to \operatorname{carr}(s))$.
\item $\operatorname{zero}(s) = \operatorname{nth}(5, s)$, a constant (\texttt{constant}): $\operatorname{zero}(s) \in \operatorname{carr}(s)$.
\item $\operatorname{one}(s) = \operatorname{nth}(6, s)$, a constant (\texttt{constant}): $\operatorname{one}(s) \in \operatorname{carr}(s)$.
\item $\operatorname{idl}(s) = \operatorname{nth}(7, s)$, a constant (\texttt{constant}): $\operatorname{idl}(s) \in \operatorname{power}(\operatorname{carr}(s))$.
\end{enumerate}
The laws (the \texttt{property} and \texttt{law} clauses):
\begin{itemize}[nosep]
\item $\operatorname{is\text{-}associative}(\operatorname{add}(s), \operatorname{carr}(s))$
\item $\operatorname{is\text{-}commutative}(\operatorname{add}(s), \operatorname{carr}(s))$
\item $\operatorname{is\text{-}identity}(\operatorname{add}(s), \operatorname{zero}(s), \operatorname{carr}(s))$
\item $\operatorname{has\text{-}inverses}(\operatorname{add}(s), \operatorname{zero}(s), \operatorname{neg}(s), \operatorname{carr}(s))$
\item $\operatorname{is\text{-}associative}(\operatorname{mul}(s), \operatorname{carr}(s))$
\item $\operatorname{is\text{-}identity}(\operatorname{mul}(s), \operatorname{one}(s), \operatorname{carr}(s))$
\item $\operatorname{is\text{-}distributive}(\operatorname{add}(s), \operatorname{mul}(s), \operatorname{carr}(s))$
\item $\operatorname{is\text{-}ideal\text{-}in}(\operatorname{idl}(s), \operatorname{carr}(s), \operatorname{add}(s), \operatorname{mul}(s), \operatorname{neg}(s), \operatorname{zero}(s))$
\end{itemize}
The characteristic law (the axiom \texttt{is-ringoid}): $\operatorname{is\text{-}ringoid}(s)$ holds if and only if $\operatorname{length}(s) = 7$, every component has the typing listed with it, every law holds, and $s \in \mathbf{Set}$ (the structure is a set).
}
\par\medskip\noindent\textbf{Definition} (\texttt{ringoid-quotient}).\;$\operatorname{ringoid\text{-}quotient}(r) :=$ $\operatorname{quotient}(\operatorname{ringoid\text{-}setoid}(r))$
\par\medskip\noindent\textbf{Definition} (\texttt{is-ideal-in}).\;For all $\operatorname{i\_}$, $\operatorname{carr\_}$, $\operatorname{add\_}$, $\operatorname{mul\_}$, $\operatorname{neg\_}$ and $\operatorname{zero\_}$, 
\begin{equation*}
\begin{array}{@{}l@{}} \operatorname{is\text{-}ideal\text{-}in}(\operatorname{i\_}, \operatorname{carr\_}, \operatorname{add\_}, \operatorname{mul\_}, \operatorname{neg\_}, \operatorname{zero\_}) \\ \Leftrightarrow\; \operatorname{i\_} \subseteq \operatorname{carr\_} \wedge \\ \operatorname{zero\_} \in \operatorname{i\_} \wedge \\ \forall a, b \in \operatorname{i\_}.\; (\operatorname{add\_}(a, b) \in \operatorname{i\_}) \wedge \\ \forall a \in \operatorname{i\_}.\; (\operatorname{neg\_}(a) \in \operatorname{i\_}) \wedge \\ \forall \operatorname{x\_} \in \operatorname{carr\_}, a \in \operatorname{i\_}.\; (\operatorname{mul\_}(\operatorname{x\_}, a) \in \operatorname{i\_}) \wedge \\ \forall \operatorname{x\_} \in \operatorname{carr\_}, a \in \operatorname{i\_}.\; (\operatorname{mul\_}(a, \operatorname{x\_}) \in \operatorname{i\_}) \end{array}
\end{equation*}

\par\medskip\noindent\textbf{Definition} (\texttt{ringoid-rel}).\;$\operatorname{ringoid\text{-}rel}(r) :=$
\begin{equation*}
\{p \in (\operatorname{carr}(r) \times \operatorname{carr}(r)) : (\operatorname{add}(r)(\operatorname{nth}(1, p), \operatorname{neg}(r)(\operatorname{nth}(2, p))) \in \operatorname{idl}(r))\}
\end{equation*}

\par\medskip\noindent\textbf{Definition} (\texttt{ringoid-setoid}).\;$\operatorname{ringoid\text{-}setoid}(r) :=$ ${[}\operatorname{carr}(r), \operatorname{ringoid\text{-}rel}(r)]$
\par\medskip\noindent\textbf{Definition} (\texttt{ringoid-quotient-ring}).\;$\operatorname{ringoid\text{-}quotient\text{-}ring}(r) :=$
\begin{equation*}
{[}\operatorname{quotient}(\operatorname{ringoid\text{-}setoid}(r)), \operatorname{descend2}(\operatorname{ringoid\text{-}setoid}(r), (\lambda {[}a, b] \in (\operatorname{carr}(r) \times \operatorname{carr}(r)).\; \operatorname{class}(\operatorname{ringoid\text{-}setoid}(r), \operatorname{add}(r)(a, b)))), \operatorname{descend2}(\operatorname{ringoid\text{-}setoid}(r), (\lambda {[}a, b] \in (\operatorname{carr}(r) \times \operatorname{carr}(r)).\; \operatorname{class}(\operatorname{ringoid\text{-}setoid}(r), \operatorname{mul}(r)(a, b)))), \operatorname{descend}(\operatorname{ringoid\text{-}setoid}(r), (\lambda a \in \operatorname{carr}(r).\; \operatorname{class}(\operatorname{ringoid\text{-}setoid}(r), \operatorname{neg}(r)(a)))), \operatorname{class}(\operatorname{ringoid\text{-}setoid}(r), \operatorname{zero}(r)), \operatorname{class}(\operatorname{ringoid\text{-}setoid}(r), \operatorname{one}(r))]
\end{equation*}

\par\medskip\noindent\textbf{Definition} (\texttt{hom-ringoid}).\;$\operatorname{hom\text{-}ringoid}(a, b) :=$
\begin{equation*}
\{\operatorname{homf\_} \in (\operatorname{carr}(a) \to \operatorname{carr}(b)) : \operatorname{is\text{-}hom\text{-}ringoid}(a, b, \operatorname{homf\_})\}
\end{equation*}

\vnbdef{rq-carr}{\vnbstmt{\begin{array}{@{}l@{}} \forall r. \\ \quad \operatorname{is\text{-}ringoid}(r) \Rightarrow  \\ \quad \quad (\operatorname{carr}(\operatorname{ringoid\text{-}quotient\text{-}ring}(r)) \simeq \operatorname{quotient}(\operatorname{ringoid\text{-}setoid}(r))) \end{array}}}
\vnbdef{rq-add}{\vnbstmt{\begin{array}{@{}l@{}} \forall r. \\ \quad \operatorname{is\text{-}ringoid}(r) \Rightarrow  \\ \quad \quad (\operatorname{add}(\operatorname{ringoid\text{-}quotient\text{-}ring}(r)) \simeq \operatorname{descend2}(\operatorname{ringoid\text{-}setoid}(r), (\lambda {[}a, b] \in (\operatorname{carr}(r) \times \operatorname{carr}(r)).\; \operatorname{class}(\operatorname{ringoid\text{-}setoid}(r), \operatorname{add}(r)(a, b))))) \end{array}}}
\vnbdef{rq-mul}{\vnbstmt{\begin{array}{@{}l@{}} \forall r. \\ \quad \operatorname{is\text{-}ringoid}(r) \Rightarrow  \\ \quad \quad (\operatorname{mul}(\operatorname{ringoid\text{-}quotient\text{-}ring}(r)) \simeq \operatorname{descend2}(\operatorname{ringoid\text{-}setoid}(r), (\lambda {[}a, b] \in (\operatorname{carr}(r) \times \operatorname{carr}(r)).\; \operatorname{class}(\operatorname{ringoid\text{-}setoid}(r), \operatorname{mul}(r)(a, b))))) \end{array}}}
\vnbdef{rq-neg}{\vnbstmt{\begin{array}{@{}l@{}} \forall r. \\ \quad \operatorname{is\text{-}ringoid}(r) \Rightarrow  \\ \quad \quad (\operatorname{neg}(\operatorname{ringoid\text{-}quotient\text{-}ring}(r)) \simeq \operatorname{descend}(\operatorname{ringoid\text{-}setoid}(r), (\lambda a \in \operatorname{carr}(r).\; \operatorname{class}(\operatorname{ringoid\text{-}setoid}(r), \operatorname{neg}(r)(a))))) \end{array}}}
\vnbdef{rq-zero}{\vnbstmt{\begin{array}{@{}l@{}} \forall r. \\ \quad \operatorname{is\text{-}ringoid}(r) \Rightarrow  \\ \quad \quad (\operatorname{zero}(\operatorname{ringoid\text{-}quotient\text{-}ring}(r)) \simeq \operatorname{class}(\operatorname{ringoid\text{-}setoid}(r), \operatorname{zero}(r))) \end{array}}}
\vnbdef{rq-one}{\vnbstmt{\begin{array}{@{}l@{}} \forall r. \\ \quad \operatorname{is\text{-}ringoid}(r) \Rightarrow  \\ \quad \quad (\operatorname{one}(\operatorname{ringoid\text{-}quotient\text{-}ring}(r)) \simeq \operatorname{class}(\operatorname{ringoid\text{-}setoid}(r), \operatorname{one}(r))) \end{array}}}
\vnbdef{ringoid-setoid-pts}{\vnbstmt{\forall r.\; (\operatorname{is\text{-}ringoid}(r) \Rightarrow (\operatorname{pts}(\operatorname{ringoid\text{-}setoid}(r)) \simeq \operatorname{carr}(r)))}}
\vnbdef{ringoid-rel-mem}{\vnbstmt{\begin{array}{@{}l@{}} \forall r, a, b. \\ \quad {[}a, b] \in \operatorname{ringoid\text{-}rel}(r) \\ \quad \Leftrightarrow\; ({[}a, b] \in (\operatorname{carr}(r) \times \operatorname{carr}(r))) \wedge (\operatorname{add}(r)(a, \operatorname{neg}(r)(b)) \in \operatorname{idl}(r)) \end{array}}}
\vnbdef{ringoid-carr-in-set}{\vnbstmt{\forall r.\; (\operatorname{is\text{-}ringoid}(r) \Rightarrow (\operatorname{carr}(r) \in \mathbf{Set}))}}
\vnbdef{ringoid-add-right-inv}{\vnbstmt{\begin{array}{@{}l@{}} \forall r, a. \\ \quad (\operatorname{is\text{-}ringoid}(r) \wedge (a \in \operatorname{carr}(r))) \Rightarrow (\operatorname{add}(r)(a, \operatorname{neg}(r)(a)) = \operatorname{zero}(r)) \end{array}}}
\vnbdef{ringoid-ideal-subset}{\vnbstmt{\forall r.\; (\operatorname{is\text{-}ringoid}(r) \Rightarrow (\operatorname{idl}(r) \subseteq \operatorname{carr}(r)))}}
\vnbdef{ringoid-ideal-zero}{\vnbstmt{\forall r.\; (\operatorname{is\text{-}ringoid}(r) \Rightarrow (\operatorname{zero}(r) \in \operatorname{idl}(r)))}}
\vnbdef{ringoid-ideal-add}{\vnbstmt{\begin{array}{@{}l@{}} \forall r, a, b. \\ \quad (\operatorname{is\text{-}ringoid}(r) \wedge ((a \in \operatorname{idl}(r)) \wedge (b \in \operatorname{idl}(r)))) \Rightarrow (\operatorname{add}(r)(a, b) \in \operatorname{idl}(r)) \end{array}}}
\vnbdef{ringoid-ideal-neg}{\vnbstmt{\forall r, a.\; ((\operatorname{is\text{-}ringoid}(r) \wedge (a \in \operatorname{idl}(r))) \Rightarrow (\operatorname{neg}(r)(a) \in \operatorname{idl}(r)))}}
\vnbdef{ringoid-ideal-mul-left}{\vnbstmt{\begin{array}{@{}l@{}} \forall r, x, a. \\ \quad (\operatorname{is\text{-}ringoid}(r) \wedge ((x \in \operatorname{carr}(r)) \wedge (a \in \operatorname{idl}(r)))) \Rightarrow (\operatorname{mul}(r)(x, a) \in \operatorname{idl}(r)) \end{array}}}
\vnbdef{ringoid-ideal-mul-right}{\vnbstmt{\begin{array}{@{}l@{}} \forall r, a, x. \\ \quad (\operatorname{is\text{-}ringoid}(r) \wedge ((a \in \operatorname{idl}(r)) \wedge (x \in \operatorname{carr}(r)))) \Rightarrow (\operatorname{mul}(r)(a, x) \in \operatorname{idl}(r)) \end{array}}}
\vnbdef{ringoid-class}{\vnbstmt{\forall s.\; ((s \in \mathit{ringoid}) \Leftrightarrow \operatorname{is\text{-}ringoid}(s))}}
\vnbdef{is-hom-ringoid-def}{\vnbstmt{\begin{array}{@{}l@{}} \forall a, b, f. \\ \quad \operatorname{is\text{-}hom\text{-}ringoid}(a, b, f) \\ \quad \Leftrightarrow\; \operatorname{is\text{-}ringoid}(a) \wedge \\ \quad \operatorname{is\text{-}ringoid}(b) \wedge \\ \quad f \in (\operatorname{carr}(a) \to \operatorname{carr}(b)) \wedge \\ \quad \operatorname{is\text{-}hom\text{-}ring}(\operatorname{ringoid\text{-}as\text{-}ring}(a), \operatorname{ringoid\text{-}as\text{-}ring}(b), f) \wedge \\ \quad \forall x \in \operatorname{idl}(a).\; (f(x) \in \operatorname{idl}(b)) \end{array}}}

\vnbbookfile{structure-library/commutative-ring}{structure-library/commutative-ring}
\noindent{\small commutative-ring.scm -- COMMUTATIVE-RING: a ring whose MUL is commutative.}\par
\par\medskip\noindent\textbf{Definition} (\texttt{hom-commutative-ring}).\;$\operatorname{hom\text{-}commutative\text{-}ring}(a, b) :=$
\begin{equation*}
\{\operatorname{homf\_} \in (\operatorname{carr}(a) \to \operatorname{carr}(b)) : \operatorname{is\text{-}hom\text{-}commutative\text{-}ring}(a, b, \operatorname{homf\_})\}
\end{equation*}

\vnbdef{is-commutative-ring-def}{\vnbstmt{\begin{array}{@{}l@{}} \forall s. \\ \quad \operatorname{is\text{-}commutative\text{-}ring}(s) \\ \quad \Leftrightarrow\; (\operatorname{is\text{-}ring}(s) \wedge \forall a, b \in \operatorname{carr}(s).\; (\operatorname{mul}(s)(a, b) = \operatorname{mul}(s)(b, a))) \end{array}}}
\vnbdef{is-hom-commutative-ring-def}{\vnbstmt{\begin{array}{@{}l@{}} \forall a, b, f. \\ \quad \operatorname{is\text{-}hom\text{-}commutative\text{-}ring}(a, b, f) \\ \quad \Leftrightarrow\; \operatorname{is\text{-}commutative\text{-}ring}(a) \wedge (\operatorname{is\text{-}commutative\text{-}ring}(b) \wedge \operatorname{is\text{-}hom\text{-}ring}(a, b, f)) \end{array}}}
\vnbdef{commutative-ring-class}{\vnbstmt{\forall s.\; ((s \in \operatorname{commutative\text{-}ring}) \Leftrightarrow \operatorname{is\text{-}commutative\text{-}ring}(s))}}

\vnbbookfile{structure-library/integral-domain}{structure-library/integral-domain}
\noindent{\small integral-domain.scm -- INTEGRAL-DOMAIN: a commutative ring with ONE /= ZERO and no zero divisors.  Reuses the RING shape and accessors.}\par
\par\medskip\noindent\textbf{Definition} (\texttt{hom-integral-domain}).\;$\operatorname{hom\text{-}integral\text{-}domain}(a, b) :=$
\begin{equation*}
\{\operatorname{homf\_} \in (\operatorname{carr}(a) \to \operatorname{carr}(b)) : \operatorname{is\text{-}hom\text{-}integral\text{-}domain}(a, b, \operatorname{homf\_})\}
\end{equation*}

\vnbdef{is-integral-domain-def}{\vnbstmt{\begin{array}{@{}l@{}} \forall s. \\ \quad \operatorname{is\text{-}integral\text{-}domain}(s) \\ \quad \Leftrightarrow\; \operatorname{is\text{-}commutative\text{-}ring}(s) \wedge \\ \quad \neg (\operatorname{one}(s) = \operatorname{zero}(s)) \wedge \\ \quad \forall a, b \in \operatorname{carr}(s). \\ \quad \quad (\operatorname{mul}(s)(a, b) = \operatorname{zero}(s)) \Rightarrow ((a = \operatorname{zero}(s)) \vee (b = \operatorname{zero}(s))) \end{array}}}
\vnbdef{is-hom-integral-domain-def}{\vnbstmt{\begin{array}{@{}l@{}} \forall a, b, f. \\ \quad \operatorname{is\text{-}hom\text{-}integral\text{-}domain}(a, b, f) \\ \quad \Leftrightarrow\; (\operatorname{is\text{-}integral\text{-}domain}(a) \wedge (\operatorname{is\text{-}integral\text{-}domain}(b) \wedge \operatorname{is\text{-}hom\text{-}commutative\text{-}ring}(a, b, f))) \end{array}}}
\vnbdef{integral-domain-class}{\vnbstmt{\forall s.\; ((s \in \operatorname{integral\text{-}domain}) \Leftrightarrow \operatorname{is\text{-}integral\text{-}domain}(s))}}

\vnbbookfile{structure-library/field}{structure-library/field}
\noindent{\small field.scm -- FIELD: a commutative ring in which every nonzero element has a multiplicative inverse.}\par
\vnbstructdef{field}{A field $s$ is an 8-tuple whose components are listed below.\par
\begin{enumerate}[nosep]
\item $\operatorname{carr}(s) = \operatorname{nth}(1, s)$, a carrier (\texttt{carriers}): $\operatorname{carr}(s) \in \mathbf{Set}$.
\item $\operatorname{add}(s) = \operatorname{nth}(2, s)$, an operation (\texttt{op}): $\operatorname{add}(s) \in ((\operatorname{carr}(s) \times \operatorname{carr}(s)) \to \operatorname{carr}(s))$.
\item $\operatorname{mul}(s) = \operatorname{nth}(3, s)$, an operation (\texttt{op}): $\operatorname{mul}(s) \in ((\operatorname{carr}(s) \times \operatorname{carr}(s)) \to \operatorname{carr}(s))$.
\item $\operatorname{neg}(s) = \operatorname{nth}(4, s)$, an operation (\texttt{op}): $\operatorname{neg}(s) \in (\operatorname{carr}(s) \to \operatorname{carr}(s))$.
\item $\operatorname{zero}(s) = \operatorname{nth}(5, s)$, a constant (\texttt{constant}): $\operatorname{zero}(s) \in \operatorname{carr}(s)$.
\item $\operatorname{one}(s) = \operatorname{nth}(6, s)$, a constant (\texttt{constant}): $\operatorname{one}(s) \in \operatorname{carr}(s)$.
\item $\operatorname{non\text{-}zero}(s) = \operatorname{nth}(7, s)$, a derived carrier, fixed by the other components (\texttt{derived}): $\operatorname{non\text{-}zero}(s) \in \mathbf{Set}$ and $\operatorname{non\text{-}zero}(s) = \operatorname{difference}(\operatorname{carr}(s), \operatorname{singleton}(\operatorname{zero}(s)))$.
\item $\operatorname{mul\text{-}inv}(s) = \operatorname{nth}(8, s)$, an operation (\texttt{op}): $\operatorname{mul\text{-}inv}(s) \in (\operatorname{non\text{-}zero}(s) \to \operatorname{non\text{-}zero}(s))$.
\end{enumerate}
The laws (the \texttt{property} and \texttt{law} clauses):
\begin{itemize}[nosep]
\item $\operatorname{is\text{-}associative}(\operatorname{add}(s), \operatorname{carr}(s))$
\item $\operatorname{is\text{-}commutative}(\operatorname{add}(s), \operatorname{carr}(s))$
\item $\operatorname{is\text{-}identity}(\operatorname{add}(s), \operatorname{zero}(s), \operatorname{carr}(s))$
\item $\operatorname{has\text{-}inverses}(\operatorname{add}(s), \operatorname{zero}(s), \operatorname{neg}(s), \operatorname{carr}(s))$
\item $\operatorname{is\text{-}associative}(\operatorname{mul}(s), \operatorname{carr}(s))$
\item $\operatorname{is\text{-}commutative}(\operatorname{mul}(s), \operatorname{carr}(s))$
\item $\operatorname{is\text{-}identity}(\operatorname{mul}(s), \operatorname{one}(s), \operatorname{carr}(s))$
\item $\operatorname{is\text{-}distributive}(\operatorname{add}(s), \operatorname{mul}(s), \operatorname{carr}(s))$
\item $\neg (\operatorname{zero}(s) = \operatorname{one}(s))$
\item $\forall \operatorname{a\_} \in \operatorname{non\text{-}zero}(s).\; (\operatorname{mul}(s)(\operatorname{a\_}, \operatorname{mul\text{-}inv}(s)(\operatorname{a\_})) = \operatorname{one}(s))$
\end{itemize}
The characteristic law (the axiom \texttt{is-field}): $\operatorname{is\text{-}field}(s)$ holds if and only if $\operatorname{length}(s) = 8$, every component has the typing listed with it, every law holds, and $s \in \mathbf{Set}$ (the structure is a set).
}
\par\medskip\noindent\textbf{Definition} (\texttt{hom-field}).\;$\operatorname{hom\text{-}field}(a, b) :=$ $\{\operatorname{homf\_} \in (\operatorname{carr}(a) \to \operatorname{carr}(b)) : \operatorname{is\text{-}hom\text{-}field}(a, b, \operatorname{homf\_})\}$
\par\medskip\noindent\textbf{Definition} (\texttt{hom-field-ring}).\;$\operatorname{hom\text{-}field\text{-}ring}(a, b) :=$
\begin{equation*}
\{\operatorname{homf\_} \in (\operatorname{carr}(a) \to \operatorname{carr}(b)) : \operatorname{is\text{-}hom\text{-}field\text{-}ring}(a, b, \operatorname{homf\_})\}
\end{equation*}

\vnbdef{field-non-zero-carrier}{\vnbstmt{\begin{array}{@{}l@{}} \forall s. \\ \quad \operatorname{is\text{-}field}(s) \Rightarrow (\operatorname{non\text{-}zero}(s) = \operatorname{difference}(\operatorname{carr}(s), \operatorname{singleton}(\operatorname{zero}(s)))) \end{array}}}
\vnbdef{field-class}{\vnbstmt{\forall s.\; ((s \in \mathit{field}) \Leftrightarrow \operatorname{is\text{-}field}(s))}}
\vnbdef{is-field-ring-def}{\vnbstmt{\begin{array}{@{}l@{}} \forall s. \\ \quad \operatorname{is\text{-}field\text{-}ring}(s) \\ \quad \Leftrightarrow\; \operatorname{is\text{-}commutative\text{-}ring}(s) \wedge \\ \quad \neg (\operatorname{one}(s) = \operatorname{zero}(s)) \wedge \\ \quad \forall a \in \operatorname{carr}(s). \\ \quad \quad \neg (a = \operatorname{zero}(s)) \Rightarrow \exists b \in \operatorname{carr}(s).\; (\operatorname{mul}(s)(a, b) = \operatorname{one}(s)) \end{array}}}
\vnbdef{field-ring-class}{\vnbstmt{\forall s.\; ((s \in \operatorname{field\text{-}ring}) \Leftrightarrow \operatorname{is\text{-}field\text{-}ring}(s))}}
\vnbdef{is-hom-field-ring-def}{\vnbstmt{\begin{array}{@{}l@{}} \forall a, b, f. \\ \quad \operatorname{is\text{-}hom\text{-}field\text{-}ring}(a, b, f) \\ \quad \Leftrightarrow\; \operatorname{is\text{-}field\text{-}ring}(a) \wedge (\operatorname{is\text{-}field\text{-}ring}(b) \wedge \operatorname{is\text{-}hom\text{-}commutative\text{-}ring}(a, b, f)) \end{array}}}
\vnbdef{is-hom-field-def}{\vnbstmt{\begin{array}{@{}l@{}} \forall a, b, f. \\ \quad \operatorname{is\text{-}hom\text{-}field}(a, b, f) \\ \quad \Leftrightarrow\; \operatorname{is\text{-}field}(a) \wedge \\ \quad \operatorname{is\text{-}field}(b) \wedge \\ \quad f \in (\operatorname{carr}(a) \to \operatorname{carr}(b)) \wedge \\ \quad \forall x_{1}, x_{2} \in \operatorname{carr}(a). \\ \quad \quad f(\operatorname{add}(a)(x_{1}, x_{2})) = \operatorname{add}(b)(f(x_{1}), f(x_{2})) \wedge \\ \quad \forall x_{1}, x_{2} \in \operatorname{carr}(a). \\ \quad \quad f(\operatorname{mul}(a)(x_{1}, x_{2})) = \operatorname{mul}(b)(f(x_{1}), f(x_{2})) \wedge \\ \quad \forall x_{1} \in \operatorname{carr}(a).\; (f(\operatorname{neg}(a)(x_{1})) = \operatorname{neg}(b)(f(x_{1}))) \wedge \\ \quad f(\operatorname{zero}(a)) = \operatorname{zero}(b) \wedge \\ \quad f(\operatorname{one}(a)) = \operatorname{one}(b) \wedge \\ \quad \forall x_{1} \in \operatorname{non\text{-}zero}(a).\; (f(\operatorname{mul\text{-}inv}(a)(x_{1})) = \operatorname{mul\text{-}inv}(b)(f(x_{1}))) \end{array}}}

\vnbbookfile{structure-library/euclidean-ring}{structure-library/euclidean-ring}
\noindent{\small euclidean-ring.scm -- EUCLIDEAN-RING: an integral domain admitting a Euclidean (degree) function with division-with-remainder.}\par
\par\medskip\noindent\textbf{Definition} (\texttt{euclidean-gauges}).\;$\operatorname{euclidean\text{-}gauges}(s) :=$ $\{\mathit{dg} \in (\operatorname{carr}(s) \to \mathbb{N}) : \operatorname{has\text{-}div\text{-}remainder}(s, \mathit{dg})\}$
\par\medskip\noindent\textbf{Definition} (\texttt{gauge}).\;$\operatorname{gauge}(s) :=$ $\varepsilon(\operatorname{euclidean\text{-}gauges}(s))$
\par\medskip\noindent\textbf{Definition} (\texttt{has-div-remainder}).\;For all $s$ and $\mathit{dg}$, 
\begin{equation*}
\begin{array}{@{}l@{}} \operatorname{has\text{-}div\text{-}remainder}(s, \mathit{dg}) \\ \Leftrightarrow\; \forall \operatorname{a\_}, b \in \operatorname{carr}(s). \\ \quad \neg (b = \operatorname{zero}(s)) \Rightarrow  \\ \quad \quad \exists q, r \in \operatorname{carr}(s). \\ \quad \quad \quad (\operatorname{a\_} = \operatorname{add}(s)(\operatorname{mul}(s)(q, b), r)) \wedge ((r = \operatorname{zero}(s)) \vee (\operatorname{succ}(\operatorname{dg}(r)) \leq \operatorname{dg}(b))) \end{array}
\end{equation*}

\par\medskip\noindent\textbf{Definition} (\texttt{hom-euclidean-ring}).\;$\operatorname{hom\text{-}euclidean\text{-}ring}(a, b) :=$
\begin{equation*}
\{\operatorname{homf\_} \in (\operatorname{carr}(a) \to \operatorname{carr}(b)) : \operatorname{is\text{-}hom\text{-}euclidean\text{-}ring}(a, b, \operatorname{homf\_})\}
\end{equation*}

\vnbdef{is-euclidean-ring-def}{\vnbstmt{\begin{array}{@{}l@{}} \forall s. \\ \quad \operatorname{is\text{-}euclidean\text{-}ring}(s) \\ \quad \Leftrightarrow\; \operatorname{is\text{-}integral\text{-}domain}(s) \wedge \\ \quad \exists g \in (\operatorname{carr}(s) \to \mathbb{N}). \forall a, b \in \operatorname{carr}(s). \\ \quad \quad \neg (b = \operatorname{zero}(s)) \Rightarrow  \\ \quad \quad \quad \exists q, r \in \operatorname{carr}(s). \\ \quad \quad \quad \quad (a = \operatorname{add}(s)(\operatorname{mul}(s)(q, b), r)) \wedge ((r = \operatorname{zero}(s)) \vee ((g(r) + 1) \leq g(b))) \end{array}}}
\vnbdef{euclidean-ring-class}{\vnbstmt{\forall s.\; ((s \in \operatorname{euclidean\text{-}ring}) \Leftrightarrow \operatorname{is\text{-}euclidean\text{-}ring}(s))}}
\vnbdef{is-hom-euclidean-ring-def}{\vnbstmt{\begin{array}{@{}l@{}} \forall a, b, f. \\ \quad \operatorname{is\text{-}hom\text{-}euclidean\text{-}ring}(a, b, f) \\ \quad \Leftrightarrow\; (\operatorname{is\text{-}euclidean\text{-}ring}(a) \wedge (\operatorname{is\text{-}euclidean\text{-}ring}(b) \wedge \operatorname{is\text{-}hom\text{-}integral\text{-}domain}(a, b, f))) \end{array}}}

\vnbbookfile{structure-library/ideal}{structure-library/ideal}
\noindent{\small RETIRED 2026-09-17 (proven): ideal-elt-in-carrier, principal-ideal-in-ideal -- theorem-library/rake-algebra.scm ideal.scm -- ideals, principal ideals, and principal-ideal domains over a commutative ring.  The vocabulary the Euclidean-ring =$>$ PID stress test needs (theorem-library/euclidean-ring-pid.scm).}\par
\par\medskip\noindent\textbf{Definition} (\texttt{principal-ideal}).\;$\operatorname{principal\text{-}ideal}(s, a) :=$ $\{x \in \operatorname{carr}(s) : \exists r \in \operatorname{carr}(s).\; (x = \operatorname{mul}(s)(r, a))\}$
\par\medskip\noindent\textbf{Definition} (\texttt{is-ideal}).\;For all $s$ and $i$, 
\begin{equation*}
\begin{array}{@{}l@{}} \operatorname{is\text{-}ideal}(s, i) \\ \Leftrightarrow\; \operatorname{is\text{-}commutative\text{-}ring}(s) \wedge \\ i \subseteq \operatorname{carr}(s) \wedge \\ \operatorname{zero}(s) \in i \wedge \\ \forall a, b \in i.\; (\operatorname{add}(s)(a, b) \in i) \wedge \\ \forall a \in i.\; (\operatorname{neg}(s)(a) \in i) \wedge \\ \forall r \in \operatorname{carr}(s), a \in i.\; (\operatorname{mul}(s)(r, a) \in i) \end{array}
\end{equation*}

\par\medskip\noindent\textbf{Definition} (\texttt{hom-pid}).\;$\operatorname{hom\text{-}pid}(a, b) :=$ $\{\operatorname{homf\_} \in (\operatorname{carr}(a) \to \operatorname{carr}(b)) : \operatorname{is\text{-}hom\text{-}pid}(a, b, \operatorname{homf\_})\}$
\vnbdef{principal-ideal-membership}{\vnbstmt{\begin{array}{@{}l@{}} \forall s, a, x. \\ \quad (x \in \operatorname{principal\text{-}ideal}(s, a)) \Leftrightarrow ((x \in \operatorname{carr}(s)) \wedge \exists r \in \operatorname{carr}(s).\; (x = \operatorname{mul}(s)(r, a))) \end{array}}}
\vnbdef{is-pid-def}{\vnbstmt{\begin{array}{@{}l@{}} \forall s. \\ \quad \operatorname{is\text{-}pid}(s) \\ \quad \Leftrightarrow\; \operatorname{is\text{-}integral\text{-}domain}(s) \wedge \\ \quad \forall i \in \mathbf{Set}. \\ \quad \quad \operatorname{is\text{-}ideal}(s, i) \Rightarrow \exists a \in \operatorname{carr}(s).\; (i = \operatorname{principal\text{-}ideal}(s, a)) \end{array}}}
\vnbdef{pid-class}{\vnbstmt{\forall s.\; ((s \in \mathit{pid}) \Leftrightarrow \operatorname{is\text{-}pid}(s))}}
\vnbdef{is-hom-pid-def}{\vnbstmt{\begin{array}{@{}l@{}} \forall a, b, f. \\ \quad \operatorname{is\text{-}hom\text{-}pid}(a, b, f) \Leftrightarrow (\operatorname{is\text{-}pid}(a) \wedge (\operatorname{is\text{-}pid}(b) \wedge \operatorname{is\text{-}hom\text{-}integral\text{-}domain}(a, b, f))) \end{array}}}

\vnbbookfile{structure-library/numeric-instances}{structure-library/numeric-instances}
\noindent{\small basic-rings.scm -- "officialize" NN/ZZ/QQ/RR/CC as algebraic structures}\par
\par\medskip\noindent\textbf{Definition} (\texttt{zz-ring}).\;
\begin{equation*}
\begin{array}{@{}l@{}} \operatorname{zz\text{-}ring} \\ =\; {[}\mathbb{Z}, (\lambda {[}\operatorname{x\_}, \operatorname{y\_}] \in (\mathbb{Z} \times \mathbb{Z}).\; +(\operatorname{x\_}, \operatorname{y\_})), (\lambda {[}\operatorname{x\_}, \operatorname{y\_}] \in (\mathbb{Z} \times \mathbb{Z}).\; \cdot(\operatorname{x\_}, \operatorname{y\_})), (\lambda \operatorname{x\_} \in \mathbb{Z}.\; -(\operatorname{x\_})), 0, 1] \end{array}
\end{equation*}

\par\medskip\noindent\textbf{Definition} (\texttt{qq-ring}).\;
\begin{equation*}
\begin{array}{@{}l@{}} \operatorname{qq\text{-}ring} \\ =\; {[}\mathbb{Q}, (\lambda {[}\operatorname{x\_}, \operatorname{y\_}] \in (\mathbb{Q} \times \mathbb{Q}).\; +(\operatorname{x\_}, \operatorname{y\_})), (\lambda {[}\operatorname{x\_}, \operatorname{y\_}] \in (\mathbb{Q} \times \mathbb{Q}).\; \cdot(\operatorname{x\_}, \operatorname{y\_})), (\lambda \operatorname{x\_} \in \mathbb{Q}.\; -(\operatorname{x\_})), 0, 1] \end{array}
\end{equation*}

\par\medskip\noindent\textbf{Definition} (\texttt{rr-normed-field}).\;
\begin{equation*}
\begin{array}{@{}l@{}} \operatorname{rr\text{-}normed\text{-}field} \\ =\; {[}\mathbb{R}, (\lambda {[}\operatorname{x\_}, \operatorname{y\_}] \in (\mathbb{R} \times \mathbb{R}).\; +(\operatorname{x\_}, \operatorname{y\_})), (\lambda {[}\operatorname{x\_}, \operatorname{y\_}] \in (\mathbb{R} \times \mathbb{R}).\; \cdot(\operatorname{x\_}, \operatorname{y\_})), (\lambda \operatorname{x\_} \in \mathbb{R}.\; -(\operatorname{x\_})), 0, 1, (\lambda \operatorname{x\_} \in \mathbb{R}.\; \lvert \operatorname{x\_} \rvert)] \end{array}
\end{equation*}

\par\medskip\noindent\textbf{Definition} (\texttt{cc-normed-field}).\;
\begin{equation*}
\begin{array}{@{}l@{}} \operatorname{cc\text{-}normed\text{-}field} \\ =\; {[}\mathbb{C}, (\lambda {[}\operatorname{x\_}, \operatorname{y\_}] \in (\mathbb{C} \times \mathbb{C}).\; +(\operatorname{x\_}, \operatorname{y\_})), (\lambda {[}\operatorname{x\_}, \operatorname{y\_}] \in (\mathbb{C} \times \mathbb{C}).\; \cdot(\operatorname{x\_}, \operatorname{y\_})), (\lambda \operatorname{x\_} \in \mathbb{C}.\; -(\operatorname{x\_})), 0, 1, (\lambda \operatorname{x\_} \in \mathbb{C}.\; \operatorname{magnitude}(\operatorname{x\_}))] \end{array}
\end{equation*}

\par\medskip\noindent\textbf{Definition} (\texttt{nn-add-monoid}).\;
\begin{equation*}
\operatorname{nn\text{-}add\text{-}monoid} = {[}\mathbb{N}, (\lambda {[}\operatorname{x\_}, \operatorname{y\_}] \in (\mathbb{N} \times \mathbb{N}).\; +(\operatorname{x\_}, \operatorname{y\_})), 0]
\end{equation*}

\par\medskip\noindent\textbf{Definition} (\texttt{rr-ms}).\;
\begin{equation*}
\operatorname{rr\text{-}ms} = {[}\mathbb{R}, (\lambda {[}x, y] \in (\mathbb{R} \times \mathbb{R}).\; \lvert -(x, y) \rvert)]
\end{equation*}

\par\medskip\noindent\textbf{Definition} (\texttt{qq-field}).\;
\begin{equation*}
\begin{array}{@{}l@{}} \operatorname{qq\text{-}field} \\ =\; {[}\mathbb{Q}, (\lambda {[}\operatorname{x\_}, \operatorname{y\_}] \in (\mathbb{Q} \times \mathbb{Q}).\; +(\operatorname{x\_}, \operatorname{y\_})), (\lambda {[}\operatorname{x\_}, \operatorname{y\_}] \in (\mathbb{Q} \times \mathbb{Q}).\; \cdot(\operatorname{x\_}, \operatorname{y\_})), (\lambda \operatorname{x\_} \in \mathbb{Q}.\; -(\operatorname{x\_})), 0, 1, \operatorname{difference}(\mathbb{Q}, \operatorname{singleton}(0)), (\lambda \operatorname{x\_} \in \operatorname{difference}(\mathbb{Q}, \operatorname{singleton}(0)).\; \operatorname{recip}(\operatorname{x\_}))] \end{array}
\end{equation*}

\vnbdef{binplus-apply}{\vnbstmt{\forall x, y.\; ((x + y) \simeq (x + y))}}
\vnbdef{bintimes-apply}{\vnbstmt{\forall x, y.\; ((x \cdot y) \simeq (x \cdot y))}}
\vnbdef{binneg-apply}{\vnbstmt{\forall x.\; (\operatorname{binneg}(x) \simeq -x)}}
\vnbdef{nary-plus-2}{\vnbstmt{\forall x, y.\; ((x + y) \simeq (x + y))}}
\vnbdef{nary-times-2}{\vnbstmt{\forall x, y.\; ((x \cdot y) \simeq (x \cdot y))}}
\vnbdef{nary-neg-1}{\vnbstmt{\forall x.\; (-x \simeq \operatorname{binneg}(x))}}
\vnbdef{nary-minus-2}{\vnbstmt{\forall x, y.\; ((x - y) \simeq (x + \operatorname{binneg}(y)))}}
\vnbdef{nary-plus-3}{\vnbstmt{\forall x, y, z.\; ((x + y + z) \simeq ((x + y) + z))}}
\vnbdef{nary-plus-4}{\vnbstmt{\forall w, x, y, z.\; ((w + x + y + z) \simeq (((w + x) + y) + z))}}
\vnbdef{nary-plus-5}{\vnbstmt{\begin{array}{@{}l@{}} \forall v, w, x, y, z. \\ \quad ((v + w + x + y + z) \simeq ((((v + w) + x) + y) + z)) \end{array}}}
\vnbdef{nary-times-3}{\vnbstmt{\forall x, y, z.\; ((x \cdot y \cdot z) \simeq ((x \cdot y) \cdot z))}}
\vnbdef{nary-times-4}{\vnbstmt{\forall w, x, y, z.\; ((w \cdot x \cdot y \cdot z) \simeq (((w \cdot x) \cdot y) \cdot z))}}
\vnbdef{nary-times-5}{\vnbstmt{\begin{array}{@{}l@{}} \forall v, w, x, y, z. \\ \quad ((v \cdot w \cdot x \cdot y \cdot z) \simeq ((((v \cdot w) \cdot x) \cdot y) \cdot z)) \end{array}}}
\vnbdef{cc-normed-field@neg}{\vnbstmt{(\operatorname{neg}(\operatorname{cc\text{-}normed\text{-}field}) \simeq (\lambda x \in \mathbb{C}.\; -x))}}
\vnbdef{rr-normed-field@fnrm}{\vnbstmt{(\operatorname{fnrm}(\operatorname{rr\text{-}normed\text{-}field}) \simeq (\lambda x \in \mathbb{R}.\; \lvert x \rvert))}}
\vnbdef{qq-field@neg}{\vnbstmt{(\operatorname{neg}(\operatorname{qq\text{-}field}) \simeq (\lambda x \in \mathbb{Q}.\; -x))}}
\vnbdef{rr-normed-field@carr}{\vnbstmt{(\operatorname{carr}(\operatorname{rr\text{-}normed\text{-}field}) \simeq \mathbb{R})}}
\vnbdef{qq-field@one}{\vnbstmt{(\operatorname{one}(\operatorname{qq\text{-}field}) \simeq 1)}}
\vnbdef{qq-field@mul}{\vnbstmt{(\operatorname{mul}(\operatorname{qq\text{-}field}) \simeq (\lambda {[}\operatorname{x\_}, \operatorname{y\_}] \in (\mathbb{Q} \times \mathbb{Q}).\; (\operatorname{x\_} \cdot \operatorname{y\_})))}}
\vnbdef{qq-ring@add}{\vnbstmt{(\operatorname{add}(\operatorname{qq\text{-}ring}) \simeq (\lambda {[}\operatorname{x\_}, \operatorname{y\_}] \in (\mathbb{Q} \times \mathbb{Q}).\; (\operatorname{x\_} + \operatorname{y\_})))}}
\vnbdef{zz-ring@neg}{\vnbstmt{(\operatorname{neg}(\operatorname{zz\text{-}ring}) \simeq (\lambda x \in \mathbb{Z}.\; -x))}}
\vnbdef{cc-normed-field@one}{\vnbstmt{(\operatorname{one}(\operatorname{cc\text{-}normed\text{-}field}) \simeq 1)}}
\vnbdef{cc-normed-field@carr}{\vnbstmt{(\operatorname{carr}(\operatorname{cc\text{-}normed\text{-}field}) \simeq \mathbb{C})}}
\vnbdef{qq-field@non-zero}{\vnbstmt{(\operatorname{non\text{-}zero}(\operatorname{qq\text{-}field}) \simeq \operatorname{difference}(\mathbb{Q}, \operatorname{singleton}(0)))}}
\vnbdef{zz-ring@add}{\vnbstmt{(\operatorname{add}(\operatorname{zz\text{-}ring}) \simeq (\lambda {[}\operatorname{x\_}, \operatorname{y\_}] \in (\mathbb{Z} \times \mathbb{Z}).\; (\operatorname{x\_} + \operatorname{y\_})))}}
\vnbdef{cc-normed-field@zero}{\vnbstmt{(\operatorname{zero}(\operatorname{cc\text{-}normed\text{-}field}) \simeq 0)}}
\vnbdef{rr-normed-field@add}{\vnbstmt{(\operatorname{add}(\operatorname{rr\text{-}normed\text{-}field}) \simeq (\lambda {[}\operatorname{x\_}, \operatorname{y\_}] \in (\mathbb{R} \times \mathbb{R}).\; (\operatorname{x\_} + \operatorname{y\_})))}}
\vnbdef{cc-normed-field@fnrm}{\vnbstmt{(\operatorname{fnrm}(\operatorname{cc\text{-}normed\text{-}field}) \simeq (\lambda x \in \mathbb{C}.\; \operatorname{magnitude}(x)))}}
\vnbdef{qq-ring@carr}{\vnbstmt{(\operatorname{carr}(\operatorname{qq\text{-}ring}) \simeq \mathbb{Q})}}
\vnbdef{zz-ring@carr}{\vnbstmt{(\operatorname{carr}(\operatorname{zz\text{-}ring}) \simeq \mathbb{Z})}}
\vnbdef{cc-normed-field@add}{\vnbstmt{(\operatorname{add}(\operatorname{cc\text{-}normed\text{-}field}) \simeq (\lambda {[}\operatorname{x\_}, \operatorname{y\_}] \in (\mathbb{C} \times \mathbb{C}).\; (\operatorname{x\_} + \operatorname{y\_})))}}
\vnbdef{zz-ring@mul}{\vnbstmt{(\operatorname{mul}(\operatorname{zz\text{-}ring}) \simeq (\lambda {[}\operatorname{x\_}, \operatorname{y\_}] \in (\mathbb{Z} \times \mathbb{Z}).\; (\operatorname{x\_} \cdot \operatorname{y\_})))}}
\vnbdef{nn-add-monoid@iden}{\vnbstmt{(\operatorname{iden}(\operatorname{nn\text{-}add\text{-}monoid}) \simeq 0)}}
\vnbdef{cc-normed-field@mul}{\vnbstmt{(\operatorname{mul}(\operatorname{cc\text{-}normed\text{-}field}) \simeq (\lambda {[}\operatorname{x\_}, \operatorname{y\_}] \in (\mathbb{C} \times \mathbb{C}).\; (\operatorname{x\_} \cdot \operatorname{y\_})))}}
\vnbdef{nn-add-monoid@opr}{\vnbstmt{(\operatorname{opr}(\operatorname{nn\text{-}add\text{-}monoid}) \simeq (\lambda {[}\operatorname{x\_}, \operatorname{y\_}] \in (\mathbb{N} \times \mathbb{N}).\; (\operatorname{x\_} + \operatorname{y\_})))}}
\vnbdef{rr-normed-field@neg}{\vnbstmt{(\operatorname{neg}(\operatorname{rr\text{-}normed\text{-}field}) \simeq (\lambda x \in \mathbb{R}.\; -x))}}
\vnbdef{qq-ring@one}{\vnbstmt{(\operatorname{one}(\operatorname{qq\text{-}ring}) \simeq 1)}}
\vnbdef{rr-normed-field@one}{\vnbstmt{(\operatorname{one}(\operatorname{rr\text{-}normed\text{-}field}) \simeq 1)}}
\vnbdef{rr-normed-field@zero}{\vnbstmt{(\operatorname{zero}(\operatorname{rr\text{-}normed\text{-}field}) \simeq 0)}}
\vnbdef{qq-ring@neg}{\vnbstmt{(\operatorname{neg}(\operatorname{qq\text{-}ring}) \simeq (\lambda x \in \mathbb{Q}.\; -x))}}
\vnbdef{qq-field@carr}{\vnbstmt{(\operatorname{carr}(\operatorname{qq\text{-}field}) \simeq \mathbb{Q})}}
\vnbdef{rr-ms@pts}{\vnbstmt{(\operatorname{pts}(\operatorname{rr\text{-}ms}) \simeq \mathbb{R})}}
\vnbdef{qq-ring@mul}{\vnbstmt{(\operatorname{mul}(\operatorname{qq\text{-}ring}) \simeq (\lambda {[}\operatorname{x\_}, \operatorname{y\_}] \in (\mathbb{Q} \times \mathbb{Q}).\; (\operatorname{x\_} \cdot \operatorname{y\_})))}}
\vnbdef{rr-ms@dist}{\vnbstmt{(\operatorname{dist}(\operatorname{rr\text{-}ms}) \simeq (\lambda {[}x, y] \in (\mathbb{R} \times \mathbb{R}).\; \lvert (x - y) \rvert))}}
\vnbdef{qq-field@zero}{\vnbstmt{(\operatorname{zero}(\operatorname{qq\text{-}field}) \simeq 0)}}
\vnbdef{rr-normed-field@mul}{\vnbstmt{(\operatorname{mul}(\operatorname{rr\text{-}normed\text{-}field}) \simeq (\lambda {[}\operatorname{x\_}, \operatorname{y\_}] \in (\mathbb{R} \times \mathbb{R}).\; (\operatorname{x\_} \cdot \operatorname{y\_})))}}
\vnbdef{qq-ring@zero}{\vnbstmt{(\operatorname{zero}(\operatorname{qq\text{-}ring}) \simeq 0)}}
\vnbdef{qq-field@add}{\vnbstmt{(\operatorname{add}(\operatorname{qq\text{-}field}) \simeq (\lambda {[}\operatorname{x\_}, \operatorname{y\_}] \in (\mathbb{Q} \times \mathbb{Q}).\; (\operatorname{x\_} + \operatorname{y\_})))}}
\vnbdef{zz-ring@one}{\vnbstmt{(\operatorname{one}(\operatorname{zz\text{-}ring}) \simeq 1)}}
\vnbdef{qq-field@mul-inv}{\vnbstmt{(\operatorname{mul\text{-}inv}(\operatorname{qq\text{-}field}) \simeq (\lambda x \in \operatorname{difference}(\mathbb{Q}, \operatorname{singleton}(0)).\; \operatorname{recip}(x)))}}
\vnbdef{zz-ring@zero}{\vnbstmt{(\operatorname{zero}(\operatorname{zz\text{-}ring}) \simeq 0)}}
\vnbdef{nn-add-monoid@carr}{\vnbstmt{(\operatorname{carr}(\operatorname{nn\text{-}add\text{-}monoid}) \simeq \mathbb{N})}}

\vnbbookfile{structure-library/zz-divisibility}{structure-library/zz-divisibility}
\noindent{\small zz-divisibility.scm -- divisibility, Bezout combinations, gcd and coprimality on the INTEGERS.  Vocabulary only; nothing is asserted here.}\par
\par\medskip\noindent\textbf{Definition} (\texttt{zz-divides}).\;For all $a$ and $b$, 
\begin{equation*}
\operatorname{zz\text{-}divides}(a, b) \Leftrightarrow ((a \in \mathbb{Z}) \wedge ((b \in \mathbb{Z}) \wedge \exists \operatorname{c\_} \in \mathbb{Z}.\; (b = \cdot(a, \operatorname{c\_}))))
\end{equation*}

\par\medskip\noindent\textbf{Definition} (\texttt{zz-bezout-set}).\;$\operatorname{zz\text{-}bezout\text{-}set}(a, b) :=$
\begin{equation*}
\{z \in \mathbb{Z} : \exists \operatorname{x\_}, \operatorname{y\_} \in \mathbb{Z}.\; (z = ((\operatorname{x\_} \cdot a) + (\operatorname{y\_} \cdot b)))\}
\end{equation*}

\par\medskip\noindent\textbf{Definition} (\texttt{zz-is-gcd}).\;For all $a$, $b$ and $d$, 
\begin{equation*}
\begin{array}{@{}l@{}} \operatorname{zz\text{-}is\text{-}gcd}(a, b, d) \\ \Leftrightarrow\; \operatorname{zz\text{-}divides}(d, a) \wedge \\ \operatorname{zz\text{-}divides}(d, b) \wedge \\ \forall \operatorname{e\_} \in \mathbb{Z}. \\ \quad \operatorname{zz\text{-}divides}(\operatorname{e\_}, a) \Rightarrow (\operatorname{zz\text{-}divides}(\operatorname{e\_}, b) \Rightarrow \operatorname{zz\text{-}divides}(\operatorname{e\_}, d)) \end{array}
\end{equation*}

\par\medskip\noindent\textbf{Definition} (\texttt{zz-coprime}).\;For all $a$ and $b$,  $\operatorname{zz\text{-}coprime}(a, b) \Leftrightarrow \operatorname{zz\text{-}is\text{-}gcd}(a, b, 1)$
\vnbdef{zz-bezout-set-membership}{\vnbstmt{\begin{array}{@{}l@{}} \forall a, b, z. \\ \quad (z \in \operatorname{zz\text{-}bezout\text{-}set}(a, b)) \Leftrightarrow ((z \in \mathbb{Z}) \wedge \exists x, y \in \mathbb{Z}.\; (z = ((x \cdot a) + (y \cdot b)))) \end{array}}}

\vnbbookfile{structure-library/monoid-power}{structure-library/monoid-power}
\noindent{\small RETIRED 2026-09-17 (proven): mpow-type -- was an unwarranted axiom; theorem-library/rake-algebra2.scm (mpow-type-ind, exponent outermost, then one fact). monoid-power.scm -- MPOW: the natural-number power x\textasciicircum{}n in a monoid.}\par
\par\medskip\noindent\textbf{Definition} (\texttt{mpow}).\;For all $m$ and $x$,  $(\operatorname{mpow}(m, x, 0) \simeq \operatorname{iden}(m))$\par
Let $m$ and $x$ be arbitrary. Let $n \in \mathbb{N}$. Then $(\operatorname{mpow}(m, x, \operatorname{succ}(n)) \simeq \operatorname{opr}(m)(x, \operatorname{mpow}(m, x, n)))$

\vnbbookfile{structure-library/zz-action}{structure-library/zz-action}
\noindent{\small zz-action.scm -- ZZ-ACT: the integer action on an abelian group.}\par
\par\medskip\noindent\textbf{Definition} (\texttt{zz-act}).\;$\operatorname{zz\text{-}act}(\operatorname{g\_}, \operatorname{k\_}, \operatorname{a\_}) :=$
\begin{equation*}
\left(\text{if }(\operatorname{k\_} \in \mathbb{N})\text{ then }\operatorname{mpow}(\operatorname{abelian\text{-}group\text{-}as\text{-}monoid}(\operatorname{g\_}), \operatorname{a\_}, \operatorname{k\_})\text{ else }\operatorname{inv}(\operatorname{g\_})(\operatorname{mpow}(\operatorname{abelian\text{-}group\text{-}as\text{-}monoid}(\operatorname{g\_}), \operatorname{a\_}, -\operatorname{k\_}))\right)
\end{equation*}


\vnbbookfile{structure-library/ring-power}{structure-library/ring-power}
\noindent{\small RETIRED 2026-09-17 (proven, rake batch P): ring-power-one, ring-power-add, ring-power-mult -- theorem-library/rake-combinatorics.scm, modulo 0 (mpow laws through the multiplicative view). RETIRED 2026-09-17 (proven): ring-power-type -- theorem-library/rake-finsum-typing.scm ring-power.scm -- RING-POWER: the ...}\par
\par\medskip\noindent\textbf{Definition} (\texttt{ring-power}).\;$\operatorname{ring\text{-}power}(r, x, n) :=$ $\operatorname{mpow}(\operatorname{commutative\text{-}ring\text{-}multiplicative\text{-}cm}(r), x, n)$

\vnbbookfile{structure-library/sequences}{structure-library/sequences}
\noindent{\small sequences.scm -- PROD-ORD, SUM, SUM-AG: products and sums over initial NN-segments}\par
\par\medskip\noindent\textbf{Definition} (\texttt{prod-ord}).\;For all $m$ and $f$,  $(\operatorname{prod\text{-}ord}(m, f, 0) \simeq \operatorname{iden}(m))$\par
Let $m$ and $f$ be arbitrary. Let $n \in \mathbb{N}$. Then $(\operatorname{prod\text{-}ord}(m, f, \operatorname{succ}(n)) \simeq \operatorname{opr}(m)(\operatorname{prod\text{-}ord}(m, f, n), f(n)))$
\par\medskip\noindent\textbf{Definition} (\texttt{sum}).\;For all $r$ and $f$,  $(\operatorname{sum}(r, f, 0) \simeq \operatorname{zero}(r))$\par
Let $r$ and $f$ be arbitrary. Let $n \in \mathbb{N}$. Then $(\operatorname{sum}(r, f, \operatorname{succ}(n)) \simeq \operatorname{add}(r)(\operatorname{sum}(r, f, n), f(n)))$
\par\medskip\noindent\textbf{Definition} (\texttt{sum-ag}).\;For all $\mathit{ag}$ and $f$,  $(\operatorname{sum\text{-}ag}(\mathit{ag}, f, 0) \simeq \operatorname{iden}(\mathit{ag}))$\par
Let $\mathit{ag}$ and $f$ be arbitrary. Let $n \in \mathbb{N}$. Then $(\operatorname{sum\text{-}ag}(\mathit{ag}, f, \operatorname{succ}(n)) \simeq \operatorname{opr}(\mathit{ag})(\operatorname{sum\text{-}ag}(\mathit{ag}, f, n), f(n)))$
\par\medskip\noindent\textbf{Definition} (\texttt{ring-prod-n}).\;For every $f$,  $(\operatorname{ring\text{-}prod\text{-}n}(f, 0) \simeq \operatorname{zero\text{-}ring})$\par
Let $f$ be arbitrary. Let $n \in \mathbb{N}$. Then
\begin{equation*}
(\operatorname{ring\text{-}prod\text{-}n}(f, \operatorname{succ}(n)) \simeq \operatorname{ring\text{-}prod}(\operatorname{ring\text{-}prod\text{-}n}(f, n), f(n)))
\end{equation*}


\vnbbookfile{structure-library/finprod}{structure-library/finprod}
\noindent{\small finprod.scm -- FINPROD / PROD-RING: finite PRODUCT of a function over a finite set.}\par
\par\medskip\noindent\textbf{Definition} (\texttt{finprod}).\;$\operatorname{finprod}(m, f, s) :=$ $\operatorname{finsum}(m, f, s)$
\par\medskip\noindent\textbf{Definition} (\texttt{prod-ring}).\;$\operatorname{prod\text{-}ring}(r, f, s) :=$ $\operatorname{finprod}(\operatorname{commutative\text{-}ring\text{-}multiplicative\text{-}cm}(r), f, s)$
\par\medskip\noindent\textbf{Definition} (\texttt{prod-set}).\;$\operatorname{prod\text{-}set}(\mathit{cm}, s, f) :=$ $\operatorname{finsum}(\mathit{cm}, f, s)$

\vnbbookfile{structure-library/polynomial}{structure-library/polynomial}
\noindent{\small polynomial.scm -- the monoid-algebra construction A[M], and polynomials as its specialisation A[NN].}\par
\par\medskip\noindent\textbf{Definition} (\texttt{supp}).\;$\operatorname{supp}(a, m, f) :=$ $\{\operatorname{x\_} \in \operatorname{carr}(m) : \neg (f(\operatorname{x\_}) = \operatorname{zero}(a))\}$
\par\medskip\noindent\textbf{Definition} (\texttt{finsupp}).\;$\operatorname{finsupp}(a, m) :=$ $\{f \in (\operatorname{carr}(m) \to \operatorname{carr}(a)) : (|\operatorname{supp}(a, m, f)| \in \mathbb{N})\}$
\par\medskip\noindent\textbf{Definition} (\texttt{monalg-add}).\;$\operatorname{monalg\text{-}add}(a, m, f, g) :=$ $(\lambda \operatorname{x\_} \in \operatorname{carr}(m).\; \operatorname{add}(a)(f(\operatorname{x\_}), g(\operatorname{x\_})))$
\par\medskip\noindent\textbf{Definition} (\texttt{monalg-neg}).\;$\operatorname{monalg\text{-}neg}(a, m, f) :=$ $(\lambda \operatorname{x\_} \in \operatorname{carr}(m).\; \operatorname{neg}(a)(f(\operatorname{x\_})))$
\par\medskip\noindent\textbf{Definition} (\texttt{monalg-zero}).\;$\operatorname{monalg\text{-}zero}(a, m) :=$ $(\lambda \operatorname{x\_} \in \operatorname{carr}(m).\; \operatorname{zero}(a))$
\par\medskip\noindent\textbf{Definition} (\texttt{monalg-one}).\;$\operatorname{monalg\text{-}one}(a, m) :=$ $(\lambda \operatorname{x\_} \in \operatorname{carr}(m).\; \left(\text{if }(\operatorname{x\_} = \operatorname{iden}(m))\text{ then }\operatorname{one}(a)\text{ else }\operatorname{zero}(a)\right))$
\par\medskip\noindent\textbf{Definition} (\texttt{monalg-mul}).\;$\operatorname{monalg\text{-}mul}(a, m, f, g) :=$
\begin{equation*}
\begin{array}{@{}l@{}} (\lambda \operatorname{x\_} \in \operatorname{carr}(m).\; \\ \quad \operatorname{finsum}(\operatorname{ring\text{-}additive\text{-}ag}(a), \\ \quad \quad (\lambda p \in \{p \in (\operatorname{supp}(a, m, f) \times \operatorname{supp}(a, m, g)) : (\operatorname{opr}(m)(\operatorname{nth}(1, p), \operatorname{nth}(2, p)) = \operatorname{x\_})\}.\; \\ \quad \quad \quad \operatorname{mul}(a)(f(\operatorname{nth}(1, p)), g(\operatorname{nth}(2, p)))), \\ \quad \quad \{p \in (\operatorname{supp}(a, m, f) \times \operatorname{supp}(a, m, g)) : (\operatorname{opr}(m)(\operatorname{nth}(1, p), \operatorname{nth}(2, p)) = \operatorname{x\_})\})) \end{array}
\end{equation*}

\par\medskip\noindent\textbf{Definition} (\texttt{monalg}).\;$\operatorname{monalg}(a, m) :=$
\begin{equation*}
{[}\operatorname{finsupp}(a, m), (\lambda {[}f, g] \in (\operatorname{finsupp}(a, m) \times \operatorname{finsupp}(a, m)).\; \operatorname{monalg\text{-}add}(a, m, f, g)), (\lambda {[}f, g] \in (\operatorname{finsupp}(a, m) \times \operatorname{finsupp}(a, m)).\; \operatorname{monalg\text{-}mul}(a, m, f, g)), (\lambda f \in \operatorname{finsupp}(a, m).\; \operatorname{monalg\text{-}neg}(a, m, f)), \operatorname{monalg\text{-}zero}(a, m), \operatorname{monalg\text{-}one}(a, m)]
\end{equation*}

\par\medskip\noindent\textbf{Definition} (\texttt{poly}).\;$\operatorname{poly}(a) :=$ $\operatorname{monalg}(a, \operatorname{nn\text{-}add\text{-}monoid})$

\vnbbookfile{structure-library/la-z6-vocab}{structure-library/la-z6-vocab}
\noindent{\small la-z6-vocab.scm -- polynomials in n indeterminates and the permutation of the indeterminates (linear algebra, the paragraph before Proposition 3.15 and 3.15 itself) (LA-Z6, 2026-10-25). No proof in this file.}\par
\par\medskip\noindent\textbf{Definition} (\texttt{is-iso-ring}).\;For all $a$, $b$ and $f$, 
\begin{equation*}
\operatorname{is\text{-}iso\text{-}ring}(a, b, f) \Leftrightarrow (\operatorname{is\text{-}hom\text{-}ring}(a, b, f) \wedge (f \in \mathrm{Bij}(\operatorname{carr}(a), \operatorname{carr}(b))))
\end{equation*}

\par\medskip\noindent\textbf{Definition} (\texttt{monalg-pullback}).\;$\operatorname{monalg\text{-}pullback}(a, m, \mathit{ph}, f) :=$ $(\lambda \operatorname{mpx\_} \in \operatorname{carr}(m).\; f(\operatorname{ph}(\operatorname{mpx\_})))$
\par\medskip\noindent\textbf{Definition} (\texttt{mi-add}).\;$\operatorname{mi\text{-}add}(x, u, v) :=$ $(\lambda \operatorname{mia\_} \in x.\; (u(\operatorname{mia\_}) + v(\operatorname{mia\_})))$
\par\medskip\noindent\textbf{Definition} (\texttt{mi-zero}).\;$\operatorname{mi\text{-}zero}(x) :=$ $(\lambda \operatorname{mia\_} \in x.\; 0)$
\par\medskip\noindent\textbf{Definition} (\texttt{multi-index-monoid}).\;$\operatorname{multi\text{-}index\text{-}monoid}(x) :=$
\begin{equation*}
{[}(x \to \mathbb{N}), (\lambda {[}\operatorname{miu\_}, \operatorname{miv\_}] \in ((x \to \mathbb{N}) \times (x \to \mathbb{N})).\; \operatorname{mi\text{-}add}(x, \operatorname{miu\_}, \operatorname{miv\_})), \operatorname{mi\text{-}zero}(x)]
\end{equation*}

\par\medskip\noindent\textbf{Definition} (\texttt{mi-reindex}).\;$\operatorname{mi\text{-}reindex}(x, s, b) :=$ $(\lambda \operatorname{mir\_} \in x.\; b(s(\operatorname{mir\_})))$
\par\medskip\noindent\textbf{Definition} (\texttt{mpoly}).\;$\operatorname{mpoly}(r, n) :=$ $\operatorname{monalg}(r, \operatorname{multi\text{-}index\text{-}monoid}(\operatorname{interval}(1, n)))$
\par\medskip\noindent\textbf{Definition} (\texttt{var-permute}).\;$\operatorname{var\text{-}permute}(r, n, s, p) :=$
\begin{equation*}
\begin{array}{@{}l@{}} \operatorname{monalg\text{-}pullback}(r, \operatorname{multi\text{-}index\text{-}monoid}(\operatorname{interval}(1, n)), \\ \quad (\lambda \operatorname{mib\_} \in (\operatorname{interval}(1, n) \to \mathbb{N}).\; \operatorname{mi\text{-}reindex}(\operatorname{interval}(1, n), s, \operatorname{mib\_})), \\ \quad p) \end{array}
\end{equation*}


\vnbbookfile{structure-library/binomial}{structure-library/binomial}
\noindent{\small MOVED 2026-09-20 (batch 12-A) from theorem-library/: this file is VOCABULARY -- definitions, notation and warranted supports, not one proof -- and every theorem stated with it had to load below it.  Its load.scm slot is unchanged. RETIRED 2026-09-17 (proven): bt-mul-comm -- theorem-library/rake-combinatorics.scm}\par
\par\medskip\noindent\textbf{Definition} (\texttt{comb-kk}).\;For all $r$, $x$ and $y$, 
\begin{equation*}
(\operatorname{comb\text{-}kk}(r, x, y, 0) \simeq (\lambda k \in \mathbb{Z}.\; \left(\text{if }(k = 0)\text{ then }\operatorname{one}(r)\text{ else }\operatorname{zero}(r)\right)))
\end{equation*}


\noindent\emph{(Here $\operatorname{comb\text{-}kk}(r, x, y, 0)$ is the internal representation of the binomial coefficient term; it is used throughout the proof.)}
\par
Let $r$, $x$ and $y$ be arbitrary. Let $m \in \mathbb{N}$. Then
\begin{equation*}
(\operatorname{comb\text{-}kk}(r, x, y, \operatorname{succ}(m)) \simeq (\lambda k \in \mathbb{Z}.\; \operatorname{add}(r)(\operatorname{mul}(r)(x, \operatorname{comb\text{-}kk}(r, x, y, m)(-(k, 1))), \operatorname{mul}(r)(y, \operatorname{comb\text{-}kk}(r, x, y, m)(k)))))
\end{equation*}


\noindent\emph{(Here $\operatorname{comb\text{-}kk}(r, x, y, \operatorname{succ}(m))$ is the internal representation of the binomial coefficient term; it is used throughout the proof.)}


\vnbbookfile{structure-library/is-set-laws}{structure-library/is-set-laws}
\noindent{\small structure-library/is-set-laws.scm -- THE GENERATED SETHOOD THEOREMS (W1I, 2026-11-06).}\par
\vnbprop{is-semigroup-is-set}{\vnbstmt{\forall s.\; (\operatorname{is\text{-}semigroup}(s) \Rightarrow (s \in \mathbf{Set}))}}
\vnbprop{is-monoid-is-set}{\vnbstmt{\forall s.\; (\operatorname{is\text{-}monoid}(s) \Rightarrow (s \in \mathbf{Set}))}}
\vnbprop{is-comm-monoid-is-set}{\vnbstmt{\forall s.\; (\operatorname{is\text{-}comm\text{-}monoid}(s) \Rightarrow (s \in \mathbf{Set}))}}
\vnbprop{is-group-is-set}{\vnbstmt{\forall s.\; (\operatorname{is\text{-}group}(s) \Rightarrow (s \in \mathbf{Set}))}}
\vnbprop{is-abelian-group-is-set}{\vnbstmt{\forall s.\; (\operatorname{is\text{-}abelian\text{-}group}(s) \Rightarrow (s \in \mathbf{Set}))}}
\vnbprop{is-ring-is-set}{\vnbstmt{\forall s.\; (\operatorname{is\text{-}ring}(s) \Rightarrow (s \in \mathbf{Set}))}}
\vnbprop{is-setoid-is-set}{\vnbstmt{\forall s.\; (\operatorname{is\text{-}setoid}(s) \Rightarrow (s \in \mathbf{Set}))}}
\vnbprop{is-ringoid-is-set}{\vnbstmt{\forall s.\; (\operatorname{is\text{-}ringoid}(s) \Rightarrow (s \in \mathbf{Set}))}}
\vnbprop{is-commutative-ring-is-set}{\vnbstmt{\forall s.\; (\operatorname{is\text{-}commutative\text{-}ring}(s) \Rightarrow (s \in \mathbf{Set}))}}
\vnbprop{is-integral-domain-is-set}{\vnbstmt{\forall s.\; (\operatorname{is\text{-}integral\text{-}domain}(s) \Rightarrow (s \in \mathbf{Set}))}}
\vnbprop{is-field-is-set}{\vnbstmt{\forall s.\; (\operatorname{is\text{-}field}(s) \Rightarrow (s \in \mathbf{Set}))}}
\vnbprop{is-field-ring-is-set}{\vnbstmt{\forall s.\; (\operatorname{is\text{-}field\text{-}ring}(s) \Rightarrow (s \in \mathbf{Set}))}}
\vnbprop{is-euclidean-ring-is-set}{\vnbstmt{\forall s.\; (\operatorname{is\text{-}euclidean\text{-}ring}(s) \Rightarrow (s \in \mathbf{Set}))}}
\vnbprop{is-pid-is-set}{\vnbstmt{\forall s.\; (\operatorname{is\text{-}pid}(s) \Rightarrow (s \in \mathbf{Set}))}}

\vnbbookfile{theorem-library/poly-membership}{theorem-library/poly-membership}
\noindent{\small poly-membership.scm -- SUPP / FINSUPP / POLY read-offs, PROVEN.}\par
\vnbprop{supp-unfold}{\vnbstmt{\forall a, m, f.\; (\operatorname{supp}(a, m, f) \simeq \{x \in \operatorname{carr}(m) : \neg (f(x) = \operatorname{zero}(a))\})}}
\vnbprop{supp-membership}{\vnbstmt{\begin{array}{@{}l@{}} \forall a, m, f, x. \\ \quad (x \in \operatorname{supp}(a, m, f)) \Leftrightarrow ((x \in \operatorname{carr}(m)) \wedge \neg (f(x) = \operatorname{zero}(a))) \end{array}}}
\vnbprop{supp-in-set}{\vnbstmt{\forall a, m, f.\; ((\operatorname{carr}(m) \in \mathbf{Set}) \Rightarrow (\operatorname{supp}(a, m, f) \in \mathbf{Set}))}}
\vnbprop{finsupp-unfold}{\vnbstmt{\begin{array}{@{}l@{}} \forall a, m. \\ \quad (\operatorname{finsupp}(a, m) \simeq \{f \in (\operatorname{carr}(m) \to \operatorname{carr}(a)) : (|\operatorname{supp}(a, m, f)| \in \mathbb{N})\}) \end{array}}}
\vnbprop{finsupp-membership}{\vnbstmt{\begin{array}{@{}l@{}} \forall a, m, f. \\ \quad (f \in \operatorname{finsupp}(a, m)) \Leftrightarrow ((f \in (\operatorname{carr}(m) \to \operatorname{carr}(a))) \wedge (|\operatorname{supp}(a, m, f)| \in \mathbb{N})) \end{array}}}
\vnbprop{poly-carrier}{\vnbstmt{\forall a.\; (\operatorname{is\text{-}ring}(a) \Rightarrow (\operatorname{carr}(\operatorname{poly}(a)) \simeq \operatorname{finsupp}(a, \operatorname{nn\text{-}add\text{-}monoid})))}}
\vnbprop{poly-membership}{\vnbstmt{\begin{array}{@{}l@{}} \forall a, f. \\ \quad \operatorname{is\text{-}ring}(a) \Rightarrow  \\ \quad \quad (f \in \operatorname{carr}(\operatorname{poly}(a))) \Leftrightarrow ((f \in \operatorname{carr}(a)^{\mathbb{N}}) \wedge (|\operatorname{supp}(a, \operatorname{nn\text{-}add\text{-}monoid}, f)| \in \mathbb{N})) \end{array}}}

\vnbbookfile{theorem-library/ag-view-read-offs}{theorem-library/ag-view-read-offs}
\noindent{\small ag-view-read-offs.scm -- THE ADDITIVE-GROUP VIEW READ-OFFS, PROVEN.}\par
\vnbprop{ras-carr}{\vnbstmt{\forall a.\; (\operatorname{is\text{-}ring}(a) \Rightarrow (\operatorname{carr}(\operatorname{ring\text{-}additive\text{-}ag}(a)) = \operatorname{carr}(a)))}}
\vnbprop{ras-op}{\vnbstmt{\forall a.\; (\operatorname{is\text{-}ring}(a) \Rightarrow (\operatorname{opr}(\operatorname{ring\text{-}additive\text{-}ag}(a)) = \operatorname{add}(a)))}}
\vnbprop{ras-id}{\vnbstmt{\forall a.\; (\operatorname{is\text{-}ring}(a) \Rightarrow (\operatorname{iden}(\operatorname{ring\text{-}additive\text{-}ag}(a)) = \operatorname{zero}(a)))}}

\vnbbookfile{theorem-library/view-typing-laws}{theorem-library/view-typing-laws}
\noindent{\small theorem-library/view-typing-laws.scm -- THE VIEW TYPING THEOREMS, PROVEN (D8, 2026-11-15).}\par
\vnbprop{ringoid-as-ring-is-ring}{\vnbstmt{\forall r.\; (\operatorname{is\text{-}ringoid}(r) \Rightarrow \operatorname{is\text{-}ring}(\operatorname{ringoid\text{-}as\text{-}ring}(r)))}}
\vnbprop{ring-additive-ag-is-abelian-group}{\vnbstmt{\forall r.\; (\operatorname{is\text{-}ring}(r) \Rightarrow \operatorname{is\text{-}abelian\text{-}group}(\operatorname{ring\text{-}additive\text{-}ag}(r)))}}
\vnbprop{ring-multiplicative-monoid-is-monoid}{\vnbstmt{\forall r.\; (\operatorname{is\text{-}ring}(r) \Rightarrow \operatorname{is\text{-}monoid}(\operatorname{ring\text{-}multiplicative\text{-}monoid}(r)))}}
\vnbprop{abelian-group-as-monoid-is-monoid}{\vnbstmt{\forall r.\; (\operatorname{is\text{-}abelian\text{-}group}(r) \Rightarrow \operatorname{is\text{-}monoid}(\operatorname{abelian\text{-}group\text{-}as\text{-}monoid}(r)))}}
\vnbprop{commutative-ring-additive-ag-is-abelian-group}{\vnbstmt{\begin{array}{@{}l@{}} \forall r. \\ \quad \operatorname{is\text{-}commutative\text{-}ring}(r) \Rightarrow \operatorname{is\text{-}abelian\text{-}group}(\operatorname{commutative\text{-}ring\text{-}additive\text{-}ag}(r)) \end{array}}}
\vnbprop{field-additive-ag-is-abelian-group}{\vnbstmt{\forall r.\; (\operatorname{is\text{-}field}(r) \Rightarrow \operatorname{is\text{-}abelian\text{-}group}(\operatorname{field\text{-}additive\text{-}ag}(r)))}}
\vnbprop{commutative-ring-multiplicative-cm-is-comm-monoid}{\vnbstmt{\begin{array}{@{}l@{}} \forall r. \\ \quad \operatorname{is\text{-}commutative\text{-}ring}(r) \Rightarrow  \\ \quad \quad \operatorname{is\text{-}comm\text{-}monoid}(\operatorname{commutative\text{-}ring\text{-}multiplicative\text{-}cm}(r)) \end{array}}}

\vnbbookfile{theorem-library/rake-setoid}{theorem-library/rake-setoid}
\noindent{\small rake-setoid.scm -- BATCH F of the 2026-09-17 rake: setoids and carriers. Seven asserted leaves of the PSS, PROVEN.  Every statement here is its support's statement UNCHANGED.}\par
\vnbprop{class-unfold}{\vnbstmt{\forall s, a.\; (\operatorname{class}(s, a) \simeq \{b \in \operatorname{pts}(s) : \operatorname{related}(s, a, b)\})}}
\vnbprop{proj-unfold}{\vnbstmt{\forall s.\; (\operatorname{proj}(s) \simeq (\lambda a \in \operatorname{pts}(s).\; \operatorname{class}(s, a)))}}
\vnbprop{quotient-unfold}{\vnbstmt{\forall s.\; (\operatorname{quotient}(s) \simeq \operatorname{image}(\operatorname{proj}(s), \operatorname{pts}(s)))}}
\vnbprop{respects-unfold}{\vnbstmt{\begin{array}{@{}l@{}} \forall s, f. \\ \quad \operatorname{respects}(s, f) \Leftrightarrow \forall a, b \in \operatorname{pts}(s).\; (\operatorname{related}(s, a, b) \Rightarrow (f(a) = f(b))) \end{array}}}
\vnbprop{respects2-unfold}{\vnbstmt{\begin{array}{@{}l@{}} \forall s, f. \\ \quad \operatorname{respects2}(s, f) \\ \quad \Leftrightarrow\; \forall a, b, c, d. \\ \quad \quad a \in \operatorname{pts}(s) \wedge \\ \quad \quad b \in \operatorname{pts}(s) \wedge \\ \quad \quad c \in \operatorname{pts}(s) \wedge \\ \quad \quad d \in \operatorname{pts}(s) \wedge \\ \quad \quad \operatorname{related}(s, a, c) \wedge \\ \quad \quad \operatorname{related}(s, b, d) \Rightarrow \\ \quad \quad \quad f(a, b) = f(c, d) \end{array}}}
\vnbprop{class-is-set}{\vnbstmt{\forall s.; \forall a \in \operatorname{pts}(s).\; (\operatorname{is\text{-}setoid}(s) \Rightarrow (\operatorname{class}(s, a) \in \mathbf{Set}))}}
\vnbprop{quotient-is-set}{\vnbstmt{\forall s.\; (\operatorname{is\text{-}setoid}(s) \Rightarrow (\operatorname{quotient}(s) \in \mathbf{Set}))}}
\vnbprop{proj-in-fun}{\vnbstmt{\forall s.\; (\operatorname{is\text{-}setoid}(s) \Rightarrow (\operatorname{proj}(s) \in (\operatorname{pts}(s) \to \operatorname{quotient}(s))))}}
\vnbprop{descend-in-fun}{\vnbstmt{\begin{array}{@{}l@{}} \forall s, z.; \forall f \in (\operatorname{pts}(s) \to z). \\ \quad (\operatorname{is\text{-}setoid}(s) \wedge \operatorname{respects}(s, f)) \Rightarrow (\operatorname{descend}(s, f) \in (\operatorname{quotient}(s) \to z)) \end{array}}}
\vnbprop{descend2-in-fun}{\vnbstmt{\begin{array}{@{}l@{}} \forall s, z.; \forall f \in ((\operatorname{pts}(s) \times \operatorname{pts}(s)) \to z). \\ \quad (\operatorname{is\text{-}setoid}(s) \wedge \operatorname{respects2}(s, f)) \Rightarrow  \\ \quad \quad \operatorname{descend2}(s, f) \in ((\operatorname{quotient}(s) \times \operatorname{quotient}(s)) \to z) \end{array}}}

\vnbbookfile{theorem-library/bt-shims}{theorem-library/bt-shims}
\noindent{\small bt-shims.scm -- the five "Bernstein typing" shims of theorem-library/binomial.scm, PROVEN.}\par
\vnbprop{bt-nn-in-zz}{\vnbstmt{\forall n \in \mathbb{N}.\; (n \in \mathbb{Z})}}
\vnbprop{bt-succ-in-nn}{\vnbstmt{\forall n \in \mathbb{N}.\; ((n + 1) \in \mathbb{N})}}
\vnbprop{bt-neg1-in-zz}{\vnbstmt{-1 \in \mathbb{Z}}}
\vnbprop{bt-lt-succ}{\vnbstmt{\forall n \in \mathbb{N}.\; (n < (n + 1))}}
\vnbprop{bt-neg1-neg}{\vnbstmt{-1 < 0}}

\vnbbookfile{structure-library/subtype-laws}{structure-library/subtype-laws}
\noindent{\small subtype-laws.scm -- the trivial subtype-subsumption laws, PROVEN.}\par
\vnbprop{abelian-group-is-group}{\vnbstmt{\forall s.\; (\operatorname{is\text{-}abelian\text{-}group}(s) \Rightarrow \operatorname{is\text{-}group}(s))}}
\vnbprop{commutative-ring-is-ring}{\vnbstmt{\forall s.\; (\operatorname{is\text{-}commutative\text{-}ring}(s) \Rightarrow \operatorname{is\text{-}ring}(s))}}
\vnbprop{integral-domain-is-commutative-ring}{\vnbstmt{\forall s.\; (\operatorname{is\text{-}integral\text{-}domain}(s) \Rightarrow \operatorname{is\text{-}commutative\text{-}ring}(s))}}
\vnbprop{euclidean-ring-is-integral-domain}{\vnbstmt{\forall s.\; (\operatorname{is\text{-}euclidean\text{-}ring}(s) \Rightarrow \operatorname{is\text{-}integral\text{-}domain}(s))}}
\vnbprop{comm-monoid-is-monoid}{\vnbstmt{\forall s.\; (\operatorname{is\text{-}comm\text{-}monoid}(s) \Rightarrow \operatorname{is\text{-}monoid}(s))}}
\vnbprop{abelian-group-opr-comm}{\vnbstmt{\begin{array}{@{}l@{}} \forall s.; \forall a, b \in \operatorname{carr}(s). \\ \quad \operatorname{is\text{-}abelian\text{-}group}(s) \Rightarrow (\operatorname{opr}(s)(a, b) = \operatorname{opr}(s)(b, a)) \end{array}}}
\vnbprop{group-assoc}{\vnbstmt{\begin{array}{@{}l@{}} \forall s.; \forall a, b, c \in \operatorname{carr}(s). \\ \quad \operatorname{is\text{-}group}(s) \Rightarrow (\operatorname{opr}(s)(\operatorname{opr}(s)(a, b), c) = \operatorname{opr}(s)(a, \operatorname{opr}(s)(b, c))) \end{array}}}
\vnbprop{group-left-id}{\vnbstmt{\forall s.; \forall a \in \operatorname{carr}(s).\; (\operatorname{is\text{-}group}(s) \Rightarrow (\operatorname{opr}(s)(\operatorname{iden}(s), a) = a))}}
\vnbprop{group-left-inv}{\vnbstmt{\begin{array}{@{}l@{}} \forall s.; \forall a \in \operatorname{carr}(s). \\ \quad \operatorname{is\text{-}group}(s) \Rightarrow (\operatorname{opr}(s)(\operatorname{inv}(s)(a), a) = \operatorname{iden}(s)) \end{array}}}
\vnbprop{group-identity-in}{\vnbstmt{\forall s.\; (\operatorname{is\text{-}group}(s) \Rightarrow (\operatorname{iden}(s) \in \operatorname{carr}(s)))}}

\vnbbookfile{theorem-library/op-typing}{theorem-library/op-typing}
\noindent{\small A file shared across books; the entries that belong to this one.}\par
\noindent{\small op-typing.scm -- the codomain typing of a structure operation, in APPLIED form, PROVEN.  One driver, seven theorems.}\par
\vnbprop{ring-add-closed}{\vnbstmt{\begin{array}{@{}l@{}} \forall s.; \forall a, b \in \operatorname{carr}(s). \\ \quad \operatorname{is\text{-}ring}(s) \Rightarrow (\operatorname{add}(s)(a, b) \in \operatorname{carr}(s)) \end{array}}}
\vnbprop{ring-carrier-closed-mul}{\vnbstmt{\begin{array}{@{}l@{}} \forall r.; \forall a, b \in \operatorname{carr}(r). \\ \quad \operatorname{is\text{-}ring}(r) \Rightarrow (\operatorname{mul}(r)(a, b) \in \operatorname{carr}(r)) \end{array}}}
\vnbprop{ring-neg-in-carr}{\vnbstmt{\forall s.; \forall a \in \operatorname{carr}(s).\; (\operatorname{is\text{-}ring}(s) \Rightarrow (\operatorname{neg}(s)(a) \in \operatorname{carr}(s)))}}
\vnbprop{bt-add-in-carr}{\vnbstmt{\begin{array}{@{}l@{}} \forall r.; \forall a, b \in \operatorname{carr}(r). \\ \quad \operatorname{is\text{-}commutative\text{-}ring}(r) \Rightarrow (\operatorname{add}(r)(a, b) \in \operatorname{carr}(r)) \end{array}}}

\vnbbookfile{theorem-library/finsum-type-proof}{theorem-library/finsum-type-proof}
\noindent{\small theorem-library/finsum-type-proof.scm -- the FINSUM typing floor, PROVEN.}\par
\vnbprop{group-carrier-closed-opr}{\vnbstmt{\begin{array}{@{}l@{}} \forall s.; \forall a, b \in \operatorname{carr}(s). \\ \quad \operatorname{is\text{-}group}(s) \Rightarrow (\operatorname{opr}(s)(a, b) \in \operatorname{carr}(s)) \end{array}}}
\vnbprop{sum-ag-type-ind}{\vnbstmt{\begin{array}{@{}l@{}} \forall n \in \mathbb{N}, g.; \forall f \in (\mathbb{N} \to \operatorname{carr}(g)). \\ \quad \operatorname{is\text{-}abelian\text{-}group}(g) \Rightarrow (\operatorname{sum\text{-}ag}(g, f, n) \in \operatorname{carr}(g)) \end{array}}}
\vnbprop{sum-ag-type}{\vnbstmt{\begin{array}{@{}l@{}} \forall g.; \forall f \in (\mathbb{N} \to \operatorname{carr}(g)), n \in \mathbb{N}. \\ \quad \operatorname{is\text{-}abelian\text{-}group}(g) \Rightarrow (\operatorname{sum\text{-}ag}(g, f, n) \in \operatorname{carr}(g)) \end{array}}}
\vnbprop{enum-fam-in-fun}{\vnbstmt{\begin{array}{@{}l@{}} \forall n \in \mathbb{N}, g, s.; \forall i \in (\mathrm{OS}(n) \to s), f \in (s \to \operatorname{carr}(g)). \\ \quad \operatorname{is\text{-}group}(g) \Rightarrow (\operatorname{enum\text{-}fam}(g, f, i, n) \in (\mathbb{N} \to \operatorname{carr}(g))) \end{array}}}
\vnbprop{finsum-type}{\vnbstmt{\begin{array}{@{}l@{}} \forall g.; \forall s \in \mathbf{Set}, f \in (s \to \operatorname{carr}(g)). \\ \quad (\operatorname{is\text{-}abelian\text{-}group}(g) \wedge (|s| \in \mathbb{N})) \Rightarrow (\operatorname{finsum}(g, f, s) \in \operatorname{carr}(g)) \end{array}}}
\vnbprop{sum-ag-all-id-ind}{\vnbstmt{\begin{array}{@{}l@{}} \forall n \in \mathbb{N}, g, a. \\ \quad (\operatorname{is\text{-}abelian\text{-}group}(g) \wedge \forall i \in \mathrm{OS}(n).\; (a(i) = \operatorname{iden}(g))) \Rightarrow  \\ \quad \quad \operatorname{sum\text{-}ag}(g, a, n) = \operatorname{iden}(g) \end{array}}}
\vnbprop{finsum-all-id}{\vnbstmt{\begin{array}{@{}l@{}} \forall g.; \forall s \in \mathbf{Set}, f \in (s \to \operatorname{carr}(g)). \\ \quad (\operatorname{is\text{-}abelian\text{-}group}(g) \wedge ((|s| \in \mathbb{N}) \wedge \forall z \in s.\; (f(z) = \operatorname{iden}(g)))) \Rightarrow  \\ \quad \quad \operatorname{finsum}(g, f, s) = \operatorname{iden}(g) \end{array}}}

\vnbbookfile{theorem-library/finsum-single-support}{theorem-library/finsum-single-support}
\noindent{\small finsum-single-support.scm -- a finite sum whose summand vanishes off a single index equals its value there.  Both theorems `modulo 0'.}\par
\vnbprop{sum-ag-single-ind}{\vnbstmt{\begin{array}{@{}l@{}} \forall n \in \mathbb{N}, g, a.; \forall p \in \mathrm{OS}(n). \\ \quad \operatorname{is\text{-}abelian\text{-}group}(g) \wedge \\ \quad a(p) \in \operatorname{carr}(g) \wedge \\ \quad \forall i \in \mathrm{OS}(n).\; (\neg (i = p) \Rightarrow (a(i) = \operatorname{iden}(g))) \Rightarrow \\ \quad \quad \operatorname{sum\text{-}ag}(g, a, n) = a(p) \end{array}}}
\vnbprop{finsum-single-support}{\vnbstmt{\begin{array}{@{}l@{}} \forall g.; \forall s \in \mathbf{Set}, f \in (s \to \operatorname{carr}(g)), a \in s. \\ \quad (\operatorname{is\text{-}abelian\text{-}group}(g) \wedge ((|s| \in \mathbb{N}) \wedge \forall j \in s.\; (\neg (j = a) \Rightarrow (f(j) = \operatorname{iden}(g))))) \Rightarrow  \\ \quad \quad \operatorname{finsum}(g, f, s) = f(a) \end{array}}}

\vnbbookfile{structure-library/integral-domain-laws}{structure-library/integral-domain-laws}
\noindent{\small integral-domain-laws.scm -- integral-domain-cancel-zero, PROVEN from the definition.}\par
\vnbprop{integral-domain-cancel-zero}{\vnbstmt{\begin{array}{@{}l@{}} \forall s.; \forall a, b \in \operatorname{carr}(s). \\ \quad (\operatorname{is\text{-}integral\text{-}domain}(s) \wedge ((\operatorname{mul}(s)(a, b) = \operatorname{zero}(s)) \wedge \neg (b = \operatorname{zero}(s)))) \Rightarrow  \\ \quad \quad a = \operatorname{zero}(s) \end{array}}}
\vnbprop{integral-domain-nontrivial}{\vnbstmt{\forall s.\; (\operatorname{is\text{-}integral\text{-}domain}(s) \Rightarrow \neg (\operatorname{one}(s) = \operatorname{zero}(s)))}}

\vnbbookfile{theorem-library/zz-ring-is-ring}{theorem-library/zz-ring-is-ring}
\noindent{\small zz-ring-is-ring.scm -- IS-RING(ZZ-RING), PROVED.}\par
\vnbprop{zz-is-ring}{\vnbstmt{\operatorname{is\text{-}ring}(\operatorname{zz\text{-}ring})}}
\vnbprop{qq-is-ring}{\vnbstmt{\operatorname{is\text{-}ring}(\operatorname{qq\text{-}ring})}}

\vnbbookfile{theorem-library/zz-integral-domain}{theorem-library/zz-integral-domain}
\noindent{\small zz-integral-domain.scm -- IS-COMMUTATIVE-RING(ZZ-RING) and IS-INTEGRAL-DOMAIN(ZZ-RING), PROVED.}\par
\vnbprop{zz-is-commutative-ring}{\vnbstmt{\operatorname{is\text{-}commutative\text{-}ring}(\operatorname{zz\text{-}ring})}}
\vnbprop{zz-no-zero-divisors}{\vnbstmt{\forall a, b \in \mathbb{Z}.\; (((a \cdot b) = 0) \Rightarrow ((a = 0) \vee (b = 0)))}}
\vnbprop{zz-is-integral-domain}{\vnbstmt{\operatorname{is\text{-}integral\text{-}domain}(\operatorname{zz\text{-}ring})}}

\vnbbookfile{theorem-library/nn-add-monoid}{theorem-library/nn-add-monoid}
\noindent{\small nn-add-monoid.scm -- IS-COMM-MONOID(NN-ADD-MONOID), PROVED.}\par
\vnbprop{nn-add-monoid-is-comm-monoid}{\vnbstmt{\operatorname{is\text{-}comm\text{-}monoid}(\operatorname{nn\text{-}add\text{-}monoid})}}

\vnbbookfile{theorem-library/cancellation}{theorem-library/cancellation}
\noindent{\small cancellation.scm -- cancellation, proved ONCE where it lives, then carried.}\par
\vnbprop{group-inv-in}{\vnbstmt{\forall s.; \forall a \in \operatorname{carr}(s).\; (\operatorname{is\text{-}group}(s) \Rightarrow (\operatorname{inv}(s)(a) \in \operatorname{carr}(s)))}}
\vnbprop{group-cancel-left}{\vnbstmt{\begin{array}{@{}l@{}} \forall s.; \forall a, b, c \in \operatorname{carr}(s). \\ \quad (\operatorname{is\text{-}group}(s) \wedge (\operatorname{opr}(s)(c, a) = \operatorname{opr}(s)(c, b))) \Rightarrow (a = b) \end{array}}}
\vnbprop{ag-cancel-right}{\vnbstmt{\begin{array}{@{}l@{}} \forall s.; \forall a, b, c \in \operatorname{carr}(s). \\ \quad (\operatorname{is\text{-}abelian\text{-}group}(s) \wedge (\operatorname{opr}(s)(a, c) = \operatorname{opr}(s)(b, c))) \Rightarrow (a = b) \end{array}}}
\vnbprop{abelian-group-idempotent-is-id}{\vnbstmt{\begin{array}{@{}l@{}} \forall s.; \forall a \in \operatorname{carr}(s). \\ \quad (\operatorname{is\text{-}abelian\text{-}group}(s) \wedge (\operatorname{opr}(s)(a, a) = a)) \Rightarrow (a = \operatorname{iden}(s)) \end{array}}}
\vnbprop{zz-add-cancel}{\vnbstmt{\forall a, b, c \in \mathbb{Z}.\; (((a + c) = (b + c)) \Rightarrow (a = b))}}
\vnbprop{qq-add-cancel}{\vnbstmt{\forall a, b, c \in \mathbb{Q}.\; (((a + c) = (b + c)) \Rightarrow (a = b))}}

\vnbbookfile{theorem-library/ring-zero-one-power}{theorem-library/ring-zero-one-power}
\noindent{\small ring-zero-one-power.scm -- four ring facts that were asserted, PROVEN.}\par
\vnbprop{ring-one-in}{\vnbstmt{\forall s.\; (\operatorname{is\text{-}ring}(s) \Rightarrow (\operatorname{one}(s) \in \operatorname{carr}(s)))}}
\vnbprop{ring-mul-zero-left}{\vnbstmt{\forall s.; \forall a \in \operatorname{carr}(s).\; (\operatorname{is\text{-}ring}(s) \Rightarrow (\operatorname{mul}(s)(\operatorname{zero}(s), a) = \operatorname{zero}(s)))}}
\vnbprop{ring-mul-zero-right}{\vnbstmt{\forall s.; \forall a \in \operatorname{carr}(s).\; (\operatorname{is\text{-}ring}(s) \Rightarrow (\operatorname{mul}(s)(a, \operatorname{zero}(s)) = \operatorname{zero}(s)))}}
\vnbprop{ring-power-zero}{\vnbstmt{\forall r, x.\; (\operatorname{is\text{-}ring}(r) \Rightarrow (\operatorname{ring\text{-}power}(r, x, 0) \simeq \operatorname{one}(r)))}}

\vnbbookfile{theorem-library/nn-integral}{theorem-library/nn-integral}
\noindent{\small nn-integral.scm -- the NN arithmetic support layer.}\par
\vnbprop{zz-mul-cancel-zero}{\vnbstmt{\forall a, b \in \mathbb{Z}.\; ((((a \cdot b) = 0) \wedge \neg (b = 0)) \Rightarrow (a = 0))}}
\vnbprop{nn-mul-nonzero}{\vnbstmt{\forall a, b \in \mathbb{N}.\; ((\neg (a = 0) \wedge \neg (b = 0)) \Rightarrow \neg ((a \cdot b) = 0))}}
\vnbprop{nn-2-cancel}{\vnbstmt{\forall x, y \in \mathbb{N}.\; (((2 \cdot x) = (2 \cdot y)) \Rightarrow (x = y))}}
\vnbprop{nn-lt-double}{\vnbstmt{\forall k \in \mathbb{N}.\; (\neg (k = 0) \Rightarrow (k < (2 \cdot k)))}}

\vnbbookfile{theorem-library/comb-kk-laws}{theorem-library/comb-kk-laws}
\noindent{\small comb-kk-laws.scm -- the six COMB-KK / Bernstein plumbing supports of theorem-library/binomial.scm, PROVEN.}\par
\vnbprop{sum-in-carr-zz-ind}{\vnbstmt{\begin{array}{@{}l@{}} \forall n \in \mathbb{N}, r.; \forall g \in (\mathbb{Z} \to \operatorname{carr}(r)). \\ \quad \operatorname{is\text{-}commutative\text{-}ring}(r) \Rightarrow (\operatorname{sum}(r, g, n) \in \operatorname{carr}(r)) \end{array}}}
\vnbprop{bt-sum-in-carr-zz}{\vnbstmt{\begin{array}{@{}l@{}} \forall r.; \forall g \in (\mathbb{Z} \to \operatorname{carr}(r)), n \in \mathbb{N}. \\ \quad \operatorname{is\text{-}commutative\text{-}ring}(r) \Rightarrow (\operatorname{sum}(r, g, n) \in \operatorname{carr}(r)) \end{array}}}
\vnbprop{comb-kk-0-0}{\vnbstmt{\begin{array}{@{}l@{}} \forall r.; \forall x, y \in \operatorname{carr}(r). \\ \quad \operatorname{is\text{-}commutative\text{-}ring}(r) \Rightarrow (\operatorname{comb\text{-}kk}(r, x, y, 0)(0) = \operatorname{one}(r)) \end{array}}}
\vnbprop{comb-kk-in-fun-ind}{\vnbstmt{\begin{array}{@{}l@{}} \forall m \in \mathbb{N}, r.; \forall x, y \in \operatorname{carr}(r). \\ \quad \operatorname{is\text{-}commutative\text{-}ring}(r) \Rightarrow (\operatorname{comb\text{-}kk}(r, x, y, m) \in (\mathbb{Z} \to \operatorname{carr}(r))) \end{array}}}
\vnbprop{comb-kk-in-fun}{\vnbstmt{\begin{array}{@{}l@{}} \forall r.; \forall x, y \in \operatorname{carr}(r), m \in \mathbb{N}. \\ \quad \operatorname{is\text{-}commutative\text{-}ring}(r) \Rightarrow (\operatorname{comb\text{-}kk}(r, x, y, m) \in (\mathbb{Z} \to \operatorname{carr}(r))) \end{array}}}
\vnbprop{comb-kk-null-ind}{\vnbstmt{\begin{array}{@{}l@{}} \forall m \in \mathbb{N}, r.; \forall x, y \in \operatorname{carr}(r), k \in \mathbb{Z}. \\ \quad (\operatorname{is\text{-}commutative\text{-}ring}(r) \wedge (k < 0)) \Rightarrow (\operatorname{comb\text{-}kk}(r, x, y, m)(k) = \operatorname{zero}(r)) \end{array}}}
\vnbprop{comb-kk-null}{\vnbstmt{\begin{array}{@{}l@{}} \forall r.; \forall x, y \in \operatorname{carr}(r), m \in \mathbb{N}, k \in \mathbb{Z}. \\ \quad (\operatorname{is\text{-}commutative\text{-}ring}(r) \wedge (k < 0)) \Rightarrow (\operatorname{comb\text{-}kk}(r, x, y, m)(k) = \operatorname{zero}(r)) \end{array}}}
\vnbprop{comb-kk-above-ind}{\vnbstmt{\begin{array}{@{}l@{}} \forall m \in \mathbb{N}, r.; \forall x, y \in \operatorname{carr}(r), k \in \mathbb{Z}. \\ \quad (\operatorname{is\text{-}commutative\text{-}ring}(r) \wedge (m < k)) \Rightarrow (\operatorname{comb\text{-}kk}(r, x, y, m)(k) = \operatorname{zero}(r)) \end{array}}}
\vnbprop{comb-kk-above}{\vnbstmt{\begin{array}{@{}l@{}} \forall r.; \forall x, y \in \operatorname{carr}(r), m \in \mathbb{N}, k \in \mathbb{Z}. \\ \quad (\operatorname{is\text{-}commutative\text{-}ring}(r) \wedge (m < k)) \Rightarrow (\operatorname{comb\text{-}kk}(r, x, y, m)(k) = \operatorname{zero}(r)) \end{array}}}

\vnbbookfile{theorem-library/zz-division}{theorem-library/zz-division}
\noindent{\small zz-division.scm -- DIVISION WITH REMAINDER, and ZZ as a Euclidean ring.}\par
\vnbprop{nn-division}{\vnbstmt{\begin{array}{@{}l@{}} \forall a, b \in \mathbb{N}. \\ \quad \neg (b = 0) \Rightarrow \exists q, r \in \mathbb{N}.\; ((a = ((q \cdot b) + r)) \wedge ((r + 1) \leq b)) \end{array}}}
\vnbprop{zz-division}{\vnbstmt{\begin{array}{@{}l@{}} \forall a, b \in \mathbb{Z}. \\ \quad \neg (b = 0) \Rightarrow  \\ \quad \quad \exists q, r \in \mathbb{Z}.\; ((a = ((q \cdot b) + r)) \wedge ((r = 0) \vee ((\lvert r \rvert + 1) \leq \lvert b \rvert))) \end{array}}}
\vnbprop{zz-is-euclidean-ring}{\vnbstmt{\operatorname{is\text{-}euclidean\text{-}ring}(\operatorname{zz\text{-}ring})}}

\vnbbookfile{theorem-library/rake-algebra}{theorem-library/rake-algebra}
\noindent{\small rake-algebra.scm -- BATCH D of the 2026-09-17 rake: seven asserted leaves of the PSS, PROVEN.  Every statement here is its support's statement UNCHANGED.}\par
\vnbprop{ring-carr-in-set}{\vnbstmt{\forall r.\; (\operatorname{is\text{-}ring}(r) \Rightarrow (\operatorname{carr}(r) \in \mathbf{Set}))}}
\vnbprop{ideal-elt-in-carrier}{\vnbstmt{\forall s, i, x.\; ((\operatorname{is\text{-}ideal}(s, i) \wedge (x \in i)) \Rightarrow (x \in \operatorname{carr}(s)))}}
\vnbprop{principal-ideal-in-ideal}{\vnbstmt{\begin{array}{@{}l@{}} \forall s, i, b.; \forall x \in \operatorname{principal\text{-}ideal}(s, b). \\ \quad (\operatorname{is\text{-}ideal}(s, i) \wedge (b \in i)) \Rightarrow (x \in i) \end{array}}}
\vnbprop{ring-add-right-id}{\vnbstmt{\forall s.; \forall a \in \operatorname{carr}(s).\; (\operatorname{is\text{-}ring}(s) \Rightarrow (\operatorname{add}(s)(a, \operatorname{zero}(s)) = a))}}

\vnbbookfile{theorem-library/rake-algebra2}{theorem-library/rake-algebra2}
\noindent{\small rake-algebra2.scm -- BATCH I of the 2026-09-17 rake: the algebra two-liners and the four unwarranted algebra axioms, PROVEN.  Every statement is its site's statement UNCHANGED (extracted from the site, not retyped).}\par
\vnbprop{abelian-group-assoc}{\vnbstmt{\begin{array}{@{}l@{}} \forall s.; \forall a, b, c \in \operatorname{carr}(s). \\ \quad \operatorname{is\text{-}abelian\text{-}group}(s) \Rightarrow  \\ \quad \quad \operatorname{opr}(s)(\operatorname{opr}(s)(a, b), c) = \operatorname{opr}(s)(a, \operatorname{opr}(s)(b, c)) \end{array}}}
\vnbprop{abelian-group-right-id}{\vnbstmt{\forall s.; \forall a \in \operatorname{carr}(s).\; (\operatorname{is\text{-}abelian\text{-}group}(s) \Rightarrow (\operatorname{opr}(s)(a, \operatorname{iden}(s)) = a))}}
\vnbprop{abelian-group-inverse-unique}{\vnbstmt{\begin{array}{@{}l@{}} \forall s.; \forall a, b \in \operatorname{carr}(s). \\ \quad (\operatorname{is\text{-}abelian\text{-}group}(s) \wedge (\operatorname{opr}(s)(a, b) = \operatorname{iden}(s))) \Rightarrow (b = \operatorname{inv}(s)(a)) \end{array}}}
\vnbprop{ring-add-right-inv}{\vnbstmt{\forall s.; \forall a \in \operatorname{carr}(s).\; (\operatorname{is\text{-}ring}(s) \Rightarrow (\operatorname{add}(s)(a, \operatorname{neg}(s)(a)) = \operatorname{zero}(s)))}}
\vnbprop{ring-add-inverse-unique}{\vnbstmt{\begin{array}{@{}l@{}} \forall s.; \forall a, b \in \operatorname{carr}(s). \\ \quad (\operatorname{is\text{-}ring}(s) \wedge (\operatorname{add}(s)(a, b) = \operatorname{zero}(s))) \Rightarrow (b = \operatorname{neg}(s)(a)) \end{array}}}
\vnbprop{ring-neg-neg}{\vnbstmt{\forall s.; \forall r \in \operatorname{carr}(s).\; (\operatorname{is\text{-}ring}(s) \Rightarrow (\operatorname{neg}(s)(\operatorname{neg}(s)(r)) = r))}}
\vnbprop{ring-neg-mul-left}{\vnbstmt{\begin{array}{@{}l@{}} \forall s.; \forall a, b \in \operatorname{carr}(s). \\ \quad \operatorname{is\text{-}ring}(s) \Rightarrow (\operatorname{mul}(s)(\operatorname{neg}(s)(a), b) = \operatorname{neg}(s)(\operatorname{mul}(s)(a, b))) \end{array}}}
\vnbprop{matrix-sethood}{\vnbstmt{\forall s \in \mathbf{Set}.\; (\operatorname{matrix}(s) \in \mathbf{Set})}}
\vnbprop{monoid-identity-in}{\vnbstmt{\forall m.\; (\operatorname{is\text{-}monoid}(m) \Rightarrow (\operatorname{iden}(m) \in \operatorname{carr}(m)))}}
\vnbprop{monoid-carrier-closed-opr}{\vnbstmt{\begin{array}{@{}l@{}} \forall m, a, b. \\ \quad (\operatorname{is\text{-}monoid}(m) \wedge ((a \in \operatorname{carr}(m)) \wedge (b \in \operatorname{carr}(m)))) \Rightarrow (\operatorname{opr}(m)(a, b) \in \operatorname{carr}(m)) \end{array}}}
\vnbprop{mpow-type-ind}{\vnbstmt{\begin{array}{@{}l@{}} \forall n \in \mathbb{N}, m.; \forall x \in \operatorname{carr}(m). \\ \quad \operatorname{is\text{-}monoid}(m) \Rightarrow (\operatorname{mpow}(m, x, n) \in \operatorname{carr}(m)) \end{array}}}
\vnbprop{mpow-type}{\vnbstmt{\begin{array}{@{}l@{}} \forall m.; \forall x \in \operatorname{carr}(m), n \in \mathbb{N}. \\ \quad \operatorname{is\text{-}monoid}(m) \Rightarrow (\operatorname{mpow}(m, x, n) \in \operatorname{carr}(m)) \end{array}}}

\vnbbookfile{theorem-library/rake-monoid-laws}{theorem-library/rake-monoid-laws}
\noindent{\small rake-monoid-laws.scm -- BATCH 5c-P of the 2026-09-18 rake: the three MONOID shape projections, PROVEN.}\par
\vnbprop{monoid-assoc}{\vnbstmt{\begin{array}{@{}l@{}} \forall s.; \forall a, b, c \in \operatorname{carr}(s). \\ \quad \operatorname{is\text{-}monoid}(s) \Rightarrow (\operatorname{opr}(s)(\operatorname{opr}(s)(a, b), c) = \operatorname{opr}(s)(a, \operatorname{opr}(s)(b, c))) \end{array}}}
\vnbprop{monoid-left-id}{\vnbstmt{\forall s.; \forall a \in \operatorname{carr}(s).\; (\operatorname{is\text{-}monoid}(s) \Rightarrow (\operatorname{opr}(s)(\operatorname{iden}(s), a) = a))}}
\vnbprop{monoid-right-id}{\vnbstmt{\forall s.; \forall a \in \operatorname{carr}(s).\; (\operatorname{is\text{-}monoid}(s) \Rightarrow (\operatorname{opr}(s)(a, \operatorname{iden}(s)) = a))}}

\vnbbookfile{theorem-library/rake-zz-act}{theorem-library/rake-zz-act}
\noindent{\small rake-zz-act.scm -- BATCH 6 of the rake: ZZ-ACT made a DEFINITION, and the five billed ZZ-ACT laws proven from it, each `modulo 0'.}\par
\vnbprop{ag-as-monoid-carr}{\vnbstmt{\begin{array}{@{}l@{}} \forall g. \\ \quad \operatorname{is\text{-}abelian\text{-}group}(g) \Rightarrow (\operatorname{carr}(g) \simeq \operatorname{carr}(\operatorname{abelian\text{-}group\text{-}as\text{-}monoid}(g))) \end{array}}}
\vnbprop{ag-as-monoid-iden}{\vnbstmt{\begin{array}{@{}l@{}} \forall g. \\ \quad \operatorname{is\text{-}abelian\text{-}group}(g) \Rightarrow (\operatorname{iden}(g) \simeq \operatorname{iden}(\operatorname{abelian\text{-}group\text{-}as\text{-}monoid}(g))) \end{array}}}
\vnbprop{ag-as-monoid-opr}{\vnbstmt{\begin{array}{@{}l@{}} \forall g. \\ \quad \operatorname{is\text{-}abelian\text{-}group}(g) \Rightarrow (\operatorname{opr}(g) \simeq \operatorname{opr}(\operatorname{abelian\text{-}group\text{-}as\text{-}monoid}(g))) \end{array}}}
\vnbprop{abelian-group-right-inv}{\vnbstmt{\begin{array}{@{}l@{}} \forall s.; \forall a \in \operatorname{carr}(s). \\ \quad \operatorname{is\text{-}abelian\text{-}group}(s) \Rightarrow (\operatorname{opr}(s)(a, \operatorname{inv}(s)(a)) = \operatorname{iden}(s)) \end{array}}}
\vnbprop{abelian-group-inv-iden}{\vnbstmt{\forall g.\; (\operatorname{is\text{-}abelian\text{-}group}(g) \Rightarrow (\operatorname{inv}(g)(\operatorname{iden}(g)) = \operatorname{iden}(g)))}}
\vnbprop{abelian-group-inv-inv}{\vnbstmt{\forall s.; \forall a \in \operatorname{carr}(s).\; (\operatorname{is\text{-}abelian\text{-}group}(s) \Rightarrow (\operatorname{inv}(s)(\operatorname{inv}(s)(a)) = a))}}
\vnbprop{abelian-group-inv-opr}{\vnbstmt{\begin{array}{@{}l@{}} \forall s.; \forall a, b \in \operatorname{carr}(s). \\ \quad \operatorname{is\text{-}abelian\text{-}group}(s) \Rightarrow  \\ \quad \quad \operatorname{inv}(s)(\operatorname{opr}(s)(a, b)) = \operatorname{opr}(s)(\operatorname{inv}(s)(a), \operatorname{inv}(s)(b)) \end{array}}}
\vnbprop{mpow-one}{\vnbstmt{\forall m.; \forall x \in \operatorname{carr}(m).\; (\operatorname{is\text{-}monoid}(m) \Rightarrow (\operatorname{mpow}(m, x, 1) = x))}}
\vnbprop{mpow-add-ind}{\vnbstmt{\begin{array}{@{}l@{}} \forall j \in \mathbb{N}, m.; \forall x \in \operatorname{carr}(m), k \in \mathbb{N}. \\ \quad \operatorname{is\text{-}monoid}(m) \Rightarrow (\operatorname{mpow}(m, x, (j + k)) = \operatorname{opr}(m)(\operatorname{mpow}(m, x, j), \operatorname{mpow}(m, x, k))) \end{array}}}
\vnbprop{mpow-add}{\vnbstmt{\begin{array}{@{}l@{}} \forall m.; \forall x \in \operatorname{carr}(m), j, k \in \mathbb{N}. \\ \quad \operatorname{is\text{-}monoid}(m) \Rightarrow (\operatorname{mpow}(m, x, (j + k)) = \operatorname{opr}(m)(\operatorname{mpow}(m, x, j), \operatorname{mpow}(m, x, k))) \end{array}}}
\vnbprop{ag-mpow-diff-nonneg}{\vnbstmt{\begin{array}{@{}l@{}} \forall g.; \forall d, n, m \in \mathbb{N}, a \in \operatorname{carr}(g). \\ \quad (\operatorname{is\text{-}abelian\text{-}group}(g) \wedge ((d + n) = m)) \Rightarrow  \\ \quad \quad \operatorname{opr}(g)(\operatorname{mpow}(\operatorname{abelian\text{-}group\text{-}as\text{-}monoid}(g), a, m), \operatorname{inv}(g)(\operatorname{mpow}(\operatorname{abelian\text{-}group\text{-}as\text{-}monoid}(g), a, n))) \\ \quad \quad =\; \operatorname{mpow}(\operatorname{abelian\text{-}group\text{-}as\text{-}monoid}(g), a, d) \end{array}}}
\vnbprop{ag-mpow-diff-neg}{\vnbstmt{\begin{array}{@{}l@{}} \forall g.; \forall p, m, n \in \mathbb{N}, a \in \operatorname{carr}(g). \\ \quad (\operatorname{is\text{-}abelian\text{-}group}(g) \wedge ((p + m) = n)) \Rightarrow  \\ \quad \quad \operatorname{opr}(g)(\operatorname{mpow}(\operatorname{abelian\text{-}group\text{-}as\text{-}monoid}(g), a, m), \operatorname{inv}(g)(\operatorname{mpow}(\operatorname{abelian\text{-}group\text{-}as\text{-}monoid}(g), a, n))) \\ \quad \quad =\; \operatorname{inv}(g)(\operatorname{mpow}(\operatorname{abelian\text{-}group\text{-}as\text{-}monoid}(g), a, p)) \end{array}}}
\vnbprop{zz-act-nonneg}{\vnbstmt{\begin{array}{@{}l@{}} \forall g.; \forall k \in \mathbb{N}, a \in \operatorname{carr}(g). \\ \quad \operatorname{is\text{-}abelian\text{-}group}(g) \Rightarrow  \\ \quad \quad \operatorname{zz\text{-}act}(g, k, a) = \operatorname{mpow}(\operatorname{abelian\text{-}group\text{-}as\text{-}monoid}(g), a, k) \end{array}}}
\vnbprop{zz-act-neg}{\vnbstmt{\begin{array}{@{}l@{}} \forall g.; \forall k \in \mathbb{N}, a \in \operatorname{carr}(g). \\ \quad \operatorname{is\text{-}abelian\text{-}group}(g) \Rightarrow  \\ \quad \quad \operatorname{zz\text{-}act}(g, -k, a) = \operatorname{inv}(g)(\operatorname{mpow}(\operatorname{abelian\text{-}group\text{-}as\text{-}monoid}(g), a, k)) \end{array}}}
\vnbprop{zz-act-zero}{\vnbstmt{\forall g.; \forall a \in \operatorname{carr}(g).\; (\operatorname{is\text{-}abelian\text{-}group}(g) \Rightarrow (\operatorname{zz\text{-}act}(g, 0, a) = \operatorname{iden}(g)))}}
\vnbprop{zz-act-one}{\vnbstmt{\forall g.; \forall a \in \operatorname{carr}(g).\; (\operatorname{is\text{-}abelian\text{-}group}(g) \Rightarrow (\operatorname{zz\text{-}act}(g, 1, a) = a))}}
\vnbprop{zz-act-type}{\vnbstmt{\begin{array}{@{}l@{}} \forall g.; \forall k \in \mathbb{Z}, a \in \operatorname{carr}(g). \\ \quad \operatorname{is\text{-}abelian\text{-}group}(g) \Rightarrow (\operatorname{zz\text{-}act}(g, k, a) \in \operatorname{carr}(g)) \end{array}}}
\vnbprop{zz-act-neg-sign}{\vnbstmt{\begin{array}{@{}l@{}} \forall g.; \forall k \in \mathbb{Z}, a \in \operatorname{carr}(g). \\ \quad \operatorname{is\text{-}abelian\text{-}group}(g) \Rightarrow (\operatorname{zz\text{-}act}(g, -k, a) = \operatorname{inv}(g)(\operatorname{zz\text{-}act}(g, k, a))) \end{array}}}
\vnbprop{zz-act-add}{\vnbstmt{\begin{array}{@{}l@{}} \forall g.; \forall j, k \in \mathbb{Z}, a \in \operatorname{carr}(g). \\ \quad \operatorname{is\text{-}abelian\text{-}group}(g) \Rightarrow  \\ \quad \quad \operatorname{zz\text{-}act}(g, (j + k), a) = \operatorname{opr}(g)(\operatorname{zz\text{-}act}(g, j, a), \operatorname{zz\text{-}act}(g, k, a)) \end{array}}}

\vnbbookfile{theorem-library/rake-ringoid-additive}{theorem-library/rake-ringoid-additive}
\noindent{\small rake-ringoid-additive.scm -- the additive group of a RINGOID, proven.}\par
\vnbprop{r6b-radd-type}{\vnbstmt{\begin{array}{@{}l@{}} \forall r, a, b. \\ \quad (\operatorname{is\text{-}ringoid}(r) \wedge ((a \in \operatorname{carr}(r)) \wedge (b \in \operatorname{carr}(r)))) \Rightarrow  \\ \quad \quad \operatorname{add}(r)(a, b) \in \operatorname{carr}(r) \end{array}}}
\vnbprop{r6b-rneg-type}{\vnbstmt{\forall r.; \forall a \in \operatorname{carr}(r).\; (\operatorname{is\text{-}ringoid}(r) \Rightarrow (\operatorname{neg}(r)(a) \in \operatorname{carr}(r)))}}
\vnbprop{r6b-radd-right-id}{\vnbstmt{\forall r.; \forall a \in \operatorname{carr}(r).\; (\operatorname{is\text{-}ringoid}(r) \Rightarrow (\operatorname{add}(r)(a, \operatorname{zero}(r)) = a))}}
\vnbprop{r6b-radd-interchange}{\vnbstmt{\begin{array}{@{}l@{}} \forall r, a, b, c, d. \\ \quad (\operatorname{is\text{-}ringoid}(r) \wedge ((a \in \operatorname{carr}(r)) \wedge ((b \in \operatorname{carr}(r)) \wedge ((c \in \operatorname{carr}(r)) \wedge (d \in \operatorname{carr}(r)))))) \Rightarrow  \\ \quad \quad \operatorname{add}(r)(\operatorname{add}(r)(a, b), \operatorname{add}(r)(c, d)) = \operatorname{add}(r)(\operatorname{add}(r)(a, c), \operatorname{add}(r)(b, d)) \end{array}}}
\vnbprop{r6b-rinv-unique}{\vnbstmt{\begin{array}{@{}l@{}} \forall r, a, b. \\ \quad (\operatorname{is\text{-}ringoid}(r) \wedge ((a \in \operatorname{carr}(r)) \wedge ((b \in \operatorname{carr}(r)) \wedge (\operatorname{add}(r)(a, b) = \operatorname{zero}(r))))) \Rightarrow  \\ \quad \quad b = \operatorname{neg}(r)(a) \end{array}}}
\vnbprop{r6b-rneg-add}{\vnbstmt{\begin{array}{@{}l@{}} \forall r, a, b. \\ \quad (\operatorname{is\text{-}ringoid}(r) \wedge ((a \in \operatorname{carr}(r)) \wedge (b \in \operatorname{carr}(r)))) \Rightarrow  \\ \quad \quad \operatorname{neg}(r)(\operatorname{add}(r)(a, b)) = \operatorname{add}(r)(\operatorname{neg}(r)(a), \operatorname{neg}(r)(b)) \end{array}}}
\vnbprop{r6b-rneg-neg}{\vnbstmt{\forall r.; \forall a \in \operatorname{carr}(r).\; (\operatorname{is\text{-}ringoid}(r) \Rightarrow (\operatorname{neg}(r)(\operatorname{neg}(r)(a)) = a))}}
\vnbprop{r6b-rdiff-add}{\vnbstmt{\begin{array}{@{}l@{}} \forall r, a, b, u, v. \\ \quad (\operatorname{is\text{-}ringoid}(r) \wedge ((a \in \operatorname{carr}(r)) \wedge ((b \in \operatorname{carr}(r)) \wedge ((u \in \operatorname{carr}(r)) \wedge (v \in \operatorname{carr}(r)))))) \Rightarrow  \\ \quad \quad \operatorname{add}(r)(\operatorname{add}(r)(a, b), \operatorname{neg}(r)(\operatorname{add}(r)(u, v))) \\ \quad \quad =\; \operatorname{add}(r)(\operatorname{add}(r)(a, \operatorname{neg}(r)(u)), \operatorname{add}(r)(b, \operatorname{neg}(r)(v))) \end{array}}}
\vnbprop{ringoid-neg-diff}{\vnbstmt{\begin{array}{@{}l@{}} \forall r, a, b. \\ \quad (\operatorname{is\text{-}ringoid}(r) \wedge ((a \in \operatorname{carr}(r)) \wedge (b \in \operatorname{carr}(r)))) \Rightarrow  \\ \quad \quad \operatorname{neg}(r)(\operatorname{add}(r)(a, \operatorname{neg}(r)(b))) = \operatorname{add}(r)(b, \operatorname{neg}(r)(a)) \end{array}}}
\vnbprop{ringoid-diff-telescope}{\vnbstmt{\begin{array}{@{}l@{}} \forall r, a, b, c. \\ \quad (\operatorname{is\text{-}ringoid}(r) \wedge ((a \in \operatorname{carr}(r)) \wedge ((b \in \operatorname{carr}(r)) \wedge (c \in \operatorname{carr}(r))))) \Rightarrow  \\ \quad \quad \operatorname{add}(r)(\operatorname{add}(r)(a, \operatorname{neg}(r)(b)), \operatorname{add}(r)(b, \operatorname{neg}(r)(c))) = \operatorname{add}(r)(a, \operatorname{neg}(r)(c)) \end{array}}}

\vnbbookfile{theorem-library/rake-finsum-typing}{theorem-library/rake-finsum-typing}
\noindent{\small rake-finsum-typing.scm -- the finsum / monoid-power typing leaves of the rake (batch G), PROVEN, together with the guarded form of the one statement in the batch that is NOT provable as stated.}\par
\vnbprop{comm-monoid-carrier-closed-opr}{\vnbstmt{\begin{array}{@{}l@{}} \forall s.; \forall a, b \in \operatorname{carr}(s). \\ \quad \operatorname{is\text{-}comm\text{-}monoid}(s) \Rightarrow (\operatorname{opr}(s)(a, b) \in \operatorname{carr}(s)) \end{array}}}
\vnbprop{comm-monoid-identity-in-carr}{\vnbstmt{\forall s.\; (\operatorname{is\text{-}comm\text{-}monoid}(s) \Rightarrow (\operatorname{iden}(s) \in \operatorname{carr}(s)))}}
\vnbprop{sum-ag-comm-monoid-type-ind}{\vnbstmt{\begin{array}{@{}l@{}} \forall n \in \mathbb{N}, m.; \forall f \in (\mathbb{N} \to \operatorname{carr}(m)). \\ \quad \operatorname{is\text{-}comm\text{-}monoid}(m) \Rightarrow (\operatorname{sum\text{-}ag}(m, f, n) \in \operatorname{carr}(m)) \end{array}}}
\vnbprop{enum-fam-comm-monoid-in-fun}{\vnbstmt{\begin{array}{@{}l@{}} \forall n \in \mathbb{N}, m, s.; \forall i \in (\mathrm{OS}(n) \to s), f \in (s \to \operatorname{carr}(m)). \\ \quad \operatorname{is\text{-}comm\text{-}monoid}(m) \Rightarrow (\operatorname{enum\text{-}fam}(m, f, i, n) \in (\mathbb{N} \to \operatorname{carr}(m))) \end{array}}}
\vnbprop{finsum-comm-monoid-type}{\vnbstmt{\begin{array}{@{}l@{}} \forall m.; \forall s \in \mathbf{Set}, f \in (s \to \operatorname{carr}(m)). \\ \quad (\operatorname{is\text{-}comm\text{-}monoid}(m) \wedge (|s| \in \mathbb{N})) \Rightarrow (\operatorname{finsum}(m, f, s) \in \operatorname{carr}(m)) \end{array}}}
\vnbprop{crmcm-carr}{\vnbstmt{\begin{array}{@{}l@{}} \forall r. \\ \quad \operatorname{is\text{-}commutative\text{-}ring}(r) \Rightarrow  \\ \quad \quad \operatorname{carr}(\operatorname{commutative\text{-}ring\text{-}multiplicative\text{-}cm}(r)) = \operatorname{carr}(r) \end{array}}}
\vnbprop{prod-ring-type}{\vnbstmt{\begin{array}{@{}l@{}} \forall r.; \forall x \in \mathbf{Set}, f \in (x \to \operatorname{carr}(r)). \\ \quad (\operatorname{is\text{-}commutative\text{-}ring}(r) \wedge (|x| \in \mathbb{N})) \Rightarrow (\operatorname{prod\text{-}ring}(r, f, x) \in \operatorname{carr}(r)) \end{array}}}
\vnbprop{mpow-comm-monoid-type-ind}{\vnbstmt{\begin{array}{@{}l@{}} \forall n \in \mathbb{N}, m.; \forall x \in \operatorname{carr}(m). \\ \quad \operatorname{is\text{-}comm\text{-}monoid}(m) \Rightarrow (\operatorname{mpow}(m, x, n) \in \operatorname{carr}(m)) \end{array}}}
\vnbprop{ring-power-type}{\vnbstmt{\begin{array}{@{}l@{}} \forall r.; \forall x \in \operatorname{carr}(r), n \in \mathbb{N}. \\ \quad \operatorname{is\text{-}commutative\text{-}ring}(r) \Rightarrow (\operatorname{ring\text{-}power}(r, x, n) \in \operatorname{carr}(r)) \end{array}}}
\vnbprop{sum-ag-in-subset-ind}{\vnbstmt{\begin{array}{@{}l@{}} \forall n \in \mathbb{N}, g, m, a. \\ \quad \operatorname{is\text{-}abelian\text{-}group}(g) \wedge \\ \quad \operatorname{iden}(g) \in m \wedge \\ \quad \forall x, y \in m.\; (\operatorname{opr}(g)(x, y) \in m) \wedge \\ \quad \forall i \in \mathrm{OS}(n).\; (a(i) \in m) \Rightarrow \\ \quad \quad \operatorname{sum\text{-}ag}(g, a, n) \in m \end{array}}}
\vnbprop{finsum-in-subset}{\vnbstmt{\begin{array}{@{}l@{}} \forall g, m.; \forall s \in \mathbf{Set}, f \in (s \to \operatorname{carr}(g)). \\ \quad \operatorname{is\text{-}abelian\text{-}group}(g) \wedge \\ \quad \operatorname{iden}(g) \in m \wedge \\ \quad \forall x, y \in m.\; (\operatorname{opr}(g)(x, y) \in m) \wedge \\ \quad |s| \in \mathbb{N} \wedge \\ \quad \forall z \in s.\; (f(z) \in m) \Rightarrow \\ \quad \quad \operatorname{finsum}(g, f, s) \in m \end{array}}}

\vnbbookfile{theorem-library/rake-finsum-welldef}{theorem-library/rake-finsum-welldef}
\noindent{\small theorem-library/rake-finsum-welldef.scm -- rake batch N (2026-09-17), part 2 of 2 --------------------------------------------------------------------------- NOTE, 2026-09-20 (batch 9-B).  This file was written while cardinality was AXIOMATISED under the name CARD and the defined constant was its companion}\par
\vnbprop{opr-rearrange-laws}{\vnbstmt{\begin{array}{@{}l@{}} \forall s.; \forall a, b, c \in \operatorname{carr}(s). \\ \quad \operatorname{iden}(s) \in \operatorname{carr}(s) \wedge \\ \quad \forall \alpha, \beta \in \operatorname{carr}(s).\; (\operatorname{opr}(s)(\alpha, \beta) \in \operatorname{carr}(s)) \wedge \\ \quad \forall \alpha, \beta, \gamma \in \operatorname{carr}(s). \\ \quad \quad \operatorname{opr}(s)(\operatorname{opr}(s)(\alpha, \beta), \gamma) = \operatorname{opr}(s)(\alpha, \operatorname{opr}(s)(\beta, \gamma)) \wedge \\ \quad \forall \alpha, \beta \in \operatorname{carr}(s). \\ \quad \quad \operatorname{opr}(s)(\alpha, \beta) = \operatorname{opr}(s)(\beta, \alpha) \Rightarrow \\ \quad \quad \operatorname{opr}(s)(\operatorname{opr}(s)(a, b), c) = \operatorname{opr}(s)(\operatorname{opr}(s)(a, c), b) \end{array}}}
\vnbprop{sum-ag-type-laws}{\vnbstmt{\begin{array}{@{}l@{}} \forall n \in \mathbb{N}, s.; \forall g \in (\mathbb{N} \to \operatorname{carr}(s)). \\ \quad \operatorname{iden}(s) \in \operatorname{carr}(s) \wedge \\ \quad \forall a, b \in \operatorname{carr}(s).\; (\operatorname{opr}(s)(a, b) \in \operatorname{carr}(s)) \wedge \\ \quad \forall a, b, c \in \operatorname{carr}(s). \\ \quad \quad \operatorname{opr}(s)(\operatorname{opr}(s)(a, b), c) = \operatorname{opr}(s)(a, \operatorname{opr}(s)(b, c)) \wedge \\ \quad \forall a, b \in \operatorname{carr}(s).\; (\operatorname{opr}(s)(a, b) = \operatorname{opr}(s)(b, a)) \Rightarrow \\ \quad \quad \operatorname{sum\text{-}ag}(s, g, n) \in \operatorname{carr}(s) \end{array}}}
\vnbprop{sum-ag-segment-congruence}{\vnbstmt{\begin{array}{@{}l@{}} \forall n \in \mathbb{N}, a, g, h. \\ \quad \forall i \in \mathrm{OS}(n).\; (g(i) \simeq h(i)) \Rightarrow  \\ \quad \quad (\operatorname{sum\text{-}ag}(a, g, n) \simeq \operatorname{sum\text{-}ag}(a, h, n)) \end{array}}}
\vnbprop{sum-ag-replace-entry-laws}{\vnbstmt{\begin{array}{@{}l@{}} \forall n \in \mathbb{N}, s.; \forall g, p \in (\mathbb{N} \to \operatorname{carr}(s)), k \in \mathrm{OS}(n). \\ \quad \operatorname{iden}(s) \in \operatorname{carr}(s) \wedge \\ \quad \forall a, b \in \operatorname{carr}(s).\; (\operatorname{opr}(s)(a, b) \in \operatorname{carr}(s)) \wedge \\ \quad \forall a, b, c \in \operatorname{carr}(s). \\ \quad \quad \operatorname{opr}(s)(\operatorname{opr}(s)(a, b), c) = \operatorname{opr}(s)(a, \operatorname{opr}(s)(b, c)) \wedge \\ \quad \forall a, b \in \operatorname{carr}(s).\; (\operatorname{opr}(s)(a, b) = \operatorname{opr}(s)(b, a)) \wedge \\ \quad \forall i \in \mathrm{OS}(n).\; (\neg (i = k) \Rightarrow (g(i) \simeq p(i))) \Rightarrow \\ \quad \quad \operatorname{opr}(s)(\operatorname{sum\text{-}ag}(s, g, n), p(k)) = \operatorname{opr}(s)(\operatorname{sum\text{-}ag}(s, p, n), g(k)) \end{array}}}
\vnbprop{seg-perm-restrict-is-bijection}{\vnbstmt{\begin{array}{@{}l@{}} \forall n \in \mathbb{N}, i \in \mathrm{Bij}(\mathrm{OS}((n + 1)), \mathrm{OS}((n + 1))), k. \\ \quad (i(n) = k) \Rightarrow  \\ \quad \quad (\lambda a \in \mathrm{OS}(n).\; \left(\text{if }(i(a) = n)\text{ then }k\text{ else }i(a)\right)) \in \mathrm{Bij}(\mathrm{OS}(n), \mathrm{OS}(n)) \end{array}}}
\vnbprop{sum-ag-permutation-invariance-laws}{\vnbstmt{\begin{array}{@{}l@{}} \forall n \in \mathbb{N}, s.; \forall g, h \in (\mathbb{N} \to \operatorname{carr}(s)), a \in \mathrm{Bij}(\mathrm{OS}(n), \mathrm{OS}(n)). \\ \quad \operatorname{iden}(s) \in \operatorname{carr}(s) \wedge \\ \quad \forall \alpha, b \in \operatorname{carr}(s).\; (\operatorname{opr}(s)(\alpha, b) \in \operatorname{carr}(s)) \wedge \\ \quad \forall \alpha, b, c \in \operatorname{carr}(s). \\ \quad \quad \operatorname{opr}(s)(\operatorname{opr}(s)(\alpha, b), c) = \operatorname{opr}(s)(\alpha, \operatorname{opr}(s)(b, c)) \wedge \\ \quad \forall \alpha, b \in \operatorname{carr}(s). \\ \quad \quad \operatorname{opr}(s)(\alpha, b) = \operatorname{opr}(s)(b, \alpha) \wedge \\ \quad \forall i \in \mathrm{OS}(n).\; (h(i) = g(a(i))) \Rightarrow \\ \quad \quad \operatorname{sum\text{-}ag}(s, g, n) = \operatorname{sum\text{-}ag}(s, h, n) \end{array}}}
\vnbprop{sum-ag-permutation-invariance}{\vnbstmt{\begin{array}{@{}l@{}} \forall n \in \mathbb{N}, a.; \forall g, h \in (\mathbb{N} \to \operatorname{carr}(a)), b \in \mathrm{Bij}(\mathrm{OS}(n), \mathrm{OS}(n)). \\ \quad (\operatorname{is\text{-}abelian\text{-}group}(a) \wedge \forall i \in \mathrm{OS}(n).\; (h(i) = g(b(i)))) \Rightarrow  \\ \quad \quad \operatorname{sum\text{-}ag}(a, g, n) = \operatorname{sum\text{-}ag}(a, h, n) \end{array}}}
\vnbprop{finsum-comm-monoid-permutation-invariance}{\vnbstmt{\begin{array}{@{}l@{}} \forall n \in \mathbb{N}, m.; \forall g, h \in (\mathbb{N} \to \operatorname{carr}(m)), a \in \mathrm{Bij}(\mathrm{OS}(n), \mathrm{OS}(n)). \\ \quad (\operatorname{is\text{-}comm\text{-}monoid}(m) \wedge \forall i \in \mathrm{OS}(n).\; (h(i) = g(a(i)))) \Rightarrow  \\ \quad \quad \operatorname{sum\text{-}ag}(m, g, n) = \operatorname{sum\text{-}ag}(m, h, n) \end{array}}}
\vnbprop{finsum-well-defined}{\vnbstmt{\begin{array}{@{}l@{}} \forall s \in \mathbf{Set}, g.; \forall f \in (s \to \operatorname{carr}(g)), m \in \mathrm{Bij}(\mathrm{OS}(|s|), s). \\ \quad ((|s| \in \mathbb{N}) \wedge \operatorname{is\text{-}abelian\text{-}group}(g)) \Rightarrow  \\ \quad \quad \operatorname{finsum}(g, f, s) = \operatorname{sum\text{-}ag}(g, \operatorname{enum\text{-}fam}(g, f, m, |s|), |s|) \end{array}}}
\vnbprop{finsum-comm-monoid-well-defined}{\vnbstmt{\begin{array}{@{}l@{}} \forall s \in \mathbf{Set}, m.; \forall f \in (s \to \operatorname{carr}(m)), a \in \mathrm{Bij}(\mathrm{OS}(|s|), s). \\ \quad ((|s| \in \mathbb{N}) \wedge \operatorname{is\text{-}comm\text{-}monoid}(m)) \Rightarrow  \\ \quad \quad \operatorname{finsum}(m, f, s) = \operatorname{sum\text{-}ag}(m, \operatorname{enum\text{-}fam}(m, f, a, |s|), |s|) \end{array}}}

\vnbbookfile{theorem-library/finsum-insert}{theorem-library/finsum-insert}
\noindent{\small theorem-library/finsum-insert.scm -- the insertion recurrence for FINSUM. --------------------------------------------------------------------------- NOTE, 2026-09-20 (batch 9-B).  This file was written while cardinality was AXIOMATISED under the name CARD and the defined constant was its companion}\par
\vnbprop{finsum-insert-from-enum}{\vnbstmt{\begin{array}{@{}l@{}} \forall m, f, x, k, a. \\ \quad |x| \in \mathbb{N} \wedge \\ \quad |(x \cup \{k\})| = (|x| + 1) \wedge \\ \quad (a(|x|) \simeq k) \wedge \\ \quad \forall i \in \mathrm{OS}(|x|).\; (a(i) \simeq \operatorname{fin\text{-}enum}(x)(i)) \wedge \\ \quad \operatorname{finsum}(m, f, (x \cup \{k\})) \\ \quad =\; \operatorname{sum\text{-}ag}(m, \operatorname{enum\text{-}fam}(m, f, a, |(x \cup \{k\})|), |(x \cup \{k\})|) \Rightarrow \\ \quad \quad \operatorname{finsum}(m, f, (x \cup \{k\})) = \operatorname{opr}(m)(\operatorname{finsum}(m, f, x), f(k)) \end{array}}}
\vnbprop{finsum-insert}{\vnbstmt{\begin{array}{@{}l@{}} \forall m.; \forall x, k \in \mathbf{Set}, f \in ((x \cup \{k\}) \to \operatorname{carr}(m)). \\ \quad (\operatorname{is\text{-}comm\text{-}monoid}(m) \wedge ((|x| \in \mathbb{N}) \wedge \neg (k \in x))) \Rightarrow  \\ \quad \quad \operatorname{finsum}(m, f, (x \cup \{k\})) = \operatorname{opr}(m)(\operatorname{finsum}(m, f, x), f(k)) \end{array}}}
\vnbprop{finsum-insert-ag}{\vnbstmt{\begin{array}{@{}l@{}} \forall g.; \forall x, k \in \mathbf{Set}, f \in ((x \cup \{k\}) \to \operatorname{carr}(g)). \\ \quad (\operatorname{is\text{-}abelian\text{-}group}(g) \wedge ((|x| \in \mathbb{N}) \wedge \neg (k \in x))) \Rightarrow  \\ \quad \quad \operatorname{finsum}(g, f, (x \cup \{k\})) = \operatorname{opr}(g)(\operatorname{finsum}(g, f, x), f(k)) \end{array}}}
\vnbprop{finsum-empty}{\vnbstmt{\forall g, f.\; (\operatorname{finsum}(g, f, \emptyset) \simeq \operatorname{iden}(g))}}

\vnbbookfile{theorem-library/field-ring-view}{theorem-library/field-ring-view}
\noindent{\small field-ring-view.scm -- EVERY FIELD IS A FIELD-RING.}\par
\vnbprop{field-id-carr}{\vnbstmt{\forall f.\; (\operatorname{is\text{-}field}(f) \Rightarrow (\operatorname{carr}(\operatorname{field\text{-}as\text{-}integral\text{-}domain}(f)) \simeq \operatorname{carr}(f)))}}
\vnbprop{field-id-mul}{\vnbstmt{\forall f.\; (\operatorname{is\text{-}field}(f) \Rightarrow (\operatorname{mul}(\operatorname{field\text{-}as\text{-}integral\text{-}domain}(f)) \simeq \operatorname{mul}(f)))}}
\vnbprop{field-id-one}{\vnbstmt{\forall f.\; (\operatorname{is\text{-}field}(f) \Rightarrow (\operatorname{one}(\operatorname{field\text{-}as\text{-}integral\text{-}domain}(f)) \simeq \operatorname{one}(f)))}}
\vnbprop{field-id-zero}{\vnbstmt{\forall f.\; (\operatorname{is\text{-}field}(f) \Rightarrow (\operatorname{zero}(\operatorname{field\text{-}as\text{-}integral\text{-}domain}(f)) \simeq \operatorname{zero}(f)))}}
\vnbprop{field-id-add}{\vnbstmt{\forall f.\; (\operatorname{is\text{-}field}(f) \Rightarrow (\operatorname{add}(\operatorname{field\text{-}as\text{-}integral\text{-}domain}(f)) \simeq \operatorname{add}(f)))}}
\vnbprop{field-id-neg}{\vnbstmt{\forall f.\; (\operatorname{is\text{-}field}(f) \Rightarrow (\operatorname{neg}(\operatorname{field\text{-}as\text{-}integral\text{-}domain}(f)) \simeq \operatorname{neg}(f)))}}
\vnbprop{field-ring-is-integral-domain}{\vnbstmt{\forall s.\; (\operatorname{is\text{-}field\text{-}ring}(s) \Rightarrow \operatorname{is\text{-}integral\text{-}domain}(s))}}
\vnbprop{field-is-field-ring}{\vnbstmt{\forall s.\; (\operatorname{is\text{-}field}(s) \Rightarrow \operatorname{is\text{-}field\text{-}ring}(\operatorname{field\text{-}as\text{-}integral\text{-}domain}(s)))}}
\vnbprop{field-as-integral-domain-is-integral-domain}{\vnbstmt{\forall r.\; (\operatorname{is\text{-}field}(r) \Rightarrow \operatorname{is\text{-}integral\text{-}domain}(\operatorname{field\text{-}as\text{-}integral\text{-}domain}(r)))}}

\vnbbookfile{theorem-library/view-typing-proofs}{theorem-library/view-typing-proofs}
\noindent{\small theorem-library/view-typing-proofs.scm -- THE DEFERRED FIELD VIEW TYPINGS, PROVEN (D8, 2026-11-15; the user's decision "Prove views + fix F1/F2").}\par
\vnbprop{field-as-euclidean-ring-is-euclidean-ring}{\vnbstmt{\forall r.\; (\operatorname{is\text{-}field}(r) \Rightarrow \operatorname{is\text{-}euclidean\text{-}ring}(\operatorname{field\text{-}as\text{-}euclidean\text{-}ring}(r)))}}
\vnbprop{field-non-zero-mul-closed}{\vnbstmt{\begin{array}{@{}l@{}} \forall r.; \forall u, v \in \operatorname{non\text{-}zero}(r). \\ \quad \operatorname{is\text{-}field}(r) \Rightarrow (\operatorname{mul}(r)(u, v) \in \operatorname{non\text{-}zero}(r)) \end{array}}}
\vnbprop{field-multiplicative-group-is-group}{\vnbstmt{\forall r.\; (\operatorname{is\text{-}field}(r) \Rightarrow \operatorname{is\text{-}group}(\operatorname{field\text{-}multiplicative\text{-}group}(r)))}}

\vnbbookfile{theorem-library/hom-laws}{theorem-library/hom-laws}
\noindent{\small hom-laws.scm -- the models of a structure form a CATEGORY: identity arrows, composition, and the hom-set as a set.}\par
\vnbprop{hom-semigroup-member-iff}{\vnbstmt{\begin{array}{@{}l@{}} \forall a, b, f. \\ \quad (f \in \operatorname{hom\text{-}semigroup}(a, b)) \Leftrightarrow ((f \in (\operatorname{carr}(a) \to \operatorname{carr}(b))) \wedge \operatorname{is\text{-}hom\text{-}semigroup}(a, b, f)) \end{array}}}
\vnbprop{hom-semigroup-id}{\vnbstmt{\forall a.\; (\operatorname{is\text{-}semigroup}(a) \Rightarrow \operatorname{is\text{-}hom\text{-}semigroup}(a, a, \operatorname{id\text{-}fun}(\operatorname{carr}(a))))}}
\vnbprop{hom-semigroup-in-set}{\vnbstmt{\begin{array}{@{}l@{}} \forall a, b. \\ \quad (\operatorname{is\text{-}semigroup}(a) \wedge \operatorname{is\text{-}semigroup}(b)) \Rightarrow (\operatorname{hom\text{-}semigroup}(a, b) \in \mathbf{Set}) \end{array}}}
\vnbprop{hom-monoid-member-iff}{\vnbstmt{\begin{array}{@{}l@{}} \forall a, b, f. \\ \quad (f \in \operatorname{hom\text{-}monoid}(a, b)) \Leftrightarrow ((f \in (\operatorname{carr}(a) \to \operatorname{carr}(b))) \wedge \operatorname{is\text{-}hom\text{-}monoid}(a, b, f)) \end{array}}}
\vnbprop{hom-monoid-id}{\vnbstmt{\forall a.\; (\operatorname{is\text{-}monoid}(a) \Rightarrow \operatorname{is\text{-}hom\text{-}monoid}(a, a, \operatorname{id\text{-}fun}(\operatorname{carr}(a))))}}
\vnbprop{hom-monoid-in-set}{\vnbstmt{\forall a, b.\; ((\operatorname{is\text{-}monoid}(a) \wedge \operatorname{is\text{-}monoid}(b)) \Rightarrow (\operatorname{hom\text{-}monoid}(a, b) \in \mathbf{Set}))}}
\vnbprop{hom-comm-monoid-member-iff}{\vnbstmt{\begin{array}{@{}l@{}} \forall a, b, f. \\ \quad (f \in \operatorname{hom\text{-}comm\text{-}monoid}(a, b)) \Leftrightarrow ((f \in (\operatorname{carr}(a) \to \operatorname{carr}(b))) \wedge \operatorname{is\text{-}hom\text{-}comm\text{-}monoid}(a, b, f)) \end{array}}}
\vnbprop{hom-comm-monoid-id}{\vnbstmt{\forall a.\; (\operatorname{is\text{-}comm\text{-}monoid}(a) \Rightarrow \operatorname{is\text{-}hom\text{-}comm\text{-}monoid}(a, a, \operatorname{id\text{-}fun}(\operatorname{carr}(a))))}}
\vnbprop{hom-comm-monoid-in-set}{\vnbstmt{\begin{array}{@{}l@{}} \forall a, b. \\ \quad (\operatorname{is\text{-}comm\text{-}monoid}(a) \wedge \operatorname{is\text{-}comm\text{-}monoid}(b)) \Rightarrow (\operatorname{hom\text{-}comm\text{-}monoid}(a, b) \in \mathbf{Set}) \end{array}}}
\vnbprop{hom-group-member-iff}{\vnbstmt{\begin{array}{@{}l@{}} \forall a, b, f. \\ \quad (f \in \operatorname{hom\text{-}group}(a, b)) \Leftrightarrow ((f \in (\operatorname{carr}(a) \to \operatorname{carr}(b))) \wedge \operatorname{is\text{-}hom\text{-}group}(a, b, f)) \end{array}}}
\vnbprop{hom-group-id}{\vnbstmt{\forall a.\; (\operatorname{is\text{-}group}(a) \Rightarrow \operatorname{is\text{-}hom\text{-}group}(a, a, \operatorname{id\text{-}fun}(\operatorname{carr}(a))))}}
\vnbprop{hom-group-in-set}{\vnbstmt{\forall a, b.\; ((\operatorname{is\text{-}group}(a) \wedge \operatorname{is\text{-}group}(b)) \Rightarrow (\operatorname{hom\text{-}group}(a, b) \in \mathbf{Set}))}}
\vnbprop{hom-abelian-group-member-iff}{\vnbstmt{\begin{array}{@{}l@{}} \forall a, b, f. \\ \quad (f \in \operatorname{hom\text{-}abelian\text{-}group}(a, b)) \Leftrightarrow ((f \in (\operatorname{carr}(a) \to \operatorname{carr}(b))) \wedge \operatorname{is\text{-}hom\text{-}abelian\text{-}group}(a, b, f)) \end{array}}}
\vnbprop{hom-abelian-group-id}{\vnbstmt{\begin{array}{@{}l@{}} \forall a. \\ \quad \operatorname{is\text{-}abelian\text{-}group}(a) \Rightarrow \operatorname{is\text{-}hom\text{-}abelian\text{-}group}(a, a, \operatorname{id\text{-}fun}(\operatorname{carr}(a))) \end{array}}}
\vnbprop{hom-abelian-group-in-set}{\vnbstmt{\begin{array}{@{}l@{}} \forall a, b. \\ \quad (\operatorname{is\text{-}abelian\text{-}group}(a) \wedge \operatorname{is\text{-}abelian\text{-}group}(b)) \Rightarrow  \\ \quad \quad \operatorname{hom\text{-}abelian\text{-}group}(a, b) \in \mathbf{Set} \end{array}}}
\vnbprop{hom-ring-member-iff}{\vnbstmt{\begin{array}{@{}l@{}} \forall a, b, f. \\ \quad (f \in \operatorname{hom\text{-}ring}(a, b)) \Leftrightarrow ((f \in (\operatorname{carr}(a) \to \operatorname{carr}(b))) \wedge \operatorname{is\text{-}hom\text{-}ring}(a, b, f)) \end{array}}}
\vnbprop{hom-ring-id}{\vnbstmt{\forall a.\; (\operatorname{is\text{-}ring}(a) \Rightarrow \operatorname{is\text{-}hom\text{-}ring}(a, a, \operatorname{id\text{-}fun}(\operatorname{carr}(a))))}}
\vnbprop{hom-ring-in-set}{\vnbstmt{\forall a, b.\; ((\operatorname{is\text{-}ring}(a) \wedge \operatorname{is\text{-}ring}(b)) \Rightarrow (\operatorname{hom\text{-}ring}(a, b) \in \mathbf{Set}))}}
\vnbprop{hom-commutative-ring-member-iff}{\vnbstmt{\begin{array}{@{}l@{}} \forall a, b, f. \\ \quad (f \in \operatorname{hom\text{-}commutative\text{-}ring}(a, b)) \Leftrightarrow ((f \in (\operatorname{carr}(a) \to \operatorname{carr}(b))) \wedge \operatorname{is\text{-}hom\text{-}commutative\text{-}ring}(a, b, f)) \end{array}}}
\vnbprop{hom-commutative-ring-id}{\vnbstmt{\begin{array}{@{}l@{}} \forall a. \\ \quad \operatorname{is\text{-}commutative\text{-}ring}(a) \Rightarrow \operatorname{is\text{-}hom\text{-}commutative\text{-}ring}(a, a, \operatorname{id\text{-}fun}(\operatorname{carr}(a))) \end{array}}}
\vnbprop{hom-commutative-ring-in-set}{\vnbstmt{\begin{array}{@{}l@{}} \forall a, b. \\ \quad (\operatorname{is\text{-}commutative\text{-}ring}(a) \wedge \operatorname{is\text{-}commutative\text{-}ring}(b)) \Rightarrow  \\ \quad \quad \operatorname{hom\text{-}commutative\text{-}ring}(a, b) \in \mathbf{Set} \end{array}}}
\vnbprop{hom-integral-domain-member-iff}{\vnbstmt{\begin{array}{@{}l@{}} \forall a, b, f. \\ \quad (f \in \operatorname{hom\text{-}integral\text{-}domain}(a, b)) \Leftrightarrow ((f \in (\operatorname{carr}(a) \to \operatorname{carr}(b))) \wedge \operatorname{is\text{-}hom\text{-}integral\text{-}domain}(a, b, f)) \end{array}}}
\vnbprop{hom-integral-domain-id}{\vnbstmt{\begin{array}{@{}l@{}} \forall a. \\ \quad \operatorname{is\text{-}integral\text{-}domain}(a) \Rightarrow \operatorname{is\text{-}hom\text{-}integral\text{-}domain}(a, a, \operatorname{id\text{-}fun}(\operatorname{carr}(a))) \end{array}}}
\vnbprop{hom-integral-domain-in-set}{\vnbstmt{\begin{array}{@{}l@{}} \forall a, b. \\ \quad (\operatorname{is\text{-}integral\text{-}domain}(a) \wedge \operatorname{is\text{-}integral\text{-}domain}(b)) \Rightarrow  \\ \quad \quad \operatorname{hom\text{-}integral\text{-}domain}(a, b) \in \mathbf{Set} \end{array}}}
\vnbprop{hom-euclidean-ring-member-iff}{\vnbstmt{\begin{array}{@{}l@{}} \forall a, b, f. \\ \quad (f \in \operatorname{hom\text{-}euclidean\text{-}ring}(a, b)) \Leftrightarrow ((f \in (\operatorname{carr}(a) \to \operatorname{carr}(b))) \wedge \operatorname{is\text{-}hom\text{-}euclidean\text{-}ring}(a, b, f)) \end{array}}}
\vnbprop{hom-euclidean-ring-id}{\vnbstmt{\begin{array}{@{}l@{}} \forall a. \\ \quad \operatorname{is\text{-}euclidean\text{-}ring}(a) \Rightarrow \operatorname{is\text{-}hom\text{-}euclidean\text{-}ring}(a, a, \operatorname{id\text{-}fun}(\operatorname{carr}(a))) \end{array}}}
\vnbprop{hom-euclidean-ring-in-set}{\vnbstmt{\begin{array}{@{}l@{}} \forall a, b. \\ \quad (\operatorname{is\text{-}euclidean\text{-}ring}(a) \wedge \operatorname{is\text{-}euclidean\text{-}ring}(b)) \Rightarrow  \\ \quad \quad \operatorname{hom\text{-}euclidean\text{-}ring}(a, b) \in \mathbf{Set} \end{array}}}
\vnbprop{hom-pid-member-iff}{\vnbstmt{\begin{array}{@{}l@{}} \forall a, b, f. \\ \quad (f \in \operatorname{hom\text{-}pid}(a, b)) \Leftrightarrow ((f \in (\operatorname{carr}(a) \to \operatorname{carr}(b))) \wedge \operatorname{is\text{-}hom\text{-}pid}(a, b, f)) \end{array}}}
\vnbprop{hom-pid-id}{\vnbstmt{\forall a.\; (\operatorname{is\text{-}pid}(a) \Rightarrow \operatorname{is\text{-}hom\text{-}pid}(a, a, \operatorname{id\text{-}fun}(\operatorname{carr}(a))))}}
\vnbprop{hom-pid-in-set}{\vnbstmt{\forall a, b.\; ((\operatorname{is\text{-}pid}(a) \wedge \operatorname{is\text{-}pid}(b)) \Rightarrow (\operatorname{hom\text{-}pid}(a, b) \in \mathbf{Set}))}}
\vnbprop{hom-field-ring-member-iff}{\vnbstmt{\begin{array}{@{}l@{}} \forall a, b, f. \\ \quad (f \in \operatorname{hom\text{-}field\text{-}ring}(a, b)) \Leftrightarrow ((f \in (\operatorname{carr}(a) \to \operatorname{carr}(b))) \wedge \operatorname{is\text{-}hom\text{-}field\text{-}ring}(a, b, f)) \end{array}}}
\vnbprop{hom-field-ring-id}{\vnbstmt{\forall a.\; (\operatorname{is\text{-}field\text{-}ring}(a) \Rightarrow \operatorname{is\text{-}hom\text{-}field\text{-}ring}(a, a, \operatorname{id\text{-}fun}(\operatorname{carr}(a))))}}
\vnbprop{hom-field-ring-in-set}{\vnbstmt{\begin{array}{@{}l@{}} \forall a, b. \\ \quad (\operatorname{is\text{-}field\text{-}ring}(a) \wedge \operatorname{is\text{-}field\text{-}ring}(b)) \Rightarrow (\operatorname{hom\text{-}field\text{-}ring}(a, b) \in \mathbf{Set}) \end{array}}}
\vnbprop{hom-field-non-zero}{\vnbstmt{\begin{array}{@{}l@{}} \forall a, b, f.; \forall x \in \operatorname{non\text{-}zero}(a). \\ \quad \operatorname{is\text{-}hom\text{-}field}(a, b, f) \Rightarrow (f(x) \in \operatorname{non\text{-}zero}(b)) \end{array}}}
\vnbprop{hom-field-member-iff}{\vnbstmt{\begin{array}{@{}l@{}} \forall a, b, f. \\ \quad (f \in \operatorname{hom\text{-}field}(a, b)) \Leftrightarrow ((f \in (\operatorname{carr}(a) \to \operatorname{carr}(b))) \wedge \operatorname{is\text{-}hom\text{-}field}(a, b, f)) \end{array}}}
\vnbprop{hom-field-id}{\vnbstmt{\forall a.\; (\operatorname{is\text{-}field}(a) \Rightarrow \operatorname{is\text{-}hom\text{-}field}(a, a, \operatorname{id\text{-}fun}(\operatorname{carr}(a))))}}
\vnbprop{hom-field-compose}{\vnbstmt{\begin{array}{@{}l@{}} \forall a, b, c, f, g. \\ \quad (\operatorname{is\text{-}hom\text{-}field}(a, b, f) \wedge \operatorname{is\text{-}hom\text{-}field}(b, c, g)) \Rightarrow  \\ \quad \quad \operatorname{is\text{-}hom\text{-}field}(a, c, \operatorname{compose}(g, f)) \end{array}}}
\vnbprop{hom-field-in-set}{\vnbstmt{\forall a, b.\; ((\operatorname{is\text{-}field}(a) \wedge \operatorname{is\text{-}field}(b)) \Rightarrow (\operatorname{hom\text{-}field}(a, b) \in \mathbf{Set}))}}
\vnbprop{hom-ringoid-member-iff}{\vnbstmt{\begin{array}{@{}l@{}} \forall a, b, f. \\ \quad (f \in \operatorname{hom\text{-}ringoid}(a, b)) \Leftrightarrow ((f \in (\operatorname{carr}(a) \to \operatorname{carr}(b))) \wedge \operatorname{is\text{-}hom\text{-}ringoid}(a, b, f)) \end{array}}}
\vnbprop{hom-ringoid-id}{\vnbstmt{\forall a.\; (\operatorname{is\text{-}ringoid}(a) \Rightarrow \operatorname{is\text{-}hom\text{-}ringoid}(a, a, \operatorname{id\text{-}fun}(\operatorname{carr}(a))))}}
\vnbprop{hom-ringoid-in-set}{\vnbstmt{\forall a, b.\; ((\operatorname{is\text{-}ringoid}(a) \wedge \operatorname{is\text{-}ringoid}(b)) \Rightarrow (\operatorname{hom\text{-}ringoid}(a, b) \in \mathbf{Set}))}}

\vnbbookfile{theorem-library/hom-kinds}{theorem-library/hom-kinds}
\noindent{\small hom-kinds.scm -- the category laws for the arrows the generic driver of hom-laws.scm does not reach: the FAMILY slot kinds, and the RELATION kind (SETOID, since batch 40; two carriers in batch 39).}\par
\vnbprop{hom-setoid-member-iff}{\vnbstmt{\begin{array}{@{}l@{}} \forall a, b, f. \\ \quad (f \in \operatorname{hom\text{-}setoid}(a, b)) \Leftrightarrow ((f \in (\operatorname{pts}(a) \to \operatorname{pts}(b))) \wedge \operatorname{is\text{-}hom\text{-}setoid}(a, b, f)) \end{array}}}
\vnbprop{hom-setoid-id}{\vnbstmt{\forall a.\; (\operatorname{is\text{-}setoid}(a) \Rightarrow \operatorname{is\text{-}hom\text{-}setoid}(a, a, \operatorname{id\text{-}fun}(\operatorname{pts}(a))))}}
\vnbprop{hom-setoid-in-set}{\vnbstmt{\forall a, b.\; ((\operatorname{is\text{-}setoid}(a) \wedge \operatorname{is\text{-}setoid}(b)) \Rightarrow (\operatorname{hom\text{-}setoid}(a, b) \in \mathbf{Set}))}}

\vnbbookfile{theorem-library/hom-functors}{theorem-library/hom-functors}
\noindent{\small hom-functors.scm -- the Hom functors: Hom(a, -) covariant, Hom(-, b) contravariant, as maps on hom-sets.}\par
\vnbprop{hom-semigroup-post-type}{\vnbstmt{\begin{array}{@{}l@{}} \forall a, b, c, h. \\ \quad (\operatorname{is\text{-}semigroup}(a) \wedge (h \in \operatorname{hom\text{-}semigroup}(b, c))) \Rightarrow  \\ \quad \quad (\lambda f \in \operatorname{hom\text{-}semigroup}(a, b).\; \operatorname{compose}(h, f)) \in (\operatorname{hom\text{-}semigroup}(a, b) \to \operatorname{hom\text{-}semigroup}(a, c)) \end{array}}}
\vnbprop{hom-semigroup-pre-type}{\vnbstmt{\begin{array}{@{}l@{}} \forall a, b, c, k. \\ \quad (\operatorname{is\text{-}semigroup}(b) \wedge (k \in \operatorname{hom\text{-}semigroup}(c, a))) \Rightarrow  \\ \quad \quad (\lambda f \in \operatorname{hom\text{-}semigroup}(a, b).\; \operatorname{compose}(f, k)) \in (\operatorname{hom\text{-}semigroup}(a, b) \to \operatorname{hom\text{-}semigroup}(c, b)) \end{array}}}
\vnbprop{hom-monoid-post-type}{\vnbstmt{\begin{array}{@{}l@{}} \forall a, b, c, h. \\ \quad (\operatorname{is\text{-}monoid}(a) \wedge (h \in \operatorname{hom\text{-}monoid}(b, c))) \Rightarrow  \\ \quad \quad (\lambda f \in \operatorname{hom\text{-}monoid}(a, b).\; \operatorname{compose}(h, f)) \in (\operatorname{hom\text{-}monoid}(a, b) \to \operatorname{hom\text{-}monoid}(a, c)) \end{array}}}
\vnbprop{hom-monoid-pre-type}{\vnbstmt{\begin{array}{@{}l@{}} \forall a, b, c, k. \\ \quad (\operatorname{is\text{-}monoid}(b) \wedge (k \in \operatorname{hom\text{-}monoid}(c, a))) \Rightarrow  \\ \quad \quad (\lambda f \in \operatorname{hom\text{-}monoid}(a, b).\; \operatorname{compose}(f, k)) \in (\operatorname{hom\text{-}monoid}(a, b) \to \operatorname{hom\text{-}monoid}(c, b)) \end{array}}}
\vnbprop{hom-comm-monoid-post-type}{\vnbstmt{\begin{array}{@{}l@{}} \forall a, b, c, h. \\ \quad (\operatorname{is\text{-}comm\text{-}monoid}(a) \wedge (h \in \operatorname{hom\text{-}comm\text{-}monoid}(b, c))) \Rightarrow  \\ \quad \quad (\lambda f \in \operatorname{hom\text{-}comm\text{-}monoid}(a, b).\; \operatorname{compose}(h, f)) \\ \quad \quad \in\; (\operatorname{hom\text{-}comm\text{-}monoid}(a, b) \to \operatorname{hom\text{-}comm\text{-}monoid}(a, c)) \end{array}}}
\vnbprop{hom-comm-monoid-pre-type}{\vnbstmt{\begin{array}{@{}l@{}} \forall a, b, c, k. \\ \quad (\operatorname{is\text{-}comm\text{-}monoid}(b) \wedge (k \in \operatorname{hom\text{-}comm\text{-}monoid}(c, a))) \Rightarrow  \\ \quad \quad (\lambda f \in \operatorname{hom\text{-}comm\text{-}monoid}(a, b).\; \operatorname{compose}(f, k)) \\ \quad \quad \in\; (\operatorname{hom\text{-}comm\text{-}monoid}(a, b) \to \operatorname{hom\text{-}comm\text{-}monoid}(c, b)) \end{array}}}
\vnbprop{hom-group-post-type}{\vnbstmt{\begin{array}{@{}l@{}} \forall a, b, c, h. \\ \quad (\operatorname{is\text{-}group}(a) \wedge (h \in \operatorname{hom\text{-}group}(b, c))) \Rightarrow  \\ \quad \quad (\lambda f \in \operatorname{hom\text{-}group}(a, b).\; \operatorname{compose}(h, f)) \in (\operatorname{hom\text{-}group}(a, b) \to \operatorname{hom\text{-}group}(a, c)) \end{array}}}
\vnbprop{hom-group-pre-type}{\vnbstmt{\begin{array}{@{}l@{}} \forall a, b, c, k. \\ \quad (\operatorname{is\text{-}group}(b) \wedge (k \in \operatorname{hom\text{-}group}(c, a))) \Rightarrow  \\ \quad \quad (\lambda f \in \operatorname{hom\text{-}group}(a, b).\; \operatorname{compose}(f, k)) \in (\operatorname{hom\text{-}group}(a, b) \to \operatorname{hom\text{-}group}(c, b)) \end{array}}}
\vnbprop{hom-abelian-group-post-type}{\vnbstmt{\begin{array}{@{}l@{}} \forall a, b, c, h. \\ \quad (\operatorname{is\text{-}abelian\text{-}group}(a) \wedge (h \in \operatorname{hom\text{-}abelian\text{-}group}(b, c))) \Rightarrow  \\ \quad \quad (\lambda f \in \operatorname{hom\text{-}abelian\text{-}group}(a, b).\; \operatorname{compose}(h, f)) \\ \quad \quad \in\; (\operatorname{hom\text{-}abelian\text{-}group}(a, b) \to \operatorname{hom\text{-}abelian\text{-}group}(a, c)) \end{array}}}
\vnbprop{hom-abelian-group-pre-type}{\vnbstmt{\begin{array}{@{}l@{}} \forall a, b, c, k. \\ \quad (\operatorname{is\text{-}abelian\text{-}group}(b) \wedge (k \in \operatorname{hom\text{-}abelian\text{-}group}(c, a))) \Rightarrow  \\ \quad \quad (\lambda f \in \operatorname{hom\text{-}abelian\text{-}group}(a, b).\; \operatorname{compose}(f, k)) \\ \quad \quad \in\; (\operatorname{hom\text{-}abelian\text{-}group}(a, b) \to \operatorname{hom\text{-}abelian\text{-}group}(c, b)) \end{array}}}
\vnbprop{hom-ring-post-type}{\vnbstmt{\begin{array}{@{}l@{}} \forall a, b, c, h. \\ \quad (\operatorname{is\text{-}ring}(a) \wedge (h \in \operatorname{hom\text{-}ring}(b, c))) \Rightarrow  \\ \quad \quad (\lambda f \in \operatorname{hom\text{-}ring}(a, b).\; \operatorname{compose}(h, f)) \in (\operatorname{hom\text{-}ring}(a, b) \to \operatorname{hom\text{-}ring}(a, c)) \end{array}}}
\vnbprop{hom-ring-pre-type}{\vnbstmt{\begin{array}{@{}l@{}} \forall a, b, c, k. \\ \quad (\operatorname{is\text{-}ring}(b) \wedge (k \in \operatorname{hom\text{-}ring}(c, a))) \Rightarrow  \\ \quad \quad (\lambda f \in \operatorname{hom\text{-}ring}(a, b).\; \operatorname{compose}(f, k)) \in (\operatorname{hom\text{-}ring}(a, b) \to \operatorname{hom\text{-}ring}(c, b)) \end{array}}}
\vnbprop{hom-commutative-ring-post-type}{\vnbstmt{\begin{array}{@{}l@{}} \forall a, b, c, h. \\ \quad (\operatorname{is\text{-}commutative\text{-}ring}(a) \wedge (h \in \operatorname{hom\text{-}commutative\text{-}ring}(b, c))) \Rightarrow  \\ \quad \quad (\lambda f \in \operatorname{hom\text{-}commutative\text{-}ring}(a, b).\; \operatorname{compose}(h, f)) \\ \quad \quad \in\; (\operatorname{hom\text{-}commutative\text{-}ring}(a, b) \to \operatorname{hom\text{-}commutative\text{-}ring}(a, c)) \end{array}}}
\vnbprop{hom-commutative-ring-pre-type}{\vnbstmt{\begin{array}{@{}l@{}} \forall a, b, c, k. \\ \quad (\operatorname{is\text{-}commutative\text{-}ring}(b) \wedge (k \in \operatorname{hom\text{-}commutative\text{-}ring}(c, a))) \Rightarrow  \\ \quad \quad (\lambda f \in \operatorname{hom\text{-}commutative\text{-}ring}(a, b).\; \operatorname{compose}(f, k)) \\ \quad \quad \in\; (\operatorname{hom\text{-}commutative\text{-}ring}(a, b) \to \operatorname{hom\text{-}commutative\text{-}ring}(c, b)) \end{array}}}
\vnbprop{hom-integral-domain-post-type}{\vnbstmt{\begin{array}{@{}l@{}} \forall a, b, c, h. \\ \quad (\operatorname{is\text{-}integral\text{-}domain}(a) \wedge (h \in \operatorname{hom\text{-}integral\text{-}domain}(b, c))) \Rightarrow  \\ \quad \quad (\lambda f \in \operatorname{hom\text{-}integral\text{-}domain}(a, b).\; \operatorname{compose}(h, f)) \\ \quad \quad \in\; (\operatorname{hom\text{-}integral\text{-}domain}(a, b) \to \operatorname{hom\text{-}integral\text{-}domain}(a, c)) \end{array}}}
\vnbprop{hom-integral-domain-pre-type}{\vnbstmt{\begin{array}{@{}l@{}} \forall a, b, c, k. \\ \quad (\operatorname{is\text{-}integral\text{-}domain}(b) \wedge (k \in \operatorname{hom\text{-}integral\text{-}domain}(c, a))) \Rightarrow  \\ \quad \quad (\lambda f \in \operatorname{hom\text{-}integral\text{-}domain}(a, b).\; \operatorname{compose}(f, k)) \\ \quad \quad \in\; (\operatorname{hom\text{-}integral\text{-}domain}(a, b) \to \operatorname{hom\text{-}integral\text{-}domain}(c, b)) \end{array}}}
\vnbprop{hom-euclidean-ring-post-type}{\vnbstmt{\begin{array}{@{}l@{}} \forall a, b, c, h. \\ \quad (\operatorname{is\text{-}euclidean\text{-}ring}(a) \wedge (h \in \operatorname{hom\text{-}euclidean\text{-}ring}(b, c))) \Rightarrow  \\ \quad \quad (\lambda f \in \operatorname{hom\text{-}euclidean\text{-}ring}(a, b).\; \operatorname{compose}(h, f)) \\ \quad \quad \in\; (\operatorname{hom\text{-}euclidean\text{-}ring}(a, b) \to \operatorname{hom\text{-}euclidean\text{-}ring}(a, c)) \end{array}}}
\vnbprop{hom-euclidean-ring-pre-type}{\vnbstmt{\begin{array}{@{}l@{}} \forall a, b, c, k. \\ \quad (\operatorname{is\text{-}euclidean\text{-}ring}(b) \wedge (k \in \operatorname{hom\text{-}euclidean\text{-}ring}(c, a))) \Rightarrow  \\ \quad \quad (\lambda f \in \operatorname{hom\text{-}euclidean\text{-}ring}(a, b).\; \operatorname{compose}(f, k)) \\ \quad \quad \in\; (\operatorname{hom\text{-}euclidean\text{-}ring}(a, b) \to \operatorname{hom\text{-}euclidean\text{-}ring}(c, b)) \end{array}}}
\vnbprop{hom-pid-post-type}{\vnbstmt{\begin{array}{@{}l@{}} \forall a, b, c, h. \\ \quad (\operatorname{is\text{-}pid}(a) \wedge (h \in \operatorname{hom\text{-}pid}(b, c))) \Rightarrow  \\ \quad \quad (\lambda f \in \operatorname{hom\text{-}pid}(a, b).\; \operatorname{compose}(h, f)) \in (\operatorname{hom\text{-}pid}(a, b) \to \operatorname{hom\text{-}pid}(a, c)) \end{array}}}
\vnbprop{hom-pid-pre-type}{\vnbstmt{\begin{array}{@{}l@{}} \forall a, b, c, k. \\ \quad (\operatorname{is\text{-}pid}(b) \wedge (k \in \operatorname{hom\text{-}pid}(c, a))) \Rightarrow  \\ \quad \quad (\lambda f \in \operatorname{hom\text{-}pid}(a, b).\; \operatorname{compose}(f, k)) \in (\operatorname{hom\text{-}pid}(a, b) \to \operatorname{hom\text{-}pid}(c, b)) \end{array}}}
\vnbprop{hom-field-ring-post-type}{\vnbstmt{\begin{array}{@{}l@{}} \forall a, b, c, h. \\ \quad (\operatorname{is\text{-}field\text{-}ring}(a) \wedge (h \in \operatorname{hom\text{-}field\text{-}ring}(b, c))) \Rightarrow  \\ \quad \quad (\lambda f \in \operatorname{hom\text{-}field\text{-}ring}(a, b).\; \operatorname{compose}(h, f)) \\ \quad \quad \in\; (\operatorname{hom\text{-}field\text{-}ring}(a, b) \to \operatorname{hom\text{-}field\text{-}ring}(a, c)) \end{array}}}
\vnbprop{hom-field-ring-pre-type}{\vnbstmt{\begin{array}{@{}l@{}} \forall a, b, c, k. \\ \quad (\operatorname{is\text{-}field\text{-}ring}(b) \wedge (k \in \operatorname{hom\text{-}field\text{-}ring}(c, a))) \Rightarrow  \\ \quad \quad (\lambda f \in \operatorname{hom\text{-}field\text{-}ring}(a, b).\; \operatorname{compose}(f, k)) \\ \quad \quad \in\; (\operatorname{hom\text{-}field\text{-}ring}(a, b) \to \operatorname{hom\text{-}field\text{-}ring}(c, b)) \end{array}}}
\vnbprop{hom-field-post-type}{\vnbstmt{\begin{array}{@{}l@{}} \forall a, b, c, h. \\ \quad (\operatorname{is\text{-}field}(a) \wedge (h \in \operatorname{hom\text{-}field}(b, c))) \Rightarrow  \\ \quad \quad (\lambda f \in \operatorname{hom\text{-}field}(a, b).\; \operatorname{compose}(h, f)) \in (\operatorname{hom\text{-}field}(a, b) \to \operatorname{hom\text{-}field}(a, c)) \end{array}}}
\vnbprop{hom-field-pre-type}{\vnbstmt{\begin{array}{@{}l@{}} \forall a, b, c, k. \\ \quad (\operatorname{is\text{-}field}(b) \wedge (k \in \operatorname{hom\text{-}field}(c, a))) \Rightarrow  \\ \quad \quad (\lambda f \in \operatorname{hom\text{-}field}(a, b).\; \operatorname{compose}(f, k)) \in (\operatorname{hom\text{-}field}(a, b) \to \operatorname{hom\text{-}field}(c, b)) \end{array}}}
\vnbprop{hom-ringoid-post-type}{\vnbstmt{\begin{array}{@{}l@{}} \forall a, b, c, h. \\ \quad (\operatorname{is\text{-}ringoid}(a) \wedge (h \in \operatorname{hom\text{-}ringoid}(b, c))) \Rightarrow  \\ \quad \quad (\lambda f \in \operatorname{hom\text{-}ringoid}(a, b).\; \operatorname{compose}(h, f)) \in (\operatorname{hom\text{-}ringoid}(a, b) \to \operatorname{hom\text{-}ringoid}(a, c)) \end{array}}}
\vnbprop{hom-ringoid-pre-type}{\vnbstmt{\begin{array}{@{}l@{}} \forall a, b, c, k. \\ \quad (\operatorname{is\text{-}ringoid}(b) \wedge (k \in \operatorname{hom\text{-}ringoid}(c, a))) \Rightarrow  \\ \quad \quad (\lambda f \in \operatorname{hom\text{-}ringoid}(a, b).\; \operatorname{compose}(f, k)) \in (\operatorname{hom\text{-}ringoid}(a, b) \to \operatorname{hom\text{-}ringoid}(c, b)) \end{array}}}
\vnbprop{hom-setoid-post-type}{\vnbstmt{\begin{array}{@{}l@{}} \forall a, b, c, h. \\ \quad (\operatorname{is\text{-}setoid}(a) \wedge (h \in \operatorname{hom\text{-}setoid}(b, c))) \Rightarrow  \\ \quad \quad (\lambda f \in \operatorname{hom\text{-}setoid}(a, b).\; \operatorname{compose}(h, f)) \in (\operatorname{hom\text{-}setoid}(a, b) \to \operatorname{hom\text{-}setoid}(a, c)) \end{array}}}
\vnbprop{hom-setoid-pre-type}{\vnbstmt{\begin{array}{@{}l@{}} \forall a, b, c, k. \\ \quad (\operatorname{is\text{-}setoid}(b) \wedge (k \in \operatorname{hom\text{-}setoid}(c, a))) \Rightarrow  \\ \quad \quad (\lambda f \in \operatorname{hom\text{-}setoid}(a, b).\; \operatorname{compose}(f, k)) \in (\operatorname{hom\text{-}setoid}(a, b) \to \operatorname{hom\text{-}setoid}(c, b)) \end{array}}}

\vnbbookfile{theorem-library/qq-field-is-field}{theorem-library/qq-field-is-field}
\noindent{\small qq-field-is-field.scm -- IS-FIELD(QQ-FIELD), PROVED.}\par
\vnbprop{qq-recip-nonzero}{\vnbstmt{\forall a \in \mathbb{Q}.\; (\neg (a = 0) \Rightarrow \neg (\operatorname{recip}(a) = 0))}}
\vnbprop{qq-field-mul-inv-type}{\vnbstmt{\begin{array}{@{}l@{}} (\lambda x \in \operatorname{difference}(\mathbb{Q}, \operatorname{singleton}(0)).\; \operatorname{recip}(x)) \\ \in\; (\operatorname{difference}(\mathbb{Q}, \operatorname{singleton}(0)) \to \operatorname{difference}(\mathbb{Q}, \operatorname{singleton}(0))) \end{array}}}
\vnbprop{qq-field-is-field}{\vnbstmt{\operatorname{is\text{-}field}(\operatorname{qq\text{-}field})}}

\vnbbookfile{theorem-library/rake-finsum-laws}{theorem-library/rake-finsum-laws}
\noindent{\small theorem-library/rake-finsum-laws.scm -- the ADDITIVE algebra of FINSUM, proven.  Rake batch J (2026-09-17), part 1 of 2; part 2 is theorem-library/rake-finsum-laws2.scm (the three distribution laws and the interval back-peel), which loads AFTER this file and cites two of its}\par
\vnbprop{enum-fam-value}{\vnbstmt{\begin{array}{@{}l@{}} \forall n \in \mathbb{N}, g, u, a.; \forall i \in \mathrm{OS}(n). \\ \quad (\operatorname{enum\text{-}fam}(g, u, a, n)(i) \simeq u(a(i))) \end{array}}}
\vnbprop{sum-ag-type-ptwise}{\vnbstmt{\begin{array}{@{}l@{}} \forall n \in \mathbb{N}, g, a. \\ \quad (\operatorname{is\text{-}abelian\text{-}group}(g) \wedge \forall i \in \mathrm{OS}(n).\; (a(i) \in \operatorname{carr}(g))) \Rightarrow  \\ \quad \quad \operatorname{sum\text{-}ag}(g, a, n) \in \operatorname{carr}(g) \end{array}}}
\vnbprop{abelian-group-opr-interchange}{\vnbstmt{\begin{array}{@{}l@{}} \forall s.; \forall a, b, c, d \in \operatorname{carr}(s). \\ \quad \operatorname{is\text{-}abelian\text{-}group}(s) \Rightarrow  \\ \quad \quad \operatorname{opr}(s)(\operatorname{opr}(s)(a, b), \operatorname{opr}(s)(c, d)) = \operatorname{opr}(s)(\operatorname{opr}(s)(a, c), \operatorname{opr}(s)(b, d)) \end{array}}}
\vnbprop{finsum-congruence-q}{\vnbstmt{\begin{array}{@{}l@{}} \forall a.; \forall s \in \mathbf{Set}, f, g. \\ \quad (\operatorname{is\text{-}abelian\text{-}group}(a) \wedge ((|s| \in \mathbb{N}) \wedge \forall z \in s.\; (f(z) = g(z)))) \Rightarrow  \\ \quad \quad (\operatorname{finsum}(a, f, s) \simeq \operatorname{finsum}(a, g, s)) \end{array}}}
\vnbprop{finsum-congruence-guarded}{\vnbstmt{\begin{array}{@{}l@{}} \forall a.; \forall s \in \mathbf{Set}, f \in (s \to \operatorname{carr}(a)), g. \\ \quad (\operatorname{is\text{-}abelian\text{-}group}(a) \wedge ((|s| \in \mathbb{N}) \wedge \forall z \in s.\; (f(z) = g(z)))) \Rightarrow  \\ \quad \quad \operatorname{finsum}(a, f, s) = \operatorname{finsum}(a, g, s) \end{array}}}
\vnbprop{finsum-type-ptwise}{\vnbstmt{\begin{array}{@{}l@{}} \forall g.; \forall s \in \mathbf{Set}, f. \\ \quad (\operatorname{is\text{-}abelian\text{-}group}(g) \wedge ((|s| \in \mathbb{N}) \wedge \forall z \in s.\; (f(z) \in \operatorname{carr}(g)))) \Rightarrow  \\ \quad \quad \operatorname{finsum}(g, f, s) \in \operatorname{carr}(g) \end{array}}}
\vnbprop{finsum-congruence}{\vnbstmt{\begin{array}{@{}l@{}} \forall a.; \forall s \in \mathbf{Set}, f, g. \\ \quad \operatorname{is\text{-}abelian\text{-}group}(a) \wedge \\ \quad |s| \in \mathbb{N} \wedge \\ \quad \forall z \in s.\; (f(z) = g(z)) \wedge \\ \quad \forall z \in s.\; (f(z) \in \operatorname{carr}(a)) \Rightarrow \\ \quad \quad \operatorname{finsum}(a, f, s) = \operatorname{finsum}(a, g, s) \end{array}}}
\vnbprop{sum-ag-add-ind}{\vnbstmt{\begin{array}{@{}l@{}} \forall n \in \mathbb{N}, g, a, b, c. \\ \quad \operatorname{is\text{-}abelian\text{-}group}(g) \wedge \\ \quad \forall i \in \mathrm{OS}(n).\; (a(i) \in \operatorname{carr}(g)) \wedge \\ \quad \forall i \in \mathrm{OS}(n).\; (b(i) \in \operatorname{carr}(g)) \wedge \\ \quad \forall i \in \mathrm{OS}(n).\; (c(i) \simeq \operatorname{opr}(g)(a(i), b(i))) \Rightarrow \\ \quad \quad \operatorname{sum\text{-}ag}(g, c, n) = \operatorname{opr}(g)(\operatorname{sum\text{-}ag}(g, a, n), \operatorname{sum\text{-}ag}(g, b, n)) \end{array}}}
\vnbprop{finsum-add-ag}{\vnbstmt{\begin{array}{@{}l@{}} \forall g.; \forall s \in \mathbf{Set}, f, h \in (s \to \operatorname{carr}(g)). \\ \quad (\operatorname{is\text{-}abelian\text{-}group}(g) \wedge (|s| \in \mathbb{N})) \Rightarrow  \\ \quad \quad \operatorname{finsum}(g, (\lambda z \in s.\; \operatorname{opr}(g)(f(z), h(z))), s) = \operatorname{opr}(g)(\operatorname{finsum}(g, f, s), \operatorname{finsum}(g, h, s)) \end{array}}}
\vnbprop{sum-ag-two-support-ind}{\vnbstmt{\begin{array}{@{}l@{}} \forall n \in \mathbb{N}, g, a, p, q. \\ \quad \operatorname{is\text{-}abelian\text{-}group}(g) \wedge \\ \quad p \in \mathrm{OS}(n) \wedge \\ \quad q \in \mathrm{OS}(n) \wedge \\ \quad \neg (p = q) \wedge \\ \quad a(p) \in \operatorname{carr}(g) \wedge \\ \quad a(q) \in \operatorname{carr}(g) \wedge \\ \quad \forall i \in \mathrm{OS}(n).\; ((\neg (i = p) \wedge \neg (i = q)) \Rightarrow (a(i) = \operatorname{iden}(g))) \Rightarrow \\ \quad \quad \operatorname{sum\text{-}ag}(g, a, n) = \operatorname{opr}(g)(a(p), a(q)) \end{array}}}
\vnbprop{finsum-two-support}{\vnbstmt{\begin{array}{@{}l@{}} \forall g.; \forall s \in \mathbf{Set}, f \in (s \to \operatorname{carr}(g)), a, b \in s. \\ \quad \operatorname{is\text{-}abelian\text{-}group}(g) \wedge \\ \quad |s| \in \mathbb{N} \wedge \\ \quad \neg (a = b) \wedge \\ \quad \forall j \in s.\; ((\neg (j = a) \wedge \neg (j = b)) \Rightarrow (f(j) = \operatorname{iden}(g))) \Rightarrow \\ \quad \quad \operatorname{finsum}(g, f, s) = \operatorname{opr}(g)(f(a), f(b)) \end{array}}}

\vnbbookfile{theorem-library/rake-finsum-laws2}{theorem-library/rake-finsum-laws2}
\noindent{\small theorem-library/rake-finsum-laws2.scm -- rake batch J, part 2: the three DISTRIBUTION laws of FINSUM (a ring element on either side, a module scalar) and the interval back-peel.  Part 1 is theorem-library/rake-finsum-laws.scm, which must load FIRST: this file cites its `enum-fam-value' and}\par
\vnbprop{sum-ag-ring-left-ind}{\vnbstmt{\begin{array}{@{}l@{}} \forall n \in \mathbb{N}, g.; \forall r \in \operatorname{carr}(g), a, b. \\ \quad \operatorname{is\text{-}ring}(g) \wedge \\ \quad \forall i \in \mathrm{OS}(n).\; (a(i) \in \operatorname{carr}(g)) \wedge \\ \quad \forall i \in \mathrm{OS}(n).\; (b(i) \simeq \operatorname{mul}(g)(r, a(i))) \Rightarrow \\ \quad \quad \operatorname{sum\text{-}ag}(\operatorname{ring\text{-}additive\text{-}ag}(g), b, n) = \operatorname{mul}(g)(r, \operatorname{sum\text{-}ag}(\operatorname{ring\text{-}additive\text{-}ag}(g), a, n)) \end{array}}}
\vnbprop{finsum-ring-distrib-left-gen}{\vnbstmt{\begin{array}{@{}l@{}} \forall g.; \forall r \in \operatorname{carr}(g), s \in \mathbf{Set}, f \in (s \to \operatorname{carr}(g)). \\ \quad (\operatorname{is\text{-}ring}(g) \wedge (|s| \in \mathbb{N})) \Rightarrow  \\ \quad \quad \operatorname{mul}(g)(r, \operatorname{finsum}(\operatorname{ring\text{-}additive\text{-}ag}(g), f, s)) \\ \quad \quad =\; \operatorname{finsum}(\operatorname{ring\text{-}additive\text{-}ag}(g), (\lambda z \in s.\; \operatorname{mul}(g)(r, f(z))), s) \end{array}}}
\vnbprop{sum-ag-ring-right-ind}{\vnbstmt{\begin{array}{@{}l@{}} \forall n \in \mathbb{N}, g.; \forall r \in \operatorname{carr}(g), a, b. \\ \quad \operatorname{is\text{-}ring}(g) \wedge \\ \quad \forall i \in \mathrm{OS}(n).\; (a(i) \in \operatorname{carr}(g)) \wedge \\ \quad \forall i \in \mathrm{OS}(n).\; (b(i) \simeq \operatorname{mul}(g)(a(i), r)) \Rightarrow \\ \quad \quad \operatorname{sum\text{-}ag}(\operatorname{ring\text{-}additive\text{-}ag}(g), b, n) = \operatorname{mul}(g)(\operatorname{sum\text{-}ag}(\operatorname{ring\text{-}additive\text{-}ag}(g), a, n), r) \end{array}}}
\vnbprop{finsum-ring-distrib-right-gen}{\vnbstmt{\begin{array}{@{}l@{}} \forall g.; \forall r \in \operatorname{carr}(g), s \in \mathbf{Set}, f \in (s \to \operatorname{carr}(g)). \\ \quad (\operatorname{is\text{-}ring}(g) \wedge (|s| \in \mathbb{N})) \Rightarrow  \\ \quad \quad \operatorname{mul}(g)(\operatorname{finsum}(\operatorname{ring\text{-}additive\text{-}ag}(g), f, s), r) \\ \quad \quad =\; \operatorname{finsum}(\operatorname{ring\text{-}additive\text{-}ag}(g), (\lambda z \in s.\; \operatorname{mul}(g)(f(z), r)), s) \end{array}}}
\vnbprop{finsum-interval-peel}{\vnbstmt{\begin{array}{@{}l@{}} \forall g.; \forall n \in \mathbb{N}, f \in (\operatorname{interval}(1, (n + 1)) \to \operatorname{carr}(g)). \\ \quad \operatorname{is\text{-}abelian\text{-}group}(g) \Rightarrow  \\ \quad \quad \operatorname{finsum}(g, f, \operatorname{interval}(1, (n + 1))) = \operatorname{opr}(g)(\operatorname{finsum}(g, f, \operatorname{interval}(1, n)), f((n + 1))) \end{array}}}

\vnbbookfile{theorem-library/rake-finsum-core}{theorem-library/rake-finsum-core}
\noindent{\small theorem-library/rake-finsum-core.scm -- the FINSUM CORE, proven. Rake batch M (2026-09-17).  Nine of the batch's ten leaves; the tenth (finsum-embed) is BLOCKED and its obstacle is stated at the end of this header.}\par
\vnbprop{finsum-ring-distrib-left}{\vnbstmt{\begin{array}{@{}l@{}} \forall g.; \forall r \in \operatorname{carr}(g), s \in \mathbf{Set}, f \in (s \to \operatorname{carr}(g)). \\ \quad (\operatorname{is\text{-}commutative\text{-}ring}(g) \wedge (|s| \in \mathbb{N})) \Rightarrow  \\ \quad \quad \operatorname{mul}(g)(r, \operatorname{finsum}(\operatorname{commutative\text{-}ring\text{-}additive\text{-}ag}(g), f, s)) \\ \quad \quad =\; \operatorname{finsum}(\operatorname{commutative\text{-}ring\text{-}additive\text{-}ag}(g), (\lambda z \in s.\; \operatorname{mul}(g)(r, f(z))), s) \end{array}}}
\vnbprop{finsum-ring-distrib-right}{\vnbstmt{\begin{array}{@{}l@{}} \forall g.; \forall r \in \operatorname{carr}(g), s \in \mathbf{Set}, f \in (s \to \operatorname{carr}(g)). \\ \quad (\operatorname{is\text{-}commutative\text{-}ring}(g) \wedge (|s| \in \mathbb{N})) \Rightarrow  \\ \quad \quad \operatorname{mul}(g)(\operatorname{finsum}(\operatorname{commutative\text{-}ring\text{-}additive\text{-}ag}(g), f, s), r) \\ \quad \quad =\; \operatorname{finsum}(\operatorname{commutative\text{-}ring\text{-}additive\text{-}ag}(g), (\lambda z \in s.\; \operatorname{mul}(g)(f(z), r)), s) \end{array}}}
\vnbprop{finsum-singleton}{\vnbstmt{\begin{array}{@{}l@{}} \forall g.; \forall x \in \mathbf{Set}, f \in (\{x\} \to \operatorname{carr}(g)). \\ \quad \operatorname{is\text{-}abelian\text{-}group}(g) \Rightarrow (\operatorname{finsum}(g, f, \{x\}) = f(x)) \end{array}}}
\vnbprop{finsum-ord-peel}{\vnbstmt{\begin{array}{@{}l@{}} \forall g.; \forall n \in \mathbb{N}, f \in (\mathrm{OS}((n + 1)) \to \operatorname{carr}(g)). \\ \quad \operatorname{is\text{-}abelian\text{-}group}(g) \Rightarrow  \\ \quad \quad \operatorname{finsum}(g, f, \mathrm{OS}((n + 1))) = \operatorname{opr}(g)(\operatorname{finsum}(g, f, \mathrm{OS}(n)), f(n)) \end{array}}}
\vnbprop{comm-monoid-opr-comm}{\vnbstmt{\begin{array}{@{}l@{}} \forall s.; \forall a, b \in \operatorname{carr}(s). \\ \quad \operatorname{is\text{-}comm\text{-}monoid}(s) \Rightarrow (\operatorname{opr}(s)(a, b) = \operatorname{opr}(s)(b, a)) \end{array}}}
\vnbprop{comm-monoid-assoc}{\vnbstmt{\begin{array}{@{}l@{}} \forall s.; \forall a, b, c \in \operatorname{carr}(s). \\ \quad \operatorname{is\text{-}comm\text{-}monoid}(s) \Rightarrow  \\ \quad \quad \operatorname{opr}(s)(\operatorname{opr}(s)(a, b), c) = \operatorname{opr}(s)(a, \operatorname{opr}(s)(b, c)) \end{array}}}
\vnbprop{comm-monoid-left-id}{\vnbstmt{\forall s.; \forall a \in \operatorname{carr}(s).\; (\operatorname{is\text{-}comm\text{-}monoid}(s) \Rightarrow (\operatorname{opr}(s)(\operatorname{iden}(s), a) = a))}}
\vnbprop{comm-monoid-opr-interchange}{\vnbstmt{\begin{array}{@{}l@{}} \forall s.; \forall a, b, c, d \in \operatorname{carr}(s). \\ \quad \operatorname{is\text{-}comm\text{-}monoid}(s) \Rightarrow  \\ \quad \quad \operatorname{opr}(s)(\operatorname{opr}(s)(a, b), \operatorname{opr}(s)(c, d)) = \operatorname{opr}(s)(\operatorname{opr}(s)(a, c), \operatorname{opr}(s)(b, d)) \end{array}}}
\vnbprop{sum-ag-comm-monoid-type-ptwise}{\vnbstmt{\begin{array}{@{}l@{}} \forall n \in \mathbb{N}, m, g. \\ \quad (\operatorname{is\text{-}comm\text{-}monoid}(m) \wedge \forall i \in \mathrm{OS}(n).\; (g(i) \in \operatorname{carr}(m))) \Rightarrow  \\ \quad \quad \operatorname{sum\text{-}ag}(m, g, n) \in \operatorname{carr}(m) \end{array}}}
\vnbprop{sum-ag-comm-monoid-add-ind}{\vnbstmt{\begin{array}{@{}l@{}} \forall n \in \mathbb{N}, m, a, b, c. \\ \quad \operatorname{is\text{-}comm\text{-}monoid}(m) \wedge \\ \quad \forall i \in \mathrm{OS}(n).\; (a(i) \in \operatorname{carr}(m)) \wedge \\ \quad \forall i \in \mathrm{OS}(n).\; (b(i) \in \operatorname{carr}(m)) \wedge \\ \quad \forall i \in \mathrm{OS}(n).\; (c(i) \simeq \operatorname{opr}(m)(a(i), b(i))) \Rightarrow \\ \quad \quad \operatorname{sum\text{-}ag}(m, c, n) = \operatorname{opr}(m)(\operatorname{sum\text{-}ag}(m, a, n), \operatorname{sum\text{-}ag}(m, b, n)) \end{array}}}
\vnbprop{finsum-add}{\vnbstmt{\begin{array}{@{}l@{}} \forall m.; \forall s \in \mathbf{Set}, f, h \in (s \to \operatorname{carr}(m)). \\ \quad (\operatorname{is\text{-}comm\text{-}monoid}(m) \wedge (|s| \in \mathbb{N})) \Rightarrow  \\ \quad \quad \operatorname{finsum}(m, (\lambda z \in s.\; \operatorname{opr}(m)(f(z), h(z))), s) = \operatorname{opr}(m)(\operatorname{finsum}(m, f, s), \operatorname{finsum}(m, h, s)) \end{array}}}
\vnbprop{sum-ag-fubini-ind}{\vnbstmt{\begin{array}{@{}l@{}} \forall m \in \mathbb{N}, g.; \forall n \in \mathbb{N}, f, \mathit{gg}, s, \mathit{cs}. \\ \quad \operatorname{is\text{-}abelian\text{-}group}(g) \wedge \\ \quad \forall i \in \mathrm{OS}(m), j \in \mathrm{OS}(n).\; (f(i)(j) \in \operatorname{carr}(g)) \wedge \\ \quad \forall i \in \mathrm{OS}(m), j \in \mathrm{OS}(n).\; (\operatorname{gg}(j)(i) \simeq f(i)(j)) \wedge \\ \quad \forall i \in \mathrm{OS}(m).\; (s(i) \simeq \operatorname{sum\text{-}ag}(g, f(i), n)) \wedge \\ \quad \mathit{cs} \in (\mathrm{OS}(n) \to \operatorname{carr}(g)) \wedge \\ \quad \forall j \in \mathrm{OS}(n).\; (\operatorname{cs}(j) \simeq \operatorname{sum\text{-}ag}(g, \operatorname{gg}(j), m)) \wedge \\ \quad \forall \mathit{fbi} \in \mathrm{OS}(m).\; (s(\mathit{fbi}) \in \operatorname{carr}(g)) \wedge \\ \quad \forall \mathit{fbj} \in \mathrm{OS}(n).\; (\operatorname{cs}(\mathit{fbj}) \in \operatorname{carr}(g)) \Rightarrow \\ \quad \quad \operatorname{sum\text{-}ag}(g, s, m) = \operatorname{sum\text{-}ag}(g, \mathit{cs}, n) \end{array}}}
\vnbprop{finsum-fubini}{\vnbstmt{\begin{array}{@{}l@{}} \forall g.; \forall x, y \in \mathbf{Set}, f \in ((x \times y) \to \operatorname{carr}(g)). \\ \quad (\operatorname{is\text{-}abelian\text{-}group}(g) \wedge ((|x| \in \mathbb{N}) \wedge (|y| \in \mathbb{N}))) \Rightarrow  \\ \quad \quad \operatorname{finsum}(g, (\lambda i \in x.\; \operatorname{finsum}(g, (\lambda j \in y.\; f({[}i, j])), y)), x) \\ \quad \quad =\; \operatorname{finsum}(g, (\lambda j \in y.\; \operatorname{finsum}(g, (\lambda i \in x.\; f({[}i, j])), x)), y) \end{array}}}
\vnbprop{ras-inv}{\vnbstmt{\forall a.\; (\operatorname{is\text{-}ring}(a) \Rightarrow (\operatorname{inv}(\operatorname{ring\text{-}additive\text{-}ag}(a)) = \operatorname{neg}(a)))}}
\vnbprop{commutative-ring-mul-comm}{\vnbstmt{\begin{array}{@{}l@{}} \forall s.; \forall a, b \in \operatorname{carr}(s). \\ \quad \operatorname{is\text{-}commutative\text{-}ring}(s) \Rightarrow (\operatorname{mul}(s)(a, b) = \operatorname{mul}(s)(b, a)) \end{array}}}
\vnbprop{ring-scalar-zz-nn}{\vnbstmt{\begin{array}{@{}l@{}} \forall k \in \mathbb{N}, g.; \forall r, a \in \operatorname{carr}(g). \\ \quad \operatorname{is\text{-}commutative\text{-}ring}(g) \Rightarrow  \\ \quad \quad \operatorname{mul}(g)(r, \operatorname{zz\text{-}act}(\operatorname{ring\text{-}additive\text{-}ag}(g), k, a)) = \operatorname{zz\text{-}act}(\operatorname{ring\text{-}additive\text{-}ag}(g), k, \operatorname{mul}(g)(r, a)) \end{array}}}
\vnbprop{finsum-ring-scalar-zz}{\vnbstmt{\begin{array}{@{}l@{}} \forall g.; \forall r \in \operatorname{carr}(g), c \in \mathbb{Z}, a \in \operatorname{carr}(g). \\ \quad \operatorname{is\text{-}commutative\text{-}ring}(g) \Rightarrow  \\ \quad \quad \operatorname{mul}(g)(r, \operatorname{zz\text{-}act}(\operatorname{commutative\text{-}ring\text{-}additive\text{-}ag}(g), c, a)) \\ \quad \quad =\; \operatorname{zz\text{-}act}(\operatorname{commutative\text{-}ring\text{-}additive\text{-}ag}(g), c, \operatorname{mul}(g)(r, a)) \end{array}}}

\vnbbookfile{theorem-library/rake-finsum-cm-ptwise}{theorem-library/rake-finsum-cm-ptwise}
\noindent{\small theorem-library/rake-finsum-cm-ptwise.scm -- finsum-comm-monoid-type-ptwise, the COMMUTATIVE-MONOID twin of `finsum-type-ptwise'.  Rake batch 5c-K.}\par
\vnbprop{finsum-comm-monoid-type-ptwise}{\vnbstmt{\begin{array}{@{}l@{}} \forall m.; \forall s \in \mathbf{Set}, f. \\ \quad (\operatorname{is\text{-}comm\text{-}monoid}(m) \wedge ((|s| \in \mathbb{N}) \wedge \forall z \in s.\; (f(z) \in \operatorname{carr}(m)))) \Rightarrow  \\ \quad \quad \operatorname{finsum}(m, f, s) \in \operatorname{carr}(m) \end{array}}}

\vnbbookfile{theorem-library/rake-finsum-fubini-c}{theorem-library/rake-finsum-fubini-c}
\noindent{\small rake-finsum-fubini-c.scm -- finsum-fubini-c, moved out of rake-finsum-laws.scm on 2026-09-17 (rake batch M) because it cites finsum-fubini, now PROVEN in rake-finsum-core.scm, which loads after rake-finsum-laws.  Window: after rake-finsum-core, before finsum-fiber (its earliest citer).  Helper prefix rkj-}\par
\vnbprop{finsum-fubini-c}{\vnbstmt{\begin{array}{@{}l@{}} \forall g.; \forall x, y \in \mathbf{Set}, f \in ((x \times y) \to \operatorname{carr}(g)). \\ \quad (\operatorname{is\text{-}abelian\text{-}group}(g) \wedge ((|x| \in \mathbb{N}) \wedge (|y| \in \mathbb{N}))) \Rightarrow  \\ \quad \quad \operatorname{finsum}(g, (\lambda i \in x.\; \operatorname{finsum}(g, (\lambda j \in y.\; f({[}i, j])), y)), x) \\ \quad \quad =\; \operatorname{finsum}(g, (\lambda j \in y.\; \operatorname{finsum}(g, (\lambda i \in x.\; f({[}i, j])), x)), y) \end{array}}}

\vnbbookfile{theorem-library/rake-combinatorics}{theorem-library/rake-combinatorics}
\noindent{\small rake-combinatorics.scm -- BATCH P of the 2026-09-17 rake: NN combinatorics, the ring powers, the powerset insert-split, and the two CHOOSE boundary values.  Seventeen asserted supports, PROVEN; every statement is copied LITERALLY from its site.}\par
\vnbprop{crmcm-op}{\vnbstmt{\begin{array}{@{}l@{}} \forall r. \\ \quad \operatorname{is\text{-}commutative\text{-}ring}(r) \Rightarrow  \\ \quad \quad \operatorname{opr}(\operatorname{commutative\text{-}ring\text{-}multiplicative\text{-}cm}(r)) = \operatorname{mul}(r) \end{array}}}
\vnbprop{crmcm-iden}{\vnbstmt{\begin{array}{@{}l@{}} \forall r. \\ \quad \operatorname{is\text{-}commutative\text{-}ring}(r) \Rightarrow  \\ \quad \quad \operatorname{iden}(\operatorname{commutative\text{-}ring\text{-}multiplicative\text{-}cm}(r)) = \operatorname{one}(r) \end{array}}}
\vnbprop{ring-power-succ-left}{\vnbstmt{\begin{array}{@{}l@{}} \forall r.; \forall x \in \operatorname{carr}(r), n \in \mathbb{N}. \\ \quad \operatorname{is\text{-}commutative\text{-}ring}(r) \Rightarrow  \\ \quad \quad \operatorname{ring\text{-}power}(r, x, (n + 1)) = \operatorname{mul}(r)(x, \operatorname{ring\text{-}power}(r, x, n)) \end{array}}}
\vnbprop{ring-power-one}{\vnbstmt{\begin{array}{@{}l@{}} \forall r.; \forall x \in \operatorname{carr}(r). \\ \quad \operatorname{is\text{-}commutative\text{-}ring}(r) \Rightarrow (\operatorname{ring\text{-}power}(r, x, 1) = x) \end{array}}}
\vnbprop{ring-power-succ}{\vnbstmt{\begin{array}{@{}l@{}} \forall r.; \forall x \in \operatorname{carr}(r), n \in \mathbb{N}. \\ \quad \operatorname{is\text{-}commutative\text{-}ring}(r) \Rightarrow  \\ \quad \quad \operatorname{ring\text{-}power}(r, x, (n + 1)) = \operatorname{mul}(r)(\operatorname{ring\text{-}power}(r, x, n), x) \end{array}}}
\vnbprop{ring-mul-interchange}{\vnbstmt{\begin{array}{@{}l@{}} \forall r.; \forall a, b, c, d \in \operatorname{carr}(r). \\ \quad \operatorname{is\text{-}commutative\text{-}ring}(r) \Rightarrow  \\ \quad \quad \operatorname{mul}(r)(\operatorname{mul}(r)(a, b), \operatorname{mul}(r)(c, d)) = \operatorname{mul}(r)(\operatorname{mul}(r)(a, c), \operatorname{mul}(r)(b, d)) \end{array}}}
\vnbprop{ring-power-add-ind}{\vnbstmt{\begin{array}{@{}l@{}} \forall k \in \mathbb{N}, r.; \forall x \in \operatorname{carr}(r), j \in \mathbb{N}. \\ \quad \operatorname{is\text{-}commutative\text{-}ring}(r) \Rightarrow  \\ \quad \quad \operatorname{ring\text{-}power}(r, x, (j + k)) = \operatorname{mul}(r)(\operatorname{ring\text{-}power}(r, x, j), \operatorname{ring\text{-}power}(r, x, k)) \end{array}}}
\vnbprop{ring-power-add}{\vnbstmt{\begin{array}{@{}l@{}} \forall r.; \forall x \in \operatorname{carr}(r), j, k \in \mathbb{N}. \\ \quad \operatorname{is\text{-}commutative\text{-}ring}(r) \Rightarrow  \\ \quad \quad \operatorname{ring\text{-}power}(r, x, (j + k)) = \operatorname{mul}(r)(\operatorname{ring\text{-}power}(r, x, j), \operatorname{ring\text{-}power}(r, x, k)) \end{array}}}
\vnbprop{ring-power-mult-ind}{\vnbstmt{\begin{array}{@{}l@{}} \forall n \in \mathbb{N}, r.; \forall x, y \in \operatorname{carr}(r). \\ \quad \operatorname{is\text{-}commutative\text{-}ring}(r) \Rightarrow  \\ \quad \quad \operatorname{ring\text{-}power}(r, \operatorname{mul}(r)(x, y), n) = \operatorname{mul}(r)(\operatorname{ring\text{-}power}(r, x, n), \operatorname{ring\text{-}power}(r, y, n)) \end{array}}}
\vnbprop{ring-power-mult}{\vnbstmt{\begin{array}{@{}l@{}} \forall r.; \forall x, y \in \operatorname{carr}(r), n \in \mathbb{N}. \\ \quad \operatorname{is\text{-}commutative\text{-}ring}(r) \Rightarrow  \\ \quad \quad \operatorname{ring\text{-}power}(r, \operatorname{mul}(r)(x, y), n) = \operatorname{mul}(r)(\operatorname{ring\text{-}power}(r, x, n), \operatorname{ring\text{-}power}(r, y, n)) \end{array}}}
\vnbprop{prod-ring-empty}{\vnbstmt{\begin{array}{@{}l@{}} \forall r, f. \\ \quad \operatorname{is\text{-}commutative\text{-}ring}(r) \Rightarrow (\operatorname{prod\text{-}ring}(r, f, \emptyset) = \operatorname{one}(r)) \end{array}}}
\vnbprop{prod-ring-insert}{\vnbstmt{\begin{array}{@{}l@{}} \forall r.; \forall x, k \in \mathbf{Set}, f \in ((x \cup \{k\}) \to \operatorname{carr}(r)). \\ \quad (\operatorname{is\text{-}commutative\text{-}ring}(r) \wedge ((|x| \in \mathbb{N}) \wedge \neg (k \in x))) \Rightarrow  \\ \quad \quad \operatorname{prod\text{-}ring}(r, f, (x \cup \{k\})) = \operatorname{mul}(r)(\operatorname{prod\text{-}ring}(r, f, x), f(k)) \end{array}}}
\vnbprop{prod-ring-singleton}{\vnbstmt{\begin{array}{@{}l@{}} \forall r.; \forall x \in \mathbf{Set}, f \in (\{x\} \to \operatorname{carr}(r)). \\ \quad \operatorname{is\text{-}commutative\text{-}ring}(r) \Rightarrow (\operatorname{prod\text{-}ring}(r, f, \{x\}) = f(x)) \end{array}}}

\vnbbookfile{theorem-library/ccr1109a2-geom}{theorem-library/ccr1109a2-geom}
\noindent{\small ccr1109a2-geom.scm -- wave of 2026-11-09/10 (agent CCR-1109-A2), wired 2026-11-10 from scratchpad/ccr1109a2/ccr1109a2-geom.scm. Library: algebra.}\par
\vnbprop{ring-geometric-partial-sum-ind}{\vnbstmt{\begin{array}{@{}l@{}} \forall n \in \mathbb{N}, r.; \forall z \in \operatorname{carr}(r). \\ \quad \operatorname{is\text{-}commutative\text{-}ring}(r) \Rightarrow  \\ \quad \quad \operatorname{mul}(r)(\operatorname{add}(r)(\operatorname{one}(r), \operatorname{neg}(r)(z)), \\ \quad \quad \quad \operatorname{sum\text{-}ag}(\operatorname{ring\text{-}additive\text{-}ag}(r), (\lambda k \in \mathbb{N}.\; \operatorname{ring\text{-}power}(r, z, k)), n)) \\ \quad \quad =\; \operatorname{add}(r)(\operatorname{one}(r), \operatorname{neg}(r)(\operatorname{ring\text{-}power}(r, z, n))) \end{array}}}
\vnbprop{ring-geometric-partial-sum}{\vnbstmt{\begin{array}{@{}l@{}} \forall r.; \forall z \in \operatorname{carr}(r), n \in \mathbb{N}. \\ \quad \operatorname{is\text{-}commutative\text{-}ring}(r) \Rightarrow  \\ \quad \quad \operatorname{mul}(r)(\operatorname{add}(r)(\operatorname{one}(r), \operatorname{neg}(r)(z)), \\ \quad \quad \quad \operatorname{sum\text{-}ag}(\operatorname{ring\text{-}additive\text{-}ag}(r), (\lambda k \in \mathbb{N}.\; \operatorname{ring\text{-}power}(r, z, k)), n)) \\ \quad \quad =\; \operatorname{add}(r)(\operatorname{one}(r), \operatorname{neg}(r)(\operatorname{ring\text{-}power}(r, z, n))) \end{array}}}

\vnbbookfile{theorem-library/rake-finsum-union}{theorem-library/rake-finsum-union}
\noindent{\small theorem-library/rake-finsum-union.scm -- the two set/card bricks that three rake batches built inline, the POINTWISE finsum recurrences they make possible, the DISJOINT-UNION law for FINSUM, and `finsum-embed'. Rake batch 5c, assignment 5c-M.  Helper prefix `r7m-'; every helper is}\par
\vnbprop{card-insert-converse}{\vnbstmt{\begin{array}{@{}l@{}} \forall a, x \in \mathbf{Set}, k \in \mathbb{N}. \\ \quad (\neg (x \in a) \wedge (|(a \cup \{x\})| = (k + 1))) \Rightarrow (|a| = k) \end{array}}}
\vnbprop{remove-restore}{\vnbstmt{\forall a \in \mathbf{Set}, x \in a.\; ((\operatorname{difference}(a, \{x\}) \cup \{x\}) = a)}}
\vnbprop{finsum-insert-ptwise}{\vnbstmt{\begin{array}{@{}l@{}} \forall g.; \forall s, k \in \mathbf{Set}, f. \\ \quad \operatorname{is\text{-}abelian\text{-}group}(g) \wedge \\ \quad |s| \in \mathbb{N} \wedge \\ \quad \neg (k \in s) \wedge \\ \quad \forall z \in (s \cup \{k\}).\; (f(z) \in \operatorname{carr}(g)) \Rightarrow \\ \quad \quad \operatorname{finsum}(g, f, (s \cup \{k\})) = \operatorname{opr}(g)(\operatorname{finsum}(g, f, s), f(k)) \end{array}}}
\vnbprop{finsum-union-disjoint}{\vnbstmt{\begin{array}{@{}l@{}} \forall g.; \forall s, t \in \mathbf{Set}, f. \\ \quad \operatorname{is\text{-}abelian\text{-}group}(g) \wedge \\ \quad |s| \in \mathbb{N} \wedge \\ \quad |t| \in \mathbb{N} \wedge \\ \quad \forall z \in t.\; \neg (z \in s) \wedge \\ \quad \forall z \in (s \cup t).\; (f(z) \in \operatorname{carr}(g)) \Rightarrow \\ \quad \quad \operatorname{finsum}(g, f, (s \cup t)) = \operatorname{opr}(g)(\operatorname{finsum}(g, f, s), \operatorname{finsum}(g, f, t)) \end{array}}}
\vnbprop{finsum-union-disjoint-fun}{\vnbstmt{\begin{array}{@{}l@{}} \forall g.; \forall s, t \in \mathbf{Set}, f \in ((s \cup t) \to \operatorname{carr}(g)). \\ \quad (\operatorname{is\text{-}abelian\text{-}group}(g) \wedge ((|s| \in \mathbb{N}) \wedge ((|t| \in \mathbb{N}) \wedge ((s \cap t) = \emptyset)))) \Rightarrow  \\ \quad \quad \operatorname{finsum}(g, f, (s \cup t)) = \operatorname{opr}(g)(\operatorname{finsum}(g, f, s), \operatorname{finsum}(g, f, t)) \end{array}}}
\vnbprop{finsum-all-id-ptwise}{\vnbstmt{\begin{array}{@{}l@{}} \forall g.; \forall s \in \mathbf{Set}, f. \\ \quad (\operatorname{is\text{-}abelian\text{-}group}(g) \wedge ((|s| \in \mathbb{N}) \wedge \forall z \in s.\; (f(z) = \operatorname{iden}(g)))) \Rightarrow  \\ \quad \quad \operatorname{finsum}(g, f, s) = \operatorname{iden}(g) \end{array}}}
\vnbprop{finsum-embed}{\vnbstmt{\begin{array}{@{}l@{}} \forall g.; \forall s, a \in \mathbf{Set}, f \in (a \to \operatorname{carr}(g)). \\ \quad \operatorname{is\text{-}abelian\text{-}group}(g) \wedge \\ \quad |s| \in \mathbb{N} \wedge \\ \quad |a| \in \mathbb{N} \wedge \\ \quad s \subseteq a \wedge \\ \quad \forall z \in a.\; (\neg (z \in s) \Rightarrow (f(z) = \operatorname{iden}(g))) \Rightarrow \\ \quad \quad \operatorname{finsum}(g, f, a) = \operatorname{finsum}(g, f, s) \end{array}}}

\vnbbookfile{theorem-library/rake-finsum-cm-union}{theorem-library/rake-finsum-cm-union}
\noindent{\small theorem-library/rake-finsum-cm-union.scm -- the COMMUTATIVE-MONOID mirror of the FINSUM congruence / insertion / disjoint-union block, plus the two set identities that a peeling induction over POWER(X) needs. Rake batch 5c, assignment 5c-T.  Helper prefix `r7t-'; every helper is}\par
\vnbprop{finsum-cm-congruence-q}{\vnbstmt{\begin{array}{@{}l@{}} \forall m.; \forall s \in \mathbf{Set}, f, g. \\ \quad (\operatorname{is\text{-}comm\text{-}monoid}(m) \wedge ((|s| \in \mathbb{N}) \wedge \forall z \in s.\; (f(z) = g(z)))) \Rightarrow  \\ \quad \quad (\operatorname{finsum}(m, f, s) \simeq \operatorname{finsum}(m, g, s)) \end{array}}}
\vnbprop{finsum-cm-insert-ptwise}{\vnbstmt{\begin{array}{@{}l@{}} \forall m.; \forall s, k \in \mathbf{Set}, f. \\ \quad \operatorname{is\text{-}comm\text{-}monoid}(m) \wedge \\ \quad |s| \in \mathbb{N} \wedge \\ \quad \neg (k \in s) \wedge \\ \quad \forall z \in (s \cup \{k\}).\; (f(z) \in \operatorname{carr}(m)) \Rightarrow \\ \quad \quad \operatorname{finsum}(m, f, (s \cup \{k\})) = \operatorname{opr}(m)(\operatorname{finsum}(m, f, s), f(k)) \end{array}}}
\vnbprop{finsum-cm-union-disjoint}{\vnbstmt{\begin{array}{@{}l@{}} \forall m.; \forall s, t \in \mathbf{Set}, f. \\ \quad \operatorname{is\text{-}comm\text{-}monoid}(m) \wedge \\ \quad |s| \in \mathbb{N} \wedge \\ \quad |t| \in \mathbb{N} \wedge \\ \quad \forall z \in t.\; \neg (z \in s) \wedge \\ \quad \forall z \in (s \cup t).\; (f(z) \in \operatorname{carr}(m)) \Rightarrow \\ \quad \quad \operatorname{finsum}(m, f, (s \cup t)) = \operatorname{opr}(m)(\operatorname{finsum}(m, f, s), \operatorname{finsum}(m, f, t)) \end{array}}}
\vnbprop{insert-difference-outside}{\vnbstmt{\begin{array}{@{}l@{}} \forall x, s.; \forall k \in \mathbf{Set}. \\ \quad \neg (k \in s) \Rightarrow  \\ \quad \quad \operatorname{difference}((x \cup \{k\}), s) = (\operatorname{difference}(x, s) \cup \{k\}) \end{array}}}
\vnbprop{insert-difference-cancel}{\vnbstmt{\begin{array}{@{}l@{}} \forall x, s.; \forall k \in \mathbf{Set}. \\ \quad \neg (k \in x) \Rightarrow  \\ \quad \quad \operatorname{difference}((x \cup \{k\}), (s \cup \{k\})) = \operatorname{difference}(x, s) \end{array}}}

\vnbbookfile{theorem-library/rake-sum-set-defined}{theorem-library/rake-sum-set-defined}
\noindent{\small theorem-library/rake-sum-set-defined.scm -- SUM-SET DEFINED, and its four characterising laws PROVEN from the definition.  Rake batch 5c, assignment 5c-S (the user's decision (B) of 2026-09-18).  Helper prefix `r7s-'.}\par
\vnbprop{sum-set-empty}{\vnbstmt{\forall r, f.\; (\operatorname{is\text{-}ring}(r) \Rightarrow (\operatorname{sum\text{-}set}(r, \emptyset, f) \simeq \operatorname{zero}(r)))}}
\vnbprop{finsum-singleton-ptwise}{\vnbstmt{\begin{array}{@{}l@{}} \forall g.; \forall x \in \mathbf{Set}, f. \\ \quad (\operatorname{is\text{-}abelian\text{-}group}(g) \wedge (f(x) \in \operatorname{carr}(g))) \Rightarrow (\operatorname{finsum}(g, f, \{x\}) = f(x)) \end{array}}}
\vnbprop{sum-set-singleton-defined}{\vnbstmt{\begin{array}{@{}l@{}} \forall r.; \forall x \in \mathbf{Set}, f. \\ \quad (\operatorname{is\text{-}ring}(r) \wedge (f(x) \in \operatorname{carr}(r))) \Rightarrow (\operatorname{sum\text{-}set}(r, \{x\}, f) = f(x)) \end{array}}}
\vnbprop{sum-set-type-defined}{\vnbstmt{\begin{array}{@{}l@{}} \forall r, x, s, f. \\ \quad (\operatorname{is\text{-}ring}(r) \wedge ((x \in \mathbf{Set}) \wedge ((s \subseteq x) \wedge ((f \in (x \to \operatorname{carr}(r))) \wedge (|s| \in \mathbb{N}))))) \Rightarrow  \\ \quad \quad \operatorname{sum\text{-}set}(r, s, f) \in \operatorname{carr}(r) \end{array}}}
\vnbprop{sum-set-disjoint-union-defined}{\vnbstmt{\begin{array}{@{}l@{}} \forall r, x, a, b, f. \\ \quad \operatorname{is\text{-}ring}(r) \wedge \\ \quad x \in \mathbf{Set} \wedge \\ \quad a \in \mathbf{Set} \wedge \\ \quad b \in \mathbf{Set} \wedge \\ \quad (a \cap b) = \emptyset \wedge \\ \quad (a \cup b) \subseteq x \wedge \\ \quad f \in (x \to \operatorname{carr}(r)) \wedge \\ \quad |a| \in \mathbb{N} \wedge \\ \quad |b| \in \mathbb{N} \Rightarrow \\ \quad \quad \operatorname{sum\text{-}set}(r, (a \cup b), f) = \operatorname{add}(r)(\operatorname{sum\text{-}set}(r, a, f), \operatorname{sum\text{-}set}(r, b, f)) \end{array}}}

\vnbbookfile{theorem-library/rake-prod-set-defined}{theorem-library/rake-prod-set-defined}
\noindent{\small theorem-library/rake-prod-set-defined.scm -- PROD-SET DEFINED, and its four characterising laws PROVEN from the definition.  Rake batch 5c, assignment 5c-P (the user's decision of 2026-09-19).  Helper prefix `psd-'.}\par
\vnbprop{prod-set-empty}{\vnbstmt{\forall m, f.\; (\operatorname{prod\text{-}set}(m, \emptyset, f) \simeq \operatorname{iden}(m))}}
\vnbprop{finsum-cm-singleton-ptwise}{\vnbstmt{\begin{array}{@{}l@{}} \forall m.; \forall x \in \mathbf{Set}, f. \\ \quad (\operatorname{is\text{-}comm\text{-}monoid}(m) \wedge (f(x) \in \operatorname{carr}(m))) \Rightarrow (\operatorname{finsum}(m, f, \{x\}) = f(x)) \end{array}}}
\vnbprop{prod-set-singleton-defined}{\vnbstmt{\begin{array}{@{}l@{}} \forall m.; \forall x \in \mathbf{Set}, f. \\ \quad (\operatorname{is\text{-}comm\text{-}monoid}(m) \wedge (f(x) \in \operatorname{carr}(m))) \Rightarrow (\operatorname{prod\text{-}set}(m, \{x\}, f) = f(x)) \end{array}}}
\vnbprop{prod-set-type-defined}{\vnbstmt{\begin{array}{@{}l@{}} \forall m, x, s, f. \\ \quad (\operatorname{is\text{-}comm\text{-}monoid}(m) \wedge ((x \in \mathbf{Set}) \wedge ((s \subseteq x) \wedge ((f \in (x \to \operatorname{carr}(m))) \wedge (|s| \in \mathbb{N}))))) \Rightarrow  \\ \quad \quad \operatorname{prod\text{-}set}(m, s, f) \in \operatorname{carr}(m) \end{array}}}
\vnbprop{prod-set-disjoint-union-defined}{\vnbstmt{\begin{array}{@{}l@{}} \forall m, a, b, f. \\ \quad \operatorname{is\text{-}comm\text{-}monoid}(m) \wedge \\ \quad a \in \mathbf{Set} \wedge \\ \quad b \in \mathbf{Set} \wedge \\ \quad (a \cap b) = \emptyset \wedge \\ \quad f \in ((a \cup b) \to \operatorname{carr}(m)) \wedge \\ \quad |a| \in \mathbb{N} \wedge \\ \quad |b| \in \mathbb{N} \Rightarrow \\ \quad \quad \operatorname{prod\text{-}set}(m, (a \cup b), f) = \operatorname{opr}(m)(\operatorname{prod\text{-}set}(m, a, f), \operatorname{prod\text{-}set}(m, b, f)) \end{array}}}

\vnbbookfile{theorem-library/rake-sum-set-scalar}{theorem-library/rake-sum-set-scalar}
\noindent{\small rake-sum-set-scalar.scm -- the two scalar pull-outs for SUM-SET, PROVEN.}\par
\vnbprop{sum-set-left-scalar}{\vnbstmt{\begin{array}{@{}l@{}} \forall s.; \forall a \in \operatorname{carr}(s), x \in \mathbf{Set}, f \in (x \to \operatorname{carr}(s)). \\ \quad (\operatorname{is\text{-}ring}(s) \wedge (|x| \in \mathbb{N})) \Rightarrow  \\ \quad \quad \operatorname{mul}(s)(a, \operatorname{sum\text{-}set}(s, x, f)) = \operatorname{sum\text{-}set}(s, x, (\lambda z \in x.\; \operatorname{mul}(s)(a, f(z)))) \end{array}}}
\vnbprop{sum-set-right-scalar}{\vnbstmt{\begin{array}{@{}l@{}} \forall s.; \forall b \in \operatorname{carr}(s), x \in \mathbf{Set}, f \in (x \to \operatorname{carr}(s)). \\ \quad (\operatorname{is\text{-}ring}(s) \wedge (|x| \in \mathbb{N})) \Rightarrow  \\ \quad \quad \operatorname{mul}(s)(\operatorname{sum\text{-}set}(s, x, f), b) = \operatorname{sum\text{-}set}(s, x, (\lambda z \in x.\; \operatorname{mul}(s)(f(z), b))) \end{array}}}

\vnbbookfile{theorem-library/fin-subset-monoid}{theorem-library/fin-subset-monoid}
\noindent{\small fin-subset-monoid.scm -- FIN-SUBSETS(a) is a commutative monoid under union.}\par
\par\medskip\noindent\textbf{Definition} (\texttt{fin-subset-monoid}).\;$\operatorname{fin\text{-}subset\text{-}monoid}(\operatorname{a\_}) :=$
\begin{equation*}
{[}\operatorname{fin\text{-}subsets}(\operatorname{a\_}), (\lambda {[}\operatorname{x\_}, \operatorname{y\_}] \in (\operatorname{fin\text{-}subsets}(\operatorname{a\_}) \times \operatorname{fin\text{-}subsets}(\operatorname{a\_})).\; (\operatorname{x\_} \cup \operatorname{y\_})), \emptyset]
\end{equation*}

\vnbprop{fsm-carr}{\vnbstmt{\forall a \in \mathbf{Set}.\; (\operatorname{carr}(\operatorname{fin\text{-}subset\text{-}monoid}(a)) \simeq \operatorname{fin\text{-}subsets}(a))}}
\vnbprop{fsm-opr}{\vnbstmt{\begin{array}{@{}l@{}} \forall a \in \mathbf{Set}. \\ \quad (\operatorname{opr}(\operatorname{fin\text{-}subset\text{-}monoid}(a)) \simeq (\lambda {[}\operatorname{x\_}, \operatorname{y\_}] \in (\operatorname{fin\text{-}subsets}(a) \times \operatorname{fin\text{-}subsets}(a)).\; (\operatorname{x\_} \cup \operatorname{y\_}))) \end{array}}}
\vnbprop{fsm-iden}{\vnbstmt{\forall a \in \mathbf{Set}.\; (\operatorname{iden}(\operatorname{fin\text{-}subset\text{-}monoid}(a)) \simeq \emptyset)}}
\vnbprop{fin-subset-monoid-is-comm-monoid}{\vnbstmt{\forall a \in \mathbf{Set}.\; \operatorname{is\text{-}comm\text{-}monoid}(\operatorname{fin\text{-}subset\text{-}monoid}(a))}}

\vnbbookfile{theorem-library/rake-algebra3}{theorem-library/rake-algebra3}
\noindent{\small theorem-library/rake-algebra3.scm -- rake batch 5, batch V (algebra leftovers).}\par
\vnbprop{qq-nonneg-is-fraction}{\vnbstmt{\begin{array}{@{}l@{}} \forall q \in \mathbb{Q}. \\ \quad (0 \leq q) \Rightarrow \exists a, b \in \mathbb{N}.\; (\neg (b = 0) \wedge ((q \cdot b) = a)) \end{array}}}
\vnbprop{card-bijection-eq}{\vnbstmt{\begin{array}{@{}l@{}} \forall m, d, i. \\ \quad ((m \in \mathbf{Set}) \wedge ((|m| \in \mathbb{N}) \wedge (i \in \mathrm{Bij}(m, d)))) \Rightarrow (|d| = |m|) \end{array}}}
\vnbprop{finsum-reindex}{\vnbstmt{\begin{array}{@{}l@{}} \forall m.; \forall s, t \in \mathbf{Set}, i \in \mathrm{Bij}(t, s), f \in (s \to \operatorname{carr}(m)). \\ \quad (\operatorname{is\text{-}comm\text{-}monoid}(m) \wedge ((|s| \in \mathbb{N}) \wedge (|t| \in \mathbb{N}))) \Rightarrow  \\ \quad \quad \operatorname{finsum}(m, f, s) = \operatorname{finsum}(m, (\lambda z \in t.\; f(i(z))), t) \end{array}}}
\vnbprop{finsum-reindex-ag}{\vnbstmt{\begin{array}{@{}l@{}} \forall g.; \forall s, t \in \mathbf{Set}, i \in \mathrm{Bij}(t, s), f \in (s \to \operatorname{carr}(g)). \\ \quad (\operatorname{is\text{-}abelian\text{-}group}(g) \wedge ((|s| \in \mathbb{N}) \wedge (|t| \in \mathbb{N}))) \Rightarrow  \\ \quad \quad \operatorname{finsum}(g, f, s) = \operatorname{finsum}(g, (\lambda z \in t.\; f(i(z))), t) \end{array}}}

\vnbbookfile{theorem-library/poly-tail-zero}{theorem-library/poly-tail-zero}
\noindent{\small poly-tail-zero.scm -- "a polynomial is a sequence that is eventually zero".}\par
\vnbprop{poly-tail-zero-fwd}{\vnbstmt{\begin{array}{@{}l@{}} \forall a, f. \\ \quad (\operatorname{is\text{-}ring}(a) \wedge (f \in \operatorname{carr}(\operatorname{poly}(a)))) \Rightarrow  \\ \quad \quad \exists n \in \mathbb{N}.\; \forall k \in \mathbb{N}.\; ((n \leq k) \Rightarrow (f(k) = \operatorname{zero}(a))) \end{array}}}
\vnbprop{poly-tail-zero-bwd}{\vnbstmt{\begin{array}{@{}l@{}} \forall a, f. \\ \quad \operatorname{is\text{-}ring}(a) \wedge \\ \quad f \in \operatorname{carr}(a)^{\mathbb{N}} \wedge \\ \quad \exists n \in \mathbb{N}.\; \forall k \in \mathbb{N}.\; ((n \leq k) \Rightarrow (f(k) = \operatorname{zero}(a))) \Rightarrow \\ \quad \quad f \in \operatorname{carr}(\operatorname{poly}(a)) \end{array}}}
\vnbprop{poly-tail-zero}{\vnbstmt{\begin{array}{@{}l@{}} \forall a, f. \\ \quad \operatorname{is\text{-}ring}(a) \Rightarrow  \\ \quad \quad f \in \operatorname{carr}(\operatorname{poly}(a)) \\ \quad \quad \Leftrightarrow\; ((f \in \operatorname{carr}(a)^{\mathbb{N}}) \wedge \exists n \in \mathbb{N}.\; \forall k \in \mathbb{N}.\; ((n \leq k) \Rightarrow (f(k) = \operatorname{zero}(a)))) \end{array}}}

\vnbbookfile{theorem-library/finsum-fiber}{theorem-library/finsum-fiber}
\noindent{\small finsum-fiber.scm -- FIBERED FUBINI: group the terms of a finite sum by the fibers of an arbitrary map, and the CARTESIAN projection brick it needs.}\par
\vnbprop{finsum-fiber}{\vnbstmt{\begin{array}{@{}l@{}} \forall g.; \forall s, t \in \mathbf{Set}, i \in (s \to t), f \in (s \to \operatorname{carr}(g)). \\ \quad (\operatorname{is\text{-}abelian\text{-}group}(g) \wedge ((|s| \in \mathbb{N}) \wedge (|t| \in \mathbb{N}))) \Rightarrow  \\ \quad \quad \operatorname{finsum}(g, f, s) = \operatorname{finsum}(g, (\lambda a \in t.\; \operatorname{finsum}(g, f, \{b \in s : (i(b) = a)\})), t) \end{array}}}

\vnbbookfile{theorem-library/monalg-laws}{theorem-library/monalg-laws}
\noindent{\small monalg-laws.scm -- the monoid-algebra laws that go through without reindexing, PROVEN.}\par
\vnbprop{monoid-carrier-is-set}{\vnbstmt{\forall m.\; (\operatorname{is\text{-}monoid}(m) \Rightarrow (\operatorname{carr}(m) \in \mathbf{Set}))}}
\vnbprop{monalg-add-apply}{\vnbstmt{\begin{array}{@{}l@{}} \forall a, m, f, g.; \forall x \in \operatorname{carr}(m). \\ \quad (\operatorname{monalg\text{-}add}(a, m, f, g)(x) \simeq \operatorname{add}(a)(f(x), g(x))) \end{array}}}
\vnbprop{monalg-add-fun}{\vnbstmt{\begin{array}{@{}l@{}} \forall a, m, f, g. \\ \quad (\operatorname{is\text{-}ring}(a) \wedge (\operatorname{is\text{-}monoid}(m) \wedge ((f \in \operatorname{finsupp}(a, m)) \wedge (g \in \operatorname{finsupp}(a, m))))) \Rightarrow  \\ \quad \quad \operatorname{monalg\text{-}add}(a, m, f, g) \in \operatorname{finsupp}(a, m) \end{array}}}

\vnbbookfile{theorem-library/poly-zero}{theorem-library/poly-zero}
\noindent{\small poly-zero.scm -- the ZERO of a monoid algebra, pointwise, and the EXTENSIONALITY BRIDGE for polynomials.  Nothing here is asserted; every statement is `proven modulo 0'.}\par
\vnbprop{monalg-zero-apply}{\vnbstmt{\forall a, m.; \forall x \in \operatorname{carr}(m).\; (\operatorname{monalg\text{-}zero}(a, m)(x) \simeq \operatorname{zero}(a))}}
\vnbprop{monalg-zero-fun}{\vnbstmt{\begin{array}{@{}l@{}} \forall a, m. \\ \quad ((\operatorname{carr}(m) \in \mathbf{Set}) \wedge (\operatorname{zero}(a) \in \operatorname{carr}(a))) \Rightarrow  \\ \quad \quad \operatorname{monalg\text{-}zero}(a, m) \in (\operatorname{carr}(m) \to \operatorname{carr}(a)) \end{array}}}
\vnbprop{poly-zero-unfold}{\vnbstmt{\forall a.\; (\operatorname{is\text{-}ring}(a) \Rightarrow (\operatorname{zero}(\operatorname{poly}(a)) \simeq \operatorname{monalg\text{-}zero}(a, \operatorname{nn\text{-}add\text{-}monoid})))}}
\vnbprop{poly-zero-apply}{\vnbstmt{\forall a, k.\; ((\operatorname{is\text{-}ring}(a) \wedge (k \in \mathbb{N})) \Rightarrow (\operatorname{zero}(\operatorname{poly}(a))(k) = \operatorname{zero}(a)))}}
\vnbprop{poly-zero-fun}{\vnbstmt{\forall a.\; (\operatorname{is\text{-}ring}(a) \Rightarrow (\operatorname{zero}(\operatorname{poly}(a)) \in (\mathbb{N} \to \operatorname{carr}(a))))}}
\vnbprop{poly-zero-iff}{\vnbstmt{\begin{array}{@{}l@{}} \forall a, p. \\ \quad (\operatorname{is\text{-}ring}(a) \wedge (p \in \operatorname{carr}(\operatorname{poly}(a)))) \Rightarrow  \\ \quad \quad (p = \operatorname{zero}(\operatorname{poly}(a))) \Leftrightarrow \forall k \in \mathbb{N}.\; (p(k) = \operatorname{zero}(a)) \end{array}}}

\vnbbookfile{theorem-library/monalg-is-ring}{theorem-library/monalg-is-ring}
\noindent{\small monalg-is-ring.scm -- the monoid algebra A[M] is a ring (structure-library/polynomial.scm's umbrella seed), PROVEN, with every additive law, both distributive laws, both unit laws and the typing of the convolution proven on the way.  The one law still chained to its seed is the associativity of convolution, ...}\par
\vnbprop{monalg-carr}{\vnbstmt{\begin{array}{@{}l@{}} \forall a, m. \\ \quad (\operatorname{is\text{-}ring}(a) \wedge \operatorname{is\text{-}monoid}(m)) \Rightarrow (\operatorname{carr}(\operatorname{monalg}(a, m)) \simeq \operatorname{finsupp}(a, m)) \end{array}}}
\vnbprop{monalg-add-op}{\vnbstmt{\begin{array}{@{}l@{}} \forall a, m. \\ \quad (\operatorname{is\text{-}ring}(a) \wedge \operatorname{is\text{-}monoid}(m)) \Rightarrow  \\ \quad \quad (\operatorname{add}(\operatorname{monalg}(a, m)) \simeq (\lambda {[}f, g] \in (\operatorname{finsupp}(a, m) \times \operatorname{finsupp}(a, m)).\; \operatorname{monalg\text{-}add}(a, m, f, g))) \end{array}}}
\vnbprop{monalg-mul-op}{\vnbstmt{\begin{array}{@{}l@{}} \forall a, m. \\ \quad (\operatorname{is\text{-}ring}(a) \wedge \operatorname{is\text{-}monoid}(m)) \Rightarrow  \\ \quad \quad (\operatorname{mul}(\operatorname{monalg}(a, m)) \simeq (\lambda {[}f, g] \in (\operatorname{finsupp}(a, m) \times \operatorname{finsupp}(a, m)).\; \operatorname{monalg\text{-}mul}(a, m, f, g))) \end{array}}}
\vnbprop{monalg-neg-op}{\vnbstmt{\begin{array}{@{}l@{}} \forall a, m. \\ \quad (\operatorname{is\text{-}ring}(a) \wedge \operatorname{is\text{-}monoid}(m)) \Rightarrow  \\ \quad \quad (\operatorname{neg}(\operatorname{monalg}(a, m)) \simeq (\lambda f \in \operatorname{finsupp}(a, m).\; \operatorname{monalg\text{-}neg}(a, m, f))) \end{array}}}
\vnbprop{monalg-zero-slot}{\vnbstmt{\begin{array}{@{}l@{}} \forall a, m. \\ \quad (\operatorname{is\text{-}ring}(a) \wedge \operatorname{is\text{-}monoid}(m)) \Rightarrow (\operatorname{zero}(\operatorname{monalg}(a, m)) \simeq \operatorname{monalg\text{-}zero}(a, m)) \end{array}}}
\vnbprop{monalg-one-slot}{\vnbstmt{\begin{array}{@{}l@{}} \forall a, m. \\ \quad (\operatorname{is\text{-}ring}(a) \wedge \operatorname{is\text{-}monoid}(m)) \Rightarrow (\operatorname{one}(\operatorname{monalg}(a, m)) \simeq \operatorname{monalg\text{-}one}(a, m)) \end{array}}}
\vnbprop{finsupp-in-set}{\vnbstmt{\forall a, m.\; ((\operatorname{is\text{-}ring}(a) \wedge \operatorname{is\text{-}monoid}(m)) \Rightarrow (\operatorname{finsupp}(a, m) \in \mathbf{Set}))}}
\vnbprop{monalg-ext}{\vnbstmt{\begin{array}{@{}l@{}} \forall a, m, f, g. \\ \quad ((f \in \operatorname{finsupp}(a, m)) \wedge ((g \in \operatorname{finsupp}(a, m)) \wedge \forall x \in \operatorname{carr}(m).\; (f(x) = g(x)))) \Rightarrow  \\ \quad \quad f = g \end{array}}}
\vnbprop{monalg-neg-apply}{\vnbstmt{\forall a, m, f.; \forall x \in \operatorname{carr}(m).\; (\operatorname{monalg\text{-}neg}(a, m, f)(x) \simeq \operatorname{neg}(a)(f(x)))}}
\vnbprop{monalg-one-apply}{\vnbstmt{\begin{array}{@{}l@{}} \forall a, m.; \forall x \in \operatorname{carr}(m). \\ \quad (\operatorname{monalg\text{-}one}(a, m)(x) \simeq \left(\text{if }(x = \operatorname{iden}(m))\text{ then }\operatorname{one}(a)\text{ else }\operatorname{zero}(a)\right)) \end{array}}}
\vnbprop{ring-neg-zero}{\vnbstmt{\forall a.\; (\operatorname{is\text{-}ring}(a) \Rightarrow (\operatorname{neg}(a)(\operatorname{zero}(a)) = \operatorname{zero}(a)))}}
\vnbprop{monalg-neg-fun}{\vnbstmt{\begin{array}{@{}l@{}} \forall a, m, f. \\ \quad (\operatorname{is\text{-}ring}(a) \wedge (\operatorname{is\text{-}monoid}(m) \wedge (f \in \operatorname{finsupp}(a, m)))) \Rightarrow  \\ \quad \quad \operatorname{monalg\text{-}neg}(a, m, f) \in \operatorname{finsupp}(a, m) \end{array}}}
\vnbprop{monalg-zero-in}{\vnbstmt{\begin{array}{@{}l@{}} \forall a, m. \\ \quad (\operatorname{is\text{-}ring}(a) \wedge \operatorname{is\text{-}monoid}(m)) \Rightarrow (\operatorname{monalg\text{-}zero}(a, m) \in \operatorname{finsupp}(a, m)) \end{array}}}
\vnbprop{monalg-one-in}{\vnbstmt{\begin{array}{@{}l@{}} \forall a, m. \\ \quad (\operatorname{is\text{-}ring}(a) \wedge \operatorname{is\text{-}monoid}(m)) \Rightarrow (\operatorname{monalg\text{-}one}(a, m) \in \operatorname{finsupp}(a, m)) \end{array}}}
\vnbprop{monalg-add-assoc}{\vnbstmt{\begin{array}{@{}l@{}} \forall a, m, f, g, h. \\ \quad (\operatorname{is\text{-}ring}(a) \wedge (\operatorname{is\text{-}monoid}(m) \wedge ((f \in \operatorname{finsupp}(a, m)) \wedge ((g \in \operatorname{finsupp}(a, m)) \wedge (h \in \operatorname{finsupp}(a, m)))))) \Rightarrow  \\ \quad \quad \operatorname{monalg\text{-}add}(a, m, \operatorname{monalg\text{-}add}(a, m, f, g), h) = \operatorname{monalg\text{-}add}(a, m, f, \operatorname{monalg\text{-}add}(a, m, g, h)) \end{array}}}
\vnbprop{monalg-add-comm}{\vnbstmt{\begin{array}{@{}l@{}} \forall a, m, f, g. \\ \quad (\operatorname{is\text{-}ring}(a) \wedge (\operatorname{is\text{-}monoid}(m) \wedge ((f \in \operatorname{finsupp}(a, m)) \wedge (g \in \operatorname{finsupp}(a, m))))) \Rightarrow  \\ \quad \quad \operatorname{monalg\text{-}add}(a, m, f, g) = \operatorname{monalg\text{-}add}(a, m, g, f) \end{array}}}
\vnbprop{monalg-add-zero-left}{\vnbstmt{\begin{array}{@{}l@{}} \forall a, m, f. \\ \quad (\operatorname{is\text{-}ring}(a) \wedge (\operatorname{is\text{-}monoid}(m) \wedge (f \in \operatorname{finsupp}(a, m)))) \Rightarrow  \\ \quad \quad \operatorname{monalg\text{-}add}(a, m, \operatorname{monalg\text{-}zero}(a, m), f) = f \end{array}}}
\vnbprop{monalg-add-zero-right}{\vnbstmt{\begin{array}{@{}l@{}} \forall a, m, f. \\ \quad (\operatorname{is\text{-}ring}(a) \wedge (\operatorname{is\text{-}monoid}(m) \wedge (f \in \operatorname{finsupp}(a, m)))) \Rightarrow  \\ \quad \quad \operatorname{monalg\text{-}add}(a, m, f, \operatorname{monalg\text{-}zero}(a, m)) = f \end{array}}}
\vnbprop{monalg-add-neg-left}{\vnbstmt{\begin{array}{@{}l@{}} \forall a, m, f. \\ \quad (\operatorname{is\text{-}ring}(a) \wedge (\operatorname{is\text{-}monoid}(m) \wedge (f \in \operatorname{finsupp}(a, m)))) \Rightarrow  \\ \quad \quad \operatorname{monalg\text{-}add}(a, m, \operatorname{monalg\text{-}neg}(a, m, f), f) = \operatorname{monalg\text{-}zero}(a, m) \end{array}}}
\vnbprop{monalg-add-neg-right}{\vnbstmt{\begin{array}{@{}l@{}} \forall a, m, f. \\ \quad (\operatorname{is\text{-}ring}(a) \wedge (\operatorname{is\text{-}monoid}(m) \wedge (f \in \operatorname{finsupp}(a, m)))) \Rightarrow  \\ \quad \quad \operatorname{monalg\text{-}add}(a, m, f, \operatorname{monalg\text{-}neg}(a, m, f)) = \operatorname{monalg\text{-}zero}(a, m) \end{array}}}
\vnbprop{raag-carr}{\vnbstmt{\forall a.\; (\operatorname{is\text{-}ring}(a) \Rightarrow (\operatorname{carr}(\operatorname{ring\text{-}additive\text{-}ag}(a)) \simeq \operatorname{carr}(a)))}}
\vnbprop{raag-opr}{\vnbstmt{\forall a.\; (\operatorname{is\text{-}ring}(a) \Rightarrow (\operatorname{opr}(\operatorname{ring\text{-}additive\text{-}ag}(a)) \simeq \operatorname{add}(a)))}}
\vnbprop{raag-iden}{\vnbstmt{\forall a.\; (\operatorname{is\text{-}ring}(a) \Rightarrow (\operatorname{iden}(\operatorname{ring\text{-}additive\text{-}ag}(a)) \simeq \operatorname{zero}(a)))}}
\vnbprop{monalg-mul-apply}{\vnbstmt{\begin{array}{@{}l@{}} \forall a, m, f, g.; \forall x \in \operatorname{carr}(m). \\ \quad (\operatorname{monalg\text{-}mul}(a, m, f, g)(x) \simeq \operatorname{finsum}(\operatorname{ring\text{-}additive\text{-}ag}(a), (\lambda p \in \{p \in (\operatorname{supp}(a, m, f) \times \operatorname{supp}(a, m, g)) : (\operatorname{opr}(m)(\operatorname{nth}(1, p), \operatorname{nth}(2, p)) = x)\}.\; \operatorname{mul}(a)(f(\operatorname{nth}(1, p)), g(\operatorname{nth}(2, p)))), \{p \in (\operatorname{supp}(a, m, f) \times \operatorname{supp}(a, m, g)) : (\operatorname{opr}(m)(\operatorname{nth}(1, p), \operatorname{nth}(2, p)) = x)\})) \end{array}}}
\vnbprop{monalg-mul-value-in-carr}{\vnbstmt{\begin{array}{@{}l@{}} \forall a, m, f, g, x. \\ \quad (\operatorname{is\text{-}ring}(a) \wedge (\operatorname{is\text{-}monoid}(m) \wedge ((f \in \operatorname{finsupp}(a, m)) \wedge ((g \in \operatorname{finsupp}(a, m)) \wedge (x \in \operatorname{carr}(m)))))) \Rightarrow  \\ \quad \quad \operatorname{monalg\text{-}mul}(a, m, f, g)(x) \in \operatorname{carr}(a) \end{array}}}
\vnbprop{monalg-mul-fun}{\vnbstmt{\begin{array}{@{}l@{}} \forall a, m, f, g. \\ \quad (\operatorname{is\text{-}ring}(a) \wedge (\operatorname{is\text{-}monoid}(m) \wedge ((f \in \operatorname{finsupp}(a, m)) \wedge (g \in \operatorname{finsupp}(a, m))))) \Rightarrow  \\ \quad \quad \operatorname{monalg\text{-}mul}(a, m, f, g) \in \operatorname{finsupp}(a, m) \end{array}}}
\vnbprop{supp-add-subset}{\vnbstmt{\begin{array}{@{}l@{}} \forall a, m, g, h, z. \\ \quad (\operatorname{is\text{-}ring}(a) \wedge ((g \in \operatorname{finsupp}(a, m)) \wedge ((h \in \operatorname{finsupp}(a, m)) \wedge (z \in \operatorname{supp}(a, m, \operatorname{monalg\text{-}add}(a, m, g, h)))))) \Rightarrow  \\ \quad \quad z \in (\operatorname{supp}(a, m, g) \cup \operatorname{supp}(a, m, h)) \end{array}}}
\vnbprop{monalg-distrib-left-at}{\vnbstmt{\begin{array}{@{}l@{}} \forall a, m, f, g, h, x. \\ \quad \operatorname{is\text{-}ring}(a) \wedge \\ \quad \operatorname{is\text{-}monoid}(m) \wedge \\ \quad f \in \operatorname{finsupp}(a, m) \wedge \\ \quad g \in \operatorname{finsupp}(a, m) \wedge \\ \quad h \in \operatorname{finsupp}(a, m) \wedge \\ \quad x \in \operatorname{carr}(m) \Rightarrow \\ \quad \quad \operatorname{finsum}(\operatorname{ring\text{-}additive\text{-}ag}(a), \\ \quad \quad \quad (\lambda p \in \{p \in (\operatorname{supp}(a, m, f) \times \operatorname{supp}(a, m, \operatorname{monalg\text{-}add}(a, m, g, h))) : (\operatorname{opr}(m)(\operatorname{nth}(1, p), \operatorname{nth}(2, p)) = x)\}.\; \\ \quad \quad \quad \quad \operatorname{mul}(a)(f(\operatorname{nth}(1, p)), \operatorname{monalg\text{-}add}(a, m, g, h)(\operatorname{nth}(2, p)))), \\ \quad \quad \quad \{p \in (\operatorname{supp}(a, m, f) \times \operatorname{supp}(a, m, \operatorname{monalg\text{-}add}(a, m, g, h))) : (\operatorname{opr}(m)(\operatorname{nth}(1, p), \operatorname{nth}(2, p)) = x)\}) \\ \quad \quad =\; \operatorname{add}(a)(\operatorname{finsum}(\operatorname{ring\text{-}additive\text{-}ag}(a), \\ \quad \quad \quad \quad (\lambda p \in \{p \in (\operatorname{supp}(a, m, f) \times \operatorname{supp}(a, m, g)) : (\operatorname{opr}(m)(\operatorname{nth}(1, p), \operatorname{nth}(2, p)) = x)\}.\; \\ \quad \quad \quad \quad \quad \operatorname{mul}(a)(f(\operatorname{nth}(1, p)), g(\operatorname{nth}(2, p)))), \\ \quad \quad \quad \quad \{p \in (\operatorname{supp}(a, m, f) \times \operatorname{supp}(a, m, g)) : (\operatorname{opr}(m)(\operatorname{nth}(1, p), \operatorname{nth}(2, p)) = x)\}), \\ \quad \quad \quad \operatorname{finsum}(\operatorname{ring\text{-}additive\text{-}ag}(a), \\ \quad \quad \quad \quad (\lambda p \in \{p \in (\operatorname{supp}(a, m, f) \times \operatorname{supp}(a, m, h)) : (\operatorname{opr}(m)(\operatorname{nth}(1, p), \operatorname{nth}(2, p)) = x)\}.\; \\ \quad \quad \quad \quad \quad \operatorname{mul}(a)(f(\operatorname{nth}(1, p)), h(\operatorname{nth}(2, p)))), \\ \quad \quad \quad \quad \{p \in (\operatorname{supp}(a, m, f) \times \operatorname{supp}(a, m, h)) : (\operatorname{opr}(m)(\operatorname{nth}(1, p), \operatorname{nth}(2, p)) = x)\})) \end{array}}}
\vnbprop{monalg-distrib-right-at}{\vnbstmt{\begin{array}{@{}l@{}} \forall a, m, f, g, h, x. \\ \quad \operatorname{is\text{-}ring}(a) \wedge \\ \quad \operatorname{is\text{-}monoid}(m) \wedge \\ \quad f \in \operatorname{finsupp}(a, m) \wedge \\ \quad g \in \operatorname{finsupp}(a, m) \wedge \\ \quad h \in \operatorname{finsupp}(a, m) \wedge \\ \quad x \in \operatorname{carr}(m) \Rightarrow \\ \quad \quad \operatorname{finsum}(\operatorname{ring\text{-}additive\text{-}ag}(a), \\ \quad \quad \quad (\lambda p \in \{p \in (\operatorname{supp}(a, m, \operatorname{monalg\text{-}add}(a, m, f, g)) \times \operatorname{supp}(a, m, h)) : (\operatorname{opr}(m)(\operatorname{nth}(1, p), \operatorname{nth}(2, p)) = x)\}.\; \\ \quad \quad \quad \quad \operatorname{mul}(a)(\operatorname{monalg\text{-}add}(a, m, f, g)(\operatorname{nth}(1, p)), h(\operatorname{nth}(2, p)))), \\ \quad \quad \quad \{p \in (\operatorname{supp}(a, m, \operatorname{monalg\text{-}add}(a, m, f, g)) \times \operatorname{supp}(a, m, h)) : (\operatorname{opr}(m)(\operatorname{nth}(1, p), \operatorname{nth}(2, p)) = x)\}) \\ \quad \quad =\; \operatorname{add}(a)(\operatorname{finsum}(\operatorname{ring\text{-}additive\text{-}ag}(a), \\ \quad \quad \quad \quad (\lambda p \in \{p \in (\operatorname{supp}(a, m, f) \times \operatorname{supp}(a, m, h)) : (\operatorname{opr}(m)(\operatorname{nth}(1, p), \operatorname{nth}(2, p)) = x)\}.\; \\ \quad \quad \quad \quad \quad \operatorname{mul}(a)(f(\operatorname{nth}(1, p)), h(\operatorname{nth}(2, p)))), \\ \quad \quad \quad \quad \{p \in (\operatorname{supp}(a, m, f) \times \operatorname{supp}(a, m, h)) : (\operatorname{opr}(m)(\operatorname{nth}(1, p), \operatorname{nth}(2, p)) = x)\}), \\ \quad \quad \quad \operatorname{finsum}(\operatorname{ring\text{-}additive\text{-}ag}(a), \\ \quad \quad \quad \quad (\lambda p \in \{p \in (\operatorname{supp}(a, m, g) \times \operatorname{supp}(a, m, h)) : (\operatorname{opr}(m)(\operatorname{nth}(1, p), \operatorname{nth}(2, p)) = x)\}.\; \\ \quad \quad \quad \quad \quad \operatorname{mul}(a)(g(\operatorname{nth}(1, p)), h(\operatorname{nth}(2, p)))), \\ \quad \quad \quad \quad \{p \in (\operatorname{supp}(a, m, g) \times \operatorname{supp}(a, m, h)) : (\operatorname{opr}(m)(\operatorname{nth}(1, p), \operatorname{nth}(2, p)) = x)\})) \end{array}}}
\vnbprop{monalg-one-supp-subset}{\vnbstmt{\begin{array}{@{}l@{}} \forall a, m, z. \\ \quad (\operatorname{is\text{-}ring}(a) \wedge (\operatorname{is\text{-}monoid}(m) \wedge (z \in \operatorname{supp}(a, m, \operatorname{monalg\text{-}one}(a, m))))) \Rightarrow  \\ \quad \quad z \in \{\operatorname{iden}(m)\} \end{array}}}
\vnbprop{monalg-one-left-at}{\vnbstmt{\begin{array}{@{}l@{}} \forall a, m, f, x. \\ \quad (\operatorname{is\text{-}ring}(a) \wedge (\operatorname{is\text{-}monoid}(m) \wedge ((f \in \operatorname{finsupp}(a, m)) \wedge (x \in \operatorname{carr}(m))))) \Rightarrow  \\ \quad \quad \operatorname{finsum}(\operatorname{ring\text{-}additive\text{-}ag}(a), \\ \quad \quad \quad (\lambda p \in \{p \in (\operatorname{supp}(a, m, \operatorname{monalg\text{-}one}(a, m)) \times \operatorname{supp}(a, m, f)) : (\operatorname{opr}(m)(\operatorname{nth}(1, p), \operatorname{nth}(2, p)) = x)\}.\; \\ \quad \quad \quad \quad \operatorname{mul}(a)(\operatorname{monalg\text{-}one}(a, m)(\operatorname{nth}(1, p)), f(\operatorname{nth}(2, p)))), \\ \quad \quad \quad \{p \in (\operatorname{supp}(a, m, \operatorname{monalg\text{-}one}(a, m)) \times \operatorname{supp}(a, m, f)) : (\operatorname{opr}(m)(\operatorname{nth}(1, p), \operatorname{nth}(2, p)) = x)\}) \\ \quad \quad =\; f(x) \end{array}}}
\vnbprop{monalg-one-right-at}{\vnbstmt{\begin{array}{@{}l@{}} \forall a, m, f, x. \\ \quad (\operatorname{is\text{-}ring}(a) \wedge (\operatorname{is\text{-}monoid}(m) \wedge ((f \in \operatorname{finsupp}(a, m)) \wedge (x \in \operatorname{carr}(m))))) \Rightarrow  \\ \quad \quad \operatorname{finsum}(\operatorname{ring\text{-}additive\text{-}ag}(a), \\ \quad \quad \quad (\lambda p \in \{p \in (\operatorname{supp}(a, m, f) \times \operatorname{supp}(a, m, \operatorname{monalg\text{-}one}(a, m))) : (\operatorname{opr}(m)(\operatorname{nth}(1, p), \operatorname{nth}(2, p)) = x)\}.\; \\ \quad \quad \quad \quad \operatorname{mul}(a)(f(\operatorname{nth}(1, p)), \operatorname{monalg\text{-}one}(a, m)(\operatorname{nth}(2, p)))), \\ \quad \quad \quad \{p \in (\operatorname{supp}(a, m, f) \times \operatorname{supp}(a, m, \operatorname{monalg\text{-}one}(a, m))) : (\operatorname{opr}(m)(\operatorname{nth}(1, p), \operatorname{nth}(2, p)) = x)\}) \\ \quad \quad =\; f(x) \end{array}}}
\vnbprop{monalg-distrib-left}{\vnbstmt{\begin{array}{@{}l@{}} \forall a, m, f, g, h. \\ \quad (\operatorname{is\text{-}ring}(a) \wedge (\operatorname{is\text{-}monoid}(m) \wedge ((f \in \operatorname{finsupp}(a, m)) \wedge ((g \in \operatorname{finsupp}(a, m)) \wedge (h \in \operatorname{finsupp}(a, m)))))) \Rightarrow  \\ \quad \quad \operatorname{monalg\text{-}mul}(a, m, f, \operatorname{monalg\text{-}add}(a, m, g, h)) \\ \quad \quad =\; \operatorname{monalg\text{-}add}(a, m, \operatorname{monalg\text{-}mul}(a, m, f, g), \operatorname{monalg\text{-}mul}(a, m, f, h)) \end{array}}}
\vnbprop{monalg-distrib-right}{\vnbstmt{\begin{array}{@{}l@{}} \forall a, m, f, g, h. \\ \quad (\operatorname{is\text{-}ring}(a) \wedge (\operatorname{is\text{-}monoid}(m) \wedge ((f \in \operatorname{finsupp}(a, m)) \wedge ((g \in \operatorname{finsupp}(a, m)) \wedge (h \in \operatorname{finsupp}(a, m)))))) \Rightarrow  \\ \quad \quad \operatorname{monalg\text{-}mul}(a, m, \operatorname{monalg\text{-}add}(a, m, f, g), h) \\ \quad \quad =\; \operatorname{monalg\text{-}add}(a, m, \operatorname{monalg\text{-}mul}(a, m, f, h), \operatorname{monalg\text{-}mul}(a, m, g, h)) \end{array}}}
\vnbprop{monalg-one-left-fn}{\vnbstmt{\begin{array}{@{}l@{}} \forall a, m, f. \\ \quad (\operatorname{is\text{-}ring}(a) \wedge (\operatorname{is\text{-}monoid}(m) \wedge (f \in \operatorname{finsupp}(a, m)))) \Rightarrow  \\ \quad \quad \operatorname{monalg\text{-}mul}(a, m, \operatorname{monalg\text{-}one}(a, m), f) = f \end{array}}}
\vnbprop{monalg-one-right-fn}{\vnbstmt{\begin{array}{@{}l@{}} \forall a, m, f. \\ \quad (\operatorname{is\text{-}ring}(a) \wedge (\operatorname{is\text{-}monoid}(m) \wedge (f \in \operatorname{finsupp}(a, m)))) \Rightarrow  \\ \quad \quad \operatorname{monalg\text{-}mul}(a, m, f, \operatorname{monalg\text{-}one}(a, m)) = f \end{array}}}
\vnbprop{monalg-one-left}{\vnbstmt{\begin{array}{@{}l@{}} \forall a, m, f, x. \\ \quad (\operatorname{is\text{-}ring}(a) \wedge (\operatorname{is\text{-}monoid}(m) \wedge ((f \in \operatorname{finsupp}(a, m)) \wedge (x \in \operatorname{carr}(m))))) \Rightarrow  \\ \quad \quad \operatorname{monalg\text{-}mul}(a, m, \operatorname{monalg\text{-}one}(a, m), f)(x) = f(x) \end{array}}}
\vnbprop{bijection-from-inverse}{\vnbstmt{\begin{array}{@{}l@{}} \forall x, y, i, a. \\ \quad i \in (x \to y) \wedge \\ \quad a \in (y \to x) \wedge \\ \quad \forall u \in x.\; (a(i(u)) = u) \wedge \\ \quad \forall v \in y.\; (i(a(v)) = v) \Rightarrow \\ \quad \quad i \in \mathrm{Bij}(x, y) \end{array}}}
\vnbprop{lambda-compose-value}{\vnbstmt{\forall f, h, m.; \forall t \in m.\; ((\lambda z \in m.\; f(h(z)))(t) \simeq f(h(t)))}}
\vnbprop{finsum-reindex-inverse}{\vnbstmt{\begin{array}{@{}l@{}} \forall g.; \forall s, t \in \mathbf{Set}, i \in (t \to s), a \in (s \to t), f \in (s \to \operatorname{carr}(g)). \\ \quad \operatorname{is\text{-}abelian\text{-}group}(g) \wedge \\ \quad |s| \in \mathbb{N} \wedge \\ \quad |t| \in \mathbb{N} \wedge \\ \quad \forall u \in t.\; (a(i(u)) = u) \wedge \\ \quad \forall v \in s.\; (i(a(v)) = v) \Rightarrow \\ \quad \quad \operatorname{finsum}(g, f, s) = \operatorname{finsum}(g, (\lambda z \in t.\; f(i(z))), t) \end{array}}}
\vnbprop{finsum-reindex-inverse-ptwise}{\vnbstmt{\begin{array}{@{}l@{}} \forall g.; \forall s, t \in \mathbf{Set}, i \in (t \to s), a \in (s \to t), f. \\ \quad \operatorname{is\text{-}abelian\text{-}group}(g) \wedge \\ \quad |s| \in \mathbb{N} \wedge \\ \quad |t| \in \mathbb{N} \wedge \\ \quad \forall u \in t.\; (a(i(u)) = u) \wedge \\ \quad \forall v \in s.\; (i(a(v)) = v) \wedge \\ \quad \forall b \in s.\; (f(b) \in \operatorname{carr}(g)) \Rightarrow \\ \quad \quad \operatorname{finsum}(g, f, s) = \operatorname{finsum}(g, (\lambda z \in t.\; f(i(z))), t) \end{array}}}
\vnbprop{finsum-fiber-value}{\vnbstmt{\begin{array}{@{}l@{}} \forall g, f, h, s, t.; \forall a \in t. \\ \quad ((\lambda b \in t.\; \operatorname{finsum}(g, f, \{c \in s : (h(c) = b)\}))(a) \simeq \operatorname{finsum}(g, f, \{c \in s : (h(c) = a)\})) \end{array}}}
\vnbprop{finsum-fiber-slice}{\vnbstmt{\begin{array}{@{}l@{}} \forall g.; \forall s, k \in \mathbf{Set}, h, p, f, l, e. \\ \quad \operatorname{is\text{-}abelian\text{-}group}(g) \wedge \\ \quad |s| \in \mathbb{N} \wedge \\ \quad |k| \in \mathbb{N} \wedge \\ \quad l \in (k \to s) \wedge \\ \quad \forall a \in k.\; (h(l(a)) = p) \wedge \\ \quad \forall w \in s.\; ((h(w) = p) \Rightarrow (e(w) \in k)) \wedge \\ \quad \forall a \in k.\; (e(l(a)) = a) \wedge \\ \quad \forall w \in s.\; ((h(w) = p) \Rightarrow (l(e(w)) = w)) \wedge \\ \quad \forall b \in s.\; (f(b) \in \operatorname{carr}(g)) \Rightarrow \\ \quad \quad \operatorname{finsum}(g, f, \{c \in s : (h(c) = p)\}) = \operatorname{finsum}(g, (\lambda z \in k.\; f(l(z))), k) \end{array}}}
\vnbprop{monalg-mul-supp-in-image}{\vnbstmt{\begin{array}{@{}l@{}} \forall a, m, f, g, z. \\ \quad \operatorname{is\text{-}ring}(a) \wedge \\ \quad \operatorname{is\text{-}monoid}(m) \wedge \\ \quad f \in \operatorname{finsupp}(a, m) \wedge \\ \quad g \in \operatorname{finsupp}(a, m) \wedge \\ \quad z \in \operatorname{supp}(a, m, \operatorname{monalg\text{-}mul}(a, m, f, g)) \Rightarrow \\ \quad \quad z \\ \quad \quad \in\; \operatorname{image}((\lambda q \in (\operatorname{supp}(a, m, f) \times \operatorname{supp}(a, m, g)).\; \operatorname{opr}(m)(\operatorname{nth}(1, q), \operatorname{nth}(2, q))), \\ \quad \quad \quad (\operatorname{supp}(a, m, f) \times \operatorname{supp}(a, m, g))) \end{array}}}
\vnbprop{monalg-mul-assoc-left-core}{\vnbstmt{\begin{array}{@{}l@{}} \forall a, m, f, g, h, x. \\ \quad \operatorname{is\text{-}ring}(a) \wedge \\ \quad \operatorname{is\text{-}monoid}(m) \wedge \\ \quad f \in \operatorname{finsupp}(a, m) \wedge \\ \quad g \in \operatorname{finsupp}(a, m) \wedge \\ \quad h \in \operatorname{finsupp}(a, m) \wedge \\ \quad x \in \operatorname{carr}(m) \Rightarrow \\ \quad \quad \operatorname{finsum}(\operatorname{ring\text{-}additive\text{-}ag}(a), \\ \quad \quad \quad (\lambda p \in \{p \in (\operatorname{image}((\lambda q \in (\operatorname{supp}(a, m, f) \times \operatorname{supp}(a, m, g)).\; \operatorname{opr}(m)(\operatorname{nth}(1, q), \operatorname{nth}(2, q))), (\operatorname{supp}(a, m, f) \times \operatorname{supp}(a, m, g))) \times \operatorname{supp}(a, m, h)) : (\operatorname{opr}(m)(\operatorname{nth}(1, p), \operatorname{nth}(2, p)) = x)\}.\; \\ \quad \quad \quad \quad \operatorname{mul}(a)(\operatorname{monalg\text{-}mul}(a, m, f, g)(\operatorname{nth}(1, p)), h(\operatorname{nth}(2, p)))), \\ \quad \quad \quad \{p \in (\operatorname{image}((\lambda q \in (\operatorname{supp}(a, m, f) \times \operatorname{supp}(a, m, g)).\; \operatorname{opr}(m)(\operatorname{nth}(1, q), \operatorname{nth}(2, q))), (\operatorname{supp}(a, m, f) \times \operatorname{supp}(a, m, g))) \times \operatorname{supp}(a, m, h)) : (\operatorname{opr}(m)(\operatorname{nth}(1, p), \operatorname{nth}(2, p)) = x)\}) \\ \quad \quad =\; \operatorname{finsum}(\operatorname{ring\text{-}additive\text{-}ag}(a), \\ \quad \quad \quad (\lambda w \in \{w \in ((\operatorname{supp}(a, m, f) \times \operatorname{supp}(a, m, g)) \times \operatorname{supp}(a, m, h)) : (\operatorname{opr}(m)(\operatorname{opr}(m)(\operatorname{nth}(1, \operatorname{nth}(1, w)), \operatorname{nth}(2, \operatorname{nth}(1, w))), \operatorname{nth}(2, w)) = x)\}.\; \\ \quad \quad \quad \quad \operatorname{mul}(a)(\operatorname{mul}(a)(f(\operatorname{nth}(1, \operatorname{nth}(1, w))), g(\operatorname{nth}(2, \operatorname{nth}(1, w)))), h(\operatorname{nth}(2, w)))), \\ \quad \quad \quad \{w \in ((\operatorname{supp}(a, m, f) \times \operatorname{supp}(a, m, g)) \times \operatorname{supp}(a, m, h)) : (\operatorname{opr}(m)(\operatorname{opr}(m)(\operatorname{nth}(1, \operatorname{nth}(1, w)), \operatorname{nth}(2, \operatorname{nth}(1, w))), \operatorname{nth}(2, w)) = x)\}) \end{array}}}
\vnbprop{monalg-mul-assoc-right-core}{\vnbstmt{\begin{array}{@{}l@{}} \forall a, m, f, g, h, x. \\ \quad \operatorname{is\text{-}ring}(a) \wedge \\ \quad \operatorname{is\text{-}monoid}(m) \wedge \\ \quad f \in \operatorname{finsupp}(a, m) \wedge \\ \quad g \in \operatorname{finsupp}(a, m) \wedge \\ \quad h \in \operatorname{finsupp}(a, m) \wedge \\ \quad x \in \operatorname{carr}(m) \Rightarrow \\ \quad \quad \operatorname{finsum}(\operatorname{ring\text{-}additive\text{-}ag}(a), \\ \quad \quad \quad (\lambda p \in \{p \in (\operatorname{supp}(a, m, f) \times \operatorname{image}((\lambda q \in (\operatorname{supp}(a, m, g) \times \operatorname{supp}(a, m, h)).\; \operatorname{opr}(m)(\operatorname{nth}(1, q), \operatorname{nth}(2, q))), (\operatorname{supp}(a, m, g) \times \operatorname{supp}(a, m, h)))) : (\operatorname{opr}(m)(\operatorname{nth}(1, p), \operatorname{nth}(2, p)) = x)\}.\; \\ \quad \quad \quad \quad \operatorname{mul}(a)(f(\operatorname{nth}(1, p)), \operatorname{monalg\text{-}mul}(a, m, g, h)(\operatorname{nth}(2, p)))), \\ \quad \quad \quad \{p \in (\operatorname{supp}(a, m, f) \times \operatorname{image}((\lambda q \in (\operatorname{supp}(a, m, g) \times \operatorname{supp}(a, m, h)).\; \operatorname{opr}(m)(\operatorname{nth}(1, q), \operatorname{nth}(2, q))), (\operatorname{supp}(a, m, g) \times \operatorname{supp}(a, m, h)))) : (\operatorname{opr}(m)(\operatorname{nth}(1, p), \operatorname{nth}(2, p)) = x)\}) \\ \quad \quad =\; \operatorname{finsum}(\operatorname{ring\text{-}additive\text{-}ag}(a), \\ \quad \quad \quad (\lambda w \in \{w \in ((\operatorname{supp}(a, m, f) \times \operatorname{supp}(a, m, g)) \times \operatorname{supp}(a, m, h)) : (\operatorname{opr}(m)(\operatorname{opr}(m)(\operatorname{nth}(1, \operatorname{nth}(1, w)), \operatorname{nth}(2, \operatorname{nth}(1, w))), \operatorname{nth}(2, w)) = x)\}.\; \\ \quad \quad \quad \quad \operatorname{mul}(a)(f(\operatorname{nth}(1, \operatorname{nth}(1, w))), \operatorname{mul}(a)(g(\operatorname{nth}(2, \operatorname{nth}(1, w))), h(\operatorname{nth}(2, w))))), \\ \quad \quad \quad \{w \in ((\operatorname{supp}(a, m, f) \times \operatorname{supp}(a, m, g)) \times \operatorname{supp}(a, m, h)) : (\operatorname{opr}(m)(\operatorname{opr}(m)(\operatorname{nth}(1, \operatorname{nth}(1, w)), \operatorname{nth}(2, \operatorname{nth}(1, w))), \operatorname{nth}(2, w)) = x)\}) \end{array}}}
\vnbprop{monalg-mul-assoc-at}{\vnbstmt{\begin{array}{@{}l@{}} \forall a, m, f, g, h, x. \\ \quad \operatorname{is\text{-}ring}(a) \wedge \\ \quad \operatorname{is\text{-}monoid}(m) \wedge \\ \quad f \in \operatorname{finsupp}(a, m) \wedge \\ \quad g \in \operatorname{finsupp}(a, m) \wedge \\ \quad h \in \operatorname{finsupp}(a, m) \wedge \\ \quad x \in \operatorname{carr}(m) \Rightarrow \\ \quad \quad \operatorname{monalg\text{-}mul}(a, m, \operatorname{monalg\text{-}mul}(a, m, f, g), h)(x) = \operatorname{monalg\text{-}mul}(a, m, f, \operatorname{monalg\text{-}mul}(a, m, g, h))(x) \end{array}}}
\vnbprop{monalg-mul-assoc}{\vnbstmt{\begin{array}{@{}l@{}} \forall a, m, f, g, h. \\ \quad (\operatorname{is\text{-}ring}(a) \wedge (\operatorname{is\text{-}monoid}(m) \wedge ((f \in \operatorname{finsupp}(a, m)) \wedge ((g \in \operatorname{finsupp}(a, m)) \wedge (h \in \operatorname{finsupp}(a, m)))))) \Rightarrow  \\ \quad \quad \operatorname{monalg\text{-}mul}(a, m, \operatorname{monalg\text{-}mul}(a, m, f, g), h) = \operatorname{monalg\text{-}mul}(a, m, f, \operatorname{monalg\text{-}mul}(a, m, g, h)) \end{array}}}
\vnbprop{monalg-is-ring}{\vnbstmt{\forall a, m.\; ((\operatorname{is\text{-}ring}(a) \wedge \operatorname{is\text{-}monoid}(m)) \Rightarrow \operatorname{is\text{-}ring}(\operatorname{monalg}(a, m)))}}

\vnbbookfile{theorem-library/la-z6-pullback}{theorem-library/la-z6-pullback}
\noindent{\small la-z6-pullback.scm -- a monoid automorphism induces a ring automorphism of the monoid algebra (LA-Z6, 2026-10-25).  The general mechanism under linear algebra 3.15.}\par
\vnbprop{monalg-pullback-apply}{\vnbstmt{\forall a, m, h, f.; \forall x \in \operatorname{carr}(m).\; (\operatorname{monalg\text{-}pullback}(a, m, h, f)(x) \simeq f(h(x)))}}
\vnbprop{monalg-pullback-fun}{\vnbstmt{\begin{array}{@{}l@{}} \forall a, m, h, f. \\ \quad (\operatorname{is\text{-}monoid}(m) \wedge ((h \in (\operatorname{carr}(m) \to \operatorname{carr}(m))) \wedge (f \in (\operatorname{carr}(m) \to \operatorname{carr}(a))))) \Rightarrow  \\ \quad \quad \operatorname{monalg\text{-}pullback}(a, m, h, f) \in (\operatorname{carr}(m) \to \operatorname{carr}(a)) \end{array}}}
\vnbprop{monalg-pullback-supp-fwd}{\vnbstmt{\begin{array}{@{}l@{}} \forall a, m, h, f, z. \\ \quad \operatorname{is\text{-}monoid}(m) \wedge \\ \quad h \in (\operatorname{carr}(m) \to \operatorname{carr}(m)) \wedge \\ \quad f \in (\operatorname{carr}(m) \to \operatorname{carr}(a)) \wedge \\ \quad z \in \operatorname{supp}(a, m, \operatorname{monalg\text{-}pullback}(a, m, h, f)) \Rightarrow \\ \quad \quad h(z) \in \operatorname{supp}(a, m, f) \end{array}}}
\vnbprop{monalg-pullback-supp-bwd}{\vnbstmt{\begin{array}{@{}l@{}} \forall a, m, h, s, f, z. \\ \quad \operatorname{is\text{-}monoid}(m) \wedge \\ \quad h \in (\operatorname{carr}(m) \to \operatorname{carr}(m)) \wedge \\ \quad s \in (\operatorname{carr}(m) \to \operatorname{carr}(m)) \wedge \\ \quad \forall u \in \operatorname{carr}(m).\; (h(s(u)) = u) \wedge \\ \quad f \in (\operatorname{carr}(m) \to \operatorname{carr}(a)) \wedge \\ \quad z \in \operatorname{supp}(a, m, f) \Rightarrow \\ \quad \quad s(z) \in \operatorname{supp}(a, m, \operatorname{monalg\text{-}pullback}(a, m, h, f)) \end{array}}}
\vnbprop{monalg-pullback-finsupp}{\vnbstmt{\begin{array}{@{}l@{}} \forall a, m, h, s, f. \\ \quad \operatorname{is\text{-}monoid}(m) \wedge \\ \quad h \in (\operatorname{carr}(m) \to \operatorname{carr}(m)) \wedge \\ \quad s \in (\operatorname{carr}(m) \to \operatorname{carr}(m)) \wedge \\ \quad \forall u \in \operatorname{carr}(m).\; (s(h(u)) = u) \wedge \\ \quad f \in \operatorname{finsupp}(a, m) \Rightarrow \\ \quad \quad \operatorname{monalg\text{-}pullback}(a, m, h, f) \in \operatorname{finsupp}(a, m) \end{array}}}
\vnbprop{monalg-pullback-add}{\vnbstmt{\begin{array}{@{}l@{}} \forall a, m, h, s, f, g. \\ \quad \operatorname{is\text{-}ring}(a) \wedge \\ \quad \operatorname{is\text{-}monoid}(m) \wedge \\ \quad h \in (\operatorname{carr}(m) \to \operatorname{carr}(m)) \wedge \\ \quad s \in (\operatorname{carr}(m) \to \operatorname{carr}(m)) \wedge \\ \quad \forall u \in \operatorname{carr}(m).\; (s(h(u)) = u) \wedge \\ \quad f \in \operatorname{finsupp}(a, m) \wedge \\ \quad g \in \operatorname{finsupp}(a, m) \Rightarrow \\ \quad \quad \operatorname{monalg\text{-}pullback}(a, m, h, \operatorname{monalg\text{-}add}(a, m, f, g)) \\ \quad \quad =\; \operatorname{monalg\text{-}add}(a, m, \operatorname{monalg\text{-}pullback}(a, m, h, f), \operatorname{monalg\text{-}pullback}(a, m, h, g)) \end{array}}}
\vnbprop{monalg-pullback-neg}{\vnbstmt{\begin{array}{@{}l@{}} \forall a, m, h, s, f. \\ \quad \operatorname{is\text{-}ring}(a) \wedge \\ \quad \operatorname{is\text{-}monoid}(m) \wedge \\ \quad h \in (\operatorname{carr}(m) \to \operatorname{carr}(m)) \wedge \\ \quad s \in (\operatorname{carr}(m) \to \operatorname{carr}(m)) \wedge \\ \quad \forall u \in \operatorname{carr}(m).\; (s(h(u)) = u) \wedge \\ \quad f \in \operatorname{finsupp}(a, m) \Rightarrow \\ \quad \quad \operatorname{monalg\text{-}pullback}(a, m, h, \operatorname{monalg\text{-}neg}(a, m, f)) = \operatorname{monalg\text{-}neg}(a, m, \operatorname{monalg\text{-}pullback}(a, m, h, f)) \end{array}}}
\vnbprop{monalg-pullback-zero}{\vnbstmt{\begin{array}{@{}l@{}} \forall a, m, h, s. \\ \quad \operatorname{is\text{-}ring}(a) \wedge \\ \quad \operatorname{is\text{-}monoid}(m) \wedge \\ \quad h \in (\operatorname{carr}(m) \to \operatorname{carr}(m)) \wedge \\ \quad s \in (\operatorname{carr}(m) \to \operatorname{carr}(m)) \wedge \\ \quad \forall u \in \operatorname{carr}(m).\; (s(h(u)) = u) \Rightarrow \\ \quad \quad \operatorname{monalg\text{-}pullback}(a, m, h, \operatorname{monalg\text{-}zero}(a, m)) = \operatorname{monalg\text{-}zero}(a, m) \end{array}}}
\vnbprop{monalg-pullback-one}{\vnbstmt{\begin{array}{@{}l@{}} \forall a, m, h, s. \\ \quad \operatorname{is\text{-}ring}(a) \wedge \\ \quad \operatorname{is\text{-}monoid}(m) \wedge \\ \quad h \in (\operatorname{carr}(m) \to \operatorname{carr}(m)) \wedge \\ \quad s \in (\operatorname{carr}(m) \to \operatorname{carr}(m)) \wedge \\ \quad \forall u \in \operatorname{carr}(m).\; (s(h(u)) = u) \wedge \\ \quad h(\operatorname{iden}(m)) = \operatorname{iden}(m) \Rightarrow \\ \quad \quad \operatorname{monalg\text{-}pullback}(a, m, h, \operatorname{monalg\text{-}one}(a, m)) = \operatorname{monalg\text{-}one}(a, m) \end{array}}}
\vnbprop{monalg-pullback-index-fun}{\vnbstmt{\begin{array}{@{}l@{}} \forall a, m, h, s, f, g, x. \\ \quad \operatorname{is\text{-}ring}(a) \wedge \\ \quad \operatorname{is\text{-}monoid}(m) \wedge \\ \quad h \in (\operatorname{carr}(m) \to \operatorname{carr}(m)) \wedge \\ \quad s \in (\operatorname{carr}(m) \to \operatorname{carr}(m)) \wedge \\ \quad \forall u \in \operatorname{carr}(m).\; (s(h(u)) = u) \wedge \\ \quad \forall u \in \operatorname{carr}(m).\; (h(s(u)) = u) \wedge \\ \quad \forall u, v \in \operatorname{carr}(m).\; (h(\operatorname{opr}(m)(u, v)) = \operatorname{opr}(m)(h(u), h(v))) \wedge \\ \quad f \in \operatorname{finsupp}(a, m) \wedge \\ \quad g \in \operatorname{finsupp}(a, m) \wedge \\ \quad x \in \operatorname{carr}(m) \Rightarrow \\ \quad \quad (\lambda q \in \{p \in (\operatorname{supp}(a, m, \operatorname{monalg\text{-}pullback}(a, m, h, f)) \times \operatorname{supp}(a, m, \operatorname{monalg\text{-}pullback}(a, m, h, g))) : (\operatorname{opr}(m)(\operatorname{nth}(1, p), \operatorname{nth}(2, p)) = x)\}.\; \\ \quad \quad \quad {[}h(\operatorname{nth}(1, q)), h(\operatorname{nth}(2, q))]) \\ \quad \quad \in\; (\{p \in (\operatorname{supp}(a, m, \operatorname{monalg\text{-}pullback}(a, m, h, f)) \times \operatorname{supp}(a, m, \operatorname{monalg\text{-}pullback}(a, m, h, g))) : (\operatorname{opr}(m)(\operatorname{nth}(1, p), \operatorname{nth}(2, p)) = x)\} \to \{p \in (\operatorname{supp}(a, m, f) \times \operatorname{supp}(a, m, g)) : (\operatorname{opr}(m)(\operatorname{nth}(1, p), \operatorname{nth}(2, p)) = h(x))\}) \end{array}}}
\vnbprop{monalg-pullback-index-inv-fun}{\vnbstmt{\begin{array}{@{}l@{}} \forall a, m, h, s, f, g, x. \\ \quad \operatorname{is\text{-}ring}(a) \wedge \\ \quad \operatorname{is\text{-}monoid}(m) \wedge \\ \quad h \in (\operatorname{carr}(m) \to \operatorname{carr}(m)) \wedge \\ \quad s \in (\operatorname{carr}(m) \to \operatorname{carr}(m)) \wedge \\ \quad \forall u \in \operatorname{carr}(m).\; (s(h(u)) = u) \wedge \\ \quad \forall u \in \operatorname{carr}(m).\; (h(s(u)) = u) \wedge \\ \quad \forall u, v \in \operatorname{carr}(m).\; (h(\operatorname{opr}(m)(u, v)) = \operatorname{opr}(m)(h(u), h(v))) \wedge \\ \quad f \in \operatorname{finsupp}(a, m) \wedge \\ \quad g \in \operatorname{finsupp}(a, m) \wedge \\ \quad x \in \operatorname{carr}(m) \Rightarrow \\ \quad \quad (\lambda q \in \{p \in (\operatorname{supp}(a, m, f) \times \operatorname{supp}(a, m, g)) : (\operatorname{opr}(m)(\operatorname{nth}(1, p), \operatorname{nth}(2, p)) = h(x))\}.\; \\ \quad \quad \quad {[}s(\operatorname{nth}(1, q)), s(\operatorname{nth}(2, q))]) \\ \quad \quad \in\; (\{p \in (\operatorname{supp}(a, m, f) \times \operatorname{supp}(a, m, g)) : (\operatorname{opr}(m)(\operatorname{nth}(1, p), \operatorname{nth}(2, p)) = h(x))\} \to \{p \in (\operatorname{supp}(a, m, \operatorname{monalg\text{-}pullback}(a, m, h, f)) \times \operatorname{supp}(a, m, \operatorname{monalg\text{-}pullback}(a, m, h, g))) : (\operatorname{opr}(m)(\operatorname{nth}(1, p), \operatorname{nth}(2, p)) = x)\}) \end{array}}}
\vnbprop{monalg-pullback-index-left}{\vnbstmt{\begin{array}{@{}l@{}} \forall a, m, h, s, f, g, x.; \forall q \in \{p \in (\operatorname{supp}(a, m, \operatorname{monalg\text{-}pullback}(a, m, h, f)) \times \operatorname{supp}(a, m, \operatorname{monalg\text{-}pullback}(a, m, h, g))) : (\operatorname{opr}(m)(\operatorname{nth}(1, p), \operatorname{nth}(2, p)) = x)\}. \\ \quad \operatorname{is\text{-}ring}(a) \wedge \\ \quad \operatorname{is\text{-}monoid}(m) \wedge \\ \quad h \in (\operatorname{carr}(m) \to \operatorname{carr}(m)) \wedge \\ \quad s \in (\operatorname{carr}(m) \to \operatorname{carr}(m)) \wedge \\ \quad \forall u \in \operatorname{carr}(m).\; (s(h(u)) = u) \wedge \\ \quad \forall u \in \operatorname{carr}(m).\; (h(s(u)) = u) \wedge \\ \quad \forall u, v \in \operatorname{carr}(m).\; (h(\operatorname{opr}(m)(u, v)) = \operatorname{opr}(m)(h(u), h(v))) \wedge \\ \quad f \in \operatorname{finsupp}(a, m) \wedge \\ \quad g \in \operatorname{finsupp}(a, m) \wedge \\ \quad x \in \operatorname{carr}(m) \Rightarrow \\ \quad \quad (\lambda b \in \{p \in (\operatorname{supp}(a, m, f) \times \operatorname{supp}(a, m, g)) : (\operatorname{opr}(m)(\operatorname{nth}(1, p), \operatorname{nth}(2, p)) = h(x))\}.\; {[}s(\operatorname{nth}(1, b)), s(\operatorname{nth}(2, b))])((\lambda b \in \{p \in (\operatorname{supp}(a, m, \operatorname{monalg\text{-}pullback}(a, m, h, f)) \times \operatorname{supp}(a, m, \operatorname{monalg\text{-}pullback}(a, m, h, g))) : (\operatorname{opr}(m)(\operatorname{nth}(1, p), \operatorname{nth}(2, p)) = x)\}.\; {[}h(\operatorname{nth}(1, b)), h(\operatorname{nth}(2, b))])(q)) \\ \quad \quad =\; q \end{array}}}
\vnbprop{monalg-pullback-index-right}{\vnbstmt{\begin{array}{@{}l@{}} \forall a, m, h, s, f, g, x.; \forall q \in \{p \in (\operatorname{supp}(a, m, f) \times \operatorname{supp}(a, m, g)) : (\operatorname{opr}(m)(\operatorname{nth}(1, p), \operatorname{nth}(2, p)) = h(x))\}. \\ \quad \operatorname{is\text{-}ring}(a) \wedge \\ \quad \operatorname{is\text{-}monoid}(m) \wedge \\ \quad h \in (\operatorname{carr}(m) \to \operatorname{carr}(m)) \wedge \\ \quad s \in (\operatorname{carr}(m) \to \operatorname{carr}(m)) \wedge \\ \quad \forall u \in \operatorname{carr}(m).\; (s(h(u)) = u) \wedge \\ \quad \forall u \in \operatorname{carr}(m).\; (h(s(u)) = u) \wedge \\ \quad \forall u, v \in \operatorname{carr}(m).\; (h(\operatorname{opr}(m)(u, v)) = \operatorname{opr}(m)(h(u), h(v))) \wedge \\ \quad f \in \operatorname{finsupp}(a, m) \wedge \\ \quad g \in \operatorname{finsupp}(a, m) \wedge \\ \quad x \in \operatorname{carr}(m) \Rightarrow \\ \quad \quad (\lambda b \in \{p \in (\operatorname{supp}(a, m, \operatorname{monalg\text{-}pullback}(a, m, h, f)) \times \operatorname{supp}(a, m, \operatorname{monalg\text{-}pullback}(a, m, h, g))) : (\operatorname{opr}(m)(\operatorname{nth}(1, p), \operatorname{nth}(2, p)) = x)\}.\; {[}h(\operatorname{nth}(1, b)), h(\operatorname{nth}(2, b))])((\lambda b \in \{p \in (\operatorname{supp}(a, m, f) \times \operatorname{supp}(a, m, g)) : (\operatorname{opr}(m)(\operatorname{nth}(1, p), \operatorname{nth}(2, p)) = h(x))\}.\; {[}s(\operatorname{nth}(1, b)), s(\operatorname{nth}(2, b))])(q)) \\ \quad \quad =\; q \end{array}}}
\vnbprop{monalg-pullback-index-bijection}{\vnbstmt{\begin{array}{@{}l@{}} \forall a, m, h, s, f, g, x. \\ \quad \operatorname{is\text{-}ring}(a) \wedge \\ \quad \operatorname{is\text{-}monoid}(m) \wedge \\ \quad h \in (\operatorname{carr}(m) \to \operatorname{carr}(m)) \wedge \\ \quad s \in (\operatorname{carr}(m) \to \operatorname{carr}(m)) \wedge \\ \quad \forall u \in \operatorname{carr}(m).\; (s(h(u)) = u) \wedge \\ \quad \forall u \in \operatorname{carr}(m).\; (h(s(u)) = u) \wedge \\ \quad \forall u, v \in \operatorname{carr}(m).\; (h(\operatorname{opr}(m)(u, v)) = \operatorname{opr}(m)(h(u), h(v))) \wedge \\ \quad f \in \operatorname{finsupp}(a, m) \wedge \\ \quad g \in \operatorname{finsupp}(a, m) \wedge \\ \quad x \in \operatorname{carr}(m) \Rightarrow \\ \quad \quad (\lambda q \in \{p \in (\operatorname{supp}(a, m, \operatorname{monalg\text{-}pullback}(a, m, h, f)) \times \operatorname{supp}(a, m, \operatorname{monalg\text{-}pullback}(a, m, h, g))) : (\operatorname{opr}(m)(\operatorname{nth}(1, p), \operatorname{nth}(2, p)) = x)\}.\; \\ \quad \quad \quad {[}h(\operatorname{nth}(1, q)), h(\operatorname{nth}(2, q))]) \\ \quad \quad \in\; \mathrm{Bij}(\{p \in (\operatorname{supp}(a, m, \operatorname{monalg\text{-}pullback}(a, m, h, f)) \times \operatorname{supp}(a, m, \operatorname{monalg\text{-}pullback}(a, m, h, g))) : (\operatorname{opr}(m)(\operatorname{nth}(1, p), \operatorname{nth}(2, p)) = x)\}, \{p \in (\operatorname{supp}(a, m, f) \times \operatorname{supp}(a, m, g)) : (\operatorname{opr}(m)(\operatorname{nth}(1, p), \operatorname{nth}(2, p)) = h(x))\}) \end{array}}}
\vnbprop{monalg-pullback-mul-at}{\vnbstmt{\begin{array}{@{}l@{}} \forall a, m, h, s, f, g, x. \\ \quad \operatorname{is\text{-}ring}(a) \wedge \\ \quad \operatorname{is\text{-}monoid}(m) \wedge \\ \quad h \in (\operatorname{carr}(m) \to \operatorname{carr}(m)) \wedge \\ \quad s \in (\operatorname{carr}(m) \to \operatorname{carr}(m)) \wedge \\ \quad \forall u \in \operatorname{carr}(m).\; (s(h(u)) = u) \wedge \\ \quad \forall u \in \operatorname{carr}(m).\; (h(s(u)) = u) \wedge \\ \quad \forall u, v \in \operatorname{carr}(m).\; (h(\operatorname{opr}(m)(u, v)) = \operatorname{opr}(m)(h(u), h(v))) \wedge \\ \quad f \in \operatorname{finsupp}(a, m) \wedge \\ \quad g \in \operatorname{finsupp}(a, m) \wedge \\ \quad x \in \operatorname{carr}(m) \Rightarrow \\ \quad \quad \operatorname{monalg\text{-}pullback}(a, m, h, \operatorname{monalg\text{-}mul}(a, m, f, g))(x) \\ \quad \quad =\; \operatorname{monalg\text{-}mul}(a, m, \operatorname{monalg\text{-}pullback}(a, m, h, f), \operatorname{monalg\text{-}pullback}(a, m, h, g))(x) \end{array}}}
\vnbprop{monalg-pullback-mul}{\vnbstmt{\begin{array}{@{}l@{}} \forall a, m, h, s, f, g. \\ \quad \operatorname{is\text{-}ring}(a) \wedge \\ \quad \operatorname{is\text{-}monoid}(m) \wedge \\ \quad h \in (\operatorname{carr}(m) \to \operatorname{carr}(m)) \wedge \\ \quad s \in (\operatorname{carr}(m) \to \operatorname{carr}(m)) \wedge \\ \quad \forall u \in \operatorname{carr}(m).\; (s(h(u)) = u) \wedge \\ \quad \forall u \in \operatorname{carr}(m).\; (h(s(u)) = u) \wedge \\ \quad \forall u, v \in \operatorname{carr}(m).\; (h(\operatorname{opr}(m)(u, v)) = \operatorname{opr}(m)(h(u), h(v))) \wedge \\ \quad f \in \operatorname{finsupp}(a, m) \wedge \\ \quad g \in \operatorname{finsupp}(a, m) \Rightarrow \\ \quad \quad \operatorname{monalg\text{-}pullback}(a, m, h, \operatorname{monalg\text{-}mul}(a, m, f, g)) \\ \quad \quad =\; \operatorname{monalg\text{-}mul}(a, m, \operatorname{monalg\text{-}pullback}(a, m, h, f), \operatorname{monalg\text{-}pullback}(a, m, h, g)) \end{array}}}
\vnbprop{monalg-pullback-inverse}{\vnbstmt{\begin{array}{@{}l@{}} \forall a, m, h, s, f. \\ \quad \operatorname{is\text{-}ring}(a) \wedge \\ \quad \operatorname{is\text{-}monoid}(m) \wedge \\ \quad h \in (\operatorname{carr}(m) \to \operatorname{carr}(m)) \wedge \\ \quad s \in (\operatorname{carr}(m) \to \operatorname{carr}(m)) \wedge \\ \quad \forall u \in \operatorname{carr}(m).\; (s(h(u)) = u) \wedge \\ \quad \forall u \in \operatorname{carr}(m).\; (h(s(u)) = u) \wedge \\ \quad f \in \operatorname{finsupp}(a, m) \Rightarrow \\ \quad \quad \operatorname{monalg\text{-}pullback}(a, m, s, \operatorname{monalg\text{-}pullback}(a, m, h, f)) = f \end{array}}}
\vnbprop{monalg-pullback-map-fun}{\vnbstmt{\begin{array}{@{}l@{}} \forall a, m, h, s. \\ \quad \operatorname{is\text{-}ring}(a) \wedge \\ \quad \operatorname{is\text{-}monoid}(m) \wedge \\ \quad h \in (\operatorname{carr}(m) \to \operatorname{carr}(m)) \wedge \\ \quad s \in (\operatorname{carr}(m) \to \operatorname{carr}(m)) \wedge \\ \quad \forall u \in \operatorname{carr}(m).\; (s(h(u)) = u) \Rightarrow \\ \quad \quad (\lambda f \in \operatorname{finsupp}(a, m).\; \operatorname{monalg\text{-}pullback}(a, m, h, f)) \in (\operatorname{finsupp}(a, m) \to \operatorname{finsupp}(a, m)) \end{array}}}
\vnbprop{monalg-pullback-map-bijection}{\vnbstmt{\begin{array}{@{}l@{}} \forall a, m, h, s. \\ \quad \operatorname{is\text{-}ring}(a) \wedge \\ \quad \operatorname{is\text{-}monoid}(m) \wedge \\ \quad h \in (\operatorname{carr}(m) \to \operatorname{carr}(m)) \wedge \\ \quad s \in (\operatorname{carr}(m) \to \operatorname{carr}(m)) \wedge \\ \quad \forall u \in \operatorname{carr}(m).\; (s(h(u)) = u) \wedge \\ \quad \forall u \in \operatorname{carr}(m).\; (h(s(u)) = u) \Rightarrow \\ \quad \quad (\lambda f \in \operatorname{finsupp}(a, m).\; \operatorname{monalg\text{-}pullback}(a, m, h, f)) \in \mathrm{Bij}(\operatorname{finsupp}(a, m), \operatorname{finsupp}(a, m)) \end{array}}}
\vnbprop{monalg-pullback-iso}{\vnbstmt{\begin{array}{@{}l@{}} \forall a, m, h, s. \\ \quad \operatorname{is\text{-}ring}(a) \wedge \\ \quad \operatorname{is\text{-}monoid}(m) \wedge \\ \quad h \in (\operatorname{carr}(m) \to \operatorname{carr}(m)) \wedge \\ \quad s \in (\operatorname{carr}(m) \to \operatorname{carr}(m)) \wedge \\ \quad \forall u \in \operatorname{carr}(m).\; (s(h(u)) = u) \wedge \\ \quad \forall u \in \operatorname{carr}(m).\; (h(s(u)) = u) \wedge \\ \quad \forall u, v \in \operatorname{carr}(m).\; (h(\operatorname{opr}(m)(u, v)) = \operatorname{opr}(m)(h(u), h(v))) \wedge \\ \quad h(\operatorname{iden}(m)) = \operatorname{iden}(m) \Rightarrow \\ \quad \quad \operatorname{is\text{-}iso\text{-}ring}(\operatorname{monalg}(a, m), \operatorname{monalg}(a, m), (\lambda f \in \operatorname{finsupp}(a, m).\; \operatorname{monalg\text{-}pullback}(a, m, h, f))) \end{array}}}

\vnbbookfile{theorem-library/rake-prod-of-sums}{theorem-library/rake-prod-of-sums}
\noindent{\small rake-prod-of-sums.scm -- the product-of-sums expansion, proven.}\par
\vnbprop{difference-empty-left}{\vnbstmt{\forall s.\; (\operatorname{difference}(\emptyset, s) = \emptyset)}}
\vnbprop{insert-remove}{\vnbstmt{\begin{array}{@{}l@{}} \forall s.; \forall k \in \mathbf{Set}. \\ \quad \neg (k \in s) \Rightarrow (\operatorname{difference}((s \cup \{k\}), \{k\}) = s) \end{array}}}
\vnbprop{prod-ring-type-ptwise}{\vnbstmt{\begin{array}{@{}l@{}} \forall r.; \forall x \in \mathbf{Set}, f. \\ \quad (\operatorname{is\text{-}commutative\text{-}ring}(r) \wedge ((|x| \in \mathbb{N}) \wedge \forall z \in x.\; (f(z) \in \operatorname{carr}(r)))) \Rightarrow  \\ \quad \quad \operatorname{prod\text{-}ring}(r, f, x) \in \operatorname{carr}(r) \end{array}}}
\vnbprop{prod-ring-insert-ptwise}{\vnbstmt{\begin{array}{@{}l@{}} \forall r.; \forall x, k \in \mathbf{Set}, f. \\ \quad \operatorname{is\text{-}commutative\text{-}ring}(r) \wedge \\ \quad |x| \in \mathbb{N} \wedge \\ \quad \neg (k \in x) \wedge \\ \quad \forall z \in (x \cup \{k\}).\; (f(z) \in \operatorname{carr}(r)) \Rightarrow \\ \quad \quad \operatorname{prod\text{-}ring}(r, f, (x \cup \{k\})) = \operatorname{mul}(r)(\operatorname{prod\text{-}ring}(r, f, x), f(k)) \end{array}}}
\vnbprop{prod-ring-congruence-q}{\vnbstmt{\begin{array}{@{}l@{}} \forall r.; \forall x \in \mathbf{Set}, f, g. \\ \quad (\operatorname{is\text{-}commutative\text{-}ring}(r) \wedge ((|x| \in \mathbb{N}) \wedge \forall z \in x.\; (f(z) = g(z)))) \Rightarrow  \\ \quad \quad (\operatorname{prod\text{-}ring}(r, f, x) \simeq \operatorname{prod\text{-}ring}(r, g, x)) \end{array}}}
\vnbprop{cra-carr}{\vnbstmt{\begin{array}{@{}l@{}} \forall r. \\ \quad \operatorname{is\text{-}commutative\text{-}ring}(r) \Rightarrow (\operatorname{carr}(\operatorname{commutative\text{-}ring\text{-}additive\text{-}ag}(r)) = \operatorname{carr}(r)) \end{array}}}
\vnbprop{cra-op}{\vnbstmt{\begin{array}{@{}l@{}} \forall r. \\ \quad \operatorname{is\text{-}commutative\text{-}ring}(r) \Rightarrow (\operatorname{opr}(\operatorname{commutative\text{-}ring\text{-}additive\text{-}ag}(r)) = \operatorname{add}(r)) \end{array}}}
\vnbprop{difference-in-power}{\vnbstmt{\forall w \in \mathbf{Set}, s.\; (\operatorname{difference}(w, s) \in \operatorname{power}(w))}}
\vnbprop{prod-ring-type-subset}{\vnbstmt{\begin{array}{@{}l@{}} \forall r.; \forall w \in \mathbf{Set}, f.; \forall s \in \operatorname{power}(w). \\ \quad (\operatorname{is\text{-}commutative\text{-}ring}(r) \wedge ((|w| \in \mathbb{N}) \wedge \forall z \in w.\; (f(z) \in \operatorname{carr}(r)))) \Rightarrow  \\ \quad \quad \operatorname{prod\text{-}ring}(r, f, s) \in \operatorname{carr}(r) \end{array}}}
\vnbprop{prod-of-sums-summand-type}{\vnbstmt{\begin{array}{@{}l@{}} \forall r.; \forall w \in \mathbf{Set}, a, b.; \forall s \in \operatorname{power}(w). \\ \quad \operatorname{is\text{-}commutative\text{-}ring}(r) \wedge \\ \quad |w| \in \mathbb{N} \wedge \\ \quad \forall z \in w.\; (a(z) \in \operatorname{carr}(r)) \wedge \\ \quad \forall z \in w.\; (b(z) \in \operatorname{carr}(r)) \Rightarrow \\ \quad \quad \operatorname{mul}(r)(\operatorname{prod\text{-}ring}(r, a, s), \operatorname{prod\text{-}ring}(r, b, \operatorname{difference}(w, s))) \in \operatorname{carr}(r) \end{array}}}
\vnbprop{power-insert-injective}{\vnbstmt{\begin{array}{@{}l@{}} \forall x, k \in \mathbf{Set}. \\ \quad \neg (k \in x) \Rightarrow  \\ \quad \quad (\lambda s \in \operatorname{power}(x).\; (s \cup \{k\})) \in \operatorname{injection}(\operatorname{power}(x), \operatorname{power}((x \cup \{k\}))) \end{array}}}
\vnbprop{lambda-eta-value}{\vnbstmt{\forall f, m.; \forall t \in m.\; ((\lambda z \in m.\; f(z))(t) \simeq f(t))}}
\vnbprop{prod-of-sums-ptwise}{\vnbstmt{\begin{array}{@{}l@{}} \forall r.; \forall w \in \mathbf{Set}, a, b. \\ \quad \operatorname{is\text{-}commutative\text{-}ring}(r) \wedge \\ \quad |w| \in \mathbb{N} \wedge \\ \quad \forall z \in w.\; (a(z) \in \operatorname{carr}(r)) \wedge \\ \quad \forall z \in w.\; (b(z) \in \operatorname{carr}(r)) \Rightarrow \\ \quad \quad \operatorname{prod\text{-}ring}(r, (\lambda v \in w.\; \operatorname{add}(r)(a(v), b(v))), w) \\ \quad \quad =\; \operatorname{finsum}(\operatorname{commutative\text{-}ring\text{-}additive\text{-}ag}(r), \\ \quad \quad \quad (\lambda \mathit{sv} \in \operatorname{power}(w).\; \operatorname{mul}(r)(\operatorname{prod\text{-}ring}(r, a, \mathit{sv}), \operatorname{prod\text{-}ring}(r, b, \operatorname{difference}(w, \mathit{sv})))), \\ \quad \quad \quad \operatorname{power}(w)) \end{array}}}
\vnbprop{prod-of-sums-expansion}{\vnbstmt{\begin{array}{@{}l@{}} \forall r.; \forall x \in \mathbf{Set}, a, b \in (x \to \operatorname{carr}(r)). \\ \quad (\operatorname{is\text{-}commutative\text{-}ring}(r) \wedge (|x| \in \mathbb{N})) \Rightarrow  \\ \quad \quad \operatorname{prod\text{-}ring}(r, (\lambda k \in x.\; \operatorname{add}(r)(a(k), b(k))), x) \\ \quad \quad =\; \operatorname{finsum}(\operatorname{commutative\text{-}ring\text{-}additive\text{-}ag}(r), \\ \quad \quad \quad (\lambda s \in \operatorname{power}(x).\; \operatorname{mul}(r)(\operatorname{prod\text{-}ring}(r, a, s), \operatorname{prod\text{-}ring}(r, b, \operatorname{difference}(x, s)))), \\ \quad \quad \quad \operatorname{power}(x)) \end{array}}}

\vnbbookfile{theorem-library/rake-monalg-comm}{theorem-library/rake-monalg-comm}
\noindent{\small theorem-library/rake-monalg-comm.scm -- rake batch 5b, assignment 5b-A.}\par
\vnbprop{monalg-index-swap-fun}{\vnbstmt{\begin{array}{@{}l@{}} \forall a, m, f, g, x. \\ \quad \operatorname{is\text{-}comm\text{-}monoid}(m) \Rightarrow  \\ \quad \quad (\lambda q \in \{p \in (\operatorname{supp}(a, m, f) \times \operatorname{supp}(a, m, g)) : (\operatorname{opr}(m)(\operatorname{nth}(1, p), \operatorname{nth}(2, p)) = x)\}.\; \\ \quad \quad \quad {[}\operatorname{nth}(2, q), \operatorname{nth}(1, q)]) \\ \quad \quad \in\; (\{p \in (\operatorname{supp}(a, m, f) \times \operatorname{supp}(a, m, g)) : (\operatorname{opr}(m)(\operatorname{nth}(1, p), \operatorname{nth}(2, p)) = x)\} \to \{p \in (\operatorname{supp}(a, m, g) \times \operatorname{supp}(a, m, f)) : (\operatorname{opr}(m)(\operatorname{nth}(1, p), \operatorname{nth}(2, p)) = x)\}) \end{array}}}
\vnbprop{monalg-index-swap-inverse}{\vnbstmt{\begin{array}{@{}l@{}} \forall a, m, f, g, x.; \forall u \in \{p \in (\operatorname{supp}(a, m, f) \times \operatorname{supp}(a, m, g)) : (\operatorname{opr}(m)(\operatorname{nth}(1, p), \operatorname{nth}(2, p)) = x)\}. \\ \quad \operatorname{is\text{-}comm\text{-}monoid}(m) \Rightarrow  \\ \quad \quad (\lambda q \in \{p \in (\operatorname{supp}(a, m, g) \times \operatorname{supp}(a, m, f)) : (\operatorname{opr}(m)(\operatorname{nth}(1, p), \operatorname{nth}(2, p)) = x)\}.\; {[}\operatorname{nth}(2, q), \operatorname{nth}(1, q)])((\lambda q \in \{p \in (\operatorname{supp}(a, m, f) \times \operatorname{supp}(a, m, g)) : (\operatorname{opr}(m)(\operatorname{nth}(1, p), \operatorname{nth}(2, p)) = x)\}.\; {[}\operatorname{nth}(2, q), \operatorname{nth}(1, q)])(u)) \\ \quad \quad =\; u \end{array}}}
\vnbprop{monalg-index-swap-bijection}{\vnbstmt{\begin{array}{@{}l@{}} \forall a, m, f, g, x. \\ \quad \operatorname{is\text{-}comm\text{-}monoid}(m) \Rightarrow  \\ \quad \quad (\lambda q \in \{p \in (\operatorname{supp}(a, m, f) \times \operatorname{supp}(a, m, g)) : (\operatorname{opr}(m)(\operatorname{nth}(1, p), \operatorname{nth}(2, p)) = x)\}.\; \\ \quad \quad \quad {[}\operatorname{nth}(2, q), \operatorname{nth}(1, q)]) \\ \quad \quad \in\; \mathrm{Bij}(\{p \in (\operatorname{supp}(a, m, f) \times \operatorname{supp}(a, m, g)) : (\operatorname{opr}(m)(\operatorname{nth}(1, p), \operatorname{nth}(2, p)) = x)\}, \{p \in (\operatorname{supp}(a, m, g) \times \operatorname{supp}(a, m, f)) : (\operatorname{opr}(m)(\operatorname{nth}(1, p), \operatorname{nth}(2, p)) = x)\}) \end{array}}}
\vnbprop{monalg-mul-comm-at}{\vnbstmt{\begin{array}{@{}l@{}} \forall a, m, f, g, x. \\ \quad \operatorname{is\text{-}commutative\text{-}ring}(a) \wedge \\ \quad \operatorname{is\text{-}comm\text{-}monoid}(m) \wedge \\ \quad f \in \operatorname{finsupp}(a, m) \wedge \\ \quad g \in \operatorname{finsupp}(a, m) \wedge \\ \quad x \in \operatorname{carr}(m) \Rightarrow \\ \quad \quad \operatorname{monalg\text{-}mul}(a, m, f, g)(x) = \operatorname{monalg\text{-}mul}(a, m, g, f)(x) \end{array}}}
\vnbprop{monalg-mul-comm}{\vnbstmt{\begin{array}{@{}l@{}} \forall a, m, f, g. \\ \quad (\operatorname{is\text{-}commutative\text{-}ring}(a) \wedge (\operatorname{is\text{-}comm\text{-}monoid}(m) \wedge ((f \in \operatorname{finsupp}(a, m)) \wedge (g \in \operatorname{finsupp}(a, m))))) \Rightarrow  \\ \quad \quad \operatorname{monalg\text{-}mul}(a, m, f, g) = \operatorname{monalg\text{-}mul}(a, m, g, f) \end{array}}}
\vnbprop{monalg-comm}{\vnbstmt{\begin{array}{@{}l@{}} \forall a, m. \\ \quad (\operatorname{is\text{-}commutative\text{-}ring}(a) \wedge \operatorname{is\text{-}comm\text{-}monoid}(m)) \Rightarrow  \\ \quad \quad \operatorname{is\text{-}commutative\text{-}ring}(\operatorname{monalg}(a, m)) \end{array}}}

\vnbbookfile{theorem-library/la-z6-mpoly}{theorem-library/la-z6-mpoly}
\noindent{\small la-z6-mpoly.scm -- R[x\_1, ..., x\_n] and linear algebra Proposition 3.15 (LA-Z6, 2026-10-25).}\par
\vnbprop{multi-index-monoid-carr}{\vnbstmt{\forall s \in \mathbf{Set}.\; (\operatorname{carr}(\operatorname{multi\text{-}index\text{-}monoid}(s)) \simeq (s \to \mathbb{N}))}}
\vnbprop{multi-index-monoid-opr}{\vnbstmt{\begin{array}{@{}l@{}} \forall s \in \mathbf{Set}. \\ \quad (\operatorname{opr}(\operatorname{multi\text{-}index\text{-}monoid}(s)) \simeq (\lambda {[}\operatorname{miu\_}, \operatorname{miv\_}] \in ((s \to \mathbb{N}) \times (s \to \mathbb{N})).\; \operatorname{mi\text{-}add}(s, \operatorname{miu\_}, \operatorname{miv\_}))) \end{array}}}
\vnbprop{multi-index-monoid-iden}{\vnbstmt{\forall s \in \mathbf{Set}.\; (\operatorname{iden}(\operatorname{multi\text{-}index\text{-}monoid}(s)) \simeq \operatorname{mi\text{-}zero}(s))}}
\vnbprop{mi-add-apply}{\vnbstmt{\forall s, u, v.; \forall j \in s.\; (\operatorname{mi\text{-}add}(s, u, v)(j) \simeq (u(j) + v(j)))}}
\vnbprop{mi-zero-apply}{\vnbstmt{\forall s.; \forall j \in s.\; (\operatorname{mi\text{-}zero}(s)(j) \simeq 0)}}
\vnbprop{mi-reindex-apply}{\vnbstmt{\forall s, a, b.; \forall j \in s.\; (\operatorname{mi\text{-}reindex}(s, a, b)(j) \simeq b(a(j)))}}
\vnbprop{mi-add-fun}{\vnbstmt{\begin{array}{@{}l@{}} \forall s, u, v. \\ \quad ((s \in \mathbf{Set}) \wedge ((u \in (s \to \mathbb{N})) \wedge (v \in (s \to \mathbb{N})))) \Rightarrow  \\ \quad \quad \operatorname{mi\text{-}add}(s, u, v) \in (s \to \mathbb{N}) \end{array}}}
\vnbprop{mi-zero-fun}{\vnbstmt{\forall s \in \mathbf{Set}.\; (\operatorname{mi\text{-}zero}(s) \in (s \to \mathbb{N}))}}
\vnbprop{mi-reindex-fun}{\vnbstmt{\begin{array}{@{}l@{}} \forall s, a, b. \\ \quad ((s \in \mathbf{Set}) \wedge ((a \in (s \to s)) \wedge (b \in (s \to \mathbb{N})))) \Rightarrow  \\ \quad \quad \operatorname{mi\text{-}reindex}(s, a, b) \in (s \to \mathbb{N}) \end{array}}}
\vnbprop{mi-add-assoc}{\vnbstmt{\begin{array}{@{}l@{}} \forall s, u, v, w. \\ \quad ((s \in \mathbf{Set}) \wedge ((u \in (s \to \mathbb{N})) \wedge ((v \in (s \to \mathbb{N})) \wedge (w \in (s \to \mathbb{N}))))) \Rightarrow  \\ \quad \quad \operatorname{mi\text{-}add}(s, \operatorname{mi\text{-}add}(s, u, v), w) = \operatorname{mi\text{-}add}(s, u, \operatorname{mi\text{-}add}(s, v, w)) \end{array}}}
\vnbprop{mi-add-comm}{\vnbstmt{\begin{array}{@{}l@{}} \forall s, u, v. \\ \quad ((s \in \mathbf{Set}) \wedge ((u \in (s \to \mathbb{N})) \wedge (v \in (s \to \mathbb{N})))) \Rightarrow  \\ \quad \quad \operatorname{mi\text{-}add}(s, u, v) = \operatorname{mi\text{-}add}(s, v, u) \end{array}}}
\vnbprop{mi-add-zero-right}{\vnbstmt{\begin{array}{@{}l@{}} \forall s, u. \\ \quad ((s \in \mathbf{Set}) \wedge (u \in (s \to \mathbb{N}))) \Rightarrow (\operatorname{mi\text{-}add}(s, u, \operatorname{mi\text{-}zero}(s)) = u) \end{array}}}
\vnbprop{mi-add-zero-left}{\vnbstmt{\begin{array}{@{}l@{}} \forall s, u. \\ \quad ((s \in \mathbf{Set}) \wedge (u \in (s \to \mathbb{N}))) \Rightarrow (\operatorname{mi\text{-}add}(s, \operatorname{mi\text{-}zero}(s), u) = u) \end{array}}}
\vnbprop{multi-index-monoid-is-comm-monoid}{\vnbstmt{\forall s \in \mathbf{Set}.\; \operatorname{is\text{-}comm\text{-}monoid}(\operatorname{multi\text{-}index\text{-}monoid}(s))}}
\vnbprop{mi-reindex-add}{\vnbstmt{\begin{array}{@{}l@{}} \forall s, a, u, v. \\ \quad ((s \in \mathbf{Set}) \wedge ((a \in (s \to s)) \wedge ((u \in (s \to \mathbb{N})) \wedge (v \in (s \to \mathbb{N}))))) \Rightarrow  \\ \quad \quad \operatorname{mi\text{-}reindex}(s, a, \operatorname{mi\text{-}add}(s, u, v)) = \operatorname{mi\text{-}add}(s, \operatorname{mi\text{-}reindex}(s, a, u), \operatorname{mi\text{-}reindex}(s, a, v)) \end{array}}}
\vnbprop{mi-reindex-zero}{\vnbstmt{\begin{array}{@{}l@{}} \forall s, a. \\ \quad ((s \in \mathbf{Set}) \wedge (a \in (s \to s))) \Rightarrow (\operatorname{mi\text{-}reindex}(s, a, \operatorname{mi\text{-}zero}(s)) = \operatorname{mi\text{-}zero}(s)) \end{array}}}
\vnbprop{mi-reindex-inverse}{\vnbstmt{\begin{array}{@{}l@{}} \forall s, a, t, b. \\ \quad ((s \in \mathbf{Set}) \wedge ((a \in (s \to s)) \wedge ((t \in (s \to s)) \wedge (\forall q \in s.\; (a(t(q)) = q) \wedge (b \in (s \to \mathbb{N})))))) \Rightarrow  \\ \quad \quad \operatorname{mi\text{-}reindex}(s, t, \operatorname{mi\text{-}reindex}(s, a, b)) = b \end{array}}}
\vnbprop{multi-index-reindex-is-ring-iso}{\vnbstmt{\begin{array}{@{}l@{}} \forall r, s, a. \\ \quad (\operatorname{is\text{-}ring}(r) \wedge ((s \in \mathbf{Set}) \wedge (a \in \mathrm{Bij}(s, s)))) \Rightarrow  \\ \quad \quad \operatorname{is\text{-}iso\text{-}ring}(\operatorname{monalg}(r, \operatorname{multi\text{-}index\text{-}monoid}(s)), \operatorname{monalg}(r, \operatorname{multi\text{-}index\text{-}monoid}(s)), \\ \quad \quad \quad (\lambda f \in \operatorname{finsupp}(r, \operatorname{multi\text{-}index\text{-}monoid}(s)).\; \\ \quad \quad \quad \quad \operatorname{monalg\text{-}pullback}(r, \operatorname{multi\text{-}index\text{-}monoid}(s), (\lambda b \in (s \to \mathbb{N}).\; \operatorname{mi\text{-}reindex}(s, a, b)), f))) \end{array}}}
\vnbprop{mpoly-is-ring}{\vnbstmt{\forall r, n.\; ((\operatorname{is\text{-}ring}(r) \wedge (n \in \mathbb{N})) \Rightarrow \operatorname{is\text{-}ring}(\operatorname{mpoly}(r, n)))}}
\vnbprop{mpoly-is-commutative-ring}{\vnbstmt{\begin{array}{@{}l@{}} \forall r, n. \\ \quad (\operatorname{is\text{-}commutative\text{-}ring}(r) \wedge (n \in \mathbb{N})) \Rightarrow \operatorname{is\text{-}commutative\text{-}ring}(\operatorname{mpoly}(r, n)) \end{array}}}
\vnbprop{var-permute-is-ring-iso}{\vnbstmt{\begin{array}{@{}l@{}} \forall r, n, s. \\ \quad (\operatorname{is\text{-}ring}(r) \wedge ((n \in \mathbb{N}) \wedge (s \in \mathrm{Bij}(\operatorname{interval}(1, n), \operatorname{interval}(1, n))))) \Rightarrow  \\ \quad \quad \operatorname{is\text{-}iso\text{-}ring}(\operatorname{mpoly}(r, n), \operatorname{mpoly}(r, n), (\lambda p \in \operatorname{carr}(\operatorname{mpoly}(r, n)).\; \operatorname{var\text{-}permute}(r, n, s, p))) \end{array}}}

\vnbbookfile{structure-library/poly-degree}{structure-library/poly-degree}
\noindent{\small poly-degree.scm -- DEG, LEADCOEF and MONOMIAL for the monoid algebra's polynomial specialisation A[NN].}\par
\par\medskip\noindent\textbf{Definition} (\texttt{deg}).\;$\operatorname{deg}(a, p) :=$
\begin{equation*}
(\iota\, \operatorname{n\_}.\;((\operatorname{n\_} \in \mathbb{N}) \wedge (\operatorname{deg\text{-}bound}(a, p, \operatorname{n\_}) \wedge \forall \operatorname{j\_} \in \mathbb{N}.\; (\operatorname{deg\text{-}bound}(a, p, \operatorname{j\_}) \Rightarrow (\operatorname{n\_} \leq \operatorname{j\_})))))
\end{equation*}

\par\medskip\noindent\textbf{Definition} (\texttt{leadcoef}).\;$\operatorname{leadcoef}(a, p) :=$ $p(\operatorname{deg}(a, p))$
\par\medskip\noindent\textbf{Definition} (\texttt{monomial}).\;$\operatorname{monomial}(a, c, n) :=$ $(\lambda \operatorname{x\_} \in \mathbb{N}.\; \left(\text{if }(\operatorname{x\_} = n)\text{ then }c\text{ else }\operatorname{zero}(a)\right))$
\par\medskip\noindent\textbf{Definition} (\texttt{deg-bound}).\;For all $a$, $p$ and $n$, 
\begin{equation*}
\operatorname{deg\text{-}bound}(a, p, n) \Leftrightarrow \forall \operatorname{k\_} \in \mathbb{N}.\; ((n < \operatorname{k\_}) \Rightarrow (p(\operatorname{k\_}) = \operatorname{zero}(a)))
\end{equation*}


\vnbbookfile{theorem-library/poly-degree-laws}{theorem-library/poly-degree-laws}
\noindent{\small poly-degree-laws.scm -- DEG / LEADCOEF / MONOMIAL, PROVEN.}\par
\vnbprop{deg-well-defined}{\vnbstmt{\begin{array}{@{}l@{}} \forall a, p. \\ \quad (\operatorname{is\text{-}ring}(a) \wedge (p \in \operatorname{carr}(\operatorname{poly}(a)))) \Rightarrow  \\ \quad \quad \operatorname{deg}(a, p) \in \mathbb{N} \wedge \\ \quad \quad \operatorname{deg\text{-}bound}(a, p, \operatorname{deg}(a, p)) \wedge \\ \quad \quad \forall j \in \mathbb{N}.\; (\operatorname{deg\text{-}bound}(a, p, j) \Rightarrow (\operatorname{deg}(a, p) \leq j)) \end{array}}}
\vnbprop{deg-bound-apply}{\vnbstmt{\begin{array}{@{}l@{}} \forall a, p, n, k. \\ \quad (\operatorname{deg\text{-}bound}(a, p, n) \wedge ((k \in \mathbb{N}) \wedge (n < k))) \Rightarrow (p(k) = \operatorname{zero}(a)) \end{array}}}
\vnbprop{deg-unfold}{\vnbstmt{\begin{array}{@{}l@{}} \forall a, p. \\ \quad (\operatorname{deg}(a, p) \simeq (\iota\, n.\;((n \in \mathbb{N}) \wedge (\operatorname{deg\text{-}bound}(a, p, n) \wedge \forall j \in \mathbb{N}.\; (\operatorname{deg\text{-}bound}(a, p, j) \Rightarrow (n \leq j)))))) \end{array}}}
\vnbprop{leadcoef-unfold}{\vnbstmt{\forall a, p.\; (\operatorname{leadcoef}(a, p) \simeq p(\operatorname{deg}(a, p)))}}
\vnbprop{leadcoef-nonzero-ptwise}{\vnbstmt{\begin{array}{@{}l@{}} \forall a, p. \\ \quad (\operatorname{is\text{-}ring}(a) \wedge ((p \in \operatorname{carr}(\operatorname{poly}(a))) \wedge \exists k \in \mathbb{N}.\; \neg (p(k) = \operatorname{zero}(a)))) \Rightarrow  \\ \quad \quad \neg (\operatorname{leadcoef}(a, p) = \operatorname{zero}(a)) \end{array}}}
\vnbprop{leadcoef-nonzero}{\vnbstmt{\begin{array}{@{}l@{}} \forall a, p. \\ \quad (\operatorname{is\text{-}ring}(a) \wedge ((p \in \operatorname{carr}(\operatorname{poly}(a))) \wedge \neg (p = \operatorname{zero}(\operatorname{poly}(a))))) \Rightarrow  \\ \quad \quad \neg (\operatorname{leadcoef}(a, p) = \operatorname{zero}(a)) \end{array}}}
\vnbprop{monomial-at}{\vnbstmt{\forall a, c, n.\; (((c \in \operatorname{carr}(a)) \wedge (n \in \mathbb{N})) \Rightarrow (\operatorname{monomial}(a, c, n)(n) = c))}}
\vnbprop{monomial-off}{\vnbstmt{\begin{array}{@{}l@{}} \forall a, c, n, k. \\ \quad (\operatorname{is\text{-}ring}(a) \wedge ((c \in \operatorname{carr}(a)) \wedge ((k \in \mathbb{N}) \wedge \neg (k = n)))) \Rightarrow  \\ \quad \quad \operatorname{monomial}(a, c, n)(k) = \operatorname{zero}(a) \end{array}}}
\vnbprop{monomial-in-poly}{\vnbstmt{\begin{array}{@{}l@{}} \forall a, c, n. \\ \quad (\operatorname{is\text{-}ring}(a) \wedge ((c \in \operatorname{carr}(a)) \wedge (n \in \mathbb{N}))) \Rightarrow  \\ \quad \quad \operatorname{monomial}(a, c, n) \in \operatorname{carr}(\operatorname{poly}(a)) \end{array}}}
\vnbprop{monomial-support}{\vnbstmt{\begin{array}{@{}l@{}} \forall a, c, n, k. \\ \quad (\operatorname{is\text{-}ring}(a) \wedge ((c \in \operatorname{carr}(a)) \wedge (\neg (c = \operatorname{zero}(a)) \wedge (n \in \mathbb{N})))) \Rightarrow  \\ \quad \quad (k \in \operatorname{supp}(a, \operatorname{nn\text{-}add\text{-}monoid}, \operatorname{monomial}(a, c, n))) \Leftrightarrow ((k \in \mathbb{N}) \wedge (k = n)) \end{array}}}
\vnbprop{deg-is}{\vnbstmt{\begin{array}{@{}l@{}} \forall a, p, m. \\ \quad \operatorname{is\text{-}ring}(a) \wedge \\ \quad p \in \operatorname{carr}(\operatorname{poly}(a)) \wedge \\ \quad m \in \mathbb{N} \wedge \\ \quad \operatorname{deg\text{-}bound}(a, p, m) \wedge \\ \quad \forall j \in \mathbb{N}.\; (\operatorname{deg\text{-}bound}(a, p, j) \Rightarrow (m \leq j)) \Rightarrow \\ \quad \quad \operatorname{deg}(a, p) = m \end{array}}}
\vnbprop{monomial-deg}{\vnbstmt{\begin{array}{@{}l@{}} \forall a, c, n. \\ \quad (\operatorname{is\text{-}ring}(a) \wedge ((c \in \operatorname{carr}(a)) \wedge (\neg (c = \operatorname{zero}(a)) \wedge (n \in \mathbb{N})))) \Rightarrow  \\ \quad \quad \operatorname{deg}(a, \operatorname{monomial}(a, c, n)) = n \end{array}}}
\vnbprop{monomial-leadcoef}{\vnbstmt{\begin{array}{@{}l@{}} \forall a, c, n. \\ \quad (\operatorname{is\text{-}ring}(a) \wedge ((c \in \operatorname{carr}(a)) \wedge (\neg (c = \operatorname{zero}(a)) \wedge (n \in \mathbb{N})))) \Rightarrow  \\ \quad \quad \operatorname{leadcoef}(a, \operatorname{monomial}(a, c, n)) = c \end{array}}}
\vnbprop{deg-add-le}{\vnbstmt{\begin{array}{@{}l@{}} \forall a, f, g, n. \\ \quad \operatorname{is\text{-}ring}(a) \wedge \\ \quad f \in \operatorname{carr}(\operatorname{poly}(a)) \wedge \\ \quad g \in \operatorname{carr}(\operatorname{poly}(a)) \wedge \\ \quad n \in \mathbb{N} \wedge \\ \quad \operatorname{deg}(a, f) \leq n \wedge \\ \quad \operatorname{deg}(a, g) \leq n \Rightarrow \\ \quad \quad \operatorname{deg}(a, \operatorname{monalg\text{-}add}(a, \operatorname{nn\text{-}add\text{-}monoid}, f, g)) \leq n \end{array}}}

\vnbbookfile{theorem-library/gauge-is-degree}{theorem-library/gauge-is-degree}
\noindent{\small gauge-is-degree.scm -- the named Euclidean gauge GAUGE(s) is a degree function, PROVEN, together with the four supports its warrant cites.}\par
\vnbprop{euclidean-gauges-unfold}{\vnbstmt{\begin{array}{@{}l@{}} \forall s. \\ \quad (\operatorname{euclidean\text{-}gauges}(s) \simeq \{g \in (\operatorname{carr}(s) \to \mathbb{N}) : \operatorname{has\text{-}div\text{-}remainder}(s, g)\}) \end{array}}}
\vnbprop{euclidean-ring-has-gauge}{\vnbstmt{\begin{array}{@{}l@{}} \forall s. \\ \quad \operatorname{is\text{-}euclidean\text{-}ring}(s) \Rightarrow  \\ \quad \quad \exists g \in (\operatorname{carr}(s) \to \mathbb{N}).\; \operatorname{has\text{-}div\text{-}remainder}(s, g) \end{array}}}
\vnbprop{gauges-mem-build}{\vnbstmt{\begin{array}{@{}l@{}} \forall s.; \forall g \in (\operatorname{carr}(s) \to \mathbb{N}). \\ \quad \operatorname{has\text{-}div\text{-}remainder}(s, g) \Rightarrow (g \in \operatorname{euclidean\text{-}gauges}(s)) \end{array}}}
\vnbprop{gauges-in-fun}{\vnbstmt{\forall s.; \forall g \in \operatorname{euclidean\text{-}gauges}(s).\; (g \in (\operatorname{carr}(s) \to \mathbb{N}))}}
\vnbprop{gauges-spec}{\vnbstmt{\forall s.; \forall g \in \operatorname{euclidean\text{-}gauges}(s).\; \operatorname{has\text{-}div\text{-}remainder}(s, g)}}
\vnbprop{gauge-is-degree}{\vnbstmt{\begin{array}{@{}l@{}} \forall s. \\ \quad \operatorname{is\text{-}euclidean\text{-}ring}(s) \Rightarrow  \\ \quad \quad (\operatorname{gauge}(s) \in (\operatorname{carr}(s) \to \mathbb{N})) \wedge \operatorname{has\text{-}div\text{-}remainder}(s, \operatorname{gauge}(s)) \end{array}}}

\vnbbookfile{theorem-library/euclidean-ideal-generator-proof}{theorem-library/euclidean-ideal-generator-proof}
\noindent{\small euclidean-ideal-generator-proof.scm -- every ideal of a Euclidean ring has a generator: some b in I with I subset (b).}\par
\vnbprop{euclidean-ideal-has-generator}{\vnbstmt{\begin{array}{@{}l@{}} \forall s, i. \\ \quad (\operatorname{is\text{-}euclidean\text{-}ring}(s) \wedge \operatorname{is\text{-}ideal}(s, i)) \Rightarrow  \\ \quad \quad \exists b \in i.\; \forall a \in i.\; (a \in \operatorname{principal\text{-}ideal}(s, b)) \end{array}}}

\vnbbookfile{structure-library/catools-gal-vocab}{structure-library/catools-gal-vocab}
\noindent{\small catools-gal-vocab.scm -- divisibility in an integral domain; Euclidean domains as in the galois-theory notes (Definition 1.2) (CA-TOOLS, galois-theory B1, 2026-10-26). DEFINITIONS ONLY.  Destination: structure-library/catools-gal-vocab.scm, library algebra. Load window: after structure-library/euclidean-ring and ...}\par
\par\medskip\noindent\textbf{Definition} (\texttt{is-euclidean-domain}).\;For every $\operatorname{ges\_}$, 
\begin{equation*}
\begin{array}{@{}l@{}} \operatorname{is\text{-}euclidean\text{-}domain}(\operatorname{ges\_}) \\ \Leftrightarrow\; \operatorname{is\text{-}integral\text{-}domain}(\operatorname{ges\_}) \wedge \exists \operatorname{gen\_}.\; \operatorname{is\text{-}degree\text{-}function}(\operatorname{ges\_}, \operatorname{gen\_}) \end{array}
\end{equation*}

\par\medskip\noindent\textbf{Definition} (\texttt{is-unit}).\;For all $\operatorname{gus\_}$ and $\operatorname{guu\_}$, 
\begin{equation*}
\begin{array}{@{}l@{}} \operatorname{is\text{-}unit}(\operatorname{gus\_}, \operatorname{guu\_}) \\ \Leftrightarrow\; ((\operatorname{guu\_} \in \operatorname{carr}(\operatorname{gus\_})) \wedge \exists \operatorname{guv\_} \in \operatorname{carr}(\operatorname{gus\_}).\; (\operatorname{mul}(\operatorname{gus\_})(\operatorname{guu\_}, \operatorname{guv\_}) = \operatorname{one}(\operatorname{gus\_}))) \end{array}
\end{equation*}

\par\medskip\noindent\textbf{Definition} (\texttt{divides}).\;For all $\operatorname{gds\_}$, $\operatorname{gda\_}$ and $\operatorname{gdb\_}$, 
\begin{equation*}
\begin{array}{@{}l@{}} \operatorname{divides}(\operatorname{gds\_}, \operatorname{gda\_}, \operatorname{gdb\_}) \\ \Leftrightarrow\; \operatorname{gda\_} \in \operatorname{carr}(\operatorname{gds\_}) \wedge \\ \operatorname{gdb\_} \in \operatorname{carr}(\operatorname{gds\_}) \wedge \\ \exists \operatorname{gdk\_} \in \operatorname{carr}(\operatorname{gds\_}).\; (\operatorname{gdb\_} = \operatorname{mul}(\operatorname{gds\_})(\operatorname{gdk\_}, \operatorname{gda\_})) \end{array}
\end{equation*}

\par\medskip\noindent\textbf{Definition} (\texttt{is-prime-element}).\;For all $\operatorname{gps\_}$ and $\operatorname{gpp\_}$, 
\begin{equation*}
\begin{array}{@{}l@{}} \operatorname{is\text{-}prime\text{-}element}(\operatorname{gps\_}, \operatorname{gpp\_}) \\ \Leftrightarrow\; \operatorname{gpp\_} \in \operatorname{carr}(\operatorname{gps\_}) \wedge \\ \neg (\operatorname{gpp\_} = \operatorname{zero}(\operatorname{gps\_})) \wedge \\ \neg \operatorname{is\text{-}unit}(\operatorname{gps\_}, \operatorname{gpp\_}) \wedge \\ \forall \operatorname{gpa\_}, \operatorname{gpb\_} \in \operatorname{carr}(\operatorname{gps\_}). \\ \quad (\operatorname{gpp\_} = \operatorname{mul}(\operatorname{gps\_})(\operatorname{gpa\_}, \operatorname{gpb\_})) \Rightarrow  \\ \quad \quad \operatorname{is\text{-}unit}(\operatorname{gps\_}, \operatorname{gpa\_}) \vee \operatorname{is\text{-}unit}(\operatorname{gps\_}, \operatorname{gpb\_}) \end{array}
\end{equation*}

\par\medskip\noindent\textbf{Definition} (\texttt{is-degree-function}).\;For all $\operatorname{gfs\_}$ and $\operatorname{gfn\_}$, 
\begin{equation*}
\begin{array}{@{}l@{}} \operatorname{is\text{-}degree\text{-}function}(\operatorname{gfs\_}, \operatorname{gfn\_}) \\ \Leftrightarrow\; \operatorname{is\text{-}integral\text{-}domain}(\operatorname{gfs\_}) \wedge \\ \operatorname{gfn\_} \in (\operatorname{carr}(\operatorname{gfs\_}) \to \mathbb{N}) \wedge \\ \forall \operatorname{gfa\_}, \operatorname{gfb\_} \in \operatorname{carr}(\operatorname{gfs\_}). \\ \quad \neg (\operatorname{gfa\_} = \operatorname{zero}(\operatorname{gfs\_})) \Rightarrow  \\ \quad \neg (\operatorname{gfb\_} = \operatorname{zero}(\operatorname{gfs\_})) \Rightarrow (\operatorname{gfn\_}(\operatorname{gfa\_}) \leq \operatorname{gfn\_}(\operatorname{mul}(\operatorname{gfs\_})(\operatorname{gfa\_}, \operatorname{gfb\_}))) \wedge \\ \forall \operatorname{gfa\_}, \operatorname{gfb\_} \in \operatorname{carr}(\operatorname{gfs\_}). \\ \quad \neg (\operatorname{gfa\_} = \operatorname{zero}(\operatorname{gfs\_})) \Rightarrow  \\ \quad \quad \exists \operatorname{gfq\_} \in \operatorname{carr}(\operatorname{gfs\_}). \\ \quad \quad \quad \operatorname{gfb\_} = \operatorname{mul}(\operatorname{gfs\_})(\operatorname{gfa\_}, \operatorname{gfq\_}) \vee \\ \quad \quad \quad \operatorname{gfn\_}(\operatorname{add}(\operatorname{gfs\_})(\operatorname{gfb\_}, \operatorname{neg}(\operatorname{gfs\_})(\operatorname{mul}(\operatorname{gfs\_})(\operatorname{gfa\_}, \operatorname{gfq\_})))) < \operatorname{gfn\_}(\operatorname{gfa\_}) \end{array}
\end{equation*}


\vnbbookfile{theorem-library/catools-gal-b1}{theorem-library/catools-gal-b1}
\noindent{\small catools-gal-b1.scm -- Euclidean domains (galois-theory notes, Lemma 1.7, Corollary 1.5) (CA-TOOLS, galois-theory B1, 2026-10-26).  Destination: theorem-library/catools-gal-b1.scm, library algebra. Vocabulary: catools-gal-vocab.scm (IS-UNIT, DIVIDES, IS-PRIME-ELEMENT, IS-DEGREE-FUNCTION, IS-EUCLIDEAN-DOMAIN).}\par
\vnbprop{degree-proper-divisor}{\vnbstmt{\begin{array}{@{}l@{}} \forall s, n.; \forall a, b, k \in \operatorname{carr}(s). \\ \quad (\operatorname{is\text{-}degree\text{-}function}(s, n) \wedge (\neg (a = \operatorname{zero}(s)) \wedge ((a = \operatorname{mul}(s)(k, b)) \wedge \neg \operatorname{is\text{-}unit}(s, k)))) \Rightarrow  \\ \quad \quad n(b) < n(a) \end{array}}}
\vnbprop{euclidean-domain-is-euclidean-ring}{\vnbstmt{\forall s.\; (\operatorname{is\text{-}euclidean\text{-}domain}(s) \Rightarrow \operatorname{is\text{-}euclidean\text{-}ring}(s))}}
\vnbprop{two-element-ideal}{\vnbstmt{\begin{array}{@{}l@{}} \forall s.; \forall p, a \in \operatorname{carr}(s). \\ \quad \operatorname{is\text{-}commutative\text{-}ring}(s) \Rightarrow  \\ \quad \quad \operatorname{is\text{-}ideal}(s, \\ \quad \quad \quad \{x \in \operatorname{carr}(s) : \exists u, v \in \operatorname{carr}(s).\; (x = \operatorname{add}(s)(\operatorname{mul}(s)(u, p), \operatorname{mul}(s)(v, a)))\}) \end{array}}}
\vnbprop{euclid-lemma}{\vnbstmt{\begin{array}{@{}l@{}} \forall s, p.; \forall a, b \in \operatorname{carr}(s). \\ \quad (\operatorname{is\text{-}euclidean\text{-}domain}(s) \wedge (\operatorname{is\text{-}prime\text{-}element}(s, p) \wedge \operatorname{divides}(s, p, \operatorname{mul}(s)(a, b)))) \Rightarrow  \\ \quad \quad \operatorname{divides}(s, p, a) \vee \operatorname{divides}(s, p, b) \end{array}}}

\vnbbookfile{theorem-library/zz-bezout-proof}{theorem-library/zz-bezout-proof}
\noindent{\small zz-bezout-proof.scm -- BEZOUT on the integers, from the Euclidean-ring core.}\par
\vnbprop{zz-bezout-set-is-ideal}{\vnbstmt{\begin{array}{@{}l@{}} \forall a, b. \\ \quad ((a \in \mathbb{Z}) \wedge (b \in \mathbb{Z})) \Rightarrow \operatorname{is\text{-}ideal}(\operatorname{zz\text{-}ring}, \operatorname{zz\text{-}bezout\text{-}set}(a, b)) \end{array}}}
\vnbprop{zz-bezout}{\vnbstmt{\begin{array}{@{}l@{}} \forall a, b. \\ \quad ((a \in \mathbb{Z}) \wedge (b \in \mathbb{Z})) \Rightarrow  \\ \quad \quad \exists d \in \operatorname{zz\text{-}bezout\text{-}set}(a, b).\; (\operatorname{zz\text{-}divides}(d, a) \wedge \operatorname{zz\text{-}divides}(d, b)) \end{array}}}

\vnbbookfile{theorem-library/euclidean-division-proof}{theorem-library/euclidean-division-proof}
\noindent{\small euclidean-division-proof.scm -- Brick B2 (LA Phase B, algebraic-numbers.pdf ch.3): division-with-remainder in usable form over any euclidean ring.  For a euclidean ring s, a in CARR, b in CARR, b/=0, there are q,r with a = q.b + r and (r = 0 or deg(r) $<$ deg(b)), where deg = GAUGE(s) (the named degree).  This}\par
\vnbprop{euclidean-division}{\vnbstmt{\begin{array}{@{}l@{}} \forall s.; \forall a, b \in \operatorname{carr}(s). \\ \quad (\operatorname{is\text{-}euclidean\text{-}ring}(s) \wedge \neg (b = \operatorname{zero}(s))) \Rightarrow  \\ \quad \quad \exists q, r \in \operatorname{carr}(s). \\ \quad \quad \quad ((a = \operatorname{add}(s)(\operatorname{mul}(s)(q, b), r)) \wedge ((r = \operatorname{zero}(s)) \vee ((\operatorname{gauge}(s)(r) + 1) \leq \operatorname{gauge}(s)(b)))) \end{array}}}

\vnbbookfile{theorem-library/field-one-ne}{theorem-library/field-one-ne}
\noindent{\small field-one-ne.scm -- lemmas MOVED 2026-10-27 (integration of the 10-26/27 wave) out of theorem-library/la-det-adjoint.scm, so that files of library algebra may cite them (the library gate). The proofs are the originals, verbatim; the helper defines are a copy of la-det-adjoint.scm's.  Loaded immediately before ...}\par
\vnbprop{ladet-field-one-ne}{\vnbstmt{\forall r.\; (\operatorname{is\text{-}field\text{-}ring}(r) \Rightarrow \neg (\operatorname{one}(r) = \operatorname{zero}(r)))}}

\vnbbookfile{theorem-library/zz-divides-basics}{theorem-library/zz-divides-basics}
\noindent{\small zz-divides-basics.scm -- lemmas MOVED 2026-10-27 (integration of the 10-26/27 wave) out of theorem-library/la-z-gcd.scm, so that files of library algebra may cite them (the library gate). The proofs are the originals, verbatim; the helper defines are a copy of la-z-gcd.scm's.  Loaded immediately before la-z-gcd.scm.}\par
\vnbprop{zz-mul-eq-one}{\vnbstmt{\forall u, v.\; (((u \in \mathbb{Z}) \wedge ((v \in \mathbb{Z}) \wedge ((u \cdot v) = 1))) \Rightarrow ((u = 1) \vee (u = -1)))}}
\vnbprop{zz-divides-antisym}{\vnbstmt{\forall a, b.\; ((\operatorname{zz\text{-}divides}(a, b) \wedge \operatorname{zz\text{-}divides}(b, a)) \Rightarrow ((a = b) \vee (a = -b)))}}

\vnbbookfile{structure-library/zz-nt-vocab}{structure-library/zz-nt-vocab}
\noindent{\small structure-library/zz-nt-vocab.scm -- vocabulary for ch. 1 of the algebraic-numbers notes (ALGNUM, 2026-10-26). DEFINITIONS ONLY.  Library: algebra.  Load window: after zz-divisibility.}\par
\par\medskip\noindent\textbf{Definition} (\texttt{euler-phi}).\;$\operatorname{euler\text{-}phi}(n) :=$ $|\{\operatorname{ephk\_} \in \operatorname{interval}(1, n) : \operatorname{zz\text{-}coprime}(\operatorname{ephk\_}, n)\}|$
\par\medskip\noindent\textbf{Definition} (\texttt{zz-residue}).\;$\operatorname{zz\text{-}residue}(k, m) :=$ $\varepsilon(\{\operatorname{zrr\_} \in \operatorname{interval}(1, m) : \operatorname{zz\text{-}divides}(m, (k - \operatorname{zrr\_}))\})$
\par\medskip\noindent\textbf{Definition} (\texttt{is-multiplicative}).\;For every $f$, 
\begin{equation*}
\begin{array}{@{}l@{}} \operatorname{is\text{-}multiplicative}(f) \\ \Leftrightarrow\; f \in (\mathbb{N} \to \mathbb{R}) \wedge \\ \forall \operatorname{mlm\_}, \operatorname{mln\_} \in \mathbb{N}. \\ \quad (1 \leq \operatorname{mlm\_}) \Rightarrow  \\ \quad (1 \leq \operatorname{mln\_}) \Rightarrow  \\ \quad \operatorname{zz\text{-}coprime}(\operatorname{mlm\_}, \operatorname{mln\_}) \Rightarrow (f(\cdot(\operatorname{mlm\_}, \operatorname{mln\_})) = \cdot(f(\operatorname{mlm\_}), f(\operatorname{mln\_}))) \end{array}
\end{equation*}


\vnbbookfile{theorem-library/zz-number-theory}{theorem-library/zz-number-theory}
\noindent{\small theorem-library/zz-number-theory.scm -- ALGEBRAIC-NUMBERS NOTES ch. 1: division with remainder, the gcd, the Chinese remainder theorem (ALGNUM, bundles B1 and B4, 2026-10-26). Vocabulary: structure-library/zz-divisibility.scm (ZZ-DIVIDES, ZZ-IS-GCD, ZZ-COPRIME, ZZ-BEZOUT-SET), structure-library/zz-nt-vocab.scm ...}\par
\vnbprop{zz-divides-lincomb}{\vnbstmt{\begin{array}{@{}l@{}} \forall d, x, y, s, t. \\ \quad (\operatorname{zz\text{-}divides}(d, x) \wedge (\operatorname{zz\text{-}divides}(d, y) \wedge ((s \in \mathbb{Z}) \wedge (t \in \mathbb{Z})))) \Rightarrow  \\ \quad \quad \operatorname{zz\text{-}divides}(d, ((s \cdot x) + (t \cdot y))) \end{array}}}
\vnbprop{nn-mul-gap}{\vnbstmt{\begin{array}{@{}l@{}} \forall a, n, m. \\ \quad ((a \in \mathbb{N}) \wedge ((n \in \mathbb{N}) \wedge ((m \in \mathbb{N}) \wedge (n < m)))) \Rightarrow (((a \cdot n) + a) \leq (a \cdot m)) \end{array}}}
\vnbprop{nn-division-remainder}{\vnbstmt{\begin{array}{@{}l@{}} \forall a, b. \\ \quad ((a \in \mathbb{N}) \wedge ((b \in \mathbb{N}) \wedge (1 \leq a))) \Rightarrow  \\ \quad \quad \exists n, r \in \mathbb{N}.\; ((b = ((a \cdot n) + r)) \wedge (r < a)) \end{array}}}
\vnbprop{nn-division-remainder-unique}{\vnbstmt{\begin{array}{@{}l@{}} \forall a, b, n, r, m, s. \\ \quad a \in \mathbb{N} \wedge \\ \quad b \in \mathbb{N} \wedge \\ \quad n \in \mathbb{N} \wedge \\ \quad r \in \mathbb{N} \wedge \\ \quad m \in \mathbb{N} \wedge \\ \quad s \in \mathbb{N} \wedge \\ \quad b = ((a \cdot n) + r) \wedge \\ \quad r < a \wedge \\ \quad b = ((a \cdot m) + s) \wedge \\ \quad s < a \Rightarrow \\ \quad \quad (n = m) \wedge (r = s) \end{array}}}
\vnbprop{nn-floor-unique}{\vnbstmt{\begin{array}{@{}l@{}} \forall a, b. \\ \quad ((a \in \mathbb{N}) \wedge ((b \in \mathbb{N}) \wedge (1 \leq a))) \Rightarrow  \\ \quad \quad \exists n \in \mathbb{N}.\; (((a \cdot n) \leq b) \wedge (b < ((a \cdot n) + a))) \wedge \\ \quad \quad \forall n, m \in \mathbb{N}. \\ \quad \quad \quad ((((a \cdot n) \leq b) \wedge (b < ((a \cdot n) + a))) \wedge (((a \cdot m) \leq b) \wedge (b < ((a \cdot m) + a)))) \Rightarrow (n = m) \end{array}}}
\vnbprop{zz-gcd-remainder}{\vnbstmt{\begin{array}{@{}l@{}} \forall a, b, n, r, d. \\ \quad ((a \in \mathbb{Z}) \wedge ((b \in \mathbb{Z}) \wedge ((n \in \mathbb{Z}) \wedge ((r \in \mathbb{Z}) \wedge (b = ((a \cdot n) + r)))))) \Rightarrow  \\ \quad \quad \operatorname{zz\text{-}is\text{-}gcd}(a, b, d) \Leftrightarrow \operatorname{zz\text{-}is\text{-}gcd}(a, r, d) \end{array}}}
\vnbprop{zz-gcd-in-bezout-set}{\vnbstmt{\begin{array}{@{}l@{}} \forall a, b, g. \\ \quad ((a \in \mathbb{Z}) \wedge ((b \in \mathbb{Z}) \wedge \operatorname{zz\text{-}is\text{-}gcd}(a, b, g))) \Rightarrow (g \in \operatorname{zz\text{-}bezout\text{-}set}(a, b)) \end{array}}}
\vnbprop{euclid-two-steps-halve}{\vnbstmt{\begin{array}{@{}l@{}} \forall a, b, c, d, f, g. \\ \quad a \in \mathbb{N} \wedge \\ \quad b \in \mathbb{N} \wedge \\ \quad c \in \mathbb{N} \wedge \\ \quad d \in \mathbb{N} \wedge \\ \quad f \in \mathbb{N} \wedge \\ \quad g \in \mathbb{N} \wedge \\ \quad b = ((a \cdot c) + d) \wedge \\ \quad d < a \wedge \\ \quad 1 \leq d \wedge \\ \quad a = ((d \cdot f) + g) \wedge \\ \quad g < d \Rightarrow \\ \quad \quad (2 \cdot g) < a \end{array}}}
\vnbprop{zz-divides-mul-right}{\vnbstmt{\forall d, x, y.\; ((\operatorname{zz\text{-}divides}(d, x) \wedge (y \in \mathbb{Z})) \Rightarrow \operatorname{zz\text{-}divides}(d, (x \cdot y)))}}
\vnbprop{zz-divides-trans}{\vnbstmt{\begin{array}{@{}l@{}} \forall d, x, y. \\ \quad (\operatorname{zz\text{-}divides}(d, x) \wedge \operatorname{zz\text{-}divides}(x, y)) \Rightarrow \operatorname{zz\text{-}divides}(d, y) \end{array}}}
\vnbprop{zz-coprime-bezout}{\vnbstmt{\begin{array}{@{}l@{}} \forall a, b. \\ \quad ((a \in \mathbb{Z}) \wedge (b \in \mathbb{Z})) \Rightarrow  \\ \quad \quad \operatorname{zz\text{-}coprime}(a, b) \Leftrightarrow \exists x, y \in \mathbb{Z}.\; (((x \cdot a) + (y \cdot b)) = 1) \end{array}}}
\vnbprop{zz-coprime-mul}{\vnbstmt{\begin{array}{@{}l@{}} \forall a, b, m. \\ \quad ((a \in \mathbb{Z}) \wedge ((b \in \mathbb{Z}) \wedge ((m \in \mathbb{Z}) \wedge (\operatorname{zz\text{-}coprime}(a, m) \wedge \operatorname{zz\text{-}coprime}(b, m))))) \Rightarrow  \\ \quad \quad \operatorname{zz\text{-}coprime}((a \cdot b), m) \end{array}}}
\vnbprop{zz-crt-strong}{\vnbstmt{\begin{array}{@{}l@{}} \forall k \in \mathbb{N}, n \in (\mathbb{N} \to \mathbb{Z}). \\ \quad \forall i, j \in \mathbb{N}.\; (((i < k) \wedge ((j < k) \wedge \neg (i = j))) \Rightarrow \operatorname{zz\text{-}coprime}(n(i), n(j))) \Rightarrow  \\ \quad \quad \exists m \in \mathbb{Z}. \\ \quad \quad \quad \forall i \in \mathbb{N}.\; ((i < k) \Rightarrow \operatorname{zz\text{-}divides}(n(i), m)) \wedge \\ \quad \quad \quad \forall q \in \mathbb{Z}. \\ \quad \quad \quad \quad \forall i \in \mathbb{N}.\; ((i < k) \Rightarrow \operatorname{zz\text{-}coprime}(n(i), q)) \Rightarrow \operatorname{zz\text{-}coprime}(m, q) \wedge \\ \quad \quad \quad \forall a \in (\mathbb{N} \to \mathbb{Z}). \exists x \in \mathbb{Z}. \forall i \in \mathbb{N}. \\ \quad \quad \quad \quad (i < k) \Rightarrow \operatorname{zz\text{-}divides}(n(i), (x - a(i))) \end{array}}}
\vnbprop{zz-chinese-remainder}{\vnbstmt{\begin{array}{@{}l@{}} \forall k, n. \\ \quad ((k \in \mathbb{N}) \wedge (n \in (\mathbb{N} \to \mathbb{Z}))) \Rightarrow  \\ \quad \quad \forall i, j \in \mathbb{N}. \\ \quad \quad \quad ((i < k) \wedge ((j < k) \wedge \neg (i = j))) \Rightarrow \operatorname{zz\text{-}coprime}(n(i), n(j)) \\ \quad \quad \Leftrightarrow\; \forall a \in (\mathbb{N} \to \mathbb{Z}). \exists x \in \mathbb{Z}. \forall i \in \mathbb{N}. \\ \quad \quad \quad (i < k) \Rightarrow \operatorname{zz\text{-}divides}(n(i), (x - a(i))) \end{array}}}
\vnbprop{zz-linear-congruences}{\vnbstmt{\begin{array}{@{}l@{}} \forall k, n, c.; \forall a \in (\mathbb{N} \to \mathbb{Z}). \\ \quad k \in \mathbb{N} \wedge \\ \quad n \in (\mathbb{N} \to \mathbb{Z}) \wedge \\ \quad c \in (\mathbb{N} \to \mathbb{Z}) \wedge \\ \quad \forall i, j \in \mathbb{N}. \\ \quad \quad ((i < k) \wedge ((j < k) \wedge \neg (i = j))) \Rightarrow \operatorname{zz\text{-}coprime}(n(i), n(j)) \wedge \\ \quad \forall i \in \mathbb{N}.\; ((i < k) \Rightarrow \operatorname{zz\text{-}coprime}(c(i), n(i))) \Rightarrow \\ \quad \quad \exists x \in \mathbb{Z}. \forall i \in \mathbb{N}. \\ \quad \quad \quad (i < k) \Rightarrow \operatorname{zz\text{-}divides}(n(i), ((c(i) \cdot x) - a(i))) \end{array}}}
\vnbprop{zz-small-multiple-zero}{\vnbstmt{\begin{array}{@{}l@{}} \forall m, x. \\ \quad ((m \in \mathbb{Z}) \wedge ((x \in \mathbb{Z}) \wedge ((1 \leq m) \wedge (\operatorname{zz\text{-}divides}(m, x) \wedge ((x < m) \wedge (-m < x)))))) \Rightarrow  \\ \quad \quad x = 0 \end{array}}}
\vnbprop{zz-residue-in}{\vnbstmt{\begin{array}{@{}l@{}} \forall m, k. \\ \quad ((m \in \mathbb{N}) \wedge ((1 \leq m) \wedge (k \in \mathbb{Z}))) \Rightarrow  \\ \quad \quad (\operatorname{zz\text{-}residue}(k, m) \in \operatorname{interval}(1, m)) \wedge \operatorname{zz\text{-}divides}(m, (k - \operatorname{zz\text{-}residue}(k, m))) \end{array}}}
\vnbprop{zz-residue-unique}{\vnbstmt{\begin{array}{@{}l@{}} \forall m, k, r. \\ \quad ((m \in \mathbb{N}) \wedge ((1 \leq m) \wedge ((k \in \mathbb{Z}) \wedge ((r \in \operatorname{interval}(1, m)) \wedge \operatorname{zz\text{-}divides}(m, (k - r)))))) \Rightarrow  \\ \quad \quad r = \operatorname{zz\text{-}residue}(k, m) \end{array}}}
\vnbprop{zz-coprime-congruent}{\vnbstmt{\begin{array}{@{}l@{}} \forall m, k, r. \\ \quad ((m \in \mathbb{Z}) \wedge ((k \in \mathbb{Z}) \wedge ((r \in \mathbb{Z}) \wedge (\operatorname{zz\text{-}divides}(m, (k - r)) \wedge \operatorname{zz\text{-}coprime}(k, m))))) \Rightarrow  \\ \quad \quad \operatorname{zz\text{-}coprime}(r, m) \end{array}}}
\vnbprop{zz-coprime-sym}{\vnbstmt{\begin{array}{@{}l@{}} \forall a, b. \\ \quad ((a \in \mathbb{Z}) \wedge ((b \in \mathbb{Z}) \wedge \operatorname{zz\text{-}coprime}(a, b))) \Rightarrow \operatorname{zz\text{-}coprime}(b, a) \end{array}}}
\vnbprop{zz-coprime-factor}{\vnbstmt{\begin{array}{@{}l@{}} \forall a, b, c. \\ \quad ((a \in \mathbb{Z}) \wedge ((b \in \mathbb{Z}) \wedge ((c \in \mathbb{Z}) \wedge \operatorname{zz\text{-}coprime}(a, (b \cdot c))))) \Rightarrow  \\ \quad \quad \operatorname{zz\text{-}coprime}(a, b) \end{array}}}
\vnbprop{zz-coprime-divides-product}{\vnbstmt{\begin{array}{@{}l@{}} \forall m, n, x. \\ \quad ((m \in \mathbb{Z}) \wedge ((n \in \mathbb{Z}) \wedge ((x \in \mathbb{Z}) \wedge (\operatorname{zz\text{-}coprime}(m, n) \wedge (\operatorname{zz\text{-}divides}(m, x) \wedge \operatorname{zz\text{-}divides}(n, x)))))) \Rightarrow  \\ \quad \quad \operatorname{zz\text{-}divides}((m \cdot n), x) \end{array}}}
\vnbprop{euler-phi-multiplicative}{\vnbstmt{\begin{array}{@{}l@{}} \forall m, n. \\ \quad ((m \in \mathbb{N}) \wedge ((n \in \mathbb{N}) \wedge ((1 \leq m) \wedge ((1 \leq n) \wedge \operatorname{zz\text{-}coprime}(m, n))))) \Rightarrow  \\ \quad \quad \operatorname{euler\text{-}phi}((m \cdot n)) = (\operatorname{euler\text{-}phi}(m) \cdot \operatorname{euler\text{-}phi}(n)) \end{array}}}

\vnbbookfile{theorem-library/prod-ring-succ}{theorem-library/prod-ring-succ}
\noindent{\small prod-ring-succ.scm -- lemmas MOVED 2026-10-27 (integration of the 10-26/27 wave) out of theorem-library/la-z2-diag.scm, so that files of library algebra may cite them (the library gate). The proofs are the originals, verbatim; the helper defines are a copy of la-z2-diag.scm's.  Loaded immediately before la-z2-diag.scm.}\par
\vnbprop{lzd-prod-succ}{\vnbstmt{\begin{array}{@{}l@{}} \forall r, s, g. \\ \quad (\operatorname{is\text{-}commutative\text{-}ring}(r) \wedge ((s \in \mathbb{N}) \wedge \forall z \in \operatorname{interval}(1, (s + 1)).\; (g(z) \in \operatorname{carr}(r)))) \Rightarrow  \\ \quad \quad \operatorname{prod\text{-}ring}(r, g, \operatorname{interval}(1, (s + 1))) = \operatorname{mul}(r)(\operatorname{prod\text{-}ring}(r, g, \operatorname{interval}(1, s)), g((s + 1))) \end{array}}}

\vnbbookfile{structure-library/gal-ufd-vocab}{structure-library/gal-ufd-vocab}
\noindent{\small gal-ufd-vocab.scm -- associates and unique factorization domains (galois-theory notes, section 1) (GALOIS, 2026-10-26).  DEFINITIONS ONLY.  Destination: structure-library/gal-ufd-vocab.scm, library algebra.  Load window: after scratchpad/catools/catools-gal-vocab (IS-UNIT, IS-PRIME-ELEMENT) and ...}\par
\par\medskip\noindent\textbf{Definition} (\texttt{is-ufd}).\;For every $\operatorname{gus\_}$, 
\begin{equation*}
\begin{array}{@{}l@{}} \operatorname{is\text{-}ufd}(\operatorname{gus\_}) \\ \Leftrightarrow\; \operatorname{is\text{-}integral\text{-}domain}(\operatorname{gus\_}) \wedge \\ \forall \operatorname{gua\_} \in \operatorname{carr}(\operatorname{gus\_}). \\ \quad \neg (\operatorname{gua\_} = \operatorname{zero}(\operatorname{gus\_})) \Rightarrow  \\ \quad \quad \exists \operatorname{guu\_}. \\ \quad \quad \quad \operatorname{is\text{-}unit}(\operatorname{gus\_}, \operatorname{guu\_}) \wedge \\ \quad \quad \quad \exists \operatorname{gun\_} \in \mathbb{N}, \operatorname{gup\_}. \\ \quad \quad \quad \quad \operatorname{is\text{-}prime\text{-}seq}(\operatorname{gus\_}, \operatorname{gup\_}, \operatorname{gun\_}) \wedge \\ \quad \quad \quad \quad \operatorname{gua\_} = \operatorname{mul}(\operatorname{gus\_})(\operatorname{guu\_}, \operatorname{prod\text{-}ring}(\operatorname{gus\_}, \operatorname{gup\_}, \operatorname{interval}(1, \operatorname{gun\_}))) \wedge \\ \forall \operatorname{gun\_} \in \mathbb{N}, \operatorname{gup\_}. \\ \quad \operatorname{is\text{-}prime\text{-}seq}(\operatorname{gus\_}, \operatorname{gup\_}, \operatorname{gun\_}) \Rightarrow  \\ \quad \quad \forall \operatorname{gum\_} \in \mathbb{N}, \operatorname{guq\_}. \\ \quad \quad \quad \operatorname{is\text{-}prime\text{-}seq}(\operatorname{gus\_}, \operatorname{guq\_}, \operatorname{gum\_}) \Rightarrow  \\ \quad \quad \quad \quad \forall \operatorname{guu\_}. \\ \quad \quad \quad \quad \quad \operatorname{is\text{-}unit}(\operatorname{gus\_}, \operatorname{guu\_}) \Rightarrow  \\ \quad \quad \quad \quad \quad \quad \forall \operatorname{guv\_}. \\ \quad \quad \quad \quad \quad \quad \quad \operatorname{is\text{-}unit}(\operatorname{gus\_}, \operatorname{guv\_}) \Rightarrow  \\ \quad \quad \quad \quad \quad \quad \quad \operatorname{mul}(\operatorname{gus\_})(\operatorname{guu\_}, \operatorname{prod\text{-}ring}(\operatorname{gus\_}, \operatorname{gup\_}, \operatorname{interval}(1, \operatorname{gun\_}))) \\ \quad \quad \quad \quad \quad \quad \quad =\; \operatorname{mul}(\operatorname{gus\_})(\operatorname{guv\_}, \operatorname{prod\text{-}ring}(\operatorname{gus\_}, \operatorname{guq\_}, \operatorname{interval}(1, \operatorname{gum\_}))) \Rightarrow \\ \quad \quad \quad \quad \quad \quad \quad \quad \operatorname{gun\_} = \operatorname{gum\_} \wedge \\ \quad \quad \quad \quad \quad \quad \quad \quad \exists \operatorname{gus2\_} \in \mathrm{Bij}(\operatorname{interval}(1, \operatorname{gun\_}), \operatorname{interval}(1, \operatorname{gum\_})). \forall \operatorname{gui\_} \in \operatorname{interval}(1, \operatorname{gun\_}). \\ \quad \quad \quad \quad \quad \quad \quad \quad \quad \operatorname{associated}(\operatorname{gus\_}, \operatorname{gup\_}(\operatorname{gui\_}), \operatorname{guq\_}(\operatorname{gus2\_}(\operatorname{gui\_}))) \end{array}
\end{equation*}

\par\medskip\noindent\textbf{Definition} (\texttt{associated}).\;For all $\operatorname{gss\_}$, $\operatorname{gsa\_}$ and $\operatorname{gsb\_}$, 
\begin{equation*}
\begin{array}{@{}l@{}} \operatorname{associated}(\operatorname{gss\_}, \operatorname{gsa\_}, \operatorname{gsb\_}) \\ \Leftrightarrow\; \operatorname{gsa\_} \in \operatorname{carr}(\operatorname{gss\_}) \wedge \\ \operatorname{gsb\_} \in \operatorname{carr}(\operatorname{gss\_}) \wedge \\ \exists \operatorname{gsw\_}.\; (\operatorname{is\text{-}unit}(\operatorname{gss\_}, \operatorname{gsw\_}) \wedge (\operatorname{gsb\_} = \operatorname{mul}(\operatorname{gss\_})(\operatorname{gsw\_}, \operatorname{gsa\_}))) \end{array}
\end{equation*}

\par\medskip\noindent\textbf{Definition} (\texttt{is-prime-seq}).\;For all $\operatorname{gqs\_}$, $\operatorname{gqp\_}$ and $\operatorname{gqn\_}$, 
\begin{equation*}
\begin{array}{@{}l@{}} \operatorname{is\text{-}prime\text{-}seq}(\operatorname{gqs\_}, \operatorname{gqp\_}, \operatorname{gqn\_}) \\ \Leftrightarrow\; \forall \operatorname{gqi\_} \in \operatorname{interval}(1, \operatorname{gqn\_}).\; \operatorname{is\text{-}prime\text{-}element}(\operatorname{gqs\_}, \operatorname{gqp\_}(\operatorname{gqi\_})) \end{array}
\end{equation*}


\vnbbookfile{theorem-library/gal-ufd}{theorem-library/gal-ufd}
\noindent{\small gal-ufd.scm -- a Euclidean domain is a unique factorization domain (galois-theory notes, Corollary 1.6 and Proposition 1.8) (GALOIS, 2026-10-26). Destination: theorem-library/gal-ufd.scm, library algebra.  Vocabulary: catools-gal-vocab.scm, gal-ufd-vocab.scm; uses catools-gal-b1 (degree-proper-divisor, euclid-lemma).}\par
\vnbprop{prime-divides-prime-product}{\vnbstmt{\begin{array}{@{}l@{}} \forall n \in \mathbb{N}, s, p, q. \\ \quad \operatorname{is\text{-}euclidean\text{-}domain}(s) \wedge \\ \quad \operatorname{is\text{-}prime\text{-}element}(s, p) \wedge \\ \quad \operatorname{is\text{-}prime\text{-}seq}(s, q, n) \wedge \\ \quad \operatorname{divides}(s, p, \operatorname{prod\text{-}ring}(s, q, \operatorname{interval}(1, n))) \Rightarrow \\ \quad \quad \exists k \in \operatorname{interval}(1, n).\; \operatorname{associated}(s, p, q(k)) \end{array}}}
\vnbprop{prime-divisor-bounded}{\vnbstmt{\begin{array}{@{}l@{}} \forall b \in \mathbb{N}, s, n.; \forall a \in \operatorname{carr}(s). \\ \quad (\operatorname{is\text{-}degree\text{-}function}(s, n) \wedge ((n(a) < b) \wedge (\neg (a = \operatorname{zero}(s)) \wedge \neg \operatorname{is\text{-}unit}(s, a)))) \Rightarrow  \\ \quad \quad \exists p. \\ \quad \quad \quad \operatorname{is\text{-}prime\text{-}element}(s, p) \wedge \exists d \in \operatorname{carr}(s).\; (a = \operatorname{mul}(s)(p, d)) \end{array}}}
\vnbprop{factorization-exists-bounded}{\vnbstmt{\begin{array}{@{}l@{}} \forall b \in \mathbb{N}, s, n.; \forall a \in \operatorname{carr}(s). \\ \quad (\operatorname{is\text{-}degree\text{-}function}(s, n) \wedge ((n(a) < b) \wedge \neg (a = \operatorname{zero}(s)))) \Rightarrow  \\ \quad \quad \exists u \in \operatorname{carr}(s). \\ \quad \quad \quad \operatorname{is\text{-}unit}(s, u) \wedge \\ \quad \quad \quad \exists \mathit{gun} \in \mathbb{N}, p \in (\operatorname{interval}(1, \mathit{gun}) \to \operatorname{carr}(s)). \\ \quad \quad \quad \quad \operatorname{is\text{-}prime\text{-}seq}(s, p, \mathit{gun}) \wedge (a = \operatorname{mul}(s)(u, \operatorname{prod\text{-}ring}(s, p, \operatorname{interval}(1, \mathit{gun})))) \end{array}}}
\vnbprop{integral-domain-cancel-right}{\vnbstmt{\begin{array}{@{}l@{}} \forall s.; \forall a, b, c \in \operatorname{carr}(s). \\ \quad (\operatorname{is\text{-}integral\text{-}domain}(s) \wedge (\neg (c = \operatorname{zero}(s)) \wedge (\operatorname{mul}(s)(a, c) = \operatorname{mul}(s)(b, c)))) \Rightarrow  \\ \quad \quad a = b \end{array}}}
\vnbprop{interval-transposition-exists}{\vnbstmt{\begin{array}{@{}l@{}} \forall m \in \mathbb{N}, k \in \operatorname{interval}(1, m). \exists h \in \mathrm{Bij}(\operatorname{interval}(1, m), \operatorname{interval}(1, m)). \\ \quad h(m) = k \end{array}}}
\vnbprop{prod-ring-permute-ptwise}{\vnbstmt{\begin{array}{@{}l@{}} \forall s.; \forall \mathit{xs} \in \mathbf{Set}, h \in \mathrm{Bij}(\mathit{xs}, \mathit{xs}), q. \\ \quad (\operatorname{is\text{-}commutative\text{-}ring}(s) \wedge ((|\mathit{xs}| \in \mathbb{N}) \wedge \forall z \in \mathit{xs}.\; (q(z) \in \operatorname{carr}(s)))) \Rightarrow  \\ \quad \quad \operatorname{prod\text{-}ring}(s, q, \mathit{xs}) = \operatorname{prod\text{-}ring}(s, (\lambda y \in \mathit{xs}.\; q(h(y))), \mathit{xs}) \end{array}}}
\vnbprop{bijection-extend-top}{\vnbstmt{\begin{array}{@{}l@{}} \forall n \in \mathbb{N}, g \in \mathrm{Bij}(\operatorname{interval}(1, n), \operatorname{interval}(1, n)). \\ \quad (\lambda e \in \operatorname{interval}(1, (n + 1)).\; \left(\text{if }(e = (n + 1))\text{ then }(n + 1)\text{ else }g(e)\right)) \\ \quad \in\; \mathrm{Bij}(\operatorname{interval}(1, (n + 1)), \operatorname{interval}(1, (n + 1))) \end{array}}}
\vnbprop{factorization-unique}{\vnbstmt{\begin{array}{@{}l@{}} \forall n \in \mathbb{N}, s, p.; \forall m \in \mathbb{N}, q \in (\operatorname{interval}(1, m) \to \operatorname{carr}(s)), u.; \forall v \in \operatorname{carr}(s). \\ \quad \operatorname{is\text{-}euclidean\text{-}domain}(s) \wedge \\ \quad \operatorname{is\text{-}prime\text{-}seq}(s, p, n) \wedge \\ \quad \operatorname{is\text{-}prime\text{-}seq}(s, q, m) \wedge \\ \quad \operatorname{is\text{-}unit}(s, u) \wedge \\ \quad \operatorname{is\text{-}unit}(s, v) \wedge \\ \quad \operatorname{mul}(s)(u, \operatorname{prod\text{-}ring}(s, p, \operatorname{interval}(1, n))) = \operatorname{mul}(s)(v, \operatorname{prod\text{-}ring}(s, q, \operatorname{interval}(1, m))) \Rightarrow \\ \quad \quad n = m \wedge \\ \quad \quad \exists s_{2} \in \mathrm{Bij}(\operatorname{interval}(1, n), \operatorname{interval}(1, m)). \forall i \in \operatorname{interval}(1, n). \\ \quad \quad \quad \operatorname{associated}(s, p(i), q(\operatorname{s2}(i))) \end{array}}}
\vnbprop{euclidean-domain-is-ufd}{\vnbstmt{\forall s.\; (\operatorname{is\text{-}euclidean\text{-}domain}(s) \Rightarrow \operatorname{is\text{-}ufd}(s))}}

\vnbbookfile{theorem-library/binomial-proof}{theorem-library/binomial-proof}
\noindent{\small binomial-proof.scm -- machine proof of the Binomial Theorem (SUM form), PROVEN by ordinary induction, IMPS-style (docs/algebra.t).}\par
\vnbprop{sum-expansion}{\vnbstmt{\begin{array}{@{}l@{}} \forall n \in \mathbb{N}, r.; \forall x, y \in \operatorname{carr}(r), g \in (\mathbb{Z} \to \operatorname{carr}(r)). \\ \quad \operatorname{is\text{-}commutative\text{-}ring}(r) \Rightarrow  \\ \quad \quad \operatorname{sum}(r, (\lambda k \in \mathbb{Z}.\; \operatorname{add}(r)(\operatorname{mul}(r)(x, g((k - 1))), \operatorname{mul}(r)(y, g(k)))), (n + 1)) \\ \quad \quad =\; \operatorname{add}(r)(\operatorname{add}(r)(\operatorname{mul}(r)(\operatorname{add}(r)(x, y), \operatorname{sum}(r, g, n)), \operatorname{mul}(r)(x, g(-1))), \operatorname{mul}(r)(y, g(n))) \end{array}}}
\vnbprop{binomial-theorem}{\vnbstmt{\begin{array}{@{}l@{}} \forall n \in \mathbb{N}, r.; \forall x, y \in \operatorname{carr}(r). \\ \quad \operatorname{is\text{-}commutative\text{-}ring}(r) \Rightarrow  \\ \quad \quad \operatorname{ring\text{-}power}(r, \operatorname{add}(r)(x, y), n) = \operatorname{sum}(r, \operatorname{comb\text{-}kk}(r, x, y, n), (n + 1)) \end{array}}}

\vnbbookfile{structure-library/gal-char-vocab}{structure-library/gal-char-vocab}
\noindent{\small gal-char-vocab.scm -- prime natural numbers and the characteristic of a field (galois-theory notes, section 3) (GALOIS B8, 2026-10-26).  DEFINITIONS ONLY.  Destination: structure-library/gal-char-vocab.scm, library algebra.  Load window: after catools-gal-vocab (IS-PRIME-ELEMENT), zz-action.scm (ZZ-ACT).}\par
\par\medskip\noindent\textbf{Definition} (\texttt{is-prime-nat}).\;For every $\operatorname{gpn\_}$, 
\begin{equation*}
\operatorname{is\text{-}prime\text{-}nat}(\operatorname{gpn\_}) \Leftrightarrow ((\operatorname{gpn\_} \in \mathbb{N}) \wedge \operatorname{is\text{-}prime\text{-}element}(\operatorname{zz\text{-}ring}, \operatorname{gpn\_}))
\end{equation*}

\par\medskip\noindent\textbf{Definition} (\texttt{has-characteristic}).\;For all $\operatorname{ghc\_}$ and $\operatorname{gcp\_}$, 
\begin{equation*}
\begin{array}{@{}l@{}} \operatorname{has\text{-}characteristic}(\operatorname{ghc\_}, \operatorname{gcp\_}) \\ \Leftrightarrow\; (\operatorname{is\text{-}prime\text{-}nat}(\operatorname{gcp\_}) \wedge (\operatorname{zz\text{-}act}(\operatorname{ring\text{-}additive\text{-}ag}(\operatorname{ghc\_}), \operatorname{gcp\_}, \operatorname{one}(\operatorname{ghc\_})) = \operatorname{zero}(\operatorname{ghc\_}))) \end{array}
\end{equation*}


\vnbbookfile{theorem-library/zz-ring-arith}{theorem-library/zz-ring-arith}
\noindent{\small zz-ring-arith.scm -- lemmas MOVED 2026-10-27 (integration of the 10-26/27 wave) out of theorem-library/la-z5-zzmod.scm, so that files of library algebra may cite them (the library gate). The proofs are the originals, verbatim; the helper defines are a copy of la-z5-zzmod.scm's.  Loaded immediately before ...}\par
\vnbprop{zz-ring-add-neg}{\vnbstmt{\begin{array}{@{}l@{}} \forall x, y. \\ \quad ((x \in \mathbb{Z}) \wedge (y \in \mathbb{Z})) \Rightarrow (\operatorname{add}(\operatorname{zz\text{-}ring})(x, \operatorname{neg}(\operatorname{zz\text{-}ring})(y)) = (x + -y)) \end{array}}}
\vnbprop{zz-ring-mul-is-mul}{\vnbstmt{\forall t, b.\; (((t \in \mathbb{Z}) \wedge (b \in \mathbb{Z})) \Rightarrow (\operatorname{mul}(\operatorname{zz\text{-}ring})(t, b) = (t \cdot b)))}}

\vnbbookfile{theorem-library/poly-is-ring-proof}{theorem-library/poly-is-ring-proof}
\noindent{\small poly-is-ring-proof.scm -- POLY(A) is a ring when A is a ring.}\par
\vnbprop{poly-is-ring}{\vnbstmt{\forall a.\; (\operatorname{is\text{-}ring}(a) \Rightarrow \operatorname{is\text{-}ring}(\operatorname{poly}(a)))}}

\vnbbookfile{structure-library/gal-poly-vocab}{structure-library/gal-poly-vocab}
\noindent{\small gal-poly-vocab.scm -- evaluation of a formal polynomial and the polynomial x - c (galois-theory notes, section 8: "the evaluation mapping F[x] x F -$>$ F, (p(x), a) |-$>$ p(a)"; "a root of a polynomial is an element a in F such that p(a) = 0") (GALOIS B2, 2026-10-26). DEFINITIONS ONLY.  Destination: ...}\par
\par\medskip\noindent\textbf{Definition} (\texttt{poly-eval}).\;$\operatorname{poly\text{-}eval}(a, p, t) :=$
\begin{equation*}
\begin{array}{@{}l@{}} \operatorname{sum\text{-}ag}(\operatorname{ring\text{-}additive\text{-}ag}(a), \\ \quad (\lambda \operatorname{gek\_} \in \mathbb{N}.\; \operatorname{mul}(a)(p(\operatorname{gek\_}), \operatorname{ring\text{-}power}(a, t, \operatorname{gek\_}))), \\ \quad (\operatorname{deg}(a, p) + 1)) \end{array}
\end{equation*}

\par\medskip\noindent\textbf{Definition} (\texttt{x-minus}).\;$\operatorname{x\text{-}minus}(a, c) :=$
\begin{equation*}
\operatorname{add}(\operatorname{poly}(a))(\operatorname{monomial}(a, \operatorname{one}(a), 1), \operatorname{monomial}(a, \operatorname{neg}(a)(c), 0))
\end{equation*}

\par\medskip\noindent\textbf{Definition} (\texttt{poly-le}).\;$\operatorname{poly\text{-}le}(a, m) :=$ $\{\operatorname{gpl\_} \in \operatorname{carr}(\operatorname{poly}(a)) : \operatorname{deg\text{-}bound}(a, \operatorname{gpl\_}, m)\}$
\par\medskip\noindent\textbf{Definition} (\texttt{eval-at-points}).\;$\operatorname{eval\text{-}at\text{-}points}(a, \mathit{th}, m) :=$
\begin{equation*}
\begin{array}{@{}l@{}} (\lambda \operatorname{gep\_} \in \operatorname{poly\text{-}le}(a, m).\; \\ \quad (\lambda \operatorname{gei\_} \in \operatorname{interval}(1, (m + 1)).\; \operatorname{poly\text{-}eval}(a, \operatorname{gep\_}, \operatorname{th}(\operatorname{gei\_})))) \end{array}
\end{equation*}


\vnbbookfile{theorem-library/gal-poly}{theorem-library/gal-poly}
\noindent{\small gal-poly.scm -- F[x]: degree of a product, division with remainder, F[x] a Euclidean domain, a PID and a UFD, evaluation and roots (galois-theory notes, Proposition 1.9, Corollary 1.10, Proposition 1.43, Corollaries 1.45, 1.46, 1.48) (GALOIS B2, 2026-10-26). Destination: theorem-library/gal-poly.scm, library algebra.  ...}\par
\vnbprop{poly-add-at}{\vnbstmt{\begin{array}{@{}l@{}} \forall a.; \forall f, g \in \operatorname{carr}(\operatorname{poly}(a)), k \in \mathbb{N}. \\ \quad \operatorname{is\text{-}ring}(a) \Rightarrow (\operatorname{add}(\operatorname{poly}(a))(f, g)(k) = \operatorname{add}(a)(f(k), g(k))) \end{array}}}
\vnbprop{poly-neg-at}{\vnbstmt{\begin{array}{@{}l@{}} \forall a.; \forall f \in \operatorname{carr}(\operatorname{poly}(a)), k \in \mathbb{N}. \\ \quad \operatorname{is\text{-}ring}(a) \Rightarrow (\operatorname{neg}(\operatorname{poly}(a))(f)(k) = \operatorname{neg}(a)(f(k))) \end{array}}}
\vnbprop{poly-mul-unfold}{\vnbstmt{\begin{array}{@{}l@{}} \forall a.; \forall f, g \in \operatorname{carr}(\operatorname{poly}(a)). \\ \quad \operatorname{is\text{-}ring}(a) \Rightarrow (\operatorname{mul}(\operatorname{poly}(a))(f, g) = \operatorname{monalg\text{-}mul}(a, \operatorname{nn\text{-}add\text{-}monoid}, f, g)) \end{array}}}
\vnbprop{poly-mul-deg-bound}{\vnbstmt{\begin{array}{@{}l@{}} \forall a.; \forall f, g \in \operatorname{carr}(\operatorname{poly}(a)), n, d \in \mathbb{N}. \\ \quad (\operatorname{is\text{-}ring}(a) \wedge (\operatorname{deg\text{-}bound}(a, f, n) \wedge \operatorname{deg\text{-}bound}(a, g, d))) \Rightarrow  \\ \quad \quad \operatorname{deg\text{-}bound}(a, \operatorname{monalg\text{-}mul}(a, \operatorname{nn\text{-}add\text{-}monoid}, f, g), (n + d)) \end{array}}}
\vnbprop{poly-mul-top-coef}{\vnbstmt{\begin{array}{@{}l@{}} \forall a.; \forall f, g \in \operatorname{carr}(\operatorname{poly}(a)), n, d \in \mathbb{N}. \\ \quad (\operatorname{is\text{-}ring}(a) \wedge (\operatorname{deg\text{-}bound}(a, f, n) \wedge \operatorname{deg\text{-}bound}(a, g, d))) \Rightarrow  \\ \quad \quad \operatorname{monalg\text{-}mul}(a, \operatorname{nn\text{-}add\text{-}monoid}, f, g)((n + d)) = \operatorname{mul}(a)(f(n), g(d)) \end{array}}}
\vnbprop{poly-mul-deg}{\vnbstmt{\begin{array}{@{}l@{}} \forall a.; \forall f, g \in \operatorname{carr}(\operatorname{poly}(a)). \\ \quad (\operatorname{is\text{-}integral\text{-}domain}(a) \wedge (\neg (f = \operatorname{zero}(\operatorname{poly}(a))) \wedge \neg (g = \operatorname{zero}(\operatorname{poly}(a))))) \Rightarrow  \\ \quad \quad (\neg (\operatorname{mul}(\operatorname{poly}(a))(f, g) = \operatorname{zero}(\operatorname{poly}(a))) \wedge (\operatorname{deg}(a, \operatorname{mul}(\operatorname{poly}(a))(f, g)) = (\operatorname{deg}(a, f) + \operatorname{deg}(a, g)))) \end{array}}}
\vnbprop{poly-is-integral-domain}{\vnbstmt{\forall a.\; (\operatorname{is\text{-}integral\text{-}domain}(a) \Rightarrow \operatorname{is\text{-}integral\text{-}domain}(\operatorname{poly}(a)))}}
\vnbprop{monomial-deg-bound}{\vnbstmt{\begin{array}{@{}l@{}} \forall a.; \forall c \in \operatorname{carr}(a), e \in \mathbb{N}. \\ \quad \operatorname{is\text{-}ring}(a) \Rightarrow \operatorname{deg\text{-}bound}(a, \operatorname{monomial}(a, c, e), e) \end{array}}}
\vnbprop{deg-drop}{\vnbstmt{\begin{array}{@{}l@{}} \forall a.; \forall h \in \operatorname{carr}(\operatorname{poly}(a)), n \in \mathbb{N}. \\ \quad (\operatorname{is\text{-}ring}(a) \wedge (\operatorname{deg\text{-}bound}(a, h, n) \wedge (h(n) = \operatorname{zero}(a)))) \Rightarrow  \\ \quad \quad (h = \operatorname{zero}(\operatorname{poly}(a))) \vee (\operatorname{deg}(a, h) < n) \end{array}}}
\vnbprop{poly-division-bounded}{\vnbstmt{\begin{array}{@{}l@{}} \forall b \in \mathbb{N}, k.; \forall a, \mathit{pb} \in \operatorname{carr}(\operatorname{poly}(k)). \\ \quad (\operatorname{is\text{-}field\text{-}ring}(k) \wedge ((\operatorname{deg}(k, a) < b) \wedge \neg (\mathit{pb} = \operatorname{zero}(\operatorname{poly}(k))))) \Rightarrow  \\ \quad \quad \exists s \in \operatorname{carr}(\operatorname{poly}(k)). \\ \quad \quad \quad a = \operatorname{mul}(\operatorname{poly}(k))(\mathit{pb}, s) \vee \\ \quad \quad \quad \operatorname{deg}(k, \operatorname{add}(\operatorname{poly}(k))(a, \operatorname{neg}(\operatorname{poly}(k))(\operatorname{mul}(\operatorname{poly}(k))(\mathit{pb}, s)))) < \operatorname{deg}(k, \mathit{pb}) \end{array}}}
\vnbprop{poly-division-with-remainder}{\vnbstmt{\begin{array}{@{}l@{}} \forall k.; \forall a, b \in \operatorname{carr}(\operatorname{poly}(k)). \\ \quad (\operatorname{is\text{-}field\text{-}ring}(k) \wedge \neg (b = \operatorname{zero}(\operatorname{poly}(k)))) \Rightarrow  \\ \quad \quad \exists s, r \in \operatorname{carr}(\operatorname{poly}(k)). \\ \quad \quad \quad ((a = \operatorname{add}(\operatorname{poly}(k))(\operatorname{mul}(\operatorname{poly}(k))(b, s), r)) \wedge ((r = \operatorname{zero}(\operatorname{poly}(k))) \vee (\operatorname{deg}(k, r) < \operatorname{deg}(k, b)))) \end{array}}}
\vnbprop{poly-degree-is-degree-function}{\vnbstmt{\begin{array}{@{}l@{}} \forall k. \\ \quad \operatorname{is\text{-}field\text{-}ring}(k) \Rightarrow  \\ \quad \quad \operatorname{is\text{-}degree\text{-}function}(\operatorname{poly}(k), (\lambda z \in \operatorname{carr}(\operatorname{poly}(k)).\; \operatorname{deg}(k, z))) \end{array}}}
\vnbprop{euclidean-domain-is-pid}{\vnbstmt{\forall s.\; (\operatorname{is\text{-}euclidean\text{-}domain}(s) \Rightarrow \operatorname{is\text{-}pid}(s))}}
\vnbprop{poly-is-euclidean-domain}{\vnbstmt{\forall k.\; (\operatorname{is\text{-}field\text{-}ring}(k) \Rightarrow \operatorname{is\text{-}euclidean\text{-}domain}(\operatorname{poly}(k)))}}
\vnbprop{poly-is-pid-and-ufd}{\vnbstmt{\forall k.\; (\operatorname{is\text{-}field\text{-}ring}(k) \Rightarrow (\operatorname{is\text{-}pid}(\operatorname{poly}(k)) \wedge \operatorname{is\text{-}ufd}(\operatorname{poly}(k))))}}
\vnbprop{poly-mul-single-coef}{\vnbstmt{\begin{array}{@{}l@{}} \forall a.; \forall f, g \in \operatorname{carr}(\operatorname{poly}(a)), n, d \in \mathbb{N}. \\ \quad (\operatorname{is\text{-}ring}(a) \wedge \forall i \in \mathbb{N}.\; (\neg (i = n) \Rightarrow (f(i) = \operatorname{zero}(a)))) \Rightarrow  \\ \quad \quad \operatorname{monalg\text{-}mul}(a, \operatorname{nn\text{-}add\text{-}monoid}, f, g)((n + d)) = \operatorname{mul}(a)(f(n), g(d)) \end{array}}}
\vnbprop{poly-eval-sum-stable}{\vnbstmt{\begin{array}{@{}l@{}} \forall j \in \mathbb{N}, a.; \forall p \in \operatorname{carr}(\operatorname{poly}(a)), t \in \operatorname{carr}(a). \\ \quad \operatorname{is\text{-}commutative\text{-}ring}(a) \Rightarrow  \\ \quad \quad \operatorname{sum\text{-}ag}(\operatorname{ring\text{-}additive\text{-}ag}(a), \\ \quad \quad \quad (\lambda k \in \mathbb{N}.\; \operatorname{mul}(a)(p(k), \operatorname{ring\text{-}power}(a, t, k))), \\ \quad \quad \quad ((\operatorname{deg}(a, p) + j) + 1)) \\ \quad \quad =\; \operatorname{sum\text{-}ag}(\operatorname{ring\text{-}additive\text{-}ag}(a), (\lambda k \in \mathbb{N}.\; \operatorname{mul}(a)(p(k), \operatorname{ring\text{-}power}(a, t, k))), (\operatorname{deg}(a, p) + 1)) \end{array}}}
\vnbprop{poly-eval-bound}{\vnbstmt{\begin{array}{@{}l@{}} \forall a.; \forall p \in \operatorname{carr}(\operatorname{poly}(a)), t \in \operatorname{carr}(a), n \in \mathbb{N}. \\ \quad (\operatorname{is\text{-}commutative\text{-}ring}(a) \wedge \operatorname{deg\text{-}bound}(a, p, n)) \Rightarrow  \\ \quad \quad \operatorname{poly\text{-}eval}(a, p, t) \\ \quad \quad =\; \operatorname{sum\text{-}ag}(\operatorname{ring\text{-}additive\text{-}ag}(a), (\lambda k \in \mathbb{N}.\; \operatorname{mul}(a)(p(k), \operatorname{ring\text{-}power}(a, t, k))), (n + 1)) \end{array}}}
\vnbprop{poly-eval-add}{\vnbstmt{\begin{array}{@{}l@{}} \forall a.; \forall p, q \in \operatorname{carr}(\operatorname{poly}(a)), t \in \operatorname{carr}(a). \\ \quad \operatorname{is\text{-}commutative\text{-}ring}(a) \Rightarrow  \\ \quad \quad \operatorname{poly\text{-}eval}(a, \operatorname{add}(\operatorname{poly}(a))(p, q), t) = \operatorname{add}(a)(\operatorname{poly\text{-}eval}(a, p, t), \operatorname{poly\text{-}eval}(a, q, t)) \end{array}}}
\vnbprop{poly-mul-single-low}{\vnbstmt{\begin{array}{@{}l@{}} \forall a.; \forall f, g \in \operatorname{carr}(\operatorname{poly}(a)), n, k \in \mathbb{N}. \\ \quad (\operatorname{is\text{-}ring}(a) \wedge (\forall i \in \mathbb{N}.\; (\neg (i = n) \Rightarrow (f(i) = \operatorname{zero}(a))) \wedge (k < n))) \Rightarrow  \\ \quad \quad \operatorname{monalg\text{-}mul}(a, \operatorname{nn\text{-}add\text{-}monoid}, f, g)(k) = \operatorname{zero}(a) \end{array}}}
\vnbprop{poly-eval-monomial-mul-partial}{\vnbstmt{\begin{array}{@{}l@{}} \forall j \in \mathbb{N}, a.; \forall c \in \operatorname{carr}(a), e \in \mathbb{N}, s \in \operatorname{carr}(\operatorname{poly}(a)), t \in \operatorname{carr}(a). \\ \quad \operatorname{is\text{-}commutative\text{-}ring}(a) \Rightarrow  \\ \quad \quad \operatorname{sum\text{-}ag}(\operatorname{ring\text{-}additive\text{-}ag}(a), \\ \quad \quad \quad (\lambda k \in \mathbb{N}.\; \\ \quad \quad \quad \quad \operatorname{mul}(a)(\operatorname{monalg\text{-}mul}(a, \operatorname{nn\text{-}add\text{-}monoid}, \operatorname{monomial}(a, c, e), s)(k), \operatorname{ring\text{-}power}(a, t, k))), \\ \quad \quad \quad (e + j)) \\ \quad \quad =\; \operatorname{mul}(a)(\operatorname{mul}(a)(c, \operatorname{ring\text{-}power}(a, t, e)), \\ \quad \quad \quad \operatorname{sum\text{-}ag}(\operatorname{ring\text{-}additive\text{-}ag}(a), (\lambda k \in \mathbb{N}.\; \operatorname{mul}(a)(s(k), \operatorname{ring\text{-}power}(a, t, k))), j)) \end{array}}}
\vnbprop{poly-eval-monomial-mul}{\vnbstmt{\begin{array}{@{}l@{}} \forall a.; \forall c \in \operatorname{carr}(a), e \in \mathbb{N}, s \in \operatorname{carr}(\operatorname{poly}(a)), t \in \operatorname{carr}(a). \\ \quad \operatorname{is\text{-}commutative\text{-}ring}(a) \Rightarrow  \\ \quad \quad \operatorname{poly\text{-}eval}(a, \operatorname{mul}(\operatorname{poly}(a))(\operatorname{monomial}(a, c, e), s), t) \\ \quad \quad =\; \operatorname{mul}(a)(\operatorname{mul}(a)(c, \operatorname{ring\text{-}power}(a, t, e)), \operatorname{poly\text{-}eval}(a, s, t)) \end{array}}}
\vnbprop{poly-eval-x-minus-mul}{\vnbstmt{\begin{array}{@{}l@{}} \forall a.; \forall l \in \operatorname{carr}(a), s \in \operatorname{carr}(\operatorname{poly}(a)), t \in \operatorname{carr}(a). \\ \quad \operatorname{is\text{-}commutative\text{-}ring}(a) \Rightarrow  \\ \quad \quad \operatorname{poly\text{-}eval}(a, \operatorname{mul}(\operatorname{poly}(a))(\operatorname{x\text{-}minus}(a, l), s), t) \\ \quad \quad =\; \operatorname{mul}(a)(\operatorname{add}(a)(t, \operatorname{neg}(a)(l)), \operatorname{poly\text{-}eval}(a, s, t)) \end{array}}}
\vnbprop{x-minus-degree-one}{\vnbstmt{\begin{array}{@{}l@{}} \forall a.; \forall l \in \operatorname{carr}(a). \\ \quad (\operatorname{is\text{-}commutative\text{-}ring}(a) \wedge \neg (\operatorname{one}(a) = \operatorname{zero}(a))) \Rightarrow  \\ \quad \quad ((\operatorname{x\text{-}minus}(a, l) \in \operatorname{carr}(\operatorname{poly}(a))) \wedge (\neg (\operatorname{x\text{-}minus}(a, l) = \operatorname{zero}(\operatorname{poly}(a))) \wedge (\operatorname{deg}(a, \operatorname{x\text{-}minus}(a, l)) = 1))) \end{array}}}
\vnbprop{poly-eval-constant}{\vnbstmt{\begin{array}{@{}l@{}} \forall a.; \forall p \in \operatorname{carr}(\operatorname{poly}(a)), t \in \operatorname{carr}(a). \\ \quad (\operatorname{is\text{-}commutative\text{-}ring}(a) \wedge \operatorname{deg\text{-}bound}(a, p, 0)) \Rightarrow (\operatorname{poly\text{-}eval}(a, p, t) = p(0)) \end{array}}}
\vnbprop{poly-eval-type}{\vnbstmt{\begin{array}{@{}l@{}} \forall a.; \forall p \in \operatorname{carr}(\operatorname{poly}(a)), t \in \operatorname{carr}(a). \\ \quad \operatorname{is\text{-}commutative\text{-}ring}(a) \Rightarrow (\operatorname{poly\text{-}eval}(a, p, t) \in \operatorname{carr}(a)) \end{array}}}
\vnbprop{poly-root-iff-x-minus-divides}{\vnbstmt{\begin{array}{@{}l@{}} \forall k.; \forall p \in \operatorname{carr}(\operatorname{poly}(k)), l \in \operatorname{carr}(k). \\ \quad \operatorname{is\text{-}field\text{-}ring}(k) \Rightarrow  \\ \quad \quad (\operatorname{poly\text{-}eval}(k, p, l) = \operatorname{zero}(k)) \Leftrightarrow \operatorname{divides}(\operatorname{poly}(k), \operatorname{x\text{-}minus}(k, l), p) \end{array}}}
\vnbprop{poly-roots-at-most-degree}{\vnbstmt{\begin{array}{@{}l@{}} \forall m \in \mathbb{N}, k.; \forall p \in \operatorname{carr}(\operatorname{poly}(k)), l. \\ \quad \operatorname{is\text{-}field\text{-}ring}(k) \wedge \\ \quad \neg (p = \operatorname{zero}(\operatorname{poly}(k))) \wedge \\ \quad \forall i \in \operatorname{interval}(1, m). \\ \quad \quad (l(i) \in \operatorname{carr}(k)) \wedge (\operatorname{poly\text{-}eval}(k, p, l(i)) = \operatorname{zero}(k)) \wedge \\ \quad \forall i, j \in \operatorname{interval}(1, m).\; ((l(i) = l(j)) \Rightarrow (i = j)) \Rightarrow \\ \quad \quad m \leq \operatorname{deg}(k, p) \end{array}}}
\vnbprop{poly-vanishing-iff-zero}{\vnbstmt{\begin{array}{@{}l@{}} \forall k.; \forall p \in \operatorname{carr}(\operatorname{poly}(k)), l \in \operatorname{injection}(\operatorname{interval}(1, (\operatorname{deg}(k, p) + 1)), \operatorname{carr}(k)). \\ \quad \operatorname{is\text{-}field\text{-}ring}(k) \Rightarrow  \\ \quad \quad \forall a \in \operatorname{carr}(k).\; (\operatorname{poly\text{-}eval}(k, p, a) = \operatorname{zero}(k)) \Leftrightarrow \forall b \in \mathbb{N}.\; (p(b) = \operatorname{zero}(k)) \end{array}}}
\vnbprop{infinite-set-finite-injections}{\vnbstmt{\begin{array}{@{}l@{}} \forall n \in \mathbb{N}, s \in \mathbf{Set}. \\ \quad \neg (|s| \in \mathbb{N}) \Rightarrow \exists \mathit{gin}.\; (\mathit{gin} \in \operatorname{injection}(\operatorname{interval}(1, n), s)) \end{array}}}
\vnbprop{poly-eval-zero}{\vnbstmt{\begin{array}{@{}l@{}} \forall a.; \forall t \in \operatorname{carr}(a). \\ \quad \operatorname{is\text{-}commutative\text{-}ring}(a) \Rightarrow (\operatorname{poly\text{-}eval}(a, \operatorname{zero}(\operatorname{poly}(a)), t) = \operatorname{zero}(a)) \end{array}}}
\vnbprop{poly-eval-neg}{\vnbstmt{\begin{array}{@{}l@{}} \forall a.; \forall q \in \operatorname{carr}(\operatorname{poly}(a)), t \in \operatorname{carr}(a). \\ \quad \operatorname{is\text{-}commutative\text{-}ring}(a) \Rightarrow  \\ \quad \quad \operatorname{poly\text{-}eval}(a, \operatorname{neg}(\operatorname{poly}(a))(q), t) = \operatorname{neg}(a)(\operatorname{poly\text{-}eval}(a, q, t)) \end{array}}}
\vnbprop{poly-determined-by-values}{\vnbstmt{\begin{array}{@{}l@{}} \forall k.; \forall p, q \in \operatorname{carr}(\operatorname{poly}(k)). \\ \quad (\operatorname{is\text{-}field\text{-}ring}(k) \wedge \neg (|\operatorname{carr}(k)| \in \mathbb{N})) \Rightarrow  \\ \quad \quad (p = q) \Leftrightarrow \forall a \in \operatorname{carr}(k).\; (\operatorname{poly\text{-}eval}(k, p, a) = \operatorname{poly\text{-}eval}(k, q, a)) \end{array}}}

\vnbbookfile{theorem-library/gal-interp}{theorem-library/gal-interp}
\noindent{\small gal-interp.scm -- Lagrange interpolation (galois-theory notes, Proposition 1.47, Corollary 1.49) and the root bound for x\textasciicircum{}d = 1 (Proposition 1.64, first sentence) (GALOIS, 2026-10-26).  Destination: theorem-library/gal-interp.scm, library algebra. Load window: after gal-poly (and so after gal-ufd and the galois ...}\par
\vnbprop{poly-add-deg-bound}{\vnbstmt{\begin{array}{@{}l@{}} \forall a.; \forall f, g \in \operatorname{carr}(\operatorname{poly}(a)), n \in \mathbb{N}. \\ \quad (\operatorname{is\text{-}commutative\text{-}ring}(a) \wedge (\operatorname{deg\text{-}bound}(a, f, n) \wedge \operatorname{deg\text{-}bound}(a, g, n))) \Rightarrow  \\ \quad \quad \operatorname{deg\text{-}bound}(a, \operatorname{add}(\operatorname{poly}(a))(f, g), n) \end{array}}}
\vnbprop{poly-neg-deg-bound}{\vnbstmt{\begin{array}{@{}l@{}} \forall a.; \forall f \in \operatorname{carr}(\operatorname{poly}(a)), n \in \mathbb{N}. \\ \quad (\operatorname{is\text{-}commutative\text{-}ring}(a) \wedge \operatorname{deg\text{-}bound}(a, f, n)) \Rightarrow  \\ \quad \quad \operatorname{deg\text{-}bound}(a, \operatorname{neg}(\operatorname{poly}(a))(f), n) \end{array}}}
\vnbprop{poly-vanishing-at-points}{\vnbstmt{\begin{array}{@{}l@{}} \forall n \in \mathbb{N}, k, h. \\ \quad (\operatorname{is\text{-}field\text{-}ring}(k) \wedge \forall i \in \operatorname{interval}(1, n).\; (h(i) \in \operatorname{carr}(k))) \Rightarrow  \\ \quad \quad \exists \mathit{gnn} \in \operatorname{carr}(\operatorname{poly}(k)). \\ \quad \quad \quad \operatorname{deg\text{-}bound}(k, \mathit{gnn}, n) \wedge \\ \quad \quad \quad \forall i \in \operatorname{interval}(1, n).\; (\operatorname{poly\text{-}eval}(k, \mathit{gnn}, h(i)) = \operatorname{zero}(k)) \wedge \\ \quad \quad \quad \forall t \in \operatorname{carr}(k). \\ \quad \quad \quad \quad \forall i \in \operatorname{interval}(1, n).\; \neg (t = h(i)) \Rightarrow  \\ \quad \quad \quad \quad \quad \neg (\operatorname{poly\text{-}eval}(k, \mathit{gnn}, t) = \operatorname{zero}(k)) \end{array}}}
\vnbprop{poly-interpolation-exists}{\vnbstmt{\begin{array}{@{}l@{}} \forall m \in \mathbb{N}, k, h, v. \\ \quad \operatorname{is\text{-}field\text{-}ring}(k) \wedge \\ \quad \forall i \in \operatorname{interval}(1, (m + 1)).\; (h(i) \in \operatorname{carr}(k)) \wedge \\ \quad \forall i, j \in \operatorname{interval}(1, (m + 1)). \\ \quad \quad (h(i) = h(j)) \Rightarrow (i = j) \wedge \\ \quad \forall i \in \operatorname{interval}(1, (m + 1)).\; (v(i) \in \operatorname{carr}(k)) \Rightarrow \\ \quad \quad \exists p \in \operatorname{carr}(\operatorname{poly}(k)). \\ \quad \quad \quad (\operatorname{deg\text{-}bound}(k, p, m) \wedge \forall i \in \operatorname{interval}(1, (m + 1)).\; (\operatorname{poly\text{-}eval}(k, p, h(i)) = v(i))) \end{array}}}
\vnbprop{poly-interpolation-unique}{\vnbstmt{\begin{array}{@{}l@{}} \forall m \in \mathbb{N}, k, h.; \forall p, q \in \operatorname{carr}(\operatorname{poly}(k)). \\ \quad \operatorname{is\text{-}field\text{-}ring}(k) \wedge \\ \quad \forall i \in \operatorname{interval}(1, (m + 1)).\; (h(i) \in \operatorname{carr}(k)) \wedge \\ \quad \forall i, j \in \operatorname{interval}(1, (m + 1)). \\ \quad \quad (h(i) = h(j)) \Rightarrow (i = j) \wedge \\ \quad \operatorname{deg\text{-}bound}(k, p, m) \wedge \\ \quad \operatorname{deg\text{-}bound}(k, q, m) \wedge \\ \quad \forall i \in \operatorname{interval}(1, (m + 1)). \\ \quad \quad \operatorname{poly\text{-}eval}(k, p, h(i)) = \operatorname{poly\text{-}eval}(k, q, h(i)) \Rightarrow \\ \quad \quad p = q \end{array}}}
\vnbprop{poly-le-unfold}{\vnbstmt{\forall a, m.\; (\operatorname{poly\text{-}le}(a, m) \simeq \{l \in \operatorname{carr}(\operatorname{poly}(a)) : \operatorname{deg\text{-}bound}(a, l, m)\})}}
\vnbprop{eval-at-points-bijective-iff-distinct}{\vnbstmt{\begin{array}{@{}l@{}} \forall k.; \forall m \in \mathbb{N}, h. \\ \quad (\operatorname{is\text{-}field\text{-}ring}(k) \wedge \forall i \in \operatorname{interval}(1, (m + 1)).\; (h(i) \in \operatorname{carr}(k))) \Rightarrow  \\ \quad \quad \operatorname{eval\text{-}at\text{-}points}(k, h, m) \in \mathrm{Bij}(\operatorname{poly\text{-}le}(k, m), (\operatorname{interval}(1, (m + 1)) \to \operatorname{carr}(k))) \\ \quad \quad \Leftrightarrow\; \forall i, j \in \operatorname{interval}(1, (m + 1)). \\ \quad \quad \quad (h(i) = h(j)) \Rightarrow (i = j) \end{array}}}
\vnbprop{segment-interval-bijection}{\vnbstmt{\begin{array}{@{}l@{}} \forall n \in \mathbb{N}. \\ \quad (\lambda k \in \mathrm{OS}(n).\; (k + 1)) \in \mathrm{Bij}(\mathrm{OS}(n), \operatorname{interval}(1, n)) \end{array}}}
\vnbprop{finite-set-enumeration}{\vnbstmt{\begin{array}{@{}l@{}} \forall s \in \mathbf{Set}, n \in \mathbb{N}. \\ \quad (|s| = n) \Rightarrow \exists h.\; (h \in \mathrm{Bij}(\operatorname{interval}(1, n), s)) \end{array}}}
\vnbprop{finite-field-maps-are-polynomials}{\vnbstmt{\begin{array}{@{}l@{}} \forall k.; \forall m \in \mathbb{N}. \\ \quad (\operatorname{is\text{-}field\text{-}ring}(k) \wedge (|\operatorname{carr}(k)| = (m + 1))) \Rightarrow  \\ \quad \quad (\lambda p \in \operatorname{poly\text{-}le}(k, m).\; (\lambda t \in \operatorname{carr}(k).\; \operatorname{poly\text{-}eval}(k, p, t))) \\ \quad \quad \in\; \mathrm{Bij}(\operatorname{poly\text{-}le}(k, m), (\operatorname{carr}(k) \to \operatorname{carr}(k))) \end{array}}}
\vnbprop{poly-one-is-monomial}{\vnbstmt{\forall a.\; (\operatorname{is\text{-}ring}(a) \Rightarrow (\operatorname{one}(\operatorname{poly}(a)) = \operatorname{monomial}(a, \operatorname{one}(a), 0)))}}
\vnbprop{poly-eval-monomial}{\vnbstmt{\begin{array}{@{}l@{}} \forall a.; \forall c \in \operatorname{carr}(a), e \in \mathbb{N}, t \in \operatorname{carr}(a). \\ \quad \operatorname{is\text{-}commutative\text{-}ring}(a) \Rightarrow  \\ \quad \quad \operatorname{poly\text{-}eval}(a, \operatorname{monomial}(a, c, e), t) = \operatorname{mul}(a)(c, \operatorname{ring\text{-}power}(a, t, e)) \end{array}}}
\vnbprop{unity-roots-at-most-d}{\vnbstmt{\begin{array}{@{}l@{}} \forall k.; \forall d, m \in \mathbb{N}, l. \\ \quad \operatorname{is\text{-}field\text{-}ring}(k) \wedge \\ \quad 1 \leq d \wedge \\ \quad \forall i \in \operatorname{interval}(1, m). \\ \quad \quad (l(i) \in \operatorname{carr}(k)) \wedge (\operatorname{ring\text{-}power}(k, l(i), d) = \operatorname{one}(k)) \wedge \\ \quad \forall i, j \in \operatorname{interval}(1, m).\; ((l(i) = l(j)) \Rightarrow (i = j)) \Rightarrow \\ \quad \quad m \leq d \end{array}}}

\vnbbookfile{structure-library/gal-gcd-vocab}{structure-library/gal-gcd-vocab}
\noindent{\small gal-gcd-vocab.scm -- a greatest common divisor of a set of ring elements (galois-theory notes, section 4, "gcd sdet A") (GALOIS, 2026-10-26).  DEFINITIONS ONLY.  Destination: structure-library/gal-gcd-vocab.scm, library algebra.  Load window: after catools-gal-vocab (DIVIDES). IS-GCD-SET(s, S, d): d divides every ...}\par
\par\medskip\noindent\textbf{Definition} (\texttt{is-gcd-set}).\;For all $\operatorname{ggs\_}$, $\operatorname{ggt\_}$ and $\operatorname{ggd\_}$, 
\begin{equation*}
\begin{array}{@{}l@{}} \operatorname{is\text{-}gcd\text{-}set}(\operatorname{ggs\_}, \operatorname{ggt\_}, \operatorname{ggd\_}) \\ \Leftrightarrow\; \operatorname{ggd\_} \in \operatorname{carr}(\operatorname{ggs\_}) \wedge \\ \forall \operatorname{ggx\_} \in \operatorname{ggt\_}.\; \operatorname{divides}(\operatorname{ggs\_}, \operatorname{ggd\_}, \operatorname{ggx\_}) \wedge \\ \forall \operatorname{gge\_} \in \operatorname{carr}(\operatorname{ggs\_}). \\ \quad \forall \operatorname{ggx\_} \in \operatorname{ggt\_}.\; \operatorname{divides}(\operatorname{ggs\_}, \operatorname{gge\_}, \operatorname{ggx\_}) \Rightarrow  \\ \quad \quad \operatorname{divides}(\operatorname{ggs\_}, \operatorname{gge\_}, \operatorname{ggd\_}) \end{array}
\end{equation*}


\vnbbookfile{theorem-library/rake-asm-2}{theorem-library/rake-asm-2}
\noindent{\small rake-asm-2.scm -- RAKE-ASM agent RA-2 (2026-10-12): elementary laws that stood as ASSUMPTIONS (add-axiom! / support), each proven here under its EXACT name and statement (copied character for character from its definition site).}\par
\vnbprop{prod-ord-singleton}{\vnbstmt{\begin{array}{@{}l@{}} \forall m.; \forall f \in (\mathbb{N} \to \operatorname{carr}(m)). \\ \quad \operatorname{is\text{-}monoid}(m) \Rightarrow (\operatorname{prod\text{-}ord}(m, f, 1) = f(0)) \end{array}}}
\vnbprop{sum-singleton}{\vnbstmt{\forall r.; \forall f \in (\mathbb{N} \to \operatorname{carr}(r)).\; (\operatorname{is\text{-}ring}(r) \Rightarrow (\operatorname{sum}(r, f, 1) = f(0)))}}
\vnbprop{sum-ag-singleton}{\vnbstmt{\begin{array}{@{}l@{}} \forall g.; \forall f \in (\mathbb{N} \to \operatorname{carr}(g)). \\ \quad \operatorname{is\text{-}abelian\text{-}group}(g) \Rightarrow (\operatorname{sum\text{-}ag}(g, f, 1) = f(0)) \end{array}}}
\vnbprop{prod-ord-type-ind}{\vnbstmt{\begin{array}{@{}l@{}} \forall n \in \mathbb{N}, m.; \forall f \in (\mathbb{N} \to \operatorname{carr}(m)). \\ \quad \operatorname{is\text{-}monoid}(m) \Rightarrow (\operatorname{prod\text{-}ord}(m, f, n) \in \operatorname{carr}(m)) \end{array}}}
\vnbprop{prod-ord-type}{\vnbstmt{\begin{array}{@{}l@{}} \forall m.; \forall f \in (\mathbb{N} \to \operatorname{carr}(m)), n \in \mathbb{N}. \\ \quad \operatorname{is\text{-}monoid}(m) \Rightarrow (\operatorname{prod\text{-}ord}(m, f, n) \in \operatorname{carr}(m)) \end{array}}}
\vnbprop{sum-type-ind}{\vnbstmt{\begin{array}{@{}l@{}} \forall n \in \mathbb{N}, r.; \forall f \in (\mathbb{N} \to \operatorname{carr}(r)). \\ \quad \operatorname{is\text{-}ring}(r) \Rightarrow (\operatorname{sum}(r, f, n) \in \operatorname{carr}(r)) \end{array}}}
\vnbprop{sum-type}{\vnbstmt{\begin{array}{@{}l@{}} \forall r.; \forall f \in (\mathbb{N} \to \operatorname{carr}(r)), n \in \mathbb{N}. \\ \quad \operatorname{is\text{-}ring}(r) \Rightarrow (\operatorname{sum}(r, f, n) \in \operatorname{carr}(r)) \end{array}}}
\vnbprop{sum-ag-as-reduce-ind}{\vnbstmt{\begin{array}{@{}l@{}} \forall k \in \mathbb{N}, g.; \forall f \in (\mathbb{N} \to \operatorname{carr}(g)). \\ \quad \operatorname{is\text{-}abelian\text{-}group}(g) \Rightarrow (\operatorname{sum\text{-}ag}(g, f, (k + 1)) = \operatorname{reduce}(\operatorname{opr}(g), f, (k + 1))) \end{array}}}
\vnbprop{sum-ag-as-reduce}{\vnbstmt{\begin{array}{@{}l@{}} \forall g.; \forall f \in (\mathbb{N} \to \operatorname{carr}(g)), n \in \mathbb{N}. \\ \quad (\operatorname{is\text{-}abelian\text{-}group}(g) \wedge (1 \leq n)) \Rightarrow (\operatorname{sum\text{-}ag}(g, f, n) = \operatorname{reduce}(\operatorname{opr}(g), f, n)) \end{array}}}
\vnbprop{mpow-mult-ind}{\vnbstmt{\begin{array}{@{}l@{}} \forall n \in \mathbb{N}, m.; \forall x, y \in \operatorname{carr}(m). \\ \quad \operatorname{is\text{-}comm\text{-}monoid}(m) \Rightarrow  \\ \quad \quad \operatorname{mpow}(m, \operatorname{opr}(m)(x, y), n) = \operatorname{opr}(m)(\operatorname{mpow}(m, x, n), \operatorname{mpow}(m, y, n)) \end{array}}}
\vnbprop{mpow-mult}{\vnbstmt{\begin{array}{@{}l@{}} \forall m.; \forall x, y \in \operatorname{carr}(m), n \in \mathbb{N}. \\ \quad \operatorname{is\text{-}comm\text{-}monoid}(m) \Rightarrow  \\ \quad \quad \operatorname{mpow}(m, \operatorname{opr}(m)(x, y), n) = \operatorname{opr}(m)(\operatorname{mpow}(m, x, n), \operatorname{mpow}(m, y, n)) \end{array}}}
\vnbprop{sum-left-scalar-ind}{\vnbstmt{\begin{array}{@{}l@{}} \forall n \in \mathbb{N}, r.; \forall a \in \operatorname{carr}(r), f \in (\mathbb{N} \to \operatorname{carr}(r)). \\ \quad \operatorname{is\text{-}ring}(r) \Rightarrow  \\ \quad \quad \operatorname{sum}(r, (\lambda i \in \mathbb{N}.\; \operatorname{mul}(r)(a, f(i))), n) = \operatorname{mul}(r)(a, \operatorname{sum}(r, f, n)) \end{array}}}
\vnbprop{sum-left-scalar}{\vnbstmt{\begin{array}{@{}l@{}} \forall r.; \forall a \in \operatorname{carr}(r), f \in (\mathbb{N} \to \operatorname{carr}(r)), n \in \mathbb{N}. \\ \quad \operatorname{is\text{-}ring}(r) \Rightarrow  \\ \quad \quad \operatorname{sum}(r, (\lambda i \in \mathbb{N}.\; \operatorname{mul}(r)(a, f(i))), n) = \operatorname{mul}(r)(a, \operatorname{sum}(r, f, n)) \end{array}}}
\vnbprop{zz-act-distrib-nn}{\vnbstmt{\begin{array}{@{}l@{}} \forall k \in \mathbb{N}, g.; \forall a, b \in \operatorname{carr}(g). \\ \quad \operatorname{is\text{-}abelian\text{-}group}(g) \Rightarrow  \\ \quad \quad \operatorname{zz\text{-}act}(g, k, \operatorname{opr}(g)(a, b)) = \operatorname{opr}(g)(\operatorname{zz\text{-}act}(g, k, a), \operatorname{zz\text{-}act}(g, k, b)) \end{array}}}
\vnbprop{zz-act-distrib}{\vnbstmt{\begin{array}{@{}l@{}} \forall g.; \forall k \in \mathbb{Z}, a, b \in \operatorname{carr}(g). \\ \quad \operatorname{is\text{-}abelian\text{-}group}(g) \Rightarrow  \\ \quad \quad \operatorname{zz\text{-}act}(g, k, \operatorname{opr}(g)(a, b)) = \operatorname{opr}(g)(\operatorname{zz\text{-}act}(g, k, a), \operatorname{zz\text{-}act}(g, k, b)) \end{array}}}
\vnbprop{zz-act-assoc-nn}{\vnbstmt{\begin{array}{@{}l@{}} \forall j \in \mathbb{N}, g.; \forall k \in \mathbb{Z}, a \in \operatorname{carr}(g). \\ \quad \operatorname{is\text{-}abelian\text{-}group}(g) \Rightarrow (\operatorname{zz\text{-}act}(g, (j \cdot k), a) = \operatorname{zz\text{-}act}(g, j, \operatorname{zz\text{-}act}(g, k, a))) \end{array}}}
\vnbprop{zz-act-assoc}{\vnbstmt{\begin{array}{@{}l@{}} \forall g.; \forall j, k \in \mathbb{Z}, a \in \operatorname{carr}(g). \\ \quad \operatorname{is\text{-}abelian\text{-}group}(g) \Rightarrow (\operatorname{zz\text{-}act}(g, (j \cdot k), a) = \operatorname{zz\text{-}act}(g, j, \operatorname{zz\text{-}act}(g, k, a))) \end{array}}}

\vnbbookfile{theorem-library/gal-char}{theorem-library/gal-char}
\noindent{\small gal-char.scm -- primes in Z dividing binomial coefficients, the characteristic of a field and the Frobenius map (galois-theory notes, Lemma 1.12, Proposition 1.14) (GALOIS B8, 2026-10-26). Destination: theorem-library/gal-char.scm, library algebra.  Load window: after catools-gal-b1 (two-element-ideal), the binomial ...}\par
\vnbprop{euclid-lemma-euclidean-ring}{\vnbstmt{\begin{array}{@{}l@{}} \forall s, p.; \forall a, b \in \operatorname{carr}(s). \\ \quad (\operatorname{is\text{-}euclidean\text{-}ring}(s) \wedge (\operatorname{is\text{-}prime\text{-}element}(s, p) \wedge \operatorname{divides}(s, p, \operatorname{mul}(s)(a, b)))) \Rightarrow  \\ \quad \quad \operatorname{divides}(s, p, a) \vee \operatorname{divides}(s, p, b) \end{array}}}
\vnbprop{zz-ring-divides-iff}{\vnbstmt{\forall p, x \in \mathbb{Z}.\; (\operatorname{divides}(\operatorname{zz\text{-}ring}, p, x) \Leftrightarrow \operatorname{zz\text{-}divides}(p, x))}}
\vnbprop{prime-nat-gt-one}{\vnbstmt{\forall p.\; (\operatorname{is\text{-}prime\text{-}nat}(p) \Rightarrow ((p \in \mathbb{N}) \wedge (1 < p)))}}
\vnbprop{prime-nat-euclid}{\vnbstmt{\begin{array}{@{}l@{}} \forall p.; \forall a, b \in \mathbb{Z}. \\ \quad (\operatorname{is\text{-}prime\text{-}nat}(p) \wedge \operatorname{zz\text{-}divides}(p, (a \cdot b))) \Rightarrow  \\ \quad \quad \operatorname{zz\text{-}divides}(p, a) \vee \operatorname{zz\text{-}divides}(p, b) \end{array}}}
\vnbprop{zz-divides-small}{\vnbstmt{\begin{array}{@{}l@{}} \forall p, m \in \mathbb{Z}. \\ \quad ((0 < p) \wedge ((0 < m) \wedge (m < p))) \Rightarrow \neg \operatorname{zz\text{-}divides}(p, m) \end{array}}}
\vnbprop{prime-nat-not-divides-factorial}{\vnbstmt{\begin{array}{@{}l@{}} \forall i \in \mathbb{N}, p. \\ \quad (\operatorname{is\text{-}prime\text{-}nat}(p) \wedge (i < p)) \Rightarrow \neg \operatorname{zz\text{-}divides}(p, \operatorname{factorial}(i)) \end{array}}}
\vnbprop{prime-divides-binomial}{\vnbstmt{\begin{array}{@{}l@{}} \forall p.; \forall i \in \mathbb{N}. \\ \quad (\operatorname{is\text{-}prime\text{-}nat}(p) \wedge ((0 < i) \wedge (i < p))) \Rightarrow \operatorname{zz\text{-}divides}(p, \operatorname{choose}(p, i)) \end{array}}}
\vnbprop{choose-above-zero}{\vnbstmt{\forall m, k \in \mathbb{N}.\; ((m < k) \Rightarrow (\operatorname{choose}(m, k) = 0))}}
\vnbprop{characteristic-kills}{\vnbstmt{\begin{array}{@{}l@{}} \forall r, p.; \forall a \in \operatorname{carr}(r). \\ \quad (\operatorname{is\text{-}commutative\text{-}ring}(r) \wedge \operatorname{has\text{-}characteristic}(r, p)) \Rightarrow  \\ \quad \quad \operatorname{zz\text{-}act}(\operatorname{ring\text{-}additive\text{-}ag}(r), p, a) = \operatorname{zero}(r) \end{array}}}
\vnbprop{characteristic-multiple-zero}{\vnbstmt{\begin{array}{@{}l@{}} \forall r, p.; \forall c \in \mathbb{Z}. \\ \quad (\operatorname{is\text{-}commutative\text{-}ring}(r) \wedge \operatorname{has\text{-}characteristic}(r, p)) \Rightarrow  \\ \quad \quad \operatorname{zz\text{-}act}(\operatorname{ring\text{-}additive\text{-}ag}(r), (p \cdot c), \operatorname{one}(r)) = \operatorname{zero}(r) \end{array}}}
\vnbprop{comb-kk-closed-form}{\vnbstmt{\begin{array}{@{}l@{}} \forall m \in \mathbb{N}, r.; \forall x, y \in \operatorname{carr}(r), j, k \in \mathbb{N}. \\ \quad (\operatorname{is\text{-}commutative\text{-}ring}(r) \wedge ((j + k) = m)) \Rightarrow  \\ \quad \quad \operatorname{comb\text{-}kk}(r, x, y, m)(k) \\ \quad \quad =\; \operatorname{mul}(r)(\operatorname{zz\text{-}act}(\operatorname{ring\text{-}additive\text{-}ag}(r), \operatorname{choose}(m, k), \operatorname{one}(r)), \\ \quad \quad \quad \operatorname{mul}(r)(\operatorname{ring\text{-}power}(r, x, k), \operatorname{ring\text{-}power}(r, y, j))) \end{array}}}
\vnbprop{choose-diagonal}{\vnbstmt{\forall n \in \mathbb{N}.\; (\operatorname{choose}(n, n) = 1)}}
\vnbprop{sum-collapse-first}{\vnbstmt{\begin{array}{@{}l@{}} \forall n \in \mathbb{N}, r.; \forall f \in (\mathbb{Z} \to \operatorname{carr}(r)). \\ \quad (\operatorname{is\text{-}commutative\text{-}ring}(r) \wedge \forall i \in \mathbb{Z}.\; (((0 < i) \wedge (i \leq n)) \Rightarrow (f(i) = \operatorname{zero}(r)))) \Rightarrow  \\ \quad \quad \operatorname{sum}(r, f, (n + 1)) = f(0) \end{array}}}
\vnbprop{frobenius-additive}{\vnbstmt{\begin{array}{@{}l@{}} \forall r, p.; \forall x, y \in \operatorname{carr}(r). \\ \quad (\operatorname{is\text{-}commutative\text{-}ring}(r) \wedge \operatorname{has\text{-}characteristic}(r, p)) \Rightarrow  \\ \quad \quad \operatorname{ring\text{-}power}(r, \operatorname{add}(r)(x, y), p) = \operatorname{add}(r)(\operatorname{ring\text{-}power}(r, x, p), \operatorname{ring\text{-}power}(r, y, p)) \end{array}}}
\vnbprop{ring-power-nilpotent-zero}{\vnbstmt{\begin{array}{@{}l@{}} \forall n \in \mathbb{N}, r.; \forall a \in \operatorname{carr}(r). \\ \quad (\operatorname{is\text{-}integral\text{-}domain}(r) \wedge (\operatorname{ring\text{-}power}(r, a, (n + 1)) = \operatorname{zero}(r))) \Rightarrow  \\ \quad \quad a = \operatorname{zero}(r) \end{array}}}
\vnbprop{frobenius-injective}{\vnbstmt{\begin{array}{@{}l@{}} \forall r, p.; \forall a, b \in \operatorname{carr}(r). \\ \quad (\operatorname{is\text{-}integral\text{-}domain}(r) \wedge (\operatorname{has\text{-}characteristic}(r, p) \wedge (\operatorname{ring\text{-}power}(r, a, p) = \operatorname{ring\text{-}power}(r, b, p)))) \Rightarrow  \\ \quad \quad a = b \end{array}}}

\vnbbookfile{theorem-library/zz-primes}{theorem-library/zz-primes}
\noindent{\small theorem-library/zz-primes.scm -- PRIME FACTORISATION AND EUCLID'S LEMMA (ALGEBRAIC-NUMBERS notes, Props. 1.9, 1.10; ALGNUM, 2026-10-26).  Primes are galois's IS-PRIME-NAT (gal-char-vocab: p in NN and a prime element of ZZ-RING). zz-is-euclidean-domain       Z is a Euclidean domain in the galois-theory sense (degree ...}\par
\vnbprop{zz-is-euclidean-domain}{\vnbstmt{\operatorname{is\text{-}euclidean\text{-}domain}(\operatorname{zz\text{-}ring})}}
\vnbprop{zz-unit-neg}{\vnbstmt{\forall x \in \mathbb{Z}.\; (\operatorname{is\text{-}unit}(\operatorname{zz\text{-}ring}, x) \Rightarrow \operatorname{is\text{-}unit}(\operatorname{zz\text{-}ring}, -x))}}
\vnbprop{zz-prime-abs-prime-nat}{\vnbstmt{\forall p.\; (\operatorname{is\text{-}prime\text{-}element}(\operatorname{zz\text{-}ring}, p) \Rightarrow \operatorname{is\text{-}prime\text{-}nat}(\lvert p \rvert))}}
\vnbprop{prime-nat-associated-equal}{\vnbstmt{\begin{array}{@{}l@{}} \forall p, q. \\ \quad (\operatorname{is\text{-}prime\text{-}nat}(p) \wedge (\operatorname{is\text{-}prime\text{-}nat}(q) \wedge \operatorname{associated}(\operatorname{zz\text{-}ring}, p, q))) \Rightarrow (p = q) \end{array}}}
\vnbprop{zz-prod-abs}{\vnbstmt{\begin{array}{@{}l@{}} \forall n \in \mathbb{N}, q. \\ \quad \forall \mathit{ntq} \in \operatorname{interval}(1, n).\; (q(\mathit{ntq}) \in \mathbb{Z}) \Rightarrow  \\ \quad \quad \operatorname{prod\text{-}ring}(\operatorname{zz\text{-}ring}, (\lambda i \in \mathbb{N}.\; \lvert q(i) \rvert), \operatorname{interval}(1, n)) \\ \quad \quad =\; \lvert \operatorname{prod\text{-}ring}(\operatorname{zz\text{-}ring}, q, \operatorname{interval}(1, n)) \rvert \end{array}}}
\vnbprop{prime-divides-product}{\vnbstmt{\begin{array}{@{}l@{}} \forall n \in \mathbb{N}, p, a. \\ \quad \operatorname{is\text{-}prime\text{-}nat}(p) \wedge \\ \quad \forall q \in \operatorname{interval}(1, n).\; (a(q) \in \mathbb{Z}) \wedge \\ \quad \operatorname{zz\text{-}divides}(p, \operatorname{prod\text{-}ring}(\operatorname{zz\text{-}ring}, a, \operatorname{interval}(1, n))) \Rightarrow \\ \quad \quad \exists i \in \operatorname{interval}(1, n).\; \operatorname{zz\text{-}divides}(p, a(i)) \end{array}}}
\vnbprop{prime-factorization-exists}{\vnbstmt{\begin{array}{@{}l@{}} \forall a \in \mathbb{N}. \\ \quad (1 \leq a) \Rightarrow  \\ \quad \quad \exists k \in \mathbb{N}, p. \\ \quad \quad \quad (\forall i \in \operatorname{interval}(1, k).\; \operatorname{is\text{-}prime\text{-}nat}(p(i)) \wedge (a = \operatorname{prod\text{-}ring}(\operatorname{zz\text{-}ring}, p, \operatorname{interval}(1, k)))) \end{array}}}
\vnbprop{prime-factorization-unique}{\vnbstmt{\begin{array}{@{}l@{}} \forall k \in \mathbb{N}, p.; \forall m \in \mathbb{N}, q. \\ \quad \forall i \in \operatorname{interval}(1, k).\; \operatorname{is\text{-}prime\text{-}nat}(p(i)) \wedge \\ \quad \forall i \in \operatorname{interval}(1, m).\; \operatorname{is\text{-}prime\text{-}nat}(q(i)) \wedge \\ \quad \operatorname{prod\text{-}ring}(\operatorname{zz\text{-}ring}, p, \operatorname{interval}(1, k)) = \operatorname{prod\text{-}ring}(\operatorname{zz\text{-}ring}, q, \operatorname{interval}(1, m)) \Rightarrow \\ \quad \quad k = m \wedge \\ \quad \quad \exists s \in \mathrm{Bij}(\operatorname{interval}(1, k), \operatorname{interval}(1, m)). \forall i \in \operatorname{interval}(1, k). \\ \quad \quad \quad q(s(i)) = p(i) \end{array}}}
\vnbprop{multiples-count}{\vnbstmt{\begin{array}{@{}l@{}} \forall d, c. \\ \quad ((d \in \mathbb{N}) \wedge ((c \in \mathbb{N}) \wedge (1 \leq d))) \Rightarrow  \\ \quad \quad |\{k \in \operatorname{interval}(1, (d \cdot c)) : \operatorname{zz\text{-}divides}(d, k)\}| = c \end{array}}}
\vnbprop{zz-prime-factor-exists}{\vnbstmt{\begin{array}{@{}l@{}} \forall d \in \mathbb{Z}. \\ \quad (\neg (d = 0) \wedge (\neg (d = 1) \wedge \neg (d = -1))) \Rightarrow  \\ \quad \quad \exists p.\; (\operatorname{is\text{-}prime\text{-}nat}(p) \wedge \operatorname{zz\text{-}divides}(p, d)) \end{array}}}
\vnbprop{zz-coprime-iff-no-common-prime}{\vnbstmt{\begin{array}{@{}l@{}} \forall k, n. \\ \quad ((k \in \mathbb{Z}) \wedge ((n \in \mathbb{Z}) \wedge \neg (n = 0))) \Rightarrow  \\ \quad \quad \operatorname{zz\text{-}coprime}(k, n) \Leftrightarrow \forall q.\; ((\operatorname{is\text{-}prime\text{-}nat}(q) \wedge \operatorname{zz\text{-}divides}(q, n)) \Rightarrow \neg \operatorname{zz\text{-}divides}(q, k)) \end{array}}}

\vnbbookfile{theorem-library/rake-setoid2}{theorem-library/rake-setoid2}
\noindent{\small A file shared across books; the entries that belong to this one.}\par
\noindent{\small rake-setoid2.scm -- BATCH O of the 2026-09-17 rake: the SETOID quotient and the METRIC COMPLETION.  Twelve of the thirteen asserted leaves of the batch are PROVEN here; the thirteenth is a STATEMENT DEFECT and is recorded below. Every statement is its support's statement UNCHANGED (copied from the}\par
\vnbprop{rko2-related-unfold}{\vnbstmt{\forall s, a, b.\; (\operatorname{related}(s, a, b) \Leftrightarrow ({[}a, b] \in \operatorname{rel}(s)))}}
\vnbprop{rko2-refl}{\vnbstmt{\forall s.; \forall a \in \operatorname{pts}(s).\; (\operatorname{is\text{-}setoid}(s) \Rightarrow \operatorname{related}(s, a, a))}}
\vnbprop{rko2-sym}{\vnbstmt{\begin{array}{@{}l@{}} \forall s, a, b. \\ \quad (\operatorname{is\text{-}setoid}(s) \wedge ((a \in \operatorname{pts}(s)) \wedge ((b \in \operatorname{pts}(s)) \wedge \operatorname{related}(s, a, b)))) \Rightarrow  \\ \quad \quad \operatorname{related}(s, b, a) \end{array}}}
\vnbprop{rko2-trans}{\vnbstmt{\begin{array}{@{}l@{}} \forall s, a, b, c. \\ \quad \operatorname{is\text{-}setoid}(s) \wedge \\ \quad a \in \operatorname{pts}(s) \wedge \\ \quad b \in \operatorname{pts}(s) \wedge \\ \quad c \in \operatorname{pts}(s) \wedge \\ \quad \operatorname{related}(s, a, b) \wedge \\ \quad \operatorname{related}(s, b, c) \Rightarrow \\ \quad \quad \operatorname{related}(s, a, c) \end{array}}}
\vnbprop{class-self}{\vnbstmt{\forall s.; \forall a \in \operatorname{pts}(s).\; (\operatorname{is\text{-}setoid}(s) \Rightarrow (a \in \operatorname{class}(s, a)))}}
\vnbprop{class-in-quotient}{\vnbstmt{\forall s.; \forall a \in \operatorname{pts}(s).\; (\operatorname{is\text{-}setoid}(s) \Rightarrow (\operatorname{class}(s, a) \in \operatorname{quotient}(s)))}}
\vnbprop{quotient-rep}{\vnbstmt{\begin{array}{@{}l@{}} \forall s.; \forall x \in \operatorname{quotient}(s). \\ \quad \operatorname{is\text{-}setoid}(s) \Rightarrow \exists a \in \operatorname{pts}(s).\; (x = \operatorname{class}(s, a)) \end{array}}}
\vnbprop{class-subset-carrier}{\vnbstmt{\forall s.; \forall a \in \operatorname{pts}(s).\; (\operatorname{is\text{-}setoid}(s) \Rightarrow (\operatorname{class}(s, a) \subseteq \operatorname{pts}(s)))}}
\vnbprop{class-eq-iff}{\vnbstmt{\begin{array}{@{}l@{}} \forall s.; \forall a, b \in \operatorname{pts}(s). \\ \quad \operatorname{is\text{-}setoid}(s) \Rightarrow ((\operatorname{class}(s, a) = \operatorname{class}(s, b)) \Leftrightarrow \operatorname{related}(s, a, b)) \end{array}}}
\vnbprop{class-disjoint}{\vnbstmt{\begin{array}{@{}l@{}} \forall s.; \forall a, b \in \operatorname{pts}(s). \\ \quad \operatorname{is\text{-}setoid}(s) \Rightarrow  \\ \quad \quad (\operatorname{class}(s, a) = \operatorname{class}(s, b)) \vee ((\operatorname{class}(s, a) \cap \operatorname{class}(s, b)) = \emptyset) \end{array}}}
\vnbprop{descend-computes}{\vnbstmt{\begin{array}{@{}l@{}} \forall s, z.; \forall f \in (\operatorname{pts}(s) \to z), a \in \operatorname{pts}(s). \\ \quad (\operatorname{is\text{-}setoid}(s) \wedge \operatorname{respects}(s, f)) \Rightarrow (\operatorname{descend}(s, f)(\operatorname{class}(s, a)) = f(a)) \end{array}}}
\vnbprop{descend2-computes}{\vnbstmt{\begin{array}{@{}l@{}} \forall s, z.; \forall f \in ((\operatorname{pts}(s) \times \operatorname{pts}(s)) \to z), a, b. \\ \quad (\operatorname{is\text{-}setoid}(s) \wedge (\operatorname{respects2}(s, f) \wedge ((a \in \operatorname{pts}(s)) \wedge (b \in \operatorname{pts}(s))))) \Rightarrow  \\ \quad \quad \operatorname{descend2}(s, f)(\operatorname{class}(s, a), \operatorname{class}(s, b)) = f(a, b) \end{array}}}
\vnbprop{quotient-universal}{\vnbstmt{\begin{array}{@{}l@{}} \forall s, z.; \forall f \in (\operatorname{pts}(s) \to z). \\ \quad (\operatorname{is\text{-}setoid}(s) \wedge \operatorname{respects}(s, f)) \Rightarrow  \\ \quad \quad \exists g \in (\operatorname{quotient}(s) \to z). \\ \quad \quad \quad \forall a \in \operatorname{pts}(s).\; (g(\operatorname{class}(s, a)) = f(a)) \wedge \\ \quad \quad \quad \forall b \in (\operatorname{quotient}(s) \to z). \\ \quad \quad \quad \quad \forall a \in \operatorname{pts}(s).\; (b(\operatorname{class}(s, a)) = f(a)) \Rightarrow (b = g) \end{array}}}

\vnbbookfile{theorem-library/ringoid-setoid-proof}{theorem-library/ringoid-setoid-proof}
\noindent{\small ringoid-setoid-proof.scm}\par
\vnbprop{ringoid-rel-is-equivalence}{\vnbstmt{\forall r.\; (\operatorname{is\text{-}ringoid}(r) \Rightarrow \operatorname{is\text{-}equivalence}(\operatorname{ringoid\text{-}rel}(r), \operatorname{carr}(r)))}}
\vnbprop{ringoid-setoid-is-setoid}{\vnbstmt{\forall r.\; (\operatorname{is\text{-}ringoid}(r) \Rightarrow \operatorname{is\text{-}setoid}(\operatorname{ringoid\text{-}setoid}(r)))}}

\vnbbookfile{theorem-library/rake-ringoid}{theorem-library/rake-ringoid}
\noindent{\small rake-ringoid.scm -- rq-add-computes: addition on R/I computes on classes.}\par
\vnbprop{r6b-setoid-rel}{\vnbstmt{\forall r.\; (\operatorname{is\text{-}ringoid}(r) \Rightarrow (\operatorname{rel}(\operatorname{ringoid\text{-}setoid}(r)) \simeq \operatorname{ringoid\text{-}rel}(r)))}}
\vnbprop{r6b-related-out}{\vnbstmt{\begin{array}{@{}l@{}} \forall r, a, b. \\ \quad (\operatorname{is\text{-}ringoid}(r) \wedge ((a \in \operatorname{carr}(r)) \wedge ((b \in \operatorname{carr}(r)) \wedge \operatorname{related}(\operatorname{ringoid\text{-}setoid}(r), a, b)))) \Rightarrow  \\ \quad \quad \operatorname{add}(r)(a, \operatorname{neg}(r)(b)) \in \operatorname{idl}(r) \end{array}}}
\vnbprop{r6b-related-in}{\vnbstmt{\begin{array}{@{}l@{}} \forall r, a, b. \\ \quad (\operatorname{is\text{-}ringoid}(r) \wedge ((a \in \operatorname{carr}(r)) \wedge ((b \in \operatorname{carr}(r)) \wedge (\operatorname{add}(r)(a, \operatorname{neg}(r)(b)) \in \operatorname{idl}(r))))) \Rightarrow  \\ \quad \quad \operatorname{related}(\operatorname{ringoid\text{-}setoid}(r), a, b) \end{array}}}
\vnbprop{r6b-rq-add-lam-type}{\vnbstmt{\begin{array}{@{}l@{}} \forall r. \\ \quad \operatorname{is\text{-}ringoid}(r) \Rightarrow  \\ \quad \quad (\lambda {[}a, b] \in (\operatorname{carr}(r) \times \operatorname{carr}(r)).\; \operatorname{class}(\operatorname{ringoid\text{-}setoid}(r), \operatorname{add}(r)(a, b))) \\ \quad \quad \in\; ((\operatorname{pts}(\operatorname{ringoid\text{-}setoid}(r)) \times \operatorname{pts}(\operatorname{ringoid\text{-}setoid}(r))) \to \operatorname{quotient}(\operatorname{ringoid\text{-}setoid}(r))) \end{array}}}
\vnbprop{r6b-add-respects2}{\vnbstmt{\begin{array}{@{}l@{}} \forall r. \\ \quad \operatorname{is\text{-}ringoid}(r) \Rightarrow  \\ \quad \quad \operatorname{respects2}(\operatorname{ringoid\text{-}setoid}(r), \\ \quad \quad \quad (\lambda {[}a, b] \in (\operatorname{carr}(r) \times \operatorname{carr}(r)).\; \operatorname{class}(\operatorname{ringoid\text{-}setoid}(r), \operatorname{add}(r)(a, b)))) \end{array}}}
\vnbprop{rq-add-computes}{\vnbstmt{\begin{array}{@{}l@{}} \forall r, a, b. \\ \quad (\operatorname{is\text{-}ringoid}(r) \wedge ((a \in \operatorname{carr}(r)) \wedge (b \in \operatorname{carr}(r)))) \Rightarrow  \\ \quad \quad \operatorname{add}(\operatorname{ringoid\text{-}quotient\text{-}ring}(r))(\operatorname{class}(\operatorname{ringoid\text{-}setoid}(r), a), \operatorname{class}(\operatorname{ringoid\text{-}setoid}(r), b)) \\ \quad \quad =\; \operatorname{class}(\operatorname{ringoid\text{-}setoid}(r), \operatorname{add}(r)(a, b)) \end{array}}}

\vnbbookfile{theorem-library/rake-ringoid-quotient}{theorem-library/rake-ringoid-quotient}
\noindent{\small rake-ringoid-quotient.scm -- MUL and NEG compute on R/I.}\par
\vnbprop{r7n-rmul-type}{\vnbstmt{\begin{array}{@{}l@{}} \forall r, a, b. \\ \quad (\operatorname{is\text{-}ringoid}(r) \wedge ((a \in \operatorname{carr}(r)) \wedge (b \in \operatorname{carr}(r)))) \Rightarrow  \\ \quad \quad \operatorname{mul}(r)(a, b) \in \operatorname{carr}(r) \end{array}}}
\vnbprop{r7n-rone-type}{\vnbstmt{\forall r.\; (\operatorname{is\text{-}ringoid}(r) \Rightarrow (\operatorname{one}(r) \in \operatorname{carr}(r)))}}
\vnbprop{r7n-radd-idem-zero}{\vnbstmt{\begin{array}{@{}l@{}} \forall r.; \forall a \in \operatorname{carr}(r). \\ \quad (\operatorname{is\text{-}ringoid}(r) \wedge (\operatorname{add}(r)(a, a) = a)) \Rightarrow (a = \operatorname{zero}(r)) \end{array}}}
\vnbprop{r7n-rzero-mul-left}{\vnbstmt{\forall r.; \forall b \in \operatorname{carr}(r).\; (\operatorname{is\text{-}ringoid}(r) \Rightarrow (\operatorname{mul}(r)(\operatorname{zero}(r), b) = \operatorname{zero}(r)))}}
\vnbprop{r7n-rmul-zero-right}{\vnbstmt{\forall r.; \forall a \in \operatorname{carr}(r).\; (\operatorname{is\text{-}ringoid}(r) \Rightarrow (\operatorname{mul}(r)(a, \operatorname{zero}(r)) = \operatorname{zero}(r)))}}
\vnbprop{r7n-rneg-mul-left}{\vnbstmt{\begin{array}{@{}l@{}} \forall r, a, b. \\ \quad (\operatorname{is\text{-}ringoid}(r) \wedge ((a \in \operatorname{carr}(r)) \wedge (b \in \operatorname{carr}(r)))) \Rightarrow  \\ \quad \quad \operatorname{mul}(r)(\operatorname{neg}(r)(a), b) = \operatorname{neg}(r)(\operatorname{mul}(r)(a, b)) \end{array}}}
\vnbprop{r7n-rmul-neg-right}{\vnbstmt{\begin{array}{@{}l@{}} \forall r, a, b. \\ \quad (\operatorname{is\text{-}ringoid}(r) \wedge ((a \in \operatorname{carr}(r)) \wedge (b \in \operatorname{carr}(r)))) \Rightarrow  \\ \quad \quad \operatorname{mul}(r)(a, \operatorname{neg}(r)(b)) = \operatorname{neg}(r)(\operatorname{mul}(r)(a, b)) \end{array}}}
\vnbprop{r7n-rdiff-mul}{\vnbstmt{\begin{array}{@{}l@{}} \forall r, a, b, u, v. \\ \quad (\operatorname{is\text{-}ringoid}(r) \wedge ((a \in \operatorname{carr}(r)) \wedge ((b \in \operatorname{carr}(r)) \wedge ((u \in \operatorname{carr}(r)) \wedge (v \in \operatorname{carr}(r)))))) \Rightarrow  \\ \quad \quad \operatorname{add}(r)(\operatorname{mul}(r)(a, b), \operatorname{neg}(r)(\operatorname{mul}(r)(u, v))) \\ \quad \quad =\; \operatorname{add}(r)(\operatorname{mul}(r)(\operatorname{add}(r)(a, \operatorname{neg}(r)(u)), b), \operatorname{mul}(r)(u, \operatorname{add}(r)(b, \operatorname{neg}(r)(v)))) \end{array}}}
\vnbprop{r7n-rq-mul-lam-type}{\vnbstmt{\begin{array}{@{}l@{}} \forall r. \\ \quad \operatorname{is\text{-}ringoid}(r) \Rightarrow  \\ \quad \quad (\lambda {[}a, b] \in (\operatorname{carr}(r) \times \operatorname{carr}(r)).\; \operatorname{class}(\operatorname{ringoid\text{-}setoid}(r), \operatorname{mul}(r)(a, b))) \\ \quad \quad \in\; ((\operatorname{pts}(\operatorname{ringoid\text{-}setoid}(r)) \times \operatorname{pts}(\operatorname{ringoid\text{-}setoid}(r))) \to \operatorname{quotient}(\operatorname{ringoid\text{-}setoid}(r))) \end{array}}}
\vnbprop{r7n-mul-respects2}{\vnbstmt{\begin{array}{@{}l@{}} \forall r. \\ \quad \operatorname{is\text{-}ringoid}(r) \Rightarrow  \\ \quad \quad \operatorname{respects2}(\operatorname{ringoid\text{-}setoid}(r), \\ \quad \quad \quad (\lambda {[}a, b] \in (\operatorname{carr}(r) \times \operatorname{carr}(r)).\; \operatorname{class}(\operatorname{ringoid\text{-}setoid}(r), \operatorname{mul}(r)(a, b)))) \end{array}}}
\vnbprop{rq-mul-computes}{\vnbstmt{\begin{array}{@{}l@{}} \forall r, a, b. \\ \quad (\operatorname{is\text{-}ringoid}(r) \wedge ((a \in \operatorname{carr}(r)) \wedge (b \in \operatorname{carr}(r)))) \Rightarrow  \\ \quad \quad \operatorname{mul}(\operatorname{ringoid\text{-}quotient\text{-}ring}(r))(\operatorname{class}(\operatorname{ringoid\text{-}setoid}(r), a), \operatorname{class}(\operatorname{ringoid\text{-}setoid}(r), b)) \\ \quad \quad =\; \operatorname{class}(\operatorname{ringoid\text{-}setoid}(r), \operatorname{mul}(r)(a, b)) \end{array}}}
\vnbprop{r7n-rq-neg-lam-type}{\vnbstmt{\begin{array}{@{}l@{}} \forall r. \\ \quad \operatorname{is\text{-}ringoid}(r) \Rightarrow  \\ \quad \quad (\lambda a \in \operatorname{carr}(r).\; \operatorname{class}(\operatorname{ringoid\text{-}setoid}(r), \operatorname{neg}(r)(a))) \\ \quad \quad \in\; (\operatorname{pts}(\operatorname{ringoid\text{-}setoid}(r)) \to \operatorname{quotient}(\operatorname{ringoid\text{-}setoid}(r))) \end{array}}}
\vnbprop{r7n-neg-respects}{\vnbstmt{\begin{array}{@{}l@{}} \forall r. \\ \quad \operatorname{is\text{-}ringoid}(r) \Rightarrow  \\ \quad \quad \operatorname{respects}(\operatorname{ringoid\text{-}setoid}(r), (\lambda a \in \operatorname{carr}(r).\; \operatorname{class}(\operatorname{ringoid\text{-}setoid}(r), \operatorname{neg}(r)(a)))) \end{array}}}
\vnbprop{rq-neg-computes}{\vnbstmt{\begin{array}{@{}l@{}} \forall r.; \forall a \in \operatorname{carr}(r). \\ \quad \operatorname{is\text{-}ringoid}(r) \Rightarrow  \\ \quad \quad \operatorname{neg}(\operatorname{ringoid\text{-}quotient\text{-}ring}(r))(\operatorname{class}(\operatorname{ringoid\text{-}setoid}(r), a)) = \operatorname{class}(\operatorname{ringoid\text{-}setoid}(r), \operatorname{neg}(r)(a)) \end{array}}}
\vnbprop{ringoid-quotient-is-ring}{\vnbstmt{\forall r.\; (\operatorname{is\text{-}ringoid}(r) \Rightarrow \operatorname{is\text{-}ring}(\operatorname{ringoid\text{-}quotient\text{-}ring}(r)))}}

\vnbbookfile{theorem-library/finsum-cm-ptwise-laws}{theorem-library/finsum-cm-ptwise-laws}
\noindent{\small finsum-cm-ptwise-laws.scm -- lemmas MOVED 2026-11-10 (integration of the 11-09/10 waves) out of theorem-library/integral-laws.scm into library algebra, for the library gate: cited by l2-families-2, ccr6-f, ccr6-g; finite sums in a commutative monoid. The proofs are the originals, verbatim; the helper defines are a copy ...}\par
\vnbprop{finsum-cm-all-iden-ptwise}{\vnbstmt{\begin{array}{@{}l@{}} \forall m.; \forall s \in \mathbf{Set}, f. \\ \quad (\operatorname{is\text{-}comm\text{-}monoid}(m) \wedge ((|s| \in \mathbb{N}) \wedge \forall z \in s.\; (f(z) = \operatorname{iden}(m)))) \Rightarrow  \\ \quad \quad \operatorname{finsum}(m, f, s) = \operatorname{iden}(m) \end{array}}}
\vnbprop{finsum-cm-add-ptwise}{\vnbstmt{\begin{array}{@{}l@{}} \forall m.; \forall s \in \mathbf{Set}, f, g, h. \\ \quad \operatorname{is\text{-}comm\text{-}monoid}(m) \wedge \\ \quad |s| \in \mathbb{N} \wedge \\ \quad \forall z \in s.\; (f(z) \in \operatorname{carr}(m)) \wedge \\ \quad \forall z \in s.\; (g(z) \in \operatorname{carr}(m)) \wedge \\ \quad \forall z \in s.\; (h(z) = \operatorname{opr}(m)(f(z), g(z))) \Rightarrow \\ \quad \quad \operatorname{finsum}(m, h, s) = \operatorname{opr}(m)(\operatorname{finsum}(m, f, s), \operatorname{finsum}(m, g, s)) \end{array}}}
\vnbprop{finsum-cm-pad-ptwise}{\vnbstmt{\begin{array}{@{}l@{}} \forall m.; \forall t \in \mathbf{Set}, r, g. \\ \quad \operatorname{is\text{-}comm\text{-}monoid}(m) \wedge \\ \quad |t| \in \mathbb{N} \wedge \\ \quad \forall z \in r.\; (z \in t) \wedge \\ \quad \forall z \in t.\; (g(z) \in \operatorname{carr}(m)) \wedge \\ \quad \forall z \in t.\; (\neg (z \in r) \Rightarrow (g(z) = \operatorname{iden}(m))) \Rightarrow \\ \quad \quad \operatorname{finsum}(m, g, t) = \operatorname{finsum}(m, g, r) \end{array}}}

\vnbbookfile{structure-library/cg2-vocab}{structure-library/cg2-vocab}
\noindent{\small cg2-vocab.scm -- set products and set powers in a group (compact-groups notes, Prop 5.8: H = union of the V\textasciicircum{}k).  DEFINITIONS ONLY.  Agent ISOM, 2026-10-26 (the user's to confirm).}\par
\par\medskip\noindent\textbf{Definition} (\texttt{set-prod}).\;$\operatorname{set\text{-}prod}(\operatorname{spg\_}, \operatorname{spa\_}, \operatorname{spb\_}) :=$
\begin{equation*}
\{\operatorname{spz\_} \in \operatorname{carr}(\operatorname{spg\_}) : \exists \operatorname{spx\_} \in \operatorname{spa\_}, \operatorname{spy\_} \in \operatorname{spb\_}.\; (\operatorname{spz\_} = \operatorname{opr}(\operatorname{spg\_})(\operatorname{spx\_}, \operatorname{spy\_}))\}
\end{equation*}

\par\medskip\noindent\textbf{Definition} (\texttt{set-power}).\;For all $\operatorname{sqg\_}$ and $\operatorname{sqv\_}$,  $(\operatorname{set\text{-}power}(\operatorname{sqg\_}, \operatorname{sqv\_}, 0) \simeq \operatorname{singleton}(\operatorname{iden}(\operatorname{sqg\_})))$\par
Let $\operatorname{sqg\_}$ and $\operatorname{sqv\_}$ be arbitrary. Let $\operatorname{sqk\_} \in \mathbb{N}$. Then
\begin{equation*}
(\operatorname{set\text{-}power}(\operatorname{sqg\_}, \operatorname{sqv\_}, \operatorname{succ}(\operatorname{sqk\_})) \simeq \operatorname{set\text{-}prod}(\operatorname{sqg\_}, \operatorname{set\text{-}power}(\operatorname{sqg\_}, \operatorname{sqv\_}, \operatorname{sqk\_}), \operatorname{sqv\_}))
\end{equation*}

\par\medskip\noindent\textbf{Definition} (\texttt{powers-union}).\;$\operatorname{powers\text{-}union}(\operatorname{pug\_}, \operatorname{puv\_}) :=$ $\operatorname{big\text{-}union}(\operatorname{puk\_}, \mathbb{N}, \operatorname{set\text{-}power}(\operatorname{pug\_}, \operatorname{puv\_}, \operatorname{puk\_}))$
\par\medskip\noindent\textbf{Definition} (\texttt{is-subgroup}).\;For all $\operatorname{sgg\_}$ and $\operatorname{sgh\_}$, 
\begin{equation*}
\begin{array}{@{}l@{}} \operatorname{is\text{-}subgroup}(\operatorname{sgg\_}, \operatorname{sgh\_}) \\ \Leftrightarrow\; \operatorname{sgh\_} \subseteq \operatorname{carr}(\operatorname{sgg\_}) \wedge \\ \operatorname{iden}(\operatorname{sgg\_}) \in \operatorname{sgh\_} \wedge \\ \forall \operatorname{sga\_}, \operatorname{sgb\_} \in \operatorname{sgh\_}.\; (\operatorname{opr}(\operatorname{sgg\_})(\operatorname{sga\_}, \operatorname{sgb\_}) \in \operatorname{sgh\_}) \wedge \\ \forall \operatorname{sga\_} \in \operatorname{sgh\_}.\; (\operatorname{inv}(\operatorname{sgg\_})(\operatorname{sga\_}) \in \operatorname{sgh\_}) \end{array}
\end{equation*}

\par\medskip\noindent\textbf{Definition} (\texttt{is-normal-subgroup}).\;For all $\operatorname{sgg\_}$ and $\operatorname{sgh\_}$, 
\begin{equation*}
\begin{array}{@{}l@{}} \operatorname{is\text{-}normal\text{-}subgroup}(\operatorname{sgg\_}, \operatorname{sgh\_}) \\ \Leftrightarrow\; \operatorname{is\text{-}subgroup}(\operatorname{sgg\_}, \operatorname{sgh\_}) \wedge \\ \forall \operatorname{sgu\_} \in \operatorname{carr}(\operatorname{sgg\_}), \operatorname{sga\_} \in \operatorname{sgh\_}. \\ \quad \operatorname{opr}(\operatorname{sgg\_})(\operatorname{opr}(\operatorname{sgg\_})(\operatorname{sgu\_}, \operatorname{sga\_}), \operatorname{inv}(\operatorname{sgg\_})(\operatorname{sgu\_})) \in \operatorname{sgh\_} \end{array}
\end{equation*}

\par\medskip\noindent\textbf{Definition} (\texttt{gpow}).\;For all $\operatorname{gpg\_}$ and $\operatorname{gpx\_}$,  $(\operatorname{gpow}(\operatorname{gpg\_}, \operatorname{gpx\_}, 0) \simeq \operatorname{iden}(\operatorname{gpg\_}))$\par
Let $\operatorname{gpg\_}$ and $\operatorname{gpx\_}$ be arbitrary. Let $\operatorname{gpn\_} \in \mathbb{N}$. Then
\begin{equation*}
(\operatorname{gpow}(\operatorname{gpg\_}, \operatorname{gpx\_}, \operatorname{succ}(\operatorname{gpn\_})) \simeq \operatorname{opr}(\operatorname{gpg\_})(\operatorname{gpx\_}, \operatorname{gpow}(\operatorname{gpg\_}, \operatorname{gpx\_}, \operatorname{gpn\_})))
\end{equation*}

\par\medskip\noindent\textbf{Definition} (\texttt{grp-monoid}).\;$\operatorname{grp\text{-}monoid}(\operatorname{gmg\_}) :=$ ${[}\operatorname{carr}(\operatorname{gmg\_}), \operatorname{opr}(\operatorname{gmg\_}), \operatorname{iden}(\operatorname{gmg\_})]$

\vnbbookfile{theorem-library/tg-group}{theorem-library/tg-group}
\noindent{\small tg-group.scm -- group algebra the topological-group work needs (right identity, right inverse, inverse of an inverse, inverse of a product) and the monoid underlying a group. Agent ISOM, 2026-10-26.  Helper prefix: tgg-.  Library: algebra (GRP-MONOID: tg-vocab).}\par
\vnbprop{group-right-inv-id}{\vnbstmt{\begin{array}{@{}l@{}} \forall s.; \forall a \in \operatorname{carr}(s). \\ \quad \operatorname{is\text{-}group}(s) \Rightarrow  \\ \quad \quad (\operatorname{opr}(s)(a, \operatorname{inv}(s)(a)) = \operatorname{iden}(s)) \wedge (\operatorname{opr}(s)(a, \operatorname{iden}(s)) = a) \end{array}}}
\vnbprop{group-inv-inv}{\vnbstmt{\forall s.; \forall a \in \operatorname{carr}(s).\; (\operatorname{is\text{-}group}(s) \Rightarrow (\operatorname{inv}(s)(\operatorname{inv}(s)(a)) = a))}}
\vnbprop{group-inv-opr}{\vnbstmt{\begin{array}{@{}l@{}} \forall s.; \forall a, b \in \operatorname{carr}(s). \\ \quad \operatorname{is\text{-}group}(s) \Rightarrow  \\ \quad \quad \operatorname{inv}(s)(\operatorname{opr}(s)(a, b)) = \operatorname{opr}(s)(\operatorname{inv}(s)(b), \operatorname{inv}(s)(a)) \end{array}}}
\vnbprop{group-inv-iden}{\vnbstmt{\forall s.\; (\operatorname{is\text{-}group}(s) \Rightarrow (\operatorname{inv}(s)(\operatorname{iden}(s)) = \operatorname{iden}(s)))}}
\vnbprop{grp-monoid-carr}{\vnbstmt{\forall s.\; (\operatorname{is\text{-}group}(s) \Rightarrow (\operatorname{carr}(\operatorname{grp\text{-}monoid}(s)) \simeq \operatorname{carr}(s)))}}
\vnbprop{grp-monoid-opr}{\vnbstmt{\forall s.\; (\operatorname{is\text{-}group}(s) \Rightarrow (\operatorname{opr}(\operatorname{grp\text{-}monoid}(s)) \simeq \operatorname{opr}(s)))}}
\vnbprop{grp-monoid-iden}{\vnbstmt{\forall s.\; (\operatorname{is\text{-}group}(s) \Rightarrow (\operatorname{iden}(\operatorname{grp\text{-}monoid}(s)) \simeq \operatorname{iden}(s)))}}
\vnbprop{grp-monoid-is-monoid}{\vnbstmt{\forall s.\; (\operatorname{is\text{-}group}(s) \Rightarrow \operatorname{is\text{-}monoid}(\operatorname{grp\text{-}monoid}(s)))}}
\vnbprop{group-carr-set}{\vnbstmt{\forall s.\; (\operatorname{is\text{-}group}(s) \Rightarrow (\operatorname{carr}(s) \in \mathbf{Set}))}}
\vnbprop{group-cancel-mid}{\vnbstmt{\begin{array}{@{}l@{}} \forall s.; \forall a, b, c \in \operatorname{carr}(s). \\ \quad \operatorname{is\text{-}group}(s) \Rightarrow  \\ \quad \quad \operatorname{opr}(s)(\operatorname{opr}(s)(\operatorname{inv}(s)(a), b), \operatorname{opr}(s)(\operatorname{inv}(s)(b), c)) = \operatorname{opr}(s)(\operatorname{inv}(s)(a), c) \end{array}}}
\vnbprop{group-inv-translate}{\vnbstmt{\begin{array}{@{}l@{}} \forall s.; \forall a, x, y \in \operatorname{carr}(s). \\ \quad \operatorname{is\text{-}group}(s) \Rightarrow  \\ \quad \quad \operatorname{opr}(s)(\operatorname{inv}(s)(\operatorname{opr}(s)(a, y)), \operatorname{opr}(s)(a, x)) = \operatorname{opr}(s)(\operatorname{inv}(s)(y), x) \end{array}}}
\vnbprop{group-eq-of-inv-eq-iden}{\vnbstmt{\begin{array}{@{}l@{}} \forall s.; \forall u, v \in \operatorname{carr}(s). \\ \quad (\operatorname{is\text{-}group}(s) \wedge (\operatorname{opr}(s)(\operatorname{inv}(s)(v), u) = \operatorname{iden}(s))) \Rightarrow (u = v) \end{array}}}
\vnbprop{group-mul-inv-cancel}{\vnbstmt{\begin{array}{@{}l@{}} \forall s.; \forall y, p \in \operatorname{carr}(s). \\ \quad \operatorname{is\text{-}group}(s) \Rightarrow (\operatorname{opr}(s)(y, \operatorname{opr}(s)(\operatorname{inv}(s)(y), p)) = p) \end{array}}}

\vnbbookfile{theorem-library/cg2-setprod}{theorem-library/cg2-setprod}
\noindent{\small cg2-setprod.scm -- set products and set powers in a group: V\textasciicircum{}k inside the carrier, V\textasciicircum{}k V inside V\textasciicircum{}(k+1), and the union H of the V\textasciicircum{}k is a subgroup when V is symmetric and contains e (compact-groups notes, Prop 5.8, the algebraic part: "for any set E such that E = E\textasciicircum{}-1, the union of the E\textasciicircum{}k is clearly a subgroup").}\par
\vnbprop{set-prod-unfold}{\vnbstmt{\begin{array}{@{}l@{}} \forall g, a, b. \\ \quad (\operatorname{set\text{-}prod}(g, a, b) \simeq \{z \in \operatorname{carr}(g) : \exists x \in a, y \in b.\; (z = \operatorname{opr}(g)(x, y))\}) \end{array}}}
\vnbprop{set-prod-intro}{\vnbstmt{\begin{array}{@{}l@{}} \forall g, a, b, x, y. \\ \quad (\operatorname{is\text{-}group}(g) \wedge ((x \in \operatorname{carr}(g)) \wedge ((y \in \operatorname{carr}(g)) \wedge ((x \in a) \wedge (y \in b))))) \Rightarrow  \\ \quad \quad \operatorname{opr}(g)(x, y) \in \operatorname{set\text{-}prod}(g, a, b) \end{array}}}
\vnbprop{set-prod-elim}{\vnbstmt{\begin{array}{@{}l@{}} \forall g, a, b.; \forall z \in \operatorname{set\text{-}prod}(g, a, b). \\ \quad (z \in \operatorname{carr}(g)) \wedge \exists x \in a, y \in b.\; (z = \operatorname{opr}(g)(x, y)) \end{array}}}
\vnbprop{set-power-zero-elim}{\vnbstmt{\forall g, v.; \forall z \in \operatorname{set\text{-}power}(g, v, 0).\; (z = \operatorname{iden}(g))}}
\vnbprop{set-power-zero-iden}{\vnbstmt{\forall g, v.\; (\operatorname{is\text{-}group}(g) \Rightarrow (\operatorname{iden}(g) \in \operatorname{set\text{-}power}(g, v, 0)))}}
\vnbprop{set-power-is-set}{\vnbstmt{\forall k \in \mathbb{N}, g, v.\; (\operatorname{is\text{-}group}(g) \Rightarrow (\operatorname{set\text{-}power}(g, v, k) \in \mathbf{Set}))}}
\vnbprop{set-power-succ-elim}{\vnbstmt{\begin{array}{@{}l@{}} \forall g, v, k, z. \\ \quad (\operatorname{is\text{-}group}(g) \wedge ((k \in \mathbb{N}) \wedge (z \in \operatorname{set\text{-}power}(g, v, (k + 1))))) \Rightarrow  \\ \quad \quad (z \in \operatorname{carr}(g)) \wedge \exists x \in \operatorname{set\text{-}power}(g, v, k), y \in v.\; (z = \operatorname{opr}(g)(x, y)) \end{array}}}
\vnbprop{set-power-succ-intro}{\vnbstmt{\begin{array}{@{}l@{}} \forall g, v, k, x, y. \\ \quad (\operatorname{is\text{-}group}(g) \wedge ((k \in \mathbb{N}) \wedge ((x \in \operatorname{carr}(g)) \wedge ((y \in \operatorname{carr}(g)) \wedge ((x \in \operatorname{set\text{-}power}(g, v, k)) \wedge (y \in v)))))) \Rightarrow  \\ \quad \quad \operatorname{opr}(g)(x, y) \in \operatorname{set\text{-}power}(g, v, (k + 1)) \end{array}}}
\vnbprop{set-power-sub-carr}{\vnbstmt{\begin{array}{@{}l@{}} \forall k \in \mathbb{N}, g, v. \\ \quad (\operatorname{is\text{-}group}(g) \wedge (v \subseteq \operatorname{carr}(g))) \Rightarrow (\operatorname{set\text{-}power}(g, v, k) \subseteq \operatorname{carr}(g)) \end{array}}}
\vnbprop{set-power-mono-succ}{\vnbstmt{\begin{array}{@{}l@{}} \forall g, v, k, z. \\ \quad (\operatorname{is\text{-}group}(g) \wedge ((v \subseteq \operatorname{carr}(g)) \wedge ((\operatorname{iden}(g) \in v) \wedge ((k \in \mathbb{N}) \wedge (z \in \operatorname{set\text{-}power}(g, v, k)))))) \Rightarrow  \\ \quad \quad z \in \operatorname{set\text{-}power}(g, v, (k + 1)) \end{array}}}
\vnbprop{powers-union-unfold}{\vnbstmt{\forall g, v.\; (\operatorname{powers\text{-}union}(g, v) \simeq \operatorname{big\text{-}union}(k, \mathbb{N}, \operatorname{set\text{-}power}(g, v, k)))}}
\vnbprop{powers-union-is-set}{\vnbstmt{\forall g, v.\; (\operatorname{is\text{-}group}(g) \Rightarrow (\operatorname{powers\text{-}union}(g, v) \in \mathbf{Set}))}}
\vnbprop{powers-union-intro}{\vnbstmt{\begin{array}{@{}l@{}} \forall g, v, k, z. \\ \quad ((k \in \mathbb{N}) \wedge (z \in \operatorname{set\text{-}power}(g, v, k))) \Rightarrow (z \in \operatorname{powers\text{-}union}(g, v)) \end{array}}}
\vnbprop{powers-union-elim}{\vnbstmt{\begin{array}{@{}l@{}} \forall g, v.; \forall z \in \operatorname{powers\text{-}union}(g, v). \exists k \in \mathbb{N}. \\ \quad z \in \operatorname{set\text{-}power}(g, v, k) \end{array}}}
\vnbprop{powers-union-sub-carr}{\vnbstmt{\begin{array}{@{}l@{}} \forall g, v. \\ \quad (\operatorname{is\text{-}group}(g) \wedge (v \subseteq \operatorname{carr}(g))) \Rightarrow (\operatorname{powers\text{-}union}(g, v) \subseteq \operatorname{carr}(g)) \end{array}}}
\vnbprop{powers-union-iden}{\vnbstmt{\forall g, v.\; (\operatorname{is\text{-}group}(g) \Rightarrow (\operatorname{iden}(g) \in \operatorname{powers\text{-}union}(g, v)))}}
\vnbprop{powers-union-contains}{\vnbstmt{\begin{array}{@{}l@{}} \forall g, v. \\ \quad (\operatorname{is\text{-}group}(g) \wedge (v \subseteq \operatorname{carr}(g))) \Rightarrow (v \subseteq \operatorname{powers\text{-}union}(g, v)) \end{array}}}
\vnbprop{powers-union-right-mult}{\vnbstmt{\begin{array}{@{}l@{}} \forall g, v, z, y. \\ \quad (\operatorname{is\text{-}group}(g) \wedge ((v \subseteq \operatorname{carr}(g)) \wedge ((z \in \operatorname{powers\text{-}union}(g, v)) \wedge (y \in v)))) \Rightarrow  \\ \quad \quad \operatorname{opr}(g)(z, y) \in \operatorname{powers\text{-}union}(g, v) \end{array}}}
\vnbprop{powers-union-prod}{\vnbstmt{\begin{array}{@{}l@{}} \forall m \in \mathbb{N}, g, v.; \forall x \in \operatorname{powers\text{-}union}(g, v), y \in \operatorname{set\text{-}power}(g, v, m). \\ \quad (\operatorname{is\text{-}group}(g) \wedge (v \subseteq \operatorname{carr}(g))) \Rightarrow (\operatorname{opr}(g)(x, y) \in \operatorname{powers\text{-}union}(g, v)) \end{array}}}
\vnbprop{powers-union-inv}{\vnbstmt{\begin{array}{@{}l@{}} \forall k \in \mathbb{N}, g, v.; \forall y \in \operatorname{set\text{-}power}(g, v, k). \\ \quad (\operatorname{is\text{-}group}(g) \wedge ((v \subseteq \operatorname{carr}(g)) \wedge \forall p \in v.\; (\operatorname{inv}(g)(p) \in v))) \Rightarrow  \\ \quad \quad \operatorname{inv}(g)(y) \in \operatorname{powers\text{-}union}(g, v) \end{array}}}
\vnbprop{powers-union-subgroup}{\vnbstmt{\begin{array}{@{}l@{}} \forall g, v. \\ \quad (\operatorname{is\text{-}group}(g) \wedge ((v \subseteq \operatorname{carr}(g)) \wedge \forall p \in v.\; (\operatorname{inv}(g)(p) \in v))) \Rightarrow  \\ \quad \quad \operatorname{is\text{-}subgroup}(g, \operatorname{powers\text{-}union}(g, v)) \end{array}}}

\vnbbookfile{theorem-library/rake-asm-1}{theorem-library/rake-asm-1}
\noindent{\small A file shared across books; the entries that belong to this one.}\par
\noindent{\small rake-asm-1.scm -- RAKE-ASM agent RA-1 (2026-10-12): elementary assumptions PROVEN.}\par
\vnbprop{semigroup-assoc}{\vnbstmt{\begin{array}{@{}l@{}} \forall s.; \forall a, b, c \in \operatorname{carr}(s). \\ \quad \operatorname{is\text{-}semigroup}(s) \Rightarrow  \\ \quad \quad \operatorname{opr}(s)(\operatorname{opr}(s)(a, b), c) = \operatorname{opr}(s)(a, \operatorname{opr}(s)(b, c)) \end{array}}}
\vnbprop{qq-is-commutative-ring}{\vnbstmt{\operatorname{is\text{-}commutative\text{-}ring}(\operatorname{qq\text{-}ring})}}
\vnbprop{qq-is-integral-domain}{\vnbstmt{\operatorname{is\text{-}integral\text{-}domain}(\operatorname{qq\text{-}ring})}}
\vnbprop{zero-ring-is-ring}{\vnbstmt{\operatorname{is\text{-}ring}(\operatorname{zero\text{-}ring})}}
\vnbprop{ring-prod-is-ring}{\vnbstmt{\forall x, y.\; ((\operatorname{is\text{-}ring}(x) \wedge \operatorname{is\text{-}ring}(y)) \Rightarrow \operatorname{is\text{-}ring}(\operatorname{ring\text{-}prod}(x, y)))}}
\vnbprop{ring-prod-n-is-ring-ind}{\vnbstmt{\begin{array}{@{}l@{}} \forall n \in \mathbb{N}, f. \\ \quad \forall i \in \mathbb{N}.\; \operatorname{is\text{-}ring}(f(i)) \Rightarrow \operatorname{is\text{-}ring}(\operatorname{ring\text{-}prod\text{-}n}(f, n)) \end{array}}}
\vnbprop{ring-prod-n-is-ring}{\vnbstmt{\begin{array}{@{}l@{}} \forall f.; \forall n \in \mathbb{N}. \\ \quad \forall i \in \mathbb{N}.\; \operatorname{is\text{-}ring}(f(i)) \Rightarrow \operatorname{is\text{-}ring}(\operatorname{ring\text{-}prod\text{-}n}(f, n)) \end{array}}}

\vnbbookfile{theorem-library/cg4-gpow}{theorem-library/cg4-gpow}
\noindent{\small cg4-gpow.scm -- powers x\textasciicircum{}n (n in NN) in a group: typing, commuting with what commutes with x, and x\textasciicircum{}(m+n) = x\textasciicircum{}m x\textasciicircum{}n.  For the extension of a one-parameter group segment (Prop. 5.36). Agent ISOM, 2026-10-26.  Helper prefix: gp4-.  Library: algebra.}\par
\vnbprop{gpow-type}{\vnbstmt{\begin{array}{@{}l@{}} \forall n \in \mathbb{N}, g, x. \\ \quad (\operatorname{is\text{-}group}(g) \wedge (x \in \operatorname{carr}(g))) \Rightarrow (\operatorname{gpow}(g, x, n) \in \operatorname{carr}(g)) \end{array}}}
\vnbprop{gpow-commute}{\vnbstmt{\begin{array}{@{}l@{}} \forall n \in \mathbb{N}, g, x, y. \\ \quad (\operatorname{is\text{-}group}(g) \wedge ((x \in \operatorname{carr}(g)) \wedge ((y \in \operatorname{carr}(g)) \wedge (\operatorname{opr}(g)(y, x) = \operatorname{opr}(g)(x, y))))) \Rightarrow  \\ \quad \quad \operatorname{opr}(g)(y, \operatorname{gpow}(g, x, n)) = \operatorname{opr}(g)(\operatorname{gpow}(g, x, n), y) \end{array}}}
\vnbprop{gpow-add}{\vnbstmt{\begin{array}{@{}l@{}} \forall n \in \mathbb{N}, g, x, m. \\ \quad (\operatorname{is\text{-}group}(g) \wedge ((x \in \operatorname{carr}(g)) \wedge (m \in \mathbb{N}))) \Rightarrow  \\ \quad \quad \operatorname{gpow}(g, x, (m + n)) = \operatorname{opr}(g)(\operatorname{gpow}(g, x, m), \operatorname{gpow}(g, x, n)) \end{array}}}

\vnbbookfile{structure-library/group-cosets-vocab}{structure-library/group-cosets-vocab}
\noindent{\small structure-library/group-cosets-vocab.scm -- cosets and the subgroup generated by an element (ALGNUM, B5, 2026-10-26).  DEFINITIONS ONLY.  Library: algebra.  Load window: after the subgroup vocabulary (IS-SUBGROUP, scratchpad/cg1/cg1-vocab.scm, unwired).  Purely algebraic: no topology.}\par
\par\medskip\noindent\textbf{Definition} (\texttt{left-coset}).\;$\operatorname{left\text{-}coset}(\operatorname{gcg\_}, \operatorname{gca\_}, \operatorname{gch\_}) :=$
\begin{equation*}
\{\operatorname{gcx\_} \in \operatorname{carr}(\operatorname{gcg\_}) : \exists \operatorname{gcy\_} \in \operatorname{gch\_}.\; (\operatorname{gcx\_} = \operatorname{opr}(\operatorname{gcg\_})(\operatorname{gca\_}, \operatorname{gcy\_}))\}
\end{equation*}

\par\medskip\noindent\textbf{Definition} (\texttt{left-cosets}).\;$\operatorname{left\text{-}cosets}(\operatorname{gcg\_}, \operatorname{gch\_}) :=$
\begin{equation*}
\{\operatorname{gcc\_} \in \operatorname{power}(\operatorname{carr}(\operatorname{gcg\_})) : \exists \operatorname{gca\_} \in \operatorname{carr}(\operatorname{gcg\_}).\; (\operatorname{gcc\_} = \operatorname{left\text{-}coset}(\operatorname{gcg\_}, \operatorname{gca\_}, \operatorname{gch\_}))\}
\end{equation*}

\par\medskip\noindent\textbf{Definition} (\texttt{generated-subgroup}).\;$\operatorname{generated\text{-}subgroup}(\operatorname{gcg\_}, \operatorname{gcx\_}) :=$
\begin{equation*}
\{\operatorname{gcy\_} \in \operatorname{carr}(\operatorname{gcg\_}) : \forall \operatorname{gcs\_} \in \operatorname{power}(\operatorname{carr}(\operatorname{gcg\_})).\; (\operatorname{is\text{-}subgroup}(\operatorname{gcg\_}, \operatorname{gcs\_}) \Rightarrow ((\operatorname{gcx\_} \in \operatorname{gcs\_}) \Rightarrow (\operatorname{gcy\_} \in \operatorname{gcs\_})))\}
\end{equation*}


\vnbbookfile{structure-library/group-zpow-vocab}{structure-library/group-zpow-vocab}
\noindent{\small group-zpow-vocab.scm -- integer powers in a group (ALGNUM, 2026-10-27).  DEFINITIONS ONLY.  Library: algebra. Built ON worker-04's GPOW (cg4-vocab.scm: x\textasciicircum{}n for n in NN, x\textasciicircum{}0 = e, x\textasciicircum{}(n+1) = x x\textasciicircum{}n), one name per notion: GROUP-ZPOW(g, x, k)   x\textasciicircum{}k for k in ZZ: x\textasciicircum{}k = GPOW(g, x, k) for k $>$= 0, the inverse of x\textasciicircum{}(-k) for k $<$ 0.}\par
\par\medskip\noindent\textbf{Definition} (\texttt{group-zpow}).\;$\operatorname{group\text{-}zpow}(g, x, k) :=$ $\left(\text{if }(0 \leq k)\text{ then }\operatorname{gpow}(g, x, k)\text{ else }\operatorname{inv}(g)(\operatorname{gpow}(g, x, (0 - k)))\right)$
\par\medskip\noindent\textbf{Definition} (\texttt{is-element-order}).\;For all $g$, $x$ and $d$, 
\begin{equation*}
\begin{array}{@{}l@{}} \operatorname{is\text{-}element\text{-}order}(g, x, d) \\ \Leftrightarrow\; d \in \mathbb{N} \wedge \\ 1 \leq d \wedge \\ \operatorname{gpow}(g, x, d) = \operatorname{iden}(g) \wedge \\ \forall \operatorname{eod\_} \in \mathbb{N}. \\ \quad (1 \leq \operatorname{eod\_}) \Rightarrow ((\operatorname{gpow}(g, x, \operatorname{eod\_}) = \operatorname{iden}(g)) \Rightarrow (d \leq \operatorname{eod\_})) \end{array}
\end{equation*}

\par\medskip\noindent\textbf{Definition} (\texttt{is-generator}).\;For all $g$ and $x$, 
\begin{equation*}
\begin{array}{@{}l@{}} \operatorname{is\text{-}generator}(g, x) \\ \Leftrightarrow\; ((x \in \operatorname{carr}(g)) \wedge \forall \operatorname{geny\_} \in \operatorname{carr}(g).\; \exists \operatorname{genk\_} \in \mathbb{Z}.\; (\operatorname{geny\_} = \operatorname{group\text{-}zpow}(g, x, \operatorname{genk\_}))) \end{array}
\end{equation*}

\par\medskip\noindent\textbf{Definition} (\texttt{group-prod}).\;$\operatorname{group\text{-}prod}(x, y) :=$
\begin{equation*}
{[}(\operatorname{carr}(x) \times \operatorname{carr}(y)), (\lambda {[}\operatorname{gpp\_}, \operatorname{gpq\_}] \in ((\operatorname{carr}(x) \times \operatorname{carr}(y)) \times (\operatorname{carr}(x) \times \operatorname{carr}(y))).\; {[}\operatorname{opr}(x)(\operatorname{nth}(1, \operatorname{gpp\_}), \operatorname{nth}(1, \operatorname{gpq\_})), \operatorname{opr}(y)(\operatorname{nth}(2, \operatorname{gpp\_}), \operatorname{nth}(2, \operatorname{gpq\_}))]), {[}\operatorname{iden}(x), \operatorname{iden}(y)], (\lambda \operatorname{gpp\_} \in (\operatorname{carr}(x) \times \operatorname{carr}(y)).\; {[}\operatorname{inv}(x)(\operatorname{nth}(1, \operatorname{gpp\_})), \operatorname{inv}(y)(\operatorname{nth}(2, \operatorname{gpp\_}))])]
\end{equation*}

\par\medskip\noindent\textbf{Definition} (\texttt{triv-group}).\;
\begin{equation*}
\begin{array}{@{}l@{}} \operatorname{triv\text{-}group} \\ =\; {[}\{0\}, (\lambda {[}\operatorname{tgp\_}, \operatorname{tgq\_}] \in (\{0\} \times \{0\}).\; 0), 0, (\lambda \operatorname{tgp\_} \in \{0\}.\; 0)] \end{array}
\end{equation*}

\par\medskip\noindent\textbf{Definition} (\texttt{group-famprod}).\;For every $\operatorname{gfam\_}$,  $(\operatorname{group\text{-}famprod}(\operatorname{gfam\_}, 0) \simeq \operatorname{triv\text{-}group})$\par
Let $\operatorname{gfam\_}$ be arbitrary. Let $\operatorname{gfn\_} \in \mathbb{N}$. Then
\begin{equation*}
(\operatorname{group\text{-}famprod}(\operatorname{gfam\_}, \operatorname{succ}(\operatorname{gfn\_})) \simeq \operatorname{group\text{-}prod}(\operatorname{group\text{-}famprod}(\operatorname{gfam\_}, \operatorname{gfn\_}), \operatorname{gfam\_}(\operatorname{succ}(\operatorname{gfn\_}))))
\end{equation*}

\par\medskip\noindent\textbf{Definition} (\texttt{group-famtuple}).\;For every $\operatorname{xfam\_}$,  $(\operatorname{group\text{-}famtuple}(\operatorname{xfam\_}, 0) \simeq 0)$\par
Let $\operatorname{xfam\_}$ be arbitrary. Let $\operatorname{xfn\_} \in \mathbb{N}$. Then
\begin{equation*}
(\operatorname{group\text{-}famtuple}(\operatorname{xfam\_}, \operatorname{succ}(\operatorname{xfn\_})) \simeq {[}\operatorname{group\text{-}famtuple}(\operatorname{xfam\_}, \operatorname{xfn\_}), \operatorname{xfam\_}(\operatorname{succ}(\operatorname{xfn\_}))])
\end{equation*}


\vnbbookfile{theorem-library/group-zpow}{theorem-library/group-zpow}
\noindent{\small theorem-library/group-zpow.scm -- THE GROUP KIT: integer powers, element orders, the direct product; PROP. 1.36 (ALGEBRAIC-NUMBERS notes, ch. 1; ALGNUM, 2026-10-27).  Built on worker-04's GPOW (natural powers).  No topology. group-right-cancel, zpow-nat, zpow-neg-nat, zpow-type, zpow-succ, zpow-add-nat, zpow-sub-nat, ...}\par
\vnbprop{group-right-cancel}{\vnbstmt{\begin{array}{@{}l@{}} \forall g.; \forall a, b, c \in \operatorname{carr}(g). \\ \quad (\operatorname{is\text{-}group}(g) \wedge (a = \operatorname{opr}(g)(b, c))) \Rightarrow (\operatorname{opr}(g)(a, \operatorname{inv}(g)(c)) = b) \end{array}}}
\vnbprop{zpow-nat}{\vnbstmt{\begin{array}{@{}l@{}} \forall n \in \mathbb{N}, g.; \forall x \in \operatorname{carr}(g). \\ \quad \operatorname{is\text{-}group}(g) \Rightarrow (\operatorname{group\text{-}zpow}(g, x, n) = \operatorname{gpow}(g, x, n)) \end{array}}}
\vnbprop{zpow-neg-nat}{\vnbstmt{\begin{array}{@{}l@{}} \forall n \in \mathbb{N}, g.; \forall x \in \operatorname{carr}(g). \\ \quad \operatorname{is\text{-}group}(g) \Rightarrow (\operatorname{group\text{-}zpow}(g, x, (0 - n)) = \operatorname{inv}(g)(\operatorname{gpow}(g, x, n))) \end{array}}}
\vnbprop{zpow-type}{\vnbstmt{\begin{array}{@{}l@{}} \forall k \in \mathbb{Z}, g.; \forall x \in \operatorname{carr}(g). \\ \quad \operatorname{is\text{-}group}(g) \Rightarrow (\operatorname{group\text{-}zpow}(g, x, k) \in \operatorname{carr}(g)) \end{array}}}
\vnbprop{zpow-succ}{\vnbstmt{\begin{array}{@{}l@{}} \forall k \in \mathbb{Z}, g.; \forall x \in \operatorname{carr}(g). \\ \quad \operatorname{is\text{-}group}(g) \Rightarrow (\operatorname{group\text{-}zpow}(g, x, (k + 1)) = \operatorname{opr}(g)(\operatorname{group\text{-}zpow}(g, x, k), x)) \end{array}}}
\vnbprop{zpow-add-nat}{\vnbstmt{\begin{array}{@{}l@{}} \forall n \in \mathbb{N}, m \in \mathbb{Z}, g.; \forall x \in \operatorname{carr}(g). \\ \quad \operatorname{is\text{-}group}(g) \Rightarrow  \\ \quad \quad \operatorname{group\text{-}zpow}(g, x, (m + n)) = \operatorname{opr}(g)(\operatorname{group\text{-}zpow}(g, x, m), \operatorname{group\text{-}zpow}(g, x, n)) \end{array}}}
\vnbprop{zpow-sub-nat}{\vnbstmt{\begin{array}{@{}l@{}} \forall n \in \mathbb{N}, m \in \mathbb{Z}, g.; \forall x \in \operatorname{carr}(g). \\ \quad \operatorname{is\text{-}group}(g) \Rightarrow  \\ \quad \quad \operatorname{group\text{-}zpow}(g, x, (m - n)) = \operatorname{opr}(g)(\operatorname{group\text{-}zpow}(g, x, m), \operatorname{group\text{-}zpow}(g, x, (0 - n))) \end{array}}}
\vnbprop{zpow-add}{\vnbstmt{\begin{array}{@{}l@{}} \forall m, n \in \mathbb{Z}, g.; \forall x \in \operatorname{carr}(g). \\ \quad \operatorname{is\text{-}group}(g) \Rightarrow  \\ \quad \quad \operatorname{group\text{-}zpow}(g, x, (m + n)) = \operatorname{opr}(g)(\operatorname{group\text{-}zpow}(g, x, m), \operatorname{group\text{-}zpow}(g, x, n)) \end{array}}}
\vnbprop{zpow-inv}{\vnbstmt{\begin{array}{@{}l@{}} \forall k \in \mathbb{Z}, g.; \forall x \in \operatorname{carr}(g). \\ \quad \operatorname{is\text{-}group}(g) \Rightarrow (\operatorname{group\text{-}zpow}(g, x, (0 - k)) = \operatorname{inv}(g)(\operatorname{group\text{-}zpow}(g, x, k))) \end{array}}}
\vnbprop{zz-ring-add-is-add}{\vnbstmt{\forall u, v.\; (((u \in \mathbb{Z}) \wedge (v \in \mathbb{Z})) \Rightarrow (\operatorname{add}(\operatorname{zz\text{-}ring})(u, v) = (u + v)))}}
\vnbprop{zz-ring-neg-is-neg}{\vnbstmt{\forall u \in \mathbb{Z}.\; (\operatorname{neg}(\operatorname{zz\text{-}ring})(u) = -u)}}
\vnbprop{zpow-hom-unique-nat}{\vnbstmt{\begin{array}{@{}l@{}} \forall n \in \mathbb{N}, g.; \forall x \in \operatorname{carr}(g), q. \\ \quad (\operatorname{is\text{-}group}(g) \wedge (\operatorname{is\text{-}hom\text{-}group}(\operatorname{ring\text{-}additive\text{-}ag}(\operatorname{zz\text{-}ring}), g, q) \wedge (q(1) = x))) \Rightarrow  \\ \quad \quad q(n) = \operatorname{group\text{-}zpow}(g, x, n) \end{array}}}
\vnbprop{zz-cyclic-universal}{\vnbstmt{\begin{array}{@{}l@{}} \forall g.; \forall x \in \operatorname{carr}(g). \\ \quad \operatorname{is\text{-}group}(g) \Rightarrow  \\ \quad \quad \exists i. \\ \quad \quad \quad \operatorname{is\text{-}hom\text{-}group}(\operatorname{ring\text{-}additive\text{-}ag}(\operatorname{zz\text{-}ring}), g, i) \wedge \\ \quad \quad \quad i(1) = x \wedge \\ \quad \quad \quad \forall f. \\ \quad \quad \quad \quad (\operatorname{is\text{-}hom\text{-}group}(\operatorname{ring\text{-}additive\text{-}ag}(\operatorname{zz\text{-}ring}), g, f) \wedge (f(1) = x)) \Rightarrow (f = i) \end{array}}}
\vnbprop{group-left-cancel-id}{\vnbstmt{\begin{array}{@{}l@{}} \forall g.; \forall p, q \in \operatorname{carr}(g). \\ \quad (\operatorname{is\text{-}group}(g) \wedge (p = \operatorname{opr}(g)(p, q))) \Rightarrow (q = \operatorname{iden}(g)) \end{array}}}
\vnbprop{gpow-difference-period}{\vnbstmt{\begin{array}{@{}l@{}} \forall a, b \in \mathbb{N}, g.; \forall x \in \operatorname{carr}(g). \\ \quad (\operatorname{is\text{-}group}(g) \wedge ((a \leq b) \wedge (\operatorname{gpow}(g, x, a) = \operatorname{gpow}(g, x, b)))) \Rightarrow  \\ \quad \quad \operatorname{gpow}(g, x, (b - a)) = \operatorname{iden}(g) \end{array}}}
\vnbprop{gpow-period-exists}{\vnbstmt{\begin{array}{@{}l@{}} \forall g.; \forall x \in \operatorname{carr}(g). \\ \quad (\operatorname{is\text{-}group}(g) \wedge (|\operatorname{carr}(g)| \in \mathbb{N})) \Rightarrow  \\ \quad \quad \exists d \in \mathbb{N}.\; ((1 \leq d) \wedge (\operatorname{gpow}(g, x, d) = \operatorname{iden}(g))) \end{array}}}
\vnbprop{element-order-exists}{\vnbstmt{\begin{array}{@{}l@{}} \forall g.; \forall x \in \operatorname{carr}(g). \\ \quad (\operatorname{is\text{-}group}(g) \wedge (|\operatorname{carr}(g)| \in \mathbb{N})) \Rightarrow  \\ \quad \quad \exists o.\; \operatorname{is\text{-}element\text{-}order}(g, x, o) \end{array}}}
\vnbprop{zpow-multiple-period-nat}{\vnbstmt{\begin{array}{@{}l@{}} \forall c \in \mathbb{N}, g.; \forall x \in \operatorname{carr}(g), d \in \mathbb{N}. \\ \quad (\operatorname{is\text{-}group}(g) \wedge (\operatorname{gpow}(g, x, d) = \operatorname{iden}(g))) \Rightarrow  \\ \quad \quad \operatorname{group\text{-}zpow}(g, x, (d \cdot c)) = \operatorname{iden}(g) \end{array}}}
\vnbprop{zpow-multiple-period}{\vnbstmt{\begin{array}{@{}l@{}} \forall c \in \mathbb{Z}, g.; \forall x \in \operatorname{carr}(g), d \in \mathbb{N}. \\ \quad (\operatorname{is\text{-}group}(g) \wedge (\operatorname{gpow}(g, x, d) = \operatorname{iden}(g))) \Rightarrow  \\ \quad \quad \operatorname{group\text{-}zpow}(g, x, (d \cdot c)) = \operatorname{iden}(g) \end{array}}}
\vnbprop{element-order-divides-nat}{\vnbstmt{\begin{array}{@{}l@{}} \forall k \in \mathbb{N}, g.; \forall x \in \operatorname{carr}(g), d. \\ \quad (\operatorname{is\text{-}group}(g) \wedge (\operatorname{is\text{-}element\text{-}order}(g, x, d) \wedge (\operatorname{group\text{-}zpow}(g, x, k) = \operatorname{iden}(g)))) \Rightarrow  \\ \quad \quad \operatorname{zz\text{-}divides}(d, k) \end{array}}}
\vnbprop{element-order-divides}{\vnbstmt{\begin{array}{@{}l@{}} \forall k \in \mathbb{Z}, g.; \forall x \in \operatorname{carr}(g), d. \\ \quad (\operatorname{is\text{-}group}(g) \wedge \operatorname{is\text{-}element\text{-}order}(g, x, d)) \Rightarrow  \\ \quad \quad (\operatorname{group\text{-}zpow}(g, x, k) = \operatorname{iden}(g)) \Leftrightarrow \operatorname{zz\text{-}divides}(d, k) \end{array}}}
\vnbprop{generator-order-card}{\vnbstmt{\begin{array}{@{}l@{}} \forall g.; \forall x \in \operatorname{carr}(g), d. \\ \quad (\operatorname{is\text{-}group}(g) \wedge ((|\operatorname{carr}(g)| \in \mathbb{N}) \wedge (\operatorname{is\text{-}generator}(g, x) \wedge \operatorname{is\text{-}element\text{-}order}(g, x, d)))) \Rightarrow  \\ \quad \quad |\operatorname{carr}(g)| = d \end{array}}}
\vnbprop{group-op-right-id}{\vnbstmt{\forall s.; \forall a \in \operatorname{carr}(s).\; (\operatorname{is\text{-}group}(s) \Rightarrow (\operatorname{opr}(s)(a, \operatorname{iden}(s)) = a))}}
\vnbprop{group-op-right-inv}{\vnbstmt{\begin{array}{@{}l@{}} \forall s.; \forall a \in \operatorname{carr}(s). \\ \quad \operatorname{is\text{-}group}(s) \Rightarrow (\operatorname{opr}(s)(a, \operatorname{inv}(s)(a)) = \operatorname{iden}(s)) \end{array}}}
\vnbprop{group-prod-is-group}{\vnbstmt{\forall x, y.\; ((\operatorname{is\text{-}group}(x) \wedge \operatorname{is\text{-}group}(y)) \Rightarrow \operatorname{is\text{-}group}(\operatorname{group\text{-}prod}(x, y)))}}
\vnbprop{group-prod-carr}{\vnbstmt{\begin{array}{@{}l@{}} \forall x, y. \\ \quad (\operatorname{is\text{-}group}(x) \wedge \operatorname{is\text{-}group}(y)) \Rightarrow  \\ \quad \quad (\operatorname{carr}(\operatorname{group\text{-}prod}(x, y)) \simeq (\operatorname{carr}(x) \times \operatorname{carr}(y))) \end{array}}}
\vnbprop{group-prod-opr}{\vnbstmt{\begin{array}{@{}l@{}} \forall x, y. \\ \quad (\operatorname{is\text{-}group}(x) \wedge \operatorname{is\text{-}group}(y)) \Rightarrow  \\ \quad \quad (\operatorname{opr}(\operatorname{group\text{-}prod}(x, y)) \simeq (\lambda {[}\operatorname{gpp\_}, \operatorname{gpq\_}] \in ((\operatorname{carr}(x) \times \operatorname{carr}(y)) \times (\operatorname{carr}(x) \times \operatorname{carr}(y))).\; {[}\operatorname{opr}(x)(\operatorname{nth}(1, \operatorname{gpp\_}), \operatorname{nth}(1, \operatorname{gpq\_})), \operatorname{opr}(y)(\operatorname{nth}(2, \operatorname{gpp\_}), \operatorname{nth}(2, \operatorname{gpq\_}))])) \end{array}}}
\vnbprop{group-prod-iden}{\vnbstmt{\begin{array}{@{}l@{}} \forall x, y. \\ \quad (\operatorname{is\text{-}group}(x) \wedge \operatorname{is\text{-}group}(y)) \Rightarrow (\operatorname{iden}(\operatorname{group\text{-}prod}(x, y)) \simeq {[}\operatorname{iden}(x), \operatorname{iden}(y)]) \end{array}}}
\vnbprop{group-prod-inv}{\vnbstmt{\begin{array}{@{}l@{}} \forall x, y. \\ \quad (\operatorname{is\text{-}group}(x) \wedge \operatorname{is\text{-}group}(y)) \Rightarrow  \\ \quad \quad (\operatorname{inv}(\operatorname{group\text{-}prod}(x, y)) \simeq (\lambda p \in (\operatorname{carr}(x) \times \operatorname{carr}(y)).\; {[}\operatorname{inv}(x)(\operatorname{nth}(1, p)), \operatorname{inv}(y)(\operatorname{nth}(2, p))])) \end{array}}}
\vnbprop{group-prod-opr-pair}{\vnbstmt{\begin{array}{@{}l@{}} \forall x, y, a, b, c, d. \\ \quad (\operatorname{is\text{-}group}(x) \wedge (\operatorname{is\text{-}group}(y) \wedge ((a \in \operatorname{carr}(x)) \wedge ((b \in \operatorname{carr}(y)) \wedge ((c \in \operatorname{carr}(x)) \wedge (d \in \operatorname{carr}(y))))))) \Rightarrow  \\ \quad \quad \operatorname{opr}(\operatorname{group\text{-}prod}(x, y))({[}a, b], {[}c, d]) = {[}\operatorname{opr}(x)(a, c), \operatorname{opr}(y)(b, d)] \end{array}}}
\vnbprop{group-prod-inv-pair}{\vnbstmt{\begin{array}{@{}l@{}} \forall x, y, a, b. \\ \quad (\operatorname{is\text{-}group}(x) \wedge (\operatorname{is\text{-}group}(y) \wedge ((a \in \operatorname{carr}(x)) \wedge (b \in \operatorname{carr}(y))))) \Rightarrow  \\ \quad \quad \operatorname{inv}(\operatorname{group\text{-}prod}(x, y))({[}a, b]) = {[}\operatorname{inv}(x)(a), \operatorname{inv}(y)(b)] \end{array}}}
\vnbprop{gpow-prod}{\vnbstmt{\begin{array}{@{}l@{}} \forall n \in \mathbb{N}, x, y, a, b. \\ \quad (\operatorname{is\text{-}group}(x) \wedge (\operatorname{is\text{-}group}(y) \wedge ((a \in \operatorname{carr}(x)) \wedge (b \in \operatorname{carr}(y))))) \Rightarrow  \\ \quad \quad \operatorname{gpow}(\operatorname{group\text{-}prod}(x, y), {[}a, b], n) = {[}\operatorname{gpow}(x, a, n), \operatorname{gpow}(y, b, n)] \end{array}}}
\vnbprop{zpow-prod}{\vnbstmt{\begin{array}{@{}l@{}} \forall k \in \mathbb{Z}, x, y, a, b. \\ \quad (\operatorname{is\text{-}group}(x) \wedge (\operatorname{is\text{-}group}(y) \wedge ((a \in \operatorname{carr}(x)) \wedge (b \in \operatorname{carr}(y))))) \Rightarrow  \\ \quad \quad \operatorname{group\text{-}zpow}(\operatorname{group\text{-}prod}(x, y), {[}a, b], k) = {[}\operatorname{group\text{-}zpow}(x, a, k), \operatorname{group\text{-}zpow}(y, b, k)] \end{array}}}

\vnbbookfile{theorem-library/ccr4c-finite-char}{theorem-library/ccr4c-finite-char}
\noindent{\small ccr4c-finite-char.scm -- wave of 2026-11-09/10 (agent CCR-1109-C4), wired 2026-11-10 from scratchpad/ccr1109c4/ccr4c-finite-char.scm. Library: algebra.}\par
\vnbprop{ag-zz-act-gpow}{\vnbstmt{\begin{array}{@{}l@{}} \forall n \in \mathbb{N}, g.; \forall b \in \operatorname{carr}(g). \\ \quad \operatorname{is\text{-}abelian\text{-}group}(g) \Rightarrow (\operatorname{zz\text{-}act}(g, n, b) = \operatorname{gpow}(g, b, n)) \end{array}}}
\vnbprop{ag-zz-act-period}{\vnbstmt{\begin{array}{@{}l@{}} \forall g.; \forall b \in \operatorname{carr}(g). \\ \quad (\operatorname{is\text{-}abelian\text{-}group}(g) \wedge (|\operatorname{carr}(g)| \in \mathbb{N})) \Rightarrow  \\ \quad \quad \exists n \in \mathbb{N}.\; ((1 \leq n) \wedge (\operatorname{zz\text{-}act}(g, n, b) = \operatorname{iden}(g))) \end{array}}}
\vnbprop{zz-act-one-mul}{\vnbstmt{\begin{array}{@{}l@{}} \forall j \in \mathbb{N}, r.; \forall k \in \mathbb{Z}. \\ \quad \operatorname{is\text{-}commutative\text{-}ring}(r) \Rightarrow  \\ \quad \quad \operatorname{zz\text{-}act}(\operatorname{ring\text{-}additive\text{-}ag}(r), (j \cdot k), \operatorname{one}(r)) \\ \quad \quad =\; \operatorname{mul}(r)(\operatorname{zz\text{-}act}(\operatorname{ring\text{-}additive\text{-}ag}(r), k, \operatorname{one}(r)), \operatorname{zz\text{-}act}(\operatorname{ring\text{-}additive\text{-}ag}(r), j, \operatorname{one}(r))) \end{array}}}
\vnbprop{zz-act-one-zero}{\vnbstmt{\forall r.\; (\operatorname{is\text{-}ring}(r) \Rightarrow (\operatorname{zz\text{-}act}(\operatorname{ring\text{-}additive\text{-}ag}(r), 0, \operatorname{one}(r)) = \operatorname{zero}(r)))}}
\vnbprop{zz-act-one-succ}{\vnbstmt{\begin{array}{@{}l@{}} \forall n \in \mathbb{N}, r. \\ \quad \operatorname{is\text{-}ring}(r) \Rightarrow  \\ \quad \quad \operatorname{zz\text{-}act}(\operatorname{ring\text{-}additive\text{-}ag}(r), (n + 1), \operatorname{one}(r)) \\ \quad \quad =\; \operatorname{add}(r)(\operatorname{zz\text{-}act}(\operatorname{ring\text{-}additive\text{-}ag}(r), n, \operatorname{one}(r)), \operatorname{one}(r)) \end{array}}}
\vnbprop{finite-domain-characteristic}{\vnbstmt{\begin{array}{@{}l@{}} \forall r. \\ \quad (\operatorname{is\text{-}integral\text{-}domain}(r) \wedge (|\operatorname{carr}(r)| \in \mathbb{N})) \Rightarrow  \\ \quad \quad \exists p.\; \operatorname{has\text{-}characteristic}(r, p) \end{array}}}
\vnbprop{finite-field-characteristic}{\vnbstmt{\begin{array}{@{}l@{}} \forall r. \\ \quad (\operatorname{is\text{-}field\text{-}ring}(r) \wedge (|\operatorname{carr}(r)| \in \mathbb{N})) \Rightarrow  \\ \quad \quad \exists p.\; \operatorname{has\text{-}characteristic}(r, p) \end{array}}}

\vnbbookfile{theorem-library/group-cosets}{theorem-library/group-cosets}
\noindent{\small theorem-library/group-cosets.scm -- COSETS AND LAGRANGE (ALGEBRAIC-NUMBERS notes, Prop. 1.28, Cor. 1.29; ALGNUM, bundle B5 / galois-theory B7, 2026-10-26; rebuilt 2026-10-27 without group-right-inv-id, which tb11/tg-group.scm proves with the same statement).  Purely algebraic. Vocabulary: IS-SUBGROUP ...}\par
\vnbprop{left-coset-unfold}{\vnbstmt{\begin{array}{@{}l@{}} \forall g, a, h. \\ \quad (\operatorname{left\text{-}coset}(g, a, h) \simeq \{x \in \operatorname{carr}(g) : \exists y \in h.\; (x = \operatorname{opr}(g)(a, y))\}) \end{array}}}
\vnbprop{left-cosets-unfold}{\vnbstmt{\begin{array}{@{}l@{}} \forall g, h. \\ \quad (\operatorname{left\text{-}cosets}(g, h) \simeq \{c \in \operatorname{power}(\operatorname{carr}(g)) : \exists a \in \operatorname{carr}(g).\; (c = \operatorname{left\text{-}coset}(g, a, h))\}) \end{array}}}
\vnbprop{group-inv-cancel}{\vnbstmt{\begin{array}{@{}l@{}} \forall g.; \forall a, b \in \operatorname{carr}(g). \\ \quad \operatorname{is\text{-}group}(g) \Rightarrow  \\ \quad \quad ((\operatorname{opr}(g)(\operatorname{inv}(g)(a), \operatorname{opr}(g)(a, b)) = b) \wedge (\operatorname{opr}(g)(a, \operatorname{opr}(g)(\operatorname{inv}(g)(a), b)) = b)) \end{array}}}
\vnbprop{left-coset-member}{\vnbstmt{\begin{array}{@{}l@{}} \forall g, a, h, y. \\ \quad (\operatorname{is\text{-}group}(g) \wedge ((a \in \operatorname{carr}(g)) \wedge ((h \subseteq \operatorname{carr}(g)) \wedge (y \in h)))) \Rightarrow  \\ \quad \quad \operatorname{opr}(g)(a, y) \in \operatorname{left\text{-}coset}(g, a, h) \end{array}}}
\vnbprop{left-coset-eq}{\vnbstmt{\begin{array}{@{}l@{}} \forall g, h, a, b. \\ \quad (\operatorname{is\text{-}group}(g) \wedge (\operatorname{is\text{-}subgroup}(g, h) \wedge ((a \in \operatorname{carr}(g)) \wedge (b \in \operatorname{left\text{-}coset}(g, a, h))))) \Rightarrow  \\ \quad \quad \operatorname{left\text{-}coset}(g, b, h) = \operatorname{left\text{-}coset}(g, a, h) \end{array}}}
\vnbprop{lagrange}{\vnbstmt{\begin{array}{@{}l@{}} \forall g, h. \\ \quad (\operatorname{is\text{-}group}(g) \wedge (\operatorname{is\text{-}subgroup}(g, h) \wedge (|\operatorname{carr}(g)| \in \mathbb{N}))) \Rightarrow  \\ \quad \quad |\operatorname{carr}(g)| = (|h| \cdot |\operatorname{left\text{-}cosets}(g, h)|) \end{array}}}
\vnbprop{generated-subgroup-unfold}{\vnbstmt{\begin{array}{@{}l@{}} \forall g, x. \\ \quad (\operatorname{generated\text{-}subgroup}(g, x) \simeq \{y \in \operatorname{carr}(g) : \forall s \in \operatorname{power}(\operatorname{carr}(g)).\; ((\operatorname{is\text{-}subgroup}(g, s) \wedge (x \in s)) \Rightarrow (y \in s))\}) \end{array}}}
\vnbprop{generated-subgroup-is-subgroup}{\vnbstmt{\begin{array}{@{}l@{}} \forall g, x. \\ \quad (\operatorname{is\text{-}group}(g) \wedge (x \in \operatorname{carr}(g))) \Rightarrow  \\ \quad \quad \operatorname{is\text{-}subgroup}(g, \operatorname{generated\text{-}subgroup}(g, x)) \wedge (x \in \operatorname{generated\text{-}subgroup}(g, x)) \end{array}}}
\vnbprop{lagrange-element-order}{\vnbstmt{\begin{array}{@{}l@{}} \forall g, x. \\ \quad (\operatorname{is\text{-}group}(g) \wedge ((|\operatorname{carr}(g)| \in \mathbb{N}) \wedge (x \in \operatorname{carr}(g)))) \Rightarrow  \\ \quad \quad \operatorname{zz\text{-}divides}(|\operatorname{generated\text{-}subgroup}(g, x)|, |\operatorname{carr}(g)|) \end{array}}}

\vnbbookfile{theorem-library/cyclic-product}{theorem-library/cyclic-product}
\noindent{\small theorem-library/cyclic-product.scm -- PROPS. 1.38 and 1.39 (two groups) (ALGEBRAIC-NUMBERS notes, ch. 1; ALGNUM, 2026-10-27). cyclic-product-generator   1.39 for two groups: (x1, x2) generates G1 x G2 iff card G1, card G2 are prime to each other (x1, x2 generators of the finite groups G1, G2) zpow-one, zpow-zpow(-nat)  ...}\par
\vnbprop{cyclic-product-generator}{\vnbstmt{\begin{array}{@{}l@{}} \forall a, b, c, d. \\ \quad \operatorname{is\text{-}group}(a) \wedge \\ \quad \operatorname{is\text{-}group}(b) \wedge \\ \quad |\operatorname{carr}(a)| \in \mathbb{N} \wedge \\ \quad |\operatorname{carr}(b)| \in \mathbb{N} \wedge \\ \quad \operatorname{is\text{-}generator}(a, c) \wedge \\ \quad \operatorname{is\text{-}generator}(b, d) \Rightarrow \\ \quad \quad \operatorname{is\text{-}generator}(\operatorname{group\text{-}prod}(a, b), {[}c, d]) \Leftrightarrow \operatorname{zz\text{-}coprime}(|\operatorname{carr}(a)|, |\operatorname{carr}(b)|) \end{array}}}
\vnbprop{zpow-one}{\vnbstmt{\forall g.; \forall x \in \operatorname{carr}(g).\; (\operatorname{is\text{-}group}(g) \Rightarrow (\operatorname{group\text{-}zpow}(g, x, 1) = x))}}
\vnbprop{zpow-zpow-nat}{\vnbstmt{\begin{array}{@{}l@{}} \forall m \in \mathbb{N}, k \in \mathbb{Z}, g.; \forall x \in \operatorname{carr}(g). \\ \quad \operatorname{is\text{-}group}(g) \Rightarrow  \\ \quad \quad \operatorname{group\text{-}zpow}(g, \operatorname{group\text{-}zpow}(g, x, k), m) = \operatorname{group\text{-}zpow}(g, x, (k \cdot m)) \end{array}}}
\vnbprop{zpow-zpow}{\vnbstmt{\begin{array}{@{}l@{}} \forall m, k \in \mathbb{Z}, g.; \forall x \in \operatorname{carr}(g). \\ \quad \operatorname{is\text{-}group}(g) \Rightarrow  \\ \quad \quad \operatorname{group\text{-}zpow}(g, \operatorname{group\text{-}zpow}(g, x, k), m) = \operatorname{group\text{-}zpow}(g, x, (k \cdot m)) \end{array}}}
\vnbprop{generator-power-iff}{\vnbstmt{\begin{array}{@{}l@{}} \forall g.; \forall x \in \operatorname{carr}(g), n.; \forall k \in \mathbb{Z}. \\ \quad (\operatorname{is\text{-}group}(g) \wedge (\operatorname{is\text{-}element\text{-}order}(g, x, n) \wedge \operatorname{is\text{-}generator}(g, x))) \Rightarrow  \\ \quad \quad \operatorname{is\text{-}generator}(g, \operatorname{group\text{-}zpow}(g, x, k)) \Leftrightarrow \operatorname{zz\text{-}coprime}(k, n) \end{array}}}
\vnbprop{zpow-residue}{\vnbstmt{\begin{array}{@{}l@{}} \forall g.; \forall x \in \operatorname{carr}(g), d.; \forall k \in \mathbb{Z}. \\ \quad (\operatorname{is\text{-}group}(g) \wedge \operatorname{is\text{-}element\text{-}order}(g, x, d)) \Rightarrow  \\ \quad \quad \exists r \in \mathbb{N}.\; ((1 \leq r) \wedge ((r \leq d) \wedge (\operatorname{group\text{-}zpow}(g, x, k) = \operatorname{gpow}(g, x, r)))) \end{array}}}
\vnbprop{gpow-injective-below-order}{\vnbstmt{\begin{array}{@{}l@{}} \forall g.; \forall x \in \operatorname{carr}(g), d, a, b. \\ \quad \operatorname{is\text{-}group}(g) \wedge \\ \quad \operatorname{is\text{-}element\text{-}order}(g, x, d) \wedge \\ \quad a \in \mathbb{N} \wedge \\ \quad b \in \mathbb{N} \wedge \\ \quad 1 \leq a \wedge \\ \quad a \leq d \wedge \\ \quad 1 \leq b \wedge \\ \quad b \leq d \wedge \\ \quad \operatorname{gpow}(g, x, a) = \operatorname{gpow}(g, x, b) \Rightarrow \\ \quad \quad a = b \end{array}}}
\vnbprop{generator-count}{\vnbstmt{\begin{array}{@{}l@{}} \forall g.; \forall x \in \operatorname{carr}(g). \\ \quad (\operatorname{is\text{-}group}(g) \wedge ((|\operatorname{carr}(g)| \in \mathbb{N}) \wedge \operatorname{is\text{-}generator}(g, x))) \Rightarrow  \\ \quad \quad |\{y \in \operatorname{carr}(g) : \operatorname{is\text{-}generator}(g, y)\}| = \operatorname{euler\text{-}phi}(|\operatorname{carr}(g)|) \end{array}}}

\vnbbookfile{theorem-library/group-family-product}{theorem-library/group-family-product}
\noindent{\small theorem-library/group-family-product.scm -- THE PRODUCT OF A FINITE FAMILY OF GROUPS and PROP. 1.39 for m groups (ALGEBRAIC-NUMBERS notes, ch. 1; ALGNUM, 2026-10-27).  G\_1 x ... x G\_m is GROUP-FAMPROD(gf, m), the iterated GROUP-PROD from TRIV-GROUP; (x\_1, ..., x\_m) is GROUP-FAMTUPLE(xf, m) (group-zpow-vocab). ...}\par
\vnbprop{triv-group-is-group}{\vnbstmt{\operatorname{is\text{-}group}(\operatorname{triv\text{-}group})}}
\vnbprop{triv-group-carr}{\vnbstmt{(\operatorname{carr}(\operatorname{triv\text{-}group}) \simeq \{0\})}}
\vnbprop{triv-group-iden}{\vnbstmt{(\operatorname{iden}(\operatorname{triv\text{-}group}) \simeq 0)}}
\vnbprop{triv-group-card}{\vnbstmt{|\operatorname{carr}(\operatorname{triv\text{-}group})| = 1}}
\vnbprop{triv-group-generator}{\vnbstmt{\operatorname{is\text{-}generator}(\operatorname{triv\text{-}group}, 0)}}
\vnbprop{group-famprod-facts}{\vnbstmt{\begin{array}{@{}l@{}} \forall m \in \mathbb{N}, f, a. \\ \quad \forall i \in \mathbb{N}. \\ \quad \quad ((1 \leq i) \wedge (i \leq m)) \Rightarrow  \\ \quad \quad \quad \operatorname{is\text{-}group}(f(i)) \wedge ((|\operatorname{carr}(f(i))| \in \mathbb{N}) \wedge \operatorname{is\text{-}generator}(f(i), a(i))) \Rightarrow \\ \quad \quad \operatorname{is\text{-}group}(\operatorname{group\text{-}famprod}(f, m)) \wedge \\ \quad \quad |\operatorname{carr}(\operatorname{group\text{-}famprod}(f, m))| \in \mathbb{N} \wedge \\ \quad \quad \operatorname{group\text{-}famtuple}(a, m) \in \operatorname{carr}(\operatorname{group\text{-}famprod}(f, m)) \end{array}}}
\vnbprop{group-prod-generator-left}{\vnbstmt{\begin{array}{@{}l@{}} \forall x, y, u, w. \\ \quad \operatorname{is\text{-}group}(x) \wedge \\ \quad \operatorname{is\text{-}group}(y) \wedge \\ \quad u \in \operatorname{carr}(x) \wedge \\ \quad w \in \operatorname{carr}(y) \wedge \\ \quad \operatorname{is\text{-}generator}(\operatorname{group\text{-}prod}(x, y), {[}u, w]) \Rightarrow \\ \quad \quad \operatorname{is\text{-}generator}(x, u) \end{array}}}
\vnbprop{group-famprod-card-divisible}{\vnbstmt{\begin{array}{@{}l@{}} \forall m \in \mathbb{N}, f, a.; \forall i \in \mathbb{N}. \\ \quad \forall \mathit{gi} \in \mathbb{N}. \\ \quad \quad ((1 \leq \mathit{gi}) \wedge (\mathit{gi} \leq m)) \Rightarrow  \\ \quad \quad \quad \operatorname{is\text{-}group}(f(\mathit{gi})) \wedge ((|\operatorname{carr}(f(\mathit{gi}))| \in \mathbb{N}) \wedge \operatorname{is\text{-}generator}(f(\mathit{gi}), a(\mathit{gi}))) \wedge \\ \quad 1 \leq i \wedge \\ \quad i \leq m \Rightarrow \\ \quad \quad \operatorname{zz\text{-}divides}(|\operatorname{carr}(f(i))|, |\operatorname{carr}(\operatorname{group\text{-}famprod}(f, m))|) \end{array}}}
\vnbprop{group-famprod-card-coprime}{\vnbstmt{\begin{array}{@{}l@{}} \forall m \in \mathbb{N}, f, a.; \forall b \in \mathbb{Z}. \\ \quad \forall i \in \mathbb{N}. \\ \quad \quad ((1 \leq i) \wedge (i \leq m)) \Rightarrow  \\ \quad \quad \quad \operatorname{is\text{-}group}(f(i)) \wedge ((|\operatorname{carr}(f(i))| \in \mathbb{N}) \wedge \operatorname{is\text{-}generator}(f(i), a(i))) \wedge \\ \quad \forall \mathit{ci} \in \mathbb{N}.\; (((1 \leq \mathit{ci}) \wedge (\mathit{ci} \leq m)) \Rightarrow \operatorname{zz\text{-}coprime}(|\operatorname{carr}(f(\mathit{ci}))|, b)) \Rightarrow \\ \quad \quad \operatorname{zz\text{-}coprime}(|\operatorname{carr}(\operatorname{group\text{-}famprod}(f, m))|, b) \end{array}}}
\vnbprop{cyclic-family-product-generator}{\vnbstmt{\begin{array}{@{}l@{}} \forall m \in \mathbb{N}, f, a. \\ \quad \forall i \in \mathbb{N}. \\ \quad \quad ((1 \leq i) \wedge (i \leq m)) \Rightarrow  \\ \quad \quad \quad \operatorname{is\text{-}group}(f(i)) \wedge ((|\operatorname{carr}(f(i))| \in \mathbb{N}) \wedge \operatorname{is\text{-}generator}(f(i), a(i))) \Rightarrow \\ \quad \quad \operatorname{is\text{-}generator}(\operatorname{group\text{-}famprod}(f, m), \operatorname{group\text{-}famtuple}(a, m)) \\ \quad \quad \Leftrightarrow\; \forall b, j \in \mathbb{N}. \\ \quad \quad \quad ((1 \leq b) \wedge ((b \leq m) \wedge ((1 \leq j) \wedge ((j \leq m) \wedge \neg (b = j))))) \Rightarrow  \\ \quad \quad \quad \quad \operatorname{zz\text{-}coprime}(|\operatorname{carr}(f(b))|, |\operatorname{carr}(f(j))|) \end{array}}}

\vnbbookfile{structure-library/functoriality}{structure-library/functoriality}
\noindent{\small functoriality.scm -- every view-as is a FUNCTOR, and here is the proof.}\par
\vnbprop{abelian-group-as-monoid-functorial}{\vnbstmt{\begin{array}{@{}l@{}} \forall a, b, f. \\ \quad \operatorname{is\text{-}hom\text{-}abelian\text{-}group}(a, b, f) \Rightarrow  \\ \quad \quad \operatorname{is\text{-}hom\text{-}monoid}(\operatorname{abelian\text{-}group\text{-}as\text{-}monoid}(a), \operatorname{abelian\text{-}group\text{-}as\text{-}monoid}(b), f) \end{array}}}
\vnbprop{commutative-ring-additive-ag-functorial}{\vnbstmt{\begin{array}{@{}l@{}} \forall a, b, f. \\ \quad \operatorname{is\text{-}hom\text{-}commutative\text{-}ring}(a, b, f) \Rightarrow  \\ \quad \quad \operatorname{is\text{-}hom\text{-}abelian\text{-}group}(\operatorname{commutative\text{-}ring\text{-}additive\text{-}ag}(a), \operatorname{commutative\text{-}ring\text{-}additive\text{-}ag}(b), f) \end{array}}}
\vnbprop{commutative-ring-multiplicative-cm-functorial}{\vnbstmt{\begin{array}{@{}l@{}} \forall a, b, f. \\ \quad \operatorname{is\text{-}hom\text{-}commutative\text{-}ring}(a, b, f) \Rightarrow  \\ \quad \quad \operatorname{is\text{-}hom\text{-}comm\text{-}monoid}(\operatorname{commutative\text{-}ring\text{-}multiplicative\text{-}cm}(a), \operatorname{commutative\text{-}ring\text{-}multiplicative\text{-}cm}(b), f) \end{array}}}
\vnbprop{field-additive-ag-functorial}{\vnbstmt{\begin{array}{@{}l@{}} \forall a, b, f. \\ \quad \operatorname{is\text{-}hom\text{-}field}(a, b, f) \Rightarrow  \\ \quad \quad \operatorname{is\text{-}hom\text{-}abelian\text{-}group}(\operatorname{field\text{-}additive\text{-}ag}(a), \operatorname{field\text{-}additive\text{-}ag}(b), f) \end{array}}}
\vnbprop{field-as-euclidean-ring-functorial}{\vnbstmt{\begin{array}{@{}l@{}} \forall a, b, f. \\ \quad \operatorname{is\text{-}hom\text{-}field}(a, b, f) \Rightarrow  \\ \quad \quad \operatorname{is\text{-}hom\text{-}euclidean\text{-}ring}(\operatorname{field\text{-}as\text{-}euclidean\text{-}ring}(a), \operatorname{field\text{-}as\text{-}euclidean\text{-}ring}(b), f) \end{array}}}
\vnbprop{field-as-integral-domain-functorial}{\vnbstmt{\begin{array}{@{}l@{}} \forall a, b, f. \\ \quad \operatorname{is\text{-}hom\text{-}field}(a, b, f) \Rightarrow  \\ \quad \quad \operatorname{is\text{-}hom\text{-}integral\text{-}domain}(\operatorname{field\text{-}as\text{-}integral\text{-}domain}(a), \operatorname{field\text{-}as\text{-}integral\text{-}domain}(b), f) \end{array}}}
\vnbprop{ring-additive-ag-functorial}{\vnbstmt{\begin{array}{@{}l@{}} \forall a, b, f. \\ \quad \operatorname{is\text{-}hom\text{-}ring}(a, b, f) \Rightarrow  \\ \quad \quad \operatorname{is\text{-}hom\text{-}abelian\text{-}group}(\operatorname{ring\text{-}additive\text{-}ag}(a), \operatorname{ring\text{-}additive\text{-}ag}(b), f) \end{array}}}
\vnbprop{ring-multiplicative-monoid-functorial}{\vnbstmt{\begin{array}{@{}l@{}} \forall a, b, f. \\ \quad \operatorname{is\text{-}hom\text{-}ring}(a, b, f) \Rightarrow  \\ \quad \quad \operatorname{is\text{-}hom\text{-}monoid}(\operatorname{ring\text{-}multiplicative\text{-}monoid}(a), \operatorname{ring\text{-}multiplicative\text{-}monoid}(b), f) \end{array}}}

\part{Linear Algebra}
\noindent Modules and vector spaces; matrices, elementary operations, determinants and the Smith normal form; spans, submodules, bases and rank; finite-dimensional spaces.
\par\noindent\textit{Builds on:} Algebra.
\par\noindent{\small 122 files.}

\vnbbookfile{structure-library/module}{structure-library/module}
\noindent{\small module.scm -- MODULE over a ring, with the scalar action ACT.}\par
\vnbstructdef{module}{A module $s$ is a 6-tuple whose components are listed below.\par
\begin{enumerate}[nosep]
\item $\operatorname{scal}(s) = \operatorname{nth}(1, s)$, a substructure (\texttt{substructure}): $\operatorname{is\text{-}ring}(\operatorname{scal}(s))$.
\item $\operatorname{vec}(s) = \operatorname{nth}(2, s)$, a carrier (\texttt{carriers}): $\operatorname{vec}(s) \in \mathbf{Set}$.
\item $\operatorname{vadd}(s) = \operatorname{nth}(3, s)$, an operation (\texttt{op}): $\operatorname{vadd}(s) \in ((\operatorname{vec}(s) \times \operatorname{vec}(s)) \to \operatorname{vec}(s))$.
\item $\operatorname{vzero}(s) = \operatorname{nth}(4, s)$, a constant (\texttt{constant}): $\operatorname{vzero}(s) \in \operatorname{vec}(s)$.
\item $\operatorname{vneg}(s) = \operatorname{nth}(5, s)$, an operation (\texttt{op}): $\operatorname{vneg}(s) \in (\operatorname{vec}(s) \to \operatorname{vec}(s))$.
\item $\operatorname{act}(s) = \operatorname{nth}(6, s)$, an operation (\texttt{op}): $\operatorname{act}(s) \in ((\operatorname{carr}(\operatorname{scal}(s)) \times \operatorname{vec}(s)) \to \operatorname{vec}(s))$.
\end{enumerate}
The laws (the \texttt{property} and \texttt{law} clauses):
\begin{itemize}[nosep]
\item $\operatorname{is\text{-}associative}(\operatorname{vadd}(s), \operatorname{vec}(s))$
\item $\operatorname{is\text{-}commutative}(\operatorname{vadd}(s), \operatorname{vec}(s))$
\item $\operatorname{is\text{-}identity}(\operatorname{vadd}(s), \operatorname{vzero}(s), \operatorname{vec}(s))$
\item $\operatorname{has\text{-}inverses}(\operatorname{vadd}(s), \operatorname{vzero}(s), \operatorname{vneg}(s), \operatorname{vec}(s))$
\item
\begin{equation*}
\begin{array}{@{}l@{}} \forall \operatorname{r\_} \in \operatorname{carr}(\operatorname{scal}(s)), \operatorname{x\_}, \operatorname{y\_} \in \operatorname{vec}(s). \\ \quad \operatorname{act}(s)(\operatorname{r\_}, \operatorname{vadd}(s)(\operatorname{x\_}, \operatorname{y\_})) = \operatorname{vadd}(s)(\operatorname{act}(s)(\operatorname{r\_}, \operatorname{x\_}), \operatorname{act}(s)(\operatorname{r\_}, \operatorname{y\_})) \end{array}
\end{equation*}

\item
\begin{equation*}
\begin{array}{@{}l@{}} \forall \operatorname{r\_}, \operatorname{s\_} \in \operatorname{carr}(\operatorname{scal}(s)), \operatorname{x\_} \in \operatorname{vec}(s). \\ \quad \operatorname{act}(s)(\operatorname{add}(\operatorname{scal}(s))(\operatorname{r\_}, \operatorname{s\_}), \operatorname{x\_}) = \operatorname{vadd}(s)(\operatorname{act}(s)(\operatorname{r\_}, \operatorname{x\_}), \operatorname{act}(s)(\operatorname{s\_}, \operatorname{x\_})) \end{array}
\end{equation*}

\item
\begin{equation*}
\begin{array}{@{}l@{}} \forall \operatorname{r\_}, \operatorname{s\_} \in \operatorname{carr}(\operatorname{scal}(s)), \operatorname{x\_} \in \operatorname{vec}(s). \\ \quad \operatorname{act}(s)(\operatorname{mul}(\operatorname{scal}(s))(\operatorname{r\_}, \operatorname{s\_}), \operatorname{x\_}) = \operatorname{act}(s)(\operatorname{r\_}, \operatorname{act}(s)(\operatorname{s\_}, \operatorname{x\_})) \end{array}
\end{equation*}

\item $\forall \operatorname{x\_} \in \operatorname{vec}(s).\; (\operatorname{act}(s)(\operatorname{one}(\operatorname{scal}(s)), \operatorname{x\_}) = \operatorname{x\_})$
\end{itemize}
The characteristic law (the axiom \texttt{is-module}): $\operatorname{is\text{-}module}(s)$ holds if and only if $\operatorname{length}(s) = 6$, every component has the typing listed with it, every law holds, and $s \in \mathbf{Set}$ (the structure is a set).
}
\par\medskip\noindent\textbf{Definition} (\texttt{hom-module}).\;$\operatorname{hom\text{-}module}(a, b) :=$ $\{\operatorname{homf\_} \in (\operatorname{vec}(a) \to \operatorname{vec}(b)) : \operatorname{is\text{-}hom\text{-}module}(a, b, \operatorname{homf\_})\}$
\vnbdef{module-scalar-ring}{\vnbstmt{\forall m.\; (\operatorname{is\text{-}module}(m) \Rightarrow \operatorname{is\text{-}ring}(\operatorname{scal}(m)))}}
\vnbdef{module-scalar-zero-in}{\vnbstmt{\forall m.\; (\operatorname{is\text{-}module}(m) \Rightarrow (\operatorname{zero}(\operatorname{scal}(m)) \in \operatorname{carr}(\operatorname{scal}(m))))}}
\vnbdef{module-vzero-in}{\vnbstmt{\forall m.\; (\operatorname{is\text{-}module}(m) \Rightarrow (\operatorname{vzero}(m) \in \operatorname{vec}(m)))}}
\vnbdef{module-act-type}{\vnbstmt{\begin{array}{@{}l@{}} \forall m.; \forall r \in \operatorname{carr}(\operatorname{scal}(m)), x \in \operatorname{vec}(m). \\ \quad \operatorname{is\text{-}module}(m) \Rightarrow (\operatorname{act}(m)(r, x) \in \operatorname{vec}(m)) \end{array}}}
\vnbdef{module-vadd-type}{\vnbstmt{\begin{array}{@{}l@{}} \forall m.; \forall x, y \in \operatorname{vec}(m). \\ \quad \operatorname{is\text{-}module}(m) \Rightarrow (\operatorname{vadd}(m)(x, y) \in \operatorname{vec}(m)) \end{array}}}
\vnbdef{module-vneg-type}{\vnbstmt{\forall m.; \forall x \in \operatorname{vec}(m).\; (\operatorname{is\text{-}module}(m) \Rightarrow (\operatorname{vneg}(m)(x) \in \operatorname{vec}(m)))}}
\vnbdef{module-act-distrib-vec}{\vnbstmt{\begin{array}{@{}l@{}} \forall m.; \forall r \in \operatorname{carr}(\operatorname{scal}(m)), x, y \in \operatorname{vec}(m). \\ \quad \operatorname{is\text{-}module}(m) \Rightarrow  \\ \quad \quad \operatorname{act}(m)(r, \operatorname{vadd}(m)(x, y)) = \operatorname{vadd}(m)(\operatorname{act}(m)(r, x), \operatorname{act}(m)(r, y)) \end{array}}}
\vnbdef{module-act-distrib-scalar}{\vnbstmt{\begin{array}{@{}l@{}} \forall m.; \forall r, s \in \operatorname{carr}(\operatorname{scal}(m)), x \in \operatorname{vec}(m). \\ \quad \operatorname{is\text{-}module}(m) \Rightarrow  \\ \quad \quad \operatorname{act}(m)(\operatorname{add}(\operatorname{scal}(m))(r, s), x) = \operatorname{vadd}(m)(\operatorname{act}(m)(r, x), \operatorname{act}(m)(s, x)) \end{array}}}
\vnbdef{module-act-mul-compat}{\vnbstmt{\begin{array}{@{}l@{}} \forall m.; \forall r, s \in \operatorname{carr}(\operatorname{scal}(m)), x \in \operatorname{vec}(m). \\ \quad \operatorname{is\text{-}module}(m) \Rightarrow  \\ \quad \quad \operatorname{act}(m)(\operatorname{mul}(\operatorname{scal}(m))(r, s), x) = \operatorname{act}(m)(r, \operatorname{act}(m)(s, x)) \end{array}}}
\vnbdef{module-act-unital}{\vnbstmt{\forall m.; \forall x \in \operatorname{vec}(m).\; (\operatorname{is\text{-}module}(m) \Rightarrow (\operatorname{act}(m)(\operatorname{one}(\operatorname{scal}(m)), x) = x))}}
\vnbdef{is-hom-module-def}{\vnbstmt{\begin{array}{@{}l@{}} \forall a, b, f. \\ \quad \operatorname{is\text{-}hom\text{-}module}(a, b, f) \\ \quad \Leftrightarrow\; \operatorname{is\text{-}module}(a) \wedge \\ \quad \operatorname{is\text{-}module}(b) \wedge \\ \quad f \in (\operatorname{vec}(a) \to \operatorname{vec}(b)) \wedge \\ \quad \operatorname{scal}(a) = \operatorname{scal}(b) \wedge \\ \quad \forall x_{1}, x_{2} \in \operatorname{vec}(a). \\ \quad \quad f(\operatorname{vadd}(a)(x_{1}, x_{2})) = \operatorname{vadd}(b)(f(x_{1}), f(x_{2})) \wedge \\ \quad f(\operatorname{vzero}(a)) = \operatorname{vzero}(b) \wedge \\ \quad \forall x_{1} \in \operatorname{vec}(a).\; (f(\operatorname{vneg}(a)(x_{1})) = \operatorname{vneg}(b)(f(x_{1}))) \wedge \\ \quad \forall x_{1} \in \operatorname{carr}(\operatorname{scal}(a)), x_{2} \in \operatorname{vec}(a). \\ \quad \quad f(\operatorname{act}(a)(x_{1}, x_{2})) = \operatorname{act}(b)(x_{1}, f(x_{2})) \end{array}}}
\vnbdef{module-class}{\vnbstmt{\forall s.\; ((s \in \mathit{module}) \Leftrightarrow \operatorname{is\text{-}module}(s))}}

\vnbbookfile{structure-library/finite-dimensional}{structure-library/finite-dimensional}
\noindent{\small finite-dimensional.scm -- finite-dimensional vector spaces, Zorn-free.}\par
\par\medskip\noindent\textbf{Definition} (\texttt{is-finite-dimensional}).\;For every $m$, 
\begin{equation*}
\operatorname{is\text{-}finite\text{-}dimensional}(m) \Leftrightarrow (\operatorname{is\text{-}vector\text{-}space}(m) \wedge \operatorname{is\text{-}noetherian}(m))
\end{equation*}

\par\medskip\noindent\textbf{Definition} (\texttt{is-submodule}).\;For all $m$ and $s$, 
\begin{equation*}
\begin{array}{@{}l@{}} \operatorname{is\text{-}submodule}(m, s) \\ \Leftrightarrow\; s \subseteq \operatorname{vec}(m) \wedge \\ \operatorname{vzero}(m) \in s \wedge \\ \forall \operatorname{x\_}, \operatorname{y\_} \in s.\; (\operatorname{vadd}(m)(\operatorname{x\_}, \operatorname{y\_}) \in s) \wedge \\ \forall \operatorname{x\_} \in s.\; (\operatorname{vneg}(m)(\operatorname{x\_}) \in s) \wedge \\ \forall \operatorname{r\_} \in \operatorname{carr}(\operatorname{scal}(m)), \operatorname{x\_} \in s.\; (\operatorname{act}(m)(\operatorname{r\_}, \operatorname{x\_}) \in s) \end{array}
\end{equation*}

\par\medskip\noindent\textbf{Definition} (\texttt{is-noetherian}).\;For every $m$, 
\begin{equation*}
\begin{array}{@{}l@{}} \operatorname{is\text{-}noetherian}(m) \\ \Leftrightarrow\; \operatorname{is\text{-}module}(m) \wedge \\ \forall \operatorname{f\_} \in (\mathbb{N} \to \operatorname{power}(\operatorname{vec}(m))). \\ \quad (\forall \operatorname{n\_} \in \mathbb{N}.\; \operatorname{is\text{-}submodule}(m, \operatorname{f\_}(\operatorname{n\_})) \wedge \forall \operatorname{n\_} \in \mathbb{N}.\; (\operatorname{f\_}(\operatorname{n\_}) \subseteq \operatorname{f\_}(\operatorname{succ}(\operatorname{n\_})))) \Rightarrow  \\ \quad \quad \exists \operatorname{k\_} \in \mathbb{N}.\; \forall \operatorname{n\_}.\; (((\operatorname{n\_} \in \mathbb{N}) \wedge (\operatorname{k\_} \leq \operatorname{n\_})) \Rightarrow (\operatorname{f\_}(\operatorname{n\_}) = \operatorname{f\_}(\operatorname{k\_}))) \end{array}
\end{equation*}

\par\medskip\noindent\textbf{Definition} (\texttt{is-subspace}).\;For all $m$ and $s$,  $\operatorname{is\text{-}subspace}(m, s) \Leftrightarrow \operatorname{is\text{-}submodule}(m, s)$
\par\medskip\noindent\textbf{Definition} (\texttt{span-add-one}).\;$\operatorname{span\text{-}add\text{-}one}(m, s, v) :=$
\begin{equation*}
\{\operatorname{y\_} \in \operatorname{vec}(m) : \exists \operatorname{x\_} \in s, \operatorname{r\_} \in \mathbb{R}.\; (\operatorname{y\_} = \operatorname{vadd}(m)(\operatorname{x\_}, \operatorname{act}(m)(\operatorname{r\_}, v)))\}
\end{equation*}

\par\medskip\noindent\textbf{Definition} (\texttt{hom-vector-space}).\;$\operatorname{hom\text{-}vector\text{-}space}(a, b) :=$
\begin{equation*}
\{\operatorname{homf\_} \in (\operatorname{vec}(a) \to \operatorname{vec}(b)) : \operatorname{is\text{-}hom\text{-}vector\text{-}space}(a, b, \operatorname{homf\_})\}
\end{equation*}

\vnbdef{submodule-subset}{\vnbstmt{\forall m, s.\; (\operatorname{is\text{-}submodule}(m, s) \Rightarrow (s \subseteq \operatorname{vec}(m)))}}
\vnbdef{submodule-vzero-in}{\vnbstmt{\forall m, s.\; (\operatorname{is\text{-}submodule}(m, s) \Rightarrow (\operatorname{vzero}(m) \in s))}}
\vnbdef{submodule-vadd-closed}{\vnbstmt{\forall m, s.; \forall x, y \in s.\; (\operatorname{is\text{-}submodule}(m, s) \Rightarrow (\operatorname{vadd}(m)(x, y) \in s))}}
\vnbdef{submodule-vneg-closed}{\vnbstmt{\forall m, s.; \forall x \in s.\; (\operatorname{is\text{-}submodule}(m, s) \Rightarrow (\operatorname{vneg}(m)(x) \in s))}}
\vnbdef{submodule-act-closed}{\vnbstmt{\begin{array}{@{}l@{}} \forall m, s.; \forall r \in \operatorname{carr}(\operatorname{scal}(m)), x \in s. \\ \quad \operatorname{is\text{-}submodule}(m, s) \Rightarrow (\operatorname{act}(m)(r, x) \in s) \end{array}}}
\vnbdef{is-vector-space-def}{\vnbstmt{\forall s.\; (\operatorname{is\text{-}vector\text{-}space}(s) \Leftrightarrow (\operatorname{is\text{-}module}(s) \wedge \operatorname{is\text{-}field\text{-}ring}(\operatorname{scal}(s))))}}
\vnbdef{is-hom-vector-space-def}{\vnbstmt{\begin{array}{@{}l@{}} \forall a, b, f. \\ \quad \operatorname{is\text{-}hom\text{-}vector\text{-}space}(a, b, f) \\ \quad \Leftrightarrow\; \operatorname{is\text{-}vector\text{-}space}(a) \wedge (\operatorname{is\text{-}vector\text{-}space}(b) \wedge \operatorname{is\text{-}hom\text{-}module}(a, b, f)) \end{array}}}
\vnbdef{vector-space-class}{\vnbstmt{\forall s.\; ((s \in \operatorname{vector\text{-}space}) \Leftrightarrow \operatorname{is\text{-}vector\text{-}space}(s))}}

\vnbbookfile{structure-library/transpose}{structure-library/transpose}
\noindent{\small transpose.scm -- the TRANSPOSE of a matrix (batch 32, 2026-09-25; the user's decision of 2026-09-25, on batch 25-C's proposal).}\par
\par\medskip\noindent\textbf{Definition} (\texttt{transpose}).\;$\operatorname{transpose}(a, m, n) :=$
\begin{equation*}
\begin{array}{@{}l@{}} \operatorname{matof}(n, m, \\ \quad (\lambda {[}\operatorname{tri\_}, \operatorname{trj\_}] \in (\operatorname{interval}(1, n) \times \operatorname{interval}(1, m)).\; \operatorname{entry}(a, \operatorname{trj\_}, \operatorname{tri\_}))) \end{array}
\end{equation*}


\vnbbookfile{structure-library/elementary-matrix}{structure-library/elementary-matrix}
\noindent{\small elementary-matrix.scm -- matrix units and elementary column operations (algebraic-numbers.pdf ch.3, Def 3.2-3.7).  The whole elementary-operation theory is over a COMMUTATIVE ring here; the euclidean assumption only enters later, at the reduction-to-diagonal (Smith) step, NOT here.}\par
\par\medskip\noindent\textbf{Definition} (\texttt{matunit}).\;$\operatorname{matunit}(a, n, k, l) :=$
\begin{equation*}
\begin{array}{@{}l@{}} \operatorname{matof}(n, n, \\ \quad (\lambda {[}i, j] \in (\operatorname{interval}(1, n) \times \operatorname{interval}(1, n)).\; \left(\text{if }((i = k) \wedge (j = l))\text{ then }\operatorname{one}(a)\text{ else }\operatorname{zero}(a)\right))) \end{array}
\end{equation*}

\par\medskip\noindent\textbf{Definition} (\texttt{elem-f}).\;$\operatorname{elem\text{-}f}(a, n, k, l) :=$
\begin{equation*}
\begin{array}{@{}l@{}} \operatorname{matof}(n, n, \\ \quad (\lambda {[}i, j] \in (\operatorname{interval}(1, n) \times \operatorname{interval}(1, n)).\; \\ \quad \quad \left(\text{if }(((i = k) \wedge (j = l)) \vee (((i = l) \wedge (j = k)) \vee ((i = j) \wedge (\neg (i = k) \wedge \neg (i = l)))))\text{ then }\operatorname{one}(a)\text{ else }\operatorname{zero}(a)\right))) \end{array}
\end{equation*}

\par\medskip\noindent\textbf{Definition} (\texttt{elem-g}).\;$\operatorname{elem\text{-}g}(a, n, r, k, l) :=$
\begin{equation*}
\begin{array}{@{}l@{}} \operatorname{matof}(n, n, \\ \quad (\lambda {[}i, j] \in (\operatorname{interval}(1, n) \times \operatorname{interval}(1, n)).\; \\ \quad \quad \left(\text{if }(i = j)\text{ then }\operatorname{one}(a)\text{ else }\left(\text{if }((i = k) \wedge (j = l))\text{ then }r\text{ else }\operatorname{zero}(a)\right)\right))) \end{array}
\end{equation*}

\par\medskip\noindent\textbf{Definition} (\texttt{elem-h}).\;$\operatorname{elem\text{-}h}(a, n, r, k) :=$
\begin{equation*}
\begin{array}{@{}l@{}} \operatorname{matof}(n, n, \\ \quad (\lambda {[}i, j] \in (\operatorname{interval}(1, n) \times \operatorname{interval}(1, n)).\; \\ \quad \quad \left(\text{if }(i = j)\text{ then }\left(\text{if }(i = k)\text{ then }r\text{ else }\operatorname{one}(a)\right)\text{ else }\operatorname{zero}(a)\right))) \end{array}
\end{equation*}


\vnbbookfile{structure-library/mat-equiv}{structure-library/mat-equiv}
\noindent{\small RETIRED 2026-09-17 (proven): ring-neg-neg -- theorem-library/rake-algebra2.scm RETIRED 2026-09-17 (proven): submat-type, border-type -- theorem-library/rake-mat-typing.scm RETIRED 2026-09-17 (proven): diagonal-off-entry -- theorem-library/rake-algebra.scm mat-equiv.scm -- the matrix EQUIVALENCE relation \textasciitilde{} ...}\par
\par\medskip\noindent\textbf{Definition} (\texttt{is-invertible-mat}).\;For all $a$, $n$ and $u$, 
\begin{equation*}
\begin{array}{@{}l@{}} \operatorname{is\text{-}invertible\text{-}mat}(a, n, u) \\ \Leftrightarrow\; u \in \operatorname{mat}(n, n, \operatorname{carr}(a)) \wedge \\ \exists v \in \operatorname{mat}(n, n, \operatorname{carr}(a)). \\ \quad (\operatorname{matmul}(a, u, v) = \operatorname{identmat}(a, n)) \wedge (\operatorname{matmul}(a, v, u) = \operatorname{identmat}(a, n)) \end{array}
\end{equation*}

\par\medskip\noindent\textbf{Definition} (\texttt{mat-equiv}).\;For all $a$, $m$, $n$, $c$ and $d$, 
\begin{equation*}
\begin{array}{@{}l@{}} \operatorname{mat\text{-}equiv}(a, m, n, c, d) \\ \Leftrightarrow\; \exists u. \\ \quad \operatorname{is\text{-}invertible\text{-}mat}(a, m, u) \wedge \\ \quad \exists v.\; (\operatorname{is\text{-}invertible\text{-}mat}(a, n, v) \wedge (d = \operatorname{matmul}(a, \operatorname{matmul}(a, u, c), v))) \end{array}
\end{equation*}

\par\medskip\noindent\textbf{Definition} (\texttt{submat}).\;$\operatorname{submat}(s, p, q) :=$
\begin{equation*}
\begin{array}{@{}l@{}} \operatorname{matof}(p, q, \\ \quad (\lambda {[}i, j] \in (\operatorname{interval}(1, p) \times \operatorname{interval}(1, q)).\; \operatorname{entry}(s, (i + 1), (j + 1)))) \end{array}
\end{equation*}

\par\medskip\noindent\textbf{Definition} (\texttt{border}).\;$\operatorname{border}(a, b, m, p, q) :=$
\begin{equation*}
\begin{array}{@{}l@{}} \operatorname{matof}((p + 1), (q + 1), \\ \quad (\lambda {[}i, j] \in (\operatorname{interval}(1, (p + 1)) \times \operatorname{interval}(1, (q + 1))).\; \\ \quad \quad \left(\text{if }(i = 1)\text{ then }\left(\text{if }(j = 1)\text{ then }b\text{ else }\operatorname{zero}(a)\right)\text{ else }\left(\text{if }(j = 1)\text{ then }\operatorname{zero}(a)\text{ else }\operatorname{entry}(m, \operatorname{nn\text{-}minus}(i, 1), \operatorname{nn\text{-}minus}(j, 1))\right)\right))) \end{array}
\end{equation*}

\par\medskip\noindent\textbf{Definition} (\texttt{is-diagonal}).\;For all $a$, $m$, $n$ and $d$, 
\begin{equation*}
\begin{array}{@{}l@{}} \operatorname{is\text{-}diagonal}(a, m, n, d) \\ \Leftrightarrow\; \forall i \in \operatorname{interval}(1, m), j \in \operatorname{interval}(1, n). \\ \quad \neg (i = j) \Rightarrow (\operatorname{entry}(d, i, j) = \operatorname{zero}(a)) \end{array}
\end{equation*}

\par\medskip\noindent\textbf{Definition} (\texttt{smith-staircase}).\;For all $a$, $m$, $n$, $d$ and $k$, 
\begin{equation*}
\begin{array}{@{}l@{}} \operatorname{smith\text{-}staircase}(a, m, n, d, k) \\ \Leftrightarrow\; \operatorname{is\text{-}diagonal}(a, m, n, d) \wedge \\ k \leq m \wedge \\ k \leq n \wedge \\ \forall \operatorname{i\_} \in \operatorname{interval}(1, k).\; \neg (\operatorname{entry}(d, \operatorname{i\_}, \operatorname{i\_}) = \operatorname{zero}(a)) \wedge \\ \forall \operatorname{i\_} \in \operatorname{interval}(1, m). \\ \quad \neg (\operatorname{i\_} \leq k) \Rightarrow \forall \operatorname{j\_} \in \operatorname{interval}(1, n).\; (\operatorname{entry}(d, \operatorname{i\_}, \operatorname{j\_}) = \operatorname{zero}(a)) \end{array}
\end{equation*}


\vnbbookfile{structure-library/determinant}{structure-library/determinant}
\noindent{\small RETIRED 2026-09-17 (proven): minor-type (modulo 0), det-in-carrier (modulo mpow-type) -- theorem-library/rake-mat-typing.scm, by det-in-carrier-ind with the NN binder outermost. determinant.scm -- the determinant of a square matrix over a ring, by RECURSIVE COFACTOR (Laplace) EXPANSION along the first row.  The ...}\par
\par\medskip\noindent\textbf{Definition} (\texttt{det}).\;For all $r$ and $a$,  $(\operatorname{det}(r, 0, a) \simeq \operatorname{one}(r))$\par
Suppose $r \text{ s.t. } \operatorname{is\text{-}ring}(r)$. Let $n \in \mathbb{N}$ and $a \in \operatorname{mat}(\operatorname{succ}(n), \operatorname{succ}(n), \operatorname{carr}(r))$. Then
\begin{equation*}
(\operatorname{det}(r, \operatorname{succ}(n), a) \simeq \operatorname{finsum}(\operatorname{ring\text{-}additive\text{-}ag}(r), (\lambda j \in \operatorname{interval}(1, \operatorname{succ}(n)).\; \operatorname{mul}(r)(\operatorname{mpow}(\operatorname{ring\text{-}multiplicative\text{-}monoid}(r), \operatorname{neg}(r)(\operatorname{one}(r)), \operatorname{succ}(j)), \operatorname{mul}(r)(\operatorname{entry}(a, 1, j), \operatorname{det}(r, n, \operatorname{minor}(a, 1, j, n))))), \operatorname{interval}(1, \operatorname{succ}(n))))
\end{equation*}

\par\medskip\noindent\textbf{Definition} (\texttt{minor}).\;$\operatorname{minor}(s, r, c, n) :=$
\begin{equation*}
\begin{array}{@{}l@{}} \operatorname{matof}(n, n, \\ \quad (\lambda {[}i, j] \in (\operatorname{interval}(1, n) \times \operatorname{interval}(1, n)).\; \\ \quad \quad \operatorname{entry}(s, \left(\text{if }(i < r)\text{ then }i\text{ else }(i + 1)\right), \left(\text{if }(j < c)\text{ then }j\text{ else }(j + 1)\right)))) \end{array}
\end{equation*}


\vnbbookfile{structure-library/la-det-vocab}{structure-library/la-det-vocab}
\noindent{\small la-det-vocab.scm -- triangular matrices and the formal adjoint (LA-DET, 2026-10-18).}\par
\par\medskip\noindent\textbf{Definition} (\texttt{is-lower-triangular}).\;For all $a$, $n$ and $d$, 
\begin{equation*}
\begin{array}{@{}l@{}} \operatorname{is\text{-}lower\text{-}triangular}(a, n, d) \\ \Leftrightarrow\; \forall \operatorname{ltu\_}, \operatorname{ltv\_} \in \operatorname{interval}(1, n). \\ \quad (\operatorname{ltu\_} < \operatorname{ltv\_}) \Rightarrow (\operatorname{entry}(d, \operatorname{ltu\_}, \operatorname{ltv\_}) = \operatorname{zero}(a)) \end{array}
\end{equation*}

\par\medskip\noindent\textbf{Definition} (\texttt{is-upper-triangular}).\;For all $a$, $n$ and $d$, 
\begin{equation*}
\begin{array}{@{}l@{}} \operatorname{is\text{-}upper\text{-}triangular}(a, n, d) \\ \Leftrightarrow\; \forall \operatorname{ltu\_}, \operatorname{ltv\_} \in \operatorname{interval}(1, n). \\ \quad (\operatorname{ltv\_} < \operatorname{ltu\_}) \Rightarrow (\operatorname{entry}(d, \operatorname{ltu\_}, \operatorname{ltv\_}) = \operatorname{zero}(a)) \end{array}
\end{equation*}

\par\medskip\noindent\textbf{Definition} (\texttt{adj}).\;$\operatorname{adj}(r, n, a) :=$
\begin{equation*}
\begin{array}{@{}l@{}} \operatorname{matof}((n + 1), (n + 1), \\ \quad (\lambda {[}\operatorname{adu\_}, \operatorname{adv\_}] \in (\operatorname{interval}(1, (n + 1)) \times \operatorname{interval}(1, (n + 1))).\; \\ \quad \quad \operatorname{mul}(r)(\operatorname{mpow}(\operatorname{ring\text{-}multiplicative\text{-}monoid}(r), \operatorname{neg}(r)(\operatorname{one}(r)), (\operatorname{adu\_} + \operatorname{adv\_})), \\ \quad \quad \quad \operatorname{det}(r, n, \operatorname{minor}(a, \operatorname{adv\_}, \operatorname{adu\_}, n))))) \end{array}
\end{equation*}


\vnbbookfile{structure-library/la-elem-vocab}{structure-library/la-elem-vocab}
\noindent{\small la-elem-vocab.scm -- elementary column matrices as a class, the subgroup they generate, column equivalence and column step matrices (LA-DET2, 2026-10-18).}\par
\par\medskip\noindent\textbf{Definition} (\texttt{is-elementary}).\;For all $r$, $n$ and $e$, 
\begin{equation*}
\begin{array}{@{}l@{}} \operatorname{is\text{-}elementary}(r, n, e) \\ \Leftrightarrow\; \exists \operatorname{elk\_}, \operatorname{ell\_} \in \operatorname{interval}(1, n). \\ \quad \neg (\operatorname{elk\_} = \operatorname{ell\_}) \wedge (e = \operatorname{elem\text{-}f}(r, n, \operatorname{elk\_}, \operatorname{ell\_})) \vee \\ \exists \operatorname{elr\_} \in \operatorname{carr}(r), \operatorname{elk\_}, \operatorname{ell\_} \in \operatorname{interval}(1, n). \\ \quad \neg (\operatorname{elk\_} = \operatorname{ell\_}) \wedge (e = \operatorname{elem\text{-}g}(r, n, \operatorname{elr\_}, \operatorname{elk\_}, \operatorname{ell\_})) \vee \\ \exists \operatorname{elr\_} \in \operatorname{carr}(r), \operatorname{elk\_}. \\ \quad \exists \operatorname{els\_} \in \operatorname{carr}(r). \\ \quad \quad (\operatorname{mul}(r)(\operatorname{elr\_}, \operatorname{els\_}) = \operatorname{one}(r)) \wedge (\operatorname{mul}(r)(\operatorname{els\_}, \operatorname{elr\_}) = \operatorname{one}(r)) \wedge \\ \quad \operatorname{elk\_} \in \operatorname{interval}(1, n) \wedge \\ \quad e = \operatorname{elem\text{-}h}(r, n, \operatorname{elr\_}, \operatorname{elk\_}) \end{array}
\end{equation*}

\par\medskip\noindent\textbf{Definition} (\texttt{elem-gen}).\;$\operatorname{elem\text{-}gen}(r, n) :=$
\begin{equation*}
\operatorname{intersection\text{-}of}(\{\operatorname{egs\_} \in \operatorname{power}(\operatorname{mat}(n, n, \operatorname{carr}(r))) : ((\operatorname{identmat}(r, n) \in \operatorname{egs\_}) \wedge \forall \operatorname{egx\_} \in \operatorname{egs\_}, \operatorname{ege\_} \in \mathbf{Set}.\; (\operatorname{is\text{-}elementary}(r, n, \operatorname{ege\_}) \Rightarrow (\operatorname{matmul}(r, \operatorname{egx\_}, \operatorname{ege\_}) \in \operatorname{egs\_})))\})
\end{equation*}

\par\medskip\noindent\textbf{Definition} (\texttt{col-equiv}).\;For all $r$, $n$, $a$ and $b$, 
\begin{equation*}
\operatorname{col\text{-}equiv}(r, n, a, b) \Leftrightarrow \exists \operatorname{cep\_} \in \operatorname{elem\text{-}gen}(r, n).\; (b = \operatorname{matmul}(r, a, \operatorname{cep\_}))
\end{equation*}

\par\medskip\noindent\textbf{Definition} (\texttt{is-column-step}).\;For all $r$, $m$, $n$ and $b$, 
\begin{equation*}
\begin{array}{@{}l@{}} \operatorname{is\text{-}column\text{-}step}(r, m, n, b) \\ \Leftrightarrow\; \forall \operatorname{csi\_} \in \operatorname{interval}(1, m), \operatorname{csj\_} \in \operatorname{interval}(1, n). \\ \quad (\operatorname{csi\_} < \operatorname{csj\_}) \Rightarrow (\operatorname{entry}(b, \operatorname{csi\_}, \operatorname{csj\_}) = \operatorname{zero}(r)) \wedge \\ \forall \operatorname{csi\_} \in \operatorname{interval}(1, m). \\ \quad (\operatorname{csi\_} \in \operatorname{interval}(1, n)) \Rightarrow  \\ \quad \quad (\operatorname{entry}(b, \operatorname{csi\_}, \operatorname{csi\_}) = \operatorname{zero}(r)) \vee (\operatorname{entry}(b, \operatorname{csi\_}, \operatorname{csi\_}) = \operatorname{one}(r)) \end{array}
\end{equation*}

\par\medskip\noindent\textbf{Definition} (\texttt{col-step-row}).\;For all $r$, $n$, $b$ and $i$, 
\begin{equation*}
\begin{array}{@{}l@{}} \operatorname{col\text{-}step\text{-}row}(r, n, b, i) \\ \Leftrightarrow\; \forall \operatorname{csj\_} \in \operatorname{interval}(1, n).\; ((i < \operatorname{csj\_}) \Rightarrow (\operatorname{entry}(b, i, \operatorname{csj\_}) = \operatorname{zero}(r))) \wedge \\ (i \in \operatorname{interval}(1, n)) \Rightarrow  \\ \quad (\operatorname{entry}(b, i, i) = \operatorname{zero}(r)) \vee (\operatorname{entry}(b, i, i) = \operatorname{one}(r)) \wedge \\ \quad (\operatorname{entry}(b, i, i) = \operatorname{one}(r)) \Rightarrow  \\ \quad \quad \forall \operatorname{csj\_} \in \operatorname{interval}(1, n). \\ \quad \quad \quad \neg (\operatorname{csj\_} = i) \Rightarrow (\operatorname{entry}(b, i, \operatorname{csj\_}) = \operatorname{zero}(r)) \end{array}
\end{equation*}


\vnbbookfile{structure-library/la-z-vocab}{structure-library/la-z-vocab}
\noindent{\small la-z-vocab.scm -- subdeterminants, equivalence by elementary operations and the gcd of a set of integers (LA-Z, 2026-10-19).}\par
\par\medskip\noindent\textbf{Definition} (\texttt{submatrix}).\;$\operatorname{submatrix}(a, s, p, q) :=$
\begin{equation*}
\begin{array}{@{}l@{}} \operatorname{matof}(s, s, \\ \quad (\lambda {[}\operatorname{smi\_}, \operatorname{smj\_}] \in (\operatorname{interval}(1, s) \times \operatorname{interval}(1, s)).\; \operatorname{entry}(a, p(\operatorname{smi\_}), q(\operatorname{smj\_})))) \end{array}
\end{equation*}

\par\medskip\noindent\textbf{Definition} (\texttt{sdet}).\;$\operatorname{sdet}(r, m, n, a, s) :=$
\begin{equation*}
\{\operatorname{sdx\_} \in \operatorname{carr}(r) : \exists \operatorname{sdp\_} \in \operatorname{injection}(\operatorname{interval}(1, s), \operatorname{interval}(1, m)), \operatorname{sdq\_} \in \operatorname{injection}(\operatorname{interval}(1, s), \operatorname{interval}(1, n)).\; (\operatorname{sdx\_} = \operatorname{det}(r, s, \operatorname{submatrix}(a, s, \operatorname{sdp\_}, \operatorname{sdq\_})))\}
\end{equation*}

\par\medskip\noindent\textbf{Definition} (\texttt{elem-equiv}).\;For all $r$, $m$, $n$, $a$ and $b$, 
\begin{equation*}
\begin{array}{@{}l@{}} \operatorname{elem\text{-}equiv}(r, m, n, a, b) \\ \Leftrightarrow\; \exists \operatorname{eqp\_} \in \operatorname{elem\text{-}gen}(r, m), \operatorname{eqq\_} \in \operatorname{elem\text{-}gen}(r, n). \\ \quad b = \operatorname{matmul}(r, \operatorname{matmul}(r, \operatorname{eqp\_}, a), \operatorname{eqq\_}) \end{array}
\end{equation*}

\par\medskip\noindent\textbf{Definition} (\texttt{zz-is-gcd-set}).\;For all $s$ and $d$, 
\begin{equation*}
\begin{array}{@{}l@{}} \operatorname{zz\text{-}is\text{-}gcd\text{-}set}(s, d) \\ \Leftrightarrow\; d \in \mathbb{Z} \wedge \\ \forall \operatorname{gsx\_} \in s.\; \operatorname{zz\text{-}divides}(d, \operatorname{gsx\_}) \wedge \\ \forall \operatorname{gse\_} \in \mathbb{Z}. \\ \quad \forall \operatorname{gsx\_} \in s.\; \operatorname{zz\text{-}divides}(\operatorname{gse\_}, \operatorname{gsx\_}) \Rightarrow \operatorname{zz\text{-}divides}(\operatorname{gse\_}, d) \end{array}
\end{equation*}


\vnbbookfile{structure-library/mod-seq}{structure-library/mod-seq}
\noindent{\small RETIRED 2026-09-17 (proven): finsum-act-collect-gen -- theorem-library/rake-finsum-core.scm RETIRED 2026-09-17 (proven): finsum-act-distrib-gen -- theorem-library/rake-finsum-laws2.scm RETIRED 2026-09-17 (proven): matact-entry -- theorem-library/rake-identmat.scm (the module twin of matmul-entry). RETIRED 2026-09-17 ...}\par
\par\medskip\noindent\textbf{Definition} (\texttt{matact}).\;$\operatorname{matact}(\mathit{md}, p, u) :=$
\begin{equation*}
\begin{array}{@{}l@{}} \operatorname{matof}(\operatorname{nth}(1, \operatorname{size}(p)), \operatorname{nth}(2, \operatorname{size}(u)), \\ \quad (\lambda {[}i, c] \in (\operatorname{interval}(1, \operatorname{nth}(1, \operatorname{size}(p))) \times \operatorname{interval}(1, \operatorname{nth}(2, \operatorname{size}(u)))).\; \\ \quad \quad \operatorname{finsum}(\operatorname{module\text{-}vector\text{-}ag}(\mathit{md}), \\ \quad \quad \quad (\lambda j \in \operatorname{interval}(1, \operatorname{nth}(2, \operatorname{size}(p))).\; \operatorname{act}(\mathit{md})(\operatorname{entry}(p, i, j), \operatorname{entry}(u, j, c))), \\ \quad \quad \quad \operatorname{interval}(1, \operatorname{nth}(2, \operatorname{size}(p)))))) \end{array}
\end{equation*}

\par\medskip\noindent\textbf{Definition} (\texttt{lincomb}).\;$\operatorname{lincomb}(\mathit{md}, n, c, u) :=$
\begin{equation*}
\begin{array}{@{}l@{}} \operatorname{finsum}(\operatorname{module\text{-}vector\text{-}ag}(\mathit{md}), \\ \quad (\lambda j \in \operatorname{interval}(1, n).\; \operatorname{act}(\mathit{md})(\operatorname{entry}(c, 1, j), \operatorname{entry}(u, j, 1))), \\ \quad \operatorname{interval}(1, n)) \end{array}
\end{equation*}

\par\medskip\noindent\textbf{Definition} (\texttt{generates}).\;For all $\mathit{md}$, $n$ and $u$, 
\begin{equation*}
\begin{array}{@{}l@{}} \operatorname{generates}(\mathit{md}, n, u) \\ \Leftrightarrow\; \forall \operatorname{x\_} \in \operatorname{vec}(\mathit{md}). \exists \operatorname{c\_} \in \operatorname{mat}(1, n, \operatorname{carr}(\operatorname{scal}(\mathit{md}))). \\ \quad \operatorname{x\_} = \operatorname{lincomb}(\mathit{md}, n, \operatorname{c\_}, u) \end{array}
\end{equation*}

\par\medskip\noindent\textbf{Definition} (\texttt{rel-free}).\;For all $\mathit{md}$, $n$ and $u$, 
\begin{equation*}
\begin{array}{@{}l@{}} \operatorname{rel\text{-}free}(\mathit{md}, n, u) \\ \Leftrightarrow\; \forall \operatorname{c\_} \in \operatorname{mat}(1, n, \operatorname{carr}(\operatorname{scal}(\mathit{md}))). \\ \quad (\operatorname{lincomb}(\mathit{md}, n, \operatorname{c\_}, u) = \operatorname{vzero}(\mathit{md})) \Rightarrow  \\ \quad \quad \forall \operatorname{j\_} \in \operatorname{interval}(1, n).\; (\operatorname{entry}(\operatorname{c\_}, 1, \operatorname{j\_}) = \operatorname{zero}(\operatorname{scal}(\mathit{md}))) \end{array}
\end{equation*}

\par\medskip\noindent\textbf{Definition} (\texttt{spans}).\;For all $\mathit{md}$, $n$, $u$ and $\mathit{sm}$, 
\begin{equation*}
\begin{array}{@{}l@{}} \operatorname{spans}(\mathit{md}, n, u, \mathit{sm}) \\ \Leftrightarrow\; \forall \operatorname{j\_} \in \operatorname{interval}(1, n).\; (\operatorname{entry}(u, \operatorname{j\_}, 1) \in \mathit{sm}) \wedge \\ \forall \operatorname{x\_} \in \mathit{sm}. \exists \operatorname{c\_} \in \operatorname{mat}(1, n, \operatorname{carr}(\operatorname{scal}(\mathit{md}))). \\ \quad \operatorname{x\_} = \operatorname{lincomb}(\mathit{md}, n, \operatorname{c\_}, u) \end{array}
\end{equation*}

\par\medskip\noindent\textbf{Definition} (\texttt{span}).\;$\operatorname{span}(\mathit{md}, n, u) :=$
\begin{equation*}
\{\operatorname{x\_} \in \operatorname{vec}(\mathit{md}) : \exists \operatorname{c\_} \in \operatorname{mat}(1, n, \operatorname{carr}(\operatorname{scal}(\mathit{md}))).\; (\operatorname{x\_} = \operatorname{lincomb}(\mathit{md}, n, \operatorname{c\_}, u))\}
\end{equation*}

\par\medskip\noindent\textbf{Definition} (\texttt{lastcoeff-set}).\;$\operatorname{lastcoeff\text{-}set}(\mathit{md}, p, u, \mathit{sm}) :=$
\begin{equation*}
\{\operatorname{r\_} \in \operatorname{carr}(\operatorname{scal}(\mathit{md})) : \exists \operatorname{c\_} \in \operatorname{mat}(1, (p + 1), \operatorname{carr}(\operatorname{scal}(\mathit{md}))).\; ((\operatorname{entry}(\operatorname{c\_}, 1, (p + 1)) = \operatorname{r\_}) \wedge (\operatorname{lincomb}(\mathit{md}, (p + 1), \operatorname{c\_}, u) \in \mathit{sm}))\}
\end{equation*}

\vnbdef{span-membership}{\vnbstmt{\begin{array}{@{}l@{}} \forall d, n, u, x. \\ \quad x \in \operatorname{span}(d, n, u) \\ \quad \Leftrightarrow\; ((x \in \operatorname{vec}(d)) \wedge \exists c \in \operatorname{mat}(1, n, \operatorname{carr}(\operatorname{scal}(d))).\; (x = \operatorname{lincomb}(d, n, c, u))) \end{array}}}
\vnbdef{lastcoeff-set-membership}{\vnbstmt{\begin{array}{@{}l@{}} \forall d, p, u, m, r. \\ \quad r \in \operatorname{lastcoeff\text{-}set}(d, p, u, m) \\ \quad \Leftrightarrow\; r \in \operatorname{carr}(\operatorname{scal}(d)) \wedge \\ \quad \exists c \in \operatorname{mat}(1, (p + 1), \operatorname{carr}(\operatorname{scal}(d))). \\ \quad \quad (\operatorname{entry}(c, 1, (p + 1)) = r) \wedge (\operatorname{lincomb}(d, (p + 1), c, u) \in m) \end{array}}}

\vnbbookfile{structure-library/la-z2-vocab}{structure-library/la-z2-vocab}
\noindent{\small la-z2-vocab.scm -- the principal sequence of a subgroup (linear algebra, the definition after Corollary 3.45) (LA-Z2, 2026-10-23).  No proof in this file.}\par
\par\medskip\noindent\textbf{Definition} (\texttt{is-principal-sequence}).\;For all $\mathit{md}$, $n$, $\mathit{sm}$, $\mathit{es}$ and $\mathit{dm}$, 
\begin{equation*}
\begin{array}{@{}l@{}} \operatorname{is\text{-}principal\text{-}sequence}(\mathit{md}, n, \mathit{sm}, \mathit{es}, \mathit{dm}) \\ \Leftrightarrow\; \mathit{es} \in \operatorname{mat}(n, 1, \operatorname{vec}(\mathit{md})) \wedge \\ \operatorname{generates}(\mathit{md}, n, \mathit{es}) \wedge \\ \operatorname{rel\text{-}free}(\mathit{md}, n, \mathit{es}) \wedge \\ \mathit{dm} \in \operatorname{mat}(n, n, \mathbb{Z}) \wedge \\ \operatorname{is\text{-}diagonal}(\operatorname{zz\text{-}ring}, n, n, \mathit{dm}) \wedge \\ \forall \operatorname{psi\_} \in \operatorname{interval}(1, n).\; (0 \leq \operatorname{entry}(\mathit{dm}, \operatorname{psi\_}, \operatorname{psi\_})) \wedge \\ \forall \operatorname{psi\_}, \operatorname{psj\_} \in \operatorname{interval}(1, n). \\ \quad (\operatorname{psi\_} \leq \operatorname{psj\_}) \Rightarrow \operatorname{zz\text{-}divides}(\operatorname{entry}(\mathit{dm}, \operatorname{psi\_}, \operatorname{psi\_}), \operatorname{entry}(\mathit{dm}, \operatorname{psj\_}, \operatorname{psj\_})) \wedge \\ \operatorname{spans}(\mathit{md}, n, \operatorname{matact}(\mathit{md}, \mathit{dm}, \mathit{es}), \mathit{sm}) \end{array}
\end{equation*}


\vnbbookfile{structure-library/la-z3-vocab}{structure-library/la-z3-vocab}
\noindent{\small la-z3-vocab.scm -- quotient modules, the ring as a module over itself, the cyclic groups Z/(a), finite direct sums of modules and module isomorphisms (LA-Z3, 2026-10-24).  No proof in this file.}\par
\par\medskip\noindent\textbf{Definition} (\texttt{qmod-rel}).\;$\operatorname{qmod\text{-}rel}(\mathit{md}, \mathit{sm}) :=$
\begin{equation*}
\{\operatorname{qmp\_} \in (\operatorname{vec}(\mathit{md}) \times \operatorname{vec}(\mathit{md})) : (\operatorname{vadd}(\mathit{md})(\operatorname{nth}(1, \operatorname{qmp\_}), \operatorname{vneg}(\mathit{md})(\operatorname{nth}(2, \operatorname{qmp\_}))) \in \mathit{sm})\}
\end{equation*}

\par\medskip\noindent\textbf{Definition} (\texttt{qmod-setoid}).\;$\operatorname{qmod\text{-}setoid}(\mathit{md}, \mathit{sm}) :=$ ${[}\operatorname{vec}(\mathit{md}), \operatorname{qmod\text{-}rel}(\mathit{md}, \mathit{sm})]$
\par\medskip\noindent\textbf{Definition} (\texttt{quotient-module}).\;$\operatorname{quotient\text{-}module}(\mathit{md}, \mathit{sm}) :=$
\begin{equation*}
{[}\operatorname{scal}(\mathit{md}), \operatorname{quotient}(\operatorname{qmod\text{-}setoid}(\mathit{md}, \mathit{sm})), \operatorname{qmod\text{-}vadd}(\mathit{md}, \mathit{sm}), \operatorname{class}(\operatorname{qmod\text{-}setoid}(\mathit{md}, \mathit{sm}), \operatorname{vzero}(\mathit{md})), \operatorname{qmod\text{-}vneg}(\mathit{md}, \mathit{sm}), \operatorname{qmod\text{-}act}(\mathit{md}, \mathit{sm})]
\end{equation*}

\par\medskip\noindent\textbf{Definition} (\texttt{qmod-scale}).\;$\operatorname{qmod\text{-}scale}(\mathit{md}, \mathit{sm}, r) :=$
\begin{equation*}
\begin{array}{@{}l@{}} \operatorname{descend}(\operatorname{qmod\text{-}setoid}(\mathit{md}, \mathit{sm}), \\ \quad (\lambda \operatorname{qma\_} \in \operatorname{vec}(\mathit{md}).\; \operatorname{class}(\operatorname{qmod\text{-}setoid}(\mathit{md}, \mathit{sm}), \operatorname{act}(\mathit{md})(r, \operatorname{qma\_})))) \end{array}
\end{equation*}

\par\medskip\noindent\textbf{Definition} (\texttt{ring-as-module}).\;$\operatorname{ring\text{-}as\text{-}module}(r) :=$ ${[}r, \operatorname{carr}(r), \operatorname{add}(r), \operatorname{zero}(r), \operatorname{neg}(r), \operatorname{mul}(r)]$
\par\medskip\noindent\textbf{Definition} (\texttt{zz-mod}).\;$\operatorname{zz\text{-}mod}(a) :=$
\begin{equation*}
\operatorname{quotient\text{-}module}(\operatorname{ring\text{-}as\text{-}module}(\operatorname{zz\text{-}ring}), \operatorname{principal\text{-}ideal}(\operatorname{zz\text{-}ring}, a))
\end{equation*}

\par\medskip\noindent\textbf{Definition} (\texttt{direct-sum}).\;$\operatorname{direct\text{-}sum}(s, n, f) :=$
\begin{equation*}
{[}s, \operatorname{dsum\text{-}vec}(n, f), \operatorname{dsum\text{-}vadd}(n, f), \operatorname{dsum\text{-}vzero}(n, f), \operatorname{dsum\text{-}vneg}(n, f), \operatorname{dsum\text{-}act}(s, n, f)]
\end{equation*}

\par\medskip\noindent\textbf{Definition} (\texttt{is-iso-module}).\;For all $a$, $b$ and $f$, 
\begin{equation*}
\operatorname{is\text{-}iso\text{-}module}(a, b, f) \Leftrightarrow (\operatorname{is\text{-}hom\text{-}module}(a, b, f) \wedge (f \in \mathrm{Bij}(\operatorname{vec}(a), \operatorname{vec}(b))))
\end{equation*}

\par\medskip\noindent\textbf{Definition} (\texttt{qmod-vadd}).\;$\operatorname{qmod\text{-}vadd}(\mathit{md}, \mathit{sm}) :=$
\begin{equation*}
\begin{array}{@{}l@{}} \operatorname{descend2}(\operatorname{qmod\text{-}setoid}(\mathit{md}, \mathit{sm}), \\ \quad (\lambda {[}\operatorname{qma\_}, \operatorname{qmb\_}] \in (\operatorname{vec}(\mathit{md}) \times \operatorname{vec}(\mathit{md})).\; \\ \quad \quad \operatorname{class}(\operatorname{qmod\text{-}setoid}(\mathit{md}, \mathit{sm}), \operatorname{vadd}(\mathit{md})(\operatorname{qma\_}, \operatorname{qmb\_})))) \end{array}
\end{equation*}

\par\medskip\noindent\textbf{Definition} (\texttt{qmod-vneg}).\;$\operatorname{qmod\text{-}vneg}(\mathit{md}, \mathit{sm}) :=$
\begin{equation*}
\begin{array}{@{}l@{}} \operatorname{descend}(\operatorname{qmod\text{-}setoid}(\mathit{md}, \mathit{sm}), \\ \quad (\lambda \operatorname{qma\_} \in \operatorname{vec}(\mathit{md}).\; \operatorname{class}(\operatorname{qmod\text{-}setoid}(\mathit{md}, \mathit{sm}), \operatorname{vneg}(\mathit{md})(\operatorname{qma\_})))) \end{array}
\end{equation*}

\par\medskip\noindent\textbf{Definition} (\texttt{qmod-act}).\;$\operatorname{qmod\text{-}act}(\mathit{md}, \mathit{sm}) :=$
\begin{equation*}
\begin{array}{@{}l@{}} (\lambda {[}\operatorname{qmr\_}, \operatorname{qmc\_}] \in (\operatorname{carr}(\operatorname{scal}(\mathit{md})) \times \operatorname{quotient}(\operatorname{qmod\text{-}setoid}(\mathit{md}, \mathit{sm}))).\; \\ \quad \operatorname{qmod\text{-}scale}(\mathit{md}, \mathit{sm}, \operatorname{qmr\_})(\operatorname{qmc\_})) \end{array}
\end{equation*}

\par\medskip\noindent\textbf{Definition} (\texttt{dsum-vec}).\;$\operatorname{dsum\text{-}vec}(n, f) :=$
\begin{equation*}
\{\operatorname{dsv\_} \in (\operatorname{interval}(1, n) \to \operatorname{big\text{-}union}(\operatorname{dsk\_}, \operatorname{interval}(1, n), \operatorname{vec}(f(\operatorname{dsk\_})))) : \forall \operatorname{dsj\_} \in \operatorname{interval}(1, n).\; (\operatorname{dsv\_}(\operatorname{dsj\_}) \in \operatorname{vec}(f(\operatorname{dsj\_})))\}
\end{equation*}

\par\medskip\noindent\textbf{Definition} (\texttt{dsum-vadd}).\;$\operatorname{dsum\text{-}vadd}(n, f) :=$
\begin{equation*}
\begin{array}{@{}l@{}} (\lambda {[}\operatorname{dsf\_}, \operatorname{dsg\_}] \in (\operatorname{dsum\text{-}vec}(n, f) \times \operatorname{dsum\text{-}vec}(n, f)).\; \\ \quad (\lambda \operatorname{dsi\_} \in \operatorname{interval}(1, n).\; \operatorname{vadd}(f(\operatorname{dsi\_}))(\operatorname{dsf\_}(\operatorname{dsi\_}), \operatorname{dsg\_}(\operatorname{dsi\_})))) \end{array}
\end{equation*}

\par\medskip\noindent\textbf{Definition} (\texttt{dsum-vzero}).\;$\operatorname{dsum\text{-}vzero}(n, f) :=$ $(\lambda \operatorname{dsi\_} \in \operatorname{interval}(1, n).\; \operatorname{vzero}(f(\operatorname{dsi\_})))$
\par\medskip\noindent\textbf{Definition} (\texttt{dsum-vneg}).\;$\operatorname{dsum\text{-}vneg}(n, f) :=$
\begin{equation*}
(\lambda \operatorname{dsf\_} \in \operatorname{dsum\text{-}vec}(n, f).\; (\lambda \operatorname{dsi\_} \in \operatorname{interval}(1, n).\; \operatorname{vneg}(f(\operatorname{dsi\_}))(\operatorname{dsf\_}(\operatorname{dsi\_}))))
\end{equation*}

\par\medskip\noindent\textbf{Definition} (\texttt{dsum-act}).\;$\operatorname{dsum\text{-}act}(s, n, f) :=$
\begin{equation*}
\begin{array}{@{}l@{}} (\lambda {[}\operatorname{dqr\_}, \operatorname{dsf\_}] \in (\operatorname{carr}(s) \times \operatorname{dsum\text{-}vec}(n, f)).\; \\ \quad (\lambda \operatorname{dsi\_} \in \operatorname{interval}(1, n).\; \operatorname{act}(f(\operatorname{dsi\_}))(\operatorname{dqr\_}, \operatorname{dsf\_}(\operatorname{dsi\_})))) \end{array}
\end{equation*}


\vnbbookfile{structure-library/la-z4-vocab}{structure-library/la-z4-vocab}
\noindent{\small la-z4-vocab.scm -- the free module Z\textasciicircum{}n, the image of a sequence under a map, the direct sum of cyclic groups of a sequence of integers, and invariant factors (LA-Z4, 2026-10-24).  No proof in this file.}\par
\par\medskip\noindent\textbf{Definition} (\texttt{zz-pow}).\;$\operatorname{zz\text{-}pow}(n) :=$
\begin{equation*}
\operatorname{direct\text{-}sum}(\operatorname{zz\text{-}ring}, n, (\lambda \operatorname{zpi\_} \in \operatorname{interval}(1, n).\; \operatorname{ring\text{-}as\text{-}module}(\operatorname{zz\text{-}ring})))
\end{equation*}

\par\medskip\noindent\textbf{Definition} (\texttt{zz-units}).\;$\operatorname{zz\text{-}units}(n) :=$
\begin{equation*}
\begin{array}{@{}l@{}} \operatorname{matof}(n, 1, \\ \quad (\lambda {[}\operatorname{zua\_}, \operatorname{zub\_}] \in (\operatorname{interval}(1, n) \times \operatorname{interval}(1, 1)).\; \\ \quad \quad (\lambda \operatorname{zut\_} \in \operatorname{interval}(1, n).\; \operatorname{entry}(\operatorname{unitrow}(\operatorname{zz\text{-}ring}, n, \operatorname{zua\_}), 1, \operatorname{zut\_})))) \end{array}
\end{equation*}

\par\medskip\noindent\textbf{Definition} (\texttt{zz-row}).\;$\operatorname{zz\text{-}row}(n, x) :=$
\begin{equation*}
\operatorname{matof}(1, n, (\lambda {[}\operatorname{zra\_}, \operatorname{zrb\_}] \in (\operatorname{interval}(1, 1) \times \operatorname{interval}(1, n)).\; x(\operatorname{zrb\_})))
\end{equation*}

\par\medskip\noindent\textbf{Definition} (\texttt{map-seq}).\;$\operatorname{map\text{-}seq}(h, n, u) :=$
\begin{equation*}
\begin{array}{@{}l@{}} \operatorname{matof}(n, 1, \\ \quad (\lambda {[}\operatorname{msa\_}, \operatorname{msb\_}] \in (\operatorname{interval}(1, n) \times \operatorname{interval}(1, 1)).\; h(\operatorname{entry}(u, \operatorname{msa\_}, \operatorname{msb\_})))) \end{array}
\end{equation*}

\par\medskip\noindent\textbf{Definition} (\texttt{cyclic-sum}).\;$\operatorname{cyclic\text{-}sum}(l, a) :=$
\begin{equation*}
\operatorname{direct\text{-}sum}(\operatorname{zz\text{-}ring}, l, (\lambda \operatorname{csi\_} \in \operatorname{interval}(1, l).\; \operatorname{zz\text{-}mod}(a(\operatorname{csi\_}))))
\end{equation*}

\par\medskip\noindent\textbf{Definition} (\texttt{is-invariant-factors}).\;For all $l$ and $a$, 
\begin{equation*}
\begin{array}{@{}l@{}} \operatorname{is\text{-}invariant\text{-}factors}(l, a) \\ \Leftrightarrow\; a \in (\operatorname{interval}(1, l) \to \mathbb{Z}) \wedge \\ \forall \operatorname{ifi\_} \in \operatorname{interval}(1, l).\; (0 \leq a(\operatorname{ifi\_})) \wedge \\ \forall \operatorname{ifi\_} \in \operatorname{interval}(1, l).\; \neg (a(\operatorname{ifi\_}) = 1) \wedge \\ \forall \operatorname{ifi\_}, \operatorname{ifj\_} \in \operatorname{interval}(1, l). \\ \quad (\operatorname{ifi\_} \leq \operatorname{ifj\_}) \Rightarrow \operatorname{zz\text{-}divides}(a(\operatorname{ifi\_}), a(\operatorname{ifj\_})) \end{array}
\end{equation*}

\par\medskip\noindent\textbf{Definition} (\texttt{diag-mat}).\;$\operatorname{diag\text{-}mat}(n, f) :=$
\begin{equation*}
\begin{array}{@{}l@{}} \operatorname{matof}(n, n, \\ \quad (\lambda {[}\operatorname{dga\_}, \operatorname{dgb\_}] \in (\operatorname{interval}(1, n) \times \operatorname{interval}(1, n)).\; \\ \quad \quad \operatorname{mul}(\operatorname{zz\text{-}ring})(f(\operatorname{dga\_}), \operatorname{entry}(\operatorname{identmat}(\operatorname{zz\text{-}ring}, n), \operatorname{dga\_}, \operatorname{dgb\_})))) \end{array}
\end{equation*}

\par\medskip\noindent\textbf{Definition} (\texttt{pad-ones}).\;$\operatorname{pad\text{-}ones}(m, k, f) :=$
\begin{equation*}
(\lambda \operatorname{poj\_} \in \operatorname{interval}(1, m).\; \left(\text{if }(\operatorname{poj\_} \leq k)\text{ then }1\text{ else }f(\operatorname{nn\text{-}minus}(\operatorname{poj\_}, k))\right))
\end{equation*}

\par\medskip\noindent\textbf{Definition} (\texttt{seq-head}).\;$\operatorname{seq\text{-}head}(n, z) :=$ $(\lambda \operatorname{shi\_} \in \operatorname{interval}(1, n).\; z(\operatorname{shi\_}))$
\par\medskip\noindent\textbf{Definition} (\texttt{seq-shift}).\;$\operatorname{seq\text{-}shift}(n, z) :=$ $(\lambda \operatorname{sfi\_} \in \operatorname{interval}(1, n).\; z((n + \operatorname{sfi\_})))$
\par\medskip\noindent\textbf{Definition} (\texttt{seq-join}).\;$\operatorname{seq\text{-}join}(m, n, \mathit{hd}, \mathit{tl}) :=$
\begin{equation*}
(\lambda \operatorname{jnj\_} \in \operatorname{interval}(1, m).\; \left(\text{if }(\operatorname{jnj\_} \leq n)\text{ then }\operatorname{hd}(\operatorname{jnj\_})\text{ else }\operatorname{tl}(\operatorname{nn\text{-}minus}(\operatorname{jnj\_}, n))\right))
\end{equation*}

\par\medskip\noindent\textbf{Definition} (\texttt{schanuel-map}).\;$\operatorname{schanuel\text{-}map}(m, n, \mathit{la}, \mathit{lb}) :=$
\begin{equation*}
\begin{array}{@{}l@{}} (\lambda \operatorname{scz\_} \in \operatorname{vec}(\operatorname{zz\text{-}pow}(m)).\; \\ \quad \operatorname{seq\text{-}join}(m, n, \\ \quad \quad \operatorname{vadd}(\operatorname{zz\text{-}pow}(n))(\operatorname{seq\text{-}shift}(n, \operatorname{scz\_}), \operatorname{vneg}(\operatorname{zz\text{-}pow}(n))(\operatorname{lb}(\operatorname{seq\text{-}head}(n, \operatorname{scz\_})))), \\ \quad \quad \operatorname{vadd}(\operatorname{zz\text{-}pow}(n))(\operatorname{seq\text{-}head}(n, \operatorname{scz\_}), \\ \quad \quad \quad \operatorname{la}(\operatorname{vadd}(\operatorname{zz\text{-}pow}(n))(\operatorname{seq\text{-}shift}(n, \operatorname{scz\_}), \operatorname{vneg}(\operatorname{zz\text{-}pow}(n))(\operatorname{lb}(\operatorname{seq\text{-}head}(n, \operatorname{scz\_})))))))) \end{array}
\end{equation*}


\vnbbookfile{structure-library/la-z5-vocab}{structure-library/la-z5-vocab}
\noindent{\small la-z5-vocab.scm -- the norm N\_E(F) of a submodule (linear algebra, the definition after Corollary 3.45, used by Corollary 3.52) (LA-Z5, 2026-10-25).  No proof in this file.}\par
\par\medskip\noindent\textbf{Definition} (\texttt{is-norm-data}).\;For all $\mathit{md}$, $\mathit{sm}$, $n$, $\mathit{xs}$, $\mathit{ys}$ and $\mathit{am}$, 
\begin{equation*}
\begin{array}{@{}l@{}} \operatorname{is\text{-}norm\text{-}data}(\mathit{md}, \mathit{sm}, n, \mathit{xs}, \mathit{ys}, \mathit{am}) \\ \Leftrightarrow\; n \in \mathbb{N} \wedge \\ \mathit{xs} \in \operatorname{mat}(n, 1, \operatorname{vec}(\mathit{md})) \wedge \\ \operatorname{generates}(\mathit{md}, n, \mathit{xs}) \wedge \\ \operatorname{rel\text{-}free}(\mathit{md}, n, \mathit{xs}) \wedge \\ \mathit{ys} \in \operatorname{mat}(n, 1, \operatorname{vec}(\mathit{md})) \wedge \\ \operatorname{spans}(\mathit{md}, n, \mathit{ys}, \mathit{sm}) \wedge \\ \operatorname{rel\text{-}free}(\mathit{md}, n, \mathit{ys}) \wedge \\ \mathit{am} \in \operatorname{mat}(n, n, \mathbb{Z}) \wedge \\ \mathit{ys} = \operatorname{matact}(\mathit{md}, \mathit{am}, \mathit{xs}) \end{array}
\end{equation*}

\par\medskip\noindent\textbf{Definition} (\texttt{submodule-norm}).\;$\operatorname{submodule\text{-}norm}(\mathit{md}, \mathit{sm}) :=$
\begin{equation*}
\left(\text{if }\exists \operatorname{nmn\_}, \operatorname{nmx\_}, \operatorname{nmy\_}, \operatorname{nma\_}.\; \operatorname{is\text{-}norm\text{-}data}(\mathit{md}, \mathit{sm}, \operatorname{nmn\_}, \operatorname{nmx\_}, \operatorname{nmy\_}, \operatorname{nma\_})\text{ then }(\iota\, \operatorname{nmd\_}.\;\exists \operatorname{nmn\_}, \operatorname{nmx\_}, \operatorname{nmy\_}, \operatorname{nma\_}.\; (\operatorname{is\text{-}norm\text{-}data}(\mathit{md}, \mathit{sm}, \operatorname{nmn\_}, \operatorname{nmx\_}, \operatorname{nmy\_}, \operatorname{nma\_}) \wedge (\operatorname{nmd\_} = \lvert \operatorname{det}(\operatorname{zz\text{-}ring}, \operatorname{nmn\_}, \operatorname{nma\_}) \rvert)))\text{ else }0\right)
\end{equation*}


\vnbbookfile{structure-library/qag-vocab}{structure-library/qag-vocab}
\noindent{\small qag-vocab.scm -- WAVE 2 (2026-10-27..30), wired 2026-10-30 from scratchpad/ccrpr/qag-vocab.scm. Library: linear-algebra.}\par
\par\medskip\noindent\textbf{Definition} (\texttt{ag-zz-module}).\;$\operatorname{ag\text{-}zz\text{-}module}(\operatorname{qag\_}) :=$
\begin{equation*}
{[}\operatorname{zz\text{-}ring}, \operatorname{carr}(\operatorname{qag\_}), \operatorname{opr}(\operatorname{qag\_}), \operatorname{iden}(\operatorname{qag\_}), \operatorname{inv}(\operatorname{qag\_}), (\lambda {[}\operatorname{azk\_}, \operatorname{aza\_}] \in (\mathbb{Z} \times \operatorname{carr}(\operatorname{qag\_})).\; \operatorname{zz\text{-}act}(\operatorname{qag\_}, \operatorname{azk\_}, \operatorname{aza\_}))]
\end{equation*}

\par\medskip\noindent\textbf{Definition} (\texttt{quotient-ag}).\;$\operatorname{quotient\text{-}ag}(\operatorname{qag\_}, \operatorname{qah\_}) :=$ $\operatorname{module\text{-}vector\text{-}ag}(\operatorname{quotient\text{-}module}(\operatorname{ag\text{-}zz\text{-}module}(\operatorname{qag\_}), \operatorname{qah\_}))$
\par\medskip\noindent\textbf{Definition} (\texttt{is-iso-group}).\;For all $\operatorname{qag\_}$, $\operatorname{qah\_}$ and $\operatorname{azk\_}$, 
\begin{equation*}
\begin{array}{@{}l@{}} \operatorname{is\text{-}iso\text{-}group}(\operatorname{qag\_}, \operatorname{qah\_}, \operatorname{azk\_}) \\ \Leftrightarrow\; \operatorname{is\text{-}hom\text{-}group}(\operatorname{qag\_}, \operatorname{qah\_}, \operatorname{azk\_}) \wedge (\operatorname{azk\_} \in \mathrm{Bij}(\operatorname{carr}(\operatorname{qag\_}), \operatorname{carr}(\operatorname{qah\_}))) \end{array}
\end{equation*}


\vnbbookfile{structure-library/ring-pow-vocab}{structure-library/ring-pow-vocab}
\noindent{\small ring-pow-vocab.scm -- wave of 2026-11-09/10 (agent CCR-1109-E), wired 2026-11-10 from scratchpad/ccr1109e/ring-pow-vocab.scm. Library: linear-algebra.}\par
\par\medskip\noindent\textbf{Definition} (\texttt{ring-pow}).\;$\operatorname{ring\text{-}pow}(r, n) :=$
\begin{equation*}
\operatorname{direct\text{-}sum}(r, n, (\lambda \operatorname{rqpi\_} \in \operatorname{interval}(1, n).\; \operatorname{ring\text{-}as\text{-}module}(r)))
\end{equation*}

\par\medskip\noindent\textbf{Definition} (\texttt{ring-units}).\;$\operatorname{ring\text{-}units}(r, n) :=$
\begin{equation*}
\begin{array}{@{}l@{}} \operatorname{matof}(n, 1, \\ \quad (\lambda {[}\operatorname{rqua\_}, \operatorname{rqub\_}] \in (\operatorname{interval}(1, n) \times \operatorname{interval}(1, 1)).\; \\ \quad \quad (\lambda \operatorname{rqut\_} \in \operatorname{interval}(1, n).\; \operatorname{entry}(\operatorname{unitrow}(r, n, \operatorname{rqua\_}), 1, \operatorname{rqut\_})))) \end{array}
\end{equation*}

\par\medskip\noindent\textbf{Definition} (\texttt{ring-row}).\;$\operatorname{ring\text{-}row}(r, n, x) :=$
\begin{equation*}
\operatorname{matof}(1, n, (\lambda {[}\operatorname{rqra\_}, \operatorname{rqrb\_}] \in (\operatorname{interval}(1, 1) \times \operatorname{interval}(1, n)).\; x(\operatorname{rqrb\_})))
\end{equation*}


\vnbbookfile{structure-library/is-set-laws}{structure-library/is-set-laws}
\noindent{\small A file shared across books; the entries that belong to this one.}\par
\noindent{\small structure-library/is-set-laws.scm -- THE GENERATED SETHOOD THEOREMS (W1I, 2026-11-06).}\par
\vnbprop{is-module-is-set}{\vnbstmt{\forall s.\; (\operatorname{is\text{-}module}(s) \Rightarrow (s \in \mathbf{Set}))}}
\vnbprop{is-vector-space-is-set}{\vnbstmt{\forall s.\; (\operatorname{is\text{-}vector\text{-}space}(s) \Rightarrow (s \in \mathbf{Set}))}}

\vnbbookfile{theorem-library/ag-view-read-offs}{theorem-library/ag-view-read-offs}
\noindent{\small A file shared across books; the entries that belong to this one.}\par
\noindent{\small ag-view-read-offs.scm -- THE ADDITIVE-GROUP VIEW READ-OFFS, PROVEN.}\par
\vnbprop{mvag-carr}{\vnbstmt{\forall d.\; (\operatorname{is\text{-}module}(d) \Rightarrow (\operatorname{carr}(\operatorname{module\text{-}vector\text{-}ag}(d)) = \operatorname{vec}(d)))}}
\vnbprop{mvag-op}{\vnbstmt{\forall d.\; (\operatorname{is\text{-}module}(d) \Rightarrow (\operatorname{opr}(\operatorname{module\text{-}vector\text{-}ag}(d)) = \operatorname{vadd}(d)))}}
\vnbprop{mvag-id}{\vnbstmt{\forall d.\; (\operatorname{is\text{-}module}(d) \Rightarrow (\operatorname{iden}(\operatorname{module\text{-}vector\text{-}ag}(d)) = \operatorname{vzero}(d)))}}

\vnbbookfile{theorem-library/view-typing-laws}{theorem-library/view-typing-laws}
\noindent{\small A file shared across books; the entries that belong to this one.}\par
\noindent{\small theorem-library/view-typing-laws.scm -- THE VIEW TYPING THEOREMS, PROVEN (D8, 2026-11-15).}\par
\vnbprop{module-vector-ag-is-abelian-group}{\vnbstmt{\forall r.\; (\operatorname{is\text{-}module}(r) \Rightarrow \operatorname{is\text{-}abelian\text{-}group}(\operatorname{module\text{-}vector\text{-}ag}(r)))}}

\vnbbookfile{theorem-library/rake-setoid}{theorem-library/rake-setoid}
\noindent{\small A file shared across books; the entries that belong to this one.}\par
\noindent{\small rake-setoid.scm -- BATCH F of the 2026-09-17 rake: setoids and carriers. Seven asserted leaves of the PSS, PROVEN.  Every statement here is its support's statement UNCHANGED.}\par
\vnbprop{is-vector-space-unfold}{\vnbstmt{\forall m.\; (\operatorname{is\text{-}vector\text{-}space}(m) \Leftrightarrow (\operatorname{is\text{-}module}(m) \wedge \operatorname{is\text{-}field\text{-}ring}(\operatorname{scal}(m))))}}
\vnbprop{vspace-vec-is-set}{\vnbstmt{\forall m.\; (\operatorname{is\text{-}vector\text{-}space}(m) \Rightarrow (\operatorname{vec}(m) \in \mathbf{Set}))}}

\vnbbookfile{theorem-library/entry-in-carrier}{theorem-library/entry-in-carrier}
\noindent{\small entry-in-carrier.scm -- the entries of an m-by-n matrix over X lie in X.}\par
\vnbprop{entry-in-carrier}{\vnbstmt{\begin{array}{@{}l@{}} \forall m, n, x, p, i, j. \\ \quad ((p \in \operatorname{mat}(m, n, x)) \wedge ((i \in \operatorname{interval}(1, m)) \wedge (j \in \operatorname{interval}(1, n)))) \Rightarrow  \\ \quad \quad \operatorname{entry}(p, i, j) \in x \end{array}}}

\vnbbookfile{theorem-library/rake-algebra}{theorem-library/rake-algebra}
\noindent{\small A file shared across books; the entries that belong to this one.}\par
\noindent{\small rake-algebra.scm -- BATCH D of the 2026-09-17 rake: seven asserted leaves of the PSS, PROVEN.  Every statement here is its support's statement UNCHANGED.}\par
\vnbprop{diagonal-off-entry}{\vnbstmt{\begin{array}{@{}l@{}} \forall a, m, n, d.; \forall i \in \operatorname{interval}(1, m), j \in \operatorname{interval}(1, n). \\ \quad (\operatorname{is\text{-}diagonal}(a, m, n, d) \wedge \neg (i = j)) \Rightarrow (\operatorname{entry}(d, i, j) = \operatorname{zero}(a)) \end{array}}}
\vnbprop{span-add-one-unfold}{\vnbstmt{\begin{array}{@{}l@{}} \forall m, s, v. \\ \quad (\operatorname{span\text{-}add\text{-}one}(m, s, v) \simeq \{y \in \operatorname{vec}(m) : \exists x \in s, r \in \mathbb{R}.\; (y = \operatorname{vadd}(m)(x, \operatorname{act}(m)(r, v)))\}) \end{array}}}
\vnbprop{span-add-one-membership}{\vnbstmt{\begin{array}{@{}l@{}} \forall m, s, v, w. \\ \quad w \in \operatorname{span\text{-}add\text{-}one}(m, s, v) \\ \quad \Leftrightarrow\; (w \in \operatorname{vec}(m)) \wedge \exists y \in s, r \in \mathbb{R}.\; (w = \operatorname{vadd}(m)(y, \operatorname{act}(m)(r, v))) \end{array}}}

\vnbbookfile{theorem-library/rake-finsum-typing}{theorem-library/rake-finsum-typing}
\noindent{\small A file shared across books; the entries that belong to this one.}\par
\noindent{\small rake-finsum-typing.scm -- the finsum / monoid-power typing leaves of the rake (batch G), PROVEN, together with the guarded form of the one statement in the batch that is NOT provable as stated.}\par
\vnbprop{submodule-finsum-closed}{\vnbstmt{\begin{array}{@{}l@{}} \forall d, m.; \forall s \in \mathbf{Set}, f \in (s \to \operatorname{carr}(\operatorname{module\text{-}vector\text{-}ag}(d))). \\ \quad (\operatorname{is\text{-}module}(d) \wedge (\operatorname{is\text{-}submodule}(d, m) \wedge ((|s| \in \mathbb{N}) \wedge \forall z \in s.\; (f(z) \in m)))) \Rightarrow  \\ \quad \quad \operatorname{finsum}(\operatorname{module\text{-}vector\text{-}ag}(d), f, s) \in m \end{array}}}

\vnbbookfile{theorem-library/hom-laws}{theorem-library/hom-laws}
\noindent{\small A file shared across books; the entries that belong to this one.}\par
\noindent{\small hom-laws.scm -- the models of a structure form a CATEGORY: identity arrows, composition, and the hom-set as a set.}\par
\vnbprop{hom-module-member-iff}{\vnbstmt{\begin{array}{@{}l@{}} \forall a, b, f. \\ \quad (f \in \operatorname{hom\text{-}module}(a, b)) \Leftrightarrow ((f \in (\operatorname{vec}(a) \to \operatorname{vec}(b))) \wedge \operatorname{is\text{-}hom\text{-}module}(a, b, f)) \end{array}}}
\vnbprop{hom-module-id}{\vnbstmt{\forall a.\; (\operatorname{is\text{-}module}(a) \Rightarrow \operatorname{is\text{-}hom\text{-}module}(a, a, \operatorname{id\text{-}fun}(\operatorname{vec}(a))))}}
\vnbprop{hom-module-in-set}{\vnbstmt{\forall a, b.\; ((\operatorname{is\text{-}module}(a) \wedge \operatorname{is\text{-}module}(b)) \Rightarrow (\operatorname{hom\text{-}module}(a, b) \in \mathbf{Set}))}}
\vnbprop{hom-vector-space-member-iff}{\vnbstmt{\begin{array}{@{}l@{}} \forall a, b, f. \\ \quad (f \in \operatorname{hom\text{-}vector\text{-}space}(a, b)) \Leftrightarrow ((f \in (\operatorname{vec}(a) \to \operatorname{vec}(b))) \wedge \operatorname{is\text{-}hom\text{-}vector\text{-}space}(a, b, f)) \end{array}}}
\vnbprop{hom-vector-space-id}{\vnbstmt{\begin{array}{@{}l@{}} \forall a. \\ \quad \operatorname{is\text{-}vector\text{-}space}(a) \Rightarrow \operatorname{is\text{-}hom\text{-}vector\text{-}space}(a, a, \operatorname{id\text{-}fun}(\operatorname{vec}(a))) \end{array}}}
\vnbprop{hom-vector-space-in-set}{\vnbstmt{\begin{array}{@{}l@{}} \forall a, b. \\ \quad (\operatorname{is\text{-}vector\text{-}space}(a) \wedge \operatorname{is\text{-}vector\text{-}space}(b)) \Rightarrow (\operatorname{hom\text{-}vector\text{-}space}(a, b) \in \mathbf{Set}) \end{array}}}

\vnbbookfile{theorem-library/hom-functors}{theorem-library/hom-functors}
\noindent{\small A file shared across books; the entries that belong to this one.}\par
\noindent{\small hom-functors.scm -- the Hom functors: Hom(a, -) covariant, Hom(-, b) contravariant, as maps on hom-sets.}\par
\vnbprop{hom-module-post-type}{\vnbstmt{\begin{array}{@{}l@{}} \forall a, b, c, h. \\ \quad (\operatorname{is\text{-}module}(a) \wedge (h \in \operatorname{hom\text{-}module}(b, c))) \Rightarrow  \\ \quad \quad (\lambda f \in \operatorname{hom\text{-}module}(a, b).\; \operatorname{compose}(h, f)) \in (\operatorname{hom\text{-}module}(a, b) \to \operatorname{hom\text{-}module}(a, c)) \end{array}}}
\vnbprop{hom-module-pre-type}{\vnbstmt{\begin{array}{@{}l@{}} \forall a, b, c, k. \\ \quad (\operatorname{is\text{-}module}(b) \wedge (k \in \operatorname{hom\text{-}module}(c, a))) \Rightarrow  \\ \quad \quad (\lambda f \in \operatorname{hom\text{-}module}(a, b).\; \operatorname{compose}(f, k)) \in (\operatorname{hom\text{-}module}(a, b) \to \operatorname{hom\text{-}module}(c, b)) \end{array}}}
\vnbprop{hom-vector-space-post-type}{\vnbstmt{\begin{array}{@{}l@{}} \forall a, b, c, h. \\ \quad (\operatorname{is\text{-}vector\text{-}space}(a) \wedge (h \in \operatorname{hom\text{-}vector\text{-}space}(b, c))) \Rightarrow  \\ \quad \quad (\lambda f \in \operatorname{hom\text{-}vector\text{-}space}(a, b).\; \operatorname{compose}(h, f)) \\ \quad \quad \in\; (\operatorname{hom\text{-}vector\text{-}space}(a, b) \to \operatorname{hom\text{-}vector\text{-}space}(a, c)) \end{array}}}
\vnbprop{hom-vector-space-pre-type}{\vnbstmt{\begin{array}{@{}l@{}} \forall a, b, c, k. \\ \quad (\operatorname{is\text{-}vector\text{-}space}(b) \wedge (k \in \operatorname{hom\text{-}vector\text{-}space}(c, a))) \Rightarrow  \\ \quad \quad (\lambda f \in \operatorname{hom\text{-}vector\text{-}space}(a, b).\; \operatorname{compose}(f, k)) \\ \quad \quad \in\; (\operatorname{hom\text{-}vector\text{-}space}(a, b) \to \operatorname{hom\text{-}vector\text{-}space}(c, b)) \end{array}}}

\vnbbookfile{theorem-library/qq-vs-exemplification}{theorem-library/qq-vs-exemplification}
\noindent{\small qq-vs-exemplification.scm -- THE FIRST WITNESS OF A VECTOR SPACE.}\par
\par\medskip\noindent\textbf{Definition} (\texttt{qq-line}).\;
\begin{equation*}
\begin{array}{@{}l@{}} \operatorname{qq\text{-}line} \\ =\; {[}\operatorname{field\text{-}as\text{-}integral\text{-}domain}(\operatorname{qq\text{-}field}), \mathbb{Q}, (\lambda {[}\operatorname{x\_}, \operatorname{y\_}] \in (\mathbb{Q} \times \mathbb{Q}).\; +(\operatorname{x\_}, \operatorname{y\_})), 0, (\lambda \operatorname{x\_} \in \mathbb{Q}.\; -(\operatorname{x\_})), (\lambda {[}\operatorname{r\_}, \operatorname{x\_}] \in (\mathbb{Q} \times \mathbb{Q}).\; \cdot(\operatorname{r\_}, \operatorname{x\_}))] \end{array}
\end{equation*}

\vnbdef{qq-line@vec}{\vnbstmt{(\operatorname{vec}(\operatorname{qq\text{-}line}) \simeq \mathbb{Q})}}
\vnbdef{qq-line@act}{\vnbstmt{(\operatorname{act}(\operatorname{qq\text{-}line}) \simeq (\lambda {[}\operatorname{r\_}, \operatorname{x\_}] \in (\mathbb{Q} \times \mathbb{Q}).\; (\operatorname{r\_} \cdot \operatorname{x\_})))}}
\vnbdef{qq-line@vadd}{\vnbstmt{(\operatorname{vadd}(\operatorname{qq\text{-}line}) \simeq (\lambda {[}\operatorname{x\_}, \operatorname{y\_}] \in (\mathbb{Q} \times \mathbb{Q}).\; (\operatorname{x\_} + \operatorname{y\_})))}}
\vnbdef{qq-line@vneg}{\vnbstmt{(\operatorname{vneg}(\operatorname{qq\text{-}line}) \simeq (\lambda x \in \mathbb{Q}.\; -x))}}
\vnbdef{qq-line@scal}{\vnbstmt{(\operatorname{scal}(\operatorname{qq\text{-}line}) \simeq \operatorname{field\text{-}as\text{-}integral\text{-}domain}(\operatorname{qq\text{-}field}))}}
\vnbdef{qq-line@vzero}{\vnbstmt{(\operatorname{vzero}(\operatorname{qq\text{-}line}) \simeq 0)}}
\vnbprop{qq-line-is-module}{\vnbstmt{\operatorname{is\text{-}module}(\operatorname{qq\text{-}line})}}
\vnbprop{qq-line-is-vector-space}{\vnbstmt{\operatorname{is\text{-}vector\text{-}space}(\operatorname{qq\text{-}line})}}

\vnbbookfile{theorem-library/rake-border-entry}{theorem-library/rake-border-entry}
\noindent{\small rake-border-entry.scm -- the BORDER block entry read-offs, PROVEN.}\par
\vnbprop{border-entries-defined}{\vnbstmt{\begin{array}{@{}l@{}} \forall a, b, m, p, q.; \forall c \in \operatorname{interval}(1, (p + 1)), d \in \operatorname{interval}(1, (q + 1)). \\ \quad (\operatorname{is\text{-}ring}(a) \wedge ((p \in \mathbb{N}) \wedge ((q \in \mathbb{N}) \wedge ((b \in \operatorname{carr}(a)) \wedge (m \in \operatorname{mat}(p, q, \operatorname{carr}(a))))))) \Rightarrow  \\ \quad \quad (\lambda {[}i, j] \in (\operatorname{interval}(1, (p + 1)) \times \operatorname{interval}(1, (q + 1))).\; \left(\text{if }(i = 1)\text{ then }\left(\text{if }(j = 1)\text{ then }b\text{ else }\operatorname{zero}(a)\right)\text{ else }\left(\text{if }(j = 1)\text{ then }\operatorname{zero}(a)\text{ else }\operatorname{entry}(m, \operatorname{nn\text{-}minus}(i, 1), \operatorname{nn\text{-}minus}(j, 1))\right)\right))(c, d) \\ \quad \quad \in\; \mathbf{Set} \end{array}}}
\vnbprop{border-entry-block2}{\vnbstmt{\begin{array}{@{}l@{}} \forall a, b, m, p, q, i, j. \\ \quad \operatorname{is\text{-}ring}(a) \wedge \\ \quad p \in \mathbb{N} \wedge \\ \quad q \in \mathbb{N} \wedge \\ \quad b \in \operatorname{carr}(a) \wedge \\ \quad m \in \operatorname{mat}(p, q, \operatorname{carr}(a)) \wedge \\ \quad i \in \operatorname{interval}(1, (p + 1)) \wedge \\ \quad \neg (i = 1) \wedge \\ \quad j \in \operatorname{interval}(1, (q + 1)) \wedge \\ \quad \neg (j = 1) \Rightarrow \\ \quad \quad \operatorname{entry}(\operatorname{border}(a, b, m, p, q), i, j) = \operatorname{entry}(m, \operatorname{nn\text{-}minus}(i, 1), \operatorname{nn\text{-}minus}(j, 1)) \end{array}}}
\vnbprop{border-entry-block}{\vnbstmt{\begin{array}{@{}l@{}} \forall a, b, m, p, q, i, j. \\ \quad \operatorname{is\text{-}ring}(a) \wedge \\ \quad p \in \mathbb{N} \wedge \\ \quad q \in \mathbb{N} \wedge \\ \quad b \in \operatorname{carr}(a) \wedge \\ \quad m \in \operatorname{mat}(p, q, \operatorname{carr}(a)) \wedge \\ \quad i \in \operatorname{interval}(1, p) \wedge \\ \quad j \in \operatorname{interval}(1, q) \Rightarrow \\ \quad \quad \operatorname{entry}(\operatorname{border}(a, b, m, p, q), (i + 1), (j + 1)) = \operatorname{entry}(m, i, j) \end{array}}}

\vnbbookfile{theorem-library/rake-border-siblings}{theorem-library/rake-border-siblings}
\noindent{\small rake-border-siblings.scm -- the three remaining BORDER entry read-offs, PROVEN: border-entry-11, border-entry-1j, border-entry-i1.}\par
\vnbprop{border-entry-11}{\vnbstmt{\begin{array}{@{}l@{}} \forall a, b, m, p, q. \\ \quad (\operatorname{is\text{-}ring}(a) \wedge ((p \in \mathbb{N}) \wedge ((q \in \mathbb{N}) \wedge ((b \in \operatorname{carr}(a)) \wedge (m \in \operatorname{mat}(p, q, \operatorname{carr}(a))))))) \Rightarrow  \\ \quad \quad \operatorname{entry}(\operatorname{border}(a, b, m, p, q), 1, 1) = b \end{array}}}
\vnbprop{border-entry-1j}{\vnbstmt{\begin{array}{@{}l@{}} \forall a, b, m, p, q, j. \\ \quad \operatorname{is\text{-}ring}(a) \wedge \\ \quad p \in \mathbb{N} \wedge \\ \quad q \in \mathbb{N} \wedge \\ \quad b \in \operatorname{carr}(a) \wedge \\ \quad m \in \operatorname{mat}(p, q, \operatorname{carr}(a)) \wedge \\ \quad j \in \operatorname{interval}(1, (q + 1)) \wedge \\ \quad \neg (j = 1) \Rightarrow \\ \quad \quad \operatorname{entry}(\operatorname{border}(a, b, m, p, q), 1, j) = \operatorname{zero}(a) \end{array}}}
\vnbprop{border-entry-i1}{\vnbstmt{\begin{array}{@{}l@{}} \forall a, b, m, p, q, i. \\ \quad \operatorname{is\text{-}ring}(a) \wedge \\ \quad p \in \mathbb{N} \wedge \\ \quad q \in \mathbb{N} \wedge \\ \quad b \in \operatorname{carr}(a) \wedge \\ \quad m \in \operatorname{mat}(p, q, \operatorname{carr}(a)) \wedge \\ \quad i \in \operatorname{interval}(1, (p + 1)) \wedge \\ \quad \neg (i = 1) \Rightarrow \\ \quad \quad \operatorname{entry}(\operatorname{border}(a, b, m, p, q), i, 1) = \operatorname{zero}(a) \end{array}}}
\vnbprop{succ-nn-minus-1}{\vnbstmt{\forall i \in \mathbb{N}.\; ((1 \leq i) \Rightarrow ((\operatorname{nn\text{-}minus}(i, 1) + 1) = i))}}

\vnbbookfile{theorem-library/rake-finsum-laws2}{theorem-library/rake-finsum-laws2}
\noindent{\small A file shared across books; the entries that belong to this one.}\par
\noindent{\small theorem-library/rake-finsum-laws2.scm -- rake batch J, part 2: the three DISTRIBUTION laws of FINSUM (a ring element on either side, a module scalar) and the interval back-peel.  Part 1 is theorem-library/rake-finsum-laws.scm, which must load FIRST: this file cites its `enum-fam-value' and}\par
\vnbprop{module-act-zero-vec}{\vnbstmt{\begin{array}{@{}l@{}} \forall d.; \forall r \in \operatorname{carr}(\operatorname{scal}(d)). \\ \quad \operatorname{is\text{-}module}(d) \Rightarrow (\operatorname{act}(d)(r, \operatorname{vzero}(d)) = \operatorname{vzero}(d)) \end{array}}}
\vnbprop{sum-ag-act-ind}{\vnbstmt{\begin{array}{@{}l@{}} \forall n \in \mathbb{N}, d.; \forall r \in \operatorname{carr}(\operatorname{scal}(d)), a, b. \\ \quad \operatorname{is\text{-}module}(d) \wedge \\ \quad \forall i \in \mathrm{OS}(n).\; (a(i) \in \operatorname{carr}(\operatorname{module\text{-}vector\text{-}ag}(d))) \wedge \\ \quad \forall i \in \mathrm{OS}(n).\; (b(i) \simeq \operatorname{act}(d)(r, a(i))) \Rightarrow \\ \quad \quad \operatorname{sum\text{-}ag}(\operatorname{module\text{-}vector\text{-}ag}(d), b, n) = \operatorname{act}(d)(r, \operatorname{sum\text{-}ag}(\operatorname{module\text{-}vector\text{-}ag}(d), a, n)) \end{array}}}
\vnbprop{finsum-act-distrib-gen}{\vnbstmt{\begin{array}{@{}l@{}} \forall d.; \forall r \in \operatorname{carr}(\operatorname{scal}(d)), s \in \mathbf{Set}, f \in (s \to \operatorname{carr}(\operatorname{module\text{-}vector\text{-}ag}(d))). \\ \quad (\operatorname{is\text{-}module}(d) \wedge (|s| \in \mathbb{N})) \Rightarrow  \\ \quad \quad \operatorname{act}(d)(r, \operatorname{finsum}(\operatorname{module\text{-}vector\text{-}ag}(d), f, s)) \\ \quad \quad =\; \operatorname{finsum}(\operatorname{module\text{-}vector\text{-}ag}(d), (\lambda z \in s.\; \operatorname{act}(d)(r, f(z))), s) \end{array}}}

\vnbbookfile{theorem-library/rake-finsum-core}{theorem-library/rake-finsum-core}
\noindent{\small A file shared across books; the entries that belong to this one.}\par
\noindent{\small theorem-library/rake-finsum-core.scm -- the FINSUM CORE, proven. Rake batch M (2026-09-17).  Nine of the batch's ten leaves; the tenth (finsum-embed) is BLOCKED and its obstacle is stated at the end of this header.}\par
\vnbprop{sum-ag-act-collect-ind}{\vnbstmt{\begin{array}{@{}l@{}} \forall n \in \mathbb{N}, d.; \forall x \in \operatorname{vec}(d), c, v. \\ \quad \operatorname{is\text{-}module}(d) \wedge \\ \quad \forall i \in \mathrm{OS}(n).\; (c(i) \in \operatorname{carr}(\operatorname{scal}(d))) \wedge \\ \quad \forall i \in \mathrm{OS}(n).\; (v(i) \simeq \operatorname{act}(d)(c(i), x)) \Rightarrow \\ \quad \quad \operatorname{sum\text{-}ag}(\operatorname{module\text{-}vector\text{-}ag}(d), v, n) = \operatorname{act}(d)(\operatorname{sum\text{-}ag}(\operatorname{ring\text{-}additive\text{-}ag}(\operatorname{scal}(d)), c, n), x) \end{array}}}
\vnbprop{finsum-act-collect-gen}{\vnbstmt{\begin{array}{@{}l@{}} \forall d.; \forall x \in \operatorname{vec}(d), s \in \mathbf{Set}, c \in (s \to \operatorname{carr}(\operatorname{scal}(d))). \\ \quad (\operatorname{is\text{-}module}(d) \wedge (|s| \in \mathbb{N})) \Rightarrow  \\ \quad \quad \operatorname{act}(d)(\operatorname{finsum}(\operatorname{ring\text{-}additive\text{-}ag}(\operatorname{scal}(d)), c, s), x) \\ \quad \quad =\; \operatorname{finsum}(\operatorname{module\text{-}vector\text{-}ag}(d), (\lambda z \in s.\; \operatorname{act}(d)(c(z), x)), s) \end{array}}}

\vnbbookfile{theorem-library/mat-typing-bundle}{theorem-library/mat-typing-bundle}
\noindent{\small mat-typing-bundle.scm -- three MATOF typings, PROVEN from the definitions.}\par
\vnbprop{mat-colcount-transfer}{\vnbstmt{\begin{array}{@{}l@{}} \forall n, k, y, q, m, x, r. \\ \quad ((q \in \operatorname{mat}(n, k, y)) \wedge (((n = 0) \Rightarrow ((m = 0) \vee (k = 0))) \wedge (r \in \operatorname{mat}(m, \operatorname{nth}(2, \operatorname{size}(q)), x)))) \Rightarrow  \\ \quad \quad r \in \operatorname{mat}(m, k, x) \end{array}}}
\vnbprop{identmat-type}{\vnbstmt{\forall a.; \forall n \in \mathbb{N}.\; (\operatorname{is\text{-}ring}(a) \Rightarrow (\operatorname{identmat}(a, n) \in \operatorname{mat}(n, n, \operatorname{carr}(a))))}}
\vnbprop{elem-g-type}{\vnbstmt{\begin{array}{@{}l@{}} \forall a, n, r, k, l. \\ \quad (\operatorname{is\text{-}ring}(a) \wedge ((n \in \mathbb{N}) \wedge (r \in \operatorname{carr}(a)))) \Rightarrow  \\ \quad \quad \operatorname{elem\text{-}g}(a, n, r, k, l) \in \operatorname{mat}(n, n, \operatorname{carr}(a)) \end{array}}}
\vnbprop{zeromat-type}{\vnbstmt{\begin{array}{@{}l@{}} \forall a, m, n. \\ \quad (\operatorname{is\text{-}ring}(a) \wedge ((m \in \mathbb{N}) \wedge (n \in \mathbb{N}))) \Rightarrow  \\ \quad \quad \operatorname{zeromat}(a, m, n) \in \operatorname{mat}(m, n, \operatorname{carr}(a)) \end{array}}}
\vnbprop{unitrow-type}{\vnbstmt{\begin{array}{@{}l@{}} \forall a, n, i. \\ \quad (\operatorname{is\text{-}ring}(a) \wedge (n \in \mathbb{N})) \Rightarrow (\operatorname{unitrow}(a, n, i) \in \operatorname{mat}(1, n, \operatorname{carr}(a))) \end{array}}}
\vnbprop{elem-h-type}{\vnbstmt{\begin{array}{@{}l@{}} \forall a, n, r, k. \\ \quad (\operatorname{is\text{-}ring}(a) \wedge ((n \in \mathbb{N}) \wedge (r \in \operatorname{carr}(a)))) \Rightarrow  \\ \quad \quad \operatorname{elem\text{-}h}(a, n, r, k) \in \operatorname{mat}(n, n, \operatorname{carr}(a)) \end{array}}}
\vnbprop{matmul-type}{\vnbstmt{\begin{array}{@{}l@{}} \forall a, m, n, k, p, q. \\ \quad \operatorname{is\text{-}ring}(a) \wedge \\ \quad p \in \operatorname{mat}(m, n, \operatorname{carr}(a)) \wedge \\ \quad q \in \operatorname{mat}(n, k, \operatorname{carr}(a)) \wedge \\ \quad (n = 0) \Rightarrow ((m = 0) \vee (k = 0)) \Rightarrow \\ \quad \quad \operatorname{matmul}(a, p, q) \in \operatorname{mat}(m, k, \operatorname{carr}(a)) \end{array}}}

\vnbbookfile{theorem-library/rake-mat-typing}{theorem-library/rake-mat-typing}
\noindent{\small rake-mat-typing.scm -- the ELEVEN matrix typing / sethood supports, PROVEN.}\par
\vnbprop{mat-is-set}{\vnbstmt{\forall x \in \mathbf{Set}, m, n.\; (\operatorname{mat}(m, n, x) \in \mathbf{Set})}}
\vnbprop{block-type}{\vnbstmt{\begin{array}{@{}l@{}} \forall m, n, x, p, k, l. \\ \quad ((k \in \mathbb{N}) \wedge ((l \in \mathbb{N}) \wedge ((p \in \operatorname{mat}(m, n, x)) \wedge ((k \leq m) \wedge (l \leq n))))) \Rightarrow  \\ \quad \quad \operatorname{block}(p, k, l) \in \operatorname{mat}(k, l, x) \end{array}}}
\vnbprop{snoc-col-type}{\vnbstmt{\begin{array}{@{}l@{}} \forall x, n, w, v. \\ \quad ((n \in \mathbb{N}) \wedge ((w \in \operatorname{mat}(n, 1, x)) \wedge (v \in x))) \Rightarrow  \\ \quad \quad \operatorname{snoc\text{-}col}(w, n, v) \in \operatorname{mat}((n + 1), 1, x) \end{array}}}
\vnbprop{snoc-row-type}{\vnbstmt{\begin{array}{@{}l@{}} \forall x, n, c, r. \\ \quad ((n \in \mathbb{N}) \wedge ((c \in \operatorname{mat}(1, n, x)) \wedge (r \in x))) \Rightarrow  \\ \quad \quad \operatorname{snoc\text{-}row}(c, n, r) \in \operatorname{mat}(1, (n + 1), x) \end{array}}}
\vnbprop{matadd-type}{\vnbstmt{\begin{array}{@{}l@{}} \forall a, m, n, p, q. \\ \quad (\operatorname{is\text{-}ring}(a) \wedge ((p \in \operatorname{mat}(m, n, \operatorname{carr}(a))) \wedge (q \in \operatorname{mat}(m, n, \operatorname{carr}(a))))) \Rightarrow  \\ \quad \quad \operatorname{matadd}(a, p, q) \in \operatorname{mat}(m, n, \operatorname{carr}(a)) \end{array}}}
\vnbprop{matscale-type}{\vnbstmt{\begin{array}{@{}l@{}} \forall a, m, n, r, p. \\ \quad (\operatorname{is\text{-}ring}(a) \wedge ((r \in \operatorname{carr}(a)) \wedge (p \in \operatorname{mat}(m, n, \operatorname{carr}(a))))) \Rightarrow  \\ \quad \quad \operatorname{matscale}(a, r, p) \in \operatorname{mat}(m, n, \operatorname{carr}(a)) \end{array}}}
\vnbprop{matneg-type}{\vnbstmt{\begin{array}{@{}l@{}} \forall a, m, n.; \forall p \in \operatorname{mat}(m, n, \operatorname{carr}(a)). \\ \quad \operatorname{is\text{-}ring}(a) \Rightarrow (\operatorname{matneg}(a, p) \in \operatorname{mat}(m, n, \operatorname{carr}(a))) \end{array}}}
\vnbprop{submat-type}{\vnbstmt{\begin{array}{@{}l@{}} \forall a, p, q, s. \\ \quad ((p \in \mathbb{N}) \wedge ((q \in \mathbb{N}) \wedge (s \in \operatorname{mat}((p + 1), (q + 1), \operatorname{carr}(a))))) \Rightarrow  \\ \quad \quad \operatorname{submat}(s, p, q) \in \operatorname{mat}(p, q, \operatorname{carr}(a)) \end{array}}}
\vnbprop{border-type}{\vnbstmt{\begin{array}{@{}l@{}} \forall a, b, m, p, q. \\ \quad (\operatorname{is\text{-}ring}(a) \wedge ((p \in \mathbb{N}) \wedge ((q \in \mathbb{N}) \wedge ((b \in \operatorname{carr}(a)) \wedge (m \in \operatorname{mat}(p, q, \operatorname{carr}(a))))))) \Rightarrow  \\ \quad \quad \operatorname{border}(a, b, m, p, q) \in \operatorname{mat}((p + 1), (q + 1), \operatorname{carr}(a)) \end{array}}}
\vnbprop{minor-type}{\vnbstmt{\begin{array}{@{}l@{}} \forall r, s, p, q, n. \\ \quad (\operatorname{is\text{-}ring}(r) \wedge ((p \in \mathbb{N}) \wedge ((q \in \mathbb{N}) \wedge ((n \in \mathbb{N}) \wedge (s \in \operatorname{mat}((n + 1), (n + 1), \operatorname{carr}(r))))))) \Rightarrow  \\ \quad \quad \operatorname{minor}(s, p, q, n) \in \operatorname{mat}(n, n, \operatorname{carr}(r)) \end{array}}}
\vnbprop{rmm-carr}{\vnbstmt{\forall r.\; (\operatorname{is\text{-}ring}(r) \Rightarrow (\operatorname{carr}(\operatorname{ring\text{-}multiplicative\text{-}monoid}(r)) = \operatorname{carr}(r)))}}
\vnbprop{det-in-carrier-ind}{\vnbstmt{\begin{array}{@{}l@{}} \forall r.; \forall n \in \mathbb{N}, a \in \operatorname{mat}(n, n, \operatorname{carr}(r)). \\ \quad \operatorname{is\text{-}ring}(r) \Rightarrow (\operatorname{det}(r, n, a) \in \operatorname{carr}(r)) \end{array}}}
\vnbprop{det-in-carrier}{\vnbstmt{\begin{array}{@{}l@{}} \forall r, n, a. \\ \quad (\operatorname{is\text{-}ring}(r) \wedge ((n \in \mathbb{N}) \wedge (a \in \operatorname{mat}(n, n, \operatorname{carr}(r))))) \Rightarrow  \\ \quad \quad \operatorname{det}(r, n, a) \in \operatorname{carr}(r) \end{array}}}

\vnbbookfile{theorem-library/rake-det-small}{theorem-library/rake-det-small}
\noindent{\small rake-det-small.scm -- the small determinants, PROVEN (rake batch 5b-C, 2026-09-18).}\par
\vnbprop{interval-1-1}{\vnbstmt{\operatorname{interval}(1, 1) = \{1\}}}
\vnbprop{rmm-op}{\vnbstmt{\forall r.\; (\operatorname{is\text{-}ring}(r) \Rightarrow (\operatorname{opr}(\operatorname{ring\text{-}multiplicative\text{-}monoid}(r)) = \operatorname{mul}(r)))}}
\vnbprop{rmm-iden}{\vnbstmt{\forall r.\; (\operatorname{is\text{-}ring}(r) \Rightarrow (\operatorname{iden}(\operatorname{ring\text{-}multiplicative\text{-}monoid}(r)) = \operatorname{one}(r)))}}
\vnbprop{ring-mpow-one}{\vnbstmt{\begin{array}{@{}l@{}} \forall r.; \forall x \in \operatorname{carr}(r). \\ \quad \operatorname{is\text{-}ring}(r) \Rightarrow (\operatorname{mpow}(\operatorname{ring\text{-}multiplicative\text{-}monoid}(r), x, 1) = x) \end{array}}}
\vnbprop{minor-entry}{\vnbstmt{\begin{array}{@{}l@{}} \forall r, s, p, q, n, u, v. \\ \quad \operatorname{is\text{-}ring}(r) \wedge \\ \quad n \in \mathbb{N} \wedge \\ \quad s \in \operatorname{mat}((n + 1), (n + 1), \operatorname{carr}(r)) \wedge \\ \quad u \in \operatorname{interval}(1, n) \wedge \\ \quad v \in \operatorname{interval}(1, n) \Rightarrow \\ \quad \quad \operatorname{entry}(\operatorname{minor}(s, p, q, n), u, v) = \operatorname{entry}(s, \left(\text{if }(u < p)\text{ then }u\text{ else }(u + 1)\right), \left(\text{if }(v < q)\text{ then }v\text{ else }(v + 1)\right)) \end{array}}}
\vnbprop{det-1x1}{\vnbstmt{\begin{array}{@{}l@{}} \forall r, a. \\ \quad (\operatorname{is\text{-}ring}(r) \wedge (a \in \operatorname{mat}(1, 1, \operatorname{carr}(r)))) \Rightarrow (\operatorname{det}(r, 1, a) = \operatorname{entry}(a, 1, 1)) \end{array}}}
\vnbprop{det-2x2}{\vnbstmt{\begin{array}{@{}l@{}} \forall r, a. \\ \quad (\operatorname{is\text{-}ring}(r) \wedge (a \in \operatorname{mat}(2, 2, \operatorname{carr}(r)))) \Rightarrow  \\ \quad \quad \operatorname{det}(r, 2, a) \\ \quad \quad =\; \operatorname{add}(r)(\operatorname{mul}(r)(\operatorname{entry}(a, 1, 1), \operatorname{entry}(a, 2, 2)), \\ \quad \quad \quad \operatorname{neg}(r)(\operatorname{mul}(r)(\operatorname{entry}(a, 1, 2), \operatorname{entry}(a, 2, 1)))) \end{array}}}
\vnbprop{identmat-minor}{\vnbstmt{\begin{array}{@{}l@{}} \forall r, n. \\ \quad (\operatorname{is\text{-}ring}(r) \wedge (n \in \mathbb{N})) \Rightarrow  \\ \quad \quad \operatorname{minor}(\operatorname{identmat}(r, (n + 1)), 1, 1, n) = \operatorname{identmat}(r, n) \end{array}}}
\vnbprop{mat-ring-one}{\vnbstmt{\begin{array}{@{}l@{}} \forall a, n. \\ \quad (\operatorname{is\text{-}ring}(a) \wedge (n \in \mathbb{N})) \Rightarrow (\operatorname{one}(\operatorname{mat\text{-}ring}(a, n)) = \operatorname{identmat}(a, n)) \end{array}}}
\vnbprop{det-identity-ind}{\vnbstmt{\forall r.; \forall n \in \mathbb{N}.\; (\operatorname{is\text{-}ring}(r) \Rightarrow (\operatorname{det}(r, n, \operatorname{identmat}(r, n)) = \operatorname{one}(r)))}}
\vnbprop{det-identity}{\vnbstmt{\begin{array}{@{}l@{}} \forall r, n. \\ \quad (\operatorname{is\text{-}ring}(r) \wedge (n \in \mathbb{N})) \Rightarrow (\operatorname{det}(r, n, \operatorname{one}(\operatorname{mat\text{-}ring}(r, n))) = \operatorname{one}(r)) \end{array}}}

\vnbbookfile{theorem-library/rake-mat-ring-readoffs}{theorem-library/rake-mat-ring-readoffs}
\noindent{\small rake-mat-ring-readoffs.scm -- the MAT-RING slot read-offs and the two curried-operation FUN typings, PROVEN (rake, 2026-09-19).}\par
\vnbprop{mat-ring-carr}{\vnbstmt{\begin{array}{@{}l@{}} \forall a, n. \\ \quad (\operatorname{is\text{-}ring}(a) \wedge (n \in \mathbb{N})) \Rightarrow (\operatorname{carr}(\operatorname{mat\text{-}ring}(a, n)) = \operatorname{mat}(n, n, \operatorname{carr}(a))) \end{array}}}
\vnbprop{mat-ring-add}{\vnbstmt{\begin{array}{@{}l@{}} \forall a, n. \\ \quad (\operatorname{is\text{-}ring}(a) \wedge (n \in \mathbb{N})) \Rightarrow  \\ \quad \quad \operatorname{add}(\operatorname{mat\text{-}ring}(a, n)) \\ \quad \quad =\; (\lambda {[}p, q] \in (\operatorname{mat}(n, n, \operatorname{carr}(a)) \times \operatorname{mat}(n, n, \operatorname{carr}(a))).\; \operatorname{matadd}(a, p, q)) \end{array}}}
\vnbprop{mat-ring-mul}{\vnbstmt{\begin{array}{@{}l@{}} \forall a, n. \\ \quad (\operatorname{is\text{-}ring}(a) \wedge (n \in \mathbb{N})) \Rightarrow  \\ \quad \quad \operatorname{mul}(\operatorname{mat\text{-}ring}(a, n)) \\ \quad \quad =\; (\lambda {[}p, q] \in (\operatorname{mat}(n, n, \operatorname{carr}(a)) \times \operatorname{mat}(n, n, \operatorname{carr}(a))).\; \operatorname{matmul}(a, p, q)) \end{array}}}
\vnbprop{mat-ring-neg}{\vnbstmt{\begin{array}{@{}l@{}} \forall a, n. \\ \quad (\operatorname{is\text{-}ring}(a) \wedge (n \in \mathbb{N})) \Rightarrow  \\ \quad \quad \operatorname{neg}(\operatorname{mat\text{-}ring}(a, n)) = (\lambda p \in \operatorname{mat}(n, n, \operatorname{carr}(a)).\; \operatorname{matneg}(a, p)) \end{array}}}
\vnbprop{mat-ring-zero}{\vnbstmt{\begin{array}{@{}l@{}} \forall a, n. \\ \quad (\operatorname{is\text{-}ring}(a) \wedge (n \in \mathbb{N})) \Rightarrow (\operatorname{zero}(\operatorname{mat\text{-}ring}(a, n)) = \operatorname{zeromat}(a, n, n)) \end{array}}}
\vnbprop{mat-ring-add-fun}{\vnbstmt{\begin{array}{@{}l@{}} \forall a.; \forall n \in \mathbb{N}. \\ \quad \operatorname{is\text{-}ring}(a) \Rightarrow  \\ \quad \quad (\lambda {[}p, q] \in (\operatorname{mat}(n, n, \operatorname{carr}(a)) \times \operatorname{mat}(n, n, \operatorname{carr}(a))).\; \operatorname{matadd}(a, p, q)) \\ \quad \quad \in\; ((\operatorname{mat}(n, n, \operatorname{carr}(a)) \times \operatorname{mat}(n, n, \operatorname{carr}(a))) \to \operatorname{mat}(n, n, \operatorname{carr}(a))) \end{array}}}
\vnbprop{mat-ring-mul-fun}{\vnbstmt{\begin{array}{@{}l@{}} \forall a.; \forall n \in \mathbb{N}. \\ \quad \operatorname{is\text{-}ring}(a) \Rightarrow  \\ \quad \quad (\lambda {[}p, q] \in (\operatorname{mat}(n, n, \operatorname{carr}(a)) \times \operatorname{mat}(n, n, \operatorname{carr}(a))).\; \operatorname{matmul}(a, p, q)) \\ \quad \quad \in\; ((\operatorname{mat}(n, n, \operatorname{carr}(a)) \times \operatorname{mat}(n, n, \operatorname{carr}(a))) \to \operatorname{mat}(n, n, \operatorname{carr}(a))) \end{array}}}

\vnbbookfile{theorem-library/matmul-entry-proof}{theorem-library/matmul-entry-proof}
\noindent{\small matmul-entry-proof.scm -- the matrix-product entry formula, PROVEN.}\par
\vnbprop{matmul-entry}{\vnbstmt{\begin{array}{@{}l@{}} \forall a, m, n, k, p, q.; \forall i \in \operatorname{interval}(1, m), c \in \operatorname{interval}(1, k). \\ \quad (\operatorname{is\text{-}ring}(a) \wedge ((p \in \operatorname{mat}(m, n, \operatorname{carr}(a))) \wedge ((q \in \operatorname{mat}(n, k, \operatorname{carr}(a))) \wedge (1 \leq n)))) \Rightarrow  \\ \quad \quad \operatorname{entry}(\operatorname{matmul}(a, p, q), i, c) \\ \quad \quad =\; \operatorname{finsum}(\operatorname{ring\text{-}additive\text{-}ag}(a), \\ \quad \quad \quad (\lambda j \in \operatorname{interval}(1, n).\; \operatorname{mul}(a)(\operatorname{entry}(p, i, j), \operatorname{entry}(q, j, c))), \\ \quad \quad \quad \operatorname{interval}(1, n)) \end{array}}}

\vnbbookfile{theorem-library/matunit-matact-type}{theorem-library/matunit-matact-type}
\noindent{\small matunit-matact-type.scm -- two MATOF typings, PROVEN from the definitions.}\par
\vnbprop{matunit-type}{\vnbstmt{\begin{array}{@{}l@{}} \forall a, n, k, l. \\ \quad (\operatorname{is\text{-}ring}(a) \wedge (n \in \mathbb{N})) \Rightarrow (\operatorname{matunit}(a, n, k, l) \in \operatorname{mat}(n, n, \operatorname{carr}(a))) \end{array}}}
\vnbprop{matact-type}{\vnbstmt{\begin{array}{@{}l@{}} \forall d, m, n, q, p, u. \\ \quad \operatorname{is\text{-}module}(d) \wedge \\ \quad p \in \operatorname{mat}(m, n, \operatorname{carr}(\operatorname{scal}(d))) \wedge \\ \quad u \in \operatorname{mat}(n, q, \operatorname{vec}(d)) \wedge \\ \quad (n = 0) \Rightarrow ((m = 0) \vee (q = 0)) \Rightarrow \\ \quad \quad \operatorname{matact}(d, p, u) \in \operatorname{mat}(m, q, \operatorname{vec}(d)) \end{array}}}

\vnbbookfile{theorem-library/lam-fun-bricks}{theorem-library/lam-fun-bricks}
\noindent{\small lam-fun-bricks.scm -- the lambda FUN-typing bricks of the matrix arc, PROVEN.}\par
\vnbprop{matprod-summand-type}{\vnbstmt{\begin{array}{@{}l@{}} \forall a, m, n, b, p, q, i, c. \\ \quad \operatorname{is\text{-}ring}(a) \wedge \\ \quad p \in \operatorname{mat}(m, n, \operatorname{carr}(a)) \wedge \\ \quad q \in \operatorname{mat}(n, b, \operatorname{carr}(a)) \wedge \\ \quad i \in \operatorname{interval}(1, m) \wedge \\ \quad c \in \operatorname{interval}(1, b) \Rightarrow \\ \quad \quad (\lambda j \in \operatorname{interval}(1, n).\; \operatorname{mul}(a)(\operatorname{entry}(p, i, j), \operatorname{entry}(q, j, c))) \\ \quad \quad \in\; (\operatorname{interval}(1, n) \to \operatorname{carr}(\operatorname{ring\text{-}additive\text{-}ag}(a))) \end{array}}}
\vnbprop{matmul-assoc-summand-type}{\vnbstmt{\begin{array}{@{}l@{}} \forall a, m, n, k, l, p, q, r, w, b. \\ \quad \operatorname{is\text{-}ring}(a) \wedge \\ \quad p \in \operatorname{mat}(m, n, \operatorname{carr}(a)) \wedge \\ \quad q \in \operatorname{mat}(n, k, \operatorname{carr}(a)) \wedge \\ \quad r \in \operatorname{mat}(k, l, \operatorname{carr}(a)) \wedge \\ \quad w \in \operatorname{interval}(1, m) \wedge \\ \quad b \in \operatorname{interval}(1, l) \Rightarrow \\ \quad \quad (\lambda z \in (\operatorname{interval}(1, k) \times \operatorname{interval}(1, n)).\; \\ \quad \quad \quad \operatorname{mul}(a)(\operatorname{mul}(a)(\operatorname{entry}(p, w, \operatorname{nth}(2, z)), \operatorname{entry}(q, \operatorname{nth}(2, z), \operatorname{nth}(1, z))), \operatorname{entry}(r, \operatorname{nth}(1, z), b))) \\ \quad \quad \in\; ((\operatorname{interval}(1, k) \times \operatorname{interval}(1, n)) \to \operatorname{carr}(\operatorname{ring\text{-}additive\text{-}ag}(a))) \end{array}}}
\vnbprop{matact-summand-type}{\vnbstmt{\begin{array}{@{}l@{}} \forall d, m, n, q, p, u, i, c. \\ \quad \operatorname{is\text{-}module}(d) \wedge \\ \quad p \in \operatorname{mat}(m, n, \operatorname{carr}(\operatorname{scal}(d))) \wedge \\ \quad u \in \operatorname{mat}(n, q, \operatorname{vec}(d)) \wedge \\ \quad i \in \operatorname{interval}(1, m) \wedge \\ \quad c \in \operatorname{interval}(1, q) \Rightarrow \\ \quad \quad (\lambda j \in \operatorname{interval}(1, n).\; \operatorname{act}(d)(\operatorname{entry}(p, i, j), \operatorname{entry}(u, j, c))) \\ \quad \quad \in\; (\operatorname{interval}(1, n) \to \operatorname{carr}(\operatorname{module\text{-}vector\text{-}ag}(d))) \end{array}}}
\vnbprop{matunit-summand-type}{\vnbstmt{\begin{array}{@{}l@{}} \forall a, m, n, p, k, l, i, c. \\ \quad (\operatorname{is\text{-}ring}(a) \wedge ((p \in \operatorname{mat}(m, n, \operatorname{carr}(a))) \wedge ((i \in \operatorname{interval}(1, m)) \wedge (c \in \operatorname{interval}(1, n))))) \Rightarrow  \\ \quad \quad (\lambda j \in \operatorname{interval}(1, n).\; \operatorname{mul}(a)(\operatorname{entry}(p, i, j), \operatorname{entry}(\operatorname{matunit}(a, n, k, l), j, c))) \\ \quad \quad \in\; (\operatorname{interval}(1, n) \to \operatorname{carr}(\operatorname{ring\text{-}additive\text{-}ag}(a))) \end{array}}}
\vnbprop{tel-outf-type}{\vnbstmt{\begin{array}{@{}l@{}} \forall a, m, n, k, l, p, q, r, w, b. \\ \quad \operatorname{is\text{-}ring}(a) \wedge \\ \quad p \in \operatorname{mat}(m, n, \operatorname{carr}(a)) \wedge \\ \quad q \in \operatorname{mat}(n, k, \operatorname{carr}(a)) \wedge \\ \quad r \in \operatorname{mat}(k, l, \operatorname{carr}(a)) \wedge \\ \quad w \in \operatorname{interval}(1, m) \wedge \\ \quad b \in \operatorname{interval}(1, l) \wedge \\ \quad 1 \leq n \Rightarrow \\ \quad \quad (\lambda j \in \operatorname{interval}(1, k).\; \operatorname{mul}(a)(\operatorname{entry}(\operatorname{matmul}(a, p, q), w, j), \operatorname{entry}(r, j, b))) \\ \quad \quad \in\; (\operatorname{interval}(1, k) \to \operatorname{carr}(\operatorname{ring\text{-}additive\text{-}ag}(a))) \end{array}}}
\vnbprop{tel-tout-type}{\vnbstmt{\begin{array}{@{}l@{}} \forall a, m, n, k, l, p, q, r, w, b. \\ \quad \operatorname{is\text{-}ring}(a) \wedge \\ \quad p \in \operatorname{mat}(m, n, \operatorname{carr}(a)) \wedge \\ \quad q \in \operatorname{mat}(n, k, \operatorname{carr}(a)) \wedge \\ \quad r \in \operatorname{mat}(k, l, \operatorname{carr}(a)) \wedge \\ \quad w \in \operatorname{interval}(1, m) \wedge \\ \quad b \in \operatorname{interval}(1, l) \Rightarrow \\ \quad \quad (\lambda c \in \operatorname{interval}(1, k).\; \\ \quad \quad \quad \operatorname{finsum}(\operatorname{ring\text{-}additive\text{-}ag}(a), \\ \quad \quad \quad \quad (\lambda j \in \operatorname{interval}(1, n).\; \\ \quad \quad \quad \quad \quad (\lambda z \in (\operatorname{interval}(1, k) \times \operatorname{interval}(1, n)).\; \operatorname{mul}(a)(\operatorname{mul}(a)(\operatorname{entry}(p, w, \operatorname{nth}(2, z)), \operatorname{entry}(q, \operatorname{nth}(2, z), \operatorname{nth}(1, z))), \operatorname{entry}(r, \operatorname{nth}(1, z), b)))({[}c, j])), \\ \quad \quad \quad \quad \operatorname{interval}(1, n))) \\ \quad \quad \in\; (\operatorname{interval}(1, k) \to \operatorname{carr}(\operatorname{ring\text{-}additive\text{-}ag}(a))) \end{array}}}
\vnbprop{tel-inf-type}{\vnbstmt{\begin{array}{@{}l@{}} \forall a, m, n, k, l, p, q, r, w, b, x. \\ \quad \operatorname{is\text{-}ring}(a) \wedge \\ \quad p \in \operatorname{mat}(m, n, \operatorname{carr}(a)) \wedge \\ \quad q \in \operatorname{mat}(n, k, \operatorname{carr}(a)) \wedge \\ \quad r \in \operatorname{mat}(k, l, \operatorname{carr}(a)) \wedge \\ \quad w \in \operatorname{interval}(1, m) \wedge \\ \quad b \in \operatorname{interval}(1, l) \wedge \\ \quad x \in \operatorname{interval}(1, k) \Rightarrow \\ \quad \quad (\lambda j \in \operatorname{interval}(1, n).\; \operatorname{mul}(a)(\operatorname{entry}(p, w, j), \operatorname{entry}(q, j, x))) \\ \quad \quad \in\; (\operatorname{interval}(1, n) \to \operatorname{carr}(a)) \end{array}}}
\vnbprop{tel-dist-type}{\vnbstmt{\begin{array}{@{}l@{}} \forall a, m, n, k, l, p, q, r, w, b, x. \\ \quad \operatorname{is\text{-}ring}(a) \wedge \\ \quad p \in \operatorname{mat}(m, n, \operatorname{carr}(a)) \wedge \\ \quad q \in \operatorname{mat}(n, k, \operatorname{carr}(a)) \wedge \\ \quad r \in \operatorname{mat}(k, l, \operatorname{carr}(a)) \wedge \\ \quad w \in \operatorname{interval}(1, m) \wedge \\ \quad b \in \operatorname{interval}(1, l) \wedge \\ \quad x \in \operatorname{interval}(1, k) \Rightarrow \\ \quad \quad (\lambda z \in \operatorname{interval}(1, n).\; \\ \quad \quad \quad \operatorname{mul}(a)((\lambda j \in \operatorname{interval}(1, n).\; \operatorname{mul}(a)(\operatorname{entry}(p, w, j), \operatorname{entry}(q, j, x)))(z), \\ \quad \quad \quad \quad \operatorname{entry}(r, x, b))) \\ \quad \quad \in\; (\operatorname{interval}(1, n) \to \operatorname{carr}(\operatorname{ring\text{-}additive\text{-}ag}(a))) \end{array}}}
\vnbprop{tel-red-type}{\vnbstmt{\begin{array}{@{}l@{}} \forall a, m, n, k, l, p, q, r, w, b, x. \\ \quad \operatorname{is\text{-}ring}(a) \wedge \\ \quad p \in \operatorname{mat}(m, n, \operatorname{carr}(a)) \wedge \\ \quad q \in \operatorname{mat}(n, k, \operatorname{carr}(a)) \wedge \\ \quad r \in \operatorname{mat}(k, l, \operatorname{carr}(a)) \wedge \\ \quad w \in \operatorname{interval}(1, m) \wedge \\ \quad b \in \operatorname{interval}(1, l) \wedge \\ \quad x \in \operatorname{interval}(1, k) \Rightarrow \\ \quad \quad (\lambda j \in \operatorname{interval}(1, n).\; \operatorname{mul}(a)(\operatorname{mul}(a)(\operatorname{entry}(p, w, j), \operatorname{entry}(q, j, x)), \operatorname{entry}(r, x, b))) \\ \quad \quad \in\; (\operatorname{interval}(1, n) \to \operatorname{carr}(\operatorname{ring\text{-}additive\text{-}ag}(a))) \end{array}}}
\vnbprop{ter-outf-type}{\vnbstmt{\begin{array}{@{}l@{}} \forall a, m, n, k, l, p, q, r, w, b. \\ \quad \operatorname{is\text{-}ring}(a) \wedge \\ \quad p \in \operatorname{mat}(m, n, \operatorname{carr}(a)) \wedge \\ \quad q \in \operatorname{mat}(n, k, \operatorname{carr}(a)) \wedge \\ \quad r \in \operatorname{mat}(k, l, \operatorname{carr}(a)) \wedge \\ \quad w \in \operatorname{interval}(1, m) \wedge \\ \quad b \in \operatorname{interval}(1, l) \wedge \\ \quad 1 \leq k \Rightarrow \\ \quad \quad (\lambda j \in \operatorname{interval}(1, n).\; \operatorname{mul}(a)(\operatorname{entry}(p, w, j), \operatorname{entry}(\operatorname{matmul}(a, q, r), j, b))) \\ \quad \quad \in\; (\operatorname{interval}(1, n) \to \operatorname{carr}(\operatorname{ring\text{-}additive\text{-}ag}(a))) \end{array}}}
\vnbprop{ter-tout-type}{\vnbstmt{\begin{array}{@{}l@{}} \forall a, m, n, k, l, p, q, r, w, b. \\ \quad \operatorname{is\text{-}ring}(a) \wedge \\ \quad p \in \operatorname{mat}(m, n, \operatorname{carr}(a)) \wedge \\ \quad q \in \operatorname{mat}(n, k, \operatorname{carr}(a)) \wedge \\ \quad r \in \operatorname{mat}(k, l, \operatorname{carr}(a)) \wedge \\ \quad w \in \operatorname{interval}(1, m) \wedge \\ \quad b \in \operatorname{interval}(1, l) \Rightarrow \\ \quad \quad (\lambda j \in \operatorname{interval}(1, n).\; \\ \quad \quad \quad \operatorname{finsum}(\operatorname{ring\text{-}additive\text{-}ag}(a), \\ \quad \quad \quad \quad (\lambda c \in \operatorname{interval}(1, k).\; \\ \quad \quad \quad \quad \quad (\lambda z \in (\operatorname{interval}(1, k) \times \operatorname{interval}(1, n)).\; \operatorname{mul}(a)(\operatorname{mul}(a)(\operatorname{entry}(p, w, \operatorname{nth}(2, z)), \operatorname{entry}(q, \operatorname{nth}(2, z), \operatorname{nth}(1, z))), \operatorname{entry}(r, \operatorname{nth}(1, z), b)))({[}c, j])), \\ \quad \quad \quad \quad \operatorname{interval}(1, k))) \\ \quad \quad \in\; (\operatorname{interval}(1, n) \to \operatorname{carr}(\operatorname{ring\text{-}additive\text{-}ag}(a))) \end{array}}}
\vnbprop{ter-gj-type}{\vnbstmt{\begin{array}{@{}l@{}} \forall a, m, n, k, l, p, q, r, w, b, x. \\ \quad \operatorname{is\text{-}ring}(a) \wedge \\ \quad p \in \operatorname{mat}(m, n, \operatorname{carr}(a)) \wedge \\ \quad q \in \operatorname{mat}(n, k, \operatorname{carr}(a)) \wedge \\ \quad r \in \operatorname{mat}(k, l, \operatorname{carr}(a)) \wedge \\ \quad w \in \operatorname{interval}(1, m) \wedge \\ \quad b \in \operatorname{interval}(1, l) \wedge \\ \quad x \in \operatorname{interval}(1, n) \Rightarrow \\ \quad \quad (\lambda j \in \operatorname{interval}(1, k).\; \operatorname{mul}(a)(\operatorname{entry}(q, x, j), \operatorname{entry}(r, j, b))) \\ \quad \quad \in\; (\operatorname{interval}(1, k) \to \operatorname{carr}(a)) \end{array}}}
\vnbprop{ter-dist-type}{\vnbstmt{\begin{array}{@{}l@{}} \forall a, m, n, k, l, p, q, r, w, b, x. \\ \quad \operatorname{is\text{-}ring}(a) \wedge \\ \quad p \in \operatorname{mat}(m, n, \operatorname{carr}(a)) \wedge \\ \quad q \in \operatorname{mat}(n, k, \operatorname{carr}(a)) \wedge \\ \quad r \in \operatorname{mat}(k, l, \operatorname{carr}(a)) \wedge \\ \quad w \in \operatorname{interval}(1, m) \wedge \\ \quad b \in \operatorname{interval}(1, l) \wedge \\ \quad x \in \operatorname{interval}(1, n) \Rightarrow \\ \quad \quad (\lambda z \in \operatorname{interval}(1, k).\; \\ \quad \quad \quad \operatorname{mul}(a)(\operatorname{entry}(p, w, x), \\ \quad \quad \quad \quad (\lambda j \in \operatorname{interval}(1, k).\; \operatorname{mul}(a)(\operatorname{entry}(q, x, j), \operatorname{entry}(r, j, b)))(z))) \\ \quad \quad \in\; (\operatorname{interval}(1, k) \to \operatorname{carr}(\operatorname{ring\text{-}additive\text{-}ag}(a))) \end{array}}}
\vnbprop{ter-red-type}{\vnbstmt{\begin{array}{@{}l@{}} \forall a, m, n, k, l, p, q, r, w, b, x. \\ \quad \operatorname{is\text{-}ring}(a) \wedge \\ \quad p \in \operatorname{mat}(m, n, \operatorname{carr}(a)) \wedge \\ \quad q \in \operatorname{mat}(n, k, \operatorname{carr}(a)) \wedge \\ \quad r \in \operatorname{mat}(k, l, \operatorname{carr}(a)) \wedge \\ \quad w \in \operatorname{interval}(1, m) \wedge \\ \quad b \in \operatorname{interval}(1, l) \wedge \\ \quad x \in \operatorname{interval}(1, n) \Rightarrow \\ \quad \quad (\lambda c \in \operatorname{interval}(1, k).\; \operatorname{mul}(a)(\operatorname{mul}(a)(\operatorname{entry}(p, w, x), \operatorname{entry}(q, x, c)), \operatorname{entry}(r, c, b))) \\ \quad \quad \in\; (\operatorname{interval}(1, k) \to \operatorname{carr}(\operatorname{ring\text{-}additive\text{-}ag}(a))) \end{array}}}
\vnbprop{mal-outf-type}{\vnbstmt{\begin{array}{@{}l@{}} \forall d, m, n, k, l, p, q, u, w, a. \\ \quad \operatorname{is\text{-}module}(d) \wedge \\ \quad p \in \operatorname{mat}(m, n, \operatorname{carr}(\operatorname{scal}(d))) \wedge \\ \quad q \in \operatorname{mat}(n, k, \operatorname{carr}(\operatorname{scal}(d))) \wedge \\ \quad u \in \operatorname{mat}(k, l, \operatorname{vec}(d)) \wedge \\ \quad w \in \operatorname{interval}(1, m) \wedge \\ \quad a \in \operatorname{interval}(1, l) \wedge \\ \quad 1 \leq n \Rightarrow \\ \quad \quad (\lambda j \in \operatorname{interval}(1, k).\; \operatorname{act}(d)(\operatorname{entry}(\operatorname{matmul}(\operatorname{scal}(d), p, q), w, j), \operatorname{entry}(u, j, a))) \\ \quad \quad \in\; (\operatorname{interval}(1, k) \to \operatorname{carr}(\operatorname{module\text{-}vector\text{-}ag}(d))) \end{array}}}
\vnbprop{mal-tout-type}{\vnbstmt{\begin{array}{@{}l@{}} \forall d, m, n, k, l, p, q, u, w, a. \\ \quad \operatorname{is\text{-}module}(d) \wedge \\ \quad p \in \operatorname{mat}(m, n, \operatorname{carr}(\operatorname{scal}(d))) \wedge \\ \quad q \in \operatorname{mat}(n, k, \operatorname{carr}(\operatorname{scal}(d))) \wedge \\ \quad u \in \operatorname{mat}(k, l, \operatorname{vec}(d)) \wedge \\ \quad w \in \operatorname{interval}(1, m) \wedge \\ \quad a \in \operatorname{interval}(1, l) \Rightarrow \\ \quad \quad (\lambda c \in \operatorname{interval}(1, k).\; \\ \quad \quad \quad \operatorname{finsum}(\operatorname{module\text{-}vector\text{-}ag}(d), \\ \quad \quad \quad \quad (\lambda j \in \operatorname{interval}(1, n).\; \\ \quad \quad \quad \quad \quad (\lambda z \in (\operatorname{interval}(1, k) \times \operatorname{interval}(1, n)).\; \operatorname{act}(d)(\operatorname{mul}(\operatorname{scal}(d))(\operatorname{entry}(p, w, \operatorname{nth}(2, z)), \operatorname{entry}(q, \operatorname{nth}(2, z), \operatorname{nth}(1, z))), \operatorname{entry}(u, \operatorname{nth}(1, z), a)))({[}c, j])), \\ \quad \quad \quad \quad \operatorname{interval}(1, n))) \\ \quad \quad \in\; (\operatorname{interval}(1, k) \to \operatorname{carr}(\operatorname{module\text{-}vector\text{-}ag}(d))) \end{array}}}
\vnbprop{mal-inf-type}{\vnbstmt{\begin{array}{@{}l@{}} \forall d, m, n, k, l, p, q, u, w, a, x. \\ \quad \operatorname{is\text{-}module}(d) \wedge \\ \quad p \in \operatorname{mat}(m, n, \operatorname{carr}(\operatorname{scal}(d))) \wedge \\ \quad q \in \operatorname{mat}(n, k, \operatorname{carr}(\operatorname{scal}(d))) \wedge \\ \quad u \in \operatorname{mat}(k, l, \operatorname{vec}(d)) \wedge \\ \quad w \in \operatorname{interval}(1, m) \wedge \\ \quad a \in \operatorname{interval}(1, l) \wedge \\ \quad x \in \operatorname{interval}(1, k) \Rightarrow \\ \quad \quad (\lambda j \in \operatorname{interval}(1, n).\; \operatorname{mul}(\operatorname{scal}(d))(\operatorname{entry}(p, w, j), \operatorname{entry}(q, j, x))) \\ \quad \quad \in\; (\operatorname{interval}(1, n) \to \operatorname{carr}(\operatorname{scal}(d))) \end{array}}}
\vnbprop{mal-dist-type}{\vnbstmt{\begin{array}{@{}l@{}} \forall d, m, n, k, l, p, q, u, w, a, x. \\ \quad \operatorname{is\text{-}module}(d) \wedge \\ \quad p \in \operatorname{mat}(m, n, \operatorname{carr}(\operatorname{scal}(d))) \wedge \\ \quad q \in \operatorname{mat}(n, k, \operatorname{carr}(\operatorname{scal}(d))) \wedge \\ \quad u \in \operatorname{mat}(k, l, \operatorname{vec}(d)) \wedge \\ \quad w \in \operatorname{interval}(1, m) \wedge \\ \quad a \in \operatorname{interval}(1, l) \wedge \\ \quad x \in \operatorname{interval}(1, k) \Rightarrow \\ \quad \quad (\lambda z \in \operatorname{interval}(1, n).\; \\ \quad \quad \quad \operatorname{act}(d)((\lambda j \in \operatorname{interval}(1, n).\; \operatorname{mul}(\operatorname{scal}(d))(\operatorname{entry}(p, w, j), \operatorname{entry}(q, j, x)))(z), \\ \quad \quad \quad \quad \operatorname{entry}(u, x, a))) \\ \quad \quad \in\; (\operatorname{interval}(1, n) \to \operatorname{carr}(\operatorname{module\text{-}vector\text{-}ag}(d))) \end{array}}}
\vnbprop{mal-red-type}{\vnbstmt{\begin{array}{@{}l@{}} \forall d, m, n, k, l, p, q, u, w, a, x. \\ \quad \operatorname{is\text{-}module}(d) \wedge \\ \quad p \in \operatorname{mat}(m, n, \operatorname{carr}(\operatorname{scal}(d))) \wedge \\ \quad q \in \operatorname{mat}(n, k, \operatorname{carr}(\operatorname{scal}(d))) \wedge \\ \quad u \in \operatorname{mat}(k, l, \operatorname{vec}(d)) \wedge \\ \quad w \in \operatorname{interval}(1, m) \wedge \\ \quad a \in \operatorname{interval}(1, l) \wedge \\ \quad x \in \operatorname{interval}(1, k) \Rightarrow \\ \quad \quad (\lambda j \in \operatorname{interval}(1, n).\; \\ \quad \quad \quad \operatorname{act}(d)(\operatorname{mul}(\operatorname{scal}(d))(\operatorname{entry}(p, w, j), \operatorname{entry}(q, j, x)), \operatorname{entry}(u, x, a))) \\ \quad \quad \in\; (\operatorname{interval}(1, n) \to \operatorname{carr}(\operatorname{module\text{-}vector\text{-}ag}(d))) \end{array}}}
\vnbprop{mar-outf-type}{\vnbstmt{\begin{array}{@{}l@{}} \forall d, m, n, k, l, p, q, u, w, a. \\ \quad \operatorname{is\text{-}module}(d) \wedge \\ \quad p \in \operatorname{mat}(m, n, \operatorname{carr}(\operatorname{scal}(d))) \wedge \\ \quad q \in \operatorname{mat}(n, k, \operatorname{carr}(\operatorname{scal}(d))) \wedge \\ \quad u \in \operatorname{mat}(k, l, \operatorname{vec}(d)) \wedge \\ \quad w \in \operatorname{interval}(1, m) \wedge \\ \quad a \in \operatorname{interval}(1, l) \wedge \\ \quad 1 \leq k \Rightarrow \\ \quad \quad (\lambda j \in \operatorname{interval}(1, n).\; \operatorname{act}(d)(\operatorname{entry}(p, w, j), \operatorname{entry}(\operatorname{matact}(d, q, u), j, a))) \\ \quad \quad \in\; (\operatorname{interval}(1, n) \to \operatorname{carr}(\operatorname{module\text{-}vector\text{-}ag}(d))) \end{array}}}
\vnbprop{mar-tout-type}{\vnbstmt{\begin{array}{@{}l@{}} \forall d, m, n, k, l, p, q, u, w, a. \\ \quad \operatorname{is\text{-}module}(d) \wedge \\ \quad p \in \operatorname{mat}(m, n, \operatorname{carr}(\operatorname{scal}(d))) \wedge \\ \quad q \in \operatorname{mat}(n, k, \operatorname{carr}(\operatorname{scal}(d))) \wedge \\ \quad u \in \operatorname{mat}(k, l, \operatorname{vec}(d)) \wedge \\ \quad w \in \operatorname{interval}(1, m) \wedge \\ \quad a \in \operatorname{interval}(1, l) \Rightarrow \\ \quad \quad (\lambda j \in \operatorname{interval}(1, n).\; \\ \quad \quad \quad \operatorname{finsum}(\operatorname{module\text{-}vector\text{-}ag}(d), \\ \quad \quad \quad \quad (\lambda c \in \operatorname{interval}(1, k).\; \\ \quad \quad \quad \quad \quad (\lambda z \in (\operatorname{interval}(1, k) \times \operatorname{interval}(1, n)).\; \operatorname{act}(d)(\operatorname{mul}(\operatorname{scal}(d))(\operatorname{entry}(p, w, \operatorname{nth}(2, z)), \operatorname{entry}(q, \operatorname{nth}(2, z), \operatorname{nth}(1, z))), \operatorname{entry}(u, \operatorname{nth}(1, z), a)))({[}c, j])), \\ \quad \quad \quad \quad \operatorname{interval}(1, k))) \\ \quad \quad \in\; (\operatorname{interval}(1, n) \to \operatorname{carr}(\operatorname{module\text{-}vector\text{-}ag}(d))) \end{array}}}
\vnbprop{mar-gj-type}{\vnbstmt{\begin{array}{@{}l@{}} \forall d, m, n, k, l, p, q, u, w, a, x. \\ \quad \operatorname{is\text{-}module}(d) \wedge \\ \quad p \in \operatorname{mat}(m, n, \operatorname{carr}(\operatorname{scal}(d))) \wedge \\ \quad q \in \operatorname{mat}(n, k, \operatorname{carr}(\operatorname{scal}(d))) \wedge \\ \quad u \in \operatorname{mat}(k, l, \operatorname{vec}(d)) \wedge \\ \quad w \in \operatorname{interval}(1, m) \wedge \\ \quad a \in \operatorname{interval}(1, l) \wedge \\ \quad x \in \operatorname{interval}(1, n) \Rightarrow \\ \quad \quad (\lambda j \in \operatorname{interval}(1, k).\; \operatorname{act}(d)(\operatorname{entry}(q, x, j), \operatorname{entry}(u, j, a))) \\ \quad \quad \in\; (\operatorname{interval}(1, k) \to \operatorname{carr}(\operatorname{module\text{-}vector\text{-}ag}(d))) \end{array}}}
\vnbprop{mar-dist-type}{\vnbstmt{\begin{array}{@{}l@{}} \forall d, m, n, k, l, p, q, u, w, a, x. \\ \quad \operatorname{is\text{-}module}(d) \wedge \\ \quad p \in \operatorname{mat}(m, n, \operatorname{carr}(\operatorname{scal}(d))) \wedge \\ \quad q \in \operatorname{mat}(n, k, \operatorname{carr}(\operatorname{scal}(d))) \wedge \\ \quad u \in \operatorname{mat}(k, l, \operatorname{vec}(d)) \wedge \\ \quad w \in \operatorname{interval}(1, m) \wedge \\ \quad a \in \operatorname{interval}(1, l) \wedge \\ \quad x \in \operatorname{interval}(1, n) \Rightarrow \\ \quad \quad (\lambda z \in \operatorname{interval}(1, k).\; \\ \quad \quad \quad \operatorname{act}(d)(\operatorname{entry}(p, w, x), \\ \quad \quad \quad \quad (\lambda j \in \operatorname{interval}(1, k).\; \operatorname{act}(d)(\operatorname{entry}(q, x, j), \operatorname{entry}(u, j, a)))(z))) \\ \quad \quad \in\; (\operatorname{interval}(1, k) \to \operatorname{carr}(\operatorname{module\text{-}vector\text{-}ag}(d))) \end{array}}}
\vnbprop{mar-red-type}{\vnbstmt{\begin{array}{@{}l@{}} \forall d, m, n, k, l, p, q, u, w, a, x. \\ \quad \operatorname{is\text{-}module}(d) \wedge \\ \quad p \in \operatorname{mat}(m, n, \operatorname{carr}(\operatorname{scal}(d))) \wedge \\ \quad q \in \operatorname{mat}(n, k, \operatorname{carr}(\operatorname{scal}(d))) \wedge \\ \quad u \in \operatorname{mat}(k, l, \operatorname{vec}(d)) \wedge \\ \quad w \in \operatorname{interval}(1, m) \wedge \\ \quad a \in \operatorname{interval}(1, l) \wedge \\ \quad x \in \operatorname{interval}(1, n) \Rightarrow \\ \quad \quad (\lambda c \in \operatorname{interval}(1, k).\; \\ \quad \quad \quad \operatorname{act}(d)(\operatorname{mul}(\operatorname{scal}(d))(\operatorname{entry}(p, w, x), \operatorname{entry}(q, x, c)), \operatorname{entry}(u, c, a))) \\ \quad \quad \in\; (\operatorname{interval}(1, k) \to \operatorname{carr}(\operatorname{module\text{-}vector\text{-}ag}(d))) \end{array}}}
\vnbprop{matact-assoc-summand-type}{\vnbstmt{\begin{array}{@{}l@{}} \forall d, m, n, k, l, p, q, u, w, a. \\ \quad \operatorname{is\text{-}module}(d) \wedge \\ \quad p \in \operatorname{mat}(m, n, \operatorname{carr}(\operatorname{scal}(d))) \wedge \\ \quad q \in \operatorname{mat}(n, k, \operatorname{carr}(\operatorname{scal}(d))) \wedge \\ \quad u \in \operatorname{mat}(k, l, \operatorname{vec}(d)) \wedge \\ \quad w \in \operatorname{interval}(1, m) \wedge \\ \quad a \in \operatorname{interval}(1, l) \Rightarrow \\ \quad \quad (\lambda z \in (\operatorname{interval}(1, k) \times \operatorname{interval}(1, n)).\; \\ \quad \quad \quad \operatorname{act}(d)(\operatorname{mul}(\operatorname{scal}(d))(\operatorname{entry}(p, w, \operatorname{nth}(2, z)), \operatorname{entry}(q, \operatorname{nth}(2, z), \operatorname{nth}(1, z))), \\ \quad \quad \quad \quad \operatorname{entry}(u, \operatorname{nth}(1, z), a))) \\ \quad \quad \in\; ((\operatorname{interval}(1, k) \times \operatorname{interval}(1, n)) \to \operatorname{carr}(\operatorname{module\text{-}vector\text{-}ag}(d))) \end{array}}}
\vnbprop{mra-combined-summand-type}{\vnbstmt{\begin{array}{@{}l@{}} \forall d, n, a, b, u. \\ \quad \operatorname{is\text{-}module}(d) \wedge \\ \quad a \in \operatorname{mat}(1, n, \operatorname{carr}(\operatorname{scal}(d))) \wedge \\ \quad b \in \operatorname{mat}(1, n, \operatorname{carr}(\operatorname{scal}(d))) \wedge \\ \quad u \in \operatorname{mat}(n, 1, \operatorname{vec}(d)) \Rightarrow \\ \quad \quad (\lambda z \in \operatorname{interval}(1, n).\; \\ \quad \quad \quad \operatorname{opr}(\operatorname{module\text{-}vector\text{-}ag}(d))((\lambda j \in \operatorname{interval}(1, n).\; \operatorname{act}(d)(\operatorname{entry}(a, 1, j), \operatorname{entry}(u, j, 1)))(z), \\ \quad \quad \quad \quad (\lambda j \in \operatorname{interval}(1, n).\; \operatorname{act}(d)(\operatorname{entry}(b, 1, j), \operatorname{entry}(u, j, 1)))(z))) \\ \quad \quad \in\; (\operatorname{interval}(1, n) \to \operatorname{carr}(\operatorname{module\text{-}vector\text{-}ag}(d))) \end{array}}}
\vnbprop{mrs-scaled-summand-type}{\vnbstmt{\begin{array}{@{}l@{}} \forall d, n, r, c, u. \\ \quad (\operatorname{is\text{-}module}(d) \wedge ((r \in \operatorname{carr}(\operatorname{scal}(d))) \wedge ((c \in \operatorname{mat}(1, n, \operatorname{carr}(\operatorname{scal}(d)))) \wedge (u \in \operatorname{mat}(n, 1, \operatorname{vec}(d)))))) \Rightarrow  \\ \quad \quad (\lambda z \in \operatorname{interval}(1, n).\; \\ \quad \quad \quad \operatorname{act}(d)(r, (\lambda j \in \operatorname{interval}(1, n).\; \operatorname{act}(d)(\operatorname{entry}(c, 1, j), \operatorname{entry}(u, j, 1)))(z))) \\ \quad \quad \in\; (\operatorname{interval}(1, n) \to \operatorname{carr}(\operatorname{module\text{-}vector\text{-}ag}(d))) \end{array}}}

\vnbbookfile{theorem-library/rake-identmat}{theorem-library/rake-identmat}
\noindent{\small rake-identmat.scm -- the two IDENTMAT identity laws and the module twin of matmul-entry, PROVEN (batch K, 2026-09-17).}\par
\vnbprop{identmat-left-identity}{\vnbstmt{\begin{array}{@{}l@{}} \forall a, m, n.; \forall p \in \operatorname{mat}(m, n, \operatorname{carr}(a)). \\ \quad \operatorname{is\text{-}ring}(a) \Rightarrow (\operatorname{matmul}(a, \operatorname{identmat}(a, m), p) = p) \end{array}}}
\vnbprop{identmat-right-identity}{\vnbstmt{\begin{array}{@{}l@{}} \forall a, m, n.; \forall p \in \operatorname{mat}(m, n, \operatorname{carr}(a)). \\ \quad \operatorname{is\text{-}ring}(a) \Rightarrow (\operatorname{matmul}(a, p, \operatorname{identmat}(a, n)) = p) \end{array}}}
\vnbprop{matact-entry}{\vnbstmt{\begin{array}{@{}l@{}} \forall d, m, n, q, p, u.; \forall i \in \operatorname{interval}(1, m), c \in \operatorname{interval}(1, q). \\ \quad (\operatorname{is\text{-}module}(d) \wedge ((p \in \operatorname{mat}(m, n, \operatorname{carr}(\operatorname{scal}(d)))) \wedge ((u \in \operatorname{mat}(n, q, \operatorname{vec}(d))) \wedge (1 \leq n)))) \Rightarrow  \\ \quad \quad \operatorname{entry}(\operatorname{matact}(d, p, u), i, c) \\ \quad \quad =\; \operatorname{finsum}(\operatorname{module\text{-}vector\text{-}ag}(d), \\ \quad \quad \quad (\lambda j \in \operatorname{interval}(1, n).\; \operatorname{act}(d)(\operatorname{entry}(p, i, j), \operatorname{entry}(u, j, c))), \\ \quad \quad \quad \operatorname{interval}(1, n)) \end{array}}}

\vnbbookfile{theorem-library/rake-generates-coeff}{theorem-library/rake-generates-coeff}
\noindent{\small rake-generates-coeff.scm -- generates-coeff-matrix, proven.}\par
\vnbprop{generates-coeff-matrix}{\vnbstmt{\begin{array}{@{}l@{}} \forall d, n, m, u, v. \\ \quad \operatorname{is\text{-}module}(d) \wedge \\ \quad u \in \operatorname{mat}(n, 1, \operatorname{vec}(d)) \wedge \\ \quad v \in \operatorname{mat}(m, 1, \operatorname{vec}(d)) \wedge \\ \quad (n = 0) \Rightarrow (m = 0) \wedge \\ \quad \operatorname{generates}(d, n, u) \Rightarrow \\ \quad \quad \exists a \in \operatorname{mat}(m, n, \operatorname{carr}(\operatorname{scal}(d))).\; (v = \operatorname{matact}(d, a, u)) \end{array}}}

\vnbbookfile{theorem-library/rake-smith-leaves}{theorem-library/rake-smith-leaves}
\noindent{\small rake-smith-leaves.scm -- finsum-interval-shift, PROVEN.}\par
\vnbprop{finsum-interval-shift}{\vnbstmt{\begin{array}{@{}l@{}} \forall g.; \forall n \in \mathbb{N}, f \in (\operatorname{interval}(1, (n + 1)) \to \operatorname{carr}(g)). \\ \quad \operatorname{is\text{-}abelian\text{-}group}(g) \Rightarrow  \\ \quad \quad \operatorname{finsum}(g, f, \operatorname{interval}(1, (n + 1))) \\ \quad \quad =\; \operatorname{opr}(g)(f(1), \operatorname{finsum}(g, (\lambda z \in \operatorname{interval}(1, n).\; f((z + 1))), \operatorname{interval}(1, n))) \end{array}}}

\vnbbookfile{theorem-library/module-zero-act}{theorem-library/module-zero-act}
\noindent{\small module-zero-act.scm -- 0 . x = 0\_V : the ring zero annihilates every vector.}\par
\vnbprop{module-zero-act}{\vnbstmt{\begin{array}{@{}l@{}} \forall s.; \forall x \in \operatorname{vec}(s). \\ \quad \operatorname{is\text{-}module}(s) \Rightarrow (\operatorname{act}(s)(\operatorname{zero}(\operatorname{scal}(s)), x) = \operatorname{vzero}(s)) \end{array}}}

\vnbbookfile{theorem-library/module-laws}{theorem-library/module-laws}
\noindent{\small theorem-library/module-laws.scm -- the generic laws of a MODULE: the vector group (associativity, commutativity, zero, inverse, the differences x - y = x + (-y)) and the scalar action on negatives and differences.  Prefix mlaw-.  Library: linear-algebra. (MODLAWS, 2026-10-24, the user's order: one early home for the ...}\par
\vnbprop{hbs-vadd-assoc}{\vnbstmt{\begin{array}{@{}l@{}} \forall m.; \forall x, y, z \in \operatorname{vec}(m). \\ \quad \operatorname{is\text{-}module}(m) \Rightarrow  \\ \quad \quad \operatorname{vadd}(m)(\operatorname{vadd}(m)(x, y), z) = \operatorname{vadd}(m)(x, \operatorname{vadd}(m)(y, z)) \end{array}}}
\vnbprop{hbs-vadd-comm}{\vnbstmt{\begin{array}{@{}l@{}} \forall m.; \forall x, y \in \operatorname{vec}(m). \\ \quad \operatorname{is\text{-}module}(m) \Rightarrow (\operatorname{vadd}(m)(x, y) = \operatorname{vadd}(m)(y, x)) \end{array}}}
\vnbprop{hbs-vzero-lr}{\vnbstmt{\begin{array}{@{}l@{}} \forall m.; \forall x \in \operatorname{vec}(m). \\ \quad \operatorname{is\text{-}module}(m) \Rightarrow ((\operatorname{vadd}(m)(\operatorname{vzero}(m), x) = x) \wedge (\operatorname{vadd}(m)(x, \operatorname{vzero}(m)) = x)) \end{array}}}
\vnbprop{hbs-vzero-r}{\vnbstmt{\forall m.; \forall x \in \operatorname{vec}(m).\; (\operatorname{is\text{-}module}(m) \Rightarrow (\operatorname{vadd}(m)(x, \operatorname{vzero}(m)) = x))}}
\vnbprop{hbs-vzero-l}{\vnbstmt{\forall m.; \forall x \in \operatorname{vec}(m).\; (\operatorname{is\text{-}module}(m) \Rightarrow (\operatorname{vadd}(m)(\operatorname{vzero}(m), x) = x))}}
\vnbprop{hbs-interchange}{\vnbstmt{\begin{array}{@{}l@{}} \forall m.; \forall a, b, c, d \in \operatorname{vec}(m). \\ \quad \operatorname{is\text{-}module}(m) \Rightarrow  \\ \quad \quad \operatorname{vadd}(m)(\operatorname{vadd}(m)(a, b), \operatorname{vadd}(m)(c, d)) = \operatorname{vadd}(m)(\operatorname{vadd}(m)(a, c), \operatorname{vadd}(m)(b, d)) \end{array}}}
\vnbprop{hbs-sub-vec}{\vnbstmt{\forall m, s.; \forall x \in s.\; ((\operatorname{is\text{-}module}(m) \wedge \operatorname{is\text{-}submodule}(m, s)) \Rightarrow (x \in \operatorname{vec}(m)))}}
\vnbprop{frm-mvag-inv}{\vnbstmt{\forall d.\; (\operatorname{is\text{-}module}(d) \Rightarrow (\operatorname{inv}(\operatorname{module\text{-}vector\text{-}ag}(d)) = \operatorname{vneg}(d)))}}
\vnbprop{frm-v-right-inv}{\vnbstmt{\begin{array}{@{}l@{}} \forall m.; \forall x \in \operatorname{vec}(m). \\ \quad \operatorname{is\text{-}module}(m) \Rightarrow (\operatorname{vadd}(m)(x, \operatorname{vneg}(m)(x)) = \operatorname{vzero}(m)) \end{array}}}
\vnbprop{frm-v-neg-add}{\vnbstmt{\begin{array}{@{}l@{}} \forall m.; \forall x, y \in \operatorname{vec}(m). \\ \quad \operatorname{is\text{-}module}(m) \Rightarrow  \\ \quad \quad \operatorname{vneg}(m)(\operatorname{vadd}(m)(x, y)) = \operatorname{vadd}(m)(\operatorname{vneg}(m)(x), \operatorname{vneg}(m)(y)) \end{array}}}
\vnbprop{frm-v-neg-neg}{\vnbstmt{\forall m.; \forall x \in \operatorname{vec}(m).\; (\operatorname{is\text{-}module}(m) \Rightarrow (\operatorname{vneg}(m)(\operatorname{vneg}(m)(x)) = x))}}
\vnbprop{frm-v-sub-anti}{\vnbstmt{\begin{array}{@{}l@{}} \forall m.; \forall x, y \in \operatorname{vec}(m). \\ \quad \operatorname{is\text{-}module}(m) \Rightarrow  \\ \quad \quad \operatorname{vneg}(m)(\operatorname{vadd}(m)(x, \operatorname{vneg}(m)(y))) = \operatorname{vadd}(m)(y, \operatorname{vneg}(m)(x)) \end{array}}}
\vnbprop{frm-v-sub-chain}{\vnbstmt{\begin{array}{@{}l@{}} \forall m.; \forall x, y, z \in \operatorname{vec}(m). \\ \quad \operatorname{is\text{-}module}(m) \Rightarrow  \\ \quad \quad \operatorname{vadd}(m)(\operatorname{vadd}(m)(x, \operatorname{vneg}(m)(y)), \operatorname{vadd}(m)(y, \operatorname{vneg}(m)(z))) = \operatorname{vadd}(m)(x, \operatorname{vneg}(m)(z)) \end{array}}}
\vnbprop{frm-v-sub-zero-eq}{\vnbstmt{\begin{array}{@{}l@{}} \forall m.; \forall x, y \in \operatorname{vec}(m). \\ \quad (\operatorname{is\text{-}module}(m) \wedge (\operatorname{vadd}(m)(x, \operatorname{vneg}(m)(y)) = \operatorname{vzero}(m))) \Rightarrow (x = y) \end{array}}}
\vnbprop{frm-v-sub-shift}{\vnbstmt{\begin{array}{@{}l@{}} \forall m.; \forall x, y, r \in \operatorname{vec}(m). \\ \quad \operatorname{is\text{-}module}(m) \Rightarrow  \\ \quad \quad \operatorname{vadd}(m)(\operatorname{vadd}(m)(x, r), \operatorname{vneg}(m)(\operatorname{vadd}(m)(y, r))) = \operatorname{vadd}(m)(x, \operatorname{vneg}(m)(y)) \end{array}}}
\vnbprop{frm-v-sub-add}{\vnbstmt{\begin{array}{@{}l@{}} \forall m.; \forall x, y, u, v \in \operatorname{vec}(m). \\ \quad \operatorname{is\text{-}module}(m) \Rightarrow  \\ \quad \quad \operatorname{vadd}(m)(\operatorname{vadd}(m)(x, y), \operatorname{vneg}(m)(\operatorname{vadd}(m)(u, v))) \\ \quad \quad =\; \operatorname{vadd}(m)(\operatorname{vadd}(m)(x, \operatorname{vneg}(m)(u)), \operatorname{vadd}(m)(y, \operatorname{vneg}(m)(v))) \end{array}}}
\vnbprop{frm-v-sub-zero}{\vnbstmt{\begin{array}{@{}l@{}} \forall m.; \forall x \in \operatorname{vec}(m). \\ \quad \operatorname{is\text{-}module}(m) \Rightarrow (\operatorname{vadd}(m)(x, \operatorname{vneg}(m)(\operatorname{vzero}(m))) = x) \end{array}}}
\vnbprop{frm-module-vec-set}{\vnbstmt{\forall m.\; (\operatorname{is\text{-}module}(m) \Rightarrow (\operatorname{vec}(m) \in \mathbf{Set}))}}
\vnbprop{frm-act-neg-vec}{\vnbstmt{\begin{array}{@{}l@{}} \forall m, r, x. \\ \quad (\operatorname{is\text{-}module}(m) \wedge ((r \in \operatorname{carr}(\operatorname{scal}(m))) \wedge (x \in \operatorname{vec}(m)))) \Rightarrow  \\ \quad \quad \operatorname{act}(m)(r, \operatorname{vneg}(m)(x)) = \operatorname{vneg}(m)(\operatorname{act}(m)(r, x)) \end{array}}}
\vnbprop{frm-act-neg-scal}{\vnbstmt{\begin{array}{@{}l@{}} \forall m, r, x. \\ \quad (\operatorname{is\text{-}module}(m) \wedge ((r \in \operatorname{carr}(\operatorname{scal}(m))) \wedge (x \in \operatorname{vec}(m)))) \Rightarrow  \\ \quad \quad \operatorname{act}(m)(\operatorname{neg}(\operatorname{scal}(m))(r), x) = \operatorname{vneg}(m)(\operatorname{act}(m)(r, x)) \end{array}}}
\vnbprop{frm-act-sub-vec}{\vnbstmt{\begin{array}{@{}l@{}} \forall m, r, x, y. \\ \quad (\operatorname{is\text{-}module}(m) \wedge ((r \in \operatorname{carr}(\operatorname{scal}(m))) \wedge ((x \in \operatorname{vec}(m)) \wedge (y \in \operatorname{vec}(m))))) \Rightarrow  \\ \quad \quad \operatorname{act}(m)(r, \operatorname{vadd}(m)(x, \operatorname{vneg}(m)(y))) = \operatorname{vadd}(m)(\operatorname{act}(m)(r, x), \operatorname{vneg}(m)(\operatorname{act}(m)(r, y))) \end{array}}}
\vnbprop{frm-act-sub-scal}{\vnbstmt{\begin{array}{@{}l@{}} \forall m, r, s, x. \\ \quad (\operatorname{is\text{-}module}(m) \wedge ((r \in \operatorname{carr}(\operatorname{scal}(m))) \wedge ((s \in \operatorname{carr}(\operatorname{scal}(m))) \wedge (x \in \operatorname{vec}(m))))) \Rightarrow  \\ \quad \quad \operatorname{act}(m)(\operatorname{add}(\operatorname{scal}(m))(r, \operatorname{neg}(\operatorname{scal}(m))(s)), x) \\ \quad \quad =\; \operatorname{vadd}(m)(\operatorname{act}(m)(r, x), \operatorname{vneg}(m)(\operatorname{act}(m)(s, x))) \end{array}}}
\vnbprop{frm-act-diff}{\vnbstmt{\begin{array}{@{}l@{}} \forall m, r, s, x, y. \\ \quad (\operatorname{is\text{-}module}(m) \wedge ((r \in \operatorname{carr}(\operatorname{scal}(m))) \wedge ((s \in \operatorname{carr}(\operatorname{scal}(m))) \wedge ((x \in \operatorname{vec}(m)) \wedge (y \in \operatorname{vec}(m)))))) \Rightarrow  \\ \quad \quad \operatorname{vadd}(m)(\operatorname{act}(m)(r, x), \operatorname{vneg}(m)(\operatorname{act}(m)(s, y))) \\ \quad \quad =\; \operatorname{vadd}(m)(\operatorname{act}(m)(r, \operatorname{vadd}(m)(x, \operatorname{vneg}(m)(y))), \\ \quad \quad \quad \operatorname{act}(m)(\operatorname{add}(\operatorname{scal}(m))(r, \operatorname{neg}(\operatorname{scal}(m))(s)), y)) \end{array}}}

\vnbbookfile{theorem-library/triple-entry-proof}{theorem-library/triple-entry-proof}
\noindent{\small triple-entry-proof.scm -- the two matrix-product entry-expansion lemmas that matmul-assoc rests on, PROVEN (were warranted in matrix.scm) from the (B) finite-sum bricks (finsum-additive.scm): finsum-congruence, and general-ring distribution of a scalar into a finite sum (finsum-ring-distrib-left/right-gen).}\par
\vnbprop{triple-entry-left}{\vnbstmt{\begin{array}{@{}l@{}} \forall a, m, n, k, l, p, q, r, w, b. \\ \quad \operatorname{is\text{-}ring}(a) \wedge \\ \quad p \in \operatorname{mat}(m, n, \operatorname{carr}(a)) \wedge \\ \quad q \in \operatorname{mat}(n, k, \operatorname{carr}(a)) \wedge \\ \quad r \in \operatorname{mat}(k, l, \operatorname{carr}(a)) \wedge \\ \quad 1 \leq n \wedge \\ \quad 1 \leq k \wedge \\ \quad w \in \operatorname{interval}(1, m) \wedge \\ \quad b \in \operatorname{interval}(1, l) \Rightarrow \\ \quad \quad \operatorname{entry}(\operatorname{matmul}(a, \operatorname{matmul}(a, p, q), r), w, b) \\ \quad \quad =\; \operatorname{finsum}(\operatorname{ring\text{-}additive\text{-}ag}(a), \\ \quad \quad \quad (\lambda c \in \operatorname{interval}(1, k).\; \\ \quad \quad \quad \quad \operatorname{finsum}(\operatorname{ring\text{-}additive\text{-}ag}(a), \\ \quad \quad \quad \quad \quad (\lambda j \in \operatorname{interval}(1, n).\; \\ \quad \quad \quad \quad \quad \quad (\lambda z \in (\operatorname{interval}(1, k) \times \operatorname{interval}(1, n)).\; \operatorname{mul}(a)(\operatorname{mul}(a)(\operatorname{entry}(p, w, \operatorname{nth}(2, z)), \operatorname{entry}(q, \operatorname{nth}(2, z), \operatorname{nth}(1, z))), \operatorname{entry}(r, \operatorname{nth}(1, z), b)))({[}c, j])), \\ \quad \quad \quad \quad \quad \operatorname{interval}(1, n))), \\ \quad \quad \quad \operatorname{interval}(1, k)) \end{array}}}
\vnbprop{triple-entry-right}{\vnbstmt{\begin{array}{@{}l@{}} \forall a, m, n, k, l, p, q, r, w, b. \\ \quad \operatorname{is\text{-}ring}(a) \wedge \\ \quad p \in \operatorname{mat}(m, n, \operatorname{carr}(a)) \wedge \\ \quad q \in \operatorname{mat}(n, k, \operatorname{carr}(a)) \wedge \\ \quad r \in \operatorname{mat}(k, l, \operatorname{carr}(a)) \wedge \\ \quad 1 \leq n \wedge \\ \quad 1 \leq k \wedge \\ \quad w \in \operatorname{interval}(1, m) \wedge \\ \quad b \in \operatorname{interval}(1, l) \Rightarrow \\ \quad \quad \operatorname{entry}(\operatorname{matmul}(a, p, \operatorname{matmul}(a, q, r)), w, b) \\ \quad \quad =\; \operatorname{finsum}(\operatorname{ring\text{-}additive\text{-}ag}(a), \\ \quad \quad \quad (\lambda j \in \operatorname{interval}(1, n).\; \\ \quad \quad \quad \quad \operatorname{finsum}(\operatorname{ring\text{-}additive\text{-}ag}(a), \\ \quad \quad \quad \quad \quad (\lambda c \in \operatorname{interval}(1, k).\; \\ \quad \quad \quad \quad \quad \quad (\lambda z \in (\operatorname{interval}(1, k) \times \operatorname{interval}(1, n)).\; \operatorname{mul}(a)(\operatorname{mul}(a)(\operatorname{entry}(p, w, \operatorname{nth}(2, z)), \operatorname{entry}(q, \operatorname{nth}(2, z), \operatorname{nth}(1, z))), \operatorname{entry}(r, \operatorname{nth}(1, z), b)))({[}c, j])), \\ \quad \quad \quad \quad \quad \operatorname{interval}(1, k))), \\ \quad \quad \quad \operatorname{interval}(1, n)) \end{array}}}

\vnbbookfile{theorem-library/matmul-assoc-proof}{theorem-library/matmul-assoc-proof}
\noindent{\small matmul-assoc-proof.scm -- matrix multiplication is associative, via Fubini.}\par
\vnbprop{matmul-assoc}{\vnbstmt{\begin{array}{@{}l@{}} \forall a, m, n, k, l, p, q, r. \\ \quad \operatorname{is\text{-}ring}(a) \wedge \\ \quad p \in \operatorname{mat}(m, n, \operatorname{carr}(a)) \wedge \\ \quad q \in \operatorname{mat}(n, k, \operatorname{carr}(a)) \wedge \\ \quad r \in \operatorname{mat}(k, l, \operatorname{carr}(a)) \wedge \\ \quad (n = 0) \Rightarrow ((m = 0) \vee ((k = 0) \wedge (l = 0))) \wedge \\ \quad (k = 0) \Rightarrow ((l = 0) \vee ((n = 0) \wedge (m = 0))) \Rightarrow \\ \quad \quad \operatorname{matmul}(a, \operatorname{matmul}(a, p, q), r) = \operatorname{matmul}(a, p, \operatorname{matmul}(a, q, r)) \end{array}}}

\vnbbookfile{theorem-library/elem-entry-readoffs}{theorem-library/elem-entry-readoffs}
\noindent{\small elem-entry-readoffs.scm -- the elementary-matrix ENTRY read-offs, PROVEN.}\par
\vnbprop{matunit-entries-defined}{\vnbstmt{\begin{array}{@{}l@{}} \forall a, n, k, l.; \forall b, c \in \operatorname{interval}(1, n). \\ \quad (\operatorname{is\text{-}ring}(a) \wedge (n \in \mathbb{N})) \Rightarrow  \\ \quad \quad (\lambda {[}i, j] \in (\operatorname{interval}(1, n) \times \operatorname{interval}(1, n)).\; \left(\text{if }((i = k) \wedge (j = l))\text{ then }\operatorname{one}(a)\text{ else }\operatorname{zero}(a)\right))(b, c) \\ \quad \quad \in\; \mathbf{Set} \end{array}}}
\vnbprop{elem-f-entries-defined}{\vnbstmt{\begin{array}{@{}l@{}} \forall a, n, k, l.; \forall b, c \in \operatorname{interval}(1, n). \\ \quad (\operatorname{is\text{-}ring}(a) \wedge (n \in \mathbb{N})) \Rightarrow  \\ \quad \quad (\lambda {[}i, j] \in (\operatorname{interval}(1, n) \times \operatorname{interval}(1, n)).\; \left(\text{if }(((i = k) \wedge (j = l)) \vee (((i = l) \wedge (j = k)) \vee ((i = j) \wedge (\neg (i = k) \wedge \neg (i = l)))))\text{ then }\operatorname{one}(a)\text{ else }\operatorname{zero}(a)\right))(b, c) \\ \quad \quad \in\; \mathbf{Set} \end{array}}}
\vnbprop{elem-g-entries-defined}{\vnbstmt{\begin{array}{@{}l@{}} \forall a, n, r, k, l.; \forall b, c \in \operatorname{interval}(1, n). \\ \quad (\operatorname{is\text{-}ring}(a) \wedge ((n \in \mathbb{N}) \wedge (r \in \operatorname{carr}(a)))) \Rightarrow  \\ \quad \quad (\lambda {[}i, j] \in (\operatorname{interval}(1, n) \times \operatorname{interval}(1, n)).\; \left(\text{if }(i = j)\text{ then }\operatorname{one}(a)\text{ else }\left(\text{if }((i = k) \wedge (j = l))\text{ then }r\text{ else }\operatorname{zero}(a)\right)\right))(b, c) \\ \quad \quad \in\; \mathbf{Set} \end{array}}}
\vnbprop{elem-h-entries-defined}{\vnbstmt{\begin{array}{@{}l@{}} \forall a, n, r, k.; \forall b, c \in \operatorname{interval}(1, n). \\ \quad (\operatorname{is\text{-}ring}(a) \wedge ((n \in \mathbb{N}) \wedge (r \in \operatorname{carr}(a)))) \Rightarrow  \\ \quad \quad (\lambda {[}i, j] \in (\operatorname{interval}(1, n) \times \operatorname{interval}(1, n)).\; \left(\text{if }(i = j)\text{ then }\left(\text{if }(i = k)\text{ then }r\text{ else }\operatorname{one}(a)\right)\text{ else }\operatorname{zero}(a)\right))(b, c) \\ \quad \quad \in\; \mathbf{Set} \end{array}}}
\vnbprop{entry-of-matunit}{\vnbstmt{\begin{array}{@{}l@{}} \forall a, n, k, l, i, j. \\ \quad (\operatorname{is\text{-}ring}(a) \wedge ((n \in \mathbb{N}) \wedge ((i \in \operatorname{interval}(1, n)) \wedge (j \in \operatorname{interval}(1, n))))) \Rightarrow  \\ \quad \quad \operatorname{entry}(\operatorname{matunit}(a, n, k, l), i, j) = \left(\text{if }((i = k) \wedge (j = l))\text{ then }\operatorname{one}(a)\text{ else }\operatorname{zero}(a)\right) \end{array}}}
\vnbprop{matunit-entry-off-row}{\vnbstmt{\begin{array}{@{}l@{}} \forall a, n, k, l, j, c. \\ \quad (\operatorname{is\text{-}ring}(a) \wedge ((n \in \mathbb{N}) \wedge ((j \in \operatorname{interval}(1, n)) \wedge ((c \in \operatorname{interval}(1, n)) \wedge \neg (j = k))))) \Rightarrow  \\ \quad \quad \operatorname{entry}(\operatorname{matunit}(a, n, k, l), j, c) = \operatorname{zero}(a) \end{array}}}
\vnbprop{matunit-entry-k-row}{\vnbstmt{\begin{array}{@{}l@{}} \forall a, n, k, l, c. \\ \quad (\operatorname{is\text{-}ring}(a) \wedge ((n \in \mathbb{N}) \wedge ((k \in \operatorname{interval}(1, n)) \wedge (c \in \operatorname{interval}(1, n))))) \Rightarrow  \\ \quad \quad \operatorname{entry}(\operatorname{matunit}(a, n, k, l), k, c) = \left(\text{if }(c = l)\text{ then }\operatorname{one}(a)\text{ else }\operatorname{zero}(a)\right) \end{array}}}
\vnbprop{entry-of-elem-f}{\vnbstmt{\begin{array}{@{}l@{}} \forall a, n, k, l, i, j. \\ \quad (\operatorname{is\text{-}ring}(a) \wedge ((n \in \mathbb{N}) \wedge ((i \in \operatorname{interval}(1, n)) \wedge (j \in \operatorname{interval}(1, n))))) \Rightarrow  \\ \quad \quad \operatorname{entry}(\operatorname{elem\text{-}f}(a, n, k, l), i, j) \\ \quad \quad =\; \left(\text{if }(((i = k) \wedge (j = l)) \vee (((i = l) \wedge (j = k)) \vee ((i = j) \wedge (\neg (i = k) \wedge \neg (i = l)))))\text{ then }\operatorname{one}(a)\text{ else }\operatorname{zero}(a)\right) \end{array}}}
\vnbprop{entry-of-elem-g}{\vnbstmt{\begin{array}{@{}l@{}} \forall a, n, r, k, l, i, j. \\ \quad (\operatorname{is\text{-}ring}(a) \wedge ((n \in \mathbb{N}) \wedge ((r \in \operatorname{carr}(a)) \wedge ((i \in \operatorname{interval}(1, n)) \wedge (j \in \operatorname{interval}(1, n)))))) \Rightarrow  \\ \quad \quad \operatorname{entry}(\operatorname{elem\text{-}g}(a, n, r, k, l), i, j) = \left(\text{if }(i = j)\text{ then }\operatorname{one}(a)\text{ else }\left(\text{if }((i = k) \wedge (j = l))\text{ then }r\text{ else }\operatorname{zero}(a)\right)\right) \end{array}}}
\vnbprop{entry-of-elem-h}{\vnbstmt{\begin{array}{@{}l@{}} \forall a, n, r, k, i, j. \\ \quad (\operatorname{is\text{-}ring}(a) \wedge ((n \in \mathbb{N}) \wedge ((r \in \operatorname{carr}(a)) \wedge ((i \in \operatorname{interval}(1, n)) \wedge (j \in \operatorname{interval}(1, n)))))) \Rightarrow  \\ \quad \quad \operatorname{entry}(\operatorname{elem\text{-}h}(a, n, r, k), i, j) = \left(\text{if }(i = j)\text{ then }\left(\text{if }(i = k)\text{ then }r\text{ else }\operatorname{one}(a)\right)\text{ else }\operatorname{zero}(a)\right) \end{array}}}
\vnbprop{elem-g-entry-off-l-off-diag}{\vnbstmt{\begin{array}{@{}l@{}} \forall a, n, r, k, l, j, c. \\ \quad \operatorname{is\text{-}ring}(a) \wedge \\ \quad n \in \mathbb{N} \wedge \\ \quad r \in \operatorname{carr}(a) \wedge \\ \quad j \in \operatorname{interval}(1, n) \wedge \\ \quad c \in \operatorname{interval}(1, n) \wedge \\ \quad \neg (c = l) \wedge \\ \quad \neg (j = c) \Rightarrow \\ \quad \quad \operatorname{entry}(\operatorname{elem\text{-}g}(a, n, r, k, l), j, c) = \operatorname{zero}(a) \end{array}}}
\vnbprop{elem-g-entry-off-l-diag}{\vnbstmt{\begin{array}{@{}l@{}} \forall a, n, r, k, l, c. \\ \quad (\operatorname{is\text{-}ring}(a) \wedge ((n \in \mathbb{N}) \wedge ((r \in \operatorname{carr}(a)) \wedge ((c \in \operatorname{interval}(1, n)) \wedge \neg (c = l))))) \Rightarrow  \\ \quad \quad \operatorname{entry}(\operatorname{elem\text{-}g}(a, n, r, k, l), c, c) = \operatorname{one}(a) \end{array}}}
\vnbprop{elem-g-entry-l-vanish}{\vnbstmt{\begin{array}{@{}l@{}} \forall a, n, r, k, l, j, c. \\ \quad \operatorname{is\text{-}ring}(a) \wedge \\ \quad n \in \mathbb{N} \wedge \\ \quad r \in \operatorname{carr}(a) \wedge \\ \quad j \in \operatorname{interval}(1, n) \wedge \\ \quad c \in \operatorname{interval}(1, n) \wedge \\ \quad c = l \wedge \\ \quad \neg (j = l) \wedge \\ \quad \neg (j = k) \Rightarrow \\ \quad \quad \operatorname{entry}(\operatorname{elem\text{-}g}(a, n, r, k, l), j, c) = \operatorname{zero}(a) \end{array}}}
\vnbprop{elem-g-entry-l-at-l}{\vnbstmt{\begin{array}{@{}l@{}} \forall a, n, r, k, l, c. \\ \quad (\operatorname{is\text{-}ring}(a) \wedge ((n \in \mathbb{N}) \wedge ((r \in \operatorname{carr}(a)) \wedge ((c \in \operatorname{interval}(1, n)) \wedge (c = l))))) \Rightarrow  \\ \quad \quad \operatorname{entry}(\operatorname{elem\text{-}g}(a, n, r, k, l), l, c) = \operatorname{one}(a) \end{array}}}
\vnbprop{elem-f-ck-off}{\vnbstmt{\begin{array}{@{}l@{}} \forall a, n, k, l, j, c. \\ \quad \operatorname{is\text{-}ring}(a) \wedge \\ \quad n \in \mathbb{N} \wedge \\ \quad j \in \operatorname{interval}(1, n) \wedge \\ \quad c \in \operatorname{interval}(1, n) \wedge \\ \quad c = k \wedge \\ \quad \neg (k = l) \wedge \\ \quad \neg (j = l) \Rightarrow \\ \quad \quad \operatorname{entry}(\operatorname{elem\text{-}f}(a, n, k, l), j, c) = \operatorname{zero}(a) \end{array}}}
\vnbprop{elem-f-cl-off}{\vnbstmt{\begin{array}{@{}l@{}} \forall a, n, k, l, j, c. \\ \quad \operatorname{is\text{-}ring}(a) \wedge \\ \quad n \in \mathbb{N} \wedge \\ \quad j \in \operatorname{interval}(1, n) \wedge \\ \quad c \in \operatorname{interval}(1, n) \wedge \\ \quad c = l \wedge \\ \quad \neg (k = l) \wedge \\ \quad \neg (j = k) \Rightarrow \\ \quad \quad \operatorname{entry}(\operatorname{elem\text{-}f}(a, n, k, l), j, c) = \operatorname{zero}(a) \end{array}}}
\vnbprop{elem-f-co-off}{\vnbstmt{\begin{array}{@{}l@{}} \forall a, n, k, l, j, c. \\ \quad \operatorname{is\text{-}ring}(a) \wedge \\ \quad n \in \mathbb{N} \wedge \\ \quad j \in \operatorname{interval}(1, n) \wedge \\ \quad c \in \operatorname{interval}(1, n) \wedge \\ \quad \neg (c = k) \wedge \\ \quad \neg (c = l) \wedge \\ \quad \neg (j = c) \Rightarrow \\ \quad \quad \operatorname{entry}(\operatorname{elem\text{-}f}(a, n, k, l), j, c) = \operatorname{zero}(a) \end{array}}}
\vnbprop{elem-f-co-at}{\vnbstmt{\begin{array}{@{}l@{}} \forall a, n, k, l, c. \\ \quad (\operatorname{is\text{-}ring}(a) \wedge ((n \in \mathbb{N}) \wedge ((c \in \operatorname{interval}(1, n)) \wedge (\neg (c = k) \wedge \neg (c = l))))) \Rightarrow  \\ \quad \quad \operatorname{entry}(\operatorname{elem\text{-}f}(a, n, k, l), c, c) = \operatorname{one}(a) \end{array}}}
\vnbprop{elem-h-entry-off-diag}{\vnbstmt{\begin{array}{@{}l@{}} \forall a, n, r, k, j, c. \\ \quad \operatorname{is\text{-}ring}(a) \wedge \\ \quad n \in \mathbb{N} \wedge \\ \quad r \in \operatorname{carr}(a) \wedge \\ \quad j \in \operatorname{interval}(1, n) \wedge \\ \quad c \in \operatorname{interval}(1, n) \wedge \\ \quad \neg (j = c) \Rightarrow \\ \quad \quad \operatorname{entry}(\operatorname{elem\text{-}h}(a, n, r, k), j, c) = \operatorname{zero}(a) \end{array}}}
\vnbprop{elem-h-entry-diag}{\vnbstmt{\begin{array}{@{}l@{}} \forall a, n, r, k, c. \\ \quad (\operatorname{is\text{-}ring}(a) \wedge ((n \in \mathbb{N}) \wedge ((r \in \operatorname{carr}(a)) \wedge (c \in \operatorname{interval}(1, n))))) \Rightarrow  \\ \quad \quad \operatorname{entry}(\operatorname{elem\text{-}h}(a, n, r, k), c, c) = \left(\text{if }(c = k)\text{ then }r\text{ else }\operatorname{one}(a)\right) \end{array}}}
\vnbprop{elem-f-rk-off}{\vnbstmt{\begin{array}{@{}l@{}} \forall a, n, k, l, i, j. \\ \quad \operatorname{is\text{-}ring}(a) \wedge \\ \quad n \in \mathbb{N} \wedge \\ \quad i \in \operatorname{interval}(1, n) \wedge \\ \quad j \in \operatorname{interval}(1, n) \wedge \\ \quad i = k \wedge \\ \quad \neg (k = l) \wedge \\ \quad \neg (j = l) \Rightarrow \\ \quad \quad \operatorname{entry}(\operatorname{elem\text{-}f}(a, n, k, l), i, j) = \operatorname{zero}(a) \end{array}}}
\vnbprop{elem-f-rl-off}{\vnbstmt{\begin{array}{@{}l@{}} \forall a, n, k, l, i, j. \\ \quad \operatorname{is\text{-}ring}(a) \wedge \\ \quad n \in \mathbb{N} \wedge \\ \quad i \in \operatorname{interval}(1, n) \wedge \\ \quad j \in \operatorname{interval}(1, n) \wedge \\ \quad i = l \wedge \\ \quad \neg (k = l) \wedge \\ \quad \neg (j = k) \Rightarrow \\ \quad \quad \operatorname{entry}(\operatorname{elem\text{-}f}(a, n, k, l), i, j) = \operatorname{zero}(a) \end{array}}}
\vnbprop{elem-f-ro-off}{\vnbstmt{\begin{array}{@{}l@{}} \forall a, n, k, l, i, j. \\ \quad \operatorname{is\text{-}ring}(a) \wedge \\ \quad n \in \mathbb{N} \wedge \\ \quad i \in \operatorname{interval}(1, n) \wedge \\ \quad j \in \operatorname{interval}(1, n) \wedge \\ \quad \neg (i = k) \wedge \\ \quad \neg (i = l) \wedge \\ \quad \neg (j = i) \Rightarrow \\ \quad \quad \operatorname{entry}(\operatorname{elem\text{-}f}(a, n, k, l), i, j) = \operatorname{zero}(a) \end{array}}}
\vnbprop{elem-g-rk-vanish}{\vnbstmt{\begin{array}{@{}l@{}} \forall a, n, r, k, l, i, j. \\ \quad \operatorname{is\text{-}ring}(a) \wedge \\ \quad n \in \mathbb{N} \wedge \\ \quad r \in \operatorname{carr}(a) \wedge \\ \quad i \in \operatorname{interval}(1, n) \wedge \\ \quad j \in \operatorname{interval}(1, n) \wedge \\ \quad i = k \wedge \\ \quad \neg (j = k) \wedge \\ \quad \neg (j = l) \Rightarrow \\ \quad \quad \operatorname{entry}(\operatorname{elem\text{-}g}(a, n, r, k, l), i, j) = \operatorname{zero}(a) \end{array}}}
\vnbprop{elem-g-ro-off}{\vnbstmt{\begin{array}{@{}l@{}} \forall a, n, r, k, l, i, j. \\ \quad \operatorname{is\text{-}ring}(a) \wedge \\ \quad n \in \mathbb{N} \wedge \\ \quad r \in \operatorname{carr}(a) \wedge \\ \quad i \in \operatorname{interval}(1, n) \wedge \\ \quad j \in \operatorname{interval}(1, n) \wedge \\ \quad \neg (i = k) \wedge \\ \quad \neg (j = i) \Rightarrow \\ \quad \quad \operatorname{entry}(\operatorname{elem\text{-}g}(a, n, r, k, l), i, j) = \operatorname{zero}(a) \end{array}}}
\vnbprop{elem-h-ro-off}{\vnbstmt{\begin{array}{@{}l@{}} \forall a, n, r, k, i, j. \\ \quad \operatorname{is\text{-}ring}(a) \wedge \\ \quad n \in \mathbb{N} \wedge \\ \quad r \in \operatorname{carr}(a) \wedge \\ \quad i \in \operatorname{interval}(1, n) \wedge \\ \quad j \in \operatorname{interval}(1, n) \wedge \\ \quad \neg (j = i) \Rightarrow \\ \quad \quad \operatorname{entry}(\operatorname{elem\text{-}h}(a, n, r, k), i, j) = \operatorname{zero}(a) \end{array}}}
\vnbprop{elem-f-ck-at}{\vnbstmt{\begin{array}{@{}l@{}} \forall a, n, k, l, c. \\ \quad (\operatorname{is\text{-}ring}(a) \wedge ((n \in \mathbb{N}) \wedge ((l \in \operatorname{interval}(1, n)) \wedge ((c \in \operatorname{interval}(1, n)) \wedge (c = k))))) \Rightarrow  \\ \quad \quad \operatorname{entry}(\operatorname{elem\text{-}f}(a, n, k, l), l, c) = \operatorname{one}(a) \end{array}}}
\vnbprop{elem-f-cl-at}{\vnbstmt{\begin{array}{@{}l@{}} \forall a, n, k, l, c. \\ \quad (\operatorname{is\text{-}ring}(a) \wedge ((n \in \mathbb{N}) \wedge ((k \in \operatorname{interval}(1, n)) \wedge ((c \in \operatorname{interval}(1, n)) \wedge (c = l))))) \Rightarrow  \\ \quad \quad \operatorname{entry}(\operatorname{elem\text{-}f}(a, n, k, l), k, c) = \operatorname{one}(a) \end{array}}}
\vnbprop{elem-f-col-k}{\vnbstmt{\begin{array}{@{}l@{}} \forall a, n, k, l, i. \\ \quad (\operatorname{is\text{-}ring}(a) \wedge ((n \in \mathbb{N}) \wedge ((i \in \operatorname{interval}(1, n)) \wedge ((k \in \operatorname{interval}(1, n)) \wedge \neg (k = l))))) \Rightarrow  \\ \quad \quad \operatorname{entry}(\operatorname{elem\text{-}f}(a, n, k, l), i, k) = \left(\text{if }(i = l)\text{ then }\operatorname{one}(a)\text{ else }\operatorname{zero}(a)\right) \end{array}}}
\vnbprop{elem-f-col-l}{\vnbstmt{\begin{array}{@{}l@{}} \forall a, n, k, l, i. \\ \quad (\operatorname{is\text{-}ring}(a) \wedge ((n \in \mathbb{N}) \wedge ((i \in \operatorname{interval}(1, n)) \wedge ((l \in \operatorname{interval}(1, n)) \wedge \neg (k = l))))) \Rightarrow  \\ \quad \quad \operatorname{entry}(\operatorname{elem\text{-}f}(a, n, k, l), i, l) = \left(\text{if }(i = k)\text{ then }\operatorname{one}(a)\text{ else }\operatorname{zero}(a)\right) \end{array}}}
\vnbprop{elem-f-col-other}{\vnbstmt{\begin{array}{@{}l@{}} \forall a, n, k, l, i, c. \\ \quad (\operatorname{is\text{-}ring}(a) \wedge ((n \in \mathbb{N}) \wedge ((i \in \operatorname{interval}(1, n)) \wedge ((c \in \operatorname{interval}(1, n)) \wedge (\neg (c = k) \wedge \neg (c = l)))))) \Rightarrow  \\ \quad \quad \operatorname{entry}(\operatorname{elem\text{-}f}(a, n, k, l), i, c) = \left(\text{if }(i = c)\text{ then }\operatorname{one}(a)\text{ else }\operatorname{zero}(a)\right) \end{array}}}
\vnbprop{elem-f-rk-at}{\vnbstmt{\begin{array}{@{}l@{}} \forall a, n, k, l, i. \\ \quad (\operatorname{is\text{-}ring}(a) \wedge ((n \in \mathbb{N}) \wedge ((i \in \operatorname{interval}(1, n)) \wedge ((l \in \operatorname{interval}(1, n)) \wedge ((i = k) \wedge \neg (k = l)))))) \Rightarrow  \\ \quad \quad \operatorname{entry}(\operatorname{elem\text{-}f}(a, n, k, l), i, l) = \operatorname{one}(a) \end{array}}}
\vnbprop{elem-f-rl-at}{\vnbstmt{\begin{array}{@{}l@{}} \forall a, n, k, l, i. \\ \quad (\operatorname{is\text{-}ring}(a) \wedge ((n \in \mathbb{N}) \wedge ((i \in \operatorname{interval}(1, n)) \wedge ((k \in \operatorname{interval}(1, n)) \wedge ((i = l) \wedge \neg (k = l)))))) \Rightarrow  \\ \quad \quad \operatorname{entry}(\operatorname{elem\text{-}f}(a, n, k, l), i, k) = \operatorname{one}(a) \end{array}}}
\vnbprop{elem-g-entry-l-at-k}{\vnbstmt{\begin{array}{@{}l@{}} \forall a, n, r, k, l, c. \\ \quad \operatorname{is\text{-}ring}(a) \wedge \\ \quad n \in \mathbb{N} \wedge \\ \quad r \in \operatorname{carr}(a) \wedge \\ \quad k \in \operatorname{interval}(1, n) \wedge \\ \quad c \in \operatorname{interval}(1, n) \wedge \\ \quad c = l \wedge \\ \quad \neg (k = l) \Rightarrow \\ \quad \quad \operatorname{entry}(\operatorname{elem\text{-}g}(a, n, r, k, l), k, c) = r \end{array}}}
\vnbprop{elem-g-col-k}{\vnbstmt{\begin{array}{@{}l@{}} \forall a, n, r, k, l, i. \\ \quad \operatorname{is\text{-}ring}(a) \wedge \\ \quad n \in \mathbb{N} \wedge \\ \quad r \in \operatorname{carr}(a) \wedge \\ \quad i \in \operatorname{interval}(1, n) \wedge \\ \quad k \in \operatorname{interval}(1, n) \wedge \\ \quad \neg (k = l) \Rightarrow \\ \quad \quad \operatorname{entry}(\operatorname{elem\text{-}g}(a, n, r, k, l), i, k) = \left(\text{if }(i = k)\text{ then }\operatorname{one}(a)\text{ else }\operatorname{zero}(a)\right) \end{array}}}
\vnbprop{elem-g-col-l}{\vnbstmt{\begin{array}{@{}l@{}} \forall a, n, r, k, l, i. \\ \quad (\operatorname{is\text{-}ring}(a) \wedge ((n \in \mathbb{N}) \wedge ((r \in \operatorname{carr}(a)) \wedge ((i \in \operatorname{interval}(1, n)) \wedge (l \in \operatorname{interval}(1, n)))))) \Rightarrow  \\ \quad \quad \operatorname{entry}(\operatorname{elem\text{-}g}(a, n, r, k, l), i, l) = \left(\text{if }(i = l)\text{ then }\operatorname{one}(a)\text{ else }\left(\text{if }(i = k)\text{ then }r\text{ else }\operatorname{zero}(a)\right)\right) \end{array}}}
\vnbprop{elem-g-col-other}{\vnbstmt{\begin{array}{@{}l@{}} \forall a, n, r, k, l, i, c. \\ \quad \operatorname{is\text{-}ring}(a) \wedge \\ \quad n \in \mathbb{N} \wedge \\ \quad r \in \operatorname{carr}(a) \wedge \\ \quad i \in \operatorname{interval}(1, n) \wedge \\ \quad c \in \operatorname{interval}(1, n) \wedge \\ \quad \neg (c = l) \Rightarrow \\ \quad \quad \operatorname{entry}(\operatorname{elem\text{-}g}(a, n, r, k, l), i, c) = \left(\text{if }(i = c)\text{ then }\operatorname{one}(a)\text{ else }\operatorname{zero}(a)\right) \end{array}}}
\vnbprop{elem-g-rk-at-k}{\vnbstmt{\begin{array}{@{}l@{}} \forall a, n, r, k, l, i. \\ \quad (\operatorname{is\text{-}ring}(a) \wedge ((n \in \mathbb{N}) \wedge ((r \in \operatorname{carr}(a)) \wedge ((i \in \operatorname{interval}(1, n)) \wedge ((k \in \operatorname{interval}(1, n)) \wedge (i = k)))))) \Rightarrow  \\ \quad \quad \operatorname{entry}(\operatorname{elem\text{-}g}(a, n, r, k, l), i, k) = \operatorname{one}(a) \end{array}}}
\vnbprop{elem-g-rk-at-l}{\vnbstmt{\begin{array}{@{}l@{}} \forall a, n, r, k, l, i. \\ \quad \operatorname{is\text{-}ring}(a) \wedge \\ \quad n \in \mathbb{N} \wedge \\ \quad r \in \operatorname{carr}(a) \wedge \\ \quad i \in \operatorname{interval}(1, n) \wedge \\ \quad l \in \operatorname{interval}(1, n) \wedge \\ \quad i = k \wedge \\ \quad \neg (k = l) \Rightarrow \\ \quad \quad \operatorname{entry}(\operatorname{elem\text{-}g}(a, n, r, k, l), i, l) = r \end{array}}}
\vnbprop{elem-g-ro-at}{\vnbstmt{\begin{array}{@{}l@{}} \forall a, n, r, k, l, i. \\ \quad (\operatorname{is\text{-}ring}(a) \wedge ((n \in \mathbb{N}) \wedge ((r \in \operatorname{carr}(a)) \wedge ((i \in \operatorname{interval}(1, n)) \wedge \neg (i = k))))) \Rightarrow  \\ \quad \quad \operatorname{entry}(\operatorname{elem\text{-}g}(a, n, r, k, l), i, i) = \operatorname{one}(a) \end{array}}}
\vnbprop{identmat-entries-defined}{\vnbstmt{\begin{array}{@{}l@{}} \forall a, n.; \forall b, c \in \operatorname{interval}(1, n). \\ \quad (\operatorname{is\text{-}ring}(a) \wedge (n \in \mathbb{N})) \Rightarrow  \\ \quad \quad (\lambda {[}i, j] \in (\operatorname{interval}(1, n) \times \operatorname{interval}(1, n)).\; \left(\text{if }(i = j)\text{ then }\operatorname{one}(a)\text{ else }\operatorname{zero}(a)\right))(b, c) \\ \quad \quad \in\; \mathbf{Set} \end{array}}}
\vnbprop{unitrow-entries-defined}{\vnbstmt{\begin{array}{@{}l@{}} \forall a, n, i.; \forall b \in \operatorname{interval}(1, 1), j \in \operatorname{interval}(1, n). \\ \quad (\operatorname{is\text{-}ring}(a) \wedge (n \in \mathbb{N})) \Rightarrow  \\ \quad \quad (\lambda {[}\mathit{rw}, \mathit{cl}] \in (\operatorname{interval}(1, 1) \times \operatorname{interval}(1, n)).\; \left(\text{if }(\mathit{cl} = i)\text{ then }\operatorname{one}(a)\text{ else }\operatorname{zero}(a)\right))(b, j) \\ \quad \quad \in\; \mathbf{Set} \end{array}}}
\vnbprop{zeromat-entries-defined}{\vnbstmt{\begin{array}{@{}l@{}} \forall a, m, n.; \forall b \in \operatorname{interval}(1, m), c \in \operatorname{interval}(1, n). \\ \quad (\operatorname{is\text{-}ring}(a) \wedge ((m \in \mathbb{N}) \wedge (n \in \mathbb{N}))) \Rightarrow  \\ \quad \quad (\lambda {[}i, j] \in (\operatorname{interval}(1, m) \times \operatorname{interval}(1, n)).\; \operatorname{zero}(a))(b, c) \in \mathbf{Set} \end{array}}}
\vnbprop{entry-of-identmat}{\vnbstmt{\begin{array}{@{}l@{}} \forall a, n, i, j. \\ \quad (\operatorname{is\text{-}ring}(a) \wedge ((n \in \mathbb{N}) \wedge ((i \in \operatorname{interval}(1, n)) \wedge (j \in \operatorname{interval}(1, n))))) \Rightarrow  \\ \quad \quad \operatorname{entry}(\operatorname{identmat}(a, n), i, j) = \left(\text{if }(i = j)\text{ then }\operatorname{one}(a)\text{ else }\operatorname{zero}(a)\right) \end{array}}}
\vnbprop{identmat-entry-diag}{\vnbstmt{\begin{array}{@{}l@{}} \forall a, n, i. \\ \quad (\operatorname{is\text{-}ring}(a) \wedge ((n \in \mathbb{N}) \wedge (i \in \operatorname{interval}(1, n)))) \Rightarrow  \\ \quad \quad \operatorname{entry}(\operatorname{identmat}(a, n), i, i) = \operatorname{one}(a) \end{array}}}
\vnbprop{identmat-entry-off}{\vnbstmt{\begin{array}{@{}l@{}} \forall a, n, i, j. \\ \quad (\operatorname{is\text{-}ring}(a) \wedge ((n \in \mathbb{N}) \wedge ((i \in \operatorname{interval}(1, n)) \wedge ((j \in \operatorname{interval}(1, n)) \wedge \neg (i = j))))) \Rightarrow  \\ \quad \quad \operatorname{entry}(\operatorname{identmat}(a, n), i, j) = \operatorname{zero}(a) \end{array}}}
\vnbprop{unitrow-entry-at}{\vnbstmt{\begin{array}{@{}l@{}} \forall a, n, i. \\ \quad (\operatorname{is\text{-}ring}(a) \wedge ((n \in \mathbb{N}) \wedge (i \in \operatorname{interval}(1, n)))) \Rightarrow  \\ \quad \quad \operatorname{entry}(\operatorname{unitrow}(a, n, i), 1, i) = \operatorname{one}(a) \end{array}}}
\vnbprop{unitrow-entry-off}{\vnbstmt{\begin{array}{@{}l@{}} \forall a, n, i, j. \\ \quad (\operatorname{is\text{-}ring}(a) \wedge ((n \in \mathbb{N}) \wedge ((i \in \operatorname{interval}(1, n)) \wedge ((j \in \operatorname{interval}(1, n)) \wedge \neg (j = i))))) \Rightarrow  \\ \quad \quad \operatorname{entry}(\operatorname{unitrow}(a, n, i), 1, j) = \operatorname{zero}(a) \end{array}}}
\vnbprop{entry-of-zeromat}{\vnbstmt{\begin{array}{@{}l@{}} \forall a, m, n, i, j. \\ \quad (\operatorname{is\text{-}ring}(a) \wedge ((m \in \mathbb{N}) \wedge ((n \in \mathbb{N}) \wedge ((i \in \operatorname{interval}(1, m)) \wedge (j \in \operatorname{interval}(1, n)))))) \Rightarrow  \\ \quad \quad \operatorname{entry}(\operatorname{zeromat}(a, m, n), i, j) = \operatorname{zero}(a) \end{array}}}
\vnbprop{elem-f-type}{\vnbstmt{\begin{array}{@{}l@{}} \forall a, n, k, l. \\ \quad (\operatorname{is\text{-}ring}(a) \wedge (n \in \mathbb{N})) \Rightarrow (\operatorname{elem\text{-}f}(a, n, k, l) \in \operatorname{mat}(n, n, \operatorname{carr}(a))) \end{array}}}
\vnbprop{elem-f-symmetric}{\vnbstmt{\begin{array}{@{}l@{}} \forall a, n, k, l. \\ \quad (\operatorname{is\text{-}ring}(a) \wedge (n \in \mathbb{N})) \Rightarrow (\operatorname{elem\text{-}f}(a, n, k, l) = \operatorname{elem\text{-}f}(a, n, l, k)) \end{array}}}
\vnbprop{entry-of-block}{\vnbstmt{\begin{array}{@{}l@{}} \forall p, k, l, i, j, m, n, x. \\ \quad k \in \mathbb{N} \wedge \\ \quad l \in \mathbb{N} \wedge \\ \quad p \in \operatorname{mat}(m, n, x) \wedge \\ \quad k \leq m \wedge \\ \quad l \leq n \wedge \\ \quad i \in \operatorname{interval}(1, k) \wedge \\ \quad j \in \operatorname{interval}(1, l) \Rightarrow \\ \quad \quad \operatorname{entry}(\operatorname{block}(p, k, l), i, j) = \operatorname{entry}(p, i, j) \end{array}}}
\vnbprop{entry-of-snoc-col}{\vnbstmt{\begin{array}{@{}l@{}} \forall w, n, v, i, x. \\ \quad ((n \in \mathbb{N}) \wedge ((w \in \operatorname{mat}(n, 1, x)) \wedge ((v \in x) \wedge (i \in \operatorname{interval}(1, n))))) \Rightarrow  \\ \quad \quad \operatorname{entry}(\operatorname{snoc\text{-}col}(w, n, v), i, 1) = \operatorname{entry}(w, i, 1) \end{array}}}
\vnbprop{snoc-col-last}{\vnbstmt{\begin{array}{@{}l@{}} \forall w, n, v, x. \\ \quad ((n \in \mathbb{N}) \wedge ((w \in \operatorname{mat}(n, 1, x)) \wedge (v \in x))) \Rightarrow  \\ \quad \quad \operatorname{entry}(\operatorname{snoc\text{-}col}(w, n, v), (n + 1), 1) = v \end{array}}}
\vnbprop{entry-of-snoc-row}{\vnbstmt{\begin{array}{@{}l@{}} \forall c, n, r, j, x. \\ \quad ((n \in \mathbb{N}) \wedge ((c \in \operatorname{mat}(1, n, x)) \wedge ((r \in x) \wedge (j \in \operatorname{interval}(1, n))))) \Rightarrow  \\ \quad \quad \operatorname{entry}(\operatorname{snoc\text{-}row}(c, n, r), 1, j) = \operatorname{entry}(c, 1, j) \end{array}}}
\vnbprop{snoc-row-last}{\vnbstmt{\begin{array}{@{}l@{}} \forall c, n, r, x. \\ \quad ((n \in \mathbb{N}) \wedge ((c \in \operatorname{mat}(1, n, x)) \wedge (r \in x))) \Rightarrow  \\ \quad \quad \operatorname{entry}(\operatorname{snoc\text{-}row}(c, n, r), 1, (n + 1)) = r \end{array}}}
\vnbprop{matadd-entry}{\vnbstmt{\begin{array}{@{}l@{}} \forall a, m, n, p, q.; \forall i \in \operatorname{interval}(1, m), j \in \operatorname{interval}(1, n). \\ \quad (\operatorname{is\text{-}ring}(a) \wedge ((p \in \operatorname{mat}(m, n, \operatorname{carr}(a))) \wedge (q \in \operatorname{mat}(m, n, \operatorname{carr}(a))))) \Rightarrow  \\ \quad \quad \operatorname{entry}(\operatorname{matadd}(a, p, q), i, j) = \operatorname{add}(a)(\operatorname{entry}(p, i, j), \operatorname{entry}(q, i, j)) \end{array}}}
\vnbprop{matscale-entry}{\vnbstmt{\begin{array}{@{}l@{}} \forall a, m, n, r, p.; \forall i \in \operatorname{interval}(1, m), j \in \operatorname{interval}(1, n). \\ \quad (\operatorname{is\text{-}ring}(a) \wedge ((r \in \operatorname{carr}(a)) \wedge (p \in \operatorname{mat}(m, n, \operatorname{carr}(a))))) \Rightarrow  \\ \quad \quad \operatorname{entry}(\operatorname{matscale}(a, r, p), i, j) = \operatorname{mul}(a)(r, \operatorname{entry}(p, i, j)) \end{array}}}
\vnbprop{entry-of-submat}{\vnbstmt{\begin{array}{@{}l@{}} \forall s, p, q, i, j, x. \\ \quad p \in \mathbb{N} \wedge \\ \quad q \in \mathbb{N} \wedge \\ \quad s \in \operatorname{mat}((p + 1), (q + 1), x) \wedge \\ \quad i \in \operatorname{interval}(1, p) \wedge \\ \quad j \in \operatorname{interval}(1, q) \Rightarrow \\ \quad \quad \operatorname{entry}(\operatorname{submat}(s, p, q), i, j) = \operatorname{entry}(s, (i + 1), (j + 1)) \end{array}}}

\vnbbookfile{theorem-library/rake-matrix-laws}{theorem-library/rake-matrix-laws}
\noindent{\small rake-matrix-laws.scm -- the eight matrix-algebra laws of `mat-ring-is-ring', PROVEN by ONE entrywise driver.}\par
\vnbprop{matneg-entry}{\vnbstmt{\begin{array}{@{}l@{}} \forall a, m, n, p.; \forall i \in \operatorname{interval}(1, m), j \in \operatorname{interval}(1, n). \\ \quad (\operatorname{is\text{-}ring}(a) \wedge (p \in \operatorname{mat}(m, n, \operatorname{carr}(a)))) \Rightarrow  \\ \quad \quad \operatorname{entry}(\operatorname{matneg}(a, p), i, j) = \operatorname{neg}(a)(\operatorname{entry}(p, i, j)) \end{array}}}
\vnbprop{matadd-comm}{\vnbstmt{\begin{array}{@{}l@{}} \forall a, m, n, p, q. \\ \quad (\operatorname{is\text{-}ring}(a) \wedge ((p \in \operatorname{mat}(m, n, \operatorname{carr}(a))) \wedge (q \in \operatorname{mat}(m, n, \operatorname{carr}(a))))) \Rightarrow  \\ \quad \quad \operatorname{matadd}(a, p, q) = \operatorname{matadd}(a, q, p) \end{array}}}
\vnbprop{matadd-assoc}{\vnbstmt{\begin{array}{@{}l@{}} \forall a, m, n, p, q, r. \\ \quad (\operatorname{is\text{-}ring}(a) \wedge ((p \in \operatorname{mat}(m, n, \operatorname{carr}(a))) \wedge ((q \in \operatorname{mat}(m, n, \operatorname{carr}(a))) \wedge (r \in \operatorname{mat}(m, n, \operatorname{carr}(a)))))) \Rightarrow  \\ \quad \quad \operatorname{matadd}(a, \operatorname{matadd}(a, p, q), r) = \operatorname{matadd}(a, p, \operatorname{matadd}(a, q, r)) \end{array}}}
\vnbprop{matadd-zero-left}{\vnbstmt{\begin{array}{@{}l@{}} \forall a, m, n.; \forall p \in \operatorname{mat}(m, n, \operatorname{carr}(a)). \\ \quad \operatorname{is\text{-}ring}(a) \Rightarrow (\operatorname{matadd}(a, \operatorname{zeromat}(a, m, n), p) = p) \end{array}}}
\vnbprop{matadd-zero-right}{\vnbstmt{\begin{array}{@{}l@{}} \forall a, m, n.; \forall p \in \operatorname{mat}(m, n, \operatorname{carr}(a)). \\ \quad \operatorname{is\text{-}ring}(a) \Rightarrow (\operatorname{matadd}(a, p, \operatorname{zeromat}(a, m, n)) = p) \end{array}}}
\vnbprop{matadd-neg-left}{\vnbstmt{\begin{array}{@{}l@{}} \forall a, m, n.; \forall p \in \operatorname{mat}(m, n, \operatorname{carr}(a)). \\ \quad \operatorname{is\text{-}ring}(a) \Rightarrow (\operatorname{matadd}(a, \operatorname{matneg}(a, p), p) = \operatorname{zeromat}(a, m, n)) \end{array}}}
\vnbprop{matadd-neg-right}{\vnbstmt{\begin{array}{@{}l@{}} \forall a, m, n.; \forall p \in \operatorname{mat}(m, n, \operatorname{carr}(a)). \\ \quad \operatorname{is\text{-}ring}(a) \Rightarrow (\operatorname{matadd}(a, p, \operatorname{matneg}(a, p)) = \operatorname{zeromat}(a, m, n)) \end{array}}}
\vnbprop{matmul-left-dist-guarded}{\vnbstmt{\begin{array}{@{}l@{}} \forall a, m, n, k, p, q, r. \\ \quad \operatorname{is\text{-}ring}(a) \wedge \\ \quad p \in \operatorname{mat}(m, n, \operatorname{carr}(a)) \wedge \\ \quad q \in \operatorname{mat}(n, k, \operatorname{carr}(a)) \wedge \\ \quad r \in \operatorname{mat}(n, k, \operatorname{carr}(a)) \wedge \\ \quad (n = 0) \Rightarrow ((m = 0) \vee (k = 0)) \Rightarrow \\ \quad \quad \operatorname{matmul}(a, p, \operatorname{matadd}(a, q, r)) = \operatorname{matadd}(a, \operatorname{matmul}(a, p, q), \operatorname{matmul}(a, p, r)) \end{array}}}
\vnbprop{matmul-right-dist-guarded}{\vnbstmt{\begin{array}{@{}l@{}} \forall a, m, n, k, p, q, r. \\ \quad \operatorname{is\text{-}ring}(a) \wedge \\ \quad p \in \operatorname{mat}(m, n, \operatorname{carr}(a)) \wedge \\ \quad q \in \operatorname{mat}(m, n, \operatorname{carr}(a)) \wedge \\ \quad r \in \operatorname{mat}(n, k, \operatorname{carr}(a)) \wedge \\ \quad (n = 0) \Rightarrow ((m = 0) \vee (k = 0)) \Rightarrow \\ \quad \quad \operatorname{matmul}(a, \operatorname{matadd}(a, p, q), r) = \operatorname{matadd}(a, \operatorname{matmul}(a, p, r), \operatorname{matmul}(a, q, r)) \end{array}}}

\vnbbookfile{theorem-library/mat-ring-proof}{theorem-library/mat-ring-proof}
\noindent{\small mat-ring-proof.scm -- the n-by-n matrices over a ring form a ring.}\par
\vnbprop{mat-ring-is-ring}{\vnbstmt{\forall a.; \forall n \in \mathbb{N}.\; (\operatorname{is\text{-}ring}(a) \Rightarrow \operatorname{is\text{-}ring}(\operatorname{mat\text{-}ring}(a, n)))}}

\vnbbookfile{theorem-library/matunit-shift-proof}{theorem-library/matunit-shift-proof}
\noindent{\small matunit-shift-proof.scm -- Lemma 3.3 (algebraic-numbers.pdf ch.3): the column-shift formula for right-multiplication by a matrix unit.}\par
\vnbprop{matunit-col-shift}{\vnbstmt{\begin{array}{@{}l@{}} \forall a, m, n, p, k, l.; \forall i \in \operatorname{interval}(1, m), c \in \operatorname{interval}(1, n). \\ \quad (\operatorname{is\text{-}ring}(a) \wedge ((p \in \operatorname{mat}(m, n, \operatorname{carr}(a))) \wedge ((k \in \operatorname{interval}(1, n)) \wedge (l \in \operatorname{interval}(1, n))))) \Rightarrow  \\ \quad \quad \operatorname{entry}(\operatorname{matmul}(a, p, \operatorname{matunit}(a, n, k, l)), i, c) = \left(\text{if }(c = l)\text{ then }\operatorname{entry}(p, i, k)\text{ else }\operatorname{zero}(a)\right) \end{array}}}

\vnbbookfile{theorem-library/elem-actions-proof}{theorem-library/elem-actions-proof}
\noindent{\small elem-actions-proof.scm -- Prop 3.5 (algebraic-numbers.pdf ch.3): the action of the elementary column matrices on a matrix by right-multiplication.  Each is the matmul-entry expansion + a finsum collapse (single-support for H, and for the off-l column of G; two-support for G's l-column), closed by an}\par
\vnbprop{elem-h-action}{\vnbstmt{\begin{array}{@{}l@{}} \forall a, m, n, p, r, k.; \forall i \in \operatorname{interval}(1, m), c \in \operatorname{interval}(1, n). \\ \quad (\operatorname{is\text{-}ring}(a) \wedge ((p \in \operatorname{mat}(m, n, \operatorname{carr}(a))) \wedge ((r \in \operatorname{carr}(a)) \wedge (k \in \operatorname{interval}(1, n))))) \Rightarrow  \\ \quad \quad \operatorname{entry}(\operatorname{matmul}(a, p, \operatorname{elem\text{-}h}(a, n, r, k)), i, c) \\ \quad \quad =\; \left(\text{if }(c = k)\text{ then }\operatorname{mul}(a)(\operatorname{entry}(p, i, k), r)\text{ else }\operatorname{entry}(p, i, c)\right) \end{array}}}
\vnbprop{elem-g-action}{\vnbstmt{\begin{array}{@{}l@{}} \forall a, m, n, p, r, k, l.; \forall i \in \operatorname{interval}(1, m), c \in \operatorname{interval}(1, n). \\ \quad \operatorname{is\text{-}ring}(a) \wedge \\ \quad p \in \operatorname{mat}(m, n, \operatorname{carr}(a)) \wedge \\ \quad r \in \operatorname{carr}(a) \wedge \\ \quad k \in \operatorname{interval}(1, n) \wedge \\ \quad l \in \operatorname{interval}(1, n) \wedge \\ \quad \neg (k = l) \Rightarrow \\ \quad \quad \operatorname{entry}(\operatorname{matmul}(a, p, \operatorname{elem\text{-}g}(a, n, r, k, l)), i, c) \\ \quad \quad =\; \left(\text{if }(c = l)\text{ then }\operatorname{add}(a)(\operatorname{entry}(p, i, l), \operatorname{mul}(a)(\operatorname{entry}(p, i, k), r))\text{ else }\operatorname{entry}(p, i, c)\right) \end{array}}}
\vnbprop{elem-f-action}{\vnbstmt{\begin{array}{@{}l@{}} \forall a, m, n, p, k, l.; \forall i \in \operatorname{interval}(1, m), c \in \operatorname{interval}(1, n). \\ \quad \operatorname{is\text{-}ring}(a) \wedge \\ \quad p \in \operatorname{mat}(m, n, \operatorname{carr}(a)) \wedge \\ \quad k \in \operatorname{interval}(1, n) \wedge \\ \quad l \in \operatorname{interval}(1, n) \wedge \\ \quad \neg (k = l) \Rightarrow \\ \quad \quad \operatorname{entry}(\operatorname{matmul}(a, p, \operatorname{elem\text{-}f}(a, n, k, l)), i, c) \\ \quad \quad =\; \left(\text{if }(c = k)\text{ then }\operatorname{entry}(p, i, l)\text{ else }\left(\text{if }(c = l)\text{ then }\operatorname{entry}(p, i, k)\text{ else }\operatorname{entry}(p, i, c)\right)\right) \end{array}}}

\vnbbookfile{theorem-library/elem-inverses-proof}{theorem-library/elem-inverses-proof}
\noindent{\small elem-inverses-proof.scm -- Cor 3.6 (algebraic-numbers.pdf ch.3): the elementary column matrices are invertible, with elementary inverses.  Each reduces, via matrix-entry-extensionality, to an entry identity whose LHS is the PROVEN Prop 3.5 action (with P = the would-be inverse) and whose RHS is entry-of-identmat;}\par
\vnbprop{elem-h-inverse}{\vnbstmt{\begin{array}{@{}l@{}} \forall a, n, r, s, k. \\ \quad \operatorname{is\text{-}ring}(a) \wedge \\ \quad n \in \mathbb{N} \wedge \\ \quad r \in \operatorname{carr}(a) \wedge \\ \quad s \in \operatorname{carr}(a) \wedge \\ \quad \operatorname{mul}(a)(r, s) = \operatorname{one}(a) \wedge \\ \quad k \in \operatorname{interval}(1, n) \Rightarrow \\ \quad \quad \operatorname{matmul}(a, \operatorname{elem\text{-}h}(a, n, r, k), \operatorname{elem\text{-}h}(a, n, s, k)) = \operatorname{identmat}(a, n) \end{array}}}
\vnbprop{elem-g-inverse}{\vnbstmt{\begin{array}{@{}l@{}} \forall a, n, r, k, l. \\ \quad \operatorname{is\text{-}ring}(a) \wedge \\ \quad n \in \mathbb{N} \wedge \\ \quad r \in \operatorname{carr}(a) \wedge \\ \quad k \in \operatorname{interval}(1, n) \wedge \\ \quad l \in \operatorname{interval}(1, n) \wedge \\ \quad \neg (k = l) \Rightarrow \\ \quad \quad \operatorname{matmul}(a, \operatorname{elem\text{-}g}(a, n, r, k, l), \operatorname{elem\text{-}g}(a, n, \operatorname{neg}(a)(r), k, l)) = \operatorname{identmat}(a, n) \end{array}}}
\vnbprop{elem-f-inverse}{\vnbstmt{\begin{array}{@{}l@{}} \forall a, n, k, l. \\ \quad (\operatorname{is\text{-}ring}(a) \wedge ((n \in \mathbb{N}) \wedge ((k \in \operatorname{interval}(1, n)) \wedge ((l \in \operatorname{interval}(1, n)) \wedge \neg (k = l))))) \Rightarrow  \\ \quad \quad \operatorname{matmul}(a, \operatorname{elem\text{-}f}(a, n, k, l), \operatorname{elem\text{-}f}(a, n, l, k)) = \operatorname{identmat}(a, n) \end{array}}}

\vnbbookfile{theorem-library/elem-row-actions-proof}{theorem-library/elem-row-actions-proof}
\noindent{\small elem-row-actions-proof.scm -- Prop 3.29 (row form, algebraic-numbers.pdf ch.3): the action of the elementary matrices by LEFT-multiplication (row operations). Each mirrors theorem-library/elem-actions-proof.scm with the elementary matrix as the LEFT factor of matmul-entry: (E.P)\_\{ic\} = FINSUM\_j E\_\{ij\}.P\_\{jc\}, so a}\par
\vnbprop{elem-h-row-action}{\vnbstmt{\begin{array}{@{}l@{}} \forall a, m, n, p, r, k.; \forall i \in \operatorname{interval}(1, m), c \in \operatorname{interval}(1, n). \\ \quad (\operatorname{is\text{-}ring}(a) \wedge ((p \in \operatorname{mat}(m, n, \operatorname{carr}(a))) \wedge ((r \in \operatorname{carr}(a)) \wedge (k \in \operatorname{interval}(1, m))))) \Rightarrow  \\ \quad \quad \operatorname{entry}(\operatorname{matmul}(a, \operatorname{elem\text{-}h}(a, m, r, k), p), i, c) \\ \quad \quad =\; \left(\text{if }(i = k)\text{ then }\operatorname{mul}(a)(r, \operatorname{entry}(p, k, c))\text{ else }\operatorname{entry}(p, i, c)\right) \end{array}}}
\vnbprop{elem-g-row-action}{\vnbstmt{\begin{array}{@{}l@{}} \forall a, m, n, p, r, k, l.; \forall i \in \operatorname{interval}(1, m), c \in \operatorname{interval}(1, n). \\ \quad \operatorname{is\text{-}ring}(a) \wedge \\ \quad p \in \operatorname{mat}(m, n, \operatorname{carr}(a)) \wedge \\ \quad r \in \operatorname{carr}(a) \wedge \\ \quad k \in \operatorname{interval}(1, m) \wedge \\ \quad l \in \operatorname{interval}(1, m) \wedge \\ \quad \neg (k = l) \Rightarrow \\ \quad \quad \operatorname{entry}(\operatorname{matmul}(a, \operatorname{elem\text{-}g}(a, m, r, k, l), p), i, c) \\ \quad \quad =\; \left(\text{if }(i = k)\text{ then }\operatorname{add}(a)(\operatorname{entry}(p, k, c), \operatorname{mul}(a)(r, \operatorname{entry}(p, l, c)))\text{ else }\operatorname{entry}(p, i, c)\right) \end{array}}}
\vnbprop{elem-f-row-action}{\vnbstmt{\begin{array}{@{}l@{}} \forall a, m, n, p, k, l.; \forall i \in \operatorname{interval}(1, m), c \in \operatorname{interval}(1, n). \\ \quad \operatorname{is\text{-}ring}(a) \wedge \\ \quad p \in \operatorname{mat}(m, n, \operatorname{carr}(a)) \wedge \\ \quad k \in \operatorname{interval}(1, m) \wedge \\ \quad l \in \operatorname{interval}(1, m) \wedge \\ \quad \neg (k = l) \Rightarrow \\ \quad \quad \operatorname{entry}(\operatorname{matmul}(a, \operatorname{elem\text{-}f}(a, m, k, l), p), i, c) \\ \quad \quad =\; \left(\text{if }(i = k)\text{ then }\operatorname{entry}(p, l, c)\text{ else }\left(\text{if }(i = l)\text{ then }\operatorname{entry}(p, k, c)\text{ else }\operatorname{entry}(p, i, c)\right)\right) \end{array}}}

\vnbbookfile{theorem-library/det-rows}{theorem-library/det-rows}
\noindent{\small det-rows.scm -- the determinant is LINEAR in each row and ALTERNATING; the interchange of two rows, the elementary row operations, and the expansion along any row (batch 25-C, 2026-09-24).  Linear-algebra 3.18 (row clauses), the row form of 3.19, and 3.22 (row half).  All over a COMMUTATIVE ring.}\par
\vnbprop{dtr-ring-lin1}{\vnbstmt{\begin{array}{@{}l@{}} \forall r, s, a, b, c, d. \\ \quad \operatorname{is\text{-}commutative\text{-}ring}(r) \wedge \\ \quad s \in \operatorname{carr}(r) \wedge \\ \quad a \in \operatorname{carr}(r) \wedge \\ \quad b \in \operatorname{carr}(r) \wedge \\ \quad c \in \operatorname{carr}(r) \wedge \\ \quad d \in \operatorname{carr}(r) \Rightarrow \\ \quad \quad \operatorname{mul}(r)(s, \operatorname{mul}(r)(\operatorname{add}(r)(\operatorname{mul}(r)(c, a), b), d)) \\ \quad \quad =\; \operatorname{add}(r)(\operatorname{mul}(r)(c, \operatorname{mul}(r)(s, \operatorname{mul}(r)(a, d))), \operatorname{mul}(r)(s, \operatorname{mul}(r)(b, d))) \end{array}}}
\vnbprop{dtr-ring-lin2}{\vnbstmt{\begin{array}{@{}l@{}} \forall r, s, a, c, x, y. \\ \quad \operatorname{is\text{-}commutative\text{-}ring}(r) \wedge \\ \quad s \in \operatorname{carr}(r) \wedge \\ \quad a \in \operatorname{carr}(r) \wedge \\ \quad c \in \operatorname{carr}(r) \wedge \\ \quad x \in \operatorname{carr}(r) \wedge \\ \quad y \in \operatorname{carr}(r) \Rightarrow \\ \quad \quad \operatorname{mul}(r)(s, \operatorname{mul}(r)(a, \operatorname{add}(r)(\operatorname{mul}(r)(c, x), y))) \\ \quad \quad =\; \operatorname{add}(r)(\operatorname{mul}(r)(c, \operatorname{mul}(r)(s, \operatorname{mul}(r)(a, x))), \operatorname{mul}(r)(s, \operatorname{mul}(r)(a, y))) \end{array}}}
\vnbprop{dtr-finsum-lin}{\vnbstmt{\begin{array}{@{}l@{}} \forall r, s, c, f, g, h. \\ \quad \operatorname{is\text{-}ring}(r) \wedge \\ \quad s \in \mathbf{Set} \wedge \\ \quad |s| \in \mathbb{N} \wedge \\ \quad c \in \operatorname{carr}(r) \wedge \\ \quad \forall z \in s.\; (f(z) \in \operatorname{carr}(r)) \wedge \\ \quad \forall z \in s.\; (g(z) \in \operatorname{carr}(r)) \wedge \\ \quad \forall z \in s.\; (h(z) = \operatorname{add}(r)(\operatorname{mul}(r)(c, f(z)), g(z))) \Rightarrow \\ \quad \quad \operatorname{finsum}(\operatorname{ring\text{-}additive\text{-}ag}(r), h, s) \\ \quad \quad =\; \operatorname{add}(r)(\operatorname{mul}(r)(c, \operatorname{finsum}(\operatorname{ring\text{-}additive\text{-}ag}(r), f, s)), \operatorname{finsum}(\operatorname{ring\text{-}additive\text{-}ag}(r), g, s)) \end{array}}}
\vnbprop{dtr-succ-in}{\vnbstmt{\begin{array}{@{}l@{}} \forall n, u. \\ \quad ((n \in \mathbb{N}) \wedge (u \in \operatorname{interval}(1, n))) \Rightarrow ((u + 1) \in \operatorname{interval}(1, (n + 1))) \end{array}}}
\vnbprop{dtr-skip-in}{\vnbstmt{\begin{array}{@{}l@{}} \forall n, j, v. \\ \quad ((n \in \mathbb{N}) \wedge (v \in \operatorname{interval}(1, n))) \Rightarrow  \\ \quad \quad \left(\text{if }(v < j)\text{ then }v\text{ else }(v + 1)\right) \in \operatorname{interval}(1, (n + 1)) \end{array}}}
\vnbprop{dtr-minor1-entry}{\vnbstmt{\begin{array}{@{}l@{}} \forall r, n, s, j, u, v. \\ \quad \operatorname{is\text{-}ring}(r) \wedge \\ \quad n \in \mathbb{N} \wedge \\ \quad s \in \operatorname{mat}((n + 1), (n + 1), \operatorname{carr}(r)) \wedge \\ \quad u \in \operatorname{interval}(1, n) \wedge \\ \quad v \in \operatorname{interval}(1, n) \Rightarrow \\ \quad \quad \operatorname{entry}(\operatorname{minor}(s, 1, j, n), u, v) = \operatorname{entry}(s, (u + 1), \left(\text{if }(v < j)\text{ then }v\text{ else }(v + 1)\right)) \end{array}}}
\vnbprop{dtr-sign-type}{\vnbstmt{\begin{array}{@{}l@{}} \forall r, j. \\ \quad (\operatorname{is\text{-}ring}(r) \wedge (j \in \mathbb{N})) \Rightarrow  \\ \quad \quad \operatorname{mpow}(\operatorname{ring\text{-}multiplicative\text{-}monoid}(r), \operatorname{neg}(r)(\operatorname{one}(r)), j) \in \operatorname{carr}(r) \end{array}}}
\vnbprop{dtr-minor1-congr}{\vnbstmt{\begin{array}{@{}l@{}} \forall r, n, x, y, p, j. \\ \quad \operatorname{is\text{-}ring}(r) \wedge \\ \quad n \in \mathbb{N} \wedge \\ \quad x \in \operatorname{mat}((n + 1), (n + 1), \operatorname{carr}(r)) \wedge \\ \quad y \in \operatorname{mat}((n + 1), (n + 1), \operatorname{carr}(r)) \wedge \\ \quad j \in \mathbb{N} \wedge \\ \quad p = 1 \wedge \\ \quad \forall i, k \in \operatorname{interval}(1, (n + 1)). \\ \quad \quad \neg (i = p) \Rightarrow (\operatorname{entry}(x, i, k) = \operatorname{entry}(y, i, k)) \Rightarrow \\ \quad \quad \operatorname{minor}(x, 1, j, n) = \operatorname{minor}(y, 1, j, n) \end{array}}}
\vnbprop{det-row-linear-ind}{\vnbstmt{\begin{array}{@{}l@{}} \forall n \in \mathbb{N}, r, a, b, c, p, \mathit{dsc}. \\ \quad \operatorname{is\text{-}commutative\text{-}ring}(r) \wedge \\ \quad a \in \operatorname{mat}(n, n, \operatorname{carr}(r)) \wedge \\ \quad b \in \operatorname{mat}(n, n, \operatorname{carr}(r)) \wedge \\ \quad c \in \operatorname{mat}(n, n, \operatorname{carr}(r)) \wedge \\ \quad p \in \operatorname{interval}(1, n) \wedge \\ \quad \mathit{dsc} \in \operatorname{carr}(r) \wedge \\ \quad \forall i, k \in \operatorname{interval}(1, n). \\ \quad \quad \neg (i = p) \Rightarrow (\operatorname{entry}(b, i, k) = \operatorname{entry}(a, i, k)) \wedge \\ \quad \forall i, k \in \operatorname{interval}(1, n). \\ \quad \quad \neg (i = p) \Rightarrow (\operatorname{entry}(c, i, k) = \operatorname{entry}(a, i, k)) \wedge \\ \quad \forall k \in \operatorname{interval}(1, n). \\ \quad \quad \operatorname{entry}(c, p, k) = \operatorname{add}(r)(\operatorname{mul}(r)(\mathit{dsc}, \operatorname{entry}(a, p, k)), \operatorname{entry}(b, p, k)) \Rightarrow \\ \quad \quad \operatorname{det}(r, n, c) = \operatorname{add}(r)(\operatorname{mul}(r)(\mathit{dsc}, \operatorname{det}(r, n, a)), \operatorname{det}(r, n, b)) \end{array}}}
\vnbprop{det-row-linear}{\vnbstmt{\begin{array}{@{}l@{}} \forall r, n, a, b, c, p, d. \\ \quad \operatorname{is\text{-}commutative\text{-}ring}(r) \wedge \\ \quad n \in \mathbb{N} \wedge \\ \quad a \in \operatorname{mat}(n, n, \operatorname{carr}(r)) \wedge \\ \quad b \in \operatorname{mat}(n, n, \operatorname{carr}(r)) \wedge \\ \quad c \in \operatorname{mat}(n, n, \operatorname{carr}(r)) \wedge \\ \quad p \in \operatorname{interval}(1, n) \wedge \\ \quad d \in \operatorname{carr}(r) \wedge \\ \quad \forall i, k \in \operatorname{interval}(1, n). \\ \quad \quad \neg (i = p) \Rightarrow (\operatorname{entry}(b, i, k) = \operatorname{entry}(a, i, k)) \wedge \\ \quad \forall i, k \in \operatorname{interval}(1, n). \\ \quad \quad \neg (i = p) \Rightarrow (\operatorname{entry}(c, i, k) = \operatorname{entry}(a, i, k)) \wedge \\ \quad \forall k \in \operatorname{interval}(1, n). \\ \quad \quad \operatorname{entry}(c, p, k) = \operatorname{add}(r)(\operatorname{mul}(r)(d, \operatorname{entry}(a, p, k)), \operatorname{entry}(b, p, k)) \Rightarrow \\ \quad \quad \operatorname{det}(r, n, c) = \operatorname{add}(r)(\operatorname{mul}(r)(d, \operatorname{det}(r, n, a)), \operatorname{det}(r, n, b)) \end{array}}}
\vnbprop{dtr-sign-succ}{\vnbstmt{\begin{array}{@{}l@{}} \forall r, x. \\ \quad (\operatorname{is\text{-}ring}(r) \wedge (x \in \mathbb{N})) \Rightarrow  \\ \quad \quad \operatorname{mpow}(\operatorname{ring\text{-}multiplicative\text{-}monoid}(r), \operatorname{neg}(r)(\operatorname{one}(r)), (x + 1)) \\ \quad \quad =\; \operatorname{mul}(r)(\operatorname{neg}(r)(\operatorname{one}(r)), \operatorname{mpow}(\operatorname{ring\text{-}multiplicative\text{-}monoid}(r), \operatorname{neg}(r)(\operatorname{one}(r)), x)) \end{array}}}
\vnbprop{dtr-minor-minor}{\vnbstmt{\begin{array}{@{}l@{}} \forall r, n, a, j, k. \\ \quad \operatorname{is\text{-}ring}(r) \wedge \\ \quad n \in \mathbb{N} \wedge \\ \quad a \in \operatorname{mat}(((n + 1) + 1), ((n + 1) + 1), \operatorname{carr}(r)) \wedge \\ \quad j \in \mathbb{N} \wedge \\ \quad k \in \mathbb{N} \wedge \\ \quad 1 \leq j \wedge \\ \quad j \leq k \wedge \\ \quad k \leq (n + 1) \Rightarrow \\ \quad \quad \operatorname{minor}(\operatorname{minor}(a, 1, j, (n + 1)), 1, k, n) = \operatorname{minor}(\operatorname{minor}(a, 1, (k + 1), (n + 1)), 1, j, n) \end{array}}}
\vnbprop{dtr-ring-zero-add}{\vnbstmt{\begin{array}{@{}l@{}} \forall r, y, z. \\ \quad (\operatorname{is\text{-}commutative\text{-}ring}(r) \wedge ((y \in \operatorname{carr}(r)) \wedge (z \in \operatorname{carr}(r)))) \Rightarrow  \\ \quad \quad \operatorname{add}(r)(\operatorname{add}(r)(\operatorname{zero}(r), y), z) = \operatorname{add}(r)(y, z) \end{array}}}
\vnbprop{dtr-dsum-anti}{\vnbstmt{\begin{array}{@{}l@{}} \forall n \in \mathbb{N}, r, f. \\ \quad \operatorname{is\text{-}commutative\text{-}ring}(r) \wedge \\ \quad \forall j \in \operatorname{interval}(1, (n + 1)), k \in \operatorname{interval}(1, n).\; (f(j, k) \in \operatorname{carr}(r)) \wedge \\ \quad \forall j, k \in \mathbb{N}. \\ \quad \quad ((1 \leq j) \wedge ((j \leq k) \wedge (k \leq n))) \Rightarrow (f((k + 1), j) = \operatorname{neg}(r)(f(j, k))) \Rightarrow \\ \quad \quad \operatorname{finsum}(\operatorname{ring\text{-}additive\text{-}ag}(r), \\ \quad \quad \quad (\lambda a \in \operatorname{interval}(1, (n + 1)).\; \\ \quad \quad \quad \quad \operatorname{finsum}(\operatorname{ring\text{-}additive\text{-}ag}(r), (\lambda b \in \operatorname{interval}(1, n).\; f(a, b)), \operatorname{interval}(1, n))), \\ \quad \quad \quad \operatorname{interval}(1, (n + 1))) \\ \quad \quad =\; \operatorname{zero}(r) \end{array}}}
\vnbprop{dtr-ring-anti}{\vnbstmt{\begin{array}{@{}l@{}} \forall r, s, t, a, b, d. \\ \quad \operatorname{is\text{-}commutative\text{-}ring}(r) \wedge \\ \quad s \in \operatorname{carr}(r) \wedge \\ \quad t \in \operatorname{carr}(r) \wedge \\ \quad a \in \operatorname{carr}(r) \wedge \\ \quad b \in \operatorname{carr}(r) \wedge \\ \quad d \in \operatorname{carr}(r) \Rightarrow \\ \quad \quad \operatorname{mul}(r)(\operatorname{mul}(r)(\operatorname{neg}(r)(\operatorname{one}(r)), t), \operatorname{mul}(r)(b, \operatorname{mul}(r)(s, \operatorname{mul}(r)(a, d)))) \\ \quad \quad =\; \operatorname{neg}(r)(\operatorname{mul}(r)(s, \operatorname{mul}(r)(a, \operatorname{mul}(r)(t, \operatorname{mul}(r)(b, d))))) \end{array}}}
\vnbprop{dtr-det-rows12}{\vnbstmt{\begin{array}{@{}l@{}} \forall n, r, a. \\ \quad n \in \mathbb{N} \wedge \\ \quad \operatorname{is\text{-}commutative\text{-}ring}(r) \wedge \\ \quad a \in \operatorname{mat}(((n + 1) + 1), ((n + 1) + 1), \operatorname{carr}(r)) \wedge \\ \quad \forall k \in \operatorname{interval}(1, ((n + 1) + 1)).\; (\operatorname{entry}(a, 1, k) = \operatorname{entry}(a, 2, k)) \Rightarrow \\ \quad \quad \operatorname{det}(r, ((n + 1) + 1), a) = \operatorname{zero}(r) \end{array}}}
\vnbprop{dtr-det-rows12b}{\vnbstmt{\begin{array}{@{}l@{}} \forall m, n, r, a. \\ \quad n \in \mathbb{N} \wedge \\ \quad m = ((n + 1) + 1) \wedge \\ \quad \operatorname{is\text{-}commutative\text{-}ring}(r) \wedge \\ \quad a \in \operatorname{mat}(m, m, \operatorname{carr}(r)) \wedge \\ \quad \forall k \in \operatorname{interval}(1, m).\; (\operatorname{entry}(a, 1, k) = \operatorname{entry}(a, 2, k)) \Rightarrow \\ \quad \quad \operatorname{det}(r, m, a) = \operatorname{zero}(r) \end{array}}}
\vnbprop{dtr-alt-low}{\vnbstmt{\begin{array}{@{}l@{}} \forall n \in \mathbb{N}, r, a, p, q. \\ \quad \forall r, a, p, q. \\ \quad \quad \operatorname{is\text{-}commutative\text{-}ring}(r) \wedge \\ \quad \quad a \in \operatorname{mat}(n, n, \operatorname{carr}(r)) \wedge \\ \quad \quad p \in \operatorname{interval}(1, n) \wedge \\ \quad \quad q \in \operatorname{interval}(1, n) \wedge \\ \quad \quad \neg (p = q) \wedge \\ \quad \quad \forall k \in \operatorname{interval}(1, n).\; (\operatorname{entry}(a, p, k) = \operatorname{entry}(a, q, k)) \Rightarrow \\ \quad \quad \quad \operatorname{det}(r, n, a) = \operatorname{zero}(r) \wedge \\ \quad \operatorname{is\text{-}commutative\text{-}ring}(r) \wedge \\ \quad a \in \operatorname{mat}((n + 1), (n + 1), \operatorname{carr}(r)) \wedge \\ \quad p \in \operatorname{interval}(1, (n + 1)) \wedge \\ \quad q \in \operatorname{interval}(1, (n + 1)) \wedge \\ \quad \neg (p = q) \wedge \\ \quad \neg (p = 1) \wedge \\ \quad \neg (q = 1) \wedge \\ \quad \forall k \in \operatorname{interval}(1, (n + 1)).\; (\operatorname{entry}(a, p, k) = \operatorname{entry}(a, q, k)) \Rightarrow \\ \quad \quad \operatorname{det}(r, (n + 1), a) = \operatorname{zero}(r) \end{array}}}
\vnbprop{dtr-ring-final}{\vnbstmt{\begin{array}{@{}l@{}} \forall r, t, u_{3}, u_{4}, u_{5}, u_{6}, u_{7}, u_{8}. \\ \quad \operatorname{is\text{-}commutative\text{-}ring}(r) \wedge \\ \quad t \in \operatorname{carr}(r) \wedge \\ \quad u_{3} \in \operatorname{carr}(r) \wedge \\ \quad u_{4} \in \operatorname{carr}(r) \wedge \\ \quad u_{5} \in \operatorname{carr}(r) \wedge \\ \quad u_{6} \in \operatorname{carr}(r) \wedge \\ \quad u_{7} \in \operatorname{carr}(r) \wedge \\ \quad u_{8} \in \operatorname{carr}(r) \wedge \\ \quad u_{3} = \operatorname{add}(r)(\operatorname{mul}(r)(\operatorname{one}(r), u_{4}), u_{5}) \wedge \\ \quad u_{4} = \operatorname{add}(r)(\operatorname{mul}(r)(\operatorname{one}(r), u_{6}), t) \wedge \\ \quad u_{5} = \operatorname{add}(r)(\operatorname{mul}(r)(\operatorname{one}(r), u_{7}), u_{8}) \wedge \\ \quad u_{3} = \operatorname{zero}(r) \wedge \\ \quad u_{6} = \operatorname{zero}(r) \wedge \\ \quad u_{7} = \operatorname{zero}(r) \wedge \\ \quad u_{8} = \operatorname{zero}(r) \Rightarrow \\ \quad \quad t = \operatorname{zero}(r) \end{array}}}
\vnbprop{dtr-alt-row1}{\vnbstmt{\begin{array}{@{}l@{}} \forall n \in \mathbb{N}, r, a, q. \\ \quad \forall r, a, p, q. \\ \quad \quad \operatorname{is\text{-}commutative\text{-}ring}(r) \wedge \\ \quad \quad a \in \operatorname{mat}(n, n, \operatorname{carr}(r)) \wedge \\ \quad \quad p \in \operatorname{interval}(1, n) \wedge \\ \quad \quad q \in \operatorname{interval}(1, n) \wedge \\ \quad \quad \neg (p = q) \wedge \\ \quad \quad \forall k \in \operatorname{interval}(1, n).\; (\operatorname{entry}(a, p, k) = \operatorname{entry}(a, q, k)) \Rightarrow \\ \quad \quad \quad \operatorname{det}(r, n, a) = \operatorname{zero}(r) \wedge \\ \quad \operatorname{is\text{-}commutative\text{-}ring}(r) \wedge \\ \quad a \in \operatorname{mat}((n + 1), (n + 1), \operatorname{carr}(r)) \wedge \\ \quad q \in \operatorname{interval}(1, (n + 1)) \wedge \\ \quad \neg (q = 1) \wedge \\ \quad \forall k \in \operatorname{interval}(1, (n + 1)).\; (\operatorname{entry}(a, 1, k) = \operatorname{entry}(a, q, k)) \Rightarrow \\ \quad \quad \operatorname{det}(r, (n + 1), a) = \operatorname{zero}(r) \end{array}}}
\vnbprop{det-alternating-ind}{\vnbstmt{\begin{array}{@{}l@{}} \forall n \in \mathbb{N}, r, a, p, q. \\ \quad \operatorname{is\text{-}commutative\text{-}ring}(r) \wedge \\ \quad a \in \operatorname{mat}(n, n, \operatorname{carr}(r)) \wedge \\ \quad p \in \operatorname{interval}(1, n) \wedge \\ \quad q \in \operatorname{interval}(1, n) \wedge \\ \quad \neg (p = q) \wedge \\ \quad \forall k \in \operatorname{interval}(1, n).\; (\operatorname{entry}(a, p, k) = \operatorname{entry}(a, q, k)) \Rightarrow \\ \quad \quad \operatorname{det}(r, n, a) = \operatorname{zero}(r) \end{array}}}
\vnbprop{det-alternating-rows}{\vnbstmt{\begin{array}{@{}l@{}} \forall r, n, a, p, q. \\ \quad \operatorname{is\text{-}commutative\text{-}ring}(r) \wedge \\ \quad n \in \mathbb{N} \wedge \\ \quad a \in \operatorname{mat}(n, n, \operatorname{carr}(r)) \wedge \\ \quad p \in \operatorname{interval}(1, n) \wedge \\ \quad q \in \operatorname{interval}(1, n) \wedge \\ \quad \neg (p = q) \wedge \\ \quad \forall k \in \operatorname{interval}(1, n).\; (\operatorname{entry}(a, p, k) = \operatorname{entry}(a, q, k)) \Rightarrow \\ \quad \quad \operatorname{det}(r, n, a) = \operatorname{zero}(r) \end{array}}}
\vnbprop{dtr-ring-swap}{\vnbstmt{\begin{array}{@{}l@{}} \forall r, a, b, u_{1}, u_{2}, u_{3}, u_{4}, u_{5}. \\ \quad \operatorname{is\text{-}commutative\text{-}ring}(r) \wedge \\ \quad a \in \operatorname{carr}(r) \wedge \\ \quad b \in \operatorname{carr}(r) \wedge \\ \quad u_{1} \in \operatorname{carr}(r) \wedge \\ \quad u_{2} \in \operatorname{carr}(r) \wedge \\ \quad u_{3} \in \operatorname{carr}(r) \wedge \\ \quad u_{4} \in \operatorname{carr}(r) \wedge \\ \quad u_{5} \in \operatorname{carr}(r) \wedge \\ \quad u_{1} = \operatorname{add}(r)(\operatorname{mul}(r)(\operatorname{one}(r), u_{2}), u_{3}) \wedge \\ \quad u_{2} = \operatorname{add}(r)(\operatorname{mul}(r)(\operatorname{one}(r), u_{4}), a) \wedge \\ \quad u_{3} = \operatorname{add}(r)(\operatorname{mul}(r)(\operatorname{one}(r), b), u_{5}) \wedge \\ \quad u_{1} = \operatorname{zero}(r) \wedge \\ \quad u_{4} = \operatorname{zero}(r) \wedge \\ \quad u_{5} = \operatorname{zero}(r) \Rightarrow \\ \quad \quad b = \operatorname{neg}(r)(a) \end{array}}}
\vnbprop{dtr-swap-s}{\vnbstmt{\begin{array}{@{}l@{}} \forall n, r, a, b, p, q. \\ \quad n \in \mathbb{N} \wedge \\ \quad \operatorname{is\text{-}commutative\text{-}ring}(r) \wedge \\ \quad a \in \operatorname{mat}((n + 1), (n + 1), \operatorname{carr}(r)) \wedge \\ \quad b \in \operatorname{mat}((n + 1), (n + 1), \operatorname{carr}(r)) \wedge \\ \quad p \in \operatorname{interval}(1, (n + 1)) \wedge \\ \quad q \in \operatorname{interval}(1, (n + 1)) \wedge \\ \quad \neg (p = q) \wedge \\ \quad \forall i, k \in \operatorname{interval}(1, (n + 1)). \\ \quad \quad (\neg (i = p) \wedge \neg (i = q)) \Rightarrow (\operatorname{entry}(b, i, k) = \operatorname{entry}(a, i, k)) \wedge \\ \quad \forall k \in \operatorname{interval}(1, (n + 1)).\; (\operatorname{entry}(b, p, k) = \operatorname{entry}(a, q, k)) \wedge \\ \quad \forall k \in \operatorname{interval}(1, (n + 1)).\; (\operatorname{entry}(b, q, k) = \operatorname{entry}(a, p, k)) \Rightarrow \\ \quad \quad \operatorname{det}(r, (n + 1), b) = \operatorname{neg}(r)(\operatorname{det}(r, (n + 1), a)) \end{array}}}
\vnbprop{det-swap-rows}{\vnbstmt{\begin{array}{@{}l@{}} \forall r, n, a, b, p, q. \\ \quad \operatorname{is\text{-}commutative\text{-}ring}(r) \wedge \\ \quad n \in \mathbb{N} \wedge \\ \quad a \in \operatorname{mat}(n, n, \operatorname{carr}(r)) \wedge \\ \quad b \in \operatorname{mat}(n, n, \operatorname{carr}(r)) \wedge \\ \quad p \in \operatorname{interval}(1, n) \wedge \\ \quad q \in \operatorname{interval}(1, n) \wedge \\ \quad \neg (p = q) \wedge \\ \quad \forall i, k \in \operatorname{interval}(1, n). \\ \quad \quad (\neg (i = p) \wedge \neg (i = q)) \Rightarrow (\operatorname{entry}(b, i, k) = \operatorname{entry}(a, i, k)) \wedge \\ \quad \forall k \in \operatorname{interval}(1, n).\; (\operatorname{entry}(b, p, k) = \operatorname{entry}(a, q, k)) \wedge \\ \quad \forall k \in \operatorname{interval}(1, n).\; (\operatorname{entry}(b, q, k) = \operatorname{entry}(a, p, k)) \Rightarrow \\ \quad \quad \operatorname{det}(r, n, b) = \operatorname{neg}(r)(\operatorname{det}(r, n, a)) \end{array}}}
\vnbprop{dtr-ring-negself}{\vnbstmt{\begin{array}{@{}l@{}} \forall r.; \forall u \in \operatorname{carr}(r). \\ \quad (\operatorname{is\text{-}commutative\text{-}ring}(r) \wedge (u = \operatorname{add}(r)(\operatorname{mul}(r)(\operatorname{neg}(r)(\operatorname{one}(r)), u), u))) \Rightarrow  \\ \quad \quad u = \operatorname{zero}(r) \end{array}}}
\vnbprop{dtr-zero-row-s}{\vnbstmt{\begin{array}{@{}l@{}} \forall n, r, a, k. \\ \quad n \in \mathbb{N} \wedge \\ \quad \operatorname{is\text{-}commutative\text{-}ring}(r) \wedge \\ \quad a \in \operatorname{mat}((n + 1), (n + 1), \operatorname{carr}(r)) \wedge \\ \quad k \in \operatorname{interval}(1, (n + 1)) \wedge \\ \quad \forall \mathit{drk} \in \operatorname{interval}(1, (n + 1)).\; (\operatorname{entry}(a, k, \mathit{drk}) = \operatorname{zero}(r)) \Rightarrow \\ \quad \quad \operatorname{det}(r, (n + 1), a) = \operatorname{zero}(r) \end{array}}}
\vnbprop{dtr-elem-h-s}{\vnbstmt{\begin{array}{@{}l@{}} \forall n, r, a, \mathit{dsr}, k. \\ \quad n \in \mathbb{N} \wedge \\ \quad \operatorname{is\text{-}commutative\text{-}ring}(r) \wedge \\ \quad a \in \operatorname{mat}((n + 1), (n + 1), \operatorname{carr}(r)) \wedge \\ \quad \mathit{dsr} \in \operatorname{carr}(r) \wedge \\ \quad k \in \operatorname{interval}(1, (n + 1)) \Rightarrow \\ \quad \quad \operatorname{det}(r, (n + 1), \operatorname{matmul}(r, \operatorname{elem\text{-}h}(r, (n + 1), \mathit{dsr}, k), a)) = \operatorname{mul}(r)(\mathit{dsr}, \operatorname{det}(r, (n + 1), a)) \end{array}}}
\vnbprop{dtr-elem-g-s}{\vnbstmt{\begin{array}{@{}l@{}} \forall n, r, a, \mathit{dsr}, k, l. \\ \quad n \in \mathbb{N} \wedge \\ \quad \operatorname{is\text{-}commutative\text{-}ring}(r) \wedge \\ \quad a \in \operatorname{mat}((n + 1), (n + 1), \operatorname{carr}(r)) \wedge \\ \quad \mathit{dsr} \in \operatorname{carr}(r) \wedge \\ \quad k \in \operatorname{interval}(1, (n + 1)) \wedge \\ \quad l \in \operatorname{interval}(1, (n + 1)) \wedge \\ \quad \neg (k = l) \Rightarrow \\ \quad \quad \operatorname{det}(r, (n + 1), \operatorname{matmul}(r, \operatorname{elem\text{-}g}(r, (n + 1), \mathit{dsr}, k, l), a)) = \operatorname{det}(r, (n + 1), a) \end{array}}}
\vnbprop{dtr-elem-f-s}{\vnbstmt{\begin{array}{@{}l@{}} \forall n, r, a, k, l. \\ \quad n \in \mathbb{N} \wedge \\ \quad \operatorname{is\text{-}commutative\text{-}ring}(r) \wedge \\ \quad a \in \operatorname{mat}((n + 1), (n + 1), \operatorname{carr}(r)) \wedge \\ \quad k \in \operatorname{interval}(1, (n + 1)) \wedge \\ \quad l \in \operatorname{interval}(1, (n + 1)) \wedge \\ \quad \neg (k = l) \Rightarrow \\ \quad \quad \operatorname{det}(r, (n + 1), \operatorname{matmul}(r, \operatorname{elem\text{-}f}(r, (n + 1), k, l), a)) = \operatorname{neg}(r)(\operatorname{det}(r, (n + 1), a)) \end{array}}}
\vnbprop{det-elem-h-row}{\vnbstmt{\begin{array}{@{}l@{}} \forall r, n, a, c, k. \\ \quad \operatorname{is\text{-}commutative\text{-}ring}(r) \wedge \\ \quad n \in \mathbb{N} \wedge \\ \quad a \in \operatorname{mat}(n, n, \operatorname{carr}(r)) \wedge \\ \quad c \in \operatorname{carr}(r) \wedge \\ \quad k \in \operatorname{interval}(1, n) \Rightarrow \\ \quad \quad \operatorname{det}(r, n, \operatorname{matmul}(r, \operatorname{elem\text{-}h}(r, n, c, k), a)) = \operatorname{mul}(r)(c, \operatorname{det}(r, n, a)) \end{array}}}
\vnbprop{det-elem-g-row}{\vnbstmt{\begin{array}{@{}l@{}} \forall r, n, a, c, k, l. \\ \quad \operatorname{is\text{-}commutative\text{-}ring}(r) \wedge \\ \quad n \in \mathbb{N} \wedge \\ \quad a \in \operatorname{mat}(n, n, \operatorname{carr}(r)) \wedge \\ \quad c \in \operatorname{carr}(r) \wedge \\ \quad k \in \operatorname{interval}(1, n) \wedge \\ \quad l \in \operatorname{interval}(1, n) \wedge \\ \quad \neg (k = l) \Rightarrow \\ \quad \quad \operatorname{det}(r, n, \operatorname{matmul}(r, \operatorname{elem\text{-}g}(r, n, c, k, l), a)) = \operatorname{det}(r, n, a) \end{array}}}
\vnbprop{det-elem-f-row}{\vnbstmt{\begin{array}{@{}l@{}} \forall r, n, a, k, l. \\ \quad \operatorname{is\text{-}commutative\text{-}ring}(r) \wedge \\ \quad n \in \mathbb{N} \wedge \\ \quad a \in \operatorname{mat}(n, n, \operatorname{carr}(r)) \wedge \\ \quad k \in \operatorname{interval}(1, n) \wedge \\ \quad l \in \operatorname{interval}(1, n) \wedge \\ \quad \neg (k = l) \Rightarrow \\ \quad \quad \operatorname{det}(r, n, \operatorname{matmul}(r, \operatorname{elem\text{-}f}(r, n, k, l), a)) = \operatorname{neg}(r)(\operatorname{det}(r, n, a)) \end{array}}}
\vnbprop{dtr-ring-rx}{\vnbstmt{\begin{array}{@{}l@{}} \forall r, d, e, t. \\ \quad \operatorname{is\text{-}commutative\text{-}ring}(r) \wedge \\ \quad d \in \operatorname{carr}(r) \wedge \\ \quad e \in \operatorname{carr}(r) \wedge \\ \quad t \in \operatorname{carr}(r) \wedge \\ \quad e = \operatorname{neg}(r)(d) \wedge \\ \quad e = t \Rightarrow \\ \quad \quad d = \operatorname{add}(r)(\operatorname{mul}(r)(\operatorname{neg}(r)(\operatorname{one}(r)), t), \operatorname{zero}(r)) \end{array}}}
\vnbprop{dtr-ring-rx2}{\vnbstmt{\begin{array}{@{}l@{}} \forall r, s, a, d. \\ \quad (\operatorname{is\text{-}commutative\text{-}ring}(r) \wedge ((s \in \operatorname{carr}(r)) \wedge ((a \in \operatorname{carr}(r)) \wedge (d \in \operatorname{carr}(r))))) \Rightarrow  \\ \quad \quad \operatorname{mul}(r)(\operatorname{mul}(r)(\operatorname{neg}(r)(\operatorname{one}(r)), s), \operatorname{mul}(r)(a, d)) \\ \quad \quad =\; \operatorname{add}(r)(\operatorname{mul}(r)(\operatorname{neg}(r)(\operatorname{one}(r)), \operatorname{mul}(r)(s, \operatorname{mul}(r)(a, d))), \operatorname{zero}(r)) \end{array}}}
\vnbprop{det-expand-row-ind}{\vnbstmt{\begin{array}{@{}l@{}} \forall i \in \mathbb{N}, n, r, a. \\ \quad n \in \mathbb{N} \wedge \\ \quad \operatorname{is\text{-}commutative\text{-}ring}(r) \wedge \\ \quad a \in \operatorname{mat}((n + 1), (n + 1), \operatorname{carr}(r)) \wedge \\ \quad 1 \leq i \wedge \\ \quad i \leq (n + 1) \Rightarrow \\ \quad \quad \operatorname{det}(r, (n + 1), a) \\ \quad \quad =\; \operatorname{finsum}(\operatorname{ring\text{-}additive\text{-}ag}(r), \\ \quad \quad \quad (\lambda j \in \operatorname{interval}(1, (n + 1)).\; \\ \quad \quad \quad \quad \operatorname{mul}(r)(\operatorname{mpow}(\operatorname{ring\text{-}multiplicative\text{-}monoid}(r), \operatorname{neg}(r)(\operatorname{one}(r)), (i + j)), \\ \quad \quad \quad \quad \quad \operatorname{mul}(r)(\operatorname{entry}(a, i, j), \operatorname{det}(r, n, \operatorname{minor}(a, i, j, n))))), \\ \quad \quad \quad \operatorname{interval}(1, (n + 1))) \end{array}}}
\vnbprop{det-expand-row}{\vnbstmt{\begin{array}{@{}l@{}} \forall r, n, a, i. \\ \quad \operatorname{is\text{-}commutative\text{-}ring}(r) \wedge \\ \quad n \in \mathbb{N} \wedge \\ \quad a \in \operatorname{mat}((n + 1), (n + 1), \operatorname{carr}(r)) \wedge \\ \quad i \in \operatorname{interval}(1, (n + 1)) \Rightarrow \\ \quad \quad \operatorname{det}(r, (n + 1), a) \\ \quad \quad =\; \operatorname{finsum}(\operatorname{ring\text{-}additive\text{-}ag}(r), \\ \quad \quad \quad (\lambda j \in \operatorname{interval}(1, (n + 1)).\; \\ \quad \quad \quad \quad \operatorname{mul}(r)(\operatorname{mpow}(\operatorname{ring\text{-}multiplicative\text{-}monoid}(r), \operatorname{neg}(r)(\operatorname{one}(r)), (i + j)), \\ \quad \quad \quad \quad \quad \operatorname{mul}(r)(\operatorname{entry}(a, i, j), \operatorname{det}(r, n, \operatorname{minor}(a, i, j, n))))), \\ \quad \quad \quad \operatorname{interval}(1, (n + 1))) \end{array}}}

\vnbbookfile{theorem-library/det-mul}{theorem-library/det-mul}
\noindent{\small det-mul.scm -- det(E A) = det(E) det(A) for an ELEMENTARY matrix E (batch 25-C, 2026-09-24).  The multiplicativity of the determinant as far as the tree carries it; see the report (scratchpad/triage/BATCH25-REPORTS.md, 25-C) for what stands between this and linear-algebra 3.20 / 3.27.}\par
\vnbprop{det-elem-h}{\vnbstmt{\begin{array}{@{}l@{}} \forall r, n, c, k. \\ \quad (\operatorname{is\text{-}commutative\text{-}ring}(r) \wedge ((n \in \mathbb{N}) \wedge ((c \in \operatorname{carr}(r)) \wedge (k \in \operatorname{interval}(1, n))))) \Rightarrow  \\ \quad \quad \operatorname{det}(r, n, \operatorname{elem\text{-}h}(r, n, c, k)) = c \end{array}}}
\vnbprop{det-elem-g}{\vnbstmt{\begin{array}{@{}l@{}} \forall r, n, c, k, l. \\ \quad \operatorname{is\text{-}commutative\text{-}ring}(r) \wedge \\ \quad n \in \mathbb{N} \wedge \\ \quad c \in \operatorname{carr}(r) \wedge \\ \quad k \in \operatorname{interval}(1, n) \wedge \\ \quad l \in \operatorname{interval}(1, n) \wedge \\ \quad \neg (k = l) \Rightarrow \\ \quad \quad \operatorname{det}(r, n, \operatorname{elem\text{-}g}(r, n, c, k, l)) = \operatorname{one}(r) \end{array}}}
\vnbprop{det-elem-f}{\vnbstmt{\begin{array}{@{}l@{}} \forall r, n, k, l. \\ \quad (\operatorname{is\text{-}commutative\text{-}ring}(r) \wedge ((n \in \mathbb{N}) \wedge ((k \in \operatorname{interval}(1, n)) \wedge ((l \in \operatorname{interval}(1, n)) \wedge \neg (k = l))))) \Rightarrow  \\ \quad \quad \operatorname{det}(r, n, \operatorname{elem\text{-}f}(r, n, k, l)) = \operatorname{neg}(r)(\operatorname{one}(r)) \end{array}}}
\vnbprop{det-mul-elem-h}{\vnbstmt{\begin{array}{@{}l@{}} \forall r, n, a, c, k. \\ \quad \operatorname{is\text{-}commutative\text{-}ring}(r) \wedge \\ \quad n \in \mathbb{N} \wedge \\ \quad a \in \operatorname{mat}(n, n, \operatorname{carr}(r)) \wedge \\ \quad c \in \operatorname{carr}(r) \wedge \\ \quad k \in \operatorname{interval}(1, n) \Rightarrow \\ \quad \quad \operatorname{det}(r, n, \operatorname{matmul}(r, \operatorname{elem\text{-}h}(r, n, c, k), a)) = \operatorname{mul}(r)(\operatorname{det}(r, n, \operatorname{elem\text{-}h}(r, n, c, k)), \operatorname{det}(r, n, a)) \end{array}}}
\vnbprop{det-mul-elem-g}{\vnbstmt{\begin{array}{@{}l@{}} \forall r, n, a, c, k, l. \\ \quad \operatorname{is\text{-}commutative\text{-}ring}(r) \wedge \\ \quad n \in \mathbb{N} \wedge \\ \quad a \in \operatorname{mat}(n, n, \operatorname{carr}(r)) \wedge \\ \quad c \in \operatorname{carr}(r) \wedge \\ \quad k \in \operatorname{interval}(1, n) \wedge \\ \quad l \in \operatorname{interval}(1, n) \wedge \\ \quad \neg (k = l) \Rightarrow \\ \quad \quad \operatorname{det}(r, n, \operatorname{matmul}(r, \operatorname{elem\text{-}g}(r, n, c, k, l), a)) \\ \quad \quad =\; \operatorname{mul}(r)(\operatorname{det}(r, n, \operatorname{elem\text{-}g}(r, n, c, k, l)), \operatorname{det}(r, n, a)) \end{array}}}
\vnbprop{det-mul-elem-f}{\vnbstmt{\begin{array}{@{}l@{}} \forall r, n, a, k, l. \\ \quad \operatorname{is\text{-}commutative\text{-}ring}(r) \wedge \\ \quad n \in \mathbb{N} \wedge \\ \quad a \in \operatorname{mat}(n, n, \operatorname{carr}(r)) \wedge \\ \quad k \in \operatorname{interval}(1, n) \wedge \\ \quad l \in \operatorname{interval}(1, n) \wedge \\ \quad \neg (k = l) \Rightarrow \\ \quad \quad \operatorname{det}(r, n, \operatorname{matmul}(r, \operatorname{elem\text{-}f}(r, n, k, l), a)) = \operatorname{mul}(r)(\operatorname{det}(r, n, \operatorname{elem\text{-}f}(r, n, k, l)), \operatorname{det}(r, n, a)) \end{array}}}

\vnbbookfile{theorem-library/pivot-col-reduce-proof}{theorem-library/pivot-col-reduce-proof}
\noindent{\small pivot-col-reduce-proof.scm -- Brick B2b (LA Phase B, algebraic-numbers.pdf ch.3): ONE elementary column operation reduces a first-row entry against the pivot.  Over a euclidean ring A, for P with nonzero pivot P\_\{1,1\} and target column j/=1, there are q,r with (P . G[-q,1,j])\_\{1,j\} = r and (r=0 or}\par
\vnbprop{pivot-col-reduce}{\vnbstmt{\begin{array}{@{}l@{}} \forall a, m, n, p, j. \\ \quad \operatorname{is\text{-}euclidean\text{-}ring}(a) \wedge \\ \quad p \in \operatorname{mat}(m, n, \operatorname{carr}(a)) \wedge \\ \quad 1 \in \operatorname{interval}(1, m) \wedge \\ \quad 1 \in \operatorname{interval}(1, n) \wedge \\ \quad j \in \operatorname{interval}(1, n) \wedge \\ \quad \neg (1 = j) \wedge \\ \quad \neg (\operatorname{entry}(p, 1, 1) = \operatorname{zero}(a)) \Rightarrow \\ \quad \quad \exists q, r \in \operatorname{carr}(a). \\ \quad \quad \quad \operatorname{entry}(\operatorname{matmul}(a, p, \operatorname{elem\text{-}g}(a, n, \operatorname{neg}(a)(q), 1, j)), 1, j) = r \wedge \\ \quad \quad \quad (r = \operatorname{zero}(a)) \vee ((\operatorname{gauge}(a)(r) + 1) \leq \operatorname{gauge}(a)(\operatorname{entry}(p, 1, 1))) \end{array}}}

\vnbbookfile{theorem-library/pivot-row-reduce-proof}{theorem-library/pivot-row-reduce-proof}
\noindent{\small pivot-row-reduce-proof.scm -- Brick B3a (LA Phase B, algebraic-numbers.pdf ch.3): the ROW-reduction step, mirror of pivot-col-reduce.  Over a euclidean ring A, for P with nonzero pivot P\_\{1,1\} and target row i/=1, there are q,r with (ELEM-G[-q,i,1] . P)\_\{i,1\} = r and (r=0 or deg(r) $<$ deg(P\_\{1,1\})) -- one}\par
\vnbprop{pivot-row-reduce}{\vnbstmt{\begin{array}{@{}l@{}} \forall a, m, n, p, i. \\ \quad \operatorname{is\text{-}euclidean\text{-}ring}(a) \wedge \\ \quad p \in \operatorname{mat}(m, n, \operatorname{carr}(a)) \wedge \\ \quad 1 \in \operatorname{interval}(1, m) \wedge \\ \quad 1 \in \operatorname{interval}(1, n) \wedge \\ \quad i \in \operatorname{interval}(1, m) \wedge \\ \quad \neg (i = 1) \wedge \\ \quad \neg (\operatorname{entry}(p, 1, 1) = \operatorname{zero}(a)) \Rightarrow \\ \quad \quad \exists q, r \in \operatorname{carr}(a). \\ \quad \quad \quad \operatorname{entry}(\operatorname{matmul}(a, \operatorname{elem\text{-}g}(a, m, \operatorname{neg}(a)(q), i, 1), p), i, 1) = r \wedge \\ \quad \quad \quad (r = \operatorname{zero}(a)) \vee ((\operatorname{gauge}(a)(r) + 1) \leq \operatorname{gauge}(a)(\operatorname{entry}(p, 1, 1))) \end{array}}}

\vnbbookfile{theorem-library/mat-equiv-proof}{theorem-library/mat-equiv-proof}
\noindent{\small mat-equiv-proof.scm -- Brick B3 infrastructure: basic laws of \textasciitilde{} (mat-equiv.scm). invertible-mat-is-mat, identmat-invertible, mat-equiv-refl, mat-equiv-right-mult, mat-equiv-left-mult, elem-f-invertible.  \textasciitilde{} is reflexive and each elementary swap (row/col) yields an equivalent matrix.  elem-g-invertible + transitivity + ...}\par
\vnbprop{invertible-mat-is-mat}{\vnbstmt{\forall a, n, u.\; (\operatorname{is\text{-}invertible\text{-}mat}(a, n, u) \Rightarrow (u \in \operatorname{mat}(n, n, \operatorname{carr}(a))))}}
\vnbprop{identmat-invertible}{\vnbstmt{\forall a.; \forall n \in \mathbb{N}.\; (\operatorname{is\text{-}ring}(a) \Rightarrow \operatorname{is\text{-}invertible\text{-}mat}(a, n, \operatorname{identmat}(a, n)))}}
\vnbprop{mat-equiv-refl}{\vnbstmt{\begin{array}{@{}l@{}} \forall a, m, n.; \forall c \in \operatorname{mat}(m, n, \operatorname{carr}(a)). \\ \quad \operatorname{is\text{-}ring}(a) \Rightarrow \operatorname{mat\text{-}equiv}(a, m, n, c, c) \end{array}}}
\vnbprop{mat-equiv-right-mult}{\vnbstmt{\begin{array}{@{}l@{}} \forall a, m, n, c, v. \\ \quad (\operatorname{is\text{-}ring}(a) \wedge ((c \in \operatorname{mat}(m, n, \operatorname{carr}(a))) \wedge \operatorname{is\text{-}invertible\text{-}mat}(a, n, v))) \Rightarrow  \\ \quad \quad \operatorname{mat\text{-}equiv}(a, m, n, c, \operatorname{matmul}(a, c, v)) \end{array}}}
\vnbprop{mat-equiv-left-mult}{\vnbstmt{\begin{array}{@{}l@{}} \forall a, m, n, c, u. \\ \quad (\operatorname{is\text{-}ring}(a) \wedge ((c \in \operatorname{mat}(m, n, \operatorname{carr}(a))) \wedge \operatorname{is\text{-}invertible\text{-}mat}(a, m, u))) \Rightarrow  \\ \quad \quad \operatorname{mat\text{-}equiv}(a, m, n, c, \operatorname{matmul}(a, u, c)) \end{array}}}
\vnbprop{elem-f-invertible}{\vnbstmt{\begin{array}{@{}l@{}} \forall a, n, k, l. \\ \quad (\operatorname{is\text{-}ring}(a) \wedge ((n \in \mathbb{N}) \wedge ((k \in \operatorname{interval}(1, n)) \wedge ((l \in \operatorname{interval}(1, n)) \wedge \neg (k = l))))) \Rightarrow  \\ \quad \quad \operatorname{is\text{-}invertible\text{-}mat}(a, n, \operatorname{elem\text{-}f}(a, n, k, l)) \end{array}}}
\vnbprop{elem-g-param-cong}{\vnbstmt{\begin{array}{@{}l@{}} \forall a, n, r, s, k, l. \\ \quad (\operatorname{is\text{-}ring}(a) \wedge ((n \in \mathbb{N}) \wedge ((s \in \operatorname{carr}(a)) \wedge (r = s)))) \Rightarrow  \\ \quad \quad \operatorname{elem\text{-}g}(a, n, r, k, l) = \operatorname{elem\text{-}g}(a, n, s, k, l) \end{array}}}
\vnbprop{elem-g-invertible}{\vnbstmt{\begin{array}{@{}l@{}} \forall a, n, r, k, l. \\ \quad \operatorname{is\text{-}ring}(a) \wedge \\ \quad n \in \mathbb{N} \wedge \\ \quad r \in \operatorname{carr}(a) \wedge \\ \quad k \in \operatorname{interval}(1, n) \wedge \\ \quad l \in \operatorname{interval}(1, n) \wedge \\ \quad \neg (k = l) \Rightarrow \\ \quad \quad \operatorname{is\text{-}invertible\text{-}mat}(a, n, \operatorname{elem\text{-}g}(a, n, r, k, l)) \end{array}}}
\vnbprop{product-of-invertibles-is-invertible}{\vnbstmt{\begin{array}{@{}l@{}} \forall a, n, b, c. \\ \quad (\operatorname{is\text{-}ring}(a) \wedge (\operatorname{is\text{-}invertible\text{-}mat}(a, n, b) \wedge \operatorname{is\text{-}invertible\text{-}mat}(a, n, c))) \Rightarrow  \\ \quad \quad \operatorname{is\text{-}invertible\text{-}mat}(a, n, \operatorname{matmul}(a, b, c)) \end{array}}}
\vnbprop{mat-equiv-trans}{\vnbstmt{\begin{array}{@{}l@{}} \forall a, m, n, c, d, e. \\ \quad (\operatorname{is\text{-}ring}(a) \wedge ((c \in \operatorname{mat}(m, n, \operatorname{carr}(a))) \wedge (\operatorname{mat\text{-}equiv}(a, m, n, c, d) \wedge \operatorname{mat\text{-}equiv}(a, m, n, d, e)))) \Rightarrow  \\ \quad \quad \operatorname{mat\text{-}equiv}(a, m, n, c, e) \end{array}}}

\vnbbookfile{theorem-library/min-degree-entry-proof}{theorem-library/min-degree-entry-proof}
\noindent{\small min-degree-entry-proof.scm -- a matrix with a nonzero entry has a nonzero entry of MINIMAL Euclidean degree.  The "mu is achieved" fact that starts the Smith reduction (Prop 3.36).}\par
\vnbprop{min-degree-entry}{\vnbstmt{\begin{array}{@{}l@{}} \forall a, m, n.; \forall p \in \operatorname{mat}(m, n, \operatorname{carr}(a)). \\ \quad \operatorname{is\text{-}euclidean\text{-}ring}(a) \wedge \\ \quad \exists b \in \operatorname{interval}(1, m), c \in \operatorname{interval}(1, n). \\ \quad \quad \neg (\operatorname{entry}(p, b, c) = \operatorname{zero}(a)) \Rightarrow \\ \quad \quad \exists s \in \operatorname{interval}(1, m), d \in \operatorname{interval}(1, n). \\ \quad \quad \quad \neg (\operatorname{entry}(p, s, d) = \operatorname{zero}(a)) \wedge \\ \quad \quad \quad \forall i, j. \\ \quad \quad \quad \quad ((i \in \operatorname{interval}(1, m)) \wedge ((j \in \operatorname{interval}(1, n)) \wedge \neg (\operatorname{entry}(p, i, j) = \operatorname{zero}(a)))) \Rightarrow  \\ \quad \quad \quad \quad \quad \operatorname{gauge}(a)(\operatorname{entry}(p, s, d)) \leq \operatorname{gauge}(a)(\operatorname{entry}(p, i, j)) \end{array}}}

\vnbbookfile{theorem-library/smith-proof}{theorem-library/smith-proof}
\noindent{\small smith-proof.scm -- Brick B3 (LA Phase B): the Smith normal-form reduction, built on the \textasciitilde{} equivalence infrastructure (mat-equiv-proof.scm). equiv-mul-both  -- A \textasciitilde{} U.A.V for U,V invertible (sandwiching preserves \textasciitilde{}); swap-to-corner  -- bring any entry (i0,j0) to position (1,1) preserving \textasciitilde{}}\par
\vnbprop{equiv-mul-both}{\vnbstmt{\begin{array}{@{}l@{}} \forall r, m, n, a, u, v. \\ \quad (\operatorname{is\text{-}ring}(r) \wedge ((a \in \operatorname{mat}(m, n, \operatorname{carr}(r))) \wedge (\operatorname{is\text{-}invertible\text{-}mat}(r, m, u) \wedge \operatorname{is\text{-}invertible\text{-}mat}(r, n, v)))) \Rightarrow  \\ \quad \quad \operatorname{mat\text{-}equiv}(r, m, n, a, \operatorname{matmul}(r, \operatorname{matmul}(r, u, a), v)) \end{array}}}
\vnbprop{swap-to-corner}{\vnbstmt{\begin{array}{@{}l@{}} \forall r, m, n, a, c, d. \\ \quad \operatorname{is\text{-}ring}(r) \wedge \\ \quad a \in \operatorname{mat}(m, n, \operatorname{carr}(r)) \wedge \\ \quad 1 \in \operatorname{interval}(1, m) \wedge \\ \quad 1 \in \operatorname{interval}(1, n) \wedge \\ \quad c \in \operatorname{interval}(1, m) \wedge \\ \quad d \in \operatorname{interval}(1, n) \wedge \\ \quad \neg (1 = c) \wedge \\ \quad \neg (1 = d) \Rightarrow \\ \quad \quad \exists b.\; (\operatorname{mat\text{-}equiv}(r, m, n, a, b) \wedge (\operatorname{entry}(b, 1, 1) = \operatorname{entry}(a, c, d))) \end{array}}}
\vnbprop{swap-to-corner-gen}{\vnbstmt{\begin{array}{@{}l@{}} \forall r, m, n, a, c, d. \\ \quad \operatorname{is\text{-}ring}(r) \wedge \\ \quad a \in \operatorname{mat}(m, n, \operatorname{carr}(r)) \wedge \\ \quad 1 \in \operatorname{interval}(1, m) \wedge \\ \quad 1 \in \operatorname{interval}(1, n) \wedge \\ \quad c \in \operatorname{interval}(1, m) \wedge \\ \quad d \in \operatorname{interval}(1, n) \Rightarrow \\ \quad \quad \exists b.\; (\operatorname{mat\text{-}equiv}(r, m, n, a, b) \wedge (\operatorname{entry}(b, 1, 1) = \operatorname{entry}(a, c, d))) \end{array}}}
\vnbprop{place-min-pivot}{\vnbstmt{\begin{array}{@{}l@{}} \forall a, m, n.; \forall p \in \operatorname{mat}(m, n, \operatorname{carr}(a)). \\ \quad \operatorname{is\text{-}euclidean\text{-}ring}(a) \wedge \\ \quad 1 \in \operatorname{interval}(1, m) \wedge \\ \quad 1 \in \operatorname{interval}(1, n) \wedge \\ \quad \exists c \in \operatorname{interval}(1, m), d \in \operatorname{interval}(1, n). \\ \quad \quad \neg (\operatorname{entry}(p, c, d) = \operatorname{zero}(a)) \Rightarrow \\ \quad \quad \exists b. \\ \quad \quad \quad \operatorname{mat\text{-}equiv}(a, m, n, p, b) \wedge \\ \quad \quad \quad \neg (\operatorname{entry}(b, 1, 1) = \operatorname{zero}(a)) \wedge \\ \quad \quad \quad \forall i, j. \\ \quad \quad \quad \quad ((i \in \operatorname{interval}(1, m)) \wedge ((j \in \operatorname{interval}(1, n)) \wedge \neg (\operatorname{entry}(p, i, j) = \operatorname{zero}(a)))) \Rightarrow  \\ \quad \quad \quad \quad \quad \operatorname{gauge}(a)(\operatorname{entry}(b, 1, 1)) \leq \operatorname{gauge}(a)(\operatorname{entry}(p, i, j)) \end{array}}}

\vnbbookfile{theorem-library/class-min-pivot-proof}{theorem-library/class-min-pivot-proof}
\noindent{\small class-min-pivot-proof.scm -- the minimal-degree pivot over the whole \textasciitilde{}-equivalence class, and the typing lemma it needs.}\par
\vnbprop{mat-equiv-target-is-mat}{\vnbstmt{\begin{array}{@{}l@{}} \forall a, m, n, c, d. \\ \quad (\operatorname{is\text{-}ring}(a) \wedge ((c \in \operatorname{mat}(m, n, \operatorname{carr}(a))) \wedge \operatorname{mat\text{-}equiv}(a, m, n, c, d))) \Rightarrow  \\ \quad \quad d \in \operatorname{mat}(m, n, \operatorname{carr}(a)) \end{array}}}
\vnbprop{class-min-pivot}{\vnbstmt{\begin{array}{@{}l@{}} \forall a, m, n.; \forall p \in \operatorname{mat}(m, n, \operatorname{carr}(a)). \\ \quad \operatorname{is\text{-}euclidean\text{-}ring}(a) \wedge \\ \quad 1 \in \operatorname{interval}(1, m) \wedge \\ \quad 1 \in \operatorname{interval}(1, n) \wedge \\ \quad \exists d \in \operatorname{interval}(1, m), f \in \operatorname{interval}(1, n). \\ \quad \quad \neg (\operatorname{entry}(p, d, f) = \operatorname{zero}(a)) \Rightarrow \\ \quad \quad \exists b. \\ \quad \quad \quad \operatorname{mat\text{-}equiv}(a, m, n, p, b) \wedge \\ \quad \quad \quad \neg (\operatorname{entry}(b, 1, 1) = \operatorname{zero}(a)) \wedge \\ \quad \quad \quad \forall c, i, j. \\ \quad \quad \quad \quad \operatorname{mat\text{-}equiv}(a, m, n, p, c) \wedge \\ \quad \quad \quad \quad i \in \operatorname{interval}(1, m) \wedge \\ \quad \quad \quad \quad j \in \operatorname{interval}(1, n) \wedge \\ \quad \quad \quad \quad \neg (\operatorname{entry}(c, i, j) = \operatorname{zero}(a)) \Rightarrow \\ \quad \quad \quad \quad \quad \operatorname{gauge}(a)(\operatorname{entry}(b, 1, 1)) \leq \operatorname{gauge}(a)(\operatorname{entry}(c, i, j)) \end{array}}}

\vnbbookfile{theorem-library/smith-clear-proof}{theorem-library/smith-clear-proof}
\noindent{\small smith-clear-proof.scm -- Brick B3 (LA Phase B): the DESCENT step of the Smith reduction.  pivot-clears-col: over a euclidean ring, if P's pivot P\_\{1,1\} is nonzero and of MINIMAL degree over the whole equivalence class of P (class-min-pivot supplies such a P), then a single elementary COLUMN op zeroes}\par
\vnbprop{pivot-clears-col}{\vnbstmt{\begin{array}{@{}l@{}} \forall a, m, n, p, j. \\ \quad \operatorname{is\text{-}euclidean\text{-}ring}(a) \wedge \\ \quad p \in \operatorname{mat}(m, n, \operatorname{carr}(a)) \wedge \\ \quad 1 \in \operatorname{interval}(1, m) \wedge \\ \quad 1 \in \operatorname{interval}(1, n) \wedge \\ \quad j \in \operatorname{interval}(1, n) \wedge \\ \quad \neg (1 = j) \wedge \\ \quad \neg (\operatorname{entry}(p, 1, 1) = \operatorname{zero}(a)) \wedge \\ \quad \forall c, i, b. \\ \quad \quad \operatorname{mat\text{-}equiv}(a, m, n, p, c) \wedge \\ \quad \quad i \in \operatorname{interval}(1, m) \wedge \\ \quad \quad b \in \operatorname{interval}(1, n) \wedge \\ \quad \quad \neg (\operatorname{entry}(c, i, b) = \operatorname{zero}(a)) \Rightarrow \\ \quad \quad \quad \operatorname{gauge}(a)(\operatorname{entry}(p, 1, 1)) \leq \operatorname{gauge}(a)(\operatorname{entry}(c, i, b)) \Rightarrow \\ \quad \quad \exists q. \\ \quad \quad \quad \operatorname{mat\text{-}equiv}(a, m, n, p, q) \wedge \\ \quad \quad \quad \operatorname{entry}(q, 1, 1) = \operatorname{entry}(p, 1, 1) \wedge \\ \quad \quad \quad \operatorname{entry}(q, 1, j) = \operatorname{zero}(a) \wedge \\ \quad \quad \quad \forall i, c. \\ \quad \quad \quad \quad ((i \in \operatorname{interval}(1, m)) \wedge ((c \in \operatorname{interval}(1, n)) \wedge \neg (c = j))) \Rightarrow  \\ \quad \quad \quad \quad \quad \operatorname{entry}(q, i, c) = \operatorname{entry}(p, i, c) \end{array}}}

\vnbbookfile{theorem-library/clear-first-row-proof}{theorem-library/clear-first-row-proof}
\noindent{\small clear-first-row-proof.scm -- Brick B3 (LA Phase B): clear the whole first row. Iterates pivot-clears-col across columns 2..n by NN-induction on the number of cleared columns, then instantiates at n.  With a class-minimal pivot at (1,1), the result is a matrix \textasciitilde{} P, same pivot, whose first row is zero off the pivot.}\par
\vnbprop{classmin-transport}{\vnbstmt{\begin{array}{@{}l@{}} \forall a, m, n, p, q, c, i, j. \\ \quad \operatorname{is\text{-}ring}(a) \wedge \\ \quad p \in \operatorname{mat}(m, n, \operatorname{carr}(a)) \wedge \\ \quad \operatorname{mat\text{-}equiv}(a, m, n, p, q) \wedge \\ \quad \operatorname{entry}(q, 1, 1) = \operatorname{entry}(p, 1, 1) \wedge \\ \quad \forall c, i, j. \\ \quad \quad \operatorname{mat\text{-}equiv}(a, m, n, p, c) \wedge \\ \quad \quad i \in \operatorname{interval}(1, m) \wedge \\ \quad \quad j \in \operatorname{interval}(1, n) \wedge \\ \quad \quad \neg (\operatorname{entry}(c, i, j) = \operatorname{zero}(a)) \Rightarrow \\ \quad \quad \quad \operatorname{gauge}(a)(\operatorname{entry}(p, 1, 1)) \leq \operatorname{gauge}(a)(\operatorname{entry}(c, i, j)) \wedge \\ \quad \operatorname{mat\text{-}equiv}(a, m, n, q, c) \wedge \\ \quad i \in \operatorname{interval}(1, m) \wedge \\ \quad j \in \operatorname{interval}(1, n) \wedge \\ \quad \neg (\operatorname{entry}(c, i, j) = \operatorname{zero}(a)) \Rightarrow \\ \quad \quad \operatorname{gauge}(a)(\operatorname{entry}(q, 1, 1)) \leq \operatorname{gauge}(a)(\operatorname{entry}(c, i, j)) \end{array}}}
\vnbprop{clear-row-upto}{\vnbstmt{\begin{array}{@{}l@{}} \forall a, m, n.; \forall p \in \operatorname{mat}(m, n, \operatorname{carr}(a)), k \in \mathbb{N}. \\ \quad \operatorname{is\text{-}euclidean\text{-}ring}(a) \wedge \\ \quad 1 \in \operatorname{interval}(1, m) \wedge \\ \quad 1 \in \operatorname{interval}(1, n) \wedge \\ \quad \neg (\operatorname{entry}(p, 1, 1) = \operatorname{zero}(a)) \wedge \\ \quad \forall c, i, b. \\ \quad \quad \operatorname{mat\text{-}equiv}(a, m, n, p, c) \wedge \\ \quad \quad i \in \operatorname{interval}(1, m) \wedge \\ \quad \quad b \in \operatorname{interval}(1, n) \wedge \\ \quad \quad \neg (\operatorname{entry}(c, i, b) = \operatorname{zero}(a)) \Rightarrow \\ \quad \quad \quad \operatorname{gauge}(a)(\operatorname{entry}(p, 1, 1)) \leq \operatorname{gauge}(a)(\operatorname{entry}(c, i, b)) \Rightarrow \\ \quad \quad \exists q \in \operatorname{mat}(m, n, \operatorname{carr}(a)). \\ \quad \quad \quad \operatorname{mat\text{-}equiv}(a, m, n, p, q) \wedge \\ \quad \quad \quad \operatorname{entry}(q, 1, 1) = \operatorname{entry}(p, 1, 1) \wedge \\ \quad \quad \quad \forall j \in \operatorname{interval}(1, n). \\ \quad \quad \quad \quad (\neg (j = 1) \wedge (j \leq k)) \Rightarrow (\operatorname{entry}(q, 1, j) = \operatorname{zero}(a)) \end{array}}}
\vnbprop{clear-first-row}{\vnbstmt{\begin{array}{@{}l@{}} \forall a, m, n.; \forall p \in \operatorname{mat}(m, n, \operatorname{carr}(a)). \\ \quad \operatorname{is\text{-}euclidean\text{-}ring}(a) \wedge \\ \quad m \in \mathbb{N} \wedge \\ \quad n \in \mathbb{N} \wedge \\ \quad 1 \in \operatorname{interval}(1, m) \wedge \\ \quad 1 \in \operatorname{interval}(1, n) \wedge \\ \quad \neg (\operatorname{entry}(p, 1, 1) = \operatorname{zero}(a)) \wedge \\ \quad \forall c, i, b. \\ \quad \quad \operatorname{mat\text{-}equiv}(a, m, n, p, c) \wedge \\ \quad \quad i \in \operatorname{interval}(1, m) \wedge \\ \quad \quad b \in \operatorname{interval}(1, n) \wedge \\ \quad \quad \neg (\operatorname{entry}(c, i, b) = \operatorname{zero}(a)) \Rightarrow \\ \quad \quad \quad \operatorname{gauge}(a)(\operatorname{entry}(p, 1, 1)) \leq \operatorname{gauge}(a)(\operatorname{entry}(c, i, b)) \Rightarrow \\ \quad \quad \exists q. \\ \quad \quad \quad \operatorname{mat\text{-}equiv}(a, m, n, p, q) \wedge \\ \quad \quad \quad \operatorname{entry}(q, 1, 1) = \operatorname{entry}(p, 1, 1) \wedge \\ \quad \quad \quad \forall j \in \operatorname{interval}(1, n).\; (\neg (j = 1) \Rightarrow (\operatorname{entry}(q, 1, j) = \operatorname{zero}(a))) \end{array}}}

\vnbbookfile{theorem-library/clear-first-col-proof}{theorem-library/clear-first-col-proof}
\noindent{\small clear-first-col-proof.scm -- Brick B3 (LA Phase B): clear the whole first column. The ROW-op mirror of clear-first-row: left-multiply by elementary row matrices. pivot-clears-row -- one row op zeroes an off-pivot column-1 entry (i,1); the euclidean remainder is forced to vanish by class-minimality}\par
\vnbprop{pivot-clears-row}{\vnbstmt{\begin{array}{@{}l@{}} \forall a, m, n, p, i. \\ \quad \operatorname{is\text{-}euclidean\text{-}ring}(a) \wedge \\ \quad p \in \operatorname{mat}(m, n, \operatorname{carr}(a)) \wedge \\ \quad 1 \in \operatorname{interval}(1, m) \wedge \\ \quad 1 \in \operatorname{interval}(1, n) \wedge \\ \quad i \in \operatorname{interval}(1, m) \wedge \\ \quad \neg (i = 1) \wedge \\ \quad \neg (\operatorname{entry}(p, 1, 1) = \operatorname{zero}(a)) \wedge \\ \quad \forall c, b, j. \\ \quad \quad \operatorname{mat\text{-}equiv}(a, m, n, p, c) \wedge \\ \quad \quad b \in \operatorname{interval}(1, m) \wedge \\ \quad \quad j \in \operatorname{interval}(1, n) \wedge \\ \quad \quad \neg (\operatorname{entry}(c, b, j) = \operatorname{zero}(a)) \Rightarrow \\ \quad \quad \quad \operatorname{gauge}(a)(\operatorname{entry}(p, 1, 1)) \leq \operatorname{gauge}(a)(\operatorname{entry}(c, b, j)) \Rightarrow \\ \quad \quad \exists q. \\ \quad \quad \quad \operatorname{mat\text{-}equiv}(a, m, n, p, q) \wedge \\ \quad \quad \quad \operatorname{entry}(q, 1, 1) = \operatorname{entry}(p, 1, 1) \wedge \\ \quad \quad \quad \operatorname{entry}(q, i, 1) = \operatorname{zero}(a) \wedge \\ \quad \quad \quad \forall w, c. \\ \quad \quad \quad \quad ((w \in \operatorname{interval}(1, m)) \wedge ((c \in \operatorname{interval}(1, n)) \wedge \neg (w = i))) \Rightarrow  \\ \quad \quad \quad \quad \quad \operatorname{entry}(q, w, c) = \operatorname{entry}(p, w, c) \end{array}}}
\vnbprop{clear-col-upto}{\vnbstmt{\begin{array}{@{}l@{}} \forall a, m, n.; \forall p \in \operatorname{mat}(m, n, \operatorname{carr}(a)), k \in \mathbb{N}. \\ \quad \operatorname{is\text{-}euclidean\text{-}ring}(a) \wedge \\ \quad 1 \in \operatorname{interval}(1, m) \wedge \\ \quad 1 \in \operatorname{interval}(1, n) \wedge \\ \quad \neg (\operatorname{entry}(p, 1, 1) = \operatorname{zero}(a)) \wedge \\ \quad \forall c, b, j. \\ \quad \quad \operatorname{mat\text{-}equiv}(a, m, n, p, c) \wedge \\ \quad \quad b \in \operatorname{interval}(1, m) \wedge \\ \quad \quad j \in \operatorname{interval}(1, n) \wedge \\ \quad \quad \neg (\operatorname{entry}(c, b, j) = \operatorname{zero}(a)) \Rightarrow \\ \quad \quad \quad \operatorname{gauge}(a)(\operatorname{entry}(p, 1, 1)) \leq \operatorname{gauge}(a)(\operatorname{entry}(c, b, j)) \Rightarrow \\ \quad \quad \exists q \in \operatorname{mat}(m, n, \operatorname{carr}(a)). \\ \quad \quad \quad \operatorname{mat\text{-}equiv}(a, m, n, p, q) \wedge \\ \quad \quad \quad \operatorname{entry}(q, 1, 1) = \operatorname{entry}(p, 1, 1) \wedge \\ \quad \quad \quad \forall i \in \operatorname{interval}(1, m). \\ \quad \quad \quad \quad (\neg (i = 1) \wedge (i \leq k)) \Rightarrow (\operatorname{entry}(q, i, 1) = \operatorname{zero}(a)) \wedge \\ \quad \quad \quad \forall c \in \operatorname{interval}(1, n).\; (\operatorname{entry}(q, 1, c) = \operatorname{entry}(p, 1, c)) \end{array}}}
\vnbprop{clear-first-col}{\vnbstmt{\begin{array}{@{}l@{}} \forall a, m, n.; \forall p \in \operatorname{mat}(m, n, \operatorname{carr}(a)). \\ \quad \operatorname{is\text{-}euclidean\text{-}ring}(a) \wedge \\ \quad m \in \mathbb{N} \wedge \\ \quad n \in \mathbb{N} \wedge \\ \quad 1 \in \operatorname{interval}(1, m) \wedge \\ \quad 1 \in \operatorname{interval}(1, n) \wedge \\ \quad \neg (\operatorname{entry}(p, 1, 1) = \operatorname{zero}(a)) \wedge \\ \quad \forall c, b, j. \\ \quad \quad \operatorname{mat\text{-}equiv}(a, m, n, p, c) \wedge \\ \quad \quad b \in \operatorname{interval}(1, m) \wedge \\ \quad \quad j \in \operatorname{interval}(1, n) \wedge \\ \quad \quad \neg (\operatorname{entry}(c, b, j) = \operatorname{zero}(a)) \Rightarrow \\ \quad \quad \quad \operatorname{gauge}(a)(\operatorname{entry}(p, 1, 1)) \leq \operatorname{gauge}(a)(\operatorname{entry}(c, b, j)) \Rightarrow \\ \quad \quad \exists q. \\ \quad \quad \quad \operatorname{mat\text{-}equiv}(a, m, n, p, q) \wedge \\ \quad \quad \quad \operatorname{entry}(q, 1, 1) = \operatorname{entry}(p, 1, 1) \wedge \\ \quad \quad \quad \forall i \in \operatorname{interval}(1, m).\; (\neg (i = 1) \Rightarrow (\operatorname{entry}(q, i, 1) = \operatorname{zero}(a))) \wedge \\ \quad \quad \quad \forall c \in \operatorname{interval}(1, n).\; (\operatorname{entry}(q, 1, c) = \operatorname{entry}(p, 1, c)) \end{array}}}

\vnbbookfile{theorem-library/matact-summand-type-le-proof}{theorem-library/matact-summand-type-le-proof}
\noindent{\small matact-summand-type-le-proof.scm -- matact-summand-type-le, GUARDED and PROVEN.}\par
\vnbprop{matact-summand-type-le-guarded}{\vnbstmt{\begin{array}{@{}l@{}} \forall d, m, n, q, p, u, i, c, k. \\ \quad \operatorname{is\text{-}module}(d) \wedge \\ \quad p \in \operatorname{mat}(m, n, \operatorname{carr}(\operatorname{scal}(d))) \wedge \\ \quad u \in \operatorname{mat}(n, q, \operatorname{vec}(d)) \wedge \\ \quad i \in \operatorname{interval}(1, m) \wedge \\ \quad c \in \operatorname{interval}(1, q) \wedge \\ \quad k \in \mathbb{N} \wedge \\ \quad k \leq n \Rightarrow \\ \quad \quad (\lambda j \in \operatorname{interval}(1, k).\; \operatorname{act}(d)(\operatorname{entry}(p, i, j), \operatorname{entry}(u, j, c))) \\ \quad \quad \in\; (\operatorname{interval}(1, k) \to \operatorname{carr}(\operatorname{module\text{-}vector\text{-}ag}(d))) \end{array}}}

\vnbbookfile{theorem-library/border-mult-proof}{theorem-library/border-mult-proof}
\noindent{\small border-mult-proof.scm -- Brick B3 (LA Phase B): block-matrix multiplication. BORDER(A,b,X,p,q) . BORDER(A,d,Y,q,r) = BORDER(A, b*d, X.Y, p, r)   [trust:none] The [[b,0],[0,X]] . [[d,0],[0,Y]] = [[bd,0],[0,XY]] direct-sum law -- the crux that lets the Smith recursion border a block equivalence back to the full}\par
\vnbprop{border-mult}{\vnbstmt{\begin{array}{@{}l@{}} \forall a, p, q, r, x, y, b, d. \\ \quad \operatorname{is\text{-}ring}(a) \wedge \\ \quad p \in \mathbb{N} \wedge \\ \quad q \in \mathbb{N} \wedge \\ \quad r \in \mathbb{N} \wedge \\ \quad x \in \operatorname{mat}(p, q, \operatorname{carr}(a)) \wedge \\ \quad y \in \operatorname{mat}(q, r, \operatorname{carr}(a)) \wedge \\ \quad b \in \operatorname{carr}(a) \wedge \\ \quad d \in \operatorname{carr}(a) \wedge \\ \quad (q = 0) \Rightarrow ((p = 0) \vee (r = 0)) \Rightarrow \\ \quad \quad \operatorname{matmul}(a, \operatorname{border}(a, b, x, p, q), \operatorname{border}(a, d, y, q, r)) \\ \quad \quad =\; \operatorname{border}(a, \operatorname{mul}(a)(b, d), \operatorname{matmul}(a, x, y), p, r) \end{array}}}

\vnbbookfile{theorem-library/border-assembly-proof}{theorem-library/border-assembly-proof}
\vnbprop{border-identity}{\vnbstmt{\begin{array}{@{}l@{}} \forall a.; \forall p \in \mathbb{N}. \\ \quad \operatorname{is\text{-}ring}(a) \Rightarrow  \\ \quad \quad \operatorname{border}(a, \operatorname{one}(a), \operatorname{identmat}(a, p), p, p) = \operatorname{identmat}(a, (p + 1)) \end{array}}}
\vnbprop{border-is-diagonal}{\vnbstmt{\begin{array}{@{}l@{}} \forall a, b, e, p, q. \\ \quad \operatorname{is\text{-}ring}(a) \wedge \\ \quad p \in \mathbb{N} \wedge \\ \quad q \in \mathbb{N} \wedge \\ \quad \operatorname{is\text{-}diagonal}(a, p, q, e) \wedge \\ \quad b \in \operatorname{carr}(a) \wedge \\ \quad e \in \operatorname{mat}(p, q, \operatorname{carr}(a)) \Rightarrow \\ \quad \quad \operatorname{is\text{-}diagonal}(a, (p + 1), (q + 1), \operatorname{border}(a, b, e, p, q)) \end{array}}}
\vnbprop{border-invertible}{\vnbstmt{\begin{array}{@{}l@{}} \forall a, n, u. \\ \quad (\operatorname{is\text{-}ring}(a) \wedge ((n \in \mathbb{N}) \wedge \operatorname{is\text{-}invertible\text{-}mat}(a, n, u))) \Rightarrow  \\ \quad \quad \operatorname{is\text{-}invertible\text{-}mat}(a, (n + 1), \operatorname{border}(a, \operatorname{one}(a), u, n, n)) \end{array}}}
\vnbprop{bordering}{\vnbstmt{\begin{array}{@{}l@{}} \forall a, p, q, b, w, g, e. \\ \quad \operatorname{is\text{-}ring}(a) \wedge \\ \quad p \in \mathbb{N} \wedge \\ \quad q \in \mathbb{N} \wedge \\ \quad b \in \operatorname{carr}(a) \wedge \\ \quad w \in \operatorname{mat}(p, q, \operatorname{carr}(a)) \wedge \\ \quad g \in \operatorname{mat}((p + 1), (q + 1), \operatorname{carr}(a)) \wedge \\ \quad g = \operatorname{border}(a, b, w, p, q) \wedge \\ \quad \operatorname{mat\text{-}equiv}(a, p, q, w, e) \Rightarrow \\ \quad \quad \operatorname{mat\text{-}equiv}(a, (p + 1), (q + 1), g, \operatorname{border}(a, b, e, p, q)) \end{array}}}

\vnbbookfile{theorem-library/bordered-eq-border-proof}{theorem-library/bordered-eq-border-proof}
\noindent{\small bordered-eq-border -- a matrix C whose first row and first column are cleared off the pivot IS literally BORDER(C\_\{1,1\}, SUBMAT(C)): C = [[C11,0],[0,C']]. The bridge from "cleared cross" (clear-first-row + clear-first-col) to the BORDER block form, so the Smith recursion can border a sub-block equivalence}\par
\vnbprop{bordered-eq-border}{\vnbstmt{\begin{array}{@{}l@{}} \forall a, p, q, c. \\ \quad \operatorname{is\text{-}ring}(a) \wedge \\ \quad p \in \mathbb{N} \wedge \\ \quad q \in \mathbb{N} \wedge \\ \quad c \in \operatorname{mat}((p + 1), (q + 1), \operatorname{carr}(a)) \wedge \\ \quad \forall j \in \operatorname{interval}(1, (q + 1)).\; (\neg (j = 1) \Rightarrow (\operatorname{entry}(c, 1, j) = \operatorname{zero}(a))) \wedge \\ \quad \forall i \in \operatorname{interval}(1, (p + 1)).\; (\neg (i = 1) \Rightarrow (\operatorname{entry}(c, i, 1) = \operatorname{zero}(a))) \Rightarrow \\ \quad \quad c = \operatorname{border}(a, \operatorname{entry}(c, 1, 1), \operatorname{submat}(c, p, q), p, q) \end{array}}}

\vnbbookfile{theorem-library/det-alternating-form}{theorem-library/det-alternating-form}
\noindent{\small det-alternating-form.scm -- Hoffman-Kunze 5.3 Theorem 2 (batch 26-B, 2026-09-24): an ALTERNATING n-LINEAR function D on the n x n matrices over a commutative ring R satisfies  D(A) = det(A) D(I)  for every A.  Proven WITHOUT permutations.}\par
\vnbprop{daf-ring-rdist}{\vnbstmt{\begin{array}{@{}l@{}} \forall r, s, x, y, z. \\ \quad (\operatorname{is\text{-}commutative\text{-}ring}(r) \wedge ((s \in \operatorname{carr}(r)) \wedge ((x \in \operatorname{carr}(r)) \wedge ((y \in \operatorname{carr}(r)) \wedge (z \in \operatorname{carr}(r)))))) \Rightarrow  \\ \quad \quad \operatorname{mul}(r)(\operatorname{add}(r)(\operatorname{mul}(r)(s, x), y), z) = \operatorname{add}(r)(\operatorname{mul}(r)(s, \operatorname{mul}(r)(x, z)), \operatorname{mul}(r)(y, z)) \end{array}}}
\vnbprop{daf-ring-cancel0}{\vnbstmt{\begin{array}{@{}l@{}} \forall r, s, x, y. \\ \quad \operatorname{is\text{-}commutative\text{-}ring}(r) \wedge \\ \quad s \in \operatorname{carr}(r) \wedge \\ \quad x \in \operatorname{carr}(r) \wedge \\ \quad y \in \operatorname{carr}(r) \wedge \\ \quad x = \operatorname{add}(r)(\operatorname{mul}(r)(s, \operatorname{zero}(r)), y) \Rightarrow \\ \quad \quad y = x \end{array}}}
\vnbprop{daf-det-fun}{\vnbstmt{\begin{array}{@{}l@{}} \forall r, n. \\ \quad (\operatorname{is\text{-}commutative\text{-}ring}(r) \wedge (n \in \mathbb{N})) \Rightarrow  \\ \quad \quad (\lambda l \in \operatorname{mat}(n, n, \operatorname{carr}(r)).\; \operatorname{det}(r, n, l)) \in (\operatorname{mat}(n, n, \operatorname{carr}(r)) \to \operatorname{carr}(r)) \end{array}}}
\vnbprop{daf-det-lin}{\vnbstmt{\begin{array}{@{}l@{}} \forall r, n, a, b, c, p, s. \\ \quad \operatorname{is\text{-}commutative\text{-}ring}(r) \wedge \\ \quad n \in \mathbb{N} \wedge \\ \quad a \in \operatorname{mat}(n, n, \operatorname{carr}(r)) \wedge \\ \quad b \in \operatorname{mat}(n, n, \operatorname{carr}(r)) \wedge \\ \quad c \in \operatorname{mat}(n, n, \operatorname{carr}(r)) \wedge \\ \quad p \in \operatorname{interval}(1, n) \wedge \\ \quad s \in \operatorname{carr}(r) \wedge \\ \quad \forall i, k \in \operatorname{interval}(1, n). \\ \quad \quad \neg (i = p) \Rightarrow (\operatorname{entry}(b, i, k) = \operatorname{entry}(a, i, k)) \wedge \\ \quad \forall i, k \in \operatorname{interval}(1, n). \\ \quad \quad \neg (i = p) \Rightarrow (\operatorname{entry}(c, i, k) = \operatorname{entry}(a, i, k)) \wedge \\ \quad \forall k \in \operatorname{interval}(1, n). \\ \quad \quad \operatorname{entry}(c, p, k) = \operatorname{add}(r)(\operatorname{mul}(r)(s, \operatorname{entry}(a, p, k)), \operatorname{entry}(b, p, k)) \Rightarrow \\ \quad \quad (\lambda l \in \operatorname{mat}(n, n, \operatorname{carr}(r)).\; \operatorname{det}(r, n, l))(c) \\ \quad \quad =\; \operatorname{add}(r)(\operatorname{mul}(r)(s, (\lambda l \in \operatorname{mat}(n, n, \operatorname{carr}(r)).\; \operatorname{det}(r, n, l))(a)), \\ \quad \quad \quad (\lambda l \in \operatorname{mat}(n, n, \operatorname{carr}(r)).\; \operatorname{det}(r, n, l))(b)) \end{array}}}
\vnbprop{daf-det-alt}{\vnbstmt{\begin{array}{@{}l@{}} \forall r, n, a, p, q. \\ \quad \operatorname{is\text{-}commutative\text{-}ring}(r) \wedge \\ \quad n \in \mathbb{N} \wedge \\ \quad a \in \operatorname{mat}(n, n, \operatorname{carr}(r)) \wedge \\ \quad p \in \operatorname{interval}(1, n) \wedge \\ \quad q \in \operatorname{interval}(1, n) \wedge \\ \quad \neg (p = q) \wedge \\ \quad \forall k \in \operatorname{interval}(1, n).\; (\operatorname{entry}(a, p, k) = \operatorname{entry}(a, q, k)) \Rightarrow \\ \quad \quad (\lambda l \in \operatorname{mat}(n, n, \operatorname{carr}(r)).\; \operatorname{det}(r, n, l))(a) = \operatorname{zero}(r) \end{array}}}
\vnbprop{daf-mm-row-eq}{\vnbstmt{\begin{array}{@{}l@{}} \forall r, n, x, y, q, i, j, k. \\ \quad \operatorname{is\text{-}commutative\text{-}ring}(r) \wedge \\ \quad n \in \mathbb{N} \wedge \\ \quad x \in \operatorname{mat}(n, n, \operatorname{carr}(r)) \wedge \\ \quad y \in \operatorname{mat}(n, n, \operatorname{carr}(r)) \wedge \\ \quad q \in \operatorname{mat}(n, n, \operatorname{carr}(r)) \wedge \\ \quad i \in \operatorname{interval}(1, n) \wedge \\ \quad j \in \operatorname{interval}(1, n) \wedge \\ \quad k \in \operatorname{interval}(1, n) \wedge \\ \quad \forall \mathit{drk} \in \operatorname{interval}(1, n).\; (\operatorname{entry}(x, i, \mathit{drk}) = \operatorname{entry}(y, j, \mathit{drk})) \Rightarrow \\ \quad \quad \operatorname{entry}(\operatorname{matmul}(r, x, q), i, k) = \operatorname{entry}(\operatorname{matmul}(r, y, q), j, k) \end{array}}}
\vnbprop{daf-mm-row-lin}{\vnbstmt{\begin{array}{@{}l@{}} \forall r, n, a, b, c, q, p, s, k. \\ \quad \operatorname{is\text{-}commutative\text{-}ring}(r) \wedge \\ \quad n \in \mathbb{N} \wedge \\ \quad a \in \operatorname{mat}(n, n, \operatorname{carr}(r)) \wedge \\ \quad b \in \operatorname{mat}(n, n, \operatorname{carr}(r)) \wedge \\ \quad c \in \operatorname{mat}(n, n, \operatorname{carr}(r)) \wedge \\ \quad q \in \operatorname{mat}(n, n, \operatorname{carr}(r)) \wedge \\ \quad p \in \operatorname{interval}(1, n) \wedge \\ \quad s \in \operatorname{carr}(r) \wedge \\ \quad k \in \operatorname{interval}(1, n) \wedge \\ \quad \forall \mathit{drk} \in \operatorname{interval}(1, n). \\ \quad \quad \operatorname{entry}(c, p, \mathit{drk}) = \operatorname{add}(r)(\operatorname{mul}(r)(s, \operatorname{entry}(a, p, \mathit{drk})), \operatorname{entry}(b, p, \mathit{drk})) \Rightarrow \\ \quad \quad \operatorname{entry}(\operatorname{matmul}(r, c, q), p, k) \\ \quad \quad =\; \operatorname{add}(r)(\operatorname{mul}(r)(s, \operatorname{entry}(\operatorname{matmul}(r, a, q), p, k)), \operatorname{entry}(\operatorname{matmul}(r, b, q), p, k)) \end{array}}}
\vnbprop{daf-rm-fun}{\vnbstmt{\begin{array}{@{}l@{}} \forall r, n, d, q. \\ \quad \operatorname{is\text{-}commutative\text{-}ring}(r) \wedge \\ \quad n \in \mathbb{N} \wedge \\ \quad d \in (\operatorname{mat}(n, n, \operatorname{carr}(r)) \to \operatorname{carr}(r)) \wedge \\ \quad q \in \operatorname{mat}(n, n, \operatorname{carr}(r)) \Rightarrow \\ \quad \quad (\lambda l \in \operatorname{mat}(n, n, \operatorname{carr}(r)).\; d(\operatorname{matmul}(r, l, q))) \in (\operatorname{mat}(n, n, \operatorname{carr}(r)) \to \operatorname{carr}(r)) \end{array}}}
\vnbprop{daf-rm-lin}{\vnbstmt{\begin{array}{@{}l@{}} \forall r, n, d, q, a, b, c, p, s. \\ \quad \operatorname{is\text{-}commutative\text{-}ring}(r) \wedge \\ \quad n \in \mathbb{N} \wedge \\ \quad d \in (\operatorname{mat}(n, n, \operatorname{carr}(r)) \to \operatorname{carr}(r)) \wedge \\ \quad \forall a, b, c, p, s. \\ \quad \quad a \in \operatorname{mat}(n, n, \operatorname{carr}(r)) \wedge \\ \quad \quad b \in \operatorname{mat}(n, n, \operatorname{carr}(r)) \wedge \\ \quad \quad c \in \operatorname{mat}(n, n, \operatorname{carr}(r)) \wedge \\ \quad \quad p \in \operatorname{interval}(1, n) \wedge \\ \quad \quad s \in \operatorname{carr}(r) \wedge \\ \quad \quad \forall i, k \in \operatorname{interval}(1, n). \\ \quad \quad \quad \neg (i = p) \Rightarrow (\operatorname{entry}(b, i, k) = \operatorname{entry}(a, i, k)) \wedge \\ \quad \quad \forall i, k \in \operatorname{interval}(1, n). \\ \quad \quad \quad \neg (i = p) \Rightarrow (\operatorname{entry}(c, i, k) = \operatorname{entry}(a, i, k)) \wedge \\ \quad \quad \forall k \in \operatorname{interval}(1, n). \\ \quad \quad \quad \operatorname{entry}(c, p, k) = \operatorname{add}(r)(\operatorname{mul}(r)(s, \operatorname{entry}(a, p, k)), \operatorname{entry}(b, p, k)) \Rightarrow \\ \quad \quad \quad d(c) = \operatorname{add}(r)(\operatorname{mul}(r)(s, d(a)), d(b)) \wedge \\ \quad q \in \operatorname{mat}(n, n, \operatorname{carr}(r)) \wedge \\ \quad a \in \operatorname{mat}(n, n, \operatorname{carr}(r)) \wedge \\ \quad b \in \operatorname{mat}(n, n, \operatorname{carr}(r)) \wedge \\ \quad c \in \operatorname{mat}(n, n, \operatorname{carr}(r)) \wedge \\ \quad p \in \operatorname{interval}(1, n) \wedge \\ \quad s \in \operatorname{carr}(r) \wedge \\ \quad \forall i, k \in \operatorname{interval}(1, n). \\ \quad \quad \neg (i = p) \Rightarrow (\operatorname{entry}(b, i, k) = \operatorname{entry}(a, i, k)) \wedge \\ \quad \forall i, k \in \operatorname{interval}(1, n). \\ \quad \quad \neg (i = p) \Rightarrow (\operatorname{entry}(c, i, k) = \operatorname{entry}(a, i, k)) \wedge \\ \quad \forall k \in \operatorname{interval}(1, n). \\ \quad \quad \operatorname{entry}(c, p, k) = \operatorname{add}(r)(\operatorname{mul}(r)(s, \operatorname{entry}(a, p, k)), \operatorname{entry}(b, p, k)) \Rightarrow \\ \quad \quad (\lambda l \in \operatorname{mat}(n, n, \operatorname{carr}(r)).\; d(\operatorname{matmul}(r, l, q)))(c) \\ \quad \quad =\; \operatorname{add}(r)(\operatorname{mul}(r)(s, (\lambda l \in \operatorname{mat}(n, n, \operatorname{carr}(r)).\; d(\operatorname{matmul}(r, l, q)))(a)), \\ \quad \quad \quad (\lambda l \in \operatorname{mat}(n, n, \operatorname{carr}(r)).\; d(\operatorname{matmul}(r, l, q)))(b)) \end{array}}}
\vnbprop{daf-rm-alt}{\vnbstmt{\begin{array}{@{}l@{}} \forall r, n, d, q, a, p, \mathit{dfq}. \\ \quad \operatorname{is\text{-}commutative\text{-}ring}(r) \wedge \\ \quad n \in \mathbb{N} \wedge \\ \quad d \in (\operatorname{mat}(n, n, \operatorname{carr}(r)) \to \operatorname{carr}(r)) \wedge \\ \quad \forall a, p, \mathit{dfq}. \\ \quad \quad a \in \operatorname{mat}(n, n, \operatorname{carr}(r)) \wedge \\ \quad \quad p \in \operatorname{interval}(1, n) \wedge \\ \quad \quad \mathit{dfq} \in \operatorname{interval}(1, n) \wedge \\ \quad \quad \neg (p = \mathit{dfq}) \wedge \\ \quad \quad \forall k \in \operatorname{interval}(1, n).\; (\operatorname{entry}(a, p, k) = \operatorname{entry}(a, \mathit{dfq}, k)) \Rightarrow \\ \quad \quad \quad d(a) = \operatorname{zero}(r) \wedge \\ \quad q \in \operatorname{mat}(n, n, \operatorname{carr}(r)) \wedge \\ \quad a \in \operatorname{mat}(n, n, \operatorname{carr}(r)) \wedge \\ \quad p \in \operatorname{interval}(1, n) \wedge \\ \quad \mathit{dfq} \in \operatorname{interval}(1, n) \wedge \\ \quad \neg (p = \mathit{dfq}) \wedge \\ \quad \forall k \in \operatorname{interval}(1, n).\; (\operatorname{entry}(a, p, k) = \operatorname{entry}(a, \mathit{dfq}, k)) \Rightarrow \\ \quad \quad (\lambda l \in \operatorname{mat}(n, n, \operatorname{carr}(r)).\; d(\operatorname{matmul}(r, l, q)))(a) = \operatorname{zero}(r) \end{array}}}
\vnbprop{daf-det-rm-fun}{\vnbstmt{\begin{array}{@{}l@{}} \forall r, n, q. \\ \quad (\operatorname{is\text{-}commutative\text{-}ring}(r) \wedge ((n \in \mathbb{N}) \wedge (q \in \operatorname{mat}(n, n, \operatorname{carr}(r))))) \Rightarrow  \\ \quad \quad (\lambda l \in \operatorname{mat}(n, n, \operatorname{carr}(r)).\; \operatorname{det}(r, n, \operatorname{matmul}(r, l, q))) \in (\operatorname{mat}(n, n, \operatorname{carr}(r)) \to \operatorname{carr}(r)) \end{array}}}
\vnbprop{daf-det-rm-lin}{\vnbstmt{\begin{array}{@{}l@{}} \forall r, n, q, a, b, c, p, s. \\ \quad \operatorname{is\text{-}commutative\text{-}ring}(r) \wedge \\ \quad n \in \mathbb{N} \wedge \\ \quad q \in \operatorname{mat}(n, n, \operatorname{carr}(r)) \wedge \\ \quad a \in \operatorname{mat}(n, n, \operatorname{carr}(r)) \wedge \\ \quad b \in \operatorname{mat}(n, n, \operatorname{carr}(r)) \wedge \\ \quad c \in \operatorname{mat}(n, n, \operatorname{carr}(r)) \wedge \\ \quad p \in \operatorname{interval}(1, n) \wedge \\ \quad s \in \operatorname{carr}(r) \wedge \\ \quad \forall i, k \in \operatorname{interval}(1, n). \\ \quad \quad \neg (i = p) \Rightarrow (\operatorname{entry}(b, i, k) = \operatorname{entry}(a, i, k)) \wedge \\ \quad \forall i, k \in \operatorname{interval}(1, n). \\ \quad \quad \neg (i = p) \Rightarrow (\operatorname{entry}(c, i, k) = \operatorname{entry}(a, i, k)) \wedge \\ \quad \forall k \in \operatorname{interval}(1, n). \\ \quad \quad \operatorname{entry}(c, p, k) = \operatorname{add}(r)(\operatorname{mul}(r)(s, \operatorname{entry}(a, p, k)), \operatorname{entry}(b, p, k)) \Rightarrow \\ \quad \quad (\lambda l \in \operatorname{mat}(n, n, \operatorname{carr}(r)).\; \operatorname{det}(r, n, \operatorname{matmul}(r, l, q)))(c) \\ \quad \quad =\; \operatorname{add}(r)(\operatorname{mul}(r)(s, (\lambda l \in \operatorname{mat}(n, n, \operatorname{carr}(r)).\; \operatorname{det}(r, n, \operatorname{matmul}(r, l, q)))(a)), \\ \quad \quad \quad (\lambda l \in \operatorname{mat}(n, n, \operatorname{carr}(r)).\; \operatorname{det}(r, n, \operatorname{matmul}(r, l, q)))(b)) \end{array}}}
\vnbprop{daf-det-rm-alt}{\vnbstmt{\begin{array}{@{}l@{}} \forall r, n, q, a, p, \mathit{dfq}. \\ \quad \operatorname{is\text{-}commutative\text{-}ring}(r) \wedge \\ \quad n \in \mathbb{N} \wedge \\ \quad q \in \operatorname{mat}(n, n, \operatorname{carr}(r)) \wedge \\ \quad a \in \operatorname{mat}(n, n, \operatorname{carr}(r)) \wedge \\ \quad p \in \operatorname{interval}(1, n) \wedge \\ \quad \mathit{dfq} \in \operatorname{interval}(1, n) \wedge \\ \quad \neg (p = \mathit{dfq}) \wedge \\ \quad \forall k \in \operatorname{interval}(1, n).\; (\operatorname{entry}(a, p, k) = \operatorname{entry}(a, \mathit{dfq}, k)) \Rightarrow \\ \quad \quad (\lambda l \in \operatorname{mat}(n, n, \operatorname{carr}(r)).\; \operatorname{det}(r, n, \operatorname{matmul}(r, l, q)))(a) = \operatorname{zero}(r) \end{array}}}
\vnbprop{daf-zero-row}{\vnbstmt{\begin{array}{@{}l@{}} \forall r, n, d, a, k. \\ \quad \operatorname{is\text{-}commutative\text{-}ring}(r) \wedge \\ \quad n \in \mathbb{N} \wedge \\ \quad d \in (\operatorname{mat}(n, n, \operatorname{carr}(r)) \to \operatorname{carr}(r)) \wedge \\ \quad \forall \mathit{dfa}, b, c, p, s. \\ \quad \quad \mathit{dfa} \in \operatorname{mat}(n, n, \operatorname{carr}(r)) \wedge \\ \quad \quad b \in \operatorname{mat}(n, n, \operatorname{carr}(r)) \wedge \\ \quad \quad c \in \operatorname{mat}(n, n, \operatorname{carr}(r)) \wedge \\ \quad \quad p \in \operatorname{interval}(1, n) \wedge \\ \quad \quad s \in \operatorname{carr}(r) \wedge \\ \quad \quad \forall i, \mathit{drk} \in \operatorname{interval}(1, n). \\ \quad \quad \quad \neg (i = p) \Rightarrow (\operatorname{entry}(b, i, \mathit{drk}) = \operatorname{entry}(\mathit{dfa}, i, \mathit{drk})) \wedge \\ \quad \quad \forall i, \mathit{drk} \in \operatorname{interval}(1, n). \\ \quad \quad \quad \neg (i = p) \Rightarrow (\operatorname{entry}(c, i, \mathit{drk}) = \operatorname{entry}(\mathit{dfa}, i, \mathit{drk})) \wedge \\ \quad \quad \forall \mathit{drk} \in \operatorname{interval}(1, n). \\ \quad \quad \quad \operatorname{entry}(c, p, \mathit{drk}) = \operatorname{add}(r)(\operatorname{mul}(r)(s, \operatorname{entry}(\mathit{dfa}, p, \mathit{drk})), \operatorname{entry}(b, p, \mathit{drk})) \Rightarrow \\ \quad \quad \quad d(c) = \operatorname{add}(r)(\operatorname{mul}(r)(s, d(\mathit{dfa})), d(b)) \wedge \\ \quad a \in \operatorname{mat}(n, n, \operatorname{carr}(r)) \wedge \\ \quad k \in \operatorname{interval}(1, n) \wedge \\ \quad \forall \mathit{drk} \in \operatorname{interval}(1, n).\; (\operatorname{entry}(a, k, \mathit{drk}) = \operatorname{zero}(r)) \Rightarrow \\ \quad \quad d(a) = \operatorname{zero}(r) \end{array}}}
\vnbprop{daf-bd-fun}{\vnbstmt{\begin{array}{@{}l@{}} \forall r, n, d. \\ \quad (\operatorname{is\text{-}commutative\text{-}ring}(r) \wedge ((n \in \mathbb{N}) \wedge (d \in (\operatorname{mat}((n + 1), (n + 1), \operatorname{carr}(r)) \to \operatorname{carr}(r))))) \Rightarrow  \\ \quad \quad (\lambda l \in \operatorname{mat}(n, n, \operatorname{carr}(r)).\; d(\operatorname{border}(r, \operatorname{one}(r), l, n, n))) \in (\operatorname{mat}(n, n, \operatorname{carr}(r)) \to \operatorname{carr}(r)) \end{array}}}
\vnbprop{daf-bd-lin}{\vnbstmt{\begin{array}{@{}l@{}} \forall r, n, d, a, b, c, p, s. \\ \quad \operatorname{is\text{-}commutative\text{-}ring}(r) \wedge \\ \quad n \in \mathbb{N} \wedge \\ \quad d \in (\operatorname{mat}((n + 1), (n + 1), \operatorname{carr}(r)) \to \operatorname{carr}(r)) \wedge \\ \quad \forall a, b, c, p, s. \\ \quad \quad a \in \operatorname{mat}((n + 1), (n + 1), \operatorname{carr}(r)) \wedge \\ \quad \quad b \in \operatorname{mat}((n + 1), (n + 1), \operatorname{carr}(r)) \wedge \\ \quad \quad c \in \operatorname{mat}((n + 1), (n + 1), \operatorname{carr}(r)) \wedge \\ \quad \quad p \in \operatorname{interval}(1, (n + 1)) \wedge \\ \quad \quad s \in \operatorname{carr}(r) \wedge \\ \quad \quad \forall i, k \in \operatorname{interval}(1, (n + 1)). \\ \quad \quad \quad \neg (i = p) \Rightarrow (\operatorname{entry}(b, i, k) = \operatorname{entry}(a, i, k)) \wedge \\ \quad \quad \forall i, k \in \operatorname{interval}(1, (n + 1)). \\ \quad \quad \quad \neg (i = p) \Rightarrow (\operatorname{entry}(c, i, k) = \operatorname{entry}(a, i, k)) \wedge \\ \quad \quad \forall k \in \operatorname{interval}(1, (n + 1)). \\ \quad \quad \quad \operatorname{entry}(c, p, k) = \operatorname{add}(r)(\operatorname{mul}(r)(s, \operatorname{entry}(a, p, k)), \operatorname{entry}(b, p, k)) \Rightarrow \\ \quad \quad \quad d(c) = \operatorname{add}(r)(\operatorname{mul}(r)(s, d(a)), d(b)) \wedge \\ \quad a \in \operatorname{mat}(n, n, \operatorname{carr}(r)) \wedge \\ \quad b \in \operatorname{mat}(n, n, \operatorname{carr}(r)) \wedge \\ \quad c \in \operatorname{mat}(n, n, \operatorname{carr}(r)) \wedge \\ \quad p \in \operatorname{interval}(1, n) \wedge \\ \quad s \in \operatorname{carr}(r) \wedge \\ \quad \forall i, k \in \operatorname{interval}(1, n). \\ \quad \quad \neg (i = p) \Rightarrow (\operatorname{entry}(b, i, k) = \operatorname{entry}(a, i, k)) \wedge \\ \quad \forall i, k \in \operatorname{interval}(1, n). \\ \quad \quad \neg (i = p) \Rightarrow (\operatorname{entry}(c, i, k) = \operatorname{entry}(a, i, k)) \wedge \\ \quad \forall k \in \operatorname{interval}(1, n). \\ \quad \quad \operatorname{entry}(c, p, k) = \operatorname{add}(r)(\operatorname{mul}(r)(s, \operatorname{entry}(a, p, k)), \operatorname{entry}(b, p, k)) \Rightarrow \\ \quad \quad (\lambda l \in \operatorname{mat}(n, n, \operatorname{carr}(r)).\; d(\operatorname{border}(r, \operatorname{one}(r), l, n, n)))(c) \\ \quad \quad =\; \operatorname{add}(r)(\operatorname{mul}(r)(s, (\lambda l \in \operatorname{mat}(n, n, \operatorname{carr}(r)).\; d(\operatorname{border}(r, \operatorname{one}(r), l, n, n)))(a)), \\ \quad \quad \quad (\lambda l \in \operatorname{mat}(n, n, \operatorname{carr}(r)).\; d(\operatorname{border}(r, \operatorname{one}(r), l, n, n)))(b)) \end{array}}}
\vnbprop{daf-bd-alt}{\vnbstmt{\begin{array}{@{}l@{}} \forall r, n, d, a, p, q. \\ \quad \operatorname{is\text{-}commutative\text{-}ring}(r) \wedge \\ \quad n \in \mathbb{N} \wedge \\ \quad d \in (\operatorname{mat}((n + 1), (n + 1), \operatorname{carr}(r)) \to \operatorname{carr}(r)) \wedge \\ \quad \forall a, p, q. \\ \quad \quad a \in \operatorname{mat}((n + 1), (n + 1), \operatorname{carr}(r)) \wedge \\ \quad \quad p \in \operatorname{interval}(1, (n + 1)) \wedge \\ \quad \quad q \in \operatorname{interval}(1, (n + 1)) \wedge \\ \quad \quad \neg (p = q) \wedge \\ \quad \quad \forall k \in \operatorname{interval}(1, (n + 1)).\; (\operatorname{entry}(a, p, k) = \operatorname{entry}(a, q, k)) \Rightarrow \\ \quad \quad \quad d(a) = \operatorname{zero}(r) \wedge \\ \quad a \in \operatorname{mat}(n, n, \operatorname{carr}(r)) \wedge \\ \quad p \in \operatorname{interval}(1, n) \wedge \\ \quad q \in \operatorname{interval}(1, n) \wedge \\ \quad \neg (p = q) \wedge \\ \quad \forall k \in \operatorname{interval}(1, n).\; (\operatorname{entry}(a, p, k) = \operatorname{entry}(a, q, k)) \Rightarrow \\ \quad \quad (\lambda l \in \operatorname{mat}(n, n, \operatorname{carr}(r)).\; d(\operatorname{border}(r, \operatorname{one}(r), l, n, n)))(a) = \operatorname{zero}(r) \end{array}}}
\vnbprop{daf-swap}{\vnbstmt{\begin{array}{@{}l@{}} \forall r, n, d, a, b, p, q. \\ \quad \operatorname{is\text{-}commutative\text{-}ring}(r) \wedge \\ \quad n \in \mathbb{N} \wedge \\ \quad d \in (\operatorname{mat}(n, n, \operatorname{carr}(r)) \to \operatorname{carr}(r)) \wedge \\ \quad \forall \mathit{dfa}, \mathit{dfb}, c, \mathit{dfp}, s. \\ \quad \quad \mathit{dfa} \in \operatorname{mat}(n, n, \operatorname{carr}(r)) \wedge \\ \quad \quad \mathit{dfb} \in \operatorname{mat}(n, n, \operatorname{carr}(r)) \wedge \\ \quad \quad c \in \operatorname{mat}(n, n, \operatorname{carr}(r)) \wedge \\ \quad \quad \mathit{dfp} \in \operatorname{interval}(1, n) \wedge \\ \quad \quad s \in \operatorname{carr}(r) \wedge \\ \quad \quad \forall i, k \in \operatorname{interval}(1, n). \\ \quad \quad \quad \neg (i = \mathit{dfp}) \Rightarrow (\operatorname{entry}(\mathit{dfb}, i, k) = \operatorname{entry}(\mathit{dfa}, i, k)) \wedge \\ \quad \quad \forall i, k \in \operatorname{interval}(1, n). \\ \quad \quad \quad \neg (i = \mathit{dfp}) \Rightarrow (\operatorname{entry}(c, i, k) = \operatorname{entry}(\mathit{dfa}, i, k)) \wedge \\ \quad \quad \forall k \in \operatorname{interval}(1, n). \\ \quad \quad \quad \operatorname{entry}(c, \mathit{dfp}, k) = \operatorname{add}(r)(\operatorname{mul}(r)(s, \operatorname{entry}(\mathit{dfa}, \mathit{dfp}, k)), \operatorname{entry}(\mathit{dfb}, \mathit{dfp}, k)) \Rightarrow \\ \quad \quad \quad d(c) = \operatorname{add}(r)(\operatorname{mul}(r)(s, d(\mathit{dfa})), d(\mathit{dfb})) \wedge \\ \quad \forall \mathit{dfa}, \mathit{dfp}, \mathit{dfq}. \\ \quad \quad \mathit{dfa} \in \operatorname{mat}(n, n, \operatorname{carr}(r)) \wedge \\ \quad \quad \mathit{dfp} \in \operatorname{interval}(1, n) \wedge \\ \quad \quad \mathit{dfq} \in \operatorname{interval}(1, n) \wedge \\ \quad \quad \neg (\mathit{dfp} = \mathit{dfq}) \wedge \\ \quad \quad \forall k \in \operatorname{interval}(1, n).\; (\operatorname{entry}(\mathit{dfa}, \mathit{dfp}, k) = \operatorname{entry}(\mathit{dfa}, \mathit{dfq}, k)) \Rightarrow \\ \quad \quad \quad d(\mathit{dfa}) = \operatorname{zero}(r) \wedge \\ \quad a \in \operatorname{mat}(n, n, \operatorname{carr}(r)) \wedge \\ \quad b \in \operatorname{mat}(n, n, \operatorname{carr}(r)) \wedge \\ \quad p \in \operatorname{interval}(1, n) \wedge \\ \quad q \in \operatorname{interval}(1, n) \wedge \\ \quad \neg (p = q) \wedge \\ \quad \forall i, k \in \operatorname{interval}(1, n). \\ \quad \quad (\neg (i = p) \wedge \neg (i = q)) \Rightarrow (\operatorname{entry}(b, i, k) = \operatorname{entry}(a, i, k)) \wedge \\ \quad \forall k \in \operatorname{interval}(1, n).\; (\operatorname{entry}(b, p, k) = \operatorname{entry}(a, q, k)) \wedge \\ \quad \forall k \in \operatorname{interval}(1, n).\; (\operatorname{entry}(b, q, k) = \operatorname{entry}(a, p, k)) \Rightarrow \\ \quad \quad d(b) = \operatorname{neg}(r)(d(a)) \end{array}}}
\vnbprop{daf-elem-f}{\vnbstmt{\begin{array}{@{}l@{}} \forall r, n, d, k, l. \\ \quad \operatorname{is\text{-}commutative\text{-}ring}(r) \wedge \\ \quad n \in \mathbb{N} \wedge \\ \quad d \in (\operatorname{mat}(n, n, \operatorname{carr}(r)) \to \operatorname{carr}(r)) \wedge \\ \quad \forall a, b, c, p, s. \\ \quad \quad a \in \operatorname{mat}(n, n, \operatorname{carr}(r)) \wedge \\ \quad \quad b \in \operatorname{mat}(n, n, \operatorname{carr}(r)) \wedge \\ \quad \quad c \in \operatorname{mat}(n, n, \operatorname{carr}(r)) \wedge \\ \quad \quad p \in \operatorname{interval}(1, n) \wedge \\ \quad \quad s \in \operatorname{carr}(r) \wedge \\ \quad \quad \forall i, \mathit{drk} \in \operatorname{interval}(1, n). \\ \quad \quad \quad \neg (i = p) \Rightarrow (\operatorname{entry}(b, i, \mathit{drk}) = \operatorname{entry}(a, i, \mathit{drk})) \wedge \\ \quad \quad \forall i, \mathit{drk} \in \operatorname{interval}(1, n). \\ \quad \quad \quad \neg (i = p) \Rightarrow (\operatorname{entry}(c, i, \mathit{drk}) = \operatorname{entry}(a, i, \mathit{drk})) \wedge \\ \quad \quad \forall \mathit{drk} \in \operatorname{interval}(1, n). \\ \quad \quad \quad \operatorname{entry}(c, p, \mathit{drk}) = \operatorname{add}(r)(\operatorname{mul}(r)(s, \operatorname{entry}(a, p, \mathit{drk})), \operatorname{entry}(b, p, \mathit{drk})) \Rightarrow \\ \quad \quad \quad d(c) = \operatorname{add}(r)(\operatorname{mul}(r)(s, d(a)), d(b)) \wedge \\ \quad \forall a, p, q. \\ \quad \quad a \in \operatorname{mat}(n, n, \operatorname{carr}(r)) \wedge \\ \quad \quad p \in \operatorname{interval}(1, n) \wedge \\ \quad \quad q \in \operatorname{interval}(1, n) \wedge \\ \quad \quad \neg (p = q) \wedge \\ \quad \quad \forall \mathit{drk} \in \operatorname{interval}(1, n).\; (\operatorname{entry}(a, p, \mathit{drk}) = \operatorname{entry}(a, q, \mathit{drk})) \Rightarrow \\ \quad \quad \quad d(a) = \operatorname{zero}(r) \wedge \\ \quad k \in \operatorname{interval}(1, n) \wedge \\ \quad l \in \operatorname{interval}(1, n) \wedge \\ \quad \neg (k = l) \Rightarrow \\ \quad \quad d(\operatorname{elem\text{-}f}(r, n, k, l)) = \operatorname{neg}(r)(d(\operatorname{identmat}(r, n))) \end{array}}}
\vnbprop{daf-clear-ind}{\vnbstmt{\begin{array}{@{}l@{}} \forall m \in \mathbb{N}, r, n, d, y. \\ \quad \operatorname{is\text{-}commutative\text{-}ring}(r) \wedge \\ \quad n \in \mathbb{N} \wedge \\ \quad d \in (\operatorname{mat}((n + 1), (n + 1), \operatorname{carr}(r)) \to \operatorname{carr}(r)) \wedge \\ \quad \forall a, b, c, p, s. \\ \quad \quad a \in \operatorname{mat}((n + 1), (n + 1), \operatorname{carr}(r)) \wedge \\ \quad \quad b \in \operatorname{mat}((n + 1), (n + 1), \operatorname{carr}(r)) \wedge \\ \quad \quad c \in \operatorname{mat}((n + 1), (n + 1), \operatorname{carr}(r)) \wedge \\ \quad \quad p \in \operatorname{interval}(1, (n + 1)) \wedge \\ \quad \quad s \in \operatorname{carr}(r) \wedge \\ \quad \quad \forall i, k \in \operatorname{interval}(1, (n + 1)). \\ \quad \quad \quad \neg (i = p) \Rightarrow (\operatorname{entry}(b, i, k) = \operatorname{entry}(a, i, k)) \wedge \\ \quad \quad \forall i, k \in \operatorname{interval}(1, (n + 1)). \\ \quad \quad \quad \neg (i = p) \Rightarrow (\operatorname{entry}(c, i, k) = \operatorname{entry}(a, i, k)) \wedge \\ \quad \quad \forall k \in \operatorname{interval}(1, (n + 1)). \\ \quad \quad \quad \operatorname{entry}(c, p, k) = \operatorname{add}(r)(\operatorname{mul}(r)(s, \operatorname{entry}(a, p, k)), \operatorname{entry}(b, p, k)) \Rightarrow \\ \quad \quad \quad d(c) = \operatorname{add}(r)(\operatorname{mul}(r)(s, d(a)), d(b)) \wedge \\ \quad \forall a, p, q. \\ \quad \quad a \in \operatorname{mat}((n + 1), (n + 1), \operatorname{carr}(r)) \wedge \\ \quad \quad p \in \operatorname{interval}(1, (n + 1)) \wedge \\ \quad \quad q \in \operatorname{interval}(1, (n + 1)) \wedge \\ \quad \quad \neg (p = q) \wedge \\ \quad \quad \forall k \in \operatorname{interval}(1, (n + 1)).\; (\operatorname{entry}(a, p, k) = \operatorname{entry}(a, q, k)) \Rightarrow \\ \quad \quad \quad d(a) = \operatorname{zero}(r) \wedge \\ \quad y \in \operatorname{mat}((n + 1), (n + 1), \operatorname{carr}(r)) \wedge \\ \quad \operatorname{entry}(y, 1, 1) = \operatorname{one}(r) \wedge \\ \quad \forall k \in \operatorname{interval}(1, (n + 1)).\; (\neg (k = 1) \Rightarrow (\operatorname{entry}(y, 1, k) = \operatorname{zero}(r))) \wedge \\ \quad m \leq n \Rightarrow \\ \quad \quad d(\operatorname{matof}((n + 1), (n + 1), \\ \quad \quad \quad \quad (\lambda {[}\operatorname{dav\_}, \operatorname{daw\_}] \in (\operatorname{interval}(1, (n + 1)) \times \operatorname{interval}(1, (n + 1))).\; \\ \quad \quad \quad \quad \quad \left(\text{if }(\operatorname{daw\_} = 1)\text{ then }\left(\text{if }(\operatorname{dav\_} = 1)\text{ then }\operatorname{entry}(y, \operatorname{dav\_}, \operatorname{daw\_})\text{ else }\left(\text{if }(\operatorname{dav\_} \leq (m + 1))\text{ then }\operatorname{zero}(r)\text{ else }\operatorname{entry}(y, \operatorname{dav\_}, \operatorname{daw\_})\right)\right)\text{ else }\operatorname{entry}(y, \operatorname{dav\_}, \operatorname{daw\_})\right)))) \\ \quad \quad =\; d(y) \end{array}}}
\vnbprop{daf-clear}{\vnbstmt{\begin{array}{@{}l@{}} \forall r, n, d, y. \\ \quad \operatorname{is\text{-}commutative\text{-}ring}(r) \wedge \\ \quad n \in \mathbb{N} \wedge \\ \quad d \in (\operatorname{mat}((n + 1), (n + 1), \operatorname{carr}(r)) \to \operatorname{carr}(r)) \wedge \\ \quad \forall a, b, c, p, s. \\ \quad \quad a \in \operatorname{mat}((n + 1), (n + 1), \operatorname{carr}(r)) \wedge \\ \quad \quad b \in \operatorname{mat}((n + 1), (n + 1), \operatorname{carr}(r)) \wedge \\ \quad \quad c \in \operatorname{mat}((n + 1), (n + 1), \operatorname{carr}(r)) \wedge \\ \quad \quad p \in \operatorname{interval}(1, (n + 1)) \wedge \\ \quad \quad s \in \operatorname{carr}(r) \wedge \\ \quad \quad \forall i, k \in \operatorname{interval}(1, (n + 1)). \\ \quad \quad \quad \neg (i = p) \Rightarrow (\operatorname{entry}(b, i, k) = \operatorname{entry}(a, i, k)) \wedge \\ \quad \quad \forall i, k \in \operatorname{interval}(1, (n + 1)). \\ \quad \quad \quad \neg (i = p) \Rightarrow (\operatorname{entry}(c, i, k) = \operatorname{entry}(a, i, k)) \wedge \\ \quad \quad \forall k \in \operatorname{interval}(1, (n + 1)). \\ \quad \quad \quad \operatorname{entry}(c, p, k) = \operatorname{add}(r)(\operatorname{mul}(r)(s, \operatorname{entry}(a, p, k)), \operatorname{entry}(b, p, k)) \Rightarrow \\ \quad \quad \quad d(c) = \operatorname{add}(r)(\operatorname{mul}(r)(s, d(a)), d(b)) \wedge \\ \quad \forall a, p, q. \\ \quad \quad a \in \operatorname{mat}((n + 1), (n + 1), \operatorname{carr}(r)) \wedge \\ \quad \quad p \in \operatorname{interval}(1, (n + 1)) \wedge \\ \quad \quad q \in \operatorname{interval}(1, (n + 1)) \wedge \\ \quad \quad \neg (p = q) \wedge \\ \quad \quad \forall k \in \operatorname{interval}(1, (n + 1)).\; (\operatorname{entry}(a, p, k) = \operatorname{entry}(a, q, k)) \Rightarrow \\ \quad \quad \quad d(a) = \operatorname{zero}(r) \wedge \\ \quad y \in \operatorname{mat}((n + 1), (n + 1), \operatorname{carr}(r)) \wedge \\ \quad \operatorname{entry}(y, 1, 1) = \operatorname{one}(r) \wedge \\ \quad \forall k \in \operatorname{interval}(1, (n + 1)).\; (\neg (k = 1) \Rightarrow (\operatorname{entry}(y, 1, k) = \operatorname{zero}(r))) \Rightarrow \\ \quad \quad d(\operatorname{matof}((n + 1), (n + 1), \\ \quad \quad \quad \quad (\lambda {[}\operatorname{dav\_}, \operatorname{daw\_}] \in (\operatorname{interval}(1, (n + 1)) \times \operatorname{interval}(1, (n + 1))).\; \\ \quad \quad \quad \quad \quad \left(\text{if }(\operatorname{daw\_} = 1)\text{ then }\left(\text{if }(\operatorname{dav\_} = 1)\text{ then }\operatorname{entry}(y, \operatorname{dav\_}, \operatorname{daw\_})\text{ else }\left(\text{if }(\operatorname{dav\_} \leq (n + 1))\text{ then }\operatorname{zero}(r)\text{ else }\operatorname{entry}(y, \operatorname{dav\_}, \operatorname{daw\_})\right)\right)\text{ else }\operatorname{entry}(y, \operatorname{dav\_}, \operatorname{daw\_})\right)))) \\ \quad \quad =\; d(y) \end{array}}}
\vnbprop{daf-bordered}{\vnbstmt{\begin{array}{@{}l@{}} \forall r, n, y. \\ \quad \operatorname{is\text{-}commutative\text{-}ring}(r) \wedge \\ \quad n \in \mathbb{N} \wedge \\ \quad y \in \operatorname{mat}((n + 1), (n + 1), \operatorname{carr}(r)) \wedge \\ \quad \operatorname{entry}(y, 1, 1) = \operatorname{one}(r) \wedge \\ \quad \forall k \in \operatorname{interval}(1, (n + 1)).\; (\neg (k = 1) \Rightarrow (\operatorname{entry}(y, 1, k) = \operatorname{zero}(r))) \Rightarrow \\ \quad \quad \operatorname{matof}((n + 1), (n + 1), \\ \quad \quad \quad (\lambda {[}\operatorname{dav\_}, \operatorname{daw\_}] \in (\operatorname{interval}(1, (n + 1)) \times \operatorname{interval}(1, (n + 1))).\; \\ \quad \quad \quad \quad \left(\text{if }(\operatorname{daw\_} = 1)\text{ then }\left(\text{if }(\operatorname{dav\_} = 1)\text{ then }\operatorname{entry}(y, \operatorname{dav\_}, \operatorname{daw\_})\text{ else }\left(\text{if }(\operatorname{dav\_} \leq (n + 1))\text{ then }\operatorname{zero}(r)\text{ else }\operatorname{entry}(y, \operatorname{dav\_}, \operatorname{daw\_})\right)\right)\text{ else }\operatorname{entry}(y, \operatorname{dav\_}, \operatorname{daw\_})\right))) \\ \quad \quad =\; \operatorname{border}(r, \operatorname{one}(r), \\ \quad \quad \quad \operatorname{submat}(\operatorname{matof}((n + 1), (n + 1), \\ \quad \quad \quad \quad \quad (\lambda {[}\operatorname{dav\_}, \operatorname{daw\_}] \in (\operatorname{interval}(1, (n + 1)) \times \operatorname{interval}(1, (n + 1))).\; \\ \quad \quad \quad \quad \quad \quad \left(\text{if }(\operatorname{daw\_} = 1)\text{ then }\left(\text{if }(\operatorname{dav\_} = 1)\text{ then }\operatorname{entry}(y, \operatorname{dav\_}, \operatorname{daw\_})\text{ else }\left(\text{if }(\operatorname{dav\_} \leq (n + 1))\text{ then }\operatorname{zero}(r)\text{ else }\operatorname{entry}(y, \operatorname{dav\_}, \operatorname{daw\_})\right)\right)\text{ else }\operatorname{entry}(y, \operatorname{dav\_}, \operatorname{daw\_})\right))), \\ \quad \quad \quad \quad n, n), \\ \quad \quad \quad n, n) \end{array}}}
\vnbprop{daf-cl}{\vnbstmt{\begin{array}{@{}l@{}} \forall n \in \mathbb{N}, r, d, y. \\ \quad \forall r, d.; \forall x \in \operatorname{mat}(n, n, \operatorname{carr}(r)). \\ \quad \quad \operatorname{is\text{-}commutative\text{-}ring}(r) \wedge \\ \quad \quad d \in (\operatorname{mat}(n, n, \operatorname{carr}(r)) \to \operatorname{carr}(r)) \wedge \\ \quad \quad \forall a, b, c, p, s. \\ \quad \quad \quad a \in \operatorname{mat}(n, n, \operatorname{carr}(r)) \wedge \\ \quad \quad \quad b \in \operatorname{mat}(n, n, \operatorname{carr}(r)) \wedge \\ \quad \quad \quad c \in \operatorname{mat}(n, n, \operatorname{carr}(r)) \wedge \\ \quad \quad \quad p \in \operatorname{interval}(1, n) \wedge \\ \quad \quad \quad s \in \operatorname{carr}(r) \wedge \\ \quad \quad \quad \forall i, k \in \operatorname{interval}(1, n). \\ \quad \quad \quad \quad \neg (i = p) \Rightarrow (\operatorname{entry}(b, i, k) = \operatorname{entry}(a, i, k)) \wedge \\ \quad \quad \quad \forall i, k \in \operatorname{interval}(1, n). \\ \quad \quad \quad \quad \neg (i = p) \Rightarrow (\operatorname{entry}(c, i, k) = \operatorname{entry}(a, i, k)) \wedge \\ \quad \quad \quad \forall k \in \operatorname{interval}(1, n). \\ \quad \quad \quad \quad \operatorname{entry}(c, p, k) = \operatorname{add}(r)(\operatorname{mul}(r)(s, \operatorname{entry}(a, p, k)), \operatorname{entry}(b, p, k)) \Rightarrow \\ \quad \quad \quad \quad d(c) = \operatorname{add}(r)(\operatorname{mul}(r)(s, d(a)), d(b)) \wedge \\ \quad \quad \forall a, p, q. \\ \quad \quad \quad a \in \operatorname{mat}(n, n, \operatorname{carr}(r)) \wedge \\ \quad \quad \quad p \in \operatorname{interval}(1, n) \wedge \\ \quad \quad \quad q \in \operatorname{interval}(1, n) \wedge \\ \quad \quad \quad \neg (p = q) \wedge \\ \quad \quad \quad \forall k \in \operatorname{interval}(1, n).\; (\operatorname{entry}(a, p, k) = \operatorname{entry}(a, q, k)) \Rightarrow \\ \quad \quad \quad \quad d(a) = \operatorname{zero}(r) \Rightarrow \\ \quad \quad \quad d(x) = \operatorname{mul}(r)(\operatorname{det}(r, n, x), d(\operatorname{identmat}(r, n))) \wedge \\ \quad \operatorname{is\text{-}commutative\text{-}ring}(r) \wedge \\ \quad d \in (\operatorname{mat}((n + 1), (n + 1), \operatorname{carr}(r)) \to \operatorname{carr}(r)) \wedge \\ \quad \forall a, b, c, p, s. \\ \quad \quad a \in \operatorname{mat}((n + 1), (n + 1), \operatorname{carr}(r)) \wedge \\ \quad \quad b \in \operatorname{mat}((n + 1), (n + 1), \operatorname{carr}(r)) \wedge \\ \quad \quad c \in \operatorname{mat}((n + 1), (n + 1), \operatorname{carr}(r)) \wedge \\ \quad \quad p \in \operatorname{interval}(1, (n + 1)) \wedge \\ \quad \quad s \in \operatorname{carr}(r) \wedge \\ \quad \quad \forall i, k \in \operatorname{interval}(1, (n + 1)). \\ \quad \quad \quad \neg (i = p) \Rightarrow (\operatorname{entry}(b, i, k) = \operatorname{entry}(a, i, k)) \wedge \\ \quad \quad \forall i, k \in \operatorname{interval}(1, (n + 1)). \\ \quad \quad \quad \neg (i = p) \Rightarrow (\operatorname{entry}(c, i, k) = \operatorname{entry}(a, i, k)) \wedge \\ \quad \quad \forall k \in \operatorname{interval}(1, (n + 1)). \\ \quad \quad \quad \operatorname{entry}(c, p, k) = \operatorname{add}(r)(\operatorname{mul}(r)(s, \operatorname{entry}(a, p, k)), \operatorname{entry}(b, p, k)) \Rightarrow \\ \quad \quad \quad d(c) = \operatorname{add}(r)(\operatorname{mul}(r)(s, d(a)), d(b)) \wedge \\ \quad \forall a, p, q. \\ \quad \quad a \in \operatorname{mat}((n + 1), (n + 1), \operatorname{carr}(r)) \wedge \\ \quad \quad p \in \operatorname{interval}(1, (n + 1)) \wedge \\ \quad \quad q \in \operatorname{interval}(1, (n + 1)) \wedge \\ \quad \quad \neg (p = q) \wedge \\ \quad \quad \forall k \in \operatorname{interval}(1, (n + 1)).\; (\operatorname{entry}(a, p, k) = \operatorname{entry}(a, q, k)) \Rightarrow \\ \quad \quad \quad d(a) = \operatorname{zero}(r) \wedge \\ \quad y \in \operatorname{mat}((n + 1), (n + 1), \operatorname{carr}(r)) \wedge \\ \quad \operatorname{entry}(y, 1, 1) = \operatorname{one}(r) \wedge \\ \quad \forall k \in \operatorname{interval}(1, (n + 1)).\; (\neg (k = 1) \Rightarrow (\operatorname{entry}(y, 1, k) = \operatorname{zero}(r))) \Rightarrow \\ \quad \quad d(y) \\ \quad \quad =\; \operatorname{mul}(r)(\operatorname{det}(r, n, \\ \quad \quad \quad \quad \operatorname{submat}(\operatorname{matof}((n + 1), (n + 1), \\ \quad \quad \quad \quad \quad \quad (\lambda {[}\operatorname{dav\_}, \operatorname{daw\_}] \in (\operatorname{interval}(1, (n + 1)) \times \operatorname{interval}(1, (n + 1))).\; \\ \quad \quad \quad \quad \quad \quad \quad \left(\text{if }(\operatorname{daw\_} = 1)\text{ then }\left(\text{if }(\operatorname{dav\_} = 1)\text{ then }\operatorname{entry}(y, \operatorname{dav\_}, \operatorname{daw\_})\text{ else }\left(\text{if }(\operatorname{dav\_} \leq (n + 1))\text{ then }\operatorname{zero}(r)\text{ else }\operatorname{entry}(y, \operatorname{dav\_}, \operatorname{daw\_})\right)\right)\text{ else }\operatorname{entry}(y, \operatorname{dav\_}, \operatorname{daw\_})\right))), \\ \quad \quad \quad \quad \quad n, n)), \\ \quad \quad \quad d(\operatorname{identmat}(r, (n + 1)))) \end{array}}}
\vnbprop{daf-basis}{\vnbstmt{\begin{array}{@{}l@{}} \forall n \in \mathbb{N}, r, d, x, j. \\ \quad \forall r, d.; \forall \mathit{dfx} \in \operatorname{mat}(n, n, \operatorname{carr}(r)). \\ \quad \quad \operatorname{is\text{-}commutative\text{-}ring}(r) \wedge \\ \quad \quad d \in (\operatorname{mat}(n, n, \operatorname{carr}(r)) \to \operatorname{carr}(r)) \wedge \\ \quad \quad \forall a, b, c, p, s. \\ \quad \quad \quad a \in \operatorname{mat}(n, n, \operatorname{carr}(r)) \wedge \\ \quad \quad \quad b \in \operatorname{mat}(n, n, \operatorname{carr}(r)) \wedge \\ \quad \quad \quad c \in \operatorname{mat}(n, n, \operatorname{carr}(r)) \wedge \\ \quad \quad \quad p \in \operatorname{interval}(1, n) \wedge \\ \quad \quad \quad s \in \operatorname{carr}(r) \wedge \\ \quad \quad \quad \forall i, k \in \operatorname{interval}(1, n). \\ \quad \quad \quad \quad \neg (i = p) \Rightarrow (\operatorname{entry}(b, i, k) = \operatorname{entry}(a, i, k)) \wedge \\ \quad \quad \quad \forall i, k \in \operatorname{interval}(1, n). \\ \quad \quad \quad \quad \neg (i = p) \Rightarrow (\operatorname{entry}(c, i, k) = \operatorname{entry}(a, i, k)) \wedge \\ \quad \quad \quad \forall k \in \operatorname{interval}(1, n). \\ \quad \quad \quad \quad \operatorname{entry}(c, p, k) = \operatorname{add}(r)(\operatorname{mul}(r)(s, \operatorname{entry}(a, p, k)), \operatorname{entry}(b, p, k)) \Rightarrow \\ \quad \quad \quad \quad d(c) = \operatorname{add}(r)(\operatorname{mul}(r)(s, d(a)), d(b)) \wedge \\ \quad \quad \forall a, p, q. \\ \quad \quad \quad a \in \operatorname{mat}(n, n, \operatorname{carr}(r)) \wedge \\ \quad \quad \quad p \in \operatorname{interval}(1, n) \wedge \\ \quad \quad \quad q \in \operatorname{interval}(1, n) \wedge \\ \quad \quad \quad \neg (p = q) \wedge \\ \quad \quad \quad \forall k \in \operatorname{interval}(1, n).\; (\operatorname{entry}(a, p, k) = \operatorname{entry}(a, q, k)) \Rightarrow \\ \quad \quad \quad \quad d(a) = \operatorname{zero}(r) \Rightarrow \\ \quad \quad \quad d(\mathit{dfx}) = \operatorname{mul}(r)(\operatorname{det}(r, n, \mathit{dfx}), d(\operatorname{identmat}(r, n))) \wedge \\ \quad \operatorname{is\text{-}commutative\text{-}ring}(r) \wedge \\ \quad d \in (\operatorname{mat}((n + 1), (n + 1), \operatorname{carr}(r)) \to \operatorname{carr}(r)) \wedge \\ \quad \forall a, b, c, p, s. \\ \quad \quad a \in \operatorname{mat}((n + 1), (n + 1), \operatorname{carr}(r)) \wedge \\ \quad \quad b \in \operatorname{mat}((n + 1), (n + 1), \operatorname{carr}(r)) \wedge \\ \quad \quad c \in \operatorname{mat}((n + 1), (n + 1), \operatorname{carr}(r)) \wedge \\ \quad \quad p \in \operatorname{interval}(1, (n + 1)) \wedge \\ \quad \quad s \in \operatorname{carr}(r) \wedge \\ \quad \quad \forall i, k \in \operatorname{interval}(1, (n + 1)). \\ \quad \quad \quad \neg (i = p) \Rightarrow (\operatorname{entry}(b, i, k) = \operatorname{entry}(a, i, k)) \wedge \\ \quad \quad \forall i, k \in \operatorname{interval}(1, (n + 1)). \\ \quad \quad \quad \neg (i = p) \Rightarrow (\operatorname{entry}(c, i, k) = \operatorname{entry}(a, i, k)) \wedge \\ \quad \quad \forall k \in \operatorname{interval}(1, (n + 1)). \\ \quad \quad \quad \operatorname{entry}(c, p, k) = \operatorname{add}(r)(\operatorname{mul}(r)(s, \operatorname{entry}(a, p, k)), \operatorname{entry}(b, p, k)) \Rightarrow \\ \quad \quad \quad d(c) = \operatorname{add}(r)(\operatorname{mul}(r)(s, d(a)), d(b)) \wedge \\ \quad \forall a, p, q. \\ \quad \quad a \in \operatorname{mat}((n + 1), (n + 1), \operatorname{carr}(r)) \wedge \\ \quad \quad p \in \operatorname{interval}(1, (n + 1)) \wedge \\ \quad \quad q \in \operatorname{interval}(1, (n + 1)) \wedge \\ \quad \quad \neg (p = q) \wedge \\ \quad \quad \forall k \in \operatorname{interval}(1, (n + 1)).\; (\operatorname{entry}(a, p, k) = \operatorname{entry}(a, q, k)) \Rightarrow \\ \quad \quad \quad d(a) = \operatorname{zero}(r) \wedge \\ \quad x \in \operatorname{mat}((n + 1), (n + 1), \operatorname{carr}(r)) \wedge \\ \quad j \in \operatorname{interval}(1, (n + 1)) \wedge \\ \quad \forall k \in \operatorname{interval}(1, (n + 1)).\; (\operatorname{entry}(x, 1, k) = \left(\text{if }(k = j)\text{ then }\operatorname{one}(r)\text{ else }\operatorname{zero}(r)\right)) \Rightarrow \\ \quad \quad d(x) = \operatorname{mul}(r)(\operatorname{det}(r, (n + 1), x), d(\operatorname{identmat}(r, (n + 1)))) \end{array}}}
\vnbprop{daf-psum-ind}{\vnbstmt{\begin{array}{@{}l@{}} \forall m \in \mathbb{N}, r, n, d, a. \\ \quad \operatorname{is\text{-}commutative\text{-}ring}(r) \wedge \\ \quad n \in \mathbb{N} \wedge \\ \quad d \in (\operatorname{mat}((n + 1), (n + 1), \operatorname{carr}(r)) \to \operatorname{carr}(r)) \wedge \\ \quad \forall \mathit{dfa}, b, c, p, s. \\ \quad \quad \mathit{dfa} \in \operatorname{mat}((n + 1), (n + 1), \operatorname{carr}(r)) \wedge \\ \quad \quad b \in \operatorname{mat}((n + 1), (n + 1), \operatorname{carr}(r)) \wedge \\ \quad \quad c \in \operatorname{mat}((n + 1), (n + 1), \operatorname{carr}(r)) \wedge \\ \quad \quad p \in \operatorname{interval}(1, (n + 1)) \wedge \\ \quad \quad s \in \operatorname{carr}(r) \wedge \\ \quad \quad \forall i, k \in \operatorname{interval}(1, (n + 1)). \\ \quad \quad \quad \neg (i = p) \Rightarrow (\operatorname{entry}(b, i, k) = \operatorname{entry}(\mathit{dfa}, i, k)) \wedge \\ \quad \quad \forall i, k \in \operatorname{interval}(1, (n + 1)). \\ \quad \quad \quad \neg (i = p) \Rightarrow (\operatorname{entry}(c, i, k) = \operatorname{entry}(\mathit{dfa}, i, k)) \wedge \\ \quad \quad \forall k \in \operatorname{interval}(1, (n + 1)). \\ \quad \quad \quad \operatorname{entry}(c, p, k) = \operatorname{add}(r)(\operatorname{mul}(r)(s, \operatorname{entry}(\mathit{dfa}, p, k)), \operatorname{entry}(b, p, k)) \Rightarrow \\ \quad \quad \quad d(c) = \operatorname{add}(r)(\operatorname{mul}(r)(s, d(\mathit{dfa})), d(b)) \wedge \\ \quad \forall \mathit{dfa}, p, q. \\ \quad \quad \mathit{dfa} \in \operatorname{mat}((n + 1), (n + 1), \operatorname{carr}(r)) \wedge \\ \quad \quad p \in \operatorname{interval}(1, (n + 1)) \wedge \\ \quad \quad q \in \operatorname{interval}(1, (n + 1)) \wedge \\ \quad \quad \neg (p = q) \wedge \\ \quad \quad \forall k \in \operatorname{interval}(1, (n + 1)).\; (\operatorname{entry}(\mathit{dfa}, p, k) = \operatorname{entry}(\mathit{dfa}, q, k)) \Rightarrow \\ \quad \quad \quad d(\mathit{dfa}) = \operatorname{zero}(r) \wedge \\ \quad a \in \operatorname{mat}((n + 1), (n + 1), \operatorname{carr}(r)) \wedge \\ \quad \forall j \in \operatorname{interval}(1, (n + 1)). \\ \quad \quad d(\operatorname{matof}((n + 1), (n + 1), \\ \quad \quad \quad \quad (\lambda {[}\operatorname{dav\_}, \operatorname{daw\_}] \in (\operatorname{interval}(1, (n + 1)) \times \operatorname{interval}(1, (n + 1))).\; \\ \quad \quad \quad \quad \quad \left(\text{if }(\operatorname{dav\_} = 1)\text{ then }\left(\text{if }(\operatorname{daw\_} = j)\text{ then }\operatorname{one}(r)\text{ else }\operatorname{zero}(r)\right)\text{ else }\operatorname{entry}(a, \operatorname{dav\_}, \operatorname{daw\_})\right)))) \\ \quad \quad =\; \operatorname{mul}(r)(\operatorname{det}(r, (n + 1), \\ \quad \quad \quad \quad \operatorname{matof}((n + 1), (n + 1), \\ \quad \quad \quad \quad \quad (\lambda {[}\operatorname{dav\_}, \operatorname{daw\_}] \in (\operatorname{interval}(1, (n + 1)) \times \operatorname{interval}(1, (n + 1))).\; \\ \quad \quad \quad \quad \quad \quad \left(\text{if }(\operatorname{dav\_} = 1)\text{ then }\left(\text{if }(\operatorname{daw\_} = j)\text{ then }\operatorname{one}(r)\text{ else }\operatorname{zero}(r)\right)\text{ else }\operatorname{entry}(a, \operatorname{dav\_}, \operatorname{daw\_})\right)))), \\ \quad \quad \quad d(\operatorname{identmat}(r, (n + 1)))) \wedge \\ \quad m \leq (n + 1) \Rightarrow \\ \quad \quad d(\operatorname{matof}((n + 1), (n + 1), \\ \quad \quad \quad \quad (\lambda {[}\operatorname{dav\_}, \operatorname{daw\_}] \in (\operatorname{interval}(1, (n + 1)) \times \operatorname{interval}(1, (n + 1))).\; \\ \quad \quad \quad \quad \quad \left(\text{if }(\operatorname{dav\_} = 1)\text{ then }\left(\text{if }(\operatorname{daw\_} \leq m)\text{ then }\operatorname{entry}(a, 1, \operatorname{daw\_})\text{ else }\operatorname{zero}(r)\right)\text{ else }\operatorname{entry}(a, \operatorname{dav\_}, \operatorname{daw\_})\right)))) \\ \quad \quad =\; \operatorname{mul}(r)(\operatorname{det}(r, (n + 1), \\ \quad \quad \quad \quad \operatorname{matof}((n + 1), (n + 1), \\ \quad \quad \quad \quad \quad (\lambda {[}\operatorname{dav\_}, \operatorname{daw\_}] \in (\operatorname{interval}(1, (n + 1)) \times \operatorname{interval}(1, (n + 1))).\; \\ \quad \quad \quad \quad \quad \quad \left(\text{if }(\operatorname{dav\_} = 1)\text{ then }\left(\text{if }(\operatorname{daw\_} \leq m)\text{ then }\operatorname{entry}(a, 1, \operatorname{daw\_})\text{ else }\operatorname{zero}(r)\right)\text{ else }\operatorname{entry}(a, \operatorname{dav\_}, \operatorname{daw\_})\right)))), \\ \quad \quad \quad d(\operatorname{identmat}(r, (n + 1)))) \end{array}}}
\vnbprop{daf-psum}{\vnbstmt{\begin{array}{@{}l@{}} \forall r, n, d, a. \\ \quad \operatorname{is\text{-}commutative\text{-}ring}(r) \wedge \\ \quad n \in \mathbb{N} \wedge \\ \quad d \in (\operatorname{mat}((n + 1), (n + 1), \operatorname{carr}(r)) \to \operatorname{carr}(r)) \wedge \\ \quad \forall \mathit{dfa}, b, c, p, s. \\ \quad \quad \mathit{dfa} \in \operatorname{mat}((n + 1), (n + 1), \operatorname{carr}(r)) \wedge \\ \quad \quad b \in \operatorname{mat}((n + 1), (n + 1), \operatorname{carr}(r)) \wedge \\ \quad \quad c \in \operatorname{mat}((n + 1), (n + 1), \operatorname{carr}(r)) \wedge \\ \quad \quad p \in \operatorname{interval}(1, (n + 1)) \wedge \\ \quad \quad s \in \operatorname{carr}(r) \wedge \\ \quad \quad \forall i, k \in \operatorname{interval}(1, (n + 1)). \\ \quad \quad \quad \neg (i = p) \Rightarrow (\operatorname{entry}(b, i, k) = \operatorname{entry}(\mathit{dfa}, i, k)) \wedge \\ \quad \quad \forall i, k \in \operatorname{interval}(1, (n + 1)). \\ \quad \quad \quad \neg (i = p) \Rightarrow (\operatorname{entry}(c, i, k) = \operatorname{entry}(\mathit{dfa}, i, k)) \wedge \\ \quad \quad \forall k \in \operatorname{interval}(1, (n + 1)). \\ \quad \quad \quad \operatorname{entry}(c, p, k) = \operatorname{add}(r)(\operatorname{mul}(r)(s, \operatorname{entry}(\mathit{dfa}, p, k)), \operatorname{entry}(b, p, k)) \Rightarrow \\ \quad \quad \quad d(c) = \operatorname{add}(r)(\operatorname{mul}(r)(s, d(\mathit{dfa})), d(b)) \wedge \\ \quad \forall \mathit{dfa}, p, q. \\ \quad \quad \mathit{dfa} \in \operatorname{mat}((n + 1), (n + 1), \operatorname{carr}(r)) \wedge \\ \quad \quad p \in \operatorname{interval}(1, (n + 1)) \wedge \\ \quad \quad q \in \operatorname{interval}(1, (n + 1)) \wedge \\ \quad \quad \neg (p = q) \wedge \\ \quad \quad \forall k \in \operatorname{interval}(1, (n + 1)).\; (\operatorname{entry}(\mathit{dfa}, p, k) = \operatorname{entry}(\mathit{dfa}, q, k)) \Rightarrow \\ \quad \quad \quad d(\mathit{dfa}) = \operatorname{zero}(r) \wedge \\ \quad a \in \operatorname{mat}((n + 1), (n + 1), \operatorname{carr}(r)) \wedge \\ \quad \forall j \in \operatorname{interval}(1, (n + 1)). \\ \quad \quad d(\operatorname{matof}((n + 1), (n + 1), \\ \quad \quad \quad \quad (\lambda {[}\operatorname{dav\_}, \operatorname{daw\_}] \in (\operatorname{interval}(1, (n + 1)) \times \operatorname{interval}(1, (n + 1))).\; \\ \quad \quad \quad \quad \quad \left(\text{if }(\operatorname{dav\_} = 1)\text{ then }\left(\text{if }(\operatorname{daw\_} = j)\text{ then }\operatorname{one}(r)\text{ else }\operatorname{zero}(r)\right)\text{ else }\operatorname{entry}(a, \operatorname{dav\_}, \operatorname{daw\_})\right)))) \\ \quad \quad =\; \operatorname{mul}(r)(\operatorname{det}(r, (n + 1), \\ \quad \quad \quad \quad \operatorname{matof}((n + 1), (n + 1), \\ \quad \quad \quad \quad \quad (\lambda {[}\operatorname{dav\_}, \operatorname{daw\_}] \in (\operatorname{interval}(1, (n + 1)) \times \operatorname{interval}(1, (n + 1))).\; \\ \quad \quad \quad \quad \quad \quad \left(\text{if }(\operatorname{dav\_} = 1)\text{ then }\left(\text{if }(\operatorname{daw\_} = j)\text{ then }\operatorname{one}(r)\text{ else }\operatorname{zero}(r)\right)\text{ else }\operatorname{entry}(a, \operatorname{dav\_}, \operatorname{daw\_})\right)))), \\ \quad \quad \quad d(\operatorname{identmat}(r, (n + 1)))) \Rightarrow \\ \quad \quad d(a) = \operatorname{mul}(r)(\operatorname{det}(r, (n + 1), a), d(\operatorname{identmat}(r, (n + 1)))) \end{array}}}
\vnbprop{det-alternating-form-ind}{\vnbstmt{\begin{array}{@{}l@{}} \forall n \in \mathbb{N}, r, d.; \forall x \in \operatorname{mat}(n, n, \operatorname{carr}(r)). \\ \quad \operatorname{is\text{-}commutative\text{-}ring}(r) \wedge \\ \quad d \in (\operatorname{mat}(n, n, \operatorname{carr}(r)) \to \operatorname{carr}(r)) \wedge \\ \quad \forall a, b, c, p, s. \\ \quad \quad a \in \operatorname{mat}(n, n, \operatorname{carr}(r)) \wedge \\ \quad \quad b \in \operatorname{mat}(n, n, \operatorname{carr}(r)) \wedge \\ \quad \quad c \in \operatorname{mat}(n, n, \operatorname{carr}(r)) \wedge \\ \quad \quad p \in \operatorname{interval}(1, n) \wedge \\ \quad \quad s \in \operatorname{carr}(r) \wedge \\ \quad \quad \forall i, k \in \operatorname{interval}(1, n). \\ \quad \quad \quad \neg (i = p) \Rightarrow (\operatorname{entry}(b, i, k) = \operatorname{entry}(a, i, k)) \wedge \\ \quad \quad \forall i, k \in \operatorname{interval}(1, n). \\ \quad \quad \quad \neg (i = p) \Rightarrow (\operatorname{entry}(c, i, k) = \operatorname{entry}(a, i, k)) \wedge \\ \quad \quad \forall k \in \operatorname{interval}(1, n). \\ \quad \quad \quad \operatorname{entry}(c, p, k) = \operatorname{add}(r)(\operatorname{mul}(r)(s, \operatorname{entry}(a, p, k)), \operatorname{entry}(b, p, k)) \Rightarrow \\ \quad \quad \quad d(c) = \operatorname{add}(r)(\operatorname{mul}(r)(s, d(a)), d(b)) \wedge \\ \quad \forall a, p, q. \\ \quad \quad a \in \operatorname{mat}(n, n, \operatorname{carr}(r)) \wedge \\ \quad \quad p \in \operatorname{interval}(1, n) \wedge \\ \quad \quad q \in \operatorname{interval}(1, n) \wedge \\ \quad \quad \neg (p = q) \wedge \\ \quad \quad \forall k \in \operatorname{interval}(1, n).\; (\operatorname{entry}(a, p, k) = \operatorname{entry}(a, q, k)) \Rightarrow \\ \quad \quad \quad d(a) = \operatorname{zero}(r) \Rightarrow \\ \quad \quad d(x) = \operatorname{mul}(r)(\operatorname{det}(r, n, x), d(\operatorname{identmat}(r, n))) \end{array}}}
\vnbprop{det-alternating-form}{\vnbstmt{\begin{array}{@{}l@{}} \forall r, n, d, a. \\ \quad \operatorname{is\text{-}commutative\text{-}ring}(r) \wedge \\ \quad n \in \mathbb{N} \wedge \\ \quad d \in (\operatorname{mat}(n, n, \operatorname{carr}(r)) \to \operatorname{carr}(r)) \wedge \\ \quad \forall \mathit{dfa}, b, c, p, s. \\ \quad \quad \mathit{dfa} \in \operatorname{mat}(n, n, \operatorname{carr}(r)) \wedge \\ \quad \quad b \in \operatorname{mat}(n, n, \operatorname{carr}(r)) \wedge \\ \quad \quad c \in \operatorname{mat}(n, n, \operatorname{carr}(r)) \wedge \\ \quad \quad p \in \operatorname{interval}(1, n) \wedge \\ \quad \quad s \in \operatorname{carr}(r) \wedge \\ \quad \quad \forall i, k \in \operatorname{interval}(1, n). \\ \quad \quad \quad \neg (i = p) \Rightarrow (\operatorname{entry}(b, i, k) = \operatorname{entry}(\mathit{dfa}, i, k)) \wedge \\ \quad \quad \forall i, k \in \operatorname{interval}(1, n). \\ \quad \quad \quad \neg (i = p) \Rightarrow (\operatorname{entry}(c, i, k) = \operatorname{entry}(\mathit{dfa}, i, k)) \wedge \\ \quad \quad \forall k \in \operatorname{interval}(1, n). \\ \quad \quad \quad \operatorname{entry}(c, p, k) = \operatorname{add}(r)(\operatorname{mul}(r)(s, \operatorname{entry}(\mathit{dfa}, p, k)), \operatorname{entry}(b, p, k)) \Rightarrow \\ \quad \quad \quad d(c) = \operatorname{add}(r)(\operatorname{mul}(r)(s, d(\mathit{dfa})), d(b)) \wedge \\ \quad \forall \mathit{dfa}, p, q. \\ \quad \quad \mathit{dfa} \in \operatorname{mat}(n, n, \operatorname{carr}(r)) \wedge \\ \quad \quad p \in \operatorname{interval}(1, n) \wedge \\ \quad \quad q \in \operatorname{interval}(1, n) \wedge \\ \quad \quad \neg (p = q) \wedge \\ \quad \quad \forall k \in \operatorname{interval}(1, n).\; (\operatorname{entry}(\mathit{dfa}, p, k) = \operatorname{entry}(\mathit{dfa}, q, k)) \Rightarrow \\ \quad \quad \quad d(\mathit{dfa}) = \operatorname{zero}(r) \wedge \\ \quad a \in \operatorname{mat}(n, n, \operatorname{carr}(r)) \Rightarrow \\ \quad \quad d(a) = \operatorname{mul}(r)(\operatorname{det}(r, n, a), d(\operatorname{identmat}(r, n))) \end{array}}}

\vnbbookfile{theorem-library/det-multiplicative}{theorem-library/det-multiplicative}
\noindent{\small det-multiplicative.scm -- det(P Q) = det(P) det(Q) over a commutative ring (Hoffman-Kunze 5.3 Theorem 3; the notes' proposition "det is multiplicative"), and the notes' corollary 3.20 (R a field, V invertible: det(A V) = det A det V). Batch 26-B, 2026-09-24.}\par
\vnbprop{det-multiplicative}{\vnbstmt{\begin{array}{@{}l@{}} \forall r, n, p, q. \\ \quad (\operatorname{is\text{-}commutative\text{-}ring}(r) \wedge ((n \in \mathbb{N}) \wedge ((p \in \operatorname{mat}(n, n, \operatorname{carr}(r))) \wedge (q \in \operatorname{mat}(n, n, \operatorname{carr}(r)))))) \Rightarrow  \\ \quad \quad \operatorname{det}(r, n, \operatorname{matmul}(r, p, q)) = \operatorname{mul}(r)(\operatorname{det}(r, n, p), \operatorname{det}(r, n, q)) \end{array}}}
\vnbprop{det-mul-invertible}{\vnbstmt{\begin{array}{@{}l@{}} \forall r, n, a, v. \\ \quad (\operatorname{is\text{-}field\text{-}ring}(r) \wedge ((n \in \mathbb{N}) \wedge ((a \in \operatorname{mat}(n, n, \operatorname{carr}(r))) \wedge \operatorname{is\text{-}invertible\text{-}mat}(r, n, v)))) \Rightarrow  \\ \quad \quad \operatorname{det}(r, n, \operatorname{matmul}(r, a, v)) = \operatorname{mul}(r)(\operatorname{det}(r, n, a), \operatorname{det}(r, n, v)) \end{array}}}

\vnbbookfile{theorem-library/det-transpose}{theorem-library/det-transpose}
\noindent{\small det-transpose.scm -- the transpose of a matrix and det(A\textasciicircum{}t) = det A; the COLUMN clauses of the linear-algebra notes' 3.18 (linearity in each column, two equal columns give 0, interchanging two columns negates det), 3.19 (det under an elementary COLUMN operation) and 3.22 (expansion along any column).  Batch 32}\par
\vnbprop{transpose-in-mat}{\vnbstmt{\begin{array}{@{}l@{}} \forall a, m, n, x. \\ \quad (a \in \operatorname{mat}(m, n, x)) \Rightarrow (\operatorname{transpose}(a, m, n) \in \operatorname{mat}(n, m, x)) \end{array}}}
\vnbprop{transpose-entry}{\vnbstmt{\begin{array}{@{}l@{}} \forall a, m, n, x, i, j. \\ \quad ((a \in \operatorname{mat}(m, n, x)) \wedge ((i \in \operatorname{interval}(1, n)) \wedge (j \in \operatorname{interval}(1, m)))) \Rightarrow  \\ \quad \quad \operatorname{entry}(\operatorname{transpose}(a, m, n), i, j) = \operatorname{entry}(a, j, i) \end{array}}}
\vnbprop{transpose-transpose}{\vnbstmt{\begin{array}{@{}l@{}} \forall a, m, n, x. \\ \quad (a \in \operatorname{mat}(m, n, x)) \Rightarrow (\operatorname{transpose}(\operatorname{transpose}(a, m, n), n, m) = a) \end{array}}}
\vnbprop{transpose-identity}{\vnbstmt{\begin{array}{@{}l@{}} \forall r, n. \\ \quad (\operatorname{is\text{-}ring}(r) \wedge (n \in \mathbb{N})) \Rightarrow (\operatorname{transpose}(\operatorname{identmat}(r, n), n, n) = \operatorname{identmat}(r, n)) \end{array}}}
\vnbprop{transpose-matmul}{\vnbstmt{\begin{array}{@{}l@{}} \forall r, m, n, c, p, q. \\ \quad \operatorname{is\text{-}commutative\text{-}ring}(r) \wedge \\ \quad p \in \operatorname{mat}(m, n, \operatorname{carr}(r)) \wedge \\ \quad q \in \operatorname{mat}(n, c, \operatorname{carr}(r)) \wedge \\ \quad (n = 0) \Rightarrow ((m = 0) \vee (c = 0)) \Rightarrow \\ \quad \quad \operatorname{transpose}(\operatorname{matmul}(r, p, q), m, c) = \operatorname{matmul}(r, \operatorname{transpose}(q, n, c), \operatorname{transpose}(p, m, n)) \end{array}}}
\vnbprop{dtt-ring-mul-neg1}{\vnbstmt{\begin{array}{@{}l@{}} \forall r.; \forall x \in \operatorname{carr}(r). \\ \quad \operatorname{is\text{-}commutative\text{-}ring}(r) \Rightarrow (\operatorname{mul}(r)(x, \operatorname{neg}(r)(\operatorname{one}(r))) = \operatorname{neg}(r)(x)) \end{array}}}
\vnbprop{dtt-ring-mul-one}{\vnbstmt{\begin{array}{@{}l@{}} \forall r.; \forall x \in \operatorname{carr}(r). \\ \quad \operatorname{is\text{-}commutative\text{-}ring}(r) \Rightarrow (\operatorname{mul}(r)(x, \operatorname{one}(r)) = x) \end{array}}}
\vnbprop{dtt-ring-neg-lin}{\vnbstmt{\begin{array}{@{}l@{}} \forall r, s, a, b. \\ \quad (\operatorname{is\text{-}commutative\text{-}ring}(r) \wedge ((s \in \operatorname{carr}(r)) \wedge ((a \in \operatorname{carr}(r)) \wedge (b \in \operatorname{carr}(r))))) \Rightarrow  \\ \quad \quad \operatorname{neg}(r)(\operatorname{add}(r)(\operatorname{mul}(r)(s, \operatorname{neg}(r)(a)), \operatorname{neg}(r)(b))) = \operatorname{add}(r)(\operatorname{mul}(r)(s, a), b) \end{array}}}
\vnbprop{dtt-ring-neg-neg}{\vnbstmt{\begin{array}{@{}l@{}} \forall r.; \forall x \in \operatorname{carr}(r). \\ \quad \operatorname{is\text{-}commutative\text{-}ring}(r) \Rightarrow (\operatorname{neg}(r)(\operatorname{neg}(r)(x)) = x) \end{array}}}
\vnbprop{dtt-ring-negx-x}{\vnbstmt{\begin{array}{@{}l@{}} \forall r.; \forall x \in \operatorname{carr}(r). \\ \quad \operatorname{is\text{-}commutative\text{-}ring}(r) \Rightarrow  \\ \quad \quad \operatorname{add}(r)(\operatorname{mul}(r)(\operatorname{neg}(r)(\operatorname{one}(r)), x), x) = \operatorname{zero}(r) \end{array}}}
\vnbprop{dtt-ring-zero-lin}{\vnbstmt{\begin{array}{@{}l@{}} \forall r. \\ \quad \operatorname{is\text{-}commutative\text{-}ring}(r) \Rightarrow  \\ \quad \quad \operatorname{zero}(r) = \operatorname{add}(r)(\operatorname{mul}(r)(\operatorname{neg}(r)(\operatorname{one}(r)), \operatorname{zero}(r)), \operatorname{zero}(r)) \end{array}}}
\vnbprop{dtt-ring-x-xneg}{\vnbstmt{\begin{array}{@{}l@{}} \forall r.; \forall x \in \operatorname{carr}(r). \\ \quad \operatorname{is\text{-}commutative\text{-}ring}(r) \Rightarrow  \\ \quad \quad \operatorname{add}(r)(x, \operatorname{mul}(r)(x, \operatorname{neg}(r)(\operatorname{one}(r)))) = \operatorname{zero}(r) \end{array}}}
\vnbprop{det-mul-elem-f-col}{\vnbstmt{\begin{array}{@{}l@{}} \forall r, n, a, k, l. \\ \quad \operatorname{is\text{-}commutative\text{-}ring}(r) \wedge \\ \quad n \in \mathbb{N} \wedge \\ \quad a \in \operatorname{mat}(n, n, \operatorname{carr}(r)) \wedge \\ \quad k \in \operatorname{interval}(1, n) \wedge \\ \quad l \in \operatorname{interval}(1, n) \wedge \\ \quad \neg (k = l) \Rightarrow \\ \quad \quad \operatorname{det}(r, n, \operatorname{matmul}(r, a, \operatorname{elem\text{-}f}(r, n, k, l))) = \operatorname{mul}(r)(\operatorname{det}(r, n, a), \operatorname{det}(r, n, \operatorname{elem\text{-}f}(r, n, k, l))) \end{array}}}
\vnbprop{det-elem-f-col}{\vnbstmt{\begin{array}{@{}l@{}} \forall r, n, a, k, l. \\ \quad \operatorname{is\text{-}commutative\text{-}ring}(r) \wedge \\ \quad n \in \mathbb{N} \wedge \\ \quad a \in \operatorname{mat}(n, n, \operatorname{carr}(r)) \wedge \\ \quad k \in \operatorname{interval}(1, n) \wedge \\ \quad l \in \operatorname{interval}(1, n) \wedge \\ \quad \neg (k = l) \Rightarrow \\ \quad \quad \operatorname{det}(r, n, \operatorname{matmul}(r, a, \operatorname{elem\text{-}f}(r, n, k, l))) = \operatorname{neg}(r)(\operatorname{det}(r, n, a)) \end{array}}}
\vnbprop{det-mul-elem-g-col}{\vnbstmt{\begin{array}{@{}l@{}} \forall r, n, a, c, k, l. \\ \quad \operatorname{is\text{-}commutative\text{-}ring}(r) \wedge \\ \quad n \in \mathbb{N} \wedge \\ \quad a \in \operatorname{mat}(n, n, \operatorname{carr}(r)) \wedge \\ \quad c \in \operatorname{carr}(r) \wedge \\ \quad k \in \operatorname{interval}(1, n) \wedge \\ \quad l \in \operatorname{interval}(1, n) \wedge \\ \quad \neg (k = l) \Rightarrow \\ \quad \quad \operatorname{det}(r, n, \operatorname{matmul}(r, a, \operatorname{elem\text{-}g}(r, n, c, k, l))) \\ \quad \quad =\; \operatorname{mul}(r)(\operatorname{det}(r, n, a), \operatorname{det}(r, n, \operatorname{elem\text{-}g}(r, n, c, k, l))) \end{array}}}
\vnbprop{det-elem-g-col}{\vnbstmt{\begin{array}{@{}l@{}} \forall r, n, a, c, k, l. \\ \quad \operatorname{is\text{-}commutative\text{-}ring}(r) \wedge \\ \quad n \in \mathbb{N} \wedge \\ \quad a \in \operatorname{mat}(n, n, \operatorname{carr}(r)) \wedge \\ \quad c \in \operatorname{carr}(r) \wedge \\ \quad k \in \operatorname{interval}(1, n) \wedge \\ \quad l \in \operatorname{interval}(1, n) \wedge \\ \quad \neg (k = l) \Rightarrow \\ \quad \quad \operatorname{det}(r, n, \operatorname{matmul}(r, a, \operatorname{elem\text{-}g}(r, n, c, k, l))) = \operatorname{det}(r, n, a) \end{array}}}
\vnbprop{det-mul-elem-h-col}{\vnbstmt{\begin{array}{@{}l@{}} \forall r, n, a, c, k. \\ \quad \operatorname{is\text{-}commutative\text{-}ring}(r) \wedge \\ \quad n \in \mathbb{N} \wedge \\ \quad a \in \operatorname{mat}(n, n, \operatorname{carr}(r)) \wedge \\ \quad c \in \operatorname{carr}(r) \wedge \\ \quad k \in \operatorname{interval}(1, n) \Rightarrow \\ \quad \quad \operatorname{det}(r, n, \operatorname{matmul}(r, a, \operatorname{elem\text{-}h}(r, n, c, k))) = \operatorname{mul}(r)(\operatorname{det}(r, n, a), \operatorname{det}(r, n, \operatorname{elem\text{-}h}(r, n, c, k))) \end{array}}}
\vnbprop{det-elem-h-col}{\vnbstmt{\begin{array}{@{}l@{}} \forall r, n, a, c, k. \\ \quad \operatorname{is\text{-}commutative\text{-}ring}(r) \wedge \\ \quad n \in \mathbb{N} \wedge \\ \quad a \in \operatorname{mat}(n, n, \operatorname{carr}(r)) \wedge \\ \quad c \in \operatorname{carr}(r) \wedge \\ \quad k \in \operatorname{interval}(1, n) \Rightarrow \\ \quad \quad \operatorname{det}(r, n, \operatorname{matmul}(r, a, \operatorname{elem\text{-}h}(r, n, c, k))) = \operatorname{mul}(r)(c, \operatorname{det}(r, n, a)) \end{array}}}
\vnbprop{dtt-minor1-colcongr}{\vnbstmt{\begin{array}{@{}l@{}} \forall r, n, x, y, p. \\ \quad \operatorname{is\text{-}ring}(r) \wedge \\ \quad n \in \mathbb{N} \wedge \\ \quad x \in \operatorname{mat}((n + 1), (n + 1), \operatorname{carr}(r)) \wedge \\ \quad y \in \operatorname{mat}((n + 1), (n + 1), \operatorname{carr}(r)) \wedge \\ \quad p = 1 \wedge \\ \quad \forall j, i \in \operatorname{interval}(1, (n + 1)). \\ \quad \quad \neg (j = p) \Rightarrow (\operatorname{entry}(x, i, j) = \operatorname{entry}(y, i, j)) \Rightarrow \\ \quad \quad \operatorname{minor}(x, 1, p, n) = \operatorname{minor}(y, 1, p, n) \end{array}}}
\vnbprop{dtt-col1-ind}{\vnbstmt{\begin{array}{@{}l@{}} \forall n \in \mathbb{N}, r, a, b, c, p, \mathit{dsc}. \\ \quad \operatorname{is\text{-}commutative\text{-}ring}(r) \wedge \\ \quad a \in \operatorname{mat}(n, n, \operatorname{carr}(r)) \wedge \\ \quad b \in \operatorname{mat}(n, n, \operatorname{carr}(r)) \wedge \\ \quad c \in \operatorname{mat}(n, n, \operatorname{carr}(r)) \wedge \\ \quad p \in \operatorname{interval}(1, n) \wedge \\ \quad p = 1 \wedge \\ \quad \mathit{dsc} \in \operatorname{carr}(r) \wedge \\ \quad \forall j, i \in \operatorname{interval}(1, n). \\ \quad \quad \neg (j = p) \Rightarrow (\operatorname{entry}(b, i, j) = \operatorname{entry}(a, i, j)) \wedge \\ \quad \forall j, i \in \operatorname{interval}(1, n). \\ \quad \quad \neg (j = p) \Rightarrow (\operatorname{entry}(c, i, j) = \operatorname{entry}(a, i, j)) \wedge \\ \quad \forall i \in \operatorname{interval}(1, n). \\ \quad \quad \operatorname{entry}(c, i, p) = \operatorname{add}(r)(\operatorname{mul}(r)(\mathit{dsc}, \operatorname{entry}(a, i, p)), \operatorname{entry}(b, i, p)) \Rightarrow \\ \quad \quad \operatorname{det}(r, n, c) = \operatorname{add}(r)(\operatorname{mul}(r)(\mathit{dsc}, \operatorname{det}(r, n, a)), \operatorname{det}(r, n, b)) \end{array}}}
\vnbprop{det-col-linear}{\vnbstmt{\begin{array}{@{}l@{}} \forall r, n, a, b, c, p, d. \\ \quad \operatorname{is\text{-}commutative\text{-}ring}(r) \wedge \\ \quad n \in \mathbb{N} \wedge \\ \quad a \in \operatorname{mat}(n, n, \operatorname{carr}(r)) \wedge \\ \quad b \in \operatorname{mat}(n, n, \operatorname{carr}(r)) \wedge \\ \quad c \in \operatorname{mat}(n, n, \operatorname{carr}(r)) \wedge \\ \quad p \in \operatorname{interval}(1, n) \wedge \\ \quad d \in \operatorname{carr}(r) \wedge \\ \quad \forall j, i \in \operatorname{interval}(1, n). \\ \quad \quad \neg (j = p) \Rightarrow (\operatorname{entry}(b, i, j) = \operatorname{entry}(a, i, j)) \wedge \\ \quad \forall j, i \in \operatorname{interval}(1, n). \\ \quad \quad \neg (j = p) \Rightarrow (\operatorname{entry}(c, i, j) = \operatorname{entry}(a, i, j)) \wedge \\ \quad \forall i \in \operatorname{interval}(1, n). \\ \quad \quad \operatorname{entry}(c, i, p) = \operatorname{add}(r)(\operatorname{mul}(r)(d, \operatorname{entry}(a, i, p)), \operatorname{entry}(b, i, p)) \Rightarrow \\ \quad \quad \operatorname{det}(r, n, c) = \operatorname{add}(r)(\operatorname{mul}(r)(d, \operatorname{det}(r, n, a)), \operatorname{det}(r, n, b)) \end{array}}}
\vnbprop{dtt-zero-col}{\vnbstmt{\begin{array}{@{}l@{}} \forall r, n, c, q. \\ \quad \operatorname{is\text{-}commutative\text{-}ring}(r) \wedge \\ \quad n \in \mathbb{N} \wedge \\ \quad c \in \operatorname{mat}(n, n, \operatorname{carr}(r)) \wedge \\ \quad q \in \operatorname{interval}(1, n) \wedge \\ \quad \forall i \in \operatorname{interval}(1, n).\; (\operatorname{entry}(c, i, q) = \operatorname{zero}(r)) \Rightarrow \\ \quad \quad \operatorname{det}(r, n, c) = \operatorname{zero}(r) \end{array}}}
\vnbprop{det-equal-cols-zero}{\vnbstmt{\begin{array}{@{}l@{}} \forall r, n, a, p, q. \\ \quad \operatorname{is\text{-}commutative\text{-}ring}(r) \wedge \\ \quad n \in \mathbb{N} \wedge \\ \quad a \in \operatorname{mat}(n, n, \operatorname{carr}(r)) \wedge \\ \quad p \in \operatorname{interval}(1, n) \wedge \\ \quad q \in \operatorname{interval}(1, n) \wedge \\ \quad \neg (p = q) \wedge \\ \quad \forall i \in \operatorname{interval}(1, n).\; (\operatorname{entry}(a, i, p) = \operatorname{entry}(a, i, q)) \Rightarrow \\ \quad \quad \operatorname{det}(r, n, a) = \operatorname{zero}(r) \end{array}}}
\vnbprop{det-swap-cols}{\vnbstmt{\begin{array}{@{}l@{}} \forall r, n, a, b, p, q. \\ \quad \operatorname{is\text{-}commutative\text{-}ring}(r) \wedge \\ \quad n \in \mathbb{N} \wedge \\ \quad a \in \operatorname{mat}(n, n, \operatorname{carr}(r)) \wedge \\ \quad b \in \operatorname{mat}(n, n, \operatorname{carr}(r)) \wedge \\ \quad p \in \operatorname{interval}(1, n) \wedge \\ \quad q \in \operatorname{interval}(1, n) \wedge \\ \quad \neg (p = q) \wedge \\ \quad \forall j, i \in \operatorname{interval}(1, n). \\ \quad \quad (\neg (j = p) \wedge \neg (j = q)) \Rightarrow (\operatorname{entry}(b, i, j) = \operatorname{entry}(a, i, j)) \wedge \\ \quad \forall i \in \operatorname{interval}(1, n).\; (\operatorname{entry}(b, i, p) = \operatorname{entry}(a, i, q)) \wedge \\ \quad \forall i \in \operatorname{interval}(1, n).\; (\operatorname{entry}(b, i, q) = \operatorname{entry}(a, i, p)) \Rightarrow \\ \quad \quad \operatorname{det}(r, n, b) = \operatorname{neg}(r)(\operatorname{det}(r, n, a)) \end{array}}}
\vnbprop{dtt-tr-fun}{\vnbstmt{\begin{array}{@{}l@{}} \forall r, n. \\ \quad (\operatorname{is\text{-}commutative\text{-}ring}(r) \wedge (n \in \mathbb{N})) \Rightarrow  \\ \quad \quad (\lambda x \in \operatorname{mat}(n, n, \operatorname{carr}(r)).\; \operatorname{det}(r, n, \operatorname{transpose}(x, n, n))) \in (\operatorname{mat}(n, n, \operatorname{carr}(r)) \to \operatorname{carr}(r)) \end{array}}}
\vnbprop{dtt-tr-lin}{\vnbstmt{\begin{array}{@{}l@{}} \forall r, n, a, b, c, p, s. \\ \quad \operatorname{is\text{-}commutative\text{-}ring}(r) \wedge \\ \quad n \in \mathbb{N} \wedge \\ \quad a \in \operatorname{mat}(n, n, \operatorname{carr}(r)) \wedge \\ \quad b \in \operatorname{mat}(n, n, \operatorname{carr}(r)) \wedge \\ \quad c \in \operatorname{mat}(n, n, \operatorname{carr}(r)) \wedge \\ \quad p \in \operatorname{interval}(1, n) \wedge \\ \quad s \in \operatorname{carr}(r) \wedge \\ \quad \forall i, k \in \operatorname{interval}(1, n). \\ \quad \quad \neg (i = p) \Rightarrow (\operatorname{entry}(b, i, k) = \operatorname{entry}(a, i, k)) \wedge \\ \quad \forall i, k \in \operatorname{interval}(1, n). \\ \quad \quad \neg (i = p) \Rightarrow (\operatorname{entry}(c, i, k) = \operatorname{entry}(a, i, k)) \wedge \\ \quad \forall k \in \operatorname{interval}(1, n). \\ \quad \quad \operatorname{entry}(c, p, k) = \operatorname{add}(r)(\operatorname{mul}(r)(s, \operatorname{entry}(a, p, k)), \operatorname{entry}(b, p, k)) \Rightarrow \\ \quad \quad (\lambda x \in \operatorname{mat}(n, n, \operatorname{carr}(r)).\; \operatorname{det}(r, n, \operatorname{transpose}(x, n, n)))(c) \\ \quad \quad =\; \operatorname{add}(r)(\operatorname{mul}(r)(s, (\lambda x \in \operatorname{mat}(n, n, \operatorname{carr}(r)).\; \operatorname{det}(r, n, \operatorname{transpose}(x, n, n)))(a)), \\ \quad \quad \quad (\lambda x \in \operatorname{mat}(n, n, \operatorname{carr}(r)).\; \operatorname{det}(r, n, \operatorname{transpose}(x, n, n)))(b)) \end{array}}}
\vnbprop{dtt-tr-alt}{\vnbstmt{\begin{array}{@{}l@{}} \forall r, n, a, p, q. \\ \quad \operatorname{is\text{-}commutative\text{-}ring}(r) \wedge \\ \quad n \in \mathbb{N} \wedge \\ \quad a \in \operatorname{mat}(n, n, \operatorname{carr}(r)) \wedge \\ \quad p \in \operatorname{interval}(1, n) \wedge \\ \quad q \in \operatorname{interval}(1, n) \wedge \\ \quad \neg (p = q) \wedge \\ \quad \forall k \in \operatorname{interval}(1, n).\; (\operatorname{entry}(a, p, k) = \operatorname{entry}(a, q, k)) \Rightarrow \\ \quad \quad (\lambda x \in \operatorname{mat}(n, n, \operatorname{carr}(r)).\; \operatorname{det}(r, n, \operatorname{transpose}(x, n, n)))(a) = \operatorname{zero}(r) \end{array}}}
\vnbprop{det-transpose}{\vnbstmt{\begin{array}{@{}l@{}} \forall r, n, a. \\ \quad (\operatorname{is\text{-}commutative\text{-}ring}(r) \wedge ((n \in \mathbb{N}) \wedge (a \in \operatorname{mat}(n, n, \operatorname{carr}(r))))) \Rightarrow  \\ \quad \quad \operatorname{det}(r, n, \operatorname{transpose}(a, n, n)) = \operatorname{det}(r, n, a) \end{array}}}
\vnbprop{dtt-minor-transpose}{\vnbstmt{\begin{array}{@{}l@{}} \forall r, n, s, p, q. \\ \quad (\operatorname{is\text{-}ring}(r) \wedge ((n \in \mathbb{N}) \wedge ((s \in \operatorname{mat}((n + 1), (n + 1), \operatorname{carr}(r))) \wedge ((p \in \mathbb{N}) \wedge (q \in \mathbb{N}))))) \Rightarrow  \\ \quad \quad \operatorname{minor}(\operatorname{transpose}(s, (n + 1), (n + 1)), q, p, n) = \operatorname{transpose}(\operatorname{minor}(s, p, q, n), n, n) \end{array}}}
\vnbprop{det-expand-col}{\vnbstmt{\begin{array}{@{}l@{}} \forall r, n, a, j. \\ \quad \operatorname{is\text{-}commutative\text{-}ring}(r) \wedge \\ \quad n \in \mathbb{N} \wedge \\ \quad a \in \operatorname{mat}((n + 1), (n + 1), \operatorname{carr}(r)) \wedge \\ \quad j \in \operatorname{interval}(1, (n + 1)) \Rightarrow \\ \quad \quad \operatorname{det}(r, (n + 1), a) \\ \quad \quad =\; \operatorname{finsum}(\operatorname{ring\text{-}additive\text{-}ag}(r), \\ \quad \quad \quad (\lambda e \in \operatorname{interval}(1, (n + 1)).\; \\ \quad \quad \quad \quad \operatorname{mul}(r)(\operatorname{mpow}(\operatorname{ring\text{-}multiplicative\text{-}monoid}(r), \operatorname{neg}(r)(\operatorname{one}(r)), (e + j)), \\ \quad \quad \quad \quad \quad \operatorname{mul}(r)(\operatorname{entry}(a, e, j), \operatorname{det}(r, n, \operatorname{minor}(a, e, j, n))))), \\ \quad \quad \quad \operatorname{interval}(1, (n + 1))) \end{array}}}

\vnbbookfile{theorem-library/clear-pivot-cross-proof}{theorem-library/clear-pivot-cross-proof}
\noindent{\small clear-pivot-cross -- over a euclidean ring, a matrix P (in MAT(succ p, succ q)) with a nonzero entry is \textasciitilde{}-equivalent to a matrix C2 in BORDER block form [[b,0],[0,C']] with b = C2\_\{1,1\} nonzero.  This is the per-level Smith step: place a class-minimal pivot at (1,1), clear its row and column, and read the}\par
\vnbprop{clear-pivot-cross}{\vnbstmt{\begin{array}{@{}l@{}} \forall a, p, q, m. \\ \quad \operatorname{is\text{-}euclidean\text{-}ring}(a) \wedge \\ \quad p \in \mathbb{N} \wedge \\ \quad q \in \mathbb{N} \wedge \\ \quad m \in \operatorname{mat}((p + 1), (q + 1), \operatorname{carr}(a)) \wedge \\ \quad \exists b \in \operatorname{interval}(1, (p + 1)), c \in \operatorname{interval}(1, (q + 1)). \\ \quad \quad \neg (\operatorname{entry}(m, b, c) = \operatorname{zero}(a)) \Rightarrow \\ \quad \quad \exists d. \\ \quad \quad \quad \operatorname{mat\text{-}equiv}(a, (p + 1), (q + 1), m, d) \wedge \\ \quad \quad \quad d = \operatorname{border}(a, \operatorname{entry}(d, 1, 1), \operatorname{submat}(d, p, q), p, q) \wedge \\ \quad \quad \quad \neg (\operatorname{entry}(d, 1, 1) = \operatorname{zero}(a)) \end{array}}}

\vnbbookfile{theorem-library/smith-diagonalization-proof}{theorem-library/smith-diagonalization-proof}
\noindent{\small smith-diagonalization -- every matrix P in MAT(m,n) over a euclidean ring is \textasciitilde{}-equivalent to a DIAGONAL matrix (Smith normal form, DIAGONALIZATION ONLY -- no divisibility chain).  The capstone of the LA arc.  Induction (ni) on the row dim k, inner statement INNER(k) = forall n P over MAT(k,n) exists diagonal}\par
\vnbprop{smith-diagonalization}{\vnbstmt{\begin{array}{@{}l@{}} \forall a.; \forall k \in \mathbb{N}, n, p. \\ \quad (\operatorname{is\text{-}euclidean\text{-}ring}(a) \wedge ((n \in \mathbb{N}) \wedge (p \in \operatorname{mat}(k, n, \operatorname{carr}(a))))) \Rightarrow  \\ \quad \quad \exists d \in \operatorname{mat}(k, n, \operatorname{carr}(a)). \\ \quad \quad \quad \operatorname{mat\text{-}equiv}(a, k, n, p, d) \wedge \operatorname{is\text{-}diagonal}(a, k, n, d) \end{array}}}

\vnbbookfile{theorem-library/smith-staircase-proof}{theorem-library/smith-staircase-proof}
\noindent{\small smith-staircase-proof.scm -- the Smith normal form with its nonzero diagonal entries recorded as an INITIAL SEGMENT.}\par
\vnbprop{border-staircase}{\vnbstmt{\begin{array}{@{}l@{}} \forall a, b, w, p, q, c. \\ \quad \operatorname{is\text{-}ring}(a) \wedge \\ \quad p \in \mathbb{N} \wedge \\ \quad q \in \mathbb{N} \wedge \\ \quad c \in \mathbb{N} \wedge \\ \quad b \in \operatorname{carr}(a) \wedge \\ \quad \neg (b = \operatorname{zero}(a)) \wedge \\ \quad \operatorname{smith\text{-}staircase}(a, p, q, w, c) \wedge \\ \quad w \in \operatorname{mat}(p, q, \operatorname{carr}(a)) \Rightarrow \\ \quad \quad \operatorname{smith\text{-}staircase}(a, (p + 1), (q + 1), \operatorname{border}(a, b, w, p, q), (c + 1)) \end{array}}}
\vnbprop{smith-normal-form}{\vnbstmt{\begin{array}{@{}l@{}} \forall a.; \forall k \in \mathbb{N}, n, p. \\ \quad (\operatorname{is\text{-}euclidean\text{-}ring}(a) \wedge ((n \in \mathbb{N}) \wedge (p \in \operatorname{mat}(k, n, \operatorname{carr}(a))))) \Rightarrow  \\ \quad \quad \exists d \in \operatorname{mat}(k, n, \operatorname{carr}(a)), b \in \mathbb{N}. \\ \quad \quad \quad \operatorname{mat\text{-}equiv}(a, k, n, p, d) \wedge \operatorname{smith\text{-}staircase}(a, k, n, d, b) \end{array}}}

\vnbbookfile{theorem-library/matact-row-linear-proof}{theorem-library/matact-row-linear-proof}
\noindent{\small matact-row-linear-proof.scm -- the coefficient row acts LINEARLY.}\par
\vnbprop{lincomb-unfold}{\vnbstmt{\begin{array}{@{}l@{}} \forall d, n, c, u. \\ \quad (\operatorname{lincomb}(d, n, c, u) \simeq \operatorname{finsum}(\operatorname{module\text{-}vector\text{-}ag}(d), (\lambda j \in \operatorname{interval}(1, n).\; \operatorname{act}(d)(\operatorname{entry}(c, 1, j), \operatorname{entry}(u, j, 1))), \operatorname{interval}(1, n))) \end{array}}}
\vnbprop{lincomb-type}{\vnbstmt{\begin{array}{@{}l@{}} \forall d, n, c, u. \\ \quad (\operatorname{is\text{-}module}(d) \wedge ((c \in \operatorname{mat}(1, n, \operatorname{carr}(\operatorname{scal}(d)))) \wedge (u \in \operatorname{mat}(n, 1, \operatorname{vec}(d))))) \Rightarrow  \\ \quad \quad \operatorname{lincomb}(d, n, c, u) \in \operatorname{vec}(d) \end{array}}}
\vnbprop{matact-lincomb}{\vnbstmt{\begin{array}{@{}l@{}} \forall d, n, c, u. \\ \quad (\operatorname{is\text{-}module}(d) \wedge ((c \in \operatorname{mat}(1, n, \operatorname{carr}(\operatorname{scal}(d)))) \wedge ((u \in \operatorname{mat}(n, 1, \operatorname{vec}(d))) \wedge (1 \leq n)))) \Rightarrow  \\ \quad \quad \operatorname{entry}(\operatorname{matact}(d, c, u), 1, 1) = \operatorname{lincomb}(d, n, c, u) \end{array}}}
\vnbprop{lincomb-row-add}{\vnbstmt{\begin{array}{@{}l@{}} \forall d, n, a, b, u. \\ \quad \operatorname{is\text{-}module}(d) \wedge \\ \quad a \in \operatorname{mat}(1, n, \operatorname{carr}(\operatorname{scal}(d))) \wedge \\ \quad b \in \operatorname{mat}(1, n, \operatorname{carr}(\operatorname{scal}(d))) \wedge \\ \quad u \in \operatorname{mat}(n, 1, \operatorname{vec}(d)) \Rightarrow \\ \quad \quad \operatorname{lincomb}(d, n, \operatorname{matadd}(\operatorname{scal}(d), a, b), u) = \operatorname{vadd}(d)(\operatorname{lincomb}(d, n, a, u), \operatorname{lincomb}(d, n, b, u)) \end{array}}}
\vnbprop{lincomb-row-scale}{\vnbstmt{\begin{array}{@{}l@{}} \forall d, n, r, c, u. \\ \quad (\operatorname{is\text{-}module}(d) \wedge ((r \in \operatorname{carr}(\operatorname{scal}(d))) \wedge ((c \in \operatorname{mat}(1, n, \operatorname{carr}(\operatorname{scal}(d)))) \wedge (u \in \operatorname{mat}(n, 1, \operatorname{vec}(d)))))) \Rightarrow  \\ \quad \quad \operatorname{lincomb}(d, n, \operatorname{matscale}(\operatorname{scal}(d), r, c), u) = \operatorname{act}(d)(r, \operatorname{lincomb}(d, n, c, u)) \end{array}}}

\vnbbookfile{theorem-library/span-bricks-proof}{theorem-library/span-bricks-proof}
\noindent{\small span-bricks-proof.scm -- the module-action bricks under spans-submodule-fg.}\par
\vnbprop{module-act-neg-one}{\vnbstmt{\begin{array}{@{}l@{}} \forall d.; \forall x \in \operatorname{vec}(d). \\ \quad \operatorname{is\text{-}module}(d) \Rightarrow (\operatorname{act}(d)(\operatorname{neg}(\operatorname{scal}(d))(\operatorname{one}(\operatorname{scal}(d))), x) = \operatorname{vneg}(d)(x)) \end{array}}}
\vnbprop{lincomb-zerorow}{\vnbstmt{\begin{array}{@{}l@{}} \forall d, n.; \forall u \in \operatorname{mat}(n, 1, \operatorname{vec}(d)). \\ \quad \operatorname{is\text{-}module}(d) \Rightarrow (\operatorname{lincomb}(d, n, \operatorname{zeromat}(\operatorname{scal}(d), 1, n), u) = \operatorname{vzero}(d)) \end{array}}}
\vnbprop{lincomb-unitrow}{\vnbstmt{\begin{array}{@{}l@{}} \forall d, n, i.; \forall u \in \operatorname{mat}(n, 1, \operatorname{vec}(d)). \\ \quad (\operatorname{is\text{-}module}(d) \wedge (i \in \operatorname{interval}(1, n))) \Rightarrow  \\ \quad \quad \operatorname{lincomb}(d, n, \operatorname{unitrow}(\operatorname{scal}(d), n, i), u) = \operatorname{entry}(u, i, 1) \end{array}}}
\vnbprop{lincomb-empty}{\vnbstmt{\forall d, c, u.\; (\operatorname{is\text{-}module}(d) \Rightarrow (\operatorname{lincomb}(d, 0, c, u) = \operatorname{vzero}(d)))}}
\vnbprop{span-is-submodule}{\vnbstmt{\begin{array}{@{}l@{}} \forall d, n.; \forall u \in \operatorname{mat}(n, 1, \operatorname{vec}(d)). \\ \quad \operatorname{is\text{-}module}(d) \Rightarrow \operatorname{is\text{-}submodule}(d, \operatorname{span}(d, n, u)) \end{array}}}
\vnbprop{spans-span}{\vnbstmt{\begin{array}{@{}l@{}} \forall d, n.; \forall u \in \operatorname{mat}(n, 1, \operatorname{vec}(d)). \\ \quad \operatorname{is\text{-}module}(d) \Rightarrow \operatorname{spans}(d, n, u, \operatorname{span}(d, n, u)) \end{array}}}

\vnbbookfile{theorem-library/span-bricks2-proof}{theorem-library/span-bricks2-proof}
\noindent{\small span-bricks2-proof.scm -- bricks 4, 5, 6 under spans-submodule-fg.}\par
\vnbprop{lincomb-row-peel}{\vnbstmt{\begin{array}{@{}l@{}} \forall d.; \forall n \in \mathbb{N}, c \in \operatorname{mat}(1, (n + 1), \operatorname{carr}(\operatorname{scal}(d))), u \in \operatorname{mat}((n + 1), 1, \operatorname{vec}(d)). \\ \quad \operatorname{is\text{-}module}(d) \Rightarrow  \\ \quad \quad \operatorname{lincomb}(d, (n + 1), c, u) \\ \quad \quad =\; \operatorname{vadd}(d)(\operatorname{lincomb}(d, n, \operatorname{block}(c, 1, n), \operatorname{block}(u, n, 1)), \\ \quad \quad \quad \operatorname{act}(d)(\operatorname{entry}(c, 1, (n + 1)), \operatorname{entry}(u, (n + 1), 1))) \end{array}}}
\vnbprop{lincomb-snoc}{\vnbstmt{\begin{array}{@{}l@{}} \forall d.; \forall n \in \mathbb{N}, c \in \operatorname{mat}(1, n, \operatorname{carr}(\operatorname{scal}(d))), u \in \operatorname{mat}(n, 1, \operatorname{vec}(d)), r \in \operatorname{carr}(\operatorname{scal}(d)), x \in \operatorname{vec}(d). \\ \quad \operatorname{is\text{-}module}(d) \Rightarrow  \\ \quad \quad \operatorname{lincomb}(d, (n + 1), \operatorname{snoc\text{-}row}(c, n, r), \operatorname{snoc\text{-}col}(u, n, x)) \\ \quad \quad =\; \operatorname{vadd}(d)(\operatorname{lincomb}(d, n, c, u), \operatorname{act}(d)(r, x)) \end{array}}}
\vnbprop{submodule-intersection}{\vnbstmt{\begin{array}{@{}l@{}} \forall d, a, b. \\ \quad (\operatorname{is\text{-}module}(d) \wedge (\operatorname{is\text{-}submodule}(d, a) \wedge \operatorname{is\text{-}submodule}(d, b))) \Rightarrow  \\ \quad \quad \operatorname{is\text{-}submodule}(d, (a \cap b)) \end{array}}}
\vnbprop{mat-1-0-nonempty}{\vnbstmt{\forall x.\; \exists p.\; (p \in \operatorname{mat}(1, 0, x))}}
\vnbprop{mat-0-1-nonempty}{\vnbstmt{\forall x.\; \exists p.\; (p \in \operatorname{mat}(0, 1, x))}}

\vnbbookfile{theorem-library/lastcoeff-ideal-proof}{theorem-library/lastcoeff-ideal-proof}
\noindent{\small lastcoeff-ideal-proof.scm -- the two self-contained lemmas the spans-submodule-fg descent leans on (no induction here).}\par
\vnbprop{lastcoeff-zero-in-span}{\vnbstmt{\begin{array}{@{}l@{}} \forall d.; \forall p \in \mathbb{N}, u \in \operatorname{mat}((p + 1), 1, \operatorname{vec}(d)), c \in \operatorname{mat}(1, (p + 1), \operatorname{carr}(\operatorname{scal}(d))). \\ \quad (\operatorname{is\text{-}module}(d) \wedge (\operatorname{entry}(c, 1, (p + 1)) = \operatorname{zero}(\operatorname{scal}(d)))) \Rightarrow  \\ \quad \quad \operatorname{lincomb}(d, (p + 1), c, u) \in \operatorname{span}(d, p, \operatorname{block}(u, p, 1)) \end{array}}}
\vnbprop{lastcoeff-set-is-ideal}{\vnbstmt{\begin{array}{@{}l@{}} \forall d.; \forall p \in \mathbb{N}, u \in \operatorname{mat}((p + 1), 1, \operatorname{vec}(d)), m. \\ \quad (\operatorname{is\text{-}module}(d) \wedge (\operatorname{is\text{-}commutative\text{-}ring}(\operatorname{scal}(d)) \wedge \operatorname{is\text{-}submodule}(d, m))) \Rightarrow  \\ \quad \quad \operatorname{is\text{-}ideal}(\operatorname{scal}(d), \operatorname{lastcoeff\text{-}set}(d, p, u, m)) \end{array}}}
\vnbprop{descent-remainder}{\vnbstmt{\begin{array}{@{}l@{}} \forall d.; \forall n \in \mathbb{N}, u \in \operatorname{mat}((n + 1), 1, \operatorname{vec}(d)), c, a \in \operatorname{mat}(1, (n + 1), \operatorname{carr}(\operatorname{scal}(d))), q, b \in \operatorname{carr}(\operatorname{scal}(d)). \\ \quad \operatorname{is\text{-}module}(d) \wedge \\ \quad \operatorname{is\text{-}commutative\text{-}ring}(\operatorname{scal}(d)) \wedge \\ \quad \operatorname{entry}(a, 1, (n + 1)) = b \wedge \\ \quad \operatorname{entry}(c, 1, (n + 1)) = \operatorname{mul}(\operatorname{scal}(d))(q, b) \Rightarrow \\ \quad \quad \exists y \in \operatorname{span}(d, n, \operatorname{block}(u, n, 1)). \\ \quad \quad \quad \operatorname{lincomb}(d, (n + 1), c, u) = \operatorname{vadd}(d)(y, \operatorname{act}(d)(q, \operatorname{lincomb}(d, (n + 1), a, u))) \wedge \\ \quad \quad \quad y = \operatorname{vadd}(d)(\operatorname{lincomb}(d, (n + 1), c, u), \operatorname{act}(d)(\operatorname{neg}(\operatorname{scal}(d))(q), \operatorname{lincomb}(d, (n + 1), a, u))) \end{array}}}

\vnbbookfile{theorem-library/spans-submodule-fg-proof}{theorem-library/spans-submodule-fg-proof}
\noindent{\small spans-submodule-fg-proof.scm -- THE DESCENT.}\par
\vnbprop{spans-fg-base}{\vnbstmt{\begin{array}{@{}l@{}} \forall d.; \forall u \in \operatorname{mat}(0, 1, \operatorname{vec}(d)), m, a \in \operatorname{power}(\operatorname{vec}(d)). \\ \quad \operatorname{is\text{-}module}(d) \wedge \\ \quad \operatorname{is\text{-}euclidean\text{-}ring}(\operatorname{scal}(d)) \wedge \\ \quad \operatorname{is\text{-}submodule}(d, m) \wedge \\ \quad \operatorname{spans}(d, 0, u, m) \wedge \\ \quad \operatorname{is\text{-}submodule}(d, a) \wedge \\ \quad a \subseteq m \Rightarrow \\ \quad \quad \exists k. \\ \quad \quad \quad ((k \in \mathbb{N}) \wedge (k \leq 0)) \wedge \exists w \in \operatorname{mat}(k, 1, \operatorname{vec}(d)).\; \operatorname{spans}(d, k, w, a) \end{array}}}
\vnbprop{spans-fg-step}{\vnbstmt{\begin{array}{@{}l@{}} \forall d.; \forall n \in \mathbb{N}, u \in \operatorname{mat}((n + 1), 1, \operatorname{vec}(d)), m, a \in \operatorname{power}(\operatorname{vec}(d)). \\ \quad \operatorname{is\text{-}module}(d) \wedge \\ \quad \operatorname{is\text{-}euclidean\text{-}ring}(\operatorname{scal}(d)) \wedge \\ \quad \forall u \in \operatorname{mat}(n, 1, \operatorname{vec}(d)), m, a \in \operatorname{power}(\operatorname{vec}(d)). \\ \quad \quad (\operatorname{is\text{-}submodule}(d, m) \wedge (\operatorname{spans}(d, n, u, m) \wedge (\operatorname{is\text{-}submodule}(d, a) \wedge (a \subseteq m)))) \Rightarrow  \\ \quad \quad \quad \exists k. \\ \quad \quad \quad \quad ((k \in \mathbb{N}) \wedge (k \leq n)) \wedge \exists w \in \operatorname{mat}(k, 1, \operatorname{vec}(d)).\; \operatorname{spans}(d, k, w, a) \wedge \\ \quad \operatorname{is\text{-}submodule}(d, m) \wedge \\ \quad \operatorname{spans}(d, (n + 1), u, m) \wedge \\ \quad \operatorname{is\text{-}submodule}(d, a) \wedge \\ \quad a \subseteq m \Rightarrow \\ \quad \quad \exists k. \\ \quad \quad \quad (((k \in \mathbb{N}) \wedge (k \leq (n + 1))) \wedge \exists w \in \operatorname{mat}(k, 1, \operatorname{vec}(d)).\; \operatorname{spans}(d, k, w, a)) \end{array}}}
\vnbprop{spans-submodule-fg}{\vnbstmt{\begin{array}{@{}l@{}} \forall d.; \forall n \in \mathbb{N}, u \in \operatorname{mat}(n, 1, \operatorname{vec}(d)), m, a \in \operatorname{power}(\operatorname{vec}(d)). \\ \quad \operatorname{is\text{-}module}(d) \wedge \\ \quad \operatorname{is\text{-}euclidean\text{-}ring}(\operatorname{scal}(d)) \wedge \\ \quad \operatorname{is\text{-}submodule}(d, m) \wedge \\ \quad \operatorname{spans}(d, n, u, m) \wedge \\ \quad \operatorname{is\text{-}submodule}(d, a) \wedge \\ \quad a \subseteq m \Rightarrow \\ \quad \quad \exists k. \\ \quad \quad \quad ((k \in \mathbb{N}) \wedge (k \leq n)) \wedge \exists w \in \operatorname{mat}(k, 1, \operatorname{vec}(d)).\; \operatorname{spans}(d, k, w, a) \end{array}}}

\vnbbookfile{theorem-library/submodule-fg-proof}{theorem-library/submodule-fg-proof}
\noindent{\small submodule-fg-proof.scm -- submodule-fg, PROVEN as a corollary of the SPANS-relativized statement it should always have had.}\par
\vnbprop{whole-module-is-submodule}{\vnbstmt{\forall d.\; (\operatorname{is\text{-}module}(d) \Rightarrow \operatorname{is\text{-}submodule}(d, \operatorname{vec}(d)))}}
\vnbprop{generates-implies-spans-vec}{\vnbstmt{\begin{array}{@{}l@{}} \forall d.; \forall n \in \mathbb{N}, u \in \operatorname{mat}(n, 1, \operatorname{vec}(d)). \\ \quad (\operatorname{is\text{-}module}(d) \wedge \operatorname{generates}(d, n, u)) \Rightarrow \operatorname{spans}(d, n, u, \operatorname{vec}(d)) \end{array}}}
\vnbprop{submodule-fg}{\vnbstmt{\begin{array}{@{}l@{}} \forall d.; \forall n \in \mathbb{N}, u \in \operatorname{mat}(n, 1, \operatorname{vec}(d)), m. \\ \quad ((\operatorname{is\text{-}module}(d) \wedge \operatorname{is\text{-}euclidean\text{-}ring}(\operatorname{scal}(d))) \wedge (\operatorname{generates}(d, n, u) \wedge \operatorname{is\text{-}submodule}(d, m))) \Rightarrow  \\ \quad \quad \exists k. \\ \quad \quad \quad ((k \in \mathbb{N}) \wedge (k \leq n)) \wedge \exists w \in \operatorname{mat}(k, 1, \operatorname{vec}(d)).\; \operatorname{spans}(d, k, w, m) \end{array}}}

\vnbbookfile{theorem-library/matact-assoc-proof}{theorem-library/matact-assoc-proof}
\noindent{\small matact-assoc-proof.scm -- Remark 3.39: the matrix action on a sequence of module elements is associative.}\par
\vnbprop{matact-triple-left}{\vnbstmt{\begin{array}{@{}l@{}} \forall d, m, n, k, l, p, q, u, w, a. \\ \quad \operatorname{is\text{-}module}(d) \wedge \\ \quad p \in \operatorname{mat}(m, n, \operatorname{carr}(\operatorname{scal}(d))) \wedge \\ \quad q \in \operatorname{mat}(n, k, \operatorname{carr}(\operatorname{scal}(d))) \wedge \\ \quad u \in \operatorname{mat}(k, l, \operatorname{vec}(d)) \wedge \\ \quad 1 \leq n \wedge \\ \quad 1 \leq k \wedge \\ \quad w \in \operatorname{interval}(1, m) \wedge \\ \quad a \in \operatorname{interval}(1, l) \Rightarrow \\ \quad \quad \operatorname{entry}(\operatorname{matact}(d, \operatorname{matmul}(\operatorname{scal}(d), p, q), u), w, a) \\ \quad \quad =\; \operatorname{finsum}(\operatorname{module\text{-}vector\text{-}ag}(d), \\ \quad \quad \quad (\lambda c \in \operatorname{interval}(1, k).\; \\ \quad \quad \quad \quad \operatorname{finsum}(\operatorname{module\text{-}vector\text{-}ag}(d), \\ \quad \quad \quad \quad \quad (\lambda j \in \operatorname{interval}(1, n).\; \\ \quad \quad \quad \quad \quad \quad (\lambda z \in (\operatorname{interval}(1, k) \times \operatorname{interval}(1, n)).\; \operatorname{act}(d)(\operatorname{mul}(\operatorname{scal}(d))(\operatorname{entry}(p, w, \operatorname{nth}(2, z)), \operatorname{entry}(q, \operatorname{nth}(2, z), \operatorname{nth}(1, z))), \operatorname{entry}(u, \operatorname{nth}(1, z), a)))({[}c, j])), \\ \quad \quad \quad \quad \quad \operatorname{interval}(1, n))), \\ \quad \quad \quad \operatorname{interval}(1, k)) \end{array}}}
\vnbprop{matact-triple-right}{\vnbstmt{\begin{array}{@{}l@{}} \forall d, m, n, k, l, p, q, u, w, a. \\ \quad \operatorname{is\text{-}module}(d) \wedge \\ \quad p \in \operatorname{mat}(m, n, \operatorname{carr}(\operatorname{scal}(d))) \wedge \\ \quad q \in \operatorname{mat}(n, k, \operatorname{carr}(\operatorname{scal}(d))) \wedge \\ \quad u \in \operatorname{mat}(k, l, \operatorname{vec}(d)) \wedge \\ \quad 1 \leq n \wedge \\ \quad 1 \leq k \wedge \\ \quad w \in \operatorname{interval}(1, m) \wedge \\ \quad a \in \operatorname{interval}(1, l) \Rightarrow \\ \quad \quad \operatorname{entry}(\operatorname{matact}(d, p, \operatorname{matact}(d, q, u)), w, a) \\ \quad \quad =\; \operatorname{finsum}(\operatorname{module\text{-}vector\text{-}ag}(d), \\ \quad \quad \quad (\lambda j \in \operatorname{interval}(1, n).\; \\ \quad \quad \quad \quad \operatorname{finsum}(\operatorname{module\text{-}vector\text{-}ag}(d), \\ \quad \quad \quad \quad \quad (\lambda c \in \operatorname{interval}(1, k).\; \\ \quad \quad \quad \quad \quad \quad (\lambda z \in (\operatorname{interval}(1, k) \times \operatorname{interval}(1, n)).\; \operatorname{act}(d)(\operatorname{mul}(\operatorname{scal}(d))(\operatorname{entry}(p, w, \operatorname{nth}(2, z)), \operatorname{entry}(q, \operatorname{nth}(2, z), \operatorname{nth}(1, z))), \operatorname{entry}(u, \operatorname{nth}(1, z), a)))({[}c, j])), \\ \quad \quad \quad \quad \quad \operatorname{interval}(1, k))), \\ \quad \quad \quad \operatorname{interval}(1, n)) \end{array}}}
\vnbprop{matact-assoc}{\vnbstmt{\begin{array}{@{}l@{}} \forall d, m, n, k, l, p, q, u. \\ \quad \operatorname{is\text{-}module}(d) \wedge \\ \quad p \in \operatorname{mat}(m, n, \operatorname{carr}(\operatorname{scal}(d))) \wedge \\ \quad q \in \operatorname{mat}(n, k, \operatorname{carr}(\operatorname{scal}(d))) \wedge \\ \quad u \in \operatorname{mat}(k, l, \operatorname{vec}(d)) \wedge \\ \quad (n = 0) \Rightarrow ((m = 0) \vee ((k = 0) \wedge (l = 0))) \wedge \\ \quad (k = 0) \Rightarrow ((l = 0) \vee ((n = 0) \wedge (m = 0))) \Rightarrow \\ \quad \quad \operatorname{matact}(d, \operatorname{matmul}(\operatorname{scal}(d), p, q), u) = \operatorname{matact}(d, p, \operatorname{matact}(d, q, u)) \end{array}}}

\vnbbookfile{theorem-library/mod-basis-proof}{theorem-library/mod-basis-proof}
\noindent{\small mod-basis-proof.scm -- Lemma 3.40 (algebraic-numbers.pdf ch.3 sec 8.2): an invertible matrix carries generating sequences to generating sequences and relation-free sequences to relation-free sequences.  Plus the identity action I . u = u that says MATACT is a genuine action.}\par
\vnbprop{matact-identmat}{\vnbstmt{\begin{array}{@{}l@{}} \forall d, n, q.; \forall u \in \operatorname{mat}(n, q, \operatorname{vec}(d)). \\ \quad \operatorname{is\text{-}module}(d) \Rightarrow (\operatorname{matact}(d, \operatorname{identmat}(\operatorname{scal}(d), n), u) = u) \end{array}}}
\vnbprop{generates-transport}{\vnbstmt{\begin{array}{@{}l@{}} \forall d, n, u, m. \\ \quad (\operatorname{is\text{-}module}(d) \wedge ((u \in \operatorname{mat}(n, 1, \operatorname{vec}(d))) \wedge (\operatorname{is\text{-}invertible\text{-}mat}(\operatorname{scal}(d), n, m) \wedge \operatorname{generates}(d, n, u)))) \Rightarrow  \\ \quad \quad \operatorname{generates}(d, n, \operatorname{matact}(d, m, u)) \end{array}}}
\vnbprop{free-transport}{\vnbstmt{\begin{array}{@{}l@{}} \forall d, m, v, a. \\ \quad (\operatorname{is\text{-}module}(d) \wedge ((v \in \operatorname{mat}(m, 1, \operatorname{vec}(d))) \wedge (\operatorname{is\text{-}invertible\text{-}mat}(\operatorname{scal}(d), m, a) \wedge \operatorname{rel\text{-}free}(d, m, v)))) \Rightarrow  \\ \quad \quad \operatorname{rel\text{-}free}(d, m, \operatorname{matact}(d, a, v)) \end{array}}}

\vnbbookfile{theorem-library/inverse-invertible-proof}{theorem-library/inverse-invertible-proof}
\noindent{\small inverse-invertible-proof.scm -- the inverse of an invertible matrix is invertible.}\par
\vnbprop{inverse-is-invertible}{\vnbstmt{\begin{array}{@{}l@{}} \forall a, n, m, b. \\ \quad \operatorname{is\text{-}ring}(a) \wedge \\ \quad m \in \operatorname{mat}(n, n, \operatorname{carr}(a)) \wedge \\ \quad b \in \operatorname{mat}(n, n, \operatorname{carr}(a)) \wedge \\ \quad \operatorname{matmul}(a, m, b) = \operatorname{identmat}(a, n) \wedge \\ \quad \operatorname{matmul}(a, b, m) = \operatorname{identmat}(a, n) \Rightarrow \\ \quad \quad \operatorname{is\text{-}invertible\text{-}mat}(a, n, b) \end{array}}}

\vnbbookfile{theorem-library/la-det-adjoint}{theorem-library/la-det-adjoint}
\noindent{\small la-det-adjoint.scm -- the triangular determinant (linear-algebra 3.23), the formal adjoint (3.25: adj(A) A = det(A) I = A adj(A)) and invertibility over a field (3.26: A invertible iff det A /= 0).  LA-DET, 2026-10-18.  The pedestrian route of the notes: cofactor expansion only (det-expand-row / det-expand-col, ...}\par
\vnbprop{ladet-nn-lt-of-le-ne}{\vnbstmt{\forall a, b.\; (((a \in \mathbb{N}) \wedge ((b \in \mathbb{N}) \wedge ((a \leq b) \wedge \neg (a = b)))) \Rightarrow (a < b))}}
\vnbprop{ladet-sign-even}{\vnbstmt{\begin{array}{@{}l@{}} \forall a \in \mathbb{N}, r. \\ \quad \operatorname{is\text{-}commutative\text{-}ring}(r) \Rightarrow  \\ \quad \quad \operatorname{mpow}(\operatorname{ring\text{-}multiplicative\text{-}monoid}(r), \operatorname{neg}(r)(\operatorname{one}(r)), (a + a)) = \operatorname{one}(r) \end{array}}}
\vnbprop{lower-triangular-entry}{\vnbstmt{\begin{array}{@{}l@{}} \forall a, n, d.; \forall i, j \in \operatorname{interval}(1, n). \\ \quad (\operatorname{is\text{-}lower\text{-}triangular}(a, n, d) \wedge (i < j)) \Rightarrow (\operatorname{entry}(d, i, j) = \operatorname{zero}(a)) \end{array}}}
\vnbprop{upper-triangular-entry}{\vnbstmt{\begin{array}{@{}l@{}} \forall a, n, d.; \forall i, j \in \operatorname{interval}(1, n). \\ \quad (\operatorname{is\text{-}upper\text{-}triangular}(a, n, d) \wedge (j < i)) \Rightarrow (\operatorname{entry}(d, i, j) = \operatorname{zero}(a)) \end{array}}}
\vnbprop{ladet-minor-last-entry}{\vnbstmt{\begin{array}{@{}l@{}} \forall r, n, s, u, v. \\ \quad \operatorname{is\text{-}ring}(r) \wedge \\ \quad n \in \mathbb{N} \wedge \\ \quad s \in \operatorname{mat}((n + 1), (n + 1), \operatorname{carr}(r)) \wedge \\ \quad u \in \operatorname{interval}(1, n) \wedge \\ \quad v \in \operatorname{interval}(1, n) \Rightarrow \\ \quad \quad \operatorname{entry}(\operatorname{minor}(s, (n + 1), (n + 1), n), u, v) = \operatorname{entry}(s, u, v) \end{array}}}
\vnbprop{ladet-minor-lower}{\vnbstmt{\begin{array}{@{}l@{}} \forall r, n, s. \\ \quad \operatorname{is\text{-}ring}(r) \wedge \\ \quad n \in \mathbb{N} \wedge \\ \quad s \in \operatorname{mat}((n + 1), (n + 1), \operatorname{carr}(r)) \wedge \\ \quad \operatorname{is\text{-}lower\text{-}triangular}(r, (n + 1), s) \Rightarrow \\ \quad \quad \operatorname{is\text{-}lower\text{-}triangular}(r, n, \operatorname{minor}(s, (n + 1), (n + 1), n)) \end{array}}}
\vnbprop{ladet-minor-upper}{\vnbstmt{\begin{array}{@{}l@{}} \forall r, n, s. \\ \quad \operatorname{is\text{-}ring}(r) \wedge \\ \quad n \in \mathbb{N} \wedge \\ \quad s \in \operatorname{mat}((n + 1), (n + 1), \operatorname{carr}(r)) \wedge \\ \quad \operatorname{is\text{-}upper\text{-}triangular}(r, (n + 1), s) \Rightarrow \\ \quad \quad \operatorname{is\text{-}upper\text{-}triangular}(r, n, \operatorname{minor}(s, (n + 1), (n + 1), n)) \end{array}}}
\vnbprop{ladet-det-lower-ind}{\vnbstmt{\begin{array}{@{}l@{}} \forall n \in \mathbb{N}, r, d. \\ \quad (\operatorname{is\text{-}commutative\text{-}ring}(r) \wedge ((d \in \operatorname{mat}(n, n, \operatorname{carr}(r))) \wedge \operatorname{is\text{-}lower\text{-}triangular}(r, n, d))) \Rightarrow  \\ \quad \quad \operatorname{det}(r, n, d) = \operatorname{prod\text{-}ring}(r, (\lambda k \in \operatorname{interval}(1, n).\; \operatorname{entry}(d, k, k)), \operatorname{interval}(1, n)) \end{array}}}
\vnbprop{ladet-det-upper-ind}{\vnbstmt{\begin{array}{@{}l@{}} \forall n \in \mathbb{N}, r, d. \\ \quad (\operatorname{is\text{-}commutative\text{-}ring}(r) \wedge ((d \in \operatorname{mat}(n, n, \operatorname{carr}(r))) \wedge \operatorname{is\text{-}upper\text{-}triangular}(r, n, d))) \Rightarrow  \\ \quad \quad \operatorname{det}(r, n, d) = \operatorname{prod\text{-}ring}(r, (\lambda k \in \operatorname{interval}(1, n).\; \operatorname{entry}(d, k, k)), \operatorname{interval}(1, n)) \end{array}}}
\vnbprop{det-lower-triangular}{\vnbstmt{\begin{array}{@{}l@{}} \forall r, n, d. \\ \quad (\operatorname{is\text{-}commutative\text{-}ring}(r) \wedge ((n \in \mathbb{N}) \wedge ((d \in \operatorname{mat}(n, n, \operatorname{carr}(r))) \wedge \operatorname{is\text{-}lower\text{-}triangular}(r, n, d)))) \Rightarrow  \\ \quad \quad \operatorname{det}(r, n, d) = \operatorname{prod\text{-}ring}(r, (\lambda k \in \operatorname{interval}(1, n).\; \operatorname{entry}(d, k, k)), \operatorname{interval}(1, n)) \end{array}}}
\vnbprop{det-upper-triangular}{\vnbstmt{\begin{array}{@{}l@{}} \forall r, n, d. \\ \quad (\operatorname{is\text{-}commutative\text{-}ring}(r) \wedge ((n \in \mathbb{N}) \wedge ((d \in \operatorname{mat}(n, n, \operatorname{carr}(r))) \wedge \operatorname{is\text{-}upper\text{-}triangular}(r, n, d)))) \Rightarrow  \\ \quad \quad \operatorname{det}(r, n, d) = \operatorname{prod\text{-}ring}(r, (\lambda k \in \operatorname{interval}(1, n).\; \operatorname{entry}(d, k, k)), \operatorname{interval}(1, n)) \end{array}}}
\vnbprop{ladet-skip-ne}{\vnbstmt{\forall w, c.\; (((w \in \mathbb{N}) \wedge (c \in \mathbb{N})) \Rightarrow \neg (\left(\text{if }(w < c)\text{ then }w\text{ else }(w + 1)\right) = c))}}
\vnbprop{ladet-minor-rowcongr}{\vnbstmt{\begin{array}{@{}l@{}} \forall r, n, x, y, p, q. \\ \quad \operatorname{is\text{-}ring}(r) \wedge \\ \quad n \in \mathbb{N} \wedge \\ \quad x \in \operatorname{mat}((n + 1), (n + 1), \operatorname{carr}(r)) \wedge \\ \quad y \in \operatorname{mat}((n + 1), (n + 1), \operatorname{carr}(r)) \wedge \\ \quad p \in \mathbb{N} \wedge \\ \quad q \in \mathbb{N} \wedge \\ \quad \forall i, k \in \operatorname{interval}(1, (n + 1)). \\ \quad \quad \neg (i = p) \Rightarrow (\operatorname{entry}(x, i, k) = \operatorname{entry}(y, i, k)) \Rightarrow \\ \quad \quad \operatorname{minor}(x, p, q, n) = \operatorname{minor}(y, p, q, n) \end{array}}}
\vnbprop{ladet-minor-colcongr}{\vnbstmt{\begin{array}{@{}l@{}} \forall r, n, x, y, p, q. \\ \quad \operatorname{is\text{-}ring}(r) \wedge \\ \quad n \in \mathbb{N} \wedge \\ \quad x \in \operatorname{mat}((n + 1), (n + 1), \operatorname{carr}(r)) \wedge \\ \quad y \in \operatorname{mat}((n + 1), (n + 1), \operatorname{carr}(r)) \wedge \\ \quad p \in \mathbb{N} \wedge \\ \quad q \in \mathbb{N} \wedge \\ \quad \forall j, i \in \operatorname{interval}(1, (n + 1)). \\ \quad \quad \neg (j = q) \Rightarrow (\operatorname{entry}(x, i, j) = \operatorname{entry}(y, i, j)) \Rightarrow \\ \quad \quad \operatorname{minor}(x, p, q, n) = \operatorname{minor}(y, p, q, n) \end{array}}}
\vnbprop{ladet-ring-r1}{\vnbstmt{\begin{array}{@{}l@{}} \forall r, s, x, d. \\ \quad (\operatorname{is\text{-}commutative\text{-}ring}(r) \wedge ((s \in \operatorname{carr}(r)) \wedge ((x \in \operatorname{carr}(r)) \wedge (d \in \operatorname{carr}(r))))) \Rightarrow  \\ \quad \quad \operatorname{mul}(r)(s, \operatorname{mul}(r)(x, d)) = \operatorname{mul}(r)(x, \operatorname{mul}(r)(s, d)) \end{array}}}
\vnbprop{ladet-ring-c1}{\vnbstmt{\begin{array}{@{}l@{}} \forall r, s, x, d. \\ \quad (\operatorname{is\text{-}commutative\text{-}ring}(r) \wedge ((s \in \operatorname{carr}(r)) \wedge ((x \in \operatorname{carr}(r)) \wedge (d \in \operatorname{carr}(r))))) \Rightarrow  \\ \quad \quad \operatorname{mul}(r)(s, \operatorname{mul}(r)(x, d)) = \operatorname{mul}(r)(\operatorname{mul}(r)(s, d), x) \end{array}}}
\vnbprop{ladet-row-diag}{\vnbstmt{\begin{array}{@{}l@{}} \forall r, n, a, k. \\ \quad \operatorname{is\text{-}commutative\text{-}ring}(r) \wedge \\ \quad n \in \mathbb{N} \wedge \\ \quad a \in \operatorname{mat}((n + 1), (n + 1), \operatorname{carr}(r)) \wedge \\ \quad k \in \operatorname{interval}(1, (n + 1)) \Rightarrow \\ \quad \quad \operatorname{finsum}(\operatorname{ring\text{-}additive\text{-}ag}(r), \\ \quad \quad \quad (\lambda j \in \operatorname{interval}(1, (n + 1)).\; \\ \quad \quad \quad \quad \operatorname{mul}(r)(\operatorname{entry}(a, k, j), \\ \quad \quad \quad \quad \quad \operatorname{mul}(r)(\operatorname{mpow}(\operatorname{ring\text{-}multiplicative\text{-}monoid}(r), \operatorname{neg}(r)(\operatorname{one}(r)), (j + k)), \operatorname{det}(r, n, \operatorname{minor}(a, k, j, n))))), \\ \quad \quad \quad \operatorname{interval}(1, (n + 1))) \\ \quad \quad =\; \operatorname{det}(r, (n + 1), a) \end{array}}}
\vnbprop{ladet-col-diag}{\vnbstmt{\begin{array}{@{}l@{}} \forall r, n, a, k. \\ \quad \operatorname{is\text{-}commutative\text{-}ring}(r) \wedge \\ \quad n \in \mathbb{N} \wedge \\ \quad a \in \operatorname{mat}((n + 1), (n + 1), \operatorname{carr}(r)) \wedge \\ \quad k \in \operatorname{interval}(1, (n + 1)) \Rightarrow \\ \quad \quad \operatorname{finsum}(\operatorname{ring\text{-}additive\text{-}ag}(r), \\ \quad \quad \quad (\lambda j \in \operatorname{interval}(1, (n + 1)).\; \\ \quad \quad \quad \quad \operatorname{mul}(r)(\operatorname{mul}(r)(\operatorname{mpow}(\operatorname{ring\text{-}multiplicative\text{-}monoid}(r), \operatorname{neg}(r)(\operatorname{one}(r)), (k + j)), \operatorname{det}(r, n, \operatorname{minor}(a, j, k, n))), \\ \quad \quad \quad \quad \quad \operatorname{entry}(a, j, k))), \\ \quad \quad \quad \operatorname{interval}(1, (n + 1))) \\ \quad \quad =\; \operatorname{det}(r, (n + 1), a) \end{array}}}
\vnbprop{ladet-row-off}{\vnbstmt{\begin{array}{@{}l@{}} \forall r, n, a, i, k. \\ \quad \operatorname{is\text{-}commutative\text{-}ring}(r) \wedge \\ \quad n \in \mathbb{N} \wedge \\ \quad a \in \operatorname{mat}((n + 1), (n + 1), \operatorname{carr}(r)) \wedge \\ \quad i \in \operatorname{interval}(1, (n + 1)) \wedge \\ \quad k \in \operatorname{interval}(1, (n + 1)) \wedge \\ \quad \neg (i = k) \Rightarrow \\ \quad \quad \operatorname{finsum}(\operatorname{ring\text{-}additive\text{-}ag}(r), \\ \quad \quad \quad (\lambda j \in \operatorname{interval}(1, (n + 1)).\; \\ \quad \quad \quad \quad \operatorname{mul}(r)(\operatorname{entry}(a, i, j), \\ \quad \quad \quad \quad \quad \operatorname{mul}(r)(\operatorname{mpow}(\operatorname{ring\text{-}multiplicative\text{-}monoid}(r), \operatorname{neg}(r)(\operatorname{one}(r)), (j + k)), \operatorname{det}(r, n, \operatorname{minor}(a, k, j, n))))), \\ \quad \quad \quad \operatorname{interval}(1, (n + 1))) \\ \quad \quad =\; \operatorname{zero}(r) \end{array}}}
\vnbprop{ladet-col-off}{\vnbstmt{\begin{array}{@{}l@{}} \forall r, n, a, i, k. \\ \quad \operatorname{is\text{-}commutative\text{-}ring}(r) \wedge \\ \quad n \in \mathbb{N} \wedge \\ \quad a \in \operatorname{mat}((n + 1), (n + 1), \operatorname{carr}(r)) \wedge \\ \quad i \in \operatorname{interval}(1, (n + 1)) \wedge \\ \quad k \in \operatorname{interval}(1, (n + 1)) \wedge \\ \quad \neg (i = k) \Rightarrow \\ \quad \quad \operatorname{finsum}(\operatorname{ring\text{-}additive\text{-}ag}(r), \\ \quad \quad \quad (\lambda j \in \operatorname{interval}(1, (n + 1)).\; \\ \quad \quad \quad \quad \operatorname{mul}(r)(\operatorname{mul}(r)(\operatorname{mpow}(\operatorname{ring\text{-}multiplicative\text{-}monoid}(r), \operatorname{neg}(r)(\operatorname{one}(r)), (i + j)), \operatorname{det}(r, n, \operatorname{minor}(a, j, i, n))), \\ \quad \quad \quad \quad \quad \operatorname{entry}(a, j, k))), \\ \quad \quad \quad \operatorname{interval}(1, (n + 1))) \\ \quad \quad =\; \operatorname{zero}(r) \end{array}}}
\vnbprop{adj-type}{\vnbstmt{\begin{array}{@{}l@{}} \forall r, n, a. \\ \quad (\operatorname{is\text{-}ring}(r) \wedge ((n \in \mathbb{N}) \wedge (a \in \operatorname{mat}((n + 1), (n + 1), \operatorname{carr}(r))))) \Rightarrow  \\ \quad \quad \operatorname{adj}(r, n, a) \in \operatorname{mat}((n + 1), (n + 1), \operatorname{carr}(r)) \end{array}}}
\vnbprop{adj-entry}{\vnbstmt{\begin{array}{@{}l@{}} \forall r, n, a, i, j. \\ \quad \operatorname{is\text{-}ring}(r) \wedge \\ \quad n \in \mathbb{N} \wedge \\ \quad a \in \operatorname{mat}((n + 1), (n + 1), \operatorname{carr}(r)) \wedge \\ \quad i \in \operatorname{interval}(1, (n + 1)) \wedge \\ \quad j \in \operatorname{interval}(1, (n + 1)) \Rightarrow \\ \quad \quad \operatorname{entry}(\operatorname{adj}(r, n, a), i, j) \\ \quad \quad =\; \operatorname{mul}(r)(\operatorname{mpow}(\operatorname{ring\text{-}multiplicative\text{-}monoid}(r), \operatorname{neg}(r)(\operatorname{one}(r)), (i + j)), \operatorname{det}(r, n, \operatorname{minor}(a, j, i, n))) \end{array}}}
\vnbprop{adj-mul}{\vnbstmt{\begin{array}{@{}l@{}} \forall r, n, a. \\ \quad (\operatorname{is\text{-}commutative\text{-}ring}(r) \wedge ((n \in \mathbb{N}) \wedge (a \in \operatorname{mat}((n + 1), (n + 1), \operatorname{carr}(r))))) \Rightarrow  \\ \quad \quad \operatorname{matmul}(r, \operatorname{adj}(r, n, a), a) = \operatorname{matscale}(r, \operatorname{det}(r, (n + 1), a), \operatorname{identmat}(r, (n + 1))) \end{array}}}
\vnbprop{mul-adj}{\vnbstmt{\begin{array}{@{}l@{}} \forall r, n, a. \\ \quad (\operatorname{is\text{-}commutative\text{-}ring}(r) \wedge ((n \in \mathbb{N}) \wedge (a \in \operatorname{mat}((n + 1), (n + 1), \operatorname{carr}(r))))) \Rightarrow  \\ \quad \quad \operatorname{matmul}(r, a, \operatorname{adj}(r, n, a)) = \operatorname{matscale}(r, \operatorname{det}(r, (n + 1), a), \operatorname{identmat}(r, (n + 1))) \end{array}}}
\vnbprop{matscale-one}{\vnbstmt{\begin{array}{@{}l@{}} \forall r, m, k, x. \\ \quad (\operatorname{is\text{-}ring}(r) \wedge (x \in \operatorname{mat}(m, k, \operatorname{carr}(r)))) \Rightarrow (\operatorname{matscale}(r, \operatorname{one}(r), x) = x) \end{array}}}
\vnbprop{matscale-matscale}{\vnbstmt{\begin{array}{@{}l@{}} \forall r, m, k, x, b, d. \\ \quad (\operatorname{is\text{-}ring}(r) \wedge ((x \in \operatorname{mat}(m, k, \operatorname{carr}(r))) \wedge ((b \in \operatorname{carr}(r)) \wedge (d \in \operatorname{carr}(r))))) \Rightarrow  \\ \quad \quad \operatorname{matscale}(r, b, \operatorname{matscale}(r, d, x)) = \operatorname{matscale}(r, \operatorname{mul}(r)(b, d), x) \end{array}}}
\vnbprop{matmul-matscale-right}{\vnbstmt{\begin{array}{@{}l@{}} \forall r, n, p, q, c. \\ \quad \operatorname{is\text{-}commutative\text{-}ring}(r) \wedge \\ \quad n \in \mathbb{N} \wedge \\ \quad p \in \operatorname{mat}((n + 1), (n + 1), \operatorname{carr}(r)) \wedge \\ \quad q \in \operatorname{mat}((n + 1), (n + 1), \operatorname{carr}(r)) \wedge \\ \quad c \in \operatorname{carr}(r) \Rightarrow \\ \quad \quad \operatorname{matmul}(r, p, \operatorname{matscale}(r, c, q)) = \operatorname{matscale}(r, c, \operatorname{matmul}(r, p, q)) \end{array}}}
\vnbprop{matmul-matscale-left}{\vnbstmt{\begin{array}{@{}l@{}} \forall r, n, p, q, c. \\ \quad \operatorname{is\text{-}commutative\text{-}ring}(r) \wedge \\ \quad n \in \mathbb{N} \wedge \\ \quad p \in \operatorname{mat}((n + 1), (n + 1), \operatorname{carr}(r)) \wedge \\ \quad q \in \operatorname{mat}((n + 1), (n + 1), \operatorname{carr}(r)) \wedge \\ \quad c \in \operatorname{carr}(r) \Rightarrow \\ \quad \quad \operatorname{matmul}(r, \operatorname{matscale}(r, c, p), q) = \operatorname{matscale}(r, c, \operatorname{matmul}(r, p, q)) \end{array}}}
\vnbprop{ladet-field-inverse}{\vnbstmt{\begin{array}{@{}l@{}} \forall r.; \forall x \in \operatorname{carr}(r). \\ \quad (\operatorname{is\text{-}field\text{-}ring}(r) \wedge \neg (x = \operatorname{zero}(r))) \Rightarrow  \\ \quad \quad \exists b \in \operatorname{carr}(r).\; (\operatorname{mul}(r)(x, b) = \operatorname{one}(r)) \end{array}}}
\vnbprop{cramer-inverse}{\vnbstmt{\begin{array}{@{}l@{}} \forall r, n, a, b. \\ \quad \operatorname{is\text{-}commutative\text{-}ring}(r) \wedge \\ \quad n \in \mathbb{N} \wedge \\ \quad a \in \operatorname{mat}((n + 1), (n + 1), \operatorname{carr}(r)) \wedge \\ \quad b \in \operatorname{carr}(r) \wedge \\ \quad \operatorname{mul}(r)(\operatorname{det}(r, (n + 1), a), b) = \operatorname{one}(r) \Rightarrow \\ \quad \quad \operatorname{matmul}(r, a, \operatorname{matscale}(r, b, \operatorname{adj}(r, n, a))) = \operatorname{identmat}(r, (n + 1)) \wedge \\ \quad \quad \operatorname{matmul}(r, \operatorname{matscale}(r, b, \operatorname{adj}(r, n, a)), a) = \operatorname{identmat}(r, (n + 1)) \end{array}}}
\vnbprop{invertible-iff-det-nonzero}{\vnbstmt{\begin{array}{@{}l@{}} \forall r, n, a. \\ \quad (\operatorname{is\text{-}field\text{-}ring}(r) \wedge ((n \in \mathbb{N}) \wedge (a \in \operatorname{mat}(n, n, \operatorname{carr}(r))))) \Rightarrow  \\ \quad \quad \operatorname{is\text{-}invertible\text{-}mat}(r, n, a) \Leftrightarrow \neg (\operatorname{det}(r, n, a) = \operatorname{zero}(r)) \end{array}}}

\vnbbookfile{theorem-library/ccr-adj-unit}{theorem-library/ccr-adj-unit}
\noindent{\small ccr-adj-unit.scm -- wave of 2026-11-09/10 (agent CCR-1109-B4), wired 2026-11-10 from scratchpad/ccr1109b4/ccr-adj-unit.scm. Library: linear-algebra.}\par
\vnbprop{det-unit-iff-invertible}{\vnbstmt{\begin{array}{@{}l@{}} \forall r, n, a. \\ \quad (\operatorname{is\text{-}commutative\text{-}ring}(r) \wedge ((n \in \mathbb{N}) \wedge (a \in \operatorname{mat}(n, n, \operatorname{carr}(r))))) \Rightarrow  \\ \quad \quad \operatorname{is\text{-}invertible\text{-}mat}(r, n, a) \Leftrightarrow \exists b \in \operatorname{carr}(r).\; (\operatorname{mul}(r)(\operatorname{det}(r, n, a), b) = \operatorname{one}(r)) \end{array}}}

\vnbbookfile{theorem-library/la-elem-gen}{theorem-library/la-elem-gen}
\noindent{\small la-elem-gen.scm -- elementary column matrices, column equivalence, the column-step proposition and linear-algebra 3.12 / 3.13 (LA-DET2, 2026-10-18).}\par
\vnbprop{elem-f-is-elementary}{\vnbstmt{\begin{array}{@{}l@{}} \forall r, n, k, l. \\ \quad (\operatorname{is\text{-}ring}(r) \wedge ((n \in \mathbb{N}) \wedge ((k \in \operatorname{interval}(1, n)) \wedge ((l \in \operatorname{interval}(1, n)) \wedge \neg (k = l))))) \Rightarrow  \\ \quad \quad \operatorname{is\text{-}elementary}(r, n, \operatorname{elem\text{-}f}(r, n, k, l)) \end{array}}}
\vnbprop{elem-g-is-elementary}{\vnbstmt{\begin{array}{@{}l@{}} \forall r, n, c, k, l. \\ \quad \operatorname{is\text{-}ring}(r) \wedge \\ \quad n \in \mathbb{N} \wedge \\ \quad c \in \operatorname{carr}(r) \wedge \\ \quad k \in \operatorname{interval}(1, n) \wedge \\ \quad l \in \operatorname{interval}(1, n) \wedge \\ \quad \neg (k = l) \Rightarrow \\ \quad \quad \operatorname{is\text{-}elementary}(r, n, \operatorname{elem\text{-}g}(r, n, c, k, l)) \end{array}}}
\vnbprop{elem-h-is-elementary}{\vnbstmt{\begin{array}{@{}l@{}} \forall r, n, c, u, k. \\ \quad \operatorname{is\text{-}ring}(r) \wedge \\ \quad n \in \mathbb{N} \wedge \\ \quad c \in \operatorname{carr}(r) \wedge \\ \quad u \in \operatorname{carr}(r) \wedge \\ \quad \operatorname{mul}(r)(c, u) = \operatorname{one}(r) \wedge \\ \quad \operatorname{mul}(r)(u, c) = \operatorname{one}(r) \wedge \\ \quad k \in \operatorname{interval}(1, n) \Rightarrow \\ \quad \quad \operatorname{is\text{-}elementary}(r, n, \operatorname{elem\text{-}h}(r, n, c, k)) \end{array}}}
\vnbprop{elementary-type}{\vnbstmt{\begin{array}{@{}l@{}} \forall r, n, e. \\ \quad (\operatorname{is\text{-}ring}(r) \wedge ((n \in \mathbb{N}) \wedge \operatorname{is\text{-}elementary}(r, n, e))) \Rightarrow  \\ \quad \quad e \in \operatorname{mat}(n, n, \operatorname{carr}(r)) \end{array}}}
\vnbprop{elementary-inverse}{\vnbstmt{\begin{array}{@{}l@{}} \forall r, n, e. \\ \quad (\operatorname{is\text{-}ring}(r) \wedge ((n \in \mathbb{N}) \wedge \operatorname{is\text{-}elementary}(r, n, e))) \Rightarrow  \\ \quad \quad \exists v. \\ \quad \quad \quad (\operatorname{is\text{-}elementary}(r, n, v) \wedge ((\operatorname{matmul}(r, e, v) = \operatorname{identmat}(r, n)) \wedge (\operatorname{matmul}(r, v, e) = \operatorname{identmat}(r, n)))) \end{array}}}
\vnbprop{elem-gen-unfold}{\vnbstmt{\begin{array}{@{}l@{}} \forall r, n. \\ \quad (\operatorname{elem\text{-}gen}(r, n) \simeq \operatorname{intersection\text{-}of}(\{s \in \operatorname{power}(\operatorname{mat}(n, n, \operatorname{carr}(r))) : ((\operatorname{identmat}(r, n) \in s) \wedge \forall x \in s, e \in \mathbf{Set}.\; (\operatorname{is\text{-}elementary}(r, n, e) \Rightarrow (\operatorname{matmul}(r, x, e) \in s)))\})) \end{array}}}
\vnbprop{laeg-fam-has-mat}{\vnbstmt{\begin{array}{@{}l@{}} \forall r, n. \\ \quad (\operatorname{is\text{-}ring}(r) \wedge (n \in \mathbb{N})) \Rightarrow  \\ \quad \quad \operatorname{mat}(n, n, \operatorname{carr}(r)) \\ \quad \quad \in\; \{s \in \operatorname{power}(\operatorname{mat}(n, n, \operatorname{carr}(r))) : ((\operatorname{identmat}(r, n) \in s) \wedge \forall x \in s, e \in \mathbf{Set}.\; (\operatorname{is\text{-}elementary}(r, n, e) \Rightarrow (\operatorname{matmul}(r, x, e) \in s)))\} \end{array}}}
\vnbprop{elem-gen-mat}{\vnbstmt{\begin{array}{@{}l@{}} \forall r, n, p. \\ \quad (\operatorname{is\text{-}ring}(r) \wedge ((n \in \mathbb{N}) \wedge (p \in \operatorname{elem\text{-}gen}(r, n)))) \Rightarrow (p \in \operatorname{mat}(n, n, \operatorname{carr}(r))) \end{array}}}
\vnbprop{elem-gen-identity}{\vnbstmt{\forall r, n.\; ((\operatorname{is\text{-}ring}(r) \wedge (n \in \mathbb{N})) \Rightarrow (\operatorname{identmat}(r, n) \in \operatorname{elem\text{-}gen}(r, n)))}}
\vnbprop{elem-gen-step}{\vnbstmt{\begin{array}{@{}l@{}} \forall r, n, p, e. \\ \quad (\operatorname{is\text{-}ring}(r) \wedge ((n \in \mathbb{N}) \wedge ((p \in \operatorname{elem\text{-}gen}(r, n)) \wedge \operatorname{is\text{-}elementary}(r, n, e)))) \Rightarrow  \\ \quad \quad \operatorname{matmul}(r, p, e) \in \operatorname{elem\text{-}gen}(r, n) \end{array}}}
\vnbprop{elem-gen-induction}{\vnbstmt{\begin{array}{@{}l@{}} \forall r, n, s.; \forall p \in \operatorname{elem\text{-}gen}(r, n). \\ \quad \operatorname{is\text{-}ring}(r) \wedge \\ \quad n \in \mathbb{N} \wedge \\ \quad \operatorname{identmat}(r, n) \in s \wedge \\ \quad \forall y \in s, f.\; (\operatorname{is\text{-}elementary}(r, n, f) \Rightarrow (\operatorname{matmul}(r, y, f) \in s)) \Rightarrow \\ \quad \quad p \in s \end{array}}}
\vnbprop{elem-gen-mul}{\vnbstmt{\begin{array}{@{}l@{}} \forall r, n, p, q. \\ \quad (\operatorname{is\text{-}ring}(r) \wedge ((n \in \mathbb{N}) \wedge ((p \in \operatorname{elem\text{-}gen}(r, n)) \wedge (q \in \operatorname{elem\text{-}gen}(r, n))))) \Rightarrow  \\ \quad \quad \operatorname{matmul}(r, p, q) \in \operatorname{elem\text{-}gen}(r, n) \end{array}}}
\vnbprop{elementary-in-elem-gen}{\vnbstmt{\begin{array}{@{}l@{}} \forall r, n, e. \\ \quad (\operatorname{is\text{-}ring}(r) \wedge ((n \in \mathbb{N}) \wedge \operatorname{is\text{-}elementary}(r, n, e))) \Rightarrow (e \in \operatorname{elem\text{-}gen}(r, n)) \end{array}}}
\vnbprop{elem-gen-inverse}{\vnbstmt{\begin{array}{@{}l@{}} \forall r, n, p. \\ \quad (\operatorname{is\text{-}ring}(r) \wedge ((n \in \mathbb{N}) \wedge (p \in \operatorname{elem\text{-}gen}(r, n)))) \Rightarrow  \\ \quad \quad \exists v \in \operatorname{elem\text{-}gen}(r, n). \\ \quad \quad \quad (\operatorname{matmul}(r, p, v) = \operatorname{identmat}(r, n)) \wedge (\operatorname{matmul}(r, v, p) = \operatorname{identmat}(r, n)) \end{array}}}
\vnbprop{elem-gen-invertible}{\vnbstmt{\begin{array}{@{}l@{}} \forall r, n, p. \\ \quad (\operatorname{is\text{-}ring}(r) \wedge ((n \in \mathbb{N}) \wedge (p \in \operatorname{elem\text{-}gen}(r, n)))) \Rightarrow  \\ \quad \quad \operatorname{is\text{-}invertible\text{-}mat}(r, n, p) \end{array}}}
\vnbprop{mat-left-inverse-eq-right}{\vnbstmt{\begin{array}{@{}l@{}} \forall r, n, a, p, v. \\ \quad \operatorname{is\text{-}ring}(r) \wedge \\ \quad n \in \mathbb{N} \wedge \\ \quad a \in \operatorname{mat}(n, n, \operatorname{carr}(r)) \wedge \\ \quad p \in \operatorname{mat}(n, n, \operatorname{carr}(r)) \wedge \\ \quad v \in \operatorname{mat}(n, n, \operatorname{carr}(r)) \wedge \\ \quad \operatorname{matmul}(r, a, p) = \operatorname{identmat}(r, n) \wedge \\ \quad \operatorname{matmul}(r, p, v) = \operatorname{identmat}(r, n) \Rightarrow \\ \quad \quad a = v \end{array}}}
\vnbprop{laeg-prod-zero-ind}{\vnbstmt{\begin{array}{@{}l@{}} \forall n \in \mathbb{N}, r, f.; \forall k \in \operatorname{interval}(1, n). \\ \quad (\operatorname{is\text{-}commutative\text{-}ring}(r) \wedge (\forall z \in \operatorname{interval}(1, n).\; (f(z) \in \operatorname{carr}(r)) \wedge (f(k) = \operatorname{zero}(r)))) \Rightarrow  \\ \quad \quad \operatorname{prod\text{-}ring}(r, f, \operatorname{interval}(1, n)) = \operatorname{zero}(r) \end{array}}}
\vnbprop{prod-ring-zero-factor}{\vnbstmt{\begin{array}{@{}l@{}} \forall r, n, f.; \forall k \in \operatorname{interval}(1, n). \\ \quad \operatorname{is\text{-}commutative\text{-}ring}(r) \wedge \\ \quad n \in \mathbb{N} \wedge \\ \quad \forall z \in \operatorname{interval}(1, n).\; (f(z) \in \operatorname{carr}(r)) \wedge \\ \quad f(k) = \operatorname{zero}(r) \Rightarrow \\ \quad \quad \operatorname{prod\text{-}ring}(r, f, \operatorname{interval}(1, n)) = \operatorname{zero}(r) \end{array}}}
\vnbprop{laeg-op-f-agree}{\vnbstmt{\begin{array}{@{}l@{}} \forall r, m, n, t, x, a, k, l.; \forall i \in \operatorname{interval}(1, t), j \in \operatorname{interval}(1, n). \\ \quad \operatorname{is\text{-}ring}(r) \wedge \\ \quad m \in \mathbb{N} \wedge \\ \quad n \in \mathbb{N} \wedge \\ \quad t \in \mathbb{N} \wedge \\ \quad t \leq m \wedge \\ \quad x \in \operatorname{mat}(m, n, \operatorname{carr}(r)) \wedge \\ \quad a \in \operatorname{mat}(m, n, \operatorname{carr}(r)) \wedge \\ \quad \forall i \in \operatorname{interval}(1, t), j \in \operatorname{interval}(1, n). \\ \quad \quad \operatorname{entry}(a, i, j) = \operatorname{entry}(x, i, j) \wedge \\ \quad \forall i \in \operatorname{interval}(1, t), j \in \operatorname{interval}(1, n). \\ \quad \quad ((t + 1) \leq j) \Rightarrow (\operatorname{entry}(x, i, j) = \operatorname{zero}(r)) \wedge \\ \quad k \in \operatorname{interval}(1, n) \wedge \\ \quad l \in \operatorname{interval}(1, n) \wedge \\ \quad \neg (k = l) \wedge \\ \quad (t + 1) \leq k \wedge \\ \quad (t + 1) \leq l \Rightarrow \\ \quad \quad \operatorname{entry}(\operatorname{matmul}(r, a, \operatorname{elem\text{-}f}(r, n, k, l)), i, j) = \operatorname{entry}(x, i, j) \end{array}}}
\vnbprop{laeg-op-h-agree}{\vnbstmt{\begin{array}{@{}l@{}} \forall r, m, n, t, x, a, c, k.; \forall i \in \operatorname{interval}(1, t), j \in \operatorname{interval}(1, n). \\ \quad \operatorname{is\text{-}ring}(r) \wedge \\ \quad m \in \mathbb{N} \wedge \\ \quad n \in \mathbb{N} \wedge \\ \quad t \in \mathbb{N} \wedge \\ \quad t \leq m \wedge \\ \quad x \in \operatorname{mat}(m, n, \operatorname{carr}(r)) \wedge \\ \quad a \in \operatorname{mat}(m, n, \operatorname{carr}(r)) \wedge \\ \quad \forall i \in \operatorname{interval}(1, t), j \in \operatorname{interval}(1, n). \\ \quad \quad \operatorname{entry}(a, i, j) = \operatorname{entry}(x, i, j) \wedge \\ \quad \forall i \in \operatorname{interval}(1, t), j \in \operatorname{interval}(1, n). \\ \quad \quad ((t + 1) \leq j) \Rightarrow (\operatorname{entry}(x, i, j) = \operatorname{zero}(r)) \wedge \\ \quad c \in \operatorname{carr}(r) \wedge \\ \quad k \in \operatorname{interval}(1, n) \wedge \\ \quad (t + 1) \leq k \Rightarrow \\ \quad \quad \operatorname{entry}(\operatorname{matmul}(r, a, \operatorname{elem\text{-}h}(r, n, c, k)), i, j) = \operatorname{entry}(x, i, j) \end{array}}}
\vnbprop{laeg-op-g-agree}{\vnbstmt{\begin{array}{@{}l@{}} \forall r, m, n, t, x, a, c, k, l.; \forall i \in \operatorname{interval}(1, t), j \in \operatorname{interval}(1, n). \\ \quad \operatorname{is\text{-}ring}(r) \wedge \\ \quad m \in \mathbb{N} \wedge \\ \quad n \in \mathbb{N} \wedge \\ \quad t \in \mathbb{N} \wedge \\ \quad t \leq m \wedge \\ \quad x \in \operatorname{mat}(m, n, \operatorname{carr}(r)) \wedge \\ \quad a \in \operatorname{mat}(m, n, \operatorname{carr}(r)) \wedge \\ \quad \forall i \in \operatorname{interval}(1, t), j \in \operatorname{interval}(1, n). \\ \quad \quad \operatorname{entry}(a, i, j) = \operatorname{entry}(x, i, j) \wedge \\ \quad \forall i \in \operatorname{interval}(1, t), j \in \operatorname{interval}(1, n). \\ \quad \quad ((t + 1) \leq j) \Rightarrow (\operatorname{entry}(x, i, j) = \operatorname{zero}(r)) \wedge \\ \quad c \in \operatorname{carr}(r) \wedge \\ \quad k \in \operatorname{interval}(1, n) \wedge \\ \quad l \in \operatorname{interval}(1, n) \wedge \\ \quad \neg (k = l) \wedge \\ \quad (t + 1) \leq k \Rightarrow \\ \quad \quad \operatorname{entry}(\operatorname{matmul}(r, a, \operatorname{elem\text{-}g}(r, n, c, k, l)), i, j) = \operatorname{entry}(x, i, j) \end{array}}}
\vnbprop{col-equiv-refl}{\vnbstmt{\begin{array}{@{}l@{}} \forall r, m, n, a. \\ \quad (\operatorname{is\text{-}ring}(r) \wedge ((m \in \mathbb{N}) \wedge ((n \in \mathbb{N}) \wedge (a \in \operatorname{mat}(m, n, \operatorname{carr}(r)))))) \Rightarrow  \\ \quad \quad \operatorname{col\text{-}equiv}(r, n, a, a) \end{array}}}
\vnbprop{col-equiv-mat}{\vnbstmt{\begin{array}{@{}l@{}} \forall r, m, n, a, b. \\ \quad (\operatorname{is\text{-}ring}(r) \wedge ((m \in \mathbb{N}) \wedge ((n \in \mathbb{N}) \wedge ((a \in \operatorname{mat}(m, n, \operatorname{carr}(r))) \wedge \operatorname{col\text{-}equiv}(r, n, a, b))))) \Rightarrow  \\ \quad \quad b \in \operatorname{mat}(m, n, \operatorname{carr}(r)) \end{array}}}
\vnbprop{col-equiv-trans}{\vnbstmt{\begin{array}{@{}l@{}} \forall r, m, n, a, b, c. \\ \quad \operatorname{is\text{-}ring}(r) \wedge \\ \quad m \in \mathbb{N} \wedge \\ \quad n \in \mathbb{N} \wedge \\ \quad a \in \operatorname{mat}(m, n, \operatorname{carr}(r)) \wedge \\ \quad \operatorname{col\text{-}equiv}(r, n, a, b) \wedge \\ \quad \operatorname{col\text{-}equiv}(r, n, b, c) \Rightarrow \\ \quad \quad \operatorname{col\text{-}equiv}(r, n, a, c) \end{array}}}
\vnbprop{col-equiv-elem}{\vnbstmt{\begin{array}{@{}l@{}} \forall r, m, n, a, e. \\ \quad (\operatorname{is\text{-}ring}(r) \wedge ((m \in \mathbb{N}) \wedge ((n \in \mathbb{N}) \wedge ((a \in \operatorname{mat}(m, n, \operatorname{carr}(r))) \wedge \operatorname{is\text{-}elementary}(r, n, e))))) \Rightarrow  \\ \quad \quad \operatorname{col\text{-}equiv}(r, n, a, \operatorname{matmul}(r, a, e)) \end{array}}}
\vnbprop{laeg-agree-refl}{\vnbstmt{\begin{array}{@{}l@{}} \forall r, m, n, t, x.; \forall i \in \operatorname{interval}(1, t), j \in \operatorname{interval}(1, n). \\ \quad (\operatorname{is\text{-}ring}(r) \wedge ((m \in \mathbb{N}) \wedge ((n \in \mathbb{N}) \wedge ((t \in \mathbb{N}) \wedge ((t \leq m) \wedge (x \in \operatorname{mat}(m, n, \operatorname{carr}(r)))))))) \Rightarrow  \\ \quad \quad \operatorname{entry}(x, i, j) = \operatorname{entry}(x, i, j) \end{array}}}
\vnbprop{laeg-rows-zr}{\vnbstmt{\begin{array}{@{}l@{}} \forall r, n, t, x.; \forall i \in \operatorname{interval}(1, t), j \in \operatorname{interval}(1, n). \\ \quad ((n \in \mathbb{N}) \wedge ((t \in \mathbb{N}) \wedge (\forall i \in \operatorname{interval}(1, t).\; \operatorname{col\text{-}step\text{-}row}(r, n, x, i) \wedge ((t + 1) \leq j)))) \Rightarrow  \\ \quad \quad \operatorname{entry}(x, i, j) = \operatorname{zero}(r) \end{array}}}
\vnbprop{laeg-csr-transfer}{\vnbstmt{\begin{array}{@{}l@{}} \forall r, m, n, t, x, a, i. \\ \quad \operatorname{is\text{-}ring}(r) \wedge \\ \quad m \in \mathbb{N} \wedge \\ \quad n \in \mathbb{N} \wedge \\ \quad t \in \mathbb{N} \wedge \\ \quad t \leq m \wedge \\ \quad x \in \operatorname{mat}(m, n, \operatorname{carr}(r)) \wedge \\ \quad a \in \operatorname{mat}(m, n, \operatorname{carr}(r)) \wedge \\ \quad \forall \mathit{agi} \in \operatorname{interval}(1, t), j \in \operatorname{interval}(1, n). \\ \quad \quad \operatorname{entry}(a, \mathit{agi}, j) = \operatorname{entry}(x, \mathit{agi}, j) \wedge \\ \quad i \in \operatorname{interval}(1, t) \wedge \\ \quad \operatorname{col\text{-}step\text{-}row}(r, n, x, i) \Rightarrow \\ \quad \quad \operatorname{col\text{-}step\text{-}row}(r, n, a, i) \end{array}}}
\vnbprop{laeg-step-pivot}{\vnbstmt{\begin{array}{@{}l@{}} \forall r, m, n, t, x, z. \\ \quad \operatorname{is\text{-}field\text{-}ring}(r) \wedge \\ \quad m \in \mathbb{N} \wedge \\ \quad n \in \mathbb{N} \wedge \\ \quad t \in \mathbb{N} \wedge \\ \quad (t + 1) \leq m \wedge \\ \quad (t + 1) \leq n \wedge \\ \quad x \in \operatorname{mat}(m, n, \operatorname{carr}(r)) \wedge \\ \quad \forall i \in \operatorname{interval}(1, t), j \in \operatorname{interval}(1, n). \\ \quad \quad ((t + 1) \leq j) \Rightarrow (\operatorname{entry}(x, i, j) = \operatorname{zero}(r)) \wedge \\ \quad z \in \operatorname{interval}(1, n) \wedge \\ \quad (t + 1) \leq z \wedge \\ \quad \neg (\operatorname{entry}(x, (t + 1), z) = \operatorname{zero}(r)) \Rightarrow \\ \quad \quad \exists c. \\ \quad \quad \quad \operatorname{col\text{-}equiv}(r, n, x, c) \wedge \\ \quad \quad \quad \forall i \in \operatorname{interval}(1, t), j \in \operatorname{interval}(1, n). \\ \quad \quad \quad \quad \operatorname{entry}(c, i, j) = \operatorname{entry}(x, i, j) \wedge \\ \quad \quad \quad \neg (\operatorname{entry}(c, (t + 1), (t + 1)) = \operatorname{zero}(r)) \end{array}}}
\vnbprop{laeg-step-norm}{\vnbstmt{\begin{array}{@{}l@{}} \forall r, m, n, t, x, a. \\ \quad \operatorname{is\text{-}field\text{-}ring}(r) \wedge \\ \quad m \in \mathbb{N} \wedge \\ \quad n \in \mathbb{N} \wedge \\ \quad t \in \mathbb{N} \wedge \\ \quad (t + 1) \leq m \wedge \\ \quad (t + 1) \leq n \wedge \\ \quad x \in \operatorname{mat}(m, n, \operatorname{carr}(r)) \wedge \\ \quad \forall i \in \operatorname{interval}(1, t), j \in \operatorname{interval}(1, n). \\ \quad \quad ((t + 1) \leq j) \Rightarrow (\operatorname{entry}(x, i, j) = \operatorname{zero}(r)) \wedge \\ \quad a \in \operatorname{mat}(m, n, \operatorname{carr}(r)) \wedge \\ \quad \forall i \in \operatorname{interval}(1, t), j \in \operatorname{interval}(1, n). \\ \quad \quad \operatorname{entry}(a, i, j) = \operatorname{entry}(x, i, j) \wedge \\ \quad \neg (\operatorname{entry}(a, (t + 1), (t + 1)) = \operatorname{zero}(r)) \Rightarrow \\ \quad \quad \exists c. \\ \quad \quad \quad \operatorname{col\text{-}equiv}(r, n, a, c) \wedge \\ \quad \quad \quad \forall i \in \operatorname{interval}(1, t), j \in \operatorname{interval}(1, n). \\ \quad \quad \quad \quad \operatorname{entry}(c, i, j) = \operatorname{entry}(x, i, j) \wedge \\ \quad \quad \quad \operatorname{entry}(c, (t + 1), (t + 1)) = \operatorname{one}(r) \end{array}}}
\vnbprop{laeg-clear-ind}{\vnbstmt{\begin{array}{@{}l@{}} \forall s \in \mathbb{N}, r, m, n, t, x, a. \\ \quad \operatorname{is\text{-}field\text{-}ring}(r) \wedge \\ \quad m \in \mathbb{N} \wedge \\ \quad n \in \mathbb{N} \wedge \\ \quad t \in \mathbb{N} \wedge \\ \quad (t + 1) \leq m \wedge \\ \quad (t + 1) \leq n \wedge \\ \quad x \in \operatorname{mat}(m, n, \operatorname{carr}(r)) \wedge \\ \quad \forall i \in \operatorname{interval}(1, t), j \in \operatorname{interval}(1, n). \\ \quad \quad ((t + 1) \leq j) \Rightarrow (\operatorname{entry}(x, i, j) = \operatorname{zero}(r)) \wedge \\ \quad a \in \operatorname{mat}(m, n, \operatorname{carr}(r)) \wedge \\ \quad \forall i \in \operatorname{interval}(1, t), j \in \operatorname{interval}(1, n). \\ \quad \quad \operatorname{entry}(a, i, j) = \operatorname{entry}(x, i, j) \wedge \\ \quad \operatorname{entry}(a, (t + 1), (t + 1)) = \operatorname{one}(r) \wedge \\ \quad s \leq n \Rightarrow \\ \quad \quad \exists c \in \operatorname{mat}(m, n, \operatorname{carr}(r)). \\ \quad \quad \quad \operatorname{col\text{-}equiv}(r, n, a, c) \wedge \\ \quad \quad \quad \forall i \in \operatorname{interval}(1, t), j \in \operatorname{interval}(1, n). \\ \quad \quad \quad \quad \operatorname{entry}(c, i, j) = \operatorname{entry}(x, i, j) \wedge \\ \quad \quad \quad \forall \mathit{csa} \in \operatorname{interval}(1, m).\; (\operatorname{entry}(c, \mathit{csa}, (t + 1)) = \operatorname{entry}(a, \mathit{csa}, (t + 1))) \wedge \\ \quad \quad \quad \forall \mathit{clj} \in \operatorname{interval}(1, n). \\ \quad \quad \quad \quad ((\mathit{clj} \leq s) \wedge \neg (\mathit{clj} = (t + 1))) \Rightarrow (\operatorname{entry}(c, (t + 1), \mathit{clj}) = \operatorname{zero}(r)) \end{array}}}
\vnbprop{laeg-row-step}{\vnbstmt{\begin{array}{@{}l@{}} \forall r, m, n, t, x. \\ \quad \operatorname{is\text{-}field\text{-}ring}(r) \wedge \\ \quad m \in \mathbb{N} \wedge \\ \quad n \in \mathbb{N} \wedge \\ \quad t \in \mathbb{N} \wedge \\ \quad (t + 1) \leq m \wedge \\ \quad x \in \operatorname{mat}(m, n, \operatorname{carr}(r)) \wedge \\ \quad \forall i \in \operatorname{interval}(1, t).\; \operatorname{col\text{-}step\text{-}row}(r, n, x, i) \Rightarrow \\ \quad \quad \exists c. \\ \quad \quad \quad (\operatorname{col\text{-}equiv}(r, n, x, c) \wedge \forall i \in \operatorname{interval}(1, (t + 1)).\; \operatorname{col\text{-}step\text{-}row}(r, n, c, i)) \end{array}}}
\vnbprop{laeg-colstep-ind}{\vnbstmt{\begin{array}{@{}l@{}} \forall t \in \mathbb{N}, r, m, n, a. \\ \quad (\operatorname{is\text{-}field\text{-}ring}(r) \wedge ((m \in \mathbb{N}) \wedge ((n \in \mathbb{N}) \wedge ((t \leq m) \wedge (a \in \operatorname{mat}(m, n, \operatorname{carr}(r))))))) \Rightarrow  \\ \quad \quad \exists c \in \operatorname{mat}(m, n, \operatorname{carr}(r)). \\ \quad \quad \quad (\operatorname{col\text{-}equiv}(r, n, a, c) \wedge \forall i \in \operatorname{interval}(1, t).\; \operatorname{col\text{-}step\text{-}row}(r, n, c, i)) \end{array}}}
\vnbprop{column-step-reduced}{\vnbstmt{\begin{array}{@{}l@{}} \forall r, m, n, a. \\ \quad (\operatorname{is\text{-}field\text{-}ring}(r) \wedge ((m \in \mathbb{N}) \wedge ((n \in \mathbb{N}) \wedge (a \in \operatorname{mat}(m, n, \operatorname{carr}(r)))))) \Rightarrow  \\ \quad \quad \exists c. \\ \quad \quad \quad (\operatorname{col\text{-}equiv}(r, n, a, c) \wedge \forall i \in \operatorname{interval}(1, m).\; \operatorname{col\text{-}step\text{-}row}(r, n, c, i)) \end{array}}}
\vnbprop{laeg-rows-colstep}{\vnbstmt{\begin{array}{@{}l@{}} \forall r, m, n, c. \\ \quad ((m \in \mathbb{N}) \wedge ((n \in \mathbb{N}) \wedge \forall i \in \operatorname{interval}(1, m).\; \operatorname{col\text{-}step\text{-}row}(r, n, c, i))) \Rightarrow  \\ \quad \quad \operatorname{is\text{-}column\text{-}step}(r, m, n, c) \end{array}}}
\vnbprop{column-step-exists}{\vnbstmt{\begin{array}{@{}l@{}} \forall r, m, n, a. \\ \quad (\operatorname{is\text{-}field\text{-}ring}(r) \wedge ((m \in \mathbb{N}) \wedge ((n \in \mathbb{N}) \wedge (a \in \operatorname{mat}(m, n, \operatorname{carr}(r)))))) \Rightarrow  \\ \quad \quad \exists e.\; (\operatorname{col\text{-}equiv}(r, n, a, e) \wedge \operatorname{is\text{-}column\text{-}step}(r, m, n, e)) \end{array}}}
\vnbprop{col-equiv-identity-in-elem-gen}{\vnbstmt{\begin{array}{@{}l@{}} \forall r, n, a. \\ \quad (\operatorname{is\text{-}ring}(r) \wedge ((n \in \mathbb{N}) \wedge ((a \in \operatorname{mat}(n, n, \operatorname{carr}(r))) \wedge \operatorname{col\text{-}equiv}(r, n, a, \operatorname{identmat}(r, n))))) \Rightarrow  \\ \quad \quad a \in \operatorname{elem\text{-}gen}(r, n) \end{array}}}
\vnbprop{laeg-reduced-invertible-identity}{\vnbstmt{\begin{array}{@{}l@{}} \forall r, n, c. \\ \quad \operatorname{is\text{-}field\text{-}ring}(r) \wedge \\ \quad n \in \mathbb{N} \wedge \\ \quad c \in \operatorname{mat}(n, n, \operatorname{carr}(r)) \wedge \\ \quad \operatorname{is\text{-}invertible\text{-}mat}(r, n, c) \wedge \\ \quad \forall i \in \operatorname{interval}(1, n).\; \operatorname{col\text{-}step\text{-}row}(r, n, c, i) \Rightarrow \\ \quad \quad c = \operatorname{identmat}(r, n) \end{array}}}
\vnbprop{invertible-iff-col-equiv-identity}{\vnbstmt{\begin{array}{@{}l@{}} \forall r, n, a. \\ \quad (\operatorname{is\text{-}field\text{-}ring}(r) \wedge ((n \in \mathbb{N}) \wedge (a \in \operatorname{mat}(n, n, \operatorname{carr}(r))))) \Rightarrow  \\ \quad \quad \operatorname{is\text{-}invertible\text{-}mat}(r, n, a) \Leftrightarrow \operatorname{col\text{-}equiv}(r, n, a, \operatorname{identmat}(r, n)) \end{array}}}
\vnbprop{invertible-iff-elem-gen}{\vnbstmt{\begin{array}{@{}l@{}} \forall r, n, a. \\ \quad (\operatorname{is\text{-}field\text{-}ring}(r) \wedge ((n \in \mathbb{N}) \wedge (a \in \operatorname{mat}(n, n, \operatorname{carr}(r))))) \Rightarrow  \\ \quad \quad \operatorname{is\text{-}invertible\text{-}mat}(r, n, a) \Leftrightarrow (a \in \operatorname{elem\text{-}gen}(r, n)) \end{array}}}

\vnbbookfile{theorem-library/la-z6-col}{theorem-library/la-z6-col}
\noindent{\small la-z6-col.scm -- linear algebra Lemma 3.8 over a field, in column equivalence (LA-Z6, 2026-10-25).}\par
\vnbprop{col-reduce-first-row}{\vnbstmt{\begin{array}{@{}l@{}} \forall r, m, n, a. \\ \quad (\operatorname{is\text{-}field\text{-}ring}(r) \wedge ((m \in \mathbb{N}) \wedge ((n \in \mathbb{N}) \wedge ((1 \leq m) \wedge ((1 \leq n) \wedge (a \in \operatorname{mat}(m, n, \operatorname{carr}(r)))))))) \Rightarrow  \\ \quad \quad \exists b. \\ \quad \quad \quad \operatorname{col\text{-}equiv}(r, n, a, b) \wedge \\ \quad \quad \quad \forall j \in \operatorname{interval}(2, n).\; (\operatorname{entry}(b, 1, j) = \operatorname{zero}(r)) \wedge \\ \quad \quad \quad (\operatorname{entry}(b, 1, 1) = \operatorname{zero}(r)) \vee (\operatorname{entry}(b, 1, 1) = \operatorname{one}(r)) \end{array}}}

\vnbbookfile{theorem-library/la-z-sdet}{theorem-library/la-z-sdet}
\noindent{\small la-z-sdet.scm -- subdeterminants under the elementary operations (linear algebra 3.29) and the invariance of their gcd over Z (3.31) (LA-Z, 2026-10-19).}\par
\vnbprop{submatrix-type}{\vnbstmt{\begin{array}{@{}l@{}} \forall r, m, n, s, a, p, q. \\ \quad \operatorname{is\text{-}ring}(r) \wedge \\ \quad m \in \mathbb{N} \wedge \\ \quad n \in \mathbb{N} \wedge \\ \quad s \in \mathbb{N} \wedge \\ \quad a \in \operatorname{mat}(m, n, \operatorname{carr}(r)) \wedge \\ \quad p \in (\operatorname{interval}(1, s) \to \operatorname{interval}(1, m)) \wedge \\ \quad q \in (\operatorname{interval}(1, s) \to \operatorname{interval}(1, n)) \Rightarrow \\ \quad \quad \operatorname{submatrix}(a, s, p, q) \in \operatorname{mat}(s, s, \operatorname{carr}(r)) \end{array}}}
\vnbprop{submatrix-entry}{\vnbstmt{\begin{array}{@{}l@{}} \forall r, m, n, s, a, p, q, i, j. \\ \quad \operatorname{is\text{-}ring}(r) \wedge \\ \quad m \in \mathbb{N} \wedge \\ \quad n \in \mathbb{N} \wedge \\ \quad s \in \mathbb{N} \wedge \\ \quad a \in \operatorname{mat}(m, n, \operatorname{carr}(r)) \wedge \\ \quad p \in (\operatorname{interval}(1, s) \to \operatorname{interval}(1, m)) \wedge \\ \quad q \in (\operatorname{interval}(1, s) \to \operatorname{interval}(1, n)) \wedge \\ \quad i \in \operatorname{interval}(1, s) \wedge \\ \quad j \in \operatorname{interval}(1, s) \Rightarrow \\ \quad \quad \operatorname{entry}(\operatorname{submatrix}(a, s, p, q), i, j) = \operatorname{entry}(a, p(i), q(j)) \end{array}}}
\vnbprop{submatrix-ext}{\vnbstmt{\begin{array}{@{}l@{}} \forall r, m, n, s, a, p, q, b, c, d. \\ \quad \operatorname{is\text{-}ring}(r) \wedge \\ \quad m \in \mathbb{N} \wedge \\ \quad n \in \mathbb{N} \wedge \\ \quad s \in \mathbb{N} \wedge \\ \quad a \in \operatorname{mat}(m, n, \operatorname{carr}(r)) \wedge \\ \quad p \in (\operatorname{interval}(1, s) \to \operatorname{interval}(1, m)) \wedge \\ \quad q \in (\operatorname{interval}(1, s) \to \operatorname{interval}(1, n)) \wedge \\ \quad b \in \operatorname{mat}(m, n, \operatorname{carr}(r)) \wedge \\ \quad c \in (\operatorname{interval}(1, s) \to \operatorname{interval}(1, m)) \wedge \\ \quad d \in (\operatorname{interval}(1, s) \to \operatorname{interval}(1, n)) \wedge \\ \quad \forall i, j \in \operatorname{interval}(1, s). \\ \quad \quad \operatorname{entry}(a, p(i), q(j)) = \operatorname{entry}(b, c(i), d(j)) \Rightarrow \\ \quad \quad \operatorname{submatrix}(a, s, p, q) = \operatorname{submatrix}(b, s, c, d) \end{array}}}
\vnbprop{sdet-unfold}{\vnbstmt{\begin{array}{@{}l@{}} \forall r, m, n, a, s. \\ \quad (\operatorname{sdet}(r, m, n, a, s) \simeq \{x \in \operatorname{carr}(r) : \exists p \in \operatorname{injection}(\operatorname{interval}(1, s), \operatorname{interval}(1, m)), q \in \operatorname{injection}(\operatorname{interval}(1, s), \operatorname{interval}(1, n)).\; (x = \operatorname{det}(r, s, \operatorname{submatrix}(a, s, p, q)))\}) \end{array}}}
\vnbprop{sdet-intro}{\vnbstmt{\begin{array}{@{}l@{}} \forall r, m, n, s, a, p, q. \\ \quad \operatorname{is\text{-}ring}(r) \wedge \\ \quad m \in \mathbb{N} \wedge \\ \quad n \in \mathbb{N} \wedge \\ \quad s \in \mathbb{N} \wedge \\ \quad a \in \operatorname{mat}(m, n, \operatorname{carr}(r)) \wedge \\ \quad p \in \operatorname{injection}(\operatorname{interval}(1, s), \operatorname{interval}(1, m)) \wedge \\ \quad q \in \operatorname{injection}(\operatorname{interval}(1, s), \operatorname{interval}(1, n)) \Rightarrow \\ \quad \quad \operatorname{det}(r, s, \operatorname{submatrix}(a, s, p, q)) \in \operatorname{sdet}(r, m, n, a, s) \end{array}}}
\vnbprop{sdet-elim}{\vnbstmt{\begin{array}{@{}l@{}} \forall r, m, n, a, s.; \forall x \in \operatorname{sdet}(r, m, n, a, s). \exists p \in \operatorname{injection}(\operatorname{interval}(1, s), \operatorname{interval}(1, m)), q \in \operatorname{injection}(\operatorname{interval}(1, s), \operatorname{interval}(1, n)). \\ \quad x = \operatorname{det}(r, s, \operatorname{submatrix}(a, s, p, q)) \end{array}}}
\vnbprop{lzs-tau-tau}{\vnbstmt{\begin{array}{@{}l@{}} \forall k, l, y. \\ \quad \neg (k = l) \Rightarrow  \\ \quad \quad \left(\text{if }(\left(\text{if }(y = k)\text{ then }l\text{ else }\left(\text{if }(y = l)\text{ then }k\text{ else }y\right)\right) = k)\text{ then }l\text{ else }\left(\text{if }(\left(\text{if }(y = k)\text{ then }l\text{ else }\left(\text{if }(y = l)\text{ then }k\text{ else }y\right)\right) = l)\text{ then }k\text{ else }\left(\text{if }(y = k)\text{ then }l\text{ else }\left(\text{if }(y = l)\text{ then }k\text{ else }y\right)\right)\right)\right) \\ \quad \quad =\; y \end{array}}}
\vnbprop{lzs-swap-injection}{\vnbstmt{\begin{array}{@{}l@{}} \forall s, n, q, k, l. \\ \quad s \in \mathbb{N} \wedge \\ \quad n \in \mathbb{N} \wedge \\ \quad q \in \operatorname{injection}(\operatorname{interval}(1, s), \operatorname{interval}(1, n)) \wedge \\ \quad k \in \operatorname{interval}(1, n) \wedge \\ \quad l \in \operatorname{interval}(1, n) \wedge \\ \quad \neg (k = l) \Rightarrow \\ \quad \quad (\lambda j \in \operatorname{interval}(1, s).\; \left(\text{if }(q(j) = k)\text{ then }l\text{ else }\left(\text{if }(q(j) = l)\text{ then }k\text{ else }q(j)\right)\right)) \\ \quad \quad \in\; \operatorname{injection}(\operatorname{interval}(1, s), \operatorname{interval}(1, n)) \end{array}}}
\vnbprop{lzs-replace-injection}{\vnbstmt{\begin{array}{@{}l@{}} \forall s, n, q, a, k. \\ \quad s \in \mathbb{N} \wedge \\ \quad n \in \mathbb{N} \wedge \\ \quad q \in \operatorname{injection}(\operatorname{interval}(1, s), \operatorname{interval}(1, n)) \wedge \\ \quad a \in \operatorname{interval}(1, s) \wedge \\ \quad k \in \operatorname{interval}(1, n) \wedge \\ \quad \forall j \in \operatorname{interval}(1, s).\; \neg (q(j) = k) \Rightarrow \\ \quad \quad (\lambda w \in \operatorname{interval}(1, s).\; \left(\text{if }(w = a)\text{ then }k\text{ else }q(w)\right)) \in \operatorname{injection}(\operatorname{interval}(1, s), \operatorname{interval}(1, n)) \end{array}}}
\vnbprop{sdet-col-swap-sub}{\vnbstmt{\begin{array}{@{}l@{}} \forall r, m, n, s, a, k, l, x. \\ \quad \operatorname{is\text{-}ring}(r) \wedge \\ \quad m \in \mathbb{N} \wedge \\ \quad n \in \mathbb{N} \wedge \\ \quad s \in \mathbb{N} \wedge \\ \quad a \in \operatorname{mat}(m, n, \operatorname{carr}(r)) \wedge \\ \quad k \in \operatorname{interval}(1, n) \wedge \\ \quad l \in \operatorname{interval}(1, n) \wedge \\ \quad \neg (k = l) \wedge \\ \quad x \in \operatorname{sdet}(r, m, n, \operatorname{matmul}(r, a, \operatorname{elem\text{-}f}(r, n, k, l)), s) \Rightarrow \\ \quad \quad x \in \operatorname{sdet}(r, m, n, a, s) \end{array}}}
\vnbprop{sdet-col-swap}{\vnbstmt{\begin{array}{@{}l@{}} \forall r, m, n, s, a, k, l. \\ \quad \operatorname{is\text{-}ring}(r) \wedge \\ \quad m \in \mathbb{N} \wedge \\ \quad n \in \mathbb{N} \wedge \\ \quad s \in \mathbb{N} \wedge \\ \quad a \in \operatorname{mat}(m, n, \operatorname{carr}(r)) \wedge \\ \quad k \in \operatorname{interval}(1, n) \wedge \\ \quad l \in \operatorname{interval}(1, n) \wedge \\ \quad \neg (k = l) \Rightarrow \\ \quad \quad \operatorname{sdet}(r, m, n, \operatorname{matmul}(r, a, \operatorname{elem\text{-}f}(r, n, k, l)), s) = \operatorname{sdet}(r, m, n, a, s) \end{array}}}
\vnbprop{lzs-ring-a}{\vnbstmt{\begin{array}{@{}l@{}} \forall r, u, v, w. \\ \quad (\operatorname{is\text{-}commutative\text{-}ring}(r) \wedge ((u \in \operatorname{carr}(r)) \wedge ((v \in \operatorname{carr}(r)) \wedge (w \in \operatorname{carr}(r))))) \Rightarrow  \\ \quad \quad \operatorname{add}(r)(u, \operatorname{mul}(r)(v, w)) = \operatorname{add}(r)(\operatorname{mul}(r)(w, v), u) \end{array}}}
\vnbprop{lzs-ring-b}{\vnbstmt{\begin{array}{@{}l@{}} \forall r, u, v, w. \\ \quad (\operatorname{is\text{-}commutative\text{-}ring}(r) \wedge ((u \in \operatorname{carr}(r)) \wedge ((v \in \operatorname{carr}(r)) \wedge (w \in \operatorname{carr}(r))))) \Rightarrow  \\ \quad \quad \operatorname{add}(r)(\operatorname{mul}(r)(w, v), u) = \operatorname{add}(r)(u, \operatorname{mul}(r)(w, v)) \end{array}}}
\vnbprop{lzs-ring-z}{\vnbstmt{\begin{array}{@{}l@{}} \forall r, u, v, w. \\ \quad (\operatorname{is\text{-}commutative\text{-}ring}(r) \wedge ((u \in \operatorname{carr}(r)) \wedge ((v \in \operatorname{carr}(r)) \wedge (w \in \operatorname{carr}(r))))) \Rightarrow  \\ \quad \quad \operatorname{add}(r)(\operatorname{mul}(r)(w, \operatorname{zero}(r)), u) = \operatorname{add}(r)(u, \operatorname{mul}(r)(\operatorname{zero}(r), u)) \end{array}}}
\vnbprop{lzs-ring-z0}{\vnbstmt{\begin{array}{@{}l@{}} \forall r, u, v, w. \\ \quad (\operatorname{is\text{-}commutative\text{-}ring}(r) \wedge ((u \in \operatorname{carr}(r)) \wedge ((v \in \operatorname{carr}(r)) \wedge (w \in \operatorname{carr}(r))))) \Rightarrow  \\ \quad \quad u = \operatorname{add}(r)(u, \operatorname{mul}(r)(\operatorname{zero}(r), u)) \end{array}}}
\vnbprop{lzs-ring-c}{\vnbstmt{\begin{array}{@{}l@{}} \forall r, u, v, w. \\ \quad (\operatorname{is\text{-}commutative\text{-}ring}(r) \wedge ((u \in \operatorname{carr}(r)) \wedge ((v \in \operatorname{carr}(r)) \wedge (w \in \operatorname{carr}(r))))) \Rightarrow  \\ \quad \quad \operatorname{add}(r)(u, \operatorname{mul}(r)(w, v)) = \operatorname{add}(r)(\operatorname{mul}(r)(w, v), u) \end{array}}}
\vnbprop{sdet-col-add}{\vnbstmt{\begin{array}{@{}l@{}} \forall r, m, n, s, a, c, k, l, x. \\ \quad \operatorname{is\text{-}commutative\text{-}ring}(r) \wedge \\ \quad m \in \mathbb{N} \wedge \\ \quad n \in \mathbb{N} \wedge \\ \quad s \in \mathbb{N} \wedge \\ \quad a \in \operatorname{mat}(m, n, \operatorname{carr}(r)) \wedge \\ \quad c \in \operatorname{carr}(r) \wedge \\ \quad k \in \operatorname{interval}(1, n) \wedge \\ \quad l \in \operatorname{interval}(1, n) \wedge \\ \quad \neg (k = l) \wedge \\ \quad x \in \operatorname{sdet}(r, m, n, \operatorname{matmul}(r, a, \operatorname{elem\text{-}g}(r, n, c, k, l)), s) \Rightarrow \\ \quad \quad \exists \mathit{sda}, b \in \operatorname{sdet}(r, m, n, a, s), t \in \operatorname{carr}(r). \\ \quad \quad \quad ((t = \operatorname{zero}(r)) \vee ((t = c) \vee (t = \operatorname{neg}(r)(c)))) \wedge (x = \operatorname{add}(r)(\mathit{sda}, \operatorname{mul}(r)(t, b))) \end{array}}}
\vnbprop{sdet-row-swap-sub}{\vnbstmt{\begin{array}{@{}l@{}} \forall r, m, n, s, a, k, l, x. \\ \quad \operatorname{is\text{-}ring}(r) \wedge \\ \quad m \in \mathbb{N} \wedge \\ \quad n \in \mathbb{N} \wedge \\ \quad s \in \mathbb{N} \wedge \\ \quad a \in \operatorname{mat}(m, n, \operatorname{carr}(r)) \wedge \\ \quad k \in \operatorname{interval}(1, m) \wedge \\ \quad l \in \operatorname{interval}(1, m) \wedge \\ \quad \neg (k = l) \wedge \\ \quad x \in \operatorname{sdet}(r, m, n, \operatorname{matmul}(r, \operatorname{elem\text{-}f}(r, m, k, l), a), s) \Rightarrow \\ \quad \quad x \in \operatorname{sdet}(r, m, n, a, s) \end{array}}}
\vnbprop{sdet-row-swap}{\vnbstmt{\begin{array}{@{}l@{}} \forall r, m, n, s, a, k, l. \\ \quad \operatorname{is\text{-}ring}(r) \wedge \\ \quad m \in \mathbb{N} \wedge \\ \quad n \in \mathbb{N} \wedge \\ \quad s \in \mathbb{N} \wedge \\ \quad a \in \operatorname{mat}(m, n, \operatorname{carr}(r)) \wedge \\ \quad k \in \operatorname{interval}(1, m) \wedge \\ \quad l \in \operatorname{interval}(1, m) \wedge \\ \quad \neg (k = l) \Rightarrow \\ \quad \quad \operatorname{sdet}(r, m, n, \operatorname{matmul}(r, \operatorname{elem\text{-}f}(r, m, k, l), a), s) = \operatorname{sdet}(r, m, n, a, s) \end{array}}}
\vnbprop{sdet-row-add}{\vnbstmt{\begin{array}{@{}l@{}} \forall r, m, n, s, a, c, k, l, x. \\ \quad \operatorname{is\text{-}commutative\text{-}ring}(r) \wedge \\ \quad m \in \mathbb{N} \wedge \\ \quad n \in \mathbb{N} \wedge \\ \quad s \in \mathbb{N} \wedge \\ \quad a \in \operatorname{mat}(m, n, \operatorname{carr}(r)) \wedge \\ \quad c \in \operatorname{carr}(r) \wedge \\ \quad k \in \operatorname{interval}(1, m) \wedge \\ \quad l \in \operatorname{interval}(1, m) \wedge \\ \quad \neg (k = l) \wedge \\ \quad x \in \operatorname{sdet}(r, m, n, \operatorname{matmul}(r, \operatorname{elem\text{-}g}(r, m, c, k, l), a), s) \Rightarrow \\ \quad \quad \exists \mathit{sda}, b \in \operatorname{sdet}(r, m, n, a, s), t \in \operatorname{carr}(r). \\ \quad \quad \quad ((t = \operatorname{zero}(r)) \vee ((t = c) \vee (t = \operatorname{neg}(r)(c)))) \wedge (x = \operatorname{add}(r)(\mathit{sda}, \operatorname{mul}(r)(t, b))) \end{array}}}
\vnbprop{sdet-col-scale}{\vnbstmt{\begin{array}{@{}l@{}} \forall r, m, n, s, a, c, k, x. \\ \quad \operatorname{is\text{-}commutative\text{-}ring}(r) \wedge \\ \quad m \in \mathbb{N} \wedge \\ \quad n \in \mathbb{N} \wedge \\ \quad s \in \mathbb{N} \wedge \\ \quad a \in \operatorname{mat}(m, n, \operatorname{carr}(r)) \wedge \\ \quad c \in \operatorname{carr}(r) \wedge \\ \quad k \in \operatorname{interval}(1, n) \wedge \\ \quad x \in \operatorname{sdet}(r, m, n, \operatorname{matmul}(r, a, \operatorname{elem\text{-}h}(r, n, c, k)), s) \Rightarrow \\ \quad \quad \exists \mathit{sda} \in \operatorname{sdet}(r, m, n, a, s).\; ((x = \mathit{sda}) \vee (x = \operatorname{mul}(r)(c, \mathit{sda}))) \end{array}}}
\vnbprop{sdet-row-scale}{\vnbstmt{\begin{array}{@{}l@{}} \forall r, m, n, s, a, c, k, x. \\ \quad \operatorname{is\text{-}commutative\text{-}ring}(r) \wedge \\ \quad m \in \mathbb{N} \wedge \\ \quad n \in \mathbb{N} \wedge \\ \quad s \in \mathbb{N} \wedge \\ \quad a \in \operatorname{mat}(m, n, \operatorname{carr}(r)) \wedge \\ \quad c \in \operatorname{carr}(r) \wedge \\ \quad k \in \operatorname{interval}(1, m) \wedge \\ \quad x \in \operatorname{sdet}(r, m, n, \operatorname{matmul}(r, \operatorname{elem\text{-}h}(r, m, c, k), a), s) \Rightarrow \\ \quad \quad \exists \mathit{sda} \in \operatorname{sdet}(r, m, n, a, s).\; ((x = \mathit{sda}) \vee (x = \operatorname{mul}(r)(c, \mathit{sda}))) \end{array}}}

\vnbbookfile{theorem-library/la-z-gcd}{theorem-library/la-z-gcd}
\noindent{\small la-z-gcd.scm -- the gcd of the s-th order subdeterminants is invariant under integral elementary operations (linear algebra 3.31) (LA-Z, 2026-10-19).}\par
\vnbprop{lzg-ring-g}{\vnbstmt{\begin{array}{@{}l@{}} \forall r, c, y, t, z. \\ \quad (\operatorname{is\text{-}commutative\text{-}ring}(r) \wedge ((c \in \operatorname{carr}(r)) \wedge ((y \in \operatorname{carr}(r)) \wedge ((t \in \operatorname{carr}(r)) \wedge (z \in \operatorname{carr}(r)))))) \Rightarrow  \\ \quad \quad \operatorname{add}(r)(\operatorname{mul}(r)(c, y), \operatorname{mul}(r)(t, \operatorname{mul}(r)(c, z))) = \operatorname{mul}(r)(c, \operatorname{add}(r)(y, \operatorname{mul}(r)(t, z))) \end{array}}}
\vnbprop{lzg-ring-h}{\vnbstmt{\begin{array}{@{}l@{}} \forall r, c, y, t. \\ \quad (\operatorname{is\text{-}commutative\text{-}ring}(r) \wedge ((c \in \operatorname{carr}(r)) \wedge ((y \in \operatorname{carr}(r)) \wedge (t \in \operatorname{carr}(r))))) \Rightarrow  \\ \quad \quad \operatorname{mul}(r)(t, \operatorname{mul}(r)(c, y)) = \operatorname{mul}(r)(c, \operatorname{mul}(r)(t, y)) \end{array}}}
\vnbprop{sdet-divisor-col-step}{\vnbstmt{\begin{array}{@{}l@{}} \forall r, m, n, s, a, c, e.; \forall x \in \operatorname{sdet}(r, m, n, \operatorname{matmul}(r, a, e), s). \\ \quad \operatorname{is\text{-}commutative\text{-}ring}(r) \wedge \\ \quad m \in \mathbb{N} \wedge \\ \quad n \in \mathbb{N} \wedge \\ \quad s \in \mathbb{N} \wedge \\ \quad a \in \operatorname{mat}(m, n, \operatorname{carr}(r)) \wedge \\ \quad c \in \operatorname{carr}(r) \wedge \\ \quad \forall x \in \operatorname{sdet}(r, m, n, a, s).\; \exists y \in \operatorname{carr}(r).\; (x = \operatorname{mul}(r)(c, y)) \wedge \\ \quad \operatorname{is\text{-}elementary}(r, n, e) \Rightarrow \\ \quad \quad \exists y \in \operatorname{carr}(r).\; (x = \operatorname{mul}(r)(c, y)) \end{array}}}
\vnbprop{sdet-divisor-row-step}{\vnbstmt{\begin{array}{@{}l@{}} \forall r, m, n, s, a, c, e.; \forall x \in \operatorname{sdet}(r, m, n, \operatorname{matmul}(r, e, a), s). \\ \quad \operatorname{is\text{-}commutative\text{-}ring}(r) \wedge \\ \quad m \in \mathbb{N} \wedge \\ \quad n \in \mathbb{N} \wedge \\ \quad s \in \mathbb{N} \wedge \\ \quad a \in \operatorname{mat}(m, n, \operatorname{carr}(r)) \wedge \\ \quad c \in \operatorname{carr}(r) \wedge \\ \quad \forall x \in \operatorname{sdet}(r, m, n, a, s).\; \exists y \in \operatorname{carr}(r).\; (x = \operatorname{mul}(r)(c, y)) \wedge \\ \quad \operatorname{is\text{-}elementary}(r, m, e) \Rightarrow \\ \quad \quad \exists y \in \operatorname{carr}(r).\; (x = \operatorname{mul}(r)(c, y)) \end{array}}}
\vnbprop{sdet-divisor-col-gen}{\vnbstmt{\begin{array}{@{}l@{}} \forall r, m, n, s, c, q.; \forall a \in \operatorname{mat}(m, n, \operatorname{carr}(r)), x \in \operatorname{sdet}(r, m, n, \operatorname{matmul}(r, a, q), s). \\ \quad \operatorname{is\text{-}commutative\text{-}ring}(r) \wedge \\ \quad m \in \mathbb{N} \wedge \\ \quad n \in \mathbb{N} \wedge \\ \quad s \in \mathbb{N} \wedge \\ \quad c \in \operatorname{carr}(r) \wedge \\ \quad q \in \operatorname{elem\text{-}gen}(r, n) \wedge \\ \quad \forall x \in \operatorname{sdet}(r, m, n, a, s).\; \exists y \in \operatorname{carr}(r).\; (x = \operatorname{mul}(r)(c, y)) \Rightarrow \\ \quad \quad \exists y \in \operatorname{carr}(r).\; (x = \operatorname{mul}(r)(c, y)) \end{array}}}
\vnbprop{sdet-divisor-row-gen}{\vnbstmt{\begin{array}{@{}l@{}} \forall r, m, n, s, c, p.; \forall a \in \operatorname{mat}(m, n, \operatorname{carr}(r)), x \in \operatorname{sdet}(r, m, n, \operatorname{matmul}(r, p, a), s). \\ \quad \operatorname{is\text{-}commutative\text{-}ring}(r) \wedge \\ \quad m \in \mathbb{N} \wedge \\ \quad n \in \mathbb{N} \wedge \\ \quad s \in \mathbb{N} \wedge \\ \quad c \in \operatorname{carr}(r) \wedge \\ \quad p \in \operatorname{elem\text{-}gen}(r, m) \wedge \\ \quad \forall x \in \operatorname{sdet}(r, m, n, a, s).\; \exists y \in \operatorname{carr}(r).\; (x = \operatorname{mul}(r)(c, y)) \Rightarrow \\ \quad \quad \exists y \in \operatorname{carr}(r).\; (x = \operatorname{mul}(r)(c, y)) \end{array}}}
\vnbprop{sdet-divisor-equiv}{\vnbstmt{\begin{array}{@{}l@{}} \forall r, m, n, s, c, a, b.; \forall x \in \operatorname{sdet}(r, m, n, b, s). \\ \quad \operatorname{is\text{-}commutative\text{-}ring}(r) \wedge \\ \quad m \in \mathbb{N} \wedge \\ \quad n \in \mathbb{N} \wedge \\ \quad s \in \mathbb{N} \wedge \\ \quad c \in \operatorname{carr}(r) \wedge \\ \quad a \in \operatorname{mat}(m, n, \operatorname{carr}(r)) \wedge \\ \quad \operatorname{elem\text{-}equiv}(r, m, n, a, b) \wedge \\ \quad \forall x \in \operatorname{sdet}(r, m, n, a, s).\; \exists y \in \operatorname{carr}(r).\; (x = \operatorname{mul}(r)(c, y)) \Rightarrow \\ \quad \quad \exists y \in \operatorname{carr}(r).\; (x = \operatorname{mul}(r)(c, y)) \end{array}}}
\vnbprop{elem-equiv-symmetric}{\vnbstmt{\begin{array}{@{}l@{}} \forall r, m, n, a, b. \\ \quad (\operatorname{is\text{-}ring}(r) \wedge ((m \in \mathbb{N}) \wedge ((n \in \mathbb{N}) \wedge ((a \in \operatorname{mat}(m, n, \operatorname{carr}(r))) \wedge \operatorname{elem\text{-}equiv}(r, m, n, a, b))))) \Rightarrow  \\ \quad \quad \operatorname{elem\text{-}equiv}(r, m, n, b, a) \end{array}}}
\vnbprop{lzg-zz-to-ring}{\vnbstmt{\begin{array}{@{}l@{}} \forall d, x. \\ \quad \operatorname{zz\text{-}divides}(d, x) \Rightarrow \exists y \in \operatorname{carr}(\operatorname{zz\text{-}ring}).\; (x = \operatorname{mul}(\operatorname{zz\text{-}ring})(d, y)) \end{array}}}
\vnbprop{lzg-ring-to-zz}{\vnbstmt{\begin{array}{@{}l@{}} \forall d, x, y. \\ \quad ((d \in \mathbb{Z}) \wedge ((y \in \operatorname{carr}(\operatorname{zz\text{-}ring})) \wedge (x = \operatorname{mul}(\operatorname{zz\text{-}ring})(d, y)))) \Rightarrow  \\ \quad \quad \operatorname{zz\text{-}divides}(d, x) \end{array}}}
\vnbprop{sdet-gcd-invariant}{\vnbstmt{\begin{array}{@{}l@{}} \forall m, n, s, a, b, d, e. \\ \quad m \in \mathbb{N} \wedge \\ \quad n \in \mathbb{N} \wedge \\ \quad s \in \mathbb{N} \wedge \\ \quad a \in \operatorname{mat}(m, n, \mathbb{Z}) \wedge \\ \quad b \in \operatorname{mat}(m, n, \mathbb{Z}) \wedge \\ \quad \operatorname{elem\text{-}equiv}(\operatorname{zz\text{-}ring}, m, n, a, b) \wedge \\ \quad \operatorname{zz\text{-}is\text{-}gcd\text{-}set}(\operatorname{sdet}(\operatorname{zz\text{-}ring}, m, n, a, s), d) \wedge \\ \quad \operatorname{zz\text{-}is\text{-}gcd\text{-}set}(\operatorname{sdet}(\operatorname{zz\text{-}ring}, m, n, b, s), e) \Rightarrow \\ \quad \quad (d = e) \vee (d = -e) \end{array}}}

\vnbbookfile{theorem-library/la-z-elem}{theorem-library/la-z-elem}
\noindent{\small la-z-elem.scm -- over a euclidean ring an invertible matrix is a product of elementary matrices; over Z this is linear algebra 3.38 (LA-Z, 2026-10-19).}\par
\vnbprop{lze-ring-neg}{\vnbstmt{\begin{array}{@{}l@{}} \forall r, a, s. \\ \quad (\operatorname{is\text{-}commutative\text{-}ring}(r) \wedge ((a \in \operatorname{carr}(r)) \wedge (s \in \operatorname{carr}(r)))) \Rightarrow  \\ \quad \quad \operatorname{mul}(r)(\operatorname{mul}(r)(a, \operatorname{neg}(r)(\operatorname{one}(r))), \operatorname{neg}(r)(s)) = \operatorname{mul}(r)(a, s) \end{array}}}
\vnbprop{lze-ring-one}{\vnbstmt{\begin{array}{@{}l@{}} \forall r, a, s. \\ \quad (\operatorname{is\text{-}commutative\text{-}ring}(r) \wedge ((a \in \operatorname{carr}(r)) \wedge (s \in \operatorname{carr}(r)))) \Rightarrow  \\ \quad \quad \operatorname{mul}(r)(\operatorname{mul}(r)(a, \operatorname{one}(r)), s) = \operatorname{mul}(r)(a, s) \end{array}}}
\vnbprop{lze-ring-four}{\vnbstmt{\begin{array}{@{}l@{}} \forall r, a, \mathit{ler}, s, t. \\ \quad (\operatorname{is\text{-}commutative\text{-}ring}(r) \wedge ((a \in \operatorname{carr}(r)) \wedge ((\mathit{ler} \in \operatorname{carr}(r)) \wedge ((s \in \operatorname{carr}(r)) \wedge (t \in \operatorname{carr}(r)))))) \Rightarrow  \\ \quad \quad \operatorname{mul}(r)(\operatorname{mul}(r)(a, \mathit{ler}), \operatorname{mul}(r)(s, t)) = \operatorname{mul}(r)(\operatorname{mul}(r)(a, s), \operatorname{mul}(r)(\mathit{ler}, t)) \end{array}}}
\vnbprop{lze-ring-assoc}{\vnbstmt{\begin{array}{@{}l@{}} \forall r, a, \mathit{ler}, s. \\ \quad (\operatorname{is\text{-}commutative\text{-}ring}(r) \wedge ((a \in \operatorname{carr}(r)) \wedge ((\mathit{ler} \in \operatorname{carr}(r)) \wedge (s \in \operatorname{carr}(r))))) \Rightarrow  \\ \quad \quad \operatorname{mul}(r)(a, \operatorname{mul}(r)(\mathit{ler}, s)) = \operatorname{mul}(r)(\operatorname{mul}(r)(a, \mathit{ler}), s) \end{array}}}
\vnbprop{lze-unit-mul}{\vnbstmt{\begin{array}{@{}l@{}} \forall r, a, \mathit{ler}, s, t. \\ \quad \operatorname{is\text{-}commutative\text{-}ring}(r) \wedge \\ \quad a \in \operatorname{carr}(r) \wedge \\ \quad \mathit{ler} \in \operatorname{carr}(r) \wedge \\ \quad s \in \operatorname{carr}(r) \wedge \\ \quad t \in \operatorname{carr}(r) \wedge \\ \quad \operatorname{mul}(r)(a, s) = \operatorname{one}(r) \wedge \\ \quad \operatorname{mul}(r)(\mathit{ler}, t) = \operatorname{one}(r) \Rightarrow \\ \quad \quad \operatorname{mul}(r)(\operatorname{mul}(r)(a, \mathit{ler}), \operatorname{mul}(r)(s, t)) = \operatorname{one}(r) \end{array}}}
\vnbprop{elem-gen-det-unit}{\vnbstmt{\begin{array}{@{}l@{}} \forall r, n, p. \\ \quad (\operatorname{is\text{-}commutative\text{-}ring}(r) \wedge ((n \in \mathbb{N}) \wedge (p \in \operatorname{elem\text{-}gen}(r, n)))) \Rightarrow  \\ \quad \quad \exists s \in \operatorname{carr}(r).\; (\operatorname{mul}(r)(\operatorname{det}(r, n, p), s) = \operatorname{one}(r)) \end{array}}}
\vnbprop{det-row-single}{\vnbstmt{\begin{array}{@{}l@{}} \forall r, n, b, w. \\ \quad \operatorname{is\text{-}commutative\text{-}ring}(r) \wedge \\ \quad n \in \mathbb{N} \wedge \\ \quad b \in \operatorname{mat}((n + 1), (n + 1), \operatorname{carr}(r)) \wedge \\ \quad w \in \operatorname{interval}(1, (n + 1)) \wedge \\ \quad \forall j \in \operatorname{interval}(1, (n + 1)).\; (\neg (j = w) \Rightarrow (\operatorname{entry}(b, w, j) = \operatorname{zero}(r))) \Rightarrow \\ \quad \quad \operatorname{det}(r, (n + 1), b) = \operatorname{mul}(r)(\operatorname{entry}(b, w, w), \operatorname{det}(r, n, \operatorname{minor}(b, w, w, n))) \end{array}}}
\vnbprop{lze-zr-identity}{\vnbstmt{\begin{array}{@{}l@{}} \forall r, n, t.; \forall i \in \operatorname{interval}(1, t), j \in \operatorname{interval}(1, n). \\ \quad (\operatorname{is\text{-}ring}(r) \wedge ((n \in \mathbb{N}) \wedge ((t \in \mathbb{N}) \wedge ((t \leq n) \wedge ((t + 1) \leq j))))) \Rightarrow  \\ \quad \quad \operatorname{entry}(\operatorname{identmat}(r, n), i, j) = \operatorname{zero}(r) \end{array}}}
\vnbprop{lze-cols-f}{\vnbstmt{\begin{array}{@{}l@{}} \forall r, n, t, b, k, l.; \forall i \in \operatorname{interval}(1, n), j \in \operatorname{interval}(1, t). \\ \quad \operatorname{is\text{-}ring}(r) \wedge \\ \quad n \in \mathbb{N} \wedge \\ \quad t \in \mathbb{N} \wedge \\ \quad t \leq n \wedge \\ \quad b \in \operatorname{mat}(n, n, \operatorname{carr}(r)) \wedge \\ \quad \forall i \in \operatorname{interval}(1, n), j \in \operatorname{interval}(1, t). \\ \quad \quad \operatorname{entry}(b, i, j) = \operatorname{entry}(\operatorname{identmat}(r, n), i, j) \wedge \\ \quad k \in \operatorname{interval}(1, n) \wedge \\ \quad l \in \operatorname{interval}(1, n) \wedge \\ \quad \neg (k = l) \wedge \\ \quad (t + 1) \leq k \wedge \\ \quad (t + 1) \leq l \Rightarrow \\ \quad \quad \operatorname{entry}(\operatorname{matmul}(r, b, \operatorname{elem\text{-}f}(r, n, k, l)), i, j) = \operatorname{entry}(\operatorname{identmat}(r, n), i, j) \end{array}}}
\vnbprop{lze-cols-g}{\vnbstmt{\begin{array}{@{}l@{}} \forall r, n, t, b, c, k, l.; \forall i \in \operatorname{interval}(1, n), j \in \operatorname{interval}(1, t). \\ \quad \operatorname{is\text{-}ring}(r) \wedge \\ \quad n \in \mathbb{N} \wedge \\ \quad t \in \mathbb{N} \wedge \\ \quad t \leq n \wedge \\ \quad b \in \operatorname{mat}(n, n, \operatorname{carr}(r)) \wedge \\ \quad \forall i \in \operatorname{interval}(1, n), j \in \operatorname{interval}(1, t). \\ \quad \quad \operatorname{entry}(b, i, j) = \operatorname{entry}(\operatorname{identmat}(r, n), i, j) \wedge \\ \quad c \in \operatorname{carr}(r) \wedge \\ \quad k \in \operatorname{interval}(1, n) \wedge \\ \quad l \in \operatorname{interval}(1, n) \wedge \\ \quad \neg (k = l) \wedge \\ \quad (t + 1) \leq l \Rightarrow \\ \quad \quad \operatorname{entry}(\operatorname{matmul}(r, b, \operatorname{elem\text{-}g}(r, n, c, k, l)), i, j) = \operatorname{entry}(\operatorname{identmat}(r, n), i, j) \end{array}}}
\vnbprop{lze-cols-h}{\vnbstmt{\begin{array}{@{}l@{}} \forall r, n, t, b, c, k.; \forall i \in \operatorname{interval}(1, n), j \in \operatorname{interval}(1, t). \\ \quad \operatorname{is\text{-}ring}(r) \wedge \\ \quad n \in \mathbb{N} \wedge \\ \quad t \in \mathbb{N} \wedge \\ \quad t \leq n \wedge \\ \quad b \in \operatorname{mat}(n, n, \operatorname{carr}(r)) \wedge \\ \quad \forall i \in \operatorname{interval}(1, n), j \in \operatorname{interval}(1, t). \\ \quad \quad \operatorname{entry}(b, i, j) = \operatorname{entry}(\operatorname{identmat}(r, n), i, j) \wedge \\ \quad c \in \operatorname{carr}(r) \wedge \\ \quad k \in \operatorname{interval}(1, n) \wedge \\ \quad (t + 1) \leq k \Rightarrow \\ \quad \quad \operatorname{entry}(\operatorname{matmul}(r, b, \operatorname{elem\text{-}h}(r, n, c, k)), i, j) = \operatorname{entry}(\operatorname{identmat}(r, n), i, j) \end{array}}}
\vnbprop{lze-ring-euc}{\vnbstmt{\begin{array}{@{}l@{}} \forall r, q, b, \mathit{ler}. \\ \quad (\operatorname{is\text{-}commutative\text{-}ring}(r) \wedge ((q \in \operatorname{carr}(r)) \wedge ((b \in \operatorname{carr}(r)) \wedge (\mathit{ler} \in \operatorname{carr}(r))))) \Rightarrow  \\ \quad \quad \operatorname{add}(r)(\operatorname{add}(r)(\operatorname{mul}(r)(q, b), \mathit{ler}), \operatorname{mul}(r)(b, \operatorname{neg}(r)(q))) = \mathit{ler} \end{array}}}
\vnbprop{lze-pair-ind}{\vnbstmt{\begin{array}{@{}l@{}} \forall \mathit{lpn} \in \mathbb{N}, r, n, t, b, j. \\ \quad \operatorname{is\text{-}euclidean\text{-}ring}(r) \wedge \\ \quad n \in \mathbb{N} \wedge \\ \quad t \in \mathbb{N} \wedge \\ \quad (t + 1) \leq n \wedge \\ \quad b \in \operatorname{mat}(n, n, \operatorname{carr}(r)) \wedge \\ \quad \forall i \in \operatorname{interval}(1, t), \mathit{agj} \in \operatorname{interval}(1, n). \\ \quad \quad \operatorname{entry}(b, i, \mathit{agj}) = \operatorname{entry}(\operatorname{identmat}(r, n), i, \mathit{agj}) \wedge \\ \quad \forall \mathit{cgi} \in \operatorname{interval}(1, n), \mathit{cgj} \in \operatorname{interval}(1, t). \\ \quad \quad \operatorname{entry}(b, \mathit{cgi}, \mathit{cgj}) = \operatorname{entry}(\operatorname{identmat}(r, n), \mathit{cgi}, \mathit{cgj}) \wedge \\ \quad j \in \operatorname{interval}(1, n) \wedge \\ \quad (t + 1) < j \wedge \\ \quad \neg (\operatorname{entry}(b, (t + 1), (t + 1)) = \operatorname{zero}(r)) \Rightarrow  \\ \quad \quad \operatorname{gauge}(r)(\operatorname{entry}(b, (t + 1), (t + 1))) \leq \mathit{lpn} \Rightarrow \\ \quad \quad \exists q \in \operatorname{elem\text{-}gen}(r, n). \\ \quad \quad \quad \forall i \in \operatorname{interval}(1, t), \mathit{agj} \in \operatorname{interval}(1, n). \\ \quad \quad \quad \quad \operatorname{entry}(\operatorname{matmul}(r, b, q), i, \mathit{agj}) = \operatorname{entry}(\operatorname{identmat}(r, n), i, \mathit{agj}) \wedge \\ \quad \quad \quad \forall \mathit{cgi} \in \operatorname{interval}(1, n), \mathit{cgj} \in \operatorname{interval}(1, t). \\ \quad \quad \quad \quad \operatorname{entry}(\operatorname{matmul}(r, b, q), \mathit{cgi}, \mathit{cgj}) = \operatorname{entry}(\operatorname{identmat}(r, n), \mathit{cgi}, \mathit{cgj}) \wedge \\ \quad \quad \quad \operatorname{entry}(\operatorname{matmul}(r, b, q), (t + 1), j) = \operatorname{zero}(r) \wedge \\ \quad \quad \quad \forall c \in \operatorname{interval}(1, n). \\ \quad \quad \quad \quad (\neg (c = (t + 1)) \wedge \neg (c = j)) \Rightarrow  \\ \quad \quad \quad \quad \quad \operatorname{entry}(\operatorname{matmul}(r, b, q), (t + 1), c) = \operatorname{entry}(b, (t + 1), c) \end{array}}}
\vnbprop{lze-row-ind}{\vnbstmt{\begin{array}{@{}l@{}} \forall k \in \mathbb{N}, r, n, t, b. \\ \quad \operatorname{is\text{-}euclidean\text{-}ring}(r) \wedge \\ \quad n \in \mathbb{N} \wedge \\ \quad t \in \mathbb{N} \wedge \\ \quad (t + 1) \leq n \wedge \\ \quad b \in \operatorname{mat}(n, n, \operatorname{carr}(r)) \wedge \\ \quad \forall i \in \operatorname{interval}(1, t), j \in \operatorname{interval}(1, n). \\ \quad \quad \operatorname{entry}(b, i, j) = \operatorname{entry}(\operatorname{identmat}(r, n), i, j) \wedge \\ \quad \forall \mathit{cgi} \in \operatorname{interval}(1, n), \mathit{cgj} \in \operatorname{interval}(1, t). \\ \quad \quad \operatorname{entry}(b, \mathit{cgi}, \mathit{cgj}) = \operatorname{entry}(\operatorname{identmat}(r, n), \mathit{cgi}, \mathit{cgj}) \Rightarrow \\ \quad \quad \exists q \in \operatorname{elem\text{-}gen}(r, n). \\ \quad \quad \quad \forall i \in \operatorname{interval}(1, t), j \in \operatorname{interval}(1, n). \\ \quad \quad \quad \quad \operatorname{entry}(\operatorname{matmul}(r, b, q), i, j) = \operatorname{entry}(\operatorname{identmat}(r, n), i, j) \wedge \\ \quad \quad \quad \forall \mathit{cgi} \in \operatorname{interval}(1, n), \mathit{cgj} \in \operatorname{interval}(1, t). \\ \quad \quad \quad \quad \operatorname{entry}(\operatorname{matmul}(r, b, q), \mathit{cgi}, \mathit{cgj}) = \operatorname{entry}(\operatorname{identmat}(r, n), \mathit{cgi}, \mathit{cgj}) \wedge \\ \quad \quad \quad \forall c \in \operatorname{interval}(1, n). \\ \quad \quad \quad \quad (((t + 1) < c) \wedge (c \leq k)) \Rightarrow (\operatorname{entry}(\operatorname{matmul}(r, b, q), (t + 1), c) = \operatorname{zero}(r)) \end{array}}}
\vnbprop{lze-ring-z1}{\vnbstmt{\begin{array}{@{}l@{}} \forall r, u, v. \\ \quad (\operatorname{is\text{-}commutative\text{-}ring}(r) \wedge ((u \in \operatorname{carr}(r)) \wedge (v \in \operatorname{carr}(r)))) \Rightarrow  \\ \quad \quad \operatorname{add}(r)(u, \operatorname{mul}(r)(v, \operatorname{zero}(r))) = u \end{array}}}
\vnbprop{lze-ring-z2}{\vnbstmt{\begin{array}{@{}l@{}} \forall r.; \forall u \in \operatorname{carr}(r). \\ \quad \operatorname{is\text{-}commutative\text{-}ring}(r) \Rightarrow  \\ \quad \quad \operatorname{add}(r)(u, \operatorname{mul}(r)(\operatorname{neg}(r)(u), \operatorname{one}(r))) = \operatorname{zero}(r) \end{array}}}
\vnbprop{lze-col-ind}{\vnbstmt{\begin{array}{@{}l@{}} \forall k \in \mathbb{N}, r, n, t, c. \\ \quad \operatorname{is\text{-}commutative\text{-}ring}(r) \wedge \\ \quad n \in \mathbb{N} \wedge \\ \quad t \in \mathbb{N} \wedge \\ \quad (t + 1) \leq n \wedge \\ \quad c \in \operatorname{mat}(n, n, \operatorname{carr}(r)) \wedge \\ \quad \forall i \in \operatorname{interval}(1, t), j \in \operatorname{interval}(1, n). \\ \quad \quad \operatorname{entry}(c, i, j) = \operatorname{entry}(\operatorname{identmat}(r, n), i, j) \wedge \\ \quad \forall \mathit{cgi} \in \operatorname{interval}(1, n), \mathit{cgj} \in \operatorname{interval}(1, t). \\ \quad \quad \operatorname{entry}(c, \mathit{cgi}, \mathit{cgj}) = \operatorname{entry}(\operatorname{identmat}(r, n), \mathit{cgi}, \mathit{cgj}) \wedge \\ \quad \operatorname{entry}(c, (t + 1), (t + 1)) = \operatorname{one}(r) \wedge \\ \quad \forall \mathit{rrc} \in \operatorname{interval}(1, n). \\ \quad \quad \neg (\mathit{rrc} = (t + 1)) \Rightarrow (\operatorname{entry}(c, (t + 1), \mathit{rrc}) = \operatorname{zero}(r)) \Rightarrow \\ \quad \quad \exists p \in \operatorname{elem\text{-}gen}(r, n). \\ \quad \quad \quad \forall i \in \operatorname{interval}(1, t), j \in \operatorname{interval}(1, n). \\ \quad \quad \quad \quad \operatorname{entry}(\operatorname{matmul}(r, p, c), i, j) = \operatorname{entry}(\operatorname{identmat}(r, n), i, j) \wedge \\ \quad \quad \quad \forall \mathit{cgi} \in \operatorname{interval}(1, n), \mathit{cgj} \in \operatorname{interval}(1, t). \\ \quad \quad \quad \quad \operatorname{entry}(\operatorname{matmul}(r, p, c), \mathit{cgi}, \mathit{cgj}) = \operatorname{entry}(\operatorname{identmat}(r, n), \mathit{cgi}, \mathit{cgj}) \wedge \\ \quad \quad \quad \operatorname{entry}(\operatorname{matmul}(r, p, c), (t + 1), (t + 1)) = \operatorname{one}(r) \wedge \\ \quad \quad \quad \forall \mathit{rrc} \in \operatorname{interval}(1, n). \\ \quad \quad \quad \quad \neg (\mathit{rrc} = (t + 1)) \Rightarrow (\operatorname{entry}(\operatorname{matmul}(r, p, c), (t + 1), \mathit{rrc}) = \operatorname{zero}(r)) \wedge \\ \quad \quad \quad \forall \mathit{lci} \in \operatorname{interval}(1, n). \\ \quad \quad \quad \quad (((t + 1) < \mathit{lci}) \wedge (\mathit{lci} \leq k)) \Rightarrow  \\ \quad \quad \quad \quad \quad \operatorname{entry}(\operatorname{matmul}(r, p, c), \mathit{lci}, (t + 1)) = \operatorname{zero}(r) \end{array}}}
\vnbprop{lze-step}{\vnbstmt{\begin{array}{@{}l@{}} \forall r, n, t, u. \\ \quad \operatorname{is\text{-}euclidean\text{-}ring}(r) \wedge \\ \quad n \in \mathbb{N} \wedge \\ \quad t \in \mathbb{N} \wedge \\ \quad (t + 1) \leq n \wedge \\ \quad u \in \operatorname{mat}(n, n, \operatorname{carr}(r)) \wedge \\ \quad \forall i \in \operatorname{interval}(1, t), j \in \operatorname{interval}(1, n). \\ \quad \quad \operatorname{entry}(u, i, j) = \operatorname{entry}(\operatorname{identmat}(r, n), i, j) \wedge \\ \quad \forall \mathit{cgi} \in \operatorname{interval}(1, n), \mathit{cgj} \in \operatorname{interval}(1, t). \\ \quad \quad \operatorname{entry}(u, \mathit{cgi}, \mathit{cgj}) = \operatorname{entry}(\operatorname{identmat}(r, n), \mathit{cgi}, \mathit{cgj}) \wedge \\ \quad \exists s \in \operatorname{carr}(r).\; (\operatorname{mul}(r)(\operatorname{det}(r, n, u), s) = \operatorname{one}(r)) \Rightarrow \\ \quad \quad \exists p, q \in \operatorname{elem\text{-}gen}(r, n). \\ \quad \quad \quad \forall i \in \operatorname{interval}(1, (t + 1)), j \in \operatorname{interval}(1, n). \\ \quad \quad \quad \quad \operatorname{entry}(\operatorname{matmul}(r, \operatorname{matmul}(r, p, u), q), i, j) = \operatorname{entry}(\operatorname{identmat}(r, n), i, j) \wedge \\ \quad \quad \quad \forall \mathit{cgi} \in \operatorname{interval}(1, n), \mathit{cgj} \in \operatorname{interval}(1, (t + 1)). \\ \quad \quad \quad \quad \operatorname{entry}(\operatorname{matmul}(r, \operatorname{matmul}(r, p, u), q), \mathit{cgi}, \mathit{cgj}) = \operatorname{entry}(\operatorname{identmat}(r, n), \mathit{cgi}, \mathit{cgj}) \wedge \\ \quad \quad \quad \exists s \in \operatorname{carr}(r). \\ \quad \quad \quad \quad \operatorname{mul}(r)(\operatorname{det}(r, n, \operatorname{matmul}(r, \operatorname{matmul}(r, p, u), q)), s) = \operatorname{one}(r) \end{array}}}
\vnbprop{lze-main-ind}{\vnbstmt{\begin{array}{@{}l@{}} \forall t \in \mathbb{N}, r, n, u. \\ \quad \operatorname{is\text{-}euclidean\text{-}ring}(r) \wedge \\ \quad n \in \mathbb{N} \wedge \\ \quad u \in \operatorname{mat}(n, n, \operatorname{carr}(r)) \wedge \\ \quad \exists s \in \operatorname{carr}(r).\; (\operatorname{mul}(r)(\operatorname{det}(r, n, u), s) = \operatorname{one}(r)) \wedge \\ \quad t \leq n \Rightarrow \\ \quad \quad \exists p, q \in \operatorname{elem\text{-}gen}(r, n). \\ \quad \quad \quad \forall i \in \operatorname{interval}(1, t), j \in \operatorname{interval}(1, n). \\ \quad \quad \quad \quad \operatorname{entry}(\operatorname{matmul}(r, \operatorname{matmul}(r, p, u), q), i, j) = \operatorname{entry}(\operatorname{identmat}(r, n), i, j) \wedge \\ \quad \quad \quad \forall \mathit{cgi} \in \operatorname{interval}(1, n), \mathit{cgj} \in \operatorname{interval}(1, t). \\ \quad \quad \quad \quad \operatorname{entry}(\operatorname{matmul}(r, \operatorname{matmul}(r, p, u), q), \mathit{cgi}, \mathit{cgj}) = \operatorname{entry}(\operatorname{identmat}(r, n), \mathit{cgi}, \mathit{cgj}) \wedge \\ \quad \quad \quad \exists s \in \operatorname{carr}(r). \\ \quad \quad \quad \quad \operatorname{mul}(r)(\operatorname{det}(r, n, \operatorname{matmul}(r, \operatorname{matmul}(r, p, u), q)), s) = \operatorname{one}(r) \end{array}}}
\vnbprop{euclidean-unit-det-elem-gen}{\vnbstmt{\begin{array}{@{}l@{}} \forall r, n, u. \\ \quad \operatorname{is\text{-}euclidean\text{-}ring}(r) \wedge \\ \quad n \in \mathbb{N} \wedge \\ \quad u \in \operatorname{mat}(n, n, \operatorname{carr}(r)) \wedge \\ \quad \exists s \in \operatorname{carr}(r).\; (\operatorname{mul}(r)(\operatorname{det}(r, n, u), s) = \operatorname{one}(r)) \Rightarrow \\ \quad \quad u \in \operatorname{elem\text{-}gen}(r, n) \end{array}}}
\vnbprop{euclidean-invertible-iff-elem-gen}{\vnbstmt{\begin{array}{@{}l@{}} \forall r, n, u. \\ \quad (\operatorname{is\text{-}euclidean\text{-}ring}(r) \wedge (n \in \mathbb{N})) \Rightarrow  \\ \quad \quad \operatorname{is\text{-}invertible\text{-}mat}(r, n, u) \Leftrightarrow (u \in \operatorname{elem\text{-}gen}(r, n)) \end{array}}}
\vnbprop{zz-invertible-iff-elementary-product}{\vnbstmt{\begin{array}{@{}l@{}} \forall n \in \mathbb{N}, u. \\ \quad \operatorname{is\text{-}invertible\text{-}mat}(\operatorname{zz\text{-}ring}, n, u) \\ \quad \Leftrightarrow\; \exists p, q \in \operatorname{elem\text{-}gen}(\operatorname{zz\text{-}ring}, n). \\ \quad \quad u = \operatorname{matmul}(\operatorname{zz\text{-}ring}, p, q) \end{array}}}

\vnbbookfile{theorem-library/spans-transport-proof}{theorem-library/spans-transport-proof}
\noindent{\small spans-transport-proof.scm -- the SPANS-relativized generates-transport.}\par
\vnbprop{spans-transport}{\vnbstmt{\begin{array}{@{}l@{}} \forall d, n, u, m, a. \\ \quad \operatorname{is\text{-}module}(d) \wedge \\ \quad u \in \operatorname{mat}(n, 1, \operatorname{vec}(d)) \wedge \\ \quad \operatorname{is\text{-}invertible\text{-}mat}(\operatorname{scal}(d), n, m) \wedge \\ \quad \operatorname{is\text{-}submodule}(d, a) \wedge \\ \quad \operatorname{spans}(d, n, u, a) \Rightarrow \\ \quad \quad \operatorname{spans}(d, n, \operatorname{matact}(d, m, u), a) \end{array}}}

\vnbbookfile{theorem-library/rake-asm-b4}{theorem-library/rake-asm-b4}
\noindent{\small theorem-library/rake-asm-b4.scm  (RB-4, rake-asm-b, 2026-10-12) ==================================================================== submodule-of-free-is-free (Cor 3.46, linear-algebra.tex cor:22195-45888-207248): over a euclidean ring, a submodule sm of a module}\par
\vnbprop{rb4-finsum-ag-pad-ptwise}{\vnbstmt{\begin{array}{@{}l@{}} \forall a.; \forall t \in \mathbf{Set}, r, g. \\ \quad \operatorname{is\text{-}abelian\text{-}group}(a) \wedge \\ \quad |t| \in \mathbb{N} \wedge \\ \quad \forall z \in r.\; (z \in t) \wedge \\ \quad \forall z \in t.\; (g(z) \in \operatorname{carr}(a)) \wedge \\ \quad \forall z \in t.\; (\neg (z \in r) \Rightarrow (g(z) = \operatorname{iden}(a))) \Rightarrow \\ \quad \quad \operatorname{finsum}(a, g, t) = \operatorname{finsum}(a, g, r) \end{array}}}
\vnbprop{rb4-lincomb-trunc}{\vnbstmt{\begin{array}{@{}l@{}} \forall d, m, k, c, y. \\ \quad \operatorname{is\text{-}module}(d) \wedge \\ \quad c \in \operatorname{mat}(1, m, \operatorname{carr}(\operatorname{scal}(d))) \wedge \\ \quad y \in \operatorname{mat}(m, 1, \operatorname{vec}(d)) \wedge \\ \quad k \in \mathbb{N} \wedge \\ \quad k \leq m \wedge \\ \quad \forall j \in \operatorname{interval}(1, m). \\ \quad \quad \neg (j \leq k) \Rightarrow (\operatorname{act}(d)(\operatorname{entry}(c, 1, j), \operatorname{entry}(y, j, 1)) = \operatorname{vzero}(d)) \Rightarrow \\ \quad \quad \operatorname{lincomb}(d, m, c, y) = \operatorname{lincomb}(d, k, c, y) \end{array}}}
\vnbprop{rb4-act-entry-single}{\vnbstmt{\begin{array}{@{}l@{}} \forall d, m, n, p, e, i. \\ \quad \operatorname{is\text{-}module}(d) \wedge \\ \quad p \in \operatorname{mat}(m, n, \operatorname{carr}(\operatorname{scal}(d))) \wedge \\ \quad e \in \operatorname{mat}(n, 1, \operatorname{vec}(d)) \wedge \\ \quad i \in \operatorname{interval}(1, m) \wedge \\ \quad i \in \operatorname{interval}(1, n) \wedge \\ \quad \forall j \in \operatorname{interval}(1, n).\; (\neg (j = i) \Rightarrow (\operatorname{entry}(p, i, j) = \operatorname{zero}(\operatorname{scal}(d)))) \Rightarrow \\ \quad \quad \operatorname{entry}(\operatorname{matact}(d, p, e), i, 1) = \operatorname{act}(d)(\operatorname{entry}(p, i, i), \operatorname{entry}(e, i, 1)) \end{array}}}
\vnbprop{rb4-act-entry-zero-row}{\vnbstmt{\begin{array}{@{}l@{}} \forall d, m, n, p, e, i. \\ \quad \operatorname{is\text{-}module}(d) \wedge \\ \quad p \in \operatorname{mat}(m, n, \operatorname{carr}(\operatorname{scal}(d))) \wedge \\ \quad e \in \operatorname{mat}(n, 1, \operatorname{vec}(d)) \wedge \\ \quad i \in \operatorname{interval}(1, m) \wedge \\ \quad 1 \leq n \wedge \\ \quad \forall j \in \operatorname{interval}(1, n).\; (\operatorname{entry}(p, i, j) = \operatorname{zero}(\operatorname{scal}(d))) \Rightarrow \\ \quad \quad \operatorname{entry}(\operatorname{matact}(d, p, e), i, 1) = \operatorname{vzero}(d) \end{array}}}
\vnbprop{rb4-matmul-entry-single}{\vnbstmt{\begin{array}{@{}l@{}} \forall a, m, n, k, p, q, w, l. \\ \quad \operatorname{is\text{-}ring}(a) \wedge \\ \quad p \in \operatorname{mat}(m, n, \operatorname{carr}(a)) \wedge \\ \quad q \in \operatorname{mat}(n, k, \operatorname{carr}(a)) \wedge \\ \quad w \in \operatorname{interval}(1, m) \wedge \\ \quad l \in \operatorname{interval}(1, n) \wedge \\ \quad l \in \operatorname{interval}(1, k) \wedge \\ \quad \forall i \in \operatorname{interval}(1, n).\; (\neg (i = l) \Rightarrow (\operatorname{entry}(q, i, l) = \operatorname{zero}(a))) \Rightarrow \\ \quad \quad \operatorname{entry}(\operatorname{matmul}(a, p, q), w, l) = \operatorname{mul}(a)(\operatorname{entry}(p, w, l), \operatorname{entry}(q, l, l)) \end{array}}}
\vnbprop{rb4-diag-rel-free}{\vnbstmt{\begin{array}{@{}l@{}} \forall a, m, n, k, d, e, w. \\ \quad \operatorname{is\text{-}module}(a) \wedge \\ \quad \operatorname{is\text{-}integral\text{-}domain}(\operatorname{scal}(a)) \wedge \\ \quad m \in \mathbb{N} \wedge \\ \quad n \in \mathbb{N} \wedge \\ \quad k \in \mathbb{N} \wedge \\ \quad k \leq m \wedge \\ \quad k \leq n \wedge \\ \quad d \in \operatorname{mat}(m, n, \operatorname{carr}(\operatorname{scal}(a))) \wedge \\ \quad \operatorname{is\text{-}diagonal}(\operatorname{scal}(a), m, n, d) \wedge \\ \quad \forall i \in \operatorname{interval}(1, k).\; \neg (\operatorname{entry}(d, i, i) = \operatorname{zero}(\operatorname{scal}(a))) \wedge \\ \quad e \in \operatorname{mat}(n, 1, \operatorname{vec}(a)) \wedge \\ \quad \operatorname{rel\text{-}free}(a, n, e) \wedge \\ \quad w \in \operatorname{mat}(k, 1, \operatorname{vec}(a)) \wedge \\ \quad \forall i \in \operatorname{interval}(1, k). \\ \quad \quad \operatorname{entry}(w, i, 1) = \operatorname{act}(a)(\operatorname{entry}(d, i, i), \operatorname{entry}(e, i, 1)) \Rightarrow \\ \quad \quad \operatorname{rel\text{-}free}(a, k, w) \end{array}}}
\vnbprop{rb4-span-trunc}{\vnbstmt{\begin{array}{@{}l@{}} \forall d, m, k, y, a. \\ \quad \operatorname{is\text{-}module}(d) \wedge \\ \quad m \in \mathbb{N} \wedge \\ \quad k \in \mathbb{N} \wedge \\ \quad k \leq m \wedge \\ \quad y \in \operatorname{mat}(m, 1, \operatorname{vec}(d)) \wedge \\ \quad \operatorname{spans}(d, m, y, a) \wedge \\ \quad \forall j \in \operatorname{interval}(1, m).\; (\neg (j \leq k) \Rightarrow (\operatorname{entry}(y, j, 1) = \operatorname{vzero}(d))) \Rightarrow \\ \quad \quad \operatorname{spans}(d, k, \operatorname{block}(y, k, 1), a) \end{array}}}
\vnbprop{rb4-equiv-act}{\vnbstmt{\begin{array}{@{}l@{}} \forall d, m, n, a, b, s. \\ \quad \operatorname{is\text{-}module}(d) \wedge \\ \quad m \in \mathbb{N} \wedge \\ \quad n \in \mathbb{N} \wedge \\ \quad 1 \leq m \wedge \\ \quad 1 \leq n \wedge \\ \quad a \in \operatorname{mat}(m, n, \operatorname{carr}(\operatorname{scal}(d))) \wedge \\ \quad s \in \operatorname{mat}(n, 1, \operatorname{vec}(d)) \wedge \\ \quad \operatorname{mat\text{-}equiv}(\operatorname{scal}(d), m, n, a, b) \Rightarrow \\ \quad \quad b \in \operatorname{mat}(m, n, \operatorname{carr}(\operatorname{scal}(d))) \wedge \\ \quad \quad \exists u. \\ \quad \quad \quad \operatorname{is\text{-}invertible\text{-}mat}(\operatorname{scal}(d), m, u) \wedge \\ \quad \quad \quad \exists w. \\ \quad \quad \quad \quad (\operatorname{is\text{-}invertible\text{-}mat}(\operatorname{scal}(d), n, w) \wedge (\operatorname{matact}(d, u, \operatorname{matact}(d, a, s)) = \operatorname{matact}(d, b, \operatorname{matact}(d, w, s)))) \end{array}}}
\vnbprop{submodule-of-free-is-free}{\vnbstmt{\begin{array}{@{}l@{}} \forall d.; \forall n \in \mathbb{N}, u \in \operatorname{mat}(n, 1, \operatorname{vec}(d)), m. \\ \quad \operatorname{is\text{-}module}(d) \wedge \\ \quad \operatorname{is\text{-}euclidean\text{-}ring}(\operatorname{scal}(d)) \wedge \\ \quad \operatorname{generates}(d, n, u) \wedge \\ \quad \operatorname{rel\text{-}free}(d, n, u) \wedge \\ \quad \operatorname{is\text{-}submodule}(d, m) \Rightarrow \\ \quad \quad \exists k. \\ \quad \quad \quad (((k \in \mathbb{N}) \wedge (k \leq n)) \wedge \exists w \in \operatorname{mat}(k, 1, \operatorname{vec}(d)).\; (\operatorname{spans}(d, k, w, m) \wedge \operatorname{rel\text{-}free}(d, k, w))) \end{array}}}

\vnbbookfile{theorem-library/la-z2-action}{theorem-library/la-z2-action}
\noindent{\small la-z2-action.scm -- linear algebra 3.43: the invertible matrices act freely and transitively on the free generating sequences (LA-Z2, 2026-10-23).}\par
\vnbprop{lza-ring-sub-zero}{\vnbstmt{\begin{array}{@{}l@{}} \forall s, a, b. \\ \quad (\operatorname{is\text{-}ring}(s) \wedge ((a \in \operatorname{carr}(s)) \wedge ((b \in \operatorname{carr}(s)) \wedge (\operatorname{add}(s)(a, \operatorname{neg}(s)(b)) = \operatorname{zero}(s))))) \Rightarrow  \\ \quad \quad a = b \end{array}}}
\vnbprop{lza-vadd-neg-one}{\vnbstmt{\begin{array}{@{}l@{}} \forall d.; \forall x \in \operatorname{vec}(d). \\ \quad \operatorname{is\text{-}module}(d) \Rightarrow  \\ \quad \quad \operatorname{vadd}(d)(x, \operatorname{act}(d)(\operatorname{neg}(\operatorname{scal}(d))(\operatorname{one}(\operatorname{scal}(d))), x)) = \operatorname{vzero}(d) \end{array}}}
\vnbprop{unitrow-matmul-entry}{\vnbstmt{\begin{array}{@{}l@{}} \forall a, m, n, b, i, j. \\ \quad (\operatorname{is\text{-}ring}(a) \wedge ((b \in \operatorname{mat}(m, n, \operatorname{carr}(a))) \wedge ((i \in \operatorname{interval}(1, m)) \wedge (j \in \operatorname{interval}(1, n))))) \Rightarrow  \\ \quad \quad \operatorname{entry}(\operatorname{matmul}(a, \operatorname{unitrow}(a, m, i), b), 1, j) = \operatorname{entry}(b, i, j) \end{array}}}
\vnbprop{row-matact-lincomb}{\vnbstmt{\begin{array}{@{}l@{}} \forall d, m, n, a, s, i. \\ \quad \operatorname{is\text{-}module}(d) \wedge \\ \quad a \in \operatorname{mat}(m, n, \operatorname{carr}(\operatorname{scal}(d))) \wedge \\ \quad s \in \operatorname{mat}(n, 1, \operatorname{vec}(d)) \wedge \\ \quad i \in \operatorname{interval}(1, m) \wedge \\ \quad 1 \leq n \Rightarrow \\ \quad \quad \operatorname{lincomb}(d, n, \operatorname{matmul}(\operatorname{scal}(d), \operatorname{unitrow}(\operatorname{scal}(d), m, i), a), s) = \operatorname{entry}(\operatorname{matact}(d, a, s), i, 1) \end{array}}}
\vnbprop{rel-free-lincomb-cancel}{\vnbstmt{\begin{array}{@{}l@{}} \forall d, n, r, a, s.; \forall j \in \operatorname{interval}(1, n). \\ \quad \operatorname{is\text{-}module}(d) \wedge \\ \quad r \in \operatorname{mat}(1, n, \operatorname{carr}(\operatorname{scal}(d))) \wedge \\ \quad a \in \operatorname{mat}(1, n, \operatorname{carr}(\operatorname{scal}(d))) \wedge \\ \quad s \in \operatorname{mat}(n, 1, \operatorname{vec}(d)) \wedge \\ \quad \operatorname{rel\text{-}free}(d, n, s) \wedge \\ \quad \operatorname{lincomb}(d, n, r, s) = \operatorname{lincomb}(d, n, a, s) \Rightarrow \\ \quad \quad \operatorname{entry}(r, 1, j) = \operatorname{entry}(a, 1, j) \end{array}}}
\vnbprop{rel-free-matact-cancel}{\vnbstmt{\begin{array}{@{}l@{}} \forall d, m, n, s, a, b. \\ \quad \operatorname{is\text{-}module}(d) \wedge \\ \quad s \in \operatorname{mat}(n, 1, \operatorname{vec}(d)) \wedge \\ \quad a \in \operatorname{mat}(m, n, \operatorname{carr}(\operatorname{scal}(d))) \wedge \\ \quad b \in \operatorname{mat}(m, n, \operatorname{carr}(\operatorname{scal}(d))) \wedge \\ \quad \operatorname{rel\text{-}free}(d, n, s) \wedge \\ \quad \operatorname{matact}(d, a, s) = \operatorname{matact}(d, b, s) \Rightarrow \\ \quad \quad a = b \end{array}}}
\vnbprop{spans-coeff-matrix}{\vnbstmt{\begin{array}{@{}l@{}} \forall d, n, m, s, a, b. \\ \quad \operatorname{is\text{-}module}(d) \wedge \\ \quad s \in \operatorname{mat}(n, 1, \operatorname{vec}(d)) \wedge \\ \quad a \in \operatorname{mat}(m, 1, \operatorname{vec}(d)) \wedge \\ \quad (n = 0) \Rightarrow (m = 0) \wedge \\ \quad \operatorname{spans}(d, n, s, b) \wedge \\ \quad \forall j \in \operatorname{interval}(1, m).\; (\operatorname{entry}(a, j, 1) \in b) \Rightarrow \\ \quad \quad \exists c \in \operatorname{mat}(m, n, \operatorname{carr}(\operatorname{scal}(d))).\; (a = \operatorname{matact}(d, c, s)) \end{array}}}
\vnbprop{spans-rel-free-transitive}{\vnbstmt{\begin{array}{@{}l@{}} \forall d, n, s, a, m. \\ \quad \operatorname{is\text{-}module}(d) \wedge \\ \quad s \in \operatorname{mat}(n, 1, \operatorname{vec}(d)) \wedge \\ \quad a \in \operatorname{mat}(n, 1, \operatorname{vec}(d)) \wedge \\ \quad \operatorname{spans}(d, n, s, m) \wedge \\ \quad \operatorname{rel\text{-}free}(d, n, s) \wedge \\ \quad \operatorname{spans}(d, n, a, m) \wedge \\ \quad \operatorname{rel\text{-}free}(d, n, a) \Rightarrow \\ \quad \quad \exists b.\; (\operatorname{is\text{-}invertible\text{-}mat}(\operatorname{scal}(d), n, b) \wedge (a = \operatorname{matact}(d, b, s))) \end{array}}}
\vnbprop{free-generating-action-closed}{\vnbstmt{\begin{array}{@{}l@{}} \forall d, n, s, m. \\ \quad \operatorname{is\text{-}module}(d) \wedge \\ \quad s \in \operatorname{mat}(n, 1, \operatorname{vec}(d)) \wedge \\ \quad \operatorname{is\text{-}invertible\text{-}mat}(\operatorname{scal}(d), n, m) \wedge \\ \quad \operatorname{generates}(d, n, s) \wedge \\ \quad \operatorname{rel\text{-}free}(d, n, s) \Rightarrow \\ \quad \quad \operatorname{generates}(d, n, \operatorname{matact}(d, m, s)) \wedge \operatorname{rel\text{-}free}(d, n, \operatorname{matact}(d, m, s)) \end{array}}}
\vnbprop{free-generating-action-free}{\vnbstmt{\begin{array}{@{}l@{}} \forall d, n, s, m, a. \\ \quad \operatorname{is\text{-}module}(d) \wedge \\ \quad s \in \operatorname{mat}(n, 1, \operatorname{vec}(d)) \wedge \\ \quad \operatorname{generates}(d, n, s) \wedge \\ \quad \operatorname{rel\text{-}free}(d, n, s) \wedge \\ \quad \operatorname{is\text{-}invertible\text{-}mat}(\operatorname{scal}(d), n, m) \wedge \\ \quad \operatorname{is\text{-}invertible\text{-}mat}(\operatorname{scal}(d), n, a) \wedge \\ \quad \operatorname{matact}(d, m, s) = \operatorname{matact}(d, a, s) \Rightarrow \\ \quad \quad m = a \end{array}}}
\vnbprop{free-generating-action-stabilizer}{\vnbstmt{\begin{array}{@{}l@{}} \forall d, n, s, m. \\ \quad \operatorname{is\text{-}module}(d) \wedge \\ \quad s \in \operatorname{mat}(n, 1, \operatorname{vec}(d)) \wedge \\ \quad \operatorname{generates}(d, n, s) \wedge \\ \quad \operatorname{rel\text{-}free}(d, n, s) \wedge \\ \quad \operatorname{is\text{-}invertible\text{-}mat}(\operatorname{scal}(d), n, m) \wedge \\ \quad \operatorname{matact}(d, m, s) = s \Rightarrow \\ \quad \quad m = \operatorname{identmat}(\operatorname{scal}(d), n) \end{array}}}
\vnbprop{free-generating-action-transitive}{\vnbstmt{\begin{array}{@{}l@{}} \forall d, n, s, a. \\ \quad \operatorname{is\text{-}module}(d) \wedge \\ \quad s \in \operatorname{mat}(n, 1, \operatorname{vec}(d)) \wedge \\ \quad a \in \operatorname{mat}(n, 1, \operatorname{vec}(d)) \wedge \\ \quad \operatorname{generates}(d, n, s) \wedge \\ \quad \operatorname{rel\text{-}free}(d, n, s) \wedge \\ \quad \operatorname{generates}(d, n, a) \wedge \\ \quad \operatorname{rel\text{-}free}(d, n, a) \Rightarrow \\ \quad \quad \exists m.\; (\operatorname{is\text{-}invertible\text{-}mat}(\operatorname{scal}(d), n, m) \wedge (a = \operatorname{matact}(d, m, s))) \end{array}}}

\vnbbookfile{theorem-library/la-z2-norm}{theorem-library/la-z2-norm}
\noindent{\small la-z2-norm.scm -- linear algebra 3.45: |det A| of a change of basis between a submodule F and the module E does not depend on the bases (LA-Z2, 2026-10-23).}\par
\vnbprop{lzn-ring-cancel-unit}{\vnbstmt{\begin{array}{@{}l@{}} \forall r, a, b, c, d, e. \\ \quad \operatorname{is\text{-}commutative\text{-}ring}(r) \wedge \\ \quad a \in \operatorname{carr}(r) \wedge \\ \quad b \in \operatorname{carr}(r) \wedge \\ \quad c \in \operatorname{carr}(r) \wedge \\ \quad d \in \operatorname{carr}(r) \wedge \\ \quad e \in \operatorname{carr}(r) \wedge \\ \quad \operatorname{mul}(r)(a, b) = \operatorname{mul}(r)(c, d) \wedge \\ \quad \operatorname{mul}(r)(b, e) = \operatorname{one}(r) \Rightarrow \\ \quad \quad a = \operatorname{mul}(r)(\operatorname{mul}(r)(c, e), d) \end{array}}}
\vnbprop{change-of-basis-det-unit}{\vnbstmt{\begin{array}{@{}l@{}} \forall d, n, s, s_{2}, a, b, m, c, m_{2}. \\ \quad \operatorname{is\text{-}module}(d) \wedge \\ \quad \operatorname{is\text{-}commutative\text{-}ring}(\operatorname{scal}(d)) \wedge \\ \quad s \in \operatorname{mat}(n, 1, \operatorname{vec}(d)) \wedge \\ \quad s_{2} \in \operatorname{mat}(n, 1, \operatorname{vec}(d)) \wedge \\ \quad a \in \operatorname{mat}(n, 1, \operatorname{vec}(d)) \wedge \\ \quad b \in \operatorname{mat}(n, 1, \operatorname{vec}(d)) \wedge \\ \quad \operatorname{generates}(d, n, s) \wedge \\ \quad \operatorname{rel\text{-}free}(d, n, s) \wedge \\ \quad \operatorname{generates}(d, n, s_{2}) \wedge \\ \quad \operatorname{rel\text{-}free}(d, n, s_{2}) \wedge \\ \quad \operatorname{spans}(d, n, a, m) \wedge \\ \quad \operatorname{rel\text{-}free}(d, n, a) \wedge \\ \quad \operatorname{spans}(d, n, b, m) \wedge \\ \quad \operatorname{rel\text{-}free}(d, n, b) \wedge \\ \quad c \in \operatorname{mat}(n, n, \operatorname{carr}(\operatorname{scal}(d))) \wedge \\ \quad m_{2} \in \operatorname{mat}(n, n, \operatorname{carr}(\operatorname{scal}(d))) \wedge \\ \quad a = \operatorname{matact}(d, c, s) \wedge \\ \quad b = \operatorname{matact}(d, m_{2}, s_{2}) \Rightarrow \\ \quad \quad \exists w \in \operatorname{carr}(\operatorname{scal}(d)). \\ \quad \quad \quad \exists v \in \operatorname{carr}(\operatorname{scal}(d)).\; (\operatorname{mul}(\operatorname{scal}(d))(w, v) = \operatorname{one}(\operatorname{scal}(d))) \wedge \\ \quad \quad \quad \operatorname{det}(\operatorname{scal}(d), n, m_{2}) = \operatorname{mul}(\operatorname{scal}(d))(w, \operatorname{det}(\operatorname{scal}(d), n, c)) \end{array}}}
\vnbprop{norm-independent}{\vnbstmt{\begin{array}{@{}l@{}} \forall d, n, s, s_{2}, a, b, m, c, m_{2}. \\ \quad \operatorname{is\text{-}module}(d) \wedge \\ \quad \operatorname{scal}(d) = \operatorname{zz\text{-}ring} \wedge \\ \quad s \in \operatorname{mat}(n, 1, \operatorname{vec}(d)) \wedge \\ \quad s_{2} \in \operatorname{mat}(n, 1, \operatorname{vec}(d)) \wedge \\ \quad a \in \operatorname{mat}(n, 1, \operatorname{vec}(d)) \wedge \\ \quad b \in \operatorname{mat}(n, 1, \operatorname{vec}(d)) \wedge \\ \quad \operatorname{generates}(d, n, s) \wedge \\ \quad \operatorname{rel\text{-}free}(d, n, s) \wedge \\ \quad \operatorname{generates}(d, n, s_{2}) \wedge \\ \quad \operatorname{rel\text{-}free}(d, n, s_{2}) \wedge \\ \quad \operatorname{spans}(d, n, a, m) \wedge \\ \quad \operatorname{rel\text{-}free}(d, n, a) \wedge \\ \quad \operatorname{spans}(d, n, b, m) \wedge \\ \quad \operatorname{rel\text{-}free}(d, n, b) \wedge \\ \quad c \in \operatorname{mat}(n, n, \mathbb{Z}) \wedge \\ \quad m_{2} \in \operatorname{mat}(n, n, \mathbb{Z}) \wedge \\ \quad a = \operatorname{matact}(d, c, s) \wedge \\ \quad b = \operatorname{matact}(d, m_{2}, s_{2}) \Rightarrow \\ \quad \quad \lvert \operatorname{det}(\operatorname{zz\text{-}ring}, n, m_{2}) \rvert = \lvert \operatorname{det}(\operatorname{zz\text{-}ring}, n, c) \rvert \end{array}}}

\vnbbookfile{theorem-library/la-z2-diag}{theorem-library/la-z2-diag}
\noindent{\small la-z2-diag.scm -- the uniqueness half of linear algebra 3.36: the gcd of the order-s subdeterminants of a diagonal matrix with the divisibility chain is b\_11 ... b\_ss, so two such matrices equivalent by integral elementary operations are equal (LA-Z2, 2026-10-23).}\par
\vnbprop{interval-injection-bound}{\vnbstmt{\begin{array}{@{}l@{}} \forall s, m, t, p. \\ \quad s \in \mathbb{N} \wedge \\ \quad m \in \mathbb{N} \wedge \\ \quad t \in \mathbb{N} \wedge \\ \quad p \in \operatorname{injection}(\operatorname{interval}(1, s), \operatorname{interval}(1, m)) \wedge \\ \quad \forall i \in \operatorname{interval}(1, s).\; (p(i) \leq t) \Rightarrow \\ \quad \quad s \leq t \end{array}}}
\vnbprop{lzd-skip-type}{\vnbstmt{\begin{array}{@{}l@{}} \forall s, a, u. \\ \quad ((s \in \mathbb{N}) \wedge ((a \in \mathbb{N}) \wedge (u \in \operatorname{interval}(1, s)))) \Rightarrow  \\ \quad \quad \left(\text{if }(u < a)\text{ then }u\text{ else }(u + 1)\right) \in \operatorname{interval}(1, (s + 1)) \end{array}}}
\vnbprop{lzd-skip-inj}{\vnbstmt{\begin{array}{@{}l@{}} \forall a, u, v. \\ \quad ((a \in \mathbb{N}) \wedge ((u \in \mathbb{N}) \wedge ((v \in \mathbb{N}) \wedge (\left(\text{if }(u < a)\text{ then }u\text{ else }(u + 1)\right) = \left(\text{if }(v < a)\text{ then }v\text{ else }(v + 1)\right))))) \Rightarrow  \\ \quad \quad u = v \end{array}}}
\vnbprop{lzd-skp-fun}{\vnbstmt{\begin{array}{@{}l@{}} \forall s, m, p, a. \\ \quad ((s \in \mathbb{N}) \wedge ((m \in \mathbb{N}) \wedge ((p \in (\operatorname{interval}(1, (s + 1)) \to \operatorname{interval}(1, m))) \wedge (a \in \mathbb{N})))) \Rightarrow  \\ \quad \quad (\lambda x \in \operatorname{interval}(1, s).\; p(\left(\text{if }(x < a)\text{ then }x\text{ else }(x + 1)\right))) \in (\operatorname{interval}(1, s) \to \operatorname{interval}(1, m)) \end{array}}}
\vnbprop{lzd-skp-inj}{\vnbstmt{\begin{array}{@{}l@{}} \forall s, m, p, a. \\ \quad ((s \in \mathbb{N}) \wedge ((m \in \mathbb{N}) \wedge ((p \in \operatorname{injection}(\operatorname{interval}(1, (s + 1)), \operatorname{interval}(1, m))) \wedge (a \in \mathbb{N})))) \Rightarrow  \\ \quad \quad (\lambda x \in \operatorname{interval}(1, s).\; p(\left(\text{if }(x < a)\text{ then }x\text{ else }(x + 1)\right))) \in \operatorname{injection}(\operatorname{interval}(1, s), \operatorname{interval}(1, m)) \end{array}}}
\vnbprop{lzd-minor-submatrix}{\vnbstmt{\begin{array}{@{}l@{}} \forall r, m, n, s, a, p, q, b, c. \\ \quad \operatorname{is\text{-}ring}(r) \wedge \\ \quad m \in \mathbb{N} \wedge \\ \quad n \in \mathbb{N} \wedge \\ \quad s \in \mathbb{N} \wedge \\ \quad a \in \operatorname{mat}(m, n, \operatorname{carr}(r)) \wedge \\ \quad p \in (\operatorname{interval}(1, (s + 1)) \to \operatorname{interval}(1, m)) \wedge \\ \quad q \in (\operatorname{interval}(1, (s + 1)) \to \operatorname{interval}(1, n)) \wedge \\ \quad b \in \operatorname{interval}(1, (s + 1)) \wedge \\ \quad c \in \operatorname{interval}(1, (s + 1)) \Rightarrow \\ \quad \quad \operatorname{minor}(\operatorname{submatrix}(a, (s + 1), p, q), b, c, s) \\ \quad \quad =\; \operatorname{submatrix}(a, s, \\ \quad \quad \quad (\lambda x \in \operatorname{interval}(1, s).\; p(\left(\text{if }(x < b)\text{ then }x\text{ else }(x + 1)\right))), \\ \quad \quad \quad (\lambda x \in \operatorname{interval}(1, s).\; q(\left(\text{if }(x < c)\text{ then }x\text{ else }(x + 1)\right)))) \end{array}}}
\vnbprop{diagonal-chain-minor-divisible}{\vnbstmt{\begin{array}{@{}l@{}} \forall s \in \mathbb{N}, r, m, n, d, p, q. \\ \quad \operatorname{is\text{-}commutative\text{-}ring}(r) \wedge \\ \quad m \in \mathbb{N} \wedge \\ \quad n \in \mathbb{N} \wedge \\ \quad d \in \operatorname{mat}(m, n, \operatorname{carr}(r)) \wedge \\ \quad \operatorname{is\text{-}diagonal}(r, m, n, d) \wedge \\ \quad \forall i, j. \\ \quad \quad i \in \operatorname{interval}(1, m) \wedge \\ \quad \quad i \in \operatorname{interval}(1, n) \wedge \\ \quad \quad j \in \operatorname{interval}(1, m) \wedge \\ \quad \quad j \in \operatorname{interval}(1, n) \wedge \\ \quad \quad i \leq j \Rightarrow \\ \quad \quad \quad \exists y \in \operatorname{carr}(r).\; (\operatorname{entry}(d, j, j) = \operatorname{mul}(r)(\operatorname{entry}(d, i, i), y)) \wedge \\ \quad p \in \operatorname{injection}(\operatorname{interval}(1, s), \operatorname{interval}(1, m)) \wedge \\ \quad q \in \operatorname{injection}(\operatorname{interval}(1, s), \operatorname{interval}(1, n)) \Rightarrow \\ \quad \quad \exists y \in \operatorname{carr}(r). \\ \quad \quad \quad \operatorname{det}(r, s, \operatorname{submatrix}(d, s, p, q)) \\ \quad \quad \quad =\; \operatorname{mul}(r)(\operatorname{prod\text{-}ring}(r, (\lambda k \in \operatorname{interval}(1, s).\; \operatorname{entry}(d, k, k)), \operatorname{interval}(1, s)), y) \end{array}}}
\vnbprop{lzd-id-injection}{\vnbstmt{\begin{array}{@{}l@{}} \forall s, m. \\ \quad ((s \in \mathbb{N}) \wedge ((m \in \mathbb{N}) \wedge (s \leq m))) \Rightarrow  \\ \quad \quad (\lambda x \in \operatorname{interval}(1, s).\; x) \in \operatorname{injection}(\operatorname{interval}(1, s), \operatorname{interval}(1, m)) \end{array}}}
\vnbprop{diagonal-leading-minor}{\vnbstmt{\begin{array}{@{}l@{}} \forall r, m, n, d, s. \\ \quad \operatorname{is\text{-}commutative\text{-}ring}(r) \wedge \\ \quad m \in \mathbb{N} \wedge \\ \quad n \in \mathbb{N} \wedge \\ \quad d \in \operatorname{mat}(m, n, \operatorname{carr}(r)) \wedge \\ \quad \operatorname{is\text{-}diagonal}(r, m, n, d) \wedge \\ \quad s \in \mathbb{N} \wedge \\ \quad s \leq m \wedge \\ \quad s \leq n \Rightarrow \\ \quad \quad \operatorname{det}(r, s, \operatorname{submatrix}(d, s, (\lambda x \in \operatorname{interval}(1, s).\; x), (\lambda x \in \operatorname{interval}(1, s).\; x))) \\ \quad \quad =\; \operatorname{prod\text{-}ring}(r, (\lambda k \in \operatorname{interval}(1, s).\; \operatorname{entry}(d, k, k)), \operatorname{interval}(1, s)) \end{array}}}
\vnbprop{diagonal-chain-sdet-gcd}{\vnbstmt{\begin{array}{@{}l@{}} \forall m, n, d, s. \\ \quad m \in \mathbb{N} \wedge \\ \quad n \in \mathbb{N} \wedge \\ \quad d \in \operatorname{mat}(m, n, \mathbb{Z}) \wedge \\ \quad \operatorname{is\text{-}diagonal}(\operatorname{zz\text{-}ring}, m, n, d) \wedge \\ \quad \forall i, j. \\ \quad \quad i \in \operatorname{interval}(1, m) \wedge \\ \quad \quad i \in \operatorname{interval}(1, n) \wedge \\ \quad \quad j \in \operatorname{interval}(1, m) \wedge \\ \quad \quad j \in \operatorname{interval}(1, n) \wedge \\ \quad \quad i \leq j \Rightarrow \\ \quad \quad \quad \operatorname{zz\text{-}divides}(\operatorname{entry}(d, i, i), \operatorname{entry}(d, j, j)) \wedge \\ \quad s \in \mathbb{N} \wedge \\ \quad s \leq m \wedge \\ \quad s \leq n \Rightarrow \\ \quad \quad \operatorname{zz\text{-}is\text{-}gcd\text{-}set}(\operatorname{sdet}(\operatorname{zz\text{-}ring}, m, n, d, s), \\ \quad \quad \quad \operatorname{prod\text{-}ring}(\operatorname{zz\text{-}ring}, (\lambda k \in \operatorname{interval}(1, s).\; \operatorname{entry}(d, k, k)), \operatorname{interval}(1, s))) \end{array}}}
\vnbprop{diagonal-product-succ}{\vnbstmt{\begin{array}{@{}l@{}} \forall r, m, n, d, s. \\ \quad \operatorname{is\text{-}commutative\text{-}ring}(r) \wedge \\ \quad m \in \mathbb{N} \wedge \\ \quad n \in \mathbb{N} \wedge \\ \quad d \in \operatorname{mat}(m, n, \operatorname{carr}(r)) \wedge \\ \quad s \in \mathbb{N} \wedge \\ \quad (s + 1) \leq m \wedge \\ \quad (s + 1) \leq n \Rightarrow \\ \quad \quad \operatorname{prod\text{-}ring}(r, (\lambda k \in \operatorname{interval}(1, (s + 1)).\; \operatorname{entry}(d, k, k)), \operatorname{interval}(1, (s + 1))) \\ \quad \quad =\; \operatorname{mul}(r)(\operatorname{prod\text{-}ring}(r, (\lambda k \in \operatorname{interval}(1, s).\; \operatorname{entry}(d, k, k)), \operatorname{interval}(1, s)), \\ \quad \quad \quad \operatorname{entry}(d, (s + 1), (s + 1))) \end{array}}}
\vnbprop{diagonal-chain-product-zero}{\vnbstmt{\begin{array}{@{}l@{}} \forall s \in \mathbb{N}, m, n, d. \\ \quad m \in \mathbb{N} \wedge \\ \quad n \in \mathbb{N} \wedge \\ \quad d \in \operatorname{mat}(m, n, \mathbb{Z}) \wedge \\ \quad \operatorname{is\text{-}diagonal}(\operatorname{zz\text{-}ring}, m, n, d) \wedge \\ \quad \forall i, j. \\ \quad \quad i \in \operatorname{interval}(1, m) \wedge \\ \quad \quad i \in \operatorname{interval}(1, n) \wedge \\ \quad \quad j \in \operatorname{interval}(1, m) \wedge \\ \quad \quad j \in \operatorname{interval}(1, n) \wedge \\ \quad \quad i \leq j \Rightarrow \\ \quad \quad \quad \operatorname{zz\text{-}divides}(\operatorname{entry}(d, i, i), \operatorname{entry}(d, j, j)) \wedge \\ \quad s \leq m \wedge \\ \quad s \leq n \wedge \\ \quad \operatorname{prod\text{-}ring}(\operatorname{zz\text{-}ring}, (\lambda k \in \operatorname{interval}(1, s).\; \operatorname{entry}(d, k, k)), \operatorname{interval}(1, s)) = 0 \Rightarrow \\ \quad \quad (1 \leq s) \wedge (\operatorname{entry}(d, s, s) = 0) \end{array}}}
\vnbprop{diagonal-chain-products-determine}{\vnbstmt{\begin{array}{@{}l@{}} \forall m, n, d, a. \\ \quad m \in \mathbb{N} \wedge \\ \quad n \in \mathbb{N} \wedge \\ \quad d \in \operatorname{mat}(m, n, \mathbb{Z}) \wedge \\ \quad a \in \operatorname{mat}(m, n, \mathbb{Z}) \wedge \\ \quad \operatorname{is\text{-}diagonal}(\operatorname{zz\text{-}ring}, m, n, d) \wedge \\ \quad \operatorname{is\text{-}diagonal}(\operatorname{zz\text{-}ring}, m, n, a) \wedge \\ \quad \forall i \in \operatorname{interval}(1, m).\; ((i \in \operatorname{interval}(1, n)) \Rightarrow (0 \leq \operatorname{entry}(d, i, i))) \wedge \\ \quad \forall i \in \operatorname{interval}(1, m).\; ((i \in \operatorname{interval}(1, n)) \Rightarrow (0 \leq \operatorname{entry}(a, i, i))) \wedge \\ \quad \forall \mathit{ci}, j. \\ \quad \quad \mathit{ci} \in \operatorname{interval}(1, m) \wedge \\ \quad \quad \mathit{ci} \in \operatorname{interval}(1, n) \wedge \\ \quad \quad j \in \operatorname{interval}(1, m) \wedge \\ \quad \quad j \in \operatorname{interval}(1, n) \wedge \\ \quad \quad \mathit{ci} \leq j \Rightarrow \\ \quad \quad \quad \operatorname{zz\text{-}divides}(\operatorname{entry}(d, \mathit{ci}, \mathit{ci}), \operatorname{entry}(d, j, j)) \wedge \\ \quad \forall \mathit{ci}, j. \\ \quad \quad \mathit{ci} \in \operatorname{interval}(1, m) \wedge \\ \quad \quad \mathit{ci} \in \operatorname{interval}(1, n) \wedge \\ \quad \quad j \in \operatorname{interval}(1, m) \wedge \\ \quad \quad j \in \operatorname{interval}(1, n) \wedge \\ \quad \quad \mathit{ci} \leq j \Rightarrow \\ \quad \quad \quad \operatorname{zz\text{-}divides}(\operatorname{entry}(a, \mathit{ci}, \mathit{ci}), \operatorname{entry}(a, j, j)) \wedge \\ \quad \forall s \in \mathbb{N}. \\ \quad \quad ((s \leq m) \wedge (s \leq n)) \Rightarrow  \\ \quad \quad \quad \operatorname{prod\text{-}ring}(\operatorname{zz\text{-}ring}, (\lambda k \in \operatorname{interval}(1, s).\; \operatorname{entry}(d, k, k)), \operatorname{interval}(1, s)) \\ \quad \quad \quad =\; \operatorname{prod\text{-}ring}(\operatorname{zz\text{-}ring}, (\lambda k \in \operatorname{interval}(1, s).\; \operatorname{entry}(a, k, k)), \operatorname{interval}(1, s)) \vee \\ \quad \quad \quad \operatorname{prod\text{-}ring}(\operatorname{zz\text{-}ring}, (\lambda k \in \operatorname{interval}(1, s).\; \operatorname{entry}(d, k, k)), \operatorname{interval}(1, s)) \\ \quad \quad \quad =\; -\operatorname{prod\text{-}ring}(\operatorname{zz\text{-}ring}, (\lambda k \in \operatorname{interval}(1, s).\; \operatorname{entry}(a, k, k)), \operatorname{interval}(1, s)) \Rightarrow \\ \quad \quad d = a \end{array}}}
\vnbprop{diagonal-normal-form-unique}{\vnbstmt{\begin{array}{@{}l@{}} \forall m, n, d, a. \\ \quad m \in \mathbb{N} \wedge \\ \quad n \in \mathbb{N} \wedge \\ \quad d \in \operatorname{mat}(m, n, \mathbb{Z}) \wedge \\ \quad a \in \operatorname{mat}(m, n, \mathbb{Z}) \wedge \\ \quad \operatorname{is\text{-}diagonal}(\operatorname{zz\text{-}ring}, m, n, d) \wedge \\ \quad \operatorname{is\text{-}diagonal}(\operatorname{zz\text{-}ring}, m, n, a) \wedge \\ \quad \forall i \in \operatorname{interval}(1, m).\; ((i \in \operatorname{interval}(1, n)) \Rightarrow (0 \leq \operatorname{entry}(d, i, i))) \wedge \\ \quad \forall i \in \operatorname{interval}(1, m).\; ((i \in \operatorname{interval}(1, n)) \Rightarrow (0 \leq \operatorname{entry}(a, i, i))) \wedge \\ \quad \forall \mathit{ci}, j. \\ \quad \quad \mathit{ci} \in \operatorname{interval}(1, m) \wedge \\ \quad \quad \mathit{ci} \in \operatorname{interval}(1, n) \wedge \\ \quad \quad j \in \operatorname{interval}(1, m) \wedge \\ \quad \quad j \in \operatorname{interval}(1, n) \wedge \\ \quad \quad \mathit{ci} \leq j \Rightarrow \\ \quad \quad \quad \operatorname{zz\text{-}divides}(\operatorname{entry}(d, \mathit{ci}, \mathit{ci}), \operatorname{entry}(d, j, j)) \wedge \\ \quad \forall \mathit{ci}, j. \\ \quad \quad \mathit{ci} \in \operatorname{interval}(1, m) \wedge \\ \quad \quad \mathit{ci} \in \operatorname{interval}(1, n) \wedge \\ \quad \quad j \in \operatorname{interval}(1, m) \wedge \\ \quad \quad j \in \operatorname{interval}(1, n) \wedge \\ \quad \quad \mathit{ci} \leq j \Rightarrow \\ \quad \quad \quad \operatorname{zz\text{-}divides}(\operatorname{entry}(a, \mathit{ci}, \mathit{ci}), \operatorname{entry}(a, j, j)) \wedge \\ \quad \operatorname{elem\text{-}equiv}(\operatorname{zz\text{-}ring}, m, n, d, a) \Rightarrow \\ \quad \quad d = a \end{array}}}
\vnbprop{submatrix-mul-diagonal}{\vnbstmt{\begin{array}{@{}l@{}} \forall r, k, n, s, x, d, p, q. \\ \quad \operatorname{is\text{-}ring}(r) \wedge \\ \quad k \in \mathbb{N} \wedge \\ \quad n \in \mathbb{N} \wedge \\ \quad s \in \mathbb{N} \wedge \\ \quad x \in \operatorname{mat}(k, n, \operatorname{carr}(r)) \wedge \\ \quad d \in \operatorname{mat}(n, n, \operatorname{carr}(r)) \wedge \\ \quad \operatorname{is\text{-}diagonal}(r, n, n, d) \wedge \\ \quad p \in (\operatorname{interval}(1, s) \to \operatorname{interval}(1, k)) \wedge \\ \quad q \in \operatorname{injection}(\operatorname{interval}(1, s), \operatorname{interval}(1, n)) \Rightarrow \\ \quad \quad \operatorname{submatrix}(\operatorname{matmul}(r, x, d), s, p, q) = \operatorname{matmul}(r, \operatorname{submatrix}(x, s, p, q), \operatorname{submatrix}(d, s, q, q)) \end{array}}}
\vnbprop{diagonal-right-factor-sdet-divisible}{\vnbstmt{\begin{array}{@{}l@{}} \forall r, k, n, s, x, d.; \forall \mathit{dvx} \in \operatorname{sdet}(r, k, n, \operatorname{matmul}(r, x, d), s). \\ \quad \operatorname{is\text{-}commutative\text{-}ring}(r) \wedge \\ \quad k \in \mathbb{N} \wedge \\ \quad n \in \mathbb{N} \wedge \\ \quad s \in \mathbb{N} \wedge \\ \quad x \in \operatorname{mat}(k, n, \operatorname{carr}(r)) \wedge \\ \quad d \in \operatorname{mat}(n, n, \operatorname{carr}(r)) \wedge \\ \quad \operatorname{is\text{-}diagonal}(r, n, n, d) \wedge \\ \quad \forall i, j. \\ \quad \quad i \in \operatorname{interval}(1, n) \wedge \\ \quad \quad i \in \operatorname{interval}(1, n) \wedge \\ \quad \quad j \in \operatorname{interval}(1, n) \wedge \\ \quad \quad j \in \operatorname{interval}(1, n) \wedge \\ \quad \quad i \leq j \Rightarrow \\ \quad \quad \quad \exists y \in \operatorname{carr}(r).\; (\operatorname{entry}(d, j, j) = \operatorname{mul}(r)(\operatorname{entry}(d, i, i), y)) \wedge \\ \quad s \leq n \Rightarrow \\ \quad \quad \exists y \in \operatorname{carr}(r). \\ \quad \quad \quad \mathit{dvx} = \operatorname{mul}(r)(\operatorname{prod\text{-}ring}(r, (\lambda \mathit{pk} \in \operatorname{interval}(1, s).\; \operatorname{entry}(d, \mathit{pk}, \mathit{pk})), \operatorname{interval}(1, s)), y) \end{array}}}
\vnbprop{diagonal-normal-form-unique-invertible}{\vnbstmt{\begin{array}{@{}l@{}} \forall m, n, d, a, b, c. \\ \quad m \in \mathbb{N} \wedge \\ \quad n \in \mathbb{N} \wedge \\ \quad d \in \operatorname{mat}(m, n, \mathbb{Z}) \wedge \\ \quad a \in \operatorname{mat}(m, n, \mathbb{Z}) \wedge \\ \quad \operatorname{is\text{-}diagonal}(\operatorname{zz\text{-}ring}, m, n, d) \wedge \\ \quad \operatorname{is\text{-}diagonal}(\operatorname{zz\text{-}ring}, m, n, a) \wedge \\ \quad \forall i \in \operatorname{interval}(1, m).\; ((i \in \operatorname{interval}(1, n)) \Rightarrow (0 \leq \operatorname{entry}(d, i, i))) \wedge \\ \quad \forall i \in \operatorname{interval}(1, m).\; ((i \in \operatorname{interval}(1, n)) \Rightarrow (0 \leq \operatorname{entry}(a, i, i))) \wedge \\ \quad \forall \mathit{ci}, j. \\ \quad \quad \mathit{ci} \in \operatorname{interval}(1, m) \wedge \\ \quad \quad \mathit{ci} \in \operatorname{interval}(1, n) \wedge \\ \quad \quad j \in \operatorname{interval}(1, m) \wedge \\ \quad \quad j \in \operatorname{interval}(1, n) \wedge \\ \quad \quad \mathit{ci} \leq j \Rightarrow \\ \quad \quad \quad \operatorname{zz\text{-}divides}(\operatorname{entry}(d, \mathit{ci}, \mathit{ci}), \operatorname{entry}(d, j, j)) \wedge \\ \quad \forall \mathit{ci}, j. \\ \quad \quad \mathit{ci} \in \operatorname{interval}(1, m) \wedge \\ \quad \quad \mathit{ci} \in \operatorname{interval}(1, n) \wedge \\ \quad \quad j \in \operatorname{interval}(1, m) \wedge \\ \quad \quad j \in \operatorname{interval}(1, n) \wedge \\ \quad \quad \mathit{ci} \leq j \Rightarrow \\ \quad \quad \quad \operatorname{zz\text{-}divides}(\operatorname{entry}(a, \mathit{ci}, \mathit{ci}), \operatorname{entry}(a, j, j)) \wedge \\ \quad \operatorname{is\text{-}invertible\text{-}mat}(\operatorname{zz\text{-}ring}, m, b) \wedge \\ \quad \operatorname{is\text{-}invertible\text{-}mat}(\operatorname{zz\text{-}ring}, n, c) \wedge \\ \quad a = \operatorname{matmul}(\operatorname{zz\text{-}ring}, \operatorname{matmul}(\operatorname{zz\text{-}ring}, b, d), c) \Rightarrow \\ \quad \quad d = a \end{array}}}

\vnbbookfile{theorem-library/la-z2-principal}{theorem-library/la-z2-principal}
\noindent{\small la-z2-principal.scm -- linear algebra 3.48: the principal sequence of a subgroup is uniquely determined (LA-Z2, 2026-10-23).}\par
\vnbprop{principal-sequence-unique}{\vnbstmt{\begin{array}{@{}l@{}} \forall d, n, m, s, s_{2}, a, m_{2}. \\ \quad \operatorname{is\text{-}module}(d) \wedge \\ \quad \operatorname{scal}(d) = \operatorname{zz\text{-}ring} \wedge \\ \quad n \in \mathbb{N} \wedge \\ \quad \operatorname{is\text{-}principal\text{-}sequence}(d, n, m, s, a) \wedge \\ \quad \operatorname{is\text{-}principal\text{-}sequence}(d, n, m, s_{2}, m_{2}) \Rightarrow \\ \quad \quad a = m_{2} \end{array}}}

\vnbbookfile{theorem-library/la-z3-smith}{theorem-library/la-z3-smith}
\noindent{\small la-z3-smith.scm -- linear algebra 3.36 (existence) and 3.37 in general: every integer matrix is equivalent, by integral elementary operations, to a diagonal matrix with non-negative diagonal entries forming a divisibility chain (LA-Z3, 2026-10-24).  Prefix lz3-.}\par
\vnbprop{smith-pivot-cross-minimal}{\vnbstmt{\begin{array}{@{}l@{}} \forall a, p, q, m. \\ \quad \operatorname{is\text{-}euclidean\text{-}ring}(a) \wedge \\ \quad p \in \mathbb{N} \wedge \\ \quad q \in \mathbb{N} \wedge \\ \quad m \in \operatorname{mat}((p + 1), (q + 1), \operatorname{carr}(a)) \wedge \\ \quad \exists b \in \operatorname{interval}(1, (p + 1)), d \in \operatorname{interval}(1, (q + 1)). \\ \quad \quad \neg (\operatorname{entry}(m, b, d) = \operatorname{zero}(a)) \Rightarrow \\ \quad \quad \exists f. \\ \quad \quad \quad \operatorname{mat\text{-}equiv}(a, (p + 1), (q + 1), m, f) \wedge \\ \quad \quad \quad f = \operatorname{border}(a, \operatorname{entry}(f, 1, 1), \operatorname{submat}(f, p, q), p, q) \wedge \\ \quad \quad \quad \neg (\operatorname{entry}(f, 1, 1) = \operatorname{zero}(a)) \wedge \\ \quad \quad \quad \forall c, i, j. \\ \quad \quad \quad \quad \operatorname{mat\text{-}equiv}(a, (p + 1), (q + 1), f, c) \wedge \\ \quad \quad \quad \quad i \in \operatorname{interval}(1, (p + 1)) \wedge \\ \quad \quad \quad \quad j \in \operatorname{interval}(1, (q + 1)) \wedge \\ \quad \quad \quad \quad \neg (\operatorname{entry}(c, i, j) = \operatorname{zero}(a)) \Rightarrow \\ \quad \quad \quad \quad \quad \operatorname{gauge}(a)(\operatorname{entry}(f, 1, 1)) \leq \operatorname{gauge}(a)(\operatorname{entry}(c, i, j)) \end{array}}}
\vnbprop{smith-pivot-divides-block}{\vnbstmt{\begin{array}{@{}l@{}} \forall r, p, q, m, i, j. \\ \quad \operatorname{is\text{-}euclidean\text{-}ring}(r) \wedge \\ \quad p \in \mathbb{N} \wedge \\ \quad q \in \mathbb{N} \wedge \\ \quad m \in \operatorname{mat}((p + 1), (q + 1), \operatorname{carr}(r)) \wedge \\ \quad m = \operatorname{border}(r, \operatorname{entry}(m, 1, 1), \operatorname{submat}(m, p, q), p, q) \wedge \\ \quad \neg (\operatorname{entry}(m, 1, 1) = \operatorname{zero}(r)) \wedge \\ \quad \forall c, a, b. \\ \quad \quad \operatorname{mat\text{-}equiv}(r, (p + 1), (q + 1), m, c) \wedge \\ \quad \quad a \in \operatorname{interval}(1, (p + 1)) \wedge \\ \quad \quad b \in \operatorname{interval}(1, (q + 1)) \wedge \\ \quad \quad \neg (\operatorname{entry}(c, a, b) = \operatorname{zero}(r)) \Rightarrow \\ \quad \quad \quad \operatorname{gauge}(r)(\operatorname{entry}(m, 1, 1)) \leq \operatorname{gauge}(r)(\operatorname{entry}(c, a, b)) \wedge \\ \quad i \in \operatorname{interval}(1, p) \wedge \\ \quad j \in \operatorname{interval}(1, q) \Rightarrow \\ \quad \quad \exists y \in \operatorname{carr}(r). \\ \quad \quad \quad \operatorname{entry}(\operatorname{submat}(m, p, q), i, j) = \operatorname{mul}(r)(\operatorname{entry}(m, 1, 1), y) \end{array}}}
\vnbprop{lz3-const-injection}{\vnbstmt{\begin{array}{@{}l@{}} \forall m, i. \\ \quad ((m \in \mathbb{N}) \wedge (i \in \operatorname{interval}(1, m))) \Rightarrow  \\ \quad \quad (\lambda x \in \operatorname{interval}(1, 1).\; i) \in \operatorname{injection}(\operatorname{interval}(1, 1), \operatorname{interval}(1, m)) \end{array}}}
\vnbprop{smith-equiv-divisor}{\vnbstmt{\begin{array}{@{}l@{}} \forall r, m, n, p, e, c, i, j. \\ \quad \operatorname{is\text{-}euclidean\text{-}ring}(r) \wedge \\ \quad m \in \mathbb{N} \wedge \\ \quad n \in \mathbb{N} \wedge \\ \quad p \in \operatorname{mat}(m, n, \operatorname{carr}(r)) \wedge \\ \quad \operatorname{mat\text{-}equiv}(r, m, n, p, e) \wedge \\ \quad c \in \operatorname{carr}(r) \wedge \\ \quad \forall \mathit{dvi}, \mathit{dvj}. \\ \quad \quad ((\mathit{dvi} \in \operatorname{interval}(1, m)) \wedge (\mathit{dvj} \in \operatorname{interval}(1, n))) \Rightarrow  \\ \quad \quad \quad \exists y \in \operatorname{carr}(r).\; (\operatorname{entry}(p, \mathit{dvi}, \mathit{dvj}) = \operatorname{mul}(r)(c, y)) \wedge \\ \quad i \in \operatorname{interval}(1, m) \wedge \\ \quad j \in \operatorname{interval}(1, n) \Rightarrow \\ \quad \quad \exists y \in \operatorname{carr}(r).\; (\operatorname{entry}(e, i, j) = \operatorname{mul}(r)(c, y)) \end{array}}}
\vnbprop{border-scale-row}{\vnbstmt{\begin{array}{@{}l@{}} \forall r, p, q, b, c, e. \\ \quad (\operatorname{is\text{-}ring}(r) \wedge ((p \in \mathbb{N}) \wedge ((q \in \mathbb{N}) \wedge ((b \in \operatorname{carr}(r)) \wedge ((c \in \operatorname{carr}(r)) \wedge (e \in \operatorname{mat}(p, q, \operatorname{carr}(r)))))))) \Rightarrow  \\ \quad \quad \operatorname{matmul}(r, \operatorname{elem\text{-}h}(r, (p + 1), c, 1), \operatorname{border}(r, b, e, p, q)) = \operatorname{border}(r, \operatorname{mul}(r)(c, b), e, p, q) \end{array}}}
\vnbprop{border-normal-form}{\vnbstmt{\begin{array}{@{}l@{}} \forall p, q, b, e. \\ \quad p \in \mathbb{N} \wedge \\ \quad q \in \mathbb{N} \wedge \\ \quad b \in \mathbb{Z} \wedge \\ \quad 0 \leq b \wedge \\ \quad e \in \operatorname{mat}(p, q, \operatorname{carr}(\operatorname{zz\text{-}ring})) \wedge \\ \quad \operatorname{is\text{-}diagonal}(\operatorname{zz\text{-}ring}, p, q, e) \wedge \\ \quad \forall i \in \operatorname{interval}(1, p).\; ((i \in \operatorname{interval}(1, q)) \Rightarrow (0 \leq \operatorname{entry}(e, i, i))) \wedge \\ \quad \forall \mathit{ci}, j. \\ \quad \quad \mathit{ci} \in \operatorname{interval}(1, p) \wedge \\ \quad \quad \mathit{ci} \in \operatorname{interval}(1, q) \wedge \\ \quad \quad j \in \operatorname{interval}(1, p) \wedge \\ \quad \quad j \in \operatorname{interval}(1, q) \wedge \\ \quad \quad \mathit{ci} \leq j \Rightarrow \\ \quad \quad \quad \operatorname{zz\text{-}divides}(\operatorname{entry}(e, \mathit{ci}, \mathit{ci}), \operatorname{entry}(e, j, j)) \wedge \\ \quad \forall \mathit{dvi}, \mathit{dvj}. \\ \quad \quad ((\mathit{dvi} \in \operatorname{interval}(1, p)) \wedge (\mathit{dvj} \in \operatorname{interval}(1, q))) \Rightarrow  \\ \quad \quad \quad \operatorname{zz\text{-}divides}(b, \operatorname{entry}(e, \mathit{dvi}, \mathit{dvj})) \Rightarrow \\ \quad \quad \operatorname{is\text{-}diagonal}(\operatorname{zz\text{-}ring}, (p + 1), (q + 1), \operatorname{border}(\operatorname{zz\text{-}ring}, b, e, p, q)) \wedge \\ \quad \quad \forall i \in \operatorname{interval}(1, (p + 1)). \\ \quad \quad \quad (i \in \operatorname{interval}(1, (q + 1))) \Rightarrow (0 \leq \operatorname{entry}(\operatorname{border}(\operatorname{zz\text{-}ring}, b, e, p, q), i, i)) \wedge \\ \quad \quad \forall \mathit{ci}, j. \\ \quad \quad \quad \mathit{ci} \in \operatorname{interval}(1, (p + 1)) \wedge \\ \quad \quad \quad \mathit{ci} \in \operatorname{interval}(1, (q + 1)) \wedge \\ \quad \quad \quad j \in \operatorname{interval}(1, (p + 1)) \wedge \\ \quad \quad \quad j \in \operatorname{interval}(1, (q + 1)) \wedge \\ \quad \quad \quad \mathit{ci} \leq j \Rightarrow \\ \quad \quad \quad \quad \operatorname{zz\text{-}divides}(\operatorname{entry}(\operatorname{border}(\operatorname{zz\text{-}ring}, b, e, p, q), \mathit{ci}, \mathit{ci}), \operatorname{entry}(\operatorname{border}(\operatorname{zz\text{-}ring}, b, e, p, q), j, j)) \end{array}}}
\vnbprop{integer-normal-form-mat-equiv}{\vnbstmt{\begin{array}{@{}l@{}} \forall k \in \mathbb{N}, n, p. \\ \quad ((n \in \mathbb{N}) \wedge (p \in \operatorname{mat}(k, n, \operatorname{carr}(\operatorname{zz\text{-}ring})))) \Rightarrow  \\ \quad \quad \exists d \in \operatorname{mat}(k, n, \operatorname{carr}(\operatorname{zz\text{-}ring})). \\ \quad \quad \quad \operatorname{mat\text{-}equiv}(\operatorname{zz\text{-}ring}, k, n, p, d) \wedge \\ \quad \quad \quad \operatorname{is\text{-}diagonal}(\operatorname{zz\text{-}ring}, k, n, d) \wedge \\ \quad \quad \quad \forall i \in \operatorname{interval}(1, k).\; ((i \in \operatorname{interval}(1, n)) \Rightarrow (0 \leq \operatorname{entry}(d, i, i))) \wedge \\ \quad \quad \quad \forall \mathit{ci}, j. \\ \quad \quad \quad \quad \mathit{ci} \in \operatorname{interval}(1, k) \wedge \\ \quad \quad \quad \quad \mathit{ci} \in \operatorname{interval}(1, n) \wedge \\ \quad \quad \quad \quad j \in \operatorname{interval}(1, k) \wedge \\ \quad \quad \quad \quad j \in \operatorname{interval}(1, n) \wedge \\ \quad \quad \quad \quad \mathit{ci} \leq j \Rightarrow \\ \quad \quad \quad \quad \quad \operatorname{zz\text{-}divides}(\operatorname{entry}(d, \mathit{ci}, \mathit{ci}), \operatorname{entry}(d, j, j)) \end{array}}}
\vnbprop{integer-normal-form-elementary}{\vnbstmt{\begin{array}{@{}l@{}} \forall m, n, a. \\ \quad ((m \in \mathbb{N}) \wedge ((n \in \mathbb{N}) \wedge (a \in \operatorname{mat}(m, n, \mathbb{Z})))) \Rightarrow  \\ \quad \quad \exists u \in \operatorname{elem\text{-}gen}(\operatorname{zz\text{-}ring}, m), v \in \operatorname{elem\text{-}gen}(\operatorname{zz\text{-}ring}, n). \\ \quad \quad \quad \operatorname{is\text{-}diagonal}(\operatorname{zz\text{-}ring}, m, n, \operatorname{matmul}(\operatorname{zz\text{-}ring}, \operatorname{matmul}(\operatorname{zz\text{-}ring}, u, a), v)) \wedge \\ \quad \quad \quad \forall i \in \operatorname{interval}(1, m). \\ \quad \quad \quad \quad (i \in \operatorname{interval}(1, n)) \Rightarrow  \\ \quad \quad \quad \quad \quad 0 \leq \operatorname{entry}(\operatorname{matmul}(\operatorname{zz\text{-}ring}, \operatorname{matmul}(\operatorname{zz\text{-}ring}, u, a), v), i, i) \wedge \\ \quad \quad \quad \forall \mathit{ci}, j. \\ \quad \quad \quad \quad \mathit{ci} \in \operatorname{interval}(1, m) \wedge \\ \quad \quad \quad \quad \mathit{ci} \in \operatorname{interval}(1, n) \wedge \\ \quad \quad \quad \quad j \in \operatorname{interval}(1, m) \wedge \\ \quad \quad \quad \quad j \in \operatorname{interval}(1, n) \wedge \\ \quad \quad \quad \quad \mathit{ci} \leq j \Rightarrow \\ \quad \quad \quad \quad \quad \operatorname{zz\text{-}divides}(\operatorname{entry}(\operatorname{matmul}(\operatorname{zz\text{-}ring}, \operatorname{matmul}(\operatorname{zz\text{-}ring}, u, a), v), \mathit{ci}, \mathit{ci}), \\ \quad \quad \quad \quad \quad \quad \operatorname{entry}(\operatorname{matmul}(\operatorname{zz\text{-}ring}, \operatorname{matmul}(\operatorname{zz\text{-}ring}, u, a), v), j, j)) \end{array}}}
\vnbprop{integer-normal-form-exists}{\vnbstmt{\begin{array}{@{}l@{}} \forall m, n, a. \\ \quad ((m \in \mathbb{N}) \wedge ((n \in \mathbb{N}) \wedge (a \in \operatorname{mat}(m, n, \mathbb{Z})))) \Rightarrow  \\ \quad \quad \exists d \in \operatorname{mat}(m, n, \mathbb{Z}). \\ \quad \quad \quad \operatorname{elem\text{-}equiv}(\operatorname{zz\text{-}ring}, m, n, a, d) \wedge \\ \quad \quad \quad \operatorname{is\text{-}diagonal}(\operatorname{zz\text{-}ring}, m, n, d) \wedge \\ \quad \quad \quad \forall i \in \operatorname{interval}(1, m).\; ((i \in \operatorname{interval}(1, n)) \Rightarrow (0 \leq \operatorname{entry}(d, i, i))) \wedge \\ \quad \quad \quad \forall \mathit{ci}, j. \\ \quad \quad \quad \quad \mathit{ci} \in \operatorname{interval}(1, m) \wedge \\ \quad \quad \quad \quad \mathit{ci} \in \operatorname{interval}(1, n) \wedge \\ \quad \quad \quad \quad j \in \operatorname{interval}(1, m) \wedge \\ \quad \quad \quad \quad j \in \operatorname{interval}(1, n) \wedge \\ \quad \quad \quad \quad \mathit{ci} \leq j \Rightarrow \\ \quad \quad \quad \quad \quad \operatorname{zz\text{-}divides}(\operatorname{entry}(d, \mathit{ci}, \mathit{ci}), \operatorname{entry}(d, j, j)) \end{array}}}
\vnbprop{elem-equiv-trans}{\vnbstmt{\begin{array}{@{}l@{}} \forall r, m, n, a, b, c. \\ \quad \operatorname{is\text{-}ring}(r) \wedge \\ \quad m \in \mathbb{N} \wedge \\ \quad n \in \mathbb{N} \wedge \\ \quad a \in \operatorname{mat}(m, n, \operatorname{carr}(r)) \wedge \\ \quad \operatorname{elem\text{-}equiv}(r, m, n, a, b) \wedge \\ \quad \operatorname{elem\text{-}equiv}(r, m, n, b, c) \Rightarrow \\ \quad \quad \operatorname{elem\text{-}equiv}(r, m, n, a, c) \end{array}}}
\vnbprop{integer-normal-form}{\vnbstmt{\begin{array}{@{}l@{}} \forall m, n, a. \\ \quad ((m \in \mathbb{N}) \wedge ((n \in \mathbb{N}) \wedge (a \in \operatorname{mat}(m, n, \mathbb{Z})))) \Rightarrow  \\ \quad \quad \exists d \in \operatorname{mat}(m, n, \mathbb{Z}). \\ \quad \quad \quad \operatorname{elem\text{-}equiv}(\operatorname{zz\text{-}ring}, m, n, a, d) \wedge \\ \quad \quad \quad \operatorname{is\text{-}diagonal}(\operatorname{zz\text{-}ring}, m, n, d) \wedge \\ \quad \quad \quad \forall i \in \operatorname{interval}(1, m).\; ((i \in \operatorname{interval}(1, n)) \Rightarrow (0 \leq \operatorname{entry}(d, i, i))) \wedge \\ \quad \quad \quad \forall \mathit{ci}, j. \\ \quad \quad \quad \quad \mathit{ci} \in \operatorname{interval}(1, m) \wedge \\ \quad \quad \quad \quad \mathit{ci} \in \operatorname{interval}(1, n) \wedge \\ \quad \quad \quad \quad j \in \operatorname{interval}(1, m) \wedge \\ \quad \quad \quad \quad j \in \operatorname{interval}(1, n) \wedge \\ \quad \quad \quad \quad \mathit{ci} \leq j \Rightarrow \\ \quad \quad \quad \quad \quad \operatorname{zz\text{-}divides}(\operatorname{entry}(d, \mathit{ci}, \mathit{ci}), \operatorname{entry}(d, j, j)) \wedge \\ \quad \quad \quad \forall b \in \operatorname{mat}(m, n, \mathbb{Z}). \\ \quad \quad \quad \quad \operatorname{elem\text{-}equiv}(\operatorname{zz\text{-}ring}, m, n, a, b) \wedge \\ \quad \quad \quad \quad \operatorname{is\text{-}diagonal}(\operatorname{zz\text{-}ring}, m, n, b) \wedge \\ \quad \quad \quad \quad \forall i \in \operatorname{interval}(1, m).\; ((i \in \operatorname{interval}(1, n)) \Rightarrow (0 \leq \operatorname{entry}(b, i, i))) \wedge \\ \quad \quad \quad \quad \forall \mathit{ci}, j. \\ \quad \quad \quad \quad \quad \mathit{ci} \in \operatorname{interval}(1, m) \wedge \\ \quad \quad \quad \quad \quad \mathit{ci} \in \operatorname{interval}(1, n) \wedge \\ \quad \quad \quad \quad \quad j \in \operatorname{interval}(1, m) \wedge \\ \quad \quad \quad \quad \quad j \in \operatorname{interval}(1, n) \wedge \\ \quad \quad \quad \quad \quad \mathit{ci} \leq j \Rightarrow \\ \quad \quad \quad \quad \quad \quad \operatorname{zz\text{-}divides}(\operatorname{entry}(b, \mathit{ci}, \mathit{ci}), \operatorname{entry}(b, j, j)) \Rightarrow \\ \quad \quad \quad \quad \quad b = d \end{array}}}
\vnbprop{spans-pad-zero}{\vnbstmt{\begin{array}{@{}l@{}} \forall j \in \mathbb{N}, d, m, k, w. \\ \quad \operatorname{is\text{-}module}(d) \wedge \\ \quad \operatorname{is\text{-}submodule}(d, m) \wedge \\ \quad k \in \mathbb{N} \wedge \\ \quad w \in \operatorname{mat}(k, 1, \operatorname{vec}(d)) \wedge \\ \quad \operatorname{spans}(d, k, w, m) \wedge \\ \quad k \leq j \Rightarrow \\ \quad \quad \exists a \in \operatorname{mat}(j, 1, \operatorname{vec}(d)).\; \operatorname{spans}(d, j, a, m) \end{array}}}
\vnbprop{principal-sequence-exists}{\vnbstmt{\begin{array}{@{}l@{}} \forall d, n, u, m. \\ \quad \operatorname{is\text{-}module}(d) \wedge \\ \quad \operatorname{scal}(d) = \operatorname{zz\text{-}ring} \wedge \\ \quad n \in \mathbb{N} \wedge \\ \quad u \in \operatorname{mat}(n, 1, \operatorname{vec}(d)) \wedge \\ \quad \operatorname{generates}(d, n, u) \wedge \\ \quad \operatorname{rel\text{-}free}(d, n, u) \wedge \\ \quad \operatorname{is\text{-}submodule}(d, m) \Rightarrow \\ \quad \quad \exists s, a.\; \operatorname{is\text{-}principal\text{-}sequence}(d, n, m, s, a) \end{array}}}
\vnbprop{principal-sequence-exists-unique}{\vnbstmt{\begin{array}{@{}l@{}} \forall d, n, u, m. \\ \quad \operatorname{is\text{-}module}(d) \wedge \\ \quad \operatorname{scal}(d) = \operatorname{zz\text{-}ring} \wedge \\ \quad n \in \mathbb{N} \wedge \\ \quad u \in \operatorname{mat}(n, 1, \operatorname{vec}(d)) \wedge \\ \quad \operatorname{generates}(d, n, u) \wedge \\ \quad \operatorname{rel\text{-}free}(d, n, u) \wedge \\ \quad \operatorname{is\text{-}submodule}(d, m) \Rightarrow \\ \quad \quad \exists s, a. \\ \quad \quad \quad \operatorname{is\text{-}principal\text{-}sequence}(d, n, m, s, a) \wedge \\ \quad \quad \quad \forall s_{2}, m_{2}.\; (\operatorname{is\text{-}principal\text{-}sequence}(d, n, m, s_{2}, m_{2}) \Rightarrow (a = m_{2})) \end{array}}}

\vnbbookfile{theorem-library/rank-bound-proof}{theorem-library/rank-bound-proof}
\noindent{\small rank-bound-proof.scm -- Proposition 3.41 (algebraic-numbers.pdf ch.3 sec 8.2), the payoff of the whole Smith arc:}\par
\vnbprop{free-length-le-generators}{\vnbstmt{\begin{array}{@{}l@{}} \forall d, n, m, u, v. \\ \quad \operatorname{is\text{-}module}(d) \wedge \\ \quad \operatorname{is\text{-}euclidean\text{-}ring}(\operatorname{scal}(d)) \wedge \\ \quad n \in \mathbb{N} \wedge \\ \quad m \in \mathbb{N} \wedge \\ \quad u \in \operatorname{mat}(n, 1, \operatorname{vec}(d)) \wedge \\ \quad v \in \operatorname{mat}(m, 1, \operatorname{vec}(d)) \wedge \\ \quad \operatorname{generates}(d, n, u) \wedge \\ \quad \operatorname{rel\text{-}free}(d, m, v) \Rightarrow \\ \quad \quad m \leq n \end{array}}}

\vnbbookfile{theorem-library/gal-smith}{theorem-library/gal-smith}
\noindent{\small gal-smith.scm -- invariance of the gcd of the s-minors (galois-theory notes, Corollary 1.24) and the uniqueness of the Smith normal form (Proposition 1.27) over a Euclidean domain (GALOIS, 2026-10-26).  Destination: theorem-library/gal-smith.scm, library algebra.  Load window: after la-z sdet files (sdet-divisor-equiv, ...}\par
\vnbprop{mutual-divides-associated}{\vnbstmt{\begin{array}{@{}l@{}} \forall s, a, b. \\ \quad (\operatorname{is\text{-}integral\text{-}domain}(s) \wedge (\operatorname{divides}(s, a, b) \wedge \operatorname{divides}(s, b, a))) \Rightarrow  \\ \quad \quad \operatorname{associated}(s, a, b) \end{array}}}
\vnbprop{sdet-gcd-associated}{\vnbstmt{\begin{array}{@{}l@{}} \forall s.; \forall m, n, g \in \mathbb{N}, a, b \in \operatorname{mat}(m, n, \operatorname{carr}(s)), d, e. \\ \quad \operatorname{is\text{-}euclidean\text{-}domain}(s) \wedge \\ \quad \operatorname{elem\text{-}equiv}(s, m, n, a, b) \wedge \\ \quad \operatorname{is\text{-}gcd\text{-}set}(s, \operatorname{sdet}(s, m, n, a, g), d) \wedge \\ \quad \operatorname{is\text{-}gcd\text{-}set}(s, \operatorname{sdet}(s, m, n, b, g), e) \Rightarrow \\ \quad \quad \operatorname{associated}(s, d, e) \end{array}}}
\vnbprop{diagonal-chain-product-gcd}{\vnbstmt{\begin{array}{@{}l@{}} \forall s.; \forall m, n \in \mathbb{N}, d \in \operatorname{mat}(m, n, \operatorname{carr}(s)), g \in \mathbb{N}. \\ \quad \operatorname{is\text{-}commutative\text{-}ring}(s) \wedge \\ \quad \operatorname{is\text{-}diagonal}(s, m, n, d) \wedge \\ \quad \forall i, j. \\ \quad \quad i \in \operatorname{interval}(1, m) \wedge \\ \quad \quad i \in \operatorname{interval}(1, n) \wedge \\ \quad \quad j \in \operatorname{interval}(1, m) \wedge \\ \quad \quad j \in \operatorname{interval}(1, n) \wedge \\ \quad \quad i \leq j \Rightarrow \\ \quad \quad \quad \operatorname{divides}(s, \operatorname{entry}(d, i, i), \operatorname{entry}(d, j, j)) \wedge \\ \quad g \leq m \wedge \\ \quad g \leq n \Rightarrow \\ \quad \quad \operatorname{is\text{-}gcd\text{-}set}(s, \operatorname{sdet}(s, m, n, d, g), \\ \quad \quad \quad \operatorname{prod\text{-}ring}(s, (\lambda k \in \operatorname{interval}(1, g).\; \operatorname{entry}(d, k, k)), \operatorname{interval}(1, g))) \end{array}}}
\vnbprop{diagonal-chain-zero-propagates}{\vnbstmt{\begin{array}{@{}l@{}} \forall k \in \mathbb{N}, s.; \forall m, n \in \mathbb{N}, d \in \operatorname{mat}(m, n, \operatorname{carr}(s)). \\ \quad \operatorname{is\text{-}integral\text{-}domain}(s) \wedge \\ \quad \forall i, j. \\ \quad \quad i \in \operatorname{interval}(1, m) \wedge \\ \quad \quad i \in \operatorname{interval}(1, n) \wedge \\ \quad \quad j \in \operatorname{interval}(1, m) \wedge \\ \quad \quad j \in \operatorname{interval}(1, n) \wedge \\ \quad \quad i \leq j \Rightarrow \\ \quad \quad \quad \operatorname{divides}(s, \operatorname{entry}(d, i, i), \operatorname{entry}(d, j, j)) \wedge \\ \quad (k + 1) \leq m \wedge \\ \quad (k + 1) \leq n \wedge \\ \quad \operatorname{prod\text{-}ring}(s, (\lambda \mathit{pk} \in \operatorname{interval}(1, k).\; \operatorname{entry}(d, \mathit{pk}, \mathit{pk})), \operatorname{interval}(1, k)) = \operatorname{zero}(s) \Rightarrow \\ \quad \quad \operatorname{entry}(d, (k + 1), (k + 1)) = \operatorname{zero}(s) \end{array}}}
\vnbprop{smith-normal-form-unique}{\vnbstmt{\begin{array}{@{}l@{}} \forall s.; \forall m, n \in \mathbb{N}, a, d, \mathit{dd} \in \operatorname{mat}(m, n, \operatorname{carr}(s)), k \in \mathbb{N}. \\ \quad \operatorname{is\text{-}euclidean\text{-}domain}(s) \wedge \\ \quad \operatorname{is\text{-}diagonal}(s, m, n, d) \wedge \\ \quad \forall i, j. \\ \quad \quad i \in \operatorname{interval}(1, m) \wedge \\ \quad \quad i \in \operatorname{interval}(1, n) \wedge \\ \quad \quad j \in \operatorname{interval}(1, m) \wedge \\ \quad \quad j \in \operatorname{interval}(1, n) \wedge \\ \quad \quad i \leq j \Rightarrow \\ \quad \quad \quad \operatorname{divides}(s, \operatorname{entry}(d, i, i), \operatorname{entry}(d, j, j)) \wedge \\ \quad \operatorname{elem\text{-}equiv}(s, m, n, a, d) \wedge \\ \quad \operatorname{is\text{-}diagonal}(s, m, n, \mathit{dd}) \wedge \\ \quad \forall i, j. \\ \quad \quad i \in \operatorname{interval}(1, m) \wedge \\ \quad \quad i \in \operatorname{interval}(1, n) \wedge \\ \quad \quad j \in \operatorname{interval}(1, m) \wedge \\ \quad \quad j \in \operatorname{interval}(1, n) \wedge \\ \quad \quad i \leq j \Rightarrow \\ \quad \quad \quad \operatorname{divides}(s, \operatorname{entry}(\mathit{dd}, i, i), \operatorname{entry}(\mathit{dd}, j, j)) \wedge \\ \quad \operatorname{elem\text{-}equiv}(s, m, n, a, \mathit{dd}) \wedge \\ \quad (k + 1) \leq m \wedge \\ \quad (k + 1) \leq n \Rightarrow \\ \quad \quad \operatorname{associated}(s, \operatorname{entry}(d, (k + 1), (k + 1)), \operatorname{entry}(\mathit{dd}, (k + 1), (k + 1))) \end{array}}}

\vnbbookfile{structure-library/gal-charpoly-vocab}{structure-library/gal-charpoly-vocab}
\noindent{\small gal-charpoly-vocab.scm -- the companion matrix and the characteristic polynomial (galois-theory notes, section 7, eq. (37), (38), Proposition 1.37) (GALOIS B6, 2026-10-26).  DEFINITIONS ONLY. Destination: structure-library/gal-charpoly-vocab.scm, library linear-algebra.  Load window: after matrix.scm (MATOF, ENTRY), ...}\par
\par\medskip\noindent\textbf{Definition} (\texttt{companion}).\;$\operatorname{companion}(f, a, n) :=$
\begin{equation*}
\begin{array}{@{}l@{}} \operatorname{matof}(n, n, \\ \quad (\lambda {[}\operatorname{gci\_}, \operatorname{gcj\_}] \in (\operatorname{interval}(1, n) \times \operatorname{interval}(1, n)).\; \\ \quad \quad \left(\text{if }(\operatorname{gcj\_} = n)\text{ then }\operatorname{neg}(f)(a(\operatorname{gci\_}))\text{ else }\left(\text{if }(\operatorname{gci\_} = (\operatorname{gcj\_} + 1))\text{ then }\operatorname{one}(f)\text{ else }\operatorname{zero}(f)\right)\right))) \end{array}
\end{equation*}

\par\medskip\noindent\textbf{Definition} (\texttt{monic-poly}).\;$\operatorname{monic\text{-}poly}(f, a, n) :=$
\begin{equation*}
(\lambda \operatorname{gmk\_} \in \mathbb{N}.\; \left(\text{if }(\operatorname{gmk\_} = n)\text{ then }\operatorname{one}(f)\text{ else }\left(\text{if }(\operatorname{gmk\_} < n)\text{ then }a((\operatorname{gmk\_} + 1))\text{ else }\operatorname{zero}(f)\right)\right))
\end{equation*}

\par\medskip\noindent\textbf{Definition} (\texttt{char-mat}).\;$\operatorname{char\text{-}mat}(f, a, n) :=$
\begin{equation*}
\begin{array}{@{}l@{}} \operatorname{matof}(n, n, \\ \quad (\lambda {[}\operatorname{gmi\_}, \operatorname{gmj\_}] \in (\operatorname{interval}(1, n) \times \operatorname{interval}(1, n)).\; \\ \quad \quad \operatorname{add}(\operatorname{poly}(f))(\left(\text{if }(\operatorname{gmi\_} = \operatorname{gmj\_})\text{ then }\operatorname{monomial}(f, \operatorname{one}(f), 1)\text{ else }\operatorname{zero}(\operatorname{poly}(f))\right), \\ \quad \quad \quad \operatorname{neg}(\operatorname{poly}(f))(\operatorname{monomial}(f, \operatorname{entry}(a, \operatorname{gmi\_}, \operatorname{gmj\_}), 0))))) \end{array}
\end{equation*}

\par\medskip\noindent\textbf{Definition} (\texttt{char-poly}).\;$\operatorname{char\text{-}poly}(f, a, n) :=$ $\operatorname{det}(\operatorname{poly}(f), n, \operatorname{char\text{-}mat}(f, a, n))$

\vnbbookfile{theorem-library/gal-charpoly}{theorem-library/gal-charpoly}
\noindent{\small gal-charpoly.scm -- the characteristic polynomial of a companion matrix (galois-theory notes, Proposition 1.37) (GALOIS B6, 2026-10-26).  Destination: theorem-library/gal-charpoly.scm, library linear-algebra.  Load window: after the determinant files (det-expand-row, det-1x1, det-upper-triangular, minor-entry), ...}\par
\vnbprop{companion-entry}{\vnbstmt{\begin{array}{@{}l@{}} \forall r.; \forall v \in \mathbb{N}, \mathit{av}.; \forall i, j \in \operatorname{interval}(1, v). \\ \quad (\operatorname{is\text{-}ring}(r) \wedge \forall i \in \operatorname{interval}(1, v).\; (\operatorname{av}(i) \in \operatorname{carr}(r))) \Rightarrow  \\ \quad \quad \operatorname{entry}(\operatorname{companion}(r, \mathit{av}, v), i, j) = \left(\text{if }(j = v)\text{ then }\operatorname{neg}(r)(\operatorname{av}(i))\text{ else }\left(\text{if }(i = (j + 1))\text{ then }\operatorname{one}(r)\text{ else }\operatorname{zero}(r)\right)\right) \end{array}}}
\vnbprop{companion-in-mat}{\vnbstmt{\begin{array}{@{}l@{}} \forall r.; \forall v \in \mathbb{N}, \mathit{av}. \\ \quad (\operatorname{is\text{-}ring}(r) \wedge \forall i \in \operatorname{interval}(1, v).\; (\operatorname{av}(i) \in \operatorname{carr}(r))) \Rightarrow  \\ \quad \quad \operatorname{companion}(r, \mathit{av}, v) \in \operatorname{mat}(v, v, \operatorname{carr}(r)) \end{array}}}
\vnbprop{char-mat-entry}{\vnbstmt{\begin{array}{@{}l@{}} \forall r.; \forall v \in \mathbb{N}, m \in \operatorname{mat}(v, v, \operatorname{carr}(r)). \\ \quad \operatorname{is\text{-}ring}(r) \Rightarrow  \\ \quad \quad \operatorname{char\text{-}mat}(r, m, v) \in \operatorname{mat}(v, v, \operatorname{carr}(\operatorname{poly}(r))) \wedge \\ \quad \quad \forall i, j \in \operatorname{interval}(1, v). \\ \quad \quad \quad \operatorname{entry}(\operatorname{char\text{-}mat}(r, m, v), i, j) \\ \quad \quad \quad =\; \operatorname{add}(\operatorname{poly}(r))(\left(\text{if }(i = j)\text{ then }\operatorname{monomial}(r, \operatorname{one}(r), 1)\text{ else }\operatorname{zero}(\operatorname{poly}(r))\right), \\ \quad \quad \quad \quad \operatorname{neg}(\operatorname{poly}(r))(\operatorname{monomial}(r, \operatorname{entry}(m, i, j), 0))) \end{array}}}
\vnbprop{poly-neg-monomial}{\vnbstmt{\begin{array}{@{}l@{}} \forall r.; \forall c \in \operatorname{carr}(r), e \in \mathbb{N}. \\ \quad \operatorname{is\text{-}commutative\text{-}ring}(r) \Rightarrow  \\ \quad \quad \operatorname{neg}(\operatorname{poly}(r))(\operatorname{monomial}(r, c, e)) = \operatorname{monomial}(r, \operatorname{neg}(r)(c), e) \end{array}}}
\vnbprop{poly-monomial-zero}{\vnbstmt{\forall r.; \forall e \in \mathbb{N}.\; (\operatorname{is\text{-}ring}(r) \Rightarrow (\operatorname{monomial}(r, \operatorname{zero}(r), e) = \operatorname{zero}(\operatorname{poly}(r))))}}
\vnbprop{monic-poly-in-poly}{\vnbstmt{\begin{array}{@{}l@{}} \forall r.; \forall v \in \mathbb{N}, \mathit{av}. \\ \quad (\operatorname{is\text{-}ring}(r) \wedge \forall i \in \operatorname{interval}(1, v).\; (\operatorname{av}(i) \in \operatorname{carr}(r))) \Rightarrow  \\ \quad \quad \operatorname{monic\text{-}poly}(r, \mathit{av}, v) \in \operatorname{carr}(\operatorname{poly}(r)) \end{array}}}
\vnbprop{monic-poly-succ}{\vnbstmt{\begin{array}{@{}l@{}} \forall r.; \forall v \in \mathbb{N}, \mathit{av}. \\ \quad (\operatorname{is\text{-}commutative\text{-}ring}(r) \wedge \forall i \in \operatorname{interval}(1, (v + 1)).\; (\operatorname{av}(i) \in \operatorname{carr}(r))) \Rightarrow  \\ \quad \quad \operatorname{monic\text{-}poly}(r, \mathit{av}, (v + 1)) \\ \quad \quad =\; \operatorname{add}(\operatorname{poly}(r))(\operatorname{mul}(\operatorname{poly}(r))(\operatorname{monomial}(r, \operatorname{one}(r), 1), \operatorname{monic\text{-}poly}(r, (\lambda h \in \mathbb{N}.\; \operatorname{av}((h + 1))), v)), \\ \quad \quad \quad \operatorname{monomial}(r, \operatorname{av}(1), 0)) \end{array}}}
\vnbprop{prod-ring-neg-one-sign}{\vnbstmt{\begin{array}{@{}l@{}} \forall v \in \mathbb{N}, r. \\ \quad \operatorname{is\text{-}commutative\text{-}ring}(r) \Rightarrow  \\ \quad \quad \operatorname{prod\text{-}ring}(r, (\lambda o \in \operatorname{interval}(1, v).\; \operatorname{neg}(r)(\operatorname{one}(r))), \operatorname{interval}(1, v)) \\ \quad \quad =\; \operatorname{mpow}(\operatorname{ring\text{-}multiplicative\text{-}monoid}(r), \operatorname{neg}(r)(\operatorname{one}(r)), v) \end{array}}}
\vnbprop{char-mat-companion-row1}{\vnbstmt{\begin{array}{@{}l@{}} \forall r.; \forall v \in \mathbb{N}, \mathit{av}. \\ \quad (\operatorname{is\text{-}commutative\text{-}ring}(r) \wedge ((1 \leq v) \wedge \forall i \in \operatorname{interval}(1, (v + 1)).\; (\operatorname{av}(i) \in \operatorname{carr}(r)))) \Rightarrow  \\ \quad \quad \operatorname{entry}(\operatorname{char\text{-}mat}(r, \operatorname{companion}(r, \mathit{av}, (v + 1)), (v + 1)), 1, 1) = \operatorname{monomial}(r, \operatorname{one}(r), 1) \wedge \\ \quad \quad \operatorname{entry}(\operatorname{char\text{-}mat}(r, \operatorname{companion}(r, \mathit{av}, (v + 1)), (v + 1)), 1, (v + 1)) = \operatorname{monomial}(r, \operatorname{av}(1), 0) \wedge \\ \quad \quad \forall j \in \operatorname{interval}(1, (v + 1)). \\ \quad \quad \quad (\neg (j = 1) \wedge \neg (j = (v + 1))) \Rightarrow  \\ \quad \quad \quad \quad \operatorname{entry}(\operatorname{char\text{-}mat}(r, \operatorname{companion}(r, \mathit{av}, (v + 1)), (v + 1)), 1, j) = \operatorname{zero}(\operatorname{poly}(r)) \end{array}}}
\vnbprop{char-mat-companion-minor11}{\vnbstmt{\begin{array}{@{}l@{}} \forall r.; \forall v \in \mathbb{N}, \mathit{av}. \\ \quad (\operatorname{is\text{-}commutative\text{-}ring}(r) \wedge ((1 \leq v) \wedge \forall i \in \operatorname{interval}(1, (v + 1)).\; (\operatorname{av}(i) \in \operatorname{carr}(r)))) \Rightarrow  \\ \quad \quad \operatorname{minor}(\operatorname{char\text{-}mat}(r, \operatorname{companion}(r, \mathit{av}, (v + 1)), (v + 1)), 1, 1, v) \\ \quad \quad =\; \operatorname{char\text{-}mat}(r, \operatorname{companion}(r, (\lambda h \in \mathbb{N}.\; \operatorname{av}((h + 1))), v), v) \end{array}}}
\vnbprop{char-mat-companion-corner-det}{\vnbstmt{\begin{array}{@{}l@{}} \forall r.; \forall v \in \mathbb{N}, \mathit{av}. \\ \quad (\operatorname{is\text{-}commutative\text{-}ring}(r) \wedge ((1 \leq v) \wedge \forall i \in \operatorname{interval}(1, (v + 1)).\; (\operatorname{av}(i) \in \operatorname{carr}(r)))) \Rightarrow  \\ \quad \quad \operatorname{det}(\operatorname{poly}(r), v, \operatorname{minor}(\operatorname{char\text{-}mat}(r, \operatorname{companion}(r, \mathit{av}, (v + 1)), (v + 1)), 1, (v + 1), v)) \\ \quad \quad =\; \operatorname{mpow}(\operatorname{ring\text{-}multiplicative\text{-}monoid}(\operatorname{poly}(r)), \operatorname{neg}(\operatorname{poly}(r))(\operatorname{one}(\operatorname{poly}(r))), v) \end{array}}}
\vnbprop{monic-poly-zero}{\vnbstmt{\begin{array}{@{}l@{}} \forall r, v. \\ \quad \operatorname{is\text{-}commutative\text{-}ring}(r) \Rightarrow (\operatorname{monic\text{-}poly}(r, v, 0) = \operatorname{one}(\operatorname{poly}(r))) \end{array}}}
\vnbprop{companion-congruence}{\vnbstmt{\begin{array}{@{}l@{}} \forall r.; \forall v \in \mathbb{N}, \mathit{av}, \mathit{bv}. \\ \quad \operatorname{is\text{-}ring}(r) \wedge \\ \quad \forall i \in \operatorname{interval}(1, v).\; (\operatorname{av}(i) \in \operatorname{carr}(r)) \wedge \\ \quad \forall i \in \operatorname{interval}(1, v).\; (\operatorname{bv}(i) \in \operatorname{carr}(r)) \wedge \\ \quad \forall i \in \operatorname{interval}(1, v).\; (\operatorname{av}(i) = \operatorname{bv}(i)) \Rightarrow \\ \quad \quad \operatorname{companion}(r, \mathit{av}, v) = \operatorname{companion}(r, \mathit{bv}, v) \end{array}}}
\vnbprop{monic-poly-congruence}{\vnbstmt{\begin{array}{@{}l@{}} \forall r.; \forall v \in \mathbb{N}, \mathit{av}, \mathit{bv}. \\ \quad \operatorname{is\text{-}ring}(r) \wedge \\ \quad \forall i \in \operatorname{interval}(1, v).\; (\operatorname{av}(i) \in \operatorname{carr}(r)) \wedge \\ \quad \forall i \in \operatorname{interval}(1, v).\; (\operatorname{bv}(i) \in \operatorname{carr}(r)) \wedge \\ \quad \forall i \in \operatorname{interval}(1, v).\; (\operatorname{av}(i) = \operatorname{bv}(i)) \Rightarrow \\ \quad \quad \operatorname{monic\text{-}poly}(r, \mathit{av}, v) = \operatorname{monic\text{-}poly}(r, \mathit{bv}, v) \end{array}}}
\vnbprop{charpoly-of-companion}{\vnbstmt{\begin{array}{@{}l@{}} \forall v \in \mathbb{N}, r.; \forall \mathit{av} \in (\operatorname{interval}(1, (v + 1)) \to \operatorname{carr}(r)). \\ \quad (\operatorname{is\text{-}commutative\text{-}ring}(r) \wedge \forall i \in \operatorname{interval}(1, (v + 1)).\; (\operatorname{av}(i) \in \operatorname{carr}(r))) \Rightarrow  \\ \quad \quad \operatorname{char\text{-}poly}(r, \operatorname{companion}(r, \mathit{av}, (v + 1)), (v + 1)) = \operatorname{monic\text{-}poly}(r, \mathit{av}, (v + 1)) \end{array}}}

\vnbbookfile{structure-library/gal-op-vocab}{structure-library/gal-op-vocab}
\noindent{\small gal-op-vocab.scm -- a linear operator T on a module E over F as an F[x]-module: the iterates T\textasciicircum{}k, the operator p(T), the T-cyclic subspace, cyclic vectors, monic, annihilating and minimum polynomials, the Krylov family v, Tv, ..., T\textasciicircum{}(n-1) v, and the matrix of an operator in a basis (galois-theory notes, section 7, ...}\par
\par\medskip\noindent\textbf{Definition} (\texttt{krylov}).\;$\operatorname{krylov}(\operatorname{gkr\text{-}md\_}, \operatorname{gkr\text{-}t\_}, \operatorname{gkr\text{-}v\_}, \operatorname{gkr\text{-}n\_}) :=$
\begin{equation*}
\begin{array}{@{}l@{}} \operatorname{matof}(\operatorname{gkr\text{-}n\_}, 1, \\ \quad (\lambda {[}\operatorname{gkri\_}, \operatorname{gkrj\_}] \in (\operatorname{interval}(1, \operatorname{gkr\text{-}n\_}) \times \operatorname{interval}(1, 1)).\; \\ \quad \quad \operatorname{op\text{-}iter}(\operatorname{gkr\text{-}md\_}, \operatorname{gkr\text{-}t\_}, \operatorname{gkr\text{-}v\_}, (\operatorname{gkri\_} - 1)))) \end{array}
\end{equation*}

\par\medskip\noindent\textbf{Definition} (\texttt{op-iter}).\;For all $\operatorname{goi\text{-}md\_}$, $\operatorname{goi\text{-}t\_}$ and $\operatorname{goi\text{-}w\_}$,  $(\operatorname{op\text{-}iter}(\operatorname{goi\text{-}md\_}, \operatorname{goi\text{-}t\_}, \operatorname{goi\text{-}w\_}, 0) \simeq \operatorname{goi\text{-}w\_})$\par
Let $\operatorname{goi\text{-}md\_}$, $\operatorname{goi\text{-}t\_}$ and $\operatorname{goi\text{-}w\_}$ be arbitrary. Let $m \in \mathbb{N}$. Then
\begin{equation*}
(\operatorname{op\text{-}iter}(\operatorname{goi\text{-}md\_}, \operatorname{goi\text{-}t\_}, \operatorname{goi\text{-}w\_}, \operatorname{succ}(m)) \simeq \operatorname{goi\text{-}t\_}(\operatorname{op\text{-}iter}(\operatorname{goi\text{-}md\_}, \operatorname{goi\text{-}t\_}, \operatorname{goi\text{-}w\_}, m)))
\end{equation*}

\par\medskip\noindent\textbf{Definition} (\texttt{poly-op}).\;$\operatorname{poly\text{-}op}(\operatorname{gpo\text{-}md\_}, \operatorname{gpo\text{-}t\_}, \operatorname{gpo\text{-}p\_}, \operatorname{gpo\text{-}w\_}) :=$
\begin{equation*}
\begin{array}{@{}l@{}} \operatorname{sum\text{-}ag}(\operatorname{module\text{-}vector\text{-}ag}(\operatorname{gpo\text{-}md\_}), \\ \quad (\lambda \operatorname{gpok\_} \in \mathbb{N}.\; \operatorname{act}(\operatorname{gpo\text{-}md\_})(\operatorname{gpo\text{-}p\_}(\operatorname{gpok\_}), \operatorname{op\text{-}iter}(\operatorname{gpo\text{-}md\_}, \operatorname{gpo\text{-}t\_}, \operatorname{gpo\text{-}w\_}, \operatorname{gpok\_}))), \\ \quad (\operatorname{deg}(\operatorname{scal}(\operatorname{gpo\text{-}md\_}), \operatorname{gpo\text{-}p\_}) + 1)) \end{array}
\end{equation*}

\par\medskip\noindent\textbf{Definition} (\texttt{cyclic-subspace}).\;$\operatorname{cyclic\text{-}subspace}(\operatorname{gcs\text{-}md\_}, \operatorname{gcs\text{-}t\_}, \operatorname{gcs\text{-}v\_}) :=$
\begin{equation*}
\{\operatorname{gcsz\_} \in \operatorname{vec}(\operatorname{gcs\text{-}md\_}) : \exists \operatorname{gcsq\_} \in \operatorname{carr}(\operatorname{poly}(\operatorname{scal}(\operatorname{gcs\text{-}md\_}))).\; (\operatorname{gcsz\_} = \operatorname{poly\text{-}op}(\operatorname{gcs\text{-}md\_}, \operatorname{gcs\text{-}t\_}, \operatorname{gcsq\_}, \operatorname{gcs\text{-}v\_}))\}
\end{equation*}

\par\medskip\noindent\textbf{Definition} (\texttt{is-cyclic-vector}).\;For all $\operatorname{gcv\text{-}md\_}$, $\operatorname{gcv\text{-}t\_}$ and $\operatorname{gcv\text{-}v\_}$, 
\begin{equation*}
\begin{array}{@{}l@{}} \operatorname{is\text{-}cyclic\text{-}vector}(\operatorname{gcv\text{-}md\_}, \operatorname{gcv\text{-}t\_}, \operatorname{gcv\text{-}v\_}) \\ \Leftrightarrow\; \operatorname{gcv\text{-}v\_} \in \operatorname{vec}(\operatorname{gcv\text{-}md\_}) \wedge \\ \forall \operatorname{gcvw\_} \in \operatorname{vec}(\operatorname{gcv\text{-}md\_}). \exists \operatorname{gcvq\_} \in \operatorname{carr}(\operatorname{poly}(\operatorname{scal}(\operatorname{gcv\text{-}md\_}))). \\ \quad \operatorname{gcvw\_} = \operatorname{poly\text{-}op}(\operatorname{gcv\text{-}md\_}, \operatorname{gcv\text{-}t\_}, \operatorname{gcvq\_}, \operatorname{gcv\text{-}v\_}) \end{array}
\end{equation*}

\par\medskip\noindent\textbf{Definition} (\texttt{is-monic}).\;For all $\operatorname{gmo\text{-}f\_}$ and $\operatorname{gmo\text{-}p\_}$, 
\begin{equation*}
\begin{array}{@{}l@{}} \operatorname{is\text{-}monic}(\operatorname{gmo\text{-}f\_}, \operatorname{gmo\text{-}p\_}) \\ \Leftrightarrow\; (\operatorname{gmo\text{-}p\_} \in \operatorname{carr}(\operatorname{poly}(\operatorname{gmo\text{-}f\_}))) \wedge (\operatorname{gmo\text{-}p\_}(\operatorname{deg}(\operatorname{gmo\text{-}f\_}, \operatorname{gmo\text{-}p\_})) = \operatorname{one}(\operatorname{gmo\text{-}f\_})) \end{array}
\end{equation*}

\par\medskip\noindent\textbf{Definition} (\texttt{is-annihilating-poly}).\;For all $\operatorname{gan\text{-}md\_}$, $\operatorname{gan\text{-}t\_}$, $\operatorname{gan\text{-}v\_}$ and $\operatorname{gan\text{-}k\_}$, 
\begin{equation*}
\begin{array}{@{}l@{}} \operatorname{is\text{-}annihilating\text{-}poly}(\operatorname{gan\text{-}md\_}, \operatorname{gan\text{-}t\_}, \operatorname{gan\text{-}v\_}, \operatorname{gan\text{-}k\_}) \\ \Leftrightarrow\; \operatorname{is\text{-}monic}(\operatorname{scal}(\operatorname{gan\text{-}md\_}), \operatorname{gan\text{-}k\_}) \wedge \\ \operatorname{poly\text{-}op}(\operatorname{gan\text{-}md\_}, \operatorname{gan\text{-}t\_}, \operatorname{gan\text{-}k\_}, \operatorname{gan\text{-}v\_}) = \operatorname{vzero}(\operatorname{gan\text{-}md\_}) \wedge \\ \forall \operatorname{ganq\_} \in \operatorname{carr}(\operatorname{poly}(\operatorname{scal}(\operatorname{gan\text{-}md\_}))). \\ \quad \neg (\operatorname{ganq\_} = \operatorname{zero}(\operatorname{poly}(\operatorname{scal}(\operatorname{gan\text{-}md\_})))) \Rightarrow  \\ \quad (\operatorname{poly\text{-}op}(\operatorname{gan\text{-}md\_}, \operatorname{gan\text{-}t\_}, \operatorname{ganq\_}, \operatorname{gan\text{-}v\_}) = \operatorname{vzero}(\operatorname{gan\text{-}md\_})) \Rightarrow  \\ \quad \quad \operatorname{deg}(\operatorname{scal}(\operatorname{gan\text{-}md\_}), \operatorname{gan\text{-}k\_}) \leq \operatorname{deg}(\operatorname{scal}(\operatorname{gan\text{-}md\_}), \operatorname{ganq\_}) \end{array}
\end{equation*}

\par\medskip\noindent\textbf{Definition} (\texttt{is-minimal-poly}).\;For all $\operatorname{gmp\text{-}md\_}$, $\operatorname{gmp\text{-}t\_}$ and $\operatorname{gmp\text{-}k\_}$, 
\begin{equation*}
\begin{array}{@{}l@{}} \operatorname{is\text{-}minimal\text{-}poly}(\operatorname{gmp\text{-}md\_}, \operatorname{gmp\text{-}t\_}, \operatorname{gmp\text{-}k\_}) \\ \Leftrightarrow\; \operatorname{is\text{-}monic}(\operatorname{scal}(\operatorname{gmp\text{-}md\_}), \operatorname{gmp\text{-}k\_}) \wedge \\ \forall \operatorname{gmpw\_} \in \operatorname{vec}(\operatorname{gmp\text{-}md\_}). \\ \quad \operatorname{poly\text{-}op}(\operatorname{gmp\text{-}md\_}, \operatorname{gmp\text{-}t\_}, \operatorname{gmp\text{-}k\_}, \operatorname{gmpw\_}) = \operatorname{vzero}(\operatorname{gmp\text{-}md\_}) \wedge \\ \forall \operatorname{gmpq\_} \in \operatorname{carr}(\operatorname{poly}(\operatorname{scal}(\operatorname{gmp\text{-}md\_}))). \\ \quad \neg (\operatorname{gmpq\_} = \operatorname{zero}(\operatorname{poly}(\operatorname{scal}(\operatorname{gmp\text{-}md\_})))) \Rightarrow  \\ \quad \forall \operatorname{gmpw\_} \in \operatorname{vec}(\operatorname{gmp\text{-}md\_}).\; (\operatorname{poly\text{-}op}(\operatorname{gmp\text{-}md\_}, \operatorname{gmp\text{-}t\_}, \operatorname{gmpq\_}, \operatorname{gmpw\_}) = \operatorname{vzero}(\operatorname{gmp\text{-}md\_})) \Rightarrow  \\ \quad \quad \operatorname{deg}(\operatorname{scal}(\operatorname{gmp\text{-}md\_}), \operatorname{gmp\text{-}k\_}) \leq \operatorname{deg}(\operatorname{scal}(\operatorname{gmp\text{-}md\_}), \operatorname{gmpq\_}) \end{array}
\end{equation*}

\par\medskip\noindent\textbf{Definition} (\texttt{col-as-row}).\;$\operatorname{col\text{-}as\text{-}row}(\operatorname{gcr\text{-}a\_}, \operatorname{gcr\text{-}n\_}, \operatorname{gcr\text{-}j\_}) :=$
\begin{equation*}
\begin{array}{@{}l@{}} \operatorname{matof}(1, \operatorname{gcr\text{-}n\_}, \\ \quad (\lambda {[}\operatorname{gcri\_}, \operatorname{gcrk\_}] \in (\operatorname{interval}(1, 1) \times \operatorname{interval}(1, \operatorname{gcr\text{-}n\_})).\; \\ \quad \quad \operatorname{entry}(\operatorname{gcr\text{-}a\_}, \operatorname{gcrk\_}, \operatorname{gcr\text{-}j\_}))) \end{array}
\end{equation*}

\par\medskip\noindent\textbf{Definition} (\texttt{is-matrix-of}).\;For all $\operatorname{gmx\text{-}md\_}$, $\operatorname{gmx\text{-}n\_}$, $\operatorname{gmx\text{-}u\_}$, $\operatorname{gmx\text{-}t\_}$ and $\operatorname{gmx\text{-}a\_}$, 
\begin{equation*}
\begin{array}{@{}l@{}} \operatorname{is\text{-}matrix\text{-}of}(\operatorname{gmx\text{-}md\_}, \operatorname{gmx\text{-}n\_}, \operatorname{gmx\text{-}u\_}, \operatorname{gmx\text{-}t\_}, \operatorname{gmx\text{-}a\_}) \\ \Leftrightarrow\; \operatorname{gmx\text{-}a\_} \in \operatorname{mat}(\operatorname{gmx\text{-}n\_}, \operatorname{gmx\text{-}n\_}, \operatorname{carr}(\operatorname{scal}(\operatorname{gmx\text{-}md\_}))) \wedge \\ \forall \operatorname{gmxj\_} \in \operatorname{interval}(1, \operatorname{gmx\text{-}n\_}). \\ \quad \operatorname{gmx\text{-}t\_}(\operatorname{entry}(\operatorname{gmx\text{-}u\_}, \operatorname{gmxj\_}, 1)) \\ \quad =\; \operatorname{lincomb}(\operatorname{gmx\text{-}md\_}, \operatorname{gmx\text{-}n\_}, \operatorname{col\text{-}as\text{-}row}(\operatorname{gmx\text{-}a\_}, \operatorname{gmx\text{-}n\_}, \operatorname{gmxj\_}), \operatorname{gmx\text{-}u\_}) \end{array}
\end{equation*}

\par\medskip\noindent\textbf{Definition} (\texttt{lin-span}).\;$\operatorname{lin\text{-}span}(\operatorname{gls\text{-}md\_}, \operatorname{gls\text{-}n\_}, \operatorname{gls\text{-}u\_}) :=$
\begin{equation*}
\{\operatorname{glsx\_} \in \operatorname{vec}(\operatorname{gls\text{-}md\_}) : \exists \operatorname{glsc\_} \in \operatorname{mat}(1, \operatorname{gls\text{-}n\_}, \operatorname{carr}(\operatorname{scal}(\operatorname{gls\text{-}md\_}))).\; (\operatorname{glsx\_} = \operatorname{lincomb}(\operatorname{gls\text{-}md\_}, \operatorname{gls\text{-}n\_}, \operatorname{glsc\_}, \operatorname{gls\text{-}u\_}))\}
\end{equation*}

\par\medskip\noindent\textbf{Definition} (\texttt{is-nilpotent-op}).\;For all $\operatorname{gnp\text{-}md\_}$ and $\operatorname{gnp\text{-}t\_}$, 
\begin{equation*}
\begin{array}{@{}l@{}} \operatorname{is\text{-}nilpotent\text{-}op}(\operatorname{gnp\text{-}md\_}, \operatorname{gnp\text{-}t\_}) \\ \Leftrightarrow\; \exists \operatorname{gnpk\_} \in \mathbb{N}. \forall \operatorname{gnpw\_} \in \operatorname{vec}(\operatorname{gnp\text{-}md\_}). \\ \quad \operatorname{op\text{-}iter}(\operatorname{gnp\text{-}md\_}, \operatorname{gnp\text{-}t\_}, \operatorname{gnpw\_}, \operatorname{gnpk\_}) = \operatorname{vzero}(\operatorname{gnp\text{-}md\_}) \end{array}
\end{equation*}


\vnbbookfile{theorem-library/gal-op}{theorem-library/gal-op}
\noindent{\small gal-op.scm -- a linear operator as an F[x]-module; cyclic vectors, the Krylov basis, the companion matrix and the minimum polynomial (galois-theory notes, section 7, Proposition 1.36) (GALOIS B6, 2026-10-26).  Destination: theorem-library/gal-op.scm, library linear-algebra. Helper prefix go-.}\par
\vnbprop{op-iter-type}{\vnbstmt{\begin{array}{@{}l@{}} \forall k \in \mathbb{N}, d, t.; \forall w \in \operatorname{vec}(d). \\ \quad \operatorname{is\text{-}hom\text{-}module}(d, d, t) \Rightarrow (\operatorname{op\text{-}iter}(d, t, w, k) \in \operatorname{vec}(d)) \end{array}}}
\vnbprop{op-iter-add}{\vnbstmt{\begin{array}{@{}l@{}} \forall k \in \mathbb{N}, d, t.; \forall a, b \in \operatorname{vec}(d). \\ \quad \operatorname{is\text{-}hom\text{-}module}(d, d, t) \Rightarrow  \\ \quad \quad \operatorname{op\text{-}iter}(d, t, \operatorname{vadd}(d)(a, b), k) = \operatorname{vadd}(d)(\operatorname{op\text{-}iter}(d, t, a, k), \operatorname{op\text{-}iter}(d, t, b, k)) \end{array}}}
\vnbprop{op-iter-act}{\vnbstmt{\begin{array}{@{}l@{}} \forall k \in \mathbb{N}, d, t.; \forall c \in \operatorname{carr}(\operatorname{scal}(d)), a \in \operatorname{vec}(d). \\ \quad \operatorname{is\text{-}hom\text{-}module}(d, d, t) \Rightarrow  \\ \quad \quad \operatorname{op\text{-}iter}(d, t, \operatorname{act}(d)(c, a), k) = \operatorname{act}(d)(c, \operatorname{op\text{-}iter}(d, t, a, k)) \end{array}}}
\vnbprop{op-iter-shift}{\vnbstmt{\begin{array}{@{}l@{}} \forall k \in \mathbb{N}, d, t.; \forall w \in \operatorname{vec}(d). \\ \quad \operatorname{is\text{-}hom\text{-}module}(d, d, t) \Rightarrow (\operatorname{op\text{-}iter}(d, t, t(w), k) = t(\operatorname{op\text{-}iter}(d, t, w, k))) \end{array}}}
\vnbprop{op-iter-vzero}{\vnbstmt{\begin{array}{@{}l@{}} \forall k \in \mathbb{N}, d, t. \\ \quad \operatorname{is\text{-}hom\text{-}module}(d, d, t) \Rightarrow (\operatorname{op\text{-}iter}(d, t, \operatorname{vzero}(d), k) = \operatorname{vzero}(d)) \end{array}}}
\vnbprop{finsum-interval-as-sum-ag}{\vnbstmt{\begin{array}{@{}l@{}} \forall n \in \mathbb{N}, g.; \forall h \in (\operatorname{interval}(1, n) \to \operatorname{carr}(g)). \\ \quad \operatorname{is\text{-}abelian\text{-}group}(g) \Rightarrow  \\ \quad \quad \operatorname{finsum}(g, h, \operatorname{interval}(1, n)) = \operatorname{sum\text{-}ag}(g, (\lambda k \in \mathbb{N}.\; h((k + 1))), n) \end{array}}}
\vnbprop{krylov-entry}{\vnbstmt{\begin{array}{@{}l@{}} \forall d, t.; \forall n \in \mathbb{N}, v \in \operatorname{vec}(d), i \in \operatorname{interval}(1, n). \\ \quad \operatorname{is\text{-}hom\text{-}module}(d, d, t) \Rightarrow  \\ \quad \quad \operatorname{entry}(\operatorname{krylov}(d, t, v, n), i, 1) = \operatorname{op\text{-}iter}(d, t, v, (i - 1)) \end{array}}}
\vnbprop{krylov-in-mat}{\vnbstmt{\begin{array}{@{}l@{}} \forall d, t.; \forall n \in \mathbb{N}, v \in \operatorname{vec}(d). \\ \quad \operatorname{is\text{-}hom\text{-}module}(d, d, t) \Rightarrow (\operatorname{krylov}(d, t, v, n) \in \operatorname{mat}(n, 1, \operatorname{vec}(d))) \end{array}}}
\vnbprop{lincomb-krylov}{\vnbstmt{\begin{array}{@{}l@{}} \forall d, t.; \forall n \in \mathbb{N}, v \in \operatorname{vec}(d), c \in \operatorname{mat}(1, n, \operatorname{carr}(\operatorname{scal}(d))). \\ \quad \operatorname{is\text{-}hom\text{-}module}(d, d, t) \Rightarrow  \\ \quad \quad \operatorname{lincomb}(d, n, c, \operatorname{krylov}(d, t, v, n)) \\ \quad \quad =\; \operatorname{sum\text{-}ag}(\operatorname{module\text{-}vector\text{-}ag}(d), \\ \quad \quad \quad (\lambda k \in \mathbb{N}.\; \operatorname{act}(d)(\operatorname{entry}(c, 1, (k + 1)), \operatorname{op\text{-}iter}(d, t, v, k))), \\ \quad \quad \quad n) \end{array}}}
\vnbprop{poly-coef-in}{\vnbstmt{\forall a.; \forall p \in \operatorname{carr}(\operatorname{poly}(a)), k \in \mathbb{N}.\; (\operatorname{is\text{-}ring}(a) \Rightarrow (p(k) \in \operatorname{carr}(a)))}}
\vnbprop{poly-op-sum-stable}{\vnbstmt{\begin{array}{@{}l@{}} \forall j \in \mathbb{N}, d, t.; \forall p \in \operatorname{carr}(\operatorname{poly}(\operatorname{scal}(d))), w \in \operatorname{vec}(d). \\ \quad \operatorname{is\text{-}hom\text{-}module}(d, d, t) \Rightarrow  \\ \quad \quad \operatorname{sum\text{-}ag}(\operatorname{module\text{-}vector\text{-}ag}(d), \\ \quad \quad \quad (\lambda k \in \mathbb{N}.\; \operatorname{act}(d)(p(k), \operatorname{op\text{-}iter}(d, t, w, k))), \\ \quad \quad \quad ((\operatorname{deg}(\operatorname{scal}(d), p) + j) + 1)) \\ \quad \quad =\; \operatorname{sum\text{-}ag}(\operatorname{module\text{-}vector\text{-}ag}(d), \\ \quad \quad \quad (\lambda k \in \mathbb{N}.\; \operatorname{act}(d)(p(k), \operatorname{op\text{-}iter}(d, t, w, k))), \\ \quad \quad \quad (\operatorname{deg}(\operatorname{scal}(d), p) + 1)) \end{array}}}
\vnbprop{poly-op-bound}{\vnbstmt{\begin{array}{@{}l@{}} \forall d, t.; \forall p \in \operatorname{carr}(\operatorname{poly}(\operatorname{scal}(d))), w \in \operatorname{vec}(d), n \in \mathbb{N}. \\ \quad (\operatorname{is\text{-}hom\text{-}module}(d, d, t) \wedge \operatorname{deg\text{-}bound}(\operatorname{scal}(d), p, n)) \Rightarrow  \\ \quad \quad \operatorname{poly\text{-}op}(d, t, p, w) \\ \quad \quad =\; \operatorname{sum\text{-}ag}(\operatorname{module\text{-}vector\text{-}ag}(d), (\lambda k \in \mathbb{N}.\; \operatorname{act}(d)(p(k), \operatorname{op\text{-}iter}(d, t, w, k))), (n + 1)) \end{array}}}
\vnbprop{poly-op-type}{\vnbstmt{\begin{array}{@{}l@{}} \forall d, t.; \forall p \in \operatorname{carr}(\operatorname{poly}(\operatorname{scal}(d))), w \in \operatorname{vec}(d). \\ \quad \operatorname{is\text{-}hom\text{-}module}(d, d, t) \Rightarrow (\operatorname{poly\text{-}op}(d, t, p, w) \in \operatorname{vec}(d)) \end{array}}}
\vnbprop{poly-op-vzero}{\vnbstmt{\begin{array}{@{}l@{}} \forall d, t.; \forall p \in \operatorname{carr}(\operatorname{poly}(\operatorname{scal}(d))). \\ \quad \operatorname{is\text{-}hom\text{-}module}(d, d, t) \Rightarrow (\operatorname{poly\text{-}op}(d, t, p, \operatorname{vzero}(d)) = \operatorname{vzero}(d)) \end{array}}}
\vnbprop{poly-op-add}{\vnbstmt{\begin{array}{@{}l@{}} \forall d, t.; \forall p \in \operatorname{carr}(\operatorname{poly}(\operatorname{scal}(d))), x, y \in \operatorname{vec}(d). \\ \quad \operatorname{is\text{-}hom\text{-}module}(d, d, t) \Rightarrow  \\ \quad \quad \operatorname{poly\text{-}op}(d, t, p, \operatorname{vadd}(d)(x, y)) = \operatorname{vadd}(d)(\operatorname{poly\text{-}op}(d, t, p, x), \operatorname{poly\text{-}op}(d, t, p, y)) \end{array}}}
\vnbprop{poly-op-act}{\vnbstmt{\begin{array}{@{}l@{}} \forall d, t.; \forall p \in \operatorname{carr}(\operatorname{poly}(\operatorname{scal}(d))), c \in \operatorname{carr}(\operatorname{scal}(d)), x \in \operatorname{vec}(d). \\ \quad (\operatorname{is\text{-}hom\text{-}module}(d, d, t) \wedge \operatorname{is\text{-}commutative\text{-}ring}(\operatorname{scal}(d))) \Rightarrow  \\ \quad \quad \operatorname{poly\text{-}op}(d, t, p, \operatorname{act}(d)(c, x)) = \operatorname{act}(d)(c, \operatorname{poly\text{-}op}(d, t, p, x)) \end{array}}}
\vnbprop{hom-sum-ag}{\vnbstmt{\begin{array}{@{}l@{}} \forall n \in \mathbb{N}, d, t, g. \\ \quad (\operatorname{is\text{-}hom\text{-}module}(d, d, t) \wedge \forall i \in \mathrm{OS}(n).\; (g(i) \in \operatorname{vec}(d))) \Rightarrow  \\ \quad \quad t(\operatorname{sum\text{-}ag}(\operatorname{module\text{-}vector\text{-}ag}(d), g, n)) \\ \quad \quad =\; \operatorname{sum\text{-}ag}(\operatorname{module\text{-}vector\text{-}ag}(d), (\lambda \mathit{gti} \in \mathbb{N}.\; t(g(\mathit{gti}))), n) \end{array}}}
\vnbprop{poly-op-commutes}{\vnbstmt{\begin{array}{@{}l@{}} \forall d, t.; \forall p \in \operatorname{carr}(\operatorname{poly}(\operatorname{scal}(d))), w \in \operatorname{vec}(d). \\ \quad \operatorname{is\text{-}hom\text{-}module}(d, d, t) \Rightarrow (\operatorname{poly\text{-}op}(d, t, p, t(w)) = t(\operatorname{poly\text{-}op}(d, t, p, w))) \end{array}}}
\vnbprop{poly-op-iter}{\vnbstmt{\begin{array}{@{}l@{}} \forall j \in \mathbb{N}, d, t.; \forall p \in \operatorname{carr}(\operatorname{poly}(\operatorname{scal}(d))), w \in \operatorname{vec}(d). \\ \quad \operatorname{is\text{-}hom\text{-}module}(d, d, t) \Rightarrow  \\ \quad \quad \operatorname{poly\text{-}op}(d, t, p, \operatorname{op\text{-}iter}(d, t, w, j)) = \operatorname{op\text{-}iter}(d, t, \operatorname{poly\text{-}op}(d, t, p, w), j) \end{array}}}
\vnbprop{poly-op-sum-kernel}{\vnbstmt{\begin{array}{@{}l@{}} \forall n \in \mathbb{N}, d, t.; \forall p \in \operatorname{carr}(\operatorname{poly}(\operatorname{scal}(d))), g. \\ \quad \operatorname{is\text{-}hom\text{-}module}(d, d, t) \wedge \\ \quad \forall i \in \mathrm{OS}(n).\; (g(i) \in \operatorname{vec}(d)) \wedge \\ \quad \forall i \in \mathrm{OS}(n).\; (\operatorname{poly\text{-}op}(d, t, p, g(i)) = \operatorname{vzero}(d)) \Rightarrow \\ \quad \quad \operatorname{poly\text{-}op}(d, t, p, \operatorname{sum\text{-}ag}(\operatorname{module\text{-}vector\text{-}ag}(d), g, n)) = \operatorname{vzero}(d) \end{array}}}
\vnbprop{poly-op-kills-krylov-span}{\vnbstmt{\begin{array}{@{}l@{}} \forall d, t.; \forall p \in \operatorname{carr}(\operatorname{poly}(\operatorname{scal}(d))), v \in \operatorname{vec}(d), n \in \mathbb{N}, c \in \operatorname{mat}(1, n, \operatorname{carr}(\operatorname{scal}(d))). \\ \quad (\operatorname{is\text{-}hom\text{-}module}(d, d, t) \wedge (\operatorname{is\text{-}commutative\text{-}ring}(\operatorname{scal}(d)) \wedge (\operatorname{poly\text{-}op}(d, t, p, v) = \operatorname{vzero}(d)))) \Rightarrow  \\ \quad \quad \operatorname{poly\text{-}op}(d, t, p, \operatorname{lincomb}(d, n, c, \operatorname{krylov}(d, t, v, n))) = \operatorname{vzero}(d) \end{array}}}
\vnbprop{krylov-rel-free}{\vnbstmt{\begin{array}{@{}l@{}} \forall d, t.; \forall v \in \operatorname{vec}(d), k. \\ \quad (\operatorname{is\text{-}hom\text{-}module}(d, d, t) \wedge \operatorname{is\text{-}annihilating\text{-}poly}(d, t, v, k)) \Rightarrow  \\ \quad \quad \operatorname{rel\text{-}free}(d, \operatorname{deg}(\operatorname{scal}(d), k), \operatorname{krylov}(d, t, v, \operatorname{deg}(\operatorname{scal}(d), k))) \end{array}}}
\vnbprop{lin-span-unfold}{\vnbstmt{\begin{array}{@{}l@{}} \forall d.; \forall n \in \mathbb{N}, u \in \operatorname{mat}(n, 1, \operatorname{vec}(d)). \\ \quad \operatorname{is\text{-}module}(d) \Rightarrow  \\ \quad \quad (\operatorname{lin\text{-}span}(d, n, u) \simeq \{x \in \operatorname{vec}(d) : \exists c \in \operatorname{mat}(1, n, \operatorname{carr}(\operatorname{scal}(d))).\; (x = \operatorname{lincomb}(d, n, c, u))\}) \end{array}}}
\vnbprop{lin-span-intro}{\vnbstmt{\begin{array}{@{}l@{}} \forall d.; \forall n \in \mathbb{N}, u \in \operatorname{mat}(n, 1, \operatorname{vec}(d)), c \in \operatorname{mat}(1, n, \operatorname{carr}(\operatorname{scal}(d))). \\ \quad \operatorname{is\text{-}module}(d) \Rightarrow (\operatorname{lincomb}(d, n, c, u) \in \operatorname{lin\text{-}span}(d, n, u)) \end{array}}}
\vnbprop{lin-span-elim}{\vnbstmt{\begin{array}{@{}l@{}} \forall d.; \forall n \in \mathbb{N}, u \in \operatorname{mat}(n, 1, \operatorname{vec}(d)), x \in \operatorname{lin\text{-}span}(d, n, u). \\ \quad \operatorname{is\text{-}module}(d) \Rightarrow  \\ \quad \quad ((x \in \operatorname{vec}(d)) \wedge \exists c \in \operatorname{mat}(1, n, \operatorname{carr}(\operatorname{scal}(d))).\; (x = \operatorname{lincomb}(d, n, c, u))) \end{array}}}
\vnbprop{lin-span-zero-add}{\vnbstmt{\begin{array}{@{}l@{}} \forall d.; \forall n \in \mathbb{N}, u \in \operatorname{mat}(n, 1, \operatorname{vec}(d)). \\ \quad \operatorname{is\text{-}module}(d) \Rightarrow  \\ \quad \quad \operatorname{vzero}(d) \in \operatorname{lin\text{-}span}(d, n, u) \wedge \\ \quad \quad \forall x, y \in \operatorname{lin\text{-}span}(d, n, u). \\ \quad \quad \quad \operatorname{vadd}(d)(x, y) \in \operatorname{lin\text{-}span}(d, n, u) \end{array}}}
\vnbprop{lin-span-act}{\vnbstmt{\begin{array}{@{}l@{}} \forall d.; \forall n \in \mathbb{N}, u \in \operatorname{mat}(n, 1, \operatorname{vec}(d)), r \in \operatorname{carr}(\operatorname{scal}(d)), x \in \operatorname{lin\text{-}span}(d, n, u). \\ \quad \operatorname{is\text{-}module}(d) \Rightarrow (\operatorname{act}(d)(r, x) \in \operatorname{lin\text{-}span}(d, n, u)) \end{array}}}
\vnbprop{lin-span-member}{\vnbstmt{\begin{array}{@{}l@{}} \forall d.; \forall n \in \mathbb{N}, u \in \operatorname{mat}(n, 1, \operatorname{vec}(d)), i \in \operatorname{interval}(1, n). \\ \quad \operatorname{is\text{-}module}(d) \Rightarrow (\operatorname{entry}(u, i, 1) \in \operatorname{lin\text{-}span}(d, n, u)) \end{array}}}
\vnbprop{lin-span-sum}{\vnbstmt{\begin{array}{@{}l@{}} \forall d.; \forall n \in \mathbb{N}, u \in \operatorname{mat}(n, 1, \operatorname{vec}(d)), m \in \mathbb{N}, g. \\ \quad (\operatorname{is\text{-}module}(d) \wedge \forall i \in \mathrm{OS}(m).\; (g(i) \in \operatorname{lin\text{-}span}(d, n, u))) \Rightarrow  \\ \quad \quad \operatorname{sum\text{-}ag}(\operatorname{module\text{-}vector\text{-}ag}(d), g, m) \in \operatorname{lin\text{-}span}(d, n, u) \end{array}}}
\vnbprop{krylov-top}{\vnbstmt{\begin{array}{@{}l@{}} \forall d, t.; \forall v \in \operatorname{vec}(d), k \in \operatorname{carr}(\operatorname{poly}(\operatorname{scal}(d))). \\ \quad \operatorname{is\text{-}hom\text{-}module}(d, d, t) \wedge \\ \quad \operatorname{is\text{-}commutative\text{-}ring}(\operatorname{scal}(d)) \wedge \\ \quad k(\operatorname{deg}(\operatorname{scal}(d), k)) = \operatorname{one}(\operatorname{scal}(d)) \wedge \\ \quad \operatorname{poly\text{-}op}(d, t, k, v) = \operatorname{vzero}(d) \Rightarrow \\ \quad \quad \operatorname{op\text{-}iter}(d, t, v, \operatorname{deg}(\operatorname{scal}(d), k)) \\ \quad \quad =\; \operatorname{sum\text{-}ag}(\operatorname{module\text{-}vector\text{-}ag}(d), \\ \quad \quad \quad (\lambda \mathit{gpk} \in \mathbb{N}.\; \operatorname{act}(d)(\operatorname{neg}(\operatorname{scal}(d))(k(\mathit{gpk})), \operatorname{op\text{-}iter}(d, t, v, \mathit{gpk}))), \\ \quad \quad \quad \operatorname{deg}(\operatorname{scal}(d), k)) \end{array}}}
\vnbprop{krylov-low-in-span}{\vnbstmt{\begin{array}{@{}l@{}} \forall d, t.; \forall v \in \operatorname{vec}(d), k \in \operatorname{carr}(\operatorname{poly}(\operatorname{scal}(d))), i \in \mathbb{N}. \\ \quad \operatorname{is\text{-}hom\text{-}module}(d, d, t) \wedge \\ \quad \operatorname{is\text{-}commutative\text{-}ring}(\operatorname{scal}(d)) \wedge \\ \quad k(\operatorname{deg}(\operatorname{scal}(d), k)) = \operatorname{one}(\operatorname{scal}(d)) \wedge \\ \quad \operatorname{poly\text{-}op}(d, t, k, v) = \operatorname{vzero}(d) \wedge \\ \quad 1 \leq \operatorname{deg}(\operatorname{scal}(d), k) \wedge \\ \quad i < \operatorname{deg}(\operatorname{scal}(d), k) \Rightarrow \\ \quad \quad \operatorname{op\text{-}iter}(d, t, v, i) \in \operatorname{lin\text{-}span}(d, \operatorname{deg}(\operatorname{scal}(d), k), \operatorname{krylov}(d, t, v, \operatorname{deg}(\operatorname{scal}(d), k))) \end{array}}}
\vnbprop{krylov-top-in-span}{\vnbstmt{\begin{array}{@{}l@{}} \forall d, t.; \forall v \in \operatorname{vec}(d), k \in \operatorname{carr}(\operatorname{poly}(\operatorname{scal}(d))). \\ \quad \operatorname{is\text{-}hom\text{-}module}(d, d, t) \wedge \\ \quad \operatorname{is\text{-}commutative\text{-}ring}(\operatorname{scal}(d)) \wedge \\ \quad k(\operatorname{deg}(\operatorname{scal}(d), k)) = \operatorname{one}(\operatorname{scal}(d)) \wedge \\ \quad \operatorname{poly\text{-}op}(d, t, k, v) = \operatorname{vzero}(d) \wedge \\ \quad 1 \leq \operatorname{deg}(\operatorname{scal}(d), k) \Rightarrow \\ \quad \quad \operatorname{op\text{-}iter}(d, t, v, \operatorname{deg}(\operatorname{scal}(d), k)) \in \operatorname{lin\text{-}span}(d, \operatorname{deg}(\operatorname{scal}(d), k), \operatorname{krylov}(d, t, v, \operatorname{deg}(\operatorname{scal}(d), k))) \end{array}}}
\vnbprop{krylov-span-invariant}{\vnbstmt{\begin{array}{@{}l@{}} \forall d, t.; \forall v \in \operatorname{vec}(d), k \in \operatorname{carr}(\operatorname{poly}(\operatorname{scal}(d))), x \in \operatorname{lin\text{-}span}(d, \operatorname{deg}(\operatorname{scal}(d), k), \operatorname{krylov}(d, t, v, \operatorname{deg}(\operatorname{scal}(d), k))). \\ \quad \operatorname{is\text{-}hom\text{-}module}(d, d, t) \wedge \\ \quad \operatorname{is\text{-}commutative\text{-}ring}(\operatorname{scal}(d)) \wedge \\ \quad k(\operatorname{deg}(\operatorname{scal}(d), k)) = \operatorname{one}(\operatorname{scal}(d)) \wedge \\ \quad \operatorname{poly\text{-}op}(d, t, k, v) = \operatorname{vzero}(d) \wedge \\ \quad 1 \leq \operatorname{deg}(\operatorname{scal}(d), k) \Rightarrow \\ \quad \quad t(x) \in \operatorname{lin\text{-}span}(d, \operatorname{deg}(\operatorname{scal}(d), k), \operatorname{krylov}(d, t, v, \operatorname{deg}(\operatorname{scal}(d), k))) \end{array}}}
\vnbprop{krylov-iterates-in-span}{\vnbstmt{\begin{array}{@{}l@{}} \forall j \in \mathbb{N}, d, t.; \forall v \in \operatorname{vec}(d), k \in \operatorname{carr}(\operatorname{poly}(\operatorname{scal}(d))). \\ \quad \operatorname{is\text{-}hom\text{-}module}(d, d, t) \wedge \\ \quad \operatorname{is\text{-}commutative\text{-}ring}(\operatorname{scal}(d)) \wedge \\ \quad k(\operatorname{deg}(\operatorname{scal}(d), k)) = \operatorname{one}(\operatorname{scal}(d)) \wedge \\ \quad \operatorname{poly\text{-}op}(d, t, k, v) = \operatorname{vzero}(d) \wedge \\ \quad 1 \leq \operatorname{deg}(\operatorname{scal}(d), k) \Rightarrow \\ \quad \quad \operatorname{op\text{-}iter}(d, t, v, j) \in \operatorname{lin\text{-}span}(d, \operatorname{deg}(\operatorname{scal}(d), k), \operatorname{krylov}(d, t, v, \operatorname{deg}(\operatorname{scal}(d), k))) \end{array}}}
\vnbprop{krylov-generates}{\vnbstmt{\begin{array}{@{}l@{}} \forall d, t.; \forall v \in \operatorname{vec}(d), k \in \operatorname{carr}(\operatorname{poly}(\operatorname{scal}(d))). \\ \quad \operatorname{is\text{-}hom\text{-}module}(d, d, t) \wedge \\ \quad \operatorname{is\text{-}commutative\text{-}ring}(\operatorname{scal}(d)) \wedge \\ \quad k(\operatorname{deg}(\operatorname{scal}(d), k)) = \operatorname{one}(\operatorname{scal}(d)) \wedge \\ \quad \operatorname{poly\text{-}op}(d, t, k, v) = \operatorname{vzero}(d) \wedge \\ \quad 1 \leq \operatorname{deg}(\operatorname{scal}(d), k) \wedge \\ \quad \operatorname{is\text{-}cyclic\text{-}vector}(d, t, v) \Rightarrow \\ \quad \quad \operatorname{generates}(d, \operatorname{deg}(\operatorname{scal}(d), k), \operatorname{krylov}(d, t, v, \operatorname{deg}(\operatorname{scal}(d), k))) \end{array}}}
\vnbprop{annihilating-is-minimal}{\vnbstmt{\begin{array}{@{}l@{}} \forall d, t, v, k. \\ \quad \operatorname{is\text{-}hom\text{-}module}(d, d, t) \wedge \\ \quad \operatorname{is\text{-}commutative\text{-}ring}(\operatorname{scal}(d)) \wedge \\ \quad \operatorname{is\text{-}cyclic\text{-}vector}(d, t, v) \wedge \\ \quad \operatorname{is\text{-}annihilating\text{-}poly}(d, t, v, k) \wedge \\ \quad 1 \leq \operatorname{deg}(\operatorname{scal}(d), k) \Rightarrow \\ \quad \quad \operatorname{is\text{-}minimal\text{-}poly}(d, t, k) \end{array}}}
\vnbprop{col-as-row-type}{\vnbstmt{\begin{array}{@{}l@{}} \forall \mathit{gcra}.; \forall n \in \mathbb{N}, x.; \forall a \in \operatorname{mat}(n, n, x), j \in \operatorname{interval}(1, n). \\ \quad \operatorname{col\text{-}as\text{-}row}(a, n, j) \in \operatorname{mat}(1, n, x) \end{array}}}
\vnbprop{col-as-row-entry}{\vnbstmt{\begin{array}{@{}l@{}} \forall \mathit{gcra}.; \forall n \in \mathbb{N}, x.; \forall a \in \operatorname{mat}(n, n, x), j, k \in \operatorname{interval}(1, n). \\ \quad \operatorname{entry}(\operatorname{col\text{-}as\text{-}row}(a, n, j), 1, k) = \operatorname{entry}(a, k, j) \end{array}}}
\vnbprop{monic-poly-top}{\vnbstmt{\begin{array}{@{}l@{}} \forall r.; \forall v \in \mathbb{N}, \mathit{av}.; \forall x \in \mathbb{N}. \\ \quad (\operatorname{is\text{-}ring}(r) \wedge (\forall i \in \operatorname{interval}(1, v).\; (\operatorname{av}(i) \in \operatorname{carr}(r)) \wedge (x = v))) \Rightarrow  \\ \quad \quad \operatorname{monic\text{-}poly}(r, \mathit{av}, v)(x) = \operatorname{one}(r) \end{array}}}
\vnbprop{monic-poly-low}{\vnbstmt{\begin{array}{@{}l@{}} \forall r.; \forall v \in \mathbb{N}, \mathit{av}.; \forall x \in \mathbb{N}. \\ \quad (\operatorname{is\text{-}ring}(r) \wedge (\forall i \in \operatorname{interval}(1, v).\; (\operatorname{av}(i) \in \operatorname{carr}(r)) \wedge (x < v))) \Rightarrow  \\ \quad \quad \operatorname{monic\text{-}poly}(r, \mathit{av}, v)(x) = \operatorname{av}((x + 1)) \end{array}}}
\vnbprop{monic-poly-high}{\vnbstmt{\begin{array}{@{}l@{}} \forall r.; \forall v \in \mathbb{N}, \mathit{av}.; \forall x \in \mathbb{N}. \\ \quad (\operatorname{is\text{-}ring}(r) \wedge (\forall i \in \operatorname{interval}(1, v).\; (\operatorname{av}(i) \in \operatorname{carr}(r)) \wedge (v < x))) \Rightarrow  \\ \quad \quad \operatorname{monic\text{-}poly}(r, \mathit{av}, v)(x) = \operatorname{zero}(r) \end{array}}}
\vnbprop{monic-poly-deg}{\vnbstmt{\begin{array}{@{}l@{}} \forall r.; \forall v \in \mathbb{N}, \mathit{av}. \\ \quad (\operatorname{is\text{-}ring}(r) \wedge (\neg (\operatorname{one}(r) = \operatorname{zero}(r)) \wedge \forall i \in \operatorname{interval}(1, v).\; (\operatorname{av}(i) \in \operatorname{carr}(r)))) \Rightarrow  \\ \quad \quad (\operatorname{deg}(r, \operatorname{monic\text{-}poly}(r, \mathit{av}, v)) = v) \wedge \operatorname{is\text{-}monic}(r, \operatorname{monic\text{-}poly}(r, \mathit{av}, v)) \end{array}}}
\vnbprop{krylov-companion}{\vnbstmt{\begin{array}{@{}l@{}} \forall d, t.; \forall \mathit{nv} \in \mathbb{N}, \mathit{av}.; \forall v \in \operatorname{vec}(d). \\ \quad \operatorname{is\text{-}hom\text{-}module}(d, d, t) \wedge \\ \quad \operatorname{is\text{-}commutative\text{-}ring}(\operatorname{scal}(d)) \wedge \\ \quad \neg (\operatorname{one}(\operatorname{scal}(d)) = \operatorname{zero}(\operatorname{scal}(d))) \wedge \\ \quad 1 \leq \mathit{nv} \wedge \\ \quad \forall i \in \operatorname{interval}(1, \mathit{nv}).\; (\operatorname{av}(i) \in \operatorname{carr}(\operatorname{scal}(d))) \wedge \\ \quad \operatorname{poly\text{-}op}(d, t, \operatorname{monic\text{-}poly}(\operatorname{scal}(d), \mathit{av}, \mathit{nv}), v) = \operatorname{vzero}(d) \Rightarrow \\ \quad \quad \operatorname{is\text{-}matrix\text{-}of}(d, \mathit{nv}, \operatorname{krylov}(d, t, v, \mathit{nv}), t, \operatorname{companion}(\operatorname{scal}(d), \mathit{av}, \mathit{nv})) \end{array}}}
\vnbprop{cyclic-vector-companion}{\vnbstmt{\begin{array}{@{}l@{}} \forall d, t, v.; \forall \mathit{nv} \in \mathbb{N}, \mathit{av}. \\ \quad \operatorname{is\text{-}vector\text{-}space}(d) \wedge \\ \quad \operatorname{is\text{-}hom\text{-}module}(d, d, t) \wedge \\ \quad \operatorname{is\text{-}cyclic\text{-}vector}(d, t, v) \wedge \\ \quad 1 \leq \mathit{nv} \wedge \\ \quad \forall i \in \operatorname{interval}(1, \mathit{nv}).\; (\operatorname{av}(i) \in \operatorname{carr}(\operatorname{scal}(d))) \wedge \\ \quad \operatorname{is\text{-}annihilating\text{-}poly}(d, t, v, \operatorname{monic\text{-}poly}(\operatorname{scal}(d), \mathit{av}, \mathit{nv})) \Rightarrow \\ \quad \quad \operatorname{generates}(d, \mathit{nv}, \operatorname{krylov}(d, t, v, \mathit{nv})) \wedge \\ \quad \quad \operatorname{rel\text{-}free}(d, \mathit{nv}, \operatorname{krylov}(d, t, v, \mathit{nv})) \wedge \\ \quad \quad \operatorname{is\text{-}matrix\text{-}of}(d, \mathit{nv}, \operatorname{krylov}(d, t, v, \mathit{nv}), t, \operatorname{companion}(\operatorname{scal}(d), \mathit{av}, \mathit{nv})) \wedge \\ \quad \quad \operatorname{is\text{-}minimal\text{-}poly}(d, t, \operatorname{monic\text{-}poly}(\operatorname{scal}(d), \mathit{av}, \mathit{nv})) \end{array}}}

\vnbbookfile{structure-library/gal-ch-vocab}{structure-library/gal-ch-vocab}
\noindent{\small gal-ch-vocab.scm -- vocabulary for the general Cayley-Hamilton theorem (GALOIS B6, 2026-10-26). DEFINITIONS ONLY.  Destination: structure-library/gal-ch-vocab.scm, library linear-algebra. Load window: after gal-op-vocab.scm and the matrix files (MATOF, ENTRY, MAT).}\par
\par\medskip\noindent\textbf{Definition} (\texttt{mat-map}).\;$\operatorname{mat\text{-}map}(\operatorname{gmm\text{-}h\_}, \operatorname{gmm\text{-}m\_}, \operatorname{gmm\text{-}n\_}, \operatorname{gmm\text{-}a\_}) :=$
\begin{equation*}
\begin{array}{@{}l@{}} \operatorname{matof}(\operatorname{gmm\text{-}m\_}, \operatorname{gmm\text{-}n\_}, \\ \quad (\lambda {[}\operatorname{gmmi\_}, \operatorname{gmmj\_}] \in (\operatorname{interval}(1, \operatorname{gmm\text{-}m\_}) \times \operatorname{interval}(1, \operatorname{gmm\text{-}n\_})).\; \\ \quad \quad \operatorname{gmm\text{-}h\_}(\operatorname{entry}(\operatorname{gmm\text{-}a\_}, \operatorname{gmmi\_}, \operatorname{gmmj\_})))) \end{array}
\end{equation*}

\par\medskip\noindent\textbf{Definition} (\texttt{op-of}).\;$\operatorname{op\text{-}of}(\operatorname{gof\text{-}md\_}, \operatorname{gof\text{-}t\_}, \operatorname{gof\text{-}p\_}) :=$
\begin{equation*}
(\lambda \operatorname{gofw\_} \in \operatorname{vec}(\operatorname{gof\text{-}md\_}).\; \operatorname{poly\text{-}op}(\operatorname{gof\text{-}md\_}, \operatorname{gof\text{-}t\_}, \operatorname{gof\text{-}p\_}, \operatorname{gofw\_}))
\end{equation*}

\par\medskip\noindent\textbf{Definition} (\texttt{ops}).\;$\operatorname{ops}(\operatorname{gos\text{-}md\_}, \operatorname{gos\text{-}t\_}) :=$
\begin{equation*}
\{\operatorname{gosf\_} \in (\operatorname{vec}(\operatorname{gos\text{-}md\_}) \to \operatorname{vec}(\operatorname{gos\text{-}md\_})) : \exists \operatorname{gosp\_} \in \operatorname{carr}(\operatorname{poly}(\operatorname{scal}(\operatorname{gos\text{-}md\_}))).\; (\operatorname{gosf\_} = \operatorname{op\text{-}of}(\operatorname{gos\text{-}md\_}, \operatorname{gos\text{-}t\_}, \operatorname{gosp\_}))\}
\end{equation*}

\par\medskip\noindent\textbf{Definition} (\texttt{op-add}).\;$\operatorname{op\text{-}add}(\operatorname{goa\text{-}md\_}, \operatorname{goa\text{-}t\_}) :=$
\begin{equation*}
\begin{array}{@{}l@{}} (\lambda {[}\operatorname{goaf\_}, \operatorname{goag\_}] \in (\operatorname{ops}(\operatorname{goa\text{-}md\_}, \operatorname{goa\text{-}t\_}) \times \operatorname{ops}(\operatorname{goa\text{-}md\_}, \operatorname{goa\text{-}t\_})).\; \\ \quad (\lambda \operatorname{goaw\_} \in \operatorname{vec}(\operatorname{goa\text{-}md\_}).\; \operatorname{vadd}(\operatorname{goa\text{-}md\_})(\operatorname{goaf\_}(\operatorname{goaw\_}), \operatorname{goag\_}(\operatorname{goaw\_})))) \end{array}
\end{equation*}

\par\medskip\noindent\textbf{Definition} (\texttt{op-mul}).\;$\operatorname{op\text{-}mul}(\operatorname{gom\text{-}md\_}, \operatorname{gom\text{-}t\_}) :=$
\begin{equation*}
\begin{array}{@{}l@{}} (\lambda {[}\operatorname{gomf\_}, \operatorname{gomg\_}] \in (\operatorname{ops}(\operatorname{gom\text{-}md\_}, \operatorname{gom\text{-}t\_}) \times \operatorname{ops}(\operatorname{gom\text{-}md\_}, \operatorname{gom\text{-}t\_})).\; \\ \quad (\lambda \operatorname{gomw\_} \in \operatorname{vec}(\operatorname{gom\text{-}md\_}).\; \operatorname{gomf\_}(\operatorname{gomg\_}(\operatorname{gomw\_})))) \end{array}
\end{equation*}

\par\medskip\noindent\textbf{Definition} (\texttt{op-neg}).\;$\operatorname{op\text{-}neg}(\operatorname{gon\text{-}md\_}, \operatorname{gon\text{-}t\_}) :=$
\begin{equation*}
\begin{array}{@{}l@{}} (\lambda \operatorname{gonf\_} \in \operatorname{ops}(\operatorname{gon\text{-}md\_}, \operatorname{gon\text{-}t\_}).\; \\ \quad (\lambda \operatorname{gonw\_} \in \operatorname{vec}(\operatorname{gon\text{-}md\_}).\; \operatorname{vneg}(\operatorname{gon\text{-}md\_})(\operatorname{gonf\_}(\operatorname{gonw\_})))) \end{array}
\end{equation*}

\par\medskip\noindent\textbf{Definition} (\texttt{op-zero}).\;$\operatorname{op\text{-}zero}(\operatorname{goz\text{-}md\_}) :=$ $(\lambda \operatorname{gozw\_} \in \operatorname{vec}(\operatorname{goz\text{-}md\_}).\; \operatorname{vzero}(\operatorname{goz\text{-}md\_}))$
\par\medskip\noindent\textbf{Definition} (\texttt{op-one}).\;$\operatorname{op\text{-}one}(\operatorname{goo\text{-}md\_}) :=$ $(\lambda \operatorname{goow\_} \in \operatorname{vec}(\operatorname{goo\text{-}md\_}).\; \operatorname{goow\_})$
\par\medskip\noindent\textbf{Definition} (\texttt{op-alg}).\;$\operatorname{op\text{-}alg}(\operatorname{gal\text{-}md\_}, \operatorname{gal\text{-}t\_}) :=$
\begin{equation*}
{[}\operatorname{ops}(\operatorname{gal\text{-}md\_}, \operatorname{gal\text{-}t\_}), \operatorname{op\text{-}add}(\operatorname{gal\text{-}md\_}, \operatorname{gal\text{-}t\_}), \operatorname{op\text{-}mul}(\operatorname{gal\text{-}md\_}, \operatorname{gal\text{-}t\_}), \operatorname{op\text{-}neg}(\operatorname{gal\text{-}md\_}, \operatorname{gal\text{-}t\_}), \operatorname{op\text{-}zero}(\operatorname{gal\text{-}md\_}), \operatorname{op\text{-}one}(\operatorname{gal\text{-}md\_})]
\end{equation*}

\par\medskip\noindent\textbf{Definition} (\texttt{ev-op}).\;$\operatorname{ev\text{-}op}(\operatorname{gev\text{-}md\_}, \operatorname{gev\text{-}t\_}) :=$
\begin{equation*}
(\lambda \operatorname{gevp\_} \in \operatorname{carr}(\operatorname{poly}(\operatorname{scal}(\operatorname{gev\text{-}md\_}))).\; \operatorname{op\text{-}of}(\operatorname{gev\text{-}md\_}, \operatorname{gev\text{-}t\_}, \operatorname{gevp\_}))
\end{equation*}

\par\medskip\noindent\textbf{Definition} (\texttt{mat-act}).\;$\operatorname{mat\text{-}act}(\operatorname{gma\text{-}md\_}, \operatorname{gma\text{-}n\_}, \operatorname{gma\text{-}p\_}, \operatorname{gma\text{-}u\_}) :=$
\begin{equation*}
\begin{array}{@{}l@{}} \operatorname{matof}(\operatorname{gma\text{-}n\_}, 1, \\ \quad (\lambda {[}\operatorname{gmai\_}, \operatorname{gmaj\_}] \in (\operatorname{interval}(1, \operatorname{gma\text{-}n\_}) \times \operatorname{interval}(1, 1)).\; \\ \quad \quad \operatorname{finsum}(\operatorname{module\text{-}vector\text{-}ag}(\operatorname{gma\text{-}md\_}), \\ \quad \quad \quad (\lambda \operatorname{gmak\_} \in \operatorname{interval}(1, \operatorname{gma\text{-}n\_}).\; \operatorname{entry}(\operatorname{gma\text{-}p\_}, \operatorname{gmai\_}, \operatorname{gmak\_})(\operatorname{entry}(\operatorname{gma\text{-}u\_}, \operatorname{gmak\_}, 1))), \\ \quad \quad \quad \operatorname{interval}(1, \operatorname{gma\text{-}n\_})))) \end{array}
\end{equation*}


\vnbbookfile{theorem-library/gal-op2}{theorem-library/gal-op2}
\noindent{\small gal-op2.scm -- the matrix of an operator in a basis is unique; (pq)(T) = p(T) q(T); nilpotent operators with a cyclic vector (galois-theory notes, Example 1.39); Cayley-Hamilton for an operator with a cyclic vector (GALOIS B6, 2026-10-26).  Destination: theorem-library/gal-op2.scm, library linear-algebra. Load window: ...}\par
\vnbprop{matrix-of-unique}{\vnbstmt{\begin{array}{@{}l@{}} \forall d.; \forall n \in \mathbb{N}, u \in \operatorname{mat}(n, 1, \operatorname{vec}(d)), t, a, b. \\ \quad \operatorname{is\text{-}module}(d) \wedge \\ \quad \operatorname{is\text{-}commutative\text{-}ring}(\operatorname{scal}(d)) \wedge \\ \quad \operatorname{rel\text{-}free}(d, n, u) \wedge \\ \quad \operatorname{is\text{-}matrix\text{-}of}(d, n, u, t, a) \wedge \\ \quad \operatorname{is\text{-}matrix\text{-}of}(d, n, u, t, b) \Rightarrow \\ \quad \quad a = b \end{array}}}
\vnbprop{cayley-hamilton-cyclic}{\vnbstmt{\begin{array}{@{}l@{}} \forall d, t, v.; \forall \mathit{nv} \in \mathbb{N}, \mathit{av}, a. \\ \quad \operatorname{is\text{-}vector\text{-}space}(d) \wedge \\ \quad \operatorname{is\text{-}hom\text{-}module}(d, d, t) \wedge \\ \quad \operatorname{is\text{-}cyclic\text{-}vector}(d, t, v) \wedge \\ \quad 1 \leq \mathit{nv} \wedge \\ \quad \forall i \in \operatorname{interval}(1, \mathit{nv}).\; (\operatorname{av}(i) \in \operatorname{carr}(\operatorname{scal}(d))) \wedge \\ \quad \operatorname{is\text{-}annihilating\text{-}poly}(d, t, v, \operatorname{monic\text{-}poly}(\operatorname{scal}(d), \mathit{av}, \mathit{nv})) \wedge \\ \quad \operatorname{is\text{-}matrix\text{-}of}(d, \mathit{nv}, \operatorname{krylov}(d, t, v, \mathit{nv}), t, a) \Rightarrow \\ \quad \quad \operatorname{char\text{-}poly}(\operatorname{scal}(d), a, \mathit{nv}) = \operatorname{monic\text{-}poly}(\operatorname{scal}(d), \mathit{av}, \mathit{nv}) \wedge \\ \quad \quad \forall w \in \operatorname{vec}(d).\; (\operatorname{poly\text{-}op}(d, t, \operatorname{char\text{-}poly}(\operatorname{scal}(d), a, \mathit{nv}), w) = \operatorname{vzero}(d)) \end{array}}}
\vnbprop{poly-op-monomial}{\vnbstmt{\begin{array}{@{}l@{}} \forall d, t.; \forall c \in \operatorname{carr}(\operatorname{scal}(d)), e \in \mathbb{N}, w \in \operatorname{vec}(d). \\ \quad (\operatorname{is\text{-}hom\text{-}module}(d, d, t) \wedge \operatorname{is\text{-}commutative\text{-}ring}(\operatorname{scal}(d))) \Rightarrow  \\ \quad \quad \operatorname{poly\text{-}op}(d, t, \operatorname{monomial}(\operatorname{scal}(d), c, e), w) = \operatorname{act}(d)(c, \operatorname{op\text{-}iter}(d, t, w, e)) \end{array}}}
\vnbprop{poly-op-padd}{\vnbstmt{\begin{array}{@{}l@{}} \forall d, t.; \forall p, q \in \operatorname{carr}(\operatorname{poly}(\operatorname{scal}(d))), w \in \operatorname{vec}(d). \\ \quad (\operatorname{is\text{-}hom\text{-}module}(d, d, t) \wedge \operatorname{is\text{-}commutative\text{-}ring}(\operatorname{scal}(d))) \Rightarrow  \\ \quad \quad \operatorname{poly\text{-}op}(d, t, \operatorname{add}(\operatorname{poly}(\operatorname{scal}(d)))(p, q), w) = \operatorname{vadd}(d)(\operatorname{poly\text{-}op}(d, t, p, w), \operatorname{poly\text{-}op}(d, t, q, w)) \end{array}}}
\vnbprop{op-iter-compose}{\vnbstmt{\begin{array}{@{}l@{}} \forall e \in \mathbb{N}, d, t.; \forall w \in \operatorname{vec}(d), j \in \mathbb{N}. \\ \quad \operatorname{is\text{-}hom\text{-}module}(d, d, t) \Rightarrow  \\ \quad \quad \operatorname{op\text{-}iter}(d, t, \operatorname{op\text{-}iter}(d, t, w, j), e) = \operatorname{op\text{-}iter}(d, t, w, (j + e)) \end{array}}}
\vnbprop{poly-op-monomial-mul-partial}{\vnbstmt{\begin{array}{@{}l@{}} \forall j \in \mathbb{N}, d, t.; \forall c \in \operatorname{carr}(\operatorname{scal}(d)), e \in \mathbb{N}, s \in \operatorname{carr}(\operatorname{poly}(\operatorname{scal}(d))), w \in \operatorname{vec}(d). \\ \quad (\operatorname{is\text{-}hom\text{-}module}(d, d, t) \wedge \operatorname{is\text{-}commutative\text{-}ring}(\operatorname{scal}(d))) \Rightarrow  \\ \quad \quad \operatorname{sum\text{-}ag}(\operatorname{module\text{-}vector\text{-}ag}(d), \\ \quad \quad \quad (\lambda k \in \mathbb{N}.\; \\ \quad \quad \quad \quad \operatorname{act}(d)(\operatorname{monalg\text{-}mul}(\operatorname{scal}(d), \operatorname{nn\text{-}add\text{-}monoid}, \operatorname{monomial}(\operatorname{scal}(d), c, e), s)(k), \operatorname{op\text{-}iter}(d, t, w, k))), \\ \quad \quad \quad (e + j)) \\ \quad \quad =\; \operatorname{act}(d)(c, \\ \quad \quad \quad \operatorname{op\text{-}iter}(d, t, \\ \quad \quad \quad \quad \operatorname{sum\text{-}ag}(\operatorname{module\text{-}vector\text{-}ag}(d), (\lambda k \in \mathbb{N}.\; \operatorname{act}(d)(s(k), \operatorname{op\text{-}iter}(d, t, w, k))), j), \\ \quad \quad \quad \quad e)) \end{array}}}
\vnbprop{poly-op-monomial-mul}{\vnbstmt{\begin{array}{@{}l@{}} \forall d, t.; \forall c \in \operatorname{carr}(\operatorname{scal}(d)), e \in \mathbb{N}, s \in \operatorname{carr}(\operatorname{poly}(\operatorname{scal}(d))), w \in \operatorname{vec}(d). \\ \quad (\operatorname{is\text{-}hom\text{-}module}(d, d, t) \wedge \operatorname{is\text{-}commutative\text{-}ring}(\operatorname{scal}(d))) \Rightarrow  \\ \quad \quad \operatorname{poly\text{-}op}(d, t, \operatorname{mul}(\operatorname{poly}(\operatorname{scal}(d)))(\operatorname{monomial}(\operatorname{scal}(d), c, e), s), w) \\ \quad \quad =\; \operatorname{act}(d)(c, \operatorname{op\text{-}iter}(d, t, \operatorname{poly\text{-}op}(d, t, s, w), e)) \end{array}}}
\vnbprop{poly-op-mul-bounded}{\vnbstmt{\begin{array}{@{}l@{}} \forall b \in \mathbb{N}, d, t.; \forall p, q \in \operatorname{carr}(\operatorname{poly}(\operatorname{scal}(d))), w \in \operatorname{vec}(d). \\ \quad (\operatorname{is\text{-}hom\text{-}module}(d, d, t) \wedge (\operatorname{is\text{-}commutative\text{-}ring}(\operatorname{scal}(d)) \wedge \operatorname{deg\text{-}bound}(\operatorname{scal}(d), p, b))) \Rightarrow  \\ \quad \quad \operatorname{poly\text{-}op}(d, t, \operatorname{mul}(\operatorname{poly}(\operatorname{scal}(d)))(p, q), w) = \operatorname{poly\text{-}op}(d, t, p, \operatorname{poly\text{-}op}(d, t, q, w)) \end{array}}}
\vnbprop{poly-op-mul}{\vnbstmt{\begin{array}{@{}l@{}} \forall d, t.; \forall p, q \in \operatorname{carr}(\operatorname{poly}(\operatorname{scal}(d))), w \in \operatorname{vec}(d). \\ \quad (\operatorname{is\text{-}hom\text{-}module}(d, d, t) \wedge \operatorname{is\text{-}commutative\text{-}ring}(\operatorname{scal}(d))) \Rightarrow  \\ \quad \quad \operatorname{poly\text{-}op}(d, t, \operatorname{mul}(\operatorname{poly}(\operatorname{scal}(d)))(p, q), w) = \operatorname{poly\text{-}op}(d, t, p, \operatorname{poly\text{-}op}(d, t, q, w)) \end{array}}}
\vnbprop{op-iter-vanishes-beyond}{\vnbstmt{\begin{array}{@{}l@{}} \forall d, t.; \forall v \in \operatorname{vec}(d), m, a \in \mathbb{N}. \\ \quad (\operatorname{is\text{-}hom\text{-}module}(d, d, t) \wedge ((\operatorname{op\text{-}iter}(d, t, v, m) = \operatorname{vzero}(d)) \wedge (m \leq a))) \Rightarrow  \\ \quad \quad \operatorname{op\text{-}iter}(d, t, v, a) = \operatorname{vzero}(d) \end{array}}}
\vnbprop{field-act-cancel}{\vnbstmt{\begin{array}{@{}l@{}} \forall d.; \forall c \in \operatorname{carr}(\operatorname{scal}(d)), x \in \operatorname{vec}(d). \\ \quad (\operatorname{is\text{-}module}(d) \wedge (\operatorname{is\text{-}field\text{-}ring}(\operatorname{scal}(d)) \wedge (\neg (c = \operatorname{zero}(\operatorname{scal}(d))) \wedge (\operatorname{act}(d)(c, x) = \operatorname{vzero}(d))))) \Rightarrow  \\ \quad \quad x = \operatorname{vzero}(d) \end{array}}}
\vnbprop{nilpotent-least-index}{\vnbstmt{\begin{array}{@{}l@{}} \forall d, t.; \forall v \in \operatorname{vec}(d), k \in \operatorname{carr}(\operatorname{poly}(\operatorname{scal}(d))), a, m \in \mathbb{N}. \\ \quad \operatorname{is\text{-}hom\text{-}module}(d, d, t) \wedge \\ \quad \operatorname{is\text{-}field\text{-}ring}(\operatorname{scal}(d)) \wedge \\ \quad \operatorname{poly\text{-}op}(d, t, k, v) = \operatorname{vzero}(d) \wedge \\ \quad \neg (k(a) = \operatorname{zero}(\operatorname{scal}(d))) \wedge \\ \quad \forall i \in \mathbb{N}.\; ((i < a) \Rightarrow (k(i) = \operatorname{zero}(\operatorname{scal}(d)))) \wedge \\ \quad \operatorname{op\text{-}iter}(d, t, v, m) = \operatorname{vzero}(d) \wedge \\ \quad \forall j \in \mathbb{N}.\; ((\operatorname{op\text{-}iter}(d, t, v, j) = \operatorname{vzero}(d)) \Rightarrow (m \leq j)) \Rightarrow \\ \quad \quad m \leq a \end{array}}}
\vnbprop{nilpotent-annihilator-low}{\vnbstmt{\begin{array}{@{}l@{}} \forall d, t.; \forall v \in \operatorname{vec}(d), k.; \forall i \in \mathbb{N}. \\ \quad \operatorname{is\text{-}hom\text{-}module}(d, d, t) \wedge \\ \quad \operatorname{is\text{-}field\text{-}ring}(\operatorname{scal}(d)) \wedge \\ \quad \operatorname{is\text{-}nilpotent\text{-}op}(d, t) \wedge \\ \quad \operatorname{is\text{-}annihilating\text{-}poly}(d, t, v, k) \wedge \\ \quad i < \operatorname{deg}(\operatorname{scal}(d), k) \Rightarrow \\ \quad \quad k(i) = \operatorname{zero}(\operatorname{scal}(d)) \end{array}}}
\vnbprop{nilpotent-annihilator-is-power}{\vnbstmt{\begin{array}{@{}l@{}} \forall d, t.; \forall v \in \operatorname{vec}(d), k. \\ \quad \operatorname{is\text{-}hom\text{-}module}(d, d, t) \wedge \\ \quad \operatorname{is\text{-}field\text{-}ring}(\operatorname{scal}(d)) \wedge \\ \quad \operatorname{is\text{-}nilpotent\text{-}op}(d, t) \wedge \\ \quad \operatorname{is\text{-}annihilating\text{-}poly}(d, t, v, k) \Rightarrow \\ \quad \quad k = \operatorname{monic\text{-}poly}(\operatorname{scal}(d), (\lambda i \in \operatorname{interval}(1, \operatorname{deg}(\operatorname{scal}(d), k)).\; \operatorname{zero}(\operatorname{scal}(d))), \operatorname{deg}(\operatorname{scal}(d), k)) \end{array}}}
\vnbprop{is-annihilating-poly-eq}{\vnbstmt{\begin{array}{@{}l@{}} \forall d, t, v, a, b. \\ \quad ((a = b) \wedge \operatorname{is\text{-}annihilating\text{-}poly}(d, t, v, a)) \Rightarrow  \\ \quad \quad \operatorname{is\text{-}annihilating\text{-}poly}(d, t, v, b) \end{array}}}
\vnbprop{nilpotent-cyclic-vector}{\vnbstmt{\begin{array}{@{}l@{}} \forall d, t, v, k. \\ \quad \operatorname{is\text{-}vector\text{-}space}(d) \wedge \\ \quad \operatorname{is\text{-}hom\text{-}module}(d, d, t) \wedge \\ \quad \operatorname{is\text{-}nilpotent\text{-}op}(d, t) \wedge \\ \quad \operatorname{is\text{-}cyclic\text{-}vector}(d, t, v) \wedge \\ \quad \operatorname{is\text{-}annihilating\text{-}poly}(d, t, v, k) \wedge \\ \quad 1 \leq \operatorname{deg}(\operatorname{scal}(d), k) \Rightarrow \\ \quad \quad k = \operatorname{monic\text{-}poly}(\operatorname{scal}(d), (\lambda i \in \operatorname{interval}(1, \operatorname{deg}(\operatorname{scal}(d), k)).\; \operatorname{zero}(\operatorname{scal}(d))), \operatorname{deg}(\operatorname{scal}(d), k)) \wedge \\ \quad \quad \operatorname{generates}(d, \operatorname{deg}(\operatorname{scal}(d), k), \operatorname{krylov}(d, t, v, \operatorname{deg}(\operatorname{scal}(d), k))) \wedge \\ \quad \quad \operatorname{rel\text{-}free}(d, \operatorname{deg}(\operatorname{scal}(d), k), \operatorname{krylov}(d, t, v, \operatorname{deg}(\operatorname{scal}(d), k))) \wedge \\ \quad \quad \operatorname{is\text{-}matrix\text{-}of}(d, \operatorname{deg}(\operatorname{scal}(d), k), \operatorname{krylov}(d, t, v, \operatorname{deg}(\operatorname{scal}(d), k)), t, \\ \quad \quad \quad \operatorname{companion}(\operatorname{scal}(d), (\lambda i \in \operatorname{interval}(1, \operatorname{deg}(\operatorname{scal}(d), k)).\; \operatorname{zero}(\operatorname{scal}(d))), \operatorname{deg}(\operatorname{scal}(d), k))) \wedge \\ \quad \quad \operatorname{is\text{-}minimal\text{-}poly}(d, t, k) \end{array}}}
\vnbprop{annihilating-poly-divides}{\vnbstmt{\begin{array}{@{}l@{}} \forall d, t.; \forall v \in \operatorname{vec}(d), k.; \forall q \in \operatorname{carr}(\operatorname{poly}(\operatorname{scal}(d))). \\ \quad \operatorname{is\text{-}hom\text{-}module}(d, d, t) \wedge \\ \quad \operatorname{is\text{-}field\text{-}ring}(\operatorname{scal}(d)) \wedge \\ \quad \operatorname{is\text{-}annihilating\text{-}poly}(d, t, v, k) \wedge \\ \quad \operatorname{poly\text{-}op}(d, t, q, v) = \operatorname{vzero}(d) \Rightarrow \\ \quad \quad \exists s \in \operatorname{carr}(\operatorname{poly}(\operatorname{scal}(d))).\; (q = \operatorname{mul}(\operatorname{poly}(\operatorname{scal}(d)))(k, s)) \end{array}}}

\vnbbookfile{theorem-library/la-z3-quotient}{theorem-library/la-z3-quotient}
\noindent{\small la-z3-quotient.scm -- the quotient module E/F of a module by a submodule (LA-Z3, 2026-10-24).  Vocabulary: structure-library/la-z3-vocab (QMOD-REL, QMOD-SETOID, QUOTIENT-MODULE).  Prefix lzq-.}\par
\vnbprop{lzq-vadd-fun}{\vnbstmt{\begin{array}{@{}l@{}} \forall d. \\ \quad \operatorname{is\text{-}module}(d) \Rightarrow (\operatorname{vadd}(d) \in ((\operatorname{vec}(d) \times \operatorname{vec}(d)) \to \operatorname{vec}(d))) \end{array}}}
\vnbprop{lzq-vneg-fun}{\vnbstmt{\forall d.\; (\operatorname{is\text{-}module}(d) \Rightarrow (\operatorname{vneg}(d) \in (\operatorname{vec}(d) \to \operatorname{vec}(d))))}}
\vnbprop{lzq-act-fun}{\vnbstmt{\begin{array}{@{}l@{}} \forall d. \\ \quad \operatorname{is\text{-}module}(d) \Rightarrow (\operatorname{act}(d) \in ((\operatorname{carr}(\operatorname{scal}(d)) \times \operatorname{vec}(d)) \to \operatorname{vec}(d))) \end{array}}}
\vnbprop{lzq-comm-prop}{\vnbstmt{\forall d.\; (\operatorname{is\text{-}module}(d) \Rightarrow \operatorname{is\text{-}commutative}(\operatorname{vadd}(d), \operatorname{vec}(d)))}}
\vnbprop{lzq-ident-prop}{\vnbstmt{\forall d.\; (\operatorname{is\text{-}module}(d) \Rightarrow \operatorname{is\text{-}identity}(\operatorname{vadd}(d), \operatorname{vzero}(d), \operatorname{vec}(d)))}}
\vnbprop{lzq-inv-prop}{\vnbstmt{\begin{array}{@{}l@{}} \forall d. \\ \quad \operatorname{is\text{-}module}(d) \Rightarrow \operatorname{has\text{-}inverses}(\operatorname{vadd}(d), \operatorname{vzero}(d), \operatorname{vneg}(d), \operatorname{vec}(d)) \end{array}}}
\vnbprop{lzq-vadd-comm}{\vnbstmt{\begin{array}{@{}l@{}} \forall d, x, y. \\ \quad (\operatorname{is\text{-}module}(d) \wedge ((x \in \operatorname{vec}(d)) \wedge (y \in \operatorname{vec}(d)))) \Rightarrow  \\ \quad \quad \operatorname{vadd}(d)(x, y) = \operatorname{vadd}(d)(y, x) \end{array}}}
\vnbprop{lzq-inv-l}{\vnbstmt{\begin{array}{@{}l@{}} \forall d.; \forall x \in \operatorname{vec}(d). \\ \quad \operatorname{is\text{-}module}(d) \Rightarrow (\operatorname{vadd}(d)(\operatorname{vneg}(d)(x), x) = \operatorname{vzero}(d)) \end{array}}}
\vnbprop{lzq-neg-sum}{\vnbstmt{\begin{array}{@{}l@{}} \forall d, a, b. \\ \quad (\operatorname{is\text{-}module}(d) \wedge ((a \in \operatorname{vec}(d)) \wedge (b \in \operatorname{vec}(d)))) \Rightarrow  \\ \quad \quad \operatorname{vneg}(d)(\operatorname{vadd}(d)(a, b)) = \operatorname{vadd}(d)(\operatorname{vneg}(d)(a), \operatorname{vneg}(d)(b)) \end{array}}}
\vnbprop{lzq-d-sym}{\vnbstmt{\begin{array}{@{}l@{}} \forall d, x, y. \\ \quad (\operatorname{is\text{-}module}(d) \wedge ((x \in \operatorname{vec}(d)) \wedge (y \in \operatorname{vec}(d)))) \Rightarrow  \\ \quad \quad \operatorname{vadd}(d)(y, \operatorname{vneg}(d)(x)) = \operatorname{vneg}(d)(\operatorname{vadd}(d)(x, \operatorname{vneg}(d)(y))) \end{array}}}
\vnbprop{lzq-d-trans}{\vnbstmt{\begin{array}{@{}l@{}} \forall d, x, y, z. \\ \quad (\operatorname{is\text{-}module}(d) \wedge ((x \in \operatorname{vec}(d)) \wedge ((y \in \operatorname{vec}(d)) \wedge (z \in \operatorname{vec}(d))))) \Rightarrow  \\ \quad \quad \operatorname{vadd}(d)(x, \operatorname{vneg}(d)(z)) = \operatorname{vadd}(d)(\operatorname{vadd}(d)(x, \operatorname{vneg}(d)(y)), \operatorname{vadd}(d)(y, \operatorname{vneg}(d)(z))) \end{array}}}
\vnbprop{lzq-d-vadd}{\vnbstmt{\begin{array}{@{}l@{}} \forall d, x, y, u, w. \\ \quad (\operatorname{is\text{-}module}(d) \wedge ((x \in \operatorname{vec}(d)) \wedge ((y \in \operatorname{vec}(d)) \wedge ((u \in \operatorname{vec}(d)) \wedge (w \in \operatorname{vec}(d)))))) \Rightarrow  \\ \quad \quad \operatorname{vadd}(d)(\operatorname{vadd}(d)(x, y), \operatorname{vneg}(d)(\operatorname{vadd}(d)(u, w))) \\ \quad \quad =\; \operatorname{vadd}(d)(\operatorname{vadd}(d)(x, \operatorname{vneg}(d)(u)), \operatorname{vadd}(d)(y, \operatorname{vneg}(d)(w))) \end{array}}}
\vnbprop{lzq-d-vneg}{\vnbstmt{\begin{array}{@{}l@{}} \forall d, x, u. \\ \quad (\operatorname{is\text{-}module}(d) \wedge ((x \in \operatorname{vec}(d)) \wedge (u \in \operatorname{vec}(d)))) \Rightarrow  \\ \quad \quad \operatorname{vadd}(d)(\operatorname{vneg}(d)(x), \operatorname{vneg}(d)(\operatorname{vneg}(d)(u))) = \operatorname{vneg}(d)(\operatorname{vadd}(d)(x, \operatorname{vneg}(d)(u))) \end{array}}}
\vnbprop{lzq-d-act}{\vnbstmt{\begin{array}{@{}l@{}} \forall d, r, x, u. \\ \quad (\operatorname{is\text{-}module}(d) \wedge ((r \in \operatorname{carr}(\operatorname{scal}(d))) \wedge ((x \in \operatorname{vec}(d)) \wedge (u \in \operatorname{vec}(d))))) \Rightarrow  \\ \quad \quad \operatorname{vadd}(d)(\operatorname{act}(d)(r, x), \operatorname{vneg}(d)(\operatorname{act}(d)(r, u))) = \operatorname{act}(d)(r, \operatorname{vadd}(d)(x, \operatorname{vneg}(d)(u))) \end{array}}}
\vnbprop{qmod-rel-unfold}{\vnbstmt{\begin{array}{@{}l@{}} \forall d, m. \\ \quad (\operatorname{qmod\text{-}rel}(d, m) \simeq \{p \in (\operatorname{vec}(d) \times \operatorname{vec}(d)) : (\operatorname{vadd}(d)(\operatorname{nth}(1, p), \operatorname{vneg}(d)(\operatorname{nth}(2, p))) \in m)\}) \end{array}}}
\vnbprop{qmod-rel-in}{\vnbstmt{\begin{array}{@{}l@{}} \forall d, m, a, b. \\ \quad \operatorname{is\text{-}module}(d) \wedge \\ \quad \operatorname{is\text{-}submodule}(d, m) \wedge \\ \quad a \in \operatorname{vec}(d) \wedge \\ \quad b \in \operatorname{vec}(d) \wedge \\ \quad \operatorname{vadd}(d)(a, \operatorname{vneg}(d)(b)) \in m \Rightarrow \\ \quad \quad {[}a, b] \in \operatorname{qmod\text{-}rel}(d, m) \end{array}}}
\vnbprop{qmod-rel-out}{\vnbstmt{\begin{array}{@{}l@{}} \forall d, m, a, b. \\ \quad (\operatorname{is\text{-}module}(d) \wedge (\operatorname{is\text{-}submodule}(d, m) \wedge ({[}a, b] \in \operatorname{qmod\text{-}rel}(d, m)))) \Rightarrow  \\ \quad \quad \operatorname{vadd}(d)(a, \operatorname{vneg}(d)(b)) \in m \end{array}}}
\vnbprop{qmod-rel-is-equivalence}{\vnbstmt{\begin{array}{@{}l@{}} \forall d, m. \\ \quad (\operatorname{is\text{-}module}(d) \wedge \operatorname{is\text{-}submodule}(d, m)) \Rightarrow  \\ \quad \quad \operatorname{is\text{-}equivalence}(\operatorname{qmod\text{-}rel}(d, m), \operatorname{vec}(d)) \end{array}}}
\vnbprop{qmod-setoid-is-setoid}{\vnbstmt{\begin{array}{@{}l@{}} \forall d, m. \\ \quad (\operatorname{is\text{-}module}(d) \wedge \operatorname{is\text{-}submodule}(d, m)) \Rightarrow \operatorname{is\text{-}setoid}(\operatorname{qmod\text{-}setoid}(d, m)) \end{array}}}
\vnbprop{qmod-setoid-pts}{\vnbstmt{\begin{array}{@{}l@{}} \forall d, m. \\ \quad (\operatorname{is\text{-}module}(d) \wedge \operatorname{is\text{-}submodule}(d, m)) \Rightarrow (\operatorname{pts}(\operatorname{qmod\text{-}setoid}(d, m)) \simeq \operatorname{vec}(d)) \end{array}}}
\vnbprop{qmod-setoid-rel}{\vnbstmt{\begin{array}{@{}l@{}} \forall d, m. \\ \quad (\operatorname{is\text{-}module}(d) \wedge \operatorname{is\text{-}submodule}(d, m)) \Rightarrow  \\ \quad \quad (\operatorname{rel}(\operatorname{qmod\text{-}setoid}(d, m)) \simeq \operatorname{qmod\text{-}rel}(d, m)) \end{array}}}
\vnbprop{qmod-related-out}{\vnbstmt{\begin{array}{@{}l@{}} \forall d, m, a, b. \\ \quad (\operatorname{is\text{-}module}(d) \wedge (\operatorname{is\text{-}submodule}(d, m) \wedge \operatorname{related}(\operatorname{qmod\text{-}setoid}(d, m), a, b))) \Rightarrow  \\ \quad \quad \operatorname{vadd}(d)(a, \operatorname{vneg}(d)(b)) \in m \end{array}}}
\vnbprop{qmod-related-in}{\vnbstmt{\begin{array}{@{}l@{}} \forall d, m, a, b. \\ \quad \operatorname{is\text{-}module}(d) \wedge \\ \quad \operatorname{is\text{-}submodule}(d, m) \wedge \\ \quad a \in \operatorname{vec}(d) \wedge \\ \quad b \in \operatorname{vec}(d) \wedge \\ \quad \operatorname{vadd}(d)(a, \operatorname{vneg}(d)(b)) \in m \Rightarrow \\ \quad \quad \operatorname{related}(\operatorname{qmod\text{-}setoid}(d, m), a, b) \end{array}}}
\vnbprop{qmod-class-eq-iff}{\vnbstmt{\begin{array}{@{}l@{}} \forall d, m, a, b. \\ \quad (\operatorname{is\text{-}module}(d) \wedge (\operatorname{is\text{-}submodule}(d, m) \wedge ((a \in \operatorname{vec}(d)) \wedge (b \in \operatorname{vec}(d))))) \Rightarrow  \\ \quad \quad (\operatorname{class}(\operatorname{qmod\text{-}setoid}(d, m), a) = \operatorname{class}(\operatorname{qmod\text{-}setoid}(d, m), b)) \Leftrightarrow (\operatorname{vadd}(d)(a, \operatorname{vneg}(d)(b)) \in m) \end{array}}}
\vnbprop{qmod-lva-fun}{\vnbstmt{\begin{array}{@{}l@{}} \forall d, m. \\ \quad (\operatorname{is\text{-}module}(d) \wedge \operatorname{is\text{-}submodule}(d, m)) \Rightarrow  \\ \quad \quad (\lambda {[}\operatorname{qma\_}, \operatorname{qmb\_}] \in (\operatorname{vec}(d) \times \operatorname{vec}(d)).\; \operatorname{class}(\operatorname{qmod\text{-}setoid}(d, m), \operatorname{vadd}(d)(\operatorname{qma\_}, \operatorname{qmb\_}))) \\ \quad \quad \in\; ((\operatorname{pts}(\operatorname{qmod\text{-}setoid}(d, m)) \times \operatorname{pts}(\operatorname{qmod\text{-}setoid}(d, m))) \to \operatorname{quotient}(\operatorname{qmod\text{-}setoid}(d, m))) \end{array}}}
\vnbprop{qmod-lvn-fun}{\vnbstmt{\begin{array}{@{}l@{}} \forall d, m. \\ \quad (\operatorname{is\text{-}module}(d) \wedge \operatorname{is\text{-}submodule}(d, m)) \Rightarrow  \\ \quad \quad (\lambda a \in \operatorname{vec}(d).\; \operatorname{class}(\operatorname{qmod\text{-}setoid}(d, m), \operatorname{vneg}(d)(a))) \\ \quad \quad \in\; (\operatorname{pts}(\operatorname{qmod\text{-}setoid}(d, m)) \to \operatorname{quotient}(\operatorname{qmod\text{-}setoid}(d, m))) \end{array}}}
\vnbprop{qmod-lsc-fun}{\vnbstmt{\begin{array}{@{}l@{}} \forall d, m.; \forall r \in \operatorname{carr}(\operatorname{scal}(d)). \\ \quad (\operatorname{is\text{-}module}(d) \wedge \operatorname{is\text{-}submodule}(d, m)) \Rightarrow  \\ \quad \quad (\lambda a \in \operatorname{vec}(d).\; \operatorname{class}(\operatorname{qmod\text{-}setoid}(d, m), \operatorname{act}(d)(r, a))) \\ \quad \quad \in\; (\operatorname{pts}(\operatorname{qmod\text{-}setoid}(d, m)) \to \operatorname{quotient}(\operatorname{qmod\text{-}setoid}(d, m))) \end{array}}}
\vnbprop{qmod-lva-respects}{\vnbstmt{\begin{array}{@{}l@{}} \forall d, m. \\ \quad (\operatorname{is\text{-}module}(d) \wedge \operatorname{is\text{-}submodule}(d, m)) \Rightarrow  \\ \quad \quad \operatorname{respects2}(\operatorname{qmod\text{-}setoid}(d, m), \\ \quad \quad \quad (\lambda {[}\operatorname{qma\_}, \operatorname{qmb\_}] \in (\operatorname{vec}(d) \times \operatorname{vec}(d)).\; \operatorname{class}(\operatorname{qmod\text{-}setoid}(d, m), \operatorname{vadd}(d)(\operatorname{qma\_}, \operatorname{qmb\_})))) \end{array}}}
\vnbprop{qmod-lvn-respects}{\vnbstmt{\begin{array}{@{}l@{}} \forall d, m. \\ \quad (\operatorname{is\text{-}module}(d) \wedge \operatorname{is\text{-}submodule}(d, m)) \Rightarrow  \\ \quad \quad \operatorname{respects}(\operatorname{qmod\text{-}setoid}(d, m), (\lambda a \in \operatorname{vec}(d).\; \operatorname{class}(\operatorname{qmod\text{-}setoid}(d, m), \operatorname{vneg}(d)(a)))) \end{array}}}
\vnbprop{qmod-lsc-respects}{\vnbstmt{\begin{array}{@{}l@{}} \forall d, m.; \forall r \in \operatorname{carr}(\operatorname{scal}(d)). \\ \quad (\operatorname{is\text{-}module}(d) \wedge \operatorname{is\text{-}submodule}(d, m)) \Rightarrow  \\ \quad \quad \operatorname{respects}(\operatorname{qmod\text{-}setoid}(d, m), (\lambda a \in \operatorname{vec}(d).\; \operatorname{class}(\operatorname{qmod\text{-}setoid}(d, m), \operatorname{act}(d)(r, a)))) \end{array}}}
\vnbprop{qmod-scale-fun}{\vnbstmt{\begin{array}{@{}l@{}} \forall d, m.; \forall r \in \operatorname{carr}(\operatorname{scal}(d)). \\ \quad (\operatorname{is\text{-}module}(d) \wedge \operatorname{is\text{-}submodule}(d, m)) \Rightarrow  \\ \quad \quad \operatorname{qmod\text{-}scale}(d, m, r) \in (\operatorname{quotient}(\operatorname{qmod\text{-}setoid}(d, m)) \to \operatorname{quotient}(\operatorname{qmod\text{-}setoid}(d, m))) \end{array}}}
\vnbprop{qmod-vadd-fun0}{\vnbstmt{\begin{array}{@{}l@{}} \forall d, m. \\ \quad (\operatorname{is\text{-}module}(d) \wedge \operatorname{is\text{-}submodule}(d, m)) \Rightarrow  \\ \quad \quad \operatorname{qmod\text{-}vadd}(d, m) \\ \quad \quad \in\; ((\operatorname{quotient}(\operatorname{qmod\text{-}setoid}(d, m)) \times \operatorname{quotient}(\operatorname{qmod\text{-}setoid}(d, m))) \to \operatorname{quotient}(\operatorname{qmod\text{-}setoid}(d, m))) \end{array}}}
\vnbprop{qmod-vneg-fun0}{\vnbstmt{\begin{array}{@{}l@{}} \forall d, m. \\ \quad (\operatorname{is\text{-}module}(d) \wedge \operatorname{is\text{-}submodule}(d, m)) \Rightarrow  \\ \quad \quad \operatorname{qmod\text{-}vneg}(d, m) \in (\operatorname{quotient}(\operatorname{qmod\text{-}setoid}(d, m)) \to \operatorname{quotient}(\operatorname{qmod\text{-}setoid}(d, m))) \end{array}}}
\vnbprop{qmod-act-fun0}{\vnbstmt{\begin{array}{@{}l@{}} \forall d, m. \\ \quad (\operatorname{is\text{-}module}(d) \wedge \operatorname{is\text{-}submodule}(d, m)) \Rightarrow  \\ \quad \quad \operatorname{qmod\text{-}act}(d, m) \\ \quad \quad \in\; ((\operatorname{carr}(\operatorname{scal}(d)) \times \operatorname{quotient}(\operatorname{qmod\text{-}setoid}(d, m))) \to \operatorname{quotient}(\operatorname{qmod\text{-}setoid}(d, m))) \end{array}}}
\vnbprop{qmod-ro-scal}{\vnbstmt{\begin{array}{@{}l@{}} \forall d, m. \\ \quad (\operatorname{is\text{-}module}(d) \wedge \operatorname{is\text{-}submodule}(d, m)) \Rightarrow  \\ \quad \quad (\operatorname{scal}(\operatorname{quotient\text{-}module}(d, m)) \simeq \operatorname{scal}(d)) \end{array}}}
\vnbprop{qmod-ro-vec}{\vnbstmt{\begin{array}{@{}l@{}} \forall d, m. \\ \quad (\operatorname{is\text{-}module}(d) \wedge \operatorname{is\text{-}submodule}(d, m)) \Rightarrow  \\ \quad \quad (\operatorname{vec}(\operatorname{quotient\text{-}module}(d, m)) \simeq \operatorname{quotient}(\operatorname{qmod\text{-}setoid}(d, m))) \end{array}}}
\vnbprop{qmod-ro-vadd}{\vnbstmt{\begin{array}{@{}l@{}} \forall d, m. \\ \quad (\operatorname{is\text{-}module}(d) \wedge \operatorname{is\text{-}submodule}(d, m)) \Rightarrow  \\ \quad \quad (\operatorname{vadd}(\operatorname{quotient\text{-}module}(d, m)) \simeq \operatorname{qmod\text{-}vadd}(d, m)) \end{array}}}
\vnbprop{qmod-ro-vzero}{\vnbstmt{\begin{array}{@{}l@{}} \forall d, m. \\ \quad (\operatorname{is\text{-}module}(d) \wedge \operatorname{is\text{-}submodule}(d, m)) \Rightarrow  \\ \quad \quad (\operatorname{vzero}(\operatorname{quotient\text{-}module}(d, m)) \simeq \operatorname{class}(\operatorname{qmod\text{-}setoid}(d, m), \operatorname{vzero}(d))) \end{array}}}
\vnbprop{qmod-ro-vneg}{\vnbstmt{\begin{array}{@{}l@{}} \forall d, m. \\ \quad (\operatorname{is\text{-}module}(d) \wedge \operatorname{is\text{-}submodule}(d, m)) \Rightarrow  \\ \quad \quad (\operatorname{vneg}(\operatorname{quotient\text{-}module}(d, m)) \simeq \operatorname{qmod\text{-}vneg}(d, m)) \end{array}}}
\vnbprop{qmod-ro-act}{\vnbstmt{\begin{array}{@{}l@{}} \forall d, m. \\ \quad (\operatorname{is\text{-}module}(d) \wedge \operatorname{is\text{-}submodule}(d, m)) \Rightarrow  \\ \quad \quad (\operatorname{act}(\operatorname{quotient\text{-}module}(d, m)) \simeq \operatorname{qmod\text{-}act}(d, m)) \end{array}}}
\vnbprop{qmod-vadd-computes}{\vnbstmt{\begin{array}{@{}l@{}} \forall d, m, a, b. \\ \quad (\operatorname{is\text{-}module}(d) \wedge (\operatorname{is\text{-}submodule}(d, m) \wedge ((a \in \operatorname{vec}(d)) \wedge (b \in \operatorname{vec}(d))))) \Rightarrow  \\ \quad \quad \operatorname{vadd}(\operatorname{quotient\text{-}module}(d, m))(\operatorname{class}(\operatorname{qmod\text{-}setoid}(d, m), a), \operatorname{class}(\operatorname{qmod\text{-}setoid}(d, m), b)) \\ \quad \quad =\; \operatorname{class}(\operatorname{qmod\text{-}setoid}(d, m), \operatorname{vadd}(d)(a, b)) \end{array}}}
\vnbprop{qmod-vneg-computes}{\vnbstmt{\begin{array}{@{}l@{}} \forall d, m.; \forall a \in \operatorname{vec}(d). \\ \quad (\operatorname{is\text{-}module}(d) \wedge \operatorname{is\text{-}submodule}(d, m)) \Rightarrow  \\ \quad \quad \operatorname{vneg}(\operatorname{quotient\text{-}module}(d, m))(\operatorname{class}(\operatorname{qmod\text{-}setoid}(d, m), a)) = \operatorname{class}(\operatorname{qmod\text{-}setoid}(d, m), \operatorname{vneg}(d)(a)) \end{array}}}
\vnbprop{qmod-scale-computes}{\vnbstmt{\begin{array}{@{}l@{}} \forall d, m, r, a. \\ \quad (\operatorname{is\text{-}module}(d) \wedge (\operatorname{is\text{-}submodule}(d, m) \wedge ((r \in \operatorname{carr}(\operatorname{scal}(d))) \wedge (a \in \operatorname{vec}(d))))) \Rightarrow  \\ \quad \quad \operatorname{qmod\text{-}scale}(d, m, r)(\operatorname{class}(\operatorname{qmod\text{-}setoid}(d, m), a)) = \operatorname{class}(\operatorname{qmod\text{-}setoid}(d, m), \operatorname{act}(d)(r, a)) \end{array}}}
\vnbprop{qmod-act-computes}{\vnbstmt{\begin{array}{@{}l@{}} \forall d, m, r, a. \\ \quad (\operatorname{is\text{-}module}(d) \wedge (\operatorname{is\text{-}submodule}(d, m) \wedge ((r \in \operatorname{carr}(\operatorname{scal}(d))) \wedge (a \in \operatorname{vec}(d))))) \Rightarrow  \\ \quad \quad \operatorname{act}(\operatorname{quotient\text{-}module}(d, m))(r, \operatorname{class}(\operatorname{qmod\text{-}setoid}(d, m), a)) \\ \quad \quad =\; \operatorname{class}(\operatorname{qmod\text{-}setoid}(d, m), \operatorname{act}(d)(r, a)) \end{array}}}
\vnbprop{qmod-length}{\vnbstmt{\begin{array}{@{}l@{}} \forall d, m. \\ \quad (\operatorname{is\text{-}module}(d) \wedge \operatorname{is\text{-}submodule}(d, m)) \Rightarrow (\operatorname{length}(\operatorname{quotient\text{-}module}(d, m)) = 6) \end{array}}}
\vnbprop{qmod-scal-ring}{\vnbstmt{\begin{array}{@{}l@{}} \forall d, m. \\ \quad (\operatorname{is\text{-}module}(d) \wedge \operatorname{is\text{-}submodule}(d, m)) \Rightarrow \operatorname{is\text{-}ring}(\operatorname{scal}(\operatorname{quotient\text{-}module}(d, m))) \end{array}}}
\vnbprop{qmod-vec-set}{\vnbstmt{\begin{array}{@{}l@{}} \forall d, m. \\ \quad (\operatorname{is\text{-}module}(d) \wedge \operatorname{is\text{-}submodule}(d, m)) \Rightarrow (\operatorname{vec}(\operatorname{quotient\text{-}module}(d, m)) \in \mathbf{Set}) \end{array}}}
\vnbprop{qmod-vzero-in}{\vnbstmt{\begin{array}{@{}l@{}} \forall d, m. \\ \quad (\operatorname{is\text{-}module}(d) \wedge \operatorname{is\text{-}submodule}(d, m)) \Rightarrow  \\ \quad \quad \operatorname{vzero}(\operatorname{quotient\text{-}module}(d, m)) \in \operatorname{vec}(\operatorname{quotient\text{-}module}(d, m)) \end{array}}}
\vnbprop{qmod-vadd-fun}{\vnbstmt{\begin{array}{@{}l@{}} \forall d, m. \\ \quad (\operatorname{is\text{-}module}(d) \wedge \operatorname{is\text{-}submodule}(d, m)) \Rightarrow  \\ \quad \quad \operatorname{vadd}(\operatorname{quotient\text{-}module}(d, m)) \\ \quad \quad \in\; ((\operatorname{vec}(\operatorname{quotient\text{-}module}(d, m)) \times \operatorname{vec}(\operatorname{quotient\text{-}module}(d, m))) \to \operatorname{vec}(\operatorname{quotient\text{-}module}(d, m))) \end{array}}}
\vnbprop{qmod-vneg-fun}{\vnbstmt{\begin{array}{@{}l@{}} \forall d, m. \\ \quad (\operatorname{is\text{-}module}(d) \wedge \operatorname{is\text{-}submodule}(d, m)) \Rightarrow  \\ \quad \quad \operatorname{vneg}(\operatorname{quotient\text{-}module}(d, m)) \in (\operatorname{vec}(\operatorname{quotient\text{-}module}(d, m)) \to \operatorname{vec}(\operatorname{quotient\text{-}module}(d, m))) \end{array}}}
\vnbprop{qmod-act-fun}{\vnbstmt{\begin{array}{@{}l@{}} \forall d, m. \\ \quad (\operatorname{is\text{-}module}(d) \wedge \operatorname{is\text{-}submodule}(d, m)) \Rightarrow  \\ \quad \quad \operatorname{act}(\operatorname{quotient\text{-}module}(d, m)) \\ \quad \quad \in\; ((\operatorname{carr}(\operatorname{scal}(\operatorname{quotient\text{-}module}(d, m))) \times \operatorname{vec}(\operatorname{quotient\text{-}module}(d, m))) \to \operatorname{vec}(\operatorname{quotient\text{-}module}(d, m))) \end{array}}}
\vnbprop{qmod-assoc}{\vnbstmt{\begin{array}{@{}l@{}} \forall d, m. \\ \quad (\operatorname{is\text{-}module}(d) \wedge \operatorname{is\text{-}submodule}(d, m)) \Rightarrow  \\ \quad \quad \operatorname{is\text{-}associative}(\operatorname{vadd}(\operatorname{quotient\text{-}module}(d, m)), \operatorname{vec}(\operatorname{quotient\text{-}module}(d, m))) \end{array}}}
\vnbprop{qmod-comm}{\vnbstmt{\begin{array}{@{}l@{}} \forall d, m. \\ \quad (\operatorname{is\text{-}module}(d) \wedge \operatorname{is\text{-}submodule}(d, m)) \Rightarrow  \\ \quad \quad \operatorname{is\text{-}commutative}(\operatorname{vadd}(\operatorname{quotient\text{-}module}(d, m)), \operatorname{vec}(\operatorname{quotient\text{-}module}(d, m))) \end{array}}}
\vnbprop{qmod-identity}{\vnbstmt{\begin{array}{@{}l@{}} \forall d, m. \\ \quad (\operatorname{is\text{-}module}(d) \wedge \operatorname{is\text{-}submodule}(d, m)) \Rightarrow  \\ \quad \quad \operatorname{is\text{-}identity}(\operatorname{vadd}(\operatorname{quotient\text{-}module}(d, m)), \operatorname{vzero}(\operatorname{quotient\text{-}module}(d, m)), \operatorname{vec}(\operatorname{quotient\text{-}module}(d, m))) \end{array}}}
\vnbprop{qmod-inverses}{\vnbstmt{\begin{array}{@{}l@{}} \forall d, m. \\ \quad (\operatorname{is\text{-}module}(d) \wedge \operatorname{is\text{-}submodule}(d, m)) \Rightarrow  \\ \quad \quad \operatorname{has\text{-}inverses}(\operatorname{vadd}(\operatorname{quotient\text{-}module}(d, m)), \operatorname{vzero}(\operatorname{quotient\text{-}module}(d, m)), \\ \quad \quad \quad \operatorname{vneg}(\operatorname{quotient\text{-}module}(d, m)), \operatorname{vec}(\operatorname{quotient\text{-}module}(d, m))) \end{array}}}
\vnbprop{qmod-act-distrib-vec}{\vnbstmt{\begin{array}{@{}l@{}} \forall d, m.; \forall r \in \operatorname{carr}(\operatorname{scal}(\operatorname{quotient\text{-}module}(d, m))), x, y \in \operatorname{vec}(\operatorname{quotient\text{-}module}(d, m)). \\ \quad (\operatorname{is\text{-}module}(d) \wedge \operatorname{is\text{-}submodule}(d, m)) \Rightarrow  \\ \quad \quad \operatorname{act}(\operatorname{quotient\text{-}module}(d, m))(r, \operatorname{vadd}(\operatorname{quotient\text{-}module}(d, m))(x, y)) \\ \quad \quad =\; \operatorname{vadd}(\operatorname{quotient\text{-}module}(d, m))(\operatorname{act}(\operatorname{quotient\text{-}module}(d, m))(r, x), \operatorname{act}(\operatorname{quotient\text{-}module}(d, m))(r, y)) \end{array}}}
\vnbprop{qmod-act-distrib-scalar}{\vnbstmt{\begin{array}{@{}l@{}} \forall d, m.; \forall r, s \in \operatorname{carr}(\operatorname{scal}(\operatorname{quotient\text{-}module}(d, m))), x \in \operatorname{vec}(\operatorname{quotient\text{-}module}(d, m)). \\ \quad (\operatorname{is\text{-}module}(d) \wedge \operatorname{is\text{-}submodule}(d, m)) \Rightarrow  \\ \quad \quad \operatorname{act}(\operatorname{quotient\text{-}module}(d, m))(\operatorname{add}(\operatorname{scal}(\operatorname{quotient\text{-}module}(d, m)))(r, s), x) \\ \quad \quad =\; \operatorname{vadd}(\operatorname{quotient\text{-}module}(d, m))(\operatorname{act}(\operatorname{quotient\text{-}module}(d, m))(r, x), \operatorname{act}(\operatorname{quotient\text{-}module}(d, m))(s, x)) \end{array}}}
\vnbprop{qmod-act-mul-compat}{\vnbstmt{\begin{array}{@{}l@{}} \forall d, m.; \forall r, s \in \operatorname{carr}(\operatorname{scal}(\operatorname{quotient\text{-}module}(d, m))), x \in \operatorname{vec}(\operatorname{quotient\text{-}module}(d, m)). \\ \quad (\operatorname{is\text{-}module}(d) \wedge \operatorname{is\text{-}submodule}(d, m)) \Rightarrow  \\ \quad \quad \operatorname{act}(\operatorname{quotient\text{-}module}(d, m))(\operatorname{mul}(\operatorname{scal}(\operatorname{quotient\text{-}module}(d, m)))(r, s), x) \\ \quad \quad =\; \operatorname{act}(\operatorname{quotient\text{-}module}(d, m))(r, \operatorname{act}(\operatorname{quotient\text{-}module}(d, m))(s, x)) \end{array}}}
\vnbprop{qmod-act-unital}{\vnbstmt{\begin{array}{@{}l@{}} \forall d, m.; \forall x \in \operatorname{vec}(\operatorname{quotient\text{-}module}(d, m)). \\ \quad (\operatorname{is\text{-}module}(d) \wedge \operatorname{is\text{-}submodule}(d, m)) \Rightarrow  \\ \quad \quad \operatorname{act}(\operatorname{quotient\text{-}module}(d, m))(\operatorname{one}(\operatorname{scal}(\operatorname{quotient\text{-}module}(d, m))), x) = x \end{array}}}
\vnbprop{quotient-module-is-module}{\vnbstmt{\begin{array}{@{}l@{}} \forall d, m. \\ \quad (\operatorname{is\text{-}module}(d) \wedge \operatorname{is\text{-}submodule}(d, m)) \Rightarrow \operatorname{is\text{-}module}(\operatorname{quotient\text{-}module}(d, m)) \end{array}}}
\vnbprop{ring-as-module-is-module}{\vnbstmt{\forall r.\; (\operatorname{is\text{-}ring}(r) \Rightarrow \operatorname{is\text{-}module}(\operatorname{ring\text{-}as\text{-}module}(r)))}}
\vnbprop{principal-ideal-is-submodule}{\vnbstmt{\begin{array}{@{}l@{}} \forall r.; \forall a \in \operatorname{carr}(r). \\ \quad \operatorname{is\text{-}commutative\text{-}ring}(r) \Rightarrow  \\ \quad \quad \operatorname{is\text{-}submodule}(\operatorname{ring\text{-}as\text{-}module}(r), \operatorname{principal\text{-}ideal}(r, a)) \end{array}}}
\vnbprop{zz-mod-is-module}{\vnbstmt{\forall a \in \mathbb{Z}.\; \operatorname{is\text{-}module}(\operatorname{zz\text{-}mod}(a))}}
\vnbprop{ram-ro-scal}{\vnbstmt{\forall r.\; (\operatorname{is\text{-}ring}(r) \Rightarrow (\operatorname{scal}(\operatorname{ring\text{-}as\text{-}module}(r)) \simeq r))}}
\vnbprop{ram-ro-vec}{\vnbstmt{\forall r.\; (\operatorname{is\text{-}ring}(r) \Rightarrow (\operatorname{vec}(\operatorname{ring\text{-}as\text{-}module}(r)) \simeq \operatorname{carr}(r)))}}
\vnbprop{ram-ro-vadd}{\vnbstmt{\forall r.\; (\operatorname{is\text{-}ring}(r) \Rightarrow (\operatorname{vadd}(\operatorname{ring\text{-}as\text{-}module}(r)) \simeq \operatorname{add}(r)))}}
\vnbprop{ram-ro-vzero}{\vnbstmt{\forall r.\; (\operatorname{is\text{-}ring}(r) \Rightarrow (\operatorname{vzero}(\operatorname{ring\text{-}as\text{-}module}(r)) \simeq \operatorname{zero}(r)))}}
\vnbprop{ram-ro-vneg}{\vnbstmt{\forall r.\; (\operatorname{is\text{-}ring}(r) \Rightarrow (\operatorname{vneg}(\operatorname{ring\text{-}as\text{-}module}(r)) \simeq \operatorname{neg}(r)))}}
\vnbprop{ram-ro-act}{\vnbstmt{\forall r.\; (\operatorname{is\text{-}ring}(r) \Rightarrow (\operatorname{act}(\operatorname{ring\text{-}as\text{-}module}(r)) \simeq \operatorname{mul}(r)))}}
\vnbprop{zz-mod-unfold}{\vnbstmt{\begin{array}{@{}l@{}} \forall a. \\ \quad (\operatorname{zz\text{-}mod}(a) \simeq \operatorname{quotient\text{-}module}(\operatorname{ring\text{-}as\text{-}module}(\operatorname{zz\text{-}ring}), \operatorname{principal\text{-}ideal}(\operatorname{zz\text{-}ring}, a))) \end{array}}}
\vnbprop{zz-mod-scal}{\vnbstmt{\forall a \in \operatorname{carr}(\operatorname{zz\text{-}ring}).\; (\operatorname{scal}(\operatorname{zz\text{-}mod}(a)) \simeq \operatorname{zz\text{-}ring})}}
\vnbprop{zz-mod-vec}{\vnbstmt{\begin{array}{@{}l@{}} \forall a \in \operatorname{carr}(\operatorname{zz\text{-}ring}). \\ \quad (\operatorname{vec}(\operatorname{zz\text{-}mod}(a)) \simeq \operatorname{quotient}(\operatorname{qmod\text{-}setoid}(\operatorname{ring\text{-}as\text{-}module}(\operatorname{zz\text{-}ring}), \operatorname{principal\text{-}ideal}(\operatorname{zz\text{-}ring}, a)))) \end{array}}}
\vnbprop{zz-mod-vzero}{\vnbstmt{\begin{array}{@{}l@{}} \forall a \in \operatorname{carr}(\operatorname{zz\text{-}ring}). \\ \quad (\operatorname{vzero}(\operatorname{zz\text{-}mod}(a)) \simeq \operatorname{class}(\operatorname{qmod\text{-}setoid}(\operatorname{ring\text{-}as\text{-}module}(\operatorname{zz\text{-}ring}), \operatorname{principal\text{-}ideal}(\operatorname{zz\text{-}ring}, a)), \operatorname{zero}(\operatorname{zz\text{-}ring}))) \end{array}}}
\vnbprop{zz-mod-vadd-computes}{\vnbstmt{\begin{array}{@{}l@{}} \forall a, x, y. \\ \quad ((a \in \operatorname{carr}(\operatorname{zz\text{-}ring})) \wedge ((x \in \operatorname{carr}(\operatorname{zz\text{-}ring})) \wedge (y \in \operatorname{carr}(\operatorname{zz\text{-}ring})))) \Rightarrow  \\ \quad \quad \operatorname{vadd}(\operatorname{zz\text{-}mod}(a))(\operatorname{class}(\operatorname{qmod\text{-}setoid}(\operatorname{ring\text{-}as\text{-}module}(\operatorname{zz\text{-}ring}), \operatorname{principal\text{-}ideal}(\operatorname{zz\text{-}ring}, a)), x), \\ \quad \quad \quad \operatorname{class}(\operatorname{qmod\text{-}setoid}(\operatorname{ring\text{-}as\text{-}module}(\operatorname{zz\text{-}ring}), \operatorname{principal\text{-}ideal}(\operatorname{zz\text{-}ring}, a)), y)) \\ \quad \quad =\; \operatorname{class}(\operatorname{qmod\text{-}setoid}(\operatorname{ring\text{-}as\text{-}module}(\operatorname{zz\text{-}ring}), \operatorname{principal\text{-}ideal}(\operatorname{zz\text{-}ring}, a)), \operatorname{add}(\operatorname{zz\text{-}ring})(x, y)) \end{array}}}
\vnbprop{zz-mod-vneg-computes}{\vnbstmt{\begin{array}{@{}l@{}} \forall a, x. \\ \quad ((a \in \operatorname{carr}(\operatorname{zz\text{-}ring})) \wedge (x \in \operatorname{carr}(\operatorname{zz\text{-}ring}))) \Rightarrow  \\ \quad \quad \operatorname{vneg}(\operatorname{zz\text{-}mod}(a))(\operatorname{class}(\operatorname{qmod\text{-}setoid}(\operatorname{ring\text{-}as\text{-}module}(\operatorname{zz\text{-}ring}), \operatorname{principal\text{-}ideal}(\operatorname{zz\text{-}ring}, a)), x)) \\ \quad \quad =\; \operatorname{class}(\operatorname{qmod\text{-}setoid}(\operatorname{ring\text{-}as\text{-}module}(\operatorname{zz\text{-}ring}), \operatorname{principal\text{-}ideal}(\operatorname{zz\text{-}ring}, a)), \operatorname{neg}(\operatorname{zz\text{-}ring})(x)) \end{array}}}
\vnbprop{zz-mod-act-computes}{\vnbstmt{\begin{array}{@{}l@{}} \forall a, r, x. \\ \quad ((a \in \operatorname{carr}(\operatorname{zz\text{-}ring})) \wedge ((r \in \operatorname{carr}(\operatorname{zz\text{-}ring})) \wedge (x \in \operatorname{carr}(\operatorname{zz\text{-}ring})))) \Rightarrow  \\ \quad \quad \operatorname{act}(\operatorname{zz\text{-}mod}(a))(r, \operatorname{class}(\operatorname{qmod\text{-}setoid}(\operatorname{ring\text{-}as\text{-}module}(\operatorname{zz\text{-}ring}), \operatorname{principal\text{-}ideal}(\operatorname{zz\text{-}ring}, a)), x)) \\ \quad \quad =\; \operatorname{class}(\operatorname{qmod\text{-}setoid}(\operatorname{ring\text{-}as\text{-}module}(\operatorname{zz\text{-}ring}), \operatorname{principal\text{-}ideal}(\operatorname{zz\text{-}ring}, a)), \operatorname{mul}(\operatorname{zz\text{-}ring})(r, x)) \end{array}}}
\vnbprop{zz-mod-class-eq-iff}{\vnbstmt{\begin{array}{@{}l@{}} \forall a, x, y. \\ \quad ((a \in \operatorname{carr}(\operatorname{zz\text{-}ring})) \wedge ((x \in \operatorname{carr}(\operatorname{zz\text{-}ring})) \wedge (y \in \operatorname{carr}(\operatorname{zz\text{-}ring})))) \Rightarrow  \\ \quad \quad \operatorname{class}(\operatorname{qmod\text{-}setoid}(\operatorname{ring\text{-}as\text{-}module}(\operatorname{zz\text{-}ring}), \operatorname{principal\text{-}ideal}(\operatorname{zz\text{-}ring}, a)), x) \\ \quad \quad =\; \operatorname{class}(\operatorname{qmod\text{-}setoid}(\operatorname{ring\text{-}as\text{-}module}(\operatorname{zz\text{-}ring}), \operatorname{principal\text{-}ideal}(\operatorname{zz\text{-}ring}, a)), y) \\ \quad \quad \Leftrightarrow\; \operatorname{add}(\operatorname{zz\text{-}ring})(x, \operatorname{neg}(\operatorname{zz\text{-}ring})(y)) \in \operatorname{principal\text{-}ideal}(\operatorname{zz\text{-}ring}, a) \end{array}}}
\vnbprop{zz-mod-class-in}{\vnbstmt{\begin{array}{@{}l@{}} \forall a, x. \\ \quad ((a \in \operatorname{carr}(\operatorname{zz\text{-}ring})) \wedge (x \in \operatorname{carr}(\operatorname{zz\text{-}ring}))) \Rightarrow  \\ \quad \quad \operatorname{class}(\operatorname{qmod\text{-}setoid}(\operatorname{ring\text{-}as\text{-}module}(\operatorname{zz\text{-}ring}), \operatorname{principal\text{-}ideal}(\operatorname{zz\text{-}ring}, a)), x) \in \operatorname{vec}(\operatorname{zz\text{-}mod}(a)) \end{array}}}

\vnbbookfile{theorem-library/la-z3-dsum}{theorem-library/la-z3-dsum}
\noindent{\small la-z3-dsum.scm -- the direct sum of finitely many modules over one scalar ring is a module (LA-Z3, 2026-10-24).  Vocabulary: structure-library/la-z3-vocab (DIRECT-SUM, DSUM-VEC, DSUM-VADD, DSUM-VZERO, DSUM-VNEG, DSUM-ACT).  Prefix lzs-.}\par
\vnbprop{dsum-vec-unfold}{\vnbstmt{\begin{array}{@{}l@{}} \forall n, f. \\ \quad (\operatorname{dsum\text{-}vec}(n, f) \simeq \{v \in (\operatorname{interval}(1, n) \to \operatorname{big\text{-}union}(k, \operatorname{interval}(1, n), \operatorname{vec}(f(k)))) : \forall j \in \operatorname{interval}(1, n).\; (v(j) \in \operatorname{vec}(f(j)))\}) \end{array}}}
\vnbprop{dsum-vec-in}{\vnbstmt{\begin{array}{@{}l@{}} \forall n, f.; \forall g \in (\operatorname{interval}(1, n) \to \operatorname{big\text{-}union}(k, \operatorname{interval}(1, n), \operatorname{vec}(f(k)))). \\ \quad \forall j \in \operatorname{interval}(1, n).\; (g(j) \in \operatorname{vec}(f(j))) \Rightarrow (g \in \operatorname{dsum\text{-}vec}(n, f)) \end{array}}}
\vnbprop{dsum-vec-out}{\vnbstmt{\begin{array}{@{}l@{}} \forall n, f.; \forall g \in \operatorname{dsum\text{-}vec}(n, f). \\ \quad g \in (\operatorname{interval}(1, n) \to \operatorname{big\text{-}union}(k, \operatorname{interval}(1, n), \operatorname{vec}(f(k)))) \wedge \\ \quad \forall j \in \operatorname{interval}(1, n).\; (g(j) \in \operatorname{vec}(f(j))) \end{array}}}
\vnbprop{dsum-in-union}{\vnbstmt{\begin{array}{@{}l@{}} \forall n, f, i, x. \\ \quad ((i \in \operatorname{interval}(1, n)) \wedge (x \in \operatorname{vec}(f(i)))) \Rightarrow  \\ \quad \quad x \in \operatorname{big\text{-}union}(k, \operatorname{interval}(1, n), \operatorname{vec}(f(k))) \end{array}}}
\vnbprop{dsum-union-set}{\vnbstmt{\begin{array}{@{}l@{}} \forall s, n, f. \\ \quad \operatorname{is\text{-}ring}(s) \wedge \\ \quad n \in \mathbb{N} \wedge \\ \quad \forall i \in \operatorname{interval}(1, n).\; \operatorname{is\text{-}module}(f(i)) \wedge \\ \quad \forall i \in \operatorname{interval}(1, n).\; (\operatorname{scal}(f(i)) = s) \Rightarrow \\ \quad \quad \operatorname{big\text{-}union}(k, \operatorname{interval}(1, n), \operatorname{vec}(f(k))) \in \mathbf{Set} \end{array}}}
\vnbprop{dsum-vec-set}{\vnbstmt{\begin{array}{@{}l@{}} \forall s, n, f. \\ \quad \operatorname{is\text{-}ring}(s) \wedge \\ \quad n \in \mathbb{N} \wedge \\ \quad \forall i \in \operatorname{interval}(1, n).\; \operatorname{is\text{-}module}(f(i)) \wedge \\ \quad \forall i \in \operatorname{interval}(1, n).\; (\operatorname{scal}(f(i)) = s) \Rightarrow \\ \quad \quad \operatorname{dsum\text{-}vec}(n, f) \in \mathbf{Set} \end{array}}}
\vnbprop{dsum-vzero-in0}{\vnbstmt{\begin{array}{@{}l@{}} \forall s, n, f. \\ \quad \operatorname{is\text{-}ring}(s) \wedge \\ \quad n \in \mathbb{N} \wedge \\ \quad \forall i \in \operatorname{interval}(1, n).\; \operatorname{is\text{-}module}(f(i)) \wedge \\ \quad \forall i \in \operatorname{interval}(1, n).\; (\operatorname{scal}(f(i)) = s) \Rightarrow \\ \quad \quad \operatorname{dsum\text{-}vzero}(n, f) \in \operatorname{dsum\text{-}vec}(n, f) \end{array}}}
\vnbprop{dsum-vadd-fun0}{\vnbstmt{\begin{array}{@{}l@{}} \forall s, n, f. \\ \quad \operatorname{is\text{-}ring}(s) \wedge \\ \quad n \in \mathbb{N} \wedge \\ \quad \forall i \in \operatorname{interval}(1, n).\; \operatorname{is\text{-}module}(f(i)) \wedge \\ \quad \forall i \in \operatorname{interval}(1, n).\; (\operatorname{scal}(f(i)) = s) \Rightarrow \\ \quad \quad \operatorname{dsum\text{-}vadd}(n, f) \in ((\operatorname{dsum\text{-}vec}(n, f) \times \operatorname{dsum\text{-}vec}(n, f)) \to \operatorname{dsum\text{-}vec}(n, f)) \end{array}}}
\vnbprop{dsum-vneg-fun0}{\vnbstmt{\begin{array}{@{}l@{}} \forall s, n, f. \\ \quad \operatorname{is\text{-}ring}(s) \wedge \\ \quad n \in \mathbb{N} \wedge \\ \quad \forall i \in \operatorname{interval}(1, n).\; \operatorname{is\text{-}module}(f(i)) \wedge \\ \quad \forall i \in \operatorname{interval}(1, n).\; (\operatorname{scal}(f(i)) = s) \Rightarrow \\ \quad \quad \operatorname{dsum\text{-}vneg}(n, f) \in (\operatorname{dsum\text{-}vec}(n, f) \to \operatorname{dsum\text{-}vec}(n, f)) \end{array}}}
\vnbprop{dsum-act-fun0}{\vnbstmt{\begin{array}{@{}l@{}} \forall s, n, f. \\ \quad \operatorname{is\text{-}ring}(s) \wedge \\ \quad n \in \mathbb{N} \wedge \\ \quad \forall i \in \operatorname{interval}(1, n).\; \operatorname{is\text{-}module}(f(i)) \wedge \\ \quad \forall i \in \operatorname{interval}(1, n).\; (\operatorname{scal}(f(i)) = s) \Rightarrow \\ \quad \quad \operatorname{dsum\text{-}act}(s, n, f) \in ((\operatorname{carr}(s) \times \operatorname{dsum\text{-}vec}(n, f)) \to \operatorname{dsum\text{-}vec}(n, f)) \end{array}}}
\vnbprop{dsum-ro-scal}{\vnbstmt{\begin{array}{@{}l@{}} \forall s, n, f. \\ \quad \operatorname{is\text{-}ring}(s) \wedge \\ \quad n \in \mathbb{N} \wedge \\ \quad \forall i \in \operatorname{interval}(1, n).\; \operatorname{is\text{-}module}(f(i)) \wedge \\ \quad \forall i \in \operatorname{interval}(1, n).\; (\operatorname{scal}(f(i)) = s) \Rightarrow \\ \quad \quad (\operatorname{scal}(\operatorname{direct\text{-}sum}(s, n, f)) \simeq s) \end{array}}}
\vnbprop{dsum-ro-vec}{\vnbstmt{\begin{array}{@{}l@{}} \forall s, n, f. \\ \quad \operatorname{is\text{-}ring}(s) \wedge \\ \quad n \in \mathbb{N} \wedge \\ \quad \forall i \in \operatorname{interval}(1, n).\; \operatorname{is\text{-}module}(f(i)) \wedge \\ \quad \forall i \in \operatorname{interval}(1, n).\; (\operatorname{scal}(f(i)) = s) \Rightarrow \\ \quad \quad (\operatorname{vec}(\operatorname{direct\text{-}sum}(s, n, f)) \simeq \operatorname{dsum\text{-}vec}(n, f)) \end{array}}}
\vnbprop{dsum-ro-vadd}{\vnbstmt{\begin{array}{@{}l@{}} \forall s, n, f. \\ \quad \operatorname{is\text{-}ring}(s) \wedge \\ \quad n \in \mathbb{N} \wedge \\ \quad \forall i \in \operatorname{interval}(1, n).\; \operatorname{is\text{-}module}(f(i)) \wedge \\ \quad \forall i \in \operatorname{interval}(1, n).\; (\operatorname{scal}(f(i)) = s) \Rightarrow \\ \quad \quad (\operatorname{vadd}(\operatorname{direct\text{-}sum}(s, n, f)) \simeq \operatorname{dsum\text{-}vadd}(n, f)) \end{array}}}
\vnbprop{dsum-ro-vzero}{\vnbstmt{\begin{array}{@{}l@{}} \forall s, n, f. \\ \quad \operatorname{is\text{-}ring}(s) \wedge \\ \quad n \in \mathbb{N} \wedge \\ \quad \forall i \in \operatorname{interval}(1, n).\; \operatorname{is\text{-}module}(f(i)) \wedge \\ \quad \forall i \in \operatorname{interval}(1, n).\; (\operatorname{scal}(f(i)) = s) \Rightarrow \\ \quad \quad (\operatorname{vzero}(\operatorname{direct\text{-}sum}(s, n, f)) \simeq \operatorname{dsum\text{-}vzero}(n, f)) \end{array}}}
\vnbprop{dsum-ro-vneg}{\vnbstmt{\begin{array}{@{}l@{}} \forall s, n, f. \\ \quad \operatorname{is\text{-}ring}(s) \wedge \\ \quad n \in \mathbb{N} \wedge \\ \quad \forall i \in \operatorname{interval}(1, n).\; \operatorname{is\text{-}module}(f(i)) \wedge \\ \quad \forall i \in \operatorname{interval}(1, n).\; (\operatorname{scal}(f(i)) = s) \Rightarrow \\ \quad \quad (\operatorname{vneg}(\operatorname{direct\text{-}sum}(s, n, f)) \simeq \operatorname{dsum\text{-}vneg}(n, f)) \end{array}}}
\vnbprop{dsum-ro-act}{\vnbstmt{\begin{array}{@{}l@{}} \forall s, n, f. \\ \quad \operatorname{is\text{-}ring}(s) \wedge \\ \quad n \in \mathbb{N} \wedge \\ \quad \forall i \in \operatorname{interval}(1, n).\; \operatorname{is\text{-}module}(f(i)) \wedge \\ \quad \forall i \in \operatorname{interval}(1, n).\; (\operatorname{scal}(f(i)) = s) \Rightarrow \\ \quad \quad (\operatorname{act}(\operatorname{direct\text{-}sum}(s, n, f)) \simeq \operatorname{dsum\text{-}act}(s, n, f)) \end{array}}}
\vnbprop{dsum-vadd-closed}{\vnbstmt{\begin{array}{@{}l@{}} \forall s, n, f, a, b. \\ \quad \operatorname{is\text{-}ring}(s) \wedge \\ \quad n \in \mathbb{N} \wedge \\ \quad \forall i \in \operatorname{interval}(1, n).\; \operatorname{is\text{-}module}(f(i)) \wedge \\ \quad \forall i \in \operatorname{interval}(1, n).\; (\operatorname{scal}(f(i)) = s) \wedge \\ \quad a \in \operatorname{dsum\text{-}vec}(n, f) \wedge \\ \quad b \in \operatorname{dsum\text{-}vec}(n, f) \Rightarrow \\ \quad \quad \operatorname{vadd}(\operatorname{direct\text{-}sum}(s, n, f))(a, b) \in \operatorname{dsum\text{-}vec}(n, f) \end{array}}}
\vnbprop{dsum-vneg-closed}{\vnbstmt{\begin{array}{@{}l@{}} \forall s, n, f.; \forall a \in \operatorname{dsum\text{-}vec}(n, f). \\ \quad \operatorname{is\text{-}ring}(s) \wedge \\ \quad n \in \mathbb{N} \wedge \\ \quad \forall i \in \operatorname{interval}(1, n).\; \operatorname{is\text{-}module}(f(i)) \wedge \\ \quad \forall i \in \operatorname{interval}(1, n).\; (\operatorname{scal}(f(i)) = s) \Rightarrow \\ \quad \quad \operatorname{vneg}(\operatorname{direct\text{-}sum}(s, n, f))(a) \in \operatorname{dsum\text{-}vec}(n, f) \end{array}}}
\vnbprop{dsum-act-closed}{\vnbstmt{\begin{array}{@{}l@{}} \forall s, n, f, r, a. \\ \quad \operatorname{is\text{-}ring}(s) \wedge \\ \quad n \in \mathbb{N} \wedge \\ \quad \forall i \in \operatorname{interval}(1, n).\; \operatorname{is\text{-}module}(f(i)) \wedge \\ \quad \forall i \in \operatorname{interval}(1, n).\; (\operatorname{scal}(f(i)) = s) \wedge \\ \quad r \in \operatorname{carr}(s) \wedge \\ \quad a \in \operatorname{dsum\text{-}vec}(n, f) \Rightarrow \\ \quad \quad \operatorname{act}(\operatorname{direct\text{-}sum}(s, n, f))(r, a) \in \operatorname{dsum\text{-}vec}(n, f) \end{array}}}
\vnbprop{dsum-vzero-closed}{\vnbstmt{\begin{array}{@{}l@{}} \forall s, n, f. \\ \quad \operatorname{is\text{-}ring}(s) \wedge \\ \quad n \in \mathbb{N} \wedge \\ \quad \forall i \in \operatorname{interval}(1, n).\; \operatorname{is\text{-}module}(f(i)) \wedge \\ \quad \forall i \in \operatorname{interval}(1, n).\; (\operatorname{scal}(f(i)) = s) \Rightarrow \\ \quad \quad \operatorname{vzero}(\operatorname{direct\text{-}sum}(s, n, f)) \in \operatorname{dsum\text{-}vec}(n, f) \end{array}}}
\vnbprop{dsum-vadd-at}{\vnbstmt{\begin{array}{@{}l@{}} \forall s, n, f, a, b, i. \\ \quad \operatorname{is\text{-}ring}(s) \wedge \\ \quad n \in \mathbb{N} \wedge \\ \quad \forall i \in \operatorname{interval}(1, n).\; \operatorname{is\text{-}module}(f(i)) \wedge \\ \quad \forall i \in \operatorname{interval}(1, n).\; (\operatorname{scal}(f(i)) = s) \wedge \\ \quad a \in \operatorname{dsum\text{-}vec}(n, f) \wedge \\ \quad b \in \operatorname{dsum\text{-}vec}(n, f) \wedge \\ \quad i \in \operatorname{interval}(1, n) \Rightarrow \\ \quad \quad \operatorname{vadd}(\operatorname{direct\text{-}sum}(s, n, f))(a, b)(i) = \operatorname{vadd}(f(i))(a(i), b(i)) \end{array}}}
\vnbprop{dsum-vneg-at}{\vnbstmt{\begin{array}{@{}l@{}} \forall s, n, f, a, i. \\ \quad \operatorname{is\text{-}ring}(s) \wedge \\ \quad n \in \mathbb{N} \wedge \\ \quad \forall i \in \operatorname{interval}(1, n).\; \operatorname{is\text{-}module}(f(i)) \wedge \\ \quad \forall i \in \operatorname{interval}(1, n).\; (\operatorname{scal}(f(i)) = s) \wedge \\ \quad a \in \operatorname{dsum\text{-}vec}(n, f) \wedge \\ \quad i \in \operatorname{interval}(1, n) \Rightarrow \\ \quad \quad \operatorname{vneg}(\operatorname{direct\text{-}sum}(s, n, f))(a)(i) = \operatorname{vneg}(f(i))(a(i)) \end{array}}}
\vnbprop{dsum-act-at}{\vnbstmt{\begin{array}{@{}l@{}} \forall s, n, f, r, a, i. \\ \quad \operatorname{is\text{-}ring}(s) \wedge \\ \quad n \in \mathbb{N} \wedge \\ \quad \forall i \in \operatorname{interval}(1, n).\; \operatorname{is\text{-}module}(f(i)) \wedge \\ \quad \forall i \in \operatorname{interval}(1, n).\; (\operatorname{scal}(f(i)) = s) \wedge \\ \quad r \in \operatorname{carr}(s) \wedge \\ \quad a \in \operatorname{dsum\text{-}vec}(n, f) \wedge \\ \quad i \in \operatorname{interval}(1, n) \Rightarrow \\ \quad \quad \operatorname{act}(\operatorname{direct\text{-}sum}(s, n, f))(r, a)(i) = \operatorname{act}(f(i))(r, a(i)) \end{array}}}
\vnbprop{dsum-vzero-at}{\vnbstmt{\begin{array}{@{}l@{}} \forall s, n, f.; \forall i \in \operatorname{interval}(1, n). \\ \quad \operatorname{is\text{-}ring}(s) \wedge \\ \quad n \in \mathbb{N} \wedge \\ \quad \forall i \in \operatorname{interval}(1, n).\; \operatorname{is\text{-}module}(f(i)) \wedge \\ \quad \forall i \in \operatorname{interval}(1, n).\; (\operatorname{scal}(f(i)) = s) \Rightarrow \\ \quad \quad \operatorname{vzero}(\operatorname{direct\text{-}sum}(s, n, f))(i) = \operatorname{vzero}(f(i)) \end{array}}}
\vnbprop{dsum-length}{\vnbstmt{\begin{array}{@{}l@{}} \forall s, n, f. \\ \quad \operatorname{is\text{-}ring}(s) \wedge \\ \quad n \in \mathbb{N} \wedge \\ \quad \forall i \in \operatorname{interval}(1, n).\; \operatorname{is\text{-}module}(f(i)) \wedge \\ \quad \forall i \in \operatorname{interval}(1, n).\; (\operatorname{scal}(f(i)) = s) \Rightarrow \\ \quad \quad \operatorname{length}(\operatorname{direct\text{-}sum}(s, n, f)) = 6 \end{array}}}
\vnbprop{dsum-scal-ring}{\vnbstmt{\begin{array}{@{}l@{}} \forall s, n, f. \\ \quad \operatorname{is\text{-}ring}(s) \wedge \\ \quad n \in \mathbb{N} \wedge \\ \quad \forall i \in \operatorname{interval}(1, n).\; \operatorname{is\text{-}module}(f(i)) \wedge \\ \quad \forall i \in \operatorname{interval}(1, n).\; (\operatorname{scal}(f(i)) = s) \Rightarrow \\ \quad \quad \operatorname{is\text{-}ring}(\operatorname{scal}(\operatorname{direct\text{-}sum}(s, n, f))) \end{array}}}
\vnbprop{dsum-vec-set2}{\vnbstmt{\begin{array}{@{}l@{}} \forall s, n, f. \\ \quad \operatorname{is\text{-}ring}(s) \wedge \\ \quad n \in \mathbb{N} \wedge \\ \quad \forall i \in \operatorname{interval}(1, n).\; \operatorname{is\text{-}module}(f(i)) \wedge \\ \quad \forall i \in \operatorname{interval}(1, n).\; (\operatorname{scal}(f(i)) = s) \Rightarrow \\ \quad \quad \operatorname{vec}(\operatorname{direct\text{-}sum}(s, n, f)) \in \mathbf{Set} \end{array}}}
\vnbprop{dsum-vzero-in}{\vnbstmt{\begin{array}{@{}l@{}} \forall s, n, f. \\ \quad \operatorname{is\text{-}ring}(s) \wedge \\ \quad n \in \mathbb{N} \wedge \\ \quad \forall i \in \operatorname{interval}(1, n).\; \operatorname{is\text{-}module}(f(i)) \wedge \\ \quad \forall i \in \operatorname{interval}(1, n).\; (\operatorname{scal}(f(i)) = s) \Rightarrow \\ \quad \quad \operatorname{vzero}(\operatorname{direct\text{-}sum}(s, n, f)) \in \operatorname{vec}(\operatorname{direct\text{-}sum}(s, n, f)) \end{array}}}
\vnbprop{dsum-vadd-fun}{\vnbstmt{\begin{array}{@{}l@{}} \forall s, n, f. \\ \quad \operatorname{is\text{-}ring}(s) \wedge \\ \quad n \in \mathbb{N} \wedge \\ \quad \forall i \in \operatorname{interval}(1, n).\; \operatorname{is\text{-}module}(f(i)) \wedge \\ \quad \forall i \in \operatorname{interval}(1, n).\; (\operatorname{scal}(f(i)) = s) \Rightarrow \\ \quad \quad \operatorname{vadd}(\operatorname{direct\text{-}sum}(s, n, f)) \\ \quad \quad \in\; ((\operatorname{vec}(\operatorname{direct\text{-}sum}(s, n, f)) \times \operatorname{vec}(\operatorname{direct\text{-}sum}(s, n, f))) \to \operatorname{vec}(\operatorname{direct\text{-}sum}(s, n, f))) \end{array}}}
\vnbprop{dsum-vneg-fun}{\vnbstmt{\begin{array}{@{}l@{}} \forall s, n, f. \\ \quad \operatorname{is\text{-}ring}(s) \wedge \\ \quad n \in \mathbb{N} \wedge \\ \quad \forall i \in \operatorname{interval}(1, n).\; \operatorname{is\text{-}module}(f(i)) \wedge \\ \quad \forall i \in \operatorname{interval}(1, n).\; (\operatorname{scal}(f(i)) = s) \Rightarrow \\ \quad \quad \operatorname{vneg}(\operatorname{direct\text{-}sum}(s, n, f)) \in (\operatorname{vec}(\operatorname{direct\text{-}sum}(s, n, f)) \to \operatorname{vec}(\operatorname{direct\text{-}sum}(s, n, f))) \end{array}}}
\vnbprop{dsum-act-fun}{\vnbstmt{\begin{array}{@{}l@{}} \forall s, n, f. \\ \quad \operatorname{is\text{-}ring}(s) \wedge \\ \quad n \in \mathbb{N} \wedge \\ \quad \forall i \in \operatorname{interval}(1, n).\; \operatorname{is\text{-}module}(f(i)) \wedge \\ \quad \forall i \in \operatorname{interval}(1, n).\; (\operatorname{scal}(f(i)) = s) \Rightarrow \\ \quad \quad \operatorname{act}(\operatorname{direct\text{-}sum}(s, n, f)) \\ \quad \quad \in\; ((\operatorname{carr}(\operatorname{scal}(\operatorname{direct\text{-}sum}(s, n, f))) \times \operatorname{vec}(\operatorname{direct\text{-}sum}(s, n, f))) \to \operatorname{vec}(\operatorname{direct\text{-}sum}(s, n, f))) \end{array}}}
\vnbprop{dsum-assoc}{\vnbstmt{\begin{array}{@{}l@{}} \forall s, n, f. \\ \quad \operatorname{is\text{-}ring}(s) \wedge \\ \quad n \in \mathbb{N} \wedge \\ \quad \forall i \in \operatorname{interval}(1, n).\; \operatorname{is\text{-}module}(f(i)) \wedge \\ \quad \forall i \in \operatorname{interval}(1, n).\; (\operatorname{scal}(f(i)) = s) \Rightarrow \\ \quad \quad \operatorname{is\text{-}associative}(\operatorname{vadd}(\operatorname{direct\text{-}sum}(s, n, f)), \operatorname{vec}(\operatorname{direct\text{-}sum}(s, n, f))) \end{array}}}
\vnbprop{dsum-comm}{\vnbstmt{\begin{array}{@{}l@{}} \forall s, n, f. \\ \quad \operatorname{is\text{-}ring}(s) \wedge \\ \quad n \in \mathbb{N} \wedge \\ \quad \forall i \in \operatorname{interval}(1, n).\; \operatorname{is\text{-}module}(f(i)) \wedge \\ \quad \forall i \in \operatorname{interval}(1, n).\; (\operatorname{scal}(f(i)) = s) \Rightarrow \\ \quad \quad \operatorname{is\text{-}commutative}(\operatorname{vadd}(\operatorname{direct\text{-}sum}(s, n, f)), \operatorname{vec}(\operatorname{direct\text{-}sum}(s, n, f))) \end{array}}}
\vnbprop{dsum-identity}{\vnbstmt{\begin{array}{@{}l@{}} \forall s, n, f. \\ \quad \operatorname{is\text{-}ring}(s) \wedge \\ \quad n \in \mathbb{N} \wedge \\ \quad \forall i \in \operatorname{interval}(1, n).\; \operatorname{is\text{-}module}(f(i)) \wedge \\ \quad \forall i \in \operatorname{interval}(1, n).\; (\operatorname{scal}(f(i)) = s) \Rightarrow \\ \quad \quad \operatorname{is\text{-}identity}(\operatorname{vadd}(\operatorname{direct\text{-}sum}(s, n, f)), \operatorname{vzero}(\operatorname{direct\text{-}sum}(s, n, f)), \operatorname{vec}(\operatorname{direct\text{-}sum}(s, n, f))) \end{array}}}
\vnbprop{dsum-inverses}{\vnbstmt{\begin{array}{@{}l@{}} \forall s, n, f. \\ \quad \operatorname{is\text{-}ring}(s) \wedge \\ \quad n \in \mathbb{N} \wedge \\ \quad \forall i \in \operatorname{interval}(1, n).\; \operatorname{is\text{-}module}(f(i)) \wedge \\ \quad \forall i \in \operatorname{interval}(1, n).\; (\operatorname{scal}(f(i)) = s) \Rightarrow \\ \quad \quad \operatorname{has\text{-}inverses}(\operatorname{vadd}(\operatorname{direct\text{-}sum}(s, n, f)), \operatorname{vzero}(\operatorname{direct\text{-}sum}(s, n, f)), \\ \quad \quad \quad \operatorname{vneg}(\operatorname{direct\text{-}sum}(s, n, f)), \operatorname{vec}(\operatorname{direct\text{-}sum}(s, n, f))) \end{array}}}
\vnbprop{dsum-act-distrib-vec}{\vnbstmt{\begin{array}{@{}l@{}} \forall s, n, f.; \forall r \in \operatorname{carr}(\operatorname{scal}(\operatorname{direct\text{-}sum}(s, n, f))), x, y \in \operatorname{vec}(\operatorname{direct\text{-}sum}(s, n, f)). \\ \quad \operatorname{is\text{-}ring}(s) \wedge \\ \quad n \in \mathbb{N} \wedge \\ \quad \forall i \in \operatorname{interval}(1, n).\; \operatorname{is\text{-}module}(f(i)) \wedge \\ \quad \forall i \in \operatorname{interval}(1, n).\; (\operatorname{scal}(f(i)) = s) \Rightarrow \\ \quad \quad \operatorname{act}(\operatorname{direct\text{-}sum}(s, n, f))(r, \operatorname{vadd}(\operatorname{direct\text{-}sum}(s, n, f))(x, y)) \\ \quad \quad =\; \operatorname{vadd}(\operatorname{direct\text{-}sum}(s, n, f))(\operatorname{act}(\operatorname{direct\text{-}sum}(s, n, f))(r, x), \operatorname{act}(\operatorname{direct\text{-}sum}(s, n, f))(r, y)) \end{array}}}
\vnbprop{dsum-act-distrib-scalar}{\vnbstmt{\begin{array}{@{}l@{}} \forall s, n, f.; \forall r, a \in \operatorname{carr}(\operatorname{scal}(\operatorname{direct\text{-}sum}(s, n, f))), x \in \operatorname{vec}(\operatorname{direct\text{-}sum}(s, n, f)). \\ \quad \operatorname{is\text{-}ring}(s) \wedge \\ \quad n \in \mathbb{N} \wedge \\ \quad \forall i \in \operatorname{interval}(1, n).\; \operatorname{is\text{-}module}(f(i)) \wedge \\ \quad \forall i \in \operatorname{interval}(1, n).\; (\operatorname{scal}(f(i)) = s) \Rightarrow \\ \quad \quad \operatorname{act}(\operatorname{direct\text{-}sum}(s, n, f))(\operatorname{add}(\operatorname{scal}(\operatorname{direct\text{-}sum}(s, n, f)))(r, a), x) \\ \quad \quad =\; \operatorname{vadd}(\operatorname{direct\text{-}sum}(s, n, f))(\operatorname{act}(\operatorname{direct\text{-}sum}(s, n, f))(r, x), \operatorname{act}(\operatorname{direct\text{-}sum}(s, n, f))(a, x)) \end{array}}}
\vnbprop{dsum-act-mul-compat}{\vnbstmt{\begin{array}{@{}l@{}} \forall s, n, f.; \forall r, a \in \operatorname{carr}(\operatorname{scal}(\operatorname{direct\text{-}sum}(s, n, f))), x \in \operatorname{vec}(\operatorname{direct\text{-}sum}(s, n, f)). \\ \quad \operatorname{is\text{-}ring}(s) \wedge \\ \quad n \in \mathbb{N} \wedge \\ \quad \forall i \in \operatorname{interval}(1, n).\; \operatorname{is\text{-}module}(f(i)) \wedge \\ \quad \forall i \in \operatorname{interval}(1, n).\; (\operatorname{scal}(f(i)) = s) \Rightarrow \\ \quad \quad \operatorname{act}(\operatorname{direct\text{-}sum}(s, n, f))(\operatorname{mul}(\operatorname{scal}(\operatorname{direct\text{-}sum}(s, n, f)))(r, a), x) \\ \quad \quad =\; \operatorname{act}(\operatorname{direct\text{-}sum}(s, n, f))(r, \operatorname{act}(\operatorname{direct\text{-}sum}(s, n, f))(a, x)) \end{array}}}
\vnbprop{dsum-act-unital}{\vnbstmt{\begin{array}{@{}l@{}} \forall s, n, f.; \forall x \in \operatorname{vec}(\operatorname{direct\text{-}sum}(s, n, f)). \\ \quad \operatorname{is\text{-}ring}(s) \wedge \\ \quad n \in \mathbb{N} \wedge \\ \quad \forall i \in \operatorname{interval}(1, n).\; \operatorname{is\text{-}module}(f(i)) \wedge \\ \quad \forall i \in \operatorname{interval}(1, n).\; (\operatorname{scal}(f(i)) = s) \Rightarrow \\ \quad \quad \operatorname{act}(\operatorname{direct\text{-}sum}(s, n, f))(\operatorname{one}(\operatorname{scal}(\operatorname{direct\text{-}sum}(s, n, f))), x) = x \end{array}}}
\vnbprop{direct-sum-is-module}{\vnbstmt{\begin{array}{@{}l@{}} \forall s, n, f. \\ \quad \operatorname{is\text{-}ring}(s) \wedge \\ \quad n \in \mathbb{N} \wedge \\ \quad \forall i \in \operatorname{interval}(1, n).\; \operatorname{is\text{-}module}(f(i)) \wedge \\ \quad \forall i \in \operatorname{interval}(1, n).\; (\operatorname{scal}(f(i)) = s) \Rightarrow \\ \quad \quad \operatorname{is\text{-}module}(\operatorname{direct\text{-}sum}(s, n, f)) \end{array}}}
\vnbprop{dsum-member-intro}{\vnbstmt{\begin{array}{@{}l@{}} \forall s, n, f.; \forall g \in \mathrm{Fun}(\operatorname{interval}(1, n)). \\ \quad \operatorname{is\text{-}ring}(s) \wedge \\ \quad n \in \mathbb{N} \wedge \\ \quad \forall i \in \operatorname{interval}(1, n).\; \operatorname{is\text{-}module}(f(i)) \wedge \\ \quad \forall i \in \operatorname{interval}(1, n).\; (\operatorname{scal}(f(i)) = s) \wedge \\ \quad \forall j \in \operatorname{interval}(1, n).\; (g(j) \in \operatorname{vec}(f(j))) \Rightarrow \\ \quad \quad g \in \operatorname{vec}(\operatorname{direct\text{-}sum}(s, n, f)) \end{array}}}
\vnbprop{dsum-member-elim}{\vnbstmt{\begin{array}{@{}l@{}} \forall s, n, f, g, i. \\ \quad \operatorname{is\text{-}ring}(s) \wedge \\ \quad n \in \mathbb{N} \wedge \\ \quad \forall i \in \operatorname{interval}(1, n).\; \operatorname{is\text{-}module}(f(i)) \wedge \\ \quad \forall i \in \operatorname{interval}(1, n).\; (\operatorname{scal}(f(i)) = s) \wedge \\ \quad g \in \operatorname{vec}(\operatorname{direct\text{-}sum}(s, n, f)) \wedge \\ \quad i \in \operatorname{interval}(1, n) \Rightarrow \\ \quad \quad g(i) \in \operatorname{vec}(f(i)) \end{array}}}
\vnbprop{dsum-ext}{\vnbstmt{\begin{array}{@{}l@{}} \forall s, n, f, g, h. \\ \quad \operatorname{is\text{-}ring}(s) \wedge \\ \quad n \in \mathbb{N} \wedge \\ \quad \forall i \in \operatorname{interval}(1, n).\; \operatorname{is\text{-}module}(f(i)) \wedge \\ \quad \forall i \in \operatorname{interval}(1, n).\; (\operatorname{scal}(f(i)) = s) \wedge \\ \quad g \in \operatorname{vec}(\operatorname{direct\text{-}sum}(s, n, f)) \wedge \\ \quad h \in \operatorname{vec}(\operatorname{direct\text{-}sum}(s, n, f)) \wedge \\ \quad \forall j \in \operatorname{interval}(1, n).\; (g(j) = h(j)) \Rightarrow \\ \quad \quad g = h \end{array}}}

\vnbbookfile{theorem-library/la-z3-struct}{theorem-library/la-z3-struct}
\noindent{\small la-z3-struct.scm -- the first isomorphism theorem for modules, and linear algebra 3.49 (LA-Z3, 2026-10-24).  Vocabulary: structure-library/la-z3-vocab.  Prefix lzt-.}\par
\vnbprop{first-isomorphism-module}{\vnbstmt{\begin{array}{@{}l@{}} \forall m, n, i, a. \\ \quad \operatorname{is\text{-}hom\text{-}module}(m, n, i) \wedge \\ \quad \operatorname{is\text{-}submodule}(m, a) \wedge \\ \quad \forall y \in \operatorname{vec}(n).\; \exists x \in \operatorname{vec}(m).\; (i(x) = y) \wedge \\ \quad \forall x \in \operatorname{vec}(m).\; ((i(x) = \operatorname{vzero}(n)) \Leftrightarrow (x \in a)) \Rightarrow \\ \quad \quad \operatorname{is\text{-}iso\text{-}module}(\operatorname{quotient\text{-}module}(m, a), n, \operatorname{descend}(\operatorname{qmod\text{-}setoid}(m, a), i)) \end{array}}}
\vnbprop{class-of-member}{\vnbstmt{\begin{array}{@{}l@{}} \forall s.; \forall q \in \operatorname{quotient}(s), a \in q. \\ \quad \operatorname{is\text{-}setoid}(s) \Rightarrow ((a \in \operatorname{pts}(s)) \wedge (\operatorname{class}(s, a) = q)) \end{array}}}
\vnbprop{quotient-member-choice}{\vnbstmt{\forall s.; \forall q \in \operatorname{quotient}(s).\; (\operatorname{is\text{-}setoid}(s) \Rightarrow (\varepsilon(q) \in q))}}
\vnbprop{lzt-lc-bijection}{\vnbstmt{\begin{array}{@{}l@{}} \forall d, n, s. \\ \quad (\operatorname{is\text{-}module}(d) \wedge ((n \in \mathbb{N}) \wedge ((s \in \operatorname{mat}(n, 1, \operatorname{vec}(d))) \wedge (\operatorname{generates}(d, n, s) \wedge \operatorname{rel\text{-}free}(d, n, s))))) \Rightarrow  \\ \quad \quad (\lambda c \in \operatorname{mat}(1, n, \operatorname{carr}(\operatorname{scal}(d))).\; \operatorname{lincomb}(d, n, c, s)) \\ \quad \quad \in\; \mathrm{Bij}(\operatorname{mat}(1, n, \operatorname{carr}(\operatorname{scal}(d))), \operatorname{vec}(d)) \end{array}}}
\vnbprop{coord-lincomb}{\vnbstmt{\begin{array}{@{}l@{}} \forall d, n, s, x. \\ \quad \operatorname{is\text{-}module}(d) \wedge \\ \quad n \in \mathbb{N} \wedge \\ \quad s \in \operatorname{mat}(n, 1, \operatorname{vec}(d)) \wedge \\ \quad \operatorname{generates}(d, n, s) \wedge \\ \quad \operatorname{rel\text{-}free}(d, n, s) \wedge \\ \quad x \in \operatorname{vec}(d) \Rightarrow \\ \quad \quad (\lambda c \in \operatorname{mat}(1, n, \operatorname{carr}(\operatorname{scal}(d))).\; \operatorname{lincomb}(d, n, c, s))^{-1}(x) \\ \quad \quad \in\; \operatorname{mat}(1, n, \operatorname{carr}(\operatorname{scal}(d))) \wedge \\ \quad \quad \operatorname{lincomb}(d, n, \\ \quad \quad \quad (\lambda c \in \operatorname{mat}(1, n, \operatorname{carr}(\operatorname{scal}(d))).\; \operatorname{lincomb}(d, n, c, s))^{-1}(x), \\ \quad \quad \quad s) \\ \quad \quad =\; x \end{array}}}
\vnbprop{lincomb-coord}{\vnbstmt{\begin{array}{@{}l@{}} \forall d, n, s, c. \\ \quad \operatorname{is\text{-}module}(d) \wedge \\ \quad n \in \mathbb{N} \wedge \\ \quad s \in \operatorname{mat}(n, 1, \operatorname{vec}(d)) \wedge \\ \quad \operatorname{generates}(d, n, s) \wedge \\ \quad \operatorname{rel\text{-}free}(d, n, s) \wedge \\ \quad c \in \operatorname{mat}(1, n, \operatorname{carr}(\operatorname{scal}(d))) \Rightarrow \\ \quad \quad (\lambda \mathit{lcc} \in \operatorname{mat}(1, n, \operatorname{carr}(\operatorname{scal}(d))).\; \operatorname{lincomb}(d, n, \mathit{lcc}, s))^{-1}(\operatorname{lincomb}(d, n, c, s)) \\ \quad \quad =\; c \end{array}}}
\vnbprop{coord-vadd}{\vnbstmt{\begin{array}{@{}l@{}} \forall d, n, s, x, y, j. \\ \quad \operatorname{is\text{-}module}(d) \wedge \\ \quad n \in \mathbb{N} \wedge \\ \quad s \in \operatorname{mat}(n, 1, \operatorname{vec}(d)) \wedge \\ \quad \operatorname{generates}(d, n, s) \wedge \\ \quad \operatorname{rel\text{-}free}(d, n, s) \wedge \\ \quad x \in \operatorname{vec}(d) \wedge \\ \quad y \in \operatorname{vec}(d) \wedge \\ \quad j \in \operatorname{interval}(1, n) \Rightarrow \\ \quad \quad \operatorname{entry}((\lambda c \in \operatorname{mat}(1, n, \operatorname{carr}(\operatorname{scal}(d))).\; \operatorname{lincomb}(d, n, c, s))^{-1}(\operatorname{vadd}(d)(x, y)), \\ \quad \quad \quad 1, j) \\ \quad \quad =\; \operatorname{add}(\operatorname{scal}(d))(\operatorname{entry}((\lambda c \in \operatorname{mat}(1, n, \operatorname{carr}(\operatorname{scal}(d))).\; \operatorname{lincomb}(d, n, c, s))^{-1}(x), \\ \quad \quad \quad \quad 1, j), \\ \quad \quad \quad \operatorname{entry}((\lambda c \in \operatorname{mat}(1, n, \operatorname{carr}(\operatorname{scal}(d))).\; \operatorname{lincomb}(d, n, c, s))^{-1}(y), \\ \quad \quad \quad \quad 1, j)) \end{array}}}
\vnbprop{coord-act}{\vnbstmt{\begin{array}{@{}l@{}} \forall d, n, s, r, x, j. \\ \quad \operatorname{is\text{-}module}(d) \wedge \\ \quad n \in \mathbb{N} \wedge \\ \quad s \in \operatorname{mat}(n, 1, \operatorname{vec}(d)) \wedge \\ \quad \operatorname{generates}(d, n, s) \wedge \\ \quad \operatorname{rel\text{-}free}(d, n, s) \wedge \\ \quad r \in \operatorname{carr}(\operatorname{scal}(d)) \wedge \\ \quad x \in \operatorname{vec}(d) \wedge \\ \quad j \in \operatorname{interval}(1, n) \Rightarrow \\ \quad \quad \operatorname{entry}((\lambda c \in \operatorname{mat}(1, n, \operatorname{carr}(\operatorname{scal}(d))).\; \operatorname{lincomb}(d, n, c, s))^{-1}(\operatorname{act}(d)(r, x)), \\ \quad \quad \quad 1, j) \\ \quad \quad =\; \operatorname{mul}(\operatorname{scal}(d))(r, \\ \quad \quad \quad \operatorname{entry}((\lambda c \in \operatorname{mat}(1, n, \operatorname{carr}(\operatorname{scal}(d))).\; \operatorname{lincomb}(d, n, c, s))^{-1}(x), \\ \quad \quad \quad \quad 1, j)) \end{array}}}
\vnbprop{coord-vzero}{\vnbstmt{\begin{array}{@{}l@{}} \forall d, n, s, j. \\ \quad \operatorname{is\text{-}module}(d) \wedge \\ \quad n \in \mathbb{N} \wedge \\ \quad s \in \operatorname{mat}(n, 1, \operatorname{vec}(d)) \wedge \\ \quad \operatorname{generates}(d, n, s) \wedge \\ \quad \operatorname{rel\text{-}free}(d, n, s) \wedge \\ \quad j \in \operatorname{interval}(1, n) \Rightarrow \\ \quad \quad \operatorname{entry}((\lambda c \in \operatorname{mat}(1, n, \operatorname{carr}(\operatorname{scal}(d))).\; \operatorname{lincomb}(d, n, c, s))^{-1}(\operatorname{vzero}(d)), \\ \quad \quad \quad 1, j) \\ \quad \quad =\; \operatorname{zero}(\operatorname{scal}(d)) \end{array}}}
\vnbprop{coord-vneg}{\vnbstmt{\begin{array}{@{}l@{}} \forall d, n, s, x, j. \\ \quad \operatorname{is\text{-}module}(d) \wedge \\ \quad n \in \mathbb{N} \wedge \\ \quad s \in \operatorname{mat}(n, 1, \operatorname{vec}(d)) \wedge \\ \quad \operatorname{generates}(d, n, s) \wedge \\ \quad \operatorname{rel\text{-}free}(d, n, s) \wedge \\ \quad x \in \operatorname{vec}(d) \wedge \\ \quad j \in \operatorname{interval}(1, n) \Rightarrow \\ \quad \quad \operatorname{entry}((\lambda c \in \operatorname{mat}(1, n, \operatorname{carr}(\operatorname{scal}(d))).\; \operatorname{lincomb}(d, n, c, s))^{-1}(\operatorname{vneg}(d)(x)), \\ \quad \quad \quad 1, j) \\ \quad \quad =\; \operatorname{neg}(\operatorname{scal}(d))(\operatorname{entry}((\lambda c \in \operatorname{mat}(1, n, \operatorname{carr}(\operatorname{scal}(d))).\; \operatorname{lincomb}(d, n, c, s))^{-1}(x), \\ \quad \quad \quad \quad 1, j)) \end{array}}}
\vnbprop{lzt-fam-modules}{\vnbstmt{\begin{array}{@{}l@{}} \forall d, n, s, m, a. \\ \quad \operatorname{is\text{-}module}(d) \wedge \\ \quad n \in \mathbb{N} \wedge \\ \quad s \in \operatorname{mat}(n, 1, \operatorname{vec}(d)) \wedge \\ \quad \operatorname{generates}(d, n, s) \wedge \\ \quad \operatorname{rel\text{-}free}(d, n, s) \wedge \\ \quad \operatorname{scal}(d) = \operatorname{zz\text{-}ring} \wedge \\ \quad m \in \operatorname{mat}(n, n, \mathbb{Z}) \wedge \\ \quad \operatorname{is\text{-}diagonal}(\operatorname{zz\text{-}ring}, n, n, m) \wedge \\ \quad \operatorname{is\text{-}submodule}(d, a) \wedge \\ \quad \operatorname{spans}(d, n, \operatorname{matact}(d, m, s), a) \Rightarrow \\ \quad \quad \forall i \in \operatorname{interval}(1, n). \\ \quad \quad \quad \operatorname{is\text{-}module}((\lambda \mathit{fmi} \in \operatorname{interval}(1, n).\; \operatorname{zz\text{-}mod}(\operatorname{entry}(m, \mathit{fmi}, \mathit{fmi})))(i)) \wedge \\ \quad \quad \forall i \in \operatorname{interval}(1, n). \\ \quad \quad \quad \operatorname{scal}((\lambda \mathit{fmi} \in \operatorname{interval}(1, n).\; \operatorname{zz\text{-}mod}(\operatorname{entry}(m, \mathit{fmi}, \mathit{fmi})))(i)) = \operatorname{zz\text{-}ring} \end{array}}}
\vnbprop{lzt-ph-at}{\vnbstmt{\begin{array}{@{}l@{}} \forall d, n, s, m, a, x, j. \\ \quad \operatorname{is\text{-}module}(d) \wedge \\ \quad n \in \mathbb{N} \wedge \\ \quad s \in \operatorname{mat}(n, 1, \operatorname{vec}(d)) \wedge \\ \quad \operatorname{generates}(d, n, s) \wedge \\ \quad \operatorname{rel\text{-}free}(d, n, s) \wedge \\ \quad \operatorname{scal}(d) = \operatorname{zz\text{-}ring} \wedge \\ \quad m \in \operatorname{mat}(n, n, \mathbb{Z}) \wedge \\ \quad \operatorname{is\text{-}diagonal}(\operatorname{zz\text{-}ring}, n, n, m) \wedge \\ \quad \operatorname{is\text{-}submodule}(d, a) \wedge \\ \quad \operatorname{spans}(d, n, \operatorname{matact}(d, m, s), a) \wedge \\ \quad x \in \operatorname{vec}(d) \wedge \\ \quad j \in \operatorname{interval}(1, n) \Rightarrow \\ \quad \quad (\lambda \mathit{phx} \in \operatorname{vec}(d).\; (\lambda \mathit{phj} \in \operatorname{interval}(1, n).\; \operatorname{class}(\operatorname{qmod\text{-}setoid}(\operatorname{ring\text{-}as\text{-}module}(\operatorname{zz\text{-}ring}), \operatorname{principal\text{-}ideal}(\operatorname{zz\text{-}ring}, \operatorname{entry}(m, \mathit{phj}, \mathit{phj}))), \operatorname{entry}((\lambda c \in \operatorname{mat}(1, n, \operatorname{carr}(\operatorname{scal}(d))).\; \operatorname{lincomb}(d, n, c, s))^{-1}(\mathit{phx}), 1, \mathit{phj}))))(x)(j) \\ \quad \quad =\; \operatorname{class}(\operatorname{qmod\text{-}setoid}(\operatorname{ring\text{-}as\text{-}module}(\operatorname{zz\text{-}ring}), \operatorname{principal\text{-}ideal}(\operatorname{zz\text{-}ring}, \operatorname{entry}(m, j, j))), \\ \quad \quad \quad \operatorname{entry}((\lambda c \in \operatorname{mat}(1, n, \operatorname{carr}(\operatorname{scal}(d))).\; \operatorname{lincomb}(d, n, c, s))^{-1}(x), \\ \quad \quad \quad \quad 1, j)) \end{array}}}
\vnbprop{lzt-ph-fun}{\vnbstmt{\begin{array}{@{}l@{}} \forall d, n, s, m, a. \\ \quad \operatorname{is\text{-}module}(d) \wedge \\ \quad n \in \mathbb{N} \wedge \\ \quad s \in \operatorname{mat}(n, 1, \operatorname{vec}(d)) \wedge \\ \quad \operatorname{generates}(d, n, s) \wedge \\ \quad \operatorname{rel\text{-}free}(d, n, s) \wedge \\ \quad \operatorname{scal}(d) = \operatorname{zz\text{-}ring} \wedge \\ \quad m \in \operatorname{mat}(n, n, \mathbb{Z}) \wedge \\ \quad \operatorname{is\text{-}diagonal}(\operatorname{zz\text{-}ring}, n, n, m) \wedge \\ \quad \operatorname{is\text{-}submodule}(d, a) \wedge \\ \quad \operatorname{spans}(d, n, \operatorname{matact}(d, m, s), a) \Rightarrow \\ \quad \quad (\lambda x \in \operatorname{vec}(d).\; \\ \quad \quad \quad (\lambda j \in \operatorname{interval}(1, n).\; \\ \quad \quad \quad \quad \operatorname{class}(\operatorname{qmod\text{-}setoid}(\operatorname{ring\text{-}as\text{-}module}(\operatorname{zz\text{-}ring}), \operatorname{principal\text{-}ideal}(\operatorname{zz\text{-}ring}, \operatorname{entry}(m, j, j))), \\ \quad \quad \quad \quad \quad \operatorname{entry}((\lambda c \in \operatorname{mat}(1, n, \operatorname{carr}(\operatorname{scal}(d))).\; \operatorname{lincomb}(d, n, c, s))^{-1}(x), \\ \quad \quad \quad \quad \quad \quad 1, j)))) \\ \quad \quad \in\; (\operatorname{vec}(d) \to \operatorname{vec}(\operatorname{direct\text{-}sum}(\operatorname{zz\text{-}ring}, n, (\lambda i \in \operatorname{interval}(1, n).\; \operatorname{zz\text{-}mod}(\operatorname{entry}(m, i, i)))))) \end{array}}}
\vnbprop{lzt-ph-hom}{\vnbstmt{\begin{array}{@{}l@{}} \forall d, n, s, m, a. \\ \quad \operatorname{is\text{-}module}(d) \wedge \\ \quad n \in \mathbb{N} \wedge \\ \quad s \in \operatorname{mat}(n, 1, \operatorname{vec}(d)) \wedge \\ \quad \operatorname{generates}(d, n, s) \wedge \\ \quad \operatorname{rel\text{-}free}(d, n, s) \wedge \\ \quad \operatorname{scal}(d) = \operatorname{zz\text{-}ring} \wedge \\ \quad m \in \operatorname{mat}(n, n, \mathbb{Z}) \wedge \\ \quad \operatorname{is\text{-}diagonal}(\operatorname{zz\text{-}ring}, n, n, m) \wedge \\ \quad \operatorname{is\text{-}submodule}(d, a) \wedge \\ \quad \operatorname{spans}(d, n, \operatorname{matact}(d, m, s), a) \Rightarrow \\ \quad \quad \operatorname{is\text{-}hom\text{-}module}(d, \\ \quad \quad \quad \operatorname{direct\text{-}sum}(\operatorname{zz\text{-}ring}, n, (\lambda i \in \operatorname{interval}(1, n).\; \operatorname{zz\text{-}mod}(\operatorname{entry}(m, i, i)))), \\ \quad \quad \quad (\lambda x \in \operatorname{vec}(d).\; \\ \quad \quad \quad \quad (\lambda j \in \operatorname{interval}(1, n).\; \\ \quad \quad \quad \quad \quad \operatorname{class}(\operatorname{qmod\text{-}setoid}(\operatorname{ring\text{-}as\text{-}module}(\operatorname{zz\text{-}ring}), \operatorname{principal\text{-}ideal}(\operatorname{zz\text{-}ring}, \operatorname{entry}(m, j, j))), \\ \quad \quad \quad \quad \quad \quad \operatorname{entry}((\lambda c \in \operatorname{mat}(1, n, \operatorname{carr}(\operatorname{scal}(d))).\; \operatorname{lincomb}(d, n, c, s))^{-1}(x), \\ \quad \quad \quad \quad \quad \quad \quad 1, j))))) \end{array}}}
\vnbprop{lincomb-in-submodule}{\vnbstmt{\begin{array}{@{}l@{}} \forall d, m, n, c, s. \\ \quad \operatorname{is\text{-}module}(d) \wedge \\ \quad \operatorname{is\text{-}submodule}(d, m) \wedge \\ \quad n \in \mathbb{N} \wedge \\ \quad c \in \operatorname{mat}(1, n, \operatorname{carr}(\operatorname{scal}(d))) \wedge \\ \quad s \in \operatorname{mat}(n, 1, \operatorname{vec}(d)) \wedge \\ \quad \forall j \in \operatorname{interval}(1, n).\; (\operatorname{entry}(s, j, 1) \in m) \Rightarrow \\ \quad \quad \operatorname{lincomb}(d, n, c, s) \in m \end{array}}}
\vnbprop{lincomb-matact}{\vnbstmt{\begin{array}{@{}l@{}} \forall d, n, c, m, s. \\ \quad \operatorname{is\text{-}module}(d) \wedge \\ \quad n \in \mathbb{N} \wedge \\ \quad c \in \operatorname{mat}(1, n, \operatorname{carr}(\operatorname{scal}(d))) \wedge \\ \quad m \in \operatorname{mat}(n, n, \operatorname{carr}(\operatorname{scal}(d))) \wedge \\ \quad s \in \operatorname{mat}(n, 1, \operatorname{vec}(d)) \Rightarrow \\ \quad \quad \operatorname{lincomb}(d, n, c, \operatorname{matact}(d, m, s)) = \operatorname{lincomb}(d, n, \operatorname{matmul}(\operatorname{scal}(d), c, m), s) \end{array}}}
\vnbprop{diagonal-right-entry}{\vnbstmt{\begin{array}{@{}l@{}} \forall r, n, c, g, j. \\ \quad \operatorname{is\text{-}ring}(r) \wedge \\ \quad n \in \mathbb{N} \wedge \\ \quad c \in \operatorname{mat}(1, n, \operatorname{carr}(r)) \wedge \\ \quad g \in \operatorname{mat}(n, n, \operatorname{carr}(r)) \wedge \\ \quad \operatorname{is\text{-}diagonal}(r, n, n, g) \wedge \\ \quad j \in \operatorname{interval}(1, n) \Rightarrow \\ \quad \quad \operatorname{entry}(\operatorname{matmul}(r, c, g), 1, j) = \operatorname{mul}(r)(\operatorname{entry}(c, 1, j), \operatorname{entry}(g, j, j)) \end{array}}}
\vnbprop{lzt-ph-kernel}{\vnbstmt{\begin{array}{@{}l@{}} \forall d, n, s, m, a, x. \\ \quad \operatorname{is\text{-}module}(d) \wedge \\ \quad n \in \mathbb{N} \wedge \\ \quad s \in \operatorname{mat}(n, 1, \operatorname{vec}(d)) \wedge \\ \quad \operatorname{generates}(d, n, s) \wedge \\ \quad \operatorname{rel\text{-}free}(d, n, s) \wedge \\ \quad \operatorname{scal}(d) = \operatorname{zz\text{-}ring} \wedge \\ \quad m \in \operatorname{mat}(n, n, \mathbb{Z}) \wedge \\ \quad \operatorname{is\text{-}diagonal}(\operatorname{zz\text{-}ring}, n, n, m) \wedge \\ \quad \operatorname{is\text{-}submodule}(d, a) \wedge \\ \quad \operatorname{spans}(d, n, \operatorname{matact}(d, m, s), a) \wedge \\ \quad x \in \operatorname{vec}(d) \Rightarrow \\ \quad \quad (\lambda \mathit{phx} \in \operatorname{vec}(d).\; (\lambda j \in \operatorname{interval}(1, n).\; \operatorname{class}(\operatorname{qmod\text{-}setoid}(\operatorname{ring\text{-}as\text{-}module}(\operatorname{zz\text{-}ring}), \operatorname{principal\text{-}ideal}(\operatorname{zz\text{-}ring}, \operatorname{entry}(m, j, j))), \operatorname{entry}((\lambda c \in \operatorname{mat}(1, n, \operatorname{carr}(\operatorname{scal}(d))).\; \operatorname{lincomb}(d, n, c, s))^{-1}(\mathit{phx}), 1, j))))(x) \\ \quad \quad =\; \operatorname{vzero}(\operatorname{direct\text{-}sum}(\operatorname{zz\text{-}ring}, n, (\lambda i \in \operatorname{interval}(1, n).\; \operatorname{zz\text{-}mod}(\operatorname{entry}(m, i, i))))) \\ \quad \quad \Leftrightarrow\; x \in a \end{array}}}
\vnbprop{lzt-ph-onto}{\vnbstmt{\begin{array}{@{}l@{}} \forall d, n, s, m, a, w. \\ \quad \operatorname{is\text{-}module}(d) \wedge \\ \quad n \in \mathbb{N} \wedge \\ \quad s \in \operatorname{mat}(n, 1, \operatorname{vec}(d)) \wedge \\ \quad \operatorname{generates}(d, n, s) \wedge \\ \quad \operatorname{rel\text{-}free}(d, n, s) \wedge \\ \quad \operatorname{scal}(d) = \operatorname{zz\text{-}ring} \wedge \\ \quad m \in \operatorname{mat}(n, n, \mathbb{Z}) \wedge \\ \quad \operatorname{is\text{-}diagonal}(\operatorname{zz\text{-}ring}, n, n, m) \wedge \\ \quad \operatorname{is\text{-}submodule}(d, a) \wedge \\ \quad \operatorname{spans}(d, n, \operatorname{matact}(d, m, s), a) \wedge \\ \quad w \in \operatorname{vec}(\operatorname{direct\text{-}sum}(\operatorname{zz\text{-}ring}, n, (\lambda i \in \operatorname{interval}(1, n).\; \operatorname{zz\text{-}mod}(\operatorname{entry}(m, i, i))))) \Rightarrow \\ \quad \quad \exists x \in \operatorname{vec}(d). \\ \quad \quad \quad (\lambda \mathit{phx} \in \operatorname{vec}(d).\; (\lambda j \in \operatorname{interval}(1, n).\; \operatorname{class}(\operatorname{qmod\text{-}setoid}(\operatorname{ring\text{-}as\text{-}module}(\operatorname{zz\text{-}ring}), \operatorname{principal\text{-}ideal}(\operatorname{zz\text{-}ring}, \operatorname{entry}(m, j, j))), \operatorname{entry}((\lambda c \in \operatorname{mat}(1, n, \operatorname{carr}(\operatorname{scal}(d))).\; \operatorname{lincomb}(d, n, c, s))^{-1}(\mathit{phx}), 1, j))))(x) \\ \quad \quad \quad =\; w \end{array}}}
\vnbprop{quotient-by-diagonal-subgroup}{\vnbstmt{\begin{array}{@{}l@{}} \forall d, n, s, m, a. \\ \quad \operatorname{is\text{-}module}(d) \wedge \\ \quad n \in \mathbb{N} \wedge \\ \quad s \in \operatorname{mat}(n, 1, \operatorname{vec}(d)) \wedge \\ \quad \operatorname{generates}(d, n, s) \wedge \\ \quad \operatorname{rel\text{-}free}(d, n, s) \wedge \\ \quad \operatorname{scal}(d) = \operatorname{zz\text{-}ring} \wedge \\ \quad m \in \operatorname{mat}(n, n, \mathbb{Z}) \wedge \\ \quad \operatorname{is\text{-}diagonal}(\operatorname{zz\text{-}ring}, n, n, m) \wedge \\ \quad \operatorname{is\text{-}submodule}(d, a) \wedge \\ \quad \operatorname{spans}(d, n, \operatorname{matact}(d, m, s), a) \Rightarrow \\ \quad \quad \exists f. \\ \quad \quad \quad \operatorname{is\text{-}iso\text{-}module}(\operatorname{quotient\text{-}module}(d, a), \\ \quad \quad \quad \quad \operatorname{direct\text{-}sum}(\operatorname{zz\text{-}ring}, n, (\lambda i \in \operatorname{interval}(1, n).\; \operatorname{zz\text{-}mod}(\operatorname{entry}(m, i, i)))), \\ \quad \quad \quad \quad f) \end{array}}}

\vnbbookfile{theorem-library/la-z4-free}{theorem-library/la-z4-free}
\noindent{\small la-z4-free.scm -- module homomorphisms and LINCOMB, isomorphisms of modules, the free module Z\textasciicircum{}n and its unit basis (LA-Z4, 2026-10-24).  Vocabulary: structure-library/la-z4-vocab (ZZ-POW, ZZ-UNITS, ZZ-ROW, MAP-SEQ) and la-z3-vocab (DIRECT-SUM, RING-AS-MODULE, IS-ISO-MODULE).  Prefix lzf-.}\par
\vnbprop{map-seq-type}{\vnbstmt{\begin{array}{@{}l@{}} \forall h, v, w, n, u. \\ \quad ((h \in (v \to w)) \wedge ((n \in \mathbb{N}) \wedge (u \in \operatorname{mat}(n, 1, v)))) \Rightarrow  \\ \quad \quad \operatorname{map\text{-}seq}(h, n, u) \in \operatorname{mat}(n, 1, w) \end{array}}}
\vnbprop{map-seq-entry}{\vnbstmt{\begin{array}{@{}l@{}} \forall h, v, w, n, u, i. \\ \quad ((h \in (v \to w)) \wedge ((n \in \mathbb{N}) \wedge ((u \in \operatorname{mat}(n, 1, v)) \wedge (i \in \operatorname{interval}(1, n))))) \Rightarrow  \\ \quad \quad \operatorname{entry}(\operatorname{map\text{-}seq}(h, n, u), i, 1) = h(\operatorname{entry}(u, i, 1)) \end{array}}}
\vnbprop{map-seq-block}{\vnbstmt{\begin{array}{@{}l@{}} \forall h, v, w, m, k, u. \\ \quad ((h \in (v \to w)) \wedge ((m \in \mathbb{N}) \wedge ((k \in \mathbb{N}) \wedge ((k \leq m) \wedge (u \in \operatorname{mat}(m, 1, v)))))) \Rightarrow  \\ \quad \quad \operatorname{block}(\operatorname{map\text{-}seq}(h, m, u), k, 1) = \operatorname{map\text{-}seq}(h, k, \operatorname{block}(u, k, 1)) \end{array}}}
\vnbprop{lzf-hom-lincomb-ind}{\vnbstmt{\begin{array}{@{}l@{}} \forall n \in \mathbb{N}, a, b, h, c, u. \\ \quad (\operatorname{is\text{-}hom\text{-}module}(a, b, h) \wedge ((c \in \operatorname{mat}(1, n, \operatorname{carr}(\operatorname{scal}(a)))) \wedge (u \in \operatorname{mat}(n, 1, \operatorname{vec}(a))))) \Rightarrow  \\ \quad \quad h(\operatorname{lincomb}(a, n, c, u)) = \operatorname{lincomb}(b, n, c, \operatorname{map\text{-}seq}(h, n, u)) \end{array}}}
\vnbprop{hom-lincomb}{\vnbstmt{\begin{array}{@{}l@{}} \forall a, b, h, n, c, u. \\ \quad (\operatorname{is\text{-}hom\text{-}module}(a, b, h) \wedge ((n \in \mathbb{N}) \wedge ((c \in \operatorname{mat}(1, n, \operatorname{carr}(\operatorname{scal}(a)))) \wedge (u \in \operatorname{mat}(n, 1, \operatorname{vec}(a)))))) \Rightarrow  \\ \quad \quad h(\operatorname{lincomb}(a, n, c, u)) = \operatorname{lincomb}(b, n, c, \operatorname{map\text{-}seq}(h, n, u)) \end{array}}}
\vnbprop{iso-module-inverse}{\vnbstmt{\begin{array}{@{}l@{}} \forall a, b, f. \\ \quad \operatorname{is\text{-}iso\text{-}module}(a, b, f) \Rightarrow  \\ \quad \quad \operatorname{is\text{-}iso\text{-}module}(b, a, f^{-1}) \end{array}}}
\vnbprop{iso-module-compose}{\vnbstmt{\begin{array}{@{}l@{}} \forall a, b, c, f, g. \\ \quad (\operatorname{is\text{-}iso\text{-}module}(a, b, f) \wedge \operatorname{is\text{-}iso\text{-}module}(b, c, g)) \Rightarrow  \\ \quad \quad \operatorname{is\text{-}iso\text{-}module}(a, c, \operatorname{compose}(g, f)) \end{array}}}
\vnbprop{lincomb-single-support}{\vnbstmt{\begin{array}{@{}l@{}} \forall d, n, c, u, i. \\ \quad \operatorname{is\text{-}module}(d) \wedge \\ \quad n \in \mathbb{N} \wedge \\ \quad c \in \operatorname{mat}(1, n, \operatorname{carr}(\operatorname{scal}(d))) \wedge \\ \quad u \in \operatorname{mat}(n, 1, \operatorname{vec}(d)) \wedge \\ \quad i \in \operatorname{interval}(1, n) \wedge \\ \quad \forall j \in \operatorname{interval}(1, n).\; (\neg (j = i) \Rightarrow (\operatorname{entry}(u, j, 1) = \operatorname{vzero}(d))) \Rightarrow \\ \quad \quad \operatorname{lincomb}(d, n, c, u) = \operatorname{act}(d)(\operatorname{entry}(c, 1, i), \operatorname{entry}(u, i, 1)) \end{array}}}
\vnbprop{dsum-eval-hom}{\vnbstmt{\begin{array}{@{}l@{}} \forall s, n, f, i. \\ \quad \operatorname{is\text{-}ring}(s) \wedge \\ \quad n \in \mathbb{N} \wedge \\ \quad \forall a \in \operatorname{interval}(1, n).\; \operatorname{is\text{-}module}(f(a)) \wedge \\ \quad \forall a \in \operatorname{interval}(1, n).\; (\operatorname{scal}(f(a)) = s) \wedge \\ \quad i \in \operatorname{interval}(1, n) \Rightarrow \\ \quad \quad \operatorname{is\text{-}hom\text{-}module}(\operatorname{direct\text{-}sum}(s, n, f), f(i), (\lambda x \in \operatorname{vec}(\operatorname{direct\text{-}sum}(s, n, f)).\; x(i))) \end{array}}}
\vnbprop{dsum-lincomb-at}{\vnbstmt{\begin{array}{@{}l@{}} \forall s, n, f, m, c, u, i. \\ \quad \operatorname{is\text{-}ring}(s) \wedge \\ \quad n \in \mathbb{N} \wedge \\ \quad \forall a \in \operatorname{interval}(1, n).\; \operatorname{is\text{-}module}(f(a)) \wedge \\ \quad \forall a \in \operatorname{interval}(1, n).\; (\operatorname{scal}(f(a)) = s) \wedge \\ \quad m \in \mathbb{N} \wedge \\ \quad c \in \operatorname{mat}(1, m, \operatorname{carr}(s)) \wedge \\ \quad u \in \operatorname{mat}(m, 1, \operatorname{vec}(\operatorname{direct\text{-}sum}(s, n, f))) \wedge \\ \quad i \in \operatorname{interval}(1, n) \Rightarrow \\ \quad \quad \operatorname{lincomb}(\operatorname{direct\text{-}sum}(s, n, f), m, c, u)(i) \\ \quad \quad =\; \operatorname{lincomb}(f(i), m, c, \operatorname{map\text{-}seq}((\lambda x \in \operatorname{vec}(\operatorname{direct\text{-}sum}(s, n, f)).\; x(i)), m, u)) \end{array}}}
\vnbprop{zz-pow-unfold}{\vnbstmt{\begin{array}{@{}l@{}} \forall n. \\ \quad (\operatorname{zz\text{-}pow}(n) \simeq \operatorname{direct\text{-}sum}(\operatorname{zz\text{-}ring}, n, (\lambda i \in \operatorname{interval}(1, n).\; \operatorname{ring\text{-}as\text{-}module}(\operatorname{zz\text{-}ring})))) \end{array}}}
\vnbprop{zz-pow-family}{\vnbstmt{\begin{array}{@{}l@{}} \forall n \in \mathbb{N}. \\ \quad \forall i \in \operatorname{interval}(1, n). \\ \quad \quad \operatorname{is\text{-}module}((\lambda \mathit{zpi} \in \operatorname{interval}(1, n).\; \operatorname{ring\text{-}as\text{-}module}(\operatorname{zz\text{-}ring}))(i)) \wedge \\ \quad \forall i \in \operatorname{interval}(1, n). \\ \quad \quad \operatorname{scal}((\lambda \mathit{zpi} \in \operatorname{interval}(1, n).\; \operatorname{ring\text{-}as\text{-}module}(\operatorname{zz\text{-}ring}))(i)) = \operatorname{zz\text{-}ring} \end{array}}}
\vnbprop{zz-pow-is-module}{\vnbstmt{\forall n \in \mathbb{N}.\; \operatorname{is\text{-}module}(\operatorname{zz\text{-}pow}(n))}}
\vnbprop{zz-pow-scal}{\vnbstmt{\forall n \in \mathbb{N}.\; (\operatorname{scal}(\operatorname{zz\text{-}pow}(n)) = \operatorname{zz\text{-}ring})}}
\vnbprop{zz-pow-vec-in}{\vnbstmt{\begin{array}{@{}l@{}} \forall n, x. \\ \quad ((n \in \mathbb{N}) \wedge (x \in (\operatorname{interval}(1, n) \to \mathbb{Z}))) \Rightarrow (x \in \operatorname{vec}(\operatorname{zz\text{-}pow}(n))) \end{array}}}
\vnbprop{zz-pow-vec-out}{\vnbstmt{\begin{array}{@{}l@{}} \forall n, x, i. \\ \quad ((n \in \mathbb{N}) \wedge ((x \in \operatorname{vec}(\operatorname{zz\text{-}pow}(n))) \wedge (i \in \operatorname{interval}(1, n)))) \Rightarrow (x(i) \in \mathbb{Z}) \end{array}}}
\vnbprop{zz-pow-ext}{\vnbstmt{\begin{array}{@{}l@{}} \forall n, x, y. \\ \quad n \in \mathbb{N} \wedge \\ \quad x \in \operatorname{vec}(\operatorname{zz\text{-}pow}(n)) \wedge \\ \quad y \in \operatorname{vec}(\operatorname{zz\text{-}pow}(n)) \wedge \\ \quad \forall i \in \operatorname{interval}(1, n).\; (x(i) = y(i)) \Rightarrow \\ \quad \quad x = y \end{array}}}
\vnbprop{zz-pow-vzero-at}{\vnbstmt{\begin{array}{@{}l@{}} \forall n, i. \\ \quad ((n \in \mathbb{N}) \wedge (i \in \operatorname{interval}(1, n))) \Rightarrow (\operatorname{vzero}(\operatorname{zz\text{-}pow}(n))(i) = \operatorname{zero}(\operatorname{zz\text{-}ring})) \end{array}}}
\vnbprop{zz-unit-vec}{\vnbstmt{\begin{array}{@{}l@{}} \forall n \in \mathbb{N}, k. \\ \quad (\lambda t \in \operatorname{interval}(1, n).\; \operatorname{entry}(\operatorname{unitrow}(\operatorname{zz\text{-}ring}, n, k), 1, t)) \in \operatorname{vec}(\operatorname{zz\text{-}pow}(n)) \end{array}}}
\vnbprop{zz-units-type}{\vnbstmt{\forall n \in \mathbb{N}.\; (\operatorname{zz\text{-}units}(n) \in \operatorname{mat}(n, 1, \operatorname{vec}(\operatorname{zz\text{-}pow}(n))))}}
\vnbprop{zz-units-entry}{\vnbstmt{\begin{array}{@{}l@{}} \forall n, k. \\ \quad ((n \in \mathbb{N}) \wedge (k \in \operatorname{interval}(1, n))) \Rightarrow  \\ \quad \quad \operatorname{entry}(\operatorname{zz\text{-}units}(n), k, 1) = (\lambda t \in \operatorname{interval}(1, n).\; \operatorname{entry}(\operatorname{unitrow}(\operatorname{zz\text{-}ring}, n, k), 1, t)) \end{array}}}
\vnbprop{zz-unit-at}{\vnbstmt{\begin{array}{@{}l@{}} \forall n, k, i. \\ \quad ((n \in \mathbb{N}) \wedge (i \in \operatorname{interval}(1, n))) \Rightarrow  \\ \quad \quad (\lambda t \in \operatorname{interval}(1, n).\; \operatorname{entry}(\operatorname{unitrow}(\operatorname{zz\text{-}ring}, n, k), 1, t))(i) \\ \quad \quad =\; \operatorname{entry}(\operatorname{unitrow}(\operatorname{zz\text{-}ring}, n, k), 1, i) \end{array}}}
\vnbprop{zz-pow-lincomb-at}{\vnbstmt{\begin{array}{@{}l@{}} \forall n, c, i. \\ \quad ((n \in \mathbb{N}) \wedge ((c \in \operatorname{mat}(1, n, \operatorname{carr}(\operatorname{zz\text{-}ring}))) \wedge (i \in \operatorname{interval}(1, n)))) \Rightarrow  \\ \quad \quad \operatorname{lincomb}(\operatorname{zz\text{-}pow}(n), n, c, \operatorname{zz\text{-}units}(n))(i) = \operatorname{entry}(c, 1, i) \end{array}}}
\vnbprop{zz-row-type}{\vnbstmt{\begin{array}{@{}l@{}} \forall n, x. \\ \quad ((n \in \mathbb{N}) \wedge (x \in \operatorname{vec}(\operatorname{zz\text{-}pow}(n)))) \Rightarrow (\operatorname{zz\text{-}row}(n, x) \in \operatorname{mat}(1, n, \operatorname{carr}(\operatorname{zz\text{-}ring}))) \end{array}}}
\vnbprop{zz-row-entry}{\vnbstmt{\begin{array}{@{}l@{}} \forall n, x, i. \\ \quad ((n \in \mathbb{N}) \wedge ((x \in \operatorname{vec}(\operatorname{zz\text{-}pow}(n))) \wedge (i \in \operatorname{interval}(1, n)))) \Rightarrow  \\ \quad \quad \operatorname{entry}(\operatorname{zz\text{-}row}(n, x), 1, i) = x(i) \end{array}}}
\vnbprop{zz-pow-row-lincomb}{\vnbstmt{\begin{array}{@{}l@{}} \forall n, x. \\ \quad ((n \in \mathbb{N}) \wedge (x \in \operatorname{vec}(\operatorname{zz\text{-}pow}(n)))) \Rightarrow  \\ \quad \quad \operatorname{lincomb}(\operatorname{zz\text{-}pow}(n), n, \operatorname{zz\text{-}row}(n, x), \operatorname{zz\text{-}units}(n)) = x \end{array}}}
\vnbprop{zz-units-generate}{\vnbstmt{\forall n \in \mathbb{N}.\; \operatorname{generates}(\operatorname{zz\text{-}pow}(n), n, \operatorname{zz\text{-}units}(n))}}
\vnbprop{zz-units-rel-free}{\vnbstmt{\forall n \in \mathbb{N}.\; \operatorname{rel\text{-}free}(\operatorname{zz\text{-}pow}(n), n, \operatorname{zz\text{-}units}(n))}}
\vnbprop{hom-kernel-submodule}{\vnbstmt{\begin{array}{@{}l@{}} \forall a, b, h. \\ \quad \operatorname{is\text{-}hom\text{-}module}(a, b, h) \Rightarrow \operatorname{is\text{-}submodule}(a, \{x \in \operatorname{vec}(a) : (h(x) = \operatorname{vzero}(b))\}) \end{array}}}
\vnbprop{free-lift-type}{\vnbstmt{\begin{array}{@{}l@{}} \forall d, n, s, m, w. \\ \quad \operatorname{is\text{-}module}(d) \wedge \\ \quad n \in \mathbb{N} \wedge \\ \quad s \in \operatorname{mat}(n, 1, \operatorname{vec}(d)) \wedge \\ \quad \operatorname{generates}(d, n, s) \wedge \\ \quad \operatorname{rel\text{-}free}(d, n, s) \wedge \\ \quad \operatorname{is\text{-}module}(m) \wedge \\ \quad \operatorname{scal}(m) = \operatorname{scal}(d) \wedge \\ \quad w \in \operatorname{mat}(n, 1, \operatorname{vec}(m)) \Rightarrow \\ \quad \quad (\lambda x \in \operatorname{vec}(d).\; \\ \quad \quad \quad \operatorname{lincomb}(m, n, \\ \quad \quad \quad \quad (\lambda c \in \operatorname{mat}(1, n, \operatorname{carr}(\operatorname{scal}(d))).\; \operatorname{lincomb}(d, n, c, s))^{-1}(x), \\ \quad \quad \quad \quad w)) \\ \quad \quad \in\; (\operatorname{vec}(d) \to \operatorname{vec}(m)) \end{array}}}
\vnbprop{free-lift-lincomb}{\vnbstmt{\begin{array}{@{}l@{}} \forall d, n, s, m, w, c. \\ \quad \operatorname{is\text{-}module}(d) \wedge \\ \quad n \in \mathbb{N} \wedge \\ \quad s \in \operatorname{mat}(n, 1, \operatorname{vec}(d)) \wedge \\ \quad \operatorname{generates}(d, n, s) \wedge \\ \quad \operatorname{rel\text{-}free}(d, n, s) \wedge \\ \quad \operatorname{is\text{-}module}(m) \wedge \\ \quad \operatorname{scal}(m) = \operatorname{scal}(d) \wedge \\ \quad w \in \operatorname{mat}(n, 1, \operatorname{vec}(m)) \wedge \\ \quad c \in \operatorname{mat}(1, n, \operatorname{carr}(\operatorname{scal}(d))) \Rightarrow \\ \quad \quad (\lambda x \in \operatorname{vec}(d).\; \operatorname{lincomb}(m, n, (\lambda \mathit{lcc} \in \operatorname{mat}(1, n, \operatorname{carr}(\operatorname{scal}(d))).\; \operatorname{lincomb}(d, n, \mathit{lcc}, s))^{-1}(x), w))(\operatorname{lincomb}(d, n, c, s)) \\ \quad \quad =\; \operatorname{lincomb}(m, n, c, w) \end{array}}}
\vnbprop{free-lift-basis}{\vnbstmt{\begin{array}{@{}l@{}} \forall d, n, s, m, w, j. \\ \quad \operatorname{is\text{-}module}(d) \wedge \\ \quad n \in \mathbb{N} \wedge \\ \quad s \in \operatorname{mat}(n, 1, \operatorname{vec}(d)) \wedge \\ \quad \operatorname{generates}(d, n, s) \wedge \\ \quad \operatorname{rel\text{-}free}(d, n, s) \wedge \\ \quad \operatorname{is\text{-}module}(m) \wedge \\ \quad \operatorname{scal}(m) = \operatorname{scal}(d) \wedge \\ \quad w \in \operatorname{mat}(n, 1, \operatorname{vec}(m)) \wedge \\ \quad j \in \operatorname{interval}(1, n) \Rightarrow \\ \quad \quad (\lambda x \in \operatorname{vec}(d).\; \operatorname{lincomb}(m, n, (\lambda c \in \operatorname{mat}(1, n, \operatorname{carr}(\operatorname{scal}(d))).\; \operatorname{lincomb}(d, n, c, s))^{-1}(x), w))(\operatorname{entry}(s, j, 1)) \\ \quad \quad =\; \operatorname{entry}(w, j, 1) \end{array}}}
\vnbprop{free-lift-add}{\vnbstmt{\begin{array}{@{}l@{}} \forall d, n, s, m, w, x, y. \\ \quad \operatorname{is\text{-}module}(d) \wedge \\ \quad n \in \mathbb{N} \wedge \\ \quad s \in \operatorname{mat}(n, 1, \operatorname{vec}(d)) \wedge \\ \quad \operatorname{generates}(d, n, s) \wedge \\ \quad \operatorname{rel\text{-}free}(d, n, s) \wedge \\ \quad \operatorname{is\text{-}module}(m) \wedge \\ \quad \operatorname{scal}(m) = \operatorname{scal}(d) \wedge \\ \quad w \in \operatorname{mat}(n, 1, \operatorname{vec}(m)) \wedge \\ \quad x \in \operatorname{vec}(d) \wedge \\ \quad y \in \operatorname{vec}(d) \Rightarrow \\ \quad \quad (\lambda \mathit{lex} \in \operatorname{vec}(d).\; \operatorname{lincomb}(m, n, (\lambda c \in \operatorname{mat}(1, n, \operatorname{carr}(\operatorname{scal}(d))).\; \operatorname{lincomb}(d, n, c, s))^{-1}(\mathit{lex}), w))(\operatorname{vadd}(d)(x, y)) \\ \quad \quad =\; \operatorname{vadd}(m)((\lambda \mathit{lex} \in \operatorname{vec}(d).\; \operatorname{lincomb}(m, n, (\lambda c \in \operatorname{mat}(1, n, \operatorname{carr}(\operatorname{scal}(d))).\; \operatorname{lincomb}(d, n, c, s))^{-1}(\mathit{lex}), w))(x), \\ \quad \quad \quad (\lambda \mathit{lex} \in \operatorname{vec}(d).\; \operatorname{lincomb}(m, n, (\lambda c \in \operatorname{mat}(1, n, \operatorname{carr}(\operatorname{scal}(d))).\; \operatorname{lincomb}(d, n, c, s))^{-1}(\mathit{lex}), w))(y)) \end{array}}}
\vnbprop{free-lift-act}{\vnbstmt{\begin{array}{@{}l@{}} \forall d, n, s, m, w, r, x. \\ \quad \operatorname{is\text{-}module}(d) \wedge \\ \quad n \in \mathbb{N} \wedge \\ \quad s \in \operatorname{mat}(n, 1, \operatorname{vec}(d)) \wedge \\ \quad \operatorname{generates}(d, n, s) \wedge \\ \quad \operatorname{rel\text{-}free}(d, n, s) \wedge \\ \quad \operatorname{is\text{-}module}(m) \wedge \\ \quad \operatorname{scal}(m) = \operatorname{scal}(d) \wedge \\ \quad w \in \operatorname{mat}(n, 1, \operatorname{vec}(m)) \wedge \\ \quad r \in \operatorname{carr}(\operatorname{scal}(d)) \wedge \\ \quad x \in \operatorname{vec}(d) \Rightarrow \\ \quad \quad (\lambda \mathit{lex} \in \operatorname{vec}(d).\; \operatorname{lincomb}(m, n, (\lambda c \in \operatorname{mat}(1, n, \operatorname{carr}(\operatorname{scal}(d))).\; \operatorname{lincomb}(d, n, c, s))^{-1}(\mathit{lex}), w))(\operatorname{act}(d)(r, x)) \\ \quad \quad =\; \operatorname{act}(m)(r, \\ \quad \quad \quad (\lambda \mathit{lex} \in \operatorname{vec}(d).\; \operatorname{lincomb}(m, n, (\lambda c \in \operatorname{mat}(1, n, \operatorname{carr}(\operatorname{scal}(d))).\; \operatorname{lincomb}(d, n, c, s))^{-1}(\mathit{lex}), w))(x)) \end{array}}}
\vnbprop{free-lift-zero}{\vnbstmt{\begin{array}{@{}l@{}} \forall d, n, s, m, w. \\ \quad \operatorname{is\text{-}module}(d) \wedge \\ \quad n \in \mathbb{N} \wedge \\ \quad s \in \operatorname{mat}(n, 1, \operatorname{vec}(d)) \wedge \\ \quad \operatorname{generates}(d, n, s) \wedge \\ \quad \operatorname{rel\text{-}free}(d, n, s) \wedge \\ \quad \operatorname{is\text{-}module}(m) \wedge \\ \quad \operatorname{scal}(m) = \operatorname{scal}(d) \wedge \\ \quad w \in \operatorname{mat}(n, 1, \operatorname{vec}(m)) \Rightarrow \\ \quad \quad (\lambda x \in \operatorname{vec}(d).\; \operatorname{lincomb}(m, n, (\lambda c \in \operatorname{mat}(1, n, \operatorname{carr}(\operatorname{scal}(d))).\; \operatorname{lincomb}(d, n, c, s))^{-1}(x), w))(\operatorname{vzero}(d)) \\ \quad \quad =\; \operatorname{vzero}(m) \end{array}}}
\vnbprop{free-lift-hom}{\vnbstmt{\begin{array}{@{}l@{}} \forall d, n, s, m, w. \\ \quad \operatorname{is\text{-}module}(d) \wedge \\ \quad n \in \mathbb{N} \wedge \\ \quad s \in \operatorname{mat}(n, 1, \operatorname{vec}(d)) \wedge \\ \quad \operatorname{generates}(d, n, s) \wedge \\ \quad \operatorname{rel\text{-}free}(d, n, s) \wedge \\ \quad \operatorname{is\text{-}module}(m) \wedge \\ \quad \operatorname{scal}(m) = \operatorname{scal}(d) \wedge \\ \quad w \in \operatorname{mat}(n, 1, \operatorname{vec}(m)) \Rightarrow \\ \quad \quad \operatorname{is\text{-}hom\text{-}module}(d, m, \\ \quad \quad \quad (\lambda x \in \operatorname{vec}(d).\; \\ \quad \quad \quad \quad \operatorname{lincomb}(m, n, \\ \quad \quad \quad \quad \quad (\lambda c \in \operatorname{mat}(1, n, \operatorname{carr}(\operatorname{scal}(d))).\; \operatorname{lincomb}(d, n, c, s))^{-1}(x), \\ \quad \quad \quad \quad \quad w))) \end{array}}}
\vnbprop{iso-free-generating}{\vnbstmt{\begin{array}{@{}l@{}} \forall a, b, g, n, s. \\ \quad \operatorname{is\text{-}iso\text{-}module}(a, b, g) \wedge \\ \quad n \in \mathbb{N} \wedge \\ \quad s \in \operatorname{mat}(n, 1, \operatorname{vec}(a)) \wedge \\ \quad \operatorname{generates}(a, n, s) \wedge \\ \quad \operatorname{rel\text{-}free}(a, n, s) \Rightarrow \\ \quad \quad \operatorname{map\text{-}seq}(g, n, s) \in \operatorname{mat}(n, 1, \operatorname{vec}(b)) \wedge \\ \quad \quad \operatorname{generates}(b, n, \operatorname{map\text{-}seq}(g, n, s)) \wedge \\ \quad \quad \operatorname{rel\text{-}free}(b, n, \operatorname{map\text{-}seq}(g, n, s)) \end{array}}}

\vnbbookfile{theorem-library/ring-pow}{theorem-library/ring-pow}
\noindent{\small ring-pow.scm -- wave of 2026-11-09/10 (agent CCR-1109-E), wired 2026-11-10 from scratchpad/ccr1109e/ring-pow.scm. Library: linear-algebra.}\par
\vnbprop{ring-pow-unfold}{\vnbstmt{\begin{array}{@{}l@{}} \forall r, n. \\ \quad (\operatorname{ring\text{-}pow}(r, n) \simeq \operatorname{direct\text{-}sum}(r, n, (\lambda i \in \operatorname{interval}(1, n).\; \operatorname{ring\text{-}as\text{-}module}(r)))) \end{array}}}
\vnbprop{ring-pow-family}{\vnbstmt{\begin{array}{@{}l@{}} \forall r, n. \\ \quad (\operatorname{is\text{-}ring}(r) \wedge (n \in \mathbb{N})) \Rightarrow  \\ \quad \quad \forall i \in \operatorname{interval}(1, n). \\ \quad \quad \quad \operatorname{is\text{-}module}((\lambda \mathit{rqpi} \in \operatorname{interval}(1, n).\; \operatorname{ring\text{-}as\text{-}module}(r))(i)) \wedge \\ \quad \quad \forall i \in \operatorname{interval}(1, n). \\ \quad \quad \quad \operatorname{scal}((\lambda \mathit{rqpi} \in \operatorname{interval}(1, n).\; \operatorname{ring\text{-}as\text{-}module}(r))(i)) = r \end{array}}}
\vnbprop{ring-pow-is-module}{\vnbstmt{\forall r, n.\; ((\operatorname{is\text{-}ring}(r) \wedge (n \in \mathbb{N})) \Rightarrow \operatorname{is\text{-}module}(\operatorname{ring\text{-}pow}(r, n)))}}
\vnbprop{ring-pow-scal}{\vnbstmt{\forall r, n.\; ((\operatorname{is\text{-}ring}(r) \wedge (n \in \mathbb{N})) \Rightarrow (\operatorname{scal}(\operatorname{ring\text{-}pow}(r, n)) = r))}}
\vnbprop{ring-pow-vec-in}{\vnbstmt{\begin{array}{@{}l@{}} \forall r, n, x. \\ \quad (\operatorname{is\text{-}ring}(r) \wedge ((n \in \mathbb{N}) \wedge (x \in (\operatorname{interval}(1, n) \to \operatorname{carr}(r))))) \Rightarrow  \\ \quad \quad x \in \operatorname{vec}(\operatorname{ring\text{-}pow}(r, n)) \end{array}}}
\vnbprop{ring-pow-vec-out}{\vnbstmt{\begin{array}{@{}l@{}} \forall r, n, x, i. \\ \quad (\operatorname{is\text{-}ring}(r) \wedge ((n \in \mathbb{N}) \wedge ((x \in \operatorname{vec}(\operatorname{ring\text{-}pow}(r, n))) \wedge (i \in \operatorname{interval}(1, n))))) \Rightarrow  \\ \quad \quad x(i) \in \operatorname{carr}(r) \end{array}}}
\vnbprop{ring-pow-ext}{\vnbstmt{\begin{array}{@{}l@{}} \forall r, n, x, y. \\ \quad \operatorname{is\text{-}ring}(r) \wedge \\ \quad n \in \mathbb{N} \wedge \\ \quad x \in \operatorname{vec}(\operatorname{ring\text{-}pow}(r, n)) \wedge \\ \quad y \in \operatorname{vec}(\operatorname{ring\text{-}pow}(r, n)) \wedge \\ \quad \forall i \in \operatorname{interval}(1, n).\; (x(i) = y(i)) \Rightarrow \\ \quad \quad x = y \end{array}}}
\vnbprop{ring-pow-vzero-at}{\vnbstmt{\begin{array}{@{}l@{}} \forall r, n, i. \\ \quad (\operatorname{is\text{-}ring}(r) \wedge ((n \in \mathbb{N}) \wedge (i \in \operatorname{interval}(1, n)))) \Rightarrow  \\ \quad \quad \operatorname{vzero}(\operatorname{ring\text{-}pow}(r, n))(i) = \operatorname{zero}(r) \end{array}}}
\vnbprop{ring-unit-vec}{\vnbstmt{\begin{array}{@{}l@{}} \forall r, n, k. \\ \quad (\operatorname{is\text{-}ring}(r) \wedge (n \in \mathbb{N})) \Rightarrow  \\ \quad \quad (\lambda t \in \operatorname{interval}(1, n).\; \operatorname{entry}(\operatorname{unitrow}(r, n, k), 1, t)) \in \operatorname{vec}(\operatorname{ring\text{-}pow}(r, n)) \end{array}}}
\vnbprop{ring-units-type}{\vnbstmt{\begin{array}{@{}l@{}} \forall r, n. \\ \quad (\operatorname{is\text{-}ring}(r) \wedge (n \in \mathbb{N})) \Rightarrow  \\ \quad \quad \operatorname{ring\text{-}units}(r, n) \in \operatorname{mat}(n, 1, \operatorname{vec}(\operatorname{ring\text{-}pow}(r, n))) \end{array}}}
\vnbprop{ring-units-entry}{\vnbstmt{\begin{array}{@{}l@{}} \forall r, n, k. \\ \quad (\operatorname{is\text{-}ring}(r) \wedge ((n \in \mathbb{N}) \wedge (k \in \operatorname{interval}(1, n)))) \Rightarrow  \\ \quad \quad \operatorname{entry}(\operatorname{ring\text{-}units}(r, n), k, 1) = (\lambda t \in \operatorname{interval}(1, n).\; \operatorname{entry}(\operatorname{unitrow}(r, n, k), 1, t)) \end{array}}}
\vnbprop{ring-unit-at}{\vnbstmt{\begin{array}{@{}l@{}} \forall r, n, k, i. \\ \quad (\operatorname{is\text{-}ring}(r) \wedge ((n \in \mathbb{N}) \wedge (i \in \operatorname{interval}(1, n)))) \Rightarrow  \\ \quad \quad (\lambda t \in \operatorname{interval}(1, n).\; \operatorname{entry}(\operatorname{unitrow}(r, n, k), 1, t))(i) = \operatorname{entry}(\operatorname{unitrow}(r, n, k), 1, i) \end{array}}}
\vnbprop{ring-pow-lincomb-at}{\vnbstmt{\begin{array}{@{}l@{}} \forall r, n, c, i. \\ \quad (\operatorname{is\text{-}ring}(r) \wedge ((n \in \mathbb{N}) \wedge ((c \in \operatorname{mat}(1, n, \operatorname{carr}(r))) \wedge (i \in \operatorname{interval}(1, n))))) \Rightarrow  \\ \quad \quad \operatorname{lincomb}(\operatorname{ring\text{-}pow}(r, n), n, c, \operatorname{ring\text{-}units}(r, n))(i) = \operatorname{entry}(c, 1, i) \end{array}}}
\vnbprop{ring-row-type}{\vnbstmt{\begin{array}{@{}l@{}} \forall r, n, x. \\ \quad (\operatorname{is\text{-}ring}(r) \wedge ((n \in \mathbb{N}) \wedge (x \in \operatorname{vec}(\operatorname{ring\text{-}pow}(r, n))))) \Rightarrow  \\ \quad \quad \operatorname{ring\text{-}row}(r, n, x) \in \operatorname{mat}(1, n, \operatorname{carr}(r)) \end{array}}}
\vnbprop{ring-row-entry}{\vnbstmt{\begin{array}{@{}l@{}} \forall r, n, x, i. \\ \quad (\operatorname{is\text{-}ring}(r) \wedge ((n \in \mathbb{N}) \wedge ((x \in \operatorname{vec}(\operatorname{ring\text{-}pow}(r, n))) \wedge (i \in \operatorname{interval}(1, n))))) \Rightarrow  \\ \quad \quad \operatorname{entry}(\operatorname{ring\text{-}row}(r, n, x), 1, i) = x(i) \end{array}}}
\vnbprop{ring-pow-row-lincomb}{\vnbstmt{\begin{array}{@{}l@{}} \forall r, n, x. \\ \quad (\operatorname{is\text{-}ring}(r) \wedge ((n \in \mathbb{N}) \wedge (x \in \operatorname{vec}(\operatorname{ring\text{-}pow}(r, n))))) \Rightarrow  \\ \quad \quad \operatorname{lincomb}(\operatorname{ring\text{-}pow}(r, n), n, \operatorname{ring\text{-}row}(r, n, x), \operatorname{ring\text{-}units}(r, n)) = x \end{array}}}
\vnbprop{ring-units-generate}{\vnbstmt{\begin{array}{@{}l@{}} \forall r, n. \\ \quad (\operatorname{is\text{-}ring}(r) \wedge (n \in \mathbb{N})) \Rightarrow \operatorname{generates}(\operatorname{ring\text{-}pow}(r, n), n, \operatorname{ring\text{-}units}(r, n)) \end{array}}}
\vnbprop{ring-units-rel-free}{\vnbstmt{\begin{array}{@{}l@{}} \forall r, n. \\ \quad (\operatorname{is\text{-}ring}(r) \wedge (n \in \mathbb{N})) \Rightarrow \operatorname{rel\text{-}free}(\operatorname{ring\text{-}pow}(r, n), n, \operatorname{ring\text{-}units}(r, n)) \end{array}}}

\vnbbookfile{theorem-library/ring-pow-form}{theorem-library/ring-pow-form}
\noindent{\small ring-pow-form.scm -- wave of 2026-11-09/10 (agent CCR-1109-E2), wired 2026-11-10 from scratchpad/ccr1109e2/ring-pow-form.scm. Library: linear-algebra.}\par
\vnbprop{matmul-matscale-left-rect}{\vnbstmt{\begin{array}{@{}l@{}} \forall r, m, n, k, p, q, c. \\ \quad \operatorname{is\text{-}ring}(r) \wedge \\ \quad m \in \mathbb{N} \wedge \\ \quad n \in \mathbb{N} \wedge \\ \quad k \in \mathbb{N} \wedge \\ \quad 1 \leq n \wedge \\ \quad p \in \operatorname{mat}(m, n, \operatorname{carr}(r)) \wedge \\ \quad q \in \operatorname{mat}(n, k, \operatorname{carr}(r)) \wedge \\ \quad c \in \operatorname{carr}(r) \Rightarrow \\ \quad \quad \operatorname{matmul}(r, \operatorname{matscale}(r, c, p), q) = \operatorname{matscale}(r, c, \operatorname{matmul}(r, p, q)) \end{array}}}
\vnbprop{matmul-matscale-right-rect}{\vnbstmt{\begin{array}{@{}l@{}} \forall r, m, n, k, p, q, c. \\ \quad \operatorname{is\text{-}commutative\text{-}ring}(r) \wedge \\ \quad m \in \mathbb{N} \wedge \\ \quad n \in \mathbb{N} \wedge \\ \quad k \in \mathbb{N} \wedge \\ \quad 1 \leq n \wedge \\ \quad p \in \operatorname{mat}(m, n, \operatorname{carr}(r)) \wedge \\ \quad q \in \operatorname{mat}(n, k, \operatorname{carr}(r)) \wedge \\ \quad c \in \operatorname{carr}(r) \Rightarrow \\ \quad \quad \operatorname{matmul}(r, p, \operatorname{matscale}(r, c, q)) = \operatorname{matscale}(r, c, \operatorname{matmul}(r, p, q)) \end{array}}}
\vnbprop{transpose-matadd}{\vnbstmt{\begin{array}{@{}l@{}} \forall r, m, n, p, q. \\ \quad (\operatorname{is\text{-}ring}(r) \wedge ((m \in \mathbb{N}) \wedge ((n \in \mathbb{N}) \wedge ((p \in \operatorname{mat}(m, n, \operatorname{carr}(r))) \wedge (q \in \operatorname{mat}(m, n, \operatorname{carr}(r))))))) \Rightarrow  \\ \quad \quad \operatorname{transpose}(\operatorname{matadd}(r, p, q), m, n) = \operatorname{matadd}(r, \operatorname{transpose}(p, m, n), \operatorname{transpose}(q, m, n)) \end{array}}}
\vnbprop{transpose-matscale}{\vnbstmt{\begin{array}{@{}l@{}} \forall r, m, n, c, p. \\ \quad (\operatorname{is\text{-}ring}(r) \wedge ((m \in \mathbb{N}) \wedge ((n \in \mathbb{N}) \wedge ((c \in \operatorname{carr}(r)) \wedge (p \in \operatorname{mat}(m, n, \operatorname{carr}(r))))))) \Rightarrow  \\ \quad \quad \operatorname{transpose}(\operatorname{matscale}(r, c, p), m, n) = \operatorname{matscale}(r, c, \operatorname{transpose}(p, m, n)) \end{array}}}
\vnbprop{ring-pow-vadd-at}{\vnbstmt{\begin{array}{@{}l@{}} \forall r, n, x, y, i. \\ \quad \operatorname{is\text{-}ring}(r) \wedge \\ \quad n \in \mathbb{N} \wedge \\ \quad x \in \operatorname{vec}(\operatorname{ring\text{-}pow}(r, n)) \wedge \\ \quad y \in \operatorname{vec}(\operatorname{ring\text{-}pow}(r, n)) \wedge \\ \quad i \in \operatorname{interval}(1, n) \Rightarrow \\ \quad \quad \operatorname{vadd}(\operatorname{ring\text{-}pow}(r, n))(x, y)(i) = \operatorname{add}(r)(x(i), y(i)) \end{array}}}
\vnbprop{ring-pow-act-at}{\vnbstmt{\begin{array}{@{}l@{}} \forall r, n, a, x, i. \\ \quad (\operatorname{is\text{-}ring}(r) \wedge ((n \in \mathbb{N}) \wedge ((a \in \operatorname{carr}(r)) \wedge ((x \in \operatorname{vec}(\operatorname{ring\text{-}pow}(r, n))) \wedge (i \in \operatorname{interval}(1, n)))))) \Rightarrow  \\ \quad \quad \operatorname{act}(\operatorname{ring\text{-}pow}(r, n))(a, x)(i) = \operatorname{mul}(r)(a, x(i)) \end{array}}}
\vnbprop{ring-row-vadd}{\vnbstmt{\begin{array}{@{}l@{}} \forall r, n, x, y. \\ \quad (\operatorname{is\text{-}ring}(r) \wedge ((n \in \mathbb{N}) \wedge ((x \in \operatorname{vec}(\operatorname{ring\text{-}pow}(r, n))) \wedge (y \in \operatorname{vec}(\operatorname{ring\text{-}pow}(r, n)))))) \Rightarrow  \\ \quad \quad \operatorname{ring\text{-}row}(r, n, \operatorname{vadd}(\operatorname{ring\text{-}pow}(r, n))(x, y)) = \operatorname{matadd}(r, \operatorname{ring\text{-}row}(r, n, x), \operatorname{ring\text{-}row}(r, n, y)) \end{array}}}
\vnbprop{ring-row-act}{\vnbstmt{\begin{array}{@{}l@{}} \forall r, n, a, x. \\ \quad (\operatorname{is\text{-}ring}(r) \wedge ((n \in \mathbb{N}) \wedge ((a \in \operatorname{carr}(r)) \wedge (x \in \operatorname{vec}(\operatorname{ring\text{-}pow}(r, n)))))) \Rightarrow  \\ \quad \quad \operatorname{ring\text{-}row}(r, n, \operatorname{act}(\operatorname{ring\text{-}pow}(r, n))(a, x)) = \operatorname{matscale}(r, a, \operatorname{ring\text{-}row}(r, n, x)) \end{array}}}
\vnbprop{ring-pow-form-type}{\vnbstmt{\begin{array}{@{}l@{}} \forall r, n, a, u, v. \\ \quad \operatorname{is\text{-}commutative\text{-}ring}(r) \wedge \\ \quad n \in \mathbb{N} \wedge \\ \quad 1 \leq n \wedge \\ \quad a \in \operatorname{mat}(n, n, \operatorname{carr}(r)) \wedge \\ \quad u \in \operatorname{vec}(\operatorname{ring\text{-}pow}(r, n)) \wedge \\ \quad v \in \operatorname{vec}(\operatorname{ring\text{-}pow}(r, n)) \Rightarrow \\ \quad \quad \operatorname{entry}(\operatorname{matmul}(r, \operatorname{ring\text{-}row}(r, n, u), \operatorname{matmul}(r, a, \operatorname{transpose}(\operatorname{ring\text{-}row}(r, n, v), 1, n))), 1, 1) \\ \quad \quad \in\; \operatorname{carr}(r) \end{array}}}
\vnbprop{ring-pow-form-add-l}{\vnbstmt{\begin{array}{@{}l@{}} \forall r, n, a, x, y, z. \\ \quad \operatorname{is\text{-}commutative\text{-}ring}(r) \wedge \\ \quad n \in \mathbb{N} \wedge \\ \quad 1 \leq n \wedge \\ \quad a \in \operatorname{mat}(n, n, \operatorname{carr}(r)) \wedge \\ \quad x \in \operatorname{vec}(\operatorname{ring\text{-}pow}(r, n)) \wedge \\ \quad y \in \operatorname{vec}(\operatorname{ring\text{-}pow}(r, n)) \wedge \\ \quad z \in \operatorname{vec}(\operatorname{ring\text{-}pow}(r, n)) \Rightarrow \\ \quad \quad \operatorname{entry}(\operatorname{matmul}(r, \\ \quad \quad \quad \quad \operatorname{ring\text{-}row}(r, n, \operatorname{vadd}(\operatorname{ring\text{-}pow}(r, n))(x, y)), \\ \quad \quad \quad \quad \operatorname{matmul}(r, a, \operatorname{transpose}(\operatorname{ring\text{-}row}(r, n, z), 1, n))), \\ \quad \quad \quad 1, 1) \\ \quad \quad =\; \operatorname{add}(r)(\operatorname{entry}(\operatorname{matmul}(r, \operatorname{ring\text{-}row}(r, n, x), \operatorname{matmul}(r, a, \operatorname{transpose}(\operatorname{ring\text{-}row}(r, n, z), 1, n))), 1, 1), \\ \quad \quad \quad \operatorname{entry}(\operatorname{matmul}(r, \operatorname{ring\text{-}row}(r, n, y), \operatorname{matmul}(r, a, \operatorname{transpose}(\operatorname{ring\text{-}row}(r, n, z), 1, n))), 1, 1)) \end{array}}}
\vnbprop{ring-pow-form-add-r}{\vnbstmt{\begin{array}{@{}l@{}} \forall r, n, a, x, y, z. \\ \quad \operatorname{is\text{-}commutative\text{-}ring}(r) \wedge \\ \quad n \in \mathbb{N} \wedge \\ \quad 1 \leq n \wedge \\ \quad a \in \operatorname{mat}(n, n, \operatorname{carr}(r)) \wedge \\ \quad x \in \operatorname{vec}(\operatorname{ring\text{-}pow}(r, n)) \wedge \\ \quad y \in \operatorname{vec}(\operatorname{ring\text{-}pow}(r, n)) \wedge \\ \quad z \in \operatorname{vec}(\operatorname{ring\text{-}pow}(r, n)) \Rightarrow \\ \quad \quad \operatorname{entry}(\operatorname{matmul}(r, \operatorname{ring\text{-}row}(r, n, z), \\ \quad \quad \quad \quad \operatorname{matmul}(r, a, \operatorname{transpose}(\operatorname{ring\text{-}row}(r, n, \operatorname{vadd}(\operatorname{ring\text{-}pow}(r, n))(x, y)), 1, n))), \\ \quad \quad \quad 1, 1) \\ \quad \quad =\; \operatorname{add}(r)(\operatorname{entry}(\operatorname{matmul}(r, \operatorname{ring\text{-}row}(r, n, z), \operatorname{matmul}(r, a, \operatorname{transpose}(\operatorname{ring\text{-}row}(r, n, x), 1, n))), 1, 1), \\ \quad \quad \quad \operatorname{entry}(\operatorname{matmul}(r, \operatorname{ring\text{-}row}(r, n, z), \operatorname{matmul}(r, a, \operatorname{transpose}(\operatorname{ring\text{-}row}(r, n, y), 1, n))), 1, 1)) \end{array}}}
\vnbprop{ring-pow-form-act-l}{\vnbstmt{\begin{array}{@{}l@{}} \forall r, n, a, c, x, y. \\ \quad \operatorname{is\text{-}commutative\text{-}ring}(r) \wedge \\ \quad n \in \mathbb{N} \wedge \\ \quad 1 \leq n \wedge \\ \quad a \in \operatorname{mat}(n, n, \operatorname{carr}(r)) \wedge \\ \quad c \in \operatorname{carr}(r) \wedge \\ \quad x \in \operatorname{vec}(\operatorname{ring\text{-}pow}(r, n)) \wedge \\ \quad y \in \operatorname{vec}(\operatorname{ring\text{-}pow}(r, n)) \Rightarrow \\ \quad \quad \operatorname{entry}(\operatorname{matmul}(r, \operatorname{ring\text{-}row}(r, n, \operatorname{act}(\operatorname{ring\text{-}pow}(r, n))(c, x)), \\ \quad \quad \quad \quad \operatorname{matmul}(r, a, \operatorname{transpose}(\operatorname{ring\text{-}row}(r, n, y), 1, n))), \\ \quad \quad \quad 1, 1) \\ \quad \quad =\; \operatorname{mul}(r)(c, \\ \quad \quad \quad \operatorname{entry}(\operatorname{matmul}(r, \operatorname{ring\text{-}row}(r, n, x), \operatorname{matmul}(r, a, \operatorname{transpose}(\operatorname{ring\text{-}row}(r, n, y), 1, n))), 1, 1)) \end{array}}}
\vnbprop{ring-pow-form-act-r}{\vnbstmt{\begin{array}{@{}l@{}} \forall r, n, a, c, x, y. \\ \quad \operatorname{is\text{-}commutative\text{-}ring}(r) \wedge \\ \quad n \in \mathbb{N} \wedge \\ \quad 1 \leq n \wedge \\ \quad a \in \operatorname{mat}(n, n, \operatorname{carr}(r)) \wedge \\ \quad c \in \operatorname{carr}(r) \wedge \\ \quad x \in \operatorname{vec}(\operatorname{ring\text{-}pow}(r, n)) \wedge \\ \quad y \in \operatorname{vec}(\operatorname{ring\text{-}pow}(r, n)) \Rightarrow \\ \quad \quad \operatorname{entry}(\operatorname{matmul}(r, \operatorname{ring\text{-}row}(r, n, x), \\ \quad \quad \quad \quad \operatorname{matmul}(r, a, \operatorname{transpose}(\operatorname{ring\text{-}row}(r, n, \operatorname{act}(\operatorname{ring\text{-}pow}(r, n))(c, y)), 1, n))), \\ \quad \quad \quad 1, 1) \\ \quad \quad =\; \operatorname{mul}(r)(c, \\ \quad \quad \quad \operatorname{entry}(\operatorname{matmul}(r, \operatorname{ring\text{-}row}(r, n, x), \operatorname{matmul}(r, a, \operatorname{transpose}(\operatorname{ring\text{-}row}(r, n, y), 1, n))), 1, 1)) \end{array}}}
\vnbprop{ring-pow-matrix-form-bilinear}{\vnbstmt{\begin{array}{@{}l@{}} \forall r, n, a. \\ \quad (\operatorname{is\text{-}commutative\text{-}ring}(r) \wedge ((n \in \mathbb{N}) \wedge ((1 \leq n) \wedge (a \in \operatorname{mat}(n, n, \operatorname{carr}(r)))))) \Rightarrow  \\ \quad \quad \forall x, y, w \in \operatorname{vec}(\operatorname{ring\text{-}pow}(r, n)). \\ \quad \quad \quad \operatorname{entry}(\operatorname{matmul}(r, \\ \quad \quad \quad \quad \quad \operatorname{ring\text{-}row}(r, n, \operatorname{vadd}(\operatorname{ring\text{-}pow}(r, n))(x, y)), \\ \quad \quad \quad \quad \quad \operatorname{matmul}(r, a, \operatorname{transpose}(\operatorname{ring\text{-}row}(r, n, w), 1, n))), \\ \quad \quad \quad \quad 1, 1) \\ \quad \quad \quad =\; \operatorname{add}(r)(\operatorname{entry}(\operatorname{matmul}(r, \operatorname{ring\text{-}row}(r, n, x), \operatorname{matmul}(r, a, \operatorname{transpose}(\operatorname{ring\text{-}row}(r, n, w), 1, n))), 1, 1), \\ \quad \quad \quad \quad \operatorname{entry}(\operatorname{matmul}(r, \operatorname{ring\text{-}row}(r, n, y), \operatorname{matmul}(r, a, \operatorname{transpose}(\operatorname{ring\text{-}row}(r, n, w), 1, n))), 1, 1)) \wedge \\ \quad \quad \quad \operatorname{entry}(\operatorname{matmul}(r, \operatorname{ring\text{-}row}(r, n, w), \\ \quad \quad \quad \quad \quad \operatorname{matmul}(r, a, \operatorname{transpose}(\operatorname{ring\text{-}row}(r, n, \operatorname{vadd}(\operatorname{ring\text{-}pow}(r, n))(x, y)), 1, n))), \\ \quad \quad \quad \quad 1, 1) \\ \quad \quad \quad =\; \operatorname{add}(r)(\operatorname{entry}(\operatorname{matmul}(r, \operatorname{ring\text{-}row}(r, n, w), \operatorname{matmul}(r, a, \operatorname{transpose}(\operatorname{ring\text{-}row}(r, n, x), 1, n))), 1, 1), \\ \quad \quad \quad \quad \operatorname{entry}(\operatorname{matmul}(r, \operatorname{ring\text{-}row}(r, n, w), \operatorname{matmul}(r, a, \operatorname{transpose}(\operatorname{ring\text{-}row}(r, n, y), 1, n))), 1, 1)) \wedge \\ \quad \quad \forall \mathit{rpfa} \in \operatorname{carr}(r), x, y \in \operatorname{vec}(\operatorname{ring\text{-}pow}(r, n)). \\ \quad \quad \quad \operatorname{entry}(\operatorname{matmul}(r, \\ \quad \quad \quad \quad \quad \operatorname{ring\text{-}row}(r, n, \operatorname{act}(\operatorname{ring\text{-}pow}(r, n))(\mathit{rpfa}, x)), \\ \quad \quad \quad \quad \quad \operatorname{matmul}(r, a, \operatorname{transpose}(\operatorname{ring\text{-}row}(r, n, y), 1, n))), \\ \quad \quad \quad \quad 1, 1) \\ \quad \quad \quad =\; \operatorname{mul}(r)(\mathit{rpfa}, \\ \quad \quad \quad \quad \operatorname{entry}(\operatorname{matmul}(r, \operatorname{ring\text{-}row}(r, n, x), \operatorname{matmul}(r, a, \operatorname{transpose}(\operatorname{ring\text{-}row}(r, n, y), 1, n))), 1, 1)) \wedge \\ \quad \quad \quad \operatorname{entry}(\operatorname{matmul}(r, \operatorname{ring\text{-}row}(r, n, x), \\ \quad \quad \quad \quad \quad \operatorname{matmul}(r, a, \operatorname{transpose}(\operatorname{ring\text{-}row}(r, n, \operatorname{act}(\operatorname{ring\text{-}pow}(r, n))(\mathit{rpfa}, y)), 1, n))), \\ \quad \quad \quad \quad 1, 1) \\ \quad \quad \quad =\; \operatorname{mul}(r)(\mathit{rpfa}, \\ \quad \quad \quad \quad \operatorname{entry}(\operatorname{matmul}(r, \operatorname{ring\text{-}row}(r, n, x), \operatorname{matmul}(r, a, \operatorname{transpose}(\operatorname{ring\text{-}row}(r, n, y), 1, n))), 1, 1)) \end{array}}}

\vnbbookfile{theorem-library/zz-pow-via-ring-pow}{theorem-library/zz-pow-via-ring-pow}
\noindent{\small zz-pow-via-ring-pow.scm -- wave of 2026-11-09/10 (agent CCR-1109-E2), wired 2026-11-10 from scratchpad/ccr1109e2/zz-pow-via-ring-pow.scm. Library: linear-algebra.}\par
\vnbprop{zz-pow-is-ring-pow}{\vnbstmt{\forall n.\; (\operatorname{zz\text{-}pow}(n) \simeq \operatorname{ring\text{-}pow}(\operatorname{zz\text{-}ring}, n))}}
\vnbprop{zz-units-is-ring-units}{\vnbstmt{\forall n.\; (\operatorname{zz\text{-}units}(n) \simeq \operatorname{ring\text{-}units}(\operatorname{zz\text{-}ring}, n))}}
\vnbprop{zz-row-is-ring-row}{\vnbstmt{\forall n, x.\; (\operatorname{zz\text{-}row}(n, x) \simeq \operatorname{ring\text{-}row}(\operatorname{zz\text{-}ring}, n, x))}}
\vnbprop{zz-pow-unfold-via-ring-pow}{\vnbstmt{\begin{array}{@{}l@{}} \forall n. \\ \quad (\operatorname{zz\text{-}pow}(n) \simeq \operatorname{direct\text{-}sum}(\operatorname{zz\text{-}ring}, n, (\lambda i \in \operatorname{interval}(1, n).\; \operatorname{ring\text{-}as\text{-}module}(\operatorname{zz\text{-}ring})))) \end{array}}}
\vnbprop{zz-pow-family-via-ring-pow}{\vnbstmt{\begin{array}{@{}l@{}} \forall n \in \mathbb{N}. \\ \quad \forall i \in \operatorname{interval}(1, n). \\ \quad \quad \operatorname{is\text{-}module}((\lambda \mathit{zpi} \in \operatorname{interval}(1, n).\; \operatorname{ring\text{-}as\text{-}module}(\operatorname{zz\text{-}ring}))(i)) \wedge \\ \quad \forall i \in \operatorname{interval}(1, n). \\ \quad \quad \operatorname{scal}((\lambda \mathit{zpi} \in \operatorname{interval}(1, n).\; \operatorname{ring\text{-}as\text{-}module}(\operatorname{zz\text{-}ring}))(i)) = \operatorname{zz\text{-}ring} \end{array}}}
\vnbprop{zz-pow-is-module-via-ring-pow}{\vnbstmt{\forall n \in \mathbb{N}.\; \operatorname{is\text{-}module}(\operatorname{zz\text{-}pow}(n))}}
\vnbprop{zz-pow-scal-via-ring-pow}{\vnbstmt{\forall n \in \mathbb{N}.\; (\operatorname{scal}(\operatorname{zz\text{-}pow}(n)) = \operatorname{zz\text{-}ring})}}
\vnbprop{zz-pow-vec-in-via-ring-pow}{\vnbstmt{\begin{array}{@{}l@{}} \forall n, x. \\ \quad ((n \in \mathbb{N}) \wedge (x \in (\operatorname{interval}(1, n) \to \mathbb{Z}))) \Rightarrow (x \in \operatorname{vec}(\operatorname{zz\text{-}pow}(n))) \end{array}}}
\vnbprop{zz-pow-vec-out-via-ring-pow}{\vnbstmt{\begin{array}{@{}l@{}} \forall n, x, i. \\ \quad ((n \in \mathbb{N}) \wedge ((x \in \operatorname{vec}(\operatorname{zz\text{-}pow}(n))) \wedge (i \in \operatorname{interval}(1, n)))) \Rightarrow (x(i) \in \mathbb{Z}) \end{array}}}
\vnbprop{zz-pow-ext-via-ring-pow}{\vnbstmt{\begin{array}{@{}l@{}} \forall n, x, y. \\ \quad n \in \mathbb{N} \wedge \\ \quad x \in \operatorname{vec}(\operatorname{zz\text{-}pow}(n)) \wedge \\ \quad y \in \operatorname{vec}(\operatorname{zz\text{-}pow}(n)) \wedge \\ \quad \forall i \in \operatorname{interval}(1, n).\; (x(i) = y(i)) \Rightarrow \\ \quad \quad x = y \end{array}}}
\vnbprop{zz-pow-vzero-at-via-ring-pow}{\vnbstmt{\begin{array}{@{}l@{}} \forall n, i. \\ \quad ((n \in \mathbb{N}) \wedge (i \in \operatorname{interval}(1, n))) \Rightarrow (\operatorname{vzero}(\operatorname{zz\text{-}pow}(n))(i) = \operatorname{zero}(\operatorname{zz\text{-}ring})) \end{array}}}
\vnbprop{zz-unit-vec-via-ring-pow}{\vnbstmt{\begin{array}{@{}l@{}} \forall n \in \mathbb{N}, k. \\ \quad (\lambda t \in \operatorname{interval}(1, n).\; \operatorname{entry}(\operatorname{unitrow}(\operatorname{zz\text{-}ring}, n, k), 1, t)) \in \operatorname{vec}(\operatorname{zz\text{-}pow}(n)) \end{array}}}
\vnbprop{zz-units-type-via-ring-pow}{\vnbstmt{\forall n \in \mathbb{N}.\; (\operatorname{zz\text{-}units}(n) \in \operatorname{mat}(n, 1, \operatorname{vec}(\operatorname{zz\text{-}pow}(n))))}}
\vnbprop{zz-units-entry-via-ring-pow}{\vnbstmt{\begin{array}{@{}l@{}} \forall n, k. \\ \quad ((n \in \mathbb{N}) \wedge (k \in \operatorname{interval}(1, n))) \Rightarrow  \\ \quad \quad \operatorname{entry}(\operatorname{zz\text{-}units}(n), k, 1) = (\lambda t \in \operatorname{interval}(1, n).\; \operatorname{entry}(\operatorname{unitrow}(\operatorname{zz\text{-}ring}, n, k), 1, t)) \end{array}}}
\vnbprop{zz-unit-at-via-ring-pow}{\vnbstmt{\begin{array}{@{}l@{}} \forall n, k, i. \\ \quad ((n \in \mathbb{N}) \wedge (i \in \operatorname{interval}(1, n))) \Rightarrow  \\ \quad \quad (\lambda t \in \operatorname{interval}(1, n).\; \operatorname{entry}(\operatorname{unitrow}(\operatorname{zz\text{-}ring}, n, k), 1, t))(i) \\ \quad \quad =\; \operatorname{entry}(\operatorname{unitrow}(\operatorname{zz\text{-}ring}, n, k), 1, i) \end{array}}}
\vnbprop{zz-pow-lincomb-at-via-ring-pow}{\vnbstmt{\begin{array}{@{}l@{}} \forall n, c, i. \\ \quad ((n \in \mathbb{N}) \wedge ((c \in \operatorname{mat}(1, n, \operatorname{carr}(\operatorname{zz\text{-}ring}))) \wedge (i \in \operatorname{interval}(1, n)))) \Rightarrow  \\ \quad \quad \operatorname{lincomb}(\operatorname{zz\text{-}pow}(n), n, c, \operatorname{zz\text{-}units}(n))(i) = \operatorname{entry}(c, 1, i) \end{array}}}
\vnbprop{zz-row-type-via-ring-pow}{\vnbstmt{\begin{array}{@{}l@{}} \forall n, x. \\ \quad ((n \in \mathbb{N}) \wedge (x \in \operatorname{vec}(\operatorname{zz\text{-}pow}(n)))) \Rightarrow (\operatorname{zz\text{-}row}(n, x) \in \operatorname{mat}(1, n, \operatorname{carr}(\operatorname{zz\text{-}ring}))) \end{array}}}
\vnbprop{zz-row-entry-via-ring-pow}{\vnbstmt{\begin{array}{@{}l@{}} \forall n, x, i. \\ \quad ((n \in \mathbb{N}) \wedge ((x \in \operatorname{vec}(\operatorname{zz\text{-}pow}(n))) \wedge (i \in \operatorname{interval}(1, n)))) \Rightarrow  \\ \quad \quad \operatorname{entry}(\operatorname{zz\text{-}row}(n, x), 1, i) = x(i) \end{array}}}
\vnbprop{zz-pow-row-lincomb-via-ring-pow}{\vnbstmt{\begin{array}{@{}l@{}} \forall n, x. \\ \quad ((n \in \mathbb{N}) \wedge (x \in \operatorname{vec}(\operatorname{zz\text{-}pow}(n)))) \Rightarrow  \\ \quad \quad \operatorname{lincomb}(\operatorname{zz\text{-}pow}(n), n, \operatorname{zz\text{-}row}(n, x), \operatorname{zz\text{-}units}(n)) = x \end{array}}}
\vnbprop{zz-units-generate-via-ring-pow}{\vnbstmt{\forall n \in \mathbb{N}.\; \operatorname{generates}(\operatorname{zz\text{-}pow}(n), n, \operatorname{zz\text{-}units}(n))}}
\vnbprop{zz-units-rel-free-via-ring-pow}{\vnbstmt{\forall n \in \mathbb{N}.\; \operatorname{rel\text{-}free}(\operatorname{zz\text{-}pow}(n), n, \operatorname{zz\text{-}units}(n))}}

\vnbbookfile{theorem-library/la-z4-dsum}{theorem-library/la-z4-dsum}
\noindent{\small la-z4-dsum.scm -- dropping trivial summands from a direct sum; Z/(1) is the zero module (LA-Z4, 2026-10-24).  Vocabulary: structure-library/la-z3-vocab (DIRECT-SUM, ZZ-MOD). Prefix lzu-.}\par
\vnbprop{lzu-shift-in}{\vnbstmt{\begin{array}{@{}l@{}} \forall m, k, l, i. \\ \quad ((m \in \mathbb{N}) \wedge ((k \in \mathbb{N}) \wedge ((l \in \mathbb{N}) \wedge ((m = (k + l)) \wedge (i \in \operatorname{interval}(1, l)))))) \Rightarrow  \\ \quad \quad ((k + i) \in \operatorname{interval}(1, m)) \wedge \neg ((k + i) \leq k) \end{array}}}
\vnbprop{lzu-shift-back}{\vnbstmt{\begin{array}{@{}l@{}} \forall m, k, l, j. \\ \quad ((m \in \mathbb{N}) \wedge ((k \in \mathbb{N}) \wedge ((l \in \mathbb{N}) \wedge ((m = (k + l)) \wedge ((j \in \operatorname{interval}(1, m)) \wedge \neg (j \leq k)))))) \Rightarrow  \\ \quad \quad (\operatorname{nn\text{-}minus}(j, k) \in \operatorname{interval}(1, l)) \wedge ((k + \operatorname{nn\text{-}minus}(j, k)) = j) \end{array}}}
\vnbprop{lzu-shift-cancel}{\vnbstmt{\forall k, i.\; (((k \in \mathbb{N}) \wedge (i \in \mathbb{N})) \Rightarrow (\operatorname{nn\text{-}minus}((k + i), k) = i))}}
\vnbprop{lzu-prefix-in}{\vnbstmt{\begin{array}{@{}l@{}} \forall m, k, j. \\ \quad ((m \in \mathbb{N}) \wedge ((k \in \mathbb{N}) \wedge ((j \in \operatorname{interval}(1, m)) \wedge (j \leq k)))) \Rightarrow  \\ \quad \quad j \in \operatorname{interval}(1, k) \end{array}}}
\vnbprop{lzu-drop-vec}{\vnbstmt{\begin{array}{@{}l@{}} \forall s, m, k, l, f, g, x. \\ \quad \operatorname{is\text{-}ring}(s) \wedge \\ \quad m \in \mathbb{N} \wedge \\ \quad k \in \mathbb{N} \wedge \\ \quad l \in \mathbb{N} \wedge \\ \quad m = (k + l) \wedge \\ \quad \forall i \in \operatorname{interval}(1, m).\; \operatorname{is\text{-}module}(f(i)) \wedge \\ \quad \forall i \in \operatorname{interval}(1, m).\; (\operatorname{scal}(f(i)) = s) \wedge \\ \quad \forall i \in \operatorname{interval}(1, l).\; \operatorname{is\text{-}module}(g(i)) \wedge \\ \quad \forall i \in \operatorname{interval}(1, l).\; (\operatorname{scal}(g(i)) = s) \wedge \\ \quad \forall i \in \operatorname{interval}(1, l).\; (g(i) = f((k + i))) \wedge \\ \quad \forall i \in \operatorname{interval}(1, k), a \in \operatorname{vec}(f(i)).\; (a = \operatorname{vzero}(f(i))) \wedge \\ \quad x \in \operatorname{vec}(\operatorname{direct\text{-}sum}(s, m, f)) \Rightarrow \\ \quad \quad (\lambda \mathit{dpx} \in \operatorname{vec}(\operatorname{direct\text{-}sum}(s, m, f)).\; (\lambda \mathit{dpi} \in \operatorname{interval}(1, l).\; \operatorname{dpx}((k + \mathit{dpi}))))(x) \\ \quad \quad \in\; \operatorname{vec}(\operatorname{direct\text{-}sum}(s, l, g)) \end{array}}}
\vnbprop{lzu-drop-at}{\vnbstmt{\begin{array}{@{}l@{}} \forall s, m, k, l, f, g, x, i. \\ \quad \operatorname{is\text{-}ring}(s) \wedge \\ \quad m \in \mathbb{N} \wedge \\ \quad k \in \mathbb{N} \wedge \\ \quad l \in \mathbb{N} \wedge \\ \quad m = (k + l) \wedge \\ \quad \forall a \in \operatorname{interval}(1, m).\; \operatorname{is\text{-}module}(f(a)) \wedge \\ \quad \forall a \in \operatorname{interval}(1, m).\; (\operatorname{scal}(f(a)) = s) \wedge \\ \quad \forall a \in \operatorname{interval}(1, l).\; \operatorname{is\text{-}module}(g(a)) \wedge \\ \quad \forall a \in \operatorname{interval}(1, l).\; (\operatorname{scal}(g(a)) = s) \wedge \\ \quad \forall a \in \operatorname{interval}(1, l).\; (g(a) = f((k + a))) \wedge \\ \quad \forall a \in \operatorname{interval}(1, k), b \in \operatorname{vec}(f(a)).\; (b = \operatorname{vzero}(f(a))) \wedge \\ \quad x \in \operatorname{vec}(\operatorname{direct\text{-}sum}(s, m, f)) \wedge \\ \quad i \in \operatorname{interval}(1, l) \Rightarrow \\ \quad \quad (\lambda \mathit{dpx} \in \operatorname{vec}(\operatorname{direct\text{-}sum}(s, m, f)).\; (\lambda \mathit{dpi} \in \operatorname{interval}(1, l).\; \operatorname{dpx}((k + \mathit{dpi}))))(x)(i) \\ \quad \quad =\; x((k + i)) \end{array}}}
\vnbprop{lzu-drop-hom}{\vnbstmt{\begin{array}{@{}l@{}} \forall s, m, k, l, f, g. \\ \quad \operatorname{is\text{-}ring}(s) \wedge \\ \quad m \in \mathbb{N} \wedge \\ \quad k \in \mathbb{N} \wedge \\ \quad l \in \mathbb{N} \wedge \\ \quad m = (k + l) \wedge \\ \quad \forall i \in \operatorname{interval}(1, m).\; \operatorname{is\text{-}module}(f(i)) \wedge \\ \quad \forall i \in \operatorname{interval}(1, m).\; (\operatorname{scal}(f(i)) = s) \wedge \\ \quad \forall i \in \operatorname{interval}(1, l).\; \operatorname{is\text{-}module}(g(i)) \wedge \\ \quad \forall i \in \operatorname{interval}(1, l).\; (\operatorname{scal}(g(i)) = s) \wedge \\ \quad \forall i \in \operatorname{interval}(1, l).\; (g(i) = f((k + i))) \wedge \\ \quad \forall i \in \operatorname{interval}(1, k), x \in \operatorname{vec}(f(i)).\; (x = \operatorname{vzero}(f(i))) \Rightarrow \\ \quad \quad \operatorname{is\text{-}hom\text{-}module}(\operatorname{direct\text{-}sum}(s, m, f), \operatorname{direct\text{-}sum}(s, l, g), \\ \quad \quad \quad (\lambda \mathit{dpx} \in \operatorname{vec}(\operatorname{direct\text{-}sum}(s, m, f)).\; (\lambda \mathit{dpi} \in \operatorname{interval}(1, l).\; \operatorname{dpx}((k + \mathit{dpi}))))) \end{array}}}
\vnbprop{lzu-pad-vec}{\vnbstmt{\begin{array}{@{}l@{}} \forall s, m, k, l, f, g, y. \\ \quad \operatorname{is\text{-}ring}(s) \wedge \\ \quad m \in \mathbb{N} \wedge \\ \quad k \in \mathbb{N} \wedge \\ \quad l \in \mathbb{N} \wedge \\ \quad m = (k + l) \wedge \\ \quad \forall i \in \operatorname{interval}(1, m).\; \operatorname{is\text{-}module}(f(i)) \wedge \\ \quad \forall i \in \operatorname{interval}(1, m).\; (\operatorname{scal}(f(i)) = s) \wedge \\ \quad \forall i \in \operatorname{interval}(1, l).\; \operatorname{is\text{-}module}(g(i)) \wedge \\ \quad \forall i \in \operatorname{interval}(1, l).\; (\operatorname{scal}(g(i)) = s) \wedge \\ \quad \forall i \in \operatorname{interval}(1, l).\; (g(i) = f((k + i))) \wedge \\ \quad \forall i \in \operatorname{interval}(1, k), x \in \operatorname{vec}(f(i)).\; (x = \operatorname{vzero}(f(i))) \wedge \\ \quad y \in \operatorname{vec}(\operatorname{direct\text{-}sum}(s, l, g)) \Rightarrow \\ \quad \quad (\lambda \mathit{dpy} \in \operatorname{vec}(\operatorname{direct\text{-}sum}(s, l, g)).\; (\lambda j \in \operatorname{interval}(1, m).\; \left(\text{if }(j \leq k)\text{ then }\operatorname{vzero}(f(j))\text{ else }\operatorname{dpy}(\operatorname{nn\text{-}minus}(j, k))\right)))(y) \\ \quad \quad \in\; \operatorname{vec}(\operatorname{direct\text{-}sum}(s, m, f)) \end{array}}}
\vnbprop{lzu-drop-pad}{\vnbstmt{\begin{array}{@{}l@{}} \forall s, m, k, l, f, g, y. \\ \quad \operatorname{is\text{-}ring}(s) \wedge \\ \quad m \in \mathbb{N} \wedge \\ \quad k \in \mathbb{N} \wedge \\ \quad l \in \mathbb{N} \wedge \\ \quad m = (k + l) \wedge \\ \quad \forall i \in \operatorname{interval}(1, m).\; \operatorname{is\text{-}module}(f(i)) \wedge \\ \quad \forall i \in \operatorname{interval}(1, m).\; (\operatorname{scal}(f(i)) = s) \wedge \\ \quad \forall i \in \operatorname{interval}(1, l).\; \operatorname{is\text{-}module}(g(i)) \wedge \\ \quad \forall i \in \operatorname{interval}(1, l).\; (\operatorname{scal}(g(i)) = s) \wedge \\ \quad \forall i \in \operatorname{interval}(1, l).\; (g(i) = f((k + i))) \wedge \\ \quad \forall i \in \operatorname{interval}(1, k), x \in \operatorname{vec}(f(i)).\; (x = \operatorname{vzero}(f(i))) \wedge \\ \quad y \in \operatorname{vec}(\operatorname{direct\text{-}sum}(s, l, g)) \Rightarrow \\ \quad \quad (\lambda \mathit{dpx} \in \operatorname{vec}(\operatorname{direct\text{-}sum}(s, m, f)).\; (\lambda \mathit{dpi} \in \operatorname{interval}(1, l).\; \operatorname{dpx}((k + \mathit{dpi}))))((\lambda \mathit{dpy} \in \operatorname{vec}(\operatorname{direct\text{-}sum}(s, l, g)).\; (\lambda j \in \operatorname{interval}(1, m).\; \left(\text{if }(j \leq k)\text{ then }\operatorname{vzero}(f(j))\text{ else }\operatorname{dpy}(\operatorname{nn\text{-}minus}(j, k))\right)))(y)) \\ \quad \quad =\; y \end{array}}}
\vnbprop{lzu-drop-inj}{\vnbstmt{\begin{array}{@{}l@{}} \forall s, m, k, l, f, g, a, b. \\ \quad \operatorname{is\text{-}ring}(s) \wedge \\ \quad m \in \mathbb{N} \wedge \\ \quad k \in \mathbb{N} \wedge \\ \quad l \in \mathbb{N} \wedge \\ \quad m = (k + l) \wedge \\ \quad \forall i \in \operatorname{interval}(1, m).\; \operatorname{is\text{-}module}(f(i)) \wedge \\ \quad \forall i \in \operatorname{interval}(1, m).\; (\operatorname{scal}(f(i)) = s) \wedge \\ \quad \forall i \in \operatorname{interval}(1, l).\; \operatorname{is\text{-}module}(g(i)) \wedge \\ \quad \forall i \in \operatorname{interval}(1, l).\; (\operatorname{scal}(g(i)) = s) \wedge \\ \quad \forall i \in \operatorname{interval}(1, l).\; (g(i) = f((k + i))) \wedge \\ \quad \forall i \in \operatorname{interval}(1, k), x \in \operatorname{vec}(f(i)).\; (x = \operatorname{vzero}(f(i))) \wedge \\ \quad a \in \operatorname{vec}(\operatorname{direct\text{-}sum}(s, m, f)) \wedge \\ \quad b \in \operatorname{vec}(\operatorname{direct\text{-}sum}(s, m, f)) \wedge \\ \quad (\lambda \mathit{dpx} \in \operatorname{vec}(\operatorname{direct\text{-}sum}(s, m, f)).\; (\lambda \mathit{dpi} \in \operatorname{interval}(1, l).\; \operatorname{dpx}((k + \mathit{dpi}))))(a) \\ \quad =\; (\lambda \mathit{dpx} \in \operatorname{vec}(\operatorname{direct\text{-}sum}(s, m, f)).\; (\lambda \mathit{dpi} \in \operatorname{interval}(1, l).\; \operatorname{dpx}((k + \mathit{dpi}))))(b) \Rightarrow \\ \quad \quad a = b \end{array}}}
\vnbprop{dsum-drop-iso}{\vnbstmt{\begin{array}{@{}l@{}} \forall s, m, k, l, f, g. \\ \quad \operatorname{is\text{-}ring}(s) \wedge \\ \quad m \in \mathbb{N} \wedge \\ \quad k \in \mathbb{N} \wedge \\ \quad l \in \mathbb{N} \wedge \\ \quad m = (k + l) \wedge \\ \quad \forall i \in \operatorname{interval}(1, m).\; \operatorname{is\text{-}module}(f(i)) \wedge \\ \quad \forall i \in \operatorname{interval}(1, m).\; (\operatorname{scal}(f(i)) = s) \wedge \\ \quad \forall i \in \operatorname{interval}(1, l).\; \operatorname{is\text{-}module}(g(i)) \wedge \\ \quad \forall i \in \operatorname{interval}(1, l).\; (\operatorname{scal}(g(i)) = s) \wedge \\ \quad \forall i \in \operatorname{interval}(1, l).\; (g(i) = f((k + i))) \wedge \\ \quad \forall i \in \operatorname{interval}(1, k), x \in \operatorname{vec}(f(i)).\; (x = \operatorname{vzero}(f(i))) \Rightarrow \\ \quad \quad \operatorname{is\text{-}iso\text{-}module}(\operatorname{direct\text{-}sum}(s, m, f), \operatorname{direct\text{-}sum}(s, l, g), \\ \quad \quad \quad (\lambda \mathit{dpx} \in \operatorname{vec}(\operatorname{direct\text{-}sum}(s, m, f)).\; (\lambda \mathit{dpi} \in \operatorname{interval}(1, l).\; \operatorname{dpx}((k + \mathit{dpi}))))) \end{array}}}
\vnbprop{zz-mod-one-trivial}{\vnbstmt{\forall x \in \operatorname{vec}(\operatorname{zz\text{-}mod}(1)).\; (x = \operatorname{vzero}(\operatorname{zz\text{-}mod}(1)))}}

\vnbbookfile{theorem-library/la-z4-exist}{theorem-library/la-z4-exist}
\noindent{\small la-z4-exist.scm -- linear algebra 3.50, existence: a finitely generated abelian group is isomorphic to Z/(a\_1) + ... + Z/(a\_l) with invariant factors a\_1 | ... | a\_l, a\_i $>$= 0, a\_i /= 1 (LA-Z4, 2026-10-24).  Vocabulary: structure-library/la-z4-vocab (ZZ-POW, ZZ-UNITS, CYCLIC-SUM, IS-INVARIANT-FACTORS).  Prefix lzx-.}\par
\vnbprop{lzx-one-divides}{\vnbstmt{\forall z \in \mathbb{Z}.\; \operatorname{zz\text{-}divides}(1, z)}}
\vnbprop{lzx-ones-prefix}{\vnbstmt{\begin{array}{@{}l@{}} \forall n, m. \\ \quad n \in \mathbb{N} \wedge \\ \quad m \in \operatorname{mat}(n, n, \mathbb{Z}) \wedge \\ \quad \forall \psi \in \operatorname{interval}(1, n).\; (0 \leq \operatorname{entry}(m, \psi, \psi)) \wedge \\ \quad \forall \psi, j. \\ \quad \quad ((\psi \in \operatorname{interval}(1, n)) \wedge ((j \in \operatorname{interval}(1, n)) \wedge (\psi \leq j))) \Rightarrow  \\ \quad \quad \quad \operatorname{zz\text{-}divides}(\operatorname{entry}(m, \psi, \psi), \operatorname{entry}(m, j, j)) \Rightarrow \\ \quad \quad \exists k \in \mathbb{N}. \\ \quad \quad \quad k \leq n \wedge \\ \quad \quad \quad \forall i \in \operatorname{interval}(1, n).\; ((i \leq k) \Rightarrow (\operatorname{entry}(m, i, i) = 1)) \wedge \\ \quad \quad \quad \forall i \in \operatorname{interval}(1, n).\; (\neg (i \leq k) \Rightarrow \neg (\operatorname{entry}(m, i, i) = 1)) \end{array}}}
\vnbprop{fg-abelian-group-structure-exists}{\vnbstmt{\begin{array}{@{}l@{}} \forall e, n, u. \\ \quad (\operatorname{is\text{-}module}(e) \wedge ((\operatorname{scal}(e) = \operatorname{zz\text{-}ring}) \wedge ((n \in \mathbb{N}) \wedge ((u \in \operatorname{mat}(n, 1, \operatorname{vec}(e))) \wedge \operatorname{generates}(e, n, u))))) \Rightarrow  \\ \quad \quad \exists l \in \mathbb{N}, a, f. \\ \quad \quad \quad (l \leq n) \wedge (\operatorname{is\text{-}invariant\text{-}factors}(l, a) \wedge \operatorname{is\text{-}iso\text{-}module}(\operatorname{cyclic\text{-}sum}(l, a), e, f)) \end{array}}}

\vnbbookfile{theorem-library/la-z4-schanuel}{theorem-library/la-z4-schanuel}
\noindent{\small la-z4-schanuel.scm -- Z\textasciicircum{}n computes componentwise; Z\textasciicircum{}(n+n) as pairs (SEQ-HEAD, SEQ-SHIFT, SEQ-JOIN); Schanuel's automorphism of Z\textasciicircum{}(n+n) (LA-Z4, 2026-10-24).  The device of the uniqueness half of 3.50 (la-z4-unique).  Vocabulary: structure-library/la-z4-vocab. Prefix lzc-.}\par
\vnbprop{zz-pow-vadd-at}{\vnbstmt{\begin{array}{@{}l@{}} \forall n, x, y, i. \\ \quad ((n \in \mathbb{N}) \wedge ((x \in \operatorname{vec}(\operatorname{zz\text{-}pow}(n))) \wedge ((y \in \operatorname{vec}(\operatorname{zz\text{-}pow}(n))) \wedge (i \in \operatorname{interval}(1, n))))) \Rightarrow  \\ \quad \quad \operatorname{vadd}(\operatorname{zz\text{-}pow}(n))(x, y)(i) = \operatorname{add}(\operatorname{zz\text{-}ring})(x(i), y(i)) \end{array}}}
\vnbprop{zz-pow-vneg-at}{\vnbstmt{\begin{array}{@{}l@{}} \forall n, x, i. \\ \quad ((n \in \mathbb{N}) \wedge ((x \in \operatorname{vec}(\operatorname{zz\text{-}pow}(n))) \wedge (i \in \operatorname{interval}(1, n)))) \Rightarrow  \\ \quad \quad \operatorname{vneg}(\operatorname{zz\text{-}pow}(n))(x)(i) = \operatorname{neg}(\operatorname{zz\text{-}ring})(x(i)) \end{array}}}
\vnbprop{zz-pow-act-at}{\vnbstmt{\begin{array}{@{}l@{}} \forall n, r, x, i. \\ \quad ((n \in \mathbb{N}) \wedge ((r \in \operatorname{carr}(\operatorname{zz\text{-}ring})) \wedge ((x \in \operatorname{vec}(\operatorname{zz\text{-}pow}(n))) \wedge (i \in \operatorname{interval}(1, n))))) \Rightarrow  \\ \quad \quad \operatorname{act}(\operatorname{zz\text{-}pow}(n))(r, x)(i) = \operatorname{mul}(\operatorname{zz\text{-}ring})(r, x(i)) \end{array}}}
\vnbprop{seq-head-vec}{\vnbstmt{\begin{array}{@{}l@{}} \forall m, n, z. \\ \quad ((m \in \mathbb{N}) \wedge ((n \in \mathbb{N}) \wedge ((m = (n + n)) \wedge (z \in \operatorname{vec}(\operatorname{zz\text{-}pow}(m)))))) \Rightarrow  \\ \quad \quad \operatorname{seq\text{-}head}(n, z) \in \operatorname{vec}(\operatorname{zz\text{-}pow}(n)) \end{array}}}
\vnbprop{seq-head-at}{\vnbstmt{\begin{array}{@{}l@{}} \forall m, n, z, i. \\ \quad ((m \in \mathbb{N}) \wedge ((n \in \mathbb{N}) \wedge ((m = (n + n)) \wedge ((z \in \operatorname{vec}(\operatorname{zz\text{-}pow}(m))) \wedge (i \in \operatorname{interval}(1, n)))))) \Rightarrow  \\ \quad \quad \operatorname{seq\text{-}head}(n, z)(i) = z(i) \end{array}}}
\vnbprop{seq-shift-vec}{\vnbstmt{\begin{array}{@{}l@{}} \forall m, n, z. \\ \quad ((m \in \mathbb{N}) \wedge ((n \in \mathbb{N}) \wedge ((m = (n + n)) \wedge (z \in \operatorname{vec}(\operatorname{zz\text{-}pow}(m)))))) \Rightarrow  \\ \quad \quad \operatorname{seq\text{-}shift}(n, z) \in \operatorname{vec}(\operatorname{zz\text{-}pow}(n)) \end{array}}}
\vnbprop{seq-shift-at}{\vnbstmt{\begin{array}{@{}l@{}} \forall m, n, z, i. \\ \quad ((m \in \mathbb{N}) \wedge ((n \in \mathbb{N}) \wedge ((m = (n + n)) \wedge ((z \in \operatorname{vec}(\operatorname{zz\text{-}pow}(m))) \wedge (i \in \operatorname{interval}(1, n)))))) \Rightarrow  \\ \quad \quad \operatorname{seq\text{-}shift}(n, z)(i) = z((n + i)) \end{array}}}
\vnbprop{seq-join-vec}{\vnbstmt{\begin{array}{@{}l@{}} \forall m, n, d, l. \\ \quad ((m \in \mathbb{N}) \wedge ((n \in \mathbb{N}) \wedge ((m = (n + n)) \wedge ((d \in \operatorname{vec}(\operatorname{zz\text{-}pow}(n))) \wedge (l \in \operatorname{vec}(\operatorname{zz\text{-}pow}(n))))))) \Rightarrow  \\ \quad \quad \operatorname{seq\text{-}join}(m, n, d, l) \in \operatorname{vec}(\operatorname{zz\text{-}pow}(m)) \end{array}}}
\vnbprop{seq-join-lo}{\vnbstmt{\begin{array}{@{}l@{}} \forall m, n, d, l, j. \\ \quad m \in \mathbb{N} \wedge \\ \quad n \in \mathbb{N} \wedge \\ \quad m = (n + n) \wedge \\ \quad d \in \operatorname{vec}(\operatorname{zz\text{-}pow}(n)) \wedge \\ \quad l \in \operatorname{vec}(\operatorname{zz\text{-}pow}(n)) \wedge \\ \quad j \in \operatorname{interval}(1, m) \wedge \\ \quad j \leq n \Rightarrow \\ \quad \quad \operatorname{seq\text{-}join}(m, n, d, l)(j) = d(j) \end{array}}}
\vnbprop{seq-join-hi}{\vnbstmt{\begin{array}{@{}l@{}} \forall m, n, d, l, j. \\ \quad m \in \mathbb{N} \wedge \\ \quad n \in \mathbb{N} \wedge \\ \quad m = (n + n) \wedge \\ \quad d \in \operatorname{vec}(\operatorname{zz\text{-}pow}(n)) \wedge \\ \quad l \in \operatorname{vec}(\operatorname{zz\text{-}pow}(n)) \wedge \\ \quad j \in \operatorname{interval}(1, m) \wedge \\ \quad \neg (j \leq n) \Rightarrow \\ \quad \quad \operatorname{seq\text{-}join}(m, n, d, l)(j) = l(\operatorname{nn\text{-}minus}(j, n)) \end{array}}}
\vnbprop{seq-head-join}{\vnbstmt{\begin{array}{@{}l@{}} \forall m, n, d, l. \\ \quad ((m \in \mathbb{N}) \wedge ((n \in \mathbb{N}) \wedge ((m = (n + n)) \wedge ((d \in \operatorname{vec}(\operatorname{zz\text{-}pow}(n))) \wedge (l \in \operatorname{vec}(\operatorname{zz\text{-}pow}(n))))))) \Rightarrow  \\ \quad \quad \operatorname{seq\text{-}head}(n, \operatorname{seq\text{-}join}(m, n, d, l)) = d \end{array}}}
\vnbprop{seq-shift-join}{\vnbstmt{\begin{array}{@{}l@{}} \forall m, n, d, l. \\ \quad ((m \in \mathbb{N}) \wedge ((n \in \mathbb{N}) \wedge ((m = (n + n)) \wedge ((d \in \operatorname{vec}(\operatorname{zz\text{-}pow}(n))) \wedge (l \in \operatorname{vec}(\operatorname{zz\text{-}pow}(n))))))) \Rightarrow  \\ \quad \quad \operatorname{seq\text{-}shift}(n, \operatorname{seq\text{-}join}(m, n, d, l)) = l \end{array}}}
\vnbprop{seq-join-ext}{\vnbstmt{\begin{array}{@{}l@{}} \forall m, n, z. \\ \quad ((m \in \mathbb{N}) \wedge ((n \in \mathbb{N}) \wedge ((m = (n + n)) \wedge (z \in \operatorname{vec}(\operatorname{zz\text{-}pow}(m)))))) \Rightarrow  \\ \quad \quad \operatorname{seq\text{-}join}(m, n, \operatorname{seq\text{-}head}(n, z), \operatorname{seq\text{-}shift}(n, z)) = z \end{array}}}
\vnbprop{seq-split-ext}{\vnbstmt{\begin{array}{@{}l@{}} \forall m, n, v, w. \\ \quad m \in \mathbb{N} \wedge \\ \quad n \in \mathbb{N} \wedge \\ \quad m = (n + n) \wedge \\ \quad v \in \operatorname{vec}(\operatorname{zz\text{-}pow}(m)) \wedge \\ \quad w \in \operatorname{vec}(\operatorname{zz\text{-}pow}(m)) \wedge \\ \quad \operatorname{seq\text{-}head}(n, v) = \operatorname{seq\text{-}head}(n, w) \wedge \\ \quad \operatorname{seq\text{-}shift}(n, v) = \operatorname{seq\text{-}shift}(n, w) \Rightarrow \\ \quad \quad v = w \end{array}}}
\vnbprop{seq-head-vadd}{\vnbstmt{\begin{array}{@{}l@{}} \forall m, n, x, y. \\ \quad ((m \in \mathbb{N}) \wedge ((n \in \mathbb{N}) \wedge ((m = (n + n)) \wedge ((x \in \operatorname{vec}(\operatorname{zz\text{-}pow}(m))) \wedge (y \in \operatorname{vec}(\operatorname{zz\text{-}pow}(m))))))) \Rightarrow  \\ \quad \quad \operatorname{seq\text{-}head}(n, \operatorname{vadd}(\operatorname{zz\text{-}pow}(m))(x, y)) = \operatorname{vadd}(\operatorname{zz\text{-}pow}(n))(\operatorname{seq\text{-}head}(n, x), \operatorname{seq\text{-}head}(n, y)) \end{array}}}
\vnbprop{seq-head-vneg}{\vnbstmt{\begin{array}{@{}l@{}} \forall m, n, x. \\ \quad ((m \in \mathbb{N}) \wedge ((n \in \mathbb{N}) \wedge ((m = (n + n)) \wedge (x \in \operatorname{vec}(\operatorname{zz\text{-}pow}(m)))))) \Rightarrow  \\ \quad \quad \operatorname{seq\text{-}head}(n, \operatorname{vneg}(\operatorname{zz\text{-}pow}(m))(x)) = \operatorname{vneg}(\operatorname{zz\text{-}pow}(n))(\operatorname{seq\text{-}head}(n, x)) \end{array}}}
\vnbprop{seq-head-act}{\vnbstmt{\begin{array}{@{}l@{}} \forall m, n, r, x. \\ \quad ((m \in \mathbb{N}) \wedge ((n \in \mathbb{N}) \wedge ((m = (n + n)) \wedge ((r \in \operatorname{carr}(\operatorname{zz\text{-}ring})) \wedge (x \in \operatorname{vec}(\operatorname{zz\text{-}pow}(m))))))) \Rightarrow  \\ \quad \quad \operatorname{seq\text{-}head}(n, \operatorname{act}(\operatorname{zz\text{-}pow}(m))(r, x)) = \operatorname{act}(\operatorname{zz\text{-}pow}(n))(r, \operatorname{seq\text{-}head}(n, x)) \end{array}}}
\vnbprop{seq-head-vzero}{\vnbstmt{\begin{array}{@{}l@{}} \forall m, n. \\ \quad ((m \in \mathbb{N}) \wedge ((n \in \mathbb{N}) \wedge (m = (n + n)))) \Rightarrow  \\ \quad \quad \operatorname{seq\text{-}head}(n, \operatorname{vzero}(\operatorname{zz\text{-}pow}(m))) = \operatorname{vzero}(\operatorname{zz\text{-}pow}(n)) \end{array}}}
\vnbprop{seq-shift-vadd}{\vnbstmt{\begin{array}{@{}l@{}} \forall m, n, x, y. \\ \quad ((m \in \mathbb{N}) \wedge ((n \in \mathbb{N}) \wedge ((m = (n + n)) \wedge ((x \in \operatorname{vec}(\operatorname{zz\text{-}pow}(m))) \wedge (y \in \operatorname{vec}(\operatorname{zz\text{-}pow}(m))))))) \Rightarrow  \\ \quad \quad \operatorname{seq\text{-}shift}(n, \operatorname{vadd}(\operatorname{zz\text{-}pow}(m))(x, y)) = \operatorname{vadd}(\operatorname{zz\text{-}pow}(n))(\operatorname{seq\text{-}shift}(n, x), \operatorname{seq\text{-}shift}(n, y)) \end{array}}}
\vnbprop{seq-shift-vneg}{\vnbstmt{\begin{array}{@{}l@{}} \forall m, n, x. \\ \quad ((m \in \mathbb{N}) \wedge ((n \in \mathbb{N}) \wedge ((m = (n + n)) \wedge (x \in \operatorname{vec}(\operatorname{zz\text{-}pow}(m)))))) \Rightarrow  \\ \quad \quad \operatorname{seq\text{-}shift}(n, \operatorname{vneg}(\operatorname{zz\text{-}pow}(m))(x)) = \operatorname{vneg}(\operatorname{zz\text{-}pow}(n))(\operatorname{seq\text{-}shift}(n, x)) \end{array}}}
\vnbprop{seq-shift-act}{\vnbstmt{\begin{array}{@{}l@{}} \forall m, n, r, x. \\ \quad ((m \in \mathbb{N}) \wedge ((n \in \mathbb{N}) \wedge ((m = (n + n)) \wedge ((r \in \operatorname{carr}(\operatorname{zz\text{-}ring})) \wedge (x \in \operatorname{vec}(\operatorname{zz\text{-}pow}(m))))))) \Rightarrow  \\ \quad \quad \operatorname{seq\text{-}shift}(n, \operatorname{act}(\operatorname{zz\text{-}pow}(m))(r, x)) = \operatorname{act}(\operatorname{zz\text{-}pow}(n))(r, \operatorname{seq\text{-}shift}(n, x)) \end{array}}}
\vnbprop{seq-shift-vzero}{\vnbstmt{\begin{array}{@{}l@{}} \forall m, n. \\ \quad ((m \in \mathbb{N}) \wedge ((n \in \mathbb{N}) \wedge (m = (n + n)))) \Rightarrow  \\ \quad \quad \operatorname{seq\text{-}shift}(n, \operatorname{vzero}(\operatorname{zz\text{-}pow}(m))) = \operatorname{vzero}(\operatorname{zz\text{-}pow}(n)) \end{array}}}
\vnbprop{schanuel-at}{\vnbstmt{\begin{array}{@{}l@{}} \forall m, n, a, b, z. \\ \quad m \in \mathbb{N} \wedge \\ \quad n \in \mathbb{N} \wedge \\ \quad m = (n + n) \wedge \\ \quad \operatorname{is\text{-}hom\text{-}module}(\operatorname{zz\text{-}pow}(n), \operatorname{zz\text{-}pow}(n), a) \wedge \\ \quad \operatorname{is\text{-}hom\text{-}module}(\operatorname{zz\text{-}pow}(n), \operatorname{zz\text{-}pow}(n), b) \wedge \\ \quad z \in \operatorname{vec}(\operatorname{zz\text{-}pow}(m)) \Rightarrow \\ \quad \quad \operatorname{schanuel\text{-}map}(m, n, a, b)(z) \\ \quad \quad =\; \operatorname{seq\text{-}join}(m, n, \\ \quad \quad \quad \operatorname{vadd}(\operatorname{zz\text{-}pow}(n))(\operatorname{seq\text{-}shift}(n, z), \operatorname{vneg}(\operatorname{zz\text{-}pow}(n))(b(\operatorname{seq\text{-}head}(n, z)))), \\ \quad \quad \quad \operatorname{vadd}(\operatorname{zz\text{-}pow}(n))(\operatorname{seq\text{-}head}(n, z), \\ \quad \quad \quad \quad a(\operatorname{vadd}(\operatorname{zz\text{-}pow}(n))(\operatorname{seq\text{-}shift}(n, z), \operatorname{vneg}(\operatorname{zz\text{-}pow}(n))(b(\operatorname{seq\text{-}head}(n, z))))))) \end{array}}}
\vnbprop{schanuel-vec}{\vnbstmt{\begin{array}{@{}l@{}} \forall m, n, a, b, z. \\ \quad m \in \mathbb{N} \wedge \\ \quad n \in \mathbb{N} \wedge \\ \quad m = (n + n) \wedge \\ \quad \operatorname{is\text{-}hom\text{-}module}(\operatorname{zz\text{-}pow}(n), \operatorname{zz\text{-}pow}(n), a) \wedge \\ \quad \operatorname{is\text{-}hom\text{-}module}(\operatorname{zz\text{-}pow}(n), \operatorname{zz\text{-}pow}(n), b) \wedge \\ \quad z \in \operatorname{vec}(\operatorname{zz\text{-}pow}(m)) \Rightarrow \\ \quad \quad \operatorname{schanuel\text{-}map}(m, n, a, b)(z) \in \operatorname{vec}(\operatorname{zz\text{-}pow}(m)) \end{array}}}
\vnbprop{schanuel-head}{\vnbstmt{\begin{array}{@{}l@{}} \forall m, n, a, b, z. \\ \quad m \in \mathbb{N} \wedge \\ \quad n \in \mathbb{N} \wedge \\ \quad m = (n + n) \wedge \\ \quad \operatorname{is\text{-}hom\text{-}module}(\operatorname{zz\text{-}pow}(n), \operatorname{zz\text{-}pow}(n), a) \wedge \\ \quad \operatorname{is\text{-}hom\text{-}module}(\operatorname{zz\text{-}pow}(n), \operatorname{zz\text{-}pow}(n), b) \wedge \\ \quad z \in \operatorname{vec}(\operatorname{zz\text{-}pow}(m)) \Rightarrow \\ \quad \quad \operatorname{seq\text{-}head}(n, \operatorname{schanuel\text{-}map}(m, n, a, b)(z)) \\ \quad \quad =\; \operatorname{vadd}(\operatorname{zz\text{-}pow}(n))(\operatorname{seq\text{-}shift}(n, z), \operatorname{vneg}(\operatorname{zz\text{-}pow}(n))(b(\operatorname{seq\text{-}head}(n, z)))) \end{array}}}
\vnbprop{schanuel-shift}{\vnbstmt{\begin{array}{@{}l@{}} \forall m, n, a, b, z. \\ \quad m \in \mathbb{N} \wedge \\ \quad n \in \mathbb{N} \wedge \\ \quad m = (n + n) \wedge \\ \quad \operatorname{is\text{-}hom\text{-}module}(\operatorname{zz\text{-}pow}(n), \operatorname{zz\text{-}pow}(n), a) \wedge \\ \quad \operatorname{is\text{-}hom\text{-}module}(\operatorname{zz\text{-}pow}(n), \operatorname{zz\text{-}pow}(n), b) \wedge \\ \quad z \in \operatorname{vec}(\operatorname{zz\text{-}pow}(m)) \Rightarrow \\ \quad \quad \operatorname{seq\text{-}shift}(n, \operatorname{schanuel\text{-}map}(m, n, a, b)(z)) \\ \quad \quad =\; \operatorname{vadd}(\operatorname{zz\text{-}pow}(n))(\operatorname{seq\text{-}head}(n, z), \\ \quad \quad \quad a(\operatorname{vadd}(\operatorname{zz\text{-}pow}(n))(\operatorname{seq\text{-}shift}(n, z), \operatorname{vneg}(\operatorname{zz\text{-}pow}(n))(b(\operatorname{seq\text{-}head}(n, z)))))) \end{array}}}
\vnbprop{schanuel-type}{\vnbstmt{\begin{array}{@{}l@{}} \forall m, n, a, b. \\ \quad m \in \mathbb{N} \wedge \\ \quad n \in \mathbb{N} \wedge \\ \quad m = (n + n) \wedge \\ \quad \operatorname{is\text{-}hom\text{-}module}(\operatorname{zz\text{-}pow}(n), \operatorname{zz\text{-}pow}(n), a) \wedge \\ \quad \operatorname{is\text{-}hom\text{-}module}(\operatorname{zz\text{-}pow}(n), \operatorname{zz\text{-}pow}(n), b) \Rightarrow \\ \quad \quad \operatorname{schanuel\text{-}map}(m, n, a, b) \in (\operatorname{vec}(\operatorname{zz\text{-}pow}(m)) \to \operatorname{vec}(\operatorname{zz\text{-}pow}(m))) \end{array}}}
\vnbprop{schanuel-inverse}{\vnbstmt{\begin{array}{@{}l@{}} \forall m, n, a, b, z. \\ \quad m \in \mathbb{N} \wedge \\ \quad n \in \mathbb{N} \wedge \\ \quad m = (n + n) \wedge \\ \quad \operatorname{is\text{-}hom\text{-}module}(\operatorname{zz\text{-}pow}(n), \operatorname{zz\text{-}pow}(n), a) \wedge \\ \quad \operatorname{is\text{-}hom\text{-}module}(\operatorname{zz\text{-}pow}(n), \operatorname{zz\text{-}pow}(n), b) \wedge \\ \quad z \in \operatorname{vec}(\operatorname{zz\text{-}pow}(m)) \Rightarrow \\ \quad \quad \operatorname{schanuel\text{-}map}(m, n, b, a)(\operatorname{schanuel\text{-}map}(m, n, a, b)(z)) = z \end{array}}}
\vnbprop{schanuel-hom}{\vnbstmt{\begin{array}{@{}l@{}} \forall m, n, a, b. \\ \quad m \in \mathbb{N} \wedge \\ \quad n \in \mathbb{N} \wedge \\ \quad m = (n + n) \wedge \\ \quad \operatorname{is\text{-}hom\text{-}module}(\operatorname{zz\text{-}pow}(n), \operatorname{zz\text{-}pow}(n), a) \wedge \\ \quad \operatorname{is\text{-}hom\text{-}module}(\operatorname{zz\text{-}pow}(n), \operatorname{zz\text{-}pow}(n), b) \Rightarrow \\ \quad \quad \operatorname{is\text{-}hom\text{-}module}(\operatorname{zz\text{-}pow}(m), \operatorname{zz\text{-}pow}(m), \operatorname{schanuel\text{-}map}(m, n, a, b)) \end{array}}}
\vnbprop{schanuel-iso}{\vnbstmt{\begin{array}{@{}l@{}} \forall m, n, a, b. \\ \quad m \in \mathbb{N} \wedge \\ \quad n \in \mathbb{N} \wedge \\ \quad m = (n + n) \wedge \\ \quad \operatorname{is\text{-}hom\text{-}module}(\operatorname{zz\text{-}pow}(n), \operatorname{zz\text{-}pow}(n), a) \wedge \\ \quad \operatorname{is\text{-}hom\text{-}module}(\operatorname{zz\text{-}pow}(n), \operatorname{zz\text{-}pow}(n), b) \Rightarrow \\ \quad \quad \operatorname{is\text{-}iso\text{-}module}(\operatorname{zz\text{-}pow}(m), \operatorname{zz\text{-}pow}(m), \operatorname{schanuel\text{-}map}(m, n, a, b)) \end{array}}}

\vnbbookfile{theorem-library/la-z4-unique}{theorem-library/la-z4-unique}
\noindent{\small la-z4-unique.scm -- linear algebra 3.50, uniqueness: the invariant factors of a finitely generated abelian group are determined by it (LA-Z4, 2026-10-24).  Vocabulary: structure-library/la-z4-vocab.  Prefix lzy-.}\par
\vnbprop{free-module-lift}{\vnbstmt{\begin{array}{@{}l@{}} \forall d, n, s, y, x, i, a. \\ \quad \operatorname{is\text{-}module}(d) \wedge \\ \quad n \in \mathbb{N} \wedge \\ \quad s \in \operatorname{mat}(n, 1, \operatorname{vec}(d)) \wedge \\ \quad \operatorname{generates}(d, n, s) \wedge \\ \quad \operatorname{rel\text{-}free}(d, n, s) \wedge \\ \quad \operatorname{is\text{-}hom\text{-}module}(y, x, i) \wedge \\ \quad \forall b \in \operatorname{vec}(x).\; \exists c \in \operatorname{vec}(y).\; (i(c) = b) \wedge \\ \quad \operatorname{is\text{-}hom\text{-}module}(d, x, a) \wedge \\ \quad \operatorname{scal}(y) = \operatorname{scal}(d) \Rightarrow \\ \quad \quad \exists m.\; (\operatorname{is\text{-}hom\text{-}module}(d, y, m) \wedge \forall b \in \operatorname{vec}(d).\; (i(m(b)) = a(b))) \end{array}}}
\vnbprop{hom-matact}{\vnbstmt{\begin{array}{@{}l@{}} \forall a, g, m, n, d, u. \\ \quad \operatorname{is\text{-}hom\text{-}module}(a, a, g) \wedge \\ \quad m \in \mathbb{N} \wedge \\ \quad n \in \mathbb{N} \wedge \\ \quad 1 \leq n \wedge \\ \quad d \in \operatorname{mat}(m, n, \operatorname{carr}(\operatorname{scal}(a))) \wedge \\ \quad u \in \operatorname{mat}(n, 1, \operatorname{vec}(a)) \Rightarrow \\ \quad \quad \operatorname{matact}(a, d, \operatorname{map\text{-}seq}(g, n, u)) = \operatorname{map\text{-}seq}(g, m, \operatorname{matact}(a, d, u)) \end{array}}}
\vnbprop{diag-mat-type}{\vnbstmt{\begin{array}{@{}l@{}} \forall n, f. \\ \quad ((n \in \mathbb{N}) \wedge (f \in (\operatorname{interval}(1, n) \to \mathbb{Z}))) \Rightarrow (\operatorname{diag\text{-}mat}(n, f) \in \operatorname{mat}(n, n, \mathbb{Z})) \end{array}}}
\vnbprop{diag-mat-entry}{\vnbstmt{\begin{array}{@{}l@{}} \forall n, f, i, j. \\ \quad ((n \in \mathbb{N}) \wedge ((f \in (\operatorname{interval}(1, n) \to \mathbb{Z})) \wedge ((i \in \operatorname{interval}(1, n)) \wedge (j \in \operatorname{interval}(1, n))))) \Rightarrow  \\ \quad \quad \operatorname{entry}(\operatorname{diag\text{-}mat}(n, f), i, j) = \operatorname{mul}(\operatorname{zz\text{-}ring})(f(i), \operatorname{entry}(\operatorname{identmat}(\operatorname{zz\text{-}ring}, n), i, j)) \end{array}}}
\vnbprop{diag-mat-diag}{\vnbstmt{\begin{array}{@{}l@{}} \forall n, f, i. \\ \quad ((n \in \mathbb{N}) \wedge ((f \in (\operatorname{interval}(1, n) \to \mathbb{Z})) \wedge (i \in \operatorname{interval}(1, n)))) \Rightarrow  \\ \quad \quad \operatorname{entry}(\operatorname{diag\text{-}mat}(n, f), i, i) = f(i) \end{array}}}
\vnbprop{diag-mat-diagonal}{\vnbstmt{\begin{array}{@{}l@{}} \forall n, f. \\ \quad ((n \in \mathbb{N}) \wedge (f \in (\operatorname{interval}(1, n) \to \mathbb{Z}))) \Rightarrow  \\ \quad \quad \operatorname{is\text{-}diagonal}(\operatorname{zz\text{-}ring}, n, n, \operatorname{diag\text{-}mat}(n, f)) \end{array}}}
\vnbprop{pad-ones-lo}{\vnbstmt{\begin{array}{@{}l@{}} \forall m, k, l, f, j. \\ \quad m \in \mathbb{N} \wedge \\ \quad k \in \mathbb{N} \wedge \\ \quad l \in \mathbb{N} \wedge \\ \quad m = (k + l) \wedge \\ \quad f \in (\operatorname{interval}(1, l) \to \mathbb{Z}) \wedge \\ \quad j \in \operatorname{interval}(1, m) \wedge \\ \quad j \leq k \Rightarrow \\ \quad \quad \operatorname{pad\text{-}ones}(m, k, f)(j) = 1 \end{array}}}
\vnbprop{pad-ones-hi}{\vnbstmt{\begin{array}{@{}l@{}} \forall m, k, l, f, j. \\ \quad m \in \mathbb{N} \wedge \\ \quad k \in \mathbb{N} \wedge \\ \quad l \in \mathbb{N} \wedge \\ \quad m = (k + l) \wedge \\ \quad f \in (\operatorname{interval}(1, l) \to \mathbb{Z}) \wedge \\ \quad j \in \operatorname{interval}(1, m) \wedge \\ \quad \neg (j \leq k) \Rightarrow \\ \quad \quad \operatorname{pad\text{-}ones}(m, k, f)(j) = f(\operatorname{nn\text{-}minus}(j, k)) \end{array}}}
\vnbprop{pad-ones-type}{\vnbstmt{\begin{array}{@{}l@{}} \forall m, k, l, f. \\ \quad ((m \in \mathbb{N}) \wedge ((k \in \mathbb{N}) \wedge ((l \in \mathbb{N}) \wedge ((m = (k + l)) \wedge (f \in (\operatorname{interval}(1, l) \to \mathbb{Z})))))) \Rightarrow  \\ \quad \quad \operatorname{pad\text{-}ones}(m, k, f) \in (\operatorname{interval}(1, m) \to \mathbb{Z}) \end{array}}}
\vnbprop{pad-ones-nonneg}{\vnbstmt{\begin{array}{@{}l@{}} \forall m, k, l, f.; \forall j \in \operatorname{interval}(1, m). \\ \quad m \in \mathbb{N} \wedge \\ \quad k \in \mathbb{N} \wedge \\ \quad l \in \mathbb{N} \wedge \\ \quad m = (k + l) \wedge \\ \quad f \in (\operatorname{interval}(1, l) \to \mathbb{Z}) \wedge \\ \quad \forall i \in \operatorname{interval}(1, l).\; (0 \leq f(i)) \Rightarrow \\ \quad \quad 0 \leq \operatorname{pad\text{-}ones}(m, k, f)(j) \end{array}}}
\vnbprop{pad-ones-chain}{\vnbstmt{\begin{array}{@{}l@{}} \forall m, k, l, f, i, j. \\ \quad m \in \mathbb{N} \wedge \\ \quad k \in \mathbb{N} \wedge \\ \quad l \in \mathbb{N} \wedge \\ \quad m = (k + l) \wedge \\ \quad f \in (\operatorname{interval}(1, l) \to \mathbb{Z}) \wedge \\ \quad \forall \mathit{ifi}, \mathit{ifj}. \\ \quad \quad ((\mathit{ifi} \in \operatorname{interval}(1, l)) \wedge ((\mathit{ifj} \in \operatorname{interval}(1, l)) \wedge (\mathit{ifi} \leq \mathit{ifj}))) \Rightarrow  \\ \quad \quad \quad \operatorname{zz\text{-}divides}(f(\mathit{ifi}), f(\mathit{ifj})) \wedge \\ \quad i \in \operatorname{interval}(1, m) \wedge \\ \quad j \in \operatorname{interval}(1, m) \wedge \\ \quad i \leq j \Rightarrow \\ \quad \quad \operatorname{zz\text{-}divides}(\operatorname{pad\text{-}ones}(m, k, f)(i), \operatorname{pad\text{-}ones}(m, k, f)(j)) \end{array}}}
\vnbprop{zz-pow-coord}{\vnbstmt{\begin{array}{@{}l@{}} \forall n, x, j. \\ \quad ((n \in \mathbb{N}) \wedge ((x \in \operatorname{vec}(\operatorname{zz\text{-}pow}(n))) \wedge (j \in \operatorname{interval}(1, n)))) \Rightarrow  \\ \quad \quad \operatorname{entry}((\lambda c \in \operatorname{mat}(1, n, \operatorname{carr}(\operatorname{scal}(\operatorname{zz\text{-}pow}(n)))).\; \operatorname{lincomb}(\operatorname{zz\text{-}pow}(n), n, c, \operatorname{zz\text{-}units}(n)))^{-1}(x), \\ \quad \quad \quad 1, j) \\ \quad \quad =\; x(j) \end{array}}}
\vnbprop{zz-pow-diag-span-iff}{\vnbstmt{\begin{array}{@{}l@{}} \forall n, f, x. \\ \quad ((n \in \mathbb{N}) \wedge ((f \in (\operatorname{interval}(1, n) \to \mathbb{Z})) \wedge (x \in \operatorname{vec}(\operatorname{zz\text{-}pow}(n))))) \Rightarrow  \\ \quad \quad x \in \operatorname{span}(\operatorname{zz\text{-}pow}(n), n, \operatorname{matact}(\operatorname{zz\text{-}pow}(n), \operatorname{diag\text{-}mat}(n, f), \operatorname{zz\text{-}units}(n))) \\ \quad \quad \Leftrightarrow\; \forall j \in \operatorname{interval}(1, n).\; (x(j) \in \operatorname{principal\text{-}ideal}(\operatorname{zz\text{-}ring}, f(j))) \end{array}}}
\vnbprop{zz-pow-pad-span-iff}{\vnbstmt{\begin{array}{@{}l@{}} \forall m, n, f, z. \\ \quad ((m \in \mathbb{N}) \wedge ((n \in \mathbb{N}) \wedge ((m = (n + n)) \wedge ((f \in (\operatorname{interval}(1, n) \to \mathbb{Z})) \wedge (z \in \operatorname{vec}(\operatorname{zz\text{-}pow}(m))))))) \Rightarrow  \\ \quad \quad z \in \operatorname{span}(\operatorname{zz\text{-}pow}(m), m, \operatorname{matact}(\operatorname{zz\text{-}pow}(m), \operatorname{diag\text{-}mat}(m, \operatorname{pad\text{-}ones}(m, n, f)), \operatorname{zz\text{-}units}(m))) \\ \quad \quad \Leftrightarrow\; \operatorname{seq\text{-}shift}(n, z) \in \operatorname{span}(\operatorname{zz\text{-}pow}(n), n, \operatorname{matact}(\operatorname{zz\text{-}pow}(n), \operatorname{diag\text{-}mat}(n, f), \operatorname{zz\text{-}units}(n))) \end{array}}}
\vnbprop{schanuel-kernel-map}{\vnbstmt{\begin{array}{@{}l@{}} \forall m, n, l, e, a, b, h, g.; \forall z \in \operatorname{vec}(\operatorname{zz\text{-}pow}(m)). \\ \quad m \in \mathbb{N} \wedge \\ \quad n \in \mathbb{N} \wedge \\ \quad m = (n + n) \wedge \\ \quad l \in (\operatorname{interval}(1, n) \to \mathbb{Z}) \wedge \\ \quad e \in (\operatorname{interval}(1, n) \to \mathbb{Z}) \wedge \\ \quad \operatorname{is\text{-}hom\text{-}module}(\operatorname{zz\text{-}pow}(n), \operatorname{zz\text{-}pow}(n), a) \wedge \\ \quad \operatorname{is\text{-}hom\text{-}module}(\operatorname{zz\text{-}pow}(n), \operatorname{zz\text{-}pow}(n), b) \wedge \\ \quad \operatorname{is\text{-}hom\text{-}module}(\operatorname{direct\text{-}sum}(\operatorname{zz\text{-}ring}, n, (\lambda i \in \operatorname{interval}(1, n).\; \operatorname{zz\text{-}mod}(\operatorname{entry}(\operatorname{diag\text{-}mat}(n, l), i, i)))), \\ \quad \quad \operatorname{direct\text{-}sum}(\operatorname{zz\text{-}ring}, n, (\lambda i \in \operatorname{interval}(1, n).\; \operatorname{zz\text{-}mod}(\operatorname{entry}(\operatorname{diag\text{-}mat}(n, e), i, i)))), \\ \quad \quad h) \wedge \\ \quad g \\ \quad \in\; (\operatorname{vec}(\operatorname{direct\text{-}sum}(\operatorname{zz\text{-}ring}, n, (\lambda i \in \operatorname{interval}(1, n).\; \operatorname{zz\text{-}mod}(\operatorname{entry}(\operatorname{diag\text{-}mat}(n, e), i, i))))) \to \operatorname{vec}(\operatorname{direct\text{-}sum}(\operatorname{zz\text{-}ring}, n, (\lambda i \in \operatorname{interval}(1, n).\; \operatorname{zz\text{-}mod}(\operatorname{entry}(\operatorname{diag\text{-}mat}(n, l), i, i)))))) \wedge \\ \quad \forall x \in \operatorname{vec}(\operatorname{zz\text{-}pow}(n)). \\ \quad \quad (\lambda \mathit{phx} \in \operatorname{vec}(\operatorname{zz\text{-}pow}(n)).\; (\lambda j \in \operatorname{interval}(1, n).\; \operatorname{class}(\operatorname{qmod\text{-}setoid}(\operatorname{ring\text{-}as\text{-}module}(\operatorname{zz\text{-}ring}), \operatorname{principal\text{-}ideal}(\operatorname{zz\text{-}ring}, \operatorname{entry}(\operatorname{diag\text{-}mat}(n, e), j, j))), \operatorname{entry}((\lambda c \in \operatorname{mat}(1, n, \operatorname{carr}(\operatorname{scal}(\operatorname{zz\text{-}pow}(n)))).\; \operatorname{lincomb}(\operatorname{zz\text{-}pow}(n), n, c, \operatorname{zz\text{-}units}(n)))^{-1}(\mathit{phx}), 1, j))))(a(x)) \\ \quad \quad =\; h((\lambda \mathit{phx} \in \operatorname{vec}(\operatorname{zz\text{-}pow}(n)).\; (\lambda j \in \operatorname{interval}(1, n).\; \operatorname{class}(\operatorname{qmod\text{-}setoid}(\operatorname{ring\text{-}as\text{-}module}(\operatorname{zz\text{-}ring}), \operatorname{principal\text{-}ideal}(\operatorname{zz\text{-}ring}, \operatorname{entry}(\operatorname{diag\text{-}mat}(n, l), j, j))), \operatorname{entry}((\lambda c \in \operatorname{mat}(1, n, \operatorname{carr}(\operatorname{scal}(\operatorname{zz\text{-}pow}(n)))).\; \operatorname{lincomb}(\operatorname{zz\text{-}pow}(n), n, c, \operatorname{zz\text{-}units}(n)))^{-1}(\mathit{phx}), 1, j))))(x)) \wedge \\ \quad \forall x \in \operatorname{vec}(\operatorname{zz\text{-}pow}(n)). \\ \quad \quad (\lambda \mathit{phx} \in \operatorname{vec}(\operatorname{zz\text{-}pow}(n)).\; (\lambda j \in \operatorname{interval}(1, n).\; \operatorname{class}(\operatorname{qmod\text{-}setoid}(\operatorname{ring\text{-}as\text{-}module}(\operatorname{zz\text{-}ring}), \operatorname{principal\text{-}ideal}(\operatorname{zz\text{-}ring}, \operatorname{entry}(\operatorname{diag\text{-}mat}(n, l), j, j))), \operatorname{entry}((\lambda c \in \operatorname{mat}(1, n, \operatorname{carr}(\operatorname{scal}(\operatorname{zz\text{-}pow}(n)))).\; \operatorname{lincomb}(\operatorname{zz\text{-}pow}(n), n, c, \operatorname{zz\text{-}units}(n)))^{-1}(\mathit{phx}), 1, j))))(b(x)) \\ \quad \quad =\; g((\lambda \mathit{phx} \in \operatorname{vec}(\operatorname{zz\text{-}pow}(n)).\; (\lambda j \in \operatorname{interval}(1, n).\; \operatorname{class}(\operatorname{qmod\text{-}setoid}(\operatorname{ring\text{-}as\text{-}module}(\operatorname{zz\text{-}ring}), \operatorname{principal\text{-}ideal}(\operatorname{zz\text{-}ring}, \operatorname{entry}(\operatorname{diag\text{-}mat}(n, e), j, j))), \operatorname{entry}((\lambda c \in \operatorname{mat}(1, n, \operatorname{carr}(\operatorname{scal}(\operatorname{zz\text{-}pow}(n)))).\; \operatorname{lincomb}(\operatorname{zz\text{-}pow}(n), n, c, \operatorname{zz\text{-}units}(n)))^{-1}(\mathit{phx}), 1, j))))(x)) \wedge \\ \quad \forall y \in \operatorname{vec}(\operatorname{direct\text{-}sum}(\operatorname{zz\text{-}ring}, n, (\lambda i \in \operatorname{interval}(1, n).\; \operatorname{zz\text{-}mod}(\operatorname{entry}(\operatorname{diag\text{-}mat}(n, e), i, i))))). \\ \quad \quad h(g(y)) = y \wedge \\ \quad z \in \operatorname{span}(\operatorname{zz\text{-}pow}(m), m, \operatorname{matact}(\operatorname{zz\text{-}pow}(m), \operatorname{diag\text{-}mat}(m, \operatorname{pad\text{-}ones}(m, n, l)), \operatorname{zz\text{-}units}(m))) \Rightarrow \\ \quad \quad \operatorname{schanuel\text{-}map}(m, n, a, b)(z) \\ \quad \quad \in\; \operatorname{span}(\operatorname{zz\text{-}pow}(m), m, \operatorname{matact}(\operatorname{zz\text{-}pow}(m), \operatorname{diag\text{-}mat}(m, \operatorname{pad\text{-}ones}(m, n, e)), \operatorname{zz\text{-}units}(m))) \end{array}}}
\vnbprop{cyclic-sum-iso-core}{\vnbstmt{\begin{array}{@{}l@{}} \forall n, l, e, h. \\ \quad n \in \mathbb{N} \wedge \\ \quad l \in (\operatorname{interval}(1, n) \to \mathbb{Z}) \wedge \\ \quad e \in (\operatorname{interval}(1, n) \to \mathbb{Z}) \wedge \\ \quad \forall i \in \operatorname{interval}(1, n).\; (0 \leq l(i)) \wedge \\ \quad \forall i \in \operatorname{interval}(1, n).\; (0 \leq e(i)) \wedge \\ \quad \forall i, j. \\ \quad \quad ((i \in \operatorname{interval}(1, n)) \wedge ((j \in \operatorname{interval}(1, n)) \wedge (i \leq j))) \Rightarrow  \\ \quad \quad \quad \operatorname{zz\text{-}divides}(l(i), l(j)) \wedge \\ \quad \forall i, j. \\ \quad \quad ((i \in \operatorname{interval}(1, n)) \wedge ((j \in \operatorname{interval}(1, n)) \wedge (i \leq j))) \Rightarrow  \\ \quad \quad \quad \operatorname{zz\text{-}divides}(e(i), e(j)) \wedge \\ \quad \operatorname{is\text{-}iso\text{-}module}(\operatorname{cyclic\text{-}sum}(n, l), \operatorname{cyclic\text{-}sum}(n, e), h) \Rightarrow \\ \quad \quad l = e \end{array}}}
\vnbprop{cyclic-sum-iso-unique}{\vnbstmt{\begin{array}{@{}l@{}} \forall l, a, c, b, h. \\ \quad l \in \mathbb{N} \wedge \\ \quad c \in \mathbb{N} \wedge \\ \quad \operatorname{is\text{-}invariant\text{-}factors}(l, a) \wedge \\ \quad \operatorname{is\text{-}invariant\text{-}factors}(c, b) \wedge \\ \quad \operatorname{is\text{-}iso\text{-}module}(\operatorname{cyclic\text{-}sum}(l, a), \operatorname{cyclic\text{-}sum}(c, b), h) \Rightarrow \\ \quad \quad (l = c) \wedge (a = b) \end{array}}}
\vnbprop{fg-abelian-group-structure-unique}{\vnbstmt{\begin{array}{@{}l@{}} \forall e, l, a, f, c, b, g. \\ \quad \operatorname{is\text{-}module}(e) \wedge \\ \quad l \in \mathbb{N} \wedge \\ \quad c \in \mathbb{N} \wedge \\ \quad \operatorname{is\text{-}invariant\text{-}factors}(l, a) \wedge \\ \quad \operatorname{is\text{-}invariant\text{-}factors}(c, b) \wedge \\ \quad \operatorname{is\text{-}iso\text{-}module}(\operatorname{cyclic\text{-}sum}(l, a), e, f) \wedge \\ \quad \operatorname{is\text{-}iso\text{-}module}(\operatorname{cyclic\text{-}sum}(c, b), e, g) \Rightarrow \\ \quad \quad (l = c) \wedge (a = b) \end{array}}}
\vnbprop{fg-abelian-group-cyclic-decomposition}{\vnbstmt{\begin{array}{@{}l@{}} \forall e, n, u. \\ \quad (\operatorname{is\text{-}module}(e) \wedge ((\operatorname{scal}(e) = \operatorname{zz\text{-}ring}) \wedge ((n \in \mathbb{N}) \wedge ((u \in \operatorname{mat}(n, 1, \operatorname{vec}(e))) \wedge \operatorname{generates}(e, n, u))))) \Rightarrow  \\ \quad \quad \exists l \in \mathbb{N}, a. \\ \quad \quad \quad \operatorname{is\text{-}invariant\text{-}factors}(l, a) \wedge \\ \quad \quad \quad \exists f.\; \operatorname{is\text{-}iso\text{-}module}(\operatorname{cyclic\text{-}sum}(l, a), e, f) \wedge \\ \quad \quad \quad \forall c, b. \\ \quad \quad \quad \quad ((c \in \mathbb{N}) \wedge (\operatorname{is\text{-}invariant\text{-}factors}(c, b) \wedge \exists g.\; \operatorname{is\text{-}iso\text{-}module}(\operatorname{cyclic\text{-}sum}(c, b), e, g))) \Rightarrow  \\ \quad \quad \quad \quad \quad (l = c) \wedge (a = b) \end{array}}}

\vnbbookfile{theorem-library/la-z5-norm}{theorem-library/la-z5-norm}
\noindent{\small la-z5-norm.scm -- linear algebra 3.41 (second sentence) and the norm N\_E(F) as an object (LA-Z5, 2026-10-25).  Prefix lz5n-.}\par
\vnbprop{free-generating-length-unique}{\vnbstmt{\begin{array}{@{}l@{}} \forall d, n, m, u, v. \\ \quad \operatorname{is\text{-}module}(d) \wedge \\ \quad \operatorname{is\text{-}euclidean\text{-}ring}(\operatorname{scal}(d)) \wedge \\ \quad n \in \mathbb{N} \wedge \\ \quad m \in \mathbb{N} \wedge \\ \quad u \in \operatorname{mat}(n, 1, \operatorname{vec}(d)) \wedge \\ \quad v \in \operatorname{mat}(m, 1, \operatorname{vec}(d)) \wedge \\ \quad \operatorname{generates}(d, n, u) \wedge \\ \quad \operatorname{rel\text{-}free}(d, n, u) \wedge \\ \quad \operatorname{generates}(d, m, v) \wedge \\ \quad \operatorname{rel\text{-}free}(d, m, v) \Rightarrow \\ \quad \quad m = n \end{array}}}
\vnbprop{norm-data-det-unique}{\vnbstmt{\begin{array}{@{}l@{}} \forall d, m, n, s, a, b, c, s_{2}, f, m_{2}. \\ \quad \operatorname{is\text{-}module}(d) \wedge \\ \quad \operatorname{scal}(d) = \operatorname{zz\text{-}ring} \wedge \\ \quad \operatorname{is\text{-}norm\text{-}data}(d, m, n, s, a, b) \wedge \\ \quad \operatorname{is\text{-}norm\text{-}data}(d, m, c, s_{2}, f, m_{2}) \Rightarrow \\ \quad \quad \lvert \operatorname{det}(\operatorname{zz\text{-}ring}, c, m_{2}) \rvert = \lvert \operatorname{det}(\operatorname{zz\text{-}ring}, n, b) \rvert \end{array}}}
\vnbprop{submodule-norm-zero}{\vnbstmt{\begin{array}{@{}l@{}} \forall d, m. \\ \quad \neg \exists n, x, y, a.\; \operatorname{is\text{-}norm\text{-}data}(d, m, n, x, y, a) \Rightarrow  \\ \quad \quad \operatorname{submodule\text{-}norm}(d, m) = 0 \end{array}}}
\vnbprop{submodule-norm-value}{\vnbstmt{\begin{array}{@{}l@{}} \forall d, m, n, s, a, b. \\ \quad (\operatorname{is\text{-}module}(d) \wedge ((\operatorname{scal}(d) = \operatorname{zz\text{-}ring}) \wedge \operatorname{is\text{-}norm\text{-}data}(d, m, n, s, a, b))) \Rightarrow  \\ \quad \quad \operatorname{submodule\text{-}norm}(d, m) = \lvert \operatorname{det}(\operatorname{zz\text{-}ring}, n, b) \rvert \end{array}}}

\vnbbookfile{theorem-library/la-z5-zzmod}{theorem-library/la-z5-zzmod}
\noindent{\small la-z5-zzmod.scm -- counting: Z/(b) has b elements for b $>$= 1 (LA-Z5, 2026-10-25). Prefix lz5z-.}\par
\vnbprop{zz-principal-ideal-out}{\vnbstmt{\begin{array}{@{}l@{}} \forall b, x. \\ \quad ((b \in \mathbb{Z}) \wedge (x \in \operatorname{principal\text{-}ideal}(\operatorname{zz\text{-}ring}, b))) \Rightarrow  \\ \quad \quad (x \in \mathbb{Z}) \wedge \exists t \in \mathbb{Z}.\; (x = (t \cdot b)) \end{array}}}
\vnbprop{zz-principal-ideal-in}{\vnbstmt{\begin{array}{@{}l@{}} \forall b, x, t. \\ \quad ((b \in \mathbb{Z}) \wedge ((x \in \mathbb{Z}) \wedge ((t \in \mathbb{Z}) \wedge (x = (t \cdot b))))) \Rightarrow  \\ \quad \quad x \in \operatorname{principal\text{-}ideal}(\operatorname{zz\text{-}ring}, b) \end{array}}}
\vnbprop{zz-mod-card}{\vnbstmt{\forall b \in \mathbb{Z}.\; ((1 \leq b) \Rightarrow (|\operatorname{vec}(\operatorname{zz\text{-}mod}(b))| = b))}}
\vnbprop{zz-mod-vec-set}{\vnbstmt{\forall b \in \mathbb{Z}.\; (\operatorname{vec}(\operatorname{zz\text{-}mod}(b)) \in \mathbf{Set})}}

\vnbbookfile{theorem-library/la-z5-dsumcard}{theorem-library/la-z5-dsumcard}
\noindent{\small la-z5-dsumcard.scm -- counting: a finite direct sum has the product of the summands' cardinalities as its cardinality (LA-Z5, 2026-10-25).  Prefix lz5d-.}\par
\vnbprop{dsum-vec-set-gen}{\vnbstmt{\begin{array}{@{}l@{}} \forall n, f. \\ \quad ((n \in \mathbb{N}) \wedge \forall i \in \operatorname{interval}(1, n).\; (\operatorname{vec}(f(i)) \in \mathbf{Set})) \Rightarrow  \\ \quad \quad \operatorname{dsum\text{-}vec}(n, f) \in \mathbf{Set} \end{array}}}
\vnbprop{dsum-succ-bijection}{\vnbstmt{\begin{array}{@{}l@{}} \forall n, f. \\ \quad ((n \in \mathbb{N}) \wedge \forall i \in \operatorname{interval}(1, (n + 1)).\; (\operatorname{vec}(f(i)) \in \mathbf{Set})) \Rightarrow  \\ \quad \quad (\lambda u \in (\operatorname{dsum\text{-}vec}(n, f) \times \operatorname{vec}(f((n + 1)))).\; \\ \quad \quad \quad (\lambda \mathit{dci} \in \operatorname{interval}(1, (n + 1)).\; \left(\text{if }(\mathit{dci} \leq n)\text{ then }\operatorname{nth}(1, u)(\mathit{dci})\text{ else }\operatorname{nth}(2, u)\right))) \\ \quad \quad \in\; \mathrm{Bij}((\operatorname{dsum\text{-}vec}(n, f) \times \operatorname{vec}(f((n + 1)))), \operatorname{dsum\text{-}vec}((n + 1), f)) \end{array}}}
\vnbprop{dsum-card-step}{\vnbstmt{\begin{array}{@{}l@{}} \forall n, f. \\ \quad n \in \mathbb{N} \wedge \\ \quad \forall i \in \operatorname{interval}(1, (n + 1)).\; (\operatorname{vec}(f(i)) \in \mathbf{Set}) \wedge \\ \quad \forall i \in \operatorname{interval}(1, (n + 1)).\; (|\operatorname{vec}(f(i))| \in \mathbb{N}) \wedge \\ \quad \forall i \in \operatorname{interval}(1, n).\; (\operatorname{vec}(f(i)) \in \mathbf{Set}) \wedge \\ \quad \forall i \in \operatorname{interval}(1, n).\; (|\operatorname{vec}(f(i))| \in \mathbb{N}) \Rightarrow \\ \quad \quad |\operatorname{dsum\text{-}vec}(n, f)| \in \mathbb{N} \wedge \\ \quad \quad |\operatorname{dsum\text{-}vec}(n, f)| \\ \quad \quad =\; \operatorname{prod\text{-}ring}(\operatorname{zz\text{-}ring}, (\lambda k \in \operatorname{interval}(1, n).\; |\operatorname{vec}(f(k))|), \operatorname{interval}(1, n)) \Rightarrow \\ \quad \quad |\operatorname{dsum\text{-}vec}((n + 1), f)| \in \mathbb{N} \wedge \\ \quad \quad |\operatorname{dsum\text{-}vec}((n + 1), f)| \\ \quad \quad =\; \operatorname{prod\text{-}ring}(\operatorname{zz\text{-}ring}, (\lambda k \in \operatorname{interval}(1, (n + 1)).\; |\operatorname{vec}(f(k))|), \operatorname{interval}(1, (n + 1))) \end{array}}}
\vnbprop{dsum-card-zero}{\vnbstmt{\begin{array}{@{}l@{}} \forall f. \\ \quad \forall i \in \operatorname{interval}(1, 0).\; (\operatorname{vec}(f(i)) \in \mathbf{Set}) \wedge \\ \quad \forall i \in \operatorname{interval}(1, 0).\; (|\operatorname{vec}(f(i))| \in \mathbb{N}) \Rightarrow \\ \quad \quad |\operatorname{dsum\text{-}vec}(0, f)| \in \mathbb{N} \wedge \\ \quad \quad |\operatorname{dsum\text{-}vec}(0, f)| \\ \quad \quad =\; \operatorname{prod\text{-}ring}(\operatorname{zz\text{-}ring}, (\lambda k \in \operatorname{interval}(1, 0).\; |\operatorname{vec}(f(k))|), \operatorname{interval}(1, 0)) \end{array}}}
\vnbprop{dsum-card-ind}{\vnbstmt{\begin{array}{@{}l@{}} \forall n \in \mathbb{N}, m. \\ \quad \forall i \in \operatorname{interval}(1, n).\; (\operatorname{vec}(m(i)) \in \mathbf{Set}) \wedge \\ \quad \forall i \in \operatorname{interval}(1, n).\; (|\operatorname{vec}(m(i))| \in \mathbb{N}) \Rightarrow \\ \quad \quad |\operatorname{dsum\text{-}vec}(n, m)| \in \mathbb{N} \wedge \\ \quad \quad |\operatorname{dsum\text{-}vec}(n, m)| \\ \quad \quad =\; \operatorname{prod\text{-}ring}(\operatorname{zz\text{-}ring}, (\lambda k \in \operatorname{interval}(1, n).\; |\operatorname{vec}(m(k))|), \operatorname{interval}(1, n)) \end{array}}}
\vnbprop{dsum-card}{\vnbstmt{\begin{array}{@{}l@{}} \forall n, f. \\ \quad n \in \mathbb{N} \wedge \\ \quad \forall i \in \operatorname{interval}(1, n).\; (\operatorname{vec}(f(i)) \in \mathbf{Set}) \wedge \\ \quad \forall i \in \operatorname{interval}(1, n).\; (|\operatorname{vec}(f(i))| \in \mathbb{N}) \Rightarrow \\ \quad \quad |\operatorname{dsum\text{-}vec}(n, f)| \in \mathbb{N} \wedge \\ \quad \quad |\operatorname{dsum\text{-}vec}(n, f)| \\ \quad \quad =\; \operatorname{prod\text{-}ring}(\operatorname{zz\text{-}ring}, (\lambda k \in \operatorname{interval}(1, n).\; |\operatorname{vec}(f(k))|), \operatorname{interval}(1, n)) \end{array}}}

\vnbbookfile{theorem-library/la-z5-card52}{theorem-library/la-z5-card52}
\noindent{\small la-z5-card52.scm -- linear algebra 3.52: if E is a finitely generated free abelian group, F a subgroup and N\_E(F) /= 0, then card(E/F) = N\_E(F) (LA-Z5, 2026-10-25). Prefix lz5q-.}\par
\vnbprop{diagonal-det-prod}{\vnbstmt{\begin{array}{@{}l@{}} \forall n, m. \\ \quad ((n \in \mathbb{N}) \wedge ((m \in \operatorname{mat}(n, n, \mathbb{Z})) \wedge \operatorname{is\text{-}diagonal}(\operatorname{zz\text{-}ring}, n, n, m))) \Rightarrow  \\ \quad \quad \operatorname{det}(\operatorname{zz\text{-}ring}, n, m) \\ \quad \quad =\; \operatorname{prod\text{-}ring}(\operatorname{zz\text{-}ring}, (\lambda k \in \operatorname{interval}(1, n).\; \operatorname{entry}(m, k, k)), \operatorname{interval}(1, n)) \end{array}}}
\vnbprop{principal-sequence-of-presentation}{\vnbstmt{\begin{array}{@{}l@{}} \forall d, m, n, u, s, a. \\ \quad (\operatorname{is\text{-}module}(d) \wedge ((\operatorname{scal}(d) = \operatorname{zz\text{-}ring}) \wedge (\operatorname{is\text{-}submodule}(d, m) \wedge \operatorname{is\text{-}norm\text{-}data}(d, m, n, u, s, a)))) \Rightarrow  \\ \quad \quad \exists b, c. \\ \quad \quad \quad \operatorname{is\text{-}principal\text{-}sequence}(d, n, m, b, c) \wedge \operatorname{rel\text{-}free}(d, n, \operatorname{matact}(d, c, b)) \end{array}}}
\vnbprop{cyclic-dsum-card}{\vnbstmt{\begin{array}{@{}l@{}} \forall n, m. \\ \quad ((n \in \mathbb{N}) \wedge ((m \in \operatorname{mat}(n, n, \mathbb{Z})) \wedge \forall i \in \operatorname{interval}(1, n).\; (1 \leq \operatorname{entry}(m, i, i)))) \Rightarrow  \\ \quad \quad \operatorname{vec}(\operatorname{direct\text{-}sum}(\operatorname{zz\text{-}ring}, n, (\lambda \mathit{fmi} \in \operatorname{interval}(1, n).\; \operatorname{zz\text{-}mod}(\operatorname{entry}(m, \mathit{fmi}, \mathit{fmi}))))) \in \mathbf{Set} \wedge \\ \quad \quad |\operatorname{vec}(\operatorname{direct\text{-}sum}(\operatorname{zz\text{-}ring}, n, (\lambda \mathit{fmi} \in \operatorname{interval}(1, n).\; \operatorname{zz\text{-}mod}(\operatorname{entry}(m, \mathit{fmi}, \mathit{fmi})))))| \in \mathbb{N} \wedge \\ \quad \quad |\operatorname{vec}(\operatorname{direct\text{-}sum}(\operatorname{zz\text{-}ring}, n, (\lambda \mathit{fmi} \in \operatorname{interval}(1, n).\; \operatorname{zz\text{-}mod}(\operatorname{entry}(m, \mathit{fmi}, \mathit{fmi})))))| \\ \quad \quad =\; \operatorname{prod\text{-}ring}(\operatorname{zz\text{-}ring}, (\lambda k \in \operatorname{interval}(1, n).\; \operatorname{entry}(m, k, k)), \operatorname{interval}(1, n)) \end{array}}}
\vnbprop{quotient-card-norm}{\vnbstmt{\begin{array}{@{}l@{}} \forall d, n, u, m. \\ \quad \operatorname{is\text{-}module}(d) \wedge \\ \quad \operatorname{scal}(d) = \operatorname{zz\text{-}ring} \wedge \\ \quad n \in \mathbb{N} \wedge \\ \quad u \in \operatorname{mat}(n, 1, \operatorname{vec}(d)) \wedge \\ \quad \operatorname{generates}(d, n, u) \wedge \\ \quad \operatorname{rel\text{-}free}(d, n, u) \wedge \\ \quad \operatorname{is\text{-}submodule}(d, m) \wedge \\ \quad \neg (\operatorname{submodule\text{-}norm}(d, m) = 0) \Rightarrow \\ \quad \quad |\operatorname{vec}(\operatorname{quotient\text{-}module}(d, m))| = \operatorname{submodule\text{-}norm}(d, m) \end{array}}}

\vnbbookfile{theorem-library/qag-laws}{theorem-library/qag-laws}
\noindent{\small qag-laws.scm -- WAVE 2 (2026-10-27..30), wired 2026-10-30 from scratchpad/ccrpr/qag-laws.scm. Library: linear-algebra.}\par
\vnbprop{agm-scal}{\vnbstmt{\forall g.\; (\operatorname{is\text{-}abelian\text{-}group}(g) \Rightarrow (\operatorname{scal}(\operatorname{ag\text{-}zz\text{-}module}(g)) \simeq \operatorname{zz\text{-}ring}))}}
\vnbprop{agm-vec}{\vnbstmt{\forall g.\; (\operatorname{is\text{-}abelian\text{-}group}(g) \Rightarrow (\operatorname{vec}(\operatorname{ag\text{-}zz\text{-}module}(g)) \simeq \operatorname{carr}(g)))}}
\vnbprop{agm-vadd}{\vnbstmt{\forall g.\; (\operatorname{is\text{-}abelian\text{-}group}(g) \Rightarrow (\operatorname{vadd}(\operatorname{ag\text{-}zz\text{-}module}(g)) \simeq \operatorname{opr}(g)))}}
\vnbprop{agm-vzero}{\vnbstmt{\forall g.\; (\operatorname{is\text{-}abelian\text{-}group}(g) \Rightarrow (\operatorname{vzero}(\operatorname{ag\text{-}zz\text{-}module}(g)) \simeq \operatorname{iden}(g)))}}
\vnbprop{agm-vneg}{\vnbstmt{\forall g.\; (\operatorname{is\text{-}abelian\text{-}group}(g) \Rightarrow (\operatorname{vneg}(\operatorname{ag\text{-}zz\text{-}module}(g)) \simeq \operatorname{inv}(g)))}}
\vnbprop{agm-act}{\vnbstmt{\begin{array}{@{}l@{}} \forall g. \\ \quad \operatorname{is\text{-}abelian\text{-}group}(g) \Rightarrow  \\ \quad \quad (\operatorname{act}(\operatorname{ag\text{-}zz\text{-}module}(g)) \simeq (\lambda {[}\operatorname{azk\_}, \operatorname{aza\_}] \in (\mathbb{Z} \times \operatorname{carr}(g)).\; \operatorname{zz\text{-}act}(g, \operatorname{azk\_}, \operatorname{aza\_}))) \end{array}}}
\vnbprop{agm-act-fun}{\vnbstmt{\begin{array}{@{}l@{}} \forall g. \\ \quad \operatorname{is\text{-}abelian\text{-}group}(g) \Rightarrow  \\ \quad \quad \operatorname{act}(\operatorname{ag\text{-}zz\text{-}module}(g)) \in ((\mathbb{Z} \times \operatorname{carr}(g)) \to \operatorname{carr}(g)) \end{array}}}
\vnbprop{agm-act-value}{\vnbstmt{\begin{array}{@{}l@{}} \forall g.; \forall k \in \mathbb{Z}, a \in \operatorname{carr}(g). \\ \quad \operatorname{is\text{-}abelian\text{-}group}(g) \Rightarrow (\operatorname{act}(\operatorname{ag\text{-}zz\text{-}module}(g))(k, a) = \operatorname{zz\text{-}act}(g, k, a)) \end{array}}}
\vnbprop{ag-zz-module-is-module}{\vnbstmt{\forall g.\; (\operatorname{is\text{-}abelian\text{-}group}(g) \Rightarrow \operatorname{is\text{-}module}(\operatorname{ag\text{-}zz\text{-}module}(g)))}}
\vnbprop{ag-subgroup-zz-act-nn}{\vnbstmt{\begin{array}{@{}l@{}} \forall n \in \mathbb{N}, g, h.; \forall a \in h. \\ \quad (\operatorname{is\text{-}abelian\text{-}group}(g) \wedge \operatorname{is\text{-}subgroup}(g, h)) \Rightarrow (\operatorname{zz\text{-}act}(g, n, a) \in h) \end{array}}}
\vnbprop{ag-subgroup-zz-act}{\vnbstmt{\begin{array}{@{}l@{}} \forall g, h.; \forall k \in \mathbb{Z}, a \in h. \\ \quad (\operatorname{is\text{-}abelian\text{-}group}(g) \wedge \operatorname{is\text{-}subgroup}(g, h)) \Rightarrow (\operatorname{zz\text{-}act}(g, k, a) \in h) \end{array}}}
\vnbprop{ag-subgroup-submodule}{\vnbstmt{\begin{array}{@{}l@{}} \forall g, h. \\ \quad (\operatorname{is\text{-}abelian\text{-}group}(g) \wedge \operatorname{is\text{-}subgroup}(g, h)) \Rightarrow  \\ \quad \quad \operatorname{is\text{-}submodule}(\operatorname{ag\text{-}zz\text{-}module}(g), h) \end{array}}}
\vnbprop{quotient-ag-is-abelian-group}{\vnbstmt{\begin{array}{@{}l@{}} \forall g, h. \\ \quad (\operatorname{is\text{-}abelian\text{-}group}(g) \wedge \operatorname{is\text{-}subgroup}(g, h)) \Rightarrow  \\ \quad \quad \operatorname{is\text{-}abelian\text{-}group}(\operatorname{quotient\text{-}ag}(g, h)) \end{array}}}
\vnbprop{qag-carr-vec}{\vnbstmt{\begin{array}{@{}l@{}} \forall g, h. \\ \quad (\operatorname{is\text{-}abelian\text{-}group}(g) \wedge \operatorname{is\text{-}subgroup}(g, h)) \Rightarrow  \\ \quad \quad (\operatorname{carr}(\operatorname{quotient\text{-}ag}(g, h)) \simeq \operatorname{vec}(\operatorname{quotient\text{-}module}(\operatorname{ag\text{-}zz\text{-}module}(g), h))) \end{array}}}
\vnbprop{qag-opr-vadd}{\vnbstmt{\begin{array}{@{}l@{}} \forall g, h. \\ \quad (\operatorname{is\text{-}abelian\text{-}group}(g) \wedge \operatorname{is\text{-}subgroup}(g, h)) \Rightarrow  \\ \quad \quad (\operatorname{opr}(\operatorname{quotient\text{-}ag}(g, h)) \simeq \operatorname{vadd}(\operatorname{quotient\text{-}module}(\operatorname{ag\text{-}zz\text{-}module}(g), h))) \end{array}}}
\vnbprop{qag-iden-vzero}{\vnbstmt{\begin{array}{@{}l@{}} \forall g, h. \\ \quad (\operatorname{is\text{-}abelian\text{-}group}(g) \wedge \operatorname{is\text{-}subgroup}(g, h)) \Rightarrow  \\ \quad \quad (\operatorname{iden}(\operatorname{quotient\text{-}ag}(g, h)) \simeq \operatorname{vzero}(\operatorname{quotient\text{-}module}(\operatorname{ag\text{-}zz\text{-}module}(g), h))) \end{array}}}
\vnbprop{qag-inv-vneg}{\vnbstmt{\begin{array}{@{}l@{}} \forall g, h. \\ \quad (\operatorname{is\text{-}abelian\text{-}group}(g) \wedge \operatorname{is\text{-}subgroup}(g, h)) \Rightarrow  \\ \quad \quad (\operatorname{inv}(\operatorname{quotient\text{-}ag}(g, h)) \simeq \operatorname{vneg}(\operatorname{quotient\text{-}module}(\operatorname{ag\text{-}zz\text{-}module}(g), h))) \end{array}}}
\vnbprop{qag-carr}{\vnbstmt{\begin{array}{@{}l@{}} \forall g, h. \\ \quad (\operatorname{is\text{-}abelian\text{-}group}(g) \wedge \operatorname{is\text{-}subgroup}(g, h)) \Rightarrow  \\ \quad \quad (\operatorname{carr}(\operatorname{quotient\text{-}ag}(g, h)) \simeq \operatorname{quotient}(\operatorname{qmod\text{-}setoid}(\operatorname{ag\text{-}zz\text{-}module}(g), h))) \end{array}}}
\vnbprop{qag-class-in}{\vnbstmt{\begin{array}{@{}l@{}} \forall g, h.; \forall a \in \operatorname{carr}(g). \\ \quad (\operatorname{is\text{-}abelian\text{-}group}(g) \wedge \operatorname{is\text{-}subgroup}(g, h)) \Rightarrow  \\ \quad \quad \operatorname{class}(\operatorname{qmod\text{-}setoid}(\operatorname{ag\text{-}zz\text{-}module}(g), h), a) \in \operatorname{carr}(\operatorname{quotient\text{-}ag}(g, h)) \end{array}}}
\vnbprop{qag-rep}{\vnbstmt{\begin{array}{@{}l@{}} \forall g, h.; \forall c \in \operatorname{carr}(\operatorname{quotient\text{-}ag}(g, h)). \\ \quad (\operatorname{is\text{-}abelian\text{-}group}(g) \wedge \operatorname{is\text{-}subgroup}(g, h)) \Rightarrow  \\ \quad \quad \exists a \in \operatorname{carr}(g).\; (c = \operatorname{class}(\operatorname{qmod\text{-}setoid}(\operatorname{ag\text{-}zz\text{-}module}(g), h), a)) \end{array}}}
\vnbprop{qag-opr-class}{\vnbstmt{\begin{array}{@{}l@{}} \forall g, h.; \forall a, b \in \operatorname{carr}(g). \\ \quad (\operatorname{is\text{-}abelian\text{-}group}(g) \wedge \operatorname{is\text{-}subgroup}(g, h)) \Rightarrow  \\ \quad \quad \operatorname{opr}(\operatorname{quotient\text{-}ag}(g, h))(\operatorname{class}(\operatorname{qmod\text{-}setoid}(\operatorname{ag\text{-}zz\text{-}module}(g), h), a), \\ \quad \quad \quad \operatorname{class}(\operatorname{qmod\text{-}setoid}(\operatorname{ag\text{-}zz\text{-}module}(g), h), b)) \\ \quad \quad =\; \operatorname{class}(\operatorname{qmod\text{-}setoid}(\operatorname{ag\text{-}zz\text{-}module}(g), h), \operatorname{opr}(g)(a, b)) \end{array}}}
\vnbprop{qag-iden}{\vnbstmt{\begin{array}{@{}l@{}} \forall g, h. \\ \quad (\operatorname{is\text{-}abelian\text{-}group}(g) \wedge \operatorname{is\text{-}subgroup}(g, h)) \Rightarrow  \\ \quad \quad (\operatorname{iden}(\operatorname{quotient\text{-}ag}(g, h)) \simeq \operatorname{class}(\operatorname{qmod\text{-}setoid}(\operatorname{ag\text{-}zz\text{-}module}(g), h), \operatorname{iden}(g))) \end{array}}}
\vnbprop{qag-inv-class}{\vnbstmt{\begin{array}{@{}l@{}} \forall g, h.; \forall a \in \operatorname{carr}(g). \\ \quad (\operatorname{is\text{-}abelian\text{-}group}(g) \wedge \operatorname{is\text{-}subgroup}(g, h)) \Rightarrow  \\ \quad \quad \operatorname{inv}(\operatorname{quotient\text{-}ag}(g, h))(\operatorname{class}(\operatorname{qmod\text{-}setoid}(\operatorname{ag\text{-}zz\text{-}module}(g), h), a)) \\ \quad \quad =\; \operatorname{class}(\operatorname{qmod\text{-}setoid}(\operatorname{ag\text{-}zz\text{-}module}(g), h), \operatorname{inv}(g)(a)) \end{array}}}
\vnbprop{qag-class-eq-iff}{\vnbstmt{\begin{array}{@{}l@{}} \forall g, h.; \forall a, b \in \operatorname{carr}(g). \\ \quad (\operatorname{is\text{-}abelian\text{-}group}(g) \wedge \operatorname{is\text{-}subgroup}(g, h)) \Rightarrow  \\ \quad \quad \operatorname{class}(\operatorname{qmod\text{-}setoid}(\operatorname{ag\text{-}zz\text{-}module}(g), h), a) = \operatorname{class}(\operatorname{qmod\text{-}setoid}(\operatorname{ag\text{-}zz\text{-}module}(g), h), b) \\ \quad \quad \Leftrightarrow\; \operatorname{opr}(g)(a, \operatorname{inv}(g)(b)) \in h \end{array}}}
\vnbprop{qag-gpow-class}{\vnbstmt{\begin{array}{@{}l@{}} \forall n \in \mathbb{N}, g, h.; \forall a \in \operatorname{carr}(g). \\ \quad (\operatorname{is\text{-}abelian\text{-}group}(g) \wedge \operatorname{is\text{-}subgroup}(g, h)) \Rightarrow  \\ \quad \quad \operatorname{gpow}(\operatorname{quotient\text{-}ag}(g, h), \operatorname{class}(\operatorname{qmod\text{-}setoid}(\operatorname{ag\text{-}zz\text{-}module}(g), h), a), n) \\ \quad \quad =\; \operatorname{class}(\operatorname{qmod\text{-}setoid}(\operatorname{ag\text{-}zz\text{-}module}(g), h), \operatorname{gpow}(g, a, n)) \end{array}}}
\vnbprop{qag-pow-iden-iff}{\vnbstmt{\begin{array}{@{}l@{}} \forall n \in \mathbb{N}, g, h.; \forall a \in \operatorname{carr}(g). \\ \quad (\operatorname{is\text{-}abelian\text{-}group}(g) \wedge \operatorname{is\text{-}subgroup}(g, h)) \Rightarrow  \\ \quad \quad \operatorname{gpow}(\operatorname{quotient\text{-}ag}(g, h), \operatorname{class}(\operatorname{qmod\text{-}setoid}(\operatorname{ag\text{-}zz\text{-}module}(g), h), a), n) = \operatorname{iden}(\operatorname{quotient\text{-}ag}(g, h)) \\ \quad \quad \Leftrightarrow\; \operatorname{gpow}(g, a, n) \in h \end{array}}}
\vnbprop{qag-class-left-coset}{\vnbstmt{\begin{array}{@{}l@{}} \forall g, h.; \forall a \in \operatorname{carr}(g). \\ \quad (\operatorname{is\text{-}abelian\text{-}group}(g) \wedge \operatorname{is\text{-}subgroup}(g, h)) \Rightarrow  \\ \quad \quad \operatorname{class}(\operatorname{qmod\text{-}setoid}(\operatorname{ag\text{-}zz\text{-}module}(g), h), a) = \operatorname{left\text{-}coset}(g, a, h) \end{array}}}
\vnbprop{qag-carr-left-cosets}{\vnbstmt{\begin{array}{@{}l@{}} \forall g, h. \\ \quad (\operatorname{is\text{-}abelian\text{-}group}(g) \wedge \operatorname{is\text{-}subgroup}(g, h)) \Rightarrow  \\ \quad \quad \operatorname{carr}(\operatorname{quotient\text{-}ag}(g, h)) = \operatorname{left\text{-}cosets}(g, h) \end{array}}}

\vnbbookfile{structure-library/functoriality}{structure-library/functoriality}
\noindent{\small A file shared across books; the entries that belong to this one.}\par
\noindent{\small functoriality.scm -- every view-as is a FUNCTOR, and here is the proof.}\par
\vnbprop{module-vector-ag-functorial}{\vnbstmt{\begin{array}{@{}l@{}} \forall a, b, f. \\ \quad \operatorname{is\text{-}hom\text{-}module}(a, b, f) \Rightarrow  \\ \quad \quad \operatorname{is\text{-}hom\text{-}abelian\text{-}group}(\operatorname{module\text{-}vector\text{-}ag}(a), \operatorname{module\text{-}vector\text{-}ag}(b), f) \end{array}}}

\vnbbookfile{theorem-library/gal-ch}{theorem-library/gal-ch}
\noindent{\small gal-ch.scm -- general Cayley-Hamilton for a linear operator on a finite-dimensional space (GALOIS B6, 2026-10-26), with its general bricks: a homomorphism of abelian groups passes through a finite sum; the determinant commutes with a ring homomorphism; the operators p(T) form a commutative ring F[T] with the evaluation ...}\par
\vnbprop{hom-ag-sum-ag}{\vnbstmt{\begin{array}{@{}l@{}} \forall n \in \mathbb{N}, a, b, h, g. \\ \quad (\operatorname{is\text{-}hom\text{-}abelian\text{-}group}(a, b, h) \wedge \forall i \in \mathrm{OS}(n).\; (g(i) \in \operatorname{carr}(a))) \Rightarrow  \\ \quad \quad h(\operatorname{sum\text{-}ag}(a, g, n)) = \operatorname{sum\text{-}ag}(b, (\lambda k \in \mathbb{N}.\; h(g(k))), n) \end{array}}}
\vnbprop{hom-ag-finsum-interval}{\vnbstmt{\begin{array}{@{}l@{}} \forall a, b, h.; \forall n \in \mathbb{N}, f \in (\operatorname{interval}(1, n) \to \operatorname{carr}(a)). \\ \quad \operatorname{is\text{-}hom\text{-}abelian\text{-}group}(a, b, h) \Rightarrow  \\ \quad \quad h(\operatorname{finsum}(a, f, \operatorname{interval}(1, n))) = \operatorname{finsum}(b, (\lambda z \in \operatorname{interval}(1, n).\; h(f(z))), \operatorname{interval}(1, n)) \end{array}}}
\vnbprop{mat-map-type}{\vnbstmt{\begin{array}{@{}l@{}} \forall m, n \in \mathbb{N}, x, y.; \forall h \in (x \to y), a \in \operatorname{mat}(m, n, x). \\ \quad \operatorname{mat\text{-}map}(h, m, n, a) \in \operatorname{mat}(m, n, y) \end{array}}}
\vnbprop{mat-map-entry}{\vnbstmt{\begin{array}{@{}l@{}} \forall m, n \in \mathbb{N}, x, y.; \forall h \in (x \to y), a \in \operatorname{mat}(m, n, x), i \in \operatorname{interval}(1, m), j \in \operatorname{interval}(1, n). \\ \quad \operatorname{entry}(\operatorname{mat\text{-}map}(h, m, n, a), i, j) = h(\operatorname{entry}(a, i, j)) \end{array}}}
\vnbprop{hom-monoid-mpow}{\vnbstmt{\begin{array}{@{}l@{}} \forall k \in \mathbb{N}, a, b, h.; \forall x \in \operatorname{carr}(a). \\ \quad \operatorname{is\text{-}hom\text{-}monoid}(a, b, h) \Rightarrow (h(\operatorname{mpow}(a, x, k)) = \operatorname{mpow}(b, h(x), k)) \end{array}}}
\vnbprop{minor-mat-map}{\vnbstmt{\begin{array}{@{}l@{}} \forall b, c.; \forall h \in (\operatorname{carr}(b) \to \operatorname{carr}(c)), n \in \mathbb{N}, a \in \operatorname{mat}((n + 1), (n + 1), \operatorname{carr}(b)), p, q \in \mathbb{N}. \\ \quad (\operatorname{is\text{-}ring}(b) \wedge \operatorname{is\text{-}ring}(c)) \Rightarrow  \\ \quad \quad \operatorname{minor}(\operatorname{mat\text{-}map}(h, (n + 1), (n + 1), a), p, q, n) = \operatorname{mat\text{-}map}(h, n, n, \operatorname{minor}(a, p, q, n)) \end{array}}}
\vnbprop{det-ring-hom}{\vnbstmt{\begin{array}{@{}l@{}} \forall n \in \mathbb{N}, b, c, h.; \forall a \in \operatorname{mat}(n, n, \operatorname{carr}(b)). \\ \quad \operatorname{is\text{-}hom\text{-}ring}(b, c, h) \Rightarrow (h(\operatorname{det}(b, n, a)) = \operatorname{det}(c, n, \operatorname{mat\text{-}map}(h, n, n, a))) \end{array}}}
\vnbprop{op-of-apply}{\vnbstmt{\begin{array}{@{}l@{}} \forall d, t.; \forall p \in \operatorname{carr}(\operatorname{poly}(\operatorname{scal}(d))), w \in \operatorname{vec}(d). \\ \quad (\operatorname{is\text{-}hom\text{-}module}(d, d, t) \wedge \operatorname{is\text{-}commutative\text{-}ring}(\operatorname{scal}(d))) \Rightarrow  \\ \quad \quad \operatorname{op\text{-}of}(d, t, p)(w) = \operatorname{poly\text{-}op}(d, t, p, w) \end{array}}}
\vnbprop{op-of-in-ops}{\vnbstmt{\begin{array}{@{}l@{}} \forall d, t.; \forall p \in \operatorname{carr}(\operatorname{poly}(\operatorname{scal}(d))). \\ \quad (\operatorname{is\text{-}hom\text{-}module}(d, d, t) \wedge \operatorname{is\text{-}commutative\text{-}ring}(\operatorname{scal}(d))) \Rightarrow  \\ \quad \quad (\operatorname{op\text{-}of}(d, t, p) \in (\operatorname{vec}(d) \to \operatorname{vec}(d))) \wedge (\operatorname{op\text{-}of}(d, t, p) \in \operatorname{ops}(d, t)) \end{array}}}
\vnbprop{ops-unfold}{\vnbstmt{\begin{array}{@{}l@{}} \forall d, t. \\ \quad (\operatorname{is\text{-}hom\text{-}module}(d, d, t) \wedge \operatorname{is\text{-}commutative\text{-}ring}(\operatorname{scal}(d))) \Rightarrow  \\ \quad \quad (\operatorname{ops}(d, t) \simeq \{f \in (\operatorname{vec}(d) \to \operatorname{vec}(d)) : \exists p \in \operatorname{carr}(\operatorname{poly}(\operatorname{scal}(d))).\; (f = \operatorname{op\text{-}of}(d, t, p))\}) \end{array}}}
\vnbprop{ops-elim}{\vnbstmt{\begin{array}{@{}l@{}} \forall d, t.; \forall f \in \operatorname{ops}(d, t). \\ \quad (\operatorname{is\text{-}hom\text{-}module}(d, d, t) \wedge \operatorname{is\text{-}commutative\text{-}ring}(\operatorname{scal}(d))) \Rightarrow  \\ \quad \quad ((f \in (\operatorname{vec}(d) \to \operatorname{vec}(d))) \wedge \exists p \in \operatorname{carr}(\operatorname{poly}(\operatorname{scal}(d))).\; (f = \operatorname{op\text{-}of}(d, t, p))) \end{array}}}
\vnbprop{ops-in-set}{\vnbstmt{\begin{array}{@{}l@{}} \forall d, t. \\ \quad (\operatorname{is\text{-}hom\text{-}module}(d, d, t) \wedge \operatorname{is\text{-}commutative\text{-}ring}(\operatorname{scal}(d))) \Rightarrow (\operatorname{ops}(d, t) \in \mathbf{Set}) \end{array}}}
\vnbprop{poly-op-zero-poly}{\vnbstmt{\begin{array}{@{}l@{}} \forall d, t.; \forall w \in \operatorname{vec}(d). \\ \quad (\operatorname{is\text{-}hom\text{-}module}(d, d, t) \wedge \operatorname{is\text{-}commutative\text{-}ring}(\operatorname{scal}(d))) \Rightarrow  \\ \quad \quad \operatorname{poly\text{-}op}(d, t, \operatorname{zero}(\operatorname{poly}(\operatorname{scal}(d))), w) = \operatorname{vzero}(d) \end{array}}}
\vnbprop{poly-op-pneg}{\vnbstmt{\begin{array}{@{}l@{}} \forall d, t.; \forall p \in \operatorname{carr}(\operatorname{poly}(\operatorname{scal}(d))), w \in \operatorname{vec}(d). \\ \quad (\operatorname{is\text{-}hom\text{-}module}(d, d, t) \wedge \operatorname{is\text{-}commutative\text{-}ring}(\operatorname{scal}(d))) \Rightarrow  \\ \quad \quad \operatorname{poly\text{-}op}(d, t, \operatorname{neg}(\operatorname{poly}(\operatorname{scal}(d)))(p), w) = \operatorname{vneg}(d)(\operatorname{poly\text{-}op}(d, t, p, w)) \end{array}}}
\vnbprop{poly-op-one-poly}{\vnbstmt{\begin{array}{@{}l@{}} \forall d, t.; \forall w \in \operatorname{vec}(d). \\ \quad (\operatorname{is\text{-}hom\text{-}module}(d, d, t) \wedge \operatorname{is\text{-}commutative\text{-}ring}(\operatorname{scal}(d))) \Rightarrow  \\ \quad \quad \operatorname{poly\text{-}op}(d, t, \operatorname{one}(\operatorname{poly}(\operatorname{scal}(d))), w) = w \end{array}}}
\vnbprop{op-add-of}{\vnbstmt{\begin{array}{@{}l@{}} \forall d, t.; \forall p, q \in \operatorname{carr}(\operatorname{poly}(\operatorname{scal}(d))). \\ \quad (\operatorname{is\text{-}hom\text{-}module}(d, d, t) \wedge \operatorname{is\text{-}commutative\text{-}ring}(\operatorname{scal}(d))) \Rightarrow  \\ \quad \quad \operatorname{op\text{-}add}(d, t)(\operatorname{op\text{-}of}(d, t, p), \operatorname{op\text{-}of}(d, t, q)) = \operatorname{op\text{-}of}(d, t, \operatorname{add}(\operatorname{poly}(\operatorname{scal}(d)))(p, q)) \end{array}}}
\vnbprop{op-mul-of}{\vnbstmt{\begin{array}{@{}l@{}} \forall d, t.; \forall p, q \in \operatorname{carr}(\operatorname{poly}(\operatorname{scal}(d))). \\ \quad (\operatorname{is\text{-}hom\text{-}module}(d, d, t) \wedge \operatorname{is\text{-}commutative\text{-}ring}(\operatorname{scal}(d))) \Rightarrow  \\ \quad \quad \operatorname{op\text{-}mul}(d, t)(\operatorname{op\text{-}of}(d, t, p), \operatorname{op\text{-}of}(d, t, q)) = \operatorname{op\text{-}of}(d, t, \operatorname{mul}(\operatorname{poly}(\operatorname{scal}(d)))(p, q)) \end{array}}}
\vnbprop{op-neg-of}{\vnbstmt{\begin{array}{@{}l@{}} \forall d, t.; \forall p \in \operatorname{carr}(\operatorname{poly}(\operatorname{scal}(d))). \\ \quad (\operatorname{is\text{-}hom\text{-}module}(d, d, t) \wedge \operatorname{is\text{-}commutative\text{-}ring}(\operatorname{scal}(d))) \Rightarrow  \\ \quad \quad \operatorname{op\text{-}neg}(d, t)(\operatorname{op\text{-}of}(d, t, p)) = \operatorname{op\text{-}of}(d, t, \operatorname{neg}(\operatorname{poly}(\operatorname{scal}(d)))(p)) \end{array}}}
\vnbprop{op-zero-of}{\vnbstmt{\begin{array}{@{}l@{}} \forall d, t. \\ \quad (\operatorname{is\text{-}hom\text{-}module}(d, d, t) \wedge \operatorname{is\text{-}commutative\text{-}ring}(\operatorname{scal}(d))) \Rightarrow  \\ \quad \quad \operatorname{op\text{-}zero}(d) = \operatorname{op\text{-}of}(d, t, \operatorname{zero}(\operatorname{poly}(\operatorname{scal}(d)))) \end{array}}}
\vnbprop{op-one-of}{\vnbstmt{\begin{array}{@{}l@{}} \forall d, t. \\ \quad (\operatorname{is\text{-}hom\text{-}module}(d, d, t) \wedge \operatorname{is\text{-}commutative\text{-}ring}(\operatorname{scal}(d))) \Rightarrow  \\ \quad \quad \operatorname{op\text{-}one}(d) = \operatorname{op\text{-}of}(d, t, \operatorname{one}(\operatorname{poly}(\operatorname{scal}(d)))) \end{array}}}
\vnbprop{op-alg-carr}{\vnbstmt{\begin{array}{@{}l@{}} \forall d, t. \\ \quad (\operatorname{is\text{-}hom\text{-}module}(d, d, t) \wedge \operatorname{is\text{-}commutative\text{-}ring}(\operatorname{scal}(d))) \Rightarrow  \\ \quad \quad (\operatorname{carr}(\operatorname{op\text{-}alg}(d, t)) \simeq \operatorname{ops}(d, t)) \end{array}}}
\vnbprop{op-alg-add}{\vnbstmt{\begin{array}{@{}l@{}} \forall d, t. \\ \quad (\operatorname{is\text{-}hom\text{-}module}(d, d, t) \wedge \operatorname{is\text{-}commutative\text{-}ring}(\operatorname{scal}(d))) \Rightarrow  \\ \quad \quad (\operatorname{add}(\operatorname{op\text{-}alg}(d, t)) \simeq \operatorname{op\text{-}add}(d, t)) \end{array}}}
\vnbprop{op-alg-mul}{\vnbstmt{\begin{array}{@{}l@{}} \forall d, t. \\ \quad (\operatorname{is\text{-}hom\text{-}module}(d, d, t) \wedge \operatorname{is\text{-}commutative\text{-}ring}(\operatorname{scal}(d))) \Rightarrow  \\ \quad \quad (\operatorname{mul}(\operatorname{op\text{-}alg}(d, t)) \simeq \operatorname{op\text{-}mul}(d, t)) \end{array}}}
\vnbprop{op-alg-neg}{\vnbstmt{\begin{array}{@{}l@{}} \forall d, t. \\ \quad (\operatorname{is\text{-}hom\text{-}module}(d, d, t) \wedge \operatorname{is\text{-}commutative\text{-}ring}(\operatorname{scal}(d))) \Rightarrow  \\ \quad \quad (\operatorname{neg}(\operatorname{op\text{-}alg}(d, t)) \simeq \operatorname{op\text{-}neg}(d, t)) \end{array}}}
\vnbprop{op-alg-zero}{\vnbstmt{\begin{array}{@{}l@{}} \forall d, t. \\ \quad (\operatorname{is\text{-}hom\text{-}module}(d, d, t) \wedge \operatorname{is\text{-}commutative\text{-}ring}(\operatorname{scal}(d))) \Rightarrow  \\ \quad \quad (\operatorname{zero}(\operatorname{op\text{-}alg}(d, t)) \simeq \operatorname{op\text{-}zero}(d)) \end{array}}}
\vnbprop{op-alg-one}{\vnbstmt{\begin{array}{@{}l@{}} \forall d, t. \\ \quad (\operatorname{is\text{-}hom\text{-}module}(d, d, t) \wedge \operatorname{is\text{-}commutative\text{-}ring}(\operatorname{scal}(d))) \Rightarrow  \\ \quad \quad (\operatorname{one}(\operatorname{op\text{-}alg}(d, t)) \simeq \operatorname{op\text{-}one}(d)) \end{array}}}
\vnbprop{ops-closed}{\vnbstmt{\begin{array}{@{}l@{}} \forall d, t.; \forall f, g \in \operatorname{ops}(d, t). \\ \quad (\operatorname{is\text{-}hom\text{-}module}(d, d, t) \wedge \operatorname{is\text{-}commutative\text{-}ring}(\operatorname{scal}(d))) \Rightarrow  \\ \quad \quad \operatorname{op\text{-}add}(d, t)(f, g) \in \operatorname{ops}(d, t) \wedge \\ \quad \quad \operatorname{op\text{-}mul}(d, t)(f, g) \in \operatorname{ops}(d, t) \wedge \\ \quad \quad \operatorname{op\text{-}neg}(d, t)(f) \in \operatorname{ops}(d, t) \end{array}}}
\vnbprop{op-alg-is-commutative-ring}{\vnbstmt{\begin{array}{@{}l@{}} \forall d, t. \\ \quad (\operatorname{is\text{-}hom\text{-}module}(d, d, t) \wedge \operatorname{is\text{-}commutative\text{-}ring}(\operatorname{scal}(d))) \Rightarrow  \\ \quad \quad \operatorname{is\text{-}commutative\text{-}ring}(\operatorname{op\text{-}alg}(d, t)) \end{array}}}
\vnbprop{ev-op-apply}{\vnbstmt{\begin{array}{@{}l@{}} \forall d, t.; \forall p \in \operatorname{carr}(\operatorname{poly}(\operatorname{scal}(d))). \\ \quad (\operatorname{is\text{-}hom\text{-}module}(d, d, t) \wedge \operatorname{is\text{-}commutative\text{-}ring}(\operatorname{scal}(d))) \Rightarrow  \\ \quad \quad \operatorname{ev\text{-}op}(d, t)(p) = \operatorname{op\text{-}of}(d, t, p) \end{array}}}
\vnbprop{ev-op-ring-hom}{\vnbstmt{\begin{array}{@{}l@{}} \forall d, t. \\ \quad (\operatorname{is\text{-}hom\text{-}module}(d, d, t) \wedge \operatorname{is\text{-}commutative\text{-}ring}(\operatorname{scal}(d))) \Rightarrow  \\ \quad \quad \operatorname{is\text{-}hom\text{-}commutative\text{-}ring}(\operatorname{poly}(\operatorname{scal}(d)), \operatorname{op\text{-}alg}(d, t), \operatorname{ev\text{-}op}(d, t)) \end{array}}}
\vnbprop{poly-op-vneg}{\vnbstmt{\begin{array}{@{}l@{}} \forall d, t.; \forall p \in \operatorname{carr}(\operatorname{poly}(\operatorname{scal}(d))), x \in \operatorname{vec}(d). \\ \quad (\operatorname{is\text{-}hom\text{-}module}(d, d, t) \wedge \operatorname{is\text{-}commutative\text{-}ring}(\operatorname{scal}(d))) \Rightarrow  \\ \quad \quad \operatorname{poly\text{-}op}(d, t, p, \operatorname{vneg}(d)(x)) = \operatorname{vneg}(d)(\operatorname{poly\text{-}op}(d, t, p, x)) \end{array}}}
\vnbprop{op-is-additive}{\vnbstmt{\begin{array}{@{}l@{}} \forall d, t.; \forall f \in \operatorname{ops}(d, t). \\ \quad (\operatorname{is\text{-}hom\text{-}module}(d, d, t) \wedge \operatorname{is\text{-}commutative\text{-}ring}(\operatorname{scal}(d))) \Rightarrow  \\ \quad \quad \operatorname{is\text{-}hom\text{-}abelian\text{-}group}(\operatorname{module\text{-}vector\text{-}ag}(d), \operatorname{module\text{-}vector\text{-}ag}(d), f) \end{array}}}
\vnbprop{op-eval-additive}{\vnbstmt{\begin{array}{@{}l@{}} \forall d, t.; \forall x \in \operatorname{vec}(d). \\ \quad (\operatorname{is\text{-}hom\text{-}module}(d, d, t) \wedge \operatorname{is\text{-}commutative\text{-}ring}(\operatorname{scal}(d))) \Rightarrow  \\ \quad \quad \operatorname{is\text{-}hom\text{-}abelian\text{-}group}(\operatorname{ring\text{-}additive\text{-}ag}(\operatorname{op\text{-}alg}(d, t)), \operatorname{module\text{-}vector\text{-}ag}(d), \\ \quad \quad \quad (\lambda \mathit{gex} \in \operatorname{ops}(d, t).\; \operatorname{gex}(x))) \end{array}}}
\vnbprop{mat-act-entry}{\vnbstmt{\begin{array}{@{}l@{}} \forall d, t.; \forall n \in \mathbb{N}, p \in \operatorname{mat}(n, n, \operatorname{ops}(d, t)), u \in \operatorname{mat}(n, 1, \operatorname{vec}(d)). \\ \quad (\operatorname{is\text{-}hom\text{-}module}(d, d, t) \wedge \operatorname{is\text{-}commutative\text{-}ring}(\operatorname{scal}(d))) \Rightarrow  \\ \quad \quad \operatorname{mat\text{-}act}(d, n, p, u) \in \operatorname{mat}(n, 1, \operatorname{vec}(d)) \wedge \\ \quad \quad \forall i \in \operatorname{interval}(1, n). \\ \quad \quad \quad \operatorname{entry}(\operatorname{mat\text{-}act}(d, n, p, u), i, 1) \\ \quad \quad \quad =\; \operatorname{finsum}(\operatorname{module\text{-}vector\text{-}ag}(d), \\ \quad \quad \quad \quad (\lambda k \in \operatorname{interval}(1, n).\; \operatorname{entry}(p, i, k)(\operatorname{entry}(u, k, 1))), \\ \quad \quad \quad \quad \operatorname{interval}(1, n)) \end{array}}}
\vnbprop{mat-act-assoc-entry}{\vnbstmt{\begin{array}{@{}l@{}} \forall d, t.; \forall n \in \mathbb{N}, p, q \in \operatorname{mat}(n, n, \operatorname{ops}(d, t)), u \in \operatorname{mat}(n, 1, \operatorname{vec}(d)), i \in \operatorname{interval}(1, n). \\ \quad (\operatorname{is\text{-}hom\text{-}module}(d, d, t) \wedge (\operatorname{is\text{-}commutative\text{-}ring}(\operatorname{scal}(d)) \wedge (1 \leq n))) \Rightarrow  \\ \quad \quad \operatorname{entry}(\operatorname{mat\text{-}act}(d, n, \operatorname{matmul}(\operatorname{op\text{-}alg}(d, t), p, q), u), i, 1) \\ \quad \quad =\; \operatorname{entry}(\operatorname{mat\text{-}act}(d, n, p, \operatorname{mat\text{-}act}(d, n, q, u)), i, 1) \end{array}}}
\vnbprop{op-iter-one}{\vnbstmt{\begin{array}{@{}l@{}} \forall d, t.; \forall w \in \operatorname{vec}(d). \\ \quad \operatorname{is\text{-}hom\text{-}module}(d, d, t) \Rightarrow (\operatorname{op\text{-}iter}(d, t, w, 1) = t(w)) \end{array}}}
\vnbprop{char-mat-kills-basis}{\vnbstmt{\begin{array}{@{}l@{}} \forall d, t.; \forall n \in \mathbb{N}, u \in \operatorname{mat}(n, 1, \operatorname{vec}(d)), a \in \operatorname{mat}(n, n, \operatorname{carr}(\operatorname{scal}(d))), k \in \operatorname{interval}(1, n). \\ \quad (\operatorname{is\text{-}hom\text{-}module}(d, d, t) \wedge (\operatorname{is\text{-}commutative\text{-}ring}(\operatorname{scal}(d)) \wedge \operatorname{is\text{-}matrix\text{-}of}(d, n, u, t, a))) \Rightarrow  \\ \quad \quad \operatorname{entry}(\operatorname{mat\text{-}act}(d, n, \operatorname{mat\text{-}map}(\operatorname{ev\text{-}op}(d, t), n, n, \operatorname{transpose}(\operatorname{char\text{-}mat}(\operatorname{scal}(d), a, n), n, n)), u), k, 1) \\ \quad \quad =\; \operatorname{vzero}(d) \end{array}}}
\vnbprop{cayley-hamilton-basis}{\vnbstmt{\begin{array}{@{}l@{}} \forall d, t.; \forall n, m \in \mathbb{N}, u \in \operatorname{mat}(n, 1, \operatorname{vec}(d)), a \in \operatorname{mat}(n, n, \operatorname{carr}(\operatorname{scal}(d))), j \in \operatorname{interval}(1, n). \\ \quad \operatorname{is\text{-}hom\text{-}module}(d, d, t) \wedge \\ \quad \operatorname{is\text{-}commutative\text{-}ring}(\operatorname{scal}(d)) \wedge \\ \quad n = (m + 1) \wedge \\ \quad \operatorname{is\text{-}matrix\text{-}of}(d, n, u, t, a) \Rightarrow \\ \quad \quad \operatorname{poly\text{-}op}(d, t, \operatorname{char\text{-}poly}(\operatorname{scal}(d), a, n), \operatorname{entry}(u, j, 1)) = \operatorname{vzero}(d) \end{array}}}
\vnbprop{cayley-hamilton}{\vnbstmt{\begin{array}{@{}l@{}} \forall d, t.; \forall n \in \mathbb{N}, u \in \operatorname{mat}(n, 1, \operatorname{vec}(d)), a \in \operatorname{mat}(n, n, \operatorname{carr}(\operatorname{scal}(d))), w \in \operatorname{vec}(d). \\ \quad \operatorname{is\text{-}hom\text{-}module}(d, d, t) \wedge \\ \quad \operatorname{is\text{-}commutative\text{-}ring}(\operatorname{scal}(d)) \wedge \\ \quad \operatorname{generates}(d, n, u) \wedge \\ \quad \operatorname{is\text{-}matrix\text{-}of}(d, n, u, t, a) \Rightarrow \\ \quad \quad \operatorname{poly\text{-}op}(d, t, \operatorname{char\text{-}poly}(\operatorname{scal}(d), a, n), w) = \operatorname{vzero}(d) \end{array}}}

\vnbbookfile{theorem-library/poly-eval-mul}{theorem-library/poly-eval-mul}
\noindent{\small poly-eval-mul.scm -- EVALUATION OF POLYNOMIALS IS MULTIPLICATIVE, in any commutative ring. Agent CCR-1111B, 2026-11-11.  Destination: theorem-library/poly-eval-mul.scm.  Helper prefix pem-. Library: algebra (POLY, POLY-EVAL, POLY-OP, RING-AS-MODULE only).}\par
\vnbprop{pem-mvag}{\vnbstmt{\begin{array}{@{}l@{}} \forall r. \\ \quad \operatorname{is\text{-}ring}(r) \Rightarrow (\operatorname{module\text{-}vector\text{-}ag}(\operatorname{ring\text{-}as\text{-}module}(r)) \simeq \operatorname{ring\text{-}additive\text{-}ag}(r)) \end{array}}}
\vnbprop{pem-mult-hom}{\vnbstmt{\begin{array}{@{}l@{}} \forall r.; \forall t \in \operatorname{carr}(r). \\ \quad \operatorname{is\text{-}commutative\text{-}ring}(r) \Rightarrow  \\ \quad \quad \operatorname{is\text{-}hom\text{-}module}(\operatorname{ring\text{-}as\text{-}module}(r), \operatorname{ring\text{-}as\text{-}module}(r), (\lambda y \in \operatorname{carr}(r).\; \operatorname{mul}(r)(t, y))) \end{array}}}
\vnbprop{pem-iter}{\vnbstmt{\begin{array}{@{}l@{}} \forall k \in \mathbb{N}, r.; \forall t, y \in \operatorname{carr}(r). \\ \quad \operatorname{is\text{-}commutative\text{-}ring}(r) \Rightarrow  \\ \quad \quad \operatorname{op\text{-}iter}(\operatorname{ring\text{-}as\text{-}module}(r), (\lambda a \in \operatorname{carr}(r).\; \operatorname{mul}(r)(t, a)), y, k) \\ \quad \quad =\; \operatorname{mul}(r)(\operatorname{ring\text{-}power}(r, t, k), y) \end{array}}}
\vnbprop{pem-poly-op-one}{\vnbstmt{\begin{array}{@{}l@{}} \forall r.; \forall t \in \operatorname{carr}(r), p \in \operatorname{carr}(\operatorname{poly}(r)). \\ \quad \operatorname{is\text{-}commutative\text{-}ring}(r) \Rightarrow  \\ \quad \quad \operatorname{poly\text{-}op}(\operatorname{ring\text{-}as\text{-}module}(r), (\lambda y \in \operatorname{carr}(r).\; \operatorname{mul}(r)(t, y)), p, \operatorname{one}(r)) = \operatorname{poly\text{-}eval}(r, p, t) \end{array}}}
\vnbprop{poly-eval-mul}{\vnbstmt{\begin{array}{@{}l@{}} \forall r.; \forall t \in \operatorname{carr}(r), p, q \in \operatorname{carr}(\operatorname{poly}(r)). \\ \quad \operatorname{is\text{-}commutative\text{-}ring}(r) \Rightarrow  \\ \quad \quad \operatorname{poly\text{-}eval}(r, \operatorname{mul}(\operatorname{poly}(r))(p, q), t) = \operatorname{mul}(r)(\operatorname{poly\text{-}eval}(r, p, t), \operatorname{poly\text{-}eval}(r, q, t)) \end{array}}}

\part{Analysis}
\noindent Metric spaces and their topology, completeness, compactness and continuity; normed fields and normed vector spaces; sequences and series of real and complex numbers; the calculus of one real variable on an interval: derivatives, the mean value theorem, Taylor, antiderivatives, regulated functions and the integral; roads and line integrals; Ascoli, Baire, Hahn-Banach, Bernstein.
\par\noindent\textit{Builds on:} Linear Algebra.
\par\noindent{\small 573 files.}

\vnbbookfile{structure-library/metric-space}{structure-library/metric-space}
\noindent{\small metric-space.scm -- METRIC-SPACE structure}\par
\vnbstructdef{metric-space}{A metric space $s$ is a 2-tuple whose components are listed below.\par
\begin{enumerate}[nosep]
\item $\operatorname{pts}(s) = \operatorname{nth}(1, s)$, a carrier (\texttt{carriers}): $\operatorname{pts}(s) \in \mathbf{Set}$.
\item $\operatorname{dist}(s) = \operatorname{nth}(2, s)$, an operation (\texttt{op}): $\operatorname{dist}(s) \in ((\operatorname{pts}(s) \times \operatorname{pts}(s)) \to \mathbb{R})$.
\end{enumerate}
The laws (the \texttt{property} and \texttt{law} clauses):
\begin{itemize}[nosep]
\item $\operatorname{is\text{-}metric}(\operatorname{dist}(s), \operatorname{pts}(s))$
\end{itemize}
The characteristic law (the axiom \texttt{is-metric-space}): $\operatorname{is\text{-}metric\text{-}space}(s)$ holds if and only if $\operatorname{length}(s) = 2$, every component has the typing listed with it, every law holds, and $s \in \mathbf{Set}$ (the structure is a set).
}
\par\medskip\noindent\textbf{Definition} (\texttt{hom-metric-space}).\;$\operatorname{hom\text{-}metric\text{-}space}(a, b) :=$
\begin{equation*}
\{\operatorname{homf\_} \in (\operatorname{pts}(a) \to \operatorname{pts}(b)) : \operatorname{is\text{-}hom\text{-}metric\text{-}space}(a, b, \operatorname{homf\_})\}
\end{equation*}

\vnbdef{metric-space-class}{\vnbstmt{\forall s.\; ((s \in \operatorname{metric\text{-}space}) \Leftrightarrow \operatorname{is\text{-}metric\text{-}space}(s))}}
\vnbdef{is-hom-metric-space-def}{\vnbstmt{\begin{array}{@{}l@{}} \forall a, b, f. \\ \quad \operatorname{is\text{-}hom\text{-}metric\text{-}space}(a, b, f) \\ \quad \Leftrightarrow\; \operatorname{is\text{-}metric\text{-}space}(a) \wedge \\ \quad \operatorname{is\text{-}metric\text{-}space}(b) \wedge \\ \quad f \in (\operatorname{pts}(a) \to \operatorname{pts}(b)) \wedge \\ \quad \forall x_{1}, x_{2} \in \operatorname{pts}(a). \\ \quad \quad \operatorname{dist}(a)(x_{1}, x_{2}) = \operatorname{dist}(b)(f(x_{1}), f(x_{2})) \end{array}}}

\vnbbookfile{structure-library/metric-topology}{structure-library/metric-topology}
\noindent{\small metric-topology.scm -- balls, r-nets, total-boundedness.}\par
\par\medskip\noindent\textbf{Definition} (\texttt{ball}).\;$\operatorname{ball}(s, c, r) :=$
\begin{equation*}
\{y \in \operatorname{pts}(s) : ((\operatorname{dist}(s)(c, y) \leq r) \wedge \neg (\operatorname{dist}(s)(c, y) = r))\}
\end{equation*}

\par\medskip\noindent\textbf{Definition} (\texttt{is-r-net}).\;For all $s$, $f$, $a$ and $r$, 
\begin{equation*}
\begin{array}{@{}l@{}} \operatorname{is\text{-}r\text{-}net}(s, f, a, r) \\ \Leftrightarrow\; f \subseteq a \wedge \\ \forall p \in a. \exists c \in f. \\ \quad (\operatorname{dist}(s)(c, p) \leq r) \wedge \neg (\operatorname{dist}(s)(c, p) = r) \end{array}
\end{equation*}

\par\medskip\noindent\textbf{Definition} (\texttt{totally-bounded}).\;For every $s$, 
\begin{equation*}
\begin{array}{@{}l@{}} \operatorname{totally\text{-}bounded}(s) \\ \Leftrightarrow\; \operatorname{is\text{-}metric\text{-}space}(s) \wedge \\ \forall r. \\ \quad ((r \in \mathbb{R}) \wedge ((0 \leq r) \wedge \neg (0 = r))) \Rightarrow  \\ \quad \quad \exists f.\; ((|f| \in \mathbb{N}) \wedge \operatorname{is\text{-}r\text{-}net}(s, f, \operatorname{pts}(s), r)) \end{array}
\end{equation*}

\par\medskip\noindent\textbf{Definition} (\texttt{closed-ball}).\;$\operatorname{closed\text{-}ball}(s, c, r) :=$ $\{y \in \operatorname{pts}(s) : (\operatorname{dist}(s)(c, y) \leq r)\}$

\vnbbookfile{structure-library/metric-completeness}{structure-library/metric-completeness}
\noindent{\small metric-completeness.scm -- Cauchy sequences, convergence, completeness on a generic METRIC-SPACE.}\par
\par\medskip\noindent\textbf{Definition} (\texttt{is-cauchy-seq}).\;For all $s$ and $f$, 
\begin{equation*}
\begin{array}{@{}l@{}} \operatorname{is\text{-}cauchy\text{-}seq}(s, f) \\ \Leftrightarrow\; \operatorname{is\text{-}metric\text{-}space}(s) \wedge \\ f \in (\mathbb{N} \to \operatorname{pts}(s)) \wedge \\ \forall \mathit{eps}. \\ \quad \operatorname{pos\text{-}rr}(\mathit{eps}) \Rightarrow  \\ \quad \quad \exists n \in \mathbb{N}. \forall m, \operatorname{n\_} \in \mathbb{N}. \\ \quad \quad \quad ((n \leq m) \wedge (n \leq \operatorname{n\_})) \Rightarrow (\operatorname{dist}(s)(f(m), f(\operatorname{n\_})) \leq \mathit{eps}) \end{array}
\end{equation*}

\par\medskip\noindent\textbf{Definition} (\texttt{converges-to}).\;For all $s$, $f$ and $l$, 
\begin{equation*}
\begin{array}{@{}l@{}} \operatorname{converges\text{-}to}(s, f, l) \\ \Leftrightarrow\; \operatorname{is\text{-}metric\text{-}space}(s) \wedge \\ f \in (\mathbb{N} \to \operatorname{pts}(s)) \wedge \\ l \in \operatorname{pts}(s) \wedge \\ \forall \mathit{eps}. \\ \quad \operatorname{pos\text{-}rr}(\mathit{eps}) \Rightarrow  \\ \quad \quad \exists n \in \mathbb{N}. \forall \operatorname{n\_} \in \mathbb{N}. \\ \quad \quad \quad (n \leq \operatorname{n\_}) \Rightarrow (\operatorname{dist}(s)(f(\operatorname{n\_}), l) \leq \mathit{eps}) \end{array}
\end{equation*}

\par\medskip\noindent\textbf{Definition} (\texttt{converges}).\;For all $s$ and $f$,  $\operatorname{converges}(s, f) \Leftrightarrow \exists l.\; \operatorname{converges\text{-}to}(s, f, l)$
\par\medskip\noindent\textbf{Definition} (\texttt{is-complete}).\;For every $s$, 
\begin{equation*}
\operatorname{is\text{-}complete}(s) \Leftrightarrow (\operatorname{is\text{-}metric\text{-}space}(s) \wedge \forall f.\; (\operatorname{is\text{-}cauchy\text{-}seq}(s, f) \Rightarrow \operatorname{converges}(s, f)))
\end{equation*}

\vnbdef{complete-cauchy-converges}{\vnbstmt{\forall s, f.\; ((\operatorname{is\text{-}complete}(s) \wedge \operatorname{is\text{-}cauchy\text{-}seq}(s, f)) \Rightarrow \operatorname{converges}(s, f))}}

\vnbbookfile{structure-library/metric-continuity}{structure-library/metric-continuity}
\noindent{\small RETIRED 2026-09-17 (proven): metric-hom-is-continuous -- theorem-library/rake-analysis2.scm metric-continuity.scm -- continuous maps between metric spaces.}\par
\par\medskip\noindent\textbf{Definition} (\texttt{is-continuous-at}).\;For all $s$, $t$, $f$ and $a$, 
\begin{equation*}
\begin{array}{@{}l@{}} \operatorname{is\text{-}continuous\text{-}at}(s, t, f, a) \\ \Leftrightarrow\; \operatorname{is\text{-}metric\text{-}space}(s) \wedge \\ \operatorname{is\text{-}metric\text{-}space}(t) \wedge \\ f \in (\operatorname{pts}(s) \to \operatorname{pts}(t)) \wedge \\ a \in \operatorname{pts}(s) \wedge \\ \forall \mathit{eps}. \\ \quad \operatorname{pos\text{-}rr}(\mathit{eps}) \Rightarrow  \\ \quad \quad \exists \delta. \\ \quad \quad \quad \operatorname{pos\text{-}rr}(\delta) \wedge \\ \quad \quad \quad \forall b \in \operatorname{pts}(s). \\ \quad \quad \quad \quad (\operatorname{dist}(s)(a, b) \leq \delta) \Rightarrow (\operatorname{dist}(t)(f(a), f(b)) \leq \mathit{eps}) \end{array}
\end{equation*}

\par\medskip\noindent\textbf{Definition} (\texttt{is-continuous}).\;For all $s$, $t$ and $f$, 
\begin{equation*}
\begin{array}{@{}l@{}} \operatorname{is\text{-}continuous}(s, t, f) \\ \Leftrightarrow\; \operatorname{is\text{-}metric\text{-}space}(s) \wedge \\ \operatorname{is\text{-}metric\text{-}space}(t) \wedge \\ f \in (\operatorname{pts}(s) \to \operatorname{pts}(t)) \wedge \\ \forall a \in \operatorname{pts}(s).\; \operatorname{is\text{-}continuous\text{-}at}(s, t, f, a) \end{array}
\end{equation*}

\par\medskip\noindent\textbf{Definition} (\texttt{is-uniformly-continuous}).\;For all $s$, $t$ and $f$, 
\begin{equation*}
\begin{array}{@{}l@{}} \operatorname{is\text{-}uniformly\text{-}continuous}(s, t, f) \\ \Leftrightarrow\; \operatorname{is\text{-}metric\text{-}space}(s) \wedge \\ \operatorname{is\text{-}metric\text{-}space}(t) \wedge \\ f \in (\operatorname{pts}(s) \to \operatorname{pts}(t)) \wedge \\ \forall \mathit{eps}. \\ \quad \operatorname{pos\text{-}rr}(\mathit{eps}) \Rightarrow  \\ \quad \quad \exists \delta. \\ \quad \quad \quad \operatorname{pos\text{-}rr}(\delta) \wedge \\ \quad \quad \quad \forall a, b \in \operatorname{pts}(s). \\ \quad \quad \quad \quad (\operatorname{dist}(s)(a, b) \leq \delta) \Rightarrow (\operatorname{dist}(t)(f(a), f(b)) \leq \mathit{eps}) \end{array}
\end{equation*}

\par\medskip\noindent\textbf{Definition} (\texttt{is-isometry}).\;For all $s$, $t$ and $f$, 
\begin{equation*}
\begin{array}{@{}l@{}} \operatorname{is\text{-}isometry}(s, t, f) \\ \Leftrightarrow\; \operatorname{is\text{-}metric\text{-}space}(s) \wedge \\ \operatorname{is\text{-}metric\text{-}space}(t) \wedge \\ f \in (\operatorname{pts}(s) \to \operatorname{pts}(t)) \wedge \\ \forall x, y \in \operatorname{pts}(s).\; (\operatorname{dist}(t)(f(x), f(y)) = \operatorname{dist}(s)(x, y)) \end{array}
\end{equation*}


\vnbbookfile{structure-library/metric-open-sets}{structure-library/metric-open-sets}
\noindent{\small RETIRED 2026-09-17 (proven): continuous-implies-open-preimage -- theorem-library/rake-cont-preimage.scm (modulo 0, via dk-halve!; this file now holds no support at all). RETIRED 2026-09-17 (proven): empty-is-open, carrier-is-open, union-of-opens-open, inter-of-opens-open, preimage-complement, ...}\par
\par\medskip\noindent\textbf{Definition} (\texttt{is-open}).\;For all $s$ and $u$, 
\begin{equation*}
\begin{array}{@{}l@{}} \operatorname{is\text{-}open}(s, u) \\ \Leftrightarrow\; \operatorname{is\text{-}metric\text{-}space}(s) \wedge \\ u \subseteq \operatorname{pts}(s) \wedge \\ \forall y \in u.\; \exists r.\; (\operatorname{pos\text{-}rr}(r) \wedge (\operatorname{ball}(s, y, r) \subseteq u)) \end{array}
\end{equation*}

\par\medskip\noindent\textbf{Definition} (\texttt{preimage}).\;$\operatorname{preimage}(s, f, v) :=$ $\{a \in \operatorname{pts}(s) : (f(a) \in v)\}$
\par\medskip\noindent\textbf{Definition} (\texttt{is-closed}).\;For all $s$ and $a$, 
\begin{equation*}
\operatorname{is\text{-}closed}(s, a) \Leftrightarrow (\operatorname{is\text{-}metric\text{-}space}(s) \wedge ((a \subseteq \operatorname{pts}(s)) \wedge \operatorname{is\text{-}open}(s, \operatorname{complement\text{-}in}(\operatorname{pts}(s), a))))
\end{equation*}

\vnbdef{preimage-membership}{\vnbstmt{\forall s, f, v, a.\; ((a \in \operatorname{preimage}(s, f, v)) \Leftrightarrow ((a \in \operatorname{pts}(s)) \wedge (f(a) \in v)))}}

\vnbbookfile{structure-library/metric-connected}{structure-library/metric-connected}
\noindent{\small metric-connected.scm -- CONNECTEDNESS OF A SUBSET OF A METRIC SPACE. DEFINITION ONLY; the laws are PROVEN in theorem-library/metric-connected-laws.scm. Agent CA-5 of week 3 of the October roadmap, 2026-10-04 (a NEW DEFINITION, the user's to confirm).}\par
\par\medskip\noindent\textbf{Definition} (\texttt{is-connected}).\;For all $\operatorname{cns\_}$ and $\operatorname{cna\_}$, 
\begin{equation*}
\begin{array}{@{}l@{}} \operatorname{is\text{-}connected}(\operatorname{cns\_}, \operatorname{cna\_}) \\ \Leftrightarrow\; \operatorname{is\text{-}metric\text{-}space}(\operatorname{cns\_}) \wedge \\ \operatorname{cna\_} \subseteq \operatorname{pts}(\operatorname{cns\_}) \wedge \\ \forall \operatorname{cnu\_}. \\ \quad \operatorname{is\text{-}open}(\operatorname{cns\_}, \operatorname{cnu\_}) \Rightarrow  \\ \quad \quad \forall \operatorname{cnv\_}. \\ \quad \quad \quad \operatorname{is\text{-}open}(\operatorname{cns\_}, \operatorname{cnv\_}) \Rightarrow  \\ \quad \quad \quad (\operatorname{cna\_} \subseteq (\operatorname{cnu\_} \cup \operatorname{cnv\_})) \Rightarrow  \\ \quad \quad \quad \forall \operatorname{cnx\_} \in \operatorname{cna\_}.\; \neg ((\operatorname{cnx\_} \in \operatorname{cnu\_}) \wedge (\operatorname{cnx\_} \in \operatorname{cnv\_})) \Rightarrow  \\ \quad \quad \quad \exists \operatorname{cnx\_} \in \operatorname{cna\_}.\; (\operatorname{cnx\_} \in \operatorname{cnu\_}) \Rightarrow  \\ \quad \quad \quad \quad \neg \exists \operatorname{cnx\_} \in \operatorname{cna\_}.\; (\operatorname{cnx\_} \in \operatorname{cnv\_}) \end{array}
\end{equation*}


\vnbbookfile{structure-library/top-space}{structure-library/top-space}
\noindent{\small top-space.scm -- TOPOLOGICAL SPACES, and the functor Met -$>$ Top}\par
\vnbstructdef{top-space}{A top space $s$ is a 2-tuple whose components are listed below.\par
\begin{enumerate}[nosep]
\item $\operatorname{pts}(s) = \operatorname{nth}(1, s)$, a carrier (\texttt{carriers}): $\operatorname{pts}(s) \in \mathbf{Set}$.
\item $\operatorname{opens}(s) = \operatorname{nth}(2, s)$, a family of subsets (\texttt{family}): $\operatorname{opens}(s) \in \operatorname{power}(\operatorname{power}(\operatorname{pts}(s)))$.
\end{enumerate}
The laws (the \texttt{property} and \texttt{law} clauses):
\begin{itemize}[nosep]
\item $\emptyset \in \operatorname{opens}(s)$
\item $\operatorname{pts}(s) \in \operatorname{opens}(s)$
\item $\forall u, v \in \operatorname{opens}(s).\; ((u \cap v) \in \operatorname{opens}(s))$
\item $\forall \mathit{fam} \in \mathbf{Set}.\; ((\mathit{fam} \subseteq \operatorname{opens}(s)) \Rightarrow (\operatorname{big\text{-}union}(u, \mathit{fam}, u) \in \operatorname{opens}(s)))$
\end{itemize}
The characteristic law (the axiom \texttt{is-top-space}): $\operatorname{is\text{-}top\text{-}space}(s)$ holds if and only if $\operatorname{length}(s) = 2$, every component has the typing listed with it, every law holds, and $s \in \mathbf{Set}$ (the structure is a set).
}
\par\medskip\noindent\textbf{Definition} (\texttt{metric-top}).\;$\operatorname{metric\text{-}top}(\mathit{md}) :=$ ${[}\operatorname{pts}(\mathit{md}), \{u \in \operatorname{power}(\operatorname{pts}(\mathit{md})) : \operatorname{is\text{-}open}(\mathit{md}, u)\}]$
\par\medskip\noindent\textbf{Definition} (\texttt{is-hausdorff}).\;For every $s$, 
\begin{equation*}
\begin{array}{@{}l@{}} \operatorname{is\text{-}hausdorff}(s) \\ \Leftrightarrow\; \operatorname{is\text{-}top\text{-}space}(s) \wedge \\ \forall a, b \in \operatorname{pts}(s). \\ \quad \neg (a \simeq b) \Rightarrow  \\ \quad \quad \exists u, v \in \operatorname{opens}(s). \\ \quad \quad \quad (a \in u) \wedge ((b \in v) \wedge ((u \cap v) \simeq \emptyset)) \end{array}
\end{equation*}

\par\medskip\noindent\textbf{Definition} (\texttt{hom-top-space}).\;$\operatorname{hom\text{-}top\text{-}space}(a, b) :=$
\begin{equation*}
\{\operatorname{homf\_} \in (\operatorname{pts}(a) \to \operatorname{pts}(b)) : \operatorname{is\text{-}hom\text{-}top\text{-}space}(a, b, \operatorname{homf\_})\}
\end{equation*}

\par\medskip\noindent\textbf{Definition} (\texttt{hom-metrizable-top-space}).\;$\operatorname{hom\text{-}metrizable\text{-}top\text{-}space}(a, b) :=$
\begin{equation*}
\{\operatorname{homf\_} \in (\operatorname{pts}(a) \to \operatorname{pts}(b)) : \operatorname{is\text{-}hom\text{-}metrizable\text{-}top\text{-}space}(a, b, \operatorname{homf\_})\}
\end{equation*}

\vnbdef{top-space-class}{\vnbstmt{\forall s.\; ((s \in \operatorname{top\text{-}space}) \Leftrightarrow \operatorname{is\text{-}top\text{-}space}(s))}}
\vnbdef{is-hom-metrizable-top-space-def}{\vnbstmt{\begin{array}{@{}l@{}} \forall s, t, f. \\ \quad \operatorname{is\text{-}hom\text{-}metrizable\text{-}top\text{-}space}(s, t, f) \\ \quad \Leftrightarrow\; \operatorname{is\text{-}metrizable\text{-}top\text{-}space}(s) \wedge \\ \quad \operatorname{is\text{-}metrizable\text{-}top\text{-}space}(t) \wedge \\ \quad f \in (\operatorname{pts}(s) \to \operatorname{pts}(t)) \wedge \\ \quad \forall u \in \operatorname{opens}(t).\; (\operatorname{preimage}(s, f, u) \in \operatorname{opens}(s)) \end{array}}}
\vnbdef{is-hom-top-space-def}{\vnbstmt{\begin{array}{@{}l@{}} \forall s, t, f. \\ \quad \operatorname{is\text{-}hom\text{-}top\text{-}space}(s, t, f) \\ \quad \Leftrightarrow\; \operatorname{is\text{-}top\text{-}space}(s) \wedge \\ \quad \operatorname{is\text{-}top\text{-}space}(t) \wedge \\ \quad f \in (\operatorname{pts}(s) \to \operatorname{pts}(t)) \wedge \\ \quad \forall u \in \operatorname{opens}(t).\; (\operatorname{preimage}(s, f, u) \in \operatorname{opens}(s)) \end{array}}}
\vnbdef{is-metrizable-top-space-def}{\vnbstmt{\begin{array}{@{}l@{}} \forall s. \\ \quad \operatorname{is\text{-}metrizable\text{-}top\text{-}space}(s) \\ \quad \Leftrightarrow\; \operatorname{is\text{-}top\text{-}space}(s) \wedge \\ \quad \exists d \in \mathbf{Set}. \\ \quad \quad \operatorname{is\text{-}metric\text{-}space}(d) \wedge (s \simeq {[}\operatorname{pts}(d), \{u \in \operatorname{power}(\operatorname{pts}(d)) : \operatorname{is\text{-}open}(d, u)\}]) \end{array}}}
\vnbdef{metrizable-top-space-class}{\vnbstmt{\forall s.\; ((s \in \operatorname{metrizable\text{-}top\text{-}space}) \Leftrightarrow \operatorname{is\text{-}metrizable\text{-}top\text{-}space}(s))}}

\vnbbookfile{structure-library/compactness}{structure-library/compactness}
\noindent{\small compactness.scm -- compactness of a metric space and the four-way characterization (calculus.pdf Prop 3.12).  Supplies the missing vocabulary (IS-COMPACT, IS-OPEN-COVER, CLUSTER-POINT, HAS-FIP) and states the characterization as warranted supports (library phase); the proof}\par
\par\medskip\noindent\textbf{Definition} (\texttt{is-compact}).\;For every $s$, 
\begin{equation*}
\begin{array}{@{}l@{}} \operatorname{is\text{-}compact}(s) \\ \Leftrightarrow\; \operatorname{is\text{-}metric\text{-}space}(s) \wedge \\ \forall c \in \mathbf{Set}. \\ \quad \operatorname{is\text{-}open\text{-}cover}(s, c) \Rightarrow  \\ \quad \quad \exists f \in \mathbf{Set}.\; ((f \subseteq c) \wedge ((|f| \in \mathbb{N}) \wedge \operatorname{is\text{-}open\text{-}cover}(s, f))) \end{array}
\end{equation*}

\par\medskip\noindent\textbf{Definition} (\texttt{is-open-cover}).\;For all $s$ and $c$, 
\begin{equation*}
\begin{array}{@{}l@{}} \operatorname{is\text{-}open\text{-}cover}(s, c) \\ \Leftrightarrow\; (\operatorname{is\text{-}metric\text{-}space}(s) \wedge (\forall u \in c.\; \operatorname{is\text{-}open}(s, u) \wedge (\operatorname{big\text{-}union}(u, c, u) \simeq \operatorname{pts}(s)))) \end{array}
\end{equation*}

\par\medskip\noindent\textbf{Definition} (\texttt{cluster-point}).\;For all $s$, $f$ and $x$, 
\begin{equation*}
\begin{array}{@{}l@{}} \operatorname{cluster\text{-}point}(s, f, x) \\ \Leftrightarrow\; \operatorname{is\text{-}metric\text{-}space}(s) \wedge \\ f \in (\mathbb{N} \to \operatorname{pts}(s)) \wedge \\ x \in \operatorname{pts}(s) \wedge \\ \forall \mathit{eps}. \\ \quad \operatorname{pos\text{-}rr}(\mathit{eps}) \Rightarrow  \\ \quad \quad \forall m \in \mathbb{N}.\; \exists n \in \mathbb{N}.\; ((m \leq n) \wedge (\operatorname{dist}(s)(f(n), x) < \mathit{eps})) \end{array}
\end{equation*}

\par\medskip\noindent\textbf{Definition} (\texttt{has-fip}).\;For all $s$ and $c$, 
\begin{equation*}
\begin{array}{@{}l@{}} \operatorname{has\text{-}fip}(s, c) \\ \Leftrightarrow\; \operatorname{is\text{-}metric\text{-}space}(s) \wedge \\ \exists a_{0}.\; (a_{0} \in c) \wedge \\ \forall a \in c.\; \operatorname{is\text{-}closed}(s, a) \wedge \\ \forall f. \\ \quad (f \subseteq c) \Rightarrow  \\ \quad (|f| \in \mathbb{N}) \Rightarrow  \\ \quad \exists a_{0}.\; (a_{0} \in f) \Rightarrow \exists p.\; (p \in \operatorname{intersection\text{-}of}(f)) \end{array}
\end{equation*}

\par\medskip\noindent\textbf{Definition} (\texttt{ball-cover}).\;$\operatorname{ball\text{-}cover}(s, r) :=$ $\operatorname{image}((\lambda c \in \operatorname{pts}(s).\; \operatorname{ball}(s, c, r)), \operatorname{pts}(s))$
\par\medskip\noindent\textbf{Definition} (\texttt{centres}).\;$\operatorname{centres}(s, b, r) :=$ $\{c \in \operatorname{pts}(s) : (\operatorname{ball}(s, c, r) = b)\}$
\par\medskip\noindent\textbf{Definition} (\texttt{centre-set}).\;$\operatorname{centre\text{-}set}(s, r, f) :=$ $\operatorname{image}((\lambda b \in f.\; \varepsilon(\operatorname{centres}(s, b, r))), f)$

\vnbbookfile{structure-library/top-continuity}{structure-library/top-continuity}
\noindent{\small top-continuity.scm -- COMPACTNESS of a topological space, and the SUBSPACE TOPOLOGY.  DEFINITIONS AND NOTATION ONLY; every law is proven in theorem-library/metrizable-compact-image.scm.}\par
\par\medskip\noindent\textbf{Definition} (\texttt{is-compact-t}).\;For every $s$, 
\begin{equation*}
\begin{array}{@{}l@{}} \operatorname{is\text{-}compact\text{-}t}(s) \\ \Leftrightarrow\; \operatorname{is\text{-}top\text{-}space}(s) \wedge \\ \forall c. \\ \quad \forall u \in c.\; (u \in \operatorname{opens}(s)) \Rightarrow  \\ \quad (\operatorname{big\text{-}union}(u, c, u) \simeq \operatorname{pts}(s)) \Rightarrow  \\ \quad \quad \exists f.\; ((f \subseteq c) \wedge ((|f| \in \mathbb{N}) \wedge (\operatorname{big\text{-}union}(u, f, u) \simeq \operatorname{pts}(s)))) \end{array}
\end{equation*}

\par\medskip\noindent\textbf{Definition} (\texttt{subspace-top}).\;$\operatorname{subspace\text{-}top}(t, a) :=$
\begin{equation*}
{[}a, \{\operatorname{stw\_} \in \operatorname{power}(a) : \exists \operatorname{stu\_} \in \operatorname{opens}(t).\; (\operatorname{stw\_} = (\operatorname{stu\_} \cap a))\}]
\end{equation*}


\vnbbookfile{structure-library/separable}{structure-library/separable}
\noindent{\small separable.scm -- separability of a metric space (a dense sequence), the vocabulary and the one support the bridge compact =$>$ separable still needs. Added 2026-07-22 for the Ascoli-Arzela arc.}\par
\par\medskip\noindent\textbf{Definition} (\texttt{is-separable}).\;For every $s$, 
\begin{equation*}
\begin{array}{@{}l@{}} \operatorname{is\text{-}separable}(s) \\ \Leftrightarrow\; \operatorname{is\text{-}metric\text{-}space}(s) \wedge \\ \exists \mathit{dseq} \in (\mathbb{N} \to \operatorname{pts}(s)). \forall x \in \operatorname{pts}(s), \mathit{eps}. \\ \quad \operatorname{pos\text{-}rr}(\mathit{eps}) \Rightarrow \exists n \in \mathbb{N}.\; (\operatorname{dist}(s)(x, \operatorname{dseq}(n)) < \mathit{eps}) \end{array}
\end{equation*}


\vnbbookfile{structure-library/normed-field}{structure-library/normed-field}
\noindent{\small normed-field.scm -- NORMED-FIELD: a commutative ring with multiplicative inverses and a structurally-carried norm into RR.}\par
\vnbstructdef{normed-field}{A normed field $s$ is a 7-tuple whose components are listed below.\par
\begin{enumerate}[nosep]
\item $\operatorname{carr}(s) = \operatorname{nth}(1, s)$, a carrier (\texttt{carriers}): $\operatorname{carr}(s) \in \mathbf{Set}$.
\item $\operatorname{add}(s) = \operatorname{nth}(2, s)$, an operation (\texttt{op}): $\operatorname{add}(s) \in ((\operatorname{carr}(s) \times \operatorname{carr}(s)) \to \operatorname{carr}(s))$.
\item $\operatorname{mul}(s) = \operatorname{nth}(3, s)$, an operation (\texttt{op}): $\operatorname{mul}(s) \in ((\operatorname{carr}(s) \times \operatorname{carr}(s)) \to \operatorname{carr}(s))$.
\item $\operatorname{neg}(s) = \operatorname{nth}(4, s)$, an operation (\texttt{op}): $\operatorname{neg}(s) \in (\operatorname{carr}(s) \to \operatorname{carr}(s))$.
\item $\operatorname{zero}(s) = \operatorname{nth}(5, s)$, a constant (\texttt{constant}): $\operatorname{zero}(s) \in \operatorname{carr}(s)$.
\item $\operatorname{one}(s) = \operatorname{nth}(6, s)$, a constant (\texttt{constant}): $\operatorname{one}(s) \in \operatorname{carr}(s)$.
\item $\operatorname{fnrm}(s) = \operatorname{nth}(7, s)$, an operation (\texttt{op}): $\operatorname{fnrm}(s) \in (\operatorname{carr}(s) \to \mathbb{R})$.
\end{enumerate}
The laws (the \texttt{property} and \texttt{law} clauses):
\begin{itemize}[nosep]
\item $\operatorname{is\text{-}associative}(\operatorname{add}(s), \operatorname{carr}(s))$
\item $\operatorname{is\text{-}commutative}(\operatorname{add}(s), \operatorname{carr}(s))$
\item $\operatorname{is\text{-}identity}(\operatorname{add}(s), \operatorname{zero}(s), \operatorname{carr}(s))$
\item $\operatorname{has\text{-}inverses}(\operatorname{add}(s), \operatorname{zero}(s), \operatorname{neg}(s), \operatorname{carr}(s))$
\item $\operatorname{is\text{-}associative}(\operatorname{mul}(s), \operatorname{carr}(s))$
\item $\operatorname{is\text{-}commutative}(\operatorname{mul}(s), \operatorname{carr}(s))$
\item $\operatorname{is\text{-}identity}(\operatorname{mul}(s), \operatorname{one}(s), \operatorname{carr}(s))$
\item $\operatorname{is\text{-}distributive}(\operatorname{add}(s), \operatorname{mul}(s), \operatorname{carr}(s))$
\item $\operatorname{is\text{-}norm}(\operatorname{fnrm}(s), \operatorname{add}(s), \operatorname{mul}(s), \operatorname{zero}(s), \operatorname{carr}(s))$
\item $\neg (\operatorname{zero}(s) = \operatorname{one}(s))$
\item
\begin{equation*}
\begin{array}{@{}l@{}} \forall \operatorname{a\_} \in \operatorname{carr}(s). \\ \quad \neg (\operatorname{a\_} = \operatorname{zero}(s)) \Rightarrow  \\ \quad \quad \exists \operatorname{b\_} \in \operatorname{carr}(s).\; (\neg (\operatorname{b\_} = \operatorname{zero}(s)) \wedge (\operatorname{mul}(s)(\operatorname{a\_}, \operatorname{b\_}) = \operatorname{one}(s))) \end{array}
\end{equation*}

\end{itemize}
The characteristic law (the axiom \texttt{is-normed-field}): $\operatorname{is\text{-}normed\text{-}field}(s)$ holds if and only if $\operatorname{length}(s) = 7$, every component has the typing listed with it, every law holds, and $s \in \mathbf{Set}$ (the structure is a set).
}
\par\medskip\noindent\textbf{Definition} (\texttt{hom-normed-field}).\;$\operatorname{hom\text{-}normed\text{-}field}(a, b) :=$
\begin{equation*}
\{\operatorname{homf\_} \in (\operatorname{carr}(a) \to \operatorname{carr}(b)) : \operatorname{is\text{-}hom\text{-}normed\text{-}field}(a, b, \operatorname{homf\_})\}
\end{equation*}

\vnbdef{normed-field-class}{\vnbstmt{\forall s.\; ((s \in \operatorname{normed\text{-}field}) \Leftrightarrow \operatorname{is\text{-}normed\text{-}field}(s))}}
\vnbdef{is-hom-normed-field-def}{\vnbstmt{\begin{array}{@{}l@{}} \forall a, b, f. \\ \quad \operatorname{is\text{-}hom\text{-}normed\text{-}field}(a, b, f) \\ \quad \Leftrightarrow\; \operatorname{is\text{-}normed\text{-}field}(a) \wedge \\ \quad \operatorname{is\text{-}normed\text{-}field}(b) \wedge \\ \quad f \in (\operatorname{carr}(a) \to \operatorname{carr}(b)) \wedge \\ \quad \forall x_{1}, x_{2} \in \operatorname{carr}(a). \\ \quad \quad f(\operatorname{add}(a)(x_{1}, x_{2})) = \operatorname{add}(b)(f(x_{1}), f(x_{2})) \wedge \\ \quad \forall x_{1}, x_{2} \in \operatorname{carr}(a). \\ \quad \quad f(\operatorname{mul}(a)(x_{1}, x_{2})) = \operatorname{mul}(b)(f(x_{1}), f(x_{2})) \wedge \\ \quad \forall x_{1} \in \operatorname{carr}(a).\; (f(\operatorname{neg}(a)(x_{1})) = \operatorname{neg}(b)(f(x_{1}))) \wedge \\ \quad f(\operatorname{zero}(a)) = \operatorname{zero}(b) \wedge \\ \quad f(\operatorname{one}(a)) = \operatorname{one}(b) \wedge \\ \quad \forall x_{1} \in \operatorname{carr}(a).\; (\operatorname{fnrm}(a)(x_{1}) = \operatorname{fnrm}(b)(f(x_{1}))) \end{array}}}

\vnbbookfile{structure-library/normed-ag}{structure-library/normed-ag}
\noindent{\small normed-ag.scm -- NORMED-AG structure (normed abelian group)}\par
\vnbstructdef{normed-ag}{A normed ag $s$ is a 5-tuple whose components are listed below.\par
\begin{enumerate}[nosep]
\item $\operatorname{carr}(s) = \operatorname{nth}(1, s)$, a carrier (\texttt{carriers}): $\operatorname{carr}(s) \in \mathbf{Set}$.
\item $\operatorname{opr}(s) = \operatorname{nth}(2, s)$, an operation (\texttt{op}): $\operatorname{opr}(s) \in ((\operatorname{carr}(s) \times \operatorname{carr}(s)) \to \operatorname{carr}(s))$.
\item $\operatorname{iden}(s) = \operatorname{nth}(3, s)$, a constant (\texttt{constant}): $\operatorname{iden}(s) \in \operatorname{carr}(s)$.
\item $\operatorname{inv}(s) = \operatorname{nth}(4, s)$, an operation (\texttt{op}): $\operatorname{inv}(s) \in (\operatorname{carr}(s) \to \operatorname{carr}(s))$.
\item $\operatorname{nrm}(s) = \operatorname{nth}(5, s)$, an operation (\texttt{op}): $\operatorname{nrm}(s) \in (\operatorname{carr}(s) \to \mathbb{R})$.
\end{enumerate}
The laws (the \texttt{property} and \texttt{law} clauses):
\begin{itemize}[nosep]
\item $\operatorname{is\text{-}associative}(\operatorname{opr}(s), \operatorname{carr}(s))$
\item $\operatorname{is\text{-}identity}(\operatorname{opr}(s), \operatorname{iden}(s), \operatorname{carr}(s))$
\item $\operatorname{has\text{-}inverses}(\operatorname{opr}(s), \operatorname{iden}(s), \operatorname{inv}(s), \operatorname{carr}(s))$
\item $\operatorname{is\text{-}commutative}(\operatorname{opr}(s), \operatorname{carr}(s))$
\item $\operatorname{is\text{-}group\text{-}norm}(\operatorname{nrm}(s), \operatorname{opr}(s), \operatorname{inv}(s), \operatorname{iden}(s), \operatorname{carr}(s))$
\end{itemize}
The characteristic law (the axiom \texttt{is-normed-ag}): $\operatorname{is\text{-}normed\text{-}ag}(s)$ holds if and only if $\operatorname{length}(s) = 5$, every component has the typing listed with it, every law holds, and $s \in \mathbf{Set}$ (the structure is a set).
}
\par\medskip\noindent\textbf{Definition} (\texttt{hom-normed-ag}).\;$\operatorname{hom\text{-}normed\text{-}ag}(a, b) :=$
\begin{equation*}
\{\operatorname{homf\_} \in (\operatorname{carr}(a) \to \operatorname{carr}(b)) : \operatorname{is\text{-}hom\text{-}normed\text{-}ag}(a, b, \operatorname{homf\_})\}
\end{equation*}

\vnbdef{is-hom-normed-ag-def}{\vnbstmt{\begin{array}{@{}l@{}} \forall a, b, f. \\ \quad \operatorname{is\text{-}hom\text{-}normed\text{-}ag}(a, b, f) \\ \quad \Leftrightarrow\; \operatorname{is\text{-}normed\text{-}ag}(a) \wedge \\ \quad \operatorname{is\text{-}normed\text{-}ag}(b) \wedge \\ \quad f \in (\operatorname{carr}(a) \to \operatorname{carr}(b)) \wedge \\ \quad \forall x_{1}, x_{2} \in \operatorname{carr}(a). \\ \quad \quad f(\operatorname{opr}(a)(x_{1}, x_{2})) = \operatorname{opr}(b)(f(x_{1}), f(x_{2})) \wedge \\ \quad f(\operatorname{iden}(a)) = \operatorname{iden}(b) \wedge \\ \quad \forall x_{1} \in \operatorname{carr}(a).\; (f(\operatorname{inv}(a)(x_{1})) = \operatorname{inv}(b)(f(x_{1}))) \wedge \\ \quad \forall x_{1} \in \operatorname{carr}(a).\; (\operatorname{nrm}(a)(x_{1}) = \operatorname{nrm}(b)(f(x_{1}))) \end{array}}}
\vnbdef{normed-ag-class}{\vnbstmt{\forall s.\; ((s \in \operatorname{normed\text{-}ag}) \Leftrightarrow \operatorname{is\text{-}normed\text{-}ag}(s))}}

\vnbbookfile{structure-library/normed-field-metric}{structure-library/normed-field-metric}
\noindent{\small normed-field-metric.scm -- the metric space underlying a normed field.}\par
\par\medskip\noindent\textbf{Definition} (\texttt{nf-metric-space}).\;$\operatorname{nf\text{-}metric\text{-}space}(\mathit{nf}) :=$
\begin{equation*}
{[}\operatorname{carr}(\mathit{nf}), (\lambda {[}x, y] \in (\operatorname{carr}(\mathit{nf}) \times \operatorname{carr}(\mathit{nf})).\; \operatorname{fnrm}(\mathit{nf})(\operatorname{add}(\mathit{nf})(x, \operatorname{neg}(\mathit{nf})(y))))]
\end{equation*}

\vnbdef{nf-metric-distance}{\vnbstmt{\begin{array}{@{}l@{}} \forall f.; \forall x, y \in \operatorname{carr}(f). \\ \quad \operatorname{is\text{-}normed\text{-}field}(f) \Rightarrow  \\ \quad \quad \operatorname{dist}(\operatorname{nf\text{-}metric\text{-}space}(f))(x, y) = \operatorname{fnrm}(f)(\operatorname{add}(f)(x, \operatorname{neg}(f)(y))) \end{array}}}
\vnbdef{nf-metric-carrier}{\vnbstmt{\forall f.\; (\operatorname{is\text{-}normed\text{-}field}(f) \Rightarrow (\operatorname{pts}(\operatorname{nf\text{-}metric\text{-}space}(f)) \simeq \operatorname{carr}(f)))}}

\vnbbookfile{structure-library/normed-ag-metric}{structure-library/normed-ag-metric}
\noindent{\small normed-ag-metric.scm -- the metric space underlying a normed abelian group.}\par
\par\medskip\noindent\textbf{Definition} (\texttt{nag-metric-space}).\;$\operatorname{nag\text{-}metric\text{-}space}(\mathit{nag}) :=$
\begin{equation*}
{[}\operatorname{carr}(\mathit{nag}), (\lambda {[}u, v] \in (\operatorname{carr}(\mathit{nag}) \times \operatorname{carr}(\mathit{nag})).\; \operatorname{nrm}(\mathit{nag})(\operatorname{opr}(\mathit{nag})(u, \operatorname{inv}(\mathit{nag})(v))))]
\end{equation*}


\vnbbookfile{structure-library/normed-vector-space}{structure-library/normed-vector-space}
\noindent{\small normed-vector-space.scm -- NORMED-VECTOR-SPACE: a real normed vector space. Vocabulary only (no proofs).  Three decisions, as agreed:}\par
\vnbstructdef{normed-vector-space}{A normed vector space $s$ is a 7-tuple whose components are listed below.\par
\begin{enumerate}[nosep]
\item $\operatorname{scal}(s) = \operatorname{nth}(1, s)$, a substructure (\texttt{substructure}): $\operatorname{is\text{-}ring}(\operatorname{scal}(s))$.
\item $\operatorname{vec}(s) = \operatorname{nth}(2, s)$, a carrier (\texttt{carriers}): $\operatorname{vec}(s) \in \mathbf{Set}$.
\item $\operatorname{vadd}(s) = \operatorname{nth}(3, s)$, an operation (\texttt{op}): $\operatorname{vadd}(s) \in ((\operatorname{vec}(s) \times \operatorname{vec}(s)) \to \operatorname{vec}(s))$.
\item $\operatorname{vzero}(s) = \operatorname{nth}(4, s)$, a constant (\texttt{constant}): $\operatorname{vzero}(s) \in \operatorname{vec}(s)$.
\item $\operatorname{vneg}(s) = \operatorname{nth}(5, s)$, an operation (\texttt{op}): $\operatorname{vneg}(s) \in (\operatorname{vec}(s) \to \operatorname{vec}(s))$.
\item $\operatorname{act}(s) = \operatorname{nth}(6, s)$, an operation (\texttt{op}): $\operatorname{act}(s) \in ((\operatorname{carr}(\operatorname{scal}(s)) \times \operatorname{vec}(s)) \to \operatorname{vec}(s))$.
\item $\operatorname{vnrm}(s) = \operatorname{nth}(7, s)$, an operation (\texttt{op}): $\operatorname{vnrm}(s) \in (\operatorname{vec}(s) \to \mathbb{R})$.
\end{enumerate}
The laws (the \texttt{property} and \texttt{law} clauses):
\begin{itemize}[nosep]
\item $\operatorname{is\text{-}associative}(\operatorname{vadd}(s), \operatorname{vec}(s))$
\item $\operatorname{is\text{-}commutative}(\operatorname{vadd}(s), \operatorname{vec}(s))$
\item $\operatorname{is\text{-}identity}(\operatorname{vadd}(s), \operatorname{vzero}(s), \operatorname{vec}(s))$
\item $\operatorname{has\text{-}inverses}(\operatorname{vadd}(s), \operatorname{vzero}(s), \operatorname{vneg}(s), \operatorname{vec}(s))$
\item $\operatorname{scal}(s) = \operatorname{normed\text{-}field\text{-}as\text{-}commutative\text{-}ring}(\operatorname{rr\text{-}normed\text{-}field})$
\item
\begin{equation*}
\begin{array}{@{}l@{}} \forall \operatorname{r\_} \in \operatorname{carr}(\operatorname{scal}(s)), \operatorname{x\_}, \operatorname{y\_} \in \operatorname{vec}(s). \\ \quad \operatorname{act}(s)(\operatorname{r\_}, \operatorname{vadd}(s)(\operatorname{x\_}, \operatorname{y\_})) = \operatorname{vadd}(s)(\operatorname{act}(s)(\operatorname{r\_}, \operatorname{x\_}), \operatorname{act}(s)(\operatorname{r\_}, \operatorname{y\_})) \end{array}
\end{equation*}

\item
\begin{equation*}
\begin{array}{@{}l@{}} \forall \operatorname{r\_}, \operatorname{s\_} \in \operatorname{carr}(\operatorname{scal}(s)), \operatorname{x\_} \in \operatorname{vec}(s). \\ \quad \operatorname{act}(s)(\operatorname{add}(\operatorname{scal}(s))(\operatorname{r\_}, \operatorname{s\_}), \operatorname{x\_}) = \operatorname{vadd}(s)(\operatorname{act}(s)(\operatorname{r\_}, \operatorname{x\_}), \operatorname{act}(s)(\operatorname{s\_}, \operatorname{x\_})) \end{array}
\end{equation*}

\item
\begin{equation*}
\begin{array}{@{}l@{}} \forall \operatorname{r\_}, \operatorname{s\_} \in \operatorname{carr}(\operatorname{scal}(s)), \operatorname{x\_} \in \operatorname{vec}(s). \\ \quad \operatorname{act}(s)(\operatorname{mul}(\operatorname{scal}(s))(\operatorname{r\_}, \operatorname{s\_}), \operatorname{x\_}) = \operatorname{act}(s)(\operatorname{r\_}, \operatorname{act}(s)(\operatorname{s\_}, \operatorname{x\_})) \end{array}
\end{equation*}

\item $\forall \operatorname{x\_} \in \operatorname{vec}(s).\; (\operatorname{act}(s)(\operatorname{one}(\operatorname{scal}(s)), \operatorname{x\_}) = \operatorname{x\_})$
\item $\forall \operatorname{x\_} \in \operatorname{vec}(s).\; (0 \leq \operatorname{vnrm}(s)(\operatorname{x\_}))$
\item $\forall \operatorname{x\_} \in \operatorname{vec}(s).\; ((\operatorname{vnrm}(s)(\operatorname{x\_}) = 0) \Leftrightarrow (\operatorname{x\_} = \operatorname{vzero}(s)))$
\item
\begin{equation*}
\begin{array}{@{}l@{}} \forall \operatorname{r\_} \in \operatorname{carr}(\operatorname{scal}(s)), \operatorname{x\_} \in \operatorname{vec}(s). \\ \quad \operatorname{vnrm}(s)(\operatorname{act}(s)(\operatorname{r\_}, \operatorname{x\_})) = (\lvert \operatorname{r\_} \rvert \cdot \operatorname{vnrm}(s)(\operatorname{x\_})) \end{array}
\end{equation*}

\item
\begin{equation*}
\begin{array}{@{}l@{}} \forall \operatorname{x\_}, \operatorname{y\_} \in \operatorname{vec}(s). \\ \quad \operatorname{vnrm}(s)(\operatorname{vadd}(s)(\operatorname{x\_}, \operatorname{y\_})) \leq (\operatorname{vnrm}(s)(\operatorname{x\_}) + \operatorname{vnrm}(s)(\operatorname{y\_})) \end{array}
\end{equation*}

\end{itemize}
The characteristic law (the axiom \texttt{is-normed-vector-space}): $\operatorname{is\text{-}normed\text{-}vector\text{-}space}(s)$ holds if and only if $\operatorname{length}(s) = 7$, every component has the typing listed with it, every law holds, and $s \in \mathbf{Set}$ (the structure is a set).
}
\par\medskip\noindent\textbf{Definition} (\texttt{hom-normed-vector-space}).\;$\operatorname{hom\text{-}normed\text{-}vector\text{-}space}(a, b) :=$
\begin{equation*}
\{\operatorname{homf\_} \in (\operatorname{vec}(a) \to \operatorname{vec}(b)) : \operatorname{is\text{-}hom\text{-}normed\text{-}vector\text{-}space}(a, b, \operatorname{homf\_})\}
\end{equation*}

\vnbdef{is-hom-normed-vector-space-def}{\vnbstmt{\begin{array}{@{}l@{}} \forall a, b, f. \\ \quad \operatorname{is\text{-}hom\text{-}normed\text{-}vector\text{-}space}(a, b, f) \\ \quad \Leftrightarrow\; \operatorname{is\text{-}normed\text{-}vector\text{-}space}(a) \wedge \\ \quad \operatorname{is\text{-}normed\text{-}vector\text{-}space}(b) \wedge \\ \quad f \in (\operatorname{vec}(a) \to \operatorname{vec}(b)) \wedge \\ \quad \operatorname{scal}(a) = \operatorname{scal}(b) \wedge \\ \quad \forall x_{1}, x_{2} \in \operatorname{vec}(a). \\ \quad \quad f(\operatorname{vadd}(a)(x_{1}, x_{2})) = \operatorname{vadd}(b)(f(x_{1}), f(x_{2})) \wedge \\ \quad f(\operatorname{vzero}(a)) = \operatorname{vzero}(b) \wedge \\ \quad \forall x_{1} \in \operatorname{vec}(a).\; (f(\operatorname{vneg}(a)(x_{1})) = \operatorname{vneg}(b)(f(x_{1}))) \wedge \\ \quad \forall x_{1} \in \operatorname{carr}(\operatorname{scal}(a)), x_{2} \in \operatorname{vec}(a). \\ \quad \quad f(\operatorname{act}(a)(x_{1}, x_{2})) = \operatorname{act}(b)(x_{1}, f(x_{2})) \wedge \\ \quad \forall x_{1} \in \operatorname{vec}(a).\; (\operatorname{vnrm}(a)(x_{1}) = \operatorname{vnrm}(b)(f(x_{1}))) \end{array}}}
\vnbdef{normed-vector-space-class}{\vnbstmt{\forall s.\; ((s \in \operatorname{normed\text{-}vector\text{-}space}) \Leftrightarrow \operatorname{is\text{-}normed\text{-}vector\text{-}space}(s))}}

\vnbbookfile{structure-library/nvs-metric}{structure-library/nvs-metric}
\noindent{\small nvs-metric.scm -- the metric induced by the norm of a real normed vector space, as VOCABULARY.}\par
\par\medskip\noindent\textbf{Definition} (\texttt{nvs-metric-space}).\;$\operatorname{nvs\text{-}metric\text{-}space}(m) :=$
\begin{equation*}
{[}\operatorname{vec}(m), (\lambda {[}x, y] \in (\operatorname{vec}(m) \times \operatorname{vec}(m)).\; \operatorname{vnrm}(m)(\operatorname{vadd}(m)(x, \operatorname{vneg}(m)(y))))]
\end{equation*}


\vnbbookfile{structure-library/linear-functional}{structure-library/linear-functional}
\noindent{\small linear-functional.scm -- bounded linear functionals and the dual norm on a real NORMED-VECTOR-SPACE.  Vocabulary only (no proofs).}\par
\par\medskip\noindent\textbf{Definition} (\texttt{is-linear-functional}).\;For all $m$ and $f$, 
\begin{equation*}
\begin{array}{@{}l@{}} \operatorname{is\text{-}linear\text{-}functional}(m, f) \\ \Leftrightarrow\; f \in (\operatorname{vec}(m) \to \mathbb{R}) \wedge \\ \forall \operatorname{x\_}, \operatorname{y\_} \in \operatorname{vec}(m).\; (f(\operatorname{vadd}(m)(\operatorname{x\_}, \operatorname{y\_})) = +(f(\operatorname{x\_}), f(\operatorname{y\_}))) \wedge \\ \forall \operatorname{r\_} \in \mathbb{R}, \operatorname{x\_} \in \operatorname{vec}(m).\; (f(\operatorname{act}(m)(\operatorname{r\_}, \operatorname{x\_})) = \cdot(\operatorname{r\_}, f(\operatorname{x\_}))) \end{array}
\end{equation*}

\par\medskip\noindent\textbf{Definition} (\texttt{is-bounded-linear-functional}).\;For all $m$ and $f$, 
\begin{equation*}
\begin{array}{@{}l@{}} \operatorname{is\text{-}bounded\text{-}linear\text{-}functional}(m, f) \\ \Leftrightarrow\; \operatorname{is\text{-}linear\text{-}functional}(m, f) \wedge \\ \exists \operatorname{c\_} \in \mathbb{R}. \\ \quad (0 \leq \operatorname{c\_}) \wedge \forall \operatorname{x\_} \in \operatorname{vec}(m).\; (\lvert f(\operatorname{x\_}) \rvert \leq \cdot(\operatorname{c\_}, \operatorname{vnrm}(m)(\operatorname{x\_}))) \end{array}
\end{equation*}

\par\medskip\noindent\textbf{Definition} (\texttt{dual-norm}).\;$\operatorname{dual\text{-}norm}(m, f) :=$
\begin{equation*}
(\iota\, \operatorname{c\_}.\;((\operatorname{c\_} \in \mathbb{R}) \wedge ((0 \leq \operatorname{c\_}) \wedge (\forall \operatorname{x\_} \in \operatorname{vec}(m).\; (\lvert f(\operatorname{x\_}) \rvert \leq (\operatorname{c\_} \cdot \operatorname{vnrm}(m)(\operatorname{x\_}))) \wedge \forall \operatorname{d\_}.\; (((\operatorname{d\_} \in \mathbb{R}) \wedge ((0 \leq \operatorname{d\_}) \wedge \forall \operatorname{x\_} \in \operatorname{vec}(m).\; (\lvert f(\operatorname{x\_}) \rvert \leq (\operatorname{d\_} \cdot \operatorname{vnrm}(m)(\operatorname{x\_}))))) \Rightarrow (\operatorname{c\_} \leq \operatorname{d\_}))))))
\end{equation*}

\par\medskip\noindent\textbf{Definition} (\texttt{is-linear-functional-on}).\;For all $m$, $s$ and $f$, 
\begin{equation*}
\begin{array}{@{}l@{}} \operatorname{is\text{-}linear\text{-}functional\text{-}on}(m, s, f) \\ \Leftrightarrow\; f \in (s \to \mathbb{R}) \wedge \\ \forall \operatorname{x\_}, \operatorname{y\_} \in s.\; (f(\operatorname{vadd}(m)(\operatorname{x\_}, \operatorname{y\_})) = +(f(\operatorname{x\_}), f(\operatorname{y\_}))) \wedge \\ \forall \operatorname{r\_} \in \mathbb{R}, \operatorname{x\_} \in s.\; (f(\operatorname{act}(m)(\operatorname{r\_}, \operatorname{x\_})) = \cdot(\operatorname{r\_}, f(\operatorname{x\_}))) \end{array}
\end{equation*}

\par\medskip\noindent\textbf{Definition} (\texttt{is-bounded-linear-functional-on}).\;For all $m$, $s$ and $f$, 
\begin{equation*}
\begin{array}{@{}l@{}} \operatorname{is\text{-}bounded\text{-}linear\text{-}functional\text{-}on}(m, s, f) \\ \Leftrightarrow\; \operatorname{is\text{-}linear\text{-}functional\text{-}on}(m, s, f) \wedge \\ \exists \operatorname{c\_} \in \mathbb{R}. \\ \quad (0 \leq \operatorname{c\_}) \wedge \forall \operatorname{x\_} \in s.\; (\lvert f(\operatorname{x\_}) \rvert \leq \cdot(\operatorname{c\_}, \operatorname{vnrm}(m)(\operatorname{x\_}))) \end{array}
\end{equation*}

\par\medskip\noindent\textbf{Definition} (\texttt{dual-norm-on}).\;$\operatorname{dual\text{-}norm\text{-}on}(m, s, f) :=$
\begin{equation*}
(\iota\, \operatorname{c\_}.\;((\operatorname{c\_} \in \mathbb{R}) \wedge ((0 \leq \operatorname{c\_}) \wedge (\forall \operatorname{x\_} \in s.\; (\lvert f(\operatorname{x\_}) \rvert \leq (\operatorname{c\_} \cdot \operatorname{vnrm}(m)(\operatorname{x\_}))) \wedge \forall \operatorname{d\_}.\; (((\operatorname{d\_} \in \mathbb{R}) \wedge ((0 \leq \operatorname{d\_}) \wedge \forall \operatorname{x\_} \in s.\; (\lvert f(\operatorname{x\_}) \rvert \leq (\operatorname{d\_} \cdot \operatorname{vnrm}(m)(\operatorname{x\_}))))) \Rightarrow (\operatorname{c\_} \leq \operatorname{d\_}))))))
\end{equation*}

\par\medskip\noindent\textbf{Definition} (\texttt{extends-on}).\;For all $s$, $g$ and $f$,  $\operatorname{extends\text{-}on}(s, g, f) \Leftrightarrow \forall \operatorname{x\_} \in s.\; (g(\operatorname{x\_}) = f(\operatorname{x\_}))$
\par\medskip\noindent\textbf{Definition} (\texttt{npe}).\;For all $m$, $s$, $f$, $t$ and $g$, 
\begin{equation*}
\begin{array}{@{}l@{}} \operatorname{npe}(m, s, f, t, g) \\ \Leftrightarrow\; \operatorname{is\text{-}submodule}(m, t) \wedge \\ s \subseteq t \wedge \\ \operatorname{is\text{-}linear\text{-}functional\text{-}on}(m, t, g) \wedge \\ \operatorname{extends\text{-}on}(s, g, f) \wedge \\ \forall \operatorname{w\_} \in t.\; (\lvert g(\operatorname{w\_}) \rvert \leq \cdot(\operatorname{dual\text{-}norm\text{-}on}(m, s, f), \operatorname{vnrm}(m)(\operatorname{w\_}))) \end{array}
\end{equation*}

\par\medskip\noindent\textbf{Definition} (\texttt{good-sub}).\;For all $m$, $s$, $f$ and $t$,  $\operatorname{good\text{-}sub}(m, s, f, t) \Leftrightarrow \exists \operatorname{g\_}.\; \operatorname{npe}(m, s, f, t, \operatorname{g\_})$
\par\medskip\noindent\textbf{Definition} (\texttt{line}).\;$\operatorname{line}(m, v) :=$ $\{\operatorname{y\_} \in \operatorname{vec}(m) : \exists \operatorname{r\_} \in \mathbb{R}.\; (\operatorname{y\_} = \operatorname{act}(m)(\operatorname{r\_}, v))\}$

\vnbbookfile{structure-library/dual-space}{structure-library/dual-space}
\noindent{\small dual-space.scm -- THE DUAL of a real normed vector space, as an OBJECT.}\par
\par\medskip\noindent\textbf{Definition} (\texttt{dual}).\;${m}^{*} :=$
\begin{equation*}
{[}\operatorname{normed\text{-}field\text{-}as\text{-}commutative\text{-}ring}(\operatorname{rr\text{-}normed\text{-}field}), \operatorname{dual\text{-}vec}(m), (\lambda {[}\operatorname{dlaf\_}, \operatorname{dlag\_}] \in (\operatorname{dual\text{-}vec}(m) \times \operatorname{dual\text{-}vec}(m)).\; (\lambda \operatorname{dlax\_} \in \operatorname{vec}(m).\; (\operatorname{dlaf\_}(\operatorname{dlax\_}) + \operatorname{dlag\_}(\operatorname{dlax\_})))), (\lambda \operatorname{dlzx\_} \in \operatorname{vec}(m).\; 0), (\lambda \operatorname{dlnf\_} \in \operatorname{dual\text{-}vec}(m).\; (\lambda \operatorname{dlnx\_} \in \operatorname{vec}(m).\; -\operatorname{dlnf\_}(\operatorname{dlnx\_}))), (\lambda {[}\operatorname{dlsr\_}, \operatorname{dlsf\_}] \in (\mathbb{R} \times \operatorname{dual\text{-}vec}(m)).\; (\lambda \operatorname{dlsx\_} \in \operatorname{vec}(m).\; (\operatorname{dlsr\_} \cdot \operatorname{dlsf\_}(\operatorname{dlsx\_})))), (\lambda \operatorname{dlvf\_} \in \operatorname{dual\text{-}vec}(m).\; \operatorname{dual\text{-}norm}(m, \operatorname{dlvf\_}))]
\end{equation*}

\par\medskip\noindent\textbf{Definition} (\texttt{dual-vec}).\;$\operatorname{dual\text{-}vec}(m) :=$ $\{\operatorname{f\_} \in (\operatorname{vec}(m) \to \mathbb{R}) : \operatorname{is\text{-}bounded\text{-}linear\text{-}functional}(m, \operatorname{f\_})\}$

\vnbbookfile{structure-library/metric-subspace}{structure-library/metric-subspace}
\noindent{\small metric-subspace.scm -- the METRIC SUBSPACE of a metric space, and the RESTRICTION of a function to a subset.  DEFINITIONS ONLY; every law is PROVEN in theorem-library/metric-subspace-laws.scm.}\par
\par\medskip\noindent\textbf{Definition} (\texttt{subspace-ms}).\;$\operatorname{subspace\text{-}ms}(s, a) :=$
\begin{equation*}
{[}a, (\lambda {[}\operatorname{sbu\_}, \operatorname{sbv\_}] \in (a \times a).\; \operatorname{dist}(s)(\operatorname{sbu\_}, \operatorname{sbv\_}))]
\end{equation*}

\par\medskip\noindent\textbf{Definition} (\texttt{restrict}).\;$\operatorname{restrict}(f, a) :=$ $(\lambda \operatorname{rsx\_} \in a.\; f(\operatorname{rsx\_}))$

\vnbbookfile{structure-library/product2-vocab}{structure-library/product2-vocab}
\noindent{\small product2-vocab.scm -- THE BINARY PRODUCT OF TWO METRIC SPACES (definitions only).}\par
\par\medskip\noindent\textbf{Definition} (\texttt{product2-ms}).\;$\operatorname{product2\text{-}ms}(s, t) :=$
\begin{equation*}
{[}(\operatorname{pts}(s) \times \operatorname{pts}(t)), (\lambda {[}\operatorname{p2u\_}, \operatorname{p2v\_}] \in ((\operatorname{pts}(s) \times \operatorname{pts}(t)) \times (\operatorname{pts}(s) \times \operatorname{pts}(t))).\; \operatorname{max}(\operatorname{dist}(s)(\operatorname{nth}(1, \operatorname{p2u\_}), \operatorname{nth}(1, \operatorname{p2v\_})), \operatorname{dist}(t)(\operatorname{nth}(2, \operatorname{p2u\_}), \operatorname{nth}(2, \operatorname{p2v\_}))))]
\end{equation*}


\vnbbookfile{structure-library/diff-on}{structure-library/diff-on}
\noindent{\small diff-on.scm -- THE DERIVATIVE ON AN OPEN SET, OVER A NORMED FIELD. DEFINITIONS ONLY; every law is PROVEN in theorem-library/diff-on-laws.scm.}\par
\par\medskip\noindent\textbf{Definition} (\texttt{is-diff-on}).\;For all $k$, $u$, $f$, $a$ and $l$, 
\begin{equation*}
\begin{array}{@{}l@{}} \operatorname{is\text{-}diff\text{-}on}(k, u, f, a, l) \\ \Leftrightarrow\; \operatorname{is\text{-}normed\text{-}field}(k) \wedge \\ \operatorname{is\text{-}open}(\operatorname{nf\text{-}metric\text{-}space}(k), u) \wedge \\ f \in (u \to \operatorname{carr}(k)) \wedge \\ a \in u \wedge \\ l \in \operatorname{carr}(k) \wedge \\ \exists \varphi \in (u \to \operatorname{carr}(k)). \\ \quad \operatorname{is\text{-}continuous\text{-}at}(\operatorname{subspace\text{-}ms}(\operatorname{nf\text{-}metric\text{-}space}(k), u), \operatorname{nf\text{-}metric\text{-}space}(k), \varphi, a) \wedge \\ \quad \varphi(a) = l \wedge \\ \quad \forall x \in u. \\ \quad \quad \operatorname{add}(k)(f(x), \operatorname{neg}(k)(f(a))) = \operatorname{mul}(k)(\varphi(x), \operatorname{add}(k)(x, \operatorname{neg}(k)(a))) \end{array}
\end{equation*}

\par\medskip\noindent\textbf{Definition} (\texttt{deriv-on}).\;$\operatorname{deriv\text{-}on}(k, u, f, a) :=$ $(\iota\, l.\;\operatorname{is\text{-}diff\text{-}on}(k, u, f, a, l))$
\par\medskip\noindent\textbf{Definition} (\texttt{holomorphic-on}).\;For all $u$ and $f$, 
\begin{equation*}
\begin{array}{@{}l@{}} \operatorname{holomorphic\text{-}on}(u, f) \\ \Leftrightarrow\; \operatorname{is\text{-}open}(\operatorname{nf\text{-}metric\text{-}space}(\operatorname{cc\text{-}normed\text{-}field}), u) \wedge \\ f \in (u \to \mathbb{C}) \wedge \\ \forall a \in u.\; \exists l.\; \operatorname{is\text{-}diff\text{-}on}(\operatorname{cc\text{-}normed\text{-}field}, u, f, a, l) \end{array}
\end{equation*}


\vnbbookfile{structure-library/extreme-value}{structure-library/extreme-value}
\noindent{\small MOVED 2026-09-20 (batch 12-A) from theorem-library/, and MOVED UP in load.scm to the calculus-vocabulary block just after numeric-instances: this file is VOCABULARY -- one def-functoid, not one proof -- while 43 files are STATED with CCINT and every one of them had to load below it.  CCINT is a SEP over RR and has no ...}\par
\par\medskip\noindent\textbf{Definition} (\texttt{ccint}).\;$\operatorname{ccint}(a, b) :=$ $\{x \in \mathbb{R} : ((a \leq x) \wedge (x \leq b))\}$

\vnbbookfile{structure-library/derivative}{structure-library/derivative}
\noindent{\small derivative.scm -- the Caratheodory derivative, as VOCABULARY.}\par
\par\medskip\noindent\textbf{Definition} (\texttt{is-diff-at}).\;For all $f$, $a$ and $l$, 
\begin{equation*}
\begin{array}{@{}l@{}} \operatorname{is\text{-}diff\text{-}at}(f, a, l) \\ \Leftrightarrow\; f \in (\mathbb{R} \to \mathbb{R}) \wedge \\ a \in \mathbb{R} \wedge \\ l \in \mathbb{R} \wedge \\ \exists \varphi \in (\mathbb{R} \to \mathbb{R}). \\ \quad \operatorname{is\text{-}continuous\text{-}at}(\operatorname{rr\text{-}ms}, \operatorname{rr\text{-}ms}, \varphi, a) \wedge \\ \quad \varphi(a) = l \wedge \\ \quad \forall x \in \mathbb{R}.\; (-(f(x), f(a)) = \cdot(\varphi(x), -(x, a))) \end{array}
\end{equation*}

\par\medskip\noindent\textbf{Definition} (\texttt{deriv}).\;$\operatorname{deriv}(f, a) :=$ $(\iota\, l.\;\operatorname{is\text{-}diff\text{-}at}(f, a, l))$

\vnbbookfile{structure-library/cc-coords}{structure-library/cc-coords}
\noindent{\small cc-coords.scm -- CC and cartesian(RR, RR), as VOCABULARY.}\par
\par\medskip\noindent\textbf{Definition} (\texttt{cc-coords}).\;
\begin{equation*}
(\operatorname{cc\text{-}coords} \simeq (\lambda \operatorname{zv\_} \in \mathbb{C}.\; {[}\operatorname{real\text{-}part}(\operatorname{zv\_}), \operatorname{imag\text{-}part}(\operatorname{zv\_})]))
\end{equation*}

\par\medskip\noindent\textbf{Definition} (\texttt{cc-of-pair}).\;
\begin{equation*}
(\operatorname{cc\text{-}of\text{-}pair} \simeq (\lambda \operatorname{pv\_} \in (\mathbb{R} \times \mathbb{R}).\; +(\operatorname{nth}(1, \operatorname{pv\_}), \cdot(\operatorname{nth}(2, \operatorname{pv\_}), +i))))
\end{equation*}

\par\medskip\noindent\textbf{Definition} (\texttt{is-cc-diff-at}).\;For all $g$, $a$ and $l$, 
\begin{equation*}
\begin{array}{@{}l@{}} \operatorname{is\text{-}cc\text{-}diff\text{-}at}(g, a, l) \\ \Leftrightarrow\; g \in (\mathbb{R} \to \mathbb{C}) \wedge \\ a \in \mathbb{R} \wedge \\ l \in \mathbb{C} \wedge \\ \operatorname{is\text{-}diff\text{-}at}((\lambda \operatorname{sv\_} \in \mathbb{R}.\; \operatorname{real\text{-}part}(g(\operatorname{sv\_}))), a, \operatorname{real\text{-}part}(l)) \wedge \\ \operatorname{is\text{-}diff\text{-}at}((\lambda \operatorname{sv\_} \in \mathbb{R}.\; \operatorname{imag\text{-}part}(g(\operatorname{sv\_}))), a, \operatorname{imag\text{-}part}(l)) \end{array}
\end{equation*}


\vnbbookfile{structure-library/antiderivative}{structure-library/antiderivative}
\noindent{\small antiderivative.scm -- Definition 4.6 (calculus.pdf), as VOCABULARY.}\par
\par\medskip\noindent\textbf{Definition} (\texttt{is-antiderivative}).\;For all $f$, $\varphi$, $a$ and $b$, 
\begin{equation*}
\begin{array}{@{}l@{}} \operatorname{is\text{-}antiderivative}(f, \varphi, a, b) \\ \Leftrightarrow\; f \in (\mathbb{R} \to \mathbb{R}) \wedge \\ \varphi \in (\mathbb{R} \to \mathbb{R}) \wedge \\ a \in \mathbb{R} \wedge \\ b \in \mathbb{R} \wedge \\ a < b \wedge \\ \forall \operatorname{x\_} \in \operatorname{ccint}(a, b).\; \operatorname{is\text{-}continuous\text{-}at}(\operatorname{rr\text{-}ms}, \operatorname{rr\text{-}ms}, f, \operatorname{x\_}) \wedge \\ \forall \operatorname{th\_}. \\ \quad ((\operatorname{th\_} \in \mathbb{R}) \wedge ((a < \operatorname{th\_}) \wedge (\operatorname{th\_} < b))) \Rightarrow \operatorname{is\text{-}diff\text{-}at}(f, \operatorname{th\_}, \varphi(\operatorname{th\_})) \end{array}
\end{equation*}

\par\medskip\noindent\textbf{Definition} (\texttt{is-antiderivable}).\;For all $\varphi$, $a$ and $b$, 
\begin{equation*}
\operatorname{is\text{-}antiderivable}(\varphi, a, b) \Leftrightarrow \exists \operatorname{f\_}.\; \operatorname{is\text{-}antiderivative}(\operatorname{f\_}, \varphi, a, b)
\end{equation*}


\vnbbookfile{structure-library/complex-inner-product}{structure-library/complex-inner-product}
\noindent{\small complex-inner-product.scm -- COMPLEX-INNER-PRODUCT-SPACE: a vector space over CC carrying a sesquilinear, conjugate-symmetric, positive-definite form IP. Vocabulary only (no proofs); the two inequalities it exists for are proved in theorem-library/inner-product-inequalities.scm.}\par
\vnbstructdef{complex-inner-product-space}{A complex inner product space $s$ is a 7-tuple whose components are listed below.\par
\begin{enumerate}[nosep]
\item $\operatorname{scal}(s) = \operatorname{nth}(1, s)$, a substructure (\texttt{substructure}): $\operatorname{is\text{-}ring}(\operatorname{scal}(s))$.
\item $\operatorname{vec}(s) = \operatorname{nth}(2, s)$, a carrier (\texttt{carriers}): $\operatorname{vec}(s) \in \mathbf{Set}$.
\item $\operatorname{vadd}(s) = \operatorname{nth}(3, s)$, an operation (\texttt{op}): $\operatorname{vadd}(s) \in ((\operatorname{vec}(s) \times \operatorname{vec}(s)) \to \operatorname{vec}(s))$.
\item $\operatorname{vzero}(s) = \operatorname{nth}(4, s)$, a constant (\texttt{constant}): $\operatorname{vzero}(s) \in \operatorname{vec}(s)$.
\item $\operatorname{vneg}(s) = \operatorname{nth}(5, s)$, an operation (\texttt{op}): $\operatorname{vneg}(s) \in (\operatorname{vec}(s) \to \operatorname{vec}(s))$.
\item $\operatorname{act}(s) = \operatorname{nth}(6, s)$, an operation (\texttt{op}): $\operatorname{act}(s) \in ((\operatorname{carr}(\operatorname{scal}(s)) \times \operatorname{vec}(s)) \to \operatorname{vec}(s))$.
\item $\langle\cdot,\cdot\rangle_{s} = \operatorname{nth}(7, s)$, an operation (\texttt{op}): $\langle\cdot,\cdot\rangle_{s} \in ((\operatorname{vec}(s) \times \operatorname{vec}(s)) \to \mathbb{C})$.
\end{enumerate}
The laws (the \texttt{property} and \texttt{law} clauses):
\begin{itemize}[nosep]
\item $\operatorname{is\text{-}associative}(\operatorname{vadd}(s), \operatorname{vec}(s))$
\item $\operatorname{is\text{-}commutative}(\operatorname{vadd}(s), \operatorname{vec}(s))$
\item $\operatorname{is\text{-}identity}(\operatorname{vadd}(s), \operatorname{vzero}(s), \operatorname{vec}(s))$
\item $\operatorname{has\text{-}inverses}(\operatorname{vadd}(s), \operatorname{vzero}(s), \operatorname{vneg}(s), \operatorname{vec}(s))$
\item $\operatorname{scal}(s) = \operatorname{normed\text{-}field\text{-}as\text{-}commutative\text{-}ring}(\operatorname{cc\text{-}normed\text{-}field})$
\item $\operatorname{carr}(\operatorname{scal}(s)) = \mathbb{C}$
\item
\begin{equation*}
\begin{array}{@{}l@{}} \forall \operatorname{r\_} \in \operatorname{carr}(\operatorname{scal}(s)), \operatorname{x\_}, \operatorname{y\_} \in \operatorname{vec}(s). \\ \quad \operatorname{act}(s)(\operatorname{r\_}, \operatorname{vadd}(s)(\operatorname{x\_}, \operatorname{y\_})) = \operatorname{vadd}(s)(\operatorname{act}(s)(\operatorname{r\_}, \operatorname{x\_}), \operatorname{act}(s)(\operatorname{r\_}, \operatorname{y\_})) \end{array}
\end{equation*}

\item
\begin{equation*}
\begin{array}{@{}l@{}} \forall \operatorname{r\_}, \operatorname{s\_} \in \operatorname{carr}(\operatorname{scal}(s)), \operatorname{x\_} \in \operatorname{vec}(s). \\ \quad \operatorname{act}(s)(\operatorname{add}(\operatorname{scal}(s))(\operatorname{r\_}, \operatorname{s\_}), \operatorname{x\_}) = \operatorname{vadd}(s)(\operatorname{act}(s)(\operatorname{r\_}, \operatorname{x\_}), \operatorname{act}(s)(\operatorname{s\_}, \operatorname{x\_})) \end{array}
\end{equation*}

\item
\begin{equation*}
\begin{array}{@{}l@{}} \forall \operatorname{r\_}, \operatorname{s\_} \in \operatorname{carr}(\operatorname{scal}(s)), \operatorname{x\_} \in \operatorname{vec}(s). \\ \quad \operatorname{act}(s)(\operatorname{mul}(\operatorname{scal}(s))(\operatorname{r\_}, \operatorname{s\_}), \operatorname{x\_}) = \operatorname{act}(s)(\operatorname{r\_}, \operatorname{act}(s)(\operatorname{s\_}, \operatorname{x\_})) \end{array}
\end{equation*}

\item $\forall \operatorname{x\_} \in \operatorname{vec}(s).\; (\operatorname{act}(s)(\operatorname{one}(\operatorname{scal}(s)), \operatorname{x\_}) = \operatorname{x\_})$
\item
\begin{equation*}
\begin{array}{@{}l@{}} \forall \operatorname{x\_}, \operatorname{y\_}, \operatorname{z\_} \in \operatorname{vec}(s). \\ \quad \langle\cdot,\cdot\rangle_{s}(\operatorname{vadd}(s)(\operatorname{x\_}, \operatorname{y\_}), \operatorname{z\_}) = (\langle\cdot,\cdot\rangle_{s}(\operatorname{x\_}, \operatorname{z\_}) + \langle\cdot,\cdot\rangle_{s}(\operatorname{y\_}, \operatorname{z\_})) \end{array}
\end{equation*}

\item
\begin{equation*}
\begin{array}{@{}l@{}} \forall \operatorname{a\_} \in \mathbb{C}, \operatorname{x\_}, \operatorname{y\_} \in \operatorname{vec}(s). \\ \quad \langle\cdot,\cdot\rangle_{s}(\operatorname{act}(s)(\operatorname{a\_}, \operatorname{x\_}), \operatorname{y\_}) = (\operatorname{a\_} \cdot \langle\cdot,\cdot\rangle_{s}(\operatorname{x\_}, \operatorname{y\_})) \end{array}
\end{equation*}

\item $\forall \operatorname{x\_}, \operatorname{y\_} \in \operatorname{vec}(s).\; (\langle\cdot,\cdot\rangle_{s}(\operatorname{y\_}, \operatorname{x\_}) = \operatorname{conjugate}(\langle\cdot,\cdot\rangle_{s}(\operatorname{x\_}, \operatorname{y\_})))$
\item $\forall \operatorname{x\_} \in \operatorname{vec}(s).\; (\langle\cdot,\cdot\rangle_{s}(\operatorname{x\_}, \operatorname{x\_}) \in \mathbb{R})$
\item $\forall \operatorname{x\_} \in \operatorname{vec}(s).\; (0 \leq \langle\cdot,\cdot\rangle_{s}(\operatorname{x\_}, \operatorname{x\_}))$
\item $\forall \operatorname{x\_} \in \operatorname{vec}(s).\; ((\langle\cdot,\cdot\rangle_{s}(\operatorname{x\_}, \operatorname{x\_}) = 0) \Leftrightarrow (\operatorname{x\_} = \operatorname{vzero}(s)))$
\end{itemize}
The characteristic law (the axiom \texttt{is-complex-inner-product-space}): $\operatorname{is\text{-}complex\text{-}inner\text{-}product\text{-}space}(s)$ holds if and only if $\operatorname{length}(s) = 7$, every component has the typing listed with it, every law holds, and $s \in \mathbf{Set}$ (the structure is a set).
}
\par\medskip\noindent\textbf{Definition} (\texttt{ip-normed-ag}).\;$\operatorname{ip\text{-}normed\text{-}ag}(v) :=$
\begin{equation*}
{[}\operatorname{vec}(v), \operatorname{vadd}(v), \operatorname{vzero}(v), \operatorname{vneg}(v), (\lambda \operatorname{x\_} \in \operatorname{vec}(v).\; \lVert \operatorname{x\_} \rVert)]
\end{equation*}

\par\medskip\noindent\textbf{Definition} (\texttt{ip-norm}).\;$\lVert x \rVert :=$ $\sqrt{\langle\cdot,\cdot\rangle_{v}(x, x)}$
\par\medskip\noindent\textbf{Definition} (\texttt{hom-complex-inner-product-space}).\;$\operatorname{hom\text{-}complex\text{-}inner\text{-}product\text{-}space}(a, b) :=$
\begin{equation*}
\{\operatorname{homf\_} \in (\operatorname{vec}(a) \to \operatorname{vec}(b)) : \operatorname{is\text{-}hom\text{-}complex\text{-}inner\text{-}product\text{-}space}(a, b, \operatorname{homf\_})\}
\end{equation*}

\vnbdef{cips-act-type}{\vnbstmt{\begin{array}{@{}l@{}} \forall v.; \forall a \in \mathbb{C}, x \in \operatorname{vec}(v). \\ \quad \operatorname{is\text{-}complex\text{-}inner\text{-}product\text{-}space}(v) \Rightarrow (\operatorname{act}(v)(a, x) \in \operatorname{vec}(v)) \end{array}}}
\vnbdef{cips-scal-carrier}{\vnbstmt{\forall v.\; (\operatorname{is\text{-}complex\text{-}inner\text{-}product\text{-}space}(v) \Rightarrow (\operatorname{carr}(\operatorname{scal}(v)) = \mathbb{C}))}}
\vnbdef{cips-vzero-in}{\vnbstmt{\forall v.\; (\operatorname{is\text{-}complex\text{-}inner\text{-}product\text{-}space}(v) \Rightarrow (\operatorname{vzero}(v) \in \operatorname{vec}(v)))}}
\vnbdef{cips-vadd-vzero}{\vnbstmt{\begin{array}{@{}l@{}} \forall v.; \forall x \in \operatorname{vec}(v). \\ \quad \operatorname{is\text{-}complex\text{-}inner\text{-}product\text{-}space}(v) \Rightarrow (\operatorname{vadd}(v)(\operatorname{vzero}(v), x) = x) \end{array}}}
\vnbdef{cips-vadd-type}{\vnbstmt{\begin{array}{@{}l@{}} \forall v.; \forall x, y \in \operatorname{vec}(v). \\ \quad \operatorname{is\text{-}complex\text{-}inner\text{-}product\text{-}space}(v) \Rightarrow (\operatorname{vadd}(v)(x, y) \in \operatorname{vec}(v)) \end{array}}}
\vnbdef{cips-ip-type}{\vnbstmt{\begin{array}{@{}l@{}} \forall v.; \forall x, y \in \operatorname{vec}(v). \\ \quad \operatorname{is\text{-}complex\text{-}inner\text{-}product\text{-}space}(v) \Rightarrow (\langle\cdot,\cdot\rangle_{v}(x, y) \in \mathbb{C}) \end{array}}}
\vnbdef{cips-ip-add-left}{\vnbstmt{\begin{array}{@{}l@{}} \forall v.; \forall x, y, z \in \operatorname{vec}(v). \\ \quad \operatorname{is\text{-}complex\text{-}inner\text{-}product\text{-}space}(v) \Rightarrow  \\ \quad \quad \langle\cdot,\cdot\rangle_{v}(\operatorname{vadd}(v)(x, y), z) = (\langle\cdot,\cdot\rangle_{v}(x, z) + \langle\cdot,\cdot\rangle_{v}(y, z)) \end{array}}}
\vnbdef{cips-ip-homog-left}{\vnbstmt{\begin{array}{@{}l@{}} \forall v.; \forall a \in \mathbb{C}, x, y \in \operatorname{vec}(v). \\ \quad \operatorname{is\text{-}complex\text{-}inner\text{-}product\text{-}space}(v) \Rightarrow  \\ \quad \quad \langle\cdot,\cdot\rangle_{v}(\operatorname{act}(v)(a, x), y) = (a \cdot \langle\cdot,\cdot\rangle_{v}(x, y)) \end{array}}}
\vnbdef{cips-ip-conj-sym}{\vnbstmt{\begin{array}{@{}l@{}} \forall v.; \forall x, y \in \operatorname{vec}(v). \\ \quad \operatorname{is\text{-}complex\text{-}inner\text{-}product\text{-}space}(v) \Rightarrow  \\ \quad \quad \langle\cdot,\cdot\rangle_{v}(y, x) = \operatorname{conjugate}(\langle\cdot,\cdot\rangle_{v}(x, y)) \end{array}}}
\vnbdef{cips-ip-self-real}{\vnbstmt{\begin{array}{@{}l@{}} \forall v.; \forall x \in \operatorname{vec}(v). \\ \quad \operatorname{is\text{-}complex\text{-}inner\text{-}product\text{-}space}(v) \Rightarrow (\langle\cdot,\cdot\rangle_{v}(x, x) \in \mathbb{R}) \end{array}}}
\vnbdef{cips-ip-self-nonneg}{\vnbstmt{\begin{array}{@{}l@{}} \forall v.; \forall x \in \operatorname{vec}(v). \\ \quad \operatorname{is\text{-}complex\text{-}inner\text{-}product\text{-}space}(v) \Rightarrow (0 \leq \langle\cdot,\cdot\rangle_{v}(x, x)) \end{array}}}
\vnbdef{cips-ip-zero-vector}{\vnbstmt{\begin{array}{@{}l@{}} \forall v.; \forall x \in \operatorname{vec}(v). \\ \quad (\operatorname{is\text{-}complex\text{-}inner\text{-}product\text{-}space}(v) \wedge (\langle\cdot,\cdot\rangle_{v}(x, x) = 0)) \Rightarrow (x = \operatorname{vzero}(v)) \end{array}}}
\vnbdef{cips-ip-definite}{\vnbstmt{\begin{array}{@{}l@{}} \forall v.; \forall x \in \operatorname{vec}(v). \\ \quad \operatorname{is\text{-}complex\text{-}inner\text{-}product\text{-}space}(v) \Rightarrow ((\langle\cdot,\cdot\rangle_{v}(x, x) = 0) \Leftrightarrow (x = \operatorname{vzero}(v))) \end{array}}}
\vnbdef{ip-normed-ag-carrier}{\vnbstmt{\begin{array}{@{}l@{}} \forall v. \\ \quad \operatorname{is\text{-}complex\text{-}inner\text{-}product\text{-}space}(v) \Rightarrow (\operatorname{carr}(\operatorname{ip\text{-}normed\text{-}ag}(v)) \simeq \operatorname{vec}(v)) \end{array}}}
\vnbdef{ip-normed-ag-norm-value}{\vnbstmt{\begin{array}{@{}l@{}} \forall v.; \forall x \in \operatorname{vec}(v). \\ \quad \operatorname{is\text{-}complex\text{-}inner\text{-}product\text{-}space}(v) \Rightarrow  \\ \quad \quad (\operatorname{nrm}(\operatorname{ip\text{-}normed\text{-}ag}(v))(x) \simeq \lVert x \rVert) \end{array}}}
\vnbdef{is-hom-complex-inner-product-space-def}{\vnbstmt{\begin{array}{@{}l@{}} \forall a, b, f. \\ \quad \operatorname{is\text{-}hom\text{-}complex\text{-}inner\text{-}product\text{-}space}(a, b, f) \\ \quad \Leftrightarrow\; \operatorname{is\text{-}complex\text{-}inner\text{-}product\text{-}space}(a) \wedge \\ \quad \operatorname{is\text{-}complex\text{-}inner\text{-}product\text{-}space}(b) \wedge \\ \quad f \in (\operatorname{vec}(a) \to \operatorname{vec}(b)) \wedge \\ \quad \operatorname{scal}(a) = \operatorname{scal}(b) \wedge \\ \quad \forall x_{1}, x_{2} \in \operatorname{vec}(a). \\ \quad \quad f(\operatorname{vadd}(a)(x_{1}, x_{2})) = \operatorname{vadd}(b)(f(x_{1}), f(x_{2})) \wedge \\ \quad f(\operatorname{vzero}(a)) = \operatorname{vzero}(b) \wedge \\ \quad \forall x_{1} \in \operatorname{vec}(a).\; (f(\operatorname{vneg}(a)(x_{1})) = \operatorname{vneg}(b)(f(x_{1}))) \wedge \\ \quad \forall x_{1} \in \operatorname{carr}(\operatorname{scal}(a)), x_{2} \in \operatorname{vec}(a). \\ \quad \quad f(\operatorname{act}(a)(x_{1}, x_{2})) = \operatorname{act}(b)(x_{1}, f(x_{2})) \wedge \\ \quad \forall x_{1}, x_{2} \in \operatorname{vec}(a).\; (\langle\cdot,\cdot\rangle_{a}(x_{1}, x_{2}) = \langle\cdot,\cdot\rangle_{b}(f(x_{1}), f(x_{2}))) \end{array}}}
\vnbdef{complex-inner-product-space-class}{\vnbstmt{\begin{array}{@{}l@{}} \forall s. \\ \quad (s \in \operatorname{complex\text{-}inner\text{-}product\text{-}space}) \Leftrightarrow \operatorname{is\text{-}complex\text{-}inner\text{-}product\text{-}space}(s) \end{array}}}

\vnbbookfile{structure-library/bounded-metric}{structure-library/bounded-metric}
\noindent{\small bounded-metric.scm -- the bounded metric  rho = d/(1+d)  of a metric space, and the fact that it is topologically equivalent to the original (the identity map is bicontinuous).  Library phase: the construction is a def-functoid, the metric/equivalence facts are warranted supports.}\par
\par\medskip\noindent\textbf{Definition} (\texttt{bdd-metric}).\;$\operatorname{bdd\text{-}metric}(s) :=$
\begin{equation*}
{[}\operatorname{pts}(s), (\lambda {[}u, v] \in (\operatorname{pts}(s) \times \operatorname{pts}(s)).\; /(\operatorname{dist}(s)(u, v), (1 + \operatorname{dist}(s)(u, v))))]
\end{equation*}

\par\medskip\noindent\textbf{Definition} (\texttt{rr-bounded-ms}).\; $(\operatorname{rr\text{-}bounded\text{-}ms} \simeq \operatorname{bdd\text{-}metric}(\operatorname{rr\text{-}ms}))$
\par\medskip\noindent\textbf{Definition} (\texttt{is-bounded-metric-space}).\;For every $s$, 
\begin{equation*}
\begin{array}{@{}l@{}} \operatorname{is\text{-}bounded\text{-}metric\text{-}space}(s) \\ \Leftrightarrow\; \operatorname{is\text{-}metric\text{-}space}(s) \wedge \\ \exists b \in \mathbb{R}. \forall x, y \in \operatorname{pts}(s). \\ \quad \operatorname{dist}(s)(x, y) \leq b \end{array}
\end{equation*}


\vnbbookfile{structure-library/trunc-metric}{structure-library/trunc-metric}
\noindent{\small trunc-metric.scm -- the TRUNCATED METRIC  d'(x,y) = min(1, d(x,y))  of a metric space, as an ENDOFUNCTOR of METRIC-SPACE.}\par
\par\medskip\noindent\textbf{Definition} (\texttt{trunc-metric}).\;$\operatorname{trunc\text{-}metric}(s) :=$
\begin{equation*}
{[}\operatorname{pts}(s), (\lambda {[}u, v] \in (\operatorname{pts}(s) \times \operatorname{pts}(s)).\; \operatorname{min}(1, \operatorname{dist}(s)(u, v)))]
\end{equation*}


\vnbbookfile{structure-library/pseudometric}{structure-library/pseudometric}
\noindent{\small pseudometric.scm -- PSEUDOMETRIC-SPACE, the topology GENERATED by a countable family of pseudometrics (the gauge topology), and the metrizability cluster T2/T3/T4 stated as warranted PSS supports.}\par
\vnbstructdef{pseudometric-space}{A pseudometric space $s$ is a 2-tuple whose components are listed below.\par
\begin{enumerate}[nosep]
\item $\operatorname{pts}(s) = \operatorname{nth}(1, s)$, a carrier (\texttt{carriers}): $\operatorname{pts}(s) \in \mathbf{Set}$.
\item $\operatorname{dist}(s) = \operatorname{nth}(2, s)$, an operation (\texttt{op}): $\operatorname{dist}(s) \in ((\operatorname{pts}(s) \times \operatorname{pts}(s)) \to \mathbb{R})$.
\end{enumerate}
The laws (the \texttt{property} and \texttt{law} clauses):
\begin{itemize}[nosep]
\item $\operatorname{is\text{-}pseudometric}(\operatorname{dist}(s), \operatorname{pts}(s))$
\end{itemize}
The characteristic law (the axiom \texttt{is-pseudometric-space}): $\operatorname{is\text{-}pseudometric\text{-}space}(s)$ holds if and only if $\operatorname{length}(s) = 2$, every component has the typing listed with it, every law holds, and $s \in \mathbf{Set}$ (the structure is a set).
}
\par\medskip\noindent\textbf{Definition} (\texttt{pseudo-gauge-top}).\;$\operatorname{pseudo\text{-}gauge\text{-}top}(\mathit{fam}, \mathit{ground}) :=$
\begin{equation*}
{[}\mathit{ground}, \{u \in \operatorname{power}(\mathit{ground}) : \operatorname{is\text{-}open\text{-}gauge}(\mathit{fam}, \mathit{ground}, u)\}]
\end{equation*}

\par\medskip\noindent\textbf{Definition} (\texttt{is-countable-pseudometric-family}).\;For all $\mathit{fam}$ and $\mathit{ground}$, 
\begin{equation*}
\begin{array}{@{}l@{}} \operatorname{is\text{-}countable\text{-}pseudometric\text{-}family}(\mathit{fam}, \mathit{ground}) \\ \Leftrightarrow\; \forall \operatorname{n\_} \in \mathbb{N}.\; (\operatorname{is\text{-}pseudometric\text{-}space}(\operatorname{fam}(\operatorname{n\_})) \wedge (\operatorname{pts}(\operatorname{fam}(\operatorname{n\_})) \simeq \mathit{ground})) \end{array}
\end{equation*}

\par\medskip\noindent\textbf{Definition} (\texttt{is-open-gauge}).\;For all $\mathit{fam}$, $\mathit{ground}$ and $u$, 
\begin{equation*}
\begin{array}{@{}l@{}} \operatorname{is\text{-}open\text{-}gauge}(\mathit{fam}, \mathit{ground}, u) \\ \Leftrightarrow\; u \subseteq \mathit{ground} \wedge \\ \forall y \in u. \exists \operatorname{k\_} \in \mathbb{N}, r. \\ \quad \operatorname{pos\text{-}rr}(r) \wedge \\ \quad \forall x \in \mathit{ground}. \\ \quad \quad \forall \operatorname{n\_} \in \mathbb{N}.\; ((\operatorname{n\_} \leq \operatorname{k\_}) \Rightarrow (\operatorname{dist}(\operatorname{fam}(\operatorname{n\_}))(y, x) < r)) \Rightarrow (x \in u) \end{array}
\end{equation*}

\par\medskip\noindent\textbf{Definition} (\texttt{is-gauge-countable}).\;For every $s$, 
\begin{equation*}
\begin{array}{@{}l@{}} \operatorname{is\text{-}gauge\text{-}countable}(s) \\ \Leftrightarrow\; \operatorname{is\text{-}top\text{-}space}(s) \wedge \\ \exists \mathit{fam}. \\ \quad (\operatorname{is\text{-}countable\text{-}pseudometric\text{-}family}(\mathit{fam}, \operatorname{pts}(s)) \wedge (\operatorname{pseudo\text{-}gauge\text{-}top}(\mathit{fam}, \operatorname{pts}(s)) \simeq s)) \end{array}
\end{equation*}

\par\medskip\noindent\textbf{Definition} (\texttt{hom-pseudometric-space}).\;$\operatorname{hom\text{-}pseudometric\text{-}space}(a, b) :=$
\begin{equation*}
\{\operatorname{homf\_} \in (\operatorname{pts}(a) \to \operatorname{pts}(b)) : \operatorname{is\text{-}hom\text{-}pseudometric\text{-}space}(a, b, \operatorname{homf\_})\}
\end{equation*}

\vnbdef{is-hom-pseudometric-space-def}{\vnbstmt{\begin{array}{@{}l@{}} \forall a, b, f. \\ \quad \operatorname{is\text{-}hom\text{-}pseudometric\text{-}space}(a, b, f) \\ \quad \Leftrightarrow\; \operatorname{is\text{-}pseudometric\text{-}space}(a) \wedge \\ \quad \operatorname{is\text{-}pseudometric\text{-}space}(b) \wedge \\ \quad f \in (\operatorname{pts}(a) \to \operatorname{pts}(b)) \wedge \\ \quad \forall x_{1}, x_{2} \in \operatorname{pts}(a). \\ \quad \quad \operatorname{dist}(a)(x_{1}, x_{2}) = \operatorname{dist}(b)(f(x_{1}), f(x_{2})) \end{array}}}
\vnbdef{pseudometric-space-class}{\vnbstmt{\forall s.\; ((s \in \operatorname{pseudometric\text{-}space}) \Leftrightarrow \operatorname{is\text{-}pseudometric\text{-}space}(s))}}

\vnbbookfile{structure-library/metric-completion}{structure-library/metric-completion}
\noindent{\small RETIRED 2026-09-17 (proven, rake batch O): cauchy-setoid-is-setoid; embed-isometry -- theorem-library/rake-setoid2.scm, modulo 0.  embed-isometry was MISSPELLED here: `(EMBED M u)' applies the one-parameter functoid EMBED to the PAIR [M,u] (apply-tupling-2), which is a function on PTS([M,u]), not a point of the ...}\par
\par\medskip\noindent\textbf{Definition} (\texttt{embed}).\;$\operatorname{embed}(m) :=$
\begin{equation*}
(\lambda u \in \operatorname{pts}(m).\; \operatorname{class}(\operatorname{cauchy\text{-}setoid}(m), \operatorname{embed\text{-}seq}(m, u)))
\end{equation*}

\par\medskip\noindent\textbf{Definition} (\texttt{completion}).\;$\operatorname{completion}(m) :=$ ${[}\operatorname{quotient}(\operatorname{cauchy\text{-}setoid}(m)), \operatorname{completion\text{-}dist}(m)]$
\par\medskip\noindent\textbf{Definition} (\texttt{dist-seq}).\;$\operatorname{dist\text{-}seq}(m, f, g) :=$ $(\lambda n \in \mathbb{N}.\; \operatorname{dist}(m)(f(n), g(n)))$
\par\medskip\noindent\textbf{Definition} (\texttt{cauchy-setoid}).\;$\operatorname{cauchy\text{-}setoid}(m) :=$ ${[}\operatorname{cseq}(m), \operatorname{crel}(m)]$
\par\medskip\noindent\textbf{Definition} (\texttt{cseq-equiv}).\;For all $m$, $f$ and $g$, 
\begin{equation*}
\operatorname{cseq\text{-}equiv}(m, f, g) \Leftrightarrow \operatorname{converges\text{-}to}(\operatorname{rr\text{-}ms}, \operatorname{dist\text{-}seq}(m, f, g), 0)
\end{equation*}

\par\medskip\noindent\textbf{Definition} (\texttt{cseq}).\;$\operatorname{cseq}(m) :=$ $\{f \in (\mathbb{N} \to \operatorname{pts}(m)) : \operatorname{is\text{-}cauchy\text{-}seq}(m, f)\}$
\par\medskip\noindent\textbf{Definition} (\texttt{crel}).\;$\operatorname{crel}(m) :=$
\begin{equation*}
\{p \in (\operatorname{cseq}(m) \times \operatorname{cseq}(m)) : \operatorname{cseq\text{-}equiv}(m, \operatorname{nth}(1, p), \operatorname{nth}(2, p))\}
\end{equation*}

\par\medskip\noindent\textbf{Definition} (\texttt{completion-dist}).\;$\operatorname{completion\text{-}dist}(m) :=$
\begin{equation*}
\begin{array}{@{}l@{}} (\lambda p \in (\operatorname{quotient}(\operatorname{cauchy\text{-}setoid}(m)) \times \operatorname{quotient}(\operatorname{cauchy\text{-}setoid}(m))).\; \\ \quad (\iota\, \mathit{dval}.\;((\mathit{dval} \in \mathbb{R}) \wedge \exists f, g \in \operatorname{cseq}(m).\; ((\operatorname{nth}(1, p) = \operatorname{class}(\operatorname{cauchy\text{-}setoid}(m), f)) \wedge ((\operatorname{nth}(2, p) = \operatorname{class}(\operatorname{cauchy\text{-}setoid}(m), g)) \wedge \operatorname{converges\text{-}to}(\operatorname{rr\text{-}ms}, \operatorname{dist\text{-}seq}(m, f, g), \mathit{dval})))))) \end{array}
\end{equation*}

\par\medskip\noindent\textbf{Definition} (\texttt{embed-seq}).\;$\operatorname{embed\text{-}seq}(m, u) :=$ $(\lambda n \in \mathbb{N}.\; u)$

\vnbbookfile{structure-library/extended-reals}{structure-library/extended-reals}
\noindent{\small extended-reals.scm -- the extended real line RR-STAR = [-inf, +inf]}\par
\vnbdef{rr-star-membership}{\vnbstmt{\forall x.\; ((x \in \operatorname{rr\text{-}star}) \Leftrightarrow ((x \in \mathbb{R}) \vee ((x = \operatorname{pos\text{-}inf}) \vee (x = \operatorname{neg\text{-}inf}))))}}
\vnbaxiom{pos-inf-upper-bound}{\vnbstmt{\forall x \in \operatorname{rr\text{-}star}.\; (x \leq \operatorname{pos\text{-}inf})}}
\vnbaxiom{neg-inf-lower-bound}{\vnbstmt{\forall x \in \operatorname{rr\text{-}star}.\; (\operatorname{neg\text{-}inf} \leq x)}}
\vnbaxiom{rr-star-leq-domain}{\vnbstmt{\forall a, b.\; ((a \leq b) \Rightarrow ((a \in \operatorname{rr\text{-}star}) \wedge (b \in \operatorname{rr\text{-}star})))}}
\vnbaxiom{rr-star-leq-antisymmetric}{\vnbstmt{\begin{array}{@{}l@{}} \forall a, b. \\ \quad (((a \in \operatorname{rr\text{-}star}) \wedge (b \in \operatorname{rr\text{-}star})) \wedge ((a \leq b) \wedge (b \leq a))) \Rightarrow (a = b) \end{array}}}

\vnbbookfile{structure-library/extended-reals-pos}{structure-library/extended-reals-pos}
\noindent{\small extended-reals-pos.scm -- RR-POS-STAR = [0, +inf], the nonnegative extended reals}\par
\par\medskip\noindent\textbf{Definition} (\texttt{esup}).\;$\operatorname{esup}(\operatorname{s\_}) :=$
\begin{equation*}
(\iota\, \operatorname{bb\_}.\;((\operatorname{bb\_} \in \operatorname{rr\text{-}pos\text{-}star}) \wedge (\forall x \in \operatorname{s\_}.\; (x \leq \operatorname{bb\_}) \wedge \forall b.\; (((b \in \operatorname{rr\text{-}pos\text{-}star}) \wedge \forall x \in \operatorname{s\_}.\; (x \leq b)) \Rightarrow (\operatorname{bb\_} \leq b)))))
\end{equation*}

\par\medskip\noindent\textbf{Definition} (\texttt{rr-pos-star-add-monoid}).\; $\operatorname{rr\text{-}pos\text{-}star\text{-}add\text{-}monoid} = {[}\operatorname{rr\text{-}pos\text{-}star}, \mathit{eplus}, 0]$
\vnbdef{rr-pos-star-membership}{\vnbstmt{\forall x.\; ((x \in \operatorname{rr\text{-}pos\text{-}star}) \Leftrightarrow (((x \in \mathbb{R}) \wedge (0 \leq x)) \vee (x = \operatorname{pos\text{-}inf})))}}
\vnbdef{eplus-def}{\vnbstmt{(\mathit{eplus} \simeq (\lambda {[}\operatorname{x\_}, \operatorname{y\_}] \in (\operatorname{rr\text{-}pos\text{-}star} \times \operatorname{rr\text{-}pos\text{-}star}).\; \left(\text{if }((\operatorname{x\_} = \operatorname{pos\text{-}inf}) \vee (\operatorname{y\_} = \operatorname{pos\text{-}inf}))\text{ then }\operatorname{pos\text{-}inf}\text{ else }(\operatorname{x\_} + \operatorname{y\_})\right)))}}
\vnbdef{rr-pos-star-add-monoid@carr}{\vnbstmt{(\operatorname{carr}(\operatorname{rr\text{-}pos\text{-}star\text{-}add\text{-}monoid}) \simeq \operatorname{rr\text{-}pos\text{-}star})}}
\vnbdef{rr-pos-star-add-monoid@opr}{\vnbstmt{(\operatorname{opr}(\operatorname{rr\text{-}pos\text{-}star\text{-}add\text{-}monoid}) \simeq \mathit{eplus})}}
\vnbdef{rr-pos-star-add-monoid@iden}{\vnbstmt{(\operatorname{iden}(\operatorname{rr\text{-}pos\text{-}star\text{-}add\text{-}monoid}) \simeq 0)}}

\vnbbookfile{structure-library/extended-reals-updown}{structure-library/extended-reals-updown}
\noindent{\small extinf-vocab.scm -- the two half-extended real lines and extended addition on each. DEFINITIONS ONLY.  Agent EXT-INF, 2026-10-26 (the user's decision of the same day: the semigroup properties of extended addition "clearly stated for both"). DESTINATION: structure-library/extended-reals-updown.scm (new file), loaded ...}\par
\par\medskip\noindent\textbf{Definition} (\texttt{rr-ext-up}).\; $(\operatorname{rr\text{-}ext\text{-}up} \simeq (\mathbb{R} \cup \operatorname{singleton}(\operatorname{pos\text{-}inf})))$
\par\medskip\noindent\textbf{Definition} (\texttt{rr-ext-down}).\; $(\operatorname{rr\text{-}ext\text{-}down} \simeq (\mathbb{R} \cup \operatorname{singleton}(\operatorname{neg\text{-}inf})))$
\par\medskip\noindent\textbf{Definition} (\texttt{eplus-up}).\;$\operatorname{eplus\text{-}up}(\operatorname{eux\_}, \operatorname{euy\_}) :=$ $\left(\text{if }((\operatorname{eux\_} = \operatorname{pos\text{-}inf}) \vee (\operatorname{euy\_} = \operatorname{pos\text{-}inf}))\text{ then }\operatorname{pos\text{-}inf}\text{ else }(\operatorname{eux\_} + \operatorname{euy\_})\right)$
\par\medskip\noindent\textbf{Definition} (\texttt{eplus-down}).\;$\operatorname{eplus\text{-}down}(\operatorname{edx\_}, \operatorname{edy\_}) :=$ $\left(\text{if }((\operatorname{edx\_} = \operatorname{neg\text{-}inf}) \vee (\operatorname{edy\_} = \operatorname{neg\text{-}inf}))\text{ then }\operatorname{neg\text{-}inf}\text{ else }(\operatorname{edx\_} + \operatorname{edy\_})\right)$

\vnbbookfile{structure-library/nth-deriv}{structure-library/nth-deriv}
\noindent{\small nth-deriv.scm -- THE ITERATED DERIVATIVE ON AN OPEN SET, OVER A NORMED FIELD. DEFINITION ONLY; the laws are PROVEN in theorem-library/nth-deriv-laws.scm. Batch 36, 2026-09-26.}\par
\par\medskip\noindent\textbf{Definition} (\texttt{nth-deriv-on}).\;For all $\operatorname{ndvk\_}$, $\operatorname{ndvu\_}$ and $\operatorname{ndvf\_}$,  $(\operatorname{ndvf\_}^{(0)} \simeq \operatorname{ndvf\_})$\par
Let $\operatorname{ndvk\_}$, $\operatorname{ndvu\_}$ and $\operatorname{ndvf\_}$ be arbitrary. Let $\operatorname{ndvn\_} \in \mathbb{N}$. Then
\begin{equation*}
(\operatorname{ndvf\_}^{(\operatorname{succ}(\operatorname{ndvn\_}))} \simeq (\lambda \operatorname{ndvz\_} \in \operatorname{ndvu\_}.\; \operatorname{deriv\text{-}on}(\operatorname{ndvk\_}, \operatorname{ndvu\_}, \operatorname{ndvf\_}^{(\operatorname{ndvn\_})}, \operatorname{ndvz\_})))
\end{equation*}


\vnbbookfile{structure-library/block-family-combinatorial}{structure-library/block-family-combinatorial}
\noindent{\small MOVED 2026-09-20 (batch 12-A) from theorem-library/: this file is VOCABULARY -- definitions, notation and warranted supports, not one proof -- and every theorem stated with it had to load below it.  Its load.scm slot is unchanged. structure-library/block-family-combinatorial.scm}\par
\par\medskip\noindent\textbf{Definition} (\texttt{is-finite-cover}).\;For all $c$ and $a$, 
\begin{equation*}
\operatorname{is\text{-}finite\text{-}cover}(c, a) \Leftrightarrow ((c \in \mathbf{Set}) \wedge ((|c| \in \mathbb{N}) \wedge (a \subseteq \operatorname{big\text{-}union}(u, c, u))))
\end{equation*}


\vnbbookfile{structure-library/subsequence-capture}{structure-library/subsequence-capture}
\noindent{\small MOVED 2026-09-20 (batch 12-A) from theorem-library/: this file is VOCABULARY -- definitions, notation and warranted supports, not one proof -- and every theorem stated with it had to load below it.  Its load.scm slot is unchanged. structure-library/subsequence-capture.scm}\par
\par\medskip\noindent\textbf{Definition} (\texttt{nn-enum}).\;$\operatorname{nn\text{-}enum}(s) :=$
\begin{equation*}
\varepsilon(\{f \in (\mathbb{N} \to s) : \forall m, n \in \mathbb{N}.\; ((m < n) \Rightarrow (f(m) < f(n)))\})
\end{equation*}


\vnbbookfile{structure-library/cauchy-subsequence}{structure-library/cauchy-subsequence}
\noindent{\small MOVED 2026-09-20 (batch 12-A) from theorem-library/: this file is VOCABULARY -- definitions, notation and warranted supports, not one proof -- and every theorem stated with it had to load below it.  Its load.scm slot is unchanged. RETIRED 2026-09-17 (proven): subseq-is-fun -- theorem-library/rake-analysis2.scm}\par
\par\medskip\noindent\textbf{Definition} (\texttt{strictly-mono-nn}).\;For every $\varphi$, 
\begin{equation*}
\begin{array}{@{}l@{}} \operatorname{strictly\text{-}mono\text{-}nn}(\varphi) \\ \Leftrightarrow\; ((\varphi \in (\mathbb{N} \to \mathbb{N})) \wedge \forall m, \operatorname{n\_} \in \mathbb{N}.\; ((m < \operatorname{n\_}) \Rightarrow (\varphi(m) < \varphi(\operatorname{n\_})))) \end{array}
\end{equation*}

\par\medskip\noindent\textbf{Definition} (\texttt{subseq}).\;$\operatorname{subseq}(f, \varphi) :=$ $(\lambda k \in \mathbb{N}.\; f(\varphi(k)))$
\par\medskip\noindent\textbf{Definition} (\texttt{is-subsequence}).\;For all $s$, $y$ and $f$, 
\begin{equation*}
\begin{array}{@{}l@{}} \operatorname{is\text{-}subsequence}(s, y, f) \\ \Leftrightarrow\; ((f \in (\mathbb{N} \to \operatorname{pts}(s))) \wedge \exists \varphi.\; (\operatorname{strictly\text{-}mono\text{-}nn}(\varphi) \wedge (y = \operatorname{subseq}(f, \varphi)))) \end{array}
\end{equation*}

\par\medskip\noindent\textbf{Definition} (\texttt{subsqn}).\;For all $y$ and $f$, 
\begin{equation*}
\operatorname{subsqn}(y, f) \Leftrightarrow \exists \varphi.\; (\operatorname{strictly\text{-}mono\text{-}nn}(\varphi) \wedge (y = \operatorname{subseq}(f, \varphi)))
\end{equation*}

\par\medskip\noindent\textbf{Definition} (\texttt{is-eps-cauchy-seq}).\;For all $s$, $\mathit{eps}$ and $y$, 
\begin{equation*}
\begin{array}{@{}l@{}} \operatorname{is\text{-}eps\text{-}cauchy\text{-}seq}(s, \mathit{eps}, y) \\ \Leftrightarrow\; \operatorname{is\text{-}metric\text{-}space}(s) \wedge \\ y \in (\mathbb{N} \to \operatorname{pts}(s)) \wedge \\ \forall m, \operatorname{n\_} \in \mathbb{N}.\; (\operatorname{dist}(s)(y(m), y(\operatorname{n\_})) \leq \mathit{eps}) \end{array}
\end{equation*}

\par\medskip\noindent\textbf{Definition} (\texttt{null-rr-seq}).\;For every $\mathit{rad}$, 
\begin{equation*}
\begin{array}{@{}l@{}} \operatorname{null\text{-}rr\text{-}seq}(\mathit{rad}) \\ \Leftrightarrow\; \mathit{rad} \in (\mathbb{N} \to \mathbb{R}) \wedge \\ \forall k \in \mathbb{N}.\; \operatorname{pos\text{-}rr}(\operatorname{rad}(k)) \wedge \\ \forall \mathit{eps}. \\ \quad \operatorname{pos\text{-}rr}(\mathit{eps}) \Rightarrow  \\ \quad \quad \exists n \in \mathbb{N}.\; \forall k.\; (((k \in \mathbb{N}) \wedge (n \leq k)) \Rightarrow (\operatorname{rad}(k) \leq \mathit{eps})) \end{array}
\end{equation*}


\vnbbookfile{structure-library/extended-sum}{structure-library/extended-sum}
\noindent{\small MOVED 2026-09-20 (batch 12-A) from theorem-library/: this file is VOCABULARY -- definitions, notation and warranted supports, not one proof -- and every theorem stated with it had to load below it.  Its load.scm slot is unchanged. structure-library/extended-sum.scm -- ESUM: the unordered sum in RR-POS-STAR}\par
\par\medskip\noindent\textbf{Definition} (\texttt{esum}).\;$\operatorname{esum}(\operatorname{f\_}) :=$
\begin{equation*}
\operatorname{esup}(\{\operatorname{v\_} \in \operatorname{rr\text{-}pos\text{-}star} : \exists \operatorname{t\_}.\; (((\operatorname{t\_} \in \mathbf{Set}) \wedge ((|\operatorname{t\_}| \in \mathbb{N}) \wedge (\operatorname{t\_} \subseteq \operatorname{dom}(\operatorname{f\_})))) \wedge (\operatorname{v\_} = \operatorname{finsum}(\operatorname{rr\text{-}pos\text{-}star\text{-}add\text{-}monoid}, \operatorname{f\_}, \operatorname{t\_})))\})
\end{equation*}


\vnbbookfile{structure-library/extended-arith}{structure-library/extended-arith}
\noindent{\small extended-arith.scm -- the rest of the arithmetic and order vocabulary of RR-POS-STAR = [0, +inf]: multiplication, infima, and the limit operations of a sequence.}\par
\par\medskip\noindent\textbf{Definition} (\texttt{einf}).\;$\operatorname{einf}(s) :=$ $\operatorname{esup}(\operatorname{elower\text{-}bounds}(s))$
\par\medskip\noindent\textbf{Definition} (\texttt{etail}).\;$\operatorname{etail}(f, n) :=$ $\operatorname{image}(f, \{\operatorname{p\_} \in \mathbb{N} : (n \leq \operatorname{p\_})\})$
\par\medskip\noindent\textbf{Definition} (\texttt{elower-bounds}).\;$\operatorname{elower\text{-}bounds}(s) :=$ $\{\operatorname{b\_} \in \operatorname{rr\text{-}pos\text{-}star} : \forall \operatorname{x\_} \in s.\; (\operatorname{b\_} \leq \operatorname{x\_})\}$
\par\medskip\noindent\textbf{Definition} (\texttt{eliminf}).\;$\operatorname{eliminf}(f) :=$ $\operatorname{esup}(\operatorname{image}((\lambda \operatorname{n\_} \in \mathbb{N}.\; \operatorname{einf}(\operatorname{etail}(f, \operatorname{n\_}))), \mathbb{N}))$
\par\medskip\noindent\textbf{Definition} (\texttt{elimsup}).\;$\operatorname{elimsup}(f) :=$ $\operatorname{einf}(\operatorname{image}((\lambda \operatorname{n\_} \in \mathbb{N}.\; \operatorname{esup}(\operatorname{etail}(f, \operatorname{n\_}))), \mathbb{N}))$
\par\medskip\noindent\textbf{Definition} (\texttt{econverges-to}).\;For all $f$ and $l$, 
\begin{equation*}
\begin{array}{@{}l@{}} \operatorname{econverges\text{-}to}(f, l) \\ \Leftrightarrow\; ((f \in (\mathbb{N} \to \operatorname{rr\text{-}pos\text{-}star})) \wedge ((l \in \operatorname{rr\text{-}pos\text{-}star}) \wedge ((\operatorname{eliminf}(f) = l) \wedge (\operatorname{elimsup}(f) = l)))) \end{array}
\end{equation*}

\vnbdef{etimes-def}{\vnbstmt{(\mathit{etimes} \simeq (\lambda {[}\operatorname{x\_}, \operatorname{y\_}] \in (\operatorname{rr\text{-}pos\text{-}star} \times \operatorname{rr\text{-}pos\text{-}star}).\; \left(\text{if }((\operatorname{x\_} = \operatorname{pos\text{-}inf}) \vee (\operatorname{y\_} = \operatorname{pos\text{-}inf}))\text{ then }\left(\text{if }((\operatorname{x\_} = 0) \vee (\operatorname{y\_} = 0))\text{ then }0\text{ else }\operatorname{pos\text{-}inf}\right)\text{ else }(\operatorname{x\_} \cdot \operatorname{y\_})\right)))}}

\vnbbookfile{structure-library/summability}{structure-library/summability}
\noindent{\small MOVED 2026-09-20 (batch 12-A) from theorem-library/: this file is VOCABULARY -- definitions, notation and warranted supports, not one proof -- and every theorem stated with it had to load below it.  Its load.scm slot is unchanged. summability.scm -- unconditional summability of a normed-AG-valued function}\par
\par\medskip\noindent\textbf{Definition} (\texttt{sums-to}).\;For all $\mathit{grp}$, $f$ and $r$, 
\begin{equation*}
\begin{array}{@{}l@{}} \operatorname{sums\text{-}to}(\mathit{grp}, f, r) \\ \Leftrightarrow\; \forall \mathit{eps}. \\ \quad \operatorname{pos\text{-}rr}(\mathit{eps}) \Rightarrow  \\ \quad \quad \exists \mathit{fin} \in \mathbf{Set}. \\ \quad \quad \quad |\mathit{fin}| \in \mathbb{N} \wedge \\ \quad \quad \quad \mathit{fin} \subseteq \operatorname{dom}(f) \wedge \\ \quad \quad \quad \forall \mathit{ext}. \\ \quad \quad \quad \quad ((\mathit{ext} \in \mathbf{Set}) \wedge ((|\mathit{ext}| \in \mathbb{N}) \wedge ((\mathit{fin} \subseteq \mathit{ext}) \wedge (\mathit{ext} \subseteq \operatorname{dom}(f))))) \Rightarrow  \\ \quad \quad \quad \quad \quad \operatorname{nrm}(\mathit{grp})(\operatorname{opr}(\mathit{grp})(r, \operatorname{inv}(\mathit{grp})(\operatorname{finsum}(\operatorname{normed\text{-}ag\text{-}as\text{-}abelian\text{-}group}(\mathit{grp}), f, \mathit{ext})))) < \mathit{eps} \end{array}
\end{equation*}

\par\medskip\noindent\textbf{Definition} (\texttt{is-summable}).\;For all $\mathit{grp}$ and $f$, 
\begin{equation*}
\operatorname{is\text{-}summable}(\mathit{grp}, f) \Leftrightarrow \exists r \in \operatorname{carr}(\mathit{grp}).\; \operatorname{sums\text{-}to}(\mathit{grp}, f, r)
\end{equation*}

\par\medskip\noindent\textbf{Definition} (\texttt{is-absolutely-summable}).\;For all $\mathit{grp}$ and $f$, 
\begin{equation*}
\operatorname{is\text{-}absolutely\text{-}summable}(\mathit{grp}, f) \Leftrightarrow (\operatorname{esum}((\lambda i \in \operatorname{dom}(f).\; \operatorname{nrm}(\mathit{grp})(f(i)))) \in \mathbb{R})
\end{equation*}


\vnbbookfile{structure-library/power-series}{structure-library/power-series}
\noindent{\small MOVED 2026-09-20 (batch 12-A) from theorem-library/: this file is VOCABULARY -- definitions, notation and warranted supports, not one proof -- and every theorem stated with it had to load below it.  Its load.scm slot is unchanged. power-series.scm -- real power series  Sum\_\{n$>$=0\} coef(n) x\textasciicircum{}n  in RR.}\par
\par\medskip\noindent\textbf{Definition} (\texttt{ps-partial-sum}).\;$\operatorname{ps\text{-}partial\text{-}sum}(\mathit{coef}, x, k) :=$
\begin{equation*}
\operatorname{sum\text{-}ag}(\operatorname{normed\text{-}field\text{-}additive\text{-}ag}(\operatorname{rr\text{-}normed\text{-}field}), (\lambda n \in \mathbb{N}.\; (\operatorname{coef}(n) \cdot \operatorname{power}(x, n))), k)
\end{equation*}

\par\medskip\noindent\textbf{Definition} (\texttt{ps-converges-to-at}).\;For all $\mathit{coef}$, $x$ and $l$, 
\begin{equation*}
\begin{array}{@{}l@{}} \operatorname{ps\text{-}converges\text{-}to\text{-}at}(\mathit{coef}, x, l) \\ \Leftrightarrow\; \operatorname{converges\text{-}to}(\operatorname{rr\text{-}ms}, (\lambda k \in \mathbb{N}.\; \operatorname{ps\text{-}partial\text{-}sum}(\mathit{coef}, x, k)), l) \end{array}
\end{equation*}

\par\medskip\noindent\textbf{Definition} (\texttt{ps-converges-at}).\;For all $\mathit{coef}$ and $x$, 
\begin{equation*}
\operatorname{ps\text{-}converges\text{-}at}(\mathit{coef}, x) \Leftrightarrow \operatorname{converges}(\operatorname{rr\text{-}ms}, (\lambda k \in \mathbb{N}.\; \operatorname{ps\text{-}partial\text{-}sum}(\mathit{coef}, x, k)))
\end{equation*}

\par\medskip\noindent\textbf{Definition} (\texttt{ps-ratio-limit}).\;For all $\mathit{coef}$ and $l$, 
\begin{equation*}
\begin{array}{@{}l@{}} \operatorname{ps\text{-}ratio\text{-}limit}(\mathit{coef}, l) \\ \Leftrightarrow\; \operatorname{converges\text{-}to}(\operatorname{rr\text{-}ms}, (\lambda n \in \mathbb{N}.\; \lvert \cdot(\operatorname{coef}(\operatorname{succ}(n)), \operatorname{recip}(\operatorname{coef}(n))) \rvert), l) \end{array}
\end{equation*}

\par\medskip\noindent\textbf{Definition} (\texttt{ps-absolutely-converges-at}).\;For all $\mathit{coef}$ and $x$, 
\begin{equation*}
\operatorname{ps\text{-}absolutely\text{-}converges\text{-}at}(\mathit{coef}, x) \Leftrightarrow \operatorname{ps\text{-}converges\text{-}at}((\lambda n \in \mathbb{N}.\; \lvert \operatorname{coef}(n) \rvert), \lvert x \rvert)
\end{equation*}

\par\medskip\noindent\textbf{Definition} (\texttt{series-partial-sum}).\;$\operatorname{series\text{-}partial\text{-}sum}(f, k) :=$ $\operatorname{sum\text{-}ag}(\operatorname{normed\text{-}field\text{-}additive\text{-}ag}(\operatorname{rr\text{-}normed\text{-}field}), f, k)$
\par\medskip\noindent\textbf{Definition} (\texttt{series-converges-to}).\;For all $f$ and $l$, 
\begin{equation*}
\operatorname{series\text{-}converges\text{-}to}(f, l) \Leftrightarrow \operatorname{converges\text{-}to}(\operatorname{rr\text{-}ms}, (\lambda k \in \mathbb{N}.\; \operatorname{series\text{-}partial\text{-}sum}(f, k)), l)
\end{equation*}

\par\medskip\noindent\textbf{Definition} (\texttt{series-converges}).\;For every $f$, 
\begin{equation*}
\operatorname{series\text{-}converges}(f) \Leftrightarrow \operatorname{converges}(\operatorname{rr\text{-}ms}, (\lambda k \in \mathbb{N}.\; \operatorname{series\text{-}partial\text{-}sum}(f, k)))
\end{equation*}


\vnbbookfile{structure-library/product-metric}{structure-library/product-metric}
\noindent{\small product-metric.scm -- the countable product of metric spaces, as a metric space carrying the PRODUCT TOPOLOGY.  Library phase: functoids for the construction, warranted supports for the metric/topology facts.}\par
\par\medskip\noindent\textbf{Definition} (\texttt{product-metric-w}).\;$\operatorname{product\text{-}metric\text{-}w}(\mathit{ms}, w) :=$
\begin{equation*}
{[}\operatorname{product\text{-}carrier}(\mathit{ms}), (\lambda {[}x, y] \in (\operatorname{product\text{-}carrier}(\mathit{ms}) \times \operatorname{product\text{-}carrier}(\mathit{ms})).\; (\iota\, l.\;((l \in \mathbb{R}) \wedge \operatorname{series\text{-}converges\text{-}to}((\lambda n \in \mathbb{N}.\; (w(n) \cdot \operatorname{dist}(\operatorname{bdd\text{-}metric}(\operatorname{ms}(n)))(x(n), y(n)))), l))))]
\end{equation*}

\par\medskip\noindent\textbf{Definition} (\texttt{product-metric}).\;$\operatorname{product\text{-}metric}(\mathit{ms}) :=$ $\operatorname{product\text{-}metric\text{-}w}(\mathit{ms}, (\lambda n \in \mathbb{N}.\; /(1, \operatorname{power}(2, (n + 1)))))$
\par\medskip\noindent\textbf{Definition} (\texttt{is-ms-sequence}).\;For every $\mathit{ms}$, 
\begin{equation*}
\operatorname{is\text{-}ms\text{-}sequence}(\mathit{ms}) \Leftrightarrow \forall n \in \mathbb{N}.\; \operatorname{is\text{-}metric\text{-}space}(\operatorname{ms}(n))
\end{equation*}

\par\medskip\noindent\textbf{Definition} (\texttt{summable-weight}).\;For every $w$, 
\begin{equation*}
\operatorname{summable\text{-}weight}(w) \Leftrightarrow ((w \in (\mathbb{N} \to \mathbb{R})) \wedge (\forall n \in \mathbb{N}.\; (0 < w(n)) \wedge \operatorname{series\text{-}converges}(w)))
\end{equation*}

\par\medskip\noindent\textbf{Definition} (\texttt{product-carrier}).\;$\operatorname{product\text{-}carrier}(\mathit{ms}) :=$
\begin{equation*}
\{x \in (\mathbb{N} \to \operatorname{big\text{-}union}(n, \mathbb{N}, \operatorname{pts}(\operatorname{ms}(n)))) : \forall n \in \mathbb{N}.\; (x(n) \in \operatorname{pts}(\operatorname{ms}(n)))\}
\end{equation*}

\par\medskip\noindent\textbf{Definition} (\texttt{product-proj}).\;$\operatorname{product\text{-}proj}(\mathit{ms}, n) :=$ $(\lambda x \in \operatorname{product\text{-}carrier}(\mathit{ms}).\; x(n))$

\vnbbookfile{structure-library/c-metric-space}{structure-library/c-metric-space}
\noindent{\small c-metric-space.scm -- C-METRIC-SPACE, a set carrying a COUNTABLE FAMILY OF PSEUDOMETRICS that separates points, and the CANONICAL METRIC it induces.}\par
\vnbstructdef{c-metric-space}{A c metric space $s$ is a 2-tuple whose components are listed below.\par
\begin{enumerate}[nosep]
\item $\operatorname{pts}(s) = \operatorname{nth}(1, s)$, a carrier (\texttt{carriers}): $\operatorname{pts}(s) \in \mathbf{Set}$.
\item $\operatorname{dists}(s) = \operatorname{nth}(2, s)$, an operation (\texttt{op}): $\operatorname{dists}(s) \in (\mathbb{N} \to ((\operatorname{pts}(s) \times \operatorname{pts}(s)) \to \mathbb{R}))$.
\end{enumerate}
The laws (the \texttt{property} and \texttt{law} clauses):
\begin{itemize}[nosep]
\item $\forall \operatorname{k\_} \in \mathbb{N}.\; \operatorname{is\text{-}pseudometric}(\operatorname{dists}(s)(\operatorname{k\_}), \operatorname{pts}(s))$
\item
\begin{equation*}
\begin{array}{@{}l@{}} \forall u, v \in \operatorname{pts}(s). \\ \quad \forall \operatorname{k\_} \in \mathbb{N}.\; (\operatorname{dists}(s)(\operatorname{k\_})(u, v) = 0) \Rightarrow (u = v) \end{array}
\end{equation*}

\end{itemize}
The characteristic law (the axiom \texttt{is-c-metric-space}): $\operatorname{is\text{-}c\text{-}metric\text{-}space}(s)$ holds if and only if $\operatorname{length}(s) = 2$, every component has the typing listed with it, every law holds, and $s \in \mathbf{Set}$ (the structure is a set).
}
\par\medskip\noindent\textbf{Definition} (\texttt{c-metric}).\;$\operatorname{c\text{-}metric}(s) :=$ $\operatorname{c\text{-}metric\text{-}w}(s, (\lambda n \in \mathbb{N}.\; /(1, \operatorname{power}(2, (n + 1)))))$
\par\medskip\noindent\textbf{Definition} (\texttt{c-metric-w}).\;$\operatorname{c\text{-}metric\text{-}w}(s, w) :=$
\begin{equation*}
{[}\operatorname{pts}(s), (\lambda {[}u, v] \in (\operatorname{pts}(s) \times \operatorname{pts}(s)).\; (\iota\, l.\;((l \in \mathbb{R}) \wedge \operatorname{series\text{-}converges\text{-}to}((\lambda k \in \mathbb{N}.\; (w(k) \cdot \operatorname{min}(1, \operatorname{dists}(s)(k)(u, v)))), l))))]
\end{equation*}

\par\medskip\noindent\textbf{Definition} (\texttt{c-pseudo-family}).\;$\operatorname{c\text{-}pseudo\text{-}family}(s) :=$ $(\lambda k \in \mathbb{N}.\; {[}\operatorname{pts}(s), \operatorname{dists}(s)(k)])$
\par\medskip\noindent\textbf{Definition} (\texttt{hom-c-metric-space}).\;$\operatorname{hom\text{-}c\text{-}metric\text{-}space}(a, b) :=$
\begin{equation*}
\{\operatorname{homf\_} \in (\operatorname{pts}(a) \to \operatorname{pts}(b)) : \operatorname{is\text{-}hom\text{-}c\text{-}metric\text{-}space}(a, b, \operatorname{homf\_})\}
\end{equation*}

\vnbdef{is-hom-c-metric-space-def}{\vnbstmt{\begin{array}{@{}l@{}} \forall a, b, f. \\ \quad \operatorname{is\text{-}hom\text{-}c\text{-}metric\text{-}space}(a, b, f) \\ \quad \Leftrightarrow\; \operatorname{is\text{-}c\text{-}metric\text{-}space}(a) \wedge \\ \quad \operatorname{is\text{-}c\text{-}metric\text{-}space}(b) \wedge \\ \quad f \in (\operatorname{pts}(a) \to \operatorname{pts}(b)) \wedge \\ \quad \forall x_{1} \in \mathbb{N}.\; (\operatorname{dists}(a)(x_{1}) = \operatorname{dists}(b)(x_{1})) \end{array}}}
\vnbdef{c-metric-space-class}{\vnbstmt{\forall s.\; ((s \in \operatorname{c\text{-}metric\text{-}space}) \Leftrightarrow \operatorname{is\text{-}c\text{-}metric\text{-}space}(s))}}

\vnbbookfile{structure-library/seq-compact-product}{structure-library/seq-compact-product}
\noindent{\small MOVED 2026-09-20 (batch 12-A) from theorem-library/: this file is VOCABULARY -- definitions, notation and warranted supports, not one proof -- and every theorem stated with it had to load below it.  Its load.scm slot is unchanged. structure-library/seq-compact-product.scm}\par
\par\medskip\noindent\textbf{Definition} (\texttt{seq-compact}).\;For every $s$, 
\begin{equation*}
\begin{array}{@{}l@{}} \operatorname{seq\text{-}compact}(s) \\ \Leftrightarrow\; \operatorname{is\text{-}metric\text{-}space}(s) \wedge \\ \forall f \in (\mathbb{N} \to \operatorname{pts}(s)). \exists \varphi. \\ \quad (\operatorname{strictly\text{-}mono\text{-}nn}(\varphi) \wedge \exists l \in \operatorname{pts}(s).\; \operatorname{converges\text{-}to}(s, \operatorname{subseq}(f, \varphi), l)) \end{array}
\end{equation*}

\par\medskip\noindent\textbf{Definition} (\texttt{converges-along}).\;For all $s$, $g$, $b$ and $p$, 
\begin{equation*}
\begin{array}{@{}l@{}} \operatorname{converges\text{-}along}(s, g, b, p) \\ \Leftrightarrow\; \operatorname{is\text{-}metric\text{-}space}(s) \wedge \\ g \in (\mathbb{N} \to \operatorname{pts}(s)) \wedge \\ p \in \operatorname{pts}(s) \wedge \\ \forall \mathit{eps}. \\ \quad \operatorname{pos\text{-}rr}(\mathit{eps}) \Rightarrow  \\ \quad \quad \exists n \in \mathbb{N}. \forall i \in \mathbb{N}. \\ \quad \quad \quad (i \in b) \Rightarrow ((n \leq i) \Rightarrow (\operatorname{dist}(s)(g(i), p) \leq \mathit{eps})) \end{array}
\end{equation*}


\vnbbookfile{structure-library/basw-vocab}{structure-library/basw-vocab}
\noindent{\small basw-vocab.scm -- wave of 2026-11-09/10 (agent BA-SW), wired 2026-11-10 from scratchpad/basw-1109/basw-vocab.scm. Library: analysis.}\par
\par\medskip\noindent\textbf{Definition} (\texttt{rcont}).\;$C_{\mathbb{R}}({\operatorname{rct\_}}) :=$
\begin{equation*}
\{\operatorname{rcf\_} \in (\operatorname{pts}(\operatorname{rct\_}) \to \mathbb{R}) : \operatorname{is\text{-}hom\text{-}top\text{-}space}(\operatorname{rct\_}, \operatorname{metric\text{-}top}(\operatorname{rr\text{-}ms}), \operatorname{rcf\_})\}
\end{equation*}

\par\medskip\noindent\textbf{Definition} (\texttt{is-rcont-subalgebra}).\;For all $\operatorname{rst\_}$ and $\operatorname{rsa\_}$, 
\begin{equation*}
\begin{array}{@{}l@{}} \operatorname{is\text{-}rcont\text{-}subalgebra}(\operatorname{rst\_}, \operatorname{rsa\_}) \\ \Leftrightarrow\; \operatorname{rsa\_} \subseteq C_{\mathbb{R}}({\operatorname{rst\_}}) \wedge \\ \forall \operatorname{rsc\_} \in \mathbb{R}.\; ((\lambda \operatorname{rsz\_} \in \operatorname{pts}(\operatorname{rst\_}).\; \operatorname{rsc\_}) \in \operatorname{rsa\_}) \wedge \\ \forall \operatorname{rsf\_}, \operatorname{rsg\_} \in \operatorname{rsa\_}. \\ \quad (\lambda \operatorname{rsz\_} \in \operatorname{pts}(\operatorname{rst\_}).\; +(\operatorname{rsf\_}(\operatorname{rsz\_}), \operatorname{rsg\_}(\operatorname{rsz\_}))) \in \operatorname{rsa\_} \wedge \\ \forall \operatorname{rsf\_}, \operatorname{rsg\_} \in \operatorname{rsa\_}. \\ \quad (\lambda \operatorname{rsz\_} \in \operatorname{pts}(\operatorname{rst\_}).\; \cdot(\operatorname{rsf\_}(\operatorname{rsz\_}), \operatorname{rsg\_}(\operatorname{rsz\_}))) \in \operatorname{rsa\_} \end{array}
\end{equation*}

\par\medskip\noindent\textbf{Definition} (\texttt{is-point-separating}).\;For all $\operatorname{rst\_}$ and $\operatorname{rsa\_}$, 
\begin{equation*}
\begin{array}{@{}l@{}} \operatorname{is\text{-}point\text{-}separating}(\operatorname{rst\_}, \operatorname{rsa\_}) \\ \Leftrightarrow\; \forall \operatorname{rsp\_}, \operatorname{rsq\_} \in \operatorname{pts}(\operatorname{rst\_}). \\ \quad \neg (\operatorname{rsp\_} = \operatorname{rsq\_}) \Rightarrow \exists \operatorname{rsf\_} \in \operatorname{rsa\_}.\; \neg (\operatorname{rsf\_}(\operatorname{rsp\_}) = \operatorname{rsf\_}(\operatorname{rsq\_})) \end{array}
\end{equation*}

\par\medskip\noindent\textbf{Definition} (\texttt{rcont-closure}).\;$\operatorname{rcont\text{-}closure}(\operatorname{rst\_}, \operatorname{rsa\_}) :=$
\begin{equation*}
\{\operatorname{rsh\_} \in C_{\mathbb{R}}({\operatorname{rst\_}}) : \forall \operatorname{rse\_}.\; (\operatorname{pos\text{-}rr}(\operatorname{rse\_}) \Rightarrow \exists \operatorname{rsg\_} \in \operatorname{rsa\_}.\; \forall \operatorname{rsy\_} \in \operatorname{pts}(\operatorname{rst\_}).\; (\lvert (\operatorname{rsh\_}(\operatorname{rsy\_}) - \operatorname{rsg\_}(\operatorname{rsy\_})) \rvert < \operatorname{rse\_}))\}
\end{equation*}

\par\medskip\noindent\textbf{Definition} (\texttt{is-ccont-star-subalgebra}).\;For all $\operatorname{rst\_}$ and $\operatorname{rsa\_}$, 
\begin{equation*}
\begin{array}{@{}l@{}} \operatorname{is\text{-}ccont\text{-}star\text{-}subalgebra}(\operatorname{rst\_}, \operatorname{rsa\_}) \\ \Leftrightarrow\; \forall \operatorname{csf\_} \in \operatorname{rsa\_}. \\ \quad (\operatorname{csf\_} \in (\operatorname{pts}(\operatorname{rst\_}) \to \mathbb{C})) \wedge \operatorname{is\text{-}hom\text{-}top\text{-}space}(\operatorname{rst\_}, \operatorname{metric\text{-}top}(\operatorname{cc\text{-}ms}), \operatorname{csf\_}) \wedge \\ \forall \operatorname{csc\_} \in \mathbb{C}.\; ((\lambda \operatorname{rsz\_} \in \operatorname{pts}(\operatorname{rst\_}).\; \operatorname{csc\_}) \in \operatorname{rsa\_}) \wedge \\ \forall \operatorname{csf\_}, \operatorname{csg\_} \in \operatorname{rsa\_}. \\ \quad (\lambda \operatorname{rsz\_} \in \operatorname{pts}(\operatorname{rst\_}).\; +(\operatorname{csf\_}(\operatorname{rsz\_}), \operatorname{csg\_}(\operatorname{rsz\_}))) \in \operatorname{rsa\_} \wedge \\ \forall \operatorname{csf\_}, \operatorname{csg\_} \in \operatorname{rsa\_}. \\ \quad (\lambda \operatorname{rsz\_} \in \operatorname{pts}(\operatorname{rst\_}).\; \cdot(\operatorname{csf\_}(\operatorname{rsz\_}), \operatorname{csg\_}(\operatorname{rsz\_}))) \in \operatorname{rsa\_} \wedge \\ \forall \operatorname{csf\_} \in \operatorname{rsa\_}. \\ \quad (\lambda \operatorname{rsz\_} \in \operatorname{pts}(\operatorname{rst\_}).\; \operatorname{conjugate}(\operatorname{csf\_}(\operatorname{rsz\_}))) \in \operatorname{rsa\_} \end{array}
\end{equation*}


\vnbbookfile{structure-library/hilbert-onb-vocab}{structure-library/hilbert-onb-vocab}
\noindent{\small hilbert-onb-vocab.scm -- wave of 2026-11-09/10 (agent BA-ONB), wired 2026-11-10 from scratchpad/baonb-1109/hilbert-onb-vocab.scm. Library: analysis.}\par
\par\medskip\noindent\textbf{Definition} (\texttt{is-orthonormal-set}).\;For all $v$ and $\mathit{os}$, 
\begin{equation*}
\begin{array}{@{}l@{}} \operatorname{is\text{-}orthonormal\text{-}set}(v, \mathit{os}) \\ \Leftrightarrow\; \mathit{os} \subseteq \operatorname{vec}(v) \wedge \\ \forall \operatorname{osa\_} \in \mathit{os}.\; (\langle\cdot,\cdot\rangle_{v}(\operatorname{osa\_}, \operatorname{osa\_}) = 1) \wedge \\ \forall \operatorname{osa\_}, \operatorname{osb\_} \in \mathit{os}. \\ \quad \neg (\operatorname{osa\_} = \operatorname{osb\_}) \Rightarrow (\langle\cdot,\cdot\rangle_{v}(\operatorname{osa\_}, \operatorname{osb\_}) = 0) \end{array}
\end{equation*}

\par\medskip\noindent\textbf{Definition} (\texttt{is-orthonormal-family}).\;For all $v$, $\mathit{ix}$ and $e$, 
\begin{equation*}
\begin{array}{@{}l@{}} \operatorname{is\text{-}orthonormal\text{-}family}(v, \mathit{ix}, e) \\ \Leftrightarrow\; e \in (\mathit{ix} \to \operatorname{vec}(v)) \wedge \\ \forall \operatorname{oni\_} \in \mathit{ix}.\; (\langle\cdot,\cdot\rangle_{v}(e(\operatorname{oni\_}), e(\operatorname{oni\_})) = 1) \wedge \\ \forall \operatorname{oni\_}, \operatorname{onj\_} \in \mathit{ix}. \\ \quad \neg (\operatorname{oni\_} = \operatorname{onj\_}) \Rightarrow (\langle\cdot,\cdot\rangle_{v}(e(\operatorname{oni\_}), e(\operatorname{onj\_})) = 0) \end{array}
\end{equation*}

\par\medskip\noindent\textbf{Definition} (\texttt{is-orthonormal-basis}).\;For all $v$, $\mathit{ix}$ and $e$, 
\begin{equation*}
\begin{array}{@{}l@{}} \operatorname{is\text{-}orthonormal\text{-}basis}(v, \mathit{ix}, e) \\ \Leftrightarrow\; \operatorname{is\text{-}orthonormal\text{-}family}(v, \mathit{ix}, e) \wedge \\ \forall \operatorname{onx\_} \in \operatorname{vec}(v). \\ \quad \forall \operatorname{oni\_} \in \mathit{ix}.\; (\langle\cdot,\cdot\rangle_{v}(\operatorname{onx\_}, e(\operatorname{oni\_})) = 0) \Rightarrow (\operatorname{onx\_} = \operatorname{vzero}(v)) \end{array}
\end{equation*}


\vnbbookfile{structure-library/l2-families-vocab}{structure-library/l2-families-vocab}
\noindent{\small l2-families-vocab.scm -- wave of 2026-11-09/10 (agent BA-ONB2), wired 2026-11-10 from scratchpad/baonb2-1109/bo2-l2-vocab.scm. Library: analysis.}\par
\par\medskip\noindent\textbf{Definition} (\texttt{cc-nag}).\;
\begin{equation*}
\begin{array}{@{}l@{}} \operatorname{cc\text{-}nag} \\ =\; {[}\mathbb{C}, (\lambda {[}\operatorname{x\_}, \operatorname{y\_}] \in (\mathbb{C} \times \mathbb{C}).\; +(\operatorname{x\_}, \operatorname{y\_})), 0, (\lambda \operatorname{x\_} \in \mathbb{C}.\; -(\operatorname{x\_})), (\lambda \operatorname{x\_} \in \mathbb{C}.\; \operatorname{magnitude}(\operatorname{x\_}))] \end{array}
\end{equation*}

\par\medskip\noindent\textbf{Definition} (\texttt{cc-family-sum}).\;$\operatorname{cc\text{-}family\text{-}sum}(f) :=$ $\varepsilon(\{\operatorname{csw\_} \in \mathbb{C} : \operatorname{sums\text{-}to}(\operatorname{cc\text{-}nag}, f, \operatorname{csw\_})\})$
\par\medskip\noindent\textbf{Definition} (\texttt{l2-families}).\;$\operatorname{l2\text{-}families}(\mathit{ix}) :=$
\begin{equation*}
\{\operatorname{l2c\_} \in (\mathit{ix} \to \mathbb{C}) : (\operatorname{esum}((\lambda \operatorname{l2k\_} \in \mathit{ix}.\; (\operatorname{magnitude}(\operatorname{l2c\_}(\operatorname{l2k\_})) \cdot \operatorname{magnitude}(\operatorname{l2c\_}(\operatorname{l2k\_}))))) \in \mathbb{R})\}
\end{equation*}

\par\medskip\noindent\textbf{Definition} (\texttt{l2-ip}).\;$\operatorname{l2\text{-}ip}(\mathit{ix}) :=$
\begin{equation*}
\begin{array}{@{}l@{}} (\lambda {[}\operatorname{l2a\_}, \operatorname{l2b\_}] \in (\operatorname{l2\text{-}families}(\mathit{ix}) \times \operatorname{l2\text{-}families}(\mathit{ix})).\; \\ \quad \operatorname{cc\text{-}family\text{-}sum}((\lambda \operatorname{l2k\_} \in \mathit{ix}.\; (\operatorname{l2a\_}(\operatorname{l2k\_}) \cdot \operatorname{conjugate}(\operatorname{l2b\_}(\operatorname{l2k\_})))))) \end{array}
\end{equation*}

\par\medskip\noindent\textbf{Definition} (\texttt{l2-cips}).\;$\operatorname{l2\text{-}cips}(\mathit{ix}) :=$
\begin{equation*}
{[}\operatorname{normed\text{-}field\text{-}as\text{-}commutative\text{-}ring}(\operatorname{cc\text{-}normed\text{-}field}), \operatorname{l2\text{-}families}(\mathit{ix}), \operatorname{l2f\text{-}vadd}(\mathit{ix}), \operatorname{l2f\text{-}vzero}(\mathit{ix}), \operatorname{l2f\text{-}vneg}(\mathit{ix}), \operatorname{l2f\text{-}act}(\mathit{ix}), \operatorname{l2\text{-}ip}(\mathit{ix})]
\end{equation*}

\par\medskip\noindent\textbf{Definition} (\texttt{l2f-vadd}).\;$\operatorname{l2f\text{-}vadd}(\mathit{ix}) :=$
\begin{equation*}
\begin{array}{@{}l@{}} (\lambda {[}\operatorname{l2a\_}, \operatorname{l2b\_}] \in (\operatorname{l2\text{-}families}(\mathit{ix}) \times \operatorname{l2\text{-}families}(\mathit{ix})).\; \\ \quad (\lambda \operatorname{l2k\_} \in \mathit{ix}.\; (\operatorname{l2a\_}(\operatorname{l2k\_}) + \operatorname{l2b\_}(\operatorname{l2k\_})))) \end{array}
\end{equation*}

\par\medskip\noindent\textbf{Definition} (\texttt{l2f-vzero}).\;$\operatorname{l2f\text{-}vzero}(\mathit{ix}) :=$ $(\lambda \operatorname{l2k\_} \in \mathit{ix}.\; 0)$
\par\medskip\noindent\textbf{Definition} (\texttt{l2f-vneg}).\;$\operatorname{l2f\text{-}vneg}(\mathit{ix}) :=$
\begin{equation*}
(\lambda \operatorname{l2a\_} \in \operatorname{l2\text{-}families}(\mathit{ix}).\; (\lambda \operatorname{l2k\_} \in \mathit{ix}.\; -\operatorname{l2a\_}(\operatorname{l2k\_})))
\end{equation*}

\par\medskip\noindent\textbf{Definition} (\texttt{l2f-act}).\;$\operatorname{l2f\text{-}act}(\mathit{ix}) :=$
\begin{equation*}
(\lambda {[}\operatorname{l2z\_}, \operatorname{l2a\_}] \in (\mathbb{C} \times \operatorname{l2\text{-}families}(\mathit{ix})).\; (\lambda \operatorname{l2k\_} \in \mathit{ix}.\; (\operatorname{l2z\_} \cdot \operatorname{l2a\_}(\operatorname{l2k\_}))))
\end{equation*}

\vnbdef{cc-nag@inv}{\vnbstmt{(\operatorname{inv}(\operatorname{cc\text{-}nag}) \simeq (\lambda x \in \mathbb{C}.\; -x))}}
\vnbdef{cc-nag@carr}{\vnbstmt{(\operatorname{carr}(\operatorname{cc\text{-}nag}) \simeq \mathbb{C})}}
\vnbdef{cc-nag@iden}{\vnbstmt{(\operatorname{iden}(\operatorname{cc\text{-}nag}) \simeq 0)}}
\vnbdef{cc-nag@nrm}{\vnbstmt{(\operatorname{nrm}(\operatorname{cc\text{-}nag}) \simeq (\lambda x \in \mathbb{C}.\; \operatorname{magnitude}(x)))}}
\vnbdef{cc-nag@opr}{\vnbstmt{(\operatorname{opr}(\operatorname{cc\text{-}nag}) \simeq (\lambda {[}\operatorname{x\_}, \operatorname{y\_}] \in (\mathbb{C} \times \mathbb{C}).\; (\operatorname{x\_} + \operatorname{y\_})))}}

\vnbbookfile{structure-library/ccr4c-subnf-vocab}{structure-library/ccr4c-subnf-vocab}
\noindent{\small ccr4c-subnf-vocab.scm -- wave of 2026-11-09/10 (agent CCR-1109-C4), wired 2026-11-10 from scratchpad/ccr1109c4/ccr4c-subnf-vocab.scm. Library: analysis.}\par
\par\medskip\noindent\textbf{Definition} (\texttt{sub-nf}).\;$\operatorname{sub\text{-}nf}(\operatorname{snk\_}, \operatorname{sns\_}) :=$
\begin{equation*}
{[}\operatorname{sns\_}, (\lambda {[}\operatorname{snu\_}, \operatorname{snv\_}] \in (\operatorname{sns\_} \times \operatorname{sns\_}).\; \operatorname{add}(\operatorname{snk\_})(\operatorname{snu\_}, \operatorname{snv\_})), (\lambda {[}\operatorname{snu\_}, \operatorname{snv\_}] \in (\operatorname{sns\_} \times \operatorname{sns\_}).\; \operatorname{mul}(\operatorname{snk\_})(\operatorname{snu\_}, \operatorname{snv\_})), (\lambda \operatorname{snu\_} \in \operatorname{sns\_}.\; \operatorname{neg}(\operatorname{snk\_})(\operatorname{snu\_})), \operatorname{zero}(\operatorname{snk\_}), \operatorname{one}(\operatorname{snk\_}), (\lambda \operatorname{snu\_} \in \operatorname{sns\_}.\; \operatorname{fnrm}(\operatorname{snk\_})(\operatorname{snu\_}))]
\end{equation*}


\vnbbookfile{structure-library/ccr4d-ext-vocab}{structure-library/ccr4d-ext-vocab}
\noindent{\small ccr4d-ext-vocab.scm -- wave of 2026-11-09/10 (agent CCR-1109-C5), wired 2026-11-10 from scratchpad/ccr1109c5/ccr4d-ext-vocab.scm. Library: analysis.}\par
\par\medskip\noindent\textbf{Definition} (\texttt{nf-subfield-module}).\;$\operatorname{nf\text{-}subfield\text{-}module}(\operatorname{fmk\_}, \operatorname{fms\_}) :=$
\begin{equation*}
{[}\operatorname{normed\text{-}field\text{-}as\text{-}commutative\text{-}ring}(\operatorname{sub\text{-}nf}(\operatorname{fmk\_}, \operatorname{fms\_})), \operatorname{carr}(\operatorname{normed\text{-}field\text{-}as\text{-}commutative\text{-}ring}(\operatorname{fmk\_})), \operatorname{add}(\operatorname{normed\text{-}field\text{-}as\text{-}commutative\text{-}ring}(\operatorname{fmk\_})), \operatorname{zero}(\operatorname{normed\text{-}field\text{-}as\text{-}commutative\text{-}ring}(\operatorname{fmk\_})), \operatorname{neg}(\operatorname{normed\text{-}field\text{-}as\text{-}commutative\text{-}ring}(\operatorname{fmk\_})), (\lambda {[}\operatorname{fmr\_}, \operatorname{fmx\_}] \in (\operatorname{fms\_} \times \operatorname{carr}(\operatorname{normed\text{-}field\text{-}as\text{-}commutative\text{-}ring}(\operatorname{fmk\_}))).\; \operatorname{mul}(\operatorname{normed\text{-}field\text{-}as\text{-}commutative\text{-}ring}(\operatorname{fmk\_}))(\operatorname{fmr\_}, \operatorname{fmx\_}))]
\end{equation*}


\vnbbookfile{structure-library/is-set-laws}{structure-library/is-set-laws}
\noindent{\small A file shared across books; the entries that belong to this one.}\par
\noindent{\small structure-library/is-set-laws.scm -- THE GENERATED SETHOOD THEOREMS (W1I, 2026-11-06).}\par
\vnbprop{is-metric-space-is-set}{\vnbstmt{\forall s.\; (\operatorname{is\text{-}metric\text{-}space}(s) \Rightarrow (s \in \mathbf{Set}))}}
\vnbprop{is-top-space-is-set}{\vnbstmt{\forall s.\; (\operatorname{is\text{-}top\text{-}space}(s) \Rightarrow (s \in \mathbf{Set}))}}
\vnbprop{is-metrizable-top-space-is-set}{\vnbstmt{\forall s.\; (\operatorname{is\text{-}metrizable\text{-}top\text{-}space}(s) \Rightarrow (s \in \mathbf{Set}))}}
\vnbprop{is-normed-field-is-set}{\vnbstmt{\forall s.\; (\operatorname{is\text{-}normed\text{-}field}(s) \Rightarrow (s \in \mathbf{Set}))}}
\vnbprop{is-normed-ag-is-set}{\vnbstmt{\forall s.\; (\operatorname{is\text{-}normed\text{-}ag}(s) \Rightarrow (s \in \mathbf{Set}))}}
\vnbprop{is-normed-vector-space-is-set}{\vnbstmt{\forall s.\; (\operatorname{is\text{-}normed\text{-}vector\text{-}space}(s) \Rightarrow (s \in \mathbf{Set}))}}
\vnbprop{is-complex-inner-product-space-is-set}{\vnbstmt{\forall s.\; (\operatorname{is\text{-}complex\text{-}inner\text{-}product\text{-}space}(s) \Rightarrow (s \in \mathbf{Set}))}}
\vnbprop{is-pseudometric-space-is-set}{\vnbstmt{\forall s.\; (\operatorname{is\text{-}pseudometric\text{-}space}(s) \Rightarrow (s \in \mathbf{Set}))}}
\vnbprop{is-c-metric-space-is-set}{\vnbstmt{\forall s.\; (\operatorname{is\text{-}c\text{-}metric\text{-}space}(s) \Rightarrow (s \in \mathbf{Set}))}}

\vnbbookfile{theorem-library/subseq-apply}{theorem-library/subseq-apply}
\noindent{\small subseq-apply.scm -- the VALUE of a subsequence at one TYPED index.}\par
\vnbprop{subseq-apply}{\vnbstmt{\forall f, i.; \forall k \in \mathbb{N}.\; (\operatorname{subseq}(f, i)(k) \simeq f(i(k)))}}

\vnbbookfile{theorem-library/bdd-metric-carrier}{theorem-library/bdd-metric-carrier}
\noindent{\small bdd-metric-carrier.scm -- BDD-METRIC keeps the point set, proven.}\par
\vnbprop{bdd-metric-carrier}{\vnbstmt{\forall s.\; (\operatorname{is\text{-}metric\text{-}space}(s) \Rightarrow (\operatorname{pts}(\operatorname{bdd\text{-}metric}(s)) \simeq \operatorname{pts}(s)))}}

\vnbbookfile{theorem-library/normed-field-ring-view}{theorem-library/normed-field-ring-view}
\noindent{\small normed-field-ring-view.scm -- THE SCALAR BRIDGE, PROVEN.}\par
\vnbprop{normed-field-ring-view-carr}{\vnbstmt{\begin{array}{@{}l@{}} \forall f. \\ \quad \operatorname{is\text{-}normed\text{-}field}(f) \Rightarrow (\operatorname{carr}(\operatorname{normed\text{-}field\text{-}as\text{-}commutative\text{-}ring}(f)) \simeq \operatorname{carr}(f)) \end{array}}}
\vnbprop{normed-field-ring-view-add}{\vnbstmt{\begin{array}{@{}l@{}} \forall f. \\ \quad \operatorname{is\text{-}normed\text{-}field}(f) \Rightarrow (\operatorname{add}(\operatorname{normed\text{-}field\text{-}as\text{-}commutative\text{-}ring}(f)) \simeq \operatorname{add}(f)) \end{array}}}
\vnbprop{normed-field-ring-view-mul}{\vnbstmt{\begin{array}{@{}l@{}} \forall f. \\ \quad \operatorname{is\text{-}normed\text{-}field}(f) \Rightarrow (\operatorname{mul}(\operatorname{normed\text{-}field\text{-}as\text{-}commutative\text{-}ring}(f)) \simeq \operatorname{mul}(f)) \end{array}}}
\vnbprop{normed-field-ring-view-neg}{\vnbstmt{\begin{array}{@{}l@{}} \forall f. \\ \quad \operatorname{is\text{-}normed\text{-}field}(f) \Rightarrow (\operatorname{neg}(\operatorname{normed\text{-}field\text{-}as\text{-}commutative\text{-}ring}(f)) \simeq \operatorname{neg}(f)) \end{array}}}
\vnbprop{normed-field-ring-view-zero}{\vnbstmt{\begin{array}{@{}l@{}} \forall f. \\ \quad \operatorname{is\text{-}normed\text{-}field}(f) \Rightarrow (\operatorname{zero}(\operatorname{normed\text{-}field\text{-}as\text{-}commutative\text{-}ring}(f)) \simeq \operatorname{zero}(f)) \end{array}}}
\vnbprop{normed-field-ring-view-one}{\vnbstmt{\begin{array}{@{}l@{}} \forall f. \\ \quad \operatorname{is\text{-}normed\text{-}field}(f) \Rightarrow (\operatorname{one}(\operatorname{normed\text{-}field\text{-}as\text{-}commutative\text{-}ring}(f)) \simeq \operatorname{one}(f)) \end{array}}}

\vnbbookfile{theorem-library/view-typing-laws}{theorem-library/view-typing-laws}
\noindent{\small A file shared across books; the entries that belong to this one.}\par
\noindent{\small theorem-library/view-typing-laws.scm -- THE VIEW TYPING THEOREMS, PROVEN (D8, 2026-11-15).}\par
\vnbprop{normed-field-additive-ag-is-abelian-group}{\vnbstmt{\begin{array}{@{}l@{}} \forall r. \\ \quad \operatorname{is\text{-}normed\text{-}field}(r) \Rightarrow \operatorname{is\text{-}abelian\text{-}group}(\operatorname{normed\text{-}field\text{-}additive\text{-}ag}(r)) \end{array}}}
\vnbprop{normed-ag-as-abelian-group-is-abelian-group}{\vnbstmt{\begin{array}{@{}l@{}} \forall r. \\ \quad \operatorname{is\text{-}normed\text{-}ag}(r) \Rightarrow \operatorname{is\text{-}abelian\text{-}group}(\operatorname{normed\text{-}ag\text{-}as\text{-}abelian\text{-}group}(r)) \end{array}}}
\vnbprop{normed-field-as-commutative-ring-is-commutative-ring}{\vnbstmt{\begin{array}{@{}l@{}} \forall r. \\ \quad \operatorname{is\text{-}normed\text{-}field}(r) \Rightarrow  \\ \quad \quad \operatorname{is\text{-}commutative\text{-}ring}(\operatorname{normed\text{-}field\text{-}as\text{-}commutative\text{-}ring}(r)) \end{array}}}
\vnbprop{normed-vector-space-as-module-is-module}{\vnbstmt{\begin{array}{@{}l@{}} \forall r. \\ \quad \operatorname{is\text{-}normed\text{-}vector\text{-}space}(r) \Rightarrow \operatorname{is\text{-}module}(\operatorname{normed\text{-}vector\text{-}space\text{-}as\text{-}module}(r)) \end{array}}}
\vnbprop{complex-inner-product-space-as-module-is-module}{\vnbstmt{\begin{array}{@{}l@{}} \forall r. \\ \quad \operatorname{is\text{-}complex\text{-}inner\text{-}product\text{-}space}(r) \Rightarrow  \\ \quad \quad \operatorname{is\text{-}module}(\operatorname{complex\text{-}inner\text{-}product\text{-}space\text{-}as\text{-}module}(r)) \end{array}}}

\vnbbookfile{theorem-library/rake-setoid}{theorem-library/rake-setoid}
\noindent{\small A file shared across books; the entries that belong to this one.}\par
\noindent{\small rake-setoid.scm -- BATCH F of the 2026-09-17 rake: setoids and carriers. Seven asserted leaves of the PSS, PROVEN.  Every statement here is its support's statement UNCHANGED.}\par
\vnbprop{vec-is-set}{\vnbstmt{\forall m.\; (\operatorname{is\text{-}normed\text{-}vector\text{-}space}(m) \Rightarrow (\operatorname{vec}(m) \in \mathbf{Set}))}}

\vnbbookfile{theorem-library/rake-analysis2}{theorem-library/rake-analysis2}
\noindent{\small rake-analysis2.scm -- rake batch L (2026-09-17): five asserted leaves of the analysis / sets area, plus the general lemma that unblocked one of them.}\par
\vnbprop{subseq-is-fun}{\vnbstmt{\begin{array}{@{}l@{}} \forall s, f, i. \\ \quad ((f \in (\mathbb{N} \to \operatorname{pts}(s))) \wedge \operatorname{strictly\text{-}mono\text{-}nn}(i)) \Rightarrow  \\ \quad \quad \operatorname{subseq}(f, i) \in (\mathbb{N} \to \operatorname{pts}(s)) \end{array}}}
\vnbprop{metric-hom-is-continuous}{\vnbstmt{\forall s, t, f.\; (\operatorname{is\text{-}hom\text{-}metric\text{-}space}(s, t, f) \Rightarrow \operatorname{is\text{-}continuous}(s, t, f))}}
\vnbprop{centre-set-contains-choice}{\vnbstmt{\begin{array}{@{}l@{}} \forall s, r, f.; \forall u \in f. \\ \quad \exists c.\; (c \in \operatorname{centres}(s, u, r)) \Rightarrow  \\ \quad \quad \varepsilon(\operatorname{centres}(s, u, r)) \in \operatorname{centre\text{-}set}(s, r, f) \end{array}}}

\vnbbookfile{theorem-library/rr-star-order-minimal}{theorem-library/rr-star-order-minimal}
\noindent{\small rr-star-order-minimal.scm -- THE MINIMAL PRIMITIVE SET of the extended line, run. Agent EXT-ORDER, 2026-10-26.  Wired 2026-10-27 (batch 0): the two stamps retired from structure-library/extended-reals.scm, the three order laws below never stamped. Library: analysis (extended-reals.scm's library).}\par
\vnbprop{pos-inf-above-reals}{\vnbstmt{\forall x \in \mathbb{R}.\; \neg (\operatorname{pos\text{-}inf} \leq x)}}
\vnbprop{neg-inf-not-in-rr}{\vnbstmt{\neg (\operatorname{neg\text{-}inf} \in \mathbb{R})}}
\vnbprop{neg-inf-below-reals}{\vnbstmt{\forall x \in \mathbb{R}.\; \neg (x \leq \operatorname{neg\text{-}inf})}}
\vnbprop{rr-star-le-pos-inf-eq}{\vnbstmt{\forall x \in \operatorname{rr\text{-}star}.\; ((\operatorname{pos\text{-}inf} \leq x) \Rightarrow (x = \operatorname{pos\text{-}inf}))}}
\vnbprop{rr-star-le-neg-inf-eq}{\vnbstmt{\forall x \in \operatorname{rr\text{-}star}.\; ((x \leq \operatorname{neg\text{-}inf}) \Rightarrow (x = \operatorname{neg\text{-}inf}))}}
\vnbprop{rr-star-leq-reflexive}{\vnbstmt{\forall a \in \operatorname{rr\text{-}star}.\; (a \leq a)}}
\vnbprop{rr-star-leq-transitive}{\vnbstmt{\begin{array}{@{}l@{}} \forall a, b, c. \\ \quad (((a \in \operatorname{rr\text{-}star}) \wedge ((b \in \operatorname{rr\text{-}star}) \wedge (c \in \operatorname{rr\text{-}star}))) \wedge ((a \leq b) \wedge (b \leq c))) \Rightarrow  \\ \quad \quad a \leq c \end{array}}}
\vnbprop{rr-star-leq-total}{\vnbstmt{\forall a, b.\; (((a \in \operatorname{rr\text{-}star}) \wedge (b \in \operatorname{rr\text{-}star})) \Rightarrow ((a \leq b) \vee (b \leq a)))}}

\vnbbookfile{theorem-library/rake-infinity-points}{theorem-library/rake-infinity-points}
\noindent{\small theorem-library/rake-infinity-points.scm -- the five DERIVABLE axioms of structure-library/extended-reals.scm, proven.  Rake "infinity points".}\par
\vnbprop{rr-subset-rr-star}{\vnbstmt{\mathbb{R} \subseteq \operatorname{rr\text{-}star}}}
\vnbprop{pos-inf-in-rr-star}{\vnbstmt{\operatorname{pos\text{-}inf} \in \operatorname{rr\text{-}star}}}
\vnbprop{neg-inf-in-rr-star}{\vnbstmt{\operatorname{neg\text{-}inf} \in \operatorname{rr\text{-}star}}}
\vnbprop{pos-inf-not-in-rr}{\vnbstmt{\neg (\operatorname{pos\text{-}inf} \in \mathbb{R})}}
\vnbprop{pos-inf-neq-neg-inf}{\vnbstmt{\neg (\operatorname{pos\text{-}inf} = \operatorname{neg\text{-}inf})}}

\vnbbookfile{theorem-library/extinf-points}{theorem-library/extinf-points}
\noindent{\small extinf-points.scm -- the NEG-INF mirrors of the theorems proven from pos-inf-above-reals about the two points at infinity.  Agent EXT-INF, 2026-10-26. DESTINATION: theorem-library/extinf-points.scm.  Library: analysis (the library of structure-library/extended-reals, whose vocabulary every statement here uses).}\par
\vnbprop{neg-inf-neq-pos-inf}{\vnbstmt{\neg (\operatorname{neg\text{-}inf} = \operatorname{pos\text{-}inf})}}
\vnbprop{neg-inf-lt-rr}{\vnbstmt{\forall x \in \mathbb{R}.\; (\operatorname{neg\text{-}inf} < x)}}

\vnbbookfile{theorem-library/rr-star-order-thms}{theorem-library/rr-star-order-thms}
\noindent{\small rr-star-order-thms.scm -- consequences of the order axioms of RR-STAR (rr-star-order.scm). Agent EXT-ORDER, 2026-10-26.  DESTINATION: theorem-library/rr-star-order-thms.scm. Library: analysis.}\par
\vnbprop{pos-inf-not-below-neg-inf}{\vnbstmt{\neg (\operatorname{pos\text{-}inf} \leq \operatorname{neg\text{-}inf})}}
\vnbprop{rr-leq-reflexive-via-star}{\vnbstmt{\forall a \in \mathbb{R}.\; (a \leq a)}}
\vnbprop{rr-leq-antisymmetric-via-star}{\vnbstmt{\forall a, b.\; ((((a \in \mathbb{R}) \wedge (b \in \mathbb{R})) \wedge ((a \leq b) \wedge (b \leq a))) \Rightarrow (a = b))}}
\vnbprop{rr-leq-transitive-via-star}{\vnbstmt{\begin{array}{@{}l@{}} \forall a, b, c. \\ \quad (((a \in \mathbb{R}) \wedge ((b \in \mathbb{R}) \wedge (c \in \mathbb{R}))) \wedge ((a \leq b) \wedge (b \leq c))) \Rightarrow (a \leq c) \end{array}}}
\vnbprop{rr-leq-total-via-star}{\vnbstmt{\forall a, b.\; (((a \in \mathbb{R}) \wedge (b \in \mathbb{R})) \Rightarrow ((a \leq b) \vee (b \leq a)))}}
\vnbprop{rr-star-leq-antisymmetric-unguarded}{\vnbstmt{\forall a, b.\; (((a \leq b) \wedge (b \leq a)) \Rightarrow (a = b))}}
\vnbprop{rr-star-leq-transitive-unguarded}{\vnbstmt{\forall a, b, c.\; (((a \leq b) \wedge (b \leq c)) \Rightarrow (a \leq c))}}

\vnbbookfile{theorem-library/rake-rr-pos-star}{theorem-library/rake-rr-pos-star}
\noindent{\small theorem-library/rake-rr-pos-star.scm -- the six CARRIER facts of RR-POS-STAR = [0, +inf], proven.  Rake batch 5c-K.}\par
\vnbprop{zero-in-rr-pos-star}{\vnbstmt{0 \in \operatorname{rr\text{-}pos\text{-}star}}}
\vnbprop{pos-inf-in-rr-pos-star}{\vnbstmt{\operatorname{pos\text{-}inf} \in \operatorname{rr\text{-}pos\text{-}star}}}
\vnbprop{neg-inf-not-in-rr-pos-star}{\vnbstmt{\neg (\operatorname{neg\text{-}inf} \in \operatorname{rr\text{-}pos\text{-}star})}}
\vnbprop{rr-pos-star-nonneg}{\vnbstmt{\forall x \in \operatorname{rr\text{-}pos\text{-}star}.\; (0 \leq x)}}
\vnbprop{rr-pos-star-below-pos-inf}{\vnbstmt{\forall x \in \operatorname{rr\text{-}pos\text{-}star}.\; (x \leq \operatorname{pos\text{-}inf})}}
\vnbprop{rr-pos-star-subset-rr-star}{\vnbstmt{\operatorname{rr\text{-}pos\text{-}star} \subseteq \operatorname{rr\text{-}star}}}

\vnbbookfile{theorem-library/rake-eplus-defined}{theorem-library/rake-eplus-defined}
\noindent{\small theorem-library/rake-eplus-defined.scm -- extended addition on [0,+inf] DEFINED, and its laws proven.  Rake batch 5c-V.}\par
\vnbprop{apply-congruence-2}{\vnbstmt{\forall f, g, u, v.\; ((f \simeq g) \Rightarrow (f(u, v) \simeq g(u, v)))}}
\vnbprop{rr-pos-star-is-set}{\vnbstmt{\operatorname{rr\text{-}pos\text{-}star} \in \mathbf{Set}}}
\vnbprop{eplus-in-fun}{\vnbstmt{\mathit{eplus} \in ((\operatorname{rr\text{-}pos\text{-}star} \times \operatorname{rr\text{-}pos\text{-}star}) \to \operatorname{rr\text{-}pos\text{-}star})}}
\vnbprop{eplus-pos-inf-left}{\vnbstmt{\forall y \in \operatorname{rr\text{-}pos\text{-}star}.\; (\operatorname{eplus}(\operatorname{pos\text{-}inf}, y) = \operatorname{pos\text{-}inf})}}
\vnbprop{eplus-pos-inf-right}{\vnbstmt{\forall x \in \operatorname{rr\text{-}pos\text{-}star}.\; (\operatorname{eplus}(x, \operatorname{pos\text{-}inf}) = \operatorname{pos\text{-}inf})}}
\vnbprop{eplus-real-defined}{\vnbstmt{\begin{array}{@{}l@{}} \forall x, y. \\ \quad ((x \in \mathbb{R}) \wedge ((0 \leq x) \wedge ((y \in \mathbb{R}) \wedge (0 \leq y)))) \Rightarrow (\operatorname{eplus}(x, y) = (x + y)) \end{array}}}

\vnbbookfile{theorem-library/rake-etimes-defined}{theorem-library/rake-etimes-defined}
\noindent{\small theorem-library/rake-etimes-defined.scm -- extended MULTIPLICATION on [0,+inf] DEFINED, and its laws proven.  Rake batch 5c-Y.}\par
\vnbprop{etimes-in-fun}{\vnbstmt{\mathit{etimes} \in ((\operatorname{rr\text{-}pos\text{-}star} \times \operatorname{rr\text{-}pos\text{-}star}) \to \operatorname{rr\text{-}pos\text{-}star})}}
\vnbprop{etimes-zero-pos-inf}{\vnbstmt{\operatorname{etimes}(0, \operatorname{pos\text{-}inf}) = 0}}
\vnbprop{etimes-pos-inf-zero}{\vnbstmt{\operatorname{etimes}(\operatorname{pos\text{-}inf}, 0) = 0}}
\vnbprop{etimes-pos-inf-left}{\vnbstmt{\forall y \in \operatorname{rr\text{-}pos\text{-}star}.\; (\neg (y = 0) \Rightarrow (\operatorname{etimes}(\operatorname{pos\text{-}inf}, y) = \operatorname{pos\text{-}inf}))}}
\vnbprop{etimes-pos-inf-right}{\vnbstmt{\forall x \in \operatorname{rr\text{-}pos\text{-}star}.\; (\neg (x = 0) \Rightarrow (\operatorname{etimes}(x, \operatorname{pos\text{-}inf}) = \operatorname{pos\text{-}inf}))}}
\vnbprop{etimes-real-defined}{\vnbstmt{\begin{array}{@{}l@{}} \forall x, y. \\ \quad ((x \in \mathbb{R}) \wedge ((0 \leq x) \wedge ((y \in \mathbb{R}) \wedge (0 \leq y)))) \Rightarrow (\operatorname{etimes}(x, y) = (x \cdot y)) \end{array}}}

\vnbbookfile{theorem-library/discrete-space}{theorem-library/discrete-space}
\noindent{\small discrete-space.scm -- THE FIRST WITNESS OF A MEASURABLE SPACE, and the first outright inhabitant of TOP-SPACE.}\par
\par\medskip\noindent\textbf{Definition} (\texttt{discrete-space}).\;$\operatorname{discrete\text{-}space}(\operatorname{a\_}) :=$ ${[}\operatorname{a\_}, \operatorname{power}(\operatorname{a\_})]$
\vnbprop{ds-pts}{\vnbstmt{\forall a \in \mathbf{Set}.\; (\operatorname{pts}(\operatorname{discrete\text{-}space}(a)) \simeq a)}}
\vnbprop{ds-opens}{\vnbstmt{\forall a \in \mathbf{Set}.\; (\operatorname{opens}(\operatorname{discrete\text{-}space}(a)) \simeq \operatorname{power}(a))}}
\vnbprop{ds-sigma}{\vnbstmt{\forall a \in \mathbf{Set}.\; (\sigma(\operatorname{discrete\text{-}space}(a)) \simeq \operatorname{power}(a))}}
\vnbprop{discrete-space-is-top-space}{\vnbstmt{\forall a \in \mathbf{Set}.\; \operatorname{is\text{-}top\text{-}space}(\operatorname{discrete\text{-}space}(a))}}

\vnbbookfile{theorem-library/op-typing}{theorem-library/op-typing}
\noindent{\small op-typing.scm -- the codomain typing of a structure operation, in APPLIED form, PROVEN.  One driver, seven theorems.}\par
\vnbprop{metric-dist-real}{\vnbstmt{\begin{array}{@{}l@{}} \forall s.; \forall x, y \in \operatorname{pts}(s). \\ \quad \operatorname{is\text{-}metric\text{-}space}(s) \Rightarrow (\operatorname{dist}(s)(x, y) \in \mathbb{R}) \end{array}}}
\vnbprop{vnrm-real}{\vnbstmt{\begin{array}{@{}l@{}} \forall m, w. \\ \quad (\operatorname{is\text{-}normed\text{-}vector\text{-}space}(m) \wedge (w \in \operatorname{vec}(m))) \Rightarrow (\operatorname{vnrm}(m)(w) \in \mathbb{R}) \end{array}}}
\vnbprop{nvs-vadd-in-vec}{\vnbstmt{\begin{array}{@{}l@{}} \forall m.; \forall x, y \in \operatorname{vec}(m). \\ \quad \operatorname{is\text{-}normed\text{-}vector\text{-}space}(m) \Rightarrow (\operatorname{vadd}(m)(x, y) \in \operatorname{vec}(m)) \end{array}}}

\vnbbookfile{theorem-library/ba1i-dominated}{theorem-library/ba1i-dominated}
\noindent{\small ba1i-dominated.scm -- wave of 2026-11-09/10 (agent BA-1I), wired 2026-11-10 from scratchpad/ba1i-1109/ba1i-dominated.scm. Library: analysis.}\par
\vnbprop{complete-of-dominated-metric}{\vnbstmt{\begin{array}{@{}l@{}} \forall s, t. \\ \quad \operatorname{is\text{-}complete}(s) \wedge \\ \quad \operatorname{is\text{-}metric\text{-}space}(t) \wedge \\ \quad \operatorname{pts}(t) = \operatorname{pts}(s) \wedge \\ \quad \forall p, q \in \operatorname{pts}(t).\; (\operatorname{dist}(s)(p, q) \leq \operatorname{dist}(t)(p, q)) \wedge \\ \quad \forall p, q \in \operatorname{pts}(t).\; (\operatorname{dist}(t)(p, q) \leq (2 \cdot \operatorname{dist}(s)(p, q))) \Rightarrow \\ \quad \quad \operatorname{is\text{-}complete}(t) \end{array}}}

\vnbbookfile{theorem-library/rr-is-normed-field}{theorem-library/rr-is-normed-field}
\noindent{\small rr-is-normed-field.scm -- IS-NORMED-FIELD(RR-NORMED-FIELD), PROVED.}\par
\vnbprop{rr-is-normed-field}{\vnbstmt{\operatorname{is\text{-}normed\text{-}field}(\operatorname{rr\text{-}normed\text{-}field})}}
\vnbprop{rr-scalar-ring-carr}{\vnbstmt{(\operatorname{carr}(\operatorname{normed\text{-}field\text{-}as\text{-}commutative\text{-}ring}(\operatorname{rr\text{-}normed\text{-}field})) \simeq \mathbb{R})}}
\vnbprop{rr-scalar-ring-add}{\vnbstmt{(\operatorname{add}(\operatorname{normed\text{-}field\text{-}as\text{-}commutative\text{-}ring}(\operatorname{rr\text{-}normed\text{-}field})) \simeq (\lambda {[}\operatorname{x\_}, \operatorname{y\_}] \in (\mathbb{R} \times \mathbb{R}).\; (\operatorname{x\_} + \operatorname{y\_})))}}
\vnbprop{rr-scalar-ring-mul}{\vnbstmt{(\operatorname{mul}(\operatorname{normed\text{-}field\text{-}as\text{-}commutative\text{-}ring}(\operatorname{rr\text{-}normed\text{-}field})) \simeq (\lambda {[}\operatorname{x\_}, \operatorname{y\_}] \in (\mathbb{R} \times \mathbb{R}).\; (\operatorname{x\_} \cdot \operatorname{y\_})))}}
\vnbprop{rr-scalar-ring-neg}{\vnbstmt{(\operatorname{neg}(\operatorname{normed\text{-}field\text{-}as\text{-}commutative\text{-}ring}(\operatorname{rr\text{-}normed\text{-}field})) \simeq (\lambda x \in \mathbb{R}.\; -x))}}
\vnbprop{rr-scalar-ring-zero}{\vnbstmt{(\operatorname{zero}(\operatorname{normed\text{-}field\text{-}as\text{-}commutative\text{-}ring}(\operatorname{rr\text{-}normed\text{-}field})) \simeq 0)}}
\vnbprop{rr-scalar-ring-one}{\vnbstmt{(\operatorname{one}(\operatorname{normed\text{-}field\text{-}as\text{-}commutative\text{-}ring}(\operatorname{rr\text{-}normed\text{-}field})) \simeq 1)}}

\vnbbookfile{theorem-library/dual-space-readoffs}{theorem-library/dual-space-readoffs}
\noindent{\small dual-stamps.scm -- the four definitional read-offs of structure-library/dual-space.scm, PROVEN. Agent DUAL-FIX, 2026-10-26.  DESTINATION: theorem-library/dual-space-readoffs.scm.  Library: analysis. Load window: right after structure-library/dual-space (and the files that define its heads); no theorem citation except ...}\par
\vnbprop{dual-vec-membership-proven}{\vnbstmt{\begin{array}{@{}l@{}} \forall m, f. \\ \quad (f \in \operatorname{dual\text{-}vec}(m)) \Leftrightarrow ((f \in (\operatorname{vec}(m) \to \mathbb{R})) \wedge \operatorname{is\text{-}bounded\text{-}linear\text{-}functional}(m, f)) \end{array}}}
\vnbprop{dual-scal-proven}{\vnbstmt{\begin{array}{@{}l@{}} \forall m. \\ \quad (\operatorname{dual\text{-}vec}(m) \in \mathbf{Set}) \Rightarrow  \\ \quad \quad (\operatorname{scal}({m}^{*}) \simeq \operatorname{normed\text{-}field\text{-}as\text{-}commutative\text{-}ring}(\operatorname{rr\text{-}normed\text{-}field})) \end{array}}}
\vnbprop{dual-carrier-proven}{\vnbstmt{\forall m.\; ((\operatorname{dual\text{-}vec}(m) \in \mathbf{Set}) \Rightarrow (\operatorname{vec}({m}^{*}) \simeq \operatorname{dual\text{-}vec}(m)))}}
\vnbprop{dual-norm-value-proven}{\vnbstmt{\begin{array}{@{}l@{}} \forall m.; \forall f \in \operatorname{dual\text{-}vec}(m). \\ \quad (\operatorname{dual\text{-}vec}(m) \in \mathbf{Set}) \Rightarrow (\operatorname{vnrm}({m}^{*})(f) \simeq \operatorname{dual\text{-}norm}(m, f)) \end{array}}}

\vnbbookfile{theorem-library/rr-nvs-exemplification}{theorem-library/rr-nvs-exemplification}
\noindent{\small rr-nvs-exemplification.scm -- THE FIRST WITNESS OF A NORMED VECTOR SPACE.}\par
\par\medskip\noindent\textbf{Definition} (\texttt{rr-nvs}).\;
\begin{equation*}
\begin{array}{@{}l@{}} \operatorname{rr\text{-}nvs} \\ =\; {[}\operatorname{normed\text{-}field\text{-}as\text{-}commutative\text{-}ring}(\operatorname{rr\text{-}normed\text{-}field}), \mathbb{R}, (\lambda {[}\operatorname{x\_}, \operatorname{y\_}] \in (\mathbb{R} \times \mathbb{R}).\; +(\operatorname{x\_}, \operatorname{y\_})), 0, (\lambda \operatorname{x\_} \in \mathbb{R}.\; -(\operatorname{x\_})), (\lambda {[}\operatorname{r\_}, \operatorname{x\_}] \in (\mathbb{R} \times \mathbb{R}).\; \cdot(\operatorname{r\_}, \operatorname{x\_})), (\lambda \operatorname{x\_} \in \mathbb{R}.\; \lvert \operatorname{x\_} \rvert)] \end{array}
\end{equation*}

\vnbdef{rr-nvs@vzero}{\vnbstmt{(\operatorname{vzero}(\operatorname{rr\text{-}nvs}) \simeq 0)}}
\vnbdef{rr-nvs@scal}{\vnbstmt{(\operatorname{scal}(\operatorname{rr\text{-}nvs}) \simeq \operatorname{normed\text{-}field\text{-}as\text{-}commutative\text{-}ring}(\operatorname{rr\text{-}normed\text{-}field}))}}
\vnbdef{rr-nvs@vadd}{\vnbstmt{(\operatorname{vadd}(\operatorname{rr\text{-}nvs}) \simeq (\lambda {[}\operatorname{x\_}, \operatorname{y\_}] \in (\mathbb{R} \times \mathbb{R}).\; (\operatorname{x\_} + \operatorname{y\_})))}}
\vnbdef{rr-nvs@vec}{\vnbstmt{(\operatorname{vec}(\operatorname{rr\text{-}nvs}) \simeq \mathbb{R})}}
\vnbdef{rr-nvs@vneg}{\vnbstmt{(\operatorname{vneg}(\operatorname{rr\text{-}nvs}) \simeq (\lambda x \in \mathbb{R}.\; -x))}}
\vnbdef{rr-nvs@vnrm}{\vnbstmt{(\operatorname{vnrm}(\operatorname{rr\text{-}nvs}) \simeq (\lambda x \in \mathbb{R}.\; \lvert x \rvert))}}
\vnbdef{rr-nvs@act}{\vnbstmt{(\operatorname{act}(\operatorname{rr\text{-}nvs}) \simeq (\lambda {[}\operatorname{r\_}, \operatorname{x\_}] \in (\mathbb{R} \times \mathbb{R}).\; (\operatorname{r\_} \cdot \operatorname{x\_})))}}
\vnbprop{rr-nvs-is-normed-vector-space}{\vnbstmt{\operatorname{is\text{-}normed\text{-}vector\text{-}space}(\operatorname{rr\text{-}nvs})}}

\vnbbookfile{theorem-library/rake-algebra}{theorem-library/rake-algebra}
\noindent{\small A file shared across books; the entries that belong to this one.}\par
\noindent{\small rake-algebra.scm -- BATCH D of the 2026-09-17 rake: seven asserted leaves of the PSS, PROVEN.  Every statement here is its support's statement UNCHANGED.}\par
\vnbprop{strictly-mono-ge-id-ind}{\vnbstmt{\forall k \in \mathbb{N}, i.\; (\operatorname{strictly\text{-}mono\text{-}nn}(i) \Rightarrow (k \leq i(k)))}}
\vnbprop{strictly-mono-ge-id}{\vnbstmt{\forall i.; \forall k \in \mathbb{N}.\; (\operatorname{strictly\text{-}mono\text{-}nn}(i) \Rightarrow (k \leq i(k)))}}

\vnbbookfile{theorem-library/rake-extended-order}{theorem-library/rake-extended-order}
\noindent{\small theorem-library/rake-extended-order.scm -- the order and addition theory of RR-POS-STAR = [0, +inf], built on the user's new axiom `pos-inf-above-reals' (structure-library/extended-reals.scm, stamped `primitive', 2026-09-18).}\par
\vnbprop{rr-pos-star-le-refl}{\vnbstmt{\forall x \in \operatorname{rr\text{-}pos\text{-}star}.\; (x \leq x)}}
\vnbprop{rr-pos-star-le-real-in-rr}{\vnbstmt{\forall x, y.\; ((((x \in \operatorname{rr\text{-}pos\text{-}star}) \wedge (y \in \mathbb{R})) \wedge (x \leq y)) \Rightarrow (x \in \mathbb{R}))}}
\vnbprop{rr-pos-star-le-trans}{\vnbstmt{\begin{array}{@{}l@{}} \forall x, y, z. \\ \quad (((x \in \operatorname{rr\text{-}pos\text{-}star}) \wedge ((y \in \operatorname{rr\text{-}pos\text{-}star}) \wedge (z \in \operatorname{rr\text{-}pos\text{-}star}))) \wedge ((x \leq y) \wedge (y \leq z))) \Rightarrow  \\ \quad \quad x \leq z \end{array}}}
\vnbprop{rr-pos-star-le-antisymm}{\vnbstmt{\begin{array}{@{}l@{}} \forall x, y. \\ \quad (((x \in \operatorname{rr\text{-}pos\text{-}star}) \wedge (y \in \operatorname{rr\text{-}pos\text{-}star})) \wedge ((x \leq y) \wedge (y \leq x))) \Rightarrow (x = y) \end{array}}}
\vnbprop{eplus-closed}{\vnbstmt{\begin{array}{@{}l@{}} \forall x, y. \\ \quad ((x \in \operatorname{rr\text{-}pos\text{-}star}) \wedge (y \in \operatorname{rr\text{-}pos\text{-}star})) \Rightarrow (\operatorname{eplus}(x, y) \in \operatorname{rr\text{-}pos\text{-}star}) \end{array}}}
\vnbprop{eplus-zero-left}{\vnbstmt{\forall x \in \operatorname{rr\text{-}pos\text{-}star}.\; (\operatorname{eplus}(0, x) = x)}}
\vnbprop{eplus-zero-right}{\vnbstmt{\forall x \in \operatorname{rr\text{-}pos\text{-}star}.\; (\operatorname{eplus}(x, 0) = x)}}
\vnbprop{eplus-comm}{\vnbstmt{\begin{array}{@{}l@{}} \forall x, y. \\ \quad ((x \in \operatorname{rr\text{-}pos\text{-}star}) \wedge (y \in \operatorname{rr\text{-}pos\text{-}star})) \Rightarrow (\operatorname{eplus}(x, y) = \operatorname{eplus}(y, x)) \end{array}}}
\vnbprop{eplus-assoc}{\vnbstmt{\begin{array}{@{}l@{}} \forall x, y, z. \\ \quad ((x \in \operatorname{rr\text{-}pos\text{-}star}) \wedge ((y \in \operatorname{rr\text{-}pos\text{-}star}) \wedge (z \in \operatorname{rr\text{-}pos\text{-}star}))) \Rightarrow  \\ \quad \quad \operatorname{eplus}(\operatorname{eplus}(x, y), z) = \operatorname{eplus}(x, \operatorname{eplus}(y, z)) \end{array}}}
\vnbprop{eplus-le-left}{\vnbstmt{\forall x, y.\; (((x \in \operatorname{rr\text{-}pos\text{-}star}) \wedge (y \in \operatorname{rr\text{-}pos\text{-}star})) \Rightarrow (x \leq \operatorname{eplus}(x, y)))}}
\vnbprop{eplus-mono}{\vnbstmt{\begin{array}{@{}l@{}} \forall x, y, u, v. \\ \quad x \in \operatorname{rr\text{-}pos\text{-}star} \wedge \\ \quad y \in \operatorname{rr\text{-}pos\text{-}star} \wedge \\ \quad u \in \operatorname{rr\text{-}pos\text{-}star} \wedge \\ \quad v \in \operatorname{rr\text{-}pos\text{-}star} \wedge \\ \quad x \leq u \wedge \\ \quad y \leq v \Rightarrow \\ \quad \quad \operatorname{eplus}(x, y) \leq \operatorname{eplus}(u, v) \end{array}}}

\vnbbookfile{theorem-library/rake-esup-defined}{theorem-library/rake-esup-defined}
\noindent{\small theorem-library/rake-esup-defined.scm -- ESUP DEFINED as a definite description, and its four characterising axioms PROVEN from the definition. Assignment "ESUP and ESUM defined" (the user's decision of 2026-09-19). Helper prefix `esd-'.}\par
\vnbprop{esup-exists}{\vnbstmt{\begin{array}{@{}l@{}} \forall s. \\ \quad (s \subseteq \operatorname{rr\text{-}pos\text{-}star}) \Rightarrow  \\ \quad \quad \exists \mathit{bb} \in \operatorname{rr\text{-}pos\text{-}star}. \\ \quad \quad \quad (\forall x \in s.\; (x \leq \mathit{bb}) \wedge \forall b \in \operatorname{rr\text{-}pos\text{-}star}.\; (\forall x \in s.\; (x \leq b) \Rightarrow (\mathit{bb} \leq b))) \end{array}}}
\vnbprop{esup-unique}{\vnbstmt{\begin{array}{@{}l@{}} \forall s, b_{1}, b_{2}. \\ \quad b_{1} \in \operatorname{rr\text{-}pos\text{-}star} \wedge \\ \quad \forall x \in s.\; (x \leq b_{1}) \wedge \\ \quad \forall b \in \operatorname{rr\text{-}pos\text{-}star}.\; (\forall x \in s.\; (x \leq b) \Rightarrow (b_{1} \leq b)) \wedge \\ \quad b_{2} \in \operatorname{rr\text{-}pos\text{-}star} \wedge \\ \quad \forall x \in s.\; (x \leq b_{2}) \wedge \\ \quad \forall b \in \operatorname{rr\text{-}pos\text{-}star}.\; (\forall x \in s.\; (x \leq b) \Rightarrow (b_{2} \leq b)) \Rightarrow \\ \quad \quad b_{1} = b_{2} \end{array}}}
\vnbprop{esup-prop}{\vnbstmt{\begin{array}{@{}l@{}} \forall s. \\ \quad (s \subseteq \operatorname{rr\text{-}pos\text{-}star}) \Rightarrow  \\ \quad \quad \operatorname{esup}(s) \in \operatorname{rr\text{-}pos\text{-}star} \wedge \\ \quad \quad \forall x \in s.\; (x \leq \operatorname{esup}(s)) \wedge \\ \quad \quad \forall b \in \operatorname{rr\text{-}pos\text{-}star}.\; (\forall x \in s.\; (x \leq b) \Rightarrow (\operatorname{esup}(s) \leq b)) \end{array}}}
\vnbprop{esup-in}{\vnbstmt{\forall s.\; ((s \subseteq \operatorname{rr\text{-}pos\text{-}star}) \Rightarrow (\operatorname{esup}(s) \in \operatorname{rr\text{-}pos\text{-}star}))}}
\vnbprop{esup-upper}{\vnbstmt{\forall s.; \forall x \in s.\; ((s \subseteq \operatorname{rr\text{-}pos\text{-}star}) \Rightarrow (x \leq \operatorname{esup}(s)))}}
\vnbprop{esup-least}{\vnbstmt{\begin{array}{@{}l@{}} \forall s.; \forall b \in \operatorname{rr\text{-}pos\text{-}star}. \\ \quad ((s \subseteq \operatorname{rr\text{-}pos\text{-}star}) \wedge \forall x \in s.\; (x \leq b)) \Rightarrow (\operatorname{esup}(s) \leq b) \end{array}}}
\vnbprop{esup-empty}{\vnbstmt{\operatorname{esup}(\emptyset) = 0}}

\vnbbookfile{theorem-library/rake-rr-pos-star-monoid}{theorem-library/rake-rr-pos-star-monoid}
\noindent{\small theorem-library/rake-rr-pos-star-monoid.scm -- IS-COMM-MONOID(RR-POS-STAR-ADD-MONOID), PROVEN.  Rake batch 6-C.  Helper prefix `rpm-'.}\par
\vnbprop{rr-pos-star-is-comm-monoid}{\vnbstmt{\operatorname{is\text{-}comm\text{-}monoid}(\operatorname{rr\text{-}pos\text{-}star\text{-}add\text{-}monoid})}}

\vnbbookfile{theorem-library/view-typing-proofs}{theorem-library/view-typing-proofs}
\noindent{\small A file shared across books; the entries that belong to this one.}\par
\noindent{\small theorem-library/view-typing-proofs.scm -- THE DEFERRED FIELD VIEW TYPINGS, PROVEN (D8, 2026-11-15; the user's decision "Prove views + fix F1/F2").}\par
\vnbprop{normed-field-is-field-ring}{\vnbstmt{\begin{array}{@{}l@{}} \forall f. \\ \quad \operatorname{is\text{-}normed\text{-}field}(f) \Rightarrow \operatorname{is\text{-}field\text{-}ring}(\operatorname{normed\text{-}field\text{-}as\text{-}commutative\text{-}ring}(f)) \end{array}}}
\vnbprop{normed-field-as-integral-domain-is-integral-domain}{\vnbstmt{\begin{array}{@{}l@{}} \forall r. \\ \quad \operatorname{is\text{-}normed\text{-}field}(r) \Rightarrow  \\ \quad \quad \operatorname{is\text{-}integral\text{-}domain}(\operatorname{normed\text{-}field\text{-}as\text{-}integral\text{-}domain}(r)) \end{array}}}

\vnbbookfile{theorem-library/extinf-semigroups}{theorem-library/extinf-semigroups}
\noindent{\small extinf-semigroups.scm -- extended addition on (-inf, +inf] and on [-inf, +inf), and the order of each line: the laws, stated and proven SYMMETRICALLY.  Agent EXT-INF, 2026-10-26 (the user's decision of the day: neg-inf-below-reals, "but the semigroup properties should be clearly stated for both").}\par
\vnbprop{rr-ext-up-membership}{\vnbstmt{\forall x.\; ((x \in \operatorname{rr\text{-}ext\text{-}up}) \Leftrightarrow ((x \in \mathbb{R}) \vee (x = \operatorname{pos\text{-}inf})))}}
\vnbprop{rr-ext-down-membership}{\vnbstmt{\forall x.\; ((x \in \operatorname{rr\text{-}ext\text{-}down}) \Leftrightarrow ((x \in \mathbb{R}) \vee (x = \operatorname{neg\text{-}inf})))}}
\vnbprop{rr-in-rr-ext-up}{\vnbstmt{\forall x \in \mathbb{R}.\; (x \in \operatorname{rr\text{-}ext\text{-}up})}}
\vnbprop{pos-inf-in-rr-ext-up}{\vnbstmt{\operatorname{pos\text{-}inf} \in \operatorname{rr\text{-}ext\text{-}up}}}
\vnbprop{rr-ext-up-in-rr-star}{\vnbstmt{\forall x \in \operatorname{rr\text{-}ext\text{-}up}.\; (x \in \operatorname{rr\text{-}star})}}
\vnbprop{rr-ext-up-below-pos-inf}{\vnbstmt{\forall x \in \operatorname{rr\text{-}ext\text{-}up}.\; (x \leq \operatorname{pos\text{-}inf})}}
\vnbprop{rr-in-rr-ext-down}{\vnbstmt{\forall x \in \mathbb{R}.\; (x \in \operatorname{rr\text{-}ext\text{-}down})}}
\vnbprop{neg-inf-in-rr-ext-down}{\vnbstmt{\operatorname{neg\text{-}inf} \in \operatorname{rr\text{-}ext\text{-}down}}}
\vnbprop{rr-ext-down-in-rr-star}{\vnbstmt{\forall x \in \operatorname{rr\text{-}ext\text{-}down}.\; (x \in \operatorname{rr\text{-}star})}}
\vnbprop{rr-ext-down-above-neg-inf}{\vnbstmt{\forall x \in \operatorname{rr\text{-}ext\text{-}down}.\; (\operatorname{neg\text{-}inf} \leq x)}}
\vnbprop{rr-ext-up-le-real-in-rr}{\vnbstmt{\forall x \in \operatorname{rr\text{-}ext\text{-}up}, y \in \mathbb{R}.\; ((x \leq y) \Rightarrow (x \in \mathbb{R}))}}
\vnbprop{rr-ext-down-ge-real-in-rr}{\vnbstmt{\forall x \in \operatorname{rr\text{-}ext\text{-}down}, y \in \mathbb{R}.\; ((y \leq x) \Rightarrow (x \in \mathbb{R}))}}
\vnbprop{rr-ext-up-le-refl}{\vnbstmt{\forall x \in \operatorname{rr\text{-}ext\text{-}up}.\; (x \leq x)}}
\vnbprop{rr-ext-down-le-refl}{\vnbstmt{\forall x \in \operatorname{rr\text{-}ext\text{-}down}.\; (x \leq x)}}
\vnbprop{rr-ext-up-le-trans}{\vnbstmt{\begin{array}{@{}l@{}} \forall x, y, z \in \operatorname{rr\text{-}ext\text{-}up}. \\ \quad ((x \leq y) \wedge (y \leq z)) \Rightarrow (x \leq z) \end{array}}}
\vnbprop{rr-ext-down-le-trans}{\vnbstmt{\begin{array}{@{}l@{}} \forall x, y, z \in \operatorname{rr\text{-}ext\text{-}down}. \\ \quad ((x \leq y) \wedge (y \leq z)) \Rightarrow (x \leq z) \end{array}}}
\vnbprop{rr-ext-up-le-antisymm}{\vnbstmt{\forall x, y \in \operatorname{rr\text{-}ext\text{-}up}.\; (((x \leq y) \wedge (y \leq x)) \Rightarrow (x = y))}}
\vnbprop{rr-ext-down-le-antisymm}{\vnbstmt{\forall x, y \in \operatorname{rr\text{-}ext\text{-}down}.\; (((x \leq y) \wedge (y \leq x)) \Rightarrow (x = y))}}
\vnbprop{eplus-up-real}{\vnbstmt{\forall x, y \in \mathbb{R}.\; (\operatorname{eplus\text{-}up}(x, y) = (x + y))}}
\vnbprop{eplus-down-real}{\vnbstmt{\forall x, y \in \mathbb{R}.\; (\operatorname{eplus\text{-}down}(x, y) = (x + y))}}
\vnbprop{eplus-up-pos-inf-left}{\vnbstmt{\forall y \in \operatorname{rr\text{-}ext\text{-}up}.\; (\operatorname{eplus\text{-}up}(\operatorname{pos\text{-}inf}, y) = \operatorname{pos\text{-}inf})}}
\vnbprop{eplus-up-pos-inf-right}{\vnbstmt{\forall y \in \operatorname{rr\text{-}ext\text{-}up}.\; (\operatorname{eplus\text{-}up}(y, \operatorname{pos\text{-}inf}) = \operatorname{pos\text{-}inf})}}
\vnbprop{eplus-down-neg-inf-left}{\vnbstmt{\forall y \in \operatorname{rr\text{-}ext\text{-}down}.\; (\operatorname{eplus\text{-}down}(\operatorname{neg\text{-}inf}, y) = \operatorname{neg\text{-}inf})}}
\vnbprop{eplus-down-neg-inf-right}{\vnbstmt{\forall y \in \operatorname{rr\text{-}ext\text{-}down}.\; (\operatorname{eplus\text{-}down}(y, \operatorname{neg\text{-}inf}) = \operatorname{neg\text{-}inf})}}
\vnbprop{eplus-up-closed}{\vnbstmt{\forall x, y \in \operatorname{rr\text{-}ext\text{-}up}.\; (\operatorname{eplus\text{-}up}(x, y) \in \operatorname{rr\text{-}ext\text{-}up})}}
\vnbprop{eplus-down-closed}{\vnbstmt{\forall x, y \in \operatorname{rr\text{-}ext\text{-}down}.\; (\operatorname{eplus\text{-}down}(x, y) \in \operatorname{rr\text{-}ext\text{-}down})}}
\vnbprop{eplus-up-comm}{\vnbstmt{\forall x, y \in \operatorname{rr\text{-}ext\text{-}up}.\; (\operatorname{eplus\text{-}up}(x, y) = \operatorname{eplus\text{-}up}(y, x))}}
\vnbprop{eplus-down-comm}{\vnbstmt{\begin{array}{@{}l@{}} \forall x, y \in \operatorname{rr\text{-}ext\text{-}down}. \\ \quad \operatorname{eplus\text{-}down}(x, y) = \operatorname{eplus\text{-}down}(y, x) \end{array}}}
\vnbprop{eplus-up-assoc}{\vnbstmt{\begin{array}{@{}l@{}} \forall x, y, z \in \operatorname{rr\text{-}ext\text{-}up}. \\ \quad \operatorname{eplus\text{-}up}(\operatorname{eplus\text{-}up}(x, y), z) = \operatorname{eplus\text{-}up}(x, \operatorname{eplus\text{-}up}(y, z)) \end{array}}}
\vnbprop{eplus-down-assoc}{\vnbstmt{\begin{array}{@{}l@{}} \forall x, y, z \in \operatorname{rr\text{-}ext\text{-}down}. \\ \quad \operatorname{eplus\text{-}down}(\operatorname{eplus\text{-}down}(x, y), z) = \operatorname{eplus\text{-}down}(x, \operatorname{eplus\text{-}down}(y, z)) \end{array}}}
\vnbprop{eplus-up-zero-left}{\vnbstmt{\forall x \in \operatorname{rr\text{-}ext\text{-}up}.\; (\operatorname{eplus\text{-}up}(0, x) = x)}}
\vnbprop{eplus-up-zero-right}{\vnbstmt{\forall x \in \operatorname{rr\text{-}ext\text{-}up}.\; (\operatorname{eplus\text{-}up}(x, 0) = x)}}
\vnbprop{eplus-down-zero-left}{\vnbstmt{\forall x \in \operatorname{rr\text{-}ext\text{-}down}.\; (\operatorname{eplus\text{-}down}(0, x) = x)}}
\vnbprop{eplus-down-zero-right}{\vnbstmt{\forall x \in \operatorname{rr\text{-}ext\text{-}down}.\; (\operatorname{eplus\text{-}down}(x, 0) = x)}}
\vnbprop{eplus-up-mono}{\vnbstmt{\begin{array}{@{}l@{}} \forall x, y, z \in \operatorname{rr\text{-}ext\text{-}up}. \\ \quad (x \leq y) \Rightarrow (\operatorname{eplus\text{-}up}(x, z) \leq \operatorname{eplus\text{-}up}(y, z)) \end{array}}}
\vnbprop{eplus-down-mono}{\vnbstmt{\begin{array}{@{}l@{}} \forall x, y, z \in \operatorname{rr\text{-}ext\text{-}down}. \\ \quad (x \leq y) \Rightarrow (\operatorname{eplus\text{-}down}(x, z) \leq \operatorname{eplus\text{-}down}(y, z)) \end{array}}}
\vnbprop{rr-pos-star-in-rr-ext-up}{\vnbstmt{\forall x \in \operatorname{rr\text{-}pos\text{-}star}.\; (x \in \operatorname{rr\text{-}ext\text{-}up})}}
\vnbprop{eplus-up-agrees-eplus}{\vnbstmt{\forall x, y \in \operatorname{rr\text{-}pos\text{-}star}.\; (\operatorname{eplus\text{-}up}(x, y) = \operatorname{eplus}(x, y))}}

\vnbbookfile{theorem-library/hom-laws}{theorem-library/hom-laws}
\noindent{\small A file shared across books; the entries that belong to this one.}\par
\noindent{\small hom-laws.scm -- the models of a structure form a CATEGORY: identity arrows, composition, and the hom-set as a set.}\par
\vnbprop{hom-normed-ag-member-iff}{\vnbstmt{\begin{array}{@{}l@{}} \forall a, b, f. \\ \quad (f \in \operatorname{hom\text{-}normed\text{-}ag}(a, b)) \Leftrightarrow ((f \in (\operatorname{carr}(a) \to \operatorname{carr}(b))) \wedge \operatorname{is\text{-}hom\text{-}normed\text{-}ag}(a, b, f)) \end{array}}}
\vnbprop{hom-normed-ag-id}{\vnbstmt{\forall a.\; (\operatorname{is\text{-}normed\text{-}ag}(a) \Rightarrow \operatorname{is\text{-}hom\text{-}normed\text{-}ag}(a, a, \operatorname{id\text{-}fun}(\operatorname{carr}(a))))}}
\vnbprop{hom-normed-ag-in-set}{\vnbstmt{\begin{array}{@{}l@{}} \forall a, b. \\ \quad (\operatorname{is\text{-}normed\text{-}ag}(a) \wedge \operatorname{is\text{-}normed\text{-}ag}(b)) \Rightarrow (\operatorname{hom\text{-}normed\text{-}ag}(a, b) \in \mathbf{Set}) \end{array}}}
\vnbprop{hom-normed-field-member-iff}{\vnbstmt{\begin{array}{@{}l@{}} \forall a, b, f. \\ \quad (f \in \operatorname{hom\text{-}normed\text{-}field}(a, b)) \Leftrightarrow ((f \in (\operatorname{carr}(a) \to \operatorname{carr}(b))) \wedge \operatorname{is\text{-}hom\text{-}normed\text{-}field}(a, b, f)) \end{array}}}
\vnbprop{hom-normed-field-id}{\vnbstmt{\begin{array}{@{}l@{}} \forall a. \\ \quad \operatorname{is\text{-}normed\text{-}field}(a) \Rightarrow \operatorname{is\text{-}hom\text{-}normed\text{-}field}(a, a, \operatorname{id\text{-}fun}(\operatorname{carr}(a))) \end{array}}}
\vnbprop{hom-normed-field-in-set}{\vnbstmt{\begin{array}{@{}l@{}} \forall a, b. \\ \quad (\operatorname{is\text{-}normed\text{-}field}(a) \wedge \operatorname{is\text{-}normed\text{-}field}(b)) \Rightarrow (\operatorname{hom\text{-}normed\text{-}field}(a, b) \in \mathbf{Set}) \end{array}}}
\vnbprop{hom-normed-vector-space-member-iff}{\vnbstmt{\begin{array}{@{}l@{}} \forall a, b, f. \\ \quad f \in \operatorname{hom\text{-}normed\text{-}vector\text{-}space}(a, b) \\ \quad \Leftrightarrow\; (f \in (\operatorname{vec}(a) \to \operatorname{vec}(b))) \wedge \operatorname{is\text{-}hom\text{-}normed\text{-}vector\text{-}space}(a, b, f) \end{array}}}
\vnbprop{hom-normed-vector-space-id}{\vnbstmt{\begin{array}{@{}l@{}} \forall a. \\ \quad \operatorname{is\text{-}normed\text{-}vector\text{-}space}(a) \Rightarrow  \\ \quad \quad \operatorname{is\text{-}hom\text{-}normed\text{-}vector\text{-}space}(a, a, \operatorname{id\text{-}fun}(\operatorname{vec}(a))) \end{array}}}
\vnbprop{hom-normed-vector-space-in-set}{\vnbstmt{\begin{array}{@{}l@{}} \forall a, b. \\ \quad (\operatorname{is\text{-}normed\text{-}vector\text{-}space}(a) \wedge \operatorname{is\text{-}normed\text{-}vector\text{-}space}(b)) \Rightarrow  \\ \quad \quad \operatorname{hom\text{-}normed\text{-}vector\text{-}space}(a, b) \in \mathbf{Set} \end{array}}}
\vnbprop{hom-complex-inner-product-space-member-iff}{\vnbstmt{\begin{array}{@{}l@{}} \forall a, b, f. \\ \quad f \in \operatorname{hom\text{-}complex\text{-}inner\text{-}product\text{-}space}(a, b) \\ \quad \Leftrightarrow\; (f \in (\operatorname{vec}(a) \to \operatorname{vec}(b))) \wedge \operatorname{is\text{-}hom\text{-}complex\text{-}inner\text{-}product\text{-}space}(a, b, f) \end{array}}}
\vnbprop{hom-complex-inner-product-space-id}{\vnbstmt{\begin{array}{@{}l@{}} \forall a. \\ \quad \operatorname{is\text{-}complex\text{-}inner\text{-}product\text{-}space}(a) \Rightarrow  \\ \quad \quad \operatorname{is\text{-}hom\text{-}complex\text{-}inner\text{-}product\text{-}space}(a, a, \operatorname{id\text{-}fun}(\operatorname{vec}(a))) \end{array}}}
\vnbprop{hom-complex-inner-product-space-in-set}{\vnbstmt{\begin{array}{@{}l@{}} \forall a, b. \\ \quad (\operatorname{is\text{-}complex\text{-}inner\text{-}product\text{-}space}(a) \wedge \operatorname{is\text{-}complex\text{-}inner\text{-}product\text{-}space}(b)) \Rightarrow  \\ \quad \quad \operatorname{hom\text{-}complex\text{-}inner\text{-}product\text{-}space}(a, b) \in \mathbf{Set} \end{array}}}
\vnbprop{hom-metric-space-member-iff}{\vnbstmt{\begin{array}{@{}l@{}} \forall a, b, f. \\ \quad (f \in \operatorname{hom\text{-}metric\text{-}space}(a, b)) \Leftrightarrow ((f \in (\operatorname{pts}(a) \to \operatorname{pts}(b))) \wedge \operatorname{is\text{-}hom\text{-}metric\text{-}space}(a, b, f)) \end{array}}}
\vnbprop{hom-metric-space-id}{\vnbstmt{\begin{array}{@{}l@{}} \forall a. \\ \quad \operatorname{is\text{-}metric\text{-}space}(a) \Rightarrow \operatorname{is\text{-}hom\text{-}metric\text{-}space}(a, a, \operatorname{id\text{-}fun}(\operatorname{pts}(a))) \end{array}}}
\vnbprop{hom-metric-space-in-set}{\vnbstmt{\begin{array}{@{}l@{}} \forall a, b. \\ \quad (\operatorname{is\text{-}metric\text{-}space}(a) \wedge \operatorname{is\text{-}metric\text{-}space}(b)) \Rightarrow (\operatorname{hom\text{-}metric\text{-}space}(a, b) \in \mathbf{Set}) \end{array}}}
\vnbprop{hom-pseudometric-space-member-iff}{\vnbstmt{\begin{array}{@{}l@{}} \forall a, b, f. \\ \quad f \in \operatorname{hom\text{-}pseudometric\text{-}space}(a, b) \\ \quad \Leftrightarrow\; (f \in (\operatorname{pts}(a) \to \operatorname{pts}(b))) \wedge \operatorname{is\text{-}hom\text{-}pseudometric\text{-}space}(a, b, f) \end{array}}}
\vnbprop{hom-pseudometric-space-id}{\vnbstmt{\begin{array}{@{}l@{}} \forall a. \\ \quad \operatorname{is\text{-}pseudometric\text{-}space}(a) \Rightarrow \operatorname{is\text{-}hom\text{-}pseudometric\text{-}space}(a, a, \operatorname{id\text{-}fun}(\operatorname{pts}(a))) \end{array}}}
\vnbprop{hom-pseudometric-space-in-set}{\vnbstmt{\begin{array}{@{}l@{}} \forall a, b. \\ \quad (\operatorname{is\text{-}pseudometric\text{-}space}(a) \wedge \operatorname{is\text{-}pseudometric\text{-}space}(b)) \Rightarrow  \\ \quad \quad \operatorname{hom\text{-}pseudometric\text{-}space}(a, b) \in \mathbf{Set} \end{array}}}
\vnbprop{hom-c-metric-space-member-iff}{\vnbstmt{\begin{array}{@{}l@{}} \forall a, b, f. \\ \quad (f \in \operatorname{hom\text{-}c\text{-}metric\text{-}space}(a, b)) \Leftrightarrow ((f \in (\operatorname{pts}(a) \to \operatorname{pts}(b))) \wedge \operatorname{is\text{-}hom\text{-}c\text{-}metric\text{-}space}(a, b, f)) \end{array}}}
\vnbprop{hom-c-metric-space-id}{\vnbstmt{\begin{array}{@{}l@{}} \forall a. \\ \quad \operatorname{is\text{-}c\text{-}metric\text{-}space}(a) \Rightarrow \operatorname{is\text{-}hom\text{-}c\text{-}metric\text{-}space}(a, a, \operatorname{id\text{-}fun}(\operatorname{pts}(a))) \end{array}}}
\vnbprop{hom-c-metric-space-in-set}{\vnbstmt{\begin{array}{@{}l@{}} \forall a, b. \\ \quad (\operatorname{is\text{-}c\text{-}metric\text{-}space}(a) \wedge \operatorname{is\text{-}c\text{-}metric\text{-}space}(b)) \Rightarrow  \\ \quad \quad \operatorname{hom\text{-}c\text{-}metric\text{-}space}(a, b) \in \mathbf{Set} \end{array}}}
\vnbprop{hom-top-space-member-iff}{\vnbstmt{\begin{array}{@{}l@{}} \forall a, b, f. \\ \quad (f \in \operatorname{hom\text{-}top\text{-}space}(a, b)) \Leftrightarrow ((f \in (\operatorname{pts}(a) \to \operatorname{pts}(b))) \wedge \operatorname{is\text{-}hom\text{-}top\text{-}space}(a, b, f)) \end{array}}}
\vnbprop{hom-top-space-id}{\vnbstmt{\forall s.\; (\operatorname{is\text{-}top\text{-}space}(s) \Rightarrow \operatorname{is\text{-}hom\text{-}top\text{-}space}(s, s, \operatorname{id\text{-}fun}(\operatorname{pts}(s))))}}
\vnbprop{hom-top-space-in-set}{\vnbstmt{\begin{array}{@{}l@{}} \forall a, b. \\ \quad (\operatorname{is\text{-}top\text{-}space}(a) \wedge \operatorname{is\text{-}top\text{-}space}(b)) \Rightarrow (\operatorname{hom\text{-}top\text{-}space}(a, b) \in \mathbf{Set}) \end{array}}}
\vnbprop{hom-metrizable-top-space-member-iff}{\vnbstmt{\begin{array}{@{}l@{}} \forall a, b, f. \\ \quad f \in \operatorname{hom\text{-}metrizable\text{-}top\text{-}space}(a, b) \\ \quad \Leftrightarrow\; (f \in (\operatorname{pts}(a) \to \operatorname{pts}(b))) \wedge \operatorname{is\text{-}hom\text{-}metrizable\text{-}top\text{-}space}(a, b, f) \end{array}}}
\vnbprop{hom-metrizable-top-space-id}{\vnbstmt{\begin{array}{@{}l@{}} \forall s. \\ \quad \operatorname{is\text{-}metrizable\text{-}top\text{-}space}(s) \Rightarrow  \\ \quad \quad \operatorname{is\text{-}hom\text{-}metrizable\text{-}top\text{-}space}(s, s, \operatorname{id\text{-}fun}(\operatorname{pts}(s))) \end{array}}}
\vnbprop{hom-metrizable-top-space-in-set}{\vnbstmt{\begin{array}{@{}l@{}} \forall a, b. \\ \quad (\operatorname{is\text{-}metrizable\text{-}top\text{-}space}(a) \wedge \operatorname{is\text{-}metrizable\text{-}top\text{-}space}(b)) \Rightarrow  \\ \quad \quad \operatorname{hom\text{-}metrizable\text{-}top\text{-}space}(a, b) \in \mathbf{Set} \end{array}}}

\vnbbookfile{theorem-library/hom-functors}{theorem-library/hom-functors}
\noindent{\small A file shared across books; the entries that belong to this one.}\par
\noindent{\small hom-functors.scm -- the Hom functors: Hom(a, -) covariant, Hom(-, b) contravariant, as maps on hom-sets.}\par
\vnbprop{hom-normed-ag-post-type}{\vnbstmt{\begin{array}{@{}l@{}} \forall a, b, c, h. \\ \quad (\operatorname{is\text{-}normed\text{-}ag}(a) \wedge (h \in \operatorname{hom\text{-}normed\text{-}ag}(b, c))) \Rightarrow  \\ \quad \quad (\lambda f \in \operatorname{hom\text{-}normed\text{-}ag}(a, b).\; \operatorname{compose}(h, f)) \in (\operatorname{hom\text{-}normed\text{-}ag}(a, b) \to \operatorname{hom\text{-}normed\text{-}ag}(a, c)) \end{array}}}
\vnbprop{hom-normed-ag-pre-type}{\vnbstmt{\begin{array}{@{}l@{}} \forall a, b, c, k. \\ \quad (\operatorname{is\text{-}normed\text{-}ag}(b) \wedge (k \in \operatorname{hom\text{-}normed\text{-}ag}(c, a))) \Rightarrow  \\ \quad \quad (\lambda f \in \operatorname{hom\text{-}normed\text{-}ag}(a, b).\; \operatorname{compose}(f, k)) \in (\operatorname{hom\text{-}normed\text{-}ag}(a, b) \to \operatorname{hom\text{-}normed\text{-}ag}(c, b)) \end{array}}}
\vnbprop{hom-normed-field-post-type}{\vnbstmt{\begin{array}{@{}l@{}} \forall a, b, c, h. \\ \quad (\operatorname{is\text{-}normed\text{-}field}(a) \wedge (h \in \operatorname{hom\text{-}normed\text{-}field}(b, c))) \Rightarrow  \\ \quad \quad (\lambda f \in \operatorname{hom\text{-}normed\text{-}field}(a, b).\; \operatorname{compose}(h, f)) \\ \quad \quad \in\; (\operatorname{hom\text{-}normed\text{-}field}(a, b) \to \operatorname{hom\text{-}normed\text{-}field}(a, c)) \end{array}}}
\vnbprop{hom-normed-field-pre-type}{\vnbstmt{\begin{array}{@{}l@{}} \forall a, b, c, k. \\ \quad (\operatorname{is\text{-}normed\text{-}field}(b) \wedge (k \in \operatorname{hom\text{-}normed\text{-}field}(c, a))) \Rightarrow  \\ \quad \quad (\lambda f \in \operatorname{hom\text{-}normed\text{-}field}(a, b).\; \operatorname{compose}(f, k)) \\ \quad \quad \in\; (\operatorname{hom\text{-}normed\text{-}field}(a, b) \to \operatorname{hom\text{-}normed\text{-}field}(c, b)) \end{array}}}
\vnbprop{hom-normed-vector-space-post-type}{\vnbstmt{\begin{array}{@{}l@{}} \forall a, b, c, h. \\ \quad (\operatorname{is\text{-}normed\text{-}vector\text{-}space}(a) \wedge (h \in \operatorname{hom\text{-}normed\text{-}vector\text{-}space}(b, c))) \Rightarrow  \\ \quad \quad (\lambda f \in \operatorname{hom\text{-}normed\text{-}vector\text{-}space}(a, b).\; \operatorname{compose}(h, f)) \\ \quad \quad \in\; (\operatorname{hom\text{-}normed\text{-}vector\text{-}space}(a, b) \to \operatorname{hom\text{-}normed\text{-}vector\text{-}space}(a, c)) \end{array}}}
\vnbprop{hom-normed-vector-space-pre-type}{\vnbstmt{\begin{array}{@{}l@{}} \forall a, b, c, k. \\ \quad (\operatorname{is\text{-}normed\text{-}vector\text{-}space}(b) \wedge (k \in \operatorname{hom\text{-}normed\text{-}vector\text{-}space}(c, a))) \Rightarrow  \\ \quad \quad (\lambda f \in \operatorname{hom\text{-}normed\text{-}vector\text{-}space}(a, b).\; \operatorname{compose}(f, k)) \\ \quad \quad \in\; (\operatorname{hom\text{-}normed\text{-}vector\text{-}space}(a, b) \to \operatorname{hom\text{-}normed\text{-}vector\text{-}space}(c, b)) \end{array}}}
\vnbprop{hom-complex-inner-product-space-post-type}{\vnbstmt{\begin{array}{@{}l@{}} \forall a, b, c, h. \\ \quad (\operatorname{is\text{-}complex\text{-}inner\text{-}product\text{-}space}(a) \wedge (h \in \operatorname{hom\text{-}complex\text{-}inner\text{-}product\text{-}space}(b, c))) \Rightarrow  \\ \quad \quad (\lambda f \in \operatorname{hom\text{-}complex\text{-}inner\text{-}product\text{-}space}(a, b).\; \operatorname{compose}(h, f)) \\ \quad \quad \in\; (\operatorname{hom\text{-}complex\text{-}inner\text{-}product\text{-}space}(a, b) \to \operatorname{hom\text{-}complex\text{-}inner\text{-}product\text{-}space}(a, c)) \end{array}}}
\vnbprop{hom-complex-inner-product-space-pre-type}{\vnbstmt{\begin{array}{@{}l@{}} \forall a, b, c, k. \\ \quad (\operatorname{is\text{-}complex\text{-}inner\text{-}product\text{-}space}(b) \wedge (k \in \operatorname{hom\text{-}complex\text{-}inner\text{-}product\text{-}space}(c, a))) \Rightarrow  \\ \quad \quad (\lambda f \in \operatorname{hom\text{-}complex\text{-}inner\text{-}product\text{-}space}(a, b).\; \operatorname{compose}(f, k)) \\ \quad \quad \in\; (\operatorname{hom\text{-}complex\text{-}inner\text{-}product\text{-}space}(a, b) \to \operatorname{hom\text{-}complex\text{-}inner\text{-}product\text{-}space}(c, b)) \end{array}}}
\vnbprop{hom-metric-space-post-type}{\vnbstmt{\begin{array}{@{}l@{}} \forall a, b, c, h. \\ \quad (\operatorname{is\text{-}metric\text{-}space}(a) \wedge (h \in \operatorname{hom\text{-}metric\text{-}space}(b, c))) \Rightarrow  \\ \quad \quad (\lambda f \in \operatorname{hom\text{-}metric\text{-}space}(a, b).\; \operatorname{compose}(h, f)) \\ \quad \quad \in\; (\operatorname{hom\text{-}metric\text{-}space}(a, b) \to \operatorname{hom\text{-}metric\text{-}space}(a, c)) \end{array}}}
\vnbprop{hom-metric-space-pre-type}{\vnbstmt{\begin{array}{@{}l@{}} \forall a, b, c, k. \\ \quad (\operatorname{is\text{-}metric\text{-}space}(b) \wedge (k \in \operatorname{hom\text{-}metric\text{-}space}(c, a))) \Rightarrow  \\ \quad \quad (\lambda f \in \operatorname{hom\text{-}metric\text{-}space}(a, b).\; \operatorname{compose}(f, k)) \\ \quad \quad \in\; (\operatorname{hom\text{-}metric\text{-}space}(a, b) \to \operatorname{hom\text{-}metric\text{-}space}(c, b)) \end{array}}}
\vnbprop{hom-pseudometric-space-post-type}{\vnbstmt{\begin{array}{@{}l@{}} \forall a, b, c, h. \\ \quad (\operatorname{is\text{-}pseudometric\text{-}space}(a) \wedge (h \in \operatorname{hom\text{-}pseudometric\text{-}space}(b, c))) \Rightarrow  \\ \quad \quad (\lambda f \in \operatorname{hom\text{-}pseudometric\text{-}space}(a, b).\; \operatorname{compose}(h, f)) \\ \quad \quad \in\; (\operatorname{hom\text{-}pseudometric\text{-}space}(a, b) \to \operatorname{hom\text{-}pseudometric\text{-}space}(a, c)) \end{array}}}
\vnbprop{hom-pseudometric-space-pre-type}{\vnbstmt{\begin{array}{@{}l@{}} \forall a, b, c, k. \\ \quad (\operatorname{is\text{-}pseudometric\text{-}space}(b) \wedge (k \in \operatorname{hom\text{-}pseudometric\text{-}space}(c, a))) \Rightarrow  \\ \quad \quad (\lambda f \in \operatorname{hom\text{-}pseudometric\text{-}space}(a, b).\; \operatorname{compose}(f, k)) \\ \quad \quad \in\; (\operatorname{hom\text{-}pseudometric\text{-}space}(a, b) \to \operatorname{hom\text{-}pseudometric\text{-}space}(c, b)) \end{array}}}
\vnbprop{hom-c-metric-space-post-type}{\vnbstmt{\begin{array}{@{}l@{}} \forall a, b, c, h. \\ \quad (\operatorname{is\text{-}c\text{-}metric\text{-}space}(a) \wedge (h \in \operatorname{hom\text{-}c\text{-}metric\text{-}space}(b, c))) \Rightarrow  \\ \quad \quad (\lambda f \in \operatorname{hom\text{-}c\text{-}metric\text{-}space}(a, b).\; \operatorname{compose}(h, f)) \\ \quad \quad \in\; (\operatorname{hom\text{-}c\text{-}metric\text{-}space}(a, b) \to \operatorname{hom\text{-}c\text{-}metric\text{-}space}(a, c)) \end{array}}}
\vnbprop{hom-c-metric-space-pre-type}{\vnbstmt{\begin{array}{@{}l@{}} \forall a, b, c, k. \\ \quad (\operatorname{is\text{-}c\text{-}metric\text{-}space}(b) \wedge (k \in \operatorname{hom\text{-}c\text{-}metric\text{-}space}(c, a))) \Rightarrow  \\ \quad \quad (\lambda f \in \operatorname{hom\text{-}c\text{-}metric\text{-}space}(a, b).\; \operatorname{compose}(f, k)) \\ \quad \quad \in\; (\operatorname{hom\text{-}c\text{-}metric\text{-}space}(a, b) \to \operatorname{hom\text{-}c\text{-}metric\text{-}space}(c, b)) \end{array}}}
\vnbprop{hom-top-space-post-type}{\vnbstmt{\begin{array}{@{}l@{}} \forall a, b, c, h. \\ \quad (\operatorname{is\text{-}top\text{-}space}(a) \wedge (h \in \operatorname{hom\text{-}top\text{-}space}(b, c))) \Rightarrow  \\ \quad \quad (\lambda f \in \operatorname{hom\text{-}top\text{-}space}(a, b).\; \operatorname{compose}(h, f)) \in (\operatorname{hom\text{-}top\text{-}space}(a, b) \to \operatorname{hom\text{-}top\text{-}space}(a, c)) \end{array}}}
\vnbprop{hom-top-space-pre-type}{\vnbstmt{\begin{array}{@{}l@{}} \forall a, b, c, k. \\ \quad (\operatorname{is\text{-}top\text{-}space}(b) \wedge (k \in \operatorname{hom\text{-}top\text{-}space}(c, a))) \Rightarrow  \\ \quad \quad (\lambda f \in \operatorname{hom\text{-}top\text{-}space}(a, b).\; \operatorname{compose}(f, k)) \in (\operatorname{hom\text{-}top\text{-}space}(a, b) \to \operatorname{hom\text{-}top\text{-}space}(c, b)) \end{array}}}
\vnbprop{hom-metrizable-top-space-post-type}{\vnbstmt{\begin{array}{@{}l@{}} \forall a, b, c, h. \\ \quad (\operatorname{is\text{-}metrizable\text{-}top\text{-}space}(a) \wedge (h \in \operatorname{hom\text{-}metrizable\text{-}top\text{-}space}(b, c))) \Rightarrow  \\ \quad \quad (\lambda f \in \operatorname{hom\text{-}metrizable\text{-}top\text{-}space}(a, b).\; \operatorname{compose}(h, f)) \\ \quad \quad \in\; (\operatorname{hom\text{-}metrizable\text{-}top\text{-}space}(a, b) \to \operatorname{hom\text{-}metrizable\text{-}top\text{-}space}(a, c)) \end{array}}}
\vnbprop{hom-metrizable-top-space-pre-type}{\vnbstmt{\begin{array}{@{}l@{}} \forall a, b, c, k. \\ \quad (\operatorname{is\text{-}metrizable\text{-}top\text{-}space}(b) \wedge (k \in \operatorname{hom\text{-}metrizable\text{-}top\text{-}space}(c, a))) \Rightarrow  \\ \quad \quad (\lambda f \in \operatorname{hom\text{-}metrizable\text{-}top\text{-}space}(a, b).\; \operatorname{compose}(f, k)) \\ \quad \quad \in\; (\operatorname{hom\text{-}metrizable\text{-}top\text{-}space}(a, b) \to \operatorname{hom\text{-}metrizable\text{-}top\text{-}space}(c, b)) \end{array}}}

\vnbbookfile{structure-library/isom2-rrvs-vocab}{structure-library/isom2-rrvs-vocab}
\noindent{\small isom2-rrvs-vocab.scm -- the real numbers as a FIELD instance, the vector space of real functions on RR in a set closed under the pointwise operations, and the set of antiderivable functions on [a, b] (calculus.pdf Prop 4.8).  DEFINITIONS ONLY; laws in theorem-library/isom2-rrfield.scm and isom2-rrvs.scm.  Agent ISOM, ...}\par
\par\medskip\noindent\textbf{Definition} (\texttt{rr-field}).\;
\begin{equation*}
\begin{array}{@{}l@{}} \operatorname{rr\text{-}field} \\ =\; {[}\mathbb{R}, (\lambda {[}\operatorname{x\_}, \operatorname{y\_}] \in (\mathbb{R} \times \mathbb{R}).\; +(\operatorname{x\_}, \operatorname{y\_})), (\lambda {[}\operatorname{x\_}, \operatorname{y\_}] \in (\mathbb{R} \times \mathbb{R}).\; \cdot(\operatorname{x\_}, \operatorname{y\_})), (\lambda \operatorname{x\_} \in \mathbb{R}.\; -(\operatorname{x\_})), 0, 1, \operatorname{difference}(\mathbb{R}, \operatorname{singleton}(0)), (\lambda \operatorname{x\_} \in \operatorname{difference}(\mathbb{R}, \operatorname{singleton}(0)).\; \operatorname{recip}(\operatorname{x\_}))] \end{array}
\end{equation*}

\par\medskip\noindent\textbf{Definition} (\texttt{rr-fn-space}).\;$\operatorname{rr\text{-}fn\text{-}space}(\operatorname{rfv\_}) :=$
\begin{equation*}
{[}\operatorname{field\text{-}as\text{-}integral\text{-}domain}(\operatorname{rr\text{-}field}), \operatorname{rfv\_}, (\lambda {[}\operatorname{rff\_}, \operatorname{rfg\_}] \in (\operatorname{rfv\_} \times \operatorname{rfv\_}).\; (\lambda \operatorname{rfx\_} \in \mathbb{R}.\; (\operatorname{rff\_}(\operatorname{rfx\_}) + \operatorname{rfg\_}(\operatorname{rfx\_})))), (\lambda \operatorname{rfx\_} \in \mathbb{R}.\; 0), (\lambda \operatorname{rff\_} \in \operatorname{rfv\_}.\; (\lambda \operatorname{rfx\_} \in \mathbb{R}.\; -\operatorname{rff\_}(\operatorname{rfx\_}))), (\lambda {[}\operatorname{rfc\_}, \operatorname{rff\_}] \in (\mathbb{R} \times \operatorname{rfv\_}).\; (\lambda \operatorname{rfx\_} \in \mathbb{R}.\; (\operatorname{rfc\_} \cdot \operatorname{rff\_}(\operatorname{rfx\_}))))]
\end{equation*}

\par\medskip\noindent\textbf{Definition} (\texttt{antiderivable-fns}).\;$\operatorname{antiderivable\text{-}fns}(\operatorname{ada\_}, \operatorname{adb\_}) :=$ $\{\operatorname{adp\_} \in (\mathbb{R} \to \mathbb{R}) : \operatorname{is\text{-}antiderivable}(\operatorname{adp\_}, \operatorname{ada\_}, \operatorname{adb\_})\}$
\vnbdef{rr-field@non-zero}{\vnbstmt{(\operatorname{non\text{-}zero}(\operatorname{rr\text{-}field}) \simeq \operatorname{difference}(\mathbb{R}, \operatorname{singleton}(0)))}}
\vnbdef{rr-field@neg}{\vnbstmt{(\operatorname{neg}(\operatorname{rr\text{-}field}) \simeq (\lambda x \in \mathbb{R}.\; -x))}}
\vnbdef{rr-field@one}{\vnbstmt{(\operatorname{one}(\operatorname{rr\text{-}field}) \simeq 1)}}
\vnbdef{rr-field@carr}{\vnbstmt{(\operatorname{carr}(\operatorname{rr\text{-}field}) \simeq \mathbb{R})}}
\vnbdef{rr-field@zero}{\vnbstmt{(\operatorname{zero}(\operatorname{rr\text{-}field}) \simeq 0)}}
\vnbdef{rr-field@add}{\vnbstmt{(\operatorname{add}(\operatorname{rr\text{-}field}) \simeq (\lambda {[}\operatorname{x\_}, \operatorname{y\_}] \in (\mathbb{R} \times \mathbb{R}).\; (\operatorname{x\_} + \operatorname{y\_})))}}
\vnbdef{rr-field@mul}{\vnbstmt{(\operatorname{mul}(\operatorname{rr\text{-}field}) \simeq (\lambda {[}\operatorname{x\_}, \operatorname{y\_}] \in (\mathbb{R} \times \mathbb{R}).\; (\operatorname{x\_} \cdot \operatorname{y\_})))}}
\vnbdef{rr-field@mul-inv}{\vnbstmt{(\operatorname{mul\text{-}inv}(\operatorname{rr\text{-}field}) \simeq (\lambda x \in \operatorname{difference}(\mathbb{R}, \operatorname{singleton}(0)).\; \operatorname{recip}(x)))}}

\vnbbookfile{theorem-library/isom2-rrfield}{theorem-library/isom2-rrfield}
\noindent{\small isom2-rrfield.scm -- IS-FIELD(RR-FIELD) and the field ring of the reals: the scalars of the real vector spaces of isom2-rrvs.scm.  Agent ISOM, 2026-10-26. The proof is theorem-library/qq-field-is-field.scm (QQ-FIELD) with QQ replaced by RR and the qq- closure laws by their rr- twins; nothing else changes.}\par
\vnbprop{rr-recip-nonzero}{\vnbstmt{\forall a \in \mathbb{R}.\; (\neg (a = 0) \Rightarrow \neg (\operatorname{recip}(a) = 0))}}
\vnbprop{rr-field-mul-inv-type}{\vnbstmt{\begin{array}{@{}l@{}} (\lambda x \in \operatorname{difference}(\mathbb{R}, \operatorname{singleton}(0)).\; \operatorname{recip}(x)) \\ \in\; (\operatorname{difference}(\mathbb{R}, \operatorname{singleton}(0)) \to \operatorname{difference}(\mathbb{R}, \operatorname{singleton}(0))) \end{array}}}
\vnbprop{rr-field-is-field}{\vnbstmt{\operatorname{is\text{-}field}(\operatorname{rr\text{-}field})}}
\vnbprop{rr-field-is-field-ring}{\vnbstmt{\operatorname{is\text{-}field\text{-}ring}(\operatorname{field\text{-}as\text{-}integral\text{-}domain}(\operatorname{rr\text{-}field}))}}

\vnbbookfile{theorem-library/rake-esum-defined}{theorem-library/rake-esum-defined}
\noindent{\small theorem-library/rake-esum-defined.scm -- ESUM DEFINED as the ESUP of the set of finite partial sums, and its three characterising axioms PROVEN from the definition.  Rake batch 6-C, the ESUP half of which is theorem-library/rake-esup-defined.scm (agent 6-B).  Helper prefix `esm-'.}\par
\vnbprop{esum-partial-sum-type}{\vnbstmt{\begin{array}{@{}l@{}} \forall f, s. \\ \quad ((f \in (\operatorname{dom}(f) \to \operatorname{rr\text{-}pos\text{-}star})) \wedge ((s \in \mathbf{Set}) \wedge ((|s| \in \mathbb{N}) \wedge (s \subseteq \operatorname{dom}(f))))) \Rightarrow  \\ \quad \quad \operatorname{finsum}(\operatorname{rr\text{-}pos\text{-}star\text{-}add\text{-}monoid}, f, s) \in \operatorname{rr\text{-}pos\text{-}star} \end{array}}}
\vnbprop{esum-in}{\vnbstmt{\forall f \in (\operatorname{dom}(f) \to \operatorname{rr\text{-}pos\text{-}star}).\; (\operatorname{esum}(f) \in \operatorname{rr\text{-}pos\text{-}star})}}
\vnbprop{esum-upper}{\vnbstmt{\begin{array}{@{}l@{}} \forall f \in (\operatorname{dom}(f) \to \operatorname{rr\text{-}pos\text{-}star}), s \in \mathbf{Set}. \\ \quad ((|s| \in \mathbb{N}) \wedge (s \subseteq \operatorname{dom}(f))) \Rightarrow  \\ \quad \quad \operatorname{finsum}(\operatorname{rr\text{-}pos\text{-}star\text{-}add\text{-}monoid}, f, s) \leq \operatorname{esum}(f) \end{array}}}
\vnbprop{esum-least}{\vnbstmt{\begin{array}{@{}l@{}} \forall f \in (\operatorname{dom}(f) \to \operatorname{rr\text{-}pos\text{-}star}), b \in \operatorname{rr\text{-}pos\text{-}star}. \\ \quad \forall s \in \mathbf{Set}. \\ \quad \quad ((|s| \in \mathbb{N}) \wedge (s \subseteq \operatorname{dom}(f))) \Rightarrow  \\ \quad \quad \quad \operatorname{finsum}(\operatorname{rr\text{-}pos\text{-}star\text{-}add\text{-}monoid}, f, s) \leq b \Rightarrow \\ \quad \quad \operatorname{esum}(f) \leq b \end{array}}}

\vnbbookfile{theorem-library/rake-esum-finite}{theorem-library/rake-esum-finite}
\noindent{\small rake-esum-finite.scm -- ESUM: the finiteness dichotomy.}\par
\vnbprop{esum-bounded-implies-finite}{\vnbstmt{\begin{array}{@{}l@{}} \forall f \in (\operatorname{dom}(f) \to \operatorname{rr\text{-}pos\text{-}star}). \\ \quad \exists m \in \mathbb{R}. \forall s \in \mathbf{Set}. \\ \quad \quad ((|s| \in \mathbb{N}) \wedge (s \subseteq \operatorname{dom}(f))) \Rightarrow  \\ \quad \quad \quad \operatorname{finsum}(\operatorname{rr\text{-}pos\text{-}star\text{-}add\text{-}monoid}, f, s) \leq m \Rightarrow \\ \quad \quad \operatorname{esum}(f) \in \mathbb{R} \end{array}}}
\vnbprop{esum-finite-iff-bounded}{\vnbstmt{\begin{array}{@{}l@{}} \forall f \in (\operatorname{dom}(f) \to \operatorname{rr\text{-}pos\text{-}star}). \\ \quad \operatorname{esum}(f) \in \mathbb{R} \\ \quad \Leftrightarrow\; \exists m \in \mathbb{R}. \forall s \in \mathbf{Set}. \\ \quad \quad ((|s| \in \mathbb{N}) \wedge (s \subseteq \operatorname{dom}(f))) \Rightarrow  \\ \quad \quad \quad \operatorname{finsum}(\operatorname{rr\text{-}pos\text{-}star\text{-}add\text{-}monoid}, f, s) \leq m \end{array}}}

\vnbbookfile{theorem-library/ball-cover-lemmas}{theorem-library/ball-cover-lemmas}
\noindent{\small ball-cover-lemmas.scm -- the membership plumbing under finite-ball-subcover-r-net, PROVEN (2026-09-14).}\par
\vnbprop{is-compact-classes}{\vnbstmt{\begin{array}{@{}l@{}} \forall s. \\ \quad \operatorname{is\text{-}compact}(s) \\ \quad \Leftrightarrow\; \operatorname{is\text{-}metric\text{-}space}(s) \wedge \\ \quad \forall c. \\ \quad \quad \operatorname{is\text{-}open\text{-}cover}(s, c) \Rightarrow  \\ \quad \quad \quad \exists f.\; ((f \subseteq c) \wedge ((|f| \in \mathbb{N}) \wedge \operatorname{is\text{-}open\text{-}cover}(s, f))) \end{array}}}
\vnbprop{centres-unfold}{\vnbstmt{\forall s, b, r.\; (\operatorname{centres}(s, b, r) \simeq \{c \in \operatorname{pts}(s) : (\operatorname{ball}(s, c, r) = b)\})}}
\vnbprop{ball-cover-unfold}{\vnbstmt{\begin{array}{@{}l@{}} \forall s, r. \\ \quad (\operatorname{ball\text{-}cover}(s, r) \simeq \operatorname{image}((\lambda c \in \operatorname{pts}(s).\; \operatorname{ball}(s, c, r)), \operatorname{pts}(s))) \end{array}}}
\vnbprop{centre-set-unfold}{\vnbstmt{\begin{array}{@{}l@{}} \forall s, r, f. \\ \quad (\operatorname{centre\text{-}set}(s, r, f) \simeq \operatorname{image}((\lambda b \in f.\; \varepsilon(\operatorname{centres}(s, b, r))), f)) \end{array}}}
\vnbprop{ball-sep-unfold}{\vnbstmt{\begin{array}{@{}l@{}} \forall s, c, r. \\ \quad (\operatorname{ball}(s, c, r) \simeq \{y \in \operatorname{pts}(s) : ((\operatorname{dist}(s)(c, y) \leq r) \wedge \neg (\operatorname{dist}(s)(c, y) = r))\}) \end{array}}}
\vnbprop{centres-mem-build}{\vnbstmt{\forall s, b, r.; \forall c \in \operatorname{pts}(s).\; ((\operatorname{ball}(s, c, r) = b) \Rightarrow (c \in \operatorname{centres}(s, b, r)))}}
\vnbprop{centres-in-carrier}{\vnbstmt{\forall s, b, r.; \forall c \in \operatorname{centres}(s, b, r).\; (c \in \operatorname{pts}(s))}}
\vnbprop{centres-ball-eq}{\vnbstmt{\forall s, b, r.; \forall c \in \operatorname{centres}(s, b, r).\; (\operatorname{ball}(s, c, r) = b)}}
\vnbprop{ball-cover-mem-fwd}{\vnbstmt{\forall s, r.; \forall u \in \operatorname{ball\text{-}cover}(s, r).\; \exists c \in \operatorname{pts}(s).\; (\operatorname{ball}(s, c, r) = u)}}
\vnbprop{open-cover-covers-point}{\vnbstmt{\begin{array}{@{}l@{}} \forall s, f.; \forall p \in \operatorname{pts}(s). \\ \quad \operatorname{is\text{-}open\text{-}cover}(s, f) \Rightarrow \exists u \in f.\; (p \in u) \end{array}}}
\vnbprop{ball-point-le}{\vnbstmt{\forall s, c, r, w.; \forall p \in w.\; ((\operatorname{ball}(s, c, r) = w) \Rightarrow (\operatorname{dist}(s)(c, p) \leq r))}}
\vnbprop{ball-point-ne}{\vnbstmt{\forall s, c, r, w.; \forall p \in w.\; ((\operatorname{ball}(s, c, r) = w) \Rightarrow \neg (\operatorname{dist}(s)(c, p) = r))}}
\vnbprop{centre-set-finite-guarded}{\vnbstmt{\forall s, r.; \forall f \in \mathbf{Set}.\; ((|f| \in \mathbb{N}) \Rightarrow (|\operatorname{centre\text{-}set}(s, r, f)| \in \mathbb{N}))}}

\vnbbookfile{structure-library/metric-laws}{structure-library/metric-laws}
\noindent{\small metric-laws.scm -- the five metric laws, PROVEN from the definition.}\par
\vnbprop{metric-pos}{\vnbstmt{\begin{array}{@{}l@{}} \forall s.; \forall x, y \in \operatorname{pts}(s). \\ \quad \operatorname{is\text{-}metric\text{-}space}(s) \Rightarrow (0 \leq \operatorname{dist}(s)(x, y)) \end{array}}}
\vnbprop{metric-self-zero}{\vnbstmt{\forall s.; \forall x \in \operatorname{pts}(s).\; (\operatorname{is\text{-}metric\text{-}space}(s) \Rightarrow (\operatorname{dist}(s)(x, x) = 0))}}
\vnbprop{metric-zero-eq}{\vnbstmt{\begin{array}{@{}l@{}} \forall s.; \forall x, y \in \operatorname{pts}(s). \\ \quad (\operatorname{is\text{-}metric\text{-}space}(s) \wedge (\operatorname{dist}(s)(x, y) = 0)) \Rightarrow (x = y) \end{array}}}
\vnbprop{metric-sym}{\vnbstmt{\begin{array}{@{}l@{}} \forall s.; \forall x, y \in \operatorname{pts}(s). \\ \quad \operatorname{is\text{-}metric\text{-}space}(s) \Rightarrow (\operatorname{dist}(s)(x, y) = \operatorname{dist}(s)(y, x)) \end{array}}}
\vnbprop{metric-triangle}{\vnbstmt{\begin{array}{@{}l@{}} \forall s.; \forall x, y, z \in \operatorname{pts}(s). \\ \quad \operatorname{is\text{-}metric\text{-}space}(s) \Rightarrow (\operatorname{dist}(s)(x, z) \leq (\operatorname{dist}(s)(x, y) + \operatorname{dist}(s)(y, z))) \end{array}}}

\vnbbookfile{theorem-library/ccr1109a2-metric}{theorem-library/ccr1109a2-metric}
\noindent{\small ccr1109a2-metric.scm -- wave of 2026-11-09/10 (agent CCR-1109-A2), wired 2026-11-10 from scratchpad/ccr1109a2/ccr1109a2-metric.scm. Library: analysis.}\par
\vnbprop{metric-converges-to-cauchy}{\vnbstmt{\forall s, f, l.\; (\operatorname{converges\text{-}to}(s, f, l) \Rightarrow \operatorname{is\text{-}cauchy\text{-}seq}(s, f))}}
\vnbprop{metric-convergent-cauchy}{\vnbstmt{\forall s, f.\; (\operatorname{converges}(s, f) \Rightarrow \operatorname{is\text{-}cauchy\text{-}seq}(s, f))}}

\vnbbookfile{theorem-library/ball-is-open}{theorem-library/ball-is-open}
\noindent{\small ball-is-open.scm -- an open ball is an open set.}\par
\vnbprop{ball-is-open}{\vnbstmt{\begin{array}{@{}l@{}} \forall s.; \forall x \in \operatorname{pts}(s), r \in \mathbb{R}. \\ \quad (\operatorname{is\text{-}metric\text{-}space}(s) \wedge ((0 \leq r) \wedge \neg (0 = r))) \Rightarrow \operatorname{is\text{-}open}(s, \operatorname{ball}(s, x, r)) \end{array}}}

\vnbbookfile{theorem-library/rake-compact-complete}{theorem-library/rake-compact-complete}
\noindent{\small rake-compact-complete.scm -- two of the three asserted leaves on the bill of compact-implies-complete (calculus/compact-complete-proof.scm), PROVEN.}\par
\vnbprop{cauchy-seq-is-fun}{\vnbstmt{\forall s, f.\; (\operatorname{is\text{-}cauchy\text{-}seq}(s, f) \Rightarrow (f \in (\mathbb{N} \to \operatorname{pts}(s))))}}
\vnbprop{cauchy-cluster-converges}{\vnbstmt{\begin{array}{@{}l@{}} \forall s, f. \\ \quad (\operatorname{is\text{-}cauchy\text{-}seq}(s, f) \wedge \exists x.\; \operatorname{cluster\text{-}point}(s, f, x)) \Rightarrow  \\ \quad \quad \operatorname{converges}(s, f) \end{array}}}

\vnbbookfile{theorem-library/rake-compact-cluster}{theorem-library/rake-compact-cluster}
\noindent{\small rake-compact-cluster.scm -- a compact metric space is sequentially compact.}\par
\vnbprop{compact-seq-has-cluster}{\vnbstmt{\begin{array}{@{}l@{}} \forall s, f. \\ \quad (\operatorname{is\text{-}compact}(s) \wedge (f \in (\mathbb{N} \to \operatorname{pts}(s)))) \Rightarrow  \\ \quad \quad \exists x.\; \operatorname{cluster\text{-}point}(s, f, x) \end{array}}}

\vnbbookfile{theorem-library/rake-balls}{theorem-library/rake-balls}
\noindent{\small rake-balls.scm -- the open-ball read-offs of structure-library/metric-topology.scm, PROVEN.  Five facts that stood there as `support' + `warrant! ... 'proof' -- a tier that claims a machine-checked proof exists -- and had none:}\par
\vnbprop{ball-membership}{\vnbstmt{\begin{array}{@{}l@{}} \forall s, x, r, y. \\ \quad (y \in \operatorname{ball}(s, x, r)) \Leftrightarrow ((y \in \operatorname{pts}(s)) \wedge ((\operatorname{dist}(s)(x, y) \leq r) \wedge \neg (\operatorname{dist}(s)(x, y) = r))) \end{array}}}
\vnbprop{ball-mem-from-le}{\vnbstmt{\begin{array}{@{}l@{}} \forall s.; \forall x \in \operatorname{pts}(s), y, d, r. \\ \quad (\operatorname{is\text{-}metric\text{-}space}(s) \wedge ((y \in \operatorname{pts}(s)) \wedge ((d \in \mathbb{R}) \wedge ((r \in \mathbb{R}) \wedge ((\operatorname{dist}(s)(x, y) \leq d) \wedge (d < r)))))) \Rightarrow  \\ \quad \quad y \in \operatorname{ball}(s, x, r) \end{array}}}
\vnbprop{ball-is-set}{\vnbstmt{\forall s, x, r.\; (\operatorname{is\text{-}metric\text{-}space}(s) \Rightarrow (\operatorname{ball}(s, x, r) \in \mathbf{Set}))}}
\vnbprop{ball-center-in}{\vnbstmt{\begin{array}{@{}l@{}} \forall s.; \forall x \in \operatorname{pts}(s), r \in \mathbb{R}. \\ \quad (\operatorname{is\text{-}metric\text{-}space}(s) \wedge ((0 \leq r) \wedge \neg (0 = r))) \Rightarrow (x \in \operatorname{ball}(s, x, r)) \end{array}}}
\vnbprop{ball-2r-triangle}{\vnbstmt{\begin{array}{@{}l@{}} \forall s.; \forall x \in \operatorname{pts}(s), r \in \mathbb{R}, y, z \in \operatorname{ball}(s, x, r). \\ \quad (\operatorname{is\text{-}metric\text{-}space}(s) \wedge ((0 \leq r) \wedge \neg (0 = r))) \Rightarrow  \\ \quad \quad (\operatorname{dist}(s)(y, z) \leq (r + r)) \wedge \neg (\operatorname{dist}(s)(y, z) = (r + r)) \end{array}}}
\vnbprop{cauchy-block-estimate}{\vnbstmt{\begin{array}{@{}l@{}} \forall s.; \forall c \in \operatorname{pts}(s), r, u.; \forall y, z \in u, d.; \forall a \in \mathbb{R}. \\ \quad (\operatorname{is\text{-}metric\text{-}space}(s) \wedge (\operatorname{pos\text{-}rr}(r) \wedge ((u = \operatorname{ball}(s, c, r)) \wedge (\operatorname{pos\text{-}rr}(d) \wedge ((r \leq d) \wedge ((d + d) = a)))))) \Rightarrow  \\ \quad \quad \operatorname{dist}(s)(y, z) \leq a \end{array}}}

\vnbbookfile{theorem-library/metric-connected-laws}{theorem-library/metric-connected-laws}
\noindent{\small metric-connected-laws.scm -- A LOCALLY CONSTANT FUNCTION ON A CONNECTED SET IS CONSTANT.  Agent CA-5 of week 3 of the October roadmap, 2026-10-04. Definition: structure-library/metric-connected.scm (IS-CONNECTED, new this batch).}\par
\vnbprop{connected-locally-constant}{\vnbstmt{\begin{array}{@{}l@{}} \forall s, a, f.; \forall p, q \in a. \\ \quad \operatorname{is\text{-}connected}(s, a) \wedge \\ \quad \forall x \in a. \exists r. \\ \quad \quad \operatorname{pos\text{-}rr}(r) \wedge \forall y \in a.\; ((y \in \operatorname{ball}(s, x, r)) \Rightarrow (f(y) = f(x))) \Rightarrow \\ \quad \quad f(p) = f(q) \end{array}}}

\vnbbookfile{theorem-library/metric-subspace-laws}{theorem-library/metric-subspace-laws}
\noindent{\small metric-subspace-laws.scm -- the laws of SUBSPACE-MS and RESTRICT (structure-library/metric-subspace.scm).  Everything here is PROVEN; the definition file asserts nothing.}\par
\vnbprop{subspace-pts}{\vnbstmt{\forall s.; \forall a \in \mathbf{Set}.\; (\operatorname{pts}(\operatorname{subspace\text{-}ms}(s, a)) \simeq a)}}
\vnbprop{subspace-dist-slot}{\vnbstmt{\begin{array}{@{}l@{}} \forall s.; \forall a \in \mathbf{Set}. \\ \quad (\operatorname{dist}(\operatorname{subspace\text{-}ms}(s, a)) \simeq (\lambda {[}\operatorname{sbu\_}, \operatorname{sbv\_}] \in (a \times a).\; \operatorname{dist}(s)(\operatorname{sbu\_}, \operatorname{sbv\_}))) \end{array}}}
\vnbprop{subspace-dist}{\vnbstmt{\begin{array}{@{}l@{}} \forall s.; \forall a \in \mathbf{Set}, u, v. \\ \quad ((u \in a) \wedge (v \in a)) \Rightarrow (\operatorname{dist}(\operatorname{subspace\text{-}ms}(s, a))(u, v) \simeq \operatorname{dist}(s)(u, v)) \end{array}}}
\vnbprop{subspace-is-metric-space}{\vnbstmt{\begin{array}{@{}l@{}} \forall s, a. \\ \quad (\operatorname{is\text{-}metric\text{-}space}(s) \wedge (a \subseteq \operatorname{pts}(s))) \Rightarrow  \\ \quad \quad \operatorname{is\text{-}metric\text{-}space}(\operatorname{subspace\text{-}ms}(s, a)) \end{array}}}
\vnbprop{restrict-apply}{\vnbstmt{\forall f, a.; \forall u \in a.\; (\operatorname{restrict}(f, a)(u) \simeq f(u))}}
\vnbprop{restrict-in-fun}{\vnbstmt{\begin{array}{@{}l@{}} \forall f, m, d, a. \\ \quad ((f \in (m \to d)) \wedge (a \subseteq m)) \Rightarrow (\operatorname{restrict}(f, a) \in (a \to d)) \end{array}}}
\vnbprop{subspace-whole}{\vnbstmt{\begin{array}{@{}l@{}} \forall s. \\ \quad \operatorname{is\text{-}metric\text{-}space}(s) \Rightarrow  \\ \quad \quad (\operatorname{pts}(\operatorname{subspace\text{-}ms}(s, \operatorname{pts}(s))) \simeq \operatorname{pts}(s)) \wedge \\ \quad \quad \forall u, v. \\ \quad \quad \quad ((u \in \operatorname{pts}(s)) \wedge (v \in \operatorname{pts}(s))) \Rightarrow  \\ \quad \quad \quad \quad (\operatorname{dist}(\operatorname{subspace\text{-}ms}(s, \operatorname{pts}(s)))(u, v) \simeq \operatorname{dist}(s)(u, v)) \end{array}}}
\vnbprop{subspace-ball-membership}{\vnbstmt{\begin{array}{@{}l@{}} \forall s, a.; \forall c \in a, r, y. \\ \quad (\operatorname{is\text{-}metric\text{-}space}(s) \wedge (a \subseteq \operatorname{pts}(s))) \Rightarrow  \\ \quad \quad (y \in \operatorname{ball}(\operatorname{subspace\text{-}ms}(s, a), c, r)) \Leftrightarrow ((y \in \operatorname{ball}(s, c, r)) \wedge (y \in a)) \end{array}}}
\vnbprop{subspace-ball}{\vnbstmt{\begin{array}{@{}l@{}} \forall s, a.; \forall c \in a, r. \\ \quad (\operatorname{is\text{-}metric\text{-}space}(s) \wedge (a \subseteq \operatorname{pts}(s))) \Rightarrow  \\ \quad \quad \operatorname{ball}(\operatorname{subspace\text{-}ms}(s, a), c, r) = (\operatorname{ball}(s, c, r) \cap a) \end{array}}}
\vnbprop{subspace-open-trace}{\vnbstmt{\begin{array}{@{}l@{}} \forall s, a, u. \\ \quad (\operatorname{is\text{-}metric\text{-}space}(s) \wedge ((a \subseteq \operatorname{pts}(s)) \wedge \operatorname{is\text{-}open}(s, u))) \Rightarrow  \\ \quad \quad \operatorname{is\text{-}open}(\operatorname{subspace\text{-}ms}(s, a), (u \cap a)) \end{array}}}
\vnbprop{restrict-continuous-at}{\vnbstmt{\begin{array}{@{}l@{}} \forall s, a, t, f.; \forall \mathit{pa} \in a. \\ \quad (\operatorname{is\text{-}metric\text{-}space}(s) \wedge ((a \subseteq \operatorname{pts}(s)) \wedge \operatorname{is\text{-}continuous\text{-}at}(s, t, f, \mathit{pa}))) \Rightarrow  \\ \quad \quad \operatorname{is\text{-}continuous\text{-}at}(\operatorname{subspace\text{-}ms}(s, a), t, \operatorname{restrict}(f, a), \mathit{pa}) \end{array}}}
\vnbprop{subspace-converges-to}{\vnbstmt{\begin{array}{@{}l@{}} \forall s, a.; \forall g \in (\mathbb{N} \to a), l \in a. \\ \quad (\operatorname{is\text{-}metric\text{-}space}(s) \wedge (a \subseteq \operatorname{pts}(s))) \Rightarrow  \\ \quad \quad \operatorname{converges\text{-}to}(\operatorname{subspace\text{-}ms}(s, a), g, l) \Leftrightarrow \operatorname{converges\text{-}to}(s, g, l) \end{array}}}
\vnbprop{subspace-open-is-trace}{\vnbstmt{\begin{array}{@{}l@{}} \forall s, a, \mathit{vs}. \\ \quad (\operatorname{is\text{-}metric\text{-}space}(s) \wedge ((a \subseteq \operatorname{pts}(s)) \wedge \operatorname{is\text{-}open}(\operatorname{subspace\text{-}ms}(s, a), \mathit{vs}))) \Rightarrow  \\ \quad \quad \exists \mathit{us}.\; (\operatorname{is\text{-}open}(s, \mathit{us}) \wedge (\mathit{vs} = (\mathit{us} \cap a))) \end{array}}}
\vnbprop{closed-limit-in}{\vnbstmt{\begin{array}{@{}l@{}} \forall s, a.; \forall f \in (\mathbb{N} \to a), l. \\ \quad (\operatorname{is\text{-}metric\text{-}space}(s) \wedge (\operatorname{is\text{-}closed}(s, a) \wedge \operatorname{converges\text{-}to}(s, f, l))) \Rightarrow (l \in a) \end{array}}}
\vnbprop{subspace-is-cauchy-seq}{\vnbstmt{\begin{array}{@{}l@{}} \forall s, a.; \forall f \in (\mathbb{N} \to a). \\ \quad (\operatorname{is\text{-}metric\text{-}space}(s) \wedge (a \subseteq \operatorname{pts}(s))) \Rightarrow  \\ \quad \quad \operatorname{is\text{-}cauchy\text{-}seq}(\operatorname{subspace\text{-}ms}(s, a), f) \Leftrightarrow \operatorname{is\text{-}cauchy\text{-}seq}(s, f) \end{array}}}
\vnbprop{subspace-complete}{\vnbstmt{\begin{array}{@{}l@{}} \forall s, a. \\ \quad (\operatorname{is\text{-}complete}(s) \wedge \operatorname{is\text{-}closed}(s, a)) \Rightarrow \operatorname{is\text{-}complete}(\operatorname{subspace\text{-}ms}(s, a)) \end{array}}}

\vnbbookfile{theorem-library/compact-subspace}{theorem-library/compact-subspace}
\noindent{\small compact-subspace.scm -- a CLOSED SUBSET of a COMPACT space is COMPACT as a subspace (structure-library/metric-subspace.scm, theorem-library/metric-subspace-laws.scm).}\par
\vnbprop{open-in-power}{\vnbstmt{\forall s, u.\; ((\operatorname{is\text{-}metric\text{-}space}(s) \wedge \operatorname{is\text{-}open}(s, u)) \Rightarrow (u \in \operatorname{power}(\operatorname{pts}(s))))}}
\vnbprop{closed-subset-of-compact-is-compact}{\vnbstmt{\begin{array}{@{}l@{}} \forall s, a. \\ \quad (\operatorname{is\text{-}compact}(s) \wedge \operatorname{is\text{-}closed}(s, a)) \Rightarrow \operatorname{is\text{-}compact}(\operatorname{subspace\text{-}ms}(s, a)) \end{array}}}

\vnbbookfile{theorem-library/rake-analysis-typing}{theorem-library/rake-analysis-typing}
\noindent{\small rake-analysis-typing.scm -- BATCH H of the leaf rake (2026-09-17): the four ANALYSIS typing leaves.  Two are proven as stated; two are UNDERDETERMINED as stated -- the same defect that made `nth-deriv-v-in-vec' false -- and are proven here with the guard their scalar twins already carry.}\par
\vnbprop{rkt-cseq-is-set}{\vnbstmt{\forall m.\; (\operatorname{is\text{-}metric\text{-}space}(m) \Rightarrow (\operatorname{cseq}(m) \in \mathbf{Set}))}}
\vnbprop{rkt-completion-pts}{\vnbstmt{\begin{array}{@{}l@{}} \forall m. \\ \quad \operatorname{is\text{-}metric\text{-}space}(m) \Rightarrow (\operatorname{pts}(\operatorname{completion}(m)) \simeq \operatorname{quotient}(\operatorname{cauchy\text{-}setoid}(m))) \end{array}}}
\vnbprop{rkt-cauchy-setoid-pts}{\vnbstmt{\forall m.\; (\operatorname{is\text{-}metric\text{-}space}(m) \Rightarrow (\operatorname{pts}(\operatorname{cauchy\text{-}setoid}(m)) \simeq \operatorname{cseq}(m)))}}
\vnbprop{rkt-class-is-set}{\vnbstmt{\forall m, w.\; (\operatorname{is\text{-}metric\text{-}space}(m) \Rightarrow (\operatorname{class}(\operatorname{cauchy\text{-}setoid}(m), w) \in \mathbf{Set}))}}
\vnbprop{rkt-embed-seq-in-fun}{\vnbstmt{\begin{array}{@{}l@{}} \forall m, u. \\ \quad (\operatorname{is\text{-}metric\text{-}space}(m) \wedge (u \in \operatorname{pts}(m))) \Rightarrow (\operatorname{embed\text{-}seq}(m, u) \in (\mathbb{N} \to \operatorname{pts}(m))) \end{array}}}
\vnbprop{rkt-embed-seq-cauchy}{\vnbstmt{\begin{array}{@{}l@{}} \forall m, u. \\ \quad (\operatorname{is\text{-}metric\text{-}space}(m) \wedge (u \in \operatorname{pts}(m))) \Rightarrow \operatorname{is\text{-}cauchy\text{-}seq}(m, \operatorname{embed\text{-}seq}(m, u)) \end{array}}}
\vnbprop{rkt-embed-seq-in-cseq}{\vnbstmt{\begin{array}{@{}l@{}} \forall m, u. \\ \quad (\operatorname{is\text{-}metric\text{-}space}(m) \wedge (u \in \operatorname{pts}(m))) \Rightarrow (\operatorname{embed\text{-}seq}(m, u) \in \operatorname{cseq}(m)) \end{array}}}
\vnbprop{embed-in-fun}{\vnbstmt{\begin{array}{@{}l@{}} \forall m. \\ \quad \operatorname{is\text{-}metric\text{-}space}(m) \Rightarrow (\operatorname{embed}(m) \in (\operatorname{pts}(m) \to \operatorname{pts}(\operatorname{completion}(m)))) \end{array}}}

\vnbbookfile{theorem-library/rake-cont-preimage}{theorem-library/rake-cont-preimage}
\noindent{\small rake-cont-preimage.scm -- continuous-implies-open-preimage, PROVEN.}\par
\vnbprop{continuous-implies-open-preimage}{\vnbstmt{\begin{array}{@{}l@{}} \forall s, t, f, v. \\ \quad (\operatorname{is\text{-}continuous}(s, t, f) \wedge \operatorname{is\text{-}open}(t, v)) \Rightarrow \operatorname{is\text{-}open}(s, \operatorname{preimage}(s, f, v)) \end{array}}}

\vnbbookfile{theorem-library/pseudometric-laws}{theorem-library/pseudometric-laws}
\noindent{\small pseudometric-laws.scm -- the FOUR PSEUDOMETRIC LAWS, proven by projecting the `is-pseudometric' property, exactly as structure-library/metric-laws.scm projects `is-metric'.}\par
\vnbprop{pseudometric-self-zero}{\vnbstmt{\forall t, r.; \forall u \in r.\; (\operatorname{is\text{-}pseudometric}(t, r) \Rightarrow (t(u, u) = 0))}}
\vnbprop{pseudometric-pos}{\vnbstmt{\forall t, r.; \forall u, v \in r.\; (\operatorname{is\text{-}pseudometric}(t, r) \Rightarrow (0 \leq t(u, v)))}}
\vnbprop{pseudometric-sym}{\vnbstmt{\forall t, r.; \forall u, v \in r.\; (\operatorname{is\text{-}pseudometric}(t, r) \Rightarrow (t(u, v) = t(v, u)))}}
\vnbprop{pseudometric-triangle}{\vnbstmt{\begin{array}{@{}l@{}} \forall t, r.; \forall u, v, w \in r. \\ \quad \operatorname{is\text{-}pseudometric}(t, r) \Rightarrow (t(u, w) \leq (t(u, v) + t(v, w))) \end{array}}}

\vnbbookfile{theorem-library/bdd-metric-distance}{theorem-library/bdd-metric-distance}
\noindent{\small bdd-metric-distance.scm -- the BDD-METRIC distance formula, proven.}\par
\vnbprop{bdd-metric-distance}{\vnbstmt{\begin{array}{@{}l@{}} \forall s.; \forall x, y \in \operatorname{pts}(s). \\ \quad \operatorname{is\text{-}metric\text{-}space}(s) \Rightarrow  \\ \quad \quad \operatorname{dist}(\operatorname{bdd\text{-}metric}(s))(x, y) = /(\operatorname{dist}(s)(x, y), (1 + \operatorname{dist}(s)(x, y))) \end{array}}}

\vnbbookfile{theorem-library/bdd-metric-basics}{theorem-library/bdd-metric-basics}
\noindent{\small bdd-metric-basics.scm -- the t/(1+t) scalar family and the two BDD-METRIC statements that read off it, PROVEN.}\par
\vnbprop{bdd-fn-nonneg}{\vnbstmt{\forall t \in \mathbb{R}.\; ((0 \leq t) \Rightarrow (0 \leq /(t, (1 + t))))}}
\vnbprop{bdd-fn-lt-one}{\vnbstmt{\forall t \in \mathbb{R}.\; ((0 \leq t) \Rightarrow (/(t, (1 + t)) < 1))}}
\vnbprop{bdd-fn-le-arg}{\vnbstmt{\forall t \in \mathbb{R}.\; ((0 \leq t) \Rightarrow (/(t, (1 + t)) \leq t))}}
\vnbprop{bdd-fn-zero}{\vnbstmt{\forall t \in \mathbb{R}.\; (((0 \leq t) \wedge (t = 0)) \Rightarrow (/(t, (1 + t)) = 0))}}
\vnbprop{bdd-fn-zero-eq}{\vnbstmt{\forall t \in \mathbb{R}.\; (((0 \leq t) \wedge (/(t, (1 + t)) = 0)) \Rightarrow (t = 0))}}
\vnbprop{bdd-fn-mono}{\vnbstmt{\begin{array}{@{}l@{}} \forall s, t \in \mathbb{R}. \\ \quad ((0 \leq s) \wedge (s \leq t)) \Rightarrow (/(s, (1 + s)) \leq /(t, (1 + t))) \end{array}}}
\vnbprop{bdd-fn-subadd}{\vnbstmt{\begin{array}{@{}l@{}} \forall a, b \in \mathbb{R}. \\ \quad ((0 \leq a) \wedge (0 \leq b)) \Rightarrow  \\ \quad \quad /((a + b), (1 + (a + b))) \leq (/(a, (1 + a)) + /(b, (1 + b))) \end{array}}}
\vnbprop{bdd-metric-bounded}{\vnbstmt{\begin{array}{@{}l@{}} \forall s.; \forall x, y \in \operatorname{pts}(s). \\ \quad \operatorname{is\text{-}metric\text{-}space}(s) \Rightarrow (\operatorname{dist}(\operatorname{bdd\text{-}metric}(s))(x, y) < 1) \end{array}}}
\vnbprop{bdd-metric-is-metric-space}{\vnbstmt{\forall s.\; (\operatorname{is\text{-}metric\text{-}space}(s) \Rightarrow \operatorname{is\text{-}metric\text{-}space}(\operatorname{bdd\text{-}metric}(s)))}}

\vnbbookfile{theorem-library/rake-subseq-leaves}{theorem-library/rake-subseq-leaves}
\noindent{\small rake-subseq-leaves.scm -- the remaining asserted leaves of the subsequence / diagonalization cluster, PROVEN.}\par
\vnbprop{nn-step-mono-ptwise}{\vnbstmt{\begin{array}{@{}l@{}} \forall n \in \mathbb{N}, g.; \forall m \in \mathbb{N}. \\ \quad (\forall k \in \mathbb{N}.\; (g(k) \in \mathbb{N}) \wedge (\forall k \in \mathbb{N}.\; (g(k) < g((k + 1))) \wedge (m < n))) \Rightarrow  \\ \quad \quad g(m) < g(n) \end{array}}}
\vnbprop{nn-step-strictly-mono}{\vnbstmt{\begin{array}{@{}l@{}} \forall g \in (\mathbb{N} \to \mathbb{N}). \\ \quad \forall k \in \mathbb{N}.\; (g(k) < g((k + 1))) \Rightarrow \operatorname{strictly\text{-}mono\text{-}nn}(g) \end{array}}}
\vnbprop{subsequence-capture}{\vnbstmt{\begin{array}{@{}l@{}} \forall s \in \mathrm{Inf}(\mathbb{N}). \exists f \in (\mathbb{N} \to s). \forall m, n \in \mathbb{N}. \\ \quad (m < n) \Rightarrow (f(m) < f(n)) \end{array}}}
\vnbprop{nn-recip-succ-small-at}{\vnbstmt{\forall s.; \forall k \in \mathbb{N}.\; ((\operatorname{pos\text{-}rr}(s) \wedge (\operatorname{recip}(s) < (k + 1))) \Rightarrow (\operatorname{recip}((k + 1)) < s))}}
\vnbprop{null-rr-seq-exists}{\vnbstmt{\exists d.\; \operatorname{null\text{-}rr\text{-}seq}(d)}}

\vnbbookfile{theorem-library/rake-dc-consumers}{theorem-library/rake-dc-consumers}
\noindent{\small rake-dc-consumers.scm -- the consumers of recursive choice (rake batch 8-C).}\par
\vnbprop{cauchy-rapid-subsequence}{\vnbstmt{\begin{array}{@{}l@{}} \forall s, f, d. \\ \quad (\operatorname{is\text{-}cauchy\text{-}seq}(s, f) \wedge ((d \in (\mathbb{N} \to \mathbb{R})) \wedge \forall k \in \mathbb{N}.\; \operatorname{pos\text{-}rr}(d(k)))) \Rightarrow  \\ \quad \quad \exists i \in (\mathbb{N} \to \mathbb{N}). \\ \quad \quad \quad \forall m, n \in \mathbb{N}.\; ((m < n) \Rightarrow (i(m) < i(n))) \wedge \\ \quad \quad \quad \forall k \in \mathbb{N}.\; (\operatorname{dist}(s)(f(i(k)), f(i((k + 1)))) \leq d(k)) \end{array}}}

\vnbbookfile{theorem-library/rake-compact-iff-seq-compact}{theorem-library/rake-compact-iff-seq-compact}
\noindent{\small rake-compact-iff-seq-compact.scm -- the FORWARD half of Prop 3.12 (1)=$>$(3'): a compact metric space is sequentially compact.}\par
\vnbprop{cluster-point-has-convergent-subseq}{\vnbstmt{\begin{array}{@{}l@{}} \forall s, f, x. \\ \quad \operatorname{cluster\text{-}point}(s, f, x) \Rightarrow  \\ \quad \quad \exists i.\; (\operatorname{strictly\text{-}mono\text{-}nn}(i) \wedge \operatorname{converges\text{-}to}(s, \operatorname{subseq}(f, i), x)) \end{array}}}
\vnbprop{compact-implies-seq-compact}{\vnbstmt{\forall s.\; ((\operatorname{is\text{-}metric\text{-}space}(s) \wedge \operatorname{is\text{-}compact}(s)) \Rightarrow \operatorname{seq\text{-}compact}(s))}}

\vnbbookfile{theorem-library/rake-seq-compact-tb}{theorem-library/rake-seq-compact-tb}
\noindent{\small rake-seq-compact-tb.scm -- a sequentially compact metric space is totally bounded (the (C1) leg of Prop 3.12's missing direction; see the closing block of theorem-library/rake-compact-iff-seq-compact.scm).}\par
\vnbprop{seq-compact-implies-totally-bounded}{\vnbstmt{\forall s.\; ((\operatorname{is\text{-}metric\text{-}space}(s) \wedge \operatorname{seq\text{-}compact}(s)) \Rightarrow \operatorname{totally\text{-}bounded}(s))}}

\vnbbookfile{theorem-library/rake-lebesgue-number}{theorem-library/rake-lebesgue-number}
\noindent{\small rake-lebesgue-number.scm -- the Lebesgue number lemma, and the CONVERSE half of Prop 3.12 (1)$<$=$>$(3'): a sequentially compact metric space is compact.}\par
\vnbprop{lebesgue-number}{\vnbstmt{\begin{array}{@{}l@{}} \forall s, v. \\ \quad (\operatorname{is\text{-}metric\text{-}space}(s) \wedge (\operatorname{seq\text{-}compact}(s) \wedge \operatorname{is\text{-}open\text{-}cover}(s, v))) \Rightarrow  \\ \quad \quad \exists l. \\ \quad \quad \quad \operatorname{pos\text{-}rr}(l) \wedge \forall p \in \operatorname{pts}(s).\; \exists u \in v.\; (\operatorname{ball}(s, p, l) \subseteq u) \end{array}}}
\vnbprop{seq-compact-implies-compact}{\vnbstmt{\forall s.\; ((\operatorname{is\text{-}metric\text{-}space}(s) \wedge \operatorname{seq\text{-}compact}(s)) \Rightarrow \operatorname{is\text{-}compact}(s))}}
\vnbprop{compact-iff-seq-compact}{\vnbstmt{\forall s.\; (\operatorname{is\text{-}metric\text{-}space}(s) \Rightarrow (\operatorname{is\text{-}compact}(s) \Leftrightarrow \operatorname{seq\text{-}compact}(s)))}}

\vnbbookfile{theorem-library/rake-nn-enum}{theorem-library/rake-nn-enum}
\noindent{\small theorem-library/rake-nn-enum.scm ==================================================================== The chosen enumeration NN-ENUM and the two strictly-monotone bricks the refinement tower (convergence-block-tower) wants.}\par
\vnbprop{nn-enum-spec}{\vnbstmt{\begin{array}{@{}l@{}} \forall s \in \mathrm{Inf}(\mathbb{N}). \\ \quad \operatorname{nn\text{-}enum}(s) \in (\mathbb{N} \to s) \wedge \\ \quad \forall m, n \in \mathbb{N}.\; ((m < n) \Rightarrow (\operatorname{nn\text{-}enum}(s)(m) < \operatorname{nn\text{-}enum}(s)(n))) \end{array}}}
\vnbprop{strictly-mono-image-infinite}{\vnbstmt{\begin{array}{@{}l@{}} \forall i, a. \\ \quad (\operatorname{strictly\text{-}mono\text{-}nn}(i) \wedge ((a \subseteq \mathbb{N}) \wedge \forall k \in \mathbb{N}.\; (i(k) \in a))) \Rightarrow  \\ \quad \quad \{m \in a : \exists j \in \mathbb{N}.\; (m = i(j))\} \in \mathrm{Inf}(\mathbb{N}) \end{array}}}
\vnbprop{strictly-mono-le-reflect}{\vnbstmt{\begin{array}{@{}l@{}} \forall i.; \forall a, j \in \mathbb{N}. \\ \quad (\operatorname{strictly\text{-}mono\text{-}nn}(i) \wedge (i(a) \leq i(j))) \Rightarrow (a \leq j) \end{array}}}

\vnbbookfile{theorem-library/rake-tb-leaves}{theorem-library/rake-tb-leaves}
\noindent{\small rake-tb-leaves.scm -- the totally-bounded / pigeonhole leaves of the Cauchy-subsequence and separability arcs.  Two theorems, both modulo 0:}\par
\vnbprop{cover-block-step}{\vnbstmt{\begin{array}{@{}l@{}} \forall v \in \mathbf{Set}, f \in (\mathbb{N} \to v), c.; \forall j \in \mathrm{Inf}(\mathbb{N}). \\ \quad \operatorname{is\text{-}finite\text{-}cover}(c, v) \Rightarrow  \\ \quad \quad \exists a \in \mathrm{Inf}(\mathbb{N}). \\ \quad \quad \quad (a \subseteq j) \wedge \exists u \in c.\; \forall i \in a.\; (f(i) \in u) \end{array}}}
\vnbprop{dist-le-implies-in-carrier}{\vnbstmt{\begin{array}{@{}l@{}} \forall s, c, p, r. \\ \quad (\operatorname{is\text{-}metric\text{-}space}(s) \wedge (\operatorname{dist}(s)(c, p) \leq r)) \Rightarrow (c \in \operatorname{pts}(s)) \end{array}}}

\vnbbookfile{theorem-library/rake-tb-leaves-2}{theorem-library/rake-tb-leaves-2}
\noindent{\small rake-tb-leaves-2.scm -- the two totally-bounded leaves that the IS-R-NET repair of 2026-09-19 unblocked.}\par
\vnbprop{tb-scale-dense-seq}{\vnbstmt{\begin{array}{@{}l@{}} \forall s.; \forall a \in \operatorname{pts}(s), r. \\ \quad (\operatorname{totally\text{-}bounded}(s) \wedge \operatorname{pos\text{-}rr}(r)) \Rightarrow  \\ \quad \quad \exists g \in (\mathbb{N} \to \operatorname{pts}(s)). \forall p \in \operatorname{pts}(s). \exists j \in \mathbb{N}. \\ \quad \quad \quad \operatorname{dist}(s)(p, g(j)) < r \end{array}}}
\vnbprop{tb-rad-ball-cover}{\vnbstmt{\begin{array}{@{}l@{}} \forall s, d. \\ \quad (\operatorname{totally\text{-}bounded}(s) \wedge \forall k \in \mathbb{N}.\; \operatorname{pos\text{-}rr}(d(k))) \Rightarrow  \\ \quad \quad \operatorname{pts}(s) \in \mathbf{Set} \wedge \\ \quad \quad \exists v. \\ \quad \quad \quad \forall k \in \mathbb{N}.\; \operatorname{is\text{-}finite\text{-}cover}(v(k), \operatorname{pts}(s)) \wedge \\ \quad \quad \quad \forall k \in \mathbb{N}, u \in v(k).\; \exists c \in \operatorname{pts}(s).\; (u = \operatorname{ball}(s, c, d(k))) \end{array}}}

\vnbbookfile{theorem-library/diagonalization}{theorem-library/diagonalization}
\noindent{\small theorem-library/diagonalization.scm}\par
\vnbprop{diagonalization}{\vnbstmt{\begin{array}{@{}l@{}} \forall s \in (\mathbb{N} \to \mathrm{Inf}(\mathbb{N})). \\ \quad \forall k \in \mathbb{N}.\; (s((k + 1)) \subseteq s(k)) \Rightarrow  \\ \quad \quad \exists f. \\ \quad \quad \quad \operatorname{strictly\text{-}mono\text{-}nn}(f) \wedge \forall k, j \in \mathbb{N}.\; ((k \leq j) \Rightarrow (f(j) \in s(k))) \end{array}}}

\vnbbookfile{theorem-library/block-family-combinatorial-proof}{theorem-library/block-family-combinatorial-proof}
\noindent{\small theorem-library/block-family-combinatorial-proof.scm}\par
\vnbprop{block-family-combinatorial}{\vnbstmt{\begin{array}{@{}l@{}} \forall v \in \mathbf{Set}, f \in (\mathbb{N} \to v), a. \\ \quad \forall k \in \mathbb{N}.\; \operatorname{is\text{-}finite\text{-}cover}(a(k), v) \Rightarrow  \\ \quad \quad \exists b \in (\mathbb{N} \to \mathrm{Inf}(\mathbb{N})). \\ \quad \quad \quad \forall k \in \mathbb{N}.\; (b((k + 1)) \subseteq b(k)) \wedge \\ \quad \quad \quad \forall k \in \mathbb{N}.\; \exists u \in a(k).\; \forall i \in b(k).\; (f(i) \in u) \end{array}}}

\vnbbookfile{theorem-library/cauchy-subseq-proof}{theorem-library/cauchy-subseq-proof}
\noindent{\small theorem-library/cauchy-subseq-proof.scm ==================================================================== totally-bounded-has-cauchy-subsequence -- PROVEN (was asserted).}\par
\vnbprop{totally-bounded-has-cauchy-subsequence}{\vnbstmt{\begin{array}{@{}l@{}} \forall s.; \forall f \in (\mathbb{N} \to \operatorname{pts}(s)). \\ \quad \operatorname{totally\text{-}bounded}(s) \Rightarrow  \\ \quad \quad \exists i.\; (\operatorname{strictly\text{-}mono\text{-}nn}(i) \wedge \operatorname{is\text{-}cauchy\text{-}seq}(s, \operatorname{subseq}(f, i))) \end{array}}}

\vnbbookfile{theorem-library/rr-ms-dist}{theorem-library/rr-ms-dist}
\noindent{\small rr-ms-dist.scm -- the distance of RR-MS, on the surface.}\par
\vnbprop{rr-ms-dist}{\vnbstmt{\forall u, v.\; (((u \in \mathbb{R}) \wedge (v \in \mathbb{R})) \Rightarrow (\operatorname{dist}(\operatorname{rr\text{-}ms})(u, v) \simeq \lvert (u - v) \rvert))}}

\vnbbookfile{theorem-library/rr-metric-space-proof}{theorem-library/rr-metric-space-proof}
\noindent{\small rr-metric-space-proof.scm -- IS-METRIC-SPACE(RR-MS), PROVEN.}\par
\vnbprop{rr-is-metric-space}{\vnbstmt{\operatorname{is\text{-}metric\text{-}space}(\operatorname{rr\text{-}ms})}}

\vnbbookfile{theorem-library/trunc-metric-proof}{theorem-library/trunc-metric-proof}
\noindent{\small trunc-metric-proof.scm -- the truncated metric d'(x,y) = min(1, d(x,y)) is a metric, is bounded by 1, and is UNIFORMLY equivalent to d.}\par
\vnbprop{rr-min-one-mono}{\vnbstmt{\forall a, b \in \mathbb{R}.\; ((a \leq b) \Rightarrow (\operatorname{min}(1, a) \leq \operatorname{min}(1, b)))}}
\vnbprop{trunc-metric-dist}{\vnbstmt{\begin{array}{@{}l@{}} \forall s, u, v. \\ \quad ((u \in \operatorname{pts}(s)) \wedge (v \in \operatorname{pts}(s))) \Rightarrow  \\ \quad \quad (\operatorname{dist}(\operatorname{trunc\text{-}metric}(s))(u, v) \simeq \operatorname{min}(1, \operatorname{dist}(s)(u, v))) \end{array}}}
\vnbprop{trunc-metric-is-metric-space}{\vnbstmt{\forall s.\; (\operatorname{is\text{-}metric\text{-}space}(s) \Rightarrow \operatorname{is\text{-}metric\text{-}space}(\operatorname{trunc\text{-}metric}(s)))}}
\vnbprop{trunc-metric-bounded}{\vnbstmt{\begin{array}{@{}l@{}} \forall s, u, v. \\ \quad (\operatorname{is\text{-}metric\text{-}space}(s) \wedge ((u \in \operatorname{pts}(s)) \wedge (v \in \operatorname{pts}(s)))) \Rightarrow  \\ \quad \quad \operatorname{dist}(\operatorname{trunc\text{-}metric}(s))(u, v) \leq 1 \end{array}}}
\vnbprop{trunc-metric-id-uniformly-continuous}{\vnbstmt{\begin{array}{@{}l@{}} \forall s. \\ \quad \operatorname{is\text{-}metric\text{-}space}(s) \Rightarrow  \\ \quad \quad \operatorname{is\text{-}uniformly\text{-}continuous}(s, \operatorname{trunc\text{-}metric}(s), (\lambda a \in \operatorname{pts}(s).\; a)) \end{array}}}
\vnbprop{trunc-metric-id-uniformly-continuous-back}{\vnbstmt{\begin{array}{@{}l@{}} \forall s. \\ \quad \operatorname{is\text{-}metric\text{-}space}(s) \Rightarrow  \\ \quad \quad \operatorname{is\text{-}uniformly\text{-}continuous}(\operatorname{trunc\text{-}metric}(s), s, (\lambda a \in \operatorname{pts}(s).\; a)) \end{array}}}

\vnbbookfile{theorem-library/rabasics-trunc}{theorem-library/rabasics-trunc}
\noindent{\small rabasics-trunc.scm -- calculus.pdf 3.27: d'(x, y) = min(d(x, y), 1) is a metric, and a sequence converges for d' exactly when it converges for d.}\par
\vnbprop{trunc-metric-converges-fwd}{\vnbstmt{\begin{array}{@{}l@{}} \forall s.; \forall f \in (\mathbb{N} \to \operatorname{pts}(s)), \ell \in \operatorname{pts}(s). \\ \quad (\operatorname{is\text{-}metric\text{-}space}(s) \wedge \operatorname{converges\text{-}to}(s, f, \ell)) \Rightarrow  \\ \quad \quad \operatorname{converges\text{-}to}(\operatorname{trunc\text{-}metric}(s), f, \ell) \end{array}}}
\vnbprop{trunc-metric-converges-bwd}{\vnbstmt{\begin{array}{@{}l@{}} \forall s.; \forall f \in (\mathbb{N} \to \operatorname{pts}(s)), \ell \in \operatorname{pts}(s). \\ \quad (\operatorname{is\text{-}metric\text{-}space}(s) \wedge \operatorname{converges\text{-}to}(\operatorname{trunc\text{-}metric}(s), f, \ell)) \Rightarrow  \\ \quad \quad \operatorname{converges\text{-}to}(s, f, \ell) \end{array}}}
\vnbprop{trunc-metric-converges-iff}{\vnbstmt{\begin{array}{@{}l@{}} \forall s.; \forall f \in (\mathbb{N} \to \operatorname{pts}(s)), \ell \in \operatorname{pts}(s). \\ \quad \operatorname{is\text{-}metric\text{-}space}(s) \Rightarrow  \\ \quad \quad \operatorname{converges\text{-}to}(\operatorname{trunc\text{-}metric}(s), f, \ell) \Leftrightarrow \operatorname{converges\text{-}to}(s, f, \ell) \end{array}}}

\vnbbookfile{theorem-library/rr-complete-proof}{theorem-library/rr-complete-proof}
\noindent{\small rr-complete-proof.scm -- RR is a COMPLETE METRIC SPACE.  Proven 2026-08-01/02.}\par
\vnbprop{rr-complete}{\vnbstmt{\operatorname{is\text{-}complete}(\operatorname{rr\text{-}ms})}}

\vnbbookfile{theorem-library/metric-limit-unique}{theorem-library/metric-limit-unique}
\noindent{\small metric-limit-unique.scm -- A SEQUENCE IN A METRIC SPACE HAS AT MOST ONE LIMIT, proven once for an arbitrary metric space, with the RR and CC cases taken as instances.}\par
\vnbprop{metric-limit-unique}{\vnbstmt{\begin{array}{@{}l@{}} \forall s, f, \ell_{1}, \ell_{2}. \\ \quad (\operatorname{converges\text{-}to}(s, f, \ell_{1}) \wedge \operatorname{converges\text{-}to}(s, f, \ell_{2})) \Rightarrow (\ell_{1} = \ell_{2}) \end{array}}}
\vnbprop{rr-limit-unique}{\vnbstmt{\begin{array}{@{}l@{}} \forall f \in (\mathbb{N} \to \mathbb{R}), \ell_{1}, \ell_{2} \in \mathbb{R}. \\ \quad (\operatorname{converges\text{-}to}(\operatorname{rr\text{-}ms}, f, \ell_{1}) \wedge \operatorname{converges\text{-}to}(\operatorname{rr\text{-}ms}, f, \ell_{2})) \Rightarrow  \\ \quad \quad \ell_{1} = \ell_{2} \end{array}}}
\vnbprop{cc-limit-unique}{\vnbstmt{\begin{array}{@{}l@{}} \forall f \in (\mathbb{N} \to \mathbb{C}), \ell_{1}, \ell_{2} \in \mathbb{C}. \\ \quad (\operatorname{converges\text{-}to}(\operatorname{cc\text{-}ms}, f, \ell_{1}) \wedge \operatorname{converges\text{-}to}(\operatorname{cc\text{-}ms}, f, \ell_{2})) \Rightarrow  \\ \quad \quad \ell_{1} = \ell_{2} \end{array}}}

\vnbbookfile{theorem-library/seq-limit-core}{theorem-library/seq-limit-core}
\noindent{\small theorem-library/seq-limit-core.scm -- THE LIMIT OF A REAL SEQUENCE AS A TERM (SEQ-LIMIT), its three characterising theorems, the TAIL bound rr-limit-tail-abs-le, and rr-cauchy-converges.}\par
\par\medskip\noindent\textbf{Definition} (\texttt{seq-limit}).\;$\operatorname{seq\text{-}limit}(f) :=$
\begin{equation*}
\left(\text{if }\operatorname{converges}(\operatorname{rr\text{-}ms}, f)\text{ then }(\iota\, \operatorname{lm\_}.\;\operatorname{converges\text{-}to}(\operatorname{rr\text{-}ms}, f, \operatorname{lm\_}))\text{ else }0\right)
\end{equation*}

\vnbprop{seq-limit-converges-to}{\vnbstmt{\forall f.\; (\operatorname{converges}(\operatorname{rr\text{-}ms}, f) \Rightarrow \operatorname{converges\text{-}to}(\operatorname{rr\text{-}ms}, f, \operatorname{seq\text{-}limit}(f)))}}
\vnbprop{seq-limit-in-rr}{\vnbstmt{\forall f.\; (\operatorname{seq\text{-}limit}(f) \in \mathbb{R})}}
\vnbprop{seq-limit-value}{\vnbstmt{\forall f, v.\; (\operatorname{converges\text{-}to}(\operatorname{rr\text{-}ms}, f, v) \Rightarrow (\operatorname{seq\text{-}limit}(f) = v))}}
\vnbprop{rr-limit-tail-abs-le}{\vnbstmt{\begin{array}{@{}l@{}} \forall f \in (\mathbb{N} \to \mathbb{R}), v, a, c \in \mathbb{R}, n \in \mathbb{N}. \\ \quad (\operatorname{converges\text{-}to}(\operatorname{rr\text{-}ms}, f, v) \wedge \forall k \in \mathbb{N}.\; ((n \leq k) \Rightarrow (\lvert (f(k) - a) \rvert \leq c))) \Rightarrow  \\ \quad \quad \lvert (v - a) \rvert \leq c \end{array}}}
\vnbprop{rr-cauchy-converges}{\vnbstmt{\begin{array}{@{}l@{}} \forall f \in (\mathbb{N} \to \mathbb{R}). \\ \quad \forall s. \\ \quad \quad \operatorname{pos\text{-}rr}(s) \Rightarrow  \\ \quad \quad \quad \exists n \in \mathbb{N}. \forall m, p \in \mathbb{N}. \\ \quad \quad \quad \quad ((n \leq m) \wedge (n \leq p)) \Rightarrow (\lvert (f(m) - f(p)) \rvert \leq s) \Rightarrow \\ \quad \quad \operatorname{converges}(\operatorname{rr\text{-}ms}, f) \end{array}}}

\vnbbookfile{structure-library/ascoli-arzela-statement}{structure-library/ascoli-arzela-statement}
\noindent{\small NOTE 2026-09-20 (batch 12-B).  Where this file says the tree has no metric SUBSPACE structure, that is history: SUBSPACE-MS(s, A) and RESTRICT(f, A) are defined in structure-library/metric-subspace.scm, their laws are in theorem-library/metric-subspace-laws.scm and compact-subspace.scm, and}\par
\par\medskip\noindent\textbf{Definition} (\texttt{pointwise-bounded}).\;For all $s$ and $\mathit{fam}$, 
\begin{equation*}
\begin{array}{@{}l@{}} \operatorname{pointwise\text{-}bounded}(s, \mathit{fam}) \\ \Leftrightarrow\; \operatorname{is\text{-}metric\text{-}space}(s) \wedge \\ \mathit{fam} \in (\mathbb{N} \to (\operatorname{pts}(s) \to \mathbb{R})) \wedge \\ \forall x \in \operatorname{pts}(s). \exists \mathit{bd} \in \mathbb{R}. \forall k \in \mathbb{N}. \\ \quad \lvert \operatorname{fam}(k)(x) \rvert \leq \mathit{bd} \end{array}
\end{equation*}

\par\medskip\noindent\textbf{Definition} (\texttt{is-equicontinuous}).\;For all $s$, $t$ and $\mathit{fam}$, 
\begin{equation*}
\begin{array}{@{}l@{}} \operatorname{is\text{-}equicontinuous}(s, t, \mathit{fam}) \\ \Leftrightarrow\; \operatorname{is\text{-}metric\text{-}space}(s) \wedge \\ \operatorname{is\text{-}metric\text{-}space}(t) \wedge \\ \mathit{fam} \in (\mathbb{N} \to (\operatorname{pts}(s) \to \operatorname{pts}(t))) \wedge \\ \forall x \in \operatorname{pts}(s), \mathit{eps}. \\ \quad \operatorname{pos\text{-}rr}(\mathit{eps}) \Rightarrow  \\ \quad \quad \exists \mathit{del}. \\ \quad \quad \quad \operatorname{pos\text{-}rr}(\mathit{del}) \wedge \\ \quad \quad \quad \forall k \in \mathbb{N}, y \in \operatorname{pts}(s). \\ \quad \quad \quad \quad (\operatorname{dist}(s)(x, y) < \mathit{del}) \Rightarrow (\operatorname{dist}(t)(\operatorname{fam}(k)(x), \operatorname{fam}(k)(y)) < \mathit{eps}) \end{array}
\end{equation*}

\par\medskip\noindent\textbf{Definition} (\texttt{converges-uniformly}).\;For all $s$, $\mathit{seq}$ and $g$, 
\begin{equation*}
\begin{array}{@{}l@{}} \operatorname{converges\text{-}uniformly}(s, \mathit{seq}, g) \\ \Leftrightarrow\; \operatorname{is\text{-}metric\text{-}space}(s) \wedge \\ \mathit{seq} \in (\mathbb{N} \to (\operatorname{pts}(s) \to \mathbb{R})) \wedge \\ g \in (\operatorname{pts}(s) \to \mathbb{R}) \wedge \\ \forall \mathit{eps}. \\ \quad \operatorname{pos\text{-}rr}(\mathit{eps}) \Rightarrow  \\ \quad \quad \exists \mathit{cap} \in \mathbb{N}. \forall k. \\ \quad \quad \quad ((k \in \mathbb{N}) \wedge (\mathit{cap} \leq k)) \Rightarrow  \\ \quad \quad \quad \quad \forall x \in \operatorname{pts}(s).\; (\lvert -(\operatorname{seq}(k)(x), g(x)) \rvert < \mathit{eps}) \end{array}
\end{equation*}


\vnbbookfile{theorem-library/ascoli-analytic-cores}{theorem-library/ascoli-analytic-cores}
\noindent{\small ascoli-analytic-cores.scm -- the analytic content of Ascoli-Arzelà, machine-proved. bridge, MACHINE-PROVED, and the one rung it still stands on.}\par
\par\medskip\noindent\textbf{Definition} (\texttt{is-unif-cauchy}).\;For all $s$ and $\mathit{fam}$, 
\begin{equation*}
\begin{array}{@{}l@{}} \operatorname{is\text{-}unif\text{-}cauchy}(s, \mathit{fam}) \\ \Leftrightarrow\; \operatorname{is\text{-}metric\text{-}space}(s) \wedge \\ \mathit{fam} \in (\mathbb{N} \to (\operatorname{pts}(s) \to \mathbb{R})) \wedge \\ \forall \mathit{eps}. \\ \quad \operatorname{pos\text{-}rr}(\mathit{eps}) \Rightarrow  \\ \quad \quad \exists \mathit{cap} \in \mathbb{N}. \forall k. \\ \quad \quad \quad ((k \in \mathbb{N}) \wedge (\mathit{cap} \leq k)) \Rightarrow  \\ \quad \quad \quad \quad \forall l. \\ \quad \quad \quad \quad \quad ((l \in \mathbb{N}) \wedge (\mathit{cap} \leq l)) \Rightarrow  \\ \quad \quad \quad \quad \quad \quad \forall x \in \operatorname{pts}(s).\; (\lvert -(\operatorname{fam}(k)(x), \operatorname{fam}(l)(x)) \rvert < \mathit{eps}) \end{array}
\end{equation*}

\par\medskip\noindent\textbf{Definition} (\texttt{is-ptwise-cauchy}).\;For all $s$ and $\mathit{fam}$, 
\begin{equation*}
\begin{array}{@{}l@{}} \operatorname{is\text{-}ptwise\text{-}cauchy}(s, \mathit{fam}) \\ \Leftrightarrow\; \operatorname{is\text{-}metric\text{-}space}(s) \wedge \\ \mathit{fam} \in (\mathbb{N} \to (\operatorname{pts}(s) \to \mathbb{R})) \wedge \\ \forall x \in \operatorname{pts}(s), \mathit{eps}. \\ \quad \operatorname{pos\text{-}rr}(\mathit{eps}) \Rightarrow  \\ \quad \quad \exists \mathit{cap} \in \mathbb{N}. \forall k. \\ \quad \quad \quad ((k \in \mathbb{N}) \wedge (\mathit{cap} \leq k)) \Rightarrow  \\ \quad \quad \quad \quad \forall l.\; (((l \in \mathbb{N}) \wedge (\mathit{cap} \leq l)) \Rightarrow (\lvert -(\operatorname{fam}(k)(x), \operatorname{fam}(l)(x)) \rvert < \mathit{eps})) \end{array}
\end{equation*}

\par\medskip\noindent\textbf{Definition} (\texttt{is-dense-seq}).\;For all $s$ and $\mathit{dseq}$, 
\begin{equation*}
\begin{array}{@{}l@{}} \operatorname{is\text{-}dense\text{-}seq}(s, \mathit{dseq}) \\ \Leftrightarrow\; \operatorname{is\text{-}metric\text{-}space}(s) \wedge \\ \mathit{dseq} \in (\mathbb{N} \to \operatorname{pts}(s)) \wedge \\ \forall x \in \operatorname{pts}(s), \mathit{eps}. \\ \quad \operatorname{pos\text{-}rr}(\mathit{eps}) \Rightarrow \exists m \in \mathbb{N}.\; (\operatorname{dist}(s)(x, \operatorname{dseq}(m)) < \mathit{eps}) \end{array}
\end{equation*}

\par\medskip\noindent\textbf{Definition} (\texttt{converges-on}).\;For all $s$, $\mathit{fam}$ and $\mathit{dseq}$, 
\begin{equation*}
\begin{array}{@{}l@{}} \operatorname{converges\text{-}on}(s, \mathit{fam}, \mathit{dseq}) \\ \Leftrightarrow\; \operatorname{is\text{-}metric\text{-}space}(s) \wedge \\ \mathit{fam} \in (\mathbb{N} \to (\operatorname{pts}(s) \to \mathbb{R})) \wedge \\ \mathit{dseq} \in (\mathbb{N} \to \operatorname{pts}(s)) \wedge \\ \forall m \in \mathbb{N}.\; \operatorname{converges}(\operatorname{rr\text{-}ms}, (\lambda k \in \mathbb{N}.\; \operatorname{fam}(k)(\operatorname{dseq}(m)))) \end{array}
\end{equation*}

\vnbprop{equicont-dense-conv-ptwise-cauchy}{\vnbstmt{\begin{array}{@{}l@{}} \forall s.; \forall m \in (\mathbb{N} \to (\operatorname{pts}(s) \to \mathbb{R})). \\ \quad (\operatorname{is\text{-}equicontinuous}(s, \operatorname{rr\text{-}ms}, m) \wedge \exists q.\; (\operatorname{is\text{-}dense\text{-}seq}(s, q) \wedge \operatorname{converges\text{-}on}(s, m, q))) \Rightarrow  \\ \quad \quad \operatorname{is\text{-}ptwise\text{-}cauchy}(s, m) \end{array}}}
\vnbprop{uniform-limit-continuous}{\vnbstmt{\begin{array}{@{}l@{}} \forall s, m, g. \\ \quad m \in (\mathbb{N} \to (\operatorname{pts}(s) \to \mathbb{R})) \wedge \\ \quad g \in (\operatorname{pts}(s) \to \mathbb{R}) \wedge \\ \quad \forall k \in \mathbb{N}.\; \operatorname{is\text{-}continuous}(s, \operatorname{rr\text{-}ms}, m(k)) \wedge \\ \quad \operatorname{converges\text{-}uniformly}(s, m, g) \Rightarrow \\ \quad \quad \operatorname{is\text{-}continuous}(s, \operatorname{rr\text{-}ms}, g) \end{array}}}

\vnbbookfile{theorem-library/rake-open-sets}{theorem-library/rake-open-sets}
\noindent{\small rake-open-sets.scm -- the metric topology axioms, PROVEN.}\par
\vnbprop{empty-is-open}{\vnbstmt{\forall s.\; (\operatorname{is\text{-}metric\text{-}space}(s) \Rightarrow \operatorname{is\text{-}open}(s, \emptyset))}}
\vnbprop{carrier-is-open}{\vnbstmt{\forall s.\; (\operatorname{is\text{-}metric\text{-}space}(s) \Rightarrow \operatorname{is\text{-}open}(s, \operatorname{pts}(s)))}}
\vnbprop{union-of-opens-open}{\vnbstmt{\begin{array}{@{}l@{}} \forall s, a, g. \\ \quad (\operatorname{is\text{-}metric\text{-}space}(s) \wedge \forall i \in a.\; \operatorname{is\text{-}open}(s, g(i))) \Rightarrow  \\ \quad \quad \operatorname{is\text{-}open}(s, \operatorname{big\text{-}union}(i, a, g(i))) \end{array}}}
\vnbprop{inter-of-opens-open}{\vnbstmt{\begin{array}{@{}l@{}} \forall s, u, w. \\ \quad (\operatorname{is\text{-}metric\text{-}space}(s) \wedge (\operatorname{is\text{-}open}(s, u) \wedge \operatorname{is\text{-}open}(s, w))) \Rightarrow  \\ \quad \quad \operatorname{is\text{-}open}(s, (u \cap w)) \end{array}}}
\vnbprop{preimage-complement}{\vnbstmt{\begin{array}{@{}l@{}} \forall s, t, f, u. \\ \quad (\operatorname{is\text{-}metric\text{-}space}(s) \wedge (\operatorname{is\text{-}metric\text{-}space}(t) \wedge (f \in (\operatorname{pts}(s) \to \operatorname{pts}(t))))) \Rightarrow  \\ \quad \quad \operatorname{preimage}(s, f, \operatorname{complement\text{-}in}(\operatorname{pts}(t), u)) = \operatorname{complement\text{-}in}(\operatorname{pts}(s), \operatorname{preimage}(s, f, u)) \end{array}}}
\vnbprop{open-preimage-implies-continuous}{\vnbstmt{\begin{array}{@{}l@{}} \forall s, t, f. \\ \quad \operatorname{is\text{-}metric\text{-}space}(s) \wedge \\ \quad \operatorname{is\text{-}metric\text{-}space}(t) \wedge \\ \quad f \in (\operatorname{pts}(s) \to \operatorname{pts}(t)) \wedge \\ \quad \forall v.\; (\operatorname{is\text{-}open}(t, v) \Rightarrow \operatorname{is\text{-}open}(s, \operatorname{preimage}(s, f, v))) \Rightarrow \\ \quad \quad \operatorname{is\text{-}continuous}(s, t, f) \end{array}}}
\vnbprop{continuous-implies-closed-preimage}{\vnbstmt{\begin{array}{@{}l@{}} \forall s, t, f, a. \\ \quad (\operatorname{is\text{-}continuous}(s, t, f) \wedge \operatorname{is\text{-}closed}(t, a)) \Rightarrow  \\ \quad \quad \operatorname{is\text{-}closed}(s, \operatorname{preimage}(s, f, a)) \end{array}}}
\vnbprop{closed-preimage-implies-continuous}{\vnbstmt{\begin{array}{@{}l@{}} \forall s, t, f. \\ \quad \operatorname{is\text{-}metric\text{-}space}(s) \wedge \\ \quad \operatorname{is\text{-}metric\text{-}space}(t) \wedge \\ \quad f \in (\operatorname{pts}(s) \to \operatorname{pts}(t)) \wedge \\ \quad \forall a.\; (\operatorname{is\text{-}closed}(t, a) \Rightarrow \operatorname{is\text{-}closed}(s, \operatorname{preimage}(s, f, a))) \Rightarrow \\ \quad \quad \operatorname{is\text{-}continuous}(s, t, f) \end{array}}}

\vnbbookfile{theorem-library/rabasics-open}{theorem-library/rabasics-open}
\noindent{\small rabasics-open.scm -- calculus.pdf 3.3, complex-analysis.pdf 1.5: the intersection of FINITELY many open sets is open (the union half is union-of-opens-open, already proven).}\par
\vnbprop{intersection-of-insert}{\vnbstmt{\begin{array}{@{}l@{}} \forall f.; \forall a \in \mathbf{Set}. \\ \quad \exists w.\; (w \in f) \Rightarrow  \\ \quad \quad \operatorname{intersection\text{-}of}((f \cup \{a\})) = (\operatorname{intersection\text{-}of}(f) \cap a) \end{array}}}
\vnbprop{intersection-of-insert-empty}{\vnbstmt{\begin{array}{@{}l@{}} \forall f.; \forall a \in \mathbf{Set}. \\ \quad \neg \exists w.\; (w \in f) \Rightarrow (\operatorname{intersection\text{-}of}((f \cup \{a\})) = a) \end{array}}}
\vnbprop{inter-of-finite-opens-open}{\vnbstmt{\begin{array}{@{}l@{}} \forall s.; \forall c \in \mathbf{Set}. \\ \quad (\operatorname{is\text{-}metric\text{-}space}(s) \wedge ((|c| \in \mathbb{N}) \wedge (\exists w.\; (w \in c) \wedge \forall u \in c.\; \operatorname{is\text{-}open}(s, u)))) \Rightarrow  \\ \quad \quad \operatorname{is\text{-}open}(s, \operatorname{intersection\text{-}of}(c)) \end{array}}}

\vnbbookfile{theorem-library/ptwise-cauchy-unif}{theorem-library/ptwise-cauchy-unif}
\noindent{\small ptwise-cauchy-unif.scm -- pointwise Cauchy + equicontinuous on a COMPACT space  =$>$  uniformly Cauchy  (the user's notes, Prop 3.33, Cauchy spelling).}\par
\vnbprop{ptwise-cauchy-compact-equicont-unif}{\vnbstmt{\begin{array}{@{}l@{}} \forall s, m. \\ \quad \operatorname{is\text{-}compact}(s) \wedge \\ \quad m \in (\mathbb{N} \to (\operatorname{pts}(s) \to \mathbb{R})) \wedge \\ \quad \operatorname{is\text{-}equicontinuous}(s, \operatorname{rr\text{-}ms}, m) \wedge \\ \quad \operatorname{is\text{-}ptwise\text{-}cauchy}(s, m) \Rightarrow \\ \quad \quad \operatorname{is\text{-}unif\text{-}cauchy}(s, m) \end{array}}}

\vnbbookfile{theorem-library/ascoli-assembly}{theorem-library/ascoli-assembly}
\noindent{\small theorem-library/ascoli-assembly.scm -- equicont-dense-conv-implies-unif-cauchy.}\par
\vnbprop{equicont-dense-conv-implies-unif-cauchy}{\vnbstmt{\begin{array}{@{}l@{}} \forall s, m. \\ \quad \operatorname{is\text{-}compact}(s) \wedge \\ \quad m \in (\mathbb{N} \to (\operatorname{pts}(s) \to \mathbb{R})) \wedge \\ \quad \operatorname{is\text{-}equicontinuous}(s, \operatorname{rr\text{-}ms}, m) \wedge \\ \quad \exists q.\; (\operatorname{is\text{-}dense\text{-}seq}(s, q) \wedge \operatorname{converges\text{-}on}(s, m, q)) \Rightarrow \\ \quad \quad \operatorname{is\text{-}unif\text{-}cauchy}(s, m) \end{array}}}

\vnbbookfile{theorem-library/unif-cauchy-limit}{theorem-library/unif-cauchy-limit}
\noindent{\small unif-cauchy-limit.scm -- unif-cauchy-has-uniform-limit, PROVEN (2026-09-14).}\par
\vnbprop{unif-cauchy-has-uniform-limit}{\vnbstmt{\begin{array}{@{}l@{}} \forall s.; \forall m \in (\mathbb{N} \to (\operatorname{pts}(s) \to \mathbb{R})). \\ \quad \operatorname{is\text{-}unif\text{-}cauchy}(s, m) \Rightarrow  \\ \quad \quad \exists g \in (\operatorname{pts}(s) \to \mathbb{R}).\; \operatorname{converges\text{-}uniformly}(s, m, g) \end{array}}}

\vnbbookfile{theorem-library/ascoli-bridge}{theorem-library/ascoli-bridge}
\noindent{\small ascoli-bridge.scm -- the dense-to-uniform bridge for Ascoli-Arzelà.}\par
\vnbprop{unif-cauchy-cont-implies-uniform-limit}{\vnbstmt{\begin{array}{@{}l@{}} \forall s.; \forall m \in (\mathbb{N} \to (\operatorname{pts}(s) \to \mathbb{R})). \\ \quad (\forall k \in \mathbb{N}.\; \operatorname{is\text{-}continuous}(s, \operatorname{rr\text{-}ms}, m(k)) \wedge \operatorname{is\text{-}unif\text{-}cauchy}(s, m)) \Rightarrow  \\ \quad \quad \exists g \in (\operatorname{pts}(s) \to \mathbb{R}). \\ \quad \quad \quad \operatorname{is\text{-}continuous}(s, \operatorname{rr\text{-}ms}, g) \wedge \operatorname{converges\text{-}uniformly}(s, m, g) \end{array}}}
\vnbprop{ascoli-dense-bridge}{\vnbstmt{\begin{array}{@{}l@{}} \forall s, m. \\ \quad \operatorname{is\text{-}compact}(s) \wedge \\ \quad m \in (\mathbb{N} \to (\operatorname{pts}(s) \to \mathbb{R})) \wedge \\ \quad \forall k \in \mathbb{N}.\; \operatorname{is\text{-}continuous}(s, \operatorname{rr\text{-}ms}, m(k)) \wedge \\ \quad \operatorname{is\text{-}equicontinuous}(s, \operatorname{rr\text{-}ms}, m) \wedge \\ \quad \exists q.\; (\operatorname{is\text{-}dense\text{-}seq}(s, q) \wedge \operatorname{converges\text{-}on}(s, m, q)) \Rightarrow \\ \quad \quad \exists g \in (\operatorname{pts}(s) \to \mathbb{R}). \\ \quad \quad \quad \operatorname{is\text{-}continuous}(s, \operatorname{rr\text{-}ms}, g) \wedge \operatorname{converges\text{-}uniformly}(s, m, g) \end{array}}}

\vnbbookfile{structure-library/top-product}{structure-library/top-product}
\noindent{\small tych-vocab.scm -- THE PRODUCT OF AN ARBITRARY FAMILY OF TOPOLOGICAL SPACES (definitions only). Agent TYCH (worker-08), 2026-10-26 (NEW DEFINITIONS, the user's to confirm). DESTINATION: structure-library/top-product.scm.  Library: analysis (top-space, cardinality).}\par
\par\medskip\noindent\textbf{Definition} (\texttt{is-top-family}).\;For all $\mathit{ix}$ and $\mathit{tf}$, 
\begin{equation*}
\operatorname{is\text{-}top\text{-}family}(\mathit{ix}, \mathit{tf}) \Leftrightarrow ((\mathit{ix} \in \mathbf{Set}) \wedge \forall \operatorname{tpi\_} \in \mathit{ix}.\; \operatorname{is\text{-}top\text{-}space}(\operatorname{tf}(\operatorname{tpi\_})))
\end{equation*}

\par\medskip\noindent\textbf{Definition} (\texttt{tprod-pts}).\;$\operatorname{tprod\text{-}pts}(\mathit{ix}, \mathit{tf}) :=$
\begin{equation*}
\{\operatorname{tpx\_} \in (\mathit{ix} \to \operatorname{big\text{-}union}(\operatorname{tpj\_}, \mathit{ix}, \operatorname{pts}(\operatorname{tf}(\operatorname{tpj\_})))) : \forall \operatorname{tpi\_} \in \mathit{ix}.\; (\operatorname{tpx\_}(\operatorname{tpi\_}) \in \operatorname{pts}(\operatorname{tf}(\operatorname{tpi\_})))\}
\end{equation*}

\par\medskip\noindent\textbf{Definition} (\texttt{tprod-constr}).\;$\operatorname{tprod\text{-}constr}(\mathit{ix}, \mathit{tf}) :=$
\begin{equation*}
\{\operatorname{tpb\_} \in (\mathit{ix} \times \operatorname{big\text{-}union}(\operatorname{tpj\_}, \mathit{ix}, \operatorname{opens}(\operatorname{tf}(\operatorname{tpj\_})))) : (\operatorname{nth}(2, \operatorname{tpb\_}) \in \operatorname{opens}(\operatorname{tf}(\operatorname{nth}(1, \operatorname{tpb\_}))))\}
\end{equation*}

\par\medskip\noindent\textbf{Definition} (\texttt{tprod-box}).\;$\operatorname{tprod\text{-}box}(\mathit{ix}, \mathit{tf}, \mathit{bb}) :=$
\begin{equation*}
\{\operatorname{tpy\_} \in \operatorname{tprod\text{-}pts}(\mathit{ix}, \mathit{tf}) : \forall \operatorname{tpc\_} \in \mathit{bb}.\; (\operatorname{tpy\_}(\operatorname{nth}(1, \operatorname{tpc\_})) \in \operatorname{nth}(2, \operatorname{tpc\_}))\}
\end{equation*}

\par\medskip\noindent\textbf{Definition} (\texttt{is-tprod-open}).\;For all $\mathit{ix}$, $\mathit{tf}$ and $u$, 
\begin{equation*}
\begin{array}{@{}l@{}} \operatorname{is\text{-}tprod\text{-}open}(\mathit{ix}, \mathit{tf}, u) \\ \Leftrightarrow\; u \subseteq \operatorname{tprod\text{-}pts}(\mathit{ix}, \mathit{tf}) \wedge \\ \forall \operatorname{tpw\_} \in u. \exists \operatorname{tpf\_} \in \mathbf{Set}. \\ \quad \operatorname{tpf\_} \subseteq \operatorname{tprod\text{-}constr}(\mathit{ix}, \mathit{tf}) \wedge \\ \quad |\operatorname{tpf\_}| \in \mathbb{N} \wedge \\ \quad \operatorname{tpw\_} \in \operatorname{tprod\text{-}box}(\mathit{ix}, \mathit{tf}, \operatorname{tpf\_}) \wedge \\ \quad \operatorname{tprod\text{-}box}(\mathit{ix}, \mathit{tf}, \operatorname{tpf\_}) \subseteq u \end{array}
\end{equation*}

\par\medskip\noindent\textbf{Definition} (\texttt{tproduct-top}).\;$\operatorname{tproduct\text{-}top}(\mathit{ix}, \mathit{tf}) :=$
\begin{equation*}
{[}\operatorname{tprod\text{-}pts}(\mathit{ix}, \mathit{tf}), \{\operatorname{tpu\_} \in \operatorname{power}(\operatorname{tprod\text{-}pts}(\mathit{ix}, \mathit{tf})) : \operatorname{is\text{-}tprod\text{-}open}(\mathit{ix}, \mathit{tf}, \operatorname{tpu\_})\}]
\end{equation*}

\par\medskip\noindent\textbf{Definition} (\texttt{tproj}).\;$\operatorname{tproj}(\mathit{ix}, \mathit{tf}, i) :=$ $(\lambda \operatorname{tpz\_} \in \operatorname{tprod\text{-}pts}(\mathit{ix}, \mathit{tf}).\; \operatorname{tpz\_}(i))$
\par\medskip\noindent\textbf{Definition} (\texttt{has-finite-subcover}).\;For all $s$ and $c$, 
\begin{equation*}
\begin{array}{@{}l@{}} \operatorname{has\text{-}finite\text{-}subcover}(s, c) \\ \Leftrightarrow\; \exists \operatorname{tpg\_} \in \mathbf{Set}. \\ \quad (\operatorname{tpg\_} \subseteq c) \wedge ((|\operatorname{tpg\_}| \in \mathbb{N}) \wedge (\operatorname{big\text{-}union}(\operatorname{tpv\_}, \operatorname{tpg\_}, \operatorname{tpv\_}) \simeq \operatorname{pts}(s))) \end{array}
\end{equation*}


\vnbbookfile{structure-library/seminorm-hahn-banach}{structure-library/seminorm-hahn-banach}
\noindent{\small MOVED 2026-09-20 (batch 12-A) from theorem-library/: this file is VOCABULARY -- definitions, notation and warranted supports, not one proof -- and every theorem stated with it had to load below it.  Its load.scm slot is unchanged. seminorm-hahn-banach.scm -- seminorms on a K-vector space and the}\par
\par\medskip\noindent\textbf{Definition} (\texttt{is-seminorm}).\;For all $\mathit{fld}$, $m$ and $p$, 
\begin{equation*}
\begin{array}{@{}l@{}} \operatorname{is\text{-}seminorm}(\mathit{fld}, m, p) \\ \Leftrightarrow\; \operatorname{is\text{-}normed\text{-}field}(\mathit{fld}) \wedge \\ \operatorname{is\text{-}module}(m) \wedge \\ \operatorname{scal}(m) = \operatorname{normed\text{-}field\text{-}as\text{-}commutative\text{-}ring}(\mathit{fld}) \wedge \\ p \in (\operatorname{vec}(m) \to \mathbb{R}) \wedge \\ \forall x \in \operatorname{vec}(m).\; (0 \leq p(x)) \wedge \\ \forall x, y \in \operatorname{vec}(m).\; (p(\operatorname{vadd}(m)(x, y)) \leq +(p(x), p(y))) \wedge \\ \forall \mathit{lam} \in \operatorname{carr}(\operatorname{scal}(m)), x \in \operatorname{vec}(m). \\ \quad p(\operatorname{act}(m)(\mathit{lam}, x)) = \cdot(\operatorname{fnrm}(\mathit{fld})(\mathit{lam}), p(x)) \end{array}
\end{equation*}

\par\medskip\noindent\textbf{Definition} (\texttt{is-k-linear-on}).\;For all $m$, $s$ and $f$, 
\begin{equation*}
\begin{array}{@{}l@{}} \operatorname{is\text{-}k\text{-}linear\text{-}on}(m, s, f) \\ \Leftrightarrow\; \operatorname{is\text{-}module}(m) \wedge \\ \operatorname{is\text{-}submodule}(m, s) \wedge \\ f \in (s \to \operatorname{carr}(\operatorname{scal}(m))) \wedge \\ \forall x, y \in s.\; (f(\operatorname{vadd}(m)(x, y)) = \operatorname{add}(\operatorname{scal}(m))(f(x), f(y))) \wedge \\ \forall \mathit{lam} \in \operatorname{carr}(\operatorname{scal}(m)), x \in s. \\ \quad f(\operatorname{act}(m)(\mathit{lam}, x)) = \operatorname{mul}(\operatorname{scal}(m))(\mathit{lam}, f(x)) \end{array}
\end{equation*}

\par\medskip\noindent\textbf{Definition} (\texttt{is-k-linear}).\;For all $m$ and $f$, 
\begin{equation*}
\begin{array}{@{}l@{}} \operatorname{is\text{-}k\text{-}linear}(m, f) \\ \Leftrightarrow\; \operatorname{is\text{-}module}(m) \wedge \\ f \in (\operatorname{vec}(m) \to \operatorname{carr}(\operatorname{scal}(m))) \wedge \\ \forall x, y \in \operatorname{vec}(m). \\ \quad f(\operatorname{vadd}(m)(x, y)) = \operatorname{add}(\operatorname{scal}(m))(f(x), f(y)) \wedge \\ \forall \mathit{lam} \in \operatorname{carr}(\operatorname{scal}(m)), x \in \operatorname{vec}(m). \\ \quad f(\operatorname{act}(m)(\mathit{lam}, x)) = \operatorname{mul}(\operatorname{scal}(m))(\mathit{lam}, f(x)) \end{array}
\end{equation*}

\par\medskip\noindent\textbf{Definition} (\texttt{is-nontrivially-valued}).\;For every $\mathit{fld}$, 
\begin{equation*}
\begin{array}{@{}l@{}} \operatorname{is\text{-}nontrivially\text{-}valued}(\mathit{fld}) \\ \Leftrightarrow\; \operatorname{is\text{-}normed\text{-}field}(\mathit{fld}) \wedge \\ \exists c \in \operatorname{carr}(\mathit{fld}). \\ \quad (1 < \operatorname{fnrm}(\mathit{fld})(c)) \wedge \exists b \in \operatorname{carr}(\mathit{fld}).\; (\operatorname{mul}(\mathit{fld})(c, b) = \operatorname{one}(\mathit{fld})) \end{array}
\end{equation*}


\vnbbookfile{structure-library/baire-category}{structure-library/baire-category}
\noindent{\small MOVED 2026-09-20 (batch 12-A) from theorem-library/: this file is VOCABULARY -- definitions, notation and warranted supports, not one proof -- and every theorem stated with it had to load below it.  Its load.scm slot is unchanged. baire-category.scm -- interior/closure, nowhere-dense, meager/nonmeager, and}\par
\par\medskip\noindent\textbf{Definition} (\texttt{interior}).\;$\operatorname{interior}(s, a) :=$
\begin{equation*}
\{x \in \operatorname{pts}(s) : \exists u \in \operatorname{power}(\operatorname{pts}(s)).\; (\operatorname{is\text{-}open}(s, u) \wedge ((x \in u) \wedge (u \subseteq a)))\}
\end{equation*}

\par\medskip\noindent\textbf{Definition} (\texttt{closure}).\;$\operatorname{closure}(s, a) :=$
\begin{equation*}
\{x \in \operatorname{pts}(s) : \forall u.\; (((u \in \operatorname{power}(\operatorname{pts}(s))) \wedge (\operatorname{is\text{-}open}(s, u) \wedge (x \in u))) \Rightarrow \exists y \in u.\; (y \in a))\}
\end{equation*}

\par\medskip\noindent\textbf{Definition} (\texttt{is-nowhere-dense}).\;For all $s$ and $a$, 
\begin{equation*}
\begin{array}{@{}l@{}} \operatorname{is\text{-}nowhere\text{-}dense}(s, a) \\ \Leftrightarrow\; (\operatorname{is\text{-}metric\text{-}space}(s) \wedge ((a \subseteq \operatorname{pts}(s)) \wedge (\operatorname{interior}(s, \operatorname{closure}(s, a)) = \emptyset))) \end{array}
\end{equation*}

\par\medskip\noindent\textbf{Definition} (\texttt{is-meager}).\;For all $s$ and $a$, 
\begin{equation*}
\begin{array}{@{}l@{}} \operatorname{is\text{-}meager}(s, a) \\ \Leftrightarrow\; \operatorname{is\text{-}metric\text{-}space}(s) \wedge \\ a \subseteq \operatorname{pts}(s) \wedge \\ \exists \mathit{ee} \in (\mathbb{N} \to \operatorname{power}(\operatorname{pts}(s))). \\ \quad (\forall n \in \mathbb{N}.\; \operatorname{is\text{-}nowhere\text{-}dense}(s, \operatorname{ee}(n)) \wedge (a \subseteq \operatorname{big\text{-}union}(n, \mathbb{N}, \operatorname{ee}(n)))) \end{array}
\end{equation*}

\par\medskip\noindent\textbf{Definition} (\texttt{is-nonmeager}).\;For all $s$ and $a$, 
\begin{equation*}
\operatorname{is\text{-}nonmeager}(s, a) \Leftrightarrow (\operatorname{is\text{-}metric\text{-}space}(s) \wedge ((a \subseteq \operatorname{pts}(s)) \wedge \neg \operatorname{is\text{-}meager}(s, a)))
\end{equation*}


\vnbbookfile{theorem-library/metric-closure-laws}{theorem-library/metric-closure-laws}
\noindent{\small metric-closure-laws.scm -- CLOSURE, INTERIOR and the CLOSED BALL.}\par
\vnbprop{closed-ball-sep-unfold}{\vnbstmt{\forall s, c, r.\; (\operatorname{closed\text{-}ball}(s, c, r) \simeq \{y \in \operatorname{pts}(s) : (\operatorname{dist}(s)(c, y) \leq r)\})}}
\vnbprop{closure-sep-unfold}{\vnbstmt{\begin{array}{@{}l@{}} \forall s, a. \\ \quad (\operatorname{closure}(s, a) \simeq \{x \in \operatorname{pts}(s) : \forall u \in \operatorname{power}(\operatorname{pts}(s)).\; ((\operatorname{is\text{-}open}(s, u) \wedge (x \in u)) \Rightarrow \exists y \in u.\; (y \in a))\}) \end{array}}}
\vnbprop{interior-sep-unfold}{\vnbstmt{\begin{array}{@{}l@{}} \forall s, a. \\ \quad (\operatorname{interior}(s, a) \simeq \{x \in \operatorname{pts}(s) : \exists u \in \operatorname{power}(\operatorname{pts}(s)).\; (\operatorname{is\text{-}open}(s, u) \wedge ((x \in u) \wedge (u \subseteq a)))\}) \end{array}}}
\vnbprop{is-open-in-power}{\vnbstmt{\forall s, u.\; (\operatorname{is\text{-}open}(s, u) \Rightarrow (u \in \operatorname{power}(\operatorname{pts}(s))))}}
\vnbprop{closed-ball-membership}{\vnbstmt{\begin{array}{@{}l@{}} \forall s, c, r, y. \\ \quad (y \in \operatorname{closed\text{-}ball}(s, c, r)) \Leftrightarrow ((y \in \operatorname{pts}(s)) \wedge (\operatorname{dist}(s)(c, y) \leq r)) \end{array}}}
\vnbprop{closure-membership}{\vnbstmt{\begin{array}{@{}l@{}} \forall s, a, x. \\ \quad x \in \operatorname{closure}(s, a) \\ \quad \Leftrightarrow\; ((x \in \operatorname{pts}(s)) \wedge \forall u.\; ((\operatorname{is\text{-}open}(s, u) \wedge (x \in u)) \Rightarrow \exists y \in u.\; (y \in a))) \end{array}}}
\vnbprop{interior-membership}{\vnbstmt{\begin{array}{@{}l@{}} \forall s, a, x. \\ \quad (x \in \operatorname{interior}(s, a)) \Leftrightarrow ((x \in \operatorname{pts}(s)) \wedge \exists u.\; (\operatorname{is\text{-}open}(s, u) \wedge ((x \in u) \wedge (u \subseteq a)))) \end{array}}}
\vnbprop{closed-ball-is-set}{\vnbstmt{\forall s, c, r.\; (\operatorname{is\text{-}metric\text{-}space}(s) \Rightarrow (\operatorname{closed\text{-}ball}(s, c, r) \in \mathbf{Set}))}}
\vnbprop{closure-in-carrier}{\vnbstmt{\forall s, a.\; (\operatorname{is\text{-}metric\text{-}space}(s) \Rightarrow (\operatorname{closure}(s, a) \subseteq \operatorname{pts}(s)))}}
\vnbprop{interior-in-carrier}{\vnbstmt{\forall s, a.\; (\operatorname{is\text{-}metric\text{-}space}(s) \Rightarrow (\operatorname{interior}(s, a) \subseteq \operatorname{pts}(s)))}}
\vnbprop{interior-in-set}{\vnbstmt{\forall s, a.\; (\operatorname{is\text{-}metric\text{-}space}(s) \Rightarrow (\operatorname{interior}(s, a) \subseteq a))}}
\vnbprop{closure-contains}{\vnbstmt{\begin{array}{@{}l@{}} \forall s, a. \\ \quad (\operatorname{is\text{-}metric\text{-}space}(s) \wedge (a \subseteq \operatorname{pts}(s))) \Rightarrow (a \subseteq \operatorname{closure}(s, a)) \end{array}}}
\vnbprop{ball-in-closed-ball}{\vnbstmt{\begin{array}{@{}l@{}} \forall s, c, r. \\ \quad \operatorname{is\text{-}metric\text{-}space}(s) \Rightarrow (\operatorname{ball}(s, c, r) \subseteq \operatorname{closed\text{-}ball}(s, c, r)) \end{array}}}
\vnbprop{closed-ball-in-ball}{\vnbstmt{\begin{array}{@{}l@{}} \forall s.; \forall c \in \operatorname{pts}(s), r, q \in \mathbb{R}. \\ \quad (\operatorname{is\text{-}metric\text{-}space}(s) \wedge (r < q)) \Rightarrow (\operatorname{closed\text{-}ball}(s, c, r) \subseteq \operatorname{ball}(s, c, q)) \end{array}}}
\vnbprop{interior-mem-intro}{\vnbstmt{\begin{array}{@{}l@{}} \forall s, a, u, x. \\ \quad (\operatorname{is\text{-}open}(s, u) \wedge ((u \subseteq a) \wedge (x \in u))) \Rightarrow (x \in \operatorname{interior}(s, a)) \end{array}}}
\vnbprop{interior-is-open}{\vnbstmt{\forall s, a.\; (\operatorname{is\text{-}metric\text{-}space}(s) \Rightarrow \operatorname{is\text{-}open}(s, \operatorname{interior}(s, a)))}}
\vnbprop{closure-is-closed}{\vnbstmt{\forall s, a.\; (\operatorname{is\text{-}metric\text{-}space}(s) \Rightarrow \operatorname{is\text{-}closed}(s, \operatorname{closure}(s, a)))}}
\vnbprop{closed-ball-is-closed}{\vnbstmt{\begin{array}{@{}l@{}} \forall s.; \forall c \in \operatorname{pts}(s), r \in \mathbb{R}. \\ \quad \operatorname{is\text{-}metric\text{-}space}(s) \Rightarrow \operatorname{is\text{-}closed}(s, \operatorname{closed\text{-}ball}(s, c, r)) \end{array}}}
\vnbprop{closure-in-closed}{\vnbstmt{\forall s, a, t.\; ((\operatorname{is\text{-}closed}(s, t) \wedge (a \subseteq t)) \Rightarrow (\operatorname{closure}(s, a) \subseteq t))}}
\vnbprop{closure-of-ball-in-closed-ball}{\vnbstmt{\begin{array}{@{}l@{}} \forall s.; \forall c \in \operatorname{pts}(s), r \in \mathbb{R}. \\ \quad \operatorname{is\text{-}metric\text{-}space}(s) \Rightarrow (\operatorname{closure}(s, \operatorname{ball}(s, c, r)) \subseteq \operatorname{closed\text{-}ball}(s, c, r)) \end{array}}}
\vnbprop{closed-set-limit-in}{\vnbstmt{\begin{array}{@{}l@{}} \forall s, a, f, l. \\ \quad (\operatorname{is\text{-}closed}(s, a) \wedge (\operatorname{converges\text{-}to}(s, f, l) \wedge \forall n \in \mathbb{N}.\; (f(n) \in a))) \Rightarrow  \\ \quad \quad l \in a \end{array}}}
\vnbprop{closure-is-set}{\vnbstmt{\forall s, a.\; (\operatorname{is\text{-}metric\text{-}space}(s) \Rightarrow (\operatorname{closure}(s, a) \in \mathbf{Set}))}}
\vnbprop{interior-is-set}{\vnbstmt{\forall s, a.\; (\operatorname{is\text{-}metric\text{-}space}(s) \Rightarrow (\operatorname{interior}(s, a) \in \mathbf{Set}))}}
\vnbprop{nowhere-dense-ball-point}{\vnbstmt{\begin{array}{@{}l@{}} \forall s, a.; \forall c \in \operatorname{pts}(s), r. \\ \quad (\operatorname{is\text{-}nowhere\text{-}dense}(s, a) \wedge \operatorname{pos\text{-}rr}(r)) \Rightarrow  \\ \quad \quad \exists x \in \operatorname{ball}(s, c, r).\; \neg (x \in \operatorname{closure}(s, a)) \end{array}}}
\vnbprop{nowhere-dense-closed-ball}{\vnbstmt{\begin{array}{@{}l@{}} \forall s, a.; \forall c \in \operatorname{pts}(s), r, d. \\ \quad (\operatorname{is\text{-}nowhere\text{-}dense}(s, a) \wedge (\operatorname{pos\text{-}rr}(r) \wedge \operatorname{pos\text{-}rr}(d))) \Rightarrow  \\ \quad \quad \exists x \in \operatorname{pts}(s), q. \\ \quad \quad \quad \operatorname{pos\text{-}rr}(q) \wedge \\ \quad \quad \quad q \leq d \wedge \\ \quad \quad \quad \operatorname{closed\text{-}ball}(s, x, q) \subseteq \operatorname{ball}(s, c, r) \wedge \\ \quad \quad \quad \forall z \in \operatorname{closed\text{-}ball}(s, x, q).\; \neg (z \in \operatorname{closure}(s, a)) \end{array}}}
\vnbprop{closure-subset-closed}{\vnbstmt{\begin{array}{@{}l@{}} \forall s, a. \\ \quad (\operatorname{is\text{-}metric\text{-}space}(s) \wedge ((a \subseteq \operatorname{pts}(s)) \wedge (\operatorname{closure}(s, a) \subseteq a))) \Rightarrow  \\ \quad \quad \operatorname{is\text{-}closed}(s, a) \end{array}}}

\vnbbookfile{theorem-library/rake-baire}{theorem-library/rake-baire}
\noindent{\small rake-baire.scm -- towards the BAIRE CATEGORY theorem (batch 12-F, 2026-09-20).}\par
\vnbprop{nested-family-monotone}{\vnbstmt{\begin{array}{@{}l@{}} \forall m \in \mathbb{N}, w.; \forall g \in (\mathbb{N} \to w), n \in \mathbb{N}. \\ \quad (\forall n \in \mathbb{N}.\; (g((n + 1)) \subseteq g(n)) \wedge (n \leq m)) \Rightarrow (g(m) \subseteq g(n)) \end{array}}}

\vnbbookfile{structure-library/frechet-open-mapping}{structure-library/frechet-open-mapping}
\noindent{\small MOVED 2026-09-20 (batch 12-A) from theorem-library/: this file is VOCABULARY -- definitions, notation and warranted supports, not one proof -- and every theorem stated with it had to load below it.  Its load.scm slot is unchanged. frechet-open-mapping.scm -- Frechet spaces, continuous linear maps, and the}\par
\par\medskip\noindent\textbf{Definition} (\texttt{fr-ball}).\;$\operatorname{fr\text{-}ball}(m, \mathit{fam}, n, \mathit{eps}) :=$
\begin{equation*}
\{x \in \operatorname{vec}(m) : \forall k.\; (((k \in \mathbb{N}) \wedge (k \leq n)) \Rightarrow (\operatorname{fam}(k)(x) < \mathit{eps}))\}
\end{equation*}
\par\noindent{\small\textit{Comment.} The pair $(m, \mathit{fam})$ is a seminormed space: $m$ a module, its vectors $\operatorname{vec}(m)$, and $\mathit{fam}(k)$ the $k$-th seminorm of a family enumerated by $\mathbb{N}$. The ball is cut out by the first $n + 1$ seminorms, radius $\varepsilon$.}\par

\par\medskip\noindent\textbf{Definition} (\texttt{has-closed-graph}).\;For all $m_{1}$, $\mathit{fam}_{1}$, $m_{2}$, $\mathit{fam}_{2}$ and $\mathit{tt}$, 
\begin{equation*}
\begin{array}{@{}l@{}} \operatorname{has\text{-}closed\text{-}graph}(m_{1}, \mathit{fam}_{1}, m_{2}, \mathit{fam}_{2}, \mathit{tt}) \\ \Leftrightarrow\; \operatorname{is\text{-}linear\text{-}map}(m_{1}, m_{2}, \mathit{tt}) \wedge \\ \forall \mathit{seq} \in (\mathbb{N} \to \operatorname{vec}(m_{1})), p_{1} \in \operatorname{vec}(m_{1}), p_{2} \in \operatorname{vec}(m_{2}). \\ \quad (\operatorname{fr\text{-}conv}(m_{1}, \mathit{fam}_{1}, \mathit{seq}, p_{1}) \wedge \operatorname{fr\text{-}conv}(m_{2}, \mathit{fam}_{2}, (\lambda i \in \mathbb{N}.\; \operatorname{tt}(\operatorname{seq}(i))), p_{2})) \Rightarrow  \\ \quad \quad p_{2} = \operatorname{tt}(p_{1}) \end{array}
\end{equation*}
\par\noindent{\small\textit{Comment.} $(m_1, \mathit{fam}_1)$ is the domain and $(m_2, \mathit{fam}_2)$ the target, each a module with an $\mathbb{N}$-enumerated seminorm family; $\mathit{tt}$ is the map.}\par

\par\medskip\noindent\textbf{Definition} (\texttt{is-seminorm-family}).\;For all $\mathit{fld}$, $m$ and $\mathit{fam}$, 
\begin{equation*}
\begin{array}{@{}l@{}} \operatorname{is\text{-}seminorm\text{-}family}(\mathit{fld}, m, \mathit{fam}) \\ \Leftrightarrow\; \operatorname{is\text{-}normed\text{-}field}(\mathit{fld}) \wedge \\ \operatorname{is\text{-}module}(m) \wedge \\ \operatorname{scal}(m) = \operatorname{normed\text{-}field\text{-}as\text{-}commutative\text{-}ring}(\mathit{fld}) \wedge \\ \mathit{fam} \in (\mathbb{N} \to (\operatorname{vec}(m) \to \mathbb{R})) \wedge \\ \forall k \in \mathbb{N}.\; \operatorname{is\text{-}seminorm}(\mathit{fld}, m, \operatorname{fam}(k)) \wedge \\ \forall x \in \operatorname{vec}(m). \\ \quad \neg (x = \operatorname{vzero}(m)) \Rightarrow \exists k \in \mathbb{N}.\; \neg (\operatorname{fam}(k)(x) = 0) \end{array}
\end{equation*}
\par\noindent{\small\textit{Comment.} $\mathit{fld}$ is the normed field of scalars, $m$ a module over it, and $\mathit{fam} : \mathbb{N} \to (\operatorname{vec}(m) \to \mathbb{R})$ the family of seminorms; separating: a nonzero vector has a nonzero seminorm.}\par

\par\medskip\noindent\textbf{Definition} (\texttt{is-frechet-structure}).\;For all $\mathit{fld}$, $m$ and $\mathit{fam}$, 
\begin{equation*}
\begin{array}{@{}l@{}} \operatorname{is\text{-}frechet\text{-}structure}(\mathit{fld}, m, \mathit{fam}) \\ \Leftrightarrow\; \operatorname{is\text{-}seminorm\text{-}family}(\mathit{fld}, m, \mathit{fam}) \wedge \\ \forall \mathit{seq} \in (\mathbb{N} \to \operatorname{vec}(m)). \\ \quad \operatorname{fr\text{-}cauchy}(m, \mathit{fam}, \mathit{seq}) \Rightarrow  \\ \quad \quad \exists \mathit{lim} \in \operatorname{vec}(m).\; \operatorname{fr\text{-}conv}(m, \mathit{fam}, \mathit{seq}, \mathit{lim}) \end{array}
\end{equation*}
\par\noindent{\small\textit{Comment.} $\mathit{fld}$ is the normed field of scalars; $(m, \mathit{fam})$ is the space, $m$ a module and $\mathit{fam}(k)$ its $k$-th seminorm.}\par

\par\medskip\noindent\textbf{Definition} (\texttt{is-linear-map}).\;For all $m_{1}$, $m_{2}$ and $\mathit{tt}$, 
\begin{equation*}
\begin{array}{@{}l@{}} \operatorname{is\text{-}linear\text{-}map}(m_{1}, m_{2}, \mathit{tt}) \\ \Leftrightarrow\; \operatorname{is\text{-}module}(m_{1}) \wedge \\ \operatorname{is\text{-}module}(m_{2}) \wedge \\ \operatorname{scal}(m_{1}) = \operatorname{scal}(m_{2}) \wedge \\ \mathit{tt} \in (\operatorname{vec}(m_{1}) \to \operatorname{vec}(m_{2})) \wedge \\ \forall x, y \in \operatorname{vec}(m_{1}). \\ \quad \operatorname{tt}(\operatorname{vadd}(m_{1})(x, y)) = \operatorname{vadd}(m_{2})(\operatorname{tt}(x), \operatorname{tt}(y)) \wedge \\ \forall \mathit{lam} \in \operatorname{carr}(\operatorname{scal}(m_{1})), x \in \operatorname{vec}(m_{1}). \\ \quad \operatorname{tt}(\operatorname{act}(m_{1})(\mathit{lam}, x)) = \operatorname{act}(m_{2})(\mathit{lam}, \operatorname{tt}(x)) \end{array}
\end{equation*}
\par\noindent{\small\textit{Comment.} $m_1$ and $m_2$ are modules over one scalar ring, $\mathit{tt} : \operatorname{vec}(m_1) \to \operatorname{vec}(m_2)$ the map.}\par

\par\medskip\noindent\textbf{Definition} (\texttt{fr-conv}).\;For all $m$, $\mathit{fam}$, $\mathit{seq}$ and $\mathit{lim}$, 
\begin{equation*}
\begin{array}{@{}l@{}} \operatorname{fr\text{-}conv}(m, \mathit{fam}, \mathit{seq}, \mathit{lim}) \\ \Leftrightarrow\; \forall k \in \mathbb{N}, \mathit{eps}. \\ \quad \operatorname{pos\text{-}rr}(\mathit{eps}) \Rightarrow  \\ \quad \quad \exists \mathit{cap} \in \mathbb{N}. \forall i. \\ \quad \quad \quad ((i \in \mathbb{N}) \wedge (\mathit{cap} \leq i)) \Rightarrow (\operatorname{fam}(k)(\operatorname{vadd}(m)(\operatorname{seq}(i), \operatorname{vneg}(m)(\mathit{lim}))) < \mathit{eps}) \end{array}
\end{equation*}
\par\noindent{\small\textit{Comment.} The pair $(m, \mathit{fam})$ is a seminormed space: $m$ a module, its vectors $\operatorname{vec}(m)$, and $\mathit{fam}(k)$ the $k$-th seminorm of a family enumerated by $\mathbb{N}$. $\mathit{seq}$ is a sequence of vectors and $\mathit{lim}$ the limit.}\par

\par\medskip\noindent\textbf{Definition} (\texttt{fr-cauchy}).\;For all $m$, $\mathit{fam}$ and $\mathit{seq}$, 
\begin{equation*}
\begin{array}{@{}l@{}} \operatorname{fr\text{-}cauchy}(m, \mathit{fam}, \mathit{seq}) \\ \Leftrightarrow\; \forall k \in \mathbb{N}, \mathit{eps}. \\ \quad \operatorname{pos\text{-}rr}(\mathit{eps}) \Rightarrow  \\ \quad \quad \exists \mathit{cap} \in \mathbb{N}. \forall i. \\ \quad \quad \quad ((i \in \mathbb{N}) \wedge (\mathit{cap} \leq i)) \Rightarrow  \\ \quad \quad \quad \quad \forall j. \\ \quad \quad \quad \quad \quad ((j \in \mathbb{N}) \wedge (\mathit{cap} \leq j)) \Rightarrow  \\ \quad \quad \quad \quad \quad \quad \operatorname{fam}(k)(\operatorname{vadd}(m)(\operatorname{seq}(i), \operatorname{vneg}(m)(\operatorname{seq}(j)))) < \mathit{eps} \end{array}
\end{equation*}
\par\noindent{\small\textit{Comment.} The pair $(m, \mathit{fam})$ is a seminormed space: $m$ a module, its vectors $\operatorname{vec}(m)$, and $\mathit{fam}(k)$ the $k$-th seminorm of a family enumerated by $\mathbb{N}$. $\mathit{seq}$ is a sequence of vectors.}\par

\par\medskip\noindent\textbf{Definition} (\texttt{is-cont-lin}).\;For all $m_{1}$, $\mathit{fam}_{1}$, $m_{2}$, $\mathit{fam}_{2}$ and $\mathit{tt}$, 
\begin{equation*}
\begin{array}{@{}l@{}} \operatorname{is\text{-}cont\text{-}lin}(m_{1}, \mathit{fam}_{1}, m_{2}, \mathit{fam}_{2}, \mathit{tt}) \\ \Leftrightarrow\; \operatorname{is\text{-}linear\text{-}map}(m_{1}, m_{2}, \mathit{tt}) \wedge \\ \forall \mathit{seq} \in (\mathbb{N} \to \operatorname{vec}(m_{1})), \mathit{lim} \in \operatorname{vec}(m_{1}). \\ \quad \operatorname{fr\text{-}conv}(m_{1}, \mathit{fam}_{1}, \mathit{seq}, \mathit{lim}) \Rightarrow  \\ \quad \quad \operatorname{fr\text{-}conv}(m_{2}, \mathit{fam}_{2}, (\lambda i \in \mathbb{N}.\; \operatorname{tt}(\operatorname{seq}(i))), \operatorname{tt}(\mathit{lim})) \end{array}
\end{equation*}
\par\noindent{\small\textit{Comment.} $(m_1, \mathit{fam}_1)$ is the domain and $(m_2, \mathit{fam}_2)$ the target, each a module with an $\mathbb{N}$-enumerated seminorm family; $\mathit{tt}$ is the map. Continuity is sequential.}\par

\par\medskip\noindent\textbf{Definition} (\texttt{is-open-lin-map}).\;For all $m_{1}$, $\mathit{fam}_{1}$, $m_{2}$, $\mathit{fam}_{2}$ and $\mathit{tt}$, 
\begin{equation*}
\begin{array}{@{}l@{}} \operatorname{is\text{-}open\text{-}lin\text{-}map}(m_{1}, \mathit{fam}_{1}, m_{2}, \mathit{fam}_{2}, \mathit{tt}) \\ \Leftrightarrow\; \operatorname{is\text{-}linear\text{-}map}(m_{1}, m_{2}, \mathit{tt}) \wedge \\ \forall n \in \mathbb{N}, \mathit{eps}. \\ \quad \operatorname{pos\text{-}rr}(\mathit{eps}) \Rightarrow  \\ \quad \quad \exists \operatorname{n\_} \in \mathbb{N}, \operatorname{eps\_}. \\ \quad \quad \quad (\operatorname{pos\text{-}rr}(\operatorname{eps\_}) \wedge (\operatorname{fr\text{-}ball}(m_{2}, \mathit{fam}_{2}, \operatorname{n\_}, \operatorname{eps\_}) \subseteq \operatorname{image}(\mathit{tt}, \operatorname{fr\text{-}ball}(m_{1}, \mathit{fam}_{1}, n, \mathit{eps})))) \end{array}
\end{equation*}
\par\noindent{\small\textit{Comment.} $(m_1, \mathit{fam}_1)$ is the domain and $(m_2, \mathit{fam}_2)$ the target, each a module with an $\mathbb{N}$-enumerated seminorm family; $\mathit{tt}$ is the map.}\par


\vnbbookfile{structure-library/frechet-metric}{structure-library/frechet-metric}
\noindent{\small frechet-metric.scm -- the METRIC of a seminormed space whose seminorm family is COUNTABLE (Thayer, "Analise Funcional e Teoria Espectral", Ch. 8 s.4, Prop. 8.24). Added 2026-10-10 (batch TVS-1).  Vocabulary only: the theorems are in theorem-library/frechet-metric.scm.}\par
\par\medskip\noindent\textbf{Definition} (\texttt{fr-c-space}).\;$\operatorname{fr\text{-}c\text{-}space}(m, \mathit{fam}) :=$
\begin{equation*}
{[}\operatorname{vec}(m), (\lambda \operatorname{frk\_} \in \mathbb{N}.\; (\lambda {[}\operatorname{fru\_}, \operatorname{frv\_}] \in (\operatorname{vec}(m) \times \operatorname{vec}(m)).\; \operatorname{fam}(\operatorname{frk\_})(\operatorname{vadd}(m)(\operatorname{fru\_}, \operatorname{vneg}(m)(\operatorname{frv\_})))))]
\end{equation*}

\par\medskip\noindent\textbf{Definition} (\texttt{fr-metric-space}).\;$\operatorname{fr\text{-}metric\text{-}space}(m, \mathit{fam}) :=$ $\operatorname{c\text{-}metric}(\operatorname{fr\text{-}c\text{-}space}(m, \mathit{fam}))$
\par\medskip\noindent\textbf{Definition} (\texttt{fr-metric}).\;$\operatorname{fr\text{-}metric}(m, \mathit{fam}) :=$ $\operatorname{dist}(\operatorname{fr\text{-}metric\text{-}space}(m, \mathit{fam}))$

\vnbbookfile{structure-library/frechet-banach-steinhaus}{structure-library/frechet-banach-steinhaus}
\noindent{\small frechet-banach-steinhaus.scm -- BOUNDED SETS and EQUICONTINUOUS FAMILIES of linear maps in a space with a countable family of seminorms (Thayer, "Analise Funcional e Teoria Espectral", Ch. 8 s.4, Def. 8.31 and the equicontinuity of Thm 8.32, p. 257). Added 2026-10-11 (batch TVS-2B).  Vocabulary only: the theorems ...}\par
\par\medskip\noindent\textbf{Definition} (\texttt{is-fr-bounded}).\;For all $m$, $\mathit{fam}$ and $\mathit{aa}$, 
\begin{equation*}
\begin{array}{@{}l@{}} \operatorname{is\text{-}fr\text{-}bounded}(m, \mathit{fam}, \mathit{aa}) \\ \Leftrightarrow\; ((\mathit{aa} \subseteq \operatorname{vec}(m)) \wedge \forall k \in \mathbb{N}.\; \exists \mathit{bd} \in \mathbb{R}.\; \forall y \in \mathit{aa}.\; (\operatorname{fam}(k)(y) \leq \mathit{bd})) \end{array}
\end{equation*}
\par\noindent{\small\textit{Comment.} $(m, \mathit{fam})$ is the seminormed space ($m$ a module, $\mathit{fam}(k)$ its $k$-th seminorm); $\mathit{aa}$ is the set, a subset of $\operatorname{vec}(m)$.}\par

\par\medskip\noindent\textbf{Definition} (\texttt{is-fr-equicont}).\;For all $m_{1}$, $\mathit{fam}_{1}$, $m_{2}$, $\mathit{fam}_{2}$, $\mathit{ix}$ and $\mathit{tf}$, 
\begin{equation*}
\begin{array}{@{}l@{}} \operatorname{is\text{-}fr\text{-}equicont}(m_{1}, \mathit{fam}_{1}, m_{2}, \mathit{fam}_{2}, \mathit{ix}, \mathit{tf}) \\ \Leftrightarrow\; \mathit{tf} \in (\mathit{ix} \to (\operatorname{vec}(m_{1}) \to \operatorname{vec}(m_{2}))) \wedge \\ \forall k \in \mathbb{N}. \exists \mathit{nb} \in \mathbb{N}, \mathit{ck}. \\ \quad \operatorname{pos\text{-}rr}(\mathit{ck}) \wedge \\ \quad \forall i \in \mathit{ix}, x \in \operatorname{vec}(m_{1}), r \in \mathbb{R}. \\ \quad \quad \forall j \in \mathbb{N}.\; ((j \leq \mathit{nb}) \Rightarrow (\operatorname{fam1}(j)(x) \leq r)) \Rightarrow  \\ \quad \quad \quad \operatorname{fam2}(k)(\operatorname{tf}(i)(x)) \leq \cdot(\mathit{ck}, r) \end{array}
\end{equation*}
\par\noindent{\small\textit{Comment.} $(m_1, \mathit{fam}_1)$ is the domain and $(m_2, \mathit{fam}_2)$ the target, each a module with an $\mathbb{N}$-enumerated seminorm family; $\mathit{ix}$ is the index class and $\mathit{tf}(i)$, for $i$ in $\mathit{ix}$, the $i$-th map of the family.}\par


\vnbbookfile{theorem-library/subseq-convergence-proof}{theorem-library/subseq-convergence-proof}
\noindent{\small theorem-library/subseq-convergence-proof.scm ==================================================================== subseq-of-convergent -- PROVEN.}\par
\vnbprop{subseq-of-convergent}{\vnbstmt{\begin{array}{@{}l@{}} \forall s, f, i, l. \\ \quad (\operatorname{converges\text{-}to}(s, f, l) \wedge \operatorname{strictly\text{-}mono\text{-}nn}(i)) \Rightarrow  \\ \quad \quad \operatorname{converges\text{-}to}(s, \operatorname{subseq}(f, i), l) \end{array}}}

\vnbbookfile{theorem-library/subsqn-basics}{theorem-library/subsqn-basics}
\noindent{\small subsqn-basics.scm -- the SUBSQN relation and what is proven about it.}\par
\vnbprop{is-subsequence-iff-typed-subsqn}{\vnbstmt{\begin{array}{@{}l@{}} \forall s, y, f. \\ \quad \operatorname{is\text{-}subsequence}(s, y, f) \Leftrightarrow ((f \in (\mathbb{N} \to \operatorname{pts}(s))) \wedge \operatorname{subsqn}(y, f)) \end{array}}}
\vnbprop{subsqn-preserves-sqn}{\vnbstmt{\forall a, f, y.\; (((f \in a^{\mathbb{N}}) \wedge \operatorname{subsqn}(y, f)) \Rightarrow (y \in a^{\mathbb{N}}))}}

\vnbbookfile{theorem-library/subsequence-principle}{theorem-library/subsequence-principle}
\noindent{\small theorem-library/subsequence-principle.scm ====================================================================== subsequence-principle -- PROVEN.}\par
\vnbprop{subsequence-principle}{\vnbstmt{\begin{array}{@{}l@{}} \forall s, f, l. \\ \quad \operatorname{is\text{-}metric\text{-}space}(s) \wedge \\ \quad f \in (\mathbb{N} \to \operatorname{pts}(s)) \wedge \\ \quad l \in \operatorname{pts}(s) \wedge \\ \quad \forall i. \\ \quad \quad \operatorname{strictly\text{-}mono\text{-}nn}(i) \Rightarrow  \\ \quad \quad \quad \exists a. \\ \quad \quad \quad \quad \operatorname{strictly\text{-}mono\text{-}nn}(a) \wedge \operatorname{converges\text{-}to}(s, \operatorname{subseq}(\operatorname{subseq}(f, i), a), l) \Rightarrow \\ \quad \quad \operatorname{converges\text{-}to}(s, f, l) \end{array}}}

\vnbbookfile{theorem-library/rake-offbill-combinatorial}{theorem-library/rake-offbill-combinatorial}
\noindent{\small theorem-library/rake-offbill-combinatorial.scm ==================================================================== Batch 8-D: the off-bill `combinatorial' PSS entries that batch 6/7 made cheap.  Three theorems, in dependency order:}\par
\vnbprop{tb-block-step}{\vnbstmt{\begin{array}{@{}l@{}} \forall s.; \forall f \in (\mathbb{N} \to \operatorname{pts}(s)), r.; \forall j \in \mathrm{Inf}(\mathbb{N}). \\ \quad (\operatorname{totally\text{-}bounded}(s) \wedge \operatorname{pos\text{-}rr}(r)) \Rightarrow  \\ \quad \quad \exists a \in \mathrm{Inf}(\mathbb{N}). \\ \quad \quad \quad (a \subseteq j) \wedge \exists c \in \operatorname{pts}(s).\; \forall i \in a.\; (f(i) \in \operatorname{ball}(s, c, r)) \end{array}}}
\vnbprop{tb-has-eps-cauchy-subseq}{\vnbstmt{\begin{array}{@{}l@{}} \forall s.; \forall f \in (\mathbb{N} \to \operatorname{pts}(s)), a. \\ \quad (\operatorname{totally\text{-}bounded}(s) \wedge \operatorname{pos\text{-}rr}(a)) \Rightarrow  \\ \quad \quad \exists i.\; (\operatorname{strictly\text{-}mono\text{-}nn}(i) \wedge \operatorname{is\text{-}eps\text{-}cauchy\text{-}seq}(s, a, \operatorname{subseq}(f, i))) \end{array}}}
\vnbprop{block-family}{\vnbstmt{\begin{array}{@{}l@{}} \forall s.; \forall f \in (\mathbb{N} \to \operatorname{pts}(s)), d. \\ \quad (\operatorname{totally\text{-}bounded}(s) \wedge \operatorname{null\text{-}rr\text{-}seq}(d)) \Rightarrow  \\ \quad \quad \exists a \in (\mathbb{N} \to \mathrm{Inf}(\mathbb{N})). \\ \quad \quad \quad \forall k \in \mathbb{N}.\; (a((k + 1)) \subseteq a(k)) \wedge \\ \quad \quad \quad \forall k \in \mathbb{N}. \exists c \in \operatorname{pts}(s). \forall i \in a(k). \\ \quad \quad \quad \quad f(i) \in \operatorname{ball}(s, c, d(k)) \end{array}}}

\vnbbookfile{theorem-library/rake-baire-2}{theorem-library/rake-baire-2}
\noindent{\small rake-baire-2.scm -- the BAIRE CATEGORY theorem (batch 13-C, 2026-09-20).}\par
\vnbprop{closed-ball-in-carrier}{\vnbstmt{\begin{array}{@{}l@{}} \forall s, v, \mathit{qv}. \\ \quad \operatorname{is\text{-}metric\text{-}space}(s) \Rightarrow (\operatorname{closed\text{-}ball}(s, v, \mathit{qv}) \subseteq \operatorname{pts}(s)) \end{array}}}
\vnbprop{closed-ball-center-in}{\vnbstmt{\begin{array}{@{}l@{}} \forall s.; \forall v \in \operatorname{pts}(s), \mathit{qv} \in \mathbb{R}. \\ \quad (\operatorname{is\text{-}metric\text{-}space}(s) \wedge (0 \leq \mathit{qv})) \Rightarrow (v \in \operatorname{closed\text{-}ball}(s, v, \mathit{qv})) \end{array}}}
\vnbprop{closed-ball-in-power}{\vnbstmt{\begin{array}{@{}l@{}} \forall s, v, \mathit{qv}. \\ \quad \operatorname{is\text{-}metric\text{-}space}(s) \Rightarrow (\operatorname{closed\text{-}ball}(s, v, \mathit{qv}) \in \operatorname{power}(\operatorname{pts}(s))) \end{array}}}
\vnbprop{closed-set-limit-in-tail}{\vnbstmt{\begin{array}{@{}l@{}} \forall s, v, \mathit{fv}, \mathit{lv}, b. \\ \quad \operatorname{is\text{-}closed}(s, v) \wedge \\ \quad \operatorname{converges\text{-}to}(s, \mathit{fv}, \mathit{lv}) \wedge \\ \quad b \in \mathbb{N} \wedge \\ \quad \forall n \in \mathbb{N}.\; ((b \leq n) \Rightarrow (\operatorname{fv}(n) \in v)) \Rightarrow \\ \quad \quad \mathit{lv} \in v \end{array}}}
\vnbprop{nested-closed-balls-point}{\vnbstmt{\begin{array}{@{}l@{}} \forall s.; \forall r \in (\mathbb{N} \to \operatorname{pts}(s)), \mathit{rds} \in (\mathbb{N} \to \mathbb{R}). \\ \quad \operatorname{is\text{-}complete}(s) \wedge \\ \quad \forall n \in \mathbb{N}.\; (0 \leq \operatorname{rds}(n)) \wedge \\ \quad \forall n \in \mathbb{N}. \\ \quad \quad \operatorname{closed\text{-}ball}(s, r((n + 1)), \operatorname{rds}((n + 1))) \subseteq \operatorname{closed\text{-}ball}(s, r(n), \operatorname{rds}(n)) \wedge \\ \quad \forall n \in \mathbb{N}.\; (\operatorname{rds}((n + 1)) \leq \operatorname{recip}((n + 1))) \Rightarrow \\ \quad \quad \exists v \in \operatorname{pts}(s).\; \forall n \in \mathbb{N}.\; (v \in \operatorname{closed\text{-}ball}(s, r(n), \operatorname{rds}(n))) \end{array}}}
\vnbprop{baire-category}{\vnbstmt{\begin{array}{@{}l@{}} \forall s. \\ \quad (\operatorname{is\text{-}complete}(s) \wedge \exists x.\; (x \in \operatorname{pts}(s))) \Rightarrow \operatorname{is\text{-}nonmeager}(s, \operatorname{pts}(s)) \end{array}}}

\vnbbookfile{theorem-library/product2-laws}{theorem-library/product2-laws}
\noindent{\small product2-laws.scm -- the laws of PRODUCT2-MS (structure-library/product2-vocab.scm), the binary product of two metric spaces with the MAX metric.  Everything PROVEN, modulo 0; nothing asserted.  Helper prefix p2-.}\par
\vnbprop{product2-pts}{\vnbstmt{\begin{array}{@{}l@{}} \forall s, t. \\ \quad (\operatorname{is\text{-}metric\text{-}space}(s) \wedge \operatorname{is\text{-}metric\text{-}space}(t)) \Rightarrow  \\ \quad \quad (\operatorname{pts}(\operatorname{product2\text{-}ms}(s, t)) \simeq (\operatorname{pts}(s) \times \operatorname{pts}(t))) \end{array}}}
\vnbprop{product2-dist-slot}{\vnbstmt{\begin{array}{@{}l@{}} \forall s, t. \\ \quad (\operatorname{is\text{-}metric\text{-}space}(s) \wedge \operatorname{is\text{-}metric\text{-}space}(t)) \Rightarrow  \\ \quad \quad (\operatorname{dist}(\operatorname{product2\text{-}ms}(s, t)) \simeq (\lambda {[}\operatorname{p2u\_}, \operatorname{p2v\_}] \in ((\operatorname{pts}(s) \times \operatorname{pts}(t)) \times (\operatorname{pts}(s) \times \operatorname{pts}(t))).\; \operatorname{max}(\operatorname{dist}(s)(\operatorname{nth}(1, \operatorname{p2u\_}), \operatorname{nth}(1, \operatorname{p2v\_})), \operatorname{dist}(t)(\operatorname{nth}(2, \operatorname{p2u\_}), \operatorname{nth}(2, \operatorname{p2v\_}))))) \end{array}}}
\vnbprop{product2-dist}{\vnbstmt{\begin{array}{@{}l@{}} \forall s, t, p, q. \\ \quad \operatorname{is\text{-}metric\text{-}space}(s) \wedge \\ \quad \operatorname{is\text{-}metric\text{-}space}(t) \wedge \\ \quad p \in (\operatorname{pts}(s) \times \operatorname{pts}(t)) \wedge \\ \quad q \in (\operatorname{pts}(s) \times \operatorname{pts}(t)) \Rightarrow \\ \quad \quad (\operatorname{dist}(\operatorname{product2\text{-}ms}(s, t))(p, q) \simeq \operatorname{max}(\operatorname{dist}(s)(\operatorname{nth}(1, p), \operatorname{nth}(1, q)), \operatorname{dist}(t)(\operatorname{nth}(2, p), \operatorname{nth}(2, q)))) \end{array}}}
\vnbprop{product2-dist-pair}{\vnbstmt{\begin{array}{@{}l@{}} \forall s, t, a, c, b, d. \\ \quad \operatorname{is\text{-}metric\text{-}space}(s) \wedge \\ \quad \operatorname{is\text{-}metric\text{-}space}(t) \wedge \\ \quad a \in \operatorname{pts}(s) \wedge \\ \quad c \in \operatorname{pts}(s) \wedge \\ \quad b \in \operatorname{pts}(t) \wedge \\ \quad d \in \operatorname{pts}(t) \Rightarrow \\ \quad \quad (\operatorname{dist}(\operatorname{product2\text{-}ms}(s, t))({[}a, b], {[}c, d]) \simeq \operatorname{max}(\operatorname{dist}(s)(a, c), \operatorname{dist}(t)(b, d))) \end{array}}}
\vnbprop{product2-is-metric-space}{\vnbstmt{\begin{array}{@{}l@{}} \forall s, t. \\ \quad (\operatorname{is\text{-}metric\text{-}space}(s) \wedge \operatorname{is\text{-}metric\text{-}space}(t)) \Rightarrow  \\ \quad \quad \operatorname{is\text{-}metric\text{-}space}(\operatorname{product2\text{-}ms}(s, t)) \end{array}}}
\vnbprop{product2-fst-continuous-at}{\vnbstmt{\begin{array}{@{}l@{}} \forall s, t.; \forall h \in ((\operatorname{pts}(s) \times \operatorname{pts}(t)) \to \operatorname{pts}(s)), x \in (\operatorname{pts}(s) \times \operatorname{pts}(t)). \\ \quad \operatorname{is\text{-}metric\text{-}space}(s) \wedge \\ \quad \operatorname{is\text{-}metric\text{-}space}(t) \wedge \\ \quad \forall z \in (\operatorname{pts}(s) \times \operatorname{pts}(t)).\; (h(z) = \operatorname{nth}(1, z)) \Rightarrow \\ \quad \quad \operatorname{is\text{-}continuous\text{-}at}(\operatorname{product2\text{-}ms}(s, t), s, h, x) \end{array}}}
\vnbprop{product2-snd-continuous-at}{\vnbstmt{\begin{array}{@{}l@{}} \forall s, t.; \forall h \in ((\operatorname{pts}(s) \times \operatorname{pts}(t)) \to \operatorname{pts}(t)), x \in (\operatorname{pts}(s) \times \operatorname{pts}(t)). \\ \quad \operatorname{is\text{-}metric\text{-}space}(s) \wedge \\ \quad \operatorname{is\text{-}metric\text{-}space}(t) \wedge \\ \quad \forall z \in (\operatorname{pts}(s) \times \operatorname{pts}(t)).\; (h(z) = \operatorname{nth}(2, z)) \Rightarrow \\ \quad \quad \operatorname{is\text{-}continuous\text{-}at}(\operatorname{product2\text{-}ms}(s, t), t, h, x) \end{array}}}
\vnbprop{product2-pair-continuous-at}{\vnbstmt{\begin{array}{@{}l@{}} \forall w, s, t, f, g, x.; \forall h \in (\operatorname{pts}(w) \to (\operatorname{pts}(s) \times \operatorname{pts}(t))). \\ \quad \operatorname{is\text{-}continuous\text{-}at}(w, s, f, x) \wedge \\ \quad \operatorname{is\text{-}continuous\text{-}at}(w, t, g, x) \wedge \\ \quad \forall z \in \operatorname{pts}(w).\; (h(z) = {[}f(z), g(z)]) \Rightarrow \\ \quad \quad \operatorname{is\text{-}continuous\text{-}at}(w, \operatorname{product2\text{-}ms}(s, t), h, x) \end{array}}}
\vnbprop{product2-continuous-at-coords}{\vnbstmt{\begin{array}{@{}l@{}} \forall w, s, t, f, g, h, x. \\ \quad \operatorname{is\text{-}metric\text{-}space}(s) \wedge \\ \quad \operatorname{is\text{-}metric\text{-}space}(t) \wedge \\ \quad f \in (\operatorname{pts}(w) \to \operatorname{pts}(s)) \wedge \\ \quad g \in (\operatorname{pts}(w) \to \operatorname{pts}(t)) \wedge \\ \quad \forall z \in \operatorname{pts}(w).\; (h(z) = {[}f(z), g(z)]) \wedge \\ \quad \operatorname{is\text{-}continuous\text{-}at}(w, \operatorname{product2\text{-}ms}(s, t), h, x) \Rightarrow \\ \quad \quad \operatorname{is\text{-}continuous\text{-}at}(w, s, f, x) \wedge \operatorname{is\text{-}continuous\text{-}at}(w, t, g, x) \end{array}}}
\vnbprop{product2-continuous-at-iff}{\vnbstmt{\begin{array}{@{}l@{}} \forall w, s, t, f, g, h, x. \\ \quad \operatorname{is\text{-}metric\text{-}space}(s) \wedge \\ \quad \operatorname{is\text{-}metric\text{-}space}(t) \wedge \\ \quad f \in (\operatorname{pts}(w) \to \operatorname{pts}(s)) \wedge \\ \quad g \in (\operatorname{pts}(w) \to \operatorname{pts}(t)) \wedge \\ \quad h \in (\operatorname{pts}(w) \to (\operatorname{pts}(s) \times \operatorname{pts}(t))) \wedge \\ \quad \forall z \in \operatorname{pts}(w).\; (h(z) = {[}f(z), g(z)]) \Rightarrow \\ \quad \quad \operatorname{is\text{-}continuous\text{-}at}(w, \operatorname{product2\text{-}ms}(s, t), h, x) \\ \quad \quad \Leftrightarrow\; \operatorname{is\text{-}continuous\text{-}at}(w, s, f, x) \wedge \operatorname{is\text{-}continuous\text{-}at}(w, t, g, x) \end{array}}}
\vnbprop{product2-converges-to-iff}{\vnbstmt{\begin{array}{@{}l@{}} \forall s, t, u, a, b, l, m. \\ \quad \operatorname{is\text{-}metric\text{-}space}(s) \wedge \\ \quad \operatorname{is\text{-}metric\text{-}space}(t) \wedge \\ \quad u \in (\mathbb{N} \to (\operatorname{pts}(s) \times \operatorname{pts}(t))) \wedge \\ \quad a \in (\mathbb{N} \to \operatorname{pts}(s)) \wedge \\ \quad b \in (\mathbb{N} \to \operatorname{pts}(t)) \wedge \\ \quad \forall n \in \mathbb{N}.\; (u(n) = {[}a(n), b(n)]) \wedge \\ \quad l \in \operatorname{pts}(s) \wedge \\ \quad m \in \operatorname{pts}(t) \Rightarrow \\ \quad \quad \operatorname{converges\text{-}to}(\operatorname{product2\text{-}ms}(s, t), u, {[}l, m]) \Leftrightarrow (\operatorname{converges\text{-}to}(s, a, l) \wedge \operatorname{converges\text{-}to}(t, b, m)) \end{array}}}

\vnbbookfile{theorem-library/coord-block-estimate-proof}{theorem-library/coord-block-estimate-proof}
\noindent{\small theorem-library/coord-block-estimate-proof.scm ==================================================================== coord-block-estimate -- PROVEN.}\par
\vnbprop{coord-block-estimate}{\vnbstmt{\begin{array}{@{}l@{}} \forall s, g, b, p, a, c. \\ \quad \operatorname{converges\text{-}along}(s, g, b, p) \wedge \\ \quad \operatorname{strictly\text{-}mono\text{-}nn}(a) \wedge \\ \quad c \in \mathbb{N} \wedge \\ \quad \forall j \in \mathbb{N}.\; ((c \leq j) \Rightarrow (a(j) \in b)) \Rightarrow \\ \quad \quad \operatorname{converges\text{-}to}(s, \operatorname{subseq}(g, a), p) \end{array}}}

\vnbbookfile{theorem-library/sqrt2-proof}{theorem-library/sqrt2-proof}
\noindent{\small sqrt2-proof.scm -- sqrt(2) is irrational, stated on the naturals:}\par
\vnbprop{sqrt2-irrational}{\vnbstmt{\forall p, q \in \mathbb{N}.\; (\neg (q = 0) \Rightarrow \neg ((p \cdot p) = (2 \cdot (q \cdot q))))}}

\vnbbookfile{theorem-library/sqrt3-proof}{theorem-library/sqrt3-proof}
\noindent{\small sqrt3-proof.scm -- sqrt(3) is irrational, on the naturals: forall p,q in NN.  q /= 0  =$>$  p*p /= 3*(q*q).}\par
\vnbprop{sqrt3-irrational}{\vnbstmt{\forall p, q \in \mathbb{N}.\; (\neg (q = 0) \Rightarrow \neg ((p \cdot p) = (3 \cdot (q \cdot q))))}}

\vnbbookfile{theorem-library/continuity-basics}{theorem-library/continuity-basics}
\noindent{\small continuity-basics.scm -- the constant and identity maps are continuous, PROVEN, with their FUN typings.}\par
\vnbprop{const-lam-in-fun}{\vnbstmt{\forall c \in \mathbb{R}.\; ((\lambda x \in \mathbb{R}.\; c) \in (\mathbb{R} \to \mathbb{R}))}}
\vnbprop{ident-lam-in-fun}{\vnbstmt{(\lambda x \in \mathbb{R}.\; x) \in (\mathbb{R} \to \mathbb{R})}}
\vnbprop{const-continuous-at}{\vnbstmt{\begin{array}{@{}l@{}} \forall c, a \in \mathbb{R}. \\ \quad \operatorname{is\text{-}continuous\text{-}at}(\operatorname{rr\text{-}ms}, \operatorname{rr\text{-}ms}, (\lambda x \in \mathbb{R}.\; c), a) \end{array}}}
\vnbprop{identity-continuous-at}{\vnbstmt{\forall a \in \mathbb{R}.\; \operatorname{is\text{-}continuous\text{-}at}(\operatorname{rr\text{-}ms}, \operatorname{rr\text{-}ms}, (\lambda x \in \mathbb{R}.\; x), a)}}

\vnbbookfile{theorem-library/neg-continuous}{theorem-library/neg-continuous}
\noindent{\small neg-continuous.scm -- NEGATION preserves membership in FUN(RR,RR) and continuity, PROVEN.}\par
\vnbprop{neg-fun-in-fun}{\vnbstmt{\forall f \in (\mathbb{R} \to \mathbb{R}).\; ((\lambda z \in \mathbb{R}.\; -f(z)) \in (\mathbb{R} \to \mathbb{R}))}}
\vnbprop{neg-continuous-at}{\vnbstmt{\begin{array}{@{}l@{}} \forall f, x. \\ \quad ((f \in (\mathbb{R} \to \mathbb{R})) \wedge ((x \in \mathbb{R}) \wedge \operatorname{is\text{-}continuous\text{-}at}(\operatorname{rr\text{-}ms}, \operatorname{rr\text{-}ms}, f, x))) \Rightarrow  \\ \quad \quad \operatorname{is\text{-}continuous\text{-}at}(\operatorname{rr\text{-}ms}, \operatorname{rr\text{-}ms}, (\lambda z \in \mathbb{R}.\; -f(z)), x) \end{array}}}

\vnbbookfile{theorem-library/continuity-sum}{theorem-library/continuity-sum}
\noindent{\small continuity-sum.scm -- the pointwise SUM of two maps continuous at a point is continuous there, PROVEN, with its FUN typing.}\par
\vnbprop{sum-lam-in-fun}{\vnbstmt{\begin{array}{@{}l@{}} \forall g, h. \\ \quad ((g \in (\mathbb{R} \to \mathbb{R})) \wedge (h \in (\mathbb{R} \to \mathbb{R}))) \Rightarrow  \\ \quad \quad (\lambda x \in \mathbb{R}.\; (g(x) + h(x))) \in (\mathbb{R} \to \mathbb{R}) \end{array}}}
\vnbprop{sum-continuous-at}{\vnbstmt{\begin{array}{@{}l@{}} \forall g, h, a. \\ \quad (\operatorname{is\text{-}continuous\text{-}at}(\operatorname{rr\text{-}ms}, \operatorname{rr\text{-}ms}, g, a) \wedge \operatorname{is\text{-}continuous\text{-}at}(\operatorname{rr\text{-}ms}, \operatorname{rr\text{-}ms}, h, a)) \Rightarrow  \\ \quad \quad \operatorname{is\text{-}continuous\text{-}at}(\operatorname{rr\text{-}ms}, \operatorname{rr\text{-}ms}, (\lambda x \in \mathbb{R}.\; (g(x) + h(x))), a) \end{array}}}

\vnbbookfile{theorem-library/continuity-product}{theorem-library/continuity-product}
\noindent{\small continuity-product.scm -- the pointwise PRODUCT of two maps continuous at a point is continuous there, PROVEN, with its FUN typing.}\par
\vnbprop{prod-lam-in-fun}{\vnbstmt{\begin{array}{@{}l@{}} \forall g, h. \\ \quad ((g \in (\mathbb{R} \to \mathbb{R})) \wedge (h \in (\mathbb{R} \to \mathbb{R}))) \Rightarrow  \\ \quad \quad (\lambda x \in \mathbb{R}.\; (g(x) \cdot h(x))) \in (\mathbb{R} \to \mathbb{R}) \end{array}}}
\vnbprop{product-continuous-at}{\vnbstmt{\begin{array}{@{}l@{}} \forall g, h, a. \\ \quad (\operatorname{is\text{-}continuous\text{-}at}(\operatorname{rr\text{-}ms}, \operatorname{rr\text{-}ms}, g, a) \wedge \operatorname{is\text{-}continuous\text{-}at}(\operatorname{rr\text{-}ms}, \operatorname{rr\text{-}ms}, h, a)) \Rightarrow  \\ \quad \quad \operatorname{is\text{-}continuous\text{-}at}(\operatorname{rr\text{-}ms}, \operatorname{rr\text{-}ms}, (\lambda x \in \mathbb{R}.\; (g(x) \cdot h(x))), a) \end{array}}}

\vnbbookfile{theorem-library/continuity-transfer}{theorem-library/continuity-transfer}
\noindent{\small continuity-transfer.scm -- continuity reads only the VALUES: a map that agrees pointwise with a map continuous at a is itself continuous at a. PROVEN.}\par
\vnbprop{cont-transfer-ptwise-eq}{\vnbstmt{\begin{array}{@{}l@{}} \forall f, g, a. \\ \quad ((f \in (\mathbb{R} \to \mathbb{R})) \wedge (\operatorname{is\text{-}continuous\text{-}at}(\operatorname{rr\text{-}ms}, \operatorname{rr\text{-}ms}, g, a) \wedge \forall x \in \mathbb{R}.\; (f(x) = g(x)))) \Rightarrow  \\ \quad \quad \operatorname{is\text{-}continuous\text{-}at}(\operatorname{rr\text{-}ms}, \operatorname{rr\text{-}ms}, f, a) \end{array}}}

\vnbbookfile{theorem-library/continuity-sub}{theorem-library/continuity-sub}
\noindent{\small continuity-sub.scm -- the pointwise DIFFERENCE of two maps continuous at a point is continuous there, PROVEN, with its FUN typing.}\par
\vnbprop{sub-lam-in-fun}{\vnbstmt{\begin{array}{@{}l@{}} \forall g, h. \\ \quad ((g \in (\mathbb{R} \to \mathbb{R})) \wedge (h \in (\mathbb{R} \to \mathbb{R}))) \Rightarrow  \\ \quad \quad (\lambda x \in \mathbb{R}.\; (g(x) - h(x))) \in (\mathbb{R} \to \mathbb{R}) \end{array}}}
\vnbprop{sub-continuous-at}{\vnbstmt{\begin{array}{@{}l@{}} \forall g, h, a. \\ \quad (\operatorname{is\text{-}continuous\text{-}at}(\operatorname{rr\text{-}ms}, \operatorname{rr\text{-}ms}, g, a) \wedge \operatorname{is\text{-}continuous\text{-}at}(\operatorname{rr\text{-}ms}, \operatorname{rr\text{-}ms}, h, a)) \Rightarrow  \\ \quad \quad \operatorname{is\text{-}continuous\text{-}at}(\operatorname{rr\text{-}ms}, \operatorname{rr\text{-}ms}, (\lambda x \in \mathbb{R}.\; (g(x) - h(x))), a) \end{array}}}

\vnbbookfile{theorem-library/continuity-scale}{theorem-library/continuity-scale}
\noindent{\small continuity-scale.scm -- the SCALAR MULTIPLE of a real map: its FUN typing and its continuity, both PROVEN `modulo 0'.}\par
\vnbprop{scale-lam-in-fun}{\vnbstmt{\begin{array}{@{}l@{}} \forall c, f. \\ \quad ((c \in \mathbb{R}) \wedge (f \in (\mathbb{R} \to \mathbb{R}))) \Rightarrow ((\lambda x \in \mathbb{R}.\; (c \cdot f(x))) \in (\mathbb{R} \to \mathbb{R})) \end{array}}}
\vnbprop{scale-continuous-at}{\vnbstmt{\begin{array}{@{}l@{}} \forall c \in \mathbb{R}, f, a. \\ \quad \operatorname{is\text{-}continuous\text{-}at}(\operatorname{rr\text{-}ms}, \operatorname{rr\text{-}ms}, f, a) \Rightarrow  \\ \quad \quad \operatorname{is\text{-}continuous\text{-}at}(\operatorname{rr\text{-}ms}, \operatorname{rr\text{-}ms}, (\lambda x \in \mathbb{R}.\; (c \cdot f(x))), a) \end{array}}}

\vnbbookfile{theorem-library/continuity-compose}{theorem-library/continuity-compose}
\noindent{\small continuity-compose.scm -- COMPOSITION OF CONTINUOUS MAPS, PROVEN.}\par
\vnbprop{compose-continuous-at}{\vnbstmt{\begin{array}{@{}l@{}} \forall g, f, a. \\ \quad (\operatorname{is\text{-}continuous\text{-}at}(\operatorname{rr\text{-}ms}, \operatorname{rr\text{-}ms}, f, a) \wedge \operatorname{is\text{-}continuous\text{-}at}(\operatorname{rr\text{-}ms}, \operatorname{rr\text{-}ms}, g, f(a))) \Rightarrow  \\ \quad \quad \operatorname{is\text{-}continuous\text{-}at}(\operatorname{rr\text{-}ms}, \operatorname{rr\text{-}ms}, \operatorname{compose}(g, f), a) \end{array}}}

\vnbbookfile{theorem-library/continuity-recip}{theorem-library/continuity-recip}
\noindent{\small continuity-recip.scm -- THE RECIPROCAL OF A NON-VANISHING CONTINUOUS MAP, PROVEN `modulo 0'.}\par
\vnbprop{recip-lam-in-fun}{\vnbstmt{\begin{array}{@{}l@{}} \forall h \in (\mathbb{R} \to \mathbb{R}). \\ \quad \forall x \in \mathbb{R}.\; \neg (h(x) = 0) \Rightarrow  \\ \quad \quad (\lambda x \in \mathbb{R}.\; \operatorname{recip}(h(x))) \in (\mathbb{R} \to \mathbb{R}) \end{array}}}
\vnbprop{recip-continuous-at}{\vnbstmt{\begin{array}{@{}l@{}} \forall h, p. \\ \quad ((h \in (\mathbb{R} \to \mathbb{R})) \wedge (\forall x \in \mathbb{R}.\; \neg (h(x) = 0) \wedge \operatorname{is\text{-}continuous\text{-}at}(\operatorname{rr\text{-}ms}, \operatorname{rr\text{-}ms}, h, p))) \Rightarrow  \\ \quad \quad \operatorname{is\text{-}continuous\text{-}at}(\operatorname{rr\text{-}ms}, \operatorname{rr\text{-}ms}, (\lambda x \in \mathbb{R}.\; \operatorname{recip}(h(x))), p) \end{array}}}

\vnbbookfile{theorem-library/cont-agree-off-pt}{theorem-library/cont-agree-off-pt}
\noindent{\small cont-agree-off-pt.scm -- two functions continuous at a point and agreeing AWAY from it agree AT it, PROVEN.}\par
\vnbprop{cont-agree-off-pt}{\vnbstmt{\begin{array}{@{}l@{}} \forall a, b.; \forall t \in \mathbb{R}. \\ \quad \operatorname{is\text{-}continuous\text{-}at}(\operatorname{rr\text{-}ms}, \operatorname{rr\text{-}ms}, a, t) \wedge \\ \quad \operatorname{is\text{-}continuous\text{-}at}(\operatorname{rr\text{-}ms}, \operatorname{rr\text{-}ms}, b, t) \wedge \\ \quad \forall x \in \mathbb{R}.\; (\neg (x = t) \Rightarrow (a(x) = b(x))) \Rightarrow \\ \quad \quad a(t) = b(t) \end{array}}}

\vnbbookfile{theorem-library/continuous-one-sided-sign}{theorem-library/continuous-one-sided-sign}
\noindent{\small continuous-one-sided-sign.scm -- sign preservation under continuity, from ONE side, PROVEN.  Two theorems, one driver parameterised by the side:}\par
\vnbprop{continuous-nonneg-left}{\vnbstmt{\begin{array}{@{}l@{}} \forall g, a, h. \\ \quad g \in (\mathbb{R} \to \mathbb{R}) \wedge \\ \quad a \in \mathbb{R} \wedge \\ \quad h \in \mathbb{R} \wedge \\ \quad a < h \wedge \\ \quad \operatorname{is\text{-}continuous\text{-}at}(\operatorname{rr\text{-}ms}, \operatorname{rr\text{-}ms}, g, h) \wedge \\ \quad \forall x \in \mathbb{R}.\; (((a < x) \wedge (x < h)) \Rightarrow (0 \leq g(x))) \Rightarrow \\ \quad \quad 0 \leq g(h) \end{array}}}
\vnbprop{continuous-nonpos-right}{\vnbstmt{\begin{array}{@{}l@{}} \forall g, h, b. \\ \quad g \in (\mathbb{R} \to \mathbb{R}) \wedge \\ \quad h \in \mathbb{R} \wedge \\ \quad b \in \mathbb{R} \wedge \\ \quad h < b \wedge \\ \quad \operatorname{is\text{-}continuous\text{-}at}(\operatorname{rr\text{-}ms}, \operatorname{rr\text{-}ms}, g, h) \wedge \\ \quad \forall x \in \mathbb{R}.\; (((h < x) \wedge (x < b)) \Rightarrow (g(x) \leq 0)) \Rightarrow \\ \quad \quad g(h) \leq 0 \end{array}}}

\vnbbookfile{theorem-library/differentiation}{theorem-library/differentiation}
\noindent{\small differentiation.scm -- Chapter 2 (Differentiation) of calculus.pdf, founded on the CARATHEODORY / o(h) ALGEBRAIC definition (user's choice 2026-06-27).}\par
\vnbprop{diff-implies-continuous}{\vnbstmt{\begin{array}{@{}l@{}} \forall f, a, l. \\ \quad \operatorname{is\text{-}diff\text{-}at}(f, a, l) \Rightarrow \operatorname{is\text{-}continuous\text{-}at}(\operatorname{rr\text{-}ms}, \operatorname{rr\text{-}ms}, f, a) \end{array}}}
\vnbprop{derivative-unique}{\vnbstmt{\forall f, a, l, m.\; ((\operatorname{is\text{-}diff\text{-}at}(f, a, l) \wedge \operatorname{is\text{-}diff\text{-}at}(f, a, m)) \Rightarrow (l = m))}}
\vnbprop{deriv-const}{\vnbstmt{\begin{array}{@{}l@{}} \forall c, a. \\ \quad ((c \in \mathbb{R}) \wedge (a \in \mathbb{R})) \Rightarrow \operatorname{is\text{-}diff\text{-}at}((\lambda x \in \mathbb{R}.\; c), a, 0) \end{array}}}
\vnbprop{deriv-identity}{\vnbstmt{\forall a \in \mathbb{R}.\; \operatorname{is\text{-}diff\text{-}at}((\lambda x \in \mathbb{R}.\; x), a, 1)}}

\vnbbookfile{theorem-library/diff-transfer}{theorem-library/diff-transfer}
\noindent{\small diff-transfer.scm -- differentiability reads only the VALUES: a map that agrees pointwise with a map differentiable at a is itself differentiable at a, with the same derivative.  PROVEN, `modulo 0'.}\par
\vnbprop{diff-transfer-ptwise-eq}{\vnbstmt{\begin{array}{@{}l@{}} \forall f, g, a, l. \\ \quad ((f \in (\mathbb{R} \to \mathbb{R})) \wedge (\operatorname{is\text{-}diff\text{-}at}(g, a, l) \wedge \forall x \in \mathbb{R}.\; (f(x) \simeq g(x)))) \Rightarrow  \\ \quad \quad \operatorname{is\text{-}diff\text{-}at}(f, a, l) \end{array}}}

\vnbbookfile{theorem-library/deriv-sum-product}{theorem-library/deriv-sum-product}
\noindent{\small deriv-sum-product.scm -- calculus.pdf Prop 2.5 (the SUM rule) and Prop 2.6 (the PRODUCT rule), PROVEN, both `modulo 0'.}\par
\vnbprop{deriv-sum}{\vnbstmt{\begin{array}{@{}l@{}} \forall f, g, a, l, m. \\ \quad (\operatorname{is\text{-}diff\text{-}at}(f, a, l) \wedge \operatorname{is\text{-}diff\text{-}at}(g, a, m)) \Rightarrow  \\ \quad \quad \operatorname{is\text{-}diff\text{-}at}((\lambda x \in \mathbb{R}.\; (f(x) + g(x))), a, (l + m)) \end{array}}}
\vnbprop{deriv-product}{\vnbstmt{\begin{array}{@{}l@{}} \forall f, g, a, l, m. \\ \quad (\operatorname{is\text{-}diff\text{-}at}(f, a, l) \wedge \operatorname{is\text{-}diff\text{-}at}(g, a, m)) \Rightarrow  \\ \quad \quad \operatorname{is\text{-}diff\text{-}at}((\lambda x \in \mathbb{R}.\; (f(x) \cdot g(x))), a, ((l \cdot g(a)) + (f(a) \cdot m))) \end{array}}}

\vnbbookfile{theorem-library/chain-rule}{theorem-library/chain-rule}
\noindent{\small chain-rule.scm -- THE CHAIN RULE (calculus.pdf Prop 2.8), PROVEN.}\par
\vnbprop{deriv-chain}{\vnbstmt{\begin{array}{@{}l@{}} \forall f, g, a, l, m. \\ \quad (\operatorname{is\text{-}diff\text{-}at}(f, a, l) \wedge \operatorname{is\text{-}diff\text{-}at}(g, f(a), m)) \Rightarrow  \\ \quad \quad \operatorname{is\text{-}diff\text{-}at}(\operatorname{compose}(g, f), a, (m \cdot l)) \end{array}}}
\vnbprop{deriv-of-is-diff-at}{\vnbstmt{\forall f, a, v.\; (\operatorname{is\text{-}diff\text{-}at}(f, a, v) \Rightarrow (\operatorname{deriv}(f, a) = v))}}
\vnbprop{deriv-chain-value}{\vnbstmt{\begin{array}{@{}l@{}} \forall f, g, a, l, m. \\ \quad (\operatorname{is\text{-}diff\text{-}at}(f, a, l) \wedge \operatorname{is\text{-}diff\text{-}at}(g, f(a), m)) \Rightarrow  \\ \quad \quad \operatorname{deriv}(\operatorname{compose}(g, f), a) = (\operatorname{deriv}(g, f(a)) \cdot \operatorname{deriv}(f, a)) \end{array}}}

\vnbbookfile{theorem-library/inverse-function}{theorem-library/inverse-function}
\noindent{\small inverse-function.scm -- THE INVERSE FUNCTION THEOREM ON THE LINE, Caratheodory form, PROVEN:}\par
\vnbprop{rr-recip-solve}{\vnbstmt{\begin{array}{@{}l@{}} \forall c, u, v. \\ \quad ((c \in \mathbb{R}) \wedge ((u \in \mathbb{R}) \wedge ((v \in \mathbb{R}) \wedge (\neg (c = 0) \wedge (u = (c \cdot v)))))) \Rightarrow  \\ \quad \quad v = (\operatorname{recip}(c) \cdot u) \end{array}}}
\vnbprop{deriv-inverse}{\vnbstmt{\begin{array}{@{}l@{}} \forall f, g, a, l. \\ \quad \operatorname{is\text{-}diff\text{-}at}(f, a, l) \wedge \\ \quad \neg (l = 0) \wedge \\ \quad g \in (\mathbb{R} \to \mathbb{R}) \wedge \\ \quad \forall x \in \mathbb{R}.\; (g(f(x)) = x) \wedge \\ \quad \forall y \in \mathbb{R}.\; (f(g(y)) = y) \wedge \\ \quad \operatorname{is\text{-}continuous\text{-}at}(\operatorname{rr\text{-}ms}, \operatorname{rr\text{-}ms}, g, f(a)) \Rightarrow \\ \quad \quad \operatorname{is\text{-}diff\text{-}at}(g, f(a), \operatorname{recip}(l)) \end{array}}}

\vnbbookfile{theorem-library/higher-derivatives}{theorem-library/higher-derivatives}
\noindent{\small higher-derivatives.scm -- calculus.pdf Def 2.2: the n-th derivative f\textasciicircum{}(n), by induction on n, built on DERIV / IS-DIFF-AT (differentiation.scm).}\par
\par\medskip\noindent\textbf{Definition} (\texttt{nth-deriv}).\;For every $f$,  $(\operatorname{nth\text{-}deriv}(f, 0) \simeq f)$\par
Let $f$ be arbitrary. Let $n \in \mathbb{N}$. Then
\begin{equation*}
(\operatorname{nth\text{-}deriv}(f, \operatorname{succ}(n)) \simeq (\lambda x \in \mathbb{R}.\; \operatorname{deriv}(\operatorname{nth\text{-}deriv}(f, n), x)))
\end{equation*}

\vnbprop{nth-deriv-one}{\vnbstmt{\forall f.\; (\operatorname{nth\text{-}deriv}(f, 1) \simeq (\lambda x \in \mathbb{R}.\; \operatorname{deriv}(f, x)))}}

\vnbbookfile{structure-library/little-o}{structure-library/little-o}
\noindent{\small MOVED 2026-09-20 (batch 12-A) from theorem-library/: this file is VOCABULARY -- definitions, notation and warranted supports, not one proof -- and every theorem stated with it had to load below it.  Its load.scm slot is unchanged. little-o.scm -- calculus.pdf Section 2.2, the calculus of o/O notation, in a}\par
\par\medskip\noindent\textbf{Definition} (\texttt{little-o-at}).\;For all $g$ and $a$, 
\begin{equation*}
\begin{array}{@{}l@{}} \operatorname{little\text{-}o\text{-}at}(g, a) \\ \Leftrightarrow\; g \in (\mathbb{R} \to \mathbb{R}) \wedge \\ a \in \mathbb{R} \wedge \\ \exists \mathit{eps} \in (\mathbb{R} \to \mathbb{R}). \\ \quad (\operatorname{is\text{-}continuous\text{-}at}(\operatorname{rr\text{-}ms}, \operatorname{rr\text{-}ms}, \mathit{eps}, a) \wedge ((\operatorname{eps}(a) = 0) \wedge \forall x \in \mathbb{R}.\; (g(x) = \cdot(\operatorname{eps}(x), -(x, a))))) \end{array}
\end{equation*}


\vnbbookfile{theorem-library/ccint-basics}{theorem-library/ccint-basics}
\noindent{\small ccint-basics.scm -- CCINT's membership law, PROVEN from its definition.}\par
\vnbprop{ccint-membership}{\vnbstmt{\forall a, b, x.\; ((x \in \operatorname{ccint}(a, b)) \Leftrightarrow ((x \in \mathbb{R}) \wedge ((a \leq x) \wedge (x \leq b))))}}
\vnbprop{ccint-elt-in-rr}{\vnbstmt{\forall a, b.; \forall x \in \operatorname{ccint}(a, b).\; (x \in \mathbb{R})}}

\vnbbookfile{theorem-library/ivt-proof}{theorem-library/ivt-proof}
\noindent{\small ivt-proof.scm -- the INTERMEDIATE VALUE THEOREM, PROVEN.}\par
\vnbprop{ivt}{\vnbstmt{\begin{array}{@{}l@{}} \forall f, a, b, c. \\ \quad f \in (\mathbb{R} \to \mathbb{R}) \wedge \\ \quad a \in \mathbb{R} \wedge \\ \quad b \in \mathbb{R} \wedge \\ \quad c \in \mathbb{R} \wedge \\ \quad a \leq b \wedge \\ \quad \forall x \in \operatorname{ccint}(a, b).\; \operatorname{is\text{-}continuous\text{-}at}(\operatorname{rr\text{-}ms}, \operatorname{rr\text{-}ms}, f, x) \wedge \\ \quad f(a) \leq c \wedge \\ \quad c \leq f(b) \Rightarrow \\ \quad \quad \exists p \in \mathbb{R}.\; ((p \in \operatorname{ccint}(a, b)) \wedge (f(p) = c)) \end{array}}}

\vnbbookfile{theorem-library/monotone-inverse}{theorem-library/monotone-inverse}
\noindent{\small monotone-inverse.scm -- MONOTONICITY VOCABULARY FOR REAL FUNCTIONS, and the EXISTENCE HALF of the inverse function theorem on a closed interval. Everything here is `modulo 0'.}\par
\par\medskip\noindent\textbf{Definition} (\texttt{is-strictly-increasing-on}).\;For all $f$ and $s$, 
\begin{equation*}
\begin{array}{@{}l@{}} \operatorname{is\text{-}strictly\text{-}increasing\text{-}on}(f, s) \\ \Leftrightarrow\; ((f \in (\mathbb{R} \to \mathbb{R})) \wedge ((s \subseteq \mathbb{R}) \wedge \forall u, v \in s.\; ((u < v) \Rightarrow (f(u) < f(v))))) \end{array}
\end{equation*}

\vnbprop{ccint-subset-rr}{\vnbstmt{\forall a, b.\; (\operatorname{ccint}(a, b) \subseteq \mathbb{R})}}
\vnbprop{strict-increase-injective}{\vnbstmt{\begin{array}{@{}l@{}} \forall f, s.; \forall u, v \in s. \\ \quad (\operatorname{is\text{-}strictly\text{-}increasing\text{-}on}(f, s) \wedge (f(u) = f(v))) \Rightarrow (u = v) \end{array}}}
\vnbprop{monotone-ccint-inverse-exists}{\vnbstmt{\begin{array}{@{}l@{}} \forall f, a, b.; \forall y \in \operatorname{ccint}(f(a), f(b)). \\ \quad f \in (\mathbb{R} \to \mathbb{R}) \wedge \\ \quad a \in \mathbb{R} \wedge \\ \quad b \in \mathbb{R} \wedge \\ \quad a \leq b \wedge \\ \quad \forall x \in \operatorname{ccint}(a, b).\; \operatorname{is\text{-}continuous\text{-}at}(\operatorname{rr\text{-}ms}, \operatorname{rr\text{-}ms}, f, x) \wedge \\ \quad \operatorname{is\text{-}strictly\text{-}increasing\text{-}on}(f, \operatorname{ccint}(a, b)) \Rightarrow \\ \quad \quad \exists p \in \operatorname{ccint}(a, b). \\ \quad \quad \quad (f(p) = y) \wedge \forall q \in \operatorname{ccint}(a, b).\; ((f(q) = y) \Rightarrow (p = q)) \end{array}}}
\vnbprop{monotone-ccint-inverse-order}{\vnbstmt{\begin{array}{@{}l@{}} \forall f, a, b.; \forall u, v \in \operatorname{ccint}(a, b). \\ \quad (\operatorname{is\text{-}strictly\text{-}increasing\text{-}on}(f, \operatorname{ccint}(a, b)) \wedge (f(u) < f(v))) \Rightarrow (u < v) \end{array}}}

\vnbbookfile{theorem-library/sq-continuous}{theorem-library/sq-continuous}
\noindent{\small sq-continuous.scm -- z |-$>$ z*z is a function RR -$>$ RR and is CONTINUOUS at every real point, PROVEN.}\par
\vnbprop{sq-fun-in-fun}{\vnbstmt{(\lambda z \in \mathbb{R}.\; (z \cdot z)) \in (\mathbb{R} \to \mathbb{R})}}
\vnbprop{sq-continuous-at}{\vnbstmt{\forall t \in \mathbb{R}.\; \operatorname{is\text{-}continuous\text{-}at}(\operatorname{rr\text{-}ms}, \operatorname{rr\text{-}ms}, (\lambda z \in \mathbb{R}.\; (z \cdot z)), t)}}

\vnbbookfile{theorem-library/sqrt-defined}{theorem-library/sqrt-defined}
\noindent{\small sqrt-defined.scm -- SQRT DEFINED by the intermediate value theorem, and the five supports that used to characterise it RETIRED as theorems.}\par
\vnbprop{sqrt-exists}{\vnbstmt{\forall a \in \mathbb{R}.\; ((0 \leq a) \Rightarrow \exists x \in \mathbb{R}.\; ((0 \leq x) \wedge ((x \cdot x) = a)))}}
\vnbprop{sqrt-unique}{\vnbstmt{\begin{array}{@{}l@{}} \forall x, y. \\ \quad (((x \in \mathbb{R}) \wedge (y \in \mathbb{R})) \wedge (((0 \leq x) \wedge (0 \leq y)) \wedge ((x \cdot x) = (y \cdot y)))) \Rightarrow (x = y) \end{array}}}
\vnbprop{sqrt-prop}{\vnbstmt{\begin{array}{@{}l@{}} \forall a \in \mathbb{R}. \\ \quad (0 \leq a) \Rightarrow ((\sqrt{a} \in \mathbb{R}) \wedge ((0 \leq \sqrt{a}) \wedge ((\sqrt{a} \cdot \sqrt{a}) = a))) \end{array}}}
\vnbprop{sqrt-char}{\vnbstmt{\begin{array}{@{}l@{}} \forall a, x. \\ \quad (((a \in \mathbb{R}) \wedge (0 \leq a)) \wedge (((x \in \mathbb{R}) \wedge (0 \leq x)) \wedge ((x \cdot x) = a))) \Rightarrow (\sqrt{a} = x) \end{array}}}
\vnbprop{sqrt-nonneg}{\vnbstmt{\forall a \in \mathbb{R}.\; ((0 \leq a) \Rightarrow ((\sqrt{a} \in \mathbb{R}) \wedge (0 \leq \sqrt{a})))}}
\vnbprop{sqrt-sq}{\vnbstmt{\forall a \in \mathbb{R}.\; ((0 \leq a) \Rightarrow ((\sqrt{a} \cdot \sqrt{a}) = a))}}
\vnbprop{sqrt-of-sq}{\vnbstmt{\forall a \in \mathbb{R}.\; (\sqrt{(a \cdot a)} = \lvert a \rvert)}}
\vnbprop{sqrt-mono}{\vnbstmt{\forall a, b \in \mathbb{R}.\; (((0 \leq a) \wedge (a \leq b)) \Rightarrow (\sqrt{a} \leq \sqrt{b}))}}
\vnbprop{sqrt-mul}{\vnbstmt{\begin{array}{@{}l@{}} \forall a, b \in \mathbb{R}. \\ \quad ((0 \leq a) \wedge (0 \leq b)) \Rightarrow (\sqrt{(a \cdot b)} = (\sqrt{a} \cdot \sqrt{b})) \end{array}}}

\vnbbookfile{theorem-library/cc-magnitude}{theorem-library/cc-magnitude}
\noindent{\small cc-magnitude.scm -- the complex modulus, PROVEN from its definition.}\par
\vnbprop{cc-magnitude-closed}{\vnbstmt{\forall z \in \mathbb{C}.\; (\operatorname{magnitude}(z) \in \mathbb{R})}}
\vnbprop{cc-magnitude-nonneg}{\vnbstmt{\forall z \in \mathbb{C}.\; (0 \leq \operatorname{magnitude}(z))}}
\vnbprop{sqrt-zero}{\vnbstmt{\sqrt{0} = 0}}
\vnbprop{rr-le-of-sq-le}{\vnbstmt{\forall t, s.\; (((t \in \mathbb{R}) \wedge ((s \in \mathbb{R}) \wedge ((0 \leq s) \wedge ((t \cdot t) \leq (s \cdot s))))) \Rightarrow (t \leq s))}}
\vnbprop{cc-magnitude-zero-iff}{\vnbstmt{\forall z \in \mathbb{C}.\; ((\operatorname{magnitude}(z) = 0) \Leftrightarrow (z = 0))}}
\vnbprop{cc-magnitude-neg}{\vnbstmt{\forall z \in \mathbb{C}.\; (\operatorname{magnitude}(-z) = \operatorname{magnitude}(z))}}
\vnbprop{rr-magnitude-is-abs}{\vnbstmt{\forall a \in \mathbb{R}.\; (\operatorname{magnitude}(a) = \lvert a \rvert)}}
\vnbprop{cc-magnitude-mul}{\vnbstmt{\begin{array}{@{}l@{}} \forall a, b. \\ \quad ((a \in \mathbb{C}) \wedge (b \in \mathbb{C})) \Rightarrow (\operatorname{magnitude}((a \cdot b)) = (\operatorname{magnitude}(a) \cdot \operatorname{magnitude}(b))) \end{array}}}
\vnbprop{cc-magnitude-triangle}{\vnbstmt{\begin{array}{@{}l@{}} \forall a, b. \\ \quad ((a \in \mathbb{C}) \wedge (b \in \mathbb{C})) \Rightarrow (\operatorname{magnitude}((a + b)) \leq (\operatorname{magnitude}(a) + \operatorname{magnitude}(b))) \end{array}}}

\vnbbookfile{theorem-library/cc-normed-field}{theorem-library/cc-normed-field}
\noindent{\small cc-normed-field.scm -- IS-NORMED-FIELD(CC-NORMED-FIELD), PROVED.}\par
\vnbprop{cc-is-normed-field}{\vnbstmt{\operatorname{is\text{-}normed\text{-}field}(\operatorname{cc\text{-}normed\text{-}field})}}

\vnbbookfile{theorem-library/cc-metric-space-proof}{theorem-library/cc-metric-space-proof}
\noindent{\small cc-metric-space-proof.scm -- IS-METRIC-SPACE(CC-MS), PROVEN.}\par
\vnbprop{cc-sub-in-cc}{\vnbstmt{\forall a, b.\; (((a \in \mathbb{C}) \wedge (b \in \mathbb{C})) \Rightarrow ((a - b) \in \mathbb{C}))}}
\vnbprop{cc-sub-zero-eq}{\vnbstmt{\forall a, b.\; (((a \in \mathbb{C}) \wedge ((b \in \mathbb{C}) \wedge ((a - b) = 0))) \Rightarrow (a = b))}}
\vnbprop{cc-is-metric-space}{\vnbstmt{\operatorname{is\text{-}metric\text{-}space}(\operatorname{cc\text{-}ms})}}

\vnbbookfile{theorem-library/complex-inner-product-laws}{theorem-library/complex-inner-product-laws}
\noindent{\small complex-inner-product-laws.scm -- what a sesquilinear form does in its SECOND argument, and the two expansion identities that follow.  All PROVEN modulo 0.}\par
\vnbprop{cc-conjugate-neg}{\vnbstmt{\forall a \in \mathbb{C}.\; (\operatorname{conjugate}(-a) = -\operatorname{conjugate}(a))}}
\vnbprop{cips-ip-add-right}{\vnbstmt{\begin{array}{@{}l@{}} \forall v.; \forall x, y, z \in \operatorname{vec}(v). \\ \quad \operatorname{is\text{-}complex\text{-}inner\text{-}product\text{-}space}(v) \Rightarrow  \\ \quad \quad \langle\cdot,\cdot\rangle_{v}(x, \operatorname{vadd}(v)(y, z)) = (\langle\cdot,\cdot\rangle_{v}(x, y) + \langle\cdot,\cdot\rangle_{v}(x, z)) \end{array}}}
\vnbprop{cips-ip-conj-homog-right}{\vnbstmt{\begin{array}{@{}l@{}} \forall v.; \forall a \in \mathbb{C}, x, y \in \operatorname{vec}(v). \\ \quad \operatorname{is\text{-}complex\text{-}inner\text{-}product\text{-}space}(v) \Rightarrow  \\ \quad \quad \langle\cdot,\cdot\rangle_{v}(x, \operatorname{act}(v)(a, y)) = (\operatorname{conjugate}(a) \cdot \langle\cdot,\cdot\rangle_{v}(x, y)) \end{array}}}
\vnbprop{cips-ip-zero-right}{\vnbstmt{\begin{array}{@{}l@{}} \forall v.; \forall x \in \operatorname{vec}(v). \\ \quad \operatorname{is\text{-}complex\text{-}inner\text{-}product\text{-}space}(v) \Rightarrow (\langle\cdot,\cdot\rangle_{v}(x, \operatorname{vzero}(v)) = 0) \end{array}}}
\vnbprop{cips-expand-sum}{\vnbstmt{\begin{array}{@{}l@{}} \forall v.; \forall x, y \in \operatorname{vec}(v). \\ \quad \operatorname{is\text{-}complex\text{-}inner\text{-}product\text{-}space}(v) \Rightarrow  \\ \quad \quad \langle\cdot,\cdot\rangle_{v}(\operatorname{vadd}(v)(x, y), \operatorname{vadd}(v)(x, y)) \\ \quad \quad =\; (\langle\cdot,\cdot\rangle_{v}(x, x) + \langle\cdot,\cdot\rangle_{v}(x, y)) + (\operatorname{conjugate}(\langle\cdot,\cdot\rangle_{v}(x, y)) + \langle\cdot,\cdot\rangle_{v}(y, y)) \end{array}}}
\vnbprop{cips-expand-two}{\vnbstmt{\begin{array}{@{}l@{}} \forall v, a, b, x, y. \\ \quad (\operatorname{is\text{-}complex\text{-}inner\text{-}product\text{-}space}(v) \wedge ((a \in \mathbb{C}) \wedge ((b \in \mathbb{C}) \wedge ((x \in \operatorname{vec}(v)) \wedge (y \in \operatorname{vec}(v)))))) \Rightarrow  \\ \quad \quad \langle\cdot,\cdot\rangle_{v}(\operatorname{vadd}(v)(\operatorname{act}(v)(a, x), \operatorname{act}(v)(b, y)), \operatorname{vadd}(v)(\operatorname{act}(v)(a, x), \operatorname{act}(v)(b, y))) \\ \quad \quad =\; (((a \cdot \operatorname{conjugate}(a)) \cdot \langle\cdot,\cdot\rangle_{v}(x, x)) + ((a \cdot \operatorname{conjugate}(b)) \cdot \langle\cdot,\cdot\rangle_{v}(x, y))) + (((b \cdot \operatorname{conjugate}(a)) \cdot \operatorname{conjugate}(\langle\cdot,\cdot\rangle_{v}(x, y))) + ((b \cdot \operatorname{conjugate}(b)) \cdot \langle\cdot,\cdot\rangle_{v}(y, y))) \end{array}}}

\vnbbookfile{theorem-library/inner-product-inequalities}{theorem-library/inner-product-inequalities}
\noindent{\small inner-product-inequalities.scm -- CAUCHY-SCHWARZ and MINKOWSKI at the INNER-PRODUCT level, over a COMPLEX-INNER-PRODUCT-SPACE.  Both PROVEN modulo 0.}\par
\vnbprop{cips-schwarz}{\vnbstmt{\begin{array}{@{}l@{}} \forall v, x, y. \\ \quad (\operatorname{is\text{-}complex\text{-}inner\text{-}product\text{-}space}(v) \wedge ((x \in \operatorname{vec}(v)) \wedge (y \in \operatorname{vec}(v)))) \Rightarrow  \\ \quad \quad (\langle\cdot,\cdot\rangle_{v}(x, y) \cdot \operatorname{conjugate}(\langle\cdot,\cdot\rangle_{v}(x, y))) \leq (\langle\cdot,\cdot\rangle_{v}(x, x) \cdot \langle\cdot,\cdot\rangle_{v}(y, y)) \end{array}}}
\vnbprop{cips-minkowski}{\vnbstmt{\begin{array}{@{}l@{}} \forall v, x, y. \\ \quad (\operatorname{is\text{-}complex\text{-}inner\text{-}product\text{-}space}(v) \wedge ((x \in \operatorname{vec}(v)) \wedge (y \in \operatorname{vec}(v)))) \Rightarrow  \\ \quad \quad \lVert \operatorname{vadd}(v)(x, y) \rVert \leq (\lVert x \rVert + \lVert y \rVert) \end{array}}}

\vnbbookfile{theorem-library/rake-ip-normed-ag}{theorem-library/rake-ip-normed-ag}
\noindent{\small rake-ip-normed-ag.scm -- rake batch 5b E: the induced norm of a complex inner product makes the additive group of its vectors a NORMED ABELIAN GROUP.  The support proved here is}\par
\vnbprop{ip-norm-unfold}{\vnbstmt{\forall v, x.\; (\lVert x \rVert \simeq \sqrt{\langle\cdot,\cdot\rangle_{v}(x, x)})}}
\vnbprop{cips-vneg-type}{\vnbstmt{\begin{array}{@{}l@{}} \forall v.; \forall x \in \operatorname{vec}(v). \\ \quad \operatorname{is\text{-}complex\text{-}inner\text{-}product\text{-}space}(v) \Rightarrow (\operatorname{vneg}(v)(x) \in \operatorname{vec}(v)) \end{array}}}
\vnbprop{cips-vadd-vneg}{\vnbstmt{\begin{array}{@{}l@{}} \forall v.; \forall x \in \operatorname{vec}(v). \\ \quad \operatorname{is\text{-}complex\text{-}inner\text{-}product\text{-}space}(v) \Rightarrow (\operatorname{vadd}(v)(x, \operatorname{vneg}(v)(x)) = \operatorname{vzero}(v)) \end{array}}}
\vnbprop{cips-ip-zero-left}{\vnbstmt{\begin{array}{@{}l@{}} \forall v.; \forall z \in \operatorname{vec}(v). \\ \quad \operatorname{is\text{-}complex\text{-}inner\text{-}product\text{-}space}(v) \Rightarrow (\langle\cdot,\cdot\rangle_{v}(\operatorname{vzero}(v), z) = 0) \end{array}}}
\vnbprop{cips-ip-neg-left}{\vnbstmt{\begin{array}{@{}l@{}} \forall v.; \forall x, z \in \operatorname{vec}(v). \\ \quad \operatorname{is\text{-}complex\text{-}inner\text{-}product\text{-}space}(v) \Rightarrow  \\ \quad \quad \langle\cdot,\cdot\rangle_{v}(\operatorname{vneg}(v)(x), z) = -\langle\cdot,\cdot\rangle_{v}(x, z) \end{array}}}
\vnbprop{cips-ip-neg-self}{\vnbstmt{\begin{array}{@{}l@{}} \forall v.; \forall x \in \operatorname{vec}(v). \\ \quad \operatorname{is\text{-}complex\text{-}inner\text{-}product\text{-}space}(v) \Rightarrow  \\ \quad \quad \langle\cdot,\cdot\rangle_{v}(\operatorname{vneg}(v)(x), \operatorname{vneg}(v)(x)) = \langle\cdot,\cdot\rangle_{v}(x, x) \end{array}}}
\vnbprop{ip-normed-ag-is-normed-ag}{\vnbstmt{\begin{array}{@{}l@{}} \forall v. \\ \quad \operatorname{is\text{-}complex\text{-}inner\text{-}product\text{-}space}(v) \Rightarrow \operatorname{is\text{-}normed\text{-}ag}(\operatorname{ip\text{-}normed\text{-}ag}(v)) \end{array}}}

\vnbbookfile{theorem-library/monotone-convergence-proof}{theorem-library/monotone-convergence-proof}
\noindent{\small monotone-convergence-proof.scm -- MONOTONE CONVERGENCE on RR, PROVEN.}\par
\vnbprop{nn-monotone-step-implies-le}{\vnbstmt{\begin{array}{@{}l@{}} \forall f \in (\mathbb{N} \to \mathbb{R}), n, m \in \mathbb{N}. \\ \quad (\forall k \in \mathbb{N}.\; (f(k) \leq f((k + 1))) \wedge (m \leq n)) \Rightarrow (f(m) \leq f(n)) \end{array}}}
\vnbprop{monotone-convergence-rr}{\vnbstmt{\begin{array}{@{}l@{}} \forall f \in (\mathbb{N} \to \mathbb{R}). \\ \quad (\forall k \in \mathbb{N}.\; (f(k) \leq f((k + 1))) \wedge \exists d \in \mathbb{R}.\; \forall k \in \mathbb{N}.\; (f(k) \leq d)) \Rightarrow  \\ \quad \quad \operatorname{converges}(\operatorname{rr\text{-}ms}, f) \end{array}}}

\vnbbookfile{theorem-library/comparison-test-proof}{theorem-library/comparison-test-proof}
\noindent{\small comparison-test-proof.scm -- THE COMPARISON TEST for real series, PROVEN, together with the two order lemmas it rests on and the readouts that make a real partial sum computable.}\par
\vnbprop{rr-additive-ag-opr}{\vnbstmt{(\operatorname{opr}(\operatorname{normed\text{-}field\text{-}additive\text{-}ag}(\operatorname{rr\text{-}normed\text{-}field})) \simeq (\lambda {[}\operatorname{x\_}, \operatorname{y\_}] \in (\mathbb{R} \times \mathbb{R}).\; (\operatorname{x\_} + \operatorname{y\_})))}}
\vnbprop{rr-additive-ag-iden}{\vnbstmt{(\operatorname{iden}(\operatorname{normed\text{-}field\text{-}additive\text{-}ag}(\operatorname{rr\text{-}normed\text{-}field})) \simeq 0)}}
\vnbprop{series-partial-sum-zero}{\vnbstmt{\forall f.\; (\operatorname{series\text{-}partial\text{-}sum}(f, 0) \simeq 0)}}
\vnbprop{sum-ag-rr-in-rr}{\vnbstmt{\begin{array}{@{}l@{}} \forall k \in \mathbb{N}, f \in (\mathbb{N} \to \mathbb{R}). \\ \quad \operatorname{sum\text{-}ag}(\operatorname{normed\text{-}field\text{-}additive\text{-}ag}(\operatorname{rr\text{-}normed\text{-}field}), f, k) \in \mathbb{R} \end{array}}}
\vnbprop{series-partial-sum-in-rr}{\vnbstmt{\forall k \in \mathbb{N}, f \in (\mathbb{N} \to \mathbb{R}).\; (\operatorname{series\text{-}partial\text{-}sum}(f, k) \in \mathbb{R})}}
\vnbprop{series-partial-sum-unfold}{\vnbstmt{\begin{array}{@{}l@{}} \forall f, k. \\ \quad (\operatorname{series\text{-}partial\text{-}sum}(f, k) \simeq \operatorname{sum\text{-}ag}(\operatorname{normed\text{-}field\text{-}additive\text{-}ag}(\operatorname{rr\text{-}normed\text{-}field}), f, k)) \end{array}}}
\vnbprop{series-partial-sum-succ}{\vnbstmt{\begin{array}{@{}l@{}} \forall f.; \forall k \in \mathbb{N}. \\ \quad ((\operatorname{series\text{-}partial\text{-}sum}(f, k) \in \mathbb{R}) \wedge (f(k) \in \mathbb{R})) \Rightarrow  \\ \quad \quad (\operatorname{series\text{-}partial\text{-}sum}(f, (k + 1)) \simeq (\operatorname{series\text{-}partial\text{-}sum}(f, k) + f(k))) \end{array}}}
\vnbprop{series-partial-sum-monotone-nonneg}{\vnbstmt{\begin{array}{@{}l@{}} \forall f \in (\mathbb{N} \to \mathbb{R}), k \in \mathbb{N}. \\ \quad \forall n \in \mathbb{N}.\; (0 \leq f(n)) \Rightarrow  \\ \quad \quad \operatorname{series\text{-}partial\text{-}sum}(f, k) \leq \operatorname{series\text{-}partial\text{-}sum}(f, (k + 1)) \end{array}}}
\vnbprop{series-partial-sum-le-termwise}{\vnbstmt{\begin{array}{@{}l@{}} \forall f, g \in (\mathbb{N} \to \mathbb{R}), k \in \mathbb{N}. \\ \quad \forall n \in \mathbb{N}.\; (f(n) \leq g(n)) \Rightarrow  \\ \quad \quad \operatorname{series\text{-}partial\text{-}sum}(f, k) \leq \operatorname{series\text{-}partial\text{-}sum}(g, k) \end{array}}}
\vnbprop{monotone-convergent-bounded-above}{\vnbstmt{\begin{array}{@{}l@{}} \forall f \in (\mathbb{N} \to \mathbb{R}). \\ \quad (\forall k \in \mathbb{N}.\; (f(k) \leq f((k + 1))) \wedge \operatorname{converges}(\operatorname{rr\text{-}ms}, f)) \Rightarrow  \\ \quad \quad \exists d \in \mathbb{R}.\; \forall k \in \mathbb{N}.\; (f(k) \leq d) \end{array}}}
\vnbprop{series-partial-sum-seq-apply}{\vnbstmt{\begin{array}{@{}l@{}} \forall f.; \forall j \in \mathbb{N}. \\ \quad ((\lambda k \in \mathbb{N}.\; \operatorname{series\text{-}partial\text{-}sum}(f, k))(j) \simeq \operatorname{series\text{-}partial\text{-}sum}(f, j)) \end{array}}}
\vnbprop{series-partial-sum-seq-in-fun}{\vnbstmt{\begin{array}{@{}l@{}} \forall f \in (\mathbb{N} \to \mathbb{R}). \\ \quad (\lambda k \in \mathbb{N}.\; \operatorname{series\text{-}partial\text{-}sum}(f, k)) \in (\mathbb{N} \to \mathbb{R}) \end{array}}}
\vnbprop{abs-seq-in-fun}{\vnbstmt{\forall f \in (\mathbb{N} \to \mathbb{R}).\; ((\lambda n \in \mathbb{N}.\; \lvert f(n) \rvert) \in (\mathbb{N} \to \mathbb{R}))}}
\vnbprop{comparison-test}{\vnbstmt{\begin{array}{@{}l@{}} \forall f, g \in (\mathbb{N} \to \mathbb{R}). \\ \quad (\forall n \in \mathbb{N}.\; ((0 \leq f(n)) \wedge (f(n) \leq g(n))) \wedge \operatorname{series\text{-}converges}(g)) \Rightarrow  \\ \quad \quad \operatorname{series\text{-}converges}(f) \end{array}}}
\vnbprop{series-partial-sum-abs-le}{\vnbstmt{\begin{array}{@{}l@{}} \forall k \in \mathbb{N}, f \in (\mathbb{N} \to \mathbb{R}). \\ \quad \lvert \operatorname{series\text{-}partial\text{-}sum}(f, k) \rvert \leq \operatorname{series\text{-}partial\text{-}sum}((\lambda n \in \mathbb{N}.\; \lvert f(n) \rvert), k) \end{array}}}

\vnbbookfile{theorem-library/series-abs-triangle}{theorem-library/series-abs-triangle}
\noindent{\small series-abs-triangle.scm -- |Sum\_\{n$<$k\} f(n)| $<$= Sum\_\{n$<$k\} |f(n)|, PROVEN, and `abs(0) = 0', which the tree did not state.}\par
\vnbprop{series-abs-triangle}{\vnbstmt{\begin{array}{@{}l@{}} \forall k \in \mathbb{N}, f, g \in (\mathbb{N} \to \mathbb{R}). \\ \quad \forall j \in \mathbb{N}.\; (g(j) = \lvert f(j) \rvert) \Rightarrow  \\ \quad \quad \lvert \operatorname{series\text{-}partial\text{-}sum}(f, k) \rvert \leq \operatorname{series\text{-}partial\text{-}sum}(g, k) \end{array}}}

\vnbbookfile{theorem-library/series-block-abs}{theorem-library/series-block-abs}
\noindent{\small series-block-abs.scm -- THE BLOCK TRIANGLE INEQUALITY, in the form a Cauchy argument can use, and the NN gap lemma under it.}\par
\vnbprop{series-abs-triangle-block}{\vnbstmt{\begin{array}{@{}l@{}} \forall d, m \in \mathbb{N}, f, g \in (\mathbb{N} \to \mathbb{R}). \\ \quad \forall j \in \mathbb{N}.\; (g(j) = \lvert f(j) \rvert) \Rightarrow  \\ \quad \quad \lvert (\operatorname{series\text{-}partial\text{-}sum}(f, (m + d)) - \operatorname{series\text{-}partial\text{-}sum}(f, m)) \rvert \\ \quad \quad \leq\; (\operatorname{series\text{-}partial\text{-}sum}(g, (m + d)) - \operatorname{series\text{-}partial\text{-}sum}(g, m)) \end{array}}}
\vnbprop{series-abs-triangle-le}{\vnbstmt{\begin{array}{@{}l@{}} \forall k, m \in \mathbb{N}, f, g \in (\mathbb{N} \to \mathbb{R}). \\ \quad (\forall j \in \mathbb{N}.\; (g(j) = \lvert f(j) \rvert) \wedge (m \leq k)) \Rightarrow  \\ \quad \quad \lvert (\operatorname{series\text{-}partial\text{-}sum}(f, k) - \operatorname{series\text{-}partial\text{-}sum}(f, m)) \rvert \\ \quad \quad \leq\; (\operatorname{series\text{-}partial\text{-}sum}(g, k) - \operatorname{series\text{-}partial\text{-}sum}(g, m)) \end{array}}}

\vnbbookfile{theorem-library/series-linearity}{theorem-library/series-linearity}
\noindent{\small series-linearity.scm -- LINEARITY OF THE REAL PARTIAL SUM, in POINTWISE form.}\par
\vnbprop{series-partial-sum-in-rr-ptwise}{\vnbstmt{\begin{array}{@{}l@{}} \forall n \in \mathbb{N}, f. \\ \quad \forall k \in \mathbb{N}.\; (f(k) \in \mathbb{R}) \Rightarrow (\operatorname{series\text{-}partial\text{-}sum}(f, n) \in \mathbb{R}) \end{array}}}
\vnbprop{series-partial-sum-add-ptwise}{\vnbstmt{\begin{array}{@{}l@{}} \forall n \in \mathbb{N}, f, g, h. \\ \quad \forall k \in \mathbb{N}.\; (f(k) \in \mathbb{R}) \wedge \\ \quad \forall k \in \mathbb{N}.\; (g(k) \in \mathbb{R}) \wedge \\ \quad \forall k \in \mathbb{N}.\; (h(k) = (f(k) + g(k))) \Rightarrow \\ \quad \quad \operatorname{series\text{-}partial\text{-}sum}(h, n) = (\operatorname{series\text{-}partial\text{-}sum}(f, n) + \operatorname{series\text{-}partial\text{-}sum}(g, n)) \end{array}}}
\vnbprop{series-partial-sum-scale-ptwise}{\vnbstmt{\begin{array}{@{}l@{}} \forall n \in \mathbb{N}, c, f, h. \\ \quad ((c \in \mathbb{R}) \wedge (\forall k \in \mathbb{N}.\; (f(k) \in \mathbb{R}) \wedge \forall k \in \mathbb{N}.\; (h(k) = (c \cdot f(k))))) \Rightarrow  \\ \quad \quad \operatorname{series\text{-}partial\text{-}sum}(h, n) = (c \cdot \operatorname{series\text{-}partial\text{-}sum}(f, n)) \end{array}}}
\vnbprop{series-partial-sum-weighted-expansion}{\vnbstmt{\begin{array}{@{}l@{}} \forall n \in \mathbb{N}, x, y, g, w, p, q, s. \\ \quad x \in \mathbb{R} \wedge \\ \quad y \in \mathbb{R} \wedge \\ \quad \forall k \in \mathbb{Z}.\; (g(k) \in \mathbb{R}) \wedge \\ \quad \forall k \in \mathbb{Z}.\; (w(k) \in \mathbb{R}) \wedge \\ \quad \forall k \in \mathbb{N}.\; (q(k) \in \mathbb{R}) \wedge \\ \quad \forall k \in \mathbb{N}.\; (s(k) \in \mathbb{R}) \wedge \\ \quad \forall k \in \mathbb{N}.\; (p(k) = (w(k) \cdot ((x \cdot g((k - 1))) + (y \cdot g(k))))) \wedge \\ \quad \forall k \in \mathbb{N}.\; (q(k) = (w((k + 1)) \cdot g(k))) \wedge \\ \quad \forall k \in \mathbb{N}.\; (s(k) = (w(k) \cdot g(k))) \Rightarrow \\ \quad \quad \operatorname{series\text{-}partial\text{-}sum}(p, (n + 1)) \\ \quad \quad =\; (x \cdot ((w(0) \cdot g(-1)) + \operatorname{series\text{-}partial\text{-}sum}(q, n))) + (y \cdot (\operatorname{series\text{-}partial\text{-}sum}(s, n) + (w(n) \cdot g(n)))) \end{array}}}
\vnbprop{series-partial-sum-le-termwise-ptwise}{\vnbstmt{\begin{array}{@{}l@{}} \forall n \in \mathbb{N}, f, g. \\ \quad (\forall k \in \mathbb{N}.\; (f(k) \in \mathbb{R}) \wedge (\forall k \in \mathbb{N}.\; (g(k) \in \mathbb{R}) \wedge \forall k \in \mathbb{N}.\; (f(k) \leq g(k)))) \Rightarrow  \\ \quad \quad \operatorname{series\text{-}partial\text{-}sum}(f, n) \leq \operatorname{series\text{-}partial\text{-}sum}(g, n) \end{array}}}

\vnbbookfile{theorem-library/bernstein-basis}{theorem-library/bernstein-basis}
\noindent{\small bernstein-basis.scm -- the Bernstein basis polynomials, and the first of the two identities Proposition 5.4 of docs/calculus.pdf asks for.}\par
\par\medskip\noindent\textbf{Definition} (\texttt{bernstein-basis}).\;$\operatorname{bernstein\text{-}basis}(n, x) :=$
\begin{equation*}
\operatorname{comb\text{-}kk}(\operatorname{normed\text{-}field\text{-}as\text{-}commutative\text{-}ring}(\operatorname{rr\text{-}normed\text{-}field}), x, (1 - x), n)
\end{equation*}

\vnbprop{ring-power-of-one}{\vnbstmt{\begin{array}{@{}l@{}} \forall n \in \mathbb{N}, r. \\ \quad \operatorname{is\text{-}commutative\text{-}ring}(r) \Rightarrow (\operatorname{ring\text{-}power}(r, \operatorname{one}(r), n) = \operatorname{one}(r)) \end{array}}}
\vnbprop{series-partial-sum-is-ring-sum}{\vnbstmt{\begin{array}{@{}l@{}} \forall n \in \mathbb{N}, f. \\ \quad (\operatorname{series\text{-}partial\text{-}sum}(f, n) \simeq \operatorname{sum}(\operatorname{normed\text{-}field\text{-}as\text{-}commutative\text{-}ring}(\operatorname{rr\text{-}normed\text{-}field}), f, n)) \end{array}}}
\vnbprop{binomial-sum-value}{\vnbstmt{\begin{array}{@{}l@{}} \forall n, r, x, y. \\ \quad ((n \in \mathbb{N}) \wedge (\operatorname{is\text{-}commutative\text{-}ring}(r) \wedge ((x \in \operatorname{carr}(r)) \wedge (y \in \operatorname{carr}(r))))) \Rightarrow  \\ \quad \quad \operatorname{sum}(r, \operatorname{comb\text{-}kk}(r, x, y, n), (n + 1)) = \operatorname{ring\text{-}power}(r, \operatorname{add}(r)(x, y), n) \end{array}}}
\vnbprop{bernstein-basis-unfold}{\vnbstmt{\begin{array}{@{}l@{}} \forall n, x. \\ \quad (\operatorname{bernstein\text{-}basis}(n, x) \simeq \operatorname{comb\text{-}kk}(\operatorname{normed\text{-}field\text{-}as\text{-}commutative\text{-}ring}(\operatorname{rr\text{-}normed\text{-}field}), x, (1 - x), n)) \end{array}}}
\vnbprop{bernstein-basis-in-fun}{\vnbstmt{\forall n \in \mathbb{N}, x \in \mathbb{R}.\; (\operatorname{bernstein\text{-}basis}(n, x) \in (\mathbb{Z} \to \mathbb{R}))}}
\vnbprop{bernstein-partition}{\vnbstmt{\begin{array}{@{}l@{}} \forall n \in \mathbb{N}, x \in \mathbb{R}. \\ \quad \operatorname{series\text{-}partial\text{-}sum}(\operatorname{bernstein\text{-}basis}(n, x), (n + 1)) = 1 \end{array}}}

\vnbbookfile{theorem-library/bernstein-moments}{theorem-library/bernstein-moments}
\noindent{\small bernstein-moments.scm -- the moments of the Bernstein basis, and identity (76) of docs/calculus.pdf Proposition 5.4.}\par
\vnbprop{bernstein-basis-ptwise-in-rr}{\vnbstmt{\begin{array}{@{}l@{}} \forall n, x, k. \\ \quad ((n \in \mathbb{N}) \wedge ((x \in \mathbb{R}) \wedge (k \in \mathbb{Z}))) \Rightarrow (\operatorname{bernstein\text{-}basis}(n, x)(k) \in \mathbb{R}) \end{array}}}
\vnbprop{bernstein-basis-succ}{\vnbstmt{\begin{array}{@{}l@{}} \forall n, x, k. \\ \quad ((n \in \mathbb{N}) \wedge ((x \in \mathbb{R}) \wedge (k \in \mathbb{N}))) \Rightarrow  \\ \quad \quad \operatorname{bernstein\text{-}basis}((n + 1), x)(k) \\ \quad \quad =\; (x \cdot \operatorname{bernstein\text{-}basis}(n, x)((k - 1))) + ((1 - x) \cdot \operatorname{bernstein\text{-}basis}(n, x)(k)) \end{array}}}
\vnbprop{bernstein-basis-null}{\vnbstmt{\begin{array}{@{}l@{}} \forall n, x, k. \\ \quad ((n \in \mathbb{N}) \wedge ((x \in \mathbb{R}) \wedge ((k \in \mathbb{Z}) \wedge (k < 0)))) \Rightarrow (\operatorname{bernstein\text{-}basis}(n, x)(k) = 0) \end{array}}}
\vnbprop{bernstein-basis-above}{\vnbstmt{\begin{array}{@{}l@{}} \forall n, x, k. \\ \quad ((n \in \mathbb{N}) \wedge ((x \in \mathbb{R}) \wedge ((k \in \mathbb{Z}) \wedge (n < k)))) \Rightarrow (\operatorname{bernstein\text{-}basis}(n, x)(k) = 0) \end{array}}}
\vnbprop{bernstein-weighted-step}{\vnbstmt{\begin{array}{@{}l@{}} \forall n, x, w, p, q, s. \\ \quad n \in \mathbb{N} \wedge \\ \quad x \in \mathbb{R} \wedge \\ \quad \forall k \in \mathbb{Z}.\; (w(k) \in \mathbb{R}) \wedge \\ \quad \forall k \in \mathbb{N}.\; (q(k) \in \mathbb{R}) \wedge \\ \quad \forall k \in \mathbb{N}.\; (s(k) \in \mathbb{R}) \wedge \\ \quad \forall k \in \mathbb{N}.\; (p(k) = (w(k) \cdot \operatorname{bernstein\text{-}basis}((n + 1), x)(k))) \wedge \\ \quad \forall k \in \mathbb{N}.\; (q(k) = (w((k + 1)) \cdot \operatorname{bernstein\text{-}basis}(n, x)(k))) \wedge \\ \quad \forall k \in \mathbb{N}.\; (s(k) = (w(k) \cdot \operatorname{bernstein\text{-}basis}(n, x)(k))) \Rightarrow \\ \quad \quad \operatorname{series\text{-}partial\text{-}sum}(p, ((n + 1) + 1)) \\ \quad \quad =\; (x \cdot \operatorname{series\text{-}partial\text{-}sum}(q, (n + 1))) + ((1 - x) \cdot \operatorname{series\text{-}partial\text{-}sum}(s, (n + 1))) \end{array}}}
\vnbprop{bernstein-moment-1}{\vnbstmt{\begin{array}{@{}l@{}} \forall n \in \mathbb{N}, x, h. \\ \quad ((x \in \mathbb{R}) \wedge ((h \in (\mathbb{N} \to \mathbb{R})) \wedge \forall k \in \mathbb{N}.\; (h(k) = (k \cdot \operatorname{bernstein\text{-}basis}(n, x)(k))))) \Rightarrow  \\ \quad \quad \operatorname{series\text{-}partial\text{-}sum}(h, (n + 1)) = (n \cdot x) \end{array}}}
\vnbprop{bernstein-moment-2}{\vnbstmt{\begin{array}{@{}l@{}} \forall n \in \mathbb{N}, x, h. \\ \quad x \in \mathbb{R} \wedge \\ \quad h \in (\mathbb{N} \to \mathbb{R}) \wedge \\ \quad \forall k \in \mathbb{N}.\; (h(k) = ((k \cdot (k - 1)) \cdot \operatorname{bernstein\text{-}basis}(n, x)(k))) \Rightarrow \\ \quad \quad \operatorname{series\text{-}partial\text{-}sum}(h, (n + 1)) = ((n \cdot (n - 1)) \cdot (x \cdot x)) \end{array}}}
\vnbprop{bernstein-variance}{\vnbstmt{\begin{array}{@{}l@{}} \forall n, x, h. \\ \quad n \in \mathbb{N} \wedge \\ \quad x \in \mathbb{R} \wedge \\ \quad \forall k \in \mathbb{N}. \\ \quad \quad h(k) = (((k - (n \cdot x)) \cdot (k - (n \cdot x))) \cdot \operatorname{bernstein\text{-}basis}(n, x)(k)) \Rightarrow \\ \quad \quad \operatorname{series\text{-}partial\text{-}sum}(h, (n + 1)) = ((n \cdot x) \cdot (1 - x)) \end{array}}}
\vnbprop{bernstein-variance-bound}{\vnbstmt{\begin{array}{@{}l@{}} \forall n, x, h. \\ \quad n \in \mathbb{N} \wedge \\ \quad \neg (n = 0) \wedge \\ \quad x \in \mathbb{R} \wedge \\ \quad \forall k \in \mathbb{N}. \\ \quad \quad h(k) = (((k - (n \cdot x)) \cdot (k - (n \cdot x))) \cdot \operatorname{bernstein\text{-}basis}(n, x)(k)) \Rightarrow \\ \quad \quad (/(1, n) \cdot \operatorname{series\text{-}partial\text{-}sum}(h, (n + 1))) \leq 1 \end{array}}}

\vnbbookfile{theorem-library/series-cauchy-proof}{theorem-library/series-cauchy-proof}
\noindent{\small series-cauchy-proof.scm -- THE CAUCHY CRITERION FOR REAL SERIES, PROVEN.}\par
\vnbprop{series-cauchy-criterion}{\vnbstmt{\begin{array}{@{}l@{}} \forall f \in (\mathbb{N} \to \mathbb{R}), s. \\ \quad (\operatorname{series\text{-}converges}(f) \wedge \operatorname{pos\text{-}rr}(s)) \Rightarrow  \\ \quad \quad \exists d \in \mathbb{N}. \forall m, n \in \mathbb{N}. \\ \quad \quad \quad ((d \leq m) \wedge (m \leq n)) \Rightarrow  \\ \quad \quad \quad \quad \lvert (\operatorname{series\text{-}partial\text{-}sum}(f, n) - \operatorname{series\text{-}partial\text{-}sum}(f, m)) \rvert \leq s \end{array}}}

\vnbbookfile{theorem-library/dominated-convergence}{theorem-library/dominated-convergence}
\noindent{\small dominated-convergence.scm -- DOMINATED CONVERGENCE, PROVEN, in two forms: for real series (the content) and in ELL-ONE (the statement the user asked for), together with the twenty limit / partial-sum / complex-magnitude lemmas the two are assembled from.}\par
\par\medskip\noindent\textbf{Definition} (\texttt{ell-one}).\;
\begin{equation*}
\operatorname{ell\text{-}one} = \{\operatorname{x\_} \in \mathbb{C}^{\mathbb{N}} : \operatorname{series\text{-}converges}((\lambda \operatorname{k\_} \in \mathbb{N}.\; \operatorname{magnitude}(\operatorname{x\_}(\operatorname{k\_}))))\}
\end{equation*}

\par\medskip\noindent\textbf{Definition} (\texttt{series-limit}).\;$\operatorname{series\text{-}limit}(f) :=$ $(\iota\, \operatorname{lm\_}.\;((\operatorname{lm\_} \in \mathbb{R}) \wedge \operatorname{series\text{-}converges\text{-}to}(f, \operatorname{lm\_})))$
\vnbprop{rr-null-sum}{\vnbstmt{\begin{array}{@{}l@{}} \forall f, g, h \in (\mathbb{N} \to \mathbb{R}). \\ \quad (\operatorname{converges\text{-}to}(\operatorname{rr\text{-}ms}, f, 0) \wedge (\operatorname{converges\text{-}to}(\operatorname{rr\text{-}ms}, g, 0) \wedge \forall j \in \mathbb{N}.\; (h(j) = (f(j) + g(j))))) \Rightarrow  \\ \quad \quad \operatorname{converges\text{-}to}(\operatorname{rr\text{-}ms}, h, 0) \end{array}}}
\vnbprop{rr-limit-abs-le}{\vnbstmt{\begin{array}{@{}l@{}} \forall f \in (\mathbb{N} \to \mathbb{R}), v, d \in \mathbb{R}. \\ \quad (\operatorname{converges\text{-}to}(\operatorname{rr\text{-}ms}, f, v) \wedge \forall n \in \mathbb{N}.\; (\lvert f(n) \rvert \leq d)) \Rightarrow  \\ \quad \quad \lvert v \rvert \leq d \end{array}}}
\vnbprop{series-block-le}{\vnbstmt{\begin{array}{@{}l@{}} \forall k, n \in \mathbb{N}, f, g \in (\mathbb{N} \to \mathbb{R}). \\ \quad (\forall i \in \mathbb{N}.\; ((0 \leq f(i)) \wedge (f(i) \leq g(i))) \wedge (n \leq k)) \Rightarrow  \\ \quad \quad (\operatorname{series\text{-}partial\text{-}sum}(f, k) - \operatorname{series\text{-}partial\text{-}sum}(f, n)) \\ \quad \quad \leq\; (\operatorname{series\text{-}partial\text{-}sum}(g, k) - \operatorname{series\text{-}partial\text{-}sum}(g, n)) \end{array}}}
\vnbprop{series-head-null}{\vnbstmt{\begin{array}{@{}l@{}} \forall n \in \mathbb{N}, u.; \forall d \in (\mathbb{N} \to \mathbb{R}). \\ \quad \forall j \in \mathbb{N}.\; (u(j) \in (\mathbb{N} \to \mathbb{R})) \wedge \\ \quad \forall i \in \mathbb{N}.\; \operatorname{converges\text{-}to}(\operatorname{rr\text{-}ms}, (\lambda j \in \mathbb{N}.\; u(j)(i)), 0) \wedge \\ \quad \forall j \in \mathbb{N}.\; (d(j) = \operatorname{series\text{-}partial\text{-}sum}(u(j), n)) \Rightarrow \\ \quad \quad \operatorname{converges\text{-}to}(\operatorname{rr\text{-}ms}, d, 0) \end{array}}}
\vnbprop{series-partial-sum-nonneg}{\vnbstmt{\begin{array}{@{}l@{}} \forall k \in \mathbb{N}, f \in (\mathbb{N} \to \mathbb{R}). \\ \quad \forall i \in \mathbb{N}.\; (0 \leq f(i)) \Rightarrow (0 \leq \operatorname{series\text{-}partial\text{-}sum}(f, k)) \end{array}}}
\vnbprop{series-partial-sum-mono}{\vnbstmt{\begin{array}{@{}l@{}} \forall k, m \in \mathbb{N}, f \in (\mathbb{N} \to \mathbb{R}). \\ \quad (\forall i \in \mathbb{N}.\; (0 \leq f(i)) \wedge (m \leq k)) \Rightarrow  \\ \quad \quad \operatorname{series\text{-}partial\text{-}sum}(f, m) \leq \operatorname{series\text{-}partial\text{-}sum}(f, k) \end{array}}}
\vnbprop{series-tail-small}{\vnbstmt{\begin{array}{@{}l@{}} \forall g \in (\mathbb{N} \to \mathbb{R}), s. \\ \quad (\operatorname{series\text{-}converges}(g) \wedge \operatorname{pos\text{-}rr}(s)) \Rightarrow  \\ \quad \quad \exists d \in \mathbb{N}. \forall k \in \mathbb{N}. \\ \quad \quad \quad (d \leq k) \Rightarrow ((\operatorname{series\text{-}partial\text{-}sum}(g, k) - \operatorname{series\text{-}partial\text{-}sum}(g, d)) \leq s) \end{array}}}
\vnbprop{dominated-null-series}{\vnbstmt{\begin{array}{@{}l@{}} \forall g, e \in (\mathbb{N} \to \mathbb{R}), u. \\ \quad \forall j \in \mathbb{N}.\; (u(j) \in (\mathbb{N} \to \mathbb{R})) \wedge \\ \quad \forall j, i \in \mathbb{N}.\; ((0 \leq u(j)(i)) \wedge (u(j)(i) \leq g(i))) \wedge \\ \quad \operatorname{series\text{-}converges}(g) \wedge \\ \quad \forall i \in \mathbb{N}.\; \operatorname{converges\text{-}to}(\operatorname{rr\text{-}ms}, (\lambda j \in \mathbb{N}.\; u(j)(i)), 0) \wedge \\ \quad \forall j \in \mathbb{N}.\; \operatorname{series\text{-}converges\text{-}to}(u(j), e(j)) \Rightarrow \\ \quad \quad \operatorname{converges\text{-}to}(\operatorname{rr\text{-}ms}, e, 0) \end{array}}}
\vnbprop{series-limit-converges-to}{\vnbstmt{\begin{array}{@{}l@{}} \forall f \in (\mathbb{N} \to \mathbb{R}). \\ \quad \operatorname{series\text{-}converges}(f) \Rightarrow \operatorname{series\text{-}converges\text{-}to}(f, \operatorname{series\text{-}limit}(f)) \end{array}}}
\vnbprop{series-limit-in-rr}{\vnbstmt{\forall f \in (\mathbb{N} \to \mathbb{R}).\; (\operatorname{series\text{-}converges}(f) \Rightarrow (\operatorname{series\text{-}limit}(f) \in \mathbb{R}))}}
\vnbprop{series-partial-sum-add}{\vnbstmt{\begin{array}{@{}l@{}} \forall k \in \mathbb{N}, f, g, h \in (\mathbb{N} \to \mathbb{R}). \\ \quad \forall i \in \mathbb{N}.\; (h(i) = (f(i) + g(i))) \Rightarrow  \\ \quad \quad \operatorname{series\text{-}partial\text{-}sum}(h, k) = (\operatorname{series\text{-}partial\text{-}sum}(f, k) + \operatorname{series\text{-}partial\text{-}sum}(g, k)) \end{array}}}
\vnbprop{series-partial-sum-bounded}{\vnbstmt{\begin{array}{@{}l@{}} \forall f \in (\mathbb{N} \to \mathbb{R}). \\ \quad (\forall i \in \mathbb{N}.\; (0 \leq f(i)) \wedge \operatorname{series\text{-}converges}(f)) \Rightarrow  \\ \quad \quad \exists d \in \mathbb{R}.\; \forall x \in \mathbb{N}.\; (\operatorname{series\text{-}partial\text{-}sum}(f, x) \leq d) \end{array}}}
\vnbprop{series-converges-sum}{\vnbstmt{\begin{array}{@{}l@{}} \forall f, g, h \in (\mathbb{N} \to \mathbb{R}). \\ \quad \forall i \in \mathbb{N}.\; (0 \leq f(i)) \wedge \\ \quad \forall i \in \mathbb{N}.\; (0 \leq g(i)) \wedge \\ \quad \forall i \in \mathbb{N}.\; (h(i) = (f(i) + g(i))) \wedge \\ \quad \operatorname{series\text{-}converges}(f) \wedge \\ \quad \operatorname{series\text{-}converges}(g) \Rightarrow \\ \quad \quad \operatorname{series\text{-}converges}(h) \end{array}}}
\vnbprop{cc-ms-dist}{\vnbstmt{\begin{array}{@{}l@{}} \forall u, v. \\ \quad ((u \in \mathbb{C}) \wedge (v \in \mathbb{C})) \Rightarrow (\operatorname{dist}(\operatorname{cc\text{-}ms})(u, v) \simeq \operatorname{magnitude}((u - v))) \end{array}}}
\vnbprop{cc-magnitude-sub-triangle}{\vnbstmt{\forall p, q \in \mathbb{C}.\; (\operatorname{magnitude}((p - q)) \leq (\operatorname{magnitude}(p) + \operatorname{magnitude}(q)))}}
\vnbprop{cc-magnitude-sub-sym}{\vnbstmt{\forall p, q \in \mathbb{C}.\; (\operatorname{magnitude}((p - q)) = \operatorname{magnitude}((q - p)))}}
\vnbprop{cc-magnitude-le-add}{\vnbstmt{\forall p, q \in \mathbb{C}.\; (\operatorname{magnitude}(p) \leq (\operatorname{magnitude}((p - q)) + \operatorname{magnitude}(q)))}}
\vnbprop{cc-limit-magnitude-le}{\vnbstmt{\begin{array}{@{}l@{}} \forall s \in (\mathbb{N} \to \mathbb{C}), v \in \mathbb{C}, d \in \mathbb{R}. \\ \quad (\operatorname{converges\text{-}to}(\operatorname{cc\text{-}ms}, s, v) \wedge \forall n \in \mathbb{N}.\; (\operatorname{magnitude}(s(n)) \leq d)) \Rightarrow  \\ \quad \quad \operatorname{magnitude}(v) \leq d \end{array}}}
\vnbprop{cc-converges-magnitude-null}{\vnbstmt{\begin{array}{@{}l@{}} \forall s \in (\mathbb{N} \to \mathbb{C}), v \in \mathbb{C}, m \in (\mathbb{N} \to \mathbb{R}). \\ \quad (\operatorname{converges\text{-}to}(\operatorname{cc\text{-}ms}, s, v) \wedge \forall j \in \mathbb{N}.\; (m(j) = \operatorname{magnitude}((s(j) - v)))) \Rightarrow  \\ \quad \quad \operatorname{converges\text{-}to}(\operatorname{rr\text{-}ms}, m, 0) \end{array}}}
\vnbprop{ell-one-membership}{\vnbstmt{\begin{array}{@{}l@{}} \forall y. \\ \quad (y \in \operatorname{ell\text{-}one}) \Leftrightarrow ((y \in \mathbb{C}^{\mathbb{N}}) \wedge \operatorname{series\text{-}converges}((\lambda k \in \mathbb{N}.\; \operatorname{magnitude}(y(k))))) \end{array}}}
\vnbprop{ell-one-dominated-convergence}{\vnbstmt{\begin{array}{@{}l@{}} \forall a \in (\mathbb{N} \to \mathbb{R}), x \in \mathbb{C}^{\mathbb{N}}, e \in (\mathbb{N} \to \mathbb{R}), s. \\ \quad \forall j \in \mathbb{N}.\; (s(j) \in \mathbb{C}^{\mathbb{N}}) \wedge \\ \quad \operatorname{series\text{-}converges}(a) \wedge \\ \quad \forall j, k \in \mathbb{N}.\; (\operatorname{magnitude}(s(j)(k)) \leq a(k)) \wedge \\ \quad \forall k \in \mathbb{N}.\; \operatorname{converges\text{-}to}(\operatorname{cc\text{-}ms}, (\lambda j \in \mathbb{N}.\; s(j)(k)), x(k)) \wedge \\ \quad \forall j \in \mathbb{N}. \\ \quad \quad e(j) = \operatorname{series\text{-}limit}((\lambda k \in \mathbb{N}.\; \operatorname{magnitude}((s(j)(k) - x(k))))) \Rightarrow \\ \quad \quad (x \in \operatorname{ell\text{-}one}) \wedge \operatorname{converges\text{-}to}(\operatorname{rr\text{-}ms}, e, 0) \end{array}}}

\vnbbookfile{theorem-library/cc-complete-proof}{theorem-library/cc-complete-proof}
\noindent{\small cc-complete-proof.scm -- CC IS A COMPLETE METRIC SPACE.}\par
\vnbprop{cc-re-sub}{\vnbstmt{\begin{array}{@{}l@{}} \forall z, w. \\ \quad ((z \in \mathbb{C}) \wedge (w \in \mathbb{C})) \Rightarrow (\operatorname{real\text{-}part}((z - w)) = (\operatorname{real\text{-}part}(z) - \operatorname{real\text{-}part}(w))) \end{array}}}
\vnbprop{cc-im-sub}{\vnbstmt{\begin{array}{@{}l@{}} \forall z, w. \\ \quad ((z \in \mathbb{C}) \wedge (w \in \mathbb{C})) \Rightarrow (\operatorname{imag\text{-}part}((z - w)) = (\operatorname{imag\text{-}part}(z) - \operatorname{imag\text{-}part}(w))) \end{array}}}
\vnbprop{cc-abs-re-le-magnitude}{\vnbstmt{\forall z \in \mathbb{C}.\; (\lvert \operatorname{real\text{-}part}(z) \rvert \leq \operatorname{magnitude}(z))}}
\vnbprop{cc-abs-im-le-magnitude}{\vnbstmt{\forall z \in \mathbb{C}.\; (\lvert \operatorname{imag\text{-}part}(z) \rvert \leq \operatorname{magnitude}(z))}}
\vnbprop{cc-magnitude-le-re-im}{\vnbstmt{\forall z \in \mathbb{C}.\; (\operatorname{magnitude}(z) \leq (\lvert \operatorname{real\text{-}part}(z) \rvert + \lvert \operatorname{imag\text{-}part}(z) \rvert))}}
\vnbprop{cc-cauchy-of-dominated}{\vnbstmt{\begin{array}{@{}l@{}} \forall f \in (\mathbb{N} \to \mathbb{C}), p \in (\mathbb{N} \to \mathbb{R}). \\ \quad (\operatorname{is\text{-}cauchy\text{-}seq}(\operatorname{cc\text{-}ms}, f) \wedge \forall m, n \in \mathbb{N}.\; (\lvert (p(m) - p(n)) \rvert \leq \operatorname{magnitude}((f(m) - f(n))))) \Rightarrow  \\ \quad \quad \operatorname{is\text{-}cauchy\text{-}seq}(\operatorname{rr\text{-}ms}, p) \end{array}}}
\vnbprop{cc-converges-of-coords}{\vnbstmt{\begin{array}{@{}l@{}} \forall f \in (\mathbb{N} \to \mathbb{C}), g, q \in (\mathbb{N} \to \mathbb{R}), v \in \mathbb{C}. \\ \quad \forall n \in \mathbb{N}.\; (\operatorname{real\text{-}part}(f(n)) = g(n)) \wedge \\ \quad \forall n \in \mathbb{N}.\; (\operatorname{imag\text{-}part}(f(n)) = q(n)) \wedge \\ \quad \operatorname{converges\text{-}to}(\operatorname{rr\text{-}ms}, g, \operatorname{real\text{-}part}(v)) \wedge \\ \quad \operatorname{converges\text{-}to}(\operatorname{rr\text{-}ms}, q, \operatorname{imag\text{-}part}(v)) \Rightarrow \\ \quad \quad \operatorname{converges\text{-}to}(\operatorname{cc\text{-}ms}, f, v) \end{array}}}
\vnbprop{cc-complete}{\vnbstmt{\operatorname{is\text{-}complete}(\operatorname{cc\text{-}ms})}}

\vnbbookfile{theorem-library/series-abs-converges}{theorem-library/series-abs-converges}
\noindent{\small series-abs-converges.scm -- ABSOLUTE CONVERGENCE IMPLIES CONVERGENCE, for real series.}\par
\vnbprop{series-abs-converges}{\vnbstmt{\begin{array}{@{}l@{}} \forall f, g \in (\mathbb{N} \to \mathbb{R}). \\ \quad (\forall j \in \mathbb{N}.\; (g(j) = \lvert f(j) \rvert) \wedge \operatorname{series\text{-}converges}(g)) \Rightarrow  \\ \quad \quad \operatorname{series\text{-}converges}(f) \end{array}}}

\vnbbookfile{theorem-library/cc-magnitude-zero}{theorem-library/cc-magnitude-zero}
\noindent{\small cc-magnitude-zero.scm -- lemmas MOVED 2026-11-10 (integration of the 11-09/10 waves) out of theorem-library/cc-series.scm into library analysis, for the library gate: cited by five analysis files; CC arithmetic. The proofs are the originals, verbatim; the helper defines are a copy of the source's. cc-magnitude-zero}\par
\vnbprop{cc-magnitude-zero}{\vnbstmt{\operatorname{magnitude}(0) = 0}}

\vnbbookfile{theorem-library/mono-le-limit}{theorem-library/mono-le-limit}
\noindent{\small mono-le-limit.scm -- A TERM OF A NONDECREASING CONVERGENT SEQUENCE IS AT MOST ITS LIMIT, and the corollary the product metric's zero law needs.}\par
\vnbprop{rr-mono-le-limit}{\vnbstmt{\begin{array}{@{}l@{}} \forall f \in (\mathbb{N} \to \mathbb{R}), \ell.; \forall m \in \mathbb{N}. \\ \quad (\forall k \in \mathbb{N}.\; (f(k) \leq f((k + 1))) \wedge \operatorname{converges\text{-}to}(\operatorname{rr\text{-}ms}, f, \ell)) \Rightarrow  \\ \quad \quad f(m) \leq \ell \end{array}}}
\vnbprop{series-term-le-sum}{\vnbstmt{\begin{array}{@{}l@{}} \forall f \in (\mathbb{N} \to \mathbb{R}), \ell.; \forall n \in \mathbb{N}. \\ \quad (\forall n \in \mathbb{N}.\; (0 \leq f(n)) \wedge \operatorname{series\text{-}converges\text{-}to}(f, \ell)) \Rightarrow (f(n) \leq \ell) \end{array}}}
\vnbprop{series-nonneg-zero-sum}{\vnbstmt{\begin{array}{@{}l@{}} \forall f \in (\mathbb{N} \to \mathbb{R}), n \in \mathbb{N}. \\ \quad (\forall n \in \mathbb{N}.\; (0 \leq f(n)) \wedge \operatorname{series\text{-}converges\text{-}to}(f, 0)) \Rightarrow (f(n) = 0) \end{array}}}

\vnbbookfile{theorem-library/limit-arithmetic}{theorem-library/limit-arithmetic}
\noindent{\small limit-arithmetic.scm -- THE TWO LIMIT LAWS THE SERIES AND PRODUCT LANES WERE MISSING: `$<$=' carried across TWO limits, and the limit of a pointwise sum.  Plus the series readings of the second.}\par
\vnbprop{rr-limit-le}{\vnbstmt{\begin{array}{@{}l@{}} \forall f, g \in (\mathbb{N} \to \mathbb{R}), v, a \in \mathbb{R}. \\ \quad (\forall n \in \mathbb{N}.\; (f(n) \leq g(n)) \wedge (\operatorname{converges\text{-}to}(\operatorname{rr\text{-}ms}, f, v) \wedge \operatorname{converges\text{-}to}(\operatorname{rr\text{-}ms}, g, a))) \Rightarrow  \\ \quad \quad v \leq a \end{array}}}
\vnbprop{rr-limit-add}{\vnbstmt{\begin{array}{@{}l@{}} \forall f, g, h \in (\mathbb{N} \to \mathbb{R}), v, a \in \mathbb{R}. \\ \quad (\forall j \in \mathbb{N}.\; (h(j) = (f(j) + g(j))) \wedge (\operatorname{converges\text{-}to}(\operatorname{rr\text{-}ms}, f, v) \wedge \operatorname{converges\text{-}to}(\operatorname{rr\text{-}ms}, g, a))) \Rightarrow  \\ \quad \quad \operatorname{converges\text{-}to}(\operatorname{rr\text{-}ms}, h, (v + a)) \end{array}}}
\vnbprop{series-converges-to-converges}{\vnbstmt{\begin{array}{@{}l@{}} \forall f \in (\mathbb{N} \to \mathbb{R}), v \in \mathbb{R}. \\ \quad \operatorname{series\text{-}converges\text{-}to}(f, v) \Rightarrow \operatorname{series\text{-}converges}(f) \end{array}}}
\vnbprop{series-converges-to-add}{\vnbstmt{\begin{array}{@{}l@{}} \forall f, g, h \in (\mathbb{N} \to \mathbb{R}), v, a \in \mathbb{R}. \\ \quad (\forall i \in \mathbb{N}.\; (h(i) = (f(i) + g(i))) \wedge (\operatorname{series\text{-}converges\text{-}to}(f, v) \wedge \operatorname{series\text{-}converges\text{-}to}(g, a))) \Rightarrow  \\ \quad \quad \operatorname{series\text{-}converges\text{-}to}(h, (v + a)) \end{array}}}
\vnbprop{series-limit-add}{\vnbstmt{\begin{array}{@{}l@{}} \forall f, g, h \in (\mathbb{N} \to \mathbb{R}). \\ \quad (\forall i \in \mathbb{N}.\; (h(i) = (f(i) + g(i))) \wedge (\operatorname{series\text{-}converges}(f) \wedge \operatorname{series\text{-}converges}(g))) \Rightarrow  \\ \quad \quad \operatorname{series\text{-}limit}(h) = (\operatorname{series\text{-}limit}(f) + \operatorname{series\text{-}limit}(g)) \end{array}}}

\vnbbookfile{theorem-library/ell-two}{theorem-library/ell-two}
\noindent{\small ell-two.scm -- ELL-TWO, the square-summable complex sequences, and the closure of that carrier under the pointwise vector operations.}\par
\par\medskip\noindent\textbf{Definition} (\texttt{ell-two}).\;
\begin{equation*}
\begin{array}{@{}l@{}} \operatorname{ell\text{-}two} \\ =\; \{\operatorname{x\_} \in \mathbb{C}^{\mathbb{N}} : \operatorname{series\text{-}converges}((\lambda \operatorname{k\_} \in \mathbb{N}.\; \cdot(\operatorname{magnitude}(\operatorname{x\_}(\operatorname{k\_})), \operatorname{magnitude}(\operatorname{x\_}(\operatorname{k\_})))))\} \end{array}
\end{equation*}

\vnbprop{rr-sq-add-le}{\vnbstmt{\begin{array}{@{}l@{}} \forall u, v \in \mathbb{R}. \\ \quad ((u + v) \cdot (u + v)) \leq (((u \cdot u) + (u \cdot u)) + ((v \cdot v) + (v \cdot v))) \end{array}}}
\vnbprop{cc-magnitude-sq-add-le}{\vnbstmt{\begin{array}{@{}l@{}} \forall p, q \in \mathbb{C}. \\ \quad \operatorname{magnitude}((p + q)) \cdot \operatorname{magnitude}((p + q)) \\ \quad \leq\; ((\operatorname{magnitude}(p) \cdot \operatorname{magnitude}(p)) + (\operatorname{magnitude}(p) \cdot \operatorname{magnitude}(p))) + ((\operatorname{magnitude}(q) \cdot \operatorname{magnitude}(q)) + (\operatorname{magnitude}(q) \cdot \operatorname{magnitude}(q))) \end{array}}}
\vnbprop{ell-two-membership}{\vnbstmt{\begin{array}{@{}l@{}} \forall y. \\ \quad y \in \operatorname{ell\text{-}two} \\ \quad \Leftrightarrow\; ((y \in \mathbb{C}^{\mathbb{N}}) \wedge \operatorname{series\text{-}converges}((\lambda k \in \mathbb{N}.\; (\operatorname{magnitude}(y(k)) \cdot \operatorname{magnitude}(y(k)))))) \end{array}}}
\vnbprop{ell-two-add-closed}{\vnbstmt{\begin{array}{@{}l@{}} \forall x, y \in \operatorname{ell\text{-}two}, z \in \mathbb{C}^{\mathbb{N}}. \\ \quad \forall k \in \mathbb{N}.\; (z(k) = (x(k) + y(k))) \Rightarrow (z \in \operatorname{ell\text{-}two}) \end{array}}}
\vnbprop{zero-seq-in-fun}{\vnbstmt{(\lambda n \in \mathbb{N}.\; 0) \in (\mathbb{N} \to \mathbb{R})}}
\vnbprop{series-partial-sum-zero-seq}{\vnbstmt{\forall k \in \mathbb{N}.\; (\operatorname{series\text{-}partial\text{-}sum}((\lambda n \in \mathbb{N}.\; 0), k) = 0)}}
\vnbprop{zero-series-converges}{\vnbstmt{\operatorname{series\text{-}converges}((\lambda n \in \mathbb{N}.\; 0))}}
\vnbprop{ell-two-zero-in}{\vnbstmt{(\lambda n \in \mathbb{N}.\; 0) \in \operatorname{ell\text{-}two}}}
\vnbprop{ell-two-neg-closed}{\vnbstmt{\begin{array}{@{}l@{}} \forall x \in \operatorname{ell\text{-}two}, z \in \mathbb{C}^{\mathbb{N}}. \\ \quad \forall k \in \mathbb{N}.\; (z(k) = -x(k)) \Rightarrow (z \in \operatorname{ell\text{-}two}) \end{array}}}

\vnbbookfile{theorem-library/dyadic-weights}{theorem-library/dyadic-weights}
\noindent{\small dyadic-weights.scm -- THE CANONICAL PRODUCT-METRIC WEIGHTS ARE SUMMABLE, proven -- retiring `product-metric-default-summable' (structure-library/product-metric.scm), the last asserted support of that file that is not blocked on a bridge somebody else is proving:}\par
\vnbprop{power-closed-at}{\vnbstmt{\forall n \in \mathbb{N}, x \in \mathbb{R}.\; (\operatorname{power}(x, n) \in \mathbb{R})}}
\vnbprop{power-real-closed}{\vnbstmt{\forall x \in \mathbb{R}, n \in \mathbb{N}.\; (\operatorname{power}(x, n) \in \mathbb{R})}}
\vnbprop{rr-zero-lt-two}{\vnbstmt{0 < 2}}
\vnbprop{power-two-pos}{\vnbstmt{\forall k \in \mathbb{N}.\; (0 < \operatorname{power}(2, k))}}
\vnbprop{rr-scaled-inverse-unique}{\vnbstmt{\begin{array}{@{}l@{}} \forall a, u, v \in \mathbb{R}. \\ \quad (((a \cdot u) = 1) \wedge (((2 \cdot a) \cdot v) = 1)) \Rightarrow (u = (2 \cdot v)) \end{array}}}
\vnbprop{recip-power-two-halves}{\vnbstmt{\forall k \in \mathbb{N}.\; (\operatorname{recip}(\operatorname{power}(2, k)) = (2 \cdot \operatorname{recip}((2 \cdot \operatorname{power}(2, k)))))}}
\vnbprop{dyadic-weight-apply}{\vnbstmt{\begin{array}{@{}l@{}} \forall k \in \mathbb{N}. \\ \quad (\lambda n \in \mathbb{N}.\; /(1, \operatorname{power}(2, (n + 1))))(k) = \operatorname{recip}(\operatorname{power}(2, (k + 1))) \end{array}}}
\vnbprop{rr-halving-identity}{\vnbstmt{\forall b \in \mathbb{R}.\; (((1 - (2 \cdot b)) + b) = (1 - b))}}
\vnbprop{dyadic-weight-in-fun}{\vnbstmt{(\lambda n \in \mathbb{N}.\; /(1, \operatorname{power}(2, (n + 1)))) \in (\mathbb{N} \to \mathbb{R})}}
\vnbprop{dyadic-weight-pos}{\vnbstmt{\forall a \in \mathbb{N}.\; (0 < (\lambda n \in \mathbb{N}.\; /(1, \operatorname{power}(2, (n + 1))))(a))}}
\vnbprop{dyadic-partial-sum}{\vnbstmt{\begin{array}{@{}l@{}} \forall k \in \mathbb{N}. \\ \quad \operatorname{series\text{-}partial\text{-}sum}((\lambda n \in \mathbb{N}.\; /(1, \operatorname{power}(2, (n + 1)))), k) = (1 - \operatorname{recip}(\operatorname{power}(2, k))) \end{array}}}
\vnbprop{product-metric-default-summable}{\vnbstmt{\operatorname{summable\text{-}weight}((\lambda n \in \mathbb{N}.\; /(1, \operatorname{power}(2, (n + 1)))))}}

\vnbbookfile{theorem-library/rake-inequalities}{theorem-library/rake-inequalities}
\noindent{\small rake-inequalities.scm -- the real-inequality leaves of the rake (batch 5, batch Q), PROVEN.  Nine asserted supports, and the bridge that four of them turned out to need.}\par
\vnbprop{rr-young-2}{\vnbstmt{\forall a, b \in \mathbb{R}.\; ((a \cdot b) \leq /(((a \cdot a) + (b \cdot b)), 2))}}
\vnbprop{rr-amgm-2}{\vnbstmt{\forall a, b \in \mathbb{R}.\; ((4 \cdot (a \cdot b)) \leq ((a + b) \cdot (a + b)))}}
\vnbprop{rr-qm-am-2}{\vnbstmt{\forall a, b \in \mathbb{R}.\; (((a + b) \cdot (a + b)) \leq (2 \cdot ((a \cdot a) + (b \cdot b))))}}
\vnbprop{rr-cauchy-schwarz-2}{\vnbstmt{\begin{array}{@{}l@{}} \forall a, b, c, d \in \mathbb{R}. \\ \quad (((a \cdot c) + (b \cdot d)) \cdot ((a \cdot c) + (b \cdot d))) \leq (((a \cdot a) + (b \cdot b)) \cdot ((c \cdot c) + (d \cdot d))) \end{array}}}
\vnbprop{rr-pos-shrink}{\vnbstmt{\forall s.\; (\operatorname{pos\text{-}rr}(s) \Rightarrow \exists d.\; (\operatorname{pos\text{-}rr}(d) \wedge (d < s)))}}
\vnbprop{rr-bernoulli-ind}{\vnbstmt{\forall n \in \mathbb{N}, x \in \mathbb{R}.\; ((-1 \leq x) \Rightarrow ((1 + (n \cdot x)) \leq \operatorname{power}((1 + x), n)))}}
\vnbprop{rr-bernoulli}{\vnbstmt{\forall x \in \mathbb{R}, n \in \mathbb{N}.\; ((-1 \leq x) \Rightarrow ((1 + (n \cdot x)) \leq \operatorname{power}((1 + x), n)))}}
\vnbprop{cra-ag-rr-carr}{\vnbstmt{(\operatorname{carr}(\operatorname{commutative\text{-}ring\text{-}additive\text{-}ag}(\operatorname{rr\text{-}normed\text{-}field})) \simeq \mathbb{R})}}
\vnbprop{cra-ag-rr-iden}{\vnbstmt{(\operatorname{iden}(\operatorname{commutative\text{-}ring\text{-}additive\text{-}ag}(\operatorname{rr\text{-}normed\text{-}field})) \simeq 0)}}
\vnbprop{cra-ag-rr-opr}{\vnbstmt{(\operatorname{opr}(\operatorname{commutative\text{-}ring\text{-}additive\text{-}ag}(\operatorname{rr\text{-}normed\text{-}field})) \simeq (\lambda {[}\operatorname{x\_}, \operatorname{y\_}] \in (\mathbb{R} \times \mathbb{R}).\; (\operatorname{x\_} + \operatorname{y\_})))}}
\vnbprop{cra-ag-is-nf-ag}{\vnbstmt{\forall r.\; (\operatorname{commutative\text{-}ring\text{-}additive\text{-}ag}(r) \simeq \operatorname{normed\text{-}field\text{-}additive\text{-}ag}(r))}}
\vnbprop{cra-ag-rr-is-abelian-group}{\vnbstmt{\operatorname{is\text{-}abelian\text{-}group}(\operatorname{commutative\text{-}ring\text{-}additive\text{-}ag}(\operatorname{rr\text{-}normed\text{-}field}))}}
\vnbprop{cra-ag-rr-opr-apply}{\vnbstmt{\begin{array}{@{}l@{}} \forall u, v \in \mathbb{R}. \\ \quad (\operatorname{opr}(\operatorname{commutative\text{-}ring\text{-}additive\text{-}ag}(\operatorname{rr\text{-}normed\text{-}field}))(u, v) \simeq (u + v)) \end{array}}}
\vnbprop{rr-nonneg-has-zero}{\vnbstmt{\operatorname{iden}(\operatorname{commutative\text{-}ring\text{-}additive\text{-}ag}(\operatorname{rr\text{-}normed\text{-}field})) \in \{t \in \mathbb{R} : (0 \leq t)\}}}
\vnbprop{rr-nonneg-opr-closed}{\vnbstmt{\begin{array}{@{}l@{}} \forall x, y \in \{t \in \mathbb{R} : (0 \leq t)\}. \\ \quad \operatorname{opr}(\operatorname{commutative\text{-}ring\text{-}additive\text{-}ag}(\operatorname{rr\text{-}normed\text{-}field}))(x, y) \in \{t \in \mathbb{R} : (0 \leq t)\} \end{array}}}
\vnbprop{finsum-rr-in-rr}{\vnbstmt{\begin{array}{@{}l@{}} \forall s \in \mathbf{Set}, f \in (s \to \mathbb{R}). \\ \quad (|s| \in \mathbb{N}) \Rightarrow  \\ \quad \quad \operatorname{finsum}(\operatorname{commutative\text{-}ring\text{-}additive\text{-}ag}(\operatorname{rr\text{-}normed\text{-}field}), f, s) \in \mathbb{R} \end{array}}}
\vnbprop{finsum-rr-nonneg}{\vnbstmt{\begin{array}{@{}l@{}} \forall s \in \mathbf{Set}, f \in (s \to \mathbb{R}). \\ \quad ((|s| \in \mathbb{N}) \wedge \forall z \in s.\; (0 \leq f(z))) \Rightarrow  \\ \quad \quad 0 \leq \operatorname{finsum}(\operatorname{commutative\text{-}ring\text{-}additive\text{-}ag}(\operatorname{rr\text{-}normed\text{-}field}), f, s) \end{array}}}
\vnbprop{finsum-sq-nonneg}{\vnbstmt{\begin{array}{@{}l@{}} \forall s \in \mathbf{Set}, a \in (s \to \mathbb{R}). \\ \quad (|s| \in \mathbb{N}) \Rightarrow  \\ \quad \quad 0 \leq \operatorname{finsum}(\operatorname{commutative\text{-}ring\text{-}additive\text{-}ag}(\operatorname{rr\text{-}normed\text{-}field}), (\lambda i \in s.\; (a(i) \cdot a(i))), s) \end{array}}}
\vnbprop{finsum-le-termwise}{\vnbstmt{\begin{array}{@{}l@{}} \forall s \in \mathbf{Set}, a, b \in (s \to \mathbb{R}). \\ \quad ((|s| \in \mathbb{N}) \wedge \forall i \in s.\; (a(i) \leq b(i))) \Rightarrow  \\ \quad \quad \operatorname{finsum}(\operatorname{commutative\text{-}ring\text{-}additive\text{-}ag}(\operatorname{rr\text{-}normed\text{-}field}), a, s) \\ \quad \quad \leq\; \operatorname{finsum}(\operatorname{commutative\text{-}ring\text{-}additive\text{-}ag}(\operatorname{rr\text{-}normed\text{-}field}), b, s) \end{array}}}
\vnbprop{finsum-abs-triangle}{\vnbstmt{\begin{array}{@{}l@{}} \forall s \in \mathbf{Set}, a \in (s \to \mathbb{R}). \\ \quad (|s| \in \mathbb{N}) \Rightarrow  \\ \quad \quad \lvert \operatorname{finsum}(\operatorname{commutative\text{-}ring\text{-}additive\text{-}ag}(\operatorname{rr\text{-}normed\text{-}field}), a, s) \rvert \\ \quad \quad \leq\; \operatorname{finsum}(\operatorname{commutative\text{-}ring\text{-}additive\text{-}ag}(\operatorname{rr\text{-}normed\text{-}field}), (\lambda i \in s.\; \lvert a(i) \rvert), s) \end{array}}}

\vnbbookfile{theorem-library/mb12b-finsum}{theorem-library/mb12b-finsum}
\noindent{\small mb12b-finsum.scm -- four general facts about real finite sums FINSUM over COMMUTATIVE-RING-ADDITIVE-AG(RR-NORMED-FIELD), each by finite-set-induction, for construction-of-measures 7.38.  Agent ISOM (brief M-B1B2, follow-up), 2026-10-26.}\par
\vnbprop{finsum-rr-const-cls}{\vnbstmt{\begin{array}{@{}l@{}} \forall s \in \mathbf{Set}. \\ \quad (|s| \in \mathbb{N}) \Rightarrow  \\ \quad \quad s \\ \quad \quad \in\; \operatorname{comp}(\mathit{mcs}, \\ \quad \quad \quad \forall c \in \mathbb{R}, f \in \mathbf{Set}. \\ \quad \quad \quad \quad \forall z \in \mathit{mcs}.\; (f(z) = c) \Rightarrow  \\ \quad \quad \quad \quad \quad \operatorname{finsum}(\operatorname{commutative\text{-}ring\text{-}additive\text{-}ag}(\operatorname{rr\text{-}normed\text{-}field}), f, \mathit{mcs}) = (c \cdot |\mathit{mcs}|)) \end{array}}}
\vnbprop{finsum-rr-const}{\vnbstmt{\begin{array}{@{}l@{}} \forall s \in \mathbf{Set}, c \in \mathbb{R}, f \in (s \to \mathbb{R}). \\ \quad ((|s| \in \mathbb{N}) \wedge \forall z \in s.\; (f(z) = c)) \Rightarrow  \\ \quad \quad \operatorname{finsum}(\operatorname{commutative\text{-}ring\text{-}additive\text{-}ag}(\operatorname{rr\text{-}normed\text{-}field}), f, s) = (c \cdot |s|) \end{array}}}
\vnbprop{finite-lower-bound-gt-cls}{\vnbstmt{\begin{array}{@{}l@{}} \forall s \in \mathbf{Set}. \\ \quad (|s| \in \mathbb{N}) \Rightarrow  \\ \quad \quad s \\ \quad \quad \in\; \operatorname{comp}(\mathit{mcs}, \\ \quad \quad \quad \forall t \in \mathbb{R}, a \in \mathbf{Set}. \\ \quad \quad \quad \quad \forall w \in \mathit{mcs}.\; ((a(w) \in \mathbb{R}) \wedge (t < a(w))) \Rightarrow  \\ \quad \quad \quad \quad \quad \exists m \in \mathbb{R}.\; ((t < m) \wedge \forall w \in \mathit{mcs}.\; (m \leq a(w)))) \end{array}}}
\vnbprop{finite-lower-bound-gt}{\vnbstmt{\begin{array}{@{}l@{}} \forall s \in \mathbf{Set}, t \in \mathbb{R}, a \in (s \to \mathbb{R}). \\ \quad ((|s| \in \mathbb{N}) \wedge \forall w \in s.\; (t < a(w))) \Rightarrow  \\ \quad \quad \exists m \in \mathbb{R}.\; ((t < m) \wedge \forall w \in s.\; (m \leq a(w))) \end{array}}}
\vnbprop{finsum-rr-lt-termwise-cls}{\vnbstmt{\begin{array}{@{}l@{}} \forall s \in \mathbf{Set}. \\ \quad (|s| \in \mathbb{N}) \Rightarrow  \\ \quad \quad s \\ \quad \quad \in\; \operatorname{comp}(\mathit{mcs}, \\ \quad \quad \quad \forall a, b \in \mathbf{Set}. \\ \quad \quad \quad \quad \forall w \in \mathit{mcs}.\; ((a(w) \in \mathbb{R}) \wedge ((b(w) \in \mathbb{R}) \wedge (a(w) < b(w)))) \Rightarrow  \\ \quad \quad \quad \quad \quad \operatorname{finsum}(\operatorname{commutative\text{-}ring\text{-}additive\text{-}ag}(\operatorname{rr\text{-}normed\text{-}field}), a, \mathit{mcs}) \in \mathbb{R} \wedge \\ \quad \quad \quad \quad \quad \operatorname{finsum}(\operatorname{commutative\text{-}ring\text{-}additive\text{-}ag}(\operatorname{rr\text{-}normed\text{-}field}), b, \mathit{mcs}) \in \mathbb{R} \wedge \\ \quad \quad \quad \quad \quad \operatorname{finsum}(\operatorname{commutative\text{-}ring\text{-}additive\text{-}ag}(\operatorname{rr\text{-}normed\text{-}field}), a, \mathit{mcs}) \\ \quad \quad \quad \quad \quad \leq\; \operatorname{finsum}(\operatorname{commutative\text{-}ring\text{-}additive\text{-}ag}(\operatorname{rr\text{-}normed\text{-}field}), b, \mathit{mcs}) \wedge \\ \quad \quad \quad \quad \quad \exists k.\; (k \in \mathit{mcs}) \Rightarrow  \\ \quad \quad \quad \quad \quad \quad \operatorname{finsum}(\operatorname{commutative\text{-}ring\text{-}additive\text{-}ag}(\operatorname{rr\text{-}normed\text{-}field}), a, \mathit{mcs}) \\ \quad \quad \quad \quad \quad \quad <\; \operatorname{finsum}(\operatorname{commutative\text{-}ring\text{-}additive\text{-}ag}(\operatorname{rr\text{-}normed\text{-}field}), b, \mathit{mcs})) \end{array}}}
\vnbprop{finsum-rr-lt-termwise}{\vnbstmt{\begin{array}{@{}l@{}} \forall s \in \mathbf{Set}, a, b \in (s \to \mathbb{R}). \\ \quad ((|s| \in \mathbb{N}) \wedge (\forall w \in s.\; (a(w) < b(w)) \wedge \exists k.\; (k \in s))) \Rightarrow  \\ \quad \quad \operatorname{finsum}(\operatorname{commutative\text{-}ring\text{-}additive\text{-}ag}(\operatorname{rr\text{-}normed\text{-}field}), a, s) \\ \quad \quad <\; \operatorname{finsum}(\operatorname{commutative\text{-}ring\text{-}additive\text{-}ag}(\operatorname{rr\text{-}normed\text{-}field}), b, s) \end{array}}}
\vnbprop{finsum-rr-continuous-at-cls}{\vnbstmt{\begin{array}{@{}l@{}} \forall s \in \mathbf{Set}. \\ \quad (|s| \in \mathbb{N}) \Rightarrow  \\ \quad \quad s \\ \quad \quad \in\; \operatorname{comp}(\mathit{mcs}, \\ \quad \quad \quad \forall h \in \mathbf{Set}, g \in (\mathbb{R} \to \mathbf{Set}), a \in \mathbb{R}. \\ \quad \quad \quad \quad \forall w \in \mathit{mcs}. \\ \quad \quad \quad \quad \quad (h(w) \in (\mathbb{R} \to \mathbb{R})) \wedge \operatorname{is\text{-}continuous\text{-}at}(\operatorname{rr\text{-}ms}, \operatorname{rr\text{-}ms}, h(w), a) \wedge \\ \quad \quad \quad \quad \forall t \in \mathbb{R}, w \in \mathit{mcs}.\; ((g(t)(w) \in \mathbb{R}) \wedge (g(t)(w) = h(w)(t))) \Rightarrow \\ \quad \quad \quad \quad \quad (\lambda t \in \mathbb{R}.\; \operatorname{finsum}(\operatorname{commutative\text{-}ring\text{-}additive\text{-}ag}(\operatorname{rr\text{-}normed\text{-}field}), g(t), \mathit{mcs})) \in (\mathbb{R} \to \mathbb{R}) \wedge \\ \quad \quad \quad \quad \quad \operatorname{is\text{-}continuous\text{-}at}(\operatorname{rr\text{-}ms}, \operatorname{rr\text{-}ms}, \\ \quad \quad \quad \quad \quad \quad (\lambda t \in \mathbb{R}.\; \operatorname{finsum}(\operatorname{commutative\text{-}ring\text{-}additive\text{-}ag}(\operatorname{rr\text{-}normed\text{-}field}), g(t), \mathit{mcs})), \\ \quad \quad \quad \quad \quad \quad a)) \end{array}}}
\vnbprop{finsum-rr-continuous-at}{\vnbstmt{\begin{array}{@{}l@{}} \forall s \in \mathbf{Set}, h \in (s \to (\mathbb{R} \to \mathbb{R})), g \in (\mathbb{R} \to (s \to \mathbb{R})), a \in \mathbb{R}. \\ \quad |s| \in \mathbb{N} \wedge \\ \quad \forall w \in s.\; \operatorname{is\text{-}continuous\text{-}at}(\operatorname{rr\text{-}ms}, \operatorname{rr\text{-}ms}, h(w), a) \wedge \\ \quad \forall t \in \mathbb{R}, w \in s.\; (g(t)(w) = h(w)(t)) \Rightarrow \\ \quad \quad \operatorname{is\text{-}continuous\text{-}at}(\operatorname{rr\text{-}ms}, \operatorname{rr\text{-}ms}, \\ \quad \quad \quad (\lambda t \in \mathbb{R}.\; \operatorname{finsum}(\operatorname{commutative\text{-}ring\text{-}additive\text{-}ag}(\operatorname{rr\text{-}normed\text{-}field}), g(t), s)), \\ \quad \quad \quad a) \end{array}}}

\vnbbookfile{theorem-library/deriv-power}{theorem-library/deriv-power}
\noindent{\small deriv-power.scm -- the POWER RULE for the real monomial, PROVEN `modulo 0'.}\par
\vnbprop{pw-step-identity}{\vnbstmt{\begin{array}{@{}l@{}} \forall c, u, v. \\ \quad ((c \in \mathbb{R}) \wedge ((u \in \mathbb{R}) \wedge (v \in \mathbb{R}))) \Rightarrow  \\ \quad \quad ((c + 1) \cdot (u \cdot v)) = ((u \cdot v) + (u \cdot (c \cdot v))) \end{array}}}
\vnbprop{pow-lam-in-fun}{\vnbstmt{\forall n \in \mathbb{N}.\; ((\lambda x \in \mathbb{R}.\; \operatorname{power}(x, n)) \in (\mathbb{R} \to \mathbb{R}))}}
\vnbprop{deriv-power}{\vnbstmt{\begin{array}{@{}l@{}} \forall n \in \mathbb{N}, a \in \mathbb{R}. \\ \quad \operatorname{is\text{-}diff\text{-}at}((\lambda x \in \mathbb{R}.\; \operatorname{power}(x, (n + 1))), a, ((n + 1) \cdot \operatorname{power}(a, n))) \end{array}}}

\vnbbookfile{theorem-library/deriv-polynomial}{theorem-library/deriv-polynomial}
\noindent{\small deriv-polynomial.scm -- rung 1 of the integration arc: the DERIVATIVE OF A POLYNOMIAL FUNCTION, in the coefficient-lambda representation.}\par
\vnbprop{poly-term-lam-in-fun}{\vnbstmt{\begin{array}{@{}l@{}} \forall f \in (\mathbb{N} \to \mathbb{R}), x \in \mathbb{R}. \\ \quad (\lambda k \in \mathbb{N}.\; (f(k) \cdot \operatorname{power}(x, k))) \in (\mathbb{N} \to \mathbb{R}) \end{array}}}
\vnbprop{poly-dseq-lam-in-fun}{\vnbstmt{\begin{array}{@{}l@{}} \forall f \in (\mathbb{N} \to \mathbb{R}), t \in \mathbb{R}. \\ \quad (\lambda j \in \mathbb{N}.\; (((j + 1) \cdot f((j + 1))) \cdot \operatorname{power}(t, j))) \in (\mathbb{N} \to \mathbb{R}) \end{array}}}
\vnbprop{poly-lam-in-fun}{\vnbstmt{\begin{array}{@{}l@{}} \forall f \in (\mathbb{N} \to \mathbb{R}), n \in \mathbb{N}. \\ \quad (\lambda x \in \mathbb{R}.\; \operatorname{series\text{-}partial\text{-}sum}((\lambda k \in \mathbb{N}.\; (f(k) \cdot \operatorname{power}(x, k))), n)) \in (\mathbb{R} \to \mathbb{R}) \end{array}}}
\vnbprop{poly-step-identity}{\vnbstmt{\begin{array}{@{}l@{}} \forall c, u, v. \\ \quad ((c \in \mathbb{R}) \wedge ((u \in \mathbb{R}) \wedge (v \in \mathbb{R}))) \Rightarrow (((c \cdot u) \cdot v) = (u \cdot (c \cdot v))) \end{array}}}
\vnbprop{deriv-scalar-mult}{\vnbstmt{\begin{array}{@{}l@{}} \forall c \in \mathbb{R}, f, t, l. \\ \quad \operatorname{is\text{-}diff\text{-}at}(f, t, l) \Rightarrow \operatorname{is\text{-}diff\text{-}at}((\lambda x \in \mathbb{R}.\; (c \cdot f(x))), t, (c \cdot l)) \end{array}}}
\vnbprop{deriv-coef-monomial}{\vnbstmt{\begin{array}{@{}l@{}} \forall c \in \mathbb{R}, g \in \mathbb{N}, t \in \mathbb{R}. \\ \quad \operatorname{is\text{-}diff\text{-}at}((\lambda x \in \mathbb{R}.\; (c \cdot \operatorname{power}(x, (g + 1)))), t, (((g + 1) \cdot c) \cdot \operatorname{power}(t, g))) \end{array}}}
\vnbprop{deriv-polynomial}{\vnbstmt{\begin{array}{@{}l@{}} \forall n \in \mathbb{N}, f \in (\mathbb{N} \to \mathbb{R}), t \in \mathbb{R}. \\ \quad \operatorname{is\text{-}diff\text{-}at}((\lambda x \in \mathbb{R}.\; \operatorname{series\text{-}partial\text{-}sum}((\lambda k \in \mathbb{N}.\; (f(k) \cdot \operatorname{power}(x, k))), (n + 1))), \\ \quad \quad t, \\ \quad \quad \operatorname{series\text{-}partial\text{-}sum}((\lambda j \in \mathbb{N}.\; (((j + 1) \cdot f((j + 1))) \cdot \operatorname{power}(t, j))), n)) \end{array}}}

\vnbbookfile{theorem-library/poly-antiderivative}{theorem-library/poly-antiderivative}
\noindent{\small poly-antiderivative.scm -- calculus.pdf EXAMPLE 4.7, "every polynomial function is antiderivable", with the antiderivative written out:}\par
\vnbprop{anti-term-in-rr}{\vnbstmt{\begin{array}{@{}l@{}} \forall f \in (\mathbb{N} \to \mathbb{R}), k \in \mathbb{N}, x \in \mathbb{R}. \\ \quad ((\operatorname{recip}((k + 1)) \cdot f(k)) \cdot \operatorname{power}(x, (k + 1))) \in \mathbb{R} \end{array}}}
\vnbprop{anti-term-lam-in-fun}{\vnbstmt{\begin{array}{@{}l@{}} \forall f \in (\mathbb{N} \to \mathbb{R}), x \in \mathbb{R}. \\ \quad (\lambda k \in \mathbb{N}.\; ((\operatorname{recip}((k + 1)) \cdot f(k)) \cdot \operatorname{power}(x, (k + 1)))) \in (\mathbb{N} \to \mathbb{R}) \end{array}}}
\vnbprop{anti-lam-in-fun}{\vnbstmt{\begin{array}{@{}l@{}} \forall f \in (\mathbb{N} \to \mathbb{R}), m \in \mathbb{N}. \\ \quad (\lambda x \in \mathbb{R}.\; \operatorname{series\text{-}partial\text{-}sum}((\lambda k \in \mathbb{N}.\; ((\operatorname{recip}((k + 1)) \cdot f(k)) \cdot \operatorname{power}(x, (k + 1)))), m)) \\ \quad \in\; (\mathbb{R} \to \mathbb{R}) \end{array}}}
\vnbprop{rr-zero-plus}{\vnbstmt{\forall t \in \mathbb{R}.\; (t = t)}}
\vnbprop{recip-succ-cancel}{\vnbstmt{\begin{array}{@{}l@{}} \forall m \in \mathbb{N}, c, v \in \mathbb{R}. \\ \quad (c \cdot v) = (((m + 1) \cdot (\operatorname{recip}((m + 1)) \cdot c)) \cdot v) \end{array}}}
\vnbprop{deriv-anti-monomial}{\vnbstmt{\begin{array}{@{}l@{}} \forall c \in \mathbb{R}, m \in \mathbb{N}, t \in \mathbb{R}. \\ \quad \operatorname{is\text{-}diff\text{-}at}((\lambda x \in \mathbb{R}.\; ((\operatorname{recip}((m + 1)) \cdot c) \cdot \operatorname{power}(x, (m + 1)))), t, (c \cdot \operatorname{power}(t, m))) \end{array}}}
\vnbprop{poly-antiderivative}{\vnbstmt{\begin{array}{@{}l@{}} \forall n \in \mathbb{N}, f \in (\mathbb{N} \to \mathbb{R}), t \in \mathbb{R}. \\ \quad \operatorname{is\text{-}diff\text{-}at}((\lambda x \in \mathbb{R}.\; \\ \quad \quad \quad \operatorname{series\text{-}partial\text{-}sum}((\lambda k \in \mathbb{N}.\; ((\operatorname{recip}((k + 1)) \cdot f(k)) \cdot \operatorname{power}(x, (k + 1)))), (n + 1))), \\ \quad \quad t, \\ \quad \quad \operatorname{series\text{-}partial\text{-}sum}((\lambda k \in \mathbb{N}.\; (f(k) \cdot \operatorname{power}(t, k))), (n + 1))) \end{array}}}

\vnbbookfile{theorem-library/product-summable}{theorem-library/product-summable}
\noindent{\small product-summable.scm -- THE DEFINING SERIES OF THE PRODUCT METRIC CONVERGES, proven -- retiring the asserted `product-weighted-summable' (structure-library/product-metric.scm), which is what makes the IOTA in PRODUCT-METRIC-W a well-defined distance, and with it the second asserted}\par
\vnbprop{product-carrier-unfold}{\vnbstmt{\begin{array}{@{}l@{}} \forall s. \\ \quad (\operatorname{product\text{-}carrier}(s) \simeq \{x \in (\mathbb{N} \to \operatorname{big\text{-}union}(n, \mathbb{N}, \operatorname{pts}(s(n)))) : \forall n \in \mathbb{N}.\; (x(n) \in \operatorname{pts}(s(n)))\}) \end{array}}}
\vnbprop{product-carrier-coord}{\vnbstmt{\forall s.; \forall x \in \operatorname{product\text{-}carrier}(s), n \in \mathbb{N}.\; (x(n) \in \operatorname{pts}(s(n)))}}
\vnbprop{bdd-metric-weight-bound}{\vnbstmt{\begin{array}{@{}l@{}} \forall s.; \forall u, v \in \operatorname{pts}(s), n \in \mathbb{R}. \\ \quad (\operatorname{is\text{-}metric\text{-}space}(s) \wedge (0 \leq n)) \Rightarrow  \\ \quad \quad (n \cdot \operatorname{dist}(\operatorname{bdd\text{-}metric}(s))(u, v)) \in \mathbb{R} \wedge \\ \quad \quad 0 \leq (n \cdot \operatorname{dist}(\operatorname{bdd\text{-}metric}(s))(u, v)) \wedge \\ \quad \quad (n \cdot \operatorname{dist}(\operatorname{bdd\text{-}metric}(s))(u, v)) \leq n \end{array}}}
\vnbprop{product-metric-carrier}{\vnbstmt{\forall s, w.\; (\operatorname{pts}(\operatorname{product\text{-}metric\text{-}w}(s, w)) \simeq \operatorname{product\text{-}carrier}(s))}}
\vnbprop{product-term-seq-in-fun}{\vnbstmt{\begin{array}{@{}l@{}} \forall s, w.; \forall x, y \in \operatorname{product\text{-}carrier}(s). \\ \quad (\operatorname{is\text{-}ms\text{-}sequence}(s) \wedge \operatorname{summable\text{-}weight}(w)) \Rightarrow  \\ \quad \quad (\lambda n \in \mathbb{N}.\; (w(n) \cdot \operatorname{dist}(\operatorname{bdd\text{-}metric}(s(n)))(x(n), y(n)))) \in (\mathbb{N} \to \mathbb{R}) \end{array}}}
\vnbprop{product-weighted-summable}{\vnbstmt{\begin{array}{@{}l@{}} \forall s, w.; \forall x, y \in \operatorname{product\text{-}carrier}(s). \\ \quad (\operatorname{is\text{-}ms\text{-}sequence}(s) \wedge \operatorname{summable\text{-}weight}(w)) \Rightarrow  \\ \quad \quad \operatorname{series\text{-}converges}((\lambda n \in \mathbb{N}.\; (w(n) \cdot \operatorname{dist}(\operatorname{bdd\text{-}metric}(s(n)))(x(n), y(n))))) \end{array}}}
\vnbprop{product-metric-distance}{\vnbstmt{\begin{array}{@{}l@{}} \forall s, w.; \forall x, y \in \operatorname{product\text{-}carrier}(s). \\ \quad (\operatorname{dist}(\operatorname{product\text{-}metric\text{-}w}(s, w))(x, y) \simeq \operatorname{series\text{-}limit}((\lambda n \in \mathbb{N}.\; (w(n) \cdot \operatorname{dist}(\operatorname{bdd\text{-}metric}(s(n)))(x(n), y(n)))))) \end{array}}}
\vnbprop{product-metric-dist-converges}{\vnbstmt{\begin{array}{@{}l@{}} \forall s, w.; \forall x, y \in \operatorname{product\text{-}carrier}(s). \\ \quad (\operatorname{is\text{-}ms\text{-}sequence}(s) \wedge \operatorname{summable\text{-}weight}(w)) \Rightarrow  \\ \quad \quad \operatorname{series\text{-}converges\text{-}to}((\lambda n \in \mathbb{N}.\; (w(n) \cdot \operatorname{dist}(\operatorname{bdd\text{-}metric}(s(n)))(x(n), y(n)))), \\ \quad \quad \quad \operatorname{dist}(\operatorname{product\text{-}metric\text{-}w}(s, w))(x, y)) \end{array}}}

\vnbbookfile{theorem-library/product-is-metric-space}{theorem-library/product-is-metric-space}
\noindent{\small product-is-metric-space.scm -- THE COUNTABLE PRODUCT OF METRIC SPACES IS A METRIC SPACE, proven -- retiring the asserted `product-is-metric-space' of structure-library/product-metric.scm (line 79).}\par
\vnbprop{series-limit-le}{\vnbstmt{\begin{array}{@{}l@{}} \forall f, g \in (\mathbb{N} \to \mathbb{R}), v, a \in \mathbb{R}. \\ \quad (\forall n \in \mathbb{N}.\; (f(n) \leq g(n)) \wedge (\operatorname{series\text{-}converges\text{-}to}(f, v) \wedge \operatorname{series\text{-}converges\text{-}to}(g, a))) \Rightarrow  \\ \quad \quad v \leq a \end{array}}}
\vnbprop{series-limit-iota}{\vnbstmt{\forall f.\; ((\iota\, l.\;((l \in \mathbb{R}) \wedge \operatorname{series\text{-}converges\text{-}to}(f, l))) \simeq \operatorname{series\text{-}limit}(f))}}
\vnbprop{product-dist-in-rr}{\vnbstmt{\begin{array}{@{}l@{}} \forall s, w.; \forall x, y \in \operatorname{product\text{-}carrier}(s). \\ \quad (\operatorname{is\text{-}ms\text{-}sequence}(s) \wedge \operatorname{summable\text{-}weight}(w)) \Rightarrow  \\ \quad \quad \operatorname{dist}(\operatorname{product\text{-}metric\text{-}w}(s, w))(x, y) \in \mathbb{R} \end{array}}}
\vnbprop{product-term-nonneg}{\vnbstmt{\begin{array}{@{}l@{}} \forall s, w.; \forall x, y \in \operatorname{product\text{-}carrier}(s), a \in \mathbb{N}. \\ \quad (\operatorname{is\text{-}ms\text{-}sequence}(s) \wedge \operatorname{summable\text{-}weight}(w)) \Rightarrow  \\ \quad \quad 0 \leq (\lambda n \in \mathbb{N}.\; (w(n) \cdot \operatorname{dist}(\operatorname{bdd\text{-}metric}(s(n)))(x(n), y(n))))(a) \end{array}}}
\vnbprop{product-term-self-zero}{\vnbstmt{\begin{array}{@{}l@{}} \forall s, w.; \forall x \in \operatorname{product\text{-}carrier}(s), a \in \mathbb{N}. \\ \quad (\operatorname{is\text{-}ms\text{-}sequence}(s) \wedge \operatorname{summable\text{-}weight}(w)) \Rightarrow  \\ \quad \quad (\lambda n \in \mathbb{N}.\; (w(n) \cdot \operatorname{dist}(\operatorname{bdd\text{-}metric}(s(n)))(x(n), x(n))))(a) = 0 \end{array}}}
\vnbprop{product-term-sym}{\vnbstmt{\begin{array}{@{}l@{}} \forall s, w.; \forall x, y \in \operatorname{product\text{-}carrier}(s), a \in \mathbb{N}. \\ \quad (\operatorname{is\text{-}ms\text{-}sequence}(s) \wedge \operatorname{summable\text{-}weight}(w)) \Rightarrow  \\ \quad \quad (\lambda n \in \mathbb{N}.\; (w(n) \cdot \operatorname{dist}(\operatorname{bdd\text{-}metric}(s(n)))(x(n), y(n))))(a) \\ \quad \quad \leq\; (\lambda n \in \mathbb{N}.\; (w(n) \cdot \operatorname{dist}(\operatorname{bdd\text{-}metric}(s(n)))(y(n), x(n))))(a) \end{array}}}
\vnbprop{product-sum-seq-in-fun}{\vnbstmt{\begin{array}{@{}l@{}} \forall s, w.; \forall x, y, z \in \operatorname{product\text{-}carrier}(s). \\ \quad (\operatorname{is\text{-}ms\text{-}sequence}(s) \wedge \operatorname{summable\text{-}weight}(w)) \Rightarrow  \\ \quad \quad (\lambda n \in \mathbb{N}.\; \\ \quad \quad \quad (w(n) \cdot \operatorname{dist}(\operatorname{bdd\text{-}metric}(s(n)))(x(n), y(n))) + (w(n) \cdot \operatorname{dist}(\operatorname{bdd\text{-}metric}(s(n)))(y(n), z(n)))) \\ \quad \quad \in\; (\mathbb{N} \to \mathbb{R}) \end{array}}}
\vnbprop{product-sum-seq-split}{\vnbstmt{\begin{array}{@{}l@{}} \forall s, w.; \forall x, y, z \in \operatorname{product\text{-}carrier}(s), i \in \mathbb{N}. \\ \quad (\operatorname{is\text{-}ms\text{-}sequence}(s) \wedge \operatorname{summable\text{-}weight}(w)) \Rightarrow  \\ \quad \quad (\lambda n \in \mathbb{N}.\; ((w(n) \cdot \operatorname{dist}(\operatorname{bdd\text{-}metric}(s(n)))(x(n), y(n))) + (w(n) \cdot \operatorname{dist}(\operatorname{bdd\text{-}metric}(s(n)))(y(n), z(n)))))(i) \\ \quad \quad =\; (\lambda n \in \mathbb{N}.\; (w(n) \cdot \operatorname{dist}(\operatorname{bdd\text{-}metric}(s(n)))(x(n), y(n))))(i) + (\lambda n \in \mathbb{N}.\; (w(n) \cdot \operatorname{dist}(\operatorname{bdd\text{-}metric}(s(n)))(y(n), z(n))))(i) \end{array}}}
\vnbprop{product-term-triangle}{\vnbstmt{\begin{array}{@{}l@{}} \forall s, w.; \forall x, y, z \in \operatorname{product\text{-}carrier}(s), a \in \mathbb{N}. \\ \quad (\operatorname{is\text{-}ms\text{-}sequence}(s) \wedge \operatorname{summable\text{-}weight}(w)) \Rightarrow  \\ \quad \quad (\lambda n \in \mathbb{N}.\; (w(n) \cdot \operatorname{dist}(\operatorname{bdd\text{-}metric}(s(n)))(x(n), z(n))))(a) \\ \quad \quad \leq\; (\lambda n \in \mathbb{N}.\; ((w(n) \cdot \operatorname{dist}(\operatorname{bdd\text{-}metric}(s(n)))(x(n), y(n))) + (w(n) \cdot \operatorname{dist}(\operatorname{bdd\text{-}metric}(s(n)))(y(n), z(n)))))(a) \end{array}}}
\vnbprop{product-is-metric-space}{\vnbstmt{\begin{array}{@{}l@{}} \forall s, w. \\ \quad (\operatorname{is\text{-}ms\text{-}sequence}(s) \wedge \operatorname{summable\text{-}weight}(w)) \Rightarrow  \\ \quad \quad \operatorname{is\text{-}metric\text{-}space}(\operatorname{product\text{-}metric\text{-}w}(s, w)) \end{array}}}

\vnbbookfile{theorem-library/rake-block-tower}{theorem-library/rake-block-tower}
\noindent{\small theorem-library/rake-block-tower.scm ==================================================================== convergence-block-tower (theorem-library/seq-compact-product.scm:91) -- the refinement-tower core of the cross-coordinate diagonal argument, and}\par
\vnbprop{block-step-converges}{\vnbstmt{\begin{array}{@{}l@{}} \forall s.; \forall h \in (\mathbb{N} \to \operatorname{pts}(s)), j \in \mathrm{Inf}(\mathbb{N}). \\ \quad \operatorname{seq\text{-}compact}(s) \Rightarrow  \\ \quad \quad \exists b \in \mathrm{Inf}(\mathbb{N}). \\ \quad \quad \quad (b \subseteq j) \wedge \exists p \in \operatorname{pts}(s).\; \operatorname{converges\text{-}along}(s, h, b, p) \end{array}}}
\vnbprop{convergence-block-tower}{\vnbstmt{\begin{array}{@{}l@{}} \forall a.; \forall q \in (\mathbb{N} \to \operatorname{product\text{-}carrier}(a)). \\ \quad (\operatorname{is\text{-}ms\text{-}sequence}(a) \wedge \forall n \in \mathbb{N}.\; \operatorname{seq\text{-}compact}(a(n))) \Rightarrow  \\ \quad \quad \exists s \in (\mathbb{N} \to \mathrm{Inf}(\mathbb{N})). \\ \quad \quad \quad \forall k \in \mathbb{N}.\; (s((k + 1)) \subseteq s(k)) \wedge \\ \quad \quad \quad \exists l \in \operatorname{product\text{-}carrier}(a). \forall n \in \mathbb{N}. \\ \quad \quad \quad \quad \operatorname{converges\text{-}along}(a(n), (\lambda i \in \mathbb{N}.\; q(i)(n)), s(n), l(n)) \end{array}}}

\vnbbookfile{theorem-library/rr-null-scale}{theorem-library/rr-null-scale}
\noindent{\small rr-null-scale.scm -- a NONNEGATIVE CONSTANT MULTIPLE OF A NULL SEQUENCE IS NULL, proven, together with the eps/delta lemma that carries it.}\par
\vnbprop{rr-scale-eps}{\vnbstmt{\begin{array}{@{}l@{}} \forall c, s \in \mathbb{R}. \\ \quad ((0 \leq c) \wedge \operatorname{pos\text{-}rr}(s)) \Rightarrow  \\ \quad \quad \exists d. \\ \quad \quad \quad \operatorname{pos\text{-}rr}(d) \wedge \forall t \in \mathbb{R}.\; (((0 \leq t) \wedge (t \leq d)) \Rightarrow ((c \cdot t) \leq s)) \end{array}}}
\vnbprop{rr-null-scale}{\vnbstmt{\begin{array}{@{}l@{}} \forall c \in \mathbb{R}, f, h \in (\mathbb{N} \to \mathbb{R}). \\ \quad ((0 \leq c) \wedge (\operatorname{converges\text{-}to}(\operatorname{rr\text{-}ms}, f, 0) \wedge \forall j \in \mathbb{N}.\; (h(j) = (c \cdot f(j))))) \Rightarrow  \\ \quad \quad \operatorname{converges\text{-}to}(\operatorname{rr\text{-}ms}, h, 0) \end{array}}}

\vnbbookfile{theorem-library/bdd-metric-convergence}{theorem-library/bdd-metric-convergence}
\noindent{\small bdd-metric-convergence.scm -- CONVERGENCE IS THE SAME IN d AND IN d/(1+d), proven: a sequence converges to L in BDD-METRIC(s) if and only if it converges to L in s.}\par
\vnbprop{bdd-fn-reflect}{\vnbstmt{\begin{array}{@{}l@{}} \forall a, b \in \mathbb{R}. \\ \quad ((0 \leq a) \wedge ((0 \leq b) \wedge (/(a, (1 + a)) \leq /(b, (1 + b))))) \Rightarrow (a \leq b) \end{array}}}
\vnbprop{bdd-metric-dist-le}{\vnbstmt{\begin{array}{@{}l@{}} \forall s.; \forall x, y \in \operatorname{pts}(s). \\ \quad \operatorname{is\text{-}metric\text{-}space}(s) \Rightarrow (\operatorname{dist}(\operatorname{bdd\text{-}metric}(s))(x, y) \leq \operatorname{dist}(s)(x, y)) \end{array}}}
\vnbprop{bdd-metric-dist-reflect}{\vnbstmt{\begin{array}{@{}l@{}} \forall s.; \forall x, y \in \operatorname{pts}(s), e \in \mathbb{R}. \\ \quad (\operatorname{is\text{-}metric\text{-}space}(s) \wedge ((0 \leq e) \wedge (\operatorname{dist}(\operatorname{bdd\text{-}metric}(s))(x, y) \leq /(e, (1 + e))))) \Rightarrow  \\ \quad \quad \operatorname{dist}(s)(x, y) \leq e \end{array}}}
\vnbprop{bdd-metric-converges-fwd}{\vnbstmt{\begin{array}{@{}l@{}} \forall s.; \forall f \in (\mathbb{N} \to \operatorname{pts}(s)), \ell \in \operatorname{pts}(s). \\ \quad (\operatorname{is\text{-}metric\text{-}space}(s) \wedge \operatorname{converges\text{-}to}(s, f, \ell)) \Rightarrow  \\ \quad \quad \operatorname{converges\text{-}to}(\operatorname{bdd\text{-}metric}(s), f, \ell) \end{array}}}
\vnbprop{bdd-metric-converges-bwd}{\vnbstmt{\begin{array}{@{}l@{}} \forall s.; \forall f \in (\mathbb{N} \to \operatorname{pts}(s)), \ell \in \operatorname{pts}(s). \\ \quad (\operatorname{is\text{-}metric\text{-}space}(s) \wedge \operatorname{converges\text{-}to}(\operatorname{bdd\text{-}metric}(s), f, \ell)) \Rightarrow  \\ \quad \quad \operatorname{converges\text{-}to}(s, f, \ell) \end{array}}}
\vnbprop{bdd-metric-converges-iff}{\vnbstmt{\begin{array}{@{}l@{}} \forall s.; \forall f \in (\mathbb{N} \to \operatorname{pts}(s)), \ell \in \operatorname{pts}(s). \\ \quad \operatorname{is\text{-}metric\text{-}space}(s) \Rightarrow  \\ \quad \quad \operatorname{converges\text{-}to}(\operatorname{bdd\text{-}metric}(s), f, \ell) \Leftrightarrow \operatorname{converges\text{-}to}(s, f, \ell) \end{array}}}

\vnbbookfile{theorem-library/rake-metric-constructions}{theorem-library/rake-metric-constructions}
\noindent{\small rake-metric-constructions.scm -- BATCH S of the 2026-09-18 rake: the metric CONSTRUCTIONS (bounded metric, ball cover, countable product).  Every statement below is its support's statement UNCHANGED, copied from the support site, except where the header says otherwise.}\par
\vnbprop{bdd-metric-id-bicontinuous}{\vnbstmt{\begin{array}{@{}l@{}} \forall s. \\ \quad \operatorname{is\text{-}metric\text{-}space}(s) \Rightarrow  \\ \quad \quad \operatorname{is\text{-}continuous}(s, \operatorname{bdd\text{-}metric}(s), (\lambda x \in \operatorname{pts}(s).\; x)) \wedge \\ \quad \quad \operatorname{is\text{-}continuous}(\operatorname{bdd\text{-}metric}(s), s, (\lambda x \in \operatorname{pts}(\operatorname{bdd\text{-}metric}(s)).\; x)) \end{array}}}
\vnbprop{bdd-rr-ms-is-metric-space}{\vnbstmt{\operatorname{is\text{-}metric\text{-}space}(\operatorname{bdd\text{-}metric}(\operatorname{rr\text{-}ms}))}}
\vnbprop{bdd-rr-ms-bounded}{\vnbstmt{\forall x, y \in \mathbb{R}.\; (\operatorname{dist}(\operatorname{bdd\text{-}metric}(\operatorname{rr\text{-}ms}))(x, y) < 1)}}
\vnbprop{bdd-rr-equivalent}{\vnbstmt{\begin{array}{@{}l@{}} \operatorname{is\text{-}continuous}(\operatorname{rr\text{-}ms}, \operatorname{bdd\text{-}metric}(\operatorname{rr\text{-}ms}), (\lambda x \in \mathbb{R}.\; x)) \wedge \\ \operatorname{is\text{-}continuous}(\operatorname{bdd\text{-}metric}(\operatorname{rr\text{-}ms}), \operatorname{rr\text{-}ms}, (\lambda x \in \mathbb{R}.\; x)) \end{array}}}
\vnbprop{ball-cover-is-open-cover}{\vnbstmt{\begin{array}{@{}l@{}} \forall s.; \forall r \in \mathbb{R}. \\ \quad (\operatorname{is\text{-}metric\text{-}space}(s) \wedge ((0 \leq r) \wedge \neg (0 = r))) \Rightarrow  \\ \quad \quad \operatorname{is\text{-}open\text{-}cover}(s, \operatorname{ball\text{-}cover}(s, r)) \end{array}}}
\vnbprop{product-projection-continuous}{\vnbstmt{\begin{array}{@{}l@{}} \forall s, w.; \forall n \in \mathbb{N}. \\ \quad (\operatorname{is\text{-}ms\text{-}sequence}(s) \wedge \operatorname{summable\text{-}weight}(w)) \Rightarrow  \\ \quad \quad \operatorname{is\text{-}continuous}(\operatorname{product\text{-}metric\text{-}w}(s, w), s(n), \operatorname{product\text{-}proj}(s, n)) \end{array}}}

\vnbbookfile{theorem-library/rake-rr-bounded-ms}{theorem-library/rake-rr-bounded-ms}
\noindent{\small rake-rr-bounded-ms.scm -- the three facts about the constant RR-BOUNDED-MS (rake batch 5, integrator, 2026-09-18).}\par
\vnbprop{rr-bounded-ms-is-metric-space}{\vnbstmt{\operatorname{is\text{-}metric\text{-}space}(\operatorname{rr\text{-}bounded\text{-}ms})}}
\vnbprop{rr-bounded-ms-bounded}{\vnbstmt{\forall x, y \in \mathbb{R}.\; (\operatorname{dist}(\operatorname{rr\text{-}bounded\text{-}ms})(x, y) < 1)}}
\vnbprop{rr-bounded-equivalent}{\vnbstmt{\begin{array}{@{}l@{}} \operatorname{is\text{-}continuous}(\operatorname{rr\text{-}ms}, \operatorname{rr\text{-}bounded\text{-}ms}, (\lambda x \in \mathbb{R}.\; x)) \wedge \\ \operatorname{is\text{-}continuous}(\operatorname{rr\text{-}bounded\text{-}ms}, \operatorname{rr\text{-}ms}, (\lambda x \in \mathbb{R}.\; x)) \end{array}}}

\vnbbookfile{theorem-library/rake-compact-equivalences}{theorem-library/rake-compact-equivalences}
\noindent{\small rake-compact-equivalences.scm -- Prop 3.12, the two remaining equivalences (batch 8, assignment 8-A; helper prefix `r9a-').}\par
\vnbprop{compact-iff-cluster-point}{\vnbstmt{\begin{array}{@{}l@{}} \forall s. \\ \quad \operatorname{is\text{-}metric\text{-}space}(s) \Rightarrow  \\ \quad \quad \operatorname{is\text{-}compact}(s) \Leftrightarrow \forall f \in (\mathbb{N} \to \operatorname{pts}(s)).\; \exists x.\; \operatorname{cluster\text{-}point}(s, f, x) \end{array}}}
\vnbprop{compact-iff-tb-complete}{\vnbstmt{\begin{array}{@{}l@{}} \forall s. \\ \quad \operatorname{is\text{-}metric\text{-}space}(s) \Rightarrow  \\ \quad \quad \operatorname{is\text{-}compact}(s) \Leftrightarrow (\operatorname{totally\text{-}bounded}(s) \wedge \operatorname{is\text{-}complete}(s)) \end{array}}}

\vnbbookfile{theorem-library/rake-compact-fip}{theorem-library/rake-compact-fip}
\noindent{\small rake-compact-fip.scm -- Prop 3.12 (1) $<$=$>$ (2), the finite-intersection characterisation of compactness, PROVEN (2026-09-20, batch 11-E; helper prefix `rcf-').}\par
\vnbprop{compact-implies-fip}{\vnbstmt{\begin{array}{@{}l@{}} \forall s, c. \\ \quad (\operatorname{is\text{-}metric\text{-}space}(s) \wedge (\operatorname{is\text{-}compact}(s) \wedge \operatorname{has\text{-}fip}(s, c))) \Rightarrow  \\ \quad \quad \exists p.\; (p \in \operatorname{intersection\text{-}of}(c)) \end{array}}}
\vnbprop{fip-implies-compact}{\vnbstmt{\begin{array}{@{}l@{}} \forall s. \\ \quad (\operatorname{is\text{-}metric\text{-}space}(s) \wedge \forall c.\; (\operatorname{has\text{-}fip}(s, c) \Rightarrow \exists p.\; (p \in \operatorname{intersection\text{-}of}(c)))) \Rightarrow  \\ \quad \quad \operatorname{is\text{-}compact}(s) \end{array}}}
\vnbprop{compact-iff-fip}{\vnbstmt{\begin{array}{@{}l@{}} \forall s. \\ \quad \operatorname{is\text{-}metric\text{-}space}(s) \Rightarrow  \\ \quad \quad \operatorname{is\text{-}compact}(s) \Leftrightarrow \forall c.\; (\operatorname{has\text{-}fip}(s, c) \Rightarrow \exists p.\; (p \in \operatorname{intersection\text{-}of}(c))) \end{array}}}

\vnbbookfile{theorem-library/compact-separable-proof}{theorem-library/compact-separable-proof}
\noindent{\small theorem-library/compact-separable-proof.scm}\par
\vnbprop{compact-metric-is-separable}{\vnbstmt{\forall s.\; ((\operatorname{is\text{-}compact}(s) \wedge \exists x.\; (x \in \operatorname{pts}(s))) \Rightarrow \operatorname{is\text{-}separable}(s))}}

\vnbbookfile{theorem-library/rake-ascoli}{theorem-library/rake-ascoli}
\noindent{\small theorem-library/rake-ascoli.scm ==================================================================== Ascoli-Arzela (sequential, RR-valued) -- batch 8-B.}\par
\vnbprop{equicont-subseq}{\vnbstmt{\begin{array}{@{}l@{}} \forall s, t, m, i. \\ \quad (\operatorname{is\text{-}equicontinuous}(s, t, m) \wedge \operatorname{strictly\text{-}mono\text{-}nn}(i)) \Rightarrow  \\ \quad \quad \operatorname{is\text{-}equicontinuous}(s, t, \operatorname{subseq}(m, i)) \end{array}}}
\vnbprop{ascoli-sequential-from-diagonal}{\vnbstmt{\begin{array}{@{}l@{}} \forall s, m. \\ \quad \operatorname{is\text{-}compact}(s) \wedge \\ \quad \exists x.\; (x \in \operatorname{pts}(s)) \wedge \\ \quad m \in (\mathbb{N} \to (\operatorname{pts}(s) \to \mathbb{R})) \wedge \\ \quad \forall k \in \mathbb{N}.\; \operatorname{is\text{-}continuous}(s, \operatorname{rr\text{-}ms}, m(k)) \wedge \\ \quad \operatorname{is\text{-}equicontinuous}(s, \operatorname{rr\text{-}ms}, m) \wedge \\ \quad \operatorname{pointwise\text{-}bounded}(s, m) \wedge \\ \quad \forall q. \\ \quad \quad \operatorname{is\text{-}dense\text{-}seq}(s, q) \Rightarrow  \\ \quad \quad \quad \exists l.\; (\operatorname{strictly\text{-}mono\text{-}nn}(l) \wedge \operatorname{converges\text{-}on}(s, \operatorname{subseq}(m, l), q)) \Rightarrow \\ \quad \quad \exists l. \\ \quad \quad \quad \operatorname{strictly\text{-}mono\text{-}nn}(l) \wedge \\ \quad \quad \quad \exists g \in (\operatorname{pts}(s) \to \mathbb{R}). \\ \quad \quad \quad \quad \operatorname{is\text{-}continuous}(s, \operatorname{rr\text{-}ms}, g) \wedge \operatorname{converges\text{-}uniformly}(s, \operatorname{subseq}(m, l), g) \end{array}}}

\vnbbookfile{theorem-library/converges-dist-null}{theorem-library/converges-dist-null}
\noindent{\small converges-dist-null.scm -- CONVERGENCE IN A METRIC SPACE IS A NULL SEQUENCE OF REAL DISTANCES, proven both ways:}\par
\vnbprop{dist-seq-apply}{\vnbstmt{\begin{array}{@{}l@{}} \forall t, f, v.; \forall n \in \mathbb{N}. \\ \quad ((\lambda j \in \mathbb{N}.\; \operatorname{dist}(t)(f(j), v))(n) \simeq \operatorname{dist}(t)(f(n), v)) \end{array}}}
\vnbprop{dist-seq-in-fun}{\vnbstmt{\begin{array}{@{}l@{}} \forall t.; \forall f \in (\mathbb{N} \to \operatorname{pts}(t)), v \in \operatorname{pts}(t). \\ \quad \operatorname{is\text{-}metric\text{-}space}(t) \Rightarrow ((\lambda j \in \mathbb{N}.\; \operatorname{dist}(t)(f(j), v)) \in (\mathbb{N} \to \mathbb{R})) \end{array}}}
\vnbprop{converges-dist-null-fwd}{\vnbstmt{\begin{array}{@{}l@{}} \forall t.; \forall f \in (\mathbb{N} \to \operatorname{pts}(t)), v \in \operatorname{pts}(t). \\ \quad (\operatorname{is\text{-}metric\text{-}space}(t) \wedge \operatorname{converges\text{-}to}(t, f, v)) \Rightarrow  \\ \quad \quad \operatorname{converges\text{-}to}(\operatorname{rr\text{-}ms}, (\lambda j \in \mathbb{N}.\; \operatorname{dist}(t)(f(j), v)), 0) \end{array}}}
\vnbprop{converges-dist-null-bwd}{\vnbstmt{\begin{array}{@{}l@{}} \forall t.; \forall f \in (\mathbb{N} \to \operatorname{pts}(t)), v \in \operatorname{pts}(t). \\ \quad (\operatorname{is\text{-}metric\text{-}space}(t) \wedge \operatorname{converges\text{-}to}(\operatorname{rr\text{-}ms}, (\lambda j \in \mathbb{N}.\; \operatorname{dist}(t)(f(j), v)), 0)) \Rightarrow  \\ \quad \quad \operatorname{converges\text{-}to}(t, f, v) \end{array}}}
\vnbprop{converges-iff-dist-null}{\vnbstmt{\begin{array}{@{}l@{}} \forall t.; \forall f \in (\mathbb{N} \to \operatorname{pts}(t)), v \in \operatorname{pts}(t). \\ \quad \operatorname{is\text{-}metric\text{-}space}(t) \Rightarrow  \\ \quad \quad \operatorname{converges\text{-}to}(t, f, v) \Leftrightarrow \operatorname{converges\text{-}to}(\operatorname{rr\text{-}ms}, (\lambda j \in \mathbb{N}.\; \operatorname{dist}(t)(f(j), v)), 0) \end{array}}}

\vnbbookfile{theorem-library/sequential-continuity}{theorem-library/sequential-continuity}
\noindent{\small sequential-continuity.scm -- THE SEQUENTIAL CHARACTERIZATION OF CONTINUITY, both directions PROVEN:}\par
\vnbprop{continuous-at-implies-sequential}{\vnbstmt{\begin{array}{@{}l@{}} \forall s, t, f, a.; \forall q \in \operatorname{pts}(s)^{\mathbb{N}}. \\ \quad \operatorname{is\text{-}metric\text{-}space}(s) \wedge \\ \quad \operatorname{is\text{-}metric\text{-}space}(t) \wedge \\ \quad f \in (\operatorname{pts}(s) \to \operatorname{pts}(t)) \wedge \\ \quad a \in \operatorname{pts}(s) \wedge \\ \quad \operatorname{is\text{-}continuous\text{-}at}(s, t, f, a) \wedge \\ \quad \operatorname{converges\text{-}to}(s, q, a) \Rightarrow \\ \quad \quad \operatorname{converges\text{-}to}(t, \operatorname{compose}(f, q), f(a)) \end{array}}}
\vnbprop{not-continuous-witness}{\vnbstmt{\begin{array}{@{}l@{}} \forall s, t, f, a, b, l. \\ \quad \neg \forall \beta \in \operatorname{pts}(s).\; ((\operatorname{dist}(s)(a, \beta) \leq l) \Rightarrow (\operatorname{dist}(t)(f(a), f(\beta)) \leq b)) \Rightarrow  \\ \quad \quad \exists \beta \in \operatorname{pts}(s). \\ \quad \quad \quad (\operatorname{dist}(s)(a, \beta) \leq l) \wedge \neg (\operatorname{dist}(t)(f(a), f(\beta)) \leq b) \end{array}}}
\vnbprop{rr-recip-antitone}{\vnbstmt{\forall u, v \in \mathbb{R}.\; (((0 < u) \wedge (u \leq v)) \Rightarrow (\operatorname{recip}(v) \leq \operatorname{recip}(u)))}}
\vnbprop{nn-recip-succ-antitone}{\vnbstmt{\forall m, n \in \mathbb{N}.\; ((m \leq n) \Rightarrow (\operatorname{recip}((n + 1)) \leq \operatorname{recip}((m + 1))))}}
\vnbprop{sequential-implies-continuous-at}{\vnbstmt{\begin{array}{@{}l@{}} \forall s, t, f, a. \\ \quad \operatorname{is\text{-}metric\text{-}space}(s) \wedge \\ \quad \operatorname{is\text{-}metric\text{-}space}(t) \wedge \\ \quad f \in (\operatorname{pts}(s) \to \operatorname{pts}(t)) \wedge \\ \quad a \in \operatorname{pts}(s) \wedge \\ \quad \forall q \in \operatorname{pts}(s)^{\mathbb{N}}. \\ \quad \quad \operatorname{converges\text{-}to}(s, q, a) \Rightarrow \operatorname{converges\text{-}to}(t, \operatorname{compose}(f, q), f(a)) \Rightarrow \\ \quad \quad \operatorname{is\text{-}continuous\text{-}at}(s, t, f, a) \end{array}}}
\vnbprop{continuous-at-iff-sequential}{\vnbstmt{\begin{array}{@{}l@{}} \forall s, t, f, a. \\ \quad (\operatorname{is\text{-}metric\text{-}space}(s) \wedge (\operatorname{is\text{-}metric\text{-}space}(t) \wedge ((f \in (\operatorname{pts}(s) \to \operatorname{pts}(t))) \wedge (a \in \operatorname{pts}(s))))) \Rightarrow  \\ \quad \quad \operatorname{is\text{-}continuous\text{-}at}(s, t, f, a) \\ \quad \quad \Leftrightarrow\; \forall q \in \operatorname{pts}(s)^{\mathbb{N}}. \\ \quad \quad \quad \operatorname{converges\text{-}to}(s, q, a) \Rightarrow \operatorname{converges\text{-}to}(t, \operatorname{compose}(f, q), f(a)) \end{array}}}

\vnbbookfile{theorem-library/recip-succ-null}{theorem-library/recip-succ-null}
\noindent{\small recip-succ-null.scm -- 1/(n+1) -$>$ 0, PROVEN.}\par
\vnbprop{rr-recip-succ-converges-to-zero}{\vnbstmt{\operatorname{converges\text{-}to}(\operatorname{rr\text{-}ms}, (\lambda k \in \mathbb{N}.\; \operatorname{recip}((k + 1))), 0)}}

\vnbbookfile{theorem-library/product-convergence}{theorem-library/product-convergence}
\noindent{\small product-convergence.scm -- CONVERGENCE IN THE COUNTABLE PRODUCT IS EXACTLY COORDINATEWISE CONVERGENCE, proven both ways, retiring the asserted `product-convergence-coordinatewise' of structure-library/product-metric.scm -- the defining property of the PRODUCT TOPOLOGY, and the rung every later}\par
\vnbprop{rr-null-squeeze}{\vnbstmt{\begin{array}{@{}l@{}} \forall f, h \in (\mathbb{N} \to \mathbb{R}). \\ \quad (\operatorname{converges\text{-}to}(\operatorname{rr\text{-}ms}, f, 0) \wedge \forall j \in \mathbb{N}.\; ((0 \leq h(j)) \wedge (h(j) \leq f(j)))) \Rightarrow  \\ \quad \quad \operatorname{converges\text{-}to}(\operatorname{rr\text{-}ms}, h, 0) \end{array}}}
\vnbprop{converges-to-transfer}{\vnbstmt{\begin{array}{@{}l@{}} \forall t.; \forall f, h \in (\mathbb{N} \to \operatorname{pts}(t)), v \in \operatorname{pts}(t). \\ \quad (\operatorname{is\text{-}metric\text{-}space}(t) \wedge (\operatorname{converges\text{-}to}(t, f, v) \wedge \forall j \in \mathbb{N}.\; (h(j) = f(j)))) \Rightarrow  \\ \quad \quad \operatorname{converges\text{-}to}(t, h, v) \end{array}}}
\vnbprop{product-coord-weighted-null}{\vnbstmt{\begin{array}{@{}l@{}} \forall s, w.; \forall q \in (\mathbb{N} \to \operatorname{product\text{-}carrier}(s)), v \in \operatorname{product\text{-}carrier}(s), n \in \mathbb{N}. \\ \quad (\operatorname{is\text{-}ms\text{-}sequence}(s) \wedge (\operatorname{summable\text{-}weight}(w) \wedge \operatorname{converges\text{-}to}(s(n), (\lambda k \in \mathbb{N}.\; q(k)(n)), v(n)))) \Rightarrow  \\ \quad \quad \operatorname{converges\text{-}to}(\operatorname{rr\text{-}ms}, (\lambda j \in \mathbb{N}.\; (w(n) \cdot \operatorname{dist}(\operatorname{bdd\text{-}metric}(s(n)))(q(j)(n), v(n)))), 0) \end{array}}}
\vnbprop{product-coord-from-weighted-null}{\vnbstmt{\begin{array}{@{}l@{}} \forall s, w.; \forall q \in (\mathbb{N} \to \operatorname{product\text{-}carrier}(s)), v \in \operatorname{product\text{-}carrier}(s), n \in \mathbb{N}. \\ \quad \operatorname{is\text{-}ms\text{-}sequence}(s) \wedge \\ \quad \operatorname{summable\text{-}weight}(w) \wedge \\ \quad \operatorname{converges\text{-}to}(\operatorname{rr\text{-}ms}, (\lambda j \in \mathbb{N}.\; (w(n) \cdot \operatorname{dist}(\operatorname{bdd\text{-}metric}(s(n)))(q(j)(n), v(n)))), 0) \Rightarrow \\ \quad \quad \operatorname{converges\text{-}to}(s(n), (\lambda k \in \mathbb{N}.\; q(k)(n)), v(n)) \end{array}}}
\vnbprop{product-convergence-coordinatewise-bwd}{\vnbstmt{\begin{array}{@{}l@{}} \forall s, w.; \forall q \in (\mathbb{N} \to \operatorname{product\text{-}carrier}(s)), v \in \operatorname{product\text{-}carrier}(s). \\ \quad \operatorname{is\text{-}ms\text{-}sequence}(s) \wedge \\ \quad \operatorname{summable\text{-}weight}(w) \wedge \\ \quad \forall n \in \mathbb{N}.\; \operatorname{converges\text{-}to}(s(n), (\lambda k \in \mathbb{N}.\; q(k)(n)), v(n)) \Rightarrow \\ \quad \quad \operatorname{converges\text{-}to}(\operatorname{product\text{-}metric\text{-}w}(s, w), q, v) \end{array}}}
\vnbprop{product-convergence-coordinatewise-fwd}{\vnbstmt{\begin{array}{@{}l@{}} \forall s, w.; \forall q \in (\mathbb{N} \to \operatorname{product\text{-}carrier}(s)), v \in \operatorname{product\text{-}carrier}(s), n \in \mathbb{N}. \\ \quad (\operatorname{is\text{-}ms\text{-}sequence}(s) \wedge (\operatorname{summable\text{-}weight}(w) \wedge \operatorname{converges\text{-}to}(\operatorname{product\text{-}metric\text{-}w}(s, w), q, v))) \Rightarrow  \\ \quad \quad \operatorname{converges\text{-}to}(s(n), (\lambda k \in \mathbb{N}.\; q(k)(n)), v(n)) \end{array}}}
\vnbprop{product-convergence-coordinatewise}{\vnbstmt{\begin{array}{@{}l@{}} \forall s, w.; \forall q \in (\mathbb{N} \to \operatorname{product\text{-}carrier}(s)), v \in \operatorname{product\text{-}carrier}(s). \\ \quad (\operatorname{is\text{-}ms\text{-}sequence}(s) \wedge \operatorname{summable\text{-}weight}(w)) \Rightarrow  \\ \quad \quad \operatorname{converges\text{-}to}(\operatorname{product\text{-}metric\text{-}w}(s, w), q, v) \\ \quad \quad \Leftrightarrow\; \forall n \in \mathbb{N}.\; \operatorname{converges\text{-}to}(s(n), (\lambda k \in \mathbb{N}.\; q(k)(n)), v(n)) \end{array}}}

\vnbbookfile{theorem-library/rake-diagonal-subseq}{theorem-library/rake-diagonal-subseq}
\noindent{\small theorem-library/rake-diagonal-subseq.scm ==================================================================== The cross-coordinate DIAGONAL argument, and the sequential Tychonoff assembly on top of it.  Two supports of theorem-library/seq-compact-product}\par
\vnbprop{coordinatewise-diagonal-subseq}{\vnbstmt{\begin{array}{@{}l@{}} \forall s.; \forall q \in (\mathbb{N} \to \operatorname{product\text{-}carrier}(s)). \\ \quad (\operatorname{is\text{-}ms\text{-}sequence}(s) \wedge \forall n \in \mathbb{N}.\; \operatorname{seq\text{-}compact}(s(n))) \Rightarrow  \\ \quad \quad \exists a. \\ \quad \quad \quad \operatorname{strictly\text{-}mono\text{-}nn}(a) \wedge \\ \quad \quad \quad \exists l \in \operatorname{product\text{-}carrier}(s). \forall n \in \mathbb{N}. \\ \quad \quad \quad \quad \operatorname{converges\text{-}to}(s(n), (\lambda k \in \mathbb{N}.\; \operatorname{subseq}(q, a)(k)(n)), l(n)) \end{array}}}
\vnbprop{seq-compact-countable-product}{\vnbstmt{\begin{array}{@{}l@{}} \forall s. \\ \quad (\operatorname{is\text{-}ms\text{-}sequence}(s) \wedge \forall n \in \mathbb{N}.\; \operatorname{seq\text{-}compact}(s(n))) \Rightarrow  \\ \quad \quad \operatorname{seq\text{-}compact}(\operatorname{product\text{-}metric}(s)) \end{array}}}

\vnbbookfile{structure-library/isom-vocab}{structure-library/isom-vocab}
\noindent{\small isom-vocab.scm -- vocabulary for calculus.pdf ch. 3, sections 6 and 8: uniform and pointwise convergence of a sequence of functions between METRIC spaces (Prop 3.33), and the group of isometric isomorphisms of a metric space with its sup metric (Props 3.37, 3.38).  DEFINITIONS ONLY; the laws are proven in the ...}\par
\par\medskip\noindent\textbf{Definition} (\texttt{converges-uniformly-to}).\;For all $\operatorname{cvus\_}$, $\operatorname{cvut\_}$, $\operatorname{cvuf\_}$ and $\operatorname{cvug\_}$, 
\begin{equation*}
\begin{array}{@{}l@{}} \operatorname{converges\text{-}uniformly\text{-}to}(\operatorname{cvus\_}, \operatorname{cvut\_}, \operatorname{cvuf\_}, \operatorname{cvug\_}) \\ \Leftrightarrow\; \operatorname{is\text{-}metric\text{-}space}(\operatorname{cvus\_}) \wedge \\ \operatorname{is\text{-}metric\text{-}space}(\operatorname{cvut\_}) \wedge \\ \operatorname{cvuf\_} \in (\mathbb{N} \to (\operatorname{pts}(\operatorname{cvus\_}) \to \operatorname{pts}(\operatorname{cvut\_}))) \wedge \\ \operatorname{cvug\_} \in (\operatorname{pts}(\operatorname{cvus\_}) \to \operatorname{pts}(\operatorname{cvut\_})) \wedge \\ \forall \operatorname{cvue\_}. \\ \quad \operatorname{pos\text{-}rr}(\operatorname{cvue\_}) \Rightarrow  \\ \quad \quad \exists \operatorname{cvun\_} \in \mathbb{N}. \forall \operatorname{cvuk\_} \in \mathbb{N}. \\ \quad \quad \quad (\operatorname{cvun\_} \leq \operatorname{cvuk\_}) \Rightarrow  \\ \quad \quad \quad \quad \forall \operatorname{cvux\_} \in \operatorname{pts}(\operatorname{cvus\_}). \\ \quad \quad \quad \quad \quad \operatorname{dist}(\operatorname{cvut\_})(\operatorname{cvuf\_}(\operatorname{cvuk\_})(\operatorname{cvux\_}), \operatorname{cvug\_}(\operatorname{cvux\_})) \leq \operatorname{cvue\_} \end{array}
\end{equation*}

\par\medskip\noindent\textbf{Definition} (\texttt{converges-ptwise-to}).\;For all $\operatorname{cvps\_}$, $\operatorname{cvpt\_}$, $\operatorname{cvpf\_}$ and $\operatorname{cvpg\_}$, 
\begin{equation*}
\begin{array}{@{}l@{}} \operatorname{converges\text{-}ptwise\text{-}to}(\operatorname{cvps\_}, \operatorname{cvpt\_}, \operatorname{cvpf\_}, \operatorname{cvpg\_}) \\ \Leftrightarrow\; \operatorname{is\text{-}metric\text{-}space}(\operatorname{cvps\_}) \wedge \\ \operatorname{is\text{-}metric\text{-}space}(\operatorname{cvpt\_}) \wedge \\ \operatorname{cvpf\_} \in (\mathbb{N} \to (\operatorname{pts}(\operatorname{cvps\_}) \to \operatorname{pts}(\operatorname{cvpt\_}))) \wedge \\ \operatorname{cvpg\_} \in (\operatorname{pts}(\operatorname{cvps\_}) \to \operatorname{pts}(\operatorname{cvpt\_})) \wedge \\ \forall \operatorname{cvpx\_} \in \operatorname{pts}(\operatorname{cvps\_}), \operatorname{cvpe\_}. \\ \quad \operatorname{pos\text{-}rr}(\operatorname{cvpe\_}) \Rightarrow  \\ \quad \quad \exists \operatorname{cvpn\_} \in \mathbb{N}. \forall \operatorname{cvpk\_} \in \mathbb{N}. \\ \quad \quad \quad (\operatorname{cvpn\_} \leq \operatorname{cvpk\_}) \Rightarrow  \\ \quad \quad \quad \quad \operatorname{dist}(\operatorname{cvpt\_})(\operatorname{cvpf\_}(\operatorname{cvpk\_})(\operatorname{cvpx\_}), \operatorname{cvpg\_}(\operatorname{cvpx\_})) \leq \operatorname{cvpe\_} \end{array}
\end{equation*}

\par\medskip\noindent\textbf{Definition} (\texttt{is-equicontinuous-set}).\;For all $\operatorname{eqss\_}$, $\operatorname{eqst\_}$ and $\operatorname{eqse\_}$, 
\begin{equation*}
\begin{array}{@{}l@{}} \operatorname{is\text{-}equicontinuous\text{-}set}(\operatorname{eqss\_}, \operatorname{eqst\_}, \operatorname{eqse\_}) \\ \Leftrightarrow\; \operatorname{is\text{-}metric\text{-}space}(\operatorname{eqss\_}) \wedge \\ \operatorname{is\text{-}metric\text{-}space}(\operatorname{eqst\_}) \wedge \\ \operatorname{eqse\_} \subseteq (\operatorname{pts}(\operatorname{eqss\_}) \to \operatorname{pts}(\operatorname{eqst\_})) \wedge \\ \forall \operatorname{eqsx\_} \in \operatorname{pts}(\operatorname{eqss\_}), \operatorname{eqsp\_}. \\ \quad \operatorname{pos\text{-}rr}(\operatorname{eqsp\_}) \Rightarrow  \\ \quad \quad \exists \operatorname{eqsd\_}. \\ \quad \quad \quad \operatorname{pos\text{-}rr}(\operatorname{eqsd\_}) \wedge \\ \quad \quad \quad \forall \operatorname{eqsf\_} \in \operatorname{eqse\_}, \operatorname{eqsy\_} \in \operatorname{pts}(\operatorname{eqss\_}). \\ \quad \quad \quad \quad (\operatorname{dist}(\operatorname{eqss\_})(\operatorname{eqsx\_}, \operatorname{eqsy\_}) < \operatorname{eqsd\_}) \Rightarrow  \\ \quad \quad \quad \quad \quad \operatorname{dist}(\operatorname{eqst\_})(\operatorname{eqsf\_}(\operatorname{eqsx\_}), \operatorname{eqsf\_}(\operatorname{eqsy\_})) < \operatorname{eqsp\_} \end{array}
\end{equation*}

\par\medskip\noindent\textbf{Definition} (\texttt{sup-dist}).\;$\operatorname{sup\text{-}dist}(\operatorname{sds\_}, \operatorname{sdt\_}, \operatorname{sdf\_}, \operatorname{sdg\_}) :=$
\begin{equation*}
\operatorname{esup}(\operatorname{image}((\lambda \operatorname{sdx\_} \in \operatorname{pts}(\operatorname{sds\_}).\; \operatorname{dist}(\operatorname{sdt\_})(\operatorname{sdf\_}(\operatorname{sdx\_}), \operatorname{sdg\_}(\operatorname{sdx\_}))), \operatorname{pts}(\operatorname{sds\_})))
\end{equation*}

\par\medskip\noindent\textbf{Definition} (\texttt{isometry-group}).\;$\operatorname{isometry\text{-}group}(\operatorname{igs\_}) :=$
\begin{equation*}
\{\operatorname{igf\_} \in \mathrm{Bij}(\operatorname{pts}(\operatorname{igs\_}), \operatorname{pts}(\operatorname{igs\_})) : \operatorname{is\text{-}isometry}(\operatorname{igs\_}, \operatorname{igs\_}, \operatorname{igf\_})\}
\end{equation*}

\par\medskip\noindent\textbf{Definition} (\texttt{isometry-group-ms}).\;$\operatorname{isometry\text{-}group\text{-}ms}(\operatorname{ims\_}) :=$
\begin{equation*}
{[}\operatorname{isometry\text{-}group}(\operatorname{ims\_}), (\lambda {[}\operatorname{imu\_}, \operatorname{imv\_}] \in (\operatorname{isometry\text{-}group}(\operatorname{ims\_}) \times \operatorname{isometry\text{-}group}(\operatorname{ims\_})).\; \operatorname{sup\text{-}dist}(\operatorname{ims\_}, \operatorname{ims\_}, \operatorname{imu\_}, \operatorname{imv\_}))]
\end{equation*}


\vnbbookfile{theorem-library/isom-dense}{theorem-library/isom-dense}
\noindent{\small isom-dense.scm -- three general facts on sequences of functions between metric spaces, the machinery of calculus.pdf Lemma 3.36 and Prop 3.37.  Agent ISOM, 2026-10-26.}\par
\vnbprop{equicont-dense-ptwise-limit}{\vnbstmt{\begin{array}{@{}l@{}} \forall s, t, a, q.; \forall x \in \operatorname{pts}(s). \\ \quad \operatorname{is\text{-}complete}(t) \wedge \\ \quad \operatorname{is\text{-}equicontinuous}(s, t, a) \wedge \\ \quad \operatorname{is\text{-}dense\text{-}seq}(s, q) \wedge \\ \quad \forall m \in \mathbb{N}. \exists y \in \operatorname{pts}(t). \\ \quad \quad \operatorname{converges\text{-}to}(t, (\lambda k \in \mathbb{N}.\; a(k)(q(m))), y) \Rightarrow \\ \quad \quad \exists y \in \operatorname{pts}(t).\; \operatorname{converges\text{-}to}(t, (\lambda k \in \mathbb{N}.\; a(k)(x)), y) \end{array}}}
\vnbprop{ptwise-limit-function}{\vnbstmt{\begin{array}{@{}l@{}} \forall s, t, m. \\ \quad \operatorname{is\text{-}metric\text{-}space}(s) \wedge \\ \quad \operatorname{is\text{-}metric\text{-}space}(t) \wedge \\ \quad m \in (\mathbb{N} \to (\operatorname{pts}(s) \to \operatorname{pts}(t))) \wedge \\ \quad \forall x \in \operatorname{pts}(s). \exists y \in \operatorname{pts}(t). \\ \quad \quad \operatorname{converges\text{-}to}(t, (\lambda k \in \mathbb{N}.\; m(k)(x)), y) \Rightarrow \\ \quad \quad \exists g.\; \operatorname{converges\text{-}ptwise\text{-}to}(s, t, m, g) \end{array}}}
\vnbprop{compact-diagonal-on-seq}{\vnbstmt{\begin{array}{@{}l@{}} \forall t, d, s, q. \\ \quad (\operatorname{is\text{-}compact}(t) \wedge ((s \in (\mathbb{N} \to (d \to \operatorname{pts}(t)))) \wedge (q \in (\mathbb{N} \to d)))) \Rightarrow  \\ \quad \quad \exists a. \\ \quad \quad \quad \operatorname{strictly\text{-}mono\text{-}nn}(a) \wedge \\ \quad \quad \quad \forall m \in \mathbb{N}. \exists y \in \operatorname{pts}(t). \\ \quad \quad \quad \quad \operatorname{converges\text{-}to}(t, (\lambda k \in \mathbb{N}.\; \operatorname{subseq}(s, a)(k)(q(m))), y) \end{array}}}

\vnbbookfile{structure-library/isom2-vocab}{structure-library/isom2-vocab}
\noindent{\small isom2-vocab.scm -- the group structure of the isometry group (calculus.pdf ch. 3 sec. 8: "In fact, Gamma is a compact metric group") and the metric group. DEFINITIONS ONLY; laws in theorem-library/isom2-group.scm.  Agent ISOM, 2026-10-26 (NEW DEFINITIONS, the user's to confirm).}\par
\vnbstructdef{metric-group}{A metric group $s$ is a 5-tuple whose components are listed below.\par
\begin{enumerate}[nosep]
\item $\operatorname{carr}(s) = \operatorname{nth}(1, s)$, a carrier (\texttt{carriers}): $\operatorname{carr}(s) \in \mathbf{Set}$.
\item $\operatorname{opr}(s) = \operatorname{nth}(2, s)$, an operation (\texttt{op}): $\operatorname{opr}(s) \in ((\operatorname{carr}(s) \times \operatorname{carr}(s)) \to \operatorname{carr}(s))$.
\item $\operatorname{iden}(s) = \operatorname{nth}(3, s)$, a constant (\texttt{constant}): $\operatorname{iden}(s) \in \operatorname{carr}(s)$.
\item $\operatorname{inv}(s) = \operatorname{nth}(4, s)$, an operation (\texttt{op}): $\operatorname{inv}(s) \in (\operatorname{carr}(s) \to \operatorname{carr}(s))$.
\item $\operatorname{gdist}(s) = \operatorname{nth}(5, s)$, an operation (\texttt{op}): $\operatorname{gdist}(s) \in ((\operatorname{carr}(s) \times \operatorname{carr}(s)) \to \mathbb{R})$.
\end{enumerate}
The laws (the \texttt{property} and \texttt{law} clauses):
\begin{itemize}[nosep]
\item $\operatorname{is\text{-}associative}(\operatorname{opr}(s), \operatorname{carr}(s))$
\item $\operatorname{is\text{-}identity}(\operatorname{opr}(s), \operatorname{iden}(s), \operatorname{carr}(s))$
\item $\operatorname{has\text{-}inverses}(\operatorname{opr}(s), \operatorname{iden}(s), \operatorname{inv}(s), \operatorname{carr}(s))$
\item $\operatorname{is\text{-}metric}(\operatorname{gdist}(s), \operatorname{carr}(s))$
\item
\begin{equation*}
\begin{array}{@{}l@{}} \forall \operatorname{mgu\_}, \operatorname{mgv\_} \in \operatorname{carr}(s), \operatorname{mge\_}. \\ \quad \operatorname{pos\text{-}rr}(\operatorname{mge\_}) \Rightarrow  \\ \quad \quad \exists \operatorname{mgd\_}. \\ \quad \quad \quad \operatorname{pos\text{-}rr}(\operatorname{mgd\_}) \wedge \\ \quad \quad \quad \forall \operatorname{mgp\_}, \operatorname{mgq\_} \in \operatorname{carr}(s). \\ \quad \quad \quad \quad (\operatorname{gdist}(s)(\operatorname{mgu\_}, \operatorname{mgp\_}) \leq \operatorname{mgd\_}) \Rightarrow  \\ \quad \quad \quad \quad (\operatorname{gdist}(s)(\operatorname{mgv\_}, \operatorname{mgq\_}) \leq \operatorname{mgd\_}) \Rightarrow  \\ \quad \quad \quad \quad \quad \operatorname{gdist}(s)(\operatorname{opr}(s)(\operatorname{mgu\_}, \operatorname{mgv\_}), \operatorname{opr}(s)(\operatorname{mgp\_}, \operatorname{mgq\_})) \leq \operatorname{mge\_} \end{array}
\end{equation*}

\item
\begin{equation*}
\begin{array}{@{}l@{}} \forall \operatorname{mga\_} \in \operatorname{carr}(s), \operatorname{mge\_}. \\ \quad \operatorname{pos\text{-}rr}(\operatorname{mge\_}) \Rightarrow  \\ \quad \quad \exists \operatorname{mgd\_}. \\ \quad \quad \quad \operatorname{pos\text{-}rr}(\operatorname{mgd\_}) \wedge \\ \quad \quad \quad \forall \operatorname{mgb\_} \in \operatorname{carr}(s). \\ \quad \quad \quad \quad (\operatorname{gdist}(s)(\operatorname{mga\_}, \operatorname{mgb\_}) \leq \operatorname{mgd\_}) \Rightarrow  \\ \quad \quad \quad \quad \quad \operatorname{gdist}(s)(\operatorname{inv}(s)(\operatorname{mga\_}), \operatorname{inv}(s)(\operatorname{mgb\_})) \leq \operatorname{mge\_} \end{array}
\end{equation*}

\end{itemize}
The characteristic law (the axiom \texttt{is-metric-group}): $\operatorname{is\text{-}metric\text{-}group}(s)$ holds if and only if $\operatorname{length}(s) = 5$, every component has the typing listed with it, every law holds, and $s \in \mathbf{Set}$ (the structure is a set).
}
\par\medskip\noindent\textbf{Definition} (\texttt{isometry-group-grp}).\;$\operatorname{isometry\text{-}group\text{-}grp}(\operatorname{igs\_}) :=$
\begin{equation*}
{[}\operatorname{isometry\text{-}group}(\operatorname{igs\_}), (\lambda {[}\operatorname{igu\_}, \operatorname{igv\_}] \in (\operatorname{isometry\text{-}group}(\operatorname{igs\_}) \times \operatorname{isometry\text{-}group}(\operatorname{igs\_})).\; (\lambda \operatorname{igx\_} \in \operatorname{pts}(\operatorname{igs\_}).\; \operatorname{igu\_}(\operatorname{igv\_}(\operatorname{igx\_})))), (\lambda \operatorname{igx\_} \in \operatorname{pts}(\operatorname{igs\_}).\; \operatorname{igx\_}), (\lambda \operatorname{igi\_} \in \operatorname{isometry\text{-}group}(\operatorname{igs\_}).\; \operatorname{igi\_}^{-1})]
\end{equation*}

\par\medskip\noindent\textbf{Definition} (\texttt{is-group-metric}).\;For all $\operatorname{mgm\_}$ and $\operatorname{mgg\_}$, 
\begin{equation*}
\begin{array}{@{}l@{}} \operatorname{is\text{-}group\text{-}metric}(\operatorname{mgm\_}, \operatorname{mgg\_}) \\ \Leftrightarrow\; \operatorname{is\text{-}metric\text{-}space}(\operatorname{mgm\_}) \wedge \\ \operatorname{is\text{-}group}(\operatorname{mgg\_}) \wedge \\ \operatorname{pts}(\operatorname{mgm\_}) = \operatorname{carr}(\operatorname{mgg\_}) \wedge \\ \forall \operatorname{mgu\_}, \operatorname{mgv\_} \in \operatorname{carr}(\operatorname{mgg\_}), \operatorname{mge\_}. \\ \quad \operatorname{pos\text{-}rr}(\operatorname{mge\_}) \Rightarrow  \\ \quad \quad \exists \operatorname{mgd\_}. \\ \quad \quad \quad \operatorname{pos\text{-}rr}(\operatorname{mgd\_}) \wedge \\ \quad \quad \quad \forall \operatorname{mgp\_}, \operatorname{mgq\_} \in \operatorname{carr}(\operatorname{mgg\_}). \\ \quad \quad \quad \quad (\operatorname{dist}(\operatorname{mgm\_})(\operatorname{mgu\_}, \operatorname{mgp\_}) \leq \operatorname{mgd\_}) \Rightarrow  \\ \quad \quad \quad \quad (\operatorname{dist}(\operatorname{mgm\_})(\operatorname{mgv\_}, \operatorname{mgq\_}) \leq \operatorname{mgd\_}) \Rightarrow  \\ \quad \quad \quad \quad \quad \operatorname{dist}(\operatorname{mgm\_})(\operatorname{opr}(\operatorname{mgg\_})(\operatorname{mgu\_}, \operatorname{mgv\_}), \operatorname{opr}(\operatorname{mgg\_})(\operatorname{mgp\_}, \operatorname{mgq\_})) \leq \operatorname{mge\_} \wedge \\ \operatorname{is\text{-}continuous}(\operatorname{mgm\_}, \operatorname{mgm\_}, \operatorname{inv}(\operatorname{mgg\_})) \end{array}
\end{equation*}

\par\medskip\noindent\textbf{Definition} (\texttt{metric-group-ms}).\;$\operatorname{metric\text{-}group\text{-}ms}(\operatorname{mgr\_}) :=$ ${[}\operatorname{carr}(\operatorname{mgr\_}), \operatorname{gdist}(\operatorname{mgr\_})]$
\par\medskip\noindent\textbf{Definition} (\texttt{metric-group-group}).\;$\operatorname{metric\text{-}group\text{-}group}(\operatorname{mgr\_}) :=$ ${[}\operatorname{carr}(\operatorname{mgr\_}), \operatorname{opr}(\operatorname{mgr\_}), \operatorname{iden}(\operatorname{mgr\_}), \operatorname{inv}(\operatorname{mgr\_})]$
\par\medskip\noindent\textbf{Definition} (\texttt{make-metric-group}).\;$\operatorname{make\text{-}metric\text{-}group}(\operatorname{mmm\_}, \operatorname{mmg\_}) :=$ ${[}\operatorname{carr}(\operatorname{mmg\_}), \operatorname{opr}(\operatorname{mmg\_}), \operatorname{iden}(\operatorname{mmg\_}), \operatorname{inv}(\operatorname{mmg\_}), \operatorname{dist}(\operatorname{mmm\_})]$
\par\medskip\noindent\textbf{Definition} (\texttt{hom-metric-group}).\;$\operatorname{hom\text{-}metric\text{-}group}(a, b) :=$
\begin{equation*}
\{\operatorname{homf\_} \in (\operatorname{carr}(a) \to \operatorname{carr}(b)) : \operatorname{is\text{-}hom\text{-}metric\text{-}group}(a, b, \operatorname{homf\_})\}
\end{equation*}

\vnbdef{is-hom-metric-group-def}{\vnbstmt{\begin{array}{@{}l@{}} \forall a, b, f. \\ \quad \operatorname{is\text{-}hom\text{-}metric\text{-}group}(a, b, f) \\ \quad \Leftrightarrow\; \operatorname{is\text{-}metric\text{-}group}(a) \wedge \\ \quad \operatorname{is\text{-}metric\text{-}group}(b) \wedge \\ \quad f \in (\operatorname{carr}(a) \to \operatorname{carr}(b)) \wedge \\ \quad \forall x_{1}, x_{2} \in \operatorname{carr}(a). \\ \quad \quad f(\operatorname{opr}(a)(x_{1}, x_{2})) = \operatorname{opr}(b)(f(x_{1}), f(x_{2})) \wedge \\ \quad f(\operatorname{iden}(a)) = \operatorname{iden}(b) \wedge \\ \quad \forall x_{1} \in \operatorname{carr}(a).\; (f(\operatorname{inv}(a)(x_{1})) = \operatorname{inv}(b)(f(x_{1}))) \wedge \\ \quad \forall x_{1}, x_{2} \in \operatorname{carr}(a). \\ \quad \quad \operatorname{gdist}(a)(x_{1}, x_{2}) = \operatorname{gdist}(b)(f(x_{1}), f(x_{2})) \end{array}}}
\vnbdef{metric-group-class}{\vnbstmt{\forall s.\; ((s \in \operatorname{metric\text{-}group}) \Leftrightarrow \operatorname{is\text{-}metric\text{-}group}(s))}}
\vnbprop{is-metric-group-is-set}{\vnbstmt{\forall s.\; (\operatorname{is\text{-}metric\text{-}group}(s) \Rightarrow (s \in \mathbf{Set}))}}

\vnbbookfile{theorem-library/isom2-dense}{theorem-library/isom2-dense}
\noindent{\small isom2-dense.scm -- calculus.pdf Remark 3.32 and Lemma 3.36 in full: the pointwise limit of an equicontinuous sequence is continuous, and an equicontinuous sequence into a COMPLETE space that converges on a dense sequence converges everywhere to a continuous limit.  Agent ISOM, 2026-10-26.}\par
\vnbprop{equicont-ptwise-limit-continuous}{\vnbstmt{\begin{array}{@{}l@{}} \forall s, t, m, g. \\ \quad (\operatorname{is\text{-}equicontinuous}(s, t, m) \wedge \operatorname{converges\text{-}ptwise\text{-}to}(s, t, m, g)) \Rightarrow  \\ \quad \quad \operatorname{is\text{-}continuous}(s, t, g) \end{array}}}
\vnbprop{equicont-dense-limit-continuous}{\vnbstmt{\begin{array}{@{}l@{}} \forall s, t, a, q. \\ \quad \operatorname{is\text{-}complete}(t) \wedge \\ \quad \operatorname{is\text{-}equicontinuous}(s, t, a) \wedge \\ \quad \operatorname{is\text{-}dense\text{-}seq}(s, q) \wedge \\ \quad \forall m \in \mathbb{N}. \exists y \in \operatorname{pts}(t). \\ \quad \quad \operatorname{converges\text{-}to}(t, (\lambda k \in \mathbb{N}.\; a(k)(q(m))), y) \Rightarrow \\ \quad \quad \exists g.\; (\operatorname{converges\text{-}ptwise\text{-}to}(s, t, a, g) \wedge \operatorname{is\text{-}continuous}(s, t, g)) \end{array}}}

\vnbbookfile{theorem-library/tychonoff-proof}{theorem-library/tychonoff-proof}
\noindent{\small theorem-library/tychonoff-proof.scm ==================================================================== compact-countable-product (countable Tychonoff headline) -- PROVEN.}\par
\vnbprop{compact-countable-product}{\vnbstmt{\begin{array}{@{}l@{}} \forall s. \\ \quad (\operatorname{is\text{-}ms\text{-}sequence}(s) \wedge \forall n \in \mathbb{N}.\; \operatorname{is\text{-}compact}(s(n))) \Rightarrow  \\ \quad \quad \operatorname{is\text{-}compact}(\operatorname{product\text{-}metric}(s)) \end{array}}}

\vnbbookfile{theorem-library/rake-setoid2}{theorem-library/rake-setoid2}
\noindent{\small rake-setoid2.scm -- BATCH O of the 2026-09-17 rake: the SETOID quotient and the METRIC COMPLETION.  Twelve of the thirteen asserted leaves of the batch are PROVEN here; the thirteenth is a STATEMENT DEFECT and is recorded below. Every statement is its support's statement UNCHANGED (copied from the}\par
\vnbprop{rko2-cseq-unfold}{\vnbstmt{\forall m.\; (\operatorname{cseq}(m) \simeq \{f \in (\mathbb{N} \to \operatorname{pts}(m)) : \operatorname{is\text{-}cauchy\text{-}seq}(m, f)\})}}
\vnbprop{rko2-cseq-in-fun}{\vnbstmt{\forall m.; \forall f \in \operatorname{cseq}(m).\; (f \in (\mathbb{N} \to \operatorname{pts}(m)))}}
\vnbprop{rko2-completion-dist}{\vnbstmt{\begin{array}{@{}l@{}} \forall m. \\ \quad \operatorname{is\text{-}metric\text{-}space}(m) \Rightarrow (\operatorname{dist}(\operatorname{completion}(m)) \simeq \operatorname{completion\text{-}dist}(m)) \end{array}}}
\vnbprop{rko2-cauchy-setoid-rel}{\vnbstmt{\forall m.\; (\operatorname{is\text{-}metric\text{-}space}(m) \Rightarrow (\operatorname{rel}(\operatorname{cauchy\text{-}setoid}(m)) \simeq \operatorname{crel}(m)))}}
\vnbprop{rko2-crel-unfold}{\vnbstmt{\begin{array}{@{}l@{}} \forall m. \\ \quad (\operatorname{crel}(m) \simeq \{p \in (\operatorname{cseq}(m) \times \operatorname{cseq}(m)) : \operatorname{cseq\text{-}equiv}(m, \operatorname{nth}(1, p), \operatorname{nth}(2, p))\}) \end{array}}}
\vnbprop{rko2-crel-intro}{\vnbstmt{\begin{array}{@{}l@{}} \forall m, p, q. \\ \quad ((p \in \operatorname{cseq}(m)) \wedge ((q \in \operatorname{cseq}(m)) \wedge \operatorname{cseq\text{-}equiv}(m, p, q))) \Rightarrow ({[}p, q] \in \operatorname{crel}(m)) \end{array}}}
\vnbprop{rko2-crel-elim}{\vnbstmt{\forall m, p, q.\; (({[}p, q] \in \operatorname{crel}(m)) \Rightarrow \operatorname{cseq\text{-}equiv}(m, p, q))}}
\vnbprop{rko2-cauchy-setoid-length}{\vnbstmt{\forall m.\; (\operatorname{is\text{-}metric\text{-}space}(m) \Rightarrow (\operatorname{length}(\operatorname{cauchy\text{-}setoid}(m)) = 2))}}
\vnbprop{rko2-crel-is-set}{\vnbstmt{\forall m.\; (\operatorname{is\text{-}metric\text{-}space}(m) \Rightarrow (\operatorname{crel}(m) \in \mathbf{Set}))}}
\vnbprop{rko2-dist-seq-unfold}{\vnbstmt{\forall m, f, g.\; (\operatorname{dist\text{-}seq}(m, f, g) \simeq (\lambda n \in \mathbb{N}.\; \operatorname{dist}(m)(f(n), g(n))))}}
\vnbprop{rko2-dist-seq-type}{\vnbstmt{\begin{array}{@{}l@{}} \forall m, f, g. \\ \quad (\operatorname{is\text{-}metric\text{-}space}(m) \wedge ((f \in (\mathbb{N} \to \operatorname{pts}(m))) \wedge (g \in (\mathbb{N} \to \operatorname{pts}(m))))) \Rightarrow  \\ \quad \quad \operatorname{dist\text{-}seq}(m, f, g) \in (\mathbb{N} \to \mathbb{R}) \end{array}}}
\vnbprop{rko2-null-refl}{\vnbstmt{\begin{array}{@{}l@{}} \forall m, f. \\ \quad (\operatorname{is\text{-}metric\text{-}space}(m) \wedge (f \in (\mathbb{N} \to \operatorname{pts}(m)))) \Rightarrow  \\ \quad \quad \operatorname{converges\text{-}to}(\operatorname{rr\text{-}ms}, \operatorname{dist\text{-}seq}(m, f, f), 0) \end{array}}}
\vnbprop{rko2-null-sym}{\vnbstmt{\begin{array}{@{}l@{}} \forall m, f, g. \\ \quad \operatorname{is\text{-}metric\text{-}space}(m) \wedge \\ \quad f \in (\mathbb{N} \to \operatorname{pts}(m)) \wedge \\ \quad g \in (\mathbb{N} \to \operatorname{pts}(m)) \wedge \\ \quad \operatorname{converges\text{-}to}(\operatorname{rr\text{-}ms}, \operatorname{dist\text{-}seq}(m, f, g), 0) \Rightarrow \\ \quad \quad \operatorname{converges\text{-}to}(\operatorname{rr\text{-}ms}, \operatorname{dist\text{-}seq}(m, g, f), 0) \end{array}}}
\vnbprop{rko2-null-trans}{\vnbstmt{\begin{array}{@{}l@{}} \forall m, f, g, h. \\ \quad \operatorname{is\text{-}metric\text{-}space}(m) \wedge \\ \quad f \in (\mathbb{N} \to \operatorname{pts}(m)) \wedge \\ \quad g \in (\mathbb{N} \to \operatorname{pts}(m)) \wedge \\ \quad h \in (\mathbb{N} \to \operatorname{pts}(m)) \wedge \\ \quad \operatorname{converges\text{-}to}(\operatorname{rr\text{-}ms}, \operatorname{dist\text{-}seq}(m, f, g), 0) \wedge \\ \quad \operatorname{converges\text{-}to}(\operatorname{rr\text{-}ms}, \operatorname{dist\text{-}seq}(m, g, h), 0) \Rightarrow \\ \quad \quad \operatorname{converges\text{-}to}(\operatorname{rr\text{-}ms}, \operatorname{dist\text{-}seq}(m, f, h), 0) \end{array}}}
\vnbprop{cauchy-setoid-is-setoid}{\vnbstmt{\forall m.\; (\operatorname{is\text{-}metric\text{-}space}(m) \Rightarrow \operatorname{is\text{-}setoid}(\operatorname{cauchy\text{-}setoid}(m)))}}
\vnbprop{rko2-quad}{\vnbstmt{\begin{array}{@{}l@{}} \forall m, a, b, c, e. \\ \quad (\operatorname{is\text{-}metric\text{-}space}(m) \wedge ((a \in \operatorname{pts}(m)) \wedge ((b \in \operatorname{pts}(m)) \wedge ((c \in \operatorname{pts}(m)) \wedge (e \in \operatorname{pts}(m)))))) \Rightarrow  \\ \quad \quad \lvert (\operatorname{dist}(m)(a, b) - \operatorname{dist}(m)(c, e)) \rvert \leq (\operatorname{dist}(m)(a, c) + \operatorname{dist}(m)(b, e)) \end{array}}}
\vnbprop{rko2-dist-converges}{\vnbstmt{\begin{array}{@{}l@{}} \forall m, f, g, a, b. \\ \quad \operatorname{is\text{-}metric\text{-}space}(m) \wedge \\ \quad f \in (\mathbb{N} \to \operatorname{pts}(m)) \wedge \\ \quad g \in (\mathbb{N} \to \operatorname{pts}(m)) \wedge \\ \quad a \in \operatorname{pts}(m) \wedge \\ \quad b \in \operatorname{pts}(m) \wedge \\ \quad \operatorname{converges\text{-}to}(m, f, a) \wedge \\ \quad \operatorname{converges\text{-}to}(m, g, b) \Rightarrow \\ \quad \quad \operatorname{converges\text{-}to}(\operatorname{rr\text{-}ms}, \operatorname{dist\text{-}seq}(m, f, g), \operatorname{dist}(m)(a, b)) \end{array}}}
\vnbprop{rko2-const-seq-converges}{\vnbstmt{\begin{array}{@{}l@{}} \forall m, u. \\ \quad (\operatorname{is\text{-}metric\text{-}space}(m) \wedge (u \in \operatorname{pts}(m))) \Rightarrow \operatorname{converges\text{-}to}(m, \operatorname{embed\text{-}seq}(m, u), u) \end{array}}}
\vnbprop{rko2-class-eq-null}{\vnbstmt{\begin{array}{@{}l@{}} \forall m, p, q. \\ \quad \operatorname{is\text{-}metric\text{-}space}(m) \wedge \\ \quad p \in \operatorname{cseq}(m) \wedge \\ \quad q \in \operatorname{cseq}(m) \wedge \\ \quad \operatorname{class}(\operatorname{cauchy\text{-}setoid}(m), p) = \operatorname{class}(\operatorname{cauchy\text{-}setoid}(m), q) \Rightarrow \\ \quad \quad \operatorname{converges\text{-}to}(\operatorname{rr\text{-}ms}, \operatorname{dist\text{-}seq}(m, p, q), 0) \end{array}}}
\vnbprop{rko2-const-equiv-converges}{\vnbstmt{\begin{array}{@{}l@{}} \forall m, u, f. \\ \quad \operatorname{is\text{-}metric\text{-}space}(m) \wedge \\ \quad u \in \operatorname{pts}(m) \wedge \\ \quad f \in (\mathbb{N} \to \operatorname{pts}(m)) \wedge \\ \quad \operatorname{converges\text{-}to}(\operatorname{rr\text{-}ms}, \operatorname{dist\text{-}seq}(m, \operatorname{embed\text{-}seq}(m, u), f), 0) \Rightarrow \\ \quad \quad \operatorname{converges\text{-}to}(m, f, u) \end{array}}}
\vnbprop{embed-isometry}{\vnbstmt{\begin{array}{@{}l@{}} \forall m.; \forall u, v \in \operatorname{pts}(m). \\ \quad \operatorname{is\text{-}metric\text{-}space}(m) \Rightarrow  \\ \quad \quad \operatorname{dist}(\operatorname{completion}(m))(\operatorname{embed}(m)(u), \operatorname{embed}(m)(v)) = \operatorname{dist}(m)(u, v) \end{array}}}

\vnbbookfile{theorem-library/rake-completion-ms}{theorem-library/rake-completion-ms}
\noindent{\small rake-completion-ms.scm -- BATCH 5c-Q of the 2026-09-18 rake: completion-is-metric-space (structure-library/metric-completion.scm:107, asserted `well-known').  PROVEN `modulo 0'.  The statement is the support's, copied from the support site and unchanged:}\par
\vnbprop{r7q-cseq-cauchy}{\vnbstmt{\forall m.; \forall f \in \operatorname{cseq}(m).\; \operatorname{is\text{-}cauchy\text{-}seq}(m, f)}}
\vnbprop{r7q-dist-seq-cauchy}{\vnbstmt{\begin{array}{@{}l@{}} \forall m.; \forall f, g \in \operatorname{cseq}(m). \\ \quad \operatorname{is\text{-}metric\text{-}space}(m) \Rightarrow \operatorname{is\text{-}cauchy\text{-}seq}(\operatorname{rr\text{-}ms}, \operatorname{dist\text{-}seq}(m, f, g)) \end{array}}}
\vnbprop{r7q-dist-seq-converges}{\vnbstmt{\begin{array}{@{}l@{}} \forall m.; \forall f, g \in \operatorname{cseq}(m). \\ \quad \operatorname{is\text{-}metric\text{-}space}(m) \Rightarrow \operatorname{converges}(\operatorname{rr\text{-}ms}, \operatorname{dist\text{-}seq}(m, f, g)) \end{array}}}
\vnbprop{r7q-dist-seq-sym}{\vnbstmt{\begin{array}{@{}l@{}} \forall m.; \forall f, g \in (\mathbb{N} \to \operatorname{pts}(m)), v. \\ \quad (\operatorname{is\text{-}metric\text{-}space}(m) \wedge \operatorname{converges\text{-}to}(\operatorname{rr\text{-}ms}, \operatorname{dist\text{-}seq}(m, f, g), v)) \Rightarrow  \\ \quad \quad \operatorname{converges\text{-}to}(\operatorname{rr\text{-}ms}, \operatorname{dist\text{-}seq}(m, g, f), v) \end{array}}}
\vnbprop{r7q-null-transfer}{\vnbstmt{\begin{array}{@{}l@{}} \forall t.; \forall f, h \in (\mathbb{N} \to \operatorname{pts}(t)), v \in \operatorname{pts}(t). \\ \quad (\operatorname{is\text{-}metric\text{-}space}(t) \wedge (\operatorname{converges\text{-}to}(t, f, v) \wedge \operatorname{converges\text{-}to}(\operatorname{rr\text{-}ms}, \operatorname{dist\text{-}seq}(t, f, h), 0))) \Rightarrow  \\ \quad \quad \operatorname{converges\text{-}to}(t, h, v) \end{array}}}
\vnbprop{r7q-quad-null}{\vnbstmt{\begin{array}{@{}l@{}} \forall m.; \forall f, g, p, q \in (\mathbb{N} \to \operatorname{pts}(m)). \\ \quad \operatorname{is\text{-}metric\text{-}space}(m) \wedge \\ \quad \operatorname{converges\text{-}to}(\operatorname{rr\text{-}ms}, \operatorname{dist\text{-}seq}(m, f, p), 0) \wedge \\ \quad \operatorname{converges\text{-}to}(\operatorname{rr\text{-}ms}, \operatorname{dist\text{-}seq}(m, g, q), 0) \Rightarrow \\ \quad \quad \operatorname{converges\text{-}to}(\operatorname{rr\text{-}ms}, \operatorname{dist\text{-}seq}(\operatorname{rr\text{-}ms}, \operatorname{dist\text{-}seq}(m, f, g), \operatorname{dist\text{-}seq}(m, p, q)), 0) \end{array}}}
\vnbprop{r7q-limit-indep}{\vnbstmt{\begin{array}{@{}l@{}} \forall m.; \forall f, g, p, q \in (\mathbb{N} \to \operatorname{pts}(m)), v. \\ \quad \operatorname{is\text{-}metric\text{-}space}(m) \wedge \\ \quad \operatorname{converges\text{-}to}(\operatorname{rr\text{-}ms}, \operatorname{dist\text{-}seq}(m, f, p), 0) \wedge \\ \quad \operatorname{converges\text{-}to}(\operatorname{rr\text{-}ms}, \operatorname{dist\text{-}seq}(m, g, q), 0) \wedge \\ \quad \operatorname{converges\text{-}to}(\operatorname{rr\text{-}ms}, \operatorname{dist\text{-}seq}(m, f, g), v) \Rightarrow \\ \quad \quad \operatorname{converges\text{-}to}(\operatorname{rr\text{-}ms}, \operatorname{dist\text{-}seq}(m, p, q), v) \end{array}}}
\vnbprop{r7q-iota-dhat}{\vnbstmt{\begin{array}{@{}l@{}} \forall m, x, y. \\ \quad (\operatorname{is\text{-}metric\text{-}space}(m) \wedge ((x \in \operatorname{quotient}(\operatorname{cauchy\text{-}setoid}(m))) \wedge (y \in \operatorname{quotient}(\operatorname{cauchy\text{-}setoid}(m))))) \Rightarrow  \\ \quad \quad ((\iota\, l.\;((l \in \mathbb{R}) \wedge \exists f, g \in \operatorname{cseq}(m).\; ((x = \operatorname{class}(\operatorname{cauchy\text{-}setoid}(m), f)) \wedge ((y = \operatorname{class}(\operatorname{cauchy\text{-}setoid}(m), g)) \wedge \operatorname{converges\text{-}to}(\operatorname{rr\text{-}ms}, \operatorname{dist\text{-}seq}(m, f, g), l))))) \simeq \operatorname{dist}(\operatorname{completion}(m))(x, y)) \end{array}}}
\vnbprop{r7q-dhat-iota}{\vnbstmt{\begin{array}{@{}l@{}} \forall m, x, y. \\ \quad (\operatorname{is\text{-}metric\text{-}space}(m) \wedge ((x \in \operatorname{quotient}(\operatorname{cauchy\text{-}setoid}(m))) \wedge (y \in \operatorname{quotient}(\operatorname{cauchy\text{-}setoid}(m))))) \Rightarrow  \\ \quad \quad (\operatorname{dist}(\operatorname{completion}(m))(x, y) \simeq (\iota\, l.\;((l \in \mathbb{R}) \wedge \exists f, g \in \operatorname{cseq}(m).\; ((x = \operatorname{class}(\operatorname{cauchy\text{-}setoid}(m), f)) \wedge ((y = \operatorname{class}(\operatorname{cauchy\text{-}setoid}(m), g)) \wedge \operatorname{converges\text{-}to}(\operatorname{rr\text{-}ms}, \operatorname{dist\text{-}seq}(m, f, g), l)))))) \end{array}}}
\vnbprop{r7q-dist-value}{\vnbstmt{\begin{array}{@{}l@{}} \forall m.; \forall f, g \in \operatorname{cseq}(m), v. \\ \quad (\operatorname{is\text{-}metric\text{-}space}(m) \wedge \operatorname{converges\text{-}to}(\operatorname{rr\text{-}ms}, \operatorname{dist\text{-}seq}(m, f, g), v)) \Rightarrow  \\ \quad \quad \operatorname{dist}(\operatorname{completion}(m))(\operatorname{class}(\operatorname{cauchy\text{-}setoid}(m), f), \operatorname{class}(\operatorname{cauchy\text{-}setoid}(m), g)) = v \end{array}}}
\vnbprop{r7q-rr-pts}{\vnbstmt{(\mathbb{R} \simeq \operatorname{pts}(\operatorname{rr\text{-}ms}))}}
\vnbprop{r7q-dhat-in-rr}{\vnbstmt{\begin{array}{@{}l@{}} \forall m.; \forall x, y \in \operatorname{quotient}(\operatorname{cauchy\text{-}setoid}(m)). \\ \quad \operatorname{is\text{-}metric\text{-}space}(m) \Rightarrow (\operatorname{dist}(\operatorname{completion}(m))(x, y) \in \mathbb{R}) \end{array}}}
\vnbprop{r7q-dhat-self-zero}{\vnbstmt{\begin{array}{@{}l@{}} \forall m.; \forall x \in \operatorname{quotient}(\operatorname{cauchy\text{-}setoid}(m)). \\ \quad \operatorname{is\text{-}metric\text{-}space}(m) \Rightarrow (\operatorname{dist}(\operatorname{completion}(m))(x, x) = 0) \end{array}}}
\vnbprop{r7q-dhat-nonneg}{\vnbstmt{\begin{array}{@{}l@{}} \forall m.; \forall x, y \in \operatorname{quotient}(\operatorname{cauchy\text{-}setoid}(m)). \\ \quad \operatorname{is\text{-}metric\text{-}space}(m) \Rightarrow (0 \leq \operatorname{dist}(\operatorname{completion}(m))(x, y)) \end{array}}}
\vnbprop{r7q-dhat-zero-eq}{\vnbstmt{\begin{array}{@{}l@{}} \forall m.; \forall x, y \in \operatorname{quotient}(\operatorname{cauchy\text{-}setoid}(m)). \\ \quad (\operatorname{is\text{-}metric\text{-}space}(m) \wedge (\operatorname{dist}(\operatorname{completion}(m))(x, y) = 0)) \Rightarrow (x = y) \end{array}}}
\vnbprop{r7q-dhat-sym}{\vnbstmt{\begin{array}{@{}l@{}} \forall m.; \forall x, y \in \operatorname{quotient}(\operatorname{cauchy\text{-}setoid}(m)). \\ \quad \operatorname{is\text{-}metric\text{-}space}(m) \Rightarrow  \\ \quad \quad \operatorname{dist}(\operatorname{completion}(m))(x, y) = \operatorname{dist}(\operatorname{completion}(m))(y, x) \end{array}}}
\vnbprop{r7q-dhat-triangle}{\vnbstmt{\begin{array}{@{}l@{}} \forall m.; \forall x, y, z \in \operatorname{quotient}(\operatorname{cauchy\text{-}setoid}(m)). \\ \quad \operatorname{is\text{-}metric\text{-}space}(m) \Rightarrow  \\ \quad \quad \operatorname{dist}(\operatorname{completion}(m))(x, z) \leq (\operatorname{dist}(\operatorname{completion}(m))(x, y) + \operatorname{dist}(\operatorname{completion}(m))(y, z)) \end{array}}}
\vnbprop{completion-is-metric-space}{\vnbstmt{\forall m.\; (\operatorname{is\text{-}metric\text{-}space}(m) \Rightarrow \operatorname{is\text{-}metric\text{-}space}(\operatorname{completion}(m)))}}

\vnbbookfile{theorem-library/rake-completion-complete}{theorem-library/rake-completion-complete}
\noindent{\small rake-completion-complete.scm -- BATCH 12-D (2026-09-20): `completion-is-complete' (structure-library/metric-completion.scm:110, asserted `well-known') and the general bricks its diagonal argument needs.}\par
\vnbprop{rr-limit-tail-le}{\vnbstmt{\begin{array}{@{}l@{}} \forall f, v, c. \\ \quad ((f \in (\mathbb{N} \to \mathbb{R})) \wedge ((v \in \mathbb{R}) \wedge ((c \in \mathbb{R}) \wedge (\operatorname{converges\text{-}to}(\operatorname{rr\text{-}ms}, f, v) \wedge \operatorname{pos\text{-}rr}(c))))) \Rightarrow  \\ \quad \quad \exists t \in \mathbb{N}.\; \forall k \in \mathbb{N}.\; ((t \leq k) \Rightarrow (f(k) \leq (v + c))) \end{array}}}
\vnbprop{rr-limit-le-of-tail}{\vnbstmt{\begin{array}{@{}l@{}} \forall f, v, c, t. \\ \quad f \in (\mathbb{N} \to \mathbb{R}) \wedge \\ \quad v \in \mathbb{R} \wedge \\ \quad c \in \mathbb{R} \wedge \\ \quad t \in \mathbb{N} \wedge \\ \quad \operatorname{converges\text{-}to}(\operatorname{rr\text{-}ms}, f, v) \wedge \\ \quad \forall k \in \mathbb{N}.\; ((t \leq k) \Rightarrow (f(k) \leq c)) \Rightarrow \\ \quad \quad v \leq c \end{array}}}
\vnbprop{r12d-repfam}{\vnbstmt{\begin{array}{@{}l@{}} \forall m.; \forall f \in (\mathbb{N} \to \operatorname{quotient}(\operatorname{cauchy\text{-}setoid}(m))). \\ \quad \operatorname{is\text{-}metric\text{-}space}(m) \Rightarrow  \\ \quad \quad \exists a \in (\mathbb{N} \to \operatorname{cseq}(m)). \forall p \in \mathbb{N}. \\ \quad \quad \quad f(p) = \operatorname{class}(\operatorname{cauchy\text{-}setoid}(m), a(p)) \end{array}}}
\vnbprop{r12d-rapid-indices}{\vnbstmt{\begin{array}{@{}l@{}} \forall m, f, d. \\ \quad \operatorname{is\text{-}metric\text{-}space}(m) \wedge \\ \quad f \in (\mathbb{N} \to \operatorname{cseq}(m)) \wedge \\ \quad d \in (\mathbb{N} \to \mathbb{R}) \wedge \\ \quad \forall p \in \mathbb{N}.\; \operatorname{pos\text{-}rr}(d(p)) \Rightarrow \\ \quad \quad \exists a \in (\mathbb{N} \to \mathbb{N}). \forall p \in \mathbb{N}, i, j. \\ \quad \quad \quad ((i \in \mathbb{N}) \wedge ((j \in \mathbb{N}) \wedge ((a(p) \leq i) \wedge (a(p) \leq j)))) \Rightarrow  \\ \quad \quad \quad \quad \operatorname{dist}(m)(f(p)(i), f(p)(j)) \leq d(p) \end{array}}}
\vnbprop{r12d-key}{\vnbstmt{\begin{array}{@{}l@{}} \forall m, f, a, p, q, b, c. \\ \quad \operatorname{is\text{-}metric\text{-}space}(m) \wedge \\ \quad f \in (\mathbb{N} \to \operatorname{quotient}(\operatorname{cauchy\text{-}setoid}(m))) \wedge \\ \quad a \in (\mathbb{N} \to \operatorname{cseq}(m)) \wedge \\ \quad \forall p \in \mathbb{N}.\; (f(p) = \operatorname{class}(\operatorname{cauchy\text{-}setoid}(m), a(p))) \wedge \\ \quad p \in \mathbb{N} \wedge \\ \quad q \in \mathbb{N} \wedge \\ \quad b \in \mathbb{N} \wedge \\ \quad c \in \mathbb{R} \wedge \\ \quad \operatorname{pos\text{-}rr}(c) \Rightarrow \\ \quad \quad \exists n \in \mathbb{N}. \\ \quad \quad \quad ((b \leq n) \wedge (\operatorname{dist}(m)(a(p)(n), a(q)(n)) \leq (\operatorname{dist}(\operatorname{completion}(m))(f(p), f(q)) + c))) \end{array}}}
\vnbprop{r12d-diag-at}{\vnbstmt{\forall f, m.; \forall p \in \mathbb{N}.\; ((\lambda v \in \mathbb{N}.\; f(v)(m(v)))(p) \simeq f(p)(m(p)))}}
\vnbprop{r12d-diag-in-fun}{\vnbstmt{\begin{array}{@{}l@{}} \forall m, f, a. \\ \quad (\operatorname{is\text{-}metric\text{-}space}(m) \wedge ((f \in (\mathbb{N} \to \operatorname{cseq}(m))) \wedge (a \in (\mathbb{N} \to \mathbb{N})))) \Rightarrow  \\ \quad \quad (\lambda v \in \mathbb{N}.\; f(v)(a(v))) \in (\mathbb{N} \to \operatorname{pts}(m)) \end{array}}}
\vnbprop{r12d-rad-tail}{\vnbstmt{\begin{array}{@{}l@{}} \forall d, v. \\ \quad (\operatorname{null\text{-}rr\text{-}seq}(d) \wedge \operatorname{pos\text{-}rr}(v)) \Rightarrow  \\ \quad \quad \exists r \in \mathbb{N}.\; \forall k \in \mathbb{N}.\; ((r \leq k) \Rightarrow (d(k) \leq v)) \end{array}}}
\vnbprop{r12d-diag-est}{\vnbstmt{\begin{array}{@{}l@{}} \forall m, f, a, d, b, p, q, c. \\ \quad \operatorname{is\text{-}metric\text{-}space}(m) \wedge \\ \quad f \in (\mathbb{N} \to \operatorname{quotient}(\operatorname{cauchy\text{-}setoid}(m))) \wedge \\ \quad a \in (\mathbb{N} \to \operatorname{cseq}(m)) \wedge \\ \quad d \in (\mathbb{N} \to \mathbb{R}) \wedge \\ \quad b \in (\mathbb{N} \to \mathbb{N}) \wedge \\ \quad \forall p \in \mathbb{N}.\; (f(p) = \operatorname{class}(\operatorname{cauchy\text{-}setoid}(m), a(p))) \wedge \\ \quad \forall p \in \mathbb{N}, i, j. \\ \quad \quad ((i \in \mathbb{N}) \wedge ((j \in \mathbb{N}) \wedge ((b(p) \leq i) \wedge (b(p) \leq j)))) \Rightarrow  \\ \quad \quad \quad \operatorname{dist}(m)(a(p)(i), a(p)(j)) \leq d(p) \wedge \\ \quad p \in \mathbb{N} \wedge \\ \quad q \in \mathbb{N} \wedge \\ \quad c \in \mathbb{R} \wedge \\ \quad \operatorname{pos\text{-}rr}(c) \Rightarrow \\ \quad \quad \operatorname{dist}(m)(a(p)(b(p)), a(q)(b(q))) \leq ((d(p) + d(q)) + (\operatorname{dist}(\operatorname{completion}(m))(f(p), f(q)) + c)) \end{array}}}
\vnbprop{r12d-cross-est}{\vnbstmt{\begin{array}{@{}l@{}} \forall m, f, a, d, b, p, q, c. \\ \quad \operatorname{is\text{-}metric\text{-}space}(m) \wedge \\ \quad f \in (\mathbb{N} \to \operatorname{quotient}(\operatorname{cauchy\text{-}setoid}(m))) \wedge \\ \quad a \in (\mathbb{N} \to \operatorname{cseq}(m)) \wedge \\ \quad d \in (\mathbb{N} \to \mathbb{R}) \wedge \\ \quad b \in (\mathbb{N} \to \mathbb{N}) \wedge \\ \quad \forall p \in \mathbb{N}.\; (f(p) = \operatorname{class}(\operatorname{cauchy\text{-}setoid}(m), a(p))) \wedge \\ \quad \forall p \in \mathbb{N}, i, j. \\ \quad \quad ((i \in \mathbb{N}) \wedge ((j \in \mathbb{N}) \wedge ((b(p) \leq i) \wedge (b(p) \leq j)))) \Rightarrow  \\ \quad \quad \quad \operatorname{dist}(m)(a(p)(i), a(p)(j)) \leq d(p) \wedge \\ \quad p \in \mathbb{N} \wedge \\ \quad q \in \mathbb{N} \wedge \\ \quad c \in \mathbb{R} \wedge \\ \quad \operatorname{pos\text{-}rr}(c) \wedge \\ \quad b(p) \leq q \Rightarrow \\ \quad \quad \operatorname{dist}(m)(a(q)(b(q)), a(p)(q)) \leq ((d(q) + d(p)) + (\operatorname{dist}(\operatorname{completion}(m))(f(q), f(p)) + c)) \end{array}}}
\vnbprop{r12d-diag-cauchy}{\vnbstmt{\begin{array}{@{}l@{}} \forall m, f, a, d, b. \\ \quad \operatorname{is\text{-}metric\text{-}space}(m) \wedge \\ \quad f \in (\mathbb{N} \to \operatorname{quotient}(\operatorname{cauchy\text{-}setoid}(m))) \wedge \\ \quad a \in (\mathbb{N} \to \operatorname{cseq}(m)) \wedge \\ \quad \operatorname{null\text{-}rr\text{-}seq}(d) \wedge \\ \quad b \in (\mathbb{N} \to \mathbb{N}) \wedge \\ \quad \forall p \in \mathbb{N}.\; (f(p) = \operatorname{class}(\operatorname{cauchy\text{-}setoid}(m), a(p))) \wedge \\ \quad \forall p \in \mathbb{N}, i, j. \\ \quad \quad ((i \in \mathbb{N}) \wedge ((j \in \mathbb{N}) \wedge ((b(p) \leq i) \wedge (b(p) \leq j)))) \Rightarrow  \\ \quad \quad \quad \operatorname{dist}(m)(a(p)(i), a(p)(j)) \leq d(p) \wedge \\ \quad \operatorname{is\text{-}cauchy\text{-}seq}(\operatorname{completion}(m), f) \Rightarrow \\ \quad \quad \operatorname{is\text{-}cauchy\text{-}seq}(m, (\lambda v \in \mathbb{N}.\; a(v)(b(v)))) \end{array}}}
\vnbprop{r12d-diag-in-cseq}{\vnbstmt{\begin{array}{@{}l@{}} \forall m, f, a, d, b. \\ \quad \operatorname{is\text{-}metric\text{-}space}(m) \wedge \\ \quad f \in (\mathbb{N} \to \operatorname{quotient}(\operatorname{cauchy\text{-}setoid}(m))) \wedge \\ \quad a \in (\mathbb{N} \to \operatorname{cseq}(m)) \wedge \\ \quad \operatorname{null\text{-}rr\text{-}seq}(d) \wedge \\ \quad b \in (\mathbb{N} \to \mathbb{N}) \wedge \\ \quad \forall p \in \mathbb{N}.\; (f(p) = \operatorname{class}(\operatorname{cauchy\text{-}setoid}(m), a(p))) \wedge \\ \quad \forall p \in \mathbb{N}, i, j. \\ \quad \quad ((i \in \mathbb{N}) \wedge ((j \in \mathbb{N}) \wedge ((b(p) \leq i) \wedge (b(p) \leq j)))) \Rightarrow  \\ \quad \quad \quad \operatorname{dist}(m)(a(p)(i), a(p)(j)) \leq d(p) \wedge \\ \quad \operatorname{is\text{-}cauchy\text{-}seq}(\operatorname{completion}(m), f) \Rightarrow \\ \quad \quad (\lambda v \in \mathbb{N}.\; a(v)(b(v))) \in \operatorname{cseq}(m) \end{array}}}
\vnbprop{r12d-diag-limit}{\vnbstmt{\begin{array}{@{}l@{}} \forall m, f, a, d, b. \\ \quad \operatorname{is\text{-}metric\text{-}space}(m) \wedge \\ \quad f \in (\mathbb{N} \to \operatorname{quotient}(\operatorname{cauchy\text{-}setoid}(m))) \wedge \\ \quad a \in (\mathbb{N} \to \operatorname{cseq}(m)) \wedge \\ \quad \operatorname{null\text{-}rr\text{-}seq}(d) \wedge \\ \quad b \in (\mathbb{N} \to \mathbb{N}) \wedge \\ \quad \forall p \in \mathbb{N}.\; (f(p) = \operatorname{class}(\operatorname{cauchy\text{-}setoid}(m), a(p))) \wedge \\ \quad \forall p \in \mathbb{N}, i, j. \\ \quad \quad ((i \in \mathbb{N}) \wedge ((j \in \mathbb{N}) \wedge ((b(p) \leq i) \wedge (b(p) \leq j)))) \Rightarrow  \\ \quad \quad \quad \operatorname{dist}(m)(a(p)(i), a(p)(j)) \leq d(p) \wedge \\ \quad \operatorname{is\text{-}cauchy\text{-}seq}(\operatorname{completion}(m), f) \Rightarrow \\ \quad \quad \operatorname{converges\text{-}to}(\operatorname{completion}(m), f, \operatorname{class}(\operatorname{cauchy\text{-}setoid}(m), (\lambda v \in \mathbb{N}.\; a(v)(b(v))))) \end{array}}}
\vnbprop{completion-is-complete}{\vnbstmt{\forall m.\; (\operatorname{is\text{-}metric\text{-}space}(m) \Rightarrow \operatorname{is\text{-}complete}(\operatorname{completion}(m)))}}

\vnbbookfile{theorem-library/product-weights}{theorem-library/product-weights}
\noindent{\small product-weights.scm -- ANY TWO SUMMABLE WEIGHT SEQUENCES GIVE THE SAME TOPOLOGY ON THE COUNTABLE PRODUCT, proven -- retiring the asserted `product-weights-equivalent' of structure-library/product-metric.scm, the last statement of that file about the product's topology.}\par
\vnbprop{product-identity-continuous}{\vnbstmt{\begin{array}{@{}l@{}} \forall s, w, a. \\ \quad (\operatorname{is\text{-}ms\text{-}sequence}(s) \wedge (\operatorname{summable\text{-}weight}(w) \wedge \operatorname{summable\text{-}weight}(a))) \Rightarrow  \\ \quad \quad \operatorname{is\text{-}continuous}(\operatorname{product\text{-}metric\text{-}w}(s, w), \operatorname{product\text{-}metric\text{-}w}(s, a), \\ \quad \quad \quad (\lambda x \in \operatorname{pts}(\operatorname{product\text{-}metric\text{-}w}(s, w)).\; x)) \end{array}}}
\vnbprop{product-weights-equivalent}{\vnbstmt{\begin{array}{@{}l@{}} \forall s, w, a. \\ \quad (\operatorname{is\text{-}ms\text{-}sequence}(s) \wedge (\operatorname{summable\text{-}weight}(w) \wedge \operatorname{summable\text{-}weight}(a))) \Rightarrow  \\ \quad \quad \operatorname{is\text{-}continuous}(\operatorname{product\text{-}metric\text{-}w}(s, w), \operatorname{product\text{-}metric\text{-}w}(s, a), \\ \quad \quad \quad (\lambda x \in \operatorname{pts}(\operatorname{product\text{-}metric\text{-}w}(s, w)).\; x)) \wedge \\ \quad \quad \operatorname{is\text{-}continuous}(\operatorname{product\text{-}metric\text{-}w}(s, a), \operatorname{product\text{-}metric\text{-}w}(s, w), \\ \quad \quad \quad (\lambda x \in \operatorname{pts}(\operatorname{product\text{-}metric\text{-}w}(s, a)).\; x)) \end{array}}}

\vnbbookfile{theorem-library/c-metric-summable}{theorem-library/c-metric-summable}
\noindent{\small c-metric-summable.scm -- THE SERIES DEFINING THE CANONICAL METRIC OF A COUNTABLY-METRISED SPACE CONVERGES, proven, together with the plumbing every later argument about C-METRIC needs.}\par
\vnbprop{c-metric-pseudometric-at}{\vnbstmt{\begin{array}{@{}l@{}} \forall s.; \forall k \in \mathbb{N}. \\ \quad \operatorname{is\text{-}c\text{-}metric\text{-}space}(s) \Rightarrow \operatorname{is\text{-}pseudometric}(\operatorname{dists}(s)(k), \operatorname{pts}(s)) \end{array}}}
\vnbprop{c-metric-dist-real}{\vnbstmt{\begin{array}{@{}l@{}} \forall s.; \forall k \in \mathbb{N}, u, v \in \operatorname{pts}(s). \\ \quad \operatorname{is\text{-}c\text{-}metric\text{-}space}(s) \Rightarrow (\operatorname{dists}(s)(k)(u, v) \in \mathbb{R}) \end{array}}}
\vnbprop{c-metric-term-bound}{\vnbstmt{\begin{array}{@{}l@{}} \forall s, k, u, v, a. \\ \quad (\operatorname{is\text{-}c\text{-}metric\text{-}space}(s) \wedge ((k \in \mathbb{N}) \wedge ((u \in \operatorname{pts}(s)) \wedge ((v \in \operatorname{pts}(s)) \wedge ((a \in \mathbb{R}) \wedge (0 \leq a)))))) \Rightarrow  \\ \quad \quad (a \cdot \operatorname{min}(1, \operatorname{dists}(s)(k)(u, v))) \in \mathbb{R} \wedge \\ \quad \quad 0 \leq (a \cdot \operatorname{min}(1, \operatorname{dists}(s)(k)(u, v))) \wedge \\ \quad \quad (a \cdot \operatorname{min}(1, \operatorname{dists}(s)(k)(u, v))) \leq a \end{array}}}
\vnbprop{c-metric-term-seq-in-fun}{\vnbstmt{\begin{array}{@{}l@{}} \forall s, w, u, v. \\ \quad (\operatorname{is\text{-}c\text{-}metric\text{-}space}(s) \wedge (\operatorname{summable\text{-}weight}(w) \wedge ((u \in \operatorname{pts}(s)) \wedge (v \in \operatorname{pts}(s))))) \Rightarrow  \\ \quad \quad (\lambda k \in \mathbb{N}.\; (w(k) \cdot \operatorname{min}(1, \operatorname{dists}(s)(k)(u, v)))) \in (\mathbb{N} \to \mathbb{R}) \end{array}}}
\vnbprop{c-metric-summable}{\vnbstmt{\begin{array}{@{}l@{}} \forall s, w, u, v. \\ \quad (\operatorname{is\text{-}c\text{-}metric\text{-}space}(s) \wedge (\operatorname{summable\text{-}weight}(w) \wedge ((u \in \operatorname{pts}(s)) \wedge (v \in \operatorname{pts}(s))))) \Rightarrow  \\ \quad \quad \operatorname{series\text{-}converges}((\lambda k \in \mathbb{N}.\; (w(k) \cdot \operatorname{min}(1, \operatorname{dists}(s)(k)(u, v))))) \end{array}}}
\vnbprop{c-metric-carrier}{\vnbstmt{\forall s, w.\; (\operatorname{is\text{-}c\text{-}metric\text{-}space}(s) \Rightarrow (\operatorname{pts}(\operatorname{c\text{-}metric\text{-}w}(s, w)) \simeq \operatorname{pts}(s)))}}
\vnbprop{c-metric-distance}{\vnbstmt{\begin{array}{@{}l@{}} \forall s, w, u, v. \\ \quad ((u \in \operatorname{pts}(s)) \wedge (v \in \operatorname{pts}(s))) \Rightarrow  \\ \quad \quad (\operatorname{dist}(\operatorname{c\text{-}metric\text{-}w}(s, w))(u, v) \simeq \operatorname{series\text{-}limit}((\lambda k \in \mathbb{N}.\; (w(k) \cdot \operatorname{min}(1, \operatorname{dists}(s)(k)(u, v)))))) \end{array}}}
\vnbprop{c-metric-dist-converges}{\vnbstmt{\begin{array}{@{}l@{}} \forall s, w, u, v. \\ \quad (\operatorname{is\text{-}c\text{-}metric\text{-}space}(s) \wedge (\operatorname{summable\text{-}weight}(w) \wedge ((u \in \operatorname{pts}(s)) \wedge (v \in \operatorname{pts}(s))))) \Rightarrow  \\ \quad \quad \operatorname{series\text{-}converges\text{-}to}((\lambda k \in \mathbb{N}.\; (w(k) \cdot \operatorname{min}(1, \operatorname{dists}(s)(k)(u, v)))), \\ \quad \quad \quad \operatorname{dist}(\operatorname{c\text{-}metric\text{-}w}(s, w))(u, v)) \end{array}}}
\vnbprop{c-metric-canonical-summable}{\vnbstmt{\begin{array}{@{}l@{}} \forall s, u, v. \\ \quad (\operatorname{is\text{-}c\text{-}metric\text{-}space}(s) \wedge ((u \in \operatorname{pts}(s)) \wedge (v \in \operatorname{pts}(s)))) \Rightarrow  \\ \quad \quad \operatorname{series\text{-}converges}((\lambda k \in \mathbb{N}.\; ((\lambda n \in \mathbb{N}.\; /(1, \operatorname{power}(2, (n + 1))))(k) \cdot \operatorname{min}(1, \operatorname{dists}(s)(k)(u, v))))) \end{array}}}

\vnbbookfile{theorem-library/c-metric-is-metric}{theorem-library/c-metric-is-metric}
\noindent{\small c-metric-is-metric.scm -- THE CANONICAL METRIC OF A COUNTABLY-METRISED SPACE IS A METRIC SPACE, proven.}\par
\vnbprop{cmim-dist-in-rr}{\vnbstmt{\begin{array}{@{}l@{}} \forall s, w.; \forall x, y \in \operatorname{pts}(s). \\ \quad (\operatorname{is\text{-}c\text{-}metric\text{-}space}(s) \wedge \operatorname{summable\text{-}weight}(w)) \Rightarrow  \\ \quad \quad \operatorname{dist}(\operatorname{c\text{-}metric\text{-}w}(s, w))(x, y) \in \mathbb{R} \end{array}}}
\vnbprop{cmim-term-nonneg}{\vnbstmt{\begin{array}{@{}l@{}} \forall s, w.; \forall x, y \in \operatorname{pts}(s), n \in \mathbb{N}. \\ \quad (\operatorname{is\text{-}c\text{-}metric\text{-}space}(s) \wedge \operatorname{summable\text{-}weight}(w)) \Rightarrow  \\ \quad \quad 0 \leq (\lambda k \in \mathbb{N}.\; (w(k) \cdot \operatorname{min}(1, \operatorname{dists}(s)(k)(x, y))))(n) \end{array}}}
\vnbprop{cmim-dist-nonneg}{\vnbstmt{\begin{array}{@{}l@{}} \forall s, w.; \forall x, y \in \operatorname{pts}(s). \\ \quad (\operatorname{is\text{-}c\text{-}metric\text{-}space}(s) \wedge \operatorname{summable\text{-}weight}(w)) \Rightarrow  \\ \quad \quad 0 \leq \operatorname{dist}(\operatorname{c\text{-}metric\text{-}w}(s, w))(x, y) \end{array}}}
\vnbprop{cmim-term-self-zero}{\vnbstmt{\begin{array}{@{}l@{}} \forall s, w.; \forall x \in \operatorname{pts}(s), n \in \mathbb{N}. \\ \quad (\operatorname{is\text{-}c\text{-}metric\text{-}space}(s) \wedge \operatorname{summable\text{-}weight}(w)) \Rightarrow  \\ \quad \quad (\lambda k \in \mathbb{N}.\; (w(k) \cdot \operatorname{min}(1, \operatorname{dists}(s)(k)(x, x))))(n) = 0 \end{array}}}
\vnbprop{cmim-dist-self-zero}{\vnbstmt{\begin{array}{@{}l@{}} \forall s, w.; \forall x \in \operatorname{pts}(s). \\ \quad (\operatorname{is\text{-}c\text{-}metric\text{-}space}(s) \wedge \operatorname{summable\text{-}weight}(w)) \Rightarrow  \\ \quad \quad \operatorname{dist}(\operatorname{c\text{-}metric\text{-}w}(s, w))(x, x) = 0 \end{array}}}
\vnbprop{cmim-term-zero}{\vnbstmt{\begin{array}{@{}l@{}} \forall s, w.; \forall x, y \in \operatorname{pts}(s), n \in \mathbb{N}. \\ \quad \operatorname{is\text{-}c\text{-}metric\text{-}space}(s) \wedge \\ \quad \operatorname{summable\text{-}weight}(w) \wedge \\ \quad (\lambda k \in \mathbb{N}.\; (w(k) \cdot \operatorname{min}(1, \operatorname{dists}(s)(k)(x, y))))(n) = 0 \Rightarrow \\ \quad \quad \operatorname{dists}(s)(n)(x, y) = 0 \end{array}}}
\vnbprop{cmim-dist-zero-eq}{\vnbstmt{\begin{array}{@{}l@{}} \forall s, w.; \forall x, y \in \operatorname{pts}(s). \\ \quad (\operatorname{is\text{-}c\text{-}metric\text{-}space}(s) \wedge (\operatorname{summable\text{-}weight}(w) \wedge (\operatorname{dist}(\operatorname{c\text{-}metric\text{-}w}(s, w))(x, y) = 0))) \Rightarrow  \\ \quad \quad x = y \end{array}}}
\vnbprop{cmim-term-sym}{\vnbstmt{\begin{array}{@{}l@{}} \forall s, w.; \forall x, y \in \operatorname{pts}(s), n \in \mathbb{N}. \\ \quad (\operatorname{is\text{-}c\text{-}metric\text{-}space}(s) \wedge \operatorname{summable\text{-}weight}(w)) \Rightarrow  \\ \quad \quad (\lambda k \in \mathbb{N}.\; (w(k) \cdot \operatorname{min}(1, \operatorname{dists}(s)(k)(x, y))))(n) \\ \quad \quad \leq\; (\lambda k \in \mathbb{N}.\; (w(k) \cdot \operatorname{min}(1, \operatorname{dists}(s)(k)(y, x))))(n) \end{array}}}
\vnbprop{cmim-dist-sym}{\vnbstmt{\begin{array}{@{}l@{}} \forall s, w.; \forall x, y \in \operatorname{pts}(s). \\ \quad (\operatorname{is\text{-}c\text{-}metric\text{-}space}(s) \wedge \operatorname{summable\text{-}weight}(w)) \Rightarrow  \\ \quad \quad \operatorname{dist}(\operatorname{c\text{-}metric\text{-}w}(s, w))(x, y) = \operatorname{dist}(\operatorname{c\text{-}metric\text{-}w}(s, w))(y, x) \end{array}}}
\vnbprop{cmim-sum-seq-in-fun}{\vnbstmt{\begin{array}{@{}l@{}} \forall s, w.; \forall x, y, z \in \operatorname{pts}(s). \\ \quad (\operatorname{is\text{-}c\text{-}metric\text{-}space}(s) \wedge \operatorname{summable\text{-}weight}(w)) \Rightarrow  \\ \quad \quad (\lambda k \in \mathbb{N}.\; ((w(k) \cdot \operatorname{min}(1, \operatorname{dists}(s)(k)(x, y))) + (w(k) \cdot \operatorname{min}(1, \operatorname{dists}(s)(k)(y, z))))) \\ \quad \quad \in\; (\mathbb{N} \to \mathbb{R}) \end{array}}}
\vnbprop{cmim-sum-seq-split}{\vnbstmt{\begin{array}{@{}l@{}} \forall s, w.; \forall x, y, z \in \operatorname{pts}(s), i \in \mathbb{N}. \\ \quad (\operatorname{is\text{-}c\text{-}metric\text{-}space}(s) \wedge \operatorname{summable\text{-}weight}(w)) \Rightarrow  \\ \quad \quad (\lambda k \in \mathbb{N}.\; ((w(k) \cdot \operatorname{min}(1, \operatorname{dists}(s)(k)(x, y))) + (w(k) \cdot \operatorname{min}(1, \operatorname{dists}(s)(k)(y, z)))))(i) \\ \quad \quad =\; (\lambda k \in \mathbb{N}.\; (w(k) \cdot \operatorname{min}(1, \operatorname{dists}(s)(k)(x, y))))(i) + (\lambda k \in \mathbb{N}.\; (w(k) \cdot \operatorname{min}(1, \operatorname{dists}(s)(k)(y, z))))(i) \end{array}}}
\vnbprop{cmim-term-triangle}{\vnbstmt{\begin{array}{@{}l@{}} \forall s, w.; \forall x, y, z \in \operatorname{pts}(s), n \in \mathbb{N}. \\ \quad (\operatorname{is\text{-}c\text{-}metric\text{-}space}(s) \wedge \operatorname{summable\text{-}weight}(w)) \Rightarrow  \\ \quad \quad (\lambda k \in \mathbb{N}.\; (w(k) \cdot \operatorname{min}(1, \operatorname{dists}(s)(k)(x, z))))(n) \\ \quad \quad \leq\; (\lambda k \in \mathbb{N}.\; ((w(k) \cdot \operatorname{min}(1, \operatorname{dists}(s)(k)(x, y))) + (w(k) \cdot \operatorname{min}(1, \operatorname{dists}(s)(k)(y, z)))))(n) \end{array}}}
\vnbprop{cmim-dist-triangle}{\vnbstmt{\begin{array}{@{}l@{}} \forall s, w.; \forall x, y, z \in \operatorname{pts}(s). \\ \quad (\operatorname{is\text{-}c\text{-}metric\text{-}space}(s) \wedge \operatorname{summable\text{-}weight}(w)) \Rightarrow  \\ \quad \quad \operatorname{dist}(\operatorname{c\text{-}metric\text{-}w}(s, w))(x, z) \leq (\operatorname{dist}(\operatorname{c\text{-}metric\text{-}w}(s, w))(x, y) + \operatorname{dist}(\operatorname{c\text{-}metric\text{-}w}(s, w))(y, z)) \end{array}}}
\vnbprop{c-metric-w-is-metric-space}{\vnbstmt{\begin{array}{@{}l@{}} \forall s, w. \\ \quad (\operatorname{is\text{-}c\text{-}metric\text{-}space}(s) \wedge \operatorname{summable\text{-}weight}(w)) \Rightarrow  \\ \quad \quad \operatorname{is\text{-}metric\text{-}space}(\operatorname{c\text{-}metric\text{-}w}(s, w)) \end{array}}}
\vnbprop{c-metric-is-metric-space}{\vnbstmt{\forall s.\; (\operatorname{is\text{-}c\text{-}metric\text{-}space}(s) \Rightarrow \operatorname{is\text{-}metric\text{-}space}(\operatorname{c\text{-}metric}(s)))}}

\vnbbookfile{theorem-library/ccint-creep}{theorem-library/ccint-creep}
\noindent{\small ccint-creep.scm -- the CREEPING PRINCIPLE on a closed interval, PROVEN.}\par
\vnbprop{ccint-creep}{\vnbstmt{\begin{array}{@{}l@{}} \forall a, b, g. \\ \quad a \in \mathbb{R} \wedge \\ \quad b \in \mathbb{R} \wedge \\ \quad a \leq b \wedge \\ \quad g \subseteq \mathbb{R} \wedge \\ \quad a \in g \wedge \\ \quad \operatorname{rr\text{-}upper\text{-}bound}(g, b) \wedge \\ \quad \forall t \in \operatorname{ccint}(a, b). \exists d \in \mathbb{R}. \\ \quad \quad ((0 < d) \wedge \forall y \in \operatorname{ccint}(a, b).\; ((\exists w \in g.\; ((t - d) < w) \wedge (y \leq (t + d))) \Rightarrow (y \in g))) \Rightarrow \\ \quad \quad b \in g \end{array}}}

\vnbbookfile{theorem-library/ccint-bounded}{theorem-library/ccint-bounded}
\noindent{\small ccint-bounded.scm -- a continuous function on [a,b] is BOUNDED ABOVE, PROVEN.}\par
\vnbprop{continuous-bounded-above-on-ccint}{\vnbstmt{\begin{array}{@{}l@{}} \forall f, a, b. \\ \quad f \in (\mathbb{R} \to \mathbb{R}) \wedge \\ \quad a \in \mathbb{R} \wedge \\ \quad b \in \mathbb{R} \wedge \\ \quad a \leq b \wedge \\ \quad \forall x \in \operatorname{ccint}(a, b).\; \operatorname{is\text{-}continuous\text{-}at}(\operatorname{rr\text{-}ms}, \operatorname{rr\text{-}ms}, f, x) \Rightarrow \\ \quad \quad \exists k \in \mathbb{R}.\; \forall z \in \operatorname{ccint}(a, b).\; (f(z) \leq k) \end{array}}}

\vnbbookfile{theorem-library/uniform-continuity-ccint}{theorem-library/uniform-continuity-ccint}
\noindent{\small uniform-continuity-ccint.scm -- a continuous function on a CLOSED BOUNDED INTERVAL is UNIFORMLY continuous there, PROVEN.}\par
\vnbprop{continuous-uniformly-continuous-on-ccint}{\vnbstmt{\begin{array}{@{}l@{}} \forall f, a, b, s. \\ \quad f \in (\mathbb{R} \to \mathbb{R}) \wedge \\ \quad a \in \mathbb{R} \wedge \\ \quad b \in \mathbb{R} \wedge \\ \quad a \leq b \wedge \\ \quad \operatorname{pos\text{-}rr}(s) \wedge \\ \quad \forall x \in \operatorname{ccint}(a, b).\; \operatorname{is\text{-}continuous\text{-}at}(\operatorname{rr\text{-}ms}, \operatorname{rr\text{-}ms}, f, x) \Rightarrow \\ \quad \quad \exists d \in \mathbb{R}. \\ \quad \quad \quad 0 < d \wedge \\ \quad \quad \quad \forall p, q \in \operatorname{ccint}(a, b). \\ \quad \quad \quad \quad (\lvert (p - q) \rvert \leq d) \Rightarrow (\lvert (f(p) - f(q)) \rvert \leq s) \end{array}}}

\vnbbookfile{theorem-library/evt-proof}{theorem-library/evt-proof}
\noindent{\small evt-proof.scm -- the EXTREME VALUE THEOREM (max form), PROVEN.}\par
\vnbprop{extreme-value-max}{\vnbstmt{\begin{array}{@{}l@{}} \forall f, a, b. \\ \quad f \in (\mathbb{R} \to \mathbb{R}) \wedge \\ \quad a \in \mathbb{R} \wedge \\ \quad b \in \mathbb{R} \wedge \\ \quad a \leq b \wedge \\ \quad \forall x \in \operatorname{ccint}(a, b).\; \operatorname{is\text{-}continuous\text{-}at}(\operatorname{rr\text{-}ms}, \operatorname{rr\text{-}ms}, f, x) \Rightarrow \\ \quad \quad \exists c \in \operatorname{ccint}(a, b).\; \forall x \in \operatorname{ccint}(a, b).\; (f(x) \leq f(c)) \end{array}}}

\vnbbookfile{theorem-library/heine-borel-baby}{theorem-library/heine-borel-baby}
\noindent{\small NOTE 2026-09-20 (batch 12-B).  Where this file says the tree has no metric SUBSPACE structure, that is history: SUBSPACE-MS(s, A) and RESTRICT(f, A) are defined in structure-library/metric-subspace.scm, their laws are in theorem-library/metric-subspace-laws.scm and compact-subspace.scm, and}\par
\vnbprop{empty-subset-any}{\vnbstmt{\forall b.\; (\emptyset \subseteq b)}}
\vnbprop{union-singleton-mem}{\vnbstmt{\forall y \in \mathbf{Set}, s.\; (y \in (\{y\} \cup s))}}
\vnbprop{union-right-subset}{\vnbstmt{\forall a, s.\; (s \subseteq (a \cup s))}}
\vnbprop{union-singleton-subset}{\vnbstmt{\forall b, y, s.\; (((y \in b) \wedge (s \subseteq b)) \Rightarrow ((\{y\} \cup s) \subseteq b))}}
\vnbprop{big-union-mono}{\vnbstmt{\forall a, b.\; ((a \subseteq b) \Rightarrow (\operatorname{big\text{-}union}(u, a, u) \subseteq \operatorname{big\text{-}union}(u, b, u)))}}
\vnbprop{card-union-singleton-nn}{\vnbstmt{\begin{array}{@{}l@{}} \forall s, y \in \mathbf{Set}. \\ \quad (|s| \in \mathbb{N}) \Rightarrow (|(\{y\} \cup s)| \in \mathbb{N}) \end{array}}}
\vnbprop{heine-borel-baby}{\vnbstmt{\begin{array}{@{}l@{}} \forall a, b, f. \\ \quad a \in \mathbb{R} \wedge \\ \quad b \in \mathbb{R} \wedge \\ \quad a \leq b \wedge \\ \quad \forall v \in f.\; \operatorname{is\text{-}open}(\operatorname{rr\text{-}ms}, v) \wedge \\ \quad \operatorname{ccint}(a, b) \subseteq \operatorname{big\text{-}union}(u, f, u) \Rightarrow \\ \quad \quad \exists g. \\ \quad \quad \quad ((g \subseteq f) \wedge ((g \in \mathbf{Set}) \wedge ((|g| \in \mathbb{N}) \wedge (\operatorname{ccint}(a, b) \subseteq \operatorname{big\text{-}union}(u, g, u))))) \end{array}}}

\vnbbookfile{theorem-library/evt-min-proof}{theorem-library/evt-min-proof}
\noindent{\small evt-min-proof.scm -- the EXTREME VALUE THEOREM (min form), PROVEN.}\par
\vnbprop{extreme-value-min}{\vnbstmt{\begin{array}{@{}l@{}} \forall f, a, b. \\ \quad f \in (\mathbb{R} \to \mathbb{R}) \wedge \\ \quad a \in \mathbb{R} \wedge \\ \quad b \in \mathbb{R} \wedge \\ \quad a \leq b \wedge \\ \quad \forall x \in \operatorname{ccint}(a, b).\; \operatorname{is\text{-}continuous\text{-}at}(\operatorname{rr\text{-}ms}, \operatorname{rr\text{-}ms}, f, x) \Rightarrow \\ \quad \quad \exists c \in \operatorname{ccint}(a, b).\; \forall x \in \operatorname{ccint}(a, b).\; (f(c) \leq f(x)) \end{array}}}

\vnbbookfile{theorem-library/ccint-abs-bounded}{theorem-library/ccint-abs-bounded}
\noindent{\small ccint-abs-bounded.scm -- a continuous function on a CLOSED BOUNDED INTERVAL is bounded in ABSOLUTE VALUE there, PROVEN modulo 0.}\par
\vnbprop{continuous-abs-bounded-on-ccint}{\vnbstmt{\begin{array}{@{}l@{}} \forall f, a, b. \\ \quad f \in (\mathbb{R} \to \mathbb{R}) \wedge \\ \quad a \in \mathbb{R} \wedge \\ \quad b \in \mathbb{R} \wedge \\ \quad a \leq b \wedge \\ \quad \forall x \in \operatorname{ccint}(a, b).\; \operatorname{is\text{-}continuous\text{-}at}(\operatorname{rr\text{-}ms}, \operatorname{rr\text{-}ms}, f, x) \Rightarrow \\ \quad \quad \exists m \in \mathbb{R}.\; ((0 \leq m) \wedge \forall z \in \operatorname{ccint}(a, b).\; (\lvert f(z) \rvert \leq m)) \end{array}}}

\vnbbookfile{theorem-library/bernstein-density}{theorem-library/bernstein-density}
\noindent{\small bernstein-density.scm -- THEOREM 5.2 of docs/calculus.pdf: a continuous function on [0,1] is the uniform limit of its Bernstein polynomials.}\par
\par\medskip\noindent\textbf{Definition} (\texttt{bernstein-poly}).\;$\operatorname{bernstein\text{-}poly}(f, n, x) :=$
\begin{equation*}
\operatorname{series\text{-}partial\text{-}sum}((\lambda \operatorname{l\_} \in \mathbb{N}.\; (f(/(\operatorname{l\_}, n)) \cdot \operatorname{bernstein\text{-}basis}(n, x)(\operatorname{l\_}))), (n + 1))
\end{equation*}

\vnbprop{bernstein-basis-nonneg}{\vnbstmt{\begin{array}{@{}l@{}} \forall n \in \mathbb{N}, x, k. \\ \quad ((x \in \mathbb{R}) \wedge ((0 \leq x) \wedge ((x \leq 1) \wedge (k \in \mathbb{N})))) \Rightarrow  \\ \quad \quad 0 \leq \operatorname{bernstein\text{-}basis}(n, x)(k) \end{array}}}
\vnbprop{bernstein-term-bound}{\vnbstmt{\begin{array}{@{}l@{}} \forall f, n, x, h, a, m, c, l. \\ \quad f \in (\mathbb{R} \to \mathbb{R}) \wedge \\ \quad n \in \mathbb{N} \wedge \\ \quad \neg (n = 0) \wedge \\ \quad x \in \operatorname{ccint}(0, 1) \wedge \\ \quad h \in \mathbb{R} \wedge \\ \quad 0 < h \wedge \\ \quad a \in \mathbb{R} \wedge \\ \quad 0 < a \wedge \\ \quad m \in \mathbb{R} \wedge \\ \quad 0 \leq m \wedge \\ \quad c \in \mathbb{R} \wedge \\ \quad 0 \leq c \wedge \\ \quad ((n \cdot n) \cdot ((a \cdot a) \cdot c)) = (m + m) \wedge \\ \quad \forall p, q \in \operatorname{ccint}(0, 1). \\ \quad \quad (\lvert (p - q) \rvert \leq a) \Rightarrow (\lvert (f(p) - f(q)) \rvert \leq h) \wedge \\ \quad \forall z \in \operatorname{ccint}(0, 1).\; (\lvert f(z) \rvert \leq m) \wedge \\ \quad l \in \mathbb{N} \Rightarrow \\ \quad \quad \operatorname{bernstein\text{-}basis}(n, x)(l) \cdot (f(x) - f(/(l, n))) \\ \quad \quad \leq\; \operatorname{bernstein\text{-}basis}(n, x)(l) \cdot (h + (c \cdot ((l - (n \cdot x)) \cdot (l - (n \cdot x))))) \wedge \\ \quad \quad \operatorname{bernstein\text{-}basis}(n, x)(l) \cdot (-1 \cdot (h + (c \cdot ((l - (n \cdot x)) \cdot (l - (n \cdot x)))))) \\ \quad \quad \leq\; \operatorname{bernstein\text{-}basis}(n, x)(l) \cdot (f(x) - f(/(l, n))) \end{array}}}
\vnbprop{bernstein-approx-bound}{\vnbstmt{\begin{array}{@{}l@{}} \forall f, n, x, h, a, m, c. \\ \quad f \in (\mathbb{R} \to \mathbb{R}) \wedge \\ \quad n \in \mathbb{N} \wedge \\ \quad \neg (n = 0) \wedge \\ \quad x \in \operatorname{ccint}(0, 1) \wedge \\ \quad h \in \mathbb{R} \wedge \\ \quad 0 < h \wedge \\ \quad a \in \mathbb{R} \wedge \\ \quad 0 < a \wedge \\ \quad m \in \mathbb{R} \wedge \\ \quad 0 \leq m \wedge \\ \quad c \in \mathbb{R} \wedge \\ \quad 0 \leq c \wedge \\ \quad ((n \cdot n) \cdot ((a \cdot a) \cdot c)) = (m + m) \wedge \\ \quad (c \cdot n) \leq h \wedge \\ \quad \forall p, q \in \operatorname{ccint}(0, 1). \\ \quad \quad (\lvert (p - q) \rvert \leq a) \Rightarrow (\lvert (f(p) - f(q)) \rvert \leq h) \wedge \\ \quad \forall z \in \operatorname{ccint}(0, 1).\; (\lvert f(z) \rvert \leq m) \Rightarrow \\ \quad \quad \lvert (f(x) - \operatorname{bernstein\text{-}poly}(f, n, x)) \rvert \leq (h + h) \end{array}}}
\vnbprop{bernstein-uniform-approximation}{\vnbstmt{\begin{array}{@{}l@{}} \forall f, s. \\ \quad ((f \in (\mathbb{R} \to \mathbb{R})) \wedge (\forall x \in \operatorname{ccint}(0, 1).\; \operatorname{is\text{-}continuous\text{-}at}(\operatorname{rr\text{-}ms}, \operatorname{rr\text{-}ms}, f, x) \wedge \operatorname{pos\text{-}rr}(s))) \Rightarrow  \\ \quad \quad \exists p \in \mathbb{N}. \forall n, a. \\ \quad \quad \quad ((n \in \mathbb{N}) \wedge ((p \leq n) \wedge (a \in \operatorname{ccint}(0, 1)))) \Rightarrow  \\ \quad \quad \quad \quad \lvert (f(a) - \operatorname{bernstein\text{-}poly}(f, n, a)) \rvert \leq s \end{array}}}

\vnbbookfile{theorem-library/mvt-cluster-readoffs}{theorem-library/mvt-cluster-readoffs}
\noindent{\small theorem-library/mvt-cluster-readoffs.scm}\par
\vnbprop{diff-value-real}{\vnbstmt{\forall f, a, l.\; (\operatorname{is\text{-}diff\text{-}at}(f, a, l) \Rightarrow (l \in \mathbb{R}))}}
\vnbprop{deriv-neg}{\vnbstmt{\begin{array}{@{}l@{}} \forall f, a, l. \\ \quad \operatorname{is\text{-}diff\text{-}at}(f, a, l) \Rightarrow \operatorname{is\text{-}diff\text{-}at}((\lambda z \in \mathbb{R}.\; -f(z)), a, -l) \end{array}}}
\vnbprop{max-val-strict-interior}{\vnbstmt{\begin{array}{@{}l@{}} \forall h, a, b.; \forall c \in \operatorname{ccint}(a, b). \\ \quad ((h(a) = h(b)) \wedge (h(a) < h(c))) \Rightarrow ((a < c) \wedge (c < b)) \end{array}}}
\vnbprop{min-val-strict-interior}{\vnbstmt{\begin{array}{@{}l@{}} \forall h, a, b.; \forall c \in \operatorname{ccint}(a, b). \\ \quad ((h(a) = h(b)) \wedge (h(c) < h(a))) \Rightarrow ((a < c) \wedge (c < b)) \end{array}}}

\vnbbookfile{theorem-library/deriv-difference}{theorem-library/deriv-difference}
\noindent{\small deriv-difference.scm -- the DIFFERENCE rule for the Caratheodory derivative IS-DIFF-AT: if f has derivative L at a and g has derivative M at a, then x |-$>$ f(x) - g(x) has derivative L - M at a.  PROVEN, `modulo 0'.}\par
\vnbprop{deriv-difference}{\vnbstmt{\begin{array}{@{}l@{}} \forall f, g, a, l, m. \\ \quad (\operatorname{is\text{-}diff\text{-}at}(f, a, l) \wedge \operatorname{is\text{-}diff\text{-}at}(g, a, m)) \Rightarrow  \\ \quad \quad \operatorname{is\text{-}diff\text{-}at}((\lambda x \in \mathbb{R}.\; (f(x) - g(x))), a, (l - m)) \end{array}}}

\vnbbookfile{theorem-library/mvt-aux-guarded}{theorem-library/mvt-aux-guarded}
\noindent{\small theorem-library/mvt-aux-guarded.scm -- the four MVT auxiliary-function supports, GUARDED and PROVEN.}\par
\vnbprop{mvt-aux-diff}{\vnbstmt{\begin{array}{@{}l@{}} \forall f, a, b, x, l. \\ \quad ((a \in \mathbb{R}) \wedge ((b \in \mathbb{R}) \wedge \operatorname{is\text{-}diff\text{-}at}(f, x, l))) \Rightarrow  \\ \quad \quad \operatorname{is\text{-}diff\text{-}at}((\lambda z \in \mathbb{R}.\; ((f(z) \cdot (b - a)) - (z \cdot (f(b) - f(a))))), x, ((l \cdot (b - a)) - (f(b) - f(a)))) \end{array}}}
\vnbprop{mvt-aux-cont}{\vnbstmt{\begin{array}{@{}l@{}} \forall f, a, b, x. \\ \quad ((a \in \mathbb{R}) \wedge ((b \in \mathbb{R}) \wedge \operatorname{is\text{-}continuous\text{-}at}(\operatorname{rr\text{-}ms}, \operatorname{rr\text{-}ms}, f, x))) \Rightarrow  \\ \quad \quad \operatorname{is\text{-}continuous\text{-}at}(\operatorname{rr\text{-}ms}, \operatorname{rr\text{-}ms}, (\lambda z \in \mathbb{R}.\; ((f(z) \cdot (b - a)) - (z \cdot (f(b) - f(a))))), x) \end{array}}}
\vnbprop{gmvt-aux-diff}{\vnbstmt{\begin{array}{@{}l@{}} \forall f, g, a, b, x, l, m. \\ \quad ((a \in \mathbb{R}) \wedge ((b \in \mathbb{R}) \wedge (\operatorname{is\text{-}diff\text{-}at}(f, x, l) \wedge \operatorname{is\text{-}diff\text{-}at}(g, x, m)))) \Rightarrow  \\ \quad \quad \operatorname{is\text{-}diff\text{-}at}((\lambda z \in \mathbb{R}.\; ((f(z) \cdot (g(b) - g(a))) - (g(z) \cdot (f(b) - f(a))))), \\ \quad \quad \quad x, ((l \cdot (g(b) - g(a))) - (m \cdot (f(b) - f(a))))) \end{array}}}
\vnbprop{gmvt-aux-cont}{\vnbstmt{\begin{array}{@{}l@{}} \forall f, g, a, b, x. \\ \quad a \in \mathbb{R} \wedge \\ \quad b \in \mathbb{R} \wedge \\ \quad \operatorname{is\text{-}continuous\text{-}at}(\operatorname{rr\text{-}ms}, \operatorname{rr\text{-}ms}, f, x) \wedge \\ \quad \operatorname{is\text{-}continuous\text{-}at}(\operatorname{rr\text{-}ms}, \operatorname{rr\text{-}ms}, g, x) \Rightarrow \\ \quad \quad \operatorname{is\text{-}continuous\text{-}at}(\operatorname{rr\text{-}ms}, \operatorname{rr\text{-}ms}, (\lambda z \in \mathbb{R}.\; ((f(z) \cdot (g(b) - g(a))) - (g(z) \cdot (f(b) - f(a))))), x) \end{array}}}

\vnbbookfile{theorem-library/interior-extremum-proof}{theorem-library/interior-extremum-proof}
\noindent{\small theorem-library/interior-extremum-proof.scm ==================================================================== Prop 2.10 (interior max): interior-max-deriv-zero, MACHINE-PROVEN.}\par
\vnbprop{interior-max-deriv-zero}{\vnbstmt{\begin{array}{@{}l@{}} \forall f, a, b, c, l. \\ \quad f \in (\mathbb{R} \to \mathbb{R}) \wedge \\ \quad a \in \mathbb{R} \wedge \\ \quad b \in \mathbb{R} \wedge \\ \quad c \in \mathbb{R} \wedge \\ \quad a < c \wedge \\ \quad c < b \wedge \\ \quad \forall x \in \operatorname{ccint}(a, b).\; (f(x) \leq f(c)) \wedge \\ \quad \operatorname{is\text{-}diff\text{-}at}(f, c, l) \Rightarrow \\ \quad \quad l = 0 \end{array}}}
\vnbprop{interior-min-deriv-zero}{\vnbstmt{\begin{array}{@{}l@{}} \forall f, a, b, c, l. \\ \quad f \in (\mathbb{R} \to \mathbb{R}) \wedge \\ \quad a \in \mathbb{R} \wedge \\ \quad b \in \mathbb{R} \wedge \\ \quad c \in \mathbb{R} \wedge \\ \quad a < c \wedge \\ \quad c < b \wedge \\ \quad \forall x \in \operatorname{ccint}(a, b).\; (f(c) \leq f(x)) \wedge \\ \quad \operatorname{is\text{-}diff\text{-}at}(f, c, l) \Rightarrow \\ \quad \quad l = 0 \end{array}}}

\vnbbookfile{theorem-library/rolle-proof}{theorem-library/rolle-proof}
\noindent{\small theorem-library/rolle-proof.scm -- Rolle's theorem, MACHINE-PROVEN. EVT gives argmax c / argmin d on [a,b]; if either is interior, the proven Fermat theorem (interior-max/min-deriv-zero) gives h'=0; else h is constant and a midpoint (rr-midpoint-between) is an interior max.}\par
\vnbprop{rolle}{\vnbstmt{\begin{array}{@{}l@{}} \forall h, a, b. \\ \quad h \in (\mathbb{R} \to \mathbb{R}) \wedge \\ \quad a \in \mathbb{R} \wedge \\ \quad b \in \mathbb{R} \wedge \\ \quad a < b \wedge \\ \quad \forall x \in \operatorname{ccint}(a, b).\; \operatorname{is\text{-}continuous\text{-}at}(\operatorname{rr\text{-}ms}, \operatorname{rr\text{-}ms}, h, x) \wedge \\ \quad \forall x \in \mathbb{R}.\; (((a < x) \wedge (x < b)) \Rightarrow \exists l.\; \operatorname{is\text{-}diff\text{-}at}(h, x, l)) \wedge \\ \quad h(a) = h(b) \Rightarrow \\ \quad \quad \exists c \in \mathbb{R}.\; ((a < c) \wedge ((c < b) \wedge \operatorname{is\text{-}diff\text{-}at}(h, c, 0))) \end{array}}}

\vnbbookfile{theorem-library/mvt-proof}{theorem-library/mvt-proof}
\noindent{\small theorem-library/mvt-proof.scm -- the Mean Value Theorem, MACHINE-PROVEN. Strategy: apply Rolle's theorem to the auxiliary h(z) = f(z)*(b-a) - z*(f(b)-f(a)) which satisfies h(a)=h(b); Rolle gives an interior theta with h'(theta)=0,}\par
\vnbprop{mvt}{\vnbstmt{\begin{array}{@{}l@{}} \forall f, a, b. \\ \quad f \in (\mathbb{R} \to \mathbb{R}) \wedge \\ \quad a \in \mathbb{R} \wedge \\ \quad b \in \mathbb{R} \wedge \\ \quad a < b \wedge \\ \quad \forall x \in \operatorname{ccint}(a, b).\; \operatorname{is\text{-}continuous\text{-}at}(\operatorname{rr\text{-}ms}, \operatorname{rr\text{-}ms}, f, x) \wedge \\ \quad \forall x \in \mathbb{R}.\; (((a < x) \wedge (x < b)) \Rightarrow \exists l.\; \operatorname{is\text{-}diff\text{-}at}(f, x, l)) \Rightarrow \\ \quad \quad \exists c \in \mathbb{R}. \\ \quad \quad \quad ((a < c) \wedge ((c < b) \wedge \exists l.\; (\operatorname{is\text{-}diff\text{-}at}(f, c, l) \wedge ((l \cdot (b - a)) = (f(b) - f(a)))))) \end{array}}}

\vnbbookfile{theorem-library/deriv-constant-proof}{theorem-library/deriv-constant-proof}
\noindent{\small theorem-library/deriv-constant-proof.scm -- Cor 2.15, MACHINE-PROVEN. f continuous on [a,b], f'=0 on (a,b)  =$>$  f constant on [a,b]. Strategy: for u$<$v in [a,b], the proven MVT on [u,v] gives an interior theta with f'(theta)(v-u)=f(v)-f(u); the zero-derivative hypothesis + derivative-}\par
\vnbprop{deriv-zero-const-up}{\vnbstmt{\begin{array}{@{}l@{}} \forall f, a, b, u, v. \\ \quad f \in (\mathbb{R} \to \mathbb{R}) \wedge \\ \quad a \in \mathbb{R} \wedge \\ \quad b \in \mathbb{R} \wedge \\ \quad a < b \wedge \\ \quad \forall x \in \operatorname{ccint}(a, b).\; \operatorname{is\text{-}continuous\text{-}at}(\operatorname{rr\text{-}ms}, \operatorname{rr\text{-}ms}, f, x) \wedge \\ \quad \forall x \in \mathbb{R}.\; (((a < x) \wedge (x < b)) \Rightarrow \operatorname{is\text{-}diff\text{-}at}(f, x, 0)) \wedge \\ \quad u \in \operatorname{ccint}(a, b) \wedge \\ \quad v \in \operatorname{ccint}(a, b) \wedge \\ \quad u < v \Rightarrow \\ \quad \quad f(u) = f(v) \end{array}}}
\vnbprop{deriv-zero-implies-constant}{\vnbstmt{\begin{array}{@{}l@{}} \forall f, a, b, u, v. \\ \quad f \in (\mathbb{R} \to \mathbb{R}) \wedge \\ \quad a \in \mathbb{R} \wedge \\ \quad b \in \mathbb{R} \wedge \\ \quad a < b \wedge \\ \quad \forall x \in \operatorname{ccint}(a, b).\; \operatorname{is\text{-}continuous\text{-}at}(\operatorname{rr\text{-}ms}, \operatorname{rr\text{-}ms}, f, x) \wedge \\ \quad \forall x \in \mathbb{R}.\; (((a < x) \wedge (x < b)) \Rightarrow \operatorname{is\text{-}diff\text{-}at}(f, x, 0)) \wedge \\ \quad u \in \operatorname{ccint}(a, b) \wedge \\ \quad v \in \operatorname{ccint}(a, b) \Rightarrow \\ \quad \quad f(u) = f(v) \end{array}}}

\vnbbookfile{theorem-library/mvt-bounds-proof}{theorem-library/mvt-bounds-proof}
\noindent{\small theorem-library/mvt-bounds-proof.scm -- Cor 2.14, MACHINE-PROVEN. f' $<$= M on (a,b)  =$>$  f(b)-f(a) $<$= M(b-a)     (mvt-upper-bound) m $<$= f' on (a,b)  =$>$  m(b-a) $<$= f(b)-f(a)     (mvt-lower-bound) Strategy: MVT gives theta with f(b)-f(a) = f'(theta)(b-a) = L(b-a); the}\par
\vnbprop{mvt-upper-bound}{\vnbstmt{\begin{array}{@{}l@{}} \forall f, a, b, m. \\ \quad f \in (\mathbb{R} \to \mathbb{R}) \wedge \\ \quad a \in \mathbb{R} \wedge \\ \quad b \in \mathbb{R} \wedge \\ \quad m \in \mathbb{R} \wedge \\ \quad a < b \wedge \\ \quad \forall x \in \operatorname{ccint}(a, b).\; \operatorname{is\text{-}continuous\text{-}at}(\operatorname{rr\text{-}ms}, \operatorname{rr\text{-}ms}, f, x) \wedge \\ \quad \forall x \in \mathbb{R}. \\ \quad \quad ((a < x) \wedge (x < b)) \Rightarrow \exists l.\; (\operatorname{is\text{-}diff\text{-}at}(f, x, l) \wedge (l \leq m)) \Rightarrow \\ \quad \quad (f(b) - f(a)) \leq (m \cdot (b - a)) \end{array}}}
\vnbprop{mvt-lower-bound}{\vnbstmt{\begin{array}{@{}l@{}} \forall f, a, b, m. \\ \quad f \in (\mathbb{R} \to \mathbb{R}) \wedge \\ \quad a \in \mathbb{R} \wedge \\ \quad b \in \mathbb{R} \wedge \\ \quad m \in \mathbb{R} \wedge \\ \quad a < b \wedge \\ \quad \forall x \in \operatorname{ccint}(a, b).\; \operatorname{is\text{-}continuous\text{-}at}(\operatorname{rr\text{-}ms}, \operatorname{rr\text{-}ms}, f, x) \wedge \\ \quad \forall x \in \mathbb{R}. \\ \quad \quad ((a < x) \wedge (x < b)) \Rightarrow \exists l.\; (\operatorname{is\text{-}diff\text{-}at}(f, x, l) \wedge (m \leq l)) \Rightarrow \\ \quad \quad (m \cdot (b - a)) \leq (f(b) - f(a)) \end{array}}}

\vnbbookfile{theorem-library/mvt-abs-bound}{theorem-library/mvt-abs-bound}
\noindent{\small mvt-abs-bound.scm -- the TWO-SIDED corollary of the Mean Value Theorem:}\par
\vnbprop{mvt-abs-bound}{\vnbstmt{\begin{array}{@{}l@{}} \forall f, a, b, m. \\ \quad f \in (\mathbb{R} \to \mathbb{R}) \wedge \\ \quad a \in \mathbb{R} \wedge \\ \quad b \in \mathbb{R} \wedge \\ \quad m \in \mathbb{R} \wedge \\ \quad a < b \wedge \\ \quad \forall x \in \operatorname{ccint}(a, b).\; \operatorname{is\text{-}continuous\text{-}at}(\operatorname{rr\text{-}ms}, \operatorname{rr\text{-}ms}, f, x) \wedge \\ \quad \forall x \in \mathbb{R}. \\ \quad \quad ((a < x) \wedge (x < b)) \Rightarrow \exists l \in \mathbb{R}.\; (\operatorname{is\text{-}diff\text{-}at}(f, x, l) \wedge (\lvert l \rvert \leq m)) \Rightarrow \\ \quad \quad \lvert (f(b) - f(a)) \rvert \leq (m \cdot (b - a)) \end{array}}}

\vnbbookfile{theorem-library/continuity-local}{theorem-library/continuity-local}
\noindent{\small continuity-local.scm -- CONTINUITY AT A POINT READS ONLY A NEIGHBOURHOOD OF IT: a map agreeing with a continuous map on a BALL about the point, and saying nothing whatever elsewhere, is continuous at that point.  PROVEN.}\par
\vnbprop{continuous-at-local}{\vnbstmt{\begin{array}{@{}l@{}} \forall f, g \in (\mathbb{R} \to \mathbb{R}), t, o \in \mathbb{R}. \\ \quad 0 < o \wedge \\ \quad \operatorname{is\text{-}continuous\text{-}at}(\operatorname{rr\text{-}ms}, \operatorname{rr\text{-}ms}, g, t) \wedge \\ \quad \forall x \in \mathbb{R}.\; ((\lvert (x - t) \rvert \leq o) \Rightarrow (f(x) = g(x))) \Rightarrow \\ \quad \quad \operatorname{is\text{-}continuous\text{-}at}(\operatorname{rr\text{-}ms}, \operatorname{rr\text{-}ms}, f, t) \end{array}}}

\vnbbookfile{theorem-library/recip-star}{theorem-library/recip-star}
\noindent{\small recip-star.scm -- THE TOTALISED RECIPROCAL, and the reason it exists.}\par
\par\medskip\noindent\textbf{Definition} (\texttt{recip-star}).\;${\operatorname{u\_}}^{\ast\!-1} :=$ $\left(\text{if }(\operatorname{u\_} = 0)\text{ then }0\text{ else }\operatorname{recip}(\operatorname{u\_})\right)$
\vnbprop{recip-star-in-rr}{\vnbstmt{\forall u \in \mathbb{R}.\; ({u}^{\ast\!-1} \in \mathbb{R})}}
\vnbprop{recip-star-lam-in-fun}{\vnbstmt{(\lambda z \in \mathbb{R}.\; {z}^{\ast\!-1}) \in (\mathbb{R} \to \mathbb{R})}}
\vnbprop{recip-star-value}{\vnbstmt{\forall u \in \mathbb{R}.\; (\neg (u = 0) \Rightarrow ({u}^{\ast\!-1} = \operatorname{recip}(u)))}}
\vnbprop{recip-star-zero}{\vnbstmt{{0}^{\ast\!-1} = 0}}
\vnbprop{recip-star-inverse}{\vnbstmt{\forall u \in \mathbb{R}.\; (\neg (u = 0) \Rightarrow ((u \cdot {u}^{\ast\!-1}) = 1))}}
\vnbprop{rr-half-ball-nonzero}{\vnbstmt{\begin{array}{@{}l@{}} \forall c, m, y \in \mathbb{R}. \\ \quad ((0 < m) \wedge (((m + m) = \lvert c \rvert) \wedge (\lvert (y - c) \rvert \leq m))) \Rightarrow \neg (y = 0) \end{array}}}
\vnbprop{nonzero-ident-lam-in-fun}{\vnbstmt{(\lambda w \in \mathbb{R}.\; \left(\text{if }(w = 0)\text{ then }1\text{ else }w\right)) \in (\mathbb{R} \to \mathbb{R})}}
\vnbprop{nonzero-ident-nowhere-zero}{\vnbstmt{\forall x \in \mathbb{R}.\; \neg ((\lambda w \in \mathbb{R}.\; \left(\text{if }(w = 0)\text{ then }1\text{ else }w\right))(x) = 0)}}
\vnbprop{nonzero-ident-continuous-at}{\vnbstmt{\begin{array}{@{}l@{}} \forall c \in \mathbb{R}. \\ \quad \neg (c = 0) \Rightarrow  \\ \quad \quad \operatorname{is\text{-}continuous\text{-}at}(\operatorname{rr\text{-}ms}, \operatorname{rr\text{-}ms}, (\lambda w \in \mathbb{R}.\; \left(\text{if }(w = 0)\text{ then }1\text{ else }w\right)), c) \end{array}}}
\vnbprop{recip-star-continuous-at}{\vnbstmt{\begin{array}{@{}l@{}} \forall c \in \mathbb{R}. \\ \quad \neg (c = 0) \Rightarrow  \\ \quad \quad \operatorname{is\text{-}continuous\text{-}at}(\operatorname{rr\text{-}ms}, \operatorname{rr\text{-}ms}, (\lambda z \in \mathbb{R}.\; {z}^{\ast\!-1}), c) \end{array}}}
\vnbprop{recip-star-continuous-on-ccint}{\vnbstmt{\begin{array}{@{}l@{}} \forall a, b, x \in \mathbb{R}. \\ \quad ((0 < a) \wedge (x \in \operatorname{ccint}(a, b))) \Rightarrow  \\ \quad \quad \operatorname{is\text{-}continuous\text{-}at}(\operatorname{rr\text{-}ms}, \operatorname{rr\text{-}ms}, (\lambda z \in \mathbb{R}.\; {z}^{\ast\!-1}), x) \end{array}}}

\vnbbookfile{theorem-library/clamp}{theorem-library/clamp}
\noindent{\small clamp.scm -- CLAMP(a,b,x) = min(max(x,a),b), the retraction of the line onto the interval [a,b], and the ONE mechanism that dissolves the endpoint problem of docs/calculus.pdf Prop 4.16.}\par
\par\medskip\noindent\textbf{Definition} (\texttt{clamp}).\;$\operatorname{clamp}(a, b, x) :=$ $\operatorname{min}(\operatorname{max}(x, a), b)$
\vnbprop{clamp-in-rr}{\vnbstmt{\forall a, b, x \in \mathbb{R}.\; (\operatorname{clamp}(a, b, x) \in \mathbb{R})}}
\vnbprop{clamp-in-ccint}{\vnbstmt{\begin{array}{@{}l@{}} \forall a, b, x \in \mathbb{R}. \\ \quad (a \leq b) \Rightarrow (\operatorname{clamp}(a, b, x) \in \operatorname{ccint}(a, b)) \end{array}}}
\vnbprop{clamp-fixes}{\vnbstmt{\begin{array}{@{}l@{}} \forall a, b, x \in \mathbb{R}. \\ \quad ((a \leq x) \wedge (x \leq b)) \Rightarrow (\operatorname{clamp}(a, b, x) = x) \end{array}}}
\vnbprop{rr-max-lipschitz}{\vnbstmt{\forall u, v, c \in \mathbb{R}.\; (\lvert (\operatorname{max}(u, c) - \operatorname{max}(v, c)) \rvert \leq \lvert (u - v) \rvert)}}
\vnbprop{rr-min-lipschitz}{\vnbstmt{\forall u, v, c \in \mathbb{R}.\; (\lvert (\operatorname{min}(u, c) - \operatorname{min}(v, c)) \rvert \leq \lvert (u - v) \rvert)}}
\vnbprop{clamp-lipschitz}{\vnbstmt{\begin{array}{@{}l@{}} \forall a, b, y, x \in \mathbb{R}. \\ \quad \lvert (\operatorname{clamp}(a, b, y) - \operatorname{clamp}(a, b, x)) \rvert \leq \lvert (y - x) \rvert \end{array}}}
\vnbprop{clamp-compose-lam-in-fun}{\vnbstmt{\begin{array}{@{}l@{}} \forall f \in (\mathbb{R} \to \mathbb{R}), a, b \in \mathbb{R}. \\ \quad (\lambda z \in \mathbb{R}.\; f(\operatorname{clamp}(a, b, z))) \in (\mathbb{R} \to \mathbb{R}) \end{array}}}
\vnbprop{clamp-compose-fixes}{\vnbstmt{\begin{array}{@{}l@{}} \forall f \in (\mathbb{R} \to \mathbb{R}), a, b, x \in \mathbb{R}. \\ \quad ((a \leq x) \wedge (x \leq b)) \Rightarrow ((\lambda z \in \mathbb{R}.\; f(\operatorname{clamp}(a, b, z)))(x) = f(x)) \end{array}}}
\vnbprop{clamp-compose-continuous-at}{\vnbstmt{\begin{array}{@{}l@{}} \forall f \in (\mathbb{R} \to \mathbb{R}), a, b, t \in \mathbb{R}. \\ \quad ((a \leq b) \wedge \forall x \in \mathbb{R}.\; ((x \in \operatorname{ccint}(a, b)) \Rightarrow \operatorname{is\text{-}continuous\text{-}at}(\operatorname{rr\text{-}ms}, \operatorname{rr\text{-}ms}, f, x))) \Rightarrow  \\ \quad \quad \operatorname{is\text{-}continuous\text{-}at}(\operatorname{rr\text{-}ms}, \operatorname{rr\text{-}ms}, (\lambda z \in \mathbb{R}.\; f(\operatorname{clamp}(a, b, z))), t) \end{array}}}

\vnbbookfile{theorem-library/seq-limit}{theorem-library/seq-limit}
\noindent{\small theorem-library/seq-limit.scm -- THE LIMIT OF A REAL SEQUENCE AS A TERM, and the two pieces of limit arithmetic the tree was missing.}\par
\vnbprop{rr-limit-scale}{\vnbstmt{\begin{array}{@{}l@{}} \forall c \in \mathbb{R}, f, h \in (\mathbb{N} \to \mathbb{R}), v \in \mathbb{R}. \\ \quad (\forall j \in \mathbb{N}.\; (h(j) = (f(j) \cdot c)) \wedge \operatorname{converges\text{-}to}(\operatorname{rr\text{-}ms}, f, v)) \Rightarrow  \\ \quad \quad \operatorname{converges\text{-}to}(\operatorname{rr\text{-}ms}, h, (v \cdot c)) \end{array}}}
\vnbprop{rr-limit-sub}{\vnbstmt{\begin{array}{@{}l@{}} \forall f, g, h \in (\mathbb{N} \to \mathbb{R}), v, a \in \mathbb{R}. \\ \quad (\forall j \in \mathbb{N}.\; (h(j) = (f(j) - g(j))) \wedge (\operatorname{converges\text{-}to}(\operatorname{rr\text{-}ms}, f, v) \wedge \operatorname{converges\text{-}to}(\operatorname{rr\text{-}ms}, g, a))) \Rightarrow  \\ \quad \quad \operatorname{converges\text{-}to}(\operatorname{rr\text{-}ms}, h, (v - a)) \end{array}}}

\vnbbookfile{structure-library/regulated}{structure-library/regulated}
\noindent{\small MOVED 2026-09-20 (batch 12-A) from theorem-library/: this file is VOCABULARY -- definitions, notation and warranted supports, not one proof -- and every theorem stated with it had to load below it.  Its load.scm slot is unchanged. regulated.scm -- docs/calculus.pdf CHAPTER 4 SECTION 1: the vocabulary of}\par
\par\medskip\noindent\textbf{Definition} (\texttt{is-right-limit}).\;For all $f$, $x$ and $l$, 
\begin{equation*}
\begin{array}{@{}l@{}} \operatorname{is\text{-}right\text{-}limit}(f, x, l) \\ \Leftrightarrow\; f \in (\mathbb{R} \to \mathbb{R}) \wedge \\ x \in \mathbb{R} \wedge \\ l \in \mathbb{R} \wedge \\ \forall \operatorname{eps\_}. \\ \quad \operatorname{pos\text{-}rr}(\operatorname{eps\_}) \Rightarrow  \\ \quad \quad \exists \operatorname{delta\_}. \\ \quad \quad \quad \operatorname{pos\text{-}rr}(\operatorname{delta\_}) \wedge \\ \quad \quad \quad \forall \operatorname{t\_}. \\ \quad \quad \quad \quad ((\operatorname{t\_} \in \mathbb{R}) \wedge ((x < \operatorname{t\_}) \wedge (\operatorname{t\_} \leq +(x, \operatorname{delta\_})))) \Rightarrow (\lvert -(f(\operatorname{t\_}), l) \rvert \leq \operatorname{eps\_}) \end{array}
\end{equation*}

\par\medskip\noindent\textbf{Definition} (\texttt{is-regulated}).\;For all $f$, $a$ and $b$, 
\begin{equation*}
\begin{array}{@{}l@{}} \operatorname{is\text{-}regulated}(f, a, b) \\ \Leftrightarrow\; f \in (\mathbb{R} \to \mathbb{R}) \wedge \\ a \in \mathbb{R} \wedge \\ b \in \mathbb{R} \wedge \\ a < b \wedge \\ \forall \operatorname{x\_} \in \operatorname{ccint}(a, b). \\ \quad (\operatorname{x\_} < b) \Rightarrow \exists \operatorname{l\_}.\; \operatorname{is\text{-}right\text{-}limit}(f, \operatorname{x\_}, \operatorname{l\_}) \wedge \\ \quad (a < \operatorname{x\_}) \Rightarrow \exists \operatorname{l\_}.\; \operatorname{is\text{-}left\text{-}limit}(f, \operatorname{x\_}, \operatorname{l\_}) \end{array}
\end{equation*}

\par\medskip\noindent\textbf{Definition} (\texttt{is-unif-limit-on}).\;For all $\mathit{fam}$, $g$, $a$ and $b$, 
\begin{equation*}
\begin{array}{@{}l@{}} \operatorname{is\text{-}unif\text{-}limit\text{-}on}(\mathit{fam}, g, a, b) \\ \Leftrightarrow\; \mathit{fam} \in (\mathbb{N} \to (\mathbb{R} \to \mathbb{R})) \wedge \\ g \in (\mathbb{R} \to \mathbb{R}) \wedge \\ a \in \mathbb{R} \wedge \\ b \in \mathbb{R} \wedge \\ a < b \wedge \\ \forall \operatorname{eps\_}. \\ \quad \operatorname{pos\text{-}rr}(\operatorname{eps\_}) \Rightarrow  \\ \quad \quad \exists \operatorname{n\_} \in \mathbb{N}. \forall \operatorname{k\_}. \\ \quad \quad \quad ((\operatorname{k\_} \in \mathbb{N}) \wedge (\operatorname{n\_} \leq \operatorname{k\_})) \Rightarrow  \\ \quad \quad \quad \quad \forall \operatorname{x\_} \in \operatorname{ccint}(a, b).\; (\lvert -(g(\operatorname{x\_}), \operatorname{fam}(\operatorname{k\_})(\operatorname{x\_})) \rvert \leq \operatorname{eps\_}) \end{array}
\end{equation*}

\par\medskip\noindent\textbf{Definition} (\texttt{is-left-limit}).\;For all $f$, $x$ and $l$, 
\begin{equation*}
\begin{array}{@{}l@{}} \operatorname{is\text{-}left\text{-}limit}(f, x, l) \\ \Leftrightarrow\; f \in (\mathbb{R} \to \mathbb{R}) \wedge \\ x \in \mathbb{R} \wedge \\ l \in \mathbb{R} \wedge \\ \forall \operatorname{eps\_}. \\ \quad \operatorname{pos\text{-}rr}(\operatorname{eps\_}) \Rightarrow  \\ \quad \quad \exists \operatorname{delta\_}. \\ \quad \quad \quad \operatorname{pos\text{-}rr}(\operatorname{delta\_}) \wedge \\ \quad \quad \quad \forall \operatorname{t\_}. \\ \quad \quad \quad \quad ((\operatorname{t\_} \in \mathbb{R}) \wedge ((-(x, \operatorname{delta\_}) \leq \operatorname{t\_}) \wedge (\operatorname{t\_} < x))) \Rightarrow (\lvert -(f(\operatorname{t\_}), l) \rvert \leq \operatorname{eps\_}) \end{array}
\end{equation*}

\par\medskip\noindent\textbf{Definition} (\texttt{is-partition}).\;For all $p$, $n$, $a$ and $b$, 
\begin{equation*}
\begin{array}{@{}l@{}} \operatorname{is\text{-}partition}(p, n, a, b) \\ \Leftrightarrow\; n \in \mathbb{N} \wedge \\ 1 \leq n \wedge \\ a \in \mathbb{R} \wedge \\ b \in \mathbb{R} \wedge \\ a < b \wedge \\ p \in (\mathbb{N} \to \mathbb{R}) \wedge \\ p(0) = a \wedge \\ p(n) = b \wedge \\ \forall \operatorname{i\_}, \operatorname{j\_}. \\ \quad ((\operatorname{i\_} \in \operatorname{interval}(0, n)) \wedge ((\operatorname{j\_} \in \operatorname{interval}(0, n)) \wedge (\operatorname{i\_} < \operatorname{j\_}))) \Rightarrow (p(\operatorname{i\_}) < p(\operatorname{j\_})) \end{array}
\end{equation*}

\par\medskip\noindent\textbf{Definition} (\texttt{is-step-fn}).\;For all $f$, $a$ and $b$, 
\begin{equation*}
\begin{array}{@{}l@{}} \operatorname{is\text{-}step\text{-}fn}(f, a, b) \\ \Leftrightarrow\; f \in (\mathbb{R} \to \mathbb{R}) \wedge \\ a \in \mathbb{R} \wedge \\ b \in \mathbb{R} \wedge \\ a < b \wedge \\ \exists \operatorname{n\_}, \operatorname{p\_}. \\ \quad \operatorname{is\text{-}partition}(\operatorname{p\_}, \operatorname{n\_}, a, b) \wedge \\ \quad \forall \operatorname{i\_}. \\ \quad \quad ((\operatorname{i\_} \in \operatorname{interval}(0, \operatorname{n\_})) \wedge (\operatorname{i\_} < \operatorname{n\_})) \Rightarrow  \\ \quad \quad \quad \exists \operatorname{c\_} \in \mathbb{R}. \forall \operatorname{t\_}. \\ \quad \quad \quad \quad ((\operatorname{t\_} \in \mathbb{R}) \wedge ((\operatorname{p\_}(\operatorname{i\_}) < \operatorname{t\_}) \wedge (\operatorname{t\_} < \operatorname{p\_}(\operatorname{succ}(\operatorname{i\_}))))) \Rightarrow (f(\operatorname{t\_}) = \operatorname{c\_}) \end{array}
\end{equation*}

\par\medskip\noindent\textbf{Definition} (\texttt{is-piecewise-continuous}).\;For all $f$, $a$ and $b$, 
\begin{equation*}
\begin{array}{@{}l@{}} \operatorname{is\text{-}piecewise\text{-}continuous}(f, a, b) \\ \Leftrightarrow\; f \in (\mathbb{R} \to \mathbb{R}) \wedge \\ a \in \mathbb{R} \wedge \\ b \in \mathbb{R} \wedge \\ a < b \wedge \\ \exists \operatorname{n\_}, \operatorname{p\_}. \\ \quad \operatorname{is\text{-}partition}(\operatorname{p\_}, \operatorname{n\_}, a, b) \wedge \\ \quad \forall \operatorname{i\_}. \\ \quad \quad ((\operatorname{i\_} \in \operatorname{interval}(0, \operatorname{n\_})) \wedge (\operatorname{i\_} < \operatorname{n\_})) \Rightarrow  \\ \quad \quad \quad \forall \operatorname{t\_}. \\ \quad \quad \quad \quad ((\operatorname{t\_} \in \mathbb{R}) \wedge ((\operatorname{p\_}(\operatorname{i\_}) < \operatorname{t\_}) \wedge (\operatorname{t\_} < \operatorname{p\_}(\operatorname{succ}(\operatorname{i\_}))))) \Rightarrow  \\ \quad \quad \quad \quad \quad \operatorname{is\text{-}continuous\text{-}at}(\operatorname{rr\text{-}ms}, \operatorname{rr\text{-}ms}, f, \operatorname{t\_}) \end{array}
\end{equation*}


\vnbbookfile{theorem-library/onesided-limits}{theorem-library/onesided-limits}
\noindent{\small onesided-limits.scm -- docs/calculus.pdf Ch 4 Sec 1. Continuity at x gives BOTH one-sided limits at x, with value f(x).  Both `modulo 0'.  The two drivers are the same text with IS-RIGHT-LIMIT swapped for IS-LEFT-LIMIT: the interval side is the only difference, and `ineq!'}\par
\vnbprop{continuous-at-right-limit}{\vnbstmt{\begin{array}{@{}l@{}} \forall f, x. \\ \quad ((f \in (\mathbb{R} \to \mathbb{R})) \wedge ((x \in \mathbb{R}) \wedge \operatorname{is\text{-}continuous\text{-}at}(\operatorname{rr\text{-}ms}, \operatorname{rr\text{-}ms}, f, x))) \Rightarrow  \\ \quad \quad \operatorname{is\text{-}right\text{-}limit}(f, x, f(x)) \end{array}}}
\vnbprop{continuous-at-left-limit}{\vnbstmt{\begin{array}{@{}l@{}} \forall f, x. \\ \quad ((f \in (\mathbb{R} \to \mathbb{R})) \wedge ((x \in \mathbb{R}) \wedge \operatorname{is\text{-}continuous\text{-}at}(\operatorname{rr\text{-}ms}, \operatorname{rr\text{-}ms}, f, x))) \Rightarrow  \\ \quad \quad \operatorname{is\text{-}left\text{-}limit}(f, x, f(x)) \end{array}}}

\vnbbookfile{theorem-library/recip-succ-small}{theorem-library/recip-succ-small}
\noindent{\small recip-succ-small.scm -- the eventual form of the Archimedean bound.}\par
\vnbprop{recip-succ-small}{\vnbstmt{\begin{array}{@{}l@{}} \forall d. \\ \quad \operatorname{pos\text{-}rr}(d) \Rightarrow  \\ \quad \quad \exists n \in \mathbb{N}.\; \forall m \in \mathbb{N}.\; ((n \leq m) \Rightarrow (\operatorname{recip}((m + 1)) \leq d)) \end{array}}}

\vnbbookfile{theorem-library/continuous-is-regulated}{theorem-library/continuous-is-regulated}
\noindent{\small continuous-is-regulated.scm -- docs/calculus.pdf Ch 4, the bridge from Section 1 to the finished Section 2/3 arc.}\par
\vnbprop{continuous-is-regulated}{\vnbstmt{\begin{array}{@{}l@{}} \forall f, a, b. \\ \quad f \in (\mathbb{R} \to \mathbb{R}) \wedge \\ \quad a \in \mathbb{R} \wedge \\ \quad b \in \mathbb{R} \wedge \\ \quad a < b \wedge \\ \quad \forall x \in \operatorname{ccint}(a, b).\; \operatorname{is\text{-}continuous\text{-}at}(\operatorname{rr\text{-}ms}, \operatorname{rr\text{-}ms}, f, x) \Rightarrow \\ \quad \quad \operatorname{is\text{-}regulated}(f, a, b) \end{array}}}

\vnbbookfile{theorem-library/cauchy-criterion-left}{theorem-library/cauchy-criterion-left}
\noindent{\small cauchy-criterion-left.scm -- docs/calculus.pdf Ch 4 Sec 1.}\par
\vnbprop{cauchy-criterion-left}{\vnbstmt{\begin{array}{@{}l@{}} \forall f, x. \\ \quad f \in (\mathbb{R} \to \mathbb{R}) \wedge \\ \quad x \in \mathbb{R} \wedge \\ \quad \forall s. \\ \quad \quad \operatorname{pos\text{-}rr}(s) \Rightarrow  \\ \quad \quad \quad \exists a. \\ \quad \quad \quad \quad \operatorname{pos\text{-}rr}(a) \wedge \\ \quad \quad \quad \quad \forall b, t. \\ \quad \quad \quad \quad \quad (((b \in \mathbb{R}) \wedge (((x - a) \leq b) \wedge (b < x))) \wedge ((t \in \mathbb{R}) \wedge (((x - a) \leq t) \wedge (t < x)))) \Rightarrow  \\ \quad \quad \quad \quad \quad \quad \lvert (f(b) - f(t)) \rvert \leq s \Rightarrow \\ \quad \quad \exists l.\; \operatorname{is\text{-}left\text{-}limit}(f, x, l) \end{array}}}

\vnbbookfile{theorem-library/cauchy-criterion-right}{theorem-library/cauchy-criterion-right}
\noindent{\small cauchy-criterion-right.scm -- docs/calculus.pdf Ch 4 Sec 1.}\par
\vnbprop{cauchy-criterion-right}{\vnbstmt{\begin{array}{@{}l@{}} \forall f, x. \\ \quad f \in (\mathbb{R} \to \mathbb{R}) \wedge \\ \quad x \in \mathbb{R} \wedge \\ \quad \forall s. \\ \quad \quad \operatorname{pos\text{-}rr}(s) \Rightarrow  \\ \quad \quad \quad \exists a. \\ \quad \quad \quad \quad \operatorname{pos\text{-}rr}(a) \wedge \\ \quad \quad \quad \quad \forall b, t. \\ \quad \quad \quad \quad \quad (((b \in \mathbb{R}) \wedge ((x < b) \wedge (b \leq (x + a)))) \wedge ((t \in \mathbb{R}) \wedge ((x < t) \wedge (t \leq (x + a))))) \Rightarrow  \\ \quad \quad \quad \quad \quad \quad \lvert (f(b) - f(t)) \rvert \leq s \Rightarrow \\ \quad \quad \exists l.\; \operatorname{is\text{-}right\text{-}limit}(f, x, l) \end{array}}}

\vnbbookfile{theorem-library/partition-locates-point}{theorem-library/partition-locates-point}
\noindent{\small partition-locates-point.scm -- docs/calculus.pdf Ch 4 Sec 1.}\par
\vnbprop{partition-locates-point}{\vnbstmt{\begin{array}{@{}l@{}} \forall n \in \mathbb{N}, p, a, b, x. \\ \quad (\operatorname{is\text{-}partition}(p, n, a, b) \wedge (x \in \operatorname{ccint}(a, b))) \Rightarrow  \\ \quad \quad \exists i \in \operatorname{interval}(0, n).\; ((i < n) \wedge ((p(i) \leq x) \wedge (x \leq p((i + 1))))) \end{array}}}

\vnbbookfile{theorem-library/uniform-limit-regulated}{theorem-library/uniform-limit-regulated}
\noindent{\small uniform-limit-regulated.scm -- a uniform limit of regulated functions is regulated.  docs/calculus.pdf Prop 4.5, the sufficiency half.}\par
\vnbprop{uniform-limit-of-regulated}{\vnbstmt{\begin{array}{@{}l@{}} \forall m, g, a, b. \\ \quad (\operatorname{is\text{-}unif\text{-}limit\text{-}on}(m, g, a, b) \wedge \forall k \in \mathbb{N}.\; \operatorname{is\text{-}regulated}(m(k), a, b)) \Rightarrow  \\ \quad \quad \operatorname{is\text{-}regulated}(g, a, b) \end{array}}}

\vnbbookfile{theorem-library/uniform-limit-local}{theorem-library/uniform-limit-local}
\noindent{\small theorem-library/uniform-limit-local.scm -- A UNIFORM LIMIT IS CONTINUOUS AT A POINT, from uniform convergence on a BALL about that point only.}\par
\vnbprop{uniform-limit-continuous-at}{\vnbstmt{\begin{array}{@{}l@{}} \forall m \in (\mathbb{N} \to (\mathbb{R} \to \mathbb{R})), g \in (\mathbb{R} \to \mathbb{R}), t, o \in \mathbb{R}. \\ \quad 0 < o \wedge \\ \quad \forall k \in \mathbb{N}.\; \operatorname{is\text{-}continuous\text{-}at}(\operatorname{rr\text{-}ms}, \operatorname{rr\text{-}ms}, m(k), t) \wedge \\ \quad \forall s. \\ \quad \quad \operatorname{pos\text{-}rr}(s) \Rightarrow  \\ \quad \quad \quad \exists p \in \mathbb{N}. \forall k \in \mathbb{N}, x \in \mathbb{R}. \\ \quad \quad \quad \quad ((p \leq k) \wedge (\lvert (x - t) \rvert \leq o)) \Rightarrow (\lvert (g(x) - m(k)(x)) \rvert \leq s) \Rightarrow \\ \quad \quad \operatorname{is\text{-}continuous\text{-}at}(\operatorname{rr\text{-}ms}, \operatorname{rr\text{-}ms}, g, t) \end{array}}}

\vnbbookfile{theorem-library/diff-at-local}{theorem-library/diff-at-local}
\noindent{\small theorem-library/diff-at-local.scm -- DIFFERENTIABILITY IS LOCAL, and the Caratheodory transfer of a POINTWISE LIMIT.}\par
\vnbprop{diff-at-local}{\vnbstmt{\begin{array}{@{}l@{}} \forall f, i \in (\mathbb{R} \to \mathbb{R}), t, v, o \in \mathbb{R}. \\ \quad 0 < o \wedge \\ \quad \operatorname{is\text{-}continuous\text{-}at}(\operatorname{rr\text{-}ms}, \operatorname{rr\text{-}ms}, i, t) \wedge \\ \quad i(t) = v \wedge \\ \quad \forall x \in \mathbb{R}.\; ((\lvert (x - t) \rvert \leq o) \Rightarrow ((f(x) - f(t)) = (i(x) \cdot (x - t)))) \Rightarrow \\ \quad \quad \operatorname{is\text{-}diff\text{-}at}(f, t, v) \end{array}}}
\vnbprop{diff-at-ptwise-limit}{\vnbstmt{\begin{array}{@{}l@{}} \forall m, a \in (\mathbb{N} \to (\mathbb{R} \to \mathbb{R})), f, i \in (\mathbb{R} \to \mathbb{R}), t, o \in \mathbb{R}. \\ \quad 0 < o \wedge \\ \quad \operatorname{is\text{-}continuous\text{-}at}(\operatorname{rr\text{-}ms}, \operatorname{rr\text{-}ms}, i, t) \wedge \\ \quad \forall k \in \mathbb{N}, y \in \mathbb{R}. \\ \quad \quad (\lvert (y - t) \rvert \leq o) \Rightarrow ((m(k)(y) - m(k)(t)) = (a(k)(y) \cdot (y - t))) \wedge \\ \quad \forall y \in \mathbb{R}. \\ \quad \quad (\lvert (y - t) \rvert \leq o) \Rightarrow \operatorname{converges\text{-}to}(\operatorname{rr\text{-}ms}, (\lambda b \in \mathbb{N}.\; m(b)(y)), f(y)) \wedge \\ \quad \forall y \in \mathbb{R}. \\ \quad \quad (\lvert (y - t) \rvert \leq o) \Rightarrow \operatorname{converges\text{-}to}(\operatorname{rr\text{-}ms}, (\lambda b \in \mathbb{N}.\; a(b)(y)), i(y)) \Rightarrow \\ \quad \quad \operatorname{is\text{-}diff\text{-}at}(f, t, i(t)) \end{array}}}

\vnbbookfile{theorem-library/generalized-mvt-proof}{theorem-library/generalized-mvt-proof}
\noindent{\small theorem-library/generalized-mvt-proof.scm -- Cauchy MVT (Thm 2.11), PROVEN. L*(g(b)-g(a)) = M*(f(b)-f(a))  for some interior theta with f'=L, g'=M. Strategy: apply Rolle to the auxiliary h(z) = f(z)*(g(b)-g(a)) - g(z)*(f(b)-f(a))}\par
\vnbprop{generalized-mvt}{\vnbstmt{\begin{array}{@{}l@{}} \forall f, g, a, b. \\ \quad f \in (\mathbb{R} \to \mathbb{R}) \wedge \\ \quad g \in (\mathbb{R} \to \mathbb{R}) \wedge \\ \quad a \in \mathbb{R} \wedge \\ \quad b \in \mathbb{R} \wedge \\ \quad a < b \wedge \\ \quad \forall x \in \operatorname{ccint}(a, b). \\ \quad \quad \operatorname{is\text{-}continuous\text{-}at}(\operatorname{rr\text{-}ms}, \operatorname{rr\text{-}ms}, f, x) \wedge \operatorname{is\text{-}continuous\text{-}at}(\operatorname{rr\text{-}ms}, \operatorname{rr\text{-}ms}, g, x) \wedge \\ \quad \forall x \in \mathbb{R}. \\ \quad \quad ((a < x) \wedge (x < b)) \Rightarrow  \\ \quad \quad \quad \exists l.\; \operatorname{is\text{-}diff\text{-}at}(f, x, l) \wedge \exists m.\; \operatorname{is\text{-}diff\text{-}at}(g, x, m) \Rightarrow \\ \quad \quad \exists c \in \mathbb{R}. \\ \quad \quad \quad a < c \wedge \\ \quad \quad \quad c < b \wedge \\ \quad \quad \quad \exists l, m. \\ \quad \quad \quad \quad (\operatorname{is\text{-}diff\text{-}at}(f, c, l) \wedge (\operatorname{is\text{-}diff\text{-}at}(g, c, m) \wedge ((l \cdot (g(b) - g(a))) = (m \cdot (f(b) - f(a)))))) \end{array}}}

\vnbbookfile{theorem-library/deriv-monotone-proof}{theorem-library/deriv-monotone-proof}
\noindent{\small theorem-library/deriv-monotone-proof.scm -- increasing function theorem. f continuous on [a,b], f' $>$ 0 on (a,b)  =$>$  f strictly increasing on [a,b]. Strategy: for u$<$v in [a,b], MVT on [u,v] gives theta with f(v)-f(u) = f'(theta)(v-u);  f'(theta) $>$ 0 and v-u $>$ 0, so f(v)-f(u) $>$ 0.}\par
\vnbprop{deriv-pos-strictly-increasing}{\vnbstmt{\begin{array}{@{}l@{}} \forall f, a, b, u, v. \\ \quad f \in (\mathbb{R} \to \mathbb{R}) \wedge \\ \quad a \in \mathbb{R} \wedge \\ \quad b \in \mathbb{R} \wedge \\ \quad a < b \wedge \\ \quad \forall x \in \operatorname{ccint}(a, b).\; \operatorname{is\text{-}continuous\text{-}at}(\operatorname{rr\text{-}ms}, \operatorname{rr\text{-}ms}, f, x) \wedge \\ \quad \forall x \in \mathbb{R}. \\ \quad \quad ((a < x) \wedge (x < b)) \Rightarrow \exists l.\; (\operatorname{is\text{-}diff\text{-}at}(f, x, l) \wedge (0 < l)) \wedge \\ \quad u \in \operatorname{ccint}(a, b) \wedge \\ \quad v \in \operatorname{ccint}(a, b) \wedge \\ \quad u < v \Rightarrow \\ \quad \quad f(u) < f(v) \end{array}}}

\vnbbookfile{theorem-library/taylor-proof}{theorem-library/taylor-proof}
\noindent{\small RETIRED 2026-09-17 (proven): power-zero-base (was an add-to-pss here) -- theorem-library/rake-combinatorics.scm theorem-library/taylor-proof.scm -- Taylor's theorem, Lagrange remainder. f\textasciicircum{}(n) C\textasciicircum{}0 on [a,x], f\textasciicircum{}(n+1) exists on (a,x), a$<$x  =$>$  there is theta in (a,x) with  (n+1)! * R\_n(x) = f\textasciicircum{}(n+1)(theta) * (x-a)\textasciicircum{}(n+1),}\par
\par\medskip\noindent\textbf{Definition} (\texttt{taylor-poly}).\;$\operatorname{taylor\text{-}poly}(f, a, n, x) :=$
\begin{equation*}
\begin{array}{@{}l@{}} \operatorname{series\text{-}partial\text{-}sum}((\lambda k \in \mathbb{N}.\; ((\operatorname{nth\text{-}deriv}(f, k)(a) \cdot \operatorname{power}((x - a), k)) \cdot \operatorname{recip}(\operatorname{factorial}(k)))), \\ \quad (n + 1)) \end{array}
\end{equation*}

\par\medskip\noindent\textbf{Definition} (\texttt{taylor-differentiable}).\;For all $f$, $a$, $x$ and $n$, 
\begin{equation*}
\begin{array}{@{}l@{}} \operatorname{taylor\text{-}differentiable}(f, a, x, n) \\ \Leftrightarrow\; \forall k. \\ \quad ((k \in \mathbb{N}) \wedge (k \leq n)) \Rightarrow  \\ \quad \quad \forall t \in \operatorname{ccint}(a, x).\; \operatorname{is\text{-}continuous\text{-}at}(\operatorname{rr\text{-}ms}, \operatorname{rr\text{-}ms}, \operatorname{nth\text{-}deriv}(f, k), t) \wedge \\ \forall k. \\ \quad ((k \in \mathbb{N}) \wedge (k \leq n)) \Rightarrow  \\ \quad \quad \forall t. \\ \quad \quad \quad ((a < t) \wedge (t < x)) \Rightarrow  \\ \quad \quad \quad \quad \operatorname{is\text{-}diff\text{-}at}(\operatorname{nth\text{-}deriv}(f, k), t, \operatorname{nth\text{-}deriv}(f, \operatorname{succ}(k))(t)) \end{array}
\end{equation*}

\vnbprop{rr-recip-one}{\vnbstmt{\operatorname{recip}(1) = 1}}
\vnbprop{rr-mul-zero-mid}{\vnbstmt{\forall u, v \in \mathbb{R}.\; (((u \cdot 0) \cdot v) = 0)}}
\vnbprop{rr-mul-one-mid}{\vnbstmt{\forall u \in \mathbb{R}.\; (u = u)}}
\vnbprop{factorial-real-pos}{\vnbstmt{\forall n \in \mathbb{N}.\; ((\operatorname{factorial}(n) \in \mathbb{R}) \wedge (0 < \operatorname{factorial}(n)))}}
\vnbprop{recip-factorial-in-rr}{\vnbstmt{\forall n \in \mathbb{N}.\; (\operatorname{recip}(\operatorname{factorial}(n)) \in \mathbb{R})}}
\vnbprop{taylor-center-term-at-zero}{\vnbstmt{\begin{array}{@{}l@{}} \forall f.; \forall x \in \mathbb{R}. \\ \quad (f(x) \in \mathbb{R}) \Rightarrow  \\ \quad \quad (\lambda k \in \mathbb{N}.\; ((\operatorname{nth\text{-}deriv}(f, k)(x) \cdot \operatorname{power}((x - x), k)) \cdot \operatorname{recip}(\operatorname{factorial}(k))))(0) = f(x) \end{array}}}
\vnbprop{taylor-center-term-at-succ}{\vnbstmt{\begin{array}{@{}l@{}} \forall f.; \forall x \in \mathbb{R}, m \in \mathbb{N}. \\ \quad (\operatorname{nth\text{-}deriv}(f, (m + 1))(x) \in \mathbb{R}) \Rightarrow  \\ \quad \quad (\lambda k \in \mathbb{N}.\; ((\operatorname{nth\text{-}deriv}(f, k)(x) \cdot \operatorname{power}((x - x), k)) \cdot \operatorname{recip}(\operatorname{factorial}(k))))((m + 1)) = 0 \end{array}}}
\vnbprop{taylor-center-partial-sum}{\vnbstmt{\begin{array}{@{}l@{}} \forall n \in \mathbb{N}, f.; \forall x \in \mathbb{R}. \\ \quad \forall k \in \mathbb{N}.\; ((k \leq n) \Rightarrow (\operatorname{nth\text{-}deriv}(f, k)(x) \in \mathbb{R})) \Rightarrow  \\ \quad \quad \operatorname{series\text{-}partial\text{-}sum}((\lambda k \in \mathbb{N}.\; ((\operatorname{nth\text{-}deriv}(f, k)(x) \cdot \operatorname{power}((x - x), k)) \cdot \operatorname{recip}(\operatorname{factorial}(k)))), \\ \quad \quad \quad (n + 1)) \\ \quad \quad =\; f(x) \end{array}}}
\vnbprop{taylor-poly-at-center}{\vnbstmt{\begin{array}{@{}l@{}} \forall f, x, n. \\ \quad ((x \in \mathbb{R}) \wedge ((n \in \mathbb{N}) \wedge \forall k \in \mathbb{N}.\; ((k \leq n) \Rightarrow (\operatorname{nth\text{-}deriv}(f, k)(x) \in \mathbb{R})))) \Rightarrow  \\ \quad \quad \operatorname{taylor\text{-}poly}(f, x, n, x) = f(x) \end{array}}}
\vnbprop{taylor-derivs-in-fun}{\vnbstmt{\begin{array}{@{}l@{}} \forall f, a, x, n.; \forall k \in \mathbb{N}. \\ \quad (\operatorname{taylor\text{-}differentiable}(f, a, x, n) \wedge ((a \in \mathbb{R}) \wedge ((x \in \mathbb{R}) \wedge ((a < x) \wedge (k \leq n))))) \Rightarrow  \\ \quad \quad \operatorname{nth\text{-}deriv}(f, k) \in (\mathbb{R} \to \mathbb{R}) \end{array}}}
\vnbprop{taylor-derivs-values-real}{\vnbstmt{\begin{array}{@{}l@{}} \forall f, n.; \forall t \in \mathbb{R}, k \in \mathbb{N}. \\ \quad (\forall k \in \mathbb{N}.\; ((k \leq n) \Rightarrow (\operatorname{nth\text{-}deriv}(f, k) \in (\mathbb{R} \to \mathbb{R}))) \wedge (k \leq n)) \Rightarrow  \\ \quad \quad \operatorname{nth\text{-}deriv}(f, k)(t) \in \mathbb{R} \end{array}}}
\vnbprop{taylor-poly-in-rr-at}{\vnbstmt{\begin{array}{@{}l@{}} \forall n \in \mathbb{N}, f.; \forall a, x \in \mathbb{R}. \\ \quad \forall k \in \mathbb{N}.\; ((k \leq n) \Rightarrow (\operatorname{nth\text{-}deriv}(f, k)(a) \in \mathbb{R})) \Rightarrow  \\ \quad \quad \operatorname{series\text{-}partial\text{-}sum}((\lambda k \in \mathbb{N}.\; ((\operatorname{nth\text{-}deriv}(f, k)(a) \cdot \operatorname{power}((x - a), k)) \cdot \operatorname{recip}(\operatorname{factorial}(k)))), \\ \quad \quad \quad (n + 1)) \\ \quad \quad \in\; \mathbb{R} \end{array}}}
\vnbprop{taylor-poly-in-rr}{\vnbstmt{\begin{array}{@{}l@{}} \forall f \in (\mathbb{R} \to \mathbb{R}), a \in \mathbb{R}, n \in \mathbb{N}, x \in \mathbb{R}. \\ \quad \forall k \in \mathbb{N}.\; ((k \leq n) \Rightarrow (\operatorname{nth\text{-}deriv}(f, k) \in (\mathbb{R} \to \mathbb{R}))) \Rightarrow  \\ \quad \quad \operatorname{taylor\text{-}poly}(f, a, n, x) \in \mathbb{R} \end{array}}}
\vnbprop{taylor-g-in-fun}{\vnbstmt{\begin{array}{@{}l@{}} \forall n \in \mathbb{N}, f \in (\mathbb{R} \to \mathbb{R}), x \in \mathbb{R}. \\ \quad \forall k \in \mathbb{N}.\; ((k \leq n) \Rightarrow (\operatorname{nth\text{-}deriv}(f, k) \in (\mathbb{R} \to \mathbb{R}))) \Rightarrow  \\ \quad \quad (\lambda z \in \mathbb{R}.\; (f(x) - \operatorname{taylor\text{-}poly}(f, z, n, x))) \in (\mathbb{R} \to \mathbb{R}) \end{array}}}
\vnbprop{taylor-h-in-fun}{\vnbstmt{\forall n \in \mathbb{N}, x \in \mathbb{R}.\; ((\lambda z \in \mathbb{R}.\; \operatorname{power}((x - z), (n + 1))) \in (\mathbb{R} \to \mathbb{R}))}}
\vnbprop{taylor-h-at-a}{\vnbstmt{\begin{array}{@{}l@{}} \forall a, x, n. \\ \quad ((a \in \mathbb{R}) \wedge ((x \in \mathbb{R}) \wedge (n \in \mathbb{N}))) \Rightarrow  \\ \quad \quad (\lambda z \in \mathbb{R}.\; \operatorname{power}((x - z), (n + 1)))(a) = \operatorname{power}((x - a), (n + 1)) \end{array}}}
\vnbprop{taylor-h-at-x}{\vnbstmt{\begin{array}{@{}l@{}} \forall x \in \mathbb{R}, n \in \mathbb{N}. \\ \quad (\lambda z \in \mathbb{R}.\; \operatorname{power}((x - z), (n + 1)))(x) = 0 \end{array}}}
\vnbprop{taylor-g-at-a}{\vnbstmt{\begin{array}{@{}l@{}} \forall f, a, x, n. \\ \quad f \in (\mathbb{R} \to \mathbb{R}) \wedge \\ \quad a \in \mathbb{R} \wedge \\ \quad x \in \mathbb{R} \wedge \\ \quad n \in \mathbb{N} \wedge \\ \quad \forall k \in \mathbb{N}.\; ((k \leq n) \Rightarrow (\operatorname{nth\text{-}deriv}(f, k) \in (\mathbb{R} \to \mathbb{R}))) \Rightarrow \\ \quad \quad (\lambda z \in \mathbb{R}.\; (f(x) - \operatorname{taylor\text{-}poly}(f, z, n, x)))(a) = (f(x) - \operatorname{taylor\text{-}poly}(f, a, n, x)) \end{array}}}
\vnbprop{taylor-g-at-x}{\vnbstmt{\begin{array}{@{}l@{}} \forall f, x.; \forall n \in \mathbb{N}. \\ \quad ((f \in (\mathbb{R} \to \mathbb{R})) \wedge ((x \in \mathbb{R}) \wedge \forall k \in \mathbb{N}.\; ((k \leq n) \Rightarrow (\operatorname{nth\text{-}deriv}(f, k) \in (\mathbb{R} \to \mathbb{R}))))) \Rightarrow  \\ \quad \quad (\lambda z \in \mathbb{R}.\; (f(x) - \operatorname{taylor\text{-}poly}(f, z, n, x)))(x) = 0 \end{array}}}
\vnbprop{taylor-deriv-real}{\vnbstmt{\begin{array}{@{}l@{}} \forall f, a, x, n, t. \\ \quad (\operatorname{taylor\text{-}differentiable}(f, a, x, n) \wedge ((a < t) \wedge ((t < x) \wedge (n \in \mathbb{N})))) \Rightarrow  \\ \quad \quad \operatorname{nth\text{-}deriv}(f, (n + 1))(t) \in \mathbb{R} \end{array}}}
\vnbprop{tg-neg-scale}{\vnbstmt{\forall s, p \in \mathbb{R}.\; ((0 - (s \cdot p)) = ((s \cdot p) \cdot -1))}}
\vnbprop{taylor-h-diff}{\vnbstmt{\begin{array}{@{}l@{}} \forall x, n, t. \\ \quad ((x \in \mathbb{R}) \wedge ((t \in \mathbb{R}) \wedge (n \in \mathbb{N}))) \Rightarrow  \\ \quad \quad \operatorname{is\text{-}diff\text{-}at}((\lambda z \in \mathbb{R}.\; \operatorname{power}((x - z), (n + 1))), t, (0 - ((n + 1) \cdot \operatorname{power}((x - t), n)))) \end{array}}}
\vnbprop{tgd-recip-factorial-succ}{\vnbstmt{\forall m \in \mathbb{N}.\; (\operatorname{recip}(\operatorname{factorial}(m)) = ((m + 1) \cdot \operatorname{recip}(\operatorname{factorial}((m + 1)))))}}
\vnbprop{tgd-sub-diff}{\vnbstmt{\forall x, t \in \mathbb{R}.\; \operatorname{is\text{-}diff\text{-}at}((\lambda z \in \mathbb{R}.\; (x - z)), t, -1)}}
\vnbprop{tgd-power-sub-diff}{\vnbstmt{\begin{array}{@{}l@{}} \forall m \in \mathbb{N}, x, t \in \mathbb{R}. \\ \quad \operatorname{is\text{-}diff\text{-}at}((\lambda z \in \mathbb{R}.\; \operatorname{power}((x - z), (m + 1))), t, (((m + 1) \cdot \operatorname{power}((x - t), m)) \cdot -1)) \end{array}}}
\vnbprop{tgd-step-identity}{\vnbstmt{\begin{array}{@{}l@{}} \forall s, r, a, b, p, q \in \mathbb{R}. \\ \quad ((((b \cdot q) + (a \cdot ((s \cdot p) \cdot -1))) \cdot r) + ((a \cdot q) \cdot 0)) = (((r \cdot b) \cdot q) - (((s \cdot r) \cdot a) \cdot p)) \end{array}}}
\vnbprop{tgd-telescope-identity}{\vnbstmt{\forall u, v \in \mathbb{R}.\; (v = (u + (v - u)))}}
\vnbprop{tgd-term-diff}{\vnbstmt{\begin{array}{@{}l@{}} \forall m \in \mathbb{N}, f.; \forall x, t \in \mathbb{R}. \\ \quad \operatorname{is\text{-}diff\text{-}at}(\operatorname{nth\text{-}deriv}(f, (m + 1)), t, \operatorname{nth\text{-}deriv}(f, ((m + 1) + 1))(t)) \Rightarrow  \\ \quad \quad \operatorname{is\text{-}diff\text{-}at}((\lambda z \in \mathbb{R}.\; ((\operatorname{nth\text{-}deriv}(f, (m + 1))(z) \cdot \operatorname{power}((x - z), (m + 1))) \cdot \operatorname{recip}(\operatorname{factorial}((m + 1))))), \\ \quad \quad \quad t, \\ \quad \quad \quad (((\operatorname{recip}(\operatorname{factorial}((m + 1))) \cdot \operatorname{nth\text{-}deriv}(f, ((m + 1) + 1))(t)) \cdot \operatorname{power}((x - t), (m + 1))) - ((\operatorname{recip}(\operatorname{factorial}(m)) \cdot \operatorname{nth\text{-}deriv}(f, (m + 1))(t)) \cdot \operatorname{power}((x - t), m)))) \end{array}}}
\vnbprop{tgd-base-sum-value}{\vnbstmt{\begin{array}{@{}l@{}} \forall f \in (\mathbb{R} \to \mathbb{R}), x, z \in \mathbb{R}. \\ \quad (\operatorname{series\text{-}partial\text{-}sum}((\lambda k \in \mathbb{N}.\; ((\operatorname{nth\text{-}deriv}(f, k)(z) \cdot \operatorname{power}((x - z), k)) \cdot \operatorname{recip}(\operatorname{factorial}(k)))), 1) \simeq f(z)) \end{array}}}
\vnbprop{tgd-sum-diff}{\vnbstmt{\begin{array}{@{}l@{}} \forall m \in \mathbb{N}, f.; \forall x, t \in \mathbb{R}. \\ \quad \forall k \in \mathbb{N}.\; ((k \leq m) \Rightarrow \operatorname{is\text{-}diff\text{-}at}(\operatorname{nth\text{-}deriv}(f, k), t, \operatorname{nth\text{-}deriv}(f, (k + 1))(t))) \Rightarrow  \\ \quad \quad \operatorname{is\text{-}diff\text{-}at}((\lambda z \in \mathbb{R}.\; \\ \quad \quad \quad \quad \operatorname{series\text{-}partial\text{-}sum}((\lambda k \in \mathbb{N}.\; ((\operatorname{nth\text{-}deriv}(f, k)(z) \cdot \operatorname{power}((x - z), k)) \cdot \operatorname{recip}(\operatorname{factorial}(k)))), \\ \quad \quad \quad \quad \quad (m + 1))), \\ \quad \quad \quad t, \\ \quad \quad \quad (\operatorname{recip}(\operatorname{factorial}(m)) \cdot \operatorname{nth\text{-}deriv}(f, (m + 1))(t)) \cdot \operatorname{power}((x - t), m)) \end{array}}}
\vnbprop{taylor-g-diff}{\vnbstmt{\begin{array}{@{}l@{}} \forall f, a, x, n, t. \\ \quad ((n \in \mathbb{N}) \wedge ((x \in \mathbb{R}) \wedge (\operatorname{taylor\text{-}differentiable}(f, a, x, n) \wedge ((a < t) \wedge (t < x))))) \Rightarrow  \\ \quad \quad \operatorname{is\text{-}diff\text{-}at}((\lambda z \in \mathbb{R}.\; (f(x) - \operatorname{taylor\text{-}poly}(f, z, n, x))), t, \\ \quad \quad \quad (0 - ((\operatorname{recip}(\operatorname{factorial}(n)) \cdot \operatorname{nth\text{-}deriv}(f, (n + 1))(t)) \cdot \operatorname{power}((x - t), n)))) \end{array}}}
\vnbprop{taylor-h-cont}{\vnbstmt{\begin{array}{@{}l@{}} \forall x, n, t. \\ \quad ((x \in \mathbb{R}) \wedge ((t \in \mathbb{R}) \wedge (n \in \mathbb{N}))) \Rightarrow  \\ \quad \quad \operatorname{is\text{-}continuous\text{-}at}(\operatorname{rr\text{-}ms}, \operatorname{rr\text{-}ms}, (\lambda z \in \mathbb{R}.\; \operatorname{power}((x - z), (n + 1))), t) \end{array}}}
\vnbprop{tgc-term-cont}{\vnbstmt{\begin{array}{@{}l@{}} \forall m \in \mathbb{N}, f.; \forall x, t \in \mathbb{R}. \\ \quad \operatorname{is\text{-}continuous\text{-}at}(\operatorname{rr\text{-}ms}, \operatorname{rr\text{-}ms}, \operatorname{nth\text{-}deriv}(f, (m + 1)), t) \Rightarrow  \\ \quad \quad \operatorname{is\text{-}continuous\text{-}at}(\operatorname{rr\text{-}ms}, \operatorname{rr\text{-}ms}, \\ \quad \quad \quad (\lambda z \in \mathbb{R}.\; ((\operatorname{nth\text{-}deriv}(f, (m + 1))(z) \cdot \operatorname{power}((x - z), (m + 1))) \cdot \operatorname{recip}(\operatorname{factorial}((m + 1))))), \\ \quad \quad \quad t) \end{array}}}
\vnbprop{tgc-sum-cont}{\vnbstmt{\begin{array}{@{}l@{}} \forall m \in \mathbb{N}, f.; \forall x, t \in \mathbb{R}. \\ \quad \forall k \in \mathbb{N}.\; ((k \leq m) \Rightarrow \operatorname{is\text{-}continuous\text{-}at}(\operatorname{rr\text{-}ms}, \operatorname{rr\text{-}ms}, \operatorname{nth\text{-}deriv}(f, k), t)) \Rightarrow  \\ \quad \quad \operatorname{is\text{-}continuous\text{-}at}(\operatorname{rr\text{-}ms}, \operatorname{rr\text{-}ms}, \\ \quad \quad \quad (\lambda z \in \mathbb{R}.\; \\ \quad \quad \quad \quad \operatorname{series\text{-}partial\text{-}sum}((\lambda k \in \mathbb{N}.\; ((\operatorname{nth\text{-}deriv}(f, k)(z) \cdot \operatorname{power}((x - z), k)) \cdot \operatorname{recip}(\operatorname{factorial}(k)))), \\ \quad \quad \quad \quad \quad (m + 1))), \\ \quad \quad \quad t) \end{array}}}
\vnbprop{taylor-g-cont}{\vnbstmt{\begin{array}{@{}l@{}} \forall f, a, x, n, t. \\ \quad ((n \in \mathbb{N}) \wedge ((x \in \mathbb{R}) \wedge (\operatorname{taylor\text{-}differentiable}(f, a, x, n) \wedge (t \in \operatorname{ccint}(a, x))))) \Rightarrow  \\ \quad \quad \operatorname{is\text{-}continuous\text{-}at}(\operatorname{rr\text{-}ms}, \operatorname{rr\text{-}ms}, (\lambda z \in \mathbb{R}.\; (f(x) - \operatorname{taylor\text{-}poly}(f, z, n, x))), t) \end{array}}}
\vnbprop{taylor-gmvt-cont}{\vnbstmt{\begin{array}{@{}l@{}} \forall f, a, x, n.; \forall t \in \operatorname{ccint}(a, x). \\ \quad (\operatorname{taylor\text{-}differentiable}(f, a, x, n) \wedge ((x \in \mathbb{R}) \wedge (n \in \mathbb{N}))) \Rightarrow  \\ \quad \quad \operatorname{is\text{-}continuous\text{-}at}(\operatorname{rr\text{-}ms}, \operatorname{rr\text{-}ms}, (\lambda z \in \mathbb{R}.\; (f(x) - \operatorname{taylor\text{-}poly}(f, z, n, x))), t) \wedge \\ \quad \quad \operatorname{is\text{-}continuous\text{-}at}(\operatorname{rr\text{-}ms}, \operatorname{rr\text{-}ms}, (\lambda z \in \mathbb{R}.\; \operatorname{power}((x - z), (n + 1))), t) \end{array}}}
\vnbprop{taylor-gmvt-diff}{\vnbstmt{\begin{array}{@{}l@{}} \forall f, a, x, n.; \forall t \in \mathbb{R}. \\ \quad (\operatorname{taylor\text{-}differentiable}(f, a, x, n) \wedge ((n \in \mathbb{N}) \wedge ((x \in \mathbb{R}) \wedge ((a < t) \wedge (t < x))))) \Rightarrow  \\ \quad \quad \exists l. \\ \quad \quad \quad \operatorname{is\text{-}diff\text{-}at}((\lambda z \in \mathbb{R}.\; (f(x) - \operatorname{taylor\text{-}poly}(f, z, n, x))), t, l) \wedge \\ \quad \quad \exists m.\; \operatorname{is\text{-}diff\text{-}at}((\lambda z \in \mathbb{R}.\; \operatorname{power}((x - z), (n + 1))), t, m) \end{array}}}
\vnbprop{rr-power-pos}{\vnbstmt{\forall d \in \mathbb{R}, n \in \mathbb{N}.\; ((0 < d) \Rightarrow (0 < \operatorname{power}(d, n)))}}
\vnbprop{rr-cancel-mul-left}{\vnbstmt{\forall c, u, v \in \mathbb{R}.\; ((\neg (c = 0) \wedge ((c \cdot u) = (c \cdot v))) \Rightarrow (u = v))}}
\vnbprop{rr-recip-factorial}{\vnbstmt{\forall n \in \mathbb{N}.\; ((\operatorname{factorial}(n) \cdot \operatorname{recip}(\operatorname{factorial}(n))) = 1)}}
\vnbprop{taylor-clear}{\vnbstmt{\begin{array}{@{}l@{}} \forall n_{1}, w, d, r \in \mathbb{R}, x, a, b, c.; \forall n \in \mathbb{N}. \\ \quad 0 < w \wedge \\ \quad x = 0 \wedge \\ \quad b = 0 \wedge \\ \quad c = d \wedge \\ \quad a = r \wedge \\ \quad ((0 - ((\operatorname{recip}(\operatorname{factorial}(n)) \cdot n_{1}) \cdot w)) \cdot (b - c)) = ((0 - ((n + 1) \cdot w)) \cdot (x - a)) \Rightarrow \\ \quad \quad (\operatorname{factorial}((n + 1)) \cdot r) = (n_{1} \cdot d) \end{array}}}
\vnbprop{taylor-lagrange}{\vnbstmt{\begin{array}{@{}l@{}} \forall f, a, x, n. \\ \quad f \in (\mathbb{R} \to \mathbb{R}) \wedge \\ \quad a \in \mathbb{R} \wedge \\ \quad x \in \mathbb{R} \wedge \\ \quad n \in \mathbb{N} \wedge \\ \quad a < x \wedge \\ \quad \operatorname{taylor\text{-}differentiable}(f, a, x, n) \Rightarrow \\ \quad \quad \exists b \in \mathbb{R}. \\ \quad \quad \quad a < b \wedge \\ \quad \quad \quad b < x \wedge \\ \quad \quad \quad \operatorname{factorial}((n + 1)) \cdot (f(x) - \operatorname{taylor\text{-}poly}(f, a, n, x)) \\ \quad \quad \quad =\; \operatorname{nth\text{-}deriv}(f, (n + 1))(b) \cdot \operatorname{power}((x - a), (n + 1)) \end{array}}}

\vnbbookfile{theorem-library/rake-hb-leaves}{theorem-library/rake-hb-leaves}
\noindent{\small theorem-library/rake-hb-leaves.scm -- nine asserted leaves of the Hahn-Banach cluster, PROVEN.}\par
\vnbprop{extends-on-trans}{\vnbstmt{\begin{array}{@{}l@{}} \forall s, t, a, b, f. \\ \quad ((s \subseteq t) \wedge (\operatorname{extends\text{-}on}(t, a, b) \wedge \operatorname{extends\text{-}on}(s, b, f))) \Rightarrow  \\ \quad \quad \operatorname{extends\text{-}on}(s, a, f) \end{array}}}
\vnbprop{vnrm-nonneg}{\vnbstmt{\begin{array}{@{}l@{}} \forall m, w. \\ \quad (\operatorname{is\text{-}normed\text{-}vector\text{-}space}(m) \wedge (w \in \operatorname{vec}(m))) \Rightarrow (0 \leq \operatorname{vnrm}(m)(w)) \end{array}}}
\vnbprop{linfun-on-vec-is-linfun}{\vnbstmt{\begin{array}{@{}l@{}} \forall m, f. \\ \quad \operatorname{is\text{-}linear\text{-}functional\text{-}on}(m, \operatorname{vec}(m), f) \Rightarrow \operatorname{is\text{-}linear\text{-}functional}(m, f) \end{array}}}
\vnbprop{blf-app-real}{\vnbstmt{\begin{array}{@{}l@{}} \forall m, g, v. \\ \quad (\operatorname{is\text{-}bounded\text{-}linear\text{-}functional}(m, g) \wedge (v \in \operatorname{vec}(m))) \Rightarrow (g(v) \in \mathbb{R}) \end{array}}}
\vnbprop{bdd-linfun-abs-real}{\vnbstmt{\begin{array}{@{}l@{}} \forall m, f, x. \\ \quad (\operatorname{is\text{-}bounded\text{-}linear\text{-}functional}(m, f) \wedge (x \in \operatorname{vec}(m))) \Rightarrow (\lvert f(x) \rvert \in \mathbb{R}) \end{array}}}
\vnbprop{linfun-app-abs-real}{\vnbstmt{\begin{array}{@{}l@{}} \forall m, t, g, w. \\ \quad (\operatorname{is\text{-}linear\text{-}functional\text{-}on}(m, t, g) \wedge (w \in t)) \Rightarrow (\lvert g(w) \rvert \in \mathbb{R}) \end{array}}}
\vnbprop{abs-nonneg-le}{\vnbstmt{\forall a, c.\; (((a \in \mathbb{R}) \wedge ((0 \leq a) \wedge (\lvert a \rvert \leq c))) \Rightarrow (a \leq c))}}
\vnbprop{proper-subset-witness}{\vnbstmt{\forall a, b.\; (((a \subseteq b) \wedge \neg (a = b)) \Rightarrow \exists x \in b.\; \neg (x \in a))}}

\vnbbookfile{theorem-library/rake-dual-norm-spec}{theorem-library/rake-dual-norm-spec}
\noindent{\small rake-dual-norm-spec.scm -- the DUAL-NORM / DUAL-NORM-ON descriptions DENOTE, and the corollaries that follow from them.  Rake batch 7, assignment 7-A. Helper prefix `r8a-'.}\par
\vnbprop{rr-le-of-le-add-all-pos}{\vnbstmt{\begin{array}{@{}l@{}} \forall a, b \in \mathbb{R}. \\ \quad \forall d.\; (\operatorname{pos\text{-}rr}(d) \Rightarrow (a \leq (b + d))) \Rightarrow (a \leq b) \end{array}}}
\vnbprop{rr-recip-cancel-right}{\vnbstmt{\forall a, n \in \mathbb{R}.\; (\neg (n = 0) \Rightarrow (a = ((a \cdot \operatorname{recip}(n)) \cdot n)))}}
\vnbprop{rr-mul-recip-cancel}{\vnbstmt{\forall a, n \in \mathbb{R}.\; (\neg (n = 0) \Rightarrow (a = ((a \cdot n) \cdot \operatorname{recip}(n))))}}
\vnbprop{linfun-is-linfun-on-vec}{\vnbstmt{\begin{array}{@{}l@{}} \forall m, f. \\ \quad \operatorname{is\text{-}linear\text{-}functional}(m, f) \Rightarrow \operatorname{is\text{-}linear\text{-}functional\text{-}on}(m, \operatorname{vec}(m), f) \end{array}}}
\vnbprop{blf-is-blf-on-vec}{\vnbstmt{\begin{array}{@{}l@{}} \forall m, f. \\ \quad \operatorname{is\text{-}bounded\text{-}linear\text{-}functional}(m, f) \Rightarrow  \\ \quad \quad \operatorname{is\text{-}bounded\text{-}linear\text{-}functional\text{-}on}(m, \operatorname{vec}(m), f) \end{array}}}
\vnbprop{dual-norm-on-lub-exists}{\vnbstmt{\begin{array}{@{}l@{}} \forall m, s, f. \\ \quad (\operatorname{is\text{-}normed\text{-}vector\text{-}space}(m) \wedge ((s \subseteq \operatorname{vec}(m)) \wedge \operatorname{is\text{-}bounded\text{-}linear\text{-}functional\text{-}on}(m, s, f))) \Rightarrow  \\ \quad \quad \exists c \in \mathbb{R}. \\ \quad \quad \quad 0 \leq c \wedge \\ \quad \quad \quad \forall x \in s.\; (\lvert f(x) \rvert \leq (c \cdot \operatorname{vnrm}(m)(x))) \wedge \\ \quad \quad \quad \forall d \in \mathbb{R}. \\ \quad \quad \quad \quad ((0 \leq d) \wedge \forall x \in s.\; (\lvert f(x) \rvert \leq (d \cdot \operatorname{vnrm}(m)(x)))) \Rightarrow (c \leq d) \end{array}}}
\vnbprop{dual-norm-on-lub-unique}{\vnbstmt{\begin{array}{@{}l@{}} \forall m, s, f, c_{1}, c_{2}. \\ \quad c_{1} \in \mathbb{R} \wedge \\ \quad 0 \leq c_{1} \wedge \\ \quad \forall x \in s.\; (\lvert f(x) \rvert \leq (c_{1} \cdot \operatorname{vnrm}(m)(x))) \wedge \\ \quad \forall d \in \mathbb{R}. \\ \quad \quad ((0 \leq d) \wedge \forall x \in s.\; (\lvert f(x) \rvert \leq (d \cdot \operatorname{vnrm}(m)(x)))) \Rightarrow (c_{1} \leq d) \wedge \\ \quad c_{2} \in \mathbb{R} \wedge \\ \quad 0 \leq c_{2} \wedge \\ \quad \forall x \in s.\; (\lvert f(x) \rvert \leq (c_{2} \cdot \operatorname{vnrm}(m)(x))) \wedge \\ \quad \forall d \in \mathbb{R}. \\ \quad \quad ((0 \leq d) \wedge \forall x \in s.\; (\lvert f(x) \rvert \leq (d \cdot \operatorname{vnrm}(m)(x)))) \Rightarrow (c_{2} \leq d) \Rightarrow \\ \quad \quad c_{1} = c_{2} \end{array}}}
\vnbprop{dual-norm-on-spec}{\vnbstmt{\begin{array}{@{}l@{}} \forall m, s, f. \\ \quad (\operatorname{is\text{-}normed\text{-}vector\text{-}space}(m) \wedge ((s \subseteq \operatorname{vec}(m)) \wedge \operatorname{is\text{-}bounded\text{-}linear\text{-}functional\text{-}on}(m, s, f))) \Rightarrow  \\ \quad \quad \operatorname{dual\text{-}norm\text{-}on}(m, s, f) \in \mathbb{R} \wedge \\ \quad \quad 0 \leq \operatorname{dual\text{-}norm\text{-}on}(m, s, f) \wedge \\ \quad \quad \forall x \in s.\; (\lvert f(x) \rvert \leq (\operatorname{dual\text{-}norm\text{-}on}(m, s, f) \cdot \operatorname{vnrm}(m)(x))) \wedge \\ \quad \quad \forall d \in \mathbb{R}. \\ \quad \quad \quad ((0 \leq d) \wedge \forall x \in s.\; (\lvert f(x) \rvert \leq (d \cdot \operatorname{vnrm}(m)(x)))) \Rightarrow  \\ \quad \quad \quad \quad \operatorname{dual\text{-}norm\text{-}on}(m, s, f) \leq d \end{array}}}
\vnbprop{dual-norm-spec}{\vnbstmt{\begin{array}{@{}l@{}} \forall m, f. \\ \quad (\operatorname{is\text{-}normed\text{-}vector\text{-}space}(m) \wedge \operatorname{is\text{-}bounded\text{-}linear\text{-}functional}(m, f)) \Rightarrow  \\ \quad \quad \operatorname{dual\text{-}norm}(m, f) \in \mathbb{R} \wedge \\ \quad \quad 0 \leq \operatorname{dual\text{-}norm}(m, f) \wedge \\ \quad \quad \forall x \in \operatorname{vec}(m).\; (\lvert f(x) \rvert \leq (\operatorname{dual\text{-}norm}(m, f) \cdot \operatorname{vnrm}(m)(x))) \wedge \\ \quad \quad \forall d \in \mathbb{R}. \\ \quad \quad \quad ((0 \leq d) \wedge \forall x \in \operatorname{vec}(m).\; (\lvert f(x) \rvert \leq (d \cdot \operatorname{vnrm}(m)(x)))) \Rightarrow  \\ \quad \quad \quad \quad \operatorname{dual\text{-}norm}(m, f) \leq d \end{array}}}
\vnbprop{dual-norm-on-nonneg-nvs}{\vnbstmt{\begin{array}{@{}l@{}} \forall m, s, f. \\ \quad (\operatorname{is\text{-}normed\text{-}vector\text{-}space}(m) \wedge ((s \subseteq \operatorname{vec}(m)) \wedge \operatorname{is\text{-}bounded\text{-}linear\text{-}functional\text{-}on}(m, s, f))) \Rightarrow  \\ \quad \quad (\operatorname{dual\text{-}norm\text{-}on}(m, s, f) \in \mathbb{R}) \wedge (0 \leq \operatorname{dual\text{-}norm\text{-}on}(m, s, f)) \end{array}}}
\vnbprop{dual-norm-on-le-bound-nvs}{\vnbstmt{\begin{array}{@{}l@{}} \forall m, s, g, c. \\ \quad \operatorname{is\text{-}normed\text{-}vector\text{-}space}(m) \wedge \\ \quad s \subseteq \operatorname{vec}(m) \wedge \\ \quad \operatorname{is\text{-}linear\text{-}functional\text{-}on}(m, s, g) \wedge \\ \quad c \in \mathbb{R} \wedge \\ \quad 0 \leq c \wedge \\ \quad \forall w \in s.\; (\lvert g(w) \rvert \leq (c \cdot \operatorname{vnrm}(m)(w))) \Rightarrow \\ \quad \quad \operatorname{dual\text{-}norm\text{-}on}(m, s, g) \leq c \end{array}}}
\vnbprop{dual-norm-on-is-bound-nvs}{\vnbstmt{\begin{array}{@{}l@{}} \forall m, s, f, x. \\ \quad \operatorname{is\text{-}normed\text{-}vector\text{-}space}(m) \wedge \\ \quad s \subseteq \operatorname{vec}(m) \wedge \\ \quad \operatorname{is\text{-}bounded\text{-}linear\text{-}functional\text{-}on}(m, s, f) \wedge \\ \quad x \in s \Rightarrow \\ \quad \quad \lvert f(x) \rvert \leq (\operatorname{dual\text{-}norm\text{-}on}(m, s, f) \cdot \operatorname{vnrm}(m)(x)) \end{array}}}
\vnbprop{dual-norm-nonneg-nvs}{\vnbstmt{\begin{array}{@{}l@{}} \forall m, f. \\ \quad (\operatorname{is\text{-}normed\text{-}vector\text{-}space}(m) \wedge \operatorname{is\text{-}bounded\text{-}linear\text{-}functional}(m, f)) \Rightarrow  \\ \quad \quad (\operatorname{dual\text{-}norm}(m, f) \in \mathbb{R}) \wedge (0 \leq \operatorname{dual\text{-}norm}(m, f)) \end{array}}}
\vnbprop{dual-norm-is-bound-nvs}{\vnbstmt{\begin{array}{@{}l@{}} \forall m, f, x. \\ \quad (\operatorname{is\text{-}normed\text{-}vector\text{-}space}(m) \wedge (\operatorname{is\text{-}bounded\text{-}linear\text{-}functional}(m, f) \wedge (x \in \operatorname{vec}(m)))) \Rightarrow  \\ \quad \quad \lvert f(x) \rvert \leq (\operatorname{dual\text{-}norm}(m, f) \cdot \operatorname{vnrm}(m)(x)) \end{array}}}
\vnbprop{dual-norm-le-bound-nvs}{\vnbstmt{\begin{array}{@{}l@{}} \forall m, f, c. \\ \quad \operatorname{is\text{-}normed\text{-}vector\text{-}space}(m) \wedge \\ \quad \operatorname{is\text{-}linear\text{-}functional}(m, f) \wedge \\ \quad c \in \mathbb{R} \wedge \\ \quad 0 \leq c \wedge \\ \quad \forall w \in \operatorname{vec}(m).\; (\lvert f(w) \rvert \leq (c \cdot \operatorname{vnrm}(m)(w))) \Rightarrow \\ \quad \quad \operatorname{dual\text{-}norm}(m, f) \leq c \end{array}}}
\vnbprop{good-sub-self-nvs}{\vnbstmt{\begin{array}{@{}l@{}} \forall m, s, f. \\ \quad (\operatorname{is\text{-}normed\text{-}vector\text{-}space}(m) \wedge (\operatorname{is\text{-}submodule}(m, s) \wedge \operatorname{is\text{-}bounded\text{-}linear\text{-}functional\text{-}on}(m, s, f))) \Rightarrow  \\ \quad \quad \operatorname{good\text{-}sub}(m, s, f, s) \end{array}}}

\vnbbookfile{theorem-library/limit-frequent}{theorem-library/limit-frequent}
\noindent{\small limit-frequent.scm -- wave of 2026-11-09/10 (agent BA-2B), wired 2026-11-10 from scratchpad/ba2b-1109/limit-frequent.scm. Library: analysis.}\par
\vnbprop{limit-of-frequent-value}{\vnbstmt{\begin{array}{@{}l@{}} \forall s, f, \ell, c. \\ \quad (\operatorname{converges\text{-}to}(s, f, \ell) \wedge \forall n \in \mathbb{N}.\; \exists k \in \mathbb{N}.\; ((n \leq k) \wedge (f(k) = c))) \Rightarrow  \\ \quad \quad \ell = c \end{array}}}
\vnbprop{rr-limit-of-frequent-value}{\vnbstmt{\begin{array}{@{}l@{}} \forall f, \ell, c. \\ \quad (\operatorname{converges\text{-}to}(\operatorname{rr\text{-}ms}, f, \ell) \wedge \forall n \in \mathbb{N}.\; \exists k \in \mathbb{N}.\; ((n \leq k) \wedge (f(k) = c))) \Rightarrow  \\ \quad \quad \ell = c \end{array}}}

\vnbbookfile{theorem-library/nvs-module-view}{theorem-library/nvs-module-view}
\noindent{\small nvs-module-view.scm -- THE MODULE BRIDGE OF A NORMED VECTOR SPACE, PROVEN.}\par
\vnbprop{nvs-module-view-scal}{\vnbstmt{\begin{array}{@{}l@{}} \forall m. \\ \quad \operatorname{is\text{-}normed\text{-}vector\text{-}space}(m) \Rightarrow  \\ \quad \quad (\operatorname{scal}(\operatorname{normed\text{-}vector\text{-}space\text{-}as\text{-}module}(m)) \simeq \operatorname{scal}(m)) \end{array}}}
\vnbprop{nvs-module-view-vec}{\vnbstmt{\begin{array}{@{}l@{}} \forall m. \\ \quad \operatorname{is\text{-}normed\text{-}vector\text{-}space}(m) \Rightarrow  \\ \quad \quad (\operatorname{vec}(\operatorname{normed\text{-}vector\text{-}space\text{-}as\text{-}module}(m)) \simeq \operatorname{vec}(m)) \end{array}}}
\vnbprop{nvs-module-view-vadd}{\vnbstmt{\begin{array}{@{}l@{}} \forall m. \\ \quad \operatorname{is\text{-}normed\text{-}vector\text{-}space}(m) \Rightarrow  \\ \quad \quad (\operatorname{vadd}(\operatorname{normed\text{-}vector\text{-}space\text{-}as\text{-}module}(m)) \simeq \operatorname{vadd}(m)) \end{array}}}
\vnbprop{nvs-module-view-vzero}{\vnbstmt{\begin{array}{@{}l@{}} \forall m. \\ \quad \operatorname{is\text{-}normed\text{-}vector\text{-}space}(m) \Rightarrow  \\ \quad \quad (\operatorname{vzero}(\operatorname{normed\text{-}vector\text{-}space\text{-}as\text{-}module}(m)) \simeq \operatorname{vzero}(m)) \end{array}}}
\vnbprop{nvs-module-view-vneg}{\vnbstmt{\begin{array}{@{}l@{}} \forall m. \\ \quad \operatorname{is\text{-}normed\text{-}vector\text{-}space}(m) \Rightarrow  \\ \quad \quad (\operatorname{vneg}(\operatorname{normed\text{-}vector\text{-}space\text{-}as\text{-}module}(m)) \simeq \operatorname{vneg}(m)) \end{array}}}
\vnbprop{nvs-module-view-act}{\vnbstmt{\begin{array}{@{}l@{}} \forall m. \\ \quad \operatorname{is\text{-}normed\text{-}vector\text{-}space}(m) \Rightarrow  \\ \quad \quad (\operatorname{act}(\operatorname{normed\text{-}vector\text{-}space\text{-}as\text{-}module}(m)) \simeq \operatorname{act}(m)) \end{array}}}
\vnbprop{submodule-nvs-module-view}{\vnbstmt{\begin{array}{@{}l@{}} \forall m, t. \\ \quad (\operatorname{is\text{-}normed\text{-}vector\text{-}space}(m) \wedge \operatorname{is\text{-}submodule}(m, t)) \Rightarrow  \\ \quad \quad \operatorname{is\text{-}submodule}(\operatorname{normed\text{-}vector\text{-}space\text{-}as\text{-}module}(m), t) \end{array}}}

\vnbbookfile{theorem-library/nvs-act-laws}{theorem-library/nvs-act-laws}
\noindent{\small nvs-act-laws.scm -- THE GROUP AND SCALAR-ACTION LAWS OF A REAL NORMED VECTOR SPACE.  Since 2026-09-19 (rake batch 8, assignment 8-K1) this file is their ONE home: section (4) below holds the laws that three to five agents had each proved separately on 2026-09-19 under their own prefixes.  The NORM}\par
\vnbprop{nvs-scal-carr}{\vnbstmt{\forall m.\; (\operatorname{is\text{-}normed\text{-}vector\text{-}space}(m) \Rightarrow (\operatorname{carr}(\operatorname{scal}(m)) \simeq \mathbb{R}))}}
\vnbprop{nvs-scal-add}{\vnbstmt{\begin{array}{@{}l@{}} \forall m. \\ \quad \operatorname{is\text{-}normed\text{-}vector\text{-}space}(m) \Rightarrow  \\ \quad \quad (\operatorname{add}(\operatorname{scal}(m)) \simeq (\lambda {[}\operatorname{x\_}, \operatorname{y\_}] \in (\mathbb{R} \times \mathbb{R}).\; (\operatorname{x\_} + \operatorname{y\_}))) \end{array}}}
\vnbprop{nvs-scal-mul}{\vnbstmt{\begin{array}{@{}l@{}} \forall m. \\ \quad \operatorname{is\text{-}normed\text{-}vector\text{-}space}(m) \Rightarrow  \\ \quad \quad (\operatorname{mul}(\operatorname{scal}(m)) \simeq (\lambda {[}\operatorname{x\_}, \operatorname{y\_}] \in (\mathbb{R} \times \mathbb{R}).\; (\operatorname{x\_} \cdot \operatorname{y\_}))) \end{array}}}
\vnbprop{nvs-scal-zero}{\vnbstmt{\forall m.\; (\operatorname{is\text{-}normed\text{-}vector\text{-}space}(m) \Rightarrow (\operatorname{zero}(\operatorname{scal}(m)) \simeq 0))}}
\vnbprop{nvs-scal-one}{\vnbstmt{\forall m.\; (\operatorname{is\text{-}normed\text{-}vector\text{-}space}(m) \Rightarrow (\operatorname{one}(\operatorname{scal}(m)) \simeq 1))}}
\vnbprop{nvs-act-distrib-scalar}{\vnbstmt{\begin{array}{@{}l@{}} \forall m.; \forall r, s \in \operatorname{carr}(\operatorname{scal}(m)), x \in \operatorname{vec}(m). \\ \quad \operatorname{is\text{-}normed\text{-}vector\text{-}space}(m) \Rightarrow  \\ \quad \quad \operatorname{act}(m)(\operatorname{add}(\operatorname{scal}(m))(r, s), x) = \operatorname{vadd}(m)(\operatorname{act}(m)(r, x), \operatorname{act}(m)(s, x)) \end{array}}}
\vnbprop{nvs-act-mul-compat}{\vnbstmt{\begin{array}{@{}l@{}} \forall m.; \forall r, s \in \operatorname{carr}(\operatorname{scal}(m)), x \in \operatorname{vec}(m). \\ \quad \operatorname{is\text{-}normed\text{-}vector\text{-}space}(m) \Rightarrow  \\ \quad \quad \operatorname{act}(m)(\operatorname{mul}(\operatorname{scal}(m))(r, s), x) = \operatorname{act}(m)(r, \operatorname{act}(m)(s, x)) \end{array}}}
\vnbprop{nvs-act-unital}{\vnbstmt{\begin{array}{@{}l@{}} \forall m.; \forall x \in \operatorname{vec}(m). \\ \quad \operatorname{is\text{-}normed\text{-}vector\text{-}space}(m) \Rightarrow (\operatorname{act}(m)(\operatorname{one}(\operatorname{scal}(m)), x) = x) \end{array}}}
\vnbprop{nvs-vadd-assoc}{\vnbstmt{\begin{array}{@{}l@{}} \forall m.; \forall u, v, w \in \operatorname{vec}(m). \\ \quad \operatorname{is\text{-}normed\text{-}vector\text{-}space}(m) \Rightarrow  \\ \quad \quad \operatorname{vadd}(m)(\operatorname{vadd}(m)(u, v), w) = \operatorname{vadd}(m)(u, \operatorname{vadd}(m)(v, w)) \end{array}}}
\vnbprop{nvs-vzero-right}{\vnbstmt{\begin{array}{@{}l@{}} \forall m.; \forall u \in \operatorname{vec}(m). \\ \quad \operatorname{is\text{-}normed\text{-}vector\text{-}space}(m) \Rightarrow (\operatorname{vadd}(m)(u, \operatorname{vzero}(m)) = u) \end{array}}}
\vnbprop{nvs-act-type}{\vnbstmt{\begin{array}{@{}l@{}} \forall m.; \forall r \in \operatorname{carr}(\operatorname{scal}(m)), x \in \operatorname{vec}(m). \\ \quad \operatorname{is\text{-}normed\text{-}vector\text{-}space}(m) \Rightarrow (\operatorname{act}(m)(r, x) \in \operatorname{vec}(m)) \end{array}}}
\vnbprop{nvs-act-in-vec}{\vnbstmt{\begin{array}{@{}l@{}} \forall m.; \forall r \in \mathbb{R}, x \in \operatorname{vec}(m). \\ \quad \operatorname{is\text{-}normed\text{-}vector\text{-}space}(m) \Rightarrow (\operatorname{act}(m)(r, x) \in \operatorname{vec}(m)) \end{array}}}
\vnbprop{nvs-act-scale-assoc}{\vnbstmt{\begin{array}{@{}l@{}} \forall m.; \forall r, s \in \mathbb{R}, x \in \operatorname{vec}(m). \\ \quad \operatorname{is\text{-}normed\text{-}vector\text{-}space}(m) \Rightarrow  \\ \quad \quad \operatorname{act}(m)(r, \operatorname{act}(m)(s, x)) = \operatorname{act}(m)((r \cdot s), x) \end{array}}}
\vnbprop{nvs-act-collect}{\vnbstmt{\begin{array}{@{}l@{}} \forall m.; \forall r, s \in \mathbb{R}, x, y \in \operatorname{vec}(m). \\ \quad \operatorname{is\text{-}normed\text{-}vector\text{-}space}(m) \Rightarrow  \\ \quad \quad \operatorname{vadd}(m)(\operatorname{vadd}(m)(y, \operatorname{act}(m)(r, x)), \operatorname{act}(m)(s, x)) = \operatorname{vadd}(m)(y, \operatorname{act}(m)((r + s), x)) \end{array}}}
\vnbprop{nvs-act-zero}{\vnbstmt{\begin{array}{@{}l@{}} \forall m.; \forall x, y \in \operatorname{vec}(m). \\ \quad \operatorname{is\text{-}normed\text{-}vector\text{-}space}(m) \Rightarrow (\operatorname{vadd}(m)(y, \operatorname{act}(m)(0, x)) = y) \end{array}}}
\vnbprop{nvs-act-one}{\vnbstmt{\begin{array}{@{}l@{}} \forall m.; \forall x, y \in \operatorname{vec}(m). \\ \quad \operatorname{is\text{-}normed\text{-}vector\text{-}space}(m) \Rightarrow (\operatorname{vadd}(m)(y, \operatorname{act}(m)(1, x)) = \operatorname{vadd}(m)(y, x)) \end{array}}}
\vnbprop{nvs-scal-neg}{\vnbstmt{\begin{array}{@{}l@{}} \forall m. \\ \quad \operatorname{is\text{-}normed\text{-}vector\text{-}space}(m) \Rightarrow (\operatorname{neg}(\operatorname{scal}(m)) \simeq (\lambda x \in \mathbb{R}.\; -x)) \end{array}}}
\vnbprop{nvs-vzero-left}{\vnbstmt{\begin{array}{@{}l@{}} \forall m.; \forall u \in \operatorname{vec}(m). \\ \quad \operatorname{is\text{-}normed\text{-}vector\text{-}space}(m) \Rightarrow (\operatorname{vadd}(m)(\operatorname{vzero}(m), u) = u) \end{array}}}
\vnbprop{nvs-vadd-comm}{\vnbstmt{\begin{array}{@{}l@{}} \forall m.; \forall u, w \in \operatorname{vec}(m). \\ \quad \operatorname{is\text{-}normed\text{-}vector\text{-}space}(m) \Rightarrow (\operatorname{vadd}(m)(u, w) = \operatorname{vadd}(m)(w, u)) \end{array}}}
\vnbprop{nvs-act-distrib-vec}{\vnbstmt{\begin{array}{@{}l@{}} \forall m.; \forall r \in \operatorname{carr}(\operatorname{scal}(m)), x, y \in \operatorname{vec}(m). \\ \quad \operatorname{is\text{-}normed\text{-}vector\text{-}space}(m) \Rightarrow  \\ \quad \quad \operatorname{act}(m)(r, \operatorname{vadd}(m)(x, y)) = \operatorname{vadd}(m)(\operatorname{act}(m)(r, x), \operatorname{act}(m)(r, y)) \end{array}}}
\vnbprop{nvs-zero-act}{\vnbstmt{\begin{array}{@{}l@{}} \forall m.; \forall x \in \operatorname{vec}(m). \\ \quad \operatorname{is\text{-}normed\text{-}vector\text{-}space}(m) \Rightarrow (\operatorname{act}(m)(0, x) = \operatorname{vzero}(m)) \end{array}}}
\vnbprop{nvs-neg-one-act}{\vnbstmt{\begin{array}{@{}l@{}} \forall m.; \forall x \in \operatorname{vec}(m). \\ \quad \operatorname{is\text{-}normed\text{-}vector\text{-}space}(m) \Rightarrow (\operatorname{act}(m)(-1, x) = \operatorname{vneg}(m)(x)) \end{array}}}
\vnbprop{nvs-act-add-scalars}{\vnbstmt{\begin{array}{@{}l@{}} \forall m.; \forall r, s \in \mathbb{R}, x \in \operatorname{vec}(m). \\ \quad \operatorname{is\text{-}normed\text{-}vector\text{-}space}(m) \Rightarrow  \\ \quad \quad \operatorname{vadd}(m)(\operatorname{act}(m)(r, x), \operatorname{act}(m)(s, x)) = \operatorname{act}(m)((r + s), x) \end{array}}}
\vnbprop{nvs-act-vneg}{\vnbstmt{\begin{array}{@{}l@{}} \forall m.; \forall r \in \mathbb{R}, x \in \operatorname{vec}(m). \\ \quad \operatorname{is\text{-}normed\text{-}vector\text{-}space}(m) \Rightarrow  \\ \quad \quad \operatorname{vneg}(m)(\operatorname{act}(m)(r, x)) = \operatorname{act}(m)((-1 \cdot r), x) \end{array}}}
\vnbprop{nvs-vneg-vadd}{\vnbstmt{\begin{array}{@{}l@{}} \forall m.; \forall x, y \in \operatorname{vec}(m). \\ \quad \operatorname{is\text{-}normed\text{-}vector\text{-}space}(m) \Rightarrow  \\ \quad \quad \operatorname{vneg}(m)(\operatorname{vadd}(m)(x, y)) = \operatorname{vadd}(m)(\operatorname{vneg}(m)(x), \operatorname{vneg}(m)(y)) \end{array}}}
\vnbprop{nvs-vadd-shuffle}{\vnbstmt{\begin{array}{@{}l@{}} \forall m.; \forall a, b, c, d \in \operatorname{vec}(m). \\ \quad \operatorname{is\text{-}normed\text{-}vector\text{-}space}(m) \Rightarrow  \\ \quad \quad \operatorname{vadd}(m)(\operatorname{vadd}(m)(a, c), \operatorname{vadd}(m)(b, d)) = \operatorname{vadd}(m)(\operatorname{vadd}(m)(a, b), \operatorname{vadd}(m)(c, d)) \end{array}}}
\vnbprop{nvs-vzero-in-vec}{\vnbstmt{\forall m.\; (\operatorname{is\text{-}normed\text{-}vector\text{-}space}(m) \Rightarrow (\operatorname{vzero}(m) \in \operatorname{vec}(m)))}}
\vnbprop{nvs-vneg-in-vec}{\vnbstmt{\begin{array}{@{}l@{}} \forall m.; \forall x \in \operatorname{vec}(m). \\ \quad \operatorname{is\text{-}normed\text{-}vector\text{-}space}(m) \Rightarrow (\operatorname{vneg}(m)(x) \in \operatorname{vec}(m)) \end{array}}}
\vnbprop{nvs-vneg-left}{\vnbstmt{\begin{array}{@{}l@{}} \forall m.; \forall u \in \operatorname{vec}(m). \\ \quad \operatorname{is\text{-}normed\text{-}vector\text{-}space}(m) \Rightarrow (\operatorname{vadd}(m)(\operatorname{vneg}(m)(u), u) = \operatorname{vzero}(m)) \end{array}}}
\vnbprop{nvs-act-distrib-vec-rr}{\vnbstmt{\begin{array}{@{}l@{}} \forall m.; \forall r \in \mathbb{R}, x, y \in \operatorname{vec}(m). \\ \quad \operatorname{is\text{-}normed\text{-}vector\text{-}space}(m) \Rightarrow  \\ \quad \quad \operatorname{act}(m)(r, \operatorname{vadd}(m)(x, y)) = \operatorname{vadd}(m)(\operatorname{act}(m)(r, x), \operatorname{act}(m)(r, y)) \end{array}}}
\vnbprop{nvs-one-act}{\vnbstmt{\forall m.; \forall x \in \operatorname{vec}(m).\; (\operatorname{is\text{-}normed\text{-}vector\text{-}space}(m) \Rightarrow (\operatorname{act}(m)(1, x) = x))}}
\vnbprop{nvs-vadd-cancel}{\vnbstmt{\begin{array}{@{}l@{}} \forall m.; \forall x, z, w \in \operatorname{vec}(m). \\ \quad (\operatorname{is\text{-}normed\text{-}vector\text{-}space}(m) \wedge (\operatorname{vadd}(m)(x, z) = \operatorname{vadd}(m)(x, w))) \Rightarrow (z = w) \end{array}}}
\vnbprop{nvs-vadd-cancel-right}{\vnbstmt{\begin{array}{@{}l@{}} \forall m.; \forall x, z, w \in \operatorname{vec}(m). \\ \quad (\operatorname{is\text{-}normed\text{-}vector\text{-}space}(m) \wedge (\operatorname{vadd}(m)(z, x) = \operatorname{vadd}(m)(w, x))) \Rightarrow (z = w) \end{array}}}
\vnbprop{nvs-vneg-right}{\vnbstmt{\begin{array}{@{}l@{}} \forall m.; \forall u \in \operatorname{vec}(m). \\ \quad \operatorname{is\text{-}normed\text{-}vector\text{-}space}(m) \Rightarrow (\operatorname{vadd}(m)(u, \operatorname{vneg}(m)(u)) = \operatorname{vzero}(m)) \end{array}}}
\vnbprop{nvs-vadd-cancel-zero}{\vnbstmt{\begin{array}{@{}l@{}} \forall m.; \forall x, z \in \operatorname{vec}(m). \\ \quad (\operatorname{is\text{-}normed\text{-}vector\text{-}space}(m) \wedge (\operatorname{vadd}(m)(x, z) = x)) \Rightarrow (z = \operatorname{vzero}(m)) \end{array}}}
\vnbprop{nvs-neg-one-act-sub}{\vnbstmt{\begin{array}{@{}l@{}} \forall m.; \forall u \in \operatorname{vec}(m). \\ \quad \operatorname{is\text{-}normed\text{-}vector\text{-}space}(m) \Rightarrow (\operatorname{act}(m)(-1, u) = \operatorname{vneg}(m)(u)) \end{array}}}

\vnbbookfile{theorem-library/nvs-normed-ag-typing}{theorem-library/nvs-normed-ag-typing}
\noindent{\small theorem-library/nvs-normed-ag-typing.scm -- THE NORMED GROUP OF A NORMED VECTOR SPACE (D8, 2026-11-15; the user's decision "Prove views").}\par
\vnbprop{normed-vector-space-as-normed-ag-is-normed-ag}{\vnbstmt{\begin{array}{@{}l@{}} \forall r. \\ \quad \operatorname{is\text{-}normed\text{-}vector\text{-}space}(r) \Rightarrow  \\ \quad \quad \operatorname{is\text{-}normed\text{-}ag}(\operatorname{normed\text{-}vector\text{-}space\text{-}as\text{-}normed\text{-}ag}(r)) \end{array}}}

\vnbbookfile{theorem-library/rake-norm-metrics}{theorem-library/rake-norm-metrics}
\noindent{\small rake-norm-metrics.scm -- rake batch 5 R: continuity projections and the norm metrics, PROVEN.}\par
\vnbprop{ag-norm-metric-is-metric-space}{\vnbstmt{\begin{array}{@{}l@{}} \forall g, m. \\ \quad (\operatorname{is\text{-}abelian\text{-}group}(g) \wedge \operatorname{is\text{-}group\text{-}norm}(m, \operatorname{opr}(g), \operatorname{inv}(g), \operatorname{iden}(g), \operatorname{carr}(g))) \Rightarrow  \\ \quad \quad \operatorname{is\text{-}metric\text{-}space}({[}\operatorname{carr}(g), (\lambda {[}u, v] \in (\operatorname{carr}(g) \times \operatorname{carr}(g)).\; m(\operatorname{opr}(g)(u, \operatorname{inv}(g)(v))))]) \end{array}}}
\vnbprop{continuous-is-continuous-at}{\vnbstmt{\begin{array}{@{}l@{}} \forall s, t, f.; \forall a \in \operatorname{pts}(s). \\ \quad \operatorname{is\text{-}continuous}(s, t, f) \Rightarrow \operatorname{is\text{-}continuous\text{-}at}(s, t, f, a) \end{array}}}
\vnbprop{uniformly-continuous-is-continuous}{\vnbstmt{\begin{array}{@{}l@{}} \forall s, t, f. \\ \quad \operatorname{is\text{-}uniformly\text{-}continuous}(s, t, f) \Rightarrow \operatorname{is\text{-}continuous}(s, t, f) \end{array}}}
\vnbprop{nag-metric-carrier}{\vnbstmt{\forall g.\; (\operatorname{is\text{-}normed\text{-}ag}(g) \Rightarrow (\operatorname{pts}(\operatorname{nag\text{-}metric\text{-}space}(g)) \simeq \operatorname{carr}(g)))}}
\vnbprop{nag-metric-distance}{\vnbstmt{\begin{array}{@{}l@{}} \forall g.; \forall u, v \in \operatorname{carr}(g). \\ \quad \operatorname{is\text{-}normed\text{-}ag}(g) \Rightarrow  \\ \quad \quad \operatorname{dist}(\operatorname{nag\text{-}metric\text{-}space}(g))(u, v) = \operatorname{nrm}(g)(\operatorname{opr}(g)(u, \operatorname{inv}(g)(v))) \end{array}}}
\vnbprop{nag-metric-space-is-metric-space}{\vnbstmt{\forall g.\; (\operatorname{is\text{-}normed\text{-}ag}(g) \Rightarrow \operatorname{is\text{-}metric\text{-}space}(\operatorname{nag\text{-}metric\text{-}space}(g)))}}
\vnbprop{nvs-metric-is-ms}{\vnbstmt{\begin{array}{@{}l@{}} \forall m. \\ \quad \operatorname{is\text{-}normed\text{-}vector\text{-}space}(m) \Rightarrow \operatorname{is\text{-}metric\text{-}space}(\operatorname{nvs\text{-}metric\text{-}space}(m)) \end{array}}}

\vnbbookfile{theorem-library/basw3-nag}{theorem-library/basw3-nag}
\noindent{\small basw3-nag.scm -- wave of 2026-11-09/10 (agent BA-SW3), wired 2026-11-10 from scratchpad/basw3-1109/basw3-nag.scm. Library: analysis.}\par
\vnbprop{nag-complete-of-norm-cauchy}{\vnbstmt{\begin{array}{@{}l@{}} \forall g. \\ \quad \operatorname{is\text{-}normed\text{-}ag}(g) \wedge \\ \quad \forall f \in (\mathbb{N} \to \operatorname{carr}(g)). \\ \quad \quad \forall s. \\ \quad \quad \quad \operatorname{pos\text{-}rr}(s) \Rightarrow  \\ \quad \quad \quad \quad \exists n \in \mathbb{N}. \forall m, a \in \mathbb{N}. \\ \quad \quad \quad \quad \quad ((n \leq m) \wedge (n \leq a)) \Rightarrow (\operatorname{nrm}(g)(\operatorname{opr}(g)(f(m), \operatorname{inv}(g)(f(a)))) \leq s) \Rightarrow \\ \quad \quad \quad \exists l \in \operatorname{carr}(g). \forall s. \\ \quad \quad \quad \quad \operatorname{pos\text{-}rr}(s) \Rightarrow  \\ \quad \quad \quad \quad \quad \exists n \in \mathbb{N}. \forall a \in \mathbb{N}. \\ \quad \quad \quad \quad \quad \quad (n \leq a) \Rightarrow (\operatorname{nrm}(g)(\operatorname{opr}(g)(f(a), \operatorname{inv}(g)(l))) \leq s) \Rightarrow \\ \quad \quad \operatorname{is\text{-}complete}(\operatorname{nag\text{-}metric\text{-}space}(g)) \end{array}}}

\vnbbookfile{theorem-library/nvs-norm-laws}{theorem-library/nvs-norm-laws}
\noindent{\small nvs-norm-laws.scm -- THE NORM LAWS OF A REAL NORMED VECTOR SPACE, in ONE place.  Rake batch 8, assignment 8-K1.}\par
\vnbprop{nvs-vnrm-triangle}{\vnbstmt{\begin{array}{@{}l@{}} \forall m.; \forall x, y \in \operatorname{vec}(m). \\ \quad \operatorname{is\text{-}normed\text{-}vector\text{-}space}(m) \Rightarrow  \\ \quad \quad \operatorname{vnrm}(m)(\operatorname{vadd}(m)(x, y)) \leq (\operatorname{vnrm}(m)(x) + \operatorname{vnrm}(m)(y)) \end{array}}}
\vnbprop{nvs-vnrm-homog-c}{\vnbstmt{\begin{array}{@{}l@{}} \forall m.; \forall r \in \operatorname{carr}(\operatorname{scal}(m)), x \in \operatorname{vec}(m). \\ \quad \operatorname{is\text{-}normed\text{-}vector\text{-}space}(m) \Rightarrow  \\ \quad \quad \operatorname{vnrm}(m)(\operatorname{act}(m)(r, x)) = (\lvert r \rvert \cdot \operatorname{vnrm}(m)(x)) \end{array}}}
\vnbprop{nvs-vnrm-homog}{\vnbstmt{\begin{array}{@{}l@{}} \forall m.; \forall r \in \mathbb{R}, x \in \operatorname{vec}(m). \\ \quad \operatorname{is\text{-}normed\text{-}vector\text{-}space}(m) \Rightarrow  \\ \quad \quad \operatorname{vnrm}(m)(\operatorname{act}(m)(r, x)) = (\lvert r \rvert \cdot \operatorname{vnrm}(m)(x)) \end{array}}}
\vnbprop{nvs-vnrm-vzero}{\vnbstmt{\forall m.\; (\operatorname{is\text{-}normed\text{-}vector\text{-}space}(m) \Rightarrow (\operatorname{vnrm}(m)(\operatorname{vzero}(m)) = 0))}}
\vnbprop{nvs-ms-pts}{\vnbstmt{\begin{array}{@{}l@{}} \forall m. \\ \quad \operatorname{is\text{-}metric\text{-}space}(\operatorname{nvs\text{-}metric\text{-}space}(m)) \Rightarrow (\operatorname{pts}(\operatorname{nvs\text{-}metric\text{-}space}(m)) \simeq \operatorname{vec}(m)) \end{array}}}
\vnbprop{nvs-metric-distance}{\vnbstmt{\begin{array}{@{}l@{}} \forall m.; \forall u, v \in \operatorname{vec}(m). \\ \quad \operatorname{is\text{-}normed\text{-}vector\text{-}space}(m) \Rightarrow  \\ \quad \quad \operatorname{dist}(\operatorname{nvs\text{-}metric\text{-}space}(m))(u, v) = \operatorname{vnrm}(m)(\operatorname{vadd}(m)(u, \operatorname{vneg}(m)(v))) \end{array}}}
\vnbprop{nvs-metric-distance-q}{\vnbstmt{\begin{array}{@{}l@{}} \forall m.; \forall u, v \in \operatorname{vec}(m). \\ \quad \operatorname{is\text{-}normed\text{-}vector\text{-}space}(m) \Rightarrow  \\ \quad \quad (\operatorname{dist}(\operatorname{nvs\text{-}metric\text{-}space}(m))(u, v) \simeq \operatorname{vnrm}(m)(\operatorname{vadd}(m)(u, \operatorname{vneg}(m)(v)))) \end{array}}}
\vnbprop{nvs-dist-shift}{\vnbstmt{\begin{array}{@{}l@{}} \forall m.; \forall u, w \in \operatorname{vec}(m). \\ \quad \operatorname{is\text{-}normed\text{-}vector\text{-}space}(m) \Rightarrow  \\ \quad \quad \operatorname{dist}(\operatorname{nvs\text{-}metric\text{-}space}(m))(u, \operatorname{vadd}(m)(u, w)) = \operatorname{vnrm}(m)(w) \end{array}}}

\vnbbookfile{theorem-library/rake-hb-submodules}{theorem-library/rake-hb-submodules}
\noindent{\small rake-hb-submodules.scm -- the two SUBMODULE leaves of the Hahn-Banach cluster, PROVEN, and the normed-vector-space action bricks they need.}\par
\vnbprop{line-unfold}{\vnbstmt{\begin{array}{@{}l@{}} \forall m, v. \\ \quad (\operatorname{line}(m, v) \simeq \{y \in \operatorname{vec}(m) : \exists r \in \mathbb{R}.\; (y = \operatorname{act}(m)(r, v))\}) \end{array}}}
\vnbprop{line-is-submodule}{\vnbstmt{\begin{array}{@{}l@{}} \forall m, v. \\ \quad (\operatorname{is\text{-}normed\text{-}vector\text{-}space}(m) \wedge (v \in \operatorname{vec}(m))) \Rightarrow \operatorname{is\text{-}submodule}(m, \operatorname{line}(m, v)) \end{array}}}
\vnbprop{span-add-one-submodule}{\vnbstmt{\begin{array}{@{}l@{}} \forall m, t, v. \\ \quad (\operatorname{is\text{-}normed\text{-}vector\text{-}space}(m) \wedge (\operatorname{is\text{-}submodule}(m, t) \wedge (v \in \operatorname{vec}(m)))) \Rightarrow  \\ \quad \quad \operatorname{is\text{-}submodule}(m, \operatorname{span\text{-}add\text{-}one}(m, t, v)) \end{array}}}

\vnbbookfile{theorem-library/rake-span-add-one-guarded}{theorem-library/rake-span-add-one-guarded}
\noindent{\small rake-span-add-one-guarded.scm -- the GUARDED forms of span-add-one-superset and span-add-one-has-v, PROVEN.}\par
\vnbprop{span-add-one-superset-nvs}{\vnbstmt{\begin{array}{@{}l@{}} \forall m, t, v. \\ \quad (\operatorname{is\text{-}normed\text{-}vector\text{-}space}(m) \wedge (\operatorname{is\text{-}submodule}(m, t) \wedge (v \in \operatorname{vec}(m)))) \Rightarrow  \\ \quad \quad t \subseteq \operatorname{span\text{-}add\text{-}one}(m, t, v) \end{array}}}
\vnbprop{span-add-one-has-v-nvs}{\vnbstmt{\begin{array}{@{}l@{}} \forall m, t, v. \\ \quad (\operatorname{is\text{-}normed\text{-}vector\text{-}space}(m) \wedge (\operatorname{is\text{-}submodule}(m, t) \wedge (v \in \operatorname{vec}(m)))) \Rightarrow  \\ \quad \quad v \in \operatorname{span\text{-}add\text{-}one}(m, t, v) \end{array}}}

\vnbbookfile{theorem-library/rake-hb-leaves-2}{theorem-library/rake-hb-leaves-2}
\noindent{\small theorem-library/rake-hb-leaves-2.scm -- three Hahn-Banach leaves, PROVEN modulo 0: good-sub-submodule, good-sub-in-power, line-has-v.}\par
\vnbprop{good-sub-submodule}{\vnbstmt{\forall m, s, f, t.\; (\operatorname{good\text{-}sub}(m, s, f, t) \Rightarrow \operatorname{is\text{-}submodule}(m, t))}}
\vnbprop{good-sub-in-power}{\vnbstmt{\begin{array}{@{}l@{}} \forall m, s, f, t. \\ \quad (\operatorname{is\text{-}normed\text{-}vector\text{-}space}(m) \wedge \operatorname{good\text{-}sub}(m, s, f, t)) \Rightarrow (t \in \operatorname{power}(\operatorname{vec}(m))) \end{array}}}
\vnbprop{line-has-v}{\vnbstmt{\begin{array}{@{}l@{}} \forall m, v. \\ \quad (\operatorname{is\text{-}normed\text{-}vector\text{-}space}(m) \wedge (v \in \operatorname{vec}(m))) \Rightarrow (v \in \operatorname{line}(m, v)) \end{array}}}

\vnbbookfile{theorem-library/rake-hb-gap}{theorem-library/rake-hb-gap}
\noindent{\small rake-hb-gap.scm -- hb-gap (theorem-library/hahn-banach-proof.scm:24), the analytic core of the one-dimension Hahn-Banach step, PROVEN.}\par
\vnbprop{hbg-scale-pos}{\vnbstmt{\begin{array}{@{}l@{}} \forall d, t, u, a, o \in \mathbb{R}. \\ \quad ((0 \leq o) \wedge (((-u - (d \cdot t)) \leq a) \wedge (a \leq ((d \cdot t) - u)))) \Rightarrow  \\ \quad \quad \lvert ((o \cdot u) + (o \cdot a)) \rvert \leq (d \cdot (o \cdot t)) \end{array}}}
\vnbprop{hbg-scale-neg}{\vnbstmt{\begin{array}{@{}l@{}} \forall d, t, u, a, o \in \mathbb{R}. \\ \quad ((o \leq 0) \wedge (((-u - (d \cdot t)) \leq a) \wedge (a \leq ((d \cdot t) - u)))) \Rightarrow  \\ \quad \quad \lvert ((o \cdot u) + (o \cdot a)) \rvert \leq (d \cdot (-o \cdot t)) \end{array}}}
\vnbprop{hbg-neg-one-cancel}{\vnbstmt{\begin{array}{@{}l@{}} \forall m.; \forall y, v \in \operatorname{vec}(m). \\ \quad \operatorname{is\text{-}normed\text{-}vector\text{-}space}(m) \Rightarrow  \\ \quad \quad \operatorname{vadd}(m)(\operatorname{vadd}(m)(y, v), \operatorname{act}(m)(-1, v)) = y \end{array}}}
\vnbprop{hbg-shift}{\vnbstmt{\begin{array}{@{}l@{}} \forall m.; \forall w, z, v \in \operatorname{vec}(m). \\ \quad \operatorname{is\text{-}normed\text{-}vector\text{-}space}(m) \Rightarrow  \\ \quad \quad \operatorname{vadd}(m)(w, \operatorname{act}(m)(-1, z)) = \operatorname{vadd}(m)(\operatorname{vadd}(m)(w, v), \operatorname{act}(m)(-1, \operatorname{vadd}(m)(z, v))) \end{array}}}
\vnbprop{hbg-diff-norm}{\vnbstmt{\begin{array}{@{}l@{}} \forall m.; \forall w, z, v \in \operatorname{vec}(m). \\ \quad \operatorname{is\text{-}normed\text{-}vector\text{-}space}(m) \Rightarrow  \\ \quad \quad \operatorname{vnrm}(m)(\operatorname{vadd}(m)(w, \operatorname{act}(m)(-1, z))) \leq (\operatorname{vnrm}(m)(\operatorname{vadd}(m)(w, v)) + \operatorname{vnrm}(m)(\operatorname{vadd}(m)(z, v))) \end{array}}}
\vnbprop{hbg-scale-back}{\vnbstmt{\begin{array}{@{}l@{}} \forall m.; \forall y \in \operatorname{vec}(m), r \in \mathbb{R}. \\ \quad (\operatorname{is\text{-}normed\text{-}vector\text{-}space}(m) \wedge \neg (r = 0)) \Rightarrow  \\ \quad \quad \operatorname{act}(m)(r, \operatorname{act}(m)(\operatorname{recip}(r), y)) = y \end{array}}}
\vnbprop{hbg-scale-decomp}{\vnbstmt{\begin{array}{@{}l@{}} \forall m.; \forall y, v \in \operatorname{vec}(m), r \in \mathbb{R}. \\ \quad (\operatorname{is\text{-}normed\text{-}vector\text{-}space}(m) \wedge \neg (r = 0)) \Rightarrow  \\ \quad \quad \operatorname{vadd}(m)(y, \operatorname{act}(m)(r, v)) = \operatorname{act}(m)(r, \operatorname{vadd}(m)(\operatorname{act}(m)(\operatorname{recip}(r), y), v)) \end{array}}}
\vnbprop{hbg-linfun-app-real}{\vnbstmt{\begin{array}{@{}l@{}} \forall m, s, f, x. \\ \quad (\operatorname{is\text{-}linear\text{-}functional\text{-}on}(m, s, f) \wedge (x \in s)) \Rightarrow (f(x) \in \mathbb{R}) \end{array}}}
\vnbprop{hbg-sep}{\vnbstmt{\begin{array}{@{}l@{}} \forall m, s, f, v, c.; \forall z, w \in s. \\ \quad \operatorname{is\text{-}normed\text{-}vector\text{-}space}(m) \wedge \\ \quad \operatorname{is\text{-}submodule}(m, s) \wedge \\ \quad \operatorname{is\text{-}linear\text{-}functional\text{-}on}(m, s, f) \wedge \\ \quad v \in \operatorname{vec}(m) \wedge \\ \quad c \in \mathbb{R} \wedge \\ \quad \forall x \in s.\; (\lvert f(x) \rvert \leq (c \cdot \operatorname{vnrm}(m)(x))) \wedge \\ \quad 0 \leq c \Rightarrow \\ \quad \quad (-f(z) - (c \cdot \operatorname{vnrm}(m)(\operatorname{vadd}(m)(z, v)))) \leq ((c \cdot \operatorname{vnrm}(m)(\operatorname{vadd}(m)(w, v))) - f(w)) \end{array}}}
\vnbprop{hbg-core}{\vnbstmt{\begin{array}{@{}l@{}} \forall m, s, f, v, c, a.; \forall y \in s, r \in \mathbb{R}. \\ \quad \operatorname{is\text{-}normed\text{-}vector\text{-}space}(m) \wedge \\ \quad \operatorname{is\text{-}submodule}(m, s) \wedge \\ \quad \operatorname{is\text{-}linear\text{-}functional\text{-}on}(m, s, f) \wedge \\ \quad v \in \operatorname{vec}(m) \wedge \\ \quad c \in \mathbb{R} \wedge \\ \quad a \in \mathbb{R} \wedge \\ \quad \forall x \in s.\; (\lvert f(x) \rvert \leq (c \cdot \operatorname{vnrm}(m)(x))) \wedge \\ \quad \forall u \in s.\; ((-f(u) - (c \cdot \operatorname{vnrm}(m)(\operatorname{vadd}(m)(u, v)))) \leq a) \wedge \\ \quad \forall u \in s.\; (a \leq ((c \cdot \operatorname{vnrm}(m)(\operatorname{vadd}(m)(u, v))) - f(u))) \Rightarrow \\ \quad \quad \lvert (f(y) + (r \cdot a)) \rvert \leq (c \cdot \operatorname{vnrm}(m)(\operatorname{vadd}(m)(y, \operatorname{act}(m)(r, v)))) \end{array}}}
\vnbprop{hb-gap-bound}{\vnbstmt{\begin{array}{@{}l@{}} \forall m, s, f, v, c. \\ \quad \operatorname{is\text{-}normed\text{-}vector\text{-}space}(m) \wedge \\ \quad \operatorname{is\text{-}submodule}(m, s) \wedge \\ \quad \operatorname{is\text{-}linear\text{-}functional\text{-}on}(m, s, f) \wedge \\ \quad v \in \operatorname{vec}(m) \wedge \\ \quad c \in \mathbb{R} \wedge \\ \quad \forall x \in s.\; (\lvert f(x) \rvert \leq (c \cdot \operatorname{vnrm}(m)(x))) \wedge \\ \quad 0 \leq c \Rightarrow \\ \quad \quad \exists a \in \mathbb{R}. \forall y \in s, r \in \mathbb{R}. \\ \quad \quad \quad \lvert (f(y) + (r \cdot a)) \rvert \leq (c \cdot \operatorname{vnrm}(m)(\operatorname{vadd}(m)(y, \operatorname{act}(m)(r, v)))) \end{array}}}
\vnbprop{hb-gap}{\vnbstmt{\begin{array}{@{}l@{}} \forall m, s, f, v. \\ \quad \operatorname{is\text{-}normed\text{-}vector\text{-}space}(m) \wedge \\ \quad \operatorname{is\text{-}submodule}(m, s) \wedge \\ \quad \operatorname{is\text{-}bounded\text{-}linear\text{-}functional\text{-}on}(m, s, f) \wedge \\ \quad v \in \operatorname{vec}(m) \wedge \\ \quad \neg (v \in s) \Rightarrow \\ \quad \quad \exists a \in \mathbb{R}. \forall y \in s, r \in \mathbb{R}. \\ \quad \quad \quad \lvert (f(y) + (r \cdot a)) \rvert \leq (\operatorname{dual\text{-}norm\text{-}on}(m, s, f) \cdot \operatorname{vnrm}(m)(\operatorname{vadd}(m)(y, \operatorname{act}(m)(r, v)))) \end{array}}}

\vnbbookfile{theorem-library/rake-hb-extend-construct}{theorem-library/rake-hb-extend-construct}
\noindent{\small rake-hb-extend-construct.scm -- hb-extend-construct (theorem-library/ hahn-banach-proof.scm:44), rake batch 7, assignment 7-C, second leaf.}\par
\vnbprop{hbx-s-in-vec}{\vnbstmt{\begin{array}{@{}l@{}} \forall m, s, p. \\ \quad (\operatorname{is\text{-}normed\text{-}vector\text{-}space}(m) \wedge (\operatorname{is\text{-}submodule}(m, s) \wedge (p \in s))) \Rightarrow (p \in \operatorname{vec}(m)) \end{array}}}
\vnbprop{hbx-span-in-vec}{\vnbstmt{\forall m, s, v.; \forall w \in \operatorname{span\text{-}add\text{-}one}(m, s, v).\; (w \in \operatorname{vec}(m))}}
\vnbprop{hbx-span-decomp}{\vnbstmt{\begin{array}{@{}l@{}} \forall m, s, v.; \forall w \in \operatorname{span\text{-}add\text{-}one}(m, s, v). \exists y \in s, r \in \mathbb{R}. \\ \quad w = \operatorname{vadd}(m)(y, \operatorname{act}(m)(r, v)) \end{array}}}
\vnbprop{hbx-span-mem-eq}{\vnbstmt{\begin{array}{@{}l@{}} \forall m, s, v, w, p, c. \\ \quad \operatorname{is\text{-}normed\text{-}vector\text{-}space}(m) \wedge \\ \quad \operatorname{is\text{-}submodule}(m, s) \wedge \\ \quad v \in \operatorname{vec}(m) \wedge \\ \quad p \in s \wedge \\ \quad c \in \mathbb{R} \wedge \\ \quad w = \operatorname{vadd}(m)(p, \operatorname{act}(m)(c, v)) \Rightarrow \\ \quad \quad w \in \operatorname{span\text{-}add\text{-}one}(m, s, v) \end{array}}}
\vnbprop{hbx-decomp-unique}{\vnbstmt{\begin{array}{@{}l@{}} \forall m, s, v, p_{1}, c_{1}, p_{2}, c_{2}. \\ \quad \operatorname{is\text{-}normed\text{-}vector\text{-}space}(m) \wedge \\ \quad \operatorname{is\text{-}submodule}(m, s) \wedge \\ \quad v \in \operatorname{vec}(m) \wedge \\ \quad \neg (v \in s) \wedge \\ \quad p_{1} \in s \wedge \\ \quad c_{1} \in \mathbb{R} \wedge \\ \quad p_{2} \in s \wedge \\ \quad c_{2} \in \mathbb{R} \wedge \\ \quad \operatorname{vadd}(m)(p_{1}, \operatorname{act}(m)(c_{1}, v)) = \operatorname{vadd}(m)(p_{2}, \operatorname{act}(m)(c_{2}, v)) \Rightarrow \\ \quad \quad (c_{1} = c_{2}) \wedge (p_{1} = p_{2}) \end{array}}}
\vnbprop{hbx-vadd-exchange}{\vnbstmt{\begin{array}{@{}l@{}} \forall m, a, b, c, d. \\ \quad (\operatorname{is\text{-}normed\text{-}vector\text{-}space}(m) \wedge ((a \in \operatorname{vec}(m)) \wedge ((b \in \operatorname{vec}(m)) \wedge ((c \in \operatorname{vec}(m)) \wedge (d \in \operatorname{vec}(m)))))) \Rightarrow  \\ \quad \quad \operatorname{vadd}(m)(\operatorname{vadd}(m)(a, b), \operatorname{vadd}(m)(c, d)) = \operatorname{vadd}(m)(\operatorname{vadd}(m)(a, c), \operatorname{vadd}(m)(b, d)) \end{array}}}
\vnbprop{hbx-add-decomp}{\vnbstmt{\begin{array}{@{}l@{}} \forall m, v, y_{1}, c_{1}, y_{2}, c_{2}. \\ \quad \operatorname{is\text{-}normed\text{-}vector\text{-}space}(m) \wedge \\ \quad v \in \operatorname{vec}(m) \wedge \\ \quad y_{1} \in \operatorname{vec}(m) \wedge \\ \quad c_{1} \in \mathbb{R} \wedge \\ \quad y_{2} \in \operatorname{vec}(m) \wedge \\ \quad c_{2} \in \mathbb{R} \Rightarrow \\ \quad \quad \operatorname{vadd}(m)(\operatorname{vadd}(m)(y_{1}, \operatorname{act}(m)(c_{1}, v)), \operatorname{vadd}(m)(y_{2}, \operatorname{act}(m)(c_{2}, v))) \\ \quad \quad =\; \operatorname{vadd}(m)(\operatorname{vadd}(m)(y_{1}, y_{2}), \operatorname{act}(m)((c_{1} + c_{2}), v)) \end{array}}}
\vnbprop{hbx-f-add}{\vnbstmt{\begin{array}{@{}l@{}} \forall m, s, f, x, y. \\ \quad (\operatorname{is\text{-}linear\text{-}functional\text{-}on}(m, s, f) \wedge ((x \in s) \wedge (y \in s))) \Rightarrow  \\ \quad \quad f(\operatorname{vadd}(m)(x, y)) = (f(x) + f(y)) \end{array}}}
\vnbprop{hbx-f-act}{\vnbstmt{\begin{array}{@{}l@{}} \forall m, s, f, r, x. \\ \quad (\operatorname{is\text{-}linear\text{-}functional\text{-}on}(m, s, f) \wedge ((r \in \mathbb{R}) \wedge (x \in s))) \Rightarrow  \\ \quad \quad f(\operatorname{act}(m)(r, x)) = (r \cdot f(x)) \end{array}}}
\vnbprop{hbx-value}{\vnbstmt{\begin{array}{@{}l@{}} \forall m, s, f, v, a, u, p, c. \\ \quad \operatorname{is\text{-}normed\text{-}vector\text{-}space}(m) \wedge \\ \quad \operatorname{is\text{-}submodule}(m, s) \wedge \\ \quad \operatorname{is\text{-}linear\text{-}functional\text{-}on}(m, s, f) \wedge \\ \quad v \in \operatorname{vec}(m) \wedge \\ \quad \neg (v \in s) \wedge \\ \quad a \in \mathbb{R} \wedge \\ \quad p \in s \wedge \\ \quad c \in \mathbb{R} \wedge \\ \quad u = \operatorname{vadd}(m)(p, \operatorname{act}(m)(c, v)) \Rightarrow \\ \quad \quad (\iota\, \mathit{tv}.\;((\mathit{tv} \in \mathbb{R}) \wedge \exists \mathit{pv} \in s, \mathit{cv} \in \mathbb{R}.\; ((u = \operatorname{vadd}(m)(\mathit{pv}, \operatorname{act}(m)(\mathit{cv}, v))) \wedge (\mathit{tv} = (f(\mathit{pv}) + (\mathit{cv} \cdot a)))))) \\ \quad \quad =\; f(p) + (c \cdot a) \end{array}}}
\vnbprop{hb-extend-construct}{\vnbstmt{\begin{array}{@{}l@{}} \forall m, s, f, v, a. \\ \quad \operatorname{is\text{-}normed\text{-}vector\text{-}space}(m) \wedge \\ \quad \operatorname{is\text{-}submodule}(m, s) \wedge \\ \quad \operatorname{is\text{-}linear\text{-}functional\text{-}on}(m, s, f) \wedge \\ \quad v \in \operatorname{vec}(m) \wedge \\ \quad \neg (v \in s) \wedge \\ \quad a \in \mathbb{R} \Rightarrow \\ \quad \quad \exists g. \\ \quad \quad \quad \operatorname{is\text{-}linear\text{-}functional\text{-}on}(m, \operatorname{span\text{-}add\text{-}one}(m, s, v), g) \wedge \\ \quad \quad \quad \operatorname{extends\text{-}on}(s, g, f) \wedge \\ \quad \quad \quad \forall y \in s, r \in \mathbb{R}.\; (g(\operatorname{vadd}(m)(y, \operatorname{act}(m)(r, v))) = (f(y) + (r \cdot a))) \end{array}}}

\vnbbookfile{theorem-library/rake-line-functional}{theorem-library/rake-line-functional}
\noindent{\small rake-line-functional.scm -- line-functional-exists (theorem-library/ norm-as-sup-proof.scm:43), rake batch 7, assignment 7-C.}\par
\vnbprop{lfe-line-membership}{\vnbstmt{\begin{array}{@{}l@{}} \forall m, v, w. \\ \quad (w \in \operatorname{line}(m, v)) \Leftrightarrow ((w \in \operatorname{vec}(m)) \wedge \exists r \in \mathbb{R}.\; (w = \operatorname{act}(m)(r, v))) \end{array}}}
\vnbprop{lfe-line-in-vec}{\vnbstmt{\forall m, v.; \forall w \in \operatorname{line}(m, v).\; (w \in \operatorname{vec}(m))}}
\vnbprop{lfe-line-subset}{\vnbstmt{\forall m, v.\; (\operatorname{is\text{-}normed\text{-}vector\text{-}space}(m) \Rightarrow (\operatorname{line}(m, v) \subseteq \operatorname{vec}(m)))}}
\vnbprop{lfe-line-coef}{\vnbstmt{\forall m, v.; \forall w \in \operatorname{line}(m, v).\; \exists r \in \mathbb{R}.\; (w = \operatorname{act}(m)(r, v))}}
\vnbprop{lfe-line-mem-eq}{\vnbstmt{\begin{array}{@{}l@{}} \forall m, v, w, r. \\ \quad (\operatorname{is\text{-}normed\text{-}vector\text{-}space}(m) \wedge ((v \in \operatorname{vec}(m)) \wedge ((r \in \mathbb{R}) \wedge (w = \operatorname{act}(m)(r, v))))) \Rightarrow  \\ \quad \quad w \in \operatorname{line}(m, v) \end{array}}}
\vnbprop{lfe-coef-unique}{\vnbstmt{\begin{array}{@{}l@{}} \forall m, v.; \forall a, b \in \mathbb{R}. \\ \quad (\operatorname{is\text{-}normed\text{-}vector\text{-}space}(m) \wedge ((v \in \operatorname{vec}(m)) \wedge (\operatorname{act}(m)(a, v) = \operatorname{act}(m)(b, v)))) \Rightarrow  \\ \quad \quad (a \cdot \operatorname{vnrm}(m)(v)) = (b \cdot \operatorname{vnrm}(m)(v)) \end{array}}}
\vnbprop{lfe-coef-value}{\vnbstmt{\begin{array}{@{}l@{}} \forall m, v, u, c. \\ \quad (\operatorname{is\text{-}normed\text{-}vector\text{-}space}(m) \wedge ((v \in \operatorname{vec}(m)) \wedge ((c \in \mathbb{R}) \wedge (u = \operatorname{act}(m)(c, v))))) \Rightarrow  \\ \quad \quad (\iota\, t.\;((t \in \mathbb{R}) \wedge \exists q \in \mathbb{R}.\; ((u = \operatorname{act}(m)(q, v)) \wedge (t = (q \cdot \operatorname{vnrm}(m)(v)))))) \\ \quad \quad =\; c \cdot \operatorname{vnrm}(m)(v) \end{array}}}
\vnbprop{line-functional-exists}{\vnbstmt{\begin{array}{@{}l@{}} \forall m, v. \\ \quad (\operatorname{is\text{-}normed\text{-}vector\text{-}space}(m) \wedge (v \in \operatorname{vec}(m))) \Rightarrow  \\ \quad \quad \exists f. \\ \quad \quad \quad \operatorname{is\text{-}bounded\text{-}linear\text{-}functional\text{-}on}(m, \operatorname{line}(m, v), f) \wedge \\ \quad \quad \quad \operatorname{dual\text{-}norm\text{-}on}(m, \operatorname{line}(m, v), f) \leq 1 \wedge \\ \quad \quad \quad f(v) = \operatorname{vnrm}(m)(v) \end{array}}}

\vnbbookfile{theorem-library/hahn-banach-proof}{theorem-library/hahn-banach-proof}
\noindent{\small hahn-banach-proof.scm -- Hahn-Banach, the one-dimension-higher extension step (real normed vector space), MACHINE-PROVEN modulo the warranted core.}\par
\vnbprop{hahn-banach-extend-one}{\vnbstmt{\begin{array}{@{}l@{}} \forall m, s, f, v. \\ \quad \operatorname{is\text{-}normed\text{-}vector\text{-}space}(m) \wedge \\ \quad \operatorname{is\text{-}submodule}(m, s) \wedge \\ \quad \operatorname{is\text{-}bounded\text{-}linear\text{-}functional\text{-}on}(m, s, f) \wedge \\ \quad v \in \operatorname{vec}(m) \wedge \\ \quad \neg (v \in s) \Rightarrow \\ \quad \quad \exists g. \\ \quad \quad \quad \operatorname{is\text{-}linear\text{-}functional\text{-}on}(m, \operatorname{span\text{-}add\text{-}one}(m, s, v), g) \wedge \\ \quad \quad \quad \operatorname{extends\text{-}on}(s, g, f) \wedge \\ \quad \quad \quad \forall w \in \operatorname{span\text{-}add\text{-}one}(m, s, v). \\ \quad \quad \quad \quad \lvert g(w) \rvert \leq (\operatorname{dual\text{-}norm\text{-}on}(m, s, f) \cdot \operatorname{vnrm}(m)(w)) \end{array}}}

\vnbbookfile{theorem-library/noetherian-maximal-proof}{theorem-library/noetherian-maximal-proof}
\noindent{\small noetherian-maximal-proof.scm -- the noetherian maximal-element principle, MACHINE-PROVEN, and its specialisation hb-good-has-maximal (the input that makes finite-dimensional Hahn-Banach terminate).}\par
\vnbprop{noetherian-set-has-maximal}{\vnbstmt{\begin{array}{@{}l@{}} \forall m.; \forall g \in \mathbf{Set}. \\ \quad \operatorname{is\text{-}noetherian}(m) \wedge \\ \quad g \subseteq \operatorname{power}(\operatorname{vec}(m)) \wedge \\ \quad \forall t \in g.\; \operatorname{is\text{-}submodule}(m, t) \wedge \\ \quad \exists a.\; (a \in g) \Rightarrow \\ \quad \quad \exists t \in g.\; \forall u \in g.\; ((t \subseteq u) \Rightarrow (t = u)) \end{array}}}
\vnbprop{hb-good-has-maximal}{\vnbstmt{\begin{array}{@{}l@{}} \forall m, s, f. \\ \quad \operatorname{is\text{-}normed\text{-}vector\text{-}space}(m) \wedge \\ \quad \operatorname{is\text{-}finite\text{-}dimensional}(\operatorname{normed\text{-}vector\text{-}space\text{-}as\text{-}module}(m)) \wedge \\ \quad \operatorname{is\text{-}submodule}(m, s) \wedge \\ \quad \operatorname{is\text{-}bounded\text{-}linear\text{-}functional\text{-}on}(m, s, f) \Rightarrow \\ \quad \quad \exists t. \\ \quad \quad \quad (\operatorname{good\text{-}sub}(m, s, f, t) \wedge \forall u.\; ((\operatorname{good\text{-}sub}(m, s, f, u) \wedge (t \subseteq u)) \Rightarrow (t = u))) \end{array}}}

\vnbbookfile{theorem-library/hahn-banach-full-proof}{theorem-library/hahn-banach-full-proof}
\noindent{\small RETIRED 2026-09-17 (proven): le-bound-mono (was an add-to-pss here) -- theorem-library/rake-analysis2.scm hahn-banach-full-proof.scm -- FULL Hahn-Banach for a finite-dimensional real normed vector space, MACHINE-PROVEN modulo a warranted plumbing core.}\par
\vnbprop{good-step}{\vnbstmt{\begin{array}{@{}l@{}} \forall m, s, f, t, x. \\ \quad \operatorname{is\text{-}normed\text{-}vector\text{-}space}(m) \wedge \\ \quad \operatorname{is\text{-}bounded\text{-}linear\text{-}functional\text{-}on}(m, s, f) \wedge \\ \quad \operatorname{good\text{-}sub}(m, s, f, t) \wedge \\ \quad x \in \operatorname{vec}(m) \wedge \\ \quad \neg (x \in t) \Rightarrow \\ \quad \quad \operatorname{good\text{-}sub}(m, s, f, \operatorname{span\text{-}add\text{-}one}(m, t, x)) \end{array}}}
\vnbprop{hahn-banach}{\vnbstmt{\begin{array}{@{}l@{}} \forall m, s, f. \\ \quad \operatorname{is\text{-}normed\text{-}vector\text{-}space}(m) \wedge \\ \quad \operatorname{is\text{-}finite\text{-}dimensional}(\operatorname{normed\text{-}vector\text{-}space\text{-}as\text{-}module}(m)) \wedge \\ \quad \operatorname{is\text{-}submodule}(m, s) \wedge \\ \quad \operatorname{is\text{-}bounded\text{-}linear\text{-}functional\text{-}on}(m, s, f) \Rightarrow \\ \quad \quad \exists g. \\ \quad \quad \quad \operatorname{is\text{-}linear\text{-}functional\text{-}on}(m, \operatorname{vec}(m), g) \wedge \\ \quad \quad \quad \operatorname{extends\text{-}on}(s, g, f) \wedge \\ \quad \quad \quad \forall w \in \operatorname{vec}(m).\; (\lvert g(w) \rvert \leq (\operatorname{dual\text{-}norm\text{-}on}(m, s, f) \cdot \operatorname{vnrm}(m)(w))) \end{array}}}

\vnbbookfile{structure-library/hahn-banach-sublinear}{structure-library/hahn-banach-sublinear}
\noindent{\small structure-library/hahn-banach-sublinear.scm -- the vocabulary of the Hahn-Banach theorem for a SUBLINEAR functional (RM-3A, 2026-10-10). DEFINITIONS ONLY; every law is PROVEN in theorem-library/hahn-banach-sublinear.scm (prefix hbs-).}\par
\par\medskip\noindent\textbf{Definition} (\texttt{is-sublinear}).\;For all $m$ and $p$, 
\begin{equation*}
\begin{array}{@{}l@{}} \operatorname{is\text{-}sublinear}(m, p) \\ \Leftrightarrow\; p \in (\operatorname{vec}(m) \to \mathbb{R}) \wedge \\ \forall \operatorname{hsx\_}, \operatorname{hsy\_} \in \operatorname{vec}(m). \\ \quad p(\operatorname{vadd}(m)(\operatorname{hsx\_}, \operatorname{hsy\_})) \leq +(p(\operatorname{hsx\_}), p(\operatorname{hsy\_})) \wedge \\ \forall \operatorname{hst\_} \in \mathbb{R}. \\ \quad (0 \leq \operatorname{hst\_}) \Rightarrow  \\ \quad \quad \forall \operatorname{hsx\_} \in \operatorname{vec}(m).\; (p(\operatorname{act}(m)(\operatorname{hst\_}, \operatorname{hsx\_})) = \cdot(\operatorname{hst\_}, p(\operatorname{hsx\_}))) \end{array}
\end{equation*}


\vnbbookfile{theorem-library/hahn-banach-sublinear}{theorem-library/hahn-banach-sublinear}
\noindent{\small theorem-library/hahn-banach-sublinear.scm -- HAHN-BANACH FOR A SUBLINEAR FUNCTIONAL on a real vector space (RM-3A, 2026-10-10).  Prefix hbs-.}\par
\vnbprop{hbs-carr}{\vnbstmt{\begin{array}{@{}l@{}} \forall m. \\ \quad (\operatorname{is\text{-}module}(m) \wedge (\operatorname{scal}(m) = \operatorname{normed\text{-}field\text{-}as\text{-}commutative\text{-}ring}(\operatorname{rr\text{-}normed\text{-}field}))) \Rightarrow  \\ \quad \quad \operatorname{carr}(\operatorname{scal}(m)) = \mathbb{R} \end{array}}}
\vnbprop{hbs-in-carr}{\vnbstmt{\begin{array}{@{}l@{}} \forall m.; \forall r \in \mathbb{R}. \\ \quad (\operatorname{is\text{-}module}(m) \wedge (\operatorname{scal}(m) = \operatorname{normed\text{-}field\text{-}as\text{-}commutative\text{-}ring}(\operatorname{rr\text{-}normed\text{-}field}))) \Rightarrow  \\ \quad \quad r \in \operatorname{carr}(\operatorname{scal}(m)) \end{array}}}
\vnbprop{hbs-add-eq}{\vnbstmt{\begin{array}{@{}l@{}} \forall m.; \forall r, s \in \mathbb{R}. \\ \quad (\operatorname{is\text{-}module}(m) \wedge (\operatorname{scal}(m) = \operatorname{normed\text{-}field\text{-}as\text{-}commutative\text{-}ring}(\operatorname{rr\text{-}normed\text{-}field}))) \Rightarrow  \\ \quad \quad \operatorname{add}(\operatorname{scal}(m))(r, s) = (r + s) \end{array}}}
\vnbprop{hbs-mul-eq}{\vnbstmt{\begin{array}{@{}l@{}} \forall m.; \forall r, s \in \mathbb{R}. \\ \quad (\operatorname{is\text{-}module}(m) \wedge (\operatorname{scal}(m) = \operatorname{normed\text{-}field\text{-}as\text{-}commutative\text{-}ring}(\operatorname{rr\text{-}normed\text{-}field}))) \Rightarrow  \\ \quad \quad \operatorname{mul}(\operatorname{scal}(m))(r, s) = (r \cdot s) \end{array}}}
\vnbprop{hbs-one-eq}{\vnbstmt{\begin{array}{@{}l@{}} \forall m. \\ \quad (\operatorname{is\text{-}module}(m) \wedge (\operatorname{scal}(m) = \operatorname{normed\text{-}field\text{-}as\text{-}commutative\text{-}ring}(\operatorname{rr\text{-}normed\text{-}field}))) \Rightarrow  \\ \quad \quad \operatorname{one}(\operatorname{scal}(m)) = 1 \end{array}}}
\vnbprop{hbs-zero-eq}{\vnbstmt{\begin{array}{@{}l@{}} \forall m. \\ \quad (\operatorname{is\text{-}module}(m) \wedge (\operatorname{scal}(m) = \operatorname{normed\text{-}field\text{-}as\text{-}commutative\text{-}ring}(\operatorname{rr\text{-}normed\text{-}field}))) \Rightarrow  \\ \quad \quad \operatorname{zero}(\operatorname{scal}(m)) = 0 \end{array}}}
\vnbprop{hbs-act-in}{\vnbstmt{\begin{array}{@{}l@{}} \forall m.; \forall r \in \mathbb{R}, x \in \operatorname{vec}(m). \\ \quad (\operatorname{is\text{-}module}(m) \wedge (\operatorname{scal}(m) = \operatorname{normed\text{-}field\text{-}as\text{-}commutative\text{-}ring}(\operatorname{rr\text{-}normed\text{-}field}))) \Rightarrow  \\ \quad \quad \operatorname{act}(m)(r, x) \in \operatorname{vec}(m) \end{array}}}
\vnbprop{hbs-act-add}{\vnbstmt{\begin{array}{@{}l@{}} \forall m.; \forall r, s \in \mathbb{R}, x \in \operatorname{vec}(m). \\ \quad (\operatorname{is\text{-}module}(m) \wedge (\operatorname{scal}(m) = \operatorname{normed\text{-}field\text{-}as\text{-}commutative\text{-}ring}(\operatorname{rr\text{-}normed\text{-}field}))) \Rightarrow  \\ \quad \quad \operatorname{act}(m)((r + s), x) = \operatorname{vadd}(m)(\operatorname{act}(m)(r, x), \operatorname{act}(m)(s, x)) \end{array}}}
\vnbprop{hbs-act-mul}{\vnbstmt{\begin{array}{@{}l@{}} \forall m.; \forall r, s \in \mathbb{R}, x \in \operatorname{vec}(m). \\ \quad (\operatorname{is\text{-}module}(m) \wedge (\operatorname{scal}(m) = \operatorname{normed\text{-}field\text{-}as\text{-}commutative\text{-}ring}(\operatorname{rr\text{-}normed\text{-}field}))) \Rightarrow  \\ \quad \quad \operatorname{act}(m)((r \cdot s), x) = \operatorname{act}(m)(r, \operatorname{act}(m)(s, x)) \end{array}}}
\vnbprop{hbs-act-one}{\vnbstmt{\begin{array}{@{}l@{}} \forall m.; \forall x \in \operatorname{vec}(m). \\ \quad (\operatorname{is\text{-}module}(m) \wedge (\operatorname{scal}(m) = \operatorname{normed\text{-}field\text{-}as\text{-}commutative\text{-}ring}(\operatorname{rr\text{-}normed\text{-}field}))) \Rightarrow  \\ \quad \quad \operatorname{act}(m)(1, x) = x \end{array}}}
\vnbprop{hbs-act-zero}{\vnbstmt{\begin{array}{@{}l@{}} \forall m.; \forall x \in \operatorname{vec}(m). \\ \quad (\operatorname{is\text{-}module}(m) \wedge (\operatorname{scal}(m) = \operatorname{normed\text{-}field\text{-}as\text{-}commutative\text{-}ring}(\operatorname{rr\text{-}normed\text{-}field}))) \Rightarrow  \\ \quad \quad \operatorname{act}(m)(0, x) = \operatorname{vzero}(m) \end{array}}}
\vnbprop{hbs-act-vadd}{\vnbstmt{\begin{array}{@{}l@{}} \forall m.; \forall r \in \mathbb{R}, x, y \in \operatorname{vec}(m). \\ \quad (\operatorname{is\text{-}module}(m) \wedge (\operatorname{scal}(m) = \operatorname{normed\text{-}field\text{-}as\text{-}commutative\text{-}ring}(\operatorname{rr\text{-}normed\text{-}field}))) \Rightarrow  \\ \quad \quad \operatorname{act}(m)(r, \operatorname{vadd}(m)(x, y)) = \operatorname{vadd}(m)(\operatorname{act}(m)(r, x), \operatorname{act}(m)(r, y)) \end{array}}}
\vnbprop{hbs-neg-cancel}{\vnbstmt{\begin{array}{@{}l@{}} \forall m.; \forall x \in \operatorname{vec}(m). \\ \quad (\operatorname{is\text{-}module}(m) \wedge (\operatorname{scal}(m) = \operatorname{normed\text{-}field\text{-}as\text{-}commutative\text{-}ring}(\operatorname{rr\text{-}normed\text{-}field}))) \Rightarrow  \\ \quad \quad \operatorname{vadd}(m)(x, \operatorname{act}(m)(-1, x)) = \operatorname{vzero}(m) \end{array}}}
\vnbprop{hbs-neg-cancel-l}{\vnbstmt{\begin{array}{@{}l@{}} \forall m.; \forall x \in \operatorname{vec}(m). \\ \quad (\operatorname{is\text{-}module}(m) \wedge (\operatorname{scal}(m) = \operatorname{normed\text{-}field\text{-}as\text{-}commutative\text{-}ring}(\operatorname{rr\text{-}normed\text{-}field}))) \Rightarrow  \\ \quad \quad \operatorname{vadd}(m)(\operatorname{act}(m)(-1, x), x) = \operatorname{vzero}(m) \end{array}}}
\vnbprop{hbs-cancel-r}{\vnbstmt{\begin{array}{@{}l@{}} \forall m.; \forall x, y \in \operatorname{vec}(m). \\ \quad (\operatorname{is\text{-}module}(m) \wedge (\operatorname{scal}(m) = \operatorname{normed\text{-}field\text{-}as\text{-}commutative\text{-}ring}(\operatorname{rr\text{-}normed\text{-}field}))) \Rightarrow  \\ \quad \quad \operatorname{vadd}(m)(\operatorname{vadd}(m)(x, y), \operatorname{act}(m)(-1, y)) = x \end{array}}}
\vnbprop{hbs-cancel-l}{\vnbstmt{\begin{array}{@{}l@{}} \forall m.; \forall x, y \in \operatorname{vec}(m). \\ \quad (\operatorname{is\text{-}module}(m) \wedge (\operatorname{scal}(m) = \operatorname{normed\text{-}field\text{-}as\text{-}commutative\text{-}ring}(\operatorname{rr\text{-}normed\text{-}field}))) \Rightarrow  \\ \quad \quad \operatorname{vadd}(m)(\operatorname{vadd}(m)(x, y), \operatorname{act}(m)(-1, x)) = y \end{array}}}
\vnbprop{hbs-inf-point}{\vnbstmt{\begin{array}{@{}l@{}} \forall i, a, x. \\ \quad \exists i_{0}.\; (i_{0} \in i) \wedge \\ \quad \forall j \in i.\; (a(x, j) \in \mathbb{R}) \wedge \\ \quad \exists b \in \mathbb{R}.\; \forall j \in i.\; (b \leq a(x, j)) \Rightarrow \\ \quad \quad -\operatorname{sup}(\{w \in \mathbb{R} : \exists k \in i.\; (w = -a(x, k))\}) \in \mathbb{R} \wedge \\ \quad \quad \forall j \in i.\; (-\operatorname{sup}(\{w \in \mathbb{R} : \exists k \in i.\; (w = -a(x, k))\}) \leq a(x, j)) \wedge \\ \quad \quad \forall b \in \mathbb{R}. \\ \quad \quad \quad \forall j \in i.\; (b \leq a(x, j)) \Rightarrow  \\ \quad \quad \quad \quad b \leq -\operatorname{sup}(\{w \in \mathbb{R} : \exists k \in i.\; (w = -a(x, k))\}) \end{array}}}
\vnbprop{hbs-inf-family}{\vnbstmt{\begin{array}{@{}l@{}} \forall d, f, i. \\ \quad d \in \mathbf{Set} \wedge \\ \quad \forall x \in d.\; \exists i_{0}.\; (i_{0} \in f(x)) \wedge \\ \quad \forall x \in d, j \in f(x).\; (i(x, j) \in \mathbb{R}) \wedge \\ \quad \forall x \in d.\; \exists b \in \mathbb{R}.\; \forall j \in f(x).\; (b \leq i(x, j)) \Rightarrow \\ \quad \quad \exists r \in (d \to \mathbb{R}). \\ \quad \quad \quad \forall x \in d, j \in f(x).\; (r(x) \leq i(x, j)) \wedge \\ \quad \quad \quad \forall x \in d, b \in \mathbb{R}.\; (\forall j \in f(x).\; (b \leq i(x, j)) \Rightarrow (b \leq r(x))) \end{array}}}
\vnbprop{hbs-inf-family-c}{\vnbstmt{\begin{array}{@{}l@{}} \forall d, i, a. \\ \quad d \in \mathbf{Set} \wedge \\ \quad \forall x \in d.\; \exists i_{0}.\; (i_{0} \in i) \wedge \\ \quad \forall x \in d, j \in i.\; (a(x, j) \in \mathbb{R}) \wedge \\ \quad \forall x \in d.\; \exists b \in \mathbb{R}.\; \forall j \in i.\; (b \leq a(x, j)) \Rightarrow \\ \quad \quad \exists r \in (d \to \mathbb{R}). \\ \quad \quad \quad \forall x \in d, j \in i.\; (r(x) \leq a(x, j)) \wedge \\ \quad \quad \quad \forall x \in d, b \in \mathbb{R}.\; (\forall j \in i.\; (b \leq a(x, j)) \Rightarrow (b \leq r(x))) \end{array}}}
\vnbprop{hbs-pos-inv}{\vnbstmt{\forall t \in \mathbb{R}.\; ((0 < t) \Rightarrow \exists u \in \mathbb{R}.\; ((0 < u) \wedge ((t \cdot u) = 1)))}}
\vnbprop{hbs-rr-div-le}{\vnbstmt{\begin{array}{@{}l@{}} \forall t, u, c, a \in \mathbb{R}. \\ \quad ((0 < t) \wedge ((0 < u) \wedge (((t \cdot u) = 1) \wedge (c \leq (t \cdot a))))) \Rightarrow ((u \cdot c) \leq a) \end{array}}}
\vnbprop{hbs-rr-mul-back}{\vnbstmt{\begin{array}{@{}l@{}} \forall t, u, c, a \in \mathbb{R}. \\ \quad ((0 < t) \wedge ((0 < u) \wedge (((t \cdot u) = 1) \wedge ((u \cdot c) \leq a)))) \Rightarrow (c \leq (t \cdot a)) \end{array}}}
\vnbprop{hbs-sl-fun}{\vnbstmt{\forall m, q.\; (\operatorname{is\text{-}sublinear}(m, q) \Rightarrow (q \in (\operatorname{vec}(m) \to \mathbb{R})))}}
\vnbprop{hbs-sl-sub}{\vnbstmt{\begin{array}{@{}l@{}} \forall m, q.; \forall x, y \in \operatorname{vec}(m). \\ \quad \operatorname{is\text{-}sublinear}(m, q) \Rightarrow (q(\operatorname{vadd}(m)(x, y)) \leq (q(x) + q(y))) \end{array}}}
\vnbprop{hbs-sl-hom}{\vnbstmt{\begin{array}{@{}l@{}} \forall m, q.; \forall t \in \mathbb{R}, x \in \operatorname{vec}(m). \\ \quad (\operatorname{is\text{-}sublinear}(m, q) \wedge (0 \leq t)) \Rightarrow (q(\operatorname{act}(m)(t, x)) = (t \cdot q(x))) \end{array}}}
\vnbprop{hbs-sl-val}{\vnbstmt{\forall m, q.; \forall x \in \operatorname{vec}(m).\; (\operatorname{is\text{-}sublinear}(m, q) \Rightarrow (q(x) \in \mathbb{R}))}}
\vnbprop{hbs-sl-zero}{\vnbstmt{\begin{array}{@{}l@{}} \forall m, q. \\ \quad (\operatorname{is\text{-}module}(m) \wedge ((\operatorname{scal}(m) = \operatorname{normed\text{-}field\text{-}as\text{-}commutative\text{-}ring}(\operatorname{rr\text{-}normed\text{-}field})) \wedge \operatorname{is\text{-}sublinear}(m, q))) \Rightarrow  \\ \quad \quad q(\operatorname{vzero}(m)) = 0 \end{array}}}
\vnbprop{hbs-sl-negsum}{\vnbstmt{\begin{array}{@{}l@{}} \forall m, q.; \forall x \in \operatorname{vec}(m). \\ \quad (\operatorname{is\text{-}module}(m) \wedge ((\operatorname{scal}(m) = \operatorname{normed\text{-}field\text{-}as\text{-}commutative\text{-}ring}(\operatorname{rr\text{-}normed\text{-}field})) \wedge \operatorname{is\text{-}sublinear}(m, q))) \Rightarrow  \\ \quad \quad 0 \leq (q(x) + q(\operatorname{act}(m)(-1, x))) \end{array}}}
\vnbprop{hbs-sl-intro-pos}{\vnbstmt{\begin{array}{@{}l@{}} \forall m.; \forall q \in (\operatorname{vec}(m) \to \mathbb{R}). \\ \quad \operatorname{is\text{-}module}(m) \wedge \\ \quad \operatorname{scal}(m) = \operatorname{normed\text{-}field\text{-}as\text{-}commutative\text{-}ring}(\operatorname{rr\text{-}normed\text{-}field}) \wedge \\ \quad \forall x, y \in \operatorname{vec}(m).\; (q(\operatorname{vadd}(m)(x, y)) \leq (q(x) + q(y))) \wedge \\ \quad \forall t \in \mathbb{R}, x \in \operatorname{vec}(m).\; ((0 < t) \Rightarrow (q(\operatorname{act}(m)(t, x)) \leq (t \cdot q(x)))) \wedge \\ \quad q(\operatorname{vzero}(m)) = 0 \Rightarrow \\ \quad \quad \operatorname{is\text{-}sublinear}(m, q) \end{array}}}
\vnbprop{hbs-sl-odd-linear}{\vnbstmt{\begin{array}{@{}l@{}} \forall m, q. \\ \quad \operatorname{is\text{-}module}(m) \wedge \\ \quad \operatorname{scal}(m) = \operatorname{normed\text{-}field\text{-}as\text{-}commutative\text{-}ring}(\operatorname{rr\text{-}normed\text{-}field}) \wedge \\ \quad \operatorname{is\text{-}sublinear}(m, q) \wedge \\ \quad \forall y \in \operatorname{vec}(m).\; (q(\operatorname{act}(m)(-1, y)) \leq -q(y)) \Rightarrow \\ \quad \quad \operatorname{is\text{-}linear\text{-}functional\text{-}on}(m, \operatorname{vec}(m), q) \end{array}}}
\vnbprop{hbs-tilt-exists}{\vnbstmt{\begin{array}{@{}l@{}} \forall m, q.; \forall y \in \operatorname{vec}(m). \\ \quad (\operatorname{is\text{-}module}(m) \wedge ((\operatorname{scal}(m) = \operatorname{normed\text{-}field\text{-}as\text{-}commutative\text{-}ring}(\operatorname{rr\text{-}normed\text{-}field})) \wedge \operatorname{is\text{-}sublinear}(m, q))) \Rightarrow  \\ \quad \quad \exists r \in (\operatorname{vec}(m) \to \mathbb{R}). \\ \quad \quad \quad \forall x \in \operatorname{vec}(m), t \in \mathbb{R}. \\ \quad \quad \quad \quad (0 \leq t) \Rightarrow (r(x) \leq (q(\operatorname{vadd}(m)(x, \operatorname{act}(m)(t, y))) - (t \cdot q(y)))) \wedge \\ \quad \quad \quad \forall x \in \operatorname{vec}(m), b \in \mathbb{R}. \\ \quad \quad \quad \quad \forall t \in \mathbb{R}.\; ((0 \leq t) \Rightarrow (b \leq (q(\operatorname{vadd}(m)(x, \operatorname{act}(m)(t, y))) - (t \cdot q(y))))) \Rightarrow  \\ \quad \quad \quad \quad \quad b \leq r(x) \end{array}}}
\vnbprop{hbs-tilt-props}{\vnbstmt{\begin{array}{@{}l@{}} \forall m, q.; \forall y \in \operatorname{vec}(m), t \in (\operatorname{vec}(m) \to \mathbb{R}). \\ \quad \operatorname{is\text{-}module}(m) \wedge \\ \quad \operatorname{scal}(m) = \operatorname{normed\text{-}field\text{-}as\text{-}commutative\text{-}ring}(\operatorname{rr\text{-}normed\text{-}field}) \wedge \\ \quad \operatorname{is\text{-}sublinear}(m, q) \wedge \\ \quad \forall x \in \operatorname{vec}(m), a \in \mathbb{R}. \\ \quad \quad (0 \leq a) \Rightarrow (t(x) \leq (q(\operatorname{vadd}(m)(x, \operatorname{act}(m)(a, y))) - (a \cdot q(y)))) \wedge \\ \quad \forall x \in \operatorname{vec}(m), b \in \mathbb{R}. \\ \quad \quad \forall a \in \mathbb{R}.\; ((0 \leq a) \Rightarrow (b \leq (q(\operatorname{vadd}(m)(x, \operatorname{act}(m)(a, y))) - (a \cdot q(y))))) \Rightarrow  \\ \quad \quad \quad b \leq t(x) \Rightarrow \\ \quad \quad (\operatorname{is\text{-}sublinear}(m, t) \wedge (\forall x \in \operatorname{vec}(m).\; (t(x) \leq q(x)) \wedge (t(\operatorname{act}(m)(-1, y)) \leq -q(y)))) \end{array}}}
\vnbprop{hbs-chain-exists}{\vnbstmt{\begin{array}{@{}l@{}} \forall m, h. \\ \quad \operatorname{is\text{-}module}(m) \wedge \\ \quad \operatorname{scal}(m) = \operatorname{normed\text{-}field\text{-}as\text{-}commutative\text{-}ring}(\operatorname{rr\text{-}normed\text{-}field}) \wedge \\ \quad \exists q_{0}.\; (q_{0} \in h) \wedge \\ \quad \forall q \in h.\; (q \in (\operatorname{vec}(m) \to \mathbb{R})) \wedge \\ \quad \forall x \in \operatorname{vec}(m).\; \exists b \in \mathbb{R}.\; \forall q \in h.\; (b \leq q(x)) \Rightarrow \\ \quad \quad \exists t \in (\operatorname{vec}(m) \to \mathbb{R}). \\ \quad \quad \quad \forall x \in \operatorname{vec}(m), q \in h.\; (t(x) \leq q(x)) \wedge \\ \quad \quad \quad \forall x \in \operatorname{vec}(m), b \in \mathbb{R}.\; (\forall q \in h.\; (b \leq q(x)) \Rightarrow (b \leq t(x))) \end{array}}}
\vnbprop{hbs-chain-props}{\vnbstmt{\begin{array}{@{}l@{}} \forall m, h.; \forall t \in (\operatorname{vec}(m) \to \mathbb{R}). \\ \quad \operatorname{is\text{-}module}(m) \wedge \\ \quad \operatorname{scal}(m) = \operatorname{normed\text{-}field\text{-}as\text{-}commutative\text{-}ring}(\operatorname{rr\text{-}normed\text{-}field}) \wedge \\ \quad \exists q_{0}.\; (q_{0} \in h) \wedge \\ \quad \forall q \in h.\; \operatorname{is\text{-}sublinear}(m, q) \wedge \\ \quad \forall q_{1}, q_{2} \in h. \\ \quad \quad \forall x \in \operatorname{vec}(m).\; (\operatorname{q1}(x) \leq \operatorname{q2}(x)) \vee \forall x \in \operatorname{vec}(m).\; (\operatorname{q2}(x) \leq \operatorname{q1}(x)) \wedge \\ \quad \forall x \in \operatorname{vec}(m), q \in h.\; (t(x) \leq q(x)) \wedge \\ \quad \forall x \in \operatorname{vec}(m), b \in \mathbb{R}.\; (\forall q \in h.\; (b \leq q(x)) \Rightarrow (b \leq t(x))) \Rightarrow \\ \quad \quad \operatorname{is\text{-}sublinear}(m, t) \end{array}}}
\vnbprop{hbs-lf-fun}{\vnbstmt{\forall m, s, f.\; (\operatorname{is\text{-}linear\text{-}functional\text{-}on}(m, s, f) \Rightarrow (f \in (s \to \mathbb{R})))}}
\vnbprop{hbs-lf-add}{\vnbstmt{\begin{array}{@{}l@{}} \forall m, s, f.; \forall x, y \in s. \\ \quad \operatorname{is\text{-}linear\text{-}functional\text{-}on}(m, s, f) \Rightarrow (f(\operatorname{vadd}(m)(x, y)) = (f(x) + f(y))) \end{array}}}
\vnbprop{hbs-lf-hom}{\vnbstmt{\begin{array}{@{}l@{}} \forall m, s, f.; \forall r \in \mathbb{R}, x \in s. \\ \quad \operatorname{is\text{-}linear\text{-}functional\text{-}on}(m, s, f) \Rightarrow (f(\operatorname{act}(m)(r, x)) = (r \cdot f(x))) \end{array}}}
\vnbprop{hbs-lf-val}{\vnbstmt{\forall m, s, f.; \forall x \in s.\; (\operatorname{is\text{-}linear\text{-}functional\text{-}on}(m, s, f) \Rightarrow (f(x) \in \mathbb{R}))}}
\vnbprop{hbs-sub-act}{\vnbstmt{\begin{array}{@{}l@{}} \forall m, s.; \forall r \in \mathbb{R}, x \in s. \\ \quad (\operatorname{is\text{-}module}(m) \wedge ((\operatorname{scal}(m) = \operatorname{normed\text{-}field\text{-}as\text{-}commutative\text{-}ring}(\operatorname{rr\text{-}normed\text{-}field})) \wedge \operatorname{is\text{-}submodule}(m, s))) \Rightarrow  \\ \quad \quad \operatorname{act}(m)(r, x) \in s \end{array}}}
\vnbprop{hbs-lf-zero}{\vnbstmt{\begin{array}{@{}l@{}} \forall m, s, f. \\ \quad (\operatorname{is\text{-}module}(m) \wedge (\operatorname{is\text{-}submodule}(m, s) \wedge \operatorname{is\text{-}linear\text{-}functional\text{-}on}(m, s, f))) \Rightarrow  \\ \quad \quad f(\operatorname{vzero}(m)) = 0 \end{array}}}
\vnbprop{hbs-red-exists}{\vnbstmt{\begin{array}{@{}l@{}} \forall m, p, s, f. \\ \quad \operatorname{is\text{-}module}(m) \wedge \\ \quad \operatorname{scal}(m) = \operatorname{normed\text{-}field\text{-}as\text{-}commutative\text{-}ring}(\operatorname{rr\text{-}normed\text{-}field}) \wedge \\ \quad \operatorname{is\text{-}sublinear}(m, p) \wedge \\ \quad \operatorname{is\text{-}submodule}(m, s) \wedge \\ \quad \operatorname{is\text{-}linear\text{-}functional\text{-}on}(m, s, f) \wedge \\ \quad \forall z \in s.\; (f(z) \leq p(z)) \Rightarrow \\ \quad \quad \exists a \in (\operatorname{vec}(m) \to \mathbb{R}). \\ \quad \quad \quad \forall x \in \operatorname{vec}(m), z \in s.\; (a(x) \leq (p(\operatorname{vadd}(m)(x, z)) - f(z))) \wedge \\ \quad \quad \quad \forall x \in \operatorname{vec}(m), b \in \mathbb{R}. \\ \quad \quad \quad \quad \forall z \in s.\; (b \leq (p(\operatorname{vadd}(m)(x, z)) - f(z))) \Rightarrow (b \leq a(x)) \end{array}}}
\vnbprop{hbs-red-props}{\vnbstmt{\begin{array}{@{}l@{}} \forall m, p, s, f.; \forall z \in (\operatorname{vec}(m) \to \mathbb{R}). \\ \quad \operatorname{is\text{-}module}(m) \wedge \\ \quad \operatorname{scal}(m) = \operatorname{normed\text{-}field\text{-}as\text{-}commutative\text{-}ring}(\operatorname{rr\text{-}normed\text{-}field}) \wedge \\ \quad \operatorname{is\text{-}sublinear}(m, p) \wedge \\ \quad \operatorname{is\text{-}submodule}(m, s) \wedge \\ \quad \operatorname{is\text{-}linear\text{-}functional\text{-}on}(m, s, f) \wedge \\ \quad \forall a \in s.\; (f(a) \leq p(a)) \wedge \\ \quad \forall x \in \operatorname{vec}(m), a \in s.\; (z(x) \leq (p(\operatorname{vadd}(m)(x, a)) - f(a))) \wedge \\ \quad \forall x \in \operatorname{vec}(m), b \in \mathbb{R}. \\ \quad \quad \forall a \in s.\; (b \leq (p(\operatorname{vadd}(m)(x, a)) - f(a))) \Rightarrow (b \leq z(x)) \Rightarrow \\ \quad \quad (\operatorname{is\text{-}sublinear}(m, z) \wedge (\forall x \in \operatorname{vec}(m).\; (z(x) \leq p(x)) \wedge \forall a \in s.\; (z(a) \leq f(a)))) \end{array}}}
\vnbprop{hbs-grd-in}{\vnbstmt{\begin{array}{@{}l@{}} \forall m, z.; \forall q \in (\operatorname{vec}(m) \to \mathbb{R}). \\ \quad (\operatorname{is\text{-}sublinear}(m, q) \wedge \forall x \in \operatorname{vec}(m).\; (q(x) \leq z(x))) \Rightarrow  \\ \quad \quad q \in \{\mathit{hgq} \in (\operatorname{vec}(m) \to \mathbb{R}) : (\operatorname{is\text{-}sublinear}(m, \mathit{hgq}) \wedge \forall \mathit{hgx} \in \operatorname{vec}(m).\; (\operatorname{hgq}(\mathit{hgx}) \leq z(\mathit{hgx})))\} \end{array}}}
\vnbprop{hbs-grd-parts}{\vnbstmt{\begin{array}{@{}l@{}} \forall m, z.; \forall q \in \{\mathit{hgq} \in (\operatorname{vec}(m) \to \mathbb{R}) : (\operatorname{is\text{-}sublinear}(m, \mathit{hgq}) \wedge \forall x \in \operatorname{vec}(m).\; (\operatorname{hgq}(x) \leq z(x)))\}. \\ \quad ((q \in (\operatorname{vec}(m) \to \mathbb{R})) \wedge (\operatorname{is\text{-}sublinear}(m, q) \wedge \forall a \in \operatorname{vec}(m).\; (q(a) \leq z(a)))) \end{array}}}
\vnbprop{hbs-porel-in}{\vnbstmt{\begin{array}{@{}l@{}} \forall m, z, a, b. \\ \quad a \in \{q \in (\operatorname{vec}(m) \to \mathbb{R}) : (\operatorname{is\text{-}sublinear}(m, q) \wedge \forall x \in \operatorname{vec}(m).\; (q(x) \leq z(x)))\} \wedge \\ \quad b \in \{q \in (\operatorname{vec}(m) \to \mathbb{R}) : (\operatorname{is\text{-}sublinear}(m, q) \wedge \forall x \in \operatorname{vec}(m).\; (q(x) \leq z(x)))\} \wedge \\ \quad \forall c \in \operatorname{vec}(m).\; (b(c) \leq a(c)) \Rightarrow \\ \quad \quad {[}a, b] \\ \quad \quad \in\; \{w \in (\{q \in (\operatorname{vec}(m) \to \mathbb{R}) : (\operatorname{is\text{-}sublinear}(m, q) \wedge \forall x \in \operatorname{vec}(m).\; (q(x) \leq z(x)))\} \times \{q \in (\operatorname{vec}(m) \to \mathbb{R}) : (\operatorname{is\text{-}sublinear}(m, q) \wedge \forall x \in \operatorname{vec}(m).\; (q(x) \leq z(x)))\}) : ((\operatorname{nth}(1, w) \in \{q \in (\operatorname{vec}(m) \to \mathbb{R}) : (\operatorname{is\text{-}sublinear}(m, q) \wedge \forall x \in \operatorname{vec}(m).\; (q(x) \leq z(x)))\}) \wedge ((\operatorname{nth}(2, w) \in \{q \in (\operatorname{vec}(m) \to \mathbb{R}) : (\operatorname{is\text{-}sublinear}(m, q) \wedge \forall x \in \operatorname{vec}(m).\; (q(x) \leq z(x)))\}) \wedge \forall x \in \operatorname{vec}(m).\; (\operatorname{nth}(2, w)(x) \leq \operatorname{nth}(1, w)(x))))\} \end{array}}}
\vnbprop{hbs-porel-parts}{\vnbstmt{\begin{array}{@{}l@{}} \forall m, z, a, b. \\ \quad {[}a, b] \\ \quad \in\; \{w \in (\{q \in (\operatorname{vec}(m) \to \mathbb{R}) : (\operatorname{is\text{-}sublinear}(m, q) \wedge \forall x \in \operatorname{vec}(m).\; (q(x) \leq z(x)))\} \times \{q \in (\operatorname{vec}(m) \to \mathbb{R}) : (\operatorname{is\text{-}sublinear}(m, q) \wedge \forall x \in \operatorname{vec}(m).\; (q(x) \leq z(x)))\}) : ((\operatorname{nth}(1, w) \in \{q \in (\operatorname{vec}(m) \to \mathbb{R}) : (\operatorname{is\text{-}sublinear}(m, q) \wedge \forall x \in \operatorname{vec}(m).\; (q(x) \leq z(x)))\}) \wedge ((\operatorname{nth}(2, w) \in \{q \in (\operatorname{vec}(m) \to \mathbb{R}) : (\operatorname{is\text{-}sublinear}(m, q) \wedge \forall x \in \operatorname{vec}(m).\; (q(x) \leq z(x)))\}) \wedge \forall x \in \operatorname{vec}(m).\; (\operatorname{nth}(2, w)(x) \leq \operatorname{nth}(1, w)(x))))\} \Rightarrow \\ \quad \quad a \in \{q \in (\operatorname{vec}(m) \to \mathbb{R}) : (\operatorname{is\text{-}sublinear}(m, q) \wedge \forall x \in \operatorname{vec}(m).\; (q(x) \leq z(x)))\} \wedge \\ \quad \quad b \in \{q \in (\operatorname{vec}(m) \to \mathbb{R}) : (\operatorname{is\text{-}sublinear}(m, q) \wedge \forall x \in \operatorname{vec}(m).\; (q(x) \leq z(x)))\} \wedge \\ \quad \quad \forall c \in \operatorname{vec}(m).\; (b(c) \leq a(c)) \end{array}}}
\vnbprop{hbs-dominated-linear}{\vnbstmt{\begin{array}{@{}l@{}} \forall m, z. \\ \quad (\operatorname{is\text{-}module}(m) \wedge ((\operatorname{scal}(m) = \operatorname{normed\text{-}field\text{-}as\text{-}commutative\text{-}ring}(\operatorname{rr\text{-}normed\text{-}field})) \wedge \operatorname{is\text{-}sublinear}(m, z))) \Rightarrow  \\ \quad \quad \exists f. \\ \quad \quad \quad \operatorname{is\text{-}linear\text{-}functional\text{-}on}(m, \operatorname{vec}(m), f) \wedge \forall x \in \operatorname{vec}(m).\; (f(x) \leq z(x)) \end{array}}}
\vnbprop{hbs-hahn-banach-sublinear}{\vnbstmt{\begin{array}{@{}l@{}} \forall m, p, s, f. \\ \quad \operatorname{is\text{-}module}(m) \wedge \\ \quad \operatorname{scal}(m) = \operatorname{normed\text{-}field\text{-}as\text{-}commutative\text{-}ring}(\operatorname{rr\text{-}normed\text{-}field}) \wedge \\ \quad \operatorname{is\text{-}sublinear}(m, p) \wedge \\ \quad \operatorname{is\text{-}submodule}(m, s) \wedge \\ \quad \operatorname{is\text{-}linear\text{-}functional\text{-}on}(m, s, f) \wedge \\ \quad \forall z \in s.\; (f(z) \leq p(z)) \Rightarrow \\ \quad \quad \exists a. \\ \quad \quad \quad \operatorname{is\text{-}linear\text{-}functional\text{-}on}(m, \operatorname{vec}(m), a) \wedge \\ \quad \quad \quad \operatorname{extends\text{-}on}(s, a, f) \wedge \\ \quad \quad \quad \forall x \in \operatorname{vec}(m).\; (a(x) \leq p(x)) \end{array}}}
\vnbprop{hbs-rr-fnrm}{\vnbstmt{\forall t \in \mathbb{R}.\; (\operatorname{fnrm}(\operatorname{rr\text{-}normed\text{-}field})(t) = \lvert t \rvert)}}
\vnbprop{hahn-banach-seminorm-rr}{\vnbstmt{\begin{array}{@{}l@{}} \forall m, p, s, f. \\ \quad \operatorname{is\text{-}module}(m) \wedge \\ \quad \operatorname{scal}(m) = \operatorname{normed\text{-}field\text{-}as\text{-}commutative\text{-}ring}(\operatorname{rr\text{-}normed\text{-}field}) \wedge \\ \quad \operatorname{is\text{-}seminorm}(\operatorname{rr\text{-}normed\text{-}field}, m, p) \wedge \\ \quad \operatorname{is\text{-}submodule}(m, s) \wedge \\ \quad \operatorname{is\text{-}k\text{-}linear\text{-}on}(m, s, f) \wedge \\ \quad \forall x \in s.\; (\operatorname{fnrm}(\operatorname{rr\text{-}normed\text{-}field})(f(x)) \leq p(x)) \Rightarrow \\ \quad \quad \exists a. \\ \quad \quad \quad \operatorname{is\text{-}k\text{-}linear}(m, a) \wedge \\ \quad \quad \quad \operatorname{extends\text{-}on}(s, a, f) \wedge \\ \quad \quad \quad \forall x \in \operatorname{vec}(m).\; (\operatorname{fnrm}(\operatorname{rr\text{-}normed\text{-}field})(a(x)) \leq p(x)) \end{array}}}

\vnbbookfile{theorem-library/norm-as-sup-proof}{theorem-library/norm-as-sup-proof}
\noindent{\small norm-as-sup-proof.scm -- the norm as a supremum of bounded functionals: for a finite-dimensional real NVS m and x in VEC(m), ||x|| = sup \{ |f(x)| : f a bounded linear functional, ||f|| $<$= 1 \}.}\par
\vnbprop{norm-bounded-by-functionals}{\vnbstmt{\begin{array}{@{}l@{}} \forall m, f, x. \\ \quad \operatorname{is\text{-}normed\text{-}vector\text{-}space}(m) \wedge \\ \quad \operatorname{is\text{-}bounded\text{-}linear\text{-}functional}(m, f) \wedge \\ \quad x \in \operatorname{vec}(m) \wedge \\ \quad \operatorname{dual\text{-}norm}(m, f) \leq 1 \Rightarrow \\ \quad \quad \lvert f(x) \rvert \leq \operatorname{vnrm}(m)(x) \end{array}}}
\vnbprop{norm-attained-by-functional}{\vnbstmt{\begin{array}{@{}l@{}} \forall m, x. \\ \quad \operatorname{is\text{-}normed\text{-}vector\text{-}space}(m) \wedge \\ \quad \operatorname{is\text{-}finite\text{-}dimensional}(\operatorname{normed\text{-}vector\text{-}space\text{-}as\text{-}module}(m)) \wedge \\ \quad x \in \operatorname{vec}(m) \Rightarrow \\ \quad \quad \exists g. \\ \quad \quad \quad (\operatorname{is\text{-}bounded\text{-}linear\text{-}functional}(m, g) \wedge ((\operatorname{dual\text{-}norm}(m, g) \leq 1) \wedge (g(x) = \operatorname{vnrm}(m)(x)))) \end{array}}}
\vnbprop{norm-as-sup}{\vnbstmt{\begin{array}{@{}l@{}} \forall m, x. \\ \quad \operatorname{is\text{-}normed\text{-}vector\text{-}space}(m) \wedge \\ \quad \operatorname{is\text{-}finite\text{-}dimensional}(\operatorname{normed\text{-}vector\text{-}space\text{-}as\text{-}module}(m)) \wedge \\ \quad x \in \operatorname{vec}(m) \Rightarrow \\ \quad \quad \forall f. \\ \quad \quad \quad (\operatorname{is\text{-}bounded\text{-}linear\text{-}functional}(m, f) \wedge (\operatorname{dual\text{-}norm}(m, f) \leq 1)) \Rightarrow  \\ \quad \quad \quad \quad \lvert f(x) \rvert \leq \operatorname{vnrm}(m)(x) \wedge \\ \quad \quad \forall d \in \mathbb{R}. \\ \quad \quad \quad \forall f.\; ((\operatorname{is\text{-}bounded\text{-}linear\text{-}functional}(m, f) \wedge (\operatorname{dual\text{-}norm}(m, f) \leq 1)) \Rightarrow (\lvert f(x) \rvert \leq d)) \Rightarrow  \\ \quad \quad \quad \quad \operatorname{vnrm}(m)(x) \leq d \end{array}}}

\vnbbookfile{theorem-library/vector-taylor-proof}{theorem-library/vector-taylor-proof}
\noindent{\small vector-taylor-proof.scm -- vector-valued Taylor with a norm remainder bound, REDUCED TO THE SCALAR CASE via a norm-attaining functional (Hahn-Banach payoff).}\par
\par\medskip\noindent\textbf{Definition} (\texttt{taylor-poly-v}).\;For all $m$, $f$, $a$ and $x$, 
\begin{equation*}
(\operatorname{taylor\text{-}poly\text{-}v}(m, f, a, x, 0) \simeq \operatorname{act}(m)(\cdot(\operatorname{power}(-(x, a), 0), \operatorname{recip}(\operatorname{factorial}(0))), \operatorname{nth\text{-}deriv\text{-}v}(m, f, 0)(a)))
\end{equation*}
\par
Let $m$, $f$, $a$ and $x$ be arbitrary. Let $n \in \mathbb{N}$. Then
\begin{equation*}
(\operatorname{taylor\text{-}poly\text{-}v}(m, f, a, x, \operatorname{succ}(n)) \simeq \operatorname{vadd}(m)(\operatorname{taylor\text{-}poly\text{-}v}(m, f, a, x, n), \operatorname{act}(m)(\cdot(\operatorname{power}(-(x, a), \operatorname{succ}(n)), \operatorname{recip}(\operatorname{factorial}(\operatorname{succ}(n)))), \operatorname{nth\text{-}deriv\text{-}v}(m, f, \operatorname{succ}(n))(a))))
\end{equation*}

\par\medskip\noindent\textbf{Definition} (\texttt{is-diff-at-v}).\;For all $m$, $f$, $a$ and $l$, 
\begin{equation*}
\begin{array}{@{}l@{}} \operatorname{is\text{-}diff\text{-}at\text{-}v}(m, f, a, l) \\ \Leftrightarrow\; \operatorname{is\text{-}normed\text{-}vector\text{-}space}(m) \wedge \\ f \in (\mathbb{R} \to \operatorname{vec}(m)) \wedge \\ a \in \mathbb{R} \wedge \\ l \in \operatorname{vec}(m) \wedge \\ \exists \varphi \in (\mathbb{R} \to \operatorname{vec}(m)). \\ \quad \operatorname{is\text{-}continuous\text{-}at}(\operatorname{rr\text{-}ms}, \operatorname{nvs\text{-}metric\text{-}space}(m), \varphi, a) \wedge \\ \quad \varphi(a) = l \wedge \\ \quad \forall \operatorname{x\_} \in \mathbb{R}. \\ \quad \quad \operatorname{vadd}(m)(f(\operatorname{x\_}), \operatorname{vneg}(m)(f(a))) = \operatorname{act}(m)(-(\operatorname{x\_}, a), \varphi(\operatorname{x\_})) \end{array}
\end{equation*}

\par\medskip\noindent\textbf{Definition} (\texttt{deriv-v}).\;$\operatorname{deriv\text{-}v}(m, f, a) :=$ $(\iota\, l.\;\operatorname{is\text{-}diff\text{-}at\text{-}v}(m, f, a, l))$
\par\medskip\noindent\textbf{Definition} (\texttt{nth-deriv-v}).\;For all $m$ and $f$,  $(\operatorname{nth\text{-}deriv\text{-}v}(m, f, 0) \simeq f)$\par
Let $m$ and $f$ be arbitrary. Let $n \in \mathbb{N}$. Then
\begin{equation*}
(\operatorname{nth\text{-}deriv\text{-}v}(m, f, \operatorname{succ}(n)) \simeq (\lambda x \in \mathbb{R}.\; \operatorname{deriv\text{-}v}(m, \operatorname{nth\text{-}deriv\text{-}v}(m, f, n), x)))
\end{equation*}

\par\medskip\noindent\textbf{Definition} (\texttt{taylor-differentiable-v}).\;For all $m$, $f$, $a$, $x$ and $n$, 
\begin{equation*}
\begin{array}{@{}l@{}} \operatorname{taylor\text{-}differentiable\text{-}v}(m, f, a, x, n) \\ \Leftrightarrow\; \forall k. \\ \quad ((k \in \mathbb{N}) \wedge (k \leq n)) \Rightarrow  \\ \quad \quad \forall t \in \operatorname{ccint}(a, x). \\ \quad \quad \quad \operatorname{is\text{-}continuous\text{-}at}(\operatorname{rr\text{-}ms}, \operatorname{nvs\text{-}metric\text{-}space}(m), \operatorname{nth\text{-}deriv\text{-}v}(m, f, k), t) \wedge \\ \forall k. \\ \quad ((k \in \mathbb{N}) \wedge (k \leq n)) \Rightarrow  \\ \quad \quad \forall t. \\ \quad \quad \quad ((a < t) \wedge (t < x)) \Rightarrow  \\ \quad \quad \quad \quad \operatorname{is\text{-}diff\text{-}at\text{-}v}(m, \operatorname{nth\text{-}deriv\text{-}v}(m, f, k), t, \operatorname{nth\text{-}deriv\text{-}v}(m, f, \operatorname{succ}(k))(t)) \end{array}
\end{equation*}

\vnbprop{diff-v-value-in-vec}{\vnbstmt{\forall m, f, a, l.\; (\operatorname{is\text{-}diff\text{-}at\text{-}v}(m, f, a, l) \Rightarrow (l \in \operatorname{vec}(m)))}}
\vnbprop{taylor-v-derivs-in-fun}{\vnbstmt{\begin{array}{@{}l@{}} \forall m, f, a, x, n.; \forall j \in \mathbb{N}. \\ \quad (\operatorname{taylor\text{-}differentiable\text{-}v}(m, f, a, x, n) \wedge ((a \in \mathbb{R}) \wedge ((x \in \mathbb{R}) \wedge ((a < x) \wedge (j \leq n))))) \Rightarrow  \\ \quad \quad \operatorname{nth\text{-}deriv\text{-}v}(m, f, j) \in (\mathbb{R} \to \operatorname{vec}(m)) \end{array}}}
\vnbprop{dfun-v-mono}{\vnbstmt{\begin{array}{@{}l@{}} \forall m, f, n, k.; \forall j \in \mathbb{N}. \\ \quad \forall j \in \mathbb{N}.\; ((j \leq n) \Rightarrow (\operatorname{nth\text{-}deriv\text{-}v}(m, f, j) \in (\mathbb{R} \to \operatorname{vec}(m)))) \wedge \\ \quad n \in \mathbb{N} \wedge \\ \quad k \in \mathbb{N} \wedge \\ \quad k \leq n \wedge \\ \quad j \leq k \Rightarrow \\ \quad \quad \operatorname{nth\text{-}deriv\text{-}v}(m, f, j) \in (\mathbb{R} \to \operatorname{vec}(m)) \end{array}}}
\vnbprop{nth-deriv-v-in-vec}{\vnbstmt{\begin{array}{@{}l@{}} \forall m, f, k, t. \\ \quad \operatorname{is\text{-}normed\text{-}vector\text{-}space}(m) \wedge \\ \quad f \in (\mathbb{R} \to \operatorname{vec}(m)) \wedge \\ \quad k \in \mathbb{N} \wedge \\ \quad t \in \mathbb{R} \wedge \\ \quad \forall j \in \mathbb{N}.\; ((j \leq k) \Rightarrow (\operatorname{nth\text{-}deriv\text{-}v}(m, f, j) \in (\mathbb{R} \to \operatorname{vec}(m)))) \Rightarrow \\ \quad \quad \operatorname{nth\text{-}deriv\text{-}v}(m, f, k)(t) \in \operatorname{vec}(m) \end{array}}}
\vnbprop{taylor-v-deriv-in-vec}{\vnbstmt{\begin{array}{@{}l@{}} \forall m, f, a, x, n, t. \\ \quad (\operatorname{taylor\text{-}differentiable\text{-}v}(m, f, a, x, n) \wedge ((a < t) \wedge ((t < x) \wedge (n \in \mathbb{N})))) \Rightarrow  \\ \quad \quad \operatorname{nth\text{-}deriv\text{-}v}(m, f, (n + 1))(t) \in \operatorname{vec}(m) \end{array}}}
\vnbprop{gof-in-fun}{\vnbstmt{\begin{array}{@{}l@{}} \forall m, f, g. \\ \quad (\operatorname{is\text{-}bounded\text{-}linear\text{-}functional}(m, g) \wedge (f \in (\mathbb{R} \to \operatorname{vec}(m)))) \Rightarrow  \\ \quad \quad \operatorname{compose}(g, f) \in (\mathbb{R} \to \mathbb{R}) \end{array}}}
\vnbprop{rkt-vtaylor-poly-in-vec-ind}{\vnbstmt{\begin{array}{@{}l@{}} \forall n \in \mathbb{N}, m, f, a, x. \\ \quad \operatorname{is\text{-}normed\text{-}vector\text{-}space}(m) \wedge \\ \quad f \in (\mathbb{R} \to \operatorname{vec}(m)) \wedge \\ \quad a \in \mathbb{R} \wedge \\ \quad x \in \mathbb{R} \wedge \\ \quad \forall j \in \mathbb{N}.\; ((j \leq n) \Rightarrow (\operatorname{nth\text{-}deriv\text{-}v}(m, f, j) \in (\mathbb{R} \to \operatorname{vec}(m)))) \Rightarrow \\ \quad \quad \operatorname{taylor\text{-}poly\text{-}v}(m, f, a, x, n) \in \operatorname{vec}(m) \end{array}}}
\vnbprop{vtaylor-poly-in-vec}{\vnbstmt{\begin{array}{@{}l@{}} \forall m, f, a, x, n. \\ \quad \operatorname{is\text{-}normed\text{-}vector\text{-}space}(m) \wedge \\ \quad f \in (\mathbb{R} \to \operatorname{vec}(m)) \wedge \\ \quad a \in \mathbb{R} \wedge \\ \quad x \in \mathbb{R} \wedge \\ \quad n \in \mathbb{N} \wedge \\ \quad \forall j \in \mathbb{N}.\; ((j \leq n) \Rightarrow (\operatorname{nth\text{-}deriv\text{-}v}(m, f, j) \in (\mathbb{R} \to \operatorname{vec}(m)))) \Rightarrow \\ \quad \quad \operatorname{taylor\text{-}poly\text{-}v}(m, f, a, x, n) \in \operatorname{vec}(m) \end{array}}}
\vnbprop{vtaylor-remainder-in-vec}{\vnbstmt{\begin{array}{@{}l@{}} \forall m, f, a, x, n. \\ \quad \operatorname{is\text{-}normed\text{-}vector\text{-}space}(m) \wedge \\ \quad f \in (\mathbb{R} \to \operatorname{vec}(m)) \wedge \\ \quad a \in \mathbb{R} \wedge \\ \quad x \in \mathbb{R} \wedge \\ \quad n \in \mathbb{N} \wedge \\ \quad \forall j \in \mathbb{N}.\; ((j \leq n) \Rightarrow (\operatorname{nth\text{-}deriv\text{-}v}(m, f, j) \in (\mathbb{R} \to \operatorname{vec}(m)))) \Rightarrow \\ \quad \quad \operatorname{vadd}(m)(f(x), \operatorname{vneg}(m)(\operatorname{taylor\text{-}poly\text{-}v}(m, f, a, x, n))) \in \operatorname{vec}(m) \end{array}}}
\vnbprop{r8e-blf-vneg}{\vnbstmt{\begin{array}{@{}l@{}} \forall m, g.; \forall u \in \operatorname{vec}(m). \\ \quad (\operatorname{is\text{-}normed\text{-}vector\text{-}space}(m) \wedge \operatorname{is\text{-}bounded\text{-}linear\text{-}functional}(m, g)) \Rightarrow  \\ \quad \quad g(\operatorname{vneg}(m)(u)) = (0 - g(u)) \end{array}}}
\vnbprop{r8e-blf-sub}{\vnbstmt{\begin{array}{@{}l@{}} \forall m, g.; \forall u, v \in \operatorname{vec}(m). \\ \quad (\operatorname{is\text{-}normed\text{-}vector\text{-}space}(m) \wedge \operatorname{is\text{-}bounded\text{-}linear\text{-}functional}(m, g)) \Rightarrow  \\ \quad \quad g(\operatorname{vadd}(m)(u, \operatorname{vneg}(m)(v))) = (g(u) - g(v)) \end{array}}}
\vnbprop{r8e-eps-over}{\vnbstmt{\begin{array}{@{}l@{}} \forall c \in \mathbb{R}, e. \\ \quad ((0 \leq c) \wedge \operatorname{pos\text{-}rr}(e)) \Rightarrow  \\ \quad \quad \exists d.\; (\operatorname{pos\text{-}rr}(d) \wedge \forall z \in \mathbb{R}.\; ((z \leq d) \Rightarrow ((c \cdot z) \leq e))) \end{array}}}
\vnbprop{r8e-blf-compose-continuous}{\vnbstmt{\begin{array}{@{}l@{}} \forall m, g, h, t. \\ \quad \operatorname{is\text{-}normed\text{-}vector\text{-}space}(m) \wedge \\ \quad \operatorname{is\text{-}bounded\text{-}linear\text{-}functional}(m, g) \wedge \\ \quad h \in (\mathbb{R} \to \operatorname{vec}(m)) \wedge \\ \quad \operatorname{is\text{-}continuous\text{-}at}(\operatorname{rr\text{-}ms}, \operatorname{nvs\text{-}metric\text{-}space}(m), h, t) \Rightarrow \\ \quad \quad \operatorname{is\text{-}continuous\text{-}at}(\operatorname{rr\text{-}ms}, \operatorname{rr\text{-}ms}, \operatorname{compose}(g, h), t) \end{array}}}
\vnbprop{r8e-blf-diff-commute}{\vnbstmt{\begin{array}{@{}l@{}} \forall m, g, h, t, l. \\ \quad (\operatorname{is\text{-}bounded\text{-}linear\text{-}functional}(m, g) \wedge \operatorname{is\text{-}diff\text{-}at\text{-}v}(m, h, t, l)) \Rightarrow  \\ \quad \quad \operatorname{is\text{-}diff\text{-}at}(\operatorname{compose}(g, h), t, g(l)) \end{array}}}
\vnbprop{r8e-nth-deriv-v-diff}{\vnbstmt{\begin{array}{@{}l@{}} \forall m, f, j, x. \\ \quad j \in \mathbb{N} \wedge \\ \quad x \in \mathbb{R} \wedge \\ \quad \operatorname{nth\text{-}deriv\text{-}v}(m, f, j) \in (\mathbb{R} \to \operatorname{vec}(m)) \wedge \\ \quad \operatorname{nth\text{-}deriv\text{-}v}(m, f, (j + 1)) \in (\mathbb{R} \to \operatorname{vec}(m)) \Rightarrow \\ \quad \quad \operatorname{is\text{-}diff\text{-}at\text{-}v}(m, \operatorname{nth\text{-}deriv\text{-}v}(m, f, j), x, \operatorname{nth\text{-}deriv\text{-}v}(m, f, (j + 1))(x)) \end{array}}}
\vnbprop{r8e-gof-nth-deriv-fn}{\vnbstmt{\begin{array}{@{}l@{}} \forall k \in \mathbb{N}, m, f, g. \\ \quad \operatorname{is\text{-}bounded\text{-}linear\text{-}functional}(m, g) \wedge \\ \quad f \in (\mathbb{R} \to \operatorname{vec}(m)) \wedge \\ \quad \forall j \in \mathbb{N}.\; ((j \leq k) \Rightarrow (\operatorname{nth\text{-}deriv\text{-}v}(m, f, j) \in (\mathbb{R} \to \operatorname{vec}(m)))) \Rightarrow \\ \quad \quad \operatorname{nth\text{-}deriv}(\operatorname{compose}(g, f), k) = \operatorname{compose}(g, \operatorname{nth\text{-}deriv\text{-}v}(m, f, k)) \end{array}}}
\vnbprop{gof-nth-deriv-dfun}{\vnbstmt{\begin{array}{@{}l@{}} \forall m, f, g, n, t. \\ \quad \operatorname{is\text{-}bounded\text{-}linear\text{-}functional}(m, g) \wedge \\ \quad f \in (\mathbb{R} \to \operatorname{vec}(m)) \wedge \\ \quad n \in \mathbb{N} \wedge \\ \quad t \in \mathbb{R} \wedge \\ \quad \forall j \in \mathbb{N}.\; ((j \leq n) \Rightarrow (\operatorname{nth\text{-}deriv\text{-}v}(m, f, j) \in (\mathbb{R} \to \operatorname{vec}(m)))) \wedge \\ \quad \operatorname{is\text{-}diff\text{-}at\text{-}v}(m, \operatorname{nth\text{-}deriv\text{-}v}(m, f, n), t, \operatorname{nth\text{-}deriv\text{-}v}(m, f, (n + 1))(t)) \Rightarrow \\ \quad \quad \operatorname{nth\text{-}deriv}(\operatorname{compose}(g, f), (n + 1))(t) = g(\operatorname{nth\text{-}deriv\text{-}v}(m, f, (n + 1))(t)) \end{array}}}
\vnbprop{gof-taylor-diff-dfun}{\vnbstmt{\begin{array}{@{}l@{}} \forall m, f, g, a, x, n. \\ \quad \operatorname{is\text{-}normed\text{-}vector\text{-}space}(m) \wedge \\ \quad \operatorname{is\text{-}bounded\text{-}linear\text{-}functional}(m, g) \wedge \\ \quad f \in (\mathbb{R} \to \operatorname{vec}(m)) \wedge \\ \quad n \in \mathbb{N} \wedge \\ \quad \operatorname{taylor\text{-}differentiable\text{-}v}(m, f, a, x, n) \wedge \\ \quad \forall j \in \mathbb{N}.\; ((j \leq n) \Rightarrow (\operatorname{nth\text{-}deriv\text{-}v}(m, f, j) \in (\mathbb{R} \to \operatorname{vec}(m)))) \Rightarrow \\ \quad \quad \operatorname{taylor\text{-}differentiable}(\operatorname{compose}(g, f), a, x, n) \end{array}}}
\vnbprop{r8e-td-v-diff-at}{\vnbstmt{\begin{array}{@{}l@{}} \forall m, f, a, x, n, t. \\ \quad (\operatorname{taylor\text{-}differentiable\text{-}v}(m, f, a, x, n) \wedge ((n \in \mathbb{N}) \wedge ((a < t) \wedge (t < x)))) \Rightarrow  \\ \quad \quad \operatorname{is\text{-}diff\text{-}at\text{-}v}(m, \operatorname{nth\text{-}deriv\text{-}v}(m, f, n), t, \operatorname{nth\text{-}deriv\text{-}v}(m, f, (n + 1))(t)) \end{array}}}
\vnbprop{r8e-tp-unfold}{\vnbstmt{\begin{array}{@{}l@{}} \forall f, a, d, x. \\ \quad (\operatorname{taylor\text{-}poly}(f, a, d, x) \simeq \operatorname{series\text{-}partial\text{-}sum}((\lambda k \in \mathbb{N}.\; ((\operatorname{nth\text{-}deriv}(f, k)(a) \cdot \operatorname{power}((x - a), k)) \cdot \operatorname{recip}(\operatorname{factorial}(k)))), (d + 1))) \end{array}}}
\vnbprop{r8e-coef-comm}{\vnbstmt{\forall p, q, v \in \mathbb{R}.\; (((p \cdot q) \cdot v) = ((v \cdot p) \cdot q))}}
\vnbprop{r8e-g-of-poly-ind}{\vnbstmt{\begin{array}{@{}l@{}} \forall n \in \mathbb{N}, m, f, g, a, x. \\ \quad \operatorname{is\text{-}normed\text{-}vector\text{-}space}(m) \wedge \\ \quad \operatorname{is\text{-}bounded\text{-}linear\text{-}functional}(m, g) \wedge \\ \quad f \in (\mathbb{R} \to \operatorname{vec}(m)) \wedge \\ \quad a \in \mathbb{R} \wedge \\ \quad x \in \mathbb{R} \wedge \\ \quad \forall j \in \mathbb{N}.\; ((j \leq n) \Rightarrow (\operatorname{nth\text{-}deriv\text{-}v}(m, f, j) \in (\mathbb{R} \to \operatorname{vec}(m)))) \Rightarrow \\ \quad \quad g(\operatorname{taylor\text{-}poly\text{-}v}(m, f, a, x, n)) = \operatorname{taylor\text{-}poly}(\operatorname{compose}(g, f), a, n, x) \end{array}}}
\vnbprop{g-of-remainder-dfun}{\vnbstmt{\begin{array}{@{}l@{}} \forall m, f, g, a, x, n. \\ \quad \operatorname{is\text{-}normed\text{-}vector\text{-}space}(m) \wedge \\ \quad \operatorname{is\text{-}bounded\text{-}linear\text{-}functional}(m, g) \wedge \\ \quad f \in (\mathbb{R} \to \operatorname{vec}(m)) \wedge \\ \quad a \in \mathbb{R} \wedge \\ \quad x \in \mathbb{R} \wedge \\ \quad n \in \mathbb{N} \wedge \\ \quad \forall j \in \mathbb{N}.\; ((j \leq n) \Rightarrow (\operatorname{nth\text{-}deriv\text{-}v}(m, f, j) \in (\mathbb{R} \to \operatorname{vec}(m)))) \Rightarrow \\ \quad \quad g(\operatorname{vadd}(m)(f(x), \operatorname{vneg}(m)(\operatorname{taylor\text{-}poly\text{-}v}(m, f, a, x, n)))) \\ \quad \quad =\; (\operatorname{compose}(g, f)(x) - \operatorname{taylor\text{-}poly}(\operatorname{compose}(g, f), a, n, x)) \end{array}}}
\vnbprop{vector-taylor-remainder-bound}{\vnbstmt{\begin{array}{@{}l@{}} \forall m, f, a, x, n. \\ \quad \operatorname{is\text{-}normed\text{-}vector\text{-}space}(m) \wedge \\ \quad \operatorname{is\text{-}finite\text{-}dimensional}(\operatorname{normed\text{-}vector\text{-}space\text{-}as\text{-}module}(m)) \wedge \\ \quad f \in (\mathbb{R} \to \operatorname{vec}(m)) \wedge \\ \quad a \in \mathbb{R} \wedge \\ \quad x \in \mathbb{R} \wedge \\ \quad n \in \mathbb{N} \wedge \\ \quad a < x \wedge \\ \quad \operatorname{taylor\text{-}differentiable\text{-}v}(m, f, a, x, n) \Rightarrow \\ \quad \quad \exists b \in \mathbb{R}. \\ \quad \quad \quad a < b \wedge \\ \quad \quad \quad b < x \wedge \\ \quad \quad \quad \operatorname{factorial}((n + 1)) \cdot \operatorname{vnrm}(m)(\operatorname{vadd}(m)(f(x), \operatorname{vneg}(m)(\operatorname{taylor\text{-}poly\text{-}v}(m, f, a, x, n)))) \\ \quad \quad \quad \leq\; \operatorname{vnrm}(m)(\operatorname{nth\text{-}deriv\text{-}v}(m, f, (n + 1))(b)) \cdot \operatorname{power}((x - a), (n + 1)) \end{array}}}

\vnbbookfile{theorem-library/rake-nf-norm}{theorem-library/rake-nf-norm}
\noindent{\small rake-nf-norm.scm -- rake batch 5b-D: the norm of a negative, and the metric space underlying a normed field.}\par
\vnbprop{nf-norm-neg}{\vnbstmt{\begin{array}{@{}l@{}} \forall f.; \forall a \in \operatorname{carr}(f). \\ \quad \operatorname{is\text{-}normed\text{-}field}(f) \Rightarrow (\operatorname{fnrm}(f)(\operatorname{neg}(f)(a)) = \operatorname{fnrm}(f)(a)) \end{array}}}
\vnbprop{nf-metric-space-is-metric-space}{\vnbstmt{\forall f.\; (\operatorname{is\text{-}normed\text{-}field}(f) \Rightarrow \operatorname{is\text{-}metric\text{-}space}(\operatorname{nf\text{-}metric\text{-}space}(f)))}}

\vnbbookfile{theorem-library/diff-on-laws}{theorem-library/diff-on-laws}
\noindent{\small diff-on-laws.scm -- the laws of IS-DIFF-ON / DERIV-ON / HOLOMORPHIC-ON (structure-library/diff-on.scm).  Everything here is PROVEN; the definition file asserts nothing.}\par
\vnbprop{diff-on-normed-field}{\vnbstmt{\forall k, u, f, a, l.\; (\operatorname{is\text{-}diff\text{-}on}(k, u, f, a, l) \Rightarrow \operatorname{is\text{-}normed\text{-}field}(k))}}
\vnbprop{diff-on-open}{\vnbstmt{\begin{array}{@{}l@{}} \forall k, u, f, a, l. \\ \quad \operatorname{is\text{-}diff\text{-}on}(k, u, f, a, l) \Rightarrow \operatorname{is\text{-}open}(\operatorname{nf\text{-}metric\text{-}space}(k), u) \end{array}}}
\vnbprop{diff-on-in-fun}{\vnbstmt{\forall k, u, f, a, l.\; (\operatorname{is\text{-}diff\text{-}on}(k, u, f, a, l) \Rightarrow (f \in (u \to \operatorname{carr}(k))))}}
\vnbprop{diff-on-pt-in}{\vnbstmt{\forall k, u, f, a, l.\; (\operatorname{is\text{-}diff\text{-}on}(k, u, f, a, l) \Rightarrow (a \in u))}}
\vnbprop{diff-on-deriv-in-carr}{\vnbstmt{\forall k, u, f, a, l.\; (\operatorname{is\text{-}diff\text{-}on}(k, u, f, a, l) \Rightarrow (l \in \operatorname{carr}(k)))}}
\vnbprop{nf-open-subset-carr}{\vnbstmt{\forall k, u.\; (\operatorname{is\text{-}open}(\operatorname{nf\text{-}metric\text{-}space}(k), u) \Rightarrow (u \subseteq \operatorname{carr}(k)))}}
\vnbprop{diff-on-value-in-carr}{\vnbstmt{\begin{array}{@{}l@{}} \forall k, u, f, a, l.; \forall x \in u. \\ \quad \operatorname{is\text{-}diff\text{-}on}(k, u, f, a, l) \Rightarrow (f(x) \in \operatorname{carr}(k)) \end{array}}}
\vnbprop{rr-nf-carr}{\vnbstmt{(\operatorname{carr}(\operatorname{rr\text{-}normed\text{-}field}) \simeq \mathbb{R})}}
\vnbprop{cc-nf-carr}{\vnbstmt{(\operatorname{carr}(\operatorname{cc\text{-}normed\text{-}field}) \simeq \mathbb{C})}}
\vnbprop{rr-nf-add}{\vnbstmt{(\operatorname{add}(\operatorname{rr\text{-}normed\text{-}field}) \simeq (\lambda {[}\operatorname{x\_}, \operatorname{y\_}] \in (\mathbb{R} \times \mathbb{R}).\; (\operatorname{x\_} + \operatorname{y\_})))}}
\vnbprop{rr-nf-mul}{\vnbstmt{(\operatorname{mul}(\operatorname{rr\text{-}normed\text{-}field}) \simeq (\lambda {[}\operatorname{x\_}, \operatorname{y\_}] \in (\mathbb{R} \times \mathbb{R}).\; (\operatorname{x\_} \cdot \operatorname{y\_})))}}
\vnbprop{rr-nf-neg}{\vnbstmt{(\operatorname{neg}(\operatorname{rr\text{-}normed\text{-}field}) \simeq (\lambda x \in \mathbb{R}.\; -x))}}
\vnbprop{rr-nf-fnrm}{\vnbstmt{(\operatorname{fnrm}(\operatorname{rr\text{-}normed\text{-}field}) \simeq (\lambda x \in \mathbb{R}.\; \lvert x \rvert))}}
\vnbprop{cc-nf-fnrm}{\vnbstmt{(\operatorname{fnrm}(\operatorname{cc\text{-}normed\text{-}field}) \simeq (\lambda x \in \mathbb{C}.\; \operatorname{magnitude}(x)))}}
\vnbprop{cc-nf-zero}{\vnbstmt{(\operatorname{zero}(\operatorname{cc\text{-}normed\text{-}field}) \simeq 0)}}
\vnbprop{rr-nf-zero}{\vnbstmt{(\operatorname{zero}(\operatorname{rr\text{-}normed\text{-}field}) \simeq 0)}}
\vnbprop{rr-nf-add-apply}{\vnbstmt{\begin{array}{@{}l@{}} \forall u, v. \\ \quad ((u \in \mathbb{R}) \wedge (v \in \mathbb{R})) \Rightarrow (\operatorname{add}(\operatorname{rr\text{-}normed\text{-}field})(u, v) \simeq (u + v)) \end{array}}}
\vnbprop{rr-nf-mul-apply}{\vnbstmt{\begin{array}{@{}l@{}} \forall u, v. \\ \quad ((u \in \mathbb{R}) \wedge (v \in \mathbb{R})) \Rightarrow (\operatorname{mul}(\operatorname{rr\text{-}normed\text{-}field})(u, v) \simeq (u \cdot v)) \end{array}}}
\vnbprop{rr-nf-neg-apply}{\vnbstmt{\forall u \in \mathbb{R}.\; (\operatorname{neg}(\operatorname{rr\text{-}normed\text{-}field})(u) \simeq -u)}}
\vnbprop{rr-nf-fnrm-apply}{\vnbstmt{\forall u \in \mathbb{R}.\; (\operatorname{fnrm}(\operatorname{rr\text{-}normed\text{-}field})(u) \simeq \lvert u \rvert)}}
\vnbprop{cc-nf-fnrm-apply}{\vnbstmt{\forall u \in \mathbb{C}.\; (\operatorname{fnrm}(\operatorname{cc\text{-}normed\text{-}field})(u) \simeq \operatorname{magnitude}(u))}}
\vnbprop{ms-agree-ball}{\vnbstmt{\begin{array}{@{}l@{}} \forall a, b.; \forall c \in \operatorname{pts}(a), r. \\ \quad \operatorname{is\text{-}metric\text{-}space}(a) \wedge \\ \quad \operatorname{is\text{-}metric\text{-}space}(b) \wedge \\ \quad (\operatorname{pts}(b) \simeq \operatorname{pts}(a)) \wedge \\ \quad \forall u, v \in \operatorname{pts}(a).\; (\operatorname{dist}(b)(u, v) \simeq \operatorname{dist}(a)(u, v)) \Rightarrow \\ \quad \quad \operatorname{ball}(b, c, r) = \operatorname{ball}(a, c, r) \end{array}}}
\vnbprop{ms-agree-open}{\vnbstmt{\begin{array}{@{}l@{}} \forall a, b, w. \\ \quad \operatorname{is\text{-}metric\text{-}space}(b) \wedge \\ \quad (\operatorname{pts}(b) \simeq \operatorname{pts}(a)) \wedge \\ \quad \forall u, v \in \operatorname{pts}(a).\; (\operatorname{dist}(b)(u, v) \simeq \operatorname{dist}(a)(u, v)) \wedge \\ \quad \operatorname{is\text{-}open}(a, w) \Rightarrow \\ \quad \quad \operatorname{is\text{-}open}(b, w) \end{array}}}
\vnbprop{ms-agree-continuous-at}{\vnbstmt{\begin{array}{@{}l@{}} \forall a, b, \mathit{mta}, \mathit{mtb}, g, p. \\ \quad \operatorname{is\text{-}metric\text{-}space}(b) \wedge \\ \quad \operatorname{is\text{-}metric\text{-}space}(\mathit{mtb}) \wedge \\ \quad (\operatorname{pts}(b) \simeq \operatorname{pts}(a)) \wedge \\ \quad (\operatorname{pts}(\mathit{mtb}) \simeq \operatorname{pts}(\mathit{mta})) \wedge \\ \quad \forall u, v \in \operatorname{pts}(a).\; (\operatorname{dist}(b)(u, v) \simeq \operatorname{dist}(a)(u, v)) \wedge \\ \quad \forall u, v \in \operatorname{pts}(\mathit{mta}).\; (\operatorname{dist}(\mathit{mtb})(u, v) \simeq \operatorname{dist}(\mathit{mta})(u, v)) \wedge \\ \quad \operatorname{is\text{-}continuous\text{-}at}(a, \mathit{mta}, g, p) \Rightarrow \\ \quad \quad \operatorname{is\text{-}continuous\text{-}at}(b, \mathit{mtb}, g, p) \end{array}}}
\vnbprop{ms-agree-pts-sym}{\vnbstmt{\forall a, b.\; ((\operatorname{pts}(b) \simeq \operatorname{pts}(a)) \Rightarrow (\operatorname{pts}(a) \simeq \operatorname{pts}(b)))}}
\vnbprop{ms-agree-dist-sym}{\vnbstmt{\begin{array}{@{}l@{}} \forall a, b.; \forall u, v \in \operatorname{pts}(b). \\ \quad ((\operatorname{pts}(b) \simeq \operatorname{pts}(a)) \wedge \forall u, v \in \operatorname{pts}(a).\; (\operatorname{dist}(b)(u, v) \simeq \operatorname{dist}(a)(u, v))) \Rightarrow  \\ \quad \quad (\operatorname{dist}(a)(u, v) \simeq \operatorname{dist}(b)(u, v)) \end{array}}}
\vnbprop{nf-rr-is-metric-space}{\vnbstmt{\operatorname{is\text{-}metric\text{-}space}(\operatorname{nf\text{-}metric\text{-}space}(\operatorname{rr\text{-}normed\text{-}field}))}}
\vnbprop{nf-rr-agree-pts}{\vnbstmt{(\operatorname{pts}(\operatorname{nf\text{-}metric\text{-}space}(\operatorname{rr\text{-}normed\text{-}field})) \simeq \operatorname{pts}(\operatorname{rr\text{-}ms}))}}
\vnbprop{nf-rr-agree-dist}{\vnbstmt{\begin{array}{@{}l@{}} \forall u, v \in \operatorname{pts}(\operatorname{rr\text{-}ms}). \\ \quad (\operatorname{dist}(\operatorname{nf\text{-}metric\text{-}space}(\operatorname{rr\text{-}normed\text{-}field}))(u, v) \simeq \operatorname{dist}(\operatorname{rr\text{-}ms})(u, v)) \end{array}}}
\vnbprop{rr-open-in-nf}{\vnbstmt{\operatorname{is\text{-}open}(\operatorname{nf\text{-}metric\text{-}space}(\operatorname{rr\text{-}normed\text{-}field}), \mathbb{R})}}
\vnbprop{rr-subset-nf-pts}{\vnbstmt{\mathbb{R} \subseteq \operatorname{pts}(\operatorname{nf\text{-}metric\text{-}space}(\operatorname{rr\text{-}normed\text{-}field}))}}
\vnbprop{sub-rr-is-metric-space}{\vnbstmt{\operatorname{is\text{-}metric\text{-}space}(\operatorname{subspace\text{-}ms}(\operatorname{nf\text{-}metric\text{-}space}(\operatorname{rr\text{-}normed\text{-}field}), \mathbb{R}))}}
\vnbprop{sub-rr-agree-pts}{\vnbstmt{(\operatorname{pts}(\operatorname{subspace\text{-}ms}(\operatorname{nf\text{-}metric\text{-}space}(\operatorname{rr\text{-}normed\text{-}field}), \mathbb{R})) \simeq \operatorname{pts}(\operatorname{rr\text{-}ms}))}}
\vnbprop{sub-rr-agree-dist}{\vnbstmt{\begin{array}{@{}l@{}} \forall u, v \in \operatorname{pts}(\operatorname{rr\text{-}ms}). \\ \quad (\operatorname{dist}(\operatorname{subspace\text{-}ms}(\operatorname{nf\text{-}metric\text{-}space}(\operatorname{rr\text{-}normed\text{-}field}), \mathbb{R}))(u, v) \simeq \operatorname{dist}(\operatorname{rr\text{-}ms})(u, v)) \end{array}}}
\vnbprop{rr-nf-sub}{\vnbstmt{\begin{array}{@{}l@{}} \forall u, v. \\ \quad ((u \in \mathbb{R}) \wedge (v \in \mathbb{R})) \Rightarrow  \\ \quad \quad (\operatorname{add}(\operatorname{rr\text{-}normed\text{-}field})(u, \operatorname{neg}(\operatorname{rr\text{-}normed\text{-}field})(v)) \simeq (u - v)) \end{array}}}
\vnbprop{rr-nf-fun-class}{\vnbstmt{((\mathbb{R} \to \operatorname{carr}(\operatorname{rr\text{-}normed\text{-}field})) \simeq (\mathbb{R} \to \mathbb{R}))}}
\vnbprop{diff-at-iff-diff-on}{\vnbstmt{\begin{array}{@{}l@{}} \forall f, a, l. \\ \quad \operatorname{is\text{-}diff\text{-}at}(f, a, l) \Leftrightarrow \operatorname{is\text{-}diff\text{-}on}(\operatorname{rr\text{-}normed\text{-}field}, \mathbb{R}, f, a, l) \end{array}}}
\vnbprop{nf-norm-in-rr}{\vnbstmt{\forall k.; \forall u \in \operatorname{carr}(k).\; (\operatorname{is\text{-}normed\text{-}field}(k) \Rightarrow (\operatorname{fnrm}(k)(u) \in \mathbb{R}))}}
\vnbprop{nf-norm-nonneg}{\vnbstmt{\forall k.; \forall u \in \operatorname{carr}(k).\; (\operatorname{is\text{-}normed\text{-}field}(k) \Rightarrow (0 \leq \operatorname{fnrm}(k)(u)))}}
\vnbprop{nf-norm-mult}{\vnbstmt{\begin{array}{@{}l@{}} \forall k, u, v. \\ \quad (\operatorname{is\text{-}normed\text{-}field}(k) \wedge ((u \in \operatorname{carr}(k)) \wedge (v \in \operatorname{carr}(k)))) \Rightarrow  \\ \quad \quad \operatorname{fnrm}(k)(\operatorname{mul}(k)(u, v)) = (\operatorname{fnrm}(k)(u) \cdot \operatorname{fnrm}(k)(v)) \end{array}}}
\vnbprop{nf-norm-subadd}{\vnbstmt{\begin{array}{@{}l@{}} \forall k, u, v. \\ \quad (\operatorname{is\text{-}normed\text{-}field}(k) \wedge ((u \in \operatorname{carr}(k)) \wedge (v \in \operatorname{carr}(k)))) \Rightarrow  \\ \quad \quad \operatorname{fnrm}(k)(\operatorname{add}(k)(u, v)) \leq (\operatorname{fnrm}(k)(u) + \operatorname{fnrm}(k)(v)) \end{array}}}
\vnbprop{nf-add-in-carr}{\vnbstmt{\begin{array}{@{}l@{}} \forall k, u, v. \\ \quad (\operatorname{is\text{-}normed\text{-}field}(k) \wedge ((u \in \operatorname{carr}(k)) \wedge (v \in \operatorname{carr}(k)))) \Rightarrow  \\ \quad \quad \operatorname{add}(k)(u, v) \in \operatorname{carr}(k) \end{array}}}
\vnbprop{nf-mul-in-carr}{\vnbstmt{\begin{array}{@{}l@{}} \forall k, u, v. \\ \quad (\operatorname{is\text{-}normed\text{-}field}(k) \wedge ((u \in \operatorname{carr}(k)) \wedge (v \in \operatorname{carr}(k)))) \Rightarrow  \\ \quad \quad \operatorname{mul}(k)(u, v) \in \operatorname{carr}(k) \end{array}}}
\vnbprop{nf-neg-in-carr}{\vnbstmt{\forall k.; \forall u \in \operatorname{carr}(k).\; (\operatorname{is\text{-}normed\text{-}field}(k) \Rightarrow (\operatorname{neg}(k)(u) \in \operatorname{carr}(k)))}}
\vnbprop{nf-sub-add-back}{\vnbstmt{\begin{array}{@{}l@{}} \forall k, u, v. \\ \quad (\operatorname{is\text{-}normed\text{-}field}(k) \wedge ((u \in \operatorname{carr}(k)) \wedge (v \in \operatorname{carr}(k)))) \Rightarrow  \\ \quad \quad \operatorname{add}(k)(\operatorname{add}(k)(u, \operatorname{neg}(k)(v)), v) = u \end{array}}}
\vnbprop{nf-norm-le-add}{\vnbstmt{\begin{array}{@{}l@{}} \forall k, u, v. \\ \quad (\operatorname{is\text{-}normed\text{-}field}(k) \wedge ((u \in \operatorname{carr}(k)) \wedge (v \in \operatorname{carr}(k)))) \Rightarrow  \\ \quad \quad \operatorname{fnrm}(k)(u) \leq (\operatorname{fnrm}(k)(\operatorname{add}(k)(u, \operatorname{neg}(k)(v))) + \operatorname{fnrm}(k)(v)) \end{array}}}
\vnbprop{diff-on-implies-continuous}{\vnbstmt{\begin{array}{@{}l@{}} \forall k, u, f, a, l. \\ \quad \operatorname{is\text{-}diff\text{-}on}(k, u, f, a, l) \Rightarrow  \\ \quad \quad \operatorname{is\text{-}continuous\text{-}at}(\operatorname{subspace\text{-}ms}(\operatorname{nf\text{-}metric\text{-}space}(k), u), \operatorname{nf\text{-}metric\text{-}space}(k), f, a) \end{array}}}
\vnbprop{deriv-on-unfold}{\vnbstmt{\forall k, u, f, a.\; (\operatorname{deriv\text{-}on}(k, u, f, a) \simeq (\iota\, l.\;\operatorname{is\text{-}diff\text{-}on}(k, u, f, a, l)))}}
\vnbprop{holomorphic-on-open}{\vnbstmt{\begin{array}{@{}l@{}} \forall u, f. \\ \quad \operatorname{holomorphic\text{-}on}(u, f) \Rightarrow \operatorname{is\text{-}open}(\operatorname{nf\text{-}metric\text{-}space}(\operatorname{cc\text{-}normed\text{-}field}), u) \end{array}}}
\vnbprop{holomorphic-on-in-fun}{\vnbstmt{\forall u, f.\; (\operatorname{holomorphic\text{-}on}(u, f) \Rightarrow (f \in (u \to \mathbb{C})))}}
\vnbprop{holomorphic-on-diff}{\vnbstmt{\begin{array}{@{}l@{}} \forall u, f.; \forall a \in u. \\ \quad \operatorname{holomorphic\text{-}on}(u, f) \Rightarrow  \\ \quad \quad \exists l.\; \operatorname{is\text{-}diff\text{-}on}(\operatorname{cc\text{-}normed\text{-}field}, u, f, a, l) \end{array}}}

\vnbbookfile{theorem-library/tvs-hyperplane}{theorem-library/tvs-hyperplane}
\noindent{\small tb-hyper.scm -- topological-vector-space.pdf Prop 2.10: two linear functionals with a common hyperplane (a common coset of their kernels) are proportional, by a non-zero factor. Agent TB-01, 2026-10-26.  DESTINATION: theorem-library/tvs-hyperplane.scm (after linear-functional, hahn-banach-sublinear (hbs-lf-*)).  ...}\par
\vnbprop{hyperplane-functionals-proportional}{\vnbstmt{\begin{array}{@{}l@{}} \forall m, f, g, a, b. \\ \quad \operatorname{is\text{-}module}(m) \wedge \\ \quad \operatorname{scal}(m) = \operatorname{normed\text{-}field\text{-}as\text{-}commutative\text{-}ring}(\operatorname{rr\text{-}normed\text{-}field}) \wedge \\ \quad \operatorname{is\text{-}linear\text{-}functional\text{-}on}(m, \operatorname{vec}(m), f) \wedge \\ \quad \operatorname{is\text{-}linear\text{-}functional\text{-}on}(m, \operatorname{vec}(m), g) \wedge \\ \quad \exists \mathit{hpa} \in \operatorname{vec}(m).\; \neg (f(\mathit{hpa}) = 0) \wedge \\ \quad \exists \mathit{hpa} \in \operatorname{vec}(m).\; \neg (g(\mathit{hpa}) = 0) \wedge \\ \quad a \in \operatorname{vec}(m) \wedge \\ \quad b \in \operatorname{vec}(m) \wedge \\ \quad \forall z \in \operatorname{vec}(m).\; ((f(z) = f(a)) \Leftrightarrow (g(z) = g(b))) \Rightarrow \\ \quad \quad \exists c \in \mathbb{R}.\; (\neg (c = 0) \wedge \forall z \in \operatorname{vec}(m).\; (f(z) = (c \cdot g(z)))) \end{array}}}

\vnbbookfile{theorem-library/ms-continuity-algebra}{theorem-library/ms-continuity-algebra}
\noindent{\small ms-continuity-algebra.scm -- THE CONTINUITY ALGEBRA OFF THE REAL LINE.}\par
\vnbprop{ms-pts-is-set}{\vnbstmt{\forall s.\; (\operatorname{is\text{-}metric\text{-}space}(s) \Rightarrow (\operatorname{pts}(s) \in \mathbf{Set}))}}
\vnbprop{ms-const-lam-in-fun}{\vnbstmt{\begin{array}{@{}l@{}} \forall s, t.; \forall c \in \operatorname{pts}(t). \\ \quad \operatorname{is\text{-}metric\text{-}space}(s) \Rightarrow ((\lambda z \in \operatorname{pts}(s).\; c) \in (\operatorname{pts}(s) \to \operatorname{pts}(t))) \end{array}}}
\vnbprop{ms-const-continuous-at}{\vnbstmt{\begin{array}{@{}l@{}} \forall s, t.; \forall c \in \operatorname{pts}(t), a \in \operatorname{pts}(s). \\ \quad (\operatorname{is\text{-}metric\text{-}space}(s) \wedge \operatorname{is\text{-}metric\text{-}space}(t)) \Rightarrow  \\ \quad \quad \operatorname{is\text{-}continuous\text{-}at}(s, t, (\lambda z \in \operatorname{pts}(s).\; c), a) \end{array}}}
\vnbprop{ms-ident-lam-in-fun}{\vnbstmt{\begin{array}{@{}l@{}} \forall s. \\ \quad \operatorname{is\text{-}metric\text{-}space}(s) \Rightarrow ((\lambda z \in \operatorname{pts}(s).\; z) \in (\operatorname{pts}(s) \to \operatorname{pts}(s))) \end{array}}}
\vnbprop{ms-identity-continuous-at}{\vnbstmt{\begin{array}{@{}l@{}} \forall s.; \forall a \in \operatorname{pts}(s). \\ \quad \operatorname{is\text{-}metric\text{-}space}(s) \Rightarrow \operatorname{is\text{-}continuous\text{-}at}(s, s, (\lambda z \in \operatorname{pts}(s).\; z), a) \end{array}}}
\vnbprop{ms-cont-transfer-ptwise-eq}{\vnbstmt{\begin{array}{@{}l@{}} \forall s, t, g, f, a. \\ \quad (\operatorname{is\text{-}continuous\text{-}at}(s, t, g, a) \wedge ((f \in (\operatorname{pts}(s) \to \operatorname{pts}(t))) \wedge \forall z \in \operatorname{pts}(s).\; (f(z) = g(z)))) \Rightarrow  \\ \quad \quad \operatorname{is\text{-}continuous\text{-}at}(s, t, f, a) \end{array}}}
\vnbprop{ms-compose-continuous-at}{\vnbstmt{\begin{array}{@{}l@{}} \forall s, t, u, f, g, a. \\ \quad (\operatorname{is\text{-}continuous\text{-}at}(s, t, f, a) \wedge \operatorname{is\text{-}continuous\text{-}at}(t, u, g, f(a))) \Rightarrow  \\ \quad \quad \operatorname{is\text{-}continuous\text{-}at}(s, u, (\lambda z \in \operatorname{pts}(s).\; g(f(z))), a) \end{array}}}
\vnbprop{ms-cont-agree-off-pt}{\vnbstmt{\begin{array}{@{}l@{}} \forall s, t, f, g, p. \\ \quad \operatorname{is\text{-}continuous\text{-}at}(s, t, f, p) \wedge \\ \quad \operatorname{is\text{-}continuous\text{-}at}(s, t, g, p) \wedge \\ \quad \forall r. \\ \quad \quad \operatorname{pos\text{-}rr}(r) \Rightarrow \exists b \in \operatorname{pts}(s).\; (\neg (b = p) \wedge (\operatorname{dist}(s)(p, b) \leq r)) \wedge \\ \quad \forall y \in \operatorname{pts}(s).\; (\neg (y = p) \Rightarrow (f(y) = g(y))) \Rightarrow \\ \quad \quad f(p) = g(p) \end{array}}}
\vnbprop{nf-sub-split}{\vnbstmt{\begin{array}{@{}l@{}} \forall k, u, v, p, q. \\ \quad (\operatorname{is\text{-}normed\text{-}field}(k) \wedge ((u \in \operatorname{carr}(k)) \wedge ((v \in \operatorname{carr}(k)) \wedge ((p \in \operatorname{carr}(k)) \wedge (q \in \operatorname{carr}(k)))))) \Rightarrow  \\ \quad \quad \operatorname{add}(k)(\operatorname{add}(k)(u, v), \operatorname{neg}(k)(\operatorname{add}(k)(p, q))) \\ \quad \quad =\; \operatorname{add}(k)(\operatorname{add}(k)(u, \operatorname{neg}(k)(p)), \operatorname{add}(k)(v, \operatorname{neg}(k)(q))) \end{array}}}
\vnbprop{nf-mul-sub-split}{\vnbstmt{\begin{array}{@{}l@{}} \forall k, u, v, p, q. \\ \quad (\operatorname{is\text{-}normed\text{-}field}(k) \wedge ((u \in \operatorname{carr}(k)) \wedge ((v \in \operatorname{carr}(k)) \wedge ((p \in \operatorname{carr}(k)) \wedge (q \in \operatorname{carr}(k)))))) \Rightarrow  \\ \quad \quad \operatorname{add}(k)(\operatorname{mul}(k)(u, v), \operatorname{neg}(k)(\operatorname{mul}(k)(p, q))) \\ \quad \quad =\; \operatorname{add}(k)(\operatorname{mul}(k)(u, \operatorname{add}(k)(v, \operatorname{neg}(k)(q))), \operatorname{mul}(k)(q, \operatorname{add}(k)(u, \operatorname{neg}(k)(p)))) \end{array}}}
\vnbprop{nf-neg-sub-split}{\vnbstmt{\begin{array}{@{}l@{}} \forall k, u, v. \\ \quad (\operatorname{is\text{-}normed\text{-}field}(k) \wedge ((u \in \operatorname{carr}(k)) \wedge (v \in \operatorname{carr}(k)))) \Rightarrow  \\ \quad \quad \operatorname{add}(k)(\operatorname{neg}(k)(u), \operatorname{neg}(k)(\operatorname{neg}(k)(v))) = \operatorname{neg}(k)(\operatorname{add}(k)(u, \operatorname{neg}(k)(v))) \end{array}}}
\vnbprop{nf-sub-anti}{\vnbstmt{\begin{array}{@{}l@{}} \forall k, u, v. \\ \quad (\operatorname{is\text{-}normed\text{-}field}(k) \wedge ((u \in \operatorname{carr}(k)) \wedge (v \in \operatorname{carr}(k)))) \Rightarrow  \\ \quad \quad \operatorname{neg}(k)(\operatorname{add}(k)(u, \operatorname{neg}(k)(v))) = \operatorname{add}(k)(v, \operatorname{neg}(k)(u)) \end{array}}}
\vnbprop{nf-norm-sub-sym}{\vnbstmt{\begin{array}{@{}l@{}} \forall k, u, v. \\ \quad (\operatorname{is\text{-}normed\text{-}field}(k) \wedge ((u \in \operatorname{carr}(k)) \wedge (v \in \operatorname{carr}(k)))) \Rightarrow  \\ \quad \quad \operatorname{fnrm}(k)(\operatorname{add}(k)(u, \operatorname{neg}(k)(v))) = \operatorname{fnrm}(k)(\operatorname{add}(k)(v, \operatorname{neg}(k)(u))) \end{array}}}
\vnbprop{nf-norm-neg-diff}{\vnbstmt{\begin{array}{@{}l@{}} \forall k, u, v. \\ \quad (\operatorname{is\text{-}normed\text{-}field}(k) \wedge ((u \in \operatorname{carr}(k)) \wedge (v \in \operatorname{carr}(k)))) \Rightarrow  \\ \quad \quad \operatorname{fnrm}(k)(\operatorname{add}(k)(\operatorname{neg}(k)(u), \operatorname{neg}(k)(\operatorname{neg}(k)(v)))) = \operatorname{fnrm}(k)(\operatorname{add}(k)(u, \operatorname{neg}(k)(v))) \end{array}}}
\vnbprop{nf-norm-sum-diff}{\vnbstmt{\begin{array}{@{}l@{}} \forall k, u, v, p, q. \\ \quad (\operatorname{is\text{-}normed\text{-}field}(k) \wedge ((u \in \operatorname{carr}(k)) \wedge ((v \in \operatorname{carr}(k)) \wedge ((p \in \operatorname{carr}(k)) \wedge (q \in \operatorname{carr}(k)))))) \Rightarrow  \\ \quad \quad \operatorname{fnrm}(k)(\operatorname{add}(k)(\operatorname{add}(k)(u, v), \operatorname{neg}(k)(\operatorname{add}(k)(p, q)))) \\ \quad \quad \leq\; \operatorname{fnrm}(k)(\operatorname{add}(k)(u, \operatorname{neg}(k)(p))) + \operatorname{fnrm}(k)(\operatorname{add}(k)(v, \operatorname{neg}(k)(q))) \end{array}}}
\vnbprop{nf-norm-prod-bound}{\vnbstmt{\begin{array}{@{}l@{}} \forall k, u, v, p, q, m, \mathit{nfk}, t, e. \\ \quad \operatorname{is\text{-}normed\text{-}field}(k) \wedge \\ \quad u \in \operatorname{carr}(k) \wedge \\ \quad v \in \operatorname{carr}(k) \wedge \\ \quad p \in \operatorname{carr}(k) \wedge \\ \quad q \in \operatorname{carr}(k) \wedge \\ \quad m \in \mathbb{R} \wedge \\ \quad \mathit{nfk} \in \mathbb{R} \wedge \\ \quad t \in \mathbb{R} \wedge \\ \quad e \in \mathbb{R} \wedge \\ \quad \operatorname{fnrm}(k)(u) \leq m \wedge \\ \quad \operatorname{fnrm}(k)(q) \leq \mathit{nfk} \wedge \\ \quad \operatorname{fnrm}(k)(\operatorname{add}(k)(u, \operatorname{neg}(k)(p))) \leq t \wedge \\ \quad \operatorname{fnrm}(k)(\operatorname{add}(k)(v, \operatorname{neg}(k)(q))) \leq t \wedge \\ \quad ((m + \mathit{nfk}) \cdot t) \leq e \Rightarrow \\ \quad \quad \operatorname{fnrm}(k)(\operatorname{add}(k)(\operatorname{mul}(k)(u, v), \operatorname{neg}(k)(\operatorname{mul}(k)(p, q)))) \leq e \end{array}}}
\vnbprop{nf-sum-lam-in-fun}{\vnbstmt{\begin{array}{@{}l@{}} \forall k, s.; \forall g, h \in (\operatorname{pts}(s) \to \operatorname{pts}(\operatorname{nf\text{-}metric\text{-}space}(k))). \\ \quad (\operatorname{is\text{-}normed\text{-}field}(k) \wedge \operatorname{is\text{-}metric\text{-}space}(s)) \Rightarrow  \\ \quad \quad (\lambda z \in \operatorname{pts}(s).\; \operatorname{add}(k)(g(z), h(z))) \in (\operatorname{pts}(s) \to \operatorname{pts}(\operatorname{nf\text{-}metric\text{-}space}(k))) \end{array}}}
\vnbprop{nf-prod-lam-in-fun}{\vnbstmt{\begin{array}{@{}l@{}} \forall k, s.; \forall g, h \in (\operatorname{pts}(s) \to \operatorname{pts}(\operatorname{nf\text{-}metric\text{-}space}(k))). \\ \quad (\operatorname{is\text{-}normed\text{-}field}(k) \wedge \operatorname{is\text{-}metric\text{-}space}(s)) \Rightarrow  \\ \quad \quad (\lambda z \in \operatorname{pts}(s).\; \operatorname{mul}(k)(g(z), h(z))) \in (\operatorname{pts}(s) \to \operatorname{pts}(\operatorname{nf\text{-}metric\text{-}space}(k))) \end{array}}}
\vnbprop{nf-neg-lam-in-fun}{\vnbstmt{\begin{array}{@{}l@{}} \forall k, s.; \forall g \in (\operatorname{pts}(s) \to \operatorname{pts}(\operatorname{nf\text{-}metric\text{-}space}(k))). \\ \quad (\operatorname{is\text{-}normed\text{-}field}(k) \wedge \operatorname{is\text{-}metric\text{-}space}(s)) \Rightarrow  \\ \quad \quad (\lambda z \in \operatorname{pts}(s).\; \operatorname{neg}(k)(g(z))) \in (\operatorname{pts}(s) \to \operatorname{pts}(\operatorname{nf\text{-}metric\text{-}space}(k))) \end{array}}}
\vnbprop{nf-sum-continuous-at}{\vnbstmt{\begin{array}{@{}l@{}} \forall k, s, g, h, a. \\ \quad \operatorname{is\text{-}normed\text{-}field}(k) \wedge \\ \quad \operatorname{is\text{-}continuous\text{-}at}(s, \operatorname{nf\text{-}metric\text{-}space}(k), g, a) \wedge \\ \quad \operatorname{is\text{-}continuous\text{-}at}(s, \operatorname{nf\text{-}metric\text{-}space}(k), h, a) \Rightarrow \\ \quad \quad \operatorname{is\text{-}continuous\text{-}at}(s, \operatorname{nf\text{-}metric\text{-}space}(k), (\lambda z \in \operatorname{pts}(s).\; \operatorname{add}(k)(g(z), h(z))), a) \end{array}}}
\vnbprop{nf-neg-continuous-at}{\vnbstmt{\begin{array}{@{}l@{}} \forall k, s, g, a. \\ \quad (\operatorname{is\text{-}normed\text{-}field}(k) \wedge \operatorname{is\text{-}continuous\text{-}at}(s, \operatorname{nf\text{-}metric\text{-}space}(k), g, a)) \Rightarrow  \\ \quad \quad \operatorname{is\text{-}continuous\text{-}at}(s, \operatorname{nf\text{-}metric\text{-}space}(k), (\lambda z \in \operatorname{pts}(s).\; \operatorname{neg}(k)(g(z))), a) \end{array}}}
\vnbprop{nf-product-continuous-at}{\vnbstmt{\begin{array}{@{}l@{}} \forall k, s, g, h, a. \\ \quad \operatorname{is\text{-}normed\text{-}field}(k) \wedge \\ \quad \operatorname{is\text{-}continuous\text{-}at}(s, \operatorname{nf\text{-}metric\text{-}space}(k), g, a) \wedge \\ \quad \operatorname{is\text{-}continuous\text{-}at}(s, \operatorname{nf\text{-}metric\text{-}space}(k), h, a) \Rightarrow \\ \quad \quad \operatorname{is\text{-}continuous\text{-}at}(s, \operatorname{nf\text{-}metric\text{-}space}(k), (\lambda z \in \operatorname{pts}(s).\; \operatorname{mul}(k)(g(z), h(z))), a) \end{array}}}

\vnbbookfile{theorem-library/ms-compose-continuous}{theorem-library/ms-compose-continuous}
\noindent{\small compose.scm -- COMPOSITION OF CONTINUOUS MAPS BETWEEN METRIC SPACES, the global form.  Agent KIT-8, 2026-10-26. Destination: theorem-library/ (library analysis), right after the file proving ms-compose-continuous-at (theorem-library/ms-continuity-algebra.scm).  The POINTWISE theorems already exist and are NOT restated: ...}\par
\vnbprop{ms-compose-continuous}{\vnbstmt{\begin{array}{@{}l@{}} \forall s, m, u, f, g. \\ \quad (\operatorname{is\text{-}continuous}(s, m, f) \wedge \operatorname{is\text{-}continuous}(m, u, g)) \Rightarrow  \\ \quad \quad \operatorname{is\text{-}continuous}(s, u, (\lambda z \in \operatorname{pts}(s).\; g(f(z)))) \end{array}}}
\vnbprop{ms-compose-continuous-dom}{\vnbstmt{\begin{array}{@{}l@{}} \forall s, d, m, u, f, g. \\ \quad ((\operatorname{pts}(s) \simeq d) \wedge (\operatorname{is\text{-}continuous}(s, m, f) \wedge \operatorname{is\text{-}continuous}(m, u, g))) \Rightarrow  \\ \quad \quad \operatorname{is\text{-}continuous}(s, u, (\lambda z \in d.\; g(f(z)))) \end{array}}}

\vnbbookfile{structure-library/riesz-markov-model}{structure-library/riesz-markov-model}
\noindent{\small structure-library/riesz-markov-model.scm -- the Cantor model of a compact metric space (Riesz-Markov, stage 1, RM-1, 2026-10-10).  DEFINITIONS ONLY; every law is PROVEN in theorem-library/riesz-markov-model.scm (prefix rmm-).}\par
\par\medskip\noindent\textbf{Definition} (\texttt{rm-ball}).\;$\operatorname{rm\text{-}ball}(s, \mathit{dseq}, n) :=$ $\operatorname{closed\text{-}ball}(s, \operatorname{dseq}(\operatorname{nnfst}(n)), \operatorname{recip}((\operatorname{nnsnd}(n) + 1)))$
\par\medskip\noindent\textbf{Definition} (\texttt{rm-shell}).\;$\operatorname{rm\text{-}shell}(s, \mathit{dseq}, n) :=$
\begin{equation*}
\{\operatorname{rsy\_} \in \operatorname{pts}(s) : (\operatorname{recip}((\operatorname{nnsnd}(n) + 1)) \leq \operatorname{dist}(s)(\operatorname{dseq}(\operatorname{nnfst}(n)), \operatorname{rsy\_}))\}
\end{equation*}

\par\medskip\noindent\textbf{Definition} (\texttt{rm-two}).\; $(\operatorname{rm\text{-}two} \simeq \{0, 1\})$
\par\medskip\noindent\textbf{Definition} (\texttt{rm-two-ms}).\; $(\operatorname{rm\text{-}two\text{-}ms} \simeq \operatorname{subspace\text{-}ms}(\operatorname{rr\text{-}ms}, \operatorname{rm\text{-}two}))$
\par\medskip\noindent\textbf{Definition} (\texttt{cantor-ms}).\; $(\operatorname{cantor\text{-}ms} \simeq \operatorname{product\text{-}metric}((\lambda \operatorname{ctk\_} \in \mathbb{N}.\; \operatorname{rm\text{-}two\text{-}ms})))$
\par\medskip\noindent\textbf{Definition} (\texttt{rm-realizes}).\;For all $s$, $\mathit{dseq}$, $\mathit{es}$, $\mathit{nb}$ and $y$, 
\begin{equation*}
\begin{array}{@{}l@{}} \operatorname{rm\text{-}realizes}(s, \mathit{dseq}, \mathit{es}, \mathit{nb}, y) \\ \Leftrightarrow\; ((y \in \operatorname{pts}(s)) \wedge \forall \operatorname{rri\_} \in \mathbb{N}.\; ((\operatorname{rri\_} < \mathit{nb}) \Rightarrow ((y \in \operatorname{rm\text{-}ball}(s, \mathit{dseq}, \operatorname{rri\_})) \Leftrightarrow (\operatorname{es}(\operatorname{rri\_}) = 1)))) \end{array}
\end{equation*}

\par\medskip\noindent\textbf{Definition} (\texttt{rm-k}).\;$\operatorname{rm\text{-}k}(s, \mathit{dseq}) :=$
\begin{equation*}
\{\operatorname{rke\_} \in \operatorname{pts}(\operatorname{cantor\text{-}ms}) : \forall \operatorname{rkn\_} \in \mathbb{N}.\; \exists \operatorname{rky\_} \in \operatorname{pts}(s).\; \operatorname{rm\text{-}realizes}(s, \mathit{dseq}, \operatorname{rke\_}, \operatorname{rkn\_}, \operatorname{rky\_})\}
\end{equation*}

\par\medskip\noindent\textbf{Definition} (\texttt{rm-kms}).\;$\operatorname{rm\text{-}kms}(s, \mathit{dseq}) :=$ $\operatorname{subspace\text{-}ms}(\operatorname{cantor\text{-}ms}, \operatorname{rm\text{-}k}(s, \mathit{dseq}))$
\par\medskip\noindent\textbf{Definition} (\texttt{rm-address}).\;For all $s$, $\mathit{dseq}$, $\mathit{es}$ and $y$, 
\begin{equation*}
\begin{array}{@{}l@{}} \operatorname{rm\text{-}address}(s, \mathit{dseq}, \mathit{es}, y) \\ \Leftrightarrow\; y \in \operatorname{pts}(s) \wedge \\ \forall \operatorname{ran\_} \in \mathbb{N}. \\ \quad (\operatorname{es}(\operatorname{ran\_}) = 1) \Rightarrow (y \in \operatorname{rm\text{-}ball}(s, \mathit{dseq}, \operatorname{ran\_})) \wedge \\ \quad \neg (\operatorname{es}(\operatorname{ran\_}) = 1) \Rightarrow (y \in \operatorname{rm\text{-}shell}(s, \mathit{dseq}, \operatorname{ran\_})) \end{array}
\end{equation*}

\par\medskip\noindent\textbf{Definition} (\texttt{rm-address-below}).\;For all $s$, $\mathit{dseq}$, $\mathit{es}$, $\mathit{nb}$ and $y$, 
\begin{equation*}
\begin{array}{@{}l@{}} \operatorname{rm\text{-}address\text{-}below}(s, \mathit{dseq}, \mathit{es}, \mathit{nb}, y) \\ \Leftrightarrow\; y \in \operatorname{pts}(s) \wedge \\ \forall \operatorname{rbn\_} \in \mathbb{N}. \\ \quad (\operatorname{rbn\_} < \mathit{nb}) \Rightarrow  \\ \quad \quad (\operatorname{es}(\operatorname{rbn\_}) = 1) \Rightarrow (y \in \operatorname{rm\text{-}ball}(s, \mathit{dseq}, \operatorname{rbn\_})) \wedge \\ \quad \quad \neg (\operatorname{es}(\operatorname{rbn\_}) = 1) \Rightarrow (y \in \operatorname{rm\text{-}shell}(s, \mathit{dseq}, \operatorname{rbn\_})) \end{array}
\end{equation*}

\par\medskip\noindent\textbf{Definition} (\texttt{rm-pi}).\;$\operatorname{rm\text{-}pi}(s, \mathit{dseq}) :=$
\begin{equation*}
(\lambda \operatorname{rpe\_} \in \operatorname{rm\text{-}k}(s, \mathit{dseq}).\; (\iota\, \operatorname{rpy\_}.\;((\operatorname{rpy\_} \in \operatorname{pts}(s)) \wedge \operatorname{rm\text{-}address}(s, \mathit{dseq}, \operatorname{rpe\_}, \operatorname{rpy\_}))))
\end{equation*}

\par\medskip\noindent\textbf{Definition} (\texttt{rm-code}).\;$\operatorname{rm\text{-}code}(s, \mathit{dseq}, x) :=$ $(\lambda \operatorname{rox\_} \in \mathbb{N}.\; \left(\text{if }(x \in \operatorname{rm\text{-}ball}(s, \mathit{dseq}, \operatorname{rox\_}))\text{ then }1\text{ else }0\right))$
\par\medskip\noindent\textbf{Definition} (\texttt{rm-cyl}).\;$\operatorname{rm\text{-}cyl}(s, \mathit{dseq}, a, \mathit{nb}) :=$
\begin{equation*}
\{\operatorname{rce\_} \in \operatorname{rm\text{-}k}(s, \mathit{dseq}) : \forall \operatorname{rci\_} \in \mathbb{N}.\; ((\operatorname{rci\_} < \mathit{nb}) \Rightarrow (\operatorname{rce\_}(\operatorname{rci\_}) = a(\operatorname{rci\_})))\}
\end{equation*}


\vnbbookfile{theorem-library/riesz-markov-model}{theorem-library/riesz-markov-model}
\noindent{\small theorem-library/riesz-markov-model.scm -- RIESZ-MARKOV, STAGE 1 (RM-1, 2026-10-10): the Cantor model K of a compact metric space and the map pi : K -$>$ s.  Vocabulary: structure-library/riesz-markov-model.scm. Source: the user's notes-63, concrete form (b).  Helper prefix: rmm-.}\par
\vnbprop{rmm-two-mem}{\vnbstmt{\forall x.\; ((x \in \operatorname{rm\text{-}two}) \Leftrightarrow ((x = 0) \vee (x = 1)))}}
\vnbprop{rmm-two-cases}{\vnbstmt{\forall x \in \operatorname{rm\text{-}two}.\; ((x = 0) \vee (x = 1))}}
\vnbprop{rmm-two-makeset}{\vnbstmt{\operatorname{rm\text{-}two} = \operatorname{make\text{-}set}({[}0, 1])}}
\vnbprop{rmm-two-sub-rr}{\vnbstmt{\operatorname{rm\text{-}two} \subseteq \mathbb{R}}}
\vnbprop{rmm-two-ms-metric}{\vnbstmt{\operatorname{is\text{-}metric\text{-}space}(\operatorname{rm\text{-}two\text{-}ms})}}
\vnbprop{rmm-two-ms-pts}{\vnbstmt{(\operatorname{pts}(\operatorname{rm\text{-}two\text{-}ms}) \simeq \operatorname{rm\text{-}two})}}
\vnbprop{rmm-two-ms-dist}{\vnbstmt{\forall u, v \in \operatorname{rm\text{-}two}.\; (\operatorname{dist}(\operatorname{rm\text{-}two\text{-}ms})(u, v) \simeq \lvert (u - v) \rvert)}}
\vnbprop{rmm-two-near-eq}{\vnbstmt{\forall u, v \in \operatorname{rm\text{-}two}.\; ((\operatorname{dist}(\operatorname{rm\text{-}two\text{-}ms})(u, v) < 1) \Rightarrow (u = v))}}
\vnbprop{rmm-two-ms-set}{\vnbstmt{\operatorname{rm\text{-}two\text{-}ms} \in \mathbf{Set}}}
\vnbprop{rmm-ms-seq}{\vnbstmt{\operatorname{is\text{-}ms\text{-}sequence}((\lambda k \in \mathbb{N}.\; \operatorname{rm\text{-}two\text{-}ms}))}}
\vnbprop{rmm-cantor-metric}{\vnbstmt{\operatorname{is\text{-}metric\text{-}space}(\operatorname{cantor\text{-}ms})}}
\vnbprop{rmm-cantor-pts}{\vnbstmt{(\operatorname{pts}(\operatorname{cantor\text{-}ms}) \simeq \operatorname{product\text{-}carrier}((\lambda k \in \mathbb{N}.\; \operatorname{rm\text{-}two\text{-}ms})))}}
\vnbprop{rmm-cantor-coord}{\vnbstmt{\forall s \in \operatorname{pts}(\operatorname{cantor\text{-}ms}), n \in \mathbb{N}.\; (s(n) \in \operatorname{rm\text{-}two})}}
\vnbprop{rmm-two-seq-compact}{\vnbstmt{\operatorname{seq\text{-}compact}(\operatorname{rm\text{-}two\text{-}ms})}}
\vnbprop{rmm-two-compact}{\vnbstmt{\operatorname{is\text{-}compact}(\operatorname{rm\text{-}two\text{-}ms})}}
\vnbprop{rmm-cantor-compact}{\vnbstmt{\operatorname{is\text{-}compact}(\operatorname{cantor\text{-}ms})}}
\vnbprop{rmm-cantor-unfold}{\vnbstmt{(\operatorname{cantor\text{-}ms} \simeq \operatorname{product\text{-}metric\text{-}w}((\lambda k \in \mathbb{N}.\; \operatorname{rm\text{-}two\text{-}ms}), (\lambda n \in \mathbb{N}.\; /(1, \operatorname{power}(2, (n + 1))))))}}
\vnbprop{rmm-proj-value}{\vnbstmt{\begin{array}{@{}l@{}} \forall n \in \mathbb{N}, s \in \operatorname{pts}(\operatorname{cantor\text{-}ms}). \\ \quad \operatorname{product\text{-}proj}((\lambda k \in \mathbb{N}.\; \operatorname{rm\text{-}two\text{-}ms}), n)(s) = s(n) \end{array}}}
\vnbprop{rmm-coord-continuous}{\vnbstmt{\begin{array}{@{}l@{}} \forall n \in \mathbb{N}. \\ \quad \operatorname{is\text{-}continuous}(\operatorname{cantor\text{-}ms}, \operatorname{rm\text{-}two\text{-}ms}, \operatorname{product\text{-}proj}((\lambda k \in \mathbb{N}.\; \operatorname{rm\text{-}two\text{-}ms}), n)) \end{array}}}
\vnbprop{rmm-coord-near}{\vnbstmt{\begin{array}{@{}l@{}} \forall n \in \mathbb{N}, s \in \operatorname{pts}(\operatorname{cantor\text{-}ms}). \exists d. \\ \quad (\operatorname{pos\text{-}rr}(d) \wedge \forall b \in \operatorname{pts}(\operatorname{cantor\text{-}ms}).\; ((\operatorname{dist}(\operatorname{cantor\text{-}ms})(s, b) \leq d) \Rightarrow (b(n) = s(n)))) \end{array}}}
\vnbprop{rmm-agree-near}{\vnbstmt{\begin{array}{@{}l@{}} \forall b \in \mathbb{N}, s \in \operatorname{pts}(\operatorname{cantor\text{-}ms}). \exists d. \\ \quad \operatorname{pos\text{-}rr}(d) \wedge \\ \quad \forall \mathit{rab} \in \operatorname{pts}(\operatorname{cantor\text{-}ms}), i \in \mathbb{N}. \\ \quad \quad ((\operatorname{dist}(\operatorname{cantor\text{-}ms})(s, \mathit{rab}) \leq d) \wedge (i < b)) \Rightarrow (\operatorname{rab}(i) = s(i)) \end{array}}}
\vnbprop{rmm-dense-parts}{\vnbstmt{\begin{array}{@{}l@{}} \forall s, q. \\ \quad \operatorname{is\text{-}dense\text{-}seq}(s, q) \Rightarrow (\operatorname{is\text{-}metric\text{-}space}(s) \wedge (q \in (\mathbb{N} \to \operatorname{pts}(s)))) \end{array}}}
\vnbprop{rmm-dense-near}{\vnbstmt{\begin{array}{@{}l@{}} \forall s, q.; \forall x \in \operatorname{pts}(s), a. \\ \quad (\operatorname{is\text{-}dense\text{-}seq}(s, q) \wedge \operatorname{pos\text{-}rr}(a)) \Rightarrow  \\ \quad \quad \exists m \in \mathbb{N}.\; (\operatorname{dist}(s)(x, q(m)) < a) \end{array}}}
\vnbprop{rmm-centre-in}{\vnbstmt{\forall s, q.; \forall n \in \mathbb{N}.\; (\operatorname{is\text{-}dense\text{-}seq}(s, q) \Rightarrow (q(\operatorname{nnfst}(n)) \in \operatorname{pts}(s)))}}
\vnbprop{rmm-rad-pos}{\vnbstmt{\forall n \in \mathbb{N}.\; \operatorname{pos\text{-}rr}(\operatorname{recip}((\operatorname{nnsnd}(n) + 1)))}}
\vnbprop{rmm-ball-mem}{\vnbstmt{\begin{array}{@{}l@{}} \forall s, q, n, y. \\ \quad (y \in \operatorname{rm\text{-}ball}(s, q, n)) \Leftrightarrow ((y \in \operatorname{pts}(s)) \wedge (\operatorname{dist}(s)(q(\operatorname{nnfst}(n)), y) \leq \operatorname{recip}((\operatorname{nnsnd}(n) + 1)))) \end{array}}}
\vnbprop{rmm-shell-mem}{\vnbstmt{\begin{array}{@{}l@{}} \forall s, q, n, y. \\ \quad (y \in \operatorname{rm\text{-}shell}(s, q, n)) \Leftrightarrow ((y \in \operatorname{pts}(s)) \wedge (\operatorname{recip}((\operatorname{nnsnd}(n) + 1)) \leq \operatorname{dist}(s)(q(\operatorname{nnfst}(n)), y))) \end{array}}}
\vnbprop{rmm-near-centre}{\vnbstmt{\begin{array}{@{}l@{}} \forall s, q.; \forall x \in \operatorname{pts}(s), j \in \mathbb{N}. \\ \quad \operatorname{is\text{-}dense\text{-}seq}(s, q) \Rightarrow  \\ \quad \quad \exists n \in \mathbb{N}. \\ \quad \quad \quad (\operatorname{nnsnd}(n) = j) \wedge (\operatorname{dist}(s)(q(\operatorname{nnfst}(n)), x) < \operatorname{recip}((\operatorname{nnsnd}(n) + 1))) \end{array}}}
\vnbprop{rmm-inside-ball}{\vnbstmt{\begin{array}{@{}l@{}} \forall s, q.; \forall n \in \mathbb{N}, x \in \operatorname{pts}(s). \\ \quad (\operatorname{is\text{-}dense\text{-}seq}(s, q) \wedge (\operatorname{dist}(s)(q(\operatorname{nnfst}(n)), x) < \operatorname{recip}((\operatorname{nnsnd}(n) + 1)))) \Rightarrow  \\ \quad \quad (x \in \operatorname{rm\text{-}ball}(s, q, n)) \wedge \neg (x \in \operatorname{rm\text{-}shell}(s, q, n)) \end{array}}}
\vnbprop{rmm-strong-separation}{\vnbstmt{\begin{array}{@{}l@{}} \forall s, q.; \forall x, y \in \operatorname{pts}(s). \\ \quad (\operatorname{is\text{-}dense\text{-}seq}(s, q) \wedge \neg (x = y)) \Rightarrow  \\ \quad \quad \exists n, m \in \mathbb{N}. \\ \quad \quad \quad x \in \operatorname{rm\text{-}ball}(s, q, n) \wedge \\ \quad \quad \quad \neg (x \in \operatorname{rm\text{-}shell}(s, q, n)) \wedge \\ \quad \quad \quad y \in \operatorname{rm\text{-}ball}(s, q, m) \wedge \\ \quad \quad \quad \neg (y \in \operatorname{rm\text{-}shell}(s, q, m)) \wedge \\ \quad \quad \quad \forall z.\; \neg ((z \in \operatorname{rm\text{-}ball}(s, q, n)) \wedge (z \in \operatorname{rm\text{-}ball}(s, q, m))) \end{array}}}
\vnbprop{rmm-k-mem}{\vnbstmt{\begin{array}{@{}l@{}} \forall s, q, \mathit{es}. \\ \quad \mathit{es} \in \operatorname{rm\text{-}k}(s, q) \\ \quad \Leftrightarrow\; ((\mathit{es} \in \operatorname{pts}(\operatorname{cantor\text{-}ms})) \wedge \forall n \in \mathbb{N}.\; \exists y.\; \operatorname{rm\text{-}realizes}(s, q, \mathit{es}, n, y)) \end{array}}}
\vnbprop{rmm-k-sub}{\vnbstmt{\forall s, q.\; (\operatorname{rm\text{-}k}(s, q) \subseteq \operatorname{pts}(\operatorname{cantor\text{-}ms}))}}
\vnbprop{rmm-realizes-agree}{\vnbstmt{\begin{array}{@{}l@{}} \forall s, q, \mathit{es}, b, a, y. \\ \quad (\forall i \in \mathbb{N}.\; ((i < a) \Rightarrow (b(i) = \operatorname{es}(i))) \wedge \operatorname{rm\text{-}realizes}(s, q, \mathit{es}, a, y)) \Rightarrow  \\ \quad \quad \operatorname{rm\text{-}realizes}(s, q, b, a, y) \end{array}}}
\vnbprop{rmm-k-out}{\vnbstmt{\begin{array}{@{}l@{}} \forall s, q.; \forall \mathit{es} \in \operatorname{pts}(\operatorname{cantor\text{-}ms}). \\ \quad \neg (\mathit{es} \in \operatorname{rm\text{-}k}(s, q)) \Rightarrow  \\ \quad \quad \exists n \in \mathbb{N}.\; \forall y.\; \neg \operatorname{rm\text{-}realizes}(s, q, \mathit{es}, n, y) \end{array}}}
\vnbprop{rmm-k-closed}{\vnbstmt{\forall s, q.\; \operatorname{is\text{-}closed}(\operatorname{cantor\text{-}ms}, \operatorname{rm\text{-}k}(s, q))}}
\vnbprop{rmm-k-set}{\vnbstmt{\forall s, q.\; (\operatorname{rm\text{-}k}(s, q) \in \mathbf{Set})}}
\vnbprop{rmm-k-compact}{\vnbstmt{\forall s, q.\; \operatorname{is\text{-}compact}(\operatorname{rm\text{-}kms}(s, q))}}
\vnbprop{rmm-kms-metric}{\vnbstmt{\forall s, q.\; \operatorname{is\text{-}metric\text{-}space}(\operatorname{rm\text{-}kms}(s, q))}}
\vnbprop{rmm-kms-pts}{\vnbstmt{\forall s, q.\; (\operatorname{pts}(\operatorname{rm\text{-}kms}(s, q)) \simeq \operatorname{rm\text{-}k}(s, q))}}
\vnbprop{rmm-nn-choice}{\vnbstmt{\begin{array}{@{}l@{}} \forall a, m. \\ \quad \forall k \in \mathbb{N}.\; \exists w \in a.\; (w \in m(k)) \Rightarrow  \\ \quad \quad \exists g \in (\mathbb{N} \to a).\; \forall k \in \mathbb{N}.\; (g(k) \in m(k)) \end{array}}}
\vnbprop{rmm-address-unique}{\vnbstmt{\begin{array}{@{}l@{}} \forall s, q, \mathit{es}, y_{1}, y_{2}. \\ \quad (\operatorname{is\text{-}dense\text{-}seq}(s, q) \wedge (\operatorname{rm\text{-}address}(s, q, \mathit{es}, y_{1}) \wedge \operatorname{rm\text{-}address}(s, q, \mathit{es}, y_{2}))) \Rightarrow  \\ \quad \quad y_{1} = y_{2} \end{array}}}
\vnbprop{rmm-shell-closed}{\vnbstmt{\begin{array}{@{}l@{}} \forall s, q.; \forall n \in \mathbb{N}. \\ \quad \operatorname{is\text{-}dense\text{-}seq}(s, q) \Rightarrow \operatorname{is\text{-}closed}(s, \operatorname{rm\text{-}shell}(s, q, n)) \end{array}}}
\vnbprop{rmm-realizes-pts}{\vnbstmt{\forall s, q, \mathit{es}, b, y.\; (\operatorname{rm\text{-}realizes}(s, q, \mathit{es}, b, y) \Rightarrow (y \in \operatorname{pts}(s)))}}
\vnbprop{rmm-ball-closed}{\vnbstmt{\begin{array}{@{}l@{}} \forall s, q.; \forall n \in \mathbb{N}. \\ \quad \operatorname{is\text{-}dense\text{-}seq}(s, q) \Rightarrow \operatorname{is\text{-}closed}(s, \operatorname{rm\text{-}ball}(s, q, n)) \end{array}}}
\vnbprop{rmm-address-exists}{\vnbstmt{\begin{array}{@{}l@{}} \forall s, q.; \forall \mathit{es} \in \operatorname{rm\text{-}k}(s, q). \\ \quad (\operatorname{is\text{-}compact}(s) \wedge \operatorname{is\text{-}dense\text{-}seq}(s, q)) \Rightarrow  \\ \quad \quad \exists y.\; \operatorname{rm\text{-}address}(s, q, \mathit{es}, y) \end{array}}}
\vnbprop{rmm-address-eu}{\vnbstmt{\begin{array}{@{}l@{}} \forall s, q.; \forall \mathit{es} \in \operatorname{rm\text{-}k}(s, q). \\ \quad (\operatorname{is\text{-}compact}(s) \wedge \operatorname{is\text{-}dense\text{-}seq}(s, q)) \Rightarrow  \\ \quad \quad \exists y. \\ \quad \quad \quad \operatorname{rm\text{-}address}(s, q, \mathit{es}, y) \wedge \forall z.\; (\operatorname{rm\text{-}address}(s, q, \mathit{es}, z) \Rightarrow (y = z)) \end{array}}}
\vnbprop{rmm-address-eu-pts}{\vnbstmt{\begin{array}{@{}l@{}} \forall s, q.; \forall \mathit{es} \in \operatorname{rm\text{-}k}(s, q). \\ \quad (\operatorname{is\text{-}compact}(s) \wedge \operatorname{is\text{-}dense\text{-}seq}(s, q)) \Rightarrow  \\ \quad \quad \exists y. \\ \quad \quad \quad y \in \operatorname{pts}(s) \wedge \\ \quad \quad \quad \operatorname{rm\text{-}address}(s, q, \mathit{es}, y) \wedge \\ \quad \quad \quad \forall z \in \operatorname{pts}(s).\; (\operatorname{rm\text{-}address}(s, q, \mathit{es}, z) \Rightarrow (y = z)) \end{array}}}
\vnbprop{rmm-pi-spec}{\vnbstmt{\begin{array}{@{}l@{}} \forall s, q.; \forall \mathit{es} \in \operatorname{rm\text{-}k}(s, q). \\ \quad (\operatorname{is\text{-}compact}(s) \wedge \operatorname{is\text{-}dense\text{-}seq}(s, q)) \Rightarrow  \\ \quad \quad \operatorname{rm\text{-}address}(s, q, \mathit{es}, \operatorname{rm\text{-}pi}(s, q)(\mathit{es})) \end{array}}}
\vnbprop{rmm-pi-value}{\vnbstmt{\begin{array}{@{}l@{}} \forall s, q.; \forall \mathit{es} \in \operatorname{rm\text{-}k}(s, q), y. \\ \quad (\operatorname{is\text{-}compact}(s) \wedge (\operatorname{is\text{-}dense\text{-}seq}(s, q) \wedge \operatorname{rm\text{-}address}(s, q, \mathit{es}, y))) \Rightarrow  \\ \quad \quad \operatorname{rm\text{-}pi}(s, q)(\mathit{es}) = y \end{array}}}
\vnbprop{rmm-pi-fun}{\vnbstmt{\begin{array}{@{}l@{}} \forall s, q. \\ \quad (\operatorname{is\text{-}compact}(s) \wedge \operatorname{is\text{-}dense\text{-}seq}(s, q)) \Rightarrow  \\ \quad \quad \operatorname{rm\text{-}pi}(s, q) \in (\operatorname{rm\text{-}k}(s, q) \to \operatorname{pts}(s)) \end{array}}}
\vnbprop{rmm-code-value}{\vnbstmt{\forall s, q, x.; \forall n \in \mathbb{N}.\; ((\operatorname{rm\text{-}code}(s, q, x)(n) = 1) \Leftrightarrow (x \in \operatorname{rm\text{-}ball}(s, q, n)))}}
\vnbprop{rmm-code-two}{\vnbstmt{\forall s, q, x.\; (\operatorname{rm\text{-}code}(s, q, x) \in (\mathbb{N} \to \operatorname{rm\text{-}two}))}}
\vnbprop{rmm-cantor-pts-intro}{\vnbstmt{\forall s \in (\mathbb{N} \to \operatorname{rm\text{-}two}).\; (s \in \operatorname{pts}(\operatorname{cantor\text{-}ms}))}}
\vnbprop{rmm-code-in-k}{\vnbstmt{\begin{array}{@{}l@{}} \forall s, q.; \forall x \in \operatorname{pts}(s). \\ \quad \operatorname{is\text{-}dense\text{-}seq}(s, q) \Rightarrow (\operatorname{rm\text{-}code}(s, q, x) \in \operatorname{rm\text{-}k}(s, q)) \end{array}}}
\vnbprop{rmm-code-address}{\vnbstmt{\begin{array}{@{}l@{}} \forall s, q.; \forall x \in \operatorname{pts}(s). \\ \quad \operatorname{is\text{-}dense\text{-}seq}(s, q) \Rightarrow \operatorname{rm\text{-}address}(s, q, \operatorname{rm\text{-}code}(s, q, x), x) \end{array}}}
\vnbprop{rmm-pi-code}{\vnbstmt{\begin{array}{@{}l@{}} \forall s, q.; \forall x \in \operatorname{pts}(s). \\ \quad (\operatorname{is\text{-}compact}(s) \wedge \operatorname{is\text{-}dense\text{-}seq}(s, q)) \Rightarrow (\operatorname{rm\text{-}pi}(s, q)(\operatorname{rm\text{-}code}(s, q, x)) = x) \end{array}}}
\vnbprop{rmm-pi-onto}{\vnbstmt{\begin{array}{@{}l@{}} \forall s, q.; \forall x \in \operatorname{pts}(s). \\ \quad (\operatorname{is\text{-}compact}(s) \wedge \operatorname{is\text{-}dense\text{-}seq}(s, q)) \Rightarrow  \\ \quad \quad \exists \mathit{es} \in \operatorname{rm\text{-}k}(s, q).\; (\operatorname{rm\text{-}pi}(s, q)(\mathit{es}) = x) \end{array}}}
\vnbprop{rmm-address-below}{\vnbstmt{\begin{array}{@{}l@{}} \forall s, q, \mathit{es}, y.; \forall b \in \mathbb{N}. \\ \quad \operatorname{rm\text{-}address}(s, q, \mathit{es}, y) \Rightarrow \operatorname{rm\text{-}address\text{-}below}(s, q, \mathit{es}, b, y) \end{array}}}
\vnbprop{rmm-below-agree}{\vnbstmt{\begin{array}{@{}l@{}} \forall s, q, \mathit{es}, b, a, y. \\ \quad (\forall i \in \mathbb{N}.\; ((i < a) \Rightarrow (b(i) = \operatorname{es}(i))) \wedge \operatorname{rm\text{-}address\text{-}below}(s, q, b, a, y)) \Rightarrow  \\ \quad \quad \operatorname{rm\text{-}address\text{-}below}(s, q, \mathit{es}, a, y) \end{array}}}
\vnbprop{rmm-shrink}{\vnbstmt{\begin{array}{@{}l@{}} \forall s, q.; \forall \mathit{es} \in \operatorname{rm\text{-}k}(s, q), a. \\ \quad (\operatorname{is\text{-}compact}(s) \wedge (\operatorname{is\text{-}dense\text{-}seq}(s, q) \wedge \operatorname{pos\text{-}rr}(a))) \Rightarrow  \\ \quad \quad \exists n \in \mathbb{N}. \forall y. \\ \quad \quad \quad \operatorname{rm\text{-}address\text{-}below}(s, q, \mathit{es}, n, y) \Rightarrow (\operatorname{dist}(s)(\operatorname{rm\text{-}pi}(s, q)(\mathit{es}), y) \leq a) \end{array}}}
\vnbprop{rmm-kms-dist}{\vnbstmt{\begin{array}{@{}l@{}} \forall s, q.; \forall u, v \in \operatorname{rm\text{-}k}(s, q). \\ \quad (\operatorname{dist}(\operatorname{rm\text{-}kms}(s, q))(u, v) \simeq \operatorname{dist}(\operatorname{cantor\text{-}ms})(u, v)) \end{array}}}
\vnbprop{rmm-pi-continuous}{\vnbstmt{\begin{array}{@{}l@{}} \forall s, q. \\ \quad (\operatorname{is\text{-}compact}(s) \wedge \operatorname{is\text{-}dense\text{-}seq}(s, q)) \Rightarrow  \\ \quad \quad \operatorname{is\text{-}continuous}(\operatorname{rm\text{-}kms}(s, q), s, \operatorname{rm\text{-}pi}(s, q)) \end{array}}}
\vnbprop{rmm-cyl-mem}{\vnbstmt{\begin{array}{@{}l@{}} \forall s, q, a, b, \mathit{es}. \\ \quad (\mathit{es} \in \operatorname{rm\text{-}cyl}(s, q, a, b)) \Leftrightarrow ((\mathit{es} \in \operatorname{rm\text{-}k}(s, q)) \wedge \forall i \in \mathbb{N}.\; ((i < b) \Rightarrow (\operatorname{es}(i) = a(i)))) \end{array}}}
\vnbprop{rmm-cyl-sub}{\vnbstmt{\forall s, q, a, b.\; (\operatorname{rm\text{-}cyl}(s, q, a, b) \subseteq \operatorname{rm\text{-}k}(s, q))}}
\vnbprop{rmm-cyl-open}{\vnbstmt{\forall s, q, a.; \forall b \in \mathbb{N}.\; \operatorname{is\text{-}open}(\operatorname{rm\text{-}kms}(s, q), \operatorname{rm\text{-}cyl}(s, q, a, b))}}
\vnbprop{rmm-cyl-closed}{\vnbstmt{\forall s, q, a.; \forall b \in \mathbb{N}.\; \operatorname{is\text{-}closed}(\operatorname{rm\text{-}kms}(s, q), \operatorname{rm\text{-}cyl}(s, q, a, b))}}
\vnbprop{rmm-kms-unfold}{\vnbstmt{\forall s, q.\; (\operatorname{rm\text{-}kms}(s, q) \simeq \operatorname{subspace\text{-}ms}(\operatorname{cantor\text{-}ms}, \operatorname{rm\text{-}k}(s, q)))}}
\vnbprop{rmm-cyl-base}{\vnbstmt{\begin{array}{@{}l@{}} \forall s, q, u.; \forall \mathit{es} \in u. \\ \quad \operatorname{is\text{-}open}(\operatorname{rm\text{-}kms}(s, q), u) \Rightarrow  \\ \quad \quad \exists k \in \mathbb{N}.\; (\operatorname{rm\text{-}cyl}(s, q, \mathit{es}, k) \subseteq u) \end{array}}}
\vnbprop{rmm-cyl-self}{\vnbstmt{\forall s, q.; \forall \mathit{es} \in \operatorname{rm\text{-}k}(s, q), b.\; (\mathit{es} \in \operatorname{rm\text{-}cyl}(s, q, \mathit{es}, b))}}
\vnbprop{rmm-dense-seq-exists}{\vnbstmt{\begin{array}{@{}l@{}} \forall s. \\ \quad (\operatorname{is\text{-}compact}(s) \wedge \exists x.\; (x \in \operatorname{pts}(s))) \Rightarrow  \\ \quad \quad \exists q.\; \operatorname{is\text{-}dense\text{-}seq}(s, q) \end{array}}}

\vnbbookfile{theorem-library/rabasics-closed}{theorem-library/rabasics-closed}
\noindent{\small rabasics-closed.scm -- calculus.pdf 3.9, complex-analysis.pdf 1.9: a set A in a metric space X is closed iff every sequence of elements of A which converges, converges to an element of A.}\par
\vnbprop{seq-closed-implies-closed}{\vnbstmt{\begin{array}{@{}l@{}} \forall s, a. \\ \quad \operatorname{is\text{-}metric\text{-}space}(s) \wedge \\ \quad a \subseteq \operatorname{pts}(s) \wedge \\ \quad \forall f \in (\mathbb{N} \to a), l.\; (\operatorname{converges\text{-}to}(s, f, l) \Rightarrow (l \in a)) \Rightarrow \\ \quad \quad \operatorname{is\text{-}closed}(s, a) \end{array}}}
\vnbprop{closed-iff-sequentially-closed}{\vnbstmt{\begin{array}{@{}l@{}} \forall s, a. \\ \quad (\operatorname{is\text{-}metric\text{-}space}(s) \wedge (a \subseteq \operatorname{pts}(s))) \Rightarrow  \\ \quad \quad \operatorname{is\text{-}closed}(s, a) \Leftrightarrow \forall f \in (\mathbb{N} \to a), l.\; (\operatorname{converges\text{-}to}(s, f, l) \Rightarrow (l \in a)) \end{array}}}

\vnbbookfile{structure-library/ca51-vocab}{structure-library/ca51-vocab}
\noindent{\small ca51-vocab.scm -- local compactness and exhaustions by compact sets (CA-51, 2026-10-26). DEFINITIONS ONLY.  Destination: structure-library/ca51-vocab.scm, library analysis. Dependencies: metric-topology (CLOSED-BALL, BALL), compactness (IS-COMPACT), metric-subspace (SUBSPACE-MS).  Load window: after ...}\par
\par\medskip\noindent\textbf{Definition} (\texttt{is-locally-compact}).\;For every $\operatorname{lcs\_}$, 
\begin{equation*}
\begin{array}{@{}l@{}} \operatorname{is\text{-}locally\text{-}compact}(\operatorname{lcs\_}) \\ \Leftrightarrow\; \operatorname{is\text{-}metric\text{-}space}(\operatorname{lcs\_}) \wedge \\ \forall \operatorname{lca\_} \in \operatorname{pts}(\operatorname{lcs\_}). \exists \operatorname{lcr\_}. \\ \quad \operatorname{pos\text{-}rr}(\operatorname{lcr\_}) \wedge \operatorname{is\text{-}compact}(\operatorname{subspace\text{-}ms}(\operatorname{lcs\_}, \operatorname{closed\text{-}ball}(\operatorname{lcs\_}, \operatorname{lca\_}, \operatorname{lcr\_}))) \end{array}
\end{equation*}

\par\medskip\noindent\textbf{Definition} (\texttt{is-compact-exhaustion}).\;For all $\operatorname{ces\_}$ and $\operatorname{cek\_}$, 
\begin{equation*}
\begin{array}{@{}l@{}} \operatorname{is\text{-}compact\text{-}exhaustion}(\operatorname{ces\_}, \operatorname{cek\_}) \\ \Leftrightarrow\; \operatorname{is\text{-}metric\text{-}space}(\operatorname{ces\_}) \wedge \\ \operatorname{cek\_} \in (\mathbb{N} \to \operatorname{power}(\operatorname{pts}(\operatorname{ces\_}))) \wedge \\ \forall \operatorname{cen\_} \in \mathbb{N}.\; \operatorname{is\text{-}compact}(\operatorname{subspace\text{-}ms}(\operatorname{ces\_}, \operatorname{cek\_}(\operatorname{cen\_}))) \wedge \\ \forall \operatorname{cen\_} \in \mathbb{N}. \exists \operatorname{cer\_}. \\ \quad (\operatorname{pos\text{-}rr}(\operatorname{cer\_}) \wedge \forall \operatorname{cey\_} \in \operatorname{cek\_}(\operatorname{cen\_}).\; (\operatorname{ball}(\operatorname{ces\_}, \operatorname{cey\_}, \operatorname{cer\_}) \subseteq \operatorname{cek\_}(+(\operatorname{cen\_}, 1)))) \wedge \\ \forall \operatorname{cex\_} \in \operatorname{pts}(\operatorname{ces\_}).\; \exists \operatorname{cen\_} \in \mathbb{N}.\; (\operatorname{cex\_} \in \operatorname{cek\_}(\operatorname{cen\_})) \wedge \\ \forall \operatorname{cea\_}. \\ \quad (\operatorname{cea\_} \subseteq \operatorname{pts}(\operatorname{ces\_})) \Rightarrow  \\ \quad \operatorname{is\text{-}compact}(\operatorname{subspace\text{-}ms}(\operatorname{ces\_}, \operatorname{cea\_})) \Rightarrow  \\ \quad \quad \exists \operatorname{cen\_} \in \mathbb{N}.\; (\operatorname{cea\_} \subseteq \operatorname{cek\_}(\operatorname{cen\_})) \end{array}
\end{equation*}


\vnbbookfile{theorem-library/ca51-compact-laws}{theorem-library/ca51-compact-laws}
\noindent{\small ca51-compact-laws.scm -- compact SUBSETS of a metric space through cluster points (CA-51, 2026-10-26).  Destination: theorem-library/ca51-compact-laws.scm, library analysis.}\par
\vnbprop{lc5-cluster-to-subspace}{\vnbstmt{\begin{array}{@{}l@{}} \forall s, a.; \forall f \in (\mathbb{N} \to a), c \in a. \\ \quad (\operatorname{is\text{-}metric\text{-}space}(s) \wedge ((a \subseteq \operatorname{pts}(s)) \wedge \operatorname{cluster\text{-}point}(s, f, c))) \Rightarrow  \\ \quad \quad \operatorname{cluster\text{-}point}(\operatorname{subspace\text{-}ms}(s, a), f, c) \end{array}}}
\vnbprop{lc5-cluster-from-subspace}{\vnbstmt{\begin{array}{@{}l@{}} \forall s, a, f, c. \\ \quad (\operatorname{is\text{-}metric\text{-}space}(s) \wedge ((a \subseteq \operatorname{pts}(s)) \wedge \operatorname{cluster\text{-}point}(\operatorname{subspace\text{-}ms}(s, a), f, c))) \Rightarrow  \\ \quad \quad (c \in a) \wedge \operatorname{cluster\text{-}point}(s, f, c) \end{array}}}
\vnbprop{lc5-compact-subset-cluster}{\vnbstmt{\begin{array}{@{}l@{}} \forall s, a.; \forall f \in (\mathbb{N} \to a). \\ \quad (\operatorname{is\text{-}metric\text{-}space}(s) \wedge ((a \subseteq \operatorname{pts}(s)) \wedge \operatorname{is\text{-}compact}(\operatorname{subspace\text{-}ms}(s, a)))) \Rightarrow  \\ \quad \quad \exists c \in a.\; \operatorname{cluster\text{-}point}(s, f, c) \end{array}}}
\vnbprop{lc5-cluster-compact-subset}{\vnbstmt{\begin{array}{@{}l@{}} \forall s, a. \\ \quad \operatorname{is\text{-}metric\text{-}space}(s) \wedge \\ \quad a \subseteq \operatorname{pts}(s) \wedge \\ \quad \forall f \in (\mathbb{N} \to a).\; \exists c \in a.\; \operatorname{cluster\text{-}point}(s, f, c) \Rightarrow \\ \quad \quad \operatorname{is\text{-}compact}(\operatorname{subspace\text{-}ms}(s, a)) \end{array}}}
\vnbprop{lc5-frequent-cluster}{\vnbstmt{\begin{array}{@{}l@{}} \forall s, a.; \forall q \in (\mathbb{N} \to \operatorname{pts}(s)). \\ \quad \operatorname{is\text{-}metric\text{-}space}(s) \wedge \\ \quad a \subseteq \operatorname{pts}(s) \wedge \\ \quad \operatorname{is\text{-}compact}(\operatorname{subspace\text{-}ms}(s, a)) \wedge \\ \quad \forall k \in \mathbb{N}.\; \exists j \in \mathbb{N}.\; ((k \leq j) \wedge (q(j) \in a)) \Rightarrow \\ \quad \quad \exists c \in a.\; \operatorname{cluster\text{-}point}(s, q, c) \end{array}}}
\vnbprop{lc5-frequent-pigeonhole}{\vnbstmt{\begin{array}{@{}l@{}} \forall n \in \mathbb{N}, y, q. \\ \quad \forall k \in \mathbb{N}.\; \exists j \in \mathbb{N}.\; ((k \leq j) \wedge \exists i \in \mathbb{N}.\; ((i \leq n) \wedge (q(j) \in y(i)))) \Rightarrow  \\ \quad \quad \exists i \in \mathbb{N}. \\ \quad \quad \quad (i \leq n) \wedge \forall k \in \mathbb{N}.\; \exists j \in \mathbb{N}.\; ((k \leq j) \wedge (q(j) \in y(i))) \end{array}}}
\vnbprop{compact-finite-union}{\vnbstmt{\begin{array}{@{}l@{}} \forall n \in \mathbb{N}, s, y. \\ \quad \operatorname{is\text{-}metric\text{-}space}(s) \wedge \\ \quad \forall i \in \mathbb{N}.\; ((i \leq n) \Rightarrow (y(i) \subseteq \operatorname{pts}(s))) \wedge \\ \quad \forall i \in \mathbb{N}.\; ((i \leq n) \Rightarrow \operatorname{is\text{-}compact}(\operatorname{subspace\text{-}ms}(s, y(i)))) \Rightarrow \\ \quad \quad \operatorname{is\text{-}compact}(\operatorname{subspace\text{-}ms}(s, \{\mathit{ucy} \in \operatorname{pts}(s) : \exists \mathit{uci} \in \mathbb{N}.\; ((\mathit{uci} \leq n) \wedge (\mathit{ucy} \in y(\mathit{uci})))\})) \end{array}}}
\vnbprop{compact-union-two}{\vnbstmt{\begin{array}{@{}l@{}} \forall s, a, b. \\ \quad \operatorname{is\text{-}metric\text{-}space}(s) \wedge \\ \quad a \subseteq \operatorname{pts}(s) \wedge \\ \quad b \subseteq \operatorname{pts}(s) \wedge \\ \quad \operatorname{is\text{-}compact}(\operatorname{subspace\text{-}ms}(s, a)) \wedge \\ \quad \operatorname{is\text{-}compact}(\operatorname{subspace\text{-}ms}(s, b)) \Rightarrow \\ \quad \quad \operatorname{is\text{-}compact}(\operatorname{subspace\text{-}ms}(s, (a \cup b))) \end{array}}}
\vnbprop{closed-in-compact-subset}{\vnbstmt{\begin{array}{@{}l@{}} \forall s, c, a. \\ \quad \operatorname{is\text{-}metric\text{-}space}(s) \wedge \\ \quad c \subseteq \operatorname{pts}(s) \wedge \\ \quad \operatorname{is\text{-}compact}(\operatorname{subspace\text{-}ms}(s, c)) \wedge \\ \quad a \subseteq c \wedge \\ \quad \operatorname{is\text{-}closed}(s, a) \Rightarrow \\ \quad \quad \operatorname{is\text{-}compact}(\operatorname{subspace\text{-}ms}(s, a)) \end{array}}}

\vnbbookfile{theorem-library/product2-compact}{theorem-library/product2-compact}
\noindent{\small ch3c-prod.scm -- the product of two compact metric spaces is compact (PRODUCT2-MS, the max metric). Prefix cx-.  Library: analysis.  Destination: theorem-library/product2-compact.scm, after product2-laws. Route: sequential compactness, a subsequence in the first coordinate, then a further one in the second ...}\par
\vnbprop{product2-compact}{\vnbstmt{\begin{array}{@{}l@{}} \forall s, t. \\ \quad (\operatorname{is\text{-}compact}(s) \wedge \operatorname{is\text{-}compact}(t)) \Rightarrow \operatorname{is\text{-}compact}(\operatorname{product2\text{-}ms}(s, t)) \end{array}}}

\vnbbookfile{theorem-library/categories}{theorem-library/categories}
\noindent{\small categories.scm -- FOUR CATEGORIES BESIDE THE DEFAULTS: Lipschitz, uniformly continuous and continuous maps of metric spaces, bounded linear maps of normed vector spaces; their laws, their Hom functors, and the inclusions.}\par
\par\medskip\noindent\textbf{Definition} (\texttt{hom-lipschitz}).\;$\operatorname{hom\text{-}lipschitz}(a, b) :=$
\begin{equation*}
\{\operatorname{homf\_} \in (\operatorname{pts}(a) \to \operatorname{pts}(b)) : \operatorname{is\text{-}lipschitz\text{-}arrow}(a, b, \operatorname{homf\_})\}
\end{equation*}

\par\medskip\noindent\textbf{Definition} (\texttt{hom-uniformly-continuous}).\;$\operatorname{hom\text{-}uniformly\text{-}continuous}(a, b) :=$
\begin{equation*}
\{\operatorname{homf\_} \in (\operatorname{pts}(a) \to \operatorname{pts}(b)) : \operatorname{is\text{-}uniformly\text{-}continuous\text{-}arrow}(a, b, \operatorname{homf\_})\}
\end{equation*}

\par\medskip\noindent\textbf{Definition} (\texttt{hom-continuous}).\;$\operatorname{hom\text{-}continuous}(a, b) :=$
\begin{equation*}
\{\operatorname{homf\_} \in (\operatorname{pts}(a) \to \operatorname{pts}(b)) : \operatorname{is\text{-}continuous\text{-}arrow}(a, b, \operatorname{homf\_})\}
\end{equation*}

\par\medskip\noindent\textbf{Definition} (\texttt{hom-bounded-linear}).\;$\operatorname{hom\text{-}bounded\text{-}linear}(a, b) :=$
\begin{equation*}
\{\operatorname{homf\_} \in (\operatorname{vec}(a) \to \operatorname{vec}(b)) : \operatorname{is\text{-}bounded\text{-}linear\text{-}arrow}(a, b, \operatorname{homf\_})\}
\end{equation*}

\vnbdef{is-lipschitz-arrow-def}{\vnbstmt{\begin{array}{@{}l@{}} \forall a, b, f. \\ \quad \operatorname{is\text{-}lipschitz\text{-}arrow}(a, b, f) \\ \quad \Leftrightarrow\; \operatorname{is\text{-}metric\text{-}space}(a) \wedge \\ \quad \operatorname{is\text{-}metric\text{-}space}(b) \wedge \\ \quad f \in (\operatorname{pts}(a) \to \operatorname{pts}(b)) \wedge \\ \quad \exists k \in \mathbb{R}. \\ \quad \quad ((0 \leq k) \wedge \forall x, y \in \operatorname{pts}(a).\; (\operatorname{dist}(b)(f(x), f(y)) \leq (k \cdot \operatorname{dist}(a)(x, y)))) \end{array}}}
\vnbdef{is-uniformly-continuous-arrow-def}{\vnbstmt{\begin{array}{@{}l@{}} \forall a, b, f. \\ \quad \operatorname{is\text{-}uniformly\text{-}continuous\text{-}arrow}(a, b, f) \\ \quad \Leftrightarrow\; \operatorname{is\text{-}metric\text{-}space}(a) \wedge \\ \quad \operatorname{is\text{-}metric\text{-}space}(b) \wedge \\ \quad f \in (\operatorname{pts}(a) \to \operatorname{pts}(b)) \wedge \\ \quad \operatorname{is\text{-}uniformly\text{-}continuous}(a, b, f) \end{array}}}
\vnbdef{is-continuous-arrow-def}{\vnbstmt{\begin{array}{@{}l@{}} \forall a, b, f. \\ \quad \operatorname{is\text{-}continuous\text{-}arrow}(a, b, f) \\ \quad \Leftrightarrow\; (\operatorname{is\text{-}metric\text{-}space}(a) \wedge (\operatorname{is\text{-}metric\text{-}space}(b) \wedge ((f \in (\operatorname{pts}(a) \to \operatorname{pts}(b))) \wedge \operatorname{is\text{-}continuous}(a, b, f)))) \end{array}}}
\vnbdef{is-bounded-linear-arrow-def}{\vnbstmt{\begin{array}{@{}l@{}} \forall a, b, f. \\ \quad \operatorname{is\text{-}bounded\text{-}linear\text{-}arrow}(a, b, f) \\ \quad \Leftrightarrow\; \operatorname{is\text{-}normed\text{-}vector\text{-}space}(a) \wedge \\ \quad \operatorname{is\text{-}normed\text{-}vector\text{-}space}(b) \wedge \\ \quad f \in (\operatorname{vec}(a) \to \operatorname{vec}(b)) \wedge \\ \quad \operatorname{scal}(a) = \operatorname{scal}(b) \wedge \\ \quad \forall x_{1}, x_{2} \in \operatorname{vec}(a). \\ \quad \quad f(\operatorname{vadd}(a)(x_{1}, x_{2})) = \operatorname{vadd}(b)(f(x_{1}), f(x_{2})) \wedge \\ \quad \forall x_{1} \in \operatorname{carr}(\operatorname{scal}(a)), x_{2} \in \operatorname{vec}(a). \\ \quad \quad f(\operatorname{act}(a)(x_{1}, x_{2})) = \operatorname{act}(b)(x_{1}, f(x_{2})) \wedge \\ \quad \exists k \in \mathbb{R}. \\ \quad \quad (0 \leq k) \wedge \forall x_{1} \in \operatorname{vec}(a).\; (\operatorname{vnrm}(b)(f(x_{1})) \leq (k \cdot \operatorname{vnrm}(a)(x_{1}))) \end{array}}}
\vnbprop{hom-lipschitz-member-iff}{\vnbstmt{\begin{array}{@{}l@{}} \forall a, b, f. \\ \quad (f \in \operatorname{hom\text{-}lipschitz}(a, b)) \Leftrightarrow ((f \in (\operatorname{pts}(a) \to \operatorname{pts}(b))) \wedge \operatorname{is\text{-}lipschitz\text{-}arrow}(a, b, f)) \end{array}}}
\vnbprop{hom-lipschitz-in-set}{\vnbstmt{\begin{array}{@{}l@{}} \forall a, b. \\ \quad (\operatorname{is\text{-}metric\text{-}space}(a) \wedge \operatorname{is\text{-}metric\text{-}space}(b)) \Rightarrow (\operatorname{hom\text{-}lipschitz}(a, b) \in \mathbf{Set}) \end{array}}}
\vnbprop{hom-uniformly-continuous-member-iff}{\vnbstmt{\begin{array}{@{}l@{}} \forall a, b, f. \\ \quad f \in \operatorname{hom\text{-}uniformly\text{-}continuous}(a, b) \\ \quad \Leftrightarrow\; (f \in (\operatorname{pts}(a) \to \operatorname{pts}(b))) \wedge \operatorname{is\text{-}uniformly\text{-}continuous\text{-}arrow}(a, b, f) \end{array}}}
\vnbprop{hom-uniformly-continuous-in-set}{\vnbstmt{\begin{array}{@{}l@{}} \forall a, b. \\ \quad (\operatorname{is\text{-}metric\text{-}space}(a) \wedge \operatorname{is\text{-}metric\text{-}space}(b)) \Rightarrow  \\ \quad \quad \operatorname{hom\text{-}uniformly\text{-}continuous}(a, b) \in \mathbf{Set} \end{array}}}
\vnbprop{hom-continuous-member-iff}{\vnbstmt{\begin{array}{@{}l@{}} \forall a, b, f. \\ \quad (f \in \operatorname{hom\text{-}continuous}(a, b)) \Leftrightarrow ((f \in (\operatorname{pts}(a) \to \operatorname{pts}(b))) \wedge \operatorname{is\text{-}continuous\text{-}arrow}(a, b, f)) \end{array}}}
\vnbprop{hom-continuous-in-set}{\vnbstmt{\begin{array}{@{}l@{}} \forall a, b. \\ \quad (\operatorname{is\text{-}metric\text{-}space}(a) \wedge \operatorname{is\text{-}metric\text{-}space}(b)) \Rightarrow (\operatorname{hom\text{-}continuous}(a, b) \in \mathbf{Set}) \end{array}}}
\vnbprop{hom-bounded-linear-member-iff}{\vnbstmt{\begin{array}{@{}l@{}} \forall a, b, f. \\ \quad (f \in \operatorname{hom\text{-}bounded\text{-}linear}(a, b)) \Leftrightarrow ((f \in (\operatorname{vec}(a) \to \operatorname{vec}(b))) \wedge \operatorname{is\text{-}bounded\text{-}linear\text{-}arrow}(a, b, f)) \end{array}}}
\vnbprop{hom-bounded-linear-in-set}{\vnbstmt{\begin{array}{@{}l@{}} \forall a, b. \\ \quad (\operatorname{is\text{-}normed\text{-}vector\text{-}space}(a) \wedge \operatorname{is\text{-}normed\text{-}vector\text{-}space}(b)) \Rightarrow  \\ \quad \quad \operatorname{hom\text{-}bounded\text{-}linear}(a, b) \in \mathbf{Set} \end{array}}}
\vnbprop{hom-metric-space-is-lipschitz-arrow}{\vnbstmt{\begin{array}{@{}l@{}} \forall a, b, f. \\ \quad \operatorname{is\text{-}hom\text{-}metric\text{-}space}(a, b, f) \Rightarrow \operatorname{is\text{-}lipschitz\text{-}arrow}(a, b, f) \end{array}}}
\vnbprop{lipschitz-arrow-is-uniformly-continuous-arrow}{\vnbstmt{\begin{array}{@{}l@{}} \forall a, b, f. \\ \quad \operatorname{is\text{-}lipschitz\text{-}arrow}(a, b, f) \Rightarrow \operatorname{is\text{-}uniformly\text{-}continuous\text{-}arrow}(a, b, f) \end{array}}}
\vnbprop{uniformly-continuous-arrow-is-continuous-arrow}{\vnbstmt{\begin{array}{@{}l@{}} \forall a, b, f. \\ \quad \operatorname{is\text{-}uniformly\text{-}continuous\text{-}arrow}(a, b, f) \Rightarrow \operatorname{is\text{-}continuous\text{-}arrow}(a, b, f) \end{array}}}
\vnbprop{hom-normed-vector-space-is-bounded-linear-arrow}{\vnbstmt{\begin{array}{@{}l@{}} \forall a, b, f. \\ \quad \operatorname{is\text{-}hom\text{-}normed\text{-}vector\text{-}space}(a, b, f) \Rightarrow \operatorname{is\text{-}bounded\text{-}linear\text{-}arrow}(a, b, f) \end{array}}}
\vnbprop{hom-lipschitz-id}{\vnbstmt{\forall a.\; (\operatorname{is\text{-}metric\text{-}space}(a) \Rightarrow \operatorname{is\text{-}lipschitz\text{-}arrow}(a, a, \operatorname{id\text{-}fun}(\operatorname{pts}(a))))}}
\vnbprop{hom-uniformly-continuous-id}{\vnbstmt{\begin{array}{@{}l@{}} \forall a. \\ \quad \operatorname{is\text{-}metric\text{-}space}(a) \Rightarrow \operatorname{is\text{-}uniformly\text{-}continuous\text{-}arrow}(a, a, \operatorname{id\text{-}fun}(\operatorname{pts}(a))) \end{array}}}
\vnbprop{hom-continuous-id}{\vnbstmt{\begin{array}{@{}l@{}} \forall a. \\ \quad \operatorname{is\text{-}metric\text{-}space}(a) \Rightarrow \operatorname{is\text{-}continuous\text{-}arrow}(a, a, \operatorname{id\text{-}fun}(\operatorname{pts}(a))) \end{array}}}
\vnbprop{hom-bounded-linear-id}{\vnbstmt{\begin{array}{@{}l@{}} \forall a. \\ \quad \operatorname{is\text{-}normed\text{-}vector\text{-}space}(a) \Rightarrow \operatorname{is\text{-}bounded\text{-}linear\text{-}arrow}(a, a, \operatorname{id\text{-}fun}(\operatorname{vec}(a))) \end{array}}}
\vnbprop{hom-lipschitz-compose}{\vnbstmt{\begin{array}{@{}l@{}} \forall a, b, c, f, g. \\ \quad (\operatorname{is\text{-}lipschitz\text{-}arrow}(a, b, f) \wedge \operatorname{is\text{-}lipschitz\text{-}arrow}(b, c, g)) \Rightarrow  \\ \quad \quad \operatorname{is\text{-}lipschitz\text{-}arrow}(a, c, \operatorname{compose}(g, f)) \end{array}}}
\vnbprop{hom-uniformly-continuous-compose}{\vnbstmt{\begin{array}{@{}l@{}} \forall a, b, c, f, g. \\ \quad (\operatorname{is\text{-}uniformly\text{-}continuous\text{-}arrow}(a, b, f) \wedge \operatorname{is\text{-}uniformly\text{-}continuous\text{-}arrow}(b, c, g)) \Rightarrow  \\ \quad \quad \operatorname{is\text{-}uniformly\text{-}continuous\text{-}arrow}(a, c, \operatorname{compose}(g, f)) \end{array}}}
\vnbprop{hom-continuous-compose}{\vnbstmt{\begin{array}{@{}l@{}} \forall a, b, c, f, g. \\ \quad (\operatorname{is\text{-}continuous\text{-}arrow}(a, b, f) \wedge \operatorname{is\text{-}continuous\text{-}arrow}(b, c, g)) \Rightarrow  \\ \quad \quad \operatorname{is\text{-}continuous\text{-}arrow}(a, c, \operatorname{compose}(g, f)) \end{array}}}
\vnbprop{hom-bounded-linear-compose}{\vnbstmt{\begin{array}{@{}l@{}} \forall a, b, c, f, g. \\ \quad (\operatorname{is\text{-}bounded\text{-}linear\text{-}arrow}(a, b, f) \wedge \operatorname{is\text{-}bounded\text{-}linear\text{-}arrow}(b, c, g)) \Rightarrow  \\ \quad \quad \operatorname{is\text{-}bounded\text{-}linear\text{-}arrow}(a, c, \operatorname{compose}(g, f)) \end{array}}}
\vnbprop{hom-lipschitz-post-type}{\vnbstmt{\begin{array}{@{}l@{}} \forall a, b, c, h. \\ \quad (\operatorname{is\text{-}metric\text{-}space}(a) \wedge (h \in \operatorname{hom\text{-}lipschitz}(b, c))) \Rightarrow  \\ \quad \quad (\lambda f \in \operatorname{hom\text{-}lipschitz}(a, b).\; \operatorname{compose}(h, f)) \in (\operatorname{hom\text{-}lipschitz}(a, b) \to \operatorname{hom\text{-}lipschitz}(a, c)) \end{array}}}
\vnbprop{hom-lipschitz-pre-type}{\vnbstmt{\begin{array}{@{}l@{}} \forall a, b, c, k. \\ \quad (\operatorname{is\text{-}metric\text{-}space}(b) \wedge (k \in \operatorname{hom\text{-}lipschitz}(c, a))) \Rightarrow  \\ \quad \quad (\lambda f \in \operatorname{hom\text{-}lipschitz}(a, b).\; \operatorname{compose}(f, k)) \in (\operatorname{hom\text{-}lipschitz}(a, b) \to \operatorname{hom\text{-}lipschitz}(c, b)) \end{array}}}
\vnbprop{hom-uniformly-continuous-post-type}{\vnbstmt{\begin{array}{@{}l@{}} \forall a, b, c, h. \\ \quad (\operatorname{is\text{-}metric\text{-}space}(a) \wedge (h \in \operatorname{hom\text{-}uniformly\text{-}continuous}(b, c))) \Rightarrow  \\ \quad \quad (\lambda f \in \operatorname{hom\text{-}uniformly\text{-}continuous}(a, b).\; \operatorname{compose}(h, f)) \\ \quad \quad \in\; (\operatorname{hom\text{-}uniformly\text{-}continuous}(a, b) \to \operatorname{hom\text{-}uniformly\text{-}continuous}(a, c)) \end{array}}}
\vnbprop{hom-uniformly-continuous-pre-type}{\vnbstmt{\begin{array}{@{}l@{}} \forall a, b, c, k. \\ \quad (\operatorname{is\text{-}metric\text{-}space}(b) \wedge (k \in \operatorname{hom\text{-}uniformly\text{-}continuous}(c, a))) \Rightarrow  \\ \quad \quad (\lambda f \in \operatorname{hom\text{-}uniformly\text{-}continuous}(a, b).\; \operatorname{compose}(f, k)) \\ \quad \quad \in\; (\operatorname{hom\text{-}uniformly\text{-}continuous}(a, b) \to \operatorname{hom\text{-}uniformly\text{-}continuous}(c, b)) \end{array}}}
\vnbprop{hom-continuous-post-type}{\vnbstmt{\begin{array}{@{}l@{}} \forall a, b, c, h. \\ \quad (\operatorname{is\text{-}metric\text{-}space}(a) \wedge (h \in \operatorname{hom\text{-}continuous}(b, c))) \Rightarrow  \\ \quad \quad (\lambda f \in \operatorname{hom\text{-}continuous}(a, b).\; \operatorname{compose}(h, f)) \\ \quad \quad \in\; (\operatorname{hom\text{-}continuous}(a, b) \to \operatorname{hom\text{-}continuous}(a, c)) \end{array}}}
\vnbprop{hom-continuous-pre-type}{\vnbstmt{\begin{array}{@{}l@{}} \forall a, b, c, k. \\ \quad (\operatorname{is\text{-}metric\text{-}space}(b) \wedge (k \in \operatorname{hom\text{-}continuous}(c, a))) \Rightarrow  \\ \quad \quad (\lambda f \in \operatorname{hom\text{-}continuous}(a, b).\; \operatorname{compose}(f, k)) \\ \quad \quad \in\; (\operatorname{hom\text{-}continuous}(a, b) \to \operatorname{hom\text{-}continuous}(c, b)) \end{array}}}
\vnbprop{hom-bounded-linear-post-type}{\vnbstmt{\begin{array}{@{}l@{}} \forall a, b, c, h. \\ \quad (\operatorname{is\text{-}normed\text{-}vector\text{-}space}(a) \wedge (h \in \operatorname{hom\text{-}bounded\text{-}linear}(b, c))) \Rightarrow  \\ \quad \quad (\lambda f \in \operatorname{hom\text{-}bounded\text{-}linear}(a, b).\; \operatorname{compose}(h, f)) \\ \quad \quad \in\; (\operatorname{hom\text{-}bounded\text{-}linear}(a, b) \to \operatorname{hom\text{-}bounded\text{-}linear}(a, c)) \end{array}}}
\vnbprop{hom-bounded-linear-pre-type}{\vnbstmt{\begin{array}{@{}l@{}} \forall a, b, c, k. \\ \quad (\operatorname{is\text{-}normed\text{-}vector\text{-}space}(b) \wedge (k \in \operatorname{hom\text{-}bounded\text{-}linear}(c, a))) \Rightarrow  \\ \quad \quad (\lambda f \in \operatorname{hom\text{-}bounded\text{-}linear}(a, b).\; \operatorname{compose}(f, k)) \\ \quad \quad \in\; (\operatorname{hom\text{-}bounded\text{-}linear}(a, b) \to \operatorname{hom\text{-}bounded\text{-}linear}(c, b)) \end{array}}}
\vnbprop{hom-metric-space-subset-hom-lipschitz}{\vnbstmt{\begin{array}{@{}l@{}} \forall a, b. \\ \quad (\operatorname{is\text{-}metric\text{-}space}(a) \wedge \operatorname{is\text{-}metric\text{-}space}(b)) \Rightarrow  \\ \quad \quad \operatorname{hom\text{-}metric\text{-}space}(a, b) \subseteq \operatorname{hom\text{-}lipschitz}(a, b) \end{array}}}
\vnbprop{hom-lipschitz-subset-hom-uniformly-continuous}{\vnbstmt{\begin{array}{@{}l@{}} \forall a, b. \\ \quad (\operatorname{is\text{-}metric\text{-}space}(a) \wedge \operatorname{is\text{-}metric\text{-}space}(b)) \Rightarrow  \\ \quad \quad \operatorname{hom\text{-}lipschitz}(a, b) \subseteq \operatorname{hom\text{-}uniformly\text{-}continuous}(a, b) \end{array}}}
\vnbprop{hom-uniformly-continuous-subset-hom-continuous}{\vnbstmt{\begin{array}{@{}l@{}} \forall a, b. \\ \quad (\operatorname{is\text{-}metric\text{-}space}(a) \wedge \operatorname{is\text{-}metric\text{-}space}(b)) \Rightarrow  \\ \quad \quad \operatorname{hom\text{-}uniformly\text{-}continuous}(a, b) \subseteq \operatorname{hom\text{-}continuous}(a, b) \end{array}}}
\vnbprop{hom-normed-vector-space-subset-hom-bounded-linear}{\vnbstmt{\begin{array}{@{}l@{}} \forall a, b. \\ \quad (\operatorname{is\text{-}normed\text{-}vector\text{-}space}(a) \wedge \operatorname{is\text{-}normed\text{-}vector\text{-}space}(b)) \Rightarrow  \\ \quad \quad \operatorname{hom\text{-}normed\text{-}vector\text{-}space}(a, b) \subseteq \operatorname{hom\text{-}bounded\text{-}linear}(a, b) \end{array}}}

\vnbbookfile{theorem-library/c-continuous-category}{theorem-library/c-continuous-category}
\noindent{\small c-continuous-category.scm -- THE CATEGORY OF COUNTABLY-METRISED SPACES AND CONTINUOUS MAPS: objects the C-METRIC-SPACEs, arrows the maps f : PTS(a) -$>$ PTS(b) continuous for the canonical metrics C-METRIC(a), C-METRIC(b).}\par
\par\medskip\noindent\textbf{Definition} (\texttt{hom-c-continuous}).\;$\operatorname{hom\text{-}c\text{-}continuous}(a, b) :=$
\begin{equation*}
\{\operatorname{homf\_} \in (\operatorname{pts}(a) \to \operatorname{pts}(b)) : \operatorname{is\text{-}c\text{-}continuous\text{-}arrow}(a, b, \operatorname{homf\_})\}
\end{equation*}

\vnbdef{is-c-continuous-arrow-def}{\vnbstmt{\begin{array}{@{}l@{}} \forall a, b, f. \\ \quad \operatorname{is\text{-}c\text{-}continuous\text{-}arrow}(a, b, f) \\ \quad \Leftrightarrow\; \operatorname{is\text{-}c\text{-}metric\text{-}space}(a) \wedge \\ \quad \operatorname{is\text{-}c\text{-}metric\text{-}space}(b) \wedge \\ \quad f \in (\operatorname{pts}(a) \to \operatorname{pts}(b)) \wedge \\ \quad \operatorname{is\text{-}continuous}(\operatorname{c\text{-}metric}(a), \operatorname{c\text{-}metric}(b), f) \end{array}}}
\vnbprop{c-metric-canonical-carrier}{\vnbstmt{\forall s.\; (\operatorname{is\text{-}c\text{-}metric\text{-}space}(s) \Rightarrow (\operatorname{pts}(\operatorname{c\text{-}metric}(s)) \simeq \operatorname{pts}(s)))}}
\vnbprop{hom-c-continuous-member-iff}{\vnbstmt{\begin{array}{@{}l@{}} \forall a, b, f. \\ \quad (f \in \operatorname{hom\text{-}c\text{-}continuous}(a, b)) \Leftrightarrow ((f \in (\operatorname{pts}(a) \to \operatorname{pts}(b))) \wedge \operatorname{is\text{-}c\text{-}continuous\text{-}arrow}(a, b, f)) \end{array}}}
\vnbprop{hom-c-continuous-in-set}{\vnbstmt{\begin{array}{@{}l@{}} \forall a, b. \\ \quad (\operatorname{is\text{-}c\text{-}metric\text{-}space}(a) \wedge \operatorname{is\text{-}c\text{-}metric\text{-}space}(b)) \Rightarrow  \\ \quad \quad \operatorname{hom\text{-}c\text{-}continuous}(a, b) \in \mathbf{Set} \end{array}}}
\vnbprop{hom-c-continuous-compose}{\vnbstmt{\begin{array}{@{}l@{}} \forall a, b, c, f, g. \\ \quad (\operatorname{is\text{-}c\text{-}continuous\text{-}arrow}(a, b, f) \wedge \operatorname{is\text{-}c\text{-}continuous\text{-}arrow}(b, c, g)) \Rightarrow  \\ \quad \quad \operatorname{is\text{-}c\text{-}continuous\text{-}arrow}(a, c, \operatorname{compose}(g, f)) \end{array}}}
\vnbprop{hom-c-continuous-id-of-metric}{\vnbstmt{\begin{array}{@{}l@{}} \forall a. \\ \quad (\operatorname{is\text{-}c\text{-}metric\text{-}space}(a) \wedge \operatorname{is\text{-}metric\text{-}space}(\operatorname{c\text{-}metric}(a))) \Rightarrow  \\ \quad \quad \operatorname{is\text{-}c\text{-}continuous\text{-}arrow}(a, a, \operatorname{id\text{-}fun}(\operatorname{pts}(a))) \end{array}}}
\vnbprop{hom-c-continuous-post-type}{\vnbstmt{\begin{array}{@{}l@{}} \forall a, b, c, h. \\ \quad (\operatorname{is\text{-}c\text{-}metric\text{-}space}(a) \wedge (h \in \operatorname{hom\text{-}c\text{-}continuous}(b, c))) \Rightarrow  \\ \quad \quad (\lambda f \in \operatorname{hom\text{-}c\text{-}continuous}(a, b).\; \operatorname{compose}(h, f)) \\ \quad \quad \in\; (\operatorname{hom\text{-}c\text{-}continuous}(a, b) \to \operatorname{hom\text{-}c\text{-}continuous}(a, c)) \end{array}}}
\vnbprop{hom-c-continuous-pre-type}{\vnbstmt{\begin{array}{@{}l@{}} \forall a, b, c, k. \\ \quad (\operatorname{is\text{-}c\text{-}metric\text{-}space}(b) \wedge (k \in \operatorname{hom\text{-}c\text{-}continuous}(c, a))) \Rightarrow  \\ \quad \quad (\lambda f \in \operatorname{hom\text{-}c\text{-}continuous}(a, b).\; \operatorname{compose}(f, k)) \\ \quad \quad \in\; (\operatorname{hom\text{-}c\text{-}continuous}(a, b) \to \operatorname{hom\text{-}c\text{-}continuous}(c, b)) \end{array}}}
\vnbprop{hom-c-continuous-id}{\vnbstmt{\begin{array}{@{}l@{}} \forall a. \\ \quad \operatorname{is\text{-}c\text{-}metric\text{-}space}(a) \Rightarrow \operatorname{is\text{-}c\text{-}continuous\text{-}arrow}(a, a, \operatorname{id\text{-}fun}(\operatorname{pts}(a))) \end{array}}}

\vnbbookfile{theorem-library/diff-on-laws-2}{theorem-library/diff-on-laws-2}
\noindent{\small diff-on-laws-2.scm -- THE DERIVATIVE LAWS ON AN OPEN SET, over a normed field: uniqueness, the constant and the identity, the sum, the product and the chain rule.  Continues theorem-library/diff-on-laws.scm (batch 12-G), which proved the read-offs, the bridge `diff-at-iff-diff-on' and}\par
\vnbprop{ms-const-continuous-on}{\vnbstmt{\begin{array}{@{}l@{}} \forall s, d, t.; \forall c \in \operatorname{pts}(t), a \in \operatorname{pts}(s). \\ \quad ((\operatorname{pts}(s) \simeq d) \wedge (\operatorname{is\text{-}metric\text{-}space}(s) \wedge \operatorname{is\text{-}metric\text{-}space}(t))) \Rightarrow  \\ \quad \quad \operatorname{is\text{-}continuous\text{-}at}(s, t, (\lambda z \in d.\; c), a) \end{array}}}
\vnbprop{ms-identity-continuous-on}{\vnbstmt{\begin{array}{@{}l@{}} \forall s, d.; \forall a \in \operatorname{pts}(s). \\ \quad ((\operatorname{pts}(s) \simeq d) \wedge \operatorname{is\text{-}metric\text{-}space}(s)) \Rightarrow  \\ \quad \quad \operatorname{is\text{-}continuous\text{-}at}(s, s, (\lambda z \in d.\; z), a) \end{array}}}
\vnbprop{nf-sum-continuous-on}{\vnbstmt{\begin{array}{@{}l@{}} \forall k, s, d, g, h, a. \\ \quad \operatorname{is\text{-}normed\text{-}field}(k) \wedge \\ \quad (\operatorname{pts}(s) \simeq d) \wedge \\ \quad \operatorname{is\text{-}continuous\text{-}at}(s, \operatorname{nf\text{-}metric\text{-}space}(k), g, a) \wedge \\ \quad \operatorname{is\text{-}continuous\text{-}at}(s, \operatorname{nf\text{-}metric\text{-}space}(k), h, a) \Rightarrow \\ \quad \quad \operatorname{is\text{-}continuous\text{-}at}(s, \operatorname{nf\text{-}metric\text{-}space}(k), (\lambda z \in d.\; \operatorname{add}(k)(g(z), h(z))), a) \end{array}}}
\vnbprop{nf-product-continuous-on}{\vnbstmt{\begin{array}{@{}l@{}} \forall k, s, d, g, h, a. \\ \quad \operatorname{is\text{-}normed\text{-}field}(k) \wedge \\ \quad (\operatorname{pts}(s) \simeq d) \wedge \\ \quad \operatorname{is\text{-}continuous\text{-}at}(s, \operatorname{nf\text{-}metric\text{-}space}(k), g, a) \wedge \\ \quad \operatorname{is\text{-}continuous\text{-}at}(s, \operatorname{nf\text{-}metric\text{-}space}(k), h, a) \Rightarrow \\ \quad \quad \operatorname{is\text{-}continuous\text{-}at}(s, \operatorname{nf\text{-}metric\text{-}space}(k), (\lambda z \in d.\; \operatorname{mul}(k)(g(z), h(z))), a) \end{array}}}
\vnbprop{ms-compose-continuous-on}{\vnbstmt{\begin{array}{@{}l@{}} \forall s, d, t, u, f, g, a. \\ \quad ((\operatorname{pts}(s) \simeq d) \wedge (\operatorname{is\text{-}continuous\text{-}at}(s, t, f, a) \wedge \operatorname{is\text{-}continuous\text{-}at}(t, u, g, f(a)))) \Rightarrow  \\ \quad \quad \operatorname{is\text{-}continuous\text{-}at}(s, u, (\lambda z \in d.\; g(f(z))), a) \end{array}}}
\vnbprop{ms-cont-transfer-on}{\vnbstmt{\begin{array}{@{}l@{}} \forall s, t, d, c, g, f, a. \\ \quad (\operatorname{pts}(s) \simeq d) \wedge \\ \quad (\operatorname{pts}(t) \simeq c) \wedge \\ \quad \operatorname{is\text{-}continuous\text{-}at}(s, t, g, a) \wedge \\ \quad f \in (d \to c) \wedge \\ \quad \forall y \in d.\; (f(y) = g(y)) \Rightarrow \\ \quad \quad \operatorname{is\text{-}continuous\text{-}at}(s, t, f, a) \end{array}}}
\vnbprop{nf-lam-const-in-fun}{\vnbstmt{\begin{array}{@{}l@{}} \forall k.; \forall u \in \mathbf{Set}, c \in \operatorname{carr}(k). \\ \quad \operatorname{is\text{-}normed\text{-}field}(k) \Rightarrow ((\lambda x \in u.\; c) \in (u \to \operatorname{carr}(k))) \end{array}}}
\vnbprop{nf-lam-id-in-fun}{\vnbstmt{\begin{array}{@{}l@{}} \forall k.; \forall u \in \mathbf{Set}. \\ \quad (\operatorname{is\text{-}normed\text{-}field}(k) \wedge (u \subseteq \operatorname{carr}(k))) \Rightarrow  \\ \quad \quad (\lambda x \in u.\; x) \in (u \to \operatorname{carr}(k)) \end{array}}}
\vnbprop{nf-lam-sum-in-fun}{\vnbstmt{\begin{array}{@{}l@{}} \forall k.; \forall u \in \mathbf{Set}, f, g \in (u \to \operatorname{carr}(k)). \\ \quad \operatorname{is\text{-}normed\text{-}field}(k) \Rightarrow  \\ \quad \quad (\lambda x \in u.\; \operatorname{add}(k)(f(x), g(x))) \in (u \to \operatorname{carr}(k)) \end{array}}}
\vnbprop{nf-lam-prod-in-fun}{\vnbstmt{\begin{array}{@{}l@{}} \forall k.; \forall u \in \mathbf{Set}, f, g \in (u \to \operatorname{carr}(k)). \\ \quad \operatorname{is\text{-}normed\text{-}field}(k) \Rightarrow  \\ \quad \quad (\lambda x \in u.\; \operatorname{mul}(k)(f(x), g(x))) \in (u \to \operatorname{carr}(k)) \end{array}}}
\vnbprop{nf-norm-zero-iff}{\vnbstmt{\begin{array}{@{}l@{}} \forall k.; \forall u \in \operatorname{carr}(k). \\ \quad \operatorname{is\text{-}normed\text{-}field}(k) \Rightarrow ((\operatorname{fnrm}(k)(u) = 0) \Leftrightarrow (u = \operatorname{zero}(k))) \end{array}}}
\vnbprop{nf-norm-zero-implies}{\vnbstmt{\begin{array}{@{}l@{}} \forall k.; \forall u \in \operatorname{carr}(k). \\ \quad (\operatorname{is\text{-}normed\text{-}field}(k) \wedge (\operatorname{fnrm}(k)(u) = 0)) \Rightarrow (u = \operatorname{zero}(k)) \end{array}}}
\vnbprop{rr-nf-small-elements}{\vnbstmt{\begin{array}{@{}l@{}} \forall e. \\ \quad \operatorname{pos\text{-}rr}(e) \Rightarrow  \\ \quad \quad \exists h \in \operatorname{carr}(\operatorname{rr\text{-}normed\text{-}field}). \\ \quad \quad \quad \neg (h = \operatorname{zero}(\operatorname{rr\text{-}normed\text{-}field})) \wedge (\operatorname{fnrm}(\operatorname{rr\text{-}normed\text{-}field})(h) < e) \end{array}}}
\vnbprop{nf-open-limit-point}{\vnbstmt{\begin{array}{@{}l@{}} \forall k, u.; \forall a \in u, r. \\ \quad \operatorname{is\text{-}normed\text{-}field}(k) \wedge \\ \quad \forall e. \\ \quad \quad \operatorname{pos\text{-}rr}(e) \Rightarrow \exists h \in \operatorname{carr}(k).\; (\neg (h = \operatorname{zero}(k)) \wedge (\operatorname{fnrm}(k)(h) < e)) \wedge \\ \quad \operatorname{is\text{-}open}(\operatorname{nf\text{-}metric\text{-}space}(k), u) \wedge \\ \quad \operatorname{pos\text{-}rr}(r) \Rightarrow \\ \quad \quad \exists b \in u. \\ \quad \quad \quad \neg (b = a) \wedge (\operatorname{dist}(\operatorname{subspace\text{-}ms}(\operatorname{nf\text{-}metric\text{-}space}(k), u))(a, b) \leq r) \end{array}}}
\vnbprop{nf-zero-in-carr}{\vnbstmt{\forall k.\; (\operatorname{is\text{-}normed\text{-}field}(k) \Rightarrow (\operatorname{zero}(k) \in \operatorname{carr}(k)))}}
\vnbprop{nf-norm-of-zero}{\vnbstmt{\forall k.\; (\operatorname{is\text{-}normed\text{-}field}(k) \Rightarrow (\operatorname{fnrm}(k)(\operatorname{zero}(k)) = 0))}}
\vnbprop{nf-sub-zero-eq}{\vnbstmt{\begin{array}{@{}l@{}} \forall k, u, v. \\ \quad (\operatorname{is\text{-}normed\text{-}field}(k) \wedge ((u \in \operatorname{carr}(k)) \wedge ((v \in \operatorname{carr}(k)) \wedge (\operatorname{add}(k)(u, \operatorname{neg}(k)(v)) = \operatorname{zero}(k))))) \Rightarrow  \\ \quad \quad u = v \end{array}}}
\vnbprop{nf-mul-cancel}{\vnbstmt{\begin{array}{@{}l@{}} \forall k, u, v, w. \\ \quad \operatorname{is\text{-}normed\text{-}field}(k) \wedge \\ \quad u \in \operatorname{carr}(k) \wedge \\ \quad v \in \operatorname{carr}(k) \wedge \\ \quad w \in \operatorname{carr}(k) \wedge \\ \quad \neg (w = \operatorname{zero}(k)) \wedge \\ \quad \operatorname{mul}(k)(u, w) = \operatorname{mul}(k)(v, w) \Rightarrow \\ \quad \quad u = v \end{array}}}
\vnbprop{diff-on-unique}{\vnbstmt{\begin{array}{@{}l@{}} \forall k, u, f, a, l, m. \\ \quad \operatorname{is\text{-}normed\text{-}field}(k) \wedge \\ \quad \forall e. \\ \quad \quad \operatorname{pos\text{-}rr}(e) \Rightarrow \exists h \in \operatorname{carr}(k).\; (\neg (h = \operatorname{zero}(k)) \wedge (\operatorname{fnrm}(k)(h) < e)) \wedge \\ \quad \operatorname{is\text{-}diff\text{-}on}(k, u, f, a, l) \wedge \\ \quad \operatorname{is\text{-}diff\text{-}on}(k, u, f, a, m) \Rightarrow \\ \quad \quad l = m \end{array}}}
\vnbprop{nf-one-in-carr}{\vnbstmt{\forall k.\; (\operatorname{is\text{-}normed\text{-}field}(k) \Rightarrow (\operatorname{one}(k) \in \operatorname{carr}(k)))}}
\vnbprop{nf-open-is-set}{\vnbstmt{\begin{array}{@{}l@{}} \forall k, u. \\ \quad (\operatorname{is\text{-}normed\text{-}field}(k) \wedge \operatorname{is\text{-}open}(\operatorname{nf\text{-}metric\text{-}space}(k), u)) \Rightarrow (u \in \mathbf{Set}) \end{array}}}
\vnbprop{diff-on-const}{\vnbstmt{\begin{array}{@{}l@{}} \forall k, u.; \forall c \in \operatorname{carr}(k), a \in u. \\ \quad (\operatorname{is\text{-}normed\text{-}field}(k) \wedge \operatorname{is\text{-}open}(\operatorname{nf\text{-}metric\text{-}space}(k), u)) \Rightarrow  \\ \quad \quad \operatorname{is\text{-}diff\text{-}on}(k, u, (\lambda x \in u.\; c), a, \operatorname{zero}(k)) \end{array}}}
\vnbprop{diff-on-identity}{\vnbstmt{\begin{array}{@{}l@{}} \forall k, u.; \forall a \in u. \\ \quad (\operatorname{is\text{-}normed\text{-}field}(k) \wedge \operatorname{is\text{-}open}(\operatorname{nf\text{-}metric\text{-}space}(k), u)) \Rightarrow  \\ \quad \quad \operatorname{is\text{-}diff\text{-}on}(k, u, (\lambda x \in u.\; x), a, \operatorname{one}(k)) \end{array}}}
\vnbprop{diff-on-sum}{\vnbstmt{\begin{array}{@{}l@{}} \forall k, u, f, g, a, l, m. \\ \quad (\operatorname{is\text{-}normed\text{-}field}(k) \wedge (\operatorname{is\text{-}diff\text{-}on}(k, u, f, a, l) \wedge \operatorname{is\text{-}diff\text{-}on}(k, u, g, a, m))) \Rightarrow  \\ \quad \quad \operatorname{is\text{-}diff\text{-}on}(k, u, (\lambda x \in u.\; \operatorname{add}(k)(f(x), g(x))), a, \operatorname{add}(k)(l, m)) \end{array}}}
\vnbprop{nf-lam-const-z-in-fun}{\vnbstmt{\begin{array}{@{}l@{}} \forall k.; \forall u \in \mathbf{Set}, c \in \operatorname{carr}(k). \\ \quad \operatorname{is\text{-}normed\text{-}field}(k) \Rightarrow ((\lambda z \in u.\; c) \in (u \to \operatorname{carr}(k))) \end{array}}}
\vnbprop{nf-lam-prod-v-in-fun}{\vnbstmt{\begin{array}{@{}l@{}} \forall k.; \forall u \in \mathbf{Set}, f, g \in (u \to \operatorname{carr}(k)). \\ \quad \operatorname{is\text{-}normed\text{-}field}(k) \Rightarrow  \\ \quad \quad (\lambda v \in u.\; \operatorname{mul}(k)(f(v), g(v))) \in (u \to \operatorname{carr}(k)) \end{array}}}
\vnbprop{nf-product-continuous-v}{\vnbstmt{\begin{array}{@{}l@{}} \forall k, s, d, g, h, a. \\ \quad \operatorname{is\text{-}normed\text{-}field}(k) \wedge \\ \quad (\operatorname{pts}(s) \simeq d) \wedge \\ \quad \operatorname{is\text{-}continuous\text{-}at}(s, \operatorname{nf\text{-}metric\text{-}space}(k), g, a) \wedge \\ \quad \operatorname{is\text{-}continuous\text{-}at}(s, \operatorname{nf\text{-}metric\text{-}space}(k), h, a) \Rightarrow \\ \quad \quad \operatorname{is\text{-}continuous\text{-}at}(s, \operatorname{nf\text{-}metric\text{-}space}(k), (\lambda v \in d.\; \operatorname{mul}(k)(g(v), h(v))), a) \end{array}}}
\vnbprop{diff-on-product}{\vnbstmt{\begin{array}{@{}l@{}} \forall k, u, f, g, a, l, m. \\ \quad (\operatorname{is\text{-}normed\text{-}field}(k) \wedge (\operatorname{is\text{-}diff\text{-}on}(k, u, f, a, l) \wedge \operatorname{is\text{-}diff\text{-}on}(k, u, g, a, m))) \Rightarrow  \\ \quad \quad \operatorname{is\text{-}diff\text{-}on}(k, u, (\lambda x \in u.\; \operatorname{mul}(k)(f(x), g(x))), a, \\ \quad \quad \quad \operatorname{add}(k)(\operatorname{mul}(k)(l, g(a)), \operatorname{mul}(k)(m, f(a)))) \end{array}}}
\vnbprop{ms-corestrict-continuous}{\vnbstmt{\begin{array}{@{}l@{}} \forall s, t, \mathit{d2s}, f, a. \\ \quad \operatorname{is\text{-}metric\text{-}space}(t) \wedge \\ \quad \mathit{d2s} \subseteq \operatorname{pts}(t) \wedge \\ \quad \operatorname{is\text{-}continuous\text{-}at}(s, t, f, a) \wedge \\ \quad \forall y \in \operatorname{pts}(s).\; (f(y) \in \mathit{d2s}) \Rightarrow \\ \quad \quad \operatorname{is\text{-}continuous\text{-}at}(s, \operatorname{subspace\text{-}ms}(t, \mathit{d2s}), f, a) \end{array}}}
\vnbprop{nf-lam-comp-z-in-fun}{\vnbstmt{\begin{array}{@{}l@{}} \forall k.; \forall u \in \mathbf{Set}, s, f.; \forall g \in (s \to \operatorname{carr}(k)). \\ \quad (\operatorname{is\text{-}normed\text{-}field}(k) \wedge \forall y \in u.\; (f(y) \in s)) \Rightarrow  \\ \quad \quad (\lambda z \in u.\; g(f(z))) \in (u \to \operatorname{carr}(k)) \end{array}}}
\vnbprop{diff-on-chain}{\vnbstmt{\begin{array}{@{}l@{}} \forall k, u, s, f, g, a, l, m. \\ \quad \operatorname{is\text{-}normed\text{-}field}(k) \wedge \\ \quad \operatorname{is\text{-}diff\text{-}on}(k, u, f, a, l) \wedge \\ \quad \operatorname{is\text{-}diff\text{-}on}(k, s, g, f(a), m) \wedge \\ \quad \forall y \in u.\; (f(y) \in s) \Rightarrow \\ \quad \quad \operatorname{is\text{-}diff\text{-}on}(k, u, (\lambda x \in u.\; g(f(x))), a, \operatorname{mul}(k)(m, l)) \end{array}}}

\vnbbookfile{theorem-library/holomorphic-basics}{theorem-library/holomorphic-basics}
\noindent{\small holomorphic-basics.scm -- the COMPLEX side of IS-DIFF-ON / HOLOMORPHIC-ON.}\par
\vnbprop{cc-nf-add-apply}{\vnbstmt{\begin{array}{@{}l@{}} \forall u, v. \\ \quad ((u \in \mathbb{C}) \wedge (v \in \mathbb{C})) \Rightarrow (\operatorname{add}(\operatorname{cc\text{-}normed\text{-}field})(u, v) \simeq (u + v)) \end{array}}}
\vnbprop{cc-nf-mul-apply}{\vnbstmt{\begin{array}{@{}l@{}} \forall u, v. \\ \quad ((u \in \mathbb{C}) \wedge (v \in \mathbb{C})) \Rightarrow (\operatorname{mul}(\operatorname{cc\text{-}normed\text{-}field})(u, v) \simeq (u \cdot v)) \end{array}}}
\vnbprop{cc-nf-neg-apply}{\vnbstmt{\forall u \in \mathbb{C}.\; (\operatorname{neg}(\operatorname{cc\text{-}normed\text{-}field})(u) \simeq -u)}}
\vnbprop{cc-nf-sub}{\vnbstmt{\begin{array}{@{}l@{}} \forall u, v. \\ \quad ((u \in \mathbb{C}) \wedge (v \in \mathbb{C})) \Rightarrow  \\ \quad \quad (\operatorname{add}(\operatorname{cc\text{-}normed\text{-}field})(u, \operatorname{neg}(\operatorname{cc\text{-}normed\text{-}field})(v)) \simeq (u - v)) \end{array}}}
\vnbprop{nf-cc-is-metric-space}{\vnbstmt{\operatorname{is\text{-}metric\text{-}space}(\operatorname{nf\text{-}metric\text{-}space}(\operatorname{cc\text{-}normed\text{-}field}))}}
\vnbprop{nf-cc-agree-pts}{\vnbstmt{(\operatorname{pts}(\operatorname{nf\text{-}metric\text{-}space}(\operatorname{cc\text{-}normed\text{-}field})) \simeq \operatorname{pts}(\operatorname{cc\text{-}ms}))}}
\vnbprop{nf-cc-agree-dist}{\vnbstmt{\begin{array}{@{}l@{}} \forall u, v \in \operatorname{pts}(\operatorname{cc\text{-}ms}). \\ \quad (\operatorname{dist}(\operatorname{nf\text{-}metric\text{-}space}(\operatorname{cc\text{-}normed\text{-}field}))(u, v) \simeq \operatorname{dist}(\operatorname{cc\text{-}ms})(u, v)) \end{array}}}
\vnbprop{cc-ms-open-iff}{\vnbstmt{\forall w.\; (\operatorname{is\text{-}open}(\operatorname{nf\text{-}metric\text{-}space}(\operatorname{cc\text{-}normed\text{-}field}), w) \Leftrightarrow \operatorname{is\text{-}open}(\operatorname{cc\text{-}ms}, w))}}
\vnbprop{cc-nf-small-elements}{\vnbstmt{\begin{array}{@{}l@{}} \forall s. \\ \quad \operatorname{pos\text{-}rr}(s) \Rightarrow  \\ \quad \quad \exists h \in \operatorname{carr}(\operatorname{cc\text{-}normed\text{-}field}). \\ \quad \quad \quad \neg (h = \operatorname{zero}(\operatorname{cc\text{-}normed\text{-}field})) \wedge (\operatorname{fnrm}(\operatorname{cc\text{-}normed\text{-}field})(h) < s) \end{array}}}
\vnbprop{holomorphic-on-open-cc-ms}{\vnbstmt{\forall u, f.\; (\operatorname{holomorphic\text{-}on}(u, f) \Rightarrow \operatorname{is\text{-}open}(\operatorname{cc\text{-}ms}, u))}}
\vnbprop{holomorphic-on-iff-cc-ms}{\vnbstmt{\begin{array}{@{}l@{}} \forall u, f. \\ \quad \operatorname{holomorphic\text{-}on}(u, f) \\ \quad \Leftrightarrow\; \operatorname{is\text{-}open}(\operatorname{cc\text{-}ms}, u) \wedge \\ \quad f \in (u \to \mathbb{C}) \wedge \\ \quad \forall a \in u.\; \exists l.\; \operatorname{is\text{-}diff\text{-}on}(\operatorname{cc\text{-}normed\text{-}field}, u, f, a, l) \end{array}}}

\vnbbookfile{theorem-library/deriv-on-read-off}{theorem-library/deriv-on-read-off}
\noindent{\small deriv-on-read-off.scm -- DERIV-ON reads off IS-DIFF-ON. Batch 27-A follow-on, 2026-09-24.}\par
\vnbprop{deriv-on-of-is-diff-on}{\vnbstmt{\begin{array}{@{}l@{}} \forall k, u, f, a, l. \\ \quad \operatorname{is\text{-}normed\text{-}field}(k) \wedge \\ \quad \forall e. \\ \quad \quad \operatorname{pos\text{-}rr}(e) \Rightarrow \exists h \in \operatorname{carr}(k).\; (\neg (h = \operatorname{zero}(k)) \wedge (\operatorname{fnrm}(k)(h) < e)) \wedge \\ \quad \operatorname{is\text{-}diff\text{-}on}(k, u, f, a, l) \Rightarrow \\ \quad \quad \operatorname{deriv\text{-}on}(k, u, f, a) = l \end{array}}}
\vnbprop{cc-deriv-on-of-is-diff-on}{\vnbstmt{\begin{array}{@{}l@{}} \forall u, f, a, l. \\ \quad \operatorname{is\text{-}diff\text{-}on}(\operatorname{cc\text{-}normed\text{-}field}, u, f, a, l) \Rightarrow  \\ \quad \quad \operatorname{deriv\text{-}on}(\operatorname{cc\text{-}normed\text{-}field}, u, f, a) = l \end{array}}}

\vnbbookfile{theorem-library/holomorphic-basics-2}{theorem-library/holomorphic-basics-2}
\noindent{\small holomorphic-basics-2.scm -- section 2.1 of the user's complex-analysis notes, finished: the constant and the identity are holomorphic on any open set, sums and products of holomorphic functions are holomorphic, and so is z |-$>$ z\textasciicircum{}n for every natural n.}\par
\vnbprop{diff-on-transfer-ptwise-eq}{\vnbstmt{\begin{array}{@{}l@{}} \forall k, u, f, g, a, l. \\ \quad ((f \in (u \to \operatorname{carr}(k))) \wedge (\operatorname{is\text{-}diff\text{-}on}(k, u, g, a, l) \wedge \forall x \in u.\; (f(x) \simeq g(x)))) \Rightarrow  \\ \quad \quad \operatorname{is\text{-}diff\text{-}on}(k, u, f, a, l) \end{array}}}
\vnbprop{holomorphic-transfer-ptwise-eq}{\vnbstmt{\begin{array}{@{}l@{}} \forall u, f, g. \\ \quad (\operatorname{holomorphic\text{-}on}(u, g) \wedge ((f \in (u \to \mathbb{C})) \wedge \forall x \in u.\; (f(x) \simeq g(x)))) \Rightarrow  \\ \quad \quad \operatorname{holomorphic\text{-}on}(u, f) \end{array}}}
\vnbprop{holomorphic-const}{\vnbstmt{\begin{array}{@{}l@{}} \forall u.; \forall c \in \mathbb{C}. \\ \quad \operatorname{is\text{-}open}(\operatorname{cc\text{-}ms}, u) \Rightarrow \operatorname{holomorphic\text{-}on}(u, (\lambda x \in u.\; c)) \end{array}}}
\vnbprop{holomorphic-identity}{\vnbstmt{\forall u.\; (\operatorname{is\text{-}open}(\operatorname{cc\text{-}ms}, u) \Rightarrow \operatorname{holomorphic\text{-}on}(u, (\lambda x \in u.\; x)))}}
\vnbprop{holomorphic-sum}{\vnbstmt{\begin{array}{@{}l@{}} \forall u, f, g, m. \\ \quad \operatorname{holomorphic\text{-}on}(u, f) \wedge \\ \quad \operatorname{holomorphic\text{-}on}(u, g) \wedge \\ \quad m \in (u \to \mathbb{C}) \wedge \\ \quad \forall x \in u.\; (m(x) \simeq (f(x) + g(x))) \Rightarrow \\ \quad \quad \operatorname{holomorphic\text{-}on}(u, m) \end{array}}}
\vnbprop{holomorphic-product}{\vnbstmt{\begin{array}{@{}l@{}} \forall u, f, g, m. \\ \quad \operatorname{holomorphic\text{-}on}(u, f) \wedge \\ \quad \operatorname{holomorphic\text{-}on}(u, g) \wedge \\ \quad m \in (u \to \mathbb{C}) \wedge \\ \quad \forall x \in u.\; (m(x) \simeq (f(x) \cdot g(x))) \Rightarrow \\ \quad \quad \operatorname{holomorphic\text{-}on}(u, m) \end{array}}}
\vnbprop{holomorphic-power}{\vnbstmt{\begin{array}{@{}l@{}} \forall n \in \mathbb{N}, u.; \forall w \in (u \to \mathbb{C}). \\ \quad (\operatorname{is\text{-}open}(\operatorname{cc\text{-}ms}, u) \wedge \forall z \in u.\; (w(z) \simeq \operatorname{power}(z, n))) \Rightarrow  \\ \quad \quad \operatorname{holomorphic\text{-}on}(u, w) \end{array}}}

\vnbbookfile{theorem-library/cc-reciprocal}{theorem-library/cc-reciprocal}
\noindent{\small cc-reciprocal.scm -- z |-$>$ 1/z: the punctured plane, the continuity of the reciprocal, its derivative, and the reciprocal and quotient rules for HOLOMORPHIC-ON.}\par
\vnbprop{cc-nonzero-membership}{\vnbstmt{\forall w.\; ((w \in \operatorname{difference}(\mathbb{C}, \operatorname{singleton}(0))) \Leftrightarrow ((w \in \mathbb{C}) \wedge \neg (w = 0)))}}
\vnbprop{cc-nonzero-subset}{\vnbstmt{\operatorname{difference}(\mathbb{C}, \operatorname{singleton}(0)) \subseteq \mathbb{C}}}
\vnbprop{cc-nonzero-is-set}{\vnbstmt{\operatorname{difference}(\mathbb{C}, \operatorname{singleton}(0)) \in \mathbf{Set}}}
\vnbprop{cc-nonzero-open}{\vnbstmt{\operatorname{is\text{-}open}(\operatorname{cc\text{-}ms}, \operatorname{difference}(\mathbb{C}, \operatorname{singleton}(0)))}}
\vnbprop{cc-nonzero-open-nf}{\vnbstmt{\operatorname{is\text{-}open}(\operatorname{nf\text{-}metric\text{-}space}(\operatorname{cc\text{-}normed\text{-}field}), \operatorname{difference}(\mathbb{C}, \operatorname{singleton}(0)))}}
\vnbprop{cc-recip-lam-in-fun}{\vnbstmt{(\lambda z \in \operatorname{difference}(\mathbb{C}, \operatorname{singleton}(0)).\; \operatorname{recip}(z)) \in (\operatorname{difference}(\mathbb{C}, \operatorname{singleton}(0)) \to \mathbb{C})}}
\vnbprop{cc-recip-continuous-at}{\vnbstmt{\begin{array}{@{}l@{}} \forall a \in \operatorname{difference}(\mathbb{C}, \operatorname{singleton}(0)). \\ \quad \operatorname{is\text{-}continuous\text{-}at}(\operatorname{subspace\text{-}ms}(\operatorname{nf\text{-}metric\text{-}space}(\operatorname{cc\text{-}normed\text{-}field}), \operatorname{difference}(\mathbb{C}, \operatorname{singleton}(0))), \\ \quad \quad \operatorname{nf\text{-}metric\text{-}space}(\operatorname{cc\text{-}normed\text{-}field}), \\ \quad \quad (\lambda z \in \operatorname{difference}(\mathbb{C}, \operatorname{singleton}(0)).\; \operatorname{recip}(z)), \\ \quad \quad a) \end{array}}}
\vnbprop{cc-mul-nonzero}{\vnbstmt{\forall u, v \in \mathbb{C}.\; ((\neg (u = 0) \wedge \neg (v = 0)) \Rightarrow \neg ((u \cdot v) = 0))}}
\vnbprop{cc-recip-mul}{\vnbstmt{\begin{array}{@{}l@{}} \forall u, v \in \mathbb{C}. \\ \quad (\neg (u = 0) \wedge \neg (v = 0)) \Rightarrow (\operatorname{recip}((u \cdot v)) = (\operatorname{recip}(u) \cdot \operatorname{recip}(v))) \end{array}}}
\vnbprop{cc-recip-diff-on}{\vnbstmt{\begin{array}{@{}l@{}} \forall a \in \operatorname{difference}(\mathbb{C}, \operatorname{singleton}(0)). \\ \quad \operatorname{is\text{-}diff\text{-}on}(\operatorname{cc\text{-}normed\text{-}field}, \operatorname{difference}(\mathbb{C}, \operatorname{singleton}(0)), \\ \quad \quad (\lambda z \in \operatorname{difference}(\mathbb{C}, \operatorname{singleton}(0)).\; \operatorname{recip}(z)), \\ \quad \quad a, -\operatorname{recip}((a \cdot a))) \end{array}}}
\vnbprop{holomorphic-recip}{\vnbstmt{\operatorname{holomorphic\text{-}on}(\operatorname{difference}(\mathbb{C}, \operatorname{singleton}(0)), (\lambda z \in \operatorname{difference}(\mathbb{C}, \operatorname{singleton}(0)).\; \operatorname{recip}(z)))}}
\vnbprop{diff-on-recip-cc}{\vnbstmt{\begin{array}{@{}l@{}} \forall u, f, a, l, q. \\ \quad \operatorname{is\text{-}diff\text{-}on}(\operatorname{cc\text{-}normed\text{-}field}, u, f, a, l) \wedge \\ \quad \forall y \in u.\; \neg (f(y) = 0) \wedge \\ \quad q \in (u \to \mathbb{C}) \wedge \\ \quad \forall y \in u.\; (q(y) \simeq \operatorname{recip}(f(y))) \Rightarrow \\ \quad \quad \operatorname{is\text{-}diff\text{-}on}(\operatorname{cc\text{-}normed\text{-}field}, u, q, a, -(l \cdot \operatorname{recip}((f(a) \cdot f(a))))) \end{array}}}
\vnbprop{holomorphic-recip-of}{\vnbstmt{\begin{array}{@{}l@{}} \forall u, f, q. \\ \quad \operatorname{holomorphic\text{-}on}(u, f) \wedge \\ \quad \forall y \in u.\; \neg (f(y) = 0) \wedge \\ \quad q \in (u \to \mathbb{C}) \wedge \\ \quad \forall y \in u.\; (q(y) \simeq \operatorname{recip}(f(y))) \Rightarrow \\ \quad \quad \operatorname{holomorphic\text{-}on}(u, q) \end{array}}}
\vnbprop{holomorphic-quotient}{\vnbstmt{\begin{array}{@{}l@{}} \forall u, f, g, q. \\ \quad \operatorname{holomorphic\text{-}on}(u, f) \wedge \\ \quad \operatorname{holomorphic\text{-}on}(u, g) \wedge \\ \quad \forall y \in u.\; \neg (g(y) = 0) \wedge \\ \quad q \in (u \to \mathbb{C}) \wedge \\ \quad \forall y \in u.\; (q(y) \simeq /(f(y), g(y))) \Rightarrow \\ \quad \quad \operatorname{holomorphic\text{-}on}(u, q) \end{array}}}

\vnbbookfile{structure-library/interval-calculus}{structure-library/interval-calculus}
\noindent{\small interval-calculus.scm -- THE VOCABULARY OF THE CALCULUS ON AN INTERVAL. DEFINITIONS ONLY; every law is PROVEN in theorem-library/interval-calculus-laws.scm and interval-mvt.scm.}\par
\par\medskip\noindent\textbf{Definition} (\texttt{ooint}).\;$\operatorname{ooint}(a, b) :=$ $\{\operatorname{icx\_} \in \mathbb{R} : ((a < \operatorname{icx\_}) \wedge (\operatorname{icx\_} < b))\}$
\par\medskip\noindent\textbf{Definition} (\texttt{is-continuous-on}).\;For all $f$ and $d$, 
\begin{equation*}
\begin{array}{@{}l@{}} \operatorname{is\text{-}continuous\text{-}on}(f, d) \\ \Leftrightarrow\; ((f \in (d \to \mathbb{R})) \wedge \forall \operatorname{icp\_} \in d.\; \operatorname{is\text{-}continuous\text{-}at}(\operatorname{subspace\text{-}ms}(\operatorname{rr\text{-}ms}, d), \operatorname{rr\text{-}ms}, f, \operatorname{icp\_})) \end{array}
\end{equation*}

\par\medskip\noindent\textbf{Definition} (\texttt{has-deriv-at}).\;For all $f$, $a$ and $l$, 
\begin{equation*}
\begin{array}{@{}l@{}} \operatorname{has\text{-}deriv\text{-}at}(f, a, l) \\ \Leftrightarrow\; \exists \operatorname{ice\_}. \\ \quad \operatorname{pos\text{-}rr}(\operatorname{ice\_}) \wedge \\ \quad \operatorname{is\text{-}diff\text{-}on}(\operatorname{rr\text{-}normed\text{-}field}, \operatorname{ooint}(-(a, \operatorname{ice\_}), +(a, \operatorname{ice\_})), \operatorname{restrict}(f, \operatorname{ooint}(-(a, \operatorname{ice\_}), +(a, \operatorname{ice\_}))), a, l) \end{array}
\end{equation*}

\par\medskip\noindent\textbf{Definition} (\texttt{extend-const}).\;$\operatorname{extend\text{-}const}(f, a, b) :=$ $(\lambda \operatorname{ecx\_} \in \mathbb{R}.\; f(\operatorname{clamp}(a, b, \operatorname{ecx\_})))$

\vnbbookfile{structure-library/regulated-primitive}{structure-library/regulated-primitive}
\noindent{\small regulated-primitive.scm -- COUNTABLE SETS, ONE-SIDED LIMITS WITHIN A SET, REGULATED FUNCTIONS ON AN INTERVAL, PRIMITIVES.  DEFINITIONS ONLY; the laws are PROVEN in theorem-library/regulated-primitive-laws.scm.}\par
\par\medskip\noindent\textbf{Definition} (\texttt{is-countable}).\;For every $\operatorname{rpc\_}$, 
\begin{equation*}
\begin{array}{@{}l@{}} \operatorname{is\text{-}countable}(\operatorname{rpc\_}) \\ \Leftrightarrow\; \operatorname{rpc\_} \in \mathbf{Set} \wedge \\ \operatorname{rpc\_} = \emptyset \vee \\ \exists \operatorname{rpen\_} \in (\mathbb{N} \to \operatorname{rpc\_}). \forall \operatorname{rpy\_} \in \operatorname{rpc\_}. \exists \operatorname{rpn\_} \in \mathbb{N}. \\ \quad \operatorname{rpy\_} = \operatorname{rpen\_}(\operatorname{rpn\_}) \end{array}
\end{equation*}

\par\medskip\noindent\textbf{Definition} (\texttt{is-right-limit-within}).\;For all $\operatorname{rpf\_}$, $\operatorname{rpw\_}$, $\operatorname{rpx\_}$ and $\operatorname{rpl\_}$, 
\begin{equation*}
\begin{array}{@{}l@{}} \operatorname{is\text{-}right\text{-}limit\text{-}within}(\operatorname{rpf\_}, \operatorname{rpw\_}, \operatorname{rpx\_}, \operatorname{rpl\_}) \\ \Leftrightarrow\; \operatorname{rpw\_} \subseteq \mathbb{R} \wedge \\ \operatorname{rpf\_} \in (\operatorname{rpw\_} \to \mathbb{R}) \wedge \\ \operatorname{rpx\_} \in \mathbb{R} \wedge \\ \operatorname{rpl\_} \in \mathbb{R} \wedge \\ \forall \operatorname{rpe\_}. \\ \quad \operatorname{pos\text{-}rr}(\operatorname{rpe\_}) \Rightarrow  \\ \quad \quad \exists \operatorname{rpdl\_}. \\ \quad \quad \quad \operatorname{pos\text{-}rr}(\operatorname{rpdl\_}) \wedge \\ \quad \quad \quad \forall \operatorname{rpt\_} \in \operatorname{rpw\_}. \\ \quad \quad \quad \quad (\operatorname{rpx\_} < \operatorname{rpt\_}) \Rightarrow ((\operatorname{rpt\_} < +(\operatorname{rpx\_}, \operatorname{rpdl\_})) \Rightarrow (\lvert -(\operatorname{rpf\_}(\operatorname{rpt\_}), \operatorname{rpl\_}) \rvert < \operatorname{rpe\_})) \end{array}
\end{equation*}

\par\medskip\noindent\textbf{Definition} (\texttt{is-left-limit-within}).\;For all $\operatorname{rpf\_}$, $\operatorname{rpw\_}$, $\operatorname{rpx\_}$ and $\operatorname{rpl\_}$, 
\begin{equation*}
\begin{array}{@{}l@{}} \operatorname{is\text{-}left\text{-}limit\text{-}within}(\operatorname{rpf\_}, \operatorname{rpw\_}, \operatorname{rpx\_}, \operatorname{rpl\_}) \\ \Leftrightarrow\; \operatorname{rpw\_} \subseteq \mathbb{R} \wedge \\ \operatorname{rpf\_} \in (\operatorname{rpw\_} \to \mathbb{R}) \wedge \\ \operatorname{rpx\_} \in \mathbb{R} \wedge \\ \operatorname{rpl\_} \in \mathbb{R} \wedge \\ \forall \operatorname{rple\_}. \\ \quad \operatorname{pos\text{-}rr}(\operatorname{rple\_}) \Rightarrow  \\ \quad \quad \exists \operatorname{rpld\_}. \\ \quad \quad \quad \operatorname{pos\text{-}rr}(\operatorname{rpld\_}) \wedge \\ \quad \quad \quad \forall \operatorname{rplt\_} \in \operatorname{rpw\_}. \\ \quad \quad \quad \quad (-(\operatorname{rpx\_}, \operatorname{rpld\_}) < \operatorname{rplt\_}) \Rightarrow  \\ \quad \quad \quad \quad (\operatorname{rplt\_} < \operatorname{rpx\_}) \Rightarrow (\lvert -(\operatorname{rpf\_}(\operatorname{rplt\_}), \operatorname{rpl\_}) \rvert < \operatorname{rple\_}) \end{array}
\end{equation*}

\par\medskip\noindent\textbf{Definition} (\texttt{is-regulated-on}).\;For all $\operatorname{rpf\_}$, $\operatorname{rpa\_}$ and $\operatorname{rpb\_}$, 
\begin{equation*}
\begin{array}{@{}l@{}} \operatorname{is\text{-}regulated\text{-}on}(\operatorname{rpf\_}, \operatorname{rpa\_}, \operatorname{rpb\_}) \\ \Leftrightarrow\; \operatorname{rpa\_} \in \mathbb{R} \wedge \\ \operatorname{rpb\_} \in \mathbb{R} \wedge \\ \operatorname{rpa\_} < \operatorname{rpb\_} \wedge \\ \operatorname{rpf\_} \in (\operatorname{ccint}(\operatorname{rpa\_}, \operatorname{rpb\_}) \to \mathbb{R}) \wedge \\ \forall \operatorname{rpx\_} \in \operatorname{ccint}(\operatorname{rpa\_}, \operatorname{rpb\_}). \\ \quad (\operatorname{rpx\_} < \operatorname{rpb\_}) \Rightarrow  \\ \quad \quad \exists \operatorname{rpl\_}.\; \operatorname{is\text{-}right\text{-}limit\text{-}within}(\operatorname{rpf\_}, \operatorname{ccint}(\operatorname{rpa\_}, \operatorname{rpb\_}), \operatorname{rpx\_}, \operatorname{rpl\_}) \wedge \\ \forall \operatorname{rpx\_} \in \operatorname{ccint}(\operatorname{rpa\_}, \operatorname{rpb\_}). \\ \quad (\operatorname{rpa\_} < \operatorname{rpx\_}) \Rightarrow  \\ \quad \quad \exists \operatorname{rpl\_}.\; \operatorname{is\text{-}left\text{-}limit\text{-}within}(\operatorname{rpf\_}, \operatorname{ccint}(\operatorname{rpa\_}, \operatorname{rpb\_}), \operatorname{rpx\_}, \operatorname{rpl\_}) \end{array}
\end{equation*}

\par\medskip\noindent\textbf{Definition} (\texttt{is-primitive}).\;For all $\operatorname{rpg\_}$, $\operatorname{rpf\_}$, $\operatorname{rpa\_}$ and $\operatorname{rpb\_}$, 
\begin{equation*}
\begin{array}{@{}l@{}} \operatorname{is\text{-}primitive}(\operatorname{rpg\_}, \operatorname{rpf\_}, \operatorname{rpa\_}, \operatorname{rpb\_}) \\ \Leftrightarrow\; \operatorname{rpa\_} \in \mathbb{R} \wedge \\ \operatorname{rpb\_} \in \mathbb{R} \wedge \\ \operatorname{rpa\_} < \operatorname{rpb\_} \wedge \\ \operatorname{rpg\_} \in (\operatorname{ccint}(\operatorname{rpa\_}, \operatorname{rpb\_}) \to \mathbb{R}) \wedge \\ \operatorname{rpf\_} \in (\operatorname{ccint}(\operatorname{rpa\_}, \operatorname{rpb\_}) \to \mathbb{R}) \wedge \\ \operatorname{is\text{-}continuous\text{-}on}(\operatorname{rpg\_}, \operatorname{ccint}(\operatorname{rpa\_}, \operatorname{rpb\_})) \wedge \\ \exists \operatorname{rps\_}. \\ \quad \operatorname{is\text{-}countable}(\operatorname{rps\_}) \wedge \\ \quad \operatorname{rps\_} \subseteq \operatorname{ccint}(\operatorname{rpa\_}, \operatorname{rpb\_}) \wedge \\ \quad \forall \operatorname{rpt\_} \in \operatorname{ooint}(\operatorname{rpa\_}, \operatorname{rpb\_}). \\ \quad \quad \neg (\operatorname{rpt\_} \in \operatorname{rps\_}) \Rightarrow \operatorname{has\text{-}deriv\text{-}at}(\operatorname{rpg\_}, \operatorname{rpt\_}, \operatorname{rpf\_}(\operatorname{rpt\_})) \end{array}
\end{equation*}


\vnbbookfile{theorem-library/interval-calculus-laws}{theorem-library/interval-calculus-laws}
\noindent{\small interval-calculus-laws.scm -- the laws of the vocabulary of structure-library/interval-calculus.scm: OOINT, IS-CONTINUOUS-ON, HAS-DERIV-AT and EXTEND-CONST.  Everything here is PROVEN; the definition file asserts nothing.}\par
\vnbprop{ooint-membership}{\vnbstmt{\forall a, b, y.\; ((y \in \operatorname{ooint}(a, b)) \Leftrightarrow ((y \in \mathbb{R}) \wedge ((a < y) \wedge (y < b))))}}
\vnbprop{ooint-elt-in-rr}{\vnbstmt{\forall a, b.; \forall y \in \operatorname{ooint}(a, b).\; (y \in \mathbb{R})}}
\vnbprop{ooint-subset-rr}{\vnbstmt{\forall a, b.\; (\operatorname{ooint}(a, b) \subseteq \mathbb{R})}}
\vnbprop{ooint-subset-ccint}{\vnbstmt{\forall a, b.\; (((a \in \mathbb{R}) \wedge (b \in \mathbb{R})) \Rightarrow (\operatorname{ooint}(a, b) \subseteq \operatorname{ccint}(a, b)))}}
\vnbprop{ooint-center}{\vnbstmt{\forall y, e.\; (((y \in \mathbb{R}) \wedge \operatorname{pos\text{-}rr}(e)) \Rightarrow (y \in \operatorname{ooint}((y - e), (y + e))))}}
\vnbprop{ooint-open}{\vnbstmt{\forall a, b.\; (((a \in \mathbb{R}) \wedge (b \in \mathbb{R})) \Rightarrow \operatorname{is\text{-}open}(\operatorname{rr\text{-}ms}, \operatorname{ooint}(a, b)))}}
\vnbprop{ooint-open-nf}{\vnbstmt{\begin{array}{@{}l@{}} \forall a, b. \\ \quad ((a \in \mathbb{R}) \wedge (b \in \mathbb{R})) \Rightarrow  \\ \quad \quad \operatorname{is\text{-}open}(\operatorname{nf\text{-}metric\text{-}space}(\operatorname{rr\text{-}normed\text{-}field}), \operatorname{ooint}(a, b)) \end{array}}}
\vnbprop{subspace-restrict-continuous-at}{\vnbstmt{\begin{array}{@{}l@{}} \forall s, u, v, t, f.; \forall z \in v. \\ \quad \operatorname{is\text{-}metric\text{-}space}(s) \wedge \\ \quad u \subseteq \operatorname{pts}(s) \wedge \\ \quad v \subseteq u \wedge \\ \quad \operatorname{is\text{-}metric\text{-}space}(t) \wedge \\ \quad \operatorname{is\text{-}continuous\text{-}at}(\operatorname{subspace\text{-}ms}(s, u), t, f, z) \Rightarrow \\ \quad \quad \operatorname{is\text{-}continuous\text{-}at}(\operatorname{subspace\text{-}ms}(s, v), t, \operatorname{restrict}(f, v), z) \end{array}}}
\vnbprop{diff-on-shrink}{\vnbstmt{\begin{array}{@{}l@{}} \forall k, u, f, a, l, v. \\ \quad (\operatorname{is\text{-}diff\text{-}on}(k, u, f, a, l) \wedge (\operatorname{is\text{-}open}(\operatorname{nf\text{-}metric\text{-}space}(k), v) \wedge ((v \subseteq u) \wedge (a \in v)))) \Rightarrow  \\ \quad \quad \operatorname{is\text{-}diff\text{-}on}(k, v, \operatorname{restrict}(f, v), a, l) \end{array}}}
\vnbprop{restrict-in-fun-of-restrict}{\vnbstmt{\begin{array}{@{}l@{}} \forall f, u, c, v. \\ \quad ((\operatorname{restrict}(f, u) \in (u \to c)) \wedge (v \subseteq u)) \Rightarrow (\operatorname{restrict}(f, v) \in (v \to c)) \end{array}}}
\vnbprop{diff-on-shrink-restrict}{\vnbstmt{\begin{array}{@{}l@{}} \forall k, u, f, a, l, v. \\ \quad (\operatorname{is\text{-}diff\text{-}on}(k, u, \operatorname{restrict}(f, u), a, l) \wedge (\operatorname{is\text{-}open}(\operatorname{nf\text{-}metric\text{-}space}(k), v) \wedge ((v \subseteq u) \wedge (a \in v)))) \Rightarrow  \\ \quad \quad \operatorname{is\text{-}diff\text{-}on}(k, v, \operatorname{restrict}(f, v), a, l) \end{array}}}
\vnbprop{diff-on-ooint-shrink}{\vnbstmt{\begin{array}{@{}l@{}} \forall a \in \mathbb{R}, e, w, f, l. \\ \quad \operatorname{pos\text{-}rr}(e) \wedge \\ \quad \operatorname{pos\text{-}rr}(w) \wedge \\ \quad w \leq e \wedge \\ \quad \operatorname{is\text{-}diff\text{-}on}(\operatorname{rr\text{-}normed\text{-}field}, \operatorname{ooint}((a - e), (a + e)), \operatorname{restrict}(f, \operatorname{ooint}((a - e), (a + e))), a, l) \Rightarrow \\ \quad \quad \operatorname{is\text{-}diff\text{-}on}(\operatorname{rr\text{-}normed\text{-}field}, \operatorname{ooint}((a - w), (a + w)), \operatorname{restrict}(f, \operatorname{ooint}((a - w), (a + w))), a, l) \end{array}}}
\vnbprop{has-deriv-at-intro}{\vnbstmt{\begin{array}{@{}l@{}} \forall f, a, l, e. \\ \quad \operatorname{pos\text{-}rr}(e) \wedge \\ \quad \operatorname{is\text{-}diff\text{-}on}(\operatorname{rr\text{-}normed\text{-}field}, \operatorname{ooint}((a - e), (a + e)), \operatorname{restrict}(f, \operatorname{ooint}((a - e), (a + e))), a, l) \Rightarrow \\ \quad \quad \operatorname{has\text{-}deriv\text{-}at}(f, a, l) \end{array}}}
\vnbprop{has-deriv-at-shrink}{\vnbstmt{\begin{array}{@{}l@{}} \forall f, a, l, r. \\ \quad (\operatorname{has\text{-}deriv\text{-}at}(f, a, l) \wedge \operatorname{pos\text{-}rr}(r)) \Rightarrow  \\ \quad \quad \exists w. \\ \quad \quad \quad \operatorname{pos\text{-}rr}(w) \wedge \\ \quad \quad \quad w \leq r \wedge \\ \quad \quad \quad \operatorname{is\text{-}diff\text{-}on}(\operatorname{rr\text{-}normed\text{-}field}, \operatorname{ooint}((a - w), (a + w)), \operatorname{restrict}(f, \operatorname{ooint}((a - w), (a + w))), a, l) \end{array}}}
\vnbprop{has-deriv-at-pt-in-rr}{\vnbstmt{\forall f, a, l.\; (\operatorname{has\text{-}deriv\text{-}at}(f, a, l) \Rightarrow (a \in \mathbb{R}))}}
\vnbprop{has-deriv-at-in-rr}{\vnbstmt{\forall f, a, l.\; (\operatorname{has\text{-}deriv\text{-}at}(f, a, l) \Rightarrow (l \in \mathbb{R}))}}
\vnbprop{has-deriv-at-unique}{\vnbstmt{\begin{array}{@{}l@{}} \forall f, a, l, m. \\ \quad (\operatorname{has\text{-}deriv\text{-}at}(f, a, l) \wedge \operatorname{has\text{-}deriv\text{-}at}(f, a, m)) \Rightarrow (l = m) \end{array}}}
\vnbprop{diff-on-nbhd-implies-diff-at}{\vnbstmt{\begin{array}{@{}l@{}} \forall f \in (\mathbb{R} \to \mathbb{R}), a \in \mathbb{R}, e, l. \\ \quad \operatorname{pos\text{-}rr}(e) \wedge \\ \quad \operatorname{is\text{-}diff\text{-}on}(\operatorname{rr\text{-}normed\text{-}field}, \operatorname{ooint}((a - e), (a + e)), \operatorname{restrict}(f, \operatorname{ooint}((a - e), (a + e))), a, l) \Rightarrow \\ \quad \quad \operatorname{is\text{-}diff\text{-}at}(f, a, l) \end{array}}}
\vnbprop{has-deriv-at-iff-diff-at}{\vnbstmt{\forall f \in (\mathbb{R} \to \mathbb{R}), a, l.\; (\operatorname{has\text{-}deriv\text{-}at}(f, a, l) \Leftrightarrow \operatorname{is\text{-}diff\text{-}at}(f, a, l))}}
\vnbprop{extend-const-value}{\vnbstmt{\forall f, a, b.; \forall z \in \mathbb{R}.\; (\operatorname{extend\text{-}const}(f, a, b)(z) \simeq f(\operatorname{clamp}(a, b, z)))}}
\vnbprop{extend-const-in-fun}{\vnbstmt{\begin{array}{@{}l@{}} \forall a, b \in \mathbb{R}, f \in (\operatorname{ccint}(a, b) \to \mathbb{R}). \\ \quad (a \leq b) \Rightarrow (\operatorname{extend\text{-}const}(f, a, b) \in (\mathbb{R} \to \mathbb{R})) \end{array}}}
\vnbprop{extend-const-fixes}{\vnbstmt{\begin{array}{@{}l@{}} \forall a, b, z \in \mathbb{R}, f. \\ \quad ((a \leq z) \wedge (z \leq b)) \Rightarrow (\operatorname{extend\text{-}const}(f, a, b)(z) \simeq f(z)) \end{array}}}
\vnbprop{extend-const-ccint-fixes}{\vnbstmt{\begin{array}{@{}l@{}} \forall a, b \in \mathbb{R}, z \in \operatorname{ccint}(a, b), f. \\ \quad (\operatorname{extend\text{-}const}(f, a, b)(z) \simeq f(z)) \end{array}}}
\vnbprop{extend-const-continuous-at}{\vnbstmt{\begin{array}{@{}l@{}} \forall a, b \in \mathbb{R}, f.; \forall t \in \mathbb{R}. \\ \quad ((a \leq b) \wedge \operatorname{is\text{-}continuous\text{-}on}(f, \operatorname{ccint}(a, b))) \Rightarrow  \\ \quad \quad \operatorname{is\text{-}continuous\text{-}at}(\operatorname{rr\text{-}ms}, \operatorname{rr\text{-}ms}, \operatorname{extend\text{-}const}(f, a, b), t) \end{array}}}
\vnbprop{extend-const-deriv-fwd}{\vnbstmt{\begin{array}{@{}l@{}} \forall a, b \in \mathbb{R}, f \in (\operatorname{ccint}(a, b) \to \mathbb{R}), x \in \operatorname{ooint}(a, b), l. \\ \quad ((a \leq b) \wedge \operatorname{has\text{-}deriv\text{-}at}(f, x, l)) \Rightarrow \operatorname{is\text{-}diff\text{-}at}(\operatorname{extend\text{-}const}(f, a, b), x, l) \end{array}}}
\vnbprop{extend-const-deriv-bwd}{\vnbstmt{\begin{array}{@{}l@{}} \forall a, b \in \mathbb{R}, f \in (\operatorname{ccint}(a, b) \to \mathbb{R}), x \in \operatorname{ooint}(a, b), l. \\ \quad ((a \leq b) \wedge \operatorname{is\text{-}diff\text{-}at}(\operatorname{extend\text{-}const}(f, a, b), x, l)) \Rightarrow \operatorname{has\text{-}deriv\text{-}at}(f, x, l) \end{array}}}
\vnbprop{extend-const-deriv-iff}{\vnbstmt{\begin{array}{@{}l@{}} \forall a, b \in \mathbb{R}, f \in (\operatorname{ccint}(a, b) \to \mathbb{R}), x \in \operatorname{ooint}(a, b), l. \\ \quad (a \leq b) \Rightarrow (\operatorname{has\text{-}deriv\text{-}at}(f, x, l) \Leftrightarrow \operatorname{is\text{-}diff\text{-}at}(\operatorname{extend\text{-}const}(f, a, b), x, l)) \end{array}}}
\vnbprop{has-deriv-at-local}{\vnbstmt{\begin{array}{@{}l@{}} \forall r, f, g, x, l. \\ \quad \operatorname{pos\text{-}rr}(r) \wedge \\ \quad \operatorname{restrict}(f, \operatorname{ooint}((x - r), (x + r))) \in (\operatorname{ooint}((x - r), (x + r)) \to \mathbb{R}) \wedge \\ \quad \forall \mathit{hbx} \in \operatorname{ooint}((x - r), (x + r)).\; (f(\mathit{hbx}) \simeq g(\mathit{hbx})) \wedge \\ \quad \operatorname{has\text{-}deriv\text{-}at}(g, x, l) \Rightarrow \\ \quad \quad \operatorname{has\text{-}deriv\text{-}at}(f, x, l) \end{array}}}
\vnbprop{has-deriv-at-sum}{\vnbstmt{\begin{array}{@{}l@{}} \forall r, f, g, h, x, l, m. \\ \quad \operatorname{pos\text{-}rr}(r) \wedge \\ \quad \operatorname{restrict}(h, \operatorname{ooint}((x - r), (x + r))) \in (\operatorname{ooint}((x - r), (x + r)) \to \mathbb{R}) \wedge \\ \quad \forall \mathit{hbx} \in \operatorname{ooint}((x - r), (x + r)).\; (h(\mathit{hbx}) \simeq (f(\mathit{hbx}) + g(\mathit{hbx}))) \wedge \\ \quad \operatorname{has\text{-}deriv\text{-}at}(f, x, l) \wedge \\ \quad \operatorname{has\text{-}deriv\text{-}at}(g, x, m) \Rightarrow \\ \quad \quad \operatorname{has\text{-}deriv\text{-}at}(h, x, (l + m)) \end{array}}}
\vnbprop{has-deriv-at-product}{\vnbstmt{\begin{array}{@{}l@{}} \forall r, f, g, h, x, l, m. \\ \quad \operatorname{pos\text{-}rr}(r) \wedge \\ \quad \operatorname{restrict}(h, \operatorname{ooint}((x - r), (x + r))) \in (\operatorname{ooint}((x - r), (x + r)) \to \mathbb{R}) \wedge \\ \quad \forall \mathit{hbx} \in \operatorname{ooint}((x - r), (x + r)).\; (h(\mathit{hbx}) \simeq (f(\mathit{hbx}) \cdot g(\mathit{hbx}))) \wedge \\ \quad \operatorname{has\text{-}deriv\text{-}at}(f, x, l) \wedge \\ \quad \operatorname{has\text{-}deriv\text{-}at}(g, x, m) \Rightarrow \\ \quad \quad \operatorname{has\text{-}deriv\text{-}at}(h, x, ((l \cdot g(x)) + (m \cdot f(x)))) \end{array}}}

\vnbbookfile{theorem-library/interval-mvt}{theorem-library/interval-mvt}
\noindent{\small interval-mvt.scm -- ROLLE, THE TWO MEAN VALUE THEOREMS AND THE THREE BOUNDS, STATED FOR A FUNCTION ON ITS INTERVAL.}\par
\vnbprop{rolle-on-interval}{\vnbstmt{\begin{array}{@{}l@{}} \forall a, b \in \mathbb{R}, h \in (\operatorname{ccint}(a, b) \to \mathbb{R}). \\ \quad a < b \wedge \\ \quad \operatorname{is\text{-}continuous\text{-}on}(h, \operatorname{ccint}(a, b)) \wedge \\ \quad \forall x \in \operatorname{ooint}(a, b).\; \exists l.\; \operatorname{has\text{-}deriv\text{-}at}(h, x, l) \wedge \\ \quad h(a) = h(b) \Rightarrow \\ \quad \quad \exists c \in \operatorname{ooint}(a, b).\; \operatorname{has\text{-}deriv\text{-}at}(h, c, 0) \end{array}}}
\vnbprop{generalized-mvt-on-interval}{\vnbstmt{\begin{array}{@{}l@{}} \forall a, b \in \mathbb{R}, f, g. \\ \quad a < b \wedge \\ \quad f \in (\operatorname{ccint}(a, b) \to \mathbb{R}) \wedge \\ \quad g \in (\operatorname{ccint}(a, b) \to \mathbb{R}) \wedge \\ \quad \operatorname{is\text{-}continuous\text{-}on}(f, \operatorname{ccint}(a, b)) \wedge \\ \quad \operatorname{is\text{-}continuous\text{-}on}(g, \operatorname{ccint}(a, b)) \wedge \\ \quad \forall x \in \operatorname{ooint}(a, b).\; \exists l.\; \operatorname{has\text{-}deriv\text{-}at}(f, x, l) \wedge \\ \quad \forall x \in \operatorname{ooint}(a, b).\; \exists m.\; \operatorname{has\text{-}deriv\text{-}at}(g, x, m) \Rightarrow \\ \quad \quad \exists c \in \operatorname{ooint}(a, b), l, m. \\ \quad \quad \quad (\operatorname{has\text{-}deriv\text{-}at}(f, c, l) \wedge (\operatorname{has\text{-}deriv\text{-}at}(g, c, m) \wedge ((l \cdot (g(b) - g(a))) = (m \cdot (f(b) - f(a)))))) \end{array}}}
\vnbprop{mvt-on-interval}{\vnbstmt{\begin{array}{@{}l@{}} \forall a, b \in \mathbb{R}, f \in (\operatorname{ccint}(a, b) \to \mathbb{R}). \\ \quad a < b \wedge \\ \quad \operatorname{is\text{-}continuous\text{-}on}(f, \operatorname{ccint}(a, b)) \wedge \\ \quad \forall x \in \operatorname{ooint}(a, b).\; \exists l.\; \operatorname{has\text{-}deriv\text{-}at}(f, x, l) \Rightarrow \\ \quad \quad \exists c \in \operatorname{ooint}(a, b), l. \\ \quad \quad \quad \operatorname{has\text{-}deriv\text{-}at}(f, c, l) \wedge (l = ((f(b) - f(a)) \cdot \operatorname{recip}((b - a)))) \end{array}}}
\vnbprop{mvt-upper-bound-on-interval}{\vnbstmt{\begin{array}{@{}l@{}} \forall a, b, m \in \mathbb{R}, f \in (\operatorname{ccint}(a, b) \to \mathbb{R}). \\ \quad a < b \wedge \\ \quad \operatorname{is\text{-}continuous\text{-}on}(f, \operatorname{ccint}(a, b)) \wedge \\ \quad \forall x \in \operatorname{ooint}(a, b).\; \exists l.\; (\operatorname{has\text{-}deriv\text{-}at}(f, x, l) \wedge (l \leq m)) \Rightarrow \\ \quad \quad (f(b) - f(a)) \leq (m \cdot (b - a)) \end{array}}}
\vnbprop{mvt-lower-bound-on-interval}{\vnbstmt{\begin{array}{@{}l@{}} \forall a, b, m \in \mathbb{R}, f \in (\operatorname{ccint}(a, b) \to \mathbb{R}). \\ \quad a < b \wedge \\ \quad \operatorname{is\text{-}continuous\text{-}on}(f, \operatorname{ccint}(a, b)) \wedge \\ \quad \forall x \in \operatorname{ooint}(a, b).\; \exists l.\; (\operatorname{has\text{-}deriv\text{-}at}(f, x, l) \wedge (m \leq l)) \Rightarrow \\ \quad \quad (m \cdot (b - a)) \leq (f(b) - f(a)) \end{array}}}
\vnbprop{mvt-abs-bound-on-interval}{\vnbstmt{\begin{array}{@{}l@{}} \forall a, b, m \in \mathbb{R}, f \in (\operatorname{ccint}(a, b) \to \mathbb{R}). \\ \quad a < b \wedge \\ \quad \operatorname{is\text{-}continuous\text{-}on}(f, \operatorname{ccint}(a, b)) \wedge \\ \quad \forall x \in \operatorname{ooint}(a, b).\; \exists l.\; (\operatorname{has\text{-}deriv\text{-}at}(f, x, l) \wedge (\lvert l \rvert \leq m)) \Rightarrow \\ \quad \quad \lvert (f(b) - f(a)) \rvert \leq (m \cdot (b - a)) \end{array}}}

\vnbbookfile{structure-library/interval-extremum}{structure-library/interval-extremum}
\noindent{\small interval-extremum.scm -- LOCAL MAXIMUM AND LOCAL MINIMUM OF A FUNCTION ON A SET.  DEFINITIONS ONLY; Proposition 2.10 is proven in theorem-library/interval-extremum.scm.}\par
\par\medskip\noindent\textbf{Definition} (\texttt{is-local-max-at}).\;For all $f$, $d$ and $\operatorname{lmt\_}$, 
\begin{equation*}
\begin{array}{@{}l@{}} \operatorname{is\text{-}local\text{-}max\text{-}at}(f, d, \operatorname{lmt\_}) \\ \Leftrightarrow\; f \in (d \to \mathbb{R}) \wedge \\ \operatorname{lmt\_} \in d \wedge \\ \exists \operatorname{lma\_}. \\ \quad (\operatorname{pos\text{-}rr}(\operatorname{lma\_}) \wedge \forall \operatorname{lmx\_} \in d.\; ((\lvert -(\operatorname{lmx\_}, \operatorname{lmt\_}) \rvert \leq \operatorname{lma\_}) \Rightarrow (f(\operatorname{lmx\_}) \leq f(\operatorname{lmt\_})))) \end{array}
\end{equation*}

\par\medskip\noindent\textbf{Definition} (\texttt{is-local-min-at}).\;For all $f$, $d$ and $\operatorname{lmt\_}$, 
\begin{equation*}
\begin{array}{@{}l@{}} \operatorname{is\text{-}local\text{-}min\text{-}at}(f, d, \operatorname{lmt\_}) \\ \Leftrightarrow\; f \in (d \to \mathbb{R}) \wedge \\ \operatorname{lmt\_} \in d \wedge \\ \exists \operatorname{lma\_}. \\ \quad (\operatorname{pos\text{-}rr}(\operatorname{lma\_}) \wedge \forall \operatorname{lmx\_} \in d.\; ((\lvert -(\operatorname{lmx\_}, \operatorname{lmt\_}) \rvert \leq \operatorname{lma\_}) \Rightarrow (f(\operatorname{lmt\_}) \leq f(\operatorname{lmx\_})))) \end{array}
\end{equation*}


\vnbbookfile{theorem-library/interval-extremum}{theorem-library/interval-extremum}
\noindent{\small interval-extremum.scm -- PROPOSITION 2.10 FOR A FUNCTION ON ITS INTERVAL.}\par
\vnbprop{interval-local-max-deriv-zero}{\vnbstmt{\begin{array}{@{}l@{}} \forall a, b \in \mathbb{R}, f.; \forall c \in \operatorname{ooint}(a, b), l. \\ \quad (\operatorname{is\text{-}local\text{-}max\text{-}at}(f, \operatorname{ccint}(a, b), c) \wedge \operatorname{has\text{-}deriv\text{-}at}(f, c, l)) \Rightarrow (l = 0) \end{array}}}
\vnbprop{interval-local-min-deriv-zero}{\vnbstmt{\begin{array}{@{}l@{}} \forall a, b \in \mathbb{R}, f.; \forall c \in \operatorname{ooint}(a, b), l. \\ \quad (\operatorname{is\text{-}local\text{-}min\text{-}at}(f, \operatorname{ccint}(a, b), c) \wedge \operatorname{has\text{-}deriv\text{-}at}(f, c, l)) \Rightarrow (l = 0) \end{array}}}

\vnbbookfile{structure-library/one-sided-derivative}{structure-library/one-sided-derivative}
\noindent{\small one-sided-derivative.scm -- THE ONE-SIDED DERIVATIVES OF THE USER'S NOTES, section 2.1.  DEFINITIONS ONLY; every law is PROVEN in theorem-library/one-sided-derivative-laws.scm.}\par
\par\medskip\noindent\textbf{Definition} (\texttt{coint}).\;$\operatorname{coint}(a, b) :=$ $\{\operatorname{osq\_} \in \mathbb{R} : ((a \leq \operatorname{osq\_}) \wedge (\operatorname{osq\_} < b))\}$
\par\medskip\noindent\textbf{Definition} (\texttt{ocint}).\;$\operatorname{ocint}(a, b) :=$ $\{\operatorname{osq\_} \in \mathbb{R} : ((a < \operatorname{osq\_}) \wedge (\operatorname{osq\_} \leq b))\}$
\par\medskip\noindent\textbf{Definition} (\texttt{has-deriv-at-within}).\;For all $f$, $w$, $x$ and $l$, 
\begin{equation*}
\begin{array}{@{}l@{}} \operatorname{has\text{-}deriv\text{-}at\text{-}within}(f, w, x, l) \\ \Leftrightarrow\; \operatorname{restrict}(f, w) \in (w \to \mathbb{R}) \wedge \\ x \in w \wedge \\ l \in \mathbb{R} \wedge \\ \exists \operatorname{osp\_} \in (w \to \mathbb{R}). \\ \quad \operatorname{is\text{-}continuous\text{-}at}(\operatorname{subspace\text{-}ms}(\operatorname{rr\text{-}ms}, w), \operatorname{rr\text{-}ms}, \operatorname{osp\_}, x) \wedge \\ \quad \operatorname{osp\_}(x) = l \wedge \\ \quad \forall \operatorname{ost\_} \in w.\; (-(f(\operatorname{ost\_}), f(x)) = \cdot(\operatorname{osp\_}(\operatorname{ost\_}), -(\operatorname{ost\_}, x))) \end{array}
\end{equation*}

\par\medskip\noindent\textbf{Definition} (\texttt{has-right-deriv-at}).\;For all $f$, $x$ and $l$, 
\begin{equation*}
\begin{array}{@{}l@{}} \operatorname{has\text{-}right\text{-}deriv\text{-}at}(f, x, l) \\ \Leftrightarrow\; \exists \operatorname{ose\_}. \\ \quad \operatorname{pos\text{-}rr}(\operatorname{ose\_}) \wedge \operatorname{has\text{-}deriv\text{-}at\text{-}within}(f, \operatorname{coint}(x, +(x, \operatorname{ose\_})), x, l) \end{array}
\end{equation*}

\par\medskip\noindent\textbf{Definition} (\texttt{has-left-deriv-at}).\;For all $f$, $x$ and $l$, 
\begin{equation*}
\begin{array}{@{}l@{}} \operatorname{has\text{-}left\text{-}deriv\text{-}at}(f, x, l) \\ \Leftrightarrow\; \exists \operatorname{ose\_}. \\ \quad \operatorname{pos\text{-}rr}(\operatorname{ose\_}) \wedge \operatorname{has\text{-}deriv\text{-}at\text{-}within}(f, \operatorname{ocint}(-(x, \operatorname{ose\_}), x), x, l) \end{array}
\end{equation*}


\vnbbookfile{structure-library/calint-vocab}{structure-library/calint-vocab}
\noindent{\small calint-vocab.scm -- WAVE 2 (2026-10-27..30), wired 2026-10-30 from scratchpad/calint/calint-vocab.scm. Library: analysis.}\par
\par\medskip\noindent\textbf{Definition} (\texttt{has-deriv-on-ccint}).\;For all $\operatorname{hdcg\_}$, $\operatorname{hdcd\_}$, $\operatorname{hdca\_}$ and $\operatorname{hdcb\_}$, 
\begin{equation*}
\begin{array}{@{}l@{}} \operatorname{has\text{-}deriv\text{-}on\text{-}ccint}(\operatorname{hdcg\_}, \operatorname{hdcd\_}, \operatorname{hdca\_}, \operatorname{hdcb\_}) \\ \Leftrightarrow\; \operatorname{hdca\_} \in \mathbb{R} \wedge \\ \operatorname{hdcb\_} \in \mathbb{R} \wedge \\ \operatorname{hdca\_} < \operatorname{hdcb\_} \wedge \\ \operatorname{hdcg\_} \in (\operatorname{ccint}(\operatorname{hdca\_}, \operatorname{hdcb\_}) \to \mathbb{R}) \wedge \\ \operatorname{hdcd\_} \in (\operatorname{ccint}(\operatorname{hdca\_}, \operatorname{hdcb\_}) \to \mathbb{R}) \wedge \\ \forall \operatorname{hdct\_} \in \operatorname{ooint}(\operatorname{hdca\_}, \operatorname{hdcb\_}).\; \operatorname{has\text{-}deriv\text{-}at}(\operatorname{hdcg\_}, \operatorname{hdct\_}, \operatorname{hdcd\_}(\operatorname{hdct\_})) \wedge \\ \operatorname{has\text{-}right\text{-}deriv\text{-}at}(\operatorname{hdcg\_}, \operatorname{hdca\_}, \operatorname{hdcd\_}(\operatorname{hdca\_})) \wedge \\ \operatorname{has\text{-}left\text{-}deriv\text{-}at}(\operatorname{hdcg\_}, \operatorname{hdcb\_}, \operatorname{hdcd\_}(\operatorname{hdcb\_})) \end{array}
\end{equation*}


\vnbbookfile{theorem-library/one-sided-derivative-laws}{theorem-library/one-sided-derivative-laws}
\noindent{\small one-sided-derivative-laws.scm -- the laws of structure-library/one-sided-derivative.scm: COINT, OCINT, HAS-DERIV-AT-WITHIN, HAS-RIGHT-DERIV-AT and HAS-LEFT-DERIV-AT.  Everything here is PROVEN; the definition file asserts nothing.}\par
\vnbprop{coint-membership}{\vnbstmt{\forall a, b, y.\; ((y \in \operatorname{coint}(a, b)) \Leftrightarrow ((y \in \mathbb{R}) \wedge ((a \leq y) \wedge (y < b))))}}
\vnbprop{ocint-membership}{\vnbstmt{\forall a, b, y.\; ((y \in \operatorname{ocint}(a, b)) \Leftrightarrow ((y \in \mathbb{R}) \wedge ((a < y) \wedge (y \leq b))))}}
\vnbprop{coint-elt-in-rr}{\vnbstmt{\forall a, b.; \forall y \in \operatorname{coint}(a, b).\; (y \in \mathbb{R})}}
\vnbprop{ocint-elt-in-rr}{\vnbstmt{\forall a, b.; \forall y \in \operatorname{ocint}(a, b).\; (y \in \mathbb{R})}}
\vnbprop{coint-subset-rr}{\vnbstmt{\forall a, b.\; (\operatorname{coint}(a, b) \subseteq \mathbb{R})}}
\vnbprop{ocint-subset-rr}{\vnbstmt{\forall a, b.\; (\operatorname{ocint}(a, b) \subseteq \mathbb{R})}}
\vnbprop{coint-left-endpoint}{\vnbstmt{\forall x, r.\; (((x \in \mathbb{R}) \wedge \operatorname{pos\text{-}rr}(r)) \Rightarrow (x \in \operatorname{coint}(x, (x + r))))}}
\vnbprop{ocint-right-endpoint}{\vnbstmt{\forall x, r.\; (((x \in \mathbb{R}) \wedge \operatorname{pos\text{-}rr}(r)) \Rightarrow (x \in \operatorname{ocint}((x - r), x)))}}
\vnbprop{coint-subset-ooint}{\vnbstmt{\begin{array}{@{}l@{}} \forall x, r. \\ \quad ((x \in \mathbb{R}) \wedge \operatorname{pos\text{-}rr}(r)) \Rightarrow (\operatorname{coint}(x, (x + r)) \subseteq \operatorname{ooint}((x - r), (x + r))) \end{array}}}
\vnbprop{ocint-subset-ooint}{\vnbstmt{\begin{array}{@{}l@{}} \forall x, r. \\ \quad ((x \in \mathbb{R}) \wedge \operatorname{pos\text{-}rr}(r)) \Rightarrow (\operatorname{ocint}((x - r), x) \subseteq \operatorname{ooint}((x - r), (x + r))) \end{array}}}
\vnbprop{coint-shrink}{\vnbstmt{\begin{array}{@{}l@{}} \forall x, r, v. \\ \quad ((x \in \mathbb{R}) \wedge (\operatorname{pos\text{-}rr}(r) \wedge (\operatorname{pos\text{-}rr}(v) \wedge (v \leq r)))) \Rightarrow  \\ \quad \quad \operatorname{coint}(x, (x + v)) \subseteq \operatorname{coint}(x, (x + r)) \end{array}}}
\vnbprop{ocint-shrink}{\vnbstmt{\begin{array}{@{}l@{}} \forall x, r, v. \\ \quad ((x \in \mathbb{R}) \wedge (\operatorname{pos\text{-}rr}(r) \wedge (\operatorname{pos\text{-}rr}(v) \wedge (v \leq r)))) \Rightarrow  \\ \quad \quad \operatorname{ocint}((x - v), x) \subseteq \operatorname{ocint}((x - r), x) \end{array}}}
\vnbprop{coint-left-limit-point}{\vnbstmt{\begin{array}{@{}l@{}} \forall x, d, r. \\ \quad ((x \in \mathbb{R}) \wedge (\operatorname{pos\text{-}rr}(d) \wedge \operatorname{pos\text{-}rr}(r))) \Rightarrow  \\ \quad \quad \exists b \in \operatorname{coint}(x, (x + d)).\; (\neg (b = x) \wedge (\lvert (x - b) \rvert \leq r)) \end{array}}}
\vnbprop{ocint-right-limit-point}{\vnbstmt{\begin{array}{@{}l@{}} \forall x, d, r. \\ \quad ((x \in \mathbb{R}) \wedge (\operatorname{pos\text{-}rr}(d) \wedge \operatorname{pos\text{-}rr}(r))) \Rightarrow  \\ \quad \quad \exists b \in \operatorname{ocint}((x - d), x).\; (\neg (b = x) \wedge (\lvert (x - b) \rvert \leq r)) \end{array}}}
\vnbprop{restrict-fun-domain-subset}{\vnbstmt{\begin{array}{@{}l@{}} \forall f, d, c, w, b. \\ \quad ((f \in (d \to c)) \wedge (\operatorname{restrict}(f, w) \in (w \to b))) \Rightarrow (w \subseteq d) \end{array}}}
\vnbprop{deriv-within-in-rr}{\vnbstmt{\forall f, w, x, l.\; (\operatorname{has\text{-}deriv\text{-}at\text{-}within}(f, w, x, l) \Rightarrow (l \in \mathbb{R}))}}
\vnbprop{deriv-within-pt-in}{\vnbstmt{\forall f, w, x, l.\; (\operatorname{has\text{-}deriv\text{-}at\text{-}within}(f, w, x, l) \Rightarrow (x \in w))}}
\vnbprop{deriv-within-restrict-in-fun}{\vnbstmt{\begin{array}{@{}l@{}} \forall f, w, x, l. \\ \quad \operatorname{has\text{-}deriv\text{-}at\text{-}within}(f, w, x, l) \Rightarrow (\operatorname{restrict}(f, w) \in (w \to \mathbb{R})) \end{array}}}
\vnbprop{deriv-within-shrink}{\vnbstmt{\begin{array}{@{}l@{}} \forall f, w, v, x, l. \\ \quad ((w \subseteq \mathbb{R}) \wedge (\operatorname{has\text{-}deriv\text{-}at\text{-}within}(f, w, x, l) \wedge ((v \subseteq w) \wedge (x \in v)))) \Rightarrow  \\ \quad \quad \operatorname{has\text{-}deriv\text{-}at\text{-}within}(f, v, x, l) \end{array}}}
\vnbprop{deriv-within-local}{\vnbstmt{\begin{array}{@{}l@{}} \forall f, g, w, x, l. \\ \quad ((\operatorname{restrict}(f, w) \in (w \to \mathbb{R})) \wedge (\forall y \in w.\; (f(y) \simeq g(y)) \wedge \operatorname{has\text{-}deriv\text{-}at\text{-}within}(g, w, x, l))) \Rightarrow  \\ \quad \quad \operatorname{has\text{-}deriv\text{-}at\text{-}within}(f, w, x, l) \end{array}}}
\vnbprop{deriv-within-unique}{\vnbstmt{\begin{array}{@{}l@{}} \forall f, w, x, l, m. \\ \quad w \subseteq \mathbb{R} \wedge \\ \quad \forall r.\; (\operatorname{pos\text{-}rr}(r) \Rightarrow \exists b \in w.\; (\neg (b = x) \wedge (\lvert (x - b) \rvert \leq r))) \wedge \\ \quad \operatorname{has\text{-}deriv\text{-}at\text{-}within}(f, w, x, l) \wedge \\ \quad \operatorname{has\text{-}deriv\text{-}at\text{-}within}(f, w, x, m) \Rightarrow \\ \quad \quad l = m \end{array}}}
\vnbprop{rr-mul-neg-neg-pos}{\vnbstmt{\forall y, b.\; (((y \in \mathbb{R}) \wedge ((b \in \mathbb{R}) \wedge ((y < 0) \wedge (b < 0)))) \Rightarrow (0 < (y \cdot b)))}}
\vnbprop{has-deriv-at-of-within-ooint}{\vnbstmt{\begin{array}{@{}l@{}} \forall x, r, f, l. \\ \quad (\operatorname{pos\text{-}rr}(r) \wedge \operatorname{has\text{-}deriv\text{-}at\text{-}within}(f, \operatorname{ooint}((x - r), (x + r)), x, l)) \Rightarrow  \\ \quad \quad \operatorname{has\text{-}deriv\text{-}at}(f, x, l) \end{array}}}
\vnbprop{has-deriv-at-to-within-ooint}{\vnbstmt{\begin{array}{@{}l@{}} \forall f, x, l, r. \\ \quad (\operatorname{has\text{-}deriv\text{-}at}(f, x, l) \wedge \operatorname{pos\text{-}rr}(r)) \Rightarrow  \\ \quad \quad \exists v. \\ \quad \quad \quad \operatorname{pos\text{-}rr}(v) \wedge ((v \leq r) \wedge \operatorname{has\text{-}deriv\text{-}at\text{-}within}(f, \operatorname{ooint}((x - v), (x + v)), x, l)) \end{array}}}
\vnbprop{has-deriv-at-iff-within-ooint}{\vnbstmt{\begin{array}{@{}l@{}} \forall f, x, l. \\ \quad \operatorname{has\text{-}deriv\text{-}at}(f, x, l) \\ \quad \Leftrightarrow\; \exists v.\; (\operatorname{pos\text{-}rr}(v) \wedge \operatorname{has\text{-}deriv\text{-}at\text{-}within}(f, \operatorname{ooint}((x - v), (x + v)), x, l)) \end{array}}}
\vnbprop{has-right-deriv-at-intro}{\vnbstmt{\begin{array}{@{}l@{}} \forall f, x, l, r. \\ \quad (\operatorname{pos\text{-}rr}(r) \wedge \operatorname{has\text{-}deriv\text{-}at\text{-}within}(f, \operatorname{coint}(x, (x + r)), x, l)) \Rightarrow  \\ \quad \quad \operatorname{has\text{-}right\text{-}deriv\text{-}at}(f, x, l) \end{array}}}
\vnbprop{has-left-deriv-at-intro}{\vnbstmt{\begin{array}{@{}l@{}} \forall f, x, l, r. \\ \quad (\operatorname{pos\text{-}rr}(r) \wedge \operatorname{has\text{-}deriv\text{-}at\text{-}within}(f, \operatorname{ocint}((x - r), x), x, l)) \Rightarrow  \\ \quad \quad \operatorname{has\text{-}left\text{-}deriv\text{-}at}(f, x, l) \end{array}}}
\vnbprop{has-right-deriv-at-pt-in-rr}{\vnbstmt{\forall f, x, l.\; (\operatorname{has\text{-}right\text{-}deriv\text{-}at}(f, x, l) \Rightarrow (x \in \mathbb{R}))}}
\vnbprop{has-right-deriv-at-in-rr}{\vnbstmt{\forall f, x, l.\; (\operatorname{has\text{-}right\text{-}deriv\text{-}at}(f, x, l) \Rightarrow (l \in \mathbb{R}))}}
\vnbprop{has-left-deriv-at-pt-in-rr}{\vnbstmt{\forall f, x, l.\; (\operatorname{has\text{-}left\text{-}deriv\text{-}at}(f, x, l) \Rightarrow (x \in \mathbb{R}))}}
\vnbprop{has-left-deriv-at-in-rr}{\vnbstmt{\forall f, x, l.\; (\operatorname{has\text{-}left\text{-}deriv\text{-}at}(f, x, l) \Rightarrow (l \in \mathbb{R}))}}
\vnbprop{has-right-deriv-at-shrink}{\vnbstmt{\begin{array}{@{}l@{}} \forall f, x, l, r. \\ \quad (\operatorname{has\text{-}right\text{-}deriv\text{-}at}(f, x, l) \wedge \operatorname{pos\text{-}rr}(r)) \Rightarrow  \\ \quad \quad \exists v. \\ \quad \quad \quad \operatorname{pos\text{-}rr}(v) \wedge ((v \leq r) \wedge \operatorname{has\text{-}deriv\text{-}at\text{-}within}(f, \operatorname{coint}(x, (x + v)), x, l)) \end{array}}}
\vnbprop{has-left-deriv-at-shrink}{\vnbstmt{\begin{array}{@{}l@{}} \forall f, x, l, r. \\ \quad (\operatorname{has\text{-}left\text{-}deriv\text{-}at}(f, x, l) \wedge \operatorname{pos\text{-}rr}(r)) \Rightarrow  \\ \quad \quad \exists v. \\ \quad \quad \quad \operatorname{pos\text{-}rr}(v) \wedge ((v \leq r) \wedge \operatorname{has\text{-}deriv\text{-}at\text{-}within}(f, \operatorname{ocint}((x - v), x), x, l)) \end{array}}}
\vnbprop{has-right-deriv-at-unique}{\vnbstmt{\begin{array}{@{}l@{}} \forall f, x, l, m. \\ \quad (\operatorname{has\text{-}right\text{-}deriv\text{-}at}(f, x, l) \wedge \operatorname{has\text{-}right\text{-}deriv\text{-}at}(f, x, m)) \Rightarrow (l = m) \end{array}}}
\vnbprop{has-left-deriv-at-unique}{\vnbstmt{\begin{array}{@{}l@{}} \forall f, x, l, m. \\ \quad (\operatorname{has\text{-}left\text{-}deriv\text{-}at}(f, x, l) \wedge \operatorname{has\text{-}left\text{-}deriv\text{-}at}(f, x, m)) \Rightarrow (l = m) \end{array}}}
\vnbprop{has-right-deriv-at-local}{\vnbstmt{\begin{array}{@{}l@{}} \forall r, f, g, x, l. \\ \quad \operatorname{pos\text{-}rr}(r) \wedge \\ \quad \operatorname{restrict}(f, \operatorname{coint}(x, (x + r))) \in (\operatorname{coint}(x, (x + r)) \to \mathbb{R}) \wedge \\ \quad \forall y \in \operatorname{coint}(x, (x + r)).\; (f(y) \simeq g(y)) \wedge \\ \quad \operatorname{has\text{-}right\text{-}deriv\text{-}at}(g, x, l) \Rightarrow \\ \quad \quad \operatorname{has\text{-}right\text{-}deriv\text{-}at}(f, x, l) \end{array}}}
\vnbprop{has-left-deriv-at-local}{\vnbstmt{\begin{array}{@{}l@{}} \forall r, f, g, x, l. \\ \quad \operatorname{pos\text{-}rr}(r) \wedge \\ \quad \operatorname{restrict}(f, \operatorname{ocint}((x - r), x)) \in (\operatorname{ocint}((x - r), x) \to \mathbb{R}) \wedge \\ \quad \forall y \in \operatorname{ocint}((x - r), x).\; (f(y) \simeq g(y)) \wedge \\ \quad \operatorname{has\text{-}left\text{-}deriv\text{-}at}(g, x, l) \Rightarrow \\ \quad \quad \operatorname{has\text{-}left\text{-}deriv\text{-}at}(f, x, l) \end{array}}}
\vnbprop{has-deriv-at-implies-right}{\vnbstmt{\forall f, x, l.\; (\operatorname{has\text{-}deriv\text{-}at}(f, x, l) \Rightarrow \operatorname{has\text{-}right\text{-}deriv\text{-}at}(f, x, l))}}
\vnbprop{has-deriv-at-implies-left}{\vnbstmt{\forall f, x, l.\; (\operatorname{has\text{-}deriv\text{-}at}(f, x, l) \Rightarrow \operatorname{has\text{-}left\text{-}deriv\text{-}at}(f, x, l))}}
\vnbprop{local-max-right-deriv-nonpos}{\vnbstmt{\begin{array}{@{}l@{}} \forall f, d, x, l. \\ \quad (\operatorname{is\text{-}local\text{-}max\text{-}at}(f, d, x) \wedge \operatorname{has\text{-}right\text{-}deriv\text{-}at}(f, x, l)) \Rightarrow (l \leq 0) \end{array}}}
\vnbprop{local-max-left-deriv-nonneg}{\vnbstmt{\begin{array}{@{}l@{}} \forall f, d, x, l. \\ \quad (\operatorname{is\text{-}local\text{-}max\text{-}at}(f, d, x) \wedge \operatorname{has\text{-}left\text{-}deriv\text{-}at}(f, x, l)) \Rightarrow (0 \leq l) \end{array}}}
\vnbprop{interval-local-max-right-deriv-nonpos}{\vnbstmt{\begin{array}{@{}l@{}} \forall a, b \in \mathbb{R}, f, x, l. \\ \quad (\operatorname{is\text{-}local\text{-}max\text{-}at}(f, \operatorname{ccint}(a, b), x) \wedge \operatorname{has\text{-}right\text{-}deriv\text{-}at}(f, x, l)) \Rightarrow  \\ \quad \quad l \leq 0 \end{array}}}
\vnbprop{interval-local-max-left-deriv-nonneg}{\vnbstmt{\begin{array}{@{}l@{}} \forall a, b \in \mathbb{R}, f, x, l. \\ \quad (\operatorname{is\text{-}local\text{-}max\text{-}at}(f, \operatorname{ccint}(a, b), x) \wedge \operatorname{has\text{-}left\text{-}deriv\text{-}at}(f, x, l)) \Rightarrow (0 \leq l) \end{array}}}
\vnbprop{rr-max-left-le}{\vnbstmt{\forall y, b.\; (((y \in \mathbb{R}) \wedge ((b \in \mathbb{R}) \wedge (b \leq y))) \Rightarrow (\operatorname{max}(y, b) = y))}}
\vnbprop{rr-min-left-le}{\vnbstmt{\forall y, b.\; (((y \in \mathbb{R}) \wedge ((b \in \mathbb{R}) \wedge (y \leq b))) \Rightarrow (\operatorname{min}(y, b) = y))}}
\vnbprop{rr-max-right-le}{\vnbstmt{\forall y, b.\; (((y \in \mathbb{R}) \wedge ((b \in \mathbb{R}) \wedge (y \leq b))) \Rightarrow (\operatorname{max}(y, b) = b))}}
\vnbprop{rr-min-right-le}{\vnbstmt{\forall y, b.\; (((y \in \mathbb{R}) \wedge ((b \in \mathbb{R}) \wedge (b \leq y))) \Rightarrow (\operatorname{min}(y, b) = b))}}
\vnbprop{rr-max-lt}{\vnbstmt{\begin{array}{@{}l@{}} \forall y, b, c. \\ \quad ((y \in \mathbb{R}) \wedge ((b \in \mathbb{R}) \wedge ((c \in \mathbb{R}) \wedge ((y < c) \wedge (b < c))))) \Rightarrow (\operatorname{max}(y, b) < c) \end{array}}}
\vnbprop{rr-lt-min}{\vnbstmt{\begin{array}{@{}l@{}} \forall c, y, b. \\ \quad ((c \in \mathbb{R}) \wedge ((y \in \mathbb{R}) \wedge ((b \in \mathbb{R}) \wedge ((c < y) \wedge (c < b))))) \Rightarrow (c < \operatorname{min}(y, b)) \end{array}}}
\vnbprop{rr-max-retract-close}{\vnbstmt{\forall y, c.\; (((y \in \mathbb{R}) \wedge (c \in \mathbb{R})) \Rightarrow (\lvert (c - \operatorname{max}(y, c)) \rvert \leq \lvert (c - y) \rvert))}}
\vnbprop{rr-min-retract-close}{\vnbstmt{\forall y, c.\; (((y \in \mathbb{R}) \wedge (c \in \mathbb{R})) \Rightarrow (\lvert (c - \operatorname{min}(y, c)) \rvert \leq \lvert (c - y) \rvert))}}
\vnbprop{coint-of-ooint}{\vnbstmt{\begin{array}{@{}l@{}} \forall x, r.; \forall y \in \operatorname{ooint}((x - r), (x + r)). \\ \quad ((x \in \mathbb{R}) \wedge (\operatorname{pos\text{-}rr}(r) \wedge (x \leq y))) \Rightarrow (y \in \operatorname{coint}(x, (x + r))) \end{array}}}
\vnbprop{ocint-of-ooint}{\vnbstmt{\begin{array}{@{}l@{}} \forall x, r.; \forall y \in \operatorname{ooint}((x - r), (x + r)). \\ \quad ((x \in \mathbb{R}) \wedge (\operatorname{pos\text{-}rr}(r) \wedge (y \leq x))) \Rightarrow (y \in \operatorname{ocint}((x - r), x)) \end{array}}}
\vnbprop{has-deriv-at-of-one-sided}{\vnbstmt{\begin{array}{@{}l@{}} \forall f, x, l. \\ \quad (\operatorname{has\text{-}right\text{-}deriv\text{-}at}(f, x, l) \wedge \operatorname{has\text{-}left\text{-}deriv\text{-}at}(f, x, l)) \Rightarrow  \\ \quad \quad \operatorname{has\text{-}deriv\text{-}at}(f, x, l) \end{array}}}
\vnbprop{has-deriv-at-iff-one-sided}{\vnbstmt{\begin{array}{@{}l@{}} \forall f, x, l. \\ \quad \operatorname{has\text{-}deriv\text{-}at}(f, x, l) \Leftrightarrow (\operatorname{has\text{-}right\text{-}deriv\text{-}at}(f, x, l) \wedge \operatorname{has\text{-}left\text{-}deriv\text{-}at}(f, x, l)) \end{array}}}
\vnbprop{deriv-within-sum}{\vnbstmt{\begin{array}{@{}l@{}} \forall f, g, h, w, x, l, m. \\ \quad w \subseteq \mathbb{R} \wedge \\ \quad \operatorname{restrict}(h, w) \in (w \to \mathbb{R}) \wedge \\ \quad \forall y \in w.\; (h(y) \simeq (f(y) + g(y))) \wedge \\ \quad \operatorname{has\text{-}deriv\text{-}at\text{-}within}(f, w, x, l) \wedge \\ \quad \operatorname{has\text{-}deriv\text{-}at\text{-}within}(g, w, x, m) \Rightarrow \\ \quad \quad \operatorname{has\text{-}deriv\text{-}at\text{-}within}(h, w, x, (l + m)) \end{array}}}
\vnbprop{has-right-deriv-at-sum}{\vnbstmt{\begin{array}{@{}l@{}} \forall r, f, g, h, x, l, m. \\ \quad \operatorname{pos\text{-}rr}(r) \wedge \\ \quad \operatorname{restrict}(h, \operatorname{coint}(x, (x + r))) \in (\operatorname{coint}(x, (x + r)) \to \mathbb{R}) \wedge \\ \quad \forall y \in \operatorname{coint}(x, (x + r)).\; (h(y) \simeq (f(y) + g(y))) \wedge \\ \quad \operatorname{has\text{-}right\text{-}deriv\text{-}at}(f, x, l) \wedge \\ \quad \operatorname{has\text{-}right\text{-}deriv\text{-}at}(g, x, m) \Rightarrow \\ \quad \quad \operatorname{has\text{-}right\text{-}deriv\text{-}at}(h, x, (l + m)) \end{array}}}
\vnbprop{has-left-deriv-at-sum}{\vnbstmt{\begin{array}{@{}l@{}} \forall r, f, g, h, x, l, m. \\ \quad \operatorname{pos\text{-}rr}(r) \wedge \\ \quad \operatorname{restrict}(h, \operatorname{ocint}((x - r), x)) \in (\operatorname{ocint}((x - r), x) \to \mathbb{R}) \wedge \\ \quad \forall y \in \operatorname{ocint}((x - r), x).\; (h(y) \simeq (f(y) + g(y))) \wedge \\ \quad \operatorname{has\text{-}left\text{-}deriv\text{-}at}(f, x, l) \wedge \\ \quad \operatorname{has\text{-}left\text{-}deriv\text{-}at}(g, x, m) \Rightarrow \\ \quad \quad \operatorname{has\text{-}left\text{-}deriv\text{-}at}(h, x, (l + m)) \end{array}}}

\vnbbookfile{structure-library/interval-taylor}{structure-library/interval-taylor}
\noindent{\small interval-taylor.scm -- THE VOCABULARY OF TAYLOR'S FORMULA ON AN INTERVAL. DEFINITIONS ONLY; every theorem is in theorem-library/interval-taylor.scm.}\par
\par\medskip\noindent\textbf{Definition} (\texttt{taylor-sum}).\;$\operatorname{taylor\text{-}sum}(d, n, b, \operatorname{tv\_}) :=$ $\operatorname{series\text{-}partial\text{-}sum}(\operatorname{taylor\text{-}term}(d, b, \operatorname{tv\_}), n)$
\par\medskip\noindent\textbf{Definition} (\texttt{taylor-aux}).\;$\operatorname{taylor\text{-}aux}(d, n, a, b) :=$ $(\lambda \operatorname{tu\_} \in \operatorname{ccint}(a, b).\; \operatorname{taylor\text{-}sum}(d, n, b, \operatorname{tu\_}))$
\par\medskip\noindent\textbf{Definition} (\texttt{taylor-rem}).\;$\operatorname{taylor\text{-}rem}(d, n, a, b) :=$ $(d(0)(b) - \operatorname{taylor\text{-}sum}(d, n, b, a))$
\par\medskip\noindent\textbf{Definition} (\texttt{taylor-term}).\;$\operatorname{taylor\text{-}term}(d, b, \operatorname{tv\_}) :=$
\begin{equation*}
(\lambda \operatorname{tk\_} \in \mathbb{N}.\; ((d(\operatorname{tk\_})(\operatorname{tv\_}) \cdot \operatorname{power}((b - \operatorname{tv\_}), \operatorname{tk\_})) \cdot \operatorname{recip}(\operatorname{factorial}(\operatorname{tk\_}))))
\end{equation*}

\par\medskip\noindent\textbf{Definition} (\texttt{is-taylor-family}).\;For all $d$, $n$, $a$ and $b$, 
\begin{equation*}
\begin{array}{@{}l@{}} \operatorname{is\text{-}taylor\text{-}family}(d, n, a, b) \\ \Leftrightarrow\; n \in \mathbb{N} \wedge \\ 1 \leq n \wedge \\ a \in \mathbb{R} \wedge \\ b \in \mathbb{R} \wedge \\ a < b \wedge \\ d(n) \in (\operatorname{ooint}(a, b) \to \mathbb{R}) \wedge \\ \forall \operatorname{tk\_} \in \mathbb{N}. \\ \quad (\operatorname{tk\_} < n) \Rightarrow  \\ \quad \quad d(\operatorname{tk\_}) \in (\operatorname{ccint}(a, b) \to \mathbb{R}) \wedge \\ \quad \quad \operatorname{is\text{-}continuous\text{-}on}(d(\operatorname{tk\_}), \operatorname{ccint}(a, b)) \wedge \\ \quad \quad \forall \operatorname{tx\_} \in \operatorname{ooint}(a, b).\; \operatorname{has\text{-}deriv\text{-}at}(d(\operatorname{tk\_}), \operatorname{tx\_}, d(\operatorname{succ}(\operatorname{tk\_}))(\operatorname{tx\_})) \end{array}
\end{equation*}


\vnbbookfile{theorem-library/directional-derivative}{theorem-library/directional-derivative}
\noindent{\small directional-derivative.scm -- calculus.pdf Chapter 2 section 8 "Multivariable Calculus", Def 2.23 through Prop 2.27, over an arbitrary real normed vector space.}\par
\par\medskip\noindent\textbf{Definition} (\texttt{seg-curve}).\;$\operatorname{seg\text{-}curve}(m, f, a, \eta) :=$ $(\lambda \operatorname{u\_} \in \mathbb{R}.\; f(\operatorname{vadd}(m)(a, \operatorname{act}(m)(\operatorname{u\_}, \eta))))$
\par\medskip\noindent\textbf{Definition} (\texttt{dir-deriv}).\;$\operatorname{dir\text{-}deriv}(m, f, a, \eta) :=$ $\operatorname{deriv}(\operatorname{seg\text{-}curve}(m, f, a, \eta), 0)$
\par\medskip\noindent\textbf{Definition} (\texttt{is-dir-diff-at}).\;For all $m$, $f$, $a$, $\eta$ and $\mathit{dl}$, 
\begin{equation*}
\operatorname{is\text{-}dir\text{-}diff\text{-}at}(m, f, a, \eta, \mathit{dl}) \Leftrightarrow \operatorname{is\text{-}diff\text{-}at}(\operatorname{seg\text{-}curve}(m, f, a, \eta), 0, \mathit{dl})
\end{equation*}

\vnbprop{seg-curve-apply}{\vnbstmt{\begin{array}{@{}l@{}} \forall m, f, a, b.; \forall t \in \mathbb{R}. \\ \quad (\operatorname{seg\text{-}curve}(m, f, a, b)(t) \simeq f(\operatorname{vadd}(m)(a, \operatorname{act}(m)(t, b)))) \end{array}}}
\vnbprop{diff-at-in-fun}{\vnbstmt{\forall f, a, l.\; (\operatorname{is\text{-}diff\text{-}at}(f, a, l) \Rightarrow (f \in (\mathbb{R} \to \mathbb{R})))}}
\vnbprop{affine-lam-in-fun}{\vnbstmt{\begin{array}{@{}l@{}} \forall c, m. \\ \quad ((c \in \mathbb{R}) \wedge (m \in \mathbb{R})) \Rightarrow ((\lambda x \in \mathbb{R}.\; (c + (m \cdot x))) \in (\mathbb{R} \to \mathbb{R})) \end{array}}}
\vnbprop{deriv-affine}{\vnbstmt{\begin{array}{@{}l@{}} \forall c, m, a. \\ \quad ((c \in \mathbb{R}) \wedge ((m \in \mathbb{R}) \wedge (a \in \mathbb{R}))) \Rightarrow  \\ \quad \quad \operatorname{is\text{-}diff\text{-}at}((\lambda x \in \mathbb{R}.\; (c + (m \cdot x))), a, m) \end{array}}}
\vnbprop{dir-diff-reparam}{\vnbstmt{\begin{array}{@{}l@{}} \forall f, g, c, m, l. \\ \quad f \in (\mathbb{R} \to \mathbb{R}) \wedge \\ \quad c \in \mathbb{R} \wedge \\ \quad m \in \mathbb{R} \wedge \\ \quad \operatorname{is\text{-}diff\text{-}at}(g, c, l) \wedge \\ \quad \forall x \in \mathbb{R}.\; (f(x) \simeq g((c + (m \cdot x)))) \Rightarrow \\ \quad \quad \operatorname{is\text{-}diff\text{-}at}(f, 0, (l \cdot m)) \end{array}}}
\vnbprop{dir-diff-scale}{\vnbstmt{\begin{array}{@{}l@{}} \forall f, g, m, l. \\ \quad ((f \in (\mathbb{R} \to \mathbb{R})) \wedge ((m \in \mathbb{R}) \wedge (\operatorname{is\text{-}diff\text{-}at}(g, 0, l) \wedge \forall x \in \mathbb{R}.\; (f(x) \simeq g((m \cdot x)))))) \Rightarrow  \\ \quad \quad \operatorname{is\text{-}diff\text{-}at}(f, 0, (m \cdot l)) \end{array}}}
\vnbprop{dir-diff-shift}{\vnbstmt{\begin{array}{@{}l@{}} \forall f, g, c, l. \\ \quad ((f \in (\mathbb{R} \to \mathbb{R})) \wedge ((c \in \mathbb{R}) \wedge (\operatorname{is\text{-}diff\text{-}at}(g, c, l) \wedge \forall x \in \mathbb{R}.\; (f(x) \simeq g((c + x)))))) \Rightarrow  \\ \quad \quad \operatorname{is\text{-}diff\text{-}at}(f, 0, l) \end{array}}}
\vnbprop{seg-curve-in-fun}{\vnbstmt{\begin{array}{@{}l@{}} \forall m, f, a, b. \\ \quad (\operatorname{is\text{-}normed\text{-}vector\text{-}space}(m) \wedge ((f \in (\operatorname{vec}(m) \to \mathbb{R})) \wedge ((a \in \operatorname{vec}(m)) \wedge (b \in \operatorname{vec}(m))))) \Rightarrow  \\ \quad \quad \operatorname{seg\text{-}curve}(m, f, a, b) \in (\mathbb{R} \to \mathbb{R}) \end{array}}}
\vnbprop{dir-deriv-homogeneous}{\vnbstmt{\begin{array}{@{}l@{}} \forall m, f, a, b, c, l. \\ \quad \operatorname{is\text{-}normed\text{-}vector\text{-}space}(m) \wedge \\ \quad f \in (\operatorname{vec}(m) \to \mathbb{R}) \wedge \\ \quad a \in \operatorname{vec}(m) \wedge \\ \quad b \in \operatorname{vec}(m) \wedge \\ \quad c \in \mathbb{R} \wedge \\ \quad \operatorname{is\text{-}dir\text{-}diff\text{-}at}(m, f, a, b, l) \Rightarrow \\ \quad \quad \operatorname{is\text{-}dir\text{-}diff\text{-}at}(m, f, a, \operatorname{act}(m)(c, b), (c \cdot l)) \end{array}}}
\vnbprop{dir-deriv-mvt}{\vnbstmt{\begin{array}{@{}l@{}} \forall m, f, a, b. \\ \quad \operatorname{is\text{-}normed\text{-}vector\text{-}space}(m) \wedge \\ \quad f \in (\operatorname{vec}(m) \to \mathbb{R}) \wedge \\ \quad a \in \operatorname{vec}(m) \wedge \\ \quad b \in \operatorname{vec}(m) \wedge \\ \quad \forall x \in \operatorname{ccint}(0, 1). \\ \quad \quad \operatorname{is\text{-}continuous\text{-}at}(\operatorname{rr\text{-}ms}, \operatorname{rr\text{-}ms}, \operatorname{seg\text{-}curve}(m, f, a, b), x) \wedge \\ \quad \forall x \in \mathbb{R}. \\ \quad \quad ((0 < x) \wedge (x < 1)) \Rightarrow \exists l.\; \operatorname{is\text{-}diff\text{-}at}(\operatorname{seg\text{-}curve}(m, f, a, b), x, l) \Rightarrow \\ \quad \quad \exists h \in \mathbb{R}. \\ \quad \quad \quad 0 < h \wedge \\ \quad \quad \quad h < 1 \wedge \\ \quad \quad \quad \exists v. \\ \quad \quad \quad \quad (\operatorname{is\text{-}dir\text{-}diff\text{-}at}(m, f, \operatorname{vadd}(m)(a, \operatorname{act}(m)(h, b)), b, v) \wedge (v = (f(\operatorname{vadd}(m)(a, b)) - f(a)))) \end{array}}}
\vnbprop{dir-deriv-value}{\vnbstmt{\begin{array}{@{}l@{}} \forall m, f, a, b, l. \\ \quad \operatorname{is\text{-}dir\text{-}diff\text{-}at}(m, f, a, b, l) \Rightarrow (\operatorname{dir\text{-}deriv}(m, f, a, b) = l) \end{array}}}
\vnbprop{diff-at-shift-back}{\vnbstmt{\begin{array}{@{}l@{}} \forall f, g, c, l. \\ \quad ((f \in (\mathbb{R} \to \mathbb{R})) \wedge ((c \in \mathbb{R}) \wedge (\operatorname{is\text{-}diff\text{-}at}(g, 0, l) \wedge \forall x \in \mathbb{R}.\; (g(x) \simeq f((c + x)))))) \Rightarrow  \\ \quad \quad \operatorname{is\text{-}diff\text{-}at}(f, c, l) \end{array}}}
\vnbprop{seg-curve-diff-at}{\vnbstmt{\begin{array}{@{}l@{}} \forall m.; \forall f \in (\operatorname{vec}(m) \to \mathbb{R}), b, a \in \operatorname{vec}(m), s \in \mathbb{R}, k. \\ \quad (\operatorname{is\text{-}normed\text{-}vector\text{-}space}(m) \wedge \operatorname{is\text{-}dir\text{-}diff\text{-}at}(m, f, \operatorname{vadd}(m)(b, \operatorname{act}(m)(s, a)), a, k)) \Rightarrow  \\ \quad \quad \operatorname{is\text{-}diff\text{-}at}(\operatorname{seg\text{-}curve}(m, f, b, a), s, k) \end{array}}}
\vnbprop{dir-deriv-increment}{\vnbstmt{\begin{array}{@{}l@{}} \forall m.; \forall f \in (\operatorname{vec}(m) \to \mathbb{R}), b, a \in \operatorname{vec}(m). \\ \quad (\operatorname{is\text{-}normed\text{-}vector\text{-}space}(m) \wedge \forall x \in \operatorname{vec}(m).\; \exists k.\; \operatorname{is\text{-}dir\text{-}diff\text{-}at}(m, f, x, a, k)) \Rightarrow  \\ \quad \quad \exists h \in \mathbb{R}. \\ \quad \quad \quad ((0 < h) \wedge ((h < 1) \wedge ((f(\operatorname{vadd}(m)(b, a)) - f(b)) = \operatorname{dir\text{-}deriv}(m, f, \operatorname{vadd}(m)(b, \operatorname{act}(m)(h, a)), a)))) \end{array}}}
\vnbprop{dir-diff-value-in-rr}{\vnbstmt{\forall m, f, a, b, l.\; (\operatorname{is\text{-}dir\text{-}diff\text{-}at}(m, f, a, b, l) \Rightarrow (l \in \mathbb{R}))}}
\vnbprop{dir-deriv-increment-scaled}{\vnbstmt{\begin{array}{@{}l@{}} \forall m.; \forall f \in (\operatorname{vec}(m) \to \mathbb{R}), b, a \in \operatorname{vec}(m), t \in \mathbb{R}. \\ \quad (\operatorname{is\text{-}normed\text{-}vector\text{-}space}(m) \wedge \forall x \in \operatorname{vec}(m).\; \exists k.\; \operatorname{is\text{-}dir\text{-}diff\text{-}at}(m, f, x, a, k)) \Rightarrow  \\ \quad \quad \exists h \in \mathbb{R}. \\ \quad \quad \quad 0 < h \wedge \\ \quad \quad \quad h < 1 \wedge \\ \quad \quad \quad (f(\operatorname{vadd}(m)(b, \operatorname{act}(m)(t, a))) - f(b)) = (t \cdot \operatorname{dir\text{-}deriv}(m, f, \operatorname{vadd}(m)(b, \operatorname{act}(m)((h \cdot t), a)), a)) \end{array}}}
\vnbprop{rr-abs-unit-scale}{\vnbstmt{\forall s, t \in \mathbb{R}.\; (((0 < s) \wedge (s < 1)) \Rightarrow (\lvert (s \cdot t) \rvert \leq \lvert t \rvert))}}
\vnbprop{nvs-norm-two-term-bound}{\vnbstmt{\begin{array}{@{}l@{}} \forall m.; \forall a, b \in \operatorname{vec}(m), c, d, f \in \mathbb{R}. \\ \quad (\operatorname{is\text{-}normed\text{-}vector\text{-}space}(m) \wedge ((0 \leq f) \wedge ((\lvert c \rvert \leq f) \wedge (\lvert d \rvert \leq f)))) \Rightarrow  \\ \quad \quad \operatorname{vnrm}(m)(\operatorname{vadd}(m)(\operatorname{act}(m)(c, a), \operatorname{act}(m)(d, b))) \leq (f \cdot (1 + (\operatorname{vnrm}(m)(b) + \operatorname{vnrm}(m)(a)))) \end{array}}}
\vnbprop{nvs-norm-one-term-bound}{\vnbstmt{\begin{array}{@{}l@{}} \forall m.; \forall a, b \in \operatorname{vec}(m), c, d \in \mathbb{R}. \\ \quad (\operatorname{is\text{-}normed\text{-}vector\text{-}space}(m) \wedge ((0 \leq d) \wedge (\lvert c \rvert \leq d))) \Rightarrow  \\ \quad \quad \operatorname{vnrm}(m)(\operatorname{act}(m)(c, a)) \leq (d \cdot (1 + (\operatorname{vnrm}(m)(b) + \operatorname{vnrm}(m)(a)))) \end{array}}}
\vnbprop{rr-diff-split-mid}{\vnbstmt{\forall x, y, z \in \mathbb{R}.\; ((x - z) = ((x - y) + (y - z)))}}
\vnbprop{dir-deriv-corner-split}{\vnbstmt{\begin{array}{@{}l@{}} \forall m.; \forall f \in (\operatorname{vec}(m) \to \mathbb{R}), a, b, i \in \operatorname{vec}(m), t \in \mathbb{R}. \\ \quad \operatorname{is\text{-}normed\text{-}vector\text{-}space}(m) \wedge \\ \quad \forall x \in \operatorname{vec}(m).\; \exists k.\; \operatorname{is\text{-}dir\text{-}diff\text{-}at}(m, f, x, b, k) \wedge \\ \quad \forall x \in \operatorname{vec}(m).\; \exists k.\; \operatorname{is\text{-}dir\text{-}diff\text{-}at}(m, f, x, i, k) \Rightarrow \\ \quad \quad \exists c, d \in \operatorname{vec}(m). \\ \quad \quad \quad \operatorname{dist}(\operatorname{nvs\text{-}metric\text{-}space}(m))(a, c) \leq (\lvert t \rvert \cdot (1 + (\operatorname{vnrm}(m)(b) + \operatorname{vnrm}(m)(i)))) \wedge \\ \quad \quad \quad \operatorname{dist}(\operatorname{nvs\text{-}metric\text{-}space}(m))(a, d) \leq (\lvert t \rvert \cdot (1 + (\operatorname{vnrm}(m)(b) + \operatorname{vnrm}(m)(i)))) \wedge \\ \quad \quad \quad (\operatorname{seg\text{-}curve}(m, f, a, \operatorname{vadd}(m)(b, i))(t) - \operatorname{seg\text{-}curve}(m, f, a, \operatorname{vadd}(m)(b, i))(0)) \\ \quad \quad \quad =\; t \cdot (\operatorname{dir\text{-}deriv}(m, f, c, b) + \operatorname{dir\text{-}deriv}(m, f, d, i)) \end{array}}}
\vnbprop{rr-factor-shift}{\vnbstmt{\forall d, u, r \in \mathbb{R}.\; (((d \cdot r) \cdot u) = (d \cdot (u \cdot r)))}}
\vnbprop{dir-quotient-in-fun}{\vnbstmt{\begin{array}{@{}l@{}} \forall m.; \forall f \in (\operatorname{vec}(m) \to \mathbb{R}), a, b, i \in \operatorname{vec}(m), l, c \in \mathbb{R}. \\ \quad \operatorname{is\text{-}normed\text{-}vector\text{-}space}(m) \Rightarrow  \\ \quad \quad (\lambda z \in \mathbb{R}.\; \\ \quad \quad \quad \left(\text{if }(z = 0)\text{ then }(l + c)\text{ else }((\operatorname{seg\text{-}curve}(m, f, a, \operatorname{vadd}(m)(b, i))(z) - \operatorname{seg\text{-}curve}(m, f, a, \operatorname{vadd}(m)(b, i))(0)) \cdot \operatorname{recip}(z))\right)) \\ \quad \quad \in\; (\mathbb{R} \to \mathbb{R}) \end{array}}}
\vnbprop{dir-quotient-at-zero}{\vnbstmt{\begin{array}{@{}l@{}} \forall m.; \forall f \in (\operatorname{vec}(m) \to \mathbb{R}), a, b, i \in \operatorname{vec}(m), l, c \in \mathbb{R}. \\ \quad \operatorname{is\text{-}normed\text{-}vector\text{-}space}(m) \Rightarrow  \\ \quad \quad (\lambda z \in \mathbb{R}.\; \left(\text{if }(z = 0)\text{ then }(l + c)\text{ else }((\operatorname{seg\text{-}curve}(m, f, a, \operatorname{vadd}(m)(b, i))(z) - \operatorname{seg\text{-}curve}(m, f, a, \operatorname{vadd}(m)(b, i))(0)) \cdot \operatorname{recip}(z))\right))(0) \\ \quad \quad =\; l + c \end{array}}}
\vnbprop{dir-quotient-identity}{\vnbstmt{\begin{array}{@{}l@{}} \forall m.; \forall f \in (\operatorname{vec}(m) \to \mathbb{R}), a, b, i \in \operatorname{vec}(m), l, c, x \in \mathbb{R}. \\ \quad \operatorname{is\text{-}normed\text{-}vector\text{-}space}(m) \Rightarrow  \\ \quad \quad (\operatorname{seg\text{-}curve}(m, f, a, \operatorname{vadd}(m)(b, i))(x) - \operatorname{seg\text{-}curve}(m, f, a, \operatorname{vadd}(m)(b, i))(0)) \\ \quad \quad =\; (\lambda z \in \mathbb{R}.\; \left(\text{if }(z = 0)\text{ then }(l + c)\text{ else }((\operatorname{seg\text{-}curve}(m, f, a, \operatorname{vadd}(m)(b, i))(z) - \operatorname{seg\text{-}curve}(m, f, a, \operatorname{vadd}(m)(b, i))(0)) \cdot \operatorname{recip}(z))\right))(x) \cdot x \end{array}}}
\vnbprop{rr-mul-rotate}{\vnbstmt{\forall s, u, r \in \mathbb{R}.\; (((u \cdot s) \cdot r) = (s \cdot (u \cdot r)))}}
\vnbprop{dir-quotient-continuous}{\vnbstmt{\begin{array}{@{}l@{}} \forall m.; \forall f \in (\operatorname{vec}(m) \to \mathbb{R}), a, b, i \in \operatorname{vec}(m), l, c. \\ \quad \operatorname{is\text{-}normed\text{-}vector\text{-}space}(m) \wedge \\ \quad \operatorname{is\text{-}dir\text{-}diff\text{-}at}(m, f, a, b, l) \wedge \\ \quad \operatorname{is\text{-}dir\text{-}diff\text{-}at}(m, f, a, i, c) \wedge \\ \quad \forall x \in \operatorname{vec}(m).\; \exists k.\; \operatorname{is\text{-}dir\text{-}diff\text{-}at}(m, f, x, b, k) \wedge \\ \quad \forall x \in \operatorname{vec}(m).\; \exists k.\; \operatorname{is\text{-}dir\text{-}diff\text{-}at}(m, f, x, i, k) \wedge \\ \quad \operatorname{is\text{-}continuous\text{-}at}(\operatorname{nvs\text{-}metric\text{-}space}(m), \operatorname{rr\text{-}ms}, (\lambda x \in \operatorname{vec}(m).\; \operatorname{dir\text{-}deriv}(m, f, x, b)), a) \wedge \\ \quad \operatorname{is\text{-}continuous\text{-}at}(\operatorname{nvs\text{-}metric\text{-}space}(m), \operatorname{rr\text{-}ms}, (\lambda x \in \operatorname{vec}(m).\; \operatorname{dir\text{-}deriv}(m, f, x, i)), a) \Rightarrow \\ \quad \quad \operatorname{is\text{-}continuous\text{-}at}(\operatorname{rr\text{-}ms}, \operatorname{rr\text{-}ms}, \\ \quad \quad \quad (\lambda z \in \mathbb{R}.\; \\ \quad \quad \quad \quad \left(\text{if }(z = 0)\text{ then }(l + c)\text{ else }((\operatorname{seg\text{-}curve}(m, f, a, \operatorname{vadd}(m)(b, i))(z) - \operatorname{seg\text{-}curve}(m, f, a, \operatorname{vadd}(m)(b, i))(0)) \cdot \operatorname{recip}(z))\right)), \\ \quad \quad \quad 0) \end{array}}}
\vnbprop{dir-deriv-additive}{\vnbstmt{\begin{array}{@{}l@{}} \forall m, f, a, b, i, l, c. \\ \quad \operatorname{is\text{-}normed\text{-}vector\text{-}space}(m) \wedge \\ \quad f \in (\operatorname{vec}(m) \to \mathbb{R}) \wedge \\ \quad a \in \operatorname{vec}(m) \wedge \\ \quad b \in \operatorname{vec}(m) \wedge \\ \quad i \in \operatorname{vec}(m) \wedge \\ \quad \operatorname{is\text{-}dir\text{-}diff\text{-}at}(m, f, a, b, l) \wedge \\ \quad \operatorname{is\text{-}dir\text{-}diff\text{-}at}(m, f, a, i, c) \wedge \\ \quad \forall x \in \operatorname{vec}(m).\; \exists k.\; \operatorname{is\text{-}dir\text{-}diff\text{-}at}(m, f, x, b, k) \wedge \\ \quad \forall x \in \operatorname{vec}(m).\; \exists k.\; \operatorname{is\text{-}dir\text{-}diff\text{-}at}(m, f, x, i, k) \wedge \\ \quad \operatorname{is\text{-}continuous\text{-}at}(\operatorname{nvs\text{-}metric\text{-}space}(m), \operatorname{rr\text{-}ms}, (\lambda x \in \operatorname{vec}(m).\; \operatorname{dir\text{-}deriv}(m, f, x, b)), a) \wedge \\ \quad \operatorname{is\text{-}continuous\text{-}at}(\operatorname{nvs\text{-}metric\text{-}space}(m), \operatorname{rr\text{-}ms}, (\lambda x \in \operatorname{vec}(m).\; \operatorname{dir\text{-}deriv}(m, f, x, i)), a) \Rightarrow \\ \quad \quad \operatorname{is\text{-}dir\text{-}diff\text{-}at}(m, f, a, \operatorname{vadd}(m)(b, i), (l + c)) \end{array}}}
\vnbprop{dir-deriv-linear}{\vnbstmt{\begin{array}{@{}l@{}} \forall m, f, a, b, i, r, s, l, c. \\ \quad \operatorname{is\text{-}normed\text{-}vector\text{-}space}(m) \wedge \\ \quad f \in (\operatorname{vec}(m) \to \mathbb{R}) \wedge \\ \quad a \in \operatorname{vec}(m) \wedge \\ \quad b \in \operatorname{vec}(m) \wedge \\ \quad i \in \operatorname{vec}(m) \wedge \\ \quad r \in \mathbb{R} \wedge \\ \quad s \in \mathbb{R} \wedge \\ \quad \operatorname{is\text{-}dir\text{-}diff\text{-}at}(m, f, a, b, l) \wedge \\ \quad \operatorname{is\text{-}dir\text{-}diff\text{-}at}(m, f, a, i, c) \wedge \\ \quad \forall x \in \operatorname{vec}(m).\; \exists k.\; \operatorname{is\text{-}dir\text{-}diff\text{-}at}(m, f, x, \operatorname{act}(m)(r, b), k) \wedge \\ \quad \forall x \in \operatorname{vec}(m).\; \exists k.\; \operatorname{is\text{-}dir\text{-}diff\text{-}at}(m, f, x, \operatorname{act}(m)(s, i), k) \wedge \\ \quad \operatorname{is\text{-}continuous\text{-}at}(\operatorname{nvs\text{-}metric\text{-}space}(m), \operatorname{rr\text{-}ms}, \\ \quad \quad (\lambda x \in \operatorname{vec}(m).\; \operatorname{dir\text{-}deriv}(m, f, x, \operatorname{act}(m)(r, b))), \\ \quad \quad a) \wedge \\ \quad \operatorname{is\text{-}continuous\text{-}at}(\operatorname{nvs\text{-}metric\text{-}space}(m), \operatorname{rr\text{-}ms}, \\ \quad \quad (\lambda x \in \operatorname{vec}(m).\; \operatorname{dir\text{-}deriv}(m, f, x, \operatorname{act}(m)(s, i))), \\ \quad \quad a) \Rightarrow \\ \quad \quad \operatorname{is\text{-}dir\text{-}diff\text{-}at}(m, f, a, \operatorname{vadd}(m)(\operatorname{act}(m)(r, b), \operatorname{act}(m)(s, i)), ((r \cdot l) + (s \cdot c))) \end{array}}}

\vnbbookfile{theorem-library/nvs-taylor-proof}{theorem-library/nvs-taylor-proof}
\noindent{\small nvs-taylor-proof.scm -- Taylor's theorem with remainder bound for a map between two finite-dimensional real normed vector spaces.  PROVED.}\par
\vnbprop{rr-power-one}{\vnbstmt{\forall k \in \mathbb{N}.\; (\operatorname{power}(1, k) = 1)}}
\vnbprop{nt-unit-factor}{\vnbstmt{\forall y, k.\; (((y \in \mathbb{R}) \wedge (k \in \mathbb{N})) \Rightarrow ((y \cdot \operatorname{power}(1, (k + 1))) = y))}}
\vnbprop{nvs-taylor-remainder-bound}{\vnbstmt{\begin{array}{@{}l@{}} \forall m, b, f, a, h, n. \\ \quad \operatorname{is\text{-}normed\text{-}vector\text{-}space}(m) \wedge \\ \quad \operatorname{is\text{-}finite\text{-}dimensional}(\operatorname{normed\text{-}vector\text{-}space\text{-}as\text{-}module}(m)) \wedge \\ \quad \operatorname{is\text{-}normed\text{-}vector\text{-}space}(b) \wedge \\ \quad \operatorname{is\text{-}finite\text{-}dimensional}(\operatorname{normed\text{-}vector\text{-}space\text{-}as\text{-}module}(b)) \wedge \\ \quad f \in (\operatorname{vec}(m) \to \operatorname{vec}(b)) \wedge \\ \quad a \in \operatorname{vec}(m) \wedge \\ \quad h \in \operatorname{vec}(m) \wedge \\ \quad n \in \mathbb{N} \wedge \\ \quad \operatorname{taylor\text{-}differentiable\text{-}v}(b, (\lambda t \in \mathbb{R}.\; f(\operatorname{vadd}(m)(a, \operatorname{act}(m)(t, h)))), 0, 1, n) \Rightarrow \\ \quad \quad \exists c. \\ \quad \quad \quad 0 < c \wedge \\ \quad \quad \quad c < 1 \wedge \\ \quad \quad \quad \operatorname{factorial}((n + 1)) \cdot \operatorname{vnrm}(b)(\operatorname{vadd}(b)(f(\operatorname{vadd}(m)(a, h)), \operatorname{vneg}(b)(\operatorname{taylor\text{-}poly\text{-}v}(b, (\lambda t \in \mathbb{R}.\; f(\operatorname{vadd}(m)(a, \operatorname{act}(m)(t, h)))), 0, 1, n)))) \\ \quad \quad \quad \leq\; \operatorname{vnrm}(b)(\operatorname{nth\text{-}deriv\text{-}v}(b, (\lambda t \in \mathbb{R}.\; f(\operatorname{vadd}(m)(a, \operatorname{act}(m)(t, h)))), (n + 1))(c)) \end{array}}}

\vnbbookfile{theorem-library/antiderivative}{theorem-library/antiderivative}
\noindent{\small antiderivative.scm -- docs/calculus.pdf CHAPTER 4 SECTION 2: Definition 4.6 (antiderivative / antiderivable), Proposition 4.8 (closed under sums; the scalar half is handed off), Proposition 4.10 (two antiderivatives of the same function differ by a constant) and Corollary 4.11 (so f(b) - f(a) does}\par
\vnbprop{antiderivative-map-in-fun}{\vnbstmt{\forall f, i, a, b.\; (\operatorname{is\text{-}antiderivative}(f, i, a, b) \Rightarrow (f \in (\mathbb{R} \to \mathbb{R})))}}
\vnbprop{antiderivative-fn-in-fun}{\vnbstmt{\forall f, i, a, b.\; (\operatorname{is\text{-}antiderivative}(f, i, a, b) \Rightarrow (i \in (\mathbb{R} \to \mathbb{R})))}}
\vnbprop{antiderivative-endpoints}{\vnbstmt{\begin{array}{@{}l@{}} \forall f, i, a, b. \\ \quad \operatorname{is\text{-}antiderivative}(f, i, a, b) \Rightarrow ((a \in \mathbb{R}) \wedge ((b \in \mathbb{R}) \wedge (a < b))) \end{array}}}
\vnbprop{deriv-sub}{\vnbstmt{\begin{array}{@{}l@{}} \forall f, g, a, l, m. \\ \quad (\operatorname{is\text{-}diff\text{-}at}(f, a, l) \wedge \operatorname{is\text{-}diff\text{-}at}(g, a, m)) \Rightarrow  \\ \quad \quad \operatorname{is\text{-}diff\text{-}at}((\lambda x \in \mathbb{R}.\; (f(x) - g(x))), a, (l - m)) \end{array}}}
\vnbprop{antiderivative-differ-by-constant}{\vnbstmt{\begin{array}{@{}l@{}} \forall f, g, i, a, b. \\ \quad (\operatorname{is\text{-}antiderivative}(f, i, a, b) \wedge \operatorname{is\text{-}antiderivative}(g, i, a, b)) \Rightarrow  \\ \quad \quad \exists t \in \mathbb{R}.\; \forall s \in \operatorname{ccint}(a, b).\; (f(s) = (g(s) + t)) \end{array}}}
\vnbprop{antiderivative-endpoint-difference}{\vnbstmt{\begin{array}{@{}l@{}} \forall f, g, i, a, b. \\ \quad (\operatorname{is\text{-}antiderivative}(f, i, a, b) \wedge \operatorname{is\text{-}antiderivative}(g, i, a, b)) \Rightarrow  \\ \quad \quad (f(b) - f(a)) = (g(b) - g(a)) \end{array}}}
\vnbprop{poly-is-antiderivable}{\vnbstmt{\begin{array}{@{}l@{}} \forall n \in \mathbb{N}, f \in (\mathbb{N} \to \mathbb{R}), a, b. \\ \quad ((a \in \mathbb{R}) \wedge ((b \in \mathbb{R}) \wedge (a < b))) \Rightarrow  \\ \quad \quad \operatorname{is\text{-}antiderivable}((\lambda x \in \mathbb{R}.\; \operatorname{series\text{-}partial\text{-}sum}((\lambda k \in \mathbb{N}.\; (f(k) \cdot \operatorname{power}(x, k))), (n + 1))), \\ \quad \quad \quad a, b) \end{array}}}
\vnbprop{antiderivative-add}{\vnbstmt{\begin{array}{@{}l@{}} \forall f, g, i, c, a, b. \\ \quad (\operatorname{is\text{-}antiderivative}(f, i, a, b) \wedge \operatorname{is\text{-}antiderivative}(g, c, a, b)) \Rightarrow  \\ \quad \quad \operatorname{is\text{-}antiderivative}((\lambda x \in \mathbb{R}.\; (f(x) + g(x))), (\lambda x \in \mathbb{R}.\; (i(x) + c(x))), a, b) \end{array}}}
\vnbprop{antiderivable-add}{\vnbstmt{\begin{array}{@{}l@{}} \forall i, c, a, b. \\ \quad (\operatorname{is\text{-}antiderivable}(i, a, b) \wedge \operatorname{is\text{-}antiderivable}(c, a, b)) \Rightarrow  \\ \quad \quad \operatorname{is\text{-}antiderivable}((\lambda x \in \mathbb{R}.\; (i(x) + c(x))), a, b) \end{array}}}
\vnbprop{antiderivative-scale}{\vnbstmt{\begin{array}{@{}l@{}} \forall c, f, i, a, b. \\ \quad ((c \in \mathbb{R}) \wedge \operatorname{is\text{-}antiderivative}(f, i, a, b)) \Rightarrow  \\ \quad \quad \operatorname{is\text{-}antiderivative}((\lambda x \in \mathbb{R}.\; (c \cdot f(x))), (\lambda x \in \mathbb{R}.\; (c \cdot i(x))), a, b) \end{array}}}
\vnbprop{antiderivable-scale}{\vnbstmt{\begin{array}{@{}l@{}} \forall c, i, a, b. \\ \quad ((c \in \mathbb{R}) \wedge \operatorname{is\text{-}antiderivable}(i, a, b)) \Rightarrow  \\ \quad \quad \operatorname{is\text{-}antiderivable}((\lambda x \in \mathbb{R}.\; (c \cdot i(x))), a, b) \end{array}}}

\vnbbookfile{theorem-library/antiderivative-transfer}{theorem-library/antiderivative-transfer}
\noindent{\small antiderivative-transfer.scm -- the POINTWISE TRANSFER for Def 4.6, and the DIFFERENCE of two antiderivatives.  All three `modulo 0'.}\par
\vnbprop{antiderivative-transfer-ptwise-eq}{\vnbstmt{\begin{array}{@{}l@{}} \forall f, i \in (\mathbb{R} \to \mathbb{R}), g, c, a, b. \\ \quad (\operatorname{is\text{-}antiderivative}(g, c, a, b) \wedge (\forall x \in \mathbb{R}.\; (f(x) = g(x)) \wedge \forall x \in \mathbb{R}.\; (i(x) = c(x)))) \Rightarrow  \\ \quad \quad \operatorname{is\text{-}antiderivative}(f, i, a, b) \end{array}}}
\vnbprop{antiderivative-sub}{\vnbstmt{\begin{array}{@{}l@{}} \forall f, g, i, c, a, b. \\ \quad (\operatorname{is\text{-}antiderivative}(f, i, a, b) \wedge \operatorname{is\text{-}antiderivative}(g, c, a, b)) \Rightarrow  \\ \quad \quad \operatorname{is\text{-}antiderivative}((\lambda x \in \mathbb{R}.\; (f(x) - g(x))), (\lambda x \in \mathbb{R}.\; (i(x) - c(x))), a, b) \end{array}}}
\vnbprop{antiderivable-sub}{\vnbstmt{\begin{array}{@{}l@{}} \forall i, c, a, b. \\ \quad (\operatorname{is\text{-}antiderivable}(i, a, b) \wedge \operatorname{is\text{-}antiderivable}(c, a, b)) \Rightarrow  \\ \quad \quad \operatorname{is\text{-}antiderivable}((\lambda x \in \mathbb{R}.\; (i(x) - c(x))), a, b) \end{array}}}

\vnbbookfile{theorem-library/antiderivative-lipschitz}{theorem-library/antiderivative-lipschitz}
\noindent{\small antiderivative-lipschitz.scm -- THE WITNESS ESTIMATE of docs/calculus.pdf Prop 4.16: a bound on the INTEGRAND bounds the CARATHEODORY WITNESS, at every point of the interval and by the same constant.}\par
\vnbprop{antiderivative-increment-bound}{\vnbstmt{\begin{array}{@{}l@{}} \forall h, i \in (\mathbb{R} \to \mathbb{R}), a, b, m, u, v \in \mathbb{R}. \\ \quad \operatorname{is\text{-}antiderivative}(h, i, a, b) \wedge \\ \quad \forall x \in \mathbb{R}.\; ((x \in \operatorname{ccint}(a, b)) \Rightarrow (\lvert i(x) \rvert \leq m)) \wedge \\ \quad u \in \operatorname{ccint}(a, b) \wedge \\ \quad v \in \operatorname{ccint}(a, b) \wedge \\ \quad u < v \Rightarrow \\ \quad \quad \lvert (h(v) - h(u)) \rvert \leq (m \cdot (v - u)) \end{array}}}
\vnbprop{antiderivative-lipschitz}{\vnbstmt{\begin{array}{@{}l@{}} \forall h, i \in (\mathbb{R} \to \mathbb{R}), a, b, m, u, v \in \mathbb{R}. \\ \quad \operatorname{is\text{-}antiderivative}(h, i, a, b) \wedge \\ \quad \forall x \in \mathbb{R}.\; ((x \in \operatorname{ccint}(a, b)) \Rightarrow (\lvert i(x) \rvert \leq m)) \wedge \\ \quad u \in \operatorname{ccint}(a, b) \wedge \\ \quad v \in \operatorname{ccint}(a, b) \Rightarrow \\ \quad \quad \lvert (h(v) - h(u)) \rvert \leq (m \cdot \lvert (v - u) \rvert) \end{array}}}
\vnbprop{caratheodory-witness-bound}{\vnbstmt{\begin{array}{@{}l@{}} \forall h, i, c \in (\mathbb{R} \to \mathbb{R}), a, b, t, m, x \in \mathbb{R}. \\ \quad \operatorname{is\text{-}antiderivative}(h, i, a, b) \wedge \\ \quad \forall x \in \mathbb{R}.\; ((x \in \operatorname{ccint}(a, b)) \Rightarrow (\lvert i(x) \rvert \leq m)) \wedge \\ \quad t \in \operatorname{ccint}(a, b) \wedge \\ \quad c(t) = i(t) \wedge \\ \quad \forall x \in \mathbb{R}.\; ((h(x) - h(t)) = (c(x) \cdot (x - t))) \wedge \\ \quad x \in \operatorname{ccint}(a, b) \Rightarrow \\ \quad \quad \lvert c(x) \rvert \leq m \end{array}}}
\vnbprop{antiderivative-pair-close}{\vnbstmt{\begin{array}{@{}l@{}} \forall f, g, i, c \in (\mathbb{R} \to \mathbb{R}), a, b, m, x \in \mathbb{R}. \\ \quad \operatorname{is\text{-}antiderivative}(f, i, a, b) \wedge \\ \quad \operatorname{is\text{-}antiderivative}(g, c, a, b) \wedge \\ \quad f(a) = g(a) \wedge \\ \quad \forall x \in \mathbb{R}.\; ((x \in \operatorname{ccint}(a, b)) \Rightarrow (\lvert (i(x) - c(x)) \rvert \leq m)) \wedge \\ \quad x \in \operatorname{ccint}(a, b) \Rightarrow \\ \quad \quad \lvert (f(x) - g(x)) \rvert \leq (m \cdot (b - a)) \end{array}}}

\vnbbookfile{theorem-library/witness-family-choice}{theorem-library/witness-family-choice}
\noindent{\small witness-family-choice.scm -- TWO FAMILIES BUILT BY CHOICE, and nothing else.}\par
\vnbprop{antiderivable-family-choice}{\vnbstmt{\begin{array}{@{}l@{}} \forall m \in (\mathbb{N} \to (\mathbb{R} \to \mathbb{R})), a, b \in \mathbb{R}. \\ \quad \forall k \in \mathbb{N}.\; \operatorname{is\text{-}antiderivable}(m(k), a, b) \Rightarrow  \\ \quad \quad \exists c \in (\mathbb{N} \to (\mathbb{R} \to \mathbb{R})). \forall k \in \mathbb{N}. \\ \quad \quad \quad \operatorname{is\text{-}antiderivative}(c(k), m(k), a, b) \end{array}}}
\vnbprop{diff-at-witness-family}{\vnbstmt{\begin{array}{@{}l@{}} \forall m, a, t. \\ \quad m \in (\mathbb{N} \to (\mathbb{R} \to \mathbb{R})) \wedge \\ \quad a \in (\mathbb{N} \to \mathbb{R}) \wedge \\ \quad t \in \mathbb{R} \wedge \\ \quad \forall k \in \mathbb{N}.\; \operatorname{is\text{-}diff\text{-}at}(m(k), t, a(k)) \Rightarrow \\ \quad \quad \exists b \in (\mathbb{N} \to (\mathbb{R} \to \mathbb{R})). \forall k \in \mathbb{N}. \\ \quad \quad \quad \operatorname{is\text{-}continuous\text{-}at}(\operatorname{rr\text{-}ms}, \operatorname{rr\text{-}ms}, b(k), t) \wedge \\ \quad \quad \quad b(k)(t) = a(k) \wedge \\ \quad \quad \quad \forall x \in \mathbb{R}.\; ((m(k)(x) - m(k)(t)) = (b(k)(x) \cdot (x - t))) \end{array}}}

\vnbbookfile{theorem-library/antiderivative-normalize}{theorem-library/antiderivative-normalize}
\noindent{\small antiderivative-normalize.scm -- SHIFTING AN ANTIDERIVATIVE BY A CONSTANT, and the NORMALISED antiderivative family that docs/calculus.pdf Cor 4.17 needs and that nothing in the tree supplied.}\par
\vnbprop{antiderivative-shift-const}{\vnbstmt{\begin{array}{@{}l@{}} \forall f, i, a, b, c. \\ \quad f \in (\mathbb{R} \to \mathbb{R}) \wedge \\ \quad i \in (\mathbb{R} \to \mathbb{R}) \wedge \\ \quad a \in \mathbb{R} \wedge \\ \quad b \in \mathbb{R} \wedge \\ \quad c \in \mathbb{R} \wedge \\ \quad \operatorname{is\text{-}antiderivative}(f, i, a, b) \Rightarrow \\ \quad \quad \operatorname{is\text{-}antiderivative}((\lambda x \in \mathbb{R}.\; (f(x) - c)), i, a, b) \end{array}}}
\vnbprop{antiderivable-family-normalized}{\vnbstmt{\begin{array}{@{}l@{}} \forall m \in (\mathbb{N} \to (\mathbb{R} \to \mathbb{R})), a, b \in \mathbb{R}. \\ \quad \forall k \in \mathbb{N}.\; \operatorname{is\text{-}antiderivable}(m(k), a, b) \Rightarrow  \\ \quad \quad \exists c \in (\mathbb{N} \to (\mathbb{R} \to \mathbb{R})). \forall k \in \mathbb{N}. \\ \quad \quad \quad \operatorname{is\text{-}antiderivative}(c(k), m(k), a, b) \wedge (c(k)(a) = 0) \end{array}}}

\vnbbookfile{theorem-library/antiderivable-uniform-limit}{theorem-library/antiderivable-uniform-limit}
\noindent{\small theorem-library/antiderivable-uniform-limit.scm -- docs/calculus.pdf PROP 4.16 and COR 4.17: THE ANTIDERIVABLE FUNCTIONS ARE CLOSED UNDER UNIFORM LIMITS.  Rung 3 of the integration arc.}\par
\vnbprop{rr-converges-to-abs}{\vnbstmt{\begin{array}{@{}l@{}} \forall f \in (\mathbb{N} \to \mathbb{R}), v \in \mathbb{R}. \\ \quad \forall s. \\ \quad \quad \operatorname{pos\text{-}rr}(s) \Rightarrow  \\ \quad \quad \quad \exists n \in \mathbb{N}.\; \forall k \in \mathbb{N}.\; ((n \leq k) \Rightarrow (\lvert (f(k) - v) \rvert \leq s)) \Rightarrow \\ \quad \quad \operatorname{converges\text{-}to}(\operatorname{rr\text{-}ms}, f, v) \end{array}}}
\vnbprop{rr-limit-ptwise-eq}{\vnbstmt{\begin{array}{@{}l@{}} \forall f, h \in (\mathbb{N} \to \mathbb{R}), v \in \mathbb{R}. \\ \quad (\forall j \in \mathbb{N}.\; (h(j) = f(j)) \wedge \operatorname{converges\text{-}to}(\operatorname{rr\text{-}ms}, f, v)) \Rightarrow  \\ \quad \quad \operatorname{converges\text{-}to}(\operatorname{rr\text{-}ms}, h, v) \end{array}}}
\vnbprop{uniform-cauchy-limit}{\vnbstmt{\begin{array}{@{}l@{}} \forall g \in (\mathbb{N} \to (\mathbb{R} \to \mathbb{R})), s. \\ \quad \forall a. \\ \quad \quad \operatorname{pos\text{-}rr}(a) \Rightarrow  \\ \quad \quad \quad \exists n \in \mathbb{N}. \forall k, l \in \mathbb{N}, x \in \mathbb{R}. \\ \quad \quad \quad \quad ((n \leq k) \wedge ((n \leq l) \wedge (x \in s))) \Rightarrow (\lvert (g(k)(x) - g(l)(x)) \rvert \leq a) \Rightarrow \\ \quad \quad \exists h \in (\mathbb{R} \to \mathbb{R}). \\ \quad \quad \quad \forall x \in \mathbb{R}. \\ \quad \quad \quad \quad (x \in s) \Rightarrow \operatorname{converges\text{-}to}(\operatorname{rr\text{-}ms}, (\lambda k \in \mathbb{N}.\; g(k)(x)), h(x)) \wedge \\ \quad \quad \quad \forall a. \\ \quad \quad \quad \quad \operatorname{pos\text{-}rr}(a) \Rightarrow  \\ \quad \quad \quad \quad \quad \exists n \in \mathbb{N}. \forall k \in \mathbb{N}, x \in \mathbb{R}. \\ \quad \quad \quad \quad \quad \quad ((n \leq k) \wedge (x \in s)) \Rightarrow (\lvert (h(x) - g(k)(x)) \rvert \leq a) \end{array}}}
\vnbprop{caratheodory-witness-pair-bound}{\vnbstmt{\begin{array}{@{}l@{}} \forall c, d, f, g, h, k \in (\mathbb{R} \to \mathbb{R}), a, b, t, m, x \in \mathbb{R}. \\ \quad \operatorname{is\text{-}antiderivative}(c, f, a, b) \wedge \\ \quad \operatorname{is\text{-}antiderivative}(d, g, a, b) \wedge \\ \quad \forall x \in \mathbb{R}.\; ((x \in \operatorname{ccint}(a, b)) \Rightarrow (\lvert (f(x) - g(x)) \rvert \leq m)) \wedge \\ \quad t \in \operatorname{ccint}(a, b) \wedge \\ \quad h(t) = f(t) \wedge \\ \quad k(t) = g(t) \wedge \\ \quad \forall x \in \mathbb{R}.\; ((c(x) - c(t)) = (h(x) \cdot (x - t))) \wedge \\ \quad \forall x \in \mathbb{R}.\; ((d(x) - d(t)) = (k(x) \cdot (x - t))) \wedge \\ \quad x \in \operatorname{ccint}(a, b) \Rightarrow \\ \quad \quad \lvert (h(x) - k(x)) \rvert \leq m \end{array}}}
\vnbprop{uniform-witness-diff-at}{\vnbstmt{\begin{array}{@{}l@{}} \forall m, c \in (\mathbb{N} \to (\mathbb{R} \to \mathbb{R})), d \in (\mathbb{R} \to \mathbb{R}), v, a, b, t, o \in \mathbb{R}. \\ \quad 0 < o \wedge \\ \quad \forall y \in \mathbb{R}.\; ((\lvert (y - t) \rvert \leq o) \Rightarrow (y \in \operatorname{ccint}(a, b))) \wedge \\ \quad t \in \operatorname{ccint}(a, b) \wedge \\ \quad \forall k \in \mathbb{N}.\; \operatorname{is\text{-}continuous\text{-}at}(\operatorname{rr\text{-}ms}, \operatorname{rr\text{-}ms}, c(k), t) \wedge \\ \quad \forall k \in \mathbb{N}, y \in \mathbb{R}. \\ \quad \quad (y \in \operatorname{ccint}(a, b)) \Rightarrow ((m(k)(y) - m(k)(t)) = (c(k)(y) \cdot (y - t))) \wedge \\ \quad \forall y \in \mathbb{R}. \\ \quad \quad (y \in \operatorname{ccint}(a, b)) \Rightarrow \operatorname{converges\text{-}to}(\operatorname{rr\text{-}ms}, (\lambda k \in \mathbb{N}.\; m(k)(y)), d(y)) \wedge \\ \quad \forall s. \\ \quad \quad \operatorname{pos\text{-}rr}(s) \Rightarrow  \\ \quad \quad \quad \exists n \in \mathbb{N}. \forall k, l \in \mathbb{N}, x \in \mathbb{R}. \\ \quad \quad \quad \quad ((n \leq k) \wedge ((n \leq l) \wedge (x \in \operatorname{ccint}(a, b)))) \Rightarrow (\lvert (c(k)(x) - c(l)(x)) \rvert \leq s) \wedge \\ \quad \operatorname{converges\text{-}to}(\operatorname{rr\text{-}ms}, (\lambda k \in \mathbb{N}.\; c(k)(t)), v) \Rightarrow \\ \quad \quad \operatorname{is\text{-}diff\text{-}at}(d, t, v) \end{array}}}
\vnbprop{antiderivative-family-uniform-limit}{\vnbstmt{\begin{array}{@{}l@{}} \forall m, c \in (\mathbb{N} \to (\mathbb{R} \to \mathbb{R})), i \in (\mathbb{R} \to \mathbb{R}), a, b \in \mathbb{R}. \\ \quad a < b \wedge \\ \quad \forall k \in \mathbb{N}.\; (\operatorname{is\text{-}antiderivative}(m(k), c(k), a, b) \wedge (m(k)(a) = 0)) \wedge \\ \quad \forall s. \\ \quad \quad \operatorname{pos\text{-}rr}(s) \Rightarrow  \\ \quad \quad \quad \exists n \in \mathbb{N}. \forall k \in \mathbb{N}, x \in \mathbb{R}. \\ \quad \quad \quad \quad ((n \leq k) \wedge (x \in \operatorname{ccint}(a, b))) \Rightarrow (\lvert (i(x) - c(k)(x)) \rvert \leq s) \Rightarrow \\ \quad \quad \exists f \in (\mathbb{R} \to \mathbb{R}). \\ \quad \quad \quad \operatorname{is\text{-}antiderivative}(f, i, a, b) \wedge \\ \quad \quad \quad \forall s. \\ \quad \quad \quad \quad \operatorname{pos\text{-}rr}(s) \Rightarrow  \\ \quad \quad \quad \quad \quad \exists n \in \mathbb{N}. \forall k \in \mathbb{N}, x \in \mathbb{R}. \\ \quad \quad \quad \quad \quad \quad ((n \leq k) \wedge (x \in \operatorname{ccint}(a, b))) \Rightarrow (\lvert (f(x) - m(k)(x)) \rvert \leq s) \end{array}}}
\vnbprop{antiderivable-uniform-limit}{\vnbstmt{\begin{array}{@{}l@{}} \forall m \in (\mathbb{N} \to (\mathbb{R} \to \mathbb{R})), i \in (\mathbb{R} \to \mathbb{R}), a, b \in \mathbb{R}. \\ \quad \forall k \in \mathbb{N}.\; \operatorname{is\text{-}antiderivable}(m(k), a, b) \wedge \\ \quad \forall s. \\ \quad \quad \operatorname{pos\text{-}rr}(s) \Rightarrow  \\ \quad \quad \quad \exists n \in \mathbb{N}. \forall k \in \mathbb{N}, x \in \mathbb{R}. \\ \quad \quad \quad \quad ((n \leq k) \wedge (x \in \operatorname{ccint}(a, b))) \Rightarrow (\lvert (i(x) - m(k)(x)) \rvert \leq s) \Rightarrow \\ \quad \quad \operatorname{is\text{-}antiderivable}(i, a, b) \end{array}}}

\vnbbookfile{theorem-library/series-antiderivable}{theorem-library/series-antiderivable}
\noindent{\small series-antiderivable.scm -- A FINITE SUM OF ANTIDERIVABLE TERMS IS ANTIDERIVABLE, and the vocabulary that makes such a statement usable.}\par
\vnbprop{antiderivable-fn-in-fun}{\vnbstmt{\forall i, a, b.\; (\operatorname{is\text{-}antiderivable}(i, a, b) \Rightarrow (i \in (\mathbb{R} \to \mathbb{R})))}}
\vnbprop{antiderivative-integrand-transfer}{\vnbstmt{\begin{array}{@{}l@{}} \forall f, i, c, a, b. \\ \quad ((c \in (\mathbb{R} \to \mathbb{R})) \wedge (\operatorname{is\text{-}antiderivative}(f, i, a, b) \wedge \forall x \in \mathbb{R}.\; (c(x) \simeq i(x)))) \Rightarrow  \\ \quad \quad \operatorname{is\text{-}antiderivative}(f, c, a, b) \end{array}}}
\vnbprop{antiderivable-integrand-transfer}{\vnbstmt{\begin{array}{@{}l@{}} \forall i, c, a, b. \\ \quad ((c \in (\mathbb{R} \to \mathbb{R})) \wedge (\operatorname{is\text{-}antiderivable}(i, a, b) \wedge \forall x \in \mathbb{R}.\; (c(x) \simeq i(x)))) \Rightarrow  \\ \quad \quad \operatorname{is\text{-}antiderivable}(c, a, b) \end{array}}}
\vnbprop{antiderivable-from-deriv}{\vnbstmt{\begin{array}{@{}l@{}} \forall f, i, a, b. \\ \quad f \in (\mathbb{R} \to \mathbb{R}) \wedge \\ \quad i \in (\mathbb{R} \to \mathbb{R}) \wedge \\ \quad a \in \mathbb{R} \wedge \\ \quad b \in \mathbb{R} \wedge \\ \quad a < b \wedge \\ \quad \forall x \in \mathbb{R}.\; \operatorname{is\text{-}diff\text{-}at}(f, x, i(x)) \Rightarrow \\ \quad \quad \operatorname{is\text{-}antiderivable}(i, a, b) \end{array}}}
\vnbprop{zero-is-antiderivable}{\vnbstmt{\begin{array}{@{}l@{}} \forall a, b. \\ \quad ((a \in \mathbb{R}) \wedge ((b \in \mathbb{R}) \wedge (a < b))) \Rightarrow  \\ \quad \quad \operatorname{is\text{-}antiderivable}((\lambda x \in \mathbb{R}.\; 0), a, b) \end{array}}}
\vnbprop{series-partial-sum-antiderivable}{\vnbstmt{\begin{array}{@{}l@{}} \forall n \in \mathbb{N}, s, h, a, b. \\ \quad a \in \mathbb{R} \wedge \\ \quad b \in \mathbb{R} \wedge \\ \quad a < b \wedge \\ \quad h \in (\mathbb{R} \to \mathbb{R}) \wedge \\ \quad \forall j \in \mathbb{N}.\; \operatorname{is\text{-}antiderivable}(s(j), a, b) \wedge \\ \quad \forall x \in \mathbb{R}.\; (h(x) \simeq \operatorname{series\text{-}partial\text{-}sum}((\lambda j \in \mathbb{N}.\; s(j)(x)), n)) \Rightarrow \\ \quad \quad \operatorname{is\text{-}antiderivable}(h, a, b) \end{array}}}
\vnbprop{power-antiderivable}{\vnbstmt{\begin{array}{@{}l@{}} \forall m, a, b. \\ \quad ((m \in \mathbb{N}) \wedge ((a \in \mathbb{R}) \wedge ((b \in \mathbb{R}) \wedge (a < b)))) \Rightarrow  \\ \quad \quad \operatorname{is\text{-}antiderivable}((\lambda x \in \mathbb{R}.\; \operatorname{power}(x, m)), a, b) \end{array}}}
\vnbprop{poly-fn-antiderivable}{\vnbstmt{\begin{array}{@{}l@{}} \forall n \in \mathbb{N}, f, h, a, b. \\ \quad f \in (\mathbb{N} \to \mathbb{R}) \wedge \\ \quad a \in \mathbb{R} \wedge \\ \quad b \in \mathbb{R} \wedge \\ \quad a < b \wedge \\ \quad h \in (\mathbb{R} \to \mathbb{R}) \wedge \\ \quad \forall x \in \mathbb{R}. \\ \quad \quad (h(x) \simeq \operatorname{series\text{-}partial\text{-}sum}((\lambda k \in \mathbb{N}.\; (f(k) \cdot \operatorname{power}(x, k))), n)) \Rightarrow \\ \quad \quad \operatorname{is\text{-}antiderivable}(h, a, b) \end{array}}}

\vnbbookfile{theorem-library/isom2-rrvs}{theorem-library/isom2-rrvs}
\noindent{\small isom2-rrvs.scm -- the real vector space of functions RR -$>$ RR in a set closed under the pointwise operations, and calculus.pdf Proposition 4.8: the antiderivable functions on [a, b] form a vector space.  Agent ISOM, 2026-10-26.}\par
\vnbprop{rr-scal-carr}{\vnbstmt{(\operatorname{carr}(\operatorname{field\text{-}as\text{-}integral\text{-}domain}(\operatorname{rr\text{-}field})) \simeq \mathbb{R})}}
\vnbprop{rr-scal-add}{\vnbstmt{(\operatorname{add}(\operatorname{field\text{-}as\text{-}integral\text{-}domain}(\operatorname{rr\text{-}field})) \simeq (\lambda {[}\operatorname{x\_}, \operatorname{y\_}] \in (\mathbb{R} \times \mathbb{R}).\; (\operatorname{x\_} + \operatorname{y\_})))}}
\vnbprop{rr-scal-mul}{\vnbstmt{(\operatorname{mul}(\operatorname{field\text{-}as\text{-}integral\text{-}domain}(\operatorname{rr\text{-}field})) \simeq (\lambda {[}\operatorname{x\_}, \operatorname{y\_}] \in (\mathbb{R} \times \mathbb{R}).\; (\operatorname{x\_} \cdot \operatorname{y\_})))}}
\vnbprop{rr-scal-one}{\vnbstmt{(\operatorname{one}(\operatorname{field\text{-}as\text{-}integral\text{-}domain}(\operatorname{rr\text{-}field})) \simeq 1)}}
\vnbprop{rr-scal-zero}{\vnbstmt{(\operatorname{zero}(\operatorname{field\text{-}as\text{-}integral\text{-}domain}(\operatorname{rr\text{-}field})) \simeq 0)}}
\vnbprop{rr-fn-space-vadd-type}{\vnbstmt{\begin{array}{@{}l@{}} \forall v \in \mathbf{Set}. \\ \quad (v \subseteq (\mathbb{R} \to \mathbb{R})) \Rightarrow  \\ \quad \quad (\lambda {[}\operatorname{rff\_}, \operatorname{rfg\_}] \in (v \times v).\; (\lambda x \in \mathbb{R}.\; (\operatorname{rff\_}(x) + \operatorname{rfg\_}(x)))) \\ \quad \quad \in\; ((v \times v) \to (\mathbb{R} \to \mathbb{R})) \end{array}}}
\vnbprop{rr-fn-space-vzero-type}{\vnbstmt{(\lambda x \in \mathbb{R}.\; 0) \in (\mathbb{R} \to \mathbb{R})}}
\vnbprop{rr-fn-space-vneg-type}{\vnbstmt{\begin{array}{@{}l@{}} \forall v \in \mathbf{Set}. \\ \quad (v \subseteq (\mathbb{R} \to \mathbb{R})) \Rightarrow  \\ \quad \quad (\lambda f \in v.\; (\lambda x \in \mathbb{R}.\; -f(x))) \in (v \to (\mathbb{R} \to \mathbb{R})) \end{array}}}
\vnbprop{rr-fn-space-act-type}{\vnbstmt{\begin{array}{@{}l@{}} \forall v \in \mathbf{Set}. \\ \quad (v \subseteq (\mathbb{R} \to \mathbb{R})) \Rightarrow  \\ \quad \quad (\lambda {[}\operatorname{rfc\_}, \operatorname{rff\_}] \in (\mathbb{R} \times v).\; (\lambda x \in \mathbb{R}.\; (\operatorname{rfc\_} \cdot \operatorname{rff\_}(x)))) \\ \quad \quad \in\; ((\mathbb{R} \times v) \to (\mathbb{R} \to \mathbb{R})) \end{array}}}
\vnbprop{rr-fn-space-scal}{\vnbstmt{\begin{array}{@{}l@{}} \forall v \in \mathbf{Set}. \\ \quad (v \subseteq (\mathbb{R} \to \mathbb{R})) \Rightarrow  \\ \quad \quad (\operatorname{scal}(\operatorname{rr\text{-}fn\text{-}space}(v)) \simeq \operatorname{field\text{-}as\text{-}integral\text{-}domain}(\operatorname{rr\text{-}field})) \end{array}}}
\vnbprop{rr-fn-space-vec}{\vnbstmt{\forall v \in \mathbf{Set}.\; ((v \subseteq (\mathbb{R} \to \mathbb{R})) \Rightarrow (\operatorname{vec}(\operatorname{rr\text{-}fn\text{-}space}(v)) \simeq v))}}
\vnbprop{rr-fn-space-vadd}{\vnbstmt{\begin{array}{@{}l@{}} \forall v \in \mathbf{Set}. \\ \quad (v \subseteq (\mathbb{R} \to \mathbb{R})) \Rightarrow  \\ \quad \quad (\operatorname{vadd}(\operatorname{rr\text{-}fn\text{-}space}(v)) \simeq (\lambda {[}\operatorname{rff\_}, \operatorname{rfg\_}] \in (v \times v).\; (\lambda x \in \mathbb{R}.\; (\operatorname{rff\_}(x) + \operatorname{rfg\_}(x))))) \end{array}}}
\vnbprop{rr-fn-space-vzero}{\vnbstmt{\begin{array}{@{}l@{}} \forall v \in \mathbf{Set}. \\ \quad (v \subseteq (\mathbb{R} \to \mathbb{R})) \Rightarrow (\operatorname{vzero}(\operatorname{rr\text{-}fn\text{-}space}(v)) \simeq (\lambda x \in \mathbb{R}.\; 0)) \end{array}}}
\vnbprop{rr-fn-space-vneg}{\vnbstmt{\begin{array}{@{}l@{}} \forall v \in \mathbf{Set}. \\ \quad (v \subseteq (\mathbb{R} \to \mathbb{R})) \Rightarrow  \\ \quad \quad (\operatorname{vneg}(\operatorname{rr\text{-}fn\text{-}space}(v)) \simeq (\lambda f \in v.\; (\lambda x \in \mathbb{R}.\; -f(x)))) \end{array}}}
\vnbprop{rr-fn-space-act}{\vnbstmt{\begin{array}{@{}l@{}} \forall v \in \mathbf{Set}. \\ \quad (v \subseteq (\mathbb{R} \to \mathbb{R})) \Rightarrow  \\ \quad \quad (\operatorname{act}(\operatorname{rr\text{-}fn\text{-}space}(v)) \simeq (\lambda {[}\operatorname{rfc\_}, \operatorname{rff\_}] \in (\mathbb{R} \times v).\; (\lambda x \in \mathbb{R}.\; (\operatorname{rfc\_} \cdot \operatorname{rff\_}(x))))) \end{array}}}
\vnbprop{rr-fn-space-is-vector-space}{\vnbstmt{\begin{array}{@{}l@{}} \forall v \in \mathbf{Set}. \\ \quad v \subseteq (\mathbb{R} \to \mathbb{R}) \wedge \\ \quad \forall f, g \in v.\; ((\lambda x \in \mathbb{R}.\; (f(x) + g(x))) \in v) \wedge \\ \quad (\lambda x \in \mathbb{R}.\; 0) \in v \wedge \\ \quad \forall f \in v.\; ((\lambda x \in \mathbb{R}.\; -f(x)) \in v) \wedge \\ \quad \forall c \in \mathbb{R}, f \in v.\; ((\lambda x \in \mathbb{R}.\; (c \cdot f(x))) \in v) \Rightarrow \\ \quad \quad \operatorname{is\text{-}vector\text{-}space}(\operatorname{rr\text{-}fn\text{-}space}(v)) \end{array}}}
\vnbprop{antiderivable-fns-unfold}{\vnbstmt{\begin{array}{@{}l@{}} \forall a, b. \\ \quad (\operatorname{antiderivable\text{-}fns}(a, b) \simeq \{p \in (\mathbb{R} \to \mathbb{R}) : \operatorname{is\text{-}antiderivable}(p, a, b)\}) \end{array}}}
\vnbprop{antiderivable-fns-membership}{\vnbstmt{\begin{array}{@{}l@{}} \forall a, b, i. \\ \quad (i \in \operatorname{antiderivable\text{-}fns}(a, b)) \Leftrightarrow ((i \in (\mathbb{R} \to \mathbb{R})) \wedge \operatorname{is\text{-}antiderivable}(i, a, b)) \end{array}}}
\vnbprop{antiderivable-fns-vector-space}{\vnbstmt{\begin{array}{@{}l@{}} \forall a, b. \\ \quad ((a \in \mathbb{R}) \wedge ((b \in \mathbb{R}) \wedge (a < b))) \Rightarrow  \\ \quad \quad \operatorname{is\text{-}vector\text{-}space}(\operatorname{rr\text{-}fn\text{-}space}(\operatorname{antiderivable\text{-}fns}(a, b))) \end{array}}}

\vnbbookfile{theorem-library/c-int}{theorem-library/c-int}
\noindent{\small c-int.scm -- docs/calculus.pdf CHAPTER 4, equation (64): the DEFINITE INTEGRAL of a continuous function, `C-INT'.  Rung 4.}\par
\par\medskip\noindent\textbf{Definition} (\texttt{c-int}).\;$\int_{a}^{b} \varphi :=$
\begin{equation*}
(\iota\, \operatorname{v\_}.\;((\operatorname{v\_} \in \mathbb{R}) \wedge \exists \operatorname{f\_}.\; (\operatorname{is\text{-}antiderivative}(\operatorname{f\_}, \varphi, a, b) \wedge (\operatorname{v\_} = (\operatorname{f\_}(b) - \operatorname{f\_}(a))))))
\end{equation*}

\vnbprop{c-int-value}{\vnbstmt{\begin{array}{@{}l@{}} \forall f, i, a, b. \\ \quad \operatorname{is\text{-}antiderivative}(f, i, a, b) \Rightarrow (\int_{a}^{b} i = (f(b) - f(a))) \end{array}}}
\vnbprop{c-int-in-rr}{\vnbstmt{\forall i, a, b.\; (\operatorname{is\text{-}antiderivable}(i, a, b) \Rightarrow (\int_{a}^{b} i \in \mathbb{R}))}}
\vnbprop{c-int-add}{\vnbstmt{\begin{array}{@{}l@{}} \forall i, c, a, b. \\ \quad (\operatorname{is\text{-}antiderivable}(i, a, b) \wedge \operatorname{is\text{-}antiderivable}(c, a, b)) \Rightarrow  \\ \quad \quad \int_{a}^{b} (\lambda x \in \mathbb{R}.\; (i(x) + c(x))) = (\int_{a}^{b} i + \int_{a}^{b} c) \end{array}}}
\vnbprop{c-int-scale}{\vnbstmt{\begin{array}{@{}l@{}} \forall c, i, a, b. \\ \quad ((c \in \mathbb{R}) \wedge \operatorname{is\text{-}antiderivable}(i, a, b)) \Rightarrow  \\ \quad \quad \int_{a}^{b} (\lambda x \in \mathbb{R}.\; (c \cdot i(x))) = (c \cdot \int_{a}^{b} i) \end{array}}}

\vnbbookfile{theorem-library/c-int-oriented}{theorem-library/c-int-oriented}
\noindent{\small c-int-oriented.scm -- THE ORIENTED DEFINITE INTEGRAL, the sub-interval restriction of Def 4.6 (docs/calculus.pdf REMARK 4.9), and ADDITIVITY IN THE BOUNDS.}\par
\par\medskip\noindent\textbf{Definition} (\texttt{c-int-or}).\;$\int_{a}^{b} \varphi :=$
\begin{equation*}
\left(\text{if }\operatorname{is\text{-}antiderivable}(\varphi, a, b)\text{ then }\int_{a}^{b} \varphi\text{ else }\left(\text{if }\operatorname{is\text{-}antiderivable}(\varphi, b, a)\text{ then }-\int_{b}^{a} \varphi\text{ else }0\right)\right)
\end{equation*}

\vnbprop{rr-lt-asymm}{\vnbstmt{\forall a, b \in \mathbb{R}.\; ((a < b) \Rightarrow \neg (b < a))}}
\vnbprop{antiderivable-endpoints}{\vnbstmt{\begin{array}{@{}l@{}} \forall i, a, b. \\ \quad \operatorname{is\text{-}antiderivable}(i, a, b) \Rightarrow ((a \in \mathbb{R}) \wedge ((b \in \mathbb{R}) \wedge (a < b))) \end{array}}}
\vnbprop{antiderivable-not-reverse}{\vnbstmt{\begin{array}{@{}l@{}} \forall i, a, b. \\ \quad \operatorname{is\text{-}antiderivable}(i, a, b) \Rightarrow \neg \operatorname{is\text{-}antiderivable}(i, b, a) \end{array}}}
\vnbprop{antiderivable-not-degenerate}{\vnbstmt{\forall i, a.\; \neg \operatorname{is\text{-}antiderivable}(i, a, a)}}
\vnbprop{c-int-or-anti}{\vnbstmt{\begin{array}{@{}l@{}} \forall i, a, b. \\ \quad \operatorname{is\text{-}antiderivable}(i, a, b) \Rightarrow (\int_{a}^{b} i = \int_{a}^{b} i) \end{array}}}
\vnbprop{c-int-or-flip}{\vnbstmt{\begin{array}{@{}l@{}} \forall i, a, b. \\ \quad \operatorname{is\text{-}antiderivable}(i, a, b) \Rightarrow (\int_{b}^{a} i = -\int_{a}^{b} i) \end{array}}}
\vnbprop{c-int-or-none}{\vnbstmt{\begin{array}{@{}l@{}} \forall i, a, b. \\ \quad (\neg \operatorname{is\text{-}antiderivable}(i, a, b) \wedge \neg \operatorname{is\text{-}antiderivable}(i, b, a)) \Rightarrow  \\ \quad \quad \int_{a}^{b} i = 0 \end{array}}}
\vnbprop{c-int-or-degenerate}{\vnbstmt{\forall i, a.\; (\int_{a}^{a} i = 0)}}
\vnbprop{c-int-or-in-rr}{\vnbstmt{\forall i, a, b.\; (\int_{a}^{b} i \in \mathbb{R})}}
\vnbprop{c-int-or-reverse}{\vnbstmt{\forall i, a, b.\; (\int_{b}^{a} i = -\int_{a}^{b} i)}}
\vnbprop{antiderivative-subinterval}{\vnbstmt{\begin{array}{@{}l@{}} \forall f, i, a, b, o, c. \\ \quad (\operatorname{is\text{-}antiderivative}(f, i, a, b) \wedge ((o \in \mathbb{R}) \wedge ((c \in \mathbb{R}) \wedge ((a \leq o) \wedge ((o < c) \wedge (c \leq b)))))) \Rightarrow  \\ \quad \quad \operatorname{is\text{-}antiderivative}(f, i, o, c) \end{array}}}
\vnbprop{antiderivable-subinterval}{\vnbstmt{\begin{array}{@{}l@{}} \forall i, a, b, o, c. \\ \quad (\operatorname{is\text{-}antiderivable}(i, a, b) \wedge ((o \in \mathbb{R}) \wedge ((c \in \mathbb{R}) \wedge ((a \leq o) \wedge ((o < c) \wedge (c \leq b)))))) \Rightarrow  \\ \quad \quad \operatorname{is\text{-}antiderivable}(i, o, c) \end{array}}}
\vnbprop{c-int-or-value}{\vnbstmt{\begin{array}{@{}l@{}} \forall f, i, p, q, a, b. \\ \quad (\operatorname{is\text{-}antiderivative}(f, i, p, q) \wedge ((a \in \operatorname{ccint}(p, q)) \wedge (b \in \operatorname{ccint}(p, q)))) \Rightarrow  \\ \quad \quad \int_{a}^{b} i = (f(b) - f(a)) \end{array}}}
\vnbprop{c-int-or-additive}{\vnbstmt{\begin{array}{@{}l@{}} \forall i, p, q, a, b, c. \\ \quad (\operatorname{is\text{-}antiderivable}(i, p, q) \wedge ((a \in \operatorname{ccint}(p, q)) \wedge ((b \in \operatorname{ccint}(p, q)) \wedge (c \in \operatorname{ccint}(p, q))))) \Rightarrow  \\ \quad \quad \int_{a}^{c} i = (\int_{a}^{b} i + \int_{b}^{c} i) \end{array}}}
\vnbprop{c-int-add-bounds}{\vnbstmt{\begin{array}{@{}l@{}} \forall i, a, b, c. \\ \quad (\operatorname{is\text{-}antiderivable}(i, a, c) \wedge ((b \in \mathbb{R}) \wedge ((a < b) \wedge (b < c)))) \Rightarrow  \\ \quad \quad \int_{a}^{c} i = (\int_{a}^{b} i + \int_{b}^{c} i) \end{array}}}

\vnbbookfile{theorem-library/bernstein-ccint}{theorem-library/bernstein-ccint}
\noindent{\small bernstein-ccint.scm -- THEOREM 5.2 of docs/calculus.pdf moved off [0,1] and onto an arbitrary nondegenerate [a,b], by the affine change of variable}\par
\vnbprop{ccint-parts}{\vnbstmt{\forall o, i.; \forall t \in \operatorname{ccint}(o, i).\; ((t \in \mathbb{R}) \wedge ((o \leq t) \wedge (t \leq i)))}}
\vnbprop{ccint-affine-in}{\vnbstmt{\begin{array}{@{}l@{}} \forall a, b, t. \\ \quad (((a \in \mathbb{R}) \wedge ((b \in \mathbb{R}) \wedge (a < b))) \wedge (t \in \operatorname{ccint}(0, 1))) \Rightarrow  \\ \quad \quad (a + ((b - a) \cdot t)) \in \operatorname{ccint}(a, b) \end{array}}}
\vnbprop{ccint-coaffine-in}{\vnbstmt{\begin{array}{@{}l@{}} \forall a, b, x. \\ \quad (((a \in \mathbb{R}) \wedge ((b \in \mathbb{R}) \wedge (a < b))) \wedge (x \in \operatorname{ccint}(a, b))) \Rightarrow  \\ \quad \quad ((x - a) \cdot \operatorname{recip}((b - a))) \in \operatorname{ccint}(0, 1) \end{array}}}
\vnbprop{ccint-affine-inverse}{\vnbstmt{\begin{array}{@{}l@{}} \forall a, b, x. \\ \quad ((a \in \mathbb{R}) \wedge ((b \in \mathbb{R}) \wedge ((a < b) \wedge (x \in \mathbb{R})))) \Rightarrow  \\ \quad \quad (a + ((b - a) \cdot ((x - a) \cdot \operatorname{recip}((b - a))))) = x \end{array}}}
\vnbprop{affine-continuous-at}{\vnbstmt{\begin{array}{@{}l@{}} \forall c, m, t. \\ \quad ((c \in \mathbb{R}) \wedge ((m \in \mathbb{R}) \wedge (t \in \mathbb{R}))) \Rightarrow  \\ \quad \quad \operatorname{is\text{-}continuous\text{-}at}(\operatorname{rr\text{-}ms}, \operatorname{rr\text{-}ms}, (\lambda x \in \mathbb{R}.\; (c + (m \cdot x))), t) \end{array}}}
\vnbprop{bernstein-uniform-approximation-ccint}{\vnbstmt{\begin{array}{@{}l@{}} \forall n, a, b, p. \\ \quad n \in (\mathbb{R} \to \mathbb{R}) \wedge \\ \quad a \in \mathbb{R} \wedge \\ \quad b \in \mathbb{R} \wedge \\ \quad a < b \wedge \\ \quad \forall c \in \operatorname{ccint}(a, b).\; \operatorname{is\text{-}continuous\text{-}at}(\operatorname{rr\text{-}ms}, \operatorname{rr\text{-}ms}, n, c) \wedge \\ \quad \operatorname{pos\text{-}rr}(p) \Rightarrow \\ \quad \quad \exists \mathit{cap} \in \mathbb{N}. \forall d \in \mathbb{N}, c \in \operatorname{ccint}(a, b). \\ \quad \quad \quad (\mathit{cap} \leq d) \Rightarrow  \\ \quad \quad \quad \quad \lvert (n(c) - \operatorname{bernstein\text{-}poly}(\operatorname{compose}(n, (\lambda x \in \mathbb{R}.\; (a + ((b - a) \cdot x)))), d, ((c - a) \cdot \operatorname{recip}((b - a))))) \rvert \\ \quad \quad \quad \quad \leq\; p \end{array}}}

\vnbbookfile{theorem-library/rake-bolzano-weierstrass}{theorem-library/rake-bolzano-weierstrass}
\noindent{\small theorem-library/rake-bolzano-weierstrass.scm ==================================================================== The eps-GRID FINITE COVER of a closed interval -- batch 8-E, item (1) -- and the scaled archimedean property it needs.}\par
\vnbprop{rr-archimedean-multiple}{\vnbstmt{\forall v.; \forall \mathit{xv} \in \mathbb{R}.\; (\operatorname{pos\text{-}rr}(v) \Rightarrow \exists \mathit{nv} \in \mathbb{N}.\; (\mathit{xv} \leq (\mathit{nv} \cdot v)))}}
\vnbprop{ccint-eps-grid-cover-at}{\vnbstmt{\begin{array}{@{}l@{}} \forall v \in \mathbb{N}, \mathit{av}, \mathit{bv} \in \mathbb{R}, \mathit{ev}. \\ \quad (\operatorname{pos\text{-}rr}(\mathit{ev}) \wedge (\mathit{bv} \leq (\mathit{av} + (v \cdot \mathit{ev})))) \Rightarrow  \\ \quad \quad \exists \mathit{cv}. \\ \quad \quad \quad \operatorname{is\text{-}finite\text{-}cover}(\mathit{cv}, \operatorname{ccint}(\mathit{av}, \mathit{bv})) \wedge \\ \quad \quad \quad \forall \mathit{uv} \in \mathit{cv}, \mathit{pv}, \mathit{qv} \in \mathit{uv}. \\ \quad \quad \quad \quad (\mathit{pv} \in \mathbb{R}) \wedge ((\mathit{qv} \in \mathbb{R}) \wedge (\lvert (\mathit{pv} - \mathit{qv}) \rvert \leq \mathit{ev})) \end{array}}}
\vnbprop{ccint-eps-grid-cover}{\vnbstmt{\begin{array}{@{}l@{}} \forall v, \mathit{bv} \in \mathbb{R}, \mathit{ev}. \\ \quad \operatorname{pos\text{-}rr}(\mathit{ev}) \Rightarrow  \\ \quad \quad \exists \mathit{cv}. \\ \quad \quad \quad \operatorname{is\text{-}finite\text{-}cover}(\mathit{cv}, \operatorname{ccint}(v, \mathit{bv})) \wedge \\ \quad \quad \quad \forall \mathit{uv} \in \mathit{cv}, \mathit{pv}, \mathit{qv} \in \mathit{uv}. \\ \quad \quad \quad \quad (\mathit{pv} \in \mathbb{R}) \wedge ((\mathit{qv} \in \mathbb{R}) \wedge (\lvert (\mathit{pv} - \mathit{qv}) \rvert \leq \mathit{ev})) \end{array}}}
\vnbprop{ccint-grid-cover-seq}{\vnbstmt{\begin{array}{@{}l@{}} \forall v, \mathit{bv} \in \mathbb{R}, d. \\ \quad \forall \mathit{kv} \in \mathbb{N}.\; \operatorname{pos\text{-}rr}(d(\mathit{kv})) \Rightarrow  \\ \quad \quad \exists w. \forall \mathit{kv} \in \mathbb{N}. \\ \quad \quad \quad \operatorname{is\text{-}finite\text{-}cover}(w(\mathit{kv}), \operatorname{ccint}(v, \mathit{bv})) \wedge \\ \quad \quad \quad \forall \mathit{uv} \in w(\mathit{kv}), \mathit{pv}, \mathit{qv} \in \mathit{uv}. \\ \quad \quad \quad \quad (\mathit{pv} \in \mathbb{R}) \wedge ((\mathit{qv} \in \mathbb{R}) \wedge (\lvert (\mathit{pv} - \mathit{qv}) \rvert \leq d(\mathit{kv}))) \end{array}}}

\vnbbookfile{theorem-library/rake-bolzano-weierstrass-2}{theorem-library/rake-bolzano-weierstrass-2}
\noindent{\small NOTE 2026-09-20 (batch 12-B).  Where this file says the tree has no metric SUBSPACE structure, that is history: SUBSPACE-MS(s, A) and RESTRICT(f, A) are defined in structure-library/metric-subspace.scm, their laws are in theorem-library/metric-subspace-laws.scm and compact-subspace.scm, and}\par
\vnbprop{rr-interval-seq-has-convergent-subseq}{\vnbstmt{\begin{array}{@{}l@{}} \forall v, \mathit{hiv} \in \mathbb{R}, \mathit{hqv} \in (\mathbb{N} \to \operatorname{ccint}(v, \mathit{hiv})). \exists \mathit{phv}. \\ \quad (\operatorname{strictly\text{-}mono\text{-}nn}(\mathit{phv}) \wedge \exists \mathit{ptv} \in \mathbb{R}.\; \operatorname{converges\text{-}to}(\operatorname{rr\text{-}ms}, \operatorname{subseq}(\mathit{hqv}, \mathit{phv}), \mathit{ptv})) \end{array}}}
\vnbprop{rr-bolzano-weierstrass}{\vnbstmt{\begin{array}{@{}l@{}} \forall v \in (\mathbb{N} \to \mathbb{R}), \mathit{bdv} \in \mathbb{R}. \\ \quad \forall \mathit{iqv} \in \mathbb{N}.\; (\lvert v(\mathit{iqv}) \rvert \leq \mathit{bdv}) \Rightarrow  \\ \quad \quad \exists \mathit{phv}. \\ \quad \quad \quad (\operatorname{strictly\text{-}mono\text{-}nn}(\mathit{phv}) \wedge \exists \mathit{ptv} \in \mathbb{R}.\; \operatorname{converges\text{-}to}(\operatorname{rr\text{-}ms}, \operatorname{subseq}(v, \mathit{phv}), \mathit{ptv})) \end{array}}}
\vnbprop{bounded-block-converges}{\vnbstmt{\begin{array}{@{}l@{}} \forall v \in (\mathbb{N} \to \mathbb{R}), \mathit{bdv} \in \mathbb{R}, \mathit{jqv} \in \mathrm{Inf}(\mathbb{N}). \\ \quad \forall \mathit{iqv} \in \mathbb{N}.\; (\lvert v(\mathit{iqv}) \rvert \leq \mathit{bdv}) \Rightarrow  \\ \quad \quad \exists \mathit{bqv} \in \mathrm{Inf}(\mathbb{N}). \\ \quad \quad \quad (\mathit{bqv} \subseteq \mathit{jqv}) \wedge \exists \mathit{ptv} \in \mathbb{R}.\; \operatorname{converges\text{-}along}(\operatorname{rr\text{-}ms}, v, \mathit{bqv}, \mathit{ptv}) \end{array}}}

\vnbbookfile{theorem-library/rake-ascoli-diagonal}{theorem-library/rake-ascoli-diagonal}
\noindent{\small theorem-library/rake-ascoli-diagonal.scm ==================================================================== ASCOLI-ARZELA, the last stretch -- batch 11-F.  Two theorems:}\par
\vnbprop{ascoli-pointwise-diagonal}{\vnbstmt{\begin{array}{@{}l@{}} \forall s, m, q. \\ \quad (\operatorname{pointwise\text{-}bounded}(s, m) \wedge \operatorname{is\text{-}dense\text{-}seq}(s, q)) \Rightarrow  \\ \quad \quad \exists l.\; (\operatorname{strictly\text{-}mono\text{-}nn}(l) \wedge \operatorname{converges\text{-}on}(s, \operatorname{subseq}(m, l), q)) \end{array}}}
\vnbprop{ascoli-arzela-sequential}{\vnbstmt{\begin{array}{@{}l@{}} \forall s, m. \\ \quad \operatorname{is\text{-}compact}(s) \wedge \\ \quad \exists x.\; (x \in \operatorname{pts}(s)) \wedge \\ \quad m \in (\mathbb{N} \to (\operatorname{pts}(s) \to \mathbb{R})) \wedge \\ \quad \forall k \in \mathbb{N}.\; \operatorname{is\text{-}continuous}(s, \operatorname{rr\text{-}ms}, m(k)) \wedge \\ \quad \operatorname{is\text{-}equicontinuous}(s, \operatorname{rr\text{-}ms}, m) \wedge \\ \quad \operatorname{pointwise\text{-}bounded}(s, m) \Rightarrow \\ \quad \quad \exists l. \\ \quad \quad \quad \operatorname{strictly\text{-}mono\text{-}nn}(l) \wedge \\ \quad \quad \quad \exists g \in (\operatorname{pts}(s) \to \mathbb{R}). \\ \quad \quad \quad \quad \operatorname{is\text{-}continuous}(s, \operatorname{rr\text{-}ms}, g) \wedge \operatorname{converges\text{-}uniformly}(s, \operatorname{subseq}(m, l), g) \end{array}}}

\vnbbookfile{theorem-library/bernstein-antiderivable}{theorem-library/bernstein-antiderivable}
\noindent{\small bernstein-antiderivable.scm -- EVERY BERNSTEIN APPROXIMANT IS ANTIDERIVABLE.}\par
\vnbprop{bernstein-basis-degree-zero-diag}{\vnbstmt{\begin{array}{@{}l@{}} \forall x, k. \\ \quad ((x \in \mathbb{R}) \wedge ((k \in \mathbb{N}) \wedge (k = 0))) \Rightarrow (\operatorname{bernstein\text{-}basis}(0, x)(k) = 1) \end{array}}}
\vnbprop{bernstein-basis-degree-zero-off}{\vnbstmt{\begin{array}{@{}l@{}} \forall x, k. \\ \quad ((x \in \mathbb{R}) \wedge ((k \in \mathbb{N}) \wedge \neg (k = 0))) \Rightarrow (\operatorname{bernstein\text{-}basis}(0, x)(k) = 0) \end{array}}}
\vnbprop{bernstein-term-antiderivable}{\vnbstmt{\begin{array}{@{}l@{}} \forall n \in \mathbb{N}, a, b, k, j. \\ \quad (((a \in \mathbb{R}) \wedge ((b \in \mathbb{R}) \wedge (a < b))) \wedge ((k \in \mathbb{N}) \wedge (j \in \mathbb{N}))) \Rightarrow  \\ \quad \quad \operatorname{is\text{-}antiderivable}((\lambda x \in \mathbb{R}.\; (\operatorname{power}(x, j) \cdot \operatorname{bernstein\text{-}basis}(n, x)(k))), a, b) \end{array}}}
\vnbprop{bernstein-basis-antiderivable}{\vnbstmt{\begin{array}{@{}l@{}} \forall n, k, a, b. \\ \quad ((n \in \mathbb{N}) \wedge ((k \in \mathbb{N}) \wedge ((a \in \mathbb{R}) \wedge ((b \in \mathbb{R}) \wedge (a < b))))) \Rightarrow  \\ \quad \quad \operatorname{is\text{-}antiderivable}((\lambda x \in \mathbb{R}.\; \operatorname{bernstein\text{-}basis}(n, x)(k)), a, b) \end{array}}}
\vnbprop{bernstein-basis-scaled-antiderivable}{\vnbstmt{\begin{array}{@{}l@{}} \forall c, n, k, a, b. \\ \quad ((c \in \mathbb{R}) \wedge ((n \in \mathbb{N}) \wedge ((k \in \mathbb{N}) \wedge ((a \in \mathbb{R}) \wedge ((b \in \mathbb{R}) \wedge (a < b)))))) \Rightarrow  \\ \quad \quad \operatorname{is\text{-}antiderivable}((\lambda x \in \mathbb{R}.\; (c \cdot \operatorname{bernstein\text{-}basis}(n, x)(k))), a, b) \end{array}}}
\vnbprop{bernstein-poly-antiderivable}{\vnbstmt{\begin{array}{@{}l@{}} \forall f, n, a, b. \\ \quad ((f \in (\mathbb{R} \to \mathbb{R})) \wedge ((n \in \mathbb{N}) \wedge (\neg (n = 0) \wedge ((a \in \mathbb{R}) \wedge ((b \in \mathbb{R}) \wedge (a < b)))))) \Rightarrow  \\ \quad \quad \operatorname{is\text{-}antiderivable}((\lambda x \in \mathbb{R}.\; \operatorname{bernstein\text{-}poly}(f, n, x)), a, b) \end{array}}}

\vnbbookfile{theorem-library/antiderivable-affine-subst}{theorem-library/antiderivable-affine-subst}
\noindent{\small antiderivable-affine-subst.scm -- ANTIDERIVABILITY SURVIVES AN AFFINE CHANGE OF VARIABLE, and with it the rung-3 form of the Bernstein approximants.}\par
\vnbprop{antiderivable-affine-subst}{\vnbstmt{\begin{array}{@{}l@{}} \forall i, k, m, p, q, a, b. \\ \quad i \in (\mathbb{R} \to \mathbb{R}) \wedge \\ \quad k \in \mathbb{R} \wedge \\ \quad m \in \mathbb{R} \wedge \\ \quad \neg (m = 0) \wedge \\ \quad p \in \mathbb{R} \wedge \\ \quad q \in \mathbb{R} \wedge \\ \quad p < q \wedge \\ \quad \operatorname{is\text{-}antiderivable}(i, a, b) \wedge \\ \quad \forall t \in \operatorname{ccint}(p, q).\; ((k + (m \cdot t)) \in \operatorname{ccint}(a, b)) \wedge \\ \quad \forall t \in \mathbb{R}.\; (((p < t) \wedge (t < q)) \Rightarrow ((a < (k + (m \cdot t))) \wedge ((k + (m \cdot t)) < b))) \Rightarrow \\ \quad \quad \operatorname{is\text{-}antiderivable}((\lambda x \in \mathbb{R}.\; i((k + (m \cdot x)))), p, q) \end{array}}}
\vnbprop{antiderivable-affine-subst-pos}{\vnbstmt{\begin{array}{@{}l@{}} \forall i, k, m, p, q. \\ \quad i \in (\mathbb{R} \to \mathbb{R}) \wedge \\ \quad k \in \mathbb{R} \wedge \\ \quad m \in \mathbb{R} \wedge \\ \quad 0 < m \wedge \\ \quad p \in \mathbb{R} \wedge \\ \quad q \in \mathbb{R} \wedge \\ \quad p < q \wedge \\ \quad \operatorname{is\text{-}antiderivable}(i, (k + (m \cdot p)), (k + (m \cdot q))) \Rightarrow \\ \quad \quad \operatorname{is\text{-}antiderivable}((\lambda x \in \mathbb{R}.\; i((k + (m \cdot x)))), p, q) \end{array}}}
\vnbprop{antiderivable-affine-subst-pos-transfer}{\vnbstmt{\begin{array}{@{}l@{}} \forall i, a, k, m, p, q. \\ \quad i \in (\mathbb{R} \to \mathbb{R}) \wedge \\ \quad a \in (\mathbb{R} \to \mathbb{R}) \wedge \\ \quad k \in \mathbb{R} \wedge \\ \quad m \in \mathbb{R} \wedge \\ \quad 0 < m \wedge \\ \quad p \in \mathbb{R} \wedge \\ \quad q \in \mathbb{R} \wedge \\ \quad p < q \wedge \\ \quad \operatorname{is\text{-}antiderivable}(i, (k + (m \cdot p)), (k + (m \cdot q))) \wedge \\ \quad \forall x \in \mathbb{R}.\; (a(x) \simeq i((k + (m \cdot x)))) \Rightarrow \\ \quad \quad \operatorname{is\text{-}antiderivable}(a, p, q) \end{array}}}
\vnbprop{bernstein-poly-unit-coordinate-antiderivable}{\vnbstmt{\begin{array}{@{}l@{}} \forall g, n, a, b. \\ \quad ((g \in (\mathbb{R} \to \mathbb{R})) \wedge ((n \in \mathbb{N}) \wedge (\neg (n = 0) \wedge ((a \in \mathbb{R}) \wedge ((b \in \mathbb{R}) \wedge (a < b)))))) \Rightarrow  \\ \quad \quad \operatorname{is\text{-}antiderivable}((\lambda x \in \mathbb{R}.\; \operatorname{bernstein\text{-}poly}(g, n, ((x - a) \cdot \operatorname{recip}((b - a))))), a, b) \end{array}}}
\vnbprop{bernstein-ccint-approximant-antiderivable}{\vnbstmt{\begin{array}{@{}l@{}} \forall f, n, a, b. \\ \quad ((f \in (\mathbb{R} \to \mathbb{R})) \wedge ((n \in \mathbb{N}) \wedge (\neg (n = 0) \wedge ((a \in \mathbb{R}) \wedge ((b \in \mathbb{R}) \wedge (a < b)))))) \Rightarrow  \\ \quad \quad \operatorname{is\text{-}antiderivable}((\lambda x \in \mathbb{R}.\; \\ \quad \quad \quad \quad \operatorname{bernstein\text{-}poly}(\operatorname{compose}(f, (\lambda c \in \mathbb{R}.\; (a + ((b - a) \cdot c)))), n, ((x - a) \cdot \operatorname{recip}((b - a))))), \\ \quad \quad \quad a, b) \end{array}}}

\vnbbookfile{theorem-library/continuous-antiderivable}{theorem-library/continuous-antiderivable}
\noindent{\small continuous-antiderivable.scm -- RUNG 3, AND THE INTEGRAL THAT FOLLOWS.}\par
\vnbprop{continuous-is-antiderivable}{\vnbstmt{\begin{array}{@{}l@{}} \forall f, a, b. \\ \quad f \in (\mathbb{R} \to \mathbb{R}) \wedge \\ \quad a \in \mathbb{R} \wedge \\ \quad b \in \mathbb{R} \wedge \\ \quad a < b \wedge \\ \quad \forall x \in \operatorname{ccint}(a, b).\; \operatorname{is\text{-}continuous\text{-}at}(\operatorname{rr\text{-}ms}, \operatorname{rr\text{-}ms}, f, x) \Rightarrow \\ \quad \quad \operatorname{is\text{-}antiderivable}(f, a, b) \end{array}}}
\vnbprop{c-int-defined-for-continuous}{\vnbstmt{\begin{array}{@{}l@{}} \forall f, a, b. \\ \quad f \in (\mathbb{R} \to \mathbb{R}) \wedge \\ \quad a \in \mathbb{R} \wedge \\ \quad b \in \mathbb{R} \wedge \\ \quad a < b \wedge \\ \quad \forall x \in \operatorname{ccint}(a, b).\; \operatorname{is\text{-}continuous\text{-}at}(\operatorname{rr\text{-}ms}, \operatorname{rr\text{-}ms}, f, x) \Rightarrow \\ \quad \quad \int_{a}^{b} f \in \mathbb{R} \end{array}}}
\vnbprop{c-int-value-continuous}{\vnbstmt{\begin{array}{@{}l@{}} \forall f, a, b. \\ \quad f \in (\mathbb{R} \to \mathbb{R}) \wedge \\ \quad a \in \mathbb{R} \wedge \\ \quad b \in \mathbb{R} \wedge \\ \quad a < b \wedge \\ \quad \forall x \in \operatorname{ccint}(a, b).\; \operatorname{is\text{-}continuous\text{-}at}(\operatorname{rr\text{-}ms}, \operatorname{rr\text{-}ms}, f, x) \Rightarrow \\ \quad \quad \exists g.\; (\operatorname{is\text{-}antiderivative}(g, f, a, b) \wedge (\int_{a}^{b} f = (g(b) - g(a)))) \end{array}}}
\vnbprop{c-int-add-continuous}{\vnbstmt{\begin{array}{@{}l@{}} \forall i, c, a, b. \\ \quad i \in (\mathbb{R} \to \mathbb{R}) \wedge \\ \quad c \in (\mathbb{R} \to \mathbb{R}) \wedge \\ \quad a \in \mathbb{R} \wedge \\ \quad b \in \mathbb{R} \wedge \\ \quad a < b \wedge \\ \quad \forall d \in \operatorname{ccint}(a, b).\; \operatorname{is\text{-}continuous\text{-}at}(\operatorname{rr\text{-}ms}, \operatorname{rr\text{-}ms}, i, d) \wedge \\ \quad \forall d \in \operatorname{ccint}(a, b).\; \operatorname{is\text{-}continuous\text{-}at}(\operatorname{rr\text{-}ms}, \operatorname{rr\text{-}ms}, c, d) \Rightarrow \\ \quad \quad \int_{a}^{b} (\lambda x \in \mathbb{R}.\; (i(x) + c(x))) = (\int_{a}^{b} i + \int_{a}^{b} c) \end{array}}}
\vnbprop{c-int-scale-continuous}{\vnbstmt{\begin{array}{@{}l@{}} \forall c, i, a, b. \\ \quad c \in \mathbb{R} \wedge \\ \quad i \in (\mathbb{R} \to \mathbb{R}) \wedge \\ \quad a \in \mathbb{R} \wedge \\ \quad b \in \mathbb{R} \wedge \\ \quad a < b \wedge \\ \quad \forall d \in \operatorname{ccint}(a, b).\; \operatorname{is\text{-}continuous\text{-}at}(\operatorname{rr\text{-}ms}, \operatorname{rr\text{-}ms}, i, d) \Rightarrow \\ \quad \quad \int_{a}^{b} (\lambda x \in \mathbb{R}.\; (c \cdot i(x))) = (c \cdot \int_{a}^{b} i) \end{array}}}

\vnbbookfile{theorem-library/c-int-change-of-variable}{theorem-library/c-int-change-of-variable}
\noindent{\small c-int-change-of-variable.scm -- CHANGE OF VARIABLES FOR `C-INT', in the INCREASING case, docs/calculus.pdf's substitution rule:}\par
\vnbprop{change-of-variable-antiderivative}{\vnbstmt{\begin{array}{@{}l@{}} \forall w, i, p, r, a, q. \\ \quad p \in (\mathbb{R} \to \mathbb{R}) \wedge \\ \quad r \in (\mathbb{R} \to \mathbb{R}) \wedge \\ \quad a \in \mathbb{R} \wedge \\ \quad q \in \mathbb{R} \wedge \\ \quad a < q \wedge \\ \quad \operatorname{is\text{-}antiderivative}(w, i, p(a), p(q)) \wedge \\ \quad \forall u \in \operatorname{ccint}(a, q).\; \operatorname{is\text{-}continuous\text{-}at}(\operatorname{rr\text{-}ms}, \operatorname{rr\text{-}ms}, p, u) \wedge \\ \quad \forall b \in \mathbb{R}.\; (((a < b) \wedge (b < q)) \Rightarrow \operatorname{is\text{-}diff\text{-}at}(p, b, r(b))) \wedge \\ \quad \forall u \in \operatorname{ccint}(a, q).\; (p(u) \in \operatorname{ccint}(p(a), p(q))) \wedge \\ \quad \forall b \in \mathbb{R}.\; (((a < b) \wedge (b < q)) \Rightarrow ((p(a) < p(b)) \wedge (p(b) < p(q)))) \Rightarrow \\ \quad \quad \operatorname{is\text{-}antiderivative}(\operatorname{compose}(w, p), (\lambda x \in \mathbb{R}.\; (i(p(x)) \cdot r(x))), a, q) \end{array}}}
\vnbprop{change-of-variable-antiderivable}{\vnbstmt{\begin{array}{@{}l@{}} \forall i, p, r, a, q. \\ \quad p \in (\mathbb{R} \to \mathbb{R}) \wedge \\ \quad r \in (\mathbb{R} \to \mathbb{R}) \wedge \\ \quad a \in \mathbb{R} \wedge \\ \quad q \in \mathbb{R} \wedge \\ \quad a < q \wedge \\ \quad \operatorname{is\text{-}antiderivable}(i, p(a), p(q)) \wedge \\ \quad \forall u \in \operatorname{ccint}(a, q).\; \operatorname{is\text{-}continuous\text{-}at}(\operatorname{rr\text{-}ms}, \operatorname{rr\text{-}ms}, p, u) \wedge \\ \quad \forall w \in \mathbb{R}.\; (((a < w) \wedge (w < q)) \Rightarrow \operatorname{is\text{-}diff\text{-}at}(p, w, r(w))) \wedge \\ \quad \forall u \in \operatorname{ccint}(a, q).\; (p(u) \in \operatorname{ccint}(p(a), p(q))) \wedge \\ \quad \forall w \in \mathbb{R}.\; (((a < w) \wedge (w < q)) \Rightarrow ((p(a) < p(w)) \wedge (p(w) < p(q)))) \Rightarrow \\ \quad \quad \operatorname{is\text{-}antiderivable}((\lambda x \in \mathbb{R}.\; (i(p(x)) \cdot r(x))), a, q) \end{array}}}
\vnbprop{c-int-change-of-variable}{\vnbstmt{\begin{array}{@{}l@{}} \forall i, p, r, a, q. \\ \quad p \in (\mathbb{R} \to \mathbb{R}) \wedge \\ \quad r \in (\mathbb{R} \to \mathbb{R}) \wedge \\ \quad a \in \mathbb{R} \wedge \\ \quad q \in \mathbb{R} \wedge \\ \quad a < q \wedge \\ \quad \operatorname{is\text{-}antiderivable}(i, p(a), p(q)) \wedge \\ \quad \forall u \in \operatorname{ccint}(a, q).\; \operatorname{is\text{-}continuous\text{-}at}(\operatorname{rr\text{-}ms}, \operatorname{rr\text{-}ms}, p, u) \wedge \\ \quad \forall w \in \mathbb{R}.\; (((a < w) \wedge (w < q)) \Rightarrow \operatorname{is\text{-}diff\text{-}at}(p, w, r(w))) \wedge \\ \quad \forall u \in \operatorname{ccint}(a, q).\; (p(u) \in \operatorname{ccint}(p(a), p(q))) \wedge \\ \quad \forall w \in \mathbb{R}.\; (((a < w) \wedge (w < q)) \Rightarrow ((p(a) < p(w)) \wedge (p(w) < p(q)))) \Rightarrow \\ \quad \quad \int_{a}^{q} (\lambda x \in \mathbb{R}.\; (i(p(x)) \cdot r(x))) = \int_{p(a)}^{p(q)} i \end{array}}}
\vnbprop{c-int-change-of-variable-transfer}{\vnbstmt{\begin{array}{@{}l@{}} \forall i, a, p, r, b, q. \\ \quad p \in (\mathbb{R} \to \mathbb{R}) \wedge \\ \quad r \in (\mathbb{R} \to \mathbb{R}) \wedge \\ \quad b \in \mathbb{R} \wedge \\ \quad q \in \mathbb{R} \wedge \\ \quad b < q \wedge \\ \quad a \in (\mathbb{R} \to \mathbb{R}) \wedge \\ \quad \operatorname{is\text{-}antiderivable}(i, p(b), p(q)) \wedge \\ \quad \forall u \in \operatorname{ccint}(b, q).\; \operatorname{is\text{-}continuous\text{-}at}(\operatorname{rr\text{-}ms}, \operatorname{rr\text{-}ms}, p, u) \wedge \\ \quad \forall w \in \mathbb{R}.\; (((b < w) \wedge (w < q)) \Rightarrow \operatorname{is\text{-}diff\text{-}at}(p, w, r(w))) \wedge \\ \quad \forall u \in \operatorname{ccint}(b, q).\; (p(u) \in \operatorname{ccint}(p(b), p(q))) \wedge \\ \quad \forall w \in \mathbb{R}.\; (((b < w) \wedge (w < q)) \Rightarrow ((p(b) < p(w)) \wedge (p(w) < p(q)))) \wedge \\ \quad \forall x \in \mathbb{R}.\; (a(x) \simeq (i(p(x)) \cdot r(x))) \Rightarrow \\ \quad \quad \int_{b}^{q} a = \int_{p(b)}^{p(q)} i \end{array}}}

\vnbbookfile{theorem-library/log}{theorem-library/log}
\noindent{\small theorem-library/log.scm -- THE LOGARITHM, DEFINED AS AN INTEGRAL, and its basic laws.  The one that costs anything is the functional equation, and what it costs is a CHANGE OF VARIABLES.}\par
\par\medskip\noindent\textbf{Definition} (\texttt{log}).\;$\log\left(\operatorname{x\_}\right) :=$ $\int_{1}^{\operatorname{x\_}} (\lambda \operatorname{z\_} \in \mathbb{R}.\; {\operatorname{z\_}}^{\ast\!-1})$
\vnbprop{log-unfold}{\vnbstmt{\forall x.\; (\log\left(x\right) \simeq \int_{1}^{x} (\lambda z \in \mathbb{R}.\; {z}^{\ast\!-1}))}}
\vnbprop{log-in-rr}{\vnbstmt{\forall x.\; (\log\left(x\right) \in \mathbb{R})}}
\vnbprop{log-lam-in-fun}{\vnbstmt{(\lambda x \in \mathbb{R}.\; \log\left(x\right)) \in (\mathbb{R} \to \mathbb{R})}}
\vnbprop{log-one}{\vnbstmt{\log\left(1\right) = 0}}
\vnbprop{recip-star-antiderivable-pos}{\vnbstmt{\begin{array}{@{}l@{}} \forall a, b. \\ \quad ((a \in \mathbb{R}) \wedge ((b \in \mathbb{R}) \wedge ((0 < a) \wedge (a < b)))) \Rightarrow  \\ \quad \quad \operatorname{is\text{-}antiderivable}((\lambda z \in \mathbb{R}.\; {z}^{\ast\!-1}), a, b) \end{array}}}
\vnbprop{recip-star-scale}{\vnbstmt{\begin{array}{@{}l@{}} \forall m, u. \\ \quad ((m \in \mathbb{R}) \wedge ((u \in \mathbb{R}) \wedge \neg (m = 0))) \Rightarrow (({(m \cdot u)}^{\ast\!-1} \cdot m) = {u}^{\ast\!-1}) \end{array}}}
\vnbprop{c-int-or-scale-ordered}{\vnbstmt{\begin{array}{@{}l@{}} \forall m, a, b. \\ \quad ((m \in \mathbb{R}) \wedge ((a \in \mathbb{R}) \wedge ((b \in \mathbb{R}) \wedge ((0 < m) \wedge ((0 < a) \wedge (a < b)))))) \Rightarrow  \\ \quad \quad \int_{(m \cdot a)}^{(m \cdot b)} (\lambda z \in \mathbb{R}.\; {z}^{\ast\!-1}) \\ \quad \quad =\; \int_{a}^{b} (\lambda z \in \mathbb{R}.\; {z}^{\ast\!-1}) \end{array}}}
\vnbprop{c-int-or-scale}{\vnbstmt{\begin{array}{@{}l@{}} \forall m, a, b. \\ \quad ((m \in \mathbb{R}) \wedge ((a \in \mathbb{R}) \wedge ((b \in \mathbb{R}) \wedge ((0 < m) \wedge ((0 < a) \wedge (0 < b)))))) \Rightarrow  \\ \quad \quad \int_{(m \cdot a)}^{(m \cdot b)} (\lambda z \in \mathbb{R}.\; {z}^{\ast\!-1}) \\ \quad \quad =\; \int_{a}^{b} (\lambda z \in \mathbb{R}.\; {z}^{\ast\!-1}) \end{array}}}
\vnbprop{recip-star-int-additive}{\vnbstmt{\begin{array}{@{}l@{}} \forall u, v, w. \\ \quad ((u \in \mathbb{R}) \wedge ((v \in \mathbb{R}) \wedge ((w \in \mathbb{R}) \wedge ((0 < u) \wedge ((0 < v) \wedge (0 < w)))))) \Rightarrow  \\ \quad \quad \int_{u}^{w} (\lambda z \in \mathbb{R}.\; {z}^{\ast\!-1}) \\ \quad \quad =\; \int_{u}^{v} (\lambda z \in \mathbb{R}.\; {z}^{\ast\!-1}) + \int_{v}^{w} (\lambda z \in \mathbb{R}.\; {z}^{\ast\!-1}) \end{array}}}
\vnbprop{log-diff}{\vnbstmt{\begin{array}{@{}l@{}} \forall u, v. \\ \quad ((u \in \mathbb{R}) \wedge ((v \in \mathbb{R}) \wedge ((0 < u) \wedge (0 < v)))) \Rightarrow  \\ \quad \quad \int_{u}^{v} (\lambda z \in \mathbb{R}.\; {z}^{\ast\!-1}) = (\log\left(v\right) - \log\left(u\right)) \end{array}}}
\vnbprop{log-mul}{\vnbstmt{\begin{array}{@{}l@{}} \forall x, y. \\ \quad ((x \in \mathbb{R}) \wedge ((y \in \mathbb{R}) \wedge ((0 < x) \wedge (0 < y)))) \Rightarrow (\log\left((x \cdot y)\right) = (\log\left(x\right) + \log\left(y\right))) \end{array}}}
\vnbprop{log-strictly-increasing}{\vnbstmt{\forall u, v.\; (((u \in \mathbb{R}) \wedge ((v \in \mathbb{R}) \wedge ((0 < u) \wedge (u < v)))) \Rightarrow (\log\left(u\right) < \log\left(v\right)))}}
\vnbprop{log-neg-below-one}{\vnbstmt{\forall x \in \mathbb{R}.\; (((0 < x) \wedge (x < 1)) \Rightarrow (\log\left(x\right) < 0))}}
\vnbprop{log-pos-above-one}{\vnbstmt{\forall x \in \mathbb{R}.\; ((1 < x) \Rightarrow (0 < \log\left(x\right)))}}
\vnbprop{log-recip}{\vnbstmt{\forall x \in \mathbb{R}.\; ((0 < x) \Rightarrow (\log\left({x}^{\ast\!-1}\right) = -\log\left(x\right)))}}
\vnbprop{log-deriv}{\vnbstmt{\begin{array}{@{}l@{}} \forall c \in \mathbb{R}. \\ \quad (0 < c) \Rightarrow \operatorname{is\text{-}diff\text{-}at}((\lambda x \in \mathbb{R}.\; \log\left(x\right)), c, {c}^{\ast\!-1}) \end{array}}}

\vnbbookfile{theorem-library/r-exp}{theorem-library/r-exp}
\noindent{\small theorem-library/r-exp.scm -- THE REAL EXPONENTIAL, DEFINED AS THE INVERSE OF THE LOGARITHM, and its laws.}\par
\par\medskip\noindent\textbf{Definition} (\texttt{r-exp}).\;$\exp\left(\operatorname{y\_}\right) :=$ $(\iota\, \operatorname{x\_}.\;((\operatorname{x\_} \in \mathbb{R}) \wedge ((0 < \operatorname{x\_}) \wedge (\log\left(\operatorname{x\_}\right) = \operatorname{y\_}))))$
\vnbprop{log-two-pos}{\vnbstmt{0 < \log\left(2\right)}}
\vnbprop{log-unbounded-above}{\vnbstmt{\forall n \in \mathbb{N}.\; \exists b \in \mathbb{R}.\; ((0 < b) \wedge ((n \cdot \log\left(2\right)) \leq \log\left(b\right)))}}
\vnbprop{log-above-any}{\vnbstmt{\forall y \in \mathbb{R}.\; \exists b \in \mathbb{R}.\; ((0 < b) \wedge (y \leq \log\left(b\right)))}}
\vnbprop{log-below-any}{\vnbstmt{\forall y \in \mathbb{R}.\; \exists a \in \mathbb{R}.\; ((0 < a) \wedge (\log\left(a\right) \leq y))}}
\vnbprop{log-surjective}{\vnbstmt{\forall y \in \mathbb{R}.\; \exists x \in \mathbb{R}.\; ((0 < x) \wedge (\log\left(x\right) = y))}}
\vnbprop{log-injective-pos}{\vnbstmt{\begin{array}{@{}l@{}} \forall u, v. \\ \quad ((u \in \mathbb{R}) \wedge ((v \in \mathbb{R}) \wedge ((0 < u) \wedge ((0 < v) \wedge (\log\left(u\right) = \log\left(v\right)))))) \Rightarrow (u = v) \end{array}}}
\vnbprop{r-exp-prop}{\vnbstmt{\forall y \in \mathbb{R}.\; ((\exp\left(y\right) \in \mathbb{R}) \wedge ((0 < \exp\left(y\right)) \wedge (\log\left(\exp\left(y\right)\right) = y)))}}
\vnbprop{r-exp-in-rr}{\vnbstmt{\forall y \in \mathbb{R}.\; (\exp\left(y\right) \in \mathbb{R})}}
\vnbprop{r-exp-pos}{\vnbstmt{\forall y \in \mathbb{R}.\; (0 < \exp\left(y\right))}}
\vnbprop{log-r-exp}{\vnbstmt{\forall y \in \mathbb{R}.\; (\log\left(\exp\left(y\right)\right) = y)}}
\vnbprop{r-exp-char}{\vnbstmt{\begin{array}{@{}l@{}} \forall y, x. \\ \quad ((y \in \mathbb{R}) \wedge ((x \in \mathbb{R}) \wedge ((0 < x) \wedge (\log\left(x\right) = y)))) \Rightarrow (\exp\left(y\right) = x) \end{array}}}
\vnbprop{r-exp-lam-in-fun}{\vnbstmt{(\lambda y \in \mathbb{R}.\; \exp\left(y\right)) \in (\mathbb{R} \to \mathbb{R})}}
\vnbprop{r-exp-zero}{\vnbstmt{\exp\left(0\right) = 1}}
\vnbprop{r-exp-log}{\vnbstmt{\forall x \in \mathbb{R}.\; ((0 < x) \Rightarrow (\exp\left(\log\left(x\right)\right) = x))}}
\vnbprop{r-exp-add}{\vnbstmt{\forall u, v.\; (((u \in \mathbb{R}) \wedge (v \in \mathbb{R})) \Rightarrow (\exp\left((u + v)\right) = (\exp\left(u\right) \cdot \exp\left(v\right))))}}
\vnbprop{r-exp-strictly-increasing}{\vnbstmt{\forall u, v.\; (((u \in \mathbb{R}) \wedge ((v \in \mathbb{R}) \wedge (u < v))) \Rightarrow (\exp\left(u\right) < \exp\left(v\right)))}}
\vnbprop{r-exp-mono}{\vnbstmt{\forall u, v.\; (((u \in \mathbb{R}) \wedge ((v \in \mathbb{R}) \wedge (u \leq v))) \Rightarrow (\exp\left(u\right) \leq \exp\left(v\right)))}}
\vnbprop{r-exp-neg}{\vnbstmt{\forall y \in \mathbb{R}.\; (\exp\left(-y\right) = {\exp\left(y\right)}^{\ast\!-1})}}
\vnbprop{r-exp-log-not-global}{\vnbstmt{\neg \forall x \in \mathbb{R}.\; (\exp\left(\log\left(x\right)\right) = x)}}
\vnbprop{r-exp-continuous-at}{\vnbstmt{\begin{array}{@{}l@{}} \forall t \in \mathbb{R}. \\ \quad \operatorname{is\text{-}continuous\text{-}at}(\operatorname{rr\text{-}ms}, \operatorname{rr\text{-}ms}, (\lambda y \in \mathbb{R}.\; \exp\left(y\right)), t) \end{array}}}
\vnbprop{rr-recip-recip-pos}{\vnbstmt{\forall c \in \mathbb{R}.\; ((0 < c) \Rightarrow (\operatorname{recip}(\operatorname{recip}(c)) = c))}}
\vnbprop{deriv-right-inverse}{\vnbstmt{\begin{array}{@{}l@{}} \forall m, \mathit{gm}, \mathit{am}, v. \\ \quad \operatorname{is\text{-}diff\text{-}at}(m, \mathit{am}, v) \wedge \\ \quad \neg (v = 0) \wedge \\ \quad \mathit{gm} \in (\mathbb{R} \to \mathbb{R}) \wedge \\ \quad \forall z \in \mathbb{R}.\; (m(\operatorname{gm}(z)) = z) \wedge \\ \quad \operatorname{gm}(m(\mathit{am})) = \mathit{am} \wedge \\ \quad \operatorname{is\text{-}continuous\text{-}at}(\operatorname{rr\text{-}ms}, \operatorname{rr\text{-}ms}, \mathit{gm}, m(\mathit{am})) \Rightarrow \\ \quad \quad \operatorname{is\text{-}diff\text{-}at}(\mathit{gm}, m(\mathit{am}), \operatorname{recip}(v)) \end{array}}}
\vnbprop{r-exp-diff-at-log}{\vnbstmt{\begin{array}{@{}l@{}} \forall c, w. \\ \quad ((c \in \mathbb{R}) \wedge ((w \in \mathbb{R}) \wedge ((0 < c) \wedge (w = \log\left(c\right))))) \Rightarrow  \\ \quad \quad \operatorname{is\text{-}diff\text{-}at}((\lambda y \in \mathbb{R}.\; \exp\left(y\right)), w, c) \end{array}}}
\vnbprop{r-exp-deriv}{\vnbstmt{\forall t \in \mathbb{R}.\; \operatorname{is\text{-}diff\text{-}at}((\lambda y \in \mathbb{R}.\; \exp\left(y\right)), t, \exp\left(t\right))}}

\vnbbookfile{theorem-library/ca51-exhaustion}{theorem-library/ca51-exhaustion}
\noindent{\small ca51-exhaustion.scm -- Lemma 5.1 of the complex analysis notes: a separable locally compact metric space has an exhaustive sequence of compact sets (CA-51, 2026-10-26). Destination: theorem-library/ca51-exhaustion.scm, library analysis.}\par
\vnbprop{closed-ball-in-closed-ball}{\vnbstmt{\begin{array}{@{}l@{}} \forall s.; \forall a, b \in \operatorname{pts}(s), c, d \in \mathbb{R}. \\ \quad (\operatorname{is\text{-}metric\text{-}space}(s) \wedge ((\operatorname{dist}(s)(a, b) + d) \leq c)) \Rightarrow  \\ \quad \quad \operatorname{closed\text{-}ball}(s, b, d) \subseteq \operatorname{closed\text{-}ball}(s, a, c) \end{array}}}
\vnbprop{nn-recip-succ-gap-pos}{\vnbstmt{\forall n \in \mathbb{N}.\; \operatorname{pos\text{-}rr}((\operatorname{recip}((n + 1)) - \operatorname{recip}(((n + 1) + 1))))}}
\vnbprop{lx5-least-radius}{\vnbstmt{\begin{array}{@{}l@{}} \forall s.; \forall q \in (\mathbb{N} \to \operatorname{pts}(s)). \\ \quad \operatorname{is\text{-}locally\text{-}compact}(s) \Rightarrow  \\ \quad \quad \exists r \in (\mathbb{N} \to \mathbb{N}). \forall k \in \mathbb{N}. \\ \quad \quad \quad \operatorname{is\text{-}compact}(\operatorname{subspace\text{-}ms}(s, \operatorname{closed\text{-}ball}(s, q(k), \operatorname{recip}((r(k) + 1))))) \wedge \\ \quad \quad \quad \forall j \in \mathbb{N}. \\ \quad \quad \quad \quad \operatorname{is\text{-}compact}(\operatorname{subspace\text{-}ms}(s, \operatorname{closed\text{-}ball}(s, q(k), \operatorname{recip}((j + 1))))) \Rightarrow (r(k) \leq j) \end{array}}}
\vnbprop{lx5-pieces}{\vnbstmt{\begin{array}{@{}l@{}} \forall s.; \forall q \in (\mathbb{N} \to \operatorname{pts}(s)), r \in (\mathbb{N} \to \mathbb{N}). \\ \quad \operatorname{is\text{-}metric\text{-}space}(s) \wedge \\ \quad \forall k \in \mathbb{N}. \\ \quad \quad \operatorname{is\text{-}compact}(\operatorname{subspace\text{-}ms}(s, \operatorname{closed\text{-}ball}(s, q(k), \operatorname{recip}((r(k) + 1))))) \Rightarrow \\ \quad \quad \exists x \in (\mathbb{N} \to \operatorname{power}(\operatorname{pts}(s))). \\ \quad \quad \quad \forall n \in \mathbb{N}.\; \operatorname{is\text{-}compact}(\operatorname{subspace\text{-}ms}(s, x(n))) \wedge \\ \quad \quad \quad \forall n \in \mathbb{N}, y. \\ \quad \quad \quad \quad y \in x(n) \\ \quad \quad \quad \quad \Leftrightarrow\; ((y \in \operatorname{pts}(s)) \wedge \exists i \in \mathbb{N}.\; ((i \leq n) \wedge (\operatorname{dist}(s)(q(i), y) \leq (\operatorname{recip}((r(i) + 1)) - \operatorname{recip}((n + 1)))))) \end{array}}}
\vnbprop{locally-compact-separable-exhaustion}{\vnbstmt{\begin{array}{@{}l@{}} \forall s. \\ \quad (\operatorname{is\text{-}locally\text{-}compact}(s) \wedge \operatorname{is\text{-}separable}(s)) \Rightarrow  \\ \quad \quad \exists x.\; \operatorname{is\text{-}compact\text{-}exhaustion}(s, x) \end{array}}}
\vnbprop{compact-subset-of-subspace-iff}{\vnbstmt{\begin{array}{@{}l@{}} \forall s, u, a. \\ \quad (\operatorname{is\text{-}metric\text{-}space}(s) \wedge ((u \subseteq \operatorname{pts}(s)) \wedge (a \subseteq u))) \Rightarrow  \\ \quad \quad \operatorname{is\text{-}compact}(\operatorname{subspace\text{-}ms}(\operatorname{subspace\text{-}ms}(s, u), a)) \Leftrightarrow \operatorname{is\text{-}compact}(\operatorname{subspace\text{-}ms}(s, a)) \end{array}}}
\vnbprop{subspace-closed-ball-eq}{\vnbstmt{\begin{array}{@{}l@{}} \forall s, u.; \forall a \in u, b. \\ \quad (\operatorname{is\text{-}metric\text{-}space}(s) \wedge ((u \subseteq \operatorname{pts}(s)) \wedge (\operatorname{closed\text{-}ball}(s, a, b) \subseteq u))) \Rightarrow  \\ \quad \quad \operatorname{closed\text{-}ball}(\operatorname{subspace\text{-}ms}(s, u), a, b) = \operatorname{closed\text{-}ball}(s, a, b) \end{array}}}
\vnbprop{open-subspace-locally-compact}{\vnbstmt{\begin{array}{@{}l@{}} \forall s, u. \\ \quad (\operatorname{is\text{-}locally\text{-}compact}(s) \wedge \operatorname{is\text{-}open}(s, u)) \Rightarrow  \\ \quad \quad \operatorname{is\text{-}locally\text{-}compact}(\operatorname{subspace\text{-}ms}(s, u)) \end{array}}}
\vnbprop{dense-seq-separable}{\vnbstmt{\forall s, q.\; (\operatorname{is\text{-}dense\text{-}seq}(s, q) \Rightarrow \operatorname{is\text{-}separable}(s))}}
\vnbprop{separable-of-full-subspace}{\vnbstmt{\begin{array}{@{}l@{}} \forall s. \\ \quad (\operatorname{is\text{-}metric\text{-}space}(s) \wedge \operatorname{is\text{-}separable}(\operatorname{subspace\text{-}ms}(s, \operatorname{pts}(s)))) \Rightarrow  \\ \quad \quad \operatorname{is\text{-}separable}(s) \end{array}}}

\vnbbookfile{theorem-library/rpow-star}{theorem-library/rpow-star}
\noindent{\small theorem-library/rpow-star.scm -- THE REAL POWER, DEFINED, with the base-zero conventions carried:}\par
\par\medskip\noindent\textbf{Definition} (\texttt{rpow-star}).\;${\operatorname{x\_}}^{\operatorname{s\_}} :=$ $\left(\text{if }(0 < \operatorname{x\_})\text{ then }\exp\left((\operatorname{s\_} \cdot \log\left(\operatorname{x\_}\right))\right)\text{ else }\left(\text{if }(0 < \operatorname{s\_})\text{ then }0\text{ else }1\right)\right)$
\vnbprop{rpow-star-unfold}{\vnbstmt{\forall x, s.\; ({x}^{s} \simeq \left(\text{if }(0 < x)\text{ then }\exp\left((s \cdot \log\left(x\right))\right)\text{ else }\left(\text{if }(0 < s)\text{ then }0\text{ else }1\right)\right))}}
\vnbprop{rpow-star-value}{\vnbstmt{\forall x.; \forall s \in \mathbb{R}.\; ((0 < x) \Rightarrow ({x}^{s} = \exp\left((s \cdot \log\left(x\right))\right)))}}
\vnbprop{rpow-star-zero-base}{\vnbstmt{\forall s \in \mathbb{R}.\; ((0 < s) \Rightarrow ({0}^{s} = 0))}}
\vnbprop{rpow-star-zero-zero}{\vnbstmt{{0}^{0} = 1}}
\vnbprop{rpow-star-in-rr}{\vnbstmt{\forall x.; \forall s \in \mathbb{R}.\; ((0 < x) \Rightarrow ({x}^{s} \in \mathbb{R}))}}
\vnbprop{rpow-star-pos}{\vnbstmt{\forall x.; \forall s \in \mathbb{R}.\; ((0 < x) \Rightarrow (0 < {x}^{s}))}}
\vnbprop{log-rpow-star}{\vnbstmt{\forall x.; \forall s \in \mathbb{R}.\; ((0 < x) \Rightarrow (\log\left({x}^{s}\right) = (s \cdot \log\left(x\right))))}}
\vnbprop{rpow-star-zero}{\vnbstmt{\forall x \in \mathbb{R}.\; ((0 < x) \Rightarrow ({x}^{0} = 1))}}
\vnbprop{rpow-star-one}{\vnbstmt{\forall x \in \mathbb{R}.\; ((0 < x) \Rightarrow ({x}^{1} = x))}}
\vnbprop{rpow-star-add}{\vnbstmt{\begin{array}{@{}l@{}} \forall x, s, t. \\ \quad ((x \in \mathbb{R}) \wedge ((s \in \mathbb{R}) \wedge ((t \in \mathbb{R}) \wedge (0 < x)))) \Rightarrow  \\ \quad \quad {x}^{(s + t)} = ({x}^{s} \cdot {x}^{t}) \end{array}}}
\vnbprop{rpow-star-mul-base}{\vnbstmt{\begin{array}{@{}l@{}} \forall x, y, s. \\ \quad ((x \in \mathbb{R}) \wedge ((y \in \mathbb{R}) \wedge ((s \in \mathbb{R}) \wedge ((0 < x) \wedge (0 < y))))) \Rightarrow  \\ \quad \quad {(x \cdot y)}^{s} = ({x}^{s} \cdot {y}^{s}) \end{array}}}
\vnbprop{rpow-star-pow}{\vnbstmt{\begin{array}{@{}l@{}} \forall x, s, t. \\ \quad ((x \in \mathbb{R}) \wedge ((s \in \mathbb{R}) \wedge ((t \in \mathbb{R}) \wedge (0 < x)))) \Rightarrow  \\ \quad \quad {{x}^{s}}^{t} = {x}^{(s \cdot t)} \end{array}}}
\vnbprop{rpow-star-neg}{\vnbstmt{\begin{array}{@{}l@{}} \forall x, s. \\ \quad ((x \in \mathbb{R}) \wedge ((s \in \mathbb{R}) \wedge (0 < x))) \Rightarrow  \\ \quad \quad {x}^{-s} = {{x}^{s}}^{\ast\!-1} \end{array}}}
\vnbprop{rpow-star-mono-exp}{\vnbstmt{\begin{array}{@{}l@{}} \forall x, s, t. \\ \quad ((x \in \mathbb{R}) \wedge ((s \in \mathbb{R}) \wedge ((t \in \mathbb{R}) \wedge ((1 < x) \wedge (s < t))))) \Rightarrow  \\ \quad \quad {x}^{s} < {x}^{t} \end{array}}}
\vnbprop{rpow-star-mono-base}{\vnbstmt{\begin{array}{@{}l@{}} \forall x, y, s. \\ \quad ((x \in \mathbb{R}) \wedge ((y \in \mathbb{R}) \wedge ((s \in \mathbb{R}) \wedge ((0 < x) \wedge ((x < y) \wedge (0 < s)))))) \Rightarrow  \\ \quad \quad {x}^{s} < {y}^{s} \end{array}}}

\vnbbookfile{structure-library/ccr1109a3-laurent-vocab}{structure-library/ccr1109a3-laurent-vocab}
\noindent{\small ccr1109a3-laurent-vocab.scm -- wave of 2026-11-09/10 (agent CCR-1109-A3), wired 2026-11-10 from scratchpad/ccr1109a3/ccr1109a3-laurent-vocab.scm. Library: analysis.}\par
\par\medskip\noindent\textbf{Definition} (\texttt{laurent-abs}).\;$\operatorname{laurent\text{-}abs}(\operatorname{lrk\_}, \operatorname{lrr\_}, \operatorname{lrf\_}) :=$
\begin{equation*}
\left(\text{if }(\operatorname{lrf\_} = \operatorname{laurent\text{-}zero}(\operatorname{lrk\_}))\text{ then }0\text{ else }{\operatorname{lrr\_}}^{\operatorname{laurent\text{-}ord}(\operatorname{lrk\_}, \operatorname{lrf\_})}\right)
\end{equation*}

\par\medskip\noindent\textbf{Definition} (\texttt{laurent-mul}).\;$\operatorname{laurent\text{-}mul}(\operatorname{lrk\_}, \operatorname{lrf\_}, \operatorname{lrg\_}) :=$ $\operatorname{monalg\text{-}mul}(\operatorname{lrk\_}, \operatorname{zz\text{-}add\text{-}monoid}, \operatorname{lrf\_}, \operatorname{lrg\_})$
\par\medskip\noindent\textbf{Definition} (\texttt{zz-add-monoid}).\;
\begin{equation*}
\operatorname{zz\text{-}add\text{-}monoid} = {[}\mathbb{Z}, (\lambda {[}\operatorname{x\_}, \operatorname{y\_}] \in (\mathbb{Z} \times \mathbb{Z}).\; +(\operatorname{x\_}, \operatorname{y\_})), 0]
\end{equation*}

\par\medskip\noindent\textbf{Definition} (\texttt{laurent-supp}).\;$\operatorname{laurent\text{-}supp}(\operatorname{lrk\_}, \operatorname{lrf\_}) :=$ $\{\operatorname{lrj\_} \in \mathbb{Z} : \neg (\operatorname{lrf\_}(\operatorname{lrj\_}) = \operatorname{zero}(\operatorname{lrk\_}))\}$
\par\medskip\noindent\textbf{Definition} (\texttt{laurent-trunc}).\;$\operatorname{laurent\text{-}trunc}(\operatorname{lrk\_}, \operatorname{lrf\_}, \operatorname{lrt\_}) :=$
\begin{equation*}
(\lambda \operatorname{lrj\_} \in \mathbb{Z}.\; \left(\text{if }(\operatorname{lrj\_} \leq \operatorname{lrt\_})\text{ then }\operatorname{lrf\_}(\operatorname{lrj\_})\text{ else }\operatorname{zero}(\operatorname{lrk\_})\right))
\end{equation*}

\par\medskip\noindent\textbf{Definition} (\texttt{is-laurent}).\;For all $\operatorname{lrk\_}$ and $\operatorname{lrf\_}$, 
\begin{equation*}
\begin{array}{@{}l@{}} \operatorname{is\text{-}laurent}(\operatorname{lrk\_}, \operatorname{lrf\_}) \\ \Leftrightarrow\; \operatorname{lrf\_} \in (\mathbb{Z} \to \operatorname{carr}(\operatorname{lrk\_})) \wedge \\ \exists \operatorname{lrb\_} \in \mathbb{Z}. \forall \operatorname{lrj\_} \in \mathbb{Z}. \\ \quad (\operatorname{lrj\_} < \operatorname{lrb\_}) \Rightarrow (\operatorname{lrf\_}(\operatorname{lrj\_}) = \operatorname{zero}(\operatorname{lrk\_})) \end{array}
\end{equation*}

\par\medskip\noindent\textbf{Definition} (\texttt{laurent-set}).\;$\operatorname{laurent\text{-}set}(\operatorname{lrk\_}) :=$ $\{\operatorname{lrf\_} \in (\mathbb{Z} \to \operatorname{carr}(\operatorname{lrk\_})) : \operatorname{is\text{-}laurent}(\operatorname{lrk\_}, \operatorname{lrf\_})\}$
\par\medskip\noindent\textbf{Definition} (\texttt{laurent-add}).\;$\operatorname{laurent\text{-}add}(\operatorname{lrk\_}, \operatorname{lrf\_}, \operatorname{lrg\_}) :=$ $(\lambda \operatorname{lrj\_} \in \mathbb{Z}.\; \operatorname{add}(\operatorname{lrk\_})(\operatorname{lrf\_}(\operatorname{lrj\_}), \operatorname{lrg\_}(\operatorname{lrj\_})))$
\par\medskip\noindent\textbf{Definition} (\texttt{laurent-neg}).\;$\operatorname{laurent\text{-}neg}(\operatorname{lrk\_}, \operatorname{lrf\_}) :=$ $(\lambda \operatorname{lrj\_} \in \mathbb{Z}.\; \operatorname{neg}(\operatorname{lrk\_})(\operatorname{lrf\_}(\operatorname{lrj\_})))$
\par\medskip\noindent\textbf{Definition} (\texttt{laurent-zero}).\;$\operatorname{laurent\text{-}zero}(\operatorname{lrk\_}) :=$ $(\lambda \operatorname{lrj\_} \in \mathbb{Z}.\; \operatorname{zero}(\operatorname{lrk\_}))$
\par\medskip\noindent\textbf{Definition} (\texttt{laurent-one}).\;$\operatorname{laurent\text{-}one}(\operatorname{lrk\_}) :=$ $(\lambda \operatorname{lrj\_} \in \mathbb{Z}.\; \left(\text{if }(\operatorname{lrj\_} = 0)\text{ then }\operatorname{one}(\operatorname{lrk\_})\text{ else }\operatorname{zero}(\operatorname{lrk\_})\right))$
\par\medskip\noindent\textbf{Definition} (\texttt{laurent-ring}).\;$\operatorname{laurent\text{-}ring}(\operatorname{lrk\_}) :=$
\begin{equation*}
{[}\operatorname{laurent\text{-}set}(\operatorname{lrk\_}), (\lambda {[}\operatorname{lrf\_}, \operatorname{lrg\_}] \in (\operatorname{laurent\text{-}set}(\operatorname{lrk\_}) \times \operatorname{laurent\text{-}set}(\operatorname{lrk\_})).\; \operatorname{laurent\text{-}add}(\operatorname{lrk\_}, \operatorname{lrf\_}, \operatorname{lrg\_})), (\lambda {[}\operatorname{lrf\_}, \operatorname{lrg\_}] \in (\operatorname{laurent\text{-}set}(\operatorname{lrk\_}) \times \operatorname{laurent\text{-}set}(\operatorname{lrk\_})).\; \operatorname{laurent\text{-}mul}(\operatorname{lrk\_}, \operatorname{lrf\_}, \operatorname{lrg\_})), (\lambda \operatorname{lrf\_} \in \operatorname{laurent\text{-}set}(\operatorname{lrk\_}).\; \operatorname{laurent\text{-}neg}(\operatorname{lrk\_}, \operatorname{lrf\_})), \operatorname{laurent\text{-}zero}(\operatorname{lrk\_}), \operatorname{laurent\text{-}one}(\operatorname{lrk\_})]
\end{equation*}

\par\medskip\noindent\textbf{Definition} (\texttt{laurent-ord}).\;$\operatorname{laurent\text{-}ord}(\operatorname{lrk\_}, \operatorname{lrf\_}) :=$
\begin{equation*}
(\iota\, \operatorname{lro\_}.\;((\operatorname{lro\_} \in \mathbb{Z}) \wedge (\neg (\operatorname{lrf\_}(\operatorname{lro\_}) = \operatorname{zero}(\operatorname{lrk\_})) \wedge \forall \operatorname{lrj\_} \in \mathbb{Z}.\; ((\operatorname{lrj\_} < \operatorname{lro\_}) \Rightarrow (\operatorname{lrf\_}(\operatorname{lrj\_}) = \operatorname{zero}(\operatorname{lrk\_}))))))
\end{equation*}

\par\medskip\noindent\textbf{Definition} (\texttt{laurent-mono}).\;$\operatorname{laurent\text{-}mono}(\operatorname{lrk\_}, \operatorname{lrc\_}, \operatorname{lrt\_}) :=$ $(\lambda \operatorname{lrj\_} \in \mathbb{Z}.\; \left(\text{if }(\operatorname{lrj\_} = \operatorname{lrt\_})\text{ then }\operatorname{lrc\_}\text{ else }\operatorname{zero}(\operatorname{lrk\_})\right))$
\par\medskip\noindent\textbf{Definition} (\texttt{laurent-inv-pre}).\;For all $\operatorname{lrk\_}$, $\operatorname{lrf\_}$, $\operatorname{lro\_}$ and $\operatorname{lra\_}$, 
\begin{equation*}
(\operatorname{laurent\text{-}inv\text{-}pre}(\operatorname{lrk\_}, \operatorname{lrf\_}, \operatorname{lro\_}, \operatorname{lra\_}, 0) \simeq \operatorname{laurent\text{-}mono}(\operatorname{lrk\_}, \operatorname{lra\_}, -(\operatorname{lro\_})))
\end{equation*}
\par
Let $\operatorname{lrk\_}$, $\operatorname{lrf\_}$, $\operatorname{lro\_}$ and $\operatorname{lra\_}$ be arbitrary. Let $\operatorname{lrm\_} \in \mathbb{N}$. Then
\begin{equation*}
(\operatorname{laurent\text{-}inv\text{-}pre}(\operatorname{lrk\_}, \operatorname{lrf\_}, \operatorname{lro\_}, \operatorname{lra\_}, \operatorname{succ}(\operatorname{lrm\_})) \simeq \operatorname{laurent\text{-}add}(\operatorname{lrk\_}, \operatorname{laurent\text{-}inv\text{-}pre}(\operatorname{lrk\_}, \operatorname{lrf\_}, \operatorname{lro\_}, \operatorname{lra\_}, \operatorname{lrm\_}), \operatorname{laurent\text{-}mono}(\operatorname{lrk\_}, \operatorname{neg}(\operatorname{lrk\_})(\operatorname{mul}(\operatorname{lrk\_})(\operatorname{lra\_}, \operatorname{laurent\text{-}mul}(\operatorname{lrk\_}, \operatorname{lrf\_}, \operatorname{laurent\text{-}inv\text{-}pre}(\operatorname{lrk\_}, \operatorname{lrf\_}, \operatorname{lro\_}, \operatorname{lra\_}, \operatorname{lrm\_}))(+(\operatorname{lrm\_}, 1)))), +(+(\operatorname{lrm\_}, 1), -(\operatorname{lro\_})))))
\end{equation*}

\par\medskip\noindent\textbf{Definition} (\texttt{laurent-inv-ser}).\;$\operatorname{laurent\text{-}inv\text{-}ser}(\operatorname{lrk\_}, \operatorname{lrf\_}, \operatorname{lro\_}, \operatorname{lra\_}) :=$
\begin{equation*}
\begin{array}{@{}l@{}} (\lambda \operatorname{lrj\_} \in \mathbb{Z}.\; \\ \quad \left(\text{if }(-\operatorname{lro\_} \leq \operatorname{lrj\_})\text{ then }\operatorname{laurent\text{-}inv\text{-}pre}(\operatorname{lrk\_}, \operatorname{lrf\_}, \operatorname{lro\_}, \operatorname{lra\_}, (\operatorname{lrj\_} + \operatorname{lro\_}))(\operatorname{lrj\_})\text{ else }\operatorname{zero}(\operatorname{lrk\_})\right)) \end{array}
\end{equation*}

\par\medskip\noindent\textbf{Definition} (\texttt{laurent-nf}).\;$\operatorname{laurent\text{-}nf}(\operatorname{lrk\_}, \operatorname{lrr\_}) :=$
\begin{equation*}
{[}\operatorname{carr}(\operatorname{laurent\text{-}ring}(\operatorname{lrk\_})), \operatorname{add}(\operatorname{laurent\text{-}ring}(\operatorname{lrk\_})), \operatorname{mul}(\operatorname{laurent\text{-}ring}(\operatorname{lrk\_})), \operatorname{neg}(\operatorname{laurent\text{-}ring}(\operatorname{lrk\_})), \operatorname{zero}(\operatorname{laurent\text{-}ring}(\operatorname{lrk\_})), \operatorname{one}(\operatorname{laurent\text{-}ring}(\operatorname{lrk\_})), (\lambda \operatorname{lrf\_} \in \operatorname{carr}(\operatorname{laurent\text{-}ring}(\operatorname{lrk\_})).\; \operatorname{laurent\text{-}abs}(\operatorname{lrk\_}, \operatorname{lrr\_}, \operatorname{lrf\_}))]
\end{equation*}

\vnbdef{zz-add-monoid@carr}{\vnbstmt{(\operatorname{carr}(\operatorname{zz\text{-}add\text{-}monoid}) \simeq \mathbb{Z})}}
\vnbdef{zz-add-monoid@iden}{\vnbstmt{(\operatorname{iden}(\operatorname{zz\text{-}add\text{-}monoid}) \simeq 0)}}
\vnbdef{zz-add-monoid@opr}{\vnbstmt{(\operatorname{opr}(\operatorname{zz\text{-}add\text{-}monoid}) \simeq (\lambda {[}\operatorname{x\_}, \operatorname{y\_}] \in (\mathbb{Z} \times \mathbb{Z}).\; (\operatorname{x\_} + \operatorname{y\_})))}}

\vnbbookfile{structure-library/ccr1109a4-laurent-lim-vocab}{structure-library/ccr1109a4-laurent-lim-vocab}
\noindent{\small ccr1109a4-laurent-lim-vocab.scm -- wave of 2026-11-09/10 (agent CCR-1109-A4), wired 2026-11-10 from scratchpad/ccr1109a4/ccr1109a4-laurent-lim-vocab.scm. Library: analysis.}\par
\par\medskip\noindent\textbf{Definition} (\texttt{laurent-limc}).\;$\operatorname{laurent\text{-}limc}(\operatorname{llk\_}, \operatorname{lls\_}, \operatorname{llj\_}) :=$
\begin{equation*}
(\iota\, \operatorname{llc\_}.\;((\operatorname{llc\_} \in \operatorname{carr}(\operatorname{llk\_})) \wedge \exists \operatorname{llm\_} \in \mathbb{N}.\; \forall \operatorname{lln\_} \in \mathbb{N}.\; ((\operatorname{llm\_} \leq \operatorname{lln\_}) \Rightarrow (\operatorname{lls\_}(\operatorname{lln\_})(\operatorname{llj\_}) = \operatorname{llc\_}))))
\end{equation*}

\par\medskip\noindent\textbf{Definition} (\texttt{laurent-lim}).\;$\operatorname{laurent\text{-}lim}(\operatorname{llk\_}, \operatorname{lls\_}) :=$ $(\lambda \operatorname{llj\_} \in \mathbb{Z}.\; \operatorname{laurent\text{-}limc}(\operatorname{llk\_}, \operatorname{lls\_}, \operatorname{llj\_}))$
\par\medskip\noindent\textbf{Definition} (\texttt{laurent-ext}).\;$\operatorname{laurent\text{-}ext}(\operatorname{llk\_}, \operatorname{llm\_}) :=$
\begin{equation*}
\begin{array}{@{}l@{}} (\lambda \operatorname{llw\_} \in (\operatorname{laurent\text{-}polys}(\operatorname{llk\_}, \operatorname{llm\_}) \times \operatorname{carr}(\operatorname{llk\_})).\; \\ \quad \operatorname{laurent\text{-}add}(\operatorname{llk\_}, \operatorname{nth}(1, \operatorname{llw\_}), \operatorname{laurent\text{-}mono}(\operatorname{llk\_}, \operatorname{nth}(2, \operatorname{llw\_}), \operatorname{llm\_}))) \end{array}
\end{equation*}

\par\medskip\noindent\textbf{Definition} (\texttt{laurent-polys}).\;$\operatorname{laurent\text{-}polys}(\operatorname{llk\_}, \operatorname{llm\_}) :=$
\begin{equation*}
\{\operatorname{llg\_} \in (\mathbb{Z} \to \operatorname{carr}(\operatorname{llk\_})) : (\forall \operatorname{llj\_} \in \mathbb{Z}.\; ((\operatorname{llj\_} < 0) \Rightarrow (\operatorname{llg\_}(\operatorname{llj\_}) = \operatorname{zero}(\operatorname{llk\_}))) \wedge \forall \operatorname{llj\_} \in \mathbb{Z}.\; ((\operatorname{llm\_} \leq \operatorname{llj\_}) \Rightarrow (\operatorname{llg\_}(\operatorname{llj\_}) = \operatorname{zero}(\operatorname{llk\_}))))\}
\end{equation*}


\vnbbookfile{structure-library/ccr1109a5-vocab}{structure-library/ccr1109a5-vocab}
\noindent{\small ccr1109a5-vocab.scm -- wave of 2026-11-09/10 (agent CCR-1109-A5), wired 2026-11-10 from scratchpad/ccr1109a5/ccr1109a5-vocab.scm. Library: analysis.}\par
\par\medskip\noindent\textbf{Definition} (\texttt{triv-nf}).\;$\operatorname{triv\text{-}nf}(\operatorname{a5k\_}) :=$
\begin{equation*}
{[}\operatorname{carr}(\operatorname{a5k\_}), \operatorname{add}(\operatorname{a5k\_}), \operatorname{mul}(\operatorname{a5k\_}), \operatorname{neg}(\operatorname{a5k\_}), \operatorname{zero}(\operatorname{a5k\_}), \operatorname{one}(\operatorname{a5k\_}), (\lambda \operatorname{a5x\_} \in \operatorname{carr}(\operatorname{a5k\_}).\; \operatorname{triv\text{-}abs}(\operatorname{a5k\_}, \operatorname{a5x\_}))]
\end{equation*}

\par\medskip\noindent\textbf{Definition} (\texttt{laurent-pairing}).\;$\operatorname{laurent\text{-}pairing}(\operatorname{a5k\_}, \operatorname{a5c\_}) :=$
\begin{equation*}
\begin{array}{@{}l@{}} (\lambda {[}\operatorname{a5g\_}, \operatorname{a5f\_}] \in (\operatorname{laurent\text{-}set}(\operatorname{a5k\_}) \times \operatorname{laurent\text{-}set}(\operatorname{a5k\_})).\; \\ \quad \operatorname{a5c\_}(\operatorname{laurent\text{-}mul}(\operatorname{a5k\_}, \operatorname{a5f\_}, \operatorname{a5g\_})(0))) \end{array}
\end{equation*}

\par\medskip\noindent\textbf{Definition} (\texttt{triv-abs}).\;$\operatorname{triv\text{-}abs}(\operatorname{a5k\_}, \operatorname{a5x\_}) :=$ $\left(\text{if }(\operatorname{a5x\_} = \operatorname{zero}(\operatorname{a5k\_}))\text{ then }0\text{ else }1\right)$

\vnbbookfile{structure-library/mb12-vocab}{structure-library/mb12-vocab}
\noindent{\small mb12-vocab.scm -- vocabulary for construction-of-measures.pdf Appendix A.3: the snowflake of a metric (A.9) and the prefix metrics on X\textasciicircum{}N (A.11-A.16). DEFINITIONS ONLY.  Agent ISOM (brief M-B1B2), 2026-10-26 (NEW DEFINITIONS, the user's to confirm).}\par
\par\medskip\noindent\textbf{Definition} (\texttt{snowflake}).\;$\operatorname{snowflake}(\operatorname{snfs\_}, \operatorname{snfe\_}) :=$
\begin{equation*}
{[}\operatorname{pts}(\operatorname{snfs\_}), (\lambda {[}\operatorname{snfu\_}, \operatorname{snfv\_}] \in (\operatorname{pts}(\operatorname{snfs\_}) \times \operatorname{pts}(\operatorname{snfs\_})).\; {\operatorname{dist}(\operatorname{snfs\_})(\operatorname{snfu\_}, \operatorname{snfv\_})}^{\operatorname{snfe\_}})]
\end{equation*}

\par\medskip\noindent\textbf{Definition} (\texttt{discrete-ms}).\;$\operatorname{discrete\text{-}ms}(\operatorname{dsx\_}) :=$
\begin{equation*}
{[}\operatorname{dsx\_}, (\lambda {[}\operatorname{dsu\_}, \operatorname{dsv\_}] \in (\operatorname{dsx\_} \times \operatorname{dsx\_}).\; \left(\text{if }(\operatorname{dsu\_} = \operatorname{dsv\_})\text{ then }0\text{ else }1\right))]
\end{equation*}

\par\medskip\noindent\textbf{Definition} (\texttt{seq-space-ms}).\;$\operatorname{seq\text{-}space\text{-}ms}(\operatorname{sqx\_}) :=$
\begin{equation*}
\operatorname{product\text{-}metric\text{-}w}((\lambda \operatorname{sqk\_} \in \mathbb{N}.\; \operatorname{discrete\text{-}ms}(\operatorname{sqx\_})), (\lambda \operatorname{sqk\_} \in \mathbb{N}.\; /(1, \operatorname{power}(2, \operatorname{sqk\_}))))
\end{equation*}

\par\medskip\noindent\textbf{Definition} (\texttt{coord-permute}).\;$\operatorname{coord\text{-}permute}(\operatorname{cpx\_}, \operatorname{cpg\_}) :=$
\begin{equation*}
(\lambda \operatorname{cpw\_} \in (\mathbb{N} \to \operatorname{cpx\_}).\; (\lambda \operatorname{cpk\_} \in \mathbb{N}.\; \operatorname{cpg\_}(\operatorname{cpk\_})(\operatorname{cpw\_}(\operatorname{cpk\_}))))
\end{equation*}

\par\medskip\noindent\textbf{Definition} (\texttt{is-ultrametric}).\;For every $\operatorname{ums\_}$, 
\begin{equation*}
\begin{array}{@{}l@{}} \operatorname{is\text{-}ultrametric}(\operatorname{ums\_}) \\ \Leftrightarrow\; \operatorname{is\text{-}metric\text{-}space}(\operatorname{ums\_}) \wedge \\ \forall \operatorname{umx\_}, \operatorname{umy\_}, \operatorname{umz\_} \in \operatorname{pts}(\operatorname{ums\_}). \\ \quad \operatorname{dist}(\operatorname{ums\_})(\operatorname{umx\_}, \operatorname{umy\_}) \leq \operatorname{max}(\operatorname{dist}(\operatorname{ums\_})(\operatorname{umx\_}, \operatorname{umz\_}), \operatorname{dist}(\operatorname{ums\_})(\operatorname{umz\_}, \operatorname{umy\_})) \end{array}
\end{equation*}

\par\medskip\noindent\textbf{Definition} (\texttt{prefix-prod}).\;For all $\operatorname{ppa\_}$ and $\operatorname{ppx\_}$,  $(\operatorname{prefix\text{-}prod}(\operatorname{ppa\_}, \operatorname{ppx\_}, 0) \simeq 1)$\par
Let $\operatorname{ppa\_}$ and $\operatorname{ppx\_}$ be arbitrary. Let $\operatorname{ppn\_} \in \mathbb{N}$. Then
\begin{equation*}
(\operatorname{prefix\text{-}prod}(\operatorname{ppa\_}, \operatorname{ppx\_}, \operatorname{succ}(\operatorname{ppn\_})) \simeq \cdot(\operatorname{prefix\text{-}prod}(\operatorname{ppa\_}, \operatorname{ppx\_}, \operatorname{ppn\_}), \operatorname{ppa\_}(\operatorname{ppx\_}(\operatorname{ppn\_}))))
\end{equation*}

\par\medskip\noindent\textbf{Definition} (\texttt{agree-length}).\;$\operatorname{agree\text{-}length}(\operatorname{alx\_}, \operatorname{aly\_}) :=$
\begin{equation*}
(\iota\, \operatorname{alk\_}.\;((\operatorname{alk\_} \in \mathbb{N}) \wedge (\neg (\operatorname{alx\_}(\operatorname{alk\_}) = \operatorname{aly\_}(\operatorname{alk\_})) \wedge \forall \operatorname{alj\_} \in \mathbb{N}.\; ((\operatorname{alj\_} < \operatorname{alk\_}) \Rightarrow (\operatorname{alx\_}(\operatorname{alj\_}) = \operatorname{aly\_}(\operatorname{alj\_}))))))
\end{equation*}

\par\medskip\noindent\textbf{Definition} (\texttt{prefix-dist}).\;$\operatorname{prefix\text{-}dist}(\operatorname{ppa\_}, \operatorname{pdx\_}, \operatorname{pdy\_}) :=$
\begin{equation*}
\left(\text{if }(\operatorname{pdx\_} = \operatorname{pdy\_})\text{ then }0\text{ else }\operatorname{recip}(\operatorname{prefix\text{-}prod}(\operatorname{ppa\_}, \operatorname{pdx\_}, \operatorname{agree\text{-}length}(\operatorname{pdx\_}, \operatorname{pdy\_})))\right)
\end{equation*}

\par\medskip\noindent\textbf{Definition} (\texttt{ultra-seq-ms}).\;$\operatorname{ultra\text{-}seq\text{-}ms}(\operatorname{usx\_}, \operatorname{usa\_}) :=$
\begin{equation*}
{[}(\mathbb{N} \to \operatorname{usx\_}), (\lambda {[}\operatorname{usu\_}, \operatorname{usv\_}] \in ((\mathbb{N} \to \operatorname{usx\_}) \times (\mathbb{N} \to \operatorname{usx\_})).\; \operatorname{prefix\text{-}dist}(\operatorname{usa\_}, \operatorname{usu\_}, \operatorname{usv\_}))]
\end{equation*}


\vnbbookfile{theorem-library/mb12-seq}{theorem-library/mb12-seq}
\noindent{\small mb12-seq.scm -- construction-of-measures.pdf Appendix A.3.3, metrics on X\textasciicircum{}N: Proposition A.11 (convergence for the metric (414) is coordinatewise eventual equality), and below.  Agent ISOM (brief M-B1B2), 2026-10-26.}\par
\vnbprop{discrete-ms-lam-type}{\vnbstmt{\begin{array}{@{}l@{}} \forall x \in \mathbf{Set}. \\ \quad (\lambda {[}\operatorname{dsu\_}, \operatorname{dsv\_}] \in (x \times x).\; \left(\text{if }(\operatorname{dsu\_} = \operatorname{dsv\_})\text{ then }0\text{ else }1\right)) \in ((x \times x) \to \mathbb{R}) \end{array}}}
\vnbprop{discrete-ms-pts}{\vnbstmt{\forall x \in \mathbf{Set}.\; (\operatorname{pts}(\operatorname{discrete\text{-}ms}(x)) \simeq x)}}
\vnbprop{discrete-ms-dist-fun}{\vnbstmt{\begin{array}{@{}l@{}} \forall x \in \mathbf{Set}. \\ \quad (\operatorname{dist}(\operatorname{discrete\text{-}ms}(x)) \simeq (\lambda {[}\operatorname{dsu\_}, \operatorname{dsv\_}] \in (x \times x).\; \left(\text{if }(\operatorname{dsu\_} = \operatorname{dsv\_})\text{ then }0\text{ else }1\right))) \end{array}}}
\vnbprop{discrete-ms-dist-values}{\vnbstmt{\begin{array}{@{}l@{}} \forall x \in \mathbf{Set}, u, v \in x. \\ \quad (u = v) \Rightarrow (\operatorname{dist}(\operatorname{discrete\text{-}ms}(x))(u, v) = 0) \wedge \\ \quad \neg (u = v) \Rightarrow (\operatorname{dist}(\operatorname{discrete\text{-}ms}(x))(u, v) = 1) \end{array}}}
\vnbprop{discrete-ms-dist-bounds}{\vnbstmt{\begin{array}{@{}l@{}} \forall x \in \mathbf{Set}, u, v \in x. \\ \quad \operatorname{dist}(\operatorname{discrete\text{-}ms}(x))(u, v) \in \mathbb{R} \wedge \\ \quad 0 \leq \operatorname{dist}(\operatorname{discrete\text{-}ms}(x))(u, v) \wedge \\ \quad \operatorname{dist}(\operatorname{discrete\text{-}ms}(x))(u, v) \leq 1 \end{array}}}
\vnbprop{discrete-ms-is-metric-space}{\vnbstmt{\forall x \in \mathbf{Set}.\; \operatorname{is\text{-}metric\text{-}space}(\operatorname{discrete\text{-}ms}(x))}}
\vnbprop{discrete-converges-iff}{\vnbstmt{\begin{array}{@{}l@{}} \forall x, f, l. \\ \quad ((x \in \mathbf{Set}) \wedge ((f \in (\mathbb{N} \to x)) \wedge (l \in x))) \Rightarrow  \\ \quad \quad \operatorname{converges\text{-}to}(\operatorname{discrete\text{-}ms}(x), f, l) \\ \quad \quad \Leftrightarrow\; \exists n \in \mathbb{N}.\; \forall i \in \mathbb{N}.\; ((n \leq i) \Rightarrow (f(i) = l)) \end{array}}}
\vnbprop{seq-weight-value}{\vnbstmt{\begin{array}{@{}l@{}} \forall k \in \mathbb{N}. \\ \quad (\lambda \mathit{sqk} \in \mathbb{N}.\; /(1, \operatorname{power}(2, \mathit{sqk})))(k) = \operatorname{recip}(\operatorname{power}(2, k)) \wedge \\ \quad \operatorname{recip}(\operatorname{power}(2, k)) \in \mathbb{R} \wedge \\ \quad 0 < \operatorname{recip}(\operatorname{power}(2, k)) \wedge \\ \quad \operatorname{recip}(\operatorname{power}(2, k)) = (2 \cdot (\lambda n \in \mathbb{N}.\; /(1, \operatorname{power}(2, (n + 1))))(k)) \end{array}}}
\vnbprop{seq-weight-summable}{\vnbstmt{\operatorname{summable\text{-}weight}((\lambda k \in \mathbb{N}.\; /(1, \operatorname{power}(2, k))))}}
\vnbprop{seq-factors-ms-sequence}{\vnbstmt{\forall x \in \mathbf{Set}.\; \operatorname{is\text{-}ms\text{-}sequence}((\lambda k \in \mathbb{N}.\; \operatorname{discrete\text{-}ms}(x)))}}
\vnbprop{seq-carrier-intro}{\vnbstmt{\begin{array}{@{}l@{}} \forall x \in \mathbf{Set}, s \in (\mathbb{N} \to x). \\ \quad s \in \operatorname{product\text{-}carrier}((\lambda k \in \mathbb{N}.\; \operatorname{discrete\text{-}ms}(x))) \end{array}}}
\vnbprop{seq-space-converges-iff}{\vnbstmt{\begin{array}{@{}l@{}} \forall x \in \mathbf{Set}, s \in (\mathbb{N} \to (\mathbb{N} \to x)), u \in (\mathbb{N} \to x). \\ \quad \operatorname{converges\text{-}to}(\operatorname{seq\text{-}space\text{-}ms}(x), s, u) \\ \quad \Leftrightarrow\; \forall k \in \mathbb{N}. \exists n \in \mathbb{N}. \forall i \in \mathbb{N}. \\ \quad \quad (n \leq i) \Rightarrow (s(i)(k) = u(k)) \end{array}}}
\vnbprop{swap-values}{\vnbstmt{\begin{array}{@{}l@{}} \forall x \in \mathbf{Set}, a, b, u \in x. \\ \quad (u = a) \Rightarrow ((\lambda w \in x.\; \left(\text{if }(w = a)\text{ then }b\text{ else }\left(\text{if }(w = b)\text{ then }a\text{ else }w\right)\right))(u) = b) \wedge \\ \quad (\neg (u = a) \wedge (u = b)) \Rightarrow  \\ \quad \quad (\lambda w \in x.\; \left(\text{if }(w = a)\text{ then }b\text{ else }\left(\text{if }(w = b)\text{ then }a\text{ else }w\right)\right))(u) = a \wedge \\ \quad (\neg (u = a) \wedge \neg (u = b)) \Rightarrow  \\ \quad \quad (\lambda w \in x.\; \left(\text{if }(w = a)\text{ then }b\text{ else }\left(\text{if }(w = b)\text{ then }a\text{ else }w\right)\right))(u) = u \end{array}}}
\vnbprop{swap-in-fun}{\vnbstmt{\begin{array}{@{}l@{}} \forall x \in \mathbf{Set}, a, b \in x. \\ \quad (\lambda w \in x.\; \left(\text{if }(w = a)\text{ then }b\text{ else }\left(\text{if }(w = b)\text{ then }a\text{ else }w\right)\right)) \in (x \to x) \end{array}}}
\vnbprop{swap-involution}{\vnbstmt{\begin{array}{@{}l@{}} \forall x \in \mathbf{Set}, a, b, u \in x. \\ \quad (\lambda w \in x.\; \left(\text{if }(w = a)\text{ then }b\text{ else }\left(\text{if }(w = b)\text{ then }a\text{ else }w\right)\right))((\lambda w \in x.\; \left(\text{if }(w = a)\text{ then }b\text{ else }\left(\text{if }(w = b)\text{ then }a\text{ else }w\right)\right))(u)) \\ \quad =\; u \end{array}}}
\vnbprop{swap-bijection}{\vnbstmt{\begin{array}{@{}l@{}} \forall x \in \mathbf{Set}, a, b \in x. \\ \quad (\lambda w \in x.\; \left(\text{if }(w = a)\text{ then }b\text{ else }\left(\text{if }(w = b)\text{ then }a\text{ else }w\right)\right)) \in \mathrm{Bij}(x, x) \end{array}}}
\vnbprop{coord-permute-transitive}{\vnbstmt{\begin{array}{@{}l@{}} \forall x \in \mathbf{Set}, s, a \in (\mathbb{N} \to x). \exists g \in (\mathbb{N} \to \mathrm{Bij}(x, x)). \\ \quad \operatorname{coord\text{-}permute}(x, g)(s) = a \end{array}}}

\vnbbookfile{theorem-library/mb12-ultra}{theorem-library/mb12-ultra}
\noindent{\small mb12-ultra.scm -- the weighted prefix ultrametric on X\textasciicircum{}N (construction-of-measures, Appendix A.3: (419), Propositions A.15 and A.16).  Helper prefix: mul-. Agent ISOM (brief M-B1B2), 2026-10-26.}\par
\vnbprop{prefix-prod-ge-one-ind}{\vnbstmt{\begin{array}{@{}l@{}} \forall n \in \mathbb{N}, x, a, s. \\ \quad ((a \in (x \to \mathbb{R})) \wedge (\forall w \in x.\; (1 < a(w)) \wedge (s \in (\mathbb{N} \to x)))) \Rightarrow  \\ \quad \quad (\operatorname{prefix\text{-}prod}(a, s, n) \in \mathbb{R}) \wedge (1 \leq \operatorname{prefix\text{-}prod}(a, s, n)) \end{array}}}
\vnbprop{prefix-prod-ge-one}{\vnbstmt{\begin{array}{@{}l@{}} \forall x, a, s.; \forall n \in \mathbb{N}. \\ \quad ((a \in (x \to \mathbb{R})) \wedge (\forall w \in x.\; (1 < a(w)) \wedge (s \in (\mathbb{N} \to x)))) \Rightarrow  \\ \quad \quad (\operatorname{prefix\text{-}prod}(a, s, n) \in \mathbb{R}) \wedge (1 \leq \operatorname{prefix\text{-}prod}(a, s, n)) \end{array}}}
\vnbprop{prefix-prod-mono}{\vnbstmt{\begin{array}{@{}l@{}} \forall x, a, s.; \forall m, n \in \mathbb{N}. \\ \quad ((a \in (x \to \mathbb{R})) \wedge (\forall w \in x.\; (1 < a(w)) \wedge ((s \in (\mathbb{N} \to x)) \wedge (m \leq n)))) \Rightarrow  \\ \quad \quad \operatorname{prefix\text{-}prod}(a, s, m) \leq \operatorname{prefix\text{-}prod}(a, s, n) \end{array}}}
\vnbprop{prefix-prod-agree-ind}{\vnbstmt{\begin{array}{@{}l@{}} \forall n \in \mathbb{N}, x, a, s, b. \\ \quad a \in (x \to \mathbb{R}) \wedge \\ \quad \forall w \in x.\; (1 < a(w)) \wedge \\ \quad s \in (\mathbb{N} \to x) \wedge \\ \quad b \in (\mathbb{N} \to x) \wedge \\ \quad \forall j \in \mathbb{N}.\; ((j < n) \Rightarrow (s(j) = b(j))) \Rightarrow \\ \quad \quad \operatorname{prefix\text{-}prod}(a, s, n) = \operatorname{prefix\text{-}prod}(a, b, n) \end{array}}}
\vnbprop{seqs-differ-somewhere}{\vnbstmt{\begin{array}{@{}l@{}} \forall x, s, a. \\ \quad ((s \in (\mathbb{N} \to x)) \wedge ((a \in (\mathbb{N} \to x)) \wedge \neg (s = a))) \Rightarrow  \\ \quad \quad \exists k \in \mathbb{N}.\; \neg (s(k) = a(k)) \end{array}}}
\vnbprop{agree-length-prop}{\vnbstmt{\begin{array}{@{}l@{}} \forall x, s, a. \\ \quad ((s \in (\mathbb{N} \to x)) \wedge ((a \in (\mathbb{N} \to x)) \wedge \neg (s = a))) \Rightarrow  \\ \quad \quad \operatorname{agree\text{-}length}(s, a) \in \mathbb{N} \wedge \\ \quad \quad \neg (s(\operatorname{agree\text{-}length}(s, a)) = a(\operatorname{agree\text{-}length}(s, a))) \wedge \\ \quad \quad \forall j \in \mathbb{N}.\; ((j < \operatorname{agree\text{-}length}(s, a)) \Rightarrow (s(j) = a(j))) \end{array}}}
\vnbprop{prefix-dist-self}{\vnbstmt{\forall a, s.\; (\operatorname{prefix\text{-}dist}(a, s, s) = 0)}}
\vnbprop{prefix-dist-value}{\vnbstmt{\begin{array}{@{}l@{}} \forall x, a, s, b. \\ \quad a \in (x \to \mathbb{R}) \wedge \\ \quad \forall w \in x.\; (1 < a(w)) \wedge \\ \quad s \in (\mathbb{N} \to x) \wedge \\ \quad b \in (\mathbb{N} \to x) \wedge \\ \quad \neg (s = b) \Rightarrow \\ \quad \quad \operatorname{prefix\text{-}dist}(a, s, b) = \operatorname{recip}(\operatorname{prefix\text{-}prod}(a, s, \operatorname{agree\text{-}length}(s, b))) \wedge \\ \quad \quad \operatorname{recip}(\operatorname{prefix\text{-}prod}(a, s, \operatorname{agree\text{-}length}(s, b))) \in \mathbb{R} \wedge \\ \quad \quad 0 < \operatorname{recip}(\operatorname{prefix\text{-}prod}(a, s, \operatorname{agree\text{-}length}(s, b))) \end{array}}}
\vnbprop{prefix-dist-nonneg}{\vnbstmt{\begin{array}{@{}l@{}} \forall x, a, s, b. \\ \quad ((a \in (x \to \mathbb{R})) \wedge (\forall w \in x.\; (1 < a(w)) \wedge ((s \in (\mathbb{N} \to x)) \wedge (b \in (\mathbb{N} \to x))))) \Rightarrow  \\ \quad \quad (\operatorname{prefix\text{-}dist}(a, s, b) \in \mathbb{R}) \wedge (0 \leq \operatorname{prefix\text{-}dist}(a, s, b)) \end{array}}}
\vnbprop{agree-length-sym}{\vnbstmt{\begin{array}{@{}l@{}} \forall x, s, a. \\ \quad ((s \in (\mathbb{N} \to x)) \wedge ((a \in (\mathbb{N} \to x)) \wedge \neg (s = a))) \Rightarrow  \\ \quad \quad \operatorname{agree\text{-}length}(s, a) = \operatorname{agree\text{-}length}(a, s) \end{array}}}
\vnbprop{prefix-dist-sym}{\vnbstmt{\begin{array}{@{}l@{}} \forall x, a, s, b. \\ \quad ((a \in (x \to \mathbb{R})) \wedge (\forall w \in x.\; (1 < a(w)) \wedge ((s \in (\mathbb{N} \to x)) \wedge (b \in (\mathbb{N} \to x))))) \Rightarrow  \\ \quad \quad \operatorname{prefix\text{-}dist}(a, s, b) = \operatorname{prefix\text{-}dist}(a, b, s) \end{array}}}
\vnbprop{prefix-dist-ultra}{\vnbstmt{\begin{array}{@{}l@{}} \forall x, a, s, b, c. \\ \quad a \in (x \to \mathbb{R}) \wedge \\ \quad \forall w \in x.\; (1 < a(w)) \wedge \\ \quad s \in (\mathbb{N} \to x) \wedge \\ \quad b \in (\mathbb{N} \to x) \wedge \\ \quad c \in (\mathbb{N} \to x) \Rightarrow \\ \quad \quad \operatorname{prefix\text{-}dist}(a, s, b) \leq \operatorname{max}(\operatorname{prefix\text{-}dist}(a, s, c), \operatorname{prefix\text{-}dist}(a, c, b)) \end{array}}}
\vnbprop{ultra-seq-lam-type}{\vnbstmt{\begin{array}{@{}l@{}} \forall x, a. \\ \quad ((x \in \mathbf{Set}) \wedge ((a \in (x \to \mathbb{R})) \wedge \forall w \in x.\; (1 < a(w)))) \Rightarrow  \\ \quad \quad (\lambda {[}\operatorname{usu\_}, \operatorname{usv\_}] \in ((\mathbb{N} \to x) \times (\mathbb{N} \to x)).\; \operatorname{prefix\text{-}dist}(a, \operatorname{usu\_}, \operatorname{usv\_})) \\ \quad \quad \in\; (((\mathbb{N} \to x) \times (\mathbb{N} \to x)) \to \mathbb{R}) \end{array}}}
\vnbprop{ultra-seq-pts}{\vnbstmt{\begin{array}{@{}l@{}} \forall x, a. \\ \quad ((x \in \mathbf{Set}) \wedge ((a \in (x \to \mathbb{R})) \wedge \forall w \in x.\; (1 < a(w)))) \Rightarrow  \\ \quad \quad (\operatorname{pts}(\operatorname{ultra\text{-}seq\text{-}ms}(x, a)) \simeq (\mathbb{N} \to x)) \end{array}}}
\vnbprop{ultra-seq-dist-fun}{\vnbstmt{\begin{array}{@{}l@{}} \forall x, a. \\ \quad ((x \in \mathbf{Set}) \wedge ((a \in (x \to \mathbb{R})) \wedge \forall w \in x.\; (1 < a(w)))) \Rightarrow  \\ \quad \quad (\operatorname{dist}(\operatorname{ultra\text{-}seq\text{-}ms}(x, a)) \simeq (\lambda {[}\operatorname{usu\_}, \operatorname{usv\_}] \in ((\mathbb{N} \to x) \times (\mathbb{N} \to x)).\; \operatorname{prefix\text{-}dist}(a, \operatorname{usu\_}, \operatorname{usv\_}))) \end{array}}}
\vnbprop{ultra-seq-is-metric-space}{\vnbstmt{\begin{array}{@{}l@{}} \forall x, a. \\ \quad ((x \in \mathbf{Set}) \wedge ((a \in (x \to \mathbb{R})) \wedge \forall w \in x.\; (1 < a(w)))) \Rightarrow  \\ \quad \quad \operatorname{is\text{-}metric\text{-}space}(\operatorname{ultra\text{-}seq\text{-}ms}(x, a)) \end{array}}}
\vnbprop{ultra-seq-is-ultrametric}{\vnbstmt{\begin{array}{@{}l@{}} \forall x, a. \\ \quad ((x \in \mathbf{Set}) \wedge ((a \in (x \to \mathbb{R})) \wedge \forall w \in x.\; (1 < a(w)))) \Rightarrow  \\ \quad \quad \operatorname{is\text{-}ultrametric}(\operatorname{ultra\text{-}seq\text{-}ms}(x, a)) \end{array}}}
\vnbprop{prefix-eventually-agrees}{\vnbstmt{\begin{array}{@{}l@{}} \forall n \in \mathbb{N}, x, s, u. \\ \quad s \in (\mathbb{N} \to (\mathbb{N} \to x)) \wedge \\ \quad u \in (\mathbb{N} \to x) \wedge \\ \quad \forall k \in \mathbb{N}. \exists \mathit{cn} \in \mathbb{N}. \forall i \in \mathbb{N}. \\ \quad \quad (\mathit{cn} \leq i) \Rightarrow (s(i)(k) = u(k)) \Rightarrow \\ \quad \quad \exists \mathit{pn} \in \mathbb{N}. \forall a, j \in \mathbb{N}. \\ \quad \quad \quad ((\mathit{pn} \leq a) \wedge (j < n)) \Rightarrow (s(a)(j) = u(j)) \end{array}}}
\vnbprop{prefix-dist-small}{\vnbstmt{\begin{array}{@{}l@{}} \forall x, a, s, b.; \forall n \in \mathbb{N}. \\ \quad a \in (x \to \mathbb{R}) \wedge \\ \quad \forall w \in x.\; (1 < a(w)) \wedge \\ \quad s \in (\mathbb{N} \to x) \wedge \\ \quad b \in (\mathbb{N} \to x) \wedge \\ \quad \forall j \in \mathbb{N}.\; ((j < n) \Rightarrow (s(j) = b(j))) \Rightarrow \\ \quad \quad \operatorname{prefix\text{-}dist}(a, s, b) \leq \operatorname{recip}(\operatorname{prefix\text{-}prod}(a, b, n)) \end{array}}}
\vnbprop{prefix-dist-large}{\vnbstmt{\begin{array}{@{}l@{}} \forall x, a, s, b.; \forall k \in \mathbb{N}. \\ \quad a \in (x \to \mathbb{R}) \wedge \\ \quad \forall w \in x.\; (1 < a(w)) \wedge \\ \quad s \in (\mathbb{N} \to x) \wedge \\ \quad b \in (\mathbb{N} \to x) \wedge \\ \quad \neg (s(k) = b(k)) \Rightarrow \\ \quad \quad \operatorname{recip}(\operatorname{prefix\text{-}prod}(a, b, k)) \leq \operatorname{prefix\text{-}dist}(a, s, b) \end{array}}}
\vnbprop{ultra-seq-dist-apply}{\vnbstmt{\begin{array}{@{}l@{}} \forall x, a.; \forall s, b \in (\mathbb{N} \to x). \\ \quad ((x \in \mathbf{Set}) \wedge ((a \in (x \to \mathbb{R})) \wedge \forall w \in x.\; (1 < a(w)))) \Rightarrow  \\ \quad \quad \operatorname{dist}(\operatorname{ultra\text{-}seq\text{-}ms}(x, a))(s, b) = \operatorname{prefix\text{-}dist}(a, s, b) \end{array}}}
\vnbprop{ultra-seq-converges-iff}{\vnbstmt{\begin{array}{@{}l@{}} \forall x, a.; \forall s \in (\mathbb{N} \to (\mathbb{N} \to x)), u \in (\mathbb{N} \to x). \\ \quad x \in \mathbf{Set} \wedge \\ \quad a \in (x \to \mathbb{R}) \wedge \\ \quad \forall w \in x.\; (1 < a(w)) \wedge \\ \quad \forall z \in (\mathbb{N} \to x), m \in \mathbb{R}.\; \exists n \in \mathbb{N}.\; (m < \operatorname{prefix\text{-}prod}(a, z, n)) \Rightarrow \\ \quad \quad \operatorname{converges\text{-}to}(\operatorname{ultra\text{-}seq\text{-}ms}(x, a), s, u) \\ \quad \quad \Leftrightarrow\; \forall k \in \mathbb{N}. \exists \mathit{cn} \in \mathbb{N}. \forall i \in \mathbb{N}. \\ \quad \quad \quad (\mathit{cn} \leq i) \Rightarrow (s(i)(k) = u(k)) \end{array}}}

\vnbbookfile{structure-library/ccr1-vocab}{structure-library/ccr1-vocab}
\noindent{\small ccr1-vocab.scm -- ULTRAMETRIC SPACES (CCR.pdf, chapter 1, Definition 1.1 and the objects of Propositions 1.3-1.7).  DEFINITIONS ONLY; the laws are in theorem-library/ccr1-ultrametric.scm. Agent CCR-B1, wave MAP-1026, 2026-10-26.  Destination: structure-library/ccr1-vocab.scm. Library: analysis.  Load window: after ...}\par
\par\medskip\noindent\textbf{Definition} (\texttt{sphere}).\;$\operatorname{sphere}(\operatorname{cuvs\_}, \operatorname{cuvc\_}, \operatorname{cuvr\_}) :=$ $\{\operatorname{cuvy\_} \in \operatorname{pts}(\operatorname{cuvs\_}) : (\operatorname{dist}(\operatorname{cuvs\_})(\operatorname{cuvc\_}, \operatorname{cuvy\_}) = \operatorname{cuvr\_})\}$
\par\medskip\noindent\textbf{Definition} (\texttt{eps-rel}).\;$\operatorname{eps\text{-}rel}(\operatorname{cuvs\_}, \operatorname{cuve\_}) :=$
\begin{equation*}
\{\operatorname{cuvp\_} \in (\operatorname{pts}(\operatorname{cuvs\_}) \times \operatorname{pts}(\operatorname{cuvs\_})) : (\operatorname{dist}(\operatorname{cuvs\_})(\operatorname{nth}(1, \operatorname{cuvp\_}), \operatorname{nth}(2, \operatorname{cuvp\_})) < \operatorname{cuve\_})\}
\end{equation*}

\par\medskip\noindent\textbf{Definition} (\texttt{eps-setoid}).\;$\operatorname{eps\text{-}setoid}(\operatorname{cuvs\_}, \operatorname{cuve\_}) :=$ ${[}\operatorname{pts}(\operatorname{cuvs\_}), \operatorname{eps\text{-}rel}(\operatorname{cuvs\_}, \operatorname{cuve\_})]$

\vnbbookfile{theorem-library/ccr1-ultrametric}{theorem-library/ccr1-ultrametric}
\noindent{\small ccr1-ultrametric.scm -- ULTRAMETRIC SPACES: CCR.pdf chapter 1, Propositions 1.2, 1.3, Corollary 1.4, Propositions 1.5 and 1.7.  Vocabulary: structure-library/ccr1-vocab.scm (SPHERE, EPS-REL, EPS-SETOID) and IS-ULTRAMETRIC (agent ISOM's, scratchpad/mb12/mb12-vocab.scm).  Agent CCR-B1, wave MAP-1026, 2026-10-26.  ...}\par
\vnbprop{ultra-ms}{\vnbstmt{\forall s.\; (\operatorname{is\text{-}ultrametric}(s) \Rightarrow \operatorname{is\text{-}metric\text{-}space}(s))}}
\vnbprop{ultra-max}{\vnbstmt{\begin{array}{@{}l@{}} \forall s.; \forall x, y, z \in \operatorname{pts}(s). \\ \quad \operatorname{is\text{-}ultrametric}(s) \Rightarrow  \\ \quad \quad \operatorname{dist}(s)(x, y) \leq \operatorname{max}(\operatorname{dist}(s)(x, z), \operatorname{dist}(s)(z, y)) \end{array}}}
\vnbprop{ultra-le}{\vnbstmt{\begin{array}{@{}l@{}} \forall s, x, y, z, r. \\ \quad \operatorname{is\text{-}ultrametric}(s) \wedge \\ \quad x \in \operatorname{pts}(s) \wedge \\ \quad y \in \operatorname{pts}(s) \wedge \\ \quad z \in \operatorname{pts}(s) \wedge \\ \quad r \in \mathbb{R} \wedge \\ \quad \operatorname{dist}(s)(x, z) \leq r \wedge \\ \quad \operatorname{dist}(s)(z, y) \leq r \Rightarrow \\ \quad \quad \operatorname{dist}(s)(x, y) \leq r \end{array}}}
\vnbprop{ultra-lt}{\vnbstmt{\begin{array}{@{}l@{}} \forall s, x, y, z, r. \\ \quad \operatorname{is\text{-}ultrametric}(s) \wedge \\ \quad x \in \operatorname{pts}(s) \wedge \\ \quad y \in \operatorname{pts}(s) \wedge \\ \quad z \in \operatorname{pts}(s) \wedge \\ \quad r \in \mathbb{R} \wedge \\ \quad \operatorname{dist}(s)(x, z) < r \wedge \\ \quad \operatorname{dist}(s)(z, y) < r \Rightarrow \\ \quad \quad \operatorname{dist}(s)(x, y) < r \end{array}}}
\vnbprop{ultra-isosceles}{\vnbstmt{\begin{array}{@{}l@{}} \forall s, x, y, z. \\ \quad \operatorname{is\text{-}ultrametric}(s) \wedge \\ \quad x \in \operatorname{pts}(s) \wedge \\ \quad y \in \operatorname{pts}(s) \wedge \\ \quad z \in \operatorname{pts}(s) \wedge \\ \quad \operatorname{dist}(s)(z, y) < \operatorname{dist}(s)(x, z) \Rightarrow \\ \quad \quad \operatorname{dist}(s)(x, y) = \operatorname{dist}(s)(x, z) \end{array}}}
\vnbprop{ultra-closed-ball-nested}{\vnbstmt{\begin{array}{@{}l@{}} \forall s, a, b, r, q. \\ \quad \operatorname{is\text{-}ultrametric}(s) \wedge \\ \quad a \in \operatorname{pts}(s) \wedge \\ \quad b \in \operatorname{pts}(s) \wedge \\ \quad r \in \mathbb{R} \wedge \\ \quad q \in \mathbb{R} \wedge \\ \quad q \leq r \wedge \\ \quad \exists x \in \operatorname{closed\text{-}ball}(s, a, r).\; (x \in \operatorname{closed\text{-}ball}(s, b, q)) \Rightarrow \\ \quad \quad \operatorname{closed\text{-}ball}(s, b, q) \subseteq \operatorname{closed\text{-}ball}(s, a, r) \end{array}}}
\vnbprop{ultra-ball-nested}{\vnbstmt{\begin{array}{@{}l@{}} \forall s, a, b, r, q. \\ \quad \operatorname{is\text{-}ultrametric}(s) \wedge \\ \quad a \in \operatorname{pts}(s) \wedge \\ \quad b \in \operatorname{pts}(s) \wedge \\ \quad r \in \mathbb{R} \wedge \\ \quad q \in \mathbb{R} \wedge \\ \quad q \leq r \wedge \\ \quad \exists x \in \operatorname{ball}(s, a, r).\; (x \in \operatorname{ball}(s, b, q)) \Rightarrow \\ \quad \quad \operatorname{ball}(s, b, q) \subseteq \operatorname{ball}(s, a, r) \end{array}}}
\vnbprop{ultra-closed-ball-equal}{\vnbstmt{\begin{array}{@{}l@{}} \forall s, a, b, r. \\ \quad \operatorname{is\text{-}ultrametric}(s) \wedge \\ \quad a \in \operatorname{pts}(s) \wedge \\ \quad b \in \operatorname{pts}(s) \wedge \\ \quad r \in \mathbb{R} \wedge \\ \quad \exists x \in \operatorname{closed\text{-}ball}(s, a, r).\; (x \in \operatorname{closed\text{-}ball}(s, b, r)) \Rightarrow \\ \quad \quad \operatorname{closed\text{-}ball}(s, a, r) = \operatorname{closed\text{-}ball}(s, b, r) \end{array}}}
\vnbprop{ultra-ball-equal}{\vnbstmt{\begin{array}{@{}l@{}} \forall s, a, b, r. \\ \quad \operatorname{is\text{-}ultrametric}(s) \wedge \\ \quad a \in \operatorname{pts}(s) \wedge \\ \quad b \in \operatorname{pts}(s) \wedge \\ \quad r \in \mathbb{R} \wedge \\ \quad \exists x \in \operatorname{ball}(s, a, r).\; (x \in \operatorname{ball}(s, b, r)) \Rightarrow \\ \quad \quad \operatorname{ball}(s, a, r) = \operatorname{ball}(s, b, r) \end{array}}}
\vnbprop{ultra-closed-ball-open}{\vnbstmt{\begin{array}{@{}l@{}} \forall s, a, r. \\ \quad (\operatorname{is\text{-}ultrametric}(s) \wedge ((a \in \operatorname{pts}(s)) \wedge ((r \in \mathbb{R}) \wedge (0 < r)))) \Rightarrow  \\ \quad \quad \operatorname{is\text{-}open}(s, \operatorname{closed\text{-}ball}(s, a, r)) \end{array}}}
\vnbprop{ultra-ball-closed}{\vnbstmt{\begin{array}{@{}l@{}} \forall s, a, r. \\ \quad (\operatorname{is\text{-}ultrametric}(s) \wedge ((a \in \operatorname{pts}(s)) \wedge ((r \in \mathbb{R}) \wedge (0 < r)))) \Rightarrow  \\ \quad \quad \operatorname{is\text{-}closed}(s, \operatorname{ball}(s, a, r)) \end{array}}}
\vnbprop{sphere-unfold}{\vnbstmt{\forall s, c, r.\; (\operatorname{sphere}(s, c, r) \simeq \{y \in \operatorname{pts}(s) : (\operatorname{dist}(s)(c, y) = r)\})}}
\vnbprop{sphere-membership}{\vnbstmt{\begin{array}{@{}l@{}} \forall s, c, r, y. \\ \quad (y \in \operatorname{sphere}(s, c, r)) \Leftrightarrow ((y \in \operatorname{pts}(s)) \wedge (\operatorname{dist}(s)(c, y) = r)) \end{array}}}
\vnbprop{ultra-sphere-open}{\vnbstmt{\begin{array}{@{}l@{}} \forall s, a, r. \\ \quad (\operatorname{is\text{-}ultrametric}(s) \wedge ((a \in \operatorname{pts}(s)) \wedge ((r \in \mathbb{R}) \wedge (0 < r)))) \Rightarrow  \\ \quad \quad \operatorname{is\text{-}open}(s, \operatorname{sphere}(s, a, r)) \end{array}}}
\vnbprop{ultra-sphere-closed}{\vnbstmt{\begin{array}{@{}l@{}} \forall s, a, r. \\ \quad (\operatorname{is\text{-}ultrametric}(s) \wedge ((a \in \operatorname{pts}(s)) \wedge ((r \in \mathbb{R}) \wedge (0 < r)))) \Rightarrow  \\ \quad \quad \operatorname{is\text{-}closed}(s, \operatorname{sphere}(s, a, r)) \end{array}}}
\vnbprop{eps-rel-unfold}{\vnbstmt{\begin{array}{@{}l@{}} \forall s, e. \\ \quad (\operatorname{eps\text{-}rel}(s, e) \simeq \{p \in (\operatorname{pts}(s) \times \operatorname{pts}(s)) : (\operatorname{dist}(s)(\operatorname{nth}(1, p), \operatorname{nth}(2, p)) < e)\}) \end{array}}}
\vnbprop{eps-rel-mem}{\vnbstmt{\begin{array}{@{}l@{}} \forall s, e, u, v. \\ \quad ((u \in \operatorname{pts}(s)) \wedge (v \in \operatorname{pts}(s))) \Rightarrow  \\ \quad \quad ({[}u, v] \in \operatorname{eps\text{-}rel}(s, e)) \Leftrightarrow (\operatorname{dist}(s)(u, v) < e) \end{array}}}
\vnbprop{eps-rel-equivalence}{\vnbstmt{\begin{array}{@{}l@{}} \forall s.; \forall e \in \mathbb{R}. \\ \quad (\operatorname{is\text{-}ultrametric}(s) \wedge (0 < e)) \Rightarrow \operatorname{is\text{-}equivalence}(\operatorname{eps\text{-}rel}(s, e), \operatorname{pts}(s)) \end{array}}}
\vnbprop{eps-rel-set}{\vnbstmt{\forall s, e.\; (\operatorname{is\text{-}metric\text{-}space}(s) \Rightarrow (\operatorname{eps\text{-}rel}(s, e) \in \mathbf{Set}))}}
\vnbprop{eps-setoid-pts}{\vnbstmt{\forall s, e.\; (\operatorname{is\text{-}metric\text{-}space}(s) \Rightarrow (\operatorname{pts}(\operatorname{eps\text{-}setoid}(s, e)) \simeq \operatorname{pts}(s)))}}
\vnbprop{eps-setoid-rel}{\vnbstmt{\begin{array}{@{}l@{}} \forall s, e. \\ \quad \operatorname{is\text{-}metric\text{-}space}(s) \Rightarrow (\operatorname{rel}(\operatorname{eps\text{-}setoid}(s, e)) \simeq \operatorname{eps\text{-}rel}(s, e)) \end{array}}}
\vnbprop{eps-setoid-is-setoid}{\vnbstmt{\begin{array}{@{}l@{}} \forall s.; \forall e \in \mathbb{R}. \\ \quad (\operatorname{is\text{-}ultrametric}(s) \wedge (0 < e)) \Rightarrow \operatorname{is\text{-}setoid}(\operatorname{eps\text{-}setoid}(s, e)) \end{array}}}
\vnbprop{ultra-totally-bounded-iff}{\vnbstmt{\begin{array}{@{}l@{}} \forall s. \\ \quad \operatorname{is\text{-}ultrametric}(s) \Rightarrow  \\ \quad \quad \operatorname{totally\text{-}bounded}(s) \Leftrightarrow \forall e \in \mathbb{R}.\; ((0 < e) \Rightarrow (|\operatorname{quotient}(\operatorname{eps\text{-}setoid}(s, e))| \in \mathbb{N})) \end{array}}}

\vnbbookfile{theorem-library/rpow-star-convexity}{theorem-library/rpow-star-convexity}
\noindent{\small theorem-library/rpow-star-convexity.scm -- the two REAL inequalities of the axiomatic RPOW, proved for the DEFINED power RPOW-STAR (rpow-star.scm):}\par
\vnbprop{r-exp-ge-one-plus}{\vnbstmt{\forall z \in \mathbb{R}.\; ((1 + z) \leq \exp\left(z\right))}}
\vnbprop{log-le-minus-one}{\vnbstmt{\forall w \in \mathbb{R}.\; ((0 < w) \Rightarrow (\log\left(w\right) \leq (w - 1)))}}
\vnbprop{log-jensen-two}{\vnbstmt{\begin{array}{@{}l@{}} \forall u, v, s, t \in \mathbb{R}. \\ \quad ((0 < u) \wedge ((0 < v) \wedge ((0 < s) \wedge ((0 < t) \wedge ((s + t) = 1))))) \Rightarrow  \\ \quad \quad ((s \cdot \log\left(u\right)) + (t \cdot \log\left(v\right))) \leq \log\left(((s \cdot u) + (t \cdot v))\right) \end{array}}}
\vnbprop{young-inequality-star}{\vnbstmt{\begin{array}{@{}l@{}} \forall a, b, p, q \in \mathbb{R}. \\ \quad ((0 < a) \wedge ((0 < b) \wedge ((1 < p) \wedge ((1 < q) \wedge ((/(1, p) + /(1, q)) = 1))))) \Rightarrow  \\ \quad \quad (a \cdot b) \leq (/({a}^{p}, p) + /({b}^{q}, q)) \end{array}}}
\vnbprop{bernoulli-rpow-star}{\vnbstmt{\begin{array}{@{}l@{}} \forall x, b \in \mathbb{R}. \\ \quad ((0 < (1 + x)) \wedge (1 \leq b)) \Rightarrow ((1 + (b \cdot x)) \leq {(1 + x)}^{b}) \end{array}}}

\vnbbookfile{theorem-library/mb12b-dim}{theorem-library/mb12b-dim}
\noindent{\small mb12b-dim.scm -- construction-of-measures Proposition 7.38 (existence and uniqueness of the self-similar dimension lambda: sum over w in X of a\_w\textasciicircum{}(-lambda) = 1).  Agent ISOM (brief M-B1B2, follow-up), 2026-10-26.  Helper prefix: mdm-.  Library: analysis.}\par
\vnbprop{rpow-exp-facts}{\vnbstmt{\begin{array}{@{}l@{}} \forall b, s \in \mathbb{R}. \\ \quad (0 < b) \Rightarrow (({b}^{s} \in \mathbb{R}) \wedge (0 < {b}^{s})) \end{array}}}
\vnbprop{rpow-exp-continuous}{\vnbstmt{\begin{array}{@{}l@{}} \forall b \in \mathbb{R}. \\ \quad (0 < b) \Rightarrow  \\ \quad \quad (\lambda s \in \mathbb{R}.\; {b}^{s}) \in (\mathbb{R} \to \mathbb{R}) \wedge \\ \quad \quad \forall x \in \mathbb{R}. \\ \quad \quad \quad \operatorname{is\text{-}continuous\text{-}at}(\operatorname{rr\text{-}ms}, \operatorname{rr\text{-}ms}, (\lambda s \in \mathbb{R}.\; {b}^{s}), x) \end{array}}}
\vnbprop{rpow-exp-strict}{\vnbstmt{\begin{array}{@{}l@{}} \forall b, a, c \in \mathbb{R}. \\ \quad ((1 < b) \wedge (a < c)) \Rightarrow ({b}^{a} < {b}^{c}) \end{array}}}
\vnbprop{rpow-exp-base-le}{\vnbstmt{\begin{array}{@{}l@{}} \forall m, b, s \in \mathbb{R}. \\ \quad ((1 < m) \wedge ((m \leq b) \wedge (s \leq 0))) \Rightarrow ({b}^{s} \leq {m}^{s}) \end{array}}}
\vnbprop{self-similar-dimension}{\vnbstmt{\begin{array}{@{}l@{}} \forall x \in \mathbf{Set}, a \in (x \to \mathbb{R}). \\ \quad ((|x| \in \mathbb{N}) \wedge ((2 \leq |x|) \wedge \forall w \in x.\; (1 < a(w)))) \Rightarrow  \\ \quad \quad \exists l \in \mathbb{R}. \\ \quad \quad \quad 0 < l \wedge \\ \quad \quad \quad \operatorname{finsum}(\operatorname{commutative\text{-}ring\text{-}additive\text{-}ag}(\operatorname{rr\text{-}normed\text{-}field}), (\lambda w \in x.\; {a(w)}^{-l}), x) = 1 \wedge \\ \quad \quad \quad \forall b \in \mathbb{R}. \\ \quad \quad \quad \quad 0 < b \wedge \\ \quad \quad \quad \quad \operatorname{finsum}(\operatorname{commutative\text{-}ring\text{-}additive\text{-}ag}(\operatorname{rr\text{-}normed\text{-}field}), (\lambda w \in x.\; {a(w)}^{-b}), x) = 1 \Rightarrow \\ \quad \quad \quad \quad \quad b = l \end{array}}}

\vnbbookfile{structure-library/functor-invariance}{structure-library/functor-invariance}
\noindent{\small functor-invariance.scm -- a construction that cannot see the functor.}\par
\vnbprop{metric-group-group@abelian-group-as-monoid}{\vnbstmt{\begin{array}{@{}l@{}} \forall r. \\ \quad \operatorname{is\text{-}metric\text{-}group}(r) \Rightarrow  \\ \quad \quad (\operatorname{abelian\text{-}group\text{-}as\text{-}monoid}(\operatorname{metric\text{-}group\text{-}group}(r)) \simeq \operatorname{abelian\text{-}group\text{-}as\text{-}monoid}(r)) \end{array}}}
\vnbprop{metric-group-group@dsum-act}{\vnbstmt{\begin{array}{@{}l@{}} \forall s, n, f. \\ \quad \operatorname{is\text{-}metric\text{-}group}(s) \Rightarrow  \\ \quad \quad (\operatorname{dsum\text{-}act}(\operatorname{metric\text{-}group\text{-}group}(s), n, f) \simeq \operatorname{dsum\text{-}act}(s, n, f)) \end{array}}}
\vnbprop{metric-group-group@enum-fam}{\vnbstmt{\begin{array}{@{}l@{}} \forall g, f, i, n. \\ \quad \operatorname{is\text{-}metric\text{-}group}(g) \Rightarrow  \\ \quad \quad (\operatorname{enum\text{-}fam}(\operatorname{metric\text{-}group\text{-}group}(g), f, i, n) \simeq \operatorname{enum\text{-}fam}(g, f, i, n)) \end{array}}}
\vnbprop{metric-group-group@laurent-limc}{\vnbstmt{\begin{array}{@{}l@{}} \forall k, s, j. \\ \quad \operatorname{is\text{-}metric\text{-}group}(k) \Rightarrow  \\ \quad \quad (\operatorname{laurent\text{-}limc}(\operatorname{metric\text{-}group\text{-}group}(k), s, j) \simeq \operatorname{laurent\text{-}limc}(k, s, j)) \end{array}}}
\vnbprop{metric-group-group@metric-group-group}{\vnbstmt{\begin{array}{@{}l@{}} \forall r. \\ \quad \operatorname{is\text{-}metric\text{-}group}(r) \Rightarrow  \\ \quad \quad (\operatorname{metric\text{-}group\text{-}group}(\operatorname{metric\text{-}group\text{-}group}(r)) \simeq \operatorname{metric\text{-}group\text{-}group}(r)) \end{array}}}
\vnbprop{metric-group-group@normed-ag-as-abelian-group}{\vnbstmt{\begin{array}{@{}l@{}} \forall r. \\ \quad \operatorname{is\text{-}metric\text{-}group}(r) \Rightarrow  \\ \quad \quad (\operatorname{normed\text{-}ag\text{-}as\text{-}abelian\text{-}group}(\operatorname{metric\text{-}group\text{-}group}(r)) \simeq \operatorname{normed\text{-}ag\text{-}as\text{-}abelian\text{-}group}(r)) \end{array}}}
\vnbprop{metric-top@preimage}{\vnbstmt{\begin{array}{@{}l@{}} \forall s, f, v. \\ \quad \operatorname{is\text{-}metric\text{-}space}(s) \Rightarrow (\operatorname{preimage}(\operatorname{metric\text{-}top}(s), f, v) \simeq \operatorname{preimage}(s, f, v)) \end{array}}}
\vnbprop{metric-top@sup-dist}{\vnbstmt{\begin{array}{@{}l@{}} \forall s, t, f, g. \\ \quad \operatorname{is\text{-}metric\text{-}space}(s) \Rightarrow  \\ \quad \quad (\operatorname{sup\text{-}dist}(\operatorname{metric\text{-}top}(s), t, f, g) \simeq \operatorname{sup\text{-}dist}(s, t, f, g)) \end{array}}}
\vnbprop{trunc-metric@preimage}{\vnbstmt{\begin{array}{@{}l@{}} \forall s, f, v. \\ \quad \operatorname{is\text{-}metric\text{-}space}(s) \Rightarrow (\operatorname{preimage}(\operatorname{trunc\text{-}metric}(s), f, v) \simeq \operatorname{preimage}(s, f, v)) \end{array}}}
\vnbprop{trunc-metric@sup-dist}{\vnbstmt{\begin{array}{@{}l@{}} \forall s, t, f, g. \\ \quad \operatorname{is\text{-}metric\text{-}space}(s) \Rightarrow  \\ \quad \quad (\operatorname{sup\text{-}dist}(\operatorname{trunc\text{-}metric}(s), t, f, g) \simeq \operatorname{sup\text{-}dist}(s, t, f, g)) \end{array}}}

\vnbbookfile{theorem-library/metric-top-proof}{theorem-library/metric-top-proof}
\noindent{\small metric-top-proof.scm -- the metric open sets form a topology.}\par
\vnbprop{metric-top-is-top-space}{\vnbstmt{\forall d.\; (\operatorname{is\text{-}metric\text{-}space}(d) \Rightarrow \operatorname{is\text{-}top\text{-}space}(\operatorname{metric\text{-}top}(d)))}}
\vnbprop{metric-space-repack}{\vnbstmt{\begin{array}{@{}l@{}} \forall d. \\ \quad \operatorname{is\text{-}metric\text{-}space}(d) \Rightarrow  \\ \quad \quad {[}\operatorname{pts}(d), \operatorname{dist}(d)] \in \mathbf{Set} \wedge \\ \quad \quad \operatorname{is\text{-}metric\text{-}space}({[}\operatorname{pts}(d), \operatorname{dist}(d)]) \wedge \\ \quad \quad (\operatorname{metric\text{-}top}({[}\operatorname{pts}(d), \operatorname{dist}(d)]) \simeq \operatorname{metric\text{-}top}(d)) \end{array}}}
\vnbprop{metric-top-is-metrizable-top-space}{\vnbstmt{\forall d.\; (\operatorname{is\text{-}metric\text{-}space}(d) \Rightarrow \operatorname{is\text{-}metrizable\text{-}top\text{-}space}(\operatorname{metric\text{-}top}(d)))}}

\vnbbookfile{theorem-library/metric-top-functorial-proof}{theorem-library/metric-top-functorial-proof}
\noindent{\small metric-top-functorial-proof.scm -- the Met -$>$ Top functor acts on morphisms.}\par
\vnbprop{metric-top-functorial}{\vnbstmt{\begin{array}{@{}l@{}} \forall a, b, f. \\ \quad \operatorname{is\text{-}hom\text{-}metric\text{-}space}(a, b, f) \Rightarrow  \\ \quad \quad \operatorname{is\text{-}hom\text{-}metrizable\text{-}top\text{-}space}(\operatorname{metric\text{-}top}(a), \operatorname{metric\text{-}top}(b), f) \end{array}}}

\vnbbookfile{theorem-library/rake-bdd-metric}{theorem-library/rake-bdd-metric}
\noindent{\small rake-bdd-metric.scm -- the three asserted leaves still on the bill of `metrizable-iff-bounded-metrizable' (theorem-library/metrizable-bounded-proof.scm), PROVEN.  Statements copied VERBATIM from their support sites:}\par
\vnbprop{metrizable-has-metric-top}{\vnbstmt{\begin{array}{@{}l@{}} \forall s. \\ \quad \operatorname{is\text{-}metrizable\text{-}top\text{-}space}(s) \Rightarrow  \\ \quad \quad \exists d.\; (\operatorname{is\text{-}metric\text{-}space}(d) \wedge (\operatorname{metric\text{-}top}(d) \simeq s)) \end{array}}}
\vnbprop{bdd-metric-is-bounded-metric-space}{\vnbstmt{\forall s.\; (\operatorname{is\text{-}metric\text{-}space}(s) \Rightarrow \operatorname{is\text{-}bounded\text{-}metric\text{-}space}(\operatorname{bdd\text{-}metric}(s)))}}
\vnbprop{bdd-metric-open-bwd}{\vnbstmt{\begin{array}{@{}l@{}} \forall s, u. \\ \quad (\operatorname{is\text{-}metric\text{-}space}(s) \wedge \operatorname{is\text{-}open}(\operatorname{bdd\text{-}metric}(s), u)) \Rightarrow \operatorname{is\text{-}open}(s, u) \end{array}}}
\vnbprop{bdd-metric-open-fwd}{\vnbstmt{\begin{array}{@{}l@{}} \forall s, u. \\ \quad (\operatorname{is\text{-}metric\text{-}space}(s) \wedge \operatorname{is\text{-}open}(s, u)) \Rightarrow \operatorname{is\text{-}open}(\operatorname{bdd\text{-}metric}(s), u) \end{array}}}
\vnbprop{bdd-metric-preserves-metric-top}{\vnbstmt{\begin{array}{@{}l@{}} \forall s. \\ \quad \operatorname{is\text{-}metric\text{-}space}(s) \Rightarrow (\operatorname{metric\text{-}top}(\operatorname{bdd\text{-}metric}(s)) \simeq \operatorname{metric\text{-}top}(s)) \end{array}}}

\vnbbookfile{theorem-library/metrizable-compact-image}{theorem-library/metrizable-compact-image}
\noindent{\small metrizable-compact-image.scm -- THE CONTINUOUS IMAGE OF A COMPACT SPACE IS COMPACT, IN THE CATEGORY OF METRIZABLE TOPOLOGICAL SPACES, PROVED THROUGH SEQUENCES (the user's request, 2026-10-01: "in that category ... it can be done using sequences.  I like using sequences to prove things whenever}\par
\vnbprop{metric-top-opens-iff}{\vnbstmt{\begin{array}{@{}l@{}} \forall d, u. \\ \quad \operatorname{is\text{-}metric\text{-}space}(d) \Rightarrow ((u \in \operatorname{opens}(\operatorname{metric\text{-}top}(d))) \Leftrightarrow \operatorname{is\text{-}open}(d, u)) \end{array}}}
\vnbprop{metric-top-compact}{\vnbstmt{\begin{array}{@{}l@{}} \forall d. \\ \quad \operatorname{is\text{-}metric\text{-}space}(d) \Rightarrow (\operatorname{is\text{-}compact\text{-}t}(\operatorname{metric\text{-}top}(d)) \Leftrightarrow \operatorname{is\text{-}compact}(d)) \end{array}}}
\vnbprop{metric-top-hom}{\vnbstmt{\begin{array}{@{}l@{}} \forall s, t, f. \\ \quad (\operatorname{is\text{-}metric\text{-}space}(s) \wedge \operatorname{is\text{-}metric\text{-}space}(t)) \Rightarrow  \\ \quad \quad \operatorname{is\text{-}hom\text{-}top\text{-}space}(\operatorname{metric\text{-}top}(s), \operatorname{metric\text{-}top}(t), f) \Leftrightarrow \operatorname{is\text{-}continuous}(s, t, f) \end{array}}}
\vnbprop{seq-compact-image}{\vnbstmt{\begin{array}{@{}l@{}} \forall d, t, f. \\ \quad (\operatorname{is\text{-}metric\text{-}space}(d) \wedge (\operatorname{is\text{-}metric\text{-}space}(t) \wedge (\operatorname{is\text{-}continuous}(d, t, f) \wedge \operatorname{seq\text{-}compact}(d)))) \Rightarrow  \\ \quad \quad \operatorname{seq\text{-}compact}(\operatorname{subspace\text{-}ms}(t, \operatorname{image}(f, \operatorname{pts}(d)))) \end{array}}}
\vnbprop{metrizable-compact-t-iff-seq-compact}{\vnbstmt{\begin{array}{@{}l@{}} \forall s, d. \\ \quad (\operatorname{is\text{-}metric\text{-}space}(d) \wedge (\operatorname{metric\text{-}top}(d) \simeq s)) \Rightarrow  \\ \quad \quad \operatorname{is\text{-}compact\text{-}t}(s) \Leftrightarrow \operatorname{seq\text{-}compact}(d) \end{array}}}
\vnbprop{metrizable-hom-iff-sequential}{\vnbstmt{\begin{array}{@{}l@{}} \forall s, t, d, b, f. \\ \quad (\operatorname{is\text{-}metric\text{-}space}(d) \wedge (\operatorname{is\text{-}metric\text{-}space}(b) \wedge ((\operatorname{metric\text{-}top}(d) \simeq s) \wedge (\operatorname{metric\text{-}top}(b) \simeq t)))) \Rightarrow  \\ \quad \quad \operatorname{is\text{-}hom\text{-}top\text{-}space}(s, t, f) \\ \quad \quad \Leftrightarrow\; f \in (\operatorname{pts}(d) \to \operatorname{pts}(b)) \wedge \\ \quad \quad \forall a \in \operatorname{pts}(d), q \in \operatorname{pts}(d)^{\mathbb{N}}. \\ \quad \quad \quad \operatorname{converges\text{-}to}(d, q, a) \Rightarrow \operatorname{converges\text{-}to}(b, \operatorname{compose}(f, q), f(a)) \end{array}}}
\vnbprop{metrizable-compact-image-ms}{\vnbstmt{\begin{array}{@{}l@{}} \forall s, t, f, a. \\ \quad \operatorname{is\text{-}metrizable\text{-}top\text{-}space}(s) \wedge \\ \quad \operatorname{is\text{-}metrizable\text{-}top\text{-}space}(t) \wedge \\ \quad \operatorname{is\text{-}hom\text{-}top\text{-}space}(s, t, f) \wedge \\ \quad \operatorname{is\text{-}compact\text{-}t}(s) \wedge \\ \quad \operatorname{is\text{-}metric\text{-}space}(a) \wedge \\ \quad (\operatorname{metric\text{-}top}(a) \simeq t) \Rightarrow \\ \quad \quad \operatorname{is\text{-}compact}(\operatorname{subspace\text{-}ms}(a, \operatorname{image}(f, \operatorname{pts}(s)))) \end{array}}}
\vnbprop{subspace-top-pts}{\vnbstmt{\forall t.; \forall a \in \mathbf{Set}.\; (\operatorname{pts}(\operatorname{subspace\text{-}top}(t, a)) \simeq a)}}
\vnbprop{subspace-top-opens}{\vnbstmt{\begin{array}{@{}l@{}} \forall t.; \forall a \in \mathbf{Set}. \\ \quad (\operatorname{opens}(\operatorname{subspace\text{-}top}(t, a)) \simeq \{w \in \operatorname{power}(a) : \exists u \in \operatorname{opens}(t).\; (w = (u \cap a))\}) \end{array}}}
\vnbprop{subspace-top-is-top-space}{\vnbstmt{\begin{array}{@{}l@{}} \forall t, a. \\ \quad (\operatorname{is\text{-}top\text{-}space}(t) \wedge (a \subseteq \operatorname{pts}(t))) \Rightarrow \operatorname{is\text{-}top\text{-}space}(\operatorname{subspace\text{-}top}(t, a)) \end{array}}}
\vnbprop{subspace-top-of-metric-top}{\vnbstmt{\begin{array}{@{}l@{}} \forall t, a. \\ \quad (\operatorname{is\text{-}metric\text{-}space}(t) \wedge (a \subseteq \operatorname{pts}(t))) \Rightarrow  \\ \quad \quad (\operatorname{subspace\text{-}top}(\operatorname{metric\text{-}top}(t), a) \simeq \operatorname{metric\text{-}top}(\operatorname{subspace\text{-}ms}(t, a))) \end{array}}}
\vnbprop{metrizable-compact-image}{\vnbstmt{\begin{array}{@{}l@{}} \forall s, t, f. \\ \quad \operatorname{is\text{-}metrizable\text{-}top\text{-}space}(s) \wedge \\ \quad \operatorname{is\text{-}metrizable\text{-}top\text{-}space}(t) \wedge \\ \quad \operatorname{is\text{-}hom\text{-}top\text{-}space}(s, t, f) \wedge \\ \quad \operatorname{is\text{-}compact\text{-}t}(s) \Rightarrow \\ \quad \quad \operatorname{is\text{-}compact\text{-}t}(\operatorname{subspace\text{-}top}(t, \operatorname{image}(f, \operatorname{pts}(s)))) \end{array}}}

\vnbbookfile{theorem-library/metrizable-bounded-proof}{theorem-library/metrizable-bounded-proof}
\noindent{\small metrizable-bounded-proof.scm -- T1: a topological space is metrizable iff it is metrizable by a BOUNDED metric.}\par
\vnbprop{metrizable-iff-bounded-metrizable}{\vnbstmt{\begin{array}{@{}l@{}} \forall s. \\ \quad \operatorname{is\text{-}top\text{-}space}(s) \Rightarrow  \\ \quad \quad \operatorname{is\text{-}metrizable\text{-}top\text{-}space}(s) \Leftrightarrow \exists d.\; (\operatorname{is\text{-}bounded\text{-}metric\text{-}space}(d) \wedge (\operatorname{metric\text{-}top}(d) \simeq s)) \end{array}}}

\vnbbookfile{theorem-library/geometric-series}{theorem-library/geometric-series}
\noindent{\small geometric-series.scm -- THE GEOMETRIC MAJORANT, PROVEN, and the four small facts about powers it rests on.}\par
\vnbprop{rr-abs-power}{\vnbstmt{\forall k \in \mathbb{N}, r \in \mathbb{R}.\; (\lvert \operatorname{power}(r, k) \rvert = \operatorname{power}(\lvert r \rvert, k))}}
\vnbprop{geometric-partial-sum-mul}{\vnbstmt{\begin{array}{@{}l@{}} \forall k \in \mathbb{N}, r \in \mathbb{R}. \\ \quad ((1 - r) \cdot \operatorname{series\text{-}partial\text{-}sum}((\lambda n \in \mathbb{N}.\; \operatorname{power}(r, n)), k)) = (1 - \operatorname{power}(r, k)) \end{array}}}
\vnbprop{rr-power-nonneg}{\vnbstmt{\forall k \in \mathbb{N}, r \in \mathbb{R}.\; ((0 \leq r) \Rightarrow (0 \leq \operatorname{power}(r, k)))}}
\vnbprop{geometric-partial-sum}{\vnbstmt{\begin{array}{@{}l@{}} \forall r \in \mathbb{R}, k \in \mathbb{N}. \\ \quad \neg (r = 1) \Rightarrow  \\ \quad \quad \operatorname{series\text{-}partial\text{-}sum}((\lambda n \in \mathbb{N}.\; \operatorname{power}(r, n)), k) = ((1 - \operatorname{power}(r, k)) \cdot \operatorname{recip}((1 - r))) \end{array}}}
\vnbprop{geometric-series-converges}{\vnbstmt{\begin{array}{@{}l@{}} \forall r \in \mathbb{R}. \\ \quad ((0 \leq r) \wedge (r < 1)) \Rightarrow \operatorname{series\text{-}converges}((\lambda n \in \mathbb{N}.\; \operatorname{power}(r, n))) \end{array}}}

\vnbbookfile{theorem-library/complete-criterion}{theorem-library/complete-criterion}
\noindent{\small complete-criterion.scm -- wave of 2026-11-09/10 (agent BA-1E), wired 2026-11-10 from scratchpad/ba1e-1109/complete-criterion.scm. Library: analysis.}\par
\vnbprop{complete-of-summable-increments}{\vnbstmt{\begin{array}{@{}l@{}} \forall s. \\ \quad \operatorname{is\text{-}metric\text{-}space}(s) \wedge \\ \quad \forall f \in (\mathbb{N} \to \operatorname{pts}(s)), d \in (\mathbb{N} \to \mathbb{R}). \\ \quad \quad (\operatorname{series\text{-}converges}(d) \wedge \forall k \in \mathbb{N}.\; (\operatorname{dist}(s)(f(k), f((k + 1))) \leq d(k))) \Rightarrow  \\ \quad \quad \quad \operatorname{converges}(s, f) \Rightarrow \\ \quad \quad \operatorname{is\text{-}complete}(s) \end{array}}}
\vnbprop{nag-complete-of-summable-increments}{\vnbstmt{\begin{array}{@{}l@{}} \forall p. \\ \quad \operatorname{is\text{-}normed\text{-}ag}(p) \wedge \\ \quad \forall f \in (\mathbb{N} \to \operatorname{carr}(p)), d \in (\mathbb{N} \to \mathbb{R}). \\ \quad \quad (\operatorname{series\text{-}converges}(d) \wedge \forall k \in \mathbb{N}.\; (\operatorname{nrm}(p)(\operatorname{opr}(p)(f(k), \operatorname{inv}(p)(f((k + 1))))) \leq d(k))) \Rightarrow  \\ \quad \quad \quad \operatorname{converges}(\operatorname{nag\text{-}metric\text{-}space}(p), f) \Rightarrow \\ \quad \quad \operatorname{is\text{-}complete}(\operatorname{nag\text{-}metric\text{-}space}(p)) \end{array}}}

\vnbbookfile{theorem-library/limsup-tests}{theorem-library/limsup-tests}
\noindent{\small limsup-tests.scm -- THE UPPER LIMIT OF A SEQUENCE IN [0, +inf], THE ROOT TEST AND THE RATIO TEST, as the notes state them. Batch 25-A, 2026-09-24.}\par
\vnbprop{einf-lower}{\vnbstmt{\forall s.; \forall x \in s.\; ((s \subseteq \operatorname{rr\text{-}pos\text{-}star}) \Rightarrow (\operatorname{einf}(s) \leq x))}}
\vnbprop{einf-greatest}{\vnbstmt{\begin{array}{@{}l@{}} \forall s.; \forall b \in \operatorname{rr\text{-}pos\text{-}star}. \\ \quad ((s \subseteq \operatorname{rr\text{-}pos\text{-}star}) \wedge \forall x \in s.\; (b \leq x)) \Rightarrow (b \leq \operatorname{einf}(s)) \end{array}}}
\vnbprop{einf-in}{\vnbstmt{\forall s.\; ((s \subseteq \operatorname{rr\text{-}pos\text{-}star}) \Rightarrow (\operatorname{einf}(s) \in \operatorname{rr\text{-}pos\text{-}star}))}}
\vnbprop{etail-subset-rr-pos-star}{\vnbstmt{\begin{array}{@{}l@{}} \forall f, n. \\ \quad (\forall k \in \mathbb{N}.\; (f(k) \in \mathbb{R}) \wedge \forall k \in \mathbb{N}.\; (0 \leq f(k))) \Rightarrow  \\ \quad \quad \operatorname{etail}(f, n) \subseteq \operatorname{rr\text{-}pos\text{-}star} \end{array}}}
\vnbprop{elimsup-eventually-below}{\vnbstmt{\begin{array}{@{}l@{}} \forall f.; \forall r \in \mathbb{R}. \\ \quad (\forall k \in \mathbb{N}.\; (f(k) \in \mathbb{R}) \wedge (\forall k \in \mathbb{N}.\; (0 \leq f(k)) \wedge (\operatorname{elimsup}(f) < r))) \Rightarrow  \\ \quad \quad \exists n \in \mathbb{N}.\; \forall k \in \mathbb{N}.\; ((n \leq k) \Rightarrow (f(k) < r)) \end{array}}}
\vnbprop{elimsup-frequently-above}{\vnbstmt{\begin{array}{@{}l@{}} \forall f.; \forall r \in \mathbb{R}, n \in \mathbb{N}. \\ \quad (\forall k \in \mathbb{N}.\; (f(k) \in \mathbb{R}) \wedge (\forall k \in \mathbb{N}.\; (0 \leq f(k)) \wedge ((0 \leq r) \wedge (r < \operatorname{elimsup}(f))))) \Rightarrow  \\ \quad \quad \exists k \in \mathbb{N}.\; ((n \leq k) \wedge (r < f(k))) \end{array}}}
\vnbprop{eliminf-eventually-above}{\vnbstmt{\begin{array}{@{}l@{}} \forall f.; \forall r \in \mathbb{R}. \\ \quad (\forall k \in \mathbb{N}.\; (f(k) \in \mathbb{R}) \wedge (\forall k \in \mathbb{N}.\; (0 \leq f(k)) \wedge ((0 \leq r) \wedge (r < \operatorname{eliminf}(f))))) \Rightarrow  \\ \quad \quad \exists n \in \mathbb{N}.\; \forall k \in \mathbb{N}.\; ((n \leq k) \Rightarrow (r < f(k))) \end{array}}}
\vnbprop{elimsup-in-rr-pos-star}{\vnbstmt{\begin{array}{@{}l@{}} \forall f. \\ \quad (\forall k \in \mathbb{N}.\; (f(k) \in \mathbb{R}) \wedge \forall k \in \mathbb{N}.\; (0 \leq f(k))) \Rightarrow  \\ \quad \quad \operatorname{elimsup}(f) \in \operatorname{rr\text{-}pos\text{-}star} \end{array}}}
\vnbprop{eliminf-in-rr-pos-star}{\vnbstmt{\begin{array}{@{}l@{}} \forall f. \\ \quad (\forall k \in \mathbb{N}.\; (f(k) \in \mathbb{R}) \wedge \forall k \in \mathbb{N}.\; (0 \leq f(k))) \Rightarrow  \\ \quad \quad \operatorname{eliminf}(f) \in \operatorname{rr\text{-}pos\text{-}star} \end{array}}}
\vnbprop{rr-power-mono-base}{\vnbstmt{\forall k \in \mathbb{N}, x, y \in \mathbb{R}.\; (((0 \leq x) \wedge (x \leq y)) \Rightarrow (\operatorname{power}(x, k) \leq \operatorname{power}(y, k)))}}
\vnbprop{rpow-star-nat}{\vnbstmt{\forall k \in \mathbb{N}, x \in \mathbb{R}.\; ((0 < x) \Rightarrow ({x}^{k} = \operatorname{power}(x, k)))}}
\vnbprop{recip-star-nn-pos}{\vnbstmt{\forall k \in \mathbb{N}.\; ((1 \leq k) \Rightarrow (0 < {k}^{\ast\!-1}))}}
\vnbprop{recip-star-nn-nonneg}{\vnbstmt{\forall k \in \mathbb{N}.\; (0 \leq {k}^{\ast\!-1})}}
\vnbprop{rpow-star-root-power}{\vnbstmt{\begin{array}{@{}l@{}} \forall k \in \mathbb{N}, x \in \mathbb{R}. \\ \quad ((1 \leq k) \wedge (0 < x)) \Rightarrow (\operatorname{power}({x}^{{k}^{\ast\!-1}}, k) = x) \end{array}}}
\vnbprop{rpow-star-nonneg-in-rr}{\vnbstmt{\forall x, s \in \mathbb{R}.\; (((0 \leq x) \wedge (0 \leq s)) \Rightarrow ({x}^{s} \in \mathbb{R}))}}
\vnbprop{rpow-star-nonneg-nonneg}{\vnbstmt{\forall x, s \in \mathbb{R}.\; (((0 \leq x) \wedge (0 \leq s)) \Rightarrow (0 \leq {x}^{s}))}}
\vnbprop{series-bounded-converges}{\vnbstmt{\begin{array}{@{}l@{}} \forall f \in (\mathbb{N} \to \mathbb{R}), b \in \mathbb{R}. \\ \quad (\forall k \in \mathbb{N}.\; (0 \leq f(k)) \wedge \forall k \in \mathbb{N}.\; (\operatorname{series\text{-}partial\text{-}sum}(f, k) \leq b)) \Rightarrow  \\ \quad \quad \operatorname{series\text{-}converges}(f) \end{array}}}
\vnbprop{series-eventually-geometric-converges}{\vnbstmt{\begin{array}{@{}l@{}} \forall f \in (\mathbb{N} \to \mathbb{R}), q, c \in \mathbb{R}, n \in \mathbb{N}. \\ \quad \forall k \in \mathbb{N}.\; (0 \leq f(k)) \wedge \\ \quad 0 \leq q \wedge \\ \quad q < 1 \wedge \\ \quad 0 \leq c \wedge \\ \quad \forall k \in \mathbb{N}.\; ((n \leq k) \Rightarrow (f(k) \leq (c \cdot \operatorname{power}(q, k)))) \Rightarrow \\ \quad \quad \operatorname{series\text{-}converges}(f) \end{array}}}
\vnbprop{series-frequently-large-diverges}{\vnbstmt{\begin{array}{@{}l@{}} \forall f \in (\mathbb{N} \to \mathbb{R}), c \in \mathbb{R}. \\ \quad ((0 < c) \wedge \forall n \in \mathbb{N}.\; \exists k \in \mathbb{N}.\; ((n \leq k) \wedge (c \leq f(k)))) \Rightarrow  \\ \quad \quad \neg \operatorname{series\text{-}converges}(f) \end{array}}}
\vnbprop{root-test-converges}{\vnbstmt{\begin{array}{@{}l@{}} \forall h \in (\mathbb{N} \to \mathbb{R}). \\ \quad (\forall k \in \mathbb{N}.\; (0 \leq h(k)) \wedge (\operatorname{elimsup}((\lambda j \in \mathbb{N}.\; {h(j)}^{{j}^{\ast\!-1}})) < 1)) \Rightarrow  \\ \quad \quad \operatorname{series\text{-}converges}(h) \end{array}}}
\vnbprop{root-test-diverges}{\vnbstmt{\begin{array}{@{}l@{}} \forall h \in (\mathbb{N} \to \mathbb{R}). \\ \quad (\forall k \in \mathbb{N}.\; (0 \leq h(k)) \wedge (1 < \operatorname{elimsup}((\lambda j \in \mathbb{N}.\; {h(j)}^{{j}^{\ast\!-1}})))) \Rightarrow  \\ \quad \quad \neg \operatorname{series\text{-}converges}(h) \end{array}}}
\vnbprop{root-test}{\vnbstmt{\begin{array}{@{}l@{}} \forall h \in (\mathbb{N} \to \mathbb{R}). \\ \quad \forall k \in \mathbb{N}.\; (0 \leq h(k)) \Rightarrow  \\ \quad \quad (\operatorname{elimsup}((\lambda j \in \mathbb{N}.\; {h(j)}^{{j}^{\ast\!-1}})) < 1) \Rightarrow  \\ \quad \quad \quad \operatorname{series\text{-}converges}(h) \wedge \\ \quad \quad (1 < \operatorname{elimsup}((\lambda j \in \mathbb{N}.\; {h(j)}^{{j}^{\ast\!-1}}))) \Rightarrow  \\ \quad \quad \quad \neg \operatorname{series\text{-}converges}(h) \end{array}}}
\vnbprop{ratio-test-converges-limsup}{\vnbstmt{\begin{array}{@{}l@{}} \forall h \in (\mathbb{N} \to \mathbb{R}). \\ \quad (\forall k \in \mathbb{N}.\; (0 < h(k)) \wedge (\operatorname{elimsup}((\lambda j \in \mathbb{N}.\; (h((j + 1)) \cdot \operatorname{recip}(h(j))))) < 1)) \Rightarrow  \\ \quad \quad \operatorname{series\text{-}converges}(h) \end{array}}}
\vnbprop{ratio-test-diverges-liminf}{\vnbstmt{\begin{array}{@{}l@{}} \forall h \in (\mathbb{N} \to \mathbb{R}). \\ \quad (\forall k \in \mathbb{N}.\; (0 < h(k)) \wedge (1 < \operatorname{eliminf}((\lambda j \in \mathbb{N}.\; (h((j + 1)) \cdot \operatorname{recip}(h(j))))))) \Rightarrow  \\ \quad \quad \neg \operatorname{series\text{-}converges}(h) \end{array}}}
\vnbprop{ratio-test}{\vnbstmt{\begin{array}{@{}l@{}} \forall h \in (\mathbb{N} \to \mathbb{R}). \\ \quad \forall k \in \mathbb{N}.\; (0 < h(k)) \Rightarrow  \\ \quad \quad (\operatorname{elimsup}((\lambda j \in \mathbb{N}.\; (h((j + 1)) \cdot \operatorname{recip}(h(j))))) < 1) \Rightarrow  \\ \quad \quad \quad \operatorname{series\text{-}converges}(h) \wedge \\ \quad \quad (1 < \operatorname{eliminf}((\lambda j \in \mathbb{N}.\; (h((j + 1)) \cdot \operatorname{recip}(h(j)))))) \Rightarrow  \\ \quad \quad \quad \neg \operatorname{series\text{-}converges}(h) \end{array}}}

\vnbbookfile{theorem-library/mb12-snow}{theorem-library/mb12-snow}
\noindent{\small mb12-snow.scm -- construction-of-measures.pdf Appendix A.3.2, the snowflake construction: Lemma A.10 ((a + b)\textasciicircum{}s $<$= a\textasciicircum{}s + b\textasciicircum{}s for 0 $<$ s $<$= 1, a, b $>$= 0) and Proposition A.9 (d\textasciicircum{}s is a metric).  Agent ISOM (brief M-B1B2), 2026-10-26.}\par
\vnbprop{rpow-one-base}{\vnbstmt{\forall c \in \mathbb{R}.\; ({1}^{c} = 1)}}
\vnbprop{rpow-le-self-gt1}{\vnbstmt{\begin{array}{@{}l@{}} \forall x, s \in \mathbb{R}. \\ \quad ((1 < x) \wedge ((0 < s) \wedge (s \leq 1))) \Rightarrow ({x}^{s} \leq x) \end{array}}}
\vnbprop{rpow-ge-self-lt1}{\vnbstmt{\begin{array}{@{}l@{}} \forall b, s \in \mathbb{R}. \\ \quad ((0 < b) \wedge ((b < 1) \wedge ((0 < s) \wedge (s \leq 1)))) \Rightarrow (b \leq {b}^{s}) \end{array}}}
\vnbprop{rpow-one-plus-small}{\vnbstmt{\begin{array}{@{}l@{}} \forall t, s \in \mathbb{R}. \\ \quad ((0 < t) \wedge ((t < 1) \wedge ((0 < s) \wedge (s \leq 1)))) \Rightarrow  \\ \quad \quad {(1 + t)}^{s} \leq (1 + {t}^{s}) \end{array}}}
\vnbprop{rpow-factor-bound}{\vnbstmt{\begin{array}{@{}l@{}} \forall x, v, s \in \mathbb{R}. \\ \quad ((0 < x) \wedge ((0 < v) \wedge ({(1 + v)}^{s} \leq (1 + {v}^{s})))) \Rightarrow  \\ \quad \quad {(x \cdot (1 + v))}^{s} \leq ({x}^{s} + {(x \cdot v)}^{s}) \end{array}}}
\vnbprop{rpow-one-plus}{\vnbstmt{\begin{array}{@{}l@{}} \forall t, s \in \mathbb{R}. \\ \quad ((0 < t) \wedge ((0 < s) \wedge (s \leq 1))) \Rightarrow ({(1 + t)}^{s} \leq (1 + {t}^{s})) \end{array}}}
\vnbprop{rpow-subadditive}{\vnbstmt{\begin{array}{@{}l@{}} \forall a, b, s \in \mathbb{R}. \\ \quad ((0 < s) \wedge ((s \leq 1) \wedge ((0 \leq a) \wedge (0 \leq b)))) \Rightarrow  \\ \quad \quad {(a + b)}^{s} \leq ({a}^{s} + {b}^{s}) \end{array}}}
\vnbprop{rpow-mono-le}{\vnbstmt{\begin{array}{@{}l@{}} \forall x, y, e \in \mathbb{R}. \\ \quad ((0 \leq x) \wedge ((x \leq y) \wedge (0 < e))) \Rightarrow ({x}^{e} \leq {y}^{e}) \end{array}}}
\vnbprop{snowflake-lam-type}{\vnbstmt{\begin{array}{@{}l@{}} \forall s, e. \\ \quad (\operatorname{is\text{-}metric\text{-}space}(s) \wedge ((e \in \mathbb{R}) \wedge (0 \leq e))) \Rightarrow  \\ \quad \quad (\lambda {[}\operatorname{snfu\_}, \operatorname{snfv\_}] \in (\operatorname{pts}(s) \times \operatorname{pts}(s)).\; {\operatorname{dist}(s)(\operatorname{snfu\_}, \operatorname{snfv\_})}^{e}) \\ \quad \quad \in\; ((\operatorname{pts}(s) \times \operatorname{pts}(s)) \to \mathbb{R}) \end{array}}}
\vnbprop{snowflake-pts}{\vnbstmt{\begin{array}{@{}l@{}} \forall s, e. \\ \quad (\operatorname{is\text{-}metric\text{-}space}(s) \wedge ((e \in \mathbb{R}) \wedge (0 \leq e))) \Rightarrow (\operatorname{pts}(\operatorname{snowflake}(s, e)) \simeq \operatorname{pts}(s)) \end{array}}}
\vnbprop{snowflake-dist-fun}{\vnbstmt{\begin{array}{@{}l@{}} \forall s, e. \\ \quad (\operatorname{is\text{-}metric\text{-}space}(s) \wedge ((e \in \mathbb{R}) \wedge (0 \leq e))) \Rightarrow  \\ \quad \quad (\operatorname{dist}(\operatorname{snowflake}(s, e)) \simeq (\lambda {[}\operatorname{snfu\_}, \operatorname{snfv\_}] \in (\operatorname{pts}(s) \times \operatorname{pts}(s)).\; {\operatorname{dist}(s)(\operatorname{snfu\_}, \operatorname{snfv\_})}^{e})) \end{array}}}
\vnbprop{snowflake-dist}{\vnbstmt{\begin{array}{@{}l@{}} \forall s, e.; \forall u, v \in \operatorname{pts}(s). \\ \quad (\operatorname{is\text{-}metric\text{-}space}(s) \wedge ((e \in \mathbb{R}) \wedge (0 \leq e))) \Rightarrow  \\ \quad \quad (\operatorname{dist}(\operatorname{snowflake}(s, e))(u, v) \simeq {\operatorname{dist}(s)(u, v)}^{e}) \end{array}}}
\vnbprop{snowflake-is-metric-space}{\vnbstmt{\begin{array}{@{}l@{}} \forall s, e. \\ \quad (\operatorname{is\text{-}metric\text{-}space}(s) \wedge ((e \in \mathbb{R}) \wedge ((0 < e) \wedge (e \leq 1)))) \Rightarrow  \\ \quad \quad \operatorname{is\text{-}metric\text{-}space}(\operatorname{snowflake}(s, e)) \end{array}}}

\vnbbookfile{theorem-library/rpow-defined}{theorem-library/rpow-defined}
\noindent{\small theorem-library/rpow-defined.scm -- THE AXIOMATIC REAL POWER RETIRED: RPOW is DEFINED as RPOW-STAR, and its laws are THEOREMS under their old names and with their old statements (the user's decision, 2026-10-03).}\par
\vnbdef{rpow-def}{\vnbstmt{\forall a, b.\; (\operatorname{rpow}(a, b) \simeq {a}^{b})}}
\vnbprop{rpow-pos}{\vnbstmt{\forall a \in \mathbb{R}, b \in \mathbb{Q}.\; ((0 < a) \Rightarrow ((\operatorname{rpow}(a, b) \in \mathbb{R}) \wedge (0 < \operatorname{rpow}(a, b))))}}
\vnbprop{rpow-zero}{\vnbstmt{\forall a \in \mathbb{R}.\; ((0 < a) \Rightarrow (\operatorname{rpow}(a, 0) = 1))}}
\vnbprop{rpow-one}{\vnbstmt{\forall a \in \mathbb{R}.\; ((0 < a) \Rightarrow (\operatorname{rpow}(a, 1) = a))}}
\vnbprop{rpow-add}{\vnbstmt{\begin{array}{@{}l@{}} \forall a \in \mathbb{R}, b, d \in \mathbb{Q}. \\ \quad (0 < a) \Rightarrow (\operatorname{rpow}(a, (b + d)) = (\operatorname{rpow}(a, b) \cdot \operatorname{rpow}(a, d))) \end{array}}}
\vnbprop{rpow-mul-base}{\vnbstmt{\begin{array}{@{}l@{}} \forall a, c \in \mathbb{R}, b \in \mathbb{Q}. \\ \quad ((0 < a) \wedge (0 < c)) \Rightarrow (\operatorname{rpow}((a \cdot c), b) = (\operatorname{rpow}(a, b) \cdot \operatorname{rpow}(c, b))) \end{array}}}
\vnbprop{rpow-pow}{\vnbstmt{\begin{array}{@{}l@{}} \forall a \in \mathbb{R}, b, d \in \mathbb{Q}. \\ \quad (0 < a) \Rightarrow (\operatorname{rpow}(\operatorname{rpow}(a, b), d) = \operatorname{rpow}(a, (b \cdot d))) \end{array}}}
\vnbprop{rpow-neg}{\vnbstmt{\forall a \in \mathbb{R}, b \in \mathbb{Q}.\; ((0 < a) \Rightarrow (\operatorname{rpow}(a, (0 - b)) = \operatorname{recip}(\operatorname{rpow}(a, b))))}}
\vnbprop{rpow-nat}{\vnbstmt{\forall a \in \mathbb{R}, n \in \mathbb{N}.\; ((0 < a) \Rightarrow (\operatorname{rpow}(a, n) = \operatorname{power}(a, n)))}}
\vnbprop{rpow-zero-base}{\vnbstmt{\forall b \in \mathbb{Q}.\; ((0 < b) \Rightarrow (\operatorname{rpow}(0, b) = 0))}}
\vnbprop{rpow-zero-zero}{\vnbstmt{\operatorname{rpow}(0, 0) = 1}}
\vnbprop{rpow-mono-base}{\vnbstmt{\begin{array}{@{}l@{}} \forall a, c \in \mathbb{R}, b \in \mathbb{Q}. \\ \quad ((0 < a) \wedge ((a \leq c) \wedge (0 \leq b))) \Rightarrow (\operatorname{rpow}(a, b) \leq \operatorname{rpow}(c, b)) \end{array}}}
\vnbprop{rpow-mono-exp-ge1}{\vnbstmt{\begin{array}{@{}l@{}} \forall a \in \mathbb{R}, b, d \in \mathbb{Q}. \\ \quad ((1 \leq a) \wedge (b \leq d)) \Rightarrow (\operatorname{rpow}(a, b) \leq \operatorname{rpow}(a, d)) \end{array}}}
\vnbprop{rpow-mono-exp-le1}{\vnbstmt{\begin{array}{@{}l@{}} \forall a \in \mathbb{R}, b, d \in \mathbb{Q}. \\ \quad (((0 < a) \wedge (a \leq 1)) \wedge (b \leq d)) \Rightarrow (\operatorname{rpow}(a, d) \leq \operatorname{rpow}(a, b)) \end{array}}}
\vnbprop{sqrt-rpow}{\vnbstmt{\forall a \in \mathbb{R}.\; ((0 < a) \Rightarrow (\sqrt{a} = \operatorname{rpow}(a, /(1, 2))))}}
\vnbprop{amgm-2-sqrt}{\vnbstmt{\begin{array}{@{}l@{}} \forall a, b \in \mathbb{R}. \\ \quad ((0 \leq a) \wedge (0 \leq b)) \Rightarrow (\sqrt{(a \cdot b)} \leq /((a + b), 2)) \end{array}}}
\vnbprop{young-inequality}{\vnbstmt{\begin{array}{@{}l@{}} \forall a, b \in \mathbb{R}, p, q \in \mathbb{Q}. \\ \quad ((0 < a) \wedge ((0 < b) \wedge ((1 < p) \wedge ((1 < q) \wedge ((/(1, p) + /(1, q)) = 1))))) \Rightarrow  \\ \quad \quad (a \cdot b) \leq (/(\operatorname{rpow}(a, p), p) + /(\operatorname{rpow}(b, q), q)) \end{array}}}
\vnbprop{bernoulli-rpow}{\vnbstmt{\begin{array}{@{}l@{}} \forall x \in \mathbb{R}, b \in \mathbb{Q}. \\ \quad ((0 < (1 + x)) \wedge (1 \leq b)) \Rightarrow ((1 + (b \cdot x)) \leq \operatorname{rpow}((1 + x), b)) \end{array}}}

\vnbbookfile{theorem-library/rake-asm-b1}{theorem-library/rake-asm-b1}
\noindent{\small theorem-library/rake-asm-b1.scm -- THE FOUR FINITE-SUM INEQUALITIES OF analysis-inequalities.scm, PROVEN under their own names and statements (RAKE-ASM-B, agent RB-1, 2026-10-12):}\par
\vnbprop{rb1-fs-empty}{\vnbstmt{\begin{array}{@{}l@{}} \forall f. \\ \quad \operatorname{finsum}(\operatorname{commutative\text{-}ring\text{-}additive\text{-}ag}(\operatorname{rr\text{-}normed\text{-}field}), f, \emptyset) = 0 \end{array}}}
\vnbprop{rb1-fs-real}{\vnbstmt{\begin{array}{@{}l@{}} \forall s \in \mathbf{Set}, f. \\ \quad ((|s| \in \mathbb{N}) \wedge \forall z \in s.\; (f(z) \in \mathbb{R})) \Rightarrow  \\ \quad \quad \operatorname{finsum}(\operatorname{commutative\text{-}ring\text{-}additive\text{-}ag}(\operatorname{rr\text{-}normed\text{-}field}), f, s) \in \mathbb{R} \end{array}}}
\vnbprop{rb1-fs-ins}{\vnbstmt{\begin{array}{@{}l@{}} \forall s, k \in \mathbf{Set}, f. \\ \quad ((|s| \in \mathbb{N}) \wedge (\neg (k \in s) \wedge \forall z \in (s \cup \{k\}).\; (f(z) \in \mathbb{R}))) \Rightarrow  \\ \quad \quad \operatorname{finsum}(\operatorname{commutative\text{-}ring\text{-}additive\text{-}ag}(\operatorname{rr\text{-}normed\text{-}field}), f, (s \cup \{k\})) \\ \quad \quad =\; \operatorname{finsum}(\operatorname{commutative\text{-}ring\text{-}additive\text{-}ag}(\operatorname{rr\text{-}normed\text{-}field}), f, s) + f(k) \end{array}}}
\vnbprop{rb1-fs-lin}{\vnbstmt{\begin{array}{@{}l@{}} \forall s \in \mathbf{Set}, a, b \in \mathbb{R}, f, g, h. \\ \quad |s| \in \mathbb{N} \wedge \\ \quad \forall z \in s.\; (f(z) \in \mathbb{R}) \wedge \\ \quad \forall z \in s.\; (g(z) \in \mathbb{R}) \wedge \\ \quad \forall z \in s.\; (h(z) = ((a \cdot f(z)) + (b \cdot g(z)))) \Rightarrow \\ \quad \quad \operatorname{finsum}(\operatorname{commutative\text{-}ring\text{-}additive\text{-}ag}(\operatorname{rr\text{-}normed\text{-}field}), h, s) \\ \quad \quad =\; (a \cdot \operatorname{finsum}(\operatorname{commutative\text{-}ring\text{-}additive\text{-}ag}(\operatorname{rr\text{-}normed\text{-}field}), f, s)) + (b \cdot \operatorname{finsum}(\operatorname{commutative\text{-}ring\text{-}additive\text{-}ag}(\operatorname{rr\text{-}normed\text{-}field}), g, s)) \end{array}}}
\vnbprop{rb1-disc}{\vnbstmt{\begin{array}{@{}l@{}} \forall a, b, c \in \mathbb{R}. \\ \quad \forall l, m \in \mathbb{R}.\; (0 \leq ((((l \cdot l) \cdot a) + (-(2 \cdot (l \cdot m)) \cdot c)) + ((m \cdot m) \cdot b))) \Rightarrow  \\ \quad \quad (c \cdot c) \leq (a \cdot b) \end{array}}}
\vnbprop{cauchy-schwarz-finite}{\vnbstmt{\begin{array}{@{}l@{}} \forall s \in \mathbf{Set}, a, b \in (s \to \mathbb{R}). \\ \quad (|s| \in \mathbb{N}) \Rightarrow  \\ \quad \quad \operatorname{finsum}(\operatorname{commutative\text{-}ring\text{-}additive\text{-}ag}(\operatorname{rr\text{-}normed\text{-}field}), (\lambda i \in s.\; (a(i) \cdot b(i))), s) \cdot \operatorname{finsum}(\operatorname{commutative\text{-}ring\text{-}additive\text{-}ag}(\operatorname{rr\text{-}normed\text{-}field}), (\lambda i \in s.\; (a(i) \cdot b(i))), s) \\ \quad \quad \leq\; \operatorname{finsum}(\operatorname{commutative\text{-}ring\text{-}additive\text{-}ag}(\operatorname{rr\text{-}normed\text{-}field}), (\lambda i \in s.\; (a(i) \cdot a(i))), s) \cdot \operatorname{finsum}(\operatorname{commutative\text{-}ring\text{-}additive\text{-}ag}(\operatorname{rr\text{-}normed\text{-}field}), (\lambda i \in s.\; (b(i) \cdot b(i))), s) \end{array}}}
\vnbprop{cauchy-schwarz-sqrt}{\vnbstmt{\begin{array}{@{}l@{}} \forall s \in \mathbf{Set}, a, b \in (s \to \mathbb{R}). \\ \quad (|s| \in \mathbb{N}) \Rightarrow  \\ \quad \quad \lvert \operatorname{finsum}(\operatorname{commutative\text{-}ring\text{-}additive\text{-}ag}(\operatorname{rr\text{-}normed\text{-}field}), (\lambda i \in s.\; (a(i) \cdot b(i))), s) \rvert \\ \quad \quad \leq\; \sqrt{\operatorname{finsum}(\operatorname{commutative\text{-}ring\text{-}additive\text{-}ag}(\operatorname{rr\text{-}normed\text{-}field}), (\lambda i \in s.\; (a(i) \cdot a(i))), s)} \cdot \sqrt{\operatorname{finsum}(\operatorname{commutative\text{-}ring\text{-}additive\text{-}ag}(\operatorname{rr\text{-}normed\text{-}field}), (\lambda i \in s.\; (b(i) \cdot b(i))), s)} \end{array}}}
\vnbprop{minkowski-l2}{\vnbstmt{\begin{array}{@{}l@{}} \forall s \in \mathbf{Set}, a, b \in (s \to \mathbb{R}). \\ \quad (|s| \in \mathbb{N}) \Rightarrow  \\ \quad \quad \sqrt{\operatorname{finsum}(\operatorname{commutative\text{-}ring\text{-}additive\text{-}ag}(\operatorname{rr\text{-}normed\text{-}field}), (\lambda i \in s.\; ((a(i) + b(i)) \cdot (a(i) + b(i)))), s)} \\ \quad \quad \leq\; \sqrt{\operatorname{finsum}(\operatorname{commutative\text{-}ring\text{-}additive\text{-}ag}(\operatorname{rr\text{-}normed\text{-}field}), (\lambda i \in s.\; (a(i) \cdot a(i))), s)} + \sqrt{\operatorname{finsum}(\operatorname{commutative\text{-}ring\text{-}additive\text{-}ag}(\operatorname{rr\text{-}normed\text{-}field}), (\lambda i \in s.\; (b(i) \cdot b(i))), s)} \end{array}}}
\vnbprop{rb1-rpow-nn}{\vnbstmt{\begin{array}{@{}l@{}} \forall x \in \mathbb{R}, r \in \mathbb{Q}. \\ \quad ((0 \leq x) \wedge (0 < r)) \Rightarrow ((\operatorname{rpow}(x, r) \in \mathbb{R}) \wedge (0 \leq \operatorname{rpow}(x, r))) \end{array}}}
\vnbprop{rb1-rpow-zero}{\vnbstmt{\forall x \in \mathbb{R}, r \in \mathbb{Q}.\; (((0 \leq x) \wedge ((0 < r) \wedge (\operatorname{rpow}(x, r) = 0))) \Rightarrow (x = 0))}}
\vnbprop{rb1-rpow-scale}{\vnbstmt{\begin{array}{@{}l@{}} \forall x, c \in \mathbb{R}, r \in \mathbb{Q}. \\ \quad ((0 \leq x) \wedge ((0 < c) \wedge (0 < r))) \Rightarrow  \\ \quad \quad \operatorname{rpow}((x \cdot \operatorname{recip}(c)), r) = (\operatorname{rpow}(x, r) \cdot \operatorname{recip}(\operatorname{rpow}(c, r))) \end{array}}}
\vnbprop{rb1-young0}{\vnbstmt{\begin{array}{@{}l@{}} \forall p, q \in \mathbb{Q}, u, v \in \mathbb{R}. \\ \quad ((1 < p) \wedge ((1 < q) \wedge (((/(1, p) + /(1, q)) = 1) \wedge ((0 \leq u) \wedge (0 \leq v))))) \Rightarrow  \\ \quad \quad (u \cdot v) \leq (/(\operatorname{rpow}(u, p), p) + /(\operatorname{rpow}(v, q), q)) \end{array}}}
\vnbprop{rb1-young-scaled}{\vnbstmt{\begin{array}{@{}l@{}} \forall p, q \in \mathbb{Q}, s, t, x, y \in \mathbb{R}. \\ \quad ((1 < p) \wedge ((1 < q) \wedge (((/(1, p) + /(1, q)) = 1) \wedge ((0 < s) \wedge ((0 < t) \wedge ((0 \leq x) \wedge (0 \leq y))))))) \Rightarrow  \\ \quad \quad x \cdot y \\ \quad \quad \leq\; ((((s \cdot t) \cdot \operatorname{recip}(\operatorname{rpow}(s, p))) \cdot \operatorname{recip}(p)) \cdot \operatorname{rpow}(x, p)) + ((((s \cdot t) \cdot \operatorname{recip}(\operatorname{rpow}(t, q))) \cdot \operatorname{recip}(q)) \cdot \operatorname{rpow}(y, q)) \end{array}}}
\vnbprop{rb1-fs-term-le}{\vnbstmt{\begin{array}{@{}l@{}} \forall s \in \mathbf{Set}, f. \\ \quad ((|s| \in \mathbb{N}) \wedge (\forall z \in s.\; (f(z) \in \mathbb{R}) \wedge \forall z \in s.\; (0 \leq f(z)))) \Rightarrow  \\ \quad \quad 0 \leq \operatorname{finsum}(\operatorname{commutative\text{-}ring\text{-}additive\text{-}ag}(\operatorname{rr\text{-}normed\text{-}field}), f, s) \wedge \\ \quad \quad \forall k \in s. \\ \quad \quad \quad f(k) \leq \operatorname{finsum}(\operatorname{commutative\text{-}ring\text{-}additive\text{-}ag}(\operatorname{rr\text{-}normed\text{-}field}), f, s) \end{array}}}
\vnbprop{holder-finite}{\vnbstmt{\begin{array}{@{}l@{}} \forall s \in \mathbf{Set}, a, b \in (s \to \mathbb{R}), p, q \in \mathbb{Q}. \\ \quad ((|s| \in \mathbb{N}) \wedge ((1 < p) \wedge ((1 < q) \wedge ((/(1, p) + /(1, q)) = 1)))) \Rightarrow  \\ \quad \quad \operatorname{finsum}(\operatorname{commutative\text{-}ring\text{-}additive\text{-}ag}(\operatorname{rr\text{-}normed\text{-}field}), (\lambda i \in s.\; \lvert (a(i) \cdot b(i)) \rvert), s) \\ \quad \quad \leq\; \operatorname{rpow}(\operatorname{finsum}(\operatorname{commutative\text{-}ring\text{-}additive\text{-}ag}(\operatorname{rr\text{-}normed\text{-}field}), (\lambda i \in s.\; \operatorname{rpow}(\lvert a(i) \rvert, p)), s), /(1, p)) \cdot \operatorname{rpow}(\operatorname{finsum}(\operatorname{commutative\text{-}ring\text{-}additive\text{-}ag}(\operatorname{rr\text{-}normed\text{-}field}), (\lambda i \in s.\; \operatorname{rpow}(\lvert b(i) \rvert, q)), s), /(1, q)) \end{array}}}

\vnbbookfile{theorem-library/rps-finsum-mono-set}{theorem-library/rps-finsum-mono-set}
\noindent{\small rps-finsum-mono-set.scm -- lemmas MOVED 2026-11-10 (integration of the 11-09/10 waves) out of theorem-library/measure-laws.scm into library analysis, for the library gate: cited by sums-to-algebra; RR-POS-STAR finite sums. The proofs are the originals, verbatim; the helper defines are a copy of the source's. ...}\par
\vnbprop{rps-finsum-mono-set}{\vnbstmt{\begin{array}{@{}l@{}} \forall f, s.; \forall t \in \mathbf{Set}. \\ \quad ((|t| \in \mathbb{N}) \wedge ((s \subseteq t) \wedge \forall z \in t.\; (f(z) \in \operatorname{rr\text{-}pos\text{-}star}))) \Rightarrow  \\ \quad \quad \operatorname{finsum}(\operatorname{rr\text{-}pos\text{-}star\text{-}add\text{-}monoid}, f, s) \leq \operatorname{finsum}(\operatorname{rr\text{-}pos\text{-}star\text{-}add\text{-}monoid}, f, t) \end{array}}}

\vnbbookfile{theorem-library/rps-finsum-le-ptwise}{theorem-library/rps-finsum-le-ptwise}
\noindent{\small rps-finsum-le-ptwise.scm -- lemmas MOVED 2026-11-10 (integration of the 11-09/10 waves) out of theorem-library/integral-laws.scm into library analysis, for the library gate: cited by l2-families-2; RR-POS-STAR finite sums. The proofs are the originals, verbatim; the helper defines are a copy of the source's. ...}\par
\vnbprop{rps-finsum-le-ptwise}{\vnbstmt{\begin{array}{@{}l@{}} \forall s \in \mathbf{Set}, f, g. \\ \quad |s| \in \mathbb{N} \wedge \\ \quad \forall z \in s.\; (f(z) \in \operatorname{rr\text{-}pos\text{-}star}) \wedge \\ \quad \forall z \in s.\; (g(z) \in \operatorname{rr\text{-}pos\text{-}star}) \wedge \\ \quad \forall z \in s.\; (f(z) \leq g(z)) \Rightarrow \\ \quad \quad \operatorname{finsum}(\operatorname{rr\text{-}pos\text{-}star\text{-}add\text{-}monoid}, f, s) \leq \operatorname{finsum}(\operatorname{rr\text{-}pos\text{-}star\text{-}add\text{-}monoid}, g, s) \end{array}}}

\vnbbookfile{theorem-library/ps-series-bridges}{theorem-library/ps-series-bridges}
\noindent{\small ps-series-bridges.scm -- THE THREE PS/SERIES BRIDGES, PROVEN, and the CONVERGES twin of rr-limit-ptwise-eq that the third one needed.}\par
\vnbprop{ps-abs-term}{\vnbstmt{\begin{array}{@{}l@{}} \forall f \in (\mathbb{N} \to \mathbb{R}), x \in \mathbb{R}, n \in \mathbb{N}. \\ \quad \lvert (f(n) \cdot \operatorname{power}(x, n)) \rvert = (\lvert f(n) \rvert \cdot \operatorname{power}(\lvert x \rvert, n)) \end{array}}}
\vnbprop{ps-partial-sum-as-series}{\vnbstmt{\begin{array}{@{}l@{}} \forall f \in (\mathbb{N} \to \mathbb{R}), x \in \mathbb{R}, k \in \mathbb{N}. \\ \quad \operatorname{ps\text{-}partial\text{-}sum}(f, x, k) = \operatorname{series\text{-}partial\text{-}sum}((\lambda n \in \mathbb{N}.\; (f(n) \cdot \operatorname{power}(x, n))), k) \end{array}}}
\vnbprop{rr-converges-ptwise-eq}{\vnbstmt{\begin{array}{@{}l@{}} \forall f, h \in (\mathbb{N} \to \mathbb{R}). \\ \quad (\forall j \in \mathbb{N}.\; (h(j) = f(j)) \wedge \operatorname{converges}(\operatorname{rr\text{-}ms}, f)) \Rightarrow  \\ \quad \quad \operatorname{converges}(\operatorname{rr\text{-}ms}, h) \end{array}}}
\vnbprop{ps-converges-as-series}{\vnbstmt{\begin{array}{@{}l@{}} \forall f \in (\mathbb{N} \to \mathbb{R}), x \in \mathbb{R}. \\ \quad \operatorname{ps\text{-}converges\text{-}at}(f, x) \Leftrightarrow \operatorname{series\text{-}converges}((\lambda n \in \mathbb{N}.\; (f(n) \cdot \operatorname{power}(x, n)))) \end{array}}}

\vnbbookfile{theorem-library/ratio-test-proof}{theorem-library/ratio-test-proof}
\noindent{\small ratio-test-proof.scm -- THE RATIO TEST FOR A REAL POWER SERIES (ratio-test-converges, structure-library/power-series.scm), PROVEN. Batch 27-A, 2026-09-24.}\par
\vnbprop{elimsup-eventually-le}{\vnbstmt{\begin{array}{@{}l@{}} \forall f.; \forall r \in \mathbb{R}, n \in \mathbb{N}. \\ \quad \forall k \in \mathbb{N}.\; (f(k) \in \mathbb{R}) \wedge \\ \quad \forall k \in \mathbb{N}.\; (0 \leq f(k)) \wedge \\ \quad \forall a \in \mathbb{N}.\; ((n \leq a) \Rightarrow (f(a) \leq r)) \Rightarrow \\ \quad \quad \operatorname{elimsup}(f) \leq r \end{array}}}
\vnbprop{ratio-test-converges}{\vnbstmt{\begin{array}{@{}l@{}} \forall f \in (\mathbb{N} \to \mathbb{R}), l, x \in \mathbb{R}. \\ \quad (((0 \leq l) \wedge (\forall n \in \mathbb{N}.\; \neg (f(n) = 0) \wedge \operatorname{ps\text{-}ratio\text{-}limit}(f, l))) \wedge ((\lvert x \rvert \cdot l) < 1)) \Rightarrow  \\ \quad \quad \operatorname{ps\text{-}converges\text{-}at}(f, x) \end{array}}}

\vnbbookfile{theorem-library/rake-series}{theorem-library/rake-series}
\noindent{\small rake-series.scm -- rake batch 5 T (2026-09-18): the little-o algebra, the power-series absolute-convergence bridge, and the Taylor clearing step.}\par
\vnbprop{vtaylor-clear}{\vnbstmt{\begin{array}{@{}l@{}} \forall v, a, b, w \in \mathbb{R}, n \in \mathbb{N}. \\ \quad (((\operatorname{factorial}((n + 1)) \cdot v) = (a \cdot w)) \wedge ((a \leq b) \wedge (0 \leq w))) \Rightarrow  \\ \quad \quad (\operatorname{factorial}((n + 1)) \cdot v) \leq (b \cdot w) \end{array}}}
\vnbprop{little-o-scalar}{\vnbstmt{\begin{array}{@{}l@{}} \forall c, g, a. \\ \quad ((c \in \mathbb{R}) \wedge \operatorname{little\text{-}o\text{-}at}(g, a)) \Rightarrow  \\ \quad \quad \operatorname{little\text{-}o\text{-}at}((\lambda x \in \mathbb{R}.\; (c \cdot g(x))), a) \end{array}}}
\vnbprop{little-o-sum}{\vnbstmt{\begin{array}{@{}l@{}} \forall g, h, a. \\ \quad (\operatorname{little\text{-}o\text{-}at}(g, a) \wedge \operatorname{little\text{-}o\text{-}at}(h, a)) \Rightarrow  \\ \quad \quad \operatorname{little\text{-}o\text{-}at}((\lambda x \in \mathbb{R}.\; (g(x) + h(x))), a) \end{array}}}
\vnbprop{shift-lam-in-fun}{\vnbstmt{\forall k \in \mathbb{R}, f \in (\mathbb{R} \to \mathbb{R}).\; ((\lambda x \in \mathbb{R}.\; (f(x) + k)) \in (\mathbb{R} \to \mathbb{R}))}}
\vnbprop{shift-continuous-at}{\vnbstmt{\begin{array}{@{}l@{}} \forall k \in \mathbb{R}, f, a. \\ \quad \operatorname{is\text{-}continuous\text{-}at}(\operatorname{rr\text{-}ms}, \operatorname{rr\text{-}ms}, f, a) \Rightarrow  \\ \quad \quad \operatorname{is\text{-}continuous\text{-}at}(\operatorname{rr\text{-}ms}, \operatorname{rr\text{-}ms}, (\lambda x \in \mathbb{R}.\; (f(x) + k)), a) \end{array}}}
\vnbprop{diff-iff-little-o}{\vnbstmt{\begin{array}{@{}l@{}} \forall f, a, l. \\ \quad \operatorname{is\text{-}diff\text{-}at}(f, a, l) \\ \quad \Leftrightarrow\; f \in (\mathbb{R} \to \mathbb{R}) \wedge \\ \quad a \in \mathbb{R} \wedge \\ \quad l \in \mathbb{R} \wedge \\ \quad \operatorname{little\text{-}o\text{-}at}((\lambda x \in \mathbb{R}.\; ((f(x) - f(a)) - (l \cdot (x - a)))), a) \end{array}}}
\vnbprop{ps-absolute-implies-convergent}{\vnbstmt{\begin{array}{@{}l@{}} \forall f \in (\mathbb{N} \to \mathbb{R}), x \in \mathbb{R}. \\ \quad \operatorname{ps\text{-}absolutely\text{-}converges\text{-}at}(f, x) \Rightarrow \operatorname{ps\text{-}converges\text{-}at}(f, x) \end{array}}}
\vnbprop{esum-finite-implies-bounded}{\vnbstmt{\begin{array}{@{}l@{}} \forall f \in (\operatorname{dom}(f) \to \operatorname{rr\text{-}pos\text{-}star}). \\ \quad (\operatorname{esum}(f) \in \mathbb{R}) \Rightarrow  \\ \quad \quad \exists m \in \mathbb{R}. \forall s \in \mathbf{Set}. \\ \quad \quad \quad ((|s| \in \mathbb{N}) \wedge (s \subseteq \operatorname{dom}(f))) \Rightarrow  \\ \quad \quad \quad \quad \operatorname{finsum}(\operatorname{rr\text{-}pos\text{-}star\text{-}add\text{-}monoid}, f, s) \leq m \end{array}}}

\vnbbookfile{theorem-library/rake-series2}{theorem-library/rake-series2}
\noindent{\small rake-series2.scm -- rake batch 5c (2026-09-18), assignment 5c-R: the two series leaves batch T left BLOCKED, and the four bricks they were missing.}\par
\vnbprop{metric-chain-bound}{\vnbstmt{\begin{array}{@{}l@{}} \forall n \in \mathbb{N}, s, f, d, m. \\ \quad \operatorname{is\text{-}metric\text{-}space}(s) \wedge \\ \quad f \in (\mathbb{N} \to \operatorname{pts}(s)) \wedge \\ \quad d \in (\mathbb{N} \to \mathbb{R}) \wedge \\ \quad \forall k \in \mathbb{N}.\; (\operatorname{dist}(s)(f(k), f((k + 1))) \leq d(k)) \wedge \\ \quad m \in \mathbb{N} \wedge \\ \quad m \leq n \Rightarrow \\ \quad \quad \operatorname{dist}(s)(f(m), f(n)) \leq (\operatorname{series\text{-}partial\text{-}sum}(d, n) - \operatorname{series\text{-}partial\text{-}sum}(d, m)) \end{array}}}
\vnbprop{metric-chain-triangle}{\vnbstmt{\begin{array}{@{}l@{}} \forall s, f.; \forall m, n \in \mathbb{N}. \\ \quad ((\operatorname{is\text{-}metric\text{-}space}(s) \wedge (f \in (\mathbb{N} \to \operatorname{pts}(s)))) \wedge (m \leq n)) \Rightarrow  \\ \quad \quad \operatorname{dist}(s)(f(m), f(n)) \\ \quad \quad \leq\; (\operatorname{series\text{-}partial\text{-}sum}((\lambda k \in \mathbb{N}.\; \operatorname{dist}(s)(f(k), f((k + 1)))), n) - \operatorname{series\text{-}partial\text{-}sum}((\lambda k \in \mathbb{N}.\; \operatorname{dist}(s)(f(k), f((k + 1)))), m)) \end{array}}}
\vnbprop{summable-bound-implies-cauchy}{\vnbstmt{\begin{array}{@{}l@{}} \forall s, f, d. \\ \quad \operatorname{is\text{-}metric\text{-}space}(s) \wedge \\ \quad f \in (\mathbb{N} \to \operatorname{pts}(s)) \wedge \\ \quad d \in (\mathbb{N} \to \mathbb{R}) \wedge \\ \quad \operatorname{series\text{-}converges}(d) \wedge \\ \quad \forall k \in \mathbb{N}.\; (\operatorname{dist}(s)(f(k), f((k + 1))) \leq d(k)) \Rightarrow \\ \quad \quad \operatorname{is\text{-}cauchy\text{-}seq}(s, f) \end{array}}}
\vnbprop{rr-power-le-one}{\vnbstmt{\forall k \in \mathbb{N}, r \in \mathbb{R}.\; (((0 \leq r) \wedge (r \leq 1)) \Rightarrow (\operatorname{power}(r, k) \leq 1))}}
\vnbprop{rr-power-shrink-bound}{\vnbstmt{\forall k \in \mathbb{N}, r \in \mathbb{R}.\; (((0 \leq r) \wedge (r \leq 1)) \Rightarrow (((k \cdot (1 - r)) \cdot \operatorname{power}(r, k)) \leq 1))}}
\vnbprop{power-small-threshold}{\vnbstmt{\begin{array}{@{}l@{}} \forall r \in \mathbb{R}, s. \\ \quad (((0 \leq r) \wedge (r < 1)) \wedge \operatorname{pos\text{-}rr}(s)) \Rightarrow  \\ \quad \quad \exists p \in \mathbb{N}.\; \forall k \in \mathbb{N}.\; ((p \leq k) \Rightarrow (\operatorname{power}(r, k) \leq s)) \end{array}}}
\vnbprop{power-tends-to-zero}{\vnbstmt{\begin{array}{@{}l@{}} \forall r \in \mathbb{R}. \\ \quad (\lvert r \rvert < 1) \Rightarrow \operatorname{converges\text{-}to}(\operatorname{rr\text{-}ms}, (\lambda n \in \mathbb{N}.\; \operatorname{power}(r, n)), 0) \end{array}}}
\vnbprop{rr-const-converges-to}{\vnbstmt{\forall c \in \mathbb{R}.\; \operatorname{converges\text{-}to}(\operatorname{rr\text{-}ms}, (\lambda n \in \mathbb{N}.\; c), c)}}
\vnbprop{geometric-series-converges-to}{\vnbstmt{\begin{array}{@{}l@{}} \forall r \in \mathbb{R}. \\ \quad (\lvert r \rvert < 1) \Rightarrow \operatorname{series\text{-}converges\text{-}to}((\lambda n \in \mathbb{N}.\; \operatorname{power}(r, n)), \operatorname{recip}((1 - r))) \end{array}}}

\vnbbookfile{theorem-library/heine-borel-interval}{theorem-library/heine-borel-interval}
\noindent{\small heine-borel-interval.scm -- HEINE-BOREL for a closed interval AS A SPACE:}\par
\vnbprop{rr-limit-in-ccint}{\vnbstmt{\begin{array}{@{}l@{}} \forall o, i \in \mathbb{R}, q \in (\mathbb{N} \to \operatorname{ccint}(o, i)), t \in \mathbb{R}. \\ \quad \operatorname{converges\text{-}to}(\operatorname{rr\text{-}ms}, q, t) \Rightarrow (t \in \operatorname{ccint}(o, i)) \end{array}}}
\vnbprop{heine-borel-ccint}{\vnbstmt{\forall v, \mathit{bv} \in \mathbb{R}.\; \operatorname{is\text{-}compact}(\operatorname{subspace\text{-}ms}(\operatorname{rr\text{-}ms}, \operatorname{ccint}(v, \mathit{bv})))}}

\vnbbookfile{theorem-library/cc-coords-laws}{theorem-library/cc-coords-laws}
\noindent{\small cc-coords-laws.scm -- CC and cartesian(RR, RR), explicitly.}\par
\vnbprop{apply-congruence-1}{\vnbstmt{\forall f, g, u.\; ((f \simeq g) \Rightarrow (f(u) \simeq g(u)))}}
\vnbprop{cc-coords-apply}{\vnbstmt{\forall z \in \mathbb{C}.\; (\operatorname{cc\text{-}coords}(z) \simeq {[}\operatorname{real\text{-}part}(z), \operatorname{imag\text{-}part}(z)])}}
\vnbprop{cc-coords-in-fun}{\vnbstmt{\operatorname{cc\text{-}coords} \in (\mathbb{C} \to (\mathbb{R} \times \mathbb{R}))}}
\vnbprop{cc-of-pair-apply}{\vnbstmt{\forall a, b \in \mathbb{R}.\; (\operatorname{cc\text{-}of\text{-}pair}({[}a, b]) \simeq (a + (b \cdot +i)))}}
\vnbprop{cc-of-pair-in-fun}{\vnbstmt{\operatorname{cc\text{-}of\text{-}pair} \in ((\mathbb{R} \times \mathbb{R}) \to \mathbb{C})}}
\vnbprop{cc-of-pair-of-coords}{\vnbstmt{\forall z \in \mathbb{C}.\; (\operatorname{cc\text{-}of\text{-}pair}(\operatorname{cc\text{-}coords}(z)) = z)}}
\vnbprop{cc-coords-of-pair}{\vnbstmt{\forall r \in (\mathbb{R} \times \mathbb{R}).\; (\operatorname{cc\text{-}coords}(\operatorname{cc\text{-}of\text{-}pair}(r)) = r)}}
\vnbprop{cc-coords-is-bijection}{\vnbstmt{\operatorname{cc\text{-}coords} \in \mathrm{Bij}(\mathbb{C}, (\mathbb{R} \times \mathbb{R}))}}
\vnbprop{cc-of-pair-is-bijection}{\vnbstmt{\operatorname{cc\text{-}of\text{-}pair} \in \mathrm{Bij}((\mathbb{R} \times \mathbb{R}), \mathbb{C})}}
\vnbprop{cc-re-add}{\vnbstmt{\forall z, w \in \mathbb{C}.\; (\operatorname{real\text{-}part}((z + w)) = (\operatorname{real\text{-}part}(z) + \operatorname{real\text{-}part}(w)))}}
\vnbprop{cc-im-add}{\vnbstmt{\forall z, w \in \mathbb{C}.\; (\operatorname{imag\text{-}part}((z + w)) = (\operatorname{imag\text{-}part}(z) + \operatorname{imag\text{-}part}(w)))}}
\vnbprop{cc-re-real-mul}{\vnbstmt{\forall c \in \mathbb{R}, z \in \mathbb{C}.\; (\operatorname{real\text{-}part}((c \cdot z)) = (c \cdot \operatorname{real\text{-}part}(z)))}}
\vnbprop{cc-im-real-mul}{\vnbstmt{\forall c \in \mathbb{R}, z \in \mathbb{C}.\; (\operatorname{imag\text{-}part}((c \cdot z)) = (c \cdot \operatorname{imag\text{-}part}(z)))}}
\vnbprop{cc-re-mul}{\vnbstmt{\begin{array}{@{}l@{}} \forall z, w \in \mathbb{C}. \\ \quad \operatorname{real\text{-}part}((z \cdot w)) = ((\operatorname{real\text{-}part}(z) \cdot \operatorname{real\text{-}part}(w)) - (\operatorname{imag\text{-}part}(z) \cdot \operatorname{imag\text{-}part}(w))) \end{array}}}
\vnbprop{cc-im-mul}{\vnbstmt{\begin{array}{@{}l@{}} \forall z, w \in \mathbb{C}. \\ \quad \operatorname{imag\text{-}part}((z \cdot w)) = ((\operatorname{real\text{-}part}(z) \cdot \operatorname{imag\text{-}part}(w)) + (\operatorname{imag\text{-}part}(z) \cdot \operatorname{real\text{-}part}(w))) \end{array}}}
\vnbprop{cc-coords-add}{\vnbstmt{\begin{array}{@{}l@{}} \forall z, w \in \mathbb{C}. \\ \quad \operatorname{cc\text{-}coords}((z + w)) = {[}(\operatorname{real\text{-}part}(z) + \operatorname{real\text{-}part}(w)), (\operatorname{imag\text{-}part}(z) + \operatorname{imag\text{-}part}(w))] \end{array}}}
\vnbprop{cc-coords-real-mul}{\vnbstmt{\begin{array}{@{}l@{}} \forall c \in \mathbb{R}, z \in \mathbb{C}. \\ \quad \operatorname{cc\text{-}coords}((c \cdot z)) = {[}(c \cdot \operatorname{real\text{-}part}(z)), (c \cdot \operatorname{imag\text{-}part}(z))] \end{array}}}
\vnbprop{cc-of-pair-add}{\vnbstmt{\begin{array}{@{}l@{}} \forall a_{1}, b_{1}, a_{2}, b_{2} \in \mathbb{R}. \\ \quad \operatorname{cc\text{-}of\text{-}pair}({[}(a_{1} + a_{2}), (b_{1} + b_{2})]) = (\operatorname{cc\text{-}of\text{-}pair}({[}a_{1}, b_{1}]) + \operatorname{cc\text{-}of\text{-}pair}({[}a_{2}, b_{2}])) \end{array}}}
\vnbprop{cc-of-pair-real-mul}{\vnbstmt{\begin{array}{@{}l@{}} \forall c, a, b \in \mathbb{R}. \\ \quad \operatorname{cc\text{-}of\text{-}pair}({[}(c \cdot a), (c \cdot b)]) = (c \cdot \operatorname{cc\text{-}of\text{-}pair}({[}a, b])) \end{array}}}
\vnbprop{cc-re-dist-le}{\vnbstmt{\forall z, w \in \mathbb{C}.\; (\lvert (\operatorname{real\text{-}part}(z) - \operatorname{real\text{-}part}(w)) \rvert \leq \operatorname{magnitude}((z - w)))}}
\vnbprop{cc-im-dist-le}{\vnbstmt{\forall z, w \in \mathbb{C}.\; (\lvert (\operatorname{imag\text{-}part}(z) - \operatorname{imag\text{-}part}(w)) \rvert \leq \operatorname{magnitude}((z - w)))}}
\vnbprop{cc-dist-le-coords}{\vnbstmt{\begin{array}{@{}l@{}} \forall z, w \in \mathbb{C}. \\ \quad \operatorname{magnitude}((z - w)) \leq (\lvert (\operatorname{real\text{-}part}(z) - \operatorname{real\text{-}part}(w)) \rvert + \lvert (\operatorname{imag\text{-}part}(z) - \operatorname{imag\text{-}part}(w)) \rvert) \end{array}}}
\vnbprop{cc-coords-of-converges}{\vnbstmt{\begin{array}{@{}l@{}} \forall f \in (\mathbb{N} \to \mathbb{C}), g, q \in (\mathbb{N} \to \mathbb{R}), v \in \mathbb{C}. \\ \quad \forall n \in \mathbb{N}.\; (\operatorname{real\text{-}part}(f(n)) = g(n)) \wedge \\ \quad \forall n \in \mathbb{N}.\; (\operatorname{imag\text{-}part}(f(n)) = q(n)) \wedge \\ \quad \operatorname{converges\text{-}to}(\operatorname{cc\text{-}ms}, f, v) \Rightarrow \\ \quad \quad \operatorname{converges\text{-}to}(\operatorname{rr\text{-}ms}, g, \operatorname{real\text{-}part}(v)) \wedge \operatorname{converges\text{-}to}(\operatorname{rr\text{-}ms}, q, \operatorname{imag\text{-}part}(v)) \end{array}}}
\vnbprop{cc-converges-iff-coords}{\vnbstmt{\begin{array}{@{}l@{}} \forall f \in (\mathbb{N} \to \mathbb{C}), g, q \in (\mathbb{N} \to \mathbb{R}), v \in \mathbb{C}. \\ \quad (\forall n \in \mathbb{N}.\; (\operatorname{real\text{-}part}(f(n)) = g(n)) \wedge \forall n \in \mathbb{N}.\; (\operatorname{imag\text{-}part}(f(n)) = q(n))) \Rightarrow  \\ \quad \quad \operatorname{converges\text{-}to}(\operatorname{cc\text{-}ms}, f, v) \\ \quad \quad \Leftrightarrow\; \operatorname{converges\text{-}to}(\operatorname{rr\text{-}ms}, g, \operatorname{real\text{-}part}(v)) \wedge \operatorname{converges\text{-}to}(\operatorname{rr\text{-}ms}, q, \operatorname{imag\text{-}part}(v)) \end{array}}}
\vnbprop{cc-continuous-at-iff-coords}{\vnbstmt{\begin{array}{@{}l@{}} \forall s, f, r, i, a. \\ \quad \operatorname{is\text{-}metric\text{-}space}(s) \wedge \\ \quad f \in (\operatorname{pts}(s) \to \mathbb{C}) \wedge \\ \quad r \in (\operatorname{pts}(s) \to \mathbb{R}) \wedge \\ \quad i \in (\operatorname{pts}(s) \to \mathbb{R}) \wedge \\ \quad a \in \operatorname{pts}(s) \wedge \\ \quad \forall u \in \operatorname{pts}(s).\; (\operatorname{real\text{-}part}(f(u)) = r(u)) \wedge \\ \quad \forall u \in \operatorname{pts}(s).\; (\operatorname{imag\text{-}part}(f(u)) = i(u)) \Rightarrow \\ \quad \quad \operatorname{is\text{-}continuous\text{-}at}(s, \operatorname{cc\text{-}ms}, f, a) \\ \quad \quad \Leftrightarrow\; \operatorname{is\text{-}continuous\text{-}at}(s, \operatorname{rr\text{-}ms}, r, a) \wedge \operatorname{is\text{-}continuous\text{-}at}(s, \operatorname{rr\text{-}ms}, i, a) \end{array}}}
\vnbprop{cc-diff-at-iff-coords}{\vnbstmt{\begin{array}{@{}l@{}} \forall g \in (\mathbb{R} \to \mathbb{C}), r, i \in (\mathbb{R} \to \mathbb{R}), v \in \mathbb{R}, a \in \mathbb{C}. \\ \quad (\forall x \in \mathbb{R}.\; (\operatorname{real\text{-}part}(g(x)) = r(x)) \wedge \forall x \in \mathbb{R}.\; (\operatorname{imag\text{-}part}(g(x)) = i(x))) \Rightarrow  \\ \quad \quad \operatorname{is\text{-}cc\text{-}diff\text{-}at}(g, v, a) \Leftrightarrow (\operatorname{is\text{-}diff\text{-}at}(r, v, \operatorname{real\text{-}part}(a)) \wedge \operatorname{is\text{-}diff\text{-}at}(i, v, \operatorname{imag\text{-}part}(a))) \end{array}}}
\vnbprop{cc-diff-at-sum}{\vnbstmt{\begin{array}{@{}l@{}} \forall g, h, m \in (\mathbb{R} \to \mathbb{C}), v \in \mathbb{R}, a, b \in \mathbb{C}. \\ \quad (\forall x \in \mathbb{R}.\; (m(x) = (g(x) + h(x))) \wedge (\operatorname{is\text{-}cc\text{-}diff\text{-}at}(g, v, a) \wedge \operatorname{is\text{-}cc\text{-}diff\text{-}at}(h, v, b))) \Rightarrow  \\ \quad \quad \operatorname{is\text{-}cc\text{-}diff\text{-}at}(m, v, (a + b)) \end{array}}}
\vnbprop{cc-diff-at-real-mul}{\vnbstmt{\begin{array}{@{}l@{}} \forall c \in \mathbb{R}, g, d \in (\mathbb{R} \to \mathbb{C}), v \in \mathbb{R}, a \in \mathbb{C}. \\ \quad (\forall x \in \mathbb{R}.\; (d(x) = (c \cdot g(x))) \wedge \operatorname{is\text{-}cc\text{-}diff\text{-}at}(g, v, a)) \Rightarrow  \\ \quad \quad \operatorname{is\text{-}cc\text{-}diff\text{-}at}(d, v, (c \cdot a)) \end{array}}}
\vnbprop{cc-diff-at-product}{\vnbstmt{\begin{array}{@{}l@{}} \forall g, h, r \in (\mathbb{R} \to \mathbb{C}), v \in \mathbb{R}, a, b \in \mathbb{C}. \\ \quad (\forall x \in \mathbb{R}.\; (r(x) = (g(x) \cdot h(x))) \wedge (\operatorname{is\text{-}cc\text{-}diff\text{-}at}(g, v, a) \wedge \operatorname{is\text{-}cc\text{-}diff\text{-}at}(h, v, b))) \Rightarrow  \\ \quad \quad \operatorname{is\text{-}cc\text{-}diff\text{-}at}(r, v, ((a \cdot h(v)) + (g(v) \cdot b))) \end{array}}}

\vnbbookfile{theorem-library/product2-cc}{theorem-library/product2-cc}
\noindent{\small product2-cc.scm -- CC and the plane PRODUCT2-MS(RR-MS, RR-MS): the gap that theorem-library/cc-coords-laws.scm recorded ("there is no binary product of metric spaces in the tree"), closed.  Helper prefix p2c-.}\par
\vnbprop{product2-rr-rr-pts}{\vnbstmt{(\operatorname{pts}(\operatorname{product2\text{-}ms}(\operatorname{rr\text{-}ms}, \operatorname{rr\text{-}ms})) \simeq (\mathbb{R} \times \mathbb{R}))}}
\vnbprop{cc-coords-dist-le}{\vnbstmt{\begin{array}{@{}l@{}} \forall z, w \in \mathbb{C}. \\ \quad \operatorname{dist}(\operatorname{product2\text{-}ms}(\operatorname{rr\text{-}ms}, \operatorname{rr\text{-}ms}))(\operatorname{cc\text{-}coords}(z), \operatorname{cc\text{-}coords}(w)) \leq \operatorname{dist}(\operatorname{cc\text{-}ms})(z, w) \end{array}}}
\vnbprop{cc-coords-dist-ge}{\vnbstmt{\begin{array}{@{}l@{}} \forall z, w \in \mathbb{C}. \\ \quad \operatorname{dist}(\operatorname{cc\text{-}ms})(z, w) \leq (2 \cdot \operatorname{dist}(\operatorname{product2\text{-}ms}(\operatorname{rr\text{-}ms}, \operatorname{rr\text{-}ms}))(\operatorname{cc\text{-}coords}(z), \operatorname{cc\text{-}coords}(w))) \end{array}}}
\vnbprop{cc-coords-continuous-at}{\vnbstmt{\begin{array}{@{}l@{}} \forall z \in \mathbb{C}. \\ \quad \operatorname{is\text{-}continuous\text{-}at}(\operatorname{cc\text{-}ms}, \operatorname{product2\text{-}ms}(\operatorname{rr\text{-}ms}, \operatorname{rr\text{-}ms}), \operatorname{cc\text{-}coords}, z) \end{array}}}
\vnbprop{cc-of-pair-continuous-at}{\vnbstmt{\begin{array}{@{}l@{}} \forall p \in (\mathbb{R} \times \mathbb{R}). \\ \quad \operatorname{is\text{-}continuous\text{-}at}(\operatorname{product2\text{-}ms}(\operatorname{rr\text{-}ms}, \operatorname{rr\text{-}ms}), \operatorname{cc\text{-}ms}, \operatorname{cc\text{-}of\text{-}pair}, p) \end{array}}}

\vnbbookfile{structure-library/path-integral}{structure-library/path-integral}
\noindent{\small path-integral.scm -- THE VOCABULARY OF THE INTEGRAL ALONG A PATH. DEFINITIONS ONLY; every law is PROVEN in theorem-library/pw-antiderivative-laws.scm, cc-int-laws.scm and road-laws.scm.}\par
\par\medskip\noindent\textbf{Definition} (\texttt{is-pw-continuous-on}).\;For all $\operatorname{pphi\_}$, $a$ and $b$, 
\begin{equation*}
\begin{array}{@{}l@{}} \operatorname{is\text{-}pw\text{-}continuous\text{-}on}(\operatorname{pphi\_}, a, b) \\ \Leftrightarrow\; a \in \mathbb{R} \wedge \\ b \in \mathbb{R} \wedge \\ a < b \wedge \\ \operatorname{pphi\_} \in (\operatorname{ccint}(a, b) \to \mathbb{R}) \wedge \\ \exists \operatorname{pas\_} \in \mathbf{Set}. \\ \quad |\operatorname{pas\_}| \in \mathbb{N} \wedge \\ \quad \operatorname{pas\_} \subseteq \operatorname{ccint}(a, b) \wedge \\ \quad \forall \operatorname{pat\_} \in \operatorname{ooint}(a, b). \\ \quad \quad \neg (\operatorname{pat\_} \in \operatorname{pas\_}) \Rightarrow  \\ \quad \quad \quad \operatorname{is\text{-}continuous\text{-}at}(\operatorname{subspace\text{-}ms}(\operatorname{rr\text{-}ms}, \operatorname{ccint}(a, b)), \operatorname{rr\text{-}ms}, \operatorname{pphi\_}, \operatorname{pat\_}) \end{array}
\end{equation*}

\par\medskip\noindent\textbf{Definition} (\texttt{pw-int}).\;$\int_{a}^{b} \operatorname{pphi\_} :=$
\begin{equation*}
(\iota\, \operatorname{pav\_}.\;((\operatorname{pav\_} \in \mathbb{R}) \wedge \exists \operatorname{paw\_}.\; (\operatorname{is\text{-}primitive}(\operatorname{paw\_}, \operatorname{pphi\_}, a, b) \wedge (\operatorname{pav\_} = (\operatorname{paw\_}(b) - \operatorname{paw\_}(a))))))
\end{equation*}

\par\medskip\noindent\textbf{Definition} (\texttt{cc-int}).\;$\int_{a}^{b} \operatorname{pphi\_} :=$
\begin{equation*}
\int_{a}^{b} (\lambda \operatorname{pat\_} \in \operatorname{ccint}(a, b).\; \operatorname{real\text{-}part}(\operatorname{pphi\_}(\operatorname{pat\_}))) + (\int_{a}^{b} (\lambda \operatorname{pat\_} \in \operatorname{ccint}(a, b).\; \operatorname{imag\text{-}part}(\operatorname{pphi\_}(\operatorname{pat\_}))) \cdot +i)
\end{equation*}

\par\medskip\noindent\textbf{Definition} (\texttt{is-path}).\;For all $\mathit{pgam}$, $a$ and $b$, 
\begin{equation*}
\begin{array}{@{}l@{}} \operatorname{is\text{-}path}(\mathit{pgam}, a, b) \\ \Leftrightarrow\; a \in \mathbb{R} \wedge \\ b \in \mathbb{R} \wedge \\ a < b \wedge \\ \mathit{pgam} \in (\operatorname{ccint}(a, b) \to \mathbb{C}) \wedge \\ \forall \operatorname{pax\_} \in \operatorname{ccint}(a, b). \\ \quad \operatorname{is\text{-}continuous\text{-}at}(\operatorname{subspace\text{-}ms}(\operatorname{rr\text{-}ms}, \operatorname{ccint}(a, b)), \operatorname{cc\text{-}ms}, \mathit{pgam}, \operatorname{pax\_}) \end{array}
\end{equation*}

\par\medskip\noindent\textbf{Definition} (\texttt{is-road}).\;For all $\mathit{pgam}$, $\mathit{dgam}$, $a$ and $b$, 
\begin{equation*}
\begin{array}{@{}l@{}} \operatorname{is\text{-}road}(\mathit{pgam}, \mathit{dgam}, a, b) \\ \Leftrightarrow\; \operatorname{is\text{-}path}(\mathit{pgam}, a, b) \wedge \\ \mathit{dgam} \in (\operatorname{ccint}(a, b) \to \mathbb{C}) \wedge \\ \operatorname{is\text{-}regulated\text{-}on}((\lambda \operatorname{pat\_} \in \operatorname{ccint}(a, b).\; \operatorname{real\text{-}part}(\operatorname{dgam}(\operatorname{pat\_}))), a, b) \wedge \\ \operatorname{is\text{-}regulated\text{-}on}((\lambda \operatorname{pat\_} \in \operatorname{ccint}(a, b).\; \operatorname{imag\text{-}part}(\operatorname{dgam}(\operatorname{pat\_}))), a, b) \wedge \\ \operatorname{is\text{-}primitive}((\lambda \operatorname{pat\_} \in \operatorname{ccint}(a, b).\; \operatorname{real\text{-}part}(\operatorname{pgam}(\operatorname{pat\_}))), \\ \quad (\lambda \operatorname{pat\_} \in \operatorname{ccint}(a, b).\; \operatorname{real\text{-}part}(\operatorname{dgam}(\operatorname{pat\_}))), \\ \quad a, b) \wedge \\ \operatorname{is\text{-}primitive}((\lambda \operatorname{pat\_} \in \operatorname{ccint}(a, b).\; \operatorname{imag\text{-}part}(\operatorname{pgam}(\operatorname{pat\_}))), \\ \quad (\lambda \operatorname{pat\_} \in \operatorname{ccint}(a, b).\; \operatorname{imag\text{-}part}(\operatorname{dgam}(\operatorname{pat\_}))), \\ \quad a, b) \end{array}
\end{equation*}

\par\medskip\noindent\textbf{Definition} (\texttt{trace}).\;$\operatorname{trace}(\mathit{pgam}, a, b) :=$ $\operatorname{image}(\mathit{pgam}, \operatorname{ccint}(a, b))$
\par\medskip\noindent\textbf{Definition} (\texttt{line-int}).\;$\operatorname{line\text{-}int}(\mathit{pf}, \mathit{pgam}, \mathit{dgam}, a, b) :=$
\begin{equation*}
\int_{a}^{b} (\lambda \operatorname{pat\_} \in \operatorname{ccint}(a, b).\; (\operatorname{pf}(\operatorname{pgam}(\operatorname{pat\_})) \cdot \operatorname{dgam}(\operatorname{pat\_})))
\end{equation*}

\par\medskip\noindent\textbf{Definition} (\texttt{is-pw-antiderivable}).\;For all $\operatorname{pphi\_}$, $a$ and $b$, 
\begin{equation*}
\operatorname{is\text{-}pw\text{-}antiderivable}(\operatorname{pphi\_}, a, b) \Leftrightarrow \exists \operatorname{paw\_}.\; \operatorname{is\text{-}primitive}(\operatorname{paw\_}, \operatorname{pphi\_}, a, b)
\end{equation*}


\vnbbookfile{structure-library/segments-triangles}{structure-library/segments-triangles}
\noindent{\small segments-triangles.scm -- THE VOCABULARY OF CAUCHY-GOURSAT (complex-analysis.pdf ch. 3, s. 3.2).  DEFINITIONS ONLY; every law is PROVEN in theorem-library/goursat.scm.  Batch 30, 2026-09-25.}\par
\par\medskip\noindent\textbf{Definition} (\texttt{seg-path}).\;$\operatorname{seg\text{-}path}(\operatorname{sgp\_}, \operatorname{sgq\_}) :=$ $(\lambda \operatorname{sgs\_} \in \operatorname{ccint}(0, 1).\; (\operatorname{sgp\_} + (\operatorname{sgs\_} \cdot (\operatorname{sgq\_} - \operatorname{sgp\_}))))$
\par\medskip\noindent\textbf{Definition} (\texttt{seg-deriv}).\;$\operatorname{seg\text{-}deriv}(\operatorname{sgp\_}, \operatorname{sgq\_}) :=$ $(\lambda \operatorname{sgs\_} \in \operatorname{ccint}(0, 1).\; (\operatorname{sgq\_} - \operatorname{sgp\_}))$
\par\medskip\noindent\textbf{Definition} (\texttt{seg-int}).\;$\operatorname{seg\text{-}int}(\operatorname{sgf\_}, \operatorname{sgp\_}, \operatorname{sgq\_}) :=$
\begin{equation*}
\operatorname{line\text{-}int}(\operatorname{sgf\_}, \operatorname{seg\text{-}path}(\operatorname{sgp\_}, \operatorname{sgq\_}), \operatorname{seg\text{-}deriv}(\operatorname{sgp\_}, \operatorname{sgq\_}), 0, 1)
\end{equation*}

\par\medskip\noindent\textbf{Definition} (\texttt{tri-int}).\;$\operatorname{tri\text{-}int}(\operatorname{sgf\_}, \operatorname{sga\_}, \operatorname{sgb\_}, \operatorname{sgc\_}) :=$
\begin{equation*}
(\operatorname{seg\text{-}int}(\operatorname{sgf\_}, \operatorname{sga\_}, \operatorname{sgb\_}) + \operatorname{seg\text{-}int}(\operatorname{sgf\_}, \operatorname{sgb\_}, \operatorname{sgc\_})) + \operatorname{seg\text{-}int}(\operatorname{sgf\_}, \operatorname{sgc\_}, \operatorname{sga\_})
\end{equation*}

\par\medskip\noindent\textbf{Definition} (\texttt{conv3}).\;$\operatorname{conv3}(\operatorname{cva\_}, \operatorname{cvb\_}, \operatorname{cvc\_}) :=$
\begin{equation*}
\{\operatorname{cvz\_} \in \mathbb{C} : \exists \operatorname{cvl\_}, \operatorname{cvm\_}, \operatorname{cvn\_} \in \operatorname{ccint}(0, 1).\; ((((\operatorname{cvl\_} + \operatorname{cvm\_}) + \operatorname{cvn\_}) = 1) \wedge (\operatorname{cvz\_} = (((\operatorname{cvl\_} \cdot \operatorname{cva\_}) + (\operatorname{cvm\_} \cdot \operatorname{cvb\_})) + (\operatorname{cvn\_} \cdot \operatorname{cvc\_}))))\}
\end{equation*}

\par\medskip\noindent\textbf{Definition} (\texttt{cc-mid}).\;$\operatorname{cc\text{-}mid}(\operatorname{cmx\_}, \operatorname{cmy\_}) :=$ $(\operatorname{cmx\_} + \operatorname{cmy\_}) \cdot \operatorname{recip}(2)$
\par\medskip\noindent\textbf{Definition} (\texttt{is-convex-cc}).\;For every $\operatorname{cvu\_}$, 
\begin{equation*}
\begin{array}{@{}l@{}} \operatorname{is\text{-}convex\text{-}cc}(\operatorname{cvu\_}) \\ \Leftrightarrow\; \operatorname{cvu\_} \subseteq \mathbb{C} \wedge \\ \forall \operatorname{cvx\_}, \operatorname{cvy\_} \in \operatorname{cvu\_}, \operatorname{cvt\_} \in \operatorname{ccint}(0, 1). \\ \quad +(\operatorname{cvx\_}, \cdot(\operatorname{cvt\_}, -(\operatorname{cvy\_}, \operatorname{cvx\_}))) \in \operatorname{cvu\_} \end{array}
\end{equation*}


\vnbbookfile{structure-library/road-juxtaposition}{structure-library/road-juxtaposition}
\noindent{\small road-juxtaposition.scm -- THE JUXTAPOSITION OF ROADS, THE TRIANGULAR ROAD AND THE LENGTH OF A ROAD.  DEFINITIONS ONLY; every law is PROVEN in theorem-library/road-juxtaposition-laws.scm.  Batch 37, 2026-09-26.}\par
\par\medskip\noindent\textbf{Definition} (\texttt{juxta}).\;$\operatorname{juxta}(\operatorname{jxp\_}, \operatorname{jxq\_}, \operatorname{jxa\_}, \operatorname{jxb\_}, \operatorname{jxc\_}) :=$
\begin{equation*}
(\lambda \operatorname{jxt\_} \in \operatorname{ccint}(\operatorname{jxa\_}, \operatorname{jxc\_}).\; \left(\text{if }(\operatorname{jxt\_} \leq \operatorname{jxb\_})\text{ then }\operatorname{jxp\_}(\operatorname{jxt\_})\text{ else }\operatorname{jxq\_}(\operatorname{jxt\_})\right))
\end{equation*}

\par\medskip\noindent\textbf{Definition} (\texttt{road-length}).\;$\operatorname{road\text{-}length}(\operatorname{rld\_}, \operatorname{rla\_}, \operatorname{rlb\_}) :=$
\begin{equation*}
\int_{\operatorname{rla\_}}^{\operatorname{rlb\_}} (\lambda \operatorname{rlt\_} \in \operatorname{ccint}(\operatorname{rla\_}, \operatorname{rlb\_}).\; \operatorname{magnitude}(\operatorname{rld\_}(\operatorname{rlt\_})))
\end{equation*}

\par\medskip\noindent\textbf{Definition} (\texttt{tri-road}).\;$\operatorname{tri\text{-}road}(\operatorname{tra\_}, \operatorname{trb\_}, \operatorname{trc\_}) :=$
\begin{equation*}
\begin{array}{@{}l@{}} \operatorname{juxta}(\operatorname{juxta}(\operatorname{seg\text{-}path}(\operatorname{tra\_}, \operatorname{trb\_}), \\ \quad \quad (\lambda \operatorname{trs\_} \in \operatorname{ccint}(1, 2).\; \operatorname{seg\text{-}path}(\operatorname{trb\_}, \operatorname{trc\_})((\operatorname{trs\_} - 1))), \\ \quad \quad 0, 1, 2), \\ \quad (\lambda \operatorname{trs\_} \in \operatorname{ccint}(2, 3).\; \operatorname{seg\text{-}path}(\operatorname{trc\_}, \operatorname{tra\_})((\operatorname{trs\_} - 2))), \\ \quad 0, 2, 3) \end{array}
\end{equation*}

\par\medskip\noindent\textbf{Definition} (\texttt{tri-road-deriv}).\;$\operatorname{tri\text{-}road\text{-}deriv}(\operatorname{tra\_}, \operatorname{trb\_}, \operatorname{trc\_}) :=$
\begin{equation*}
\begin{array}{@{}l@{}} \operatorname{juxta}(\operatorname{juxta}(\operatorname{seg\text{-}deriv}(\operatorname{tra\_}, \operatorname{trb\_}), \\ \quad \quad (\lambda \operatorname{trs\_} \in \operatorname{ccint}(1, 2).\; \operatorname{seg\text{-}deriv}(\operatorname{trb\_}, \operatorname{trc\_})((\operatorname{trs\_} - 1))), \\ \quad \quad 0, 1, 2), \\ \quad (\lambda \operatorname{trs\_} \in \operatorname{ccint}(2, 3).\; \operatorname{seg\text{-}deriv}(\operatorname{trc\_}, \operatorname{tra\_})((\operatorname{trs\_} - 2))), \\ \quad 0, 2, 3) \end{array}
\end{equation*}


\vnbbookfile{theorem-library/regulated-primitive-laws}{theorem-library/regulated-primitive-laws}
\noindent{\small regulated-primitive-laws.scm -- the laws of IS-COUNTABLE, IS-RIGHT-LIMIT-WITHIN / IS-LEFT-LIMIT-WITHIN, IS-REGULATED-ON and IS-PRIMITIVE (structure-library/regulated-primitive.scm).  Batch 20-D, 2026-09-23.  THE SPECIFICATION is docs/roads-regulated-primitives-2026-09-23.md.}\par
\vnbprop{countable-of-enum}{\vnbstmt{\begin{array}{@{}l@{}} \forall d, e. \\ \quad ((d \in \mathbf{Set}) \wedge ((e \in (\mathbb{N} \to d)) \wedge \forall y \in d.\; \exists n \in \mathbb{N}.\; (y = e(n)))) \Rightarrow  \\ \quad \quad \operatorname{is\text{-}countable}(d) \end{array}}}
\vnbprop{countable-empty}{\vnbstmt{\operatorname{is\text{-}countable}(\emptyset)}}
\vnbprop{rr-abs-lt-parts}{\vnbstmt{\forall x, c \in \mathbb{R}.\; ((\lvert x \rvert < c) \Rightarrow ((-c < x) \wedge (x < c)))}}
\vnbprop{countable-singleton}{\vnbstmt{\forall x \in \mathbf{Set}.\; \operatorname{is\text{-}countable}(\{x\})}}
\vnbprop{countable-subset}{\vnbstmt{\forall a, d.\; ((\operatorname{is\text{-}countable}(d) \wedge (a \subseteq d)) \Rightarrow \operatorname{is\text{-}countable}(a))}}
\vnbprop{countable-cases}{\vnbstmt{\begin{array}{@{}l@{}} \forall d. \\ \quad \operatorname{is\text{-}countable}(d) \Rightarrow  \\ \quad \quad d \in \mathbf{Set} \wedge \\ \quad \quad ((d = \emptyset) \vee \exists e \in (\mathbb{N} \to d).\; \forall y \in d.\; \exists n \in \mathbb{N}.\; (y = e(n))) \end{array}}}
\vnbprop{countable-of-two-enum}{\vnbstmt{\begin{array}{@{}l@{}} \forall u, f, g. \\ \quad u \in \mathbf{Set} \wedge \\ \quad f \in (\mathbb{N} \to u) \wedge \\ \quad g \in (\mathbb{N} \to u) \wedge \\ \quad \forall y \in u.\; (\exists j \in \mathbb{N}.\; (y = f(j)) \vee \exists j \in \mathbb{N}.\; (y = g(j))) \Rightarrow \\ \quad \quad \operatorname{is\text{-}countable}(u) \end{array}}}
\vnbprop{countable-union-2}{\vnbstmt{\begin{array}{@{}l@{}} \forall a, b. \\ \quad (\operatorname{is\text{-}countable}(a) \wedge \operatorname{is\text{-}countable}(b)) \Rightarrow \operatorname{is\text{-}countable}((a \cup b)) \end{array}}}
\vnbprop{finite-implies-countable}{\vnbstmt{\forall d \in \mathbf{Set}.\; ((|d| \in \mathbb{N}) \Rightarrow \operatorname{is\text{-}countable}(d))}}
\vnbprop{countable-image}{\vnbstmt{\begin{array}{@{}l@{}} \forall d, h, a, b. \\ \quad (\operatorname{is\text{-}countable}(d) \wedge ((h \in (a \to b)) \wedge (d \subseteq a))) \Rightarrow  \\ \quad \quad \operatorname{is\text{-}countable}(\operatorname{image}(h, d)) \end{array}}}
\vnbprop{rr-abs-le-parts}{\vnbstmt{\forall x, c \in \mathbb{R}.\; ((\lvert x \rvert \leq c) \Rightarrow ((-c \leq x) \wedge (x \leq c)))}}
\vnbprop{rr-abs-le-of-parts}{\vnbstmt{\forall x, c \in \mathbb{R}.\; (((-c \leq x) \wedge (x \leq c)) \Rightarrow (\lvert x \rvert \leq c))}}
\vnbprop{rr-abs-lt-of-parts}{\vnbstmt{\forall x, c \in \mathbb{R}.\; (((-c < x) \wedge (x < c)) \Rightarrow (\lvert x \rvert < c))}}
\vnbprop{continuous-on-right-limit-within}{\vnbstmt{\begin{array}{@{}l@{}} \forall f, w.; \forall x \in w. \\ \quad ((w \subseteq \mathbb{R}) \wedge \operatorname{is\text{-}continuous\text{-}on}(f, w)) \Rightarrow  \\ \quad \quad \operatorname{is\text{-}right\text{-}limit\text{-}within}(f, w, x, f(x)) \end{array}}}
\vnbprop{continuous-on-left-limit-within}{\vnbstmt{\begin{array}{@{}l@{}} \forall f, w.; \forall x \in w. \\ \quad ((w \subseteq \mathbb{R}) \wedge \operatorname{is\text{-}continuous\text{-}on}(f, w)) \Rightarrow  \\ \quad \quad \operatorname{is\text{-}left\text{-}limit\text{-}within}(f, w, x, f(x)) \end{array}}}
\vnbprop{continuous-on-implies-regulated-on}{\vnbstmt{\begin{array}{@{}l@{}} \forall f, a, b. \\ \quad ((a \in \mathbb{R}) \wedge ((b \in \mathbb{R}) \wedge ((a < b) \wedge \operatorname{is\text{-}continuous\text{-}on}(f, \operatorname{ccint}(a, b))))) \Rightarrow  \\ \quad \quad \operatorname{is\text{-}regulated\text{-}on}(f, a, b) \end{array}}}
\vnbprop{right-limit-within-le}{\vnbstmt{\begin{array}{@{}l@{}} \forall f, w, x, l, m. \\ \quad \operatorname{is\text{-}right\text{-}limit\text{-}within}(f, w, x, l) \wedge \\ \quad \operatorname{is\text{-}right\text{-}limit\text{-}within}(f, w, x, m) \wedge \\ \quad \forall r.\; (\operatorname{pos\text{-}rr}(r) \Rightarrow \exists t \in w.\; ((x < t) \wedge (t < (x + r)))) \Rightarrow \\ \quad \quad (l - m) \leq 0 \end{array}}}
\vnbprop{left-limit-within-le}{\vnbstmt{\begin{array}{@{}l@{}} \forall f, w, x, l, m. \\ \quad \operatorname{is\text{-}left\text{-}limit\text{-}within}(f, w, x, l) \wedge \\ \quad \operatorname{is\text{-}left\text{-}limit\text{-}within}(f, w, x, m) \wedge \\ \quad \forall r.\; (\operatorname{pos\text{-}rr}(r) \Rightarrow \exists t \in w.\; (((x - r) < t) \wedge (t < x))) \Rightarrow \\ \quad \quad (l - m) \leq 0 \end{array}}}
\vnbprop{right-limit-within-unique}{\vnbstmt{\begin{array}{@{}l@{}} \forall f, w, x, l, m. \\ \quad \operatorname{is\text{-}right\text{-}limit\text{-}within}(f, w, x, l) \wedge \\ \quad \operatorname{is\text{-}right\text{-}limit\text{-}within}(f, w, x, m) \wedge \\ \quad \forall r.\; (\operatorname{pos\text{-}rr}(r) \Rightarrow \exists t \in w.\; ((x < t) \wedge (t < (x + r)))) \Rightarrow \\ \quad \quad l = m \end{array}}}
\vnbprop{left-limit-within-unique}{\vnbstmt{\begin{array}{@{}l@{}} \forall f, w, x, l, m. \\ \quad \operatorname{is\text{-}left\text{-}limit\text{-}within}(f, w, x, l) \wedge \\ \quad \operatorname{is\text{-}left\text{-}limit\text{-}within}(f, w, x, m) \wedge \\ \quad \forall r.\; (\operatorname{pos\text{-}rr}(r) \Rightarrow \exists t \in w.\; (((x - r) < t) \wedge (t < x))) \Rightarrow \\ \quad \quad l = m \end{array}}}
\vnbprop{ccint-right-approach}{\vnbstmt{\begin{array}{@{}l@{}} \forall a, b, x \in \mathbb{R}, r. \\ \quad ((a \leq x) \wedge ((x < b) \wedge \operatorname{pos\text{-}rr}(r))) \Rightarrow  \\ \quad \quad \exists t \in \operatorname{ccint}(a, b).\; ((x < t) \wedge (t < (x + r))) \end{array}}}
\vnbprop{ccint-left-approach}{\vnbstmt{\begin{array}{@{}l@{}} \forall a, b, x \in \mathbb{R}, r. \\ \quad ((a < x) \wedge ((x \leq b) \wedge \operatorname{pos\text{-}rr}(r))) \Rightarrow  \\ \quad \quad \exists t \in \operatorname{ccint}(a, b).\; (((x - r) < t) \wedge (t < x)) \end{array}}}
\vnbprop{right-limit-within-ccint-unique}{\vnbstmt{\begin{array}{@{}l@{}} \forall f.; \forall a, b, x \in \mathbb{R}, l, m. \\ \quad a \leq x \wedge \\ \quad x < b \wedge \\ \quad \operatorname{is\text{-}right\text{-}limit\text{-}within}(f, \operatorname{ccint}(a, b), x, l) \wedge \\ \quad \operatorname{is\text{-}right\text{-}limit\text{-}within}(f, \operatorname{ccint}(a, b), x, m) \Rightarrow \\ \quad \quad l = m \end{array}}}
\vnbprop{left-limit-within-ccint-unique}{\vnbstmt{\begin{array}{@{}l@{}} \forall f.; \forall a, b, x \in \mathbb{R}, l, m. \\ \quad a < x \wedge \\ \quad x \leq b \wedge \\ \quad \operatorname{is\text{-}left\text{-}limit\text{-}within}(f, \operatorname{ccint}(a, b), x, l) \wedge \\ \quad \operatorname{is\text{-}left\text{-}limit\text{-}within}(f, \operatorname{ccint}(a, b), x, m) \Rightarrow \\ \quad \quad l = m \end{array}}}
\vnbprop{right-limit-within-subset}{\vnbstmt{\begin{array}{@{}l@{}} \forall f, w, x, l, g, v. \\ \quad \operatorname{is\text{-}right\text{-}limit\text{-}within}(f, w, x, l) \wedge \\ \quad v \subseteq w \wedge \\ \quad g \in (v \to \mathbb{R}) \wedge \\ \quad \forall z \in v.\; (g(z) = f(z)) \Rightarrow \\ \quad \quad \operatorname{is\text{-}right\text{-}limit\text{-}within}(g, v, x, l) \end{array}}}
\vnbprop{left-limit-within-subset}{\vnbstmt{\begin{array}{@{}l@{}} \forall f, w, x, l, g, v. \\ \quad \operatorname{is\text{-}left\text{-}limit\text{-}within}(f, w, x, l) \wedge \\ \quad v \subseteq w \wedge \\ \quad g \in (v \to \mathbb{R}) \wedge \\ \quad \forall z \in v.\; (g(z) = f(z)) \Rightarrow \\ \quad \quad \operatorname{is\text{-}left\text{-}limit\text{-}within}(g, v, x, l) \end{array}}}
\vnbprop{regulated-on-restrict}{\vnbstmt{\begin{array}{@{}l@{}} \forall f, a, b, c, d, g. \\ \quad \operatorname{is\text{-}regulated\text{-}on}(f, a, b) \wedge \\ \quad c \in \mathbb{R} \wedge \\ \quad d \in \mathbb{R} \wedge \\ \quad a \leq c \wedge \\ \quad c < d \wedge \\ \quad d \leq b \wedge \\ \quad g \in (\operatorname{ccint}(c, d) \to \mathbb{R}) \wedge \\ \quad \forall z \in \operatorname{ccint}(c, d).\; (g(z) = f(z)) \Rightarrow \\ \quad \quad \operatorname{is\text{-}regulated\text{-}on}(g, c, d) \end{array}}}
\vnbprop{primitive-endpoints}{\vnbstmt{\begin{array}{@{}l@{}} \forall f, i, a, b. \\ \quad \operatorname{is\text{-}primitive}(f, i, a, b) \Rightarrow ((a \in \mathbb{R}) \wedge ((b \in \mathbb{R}) \wedge (a < b))) \end{array}}}
\vnbprop{primitive-in-fun}{\vnbstmt{\forall f, i, a, b.\; (\operatorname{is\text{-}primitive}(f, i, a, b) \Rightarrow (f \in (\operatorname{ccint}(a, b) \to \mathbb{R})))}}
\vnbprop{primitive-integrand-in-fun}{\vnbstmt{\forall f, i, a, b.\; (\operatorname{is\text{-}primitive}(f, i, a, b) \Rightarrow (i \in (\operatorname{ccint}(a, b) \to \mathbb{R})))}}
\vnbprop{primitive-continuous}{\vnbstmt{\begin{array}{@{}l@{}} \forall f, i, a, b. \\ \quad \operatorname{is\text{-}primitive}(f, i, a, b) \Rightarrow \operatorname{is\text{-}continuous\text{-}on}(f, \operatorname{ccint}(a, b)) \end{array}}}
\vnbprop{primitive-exceptional-set}{\vnbstmt{\begin{array}{@{}l@{}} \forall f, i, a, b. \\ \quad \operatorname{is\text{-}primitive}(f, i, a, b) \Rightarrow  \\ \quad \quad \exists s. \\ \quad \quad \quad \operatorname{is\text{-}countable}(s) \wedge \\ \quad \quad \quad s \subseteq \operatorname{ccint}(a, b) \wedge \\ \quad \quad \quad \forall t \in \operatorname{ooint}(a, b).\; (\neg (t \in s) \Rightarrow \operatorname{has\text{-}deriv\text{-}at}(f, t, i(t))) \end{array}}}

\vnbbookfile{theorem-library/countable-products}{theorem-library/countable-products}
\noindent{\small countable-products.scm -- wave of 2026-11-09/10 (agent BA-ONB3), wired 2026-11-10 from scratchpad/baonb3-1109/b3-count.scm. Library: analysis.}\par
\vnbprop{countable-enum-nonempty}{\vnbstmt{\begin{array}{@{}l@{}} \forall t. \\ \quad (\operatorname{is\text{-}countable}(t) \wedge \exists a.\; (a \in t)) \Rightarrow  \\ \quad \quad \exists e \in (\mathbb{N} \to t).\; \forall y \in t.\; \exists n \in \mathbb{N}.\; (y = e(n)) \end{array}}}
\vnbprop{countable-cartesian}{\vnbstmt{\begin{array}{@{}l@{}} \forall a, b. \\ \quad (\operatorname{is\text{-}countable}(a) \wedge \operatorname{is\text{-}countable}(b)) \Rightarrow \operatorname{is\text{-}countable}((a \times b)) \end{array}}}

\vnbbookfile{theorem-library/rabasics-extline}{theorem-library/rabasics-extline}
\noindent{\small rabasics-extline.scm -- calculus.pdf 1.13: x |-$>$ x / (1 - |x|) is an order preserving bijection of (-1, 1) onto R, with inverse s |-$>$ s / (1 + |s|), and it is odd.}\par
\vnbprop{rr-quot-abs}{\vnbstmt{\forall a, c \in \mathbb{R}.\; ((0 < c) \Rightarrow (\lvert (a \cdot \operatorname{recip}(c)) \rvert = (\lvert a \rvert \cdot \operatorname{recip}(c))))}}
\vnbprop{rr-pos-of-mul-eq-one}{\vnbstmt{\forall d, a \in \mathbb{R}.\; (((0 < a) \wedge ((d \cdot a) = 1)) \Rightarrow (0 < d))}}
\vnbprop{ooint-rr-iso-g-den}{\vnbstmt{\forall s \in \mathbb{R}.\; (((1 - \lvert (s \cdot \operatorname{recip}((1 + \lvert s \rvert))) \rvert) \cdot (1 + \lvert s \rvert)) = 1)}}
\vnbprop{ooint-rr-iso-f-den}{\vnbstmt{\begin{array}{@{}l@{}} \forall x \in \mathbb{R}. \\ \quad (\lvert x \rvert < 1) \Rightarrow (((1 + \lvert (x \cdot \operatorname{recip}((1 - \lvert x \rvert))) \rvert) \cdot (1 - \lvert x \rvert)) = 1) \end{array}}}
\vnbprop{ooint-rr-iso-g-small}{\vnbstmt{\forall s \in \mathbb{R}.\; (\lvert (s \cdot \operatorname{recip}((1 + \lvert s \rvert))) \rvert < 1)}}
\vnbprop{ooint-rr-iso-f-real}{\vnbstmt{\forall x \in \mathbb{R}.\; ((\lvert x \rvert < 1) \Rightarrow ((x \cdot \operatorname{recip}((1 - \lvert x \rvert))) \in \mathbb{R}))}}
\vnbprop{ooint-rr-iso-g-inv}{\vnbstmt{\begin{array}{@{}l@{}} \forall s \in \mathbb{R}. \\ \quad ((s \cdot \operatorname{recip}((1 + \lvert s \rvert))) \cdot \operatorname{recip}((1 - \lvert (s \cdot \operatorname{recip}((1 + \lvert s \rvert))) \rvert))) = s \end{array}}}
\vnbprop{ooint-rr-iso-f-inv}{\vnbstmt{\begin{array}{@{}l@{}} \forall x \in \mathbb{R}. \\ \quad (\lvert x \rvert < 1) \Rightarrow  \\ \quad \quad ((x \cdot \operatorname{recip}((1 - \lvert x \rvert))) \cdot \operatorname{recip}((1 + \lvert (x \cdot \operatorname{recip}((1 - \lvert x \rvert))) \rvert))) = x \end{array}}}
\vnbprop{rr-recip-form-nonneg}{\vnbstmt{\begin{array}{@{}l@{}} \forall s, a, r \in \mathbb{R}. \\ \quad (((a \cdot r) = 1) \wedge (a = (1 + s))) \Rightarrow ((s \cdot r) = (1 - r)) \end{array}}}
\vnbprop{rr-recip-form-neg}{\vnbstmt{\begin{array}{@{}l@{}} \forall s, a, r \in \mathbb{R}. \\ \quad (((a \cdot r) = 1) \wedge (a = (1 + -s))) \Rightarrow ((s \cdot r) = (r - 1)) \end{array}}}
\vnbprop{ooint-rr-iso-g-mono}{\vnbstmt{\begin{array}{@{}l@{}} \forall s, t \in \mathbb{R}. \\ \quad (s \leq t) \Rightarrow ((s \cdot \operatorname{recip}((1 + \lvert s \rvert))) \leq (t \cdot \operatorname{recip}((1 + \lvert t \rvert)))) \end{array}}}
\vnbprop{ooint-rr-iso-strict}{\vnbstmt{\begin{array}{@{}l@{}} \forall x, y \in \operatorname{ooint}(-1, 1). \\ \quad (x < y) \Rightarrow (/(x, (1 - \lvert x \rvert)) < /(y, (1 - \lvert y \rvert))) \end{array}}}
\vnbprop{ooint-rr-iso-odd}{\vnbstmt{\forall x \in \operatorname{ooint}(-1, 1).\; (/(-x, (1 - \lvert -x \rvert)) = -/(x, (1 - \lvert x \rvert)))}}
\vnbprop{ooint-rr-iso-left-inverse}{\vnbstmt{\begin{array}{@{}l@{}} \forall x \in \operatorname{ooint}(-1, 1). \\ \quad /(/(x, (1 - \lvert x \rvert)), (1 + \lvert /(x, (1 - \lvert x \rvert)) \rvert)) = x \end{array}}}
\vnbprop{ooint-rr-iso-right-inverse}{\vnbstmt{\begin{array}{@{}l@{}} \forall s \in \mathbb{R}. \\ \quad ((/(s, (1 + \lvert s \rvert)) \in \operatorname{ooint}(-1, 1)) \wedge (/(/(s, (1 + \lvert s \rvert)), (1 - \lvert /(s, (1 + \lvert s \rvert)) \rvert)) = s)) \end{array}}}
\vnbprop{ooint-rr-iso-bijection}{\vnbstmt{(\lambda x \in \operatorname{ooint}(-1, 1).\; /(x, (1 - \lvert x \rvert))) \in \mathrm{Bij}(\operatorname{ooint}(-1, 1), \mathbb{R})}}
\vnbprop{ooint-rr-order-iso}{\vnbstmt{\operatorname{is\text{-}order\text{-}iso}((\lambda x \in \operatorname{ooint}(-1, 1).\; /(x, (1 - \lvert x \rvert))), \operatorname{ooint}(-1, 1), \mathbb{R})}}

\vnbbookfile{theorem-library/rabasics-extline2}{theorem-library/rabasics-extline2}
\noindent{\small rabasics-extline2.scm -- calculus.pdf 1.13, last sentence: x |-$>$ x / (1 - |x|) extends to an order preserving bijection of [-1, 1] onto the extended line RR-STAR, sending 1 to POS-INF and -1 to NEG-INF.}\par
\vnbprop{rr-star-is-set}{\vnbstmt{\operatorname{rr\text{-}star} \in \mathbf{Set}}}
\vnbprop{ccint-unit-cases}{\vnbstmt{\forall x \in \operatorname{ccint}(-1, 1).\; ((x = 1) \vee ((x = -1) \vee (x \in \operatorname{ooint}(-1, 1))))}}
\vnbprop{ccint-rr-star-iso-extends}{\vnbstmt{\begin{array}{@{}l@{}} \forall u \in \operatorname{ooint}(-1, 1). \\ \quad (\lambda x \in \operatorname{ccint}(-1, 1).\; \left(\text{if }(x = 1)\text{ then }\operatorname{pos\text{-}inf}\text{ else }\left(\text{if }(x = -1)\text{ then }\operatorname{neg\text{-}inf}\text{ else }/(x, (1 - \lvert x \rvert))\right)\right))(u) \\ \quad =\; /(u, (1 - \lvert u \rvert)) \end{array}}}
\vnbprop{ccint-rr-star-iso-bijection}{\vnbstmt{\begin{array}{@{}l@{}} (\lambda x \in \operatorname{ccint}(-1, 1).\; \left(\text{if }(x = 1)\text{ then }\operatorname{pos\text{-}inf}\text{ else }\left(\text{if }(x = -1)\text{ then }\operatorname{neg\text{-}inf}\text{ else }/(x, (1 - \lvert x \rvert))\right)\right)) \\ \in\; \mathrm{Bij}(\operatorname{ccint}(-1, 1), \operatorname{rr\text{-}star}) \end{array}}}
\vnbprop{ccint-rr-star-iso-strict}{\vnbstmt{\begin{array}{@{}l@{}} \forall u, v \in \operatorname{ccint}(-1, 1). \\ \quad (u < v) \Rightarrow  \\ \quad \quad (\lambda x \in \operatorname{ccint}(-1, 1).\; \left(\text{if }(x = 1)\text{ then }\operatorname{pos\text{-}inf}\text{ else }\left(\text{if }(x = -1)\text{ then }\operatorname{neg\text{-}inf}\text{ else }/(x, (1 - \lvert x \rvert))\right)\right))(u) \\ \quad \quad <\; (\lambda x \in \operatorname{ccint}(-1, 1).\; \left(\text{if }(x = 1)\text{ then }\operatorname{pos\text{-}inf}\text{ else }\left(\text{if }(x = -1)\text{ then }\operatorname{neg\text{-}inf}\text{ else }/(x, (1 - \lvert x \rvert))\right)\right))(v) \end{array}}}
\vnbprop{ccint-rr-star-order-iso}{\vnbstmt{\begin{array}{@{}l@{}} \operatorname{is\text{-}order\text{-}iso}((\lambda x \in \operatorname{ccint}(-1, 1).\; \left(\text{if }(x = 1)\text{ then }\operatorname{pos\text{-}inf}\text{ else }\left(\text{if }(x = -1)\text{ then }\operatorname{neg\text{-}inf}\text{ else }/(x, (1 - \lvert x \rvert))\right)\right)), \\ \quad \operatorname{ccint}(-1, 1), \operatorname{rr\text{-}star}) \end{array}}}

\vnbbookfile{theorem-library/extinf-iso}{theorem-library/extinf-iso}
\noindent{\small extinf-iso.scm -- calculus.pdf 1.13 with $<$= REFLECTED: what neg-inf-below-reals buys, and the one comparison it does not decide.  Agent EXT-INF, 2026-10-26. DESTINATION: theorem-library/extinf-iso.scm.  Library: analysis. LOAD WINDOW: after rabasics-extline2 (ccint-rr-star-iso-bijection, ccint-rr-star-iso-strict)}\par
\vnbprop{rr-star-le-antisymm-given-ends}{\vnbstmt{\begin{array}{@{}l@{}} \forall x, y \in \operatorname{rr\text{-}star}. \\ \quad (\neg (\operatorname{pos\text{-}inf} \leq \operatorname{neg\text{-}inf}) \wedge ((x \leq y) \wedge (y \leq x))) \Rightarrow (x = y) \end{array}}}
\vnbprop{ccint-rr-star-iso-mono}{\vnbstmt{\begin{array}{@{}l@{}} \forall u, v \in \operatorname{ccint}(-1, 1). \\ \quad (u \leq v) \Rightarrow  \\ \quad \quad (\lambda x \in \operatorname{ccint}(-1, 1).\; \left(\text{if }(x = 1)\text{ then }\operatorname{pos\text{-}inf}\text{ else }\left(\text{if }(x = -1)\text{ then }\operatorname{neg\text{-}inf}\text{ else }/(x, (1 - \lvert x \rvert))\right)\right))(u) \\ \quad \quad \leq\; (\lambda x \in \operatorname{ccint}(-1, 1).\; \left(\text{if }(x = 1)\text{ then }\operatorname{pos\text{-}inf}\text{ else }\left(\text{if }(x = -1)\text{ then }\operatorname{neg\text{-}inf}\text{ else }/(x, (1 - \lvert x \rvert))\right)\right))(v) \end{array}}}
\vnbprop{ccint-rr-star-iso-reflects-given-ends}{\vnbstmt{\begin{array}{@{}l@{}} \forall u, v \in \operatorname{ccint}(-1, 1). \\ \quad \neg (\operatorname{pos\text{-}inf} \leq \operatorname{neg\text{-}inf}) \wedge \\ \quad (\lambda x \in \operatorname{ccint}(-1, 1).\; \left(\text{if }(x = 1)\text{ then }\operatorname{pos\text{-}inf}\text{ else }\left(\text{if }(x = -1)\text{ then }\operatorname{neg\text{-}inf}\text{ else }/(x, (1 - \lvert x \rvert))\right)\right))(u) \\ \quad \leq\; (\lambda x \in \operatorname{ccint}(-1, 1).\; \left(\text{if }(x = 1)\text{ then }\operatorname{pos\text{-}inf}\text{ else }\left(\text{if }(x = -1)\text{ then }\operatorname{neg\text{-}inf}\text{ else }/(x, (1 - \lvert x \rvert))\right)\right))(v) \Rightarrow \\ \quad \quad u \leq v \end{array}}}

\vnbbookfile{theorem-library/extorder-iso}{theorem-library/extorder-iso}
\noindent{\small extorder-iso.scm -- calculus.pdf 1.13 in full: x |-$>$ x / (1 - |x|) extends to an order isomorphism of [-1, 1] onto the extended line RR-STAR, $<$= PRESERVED AND REFLECTED. Agent EXT-ORDER, 2026-10-26.  DESTINATION: theorem-library/extorder-iso.scm (or appended to theorem-library/extinf-iso.scm).  Library: analysis.}\par
\vnbprop{rr-star-le-antisymm}{\vnbstmt{\forall x, y \in \operatorname{rr\text{-}star}.\; (((x \leq y) \wedge (y \leq x)) \Rightarrow (x = y))}}
\vnbprop{ccint-rr-star-iso-reflects}{\vnbstmt{\begin{array}{@{}l@{}} \forall u, v \in \operatorname{ccint}(-1, 1). \\ \quad (\lambda x \in \operatorname{ccint}(-1, 1).\; \left(\text{if }(x = 1)\text{ then }\operatorname{pos\text{-}inf}\text{ else }\left(\text{if }(x = -1)\text{ then }\operatorname{neg\text{-}inf}\text{ else }/(x, (1 - \lvert x \rvert))\right)\right))(u) \\ \quad \leq\; (\lambda x \in \operatorname{ccint}(-1, 1).\; \left(\text{if }(x = 1)\text{ then }\operatorname{pos\text{-}inf}\text{ else }\left(\text{if }(x = -1)\text{ then }\operatorname{neg\text{-}inf}\text{ else }/(x, (1 - \lvert x \rvert))\right)\right))(v) \Rightarrow \\ \quad \quad u \leq v \end{array}}}
\vnbprop{ccint-rr-star-iso-le-iff}{\vnbstmt{\begin{array}{@{}l@{}} \forall u, v \in \operatorname{ccint}(-1, 1). \\ \quad u \leq v \\ \quad \Leftrightarrow\; (\lambda x \in \operatorname{ccint}(-1, 1).\; \left(\text{if }(x = 1)\text{ then }\operatorname{pos\text{-}inf}\text{ else }\left(\text{if }(x = -1)\text{ then }\operatorname{neg\text{-}inf}\text{ else }/(x, (1 - \lvert x \rvert))\right)\right))(u) \\ \quad \leq\; (\lambda x \in \operatorname{ccint}(-1, 1).\; \left(\text{if }(x = 1)\text{ then }\operatorname{pos\text{-}inf}\text{ else }\left(\text{if }(x = -1)\text{ then }\operatorname{neg\text{-}inf}\text{ else }/(x, (1 - \lvert x \rvert))\right)\right))(v) \end{array}}}
\vnbprop{ccint-rr-star-order-iso-le}{\vnbstmt{\begin{array}{@{}l@{}} (\lambda x \in \operatorname{ccint}(-1, 1).\; \left(\text{if }(x = 1)\text{ then }\operatorname{pos\text{-}inf}\text{ else }\left(\text{if }(x = -1)\text{ then }\operatorname{neg\text{-}inf}\text{ else }/(x, (1 - \lvert x \rvert))\right)\right)) \\ \in\; \mathrm{Bij}(\operatorname{ccint}(-1, 1), \operatorname{rr\text{-}star}) \wedge \\ \forall u, v \in \operatorname{ccint}(-1, 1). \\ \quad u \leq v \\ \quad \Leftrightarrow\; (\lambda x \in \operatorname{ccint}(-1, 1).\; \left(\text{if }(x = 1)\text{ then }\operatorname{pos\text{-}inf}\text{ else }\left(\text{if }(x = -1)\text{ then }\operatorname{neg\text{-}inf}\text{ else }/(x, (1 - \lvert x \rvert))\right)\right))(u) \\ \quad \leq\; (\lambda x \in \operatorname{ccint}(-1, 1).\; \left(\text{if }(x = 1)\text{ then }\operatorname{pos\text{-}inf}\text{ else }\left(\text{if }(x = -1)\text{ then }\operatorname{neg\text{-}inf}\text{ else }/(x, (1 - \lvert x \rvert))\right)\right))(v) \end{array}}}

\vnbbookfile{theorem-library/qq-countable}{theorem-library/qq-countable}
\noindent{\small theorem-library/qq-countable.scm -- THE RATIONALS ARE COUNTABLE (2026-10-26, after M-A / QQ-DEBT). Helper prefix: ma- (kit), mv- (the group D).  Library: analysis (IS-COUNTABLE is regulated-primitive.scm's).}\par
\vnbprop{nn-fractions-enumerated}{\vnbstmt{\begin{array}{@{}l@{}} \forall d \in \{\mathit{mvd} \in \mathbb{R} : \exists a, b, c \in \mathbb{N}.\; ((a - b) = ((c + 1) \cdot \mathit{mvd}))\}. \exists m \in \mathbb{N}. \\ \quad (\operatorname{recip}((\operatorname{nnsnd}(\operatorname{nnsnd}(m)) + 1)) \cdot (\operatorname{nnfst}(m) - \operatorname{nnfst}(\operatorname{nnsnd}(m)))) = d \end{array}}}
\vnbprop{qq-in-nn-fractions}{\vnbstmt{\begin{array}{@{}l@{}} \forall q \in \mathbb{Q}. \\ \quad q \in \{d \in \mathbb{R} : \exists a, b, c \in \mathbb{N}.\; ((a - b) = ((c + 1) \cdot d))\} \end{array}}}
\vnbprop{qq-enumerated}{\vnbstmt{\begin{array}{@{}l@{}} \forall q \in \mathbb{Q}. \exists m \in \mathbb{N}. \\ \quad (\operatorname{recip}((\operatorname{nnsnd}(\operatorname{nnsnd}(m)) + 1)) \cdot (\operatorname{nnfst}(m) - \operatorname{nnfst}(\operatorname{nnsnd}(m)))) = q \end{array}}}
\vnbprop{qq-enumeration-type}{\vnbstmt{\begin{array}{@{}l@{}} \forall m \in \mathbb{N}. \\ \quad (\operatorname{recip}((\operatorname{nnsnd}(\operatorname{nnsnd}(m)) + 1)) \cdot (\operatorname{nnfst}(m) - \operatorname{nnfst}(\operatorname{nnsnd}(m)))) \in \mathbb{Q} \end{array}}}
\vnbprop{qq-countable}{\vnbstmt{\operatorname{is\text{-}countable}(\mathbb{Q})}}

\vnbbookfile{theorem-library/catools-evt}{theorem-library/catools-evt}
\noindent{\small catools-evt.scm -- the extreme value theorem on a compact metric space (CA-TOOLS, 2026-10-26). Destination: theorem-library/catools-evt.scm, library analysis.}\par
\vnbprop{continuous-real-near}{\vnbstmt{\begin{array}{@{}l@{}} \forall s, f.; \forall c \in \operatorname{pts}(s), a. \\ \quad (\operatorname{is\text{-}continuous}(s, \operatorname{rr\text{-}ms}, f) \wedge \operatorname{pos\text{-}rr}(a)) \Rightarrow  \\ \quad \quad \exists l. \\ \quad \quad \quad \operatorname{pos\text{-}rr}(l) \wedge \\ \quad \quad \quad \forall b \in \operatorname{pts}(s). \\ \quad \quad \quad \quad (\operatorname{dist}(s)(b, c) < l) \Rightarrow ((f(b) \leq (f(c) + a)) \wedge (f(c) \leq (f(b) + a))) \end{array}}}
\vnbprop{compact-continuous-max}{\vnbstmt{\begin{array}{@{}l@{}} \forall s, f.; \forall a \in \operatorname{pts}(s). \\ \quad (\operatorname{is\text{-}compact}(s) \wedge \operatorname{is\text{-}continuous}(s, \operatorname{rr\text{-}ms}, f)) \Rightarrow  \\ \quad \quad \exists p \in \operatorname{pts}(s).\; \forall z \in \operatorname{pts}(s).\; (f(z) \leq f(p)) \end{array}}}
\vnbprop{neg-continuous-real}{\vnbstmt{\begin{array}{@{}l@{}} \forall s, f. \\ \quad \operatorname{is\text{-}continuous}(s, \operatorname{rr\text{-}ms}, f) \Rightarrow  \\ \quad \quad \operatorname{is\text{-}continuous}(s, \operatorname{rr\text{-}ms}, (\lambda n \in \operatorname{pts}(s).\; -f(n))) \end{array}}}
\vnbprop{compact-continuous-min}{\vnbstmt{\begin{array}{@{}l@{}} \forall s, f.; \forall a \in \operatorname{pts}(s). \\ \quad (\operatorname{is\text{-}compact}(s) \wedge \operatorname{is\text{-}continuous}(s, \operatorname{rr\text{-}ms}, f)) \Rightarrow  \\ \quad \quad \exists p \in \operatorname{pts}(s).\; \forall z \in \operatorname{pts}(s).\; (f(p) \leq f(z)) \end{array}}}

\vnbbookfile{theorem-library/regulated-step-approx}{theorem-library/regulated-step-approx}
\noindent{\small regulated-step-approx.scm -- Dieudonne (7.6.1), NECESSITY: a function regulated on [a, b] is the uniform limit on [a, b] of step functions. (Batch 22-A, 2026-09-23.)}\par
\vnbprop{strict-seq-append-point}{\vnbstmt{\begin{array}{@{}l@{}} \forall p, n, d. \\ \quad p \in (\mathbb{N} \to \mathbb{R}) \wedge \\ \quad n \in \mathbb{N} \wedge \\ \quad d \in \mathbb{R} \wedge \\ \quad \forall i, j. \\ \quad \quad ((i \in \operatorname{interval}(0, n)) \wedge ((j \in \operatorname{interval}(0, n)) \wedge (i < j))) \Rightarrow (p(i) < p(j)) \wedge \\ \quad p(n) < d \Rightarrow \\ \quad \quad \exists q \in (\mathbb{N} \to \mathbb{R}). \\ \quad \quad \quad q((n + 1)) = d \wedge \\ \quad \quad \quad \forall k \in \mathbb{N}.\; ((k \leq n) \Rightarrow (q(k) = p(k))) \wedge \\ \quad \quad \quad \forall i, j. \\ \quad \quad \quad \quad ((i \in \operatorname{interval}(0, (n + 1))) \wedge ((j \in \operatorname{interval}(0, (n + 1))) \wedge (i < j))) \Rightarrow  \\ \quad \quad \quad \quad \quad q(i) < q(j) \end{array}}}
\vnbprop{fun-append-constant-piece}{\vnbstmt{\begin{array}{@{}l@{}} \forall h \in (\mathbb{R} \to \mathbb{R}), c, d, k, v \in \mathbb{R}. \\ \quad (c < d) \Rightarrow  \\ \quad \quad \exists g \in (\mathbb{R} \to \mathbb{R}). \\ \quad \quad \quad \forall t \in \mathbb{R}.\; ((t \leq c) \Rightarrow (g(t) = h(t))) \wedge \\ \quad \quad \quad \forall t \in \mathbb{R}.\; (((c < t) \wedge (t < d)) \Rightarrow (g(t) = k)) \wedge \\ \quad \quad \quad g(d) = v \end{array}}}
\vnbprop{step-pieces-append}{\vnbstmt{\begin{array}{@{}l@{}} \forall p, n, h, q, g, c, d, k.; \forall i \in \operatorname{interval}(0, (n + 1)). \\ \quad p \in (\mathbb{N} \to \mathbb{R}) \wedge \\ \quad n \in \mathbb{N} \wedge \\ \quad h \in (\mathbb{R} \to \mathbb{R}) \wedge \\ \quad q \in (\mathbb{N} \to \mathbb{R}) \wedge \\ \quad g \in (\mathbb{R} \to \mathbb{R}) \wedge \\ \quad c \in \mathbb{R} \wedge \\ \quad d \in \mathbb{R} \wedge \\ \quad k \in \mathbb{R} \wedge \\ \quad \forall \mathit{rsi}, j. \\ \quad \quad ((\mathit{rsi} \in \operatorname{interval}(0, n)) \wedge ((j \in \operatorname{interval}(0, n)) \wedge (\mathit{rsi} < j))) \Rightarrow (p(\mathit{rsi}) < p(j)) \wedge \\ \quad \forall i \in \operatorname{interval}(0, n). \\ \quad \quad (i < n) \Rightarrow  \\ \quad \quad \quad \exists \mathit{rspc} \in \mathbb{R}. \forall t \in \mathbb{R}. \\ \quad \quad \quad \quad ((p(i) < t) \wedge (t < p((i + 1)))) \Rightarrow (h(t) = \mathit{rspc}) \wedge \\ \quad p(n) = c \wedge \\ \quad \forall j \in \mathbb{N}.\; ((j \leq n) \Rightarrow (q(j) = p(j))) \wedge \\ \quad q((n + 1)) = d \wedge \\ \quad \forall \mathit{rst} \in \mathbb{R}.\; ((\mathit{rst} \leq c) \Rightarrow (g(\mathit{rst}) = h(\mathit{rst}))) \wedge \\ \quad \forall \mathit{rst} \in \mathbb{R}.\; (((c < \mathit{rst}) \wedge (\mathit{rst} < d)) \Rightarrow (g(\mathit{rst}) = k)) \wedge \\ \quad i < (n + 1) \Rightarrow \\ \quad \quad \exists \mathit{rspc} \in \mathbb{R}. \forall t \in \mathbb{R}. \\ \quad \quad \quad ((q(i) < t) \wedge (t < q((i + 1)))) \Rightarrow (g(t) = \mathit{rspc}) \end{array}}}
\vnbprop{step-approx-start}{\vnbstmt{\begin{array}{@{}l@{}} \forall a, f, e. \\ \quad ((a \in \mathbb{R}) \wedge ((f \in (\mathbb{R} \to \mathbb{R})) \wedge ((e \in \mathbb{R}) \wedge (0 \leq e)))) \Rightarrow  \\ \quad \quad \exists n \in \mathbb{N}, p \in (\mathbb{N} \to \mathbb{R}), h \in (\mathbb{R} \to \mathbb{R}). \\ \quad \quad \quad p(0) = a \wedge \\ \quad \quad \quad p(n) = a \wedge \\ \quad \quad \quad \forall i, j. \\ \quad \quad \quad \quad ((i \in \operatorname{interval}(0, n)) \wedge ((j \in \operatorname{interval}(0, n)) \wedge (i < j))) \Rightarrow (p(i) < p(j)) \wedge \\ \quad \quad \quad \forall \mathit{rspi} \in \operatorname{interval}(0, n). \\ \quad \quad \quad \quad (\mathit{rspi} < n) \Rightarrow  \\ \quad \quad \quad \quad \quad \exists c \in \mathbb{R}. \forall t \in \mathbb{R}. \\ \quad \quad \quad \quad \quad \quad ((p(\mathit{rspi}) < t) \wedge (t < p((\mathit{rspi} + 1)))) \Rightarrow (h(t) = c) \wedge \\ \quad \quad \quad \forall z \in \operatorname{ccint}(a, a).\; (\lvert (f(z) - h(z)) \rvert \leq e) \end{array}}}
\vnbprop{step-approx-append}{\vnbstmt{\begin{array}{@{}l@{}} \forall a, x, y, f, e, k. \\ \quad a \in \mathbb{R} \wedge \\ \quad x \in \mathbb{R} \wedge \\ \quad y \in \mathbb{R} \wedge \\ \quad f \in (\mathbb{R} \to \mathbb{R}) \wedge \\ \quad e \in \mathbb{R} \wedge \\ \quad k \in \mathbb{R} \wedge \\ \quad 0 \leq e \wedge \\ \quad \exists n \in \mathbb{N}, p \in (\mathbb{N} \to \mathbb{R}), h \in (\mathbb{R} \to \mathbb{R}). \\ \quad \quad p(0) = a \wedge \\ \quad \quad p(n) = x \wedge \\ \quad \quad \forall i, j. \\ \quad \quad \quad ((i \in \operatorname{interval}(0, n)) \wedge ((j \in \operatorname{interval}(0, n)) \wedge (i < j))) \Rightarrow (p(i) < p(j)) \wedge \\ \quad \quad \forall \mathit{rspi} \in \operatorname{interval}(0, n). \\ \quad \quad \quad (\mathit{rspi} < n) \Rightarrow  \\ \quad \quad \quad \quad \exists c \in \mathbb{R}. \forall t \in \mathbb{R}. \\ \quad \quad \quad \quad \quad ((p(\mathit{rspi}) < t) \wedge (t < p((\mathit{rspi} + 1)))) \Rightarrow (h(t) = c) \wedge \\ \quad \quad \forall z \in \operatorname{ccint}(a, x).\; (\lvert (f(z) - h(z)) \rvert \leq e) \wedge \\ \quad x < y \wedge \\ \quad \forall s \in \mathbb{R}.\; (((x < s) \wedge (s < y)) \Rightarrow (\lvert (f(s) - k) \rvert \leq e)) \Rightarrow \\ \quad \quad \exists n \in \mathbb{N}, p \in (\mathbb{N} \to \mathbb{R}), h \in (\mathbb{R} \to \mathbb{R}). \\ \quad \quad \quad p(0) = a \wedge \\ \quad \quad \quad p(n) = y \wedge \\ \quad \quad \quad \forall i, j. \\ \quad \quad \quad \quad ((i \in \operatorname{interval}(0, n)) \wedge ((j \in \operatorname{interval}(0, n)) \wedge (i < j))) \Rightarrow (p(i) < p(j)) \wedge \\ \quad \quad \quad \forall \mathit{rspi} \in \operatorname{interval}(0, n). \\ \quad \quad \quad \quad (\mathit{rspi} < n) \Rightarrow  \\ \quad \quad \quad \quad \quad \exists c \in \mathbb{R}. \forall t \in \mathbb{R}. \\ \quad \quad \quad \quad \quad \quad ((p(\mathit{rspi}) < t) \wedge (t < p((\mathit{rspi} + 1)))) \Rightarrow (h(t) = c) \wedge \\ \quad \quad \quad \forall z \in \operatorname{ccint}(a, y).\; (\lvert (f(z) - h(z)) \rvert \leq e) \end{array}}}
\vnbprop{regulated-on-right-window}{\vnbstmt{\begin{array}{@{}l@{}} \forall a, b, f, g, e, t. \\ \quad \operatorname{is\text{-}regulated\text{-}on}(f, a, b) \wedge \\ \quad g \in (\mathbb{R} \to \mathbb{R}) \wedge \\ \quad \forall z \in \operatorname{ccint}(a, b).\; (g(z) = f(z)) \wedge \\ \quad \operatorname{pos\text{-}rr}(e) \wedge \\ \quad t \in \operatorname{ccint}(a, b) \Rightarrow \\ \quad \quad \exists d \in \mathbb{R}. \\ \quad \quad \quad 0 < d \wedge \\ \quad \quad \quad \exists l \in \mathbb{R}. \forall s \in \mathbb{R}. \\ \quad \quad \quad \quad ((t < s) \wedge ((s < (t + d)) \wedge (s \leq b))) \Rightarrow (\lvert (g(s) - l) \rvert \leq e) \end{array}}}
\vnbprop{regulated-on-left-window}{\vnbstmt{\begin{array}{@{}l@{}} \forall a, b, f, g, e, t. \\ \quad \operatorname{is\text{-}regulated\text{-}on}(f, a, b) \wedge \\ \quad g \in (\mathbb{R} \to \mathbb{R}) \wedge \\ \quad \forall z \in \operatorname{ccint}(a, b).\; (g(z) = f(z)) \wedge \\ \quad \operatorname{pos\text{-}rr}(e) \wedge \\ \quad t \in \operatorname{ccint}(a, b) \Rightarrow \\ \quad \quad \exists d \in \mathbb{R}. \\ \quad \quad \quad 0 < d \wedge \\ \quad \quad \quad \exists l \in \mathbb{R}. \forall s \in \mathbb{R}. \\ \quad \quad \quad \quad ((a \leq s) \wedge (((t - d) < s) \wedge (s < t))) \Rightarrow (\lvert (g(s) - l) \rvert \leq e) \end{array}}}
\vnbprop{regulated-on-step-approx-creep}{\vnbstmt{\begin{array}{@{}l@{}} \forall a, b, f, g, e. \\ \quad \operatorname{is\text{-}regulated\text{-}on}(f, a, b) \wedge \\ \quad g \in (\mathbb{R} \to \mathbb{R}) \wedge \\ \quad \forall z \in \operatorname{ccint}(a, b).\; (g(z) = f(z)) \wedge \\ \quad \operatorname{pos\text{-}rr}(e) \Rightarrow \\ \quad \quad \exists n \in \mathbb{N}, p \in (\mathbb{N} \to \mathbb{R}), h \in (\mathbb{R} \to \mathbb{R}). \\ \quad \quad \quad p(0) = a \wedge \\ \quad \quad \quad p(n) = b \wedge \\ \quad \quad \quad \forall i, j. \\ \quad \quad \quad \quad ((i \in \operatorname{interval}(0, n)) \wedge ((j \in \operatorname{interval}(0, n)) \wedge (i < j))) \Rightarrow (p(i) < p(j)) \wedge \\ \quad \quad \quad \forall \mathit{rspi} \in \operatorname{interval}(0, n). \\ \quad \quad \quad \quad (\mathit{rspi} < n) \Rightarrow  \\ \quad \quad \quad \quad \quad \exists c \in \mathbb{R}. \forall t \in \mathbb{R}. \\ \quad \quad \quad \quad \quad \quad ((p(\mathit{rspi}) < t) \wedge (t < p((\mathit{rspi} + 1)))) \Rightarrow (h(t) = c) \wedge \\ \quad \quad \quad \forall \mathit{rsqz} \in \operatorname{ccint}(a, b).\; (\lvert (g(\mathit{rsqz}) - h(\mathit{rsqz})) \rvert \leq e) \end{array}}}
\vnbprop{step-approx-q-step-fn}{\vnbstmt{\begin{array}{@{}l@{}} \forall a, b, f, e. \\ \quad a \in \mathbb{R} \wedge \\ \quad b \in \mathbb{R} \wedge \\ \quad a < b \wedge \\ \quad \exists n \in \mathbb{N}, p \in (\mathbb{N} \to \mathbb{R}), h \in (\mathbb{R} \to \mathbb{R}). \\ \quad \quad p(0) = a \wedge \\ \quad \quad p(n) = b \wedge \\ \quad \quad \forall i, j. \\ \quad \quad \quad ((i \in \operatorname{interval}(0, n)) \wedge ((j \in \operatorname{interval}(0, n)) \wedge (i < j))) \Rightarrow (p(i) < p(j)) \wedge \\ \quad \quad \forall \mathit{rspi} \in \operatorname{interval}(0, n). \\ \quad \quad \quad (\mathit{rspi} < n) \Rightarrow  \\ \quad \quad \quad \quad \exists c \in \mathbb{R}. \forall t \in \mathbb{R}. \\ \quad \quad \quad \quad \quad ((p(\mathit{rspi}) < t) \wedge (t < p((\mathit{rspi} + 1)))) \Rightarrow (h(t) = c) \wedge \\ \quad \quad \forall z \in \operatorname{ccint}(a, b).\; (\lvert (f(z) - h(z)) \rvert \leq e) \Rightarrow \\ \quad \quad \exists \mathit{rsth} \in (\mathbb{R} \to \mathbb{R}). \\ \quad \quad \quad \operatorname{is\text{-}step\text{-}fn}(\mathit{rsth}, a, b) \wedge \forall z \in \operatorname{ccint}(a, b).\; (\lvert (f(z) - \operatorname{rsth}(z)) \rvert \leq e) \end{array}}}
\vnbprop{regulated-on-step-approx}{\vnbstmt{\begin{array}{@{}l@{}} \forall a, b, f, e. \\ \quad (\operatorname{is\text{-}regulated\text{-}on}(f, a, b) \wedge \operatorname{pos\text{-}rr}(e)) \Rightarrow  \\ \quad \quad \exists h \in (\mathbb{R} \to \mathbb{R}). \\ \quad \quad \quad (\operatorname{is\text{-}step\text{-}fn}(h, a, b) \wedge \forall z \in \operatorname{ccint}(a, b).\; (\lvert (\operatorname{extend\text{-}const}(f, a, b)(z) - h(z)) \rvert \leq e)) \end{array}}}
\vnbprop{regulated-on-is-step-limit}{\vnbstmt{\begin{array}{@{}l@{}} \forall a, b \in \mathbb{R}, f. \\ \quad \operatorname{is\text{-}regulated\text{-}on}(f, a, b) \Rightarrow  \\ \quad \quad \exists m \in (\mathbb{N} \to (\mathbb{R} \to \mathbb{R})). \\ \quad \quad \quad (\forall k \in \mathbb{N}.\; \operatorname{is\text{-}step\text{-}fn}(m(k), a, b) \wedge \operatorname{is\text{-}unif\text{-}limit\text{-}on}(m, \operatorname{extend\text{-}const}(f, a, b), a, b)) \end{array}}}

\vnbbookfile{theorem-library/has-deriv-at-more}{theorem-library/has-deriv-at-more}
\noindent{\small has-deriv-at-more.scm -- the four HAS-DERIV-AT laws that theorem-library/interval-calculus-laws.scm did not prove: the VALUE typing, the AFFINE function (and with it the constant and the identity), the REAL MULTIPLE and the DIFFERENCE.}\par
\vnbprop{has-deriv-at-value-in-rr}{\vnbstmt{\forall f, x, l.\; (\operatorname{has\text{-}deriv\text{-}at}(f, x, l) \Rightarrow (f(x) \in \mathbb{R}))}}
\vnbprop{has-deriv-at-affine}{\vnbstmt{\begin{array}{@{}l@{}} \forall r, f.; \forall c, m, x \in \mathbb{R}. \\ \quad \operatorname{pos\text{-}rr}(r) \wedge \\ \quad \operatorname{restrict}(f, \operatorname{ooint}((x - r), (x + r))) \in (\operatorname{ooint}((x - r), (x + r)) \to \mathbb{R}) \wedge \\ \quad \forall \mathit{hbx} \in \operatorname{ooint}((x - r), (x + r)).\; (f(\mathit{hbx}) \simeq (c + (m \cdot \mathit{hbx}))) \Rightarrow \\ \quad \quad \operatorname{has\text{-}deriv\text{-}at}(f, x, m) \end{array}}}
\vnbprop{has-deriv-at-const}{\vnbstmt{\begin{array}{@{}l@{}} \forall r, f.; \forall c, x \in \mathbb{R}. \\ \quad \operatorname{pos\text{-}rr}(r) \wedge \\ \quad \operatorname{restrict}(f, \operatorname{ooint}((x - r), (x + r))) \in (\operatorname{ooint}((x - r), (x + r)) \to \mathbb{R}) \wedge \\ \quad \forall \mathit{hbx} \in \operatorname{ooint}((x - r), (x + r)).\; (f(\mathit{hbx}) \simeq c) \Rightarrow \\ \quad \quad \operatorname{has\text{-}deriv\text{-}at}(f, x, 0) \end{array}}}
\vnbprop{has-deriv-at-identity}{\vnbstmt{\begin{array}{@{}l@{}} \forall r, f.; \forall x \in \mathbb{R}. \\ \quad \operatorname{pos\text{-}rr}(r) \wedge \\ \quad \operatorname{restrict}(f, \operatorname{ooint}((x - r), (x + r))) \in (\operatorname{ooint}((x - r), (x + r)) \to \mathbb{R}) \wedge \\ \quad \forall \mathit{hbx} \in \operatorname{ooint}((x - r), (x + r)).\; (f(\mathit{hbx}) \simeq \mathit{hbx}) \Rightarrow \\ \quad \quad \operatorname{has\text{-}deriv\text{-}at}(f, x, 1) \end{array}}}
\vnbprop{has-deriv-at-real-mul}{\vnbstmt{\begin{array}{@{}l@{}} \forall r, g, h.; \forall c \in \mathbb{R}, x, m. \\ \quad \operatorname{pos\text{-}rr}(r) \wedge \\ \quad \operatorname{restrict}(h, \operatorname{ooint}((x - r), (x + r))) \in (\operatorname{ooint}((x - r), (x + r)) \to \mathbb{R}) \wedge \\ \quad \forall \mathit{hbx} \in \operatorname{ooint}((x - r), (x + r)).\; (h(\mathit{hbx}) \simeq (c \cdot g(\mathit{hbx}))) \wedge \\ \quad \operatorname{has\text{-}deriv\text{-}at}(g, x, m) \Rightarrow \\ \quad \quad \operatorname{has\text{-}deriv\text{-}at}(h, x, (c \cdot m)) \end{array}}}
\vnbprop{has-deriv-at-difference}{\vnbstmt{\begin{array}{@{}l@{}} \forall r, f, g, h, x, l, m. \\ \quad \operatorname{pos\text{-}rr}(r) \wedge \\ \quad \operatorname{restrict}(h, \operatorname{ooint}((x - r), (x + r))) \in (\operatorname{ooint}((x - r), (x + r)) \to \mathbb{R}) \wedge \\ \quad \operatorname{restrict}(g, \operatorname{ooint}((x - r), (x + r))) \in (\operatorname{ooint}((x - r), (x + r)) \to \mathbb{R}) \wedge \\ \quad \forall \mathit{hbx} \in \operatorname{ooint}((x - r), (x + r)).\; (h(\mathit{hbx}) \simeq (f(\mathit{hbx}) - g(\mathit{hbx}))) \wedge \\ \quad \operatorname{has\text{-}deriv\text{-}at}(f, x, l) \wedge \\ \quad \operatorname{has\text{-}deriv\text{-}at}(g, x, m) \Rightarrow \\ \quad \quad \operatorname{has\text{-}deriv\text{-}at}(h, x, (l - m)) \end{array}}}

\vnbbookfile{theorem-library/has-deriv-at-chain}{theorem-library/has-deriv-at-chain}
\noindent{\small has-deriv-at-chain.scm -- THE CHAIN RULE FOR THE LOCAL DERIVATIVE, and the corollary at an affine inner map.}\par
\vnbprop{has-deriv-at-maps-into}{\vnbstmt{\begin{array}{@{}l@{}} \forall f, x, l, d, b. \\ \quad (\operatorname{has\text{-}deriv\text{-}at}(f, x, l) \wedge (\operatorname{pos\text{-}rr}(d) \wedge \operatorname{pos\text{-}rr}(b))) \Rightarrow  \\ \quad \quad \exists w. \\ \quad \quad \quad (\operatorname{pos\text{-}rr}(w) \wedge ((w \leq b) \wedge \forall y \in \operatorname{ooint}((x - w), (x + w)).\; (f(y) \in \operatorname{ooint}((f(x) - d), (f(x) + d))))) \end{array}}}
\vnbprop{has-deriv-at-chain}{\vnbstmt{\begin{array}{@{}l@{}} \forall r, f, g, h, x, l, m. \\ \quad \operatorname{pos\text{-}rr}(r) \wedge \\ \quad \operatorname{restrict}(h, \operatorname{ooint}((x - r), (x + r))) \in (\operatorname{ooint}((x - r), (x + r)) \to \mathbb{R}) \wedge \\ \quad \forall y \in \operatorname{ooint}((x - r), (x + r)).\; (h(y) \simeq g(f(y))) \wedge \\ \quad \operatorname{has\text{-}deriv\text{-}at}(f, x, l) \wedge \\ \quad \operatorname{has\text{-}deriv\text{-}at}(g, f(x), m) \Rightarrow \\ \quad \quad \operatorname{has\text{-}deriv\text{-}at}(h, x, (m \cdot l)) \end{array}}}
\vnbprop{has-deriv-at-affine-chain}{\vnbstmt{\begin{array}{@{}l@{}} \forall r, g, h.; \forall p, q, x \in \mathbb{R}, m. \\ \quad \operatorname{pos\text{-}rr}(r) \wedge \\ \quad \operatorname{restrict}(h, \operatorname{ooint}((x - r), (x + r))) \in (\operatorname{ooint}((x - r), (x + r)) \to \mathbb{R}) \wedge \\ \quad \forall y \in \operatorname{ooint}((x - r), (x + r)).\; (h(y) \simeq g((p + (q \cdot y)))) \wedge \\ \quad \operatorname{has\text{-}deriv\text{-}at}(g, (p + (q \cdot x)), m) \Rightarrow \\ \quad \quad \operatorname{has\text{-}deriv\text{-}at}(h, x, (m \cdot q)) \end{array}}}

\vnbbookfile{theorem-library/interval-taylor}{theorem-library/interval-taylor}
\noindent{\small interval-taylor.scm -- TAYLOR'S FORMULA (2.18) FOR A FUNCTION ON ITS INTERVAL, AND ITS TWO COROLLARIES.}\par
\vnbprop{gmvt-with-known-derivative}{\vnbstmt{\begin{array}{@{}l@{}} \forall a, b \in \mathbb{R}, f, g.; \forall v \in (\operatorname{ooint}(a, b) \to \mathbb{R}). \\ \quad a < b \wedge \\ \quad f \in (\operatorname{ccint}(a, b) \to \mathbb{R}) \wedge \\ \quad \operatorname{is\text{-}continuous\text{-}on}(f, \operatorname{ccint}(a, b)) \wedge \\ \quad \forall x \in \operatorname{ooint}(a, b).\; \operatorname{has\text{-}deriv\text{-}at}(f, x, v(x)) \wedge \\ \quad g \in (\operatorname{ccint}(a, b) \to \mathbb{R}) \wedge \\ \quad \operatorname{is\text{-}continuous\text{-}on}(g, \operatorname{ccint}(a, b)) \wedge \\ \quad \forall x \in \operatorname{ooint}(a, b).\; \exists m.\; \operatorname{has\text{-}deriv\text{-}at}(g, x, m) \Rightarrow \\ \quad \quad \exists i \in \operatorname{ooint}(a, b), m. \\ \quad \quad \quad \operatorname{has\text{-}deriv\text{-}at}(g, i, m) \wedge ((v(i) \cdot (g(b) - g(a))) = (m \cdot (f(b) - f(a)))) \end{array}}}
\vnbprop{taylor-family-order-in-nn}{\vnbstmt{\forall d, n, a, b.\; (\operatorname{is\text{-}taylor\text{-}family}(d, n, a, b) \Rightarrow (n \in \mathbb{N}))}}
\vnbprop{taylor-family-order-pos}{\vnbstmt{\forall d, n, a, b.\; (\operatorname{is\text{-}taylor\text{-}family}(d, n, a, b) \Rightarrow (1 \leq n))}}
\vnbprop{taylor-family-left-in-rr}{\vnbstmt{\forall d, n, a, b.\; (\operatorname{is\text{-}taylor\text{-}family}(d, n, a, b) \Rightarrow (a \in \mathbb{R}))}}
\vnbprop{taylor-family-right-in-rr}{\vnbstmt{\forall d, n, a, b.\; (\operatorname{is\text{-}taylor\text{-}family}(d, n, a, b) \Rightarrow (b \in \mathbb{R}))}}
\vnbprop{taylor-family-endpoints-lt}{\vnbstmt{\forall d, n, a, b.\; (\operatorname{is\text{-}taylor\text{-}family}(d, n, a, b) \Rightarrow (a < b))}}
\vnbprop{taylor-family-top-in-fun}{\vnbstmt{\begin{array}{@{}l@{}} \forall d, n, a, b. \\ \quad \operatorname{is\text{-}taylor\text{-}family}(d, n, a, b) \Rightarrow (d(n) \in (\operatorname{ooint}(a, b) \to \mathbb{R})) \end{array}}}
\vnbprop{taylor-family-in-fun}{\vnbstmt{\begin{array}{@{}l@{}} \forall d, n, a, b.; \forall k \in \mathbb{N}. \\ \quad (\operatorname{is\text{-}taylor\text{-}family}(d, n, a, b) \wedge (k < n)) \Rightarrow (d(k) \in (\operatorname{ccint}(a, b) \to \mathbb{R})) \end{array}}}
\vnbprop{taylor-family-continuous}{\vnbstmt{\begin{array}{@{}l@{}} \forall d, n, a, b.; \forall k \in \mathbb{N}. \\ \quad (\operatorname{is\text{-}taylor\text{-}family}(d, n, a, b) \wedge (k < n)) \Rightarrow  \\ \quad \quad \operatorname{is\text{-}continuous\text{-}on}(d(k), \operatorname{ccint}(a, b)) \end{array}}}
\vnbprop{taylor-family-deriv}{\vnbstmt{\begin{array}{@{}l@{}} \forall d, n, a, b.; \forall k \in \mathbb{N}, x \in \operatorname{ooint}(a, b). \\ \quad (\operatorname{is\text{-}taylor\text{-}family}(d, n, a, b) \wedge (k < n)) \Rightarrow  \\ \quad \quad \operatorname{has\text{-}deriv\text{-}at}(d(k), x, d((k + 1))(x)) \end{array}}}
\vnbprop{taylor-aux-value}{\vnbstmt{\begin{array}{@{}l@{}} \forall d, n, a, b.; \forall z \in \operatorname{ccint}(a, b). \\ \quad (\operatorname{taylor\text{-}aux}(d, n, a, b)(z) \simeq \operatorname{taylor\text{-}sum}(d, n, b, z)) \end{array}}}
\vnbprop{taylor-clear-recip}{\vnbstmt{\begin{array}{@{}l@{}} \forall p, q, r, s, u, w \in \mathbb{R}. \\ \quad (((s \cdot r) = 1) \wedge ((((p \cdot q) \cdot r) \cdot w) = u)) \Rightarrow  \\ \quad \quad ((u \cdot s) = ((p \cdot q) \cdot w)) \wedge (u = (((p \cdot r) \cdot q) \cdot w)) \end{array}}}
\vnbprop{series-partial-sum-in-rr-below}{\vnbstmt{\begin{array}{@{}l@{}} \forall n \in \mathbb{N}, f. \\ \quad \forall k \in \mathbb{N}.\; ((k < n) \Rightarrow (f(k) \in \mathbb{R})) \Rightarrow  \\ \quad \quad \operatorname{series\text{-}partial\text{-}sum}(f, n) \in \mathbb{R} \end{array}}}
\vnbprop{taylor-term-value}{\vnbstmt{\begin{array}{@{}l@{}} \forall d, b, v.; \forall k \in \mathbb{N}. \\ \quad (\operatorname{taylor\text{-}term}(d, b, v)(k) \simeq ((d(k)(v) \cdot \operatorname{power}((b - v), k)) \cdot \operatorname{recip}(\operatorname{factorial}(k)))) \end{array}}}
\vnbprop{taylor-sum-unfold}{\vnbstmt{\begin{array}{@{}l@{}} \forall d, n, b, v. \\ \quad (\operatorname{taylor\text{-}sum}(d, n, b, v) \simeq \operatorname{series\text{-}partial\text{-}sum}(\operatorname{taylor\text{-}term}(d, b, v), n)) \end{array}}}
\vnbprop{taylor-sum-in-rr}{\vnbstmt{\begin{array}{@{}l@{}} \forall d, n, a, b.; \forall j \in \mathbb{N}, z \in \operatorname{ccint}(a, b). \\ \quad (\operatorname{is\text{-}taylor\text{-}family}(d, n, a, b) \wedge (j \leq n)) \Rightarrow (\operatorname{taylor\text{-}sum}(d, j, b, z) \in \mathbb{R}) \end{array}}}
\vnbprop{taylor-aux-in-fun}{\vnbstmt{\begin{array}{@{}l@{}} \forall d, n, a, b.; \forall j \in \mathbb{N}. \\ \quad (\operatorname{is\text{-}taylor\text{-}family}(d, n, a, b) \wedge (j \leq n)) \Rightarrow  \\ \quad \quad \operatorname{taylor\text{-}aux}(d, j, a, b) \in (\operatorname{ccint}(a, b) \to \mathbb{R}) \end{array}}}
\vnbprop{taylor-sum-one}{\vnbstmt{\begin{array}{@{}l@{}} \forall d.; \forall b, v \in \mathbb{R}. \\ \quad (d(0)(v) \in \mathbb{R}) \Rightarrow (\operatorname{taylor\text{-}sum}(d, 1, b, v) = d(0)(v)) \end{array}}}
\vnbprop{taylor-sum-at-endpoint}{\vnbstmt{\begin{array}{@{}l@{}} \forall m \in \mathbb{N}, d.; \forall b \in \mathbb{R}. \\ \quad \forall k \in \mathbb{N}.\; ((k \leq m) \Rightarrow (d(k)(b) \in \mathbb{R})) \Rightarrow  \\ \quad \quad \operatorname{taylor\text{-}sum}(d, (m + 1), b, b) = d(0)(b) \end{array}}}
\vnbprop{taylor-family-sum-at-endpoint}{\vnbstmt{\begin{array}{@{}l@{}} \forall d, n, a, b.; \forall m \in \mathbb{N}. \\ \quad (\operatorname{is\text{-}taylor\text{-}family}(d, n, a, b) \wedge ((m + 1) \leq n)) \Rightarrow  \\ \quad \quad \operatorname{taylor\text{-}sum}(d, (m + 1), b, b) = d(0)(b) \end{array}}}
\vnbprop{cont-on-in-fun}{\vnbstmt{\forall f, m.\; (\operatorname{is\text{-}continuous\text{-}on}(f, m) \Rightarrow (f \in (m \to \mathbb{R})))}}
\vnbprop{cont-on-transfer-ptwise-eq}{\vnbstmt{\begin{array}{@{}l@{}} \forall m, f, g. \\ \quad ((m \subseteq \mathbb{R}) \wedge ((f \in (m \to \mathbb{R})) \wedge (\operatorname{is\text{-}continuous\text{-}on}(g, m) \wedge \forall z \in m.\; (f(z) = g(z))))) \Rightarrow  \\ \quad \quad \operatorname{is\text{-}continuous\text{-}on}(f, m) \end{array}}}
\vnbprop{cont-on-of-total}{\vnbstmt{\begin{array}{@{}l@{}} \forall m.; \forall g \in (\mathbb{R} \to \mathbb{R}), f \in (m \to \mathbb{R}). \\ \quad m \subseteq \mathbb{R} \wedge \\ \quad \forall y \in \mathbb{R}.\; \operatorname{is\text{-}continuous\text{-}at}(\operatorname{rr\text{-}ms}, \operatorname{rr\text{-}ms}, g, y) \wedge \\ \quad \forall z \in m.\; (f(z) = g(z)) \Rightarrow \\ \quad \quad \operatorname{is\text{-}continuous\text{-}on}(f, m) \end{array}}}
\vnbprop{power-sub-succ-continuous}{\vnbstmt{\begin{array}{@{}l@{}} \forall b \in \mathbb{R}, k \in \mathbb{N}, y \in \mathbb{R}. \\ \quad \operatorname{is\text{-}continuous\text{-}at}(\operatorname{rr\text{-}ms}, \operatorname{rr\text{-}ms}, (\lambda z \in \mathbb{R}.\; \operatorname{power}((b - z), (k + 1))), y) \end{array}}}
\vnbprop{taylor-step-continuous}{\vnbstmt{\begin{array}{@{}l@{}} \forall b, c \in \mathbb{R}, k \in \mathbb{N}, g, f \in (\mathbb{R} \to \mathbb{R}), y \in \mathbb{R}. \\ \quad (\operatorname{is\text{-}continuous\text{-}at}(\operatorname{rr\text{-}ms}, \operatorname{rr\text{-}ms}, g, y) \wedge \operatorname{is\text{-}continuous\text{-}at}(\operatorname{rr\text{-}ms}, \operatorname{rr\text{-}ms}, f, y)) \Rightarrow  \\ \quad \quad \operatorname{is\text{-}continuous\text{-}at}(\operatorname{rr\text{-}ms}, \operatorname{rr\text{-}ms}, (\lambda w \in \mathbb{R}.\; (g(w) + ((f(w) \cdot \operatorname{power}((b - w), (k + 1))) \cdot c))), y) \end{array}}}
\vnbprop{taylor-aux-step}{\vnbstmt{\begin{array}{@{}l@{}} \forall d, n, a, b.; \forall j \in \mathbb{N}, z \in \operatorname{ccint}(a, b). \\ \quad (\operatorname{is\text{-}taylor\text{-}family}(d, n, a, b) \wedge ((j + 1) \leq n)) \Rightarrow  \\ \quad \quad \operatorname{taylor\text{-}aux}(d, (j + 1), a, b)(z) \\ \quad \quad =\; \operatorname{taylor\text{-}aux}(d, j, a, b)(z) + ((d(j)(z) \cdot \operatorname{power}((b - z), j)) \cdot \operatorname{recip}(\operatorname{factorial}(j))) \end{array}}}
\vnbprop{taylor-aux-continuous-on}{\vnbstmt{\begin{array}{@{}l@{}} \forall m \in \mathbb{N}, d, n, a, b. \\ \quad (\operatorname{is\text{-}taylor\text{-}family}(d, n, a, b) \wedge ((m + 1) \leq n)) \Rightarrow  \\ \quad \quad \operatorname{is\text{-}continuous\text{-}on}(\operatorname{taylor\text{-}aux}(d, (m + 1), a, b), \operatorname{ccint}(a, b)) \end{array}}}
\vnbprop{ooint-interior-radius}{\vnbstmt{\begin{array}{@{}l@{}} \forall a, b \in \mathbb{R}, x \in \operatorname{ooint}(a, b). \exists r. \\ \quad \operatorname{pos\text{-}rr}(r) \wedge (\operatorname{ooint}((x - r), (x + r)) \subseteq \operatorname{ooint}(a, b)) \end{array}}}
\vnbprop{power-sub-succ-has-deriv}{\vnbstmt{\begin{array}{@{}l@{}} \forall b \in \mathbb{R}, k \in \mathbb{N}, x \in \mathbb{R}. \\ \quad \operatorname{has\text{-}deriv\text{-}at}((\lambda z \in \mathbb{R}.\; \operatorname{power}((b - z), (k + 1))), x, (((k + 1) \cdot \operatorname{power}((b - x), k)) \cdot -1)) \end{array}}}
\vnbprop{taylor-step-identity}{\vnbstmt{\begin{array}{@{}l@{}} \forall a, b, p, q, r, s \in \mathbb{R}. \\ \quad (((a \cdot p) \cdot (s \cdot r)) + (r \cdot ((b \cdot q) + (((s \cdot p) \cdot -1) \cdot a)))) = ((b \cdot q) \cdot r) \end{array}}}
\vnbprop{taylor-aux-has-deriv}{\vnbstmt{\begin{array}{@{}l@{}} \forall m \in \mathbb{N}, d, n, a, b.; \forall x \in \operatorname{ooint}(a, b). \\ \quad (\operatorname{is\text{-}taylor\text{-}family}(d, n, a, b) \wedge ((m + 1) \leq n)) \Rightarrow  \\ \quad \quad \operatorname{has\text{-}deriv\text{-}at}(\operatorname{taylor\text{-}aux}(d, (m + 1), a, b), x, ((d((m + 1))(x) \cdot \operatorname{power}((b - x), m)) \cdot \operatorname{recip}(\operatorname{factorial}(m)))) \end{array}}}
\vnbprop{taylor-formula-on-interval-from-aux}{\vnbstmt{\begin{array}{@{}l@{}} \forall a, b \in \mathbb{R}, d.; \forall m \in \mathbb{N}, g \in (\operatorname{ccint}(a, b) \to \mathbb{R}). \\ \quad \operatorname{is\text{-}taylor\text{-}family}(d, (m + 1), a, b) \wedge \\ \quad \operatorname{is\text{-}continuous\text{-}on}(\operatorname{taylor\text{-}aux}(d, (m + 1), a, b), \operatorname{ccint}(a, b)) \wedge \\ \quad \forall x \in \operatorname{ooint}(a, b). \\ \quad \quad \operatorname{has\text{-}deriv\text{-}at}(\operatorname{taylor\text{-}aux}(d, (m + 1), a, b), x, ((d((m + 1))(x) \cdot \operatorname{power}((b - x), m)) \cdot \operatorname{recip}(\operatorname{factorial}(m)))) \wedge \\ \quad \operatorname{taylor\text{-}sum}(d, (m + 1), b, b) = d(0)(b) \wedge \\ \quad \operatorname{is\text{-}continuous\text{-}on}(g, \operatorname{ccint}(a, b)) \wedge \\ \quad \forall x \in \operatorname{ooint}(a, b).\; \exists l.\; \operatorname{has\text{-}deriv\text{-}at}(g, x, l) \Rightarrow \\ \quad \quad \exists i \in \operatorname{ooint}(a, b), u. \\ \quad \quad \quad \operatorname{has\text{-}deriv\text{-}at}(g, i, u) \wedge \\ \quad \quad \quad ((u \cdot \operatorname{taylor\text{-}rem}(d, (m + 1), a, b)) \cdot \operatorname{factorial}(m)) = ((d((m + 1))(i) \cdot \operatorname{power}((b - i), m)) \cdot (g(b) - g(a))) \wedge \\ \quad \quad \quad u \cdot \operatorname{taylor\text{-}rem}(d, (m + 1), a, b) \\ \quad \quad \quad =\; ((d((m + 1))(i) \cdot \operatorname{recip}(\operatorname{factorial}(m))) \cdot \operatorname{power}((b - i), m)) \cdot (g(b) - g(a)) \end{array}}}
\vnbprop{taylor-formula-on-interval}{\vnbstmt{\begin{array}{@{}l@{}} \forall a, b \in \mathbb{R}, d.; \forall m \in \mathbb{N}, g \in (\operatorname{ccint}(a, b) \to \mathbb{R}). \\ \quad \operatorname{is\text{-}taylor\text{-}family}(d, (m + 1), a, b) \wedge \\ \quad \operatorname{is\text{-}continuous\text{-}on}(g, \operatorname{ccint}(a, b)) \wedge \\ \quad \forall x \in \operatorname{ooint}(a, b).\; \exists l.\; \operatorname{has\text{-}deriv\text{-}at}(g, x, l) \Rightarrow \\ \quad \quad \exists i \in \operatorname{ooint}(a, b), u. \\ \quad \quad \quad \operatorname{has\text{-}deriv\text{-}at}(g, i, u) \wedge \\ \quad \quad \quad ((u \cdot \operatorname{taylor\text{-}rem}(d, (m + 1), a, b)) \cdot \operatorname{factorial}(m)) = ((d((m + 1))(i) \cdot \operatorname{power}((b - i), m)) \cdot (g(b) - g(a))) \wedge \\ \quad \quad \quad u \cdot \operatorname{taylor\text{-}rem}(d, (m + 1), a, b) \\ \quad \quad \quad =\; ((d((m + 1))(i) \cdot \operatorname{recip}(\operatorname{factorial}(m))) \cdot \operatorname{power}((b - i), m)) \cdot (g(b) - g(a)) \end{array}}}
\vnbprop{taylor-lagrange-on-interval}{\vnbstmt{\begin{array}{@{}l@{}} \forall a, b \in \mathbb{R}, d.; \forall m \in \mathbb{N}. \\ \quad \operatorname{is\text{-}taylor\text{-}family}(d, (m + 1), a, b) \Rightarrow  \\ \quad \quad \exists i \in \operatorname{ooint}(a, b). \\ \quad \quad \quad \operatorname{taylor\text{-}rem}(d, (m + 1), a, b) = ((d((m + 1))(i) \cdot \operatorname{recip}(\operatorname{factorial}((m + 1)))) \cdot \operatorname{power}((b - a), (m + 1))) \end{array}}}

\vnbbookfile{theorem-library/zero-deriv-off-countable}{theorem-library/zero-deriv-off-countable}
\noindent{\small zero-deriv-off-countable.scm -- CONSTANCY AND THE MEAN VALUE INEQUALITY WITH A COUNTABLE EXCEPTIONAL SET (batch 20-A, 2026-09-23).}\par
\vnbprop{pw-lt-ne}{\vnbstmt{\forall u, v \in \mathbb{R}.\; ((u < v) \Rightarrow (\neg (u = v) \wedge \neg (v = u)))}}
\vnbprop{nn-antitone-step-implies-le}{\vnbstmt{\begin{array}{@{}l@{}} \forall f \in (\mathbb{N} \to \mathbb{R}), n, m \in \mathbb{N}. \\ \quad (\forall k \in \mathbb{N}.\; (f((k + 1)) \leq f(k)) \wedge (m \leq n)) \Rightarrow (f(n) \leq f(m)) \end{array}}}
\vnbprop{rr-nested-intervals}{\vnbstmt{\begin{array}{@{}l@{}} \forall o, i \in (\mathbb{N} \to \mathbb{R}). \\ \quad \forall k \in \mathbb{N}.\; (o(k) \leq o((k + 1))) \wedge \\ \quad \forall k \in \mathbb{N}.\; (i((k + 1)) \leq i(k)) \wedge \\ \quad \forall k \in \mathbb{N}.\; (o(k) \leq i(k)) \Rightarrow \\ \quad \quad \exists y \in \mathbb{R}.\; \forall n \in \mathbb{N}.\; ((o(n) \leq y) \wedge (y \leq i(n))) \end{array}}}
\vnbprop{rr-interval-avoids-sequence}{\vnbstmt{\begin{array}{@{}l@{}} \forall e \in (\mathbb{N} \to \mathbb{R}), p, q \in \mathbb{R}. \\ \quad (p < q) \Rightarrow  \\ \quad \quad \exists y \in \mathbb{R}.\; ((p < y) \wedge ((y < q) \wedge \forall n \in \mathbb{N}.\; \neg (e(n) = y))) \end{array}}}
\vnbprop{countable-image-avoids-interval}{\vnbstmt{\begin{array}{@{}l@{}} \forall c, m, h.; \forall p, q \in \mathbb{R}. \\ \quad (\operatorname{is\text{-}countable}(c) \wedge ((c \subseteq m) \wedge ((h \in (m \to \mathbb{R})) \wedge (p < q)))) \Rightarrow  \\ \quad \quad \exists y \in \mathbb{R}.\; ((p < y) \wedge ((y < q) \wedge \forall x \in c.\; \neg (h(x) = y))) \end{array}}}
\vnbprop{neg-deriv-off-countable-nonincreasing}{\vnbstmt{\begin{array}{@{}l@{}} \forall h \in (\mathbb{R} \to \mathbb{R}), u, v \in \mathbb{R}, c. \\ \quad u < v \wedge \\ \quad \forall x \in \operatorname{ccint}(u, v).\; \operatorname{is\text{-}continuous\text{-}at}(\operatorname{rr\text{-}ms}, \operatorname{rr\text{-}ms}, h, x) \wedge \\ \quad \operatorname{is\text{-}countable}(c) \wedge \\ \quad c \subseteq \mathbb{R} \wedge \\ \quad \forall t \in \operatorname{ooint}(u, v). \\ \quad \quad \neg (t \in c) \Rightarrow \exists l.\; (\operatorname{is\text{-}diff\text{-}at}(h, t, l) \wedge (l < 0)) \Rightarrow \\ \quad \quad h(v) \leq h(u) \end{array}}}
\vnbprop{mvi-off-countable}{\vnbstmt{\begin{array}{@{}l@{}} \forall h \in (\mathbb{R} \to \mathbb{R}), m, u, v \in \mathbb{R}, c. \\ \quad u < v \wedge \\ \quad \forall x \in \operatorname{ccint}(u, v).\; \operatorname{is\text{-}continuous\text{-}at}(\operatorname{rr\text{-}ms}, \operatorname{rr\text{-}ms}, h, x) \wedge \\ \quad \operatorname{is\text{-}countable}(c) \wedge \\ \quad c \subseteq \mathbb{R} \wedge \\ \quad \forall t \in \operatorname{ooint}(u, v). \\ \quad \quad \neg (t \in c) \Rightarrow \exists l.\; (\operatorname{is\text{-}diff\text{-}at}(h, t, l) \wedge (l \leq m)) \Rightarrow \\ \quad \quad (h(v) - h(u)) \leq (m \cdot (v - u)) \end{array}}}
\vnbprop{mvi-abs-off-countable}{\vnbstmt{\begin{array}{@{}l@{}} \forall h \in (\mathbb{R} \to \mathbb{R}), m, u, v \in \mathbb{R}, c. \\ \quad u < v \wedge \\ \quad \forall x \in \operatorname{ccint}(u, v).\; \operatorname{is\text{-}continuous\text{-}at}(\operatorname{rr\text{-}ms}, \operatorname{rr\text{-}ms}, h, x) \wedge \\ \quad \operatorname{is\text{-}countable}(c) \wedge \\ \quad c \subseteq \mathbb{R} \wedge \\ \quad \forall t \in \operatorname{ooint}(u, v). \\ \quad \quad \neg (t \in c) \Rightarrow \exists l.\; (\operatorname{is\text{-}diff\text{-}at}(h, t, l) \wedge (\lvert l \rvert \leq m)) \Rightarrow \\ \quad \quad \lvert (h(v) - h(u)) \rvert \leq (m \cdot (v - u)) \end{array}}}
\vnbprop{zero-deriv-off-countable}{\vnbstmt{\begin{array}{@{}l@{}} \forall s.; \forall h \in (\mathbb{R} \to \mathbb{R}), u, v \in \mathbb{R}. \\ \quad \operatorname{is\text{-}countable}(s) \wedge \\ \quad s \subseteq \mathbb{R} \wedge \\ \quad \forall x \in \mathbb{R}.\; \operatorname{is\text{-}continuous\text{-}at}(\operatorname{rr\text{-}ms}, \operatorname{rr\text{-}ms}, h, x) \wedge \\ \quad u < v \wedge \\ \quad \forall t \in \operatorname{ooint}(u, v).\; (\neg (t \in s) \Rightarrow \operatorname{is\text{-}diff\text{-}at}(h, t, 0)) \Rightarrow \\ \quad \quad h(u) = h(v) \end{array}}}
\vnbprop{mvi-abs-off-countable-on-interval}{\vnbstmt{\begin{array}{@{}l@{}} \forall a, b, m \in \mathbb{R}, f \in (\operatorname{ccint}(a, b) \to \mathbb{R}), d. \\ \quad a < b \wedge \\ \quad \operatorname{is\text{-}continuous\text{-}on}(f, \operatorname{ccint}(a, b)) \wedge \\ \quad \operatorname{is\text{-}countable}(d) \wedge \\ \quad d \subseteq \operatorname{ccint}(a, b) \wedge \\ \quad \forall t \in \operatorname{ooint}(a, b). \\ \quad \quad \neg (t \in d) \Rightarrow \exists l.\; (\operatorname{has\text{-}deriv\text{-}at}(f, t, l) \wedge (\lvert l \rvert \leq m)) \Rightarrow \\ \quad \quad \lvert (f(b) - f(a)) \rvert \leq (m \cdot (b - a)) \end{array}}}
\vnbprop{zero-deriv-off-countable-constant}{\vnbstmt{\begin{array}{@{}l@{}} \forall a, b \in \mathbb{R}, g \in (\operatorname{ccint}(a, b) \to \mathbb{R}), d.; \forall s, w \in \operatorname{ccint}(a, b). \\ \quad a < b \wedge \\ \quad \operatorname{is\text{-}continuous\text{-}on}(g, \operatorname{ccint}(a, b)) \wedge \\ \quad \operatorname{is\text{-}countable}(d) \wedge \\ \quad d \subseteq \operatorname{ccint}(a, b) \wedge \\ \quad \forall t \in \operatorname{ooint}(a, b).\; (\neg (t \in d) \Rightarrow \operatorname{has\text{-}deriv\text{-}at}(g, t, 0)) \Rightarrow \\ \quad \quad g(s) = g(w) \end{array}}}

\vnbbookfile{theorem-library/rabasics-monotone}{theorem-library/rabasics-monotone}
\noindent{\small rabasics-monotone.scm -- calculus.pdf 3.7, complex-analysis.pdf 1.7: a monotone bounded real sequence converges to its sup (resp. inf).}\par
\vnbprop{monotone-converges-to-sup}{\vnbstmt{\begin{array}{@{}l@{}} \forall f \in (\mathbb{N} \to \mathbb{R}). \\ \quad (\forall k \in \mathbb{N}.\; (f(k) \leq f((k + 1))) \wedge \exists d \in \mathbb{R}.\; \forall k \in \mathbb{N}.\; (f(k) \leq d)) \Rightarrow  \\ \quad \quad \operatorname{converges\text{-}to}(\operatorname{rr\text{-}ms}, f, \operatorname{sup}(\{y \in \mathbb{R} : \exists a \in \mathbb{N}.\; (y = f(a))\})) \end{array}}}
\vnbprop{monotone-converges-to-inf}{\vnbstmt{\begin{array}{@{}l@{}} \forall f \in (\mathbb{N} \to \mathbb{R}). \\ \quad (\forall k \in \mathbb{N}.\; (f((k + 1)) \leq f(k)) \wedge \exists d \in \mathbb{R}.\; \forall k \in \mathbb{N}.\; (d \leq f(k))) \Rightarrow  \\ \quad \quad \operatorname{converges\text{-}to}(\operatorname{rr\text{-}ms}, f, -\operatorname{sup}(\{y \in \mathbb{R} : \exists a \in \mathbb{N}.\; (y = -f(a))\})) \end{array}}}
\vnbprop{rr-neg-sup-neg-is-glb}{\vnbstmt{\begin{array}{@{}l@{}} \forall f \in (\mathbb{N} \to \mathbb{R}). \\ \quad \exists d \in \mathbb{R}.\; \forall k \in \mathbb{N}.\; (d \leq f(k)) \Rightarrow  \\ \quad \quad \forall k \in \mathbb{N}.\; (-\operatorname{sup}(\{y \in \mathbb{R} : \exists a \in \mathbb{N}.\; (y = -f(a))\}) \leq f(k)) \wedge \\ \quad \quad \forall b \in \mathbb{R}. \\ \quad \quad \quad \forall k \in \mathbb{N}.\; (b \leq f(k)) \Rightarrow  \\ \quad \quad \quad \quad b \leq -\operatorname{sup}(\{y \in \mathbb{R} : \exists a \in \mathbb{N}.\; (y = -f(a))\}) \end{array}}}
\vnbprop{rr-inf-is-glb}{\vnbstmt{\begin{array}{@{}l@{}} \forall s. \\ \quad ((s \subseteq \mathbb{R}) \wedge (\exists a.\; (a \in s) \wedge \exists d \in \mathbb{R}.\; \forall x \in s.\; (d \leq x))) \Rightarrow  \\ \quad \quad \operatorname{inf}(s) \in \mathbb{R} \wedge \\ \quad \quad \forall x \in s.\; (\operatorname{inf}(s) \leq x) \wedge \\ \quad \quad \forall b \in \mathbb{R}.\; (\forall x \in s.\; (b \leq x) \Rightarrow (b \leq \operatorname{inf}(s))) \end{array}}}
\vnbprop{monotone-converges-to-inf-literal}{\vnbstmt{\begin{array}{@{}l@{}} \forall f \in (\mathbb{N} \to \mathbb{R}). \\ \quad (\forall k \in \mathbb{N}.\; (f((k + 1)) \leq f(k)) \wedge \exists d \in \mathbb{R}.\; \forall k \in \mathbb{N}.\; (d \leq f(k))) \Rightarrow  \\ \quad \quad \operatorname{converges\text{-}to}(\operatorname{rr\text{-}ms}, f, \operatorname{inf}(\{y \in \mathbb{R} : \exists a \in \mathbb{N}.\; (y = f(a))\})) \end{array}}}

\vnbbookfile{theorem-library/pw-antiderivative-laws}{theorem-library/pw-antiderivative-laws}
\noindent{\small pw-antiderivative-laws.scm -- the laws of IS-PW-ANTIDERIVATIVE and PW-INT (structure-library/path-integral.scm), with the exceptional set a FINITE SET.}\par
\vnbprop{pw-antiderivative-value-in-rr}{\vnbstmt{\begin{array}{@{}l@{}} \forall f, i, a, b.; \forall y \in \operatorname{ccint}(a, b). \\ \quad \operatorname{is\text{-}primitive}(f, i, a, b) \Rightarrow (f(y) \in \mathbb{R}) \end{array}}}
\vnbprop{pw-not-in-insert}{\vnbstmt{\begin{array}{@{}l@{}} \forall y \in \mathbf{Set}, w, t. \\ \quad (\neg (t \in w) \wedge \neg (t = y)) \Rightarrow \neg (t \in (w \cup \{y\})) \end{array}}}
\vnbprop{ooint-inner-radius}{\vnbstmt{\begin{array}{@{}l@{}} \forall u, v \in \mathbb{R}, t \in \operatorname{ooint}(u, v). \exists r. \\ \quad \operatorname{pos\text{-}rr}(r) \wedge (\operatorname{ooint}((t - r), (t + r)) \subseteq \operatorname{ccint}(u, v)) \end{array}}}
\vnbprop{pw-zero-deriv-total}{\vnbstmt{\begin{array}{@{}l@{}} \forall h \in (\mathbb{R} \to \mathbb{R}), u, v \in \mathbb{R}. \\ \quad \forall x \in \mathbb{R}.\; \operatorname{is\text{-}continuous\text{-}at}(\operatorname{rr\text{-}ms}, \operatorname{rr\text{-}ms}, h, x) \wedge \\ \quad u < v \wedge \\ \quad \forall t \in \operatorname{ooint}(u, v).\; \operatorname{is\text{-}diff\text{-}at}(h, t, 0) \Rightarrow \\ \quad \quad h(u) = h(v) \end{array}}}
\vnbprop{deriv-zero-constant-on-interval}{\vnbstmt{\begin{array}{@{}l@{}} \forall a, b \in \mathbb{R}, h \in (\operatorname{ccint}(a, b) \to \mathbb{R}), u, v \in \operatorname{ccint}(a, b). \\ \quad ((a < b) \wedge (\operatorname{is\text{-}continuous\text{-}on}(h, \operatorname{ccint}(a, b)) \wedge \forall t \in \operatorname{ooint}(a, b).\; \operatorname{has\text{-}deriv\text{-}at}(h, t, 0))) \Rightarrow  \\ \quad \quad h(u) = h(v) \end{array}}}
\vnbprop{pw-zero-deriv-off-finite-set}{\vnbstmt{\begin{array}{@{}l@{}} \forall s \in \mathbf{Set}, h \in (\mathbb{R} \to \mathbb{R}), u, v \in \mathbb{R}. \\ \quad |s| \in \mathbb{N} \wedge \\ \quad \forall x \in \mathbb{R}.\; \operatorname{is\text{-}continuous\text{-}at}(\operatorname{rr\text{-}ms}, \operatorname{rr\text{-}ms}, h, x) \wedge \\ \quad u < v \wedge \\ \quad \forall t \in \operatorname{ooint}(u, v).\; (\neg (t \in s) \Rightarrow \operatorname{is\text{-}diff\text{-}at}(h, t, 0)) \Rightarrow \\ \quad \quad h(u) = h(v) \end{array}}}
\vnbprop{pw-antiderivative-endpoint-difference}{\vnbstmt{\begin{array}{@{}l@{}} \forall f, w, i, a, b. \\ \quad (\operatorname{is\text{-}primitive}(f, i, a, b) \wedge \operatorname{is\text{-}primitive}(w, i, a, b)) \Rightarrow  \\ \quad \quad (f(b) - f(a)) = (w(b) - w(a)) \end{array}}}
\vnbprop{pw-int-value}{\vnbstmt{\begin{array}{@{}l@{}} \forall f, i, a, b. \\ \quad \operatorname{is\text{-}primitive}(f, i, a, b) \Rightarrow (\int_{a}^{b} i = (f(b) - f(a))) \end{array}}}
\vnbprop{pw-int-in-rr}{\vnbstmt{\forall f, i, a, b.\; (\operatorname{is\text{-}primitive}(f, i, a, b) \Rightarrow (\int_{a}^{b} i \in \mathbb{R}))}}
\vnbprop{pw-antiderivative-integrand-congruence}{\vnbstmt{\begin{array}{@{}l@{}} \forall f, i, \mathit{ppsi}, a, b. \\ \quad \operatorname{is\text{-}primitive}(f, i, a, b) \wedge \\ \quad \mathit{ppsi} \in (\operatorname{ccint}(a, b) \to \mathbb{R}) \wedge \\ \quad \forall y \in \operatorname{ccint}(a, b).\; (i(y) \simeq \operatorname{ppsi}(y)) \Rightarrow \\ \quad \quad \operatorname{is\text{-}primitive}(f, \mathit{ppsi}, a, b) \end{array}}}
\vnbprop{pw-int-congruence}{\vnbstmt{\begin{array}{@{}l@{}} \forall f, i, \mathit{ppsi}, a, b. \\ \quad \operatorname{is\text{-}primitive}(f, i, a, b) \wedge \\ \quad \mathit{ppsi} \in (\operatorname{ccint}(a, b) \to \mathbb{R}) \wedge \\ \quad \forall y \in \operatorname{ccint}(a, b).\; (i(y) \simeq \operatorname{ppsi}(y)) \Rightarrow \\ \quad \quad \int_{a}^{b} i = \int_{a}^{b} \mathit{ppsi} \end{array}}}
\vnbprop{pw-antiderivative-one-piece}{\vnbstmt{\begin{array}{@{}l@{}} \forall a, b \in \mathbb{R}, f, i. \\ \quad a < b \wedge \\ \quad f \in (\operatorname{ccint}(a, b) \to \mathbb{R}) \wedge \\ \quad i \in (\operatorname{ccint}(a, b) \to \mathbb{R}) \wedge \\ \quad \operatorname{is\text{-}continuous\text{-}on}(f, \operatorname{ccint}(a, b)) \wedge \\ \quad \forall t \in \operatorname{ooint}(a, b).\; \operatorname{has\text{-}deriv\text{-}at}(f, t, i(t)) \Rightarrow \\ \quad \quad \operatorname{is\text{-}primitive}(f, i, a, b) \end{array}}}
\vnbprop{pw-restrict-continuous-on}{\vnbstmt{\begin{array}{@{}l@{}} \forall a, b \in \mathbb{R}, f, \mathit{pwf}. \\ \quad a < b \wedge \\ \quad f \in (\mathbb{R} \to \mathbb{R}) \wedge \\ \quad \forall y \in \mathbb{R}.\; \operatorname{is\text{-}continuous\text{-}at}(\operatorname{rr\text{-}ms}, \operatorname{rr\text{-}ms}, f, y) \wedge \\ \quad \mathit{pwf} \in (\operatorname{ccint}(a, b) \to \mathbb{R}) \wedge \\ \quad \forall y \in \operatorname{ccint}(a, b).\; (\operatorname{pwf}(y) \simeq f(y)) \Rightarrow \\ \quad \quad \operatorname{is\text{-}continuous\text{-}on}(\mathit{pwf}, \operatorname{ccint}(a, b)) \end{array}}}
\vnbprop{pw-affine-antiderivative}{\vnbstmt{\begin{array}{@{}l@{}} \forall a, b, c, m \in \mathbb{R}, f, i. \\ \quad a < b \wedge \\ \quad f \in (\operatorname{ccint}(a, b) \to \mathbb{R}) \wedge \\ \quad i \in (\operatorname{ccint}(a, b) \to \mathbb{R}) \wedge \\ \quad \forall y \in \operatorname{ccint}(a, b).\; (f(y) \simeq (c + (m \cdot y))) \wedge \\ \quad \forall y \in \operatorname{ccint}(a, b).\; (i(y) \simeq m) \Rightarrow \\ \quad \quad \operatorname{is\text{-}primitive}(f, i, a, b) \end{array}}}
\vnbprop{pw-int-const}{\vnbstmt{\begin{array}{@{}l@{}} \forall a, b, c \in \mathbb{R}, i \in (\operatorname{ccint}(a, b) \to \mathbb{R}). \\ \quad ((a < b) \wedge \forall y \in \operatorname{ccint}(a, b).\; (i(y) \simeq c)) \Rightarrow  \\ \quad \quad \int_{a}^{b} i = (c \cdot (b - a)) \end{array}}}
\vnbprop{pw-antiderivative-sum}{\vnbstmt{\begin{array}{@{}l@{}} \forall a, b \in \mathbb{R}, f, w, i, \mathit{ppsi}.; \forall h, \mathit{pchi} \in (\operatorname{ccint}(a, b) \to \mathbb{R}). \\ \quad \operatorname{is\text{-}primitive}(f, i, a, b) \wedge \\ \quad \operatorname{is\text{-}primitive}(w, \mathit{ppsi}, a, b) \wedge \\ \quad \forall y \in \operatorname{ccint}(a, b).\; (h(y) \simeq (f(y) + w(y))) \wedge \\ \quad \forall y \in \operatorname{ccint}(a, b).\; (\operatorname{pchi}(y) \simeq (i(y) + \operatorname{ppsi}(y))) \Rightarrow \\ \quad \quad (\operatorname{is\text{-}primitive}(h, \mathit{pchi}, a, b) \wedge (\int_{a}^{b} \mathit{pchi} = (\int_{a}^{b} i + \int_{a}^{b} \mathit{ppsi}))) \end{array}}}

\vnbbookfile{theorem-library/cc-int-laws}{theorem-library/cc-int-laws}
\noindent{\small cc-int-laws.scm -- the read-offs of CC-INT, TRACE and LINE-INT (structure-library/path-integral.scm), sections 3.3 and 3.4 of docs/paths-and-line-integrals-2026-09-21.md.}\par
\vnbprop{cc-int-unfold}{\vnbstmt{\begin{array}{@{}l@{}} \forall i, a, b. \\ \quad (\int_{a}^{b} i \simeq (\int_{a}^{b} (\lambda t \in \operatorname{ccint}(a, b).\; \operatorname{real\text{-}part}(i(t))) + (\int_{a}^{b} (\lambda t \in \operatorname{ccint}(a, b).\; \operatorname{imag\text{-}part}(i(t))) \cdot +i))) \end{array}}}
\vnbprop{line-int-unfold}{\vnbstmt{\begin{array}{@{}l@{}} \forall f, m, c, a, b. \\ \quad (\operatorname{line\text{-}int}(f, m, c, a, b) \simeq \int_{a}^{b} (\lambda t \in \operatorname{ccint}(a, b).\; (f(m(t)) \cdot c(t)))) \end{array}}}
\vnbprop{trace-unfold}{\vnbstmt{\forall m, a, b.\; (\operatorname{trace}(m, a, b) \simeq \operatorname{image}(m, \operatorname{ccint}(a, b)))}}
\vnbprop{trace-value-in}{\vnbstmt{\begin{array}{@{}l@{}} \forall m, a, b.; \forall t \in \operatorname{ccint}(a, b). \\ \quad (m \in (\operatorname{ccint}(a, b) \to \mathbb{C})) \Rightarrow (m(t) \in \operatorname{trace}(m, a, b)) \end{array}}}
\vnbprop{trace-subset-cc}{\vnbstmt{\forall m, a, b.\; ((m \in (\operatorname{ccint}(a, b) \to \mathbb{C})) \Rightarrow (\operatorname{trace}(m, a, b) \subseteq \mathbb{C}))}}
\vnbprop{trace-is-set}{\vnbstmt{\forall m, a, b.\; (((a \in \mathbb{R}) \wedge (b \in \mathbb{R})) \Rightarrow (\operatorname{trace}(m, a, b) \in \mathbf{Set}))}}
\vnbprop{line-int-integrand-in-fun}{\vnbstmt{\begin{array}{@{}l@{}} \forall a, b, m, c, f, d. \\ \quad a \in \mathbb{R} \wedge \\ \quad b \in \mathbb{R} \wedge \\ \quad m \in (\operatorname{ccint}(a, b) \to \mathbb{C}) \wedge \\ \quad c \in (\operatorname{ccint}(a, b) \to \mathbb{C}) \wedge \\ \quad f \in (d \to \mathbb{C}) \wedge \\ \quad \operatorname{trace}(m, a, b) \subseteq d \Rightarrow \\ \quad \quad (\lambda t \in \operatorname{ccint}(a, b).\; (f(m(t)) \cdot c(t))) \in (\operatorname{ccint}(a, b) \to \mathbb{C}) \end{array}}}

\vnbbookfile{theorem-library/road-laws}{theorem-library/road-laws}
\noindent{\small road-laws.scm -- the laws of IS-PW-CONTINUOUS-ON, IS-PATH and IS-ROAD (structure-library/path-integral.scm), section 3.4 of docs/paths-and-line-integrals-2026-09-21.md and Dieudonne 9.6.}\par
\vnbprop{pw-continuous-endpoints}{\vnbstmt{\begin{array}{@{}l@{}} \forall i, a, b. \\ \quad \operatorname{is\text{-}pw\text{-}continuous\text{-}on}(i, a, b) \Rightarrow ((a \in \mathbb{R}) \wedge ((b \in \mathbb{R}) \wedge (a < b))) \end{array}}}
\vnbprop{pw-continuous-in-fun}{\vnbstmt{\begin{array}{@{}l@{}} \forall i, a, b. \\ \quad \operatorname{is\text{-}pw\text{-}continuous\text{-}on}(i, a, b) \Rightarrow (i \in (\operatorname{ccint}(a, b) \to \mathbb{R})) \end{array}}}
\vnbprop{is-path-endpoints}{\vnbstmt{\forall m, a, b.\; (\operatorname{is\text{-}path}(m, a, b) \Rightarrow ((a \in \mathbb{R}) \wedge ((b \in \mathbb{R}) \wedge (a < b))))}}
\vnbprop{is-path-in-fun}{\vnbstmt{\forall m, a, b.\; (\operatorname{is\text{-}path}(m, a, b) \Rightarrow (m \in (\operatorname{ccint}(a, b) \to \mathbb{C})))}}
\vnbprop{is-path-continuous}{\vnbstmt{\begin{array}{@{}l@{}} \forall m, a, b.; \forall x \in \operatorname{ccint}(a, b). \\ \quad \operatorname{is\text{-}path}(m, a, b) \Rightarrow  \\ \quad \quad \operatorname{is\text{-}continuous\text{-}at}(\operatorname{subspace\text{-}ms}(\operatorname{rr\text{-}ms}, \operatorname{ccint}(a, b)), \operatorname{cc\text{-}ms}, m, x) \end{array}}}
\vnbprop{is-road-is-path}{\vnbstmt{\forall m, c, a, b.\; (\operatorname{is\text{-}road}(m, c, a, b) \Rightarrow \operatorname{is\text{-}path}(m, a, b))}}
\vnbprop{is-road-dgam-in-fun}{\vnbstmt{\forall m, c, a, b.\; (\operatorname{is\text{-}road}(m, c, a, b) \Rightarrow (c \in (\operatorname{ccint}(a, b) \to \mathbb{C})))}}
\vnbprop{is-road-re-antiderivative}{\vnbstmt{\begin{array}{@{}l@{}} \forall m, c, a, b. \\ \quad \operatorname{is\text{-}road}(m, c, a, b) \Rightarrow  \\ \quad \quad \operatorname{is\text{-}primitive}((\lambda t \in \operatorname{ccint}(a, b).\; \operatorname{real\text{-}part}(m(t))), \\ \quad \quad \quad (\lambda t \in \operatorname{ccint}(a, b).\; \operatorname{real\text{-}part}(c(t))), a, b) \end{array}}}
\vnbprop{is-road-im-antiderivative}{\vnbstmt{\begin{array}{@{}l@{}} \forall m, c, a, b. \\ \quad \operatorname{is\text{-}road}(m, c, a, b) \Rightarrow  \\ \quad \quad \operatorname{is\text{-}primitive}((\lambda t \in \operatorname{ccint}(a, b).\; \operatorname{imag\text{-}part}(m(t))), \\ \quad \quad \quad (\lambda t \in \operatorname{ccint}(a, b).\; \operatorname{imag\text{-}part}(c(t))), a, b) \end{array}}}
\vnbprop{is-road-re-regulated}{\vnbstmt{\begin{array}{@{}l@{}} \forall m, c, a, b. \\ \quad \operatorname{is\text{-}road}(m, c, a, b) \Rightarrow  \\ \quad \quad \operatorname{is\text{-}regulated\text{-}on}((\lambda t \in \operatorname{ccint}(a, b).\; \operatorname{real\text{-}part}(c(t))), a, b) \end{array}}}
\vnbprop{is-road-im-regulated}{\vnbstmt{\begin{array}{@{}l@{}} \forall m, c, a, b. \\ \quad \operatorname{is\text{-}road}(m, c, a, b) \Rightarrow  \\ \quad \quad \operatorname{is\text{-}regulated\text{-}on}((\lambda t \in \operatorname{ccint}(a, b).\; \operatorname{imag\text{-}part}(c(t))), a, b) \end{array}}}
\vnbprop{pw-continuous-one-piece}{\vnbstmt{\begin{array}{@{}l@{}} \forall a, b \in \mathbb{R}, i \in (\operatorname{ccint}(a, b) \to \mathbb{R}). \\ \quad ((a < b) \wedge \operatorname{is\text{-}continuous\text{-}on}(i, \operatorname{ccint}(a, b))) \Rightarrow \operatorname{is\text{-}pw\text{-}continuous\text{-}on}(i, a, b) \end{array}}}
\vnbprop{is-path-of-coords}{\vnbstmt{\begin{array}{@{}l@{}} \forall a, b \in \mathbb{R}, m, r, i. \\ \quad a < b \wedge \\ \quad m \in (\operatorname{ccint}(a, b) \to \mathbb{C}) \wedge \\ \quad r \in (\operatorname{ccint}(a, b) \to \mathbb{R}) \wedge \\ \quad i \in (\operatorname{ccint}(a, b) \to \mathbb{R}) \wedge \\ \quad \forall y \in \operatorname{ccint}(a, b).\; (\operatorname{real\text{-}part}(m(y)) = r(y)) \wedge \\ \quad \forall y \in \operatorname{ccint}(a, b).\; (\operatorname{imag\text{-}part}(m(y)) = i(y)) \wedge \\ \quad \operatorname{is\text{-}continuous\text{-}on}(r, \operatorname{ccint}(a, b)) \wedge \\ \quad \operatorname{is\text{-}continuous\text{-}on}(i, \operatorname{ccint}(a, b)) \Rightarrow \\ \quad \quad \operatorname{is\text{-}path}(m, a, b) \end{array}}}
\vnbprop{segment-is-road}{\vnbstmt{\begin{array}{@{}l@{}} \forall z_{0}, w. \\ \quad ((z_{0} \in \mathbb{C}) \wedge (w \in \mathbb{C})) \Rightarrow  \\ \quad \quad \operatorname{is\text{-}road}((\lambda x \in \operatorname{ccint}(0, 1).\; (z_{0} + (x \cdot w))), (\lambda x \in \operatorname{ccint}(0, 1).\; w), 0, 1) \end{array}}}
\vnbprop{line-int-of-one}{\vnbstmt{\begin{array}{@{}l@{}} \forall m, c, a, b, f, d. \\ \quad (\operatorname{is\text{-}road}(m, c, a, b) \wedge ((f \in (d \to \mathbb{C})) \wedge ((\operatorname{trace}(m, a, b) \subseteq d) \wedge \forall y \in d.\; (f(y) = 1)))) \Rightarrow  \\ \quad \quad \operatorname{line\text{-}int}(f, m, c, a, b) = (m(b) - m(a)) \end{array}}}

\vnbbookfile{theorem-library/pw-int-laws-2}{theorem-library/pw-int-laws-2}
\noindent{\small pw-int-laws-2.scm -- the LINEARITY and the ADDITIVITY of the integral on [a, b], continuing theorem-library/pw-antiderivative-laws.scm.}\par
\vnbprop{pw-antiderivative-real-mul}{\vnbstmt{\begin{array}{@{}l@{}} \forall a, b \in \mathbb{R}, f, i.; \forall c \in \mathbb{R}, h, \mathit{pchi} \in (\operatorname{ccint}(a, b) \to \mathbb{R}). \\ \quad \operatorname{is\text{-}primitive}(f, i, a, b) \wedge \\ \quad \forall y \in \operatorname{ccint}(a, b).\; (h(y) \simeq (c \cdot f(y))) \wedge \\ \quad \forall y \in \operatorname{ccint}(a, b).\; (\operatorname{pchi}(y) \simeq (c \cdot i(y))) \Rightarrow \\ \quad \quad \operatorname{is\text{-}primitive}(h, \mathit{pchi}, a, b) \wedge (\int_{a}^{b} \mathit{pchi} = (c \cdot \int_{a}^{b} i)) \end{array}}}
\vnbprop{pw-antiderivative-differ-by-constant}{\vnbstmt{\begin{array}{@{}l@{}} \forall f, w, i, a, b.; \forall s \in \operatorname{ccint}(a, b). \\ \quad (\operatorname{is\text{-}primitive}(f, i, a, b) \wedge \operatorname{is\text{-}primitive}(w, i, a, b)) \Rightarrow  \\ \quad \quad f(s) = (w(s) + (f(a) - w(a))) \end{array}}}
\vnbprop{pw-antiderivative-differ-by-some-constant}{\vnbstmt{\begin{array}{@{}l@{}} \forall f, w, i, a, b. \\ \quad (\operatorname{is\text{-}primitive}(f, i, a, b) \wedge \operatorname{is\text{-}primitive}(w, i, a, b)) \Rightarrow  \\ \quad \quad \exists c \in \mathbb{R}.\; \forall s \in \operatorname{ccint}(a, b).\; (f(s) = (w(s) + c)) \end{array}}}
\vnbprop{ccint-subset-ccint}{\vnbstmt{\begin{array}{@{}l@{}} \forall a, b, v, \mathit{pdv} \in \mathbb{R}. \\ \quad ((a \leq v) \wedge (\mathit{pdv} \leq b)) \Rightarrow (\operatorname{ccint}(v, \mathit{pdv}) \subseteq \operatorname{ccint}(a, b)) \end{array}}}
\vnbprop{pw-antiderivative-restrict}{\vnbstmt{\begin{array}{@{}l@{}} \forall a, b, v, \mathit{pdv} \in \mathbb{R}, f, i. \\ \quad (\operatorname{is\text{-}primitive}(f, i, a, b) \wedge ((a \leq v) \wedge ((v < \mathit{pdv}) \wedge (\mathit{pdv} \leq b)))) \Rightarrow  \\ \quad \quad \operatorname{is\text{-}primitive}(\operatorname{restrict}(f, \operatorname{ccint}(v, \mathit{pdv})), \operatorname{restrict}(i, \operatorname{ccint}(v, \mathit{pdv})), v, \mathit{pdv}) \end{array}}}
\vnbprop{pw-int-adjacent}{\vnbstmt{\begin{array}{@{}l@{}} \forall a, b \in \mathbb{R}, v \in \operatorname{ooint}(a, b), f, i. \\ \quad \operatorname{is\text{-}primitive}(f, i, a, b) \Rightarrow  \\ \quad \quad \int_{a}^{b} i = (\int_{a}^{v} \operatorname{restrict}(i, \operatorname{ccint}(a, v)) + \int_{v}^{b} \operatorname{restrict}(i, \operatorname{ccint}(v, b))) \end{array}}}
\vnbprop{cc-int-value}{\vnbstmt{\begin{array}{@{}l@{}} \forall a, b, i, f, w. \\ \quad a \in \mathbb{R} \wedge \\ \quad b \in \mathbb{R} \wedge \\ \quad \operatorname{is\text{-}primitive}(f, (\lambda t \in \operatorname{ccint}(a, b).\; \operatorname{real\text{-}part}(i(t))), a, b) \wedge \\ \quad \operatorname{is\text{-}primitive}(w, (\lambda t \in \operatorname{ccint}(a, b).\; \operatorname{imag\text{-}part}(i(t))), a, b) \Rightarrow \\ \quad \quad \int_{a}^{b} i = ((f(b) - f(a)) + ((w(b) - w(a)) \cdot +i)) \end{array}}}
\vnbprop{cc-int-in-cc}{\vnbstmt{\begin{array}{@{}l@{}} \forall a, b, i, f, w. \\ \quad a \in \mathbb{R} \wedge \\ \quad b \in \mathbb{R} \wedge \\ \quad \operatorname{is\text{-}primitive}(f, (\lambda t \in \operatorname{ccint}(a, b).\; \operatorname{real\text{-}part}(i(t))), a, b) \wedge \\ \quad \operatorname{is\text{-}primitive}(w, (\lambda t \in \operatorname{ccint}(a, b).\; \operatorname{imag\text{-}part}(i(t))), a, b) \Rightarrow \\ \quad \quad \int_{a}^{b} i \in \mathbb{C} \end{array}}}
\vnbprop{cc-int-sum}{\vnbstmt{\begin{array}{@{}l@{}} \forall a, b, i, \mathit{ppsi}, \mathit{pchi}, f, w, h, \mathit{pbw}. \\ \quad a \in \mathbb{R} \wedge \\ \quad b \in \mathbb{R} \wedge \\ \quad i \in (\operatorname{ccint}(a, b) \to \mathbb{C}) \wedge \\ \quad \mathit{ppsi} \in (\operatorname{ccint}(a, b) \to \mathbb{C}) \wedge \\ \quad \mathit{pchi} \in (\operatorname{ccint}(a, b) \to \mathbb{C}) \wedge \\ \quad \forall y \in \operatorname{ccint}(a, b).\; (\operatorname{pchi}(y) \simeq (i(y) + \operatorname{ppsi}(y))) \wedge \\ \quad \operatorname{is\text{-}primitive}(f, (\lambda t \in \operatorname{ccint}(a, b).\; \operatorname{real\text{-}part}(i(t))), a, b) \wedge \\ \quad \operatorname{is\text{-}primitive}(w, (\lambda t \in \operatorname{ccint}(a, b).\; \operatorname{imag\text{-}part}(i(t))), a, b) \wedge \\ \quad \operatorname{is\text{-}primitive}(h, (\lambda t \in \operatorname{ccint}(a, b).\; \operatorname{real\text{-}part}(\operatorname{ppsi}(t))), a, b) \wedge \\ \quad \operatorname{is\text{-}primitive}(\mathit{pbw}, (\lambda t \in \operatorname{ccint}(a, b).\; \operatorname{imag\text{-}part}(\operatorname{ppsi}(t))), a, b) \Rightarrow \\ \quad \quad \int_{a}^{b} \mathit{pchi} = (\int_{a}^{b} i + \int_{a}^{b} \mathit{ppsi}) \end{array}}}
\vnbprop{cc-int-real-mul}{\vnbstmt{\begin{array}{@{}l@{}} \forall a, b, i, \mathit{pchi}.; \forall c \in \mathbb{R}, f, w. \\ \quad a \in \mathbb{R} \wedge \\ \quad b \in \mathbb{R} \wedge \\ \quad i \in (\operatorname{ccint}(a, b) \to \mathbb{C}) \wedge \\ \quad \mathit{pchi} \in (\operatorname{ccint}(a, b) \to \mathbb{C}) \wedge \\ \quad \forall y \in \operatorname{ccint}(a, b).\; (\operatorname{pchi}(y) \simeq (c \cdot i(y))) \wedge \\ \quad \operatorname{is\text{-}primitive}(f, (\lambda t \in \operatorname{ccint}(a, b).\; \operatorname{real\text{-}part}(i(t))), a, b) \wedge \\ \quad \operatorname{is\text{-}primitive}(w, (\lambda t \in \operatorname{ccint}(a, b).\; \operatorname{imag\text{-}part}(i(t))), a, b) \Rightarrow \\ \quad \quad \int_{a}^{b} \mathit{pchi} = (c \cdot \int_{a}^{b} i) \end{array}}}

\vnbbookfile{theorem-library/taylor-cluster}{theorem-library/taylor-cluster}
\noindent{\small taylor-cluster.scm -- THE TAYLOR CLUSTER OF THE CALCULUS NOTES (calculus.pdf, ch. 2, sections 1, 6 and 7): 2.3, 2.16, 2.20, 2.21, 2.22.  Batch 25-B, 2026-09-24.}\par
\vnbprop{nth-deriv-succ-inner}{\vnbstmt{\begin{array}{@{}l@{}} \forall k \in \mathbb{N}, f. \\ \quad (\operatorname{nth\text{-}deriv}(f, (k + 1)) \simeq \operatorname{nth\text{-}deriv}((\lambda x \in \mathbb{R}.\; \operatorname{deriv}(f, x)), k)) \end{array}}}
\vnbprop{nth-deriv-monomial-core}{\vnbstmt{\begin{array}{@{}l@{}} \forall k, m \in \mathbb{N}. \\ \quad (\operatorname{nth\text{-}deriv}((\lambda x \in \mathbb{R}.\; \operatorname{power}(x, (m + k))), k) \simeq (\lambda x \in \mathbb{R}.\; ((\operatorname{factorial}((m + k)) \cdot \operatorname{recip}(\operatorname{factorial}(m))) \cdot \operatorname{power}(x, m)))) \end{array}}}
\vnbprop{nth-deriv-monomial-gap}{\vnbstmt{\begin{array}{@{}l@{}} \forall k, m \in \mathbb{N}, p. \\ \quad (p = (m + k)) \Rightarrow  \\ \quad \quad (\operatorname{nth\text{-}deriv}((\lambda x \in \mathbb{R}.\; \operatorname{power}(x, p)), k) \simeq (\lambda x \in \mathbb{R}.\; ((\operatorname{factorial}(p) \cdot \operatorname{recip}(\operatorname{factorial}(m))) \cdot \operatorname{power}(x, m)))) \end{array}}}
\vnbprop{nth-deriv-monomial-beyond}{\vnbstmt{\begin{array}{@{}l@{}} \forall j, n \in \mathbb{N}. \\ \quad (\operatorname{nth\text{-}deriv}((\lambda x \in \mathbb{R}.\; \operatorname{power}(x, n)), ((n + j) + 1)) \simeq (\lambda x \in \mathbb{R}.\; 0)) \end{array}}}
\vnbprop{nth-deriv-monomial-vanish}{\vnbstmt{\begin{array}{@{}l@{}} \forall n, k \in \mathbb{N}. \\ \quad (n < k) \Rightarrow (\operatorname{nth\text{-}deriv}((\lambda x \in \mathbb{R}.\; \operatorname{power}(x, n)), k) \simeq (\lambda x \in \mathbb{R}.\; 0)) \end{array}}}
\vnbprop{nth-deriv-monomial}{\vnbstmt{\begin{array}{@{}l@{}} \forall n, k \in \mathbb{N}. \\ \quad (k \leq n) \Rightarrow  \\ \quad \quad (\operatorname{nth\text{-}deriv}((\lambda x \in \mathbb{R}.\; \operatorname{power}(x, n)), k) \simeq (\lambda x \in \mathbb{R}.\; ((\operatorname{factorial}(n) \cdot \operatorname{recip}(\operatorname{factorial}(\operatorname{nn\text{-}minus}(n, k)))) \cdot \operatorname{power}(x, \operatorname{nn\text{-}minus}(n, k))))) \end{array}}}
\vnbprop{poly-deriv-fn}{\vnbstmt{\begin{array}{@{}l@{}} \forall n \in \mathbb{N}, f, d \in (\mathbb{N} \to \mathbb{R}). \\ \quad \forall j \in \mathbb{N}.\; (d(j) = ((j + 1) \cdot f((j + 1)))) \Rightarrow  \\ \quad \quad ((\lambda \mathit{tcx} \in \mathbb{R}.\; \operatorname{deriv}((\lambda x \in \mathbb{R}.\; \operatorname{series\text{-}partial\text{-}sum}((\lambda k \in \mathbb{N}.\; (f(k) \cdot \operatorname{power}(x, k))), (n + 1))), \mathit{tcx})) \simeq (\lambda x \in \mathbb{R}.\; \operatorname{series\text{-}partial\text{-}sum}((\lambda k \in \mathbb{N}.\; (d(k) \cdot \operatorname{power}(x, k))), n))) \end{array}}}
\vnbprop{poly-deriv-coef-exists}{\vnbstmt{\begin{array}{@{}l@{}} \forall f \in (\mathbb{N} \to \mathbb{R}). \exists d \in (\mathbb{N} \to \mathbb{R}). \forall j \in \mathbb{N}. \\ \quad d(j) = ((j + 1) \cdot f((j + 1))) \end{array}}}
\vnbprop{poly-deriv-fn-zero}{\vnbstmt{\begin{array}{@{}l@{}} \forall f \in (\mathbb{N} \to \mathbb{R}). \\ \quad ((\lambda \mathit{tcx} \in \mathbb{R}.\; \operatorname{deriv}((\lambda x \in \mathbb{R}.\; \operatorname{series\text{-}partial\text{-}sum}((\lambda k \in \mathbb{N}.\; (f(k) \cdot \operatorname{power}(x, k))), 0)), \mathit{tcx})) \simeq (\lambda x \in \mathbb{R}.\; \operatorname{series\text{-}partial\text{-}sum}((\lambda k \in \mathbb{N}.\; (f(k) \cdot \operatorname{power}(x, k))), 0))) \end{array}}}
\vnbprop{poly-nth-deriv-in-fun}{\vnbstmt{\begin{array}{@{}l@{}} \forall \mathit{tck}, n \in \mathbb{N}, f \in (\mathbb{N} \to \mathbb{R}). \\ \quad \operatorname{nth\text{-}deriv}((\lambda x \in \mathbb{R}.\; \operatorname{series\text{-}partial\text{-}sum}((\lambda k \in \mathbb{N}.\; (f(k) \cdot \operatorname{power}(x, k))), n)), \mathit{tck}) \\ \quad \in\; (\mathbb{R} \to \mathbb{R}) \end{array}}}
\vnbprop{poly-value-at-zero}{\vnbstmt{\begin{array}{@{}l@{}} \forall n \in \mathbb{N}, f \in (\mathbb{N} \to \mathbb{R}). \\ \quad \operatorname{series\text{-}partial\text{-}sum}((\lambda k \in \mathbb{N}.\; (f(k) \cdot \operatorname{power}(0, k))), (n + 1)) = f(0) \end{array}}}
\vnbprop{poly-nth-deriv-at-zero}{\vnbstmt{\begin{array}{@{}l@{}} \forall \mathit{tck}, n \in \mathbb{N}, f \in (\mathbb{N} \to \mathbb{R}). \\ \quad (\mathit{tck} \leq n) \Rightarrow  \\ \quad \quad \operatorname{nth\text{-}deriv}((\lambda x \in \mathbb{R}.\; \operatorname{series\text{-}partial\text{-}sum}((\lambda k \in \mathbb{N}.\; (f(k) \cdot \operatorname{power}(x, k))), (n + 1))), \mathit{tck})(0) \\ \quad \quad =\; \operatorname{factorial}(\mathit{tck}) \cdot f(\mathit{tck}) \end{array}}}
\vnbprop{taylor-coef-from-derivs}{\vnbstmt{\begin{array}{@{}l@{}} \forall n \in \mathbb{N}, f.; \forall a \in (\mathbb{N} \to \mathbb{R}), \mathit{tck} \in \mathbb{N}. \\ \quad \forall \mathit{tck} \in \mathbb{N}. \\ \quad \quad (\mathit{tck} \leq n) \Rightarrow  \\ \quad \quad \quad \operatorname{nth\text{-}deriv}((\lambda x \in \mathbb{R}.\; \operatorname{series\text{-}partial\text{-}sum}((\lambda k \in \mathbb{N}.\; (a(k) \cdot \operatorname{power}(x, k))), (n + 1))), \mathit{tck})(0) \\ \quad \quad \quad =\; \operatorname{nth\text{-}deriv}(f, \mathit{tck})(0) \wedge \\ \quad \mathit{tck} \leq n \Rightarrow \\ \quad \quad a(\mathit{tck}) = (\operatorname{nth\text{-}deriv}(f, \mathit{tck})(0) \cdot \operatorname{recip}(\operatorname{factorial}(\mathit{tck}))) \end{array}}}
\vnbprop{taylor-derivs-from-coef}{\vnbstmt{\begin{array}{@{}l@{}} \forall n \in \mathbb{N}, f.; \forall a \in (\mathbb{N} \to \mathbb{R}), \mathit{tck} \in \mathbb{N}. \\ \quad \forall \mathit{tck} \in \mathbb{N}.\; ((\mathit{tck} \leq n) \Rightarrow (\operatorname{nth\text{-}deriv}(f, \mathit{tck})(0) \in \mathbb{R})) \wedge \\ \quad \forall \mathit{tck} \in \mathbb{N}. \\ \quad \quad (\mathit{tck} \leq n) \Rightarrow (a(\mathit{tck}) = (\operatorname{nth\text{-}deriv}(f, \mathit{tck})(0) \cdot \operatorname{recip}(\operatorname{factorial}(\mathit{tck})))) \wedge \\ \quad \mathit{tck} \leq n \Rightarrow \\ \quad \quad \operatorname{nth\text{-}deriv}((\lambda x \in \mathbb{R}.\; \operatorname{series\text{-}partial\text{-}sum}((\lambda k \in \mathbb{N}.\; (a(k) \cdot \operatorname{power}(x, k))), (n + 1))), \mathit{tck})(0) \\ \quad \quad =\; \operatorname{nth\text{-}deriv}(f, \mathit{tck})(0) \end{array}}}
\vnbprop{taylor-interpolant-at-zero}{\vnbstmt{\begin{array}{@{}l@{}} \forall n \in \mathbb{N}, f. \\ \quad \forall \mathit{tck} \in \mathbb{N}.\; ((\mathit{tck} \leq n) \Rightarrow (\operatorname{nth\text{-}deriv}(f, \mathit{tck})(0) \in \mathbb{R})) \Rightarrow  \\ \quad \quad \exists a \in (\mathbb{N} \to \mathbb{R}). \\ \quad \quad \quad \forall \mathit{tck} \in \mathbb{N}. \\ \quad \quad \quad \quad (\mathit{tck} \leq n) \Rightarrow (a(\mathit{tck}) = (\operatorname{nth\text{-}deriv}(f, \mathit{tck})(0) \cdot \operatorname{recip}(\operatorname{factorial}(\mathit{tck})))) \wedge \\ \quad \quad \quad \forall \mathit{tck} \in \mathbb{N}. \\ \quad \quad \quad \quad (\mathit{tck} \leq n) \Rightarrow  \\ \quad \quad \quad \quad \quad \operatorname{nth\text{-}deriv}((\lambda x \in \mathbb{R}.\; \operatorname{series\text{-}partial\text{-}sum}((\lambda k \in \mathbb{N}.\; (a(k) \cdot \operatorname{power}(x, k))), (n + 1))), \mathit{tck})(0) \\ \quad \quad \quad \quad \quad =\; \operatorname{nth\text{-}deriv}(f, \mathit{tck})(0) \wedge \\ \quad \quad \quad \forall g \in (\mathbb{N} \to \mathbb{R}), \mathit{tck} \in \mathbb{N}. \\ \quad \quad \quad \quad \forall \mathit{tck} \in \mathbb{N}. \\ \quad \quad \quad \quad \quad (\mathit{tck} \leq n) \Rightarrow  \\ \quad \quad \quad \quad \quad \quad \operatorname{nth\text{-}deriv}((\lambda x \in \mathbb{R}.\; \operatorname{series\text{-}partial\text{-}sum}((\lambda k \in \mathbb{N}.\; (g(k) \cdot \operatorname{power}(x, k))), (n + 1))), \mathit{tck})(0) \\ \quad \quad \quad \quad \quad \quad =\; \operatorname{nth\text{-}deriv}(f, \mathit{tck})(0) \wedge \\ \quad \quad \quad \quad \mathit{tck} \leq n \Rightarrow \\ \quad \quad \quad \quad \quad g(\mathit{tck}) = a(\mathit{tck}) \end{array}}}
\vnbprop{series-partial-sum-shift}{\vnbstmt{\begin{array}{@{}l@{}} \forall n \in \mathbb{N}, u \in (\mathbb{N} \to \mathbb{R}). \\ \quad \operatorname{series\text{-}partial\text{-}sum}(u, (n + 1)) = (u(0) + \operatorname{series\text{-}partial\text{-}sum}((\lambda i \in \mathbb{N}.\; u((i + 1))), n)) \end{array}}}
\vnbprop{poly-degree-zero-value}{\vnbstmt{\begin{array}{@{}l@{}} \forall f \in (\mathbb{N} \to \mathbb{R}), v \in \mathbb{R}. \\ \quad \operatorname{series\text{-}partial\text{-}sum}((\lambda k \in \mathbb{N}.\; (f(k) \cdot \operatorname{power}(v, k))), 1) = f(0) \end{array}}}
\vnbprop{taylor-exact-polynomial}{\vnbstmt{\begin{array}{@{}l@{}} \forall m \in \mathbb{N}, f \in (\mathbb{N} \to \mathbb{R}), a, h \in \mathbb{R}. \\ \quad (\lambda x \in \mathbb{R}.\; \operatorname{series\text{-}partial\text{-}sum}((\lambda k \in \mathbb{N}.\; (f(k) \cdot \operatorname{power}(x, k))), (m + 1)))((a + h)) \\ \quad =\; \operatorname{series\text{-}partial\text{-}sum}((\lambda i \in \mathbb{N}.\; \\ \quad \quad \quad (\operatorname{nth\text{-}deriv}((\lambda x \in \mathbb{R}.\; \operatorname{series\text{-}partial\text{-}sum}((\lambda k \in \mathbb{N}.\; (f(k) \cdot \operatorname{power}(x, k))), (m + 1))), i)(a) \cdot \operatorname{power}(h, i)) \cdot \operatorname{recip}(\operatorname{factorial}(i))), \\ \quad \quad (m + 1)) \end{array}}}
\vnbprop{taylor-remainder-sup-bound}{\vnbstmt{\begin{array}{@{}l@{}} \forall a, b \in \mathbb{R}, f.; \forall m \in \mathbb{N}. \\ \quad \operatorname{is\text{-}taylor\text{-}family}(f, (m + 1), a, b) \wedge \\ \quad \operatorname{rr\text{-}bounded\text{-}above}(\operatorname{image}((\lambda v \in \operatorname{ooint}(a, b).\; \lvert f((m + 1))(v) \rvert), \operatorname{ooint}(a, b))) \Rightarrow \\ \quad \quad \lvert \operatorname{taylor\text{-}rem}(f, (m + 1), a, b) \rvert \\ \quad \quad \leq\; (\operatorname{power}((b - a), (m + 1)) \cdot \operatorname{recip}(\operatorname{factorial}((m + 1)))) \cdot \operatorname{sup}(\operatorname{image}((\lambda v \in \operatorname{ooint}(a, b).\; \lvert f((m + 1))(v) \rvert), \operatorname{ooint}(a, b))) \end{array}}}
\vnbprop{series-partial-sum-agree-below}{\vnbstmt{\begin{array}{@{}l@{}} \forall n \in \mathbb{N}, u, g. \\ \quad (\forall k \in \mathbb{N}.\; ((k < n) \Rightarrow (g(k) \in \mathbb{R})) \wedge \forall k \in \mathbb{N}.\; ((k < n) \Rightarrow (u(k) = g(k)))) \Rightarrow  \\ \quad \quad \operatorname{series\text{-}partial\text{-}sum}(u, n) = \operatorname{series\text{-}partial\text{-}sum}(g, n) \end{array}}}
\vnbprop{taylor-lagrange-subinterval}{\vnbstmt{\begin{array}{@{}l@{}} \forall a, b \in \mathbb{R}, f.; \forall m \in \mathbb{N}, t \in \mathbb{R}. \\ \quad \forall k \in \mathbb{N}.\; ((k \leq (m + 1)) \Rightarrow (f(k) \in (\operatorname{ccint}(a, b) \to \mathbb{R}))) \wedge \\ \quad \forall k \in \mathbb{N}.\; ((k < (m + 1)) \Rightarrow \operatorname{is\text{-}continuous\text{-}on}(f(k), \operatorname{ccint}(a, b))) \wedge \\ \quad \forall k \in \mathbb{N}, x \in \operatorname{ooint}(a, b). \\ \quad \quad (k < (m + 1)) \Rightarrow \operatorname{has\text{-}deriv\text{-}at}(f(k), x, f((k + 1))(x)) \wedge \\ \quad a < t \wedge \\ \quad t \leq b \Rightarrow \\ \quad \quad \exists i \in \operatorname{ooint}(a, t). \\ \quad \quad \quad \operatorname{taylor\text{-}rem}(f, (m + 1), a, t) = ((f((m + 1))(i) \cdot \operatorname{recip}(\operatorname{factorial}((m + 1)))) \cdot \operatorname{power}((t - a), (m + 1))) \end{array}}}
\vnbprop{taylor-sum-at-base}{\vnbstmt{\begin{array}{@{}l@{}} \forall j \in \mathbb{N}, f.; \forall a \in \mathbb{R}. \\ \quad \forall k \in \mathbb{N}.\; ((k \leq j) \Rightarrow (f(k)(a) \in \mathbb{R})) \Rightarrow  \\ \quad \quad \operatorname{series\text{-}partial\text{-}sum}((\lambda \mathit{tk} \in \mathbb{N}.\; ((f(\mathit{tk})(a) \cdot \operatorname{power}((a - a), \mathit{tk})) \cdot \operatorname{recip}(\operatorname{factorial}(\mathit{tk})))), (j + 1)) \\ \quad \quad =\; f(0)(a) \end{array}}}
\vnbprop{taylor-factor-value}{\vnbstmt{\begin{array}{@{}l@{}} \forall a, b \in \mathbb{R}, f.; \forall m \in \mathbb{N}, t \in \mathbb{R}. \\ \quad \forall k \in \mathbb{N}.\; ((k \leq (m + 1)) \Rightarrow (f(k) \in (\operatorname{ccint}(a, b) \to \mathbb{R}))) \wedge \\ \quad \forall k \in \mathbb{N}.\; ((k < (m + 1)) \Rightarrow \operatorname{is\text{-}continuous\text{-}on}(f(k), \operatorname{ccint}(a, b))) \wedge \\ \quad \forall k \in \mathbb{N}, x \in \operatorname{ooint}(a, b). \\ \quad \quad (k < (m + 1)) \Rightarrow \operatorname{has\text{-}deriv\text{-}at}(f(k), x, f((k + 1))(x)) \wedge \\ \quad a < t \wedge \\ \quad t \leq b \Rightarrow \\ \quad \quad \exists i \in \operatorname{ooint}(a, t). \\ \quad \quad \quad ((\operatorname{taylor\text{-}rem}(f, (m + 1), a, t) \cdot \operatorname{factorial}((m + 1))) \cdot \operatorname{recip}(\operatorname{power}((t - a), (m + 1)))) = f((m + 1))(i) \end{array}}}
\vnbprop{taylor-sum-continuous}{\vnbstmt{\begin{array}{@{}l@{}} \forall a \in \mathbb{R}, f.; \forall m \in \mathbb{N}, p \in \mathbb{R}. \\ \quad \forall k \in \mathbb{N}.\; ((k < (m + 1)) \Rightarrow (f(k)(a) \in \mathbb{R})) \Rightarrow  \\ \quad \quad \operatorname{is\text{-}continuous\text{-}at}(\operatorname{rr\text{-}ms}, \operatorname{rr\text{-}ms}, (\lambda s \in \mathbb{R}.\; \operatorname{taylor\text{-}sum}(f, (m + 1), s, a)), p) \end{array}}}
\vnbprop{taylor-factor-fn-continuous}{\vnbstmt{\begin{array}{@{}l@{}} \forall a, b \in \mathbb{R}, f.; \forall m \in \mathbb{N}, p \in \mathbb{R}. \\ \quad a \leq b \wedge \\ \quad \operatorname{is\text{-}continuous\text{-}on}(f(0), \operatorname{ccint}(a, b)) \wedge \\ \quad \forall k \in \mathbb{N}.\; ((k < (m + 1)) \Rightarrow (f(k)(a) \in \mathbb{R})) \wedge \\ \quad a < p \Rightarrow \\ \quad \quad \operatorname{is\text{-}continuous\text{-}at}(\operatorname{rr\text{-}ms}, \operatorname{rr\text{-}ms}, \\ \quad \quad \quad (\lambda s \in \mathbb{R}.\; \\ \quad \quad \quad \quad ((\operatorname{extend\text{-}const}(f(0), a, b)(s) - \operatorname{taylor\text{-}sum}(f, (m + 1), s, a)) \cdot \operatorname{factorial}((m + 1))) \cdot {\operatorname{power}((s - a), (m + 1))}^{\ast\!-1}), \\ \quad \quad \quad p) \end{array}}}
\vnbprop{taylor-factor-fn-value}{\vnbstmt{\begin{array}{@{}l@{}} \forall a, b \in \mathbb{R}, f.; \forall m \in \mathbb{N}, t \in \mathbb{R}. \\ \quad \forall k \in \mathbb{N}.\; ((k \leq (m + 1)) \Rightarrow (f(k) \in (\operatorname{ccint}(a, b) \to \mathbb{R}))) \wedge \\ \quad \forall k \in \mathbb{N}.\; ((k < (m + 1)) \Rightarrow \operatorname{is\text{-}continuous\text{-}on}(f(k), \operatorname{ccint}(a, b))) \wedge \\ \quad \forall k \in \mathbb{N}, x \in \operatorname{ooint}(a, b). \\ \quad \quad (k < (m + 1)) \Rightarrow \operatorname{has\text{-}deriv\text{-}at}(f(k), x, f((k + 1))(x)) \wedge \\ \quad a < t \wedge \\ \quad t \leq b \Rightarrow \\ \quad \quad \exists i \in \operatorname{ooint}(a, t). \\ \quad \quad \quad (\lambda s \in \mathbb{R}.\; (((\operatorname{extend\text{-}const}(f(0), a, b)(s) - \operatorname{taylor\text{-}sum}(f, (m + 1), s, a)) \cdot \operatorname{factorial}((m + 1))) \cdot {\operatorname{power}((s - a), (m + 1))}^{\ast\!-1}))(t) \\ \quad \quad \quad =\; f((m + 1))(i) \end{array}}}
\vnbprop{taylor-factor-glue-continuous}{\vnbstmt{\begin{array}{@{}l@{}} \forall a, b \in \mathbb{R}, f.; \forall m \in \mathbb{N}, p \in \mathbb{R}. \\ \quad a < b \wedge \\ \quad \forall k \in \mathbb{N}.\; ((k \leq (m + 1)) \Rightarrow (f(k) \in (\operatorname{ccint}(a, b) \to \mathbb{R}))) \wedge \\ \quad \forall k \in \mathbb{N}.\; ((k < (m + 1)) \Rightarrow \operatorname{is\text{-}continuous\text{-}on}(f(k), \operatorname{ccint}(a, b))) \wedge \\ \quad \forall k \in \mathbb{N}, x \in \operatorname{ooint}(a, b). \\ \quad \quad (k < (m + 1)) \Rightarrow \operatorname{has\text{-}deriv\text{-}at}(f(k), x, f((k + 1))(x)) \wedge \\ \quad \operatorname{is\text{-}continuous\text{-}on}(f((m + 1)), \operatorname{ccint}(a, b)) \Rightarrow \\ \quad \quad \operatorname{is\text{-}continuous\text{-}at}(\operatorname{rr\text{-}ms}, \operatorname{rr\text{-}ms}, \\ \quad \quad \quad (\lambda s \in \mathbb{R}.\; \\ \quad \quad \quad \quad \left(\text{if }(s \leq a)\text{ then }f((m + 1))(a)\text{ else }(((\operatorname{extend\text{-}const}(f(0), a, b)(s) - \operatorname{taylor\text{-}sum}(f, (m + 1), s, a)) \cdot \operatorname{factorial}((m + 1))) \cdot {\operatorname{power}((s - a), (m + 1))}^{\ast\!-1})\right)), \\ \quad \quad \quad p) \end{array}}}
\vnbprop{taylor-remainder-factor}{\vnbstmt{\begin{array}{@{}l@{}} \forall a, b \in \mathbb{R}, f.; \forall m \in \mathbb{N}. \\ \quad a < b \wedge \\ \quad \forall k \in \mathbb{N}.\; ((k \leq (m + 1)) \Rightarrow (f(k) \in (\operatorname{ccint}(a, b) \to \mathbb{R}))) \wedge \\ \quad \forall k \in \mathbb{N}.\; ((k < (m + 1)) \Rightarrow \operatorname{is\text{-}continuous\text{-}on}(f(k), \operatorname{ccint}(a, b))) \wedge \\ \quad \forall k \in \mathbb{N}, x \in \operatorname{ooint}(a, b). \\ \quad \quad (k < (m + 1)) \Rightarrow \operatorname{has\text{-}deriv\text{-}at}(f(k), x, f((k + 1))(x)) \wedge \\ \quad \operatorname{is\text{-}continuous\text{-}on}(f((m + 1)), \operatorname{ccint}(a, b)) \Rightarrow \\ \quad \quad \exists s \in (\operatorname{ccint}(a, b) \to \mathbb{R}). \\ \quad \quad \quad \operatorname{is\text{-}continuous\text{-}on}(s, \operatorname{ccint}(a, b)) \wedge \\ \quad \quad \quad s(a) = f((m + 1))(a) \wedge \\ \quad \quad \quad \forall t \in \operatorname{ccint}(a, b). \\ \quad \quad \quad \quad \operatorname{taylor\text{-}rem}(f, (m + 1), a, t) = ((s(t) \cdot \operatorname{power}((t - a), (m + 1))) \cdot \operatorname{recip}(\operatorname{factorial}((m + 1)))) \wedge \\ \quad \quad \quad \forall t \in \operatorname{ccint}(a, b). \\ \quad \quad \quad \quad \lvert s(t) \rvert \leq \operatorname{sup}(\operatorname{image}((\lambda v \in \operatorname{ccint}(a, b).\; \lvert f((m + 1))(v) \rvert), \operatorname{ccint}(a, b))) \end{array}}}

\vnbbookfile{structure-library/tb11-vocab}{structure-library/tb11-vocab}
\noindent{\small tb11-vocab.scm -- vocabulary for topological-vector-spaces Appendix A.3 (TB-11): finite products over a sequence in a monoid, and the subadditive regularisation of Proposition A.5.  DEFINITIONS ONLY.  Agent ISOM, 2026-10-26 (new definitions, the user's to confirm).}\par
\par\medskip\noindent\textbf{Definition} (\texttt{mprod}).\;For all $\operatorname{mpm\_}$ and $\operatorname{mpu\_}$,  $(\operatorname{mprod}(\operatorname{mpm\_}, \operatorname{mpu\_}, 0) \simeq \operatorname{iden}(\operatorname{mpm\_}))$\par
Let $\operatorname{mpm\_}$ and $\operatorname{mpu\_}$ be arbitrary. Let $\operatorname{mpn\_} \in \mathbb{N}$. Then
\begin{equation*}
(\operatorname{mprod}(\operatorname{mpm\_}, \operatorname{mpu\_}, \operatorname{succ}(\operatorname{mpn\_})) \simeq \operatorname{opr}(\operatorname{mpm\_})(\operatorname{mprod}(\operatorname{mpm\_}, \operatorname{mpu\_}, \operatorname{mpn\_}), \operatorname{mpu\_}(\operatorname{mpn\_})))
\end{equation*}

\par\medskip\noindent\textbf{Definition} (\texttt{fact-sums}).\;$\operatorname{fact\text{-}sums}(\operatorname{fsm\_}, \operatorname{fsg\_}, \operatorname{fsx\_}) :=$
\begin{equation*}
\{\operatorname{fsr\_} \in \mathbb{R} : \exists \operatorname{fsn\_} \in \mathbb{N}, \operatorname{fsu\_} \in (\mathbb{N} \to \operatorname{carr}(\operatorname{fsm\_})).\; ((\operatorname{mprod}(\operatorname{fsm\_}, \operatorname{fsu\_}, \operatorname{fsn\_}) = \operatorname{fsx\_}) \wedge (\operatorname{fsr\_} = \operatorname{series\text{-}partial\text{-}sum}((\lambda \operatorname{fsi\_} \in \mathbb{N}.\; \operatorname{fsg\_}(\operatorname{fsu\_}(\operatorname{fsi\_}))), \operatorname{fsn\_})))\}
\end{equation*}

\par\medskip\noindent\textbf{Definition} (\texttt{subadd-reg}).\;$\operatorname{subadd\text{-}reg}(\operatorname{srm\_}, \operatorname{srg\_}, \operatorname{srx\_}) :=$ $\operatorname{einf}(\operatorname{fact\text{-}sums}(\operatorname{srm\_}, \operatorname{srg\_}, \operatorname{srx\_}))$

\vnbbookfile{theorem-library/tb11-a5}{theorem-library/tb11-a5}
\noindent{\small tb11-a5.scm -- topological-vector-spaces Proposition A.5 (subadditive regularisation on a monoid) and the finite-product laws it needs.  Agent ISOM, 2026-10-26.  Helper prefix: ta5-. Library: analysis (EINF, partial sums); vocabulary tb11-vocab.scm.}\par
\vnbprop{mprod-in-carr-ind}{\vnbstmt{\begin{array}{@{}l@{}} \forall n \in \mathbb{N}, m.; \forall u \in (\mathbb{N} \to \operatorname{carr}(m)). \\ \quad \operatorname{is\text{-}monoid}(m) \Rightarrow (\operatorname{mprod}(m, u, n) \in \operatorname{carr}(m)) \end{array}}}
\vnbprop{mprod-in-carr}{\vnbstmt{\begin{array}{@{}l@{}} \forall m.; \forall u \in (\mathbb{N} \to \operatorname{carr}(m)), n \in \mathbb{N}. \\ \quad \operatorname{is\text{-}monoid}(m) \Rightarrow (\operatorname{mprod}(m, u, n) \in \operatorname{carr}(m)) \end{array}}}
\vnbprop{mprod-shift-split-ind}{\vnbstmt{\begin{array}{@{}l@{}} \forall k \in \mathbb{N}, m.; \forall u \in (\mathbb{N} \to \operatorname{carr}(m)), a \in \mathbb{N}. \\ \quad \operatorname{is\text{-}monoid}(m) \Rightarrow  \\ \quad \quad \operatorname{mprod}(m, u, (a + k)) = \operatorname{opr}(m)(\operatorname{mprod}(m, u, a), \operatorname{mprod}(m, (\lambda i \in \mathbb{N}.\; u((a + i))), k)) \end{array}}}
\vnbprop{partial-sum-shift-split-ind}{\vnbstmt{\begin{array}{@{}l@{}} \forall k \in \mathbb{N}, f \in (\mathbb{N} \to \mathbb{R}), a \in \mathbb{N}. \\ \quad \operatorname{series\text{-}partial\text{-}sum}(f, (a + k)) \\ \quad =\; \operatorname{series\text{-}partial\text{-}sum}(f, a) + \operatorname{series\text{-}partial\text{-}sum}((\lambda i \in \mathbb{N}.\; f((a + i))), k) \end{array}}}
\vnbprop{partial-sum-half-split}{\vnbstmt{\begin{array}{@{}l@{}} \forall f \in (\mathbb{N} \to \mathbb{R}), n \in \mathbb{N}. \\ \quad \forall i \in \mathbb{N}.\; (0 \leq f(i)) \Rightarrow  \\ \quad \quad \exists q \in \mathbb{N}. \\ \quad \quad \quad q \leq n \wedge \\ \quad \quad \quad (2 \cdot \operatorname{series\text{-}partial\text{-}sum}(f, q)) \leq \operatorname{series\text{-}partial\text{-}sum}(f, (n + 1)) \wedge \\ \quad \quad \quad 2 \cdot (\operatorname{series\text{-}partial\text{-}sum}(f, (n + 1)) - \operatorname{series\text{-}partial\text{-}sum}(f, (q + 1))) \\ \quad \quad \quad \leq\; \operatorname{series\text{-}partial\text{-}sum}(f, (n + 1)) \end{array}}}
\vnbprop{gauge-product-bound-ind}{\vnbstmt{\begin{array}{@{}l@{}} \forall b \in \mathbb{N}, m.; \forall g \in (\operatorname{carr}(m) \to \mathbb{R}), n \in \mathbb{N}, u \in (\mathbb{N} \to \operatorname{carr}(m)). \\ \quad \operatorname{is\text{-}monoid}(m) \wedge \\ \quad \forall x \in \operatorname{carr}(m).\; (0 \leq g(x)) \wedge \\ \quad g(\operatorname{iden}(m)) = 0 \wedge \\ \quad \forall x, y, z \in \operatorname{carr}(m). \\ \quad \quad g(\operatorname{opr}(m)(\operatorname{opr}(m)(x, y), z)) \leq (2 \cdot \operatorname{max}(g(x), \operatorname{max}(g(y), g(z)))) \wedge \\ \quad n \leq b \Rightarrow \\ \quad \quad g(\operatorname{mprod}(m, u, n)) \leq (2 \cdot \operatorname{series\text{-}partial\text{-}sum}((\lambda i \in \mathbb{N}.\; g(u(i))), n)) \end{array}}}
\vnbprop{fact-sums-unfold}{\vnbstmt{\begin{array}{@{}l@{}} \forall m, g, x. \\ \quad (\operatorname{fact\text{-}sums}(m, g, x) \simeq \{r \in \mathbb{R} : \exists n \in \mathbb{N}, u \in (\mathbb{N} \to \operatorname{carr}(m)).\; ((\operatorname{mprod}(m, u, n) = x) \wedge (r = \operatorname{series\text{-}partial\text{-}sum}((\lambda i \in \mathbb{N}.\; g(u(i))), n)))\}) \end{array}}}
\vnbprop{subadd-reg-unfold}{\vnbstmt{\forall m, g, x.\; (\operatorname{subadd\text{-}reg}(m, g, x) \simeq \operatorname{einf}(\operatorname{fact\text{-}sums}(m, g, x)))}}
\vnbprop{fact-sums-intro}{\vnbstmt{\begin{array}{@{}l@{}} \forall m.; \forall g \in (\operatorname{carr}(m) \to \mathbb{R}), u \in (\mathbb{N} \to \operatorname{carr}(m)), n \in \mathbb{N}. \\ \quad \operatorname{is\text{-}monoid}(m) \wedge \\ \quad \forall x \in \operatorname{carr}(m).\; (0 \leq g(x)) \wedge \\ \quad g(\operatorname{iden}(m)) = 0 \wedge \\ \quad \forall x, y, z \in \operatorname{carr}(m). \\ \quad \quad g(\operatorname{opr}(m)(\operatorname{opr}(m)(x, y), z)) \leq (2 \cdot \operatorname{max}(g(x), \operatorname{max}(g(y), g(z)))) \Rightarrow \\ \quad \quad \operatorname{series\text{-}partial\text{-}sum}((\lambda i \in \mathbb{N}.\; g(u(i))), n) \in \operatorname{fact\text{-}sums}(m, g, \operatorname{mprod}(m, u, n)) \end{array}}}
\vnbprop{fact-sums-nonneg}{\vnbstmt{\begin{array}{@{}l@{}} \forall m.; \forall g \in (\operatorname{carr}(m) \to \mathbb{R}), x \in \operatorname{carr}(m). \\ \quad \operatorname{is\text{-}monoid}(m) \wedge \\ \quad \forall x \in \operatorname{carr}(m).\; (0 \leq g(x)) \wedge \\ \quad g(\operatorname{iden}(m)) = 0 \wedge \\ \quad \forall x, y, z \in \operatorname{carr}(m). \\ \quad \quad g(\operatorname{opr}(m)(\operatorname{opr}(m)(x, y), z)) \leq (2 \cdot \operatorname{max}(g(x), \operatorname{max}(g(y), g(z)))) \Rightarrow \\ \quad \quad \operatorname{fact\text{-}sums}(m, g, x) \subseteq \operatorname{rr\text{-}pos\text{-}star} \end{array}}}
\vnbprop{subadd-reg-bounds}{\vnbstmt{\begin{array}{@{}l@{}} \forall m.; \forall g \in (\operatorname{carr}(m) \to \mathbb{R}), x \in \operatorname{carr}(m). \\ \quad \operatorname{is\text{-}monoid}(m) \wedge \\ \quad \forall x \in \operatorname{carr}(m).\; (0 \leq g(x)) \wedge \\ \quad g(\operatorname{iden}(m)) = 0 \wedge \\ \quad \forall x, y, z \in \operatorname{carr}(m). \\ \quad \quad g(\operatorname{opr}(m)(\operatorname{opr}(m)(x, y), z)) \leq (2 \cdot \operatorname{max}(g(x), \operatorname{max}(g(y), g(z)))) \Rightarrow \\ \quad \quad ((\operatorname{subadd\text{-}reg}(m, g, x) \in \mathbb{R}) \wedge (((1/2 \cdot g(x)) \leq \operatorname{subadd\text{-}reg}(m, g, x)) \wedge (\operatorname{subadd\text{-}reg}(m, g, x) \leq g(x)))) \end{array}}}
\vnbprop{mprod-agree-below}{\vnbstmt{\begin{array}{@{}l@{}} \forall k \in \mathbb{N}, m.; \forall u, w \in (\mathbb{N} \to \operatorname{carr}(m)). \\ \quad (\operatorname{is\text{-}monoid}(m) \wedge \forall i \in \mathbb{N}.\; ((i < k) \Rightarrow (u(i) = w(i)))) \Rightarrow  \\ \quad \quad \operatorname{mprod}(m, u, k) = \operatorname{mprod}(m, w, k) \end{array}}}
\vnbprop{fact-sums-concat}{\vnbstmt{\begin{array}{@{}l@{}} \forall m.; \forall g \in (\operatorname{carr}(m) \to \mathbb{R}), a \in \mathbb{N}, b \in (\mathbb{N} \to \operatorname{carr}(m)), c \in \mathbb{N}, d \in (\mathbb{N} \to \operatorname{carr}(m)). \\ \quad \operatorname{is\text{-}monoid}(m) \wedge \\ \quad \forall x \in \operatorname{carr}(m).\; (0 \leq g(x)) \wedge \\ \quad g(\operatorname{iden}(m)) = 0 \wedge \\ \quad \forall x, y, z \in \operatorname{carr}(m). \\ \quad \quad g(\operatorname{opr}(m)(\operatorname{opr}(m)(x, y), z)) \leq (2 \cdot \operatorname{max}(g(x), \operatorname{max}(g(y), g(z)))) \Rightarrow \\ \quad \quad \operatorname{series\text{-}partial\text{-}sum}((\lambda i \in \mathbb{N}.\; g(b(i))), a) + \operatorname{series\text{-}partial\text{-}sum}((\lambda i \in \mathbb{N}.\; g(d(i))), c) \\ \quad \quad \in\; \operatorname{fact\text{-}sums}(m, g, \operatorname{opr}(m)(\operatorname{mprod}(m, b, a), \operatorname{mprod}(m, d, c))) \end{array}}}
\vnbprop{subadd-reg-subadditive}{\vnbstmt{\begin{array}{@{}l@{}} \forall m.; \forall g \in (\operatorname{carr}(m) \to \mathbb{R}), x, y \in \operatorname{carr}(m). \\ \quad \operatorname{is\text{-}monoid}(m) \wedge \\ \quad \forall x \in \operatorname{carr}(m).\; (0 \leq g(x)) \wedge \\ \quad g(\operatorname{iden}(m)) = 0 \wedge \\ \quad \forall x, y, z \in \operatorname{carr}(m). \\ \quad \quad g(\operatorname{opr}(m)(\operatorname{opr}(m)(x, y), z)) \leq (2 \cdot \operatorname{max}(g(x), \operatorname{max}(g(y), g(z)))) \Rightarrow \\ \quad \quad \operatorname{subadd\text{-}reg}(m, g, \operatorname{opr}(m)(x, y)) \leq (\operatorname{subadd\text{-}reg}(m, g, x) + \operatorname{subadd\text{-}reg}(m, g, y)) \end{array}}}

\vnbbookfile{structure-library/tg-vocab}{structure-library/tg-vocab}
\noindent{\small tg-vocab.scm -- topological groups (topological-vector-space.pdf Appendix A.3, TB-11). DEFINITIONS ONLY.  Agent ISOM, 2026-10-26 (new definitions, the user's to confirm).}\par
\vnbstructdef{top-group}{A top group $s$ is a 5-tuple whose components are listed below.\par
\begin{enumerate}[nosep]
\item $\operatorname{carr}(s) = \operatorname{nth}(1, s)$, a carrier (\texttt{carriers}): $\operatorname{carr}(s) \in \mathbf{Set}$.
\item $\operatorname{opr}(s) = \operatorname{nth}(2, s)$, an operation (\texttt{op}): $\operatorname{opr}(s) \in ((\operatorname{carr}(s) \times \operatorname{carr}(s)) \to \operatorname{carr}(s))$.
\item $\operatorname{iden}(s) = \operatorname{nth}(3, s)$, a constant (\texttt{constant}): $\operatorname{iden}(s) \in \operatorname{carr}(s)$.
\item $\operatorname{inv}(s) = \operatorname{nth}(4, s)$, an operation (\texttt{op}): $\operatorname{inv}(s) \in (\operatorname{carr}(s) \to \operatorname{carr}(s))$.
\item $\operatorname{topology}(s) = \operatorname{nth}(5, s)$, a family of subsets (\texttt{family}): $\operatorname{topology}(s) \in \operatorname{power}(\operatorname{power}(\operatorname{carr}(s)))$.
\end{enumerate}
The laws (the \texttt{property} and \texttt{law} clauses):
\begin{itemize}[nosep]
\item $\operatorname{is\text{-}associative}(\operatorname{opr}(s), \operatorname{carr}(s))$
\item $\operatorname{is\text{-}identity}(\operatorname{opr}(s), \operatorname{iden}(s), \operatorname{carr}(s))$
\item $\operatorname{has\text{-}inverses}(\operatorname{opr}(s), \operatorname{iden}(s), \operatorname{inv}(s), \operatorname{carr}(s))$
\item $\emptyset \in \operatorname{topology}(s)$
\item $\operatorname{carr}(s) \in \operatorname{topology}(s)$
\item $\forall u, v \in \operatorname{topology}(s).\; ((u \cap v) \in \operatorname{topology}(s))$
\item $\forall \mathit{fam} \in \mathbf{Set}.\; ((\mathit{fam} \subseteq \operatorname{topology}(s)) \Rightarrow (\operatorname{big\text{-}union}(u, \mathit{fam}, u) \in \operatorname{topology}(s)))$
\item
\begin{equation*}
\begin{array}{@{}l@{}} \forall \operatorname{tga\_}, \operatorname{tgb\_} \in \operatorname{carr}(s), \operatorname{tgw\_} \in \operatorname{topology}(s). \\ \quad (\operatorname{opr}(s)(\operatorname{tga\_}, \operatorname{inv}(s)(\operatorname{tgb\_})) \in \operatorname{tgw\_}) \Rightarrow  \\ \quad \quad \exists \operatorname{tgu\_} \in \operatorname{topology}(s). \\ \quad \quad \quad \operatorname{tga\_} \in \operatorname{tgu\_} \wedge \\ \quad \quad \quad \exists \operatorname{tgv\_} \in \operatorname{topology}(s). \\ \quad \quad \quad \quad ((\operatorname{tgb\_} \in \operatorname{tgv\_}) \wedge \forall \operatorname{tgp\_} \in \operatorname{tgu\_}, \operatorname{tgq\_} \in \operatorname{tgv\_}.\; (\operatorname{opr}(s)(\operatorname{tgp\_}, \operatorname{inv}(s)(\operatorname{tgq\_})) \in \operatorname{tgw\_})) \end{array}
\end{equation*}

\end{itemize}
The characteristic law (the axiom \texttt{is-top-group}): $\operatorname{is\text{-}top\text{-}group}(s)$ holds if and only if $\operatorname{length}(s) = 5$, every component has the typing listed with it, every law holds, and $s \in \mathbf{Set}$ (the structure is a set).
}
\par\medskip\noindent\textbf{Definition} (\texttt{is-group-topology}).\;For all $\operatorname{tgs\_}$ and $\operatorname{tgg\_}$, 
\begin{equation*}
\begin{array}{@{}l@{}} \operatorname{is\text{-}group\text{-}topology}(\operatorname{tgs\_}, \operatorname{tgg\_}) \\ \Leftrightarrow\; \operatorname{is\text{-}top\text{-}space}(\operatorname{tgs\_}) \wedge \\ \operatorname{is\text{-}group}(\operatorname{tgg\_}) \wedge \\ \operatorname{pts}(\operatorname{tgs\_}) = \operatorname{carr}(\operatorname{tgg\_}) \wedge \\ \forall \operatorname{tga\_}, \operatorname{tgb\_} \in \operatorname{carr}(\operatorname{tgg\_}), \operatorname{tgw\_} \in \operatorname{opens}(\operatorname{tgs\_}). \\ \quad (\operatorname{opr}(\operatorname{tgg\_})(\operatorname{tga\_}, \operatorname{inv}(\operatorname{tgg\_})(\operatorname{tgb\_})) \in \operatorname{tgw\_}) \Rightarrow  \\ \quad \quad \exists \operatorname{tgu\_} \in \operatorname{opens}(\operatorname{tgs\_}). \\ \quad \quad \quad \operatorname{tga\_} \in \operatorname{tgu\_} \wedge \\ \quad \quad \quad \exists \operatorname{tgv\_} \in \operatorname{opens}(\operatorname{tgs\_}). \\ \quad \quad \quad \quad ((\operatorname{tgb\_} \in \operatorname{tgv\_}) \wedge \forall \operatorname{tgp\_} \in \operatorname{tgu\_}, \operatorname{tgq\_} \in \operatorname{tgv\_}.\; (\operatorname{opr}(\operatorname{tgg\_})(\operatorname{tgp\_}, \operatorname{inv}(\operatorname{tgg\_})(\operatorname{tgq\_})) \in \operatorname{tgw\_})) \end{array}
\end{equation*}

\par\medskip\noindent\textbf{Definition} (\texttt{has-countable-nbhd-base}).\;For all $\operatorname{cbs\_}$ and $\operatorname{cbx\_}$, 
\begin{equation*}
\begin{array}{@{}l@{}} \operatorname{has\text{-}countable\text{-}nbhd\text{-}base}(\operatorname{cbs\_}, \operatorname{cbx\_}) \\ \Leftrightarrow\; \exists \operatorname{cbn\_} \in (\mathbb{N} \to \operatorname{opens}(\operatorname{cbs\_})). \\ \quad \forall \operatorname{cbk\_} \in \mathbb{N}.\; (\operatorname{cbx\_} \in \operatorname{cbn\_}(\operatorname{cbk\_})) \wedge \\ \quad \forall \operatorname{cbw\_} \in \operatorname{opens}(\operatorname{cbs\_}). \\ \quad \quad (\operatorname{cbx\_} \in \operatorname{cbw\_}) \Rightarrow \exists \operatorname{cbk\_} \in \mathbb{N}.\; (\operatorname{cbn\_}(\operatorname{cbk\_}) \subseteq \operatorname{cbw\_}) \end{array}
\end{equation*}

\par\medskip\noindent\textbf{Definition} (\texttt{is-left-invariant}).\;For all $\operatorname{lim\_}$ and $\operatorname{lig\_}$, 
\begin{equation*}
\begin{array}{@{}l@{}} \operatorname{is\text{-}left\text{-}invariant}(\operatorname{lim\_}, \operatorname{lig\_}) \\ \Leftrightarrow\; \forall \operatorname{lia\_}, \operatorname{lix\_}, \operatorname{liy\_} \in \operatorname{carr}(\operatorname{lig\_}). \\ \quad \operatorname{dist}(\operatorname{lim\_})(\operatorname{opr}(\operatorname{lig\_})(\operatorname{lia\_}, \operatorname{lix\_}), \operatorname{opr}(\operatorname{lig\_})(\operatorname{lia\_}, \operatorname{liy\_})) = \operatorname{dist}(\operatorname{lim\_})(\operatorname{lix\_}, \operatorname{liy\_}) \end{array}
\end{equation*}

\par\medskip\noindent\textbf{Definition} (\texttt{is-bk-sequence}).\;For all $\operatorname{bss\_}$, $\operatorname{bsg\_}$ and $\operatorname{bsv\_}$, 
\begin{equation*}
\begin{array}{@{}l@{}} \operatorname{is\text{-}bk\text{-}sequence}(\operatorname{bss\_}, \operatorname{bsg\_}, \operatorname{bsv\_}) \\ \Leftrightarrow\; \operatorname{bsv\_} \in (\mathbb{N} \to \operatorname{opens}(\operatorname{bss\_})) \wedge \\ \operatorname{bsv\_}(0) = \operatorname{pts}(\operatorname{bss\_}) \wedge \\ \forall \operatorname{bsk\_} \in \mathbb{N}. \\ \quad \operatorname{iden}(\operatorname{bsg\_}) \in \operatorname{bsv\_}(\operatorname{bsk\_}) \wedge \\ \quad \forall \operatorname{bsp\_} \in \operatorname{bsv\_}(\operatorname{bsk\_}).\; (\operatorname{inv}(\operatorname{bsg\_})(\operatorname{bsp\_}) \in \operatorname{bsv\_}(\operatorname{bsk\_})) \wedge \\ \quad \forall \operatorname{bsp\_}, \operatorname{bsq\_}, \operatorname{bsr\_} \in \operatorname{bsv\_}(\operatorname{succ}(\operatorname{bsk\_})). \\ \quad \quad \operatorname{opr}(\operatorname{bsg\_})(\operatorname{opr}(\operatorname{bsg\_})(\operatorname{bsp\_}, \operatorname{bsq\_}), \operatorname{bsr\_}) \in \operatorname{bsv\_}(\operatorname{bsk\_}) \end{array}
\end{equation*}

\par\medskip\noindent\textbf{Definition} (\texttt{bk-gauge}).\;$\operatorname{bk\text{-}gauge}(\operatorname{bgv\_}, \operatorname{bgx\_}) :=$
\begin{equation*}
\operatorname{einf}(\{\operatorname{bgr\_} \in \mathbb{R} : \exists \operatorname{bgk\_} \in \mathbb{N}.\; ((\operatorname{bgx\_} \in \operatorname{bgv\_}(\operatorname{bgk\_})) \wedge (\operatorname{bgr\_} = \operatorname{power}(1/2, \operatorname{bgk\_})))\})
\end{equation*}

\par\medskip\noindent\textbf{Definition} (\texttt{bk-f}).\;$\operatorname{bk\text{-}f}(\operatorname{bfg\_}, \operatorname{bfv\_}, \operatorname{bfx\_}) :=$
\begin{equation*}
\operatorname{subadd\text{-}reg}(\operatorname{grp\text{-}monoid}(\operatorname{bfg\_}), (\lambda \operatorname{bfz\_} \in \operatorname{carr}(\operatorname{bfg\_}).\; \operatorname{bk\text{-}gauge}(\operatorname{bfv\_}, \operatorname{bfz\_})), \operatorname{bfx\_})
\end{equation*}

\par\medskip\noindent\textbf{Definition} (\texttt{bk-metric}).\;$\operatorname{bk\text{-}metric}(\operatorname{bmg\_}, \operatorname{bmv\_}) :=$
\begin{equation*}
{[}\operatorname{carr}(\operatorname{bmg\_}), (\lambda {[}\operatorname{bmx\_}, \operatorname{bmy\_}] \in (\operatorname{carr}(\operatorname{bmg\_}) \times \operatorname{carr}(\operatorname{bmg\_})).\; (\operatorname{bk\text{-}f}(\operatorname{bmg\_}, \operatorname{bmv\_}, \operatorname{opr}(\operatorname{bmg\_})(\operatorname{inv}(\operatorname{bmg\_})(\operatorname{bmy\_}), \operatorname{bmx\_})) + \operatorname{bk\text{-}f}(\operatorname{bmg\_}, \operatorname{bmv\_}, \operatorname{opr}(\operatorname{bmg\_})(\operatorname{inv}(\operatorname{bmg\_})(\operatorname{bmx\_}), \operatorname{bmy\_}))))]
\end{equation*}

\par\medskip\noindent\textbf{Definition} (\texttt{top-group-space}).\;$\operatorname{top\text{-}group\text{-}space}(\operatorname{tgr\_}) :=$ ${[}\operatorname{carr}(\operatorname{tgr\_}), \operatorname{topology}(\operatorname{tgr\_})]$
\par\medskip\noindent\textbf{Definition} (\texttt{top-group-group}).\;$\operatorname{top\text{-}group\text{-}group}(\operatorname{tgr\_}) :=$ ${[}\operatorname{carr}(\operatorname{tgr\_}), \operatorname{opr}(\operatorname{tgr\_}), \operatorname{iden}(\operatorname{tgr\_}), \operatorname{inv}(\operatorname{tgr\_})]$
\par\medskip\noindent\textbf{Definition} (\texttt{make-top-group}).\;$\operatorname{make\text{-}top\text{-}group}(\operatorname{mks\_}, \operatorname{mkg\_}) :=$ ${[}\operatorname{carr}(\operatorname{mkg\_}), \operatorname{opr}(\operatorname{mkg\_}), \operatorname{iden}(\operatorname{mkg\_}), \operatorname{inv}(\operatorname{mkg\_}), \operatorname{opens}(\operatorname{mks\_})]$
\par\medskip\noindent\textbf{Definition} (\texttt{hom-top-group}).\;$\operatorname{hom\text{-}top\text{-}group}(a, b) :=$
\begin{equation*}
\{\operatorname{homf\_} \in (\operatorname{carr}(a) \to \operatorname{carr}(b)) : \operatorname{is\text{-}hom\text{-}top\text{-}group}(a, b, \operatorname{homf\_})\}
\end{equation*}

\vnbdef{is-hom-top-group-def}{\vnbstmt{\begin{array}{@{}l@{}} \forall s, t, f. \\ \quad \operatorname{is\text{-}hom\text{-}top\text{-}group}(s, t, f) \\ \quad \Leftrightarrow\; \operatorname{is\text{-}top\text{-}group}(s) \wedge \\ \quad \operatorname{is\text{-}top\text{-}group}(t) \wedge \\ \quad f \in (\operatorname{carr}(s) \to \operatorname{carr}(t)) \wedge \\ \quad \forall x_{1}, x_{2} \in \operatorname{carr}(s). \\ \quad \quad f(\operatorname{opr}(s)(x_{1}, x_{2})) = \operatorname{opr}(t)(f(x_{1}), f(x_{2})) \wedge \\ \quad f(\operatorname{iden}(s)) = \operatorname{iden}(t) \wedge \\ \quad \forall x_{1} \in \operatorname{carr}(s).\; (f(\operatorname{inv}(s)(x_{1})) = \operatorname{inv}(t)(f(x_{1}))) \wedge \\ \quad \forall u \in \operatorname{topology}(t).\; (\{x \in \operatorname{carr}(s) : (f(x) \in u)\} \in \operatorname{topology}(s)) \end{array}}}
\vnbdef{top-group-class}{\vnbstmt{\forall s.\; ((s \in \operatorname{top\text{-}group}) \Leftrightarrow \operatorname{is\text{-}top\text{-}group}(s))}}
\vnbprop{is-top-group-is-set}{\vnbstmt{\forall s.\; (\operatorname{is\text{-}top\text{-}group}(s) \Rightarrow (s \in \mathbf{Set}))}}

\vnbbookfile{theorem-library/tg-struct}{theorem-library/tg-struct}
\noindent{\small tg-struct.scm -- the structure TOP-GROUP and the group topologies: the underlying space and group of a topological group (the typing obligations of the two constructed functors), the bridge IS-TOP-GROUP(G) =$>$ IS-GROUP-TOPOLOGY(space, group), and the way back, IS-GROUP-TOPOLOGY(s, g) =$>$ IS-TOP-GROUP(MAKE-TOP-GROUP(s, ...}\par
\vnbprop{top-group-space-is-top-space}{\vnbstmt{\forall g.\; (\operatorname{is\text{-}top\text{-}group}(g) \Rightarrow \operatorname{is\text{-}top\text{-}space}(\operatorname{top\text{-}group\text{-}space}(g)))}}
\vnbprop{top-group-group-is-group}{\vnbstmt{\forall g.\; (\operatorname{is\text{-}top\text{-}group}(g) \Rightarrow \operatorname{is\text{-}group}(\operatorname{top\text{-}group\text{-}group}(g)))}}
\vnbprop{top-group-group-topology}{\vnbstmt{\begin{array}{@{}l@{}} \forall g. \\ \quad \operatorname{is\text{-}top\text{-}group}(g) \Rightarrow  \\ \quad \quad \operatorname{is\text{-}group\text{-}topology}(\operatorname{top\text{-}group\text{-}space}(g), \operatorname{top\text{-}group\text{-}group}(g)) \end{array}}}
\vnbprop{make-top-group-carr}{\vnbstmt{\begin{array}{@{}l@{}} \forall s, g. \\ \quad (\operatorname{is\text{-}top\text{-}space}(s) \wedge \operatorname{is\text{-}group}(g)) \Rightarrow (\operatorname{carr}(\operatorname{make\text{-}top\text{-}group}(s, g)) \simeq \operatorname{carr}(g)) \end{array}}}
\vnbprop{make-top-group-opr}{\vnbstmt{\begin{array}{@{}l@{}} \forall s, g. \\ \quad (\operatorname{is\text{-}top\text{-}space}(s) \wedge \operatorname{is\text{-}group}(g)) \Rightarrow (\operatorname{opr}(\operatorname{make\text{-}top\text{-}group}(s, g)) \simeq \operatorname{opr}(g)) \end{array}}}
\vnbprop{make-top-group-iden}{\vnbstmt{\begin{array}{@{}l@{}} \forall s, g. \\ \quad (\operatorname{is\text{-}top\text{-}space}(s) \wedge \operatorname{is\text{-}group}(g)) \Rightarrow (\operatorname{iden}(\operatorname{make\text{-}top\text{-}group}(s, g)) \simeq \operatorname{iden}(g)) \end{array}}}
\vnbprop{make-top-group-inv}{\vnbstmt{\begin{array}{@{}l@{}} \forall s, g. \\ \quad (\operatorname{is\text{-}top\text{-}space}(s) \wedge \operatorname{is\text{-}group}(g)) \Rightarrow (\operatorname{inv}(\operatorname{make\text{-}top\text{-}group}(s, g)) \simeq \operatorname{inv}(g)) \end{array}}}
\vnbprop{make-top-group-topology}{\vnbstmt{\begin{array}{@{}l@{}} \forall s, g. \\ \quad (\operatorname{is\text{-}top\text{-}space}(s) \wedge \operatorname{is\text{-}group}(g)) \Rightarrow  \\ \quad \quad (\operatorname{topology}(\operatorname{make\text{-}top\text{-}group}(s, g)) \simeq \operatorname{opens}(s)) \end{array}}}
\vnbprop{make-top-group-is-top-group}{\vnbstmt{\begin{array}{@{}l@{}} \forall s, g. \\ \quad \operatorname{is\text{-}group\text{-}topology}(s, g) \Rightarrow \operatorname{is\text{-}top\text{-}group}(\operatorname{make\text{-}top\text{-}group}(s, g)) \end{array}}}
\vnbprop{top-group-space-functorial}{\vnbstmt{\begin{array}{@{}l@{}} \forall a, b, f. \\ \quad \operatorname{is\text{-}hom\text{-}top\text{-}group}(a, b, f) \Rightarrow  \\ \quad \quad \operatorname{is\text{-}hom\text{-}top\text{-}space}(\operatorname{top\text{-}group\text{-}space}(a), \operatorname{top\text{-}group\text{-}space}(b), f) \end{array}}}
\vnbprop{top-group-group-functorial}{\vnbstmt{\begin{array}{@{}l@{}} \forall a, b, f. \\ \quad \operatorname{is\text{-}hom\text{-}top\text{-}group}(a, b, f) \Rightarrow  \\ \quad \quad \operatorname{is\text{-}hom\text{-}group}(\operatorname{top\text{-}group\text{-}group}(a), \operatorname{top\text{-}group\text{-}group}(b), f) \end{array}}}

\vnbbookfile{theorem-library/top-groups}{theorem-library/top-groups}
\noindent{\small tg-top.scm -- what continuity of (x, y) |-$>$ x y\textasciicircum{}-1 gives in a topological group: the inverse and the left translates of an open set are open, and an open neighbourhood of e contains a symmetric open V with V V V inside it.  Plus the local criterion for openness. Agent ISOM, 2026-10-26.  Helper prefix: tgt-.  Library: ...}\par
\vnbprop{top-open-of-local}{\vnbstmt{\begin{array}{@{}l@{}} \forall s, u. \\ \quad \operatorname{is\text{-}top\text{-}space}(s) \wedge \\ \quad u \subseteq \operatorname{pts}(s) \wedge \\ \quad \forall a \in u.\; \exists w \in \operatorname{opens}(s).\; ((a \in w) \wedge (w \subseteq u)) \Rightarrow \\ \quad \quad u \in \operatorname{opens}(s) \end{array}}}
\vnbprop{top-open-subset}{\vnbstmt{\forall s.; \forall v \in \operatorname{opens}(s).\; (\operatorname{is\text{-}top\text{-}space}(s) \Rightarrow (v \subseteq \operatorname{pts}(s)))}}
\vnbprop{top-group-parts-gt}{\vnbstmt{\begin{array}{@{}l@{}} \forall s, g. \\ \quad \operatorname{is\text{-}group\text{-}topology}(s, g) \Rightarrow  \\ \quad \quad \operatorname{is\text{-}top\text{-}space}(s) \wedge \\ \quad \quad \operatorname{is\text{-}group}(g) \wedge \\ \quad \quad \operatorname{pts}(s) = \operatorname{carr}(g) \wedge \\ \quad \quad \forall a, b \in \operatorname{carr}(g), w \in \operatorname{opens}(s). \\ \quad \quad \quad (\operatorname{opr}(g)(a, \operatorname{inv}(g)(b)) \in w) \Rightarrow  \\ \quad \quad \quad \quad \exists u \in \operatorname{opens}(s). \\ \quad \quad \quad \quad \quad a \in u \wedge \\ \quad \quad \quad \quad \quad \exists v \in \operatorname{opens}(s). \\ \quad \quad \quad \quad \quad \quad (b \in v) \wedge \forall p \in u, q \in v.\; (\operatorname{opr}(g)(p, \operatorname{inv}(g)(q)) \in w) \end{array}}}
\vnbprop{top-group-parts}{\vnbstmt{\begin{array}{@{}l@{}} \forall g. \\ \quad \operatorname{is\text{-}top\text{-}group}(g) \Rightarrow  \\ \quad \quad \operatorname{is\text{-}top\text{-}space}(\operatorname{top\text{-}group\text{-}space}(g)) \wedge \\ \quad \quad \operatorname{is\text{-}group}(\operatorname{top\text{-}group\text{-}group}(g)) \wedge \\ \quad \quad \operatorname{pts}(\operatorname{top\text{-}group\text{-}space}(g)) = \operatorname{carr}(\operatorname{top\text{-}group\text{-}group}(g)) \wedge \\ \quad \quad \forall a, b \in \operatorname{carr}(\operatorname{top\text{-}group\text{-}group}(g)), w \in \operatorname{opens}(\operatorname{top\text{-}group\text{-}space}(g)). \\ \quad \quad \quad (\operatorname{opr}(\operatorname{top\text{-}group\text{-}group}(g))(a, \operatorname{inv}(\operatorname{top\text{-}group\text{-}group}(g))(b)) \in w) \Rightarrow  \\ \quad \quad \quad \quad \exists u \in \operatorname{opens}(\operatorname{top\text{-}group\text{-}space}(g)). \\ \quad \quad \quad \quad \quad a \in u \wedge \\ \quad \quad \quad \quad \quad \exists v \in \operatorname{opens}(\operatorname{top\text{-}group\text{-}space}(g)). \\ \quad \quad \quad \quad \quad \quad ((b \in v) \wedge \forall p \in u, q \in v.\; (\operatorname{opr}(\operatorname{top\text{-}group\text{-}group}(g))(p, \operatorname{inv}(\operatorname{top\text{-}group\text{-}group}(g))(q)) \in w)) \end{array}}}
\vnbprop{top-group-inv-open-gt}{\vnbstmt{\begin{array}{@{}l@{}} \forall s, g.; \forall w \in \operatorname{opens}(s). \\ \quad \operatorname{is\text{-}group\text{-}topology}(s, g) \Rightarrow (\{z \in \operatorname{pts}(s) : (\operatorname{inv}(g)(z) \in w)\} \in \operatorname{opens}(s)) \end{array}}}
\vnbprop{top-group-inv-open}{\vnbstmt{\begin{array}{@{}l@{}} \forall g.; \forall w \in \operatorname{opens}(\operatorname{top\text{-}group\text{-}space}(g)). \\ \quad \operatorname{is\text{-}top\text{-}group}(g) \Rightarrow  \\ \quad \quad \{z \in \operatorname{pts}(\operatorname{top\text{-}group\text{-}space}(g)) : (\operatorname{inv}(\operatorname{top\text{-}group\text{-}group}(g))(z) \in w)\} \in \operatorname{opens}(\operatorname{top\text{-}group\text{-}space}(g)) \end{array}}}
\vnbprop{top-group-translate-open-gt}{\vnbstmt{\begin{array}{@{}l@{}} \forall s, g.; \forall a \in \operatorname{carr}(g), w \in \operatorname{opens}(s). \\ \quad \operatorname{is\text{-}group\text{-}topology}(s, g) \Rightarrow  \\ \quad \quad \{z \in \operatorname{pts}(s) : (\operatorname{opr}(g)(\operatorname{inv}(g)(a), z) \in w)\} \in \operatorname{opens}(s) \end{array}}}
\vnbprop{top-group-translate-open}{\vnbstmt{\begin{array}{@{}l@{}} \forall g.; \forall a \in \operatorname{carr}(\operatorname{top\text{-}group\text{-}group}(g)), w \in \operatorname{opens}(\operatorname{top\text{-}group\text{-}space}(g)). \\ \quad \operatorname{is\text{-}top\text{-}group}(g) \Rightarrow  \\ \quad \quad \{z \in \operatorname{pts}(\operatorname{top\text{-}group\text{-}space}(g)) : (\operatorname{opr}(\operatorname{top\text{-}group\text{-}group}(g))(\operatorname{inv}(\operatorname{top\text{-}group\text{-}group}(g))(a), z) \in w)\} \\ \quad \quad \in\; \operatorname{opens}(\operatorname{top\text{-}group\text{-}space}(g)) \end{array}}}
\vnbprop{top-group-sym-pair-gt}{\vnbstmt{\begin{array}{@{}l@{}} \forall s, g.; \forall w \in \operatorname{opens}(s). \\ \quad (\operatorname{is\text{-}group\text{-}topology}(s, g) \wedge (\operatorname{iden}(g) \in w)) \Rightarrow  \\ \quad \quad \exists v \in \operatorname{opens}(s). \\ \quad \quad \quad \operatorname{iden}(g) \in v \wedge \\ \quad \quad \quad \forall p \in v.\; (\operatorname{inv}(g)(p) \in v) \wedge \\ \quad \quad \quad \forall \mathit{pp}, q \in v.\; (\operatorname{opr}(g)(\mathit{pp}, q) \in w) \end{array}}}
\vnbprop{top-group-sym-pair}{\vnbstmt{\begin{array}{@{}l@{}} \forall g.; \forall w \in \operatorname{opens}(\operatorname{top\text{-}group\text{-}space}(g)). \\ \quad (\operatorname{is\text{-}top\text{-}group}(g) \wedge (\operatorname{iden}(\operatorname{top\text{-}group\text{-}group}(g)) \in w)) \Rightarrow  \\ \quad \quad \exists v \in \operatorname{opens}(\operatorname{top\text{-}group\text{-}space}(g)). \\ \quad \quad \quad \operatorname{iden}(\operatorname{top\text{-}group\text{-}group}(g)) \in v \wedge \\ \quad \quad \quad \forall p \in v.\; (\operatorname{inv}(\operatorname{top\text{-}group\text{-}group}(g))(p) \in v) \wedge \\ \quad \quad \quad \forall \mathit{pp}, q \in v.\; (\operatorname{opr}(\operatorname{top\text{-}group\text{-}group}(g))(\mathit{pp}, q) \in w) \end{array}}}
\vnbprop{top-group-sym-triple-gt}{\vnbstmt{\begin{array}{@{}l@{}} \forall s, g.; \forall w \in \operatorname{opens}(s). \\ \quad (\operatorname{is\text{-}group\text{-}topology}(s, g) \wedge (\operatorname{iden}(g) \in w)) \Rightarrow  \\ \quad \quad \exists v \in \operatorname{opens}(s). \\ \quad \quad \quad \operatorname{iden}(g) \in v \wedge \\ \quad \quad \quad \forall p \in v.\; (\operatorname{inv}(g)(p) \in v) \wedge \\ \quad \quad \quad \forall \mathit{tp}, q, r \in v.\; (\operatorname{opr}(g)(\operatorname{opr}(g)(\mathit{tp}, q), r) \in w) \end{array}}}
\vnbprop{top-group-sym-triple}{\vnbstmt{\begin{array}{@{}l@{}} \forall g.; \forall w \in \operatorname{opens}(\operatorname{top\text{-}group\text{-}space}(g)). \\ \quad (\operatorname{is\text{-}top\text{-}group}(g) \wedge (\operatorname{iden}(\operatorname{top\text{-}group\text{-}group}(g)) \in w)) \Rightarrow  \\ \quad \quad \exists v \in \operatorname{opens}(\operatorname{top\text{-}group\text{-}space}(g)). \\ \quad \quad \quad \operatorname{iden}(\operatorname{top\text{-}group\text{-}group}(g)) \in v \wedge \\ \quad \quad \quad \forall p \in v.\; (\operatorname{inv}(\operatorname{top\text{-}group\text{-}group}(g))(p) \in v) \wedge \\ \quad \quad \quad \forall \mathit{tp}, q, r \in v. \\ \quad \quad \quad \quad \operatorname{opr}(\operatorname{top\text{-}group\text{-}group}(g))(\operatorname{opr}(\operatorname{top\text{-}group\text{-}group}(g))(\mathit{tp}, q), r) \in w \end{array}}}
\vnbprop{top-group-nbhd-sequence-gt}{\vnbstmt{\begin{array}{@{}l@{}} \forall s, g.; \forall q \in (\mathbb{N} \to \operatorname{opens}(s)). \\ \quad (\operatorname{is\text{-}group\text{-}topology}(s, g) \wedge \forall k \in \mathbb{N}.\; (\operatorname{iden}(g) \in q(k))) \Rightarrow  \\ \quad \quad \exists a \in (\mathbb{N} \to \operatorname{opens}(s)). \\ \quad \quad \quad a(0) = \operatorname{pts}(s) \wedge \\ \quad \quad \quad \forall \mathit{qk} \in \mathbb{N}. \\ \quad \quad \quad \quad \operatorname{iden}(g) \in a(\mathit{qk}) \wedge \\ \quad \quad \quad \quad \forall p \in a(\mathit{qk}).\; (\operatorname{inv}(g)(p) \in a(\mathit{qk})) \wedge \\ \quad \quad \quad \quad \forall \mathit{tp}, \mathit{tq}, r \in a((\mathit{qk} + 1)). \\ \quad \quad \quad \quad \quad \operatorname{opr}(g)(\operatorname{opr}(g)(\mathit{tp}, \mathit{tq}), r) \in a(\mathit{qk}) \wedge \\ \quad \quad \quad \quad a((\mathit{qk} + 1)) \subseteq q(\mathit{qk}) \end{array}}}
\vnbprop{top-group-nbhd-sequence}{\vnbstmt{\begin{array}{@{}l@{}} \forall g.; \forall q \in (\mathbb{N} \to \operatorname{opens}(\operatorname{top\text{-}group\text{-}space}(g))). \\ \quad (\operatorname{is\text{-}top\text{-}group}(g) \wedge \forall k \in \mathbb{N}.\; (\operatorname{iden}(\operatorname{top\text{-}group\text{-}group}(g)) \in q(k))) \Rightarrow  \\ \quad \quad \exists s \in (\mathbb{N} \to \operatorname{opens}(\operatorname{top\text{-}group\text{-}space}(g))). \\ \quad \quad \quad s(0) = \operatorname{pts}(\operatorname{top\text{-}group\text{-}space}(g)) \wedge \\ \quad \quad \quad \forall \mathit{qk} \in \mathbb{N}. \\ \quad \quad \quad \quad \operatorname{iden}(\operatorname{top\text{-}group\text{-}group}(g)) \in s(\mathit{qk}) \wedge \\ \quad \quad \quad \quad \forall p \in s(\mathit{qk}).\; (\operatorname{inv}(\operatorname{top\text{-}group\text{-}group}(g))(p) \in s(\mathit{qk})) \wedge \\ \quad \quad \quad \quad \forall \mathit{tp}, \mathit{tq}, r \in s((\mathit{qk} + 1)). \\ \quad \quad \quad \quad \quad \operatorname{opr}(\operatorname{top\text{-}group\text{-}group}(g))(\operatorname{opr}(\operatorname{top\text{-}group\text{-}group}(g))(\mathit{tp}, \mathit{tq}), r) \in s(\mathit{qk}) \wedge \\ \quad \quad \quad \quad s((\mathit{qk} + 1)) \subseteq q(\mathit{qk}) \end{array}}}
\vnbprop{half-power-facts}{\vnbstmt{\begin{array}{@{}l@{}} \forall k \in \mathbb{N}. \\ \quad ((\operatorname{power}(1/2, k) \in \mathbb{R}) \wedge ((0 < \operatorname{power}(1/2, k)) \wedge (\operatorname{power}(1/2, (k + 1)) = (1/2 \cdot \operatorname{power}(1/2, k))))) \end{array}}}
\vnbprop{half-power-anti}{\vnbstmt{\forall j, k \in \mathbb{N}.\; ((j \leq k) \Rightarrow (\operatorname{power}(1/2, k) \leq \operatorname{power}(1/2, j)))}}
\vnbprop{half-power-small}{\vnbstmt{\forall a \in \mathbb{R}.\; (\forall k \in \mathbb{N}.\; (a \leq \operatorname{power}(1/2, k)) \Rightarrow (a \leq 0))}}
\vnbprop{bk-sequence-parts}{\vnbstmt{\begin{array}{@{}l@{}} \forall s, g, a. \\ \quad \operatorname{is\text{-}bk\text{-}sequence}(s, g, a) \Rightarrow  \\ \quad \quad a \in (\mathbb{N} \to \operatorname{opens}(s)) \wedge \\ \quad \quad a(0) = \operatorname{pts}(s) \wedge \\ \quad \quad \forall k \in \mathbb{N}. \\ \quad \quad \quad \operatorname{iden}(g) \in a(k) \wedge \\ \quad \quad \quad \forall p \in a(k).\; (\operatorname{inv}(g)(p) \in a(k)) \wedge \\ \quad \quad \quad \forall p, q, r \in a((k + 1)). \\ \quad \quad \quad \quad \operatorname{opr}(g)(\operatorname{opr}(g)(p, q), r) \in a(k) \end{array}}}
\vnbprop{bk-gauge-unfold}{\vnbstmt{\begin{array}{@{}l@{}} \forall s, x. \\ \quad (\operatorname{bk\text{-}gauge}(s, x) \simeq \operatorname{einf}(\{r \in \mathbb{R} : \exists k \in \mathbb{N}.\; ((x \in s(k)) \wedge (r = \operatorname{power}(1/2, k)))\})) \end{array}}}
\vnbprop{bk-sequence-decreasing}{\vnbstmt{\begin{array}{@{}l@{}} \forall s, g, a.; \forall k \in \mathbb{N}. \\ \quad (\operatorname{is\text{-}group\text{-}topology}(s, g) \wedge \operatorname{is\text{-}bk\text{-}sequence}(s, g, a)) \Rightarrow  \\ \quad \quad a((k + 1)) \subseteq a(k) \end{array}}}
\vnbprop{bk-sequence-nested}{\vnbstmt{\begin{array}{@{}l@{}} \forall s, g, a.; \forall k, j \in \mathbb{N}. \\ \quad (\operatorname{is\text{-}group\text{-}topology}(s, g) \wedge (\operatorname{is\text{-}bk\text{-}sequence}(s, g, a) \wedge (k \leq j))) \Rightarrow  \\ \quad \quad a(j) \subseteq a(k) \end{array}}}
\vnbprop{bk-gauge-set-nonneg}{\vnbstmt{\begin{array}{@{}l@{}} \forall s, x. \\ \quad \{r \in \mathbb{R} : \exists k \in \mathbb{N}.\; ((x \in s(k)) \wedge (r = \operatorname{power}(1/2, k)))\} \subseteq \operatorname{rr\text{-}pos\text{-}star} \end{array}}}
\vnbprop{bk-gauge-le}{\vnbstmt{\begin{array}{@{}l@{}} \forall s, x.; \forall k \in \mathbb{N}. \\ \quad (x \in s(k)) \Rightarrow  \\ \quad \quad (\operatorname{bk\text{-}gauge}(s, x) \in \mathbb{R}) \wedge ((0 \leq \operatorname{bk\text{-}gauge}(s, x)) \wedge (\operatorname{bk\text{-}gauge}(s, x) \leq \operatorname{power}(1/2, k))) \end{array}}}
\vnbprop{bk-last-nbhd}{\vnbstmt{\begin{array}{@{}l@{}} \forall s, g, a.; \forall x \in \operatorname{pts}(s), k \in \mathbb{N}. \\ \quad (\operatorname{is\text{-}group\text{-}topology}(s, g) \wedge (\operatorname{is\text{-}bk\text{-}sequence}(s, g, a) \wedge \neg (x \in a(k)))) \Rightarrow  \\ \quad \quad \exists q \in \mathbb{N}.\; ((x \in a(q)) \wedge (\neg (x \in a((q + 1))) \wedge ((q + 1) \leq k))) \end{array}}}
\vnbprop{bk-gauge-exact}{\vnbstmt{\begin{array}{@{}l@{}} \forall s, g, a.; \forall x \in \operatorname{pts}(s), q \in \mathbb{N}. \\ \quad (\operatorname{is\text{-}group\text{-}topology}(s, g) \wedge (\operatorname{is\text{-}bk\text{-}sequence}(s, g, a) \wedge ((x \in a(q)) \wedge \neg (x \in a((q + 1)))))) \Rightarrow  \\ \quad \quad \operatorname{bk\text{-}gauge}(a, x) = \operatorname{power}(1/2, q) \end{array}}}
\vnbprop{bk-gauge-char}{\vnbstmt{\begin{array}{@{}l@{}} \forall s, g, a.; \forall x \in \operatorname{pts}(s), k \in \mathbb{N}. \\ \quad (\operatorname{is\text{-}group\text{-}topology}(s, g) \wedge (\operatorname{is\text{-}bk\text{-}sequence}(s, g, a) \wedge (\operatorname{bk\text{-}gauge}(a, x) \leq \operatorname{power}(1/2, k)))) \Rightarrow  \\ \quad \quad x \in a(k) \end{array}}}
\vnbprop{bk-gauge-iden}{\vnbstmt{\begin{array}{@{}l@{}} \forall s, g, a. \\ \quad (\operatorname{is\text{-}group\text{-}topology}(s, g) \wedge \operatorname{is\text{-}bk\text{-}sequence}(s, g, a)) \Rightarrow  \\ \quad \quad \operatorname{bk\text{-}gauge}(a, \operatorname{iden}(g)) = 0 \end{array}}}
\vnbprop{bk-gauge-power}{\vnbstmt{\begin{array}{@{}l@{}} \forall s, g, a.; \forall x \in \operatorname{pts}(s). \\ \quad (\operatorname{is\text{-}group\text{-}topology}(s, g) \wedge (\operatorname{is\text{-}bk\text{-}sequence}(s, g, a) \wedge (0 < \operatorname{bk\text{-}gauge}(a, x)))) \Rightarrow  \\ \quad \quad \exists q \in \mathbb{N}.\; (\operatorname{bk\text{-}gauge}(a, x) = \operatorname{power}(1/2, q)) \end{array}}}
\vnbprop{bk-gauge-triple}{\vnbstmt{\begin{array}{@{}l@{}} \forall s, g, a.; \forall x, y, z \in \operatorname{carr}(g). \\ \quad (\operatorname{is\text{-}group\text{-}topology}(s, g) \wedge \operatorname{is\text{-}bk\text{-}sequence}(s, g, a)) \Rightarrow  \\ \quad \quad \operatorname{bk\text{-}gauge}(a, \operatorname{opr}(g)(\operatorname{opr}(g)(x, y), z)) \\ \quad \quad \leq\; 2 \cdot \operatorname{max}(\operatorname{bk\text{-}gauge}(a, x), \operatorname{max}(\operatorname{bk\text{-}gauge}(a, y), \operatorname{bk\text{-}gauge}(a, z))) \end{array}}}
\vnbprop{bk-f-unfold}{\vnbstmt{\begin{array}{@{}l@{}} \forall g, s, x. \\ \quad (\operatorname{bk\text{-}f}(g, s, x) \simeq \operatorname{subadd\text{-}reg}(\operatorname{grp\text{-}monoid}(g), (\lambda z \in \operatorname{carr}(g).\; \operatorname{bk\text{-}gauge}(s, z)), x)) \end{array}}}
\vnbprop{bk-f-bounds}{\vnbstmt{\begin{array}{@{}l@{}} \forall s, g, a.; \forall x \in \operatorname{carr}(g). \\ \quad (\operatorname{is\text{-}group\text{-}topology}(s, g) \wedge \operatorname{is\text{-}bk\text{-}sequence}(s, g, a)) \Rightarrow  \\ \quad \quad \operatorname{bk\text{-}f}(g, a, x) \in \mathbb{R} \wedge \\ \quad \quad (1/2 \cdot \operatorname{bk\text{-}gauge}(a, x)) \leq \operatorname{bk\text{-}f}(g, a, x) \wedge \\ \quad \quad \operatorname{bk\text{-}f}(g, a, x) \leq \operatorname{bk\text{-}gauge}(a, x) \end{array}}}
\vnbprop{bk-f-subadditive}{\vnbstmt{\begin{array}{@{}l@{}} \forall s, g, c.; \forall a, b \in \operatorname{carr}(g). \\ \quad (\operatorname{is\text{-}group\text{-}topology}(s, g) \wedge \operatorname{is\text{-}bk\text{-}sequence}(s, g, c)) \Rightarrow  \\ \quad \quad \operatorname{bk\text{-}f}(g, c, \operatorname{opr}(g)(a, b)) \leq (\operatorname{bk\text{-}f}(g, c, a) + \operatorname{bk\text{-}f}(g, c, b)) \end{array}}}
\vnbprop{bk-metric-lam-type}{\vnbstmt{\begin{array}{@{}l@{}} \forall s, g, a. \\ \quad (\operatorname{is\text{-}group\text{-}topology}(s, g) \wedge \operatorname{is\text{-}bk\text{-}sequence}(s, g, a)) \Rightarrow  \\ \quad \quad (\lambda {[}\operatorname{bmx\_}, \operatorname{bmy\_}] \in (\operatorname{carr}(g) \times \operatorname{carr}(g)).\; \\ \quad \quad \quad \operatorname{bk\text{-}f}(g, a, \operatorname{opr}(g)(\operatorname{inv}(g)(\operatorname{bmy\_}), \operatorname{bmx\_})) + \operatorname{bk\text{-}f}(g, a, \operatorname{opr}(g)(\operatorname{inv}(g)(\operatorname{bmx\_}), \operatorname{bmy\_}))) \\ \quad \quad \in\; ((\operatorname{carr}(g) \times \operatorname{carr}(g)) \to \mathbb{R}) \end{array}}}
\vnbprop{bk-metric-pts}{\vnbstmt{\begin{array}{@{}l@{}} \forall s, g, a. \\ \quad (\operatorname{is\text{-}group\text{-}topology}(s, g) \wedge \operatorname{is\text{-}bk\text{-}sequence}(s, g, a)) \Rightarrow  \\ \quad \quad (\operatorname{pts}(\operatorname{bk\text{-}metric}(g, a)) \simeq \operatorname{carr}(g)) \end{array}}}
\vnbprop{bk-metric-dist-fun}{\vnbstmt{\begin{array}{@{}l@{}} \forall s, g, a. \\ \quad (\operatorname{is\text{-}group\text{-}topology}(s, g) \wedge \operatorname{is\text{-}bk\text{-}sequence}(s, g, a)) \Rightarrow  \\ \quad \quad (\operatorname{dist}(\operatorname{bk\text{-}metric}(g, a)) \simeq (\lambda {[}\operatorname{bmx\_}, \operatorname{bmy\_}] \in (\operatorname{carr}(g) \times \operatorname{carr}(g)).\; (\operatorname{bk\text{-}f}(g, a, \operatorname{opr}(g)(\operatorname{inv}(g)(\operatorname{bmy\_}), \operatorname{bmx\_})) + \operatorname{bk\text{-}f}(g, a, \operatorname{opr}(g)(\operatorname{inv}(g)(\operatorname{bmx\_}), \operatorname{bmy\_}))))) \end{array}}}
\vnbprop{bk-metric-dist}{\vnbstmt{\begin{array}{@{}l@{}} \forall s, g, a.; \forall x, y \in \operatorname{carr}(g). \\ \quad (\operatorname{is\text{-}group\text{-}topology}(s, g) \wedge \operatorname{is\text{-}bk\text{-}sequence}(s, g, a)) \Rightarrow  \\ \quad \quad \operatorname{dist}(\operatorname{bk\text{-}metric}(g, a))(x, y) \\ \quad \quad =\; \operatorname{bk\text{-}f}(g, a, \operatorname{opr}(g)(\operatorname{inv}(g)(y), x)) + \operatorname{bk\text{-}f}(g, a, \operatorname{opr}(g)(\operatorname{inv}(g)(x), y)) \end{array}}}
\vnbprop{bk-f-iden}{\vnbstmt{\begin{array}{@{}l@{}} \forall s, g, a. \\ \quad (\operatorname{is\text{-}group\text{-}topology}(s, g) \wedge \operatorname{is\text{-}bk\text{-}sequence}(s, g, a)) \Rightarrow  \\ \quad \quad \operatorname{bk\text{-}f}(g, a, \operatorname{iden}(g)) = 0 \end{array}}}
\vnbprop{bk-metric-is-metric-space}{\vnbstmt{\begin{array}{@{}l@{}} \forall s, g, a. \\ \quad \operatorname{is\text{-}group\text{-}topology}(s, g) \wedge \\ \quad \operatorname{is\text{-}bk\text{-}sequence}(s, g, a) \wedge \\ \quad \forall z \in \operatorname{carr}(g).\; (\forall k \in \mathbb{N}.\; (z \in a(k)) \Rightarrow (z = \operatorname{iden}(g))) \Rightarrow \\ \quad \quad \operatorname{is\text{-}metric\text{-}space}(\operatorname{bk\text{-}metric}(g, a)) \end{array}}}
\vnbprop{bk-metric-left-invariant}{\vnbstmt{\begin{array}{@{}l@{}} \forall s, g, a. \\ \quad (\operatorname{is\text{-}group\text{-}topology}(s, g) \wedge \operatorname{is\text{-}bk\text{-}sequence}(s, g, a)) \Rightarrow  \\ \quad \quad \operatorname{is\text{-}left\text{-}invariant}(\operatorname{bk\text{-}metric}(g, a), g) \end{array}}}
\vnbprop{bk-ball-of-nbhd}{\vnbstmt{\begin{array}{@{}l@{}} \forall s, g, b.; \forall k \in \mathbb{N}, a, z \in \operatorname{carr}(g). \\ \quad (\operatorname{is\text{-}group\text{-}topology}(s, g) \wedge (\operatorname{is\text{-}bk\text{-}sequence}(s, g, b) \wedge (\operatorname{opr}(g)(\operatorname{inv}(g)(a), z) \in b(((k + 1) + 1))))) \Rightarrow  \\ \quad \quad \operatorname{dist}(\operatorname{bk\text{-}metric}(g, b))(a, z) \leq \operatorname{power}(1/2, (k + 1)) \end{array}}}
\vnbprop{bk-nbhd-of-ball}{\vnbstmt{\begin{array}{@{}l@{}} \forall s, g, b.; \forall k \in \mathbb{N}, a, z \in \operatorname{carr}(g). \\ \quad \operatorname{is\text{-}group\text{-}topology}(s, g) \wedge \\ \quad \operatorname{is\text{-}bk\text{-}sequence}(s, g, b) \wedge \\ \quad \operatorname{dist}(\operatorname{bk\text{-}metric}(g, b))(a, z) < \operatorname{power}(1/2, (k + 1)) \Rightarrow \\ \quad \quad \operatorname{opr}(g)(\operatorname{inv}(g)(a), z) \in b(k) \end{array}}}
\vnbprop{bk-metric-topology}{\vnbstmt{\begin{array}{@{}l@{}} \forall s, g, a, u. \\ \quad \operatorname{is\text{-}group\text{-}topology}(s, g) \wedge \\ \quad \operatorname{is\text{-}bk\text{-}sequence}(s, g, a) \wedge \\ \quad \forall z \in \operatorname{carr}(g).\; (\forall k \in \mathbb{N}.\; (z \in a(k)) \Rightarrow (z = \operatorname{iden}(g))) \wedge \\ \quad \forall w \in \operatorname{opens}(s). \\ \quad \quad (\operatorname{iden}(g) \in w) \Rightarrow \exists \mathit{fk} \in \mathbb{N}.\; (a(\mathit{fk}) \subseteq w) \Rightarrow \\ \quad \quad (u \in \operatorname{opens}(s)) \Leftrightarrow \operatorname{is\text{-}open}(\operatorname{bk\text{-}metric}(g, a), u) \end{array}}}
\vnbprop{birkhoff-kakutani-gt}{\vnbstmt{\begin{array}{@{}l@{}} \forall s, g. \\ \quad (\operatorname{is\text{-}group\text{-}topology}(s, g) \wedge (\operatorname{is\text{-}hausdorff}(s) \wedge \operatorname{has\text{-}countable\text{-}nbhd\text{-}base}(s, \operatorname{iden}(g)))) \Rightarrow  \\ \quad \quad \exists d. \\ \quad \quad \quad \operatorname{is\text{-}metric\text{-}space}(d) \wedge \\ \quad \quad \quad \operatorname{pts}(d) = \operatorname{pts}(s) \wedge \\ \quad \quad \quad \operatorname{is\text{-}left\text{-}invariant}(d, g) \wedge \\ \quad \quad \quad \operatorname{opens}(\operatorname{metric\text{-}top}(d)) = \operatorname{opens}(s) \end{array}}}
\vnbprop{birkhoff-kakutani}{\vnbstmt{\begin{array}{@{}l@{}} \forall g. \\ \quad \operatorname{is\text{-}top\text{-}group}(g) \wedge \\ \quad \operatorname{is\text{-}hausdorff}(\operatorname{top\text{-}group\text{-}space}(g)) \wedge \\ \quad \operatorname{has\text{-}countable\text{-}nbhd\text{-}base}(\operatorname{top\text{-}group\text{-}space}(g), \operatorname{iden}(\operatorname{top\text{-}group\text{-}group}(g))) \Rightarrow \\ \quad \quad \exists d. \\ \quad \quad \quad \operatorname{is\text{-}metric\text{-}space}(d) \wedge \\ \quad \quad \quad \operatorname{pts}(d) = \operatorname{pts}(\operatorname{top\text{-}group\text{-}space}(g)) \wedge \\ \quad \quad \quad \operatorname{is\text{-}left\text{-}invariant}(d, \operatorname{top\text{-}group\text{-}group}(g)) \wedge \\ \quad \quad \quad \operatorname{opens}(\operatorname{metric\text{-}top}(d)) = \operatorname{opens}(\operatorname{top\text{-}group\text{-}space}(g)) \end{array}}}
\vnbprop{bk-metric-pseudometric}{\vnbstmt{\begin{array}{@{}l@{}} \forall s, g, a. \\ \quad (\operatorname{is\text{-}group\text{-}topology}(s, g) \wedge \operatorname{is\text{-}bk\text{-}sequence}(s, g, a)) \Rightarrow  \\ \quad \quad \operatorname{is\text{-}pseudometric\text{-}space}(\operatorname{bk\text{-}metric}(g, a)) \end{array}}}
\vnbprop{bk-ball-open}{\vnbstmt{\begin{array}{@{}l@{}} \forall s, g, b.; \forall a \in \operatorname{carr}(g), r. \\ \quad (\operatorname{is\text{-}group\text{-}topology}(s, g) \wedge (\operatorname{is\text{-}bk\text{-}sequence}(s, g, b) \wedge \operatorname{pos\text{-}rr}(r))) \Rightarrow  \\ \quad \quad \operatorname{ball}(\operatorname{bk\text{-}metric}(g, b), a, r) \in \operatorname{opens}(s) \end{array}}}
\vnbprop{bk-metric-self}{\vnbstmt{\begin{array}{@{}l@{}} \forall s, g, b.; \forall a \in \operatorname{carr}(g). \\ \quad (\operatorname{is\text{-}group\text{-}topology}(s, g) \wedge \operatorname{is\text{-}bk\text{-}sequence}(s, g, b)) \Rightarrow  \\ \quad \quad \operatorname{dist}(\operatorname{bk\text{-}metric}(g, b))(a, a) = 0 \end{array}}}
\vnbprop{bk-ball-in-nbhd}{\vnbstmt{\begin{array}{@{}l@{}} \forall s, g, b.; \forall k \in \mathbb{N}, a \in \operatorname{carr}(g), z \in \operatorname{ball}(\operatorname{bk\text{-}metric}(g, b), a, \operatorname{power}(1/2, (k + 1))). \\ \quad (\operatorname{is\text{-}group\text{-}topology}(s, g) \wedge \operatorname{is\text{-}bk\text{-}sequence}(s, g, b)) \Rightarrow  \\ \quad \quad (z \in \operatorname{carr}(g)) \wedge (\operatorname{opr}(g)(\operatorname{inv}(g)(a), z) \in b(k)) \end{array}}}
\vnbprop{top-group-pseudometrics-define-topology-gt}{\vnbstmt{\begin{array}{@{}l@{}} \forall s, g, u. \\ \quad (\operatorname{is\text{-}group\text{-}topology}(s, g) \wedge (u \subseteq \operatorname{pts}(s))) \Rightarrow  \\ \quad \quad u \in \operatorname{opens}(s) \\ \quad \quad \Leftrightarrow\; \forall a \in u. \exists b. \\ \quad \quad \quad (\operatorname{is\text{-}bk\text{-}sequence}(s, g, b) \wedge \exists r.\; (\operatorname{pos\text{-}rr}(r) \wedge (\operatorname{ball}(\operatorname{bk\text{-}metric}(g, b), a, r) \subseteq u))) \end{array}}}
\vnbprop{top-group-pseudometrics-define-topology}{\vnbstmt{\begin{array}{@{}l@{}} \forall g, u. \\ \quad (\operatorname{is\text{-}top\text{-}group}(g) \wedge (u \subseteq \operatorname{pts}(\operatorname{top\text{-}group\text{-}space}(g)))) \Rightarrow  \\ \quad \quad u \in \operatorname{opens}(\operatorname{top\text{-}group\text{-}space}(g)) \\ \quad \quad \Leftrightarrow\; \forall a \in u. \exists s. \\ \quad \quad \quad \operatorname{is\text{-}bk\text{-}sequence}(\operatorname{top\text{-}group\text{-}space}(g), \operatorname{top\text{-}group\text{-}group}(g), s) \wedge \\ \quad \quad \quad \exists r. \\ \quad \quad \quad \quad \operatorname{pos\text{-}rr}(r) \wedge (\operatorname{ball}(\operatorname{bk\text{-}metric}(\operatorname{top\text{-}group\text{-}group}(g), s), a, r) \subseteq u) \end{array}}}

\vnbbookfile{theorem-library/top-product}{theorem-library/top-product}
\noindent{\small tych.scm -- THE PRODUCT TOPOLOGY OF AN ARBITRARY FAMILY, its projections and universal property, and TYCHONOFF'S THEOREM for arbitrary products, from zorn-lemma (Alexander's argument, specialised to the subbasic constraints of the product).  Agent TYCH (worker-08), 2026-10-26. DESTINATION: ...}\par
\vnbprop{tprod-pts-unfold}{\vnbstmt{\begin{array}{@{}l@{}} \forall x, f. \\ \quad (\operatorname{tprod\text{-}pts}(x, f) \simeq \{\mathit{tpx} \in (x \to \operatorname{big\text{-}union}(j, x, \operatorname{pts}(f(j)))) : \forall i \in x.\; (\operatorname{tpx}(i) \in \operatorname{pts}(f(i)))\}) \end{array}}}
\vnbprop{tprod-constr-unfold}{\vnbstmt{\begin{array}{@{}l@{}} \forall x, f. \\ \quad (\operatorname{tprod\text{-}constr}(x, f) \simeq \{b \in (x \times \operatorname{big\text{-}union}(j, x, \operatorname{opens}(f(j)))) : (\operatorname{nth}(2, b) \in \operatorname{opens}(f(\operatorname{nth}(1, b))))\}) \end{array}}}
\vnbprop{tprod-box-unfold}{\vnbstmt{\begin{array}{@{}l@{}} \forall x, f, b. \\ \quad (\operatorname{tprod\text{-}box}(x, f, b) \simeq \{y \in \operatorname{tprod\text{-}pts}(x, f) : \forall c \in b.\; (y(\operatorname{nth}(1, c)) \in \operatorname{nth}(2, c))\}) \end{array}}}
\vnbprop{tprod-pts-mem}{\vnbstmt{\begin{array}{@{}l@{}} \forall x, f, z. \\ \quad z \in \operatorname{tprod\text{-}pts}(x, f) \\ \quad \Leftrightarrow\; (z \in (x \to \operatorname{big\text{-}union}(j, x, \operatorname{pts}(f(j))))) \wedge \forall i \in x.\; (z(i) \in \operatorname{pts}(f(i))) \end{array}}}
\vnbprop{tprod-constr-mem}{\vnbstmt{\begin{array}{@{}l@{}} \forall x, f, z. \\ \quad z \in \operatorname{tprod\text{-}constr}(x, f) \\ \quad \Leftrightarrow\; ((z \in (x \times \operatorname{big\text{-}union}(j, x, \operatorname{opens}(f(j))))) \wedge (\operatorname{nth}(2, z) \in \operatorname{opens}(f(\operatorname{nth}(1, z))))) \end{array}}}
\vnbprop{tprod-box-mem}{\vnbstmt{\begin{array}{@{}l@{}} \forall x, f, b, z. \\ \quad (z \in \operatorname{tprod\text{-}box}(x, f, b)) \Leftrightarrow ((z \in \operatorname{tprod\text{-}pts}(x, f)) \wedge \forall c \in b.\; (z(\operatorname{nth}(1, c)) \in \operatorname{nth}(2, c))) \end{array}}}
\vnbprop{tprod-sets}{\vnbstmt{\begin{array}{@{}l@{}} \forall x, f. \\ \quad \operatorname{is\text{-}top\text{-}family}(x, f) \Rightarrow ((\operatorname{tprod\text{-}pts}(x, f) \in \mathbf{Set}) \wedge (\operatorname{tprod\text{-}constr}(x, f) \in \mathbf{Set})) \end{array}}}
\vnbprop{tprod-top-pts}{\vnbstmt{\begin{array}{@{}l@{}} \forall x, f. \\ \quad \operatorname{is\text{-}top\text{-}family}(x, f) \Rightarrow (\operatorname{pts}(\operatorname{tproduct\text{-}top}(x, f)) \simeq \operatorname{tprod\text{-}pts}(x, f)) \end{array}}}
\vnbprop{tprod-top-opens}{\vnbstmt{\begin{array}{@{}l@{}} \forall x, f. \\ \quad \operatorname{is\text{-}top\text{-}family}(x, f) \Rightarrow  \\ \quad \quad (\operatorname{opens}(\operatorname{tproduct\text{-}top}(x, f)) \simeq \{u \in \operatorname{power}(\operatorname{tprod\text{-}pts}(x, f)) : \operatorname{is\text{-}tprod\text{-}open}(x, f, u)\}) \end{array}}}
\vnbprop{tprod-box-anti}{\vnbstmt{\begin{array}{@{}l@{}} \forall x, f, a, b. \\ \quad (a \subseteq b) \Rightarrow (\operatorname{tprod\text{-}box}(x, f, b) \subseteq \operatorname{tprod\text{-}box}(x, f, a)) \end{array}}}
\vnbprop{card-nn-is-set}{\vnbstmt{\forall x.\; ((|x| \in \mathbb{N}) \Rightarrow (x \in \mathbf{Set}))}}
\vnbprop{is-tprod-open-classes}{\vnbstmt{\begin{array}{@{}l@{}} \forall x, f, u. \\ \quad \operatorname{is\text{-}tprod\text{-}open}(x, f, u) \\ \quad \Leftrightarrow\; u \subseteq \operatorname{tprod\text{-}pts}(x, f) \wedge \\ \quad \forall w \in u. \exists \mathit{tpf}. \\ \quad \quad \mathit{tpf} \subseteq \operatorname{tprod\text{-}constr}(x, f) \wedge \\ \quad \quad |\mathit{tpf}| \in \mathbb{N} \wedge \\ \quad \quad w \in \operatorname{tprod\text{-}box}(x, f, \mathit{tpf}) \wedge \\ \quad \quad \operatorname{tprod\text{-}box}(x, f, \mathit{tpf}) \subseteq u \end{array}}}
\vnbprop{tprod-open-empty}{\vnbstmt{\forall x, f.\; (\operatorname{is\text{-}top\text{-}family}(x, f) \Rightarrow \operatorname{is\text{-}tprod\text{-}open}(x, f, \emptyset))}}
\vnbprop{tprod-open-whole}{\vnbstmt{\forall x, f.\; (\operatorname{is\text{-}top\text{-}family}(x, f) \Rightarrow \operatorname{is\text{-}tprod\text{-}open}(x, f, \operatorname{tprod\text{-}pts}(x, f)))}}
\vnbprop{tprod-open-inter}{\vnbstmt{\begin{array}{@{}l@{}} \forall x, f, u, v. \\ \quad (\operatorname{is\text{-}top\text{-}family}(x, f) \wedge (\operatorname{is\text{-}tprod\text{-}open}(x, f, u) \wedge \operatorname{is\text{-}tprod\text{-}open}(x, f, v))) \Rightarrow  \\ \quad \quad \operatorname{is\text{-}tprod\text{-}open}(x, f, (u \cap v)) \end{array}}}
\vnbprop{tprod-open-union}{\vnbstmt{\begin{array}{@{}l@{}} \forall x, f, m. \\ \quad (\operatorname{is\text{-}top\text{-}family}(x, f) \wedge \forall o \in m.\; \operatorname{is\text{-}tprod\text{-}open}(x, f, o)) \Rightarrow  \\ \quad \quad \operatorname{is\text{-}tprod\text{-}open}(x, f, \operatorname{big\text{-}union}(u, m, u)) \end{array}}}
\vnbprop{tprod-top-is-top-space}{\vnbstmt{\forall x, f.\; (\operatorname{is\text{-}top\text{-}family}(x, f) \Rightarrow \operatorname{is\text{-}top\text{-}space}(\operatorname{tproduct\text{-}top}(x, f)))}}
\vnbprop{tprod-box-open}{\vnbstmt{\begin{array}{@{}l@{}} \forall x, f, b. \\ \quad (\operatorname{is\text{-}top\text{-}family}(x, f) \wedge ((b \subseteq \operatorname{tprod\text{-}constr}(x, f)) \wedge (|b| \in \mathbb{N}))) \Rightarrow  \\ \quad \quad \operatorname{tprod\text{-}box}(x, f, b) \in \operatorname{opens}(\operatorname{tproduct\text{-}top}(x, f)) \end{array}}}
\vnbprop{tproj-value}{\vnbstmt{\begin{array}{@{}l@{}} \forall x, f, i, y. \\ \quad (\operatorname{is\text{-}top\text{-}family}(x, f) \wedge ((i \in x) \wedge (y \in \operatorname{tprod\text{-}pts}(x, f)))) \Rightarrow  \\ \quad \quad \operatorname{tproj}(x, f, i)(y) = y(i) \end{array}}}
\vnbprop{tproj-continuous}{\vnbstmt{\begin{array}{@{}l@{}} \forall x, f, i. \\ \quad (\operatorname{is\text{-}top\text{-}family}(x, f) \wedge (i \in x)) \Rightarrow  \\ \quad \quad \operatorname{is\text{-}hom\text{-}top\text{-}space}(\operatorname{tproduct\text{-}top}(x, f), f(i), \operatorname{tproj}(x, f, i)) \end{array}}}
\vnbprop{tprod-coord-nbhd}{\vnbstmt{\begin{array}{@{}l@{}} \forall x, b, z, f, a, c. \\ \quad \operatorname{is\text{-}top\text{-}family}(x, b) \wedge \\ \quad \operatorname{is\text{-}top\text{-}space}(z) \wedge \\ \quad f \in (\operatorname{pts}(z) \to \operatorname{tprod\text{-}pts}(x, b)) \wedge \\ \quad \forall k \in x.\; \operatorname{is\text{-}hom\text{-}top\text{-}space}(z, b(k), (\lambda q \in \operatorname{pts}(z).\; f(q)(k))) \wedge \\ \quad a \in \operatorname{pts}(z) \wedge \\ \quad c \subseteq \operatorname{tprod\text{-}constr}(x, b) \wedge \\ \quad c \in \mathbf{Set} \wedge \\ \quad |c| \in \mathbb{N} \wedge \\ \quad \forall \mathit{tpc} \in c.\; (f(a)(\operatorname{nth}(1, \mathit{tpc})) \in \operatorname{nth}(2, \mathit{tpc})) \Rightarrow \\ \quad \quad \exists w \in \operatorname{opens}(z). \\ \quad \quad \quad (a \in w) \wedge \forall p \in w, \mathit{tpc} \in c.\; (f(p)(\operatorname{nth}(1, \mathit{tpc})) \in \operatorname{nth}(2, \mathit{tpc})) \end{array}}}
\vnbprop{tprod-hom-iff}{\vnbstmt{\begin{array}{@{}l@{}} \forall x, a, z, f. \\ \quad (\operatorname{is\text{-}top\text{-}family}(x, a) \wedge (\operatorname{is\text{-}top\text{-}space}(z) \wedge (f \in (\operatorname{pts}(z) \to \operatorname{tprod\text{-}pts}(x, a))))) \Rightarrow  \\ \quad \quad \operatorname{is\text{-}hom\text{-}top\text{-}space}(z, \operatorname{tproduct\text{-}top}(x, a), f) \\ \quad \quad \Leftrightarrow\; \forall k \in x.\; \operatorname{is\text{-}hom\text{-}top\text{-}space}(z, a(k), (\lambda q \in \operatorname{pts}(z).\; f(q)(k))) \end{array}}}
\vnbprop{hom-top-space-ptwise-eq}{\vnbstmt{\begin{array}{@{}l@{}} \forall s, t, g, h. \\ \quad (\operatorname{is\text{-}hom\text{-}top\text{-}space}(s, t, g) \wedge ((h \in (\operatorname{pts}(s) \to \operatorname{pts}(t))) \wedge \forall q \in \operatorname{pts}(s).\; (h(q) = g(q)))) \Rightarrow  \\ \quad \quad \operatorname{is\text{-}hom\text{-}top\text{-}space}(s, t, h) \end{array}}}
\vnbprop{tprod-induced-map}{\vnbstmt{\begin{array}{@{}l@{}} \forall x, f, z, a. \\ \quad \operatorname{is\text{-}top\text{-}family}(x, f) \wedge \\ \quad \operatorname{is\text{-}top\text{-}space}(z) \wedge \\ \quad \forall k \in x.\; (a(k) \in (\operatorname{pts}(z) \to \operatorname{pts}(f(k)))) \wedge \\ \quad \forall k \in x.\; \operatorname{is\text{-}hom\text{-}top\text{-}space}(z, f(k), a(k)) \Rightarrow \\ \quad \quad \operatorname{is\text{-}hom\text{-}top\text{-}space}(z, \operatorname{tproduct\text{-}top}(x, f), (\lambda q \in \operatorname{pts}(z).\; (\lambda m \in x.\; a(m)(q)))) \wedge \\ \quad \quad \forall k \in x, p \in \operatorname{pts}(z). \\ \quad \quad \quad \operatorname{tproj}(x, f, k)((\lambda q \in \operatorname{pts}(z).\; (\lambda m \in x.\; a(m)(q)))(p)) = a(k)(p) \end{array}}}
\vnbprop{tprod-induced-map-unique}{\vnbstmt{\begin{array}{@{}l@{}} \forall x, f, z, a, h. \\ \quad \operatorname{is\text{-}top\text{-}family}(x, f) \wedge \\ \quad \operatorname{is\text{-}top\text{-}space}(z) \wedge \\ \quad \forall k \in x.\; (a(k) \in (\operatorname{pts}(z) \to \operatorname{pts}(f(k)))) \wedge \\ \quad h \in (\operatorname{pts}(z) \to \operatorname{tprod\text{-}pts}(x, f)) \wedge \\ \quad \forall k \in x, p \in \operatorname{pts}(z).\; (\operatorname{tproj}(x, f, k)(h(p)) = a(k)(p)) \Rightarrow \\ \quad \quad h = (\lambda q \in \operatorname{pts}(z).\; (\lambda m \in x.\; a(m)(q))) \end{array}}}
\vnbprop{chain-finite-in-member}{\vnbstmt{\begin{array}{@{}l@{}} \forall h.; \forall f \in \mathbf{Set}. \\ \quad \exists \mathit{tph}.\; (\mathit{tph} \in h) \wedge \\ \quad \forall a, b \in h.\; ((a \subseteq b) \vee (b \subseteq a)) \wedge \\ \quad |f| \in \mathbb{N} \wedge \\ \quad f \subseteq \operatorname{big\text{-}union}(v, h, v) \Rightarrow \\ \quad \quad \exists d \in h.\; (f \subseteq d) \end{array}}}
\vnbprop{has-finite-subcover-classes}{\vnbstmt{\begin{array}{@{}l@{}} \forall s, c. \\ \quad \operatorname{has\text{-}finite\text{-}subcover}(s, c) \\ \quad \Leftrightarrow\; \exists g.\; ((g \subseteq c) \wedge ((|g| \in \mathbb{N}) \wedge (\operatorname{big\text{-}union}(v, g, v) \simeq \operatorname{pts}(s)))) \end{array}}}
\vnbprop{alexander-maximal-cover}{\vnbstmt{\begin{array}{@{}l@{}} \forall s, a. \\ \quad (\operatorname{is\text{-}top\text{-}space}(s) \wedge (\forall u \in a.\; (u \in \operatorname{opens}(s)) \wedge \neg \operatorname{has\text{-}finite\text{-}subcover}(s, a))) \Rightarrow  \\ \quad \quad \exists m. \\ \quad \quad \quad a \subseteq m \wedge \\ \quad \quad \quad m \subseteq \operatorname{opens}(s) \wedge \\ \quad \quad \quad \neg \operatorname{has\text{-}finite\text{-}subcover}(s, m) \wedge \\ \quad \quad \quad \forall w \in \operatorname{opens}(s). \\ \quad \quad \quad \quad \neg (w \in m) \Rightarrow \operatorname{has\text{-}finite\text{-}subcover}(s, (m \cup \{w\})) \end{array}}}
\vnbprop{finite-cover-outside}{\vnbstmt{\begin{array}{@{}l@{}} \forall s, m, w. \\ \quad (\operatorname{is\text{-}top\text{-}space}(s) \wedge ((m \subseteq \operatorname{opens}(s)) \wedge ((w \in \mathbf{Set}) \wedge \operatorname{has\text{-}finite\text{-}subcover}(s, (m \cup \{w\}))))) \Rightarrow  \\ \quad \quad \exists k_{2}. \\ \quad \quad \quad k_{2} \subseteq m \wedge \\ \quad \quad \quad |k_{2}| \in \mathbb{N} \wedge \\ \quad \quad \quad \forall y \in \operatorname{pts}(s).\; (\neg (y \in w) \Rightarrow (y \in \operatorname{big\text{-}union}(v, k_{2}, v))) \end{array}}}
\vnbprop{tychonoff-factor-point}{\vnbstmt{\begin{array}{@{}l@{}} \forall x, f, m, i. \\ \quad \operatorname{is\text{-}top\text{-}family}(x, f) \wedge \\ \quad \forall \mathit{tpi} \in x.\; \operatorname{is\text{-}compact\text{-}t}(f(\mathit{tpi})) \wedge \\ \quad m \subseteq \operatorname{opens}(\operatorname{tproduct\text{-}top}(x, f)) \wedge \\ \quad \neg \operatorname{has\text{-}finite\text{-}subcover}(\operatorname{tproduct\text{-}top}(x, f), m) \wedge \\ \quad i \in x \Rightarrow \\ \quad \quad \exists e \in \operatorname{pts}(f(i)). \forall v_{5} \in \operatorname{opens}(f(i)). \\ \quad \quad \quad (\operatorname{tprod\text{-}box}(x, f, \{{[}i, v_{5}]\}) \in m) \Rightarrow \neg (e \in v_{5}) \end{array}}}
\vnbprop{tprod-finite-constraint-cover}{\vnbstmt{\begin{array}{@{}l@{}} \forall x, f, m, b. \\ \quad \operatorname{is\text{-}top\text{-}family}(x, f) \wedge \\ \quad m \subseteq \operatorname{opens}(\operatorname{tproduct\text{-}top}(x, f)) \wedge \\ \quad b \subseteq \operatorname{tprod\text{-}constr}(x, f) \wedge \\ \quad b \in \mathbf{Set} \wedge \\ \quad |b| \in \mathbb{N} \wedge \\ \quad \forall c_{8} \in b. \exists k_{2}. \\ \quad \quad k_{2} \subseteq m \wedge \\ \quad \quad |k_{2}| \in \mathbb{N} \wedge \\ \quad \quad \forall y \in \operatorname{pts}(\operatorname{tproduct\text{-}top}(x, f)). \\ \quad \quad \quad \neg (y \in \operatorname{tprod\text{-}box}(x, f, \{c_{8}\})) \Rightarrow (y \in \operatorname{big\text{-}union}(v, k_{2}, v)) \Rightarrow \\ \quad \quad \exists k_{3} \in \operatorname{power}(m). \\ \quad \quad \quad |k_{3}| \in \mathbb{N} \wedge \\ \quad \quad \quad \forall y \in \operatorname{pts}(\operatorname{tproduct\text{-}top}(x, f)). \\ \quad \quad \quad \quad \neg (y \in \operatorname{tprod\text{-}box}(x, f, b)) \Rightarrow (y \in \operatorname{big\text{-}union}(v, k_{3}, v)) \end{array}}}
\vnbprop{tychonoff}{\vnbstmt{\begin{array}{@{}l@{}} \forall x, f. \\ \quad (\operatorname{is\text{-}top\text{-}family}(x, f) \wedge \forall i \in x.\; \operatorname{is\text{-}compact\text{-}t}(f(i))) \Rightarrow  \\ \quad \quad \operatorname{is\text{-}compact\text{-}t}(\operatorname{tproduct\text{-}top}(x, f)) \end{array}}}
\vnbprop{compact-t-continuous-onto}{\vnbstmt{\begin{array}{@{}l@{}} \forall s, t, f. \\ \quad \operatorname{is\text{-}hom\text{-}top\text{-}space}(s, t, f) \wedge \\ \quad \operatorname{is\text{-}compact\text{-}t}(s) \wedge \\ \quad \forall y \in \operatorname{pts}(t).\; \exists x \in \operatorname{pts}(s).\; (f(x) = y) \Rightarrow \\ \quad \quad \operatorname{is\text{-}compact\text{-}t}(t) \end{array}}}

\vnbbookfile{structure-library/cg1-vocab}{structure-library/cg1-vocab}
\noindent{\small cg1-vocab.scm -- vocabulary for the compact-groups notes, chapter 5 section 2 (CG-1). DEFINITIONS ONLY.  Agent ISOM, 2026-10-26 (new definitions, the user's to confirm).}\par
\par\medskip\noindent\textbf{Definition} (\texttt{is-connected-t}).\;For all $\operatorname{cts\_}$ and $\operatorname{cta\_}$, 
\begin{equation*}
\begin{array}{@{}l@{}} \operatorname{is\text{-}connected\text{-}t}(\operatorname{cts\_}, \operatorname{cta\_}) \\ \Leftrightarrow\; \operatorname{cta\_} \subseteq \operatorname{pts}(\operatorname{cts\_}) \wedge \\ \forall \operatorname{ctu\_}, \operatorname{ctv\_} \in \operatorname{opens}(\operatorname{cts\_}). \\ \quad (\operatorname{cta\_} \subseteq (\operatorname{ctu\_} \cup \operatorname{ctv\_})) \Rightarrow  \\ \quad \forall \operatorname{ctx\_} \in \operatorname{cta\_}.\; \neg ((\operatorname{ctx\_} \in \operatorname{ctu\_}) \wedge (\operatorname{ctx\_} \in \operatorname{ctv\_})) \Rightarrow  \\ \quad \exists \operatorname{ctx\_} \in \operatorname{cta\_}.\; (\operatorname{ctx\_} \in \operatorname{ctu\_}) \Rightarrow  \\ \quad \quad \neg \exists \operatorname{ctx\_} \in \operatorname{cta\_}.\; (\operatorname{ctx\_} \in \operatorname{ctv\_}) \end{array}
\end{equation*}

\par\medskip\noindent\textbf{Definition} (\texttt{connected-component}).\;$\operatorname{connected\text{-}component}(\operatorname{ccs\_}, \operatorname{ccx\_}) :=$
\begin{equation*}
\operatorname{big\text{-}union}(\operatorname{ccc\_}, \{\operatorname{ccc\_} \in \operatorname{power}(\operatorname{pts}(\operatorname{ccs\_})) : (\operatorname{is\text{-}connected\text{-}t}(\operatorname{ccs\_}, \operatorname{ccc\_}) \wedge (\operatorname{ccx\_} \in \operatorname{ccc\_}))\}, \operatorname{ccc\_})
\end{equation*}


\vnbbookfile{theorem-library/frechet-metric}{theorem-library/frechet-metric}
\noindent{\small frechet-metric.scm -- a seminormed space with a COUNTABLE separating family of seminorms as a METRIC space: Thayer, "Analise Funcional e Teoria Espectral", Ch. 8 s.4, Props. 8.24-8.26 (pp. 254-256).  Batch TVS-1, 2026-10-10. Vocabulary: structure-library/frechet-metric.scm (FR-C-SPACE, FR-METRIC-SPACE,}\par
\vnbprop{frm-c-space-pts}{\vnbstmt{\forall m, a.\; ((\operatorname{vec}(m) \in \mathbf{Set}) \Rightarrow (\operatorname{pts}(\operatorname{fr\text{-}c\text{-}space}(m, a)) \simeq \operatorname{vec}(m)))}}
\vnbprop{frm-c-space-dists}{\vnbstmt{\begin{array}{@{}l@{}} \forall m, a. \\ \quad (\operatorname{vec}(m) \in \mathbf{Set}) \Rightarrow  \\ \quad \quad (\operatorname{dists}(\operatorname{fr\text{-}c\text{-}space}(m, a)) \simeq (\lambda k \in \mathbb{N}.\; (\lambda {[}\operatorname{fru\_}, \operatorname{frv\_}] \in (\operatorname{vec}(m) \times \operatorname{vec}(m)).\; a(k)(\operatorname{vadd}(m)(\operatorname{fru\_}, \operatorname{vneg}(m)(\operatorname{frv\_})))))) \end{array}}}
\vnbprop{frm-c-space-dists-apply}{\vnbstmt{\begin{array}{@{}l@{}} \forall m, a, k, x, y. \\ \quad ((\operatorname{vec}(m) \in \mathbf{Set}) \wedge ((k \in \mathbb{N}) \wedge ((x \in \operatorname{vec}(m)) \wedge (y \in \operatorname{vec}(m))))) \Rightarrow  \\ \quad \quad (\operatorname{dists}(\operatorname{fr\text{-}c\text{-}space}(m, a))(k)(x, y) \simeq a(k)(\operatorname{vadd}(m)(x, \operatorname{vneg}(m)(y)))) \end{array}}}
\vnbprop{frm-sn-module}{\vnbstmt{\forall d, m, p.\; (\operatorname{is\text{-}seminorm}(d, m, p) \Rightarrow \operatorname{is\text{-}module}(m))}}
\vnbprop{frm-sn-type}{\vnbstmt{\forall d, m, p.\; (\operatorname{is\text{-}seminorm}(d, m, p) \Rightarrow (p \in (\operatorname{vec}(m) \to \mathbb{R})))}}
\vnbprop{frm-sn-nonneg}{\vnbstmt{\forall d, m, p.; \forall x \in \operatorname{vec}(m).\; (\operatorname{is\text{-}seminorm}(d, m, p) \Rightarrow (0 \leq p(x)))}}
\vnbprop{frm-sn-subadd}{\vnbstmt{\begin{array}{@{}l@{}} \forall d, m, p.; \forall x, y \in \operatorname{vec}(m). \\ \quad \operatorname{is\text{-}seminorm}(d, m, p) \Rightarrow (p(\operatorname{vadd}(m)(x, y)) \leq (p(x) + p(y))) \end{array}}}
\vnbprop{frm-sn-zero}{\vnbstmt{\forall d, m, p.\; (\operatorname{is\text{-}seminorm}(d, m, p) \Rightarrow (p(\operatorname{vzero}(m)) = 0))}}
\vnbprop{frm-sn-neg}{\vnbstmt{\begin{array}{@{}l@{}} \forall d, m, p.; \forall x \in \operatorname{vec}(m). \\ \quad \operatorname{is\text{-}seminorm}(d, m, p) \Rightarrow (p(\operatorname{vneg}(m)(x)) = p(x)) \end{array}}}
\vnbprop{frm-fam-module}{\vnbstmt{\forall d, m, a.\; (\operatorname{is\text{-}seminorm\text{-}family}(d, m, a) \Rightarrow \operatorname{is\text{-}module}(m))}}
\vnbprop{frm-fam-type}{\vnbstmt{\begin{array}{@{}l@{}} \forall d, m, a. \\ \quad \operatorname{is\text{-}seminorm\text{-}family}(d, m, a) \Rightarrow (a \in (\mathbb{N} \to (\operatorname{vec}(m) \to \mathbb{R}))) \end{array}}}
\vnbprop{frm-fam-seminorm}{\vnbstmt{\begin{array}{@{}l@{}} \forall d, m, a.; \forall k \in \mathbb{N}. \\ \quad \operatorname{is\text{-}seminorm\text{-}family}(d, m, a) \Rightarrow \operatorname{is\text{-}seminorm}(d, m, a(k)) \end{array}}}
\vnbprop{frm-fam-separates}{\vnbstmt{\begin{array}{@{}l@{}} \forall d, m, a.; \forall x \in \operatorname{vec}(m). \\ \quad (\operatorname{is\text{-}seminorm\text{-}family}(d, m, a) \wedge \forall k \in \mathbb{N}.\; (a(k)(x) = 0)) \Rightarrow  \\ \quad \quad x = \operatorname{vzero}(m) \end{array}}}
\vnbprop{frm-c-space-is-c-metric-space}{\vnbstmt{\begin{array}{@{}l@{}} \forall d, m, a. \\ \quad \operatorname{is\text{-}seminorm\text{-}family}(d, m, a) \Rightarrow \operatorname{is\text{-}c\text{-}metric\text{-}space}(\operatorname{fr\text{-}c\text{-}space}(m, a)) \end{array}}}
\vnbprop{frm-metric-space-is-metric-space}{\vnbstmt{\begin{array}{@{}l@{}} \forall d, m, a. \\ \quad \operatorname{is\text{-}seminorm\text{-}family}(d, m, a) \Rightarrow \operatorname{is\text{-}metric\text{-}space}(\operatorname{fr\text{-}metric\text{-}space}(m, a)) \end{array}}}
\vnbprop{frm-metric-space-pts}{\vnbstmt{\begin{array}{@{}l@{}} \forall d, m, a. \\ \quad \operatorname{is\text{-}seminorm\text{-}family}(d, m, a) \Rightarrow (\operatorname{pts}(\operatorname{fr\text{-}metric\text{-}space}(m, a)) \simeq \operatorname{vec}(m)) \end{array}}}
\vnbprop{frm-c-ball-prefix}{\vnbstmt{\begin{array}{@{}l@{}} \forall s.; \forall k \in \mathbb{N}, a. \\ \quad (\operatorname{is\text{-}c\text{-}metric\text{-}space}(s) \wedge \operatorname{pos\text{-}rr}(a)) \Rightarrow  \\ \quad \quad \exists l. \\ \quad \quad \quad \operatorname{pos\text{-}rr}(l) \wedge \\ \quad \quad \quad \forall x, y \in \operatorname{pts}(s). \\ \quad \quad \quad \quad (\operatorname{dist}(\operatorname{c\text{-}metric}(s))(x, y) < l) \Rightarrow (\operatorname{dists}(s)(k)(x, y) < a) \end{array}}}
\vnbprop{frm-c-head-small}{\vnbstmt{\begin{array}{@{}l@{}} \forall s.; \forall x, y \in \operatorname{pts}(s), b \in \mathbb{N}, e \in \mathbb{R}. \\ \quad (\operatorname{is\text{-}c\text{-}metric\text{-}space}(s) \wedge ((0 \leq e) \wedge \forall j \in \mathbb{N}.\; ((j \leq b) \Rightarrow (\operatorname{dists}(s)(j)(x, y) < e)))) \Rightarrow  \\ \quad \quad \operatorname{series\text{-}partial\text{-}sum}((\lambda k \in \mathbb{N}.\; ((\lambda n \in \mathbb{N}.\; /(1, \operatorname{power}(2, (n + 1))))(k) \cdot \operatorname{min}(1, \operatorname{dists}(s)(k)(x, y)))), \\ \quad \quad \quad b) \\ \quad \quad \leq\; e \end{array}}}
\vnbprop{frm-c-term-bounds}{\vnbstmt{\begin{array}{@{}l@{}} \forall s, x, y.; \forall i \in \mathbb{N}. \\ \quad (\operatorname{is\text{-}c\text{-}metric\text{-}space}(s) \wedge ((x \in \operatorname{pts}(s)) \wedge (y \in \operatorname{pts}(s)))) \Rightarrow  \\ \quad \quad 0 \\ \quad \quad \leq\; (\lambda k \in \mathbb{N}.\; ((\lambda n \in \mathbb{N}.\; /(1, \operatorname{power}(2, (n + 1))))(k) \cdot \operatorname{min}(1, \operatorname{dists}(s)(k)(x, y))))(i) \wedge \\ \quad \quad (\lambda k \in \mathbb{N}.\; ((\lambda n \in \mathbb{N}.\; /(1, \operatorname{power}(2, (n + 1))))(k) \cdot \operatorname{min}(1, \operatorname{dists}(s)(k)(x, y))))(i) \\ \quad \quad \leq\; (\lambda n \in \mathbb{N}.\; /(1, \operatorname{power}(2, (n + 1))))(i) \end{array}}}
\vnbprop{frm-c-prefix-ball}{\vnbstmt{\begin{array}{@{}l@{}} \forall s, a. \\ \quad (\operatorname{is\text{-}c\text{-}metric\text{-}space}(s) \wedge \operatorname{pos\text{-}rr}(a)) \Rightarrow  \\ \quad \quad \exists b \in \mathbb{N}, l. \\ \quad \quad \quad \operatorname{pos\text{-}rr}(l) \wedge \\ \quad \quad \quad \forall x, y \in \operatorname{pts}(s). \\ \quad \quad \quad \quad \forall j \in \mathbb{N}.\; ((j \leq b) \Rightarrow (\operatorname{dists}(s)(j)(x, y) < l)) \Rightarrow  \\ \quad \quad \quad \quad \quad \operatorname{dist}(\operatorname{c\text{-}metric}(s))(x, y) \leq a \end{array}}}
\vnbprop{frm-ball-seminorm}{\vnbstmt{\begin{array}{@{}l@{}} \forall d, m, a.; \forall k \in \mathbb{N}, s. \\ \quad (\operatorname{is\text{-}seminorm\text{-}family}(d, m, a) \wedge \operatorname{pos\text{-}rr}(s)) \Rightarrow  \\ \quad \quad \exists l. \\ \quad \quad \quad \operatorname{pos\text{-}rr}(l) \wedge \\ \quad \quad \quad \forall x, y. \\ \quad \quad \quad \quad ((x \in \operatorname{vec}(m)) \wedge ((y \in \operatorname{vec}(m)) \wedge (\operatorname{dist}(\operatorname{fr\text{-}metric\text{-}space}(m, a))(x, y) < l))) \Rightarrow  \\ \quad \quad \quad \quad \quad a(k)(\operatorname{vadd}(m)(x, \operatorname{vneg}(m)(y))) < s \end{array}}}
\vnbprop{frm-seminorm-ball}{\vnbstmt{\begin{array}{@{}l@{}} \forall d, m, a, s. \\ \quad (\operatorname{is\text{-}seminorm\text{-}family}(d, m, a) \wedge \operatorname{pos\text{-}rr}(s)) \Rightarrow  \\ \quad \quad \exists b \in \mathbb{N}, l. \\ \quad \quad \quad \operatorname{pos\text{-}rr}(l) \wedge \\ \quad \quad \quad \forall x, y. \\ \quad \quad \quad \quad x \in \operatorname{vec}(m) \wedge \\ \quad \quad \quad \quad y \in \operatorname{vec}(m) \wedge \\ \quad \quad \quad \quad \forall j \in \mathbb{N}.\; ((j \leq b) \Rightarrow (a(j)(\operatorname{vadd}(m)(x, \operatorname{vneg}(m)(y))) < l)) \Rightarrow \\ \quad \quad \quad \quad \quad \operatorname{dist}(\operatorname{fr\text{-}metric\text{-}space}(m, a))(x, y) \leq s \end{array}}}
\vnbprop{frm-prefix-eventually}{\vnbstmt{\begin{array}{@{}l@{}} \forall b \in \mathbb{N}, u, l. \\ \quad \forall j \in \mathbb{N}. \exists p \in \mathbb{N}. \forall i, a \in \mathbb{N}. \\ \quad \quad ((p \leq i) \wedge (p \leq a)) \Rightarrow (u(j)(i)(a) < l) \Rightarrow \\ \quad \quad \exists t \in \mathbb{N}. \forall j, i, a \in \mathbb{N}. \\ \quad \quad \quad ((j \leq b) \wedge ((t \leq i) \wedge (t \leq a))) \Rightarrow (u(j)(i)(a) < l) \end{array}}}
\vnbprop{frm-conv-of-converges}{\vnbstmt{\begin{array}{@{}l@{}} \forall d, m, a, q, b. \\ \quad \operatorname{is\text{-}seminorm\text{-}family}(d, m, a) \wedge \\ \quad q \in (\mathbb{N} \to \operatorname{vec}(m)) \wedge \\ \quad b \in \operatorname{vec}(m) \wedge \\ \quad \operatorname{converges\text{-}to}(\operatorname{fr\text{-}metric\text{-}space}(m, a), q, b) \Rightarrow \\ \quad \quad \operatorname{fr\text{-}conv}(m, a, q, b) \end{array}}}
\vnbprop{frm-converges-of-conv}{\vnbstmt{\begin{array}{@{}l@{}} \forall d, m, a, q, b. \\ \quad (\operatorname{is\text{-}seminorm\text{-}family}(d, m, a) \wedge ((q \in (\mathbb{N} \to \operatorname{vec}(m))) \wedge ((b \in \operatorname{vec}(m)) \wedge \operatorname{fr\text{-}conv}(m, a, q, b)))) \Rightarrow  \\ \quad \quad \operatorname{converges\text{-}to}(\operatorname{fr\text{-}metric\text{-}space}(m, a), q, b) \end{array}}}
\vnbprop{frm-cauchy-of-metric}{\vnbstmt{\begin{array}{@{}l@{}} \forall d, m, a.; \forall q \in (\mathbb{N} \to \operatorname{vec}(m)). \\ \quad (\operatorname{is\text{-}seminorm\text{-}family}(d, m, a) \wedge \operatorname{is\text{-}cauchy\text{-}seq}(\operatorname{fr\text{-}metric\text{-}space}(m, a), q)) \Rightarrow  \\ \quad \quad \operatorname{fr\text{-}cauchy}(m, a, q) \end{array}}}
\vnbprop{frm-metric-cauchy-of-cauchy}{\vnbstmt{\begin{array}{@{}l@{}} \forall d, m, a.; \forall q \in (\mathbb{N} \to \operatorname{vec}(m)). \\ \quad (\operatorname{is\text{-}seminorm\text{-}family}(d, m, a) \wedge \operatorname{fr\text{-}cauchy}(m, a, q)) \Rightarrow  \\ \quad \quad \operatorname{is\text{-}cauchy\text{-}seq}(\operatorname{fr\text{-}metric\text{-}space}(m, a), q) \end{array}}}
\vnbprop{frm-converges-iff}{\vnbstmt{\begin{array}{@{}l@{}} \forall d, m, a, q, b. \\ \quad (\operatorname{is\text{-}seminorm\text{-}family}(d, m, a) \wedge ((q \in (\mathbb{N} \to \operatorname{vec}(m))) \wedge (b \in \operatorname{vec}(m)))) \Rightarrow  \\ \quad \quad \operatorname{converges\text{-}to}(\operatorname{fr\text{-}metric\text{-}space}(m, a), q, b) \Leftrightarrow \operatorname{fr\text{-}conv}(m, a, q, b) \end{array}}}
\vnbprop{frm-cauchy-iff}{\vnbstmt{\begin{array}{@{}l@{}} \forall d, m, a.; \forall q \in (\mathbb{N} \to \operatorname{vec}(m)). \\ \quad \operatorname{is\text{-}seminorm\text{-}family}(d, m, a) \Rightarrow  \\ \quad \quad \operatorname{is\text{-}cauchy\text{-}seq}(\operatorname{fr\text{-}metric\text{-}space}(m, a), q) \Leftrightarrow \operatorname{fr\text{-}cauchy}(m, a, q) \end{array}}}
\vnbprop{frm-frechet-iff-complete}{\vnbstmt{\begin{array}{@{}l@{}} \forall d, m, a. \\ \quad \operatorname{is\text{-}frechet\text{-}structure}(d, m, a) \Leftrightarrow (\operatorname{is\text{-}seminorm\text{-}family}(d, m, a) \wedge \operatorname{is\text{-}complete}(\operatorname{fr\text{-}metric\text{-}space}(m, a))) \end{array}}}
\vnbprop{frm-c-dist-congr}{\vnbstmt{\begin{array}{@{}l@{}} \forall s, x, y, a, b. \\ \quad \operatorname{is\text{-}c\text{-}metric\text{-}space}(s) \wedge \\ \quad x \in \operatorname{pts}(s) \wedge \\ \quad y \in \operatorname{pts}(s) \wedge \\ \quad a \in \operatorname{pts}(s) \wedge \\ \quad b \in \operatorname{pts}(s) \wedge \\ \quad \forall k \in \mathbb{N}.\; (\operatorname{dists}(s)(k)(x, y) = \operatorname{dists}(s)(k)(a, b)) \Rightarrow \\ \quad \quad \operatorname{dist}(\operatorname{c\text{-}metric}(s))(x, y) = \operatorname{dist}(\operatorname{c\text{-}metric}(s))(a, b) \end{array}}}
\vnbprop{frm-metric-translation-invariant}{\vnbstmt{\begin{array}{@{}l@{}} \forall d, m, a, x, y, r. \\ \quad (\operatorname{is\text{-}seminorm\text{-}family}(d, m, a) \wedge ((x \in \operatorname{vec}(m)) \wedge ((y \in \operatorname{vec}(m)) \wedge (r \in \operatorname{vec}(m))))) \Rightarrow  \\ \quad \quad \operatorname{dist}(\operatorname{fr\text{-}metric\text{-}space}(m, a))(\operatorname{vadd}(m)(x, r), \operatorname{vadd}(m)(y, r)) \\ \quad \quad =\; \operatorname{dist}(\operatorname{fr\text{-}metric\text{-}space}(m, a))(x, y) \end{array}}}
\vnbprop{frm-metric-value}{\vnbstmt{\begin{array}{@{}l@{}} \forall d, m, a, x, y. \\ \quad (\operatorname{is\text{-}seminorm\text{-}family}(d, m, a) \wedge ((x \in \operatorname{vec}(m)) \wedge (y \in \operatorname{vec}(m)))) \Rightarrow  \\ \quad \quad (\operatorname{fr\text{-}metric}(m, a)(x, y) \simeq \operatorname{series\text{-}limit}((\lambda k \in \mathbb{N}.\; (/(1, \operatorname{power}(2, (k + 1))) \cdot \operatorname{min}(1, a(k)(\operatorname{vadd}(m)(x, \operatorname{vneg}(m)(y)))))))) \end{array}}}
\vnbprop{frm-add-continuous}{\vnbstmt{\begin{array}{@{}l@{}} \forall d, m, c, f, g, a, b. \\ \quad \operatorname{is\text{-}seminorm\text{-}family}(d, m, c) \wedge \\ \quad f \in (\mathbb{N} \to \operatorname{vec}(m)) \wedge \\ \quad g \in (\mathbb{N} \to \operatorname{vec}(m)) \wedge \\ \quad a \in \operatorname{vec}(m) \wedge \\ \quad b \in \operatorname{vec}(m) \wedge \\ \quad \operatorname{converges\text{-}to}(\operatorname{fr\text{-}metric\text{-}space}(m, c), f, a) \wedge \\ \quad \operatorname{converges\text{-}to}(\operatorname{fr\text{-}metric\text{-}space}(m, c), g, b) \Rightarrow \\ \quad \quad \operatorname{converges\text{-}to}(\operatorname{fr\text{-}metric\text{-}space}(m, c), (\lambda n \in \mathbb{N}.\; \operatorname{vadd}(m)(f(n), g(n))), \operatorname{vadd}(m)(a, b)) \end{array}}}
\vnbprop{frm-linear-sub}{\vnbstmt{\begin{array}{@{}l@{}} \forall c, d, t, a, b. \\ \quad (\operatorname{is\text{-}linear\text{-}map}(c, d, t) \wedge ((a \in \operatorname{vec}(c)) \wedge (b \in \operatorname{vec}(c)))) \Rightarrow  \\ \quad \quad t(\operatorname{vadd}(c)(a, \operatorname{vneg}(c)(b))) = \operatorname{vadd}(d)(t(a), \operatorname{vneg}(d)(t(b))) \end{array}}}
\vnbprop{frm-cont-lin-of-bound}{\vnbstmt{\begin{array}{@{}l@{}} \forall d, a, b, c, f, t. \\ \quad \operatorname{is\text{-}seminorm\text{-}family}(d, a, b) \wedge \\ \quad \operatorname{is\text{-}seminorm\text{-}family}(d, c, f) \wedge \\ \quad \operatorname{is\text{-}linear\text{-}map}(a, c, t) \wedge \\ \quad \forall k \in \mathbb{N}. \exists g \in \mathbb{N}, h. \\ \quad \quad \operatorname{pos\text{-}rr}(h) \wedge \\ \quad \quad \forall x \in \operatorname{vec}(a), r \in \mathbb{R}. \\ \quad \quad \quad \forall j \in \mathbb{N}.\; ((j \leq g) \Rightarrow (b(j)(x) \leq r)) \Rightarrow (f(k)(t(x)) \leq (h \cdot r)) \Rightarrow \\ \quad \quad \operatorname{is\text{-}cont\text{-}lin}(a, b, c, f, t) \end{array}}}
\vnbprop{frm-ti-uniform}{\vnbstmt{\begin{array}{@{}l@{}} \forall m, a, b. \\ \quad \operatorname{is\text{-}module}(m) \wedge \\ \quad \operatorname{is\text{-}metric\text{-}space}(a) \wedge \\ \quad \operatorname{is\text{-}metric\text{-}space}(b) \wedge \\ \quad (\operatorname{pts}(a) \simeq \operatorname{vec}(m)) \wedge \\ \quad (\operatorname{pts}(b) \simeq \operatorname{vec}(m)) \wedge \\ \quad \forall x, y, r. \\ \quad \quad ((x \in \operatorname{vec}(m)) \wedge ((y \in \operatorname{vec}(m)) \wedge (r \in \operatorname{vec}(m)))) \Rightarrow  \\ \quad \quad \quad \operatorname{dist}(a)(\operatorname{vadd}(m)(x, r), \operatorname{vadd}(m)(y, r)) = \operatorname{dist}(a)(x, y) \wedge \\ \quad \forall x, y, r. \\ \quad \quad ((x \in \operatorname{vec}(m)) \wedge ((y \in \operatorname{vec}(m)) \wedge (r \in \operatorname{vec}(m)))) \Rightarrow  \\ \quad \quad \quad \operatorname{dist}(b)(\operatorname{vadd}(m)(x, r), \operatorname{vadd}(m)(y, r)) = \operatorname{dist}(b)(x, y) \wedge \\ \quad \forall q \in (\mathbb{N} \to \operatorname{vec}(m)). \\ \quad \quad \operatorname{converges\text{-}to}(a, q, \operatorname{vzero}(m)) \Rightarrow \operatorname{converges\text{-}to}(b, q, \operatorname{vzero}(m)) \Rightarrow \\ \quad \quad \operatorname{is\text{-}uniformly\text{-}continuous}(a, b, \operatorname{id\text{-}fun}(\operatorname{vec}(m))) \end{array}}}
\vnbprop{frm-metrics-uniformly-equivalent}{\vnbstmt{\begin{array}{@{}l@{}} \forall d, m, a, t. \\ \quad \operatorname{is\text{-}seminorm\text{-}family}(d, m, a) \wedge \\ \quad \operatorname{is\text{-}metric\text{-}space}(t) \wedge \\ \quad (\operatorname{pts}(t) \simeq \operatorname{vec}(m)) \wedge \\ \quad \forall x, y, r. \\ \quad \quad ((x \in \operatorname{vec}(m)) \wedge ((y \in \operatorname{vec}(m)) \wedge (r \in \operatorname{vec}(m)))) \Rightarrow  \\ \quad \quad \quad \operatorname{dist}(t)(\operatorname{vadd}(m)(x, r), \operatorname{vadd}(m)(y, r)) = \operatorname{dist}(t)(x, y) \wedge \\ \quad \forall q, b. \\ \quad \quad ((q \in (\mathbb{N} \to \operatorname{vec}(m))) \wedge (b \in \operatorname{vec}(m))) \Rightarrow  \\ \quad \quad \quad \operatorname{converges\text{-}to}(t, q, b) \Leftrightarrow \operatorname{fr\text{-}conv}(m, a, q, b) \Rightarrow \\ \quad \quad \operatorname{is\text{-}uniformly\text{-}continuous}(\operatorname{fr\text{-}metric\text{-}space}(m, a), t, \operatorname{id\text{-}fun}(\operatorname{vec}(m))) \wedge \\ \quad \quad \operatorname{is\text{-}uniformly\text{-}continuous}(t, \operatorname{fr\text{-}metric\text{-}space}(m, a), \operatorname{id\text{-}fun}(\operatorname{vec}(m))) \end{array}}}
\vnbprop{frm-sn-homog}{\vnbstmt{\begin{array}{@{}l@{}} \forall d, m, p.; \forall a \in \operatorname{carr}(\operatorname{scal}(m)), x \in \operatorname{vec}(m). \\ \quad \operatorname{is\text{-}seminorm}(d, m, p) \Rightarrow (p(\operatorname{act}(m)(a, x)) = (\operatorname{fnrm}(d)(a) \cdot p(x))) \end{array}}}
\vnbprop{frm-fam-field}{\vnbstmt{\forall d, m, a.\; (\operatorname{is\text{-}seminorm\text{-}family}(d, m, a) \Rightarrow \operatorname{is\text{-}normed\text{-}field}(d))}}
\vnbprop{frm-fam-scal}{\vnbstmt{\begin{array}{@{}l@{}} \forall d, m, a. \\ \quad \operatorname{is\text{-}seminorm\text{-}family}(d, m, a) \Rightarrow  \\ \quad \quad \operatorname{scal}(m) = \operatorname{normed\text{-}field\text{-}as\text{-}commutative\text{-}ring}(d) \end{array}}}
\vnbprop{frm-scal-continuous}{\vnbstmt{\begin{array}{@{}l@{}} \forall d, m, a.; \forall s \in (\mathbb{N} \to \operatorname{carr}(d)), b \in (\mathbb{N} \to \operatorname{vec}(m)), l \in \operatorname{carr}(d), x \in \operatorname{vec}(m). \\ \quad \operatorname{is\text{-}seminorm\text{-}family}(d, m, a) \wedge \\ \quad \operatorname{converges\text{-}to}(\operatorname{nf\text{-}metric\text{-}space}(d), s, l) \wedge \\ \quad \operatorname{converges\text{-}to}(\operatorname{fr\text{-}metric\text{-}space}(m, a), b, x) \Rightarrow \\ \quad \quad \operatorname{converges\text{-}to}(\operatorname{fr\text{-}metric\text{-}space}(m, a), (\lambda n \in \mathbb{N}.\; \operatorname{act}(m)(s(n), b(n))), \operatorname{act}(m)(l, x)) \end{array}}}
\vnbprop{frm-linear-zero}{\vnbstmt{\forall a, b, t.\; (\operatorname{is\text{-}linear\text{-}map}(a, b, t) \Rightarrow (t(\operatorname{vzero}(a)) = \operatorname{vzero}(b)))}}
\vnbprop{frm-linear-homog}{\vnbstmt{\begin{array}{@{}l@{}} \forall a, b, t.; \forall m \in \operatorname{carr}(\operatorname{scal}(a)), x \in \operatorname{vec}(a). \\ \quad \operatorname{is\text{-}linear\text{-}map}(a, b, t) \Rightarrow (t(\operatorname{act}(a)(m, x)) = \operatorname{act}(b)(m, t(x))) \end{array}}}
\vnbprop{frm-cont-lin-seq}{\vnbstmt{\begin{array}{@{}l@{}} \forall a, b, c, d, t.; \forall q \in (\mathbb{N} \to \operatorname{vec}(a)), m \in \operatorname{vec}(a). \\ \quad (\operatorname{is\text{-}cont\text{-}lin}(a, b, c, d, t) \wedge \operatorname{fr\text{-}conv}(a, b, q, m)) \Rightarrow  \\ \quad \quad \operatorname{fr\text{-}conv}(c, d, (\lambda i \in \mathbb{N}.\; t(q(i))), t(m)) \end{array}}}
\vnbprop{frm-cont-lin-nbhd}{\vnbstmt{\begin{array}{@{}l@{}} \forall d, a, b, c, f, t.; \forall k \in \mathbb{N}. \\ \quad \operatorname{is\text{-}seminorm\text{-}family}(d, a, b) \wedge \\ \quad \operatorname{is\text{-}seminorm\text{-}family}(d, c, f) \wedge \\ \quad \operatorname{is\text{-}linear\text{-}map}(a, c, t) \wedge \\ \quad \operatorname{is\text{-}cont\text{-}lin}(a, b, c, f, t) \Rightarrow \\ \quad \quad \exists g \in \mathbb{N}, l. \\ \quad \quad \quad \operatorname{pos\text{-}rr}(l) \wedge \\ \quad \quad \quad \forall x \in \operatorname{vec}(a). \\ \quad \quad \quad \quad \forall j \in \mathbb{N}.\; ((j \leq g) \Rightarrow (b(j)(x) < l)) \Rightarrow (f(k)(t(x)) < 1) \end{array}}}
\vnbprop{frm-nf-norm-one}{\vnbstmt{\begin{array}{@{}l@{}} \forall d.; \forall c \in \operatorname{carr}(d). \\ \quad (\operatorname{is\text{-}normed\text{-}field}(d) \wedge (0 < \operatorname{fnrm}(d)(c))) \Rightarrow (\operatorname{fnrm}(d)(\operatorname{one}(d)) = 1) \end{array}}}
\vnbprop{frm-nf-ladder}{\vnbstmt{\begin{array}{@{}l@{}} \forall n \in \mathbb{N}, d.; \forall c, b \in \operatorname{carr}(d). \\ \quad (\operatorname{is\text{-}normed\text{-}field}(d) \wedge ((1 < \operatorname{fnrm}(d)(c)) \wedge (\operatorname{mul}(d)(c, b) = \operatorname{one}(d)))) \Rightarrow  \\ \quad \quad \exists \mathit{fsc}, \mathit{fsb} \in \operatorname{carr}(d). \\ \quad \quad \quad (\operatorname{fnrm}(d)(\mathit{fsc}) \cdot \operatorname{fnrm}(d)(\mathit{fsb})) = 1 \wedge \\ \quad \quad \quad 1 \leq \operatorname{fnrm}(d)(\mathit{fsc}) \wedge \\ \quad \quad \quad (1 + ((n \cdot \operatorname{fnrm}(d)(c)) - n)) \leq \operatorname{fnrm}(d)(\mathit{fsc}) \wedge \\ \quad \quad \quad \forall u, v \in \mathbb{R}. \\ \quad \quad \quad \quad ((0 < u) \wedge ((0 < v) \wedge ((v < u) \wedge (u \leq (\operatorname{fnrm}(d)(\mathit{fsc}) \cdot v))))) \Rightarrow  \\ \quad \quad \quad \quad \quad \exists s \in \operatorname{carr}(d). \\ \quad \quad \quad \quad \quad \quad ((\operatorname{fnrm}(d)(s) \cdot v) < u) \wedge (u \leq (\operatorname{fnrm}(d)(c) \cdot (\operatorname{fnrm}(d)(s) \cdot v))) \wedge \\ \quad \quad \quad \forall u, v \in \mathbb{R}. \\ \quad \quad \quad \quad ((0 < u) \wedge ((0 < v) \wedge (((\operatorname{fnrm}(d)(\mathit{fsb}) \cdot v) < u) \wedge (u \leq v)))) \Rightarrow  \\ \quad \quad \quad \quad \quad \exists s \in \operatorname{carr}(d). \\ \quad \quad \quad \quad \quad \quad ((\operatorname{fnrm}(d)(s) \cdot v) < u) \wedge (u \leq (\operatorname{fnrm}(d)(c) \cdot (\operatorname{fnrm}(d)(s) \cdot v))) \end{array}}}
\vnbprop{frm-nf-bracket}{\vnbstmt{\begin{array}{@{}l@{}} \forall d.; \forall c, b \in \operatorname{carr}(d), u, v \in \mathbb{R}. \\ \quad (\operatorname{is\text{-}normed\text{-}field}(d) \wedge ((1 < \operatorname{fnrm}(d)(c)) \wedge ((\operatorname{mul}(d)(c, b) = \operatorname{one}(d)) \wedge ((0 < u) \wedge (0 < v))))) \Rightarrow  \\ \quad \quad \exists s \in \operatorname{carr}(d). \\ \quad \quad \quad ((\operatorname{fnrm}(d)(s) \cdot v) < u) \wedge (u \leq (\operatorname{fnrm}(d)(c) \cdot (\operatorname{fnrm}(d)(s) \cdot v))) \end{array}}}
\vnbprop{frm-bound-of-cont-lin}{\vnbstmt{\begin{array}{@{}l@{}} \forall d, a, b, c, f, t.; \forall k \in \mathbb{N}. \\ \quad \operatorname{is\text{-}seminorm\text{-}family}(d, a, b) \wedge \\ \quad \operatorname{is\text{-}seminorm\text{-}family}(d, c, f) \wedge \\ \quad \operatorname{is\text{-}linear\text{-}map}(a, c, t) \wedge \\ \quad \operatorname{is\text{-}nontrivially\text{-}valued}(d) \wedge \\ \quad \operatorname{is\text{-}cont\text{-}lin}(a, b, c, f, t) \Rightarrow \\ \quad \quad \exists g \in \mathbb{N}, h. \\ \quad \quad \quad \operatorname{pos\text{-}rr}(h) \wedge \\ \quad \quad \quad \forall x \in \operatorname{vec}(a), r \in \mathbb{R}. \\ \quad \quad \quad \quad \forall j \in \mathbb{N}.\; ((j \leq g) \Rightarrow (b(j)(x) \leq r)) \Rightarrow (f(k)(t(x)) \leq (h \cdot r)) \end{array}}}
\vnbprop{frm-cont-lin-iff-bound}{\vnbstmt{\begin{array}{@{}l@{}} \forall d, a, b, c, f, t. \\ \quad \operatorname{is\text{-}seminorm\text{-}family}(d, a, b) \wedge \\ \quad \operatorname{is\text{-}seminorm\text{-}family}(d, c, f) \wedge \\ \quad \operatorname{is\text{-}linear\text{-}map}(a, c, t) \wedge \\ \quad \operatorname{is\text{-}nontrivially\text{-}valued}(d) \Rightarrow \\ \quad \quad \operatorname{is\text{-}cont\text{-}lin}(a, b, c, f, t) \\ \quad \quad \Leftrightarrow\; \forall k \in \mathbb{N}. \exists g \in \mathbb{N}, h. \\ \quad \quad \quad \operatorname{pos\text{-}rr}(h) \wedge \\ \quad \quad \quad \forall x \in \operatorname{vec}(a), r \in \mathbb{R}. \\ \quad \quad \quad \quad \forall j \in \mathbb{N}.\; ((j \leq g) \Rightarrow (b(j)(x) \leq r)) \Rightarrow (f(k)(t(x)) \leq (h \cdot r)) \end{array}}}
\vnbprop{frm-rr-nontrivially-valued}{\vnbstmt{\operatorname{is\text{-}nontrivially\text{-}valued}(\operatorname{rr\text{-}normed\text{-}field})}}

\vnbbookfile{theorem-library/frechet-open-mapping-proof}{theorem-library/frechet-open-mapping-proof}
\noindent{\small frechet-open-mapping-proof.scm -- the OPEN MAPPING THEOREM for Frechet spaces and its bounded-inverse corollary: Thayer, "Analise Funcional e Teoria Espectral", Ch. 8 s.4, Thm 8.30, proved by the roadmap of exercise 8.9 (a)-(i).  Batch TVS-2A, 2026-10-11/12.  Helper prefix: fom-.  Every theorem here is proven modulo 0; ...}\par
\vnbprop{fom-prefix-bound}{\vnbstmt{\begin{array}{@{}l@{}} \forall b \in \mathbb{N}, m, x. \\ \quad \forall i \in \mathbb{N}.\; (m(i)(x) \in \mathbb{R}) \Rightarrow  \\ \quad \quad \exists p \in \mathbb{N}.\; \forall i \in \mathbb{N}.\; ((i \leq b) \Rightarrow (m(i)(x) < p)) \end{array}}}
\vnbprop{fom-dist-prefix}{\vnbstmt{\begin{array}{@{}l@{}} \forall k \in \mathbb{N}, d, m, a, s. \\ \quad (\operatorname{is\text{-}seminorm\text{-}family}(d, m, a) \wedge \operatorname{pos\text{-}rr}(s)) \Rightarrow  \\ \quad \quad \exists l. \\ \quad \quad \quad \operatorname{pos\text{-}rr}(l) \wedge \\ \quad \quad \quad \forall x, y.; \forall i \in \mathbb{N}. \\ \quad \quad \quad \quad ((x \in \operatorname{vec}(m)) \wedge ((y \in \operatorname{vec}(m)) \wedge ((\operatorname{dist}(\operatorname{fr\text{-}metric\text{-}space}(m, a))(x, y) < l) \wedge (i \leq k)))) \Rightarrow  \\ \quad \quad \quad \quad \quad a(i)(\operatorname{vadd}(m)(x, \operatorname{vneg}(m)(y))) < s \end{array}}}
\vnbprop{fom-unit-ladder}{\vnbstmt{\begin{array}{@{}l@{}} \forall n \in \mathbb{N}, d.; \forall c, b \in \operatorname{carr}(d). \\ \quad (\operatorname{is\text{-}normed\text{-}field}(d) \wedge ((1 < \operatorname{fnrm}(d)(c)) \wedge (\operatorname{mul}(d)(c, b) = \operatorname{one}(d)))) \Rightarrow  \\ \quad \quad \exists \mathit{fsc}, \mathit{fsb} \in \operatorname{carr}(d). \\ \quad \quad \quad ((\operatorname{mul}(d)(\mathit{fsc}, \mathit{fsb}) = \operatorname{one}(d)) \wedge ((1 \leq \operatorname{fnrm}(d)(\mathit{fsc})) \wedge ((1 + ((n \cdot \operatorname{fnrm}(d)(c)) - n)) \leq \operatorname{fnrm}(d)(\mathit{fsc})))) \end{array}}}
\vnbprop{fom-small-unit}{\vnbstmt{\begin{array}{@{}l@{}} \forall d.; \forall r \in \mathbb{R}, l. \\ \quad (\operatorname{is\text{-}nontrivially\text{-}valued}(d) \wedge ((0 \leq r) \wedge \operatorname{pos\text{-}rr}(l))) \Rightarrow  \\ \quad \quad \exists g, a \in \operatorname{carr}(d). \\ \quad \quad \quad \operatorname{mul}(d)(g, a) = \operatorname{one}(d) \wedge \\ \quad \quad \quad (\operatorname{fnrm}(d)(g) \cdot \operatorname{fnrm}(d)(a)) = 1 \wedge \\ \quad \quad \quad 0 < \operatorname{fnrm}(d)(g) \wedge \\ \quad \quad \quad 0 < \operatorname{fnrm}(d)(a) \wedge \\ \quad \quad \quad (\operatorname{fnrm}(d)(g) \cdot r) < l \end{array}}}
\vnbprop{fom-closure-near}{\vnbstmt{\begin{array}{@{}l@{}} \forall s, a.; \forall z \in \operatorname{closure}(s, a), o. \\ \quad (\operatorname{is\text{-}metric\text{-}space}(s) \wedge \operatorname{pos\text{-}rr}(o)) \Rightarrow \exists t \in a.\; (\operatorname{dist}(s)(z, t) < o) \end{array}}}
\vnbprop{fom-baire-ball}{\vnbstmt{\begin{array}{@{}l@{}} \forall s.; \forall a \in \operatorname{pts}(s), e \in (\mathbb{N} \to \operatorname{power}(\operatorname{pts}(s))). \\ \quad (\operatorname{is\text{-}complete}(s) \wedge \forall y \in \operatorname{pts}(s).\; \exists j \in \mathbb{N}.\; (y \in e(j))) \Rightarrow  \\ \quad \quad \exists j \in \mathbb{N}, c \in \operatorname{pts}(s), r. \\ \quad \quad \quad (\operatorname{pos\text{-}rr}(r) \wedge \forall z \in \operatorname{pts}(s).\; ((\operatorname{dist}(s)(c, z) < r) \Rightarrow (z \in \operatorname{closure}(s, e(j))))) \end{array}}}
\vnbprop{fom-v-shift-back}{\vnbstmt{\begin{array}{@{}l@{}} \forall m.; \forall w, c \in \operatorname{vec}(m). \\ \quad \operatorname{is\text{-}module}(m) \Rightarrow (\operatorname{vadd}(m)(\operatorname{vadd}(m)(w, c), \operatorname{vneg}(m)(c)) = w) \end{array}}}
\vnbprop{fom-v-quad}{\vnbstmt{\begin{array}{@{}l@{}} \forall m.; \forall p, q, a, b \in \operatorname{vec}(m). \\ \quad \operatorname{is\text{-}module}(m) \Rightarrow  \\ \quad \quad \operatorname{vadd}(m)(\operatorname{vadd}(m)(p, \operatorname{vneg}(m)(q)), \operatorname{vneg}(m)(\operatorname{vadd}(m)(a, \operatorname{vneg}(m)(b)))) \\ \quad \quad =\; \operatorname{vadd}(m)(\operatorname{vadd}(m)(p, \operatorname{vneg}(m)(a)), \operatorname{vneg}(m)(\operatorname{vadd}(m)(q, \operatorname{vneg}(m)(b)))) \end{array}}}
\vnbprop{fom-v-sub-sum}{\vnbstmt{\begin{array}{@{}l@{}} \forall m.; \forall y, a, b \in \operatorname{vec}(m). \\ \quad \operatorname{is\text{-}module}(m) \Rightarrow  \\ \quad \quad \operatorname{vadd}(m)(y, \operatorname{vneg}(m)(\operatorname{vadd}(m)(a, b))) = \operatorname{vadd}(m)(\operatorname{vadd}(m)(y, \operatorname{vneg}(m)(b)), \operatorname{vneg}(m)(a)) \end{array}}}
\vnbprop{fom-lin-m1}{\vnbstmt{\forall a, b, t.\; (\operatorname{is\text{-}linear\text{-}map}(a, b, t) \Rightarrow \operatorname{is\text{-}module}(a))}}
\vnbprop{fom-lin-m2}{\vnbstmt{\forall a, b, t.\; (\operatorname{is\text{-}linear\text{-}map}(a, b, t) \Rightarrow \operatorname{is\text{-}module}(b))}}
\vnbprop{fom-lin-scal}{\vnbstmt{\forall a, b, t.\; (\operatorname{is\text{-}linear\text{-}map}(a, b, t) \Rightarrow (\operatorname{scal}(a) = \operatorname{scal}(b)))}}
\vnbprop{fom-lin-type}{\vnbstmt{\forall a, b, t.\; (\operatorname{is\text{-}linear\text{-}map}(a, b, t) \Rightarrow (t \in (\operatorname{vec}(a) \to \operatorname{vec}(b))))}}
\vnbprop{fom-lin-add}{\vnbstmt{\begin{array}{@{}l@{}} \forall a, b, t.; \forall x, y \in \operatorname{vec}(a). \\ \quad \operatorname{is\text{-}linear\text{-}map}(a, b, t) \Rightarrow (t(\operatorname{vadd}(a)(x, y)) = \operatorname{vadd}(b)(t(x), t(y))) \end{array}}}
\vnbprop{fom-sn-sub-le}{\vnbstmt{\begin{array}{@{}l@{}} \forall d, m, p.; \forall u, v \in \operatorname{vec}(m). \\ \quad \operatorname{is\text{-}seminorm}(d, m, p) \Rightarrow (p(\operatorname{vadd}(m)(u, \operatorname{vneg}(m)(v))) \leq (p(u) + p(v))) \end{array}}}
\vnbprop{fom-act-cancel}{\vnbstmt{\begin{array}{@{}l@{}} \forall d, m.; \forall g, a \in \operatorname{carr}(d), z \in \operatorname{vec}(m). \\ \quad \operatorname{is\text{-}normed\text{-}field}(d) \wedge \\ \quad \operatorname{is\text{-}module}(m) \wedge \\ \quad \operatorname{scal}(m) = \operatorname{normed\text{-}field\text{-}as\text{-}commutative\text{-}ring}(d) \wedge \\ \quad \operatorname{mul}(d)(g, a) = \operatorname{one}(d) \Rightarrow \\ \quad \quad \operatorname{act}(m)(g, \operatorname{act}(m)(a, z)) = z \end{array}}}
\vnbprop{fom-mul-lt-left}{\vnbstmt{\forall a, b, c \in \mathbb{R}.\; (((0 < c) \wedge (a < b)) \Rightarrow ((c \cdot a) < (c \cdot b)))}}
\vnbprop{fom-baire-approx}{\vnbstmt{\begin{array}{@{}l@{}} \forall d, a, b, c, f, t.; \forall n \in \mathbb{N}, l. \\ \quad \operatorname{is\text{-}nontrivially\text{-}valued}(d) \wedge \\ \quad \operatorname{is\text{-}seminorm\text{-}family}(d, a, b) \wedge \\ \quad \operatorname{is\text{-}frechet\text{-}structure}(d, c, f) \wedge \\ \quad \operatorname{is\text{-}linear\text{-}map}(a, c, t) \wedge \\ \quad \forall y \in \operatorname{vec}(c).\; \exists x \in \operatorname{vec}(a).\; (y = t(x)) \wedge \\ \quad \operatorname{pos\text{-}rr}(l) \Rightarrow \\ \quad \quad \exists \mathit{fnb} \in \mathbb{N}, \mathit{feta}. \\ \quad \quad \quad \operatorname{pos\text{-}rr}(\mathit{feta}) \wedge \\ \quad \quad \quad \forall z \in \operatorname{vec}(c), k \in \mathbb{N}, h. \\ \quad \quad \quad \quad (\forall i \in \mathbb{N}.\; ((i \leq \mathit{fnb}) \Rightarrow (f(i)(z) < \mathit{feta})) \wedge \operatorname{pos\text{-}rr}(h)) \Rightarrow  \\ \quad \quad \quad \quad \quad \exists \mathit{fx} \in \operatorname{vec}(a). \\ \quad \quad \quad \quad \quad \quad \forall j \in \mathbb{N}.\; ((j \leq n) \Rightarrow (b(j)(\mathit{fx}) < l)) \wedge \\ \quad \quad \quad \quad \quad \quad \forall i \in \mathbb{N}.\; ((i \leq k) \Rightarrow (f(i)(\operatorname{vadd}(c)(z, \operatorname{vneg}(c)(t(\mathit{fx})))) < h)) \end{array}}}
\vnbprop{fom-a-facts}{\vnbstmt{\begin{array}{@{}l@{}} \forall a, k \in \mathbb{N}. \\ \quad \operatorname{recip}(((a + k) + 1)) \in \mathbb{R} \wedge \\ \quad 0 < \operatorname{recip}(((a + k) + 1)) \wedge \\ \quad \operatorname{recip}(((a + (k + 1)) + 1)) < \operatorname{recip}(((a + k) + 1)) \end{array}}}
\vnbprop{fom-a-small}{\vnbstmt{\begin{array}{@{}l@{}} \forall d. \\ \quad \operatorname{pos\text{-}rr}(d) \Rightarrow  \\ \quad \quad \exists p \in \mathbb{N}. \forall a, k \in \mathbb{N}. \\ \quad \quad \quad (p \leq k) \Rightarrow (\operatorname{recip}(((a + k) + 1)) \leq d) \end{array}}}
\vnbprop{fom-telescope}{\vnbstmt{\begin{array}{@{}l@{}} \forall l \in \mathbb{N}, d, m, a, f, v, b.; \forall k, j \in \mathbb{N}. \\ \quad \operatorname{is\text{-}seminorm\text{-}family}(d, m, a) \wedge \\ \quad f \in (\mathbb{N} \to \operatorname{vec}(m)) \wedge \\ \quad v \in \mathbb{N} \wedge \\ \quad b \in \mathbb{N} \wedge \\ \quad \forall i, j \in \mathbb{N}. \\ \quad \quad (j \leq (v + i)) \Rightarrow  \\ \quad \quad \quad a(j)(\operatorname{vadd}(m)(f((i + 1)), \operatorname{vneg}(m)(f(i)))) < (\operatorname{recip}(((b + i) + 1)) - \operatorname{recip}(((b + (i + 1)) + 1))) \wedge \\ \quad k \leq l \wedge \\ \quad j \leq (v + k) \Rightarrow \\ \quad \quad a(j)(\operatorname{vadd}(m)(f(l), \operatorname{vneg}(m)(f(k)))) \leq (\operatorname{recip}(((b + k) + 1)) - \operatorname{recip}(((b + l) + 1))) \end{array}}}
\vnbprop{fom-fr-ball-unfold}{\vnbstmt{\begin{array}{@{}l@{}} \forall m, a, n, s. \\ \quad (\operatorname{fr\text{-}ball}(m, a, n, s) \simeq \{x \in \operatorname{vec}(m) : \forall k \in \mathbb{N}.\; ((k \leq n) \Rightarrow (a(k)(x) < s))\}) \end{array}}}
\vnbprop{fom-sn-swap}{\vnbstmt{\begin{array}{@{}l@{}} \forall d, m, p.; \forall a, b \in \operatorname{vec}(m). \\ \quad \operatorname{is\text{-}seminorm}(d, m, p) \Rightarrow  \\ \quad \quad p(\operatorname{vadd}(m)(a, \operatorname{vneg}(m)(b))) = p(\operatorname{vadd}(m)(b, \operatorname{vneg}(m)(a))) \end{array}}}
\vnbprop{fom-open-of-approx}{\vnbstmt{\begin{array}{@{}l@{}} \forall d, a, b, c, f, t. \\ \quad \operatorname{is\text{-}frechet\text{-}structure}(d, a, b) \wedge \\ \quad \operatorname{is\text{-}seminorm\text{-}family}(d, c, f) \wedge \\ \quad \operatorname{is\text{-}cont\text{-}lin}(a, b, c, f, t) \wedge \\ \quad \forall n \in \mathbb{N}, l. \\ \quad \quad \operatorname{pos\text{-}rr}(l) \Rightarrow  \\ \quad \quad \quad \exists \mathit{fnb} \in \mathbb{N}, \mathit{feta}. \\ \quad \quad \quad \quad \operatorname{pos\text{-}rr}(\mathit{feta}) \wedge \\ \quad \quad \quad \quad \forall z \in \operatorname{vec}(c), k \in \mathbb{N}, h. \\ \quad \quad \quad \quad \quad (\forall i \in \mathbb{N}.\; ((i \leq \mathit{fnb}) \Rightarrow (f(i)(z) < \mathit{feta})) \wedge \operatorname{pos\text{-}rr}(h)) \Rightarrow  \\ \quad \quad \quad \quad \quad \quad \exists x \in \operatorname{vec}(a). \\ \quad \quad \quad \quad \quad \quad \quad \forall j \in \mathbb{N}.\; ((j \leq n) \Rightarrow (b(j)(x) < l)) \wedge \\ \quad \quad \quad \quad \quad \quad \quad \forall i \in \mathbb{N}.\; ((i \leq k) \Rightarrow (f(i)(\operatorname{vadd}(c)(z, \operatorname{vneg}(c)(t(x)))) < h)) \Rightarrow \\ \quad \quad \operatorname{is\text{-}open\text{-}lin\text{-}map}(a, b, c, f, t) \end{array}}}
\vnbprop{open-mapping-theorem}{\vnbstmt{\begin{array}{@{}l@{}} \forall d, a, b, c, f, t. \\ \quad \operatorname{is\text{-}nontrivially\text{-}valued}(d) \wedge \\ \quad \operatorname{is\text{-}frechet\text{-}structure}(d, a, b) \wedge \\ \quad \operatorname{is\text{-}frechet\text{-}structure}(d, c, f) \wedge \\ \quad \operatorname{is\text{-}cont\text{-}lin}(a, b, c, f, t) \wedge \\ \quad \forall y \in \operatorname{vec}(c).\; \exists x \in \operatorname{vec}(a).\; (y = t(x)) \Rightarrow \\ \quad \quad \operatorname{is\text{-}open\text{-}lin\text{-}map}(a, b, c, f, t) \end{array}}}
\vnbprop{open-mapping-bounded-inverse}{\vnbstmt{\begin{array}{@{}l@{}} \forall d, a, b, c, f, t, g. \\ \quad \operatorname{is\text{-}nontrivially\text{-}valued}(d) \wedge \\ \quad \operatorname{is\text{-}frechet\text{-}structure}(d, a, b) \wedge \\ \quad \operatorname{is\text{-}frechet\text{-}structure}(d, c, f) \wedge \\ \quad \operatorname{is\text{-}cont\text{-}lin}(a, b, c, f, t) \wedge \\ \quad \operatorname{is\text{-}linear\text{-}map}(c, a, g) \wedge \\ \quad \forall x \in \operatorname{vec}(a).\; (g(t(x)) = x) \wedge \\ \quad \forall y \in \operatorname{vec}(c).\; (t(g(y)) = y) \Rightarrow \\ \quad \quad \operatorname{is\text{-}cont\text{-}lin}(c, f, a, b, g) \end{array}}}

\vnbbookfile{theorem-library/ba-metric-lemmas}{theorem-library/ba-metric-lemmas}
\noindent{\small ba-metric-lemmas.scm -- wave of 2026-11-09/10 (agent BA-1B), wired 2026-11-10 from scratchpad/ba1b-1109/ba-metric-lemmas.scm. Library: analysis.}\par
\vnbprop{metric-limit-dist-le}{\vnbstmt{\begin{array}{@{}l@{}} \forall s.; \forall c \in \operatorname{pts}(s), d \in \mathbb{R}, f, l.; \forall b \in \mathbb{N}. \\ \quad \operatorname{is\text{-}metric\text{-}space}(s) \wedge \\ \quad \operatorname{converges\text{-}to}(s, f, l) \wedge \\ \quad \forall n \in \mathbb{N}.\; ((b \leq n) \Rightarrow (\operatorname{dist}(s)(c, f(n)) \leq d)) \Rightarrow \\ \quad \quad \operatorname{dist}(s)(c, l) \leq d \end{array}}}
\vnbprop{closure-near-intro}{\vnbstmt{\begin{array}{@{}l@{}} \forall s, a.; \forall x \in \operatorname{pts}(s). \\ \quad \operatorname{is\text{-}metric\text{-}space}(s) \wedge \\ \quad a \subseteq \operatorname{pts}(s) \wedge \\ \quad \forall o.\; (\operatorname{pos\text{-}rr}(o) \Rightarrow \exists y \in a.\; (\operatorname{dist}(s)(x, y) < o)) \Rightarrow \\ \quad \quad x \in \operatorname{closure}(s, a) \end{array}}}
\vnbprop{lipschitz-closure-image}{\vnbstmt{\begin{array}{@{}l@{}} \forall s.; \forall f \in (\operatorname{pts}(s) \to \operatorname{pts}(s)), c \in \mathbb{R}, a, b.; \forall x \in \operatorname{closure}(s, a). \\ \quad \operatorname{is\text{-}metric\text{-}space}(s) \wedge \\ \quad 0 \leq c \wedge \\ \quad \forall x, p \in \operatorname{pts}(s).\; (\operatorname{dist}(s)(f(x), f(p)) \leq (c \cdot \operatorname{dist}(s)(x, p))) \wedge \\ \quad a \subseteq \operatorname{pts}(s) \wedge \\ \quad b \subseteq \operatorname{pts}(s) \wedge \\ \quad \forall p \in a.\; (f(p) \in \operatorname{closure}(s, b)) \Rightarrow \\ \quad \quad f(x) \in \operatorname{closure}(s, b) \end{array}}}

\vnbbookfile{theorem-library/frechet-banach-steinhaus}{theorem-library/frechet-banach-steinhaus}
\noindent{\small frechet-banach-steinhaus.scm -- Thayer, "Analise Funcional e Teoria Espectral", Ch. 8 s.4, pp. 256-257: Thm 8.29 (closed graph bis) and Thm 8.32 (Banach-Steinhaus bis, exercise 8.10).  Batch TVS-2B, 2026-10-11/12.  Vocabulary (Def. 8.31, the equicontinuity of 8.32): structure-library/frechet-banach-steinhaus.scm.}\par
\vnbprop{fbs-conv-dom}{\vnbstmt{\begin{array}{@{}l@{}} \forall a, b, f, c, s, d, g, h. \\ \quad \operatorname{is\text{-}module}(a) \wedge \\ \quad \operatorname{is\text{-}module}(b) \wedge \\ \quad f \in (\mathbb{N} \to (\operatorname{vec}(a) \to \mathbb{R})) \wedge \\ \quad c \in (\mathbb{N} \to (\operatorname{vec}(b) \to \mathbb{R})) \wedge \\ \quad s \in (\mathbb{N} \to \operatorname{vec}(a)) \wedge \\ \quad d \in (\mathbb{N} \to \operatorname{vec}(b)) \wedge \\ \quad g \in \operatorname{vec}(a) \wedge \\ \quad h \in \operatorname{vec}(b) \wedge \\ \quad \forall k, i. \\ \quad \quad ((k \in \mathbb{N}) \wedge (i \in \mathbb{N})) \Rightarrow  \\ \quad \quad \quad c(k)(\operatorname{vadd}(b)(d(i), \operatorname{vneg}(b)(h))) \leq f(k)(\operatorname{vadd}(a)(s(i), \operatorname{vneg}(a)(g))) \wedge \\ \quad \operatorname{fr\text{-}conv}(a, f, s, g) \Rightarrow \\ \quad \quad \operatorname{fr\text{-}conv}(b, c, d, h) \end{array}}}
\vnbprop{fbs-cauchy-dom}{\vnbstmt{\begin{array}{@{}l@{}} \forall a, b, f, c, s, d. \\ \quad \operatorname{is\text{-}module}(a) \wedge \\ \quad \operatorname{is\text{-}module}(b) \wedge \\ \quad f \in (\mathbb{N} \to (\operatorname{vec}(a) \to \mathbb{R})) \wedge \\ \quad c \in (\mathbb{N} \to (\operatorname{vec}(b) \to \mathbb{R})) \wedge \\ \quad s \in (\mathbb{N} \to \operatorname{vec}(a)) \wedge \\ \quad d \in (\mathbb{N} \to \operatorname{vec}(b)) \wedge \\ \quad \forall k, i, j. \\ \quad \quad ((k \in \mathbb{N}) \wedge ((i \in \mathbb{N}) \wedge (j \in \mathbb{N}))) \Rightarrow  \\ \quad \quad \quad c(k)(\operatorname{vadd}(b)(d(i), \operatorname{vneg}(b)(d(j)))) \leq f(k)(\operatorname{vadd}(a)(s(i), \operatorname{vneg}(a)(s(j)))) \wedge \\ \quad \operatorname{fr\text{-}cauchy}(a, f, s) \Rightarrow \\ \quad \quad \operatorname{fr\text{-}cauchy}(b, c, d) \end{array}}}
\vnbprop{fbs-conv-sum}{\vnbstmt{\begin{array}{@{}l@{}} \forall a, b, f, c, d, s, g, h, m. \\ \quad \operatorname{is\text{-}module}(a) \wedge \\ \quad \operatorname{is\text{-}module}(b) \wedge \\ \quad f \in (\mathbb{N} \to (\operatorname{vec}(a) \to \mathbb{R})) \wedge \\ \quad c \in (\mathbb{N} \to (\operatorname{vec}(a) \to \mathbb{R})) \wedge \\ \quad d \in (\mathbb{N} \to (\operatorname{vec}(b) \to \mathbb{R})) \wedge \\ \quad s \in (\mathbb{N} \to \operatorname{vec}(a)) \wedge \\ \quad g \in (\mathbb{N} \to \operatorname{vec}(b)) \wedge \\ \quad h \in \operatorname{vec}(a) \wedge \\ \quad m \in \operatorname{vec}(b) \wedge \\ \quad \forall k, i. \\ \quad \quad ((k \in \mathbb{N}) \wedge (i \in \mathbb{N})) \Rightarrow  \\ \quad \quad \quad f(k)(\operatorname{vadd}(a)(s(i), \operatorname{vneg}(a)(h))) \\ \quad \quad \quad \leq\; c(k)(\operatorname{vadd}(a)(s(i), \operatorname{vneg}(a)(h))) + d(k)(\operatorname{vadd}(b)(g(i), \operatorname{vneg}(b)(m))) \wedge \\ \quad \operatorname{fr\text{-}conv}(a, c, s, h) \wedge \\ \quad \operatorname{fr\text{-}conv}(b, d, g, m) \Rightarrow \\ \quad \quad \operatorname{fr\text{-}conv}(a, f, s, h) \end{array}}}
\vnbprop{fbs-graph-fam-exists}{\vnbstmt{\begin{array}{@{}l@{}} \forall d, a, b, c, f, t. \\ \quad (\operatorname{is\text{-}seminorm\text{-}family}(d, a, b) \wedge (\operatorname{is\text{-}seminorm\text{-}family}(d, c, f) \wedge (t \in (\operatorname{vec}(a) \to \operatorname{vec}(c))))) \Rightarrow  \\ \quad \quad \exists g \in (\mathbb{N} \to (\operatorname{vec}(a) \to \mathbb{R})). \forall k \in \mathbb{N}, x \in \operatorname{vec}(a). \\ \quad \quad \quad g(k)(x) = (b(k)(x) + f(k)(t(x))) \end{array}}}
\vnbprop{fbs-graph-seminorm}{\vnbstmt{\begin{array}{@{}l@{}} \forall d, a, b, c, f, t, g.; \forall k \in \mathbb{N}. \\ \quad \operatorname{is\text{-}seminorm\text{-}family}(d, a, b) \wedge \\ \quad \operatorname{is\text{-}seminorm\text{-}family}(d, c, f) \wedge \\ \quad \operatorname{is\text{-}linear\text{-}map}(a, c, t) \wedge \\ \quad g \in (\mathbb{N} \to (\operatorname{vec}(a) \to \mathbb{R})) \wedge \\ \quad \forall k \in \mathbb{N}, x \in \operatorname{vec}(a).\; (g(k)(x) = (b(k)(x) + f(k)(t(x)))) \Rightarrow \\ \quad \quad \operatorname{is\text{-}seminorm}(d, a, g(k)) \end{array}}}
\vnbprop{fbs-graph-family}{\vnbstmt{\begin{array}{@{}l@{}} \forall d, a, b, c, f, t, g. \\ \quad \operatorname{is\text{-}seminorm\text{-}family}(d, a, b) \wedge \\ \quad \operatorname{is\text{-}seminorm\text{-}family}(d, c, f) \wedge \\ \quad \operatorname{is\text{-}linear\text{-}map}(a, c, t) \wedge \\ \quad g \in (\mathbb{N} \to (\operatorname{vec}(a) \to \mathbb{R})) \wedge \\ \quad \forall k \in \mathbb{N}, x \in \operatorname{vec}(a).\; (g(k)(x) = (b(k)(x) + f(k)(t(x)))) \Rightarrow \\ \quad \quad \operatorname{is\text{-}seminorm\text{-}family}(d, a, g) \end{array}}}
\vnbprop{fbs-graph-frechet}{\vnbstmt{\begin{array}{@{}l@{}} \forall d, a, b, c, f, t, g. \\ \quad \operatorname{is\text{-}frechet\text{-}structure}(d, a, b) \wedge \\ \quad \operatorname{is\text{-}frechet\text{-}structure}(d, c, f) \wedge \\ \quad \operatorname{has\text{-}closed\text{-}graph}(a, b, c, f, t) \wedge \\ \quad g \in (\mathbb{N} \to (\operatorname{vec}(a) \to \mathbb{R})) \wedge \\ \quad \forall k \in \mathbb{N}, x \in \operatorname{vec}(a).\; (g(k)(x) = (b(k)(x) + f(k)(t(x)))) \Rightarrow \\ \quad \quad \operatorname{is\text{-}frechet\text{-}structure}(d, a, g) \end{array}}}
\vnbprop{fbs-id-linear}{\vnbstmt{\forall a.\; (\operatorname{is\text{-}module}(a) \Rightarrow \operatorname{is\text{-}linear\text{-}map}(a, a, (\lambda v \in \operatorname{vec}(a).\; v)))}}
\vnbprop{closed-graph-theorem}{\vnbstmt{\begin{array}{@{}l@{}} \forall d, a, b, c, f, t. \\ \quad \operatorname{is\text{-}nontrivially\text{-}valued}(d) \wedge \\ \quad \operatorname{is\text{-}frechet\text{-}structure}(d, a, b) \wedge \\ \quad \operatorname{is\text{-}frechet\text{-}structure}(d, c, f) \wedge \\ \quad \operatorname{has\text{-}closed\text{-}graph}(a, b, c, f, t) \Rightarrow \\ \quad \quad \operatorname{is\text{-}cont\text{-}lin}(a, b, c, f, t) \end{array}}}
\vnbprop{fbs-baire-ball}{\vnbstmt{\begin{array}{@{}l@{}} \forall d, a, b, c, f, g, h.; \forall k \in \mathbb{N}. \\ \quad \operatorname{is\text{-}frechet\text{-}structure}(d, a, b) \wedge \\ \quad \operatorname{is\text{-}seminorm\text{-}family}(d, c, f) \wedge \\ \quad h \in (g \to (\operatorname{vec}(a) \to \operatorname{vec}(c))) \wedge \\ \quad \forall x \in \operatorname{vec}(a). \\ \quad \quad \operatorname{is\text{-}fr\text{-}bounded}(c, f, \{y \in \operatorname{vec}(c) : \exists i \in g.\; (y = h(i)(x))\}) \Rightarrow \\ \quad \quad \exists n \in \mathbb{N}, m \in \operatorname{vec}(a), p. \\ \quad \quad \quad \operatorname{pos\text{-}rr}(p) \wedge \\ \quad \quad \quad \operatorname{ball}(\operatorname{fr\text{-}metric\text{-}space}(a, b), m, p) \\ \quad \quad \quad \subseteq\; \operatorname{closure}(\operatorname{fr\text{-}metric\text{-}space}(a, b), \{\mathit{bx} \in \operatorname{vec}(a) : \forall \mathit{bi} \in g.\; (f(k)(h(\mathit{bi})(\mathit{bx})) \leq n)\}) \end{array}}}
\vnbprop{fbs-v-sub-add}{\vnbstmt{\begin{array}{@{}l@{}} \forall m.; \forall x, y \in \operatorname{vec}(m). \\ \quad \operatorname{is\text{-}module}(m) \Rightarrow (\operatorname{vadd}(m)(\operatorname{vadd}(m)(x, \operatorname{vneg}(m)(y)), y) = x) \end{array}}}
\vnbprop{fbs-q-cont}{\vnbstmt{\begin{array}{@{}l@{}} \forall d, a, b, c, f, t, k, y, s. \\ \quad \operatorname{is\text{-}seminorm\text{-}family}(d, a, b) \wedge \\ \quad \operatorname{is\text{-}seminorm\text{-}family}(d, c, f) \wedge \\ \quad \operatorname{is\text{-}cont\text{-}lin}(a, b, c, f, t) \wedge \\ \quad k \in \mathbb{N} \wedge \\ \quad y \in \operatorname{vec}(a) \wedge \\ \quad \operatorname{pos\text{-}rr}(s) \Rightarrow \\ \quad \quad \exists e. \\ \quad \quad \quad \operatorname{pos\text{-}rr}(e) \wedge \\ \quad \quad \quad \forall z \in \operatorname{vec}(a). \\ \quad \quad \quad \quad (\operatorname{dist}(\operatorname{fr\text{-}metric\text{-}space}(a, b))(y, z) \leq e) \Rightarrow  \\ \quad \quad \quad \quad \quad f(k)(t(y)) \leq (f(k)(t(z)) + s) \end{array}}}
\vnbprop{fbs-closure-bound}{\vnbstmt{\begin{array}{@{}l@{}} \forall d, a, b, c, f, x, g, k, n, z.; \forall j \in x. \\ \quad \operatorname{is\text{-}seminorm\text{-}family}(d, a, b) \wedge \\ \quad \operatorname{is\text{-}seminorm\text{-}family}(d, c, f) \wedge \\ \quad g \in (x \to (\operatorname{vec}(a) \to \operatorname{vec}(c))) \wedge \\ \quad \forall i \in x.\; \operatorname{is\text{-}cont\text{-}lin}(a, b, c, f, g(i)) \wedge \\ \quad k \in \mathbb{N} \wedge \\ \quad n \in \mathbb{N} \wedge \\ \quad z \in \operatorname{vec}(a) \wedge \\ \quad z \in \operatorname{closure}(\operatorname{fr\text{-}metric\text{-}space}(a, b), \{\mathit{bx} \in \operatorname{vec}(a) : \forall \mathit{bi} \in x.\; (f(k)(g(\mathit{bi})(\mathit{bx})) \leq n)\}) \Rightarrow \\ \quad \quad f(k)(g(j)(z)) \leq n \end{array}}}
\vnbprop{fbs-ball-bound}{\vnbstmt{\begin{array}{@{}l@{}} \forall d, a, f, g, h, t.; \forall c, b \in \operatorname{carr}(d), k, m, l, n.; \forall x \in \operatorname{vec}(a), r \in \mathbb{R}. \\ \quad \operatorname{is\text{-}seminorm\text{-}family}(d, a, f) \wedge \\ \quad \operatorname{is\text{-}seminorm\text{-}family}(d, g, h) \wedge \\ \quad \operatorname{is\text{-}linear\text{-}map}(a, g, t) \wedge \\ \quad 1 < \operatorname{fnrm}(d)(c) \wedge \\ \quad \operatorname{mul}(d)(c, b) = \operatorname{one}(d) \wedge \\ \quad k \in \mathbb{N} \wedge \\ \quad m \in \mathbb{N} \wedge \\ \quad \operatorname{pos\text{-}rr}(l) \wedge \\ \quad \operatorname{pos\text{-}rr}(n) \wedge \\ \quad \forall x \in \operatorname{vec}(a). \\ \quad \quad \forall j \in \mathbb{N}.\; ((j \leq m) \Rightarrow (f(j)(x) < l)) \Rightarrow (h(k)(t(x)) \leq n) \wedge \\ \quad \forall j \in \mathbb{N}.\; ((j \leq m) \Rightarrow (f(j)(x) \leq r)) \Rightarrow \\ \quad \quad h(k)(t(x)) \leq (((\operatorname{fnrm}(d)(c) \cdot n) \cdot \operatorname{recip}(l)) \cdot r) \end{array}}}
\vnbprop{fbs-uniform-ball}{\vnbstmt{\begin{array}{@{}l@{}} \forall d, a, b, c, f, g, h.; \forall k \in \mathbb{N}. \\ \quad \operatorname{is\text{-}frechet\text{-}structure}(d, a, b) \wedge \\ \quad \operatorname{is\text{-}seminorm\text{-}family}(d, c, f) \wedge \\ \quad h \in (g \to (\operatorname{vec}(a) \to \operatorname{vec}(c))) \wedge \\ \quad \forall i \in g.\; \operatorname{is\text{-}cont\text{-}lin}(a, b, c, f, h(i)) \wedge \\ \quad \forall x \in \operatorname{vec}(a). \\ \quad \quad \operatorname{is\text{-}fr\text{-}bounded}(c, f, \{y \in \operatorname{vec}(c) : \exists i \in g.\; (y = h(i)(x))\}) \Rightarrow \\ \quad \quad \exists m \in \mathbb{N}, l. \\ \quad \quad \quad \operatorname{pos\text{-}rr}(l) \wedge \\ \quad \quad \quad \exists n. \\ \quad \quad \quad \quad \operatorname{pos\text{-}rr}(n) \wedge \\ \quad \quad \quad \quad \forall i \in g, x \in \operatorname{vec}(a). \\ \quad \quad \quad \quad \quad \forall j \in \mathbb{N}.\; ((j \leq m) \Rightarrow (b(j)(x) < l)) \Rightarrow (f(k)(h(i)(x)) \leq n) \end{array}}}
\vnbprop{frechet-banach-steinhaus}{\vnbstmt{\begin{array}{@{}l@{}} \forall d, a, b, c, f, g, h. \\ \quad \operatorname{is\text{-}nontrivially\text{-}valued}(d) \wedge \\ \quad \operatorname{is\text{-}frechet\text{-}structure}(d, a, b) \wedge \\ \quad \operatorname{is\text{-}seminorm\text{-}family}(d, c, f) \wedge \\ \quad h \in (g \to (\operatorname{vec}(a) \to \operatorname{vec}(c))) \wedge \\ \quad \forall i \in g.\; \operatorname{is\text{-}cont\text{-}lin}(a, b, c, f, h(i)) \wedge \\ \quad \forall x \in \operatorname{vec}(a). \\ \quad \quad \operatorname{is\text{-}fr\text{-}bounded}(c, f, \{y \in \operatorname{vec}(c) : \exists i \in g.\; (y = h(i)(x))\}) \Rightarrow \\ \quad \quad \operatorname{is\text{-}fr\text{-}equicont}(a, b, c, f, g, h) \end{array}}}

\vnbbookfile{theorem-library/rake-asm-1}{theorem-library/rake-asm-1}
\noindent{\small rake-asm-1.scm -- RAKE-ASM agent RA-1 (2026-10-12): elementary assumptions PROVEN.}\par
\vnbprop{normed-ag-nrm-type}{\vnbstmt{\forall s.\; (\operatorname{is\text{-}normed\text{-}ag}(s) \Rightarrow (\operatorname{nrm}(s) \in (\operatorname{carr}(s) \to \mathbb{R})))}}
\vnbprop{normed-ag-nrm-nonneg}{\vnbstmt{\forall s.; \forall a \in \operatorname{carr}(s).\; (\operatorname{is\text{-}normed\text{-}ag}(s) \Rightarrow (0 \leq \operatorname{nrm}(s)(a)))}}
\vnbprop{normed-ag-nrm-definite}{\vnbstmt{\begin{array}{@{}l@{}} \forall s.; \forall a \in \operatorname{carr}(s). \\ \quad \operatorname{is\text{-}normed\text{-}ag}(s) \Rightarrow ((\operatorname{nrm}(s)(a) = 0) \Leftrightarrow (a = \operatorname{iden}(s))) \end{array}}}
\vnbprop{normed-ag-nrm-subadditive}{\vnbstmt{\begin{array}{@{}l@{}} \forall s.; \forall a, b \in \operatorname{carr}(s). \\ \quad \operatorname{is\text{-}normed\text{-}ag}(s) \Rightarrow (\operatorname{nrm}(s)(\operatorname{opr}(s)(a, b)) \leq (\operatorname{nrm}(s)(a) + \operatorname{nrm}(s)(b))) \end{array}}}
\vnbprop{rr-archimedean-floor}{\vnbstmt{\begin{array}{@{}l@{}} \forall r.; \forall s \in \mathbb{R}. \\ \quad (\operatorname{pos\text{-}rr}(r) \wedge (0 \leq s)) \Rightarrow \exists w \in \mathbb{N}.\; ((s < (w \cdot r)) \wedge (((w \cdot r) - r) \leq s)) \end{array}}}
\vnbprop{qq-dense-in-rr}{\vnbstmt{\begin{array}{@{}l@{}} \forall s, x. \\ \quad (\operatorname{pos\text{-}rr}(s) \wedge (x \in \mathbb{R})) \Rightarrow \exists a \in \mathbb{Q}.\; (((a - s) < x) \wedge (x < (a + s))) \end{array}}}

\vnbbookfile{theorem-library/rake-asm-b2}{theorem-library/rake-asm-b2}
\noindent{\small rake-asm-b2.scm -- RAKE-ASM-B agent RB-2 (2026-10-12): four assumptions PROVEN.}\par
\vnbprop{totally-bounded-has-cauchy-subseq-rad}{\vnbstmt{\begin{array}{@{}l@{}} \forall s.; \forall f \in (\mathbb{N} \to \operatorname{pts}(s)), d. \\ \quad (\operatorname{totally\text{-}bounded}(s) \wedge \operatorname{null\text{-}rr\text{-}seq}(d)) \Rightarrow  \\ \quad \quad \exists i.\; (\operatorname{strictly\text{-}mono\text{-}nn}(i) \wedge \operatorname{is\text{-}cauchy\text{-}seq}(s, \operatorname{subseq}(f, i))) \end{array}}}
\vnbprop{summable-bound-converges}{\vnbstmt{\begin{array}{@{}l@{}} \forall s, f, d. \\ \quad \operatorname{is\text{-}complete}(s) \wedge \\ \quad f \in (\mathbb{N} \to \operatorname{pts}(s)) \wedge \\ \quad d \in (\mathbb{N} \to \mathbb{R}) \wedge \\ \quad \operatorname{series\text{-}converges}(d) \wedge \\ \quad \forall k \in \mathbb{N}.\; (\operatorname{dist}(s)(f(k), f((k + 1))) \leq d(k)) \Rightarrow \\ \quad \quad \operatorname{converges}(s, f) \end{array}}}
\vnbprop{sums-to-unique}{\vnbstmt{\begin{array}{@{}l@{}} \forall p.; \forall f \in (\operatorname{dom}(f) \to \operatorname{carr}(p)), a, b \in \operatorname{carr}(p). \\ \quad (\operatorname{is\text{-}normed\text{-}ag}(p) \wedge (\operatorname{sums\text{-}to}(p, f, a) \wedge \operatorname{sums\text{-}to}(p, f, b))) \Rightarrow (a = b) \end{array}}}
\vnbprop{rb2-nrm-finsum-le}{\vnbstmt{\begin{array}{@{}l@{}} \forall p, f, g, d.; \forall w \in \mathbf{Set}. \\ \quad \operatorname{is\text{-}normed\text{-}ag}(p) \wedge \\ \quad \forall z \in d.\; (f(z) \in \operatorname{carr}(p)) \wedge \\ \quad \forall z \in d.\; (g(z) = \operatorname{nrm}(p)(f(z))) \wedge \\ \quad |w| \in \mathbb{N} \wedge \\ \quad w \subseteq d \Rightarrow \\ \quad \quad \operatorname{finsum}(\operatorname{rr\text{-}pos\text{-}star\text{-}add\text{-}monoid}, g, w) \in \mathbb{R} \wedge \\ \quad \quad \operatorname{nrm}(p)(\operatorname{finsum}(\operatorname{normed\text{-}ag\text{-}as\text{-}abelian\text{-}group}(p), f, w)) \leq \operatorname{finsum}(\operatorname{rr\text{-}pos\text{-}star\text{-}add\text{-}monoid}, g, w) \end{array}}}
\vnbprop{rb2-dist-union}{\vnbstmt{\begin{array}{@{}l@{}} \forall p, f, g, d.; \forall a, c \in \mathbf{Set}. \\ \quad \operatorname{is\text{-}normed\text{-}ag}(p) \wedge \\ \quad \forall z \in d.\; (f(z) \in \operatorname{carr}(p)) \wedge \\ \quad \forall z \in d.\; (g(z) = \operatorname{nrm}(p)(f(z))) \wedge \\ \quad |a| \in \mathbb{N} \wedge \\ \quad |c| \in \mathbb{N} \wedge \\ \quad \forall z \in c.\; \neg (z \in a) \wedge \\ \quad a \subseteq d \wedge \\ \quad c \subseteq d \Rightarrow \\ \quad \quad \operatorname{finsum}(\operatorname{rr\text{-}pos\text{-}star\text{-}add\text{-}monoid}, g, (a \cup c)) \\ \quad \quad =\; \operatorname{finsum}(\operatorname{rr\text{-}pos\text{-}star\text{-}add\text{-}monoid}, g, a) + \operatorname{finsum}(\operatorname{rr\text{-}pos\text{-}star\text{-}add\text{-}monoid}, g, c) \wedge \\ \quad \quad \operatorname{dist}(\operatorname{nag\text{-}metric\text{-}space}(p))(\operatorname{finsum}(\operatorname{normed\text{-}ag\text{-}as\text{-}abelian\text{-}group}(p), f, a), \\ \quad \quad \quad \operatorname{finsum}(\operatorname{normed\text{-}ag\text{-}as\text{-}abelian\text{-}group}(p), f, (a \cup c))) \\ \quad \quad \leq\; \operatorname{finsum}(\operatorname{rr\text{-}pos\text{-}star\text{-}add\text{-}monoid}, g, c) \end{array}}}
\vnbprop{rb2-dist-subset}{\vnbstmt{\begin{array}{@{}l@{}} \forall p, f, g, d.; \forall a, b \in \mathbf{Set}. \\ \quad \operatorname{is\text{-}normed\text{-}ag}(p) \wedge \\ \quad \forall z \in d.\; (f(z) \in \operatorname{carr}(p)) \wedge \\ \quad \forall z \in d.\; (g(z) = \operatorname{nrm}(p)(f(z))) \wedge \\ \quad |a| \in \mathbb{N} \wedge \\ \quad |b| \in \mathbb{N} \wedge \\ \quad a \subseteq b \wedge \\ \quad b \subseteq d \Rightarrow \\ \quad \quad \operatorname{dist}(\operatorname{nag\text{-}metric\text{-}space}(p))(\operatorname{finsum}(\operatorname{normed\text{-}ag\text{-}as\text{-}abelian\text{-}group}(p), f, a), \\ \quad \quad \quad \operatorname{finsum}(\operatorname{normed\text{-}ag\text{-}as\text{-}abelian\text{-}group}(p), f, b)) \\ \quad \quad \leq\; (\operatorname{finsum}(\operatorname{rr\text{-}pos\text{-}star\text{-}add\text{-}monoid}, g, b) - \operatorname{finsum}(\operatorname{rr\text{-}pos\text{-}star\text{-}add\text{-}monoid}, g, a)) \end{array}}}
\vnbprop{rb2-esum-approx}{\vnbstmt{\begin{array}{@{}l@{}} \forall g \in (\operatorname{dom}(g) \to \operatorname{rr\text{-}pos\text{-}star}), l. \\ \quad ((\operatorname{esum}(g) \in \mathbb{R}) \wedge \operatorname{pos\text{-}rr}(l)) \Rightarrow  \\ \quad \quad \exists t \in \mathbf{Set}. \\ \quad \quad \quad ((|t| \in \mathbb{N}) \wedge ((t \subseteq \operatorname{dom}(g)) \wedge (\operatorname{esum}(g) < (\operatorname{finsum}(\operatorname{rr\text{-}pos\text{-}star\text{-}add\text{-}monoid}, g, t) + l)))) \end{array}}}
\vnbprop{rb2-choice-seq}{\vnbstmt{\begin{array}{@{}l@{}} \forall g \in (\operatorname{dom}(g) \to \operatorname{rr\text{-}pos\text{-}star}), d. \\ \quad ((\operatorname{esum}(g) \in \mathbb{R}) \wedge \forall k \in \mathbb{N}.\; \operatorname{pos\text{-}rr}(d(k))) \Rightarrow  \\ \quad \quad \exists s. \forall n \in \mathbb{N}. \\ \quad \quad \quad s(n) \in \mathbf{Set} \wedge \\ \quad \quad \quad |s(n)| \in \mathbb{N} \wedge \\ \quad \quad \quad s(n) \subseteq \operatorname{dom}(g) \wedge \\ \quad \quad \quad \operatorname{esum}(g) < (\operatorname{finsum}(\operatorname{rr\text{-}pos\text{-}star\text{-}add\text{-}monoid}, g, s(n)) + d(n)) \end{array}}}
\vnbprop{absolute-summable-implies-summable}{\vnbstmt{\begin{array}{@{}l@{}} \forall p.; \forall f \in (\operatorname{dom}(f) \to \operatorname{carr}(p)). \\ \quad ((\operatorname{is\text{-}normed\text{-}ag}(p) \wedge \operatorname{is\text{-}complete}(\operatorname{nag\text{-}metric\text{-}space}(p))) \wedge \operatorname{is\text{-}absolutely\text{-}summable}(p, f)) \Rightarrow  \\ \quad \quad \operatorname{is\text{-}summable}(p, f) \end{array}}}

\vnbbookfile{theorem-library/primitive-glue}{theorem-library/primitive-glue}
\noindent{\small primitive-glue.scm -- PRIMITIVES ON ADJACENT INTERVALS GLUE (batch 22-C).}\par
\vnbprop{clamp-left-const}{\vnbstmt{\begin{array}{@{}l@{}} \forall a, b, x \in \mathbb{R}. \\ \quad ((x \leq a) \wedge (a \leq b)) \Rightarrow (\operatorname{clamp}(a, b, x) = a) \end{array}}}
\vnbprop{clamp-right-const}{\vnbstmt{\begin{array}{@{}l@{}} \forall a, b, x \in \mathbb{R}. \\ \quad ((a \leq b) \wedge (b \leq x)) \Rightarrow (\operatorname{clamp}(a, b, x) = b) \end{array}}}
\vnbprop{primitive-integrand-interior-congruence}{\vnbstmt{\begin{array}{@{}l@{}} \forall f, i, \mathit{ppsi}, a, b. \\ \quad \operatorname{is\text{-}primitive}(f, i, a, b) \wedge \\ \quad \mathit{ppsi} \in (\operatorname{ccint}(a, b) \to \mathbb{R}) \wedge \\ \quad \forall y \in \operatorname{ooint}(a, b).\; (i(y) \simeq \operatorname{ppsi}(y)) \Rightarrow \\ \quad \quad \operatorname{is\text{-}primitive}(f, \mathit{ppsi}, a, b) \end{array}}}
\vnbprop{primitive-glue}{\vnbstmt{\begin{array}{@{}l@{}} \forall a, c, b \in \mathbb{R}, d, h, k, m.; \forall g, f \in (\operatorname{ccint}(a, b) \to \mathbb{R}). \\ \quad a < c \wedge \\ \quad c < b \wedge \\ \quad \operatorname{is\text{-}primitive}(d, h, a, c) \wedge \\ \quad \operatorname{is\text{-}primitive}(k, m, c, b) \wedge \\ \quad d(c) = k(c) \wedge \\ \quad \forall t \in \operatorname{ccint}(a, c).\; (g(t) \simeq d(t)) \wedge \\ \quad \forall t \in \operatorname{ccint}(c, b).\; (g(t) \simeq k(t)) \wedge \\ \quad \forall t \in \operatorname{ooint}(a, c).\; (f(t) \simeq h(t)) \wedge \\ \quad \forall t \in \operatorname{ooint}(c, b).\; (f(t) \simeq m(t)) \Rightarrow \\ \quad \quad \operatorname{is\text{-}primitive}(g, f, a, b) \end{array}}}
\vnbprop{primitive-of-constant-piece}{\vnbstmt{\begin{array}{@{}l@{}} \forall u, v \in \mathbb{R}, f \in (\mathbb{R} \to \mathbb{R}), a, k \in \mathbb{R}. \\ \quad ((u < v) \wedge \forall t \in \mathbb{R}.\; (((u < t) \wedge (t < v)) \Rightarrow (f(t) = a))) \Rightarrow  \\ \quad \quad \operatorname{is\text{-}primitive}((\lambda s \in \operatorname{ccint}(u, v).\; (k + (a \cdot (s - u)))), \operatorname{restrict}(f, \operatorname{ccint}(u, v)), u, v) \end{array}}}
\vnbprop{step-fn-has-primitive-ind}{\vnbstmt{\begin{array}{@{}l@{}} \forall n \in \mathbb{N}, p, a, b.; \forall f \in (\mathbb{R} \to \mathbb{R}). \\ \quad \operatorname{is\text{-}partition}(p, n, a, b) \wedge \\ \quad \forall i \in \operatorname{interval}(0, n). \\ \quad \quad (i < n) \Rightarrow  \\ \quad \quad \quad \exists c \in \mathbb{R}. \forall t \in \mathbb{R}. \\ \quad \quad \quad \quad ((p(i) < t) \wedge (t < p((i + 1)))) \Rightarrow (f(t) = c) \Rightarrow \\ \quad \quad \exists g.\; \operatorname{is\text{-}primitive}(g, \operatorname{restrict}(f, \operatorname{ccint}(a, b)), a, b) \end{array}}}
\vnbprop{step-fn-has-primitive}{\vnbstmt{\begin{array}{@{}l@{}} \forall a, b \in \mathbb{R}, f. \\ \quad \operatorname{is\text{-}step\text{-}fn}(f, a, b) \Rightarrow  \\ \quad \quad \exists g.\; \operatorname{is\text{-}primitive}(g, \operatorname{restrict}(f, \operatorname{ccint}(a, b)), a, b) \end{array}}}

\vnbbookfile{theorem-library/pw-int-laws-3}{theorem-library/pw-int-laws-3}
\noindent{\small pw-int-laws-3.scm -- LINEARITY over the reals as ONE law, the COMPLEX scalar, and the additivity of the complex integral over adjacent intervals.  Continues theorem-library/pw-int-laws-2.scm.}\par
\vnbprop{pw-antiderivative-linear}{\vnbstmt{\begin{array}{@{}l@{}} \forall a, b, c, d \in \mathbb{R}, f, w, i, \mathit{ppsi}.; \forall h, \mathit{pchi} \in (\operatorname{ccint}(a, b) \to \mathbb{R}). \\ \quad \operatorname{is\text{-}primitive}(f, i, a, b) \wedge \\ \quad \operatorname{is\text{-}primitive}(w, \mathit{ppsi}, a, b) \wedge \\ \quad \forall y \in \operatorname{ccint}(a, b).\; (h(y) \simeq ((c \cdot f(y)) + (d \cdot w(y)))) \wedge \\ \quad \forall y \in \operatorname{ccint}(a, b).\; (\operatorname{pchi}(y) \simeq ((c \cdot i(y)) + (d \cdot \operatorname{ppsi}(y)))) \Rightarrow \\ \quad \quad (\operatorname{is\text{-}primitive}(h, \mathit{pchi}, a, b) \wedge (\int_{a}^{b} \mathit{pchi} = ((c \cdot \int_{a}^{b} i) + (d \cdot \int_{a}^{b} \mathit{ppsi})))) \end{array}}}
\vnbprop{cc-int-complex-mul}{\vnbstmt{\begin{array}{@{}l@{}} \forall a, b, i, \mathit{pchi}.; \forall w \in \mathbb{C}, f, \mathit{paw}. \\ \quad a \in \mathbb{R} \wedge \\ \quad b \in \mathbb{R} \wedge \\ \quad i \in (\operatorname{ccint}(a, b) \to \mathbb{C}) \wedge \\ \quad \mathit{pchi} \in (\operatorname{ccint}(a, b) \to \mathbb{C}) \wedge \\ \quad \forall y \in \operatorname{ccint}(a, b).\; (\operatorname{pchi}(y) \simeq (w \cdot i(y))) \wedge \\ \quad \operatorname{is\text{-}primitive}(f, (\lambda t \in \operatorname{ccint}(a, b).\; \operatorname{real\text{-}part}(i(t))), a, b) \wedge \\ \quad \operatorname{is\text{-}primitive}(\mathit{paw}, (\lambda t \in \operatorname{ccint}(a, b).\; \operatorname{imag\text{-}part}(i(t))), a, b) \Rightarrow \\ \quad \quad \int_{a}^{b} \mathit{pchi} = (w \cdot \int_{a}^{b} i) \end{array}}}
\vnbprop{cc-int-adjacent}{\vnbstmt{\begin{array}{@{}l@{}} \forall a, b.; \forall v \in \operatorname{ooint}(a, b), i, f, w. \\ \quad a \in \mathbb{R} \wedge \\ \quad b \in \mathbb{R} \wedge \\ \quad i \in (\operatorname{ccint}(a, b) \to \mathbb{C}) \wedge \\ \quad \operatorname{is\text{-}primitive}(f, (\lambda t \in \operatorname{ccint}(a, b).\; \operatorname{real\text{-}part}(i(t))), a, b) \wedge \\ \quad \operatorname{is\text{-}primitive}(w, (\lambda t \in \operatorname{ccint}(a, b).\; \operatorname{imag\text{-}part}(i(t))), a, b) \Rightarrow \\ \quad \quad \int_{a}^{b} i = (\int_{a}^{v} \operatorname{restrict}(i, \operatorname{ccint}(a, v)) + \int_{v}^{b} \operatorname{restrict}(i, \operatorname{ccint}(v, b))) \end{array}}}

\vnbbookfile{theorem-library/primitive-uniform-limit}{theorem-library/primitive-uniform-limit}
\noindent{\small theorem-library/primitive-uniform-limit.scm -- Dieudonne (8.6.4) on a compact interval [a,b], with COUNTABLE exceptional sets: A UNIFORM LIMIT OF PRIMITIVES IS A PRIMITIVE.  Batch 22-B, 2026-09-23.}\par
\vnbprop{countable-union-nn}{\vnbstmt{\begin{array}{@{}l@{}} \forall d.; \forall u \in \mathbf{Set}. \\ \quad (\forall k \in \mathbb{N}.\; \operatorname{is\text{-}countable}(d(k)) \wedge \forall y \in u.\; \exists k \in \mathbb{N}.\; (y \in d(k))) \Rightarrow  \\ \quad \quad \operatorname{is\text{-}countable}(u) \end{array}}}
\vnbprop{primitive-pair-lipschitz}{\vnbstmt{\begin{array}{@{}l@{}} \forall a, b, g_{1}, f_{1}, g_{2}, f_{2}.; \forall m \in \mathbb{R}, x, y \in \operatorname{ccint}(a, b). \\ \quad \operatorname{is\text{-}primitive}(g_{1}, f_{1}, a, b) \wedge \\ \quad \operatorname{is\text{-}primitive}(g_{2}, f_{2}, a, b) \wedge \\ \quad \forall x \in \operatorname{ccint}(a, b).\; (\lvert (\operatorname{f1}(x) - \operatorname{f2}(x)) \rvert \leq m) \wedge \\ \quad x < y \Rightarrow \\ \quad \quad \lvert ((\operatorname{g1}(y) - \operatorname{g2}(y)) - (\operatorname{g1}(x) - \operatorname{g2}(x))) \rvert \leq (m \cdot (y - x)) \end{array}}}
\vnbprop{primitive-pair-lipschitz-abs}{\vnbstmt{\begin{array}{@{}l@{}} \forall a, b, g_{1}, f_{1}, g_{2}, f_{2}.; \forall m \in \mathbb{R}, x, y \in \operatorname{ccint}(a, b). \\ \quad \operatorname{is\text{-}primitive}(g_{1}, f_{1}, a, b) \wedge \\ \quad \operatorname{is\text{-}primitive}(g_{2}, f_{2}, a, b) \wedge \\ \quad \forall x \in \operatorname{ccint}(a, b).\; (\lvert (\operatorname{f1}(x) - \operatorname{f2}(x)) \rvert \leq m) \Rightarrow \\ \quad \quad \lvert ((\operatorname{g1}(y) - \operatorname{g2}(y)) - (\operatorname{g1}(x) - \operatorname{g2}(x))) \rvert \leq (m \cdot \lvert (y - x) \rvert) \end{array}}}
\vnbprop{primitive-family-limit}{\vnbstmt{\begin{array}{@{}l@{}} \forall a, b \in \mathbb{R}, m, c \in (\mathbb{N} \to (\operatorname{ccint}(a, b) \to \mathbb{R})). \\ \quad a < b \wedge \\ \quad \forall k \in \mathbb{N}.\; \operatorname{is\text{-}primitive}(m(k), c(k), a, b) \wedge \\ \quad \forall s. \\ \quad \quad \operatorname{pos\text{-}rr}(s) \Rightarrow  \\ \quad \quad \quad \exists n \in \mathbb{N}. \forall k, l \in \mathbb{N}, x \in \operatorname{ccint}(a, b). \\ \quad \quad \quad \quad ((n \leq k) \wedge (n \leq l)) \Rightarrow (\lvert (c(k)(x) - c(l)(x)) \rvert \leq s) \wedge \\ \quad \forall k \in \mathbb{N}.\; (m(k)(a) = 0) \Rightarrow \\ \quad \quad \exists \mathit{pufc} \in (\mathbb{R} \to \mathbb{R}). \\ \quad \quad \quad \forall x \in \operatorname{ccint}(a, b). \\ \quad \quad \quad \quad \operatorname{converges\text{-}to}(\operatorname{rr\text{-}ms}, (\lambda k \in \mathbb{N}.\; m(k)(x)), \operatorname{pufc}(x)) \wedge \\ \quad \quad \quad \forall t \in \mathbb{R}.\; \operatorname{is\text{-}continuous\text{-}at}(\operatorname{rr\text{-}ms}, \operatorname{rr\text{-}ms}, \mathit{pufc}, t) \end{array}}}
\vnbprop{primitive-family-deriv-at}{\vnbstmt{\begin{array}{@{}l@{}} \forall a, b \in \mathbb{R}, m, c \in (\mathbb{N} \to (\operatorname{ccint}(a, b) \to \mathbb{R})), \mathit{pufc} \in (\mathbb{R} \to \mathbb{R}), t \in \operatorname{ooint}(a, b), v \in \mathbb{R}. \\ \quad a < b \wedge \\ \quad \forall k \in \mathbb{N}.\; \operatorname{is\text{-}primitive}(m(k), c(k), a, b) \wedge \\ \quad \forall s. \\ \quad \quad \operatorname{pos\text{-}rr}(s) \Rightarrow  \\ \quad \quad \quad \exists n \in \mathbb{N}. \forall k, l \in \mathbb{N}, x \in \operatorname{ccint}(a, b). \\ \quad \quad \quad \quad ((n \leq k) \wedge (n \leq l)) \Rightarrow (\lvert (c(k)(x) - c(l)(x)) \rvert \leq s) \wedge \\ \quad \forall x \in \operatorname{ccint}(a, b). \\ \quad \quad \operatorname{converges\text{-}to}(\operatorname{rr\text{-}ms}, (\lambda k \in \mathbb{N}.\; m(k)(x)), \operatorname{pufc}(x)) \wedge \\ \quad \forall k \in \mathbb{N}.\; \operatorname{has\text{-}deriv\text{-}at}(m(k), t, c(k)(t)) \wedge \\ \quad \operatorname{converges\text{-}to}(\operatorname{rr\text{-}ms}, (\lambda k \in \mathbb{N}.\; c(k)(t)), v) \Rightarrow \\ \quad \quad \operatorname{is\text{-}diff\text{-}at}(\mathit{pufc}, t, v) \end{array}}}
\vnbprop{primitive-uniform-limit}{\vnbstmt{\begin{array}{@{}l@{}} \forall a, b \in \mathbb{R}, m, c \in (\mathbb{N} \to (\operatorname{ccint}(a, b) \to \mathbb{R})), f \in (\operatorname{ccint}(a, b) \to \mathbb{R}). \\ \quad a < b \wedge \\ \quad \forall k \in \mathbb{N}.\; \operatorname{is\text{-}primitive}(m(k), c(k), a, b) \wedge \\ \quad \forall s. \\ \quad \quad \operatorname{pos\text{-}rr}(s) \Rightarrow  \\ \quad \quad \quad \exists n \in \mathbb{N}. \forall k \in \mathbb{N}, x \in \operatorname{ccint}(a, b). \\ \quad \quad \quad \quad (n \leq k) \Rightarrow (\lvert (c(k)(x) - f(x)) \rvert \leq s) \wedge \\ \quad \forall k \in \mathbb{N}.\; (m(k)(a) = 0) \Rightarrow \\ \quad \quad \exists g. \\ \quad \quad \quad \operatorname{is\text{-}primitive}(g, f, a, b) \wedge \\ \quad \quad \quad \forall x \in \operatorname{ccint}(a, b). \\ \quad \quad \quad \quad \operatorname{converges\text{-}to}(\operatorname{rr\text{-}ms}, (\lambda k \in \mathbb{N}.\; m(k)(x)), g(x)) \end{array}}}

\vnbbookfile{theorem-library/dual-nvs}{theorem-library/dual-nvs}
\noindent{\small dual-nvs.scm -- THE DUAL OF A REAL NORMED VECTOR SPACE IS A NORMED VECTOR SPACE. Agent DUAL-FIX, 2026-10-26.  DESTINATION: theorem-library/dual-nvs.scm.  Library: analysis. Helper prefix: dnv-.}\par
\vnbprop{dual-vec-in-set}{\vnbstmt{\forall m.\; (\operatorname{is\text{-}normed\text{-}vector\text{-}space}(m) \Rightarrow (\operatorname{dual\text{-}vec}(m) \in \mathbf{Set}))}}
\vnbprop{dual-vec-fun}{\vnbstmt{\forall m.; \forall f \in \operatorname{dual\text{-}vec}(m).\; (f \in (\operatorname{vec}(m) \to \mathbb{R}))}}
\vnbprop{dual-vec-blf}{\vnbstmt{\forall m.; \forall f \in \operatorname{dual\text{-}vec}(m).\; \operatorname{is\text{-}bounded\text{-}linear\text{-}functional}(m, f)}}
\vnbprop{dual-vec-intro}{\vnbstmt{\begin{array}{@{}l@{}} \forall m.; \forall f \in (\operatorname{vec}(m) \to \mathbb{R}). \\ \quad \operatorname{is\text{-}bounded\text{-}linear\text{-}functional}(m, f) \Rightarrow (f \in \operatorname{dual\text{-}vec}(m)) \end{array}}}
\vnbprop{dual-vec-additive}{\vnbstmt{\begin{array}{@{}l@{}} \forall m, f, x, y. \\ \quad ((f \in \operatorname{dual\text{-}vec}(m)) \wedge ((x \in \operatorname{vec}(m)) \wedge (y \in \operatorname{vec}(m)))) \Rightarrow  \\ \quad \quad f(\operatorname{vadd}(m)(x, y)) = (f(x) + f(y)) \end{array}}}
\vnbprop{dual-vec-homogeneous}{\vnbstmt{\begin{array}{@{}l@{}} \forall m, f, r, x. \\ \quad ((f \in \operatorname{dual\text{-}vec}(m)) \wedge ((r \in \mathbb{R}) \wedge (x \in \operatorname{vec}(m)))) \Rightarrow  \\ \quad \quad f(\operatorname{act}(m)(r, x)) = (r \cdot f(x)) \end{array}}}
\vnbprop{dual-vadd-readoff}{\vnbstmt{\begin{array}{@{}l@{}} \forall m, f, g. \\ \quad (\operatorname{is\text{-}normed\text{-}vector\text{-}space}(m) \wedge ((f \in \operatorname{dual\text{-}vec}(m)) \wedge (g \in \operatorname{dual\text{-}vec}(m)))) \Rightarrow  \\ \quad \quad (\operatorname{vadd}({m}^{*})(f, g) \simeq (\lambda x \in \operatorname{vec}(m).\; (f(x) + g(x)))) \end{array}}}
\vnbprop{dual-act-readoff}{\vnbstmt{\begin{array}{@{}l@{}} \forall m, r, f. \\ \quad (\operatorname{is\text{-}normed\text{-}vector\text{-}space}(m) \wedge ((r \in \mathbb{R}) \wedge (f \in \operatorname{dual\text{-}vec}(m)))) \Rightarrow  \\ \quad \quad (\operatorname{act}({m}^{*})(r, f) \simeq (\lambda x \in \operatorname{vec}(m).\; (r \cdot f(x)))) \end{array}}}
\vnbprop{dual-vneg-readoff}{\vnbstmt{\begin{array}{@{}l@{}} \forall m, f. \\ \quad (\operatorname{is\text{-}normed\text{-}vector\text{-}space}(m) \wedge (f \in \operatorname{dual\text{-}vec}(m))) \Rightarrow  \\ \quad \quad (\operatorname{vneg}({m}^{*})(f) \simeq (\lambda x \in \operatorname{vec}(m).\; -f(x))) \end{array}}}
\vnbprop{dual-vzero-readoff}{\vnbstmt{\begin{array}{@{}l@{}} \forall m. \\ \quad \operatorname{is\text{-}normed\text{-}vector\text{-}space}(m) \Rightarrow (\operatorname{vzero}({m}^{*}) \simeq (\lambda x \in \operatorname{vec}(m).\; 0)) \end{array}}}
\vnbprop{dual-lam-add-fun}{\vnbstmt{\begin{array}{@{}l@{}} \forall m, f, g. \\ \quad (\operatorname{is\text{-}normed\text{-}vector\text{-}space}(m) \wedge ((f \in \operatorname{dual\text{-}vec}(m)) \wedge (g \in \operatorname{dual\text{-}vec}(m)))) \Rightarrow  \\ \quad \quad (\lambda x \in \operatorname{vec}(m).\; (f(x) + g(x))) \in (\operatorname{vec}(m) \to \mathbb{R}) \end{array}}}
\vnbprop{dual-lam-act-fun}{\vnbstmt{\begin{array}{@{}l@{}} \forall m, r, f. \\ \quad (\operatorname{is\text{-}normed\text{-}vector\text{-}space}(m) \wedge ((r \in \mathbb{R}) \wedge (f \in \operatorname{dual\text{-}vec}(m)))) \Rightarrow  \\ \quad \quad (\lambda x \in \operatorname{vec}(m).\; (r \cdot f(x))) \in (\operatorname{vec}(m) \to \mathbb{R}) \end{array}}}
\vnbprop{dual-lam-neg-fun}{\vnbstmt{\begin{array}{@{}l@{}} \forall m, f. \\ \quad (\operatorname{is\text{-}normed\text{-}vector\text{-}space}(m) \wedge (f \in \operatorname{dual\text{-}vec}(m))) \Rightarrow  \\ \quad \quad (\lambda x \in \operatorname{vec}(m).\; -f(x)) \in (\operatorname{vec}(m) \to \mathbb{R}) \end{array}}}
\vnbprop{dual-lam-zero-fun}{\vnbstmt{\begin{array}{@{}l@{}} \forall m. \\ \quad \operatorname{is\text{-}normed\text{-}vector\text{-}space}(m) \Rightarrow ((\lambda x \in \operatorname{vec}(m).\; 0) \in (\operatorname{vec}(m) \to \mathbb{R})) \end{array}}}
\vnbprop{dual-lam-add-linear}{\vnbstmt{\begin{array}{@{}l@{}} \forall m, f, g. \\ \quad (\operatorname{is\text{-}normed\text{-}vector\text{-}space}(m) \wedge ((f \in \operatorname{dual\text{-}vec}(m)) \wedge (g \in \operatorname{dual\text{-}vec}(m)))) \Rightarrow  \\ \quad \quad \operatorname{is\text{-}linear\text{-}functional}(m, (\lambda x \in \operatorname{vec}(m).\; (f(x) + g(x)))) \end{array}}}
\vnbprop{dual-lam-act-linear}{\vnbstmt{\begin{array}{@{}l@{}} \forall m, r, f. \\ \quad (\operatorname{is\text{-}normed\text{-}vector\text{-}space}(m) \wedge ((r \in \mathbb{R}) \wedge (f \in \operatorname{dual\text{-}vec}(m)))) \Rightarrow  \\ \quad \quad \operatorname{is\text{-}linear\text{-}functional}(m, (\lambda x \in \operatorname{vec}(m).\; (r \cdot f(x)))) \end{array}}}
\vnbprop{dual-lam-neg-linear}{\vnbstmt{\begin{array}{@{}l@{}} \forall m, f. \\ \quad (\operatorname{is\text{-}normed\text{-}vector\text{-}space}(m) \wedge (f \in \operatorname{dual\text{-}vec}(m))) \Rightarrow  \\ \quad \quad \operatorname{is\text{-}linear\text{-}functional}(m, (\lambda x \in \operatorname{vec}(m).\; -f(x))) \end{array}}}
\vnbprop{dual-lam-zero-linear}{\vnbstmt{\begin{array}{@{}l@{}} \forall m. \\ \quad \operatorname{is\text{-}normed\text{-}vector\text{-}space}(m) \Rightarrow  \\ \quad \quad \operatorname{is\text{-}linear\text{-}functional}(m, (\lambda x \in \operatorname{vec}(m).\; 0)) \end{array}}}
\vnbprop{dual-rr-le-sum-of-bounds}{\vnbstmt{\begin{array}{@{}l@{}} \forall a, b, c, p, q. \\ \quad ((a \in \mathbb{R}) \wedge ((b \in \mathbb{R}) \wedge ((c \in \mathbb{R}) \wedge ((p \in \mathbb{R}) \wedge ((q \in \mathbb{R}) \wedge ((c \leq (a + b)) \wedge ((a \leq p) \wedge (b \leq q)))))))) \Rightarrow  \\ \quad \quad c \leq (p + q) \end{array}}}
\vnbprop{dual-lam-add-bound}{\vnbstmt{\begin{array}{@{}l@{}} \forall m, f, g.; \forall w \in \operatorname{vec}(m). \\ \quad (\operatorname{is\text{-}normed\text{-}vector\text{-}space}(m) \wedge ((f \in \operatorname{dual\text{-}vec}(m)) \wedge (g \in \operatorname{dual\text{-}vec}(m)))) \Rightarrow  \\ \quad \quad \lvert (\lambda x \in \operatorname{vec}(m).\; (f(x) + g(x)))(w) \rvert \leq ((\operatorname{dual\text{-}norm}(m, f) + \operatorname{dual\text{-}norm}(m, g)) \cdot \operatorname{vnrm}(m)(w)) \end{array}}}
\vnbprop{dual-lam-neg-bound}{\vnbstmt{\begin{array}{@{}l@{}} \forall m, f.; \forall w \in \operatorname{vec}(m). \\ \quad (\operatorname{is\text{-}normed\text{-}vector\text{-}space}(m) \wedge (f \in \operatorname{dual\text{-}vec}(m))) \Rightarrow  \\ \quad \quad \lvert (\lambda x \in \operatorname{vec}(m).\; -f(x))(w) \rvert \leq (\operatorname{dual\text{-}norm}(m, f) \cdot \operatorname{vnrm}(m)(w)) \end{array}}}
\vnbprop{dual-lam-zero-bound}{\vnbstmt{\begin{array}{@{}l@{}} \forall m.; \forall w \in \operatorname{vec}(m). \\ \quad \operatorname{is\text{-}normed\text{-}vector\text{-}space}(m) \Rightarrow  \\ \quad \quad \lvert (\lambda x \in \operatorname{vec}(m).\; 0)(w) \rvert \leq (0 \cdot \operatorname{vnrm}(m)(w)) \end{array}}}
\vnbprop{dual-lam-act-bound}{\vnbstmt{\begin{array}{@{}l@{}} \forall m, r, f.; \forall w \in \operatorname{vec}(m). \\ \quad (\operatorname{is\text{-}normed\text{-}vector\text{-}space}(m) \wedge ((r \in \mathbb{R}) \wedge (f \in \operatorname{dual\text{-}vec}(m)))) \Rightarrow  \\ \quad \quad \lvert (\lambda x \in \operatorname{vec}(m).\; (r \cdot f(x)))(w) \rvert \leq ((\lvert r \rvert \cdot \operatorname{dual\text{-}norm}(m, f)) \cdot \operatorname{vnrm}(m)(w)) \end{array}}}
\vnbprop{dual-lam-add-in-dual}{\vnbstmt{\begin{array}{@{}l@{}} \forall m, f, g. \\ \quad (\operatorname{is\text{-}normed\text{-}vector\text{-}space}(m) \wedge ((f \in \operatorname{dual\text{-}vec}(m)) \wedge (g \in \operatorname{dual\text{-}vec}(m)))) \Rightarrow  \\ \quad \quad (\lambda x \in \operatorname{vec}(m).\; (f(x) + g(x))) \in \operatorname{dual\text{-}vec}(m) \end{array}}}
\vnbprop{dual-lam-neg-in-dual}{\vnbstmt{\begin{array}{@{}l@{}} \forall m, f. \\ \quad (\operatorname{is\text{-}normed\text{-}vector\text{-}space}(m) \wedge (f \in \operatorname{dual\text{-}vec}(m))) \Rightarrow  \\ \quad \quad (\lambda x \in \operatorname{vec}(m).\; -f(x)) \in \operatorname{dual\text{-}vec}(m) \end{array}}}
\vnbprop{dual-lam-zero-in-dual}{\vnbstmt{\begin{array}{@{}l@{}} \forall m. \\ \quad \operatorname{is\text{-}normed\text{-}vector\text{-}space}(m) \Rightarrow ((\lambda x \in \operatorname{vec}(m).\; 0) \in \operatorname{dual\text{-}vec}(m)) \end{array}}}
\vnbprop{dual-lam-act-in-dual}{\vnbstmt{\begin{array}{@{}l@{}} \forall m, r, f. \\ \quad (\operatorname{is\text{-}normed\text{-}vector\text{-}space}(m) \wedge ((r \in \mathbb{R}) \wedge (f \in \operatorname{dual\text{-}vec}(m)))) \Rightarrow  \\ \quad \quad (\lambda x \in \operatorname{vec}(m).\; (r \cdot f(x))) \in \operatorname{dual\text{-}vec}(m) \end{array}}}
\vnbprop{dual-vadd-in-dual}{\vnbstmt{\begin{array}{@{}l@{}} \forall m, f, g. \\ \quad (\operatorname{is\text{-}normed\text{-}vector\text{-}space}(m) \wedge ((f \in \operatorname{dual\text{-}vec}(m)) \wedge (g \in \operatorname{dual\text{-}vec}(m)))) \Rightarrow  \\ \quad \quad \operatorname{vadd}({m}^{*})(f, g) \in \operatorname{dual\text{-}vec}(m) \end{array}}}
\vnbprop{dual-act-in-dual}{\vnbstmt{\begin{array}{@{}l@{}} \forall m, r, f. \\ \quad (\operatorname{is\text{-}normed\text{-}vector\text{-}space}(m) \wedge ((r \in \mathbb{R}) \wedge (f \in \operatorname{dual\text{-}vec}(m)))) \Rightarrow  \\ \quad \quad \operatorname{act}({m}^{*})(r, f) \in \operatorname{dual\text{-}vec}(m) \end{array}}}
\vnbprop{dual-vneg-in-dual}{\vnbstmt{\begin{array}{@{}l@{}} \forall m, f. \\ \quad (\operatorname{is\text{-}normed\text{-}vector\text{-}space}(m) \wedge (f \in \operatorname{dual\text{-}vec}(m))) \Rightarrow  \\ \quad \quad \operatorname{vneg}({m}^{*})(f) \in \operatorname{dual\text{-}vec}(m) \end{array}}}
\vnbprop{dual-vzero-in-dual}{\vnbstmt{\forall m.\; (\operatorname{is\text{-}normed\text{-}vector\text{-}space}(m) \Rightarrow (\operatorname{vzero}({m}^{*}) \in \operatorname{dual\text{-}vec}(m)))}}
\vnbprop{dual-vadd-apply}{\vnbstmt{\begin{array}{@{}l@{}} \forall m, f, g, x. \\ \quad (\operatorname{is\text{-}normed\text{-}vector\text{-}space}(m) \wedge ((f \in \operatorname{dual\text{-}vec}(m)) \wedge ((g \in \operatorname{dual\text{-}vec}(m)) \wedge (x \in \operatorname{vec}(m))))) \Rightarrow  \\ \quad \quad \operatorname{vadd}({m}^{*})(f, g)(x) = (f(x) + g(x)) \end{array}}}
\vnbprop{dual-act-apply}{\vnbstmt{\begin{array}{@{}l@{}} \forall m, r, f, x. \\ \quad (\operatorname{is\text{-}normed\text{-}vector\text{-}space}(m) \wedge ((r \in \mathbb{R}) \wedge ((f \in \operatorname{dual\text{-}vec}(m)) \wedge (x \in \operatorname{vec}(m))))) \Rightarrow  \\ \quad \quad \operatorname{act}({m}^{*})(r, f)(x) = (r \cdot f(x)) \end{array}}}
\vnbprop{dual-vneg-apply}{\vnbstmt{\begin{array}{@{}l@{}} \forall m, f, x. \\ \quad (\operatorname{is\text{-}normed\text{-}vector\text{-}space}(m) \wedge ((f \in \operatorname{dual\text{-}vec}(m)) \wedge (x \in \operatorname{vec}(m)))) \Rightarrow  \\ \quad \quad \operatorname{vneg}({m}^{*})(f)(x) = -f(x) \end{array}}}
\vnbprop{dual-vzero-apply}{\vnbstmt{\begin{array}{@{}l@{}} \forall m, x. \\ \quad (\operatorname{is\text{-}normed\text{-}vector\text{-}space}(m) \wedge (x \in \operatorname{vec}(m))) \Rightarrow (\operatorname{vzero}({m}^{*})(x) = 0) \end{array}}}
\vnbprop{dual-vnrm-nonneg}{\vnbstmt{\begin{array}{@{}l@{}} \forall m, f. \\ \quad (\operatorname{is\text{-}normed\text{-}vector\text{-}space}(m) \wedge (f \in \operatorname{dual\text{-}vec}(m))) \Rightarrow (0 \leq \operatorname{vnrm}({m}^{*})(f)) \end{array}}}
\vnbprop{dual-vnrm-triangle}{\vnbstmt{\begin{array}{@{}l@{}} \forall m, f, g. \\ \quad (\operatorname{is\text{-}normed\text{-}vector\text{-}space}(m) \wedge ((f \in \operatorname{dual\text{-}vec}(m)) \wedge (g \in \operatorname{dual\text{-}vec}(m)))) \Rightarrow  \\ \quad \quad \operatorname{vnrm}({m}^{*})(\operatorname{vadd}({m}^{*})(f, g)) \leq (\operatorname{vnrm}({m}^{*})(f) + \operatorname{vnrm}({m}^{*})(g)) \end{array}}}
\vnbprop{dual-norm-zero-is-vzero}{\vnbstmt{\begin{array}{@{}l@{}} \forall m, f. \\ \quad (\operatorname{is\text{-}normed\text{-}vector\text{-}space}(m) \wedge ((f \in \operatorname{dual\text{-}vec}(m)) \wedge (\operatorname{dual\text{-}norm}(m, f) = 0))) \Rightarrow  \\ \quad \quad f = \operatorname{vzero}({m}^{*}) \end{array}}}
\vnbprop{dual-vzero-norm-zero}{\vnbstmt{\begin{array}{@{}l@{}} \forall m, f. \\ \quad (\operatorname{is\text{-}normed\text{-}vector\text{-}space}(m) \wedge ((f \in \operatorname{dual\text{-}vec}(m)) \wedge (f = \operatorname{vzero}({m}^{*})))) \Rightarrow  \\ \quad \quad \operatorname{dual\text{-}norm}(m, f) = 0 \end{array}}}
\vnbprop{dual-vnrm-definite}{\vnbstmt{\begin{array}{@{}l@{}} \forall m, f. \\ \quad (\operatorname{is\text{-}normed\text{-}vector\text{-}space}(m) \wedge (f \in \operatorname{dual\text{-}vec}(m))) \Rightarrow  \\ \quad \quad (\operatorname{vnrm}({m}^{*})(f) = 0) \Leftrightarrow (f = \operatorname{vzero}({m}^{*})) \end{array}}}
\vnbprop{dual-rr-pos-of-nonneg-nonzero}{\vnbstmt{\forall a \in \mathbb{R}.\; (((0 \leq a) \wedge \neg (a = 0)) \Rightarrow (0 < a))}}
\vnbprop{dual-lam-act-norm-ge}{\vnbstmt{\begin{array}{@{}l@{}} \forall m, r, f. \\ \quad (\operatorname{is\text{-}normed\text{-}vector\text{-}space}(m) \wedge ((r \in \mathbb{R}) \wedge (f \in \operatorname{dual\text{-}vec}(m)))) \Rightarrow  \\ \quad \quad (\lvert r \rvert \cdot \operatorname{dual\text{-}norm}(m, f)) \leq \operatorname{dual\text{-}norm}(m, (\lambda x \in \operatorname{vec}(m).\; (r \cdot f(x)))) \end{array}}}
\vnbprop{dual-vnrm-homogeneous}{\vnbstmt{\begin{array}{@{}l@{}} \forall m, r, f. \\ \quad (\operatorname{is\text{-}normed\text{-}vector\text{-}space}(m) \wedge ((r \in \mathbb{R}) \wedge (f \in \operatorname{dual\text{-}vec}(m)))) \Rightarrow  \\ \quad \quad \operatorname{vnrm}({m}^{*})(\operatorname{act}({m}^{*})(r, f)) = (\lvert r \rvert \cdot \operatorname{vnrm}({m}^{*})(f)) \end{array}}}
\vnbprop{dual-vadd-typing}{\vnbstmt{\begin{array}{@{}l@{}} \forall m. \\ \quad \operatorname{is\text{-}normed\text{-}vector\text{-}space}(m) \Rightarrow  \\ \quad \quad \operatorname{vadd}({m}^{*}) \in ((\operatorname{dual\text{-}vec}(m) \times \operatorname{dual\text{-}vec}(m)) \to \operatorname{dual\text{-}vec}(m)) \end{array}}}
\vnbprop{dual-act-typing}{\vnbstmt{\begin{array}{@{}l@{}} \forall m. \\ \quad \operatorname{is\text{-}normed\text{-}vector\text{-}space}(m) \Rightarrow  \\ \quad \quad \operatorname{act}({m}^{*}) \in ((\mathbb{R} \times \operatorname{dual\text{-}vec}(m)) \to \operatorname{dual\text{-}vec}(m)) \end{array}}}
\vnbprop{dual-vneg-typing}{\vnbstmt{\begin{array}{@{}l@{}} \forall m. \\ \quad \operatorname{is\text{-}normed\text{-}vector\text{-}space}(m) \Rightarrow (\operatorname{vneg}({m}^{*}) \in (\operatorname{dual\text{-}vec}(m) \to \operatorname{dual\text{-}vec}(m))) \end{array}}}
\vnbprop{dual-vnrm-typing}{\vnbstmt{\begin{array}{@{}l@{}} \forall m. \\ \quad \operatorname{is\text{-}normed\text{-}vector\text{-}space}(m) \Rightarrow (\operatorname{vnrm}({m}^{*}) \in (\operatorname{dual\text{-}vec}(m) \to \mathbb{R})) \end{array}}}
\vnbprop{dual-is-normed-vector-space}{\vnbstmt{\forall m.\; (\operatorname{is\text{-}normed\text{-}vector\text{-}space}(m) \Rightarrow \operatorname{is\text{-}normed\text{-}vector\text{-}space}({m}^{*}))}}

\vnbbookfile{structure-library/bidual}{structure-library/bidual}
\noindent{\small bidual-vocab.scm -- THE CANONICAL MAP OF A NORMED SPACE INTO ITS BIDUAL.  Vocabulary only. Agent DUAL-FIX, 2026-10-26.  DESTINATION: structure-library/bidual.scm.  Library: analysis. Load window: after structure-library/dual-space.}\par
\par\medskip\noindent\textbf{Definition} (\texttt{bidual-eval}).\;$\iota_{m}({x}) :=$ $(\lambda \operatorname{bdef\_} \in \operatorname{dual\text{-}vec}(m).\; \operatorname{bdef\_}(x))$

\vnbbookfile{theorem-library/tvs-bidual}{theorem-library/tvs-bidual}
\noindent{\small bidual.scm -- topological-vector-space.pdf Prop 3.31: for a normed space E the canonical map iota\_E : E -$>$ E** is an isometry.  Agent DUAL-FIX, 2026-10-26. DESTINATION: theorem-library/tvs-bidual.scm.  Library: analysis.  Helper prefix: bdl-.}\par
\vnbprop{bidual-rr-le-chain}{\vnbstmt{\begin{array}{@{}l@{}} \forall a, b, c, d. \\ \quad ((a \in \mathbb{R}) \wedge ((b \in \mathbb{R}) \wedge ((c \in \mathbb{R}) \wedge ((d \in \mathbb{R}) \wedge ((a \leq b) \wedge ((b \leq c) \wedge (c \leq d))))))) \Rightarrow  \\ \quad \quad a \leq d \end{array}}}
\vnbprop{bidual-norm-sublinear-view}{\vnbstmt{\begin{array}{@{}l@{}} \forall m. \\ \quad \operatorname{is\text{-}normed\text{-}vector\text{-}space}(m) \Rightarrow  \\ \quad \quad \operatorname{is\text{-}sublinear}(\operatorname{normed\text{-}vector\text{-}space\text{-}as\text{-}module}(m), \operatorname{vnrm}(m)) \end{array}}}
\vnbprop{bidual-view-scal}{\vnbstmt{\begin{array}{@{}l@{}} \forall m. \\ \quad \operatorname{is\text{-}normed\text{-}vector\text{-}space}(m) \Rightarrow  \\ \quad \quad \operatorname{scal}(\operatorname{normed\text{-}vector\text{-}space\text{-}as\text{-}module}(m)) = \operatorname{normed\text{-}field\text{-}as\text{-}commutative\text{-}ring}(\operatorname{rr\text{-}normed\text{-}field}) \end{array}}}
\vnbprop{bidual-lf-on-to-view}{\vnbstmt{\begin{array}{@{}l@{}} \forall m, s, f. \\ \quad (\operatorname{is\text{-}normed\text{-}vector\text{-}space}(m) \wedge \operatorname{is\text{-}linear\text{-}functional\text{-}on}(m, s, f)) \Rightarrow  \\ \quad \quad \operatorname{is\text{-}linear\text{-}functional\text{-}on}(\operatorname{normed\text{-}vector\text{-}space\text{-}as\text{-}module}(m), s, f) \end{array}}}
\vnbprop{bidual-lf-from-view}{\vnbstmt{\begin{array}{@{}l@{}} \forall m, f. \\ \quad (\operatorname{is\text{-}normed\text{-}vector\text{-}space}(m) \wedge \operatorname{is\text{-}linear\text{-}functional\text{-}on}(\operatorname{normed\text{-}vector\text{-}space\text{-}as\text{-}module}(m), \operatorname{vec}(m), f)) \Rightarrow  \\ \quad \quad \operatorname{is\text{-}linear\text{-}functional}(m, f) \end{array}}}
\vnbprop{bidual-line-seed-dominated}{\vnbstmt{\begin{array}{@{}l@{}} \forall m, v, f.; \forall z \in \operatorname{line}(m, v). \\ \quad \operatorname{is\text{-}normed\text{-}vector\text{-}space}(m) \wedge \\ \quad v \in \operatorname{vec}(m) \wedge \\ \quad \operatorname{is\text{-}bounded\text{-}linear\text{-}functional\text{-}on}(m, \operatorname{line}(m, v), f) \wedge \\ \quad \operatorname{dual\text{-}norm\text{-}on}(m, \operatorname{line}(m, v), f) \leq 1 \Rightarrow \\ \quad \quad f(z) \leq \operatorname{vnrm}(m)(z) \end{array}}}
\vnbprop{bidual-two-sided-bound}{\vnbstmt{\begin{array}{@{}l@{}} \forall m, g.; \forall w \in \operatorname{vec}(m). \\ \quad \operatorname{is\text{-}normed\text{-}vector\text{-}space}(m) \wedge \\ \quad \operatorname{is\text{-}linear\text{-}functional}(m, g) \wedge \\ \quad \forall a \in \operatorname{vec}(m).\; (g(a) \leq \operatorname{vnrm}(m)(a)) \Rightarrow \\ \quad \quad \lvert g(w) \rvert \leq \operatorname{vnrm}(m)(w) \end{array}}}
\vnbprop{norming-functional-exists}{\vnbstmt{\begin{array}{@{}l@{}} \forall m, x. \\ \quad (\operatorname{is\text{-}normed\text{-}vector\text{-}space}(m) \wedge (x \in \operatorname{vec}(m))) \Rightarrow  \\ \quad \quad \exists g. \\ \quad \quad \quad (\operatorname{is\text{-}bounded\text{-}linear\text{-}functional}(m, g) \wedge ((\operatorname{dual\text{-}norm}(m, g) \leq 1) \wedge (g(x) = \operatorname{vnrm}(m)(x)))) \end{array}}}
\vnbprop{bidual-eval-apply}{\vnbstmt{\begin{array}{@{}l@{}} \forall m, x, f. \\ \quad (\operatorname{is\text{-}normed\text{-}vector\text{-}space}(m) \wedge ((x \in \operatorname{vec}(m)) \wedge (f \in \operatorname{dual\text{-}vec}(m)))) \Rightarrow  \\ \quad \quad \iota_{m}({x})(f) = f(x) \end{array}}}
\vnbprop{bidual-eval-fun}{\vnbstmt{\begin{array}{@{}l@{}} \forall m, x. \\ \quad (\operatorname{is\text{-}normed\text{-}vector\text{-}space}(m) \wedge (x \in \operatorname{vec}(m))) \Rightarrow  \\ \quad \quad \iota_{m}({x}) \in (\operatorname{dual\text{-}vec}(m) \to \mathbb{R}) \end{array}}}
\vnbprop{bidual-eval-linear}{\vnbstmt{\begin{array}{@{}l@{}} \forall m, x. \\ \quad (\operatorname{is\text{-}normed\text{-}vector\text{-}space}(m) \wedge (x \in \operatorname{vec}(m))) \Rightarrow  \\ \quad \quad \operatorname{is\text{-}linear\text{-}functional}({m}^{*}, \iota_{m}({x})) \end{array}}}
\vnbprop{bidual-eval-bound}{\vnbstmt{\begin{array}{@{}l@{}} \forall m, x.; \forall w \in \operatorname{vec}({m}^{*}). \\ \quad (\operatorname{is\text{-}normed\text{-}vector\text{-}space}(m) \wedge (x \in \operatorname{vec}(m))) \Rightarrow  \\ \quad \quad \lvert \iota_{m}({x})(w) \rvert \leq (\operatorname{vnrm}(m)(x) \cdot \operatorname{vnrm}({m}^{*})(w)) \end{array}}}
\vnbprop{bidual-eval-in-bidual}{\vnbstmt{\begin{array}{@{}l@{}} \forall m, x. \\ \quad (\operatorname{is\text{-}normed\text{-}vector\text{-}space}(m) \wedge (x \in \operatorname{vec}(m))) \Rightarrow  \\ \quad \quad \iota_{m}({x}) \in \operatorname{vec}({{m}^{*}}^{*}) \end{array}}}
\vnbprop{bidual-eval-isometry}{\vnbstmt{\begin{array}{@{}l@{}} \forall m, x. \\ \quad (\operatorname{is\text{-}normed\text{-}vector\text{-}space}(m) \wedge (x \in \operatorname{vec}(m))) \Rightarrow  \\ \quad \quad \operatorname{vnrm}({{m}^{*}}^{*})(\iota_{m}({x})) = \operatorname{vnrm}(m)(x) \end{array}}}
\vnbprop{bidual-eval-additive}{\vnbstmt{\begin{array}{@{}l@{}} \forall m, x, y. \\ \quad (\operatorname{is\text{-}normed\text{-}vector\text{-}space}(m) \wedge ((x \in \operatorname{vec}(m)) \wedge (y \in \operatorname{vec}(m)))) \Rightarrow  \\ \quad \quad \iota_{m}({\operatorname{vadd}(m)(x, y)}) = \operatorname{vadd}({{m}^{*}}^{*})(\iota_{m}({x}), \iota_{m}({y})) \end{array}}}
\vnbprop{bidual-eval-homogeneous}{\vnbstmt{\begin{array}{@{}l@{}} \forall m, r, x. \\ \quad (\operatorname{is\text{-}normed\text{-}vector\text{-}space}(m) \wedge ((r \in \mathbb{R}) \wedge (x \in \operatorname{vec}(m)))) \Rightarrow  \\ \quad \quad \iota_{m}({\operatorname{act}(m)(r, x)}) = \operatorname{act}({{m}^{*}}^{*})(r, \iota_{m}({x})) \end{array}}}

\vnbbookfile{theorem-library/norm-attained-general}{theorem-library/norm-attained-general}
\noindent{\small norm-attained-general.scm -- the norm-attaining functional WITHOUT a finite-dimension hypothesis, and the two theorems that cited the finite-dimensional one, restated without it as NEW theorems beside the old. Agent DUAL-FIX, 2026-10-26.  DESTINATION: theorem-library/norm-attained-general.scm.  Library: analysis. Load ...}\par
\vnbprop{norm-attained-by-functional-nvs}{\vnbstmt{\begin{array}{@{}l@{}} \forall m, x. \\ \quad (\operatorname{is\text{-}normed\text{-}vector\text{-}space}(m) \wedge (x \in \operatorname{vec}(m))) \Rightarrow  \\ \quad \quad \exists g. \\ \quad \quad \quad (\operatorname{is\text{-}bounded\text{-}linear\text{-}functional}(m, g) \wedge ((\operatorname{dual\text{-}norm}(m, g) \leq 1) \wedge (g(x) = \operatorname{vnrm}(m)(x)))) \end{array}}}
\vnbprop{norm-as-sup-nvs}{\vnbstmt{\begin{array}{@{}l@{}} \forall m, x. \\ \quad (\operatorname{is\text{-}normed\text{-}vector\text{-}space}(m) \wedge (x \in \operatorname{vec}(m))) \Rightarrow  \\ \quad \quad \forall f. \\ \quad \quad \quad (\operatorname{is\text{-}bounded\text{-}linear\text{-}functional}(m, f) \wedge (\operatorname{dual\text{-}norm}(m, f) \leq 1)) \Rightarrow  \\ \quad \quad \quad \quad \lvert f(x) \rvert \leq \operatorname{vnrm}(m)(x) \wedge \\ \quad \quad \forall d \in \mathbb{R}. \\ \quad \quad \quad \forall f.\; ((\operatorname{is\text{-}bounded\text{-}linear\text{-}functional}(m, f) \wedge (\operatorname{dual\text{-}norm}(m, f) \leq 1)) \Rightarrow (\lvert f(x) \rvert \leq d)) \Rightarrow  \\ \quad \quad \quad \quad \operatorname{vnrm}(m)(x) \leq d \end{array}}}
\vnbprop{vector-taylor-remainder-bound-nvs}{\vnbstmt{\begin{array}{@{}l@{}} \forall m, f, a, x, n. \\ \quad \operatorname{is\text{-}normed\text{-}vector\text{-}space}(m) \wedge \\ \quad f \in (\mathbb{R} \to \operatorname{vec}(m)) \wedge \\ \quad a \in \mathbb{R} \wedge \\ \quad x \in \mathbb{R} \wedge \\ \quad n \in \mathbb{N} \wedge \\ \quad a < x \wedge \\ \quad \operatorname{taylor\text{-}differentiable\text{-}v}(m, f, a, x, n) \Rightarrow \\ \quad \quad \exists b \in \mathbb{R}. \\ \quad \quad \quad a < b \wedge \\ \quad \quad \quad b < x \wedge \\ \quad \quad \quad \operatorname{factorial}((n + 1)) \cdot \operatorname{vnrm}(m)(\operatorname{vadd}(m)(f(x), \operatorname{vneg}(m)(\operatorname{taylor\text{-}poly\text{-}v}(m, f, a, x, n)))) \\ \quad \quad \quad \leq\; \operatorname{vnrm}(m)(\operatorname{nth\text{-}deriv\text{-}v}(m, f, (n + 1))(b)) \cdot \operatorname{power}((x - a), (n + 1)) \end{array}}}

\vnbbookfile{structure-library/rn-vocab}{structure-library/rn-vocab}
\noindent{\small rn-vocab.scm -- E = R\textasciicircum{}n WITH THE SUP NORM, calculus.pdf ch. 2 s. 8 (32), (36)-(37), (44). DEFINITIONS ONLY.  CALC-2, 2026-10-19: OPTION A of the representation question, chosen by the user 2026-10-19 (a point of R\textasciicircum{}n is a function on INTERVAL(1, n)).}\par
\par\medskip\noindent\textbf{Definition} (\texttt{rn}).\;$\operatorname{rn}(n) :=$
\begin{equation*}
{[}\operatorname{normed\text{-}field\text{-}as\text{-}commutative\text{-}ring}(\operatorname{rr\text{-}normed\text{-}field}), (\operatorname{interval}(1, n) \to \mathbb{R}), (\lambda {[}\operatorname{rnax\_}, \operatorname{rnay\_}] \in ((\operatorname{interval}(1, n) \to \mathbb{R}) \times (\operatorname{interval}(1, n) \to \mathbb{R})).\; (\lambda \operatorname{rnai\_} \in \operatorname{interval}(1, n).\; (\operatorname{rnax\_}(\operatorname{rnai\_}) + \operatorname{rnay\_}(\operatorname{rnai\_})))), (\lambda \operatorname{rnzi\_} \in \operatorname{interval}(1, n).\; 0), (\lambda \operatorname{rnnx\_} \in (\operatorname{interval}(1, n) \to \mathbb{R}).\; (\lambda \operatorname{rnni\_} \in \operatorname{interval}(1, n).\; -\operatorname{rnnx\_}(\operatorname{rnni\_}))), (\lambda {[}\operatorname{rnsr\_}, \operatorname{rnsx\_}] \in (\mathbb{R} \times (\operatorname{interval}(1, n) \to \mathbb{R})).\; (\lambda \operatorname{rnsi\_} \in \operatorname{interval}(1, n).\; (\operatorname{rnsr\_} \cdot \operatorname{rnsx\_}(\operatorname{rnsi\_})))), (\lambda \operatorname{rnvx\_} \in (\operatorname{interval}(1, n) \to \mathbb{R}).\; \operatorname{rn\text{-}maxabs}(\operatorname{rnvx\_}, n))]
\end{equation*}

\par\medskip\noindent\textbf{Definition} (\texttt{rn-maxabs}).\;For every $\operatorname{rnmx\_}$,  $(\operatorname{rn\text{-}maxabs}(\operatorname{rnmx\_}, 0) \simeq 0)$\par
Let $\operatorname{rnmx\_}$ be arbitrary. Let $\operatorname{rnmk\_} \in \mathbb{N}$. Then
\begin{equation*}
(\operatorname{rn\text{-}maxabs}(\operatorname{rnmx\_}, \operatorname{succ}(\operatorname{rnmk\_})) \simeq \operatorname{max}(\operatorname{rn\text{-}maxabs}(\operatorname{rnmx\_}, \operatorname{rnmk\_}), \lvert \operatorname{rnmx\_}(\operatorname{succ}(\operatorname{rnmk\_})) \rvert))
\end{equation*}

\par\medskip\noindent\textbf{Definition} (\texttt{rn-unit}).\;$\operatorname{rn\text{-}unit}(n, i) :=$ $(\lambda \operatorname{rnuj\_} \in \operatorname{interval}(1, n).\; \left(\text{if }(\operatorname{rnuj\_} = i)\text{ then }1\text{ else }0\right))$
\par\medskip\noindent\textbf{Definition} (\texttt{rn-trunc}).\;$\operatorname{rn\text{-}trunc}(n, x, k) :=$
\begin{equation*}
(\lambda \operatorname{rntj\_} \in \operatorname{interval}(1, n).\; \left(\text{if }(\operatorname{rntj\_} \leq k)\text{ then }x(\operatorname{rntj\_})\text{ else }0\right))
\end{equation*}

\par\medskip\noindent\textbf{Definition} (\texttt{has-dir-deriv}).\;For all $m$, $f$, $a$, $\eta$ and $l$, 
\begin{equation*}
\operatorname{has\text{-}dir\text{-}deriv}(m, f, a, \eta, l) \Leftrightarrow \operatorname{has\text{-}deriv\text{-}at}(\operatorname{seg\text{-}curve}(m, f, a, \eta), 0, l)
\end{equation*}


\vnbbookfile{theorem-library/calc2-rn-basics}{theorem-library/calc2-rn-basics}
\noindent{\small calc2-rn-basics.scm -- the read-offs of RN(n) (structure-library/rn-vocab.scm). CALC-2, 2026-10-19.  Helper prefix: c2b-.}\par
\vnbprop{rn-vec}{\vnbstmt{\forall n \in \mathbb{N}.\; (\operatorname{vec}(\operatorname{rn}(n)) \simeq (\operatorname{interval}(1, n) \to \mathbb{R}))}}
\vnbprop{rn-vadd-apply}{\vnbstmt{\begin{array}{@{}l@{}} \forall n, x, y, i. \\ \quad ((n \in \mathbb{N}) \wedge ((x \in (\operatorname{interval}(1, n) \to \mathbb{R})) \wedge ((y \in (\operatorname{interval}(1, n) \to \mathbb{R})) \wedge (i \in \operatorname{interval}(1, n))))) \Rightarrow  \\ \quad \quad \operatorname{vadd}(\operatorname{rn}(n))(x, y)(i) = (x(i) + y(i)) \end{array}}}
\vnbprop{rn-act-apply}{\vnbstmt{\begin{array}{@{}l@{}} \forall n, r, x, i. \\ \quad ((n \in \mathbb{N}) \wedge ((r \in \mathbb{R}) \wedge ((x \in (\operatorname{interval}(1, n) \to \mathbb{R})) \wedge (i \in \operatorname{interval}(1, n))))) \Rightarrow  \\ \quad \quad \operatorname{act}(\operatorname{rn}(n))(r, x)(i) = (r \cdot x(i)) \end{array}}}
\vnbprop{rn-vneg-apply}{\vnbstmt{\begin{array}{@{}l@{}} \forall n, x, i. \\ \quad ((n \in \mathbb{N}) \wedge ((x \in (\operatorname{interval}(1, n) \to \mathbb{R})) \wedge (i \in \operatorname{interval}(1, n)))) \Rightarrow  \\ \quad \quad \operatorname{vneg}(\operatorname{rn}(n))(x)(i) = -x(i) \end{array}}}
\vnbprop{rn-vzero-apply}{\vnbstmt{\forall n, i.\; (((n \in \mathbb{N}) \wedge (i \in \operatorname{interval}(1, n))) \Rightarrow (\operatorname{vzero}(\operatorname{rn}(n))(i) = 0))}}
\vnbprop{rn-vnrm-apply}{\vnbstmt{\begin{array}{@{}l@{}} \forall n, x. \\ \quad ((n \in \mathbb{N}) \wedge (x \in (\operatorname{interval}(1, n) \to \mathbb{R}))) \Rightarrow  \\ \quad \quad (\operatorname{vnrm}(\operatorname{rn}(n))(x) \simeq \operatorname{rn\text{-}maxabs}(x, n)) \end{array}}}
\vnbprop{rn-scal}{\vnbstmt{\begin{array}{@{}l@{}} \forall n \in \mathbb{N}. \\ \quad (\operatorname{scal}(\operatorname{rn}(n)) \simeq \operatorname{normed\text{-}field\text{-}as\text{-}commutative\text{-}ring}(\operatorname{rr\text{-}normed\text{-}field})) \end{array}}}
\vnbprop{rn-unit-at-same}{\vnbstmt{\forall n, i.\; (((n \in \mathbb{N}) \wedge (i \in \operatorname{interval}(1, n))) \Rightarrow (\operatorname{rn\text{-}unit}(n, i)(i) = 1))}}
\vnbprop{rn-unit-at-other}{\vnbstmt{\begin{array}{@{}l@{}} \forall n, i, j. \\ \quad ((n \in \mathbb{N}) \wedge ((i \in \operatorname{interval}(1, n)) \wedge ((j \in \operatorname{interval}(1, n)) \wedge \neg (j = i)))) \Rightarrow  \\ \quad \quad \operatorname{rn\text{-}unit}(n, i)(j) = 0 \end{array}}}
\vnbprop{rn-vadd-in-vec}{\vnbstmt{\begin{array}{@{}l@{}} \forall n, x, y. \\ \quad ((n \in \mathbb{N}) \wedge ((x \in (\operatorname{interval}(1, n) \to \mathbb{R})) \wedge (y \in (\operatorname{interval}(1, n) \to \mathbb{R})))) \Rightarrow  \\ \quad \quad \operatorname{vadd}(\operatorname{rn}(n))(x, y) \in (\operatorname{interval}(1, n) \to \mathbb{R}) \end{array}}}
\vnbprop{rn-act-in-vec}{\vnbstmt{\begin{array}{@{}l@{}} \forall n, r, x. \\ \quad ((n \in \mathbb{N}) \wedge ((r \in \mathbb{R}) \wedge (x \in (\operatorname{interval}(1, n) \to \mathbb{R})))) \Rightarrow  \\ \quad \quad \operatorname{act}(\operatorname{rn}(n))(r, x) \in (\operatorname{interval}(1, n) \to \mathbb{R}) \end{array}}}
\vnbprop{rn-vneg-in-vec}{\vnbstmt{\begin{array}{@{}l@{}} \forall n, x. \\ \quad ((n \in \mathbb{N}) \wedge (x \in (\operatorname{interval}(1, n) \to \mathbb{R}))) \Rightarrow  \\ \quad \quad \operatorname{vneg}(\operatorname{rn}(n))(x) \in (\operatorname{interval}(1, n) \to \mathbb{R}) \end{array}}}
\vnbprop{rn-vzero-in-vec}{\vnbstmt{\forall n \in \mathbb{N}.\; (\operatorname{vzero}(\operatorname{rn}(n)) \in (\operatorname{interval}(1, n) \to \mathbb{R}))}}
\vnbprop{rn-unit-in-vec}{\vnbstmt{\begin{array}{@{}l@{}} \forall n, i. \\ \quad ((n \in \mathbb{N}) \wedge (i \in \operatorname{interval}(1, n))) \Rightarrow (\operatorname{rn\text{-}unit}(n, i) \in (\operatorname{interval}(1, n) \to \mathbb{R})) \end{array}}}

\vnbbookfile{theorem-library/calc2-rn-norm}{theorem-library/calc2-rn-norm}
\noindent{\small calc2-rn-norm.scm -- the sup norm of RN(n): RN-MAXABS laws.  CALC-2, 2026-10-19.  Prefix c2n-.}\par
\vnbprop{rn-maxabs-typing}{\vnbstmt{\begin{array}{@{}l@{}} \forall k \in \mathbb{N}, n, v. \\ \quad ((n \in \mathbb{N}) \wedge ((v \in (\operatorname{interval}(1, n) \to \mathbb{R})) \wedge (k \leq n))) \Rightarrow  \\ \quad \quad (\operatorname{rn\text{-}maxabs}(v, k) \in \mathbb{R}) \wedge (0 \leq \operatorname{rn\text{-}maxabs}(v, k)) \end{array}}}
\vnbprop{rn-maxabs-ge}{\vnbstmt{\begin{array}{@{}l@{}} \forall k \in \mathbb{N}, n, v, i. \\ \quad ((n \in \mathbb{N}) \wedge ((v \in (\operatorname{interval}(1, n) \to \mathbb{R})) \wedge ((k \leq n) \wedge ((i \in \operatorname{interval}(1, n)) \wedge (i \leq k))))) \Rightarrow  \\ \quad \quad \lvert v(i) \rvert \leq \operatorname{rn\text{-}maxabs}(v, k) \end{array}}}
\vnbprop{rn-maxabs-le}{\vnbstmt{\begin{array}{@{}l@{}} \forall k \in \mathbb{N}, n, v, c. \\ \quad n \in \mathbb{N} \wedge \\ \quad v \in (\operatorname{interval}(1, n) \to \mathbb{R}) \wedge \\ \quad k \leq n \wedge \\ \quad c \in \mathbb{R} \wedge \\ \quad 0 \leq c \wedge \\ \quad \forall j \in \operatorname{interval}(1, n).\; (\lvert v(j) \rvert \leq c) \Rightarrow \\ \quad \quad \operatorname{rn\text{-}maxabs}(v, k) \leq c \end{array}}}
\vnbprop{rn-norm-typing}{\vnbstmt{\begin{array}{@{}l@{}} \forall n, v. \\ \quad ((n \in \mathbb{N}) \wedge (v \in (\operatorname{interval}(1, n) \to \mathbb{R}))) \Rightarrow  \\ \quad \quad (\operatorname{vnrm}(\operatorname{rn}(n))(v) \in \mathbb{R}) \wedge (0 \leq \operatorname{vnrm}(\operatorname{rn}(n))(v)) \end{array}}}
\vnbprop{rn-norm-ge-coord}{\vnbstmt{\begin{array}{@{}l@{}} \forall n, v, i. \\ \quad ((n \in \mathbb{N}) \wedge ((v \in (\operatorname{interval}(1, n) \to \mathbb{R})) \wedge (i \in \operatorname{interval}(1, n)))) \Rightarrow  \\ \quad \quad \lvert v(i) \rvert \leq \operatorname{vnrm}(\operatorname{rn}(n))(v) \end{array}}}
\vnbprop{rn-norm-le}{\vnbstmt{\begin{array}{@{}l@{}} \forall n, v, c. \\ \quad n \in \mathbb{N} \wedge \\ \quad v \in (\operatorname{interval}(1, n) \to \mathbb{R}) \wedge \\ \quad c \in \mathbb{R} \wedge \\ \quad 0 \leq c \wedge \\ \quad \forall j \in \operatorname{interval}(1, n).\; (\lvert v(j) \rvert \leq c) \Rightarrow \\ \quad \quad \operatorname{vnrm}(\operatorname{rn}(n))(v) \leq c \end{array}}}

\vnbbookfile{theorem-library/calc2-rn-nvs}{theorem-library/calc2-rn-nvs}
\noindent{\small calc2-rn-nvs.scm -- RN(n) is a real normed vector space.  CALC-2, 2026-10-19.  Prefix c2v-.}\par
\vnbprop{rn-vadd-assoc}{\vnbstmt{\begin{array}{@{}l@{}} \forall n, u, w, z. \\ \quad n \in \mathbb{N} \wedge \\ \quad u \in (\operatorname{interval}(1, n) \to \mathbb{R}) \wedge \\ \quad w \in (\operatorname{interval}(1, n) \to \mathbb{R}) \wedge \\ \quad z \in (\operatorname{interval}(1, n) \to \mathbb{R}) \Rightarrow \\ \quad \quad \operatorname{vadd}(\operatorname{rn}(n))(\operatorname{vadd}(\operatorname{rn}(n))(u, w), z) = \operatorname{vadd}(\operatorname{rn}(n))(u, \operatorname{vadd}(\operatorname{rn}(n))(w, z)) \end{array}}}
\vnbprop{rn-vadd-comm}{\vnbstmt{\begin{array}{@{}l@{}} \forall n, u, w. \\ \quad ((n \in \mathbb{N}) \wedge ((u \in (\operatorname{interval}(1, n) \to \mathbb{R})) \wedge (w \in (\operatorname{interval}(1, n) \to \mathbb{R})))) \Rightarrow  \\ \quad \quad \operatorname{vadd}(\operatorname{rn}(n))(u, w) = \operatorname{vadd}(\operatorname{rn}(n))(w, u) \end{array}}}
\vnbprop{rn-vzero-left}{\vnbstmt{\begin{array}{@{}l@{}} \forall n, u. \\ \quad ((n \in \mathbb{N}) \wedge (u \in (\operatorname{interval}(1, n) \to \mathbb{R}))) \Rightarrow  \\ \quad \quad \operatorname{vadd}(\operatorname{rn}(n))(\operatorname{vzero}(\operatorname{rn}(n)), u) = u \end{array}}}
\vnbprop{rn-vzero-right}{\vnbstmt{\begin{array}{@{}l@{}} \forall n, u. \\ \quad ((n \in \mathbb{N}) \wedge (u \in (\operatorname{interval}(1, n) \to \mathbb{R}))) \Rightarrow  \\ \quad \quad \operatorname{vadd}(\operatorname{rn}(n))(u, \operatorname{vzero}(\operatorname{rn}(n))) = u \end{array}}}
\vnbprop{rn-vneg-left}{\vnbstmt{\begin{array}{@{}l@{}} \forall n, u. \\ \quad ((n \in \mathbb{N}) \wedge (u \in (\operatorname{interval}(1, n) \to \mathbb{R}))) \Rightarrow  \\ \quad \quad \operatorname{vadd}(\operatorname{rn}(n))(\operatorname{vneg}(\operatorname{rn}(n))(u), u) = \operatorname{vzero}(\operatorname{rn}(n)) \end{array}}}
\vnbprop{rn-vneg-right}{\vnbstmt{\begin{array}{@{}l@{}} \forall n, u. \\ \quad ((n \in \mathbb{N}) \wedge (u \in (\operatorname{interval}(1, n) \to \mathbb{R}))) \Rightarrow  \\ \quad \quad \operatorname{vadd}(\operatorname{rn}(n))(u, \operatorname{vneg}(\operatorname{rn}(n))(u)) = \operatorname{vzero}(\operatorname{rn}(n)) \end{array}}}
\vnbprop{rn-act-vadd}{\vnbstmt{\begin{array}{@{}l@{}} \forall n, r, u, w. \\ \quad ((n \in \mathbb{N}) \wedge ((r \in \mathbb{R}) \wedge ((u \in (\operatorname{interval}(1, n) \to \mathbb{R})) \wedge (w \in (\operatorname{interval}(1, n) \to \mathbb{R}))))) \Rightarrow  \\ \quad \quad \operatorname{act}(\operatorname{rn}(n))(r, \operatorname{vadd}(\operatorname{rn}(n))(u, w)) = \operatorname{vadd}(\operatorname{rn}(n))(\operatorname{act}(\operatorname{rn}(n))(r, u), \operatorname{act}(\operatorname{rn}(n))(r, w)) \end{array}}}
\vnbprop{rn-act-add}{\vnbstmt{\begin{array}{@{}l@{}} \forall n, r, s, u. \\ \quad ((n \in \mathbb{N}) \wedge ((r \in \mathbb{R}) \wedge ((s \in \mathbb{R}) \wedge (u \in (\operatorname{interval}(1, n) \to \mathbb{R}))))) \Rightarrow  \\ \quad \quad \operatorname{act}(\operatorname{rn}(n))((r + s), u) = \operatorname{vadd}(\operatorname{rn}(n))(\operatorname{act}(\operatorname{rn}(n))(r, u), \operatorname{act}(\operatorname{rn}(n))(s, u)) \end{array}}}
\vnbprop{rn-act-mul}{\vnbstmt{\begin{array}{@{}l@{}} \forall n, r, s, u. \\ \quad ((n \in \mathbb{N}) \wedge ((r \in \mathbb{R}) \wedge ((s \in \mathbb{R}) \wedge (u \in (\operatorname{interval}(1, n) \to \mathbb{R}))))) \Rightarrow  \\ \quad \quad \operatorname{act}(\operatorname{rn}(n))((r \cdot s), u) = \operatorname{act}(\operatorname{rn}(n))(r, \operatorname{act}(\operatorname{rn}(n))(s, u)) \end{array}}}
\vnbprop{rn-act-one}{\vnbstmt{\begin{array}{@{}l@{}} \forall n, u. \\ \quad ((n \in \mathbb{N}) \wedge (u \in (\operatorname{interval}(1, n) \to \mathbb{R}))) \Rightarrow (\operatorname{act}(\operatorname{rn}(n))(1, u) = u) \end{array}}}
\vnbprop{rr-max-nonneg-scale}{\vnbstmt{\begin{array}{@{}l@{}} \forall a, b, c. \\ \quad ((a \in \mathbb{R}) \wedge ((b \in \mathbb{R}) \wedge ((c \in \mathbb{R}) \wedge (0 \leq c)))) \Rightarrow  \\ \quad \quad \operatorname{max}((c \cdot a), (c \cdot b)) = (c \cdot \operatorname{max}(a, b)) \end{array}}}
\vnbprop{rn-maxabs-scale}{\vnbstmt{\begin{array}{@{}l@{}} \forall k \in \mathbb{N}, n, v, r. \\ \quad ((n \in \mathbb{N}) \wedge ((v \in (\operatorname{interval}(1, n) \to \mathbb{R})) \wedge ((r \in \mathbb{R}) \wedge (k \leq n)))) \Rightarrow  \\ \quad \quad \operatorname{rn\text{-}maxabs}(\operatorname{act}(\operatorname{rn}(n))(r, v), k) = (\lvert r \rvert \cdot \operatorname{rn\text{-}maxabs}(v, k)) \end{array}}}
\vnbprop{rn-norm-homogeneous}{\vnbstmt{\begin{array}{@{}l@{}} \forall n, r, u. \\ \quad ((n \in \mathbb{N}) \wedge ((r \in \mathbb{R}) \wedge (u \in (\operatorname{interval}(1, n) \to \mathbb{R})))) \Rightarrow  \\ \quad \quad \operatorname{vnrm}(\operatorname{rn}(n))(\operatorname{act}(\operatorname{rn}(n))(r, u)) = (\lvert r \rvert \cdot \operatorname{vnrm}(\operatorname{rn}(n))(u)) \end{array}}}
\vnbprop{rn-norm-triangle}{\vnbstmt{\begin{array}{@{}l@{}} \forall n, u, w. \\ \quad ((n \in \mathbb{N}) \wedge ((u \in (\operatorname{interval}(1, n) \to \mathbb{R})) \wedge (w \in (\operatorname{interval}(1, n) \to \mathbb{R})))) \Rightarrow  \\ \quad \quad \operatorname{vnrm}(\operatorname{rn}(n))(\operatorname{vadd}(\operatorname{rn}(n))(u, w)) \leq (\operatorname{vnrm}(\operatorname{rn}(n))(u) + \operatorname{vnrm}(\operatorname{rn}(n))(w)) \end{array}}}
\vnbprop{rn-norm-definite}{\vnbstmt{\begin{array}{@{}l@{}} \forall n, u. \\ \quad ((n \in \mathbb{N}) \wedge (u \in (\operatorname{interval}(1, n) \to \mathbb{R}))) \Rightarrow  \\ \quad \quad (\operatorname{vnrm}(\operatorname{rn}(n))(u) = 0) \Leftrightarrow (u = \operatorname{vzero}(\operatorname{rn}(n))) \end{array}}}
\vnbprop{rn-is-normed-vector-space}{\vnbstmt{\forall n \in \mathbb{N}.\; \operatorname{is\text{-}normed\text{-}vector\text{-}space}(\operatorname{rn}(n))}}
\vnbprop{rn-trunc-in-vec}{\vnbstmt{\begin{array}{@{}l@{}} \forall n, x, k. \\ \quad ((n \in \mathbb{N}) \wedge ((x \in (\operatorname{interval}(1, n) \to \mathbb{R})) \wedge (k \in \mathbb{N}))) \Rightarrow  \\ \quad \quad \operatorname{rn\text{-}trunc}(n, x, k) \in (\operatorname{interval}(1, n) \to \mathbb{R}) \end{array}}}
\vnbprop{rn-trunc-apply-le}{\vnbstmt{\begin{array}{@{}l@{}} \forall n, x, k, j. \\ \quad ((n \in \mathbb{N}) \wedge ((x \in (\operatorname{interval}(1, n) \to \mathbb{R})) \wedge ((k \in \mathbb{N}) \wedge ((j \in \operatorname{interval}(1, n)) \wedge (j \leq k))))) \Rightarrow  \\ \quad \quad \operatorname{rn\text{-}trunc}(n, x, k)(j) = x(j) \end{array}}}
\vnbprop{rn-trunc-apply-gt}{\vnbstmt{\begin{array}{@{}l@{}} \forall n, x, k, j. \\ \quad ((n \in \mathbb{N}) \wedge ((x \in (\operatorname{interval}(1, n) \to \mathbb{R})) \wedge ((k \in \mathbb{N}) \wedge ((j \in \operatorname{interval}(1, n)) \wedge \neg (j \leq k))))) \Rightarrow  \\ \quad \quad \operatorname{rn\text{-}trunc}(n, x, k)(j) = 0 \end{array}}}
\vnbprop{rn-trunc-zero}{\vnbstmt{\begin{array}{@{}l@{}} \forall n, x. \\ \quad ((n \in \mathbb{N}) \wedge (x \in (\operatorname{interval}(1, n) \to \mathbb{R}))) \Rightarrow  \\ \quad \quad \operatorname{rn\text{-}trunc}(n, x, 0) = \operatorname{vzero}(\operatorname{rn}(n)) \end{array}}}
\vnbprop{rn-trunc-full}{\vnbstmt{\begin{array}{@{}l@{}} \forall n, x. \\ \quad ((n \in \mathbb{N}) \wedge (x \in (\operatorname{interval}(1, n) \to \mathbb{R}))) \Rightarrow (\operatorname{rn\text{-}trunc}(n, x, n) = x) \end{array}}}
\vnbprop{rn-trunc-succ}{\vnbstmt{\begin{array}{@{}l@{}} \forall n, x, k. \\ \quad ((n \in \mathbb{N}) \wedge ((x \in (\operatorname{interval}(1, n) \to \mathbb{R})) \wedge ((k \in \mathbb{N}) \wedge ((k + 1) \leq n)))) \Rightarrow  \\ \quad \quad \operatorname{rn\text{-}trunc}(n, x, (k + 1)) \\ \quad \quad =\; \operatorname{vadd}(\operatorname{rn}(n))(\operatorname{rn\text{-}trunc}(n, x, k), \operatorname{act}(\operatorname{rn}(n))(x((k + 1)), \operatorname{rn\text{-}unit}(n, (k + 1)))) \end{array}}}
\vnbprop{rn-trunc-step-norm}{\vnbstmt{\begin{array}{@{}l@{}} \forall n, x, k, t. \\ \quad n \in \mathbb{N} \wedge \\ \quad x \in (\operatorname{interval}(1, n) \to \mathbb{R}) \wedge \\ \quad k \in \mathbb{N} \wedge \\ \quad (k + 1) \leq n \wedge \\ \quad t \in \mathbb{R} \wedge \\ \quad \lvert t \rvert \leq \lvert x((k + 1)) \rvert \Rightarrow \\ \quad \quad \operatorname{vnrm}(\operatorname{rn}(n))(\operatorname{vadd}(\operatorname{rn}(n))(\operatorname{rn\text{-}trunc}(n, x, k), \operatorname{act}(\operatorname{rn}(n))(t, \operatorname{rn\text{-}unit}(n, (k + 1))))) \\ \quad \quad \leq\; \operatorname{vnrm}(\operatorname{rn}(n))(x) \end{array}}}

\vnbbookfile{theorem-library/calc2-segment}{theorem-library/calc2-segment}
\noindent{\small calc2-segment.scm -- the one-variable facts the coordinate segments need.  CALC-2, 2026-10-19. Prefix c2s-.}\par
\vnbprop{continuous-on-of-derivs-on-ccint}{\vnbstmt{\begin{array}{@{}l@{}} \forall o, i, f, h. \\ \quad o \in \mathbb{R} \wedge \\ \quad i \in \mathbb{R} \wedge \\ \quad h \in (\operatorname{ccint}(o, i) \to \mathbb{R}) \wedge \\ \quad \forall y \in \operatorname{ccint}(o, i).\; (h(y) = f(y)) \wedge \\ \quad \forall x \in \operatorname{ccint}(o, i).\; \exists l.\; \operatorname{has\text{-}deriv\text{-}at}(f, x, l) \Rightarrow \\ \quad \quad \operatorname{is\text{-}continuous\text{-}on}(h, \operatorname{ccint}(o, i)) \end{array}}}
\vnbprop{mean-value-inequality-on-interval}{\vnbstmt{\begin{array}{@{}l@{}} \forall o, i, f, l, d. \\ \quad o \in \mathbb{R} \wedge \\ \quad i \in \mathbb{R} \wedge \\ \quad o < i \wedge \\ \quad l \in \mathbb{R} \wedge \\ \quad d \in \mathbb{R} \wedge \\ \quad \forall x \in \operatorname{ccint}(o, i). \exists a. \\ \quad \quad \operatorname{has\text{-}deriv\text{-}at}(f, x, a) \wedge (\lvert (a - l) \rvert \leq d) \Rightarrow \\ \quad \quad \lvert ((f(i) - f(o)) - (l \cdot (i - o))) \rvert \leq (d \cdot (i - o)) \end{array}}}
\vnbprop{has-deriv-at-of-eps-bound}{\vnbstmt{\begin{array}{@{}l@{}} \forall g, c, r, l. \\ \quad c \in \mathbb{R} \wedge \\ \quad \operatorname{pos\text{-}rr}(r) \wedge \\ \quad l \in \mathbb{R} \wedge \\ \quad \forall t \in \operatorname{ooint}((c - r), (c + r)).\; (g(t) \in \mathbb{R}) \wedge \\ \quad \forall e. \\ \quad \quad \operatorname{pos\text{-}rr}(e) \Rightarrow  \\ \quad \quad \quad \exists d. \\ \quad \quad \quad \quad \operatorname{pos\text{-}rr}(d) \wedge \\ \quad \quad \quad \quad \forall t \in \operatorname{ooint}((c - r), (c + r)). \\ \quad \quad \quad \quad \quad (\lvert (t - c) \rvert < d) \Rightarrow (\lvert ((g(t) - g(c)) - ((t - c) \cdot l)) \rvert \leq (e \cdot \lvert (t - c) \rvert)) \Rightarrow \\ \quad \quad \operatorname{has\text{-}deriv\text{-}at}(g, c, l) \end{array}}}

\vnbbookfile{theorem-library/calc2-dirderiv}{theorem-library/calc2-dirderiv}
\noindent{\small calc2-dirderiv.scm -- the local directional derivative HAS-DIR-DERIV over a normed vector space: moving the base point along the direction.  CALC-2, 2026-10-19.  Prefix c2d-.}\par
\vnbprop{has-dir-deriv-shift}{\vnbstmt{\begin{array}{@{}l@{}} \forall m, f, b, e, t, l. \\ \quad \operatorname{is\text{-}normed\text{-}vector\text{-}space}(m) \wedge \\ \quad b \in \operatorname{vec}(m) \wedge \\ \quad e \in \operatorname{vec}(m) \wedge \\ \quad t \in \mathbb{R} \wedge \\ \quad \operatorname{has\text{-}dir\text{-}deriv}(m, f, \operatorname{vadd}(m)(b, \operatorname{act}(m)(t, e)), e, l) \Rightarrow \\ \quad \quad \operatorname{has\text{-}deriv\text{-}at}(\operatorname{seg\text{-}curve}(m, f, b, e), t, l) \end{array}}}
\vnbprop{has-dir-deriv-step}{\vnbstmt{\begin{array}{@{}l@{}} \forall m, f, b, e, \mathit{dlam}, l, c. \\ \quad \operatorname{is\text{-}normed\text{-}vector\text{-}space}(m) \wedge \\ \quad b \in \operatorname{vec}(m) \wedge \\ \quad e \in \operatorname{vec}(m) \wedge \\ \quad \mathit{dlam} \in \mathbb{R} \wedge \\ \quad l \in \mathbb{R} \wedge \\ \quad c \in \mathbb{R} \wedge \\ \quad \forall t \in \mathbb{R}. \\ \quad \quad (\lvert t \rvert \leq \lvert \mathit{dlam} \rvert) \Rightarrow  \\ \quad \quad \quad \exists d. \\ \quad \quad \quad \quad \operatorname{has\text{-}dir\text{-}deriv}(m, f, \operatorname{vadd}(m)(b, \operatorname{act}(m)(t, e)), e, d) \wedge (\lvert (d - l) \rvert \leq c) \Rightarrow \\ \quad \quad \lvert ((f(\operatorname{vadd}(m)(b, \operatorname{act}(m)(\mathit{dlam}, e))) - f(b)) - (\mathit{dlam} \cdot l)) \rvert \leq (c \cdot \lvert \mathit{dlam} \rvert) \end{array}}}

\vnbbookfile{theorem-library/calc2-prop229}{theorem-library/calc2-prop229}
\noindent{\small calc2-prop229.scm -- calculus.pdf Prop. 2.29: continuous partial derivatives give the first-order estimate on E = R\textasciicircum{}n with the sup norm.  CALC-2, 2026-10-19.  Prefix c29-.}\par
\vnbprop{rr-small-factor-real}{\vnbstmt{\begin{array}{@{}l@{}} \forall e, n. \\ \quad (\operatorname{pos\text{-}rr}(e) \wedge ((n \in \mathbb{R}) \wedge (0 \leq n))) \Rightarrow \exists c.\; (\operatorname{pos\text{-}rr}(c) \wedge ((c \cdot n) \leq e)) \end{array}}}
\vnbprop{rn-partials-estimate-trunc}{\vnbstmt{\begin{array}{@{}l@{}} \forall k \in \mathbb{N}, n, u, f, p, a, c, r. \\ \quad n \in \mathbb{N} \wedge \\ \quad \operatorname{is\text{-}open}(\operatorname{nvs\text{-}metric\text{-}space}(\operatorname{rn}(n)), u) \wedge \\ \quad f \in (u \to \mathbb{R}) \wedge \\ \quad a \in u \wedge \\ \quad \forall i \in \operatorname{interval}(1, n), x \in u. \\ \quad \quad \operatorname{has\text{-}dir\text{-}deriv}(\operatorname{rn}(n), f, x, \operatorname{rn\text{-}unit}(n, i), p(i)(x)) \wedge \\ \quad \forall i \in \operatorname{interval}(1, n), x \in u. \\ \quad \quad \operatorname{is\text{-}continuous\text{-}at}(\operatorname{subspace\text{-}ms}(\operatorname{nvs\text{-}metric\text{-}space}(\operatorname{rn}(n)), u), \operatorname{rr\text{-}ms}, p(i), x) \wedge \\ \quad \operatorname{pos\text{-}rr}(c) \wedge \\ \quad \operatorname{pos\text{-}rr}(r) \wedge \\ \quad \operatorname{ball}(\operatorname{nvs\text{-}metric\text{-}space}(\operatorname{rn}(n)), a, r) \subseteq u \wedge \\ \quad k \leq n \Rightarrow \\ \quad \quad \exists d. \\ \quad \quad \quad \operatorname{pos\text{-}rr}(d) \wedge \\ \quad \quad \quad d < r \wedge \\ \quad \quad \quad \forall e \in (\operatorname{interval}(1, n) \to \mathbb{R}). \\ \quad \quad \quad \quad (\operatorname{vnrm}(\operatorname{rn}(n))(e) \leq d) \Rightarrow  \\ \quad \quad \quad \quad \quad \lvert ((f(\operatorname{vadd}(\operatorname{rn}(n))(a, \operatorname{rn\text{-}trunc}(n, e, k))) - f(a)) - \operatorname{finsum}(\operatorname{commutative\text{-}ring\text{-}additive\text{-}ag}(\operatorname{rr\text{-}normed\text{-}field}), (\lambda \mathit{cfi} \in \operatorname{interval}(1, n).\; (e(\mathit{cfi}) \cdot p(\mathit{cfi})(a))), \operatorname{interval}(1, k))) \rvert \\ \quad \quad \quad \quad \quad \leq\; k \cdot (c \cdot \operatorname{vnrm}(\operatorname{rn}(n))(e)) \end{array}}}
\vnbprop{continuous-partials-estimate}{\vnbstmt{\begin{array}{@{}l@{}} \forall n, u, f, p, a, s. \\ \quad n \in \mathbb{N} \wedge \\ \quad \operatorname{is\text{-}open}(\operatorname{nvs\text{-}metric\text{-}space}(\operatorname{rn}(n)), u) \wedge \\ \quad f \in (u \to \mathbb{R}) \wedge \\ \quad a \in u \wedge \\ \quad \forall i \in \operatorname{interval}(1, n), x \in u. \\ \quad \quad \operatorname{has\text{-}dir\text{-}deriv}(\operatorname{rn}(n), f, x, \operatorname{rn\text{-}unit}(n, i), p(i)(x)) \wedge \\ \quad \forall i \in \operatorname{interval}(1, n), x \in u. \\ \quad \quad \operatorname{is\text{-}continuous\text{-}at}(\operatorname{subspace\text{-}ms}(\operatorname{nvs\text{-}metric\text{-}space}(\operatorname{rn}(n)), u), \operatorname{rr\text{-}ms}, p(i), x) \wedge \\ \quad \operatorname{pos\text{-}rr}(s) \Rightarrow \\ \quad \quad \exists b. \\ \quad \quad \quad \operatorname{pos\text{-}rr}(b) \wedge \\ \quad \quad \quad \forall c \in \operatorname{vec}(\operatorname{rn}(n)). \\ \quad \quad \quad \quad (\operatorname{vnrm}(\operatorname{rn}(n))(c) \leq b) \Rightarrow  \\ \quad \quad \quad \quad \quad \lvert ((f(\operatorname{vadd}(\operatorname{rn}(n))(a, c)) - f(a)) - \operatorname{finsum}(\operatorname{commutative\text{-}ring\text{-}additive\text{-}ag}(\operatorname{rr\text{-}normed\text{-}field}), (\lambda i \in \operatorname{interval}(1, n).\; (c(i) \cdot p(i)(a))), \operatorname{interval}(1, n))) \rvert \\ \quad \quad \quad \quad \quad \leq\; \operatorname{vnrm}(\operatorname{rn}(n))(c) \cdot s \end{array}}}

\vnbbookfile{theorem-library/calc2-prop228}{theorem-library/calc2-prop228}
\noindent{\small calc2-prop228.scm -- calculus.pdf Prop. 2.28 on E = R\textasciicircum{}n with the sup norm: partial derivatives defined and continuous on the open set U give the directional derivative in every direction eta at every x in U, d\_eta f(x) = sum\_i eta\_i d\_(e\_i) f(x), linear in eta and continuous in x.  Derived from Prop. 2.29 ...}\par
\vnbprop{rn-partials-sum-scale}{\vnbstmt{\begin{array}{@{}l@{}} \forall n, p, x, e, t. \\ \quad n \in \mathbb{N} \wedge \\ \quad e \in (\operatorname{interval}(1, n) \to \mathbb{R}) \wedge \\ \quad t \in \mathbb{R} \wedge \\ \quad \forall i \in \operatorname{interval}(1, n).\; (p(i)(x) \in \mathbb{R}) \Rightarrow \\ \quad \quad \operatorname{finsum}(\operatorname{commutative\text{-}ring\text{-}additive\text{-}ag}(\operatorname{rr\text{-}normed\text{-}field}), \\ \quad \quad \quad (\lambda i \in \operatorname{interval}(1, n).\; (\operatorname{act}(\operatorname{rn}(n))(t, e)(i) \cdot p(i)(x))), \\ \quad \quad \quad \operatorname{interval}(1, n)) \\ \quad \quad =\; t \cdot \operatorname{finsum}(\operatorname{commutative\text{-}ring\text{-}additive\text{-}ag}(\operatorname{rr\text{-}normed\text{-}field}), (\lambda i \in \operatorname{interval}(1, n).\; (e(i) \cdot p(i)(x))), \operatorname{interval}(1, n)) \end{array}}}
\vnbprop{continuous-partials-dir-deriv}{\vnbstmt{\begin{array}{@{}l@{}} \forall n, u, f, p.; \forall x \in u, a \in \operatorname{vec}(\operatorname{rn}(n)). \\ \quad n \in \mathbb{N} \wedge \\ \quad \operatorname{is\text{-}open}(\operatorname{nvs\text{-}metric\text{-}space}(\operatorname{rn}(n)), u) \wedge \\ \quad f \in (u \to \mathbb{R}) \wedge \\ \quad \forall i \in \operatorname{interval}(1, n), x \in u. \\ \quad \quad \operatorname{has\text{-}dir\text{-}deriv}(\operatorname{rn}(n), f, x, \operatorname{rn\text{-}unit}(n, i), p(i)(x)) \wedge \\ \quad \forall i \in \operatorname{interval}(1, n), x \in u. \\ \quad \quad \operatorname{is\text{-}continuous\text{-}at}(\operatorname{subspace\text{-}ms}(\operatorname{nvs\text{-}metric\text{-}space}(\operatorname{rn}(n)), u), \operatorname{rr\text{-}ms}, p(i), x) \Rightarrow \\ \quad \quad \operatorname{has\text{-}dir\text{-}deriv}(\operatorname{rn}(n), f, x, a, \\ \quad \quad \quad \operatorname{finsum}(\operatorname{commutative\text{-}ring\text{-}additive\text{-}ag}(\operatorname{rr\text{-}normed\text{-}field}), \\ \quad \quad \quad \quad (\lambda i \in \operatorname{interval}(1, n).\; (a(i) \cdot p(i)(x))), \operatorname{interval}(1, n))) \end{array}}}
\vnbprop{rn-partials-sum-linear}{\vnbstmt{\begin{array}{@{}l@{}} \forall n, p, x, e, z, r, s. \\ \quad n \in \mathbb{N} \wedge \\ \quad e \in (\operatorname{interval}(1, n) \to \mathbb{R}) \wedge \\ \quad z \in (\operatorname{interval}(1, n) \to \mathbb{R}) \wedge \\ \quad r \in \mathbb{R} \wedge \\ \quad s \in \mathbb{R} \wedge \\ \quad \forall i \in \operatorname{interval}(1, n).\; (p(i)(x) \in \mathbb{R}) \Rightarrow \\ \quad \quad \operatorname{finsum}(\operatorname{commutative\text{-}ring\text{-}additive\text{-}ag}(\operatorname{rr\text{-}normed\text{-}field}), \\ \quad \quad \quad (\lambda i \in \operatorname{interval}(1, n).\; \\ \quad \quad \quad \quad \operatorname{vadd}(\operatorname{rn}(n))(\operatorname{act}(\operatorname{rn}(n))(r, e), \operatorname{act}(\operatorname{rn}(n))(s, z))(i) \cdot p(i)(x)), \\ \quad \quad \quad \operatorname{interval}(1, n)) \\ \quad \quad =\; (r \cdot \operatorname{finsum}(\operatorname{commutative\text{-}ring\text{-}additive\text{-}ag}(\operatorname{rr\text{-}normed\text{-}field}), (\lambda i \in \operatorname{interval}(1, n).\; (e(i) \cdot p(i)(x))), \operatorname{interval}(1, n))) + (s \cdot \operatorname{finsum}(\operatorname{commutative\text{-}ring\text{-}additive\text{-}ag}(\operatorname{rr\text{-}normed\text{-}field}), (\lambda i \in \operatorname{interval}(1, n).\; (z(i) \cdot p(i)(x))), \operatorname{interval}(1, n))) \end{array}}}
\vnbprop{ms-rr-add-continuous-at}{\vnbstmt{\begin{array}{@{}l@{}} \forall s, f, g, a. \\ \quad \operatorname{is\text{-}metric\text{-}space}(s) \wedge \\ \quad f \in (\operatorname{pts}(s) \to \mathbb{R}) \wedge \\ \quad g \in (\operatorname{pts}(s) \to \mathbb{R}) \wedge \\ \quad a \in \operatorname{pts}(s) \wedge \\ \quad \operatorname{is\text{-}continuous\text{-}at}(s, \operatorname{rr\text{-}ms}, f, a) \wedge \\ \quad \operatorname{is\text{-}continuous\text{-}at}(s, \operatorname{rr\text{-}ms}, g, a) \Rightarrow \\ \quad \quad \operatorname{is\text{-}continuous\text{-}at}(s, \operatorname{rr\text{-}ms}, (\lambda x \in \operatorname{pts}(s).\; (f(x) + g(x))), a) \end{array}}}
\vnbprop{ms-rr-scale-continuous-at}{\vnbstmt{\begin{array}{@{}l@{}} \forall s, c, f, a. \\ \quad \operatorname{is\text{-}metric\text{-}space}(s) \wedge \\ \quad c \in \mathbb{R} \wedge \\ \quad f \in (\operatorname{pts}(s) \to \mathbb{R}) \wedge \\ \quad a \in \operatorname{pts}(s) \wedge \\ \quad \operatorname{is\text{-}continuous\text{-}at}(s, \operatorname{rr\text{-}ms}, f, a) \Rightarrow \\ \quad \quad \operatorname{is\text{-}continuous\text{-}at}(s, \operatorname{rr\text{-}ms}, (\lambda x \in \operatorname{pts}(s).\; (c \cdot f(x))), a) \end{array}}}
\vnbprop{finite-sum-continuous}{\vnbstmt{\begin{array}{@{}l@{}} \forall k \in \mathbb{N}, s, n, h, t, a. \\ \quad \operatorname{is\text{-}metric\text{-}space}(s) \wedge \\ \quad n \in \mathbb{N} \wedge \\ \quad k \leq n \wedge \\ \quad t \in (\operatorname{interval}(1, n) \to \mathbb{R}) \wedge \\ \quad a \in \operatorname{pts}(s) \wedge \\ \quad \forall i \in \operatorname{interval}(1, n).\; (h(i) \in (\operatorname{pts}(s) \to \mathbb{R})) \wedge \\ \quad \forall i \in \operatorname{interval}(1, n).\; \operatorname{is\text{-}continuous\text{-}at}(s, \operatorname{rr\text{-}ms}, h(i), a) \Rightarrow \\ \quad \quad \operatorname{is\text{-}continuous\text{-}at}(s, \operatorname{rr\text{-}ms}, \\ \quad \quad \quad (\lambda y \in \operatorname{pts}(s).\; \\ \quad \quad \quad \quad \operatorname{finsum}(\operatorname{commutative\text{-}ring\text{-}additive\text{-}ag}(\operatorname{rr\text{-}normed\text{-}field}), \\ \quad \quad \quad \quad \quad (\lambda i \in \operatorname{interval}(1, n).\; (t(i) \cdot h(i)(y))), \\ \quad \quad \quad \quad \quad \operatorname{interval}(1, k))), \\ \quad \quad \quad a) \end{array}}}
\vnbprop{continuous-partials-dir-deriv-continuous}{\vnbstmt{\begin{array}{@{}l@{}} \forall n, u, f, p.; \forall a \in \operatorname{vec}(\operatorname{rn}(n)), x \in u. \\ \quad n \in \mathbb{N} \wedge \\ \quad \operatorname{is\text{-}open}(\operatorname{nvs\text{-}metric\text{-}space}(\operatorname{rn}(n)), u) \wedge \\ \quad f \in (u \to \mathbb{R}) \wedge \\ \quad \forall i \in \operatorname{interval}(1, n), x \in u. \\ \quad \quad \operatorname{has\text{-}dir\text{-}deriv}(\operatorname{rn}(n), f, x, \operatorname{rn\text{-}unit}(n, i), p(i)(x)) \wedge \\ \quad \forall i \in \operatorname{interval}(1, n), x \in u. \\ \quad \quad \operatorname{is\text{-}continuous\text{-}at}(\operatorname{subspace\text{-}ms}(\operatorname{nvs\text{-}metric\text{-}space}(\operatorname{rn}(n)), u), \operatorname{rr\text{-}ms}, p(i), x) \Rightarrow \\ \quad \quad \operatorname{is\text{-}continuous\text{-}at}(\operatorname{subspace\text{-}ms}(\operatorname{nvs\text{-}metric\text{-}space}(\operatorname{rn}(n)), u), \\ \quad \quad \quad \operatorname{rr\text{-}ms}, \\ \quad \quad \quad (\lambda y \in u.\; \\ \quad \quad \quad \quad \operatorname{finsum}(\operatorname{commutative\text{-}ring\text{-}additive\text{-}ag}(\operatorname{rr\text{-}normed\text{-}field}), \\ \quad \quad \quad \quad \quad (\lambda i \in \operatorname{interval}(1, n).\; (a(i) \cdot p(i)(y))), \\ \quad \quad \quad \quad \quad \operatorname{interval}(1, n))), \\ \quad \quad \quad x) \end{array}}}

\vnbbookfile{structure-library/rb2-vocab}{structure-library/rb2-vocab}
\noindent{\small rb2-vocab.scm -- vocabulary for calculus.pdf 2.27 and for path-connectedness. DEFINITIONS ONLY (destination: structure-library/).  Agent RA-BASICS, 2026-10-26 (NEW DEFINITIONS, the user's to confirm; "missing vocabulary is not a stopping point").}\par
\par\medskip\noindent\textbf{Definition} (\texttt{ldir-deriv}).\;$\operatorname{ldir\text{-}deriv}(\operatorname{ldm\_}, \operatorname{ldf\_}, \operatorname{lda\_}, \operatorname{lde\_}) :=$ $(\iota\, \operatorname{ldl\_}.\;\operatorname{has\text{-}dir\text{-}deriv}(\operatorname{ldm\_}, \operatorname{ldf\_}, \operatorname{lda\_}, \operatorname{lde\_}, \operatorname{ldl\_}))$
\par\medskip\noindent\textbf{Definition} (\texttt{span2}).\;$\operatorname{span2}(\operatorname{s2m\_}, \operatorname{s2e\_}, \operatorname{s2x\_}) :=$
\begin{equation*}
\{\operatorname{s2v\_} \in \operatorname{vec}(\operatorname{s2m\_}) : \exists \operatorname{s2r\_}, \operatorname{s2s\_} \in \mathbb{R}.\; (\operatorname{s2v\_} = \operatorname{vadd}(\operatorname{s2m\_})(\operatorname{act}(\operatorname{s2m\_})(\operatorname{s2r\_}, \operatorname{s2e\_}), \operatorname{act}(\operatorname{s2m\_})(\operatorname{s2s\_}, \operatorname{s2x\_})))\}
\end{equation*}

\par\medskip\noindent\textbf{Definition} (\texttt{is-path-connected}).\;For all $\operatorname{pcs\_}$ and $\operatorname{pca\_}$, 
\begin{equation*}
\begin{array}{@{}l@{}} \operatorname{is\text{-}path\text{-}connected}(\operatorname{pcs\_}, \operatorname{pca\_}) \\ \Leftrightarrow\; \operatorname{is\text{-}metric\text{-}space}(\operatorname{pcs\_}) \wedge \\ \operatorname{pca\_} \subseteq \operatorname{pts}(\operatorname{pcs\_}) \wedge \\ \forall \operatorname{pcx\_}, \operatorname{pcy\_} \in \operatorname{pca\_}. \exists \operatorname{pcg\_} \in (\operatorname{ccint}(0, 1) \to \operatorname{pca\_}). \\ \quad (\operatorname{is\text{-}continuous}(\operatorname{subspace\text{-}ms}(\operatorname{rr\text{-}ms}, \operatorname{ccint}(0, 1)), \operatorname{pcs\_}, \operatorname{pcg\_}) \wedge ((\operatorname{pcg\_}(0) = \operatorname{pcx\_}) \wedge (\operatorname{pcg\_}(1) = \operatorname{pcy\_}))) \end{array}
\end{equation*}

\par\medskip\noindent\textbf{Definition} (\texttt{is-convex-rr}).\;For every $\operatorname{cvra\_}$, 
\begin{equation*}
\begin{array}{@{}l@{}} \operatorname{is\text{-}convex\text{-}rr}(\operatorname{cvra\_}) \\ \Leftrightarrow\; \operatorname{cvra\_} \subseteq \mathbb{R} \wedge \\ \forall \operatorname{cvrx\_}, \operatorname{cvry\_} \in \operatorname{cvra\_}, \operatorname{cvrt\_} \in \operatorname{ccint}(0, 1). \\ \quad +(\operatorname{cvrx\_}, \cdot(\operatorname{cvrt\_}, -(\operatorname{cvry\_}, \operatorname{cvrx\_}))) \in \operatorname{cvra\_} \end{array}
\end{equation*}


\vnbbookfile{theorem-library/rb2-dir}{theorem-library/rb2-dir}
\noindent{\small rb2-dir.scm -- calculus.pdf Lemmas 2.25, 2.26 and Prop. 2.27 for the LOCAL directional derivative HAS-DIR-DERIV over a normed vector space, f defined on an open set U.  Helper prefix r2d-.}\par
\vnbprop{ldir-deriv-value}{\vnbstmt{\begin{array}{@{}l@{}} \forall m, f, a, b, l. \\ \quad \operatorname{has\text{-}dir\text{-}deriv}(m, f, a, b, l) \Rightarrow (\operatorname{ldir\text{-}deriv}(m, f, a, b) = l) \end{array}}}
\vnbprop{has-dir-deriv-scale}{\vnbstmt{\begin{array}{@{}l@{}} \forall m, f, a, b, k, l. \\ \quad \operatorname{is\text{-}normed\text{-}vector\text{-}space}(m) \wedge \\ \quad a \in \operatorname{vec}(m) \wedge \\ \quad b \in \operatorname{vec}(m) \wedge \\ \quad k \in \mathbb{R} \wedge \\ \quad \operatorname{has\text{-}dir\text{-}deriv}(m, f, a, b, l) \Rightarrow \\ \quad \quad \operatorname{has\text{-}dir\text{-}deriv}(m, f, a, \operatorname{act}(m)(k, b), (k \cdot l)) \end{array}}}
\vnbprop{has-dir-deriv-add}{\vnbstmt{\begin{array}{@{}l@{}} \forall m, u, f, b, i, p, q, a. \\ \quad \operatorname{is\text{-}normed\text{-}vector\text{-}space}(m) \wedge \\ \quad \operatorname{is\text{-}open}(\operatorname{nvs\text{-}metric\text{-}space}(m), u) \wedge \\ \quad f \in (u \to \mathbb{R}) \wedge \\ \quad b \in \operatorname{vec}(m) \wedge \\ \quad i \in \operatorname{vec}(m) \wedge \\ \quad \forall x \in u.\; \operatorname{has\text{-}dir\text{-}deriv}(m, f, x, b, p(x)) \wedge \\ \quad \forall x \in u.\; \operatorname{has\text{-}dir\text{-}deriv}(m, f, x, i, q(x)) \wedge \\ \quad a \in u \wedge \\ \quad \operatorname{is\text{-}continuous\text{-}at}(\operatorname{subspace\text{-}ms}(\operatorname{nvs\text{-}metric\text{-}space}(m), u), \operatorname{rr\text{-}ms}, p, a) \wedge \\ \quad \operatorname{is\text{-}continuous\text{-}at}(\operatorname{subspace\text{-}ms}(\operatorname{nvs\text{-}metric\text{-}space}(m), u), \operatorname{rr\text{-}ms}, q, a) \Rightarrow \\ \quad \quad \operatorname{has\text{-}dir\text{-}deriv}(m, f, a, \operatorname{vadd}(m)(b, i), (p(a) + q(a))) \end{array}}}
\vnbprop{dir-deriv-span-formula}{\vnbstmt{\begin{array}{@{}l@{}} \forall m, u, f, a, i, p, q.; \forall x \in u, r, s \in \mathbb{R}. \\ \quad \operatorname{is\text{-}normed\text{-}vector\text{-}space}(m) \wedge \\ \quad \operatorname{is\text{-}open}(\operatorname{nvs\text{-}metric\text{-}space}(m), u) \wedge \\ \quad f \in (u \to \mathbb{R}) \wedge \\ \quad a \in \operatorname{vec}(m) \wedge \\ \quad i \in \operatorname{vec}(m) \wedge \\ \quad \forall x \in u.\; \operatorname{has\text{-}dir\text{-}deriv}(m, f, x, a, p(x)) \wedge \\ \quad \forall x \in u.\; \operatorname{has\text{-}dir\text{-}deriv}(m, f, x, i, q(x)) \wedge \\ \quad \forall x \in u. \\ \quad \quad \operatorname{is\text{-}continuous\text{-}at}(\operatorname{subspace\text{-}ms}(\operatorname{nvs\text{-}metric\text{-}space}(m), u), \operatorname{rr\text{-}ms}, p, x) \wedge \\ \quad \forall x \in u. \\ \quad \quad \operatorname{is\text{-}continuous\text{-}at}(\operatorname{subspace\text{-}ms}(\operatorname{nvs\text{-}metric\text{-}space}(m), u), \operatorname{rr\text{-}ms}, q, x) \Rightarrow \\ \quad \quad \operatorname{has\text{-}dir\text{-}deriv}(m, f, x, \operatorname{vadd}(m)(\operatorname{act}(m)(r, a), \operatorname{act}(m)(s, i)), ((r \cdot p(x)) + (s \cdot q(x)))) \end{array}}}
\vnbprop{span2-unfold}{\vnbstmt{\begin{array}{@{}l@{}} \forall m, a, i. \\ \quad (\operatorname{span2}(m, a, i) \simeq \{v \in \operatorname{vec}(m) : \exists r, s \in \mathbb{R}.\; (v = \operatorname{vadd}(m)(\operatorname{act}(m)(r, a), \operatorname{act}(m)(s, i)))\}) \end{array}}}
\vnbprop{span2-add-coeffs}{\vnbstmt{\begin{array}{@{}l@{}} \forall m, a, i, b, c, d, f. \\ \quad \operatorname{is\text{-}normed\text{-}vector\text{-}space}(m) \wedge \\ \quad a \in \operatorname{vec}(m) \wedge \\ \quad i \in \operatorname{vec}(m) \wedge \\ \quad b \in \mathbb{R} \wedge \\ \quad c \in \mathbb{R} \wedge \\ \quad d \in \mathbb{R} \wedge \\ \quad f \in \mathbb{R} \Rightarrow \\ \quad \quad \operatorname{vadd}(m)(\operatorname{vadd}(m)(\operatorname{act}(m)(b, a), \operatorname{act}(m)(c, i)), \operatorname{vadd}(m)(\operatorname{act}(m)(d, a), \operatorname{act}(m)(f, i))) \\ \quad \quad =\; \operatorname{vadd}(m)(\operatorname{act}(m)((b + d), a), \operatorname{act}(m)((c + f), i)) \end{array}}}
\vnbprop{span2-scale-coeffs}{\vnbstmt{\begin{array}{@{}l@{}} \forall m, a, i, c, r, s. \\ \quad (\operatorname{is\text{-}normed\text{-}vector\text{-}space}(m) \wedge ((a \in \operatorname{vec}(m)) \wedge ((i \in \operatorname{vec}(m)) \wedge ((c \in \mathbb{R}) \wedge ((r \in \mathbb{R}) \wedge (s \in \mathbb{R})))))) \Rightarrow  \\ \quad \quad \operatorname{act}(m)(c, \operatorname{vadd}(m)(\operatorname{act}(m)(r, a), \operatorname{act}(m)(s, i))) \\ \quad \quad =\; \operatorname{vadd}(m)(\operatorname{act}(m)((c \cdot r), a), \operatorname{act}(m)((c \cdot s), i)) \end{array}}}
\vnbprop{dir-deriv-span-linear}{\vnbstmt{\begin{array}{@{}l@{}} \forall m, u, f, a, i, p, q.; \forall x \in u. \\ \quad \operatorname{is\text{-}normed\text{-}vector\text{-}space}(m) \wedge \\ \quad \operatorname{is\text{-}open}(\operatorname{nvs\text{-}metric\text{-}space}(m), u) \wedge \\ \quad f \in (u \to \mathbb{R}) \wedge \\ \quad a \in \operatorname{vec}(m) \wedge \\ \quad i \in \operatorname{vec}(m) \wedge \\ \quad \forall x \in u.\; \operatorname{has\text{-}dir\text{-}deriv}(m, f, x, a, p(x)) \wedge \\ \quad \forall x \in u.\; \operatorname{has\text{-}dir\text{-}deriv}(m, f, x, i, q(x)) \wedge \\ \quad \forall x \in u. \\ \quad \quad \operatorname{is\text{-}continuous\text{-}at}(\operatorname{subspace\text{-}ms}(\operatorname{nvs\text{-}metric\text{-}space}(m), u), \operatorname{rr\text{-}ms}, p, x) \wedge \\ \quad \forall x \in u. \\ \quad \quad \operatorname{is\text{-}continuous\text{-}at}(\operatorname{subspace\text{-}ms}(\operatorname{nvs\text{-}metric\text{-}space}(m), u), \operatorname{rr\text{-}ms}, q, x) \Rightarrow \\ \quad \quad \forall b \in \operatorname{span2}(m, a, i).\; \operatorname{has\text{-}dir\text{-}deriv}(m, f, x, b, \operatorname{ldir\text{-}deriv}(m, f, x, b)) \wedge \\ \quad \quad \operatorname{is\text{-}linear\text{-}functional\text{-}on}(m, \operatorname{span2}(m, a, i), (\lambda b \in \operatorname{span2}(m, a, i).\; \operatorname{ldir\text{-}deriv}(m, f, x, b))) \end{array}}}

\vnbbookfile{structure-library/hilbert-space}{structure-library/hilbert-space}
\noindent{\small hilbert-space.scm -- HILBERT SPACES: the metric space of a complex inner product space, completeness, and the orthogonal complement.  Vocabulary only: no proof, no support.  Agent HILBERT-1, 2026-10-10 (scratchpad/triage/HILBERT-1.md).  The proofs are}\par
\par\medskip\noindent\textbf{Definition} (\texttt{perp}).\;${\mathit{ks}}^{\perp} :=$
\begin{equation*}
\{\operatorname{prp\_} \in \operatorname{vec}(v) : \forall \operatorname{prk\_} \in \mathit{ks}.\; (\langle\cdot,\cdot\rangle_{v}(\operatorname{prp\_}, \operatorname{prk\_}) = 0)\}
\end{equation*}

\par\medskip\noindent\textbf{Definition} (\texttt{ip-metric-space}).\;$\operatorname{ip\text{-}metric\text{-}space}(v) :=$ $\operatorname{nag\text{-}metric\text{-}space}(\operatorname{ip\text{-}normed\text{-}ag}(v))$
\par\medskip\noindent\textbf{Definition} (\texttt{is-hilbert-space}).\;For every $v$, 
\begin{equation*}
\operatorname{is\text{-}hilbert\text{-}space}(v) \Leftrightarrow (\operatorname{is\text{-}complex\text{-}inner\text{-}product\text{-}space}(v) \wedge \operatorname{is\text{-}complete}(\operatorname{ip\text{-}metric\text{-}space}(v)))
\end{equation*}

\par\medskip\noindent\textbf{Definition} (\texttt{is-orthogonal-projection}).\;For all $v$, $\mathit{ks}$ and $\mathit{pk}$, 
\begin{equation*}
\begin{array}{@{}l@{}} \operatorname{is\text{-}orthogonal\text{-}projection}(v, \mathit{ks}, \mathit{pk}) \\ \Leftrightarrow\; \mathit{pk} \in (\operatorname{vec}(v) \to \operatorname{vec}(v)) \wedge \\ \forall \operatorname{opz\_} \in \operatorname{vec}(v). \\ \quad (\operatorname{pk}(\operatorname{opz\_}) \in \mathit{ks}) \wedge (\operatorname{vadd}(v)(\operatorname{opz\_}, \operatorname{vneg}(v)(\operatorname{pk}(\operatorname{opz\_}))) \in {\mathit{ks}}^{\perp}) \end{array}
\end{equation*}

\par\medskip\noindent\textbf{Definition} (\texttt{is-cc-linear-functional}).\;For all $v$ and $f$, 
\begin{equation*}
\begin{array}{@{}l@{}} \operatorname{is\text{-}cc\text{-}linear\text{-}functional}(v, f) \\ \Leftrightarrow\; f \in (\operatorname{vec}(v) \to \mathbb{C}) \wedge \\ \forall \operatorname{hlx\_}, \operatorname{hly\_} \in \operatorname{vec}(v). \\ \quad f(\operatorname{vadd}(v)(\operatorname{hlx\_}, \operatorname{hly\_})) = +(f(\operatorname{hlx\_}), f(\operatorname{hly\_})) \wedge \\ \forall \operatorname{hlc\_} \in \mathbb{C}, \operatorname{hlx\_} \in \operatorname{vec}(v).\; (f(\operatorname{act}(v)(\operatorname{hlc\_}, \operatorname{hlx\_})) = \cdot(\operatorname{hlc\_}, f(\operatorname{hlx\_}))) \end{array}
\end{equation*}


\vnbbookfile{theorem-library/hilbert-projection}{theorem-library/hilbert-projection}
\noindent{\small hilbert-projection.scm -- HILBERT SPACES: the parallelogram law, K-perp, the orthogonal projection theorem (Prop. 2.19), the Riesz representation theorem (Thm. 2.22) and the density criterion (Cor. 2.21).  All PROVEN modulo 0. Agent HILBERT-1, 2026-10-10 (scratchpad/triage/HILBERT-1.md).  Vocabulary:}\par
\vnbprop{ip-metric-space-unfold}{\vnbstmt{\forall v.\; (\operatorname{ip\text{-}metric\text{-}space}(v) \simeq \operatorname{nag\text{-}metric\text{-}space}(\operatorname{ip\text{-}normed\text{-}ag}(v)))}}
\vnbprop{perp-unfold}{\vnbstmt{\forall v, s.\; ({s}^{\perp} \simeq \{p \in \operatorname{vec}(v) : \forall k \in s.\; (\langle\cdot,\cdot\rangle_{v}(p, k) = 0)\})}}
\vnbprop{perp-membership}{\vnbstmt{\begin{array}{@{}l@{}} \forall v, s, z. \\ \quad (z \in {s}^{\perp}) \Leftrightarrow ((z \in \operatorname{vec}(v)) \wedge \forall k \in s.\; (\langle\cdot,\cdot\rangle_{v}(z, k) = 0)) \end{array}}}
\vnbprop{hilbert-space-is-cips}{\vnbstmt{\forall v.\; (\operatorname{is\text{-}hilbert\text{-}space}(v) \Rightarrow \operatorname{is\text{-}complex\text{-}inner\text{-}product\text{-}space}(v))}}
\vnbprop{hilbert-space-complete}{\vnbstmt{\forall v.\; (\operatorname{is\text{-}hilbert\text{-}space}(v) \Rightarrow \operatorname{is\text{-}complete}(\operatorname{ip\text{-}metric\text{-}space}(v)))}}
\vnbprop{ip-metric-space-is-metric-space}{\vnbstmt{\begin{array}{@{}l@{}} \forall v. \\ \quad \operatorname{is\text{-}complex\text{-}inner\text{-}product\text{-}space}(v) \Rightarrow \operatorname{is\text{-}metric\text{-}space}(\operatorname{ip\text{-}metric\text{-}space}(v)) \end{array}}}
\vnbprop{ip-metric-space-pts}{\vnbstmt{\begin{array}{@{}l@{}} \forall v. \\ \quad \operatorname{is\text{-}complex\text{-}inner\text{-}product\text{-}space}(v) \Rightarrow (\operatorname{pts}(\operatorname{ip\text{-}metric\text{-}space}(v)) \simeq \operatorname{vec}(v)) \end{array}}}
\vnbprop{ip-normed-ag-opr}{\vnbstmt{\begin{array}{@{}l@{}} \forall v. \\ \quad \operatorname{is\text{-}complex\text{-}inner\text{-}product\text{-}space}(v) \Rightarrow (\operatorname{opr}(\operatorname{ip\text{-}normed\text{-}ag}(v)) \simeq \operatorname{vadd}(v)) \end{array}}}
\vnbprop{ip-normed-ag-inv}{\vnbstmt{\begin{array}{@{}l@{}} \forall v. \\ \quad \operatorname{is\text{-}complex\text{-}inner\text{-}product\text{-}space}(v) \Rightarrow (\operatorname{inv}(\operatorname{ip\text{-}normed\text{-}ag}(v)) \simeq \operatorname{vneg}(v)) \end{array}}}
\vnbprop{ip-metric-space-dist}{\vnbstmt{\begin{array}{@{}l@{}} \forall v.; \forall x, y \in \operatorname{vec}(v). \\ \quad \operatorname{is\text{-}complex\text{-}inner\text{-}product\text{-}space}(v) \Rightarrow  \\ \quad \quad \operatorname{dist}(\operatorname{ip\text{-}metric\text{-}space}(v))(x, y) \\ \quad \quad =\; \sqrt{\langle\cdot,\cdot\rangle_{v}(\operatorname{vadd}(v)(x, \operatorname{vneg}(v)(y)), \operatorname{vadd}(v)(x, \operatorname{vneg}(v)(y)))} \end{array}}}
\vnbprop{cips-ip-neg-right}{\vnbstmt{\begin{array}{@{}l@{}} \forall v.; \forall x, y \in \operatorname{vec}(v). \\ \quad \operatorname{is\text{-}complex\text{-}inner\text{-}product\text{-}space}(v) \Rightarrow  \\ \quad \quad \langle\cdot,\cdot\rangle_{v}(x, \operatorname{vneg}(v)(y)) = -\langle\cdot,\cdot\rangle_{v}(x, y) \end{array}}}
\vnbprop{cips-parallelogram-ip}{\vnbstmt{\begin{array}{@{}l@{}} \forall v.; \forall x, y \in \operatorname{vec}(v). \\ \quad \operatorname{is\text{-}complex\text{-}inner\text{-}product\text{-}space}(v) \Rightarrow  \\ \quad \quad \langle\cdot,\cdot\rangle_{v}(\operatorname{vadd}(v)(x, y), \operatorname{vadd}(v)(x, y)) + \langle\cdot,\cdot\rangle_{v}(\operatorname{vadd}(v)(x, \operatorname{vneg}(v)(y)), \operatorname{vadd}(v)(x, \operatorname{vneg}(v)(y))) \\ \quad \quad =\; (2 \cdot \langle\cdot,\cdot\rangle_{v}(x, x)) + (2 \cdot \langle\cdot,\cdot\rangle_{v}(y, y)) \end{array}}}
\vnbprop{cips-parallelogram}{\vnbstmt{\begin{array}{@{}l@{}} \forall v, x, y. \\ \quad (\operatorname{is\text{-}complex\text{-}inner\text{-}product\text{-}space}(v) \wedge ((x \in \operatorname{vec}(v)) \wedge (y \in \operatorname{vec}(v)))) \Rightarrow  \\ \quad \quad (\lVert \operatorname{vadd}(v)(x, y) \rVert \cdot \lVert \operatorname{vadd}(v)(x, y) \rVert) + (\lVert \operatorname{vadd}(v)(x, \operatorname{vneg}(v)(y)) \rVert \cdot \lVert \operatorname{vadd}(v)(x, \operatorname{vneg}(v)(y)) \rVert) \\ \quad \quad =\; (2 \cdot (\lVert x \rVert \cdot \lVert x \rVert)) + (2 \cdot (\lVert y \rVert \cdot \lVert y \rVert)) \end{array}}}
\vnbprop{cips-pythagoras}{\vnbstmt{\begin{array}{@{}l@{}} \forall v, x, y. \\ \quad (\operatorname{is\text{-}complex\text{-}inner\text{-}product\text{-}space}(v) \wedge ((x \in \operatorname{vec}(v)) \wedge ((y \in \operatorname{vec}(v)) \wedge (\langle\cdot,\cdot\rangle_{v}(x, y) = 0)))) \Rightarrow  \\ \quad \quad \lVert \operatorname{vadd}(v)(x, y) \rVert \cdot \lVert \operatorname{vadd}(v)(x, y) \rVert \\ \quad \quad =\; (\lVert x \rVert \cdot \lVert x \rVert) + (\lVert y \rVert \cdot \lVert y \rVert) \end{array}}}
\vnbprop{cc-zero-of-self-conj-zero}{\vnbstmt{\forall a \in \mathbb{C}.\; (((a \cdot \operatorname{conjugate}(a)) = 0) \Rightarrow (a = 0))}}
\vnbprop{cc-magnitude-sqrt-self-conj}{\vnbstmt{\forall z \in \mathbb{C}.\; (\operatorname{magnitude}(z) = \sqrt{(z \cdot \operatorname{conjugate}(z))})}}
\vnbprop{ip-orth-closure}{\vnbstmt{\begin{array}{@{}l@{}} \forall v, a, w, x. \\ \quad \operatorname{is\text{-}complex\text{-}inner\text{-}product\text{-}space}(v) \wedge \\ \quad w \in \operatorname{vec}(v) \wedge \\ \quad x \in \operatorname{closure}(\operatorname{ip\text{-}metric\text{-}space}(v), a) \wedge \\ \quad \forall y \in a.\; (\langle\cdot,\cdot\rangle_{v}(y, w) = 0) \Rightarrow \\ \quad \quad \langle\cdot,\cdot\rangle_{v}(x, w) = 0 \end{array}}}
\vnbprop{perp-orth}{\vnbstmt{\forall v, s, z, k.\; (((z \in {s}^{\perp}) \wedge (k \in s)) \Rightarrow (\langle\cdot,\cdot\rangle_{v}(z, k) = 0))}}
\vnbprop{perp-subset}{\vnbstmt{\begin{array}{@{}l@{}} \forall v, s. \\ \quad \operatorname{is\text{-}complex\text{-}inner\text{-}product\text{-}space}(v) \Rightarrow ({s}^{\perp} \subseteq \operatorname{vec}(v)) \end{array}}}
\vnbprop{perp-is-submodule}{\vnbstmt{\begin{array}{@{}l@{}} \forall v, s. \\ \quad (\operatorname{is\text{-}complex\text{-}inner\text{-}product\text{-}space}(v) \wedge (s \subseteq \operatorname{vec}(v))) \Rightarrow  \\ \quad \quad \operatorname{is\text{-}submodule}(v, {s}^{\perp}) \end{array}}}
\vnbprop{perp-is-closed}{\vnbstmt{\begin{array}{@{}l@{}} \forall v, s. \\ \quad (\operatorname{is\text{-}complex\text{-}inner\text{-}product\text{-}space}(v) \wedge (s \subseteq \operatorname{vec}(v))) \Rightarrow  \\ \quad \quad \operatorname{is\text{-}closed}(\operatorname{ip\text{-}metric\text{-}space}(v), {s}^{\perp}) \end{array}}}
\vnbprop{cips-dist-sq-inf-nonempty}{\vnbstmt{\begin{array}{@{}l@{}} \forall v, k, i. \\ \quad (\operatorname{is\text{-}complex\text{-}inner\text{-}product\text{-}space}(v) \wedge ((k \subseteq \operatorname{vec}(v)) \wedge (\exists e.\; (e \in k) \wedge (i \in \operatorname{vec}(v))))) \Rightarrow  \\ \quad \quad \exists a \in \mathbb{R}. \\ \quad \quad \quad 0 \leq a \wedge \\ \quad \quad \quad \forall p \in k. \\ \quad \quad \quad \quad a \leq \langle\cdot,\cdot\rangle_{v}(\operatorname{vadd}(v)(p, \operatorname{vneg}(v)(i)), \operatorname{vadd}(v)(p, \operatorname{vneg}(v)(i))) \wedge \\ \quad \quad \quad \forall n \in \mathbb{N}. \exists q \in k. \\ \quad \quad \quad \quad \langle\cdot,\cdot\rangle_{v}(\operatorname{vadd}(v)(q, \operatorname{vneg}(v)(i)), \operatorname{vadd}(v)(q, \operatorname{vneg}(v)(i))) \leq (a + \operatorname{recip}((n + 1))) \end{array}}}
\vnbprop{cips-dist-sq-inf}{\vnbstmt{\begin{array}{@{}l@{}} \forall v, k, i. \\ \quad (\operatorname{is\text{-}complex\text{-}inner\text{-}product\text{-}space}(v) \wedge (\operatorname{is\text{-}submodule}(v, k) \wedge (i \in \operatorname{vec}(v)))) \Rightarrow  \\ \quad \quad \exists a \in \mathbb{R}. \\ \quad \quad \quad 0 \leq a \wedge \\ \quad \quad \quad \forall p \in k. \\ \quad \quad \quad \quad a \leq \langle\cdot,\cdot\rangle_{v}(\operatorname{vadd}(v)(p, \operatorname{vneg}(v)(i)), \operatorname{vadd}(v)(p, \operatorname{vneg}(v)(i))) \wedge \\ \quad \quad \quad \forall n \in \mathbb{N}. \exists q \in k. \\ \quad \quad \quad \quad \langle\cdot,\cdot\rangle_{v}(\operatorname{vadd}(v)(q, \operatorname{vneg}(v)(i)), \operatorname{vadd}(v)(q, \operatorname{vneg}(v)(i))) \leq (a + \operatorname{recip}((n + 1))) \end{array}}}
\vnbprop{cips-midpoint-estimate-midpoint}{\vnbstmt{\begin{array}{@{}l@{}} \forall v, k, i, a, p_{1}, p_{2}. \\ \quad \operatorname{is\text{-}complex\text{-}inner\text{-}product\text{-}space}(v) \wedge \\ \quad k \subseteq \operatorname{vec}(v) \wedge \\ \quad \forall m_{1}, m_{2} \in k.\; (\operatorname{act}(v)(1/2, \operatorname{vadd}(v)(m_{1}, m_{2})) \in k) \wedge \\ \quad i \in \operatorname{vec}(v) \wedge \\ \quad a \in \mathbb{R} \wedge \\ \quad \forall p \in k. \\ \quad \quad a \leq \langle\cdot,\cdot\rangle_{v}(\operatorname{vadd}(v)(p, \operatorname{vneg}(v)(i)), \operatorname{vadd}(v)(p, \operatorname{vneg}(v)(i))) \wedge \\ \quad p_{1} \in k \wedge \\ \quad p_{2} \in k \Rightarrow \\ \quad \quad \langle\cdot,\cdot\rangle_{v}(\operatorname{vadd}(v)(p_{1}, \operatorname{vneg}(v)(p_{2})), \operatorname{vadd}(v)(p_{1}, \operatorname{vneg}(v)(p_{2}))) \\ \quad \quad \leq\; (((2 \cdot \langle\cdot,\cdot\rangle_{v}(\operatorname{vadd}(v)(p_{1}, \operatorname{vneg}(v)(i)), \operatorname{vadd}(v)(p_{1}, \operatorname{vneg}(v)(i)))) + (2 \cdot \langle\cdot,\cdot\rangle_{v}(\operatorname{vadd}(v)(p_{2}, \operatorname{vneg}(v)(i)), \operatorname{vadd}(v)(p_{2}, \operatorname{vneg}(v)(i))))) - (4 \cdot a)) \end{array}}}
\vnbprop{cips-midpoint-estimate}{\vnbstmt{\begin{array}{@{}l@{}} \forall v, k, i, a, p_{1}, p_{2}. \\ \quad \operatorname{is\text{-}complex\text{-}inner\text{-}product\text{-}space}(v) \wedge \\ \quad \operatorname{is\text{-}submodule}(v, k) \wedge \\ \quad i \in \operatorname{vec}(v) \wedge \\ \quad a \in \mathbb{R} \wedge \\ \quad \forall p \in k. \\ \quad \quad a \leq \langle\cdot,\cdot\rangle_{v}(\operatorname{vadd}(v)(p, \operatorname{vneg}(v)(i)), \operatorname{vadd}(v)(p, \operatorname{vneg}(v)(i))) \wedge \\ \quad p_{1} \in k \wedge \\ \quad p_{2} \in k \Rightarrow \\ \quad \quad \langle\cdot,\cdot\rangle_{v}(\operatorname{vadd}(v)(p_{1}, \operatorname{vneg}(v)(p_{2})), \operatorname{vadd}(v)(p_{1}, \operatorname{vneg}(v)(p_{2}))) \\ \quad \quad \leq\; (((2 \cdot \langle\cdot,\cdot\rangle_{v}(\operatorname{vadd}(v)(p_{1}, \operatorname{vneg}(v)(i)), \operatorname{vadd}(v)(p_{1}, \operatorname{vneg}(v)(i)))) + (2 \cdot \langle\cdot,\cdot\rangle_{v}(\operatorname{vadd}(v)(p_{2}, \operatorname{vneg}(v)(i)), \operatorname{vadd}(v)(p_{2}, \operatorname{vneg}(v)(i))))) - (4 \cdot a)) \end{array}}}
\vnbprop{cips-eq-of-dist-sq-zero}{\vnbstmt{\begin{array}{@{}l@{}} \forall v, x, y. \\ \quad \operatorname{is\text{-}complex\text{-}inner\text{-}product\text{-}space}(v) \wedge \\ \quad x \in \operatorname{vec}(v) \wedge \\ \quad y \in \operatorname{vec}(v) \wedge \\ \quad \langle\cdot,\cdot\rangle_{v}(\operatorname{vadd}(v)(x, \operatorname{vneg}(v)(y)), \operatorname{vadd}(v)(x, \operatorname{vneg}(v)(y))) = 0 \Rightarrow \\ \quad \quad x = y \end{array}}}
\vnbprop{hilbert-closest-point-midpoint}{\vnbstmt{\begin{array}{@{}l@{}} \forall v, k, i. \\ \quad \operatorname{is\text{-}hilbert\text{-}space}(v) \wedge \\ \quad k \subseteq \operatorname{vec}(v) \wedge \\ \quad \exists e.\; (e \in k) \wedge \\ \quad \forall m_{1}, m_{2} \in k.\; (\operatorname{act}(v)(1/2, \operatorname{vadd}(v)(m_{1}, m_{2})) \in k) \wedge \\ \quad \operatorname{is\text{-}closed}(\operatorname{ip\text{-}metric\text{-}space}(v), k) \wedge \\ \quad i \in \operatorname{vec}(v) \Rightarrow \\ \quad \quad \exists p \in k. \forall \mathit{hdp} \in k. \\ \quad \quad \quad \langle\cdot,\cdot\rangle_{v}(\operatorname{vadd}(v)(p, \operatorname{vneg}(v)(i)), \operatorname{vadd}(v)(p, \operatorname{vneg}(v)(i))) \\ \quad \quad \quad \leq\; \langle\cdot,\cdot\rangle_{v}(\operatorname{vadd}(v)(\mathit{hdp}, \operatorname{vneg}(v)(i)), \operatorname{vadd}(v)(\mathit{hdp}, \operatorname{vneg}(v)(i))) \end{array}}}
\vnbprop{hilbert-closest-point}{\vnbstmt{\begin{array}{@{}l@{}} \forall v, k, i. \\ \quad (\operatorname{is\text{-}hilbert\text{-}space}(v) \wedge (\operatorname{is\text{-}submodule}(v, k) \wedge (\operatorname{is\text{-}closed}(\operatorname{ip\text{-}metric\text{-}space}(v), k) \wedge (i \in \operatorname{vec}(v))))) \Rightarrow  \\ \quad \quad \exists p \in k. \forall \mathit{hdp} \in k. \\ \quad \quad \quad \langle\cdot,\cdot\rangle_{v}(\operatorname{vadd}(v)(p, \operatorname{vneg}(v)(i)), \operatorname{vadd}(v)(p, \operatorname{vneg}(v)(i))) \\ \quad \quad \quad \leq\; \langle\cdot,\cdot\rangle_{v}(\operatorname{vadd}(v)(\mathit{hdp}, \operatorname{vneg}(v)(i)), \operatorname{vadd}(v)(\mathit{hdp}, \operatorname{vneg}(v)(i))) \end{array}}}
\vnbprop{cips-variational}{\vnbstmt{\begin{array}{@{}l@{}} \forall v, w, e. \\ \quad \operatorname{is\text{-}complex\text{-}inner\text{-}product\text{-}space}(v) \wedge \\ \quad w \in \operatorname{vec}(v) \wedge \\ \quad e \in \operatorname{vec}(v) \wedge \\ \quad \forall t \in \mathbb{C}. \\ \quad \quad \langle\cdot,\cdot\rangle_{v}(w, w) \leq \langle\cdot,\cdot\rangle_{v}(\operatorname{vadd}(v)(w, \operatorname{act}(v)(t, e)), \operatorname{vadd}(v)(w, \operatorname{act}(v)(t, e))) \Rightarrow \\ \quad \quad \langle\cdot,\cdot\rangle_{v}(w, e) = 0 \end{array}}}
\vnbprop{hilbert-projection-exists}{\vnbstmt{\begin{array}{@{}l@{}} \forall v, k, i. \\ \quad (\operatorname{is\text{-}hilbert\text{-}space}(v) \wedge (\operatorname{is\text{-}submodule}(v, k) \wedge (\operatorname{is\text{-}closed}(\operatorname{ip\text{-}metric\text{-}space}(v), k) \wedge (i \in \operatorname{vec}(v))))) \Rightarrow  \\ \quad \quad \exists a \in k, b \in {k}^{\perp}.\; (i = \operatorname{vadd}(v)(a, b)) \end{array}}}
\vnbprop{hilbert-decomposition-unique}{\vnbstmt{\begin{array}{@{}l@{}} \forall v, k, a_{1}, a_{2}, b_{1}, b_{2}. \\ \quad \operatorname{is\text{-}complex\text{-}inner\text{-}product\text{-}space}(v) \wedge \\ \quad \operatorname{is\text{-}submodule}(v, k) \wedge \\ \quad a_{1} \in k \wedge \\ \quad a_{2} \in k \wedge \\ \quad b_{1} \in {k}^{\perp} \wedge \\ \quad b_{2} \in {k}^{\perp} \wedge \\ \quad \operatorname{vadd}(v)(a_{1}, b_{1}) = \operatorname{vadd}(v)(a_{2}, b_{2}) \Rightarrow \\ \quad \quad (a_{1} = a_{2}) \wedge (b_{1} = b_{2}) \end{array}}}
\vnbprop{hilbert-projection-point}{\vnbstmt{\begin{array}{@{}l@{}} \forall v, k, i. \\ \quad (\operatorname{is\text{-}hilbert\text{-}space}(v) \wedge (\operatorname{is\text{-}submodule}(v, k) \wedge (\operatorname{is\text{-}closed}(\operatorname{ip\text{-}metric\text{-}space}(v), k) \wedge (i \in \operatorname{vec}(v))))) \Rightarrow  \\ \quad \quad \exists a \in k.\; (\operatorname{vadd}(v)(i, \operatorname{vneg}(v)(a)) \in {k}^{\perp}) \end{array}}}
\vnbprop{orthogonal-projection-apply}{\vnbstmt{\begin{array}{@{}l@{}} \forall v, k, a, x. \\ \quad (\operatorname{is\text{-}orthogonal\text{-}projection}(v, k, a) \wedge (x \in \operatorname{vec}(v))) \Rightarrow  \\ \quad \quad (a(x) \in k) \wedge (\operatorname{vadd}(v)(x, \operatorname{vneg}(v)(a(x))) \in {k}^{\perp}) \end{array}}}
\vnbprop{orthogonal-projection-type}{\vnbstmt{\begin{array}{@{}l@{}} \forall v, k, a. \\ \quad \operatorname{is\text{-}orthogonal\text{-}projection}(v, k, a) \Rightarrow (a \in (\operatorname{vec}(v) \to \operatorname{vec}(v))) \end{array}}}
\vnbprop{orthogonal-projection-exists}{\vnbstmt{\begin{array}{@{}l@{}} \forall v, k. \\ \quad (\operatorname{is\text{-}hilbert\text{-}space}(v) \wedge (\operatorname{is\text{-}submodule}(v, k) \wedge \operatorname{is\text{-}closed}(\operatorname{ip\text{-}metric\text{-}space}(v), k))) \Rightarrow  \\ \quad \quad \exists \mathit{hpk}.\; \operatorname{is\text{-}orthogonal\text{-}projection}(v, k, \mathit{hpk}) \end{array}}}
\vnbprop{orthogonal-projection-unique}{\vnbstmt{\begin{array}{@{}l@{}} \forall v, k, a, b. \\ \quad \operatorname{is\text{-}complex\text{-}inner\text{-}product\text{-}space}(v) \wedge \\ \quad \operatorname{is\text{-}submodule}(v, k) \wedge \\ \quad \operatorname{is\text{-}orthogonal\text{-}projection}(v, k, a) \wedge \\ \quad \operatorname{is\text{-}orthogonal\text{-}projection}(v, k, b) \Rightarrow \\ \quad \quad a = b \end{array}}}
\vnbprop{orthogonal-projection-additive}{\vnbstmt{\begin{array}{@{}l@{}} \forall v, k, a, x, y. \\ \quad \operatorname{is\text{-}complex\text{-}inner\text{-}product\text{-}space}(v) \wedge \\ \quad \operatorname{is\text{-}submodule}(v, k) \wedge \\ \quad \operatorname{is\text{-}orthogonal\text{-}projection}(v, k, a) \wedge \\ \quad x \in \operatorname{vec}(v) \wedge \\ \quad y \in \operatorname{vec}(v) \Rightarrow \\ \quad \quad a(\operatorname{vadd}(v)(x, y)) = \operatorname{vadd}(v)(a(x), a(y)) \end{array}}}
\vnbprop{orthogonal-projection-homogeneous}{\vnbstmt{\begin{array}{@{}l@{}} \forall v, k, a, c, x. \\ \quad \operatorname{is\text{-}complex\text{-}inner\text{-}product\text{-}space}(v) \wedge \\ \quad \operatorname{is\text{-}submodule}(v, k) \wedge \\ \quad \operatorname{is\text{-}orthogonal\text{-}projection}(v, k, a) \wedge \\ \quad c \in \mathbb{C} \wedge \\ \quad x \in \operatorname{vec}(v) \Rightarrow \\ \quad \quad a(\operatorname{act}(v)(c, x)) = \operatorname{act}(v)(c, a(x)) \end{array}}}
\vnbprop{orthogonal-projection-norm-le}{\vnbstmt{\begin{array}{@{}l@{}} \forall v, k, a, x. \\ \quad \operatorname{is\text{-}complex\text{-}inner\text{-}product\text{-}space}(v) \wedge \\ \quad \operatorname{is\text{-}submodule}(v, k) \wedge \\ \quad \operatorname{is\text{-}orthogonal\text{-}projection}(v, k, a) \wedge \\ \quad x \in \operatorname{vec}(v) \Rightarrow \\ \quad \quad \lVert a(x) \rVert \leq \lVert x \rVert \end{array}}}
\vnbprop{orthogonal-projection-fixes}{\vnbstmt{\begin{array}{@{}l@{}} \forall v, k, a, y. \\ \quad \operatorname{is\text{-}complex\text{-}inner\text{-}product\text{-}space}(v) \wedge \\ \quad \operatorname{is\text{-}submodule}(v, k) \wedge \\ \quad \operatorname{is\text{-}orthogonal\text{-}projection}(v, k, a) \wedge \\ \quad y \in k \Rightarrow \\ \quad \quad a(y) = y \end{array}}}
\vnbprop{orthogonal-projection-idempotent}{\vnbstmt{\begin{array}{@{}l@{}} \forall v, k, a, x. \\ \quad \operatorname{is\text{-}complex\text{-}inner\text{-}product\text{-}space}(v) \wedge \\ \quad \operatorname{is\text{-}submodule}(v, k) \wedge \\ \quad \operatorname{is\text{-}orthogonal\text{-}projection}(v, k, a) \wedge \\ \quad x \in \operatorname{vec}(v) \Rightarrow \\ \quad \quad a(a(x)) = a(x) \end{array}}}
\vnbprop{cc-linear-functional-laws}{\vnbstmt{\begin{array}{@{}l@{}} \forall v, f. \\ \quad \operatorname{is\text{-}cc\text{-}linear\text{-}functional}(v, f) \Rightarrow  \\ \quad \quad f \in (\operatorname{vec}(v) \to \mathbb{C}) \wedge \\ \quad \quad \forall x, y \in \operatorname{vec}(v).\; (f(\operatorname{vadd}(v)(x, y)) = (f(x) + f(y))) \wedge \\ \quad \quad \forall c \in \mathbb{C}, x \in \operatorname{vec}(v).\; (f(\operatorname{act}(v)(c, x)) = (c \cdot f(x))) \end{array}}}
\vnbprop{cc-linear-functional-zero-neg}{\vnbstmt{\begin{array}{@{}l@{}} \forall v, f. \\ \quad (\operatorname{is\text{-}complex\text{-}inner\text{-}product\text{-}space}(v) \wedge \operatorname{is\text{-}cc\text{-}linear\text{-}functional}(v, f)) \Rightarrow  \\ \quad \quad (f(\operatorname{vzero}(v)) = 0) \wedge \forall x \in \operatorname{vec}(v).\; (f(\operatorname{vneg}(v)(x)) = -f(x)) \end{array}}}
\vnbprop{cc-linear-functional-kernel-membership}{\vnbstmt{\begin{array}{@{}l@{}} \forall v, f, x. \\ \quad (\operatorname{is\text{-}complex\text{-}inner\text{-}product\text{-}space}(v) \wedge (\operatorname{is\text{-}cc\text{-}linear\text{-}functional}(v, f) \wedge (x \in \operatorname{vec}(v)))) \Rightarrow  \\ \quad \quad (x \in \operatorname{preimage}(\operatorname{ip\text{-}metric\text{-}space}(v), f, \operatorname{closed\text{-}ball}(\operatorname{cc\text{-}ms}, 0, 0))) \Leftrightarrow (f(x) = 0) \end{array}}}
\vnbprop{cc-linear-functional-kernel-submodule}{\vnbstmt{\begin{array}{@{}l@{}} \forall v, f. \\ \quad (\operatorname{is\text{-}complex\text{-}inner\text{-}product\text{-}space}(v) \wedge \operatorname{is\text{-}cc\text{-}linear\text{-}functional}(v, f)) \Rightarrow  \\ \quad \quad \operatorname{is\text{-}submodule}(v, \operatorname{preimage}(\operatorname{ip\text{-}metric\text{-}space}(v), f, \operatorname{closed\text{-}ball}(\operatorname{cc\text{-}ms}, 0, 0))) \end{array}}}
\vnbprop{cc-linear-functional-kernel-closed}{\vnbstmt{\begin{array}{@{}l@{}} \forall v, f. \\ \quad (\operatorname{is\text{-}complex\text{-}inner\text{-}product\text{-}space}(v) \wedge \operatorname{is\text{-}continuous}(\operatorname{ip\text{-}metric\text{-}space}(v), \operatorname{cc\text{-}ms}, f)) \Rightarrow  \\ \quad \quad \operatorname{is\text{-}closed}(\operatorname{ip\text{-}metric\text{-}space}(v), \operatorname{preimage}(\operatorname{ip\text{-}metric\text{-}space}(v), f, \operatorname{closed\text{-}ball}(\operatorname{cc\text{-}ms}, 0, 0))) \end{array}}}
\vnbprop{cips-ip-separates}{\vnbstmt{\begin{array}{@{}l@{}} \forall v, e_{1}, e_{2}. \\ \quad \operatorname{is\text{-}complex\text{-}inner\text{-}product\text{-}space}(v) \wedge \\ \quad e_{1} \in \operatorname{vec}(v) \wedge \\ \quad e_{2} \in \operatorname{vec}(v) \wedge \\ \quad \forall x \in \operatorname{vec}(v).\; (\langle\cdot,\cdot\rangle_{v}(x, e_{1}) = \langle\cdot,\cdot\rangle_{v}(x, e_{2})) \Rightarrow \\ \quad \quad e_{1} = e_{2} \end{array}}}
\vnbprop{cips-ip-magnitude-le}{\vnbstmt{\begin{array}{@{}l@{}} \forall v, x, e. \\ \quad (\operatorname{is\text{-}complex\text{-}inner\text{-}product\text{-}space}(v) \wedge ((x \in \operatorname{vec}(v)) \wedge (e \in \operatorname{vec}(v)))) \Rightarrow  \\ \quad \quad \operatorname{magnitude}(\langle\cdot,\cdot\rangle_{v}(x, e)) \leq (\lVert x \rVert \cdot \lVert e \rVert) \end{array}}}
\vnbprop{cips-ip-bound-least}{\vnbstmt{\begin{array}{@{}l@{}} \forall v, e, c. \\ \quad \operatorname{is\text{-}complex\text{-}inner\text{-}product\text{-}space}(v) \wedge \\ \quad e \in \operatorname{vec}(v) \wedge \\ \quad c \in \mathbb{R} \wedge \\ \quad 0 \leq c \wedge \\ \quad \forall x \in \operatorname{vec}(v).\; (\operatorname{magnitude}(\langle\cdot,\cdot\rangle_{v}(x, e)) \leq (c \cdot \lVert x \rVert)) \Rightarrow \\ \quad \quad \lVert e \rVert \leq c \end{array}}}
\vnbprop{riesz-representation-of-kernel}{\vnbstmt{\begin{array}{@{}l@{}} \forall v, f, k. \\ \quad \operatorname{is\text{-}hilbert\text{-}space}(v) \wedge \\ \quad \operatorname{is\text{-}cc\text{-}linear\text{-}functional}(v, f) \wedge \\ \quad \operatorname{is\text{-}submodule}(v, k) \wedge \\ \quad \operatorname{is\text{-}closed}(\operatorname{ip\text{-}metric\text{-}space}(v), k) \wedge \\ \quad \forall x \in \operatorname{vec}(v).\; ((x \in k) \Leftrightarrow (f(x) = 0)) \Rightarrow \\ \quad \quad \exists e \in \operatorname{vec}(v).\; \forall \mathit{hrx} \in \operatorname{vec}(v).\; (f(\mathit{hrx}) = \langle\cdot,\cdot\rangle_{v}(\mathit{hrx}, e)) \end{array}}}
\vnbprop{riesz-representation-exists}{\vnbstmt{\begin{array}{@{}l@{}} \forall v, f. \\ \quad \operatorname{is\text{-}hilbert\text{-}space}(v) \wedge \\ \quad \operatorname{is\text{-}cc\text{-}linear\text{-}functional}(v, f) \wedge \\ \quad \operatorname{is\text{-}continuous}(\operatorname{ip\text{-}metric\text{-}space}(v), \operatorname{cc\text{-}ms}, f) \Rightarrow \\ \quad \quad \exists e \in \operatorname{vec}(v).\; \forall x \in \operatorname{vec}(v).\; (f(x) = \langle\cdot,\cdot\rangle_{v}(x, e)) \end{array}}}
\vnbprop{riesz-representation}{\vnbstmt{\begin{array}{@{}l@{}} \forall v, f. \\ \quad \operatorname{is\text{-}hilbert\text{-}space}(v) \wedge \\ \quad \operatorname{is\text{-}cc\text{-}linear\text{-}functional}(v, f) \wedge \\ \quad \operatorname{is\text{-}continuous}(\operatorname{ip\text{-}metric\text{-}space}(v), \operatorname{cc\text{-}ms}, f) \Rightarrow \\ \quad \quad \exists e \in \operatorname{vec}(v). \\ \quad \quad \quad \forall x \in \operatorname{vec}(v).\; (f(x) = \langle\cdot,\cdot\rangle_{v}(x, e)) \wedge \\ \quad \quad \quad \forall x \in \operatorname{vec}(v).\; (\operatorname{magnitude}(f(x)) \leq (\lVert x \rVert \cdot \lVert e \rVert)) \wedge \\ \quad \quad \quad \forall c \in \mathbb{R}. \\ \quad \quad \quad \quad ((0 \leq c) \wedge \forall x \in \operatorname{vec}(v).\; (\operatorname{magnitude}(f(x)) \leq (c \cdot \lVert x \rVert))) \Rightarrow  \\ \quad \quad \quad \quad \quad \lVert e \rVert \leq c \end{array}}}
\vnbprop{riesz-representation-unique}{\vnbstmt{\begin{array}{@{}l@{}} \forall v, f, e_{1}, e_{2}. \\ \quad \operatorname{is\text{-}complex\text{-}inner\text{-}product\text{-}space}(v) \wedge \\ \quad e_{1} \in \operatorname{vec}(v) \wedge \\ \quad e_{2} \in \operatorname{vec}(v) \wedge \\ \quad \forall x \in \operatorname{vec}(v).\; (f(x) = \langle\cdot,\cdot\rangle_{v}(x, e_{1})) \wedge \\ \quad \forall x \in \operatorname{vec}(v).\; (f(x) = \langle\cdot,\cdot\rangle_{v}(x, e_{2})) \Rightarrow \\ \quad \quad e_{1} = e_{2} \end{array}}}
\vnbprop{ip-closure-elim}{\vnbstmt{\begin{array}{@{}l@{}} \forall v, a, x, e. \\ \quad (\operatorname{is\text{-}complex\text{-}inner\text{-}product\text{-}space}(v) \wedge ((x \in \operatorname{closure}(\operatorname{ip\text{-}metric\text{-}space}(v), a)) \wedge \operatorname{pos\text{-}rr}(e))) \Rightarrow  \\ \quad \quad \exists y \in a. \\ \quad \quad \quad \langle\cdot,\cdot\rangle_{v}(\operatorname{vadd}(v)(x, \operatorname{vneg}(v)(y)), \operatorname{vadd}(v)(x, \operatorname{vneg}(v)(y))) \leq e \end{array}}}
\vnbprop{ip-closure-intro}{\vnbstmt{\begin{array}{@{}l@{}} \forall v, a, x. \\ \quad \operatorname{is\text{-}complex\text{-}inner\text{-}product\text{-}space}(v) \wedge \\ \quad a \subseteq \operatorname{vec}(v) \wedge \\ \quad x \in \operatorname{vec}(v) \wedge \\ \quad \forall e. \\ \quad \quad \operatorname{pos\text{-}rr}(e) \Rightarrow  \\ \quad \quad \quad \exists y \in a. \\ \quad \quad \quad \quad \langle\cdot,\cdot\rangle_{v}(\operatorname{vadd}(v)(x, \operatorname{vneg}(v)(y)), \operatorname{vadd}(v)(x, \operatorname{vneg}(v)(y))) \leq e \Rightarrow \\ \quad \quad x \in \operatorname{closure}(\operatorname{ip\text{-}metric\text{-}space}(v), a) \end{array}}}
\vnbprop{ip-closure-submodule}{\vnbstmt{\begin{array}{@{}l@{}} \forall v, k. \\ \quad (\operatorname{is\text{-}complex\text{-}inner\text{-}product\text{-}space}(v) \wedge \operatorname{is\text{-}submodule}(v, k)) \Rightarrow  \\ \quad \quad \operatorname{is\text{-}submodule}(v, \operatorname{closure}(\operatorname{ip\text{-}metric\text{-}space}(v), k)) \end{array}}}
\vnbprop{hilbert-dense-iff-perp-zero}{\vnbstmt{\begin{array}{@{}l@{}} \forall v, k. \\ \quad (\operatorname{is\text{-}hilbert\text{-}space}(v) \wedge \operatorname{is\text{-}submodule}(v, k)) \Rightarrow  \\ \quad \quad (\operatorname{closure}(\operatorname{ip\text{-}metric\text{-}space}(v), k) = \operatorname{vec}(v)) \Leftrightarrow \forall z \in {k}^{\perp}.\; (z = \operatorname{vzero}(v)) \end{array}}}

\vnbbookfile{theorem-library/rake-asm-b3}{theorem-library/rake-asm-b3}
\noindent{\small theorem-library/rake-asm-b3.scm -- RAKE-ASM-B, agent RB-3 (2026-10-12).  Prefix rb3-.}\par
\vnbprop{cc-scalar-ring-carr}{\vnbstmt{(\operatorname{carr}(\operatorname{normed\text{-}field\text{-}as\text{-}commutative\text{-}ring}(\operatorname{cc\text{-}normed\text{-}field})) \simeq \mathbb{C})}}
\vnbprop{cc-scalar-ring-add}{\vnbstmt{(\operatorname{add}(\operatorname{normed\text{-}field\text{-}as\text{-}commutative\text{-}ring}(\operatorname{cc\text{-}normed\text{-}field})) \simeq (\lambda {[}\operatorname{x\_}, \operatorname{y\_}] \in (\mathbb{C} \times \mathbb{C}).\; (\operatorname{x\_} + \operatorname{y\_})))}}
\vnbprop{cc-scalar-ring-mul}{\vnbstmt{(\operatorname{mul}(\operatorname{normed\text{-}field\text{-}as\text{-}commutative\text{-}ring}(\operatorname{cc\text{-}normed\text{-}field})) \simeq (\lambda {[}\operatorname{x\_}, \operatorname{y\_}] \in (\mathbb{C} \times \mathbb{C}).\; (\operatorname{x\_} \cdot \operatorname{y\_})))}}
\vnbprop{cc-scalar-ring-one}{\vnbstmt{(\operatorname{one}(\operatorname{normed\text{-}field\text{-}as\text{-}commutative\text{-}ring}(\operatorname{cc\text{-}normed\text{-}field})) \simeq 1)}}
\vnbprop{cc-scalar-ring-zero}{\vnbstmt{(\operatorname{zero}(\operatorname{normed\text{-}field\text{-}as\text{-}commutative\text{-}ring}(\operatorname{cc\text{-}normed\text{-}field})) \simeq 0)}}
\vnbprop{rb3-c-carr}{\vnbstmt{\begin{array}{@{}l@{}} \forall m. \\ \quad (\operatorname{is\text{-}module}(m) \wedge (\operatorname{scal}(m) = \operatorname{normed\text{-}field\text{-}as\text{-}commutative\text{-}ring}(\operatorname{cc\text{-}normed\text{-}field}))) \Rightarrow  \\ \quad \quad \operatorname{carr}(\operatorname{scal}(m)) = \mathbb{C} \end{array}}}
\vnbprop{rb3-c-in-carr}{\vnbstmt{\begin{array}{@{}l@{}} \forall m.; \forall r \in \mathbb{C}. \\ \quad (\operatorname{is\text{-}module}(m) \wedge (\operatorname{scal}(m) = \operatorname{normed\text{-}field\text{-}as\text{-}commutative\text{-}ring}(\operatorname{cc\text{-}normed\text{-}field}))) \Rightarrow  \\ \quad \quad r \in \operatorname{carr}(\operatorname{scal}(m)) \end{array}}}
\vnbprop{rb3-c-add-eq}{\vnbstmt{\begin{array}{@{}l@{}} \forall m.; \forall r, s \in \mathbb{C}. \\ \quad (\operatorname{is\text{-}module}(m) \wedge (\operatorname{scal}(m) = \operatorname{normed\text{-}field\text{-}as\text{-}commutative\text{-}ring}(\operatorname{cc\text{-}normed\text{-}field}))) \Rightarrow  \\ \quad \quad \operatorname{add}(\operatorname{scal}(m))(r, s) = (r + s) \end{array}}}
\vnbprop{rb3-c-mul-eq}{\vnbstmt{\begin{array}{@{}l@{}} \forall m.; \forall r, s \in \mathbb{C}. \\ \quad (\operatorname{is\text{-}module}(m) \wedge (\operatorname{scal}(m) = \operatorname{normed\text{-}field\text{-}as\text{-}commutative\text{-}ring}(\operatorname{cc\text{-}normed\text{-}field}))) \Rightarrow  \\ \quad \quad \operatorname{mul}(\operatorname{scal}(m))(r, s) = (r \cdot s) \end{array}}}
\vnbprop{rb3-c-one-eq}{\vnbstmt{\begin{array}{@{}l@{}} \forall m. \\ \quad (\operatorname{is\text{-}module}(m) \wedge (\operatorname{scal}(m) = \operatorname{normed\text{-}field\text{-}as\text{-}commutative\text{-}ring}(\operatorname{cc\text{-}normed\text{-}field}))) \Rightarrow  \\ \quad \quad \operatorname{one}(\operatorname{scal}(m)) = 1 \end{array}}}
\vnbprop{rb3-c-act-in}{\vnbstmt{\begin{array}{@{}l@{}} \forall m.; \forall r \in \mathbb{C}, x \in \operatorname{vec}(m). \\ \quad (\operatorname{is\text{-}module}(m) \wedge (\operatorname{scal}(m) = \operatorname{normed\text{-}field\text{-}as\text{-}commutative\text{-}ring}(\operatorname{cc\text{-}normed\text{-}field}))) \Rightarrow  \\ \quad \quad \operatorname{act}(m)(r, x) \in \operatorname{vec}(m) \end{array}}}
\vnbprop{rb3-c-act-add}{\vnbstmt{\begin{array}{@{}l@{}} \forall m.; \forall r, s \in \mathbb{C}, x \in \operatorname{vec}(m). \\ \quad (\operatorname{is\text{-}module}(m) \wedge (\operatorname{scal}(m) = \operatorname{normed\text{-}field\text{-}as\text{-}commutative\text{-}ring}(\operatorname{cc\text{-}normed\text{-}field}))) \Rightarrow  \\ \quad \quad \operatorname{act}(m)((r + s), x) = \operatorname{vadd}(m)(\operatorname{act}(m)(r, x), \operatorname{act}(m)(s, x)) \end{array}}}
\vnbprop{rb3-c-act-mul}{\vnbstmt{\begin{array}{@{}l@{}} \forall m.; \forall r, s \in \mathbb{C}, x \in \operatorname{vec}(m). \\ \quad (\operatorname{is\text{-}module}(m) \wedge (\operatorname{scal}(m) = \operatorname{normed\text{-}field\text{-}as\text{-}commutative\text{-}ring}(\operatorname{cc\text{-}normed\text{-}field}))) \Rightarrow  \\ \quad \quad \operatorname{act}(m)((r \cdot s), x) = \operatorname{act}(m)(r, \operatorname{act}(m)(s, x)) \end{array}}}
\vnbprop{rb3-c-act-vadd}{\vnbstmt{\begin{array}{@{}l@{}} \forall m.; \forall r \in \mathbb{C}, x, y \in \operatorname{vec}(m). \\ \quad (\operatorname{is\text{-}module}(m) \wedge (\operatorname{scal}(m) = \operatorname{normed\text{-}field\text{-}as\text{-}commutative\text{-}ring}(\operatorname{cc\text{-}normed\text{-}field}))) \Rightarrow  \\ \quad \quad \operatorname{act}(m)(r, \operatorname{vadd}(m)(x, y)) = \operatorname{vadd}(m)(\operatorname{act}(m)(r, x), \operatorname{act}(m)(r, y)) \end{array}}}
\vnbprop{rb3-tuple6}{\vnbstmt{\begin{array}{@{}l@{}} \forall a_{2}, a_{3}, a_{4}, a_{5}, a_{6}. \\ \quad ((a_{2} \in \mathbf{Set}) \wedge ((a_{3} \in \mathbf{Set}) \wedge ((a_{4} \in \mathbf{Set}) \wedge ((a_{5} \in \mathbf{Set}) \wedge (a_{6} \in \mathbf{Set}))))) \Rightarrow  \\ \quad \quad (\operatorname{scal}({[}\operatorname{normed\text{-}field\text{-}as\text{-}commutative\text{-}ring}(\operatorname{rr\text{-}normed\text{-}field}), a_{2}, a_{3}, a_{4}, a_{5}, a_{6}]) \simeq \operatorname{normed\text{-}field\text{-}as\text{-}commutative\text{-}ring}(\operatorname{rr\text{-}normed\text{-}field})) \wedge \\ \quad \quad (\operatorname{vec}({[}\operatorname{normed\text{-}field\text{-}as\text{-}commutative\text{-}ring}(\operatorname{rr\text{-}normed\text{-}field}), a_{2}, a_{3}, a_{4}, a_{5}, a_{6}]) \simeq a_{2}) \wedge \\ \quad \quad (\operatorname{vadd}({[}\operatorname{normed\text{-}field\text{-}as\text{-}commutative\text{-}ring}(\operatorname{rr\text{-}normed\text{-}field}), a_{2}, a_{3}, a_{4}, a_{5}, a_{6}]) \simeq a_{3}) \wedge \\ \quad \quad (\operatorname{vzero}({[}\operatorname{normed\text{-}field\text{-}as\text{-}commutative\text{-}ring}(\operatorname{rr\text{-}normed\text{-}field}), a_{2}, a_{3}, a_{4}, a_{5}, a_{6}]) \simeq a_{4}) \wedge \\ \quad \quad (\operatorname{vneg}({[}\operatorname{normed\text{-}field\text{-}as\text{-}commutative\text{-}ring}(\operatorname{rr\text{-}normed\text{-}field}), a_{2}, a_{3}, a_{4}, a_{5}, a_{6}]) \simeq a_{5}) \wedge \\ \quad \quad (\operatorname{act}({[}\operatorname{normed\text{-}field\text{-}as\text{-}commutative\text{-}ring}(\operatorname{rr\text{-}normed\text{-}field}), a_{2}, a_{3}, a_{4}, a_{5}, a_{6}]) \simeq a_{6}) \wedge \\ \quad \quad \operatorname{length}({[}\operatorname{normed\text{-}field\text{-}as\text{-}commutative\text{-}ring}(\operatorname{rr\text{-}normed\text{-}field}), a_{2}, a_{3}, a_{4}, a_{5}, a_{6}]) = 6 \end{array}}}
\vnbprop{rb3-mr-exists}{\vnbstmt{\begin{array}{@{}l@{}} \forall m. \\ \quad (\operatorname{is\text{-}module}(m) \wedge (\operatorname{scal}(m) = \operatorname{normed\text{-}field\text{-}as\text{-}commutative\text{-}ring}(\operatorname{cc\text{-}normed\text{-}field}))) \Rightarrow  \\ \quad \quad \exists r. \\ \quad \quad \quad \operatorname{is\text{-}module}(r) \wedge \\ \quad \quad \quad \operatorname{scal}(r) = \operatorname{normed\text{-}field\text{-}as\text{-}commutative\text{-}ring}(\operatorname{rr\text{-}normed\text{-}field}) \wedge \\ \quad \quad \quad \operatorname{vec}(r) = \operatorname{vec}(m) \wedge \\ \quad \quad \quad \operatorname{vadd}(r) = \operatorname{vadd}(m) \wedge \\ \quad \quad \quad \operatorname{vzero}(r) = \operatorname{vzero}(m) \wedge \\ \quad \quad \quad \operatorname{vneg}(r) = \operatorname{vneg}(m) \wedge \\ \quad \quad \quad \forall a \in \mathbb{R}, x \in \operatorname{vec}(m).\; (\operatorname{act}(r)(a, x) = \operatorname{act}(m)(a, x)) \end{array}}}
\vnbprop{rb3-mr-sublinear}{\vnbstmt{\begin{array}{@{}l@{}} \forall m, r, p. \\ \quad \operatorname{is\text{-}module}(m) \wedge \\ \quad \operatorname{scal}(m) = \operatorname{normed\text{-}field\text{-}as\text{-}commutative\text{-}ring}(\operatorname{cc\text{-}normed\text{-}field}) \wedge \\ \quad \operatorname{is\text{-}module}(r) \wedge \\ \quad \operatorname{scal}(r) = \operatorname{normed\text{-}field\text{-}as\text{-}commutative\text{-}ring}(\operatorname{rr\text{-}normed\text{-}field}) \wedge \\ \quad \operatorname{vec}(r) = \operatorname{vec}(m) \wedge \\ \quad \operatorname{vadd}(r) = \operatorname{vadd}(m) \wedge \\ \quad \operatorname{vzero}(r) = \operatorname{vzero}(m) \wedge \\ \quad \operatorname{vneg}(r) = \operatorname{vneg}(m) \wedge \\ \quad \forall a \in \mathbb{R}, x \in \operatorname{vec}(m).\; (\operatorname{act}(r)(a, x) = \operatorname{act}(m)(a, x)) \wedge \\ \quad \operatorname{is\text{-}seminorm}(\operatorname{cc\text{-}normed\text{-}field}, m, p) \Rightarrow \\ \quad \quad \operatorname{is\text{-}sublinear}(r, p) \end{array}}}
\vnbprop{rb3-mr-submodule}{\vnbstmt{\begin{array}{@{}l@{}} \forall m, r, s. \\ \quad \operatorname{is\text{-}module}(m) \wedge \\ \quad \operatorname{scal}(m) = \operatorname{normed\text{-}field\text{-}as\text{-}commutative\text{-}ring}(\operatorname{cc\text{-}normed\text{-}field}) \wedge \\ \quad \operatorname{is\text{-}module}(r) \wedge \\ \quad \operatorname{scal}(r) = \operatorname{normed\text{-}field\text{-}as\text{-}commutative\text{-}ring}(\operatorname{rr\text{-}normed\text{-}field}) \wedge \\ \quad \operatorname{vec}(r) = \operatorname{vec}(m) \wedge \\ \quad \operatorname{vadd}(r) = \operatorname{vadd}(m) \wedge \\ \quad \operatorname{vzero}(r) = \operatorname{vzero}(m) \wedge \\ \quad \operatorname{vneg}(r) = \operatorname{vneg}(m) \wedge \\ \quad \forall a \in \mathbb{R}, x \in \operatorname{vec}(m).\; (\operatorname{act}(r)(a, x) = \operatorname{act}(m)(a, x)) \wedge \\ \quad \operatorname{is\text{-}submodule}(m, s) \Rightarrow \\ \quad \quad \operatorname{is\text{-}submodule}(r, s) \end{array}}}
\vnbprop{rb3-mr-re-linear}{\vnbstmt{\begin{array}{@{}l@{}} \forall m, r, s, f.; \forall a \in (s \to \mathbb{R}). \\ \quad \operatorname{is\text{-}module}(m) \wedge \\ \quad \operatorname{scal}(m) = \operatorname{normed\text{-}field\text{-}as\text{-}commutative\text{-}ring}(\operatorname{cc\text{-}normed\text{-}field}) \wedge \\ \quad \operatorname{is\text{-}module}(r) \wedge \\ \quad \operatorname{scal}(r) = \operatorname{normed\text{-}field\text{-}as\text{-}commutative\text{-}ring}(\operatorname{rr\text{-}normed\text{-}field}) \wedge \\ \quad \operatorname{vec}(r) = \operatorname{vec}(m) \wedge \\ \quad \operatorname{vadd}(r) = \operatorname{vadd}(m) \wedge \\ \quad \operatorname{vzero}(r) = \operatorname{vzero}(m) \wedge \\ \quad \operatorname{vneg}(r) = \operatorname{vneg}(m) \wedge \\ \quad \forall b \in \mathbb{R}, x \in \operatorname{vec}(m).\; (\operatorname{act}(r)(b, x) = \operatorname{act}(m)(b, x)) \wedge \\ \quad \operatorname{is\text{-}submodule}(m, s) \wedge \\ \quad \operatorname{is\text{-}k\text{-}linear\text{-}on}(m, s, f) \wedge \\ \quad \forall z \in s.\; (a(z) \simeq \operatorname{real\text{-}part}(f(z))) \Rightarrow \\ \quad \quad \operatorname{is\text{-}linear\text{-}functional\text{-}on}(r, s, a) \end{array}}}
\vnbprop{rb3-cf-re}{\vnbstmt{\begin{array}{@{}l@{}} \forall m.; \forall n \in (\operatorname{vec}(m) \to \mathbb{R}), a \in (\operatorname{vec}(m) \to \mathbb{C}), y \in \operatorname{vec}(m). \\ \quad \operatorname{is\text{-}module}(m) \wedge \\ \quad \operatorname{scal}(m) = \operatorname{normed\text{-}field\text{-}as\text{-}commutative\text{-}ring}(\operatorname{cc\text{-}normed\text{-}field}) \wedge \\ \quad \forall x, \mathit{ray} \in \operatorname{vec}(m).\; (n(\operatorname{vadd}(m)(x, \mathit{ray})) = (n(x) + n(\mathit{ray}))) \wedge \\ \quad \forall r \in \mathbb{R}, \mathit{rhx} \in \operatorname{vec}(m).\; (n(\operatorname{act}(m)(r, \mathit{rhx})) = (r \cdot n(\mathit{rhx}))) \wedge \\ \quad \forall \mathit{cvy} \in \operatorname{vec}(m).\; (a(\mathit{cvy}) \simeq (n(\mathit{cvy}) + (-n(\operatorname{act}(m)(+i, \mathit{cvy})) \cdot +i))) \Rightarrow \\ \quad \quad \operatorname{real\text{-}part}(a(y)) = n(y) \end{array}}}
\vnbprop{rb3-cf-add}{\vnbstmt{\begin{array}{@{}l@{}} \forall m.; \forall n \in (\operatorname{vec}(m) \to \mathbb{R}), a \in (\operatorname{vec}(m) \to \mathbb{C}), x, y \in \operatorname{vec}(m). \\ \quad \operatorname{is\text{-}module}(m) \wedge \\ \quad \operatorname{scal}(m) = \operatorname{normed\text{-}field\text{-}as\text{-}commutative\text{-}ring}(\operatorname{cc\text{-}normed\text{-}field}) \wedge \\ \quad \forall \mathit{rax}, \mathit{ray} \in \operatorname{vec}(m).\; (n(\operatorname{vadd}(m)(\mathit{rax}, \mathit{ray})) = (n(\mathit{rax}) + n(\mathit{ray}))) \wedge \\ \quad \forall r \in \mathbb{R}, \mathit{rhx} \in \operatorname{vec}(m).\; (n(\operatorname{act}(m)(r, \mathit{rhx})) = (r \cdot n(\mathit{rhx}))) \wedge \\ \quad \forall \mathit{cvy} \in \operatorname{vec}(m).\; (a(\mathit{cvy}) \simeq (n(\mathit{cvy}) + (-n(\operatorname{act}(m)(+i, \mathit{cvy})) \cdot +i))) \Rightarrow \\ \quad \quad a(\operatorname{vadd}(m)(x, y)) = (a(x) + a(y)) \end{array}}}
\vnbprop{rb3-cf-rhom}{\vnbstmt{\begin{array}{@{}l@{}} \forall m.; \forall n \in (\operatorname{vec}(m) \to \mathbb{R}), a \in (\operatorname{vec}(m) \to \mathbb{C}), r \in \mathbb{R}, x \in \operatorname{vec}(m). \\ \quad \operatorname{is\text{-}module}(m) \wedge \\ \quad \operatorname{scal}(m) = \operatorname{normed\text{-}field\text{-}as\text{-}commutative\text{-}ring}(\operatorname{cc\text{-}normed\text{-}field}) \wedge \\ \quad \forall \mathit{rax}, y \in \operatorname{vec}(m).\; (n(\operatorname{vadd}(m)(\mathit{rax}, y)) = (n(\mathit{rax}) + n(y))) \wedge \\ \quad \forall \mathit{rhr} \in \mathbb{R}, \mathit{rhx} \in \operatorname{vec}(m).\; (n(\operatorname{act}(m)(\mathit{rhr}, \mathit{rhx})) = (\mathit{rhr} \cdot n(\mathit{rhx}))) \wedge \\ \quad \forall \mathit{cvy} \in \operatorname{vec}(m).\; (a(\mathit{cvy}) \simeq (n(\mathit{cvy}) + (-n(\operatorname{act}(m)(+i, \mathit{cvy})) \cdot +i))) \Rightarrow \\ \quad \quad a(\operatorname{act}(m)(r, x)) = (r \cdot a(x)) \end{array}}}
\vnbprop{rb3-cf-ihom}{\vnbstmt{\begin{array}{@{}l@{}} \forall m.; \forall n \in (\operatorname{vec}(m) \to \mathbb{R}), a \in (\operatorname{vec}(m) \to \mathbb{C}), x \in \operatorname{vec}(m). \\ \quad \operatorname{is\text{-}module}(m) \wedge \\ \quad \operatorname{scal}(m) = \operatorname{normed\text{-}field\text{-}as\text{-}commutative\text{-}ring}(\operatorname{cc\text{-}normed\text{-}field}) \wedge \\ \quad \forall \mathit{rax}, y \in \operatorname{vec}(m).\; (n(\operatorname{vadd}(m)(\mathit{rax}, y)) = (n(\mathit{rax}) + n(y))) \wedge \\ \quad \forall r \in \mathbb{R}, \mathit{rhx} \in \operatorname{vec}(m).\; (n(\operatorname{act}(m)(r, \mathit{rhx})) = (r \cdot n(\mathit{rhx}))) \wedge \\ \quad \forall \mathit{cvy} \in \operatorname{vec}(m).\; (a(\mathit{cvy}) \simeq (n(\mathit{cvy}) + (-n(\operatorname{act}(m)(+i, \mathit{cvy})) \cdot +i))) \Rightarrow \\ \quad \quad a(\operatorname{act}(m)(+i, x)) = (+i \cdot a(x)) \end{array}}}
\vnbprop{rb3-cf-abhom}{\vnbstmt{\begin{array}{@{}l@{}} \forall m.; \forall n \in (\operatorname{vec}(m) \to \mathbb{R}), a \in (\operatorname{vec}(m) \to \mathbb{C}), \alpha, b \in \mathbb{R}, x \in \operatorname{vec}(m). \\ \quad \operatorname{is\text{-}module}(m) \wedge \\ \quad \operatorname{scal}(m) = \operatorname{normed\text{-}field\text{-}as\text{-}commutative\text{-}ring}(\operatorname{cc\text{-}normed\text{-}field}) \wedge \\ \quad \forall \mathit{rax}, y \in \operatorname{vec}(m).\; (n(\operatorname{vadd}(m)(\mathit{rax}, y)) = (n(\mathit{rax}) + n(y))) \wedge \\ \quad \forall r \in \mathbb{R}, \mathit{rhx} \in \operatorname{vec}(m).\; (n(\operatorname{act}(m)(r, \mathit{rhx})) = (r \cdot n(\mathit{rhx}))) \wedge \\ \quad \forall \mathit{cvy} \in \operatorname{vec}(m).\; (a(\mathit{cvy}) \simeq (n(\mathit{cvy}) + (-n(\operatorname{act}(m)(+i, \mathit{cvy})) \cdot +i))) \Rightarrow \\ \quad \quad a(\operatorname{act}(m)((\alpha + (+i \cdot b)), x)) = ((\alpha + (+i \cdot b)) \cdot a(x)) \end{array}}}
\vnbprop{rb3-cf-chom}{\vnbstmt{\begin{array}{@{}l@{}} \forall m.; \forall n \in (\operatorname{vec}(m) \to \mathbb{R}), a \in (\operatorname{vec}(m) \to \mathbb{C}), b \in \mathbb{C}, x \in \operatorname{vec}(m). \\ \quad \operatorname{is\text{-}module}(m) \wedge \\ \quad \operatorname{scal}(m) = \operatorname{normed\text{-}field\text{-}as\text{-}commutative\text{-}ring}(\operatorname{cc\text{-}normed\text{-}field}) \wedge \\ \quad \forall \mathit{rax}, y \in \operatorname{vec}(m).\; (n(\operatorname{vadd}(m)(\mathit{rax}, y)) = (n(\mathit{rax}) + n(y))) \wedge \\ \quad \forall r \in \mathbb{R}, \mathit{rhx} \in \operatorname{vec}(m).\; (n(\operatorname{act}(m)(r, \mathit{rhx})) = (r \cdot n(\mathit{rhx}))) \wedge \\ \quad \forall \mathit{cvy} \in \operatorname{vec}(m).\; (a(\mathit{cvy}) \simeq (n(\mathit{cvy}) + (-n(\operatorname{act}(m)(+i, \mathit{cvy})) \cdot +i))) \Rightarrow \\ \quad \quad a(\operatorname{act}(m)(b, x)) = (b \cdot a(x)) \end{array}}}
\vnbprop{rb3-cf-klinear}{\vnbstmt{\begin{array}{@{}l@{}} \forall m.; \forall n \in (\operatorname{vec}(m) \to \mathbb{R}), a \in (\operatorname{vec}(m) \to \mathbb{C}). \\ \quad \operatorname{is\text{-}module}(m) \wedge \\ \quad \operatorname{scal}(m) = \operatorname{normed\text{-}field\text{-}as\text{-}commutative\text{-}ring}(\operatorname{cc\text{-}normed\text{-}field}) \wedge \\ \quad \forall x, y \in \operatorname{vec}(m).\; (n(\operatorname{vadd}(m)(x, y)) = (n(x) + n(y))) \wedge \\ \quad \forall r \in \mathbb{R}, \mathit{rhx} \in \operatorname{vec}(m).\; (n(\operatorname{act}(m)(r, \mathit{rhx})) = (r \cdot n(\mathit{rhx}))) \wedge \\ \quad \forall \mathit{cvy} \in \operatorname{vec}(m).\; (a(\mathit{cvy}) \simeq (n(\mathit{cvy}) + (-n(\operatorname{act}(m)(+i, \mathit{cvy})) \cdot +i))) \Rightarrow \\ \quad \quad \operatorname{is\text{-}k\text{-}linear}(m, a) \end{array}}}
\vnbprop{rb3-eq-decomp}{\vnbstmt{\forall z \in \mathbb{C}, w.\; ((w = (\operatorname{real\text{-}part}(z) + (+i \cdot \operatorname{imag\text{-}part}(z)))) \Rightarrow (w = z))}}
\vnbprop{rb3-re-im-i}{\vnbstmt{(\operatorname{real\text{-}part}(+i) = 0) \wedge (\operatorname{imag\text{-}part}(+i) = 1)}}
\vnbprop{rb3-cf-extends}{\vnbstmt{\begin{array}{@{}l@{}} \forall m.; \forall n \in (\operatorname{vec}(m) \to \mathbb{R}), a \in (\operatorname{vec}(m) \to \mathbb{C}), s.; \forall f \in (s \to \mathbb{C}). \\ \quad \operatorname{is\text{-}module}(m) \wedge \\ \quad \operatorname{scal}(m) = \operatorname{normed\text{-}field\text{-}as\text{-}commutative\text{-}ring}(\operatorname{cc\text{-}normed\text{-}field}) \wedge \\ \quad \forall x, y \in \operatorname{vec}(m).\; (n(\operatorname{vadd}(m)(x, y)) = (n(x) + n(y))) \wedge \\ \quad \forall r \in \mathbb{R}, \mathit{rhx} \in \operatorname{vec}(m).\; (n(\operatorname{act}(m)(r, \mathit{rhx})) = (r \cdot n(\mathit{rhx}))) \wedge \\ \quad \forall \mathit{cvy} \in \operatorname{vec}(m).\; (a(\mathit{cvy}) \simeq (n(\mathit{cvy}) + (-n(\operatorname{act}(m)(+i, \mathit{cvy})) \cdot +i))) \wedge \\ \quad \operatorname{is\text{-}submodule}(m, s) \wedge \\ \quad \forall \mathit{eix} \in s.\; (f(\operatorname{act}(m)(+i, \mathit{eix})) = (+i \cdot f(\mathit{eix}))) \wedge \\ \quad \forall \mathit{erx} \in s.\; (n(\mathit{erx}) = \operatorname{real\text{-}part}(f(\mathit{erx}))) \Rightarrow \\ \quad \quad \operatorname{extends\text{-}on}(s, a, f) \end{array}}}
\vnbprop{cc-magnitude-sq-conjugate}{\vnbstmt{\forall z \in \mathbb{C}.\; ((\operatorname{magnitude}(z) \cdot \operatorname{magnitude}(z)) = (z \cdot \operatorname{conjugate}(z)))}}
\vnbprop{cc-magnitude-conjugate}{\vnbstmt{\forall z \in \mathbb{C}.\; (\operatorname{magnitude}(\operatorname{conjugate}(z)) = \operatorname{magnitude}(z))}}
\vnbprop{rb3-cf-bound}{\vnbstmt{\begin{array}{@{}l@{}} \forall m.; \forall n \in (\operatorname{vec}(m) \to \mathbb{R}), a \in (\operatorname{vec}(m) \to \mathbb{C}), p.; \forall y \in \operatorname{vec}(m). \\ \quad \operatorname{is\text{-}module}(m) \wedge \\ \quad \operatorname{scal}(m) = \operatorname{normed\text{-}field\text{-}as\text{-}commutative\text{-}ring}(\operatorname{cc\text{-}normed\text{-}field}) \wedge \\ \quad \forall x, \mathit{ray} \in \operatorname{vec}(m).\; (n(\operatorname{vadd}(m)(x, \mathit{ray})) = (n(x) + n(\mathit{ray}))) \wedge \\ \quad \forall r \in \mathbb{R}, \mathit{rhx} \in \operatorname{vec}(m).\; (n(\operatorname{act}(m)(r, \mathit{rhx})) = (r \cdot n(\mathit{rhx}))) \wedge \\ \quad \forall \mathit{cvy} \in \operatorname{vec}(m).\; (a(\mathit{cvy}) \simeq (n(\mathit{cvy}) + (-n(\operatorname{act}(m)(+i, \mathit{cvy})) \cdot +i))) \wedge \\ \quad \operatorname{is\text{-}seminorm}(\operatorname{cc\text{-}normed\text{-}field}, m, p) \wedge \\ \quad \forall \mathit{bx} \in \operatorname{vec}(m).\; (n(\mathit{bx}) \leq p(\mathit{bx})) \Rightarrow \\ \quad \quad \operatorname{fnrm}(\operatorname{cc\text{-}normed\text{-}field})(a(y)) \leq p(y) \end{array}}}
\vnbprop{hahn-banach-seminorm}{\vnbstmt{\begin{array}{@{}l@{}} \forall m, p, s, f, d. \\ \quad \operatorname{is\text{-}module}(m) \wedge \\ \quad \operatorname{is\text{-}normed\text{-}field}(d) \wedge \\ \quad (d = \operatorname{rr\text{-}normed\text{-}field}) \vee (d = \operatorname{cc\text{-}normed\text{-}field}) \wedge \\ \quad \operatorname{scal}(m) = \operatorname{normed\text{-}field\text{-}as\text{-}commutative\text{-}ring}(d) \wedge \\ \quad \operatorname{is\text{-}seminorm}(d, m, p) \wedge \\ \quad \operatorname{is\text{-}submodule}(m, s) \wedge \\ \quad \operatorname{is\text{-}k\text{-}linear\text{-}on}(m, s, f) \wedge \\ \quad \forall x \in s.\; (\operatorname{fnrm}(d)(f(x)) \leq p(x)) \Rightarrow \\ \quad \quad \exists a. \\ \quad \quad \quad (\operatorname{is\text{-}k\text{-}linear}(m, a) \wedge (\operatorname{extends\text{-}on}(s, a, f) \wedge \forall x \in \operatorname{vec}(m).\; (\operatorname{fnrm}(d)(a(x)) \leq p(x)))) \end{array}}}
\vnbprop{rb3-pgt-pts}{\vnbstmt{\forall m.; \forall g \in \mathbf{Set}.\; (\operatorname{pts}(\operatorname{pseudo\text{-}gauge\text{-}top}(m, g)) \simeq g)}}
\vnbprop{rb3-pgt-opens}{\vnbstmt{\begin{array}{@{}l@{}} \forall m.; \forall g \in \mathbf{Set}. \\ \quad (\operatorname{opens}(\operatorname{pseudo\text{-}gauge\text{-}top}(m, g)) \simeq \{u \in \operatorname{power}(g) : \operatorname{is\text{-}open\text{-}gauge}(m, g, u)\}) \end{array}}}
\vnbprop{rb3-fam-dist-real}{\vnbstmt{\begin{array}{@{}l@{}} \forall m, g.; \forall n \in \mathbb{N}, y, x \in g. \\ \quad \operatorname{is\text{-}countable\text{-}pseudometric\text{-}family}(m, g) \Rightarrow (\operatorname{dist}(m(n))(y, x) \in \mathbb{R}) \end{array}}}
\vnbprop{rb3-fam-pseudo}{\vnbstmt{\begin{array}{@{}l@{}} \forall m, g.; \forall n \in \mathbb{N}. \\ \quad \operatorname{is\text{-}countable\text{-}pseudometric\text{-}family}(m, g) \Rightarrow \operatorname{is\text{-}pseudometric}(\operatorname{dist}(m(n)), g) \end{array}}}
\vnbprop{rb3-og-inter}{\vnbstmt{\begin{array}{@{}l@{}} \forall m, g, u, v. \\ \quad (\operatorname{is\text{-}countable\text{-}pseudometric\text{-}family}(m, g) \wedge (\operatorname{is\text{-}open\text{-}gauge}(m, g, u) \wedge \operatorname{is\text{-}open\text{-}gauge}(m, g, v))) \Rightarrow  \\ \quad \quad \operatorname{is\text{-}open\text{-}gauge}(m, g, (u \cap v)) \end{array}}}
\vnbprop{rb3-og-union}{\vnbstmt{\begin{array}{@{}l@{}} \forall m, g, w. \\ \quad \operatorname{is\text{-}countable\text{-}pseudometric\text{-}family}(m, g) \wedge \\ \quad \forall a \in w.\; \operatorname{is\text{-}open\text{-}gauge}(m, g, a) \wedge \\ \quad \forall a \in w.\; (a \subseteq g) \Rightarrow \\ \quad \quad \operatorname{is\text{-}open\text{-}gauge}(m, g, \operatorname{big\text{-}union}(u, w, u)) \end{array}}}
\vnbprop{rb3-og-empty}{\vnbstmt{\begin{array}{@{}l@{}} \forall m, g. \\ \quad \operatorname{is\text{-}countable\text{-}pseudometric\text{-}family}(m, g) \Rightarrow \operatorname{is\text{-}open\text{-}gauge}(m, g, \emptyset) \end{array}}}
\vnbprop{rb3-og-full}{\vnbstmt{\begin{array}{@{}l@{}} \forall m, g. \\ \quad \operatorname{is\text{-}countable\text{-}pseudometric\text{-}family}(m, g) \Rightarrow \operatorname{is\text{-}open\text{-}gauge}(m, g, g) \end{array}}}
\vnbprop{rb3-gauge-top-is-top-space}{\vnbstmt{\begin{array}{@{}l@{}} \forall m, g. \\ \quad (\operatorname{is\text{-}countable\text{-}pseudometric\text{-}family}(m, g) \wedge (g \in \mathbf{Set})) \Rightarrow  \\ \quad \quad \operatorname{is\text{-}top\text{-}space}(\operatorname{pseudo\text{-}gauge\text{-}top}(m, g)) \end{array}}}
\vnbprop{rb3-og-zero-dist}{\vnbstmt{\begin{array}{@{}l@{}} \forall m, g, u.; \forall y \in u, x \in g. \\ \quad \operatorname{is\text{-}countable\text{-}pseudometric\text{-}family}(m, g) \wedge \\ \quad \operatorname{is\text{-}open\text{-}gauge}(m, g, u) \wedge \\ \quad \forall n \in \mathbb{N}.\; (\operatorname{dist}(m(n))(y, x) = 0) \Rightarrow \\ \quad \quad x \in u \end{array}}}
\vnbprop{rb3-metrizable-hausdorff}{\vnbstmt{\forall s.\; (\operatorname{is\text{-}metrizable\text{-}top\text{-}space}(s) \Rightarrow \operatorname{is\text{-}hausdorff}(s))}}
\vnbprop{rb3-zero-family}{\vnbstmt{\begin{array}{@{}l@{}} \operatorname{is\text{-}countable\text{-}pseudometric\text{-}family}((\lambda n \in \mathbb{N}.\; {[}\mathbb{N}, (\lambda {[}\operatorname{za\_}, \operatorname{zb\_}] \in (\mathbb{N} \times \mathbb{N}).\; 0)]), \\ \quad \mathbb{N}) \end{array}}}
\vnbprop{rb3-zero-gauge-witness}{\vnbstmt{\begin{array}{@{}l@{}} \operatorname{is\text{-}gauge\text{-}countable}(\operatorname{pseudo\text{-}gauge\text{-}top}((\lambda n \in \mathbb{N}.\; {[}\mathbb{N}, (\lambda {[}\operatorname{za\_}, \operatorname{zb\_}] \in (\mathbb{N} \times \mathbb{N}).\; 0)]), \mathbb{N})) \wedge \\ \neg \operatorname{is\text{-}metrizable\text{-}top\text{-}space}(\operatorname{pseudo\text{-}gauge\text{-}top}((\lambda n \in \mathbb{N}.\; {[}\mathbb{N}, (\lambda {[}\operatorname{za\_}, \operatorname{zb\_}] \in (\mathbb{N} \times \mathbb{N}).\; 0)]), \mathbb{N})) \end{array}}}
\vnbprop{gauge-countable-not-metrizable-witness}{\vnbstmt{\exists s.\; (\operatorname{is\text{-}gauge\text{-}countable}(s) \wedge \neg \operatorname{is\text{-}metrizable\text{-}top\text{-}space}(s))}}
\vnbprop{rb3-tuple2}{\vnbstmt{\begin{array}{@{}l@{}} \forall a, b. \\ \quad ((a \in \mathbf{Set}) \wedge (b \in \mathbf{Set})) \Rightarrow ((\operatorname{pts}({[}a, b]) \simeq a) \wedge (\operatorname{dist}({[}a, b]) \simeq b)) \end{array}}}
\vnbprop{rb3-pms-of}{\vnbstmt{\begin{array}{@{}l@{}} \forall q, \mathit{dq}. \\ \quad ((q \in \mathbf{Set}) \wedge ((\mathit{dq} \in ((q \times q) \to \mathbb{R})) \wedge \operatorname{is\text{-}pseudometric}(\mathit{dq}, q))) \Rightarrow  \\ \quad \quad \operatorname{is\text{-}pseudometric\text{-}space}({[}q, \mathit{dq}]) \end{array}}}
\vnbprop{rb3-metric-pseudo}{\vnbstmt{\forall d.\; (\operatorname{is\text{-}metric\text{-}space}(d) \Rightarrow \operatorname{is\text{-}pseudometric}(\operatorname{dist}(d), \operatorname{pts}(d)))}}
\vnbprop{rb3-fam1-dist}{\vnbstmt{\begin{array}{@{}l@{}} \forall d.; \forall n \in \mathbb{N}. \\ \quad \operatorname{is\text{-}metric\text{-}space}(d) \Rightarrow  \\ \quad \quad (\operatorname{dist}((\lambda \mathit{zn} \in \mathbb{N}.\; {[}\operatorname{pts}(d), \operatorname{dist}(d)])(n)) \simeq \operatorname{dist}(d)) \end{array}}}
\vnbprop{rb3-metric-gauge-iff}{\vnbstmt{\begin{array}{@{}l@{}} \forall d, u. \\ \quad (\operatorname{is\text{-}metric\text{-}space}(d) \wedge (u \subseteq \operatorname{pts}(d))) \Rightarrow  \\ \quad \quad \operatorname{is\text{-}open}(d, u) \Leftrightarrow \operatorname{is\text{-}open\text{-}gauge}((\lambda n \in \mathbb{N}.\; {[}\operatorname{pts}(d), \operatorname{dist}(d)]), \operatorname{pts}(d), u) \end{array}}}
\vnbprop{rb3-mgc-fwd}{\vnbstmt{\begin{array}{@{}l@{}} \forall s. \\ \quad (\operatorname{is\text{-}top\text{-}space}(s) \wedge \operatorname{is\text{-}metrizable\text{-}top\text{-}space}(s)) \Rightarrow  \\ \quad \quad \operatorname{is\text{-}hausdorff}(s) \wedge \operatorname{is\text{-}gauge\text{-}countable}(s) \end{array}}}
\vnbprop{rb3-cm-term-le}{\vnbstmt{\begin{array}{@{}l@{}} \forall s, w, x, y.; \forall n \in \mathbb{N}. \\ \quad (\operatorname{is\text{-}c\text{-}metric\text{-}space}(s) \wedge (\operatorname{summable\text{-}weight}(w) \wedge ((x \in \operatorname{pts}(s)) \wedge (y \in \operatorname{pts}(s))))) \Rightarrow  \\ \quad \quad (w(n) \cdot \operatorname{min}(1, \operatorname{dists}(s)(n)(x, y))) \leq \operatorname{dist}(\operatorname{c\text{-}metric\text{-}w}(s, w))(x, y) \end{array}}}
\vnbprop{rb3-cm-dist-le}{\vnbstmt{\begin{array}{@{}l@{}} \forall s, w, e. \\ \quad \operatorname{is\text{-}c\text{-}metric\text{-}space}(s) \wedge \\ \quad \operatorname{summable\text{-}weight}(w) \wedge \\ \quad \forall k \in \mathbb{N}.\; (\operatorname{series\text{-}partial\text{-}sum}(w, k) \leq 1) \wedge \\ \quad \operatorname{pos\text{-}rr}(e) \Rightarrow \\ \quad \quad \exists d \in \mathbb{N}. \forall x, y \in \operatorname{pts}(s), r \in \mathbb{R}. \\ \quad \quad \quad ((0 < r) \wedge \forall n \in \mathbb{N}.\; ((n \leq d) \Rightarrow (\operatorname{dists}(s)(n)(x, y) < r))) \Rightarrow  \\ \quad \quad \quad \quad \operatorname{dist}(\operatorname{c\text{-}metric\text{-}w}(s, w))(x, y) \leq (r + e) \end{array}}}
\vnbprop{rb3-tuple2d}{\vnbstmt{\begin{array}{@{}l@{}} \forall a, b. \\ \quad ((a \in \mathbf{Set}) \wedge (b \in \mathbf{Set})) \Rightarrow ((\operatorname{pts}({[}a, b]) \simeq a) \wedge (\operatorname{dists}({[}a, b]) \simeq b)) \end{array}}}
\vnbprop{rb3-cs-is-cms}{\vnbstmt{\begin{array}{@{}l@{}} \forall m, g. \\ \quad \operatorname{is\text{-}countable\text{-}pseudometric\text{-}family}(m, g) \wedge \\ \quad g \in \mathbf{Set} \wedge \\ \quad \forall u, v \in g. \\ \quad \quad \forall k \in \mathbb{N}.\; (\operatorname{dist}(m(k))(u, v) = 0) \Rightarrow (u = v) \Rightarrow \\ \quad \quad \operatorname{is\text{-}c\text{-}metric\text{-}space}({[}g, (\lambda \mathit{zk} \in \mathbb{N}.\; \operatorname{dist}(m(\mathit{zk})))]) \end{array}}}
\vnbprop{rb3-haus-sep}{\vnbstmt{\begin{array}{@{}l@{}} \forall s, m, g.; \forall u, v \in g. \\ \quad \operatorname{is\text{-}hausdorff}(s) \wedge \\ \quad \operatorname{is\text{-}countable\text{-}pseudometric\text{-}family}(m, g) \wedge \\ \quad (\operatorname{pseudo\text{-}gauge\text{-}top}(m, g) \simeq s) \wedge \\ \quad \forall k \in \mathbb{N}.\; (\operatorname{dist}(m(k))(u, v) = 0) \Rightarrow \\ \quad \quad u = v \end{array}}}
\vnbprop{rb3-cm-open-iff}{\vnbstmt{\begin{array}{@{}l@{}} \forall s, w, m, g, u. \\ \quad \operatorname{is\text{-}c\text{-}metric\text{-}space}(s) \wedge \\ \quad \operatorname{summable\text{-}weight}(w) \wedge \\ \quad \forall k \in \mathbb{N}.\; (\operatorname{series\text{-}partial\text{-}sum}(w, k) \leq 1) \wedge \\ \quad \forall n, k \in \mathbb{N}.\; ((n \leq k) \Rightarrow (w(k) \leq w(n))) \wedge \\ \quad (\operatorname{pts}(s) \simeq g) \wedge \\ \quad \forall n \in \mathbb{N}.\; (\operatorname{dists}(s)(n) \simeq \operatorname{dist}(m(n))) \wedge \\ \quad u \subseteq g \Rightarrow \\ \quad \quad \operatorname{is\text{-}open}(\operatorname{c\text{-}metric\text{-}w}(s, w), u) \Leftrightarrow \operatorname{is\text{-}open\text{-}gauge}(m, g, u) \end{array}}}
\vnbprop{rb3-dyadic-le-one}{\vnbstmt{\begin{array}{@{}l@{}} \forall k \in \mathbb{N}. \\ \quad \operatorname{series\text{-}partial\text{-}sum}((\lambda n \in \mathbb{N}.\; /(1, \operatorname{power}(2, (n + 1)))), k) \leq 1 \end{array}}}
\vnbprop{rb3-dyadic-antitone}{\vnbstmt{\begin{array}{@{}l@{}} \forall a, k \in \mathbb{N}. \\ \quad (a \leq k) \Rightarrow  \\ \quad \quad (\lambda n \in \mathbb{N}.\; /(1, \operatorname{power}(2, (n + 1))))(k) \leq (\lambda n \in \mathbb{N}.\; /(1, \operatorname{power}(2, (n + 1))))(a) \end{array}}}
\vnbprop{rb3-mgc-bwd}{\vnbstmt{\begin{array}{@{}l@{}} \forall s, m, g. \\ \quad \operatorname{is\text{-}top\text{-}space}(s) \wedge \\ \quad \operatorname{is\text{-}hausdorff}(s) \wedge \\ \quad \operatorname{is\text{-}countable\text{-}pseudometric\text{-}family}(m, g) \wedge \\ \quad (\operatorname{pseudo\text{-}gauge\text{-}top}(m, g) \simeq s) \Rightarrow \\ \quad \quad \operatorname{is\text{-}metrizable\text{-}top\text{-}space}(s) \end{array}}}
\vnbprop{metrizable-iff-gauge-countable}{\vnbstmt{\begin{array}{@{}l@{}} \forall s. \\ \quad \operatorname{is\text{-}top\text{-}space}(s) \Rightarrow  \\ \quad \quad \operatorname{is\text{-}metrizable\text{-}top\text{-}space}(s) \Leftrightarrow (\operatorname{is\text{-}hausdorff}(s) \wedge \operatorname{is\text{-}gauge\text{-}countable}(s)) \end{array}}}
\vnbprop{bounded-metrizable-iff-gauge-countable}{\vnbstmt{\begin{array}{@{}l@{}} \forall s. \\ \quad \operatorname{is\text{-}top\text{-}space}(s) \Rightarrow  \\ \quad \quad \exists d.\; (\operatorname{is\text{-}bounded\text{-}metric\text{-}space}(d) \wedge (\operatorname{metric\text{-}top}(d) \simeq s)) \\ \quad \quad \Leftrightarrow\; \operatorname{is\text{-}hausdorff}(s) \wedge \operatorname{is\text{-}gauge\text{-}countable}(s) \end{array}}}

\vnbbookfile{theorem-library/cc-norming-functional}{theorem-library/cc-norming-functional}
\noindent{\small cc-norming-functional.scm -- wave of 2026-11-09/10 (agent BA-1D), wired 2026-11-10 from scratchpad/ba1d-1109/cc-norming-functional.scm. Library: analysis.}\par
\vnbprop{cnf-ring-neg}{\vnbstmt{(\operatorname{neg}(\operatorname{normed\text{-}field\text{-}as\text{-}commutative\text{-}ring}(\operatorname{cc\text{-}normed\text{-}field})) \simeq \operatorname{neg}(\operatorname{cc\text{-}normed\text{-}field}))}}
\vnbprop{cnf-c-neg-eq}{\vnbstmt{\begin{array}{@{}l@{}} \forall m.; \forall a \in \mathbb{C}. \\ \quad (\operatorname{is\text{-}module}(m) \wedge (\operatorname{scal}(m) = \operatorname{normed\text{-}field\text{-}as\text{-}commutative\text{-}ring}(\operatorname{cc\text{-}normed\text{-}field}))) \Rightarrow  \\ \quad \quad \operatorname{neg}(\operatorname{scal}(m))(a) = -a \end{array}}}
\vnbprop{cnf-coef-unique}{\vnbstmt{\begin{array}{@{}l@{}} \forall m, p.; \forall y \in \operatorname{vec}(m), a, b \in \mathbb{C}. \\ \quad (\operatorname{is\text{-}seminorm}(\operatorname{cc\text{-}normed\text{-}field}, m, p) \wedge (\operatorname{act}(m)(a, y) = \operatorname{act}(m)(b, y))) \Rightarrow  \\ \quad \quad (a \cdot p(y)) = (b \cdot p(y)) \end{array}}}
\vnbprop{cnf-coef-value}{\vnbstmt{\begin{array}{@{}l@{}} \forall m, p, y, u, c. \\ \quad (\operatorname{is\text{-}seminorm}(\operatorname{cc\text{-}normed\text{-}field}, m, p) \wedge ((y \in \operatorname{vec}(m)) \wedge ((c \in \mathbb{C}) \wedge (u = \operatorname{act}(m)(c, y))))) \Rightarrow  \\ \quad \quad (\iota\, t.\;((t \in \mathbb{C}) \wedge \exists q \in \mathbb{C}.\; ((u = \operatorname{act}(m)(q, y)) \wedge (t = (q \cdot p(y)))))) = (c \cdot p(y)) \end{array}}}
\vnbprop{cnf-line-mem}{\vnbstmt{\begin{array}{@{}l@{}} \forall m, p, y, u, c. \\ \quad (\operatorname{is\text{-}seminorm}(\operatorname{cc\text{-}normed\text{-}field}, m, p) \wedge ((y \in \operatorname{vec}(m)) \wedge ((c \in \mathbb{C}) \wedge (u = \operatorname{act}(m)(c, y))))) \Rightarrow  \\ \quad \quad u \in \{w \in \operatorname{vec}(m) : \exists r \in \mathbb{C}.\; (w = \operatorname{act}(m)(r, y))\} \end{array}}}
\vnbprop{cnf-line-coef}{\vnbstmt{\begin{array}{@{}l@{}} \forall m, y.; \forall u \in \{w \in \operatorname{vec}(m) : \exists r \in \mathbb{C}.\; (w = \operatorname{act}(m)(r, y))\}. \\ \quad (u \in \operatorname{vec}(m)) \wedge \exists r \in \mathbb{C}.\; (u = \operatorname{act}(m)(r, y)) \end{array}}}
\vnbprop{cnf-line-submodule}{\vnbstmt{\begin{array}{@{}l@{}} \forall m, p.; \forall y \in \operatorname{vec}(m). \\ \quad \operatorname{is\text{-}seminorm}(\operatorname{cc\text{-}normed\text{-}field}, m, p) \Rightarrow  \\ \quad \quad \operatorname{is\text{-}submodule}(m, \{w \in \operatorname{vec}(m) : \exists r \in \mathbb{C}.\; (w = \operatorname{act}(m)(r, y))\}) \end{array}}}
\vnbprop{cnf-seed-linear}{\vnbstmt{\begin{array}{@{}l@{}} \forall m, p.; \forall y \in \operatorname{vec}(m). \\ \quad \operatorname{is\text{-}seminorm}(\operatorname{cc\text{-}normed\text{-}field}, m, p) \Rightarrow  \\ \quad \quad \operatorname{is\text{-}k\text{-}linear\text{-}on}(m, \{w \in \operatorname{vec}(m) : \exists r \in \mathbb{C}.\; (w = \operatorname{act}(m)(r, y))\}, \\ \quad \quad \quad (\lambda u \in \{w \in \operatorname{vec}(m) : \exists r \in \mathbb{C}.\; (w = \operatorname{act}(m)(r, y))\}.\; \\ \quad \quad \quad \quad (\iota\, t.\;((t \in \mathbb{C}) \wedge \exists q \in \mathbb{C}.\; ((u = \operatorname{act}(m)(q, y)) \wedge (t = (q \cdot p(y)))))))) \end{array}}}
\vnbprop{cnf-seed-bound}{\vnbstmt{\begin{array}{@{}l@{}} \forall m, p.; \forall y \in \operatorname{vec}(m), x \in \{w \in \operatorname{vec}(m) : \exists r \in \mathbb{C}.\; (w = \operatorname{act}(m)(r, y))\}. \\ \quad \operatorname{is\text{-}seminorm}(\operatorname{cc\text{-}normed\text{-}field}, m, p) \Rightarrow  \\ \quad \quad \operatorname{fnrm}(\operatorname{cc\text{-}normed\text{-}field})((\lambda u \in \{w \in \operatorname{vec}(m) : \exists r \in \mathbb{C}.\; (w = \operatorname{act}(m)(r, y))\}.\; (\iota\, t.\;((t \in \mathbb{C}) \wedge \exists q \in \mathbb{C}.\; ((u = \operatorname{act}(m)(q, y)) \wedge (t = (q \cdot p(y)))))))(x)) \\ \quad \quad \leq\; p(x) \end{array}}}
\vnbprop{cc-norming-functional}{\vnbstmt{\begin{array}{@{}l@{}} \forall m, p.; \forall y \in \operatorname{vec}(m). \\ \quad \operatorname{is\text{-}seminorm}(\operatorname{cc\text{-}normed\text{-}field}, m, p) \Rightarrow  \\ \quad \quad \exists f. \\ \quad \quad \quad (\operatorname{is\text{-}k\text{-}linear}(m, f) \wedge (\forall x \in \operatorname{vec}(m).\; (\operatorname{magnitude}(f(x)) \leq p(x)) \wedge (f(y) = p(y)))) \end{array}}}

\vnbbookfile{theorem-library/inclusion-exclusion}{theorem-library/inclusion-exclusion}
\noindent{\small theorem-library/inclusion-exclusion.scm -- THE INCLUSION-EXCLUSION PRINCIPLE (ALGEBRAIC-NUMBERS notes, Prop. 1.20 and Cor. 1.21; ALGNUM, bundle B2, 2026-10-26). inclusion-exclusion          PROP. 1.20: for a finite index set I, a finite set X and sets A\_k (k in I): card\{x in X : x in no A\_k\} = sum over F in POWER(I) of ...}\par
\vnbprop{inclusion-exclusion-base}{\vnbstmt{\begin{array}{@{}l@{}} \forall f.; \forall s \in \mathbf{Set}. \\ \quad (|s| \in \mathbb{N}) \Rightarrow  \\ \quad \quad |\{x \in s : \forall k \in \emptyset.\; \neg (x \in f(k))\}| \\ \quad \quad =\; \operatorname{finsum}(\operatorname{ring\text{-}additive\text{-}ag}(\operatorname{normed\text{-}field\text{-}as\text{-}commutative\text{-}ring}(\operatorname{cc\text{-}normed\text{-}field})), \\ \quad \quad \quad (\lambda \mathit{ief} \in \operatorname{power}(\emptyset).\; (\operatorname{power}(-1, |\mathit{ief}|) \cdot |\{x \in s : \forall k \in \mathit{ief}.\; (x \in f(k))\}|)), \\ \quad \quad \quad \operatorname{power}(\emptyset)) \end{array}}}
\vnbprop{inclusion-exclusion-split-card}{\vnbstmt{\begin{array}{@{}l@{}} \forall f, a, s, v. \\ \quad ((a \in \mathbf{Set}) \wedge ((|a| \in \mathbb{N}) \wedge (v \in \mathbf{Set}))) \Rightarrow  \\ \quad \quad |\{x \in a : \forall k \in s.\; \neg (x \in f(k))\}| \\ \quad \quad =\; |\{x \in a : \forall k \in (s \cup \{v\}).\; \neg (x \in f(k))\}| + |\{x \in \{y \in a : (y \in f(v))\} : \forall k \in s.\; \neg (x \in f(k))\}| \end{array}}}
\vnbprop{inclusion-exclusion-split-sum}{\vnbstmt{\begin{array}{@{}l@{}} \forall f, s, i, v. \\ \quad ((s \in \mathbf{Set}) \wedge ((|s| \in \mathbb{N}) \wedge ((i \in \mathbf{Set}) \wedge ((|i| \in \mathbb{N}) \wedge ((v \in \mathbf{Set}) \wedge \neg (v \in i)))))) \Rightarrow  \\ \quad \quad \operatorname{finsum}(\operatorname{ring\text{-}additive\text{-}ag}(\operatorname{normed\text{-}field\text{-}as\text{-}commutative\text{-}ring}(\operatorname{cc\text{-}normed\text{-}field})), \\ \quad \quad \quad (\lambda \mathit{ief} \in \operatorname{power}((i \cup \{v\})).\; \\ \quad \quad \quad \quad \operatorname{power}(-1, |\mathit{ief}|) \cdot |\{x \in s : \forall k \in \mathit{ief}.\; (x \in f(k))\}|), \\ \quad \quad \quad \operatorname{power}((i \cup \{v\}))) \\ \quad \quad =\; (\operatorname{finsum}(\operatorname{ring\text{-}additive\text{-}ag}(\operatorname{normed\text{-}field\text{-}as\text{-}commutative\text{-}ring}(\operatorname{cc\text{-}normed\text{-}field})), (\lambda \mathit{ief} \in \operatorname{power}(i).\; (\operatorname{power}(-1, |\mathit{ief}|) \cdot |\{x \in s : \forall k \in \mathit{ief}.\; (x \in f(k))\}|)), \operatorname{power}(i)) - \operatorname{finsum}(\operatorname{ring\text{-}additive\text{-}ag}(\operatorname{normed\text{-}field\text{-}as\text{-}commutative\text{-}ring}(\operatorname{cc\text{-}normed\text{-}field})), (\lambda \mathit{ief} \in \operatorname{power}(i).\; (\operatorname{power}(-1, |\mathit{ief}|) \cdot |\{x \in \{y \in s : (y \in f(v))\} : \forall k \in \mathit{ief}.\; (x \in f(k))\}|)), \operatorname{power}(i))) \end{array}}}
\vnbprop{inclusion-exclusion}{\vnbstmt{\begin{array}{@{}l@{}} \forall i \in \mathbf{Set}, f.; \forall s \in \mathbf{Set}. \\ \quad ((|i| \in \mathbb{N}) \wedge (|s| \in \mathbb{N})) \Rightarrow  \\ \quad \quad |\{x \in s : \forall k \in i.\; \neg (x \in f(k))\}| \\ \quad \quad =\; \operatorname{finsum}(\operatorname{ring\text{-}additive\text{-}ag}(\operatorname{normed\text{-}field\text{-}as\text{-}commutative\text{-}ring}(\operatorname{cc\text{-}normed\text{-}field})), \\ \quad \quad \quad (\lambda \mathit{ief} \in \operatorname{power}(i).\; (\operatorname{power}(-1, |\mathit{ief}|) \cdot |\{x \in s : \forall k \in \mathit{ief}.\; (x \in f(k))\}|)), \\ \quad \quad \quad \operatorname{power}(i)) \end{array}}}
\vnbprop{inclusion-exclusion-union}{\vnbstmt{\begin{array}{@{}l@{}} \forall i \in \mathbf{Set}, f.; \forall s \in \mathbf{Set}. \\ \quad ((|i| \in \mathbb{N}) \wedge (|s| \in \mathbb{N})) \Rightarrow  \\ \quad \quad |\{x \in s : \exists k \in i.\; (x \in f(k))\}| \\ \quad \quad =\; \operatorname{finsum}(\operatorname{ring\text{-}additive\text{-}ag}(\operatorname{normed\text{-}field\text{-}as\text{-}commutative\text{-}ring}(\operatorname{cc\text{-}normed\text{-}field})), \\ \quad \quad \quad (\lambda \mathit{ief} \in \operatorname{difference}(\operatorname{power}(i), \{\emptyset\}).\; \\ \quad \quad \quad \quad \operatorname{power}(-1, \operatorname{nn\text{-}minus}(|\mathit{ief}|, 1)) \cdot |\{x \in s : \forall k \in \mathit{ief}.\; (x \in f(k))\}|), \\ \quad \quad \quad \operatorname{difference}(\operatorname{power}(i), \{\emptyset\})) \end{array}}}

\vnbbookfile{theorem-library/euler-phi-product}{theorem-library/euler-phi-product}
\noindent{\small theorem-library/euler-phi-product.scm -- PROP. 1.19 (EULER'S PRODUCT FORMULA) (ALGEBRAIC-NUMBERS notes, ch. 1; ALGNUM, 2026-10-27).  phi(n) = n prod\_\{p | n\} (1 - 1/p), by the notes' route through inclusion-exclusion (1.20). prime-nat-divides-prime, prime-set-product      P\_F = prod\_\{p in F\} p for a finite set F of ...}\par
\vnbprop{prime-nat-divides-prime}{\vnbstmt{\begin{array}{@{}l@{}} \forall q, p. \\ \quad (\operatorname{is\text{-}prime\text{-}nat}(q) \wedge (\operatorname{is\text{-}prime\text{-}nat}(p) \wedge \operatorname{zz\text{-}divides}(q, p))) \Rightarrow (q = p) \end{array}}}
\vnbprop{prime-set-product}{\vnbstmt{\begin{array}{@{}l@{}} \forall f \in \mathbf{Set}. \\ \quad ((|f| \in \mathbb{N}) \wedge \forall p \in f.\; \operatorname{is\text{-}prime\text{-}nat}(p)) \Rightarrow  \\ \quad \quad \operatorname{prod\text{-}ring}(\operatorname{zz\text{-}ring}, (\lambda \mathit{nzp} \in \mathbb{Z}.\; \mathit{nzp}), f) \in \mathbb{N} \wedge \\ \quad \quad 1 \leq \operatorname{prod\text{-}ring}(\operatorname{zz\text{-}ring}, (\lambda \mathit{nzp} \in \mathbb{Z}.\; \mathit{nzp}), f) \wedge \\ \quad \quad \forall p \in f.\; \operatorname{zz\text{-}divides}(p, \operatorname{prod\text{-}ring}(\operatorname{zz\text{-}ring}, (\lambda \mathit{nzp} \in \mathbb{Z}.\; \mathit{nzp}), f)) \wedge \\ \quad \quad \forall q \in \mathbb{N}. \\ \quad \quad \quad (\operatorname{is\text{-}prime\text{-}nat}(q) \wedge \operatorname{zz\text{-}divides}(q, \operatorname{prod\text{-}ring}(\operatorname{zz\text{-}ring}, (\lambda \mathit{nzp} \in \mathbb{Z}.\; \mathit{nzp}), f))) \Rightarrow  \\ \quad \quad \quad \quad q \in f \wedge \\ \quad \quad \forall k \in \mathbb{Z}. \\ \quad \quad \quad \forall p \in f.\; \operatorname{zz\text{-}divides}(p, k) \Rightarrow  \\ \quad \quad \quad \quad \operatorname{zz\text{-}divides}(\operatorname{prod\text{-}ring}(\operatorname{zz\text{-}ring}, (\lambda \mathit{nzp} \in \mathbb{Z}.\; \mathit{nzp}), f), k) \end{array}}}
\vnbprop{prod-ring-cc-ones}{\vnbstmt{\begin{array}{@{}l@{}} \forall f \in \mathbf{Set}. \\ \quad ((|f| \in \mathbb{N}) \wedge \forall z \in f.\; (z \in \mathbb{N})) \Rightarrow  \\ \quad \quad \operatorname{prod\text{-}ring}(\operatorname{normed\text{-}field\text{-}as\text{-}commutative\text{-}ring}(\operatorname{cc\text{-}normed\text{-}field}), (\lambda o \in \mathbb{N}.\; 1), f) = 1 \end{array}}}
\vnbprop{prime-set-recip-product}{\vnbstmt{\begin{array}{@{}l@{}} \forall f \in \mathbf{Set}. \\ \quad ((|f| \in \mathbb{N}) \wedge \forall p \in f.\; \operatorname{is\text{-}prime\text{-}nat}(p)) \Rightarrow  \\ \quad \quad \operatorname{prod\text{-}ring}(\operatorname{normed\text{-}field\text{-}as\text{-}commutative\text{-}ring}(\operatorname{cc\text{-}normed\text{-}field}), (\lambda a \in \mathbb{N}.\; -\operatorname{recip}(a)), f) \\ \quad \quad =\; \operatorname{power}(-1, |f|) \cdot \operatorname{recip}(\operatorname{prod\text{-}ring}(\operatorname{zz\text{-}ring}, (\lambda \mathit{nzp} \in \mathbb{Z}.\; \mathit{nzp}), f)) \end{array}}}
\vnbprop{prime-set-multiples-count}{\vnbstmt{\begin{array}{@{}l@{}} \forall n \in \mathbb{N}, f \in \mathbf{Set}. \\ \quad 1 \leq n \wedge \\ \quad |f| \in \mathbb{N} \wedge \\ \quad \forall p \in f.\; \operatorname{is\text{-}prime\text{-}nat}(p) \wedge \\ \quad \forall p \in f.\; \operatorname{zz\text{-}divides}(p, n) \Rightarrow \\ \quad \quad \operatorname{prod\text{-}ring}(\operatorname{zz\text{-}ring}, (\lambda \mathit{nzp} \in \mathbb{Z}.\; \mathit{nzp}), f) \cdot |\{x \in \operatorname{interval}(1, n) : \forall k \in f.\; (x \in (\lambda \mathit{nzf} \in \mathbb{N}.\; \{m \in \mathbb{Z} : \operatorname{zz\text{-}divides}(\mathit{nzf}, m)\})(k))\}| \\ \quad \quad =\; n \end{array}}}
\vnbprop{euler-phi-coprime-out}{\vnbstmt{\begin{array}{@{}l@{}} \forall n \in \mathbb{N}. \\ \quad (1 \leq n) \Rightarrow  \\ \quad \quad \{k \in \operatorname{interval}(1, n) : \operatorname{zz\text{-}coprime}(k, n)\} \\ \quad \quad =\; \{x \in \operatorname{interval}(1, n) : \forall \mathit{iek} \in \{d \in \operatorname{interval}(1, n) : (\operatorname{is\text{-}prime\text{-}nat}(d) \wedge \operatorname{zz\text{-}divides}(d, n))\}.\; \neg (x \in (\lambda f \in \mathbb{N}.\; \{m \in \mathbb{Z} : \operatorname{zz\text{-}divides}(f, m)\})(\mathit{iek}))\} \end{array}}}
\vnbprop{euler-phi-product}{\vnbstmt{\begin{array}{@{}l@{}} \forall n \in \mathbb{N}. \\ \quad (1 \leq n) \Rightarrow  \\ \quad \quad \operatorname{euler\text{-}phi}(n) \\ \quad \quad =\; n \cdot \operatorname{prod\text{-}ring}(\operatorname{normed\text{-}field\text{-}as\text{-}commutative\text{-}ring}(\operatorname{cc\text{-}normed\text{-}field}), (\lambda q \in \{d \in \operatorname{interval}(1, n) : (\operatorname{is\text{-}prime\text{-}nat}(d) \wedge \operatorname{zz\text{-}divides}(d, n))\}.\; (1 - \operatorname{recip}(q))), \{d \in \operatorname{interval}(1, n) : (\operatorname{is\text{-}prime\text{-}nat}(d) \wedge \operatorname{zz\text{-}divides}(d, n))\}) \end{array}}}

\vnbbookfile{theorem-library/cc-power-add}{theorem-library/cc-power-add}
\noindent{\small cc-power-add.scm -- lemmas MOVED 2026-10-27 (integration of the 10-26/27 wave) out of theorem-library/cc-exp.scm, so that files of library analysis may cite them (the library gate). The proofs are the originals, verbatim; the helper defines are a copy of cc-exp.scm's.  Loaded immediately before cc-exp.scm.}\par
\vnbprop{cc-power-add}{\vnbstmt{\forall n, m \in \mathbb{N}, z \in \mathbb{C}.\; (\operatorname{power}(z, (m + n)) = (\operatorname{power}(z, m) \cdot \operatorname{power}(z, n)))}}

\vnbbookfile{theorem-library/fekete}{theorem-library/fekete}
\noindent{\small fekete.scm -- wave of 2026-11-09/10 (agent BA-1C), wired 2026-11-10 from scratchpad/ba1c-1109/fekete.scm. Library: analysis.}\par
\vnbprop{fek-root-power}{\vnbstmt{\begin{array}{@{}l@{}} \forall k \in \mathbb{R}, n \in \mathbb{N}. \\ \quad ((0 < k) \wedge (1 \leq n)) \Rightarrow ({\operatorname{power}(k, n)}^{{n}^{\ast\!-1}} = k) \end{array}}}
\vnbprop{fek-root-le}{\vnbstmt{\begin{array}{@{}l@{}} \forall x, k \in \mathbb{R}, n \in \mathbb{N}. \\ \quad ((0 \leq x) \wedge ((0 < k) \wedge ((1 \leq n) \wedge (x \leq \operatorname{power}(k, n))))) \Rightarrow  \\ \quad \quad {x}^{{n}^{\ast\!-1}} \leq k \end{array}}}
\vnbprop{fek-root-lt}{\vnbstmt{\begin{array}{@{}l@{}} \forall x, k \in \mathbb{R}, n \in \mathbb{N}. \\ \quad ((0 \leq x) \wedge ((0 < k) \wedge ((1 \leq n) \wedge ({x}^{{n}^{\ast\!-1}} < k)))) \Rightarrow  \\ \quad \quad x < \operatorname{power}(k, n) \end{array}}}
\vnbprop{fek-pow-mult}{\vnbstmt{\begin{array}{@{}l@{}} \forall p \in \mathbb{N}, l \in (\mathbb{N} \to \mathbb{R}), m \in \mathbb{N}. \\ \quad \forall n \in \mathbb{N}.\; (0 \leq l(n)) \wedge \\ \quad \forall a, n \in \mathbb{N}.\; (((1 \leq a) \wedge (1 \leq n)) \Rightarrow (l((a + n)) \leq (l(a) \cdot l(n)))) \wedge \\ \quad 1 \leq m \Rightarrow \\ \quad \quad l((m \cdot (p + 1))) \leq \operatorname{power}(l(m), (p + 1)) \end{array}}}
\vnbprop{fek-bound}{\vnbstmt{\begin{array}{@{}l@{}} \forall m \in \mathbb{N}, l \in (\mathbb{N} \to \mathbb{R}). \\ \quad \forall n \in \mathbb{N}.\; (0 \leq l(n)) \Rightarrow  \\ \quad \quad \exists d \in \mathbb{R}.\; ((1 \leq d) \wedge \forall r \in \mathbb{N}.\; ((r \leq m) \Rightarrow (l(r) \leq d))) \end{array}}}
\vnbprop{fek-eventually}{\vnbstmt{\begin{array}{@{}l@{}} \forall l \in (\mathbb{N} \to \mathbb{R}), m \in \mathbb{N}. \\ \quad \forall a \in \mathbb{N}.\; (0 \leq l(a)) \wedge \\ \quad \forall b, a \in \mathbb{N}.\; (((1 \leq b) \wedge (1 \leq a)) \Rightarrow (l((b + a)) \leq (l(b) \cdot l(a)))) \wedge \\ \quad 1 \leq m \wedge \\ \quad l(m) < 1 \Rightarrow \\ \quad \quad \exists n_{0} \in \mathbb{N}.\; \forall n \in \mathbb{N}.\; ((n_{0} \leq n) \Rightarrow (l(n) \leq 1)) \end{array}}}
\vnbprop{fek-upper}{\vnbstmt{\begin{array}{@{}l@{}} \forall l \in (\mathbb{N} \to \mathbb{R}), k \in \mathbb{R}, m \in \mathbb{N}. \\ \quad \forall a \in \mathbb{N}.\; (0 \leq l(a)) \wedge \\ \quad \forall b, a \in \mathbb{N}.\; (((1 \leq b) \wedge (1 \leq a)) \Rightarrow (l((b + a)) \leq (l(b) \cdot l(a)))) \wedge \\ \quad 0 < k \wedge \\ \quad 1 \leq m \wedge \\ \quad l(m) < \operatorname{power}(k, m) \Rightarrow \\ \quad \quad \exists n_{0} \in \mathbb{N}.\; \forall n \in \mathbb{N}.\; ((n_{0} \leq n) \Rightarrow (l(n) \leq \operatorname{power}(k, n))) \end{array}}}
\vnbprop{fekete-lemma}{\vnbstmt{\begin{array}{@{}l@{}} \forall l \in (\mathbb{N} \to \mathbb{R}), v \in \mathbb{R}, g \in (\mathbb{N} \to \mathbb{R}). \\ \quad \forall a \in \mathbb{N}.\; (0 \leq l(a)) \wedge \\ \quad \forall m, a \in \mathbb{N}.\; (((1 \leq m) \wedge (1 \leq a)) \Rightarrow (l((m + a)) \leq (l(m) \cdot l(a)))) \wedge \\ \quad \forall k \in \mathbb{N}.\; ((1 \leq k) \Rightarrow (v \leq {l(k)}^{{k}^{\ast\!-1}})) \wedge \\ \quad \forall s. \\ \quad \quad \operatorname{pos\text{-}rr}(s) \Rightarrow  \\ \quad \quad \quad \exists k \in \mathbb{N}.\; ((1 \leq k) \wedge ({l(k)}^{{k}^{\ast\!-1}} < (v + s))) \wedge \\ \quad \forall n \in \mathbb{N}.\; (g(n) = {l((n + 1))}^{{(n + 1)}^{\ast\!-1}}) \Rightarrow \\ \quad \quad \operatorname{converges\text{-}to}(\operatorname{rr\text{-}ms}, g, v) \end{array}}}

\vnbbookfile{theorem-library/phi-growth}{theorem-library/phi-growth}
\noindent{\small theorem-library/phi-growth.scm -- LEMMA 1.24, PROP. 1.25, COR. 1.26 (ALGEBRAIC-NUMBERS notes, ch. 1; ALGNUM, 2026-10-27): the growth of phi. nn-power-closed, rr-power-ge-one, rr-power-mono-exp, rr-power-mono-le, rr-power-ge-base   natural powers prime-divides-prime-power, prime-decomposition(-bounded), ...}\par
\vnbprop{nn-power-closed}{\vnbstmt{\forall k, x \in \mathbb{N}.\; (\operatorname{power}(x, k) \in \mathbb{N})}}
\vnbprop{rr-power-ge-one}{\vnbstmt{\forall k \in \mathbb{N}, x \in \mathbb{R}.\; ((1 \leq x) \Rightarrow (1 \leq \operatorname{power}(x, k)))}}
\vnbprop{rr-power-mono-exp}{\vnbstmt{\forall d, k \in \mathbb{N}, x \in \mathbb{R}.\; ((1 \leq x) \Rightarrow (\operatorname{power}(x, k) \leq \operatorname{power}(x, (k + d))))}}
\vnbprop{prime-divides-prime-power}{\vnbstmt{\begin{array}{@{}l@{}} \forall a \in \mathbb{N}, q, r. \\ \quad (\operatorname{is\text{-}prime\text{-}nat}(q) \wedge (\operatorname{is\text{-}prime\text{-}nat}(r) \wedge \operatorname{zz\text{-}divides}(r, \operatorname{power}(q, a)))) \Rightarrow (r = q) \end{array}}}
\vnbprop{prime-decomposition-bounded}{\vnbstmt{\begin{array}{@{}l@{}} \forall k \in \mathbb{N}, q.; \forall n \in \mathbb{N}. \\ \quad (\operatorname{is\text{-}prime\text{-}nat}(q) \wedge ((1 \leq n) \wedge (n \leq k))) \Rightarrow  \\ \quad \quad \exists a, m \in \mathbb{N}.\; ((1 \leq m) \wedge ((n = (\operatorname{power}(q, a) \cdot m)) \wedge \neg \operatorname{zz\text{-}divides}(q, m))) \end{array}}}
\vnbprop{prime-decomposition}{\vnbstmt{\begin{array}{@{}l@{}} \forall q.; \forall n \in \mathbb{N}. \\ \quad (\operatorname{is\text{-}prime\text{-}nat}(q) \wedge (1 \leq n)) \Rightarrow  \\ \quad \quad \exists a, m \in \mathbb{N}.\; ((1 \leq m) \wedge ((n = (\operatorname{power}(q, a) \cdot m)) \wedge \neg \operatorname{zz\text{-}divides}(q, m))) \end{array}}}
\vnbprop{coprime-prime-power}{\vnbstmt{\begin{array}{@{}l@{}} \forall a \in \mathbb{N}, q.; \forall m \in \mathbb{Z}. \\ \quad (\operatorname{is\text{-}prime\text{-}nat}(q) \wedge \neg \operatorname{zz\text{-}divides}(q, m)) \Rightarrow \operatorname{zz\text{-}coprime}(m, \operatorname{power}(q, a)) \end{array}}}
\vnbprop{fun-abs-bounded-initial}{\vnbstmt{\begin{array}{@{}l@{}} \forall k \in \mathbb{N}, f \in (\mathbb{N} \to \mathbb{R}). \exists m \in \mathbb{R}. \\ \quad (1 \leq m) \wedge \forall l \in \mathbb{N}.\; (((1 \leq l) \wedge (l \leq k)) \Rightarrow (\lvert f(l) \rvert \leq m)) \end{array}}}
\vnbprop{rr-power-ge-base}{\vnbstmt{\forall a \in \mathbb{N}, x \in \mathbb{R}.\; (((1 \leq x) \wedge (1 \leq a)) \Rightarrow (x \leq \operatorname{power}(x, a)))}}
\vnbprop{rr-power-mono-le}{\vnbstmt{\forall j, k \in \mathbb{N}, x \in \mathbb{R}.\; (((1 \leq x) \wedge (k \leq j)) \Rightarrow (\operatorname{power}(x, k) \leq \operatorname{power}(x, j)))}}
\vnbprop{multiplicative-bound-step}{\vnbstmt{\begin{array}{@{}l@{}} \forall k \in \mathbb{N}, f.; \forall b \in \mathbb{N}, a \in \mathbb{R}, n \in \mathbb{N}. \\ \quad \operatorname{is\text{-}multiplicative}(f) \wedge \\ \quad \forall p.; \forall \mathit{mbk} \in \mathbb{N}. \\ \quad \quad (\operatorname{is\text{-}prime\text{-}nat}(p) \wedge (b < \operatorname{power}(p, \mathit{mbk}))) \Rightarrow (\lvert f(\operatorname{power}(p, \mathit{mbk})) \rvert \leq 1) \wedge \\ \quad 1 \leq a \wedge \\ \quad \forall l \in \mathbb{N}.\; (((1 \leq l) \wedge (l \leq (b + 1))) \Rightarrow (\lvert f(l) \rvert \leq a)) \wedge \\ \quad 1 \leq n \wedge \\ \quad n \leq k \Rightarrow \\ \quad \quad \lvert f(n) \rvert \leq \operatorname{power}(a, (|\{t \in \operatorname{interval}(1, b) : (\operatorname{is\text{-}prime\text{-}nat}(t) \wedge \operatorname{zz\text{-}divides}(t, n))\}| + 1)) \end{array}}}
\vnbprop{multiplicative-bounded}{\vnbstmt{\begin{array}{@{}l@{}} \forall f.; \forall b \in \mathbb{N}. \\ \quad \operatorname{is\text{-}multiplicative}(f) \wedge \\ \quad \forall p.; \forall k \in \mathbb{N}.\; ((\operatorname{is\text{-}prime\text{-}nat}(p) \wedge (b < \operatorname{power}(p, k))) \Rightarrow (\lvert f(\operatorname{power}(p, k)) \rvert \leq 1)) \Rightarrow \\ \quad \quad \exists \mathit{mbb} \in \mathbb{R}.\; \forall n \in \mathbb{N}.\; ((1 \leq n) \Rightarrow (\lvert f(n) \rvert \leq \mathit{mbb})) \end{array}}}
\vnbprop{euler-phi-pos}{\vnbstmt{\forall n \in \mathbb{N}.\; ((1 \leq n) \Rightarrow ((\operatorname{euler\text{-}phi}(n) \in \mathbb{N}) \wedge (1 \leq \operatorname{euler\text{-}phi}(n))))}}
\vnbprop{euler-phi-prime-power}{\vnbstmt{\begin{array}{@{}l@{}} \forall j \in \mathbb{N}, p. \\ \quad \operatorname{is\text{-}prime\text{-}nat}(p) \Rightarrow (\operatorname{euler\text{-}phi}(\operatorname{power}(p, (j + 1))) = (\operatorname{power}(p, (j + 1)) \cdot (1 - \operatorname{recip}(p)))) \end{array}}}
\vnbprop{euler-phi-ratio-multiplicative}{\vnbstmt{\begin{array}{@{}l@{}} \forall e \in \mathbb{Q}. \\ \quad \operatorname{is\text{-}multiplicative}((\lambda k \in \mathbb{N}.\; \left(\text{if }(k = 0)\text{ then }0\text{ else }(\operatorname{rpow}(k, e) \cdot \operatorname{recip}(\operatorname{euler\text{-}phi}(k)))\right))) \end{array}}}
\vnbprop{euler-phi-ratio-prime-power}{\vnbstmt{\begin{array}{@{}l@{}} \forall t \in \mathbb{Q}, y \in \mathbb{R}, j \in \mathbb{N}, p. \\ \quad ((0 < t) \wedge ((0 < y) \wedge ((\operatorname{rpow}(y, t) = 2) \wedge (\operatorname{is\text{-}prime\text{-}nat}(p) \wedge (y \leq \operatorname{power}(p, (j + 1))))))) \Rightarrow  \\ \quad \quad (\operatorname{rpow}(\operatorname{power}(p, (j + 1)), (1 - t)) \cdot \operatorname{recip}(\operatorname{euler\text{-}phi}(\operatorname{power}(p, (j + 1))))) \leq 1 \end{array}}}
\vnbprop{euler-phi-ratio-bounded}{\vnbstmt{\begin{array}{@{}l@{}} \forall t \in \mathbb{Q}. \\ \quad (0 < t) \Rightarrow  \\ \quad \quad \exists b \in \mathbb{R}. \forall n \in \mathbb{N}. \\ \quad \quad \quad (1 \leq n) \Rightarrow ((\operatorname{rpow}(n, (1 - t)) \cdot \operatorname{recip}(\operatorname{euler\text{-}phi}(n))) \leq b) \end{array}}}
\vnbprop{euler-phi-lower-bound}{\vnbstmt{\begin{array}{@{}l@{}} \forall d \in \mathbb{Q}. \\ \quad (0 < d) \Rightarrow  \\ \quad \quad \exists c \in \mathbb{R}. \\ \quad \quad \quad (0 < c) \wedge \forall n \in \mathbb{N}.\; ((1 \leq n) \Rightarrow ((c \cdot \operatorname{rpow}(n, (1 - d))) \leq \operatorname{euler\text{-}phi}(n))) \end{array}}}
\vnbprop{euler-phi-ratio-limit}{\vnbstmt{\begin{array}{@{}l@{}} \forall d \in \mathbb{Q}. \\ \quad (0 < d) \Rightarrow  \\ \quad \quad \operatorname{converges\text{-}to}(\operatorname{rr\text{-}ms}, (\lambda k \in \mathbb{N}.\; (\operatorname{rpow}((k + 1), (1 - d)) \cdot \operatorname{recip}(\operatorname{euler\text{-}phi}((k + 1))))), 0) \end{array}}}

\vnbbookfile{theorem-library/rr-pos-star-order-lemmas}{theorem-library/rr-pos-star-order-lemmas}
\noindent{\small rr-pos-star-order-lemmas.scm -- lemmas MOVED 2026-10-27 (integration of the 10-26/27 wave) out of theorem-library/integral-convergence.scm, so that files of library analysis may cite them (the library gate). The proofs are the originals, verbatim; the helper defines are a copy of integral-convergence.scm's.  Loaded ...}\par
\vnbprop{rps-le-total}{\vnbstmt{\forall x, y \in \operatorname{rr\text{-}pos\text{-}star}.\; ((x \leq y) \vee (y \leq x))}}
\vnbprop{rps-not-lt-le}{\vnbstmt{\forall x, y \in \operatorname{rr\text{-}pos\text{-}star}.\; (\neg (x < y) \Rightarrow (y \leq x))}}
\vnbprop{rr-nn-floor-mul}{\vnbstmt{\begin{array}{@{}l@{}} \forall c \in \mathbb{R}, r. \\ \quad ((0 \leq c) \wedge \operatorname{pos\text{-}rr}(r)) \Rightarrow \exists f \in \mathbb{N}.\; (((f \cdot r) \leq c) \wedge (c < ((f \cdot r) + r))) \end{array}}}
\vnbprop{rr-recip-swap-lt}{\vnbstmt{\begin{array}{@{}l@{}} \forall l, q, g \in \mathbb{R}. \\ \quad ((0 < g) \wedge ((0 < q) \wedge ((l \cdot \operatorname{recip}(g)) < q))) \Rightarrow ((l \cdot \operatorname{recip}(q)) < g) \end{array}}}
\vnbprop{rr-lt-mul-of-div-lt}{\vnbstmt{\begin{array}{@{}l@{}} \forall l, q, g \in \mathbb{R}. \\ \quad ((0 < g) \wedge ((l \cdot \operatorname{recip}(g)) < q)) \Rightarrow (l < (q \cdot g)) \end{array}}}

\vnbbookfile{theorem-library/rps-lt-le-trans}{theorem-library/rps-lt-le-trans}
\noindent{\small rps-lt-le-trans.scm -- lemmas MOVED 2026-10-30 (wave-2 integration) out of theorem-library/integral-convergence.scm into library analysis, for the library gate: cited by ch4a-coef (complex-analysis); an RR-POS-STAR order fact. The proofs are the originals, verbatim; the helper defines are a copy of the source's. ...}\par
\vnbprop{rps-lt-le-trans}{\vnbstmt{\begin{array}{@{}l@{}} \forall x, y, z \in \operatorname{rr\text{-}pos\text{-}star}. \\ \quad ((x < y) \wedge (y \leq z)) \Rightarrow (x < z) \end{array}}}

\vnbbookfile{theorem-library/esup-lt-witness}{theorem-library/esup-lt-witness}
\noindent{\small esup-lt-witness.scm -- lemmas MOVED 2026-11-10 (integration of the 11-09/10 waves) out of theorem-library/integral-convergence.scm into library analysis, for the library gate: cited by haar-fixed-point; ESUP. The proofs are the originals, verbatim; the helper defines are a copy of the source's. esup-lt-witness}\par
\vnbprop{esup-lt-witness}{\vnbstmt{\begin{array}{@{}l@{}} \forall s.; \forall a \in \operatorname{rr\text{-}pos\text{-}star}. \\ \quad ((s \subseteq \operatorname{rr\text{-}pos\text{-}star}) \wedge (a < \operatorname{esup}(s))) \Rightarrow \exists y \in s.\; (a < y) \end{array}}}

\vnbbookfile{theorem-library/subset-rr-power}{theorem-library/subset-rr-power}
\noindent{\small subset-rr-power.scm -- lemmas MOVED 2026-10-27 (integration of the 10-26/27 wave) out of theorem-library/lebesgue-interval-length.scm, so that files of library analysis may cite them (the library gate). The proofs are the originals, verbatim; the helper defines are a copy of lebesgue-interval-length.scm's.  Loaded ...}\par
\vnbprop{subset-rr-in-power-rr}{\vnbstmt{\forall s.\; ((s \subseteq \mathbb{R}) \Rightarrow (s \in \operatorname{power}(\mathbb{R})))}}

\vnbbookfile{theorem-library/metric-finite-subcover}{theorem-library/metric-finite-subcover}
\noindent{\small theorem-library/metric-finite-subcover.scm -- a closed set of a compact metric space covered by countably many open sets is covered by finitely many of them. Library: analysis.  MOVED 2026-10-20 (the wiring of CALC-3) from theorem-library/riesz-markov-cantor.scm (library measure), where RM-2 proved it as}\par
\vnbprop{compact-closed-countable-subcover}{\vnbstmt{\begin{array}{@{}l@{}} \forall m, a, q. \\ \quad \operatorname{is\text{-}compact}(m) \wedge \\ \quad \operatorname{is\text{-}closed}(m, a) \wedge \\ \quad \forall n \in \mathbb{N}.\; \operatorname{is\text{-}open}(m, q(n)) \wedge \\ \quad \forall x \in a.\; \exists n \in \mathbb{N}.\; (x \in q(n)) \Rightarrow \\ \quad \quad \exists f \in \mathbb{N}.\; \forall x \in a.\; \exists n \in \mathbb{N}.\; ((n < f) \wedge (x \in q(n))) \end{array}}}

\vnbbookfile{structure-library/positive-functional}{structure-library/positive-functional}
\noindent{\small structure-library/positive-functional.scm -- continuous real functions on a metric space, positive linear functionals on them, clopen sets.  DEFINITIONS ONLY (Riesz-Markov, stage 2, RM-2, 2026-10-10); the laws are PROVEN in theorem-library/riesz-markov-cantor.scm (prefix rmc-).}\par
\par\medskip\noindent\textbf{Definition} (\texttt{cont-fns}).\;$\operatorname{cont\text{-}fns}(s) :=$ $\{\operatorname{cfh\_} \in (\operatorname{pts}(s) \to \mathbb{R}) : \operatorname{is\text{-}continuous}(s, \operatorname{rr\text{-}ms}, \operatorname{cfh\_})\}$
\par\medskip\noindent\textbf{Definition} (\texttt{is-positive-functional}).\;For all $s$ and $\mathit{lam}$, 
\begin{equation*}
\begin{array}{@{}l@{}} \operatorname{is\text{-}positive\text{-}functional}(s, \mathit{lam}) \\ \Leftrightarrow\; \mathit{lam} \in (\operatorname{cont\text{-}fns}(s) \to \mathbb{R}) \wedge \\ \forall \operatorname{pff\_}, \operatorname{pfg\_} \in \operatorname{cont\text{-}fns}(s). \\ \quad \operatorname{lam}((\lambda \operatorname{sx\_} \in \operatorname{pts}(s).\; +(\operatorname{pff\_}(\operatorname{sx\_}), \operatorname{pfg\_}(\operatorname{sx\_})))) = +(\operatorname{lam}(\operatorname{pff\_}), \operatorname{lam}(\operatorname{pfg\_})) \wedge \\ \forall \operatorname{pfc\_} \in \mathbb{R}, \operatorname{pff\_} \in \operatorname{cont\text{-}fns}(s). \\ \quad \operatorname{lam}((\lambda \operatorname{sx\_} \in \operatorname{pts}(s).\; \cdot(\operatorname{pfc\_}, \operatorname{pff\_}(\operatorname{sx\_})))) = \cdot(\operatorname{pfc\_}, \operatorname{lam}(\operatorname{pff\_})) \wedge \\ \forall \operatorname{pff\_} \in \operatorname{cont\text{-}fns}(s). \\ \quad \forall \operatorname{pfx\_} \in \operatorname{pts}(s).\; (0 \leq \operatorname{pff\_}(\operatorname{pfx\_})) \Rightarrow (0 \leq \operatorname{lam}(\operatorname{pff\_})) \end{array}
\end{equation*}

\par\medskip\noindent\textbf{Definition} (\texttt{clopen-sets}).\;$\operatorname{clopen\text{-}sets}(s) :=$
\begin{equation*}
\{\operatorname{cou\_} \in \operatorname{power}(\operatorname{pts}(s)) : (\operatorname{is\text{-}open}(s, \operatorname{cou\_}) \wedge \operatorname{is\text{-}closed}(s, \operatorname{cou\_}))\}
\end{equation*}


\vnbbookfile{structure-library/hc2-vocab}{structure-library/hc2-vocab}
\noindent{\small hc2-vocab.scm -- wave of 2026-11-09/10 (agent HAAR-C2), wired 2026-11-10 from scratchpad/haarc2-1109/hc2-vocab.scm. Library: analysis.}\par
\par\medskip\noindent\textbf{Definition} (\texttt{ck-norm}).\;$\lVert \operatorname{ckf\_} \rVert_{\infty} :=$
\begin{equation*}
\operatorname{esup}(\operatorname{image}((\lambda \operatorname{ckp\_} \in \operatorname{pts}(\operatorname{cks\_}).\; \lvert \operatorname{ckf\_}(\operatorname{ckp\_}) \rvert), \operatorname{pts}(\operatorname{cks\_})))
\end{equation*}

\par\medskip\noindent\textbf{Definition} (\texttt{ck-nvs}).\;$C({\operatorname{cks\_}}) :=$
\begin{equation*}
{[}\operatorname{normed\text{-}field\text{-}as\text{-}commutative\text{-}ring}(\operatorname{rr\text{-}normed\text{-}field}), \operatorname{cont\text{-}fns}(\operatorname{cks\_}), (\lambda {[}\operatorname{ckaf\_}, \operatorname{ckag\_}] \in (\operatorname{cont\text{-}fns}(\operatorname{cks\_}) \times \operatorname{cont\text{-}fns}(\operatorname{cks\_})).\; (\lambda \operatorname{sx\_} \in \operatorname{pts}(\operatorname{cks\_}).\; (\operatorname{ckaf\_}(\operatorname{sx\_}) + \operatorname{ckag\_}(\operatorname{sx\_})))), (\lambda \operatorname{sx\_} \in \operatorname{pts}(\operatorname{cks\_}).\; 0), (\lambda \operatorname{cknf\_} \in \operatorname{cont\text{-}fns}(\operatorname{cks\_}).\; (\lambda \operatorname{sx\_} \in \operatorname{pts}(\operatorname{cks\_}).\; -\operatorname{cknf\_}(\operatorname{sx\_}))), (\lambda {[}\operatorname{cksr\_}, \operatorname{cksf\_}] \in (\mathbb{R} \times \operatorname{cont\text{-}fns}(\operatorname{cks\_})).\; (\lambda \operatorname{sx\_} \in \operatorname{pts}(\operatorname{cks\_}).\; (\operatorname{cksr\_} \cdot \operatorname{cksf\_}(\operatorname{sx\_})))), (\lambda \operatorname{ckvf\_} \in \operatorname{cont\text{-}fns}(\operatorname{cks\_}).\; \lVert \operatorname{ckvf\_} \rVert_{\infty})]
\end{equation*}


\vnbbookfile{theorem-library/cont-fns-laws}{theorem-library/cont-fns-laws}
\noindent{\small theorem-library/cont-fns-laws.scm -- CONT-FNS(s) and CLOPEN-SETS(s) of a metric space s: the membership laws, closure of CONT-FNS under pointwise sums and scalar multiples, and a function constant on an open set around a is continuous at a.  Library: analysis.}\par
\vnbprop{rmc-cont-fns-mem}{\vnbstmt{\begin{array}{@{}l@{}} \forall s, h. \\ \quad (h \in \operatorname{cont\text{-}fns}(s)) \Leftrightarrow ((h \in (\operatorname{pts}(s) \to \mathbb{R})) \wedge \operatorname{is\text{-}continuous}(s, \operatorname{rr\text{-}ms}, h)) \end{array}}}
\vnbprop{rmc-clopen-mem}{\vnbstmt{\begin{array}{@{}l@{}} \forall s, u. \\ \quad (u \in \operatorname{clopen\text{-}sets}(s)) \Leftrightarrow ((u \in \operatorname{power}(\operatorname{pts}(s))) \wedge (\operatorname{is\text{-}open}(s, u) \wedge \operatorname{is\text{-}closed}(s, u))) \end{array}}}
\vnbprop{rmc-clopen-intro}{\vnbstmt{\forall s, u.\; ((\operatorname{is\text{-}open}(s, u) \wedge \operatorname{is\text{-}closed}(s, u)) \Rightarrow (u \in \operatorname{clopen\text{-}sets}(s)))}}
\vnbprop{rmc-clopen-parts}{\vnbstmt{\begin{array}{@{}l@{}} \forall s.; \forall u \in \operatorname{clopen\text{-}sets}(s). \\ \quad \operatorname{is\text{-}open}(s, u) \wedge (\operatorname{is\text{-}closed}(s, u) \wedge (u \subseteq \operatorname{pts}(s))) \end{array}}}
\vnbprop{rmc-cont-fns-parts}{\vnbstmt{\begin{array}{@{}l@{}} \forall s.; \forall h \in \operatorname{cont\text{-}fns}(s). \\ \quad h \in (\operatorname{pts}(s) \to \mathbb{R}) \wedge \\ \quad \operatorname{is\text{-}metric\text{-}space}(s) \wedge \\ \quad \forall a \in \operatorname{pts}(s).\; \operatorname{is\text{-}continuous\text{-}at}(s, \operatorname{rr\text{-}ms}, h, a) \end{array}}}
\vnbprop{rmc-cont-fns-intro}{\vnbstmt{\begin{array}{@{}l@{}} \forall s, h. \\ \quad \operatorname{is\text{-}metric\text{-}space}(s) \wedge \\ \quad h \in (\operatorname{pts}(s) \to \mathbb{R}) \wedge \\ \quad \forall a \in \operatorname{pts}(s).\; \operatorname{is\text{-}continuous\text{-}at}(s, \operatorname{rr\text{-}ms}, h, a) \Rightarrow \\ \quad \quad h \in \operatorname{cont\text{-}fns}(s) \end{array}}}
\vnbprop{rmc-cont-fns-scale}{\vnbstmt{\begin{array}{@{}l@{}} \forall s, c, f. \\ \quad ((c \in \mathbb{R}) \wedge (f \in \operatorname{cont\text{-}fns}(s))) \Rightarrow  \\ \quad \quad (\lambda x \in \operatorname{pts}(s).\; (c \cdot f(x))) \in \operatorname{cont\text{-}fns}(s) \end{array}}}
\vnbprop{rmc-cont-fns-add}{\vnbstmt{\begin{array}{@{}l@{}} \forall s, f, g. \\ \quad ((f \in \operatorname{cont\text{-}fns}(s)) \wedge (g \in \operatorname{cont\text{-}fns}(s))) \Rightarrow  \\ \quad \quad (\lambda x \in \operatorname{pts}(s).\; (f(x) + g(x))) \in \operatorname{cont\text{-}fns}(s) \end{array}}}
\vnbprop{rmc-cont-at-locally-const}{\vnbstmt{\begin{array}{@{}l@{}} \forall s, h, a, w. \\ \quad (\operatorname{is\text{-}open}(s, w) \wedge ((a \in w) \wedge ((h \in (\operatorname{pts}(s) \to \mathbb{R})) \wedge \forall y \in w.\; (h(y) = h(a))))) \Rightarrow  \\ \quad \quad \operatorname{is\text{-}continuous\text{-}at}(s, \operatorname{rr\text{-}ms}, h, a) \end{array}}}

\vnbbookfile{structure-library/m-riesz-vocab}{structure-library/m-riesz-vocab}
\noindent{\small m-riesz-vocab.scm -- ORDERED VECTOR SPACES: BanachAlgebras.pdf section 4.6, before Proposition 4.32 (the M. Riesz extension).  Agent BA-1117-GNS2, 2026-11-17. Destination: structure-library/m-riesz-vocab.scm.  Library: analysis. DEFINITIONS ONLY.  Load window: after structure-library/module (VEC, VADD, VNEG, ACT, ...}\par
\par\medskip\noindent\textbf{Definition} (\texttt{is-pointed-cone}).\;For all $\operatorname{g4rm\_}$ and $\operatorname{g4rc\_}$, 
\begin{equation*}
\begin{array}{@{}l@{}} \operatorname{is\text{-}pointed\text{-}cone}(\operatorname{g4rm\_}, \operatorname{g4rc\_}) \\ \Leftrightarrow\; \operatorname{g4rc\_} \subseteq \operatorname{vec}(\operatorname{g4rm\_}) \wedge \\ \operatorname{vzero}(\operatorname{g4rm\_}) \in \operatorname{g4rc\_} \wedge \\ \forall \operatorname{g4rx\_}, \operatorname{g4ry\_} \in \operatorname{g4rc\_}.\; (\operatorname{vadd}(\operatorname{g4rm\_})(\operatorname{g4rx\_}, \operatorname{g4ry\_}) \in \operatorname{g4rc\_}) \wedge \\ \forall \operatorname{g4rt\_} \in \mathbb{R}. \\ \quad (0 \leq \operatorname{g4rt\_}) \Rightarrow \forall \operatorname{g4rx\_} \in \operatorname{g4rc\_}.\; (\operatorname{act}(\operatorname{g4rm\_})(\operatorname{g4rt\_}, \operatorname{g4rx\_}) \in \operatorname{g4rc\_}) \end{array}
\end{equation*}

\par\medskip\noindent\textbf{Definition} (\texttt{cone-le}).\;For all $\operatorname{g4rm\_}$, $\operatorname{g4rc\_}$, $\operatorname{g4rx\_}$ and $\operatorname{g4ry\_}$, 
\begin{equation*}
{\operatorname{g4rx\_}} \preceq_{\operatorname{g4rc\_}} {\operatorname{g4ry\_}} \Leftrightarrow (\operatorname{vadd}(\operatorname{g4rm\_})(\operatorname{g4ry\_}, \operatorname{vneg}(\operatorname{g4rm\_})(\operatorname{g4rx\_})) \in \operatorname{g4rc\_})
\end{equation*}

\par\medskip\noindent\textbf{Definition} (\texttt{is-cofinal}).\;For all $\operatorname{g4rm\_}$, $\operatorname{g4rc\_}$ and $\operatorname{g4rv\_}$, 
\begin{equation*}
\begin{array}{@{}l@{}} \operatorname{is\text{-}cofinal}(\operatorname{g4rm\_}, \operatorname{g4rc\_}, \operatorname{g4rv\_}) \\ \Leftrightarrow\; \forall \operatorname{g4ra\_} \in \operatorname{vec}(\operatorname{g4rm\_}). \exists \operatorname{g4rx\_} \in \operatorname{g4rv\_}. \\ \quad {\operatorname{g4ra\_}} \preceq_{\operatorname{g4rc\_}} {\operatorname{g4rx\_}} \end{array}
\end{equation*}


\vnbbookfile{theorem-library/m-riesz}{theorem-library/m-riesz}
\noindent{\small m-riesz.scm -- PROPOSITION 4.32, THE M. RIESZ EXTENSION: E a real vector space ordered by a pointed cone C, V a cofinal subspace, f linear on V and non-negative on C n V.  Then f extends to a linear functional on E non-negative on C.  F. J. Thayer, "Banach Algebras and Applications", Prop. 4.32. Agent BA-1117-GNS2, ...}\par
\vnbprop{m-riesz-f-mono}{\vnbstmt{\begin{array}{@{}l@{}} \forall m, c, v, f.; \forall z, w \in v. \\ \quad \operatorname{is\text{-}module}(m) \wedge \\ \quad \operatorname{scal}(m) = \operatorname{normed\text{-}field\text{-}as\text{-}commutative\text{-}ring}(\operatorname{rr\text{-}normed\text{-}field}) \wedge \\ \quad \operatorname{is\text{-}pointed\text{-}cone}(m, c) \wedge \\ \quad \operatorname{is\text{-}submodule}(m, v) \wedge \\ \quad \operatorname{is\text{-}cofinal}(m, c, v) \wedge \\ \quad \operatorname{is\text{-}linear\text{-}functional\text{-}on}(m, v, f) \wedge \\ \quad \forall z \in v.\; ((z \in c) \Rightarrow (0 \leq f(z))) \wedge \\ \quad \operatorname{vadd}(m)(w, \operatorname{vneg}(m)(z)) \in c \Rightarrow \\ \quad \quad f(z) \leq f(w) \end{array}}}
\vnbprop{m-riesz-iif-mem}{\vnbstmt{\begin{array}{@{}l@{}} \forall m, c, v, f.; \forall g \in \operatorname{vec}(m), w. \\ \quad \operatorname{is\text{-}module}(m) \wedge \\ \quad \operatorname{scal}(m) = \operatorname{normed\text{-}field\text{-}as\text{-}commutative\text{-}ring}(\operatorname{rr\text{-}normed\text{-}field}) \wedge \\ \quad \operatorname{is\text{-}pointed\text{-}cone}(m, c) \wedge \\ \quad \operatorname{is\text{-}submodule}(m, v) \wedge \\ \quad \operatorname{is\text{-}cofinal}(m, c, v) \wedge \\ \quad \operatorname{is\text{-}linear\text{-}functional\text{-}on}(m, v, f) \wedge \\ \quad \forall z \in v.\; ((z \in c) \Rightarrow (0 \leq f(z))) \Rightarrow \\ \quad \quad w \in (\lambda g \in \operatorname{vec}(m).\; \{s \in v : (\operatorname{vadd}(m)(s, \operatorname{vneg}(m)(g)) \in c)\})(g) \\ \quad \quad \Leftrightarrow\; w \in \{s \in v : (\operatorname{vadd}(m)(s, \operatorname{vneg}(m)(g)) \in c)\} \end{array}}}
\vnbprop{m-riesz-p-exists}{\vnbstmt{\begin{array}{@{}l@{}} \forall m, c, v, f. \\ \quad \operatorname{is\text{-}module}(m) \wedge \\ \quad \operatorname{scal}(m) = \operatorname{normed\text{-}field\text{-}as\text{-}commutative\text{-}ring}(\operatorname{rr\text{-}normed\text{-}field}) \wedge \\ \quad \operatorname{is\text{-}pointed\text{-}cone}(m, c) \wedge \\ \quad \operatorname{is\text{-}submodule}(m, v) \wedge \\ \quad \operatorname{is\text{-}cofinal}(m, c, v) \wedge \\ \quad \operatorname{is\text{-}linear\text{-}functional\text{-}on}(m, v, f) \wedge \\ \quad \forall z \in v.\; ((z \in c) \Rightarrow (0 \leq f(z))) \Rightarrow \\ \quad \quad \exists p \in (\operatorname{vec}(m) \to \mathbb{R}). \\ \quad \quad \quad \forall k \in \operatorname{vec}(m), u \in v. \\ \quad \quad \quad \quad (\operatorname{vadd}(m)(u, \operatorname{vneg}(m)(k)) \in c) \Rightarrow (p(k) \leq f(u)) \wedge \\ \quad \quad \quad \forall k \in \operatorname{vec}(m), b \in \mathbb{R}. \\ \quad \quad \quad \quad \forall u \in v.\; ((\operatorname{vadd}(m)(u, \operatorname{vneg}(m)(k)) \in c) \Rightarrow (b \leq f(u))) \Rightarrow  \\ \quad \quad \quad \quad \quad b \leq p(k) \end{array}}}
\vnbprop{m-riesz-p-props}{\vnbstmt{\begin{array}{@{}l@{}} \forall m, c, v, f.; \forall p \in (\operatorname{vec}(m) \to \mathbb{R}). \\ \quad \operatorname{is\text{-}module}(m) \wedge \\ \quad \operatorname{scal}(m) = \operatorname{normed\text{-}field\text{-}as\text{-}commutative\text{-}ring}(\operatorname{rr\text{-}normed\text{-}field}) \wedge \\ \quad \operatorname{is\text{-}pointed\text{-}cone}(m, c) \wedge \\ \quad \operatorname{is\text{-}submodule}(m, v) \wedge \\ \quad \operatorname{is\text{-}cofinal}(m, c, v) \wedge \\ \quad \operatorname{is\text{-}linear\text{-}functional\text{-}on}(m, v, f) \wedge \\ \quad \forall z \in v.\; ((z \in c) \Rightarrow (0 \leq f(z))) \wedge \\ \quad \forall k \in \operatorname{vec}(m), u \in v. \\ \quad \quad (\operatorname{vadd}(m)(u, \operatorname{vneg}(m)(k)) \in c) \Rightarrow (p(k) \leq f(u)) \wedge \\ \quad \forall k \in \operatorname{vec}(m), b \in \mathbb{R}. \\ \quad \quad \forall u \in v.\; ((\operatorname{vadd}(m)(u, \operatorname{vneg}(m)(k)) \in c) \Rightarrow (b \leq f(u))) \Rightarrow  \\ \quad \quad \quad b \leq p(k) \Rightarrow \\ \quad \quad (\operatorname{is\text{-}sublinear}(m, p) \wedge (\forall z \in v.\; (f(z) \leq p(z)) \wedge \forall x \in c.\; (p(\operatorname{vneg}(m)(x)) \leq 0))) \end{array}}}
\vnbprop{m-riesz-extension}{\vnbstmt{\begin{array}{@{}l@{}} \forall m, c, v, f. \\ \quad \operatorname{is\text{-}module}(m) \wedge \\ \quad \operatorname{scal}(m) = \operatorname{normed\text{-}field\text{-}as\text{-}commutative\text{-}ring}(\operatorname{rr\text{-}normed\text{-}field}) \wedge \\ \quad \operatorname{is\text{-}pointed\text{-}cone}(m, c) \wedge \\ \quad \operatorname{is\text{-}submodule}(m, v) \wedge \\ \quad \operatorname{is\text{-}cofinal}(m, c, v) \wedge \\ \quad \operatorname{is\text{-}linear\text{-}functional\text{-}on}(m, v, f) \wedge \\ \quad \forall z \in v.\; ((z \in c) \Rightarrow (0 \leq f(z))) \Rightarrow \\ \quad \quad \exists a. \\ \quad \quad \quad (\operatorname{is\text{-}linear\text{-}functional\text{-}on}(m, \operatorname{vec}(m), a) \wedge (\operatorname{extends\text{-}on}(v, a, f) \wedge \forall x \in c.\; (0 \leq a(x)))) \end{array}}}

\vnbbookfile{theorem-library/rr-lt-pos-inf}{theorem-library/rr-lt-pos-inf}
\noindent{\small rr-lt-pos-inf.scm -- lemmas MOVED 2026-10-30 (wave-2 integration) out of theorem-library/lp-inequalities.scm into library analysis, for the library gate: cited by ch4a-coef (complex-analysis); an extended-real order fact. The proofs are the originals, verbatim; the helper defines are a copy of the source's. ...}\par
\vnbprop{rr-lt-pos-inf}{\vnbstmt{\forall x \in \mathbb{R}.\; (x < \operatorname{pos\text{-}inf})}}

\vnbbookfile{theorem-library/etimes-real-in}{theorem-library/etimes-real-in}
\noindent{\small etimes-real-in.scm -- lemmas MOVED 2026-11-10 (integration of the 11-09/10 waves) out of theorem-library/lp-convergence.scm into library analysis, for the library gate: cited by l2-families-2; extended reals. The proofs are the originals, verbatim; the helper defines are a copy of the source's. etimes-real-in}\par
\vnbprop{etimes-real-in}{\vnbstmt{\begin{array}{@{}l@{}} \forall x, y. \\ \quad ((x \in \mathbb{R}) \wedge ((0 \leq x) \wedge ((y \in \mathbb{R}) \wedge (0 \leq y)))) \Rightarrow  \\ \quad \quad (\operatorname{etimes}(x, y) \in \mathbb{R}) \wedge (0 \leq \operatorname{etimes}(x, y)) \end{array}}}

\vnbbookfile{structure-library/sip-vocab}{structure-library/sip-vocab}
\noindent{\small sip-vocab.scm -- SEMI-INNER-PRODUCT SPACES (SIP-1, 2026-10-18).  Vocabulary only.}\par
\vnbstructdef{semi-inner-product-space}{A semi inner product space $s$ is a 7-tuple whose components are listed below.\par
\begin{enumerate}[nosep]
\item $\operatorname{scal}(s) = \operatorname{nth}(1, s)$, a substructure (\texttt{substructure}): $\operatorname{is\text{-}ring}(\operatorname{scal}(s))$.
\item $\operatorname{vec}(s) = \operatorname{nth}(2, s)$, a carrier (\texttt{carriers}): $\operatorname{vec}(s) \in \mathbf{Set}$.
\item $\operatorname{vadd}(s) = \operatorname{nth}(3, s)$, an operation (\texttt{op}): $\operatorname{vadd}(s) \in ((\operatorname{vec}(s) \times \operatorname{vec}(s)) \to \operatorname{vec}(s))$.
\item $\operatorname{vzero}(s) = \operatorname{nth}(4, s)$, a constant (\texttt{constant}): $\operatorname{vzero}(s) \in \operatorname{vec}(s)$.
\item $\operatorname{vneg}(s) = \operatorname{nth}(5, s)$, an operation (\texttt{op}): $\operatorname{vneg}(s) \in (\operatorname{vec}(s) \to \operatorname{vec}(s))$.
\item $\operatorname{act}(s) = \operatorname{nth}(6, s)$, an operation (\texttt{op}): $\operatorname{act}(s) \in ((\operatorname{carr}(\operatorname{scal}(s)) \times \operatorname{vec}(s)) \to \operatorname{vec}(s))$.
\item $\langle\cdot\mid\cdot\rangle_{s} = \operatorname{nth}(7, s)$, an operation (\texttt{op}): $\langle\cdot\mid\cdot\rangle_{s} \in ((\operatorname{vec}(s) \times \operatorname{vec}(s)) \to \mathbb{R})$.
\end{enumerate}
The laws (the \texttt{property} and \texttt{law} clauses):
\begin{itemize}[nosep]
\item $\operatorname{is\text{-}associative}(\operatorname{vadd}(s), \operatorname{vec}(s))$
\item $\operatorname{is\text{-}commutative}(\operatorname{vadd}(s), \operatorname{vec}(s))$
\item $\operatorname{is\text{-}identity}(\operatorname{vadd}(s), \operatorname{vzero}(s), \operatorname{vec}(s))$
\item $\operatorname{has\text{-}inverses}(\operatorname{vadd}(s), \operatorname{vzero}(s), \operatorname{vneg}(s), \operatorname{vec}(s))$
\item $\operatorname{scal}(s) = \operatorname{normed\text{-}field\text{-}as\text{-}commutative\text{-}ring}(\operatorname{rr\text{-}normed\text{-}field})$
\item $\operatorname{carr}(\operatorname{scal}(s)) = \mathbb{R}$
\item
\begin{equation*}
\begin{array}{@{}l@{}} \forall \operatorname{r\_} \in \operatorname{carr}(\operatorname{scal}(s)), \operatorname{x\_}, \operatorname{y\_} \in \operatorname{vec}(s). \\ \quad \operatorname{act}(s)(\operatorname{r\_}, \operatorname{vadd}(s)(\operatorname{x\_}, \operatorname{y\_})) = \operatorname{vadd}(s)(\operatorname{act}(s)(\operatorname{r\_}, \operatorname{x\_}), \operatorname{act}(s)(\operatorname{r\_}, \operatorname{y\_})) \end{array}
\end{equation*}

\item
\begin{equation*}
\begin{array}{@{}l@{}} \forall \operatorname{r\_}, \operatorname{s\_} \in \operatorname{carr}(\operatorname{scal}(s)), \operatorname{x\_} \in \operatorname{vec}(s). \\ \quad \operatorname{act}(s)(\operatorname{add}(\operatorname{scal}(s))(\operatorname{r\_}, \operatorname{s\_}), \operatorname{x\_}) = \operatorname{vadd}(s)(\operatorname{act}(s)(\operatorname{r\_}, \operatorname{x\_}), \operatorname{act}(s)(\operatorname{s\_}, \operatorname{x\_})) \end{array}
\end{equation*}

\item
\begin{equation*}
\begin{array}{@{}l@{}} \forall \operatorname{r\_}, \operatorname{s\_} \in \operatorname{carr}(\operatorname{scal}(s)), \operatorname{x\_} \in \operatorname{vec}(s). \\ \quad \operatorname{act}(s)(\operatorname{mul}(\operatorname{scal}(s))(\operatorname{r\_}, \operatorname{s\_}), \operatorname{x\_}) = \operatorname{act}(s)(\operatorname{r\_}, \operatorname{act}(s)(\operatorname{s\_}, \operatorname{x\_})) \end{array}
\end{equation*}

\item $\forall \operatorname{x\_} \in \operatorname{vec}(s).\; (\operatorname{act}(s)(\operatorname{one}(\operatorname{scal}(s)), \operatorname{x\_}) = \operatorname{x\_})$
\item $\forall \operatorname{x\_}, \operatorname{y\_} \in \operatorname{vec}(s).\; (\langle\cdot\mid\cdot\rangle_{s}(\operatorname{x\_}, \operatorname{y\_}) = \langle\cdot\mid\cdot\rangle_{s}(\operatorname{y\_}, \operatorname{x\_}))$
\item
\begin{equation*}
\begin{array}{@{}l@{}} \forall \operatorname{x\_}, \operatorname{y\_}, \operatorname{z\_} \in \operatorname{vec}(s). \\ \quad \langle\cdot\mid\cdot\rangle_{s}(\operatorname{vadd}(s)(\operatorname{x\_}, \operatorname{y\_}), \operatorname{z\_}) = (\langle\cdot\mid\cdot\rangle_{s}(\operatorname{x\_}, \operatorname{z\_}) + \langle\cdot\mid\cdot\rangle_{s}(\operatorname{y\_}, \operatorname{z\_})) \end{array}
\end{equation*}

\item
\begin{equation*}
\begin{array}{@{}l@{}} \forall \operatorname{a\_} \in \mathbb{R}, \operatorname{x\_}, \operatorname{y\_} \in \operatorname{vec}(s). \\ \quad \langle\cdot\mid\cdot\rangle_{s}(\operatorname{act}(s)(\operatorname{a\_}, \operatorname{x\_}), \operatorname{y\_}) = (\operatorname{a\_} \cdot \langle\cdot\mid\cdot\rangle_{s}(\operatorname{x\_}, \operatorname{y\_})) \end{array}
\end{equation*}

\item $\forall \operatorname{x\_} \in \operatorname{vec}(s).\; (0 \leq \langle\cdot\mid\cdot\rangle_{s}(\operatorname{x\_}, \operatorname{x\_}))$
\end{itemize}
The characteristic law (the axiom \texttt{is-semi-inner-product-space}): $\operatorname{is\text{-}semi\text{-}inner\text{-}product\text{-}space}(s)$ holds if and only if $\operatorname{length}(s) = 7$, every component has the typing listed with it, every law holds, and $s \in \mathbf{Set}$ (the structure is a set).
}
\par\medskip\noindent\textbf{Definition} (\texttt{sip-null}).\;$\operatorname{sip\text{-}null}(v) :=$ $\{\operatorname{sqn\_} \in \operatorname{vec}(v) : (\langle\cdot\mid\cdot\rangle_{v}(\operatorname{sqn\_}, \operatorname{sqn\_}) = 0)\}$
\par\medskip\noindent\textbf{Definition} (\texttt{sip-rel}).\;$\operatorname{sip\text{-}rel}(v) :=$
\begin{equation*}
\{\operatorname{sqp\_} \in (\operatorname{vec}(v) \times \operatorname{vec}(v)) : (\langle\cdot\mid\cdot\rangle_{v}(\operatorname{vadd}(v)(\operatorname{nth}(1, \operatorname{sqp\_}), \operatorname{vneg}(v)(\operatorname{nth}(2, \operatorname{sqp\_}))), \operatorname{vadd}(v)(\operatorname{nth}(1, \operatorname{sqp\_}), \operatorname{vneg}(v)(\operatorname{nth}(2, \operatorname{sqp\_})))) = 0)\}
\end{equation*}

\par\medskip\noindent\textbf{Definition} (\texttt{sip-setoid}).\;$\operatorname{sip\text{-}setoid}(v) :=$ ${[}\operatorname{vec}(v), \operatorname{sip\text{-}rel}(v)]$
\par\medskip\noindent\textbf{Definition} (\texttt{sip-qvadd}).\;$\operatorname{sip\text{-}qvadd}(v) :=$
\begin{equation*}
\begin{array}{@{}l@{}} \operatorname{descend2}(\operatorname{sip\text{-}setoid}(v), \\ \quad (\lambda {[}\operatorname{sqa\_}, \operatorname{sqb\_}] \in (\operatorname{vec}(v) \times \operatorname{vec}(v)).\; \operatorname{class}(\operatorname{sip\text{-}setoid}(v), \operatorname{vadd}(v)(\operatorname{sqa\_}, \operatorname{sqb\_})))) \end{array}
\end{equation*}

\par\medskip\noindent\textbf{Definition} (\texttt{sip-qvneg}).\;$\operatorname{sip\text{-}qvneg}(v) :=$
\begin{equation*}
\operatorname{descend}(\operatorname{sip\text{-}setoid}(v), (\lambda \operatorname{sqa\_} \in \operatorname{vec}(v).\; \operatorname{class}(\operatorname{sip\text{-}setoid}(v), \operatorname{vneg}(v)(\operatorname{sqa\_}))))
\end{equation*}

\par\medskip\noindent\textbf{Definition} (\texttt{sip-qscale}).\;$\operatorname{sip\text{-}qscale}(v, r) :=$
\begin{equation*}
\operatorname{descend}(\operatorname{sip\text{-}setoid}(v), (\lambda \operatorname{sqa\_} \in \operatorname{vec}(v).\; \operatorname{class}(\operatorname{sip\text{-}setoid}(v), \operatorname{act}(v)(r, \operatorname{sqa\_}))))
\end{equation*}

\par\medskip\noindent\textbf{Definition} (\texttt{sip-qact}).\;$\operatorname{sip\text{-}qact}(v) :=$
\begin{equation*}
(\lambda {[}\operatorname{sqr\_}, \operatorname{sqc\_}] \in (\mathbb{R} \times \operatorname{quotient}(\operatorname{sip\text{-}setoid}(v))).\; \operatorname{sip\text{-}qscale}(v, \operatorname{sqr\_})(\operatorname{sqc\_}))
\end{equation*}

\par\medskip\noindent\textbf{Definition} (\texttt{sip-qrip}).\;$\operatorname{sip\text{-}qrip}(v) :=$
\begin{equation*}
\operatorname{descend2}(\operatorname{sip\text{-}setoid}(v), (\lambda {[}\operatorname{sqa\_}, \operatorname{sqb\_}] \in (\operatorname{vec}(v) \times \operatorname{vec}(v)).\; \langle\cdot\mid\cdot\rangle_{v}(\operatorname{sqa\_}, \operatorname{sqb\_})))
\end{equation*}

\par\medskip\noindent\textbf{Definition} (\texttt{sip-quotient}).\;$\operatorname{sip\text{-}quotient}(v) :=$
\begin{equation*}
{[}\operatorname{scal}(v), \operatorname{quotient}(\operatorname{sip\text{-}setoid}(v)), \operatorname{sip\text{-}qvadd}(v), \operatorname{class}(\operatorname{sip\text{-}setoid}(v), \operatorname{vzero}(v)), \operatorname{sip\text{-}qvneg}(v), \operatorname{sip\text{-}qact}(v), \operatorname{sip\text{-}qrip}(v)]
\end{equation*}

\par\medskip\noindent\textbf{Definition} (\texttt{hom-pre-hilbert-space}).\;$\operatorname{hom\text{-}pre\text{-}hilbert\text{-}space}(a, b) :=$
\begin{equation*}
\{\operatorname{homf\_} \in (\operatorname{vec}(a) \to \operatorname{vec}(b)) : \operatorname{is\text{-}hom\text{-}pre\text{-}hilbert\text{-}space}(a, b, \operatorname{homf\_})\}
\end{equation*}

\par\medskip\noindent\textbf{Definition} (\texttt{hom-semi-inner-product-space}).\;$\operatorname{hom\text{-}semi\text{-}inner\text{-}product\text{-}space}(a, b) :=$
\begin{equation*}
\{\operatorname{homf\_} \in (\operatorname{vec}(a) \to \operatorname{vec}(b)) : \operatorname{is\text{-}hom\text{-}semi\text{-}inner\text{-}product\text{-}space}(a, b, \operatorname{homf\_})\}
\end{equation*}

\vnbdef{is-pre-hilbert-space-def}{\vnbstmt{\begin{array}{@{}l@{}} \forall s. \\ \quad \operatorname{is\text{-}pre\text{-}hilbert\text{-}space}(s) \\ \quad \Leftrightarrow\; (\operatorname{is\text{-}semi\text{-}inner\text{-}product\text{-}space}(s) \wedge \forall x \in \operatorname{vec}(s).\; ((\langle\cdot\mid\cdot\rangle_{s}(x, x) = 0) \Rightarrow (x = \operatorname{vzero}(s)))) \end{array}}}
\vnbdef{pre-hilbert-space-class}{\vnbstmt{\forall s.\; ((s \in \operatorname{pre\text{-}hilbert\text{-}space}) \Leftrightarrow \operatorname{is\text{-}pre\text{-}hilbert\text{-}space}(s))}}
\vnbdef{semi-inner-product-space-class}{\vnbstmt{\forall s.\; ((s \in \operatorname{semi\text{-}inner\text{-}product\text{-}space}) \Leftrightarrow \operatorname{is\text{-}semi\text{-}inner\text{-}product\text{-}space}(s))}}
\vnbdef{is-hom-pre-hilbert-space-def}{\vnbstmt{\begin{array}{@{}l@{}} \forall a, b, f. \\ \quad \operatorname{is\text{-}hom\text{-}pre\text{-}hilbert\text{-}space}(a, b, f) \\ \quad \Leftrightarrow\; (\operatorname{is\text{-}pre\text{-}hilbert\text{-}space}(a) \wedge (\operatorname{is\text{-}pre\text{-}hilbert\text{-}space}(b) \wedge \operatorname{is\text{-}hom\text{-}semi\text{-}inner\text{-}product\text{-}space}(a, b, f))) \end{array}}}
\vnbdef{is-hom-semi-inner-product-space-def}{\vnbstmt{\begin{array}{@{}l@{}} \forall a, b, f. \\ \quad \operatorname{is\text{-}hom\text{-}semi\text{-}inner\text{-}product\text{-}space}(a, b, f) \\ \quad \Leftrightarrow\; \operatorname{is\text{-}semi\text{-}inner\text{-}product\text{-}space}(a) \wedge \\ \quad \operatorname{is\text{-}semi\text{-}inner\text{-}product\text{-}space}(b) \wedge \\ \quad f \in (\operatorname{vec}(a) \to \operatorname{vec}(b)) \wedge \\ \quad \operatorname{scal}(a) = \operatorname{scal}(b) \wedge \\ \quad \forall x_{1}, x_{2} \in \operatorname{vec}(a). \\ \quad \quad f(\operatorname{vadd}(a)(x_{1}, x_{2})) = \operatorname{vadd}(b)(f(x_{1}), f(x_{2})) \wedge \\ \quad f(\operatorname{vzero}(a)) = \operatorname{vzero}(b) \wedge \\ \quad \forall x_{1} \in \operatorname{vec}(a).\; (f(\operatorname{vneg}(a)(x_{1})) = \operatorname{vneg}(b)(f(x_{1}))) \wedge \\ \quad \forall x_{1} \in \operatorname{carr}(\operatorname{scal}(a)), x_{2} \in \operatorname{vec}(a). \\ \quad \quad f(\operatorname{act}(a)(x_{1}, x_{2})) = \operatorname{act}(b)(x_{1}, f(x_{2})) \wedge \\ \quad \forall x_{1}, x_{2} \in \operatorname{vec}(a).\; (\langle\cdot\mid\cdot\rangle_{a}(x_{1}, x_{2}) = \langle\cdot\mid\cdot\rangle_{b}(f(x_{1}), f(x_{2}))) \end{array}}}
\vnbprop{is-semi-inner-product-space-is-set}{\vnbstmt{\forall s.\; (\operatorname{is\text{-}semi\text{-}inner\text{-}product\text{-}space}(s) \Rightarrow (s \in \mathbf{Set}))}}
\vnbprop{is-pre-hilbert-space-is-set}{\vnbstmt{\forall s.\; (\operatorname{is\text{-}pre\text{-}hilbert\text{-}space}(s) \Rightarrow (s \in \mathbf{Set}))}}

\vnbbookfile{theorem-library/sip-a}{theorem-library/sip-a}
\noindent{\small sip-a.scm -- SEMI-INNER-PRODUCT SPACES: the laws, Cauchy-Schwarz for a semidefinite form, the null vectors, the setoid, and the quotient pre-Hilbert space. (SIP-1, 2026-10-18; STOP POINT A of scratchpad/triage/SIP-1.md.)  All PROVEN modulo 0.}\par
\vnbprop{sip-as-module-is-module}{\vnbstmt{\forall r.\; (\operatorname{is\text{-}semi\text{-}inner\text{-}product\text{-}space}(r) \Rightarrow \operatorname{is\text{-}module}(\operatorname{sip\text{-}as\text{-}module}(r)))}}
\vnbprop{sip-scal-pin}{\vnbstmt{\begin{array}{@{}l@{}} \forall v. \\ \quad \operatorname{is\text{-}semi\text{-}inner\text{-}product\text{-}space}(v) \Rightarrow  \\ \quad \quad \operatorname{scal}(v) = \operatorname{normed\text{-}field\text{-}as\text{-}commutative\text{-}ring}(\operatorname{rr\text{-}normed\text{-}field}) \end{array}}}
\vnbprop{sip-scal-carr}{\vnbstmt{\forall v.\; (\operatorname{is\text{-}semi\text{-}inner\text{-}product\text{-}space}(v) \Rightarrow (\operatorname{carr}(\operatorname{scal}(v)) = \mathbb{R}))}}
\vnbprop{sip-vec-set}{\vnbstmt{\forall v.\; (\operatorname{is\text{-}semi\text{-}inner\text{-}product\text{-}space}(v) \Rightarrow (\operatorname{vec}(v) \in \mathbf{Set}))}}
\vnbprop{sip-rip-fun}{\vnbstmt{\begin{array}{@{}l@{}} \forall v. \\ \quad \operatorname{is\text{-}semi\text{-}inner\text{-}product\text{-}space}(v) \Rightarrow  \\ \quad \quad \langle\cdot\mid\cdot\rangle_{v} \in ((\operatorname{vec}(v) \times \operatorname{vec}(v)) \to \mathbb{R}) \end{array}}}
\vnbprop{sip-rip-sym}{\vnbstmt{\begin{array}{@{}l@{}} \forall v.; \forall x, y \in \operatorname{vec}(v). \\ \quad \operatorname{is\text{-}semi\text{-}inner\text{-}product\text{-}space}(v) \Rightarrow (\langle\cdot\mid\cdot\rangle_{v}(x, y) = \langle\cdot\mid\cdot\rangle_{v}(y, x)) \end{array}}}
\vnbprop{sip-rip-add-left}{\vnbstmt{\begin{array}{@{}l@{}} \forall v.; \forall x, y, z \in \operatorname{vec}(v). \\ \quad \operatorname{is\text{-}semi\text{-}inner\text{-}product\text{-}space}(v) \Rightarrow  \\ \quad \quad \langle\cdot\mid\cdot\rangle_{v}(\operatorname{vadd}(v)(x, y), z) = (\langle\cdot\mid\cdot\rangle_{v}(x, z) + \langle\cdot\mid\cdot\rangle_{v}(y, z)) \end{array}}}
\vnbprop{sip-rip-homog-left}{\vnbstmt{\begin{array}{@{}l@{}} \forall v.; \forall a \in \mathbb{R}, x, y \in \operatorname{vec}(v). \\ \quad \operatorname{is\text{-}semi\text{-}inner\text{-}product\text{-}space}(v) \Rightarrow  \\ \quad \quad \langle\cdot\mid\cdot\rangle_{v}(\operatorname{act}(v)(a, x), y) = (a \cdot \langle\cdot\mid\cdot\rangle_{v}(x, y)) \end{array}}}
\vnbprop{sip-rip-self-nonneg}{\vnbstmt{\begin{array}{@{}l@{}} \forall v.; \forall x \in \operatorname{vec}(v). \\ \quad \operatorname{is\text{-}semi\text{-}inner\text{-}product\text{-}space}(v) \Rightarrow (0 \leq \langle\cdot\mid\cdot\rangle_{v}(x, x)) \end{array}}}
\vnbprop{sip-rip-type}{\vnbstmt{\begin{array}{@{}l@{}} \forall v.; \forall x, y \in \operatorname{vec}(v). \\ \quad \operatorname{is\text{-}semi\text{-}inner\text{-}product\text{-}space}(v) \Rightarrow (\langle\cdot\mid\cdot\rangle_{v}(x, y) \in \mathbb{R}) \end{array}}}
\vnbprop{sip-rip-add-right}{\vnbstmt{\begin{array}{@{}l@{}} \forall v.; \forall x, y, z \in \operatorname{vec}(v). \\ \quad \operatorname{is\text{-}semi\text{-}inner\text{-}product\text{-}space}(v) \Rightarrow  \\ \quad \quad \langle\cdot\mid\cdot\rangle_{v}(x, \operatorname{vadd}(v)(y, z)) = (\langle\cdot\mid\cdot\rangle_{v}(x, y) + \langle\cdot\mid\cdot\rangle_{v}(x, z)) \end{array}}}
\vnbprop{sip-rip-homog-right}{\vnbstmt{\begin{array}{@{}l@{}} \forall v.; \forall a \in \mathbb{R}, x, y \in \operatorname{vec}(v). \\ \quad \operatorname{is\text{-}semi\text{-}inner\text{-}product\text{-}space}(v) \Rightarrow  \\ \quad \quad \langle\cdot\mid\cdot\rangle_{v}(x, \operatorname{act}(v)(a, y)) = (a \cdot \langle\cdot\mid\cdot\rangle_{v}(x, y)) \end{array}}}
\vnbprop{sip-rip-zero-left}{\vnbstmt{\begin{array}{@{}l@{}} \forall v.; \forall y \in \operatorname{vec}(v). \\ \quad \operatorname{is\text{-}semi\text{-}inner\text{-}product\text{-}space}(v) \Rightarrow (\langle\cdot\mid\cdot\rangle_{v}(\operatorname{vzero}(v), y) = 0) \end{array}}}
\vnbprop{sip-rip-zero-right}{\vnbstmt{\begin{array}{@{}l@{}} \forall v.; \forall x \in \operatorname{vec}(v). \\ \quad \operatorname{is\text{-}semi\text{-}inner\text{-}product\text{-}space}(v) \Rightarrow (\langle\cdot\mid\cdot\rangle_{v}(x, \operatorname{vzero}(v)) = 0) \end{array}}}
\vnbprop{sip-rip-neg-left}{\vnbstmt{\begin{array}{@{}l@{}} \forall v.; \forall x, y \in \operatorname{vec}(v). \\ \quad \operatorname{is\text{-}semi\text{-}inner\text{-}product\text{-}space}(v) \Rightarrow  \\ \quad \quad \langle\cdot\mid\cdot\rangle_{v}(\operatorname{vneg}(v)(x), y) = -\langle\cdot\mid\cdot\rangle_{v}(x, y) \end{array}}}
\vnbprop{sip-rip-neg-right}{\vnbstmt{\begin{array}{@{}l@{}} \forall v.; \forall x, y \in \operatorname{vec}(v). \\ \quad \operatorname{is\text{-}semi\text{-}inner\text{-}product\text{-}space}(v) \Rightarrow  \\ \quad \quad \langle\cdot\mid\cdot\rangle_{v}(x, \operatorname{vneg}(v)(y)) = -\langle\cdot\mid\cdot\rangle_{v}(x, y) \end{array}}}
\vnbprop{sip-expand-two}{\vnbstmt{\begin{array}{@{}l@{}} \forall v.; \forall x, y \in \operatorname{vec}(v), a, b \in \mathbb{R}. \\ \quad \operatorname{is\text{-}semi\text{-}inner\text{-}product\text{-}space}(v) \Rightarrow  \\ \quad \quad \langle\cdot\mid\cdot\rangle_{v}(\operatorname{vadd}(v)(\operatorname{act}(v)(a, x), \operatorname{act}(v)(b, y)), \operatorname{vadd}(v)(\operatorname{act}(v)(a, x), \operatorname{act}(v)(b, y))) \\ \quad \quad =\; (((a \cdot a) \cdot \langle\cdot\mid\cdot\rangle_{v}(x, x)) + ((2 \cdot (a \cdot b)) \cdot \langle\cdot\mid\cdot\rangle_{v}(x, y))) + ((b \cdot b) \cdot \langle\cdot\mid\cdot\rangle_{v}(y, y)) \end{array}}}
\vnbprop{sip-schwarz}{\vnbstmt{\begin{array}{@{}l@{}} \forall v.; \forall x, y \in \operatorname{vec}(v). \\ \quad \operatorname{is\text{-}semi\text{-}inner\text{-}product\text{-}space}(v) \Rightarrow  \\ \quad \quad (\langle\cdot\mid\cdot\rangle_{v}(x, y) \cdot \langle\cdot\mid\cdot\rangle_{v}(x, y)) \leq (\langle\cdot\mid\cdot\rangle_{v}(x, x) \cdot \langle\cdot\mid\cdot\rangle_{v}(y, y)) \end{array}}}
\vnbprop{sip-null-orth}{\vnbstmt{\begin{array}{@{}l@{}} \forall v.; \forall x, y \in \operatorname{vec}(v). \\ \quad (\operatorname{is\text{-}semi\text{-}inner\text{-}product\text{-}space}(v) \wedge (\langle\cdot\mid\cdot\rangle_{v}(x, x) = 0)) \Rightarrow (\langle\cdot\mid\cdot\rangle_{v}(x, y) = 0) \end{array}}}
\vnbprop{sip-null-vadd}{\vnbstmt{\begin{array}{@{}l@{}} \forall v.; \forall x, y \in \operatorname{vec}(v). \\ \quad (\operatorname{is\text{-}semi\text{-}inner\text{-}product\text{-}space}(v) \wedge ((\langle\cdot\mid\cdot\rangle_{v}(x, x) = 0) \wedge (\langle\cdot\mid\cdot\rangle_{v}(y, y) = 0))) \Rightarrow  \\ \quad \quad \langle\cdot\mid\cdot\rangle_{v}(\operatorname{vadd}(v)(x, y), \operatorname{vadd}(v)(x, y)) = 0 \end{array}}}
\vnbprop{sip-null-vneg}{\vnbstmt{\begin{array}{@{}l@{}} \forall v.; \forall x \in \operatorname{vec}(v). \\ \quad (\operatorname{is\text{-}semi\text{-}inner\text{-}product\text{-}space}(v) \wedge (\langle\cdot\mid\cdot\rangle_{v}(x, x) = 0)) \Rightarrow  \\ \quad \quad \langle\cdot\mid\cdot\rangle_{v}(\operatorname{vneg}(v)(x), \operatorname{vneg}(v)(x)) = 0 \end{array}}}
\vnbprop{sip-null-act}{\vnbstmt{\begin{array}{@{}l@{}} \forall v.; \forall a \in \mathbb{R}, x \in \operatorname{vec}(v). \\ \quad (\operatorname{is\text{-}semi\text{-}inner\text{-}product\text{-}space}(v) \wedge (\langle\cdot\mid\cdot\rangle_{v}(x, x) = 0)) \Rightarrow  \\ \quad \quad \langle\cdot\mid\cdot\rangle_{v}(\operatorname{act}(v)(a, x), \operatorname{act}(v)(a, x)) = 0 \end{array}}}
\vnbprop{sip-null-unfold}{\vnbstmt{\forall v.\; (\operatorname{sip\text{-}null}(v) \simeq \{n \in \operatorname{vec}(v) : (\langle\cdot\mid\cdot\rangle_{v}(n, n) = 0)\})}}
\vnbprop{sip-null-in}{\vnbstmt{\begin{array}{@{}l@{}} \forall v.; \forall x \in \operatorname{vec}(v). \\ \quad (\operatorname{is\text{-}semi\text{-}inner\text{-}product\text{-}space}(v) \wedge (\langle\cdot\mid\cdot\rangle_{v}(x, x) = 0)) \Rightarrow (x \in \operatorname{sip\text{-}null}(v)) \end{array}}}
\vnbprop{sip-null-out}{\vnbstmt{\begin{array}{@{}l@{}} \forall v.; \forall x \in \operatorname{sip\text{-}null}(v). \\ \quad \operatorname{is\text{-}semi\text{-}inner\text{-}product\text{-}space}(v) \Rightarrow ((x \in \operatorname{vec}(v)) \wedge (\langle\cdot\mid\cdot\rangle_{v}(x, x) = 0)) \end{array}}}
\vnbprop{sip-null-subspace}{\vnbstmt{\forall v.\; (\operatorname{is\text{-}semi\text{-}inner\text{-}product\text{-}space}(v) \Rightarrow \operatorname{is\text{-}subspace}(v, \operatorname{sip\text{-}null}(v)))}}
\vnbprop{sip-rel-unfold}{\vnbstmt{\begin{array}{@{}l@{}} \forall v. \\ \quad (\operatorname{sip\text{-}rel}(v) \simeq \{p \in (\operatorname{vec}(v) \times \operatorname{vec}(v)) : (\langle\cdot\mid\cdot\rangle_{v}(\operatorname{vadd}(v)(\operatorname{nth}(1, p), \operatorname{vneg}(v)(\operatorname{nth}(2, p))), \operatorname{vadd}(v)(\operatorname{nth}(1, p), \operatorname{vneg}(v)(\operatorname{nth}(2, p)))) = 0)\}) \end{array}}}
\vnbprop{sip-rel-in}{\vnbstmt{\begin{array}{@{}l@{}} \forall v.; \forall a, b \in \operatorname{vec}(v). \\ \quad \operatorname{is\text{-}semi\text{-}inner\text{-}product\text{-}space}(v) \wedge \\ \quad \langle\cdot\mid\cdot\rangle_{v}(\operatorname{vadd}(v)(a, \operatorname{vneg}(v)(b)), \operatorname{vadd}(v)(a, \operatorname{vneg}(v)(b))) = 0 \Rightarrow \\ \quad \quad {[}a, b] \in \operatorname{sip\text{-}rel}(v) \end{array}}}
\vnbprop{sip-rel-out}{\vnbstmt{\begin{array}{@{}l@{}} \forall v, a, b. \\ \quad (\operatorname{is\text{-}semi\text{-}inner\text{-}product\text{-}space}(v) \wedge ({[}a, b] \in \operatorname{sip\text{-}rel}(v))) \Rightarrow  \\ \quad \quad \langle\cdot\mid\cdot\rangle_{v}(\operatorname{vadd}(v)(a, \operatorname{vneg}(v)(b)), \operatorname{vadd}(v)(a, \operatorname{vneg}(v)(b))) = 0 \end{array}}}
\vnbprop{sip-rel-is-equivalence}{\vnbstmt{\begin{array}{@{}l@{}} \forall v. \\ \quad \operatorname{is\text{-}semi\text{-}inner\text{-}product\text{-}space}(v) \Rightarrow \operatorname{is\text{-}equivalence}(\operatorname{sip\text{-}rel}(v), \operatorname{vec}(v)) \end{array}}}
\vnbprop{sip-setoid-is-setoid}{\vnbstmt{\forall v.\; (\operatorname{is\text{-}semi\text{-}inner\text{-}product\text{-}space}(v) \Rightarrow \operatorname{is\text{-}setoid}(\operatorname{sip\text{-}setoid}(v)))}}
\vnbprop{sip-setoid-pts}{\vnbstmt{\forall v.\; (\operatorname{is\text{-}semi\text{-}inner\text{-}product\text{-}space}(v) \Rightarrow (\operatorname{pts}(\operatorname{sip\text{-}setoid}(v)) \simeq \operatorname{vec}(v)))}}
\vnbprop{sip-setoid-rel}{\vnbstmt{\begin{array}{@{}l@{}} \forall v. \\ \quad \operatorname{is\text{-}semi\text{-}inner\text{-}product\text{-}space}(v) \Rightarrow (\operatorname{rel}(\operatorname{sip\text{-}setoid}(v)) \simeq \operatorname{sip\text{-}rel}(v)) \end{array}}}
\vnbprop{sip-related-out}{\vnbstmt{\begin{array}{@{}l@{}} \forall v, a, b. \\ \quad (\operatorname{is\text{-}semi\text{-}inner\text{-}product\text{-}space}(v) \wedge \operatorname{related}(\operatorname{sip\text{-}setoid}(v), a, b)) \Rightarrow  \\ \quad \quad \langle\cdot\mid\cdot\rangle_{v}(\operatorname{vadd}(v)(a, \operatorname{vneg}(v)(b)), \operatorname{vadd}(v)(a, \operatorname{vneg}(v)(b))) = 0 \end{array}}}
\vnbprop{sip-related-in}{\vnbstmt{\begin{array}{@{}l@{}} \forall v.; \forall a, b \in \operatorname{vec}(v). \\ \quad \operatorname{is\text{-}semi\text{-}inner\text{-}product\text{-}space}(v) \wedge \\ \quad \langle\cdot\mid\cdot\rangle_{v}(\operatorname{vadd}(v)(a, \operatorname{vneg}(v)(b)), \operatorname{vadd}(v)(a, \operatorname{vneg}(v)(b))) = 0 \Rightarrow \\ \quad \quad \operatorname{related}(\operatorname{sip\text{-}setoid}(v), a, b) \end{array}}}
\vnbprop{sip-class-eq-iff}{\vnbstmt{\begin{array}{@{}l@{}} \forall v.; \forall a, b \in \operatorname{vec}(v). \\ \quad \operatorname{is\text{-}semi\text{-}inner\text{-}product\text{-}space}(v) \Rightarrow  \\ \quad \quad \operatorname{class}(\operatorname{sip\text{-}setoid}(v), a) = \operatorname{class}(\operatorname{sip\text{-}setoid}(v), b) \\ \quad \quad \Leftrightarrow\; \langle\cdot\mid\cdot\rangle_{v}(\operatorname{vadd}(v)(a, \operatorname{vneg}(v)(b)), \operatorname{vadd}(v)(a, \operatorname{vneg}(v)(b))) = 0 \end{array}}}
\vnbprop{sipq-scal}{\vnbstmt{\begin{array}{@{}l@{}} \forall v. \\ \quad \operatorname{is\text{-}semi\text{-}inner\text{-}product\text{-}space}(v) \Rightarrow (\operatorname{scal}(\operatorname{sip\text{-}quotient}(v)) \simeq \operatorname{scal}(v)) \end{array}}}
\vnbprop{sipq-vec}{\vnbstmt{\begin{array}{@{}l@{}} \forall v. \\ \quad \operatorname{is\text{-}semi\text{-}inner\text{-}product\text{-}space}(v) \Rightarrow  \\ \quad \quad (\operatorname{vec}(\operatorname{sip\text{-}quotient}(v)) \simeq \operatorname{quotient}(\operatorname{sip\text{-}setoid}(v))) \end{array}}}
\vnbprop{sipq-vadd}{\vnbstmt{\begin{array}{@{}l@{}} \forall v. \\ \quad \operatorname{is\text{-}semi\text{-}inner\text{-}product\text{-}space}(v) \Rightarrow (\operatorname{vadd}(\operatorname{sip\text{-}quotient}(v)) \simeq \operatorname{sip\text{-}qvadd}(v)) \end{array}}}
\vnbprop{sipq-vzero}{\vnbstmt{\begin{array}{@{}l@{}} \forall v. \\ \quad \operatorname{is\text{-}semi\text{-}inner\text{-}product\text{-}space}(v) \Rightarrow  \\ \quad \quad (\operatorname{vzero}(\operatorname{sip\text{-}quotient}(v)) \simeq \operatorname{class}(\operatorname{sip\text{-}setoid}(v), \operatorname{vzero}(v))) \end{array}}}
\vnbprop{sipq-vneg}{\vnbstmt{\begin{array}{@{}l@{}} \forall v. \\ \quad \operatorname{is\text{-}semi\text{-}inner\text{-}product\text{-}space}(v) \Rightarrow (\operatorname{vneg}(\operatorname{sip\text{-}quotient}(v)) \simeq \operatorname{sip\text{-}qvneg}(v)) \end{array}}}
\vnbprop{sipq-act}{\vnbstmt{\begin{array}{@{}l@{}} \forall v. \\ \quad \operatorname{is\text{-}semi\text{-}inner\text{-}product\text{-}space}(v) \Rightarrow (\operatorname{act}(\operatorname{sip\text{-}quotient}(v)) \simeq \operatorname{sip\text{-}qact}(v)) \end{array}}}
\vnbprop{sipq-rip}{\vnbstmt{\begin{array}{@{}l@{}} \forall v. \\ \quad \operatorname{is\text{-}semi\text{-}inner\text{-}product\text{-}space}(v) \Rightarrow (\langle\cdot\mid\cdot\rangle_{\operatorname{sip\text{-}quotient}(v)} \simeq \operatorname{sip\text{-}qrip}(v)) \end{array}}}
\vnbprop{sipq-lva-fun}{\vnbstmt{\begin{array}{@{}l@{}} \forall v. \\ \quad \operatorname{is\text{-}semi\text{-}inner\text{-}product\text{-}space}(v) \Rightarrow  \\ \quad \quad (\lambda {[}\operatorname{sqa\_}, \operatorname{sqb\_}] \in (\operatorname{vec}(v) \times \operatorname{vec}(v)).\; \operatorname{class}(\operatorname{sip\text{-}setoid}(v), \operatorname{vadd}(v)(\operatorname{sqa\_}, \operatorname{sqb\_}))) \\ \quad \quad \in\; ((\operatorname{pts}(\operatorname{sip\text{-}setoid}(v)) \times \operatorname{pts}(\operatorname{sip\text{-}setoid}(v))) \to \operatorname{quotient}(\operatorname{sip\text{-}setoid}(v))) \end{array}}}
\vnbprop{sipq-lvn-fun}{\vnbstmt{\begin{array}{@{}l@{}} \forall v. \\ \quad \operatorname{is\text{-}semi\text{-}inner\text{-}product\text{-}space}(v) \Rightarrow  \\ \quad \quad (\lambda a \in \operatorname{vec}(v).\; \operatorname{class}(\operatorname{sip\text{-}setoid}(v), \operatorname{vneg}(v)(a))) \\ \quad \quad \in\; (\operatorname{pts}(\operatorname{sip\text{-}setoid}(v)) \to \operatorname{quotient}(\operatorname{sip\text{-}setoid}(v))) \end{array}}}
\vnbprop{sipq-lrp-fun}{\vnbstmt{\begin{array}{@{}l@{}} \forall v. \\ \quad \operatorname{is\text{-}semi\text{-}inner\text{-}product\text{-}space}(v) \Rightarrow  \\ \quad \quad (\lambda {[}\operatorname{sqa\_}, \operatorname{sqb\_}] \in (\operatorname{vec}(v) \times \operatorname{vec}(v)).\; \langle\cdot\mid\cdot\rangle_{v}(\operatorname{sqa\_}, \operatorname{sqb\_})) \\ \quad \quad \in\; ((\operatorname{pts}(\operatorname{sip\text{-}setoid}(v)) \times \operatorname{pts}(\operatorname{sip\text{-}setoid}(v))) \to \mathbb{R}) \end{array}}}
\vnbprop{sipq-lsc-fun}{\vnbstmt{\begin{array}{@{}l@{}} \forall v.; \forall r \in \mathbb{R}. \\ \quad \operatorname{is\text{-}semi\text{-}inner\text{-}product\text{-}space}(v) \Rightarrow  \\ \quad \quad (\lambda a \in \operatorname{vec}(v).\; \operatorname{class}(\operatorname{sip\text{-}setoid}(v), \operatorname{act}(v)(r, a))) \\ \quad \quad \in\; (\operatorname{pts}(\operatorname{sip\text{-}setoid}(v)) \to \operatorname{quotient}(\operatorname{sip\text{-}setoid}(v))) \end{array}}}
\vnbprop{sip-rip-respects-left}{\vnbstmt{\begin{array}{@{}l@{}} \forall v.; \forall x, u, y \in \operatorname{vec}(v). \\ \quad \operatorname{is\text{-}semi\text{-}inner\text{-}product\text{-}space}(v) \wedge \\ \quad \langle\cdot\mid\cdot\rangle_{v}(\operatorname{vadd}(v)(x, \operatorname{vneg}(v)(u)), \operatorname{vadd}(v)(x, \operatorname{vneg}(v)(u))) = 0 \Rightarrow \\ \quad \quad \langle\cdot\mid\cdot\rangle_{v}(x, y) = \langle\cdot\mid\cdot\rangle_{v}(u, y) \end{array}}}
\vnbprop{sipq-lva-respects}{\vnbstmt{\begin{array}{@{}l@{}} \forall v. \\ \quad \operatorname{is\text{-}semi\text{-}inner\text{-}product\text{-}space}(v) \Rightarrow  \\ \quad \quad \operatorname{respects2}(\operatorname{sip\text{-}setoid}(v), \\ \quad \quad \quad (\lambda {[}\operatorname{sqa\_}, \operatorname{sqb\_}] \in (\operatorname{vec}(v) \times \operatorname{vec}(v)).\; \operatorname{class}(\operatorname{sip\text{-}setoid}(v), \operatorname{vadd}(v)(\operatorname{sqa\_}, \operatorname{sqb\_})))) \end{array}}}
\vnbprop{sipq-lvn-respects}{\vnbstmt{\begin{array}{@{}l@{}} \forall v. \\ \quad \operatorname{is\text{-}semi\text{-}inner\text{-}product\text{-}space}(v) \Rightarrow  \\ \quad \quad \operatorname{respects}(\operatorname{sip\text{-}setoid}(v), (\lambda a \in \operatorname{vec}(v).\; \operatorname{class}(\operatorname{sip\text{-}setoid}(v), \operatorname{vneg}(v)(a)))) \end{array}}}
\vnbprop{sipq-lsc-respects}{\vnbstmt{\begin{array}{@{}l@{}} \forall v.; \forall r \in \mathbb{R}. \\ \quad \operatorname{is\text{-}semi\text{-}inner\text{-}product\text{-}space}(v) \Rightarrow  \\ \quad \quad \operatorname{respects}(\operatorname{sip\text{-}setoid}(v), (\lambda a \in \operatorname{vec}(v).\; \operatorname{class}(\operatorname{sip\text{-}setoid}(v), \operatorname{act}(v)(r, a)))) \end{array}}}
\vnbprop{sipq-lrp-respects}{\vnbstmt{\begin{array}{@{}l@{}} \forall v. \\ \quad \operatorname{is\text{-}semi\text{-}inner\text{-}product\text{-}space}(v) \Rightarrow  \\ \quad \quad \operatorname{respects2}(\operatorname{sip\text{-}setoid}(v), (\lambda {[}\operatorname{sqa\_}, \operatorname{sqb\_}] \in (\operatorname{vec}(v) \times \operatorname{vec}(v)).\; \langle\cdot\mid\cdot\rangle_{v}(\operatorname{sqa\_}, \operatorname{sqb\_}))) \end{array}}}
\vnbprop{sipq-vadd-computes}{\vnbstmt{\begin{array}{@{}l@{}} \forall v.; \forall a, b \in \operatorname{vec}(v). \\ \quad \operatorname{is\text{-}semi\text{-}inner\text{-}product\text{-}space}(v) \Rightarrow  \\ \quad \quad \operatorname{vadd}(\operatorname{sip\text{-}quotient}(v))(\operatorname{class}(\operatorname{sip\text{-}setoid}(v), a), \operatorname{class}(\operatorname{sip\text{-}setoid}(v), b)) \\ \quad \quad =\; \operatorname{class}(\operatorname{sip\text{-}setoid}(v), \operatorname{vadd}(v)(a, b)) \end{array}}}
\vnbprop{sipq-vneg-computes}{\vnbstmt{\begin{array}{@{}l@{}} \forall v.; \forall a \in \operatorname{vec}(v). \\ \quad \operatorname{is\text{-}semi\text{-}inner\text{-}product\text{-}space}(v) \Rightarrow  \\ \quad \quad \operatorname{vneg}(\operatorname{sip\text{-}quotient}(v))(\operatorname{class}(\operatorname{sip\text{-}setoid}(v), a)) = \operatorname{class}(\operatorname{sip\text{-}setoid}(v), \operatorname{vneg}(v)(a)) \end{array}}}
\vnbprop{sipq-scale-computes}{\vnbstmt{\begin{array}{@{}l@{}} \forall v.; \forall r \in \mathbb{R}, a \in \operatorname{vec}(v). \\ \quad \operatorname{is\text{-}semi\text{-}inner\text{-}product\text{-}space}(v) \Rightarrow  \\ \quad \quad \operatorname{sip\text{-}qscale}(v, r)(\operatorname{class}(\operatorname{sip\text{-}setoid}(v), a)) = \operatorname{class}(\operatorname{sip\text{-}setoid}(v), \operatorname{act}(v)(r, a)) \end{array}}}
\vnbprop{sipq-rip-computes}{\vnbstmt{\begin{array}{@{}l@{}} \forall v.; \forall a, b \in \operatorname{vec}(v). \\ \quad \operatorname{is\text{-}semi\text{-}inner\text{-}product\text{-}space}(v) \Rightarrow  \\ \quad \quad \langle\cdot\mid\cdot\rangle_{\operatorname{sip\text{-}quotient}(v)}(\operatorname{class}(\operatorname{sip\text{-}setoid}(v), a), \operatorname{class}(\operatorname{sip\text{-}setoid}(v), b)) = \langle\cdot\mid\cdot\rangle_{v}(a, b) \end{array}}}
\vnbprop{sipq-scale-fun}{\vnbstmt{\begin{array}{@{}l@{}} \forall v.; \forall r \in \mathbb{R}. \\ \quad \operatorname{is\text{-}semi\text{-}inner\text{-}product\text{-}space}(v) \Rightarrow  \\ \quad \quad \operatorname{sip\text{-}qscale}(v, r) \in (\operatorname{quotient}(\operatorname{sip\text{-}setoid}(v)) \to \operatorname{quotient}(\operatorname{sip\text{-}setoid}(v))) \end{array}}}
\vnbprop{sipq-act-computes}{\vnbstmt{\begin{array}{@{}l@{}} \forall v.; \forall r \in \mathbb{R}, a \in \operatorname{vec}(v). \\ \quad \operatorname{is\text{-}semi\text{-}inner\text{-}product\text{-}space}(v) \Rightarrow  \\ \quad \quad \operatorname{act}(\operatorname{sip\text{-}quotient}(v))(r, \operatorname{class}(\operatorname{sip\text{-}setoid}(v), a)) = \operatorname{class}(\operatorname{sip\text{-}setoid}(v), \operatorname{act}(v)(r, a)) \end{array}}}
\vnbprop{sipq-vadd-fun}{\vnbstmt{\begin{array}{@{}l@{}} \forall v. \\ \quad \operatorname{is\text{-}semi\text{-}inner\text{-}product\text{-}space}(v) \Rightarrow  \\ \quad \quad \operatorname{vadd}(\operatorname{sip\text{-}quotient}(v)) \\ \quad \quad \in\; ((\operatorname{vec}(\operatorname{sip\text{-}quotient}(v)) \times \operatorname{vec}(\operatorname{sip\text{-}quotient}(v))) \to \operatorname{vec}(\operatorname{sip\text{-}quotient}(v))) \end{array}}}
\vnbprop{sipq-vneg-fun}{\vnbstmt{\begin{array}{@{}l@{}} \forall v. \\ \quad \operatorname{is\text{-}semi\text{-}inner\text{-}product\text{-}space}(v) \Rightarrow  \\ \quad \quad \operatorname{vneg}(\operatorname{sip\text{-}quotient}(v)) \in (\operatorname{vec}(\operatorname{sip\text{-}quotient}(v)) \to \operatorname{vec}(\operatorname{sip\text{-}quotient}(v))) \end{array}}}
\vnbprop{sipq-rip-fun}{\vnbstmt{\begin{array}{@{}l@{}} \forall v. \\ \quad \operatorname{is\text{-}semi\text{-}inner\text{-}product\text{-}space}(v) \Rightarrow  \\ \quad \quad \langle\cdot\mid\cdot\rangle_{\operatorname{sip\text{-}quotient}(v)} \in ((\operatorname{vec}(\operatorname{sip\text{-}quotient}(v)) \times \operatorname{vec}(\operatorname{sip\text{-}quotient}(v))) \to \mathbb{R}) \end{array}}}
\vnbprop{sipq-act-fun}{\vnbstmt{\begin{array}{@{}l@{}} \forall v. \\ \quad \operatorname{is\text{-}semi\text{-}inner\text{-}product\text{-}space}(v) \Rightarrow  \\ \quad \quad \operatorname{act}(\operatorname{sip\text{-}quotient}(v)) \\ \quad \quad \in\; ((\operatorname{carr}(\operatorname{scal}(\operatorname{sip\text{-}quotient}(v))) \times \operatorname{vec}(\operatorname{sip\text{-}quotient}(v))) \to \operatorname{vec}(\operatorname{sip\text{-}quotient}(v))) \end{array}}}
\vnbprop{sip-scal-ring}{\vnbstmt{\forall v.\; (\operatorname{is\text{-}semi\text{-}inner\text{-}product\text{-}space}(v) \Rightarrow \operatorname{is\text{-}ring}(\operatorname{scal}(v)))}}
\vnbprop{sipq-length}{\vnbstmt{\forall v.\; (\operatorname{is\text{-}semi\text{-}inner\text{-}product\text{-}space}(v) \Rightarrow (\operatorname{length}(\operatorname{sip\text{-}quotient}(v)) = 7))}}
\vnbprop{sipq-scal-ring}{\vnbstmt{\begin{array}{@{}l@{}} \forall v. \\ \quad \operatorname{is\text{-}semi\text{-}inner\text{-}product\text{-}space}(v) \Rightarrow \operatorname{is\text{-}ring}(\operatorname{scal}(\operatorname{sip\text{-}quotient}(v))) \end{array}}}
\vnbprop{sipq-scal-pin}{\vnbstmt{\begin{array}{@{}l@{}} \forall v. \\ \quad \operatorname{is\text{-}semi\text{-}inner\text{-}product\text{-}space}(v) \Rightarrow  \\ \quad \quad \operatorname{scal}(\operatorname{sip\text{-}quotient}(v)) = \operatorname{normed\text{-}field\text{-}as\text{-}commutative\text{-}ring}(\operatorname{rr\text{-}normed\text{-}field}) \end{array}}}
\vnbprop{sipq-scal-carr}{\vnbstmt{\begin{array}{@{}l@{}} \forall v. \\ \quad \operatorname{is\text{-}semi\text{-}inner\text{-}product\text{-}space}(v) \Rightarrow (\operatorname{carr}(\operatorname{scal}(\operatorname{sip\text{-}quotient}(v))) = \mathbb{R}) \end{array}}}
\vnbprop{sipq-vec-set}{\vnbstmt{\forall v.\; (\operatorname{is\text{-}semi\text{-}inner\text{-}product\text{-}space}(v) \Rightarrow (\operatorname{vec}(\operatorname{sip\text{-}quotient}(v)) \in \mathbf{Set}))}}
\vnbprop{sipq-vzero-in}{\vnbstmt{\begin{array}{@{}l@{}} \forall v. \\ \quad \operatorname{is\text{-}semi\text{-}inner\text{-}product\text{-}space}(v) \Rightarrow  \\ \quad \quad \operatorname{vzero}(\operatorname{sip\text{-}quotient}(v)) \in \operatorname{vec}(\operatorname{sip\text{-}quotient}(v)) \end{array}}}
\vnbprop{sipq-assoc}{\vnbstmt{\begin{array}{@{}l@{}} \forall v. \\ \quad \operatorname{is\text{-}semi\text{-}inner\text{-}product\text{-}space}(v) \Rightarrow  \\ \quad \quad \operatorname{is\text{-}associative}(\operatorname{vadd}(\operatorname{sip\text{-}quotient}(v)), \operatorname{vec}(\operatorname{sip\text{-}quotient}(v))) \end{array}}}
\vnbprop{sipq-comm}{\vnbstmt{\begin{array}{@{}l@{}} \forall v. \\ \quad \operatorname{is\text{-}semi\text{-}inner\text{-}product\text{-}space}(v) \Rightarrow  \\ \quad \quad \operatorname{is\text{-}commutative}(\operatorname{vadd}(\operatorname{sip\text{-}quotient}(v)), \operatorname{vec}(\operatorname{sip\text{-}quotient}(v))) \end{array}}}
\vnbprop{sipq-identity}{\vnbstmt{\begin{array}{@{}l@{}} \forall v. \\ \quad \operatorname{is\text{-}semi\text{-}inner\text{-}product\text{-}space}(v) \Rightarrow  \\ \quad \quad \operatorname{is\text{-}identity}(\operatorname{vadd}(\operatorname{sip\text{-}quotient}(v)), \operatorname{vzero}(\operatorname{sip\text{-}quotient}(v)), \operatorname{vec}(\operatorname{sip\text{-}quotient}(v))) \end{array}}}
\vnbprop{sipq-inverses}{\vnbstmt{\begin{array}{@{}l@{}} \forall v. \\ \quad \operatorname{is\text{-}semi\text{-}inner\text{-}product\text{-}space}(v) \Rightarrow  \\ \quad \quad \operatorname{has\text{-}inverses}(\operatorname{vadd}(\operatorname{sip\text{-}quotient}(v)), \operatorname{vzero}(\operatorname{sip\text{-}quotient}(v)), \\ \quad \quad \quad \operatorname{vneg}(\operatorname{sip\text{-}quotient}(v)), \operatorname{vec}(\operatorname{sip\text{-}quotient}(v))) \end{array}}}
\vnbprop{sipq-act-distrib-vec}{\vnbstmt{\begin{array}{@{}l@{}} \forall v.; \forall r \in \operatorname{carr}(\operatorname{scal}(\operatorname{sip\text{-}quotient}(v))), x, y \in \operatorname{vec}(\operatorname{sip\text{-}quotient}(v)). \\ \quad \operatorname{is\text{-}semi\text{-}inner\text{-}product\text{-}space}(v) \Rightarrow  \\ \quad \quad \operatorname{act}(\operatorname{sip\text{-}quotient}(v))(r, \operatorname{vadd}(\operatorname{sip\text{-}quotient}(v))(x, y)) \\ \quad \quad =\; \operatorname{vadd}(\operatorname{sip\text{-}quotient}(v))(\operatorname{act}(\operatorname{sip\text{-}quotient}(v))(r, x), \operatorname{act}(\operatorname{sip\text{-}quotient}(v))(r, y)) \end{array}}}
\vnbprop{sipq-act-distrib-scalar}{\vnbstmt{\begin{array}{@{}l@{}} \forall v.; \forall r, s \in \operatorname{carr}(\operatorname{scal}(\operatorname{sip\text{-}quotient}(v))), x \in \operatorname{vec}(\operatorname{sip\text{-}quotient}(v)). \\ \quad \operatorname{is\text{-}semi\text{-}inner\text{-}product\text{-}space}(v) \Rightarrow  \\ \quad \quad \operatorname{act}(\operatorname{sip\text{-}quotient}(v))(\operatorname{add}(\operatorname{scal}(\operatorname{sip\text{-}quotient}(v)))(r, s), x) \\ \quad \quad =\; \operatorname{vadd}(\operatorname{sip\text{-}quotient}(v))(\operatorname{act}(\operatorname{sip\text{-}quotient}(v))(r, x), \operatorname{act}(\operatorname{sip\text{-}quotient}(v))(s, x)) \end{array}}}
\vnbprop{sipq-act-mul-compat}{\vnbstmt{\begin{array}{@{}l@{}} \forall v.; \forall r, s \in \operatorname{carr}(\operatorname{scal}(\operatorname{sip\text{-}quotient}(v))), x \in \operatorname{vec}(\operatorname{sip\text{-}quotient}(v)). \\ \quad \operatorname{is\text{-}semi\text{-}inner\text{-}product\text{-}space}(v) \Rightarrow  \\ \quad \quad \operatorname{act}(\operatorname{sip\text{-}quotient}(v))(\operatorname{mul}(\operatorname{scal}(\operatorname{sip\text{-}quotient}(v)))(r, s), x) \\ \quad \quad =\; \operatorname{act}(\operatorname{sip\text{-}quotient}(v))(r, \operatorname{act}(\operatorname{sip\text{-}quotient}(v))(s, x)) \end{array}}}
\vnbprop{sipq-act-unital}{\vnbstmt{\begin{array}{@{}l@{}} \forall v.; \forall x \in \operatorname{vec}(\operatorname{sip\text{-}quotient}(v)). \\ \quad \operatorname{is\text{-}semi\text{-}inner\text{-}product\text{-}space}(v) \Rightarrow  \\ \quad \quad \operatorname{act}(\operatorname{sip\text{-}quotient}(v))(\operatorname{one}(\operatorname{scal}(\operatorname{sip\text{-}quotient}(v))), x) = x \end{array}}}
\vnbprop{sipq-rip-sym}{\vnbstmt{\begin{array}{@{}l@{}} \forall v.; \forall x, y \in \operatorname{vec}(\operatorname{sip\text{-}quotient}(v)). \\ \quad \operatorname{is\text{-}semi\text{-}inner\text{-}product\text{-}space}(v) \Rightarrow  \\ \quad \quad \langle\cdot\mid\cdot\rangle_{\operatorname{sip\text{-}quotient}(v)}(x, y) = \langle\cdot\mid\cdot\rangle_{\operatorname{sip\text{-}quotient}(v)}(y, x) \end{array}}}
\vnbprop{sipq-rip-add-left}{\vnbstmt{\begin{array}{@{}l@{}} \forall v.; \forall x, y, z \in \operatorname{vec}(\operatorname{sip\text{-}quotient}(v)). \\ \quad \operatorname{is\text{-}semi\text{-}inner\text{-}product\text{-}space}(v) \Rightarrow  \\ \quad \quad \langle\cdot\mid\cdot\rangle_{\operatorname{sip\text{-}quotient}(v)}(\operatorname{vadd}(\operatorname{sip\text{-}quotient}(v))(x, y), z) \\ \quad \quad =\; \langle\cdot\mid\cdot\rangle_{\operatorname{sip\text{-}quotient}(v)}(x, z) + \langle\cdot\mid\cdot\rangle_{\operatorname{sip\text{-}quotient}(v)}(y, z) \end{array}}}
\vnbprop{sipq-rip-homog-left}{\vnbstmt{\begin{array}{@{}l@{}} \forall v.; \forall a \in \mathbb{R}, x, y \in \operatorname{vec}(\operatorname{sip\text{-}quotient}(v)). \\ \quad \operatorname{is\text{-}semi\text{-}inner\text{-}product\text{-}space}(v) \Rightarrow  \\ \quad \quad \langle\cdot\mid\cdot\rangle_{\operatorname{sip\text{-}quotient}(v)}(\operatorname{act}(\operatorname{sip\text{-}quotient}(v))(a, x), y) = (a \cdot \langle\cdot\mid\cdot\rangle_{\operatorname{sip\text{-}quotient}(v)}(x, y)) \end{array}}}
\vnbprop{sipq-rip-self-nonneg}{\vnbstmt{\begin{array}{@{}l@{}} \forall v.; \forall x \in \operatorname{vec}(\operatorname{sip\text{-}quotient}(v)). \\ \quad \operatorname{is\text{-}semi\text{-}inner\text{-}product\text{-}space}(v) \Rightarrow (0 \leq \langle\cdot\mid\cdot\rangle_{\operatorname{sip\text{-}quotient}(v)}(x, x)) \end{array}}}
\vnbprop{sipq-definite}{\vnbstmt{\begin{array}{@{}l@{}} \forall v.; \forall x \in \operatorname{vec}(\operatorname{sip\text{-}quotient}(v)). \\ \quad (\operatorname{is\text{-}semi\text{-}inner\text{-}product\text{-}space}(v) \wedge (\langle\cdot\mid\cdot\rangle_{\operatorname{sip\text{-}quotient}(v)}(x, x) = 0)) \Rightarrow  \\ \quad \quad x = \operatorname{vzero}(\operatorname{sip\text{-}quotient}(v)) \end{array}}}
\vnbprop{sip-quotient-is-pre-hilbert}{\vnbstmt{\begin{array}{@{}l@{}} \forall v. \\ \quad \operatorname{is\text{-}semi\text{-}inner\text{-}product\text{-}space}(v) \Rightarrow \operatorname{is\text{-}pre\text{-}hilbert\text{-}space}(\operatorname{sip\text{-}quotient}(v)) \end{array}}}
\vnbprop{phs-is-sip}{\vnbstmt{\forall v.\; (\operatorname{is\text{-}pre\text{-}hilbert\text{-}space}(v) \Rightarrow \operatorname{is\text{-}semi\text{-}inner\text{-}product\text{-}space}(v))}}
\vnbprop{phs-definite}{\vnbstmt{\begin{array}{@{}l@{}} \forall v.; \forall x \in \operatorname{vec}(v). \\ \quad (\operatorname{is\text{-}pre\text{-}hilbert\text{-}space}(v) \wedge (\langle\cdot\mid\cdot\rangle_{v}(x, x) = 0)) \Rightarrow (x = \operatorname{vzero}(v)) \end{array}}}
\vnbprop{sip-act-neg-scalar}{\vnbstmt{\begin{array}{@{}l@{}} \forall v.; \forall q \in \mathbb{R}, x \in \operatorname{vec}(v). \\ \quad \operatorname{is\text{-}semi\text{-}inner\text{-}product\text{-}space}(v) \Rightarrow  \\ \quad \quad \operatorname{act}(v)(-q, x) = \operatorname{vneg}(v)(\operatorname{act}(v)(q, x)) \end{array}}}
\vnbprop{sip-vadd-fun}{\vnbstmt{\begin{array}{@{}l@{}} \forall v. \\ \quad \operatorname{is\text{-}semi\text{-}inner\text{-}product\text{-}space}(v) \Rightarrow  \\ \quad \quad \operatorname{vadd}(v) \in ((\operatorname{vec}(v) \times \operatorname{vec}(v)) \to \operatorname{vec}(v)) \end{array}}}
\vnbprop{sip-vzero-in}{\vnbstmt{\forall v.\; (\operatorname{is\text{-}semi\text{-}inner\text{-}product\text{-}space}(v) \Rightarrow (\operatorname{vzero}(v) \in \operatorname{vec}(v)))}}
\vnbprop{sip-vneg-fun}{\vnbstmt{\begin{array}{@{}l@{}} \forall v. \\ \quad \operatorname{is\text{-}semi\text{-}inner\text{-}product\text{-}space}(v) \Rightarrow (\operatorname{vneg}(v) \in (\operatorname{vec}(v) \to \operatorname{vec}(v))) \end{array}}}
\vnbprop{sip-act-fun}{\vnbstmt{\begin{array}{@{}l@{}} \forall v. \\ \quad \operatorname{is\text{-}semi\text{-}inner\text{-}product\text{-}space}(v) \Rightarrow  \\ \quad \quad \operatorname{act}(v) \in ((\operatorname{carr}(\operatorname{scal}(v)) \times \operatorname{vec}(v)) \to \operatorname{vec}(v)) \end{array}}}
\vnbprop{sip-vadd-assoc-prop}{\vnbstmt{\begin{array}{@{}l@{}} \forall v. \\ \quad \operatorname{is\text{-}semi\text{-}inner\text{-}product\text{-}space}(v) \Rightarrow \operatorname{is\text{-}associative}(\operatorname{vadd}(v), \operatorname{vec}(v)) \end{array}}}
\vnbprop{sip-vadd-comm-prop}{\vnbstmt{\begin{array}{@{}l@{}} \forall v. \\ \quad \operatorname{is\text{-}semi\text{-}inner\text{-}product\text{-}space}(v) \Rightarrow \operatorname{is\text{-}commutative}(\operatorname{vadd}(v), \operatorname{vec}(v)) \end{array}}}
\vnbprop{sip-vadd-identity-prop}{\vnbstmt{\begin{array}{@{}l@{}} \forall v. \\ \quad \operatorname{is\text{-}semi\text{-}inner\text{-}product\text{-}space}(v) \Rightarrow \operatorname{is\text{-}identity}(\operatorname{vadd}(v), \operatorname{vzero}(v), \operatorname{vec}(v)) \end{array}}}
\vnbprop{sip-vadd-inverses-prop}{\vnbstmt{\begin{array}{@{}l@{}} \forall v. \\ \quad \operatorname{is\text{-}semi\text{-}inner\text{-}product\text{-}space}(v) \Rightarrow  \\ \quad \quad \operatorname{has\text{-}inverses}(\operatorname{vadd}(v), \operatorname{vzero}(v), \operatorname{vneg}(v), \operatorname{vec}(v)) \end{array}}}
\vnbprop{sipq-unfold}{\vnbstmt{\begin{array}{@{}l@{}} \forall v. \\ \quad \operatorname{is\text{-}semi\text{-}inner\text{-}product\text{-}space}(v) \Rightarrow  \\ \quad \quad (\operatorname{sip\text{-}quotient}(v) \simeq {[}\operatorname{scal}(v), \operatorname{quotient}(\operatorname{sip\text{-}setoid}(v)), \operatorname{sip\text{-}qvadd}(v), \operatorname{class}(\operatorname{sip\text{-}setoid}(v), \operatorname{vzero}(v)), \operatorname{sip\text{-}qvneg}(v), \operatorname{sip\text{-}qact}(v), \operatorname{sip\text{-}qrip}(v)]) \end{array}}}

\vnbbookfile{structure-library/sip-complex-vocab}{structure-library/sip-complex-vocab}
\noindent{\small sip-complex-vocab.scm -- THE COMPLEX STRUCTURE J on a semi-inner-product space (SIP-1, 2026-10-18).  Vocabulary only.}\par
\vnbstructdef{complex-semi-inner-product-space}{A complex semi inner product space $s$ is an 8-tuple whose components are listed below.\par
\begin{enumerate}[nosep]
\item $\operatorname{scal}(s) = \operatorname{nth}(1, s)$, a substructure (\texttt{substructure}): $\operatorname{is\text{-}ring}(\operatorname{scal}(s))$.
\item $\operatorname{vec}(s) = \operatorname{nth}(2, s)$, a carrier (\texttt{carriers}): $\operatorname{vec}(s) \in \mathbf{Set}$.
\item $\operatorname{vadd}(s) = \operatorname{nth}(3, s)$, an operation (\texttt{op}): $\operatorname{vadd}(s) \in ((\operatorname{vec}(s) \times \operatorname{vec}(s)) \to \operatorname{vec}(s))$.
\item $\operatorname{vzero}(s) = \operatorname{nth}(4, s)$, a constant (\texttt{constant}): $\operatorname{vzero}(s) \in \operatorname{vec}(s)$.
\item $\operatorname{vneg}(s) = \operatorname{nth}(5, s)$, an operation (\texttt{op}): $\operatorname{vneg}(s) \in (\operatorname{vec}(s) \to \operatorname{vec}(s))$.
\item $\operatorname{act}(s) = \operatorname{nth}(6, s)$, an operation (\texttt{op}): $\operatorname{act}(s) \in ((\operatorname{carr}(\operatorname{scal}(s)) \times \operatorname{vec}(s)) \to \operatorname{vec}(s))$.
\item $\langle\cdot\mid\cdot\rangle_{s} = \operatorname{nth}(7, s)$, an operation (\texttt{op}): $\langle\cdot\mid\cdot\rangle_{s} \in ((\operatorname{vec}(s) \times \operatorname{vec}(s)) \to \mathbb{R})$.
\item $J_{s} = \operatorname{nth}(8, s)$, an operation (\texttt{op}): $J_{s} \in (\operatorname{vec}(s) \to \operatorname{vec}(s))$.
\end{enumerate}
The laws (the \texttt{property} and \texttt{law} clauses):
\begin{itemize}[nosep]
\item $\operatorname{is\text{-}semi\text{-}inner\text{-}product\text{-}space}({[}\operatorname{scal}(s), \operatorname{vec}(s), \operatorname{vadd}(s), \operatorname{vzero}(s), \operatorname{vneg}(s), \operatorname{act}(s), \langle\cdot\mid\cdot\rangle_{s}])$
\item $\operatorname{is\text{-}real\text{-}linear\text{-}op}(s, J_{s})$
\item $\operatorname{is\text{-}square\text{-}minus\text{-}id}(s, J_{s})$
\item $\operatorname{is\text{-}antisymmetric\text{-}op}(s, J_{s})$
\end{itemize}
The characteristic law (the axiom \texttt{is-complex-semi-inner-product-space}): $\operatorname{is\text{-}complex\text{-}semi\text{-}inner\text{-}product\text{-}space}(s)$ holds if and only if $\operatorname{length}(s) = 8$, every component has the typing listed with it, every law holds, and $s \in \mathbf{Set}$ (the structure is a set).
}
\par\medskip\noindent\textbf{Definition} (\texttt{cact}).\;$\operatorname{cact}(v) :=$
\begin{equation*}
\begin{array}{@{}l@{}} (\lambda {[}\operatorname{cxz\_}, \operatorname{cxx\_}] \in (\mathbb{C} \times \operatorname{vec}(v)).\; \\ \quad \operatorname{vadd}(v)(\operatorname{act}(v)(\operatorname{real\text{-}part}(\operatorname{cxz\_}), \operatorname{cxx\_}), \operatorname{act}(v)(\operatorname{imag\text{-}part}(\operatorname{cxz\_}), J_{v}(\operatorname{cxx\_})))) \end{array}
\end{equation*}

\par\medskip\noindent\textbf{Definition} (\texttt{crip}).\;$\operatorname{crip}(v) :=$
\begin{equation*}
\begin{array}{@{}l@{}} (\lambda {[}\operatorname{cxx\_}, \operatorname{cxy\_}] \in (\operatorname{vec}(v) \times \operatorname{vec}(v)).\; \\ \quad (\langle\cdot\mid\cdot\rangle_{v}(\operatorname{cxx\_}, \operatorname{cxy\_}) - (+i \cdot \langle\cdot\mid\cdot\rangle_{v}(J_{v}(\operatorname{cxx\_}), \operatorname{cxy\_})))) \end{array}
\end{equation*}

\par\medskip\noindent\textbf{Definition} (\texttt{csip-cips}).\;$\operatorname{csip\text{-}cips}(v) :=$
\begin{equation*}
{[}\operatorname{normed\text{-}field\text{-}as\text{-}commutative\text{-}ring}(\operatorname{cc\text{-}normed\text{-}field}), \operatorname{vec}(v), \operatorname{vadd}(v), \operatorname{vzero}(v), \operatorname{vneg}(v), \operatorname{cact}(v), \operatorname{crip}(v)]
\end{equation*}

\par\medskip\noindent\textbf{Definition} (\texttt{csip-quotient}).\;$\operatorname{csip\text{-}quotient}(v) :=$
\begin{equation*}
{[}\operatorname{scal}(v), \operatorname{quotient}(\operatorname{sip\text{-}setoid}(\operatorname{csip\text{-}as\text{-}sip}(v))), \operatorname{sip\text{-}qvadd}(\operatorname{csip\text{-}as\text{-}sip}(v)), \operatorname{class}(\operatorname{sip\text{-}setoid}(\operatorname{csip\text{-}as\text{-}sip}(v)), \operatorname{vzero}(v)), \operatorname{sip\text{-}qvneg}(\operatorname{csip\text{-}as\text{-}sip}(v)), \operatorname{sip\text{-}qact}(\operatorname{csip\text{-}as\text{-}sip}(v)), \operatorname{sip\text{-}qrip}(\operatorname{csip\text{-}as\text{-}sip}(v)), \operatorname{csip\text{-}qj}(v)]
\end{equation*}

\par\medskip\noindent\textbf{Definition} (\texttt{is-real-linear-op}).\;For all $v$ and $\mathit{jo}$, 
\begin{equation*}
\begin{array}{@{}l@{}} \operatorname{is\text{-}real\text{-}linear\text{-}op}(v, \mathit{jo}) \\ \Leftrightarrow\; \mathit{jo} \in (\operatorname{vec}(v) \to \operatorname{vec}(v)) \wedge \\ \forall \operatorname{x\_}, \operatorname{y\_} \in \operatorname{vec}(v). \\ \quad \operatorname{jo}(\operatorname{vadd}(v)(\operatorname{x\_}, \operatorname{y\_})) = \operatorname{vadd}(v)(\operatorname{jo}(\operatorname{x\_}), \operatorname{jo}(\operatorname{y\_})) \wedge \\ \forall \operatorname{a\_} \in \mathbb{R}, \operatorname{x\_} \in \operatorname{vec}(v).\; (\operatorname{jo}(\operatorname{act}(v)(\operatorname{a\_}, \operatorname{x\_})) = \operatorname{act}(v)(\operatorname{a\_}, \operatorname{jo}(\operatorname{x\_}))) \end{array}
\end{equation*}

\par\medskip\noindent\textbf{Definition} (\texttt{is-antisymmetric-op}).\;For all $v$ and $\mathit{jo}$, 
\begin{equation*}
\begin{array}{@{}l@{}} \operatorname{is\text{-}antisymmetric\text{-}op}(v, \mathit{jo}) \\ \Leftrightarrow\; \forall \operatorname{x\_}, \operatorname{y\_} \in \operatorname{vec}(v). \\ \quad \langle\cdot\mid\cdot\rangle_{v}(\operatorname{jo}(\operatorname{x\_}), \operatorname{y\_}) = -(\langle\cdot\mid\cdot\rangle_{v}(\operatorname{x\_}, \operatorname{jo}(\operatorname{y\_}))) \end{array}
\end{equation*}

\par\medskip\noindent\textbf{Definition} (\texttt{is-isometric-op}).\;For all $v$ and $\mathit{jo}$, 
\begin{equation*}
\begin{array}{@{}l@{}} \operatorname{is\text{-}isometric\text{-}op}(v, \mathit{jo}) \\ \Leftrightarrow\; \forall \operatorname{x\_}, \operatorname{y\_} \in \operatorname{vec}(v). \\ \quad \langle\cdot\mid\cdot\rangle_{v}(\operatorname{jo}(\operatorname{x\_}), \operatorname{jo}(\operatorname{y\_})) = \langle\cdot\mid\cdot\rangle_{v}(\operatorname{x\_}, \operatorname{y\_}) \end{array}
\end{equation*}

\par\medskip\noindent\textbf{Definition} (\texttt{is-square-minus-id}).\;For all $v$ and $\mathit{jo}$, 
\begin{equation*}
\operatorname{is\text{-}square\text{-}minus\text{-}id}(v, \mathit{jo}) \Leftrightarrow \forall \operatorname{x\_} \in \operatorname{vec}(v).\; (\operatorname{jo}(\operatorname{jo}(\operatorname{x\_})) = \operatorname{vneg}(v)(\operatorname{x\_}))
\end{equation*}

\par\medskip\noindent\textbf{Definition} (\texttt{is-sesquilinear-form}).\;For all $v$ and $\mathit{sqf}$, 
\begin{equation*}
\begin{array}{@{}l@{}} \operatorname{is\text{-}sesquilinear\text{-}form}(v, \mathit{sqf}) \\ \Leftrightarrow\; \mathit{sqf} \in ((\operatorname{vec}(v) \times \operatorname{vec}(v)) \to \mathbb{C}) \wedge \\ \forall \operatorname{x\_}, \operatorname{y\_}, \operatorname{z\_} \in \operatorname{vec}(v). \\ \quad \operatorname{sqf}(\operatorname{vadd}(v)(\operatorname{x\_}, \operatorname{y\_}), \operatorname{z\_}) = +(\operatorname{sqf}(\operatorname{x\_}, \operatorname{z\_}), \operatorname{sqf}(\operatorname{y\_}, \operatorname{z\_})) \wedge \\ \forall \operatorname{a\_} \in \mathbb{C}, \operatorname{x\_}, \operatorname{y\_} \in \operatorname{vec}(v). \\ \quad \operatorname{sqf}(\operatorname{cact}(v)(\operatorname{a\_}, \operatorname{x\_}), \operatorname{y\_}) = \cdot(\operatorname{a\_}, \operatorname{sqf}(\operatorname{x\_}, \operatorname{y\_})) \wedge \\ \forall \operatorname{x\_}, \operatorname{y\_} \in \operatorname{vec}(v).\; (\operatorname{sqf}(\operatorname{y\_}, \operatorname{x\_}) = \operatorname{conjugate}(\operatorname{sqf}(\operatorname{x\_}, \operatorname{y\_}))) \end{array}
\end{equation*}

\par\medskip\noindent\textbf{Definition} (\texttt{csip-qj}).\;$\operatorname{csip\text{-}qj}(v) :=$
\begin{equation*}
\begin{array}{@{}l@{}} \operatorname{descend}(\operatorname{sip\text{-}setoid}(\operatorname{csip\text{-}as\text{-}sip}(v)), \\ \quad (\lambda \operatorname{sqa\_} \in \operatorname{vec}(v).\; \operatorname{class}(\operatorname{sip\text{-}setoid}(\operatorname{csip\text{-}as\text{-}sip}(v)), J_{v}(\operatorname{sqa\_})))) \end{array}
\end{equation*}

\par\medskip\noindent\textbf{Definition} (\texttt{cips-ract}).\;$\operatorname{cips\text{-}ract}(w) :=$
\begin{equation*}
(\lambda {[}\operatorname{cxr\_}, \operatorname{cxx\_}] \in (\mathbb{R} \times \operatorname{vec}(w)).\; \operatorname{act}(w)(\operatorname{cxr\_}, \operatorname{cxx\_}))
\end{equation*}

\par\medskip\noindent\textbf{Definition} (\texttt{cips-rip}).\;$\operatorname{cips\text{-}rip}(w) :=$
\begin{equation*}
(\lambda {[}\operatorname{cxx\_}, \operatorname{cxy\_}] \in (\operatorname{vec}(w) \times \operatorname{vec}(w)).\; \operatorname{real\text{-}part}(\langle\cdot,\cdot\rangle_{w}(\operatorname{cxx\_}, \operatorname{cxy\_})))
\end{equation*}

\par\medskip\noindent\textbf{Definition} (\texttt{cips-j}).\;$\operatorname{cips\text{-}j}(w) :=$ $(\lambda \operatorname{cxx\_} \in \operatorname{vec}(w).\; \operatorname{act}(w)(+i, \operatorname{cxx\_}))$
\par\medskip\noindent\textbf{Definition} (\texttt{cips-rsip}).\;$\operatorname{cips\text{-}rsip}(w) :=$
\begin{equation*}
{[}\operatorname{normed\text{-}field\text{-}as\text{-}commutative\text{-}ring}(\operatorname{rr\text{-}normed\text{-}field}), \operatorname{vec}(w), \operatorname{vadd}(w), \operatorname{vzero}(w), \operatorname{vneg}(w), \operatorname{cips\text{-}ract}(w), \operatorname{cips\text{-}rip}(w)]
\end{equation*}

\par\medskip\noindent\textbf{Definition} (\texttt{cips-real}).\;$\operatorname{cips\text{-}real}(w) :=$
\begin{equation*}
{[}\operatorname{normed\text{-}field\text{-}as\text{-}commutative\text{-}ring}(\operatorname{rr\text{-}normed\text{-}field}), \operatorname{vec}(w), \operatorname{vadd}(w), \operatorname{vzero}(w), \operatorname{vneg}(w), \operatorname{cips\text{-}ract}(w), \operatorname{cips\text{-}rip}(w), \operatorname{cips\text{-}j}(w)]
\end{equation*}

\par\medskip\noindent\textbf{Definition} (\texttt{complexify}).\;$\operatorname{complexify}(v) :=$
\begin{equation*}
{[}\operatorname{scal}(v), (\operatorname{vec}(v) \times \operatorname{vec}(v)), \operatorname{cx\text{-}vadd}(v), {[}\operatorname{vzero}(v), \operatorname{vzero}(v)], \operatorname{cx\text{-}vneg}(v), \operatorname{cx\text{-}act}(v), \operatorname{cx\text{-}rip}(v), \operatorname{cx\text{-}j}(v)]
\end{equation*}

\par\medskip\noindent\textbf{Definition} (\texttt{cx-vadd}).\;$\operatorname{cx\text{-}vadd}(v) :=$
\begin{equation*}
\begin{array}{@{}l@{}} (\lambda {[}\operatorname{cp\_}, \operatorname{cq\_}] \in ((\operatorname{vec}(v) \times \operatorname{vec}(v)) \times (\operatorname{vec}(v) \times \operatorname{vec}(v))).\; \\ \quad {[}\operatorname{vadd}(v)(\operatorname{nth}(1, \operatorname{cp\_}), \operatorname{nth}(1, \operatorname{cq\_})), \operatorname{vadd}(v)(\operatorname{nth}(2, \operatorname{cp\_}), \operatorname{nth}(2, \operatorname{cq\_}))]) \end{array}
\end{equation*}

\par\medskip\noindent\textbf{Definition} (\texttt{cx-vneg}).\;$\operatorname{cx\text{-}vneg}(v) :=$
\begin{equation*}
(\lambda \operatorname{cp\_} \in (\operatorname{vec}(v) \times \operatorname{vec}(v)).\; {[}\operatorname{vneg}(v)(\operatorname{nth}(1, \operatorname{cp\_})), \operatorname{vneg}(v)(\operatorname{nth}(2, \operatorname{cp\_}))])
\end{equation*}

\par\medskip\noindent\textbf{Definition} (\texttt{cx-act}).\;$\operatorname{cx\text{-}act}(v) :=$
\begin{equation*}
\begin{array}{@{}l@{}} (\lambda {[}\operatorname{cr\_}, \operatorname{cp\_}] \in (\operatorname{carr}(\operatorname{scal}(v)) \times (\operatorname{vec}(v) \times \operatorname{vec}(v))).\; \\ \quad {[}\operatorname{act}(v)(\operatorname{cr\_}, \operatorname{nth}(1, \operatorname{cp\_})), \operatorname{act}(v)(\operatorname{cr\_}, \operatorname{nth}(2, \operatorname{cp\_}))]) \end{array}
\end{equation*}

\par\medskip\noindent\textbf{Definition} (\texttt{cx-rip}).\;$\operatorname{cx\text{-}rip}(v) :=$
\begin{equation*}
\begin{array}{@{}l@{}} (\lambda {[}\operatorname{cp\_}, \operatorname{cq\_}] \in ((\operatorname{vec}(v) \times \operatorname{vec}(v)) \times (\operatorname{vec}(v) \times \operatorname{vec}(v))).\; \\ \quad \langle\cdot\mid\cdot\rangle_{v}(\operatorname{nth}(1, \operatorname{cp\_}), \operatorname{nth}(1, \operatorname{cq\_})) + \langle\cdot\mid\cdot\rangle_{v}(\operatorname{nth}(2, \operatorname{cp\_}), \operatorname{nth}(2, \operatorname{cq\_}))) \end{array}
\end{equation*}

\par\medskip\noindent\textbf{Definition} (\texttt{cx-j}).\;$\operatorname{cx\text{-}j}(v) :=$
\begin{equation*}
(\lambda \operatorname{cp\_} \in (\operatorname{vec}(v) \times \operatorname{vec}(v)).\; {[}\operatorname{vneg}(v)(\operatorname{nth}(2, \operatorname{cp\_})), \operatorname{nth}(1, \operatorname{cp\_})])
\end{equation*}

\par\medskip\noindent\textbf{Definition} (\texttt{cx-sip}).\;$\operatorname{cx\text{-}sip}(v) :=$
\begin{equation*}
{[}\operatorname{scal}(v), (\operatorname{vec}(v) \times \operatorname{vec}(v)), \operatorname{cx\text{-}vadd}(v), {[}\operatorname{vzero}(v), \operatorname{vzero}(v)], \operatorname{cx\text{-}vneg}(v), \operatorname{cx\text{-}act}(v), \operatorname{cx\text{-}rip}(v)]
\end{equation*}

\par\medskip\noindent\textbf{Definition} (\texttt{hom-complex-pre-hilbert-space}).\;$\operatorname{hom\text{-}complex\text{-}pre\text{-}hilbert\text{-}space}(a, b) :=$
\begin{equation*}
\{\operatorname{homf\_} \in (\operatorname{vec}(a) \to \operatorname{vec}(b)) : \operatorname{is\text{-}hom\text{-}complex\text{-}pre\text{-}hilbert\text{-}space}(a, b, \operatorname{homf\_})\}
\end{equation*}

\par\medskip\noindent\textbf{Definition} (\texttt{hom-complex-semi-inner-product-space}).\;$\operatorname{hom\text{-}complex\text{-}semi\text{-}inner\text{-}product\text{-}space}(a, b) :=$
\begin{equation*}
\{\operatorname{homf\_} \in (\operatorname{vec}(a) \to \operatorname{vec}(b)) : \operatorname{is\text{-}hom\text{-}complex\text{-}semi\text{-}inner\text{-}product\text{-}space}(a, b, \operatorname{homf\_})\}
\end{equation*}

\vnbdef{is-hom-complex-pre-hilbert-space-def}{\vnbstmt{\begin{array}{@{}l@{}} \forall a, b, f. \\ \quad \operatorname{is\text{-}hom\text{-}complex\text{-}pre\text{-}hilbert\text{-}space}(a, b, f) \\ \quad \Leftrightarrow\; \operatorname{is\text{-}complex\text{-}pre\text{-}hilbert\text{-}space}(a) \wedge \\ \quad \operatorname{is\text{-}complex\text{-}pre\text{-}hilbert\text{-}space}(b) \wedge \\ \quad \operatorname{is\text{-}hom\text{-}complex\text{-}semi\text{-}inner\text{-}product\text{-}space}(a, b, f) \end{array}}}
\vnbdef{is-hom-complex-semi-inner-product-space-def}{\vnbstmt{\begin{array}{@{}l@{}} \forall a, b, f. \\ \quad \operatorname{is\text{-}hom\text{-}complex\text{-}semi\text{-}inner\text{-}product\text{-}space}(a, b, f) \\ \quad \Leftrightarrow\; \operatorname{is\text{-}complex\text{-}semi\text{-}inner\text{-}product\text{-}space}(a) \wedge \\ \quad \operatorname{is\text{-}complex\text{-}semi\text{-}inner\text{-}product\text{-}space}(b) \wedge \\ \quad f \in (\operatorname{vec}(a) \to \operatorname{vec}(b)) \wedge \\ \quad \operatorname{scal}(a) = \operatorname{scal}(b) \wedge \\ \quad \forall x_{1}, x_{2} \in \operatorname{vec}(a). \\ \quad \quad f(\operatorname{vadd}(a)(x_{1}, x_{2})) = \operatorname{vadd}(b)(f(x_{1}), f(x_{2})) \wedge \\ \quad f(\operatorname{vzero}(a)) = \operatorname{vzero}(b) \wedge \\ \quad \forall x_{1} \in \operatorname{vec}(a).\; (f(\operatorname{vneg}(a)(x_{1})) = \operatorname{vneg}(b)(f(x_{1}))) \wedge \\ \quad \forall x_{1} \in \operatorname{carr}(\operatorname{scal}(a)), x_{2} \in \operatorname{vec}(a). \\ \quad \quad f(\operatorname{act}(a)(x_{1}, x_{2})) = \operatorname{act}(b)(x_{1}, f(x_{2})) \wedge \\ \quad \forall x_{1}, x_{2} \in \operatorname{vec}(a).\; (\langle\cdot\mid\cdot\rangle_{a}(x_{1}, x_{2}) = \langle\cdot\mid\cdot\rangle_{b}(f(x_{1}), f(x_{2}))) \wedge \\ \quad \forall x_{1} \in \operatorname{vec}(a).\; (f(J_{a}(x_{1})) = J_{b}(f(x_{1}))) \end{array}}}
\vnbdef{complex-semi-inner-product-space-class}{\vnbstmt{\begin{array}{@{}l@{}} \forall s. \\ \quad (s \in \operatorname{complex\text{-}semi\text{-}inner\text{-}product\text{-}space}) \Leftrightarrow \operatorname{is\text{-}complex\text{-}semi\text{-}inner\text{-}product\text{-}space}(s) \end{array}}}
\vnbdef{complex-pre-hilbert-space-class}{\vnbstmt{\forall s.\; ((s \in \operatorname{complex\text{-}pre\text{-}hilbert\text{-}space}) \Leftrightarrow \operatorname{is\text{-}complex\text{-}pre\text{-}hilbert\text{-}space}(s))}}
\vnbdef{is-complex-pre-hilbert-space-def}{\vnbstmt{\begin{array}{@{}l@{}} \forall s. \\ \quad \operatorname{is\text{-}complex\text{-}pre\text{-}hilbert\text{-}space}(s) \\ \quad \Leftrightarrow\; \operatorname{is\text{-}complex\text{-}semi\text{-}inner\text{-}product\text{-}space}(s) \wedge \\ \quad \forall x \in \operatorname{vec}(s).\; ((\langle\cdot\mid\cdot\rangle_{s}(x, x) = 0) \Rightarrow (x = \operatorname{vzero}(s))) \end{array}}}
\vnbprop{is-complex-semi-inner-product-space-is-set}{\vnbstmt{\forall s.\; (\operatorname{is\text{-}complex\text{-}semi\text{-}inner\text{-}product\text{-}space}(s) \Rightarrow (s \in \mathbf{Set}))}}
\vnbprop{is-complex-pre-hilbert-space-is-set}{\vnbstmt{\forall s.\; (\operatorname{is\text{-}complex\text{-}pre\text{-}hilbert\text{-}space}(s) \Rightarrow (s \in \mathbf{Set}))}}

\vnbbookfile{theorem-library/sip-b}{theorem-library/sip-b}
\noindent{\small sip-b.scm -- THE COMPLEX STRUCTURE BY J (SIP-1, 2026-10-18; STOP POINT B of scratchpad/triage/SIP-1.md).  All PROVEN modulo 0.}\par
\vnbprop{csip-as-sip-is-semi-inner-product-space}{\vnbstmt{\begin{array}{@{}l@{}} \forall r. \\ \quad \operatorname{is\text{-}complex\text{-}semi\text{-}inner\text{-}product\text{-}space}(r) \Rightarrow  \\ \quad \quad \operatorname{is\text{-}semi\text{-}inner\text{-}product\text{-}space}(\operatorname{csip\text{-}as\text{-}sip}(r)) \end{array}}}
\vnbprop{sip-isometric-of-antisym-sq}{\vnbstmt{\begin{array}{@{}l@{}} \forall v.; \forall o \in (\operatorname{vec}(v) \to \operatorname{vec}(v)). \\ \quad (\operatorname{is\text{-}semi\text{-}inner\text{-}product\text{-}space}(v) \wedge (\operatorname{is\text{-}antisymmetric\text{-}op}(v, o) \wedge \operatorname{is\text{-}square\text{-}minus\text{-}id}(v, o))) \Rightarrow  \\ \quad \quad \operatorname{is\text{-}isometric\text{-}op}(v, o) \end{array}}}
\vnbprop{sip-antisym-of-isometric-sq}{\vnbstmt{\begin{array}{@{}l@{}} \forall v.; \forall o \in (\operatorname{vec}(v) \to \operatorname{vec}(v)). \\ \quad (\operatorname{is\text{-}semi\text{-}inner\text{-}product\text{-}space}(v) \wedge (\operatorname{is\text{-}isometric\text{-}op}(v, o) \wedge \operatorname{is\text{-}square\text{-}minus\text{-}id}(v, o))) \Rightarrow  \\ \quad \quad \operatorname{is\text{-}antisymmetric\text{-}op}(v, o) \end{array}}}
\vnbprop{phs-square-of-antisym-isometric}{\vnbstmt{\begin{array}{@{}l@{}} \forall v.; \forall o \in (\operatorname{vec}(v) \to \operatorname{vec}(v)). \\ \quad (\operatorname{is\text{-}pre\text{-}hilbert\text{-}space}(v) \wedge (\operatorname{is\text{-}antisymmetric\text{-}op}(v, o) \wedge \operatorname{is\text{-}isometric\text{-}op}(v, o))) \Rightarrow  \\ \quad \quad \operatorname{is\text{-}square\text{-}minus\text{-}id}(v, o) \end{array}}}
\vnbprop{csip-jop-fun}{\vnbstmt{\begin{array}{@{}l@{}} \forall v. \\ \quad \operatorname{is\text{-}complex\text{-}semi\text{-}inner\text{-}product\text{-}space}(v) \Rightarrow (J_{v} \in (\operatorname{vec}(v) \to \operatorname{vec}(v))) \end{array}}}
\vnbprop{csip-real-linear}{\vnbstmt{\begin{array}{@{}l@{}} \forall v. \\ \quad \operatorname{is\text{-}complex\text{-}semi\text{-}inner\text{-}product\text{-}space}(v) \Rightarrow \operatorname{is\text{-}real\text{-}linear\text{-}op}(v, J_{v}) \end{array}}}
\vnbprop{csip-square}{\vnbstmt{\begin{array}{@{}l@{}} \forall v. \\ \quad \operatorname{is\text{-}complex\text{-}semi\text{-}inner\text{-}product\text{-}space}(v) \Rightarrow \operatorname{is\text{-}square\text{-}minus\text{-}id}(v, J_{v}) \end{array}}}
\vnbprop{csip-antisym}{\vnbstmt{\begin{array}{@{}l@{}} \forall v. \\ \quad \operatorname{is\text{-}complex\text{-}semi\text{-}inner\text{-}product\text{-}space}(v) \Rightarrow \operatorname{is\text{-}antisymmetric\text{-}op}(v, J_{v}) \end{array}}}
\vnbprop{csip-jop-add}{\vnbstmt{\begin{array}{@{}l@{}} \forall v.; \forall x, y \in \operatorname{vec}(v). \\ \quad \operatorname{is\text{-}complex\text{-}semi\text{-}inner\text{-}product\text{-}space}(v) \Rightarrow  \\ \quad \quad J_{v}(\operatorname{vadd}(v)(x, y)) = \operatorname{vadd}(v)(J_{v}(x), J_{v}(y)) \end{array}}}
\vnbprop{csip-jop-act}{\vnbstmt{\begin{array}{@{}l@{}} \forall v.; \forall a \in \mathbb{R}, x \in \operatorname{vec}(v). \\ \quad \operatorname{is\text{-}complex\text{-}semi\text{-}inner\text{-}product\text{-}space}(v) \Rightarrow  \\ \quad \quad J_{v}(\operatorname{act}(v)(a, x)) = \operatorname{act}(v)(a, J_{v}(x)) \end{array}}}
\vnbprop{csip-jop-sq}{\vnbstmt{\begin{array}{@{}l@{}} \forall v.; \forall x \in \operatorname{vec}(v). \\ \quad \operatorname{is\text{-}complex\text{-}semi\text{-}inner\text{-}product\text{-}space}(v) \Rightarrow  \\ \quad \quad J_{v}(J_{v}(x)) = \operatorname{vneg}(v)(x) \end{array}}}
\vnbprop{csip-jop-antisym}{\vnbstmt{\begin{array}{@{}l@{}} \forall v.; \forall x, y \in \operatorname{vec}(v). \\ \quad \operatorname{is\text{-}complex\text{-}semi\text{-}inner\text{-}product\text{-}space}(v) \Rightarrow  \\ \quad \quad \langle\cdot\mid\cdot\rangle_{v}(J_{v}(x), y) = -\langle\cdot\mid\cdot\rangle_{v}(x, J_{v}(y)) \end{array}}}
\vnbprop{csip-jop-type}{\vnbstmt{\begin{array}{@{}l@{}} \forall v.; \forall x \in \operatorname{vec}(v). \\ \quad \operatorname{is\text{-}complex\text{-}semi\text{-}inner\text{-}product\text{-}space}(v) \Rightarrow (J_{v}(x) \in \operatorname{vec}(v)) \end{array}}}
\vnbprop{csip-nf-add}{\vnbstmt{\begin{array}{@{}l@{}} \forall v, p, q, r, s, x. \\ \quad \operatorname{is\text{-}complex\text{-}semi\text{-}inner\text{-}product\text{-}space}(v) \wedge \\ \quad p \in \mathbb{R} \wedge \\ \quad q \in \mathbb{R} \wedge \\ \quad r \in \mathbb{R} \wedge \\ \quad s \in \mathbb{R} \wedge \\ \quad x \in \operatorname{vec}(v) \Rightarrow \\ \quad \quad \operatorname{vadd}(v)(\operatorname{vadd}(v)(\operatorname{act}(v)(p, x), \operatorname{act}(v)(q, J_{v}(x))), \\ \quad \quad \quad \operatorname{vadd}(v)(\operatorname{act}(v)(r, x), \operatorname{act}(v)(s, J_{v}(x)))) \\ \quad \quad =\; \operatorname{vadd}(v)(\operatorname{act}(v)((p + r), x), \operatorname{act}(v)((q + s), J_{v}(x))) \end{array}}}
\vnbprop{csip-nf-scale}{\vnbstmt{\begin{array}{@{}l@{}} \forall v, t, p, q, x. \\ \quad (\operatorname{is\text{-}complex\text{-}semi\text{-}inner\text{-}product\text{-}space}(v) \wedge ((t \in \mathbb{R}) \wedge ((p \in \mathbb{R}) \wedge ((q \in \mathbb{R}) \wedge (x \in \operatorname{vec}(v)))))) \Rightarrow  \\ \quad \quad \operatorname{act}(v)(t, \operatorname{vadd}(v)(\operatorname{act}(v)(p, x), \operatorname{act}(v)(q, J_{v}(x)))) \\ \quad \quad =\; \operatorname{vadd}(v)(\operatorname{act}(v)((t \cdot p), x), \operatorname{act}(v)((t \cdot q), J_{v}(x))) \end{array}}}
\vnbprop{csip-nf-jop}{\vnbstmt{\begin{array}{@{}l@{}} \forall v, p, q, x. \\ \quad (\operatorname{is\text{-}complex\text{-}semi\text{-}inner\text{-}product\text{-}space}(v) \wedge ((p \in \mathbb{R}) \wedge ((q \in \mathbb{R}) \wedge (x \in \operatorname{vec}(v))))) \Rightarrow  \\ \quad \quad J_{v}(\operatorname{vadd}(v)(\operatorname{act}(v)(p, x), \operatorname{act}(v)(q, J_{v}(x)))) \\ \quad \quad =\; \operatorname{vadd}(v)(\operatorname{act}(v)(-q, x), \operatorname{act}(v)(p, J_{v}(x))) \end{array}}}
\vnbprop{csip-cact-value}{\vnbstmt{\begin{array}{@{}l@{}} \forall v.; \forall z \in \mathbb{C}, x \in \operatorname{vec}(v). \\ \quad \operatorname{is\text{-}complex\text{-}semi\text{-}inner\text{-}product\text{-}space}(v) \Rightarrow  \\ \quad \quad \operatorname{cact}(v)(z, x) = \operatorname{vadd}(v)(\operatorname{act}(v)(\operatorname{real\text{-}part}(z), x), \operatorname{act}(v)(\operatorname{imag\text{-}part}(z), J_{v}(x))) \end{array}}}
\vnbprop{csip-cact-type-pt}{\vnbstmt{\begin{array}{@{}l@{}} \forall v.; \forall z \in \mathbb{C}, x \in \operatorname{vec}(v). \\ \quad \operatorname{is\text{-}complex\text{-}semi\text{-}inner\text{-}product\text{-}space}(v) \Rightarrow (\operatorname{cact}(v)(z, x) \in \operatorname{vec}(v)) \end{array}}}
\vnbprop{csip-cact-fun}{\vnbstmt{\begin{array}{@{}l@{}} \forall v. \\ \quad \operatorname{is\text{-}complex\text{-}semi\text{-}inner\text{-}product\text{-}space}(v) \Rightarrow  \\ \quad \quad \operatorname{cact}(v) \in ((\mathbb{C} \times \operatorname{vec}(v)) \to \operatorname{vec}(v)) \end{array}}}
\vnbprop{csip-cact-vadd}{\vnbstmt{\begin{array}{@{}l@{}} \forall v.; \forall z \in \mathbb{C}, x, y \in \operatorname{vec}(v). \\ \quad \operatorname{is\text{-}complex\text{-}semi\text{-}inner\text{-}product\text{-}space}(v) \Rightarrow  \\ \quad \quad \operatorname{cact}(v)(z, \operatorname{vadd}(v)(x, y)) = \operatorname{vadd}(v)(\operatorname{cact}(v)(z, x), \operatorname{cact}(v)(z, y)) \end{array}}}
\vnbprop{csip-cact-add}{\vnbstmt{\begin{array}{@{}l@{}} \forall v.; \forall z, w \in \mathbb{C}, x \in \operatorname{vec}(v). \\ \quad \operatorname{is\text{-}complex\text{-}semi\text{-}inner\text{-}product\text{-}space}(v) \Rightarrow  \\ \quad \quad \operatorname{cact}(v)((z + w), x) = \operatorname{vadd}(v)(\operatorname{cact}(v)(z, x), \operatorname{cact}(v)(w, x)) \end{array}}}
\vnbprop{csip-cact-mul}{\vnbstmt{\begin{array}{@{}l@{}} \forall v.; \forall z, w \in \mathbb{C}, x \in \operatorname{vec}(v). \\ \quad \operatorname{is\text{-}complex\text{-}semi\text{-}inner\text{-}product\text{-}space}(v) \Rightarrow  \\ \quad \quad \operatorname{cact}(v)((z \cdot w), x) = \operatorname{cact}(v)(z, \operatorname{cact}(v)(w, x)) \end{array}}}
\vnbprop{csip-cact-one}{\vnbstmt{\begin{array}{@{}l@{}} \forall v.; \forall x \in \operatorname{vec}(v). \\ \quad \operatorname{is\text{-}complex\text{-}semi\text{-}inner\text{-}product\text{-}space}(v) \Rightarrow (\operatorname{cact}(v)(1, x) = x) \end{array}}}
\vnbprop{csip-crip-value}{\vnbstmt{\begin{array}{@{}l@{}} \forall v, x, y. \\ \quad (\operatorname{is\text{-}complex\text{-}semi\text{-}inner\text{-}product\text{-}space}(v) \wedge ((x \in \operatorname{vec}(v)) \wedge (y \in \operatorname{vec}(v)))) \Rightarrow  \\ \quad \quad \operatorname{crip}(v)(x, y) = (\langle\cdot\mid\cdot\rangle_{v}(x, y) - (+i \cdot \langle\cdot\mid\cdot\rangle_{v}(J_{v}(x), y))) \end{array}}}
\vnbprop{csip-crip-type-pt}{\vnbstmt{\begin{array}{@{}l@{}} \forall v.; \forall x, y \in \operatorname{vec}(v). \\ \quad \operatorname{is\text{-}complex\text{-}semi\text{-}inner\text{-}product\text{-}space}(v) \Rightarrow (\operatorname{crip}(v)(x, y) \in \mathbb{C}) \end{array}}}
\vnbprop{csip-crip-fun}{\vnbstmt{\begin{array}{@{}l@{}} \forall v. \\ \quad \operatorname{is\text{-}complex\text{-}semi\text{-}inner\text{-}product\text{-}space}(v) \Rightarrow  \\ \quad \quad \operatorname{crip}(v) \in ((\operatorname{vec}(v) \times \operatorname{vec}(v)) \to \mathbb{C}) \end{array}}}
\vnbprop{csip-crip-add-left}{\vnbstmt{\begin{array}{@{}l@{}} \forall v.; \forall x, y, z \in \operatorname{vec}(v). \\ \quad \operatorname{is\text{-}complex\text{-}semi\text{-}inner\text{-}product\text{-}space}(v) \Rightarrow  \\ \quad \quad \operatorname{crip}(v)(\operatorname{vadd}(v)(x, y), z) = (\operatorname{crip}(v)(x, z) + \operatorname{crip}(v)(y, z)) \end{array}}}
\vnbprop{csip-crip-homog-left}{\vnbstmt{\begin{array}{@{}l@{}} \forall v.; \forall z \in \mathbb{C}, x, y \in \operatorname{vec}(v). \\ \quad \operatorname{is\text{-}complex\text{-}semi\text{-}inner\text{-}product\text{-}space}(v) \Rightarrow  \\ \quad \quad \operatorname{crip}(v)(\operatorname{cact}(v)(z, x), y) = (z \cdot \operatorname{crip}(v)(x, y)) \end{array}}}
\vnbprop{csip-crip-conj-sym}{\vnbstmt{\begin{array}{@{}l@{}} \forall v.; \forall x, y \in \operatorname{vec}(v). \\ \quad \operatorname{is\text{-}complex\text{-}semi\text{-}inner\text{-}product\text{-}space}(v) \Rightarrow  \\ \quad \quad \operatorname{crip}(v)(y, x) = \operatorname{conjugate}(\operatorname{crip}(v)(x, y)) \end{array}}}
\vnbprop{csip-crip-self}{\vnbstmt{\begin{array}{@{}l@{}} \forall v.; \forall x \in \operatorname{vec}(v). \\ \quad \operatorname{is\text{-}complex\text{-}semi\text{-}inner\text{-}product\text{-}space}(v) \Rightarrow (\operatorname{crip}(v)(x, x) = \langle\cdot\mid\cdot\rangle_{v}(x, x)) \end{array}}}
\vnbprop{sip-cc-ring-carr}{\vnbstmt{\operatorname{carr}(\operatorname{normed\text{-}field\text{-}as\text{-}commutative\text{-}ring}(\operatorname{cc\text{-}normed\text{-}field})) = \mathbb{C}}}
\vnbprop{sip-cc-ring-one}{\vnbstmt{\operatorname{one}(\operatorname{normed\text{-}field\text{-}as\text{-}commutative\text{-}ring}(\operatorname{cc\text{-}normed\text{-}field})) = 1}}
\vnbprop{sip-cc-ring-add-value}{\vnbstmt{\begin{array}{@{}l@{}} \forall a, b. \\ \quad ((a \in \mathbb{C}) \wedge (b \in \mathbb{C})) \Rightarrow  \\ \quad \quad \operatorname{add}(\operatorname{normed\text{-}field\text{-}as\text{-}commutative\text{-}ring}(\operatorname{cc\text{-}normed\text{-}field}))(a, b) = (a + b) \end{array}}}
\vnbprop{sip-cc-ring-mul-value}{\vnbstmt{\begin{array}{@{}l@{}} \forall a, b. \\ \quad ((a \in \mathbb{C}) \wedge (b \in \mathbb{C})) \Rightarrow  \\ \quad \quad \operatorname{mul}(\operatorname{normed\text{-}field\text{-}as\text{-}commutative\text{-}ring}(\operatorname{cc\text{-}normed\text{-}field}))(a, b) = (a \cdot b) \end{array}}}
\vnbprop{sip-cc-ring-neg-value}{\vnbstmt{\begin{array}{@{}l@{}} \forall a \in \mathbb{C}. \\ \quad \operatorname{neg}(\operatorname{normed\text{-}field\text{-}as\text{-}commutative\text{-}ring}(\operatorname{cc\text{-}normed\text{-}field}))(a) = -a \end{array}}}
\vnbprop{cphs-is-csip}{\vnbstmt{\begin{array}{@{}l@{}} \forall v. \\ \quad \operatorname{is\text{-}complex\text{-}pre\text{-}hilbert\text{-}space}(v) \Rightarrow \operatorname{is\text{-}complex\text{-}semi\text{-}inner\text{-}product\text{-}space}(v) \end{array}}}
\vnbprop{cphs-definite}{\vnbstmt{\begin{array}{@{}l@{}} \forall v.; \forall x \in \operatorname{vec}(v). \\ \quad (\operatorname{is\text{-}complex\text{-}pre\text{-}hilbert\text{-}space}(v) \wedge (\langle\cdot\mid\cdot\rangle_{v}(x, x) = 0)) \Rightarrow (x = \operatorname{vzero}(v)) \end{array}}}
\vnbprop{csip-cips-scal}{\vnbstmt{\begin{array}{@{}l@{}} \forall v. \\ \quad \operatorname{is\text{-}complex\text{-}semi\text{-}inner\text{-}product\text{-}space}(v) \Rightarrow  \\ \quad \quad (\operatorname{scal}(\operatorname{csip\text{-}cips}(v)) \simeq \operatorname{normed\text{-}field\text{-}as\text{-}commutative\text{-}ring}(\operatorname{cc\text{-}normed\text{-}field})) \end{array}}}
\vnbprop{csip-cips-vec}{\vnbstmt{\begin{array}{@{}l@{}} \forall v. \\ \quad \operatorname{is\text{-}complex\text{-}semi\text{-}inner\text{-}product\text{-}space}(v) \Rightarrow (\operatorname{vec}(\operatorname{csip\text{-}cips}(v)) \simeq \operatorname{vec}(v)) \end{array}}}
\vnbprop{csip-cips-vadd}{\vnbstmt{\begin{array}{@{}l@{}} \forall v. \\ \quad \operatorname{is\text{-}complex\text{-}semi\text{-}inner\text{-}product\text{-}space}(v) \Rightarrow (\operatorname{vadd}(\operatorname{csip\text{-}cips}(v)) \simeq \operatorname{vadd}(v)) \end{array}}}
\vnbprop{csip-cips-vzero}{\vnbstmt{\begin{array}{@{}l@{}} \forall v. \\ \quad \operatorname{is\text{-}complex\text{-}semi\text{-}inner\text{-}product\text{-}space}(v) \Rightarrow (\operatorname{vzero}(\operatorname{csip\text{-}cips}(v)) \simeq \operatorname{vzero}(v)) \end{array}}}
\vnbprop{csip-cips-vneg}{\vnbstmt{\begin{array}{@{}l@{}} \forall v. \\ \quad \operatorname{is\text{-}complex\text{-}semi\text{-}inner\text{-}product\text{-}space}(v) \Rightarrow (\operatorname{vneg}(\operatorname{csip\text{-}cips}(v)) \simeq \operatorname{vneg}(v)) \end{array}}}
\vnbprop{csip-cips-act}{\vnbstmt{\begin{array}{@{}l@{}} \forall v. \\ \quad \operatorname{is\text{-}complex\text{-}semi\text{-}inner\text{-}product\text{-}space}(v) \Rightarrow (\operatorname{act}(\operatorname{csip\text{-}cips}(v)) \simeq \operatorname{cact}(v)) \end{array}}}
\vnbprop{csip-cips-ip}{\vnbstmt{\begin{array}{@{}l@{}} \forall v. \\ \quad \operatorname{is\text{-}complex\text{-}semi\text{-}inner\text{-}product\text{-}space}(v) \Rightarrow (\langle\cdot,\cdot\rangle_{\operatorname{csip\text{-}cips}(v)} \simeq \operatorname{crip}(v)) \end{array}}}
\vnbprop{csip-cips-length}{\vnbstmt{\forall v.\; (\operatorname{is\text{-}complex\text{-}pre\text{-}hilbert\text{-}space}(v) \Rightarrow (\operatorname{length}(\operatorname{csip\text{-}cips}(v)) = 7))}}
\vnbprop{csip-cips-scal-ring}{\vnbstmt{\forall v.\; (\operatorname{is\text{-}complex\text{-}pre\text{-}hilbert\text{-}space}(v) \Rightarrow \operatorname{is\text{-}ring}(\operatorname{scal}(\operatorname{csip\text{-}cips}(v))))}}
\vnbprop{csip-cips-vec-set}{\vnbstmt{\forall v.\; (\operatorname{is\text{-}complex\text{-}pre\text{-}hilbert\text{-}space}(v) \Rightarrow (\operatorname{vec}(\operatorname{csip\text{-}cips}(v)) \in \mathbf{Set}))}}
\vnbprop{csip-cips-vadd-fun}{\vnbstmt{\begin{array}{@{}l@{}} \forall v. \\ \quad \operatorname{is\text{-}complex\text{-}pre\text{-}hilbert\text{-}space}(v) \Rightarrow  \\ \quad \quad \operatorname{vadd}(\operatorname{csip\text{-}cips}(v)) \in ((\operatorname{vec}(\operatorname{csip\text{-}cips}(v)) \times \operatorname{vec}(\operatorname{csip\text{-}cips}(v))) \to \operatorname{vec}(\operatorname{csip\text{-}cips}(v))) \end{array}}}
\vnbprop{csip-cips-vzero-in}{\vnbstmt{\begin{array}{@{}l@{}} \forall v. \\ \quad \operatorname{is\text{-}complex\text{-}pre\text{-}hilbert\text{-}space}(v) \Rightarrow (\operatorname{vzero}(\operatorname{csip\text{-}cips}(v)) \in \operatorname{vec}(\operatorname{csip\text{-}cips}(v))) \end{array}}}
\vnbprop{csip-cips-vneg-fun}{\vnbstmt{\begin{array}{@{}l@{}} \forall v. \\ \quad \operatorname{is\text{-}complex\text{-}pre\text{-}hilbert\text{-}space}(v) \Rightarrow  \\ \quad \quad \operatorname{vneg}(\operatorname{csip\text{-}cips}(v)) \in (\operatorname{vec}(\operatorname{csip\text{-}cips}(v)) \to \operatorname{vec}(\operatorname{csip\text{-}cips}(v))) \end{array}}}
\vnbprop{csip-cips-act-fun}{\vnbstmt{\begin{array}{@{}l@{}} \forall v. \\ \quad \operatorname{is\text{-}complex\text{-}pre\text{-}hilbert\text{-}space}(v) \Rightarrow  \\ \quad \quad \operatorname{act}(\operatorname{csip\text{-}cips}(v)) \in ((\operatorname{carr}(\operatorname{scal}(\operatorname{csip\text{-}cips}(v))) \times \operatorname{vec}(\operatorname{csip\text{-}cips}(v))) \to \operatorname{vec}(\operatorname{csip\text{-}cips}(v))) \end{array}}}
\vnbprop{csip-cips-ip-fun}{\vnbstmt{\begin{array}{@{}l@{}} \forall v. \\ \quad \operatorname{is\text{-}complex\text{-}pre\text{-}hilbert\text{-}space}(v) \Rightarrow  \\ \quad \quad \langle\cdot,\cdot\rangle_{\operatorname{csip\text{-}cips}(v)} \in ((\operatorname{vec}(\operatorname{csip\text{-}cips}(v)) \times \operatorname{vec}(\operatorname{csip\text{-}cips}(v))) \to \mathbb{C}) \end{array}}}
\vnbprop{csip-cips-assoc}{\vnbstmt{\begin{array}{@{}l@{}} \forall v. \\ \quad \operatorname{is\text{-}complex\text{-}pre\text{-}hilbert\text{-}space}(v) \Rightarrow  \\ \quad \quad \operatorname{is\text{-}associative}(\operatorname{vadd}(\operatorname{csip\text{-}cips}(v)), \operatorname{vec}(\operatorname{csip\text{-}cips}(v))) \end{array}}}
\vnbprop{csip-cips-comm}{\vnbstmt{\begin{array}{@{}l@{}} \forall v. \\ \quad \operatorname{is\text{-}complex\text{-}pre\text{-}hilbert\text{-}space}(v) \Rightarrow  \\ \quad \quad \operatorname{is\text{-}commutative}(\operatorname{vadd}(\operatorname{csip\text{-}cips}(v)), \operatorname{vec}(\operatorname{csip\text{-}cips}(v))) \end{array}}}
\vnbprop{csip-cips-identity}{\vnbstmt{\begin{array}{@{}l@{}} \forall v. \\ \quad \operatorname{is\text{-}complex\text{-}pre\text{-}hilbert\text{-}space}(v) \Rightarrow  \\ \quad \quad \operatorname{is\text{-}identity}(\operatorname{vadd}(\operatorname{csip\text{-}cips}(v)), \operatorname{vzero}(\operatorname{csip\text{-}cips}(v)), \operatorname{vec}(\operatorname{csip\text{-}cips}(v))) \end{array}}}
\vnbprop{csip-cips-inverses}{\vnbstmt{\begin{array}{@{}l@{}} \forall v. \\ \quad \operatorname{is\text{-}complex\text{-}pre\text{-}hilbert\text{-}space}(v) \Rightarrow  \\ \quad \quad \operatorname{has\text{-}inverses}(\operatorname{vadd}(\operatorname{csip\text{-}cips}(v)), \operatorname{vzero}(\operatorname{csip\text{-}cips}(v)), \operatorname{vneg}(\operatorname{csip\text{-}cips}(v)), \operatorname{vec}(\operatorname{csip\text{-}cips}(v))) \end{array}}}
\vnbprop{csip-cips-scal-pin}{\vnbstmt{\begin{array}{@{}l@{}} \forall v. \\ \quad \operatorname{is\text{-}complex\text{-}pre\text{-}hilbert\text{-}space}(v) \Rightarrow  \\ \quad \quad \operatorname{scal}(\operatorname{csip\text{-}cips}(v)) = \operatorname{normed\text{-}field\text{-}as\text{-}commutative\text{-}ring}(\operatorname{cc\text{-}normed\text{-}field}) \end{array}}}
\vnbprop{csip-cips-scal-carr}{\vnbstmt{\begin{array}{@{}l@{}} \forall v. \\ \quad \operatorname{is\text{-}complex\text{-}pre\text{-}hilbert\text{-}space}(v) \Rightarrow (\operatorname{carr}(\operatorname{scal}(\operatorname{csip\text{-}cips}(v))) = \mathbb{C}) \end{array}}}
\vnbprop{csip-cips-act-vadd}{\vnbstmt{\begin{array}{@{}l@{}} \forall v.; \forall r \in \operatorname{carr}(\operatorname{scal}(\operatorname{csip\text{-}cips}(v))), x, y \in \operatorname{vec}(\operatorname{csip\text{-}cips}(v)). \\ \quad \operatorname{is\text{-}complex\text{-}pre\text{-}hilbert\text{-}space}(v) \Rightarrow  \\ \quad \quad \operatorname{act}(\operatorname{csip\text{-}cips}(v))(r, \operatorname{vadd}(\operatorname{csip\text{-}cips}(v))(x, y)) \\ \quad \quad =\; \operatorname{vadd}(\operatorname{csip\text{-}cips}(v))(\operatorname{act}(\operatorname{csip\text{-}cips}(v))(r, x), \operatorname{act}(\operatorname{csip\text{-}cips}(v))(r, y)) \end{array}}}
\vnbprop{csip-cips-act-add}{\vnbstmt{\begin{array}{@{}l@{}} \forall v.; \forall r, s \in \operatorname{carr}(\operatorname{scal}(\operatorname{csip\text{-}cips}(v))), x \in \operatorname{vec}(\operatorname{csip\text{-}cips}(v)). \\ \quad \operatorname{is\text{-}complex\text{-}pre\text{-}hilbert\text{-}space}(v) \Rightarrow  \\ \quad \quad \operatorname{act}(\operatorname{csip\text{-}cips}(v))(\operatorname{add}(\operatorname{scal}(\operatorname{csip\text{-}cips}(v)))(r, s), x) \\ \quad \quad =\; \operatorname{vadd}(\operatorname{csip\text{-}cips}(v))(\operatorname{act}(\operatorname{csip\text{-}cips}(v))(r, x), \operatorname{act}(\operatorname{csip\text{-}cips}(v))(s, x)) \end{array}}}
\vnbprop{csip-cips-act-mul}{\vnbstmt{\begin{array}{@{}l@{}} \forall v.; \forall r, s \in \operatorname{carr}(\operatorname{scal}(\operatorname{csip\text{-}cips}(v))), x \in \operatorname{vec}(\operatorname{csip\text{-}cips}(v)). \\ \quad \operatorname{is\text{-}complex\text{-}pre\text{-}hilbert\text{-}space}(v) \Rightarrow  \\ \quad \quad \operatorname{act}(\operatorname{csip\text{-}cips}(v))(\operatorname{mul}(\operatorname{scal}(\operatorname{csip\text{-}cips}(v)))(r, s), x) \\ \quad \quad =\; \operatorname{act}(\operatorname{csip\text{-}cips}(v))(r, \operatorname{act}(\operatorname{csip\text{-}cips}(v))(s, x)) \end{array}}}
\vnbprop{csip-cips-act-one}{\vnbstmt{\begin{array}{@{}l@{}} \forall v.; \forall x \in \operatorname{vec}(\operatorname{csip\text{-}cips}(v)). \\ \quad \operatorname{is\text{-}complex\text{-}pre\text{-}hilbert\text{-}space}(v) \Rightarrow  \\ \quad \quad \operatorname{act}(\operatorname{csip\text{-}cips}(v))(\operatorname{one}(\operatorname{scal}(\operatorname{csip\text{-}cips}(v))), x) = x \end{array}}}
\vnbprop{csip-cips-ip-add-left}{\vnbstmt{\begin{array}{@{}l@{}} \forall v.; \forall x, y, z \in \operatorname{vec}(\operatorname{csip\text{-}cips}(v)). \\ \quad \operatorname{is\text{-}complex\text{-}pre\text{-}hilbert\text{-}space}(v) \Rightarrow  \\ \quad \quad \langle\cdot,\cdot\rangle_{\operatorname{csip\text{-}cips}(v)}(\operatorname{vadd}(\operatorname{csip\text{-}cips}(v))(x, y), z) \\ \quad \quad =\; \langle\cdot,\cdot\rangle_{\operatorname{csip\text{-}cips}(v)}(x, z) + \langle\cdot,\cdot\rangle_{\operatorname{csip\text{-}cips}(v)}(y, z) \end{array}}}
\vnbprop{csip-cips-ip-homog-left}{\vnbstmt{\begin{array}{@{}l@{}} \forall v.; \forall a \in \mathbb{C}, x, y \in \operatorname{vec}(\operatorname{csip\text{-}cips}(v)). \\ \quad \operatorname{is\text{-}complex\text{-}pre\text{-}hilbert\text{-}space}(v) \Rightarrow  \\ \quad \quad \langle\cdot,\cdot\rangle_{\operatorname{csip\text{-}cips}(v)}(\operatorname{act}(\operatorname{csip\text{-}cips}(v))(a, x), y) = (a \cdot \langle\cdot,\cdot\rangle_{\operatorname{csip\text{-}cips}(v)}(x, y)) \end{array}}}
\vnbprop{csip-cips-ip-conj-sym}{\vnbstmt{\begin{array}{@{}l@{}} \forall v.; \forall x, y \in \operatorname{vec}(\operatorname{csip\text{-}cips}(v)). \\ \quad \operatorname{is\text{-}complex\text{-}pre\text{-}hilbert\text{-}space}(v) \Rightarrow  \\ \quad \quad \langle\cdot,\cdot\rangle_{\operatorname{csip\text{-}cips}(v)}(y, x) = \operatorname{conjugate}(\langle\cdot,\cdot\rangle_{\operatorname{csip\text{-}cips}(v)}(x, y)) \end{array}}}
\vnbprop{csip-cips-ip-self-real}{\vnbstmt{\begin{array}{@{}l@{}} \forall v.; \forall x \in \operatorname{vec}(\operatorname{csip\text{-}cips}(v)). \\ \quad \operatorname{is\text{-}complex\text{-}pre\text{-}hilbert\text{-}space}(v) \Rightarrow (\langle\cdot,\cdot\rangle_{\operatorname{csip\text{-}cips}(v)}(x, x) \in \mathbb{R}) \end{array}}}
\vnbprop{csip-cips-ip-self-nonneg}{\vnbstmt{\begin{array}{@{}l@{}} \forall v.; \forall x \in \operatorname{vec}(\operatorname{csip\text{-}cips}(v)). \\ \quad \operatorname{is\text{-}complex\text{-}pre\text{-}hilbert\text{-}space}(v) \Rightarrow (0 \leq \langle\cdot,\cdot\rangle_{\operatorname{csip\text{-}cips}(v)}(x, x)) \end{array}}}
\vnbprop{csip-cips-ip-definite}{\vnbstmt{\begin{array}{@{}l@{}} \forall v.; \forall x \in \operatorname{vec}(\operatorname{csip\text{-}cips}(v)). \\ \quad \operatorname{is\text{-}complex\text{-}pre\text{-}hilbert\text{-}space}(v) \Rightarrow  \\ \quad \quad (\langle\cdot,\cdot\rangle_{\operatorname{csip\text{-}cips}(v)}(x, x) = 0) \Leftrightarrow (x = \operatorname{vzero}(\operatorname{csip\text{-}cips}(v))) \end{array}}}
\vnbprop{cphs-is-complex-inner-product-space}{\vnbstmt{\begin{array}{@{}l@{}} \forall v. \\ \quad \operatorname{is\text{-}complex\text{-}pre\text{-}hilbert\text{-}space}(v) \Rightarrow  \\ \quad \quad \operatorname{is\text{-}complex\text{-}inner\text{-}product\text{-}space}(\operatorname{csip\text{-}cips}(v)) \end{array}}}
\vnbprop{cc-re-of-real}{\vnbstmt{\forall z \in \mathbb{R}.\; (\operatorname{real\text{-}part}(z) = z)}}
\vnbprop{cc-re-im-neg-i}{\vnbstmt{(\operatorname{real\text{-}part}(-+i) = 0) \wedge (\operatorname{imag\text{-}part}(-+i) = -1)}}
\vnbprop{cc-re-i-mul}{\vnbstmt{\forall z \in \mathbb{C}.\; (\operatorname{real\text{-}part}((+i \cdot z)) = -\operatorname{imag\text{-}part}(z))}}
\vnbprop{cc-re-negi-mul}{\vnbstmt{\forall z \in \mathbb{C}.\; (\operatorname{real\text{-}part}((-+i \cdot z)) = \operatorname{imag\text{-}part}(z))}}
\vnbprop{cc-re-conj}{\vnbstmt{\forall z \in \mathbb{C}.\; (\operatorname{real\text{-}part}(\operatorname{conjugate}(z)) = \operatorname{real\text{-}part}(z))}}
\vnbprop{cips-vec-set}{\vnbstmt{\forall w.\; (\operatorname{is\text{-}complex\text{-}inner\text{-}product\text{-}space}(w) \Rightarrow (\operatorname{vec}(w) \in \mathbf{Set}))}}
\vnbprop{cips-vadd-fun}{\vnbstmt{\begin{array}{@{}l@{}} \forall w. \\ \quad \operatorname{is\text{-}complex\text{-}inner\text{-}product\text{-}space}(w) \Rightarrow  \\ \quad \quad \operatorname{vadd}(w) \in ((\operatorname{vec}(w) \times \operatorname{vec}(w)) \to \operatorname{vec}(w)) \end{array}}}
\vnbprop{cips-vneg-fun}{\vnbstmt{\begin{array}{@{}l@{}} \forall w. \\ \quad \operatorname{is\text{-}complex\text{-}inner\text{-}product\text{-}space}(w) \Rightarrow (\operatorname{vneg}(w) \in (\operatorname{vec}(w) \to \operatorname{vec}(w))) \end{array}}}
\vnbprop{cips-assoc-prop}{\vnbstmt{\begin{array}{@{}l@{}} \forall w. \\ \quad \operatorname{is\text{-}complex\text{-}inner\text{-}product\text{-}space}(w) \Rightarrow \operatorname{is\text{-}associative}(\operatorname{vadd}(w), \operatorname{vec}(w)) \end{array}}}
\vnbprop{cips-comm-prop}{\vnbstmt{\begin{array}{@{}l@{}} \forall w. \\ \quad \operatorname{is\text{-}complex\text{-}inner\text{-}product\text{-}space}(w) \Rightarrow \operatorname{is\text{-}commutative}(\operatorname{vadd}(w), \operatorname{vec}(w)) \end{array}}}
\vnbprop{cips-identity-prop}{\vnbstmt{\begin{array}{@{}l@{}} \forall w. \\ \quad \operatorname{is\text{-}complex\text{-}inner\text{-}product\text{-}space}(w) \Rightarrow \operatorname{is\text{-}identity}(\operatorname{vadd}(w), \operatorname{vzero}(w), \operatorname{vec}(w)) \end{array}}}
\vnbprop{cips-inverses-prop}{\vnbstmt{\begin{array}{@{}l@{}} \forall w. \\ \quad \operatorname{is\text{-}complex\text{-}inner\text{-}product\text{-}space}(w) \Rightarrow  \\ \quad \quad \operatorname{has\text{-}inverses}(\operatorname{vadd}(w), \operatorname{vzero}(w), \operatorname{vneg}(w), \operatorname{vec}(w)) \end{array}}}
\vnbprop{cips-scal-pin}{\vnbstmt{\begin{array}{@{}l@{}} \forall w. \\ \quad \operatorname{is\text{-}complex\text{-}inner\text{-}product\text{-}space}(w) \Rightarrow  \\ \quad \quad \operatorname{scal}(w) = \operatorname{normed\text{-}field\text{-}as\text{-}commutative\text{-}ring}(\operatorname{cc\text{-}normed\text{-}field}) \end{array}}}
\vnbprop{cips-module-view-scal}{\vnbstmt{\begin{array}{@{}l@{}} \forall w. \\ \quad \operatorname{is\text{-}complex\text{-}inner\text{-}product\text{-}space}(w) \Rightarrow  \\ \quad \quad (\operatorname{scal}(\operatorname{complex\text{-}inner\text{-}product\text{-}space\text{-}as\text{-}module}(w)) \simeq \operatorname{scal}(w)) \end{array}}}
\vnbprop{cips-module-view-vec}{\vnbstmt{\begin{array}{@{}l@{}} \forall w. \\ \quad \operatorname{is\text{-}complex\text{-}inner\text{-}product\text{-}space}(w) \Rightarrow  \\ \quad \quad (\operatorname{vec}(\operatorname{complex\text{-}inner\text{-}product\text{-}space\text{-}as\text{-}module}(w)) \simeq \operatorname{vec}(w)) \end{array}}}
\vnbprop{cips-module-view-vneg}{\vnbstmt{\begin{array}{@{}l@{}} \forall w. \\ \quad \operatorname{is\text{-}complex\text{-}inner\text{-}product\text{-}space}(w) \Rightarrow  \\ \quad \quad (\operatorname{vneg}(\operatorname{complex\text{-}inner\text{-}product\text{-}space\text{-}as\text{-}module}(w)) \simeq \operatorname{vneg}(w)) \end{array}}}
\vnbprop{cips-module-view-act}{\vnbstmt{\begin{array}{@{}l@{}} \forall w. \\ \quad \operatorname{is\text{-}complex\text{-}inner\text{-}product\text{-}space}(w) \Rightarrow  \\ \quad \quad (\operatorname{act}(\operatorname{complex\text{-}inner\text{-}product\text{-}space\text{-}as\text{-}module}(w)) \simeq \operatorname{act}(w)) \end{array}}}
\vnbprop{cips-act-neg-one}{\vnbstmt{\begin{array}{@{}l@{}} \forall w.; \forall x \in \operatorname{vec}(w). \\ \quad \operatorname{is\text{-}complex\text{-}inner\text{-}product\text{-}space}(w) \Rightarrow (\operatorname{act}(w)(-1, x) = \operatorname{vneg}(w)(x)) \end{array}}}
\vnbprop{cips-j-add}{\vnbstmt{\begin{array}{@{}l@{}} \forall w.; \forall a, b \in \operatorname{vec}(w). \\ \quad \operatorname{is\text{-}complex\text{-}inner\text{-}product\text{-}space}(w) \Rightarrow  \\ \quad \quad \operatorname{act}(w)(+i, \operatorname{vadd}(w)(a, b)) = \operatorname{vadd}(w)(\operatorname{act}(w)(+i, a), \operatorname{act}(w)(+i, b)) \end{array}}}
\vnbprop{cips-j-act}{\vnbstmt{\begin{array}{@{}l@{}} \forall w.; \forall a \in \mathbb{R}, \mathit{ja} \in \operatorname{vec}(w). \\ \quad \operatorname{is\text{-}complex\text{-}inner\text{-}product\text{-}space}(w) \Rightarrow  \\ \quad \quad \operatorname{act}(w)(+i, \operatorname{act}(w)(a, \mathit{ja})) = \operatorname{act}(w)(a, \operatorname{act}(w)(+i, \mathit{ja})) \end{array}}}
\vnbprop{cips-j-sq}{\vnbstmt{\begin{array}{@{}l@{}} \forall w.; \forall a \in \operatorname{vec}(w). \\ \quad \operatorname{is\text{-}complex\text{-}inner\text{-}product\text{-}space}(w) \Rightarrow  \\ \quad \quad \operatorname{act}(w)(+i, \operatorname{act}(w)(+i, a)) = \operatorname{vneg}(w)(a) \end{array}}}
\vnbprop{cips-j-antisym}{\vnbstmt{\begin{array}{@{}l@{}} \forall w.; \forall a, b \in \operatorname{vec}(w). \\ \quad \operatorname{is\text{-}complex\text{-}inner\text{-}product\text{-}space}(w) \Rightarrow  \\ \quad \quad \operatorname{real\text{-}part}(\langle\cdot,\cdot\rangle_{w}(\operatorname{act}(w)(+i, a), b)) = -\operatorname{real\text{-}part}(\langle\cdot,\cdot\rangle_{w}(a, \operatorname{act}(w)(+i, b))) \end{array}}}
\vnbprop{cips-re-self}{\vnbstmt{\begin{array}{@{}l@{}} \forall w.; \forall a \in \operatorname{vec}(w). \\ \quad \operatorname{is\text{-}complex\text{-}inner\text{-}product\text{-}space}(w) \Rightarrow  \\ \quad \quad \operatorname{real\text{-}part}(\langle\cdot,\cdot\rangle_{w}(a, a)) = \langle\cdot,\cdot\rangle_{w}(a, a) \end{array}}}
\vnbprop{cips-re-sym}{\vnbstmt{\begin{array}{@{}l@{}} \forall w.; \forall a, b \in \operatorname{vec}(w). \\ \quad \operatorname{is\text{-}complex\text{-}inner\text{-}product\text{-}space}(w) \Rightarrow  \\ \quad \quad \operatorname{real\text{-}part}(\langle\cdot,\cdot\rangle_{w}(a, b)) = \operatorname{real\text{-}part}(\langle\cdot,\cdot\rangle_{w}(b, a)) \end{array}}}
\vnbprop{cips-re-add-left}{\vnbstmt{\begin{array}{@{}l@{}} \forall w.; \forall a, b, c \in \operatorname{vec}(w). \\ \quad \operatorname{is\text{-}complex\text{-}inner\text{-}product\text{-}space}(w) \Rightarrow  \\ \quad \quad \operatorname{real\text{-}part}(\langle\cdot,\cdot\rangle_{w}(\operatorname{vadd}(w)(a, b), c)) = (\operatorname{real\text{-}part}(\langle\cdot,\cdot\rangle_{w}(a, c)) + \operatorname{real\text{-}part}(\langle\cdot,\cdot\rangle_{w}(b, c))) \end{array}}}
\vnbprop{cips-re-homog-left}{\vnbstmt{\begin{array}{@{}l@{}} \forall w.; \forall a \in \mathbb{R}, \mathit{ja}, b \in \operatorname{vec}(w). \\ \quad \operatorname{is\text{-}complex\text{-}inner\text{-}product\text{-}space}(w) \Rightarrow  \\ \quad \quad \operatorname{real\text{-}part}(\langle\cdot,\cdot\rangle_{w}(\operatorname{act}(w)(a, \mathit{ja}), b)) = (a \cdot \operatorname{real\text{-}part}(\langle\cdot,\cdot\rangle_{w}(\mathit{ja}, b))) \end{array}}}
\vnbprop{cips-ract-add}{\vnbstmt{\begin{array}{@{}l@{}} \forall w.; \forall a, \mathit{sa} \in \mathbb{R}, \mathit{ja} \in \operatorname{vec}(w). \\ \quad \operatorname{is\text{-}complex\text{-}inner\text{-}product\text{-}space}(w) \Rightarrow  \\ \quad \quad \operatorname{act}(w)((a + \mathit{sa}), \mathit{ja}) = \operatorname{vadd}(w)(\operatorname{act}(w)(a, \mathit{ja}), \operatorname{act}(w)(\mathit{sa}, \mathit{ja})) \end{array}}}
\vnbprop{cips-ract-mul}{\vnbstmt{\begin{array}{@{}l@{}} \forall w.; \forall a, \mathit{sa} \in \mathbb{R}, \mathit{ja} \in \operatorname{vec}(w). \\ \quad \operatorname{is\text{-}complex\text{-}inner\text{-}product\text{-}space}(w) \Rightarrow  \\ \quad \quad \operatorname{act}(w)((a \cdot \mathit{sa}), \mathit{ja}) = \operatorname{act}(w)(a, \operatorname{act}(w)(\mathit{sa}, \mathit{ja})) \end{array}}}
\vnbprop{cips-ract-one}{\vnbstmt{\begin{array}{@{}l@{}} \forall w.; \forall a \in \operatorname{vec}(w). \\ \quad \operatorname{is\text{-}complex\text{-}inner\text{-}product\text{-}space}(w) \Rightarrow (\operatorname{act}(w)(1, a) = a) \end{array}}}
\vnbprop{sip-rr-ring-add-value}{\vnbstmt{\begin{array}{@{}l@{}} \forall p, q \in \mathbb{R}. \\ \quad \operatorname{add}(\operatorname{normed\text{-}field\text{-}as\text{-}commutative\text{-}ring}(\operatorname{rr\text{-}normed\text{-}field}))(p, q) = (p + q) \end{array}}}
\vnbprop{sip-rr-ring-mul-value}{\vnbstmt{\begin{array}{@{}l@{}} \forall p, q \in \mathbb{R}. \\ \quad \operatorname{mul}(\operatorname{normed\text{-}field\text{-}as\text{-}commutative\text{-}ring}(\operatorname{rr\text{-}normed\text{-}field}))(p, q) = (p \cdot q) \end{array}}}
\vnbprop{cips-ract-fun}{\vnbstmt{\begin{array}{@{}l@{}} \forall w. \\ \quad \operatorname{is\text{-}complex\text{-}inner\text{-}product\text{-}space}(w) \Rightarrow  \\ \quad \quad \operatorname{cips\text{-}ract}(w) \in ((\mathbb{R} \times \operatorname{vec}(w)) \to \operatorname{vec}(w)) \end{array}}}
\vnbprop{cips-rip-fun}{\vnbstmt{\begin{array}{@{}l@{}} \forall w. \\ \quad \operatorname{is\text{-}complex\text{-}inner\text{-}product\text{-}space}(w) \Rightarrow  \\ \quad \quad \operatorname{cips\text{-}rip}(w) \in ((\operatorname{vec}(w) \times \operatorname{vec}(w)) \to \mathbb{R}) \end{array}}}
\vnbprop{cips-j-fun}{\vnbstmt{\begin{array}{@{}l@{}} \forall w. \\ \quad \operatorname{is\text{-}complex\text{-}inner\text{-}product\text{-}space}(w) \Rightarrow (\operatorname{cips\text{-}j}(w) \in (\operatorname{vec}(w) \to \operatorname{vec}(w))) \end{array}}}
\vnbprop{cips-rsip-scal}{\vnbstmt{\begin{array}{@{}l@{}} \forall w. \\ \quad \operatorname{is\text{-}complex\text{-}inner\text{-}product\text{-}space}(w) \Rightarrow  \\ \quad \quad (\operatorname{scal}(\operatorname{cips\text{-}rsip}(w)) \simeq \operatorname{normed\text{-}field\text{-}as\text{-}commutative\text{-}ring}(\operatorname{rr\text{-}normed\text{-}field})) \end{array}}}
\vnbprop{cips-real-scal}{\vnbstmt{\begin{array}{@{}l@{}} \forall w. \\ \quad \operatorname{is\text{-}complex\text{-}inner\text{-}product\text{-}space}(w) \Rightarrow  \\ \quad \quad (\operatorname{scal}(\operatorname{cips\text{-}real}(w)) \simeq \operatorname{normed\text{-}field\text{-}as\text{-}commutative\text{-}ring}(\operatorname{rr\text{-}normed\text{-}field})) \end{array}}}
\vnbprop{cips-rsip-vec}{\vnbstmt{\begin{array}{@{}l@{}} \forall w. \\ \quad \operatorname{is\text{-}complex\text{-}inner\text{-}product\text{-}space}(w) \Rightarrow (\operatorname{vec}(\operatorname{cips\text{-}rsip}(w)) \simeq \operatorname{vec}(w)) \end{array}}}
\vnbprop{cips-real-vec}{\vnbstmt{\begin{array}{@{}l@{}} \forall w. \\ \quad \operatorname{is\text{-}complex\text{-}inner\text{-}product\text{-}space}(w) \Rightarrow (\operatorname{vec}(\operatorname{cips\text{-}real}(w)) \simeq \operatorname{vec}(w)) \end{array}}}
\vnbprop{cips-rsip-vadd}{\vnbstmt{\begin{array}{@{}l@{}} \forall w. \\ \quad \operatorname{is\text{-}complex\text{-}inner\text{-}product\text{-}space}(w) \Rightarrow (\operatorname{vadd}(\operatorname{cips\text{-}rsip}(w)) \simeq \operatorname{vadd}(w)) \end{array}}}
\vnbprop{cips-real-vadd}{\vnbstmt{\begin{array}{@{}l@{}} \forall w. \\ \quad \operatorname{is\text{-}complex\text{-}inner\text{-}product\text{-}space}(w) \Rightarrow (\operatorname{vadd}(\operatorname{cips\text{-}real}(w)) \simeq \operatorname{vadd}(w)) \end{array}}}
\vnbprop{cips-rsip-vzero}{\vnbstmt{\begin{array}{@{}l@{}} \forall w. \\ \quad \operatorname{is\text{-}complex\text{-}inner\text{-}product\text{-}space}(w) \Rightarrow (\operatorname{vzero}(\operatorname{cips\text{-}rsip}(w)) \simeq \operatorname{vzero}(w)) \end{array}}}
\vnbprop{cips-real-vzero}{\vnbstmt{\begin{array}{@{}l@{}} \forall w. \\ \quad \operatorname{is\text{-}complex\text{-}inner\text{-}product\text{-}space}(w) \Rightarrow (\operatorname{vzero}(\operatorname{cips\text{-}real}(w)) \simeq \operatorname{vzero}(w)) \end{array}}}
\vnbprop{cips-rsip-vneg}{\vnbstmt{\begin{array}{@{}l@{}} \forall w. \\ \quad \operatorname{is\text{-}complex\text{-}inner\text{-}product\text{-}space}(w) \Rightarrow (\operatorname{vneg}(\operatorname{cips\text{-}rsip}(w)) \simeq \operatorname{vneg}(w)) \end{array}}}
\vnbprop{cips-real-vneg}{\vnbstmt{\begin{array}{@{}l@{}} \forall w. \\ \quad \operatorname{is\text{-}complex\text{-}inner\text{-}product\text{-}space}(w) \Rightarrow (\operatorname{vneg}(\operatorname{cips\text{-}real}(w)) \simeq \operatorname{vneg}(w)) \end{array}}}
\vnbprop{cips-rsip-act}{\vnbstmt{\begin{array}{@{}l@{}} \forall w. \\ \quad \operatorname{is\text{-}complex\text{-}inner\text{-}product\text{-}space}(w) \Rightarrow (\operatorname{act}(\operatorname{cips\text{-}rsip}(w)) \simeq \operatorname{cips\text{-}ract}(w)) \end{array}}}
\vnbprop{cips-real-act}{\vnbstmt{\begin{array}{@{}l@{}} \forall w. \\ \quad \operatorname{is\text{-}complex\text{-}inner\text{-}product\text{-}space}(w) \Rightarrow (\operatorname{act}(\operatorname{cips\text{-}real}(w)) \simeq \operatorname{cips\text{-}ract}(w)) \end{array}}}
\vnbprop{cips-rsip-rip}{\vnbstmt{\begin{array}{@{}l@{}} \forall w. \\ \quad \operatorname{is\text{-}complex\text{-}inner\text{-}product\text{-}space}(w) \Rightarrow (\langle\cdot\mid\cdot\rangle_{\operatorname{cips\text{-}rsip}(w)} \simeq \operatorname{cips\text{-}rip}(w)) \end{array}}}
\vnbprop{cips-real-rip}{\vnbstmt{\begin{array}{@{}l@{}} \forall w. \\ \quad \operatorname{is\text{-}complex\text{-}inner\text{-}product\text{-}space}(w) \Rightarrow (\langle\cdot\mid\cdot\rangle_{\operatorname{cips\text{-}real}(w)} \simeq \operatorname{cips\text{-}rip}(w)) \end{array}}}
\vnbprop{cips-real-jop}{\vnbstmt{\begin{array}{@{}l@{}} \forall w. \\ \quad \operatorname{is\text{-}complex\text{-}inner\text{-}product\text{-}space}(w) \Rightarrow (J_{\operatorname{cips\text{-}real}(w)} \simeq \operatorname{cips\text{-}j}(w)) \end{array}}}
\vnbprop{cips-rsip-length}{\vnbstmt{\forall w.\; (\operatorname{is\text{-}complex\text{-}inner\text{-}product\text{-}space}(w) \Rightarrow (\operatorname{length}(\operatorname{cips\text{-}rsip}(w)) = 7))}}
\vnbprop{cips-rsip-scal-ring}{\vnbstmt{\begin{array}{@{}l@{}} \forall w. \\ \quad \operatorname{is\text{-}complex\text{-}inner\text{-}product\text{-}space}(w) \Rightarrow \operatorname{is\text{-}ring}(\operatorname{scal}(\operatorname{cips\text{-}rsip}(w))) \end{array}}}
\vnbprop{cips-rsip-vec-set}{\vnbstmt{\forall w.\; (\operatorname{is\text{-}complex\text{-}inner\text{-}product\text{-}space}(w) \Rightarrow (\operatorname{vec}(\operatorname{cips\text{-}rsip}(w)) \in \mathbf{Set}))}}
\vnbprop{cips-rsip-vadd-fun}{\vnbstmt{\begin{array}{@{}l@{}} \forall w. \\ \quad \operatorname{is\text{-}complex\text{-}inner\text{-}product\text{-}space}(w) \Rightarrow  \\ \quad \quad \operatorname{vadd}(\operatorname{cips\text{-}rsip}(w)) \in ((\operatorname{vec}(\operatorname{cips\text{-}rsip}(w)) \times \operatorname{vec}(\operatorname{cips\text{-}rsip}(w))) \to \operatorname{vec}(\operatorname{cips\text{-}rsip}(w))) \end{array}}}
\vnbprop{cips-rsip-vzero-in}{\vnbstmt{\begin{array}{@{}l@{}} \forall w. \\ \quad \operatorname{is\text{-}complex\text{-}inner\text{-}product\text{-}space}(w) \Rightarrow  \\ \quad \quad \operatorname{vzero}(\operatorname{cips\text{-}rsip}(w)) \in \operatorname{vec}(\operatorname{cips\text{-}rsip}(w)) \end{array}}}
\vnbprop{cips-rsip-vneg-fun}{\vnbstmt{\begin{array}{@{}l@{}} \forall w. \\ \quad \operatorname{is\text{-}complex\text{-}inner\text{-}product\text{-}space}(w) \Rightarrow  \\ \quad \quad \operatorname{vneg}(\operatorname{cips\text{-}rsip}(w)) \in (\operatorname{vec}(\operatorname{cips\text{-}rsip}(w)) \to \operatorname{vec}(\operatorname{cips\text{-}rsip}(w))) \end{array}}}
\vnbprop{cips-rsip-act-fun}{\vnbstmt{\begin{array}{@{}l@{}} \forall w. \\ \quad \operatorname{is\text{-}complex\text{-}inner\text{-}product\text{-}space}(w) \Rightarrow  \\ \quad \quad \operatorname{act}(\operatorname{cips\text{-}rsip}(w)) \in ((\operatorname{carr}(\operatorname{scal}(\operatorname{cips\text{-}rsip}(w))) \times \operatorname{vec}(\operatorname{cips\text{-}rsip}(w))) \to \operatorname{vec}(\operatorname{cips\text{-}rsip}(w))) \end{array}}}
\vnbprop{cips-rsip-rip-fun}{\vnbstmt{\begin{array}{@{}l@{}} \forall w. \\ \quad \operatorname{is\text{-}complex\text{-}inner\text{-}product\text{-}space}(w) \Rightarrow  \\ \quad \quad \langle\cdot\mid\cdot\rangle_{\operatorname{cips\text{-}rsip}(w)} \in ((\operatorname{vec}(\operatorname{cips\text{-}rsip}(w)) \times \operatorname{vec}(\operatorname{cips\text{-}rsip}(w))) \to \mathbb{R}) \end{array}}}
\vnbprop{cips-rsip-assoc}{\vnbstmt{\begin{array}{@{}l@{}} \forall w. \\ \quad \operatorname{is\text{-}complex\text{-}inner\text{-}product\text{-}space}(w) \Rightarrow  \\ \quad \quad \operatorname{is\text{-}associative}(\operatorname{vadd}(\operatorname{cips\text{-}rsip}(w)), \operatorname{vec}(\operatorname{cips\text{-}rsip}(w))) \end{array}}}
\vnbprop{cips-rsip-comm}{\vnbstmt{\begin{array}{@{}l@{}} \forall w. \\ \quad \operatorname{is\text{-}complex\text{-}inner\text{-}product\text{-}space}(w) \Rightarrow  \\ \quad \quad \operatorname{is\text{-}commutative}(\operatorname{vadd}(\operatorname{cips\text{-}rsip}(w)), \operatorname{vec}(\operatorname{cips\text{-}rsip}(w))) \end{array}}}
\vnbprop{cips-rsip-identity}{\vnbstmt{\begin{array}{@{}l@{}} \forall w. \\ \quad \operatorname{is\text{-}complex\text{-}inner\text{-}product\text{-}space}(w) \Rightarrow  \\ \quad \quad \operatorname{is\text{-}identity}(\operatorname{vadd}(\operatorname{cips\text{-}rsip}(w)), \operatorname{vzero}(\operatorname{cips\text{-}rsip}(w)), \operatorname{vec}(\operatorname{cips\text{-}rsip}(w))) \end{array}}}
\vnbprop{cips-rsip-inverses}{\vnbstmt{\begin{array}{@{}l@{}} \forall w. \\ \quad \operatorname{is\text{-}complex\text{-}inner\text{-}product\text{-}space}(w) \Rightarrow  \\ \quad \quad \operatorname{has\text{-}inverses}(\operatorname{vadd}(\operatorname{cips\text{-}rsip}(w)), \operatorname{vzero}(\operatorname{cips\text{-}rsip}(w)), \operatorname{vneg}(\operatorname{cips\text{-}rsip}(w)), \operatorname{vec}(\operatorname{cips\text{-}rsip}(w))) \end{array}}}
\vnbprop{cips-rsip-scal-pin}{\vnbstmt{\begin{array}{@{}l@{}} \forall w. \\ \quad \operatorname{is\text{-}complex\text{-}inner\text{-}product\text{-}space}(w) \Rightarrow  \\ \quad \quad \operatorname{scal}(\operatorname{cips\text{-}rsip}(w)) = \operatorname{normed\text{-}field\text{-}as\text{-}commutative\text{-}ring}(\operatorname{rr\text{-}normed\text{-}field}) \end{array}}}
\vnbprop{cips-rsip-scal-carr}{\vnbstmt{\begin{array}{@{}l@{}} \forall w. \\ \quad \operatorname{is\text{-}complex\text{-}inner\text{-}product\text{-}space}(w) \Rightarrow (\operatorname{carr}(\operatorname{scal}(\operatorname{cips\text{-}rsip}(w))) = \mathbb{R}) \end{array}}}
\vnbprop{cips-rsip-act-vadd}{\vnbstmt{\begin{array}{@{}l@{}} \forall w.; \forall r \in \operatorname{carr}(\operatorname{scal}(\operatorname{cips\text{-}rsip}(w))), x, y \in \operatorname{vec}(\operatorname{cips\text{-}rsip}(w)). \\ \quad \operatorname{is\text{-}complex\text{-}inner\text{-}product\text{-}space}(w) \Rightarrow  \\ \quad \quad \operatorname{act}(\operatorname{cips\text{-}rsip}(w))(r, \operatorname{vadd}(\operatorname{cips\text{-}rsip}(w))(x, y)) \\ \quad \quad =\; \operatorname{vadd}(\operatorname{cips\text{-}rsip}(w))(\operatorname{act}(\operatorname{cips\text{-}rsip}(w))(r, x), \operatorname{act}(\operatorname{cips\text{-}rsip}(w))(r, y)) \end{array}}}
\vnbprop{cips-rsip-act-add}{\vnbstmt{\begin{array}{@{}l@{}} \forall w.; \forall r, s \in \operatorname{carr}(\operatorname{scal}(\operatorname{cips\text{-}rsip}(w))), x \in \operatorname{vec}(\operatorname{cips\text{-}rsip}(w)). \\ \quad \operatorname{is\text{-}complex\text{-}inner\text{-}product\text{-}space}(w) \Rightarrow  \\ \quad \quad \operatorname{act}(\operatorname{cips\text{-}rsip}(w))(\operatorname{add}(\operatorname{scal}(\operatorname{cips\text{-}rsip}(w)))(r, s), x) \\ \quad \quad =\; \operatorname{vadd}(\operatorname{cips\text{-}rsip}(w))(\operatorname{act}(\operatorname{cips\text{-}rsip}(w))(r, x), \operatorname{act}(\operatorname{cips\text{-}rsip}(w))(s, x)) \end{array}}}
\vnbprop{cips-rsip-act-mul}{\vnbstmt{\begin{array}{@{}l@{}} \forall w.; \forall r, s \in \operatorname{carr}(\operatorname{scal}(\operatorname{cips\text{-}rsip}(w))), x \in \operatorname{vec}(\operatorname{cips\text{-}rsip}(w)). \\ \quad \operatorname{is\text{-}complex\text{-}inner\text{-}product\text{-}space}(w) \Rightarrow  \\ \quad \quad \operatorname{act}(\operatorname{cips\text{-}rsip}(w))(\operatorname{mul}(\operatorname{scal}(\operatorname{cips\text{-}rsip}(w)))(r, s), x) \\ \quad \quad =\; \operatorname{act}(\operatorname{cips\text{-}rsip}(w))(r, \operatorname{act}(\operatorname{cips\text{-}rsip}(w))(s, x)) \end{array}}}
\vnbprop{cips-rsip-act-one}{\vnbstmt{\begin{array}{@{}l@{}} \forall w.; \forall x \in \operatorname{vec}(\operatorname{cips\text{-}rsip}(w)). \\ \quad \operatorname{is\text{-}complex\text{-}inner\text{-}product\text{-}space}(w) \Rightarrow  \\ \quad \quad \operatorname{act}(\operatorname{cips\text{-}rsip}(w))(\operatorname{one}(\operatorname{scal}(\operatorname{cips\text{-}rsip}(w))), x) = x \end{array}}}
\vnbprop{cips-rsip-rip-sym}{\vnbstmt{\begin{array}{@{}l@{}} \forall w.; \forall x, y \in \operatorname{vec}(\operatorname{cips\text{-}rsip}(w)). \\ \quad \operatorname{is\text{-}complex\text{-}inner\text{-}product\text{-}space}(w) \Rightarrow  \\ \quad \quad \langle\cdot\mid\cdot\rangle_{\operatorname{cips\text{-}rsip}(w)}(x, y) = \langle\cdot\mid\cdot\rangle_{\operatorname{cips\text{-}rsip}(w)}(y, x) \end{array}}}
\vnbprop{cips-rsip-rip-add-left}{\vnbstmt{\begin{array}{@{}l@{}} \forall w.; \forall x, y, z \in \operatorname{vec}(\operatorname{cips\text{-}rsip}(w)). \\ \quad \operatorname{is\text{-}complex\text{-}inner\text{-}product\text{-}space}(w) \Rightarrow  \\ \quad \quad \langle\cdot\mid\cdot\rangle_{\operatorname{cips\text{-}rsip}(w)}(\operatorname{vadd}(\operatorname{cips\text{-}rsip}(w))(x, y), z) \\ \quad \quad =\; \langle\cdot\mid\cdot\rangle_{\operatorname{cips\text{-}rsip}(w)}(x, z) + \langle\cdot\mid\cdot\rangle_{\operatorname{cips\text{-}rsip}(w)}(y, z) \end{array}}}
\vnbprop{cips-rsip-rip-homog-left}{\vnbstmt{\begin{array}{@{}l@{}} \forall w.; \forall a \in \mathbb{R}, x, y \in \operatorname{vec}(\operatorname{cips\text{-}rsip}(w)). \\ \quad \operatorname{is\text{-}complex\text{-}inner\text{-}product\text{-}space}(w) \Rightarrow  \\ \quad \quad \langle\cdot\mid\cdot\rangle_{\operatorname{cips\text{-}rsip}(w)}(\operatorname{act}(\operatorname{cips\text{-}rsip}(w))(a, x), y) = (a \cdot \langle\cdot\mid\cdot\rangle_{\operatorname{cips\text{-}rsip}(w)}(x, y)) \end{array}}}
\vnbprop{cips-rsip-rip-self-nonneg}{\vnbstmt{\begin{array}{@{}l@{}} \forall w.; \forall x \in \operatorname{vec}(\operatorname{cips\text{-}rsip}(w)). \\ \quad \operatorname{is\text{-}complex\text{-}inner\text{-}product\text{-}space}(w) \Rightarrow (0 \leq \langle\cdot\mid\cdot\rangle_{\operatorname{cips\text{-}rsip}(w)}(x, x)) \end{array}}}
\vnbprop{cips-rsip-is-sip}{\vnbstmt{\begin{array}{@{}l@{}} \forall w. \\ \quad \operatorname{is\text{-}complex\text{-}inner\text{-}product\text{-}space}(w) \Rightarrow  \\ \quad \quad \operatorname{is\text{-}semi\text{-}inner\text{-}product\text{-}space}(\operatorname{cips\text{-}rsip}(w)) \end{array}}}
\vnbprop{cips-rsip-unfold}{\vnbstmt{\begin{array}{@{}l@{}} \forall w. \\ \quad (\operatorname{cips\text{-}rsip}(w) \simeq {[}\operatorname{normed\text{-}field\text{-}as\text{-}commutative\text{-}ring}(\operatorname{rr\text{-}normed\text{-}field}), \operatorname{vec}(w), \operatorname{vadd}(w), \operatorname{vzero}(w), \operatorname{vneg}(w), \operatorname{cips\text{-}ract}(w), \operatorname{cips\text{-}rip}(w)]) \end{array}}}
\vnbprop{cips-rip-unfold}{\vnbstmt{\begin{array}{@{}l@{}} \forall w. \\ \quad (\operatorname{cips\text{-}rip}(w) \simeq (\lambda {[}\operatorname{cxx\_}, \operatorname{cxy\_}] \in (\operatorname{vec}(w) \times \operatorname{vec}(w)).\; \operatorname{real\text{-}part}(\langle\cdot,\cdot\rangle_{w}(\operatorname{cxx\_}, \operatorname{cxy\_})))) \end{array}}}
\vnbprop{cips-real-length}{\vnbstmt{\forall w.\; (\operatorname{is\text{-}complex\text{-}inner\text{-}product\text{-}space}(w) \Rightarrow (\operatorname{length}(\operatorname{cips\text{-}real}(w)) = 8))}}
\vnbprop{cips-real-scal-ring}{\vnbstmt{\begin{array}{@{}l@{}} \forall w. \\ \quad \operatorname{is\text{-}complex\text{-}inner\text{-}product\text{-}space}(w) \Rightarrow \operatorname{is\text{-}ring}(\operatorname{scal}(\operatorname{cips\text{-}real}(w))) \end{array}}}
\vnbprop{cips-real-vec-set}{\vnbstmt{\forall w.\; (\operatorname{is\text{-}complex\text{-}inner\text{-}product\text{-}space}(w) \Rightarrow (\operatorname{vec}(\operatorname{cips\text{-}real}(w)) \in \mathbf{Set}))}}
\vnbprop{cips-real-vadd-fun}{\vnbstmt{\begin{array}{@{}l@{}} \forall w. \\ \quad \operatorname{is\text{-}complex\text{-}inner\text{-}product\text{-}space}(w) \Rightarrow  \\ \quad \quad \operatorname{vadd}(\operatorname{cips\text{-}real}(w)) \in ((\operatorname{vec}(\operatorname{cips\text{-}real}(w)) \times \operatorname{vec}(\operatorname{cips\text{-}real}(w))) \to \operatorname{vec}(\operatorname{cips\text{-}real}(w))) \end{array}}}
\vnbprop{cips-real-vzero-in}{\vnbstmt{\begin{array}{@{}l@{}} \forall w. \\ \quad \operatorname{is\text{-}complex\text{-}inner\text{-}product\text{-}space}(w) \Rightarrow  \\ \quad \quad \operatorname{vzero}(\operatorname{cips\text{-}real}(w)) \in \operatorname{vec}(\operatorname{cips\text{-}real}(w)) \end{array}}}
\vnbprop{cips-real-vneg-fun}{\vnbstmt{\begin{array}{@{}l@{}} \forall w. \\ \quad \operatorname{is\text{-}complex\text{-}inner\text{-}product\text{-}space}(w) \Rightarrow  \\ \quad \quad \operatorname{vneg}(\operatorname{cips\text{-}real}(w)) \in (\operatorname{vec}(\operatorname{cips\text{-}real}(w)) \to \operatorname{vec}(\operatorname{cips\text{-}real}(w))) \end{array}}}
\vnbprop{cips-real-act-fun}{\vnbstmt{\begin{array}{@{}l@{}} \forall w. \\ \quad \operatorname{is\text{-}complex\text{-}inner\text{-}product\text{-}space}(w) \Rightarrow  \\ \quad \quad \operatorname{act}(\operatorname{cips\text{-}real}(w)) \in ((\operatorname{carr}(\operatorname{scal}(\operatorname{cips\text{-}real}(w))) \times \operatorname{vec}(\operatorname{cips\text{-}real}(w))) \to \operatorname{vec}(\operatorname{cips\text{-}real}(w))) \end{array}}}
\vnbprop{cips-real-rip-fun}{\vnbstmt{\begin{array}{@{}l@{}} \forall w. \\ \quad \operatorname{is\text{-}complex\text{-}inner\text{-}product\text{-}space}(w) \Rightarrow  \\ \quad \quad \langle\cdot\mid\cdot\rangle_{\operatorname{cips\text{-}real}(w)} \in ((\operatorname{vec}(\operatorname{cips\text{-}real}(w)) \times \operatorname{vec}(\operatorname{cips\text{-}real}(w))) \to \mathbb{R}) \end{array}}}
\vnbprop{cips-real-jop-fun}{\vnbstmt{\begin{array}{@{}l@{}} \forall w. \\ \quad \operatorname{is\text{-}complex\text{-}inner\text{-}product\text{-}space}(w) \Rightarrow  \\ \quad \quad J_{\operatorname{cips\text{-}real}(w)} \in (\operatorname{vec}(\operatorname{cips\text{-}real}(w)) \to \operatorname{vec}(\operatorname{cips\text{-}real}(w))) \end{array}}}
\vnbprop{cips-real-sip}{\vnbstmt{\begin{array}{@{}l@{}} \forall w. \\ \quad \operatorname{is\text{-}complex\text{-}inner\text{-}product\text{-}space}(w) \Rightarrow  \\ \quad \quad \operatorname{is\text{-}semi\text{-}inner\text{-}product\text{-}space}({[}\operatorname{scal}(\operatorname{cips\text{-}real}(w)), \operatorname{vec}(\operatorname{cips\text{-}real}(w)), \operatorname{vadd}(\operatorname{cips\text{-}real}(w)), \operatorname{vzero}(\operatorname{cips\text{-}real}(w)), \operatorname{vneg}(\operatorname{cips\text{-}real}(w)), \operatorname{act}(\operatorname{cips\text{-}real}(w)), \langle\cdot\mid\cdot\rangle_{\operatorname{cips\text{-}real}(w)}]) \end{array}}}
\vnbprop{cips-real-linear}{\vnbstmt{\begin{array}{@{}l@{}} \forall w. \\ \quad \operatorname{is\text{-}complex\text{-}inner\text{-}product\text{-}space}(w) \Rightarrow  \\ \quad \quad \operatorname{is\text{-}real\text{-}linear\text{-}op}(\operatorname{cips\text{-}real}(w), J_{\operatorname{cips\text{-}real}(w)}) \end{array}}}
\vnbprop{cips-real-square}{\vnbstmt{\begin{array}{@{}l@{}} \forall w. \\ \quad \operatorname{is\text{-}complex\text{-}inner\text{-}product\text{-}space}(w) \Rightarrow  \\ \quad \quad \operatorname{is\text{-}square\text{-}minus\text{-}id}(\operatorname{cips\text{-}real}(w), J_{\operatorname{cips\text{-}real}(w)}) \end{array}}}
\vnbprop{cips-real-antisym}{\vnbstmt{\begin{array}{@{}l@{}} \forall w. \\ \quad \operatorname{is\text{-}complex\text{-}inner\text{-}product\text{-}space}(w) \Rightarrow  \\ \quad \quad \operatorname{is\text{-}antisymmetric\text{-}op}(\operatorname{cips\text{-}real}(w), J_{\operatorname{cips\text{-}real}(w)}) \end{array}}}
\vnbprop{cips-real-definite}{\vnbstmt{\begin{array}{@{}l@{}} \forall w.; \forall x \in \operatorname{vec}(\operatorname{cips\text{-}real}(w)). \\ \quad (\operatorname{is\text{-}complex\text{-}inner\text{-}product\text{-}space}(w) \wedge (\langle\cdot\mid\cdot\rangle_{\operatorname{cips\text{-}real}(w)}(x, x) = 0)) \Rightarrow  \\ \quad \quad x = \operatorname{vzero}(\operatorname{cips\text{-}real}(w)) \end{array}}}
\vnbprop{cips-real-is-cphs}{\vnbstmt{\begin{array}{@{}l@{}} \forall w. \\ \quad \operatorname{is\text{-}complex\text{-}inner\text{-}product\text{-}space}(w) \Rightarrow  \\ \quad \quad \operatorname{is\text{-}complex\text{-}pre\text{-}hilbert\text{-}space}(\operatorname{cips\text{-}real}(w)) \end{array}}}
\vnbprop{csip-jop-isometric}{\vnbstmt{\begin{array}{@{}l@{}} \forall v.; \forall x, y \in \operatorname{vec}(v). \\ \quad \operatorname{is\text{-}complex\text{-}semi\text{-}inner\text{-}product\text{-}space}(v) \Rightarrow  \\ \quad \quad \langle\cdot\mid\cdot\rangle_{v}(J_{v}(x), J_{v}(y)) = \langle\cdot\mid\cdot\rangle_{v}(x, y) \end{array}}}
\vnbprop{csip-jop-sub}{\vnbstmt{\begin{array}{@{}l@{}} \forall v.; \forall x, y \in \operatorname{vec}(v). \\ \quad \operatorname{is\text{-}complex\text{-}semi\text{-}inner\text{-}product\text{-}space}(v) \Rightarrow  \\ \quad \quad J_{v}(\operatorname{vadd}(v)(x, \operatorname{vneg}(v)(y))) = \operatorname{vadd}(v)(J_{v}(x), \operatorname{vneg}(v)(J_{v}(y))) \end{array}}}
\vnbprop{csip-qj-lam-fun}{\vnbstmt{\begin{array}{@{}l@{}} \forall v. \\ \quad \operatorname{is\text{-}complex\text{-}semi\text{-}inner\text{-}product\text{-}space}(v) \Rightarrow  \\ \quad \quad (\lambda a \in \operatorname{vec}(v).\; \operatorname{class}(\operatorname{sip\text{-}setoid}(\operatorname{csip\text{-}as\text{-}sip}(v)), J_{v}(a))) \\ \quad \quad \in\; (\operatorname{pts}(\operatorname{sip\text{-}setoid}(\operatorname{csip\text{-}as\text{-}sip}(v))) \to \operatorname{quotient}(\operatorname{sip\text{-}setoid}(\operatorname{csip\text{-}as\text{-}sip}(v)))) \end{array}}}
\vnbprop{csip-qj-respects}{\vnbstmt{\begin{array}{@{}l@{}} \forall v. \\ \quad \operatorname{is\text{-}complex\text{-}semi\text{-}inner\text{-}product\text{-}space}(v) \Rightarrow  \\ \quad \quad \operatorname{respects}(\operatorname{sip\text{-}setoid}(\operatorname{csip\text{-}as\text{-}sip}(v)), \\ \quad \quad \quad (\lambda a \in \operatorname{vec}(v).\; \operatorname{class}(\operatorname{sip\text{-}setoid}(\operatorname{csip\text{-}as\text{-}sip}(v)), J_{v}(a)))) \end{array}}}
\vnbprop{csip-qj-computes}{\vnbstmt{\begin{array}{@{}l@{}} \forall v.; \forall a \in \operatorname{vec}(v). \\ \quad \operatorname{is\text{-}complex\text{-}semi\text{-}inner\text{-}product\text{-}space}(v) \Rightarrow  \\ \quad \quad \operatorname{csip\text{-}qj}(v)(\operatorname{class}(\operatorname{sip\text{-}setoid}(\operatorname{csip\text{-}as\text{-}sip}(v)), a)) = \operatorname{class}(\operatorname{sip\text{-}setoid}(\operatorname{csip\text{-}as\text{-}sip}(v)), J_{v}(a)) \end{array}}}
\vnbprop{csip-qj-fun}{\vnbstmt{\begin{array}{@{}l@{}} \forall v. \\ \quad \operatorname{is\text{-}complex\text{-}semi\text{-}inner\text{-}product\text{-}space}(v) \Rightarrow  \\ \quad \quad \operatorname{csip\text{-}qj}(v) \in (\operatorname{quotient}(\operatorname{sip\text{-}setoid}(\operatorname{csip\text{-}as\text{-}sip}(v))) \to \operatorname{quotient}(\operatorname{sip\text{-}setoid}(\operatorname{csip\text{-}as\text{-}sip}(v)))) \end{array}}}
\vnbprop{csip-quotient-scal}{\vnbstmt{\begin{array}{@{}l@{}} \forall v. \\ \quad \operatorname{is\text{-}complex\text{-}semi\text{-}inner\text{-}product\text{-}space}(v) \Rightarrow (\operatorname{scal}(\operatorname{csip\text{-}quotient}(v)) \simeq \operatorname{scal}(v)) \end{array}}}
\vnbprop{csip-quotient-vec}{\vnbstmt{\begin{array}{@{}l@{}} \forall v. \\ \quad \operatorname{is\text{-}complex\text{-}semi\text{-}inner\text{-}product\text{-}space}(v) \Rightarrow  \\ \quad \quad (\operatorname{vec}(\operatorname{csip\text{-}quotient}(v)) \simeq \operatorname{quotient}(\operatorname{sip\text{-}setoid}(\operatorname{csip\text{-}as\text{-}sip}(v)))) \end{array}}}
\vnbprop{csip-quotient-vadd}{\vnbstmt{\begin{array}{@{}l@{}} \forall v. \\ \quad \operatorname{is\text{-}complex\text{-}semi\text{-}inner\text{-}product\text{-}space}(v) \Rightarrow  \\ \quad \quad (\operatorname{vadd}(\operatorname{csip\text{-}quotient}(v)) \simeq \operatorname{sip\text{-}qvadd}(\operatorname{csip\text{-}as\text{-}sip}(v))) \end{array}}}
\vnbprop{csip-quotient-vzero}{\vnbstmt{\begin{array}{@{}l@{}} \forall v. \\ \quad \operatorname{is\text{-}complex\text{-}semi\text{-}inner\text{-}product\text{-}space}(v) \Rightarrow  \\ \quad \quad (\operatorname{vzero}(\operatorname{csip\text{-}quotient}(v)) \simeq \operatorname{class}(\operatorname{sip\text{-}setoid}(\operatorname{csip\text{-}as\text{-}sip}(v)), \operatorname{vzero}(v))) \end{array}}}
\vnbprop{csip-quotient-vneg}{\vnbstmt{\begin{array}{@{}l@{}} \forall v. \\ \quad \operatorname{is\text{-}complex\text{-}semi\text{-}inner\text{-}product\text{-}space}(v) \Rightarrow  \\ \quad \quad (\operatorname{vneg}(\operatorname{csip\text{-}quotient}(v)) \simeq \operatorname{sip\text{-}qvneg}(\operatorname{csip\text{-}as\text{-}sip}(v))) \end{array}}}
\vnbprop{csip-quotient-act}{\vnbstmt{\begin{array}{@{}l@{}} \forall v. \\ \quad \operatorname{is\text{-}complex\text{-}semi\text{-}inner\text{-}product\text{-}space}(v) \Rightarrow  \\ \quad \quad (\operatorname{act}(\operatorname{csip\text{-}quotient}(v)) \simeq \operatorname{sip\text{-}qact}(\operatorname{csip\text{-}as\text{-}sip}(v))) \end{array}}}
\vnbprop{csip-quotient-rip}{\vnbstmt{\begin{array}{@{}l@{}} \forall v. \\ \quad \operatorname{is\text{-}complex\text{-}semi\text{-}inner\text{-}product\text{-}space}(v) \Rightarrow  \\ \quad \quad (\langle\cdot\mid\cdot\rangle_{\operatorname{csip\text{-}quotient}(v)} \simeq \operatorname{sip\text{-}qrip}(\operatorname{csip\text{-}as\text{-}sip}(v))) \end{array}}}
\vnbprop{csip-quotient-jop}{\vnbstmt{\begin{array}{@{}l@{}} \forall v. \\ \quad \operatorname{is\text{-}complex\text{-}semi\text{-}inner\text{-}product\text{-}space}(v) \Rightarrow  \\ \quad \quad (J_{\operatorname{csip\text{-}quotient}(v)} \simeq \operatorname{csip\text{-}qj}(v)) \end{array}}}
\vnbprop{csip-quotient-vadd-computes}{\vnbstmt{\begin{array}{@{}l@{}} \forall v, a, b. \\ \quad (\operatorname{is\text{-}complex\text{-}semi\text{-}inner\text{-}product\text{-}space}(v) \wedge ((a \in \operatorname{vec}(v)) \wedge (b \in \operatorname{vec}(v)))) \Rightarrow  \\ \quad \quad \operatorname{vadd}(\operatorname{csip\text{-}quotient}(v))(\operatorname{class}(\operatorname{sip\text{-}setoid}(\operatorname{csip\text{-}as\text{-}sip}(v)), a), \operatorname{class}(\operatorname{sip\text{-}setoid}(\operatorname{csip\text{-}as\text{-}sip}(v)), b)) \\ \quad \quad =\; \operatorname{class}(\operatorname{sip\text{-}setoid}(\operatorname{csip\text{-}as\text{-}sip}(v)), \operatorname{vadd}(v)(a, b)) \end{array}}}
\vnbprop{csip-quotient-vneg-computes}{\vnbstmt{\begin{array}{@{}l@{}} \forall v.; \forall a \in \operatorname{vec}(v). \\ \quad \operatorname{is\text{-}complex\text{-}semi\text{-}inner\text{-}product\text{-}space}(v) \Rightarrow  \\ \quad \quad \operatorname{vneg}(\operatorname{csip\text{-}quotient}(v))(\operatorname{class}(\operatorname{sip\text{-}setoid}(\operatorname{csip\text{-}as\text{-}sip}(v)), a)) \\ \quad \quad =\; \operatorname{class}(\operatorname{sip\text{-}setoid}(\operatorname{csip\text{-}as\text{-}sip}(v)), \operatorname{vneg}(v)(a)) \end{array}}}
\vnbprop{csip-quotient-act-computes}{\vnbstmt{\begin{array}{@{}l@{}} \forall v, r, a. \\ \quad (\operatorname{is\text{-}complex\text{-}semi\text{-}inner\text{-}product\text{-}space}(v) \wedge ((r \in \mathbb{R}) \wedge (a \in \operatorname{vec}(v)))) \Rightarrow  \\ \quad \quad \operatorname{act}(\operatorname{csip\text{-}quotient}(v))(r, \operatorname{class}(\operatorname{sip\text{-}setoid}(\operatorname{csip\text{-}as\text{-}sip}(v)), a)) \\ \quad \quad =\; \operatorname{class}(\operatorname{sip\text{-}setoid}(\operatorname{csip\text{-}as\text{-}sip}(v)), \operatorname{act}(v)(r, a)) \end{array}}}
\vnbprop{csip-quotient-rip-computes}{\vnbstmt{\begin{array}{@{}l@{}} \forall v, a, b. \\ \quad (\operatorname{is\text{-}complex\text{-}semi\text{-}inner\text{-}product\text{-}space}(v) \wedge ((a \in \operatorname{vec}(v)) \wedge (b \in \operatorname{vec}(v)))) \Rightarrow  \\ \quad \quad \langle\cdot\mid\cdot\rangle_{\operatorname{csip\text{-}quotient}(v)}(\operatorname{class}(\operatorname{sip\text{-}setoid}(\operatorname{csip\text{-}as\text{-}sip}(v)), a), \operatorname{class}(\operatorname{sip\text{-}setoid}(\operatorname{csip\text{-}as\text{-}sip}(v)), b)) \\ \quad \quad =\; \langle\cdot\mid\cdot\rangle_{v}(a, b) \end{array}}}
\vnbprop{csip-quotient-jop-computes}{\vnbstmt{\begin{array}{@{}l@{}} \forall v.; \forall a \in \operatorname{vec}(v). \\ \quad \operatorname{is\text{-}complex\text{-}semi\text{-}inner\text{-}product\text{-}space}(v) \Rightarrow  \\ \quad \quad J_{\operatorname{csip\text{-}quotient}(v)}(\operatorname{class}(\operatorname{sip\text{-}setoid}(\operatorname{csip\text{-}as\text{-}sip}(v)), a)) \\ \quad \quad =\; \operatorname{class}(\operatorname{sip\text{-}setoid}(\operatorname{csip\text{-}as\text{-}sip}(v)), J_{v}(a)) \end{array}}}
\vnbprop{csip-quotient-length}{\vnbstmt{\begin{array}{@{}l@{}} \forall v. \\ \quad \operatorname{is\text{-}complex\text{-}semi\text{-}inner\text{-}product\text{-}space}(v) \Rightarrow (\operatorname{length}(\operatorname{csip\text{-}quotient}(v)) = 8) \end{array}}}
\vnbprop{csip-quotient-scal-ring}{\vnbstmt{\begin{array}{@{}l@{}} \forall v. \\ \quad \operatorname{is\text{-}complex\text{-}semi\text{-}inner\text{-}product\text{-}space}(v) \Rightarrow \operatorname{is\text{-}ring}(\operatorname{scal}(\operatorname{csip\text{-}quotient}(v))) \end{array}}}
\vnbprop{csip-quotient-vec-set}{\vnbstmt{\begin{array}{@{}l@{}} \forall v. \\ \quad \operatorname{is\text{-}complex\text{-}semi\text{-}inner\text{-}product\text{-}space}(v) \Rightarrow (\operatorname{vec}(\operatorname{csip\text{-}quotient}(v)) \in \mathbf{Set}) \end{array}}}
\vnbprop{csip-quotient-vadd-fun}{\vnbstmt{\begin{array}{@{}l@{}} \forall v. \\ \quad \operatorname{is\text{-}complex\text{-}semi\text{-}inner\text{-}product\text{-}space}(v) \Rightarrow  \\ \quad \quad \operatorname{vadd}(\operatorname{csip\text{-}quotient}(v)) \\ \quad \quad \in\; ((\operatorname{vec}(\operatorname{csip\text{-}quotient}(v)) \times \operatorname{vec}(\operatorname{csip\text{-}quotient}(v))) \to \operatorname{vec}(\operatorname{csip\text{-}quotient}(v))) \end{array}}}
\vnbprop{csip-quotient-vzero-in}{\vnbstmt{\begin{array}{@{}l@{}} \forall v. \\ \quad \operatorname{is\text{-}complex\text{-}semi\text{-}inner\text{-}product\text{-}space}(v) \Rightarrow  \\ \quad \quad \operatorname{vzero}(\operatorname{csip\text{-}quotient}(v)) \in \operatorname{vec}(\operatorname{csip\text{-}quotient}(v)) \end{array}}}
\vnbprop{csip-quotient-vneg-fun}{\vnbstmt{\begin{array}{@{}l@{}} \forall v. \\ \quad \operatorname{is\text{-}complex\text{-}semi\text{-}inner\text{-}product\text{-}space}(v) \Rightarrow  \\ \quad \quad \operatorname{vneg}(\operatorname{csip\text{-}quotient}(v)) \in (\operatorname{vec}(\operatorname{csip\text{-}quotient}(v)) \to \operatorname{vec}(\operatorname{csip\text{-}quotient}(v))) \end{array}}}
\vnbprop{csip-quotient-act-fun}{\vnbstmt{\begin{array}{@{}l@{}} \forall v. \\ \quad \operatorname{is\text{-}complex\text{-}semi\text{-}inner\text{-}product\text{-}space}(v) \Rightarrow  \\ \quad \quad \operatorname{act}(\operatorname{csip\text{-}quotient}(v)) \\ \quad \quad \in\; ((\operatorname{carr}(\operatorname{scal}(\operatorname{csip\text{-}quotient}(v))) \times \operatorname{vec}(\operatorname{csip\text{-}quotient}(v))) \to \operatorname{vec}(\operatorname{csip\text{-}quotient}(v))) \end{array}}}
\vnbprop{csip-quotient-rip-fun}{\vnbstmt{\begin{array}{@{}l@{}} \forall v. \\ \quad \operatorname{is\text{-}complex\text{-}semi\text{-}inner\text{-}product\text{-}space}(v) \Rightarrow  \\ \quad \quad \langle\cdot\mid\cdot\rangle_{\operatorname{csip\text{-}quotient}(v)} \in ((\operatorname{vec}(\operatorname{csip\text{-}quotient}(v)) \times \operatorname{vec}(\operatorname{csip\text{-}quotient}(v))) \to \mathbb{R}) \end{array}}}
\vnbprop{csip-quotient-jop-fun}{\vnbstmt{\begin{array}{@{}l@{}} \forall v. \\ \quad \operatorname{is\text{-}complex\text{-}semi\text{-}inner\text{-}product\text{-}space}(v) \Rightarrow  \\ \quad \quad J_{\operatorname{csip\text{-}quotient}(v)} \in (\operatorname{vec}(\operatorname{csip\text{-}quotient}(v)) \to \operatorname{vec}(\operatorname{csip\text{-}quotient}(v))) \end{array}}}
\vnbprop{csip-quotient-sip}{\vnbstmt{\begin{array}{@{}l@{}} \forall v. \\ \quad \operatorname{is\text{-}complex\text{-}semi\text{-}inner\text{-}product\text{-}space}(v) \Rightarrow  \\ \quad \quad \operatorname{is\text{-}semi\text{-}inner\text{-}product\text{-}space}({[}\operatorname{scal}(\operatorname{csip\text{-}quotient}(v)), \operatorname{vec}(\operatorname{csip\text{-}quotient}(v)), \operatorname{vadd}(\operatorname{csip\text{-}quotient}(v)), \operatorname{vzero}(\operatorname{csip\text{-}quotient}(v)), \operatorname{vneg}(\operatorname{csip\text{-}quotient}(v)), \operatorname{act}(\operatorname{csip\text{-}quotient}(v)), \langle\cdot\mid\cdot\rangle_{\operatorname{csip\text{-}quotient}(v)}]) \end{array}}}
\vnbprop{csip-quotient-linear}{\vnbstmt{\begin{array}{@{}l@{}} \forall v. \\ \quad \operatorname{is\text{-}complex\text{-}semi\text{-}inner\text{-}product\text{-}space}(v) \Rightarrow  \\ \quad \quad \operatorname{is\text{-}real\text{-}linear\text{-}op}(\operatorname{csip\text{-}quotient}(v), J_{\operatorname{csip\text{-}quotient}(v)}) \end{array}}}
\vnbprop{csip-quotient-square}{\vnbstmt{\begin{array}{@{}l@{}} \forall v. \\ \quad \operatorname{is\text{-}complex\text{-}semi\text{-}inner\text{-}product\text{-}space}(v) \Rightarrow  \\ \quad \quad \operatorname{is\text{-}square\text{-}minus\text{-}id}(\operatorname{csip\text{-}quotient}(v), J_{\operatorname{csip\text{-}quotient}(v)}) \end{array}}}
\vnbprop{csip-quotient-antisym}{\vnbstmt{\begin{array}{@{}l@{}} \forall v. \\ \quad \operatorname{is\text{-}complex\text{-}semi\text{-}inner\text{-}product\text{-}space}(v) \Rightarrow  \\ \quad \quad \operatorname{is\text{-}antisymmetric\text{-}op}(\operatorname{csip\text{-}quotient}(v), J_{\operatorname{csip\text{-}quotient}(v)}) \end{array}}}
\vnbprop{csip-quotient-definite}{\vnbstmt{\begin{array}{@{}l@{}} \forall v.; \forall x \in \operatorname{vec}(\operatorname{csip\text{-}quotient}(v)). \\ \quad (\operatorname{is\text{-}complex\text{-}semi\text{-}inner\text{-}product\text{-}space}(v) \wedge (\langle\cdot\mid\cdot\rangle_{\operatorname{csip\text{-}quotient}(v)}(x, x) = 0)) \Rightarrow  \\ \quad \quad x = \operatorname{vzero}(\operatorname{csip\text{-}quotient}(v)) \end{array}}}
\vnbprop{csip-quotient-is-cphs}{\vnbstmt{\begin{array}{@{}l@{}} \forall v. \\ \quad \operatorname{is\text{-}complex\text{-}semi\text{-}inner\text{-}product\text{-}space}(v) \Rightarrow  \\ \quad \quad \operatorname{is\text{-}complex\text{-}pre\text{-}hilbert\text{-}space}(\operatorname{csip\text{-}quotient}(v)) \end{array}}}
\vnbprop{csip-cact-neg-one}{\vnbstmt{\begin{array}{@{}l@{}} \forall v.; \forall x \in \operatorname{vec}(v). \\ \quad \operatorname{is\text{-}complex\text{-}semi\text{-}inner\text{-}product\text{-}space}(v) \Rightarrow  \\ \quad \quad \operatorname{cact}(v)(-1, x) = \operatorname{vneg}(v)(x) \end{array}}}
\vnbprop{csip-cact-neg-i}{\vnbstmt{\begin{array}{@{}l@{}} \forall v.; \forall x \in \operatorname{vec}(v). \\ \quad \operatorname{is\text{-}complex\text{-}semi\text{-}inner\text{-}product\text{-}space}(v) \Rightarrow  \\ \quad \quad \operatorname{cact}(v)(-+i, x) = \operatorname{vneg}(v)(\operatorname{cact}(v)(+i, x)) \end{array}}}
\vnbprop{ses-type}{\vnbstmt{\begin{array}{@{}l@{}} \forall v, f.; \forall x, y \in \operatorname{vec}(v). \\ \quad (\operatorname{is\text{-}complex\text{-}semi\text{-}inner\text{-}product\text{-}space}(v) \wedge \operatorname{is\text{-}sesquilinear\text{-}form}(v, f)) \Rightarrow  \\ \quad \quad f(x, y) \in \mathbb{C} \end{array}}}
\vnbprop{ses-add-left}{\vnbstmt{\begin{array}{@{}l@{}} \forall v, f.; \forall x, y, z \in \operatorname{vec}(v). \\ \quad (\operatorname{is\text{-}complex\text{-}semi\text{-}inner\text{-}product\text{-}space}(v) \wedge \operatorname{is\text{-}sesquilinear\text{-}form}(v, f)) \Rightarrow  \\ \quad \quad f(\operatorname{vadd}(v)(x, y), z) = (f(x, z) + f(y, z)) \end{array}}}
\vnbprop{ses-homog-left}{\vnbstmt{\begin{array}{@{}l@{}} \forall v, f.; \forall a \in \mathbb{C}, x, y \in \operatorname{vec}(v). \\ \quad (\operatorname{is\text{-}complex\text{-}semi\text{-}inner\text{-}product\text{-}space}(v) \wedge \operatorname{is\text{-}sesquilinear\text{-}form}(v, f)) \Rightarrow  \\ \quad \quad f(\operatorname{cact}(v)(a, x), y) = (a \cdot f(x, y)) \end{array}}}
\vnbprop{ses-conj-sym}{\vnbstmt{\begin{array}{@{}l@{}} \forall v, f.; \forall x, y \in \operatorname{vec}(v). \\ \quad (\operatorname{is\text{-}complex\text{-}semi\text{-}inner\text{-}product\text{-}space}(v) \wedge \operatorname{is\text{-}sesquilinear\text{-}form}(v, f)) \Rightarrow  \\ \quad \quad f(y, x) = \operatorname{conjugate}(f(x, y)) \end{array}}}
\vnbprop{ses-add-right}{\vnbstmt{\begin{array}{@{}l@{}} \forall v, f.; \forall x, y, z \in \operatorname{vec}(v). \\ \quad (\operatorname{is\text{-}complex\text{-}semi\text{-}inner\text{-}product\text{-}space}(v) \wedge \operatorname{is\text{-}sesquilinear\text{-}form}(v, f)) \Rightarrow  \\ \quad \quad f(x, \operatorname{vadd}(v)(y, z)) = (f(x, y) + f(x, z)) \end{array}}}
\vnbprop{ses-homog-right}{\vnbstmt{\begin{array}{@{}l@{}} \forall v, f.; \forall k \in \mathbb{C}, x, y \in \operatorname{vec}(v). \\ \quad (\operatorname{is\text{-}complex\text{-}semi\text{-}inner\text{-}product\text{-}space}(v) \wedge \operatorname{is\text{-}sesquilinear\text{-}form}(v, f)) \Rightarrow  \\ \quad \quad f(x, \operatorname{cact}(v)(k, y)) = (\operatorname{conjugate}(k) \cdot f(x, y)) \end{array}}}
\vnbprop{ses-expand}{\vnbstmt{\begin{array}{@{}l@{}} \forall v, f.; \forall c \in \mathbb{C}, x, y \in \operatorname{vec}(v). \\ \quad (\operatorname{is\text{-}complex\text{-}semi\text{-}inner\text{-}product\text{-}space}(v) \wedge \operatorname{is\text{-}sesquilinear\text{-}form}(v, f)) \Rightarrow  \\ \quad \quad f(\operatorname{vadd}(v)(x, \operatorname{cact}(v)(c, y)), \operatorname{vadd}(v)(x, \operatorname{cact}(v)(c, y))) \\ \quad \quad =\; ((f(x, x) + (\operatorname{conjugate}(c) \cdot f(x, y))) + (c \cdot f(y, x))) + ((c \cdot \operatorname{conjugate}(c)) \cdot f(y, y)) \end{array}}}
\vnbprop{ses-polarization}{\vnbstmt{\begin{array}{@{}l@{}} \forall v, f.; \forall x, y \in \operatorname{vec}(v). \\ \quad (\operatorname{is\text{-}complex\text{-}semi\text{-}inner\text{-}product\text{-}space}(v) \wedge \operatorname{is\text{-}sesquilinear\text{-}form}(v, f)) \Rightarrow  \\ \quad \quad 4 \cdot f(x, y) \\ \quad \quad =\; (((f(\operatorname{vadd}(v)(x, y), \operatorname{vadd}(v)(x, y)) - f(\operatorname{vadd}(v)(x, \operatorname{vneg}(v)(y)), \operatorname{vadd}(v)(x, \operatorname{vneg}(v)(y)))) + (+i \cdot f(\operatorname{vadd}(v)(x, \operatorname{cact}(v)(+i, y)), \operatorname{vadd}(v)(x, \operatorname{cact}(v)(+i, y))))) - (+i \cdot f(\operatorname{vadd}(v)(x, \operatorname{vneg}(v)(\operatorname{cact}(v)(+i, y))), \operatorname{vadd}(v)(x, \operatorname{vneg}(v)(\operatorname{cact}(v)(+i, y)))))) \end{array}}}
\vnbprop{csip-crip-sesquilinear}{\vnbstmt{\begin{array}{@{}l@{}} \forall v. \\ \quad \operatorname{is\text{-}complex\text{-}semi\text{-}inner\text{-}product\text{-}space}(v) \Rightarrow \operatorname{is\text{-}sesquilinear\text{-}form}(v, \operatorname{crip}(v)) \end{array}}}
\vnbprop{cc-four-cancel}{\vnbstmt{\forall a, b \in \mathbb{C}.\; (((4 \cdot a) = (4 \cdot b)) \Rightarrow (a = b))}}
\vnbprop{csip-crip-unique}{\vnbstmt{\begin{array}{@{}l@{}} \forall v, f.; \forall x, y \in \operatorname{vec}(v). \\ \quad \operatorname{is\text{-}complex\text{-}semi\text{-}inner\text{-}product\text{-}space}(v) \wedge \\ \quad \operatorname{is\text{-}sesquilinear\text{-}form}(v, f) \wedge \\ \quad \forall \mathit{ux} \in \operatorname{vec}(v).\; (f(\mathit{ux}, \mathit{ux}) = \langle\cdot\mid\cdot\rangle_{v}(\mathit{ux}, \mathit{ux})) \Rightarrow \\ \quad \quad f(x, y) = \operatorname{crip}(v)(x, y) \end{array}}}
\vnbprop{cx-vadd-value}{\vnbstmt{\begin{array}{@{}l@{}} \forall v, x, \mathit{qx}. \\ \quad \operatorname{is\text{-}semi\text{-}inner\text{-}product\text{-}space}(v) \wedge \\ \quad x \in (\operatorname{vec}(v) \times \operatorname{vec}(v)) \wedge \\ \quad \mathit{qx} \in (\operatorname{vec}(v) \times \operatorname{vec}(v)) \Rightarrow \\ \quad \quad \operatorname{cx\text{-}vadd}(v)(x, \mathit{qx}) = {[}\operatorname{vadd}(v)(\operatorname{nth}(1, x), \operatorname{nth}(1, \mathit{qx})), \operatorname{vadd}(v)(\operatorname{nth}(2, x), \operatorname{nth}(2, \mathit{qx}))] \end{array}}}
\vnbprop{cx-vneg-value}{\vnbstmt{\begin{array}{@{}l@{}} \forall v.; \forall x \in (\operatorname{vec}(v) \times \operatorname{vec}(v)). \\ \quad \operatorname{is\text{-}semi\text{-}inner\text{-}product\text{-}space}(v) \Rightarrow  \\ \quad \quad \operatorname{cx\text{-}vneg}(v)(x) = {[}\operatorname{vneg}(v)(\operatorname{nth}(1, x)), \operatorname{vneg}(v)(\operatorname{nth}(2, x))] \end{array}}}
\vnbprop{cx-act-value}{\vnbstmt{\begin{array}{@{}l@{}} \forall v, x, \mathit{px}. \\ \quad (\operatorname{is\text{-}semi\text{-}inner\text{-}product\text{-}space}(v) \wedge ((x \in \mathbb{R}) \wedge (\mathit{px} \in (\operatorname{vec}(v) \times \operatorname{vec}(v))))) \Rightarrow  \\ \quad \quad \operatorname{cx\text{-}act}(v)(x, \mathit{px}) = {[}\operatorname{act}(v)(x, \operatorname{nth}(1, \mathit{px})), \operatorname{act}(v)(x, \operatorname{nth}(2, \mathit{px}))] \end{array}}}
\vnbprop{cx-rip-value}{\vnbstmt{\begin{array}{@{}l@{}} \forall v, x, \mathit{qx}. \\ \quad \operatorname{is\text{-}semi\text{-}inner\text{-}product\text{-}space}(v) \wedge \\ \quad x \in (\operatorname{vec}(v) \times \operatorname{vec}(v)) \wedge \\ \quad \mathit{qx} \in (\operatorname{vec}(v) \times \operatorname{vec}(v)) \Rightarrow \\ \quad \quad \operatorname{cx\text{-}rip}(v)(x, \mathit{qx}) = (\langle\cdot\mid\cdot\rangle_{v}(\operatorname{nth}(1, x), \operatorname{nth}(1, \mathit{qx})) + \langle\cdot\mid\cdot\rangle_{v}(\operatorname{nth}(2, x), \operatorname{nth}(2, \mathit{qx}))) \end{array}}}
\vnbprop{cx-j-value}{\vnbstmt{\begin{array}{@{}l@{}} \forall v.; \forall x \in (\operatorname{vec}(v) \times \operatorname{vec}(v)). \\ \quad \operatorname{is\text{-}semi\text{-}inner\text{-}product\text{-}space}(v) \Rightarrow  \\ \quad \quad \operatorname{cx\text{-}j}(v)(x) = {[}\operatorname{vneg}(v)(\operatorname{nth}(2, x)), \operatorname{nth}(1, x)] \end{array}}}
\vnbprop{cx-sip-scal}{\vnbstmt{\forall v.\; (\operatorname{is\text{-}semi\text{-}inner\text{-}product\text{-}space}(v) \Rightarrow (\operatorname{scal}(\operatorname{cx\text{-}sip}(v)) \simeq \operatorname{scal}(v)))}}
\vnbprop{complexify-scal}{\vnbstmt{\begin{array}{@{}l@{}} \forall v. \\ \quad \operatorname{is\text{-}semi\text{-}inner\text{-}product\text{-}space}(v) \Rightarrow (\operatorname{scal}(\operatorname{complexify}(v)) \simeq \operatorname{scal}(v)) \end{array}}}
\vnbprop{cx-sip-vec}{\vnbstmt{\begin{array}{@{}l@{}} \forall v. \\ \quad \operatorname{is\text{-}semi\text{-}inner\text{-}product\text{-}space}(v) \Rightarrow  \\ \quad \quad (\operatorname{vec}(\operatorname{cx\text{-}sip}(v)) \simeq (\operatorname{vec}(v) \times \operatorname{vec}(v))) \end{array}}}
\vnbprop{complexify-vec}{\vnbstmt{\begin{array}{@{}l@{}} \forall v. \\ \quad \operatorname{is\text{-}semi\text{-}inner\text{-}product\text{-}space}(v) \Rightarrow  \\ \quad \quad (\operatorname{vec}(\operatorname{complexify}(v)) \simeq (\operatorname{vec}(v) \times \operatorname{vec}(v))) \end{array}}}
\vnbprop{cx-sip-vadd}{\vnbstmt{\begin{array}{@{}l@{}} \forall v. \\ \quad \operatorname{is\text{-}semi\text{-}inner\text{-}product\text{-}space}(v) \Rightarrow (\operatorname{vadd}(\operatorname{cx\text{-}sip}(v)) \simeq \operatorname{cx\text{-}vadd}(v)) \end{array}}}
\vnbprop{complexify-vadd}{\vnbstmt{\begin{array}{@{}l@{}} \forall v. \\ \quad \operatorname{is\text{-}semi\text{-}inner\text{-}product\text{-}space}(v) \Rightarrow (\operatorname{vadd}(\operatorname{complexify}(v)) \simeq \operatorname{cx\text{-}vadd}(v)) \end{array}}}
\vnbprop{cx-sip-vzero}{\vnbstmt{\begin{array}{@{}l@{}} \forall v. \\ \quad \operatorname{is\text{-}semi\text{-}inner\text{-}product\text{-}space}(v) \Rightarrow (\operatorname{vzero}(\operatorname{cx\text{-}sip}(v)) \simeq {[}\operatorname{vzero}(v), \operatorname{vzero}(v)]) \end{array}}}
\vnbprop{complexify-vzero}{\vnbstmt{\begin{array}{@{}l@{}} \forall v. \\ \quad \operatorname{is\text{-}semi\text{-}inner\text{-}product\text{-}space}(v) \Rightarrow  \\ \quad \quad (\operatorname{vzero}(\operatorname{complexify}(v)) \simeq {[}\operatorname{vzero}(v), \operatorname{vzero}(v)]) \end{array}}}
\vnbprop{cx-sip-vneg}{\vnbstmt{\begin{array}{@{}l@{}} \forall v. \\ \quad \operatorname{is\text{-}semi\text{-}inner\text{-}product\text{-}space}(v) \Rightarrow (\operatorname{vneg}(\operatorname{cx\text{-}sip}(v)) \simeq \operatorname{cx\text{-}vneg}(v)) \end{array}}}
\vnbprop{complexify-vneg}{\vnbstmt{\begin{array}{@{}l@{}} \forall v. \\ \quad \operatorname{is\text{-}semi\text{-}inner\text{-}product\text{-}space}(v) \Rightarrow (\operatorname{vneg}(\operatorname{complexify}(v)) \simeq \operatorname{cx\text{-}vneg}(v)) \end{array}}}
\vnbprop{cx-sip-act}{\vnbstmt{\forall v.\; (\operatorname{is\text{-}semi\text{-}inner\text{-}product\text{-}space}(v) \Rightarrow (\operatorname{act}(\operatorname{cx\text{-}sip}(v)) \simeq \operatorname{cx\text{-}act}(v)))}}
\vnbprop{complexify-act}{\vnbstmt{\begin{array}{@{}l@{}} \forall v. \\ \quad \operatorname{is\text{-}semi\text{-}inner\text{-}product\text{-}space}(v) \Rightarrow (\operatorname{act}(\operatorname{complexify}(v)) \simeq \operatorname{cx\text{-}act}(v)) \end{array}}}
\vnbprop{cx-sip-rip}{\vnbstmt{\forall v.\; (\operatorname{is\text{-}semi\text{-}inner\text{-}product\text{-}space}(v) \Rightarrow (\langle\cdot\mid\cdot\rangle_{\operatorname{cx\text{-}sip}(v)} \simeq \operatorname{cx\text{-}rip}(v)))}}
\vnbprop{complexify-rip}{\vnbstmt{\begin{array}{@{}l@{}} \forall v. \\ \quad \operatorname{is\text{-}semi\text{-}inner\text{-}product\text{-}space}(v) \Rightarrow (\langle\cdot\mid\cdot\rangle_{\operatorname{complexify}(v)} \simeq \operatorname{cx\text{-}rip}(v)) \end{array}}}
\vnbprop{complexify-jop}{\vnbstmt{\begin{array}{@{}l@{}} \forall v. \\ \quad \operatorname{is\text{-}semi\text{-}inner\text{-}product\text{-}space}(v) \Rightarrow (J_{\operatorname{complexify}(v)} \simeq \operatorname{cx\text{-}j}(v)) \end{array}}}
\vnbprop{cx-j-fun}{\vnbstmt{\begin{array}{@{}l@{}} \forall v. \\ \quad \operatorname{is\text{-}semi\text{-}inner\text{-}product\text{-}space}(v) \Rightarrow  \\ \quad \quad \operatorname{cx\text{-}j}(v) \in ((\operatorname{vec}(v) \times \operatorname{vec}(v)) \to (\operatorname{vec}(v) \times \operatorname{vec}(v))) \end{array}}}
\vnbprop{cx-vadd-fun}{\vnbstmt{\begin{array}{@{}l@{}} \forall v. \\ \quad \operatorname{is\text{-}semi\text{-}inner\text{-}product\text{-}space}(v) \Rightarrow  \\ \quad \quad \operatorname{cx\text{-}vadd}(v) \\ \quad \quad \in\; (((\operatorname{vec}(v) \times \operatorname{vec}(v)) \times (\operatorname{vec}(v) \times \operatorname{vec}(v))) \to (\operatorname{vec}(v) \times \operatorname{vec}(v))) \end{array}}}
\vnbprop{cx-vneg-fun}{\vnbstmt{\begin{array}{@{}l@{}} \forall v. \\ \quad \operatorname{is\text{-}semi\text{-}inner\text{-}product\text{-}space}(v) \Rightarrow  \\ \quad \quad \operatorname{cx\text{-}vneg}(v) \in ((\operatorname{vec}(v) \times \operatorname{vec}(v)) \to (\operatorname{vec}(v) \times \operatorname{vec}(v))) \end{array}}}
\vnbprop{cx-act-fun}{\vnbstmt{\begin{array}{@{}l@{}} \forall v. \\ \quad \operatorname{is\text{-}semi\text{-}inner\text{-}product\text{-}space}(v) \Rightarrow  \\ \quad \quad \operatorname{cx\text{-}act}(v) \in ((\mathbb{R} \times (\operatorname{vec}(v) \times \operatorname{vec}(v))) \to (\operatorname{vec}(v) \times \operatorname{vec}(v))) \end{array}}}
\vnbprop{cx-rip-fun}{\vnbstmt{\begin{array}{@{}l@{}} \forall v. \\ \quad \operatorname{is\text{-}semi\text{-}inner\text{-}product\text{-}space}(v) \Rightarrow  \\ \quad \quad \operatorname{cx\text{-}rip}(v) \in (((\operatorname{vec}(v) \times \operatorname{vec}(v)) \times (\operatorname{vec}(v) \times \operatorname{vec}(v))) \to \mathbb{R}) \end{array}}}
\vnbprop{cx-sip-length}{\vnbstmt{\forall v.\; (\operatorname{is\text{-}semi\text{-}inner\text{-}product\text{-}space}(v) \Rightarrow (\operatorname{length}(\operatorname{cx\text{-}sip}(v)) = 7))}}
\vnbprop{cx-sip-scal-ring}{\vnbstmt{\forall v.\; (\operatorname{is\text{-}semi\text{-}inner\text{-}product\text{-}space}(v) \Rightarrow \operatorname{is\text{-}ring}(\operatorname{scal}(\operatorname{cx\text{-}sip}(v))))}}
\vnbprop{cx-sip-vec-set}{\vnbstmt{\forall v.\; (\operatorname{is\text{-}semi\text{-}inner\text{-}product\text{-}space}(v) \Rightarrow (\operatorname{vec}(\operatorname{cx\text{-}sip}(v)) \in \mathbf{Set}))}}
\vnbprop{cx-sip-vadd-fun}{\vnbstmt{\begin{array}{@{}l@{}} \forall v. \\ \quad \operatorname{is\text{-}semi\text{-}inner\text{-}product\text{-}space}(v) \Rightarrow  \\ \quad \quad \operatorname{vadd}(\operatorname{cx\text{-}sip}(v)) \in ((\operatorname{vec}(\operatorname{cx\text{-}sip}(v)) \times \operatorname{vec}(\operatorname{cx\text{-}sip}(v))) \to \operatorname{vec}(\operatorname{cx\text{-}sip}(v))) \end{array}}}
\vnbprop{cx-sip-vzero-in}{\vnbstmt{\begin{array}{@{}l@{}} \forall v. \\ \quad \operatorname{is\text{-}semi\text{-}inner\text{-}product\text{-}space}(v) \Rightarrow (\operatorname{vzero}(\operatorname{cx\text{-}sip}(v)) \in \operatorname{vec}(\operatorname{cx\text{-}sip}(v))) \end{array}}}
\vnbprop{cx-sip-vneg-fun}{\vnbstmt{\begin{array}{@{}l@{}} \forall v. \\ \quad \operatorname{is\text{-}semi\text{-}inner\text{-}product\text{-}space}(v) \Rightarrow  \\ \quad \quad \operatorname{vneg}(\operatorname{cx\text{-}sip}(v)) \in (\operatorname{vec}(\operatorname{cx\text{-}sip}(v)) \to \operatorname{vec}(\operatorname{cx\text{-}sip}(v))) \end{array}}}
\vnbprop{cx-sip-act-fun}{\vnbstmt{\begin{array}{@{}l@{}} \forall v. \\ \quad \operatorname{is\text{-}semi\text{-}inner\text{-}product\text{-}space}(v) \Rightarrow  \\ \quad \quad \operatorname{act}(\operatorname{cx\text{-}sip}(v)) \in ((\operatorname{carr}(\operatorname{scal}(\operatorname{cx\text{-}sip}(v))) \times \operatorname{vec}(\operatorname{cx\text{-}sip}(v))) \to \operatorname{vec}(\operatorname{cx\text{-}sip}(v))) \end{array}}}
\vnbprop{cx-sip-rip-fun}{\vnbstmt{\begin{array}{@{}l@{}} \forall v. \\ \quad \operatorname{is\text{-}semi\text{-}inner\text{-}product\text{-}space}(v) \Rightarrow  \\ \quad \quad \langle\cdot\mid\cdot\rangle_{\operatorname{cx\text{-}sip}(v)} \in ((\operatorname{vec}(\operatorname{cx\text{-}sip}(v)) \times \operatorname{vec}(\operatorname{cx\text{-}sip}(v))) \to \mathbb{R}) \end{array}}}
\vnbprop{cx-sip-assoc}{\vnbstmt{\begin{array}{@{}l@{}} \forall v. \\ \quad \operatorname{is\text{-}semi\text{-}inner\text{-}product\text{-}space}(v) \Rightarrow  \\ \quad \quad \operatorname{is\text{-}associative}(\operatorname{vadd}(\operatorname{cx\text{-}sip}(v)), \operatorname{vec}(\operatorname{cx\text{-}sip}(v))) \end{array}}}
\vnbprop{cx-sip-comm}{\vnbstmt{\begin{array}{@{}l@{}} \forall v. \\ \quad \operatorname{is\text{-}semi\text{-}inner\text{-}product\text{-}space}(v) \Rightarrow  \\ \quad \quad \operatorname{is\text{-}commutative}(\operatorname{vadd}(\operatorname{cx\text{-}sip}(v)), \operatorname{vec}(\operatorname{cx\text{-}sip}(v))) \end{array}}}
\vnbprop{cx-sip-identity}{\vnbstmt{\begin{array}{@{}l@{}} \forall v. \\ \quad \operatorname{is\text{-}semi\text{-}inner\text{-}product\text{-}space}(v) \Rightarrow  \\ \quad \quad \operatorname{is\text{-}identity}(\operatorname{vadd}(\operatorname{cx\text{-}sip}(v)), \operatorname{vzero}(\operatorname{cx\text{-}sip}(v)), \operatorname{vec}(\operatorname{cx\text{-}sip}(v))) \end{array}}}
\vnbprop{cx-sip-inverses}{\vnbstmt{\begin{array}{@{}l@{}} \forall v. \\ \quad \operatorname{is\text{-}semi\text{-}inner\text{-}product\text{-}space}(v) \Rightarrow  \\ \quad \quad \operatorname{has\text{-}inverses}(\operatorname{vadd}(\operatorname{cx\text{-}sip}(v)), \operatorname{vzero}(\operatorname{cx\text{-}sip}(v)), \operatorname{vneg}(\operatorname{cx\text{-}sip}(v)), \operatorname{vec}(\operatorname{cx\text{-}sip}(v))) \end{array}}}
\vnbprop{cx-sip-scal-pin}{\vnbstmt{\begin{array}{@{}l@{}} \forall v. \\ \quad \operatorname{is\text{-}semi\text{-}inner\text{-}product\text{-}space}(v) \Rightarrow  \\ \quad \quad \operatorname{scal}(\operatorname{cx\text{-}sip}(v)) = \operatorname{normed\text{-}field\text{-}as\text{-}commutative\text{-}ring}(\operatorname{rr\text{-}normed\text{-}field}) \end{array}}}
\vnbprop{cx-sip-scal-carr}{\vnbstmt{\forall v.\; (\operatorname{is\text{-}semi\text{-}inner\text{-}product\text{-}space}(v) \Rightarrow (\operatorname{carr}(\operatorname{scal}(\operatorname{cx\text{-}sip}(v))) = \mathbb{R}))}}
\vnbprop{cx-sip-act-vadd}{\vnbstmt{\begin{array}{@{}l@{}} \forall v.; \forall r \in \operatorname{carr}(\operatorname{scal}(\operatorname{cx\text{-}sip}(v))), x, y \in \operatorname{vec}(\operatorname{cx\text{-}sip}(v)). \\ \quad \operatorname{is\text{-}semi\text{-}inner\text{-}product\text{-}space}(v) \Rightarrow  \\ \quad \quad \operatorname{act}(\operatorname{cx\text{-}sip}(v))(r, \operatorname{vadd}(\operatorname{cx\text{-}sip}(v))(x, y)) \\ \quad \quad =\; \operatorname{vadd}(\operatorname{cx\text{-}sip}(v))(\operatorname{act}(\operatorname{cx\text{-}sip}(v))(r, x), \operatorname{act}(\operatorname{cx\text{-}sip}(v))(r, y)) \end{array}}}
\vnbprop{cx-sip-act-add}{\vnbstmt{\begin{array}{@{}l@{}} \forall v.; \forall r, s \in \operatorname{carr}(\operatorname{scal}(\operatorname{cx\text{-}sip}(v))), x \in \operatorname{vec}(\operatorname{cx\text{-}sip}(v)). \\ \quad \operatorname{is\text{-}semi\text{-}inner\text{-}product\text{-}space}(v) \Rightarrow  \\ \quad \quad \operatorname{act}(\operatorname{cx\text{-}sip}(v))(\operatorname{add}(\operatorname{scal}(\operatorname{cx\text{-}sip}(v)))(r, s), x) \\ \quad \quad =\; \operatorname{vadd}(\operatorname{cx\text{-}sip}(v))(\operatorname{act}(\operatorname{cx\text{-}sip}(v))(r, x), \operatorname{act}(\operatorname{cx\text{-}sip}(v))(s, x)) \end{array}}}
\vnbprop{cx-sip-act-mul}{\vnbstmt{\begin{array}{@{}l@{}} \forall v.; \forall r, s \in \operatorname{carr}(\operatorname{scal}(\operatorname{cx\text{-}sip}(v))), x \in \operatorname{vec}(\operatorname{cx\text{-}sip}(v)). \\ \quad \operatorname{is\text{-}semi\text{-}inner\text{-}product\text{-}space}(v) \Rightarrow  \\ \quad \quad \operatorname{act}(\operatorname{cx\text{-}sip}(v))(\operatorname{mul}(\operatorname{scal}(\operatorname{cx\text{-}sip}(v)))(r, s), x) = \operatorname{act}(\operatorname{cx\text{-}sip}(v))(r, \operatorname{act}(\operatorname{cx\text{-}sip}(v))(s, x)) \end{array}}}
\vnbprop{cx-sip-act-one}{\vnbstmt{\begin{array}{@{}l@{}} \forall v.; \forall x \in \operatorname{vec}(\operatorname{cx\text{-}sip}(v)). \\ \quad \operatorname{is\text{-}semi\text{-}inner\text{-}product\text{-}space}(v) \Rightarrow  \\ \quad \quad \operatorname{act}(\operatorname{cx\text{-}sip}(v))(\operatorname{one}(\operatorname{scal}(\operatorname{cx\text{-}sip}(v))), x) = x \end{array}}}
\vnbprop{cx-sip-rip-sym}{\vnbstmt{\begin{array}{@{}l@{}} \forall v.; \forall x, y \in \operatorname{vec}(\operatorname{cx\text{-}sip}(v)). \\ \quad \operatorname{is\text{-}semi\text{-}inner\text{-}product\text{-}space}(v) \Rightarrow  \\ \quad \quad \langle\cdot\mid\cdot\rangle_{\operatorname{cx\text{-}sip}(v)}(x, y) = \langle\cdot\mid\cdot\rangle_{\operatorname{cx\text{-}sip}(v)}(y, x) \end{array}}}
\vnbprop{cx-sip-rip-add-left}{\vnbstmt{\begin{array}{@{}l@{}} \forall v.; \forall x, y, z \in \operatorname{vec}(\operatorname{cx\text{-}sip}(v)). \\ \quad \operatorname{is\text{-}semi\text{-}inner\text{-}product\text{-}space}(v) \Rightarrow  \\ \quad \quad \langle\cdot\mid\cdot\rangle_{\operatorname{cx\text{-}sip}(v)}(\operatorname{vadd}(\operatorname{cx\text{-}sip}(v))(x, y), z) = (\langle\cdot\mid\cdot\rangle_{\operatorname{cx\text{-}sip}(v)}(x, z) + \langle\cdot\mid\cdot\rangle_{\operatorname{cx\text{-}sip}(v)}(y, z)) \end{array}}}
\vnbprop{cx-sip-rip-homog-left}{\vnbstmt{\begin{array}{@{}l@{}} \forall v.; \forall a \in \mathbb{R}, x, y \in \operatorname{vec}(\operatorname{cx\text{-}sip}(v)). \\ \quad \operatorname{is\text{-}semi\text{-}inner\text{-}product\text{-}space}(v) \Rightarrow  \\ \quad \quad \langle\cdot\mid\cdot\rangle_{\operatorname{cx\text{-}sip}(v)}(\operatorname{act}(\operatorname{cx\text{-}sip}(v))(a, x), y) = (a \cdot \langle\cdot\mid\cdot\rangle_{\operatorname{cx\text{-}sip}(v)}(x, y)) \end{array}}}
\vnbprop{cx-sip-rip-self-nonneg}{\vnbstmt{\begin{array}{@{}l@{}} \forall v.; \forall x \in \operatorname{vec}(\operatorname{cx\text{-}sip}(v)). \\ \quad \operatorname{is\text{-}semi\text{-}inner\text{-}product\text{-}space}(v) \Rightarrow (0 \leq \langle\cdot\mid\cdot\rangle_{\operatorname{cx\text{-}sip}(v)}(x, x)) \end{array}}}
\vnbprop{cx-sip-is-sip}{\vnbstmt{\begin{array}{@{}l@{}} \forall v. \\ \quad \operatorname{is\text{-}semi\text{-}inner\text{-}product\text{-}space}(v) \Rightarrow \operatorname{is\text{-}semi\text{-}inner\text{-}product\text{-}space}(\operatorname{cx\text{-}sip}(v)) \end{array}}}
\vnbprop{cx-sip-unfold}{\vnbstmt{\begin{array}{@{}l@{}} \forall v. \\ \quad \operatorname{is\text{-}semi\text{-}inner\text{-}product\text{-}space}(v) \Rightarrow  \\ \quad \quad (\operatorname{cx\text{-}sip}(v) \simeq {[}\operatorname{scal}(v), (\operatorname{vec}(v) \times \operatorname{vec}(v)), \operatorname{cx\text{-}vadd}(v), {[}\operatorname{vzero}(v), \operatorname{vzero}(v)], \operatorname{cx\text{-}vneg}(v), \operatorname{cx\text{-}act}(v), \operatorname{cx\text{-}rip}(v)]) \end{array}}}
\vnbprop{complexify-length}{\vnbstmt{\forall v.\; (\operatorname{is\text{-}semi\text{-}inner\text{-}product\text{-}space}(v) \Rightarrow (\operatorname{length}(\operatorname{complexify}(v)) = 8))}}
\vnbprop{complexify-scal-ring}{\vnbstmt{\forall v.\; (\operatorname{is\text{-}semi\text{-}inner\text{-}product\text{-}space}(v) \Rightarrow \operatorname{is\text{-}ring}(\operatorname{scal}(\operatorname{complexify}(v))))}}
\vnbprop{complexify-vec-set}{\vnbstmt{\forall v.\; (\operatorname{is\text{-}semi\text{-}inner\text{-}product\text{-}space}(v) \Rightarrow (\operatorname{vec}(\operatorname{complexify}(v)) \in \mathbf{Set}))}}
\vnbprop{complexify-vadd-fun}{\vnbstmt{\begin{array}{@{}l@{}} \forall v. \\ \quad \operatorname{is\text{-}semi\text{-}inner\text{-}product\text{-}space}(v) \Rightarrow  \\ \quad \quad \operatorname{vadd}(\operatorname{complexify}(v)) \in ((\operatorname{vec}(\operatorname{complexify}(v)) \times \operatorname{vec}(\operatorname{complexify}(v))) \to \operatorname{vec}(\operatorname{complexify}(v))) \end{array}}}
\vnbprop{complexify-vzero-in}{\vnbstmt{\begin{array}{@{}l@{}} \forall v. \\ \quad \operatorname{is\text{-}semi\text{-}inner\text{-}product\text{-}space}(v) \Rightarrow  \\ \quad \quad \operatorname{vzero}(\operatorname{complexify}(v)) \in \operatorname{vec}(\operatorname{complexify}(v)) \end{array}}}
\vnbprop{complexify-vneg-fun}{\vnbstmt{\begin{array}{@{}l@{}} \forall v. \\ \quad \operatorname{is\text{-}semi\text{-}inner\text{-}product\text{-}space}(v) \Rightarrow  \\ \quad \quad \operatorname{vneg}(\operatorname{complexify}(v)) \in (\operatorname{vec}(\operatorname{complexify}(v)) \to \operatorname{vec}(\operatorname{complexify}(v))) \end{array}}}
\vnbprop{complexify-act-fun}{\vnbstmt{\begin{array}{@{}l@{}} \forall v. \\ \quad \operatorname{is\text{-}semi\text{-}inner\text{-}product\text{-}space}(v) \Rightarrow  \\ \quad \quad \operatorname{act}(\operatorname{complexify}(v)) \\ \quad \quad \in\; ((\operatorname{carr}(\operatorname{scal}(\operatorname{complexify}(v))) \times \operatorname{vec}(\operatorname{complexify}(v))) \to \operatorname{vec}(\operatorname{complexify}(v))) \end{array}}}
\vnbprop{complexify-rip-fun}{\vnbstmt{\begin{array}{@{}l@{}} \forall v. \\ \quad \operatorname{is\text{-}semi\text{-}inner\text{-}product\text{-}space}(v) \Rightarrow  \\ \quad \quad \langle\cdot\mid\cdot\rangle_{\operatorname{complexify}(v)} \in ((\operatorname{vec}(\operatorname{complexify}(v)) \times \operatorname{vec}(\operatorname{complexify}(v))) \to \mathbb{R}) \end{array}}}
\vnbprop{complexify-jop-fun}{\vnbstmt{\begin{array}{@{}l@{}} \forall v. \\ \quad \operatorname{is\text{-}semi\text{-}inner\text{-}product\text{-}space}(v) \Rightarrow  \\ \quad \quad J_{\operatorname{complexify}(v)} \in (\operatorname{vec}(\operatorname{complexify}(v)) \to \operatorname{vec}(\operatorname{complexify}(v))) \end{array}}}
\vnbprop{complexify-sip}{\vnbstmt{\begin{array}{@{}l@{}} \forall v. \\ \quad \operatorname{is\text{-}semi\text{-}inner\text{-}product\text{-}space}(v) \Rightarrow  \\ \quad \quad \operatorname{is\text{-}semi\text{-}inner\text{-}product\text{-}space}({[}\operatorname{scal}(\operatorname{complexify}(v)), \operatorname{vec}(\operatorname{complexify}(v)), \operatorname{vadd}(\operatorname{complexify}(v)), \operatorname{vzero}(\operatorname{complexify}(v)), \operatorname{vneg}(\operatorname{complexify}(v)), \operatorname{act}(\operatorname{complexify}(v)), \langle\cdot\mid\cdot\rangle_{\operatorname{complexify}(v)}]) \end{array}}}
\vnbprop{complexify-linear}{\vnbstmt{\begin{array}{@{}l@{}} \forall v. \\ \quad \operatorname{is\text{-}semi\text{-}inner\text{-}product\text{-}space}(v) \Rightarrow  \\ \quad \quad \operatorname{is\text{-}real\text{-}linear\text{-}op}(\operatorname{complexify}(v), J_{\operatorname{complexify}(v)}) \end{array}}}
\vnbprop{complexify-square}{\vnbstmt{\begin{array}{@{}l@{}} \forall v. \\ \quad \operatorname{is\text{-}semi\text{-}inner\text{-}product\text{-}space}(v) \Rightarrow  \\ \quad \quad \operatorname{is\text{-}square\text{-}minus\text{-}id}(\operatorname{complexify}(v), J_{\operatorname{complexify}(v)}) \end{array}}}
\vnbprop{complexify-antisym}{\vnbstmt{\begin{array}{@{}l@{}} \forall v. \\ \quad \operatorname{is\text{-}semi\text{-}inner\text{-}product\text{-}space}(v) \Rightarrow  \\ \quad \quad \operatorname{is\text{-}antisymmetric\text{-}op}(\operatorname{complexify}(v), J_{\operatorname{complexify}(v)}) \end{array}}}
\vnbprop{complexify-is-csip}{\vnbstmt{\begin{array}{@{}l@{}} \forall v. \\ \quad \operatorname{is\text{-}semi\text{-}inner\text{-}product\text{-}space}(v) \Rightarrow  \\ \quad \quad \operatorname{is\text{-}complex\text{-}semi\text{-}inner\text{-}product\text{-}space}(\operatorname{complexify}(v)) \end{array}}}

\vnbbookfile{theorem-library/hilbert-onb}{theorem-library/hilbert-onb}
\noindent{\small hilbert-onb.scm -- wave of 2026-11-09/10 (agent BA-ONB), wired 2026-11-10 from scratchpad/baonb-1109/hilbert-onb.scm. Library: analysis.}\par
\vnbprop{onb-unit-scale-identity}{\vnbstmt{\begin{array}{@{}l@{}} \forall a, b, t \in \mathbb{R}. \\ \quad (((a \cdot a) = t) \wedge ((a \cdot b) = 1)) \Rightarrow ((b \cdot (b \cdot t)) = 1) \end{array}}}
\vnbprop{onb-unit-multiple}{\vnbstmt{\begin{array}{@{}l@{}} \forall v.; \forall x \in \operatorname{vec}(v). \\ \quad (\operatorname{is\text{-}complex\text{-}inner\text{-}product\text{-}space}(v) \wedge \neg (x = \operatorname{vzero}(v))) \Rightarrow  \\ \quad \quad \exists c \in \mathbb{C}.\; (\langle\cdot,\cdot\rangle_{v}(\operatorname{act}(v)(c, x), \operatorname{act}(v)(c, x)) = 1) \end{array}}}
\vnbprop{onb-orthonormal-chain-union}{\vnbstmt{\begin{array}{@{}l@{}} \forall v, h. \\ \quad \operatorname{is\text{-}complex\text{-}inner\text{-}product\text{-}space}(v) \wedge \\ \quad \forall a \in h.\; \operatorname{is\text{-}orthonormal\text{-}set}(v, a) \wedge \\ \quad \forall a, b \in h.\; ((a \subseteq b) \vee (b \subseteq a)) \Rightarrow \\ \quad \quad \operatorname{is\text{-}orthonormal\text{-}set}(v, \operatorname{big\text{-}union}(w, h, w)) \end{array}}}
\vnbprop{onb-orthonormal-insert}{\vnbstmt{\begin{array}{@{}l@{}} \forall v, m.; \forall u \in \operatorname{vec}(v). \\ \quad \operatorname{is\text{-}complex\text{-}inner\text{-}product\text{-}space}(v) \wedge \\ \quad \operatorname{is\text{-}orthonormal\text{-}set}(v, m) \wedge \\ \quad \langle\cdot,\cdot\rangle_{v}(u, u) = 1 \wedge \\ \quad \forall a \in m.\; (\langle\cdot,\cdot\rangle_{v}(u, a) = 0) \Rightarrow \\ \quad \quad \operatorname{is\text{-}orthonormal\text{-}set}(v, (m \cup \{u\})) \end{array}}}
\vnbprop{onb-complete-orthonormal-set-exists}{\vnbstmt{\begin{array}{@{}l@{}} \forall v, a. \\ \quad (\operatorname{is\text{-}complex\text{-}inner\text{-}product\text{-}space}(v) \wedge \operatorname{is\text{-}orthonormal\text{-}set}(v, a)) \Rightarrow  \\ \quad \quad \exists m. \\ \quad \quad \quad a \subseteq m \wedge \\ \quad \quad \quad \operatorname{is\text{-}orthonormal\text{-}set}(v, m) \wedge \\ \quad \quad \quad \forall x \in \operatorname{vec}(v).\; (\forall u \in m.\; (\langle\cdot,\cdot\rangle_{v}(x, u) = 0) \Rightarrow (x = \operatorname{vzero}(v))) \end{array}}}
\vnbprop{onb-image-orthonormal}{\vnbstmt{\begin{array}{@{}l@{}} \forall v, x, e. \\ \quad (\operatorname{is\text{-}complex\text{-}inner\text{-}product\text{-}space}(v) \wedge \operatorname{is\text{-}orthonormal\text{-}family}(v, x, e)) \Rightarrow  \\ \quad \quad \operatorname{is\text{-}orthonormal\text{-}set}(v, \operatorname{image}(e, x)) \end{array}}}
\vnbprop{onb-of-complete-orthonormal-set}{\vnbstmt{\begin{array}{@{}l@{}} \forall v, m. \\ \quad \operatorname{is\text{-}complex\text{-}inner\text{-}product\text{-}space}(v) \wedge \\ \quad \operatorname{is\text{-}orthonormal\text{-}set}(v, m) \wedge \\ \quad \forall x \in \operatorname{vec}(v).\; (\forall u \in m.\; (\langle\cdot,\cdot\rangle_{v}(x, u) = 0) \Rightarrow (x = \operatorname{vzero}(v))) \Rightarrow \\ \quad \quad \operatorname{is\text{-}orthonormal\text{-}basis}(v, m, (\lambda k \in m.\; k)) \end{array}}}
\vnbprop{onb-exists}{\vnbstmt{\begin{array}{@{}l@{}} \forall v. \\ \quad \operatorname{is\text{-}complex\text{-}inner\text{-}product\text{-}space}(v) \Rightarrow  \\ \quad \quad \exists x, e.\; \operatorname{is\text{-}orthonormal\text{-}basis}(v, x, e) \end{array}}}
\vnbprop{onb-extension}{\vnbstmt{\begin{array}{@{}l@{}} \forall v, x, e. \\ \quad (\operatorname{is\text{-}complex\text{-}inner\text{-}product\text{-}space}(v) \wedge \operatorname{is\text{-}orthonormal\text{-}family}(v, x, e)) \Rightarrow  \\ \quad \quad \exists a, f. \\ \quad \quad \quad (\operatorname{is\text{-}orthonormal\text{-}basis}(v, a, f) \wedge \forall i \in x.\; \exists j \in a.\; (f(j) = e(i))) \end{array}}}
\vnbprop{onb-ip-normed-ag-iden}{\vnbstmt{\begin{array}{@{}l@{}} \forall v. \\ \quad \operatorname{is\text{-}complex\text{-}inner\text{-}product\text{-}space}(v) \Rightarrow (\operatorname{iden}(\operatorname{ip\text{-}normed\text{-}ag}(v)) \simeq \operatorname{vzero}(v)) \end{array}}}
\vnbprop{onb-partial-sums}{\vnbstmt{\begin{array}{@{}l@{}} \forall v, a, e.; \forall x \in \operatorname{vec}(v), w \in \mathbf{Set}. \\ \quad \operatorname{is\text{-}complex\text{-}inner\text{-}product\text{-}space}(v) \wedge \\ \quad \operatorname{is\text{-}orthonormal\text{-}family}(v, a, e) \wedge \\ \quad |w| \in \mathbb{N} \wedge \\ \quad w \subseteq a \Rightarrow \\ \quad \quad \operatorname{finsum}(\operatorname{normed\text{-}ag\text{-}as\text{-}abelian\text{-}group}(\operatorname{ip\text{-}normed\text{-}ag}(v)), \\ \quad \quad \quad (\lambda k \in a.\; \operatorname{act}(v)(\langle\cdot,\cdot\rangle_{v}(x, e(k)), e(k))), w) \\ \quad \quad \in\; \operatorname{vec}(v) \wedge \\ \quad \quad \forall q \in a. \\ \quad \quad \quad (q \in w) \Rightarrow  \\ \quad \quad \quad \quad \langle\cdot,\cdot\rangle_{v}(\operatorname{finsum}(\operatorname{normed\text{-}ag\text{-}as\text{-}abelian\text{-}group}(\operatorname{ip\text{-}normed\text{-}ag}(v)), \\ \quad \quad \quad \quad \quad \quad (\lambda k \in a.\; \operatorname{act}(v)(\langle\cdot,\cdot\rangle_{v}(x, e(k)), e(k))), w), \\ \quad \quad \quad \quad \quad e(q)) \\ \quad \quad \quad \quad =\; \langle\cdot,\cdot\rangle_{v}(x, e(q)) \wedge \\ \quad \quad \forall q \in a. \\ \quad \quad \quad \neg (q \in w) \Rightarrow  \\ \quad \quad \quad \quad \langle\cdot,\cdot\rangle_{v}(\operatorname{finsum}(\operatorname{normed\text{-}ag\text{-}as\text{-}abelian\text{-}group}(\operatorname{ip\text{-}normed\text{-}ag}(v)), \\ \quad \quad \quad \quad \quad \quad (\lambda k \in a.\; \operatorname{act}(v)(\langle\cdot,\cdot\rangle_{v}(x, e(k)), e(k))), w), \\ \quad \quad \quad \quad \quad e(q)) \\ \quad \quad \quad \quad =\; 0 \wedge \\ \quad \quad \operatorname{finsum}(\operatorname{rr\text{-}pos\text{-}star\text{-}add\text{-}monoid}, \\ \quad \quad \quad (\lambda k \in a.\; (\operatorname{magnitude}(\langle\cdot,\cdot\rangle_{v}(x, e(k))) \cdot \operatorname{magnitude}(\langle\cdot,\cdot\rangle_{v}(x, e(k))))), \\ \quad \quad \quad w) \\ \quad \quad \in\; \mathbb{R} \wedge \\ \quad \quad \langle\cdot,\cdot\rangle_{v}(\operatorname{finsum}(\operatorname{normed\text{-}ag\text{-}as\text{-}abelian\text{-}group}(\operatorname{ip\text{-}normed\text{-}ag}(v)), \\ \quad \quad \quad \quad (\lambda k \in a.\; \operatorname{act}(v)(\langle\cdot,\cdot\rangle_{v}(x, e(k)), e(k))), w), \\ \quad \quad \quad \operatorname{finsum}(\operatorname{normed\text{-}ag\text{-}as\text{-}abelian\text{-}group}(\operatorname{ip\text{-}normed\text{-}ag}(v)), \\ \quad \quad \quad \quad (\lambda k \in a.\; \operatorname{act}(v)(\langle\cdot,\cdot\rangle_{v}(x, e(k)), e(k))), w)) \\ \quad \quad =\; \operatorname{finsum}(\operatorname{rr\text{-}pos\text{-}star\text{-}add\text{-}monoid}, \\ \quad \quad \quad (\lambda k \in a.\; (\operatorname{magnitude}(\langle\cdot,\cdot\rangle_{v}(x, e(k))) \cdot \operatorname{magnitude}(\langle\cdot,\cdot\rangle_{v}(x, e(k))))), \\ \quad \quad \quad w) \wedge \\ \quad \quad \langle\cdot,\cdot\rangle_{v}(\operatorname{finsum}(\operatorname{normed\text{-}ag\text{-}as\text{-}abelian\text{-}group}(\operatorname{ip\text{-}normed\text{-}ag}(v)), \\ \quad \quad \quad \quad (\lambda k \in a.\; \operatorname{act}(v)(\langle\cdot,\cdot\rangle_{v}(x, e(k)), e(k))), w), \\ \quad \quad \quad x) \\ \quad \quad =\; \operatorname{finsum}(\operatorname{rr\text{-}pos\text{-}star\text{-}add\text{-}monoid}, \\ \quad \quad \quad (\lambda k \in a.\; (\operatorname{magnitude}(\langle\cdot,\cdot\rangle_{v}(x, e(k))) \cdot \operatorname{magnitude}(\langle\cdot,\cdot\rangle_{v}(x, e(k))))), \\ \quad \quad \quad w) \end{array}}}
\vnbprop{onb-bessel-partial}{\vnbstmt{\begin{array}{@{}l@{}} \forall v, a, e.; \forall x \in \operatorname{vec}(v), w \in \mathbf{Set}. \\ \quad \operatorname{is\text{-}complex\text{-}inner\text{-}product\text{-}space}(v) \wedge \\ \quad \operatorname{is\text{-}orthonormal\text{-}family}(v, a, e) \wedge \\ \quad |w| \in \mathbb{N} \wedge \\ \quad w \subseteq a \Rightarrow \\ \quad \quad \operatorname{finsum}(\operatorname{rr\text{-}pos\text{-}star\text{-}add\text{-}monoid}, \\ \quad \quad \quad (\lambda k \in a.\; (\operatorname{magnitude}(\langle\cdot,\cdot\rangle_{v}(x, e(k))) \cdot \operatorname{magnitude}(\langle\cdot,\cdot\rangle_{v}(x, e(k))))), \\ \quad \quad \quad w) \\ \quad \quad \leq\; \langle\cdot,\cdot\rangle_{v}(x, x) \end{array}}}
\vnbprop{bessel-inequality}{\vnbstmt{\begin{array}{@{}l@{}} \forall v, a, e.; \forall x \in \operatorname{vec}(v). \\ \quad (\operatorname{is\text{-}complex\text{-}inner\text{-}product\text{-}space}(v) \wedge \operatorname{is\text{-}orthonormal\text{-}family}(v, a, e)) \Rightarrow  \\ \quad \quad \operatorname{esum}((\lambda k \in a.\; (\operatorname{magnitude}(\langle\cdot,\cdot\rangle_{v}(x, e(k))) \cdot \operatorname{magnitude}(\langle\cdot,\cdot\rangle_{v}(x, e(k)))))) \leq \langle\cdot,\cdot\rangle_{v}(x, x) \end{array}}}
\vnbprop{onb-cauchy-net-choice-seq}{\vnbstmt{\begin{array}{@{}l@{}} \forall p, f, d. \\ \quad \forall s. \\ \quad \quad \operatorname{pos\text{-}rr}(s) \Rightarrow  \\ \quad \quad \quad \exists a \in \mathbf{Set}. \\ \quad \quad \quad \quad |a| \in \mathbb{N} \wedge \\ \quad \quad \quad \quad a \subseteq \operatorname{dom}(f) \wedge \\ \quad \quad \quad \quad \forall \mathit{cna}, b \in \mathbf{Set}. \\ \quad \quad \quad \quad \quad |\mathit{cna}| \in \mathbb{N} \wedge \\ \quad \quad \quad \quad \quad a \subseteq \mathit{cna} \wedge \\ \quad \quad \quad \quad \quad \mathit{cna} \subseteq \operatorname{dom}(f) \wedge \\ \quad \quad \quad \quad \quad |b| \in \mathbb{N} \wedge \\ \quad \quad \quad \quad \quad a \subseteq b \wedge \\ \quad \quad \quad \quad \quad b \subseteq \operatorname{dom}(f) \Rightarrow \\ \quad \quad \quad \quad \quad \quad \operatorname{dist}(\operatorname{nag\text{-}metric\text{-}space}(p))(\operatorname{finsum}(\operatorname{normed\text{-}ag\text{-}as\text{-}abelian\text{-}group}(p), f, \mathit{cna}), \\ \quad \quad \quad \quad \quad \quad \quad \operatorname{finsum}(\operatorname{normed\text{-}ag\text{-}as\text{-}abelian\text{-}group}(p), f, b)) \\ \quad \quad \quad \quad \quad \quad <\; s \wedge \\ \quad \forall k \in \mathbb{N}.\; \operatorname{pos\text{-}rr}(d(k)) \Rightarrow \\ \quad \quad \exists c. \forall n \in \mathbb{N}. \\ \quad \quad \quad c(n) \in \mathbf{Set} \wedge \\ \quad \quad \quad |c(n)| \in \mathbb{N} \wedge \\ \quad \quad \quad c(n) \subseteq \operatorname{dom}(f) \wedge \\ \quad \quad \quad \forall \mathit{cna}, b \in \mathbf{Set}. \\ \quad \quad \quad \quad |\mathit{cna}| \in \mathbb{N} \wedge \\ \quad \quad \quad \quad c(n) \subseteq \mathit{cna} \wedge \\ \quad \quad \quad \quad \mathit{cna} \subseteq \operatorname{dom}(f) \wedge \\ \quad \quad \quad \quad |b| \in \mathbb{N} \wedge \\ \quad \quad \quad \quad c(n) \subseteq b \wedge \\ \quad \quad \quad \quad b \subseteq \operatorname{dom}(f) \Rightarrow \\ \quad \quad \quad \quad \quad \operatorname{dist}(\operatorname{nag\text{-}metric\text{-}space}(p))(\operatorname{finsum}(\operatorname{normed\text{-}ag\text{-}as\text{-}abelian\text{-}group}(p), f, \mathit{cna}), \\ \quad \quad \quad \quad \quad \quad \operatorname{finsum}(\operatorname{normed\text{-}ag\text{-}as\text{-}abelian\text{-}group}(p), f, b)) \\ \quad \quad \quad \quad \quad <\; d(n) \end{array}}}
\vnbprop{summable-of-cauchy-net}{\vnbstmt{\begin{array}{@{}l@{}} \forall p.; \forall f \in (\operatorname{dom}(f) \to \operatorname{carr}(p)). \\ \quad \operatorname{is\text{-}normed\text{-}ag}(p) \wedge \\ \quad \operatorname{is\text{-}complete}(\operatorname{nag\text{-}metric\text{-}space}(p)) \wedge \\ \quad \forall s. \\ \quad \quad \operatorname{pos\text{-}rr}(s) \Rightarrow  \\ \quad \quad \quad \exists n \in \mathbf{Set}. \\ \quad \quad \quad \quad |n| \in \mathbb{N} \wedge \\ \quad \quad \quad \quad n \subseteq \operatorname{dom}(f) \wedge \\ \quad \quad \quad \quad \forall a, b \in \mathbf{Set}. \\ \quad \quad \quad \quad \quad |a| \in \mathbb{N} \wedge \\ \quad \quad \quad \quad \quad n \subseteq a \wedge \\ \quad \quad \quad \quad \quad a \subseteq \operatorname{dom}(f) \wedge \\ \quad \quad \quad \quad \quad |b| \in \mathbb{N} \wedge \\ \quad \quad \quad \quad \quad n \subseteq b \wedge \\ \quad \quad \quad \quad \quad b \subseteq \operatorname{dom}(f) \Rightarrow \\ \quad \quad \quad \quad \quad \quad \operatorname{dist}(\operatorname{nag\text{-}metric\text{-}space}(p))(\operatorname{finsum}(\operatorname{normed\text{-}ag\text{-}as\text{-}abelian\text{-}group}(p), f, a), \\ \quad \quad \quad \quad \quad \quad \quad \operatorname{finsum}(\operatorname{normed\text{-}ag\text{-}as\text{-}abelian\text{-}group}(p), f, b)) \\ \quad \quad \quad \quad \quad \quad <\; s \Rightarrow \\ \quad \quad \operatorname{is\text{-}summable}(p, f) \end{array}}}
\vnbprop{onb-partial-nested}{\vnbstmt{\begin{array}{@{}l@{}} \forall v, a, e.; \forall x \in \operatorname{vec}(v), \mathit{ona}, w \in \mathbf{Set}. \\ \quad \operatorname{is\text{-}complex\text{-}inner\text{-}product\text{-}space}(v) \wedge \\ \quad \operatorname{is\text{-}orthonormal\text{-}family}(v, a, e) \wedge \\ \quad |\mathit{ona}| \in \mathbb{N} \wedge \\ \quad \mathit{ona} \subseteq a \wedge \\ \quad |w| \in \mathbb{N} \wedge \\ \quad w \subseteq \mathit{ona} \Rightarrow \\ \quad \quad \langle\cdot,\cdot\rangle_{v}(\operatorname{finsum}(\operatorname{normed\text{-}ag\text{-}as\text{-}abelian\text{-}group}(\operatorname{ip\text{-}normed\text{-}ag}(v)), \\ \quad \quad \quad \quad (\lambda k \in a.\; \operatorname{act}(v)(\langle\cdot,\cdot\rangle_{v}(x, e(k)), e(k))), w), \\ \quad \quad \quad \operatorname{finsum}(\operatorname{normed\text{-}ag\text{-}as\text{-}abelian\text{-}group}(\operatorname{ip\text{-}normed\text{-}ag}(v)), \\ \quad \quad \quad \quad (\lambda k \in a.\; \operatorname{act}(v)(\langle\cdot,\cdot\rangle_{v}(x, e(k)), e(k))), \mathit{ona})) \\ \quad \quad =\; \operatorname{finsum}(\operatorname{rr\text{-}pos\text{-}star\text{-}add\text{-}monoid}, \\ \quad \quad \quad (\lambda k \in a.\; (\operatorname{magnitude}(\langle\cdot,\cdot\rangle_{v}(x, e(k))) \cdot \operatorname{magnitude}(\langle\cdot,\cdot\rangle_{v}(x, e(k))))), \\ \quad \quad \quad w) \end{array}}}
\vnbprop{onb-partial-diff}{\vnbstmt{\begin{array}{@{}l@{}} \forall v, a, e.; \forall x \in \operatorname{vec}(v), \mathit{ona}, w \in \mathbf{Set}. \\ \quad \operatorname{is\text{-}complex\text{-}inner\text{-}product\text{-}space}(v) \wedge \\ \quad \operatorname{is\text{-}orthonormal\text{-}family}(v, a, e) \wedge \\ \quad |\mathit{ona}| \in \mathbb{N} \wedge \\ \quad \mathit{ona} \subseteq a \wedge \\ \quad |w| \in \mathbb{N} \wedge \\ \quad w \subseteq \mathit{ona} \Rightarrow \\ \quad \quad \langle\cdot,\cdot\rangle_{v}(\operatorname{vadd}(v)(\operatorname{finsum}(\operatorname{normed\text{-}ag\text{-}as\text{-}abelian\text{-}group}(\operatorname{ip\text{-}normed\text{-}ag}(v)), \\ \quad \quad \quad \quad \quad (\lambda k \in a.\; \operatorname{act}(v)(\langle\cdot,\cdot\rangle_{v}(x, e(k)), e(k))), \mathit{ona}), \\ \quad \quad \quad \quad \operatorname{vneg}(v)(\operatorname{finsum}(\operatorname{normed\text{-}ag\text{-}as\text{-}abelian\text{-}group}(\operatorname{ip\text{-}normed\text{-}ag}(v)), \\ \quad \quad \quad \quad \quad \quad (\lambda k \in a.\; \operatorname{act}(v)(\langle\cdot,\cdot\rangle_{v}(x, e(k)), e(k))), w))), \\ \quad \quad \quad \operatorname{vadd}(v)(\operatorname{finsum}(\operatorname{normed\text{-}ag\text{-}as\text{-}abelian\text{-}group}(\operatorname{ip\text{-}normed\text{-}ag}(v)), \\ \quad \quad \quad \quad \quad (\lambda k \in a.\; \operatorname{act}(v)(\langle\cdot,\cdot\rangle_{v}(x, e(k)), e(k))), \mathit{ona}), \\ \quad \quad \quad \quad \operatorname{vneg}(v)(\operatorname{finsum}(\operatorname{normed\text{-}ag\text{-}as\text{-}abelian\text{-}group}(\operatorname{ip\text{-}normed\text{-}ag}(v)), \\ \quad \quad \quad \quad \quad \quad (\lambda k \in a.\; \operatorname{act}(v)(\langle\cdot,\cdot\rangle_{v}(x, e(k)), e(k))), w)))) \\ \quad \quad =\; \operatorname{finsum}(\operatorname{rr\text{-}pos\text{-}star\text{-}add\text{-}monoid}, (\lambda k \in a.\; (\operatorname{magnitude}(\langle\cdot,\cdot\rangle_{v}(x, e(k))) \cdot \operatorname{magnitude}(\langle\cdot,\cdot\rangle_{v}(x, e(k))))), \mathit{ona}) + -\operatorname{finsum}(\operatorname{rr\text{-}pos\text{-}star\text{-}add\text{-}monoid}, (\lambda k \in a.\; (\operatorname{magnitude}(\langle\cdot,\cdot\rangle_{v}(x, e(k))) \cdot \operatorname{magnitude}(\langle\cdot,\cdot\rangle_{v}(x, e(k))))), w) \end{array}}}
\vnbprop{onb-expansion-cauchy-net}{\vnbstmt{\begin{array}{@{}l@{}} \forall v, a, e.; \forall x \in \operatorname{vec}(v), s. \\ \quad (\operatorname{is\text{-}complex\text{-}inner\text{-}product\text{-}space}(v) \wedge (\operatorname{is\text{-}orthonormal\text{-}family}(v, a, e) \wedge \operatorname{pos\text{-}rr}(s))) \Rightarrow  \\ \quad \quad \exists n \in \mathbf{Set}. \\ \quad \quad \quad |n| \in \mathbb{N} \wedge \\ \quad \quad \quad n \subseteq \operatorname{dom}((\lambda k \in a.\; \operatorname{act}(v)(\langle\cdot,\cdot\rangle_{v}(x, e(k)), e(k)))) \wedge \\ \quad \quad \quad \forall \mathit{cna}, b \in \mathbf{Set}. \\ \quad \quad \quad \quad |\mathit{cna}| \in \mathbb{N} \wedge \\ \quad \quad \quad \quad n \subseteq \mathit{cna} \wedge \\ \quad \quad \quad \quad \mathit{cna} \subseteq \operatorname{dom}((\lambda k \in a.\; \operatorname{act}(v)(\langle\cdot,\cdot\rangle_{v}(x, e(k)), e(k)))) \wedge \\ \quad \quad \quad \quad |b| \in \mathbb{N} \wedge \\ \quad \quad \quad \quad n \subseteq b \wedge \\ \quad \quad \quad \quad b \subseteq \operatorname{dom}((\lambda k \in a.\; \operatorname{act}(v)(\langle\cdot,\cdot\rangle_{v}(x, e(k)), e(k)))) \Rightarrow \\ \quad \quad \quad \quad \quad \operatorname{dist}(\operatorname{nag\text{-}metric\text{-}space}(\operatorname{ip\text{-}normed\text{-}ag}(v)))(\operatorname{finsum}(\operatorname{normed\text{-}ag\text{-}as\text{-}abelian\text{-}group}(\operatorname{ip\text{-}normed\text{-}ag}(v)), \\ \quad \quad \quad \quad \quad \quad \quad (\lambda k \in a.\; \operatorname{act}(v)(\langle\cdot,\cdot\rangle_{v}(x, e(k)), e(k))), \mathit{cna}), \\ \quad \quad \quad \quad \quad \quad \operatorname{finsum}(\operatorname{normed\text{-}ag\text{-}as\text{-}abelian\text{-}group}(\operatorname{ip\text{-}normed\text{-}ag}(v)), \\ \quad \quad \quad \quad \quad \quad \quad (\lambda k \in a.\; \operatorname{act}(v)(\langle\cdot,\cdot\rangle_{v}(x, e(k)), e(k))), b)) \\ \quad \quad \quad \quad \quad <\; s \end{array}}}
\vnbprop{onb-coef-limit}{\vnbstmt{\begin{array}{@{}l@{}} \forall v, a, e.; \forall x, l \in \operatorname{vec}(v), p \in a. \\ \quad \operatorname{is\text{-}complex\text{-}inner\text{-}product\text{-}space}(v) \wedge \\ \quad \operatorname{is\text{-}orthonormal\text{-}family}(v, a, e) \wedge \\ \quad \operatorname{sums\text{-}to}(\operatorname{ip\text{-}normed\text{-}ag}(v), (\lambda k \in a.\; \operatorname{act}(v)(\langle\cdot,\cdot\rangle_{v}(x, e(k)), e(k))), l) \Rightarrow \\ \quad \quad \langle\cdot,\cdot\rangle_{v}(l, e(p)) = \langle\cdot,\cdot\rangle_{v}(x, e(p)) \end{array}}}
\vnbprop{onb-expansion}{\vnbstmt{\begin{array}{@{}l@{}} \forall v, a, e.; \forall x \in \operatorname{vec}(v). \\ \quad (\operatorname{is\text{-}hilbert\text{-}space}(v) \wedge \operatorname{is\text{-}orthonormal\text{-}basis}(v, a, e)) \Rightarrow  \\ \quad \quad \operatorname{sums\text{-}to}(\operatorname{ip\text{-}normed\text{-}ag}(v), (\lambda k \in a.\; \operatorname{act}(v)(\langle\cdot,\cdot\rangle_{v}(x, e(k)), e(k))), x) \end{array}}}
\vnbprop{onb-residual}{\vnbstmt{\begin{array}{@{}l@{}} \forall v, a, e.; \forall x \in \operatorname{vec}(v), w \in \mathbf{Set}. \\ \quad \operatorname{is\text{-}complex\text{-}inner\text{-}product\text{-}space}(v) \wedge \\ \quad \operatorname{is\text{-}orthonormal\text{-}family}(v, a, e) \wedge \\ \quad |w| \in \mathbb{N} \wedge \\ \quad w \subseteq a \Rightarrow \\ \quad \quad \langle\cdot,\cdot\rangle_{v}(\operatorname{vadd}(v)(x, \\ \quad \quad \quad \quad \operatorname{vneg}(v)(\operatorname{finsum}(\operatorname{normed\text{-}ag\text{-}as\text{-}abelian\text{-}group}(\operatorname{ip\text{-}normed\text{-}ag}(v)), \\ \quad \quad \quad \quad \quad \quad (\lambda k \in a.\; \operatorname{act}(v)(\langle\cdot,\cdot\rangle_{v}(x, e(k)), e(k))), w))), \\ \quad \quad \quad \operatorname{vadd}(v)(x, \\ \quad \quad \quad \quad \operatorname{vneg}(v)(\operatorname{finsum}(\operatorname{normed\text{-}ag\text{-}as\text{-}abelian\text{-}group}(\operatorname{ip\text{-}normed\text{-}ag}(v)), \\ \quad \quad \quad \quad \quad \quad (\lambda k \in a.\; \operatorname{act}(v)(\langle\cdot,\cdot\rangle_{v}(x, e(k)), e(k))), w)))) \\ \quad \quad =\; \langle\cdot,\cdot\rangle_{v}(x, x) + -\operatorname{finsum}(\operatorname{rr\text{-}pos\text{-}star\text{-}add\text{-}monoid}, (\lambda k \in a.\; (\operatorname{magnitude}(\langle\cdot,\cdot\rangle_{v}(x, e(k))) \cdot \operatorname{magnitude}(\langle\cdot,\cdot\rangle_{v}(x, e(k))))), w) \end{array}}}
\vnbprop{parseval-identity}{\vnbstmt{\begin{array}{@{}l@{}} \forall v, a, e.; \forall x \in \operatorname{vec}(v). \\ \quad (\operatorname{is\text{-}hilbert\text{-}space}(v) \wedge \operatorname{is\text{-}orthonormal\text{-}basis}(v, a, e)) \Rightarrow  \\ \quad \quad \operatorname{esum}((\lambda k \in a.\; (\operatorname{magnitude}(\langle\cdot,\cdot\rangle_{v}(x, e(k))) \cdot \operatorname{magnitude}(\langle\cdot,\cdot\rangle_{v}(x, e(k)))))) = \langle\cdot,\cdot\rangle_{v}(x, x) \end{array}}}

\vnbbookfile{theorem-library/abs-summable-cn}{theorem-library/abs-summable-cn}
\noindent{\small abs-summable-cn.scm -- wave of 2026-11-09/10 (agent BA-ONB2), wired 2026-11-10 from scratchpad/baonb2-1109/bo2-abs.scm. Library: analysis.}\par
\vnbprop{abs-summable-implies-summable-cn}{\vnbstmt{\begin{array}{@{}l@{}} \forall p.; \forall f \in (\operatorname{dom}(f) \to \operatorname{carr}(p)). \\ \quad ((\operatorname{is\text{-}normed\text{-}ag}(p) \wedge \operatorname{is\text{-}complete}(\operatorname{nag\text{-}metric\text{-}space}(p))) \wedge \operatorname{is\text{-}absolutely\text{-}summable}(p, f)) \Rightarrow  \\ \quad \quad \operatorname{is\text{-}summable}(p, f) \end{array}}}

\vnbbookfile{structure-library/sip-hilbert-vocab}{structure-library/sip-hilbert-vocab}
\noindent{\small sip-hilbert-vocab.scm -- REAL HILBERT SPACES: the norm and the metric of a semi-inner-product space, completeness, the orthogonal complement and the orthogonal projection.  Vocabulary only: no proof, no support.  (SIP-D, 2026-10-19; scratchpad/triage/SIP-D.md.)  The proofs are theorem-library/sip-d1.scm}\par
\par\medskip\noindent\textbf{Definition} (\texttt{sip-normed-ag}).\;$\operatorname{sip\text{-}normed\text{-}ag}(v) :=$
\begin{equation*}
{[}\operatorname{vec}(v), \operatorname{vadd}(v), \operatorname{vzero}(v), \operatorname{vneg}(v), (\lambda \operatorname{shn\_} \in \operatorname{vec}(v).\; \sqrt{\langle\cdot\mid\cdot\rangle_{v}(\operatorname{shn\_}, \operatorname{shn\_})})]
\end{equation*}

\par\medskip\noindent\textbf{Definition} (\texttt{sip-metric-space}).\;$\operatorname{sip\text{-}metric\text{-}space}(v) :=$ $\operatorname{nag\text{-}metric\text{-}space}(\operatorname{sip\text{-}normed\text{-}ag}(v))$
\par\medskip\noindent\textbf{Definition} (\texttt{is-real-hilbert-space}).\;For every $v$, 
\begin{equation*}
\operatorname{is\text{-}real\text{-}hilbert\text{-}space}(v) \Leftrightarrow (\operatorname{is\text{-}pre\text{-}hilbert\text{-}space}(v) \wedge \operatorname{is\text{-}complete}(\operatorname{sip\text{-}metric\text{-}space}(v)))
\end{equation*}

\par\medskip\noindent\textbf{Definition} (\texttt{sip-perp}).\;${\mathit{ks}}^{\perp} :=$
\begin{equation*}
\{\operatorname{spp\_} \in \operatorname{vec}(v) : \forall \operatorname{spk\_} \in \mathit{ks}.\; (\langle\cdot\mid\cdot\rangle_{v}(\operatorname{spp\_}, \operatorname{spk\_}) = 0)\}
\end{equation*}

\par\medskip\noindent\textbf{Definition} (\texttt{is-convex-subset}).\;For all $v$ and $\mathit{cs}$, 
\begin{equation*}
\begin{array}{@{}l@{}} \operatorname{is\text{-}convex\text{-}subset}(v, \mathit{cs}) \\ \Leftrightarrow\; \mathit{cs} \subseteq \operatorname{vec}(v) \wedge \\ \forall \operatorname{sxa\_}, \operatorname{sxb\_} \in \mathit{cs}, \operatorname{sxt\_} \in \mathbb{R}. \\ \quad (0 \leq \operatorname{sxt\_}) \Rightarrow  \\ \quad (\operatorname{sxt\_} \leq 1) \Rightarrow  \\ \quad \quad \operatorname{vadd}(v)(\operatorname{act}(v)(-(1, \operatorname{sxt\_}), \operatorname{sxa\_}), \operatorname{act}(v)(\operatorname{sxt\_}, \operatorname{sxb\_})) \in \mathit{cs} \end{array}
\end{equation*}

\par\medskip\noindent\textbf{Definition} (\texttt{is-real-orthogonal-projection}).\;For all $v$, $\mathit{ks}$ and $\mathit{pk}$, 
\begin{equation*}
\begin{array}{@{}l@{}} \operatorname{is\text{-}real\text{-}orthogonal\text{-}projection}(v, \mathit{ks}, \mathit{pk}) \\ \Leftrightarrow\; \mathit{pk} \in (\operatorname{vec}(v) \to \operatorname{vec}(v)) \wedge \\ \forall \operatorname{sop\_} \in \operatorname{vec}(v). \\ \quad (\operatorname{pk}(\operatorname{sop\_}) \in \mathit{ks}) \wedge (\operatorname{vadd}(v)(\operatorname{sop\_}, \operatorname{vneg}(v)(\operatorname{pk}(\operatorname{sop\_}))) \in {\mathit{ks}}^{\perp}) \end{array}
\end{equation*}


\vnbbookfile{structure-library/minkowski-functional}{structure-library/minkowski-functional}
\noindent{\small mink-vocab.scm -- absorbing and symmetric sets, the Minkowski functional (topological-vector-space.pdf section 3.3, Def 2.28).  Vocabulary only.  Agent DUAL-FIX, 2026-10-26 (NEW DEFINITIONS, the user's to confirm; added under the rule "missing vocabulary is not a stopping point"). DESTINATION: ...}\par
\par\medskip\noindent\textbf{Definition} (\texttt{is-absorbing-subset}).\;For all $m$ and $\mathit{cs}$, 
\begin{equation*}
\begin{array}{@{}l@{}} \operatorname{is\text{-}absorbing\text{-}subset}(m, \mathit{cs}) \\ \Leftrightarrow\; \mathit{cs} \subseteq \operatorname{vec}(m) \wedge \\ \forall \operatorname{abx\_} \in \operatorname{vec}(m). \exists \operatorname{abr\_} \in \mathbb{R}. \\ \quad (0 < \operatorname{abr\_}) \wedge \exists \operatorname{aby\_} \in \mathit{cs}.\; (\operatorname{abx\_} = \operatorname{act}(m)(\operatorname{abr\_}, \operatorname{aby\_})) \end{array}
\end{equation*}

\par\medskip\noindent\textbf{Definition} (\texttt{is-symmetric-subset}).\;For all $m$ and $\mathit{cs}$, 
\begin{equation*}
\operatorname{is\text{-}symmetric\text{-}subset}(m, \mathit{cs}) \Leftrightarrow ((\mathit{cs} \subseteq \operatorname{vec}(m)) \wedge \forall \operatorname{syx\_} \in \mathit{cs}.\; (\operatorname{act}(m)(-(1), \operatorname{syx\_}) \in \mathit{cs}))
\end{equation*}

\par\medskip\noindent\textbf{Definition} (\texttt{mink-set}).\;$\operatorname{mink\text{-}set}(m, \mathit{cs}, x) :=$
\begin{equation*}
\{\operatorname{mkr\_} \in \mathbb{R} : ((0 < \operatorname{mkr\_}) \wedge \exists \operatorname{mky\_} \in \mathit{cs}.\; (x = \operatorname{act}(m)(\operatorname{mkr\_}, \operatorname{mky\_})))\}
\end{equation*}

\par\medskip\noindent\textbf{Definition} (\texttt{minkowski-functional}).\;$p_{\mathit{cs}} :=$ $(\lambda \operatorname{mkx\_} \in \operatorname{vec}(m).\; \operatorname{inf}(\operatorname{mink\text{-}set}(m, \mathit{cs}, \operatorname{mkx\_})))$

\vnbbookfile{theorem-library/tvs-minkowski}{theorem-library/tvs-minkowski}
\noindent{\small minkowski.scm -- topological-vector-space.pdf Prop 2.29: the Minkowski functional of a symmetric, absorbing, convex subset of a real vector space is a seminorm.  Agent DUAL-FIX, 2026-10-26. DESTINATION: theorem-library/tvs-minkowski.scm.  Library: analysis.  Helper prefix: mkw-. Needs: ...}\par
\vnbprop{rmod-act-type}{\vnbstmt{\begin{array}{@{}l@{}} \forall m, r, x. \\ \quad \operatorname{is\text{-}module}(m) \wedge \\ \quad \operatorname{scal}(m) = \operatorname{normed\text{-}field\text{-}as\text{-}commutative\text{-}ring}(\operatorname{rr\text{-}normed\text{-}field}) \wedge \\ \quad r \in \mathbb{R} \wedge \\ \quad x \in \operatorname{vec}(m) \Rightarrow \\ \quad \quad \operatorname{act}(m)(r, x) \in \operatorname{vec}(m) \end{array}}}
\vnbprop{rmod-act-mul}{\vnbstmt{\begin{array}{@{}l@{}} \forall m, r, s, x. \\ \quad \operatorname{is\text{-}module}(m) \wedge \\ \quad \operatorname{scal}(m) = \operatorname{normed\text{-}field\text{-}as\text{-}commutative\text{-}ring}(\operatorname{rr\text{-}normed\text{-}field}) \wedge \\ \quad r \in \mathbb{R} \wedge \\ \quad s \in \mathbb{R} \wedge \\ \quad x \in \operatorname{vec}(m) \Rightarrow \\ \quad \quad \operatorname{act}(m)((r \cdot s), x) = \operatorname{act}(m)(r, \operatorname{act}(m)(s, x)) \end{array}}}
\vnbprop{rmod-act-add}{\vnbstmt{\begin{array}{@{}l@{}} \forall m, r, s, x. \\ \quad \operatorname{is\text{-}module}(m) \wedge \\ \quad \operatorname{scal}(m) = \operatorname{normed\text{-}field\text{-}as\text{-}commutative\text{-}ring}(\operatorname{rr\text{-}normed\text{-}field}) \wedge \\ \quad r \in \mathbb{R} \wedge \\ \quad s \in \mathbb{R} \wedge \\ \quad x \in \operatorname{vec}(m) \Rightarrow \\ \quad \quad \operatorname{act}(m)((r + s), x) = \operatorname{vadd}(m)(\operatorname{act}(m)(r, x), \operatorname{act}(m)(s, x)) \end{array}}}
\vnbprop{rmod-act-one}{\vnbstmt{\begin{array}{@{}l@{}} \forall m, x. \\ \quad (\operatorname{is\text{-}module}(m) \wedge ((\operatorname{scal}(m) = \operatorname{normed\text{-}field\text{-}as\text{-}commutative\text{-}ring}(\operatorname{rr\text{-}normed\text{-}field})) \wedge (x \in \operatorname{vec}(m)))) \Rightarrow  \\ \quad \quad \operatorname{act}(m)(1, x) = x \end{array}}}
\vnbprop{rmod-act-zero}{\vnbstmt{\begin{array}{@{}l@{}} \forall m, x. \\ \quad (\operatorname{is\text{-}module}(m) \wedge ((\operatorname{scal}(m) = \operatorname{normed\text{-}field\text{-}as\text{-}commutative\text{-}ring}(\operatorname{rr\text{-}normed\text{-}field})) \wedge (x \in \operatorname{vec}(m)))) \Rightarrow  \\ \quad \quad \operatorname{act}(m)(0, x) = \operatorname{vzero}(m) \end{array}}}
\vnbprop{rmod-act-vzero}{\vnbstmt{\begin{array}{@{}l@{}} \forall m, r. \\ \quad (\operatorname{is\text{-}module}(m) \wedge ((\operatorname{scal}(m) = \operatorname{normed\text{-}field\text{-}as\text{-}commutative\text{-}ring}(\operatorname{rr\text{-}normed\text{-}field})) \wedge (r \in \mathbb{R}))) \Rightarrow  \\ \quad \quad \operatorname{act}(m)(r, \operatorname{vzero}(m)) = \operatorname{vzero}(m) \end{array}}}
\vnbprop{rmod-act-distrib}{\vnbstmt{\begin{array}{@{}l@{}} \forall m, r, x, y. \\ \quad \operatorname{is\text{-}module}(m) \wedge \\ \quad \operatorname{scal}(m) = \operatorname{normed\text{-}field\text{-}as\text{-}commutative\text{-}ring}(\operatorname{rr\text{-}normed\text{-}field}) \wedge \\ \quad r \in \mathbb{R} \wedge \\ \quad x \in \operatorname{vec}(m) \wedge \\ \quad y \in \operatorname{vec}(m) \Rightarrow \\ \quad \quad \operatorname{act}(m)(r, \operatorname{vadd}(m)(x, y)) = \operatorname{vadd}(m)(\operatorname{act}(m)(r, x), \operatorname{act}(m)(r, y)) \end{array}}}
\vnbprop{rmod-vzero-left}{\vnbstmt{\forall m, x.\; ((\operatorname{is\text{-}module}(m) \wedge (x \in \operatorname{vec}(m))) \Rightarrow (\operatorname{vadd}(m)(\operatorname{vzero}(m), x) = x))}}
\vnbprop{mink-vzero-in}{\vnbstmt{\begin{array}{@{}l@{}} \forall m, s. \\ \quad \operatorname{is\text{-}module}(m) \wedge \\ \quad \operatorname{scal}(m) = \operatorname{normed\text{-}field\text{-}as\text{-}commutative\text{-}ring}(\operatorname{rr\text{-}normed\text{-}field}) \wedge \\ \quad \operatorname{is\text{-}convex\text{-}subset}(m, s) \wedge \\ \quad \operatorname{is\text{-}symmetric\text{-}subset}(m, s) \wedge \\ \quad \operatorname{is\text{-}absorbing\text{-}subset}(m, s) \Rightarrow \\ \quad \quad \operatorname{vzero}(m) \in s \end{array}}}
\vnbprop{mink-scale-in}{\vnbstmt{\begin{array}{@{}l@{}} \forall m, s, y, t. \\ \quad \operatorname{is\text{-}module}(m) \wedge \\ \quad \operatorname{scal}(m) = \operatorname{normed\text{-}field\text{-}as\text{-}commutative\text{-}ring}(\operatorname{rr\text{-}normed\text{-}field}) \wedge \\ \quad \operatorname{is\text{-}convex\text{-}subset}(m, s) \wedge \\ \quad \operatorname{is\text{-}absorbing\text{-}subset}(m, s) \wedge \\ \quad \operatorname{vzero}(m) \in s \wedge \\ \quad y \in s \wedge \\ \quad t \in \mathbb{R} \wedge \\ \quad 0 \leq t \wedge \\ \quad t \leq 1 \Rightarrow \\ \quad \quad \operatorname{act}(m)(t, y) \in s \end{array}}}
\vnbprop{mink-set-unfold}{\vnbstmt{\begin{array}{@{}l@{}} \forall m, s, x. \\ \quad (\operatorname{mink\text{-}set}(m, s, x) \simeq \{r \in \mathbb{R} : ((0 < r) \wedge \exists y \in s.\; (x = \operatorname{act}(m)(r, y)))\}) \end{array}}}
\vnbprop{mink-set-intro}{\vnbstmt{\begin{array}{@{}l@{}} \forall m, s, x, r, y. \\ \quad ((r \in \mathbb{R}) \wedge ((0 < r) \wedge ((y \in s) \wedge (x = \operatorname{act}(m)(r, y))))) \Rightarrow  \\ \quad \quad r \in \operatorname{mink\text{-}set}(m, s, x) \end{array}}}
\vnbprop{mink-set-elim}{\vnbstmt{\begin{array}{@{}l@{}} \forall m, s, x.; \forall r \in \operatorname{mink\text{-}set}(m, s, x). \\ \quad (r \in \mathbb{R}) \wedge ((0 < r) \wedge \exists y \in s.\; (x = \operatorname{act}(m)(r, y))) \end{array}}}
\vnbprop{mink-set-glb}{\vnbstmt{\begin{array}{@{}l@{}} \forall m, s, z. \\ \quad \operatorname{is\text{-}module}(m) \wedge \\ \quad \operatorname{scal}(m) = \operatorname{normed\text{-}field\text{-}as\text{-}commutative\text{-}ring}(\operatorname{rr\text{-}normed\text{-}field}) \wedge \\ \quad \operatorname{is\text{-}convex\text{-}subset}(m, s) \wedge \\ \quad \operatorname{is\text{-}absorbing\text{-}subset}(m, s) \wedge \\ \quad \operatorname{vzero}(m) \in s \wedge \\ \quad z \in \operatorname{vec}(m) \Rightarrow \\ \quad \quad \operatorname{inf}(\operatorname{mink\text{-}set}(m, s, z)) \in \mathbb{R} \wedge \\ \quad \quad \forall w \in \operatorname{mink\text{-}set}(m, s, z).\; (\operatorname{inf}(\operatorname{mink\text{-}set}(m, s, z)) \leq w) \wedge \\ \quad \quad \forall b \in \mathbb{R}. \\ \quad \quad \quad \forall w \in \operatorname{mink\text{-}set}(m, s, z).\; (b \leq w) \Rightarrow (b \leq \operatorname{inf}(\operatorname{mink\text{-}set}(m, s, z))) \end{array}}}
\vnbprop{mink-value}{\vnbstmt{\begin{array}{@{}l@{}} \forall m, s.; \forall z \in \operatorname{vec}(m). \\ \quad (p_{s}(z) \simeq \operatorname{inf}(\operatorname{mink\text{-}set}(m, s, z))) \end{array}}}
\vnbprop{mink-rr-scale-le}{\vnbstmt{\begin{array}{@{}l@{}} \forall a, b, c. \\ \quad ((a \in \mathbb{R}) \wedge ((b \in \mathbb{R}) \wedge ((c \in \mathbb{R}) \wedge ((0 \leq c) \wedge (a \leq b))))) \Rightarrow ((c \cdot a) \leq (c \cdot b)) \end{array}}}
\vnbprop{mink-rr-mul-pos}{\vnbstmt{\forall a, b.\; (((a \in \mathbb{R}) \wedge ((b \in \mathbb{R}) \wedge ((0 < a) \wedge (0 < b)))) \Rightarrow (0 < (a \cdot b)))}}
\vnbprop{mink-rr-div-le}{\vnbstmt{\begin{array}{@{}l@{}} \forall a, b, c. \\ \quad ((a \in \mathbb{R}) \wedge ((b \in \mathbb{R}) \wedge ((c \in \mathbb{R}) \wedge ((0 < c) \wedge ((c \cdot a) \leq b))))) \Rightarrow  \\ \quad \quad a \leq (\operatorname{recip}(c) \cdot b) \end{array}}}
\vnbprop{mink-rr-le-div}{\vnbstmt{\begin{array}{@{}l@{}} \forall a, b, c. \\ \quad ((a \in \mathbb{R}) \wedge ((b \in \mathbb{R}) \wedge ((c \in \mathbb{R}) \wedge ((0 < c) \wedge (a \leq (\operatorname{recip}(c) \cdot b)))))) \Rightarrow  \\ \quad \quad (c \cdot a) \leq b \end{array}}}
\vnbprop{mink-rr-div-ge}{\vnbstmt{\begin{array}{@{}l@{}} \forall a, b, c. \\ \quad ((a \in \mathbb{R}) \wedge ((b \in \mathbb{R}) \wedge ((c \in \mathbb{R}) \wedge ((0 < c) \wedge (b \leq (c \cdot a)))))) \Rightarrow  \\ \quad \quad (\operatorname{recip}(c) \cdot b) \leq a \end{array}}}
\vnbprop{mink-rr-ge-div}{\vnbstmt{\begin{array}{@{}l@{}} \forall a, b, c. \\ \quad ((a \in \mathbb{R}) \wedge ((b \in \mathbb{R}) \wedge ((c \in \mathbb{R}) \wedge ((0 < c) \wedge ((\operatorname{recip}(c) \cdot b) \leq a))))) \Rightarrow  \\ \quad \quad b \leq (c \cdot a) \end{array}}}
\vnbprop{rmod-act-recip-cancel}{\vnbstmt{\begin{array}{@{}l@{}} \forall m, c, x. \\ \quad \operatorname{is\text{-}module}(m) \wedge \\ \quad \operatorname{scal}(m) = \operatorname{normed\text{-}field\text{-}as\text{-}commutative\text{-}ring}(\operatorname{rr\text{-}normed\text{-}field}) \wedge \\ \quad c \in \mathbb{R} \wedge \\ \quad 0 < c \wedge \\ \quad x \in \operatorname{vec}(m) \Rightarrow \\ \quad \quad \operatorname{act}(m)(\operatorname{recip}(c), \operatorname{act}(m)(c, x)) = x \end{array}}}
\vnbprop{mink-set-scale}{\vnbstmt{\begin{array}{@{}l@{}} \forall m, s, x, c, r. \\ \quad \operatorname{is\text{-}module}(m) \wedge \\ \quad \operatorname{scal}(m) = \operatorname{normed\text{-}field\text{-}as\text{-}commutative\text{-}ring}(\operatorname{rr\text{-}normed\text{-}field}) \wedge \\ \quad \operatorname{is\text{-}convex\text{-}subset}(m, s) \wedge \\ \quad \operatorname{is\text{-}absorbing\text{-}subset}(m, s) \wedge \\ \quad \operatorname{vzero}(m) \in s \wedge \\ \quad r \in \operatorname{mink\text{-}set}(m, s, x) \wedge \\ \quad c \in \mathbb{R} \wedge \\ \quad 0 < c \Rightarrow \\ \quad \quad (c \cdot r) \in \operatorname{mink\text{-}set}(m, s, \operatorname{act}(m)(c, x)) \end{array}}}
\vnbprop{mink-set-neg}{\vnbstmt{\begin{array}{@{}l@{}} \forall m, s, x, r. \\ \quad \operatorname{is\text{-}module}(m) \wedge \\ \quad \operatorname{scal}(m) = \operatorname{normed\text{-}field\text{-}as\text{-}commutative\text{-}ring}(\operatorname{rr\text{-}normed\text{-}field}) \wedge \\ \quad \operatorname{is\text{-}convex\text{-}subset}(m, s) \wedge \\ \quad \operatorname{is\text{-}symmetric\text{-}subset}(m, s) \wedge \\ \quad \operatorname{is\text{-}absorbing\text{-}subset}(m, s) \wedge \\ \quad r \in \operatorname{mink\text{-}set}(m, s, x) \Rightarrow \\ \quad \quad r \in \operatorname{mink\text{-}set}(m, s, \operatorname{act}(m)(-1, x)) \end{array}}}
\vnbprop{mink-set-zero}{\vnbstmt{\begin{array}{@{}l@{}} \forall m, s, r. \\ \quad \operatorname{is\text{-}module}(m) \wedge \\ \quad \operatorname{scal}(m) = \operatorname{normed\text{-}field\text{-}as\text{-}commutative\text{-}ring}(\operatorname{rr\text{-}normed\text{-}field}) \wedge \\ \quad \operatorname{is\text{-}convex\text{-}subset}(m, s) \wedge \\ \quad \operatorname{is\text{-}absorbing\text{-}subset}(m, s) \wedge \\ \quad \operatorname{vzero}(m) \in s \wedge \\ \quad r \in \mathbb{R} \wedge \\ \quad 0 < r \Rightarrow \\ \quad \quad r \in \operatorname{mink\text{-}set}(m, s, \operatorname{vzero}(m)) \end{array}}}
\vnbprop{mink-set-sum}{\vnbstmt{\begin{array}{@{}l@{}} \forall m, a, x, y, r, s. \\ \quad \operatorname{is\text{-}module}(m) \wedge \\ \quad \operatorname{scal}(m) = \operatorname{normed\text{-}field\text{-}as\text{-}commutative\text{-}ring}(\operatorname{rr\text{-}normed\text{-}field}) \wedge \\ \quad \operatorname{is\text{-}convex\text{-}subset}(m, a) \wedge \\ \quad \operatorname{is\text{-}absorbing\text{-}subset}(m, a) \wedge \\ \quad \operatorname{vzero}(m) \in a \wedge \\ \quad r \in \operatorname{mink\text{-}set}(m, a, x) \wedge \\ \quad s \in \operatorname{mink\text{-}set}(m, a, y) \Rightarrow \\ \quad \quad (r + s) \in \operatorname{mink\text{-}set}(m, a, \operatorname{vadd}(m)(x, y)) \end{array}}}
\vnbprop{mink-inf-nonneg}{\vnbstmt{\begin{array}{@{}l@{}} \forall m, s, z. \\ \quad \operatorname{is\text{-}module}(m) \wedge \\ \quad \operatorname{scal}(m) = \operatorname{normed\text{-}field\text{-}as\text{-}commutative\text{-}ring}(\operatorname{rr\text{-}normed\text{-}field}) \wedge \\ \quad \operatorname{is\text{-}convex\text{-}subset}(m, s) \wedge \\ \quad \operatorname{is\text{-}absorbing\text{-}subset}(m, s) \wedge \\ \quad \operatorname{vzero}(m) \in s \wedge \\ \quad z \in \operatorname{vec}(m) \Rightarrow \\ \quad \quad 0 \leq \operatorname{inf}(\operatorname{mink\text{-}set}(m, s, z)) \end{array}}}
\vnbprop{mink-inf-subadditive}{\vnbstmt{\begin{array}{@{}l@{}} \forall m, s, x, y. \\ \quad \operatorname{is\text{-}module}(m) \wedge \\ \quad \operatorname{scal}(m) = \operatorname{normed\text{-}field\text{-}as\text{-}commutative\text{-}ring}(\operatorname{rr\text{-}normed\text{-}field}) \wedge \\ \quad \operatorname{is\text{-}convex\text{-}subset}(m, s) \wedge \\ \quad \operatorname{is\text{-}absorbing\text{-}subset}(m, s) \wedge \\ \quad \operatorname{vzero}(m) \in s \wedge \\ \quad x \in \operatorname{vec}(m) \wedge \\ \quad y \in \operatorname{vec}(m) \Rightarrow \\ \quad \quad \operatorname{inf}(\operatorname{mink\text{-}set}(m, s, \operatorname{vadd}(m)(x, y))) \leq (\operatorname{inf}(\operatorname{mink\text{-}set}(m, s, x)) + \operatorname{inf}(\operatorname{mink\text{-}set}(m, s, y))) \end{array}}}
\vnbprop{mink-inf-scale}{\vnbstmt{\begin{array}{@{}l@{}} \forall m, s, x, c. \\ \quad \operatorname{is\text{-}module}(m) \wedge \\ \quad \operatorname{scal}(m) = \operatorname{normed\text{-}field\text{-}as\text{-}commutative\text{-}ring}(\operatorname{rr\text{-}normed\text{-}field}) \wedge \\ \quad \operatorname{is\text{-}convex\text{-}subset}(m, s) \wedge \\ \quad \operatorname{is\text{-}absorbing\text{-}subset}(m, s) \wedge \\ \quad \operatorname{vzero}(m) \in s \wedge \\ \quad x \in \operatorname{vec}(m) \wedge \\ \quad c \in \mathbb{R} \wedge \\ \quad 0 < c \Rightarrow \\ \quad \quad \operatorname{inf}(\operatorname{mink\text{-}set}(m, s, \operatorname{act}(m)(c, x))) = (c \cdot \operatorname{inf}(\operatorname{mink\text{-}set}(m, s, x))) \end{array}}}
\vnbprop{mink-inf-zero}{\vnbstmt{\begin{array}{@{}l@{}} \forall m, s. \\ \quad \operatorname{is\text{-}module}(m) \wedge \\ \quad \operatorname{scal}(m) = \operatorname{normed\text{-}field\text{-}as\text{-}commutative\text{-}ring}(\operatorname{rr\text{-}normed\text{-}field}) \wedge \\ \quad \operatorname{is\text{-}convex\text{-}subset}(m, s) \wedge \\ \quad \operatorname{is\text{-}absorbing\text{-}subset}(m, s) \wedge \\ \quad \operatorname{vzero}(m) \in s \Rightarrow \\ \quad \quad \operatorname{inf}(\operatorname{mink\text{-}set}(m, s, \operatorname{vzero}(m))) = 0 \end{array}}}
\vnbprop{mink-inf-neg-le}{\vnbstmt{\begin{array}{@{}l@{}} \forall m, s, x. \\ \quad \operatorname{is\text{-}module}(m) \wedge \\ \quad \operatorname{scal}(m) = \operatorname{normed\text{-}field\text{-}as\text{-}commutative\text{-}ring}(\operatorname{rr\text{-}normed\text{-}field}) \wedge \\ \quad \operatorname{is\text{-}convex\text{-}subset}(m, s) \wedge \\ \quad \operatorname{is\text{-}symmetric\text{-}subset}(m, s) \wedge \\ \quad \operatorname{is\text{-}absorbing\text{-}subset}(m, s) \wedge \\ \quad x \in \operatorname{vec}(m) \Rightarrow \\ \quad \quad \operatorname{inf}(\operatorname{mink\text{-}set}(m, s, \operatorname{act}(m)(-1, x))) \leq \operatorname{inf}(\operatorname{mink\text{-}set}(m, s, x)) \end{array}}}
\vnbprop{mink-inf-neg}{\vnbstmt{\begin{array}{@{}l@{}} \forall m, s, x. \\ \quad \operatorname{is\text{-}module}(m) \wedge \\ \quad \operatorname{scal}(m) = \operatorname{normed\text{-}field\text{-}as\text{-}commutative\text{-}ring}(\operatorname{rr\text{-}normed\text{-}field}) \wedge \\ \quad \operatorname{is\text{-}convex\text{-}subset}(m, s) \wedge \\ \quad \operatorname{is\text{-}symmetric\text{-}subset}(m, s) \wedge \\ \quad \operatorname{is\text{-}absorbing\text{-}subset}(m, s) \wedge \\ \quad x \in \operatorname{vec}(m) \Rightarrow \\ \quad \quad \operatorname{inf}(\operatorname{mink\text{-}set}(m, s, \operatorname{act}(m)(-1, x))) = \operatorname{inf}(\operatorname{mink\text{-}set}(m, s, x)) \end{array}}}
\vnbprop{mink-inf-homogeneous}{\vnbstmt{\begin{array}{@{}l@{}} \forall m, s, x, \ell. \\ \quad \operatorname{is\text{-}module}(m) \wedge \\ \quad \operatorname{scal}(m) = \operatorname{normed\text{-}field\text{-}as\text{-}commutative\text{-}ring}(\operatorname{rr\text{-}normed\text{-}field}) \wedge \\ \quad \operatorname{is\text{-}convex\text{-}subset}(m, s) \wedge \\ \quad \operatorname{is\text{-}symmetric\text{-}subset}(m, s) \wedge \\ \quad \operatorname{is\text{-}absorbing\text{-}subset}(m, s) \wedge \\ \quad x \in \operatorname{vec}(m) \wedge \\ \quad \ell \in \mathbb{R} \Rightarrow \\ \quad \quad \operatorname{inf}(\operatorname{mink\text{-}set}(m, s, \operatorname{act}(m)(\ell, x))) = (\lvert \ell \rvert \cdot \operatorname{inf}(\operatorname{mink\text{-}set}(m, s, x))) \end{array}}}
\vnbprop{minkowski-functional-is-seminorm}{\vnbstmt{\begin{array}{@{}l@{}} \forall m, s. \\ \quad \operatorname{is\text{-}module}(m) \wedge \\ \quad \operatorname{scal}(m) = \operatorname{normed\text{-}field\text{-}as\text{-}commutative\text{-}ring}(\operatorname{rr\text{-}normed\text{-}field}) \wedge \\ \quad \operatorname{is\text{-}convex\text{-}subset}(m, s) \wedge \\ \quad \operatorname{is\text{-}symmetric\text{-}subset}(m, s) \wedge \\ \quad \operatorname{is\text{-}absorbing\text{-}subset}(m, s) \Rightarrow \\ \quad \quad \operatorname{is\text{-}seminorm}(\operatorname{rr\text{-}normed\text{-}field}, m, p_{s}) \end{array}}}

\vnbbookfile{structure-library/seminorm-extension}{structure-library/seminorm-extension}
\noindent{\small sxt-vocab.scm -- a seminorm on a subspace, and the extension of topological-vector-space.pdf Lemma 4.20, Eq. (113).  Vocabulary only.  Agent DUAL-FIX, 2026-10-26 (NEW DEFINITIONS, the user's to confirm). DESTINATION: structure-library/seminorm-extension.scm.  Library: analysis. Load window: after the INF vocabulary ...}\par
\par\medskip\noindent\textbf{Definition} (\texttt{is-seminorm-on}).\;For all $m$, $s$ and $p$, 
\begin{equation*}
\begin{array}{@{}l@{}} \operatorname{is\text{-}seminorm\text{-}on}(m, s, p) \\ \Leftrightarrow\; p \in (s \to \mathbb{R}) \wedge \\ \forall \operatorname{snx\_} \in s.\; (0 \leq p(\operatorname{snx\_})) \wedge \\ \forall \operatorname{snx\_}, \operatorname{sny\_} \in s.\; (p(\operatorname{vadd}(m)(\operatorname{snx\_}, \operatorname{sny\_})) \leq +(p(\operatorname{snx\_}), p(\operatorname{sny\_}))) \wedge \\ \forall \operatorname{snl\_} \in \mathbb{R}, \operatorname{snx\_} \in s.\; (p(\operatorname{act}(m)(\operatorname{snl\_}, \operatorname{snx\_})) = \cdot(\lvert \operatorname{snl\_} \rvert, p(\operatorname{snx\_}))) \end{array}
\end{equation*}

\par\medskip\noindent\textbf{Definition} (\texttt{seminorm-ext-set}).\;$\operatorname{seminorm\text{-}ext\text{-}set}(m, s, p, q, y) :=$
\begin{equation*}
\{\operatorname{sxt\_} \in \mathbb{R} : \exists \operatorname{sev\_} \in s, \operatorname{sew\_} \in \operatorname{vec}(m).\; ((y = \operatorname{vadd}(m)(\operatorname{sev\_}, \operatorname{sew\_})) \wedge (\operatorname{sxt\_} = (p(\operatorname{sev\_}) + q(\operatorname{sew\_}))))\}
\end{equation*}

\par\medskip\noindent\textbf{Definition} (\texttt{seminorm-ext}).\;$\operatorname{seminorm\text{-}ext}(m, s, p, q) :=$
\begin{equation*}
(\lambda \operatorname{sxy\_} \in \operatorname{vec}(m).\; \operatorname{inf}(\operatorname{seminorm\text{-}ext\text{-}set}(m, s, p, q, \operatorname{sxy\_})))
\end{equation*}


\vnbbookfile{theorem-library/tvs-seminorm-extension}{theorem-library/tvs-seminorm-extension}
\noindent{\small sxt.scm -- topological-vector-space.pdf Lemma 4.20: a seminorm p on a subspace E of F, dominated by a seminorm q of F, extends to a seminorm p\textasciitilde{} of F with p\textasciitilde{} $<$= q, namely p\textasciitilde{}(y) = inf \{ p(v) + q(w) : v in E, w in F, y = v + w \}. Agent DUAL-FIX, 2026-10-26.  DESTINATION: theorem-library/tvs-seminorm-extension.scm.  ...}\par
\vnbprop{rmod-vadd-assoc}{\vnbstmt{\begin{array}{@{}l@{}} \forall m, x, y, z. \\ \quad (\operatorname{is\text{-}module}(m) \wedge ((x \in \operatorname{vec}(m)) \wedge ((y \in \operatorname{vec}(m)) \wedge (z \in \operatorname{vec}(m))))) \Rightarrow  \\ \quad \quad \operatorname{vadd}(m)(\operatorname{vadd}(m)(x, y), z) = \operatorname{vadd}(m)(x, \operatorname{vadd}(m)(y, z)) \end{array}}}
\vnbprop{rmod-vadd-comm}{\vnbstmt{\begin{array}{@{}l@{}} \forall m, x, y. \\ \quad (\operatorname{is\text{-}module}(m) \wedge ((x \in \operatorname{vec}(m)) \wedge (y \in \operatorname{vec}(m)))) \Rightarrow  \\ \quad \quad \operatorname{vadd}(m)(x, y) = \operatorname{vadd}(m)(y, x) \end{array}}}
\vnbprop{rmod-vzero-right}{\vnbstmt{\forall m, x.\; ((\operatorname{is\text{-}module}(m) \wedge (x \in \operatorname{vec}(m))) \Rightarrow (\operatorname{vadd}(m)(x, \operatorname{vzero}(m)) = x))}}
\vnbprop{rmod-vneg-left}{\vnbstmt{\begin{array}{@{}l@{}} \forall m, x. \\ \quad (\operatorname{is\text{-}module}(m) \wedge (x \in \operatorname{vec}(m))) \Rightarrow (\operatorname{vadd}(m)(\operatorname{vneg}(m)(x), x) = \operatorname{vzero}(m)) \end{array}}}
\vnbprop{rmod-vadd-swap}{\vnbstmt{\begin{array}{@{}l@{}} \forall m, a, b, c, d. \\ \quad (\operatorname{is\text{-}module}(m) \wedge ((a \in \operatorname{vec}(m)) \wedge ((b \in \operatorname{vec}(m)) \wedge ((c \in \operatorname{vec}(m)) \wedge (d \in \operatorname{vec}(m)))))) \Rightarrow  \\ \quad \quad \operatorname{vadd}(m)(\operatorname{vadd}(m)(a, b), \operatorname{vadd}(m)(c, d)) = \operatorname{vadd}(m)(\operatorname{vadd}(m)(a, c), \operatorname{vadd}(m)(b, d)) \end{array}}}
\vnbprop{rmod-vneg-cancel}{\vnbstmt{\begin{array}{@{}l@{}} \forall m, v, w. \\ \quad (\operatorname{is\text{-}module}(m) \wedge ((v \in \operatorname{vec}(m)) \wedge (w \in \operatorname{vec}(m)))) \Rightarrow  \\ \quad \quad \operatorname{vadd}(m)(\operatorname{vneg}(m)(v), \operatorname{vadd}(m)(v, w)) = w \end{array}}}
\vnbprop{rmod-act-recip-cancel-nz}{\vnbstmt{\begin{array}{@{}l@{}} \forall m, c, x. \\ \quad \operatorname{is\text{-}module}(m) \wedge \\ \quad \operatorname{scal}(m) = \operatorname{normed\text{-}field\text{-}as\text{-}commutative\text{-}ring}(\operatorname{rr\text{-}normed\text{-}field}) \wedge \\ \quad c \in \mathbb{R} \wedge \\ \quad \neg (c = 0) \wedge \\ \quad x \in \operatorname{vec}(m) \Rightarrow \\ \quad \quad \operatorname{act}(m)(\operatorname{recip}(c), \operatorname{act}(m)(c, x)) = x \end{array}}}
\vnbprop{sxt-q-nonneg}{\vnbstmt{\begin{array}{@{}l@{}} \forall m, q, x. \\ \quad (\operatorname{is\text{-}seminorm}(\operatorname{rr\text{-}normed\text{-}field}, m, q) \wedge (x \in \operatorname{vec}(m))) \Rightarrow (0 \leq q(x)) \end{array}}}
\vnbprop{sxt-q-subadd}{\vnbstmt{\begin{array}{@{}l@{}} \forall m, q, x, y. \\ \quad (\operatorname{is\text{-}seminorm}(\operatorname{rr\text{-}normed\text{-}field}, m, q) \wedge ((x \in \operatorname{vec}(m)) \wedge (y \in \operatorname{vec}(m)))) \Rightarrow  \\ \quad \quad q(\operatorname{vadd}(m)(x, y)) \leq (q(x) + q(y)) \end{array}}}
\vnbprop{sxt-q-homog}{\vnbstmt{\begin{array}{@{}l@{}} \forall m, q, \ell, x. \\ \quad \operatorname{is\text{-}module}(m) \wedge \\ \quad \operatorname{scal}(m) = \operatorname{normed\text{-}field\text{-}as\text{-}commutative\text{-}ring}(\operatorname{rr\text{-}normed\text{-}field}) \wedge \\ \quad \operatorname{is\text{-}seminorm}(\operatorname{rr\text{-}normed\text{-}field}, m, q) \wedge \\ \quad \ell \in \mathbb{R} \wedge \\ \quad x \in \operatorname{vec}(m) \Rightarrow \\ \quad \quad q(\operatorname{act}(m)(\ell, x)) = (\lvert \ell \rvert \cdot q(x)) \end{array}}}
\vnbprop{sxt-q-zero}{\vnbstmt{\begin{array}{@{}l@{}} \forall m, q. \\ \quad \operatorname{is\text{-}module}(m) \wedge \\ \quad \operatorname{scal}(m) = \operatorname{normed\text{-}field\text{-}as\text{-}commutative\text{-}ring}(\operatorname{rr\text{-}normed\text{-}field}) \wedge \\ \quad \operatorname{is\text{-}seminorm}(\operatorname{rr\text{-}normed\text{-}field}, m, q) \Rightarrow \\ \quad \quad q(\operatorname{vzero}(m)) = 0 \end{array}}}
\vnbprop{sxt-p-nonneg}{\vnbstmt{\forall m, s, p, x.\; ((\operatorname{is\text{-}seminorm\text{-}on}(m, s, p) \wedge (x \in s)) \Rightarrow (0 \leq p(x)))}}
\vnbprop{sxt-p-subadd}{\vnbstmt{\begin{array}{@{}l@{}} \forall m, s, p, x, y. \\ \quad (\operatorname{is\text{-}seminorm\text{-}on}(m, s, p) \wedge ((x \in s) \wedge (y \in s))) \Rightarrow  \\ \quad \quad p(\operatorname{vadd}(m)(x, y)) \leq (p(x) + p(y)) \end{array}}}
\vnbprop{sxt-p-homog}{\vnbstmt{\begin{array}{@{}l@{}} \forall m, s, p, \ell, x. \\ \quad (\operatorname{is\text{-}seminorm\text{-}on}(m, s, p) \wedge ((\ell \in \mathbb{R}) \wedge (x \in s))) \Rightarrow  \\ \quad \quad p(\operatorname{act}(m)(\ell, x)) = (\lvert \ell \rvert \cdot p(x)) \end{array}}}
\vnbprop{sxt-p-zero}{\vnbstmt{\begin{array}{@{}l@{}} \forall m, s, p. \\ \quad \operatorname{is\text{-}module}(m) \wedge \\ \quad \operatorname{scal}(m) = \operatorname{normed\text{-}field\text{-}as\text{-}commutative\text{-}ring}(\operatorname{rr\text{-}normed\text{-}field}) \wedge \\ \quad \operatorname{is\text{-}submodule}(m, s) \wedge \\ \quad \operatorname{is\text{-}seminorm\text{-}on}(m, s, p) \Rightarrow \\ \quad \quad p(\operatorname{vzero}(m)) = 0 \end{array}}}
\vnbprop{sxt-set-unfold}{\vnbstmt{\begin{array}{@{}l@{}} \forall m, s, p, q, y. \\ \quad (\operatorname{seminorm\text{-}ext\text{-}set}(m, s, p, q, y) \simeq \{t \in \mathbb{R} : \exists v \in s, w \in \operatorname{vec}(m).\; ((y = \operatorname{vadd}(m)(v, w)) \wedge (t = (p(v) + q(w))))\}) \end{array}}}
\vnbprop{sxt-set-intro}{\vnbstmt{\begin{array}{@{}l@{}} \forall m, s, p, q, y, v, w. \\ \quad ((p \in (s \to \mathbb{R})) \wedge ((q \in (\operatorname{vec}(m) \to \mathbb{R})) \wedge ((v \in s) \wedge ((w \in \operatorname{vec}(m)) \wedge (y = \operatorname{vadd}(m)(v, w)))))) \Rightarrow  \\ \quad \quad (p(v) + q(w)) \in \operatorname{seminorm\text{-}ext\text{-}set}(m, s, p, q, y) \end{array}}}
\vnbprop{sxt-set-elim}{\vnbstmt{\begin{array}{@{}l@{}} \forall m, s, p, q, y.; \forall t \in \operatorname{seminorm\text{-}ext\text{-}set}(m, s, p, q, y). \\ \quad ((t \in \mathbb{R}) \wedge \exists v \in s, w \in \operatorname{vec}(m).\; ((y = \operatorname{vadd}(m)(v, w)) \wedge (t = (p(v) + q(w))))) \end{array}}}
\vnbprop{sxt-set-glb}{\vnbstmt{\begin{array}{@{}l@{}} \forall m, s, p, q, z. \\ \quad \operatorname{is\text{-}module}(m) \wedge \\ \quad \operatorname{scal}(m) = \operatorname{normed\text{-}field\text{-}as\text{-}commutative\text{-}ring}(\operatorname{rr\text{-}normed\text{-}field}) \wedge \\ \quad \operatorname{is\text{-}submodule}(m, s) \wedge \\ \quad \operatorname{is\text{-}seminorm\text{-}on}(m, s, p) \wedge \\ \quad \operatorname{is\text{-}seminorm}(\operatorname{rr\text{-}normed\text{-}field}, m, q) \wedge \\ \quad \forall x \in s.\; (p(x) \leq q(x)) \wedge \\ \quad z \in \operatorname{vec}(m) \Rightarrow \\ \quad \quad \operatorname{inf}(\operatorname{seminorm\text{-}ext\text{-}set}(m, s, p, q, z)) \in \mathbb{R} \wedge \\ \quad \quad \forall w \in \operatorname{seminorm\text{-}ext\text{-}set}(m, s, p, q, z). \\ \quad \quad \quad \operatorname{inf}(\operatorname{seminorm\text{-}ext\text{-}set}(m, s, p, q, z)) \leq w \wedge \\ \quad \quad \forall b \in \mathbb{R}. \\ \quad \quad \quad \forall w \in \operatorname{seminorm\text{-}ext\text{-}set}(m, s, p, q, z).\; (b \leq w) \Rightarrow  \\ \quad \quad \quad \quad b \leq \operatorname{inf}(\operatorname{seminorm\text{-}ext\text{-}set}(m, s, p, q, z)) \end{array}}}
\vnbprop{sxt-value}{\vnbstmt{\begin{array}{@{}l@{}} \forall m, s, p, q.; \forall z \in \operatorname{vec}(m). \\ \quad (\operatorname{seminorm\text{-}ext}(m, s, p, q)(z) \simeq \operatorname{inf}(\operatorname{seminorm\text{-}ext\text{-}set}(m, s, p, q, z))) \end{array}}}
\vnbprop{sxt-inf-le-q}{\vnbstmt{\begin{array}{@{}l@{}} \forall m, s, p, q, y. \\ \quad \operatorname{is\text{-}module}(m) \wedge \\ \quad \operatorname{scal}(m) = \operatorname{normed\text{-}field\text{-}as\text{-}commutative\text{-}ring}(\operatorname{rr\text{-}normed\text{-}field}) \wedge \\ \quad \operatorname{is\text{-}submodule}(m, s) \wedge \\ \quad \operatorname{is\text{-}seminorm\text{-}on}(m, s, p) \wedge \\ \quad \operatorname{is\text{-}seminorm}(\operatorname{rr\text{-}normed\text{-}field}, m, q) \wedge \\ \quad \forall x \in s.\; (p(x) \leq q(x)) \wedge \\ \quad y \in \operatorname{vec}(m) \Rightarrow \\ \quad \quad \operatorname{inf}(\operatorname{seminorm\text{-}ext\text{-}set}(m, s, p, q, y)) \leq q(y) \end{array}}}
\vnbprop{sxt-inf-extends}{\vnbstmt{\begin{array}{@{}l@{}} \forall m, s, p, q, x. \\ \quad \operatorname{is\text{-}module}(m) \wedge \\ \quad \operatorname{scal}(m) = \operatorname{normed\text{-}field\text{-}as\text{-}commutative\text{-}ring}(\operatorname{rr\text{-}normed\text{-}field}) \wedge \\ \quad \operatorname{is\text{-}submodule}(m, s) \wedge \\ \quad \operatorname{is\text{-}seminorm\text{-}on}(m, s, p) \wedge \\ \quad \operatorname{is\text{-}seminorm}(\operatorname{rr\text{-}normed\text{-}field}, m, q) \wedge \\ \quad \forall \mathit{dmx} \in s.\; (p(\mathit{dmx}) \leq q(\mathit{dmx})) \wedge \\ \quad x \in s \Rightarrow \\ \quad \quad \operatorname{inf}(\operatorname{seminorm\text{-}ext\text{-}set}(m, s, p, q, x)) = p(x) \end{array}}}
\vnbprop{sxt-inf-nonneg}{\vnbstmt{\begin{array}{@{}l@{}} \forall m, s, p, q, z. \\ \quad \operatorname{is\text{-}module}(m) \wedge \\ \quad \operatorname{scal}(m) = \operatorname{normed\text{-}field\text{-}as\text{-}commutative\text{-}ring}(\operatorname{rr\text{-}normed\text{-}field}) \wedge \\ \quad \operatorname{is\text{-}submodule}(m, s) \wedge \\ \quad \operatorname{is\text{-}seminorm\text{-}on}(m, s, p) \wedge \\ \quad \operatorname{is\text{-}seminorm}(\operatorname{rr\text{-}normed\text{-}field}, m, q) \wedge \\ \quad \forall x \in s.\; (p(x) \leq q(x)) \wedge \\ \quad z \in \operatorname{vec}(m) \Rightarrow \\ \quad \quad 0 \leq \operatorname{inf}(\operatorname{seminorm\text{-}ext\text{-}set}(m, s, p, q, z)) \end{array}}}
\vnbprop{sxt-elem-sum}{\vnbstmt{\begin{array}{@{}l@{}} \forall m, s, p, q, x, y, a, b. \\ \quad \operatorname{is\text{-}module}(m) \wedge \\ \quad \operatorname{scal}(m) = \operatorname{normed\text{-}field\text{-}as\text{-}commutative\text{-}ring}(\operatorname{rr\text{-}normed\text{-}field}) \wedge \\ \quad \operatorname{is\text{-}submodule}(m, s) \wedge \\ \quad \operatorname{is\text{-}seminorm\text{-}on}(m, s, p) \wedge \\ \quad \operatorname{is\text{-}seminorm}(\operatorname{rr\text{-}normed\text{-}field}, m, q) \wedge \\ \quad \forall \mathit{dmx} \in s.\; (p(\mathit{dmx}) \leq q(\mathit{dmx})) \wedge \\ \quad x \in \operatorname{vec}(m) \wedge \\ \quad y \in \operatorname{vec}(m) \wedge \\ \quad a \in \operatorname{seminorm\text{-}ext\text{-}set}(m, s, p, q, x) \wedge \\ \quad b \in \operatorname{seminorm\text{-}ext\text{-}set}(m, s, p, q, y) \Rightarrow \\ \quad \quad \operatorname{inf}(\operatorname{seminorm\text{-}ext\text{-}set}(m, s, p, q, \operatorname{vadd}(m)(x, y))) \leq (a + b) \end{array}}}
\vnbprop{sxt-inf-subadditive}{\vnbstmt{\begin{array}{@{}l@{}} \forall m, s, p, q, x, y. \\ \quad \operatorname{is\text{-}module}(m) \wedge \\ \quad \operatorname{scal}(m) = \operatorname{normed\text{-}field\text{-}as\text{-}commutative\text{-}ring}(\operatorname{rr\text{-}normed\text{-}field}) \wedge \\ \quad \operatorname{is\text{-}submodule}(m, s) \wedge \\ \quad \operatorname{is\text{-}seminorm\text{-}on}(m, s, p) \wedge \\ \quad \operatorname{is\text{-}seminorm}(\operatorname{rr\text{-}normed\text{-}field}, m, q) \wedge \\ \quad \forall \mathit{dmx} \in s.\; (p(\mathit{dmx}) \leq q(\mathit{dmx})) \wedge \\ \quad x \in \operatorname{vec}(m) \wedge \\ \quad y \in \operatorname{vec}(m) \Rightarrow \\ \quad \quad \operatorname{inf}(\operatorname{seminorm\text{-}ext\text{-}set}(m, s, p, q, \operatorname{vadd}(m)(x, y))) \\ \quad \quad \leq\; \operatorname{inf}(\operatorname{seminorm\text{-}ext\text{-}set}(m, s, p, q, x)) + \operatorname{inf}(\operatorname{seminorm\text{-}ext\text{-}set}(m, s, p, q, y)) \end{array}}}
\vnbprop{sxt-elem-scale}{\vnbstmt{\begin{array}{@{}l@{}} \forall m, s, p, q, x, \ell, t. \\ \quad \operatorname{is\text{-}module}(m) \wedge \\ \quad \operatorname{scal}(m) = \operatorname{normed\text{-}field\text{-}as\text{-}commutative\text{-}ring}(\operatorname{rr\text{-}normed\text{-}field}) \wedge \\ \quad \operatorname{is\text{-}submodule}(m, s) \wedge \\ \quad \operatorname{is\text{-}seminorm\text{-}on}(m, s, p) \wedge \\ \quad \operatorname{is\text{-}seminorm}(\operatorname{rr\text{-}normed\text{-}field}, m, q) \wedge \\ \quad \forall \mathit{dmx} \in s.\; (p(\mathit{dmx}) \leq q(\mathit{dmx})) \wedge \\ \quad x \in \operatorname{vec}(m) \wedge \\ \quad \ell \in \mathbb{R} \wedge \\ \quad t \in \operatorname{seminorm\text{-}ext\text{-}set}(m, s, p, q, x) \Rightarrow \\ \quad \quad (\lvert \ell \rvert \cdot t) \in \operatorname{seminorm\text{-}ext\text{-}set}(m, s, p, q, \operatorname{act}(m)(\ell, x)) \end{array}}}
\vnbprop{sxt-inf-homog-le}{\vnbstmt{\begin{array}{@{}l@{}} \forall m, s, p, q, x, \ell. \\ \quad \operatorname{is\text{-}module}(m) \wedge \\ \quad \operatorname{scal}(m) = \operatorname{normed\text{-}field\text{-}as\text{-}commutative\text{-}ring}(\operatorname{rr\text{-}normed\text{-}field}) \wedge \\ \quad \operatorname{is\text{-}submodule}(m, s) \wedge \\ \quad \operatorname{is\text{-}seminorm\text{-}on}(m, s, p) \wedge \\ \quad \operatorname{is\text{-}seminorm}(\operatorname{rr\text{-}normed\text{-}field}, m, q) \wedge \\ \quad \forall \mathit{dmx} \in s.\; (p(\mathit{dmx}) \leq q(\mathit{dmx})) \wedge \\ \quad x \in \operatorname{vec}(m) \wedge \\ \quad \ell \in \mathbb{R} \Rightarrow \\ \quad \quad \operatorname{inf}(\operatorname{seminorm\text{-}ext\text{-}set}(m, s, p, q, \operatorname{act}(m)(\ell, x))) \\ \quad \quad \leq\; \lvert \ell \rvert \cdot \operatorname{inf}(\operatorname{seminorm\text{-}ext\text{-}set}(m, s, p, q, x)) \end{array}}}
\vnbprop{sxt-inf-homogeneous}{\vnbstmt{\begin{array}{@{}l@{}} \forall m, s, p, q, x, \ell. \\ \quad \operatorname{is\text{-}module}(m) \wedge \\ \quad \operatorname{scal}(m) = \operatorname{normed\text{-}field\text{-}as\text{-}commutative\text{-}ring}(\operatorname{rr\text{-}normed\text{-}field}) \wedge \\ \quad \operatorname{is\text{-}submodule}(m, s) \wedge \\ \quad \operatorname{is\text{-}seminorm\text{-}on}(m, s, p) \wedge \\ \quad \operatorname{is\text{-}seminorm}(\operatorname{rr\text{-}normed\text{-}field}, m, q) \wedge \\ \quad \forall \mathit{dmx} \in s.\; (p(\mathit{dmx}) \leq q(\mathit{dmx})) \wedge \\ \quad x \in \operatorname{vec}(m) \wedge \\ \quad \ell \in \mathbb{R} \Rightarrow \\ \quad \quad \operatorname{inf}(\operatorname{seminorm\text{-}ext\text{-}set}(m, s, p, q, \operatorname{act}(m)(\ell, x))) \\ \quad \quad =\; \lvert \ell \rvert \cdot \operatorname{inf}(\operatorname{seminorm\text{-}ext\text{-}set}(m, s, p, q, x)) \end{array}}}
\vnbprop{seminorm-extension}{\vnbstmt{\begin{array}{@{}l@{}} \forall m, s, p, q. \\ \quad \operatorname{is\text{-}module}(m) \wedge \\ \quad \operatorname{scal}(m) = \operatorname{normed\text{-}field\text{-}as\text{-}commutative\text{-}ring}(\operatorname{rr\text{-}normed\text{-}field}) \wedge \\ \quad \operatorname{is\text{-}submodule}(m, s) \wedge \\ \quad \operatorname{is\text{-}seminorm\text{-}on}(m, s, p) \wedge \\ \quad \operatorname{is\text{-}seminorm}(\operatorname{rr\text{-}normed\text{-}field}, m, q) \wedge \\ \quad \forall x \in s.\; (p(x) \leq q(x)) \Rightarrow \\ \quad \quad \operatorname{is\text{-}seminorm}(\operatorname{rr\text{-}normed\text{-}field}, m, \operatorname{seminorm\text{-}ext}(m, s, p, q)) \wedge \\ \quad \quad \operatorname{extends\text{-}on}(s, \operatorname{seminorm\text{-}ext}(m, s, p, q), p) \wedge \\ \quad \quad \forall y \in \operatorname{vec}(m).\; (\operatorname{seminorm\text{-}ext}(m, s, p, q)(y) \leq q(y)) \end{array}}}

\vnbbookfile{structure-library/tvs-vocab}{structure-library/tvs-vocab}
\noindent{\small tb-vocab.scm -- vocabulary for TB-0 (topological-vector-space.pdf ch. 1 s. 1-2): topological vector spaces and the topology of an ARBITRARY separating family of seminorms.  DEFINITIONS ONLY. Agent TB-01, 2026-10-26 (new definitions under the rule "missing vocabulary is not a stopping point"). DESTINATION: ...}\par
\par\medskip\noindent\textbf{Definition} (\texttt{is-seminorm-family-on}).\;For all $\mathit{fld}$, $m$, $\mathit{ix}$ and $\mathit{pf}$, 
\begin{equation*}
\begin{array}{@{}l@{}} \operatorname{is\text{-}seminorm\text{-}family\text{-}on}(\mathit{fld}, m, \mathit{ix}, \mathit{pf}) \\ \Leftrightarrow\; \operatorname{is\text{-}normed\text{-}field}(\mathit{fld}) \wedge \\ \operatorname{is\text{-}module}(m) \wedge \\ \operatorname{scal}(m) = \operatorname{normed\text{-}field\text{-}as\text{-}commutative\text{-}ring}(\mathit{fld}) \wedge \\ \mathit{ix} \in \mathbf{Set} \wedge \\ \mathit{pf} \in (\mathit{ix} \to (\operatorname{vec}(m) \to \mathbb{R})) \wedge \\ \forall \operatorname{sni\_} \in \mathit{ix}.\; \operatorname{is\text{-}seminorm}(\mathit{fld}, m, \operatorname{pf}(\operatorname{sni\_})) \wedge \\ \forall \operatorname{snx\_} \in \operatorname{vec}(m). \\ \quad \forall \operatorname{sni\_} \in \mathit{ix}.\; (\operatorname{pf}(\operatorname{sni\_})(\operatorname{snx\_}) = 0) \Rightarrow (\operatorname{snx\_} = \operatorname{vzero}(m)) \end{array}
\end{equation*}

\par\medskip\noindent\textbf{Definition} (\texttt{sn-ball}).\;$\operatorname{sn\text{-}ball}(m, \mathit{pf}, \mathit{ff}, x, \mathit{eps}) :=$
\begin{equation*}
\{\operatorname{snby\_} \in \operatorname{vec}(m) : \forall \operatorname{sni\_} \in \mathit{ff}.\; (\operatorname{pf}(\operatorname{sni\_})(\operatorname{vadd}(m)(\operatorname{snby\_}, \operatorname{vneg}(m)(x))) \leq \mathit{eps})\}
\end{equation*}

\par\medskip\noindent\textbf{Definition} (\texttt{is-sn-open}).\;For all $m$, $\mathit{ix}$, $\mathit{pf}$ and $u$, 
\begin{equation*}
\begin{array}{@{}l@{}} \operatorname{is\text{-}sn\text{-}open}(m, \mathit{ix}, \mathit{pf}, u) \\ \Leftrightarrow\; u \subseteq \operatorname{vec}(m) \wedge \\ \forall \operatorname{snx\_} \in u. \exists \operatorname{snf\_} \in \mathbf{Set}. \\ \quad |\operatorname{snf\_}| \in \mathbb{N} \wedge \\ \quad \operatorname{snf\_} \subseteq \mathit{ix} \wedge \\ \quad \exists \operatorname{sne\_}.\; (\operatorname{pos\text{-}rr}(\operatorname{sne\_}) \wedge (\operatorname{sn\text{-}ball}(m, \mathit{pf}, \operatorname{snf\_}, \operatorname{snx\_}, \operatorname{sne\_}) \subseteq u)) \end{array}
\end{equation*}

\par\medskip\noindent\textbf{Definition} (\texttt{seminorm-top}).\;$\operatorname{seminorm\text{-}top}(m, \mathit{ix}, \mathit{pf}) :=$
\begin{equation*}
{[}\operatorname{vec}(m), \{\operatorname{snu\_} \in \operatorname{power}(\operatorname{vec}(m)) : \operatorname{is\text{-}sn\text{-}open}(m, \mathit{ix}, \mathit{pf}, \operatorname{snu\_})\}]
\end{equation*}

\par\medskip\noindent\textbf{Definition} (\texttt{is-tvs}).\;For all $\mathit{fld}$, $m$ and $t$, 
\begin{equation*}
\begin{array}{@{}l@{}} \operatorname{is\text{-}tvs}(\mathit{fld}, m, t) \\ \Leftrightarrow\; \operatorname{is\text{-}normed\text{-}field}(\mathit{fld}) \wedge \\ \operatorname{is\text{-}module}(m) \wedge \\ \operatorname{scal}(m) = \operatorname{normed\text{-}field\text{-}as\text{-}commutative\text{-}ring}(\mathit{fld}) \wedge \\ \operatorname{is\text{-}hausdorff}(t) \wedge \\ (\operatorname{pts}(t) \simeq \operatorname{vec}(m)) \wedge \\ \forall \operatorname{tvx\_}, \operatorname{tvy\_} \in \operatorname{vec}(m), \operatorname{tvu\_} \in \operatorname{opens}(t). \\ \quad (\operatorname{vadd}(m)(\operatorname{tvx\_}, \operatorname{tvy\_}) \in \operatorname{tvu\_}) \Rightarrow  \\ \quad \quad \exists \operatorname{tvv\_} \in \operatorname{opens}(t). \\ \quad \quad \quad \operatorname{tvx\_} \in \operatorname{tvv\_} \wedge \\ \quad \quad \quad \exists \operatorname{tvw\_} \in \operatorname{opens}(t). \\ \quad \quad \quad \quad ((\operatorname{tvy\_} \in \operatorname{tvw\_}) \wedge \forall \operatorname{tva\_} \in \operatorname{tvv\_}, \operatorname{tvb\_} \in \operatorname{tvw\_}.\; (\operatorname{vadd}(m)(\operatorname{tva\_}, \operatorname{tvb\_}) \in \operatorname{tvu\_})) \wedge \\ \forall \operatorname{tvr\_} \in \operatorname{carr}(\operatorname{scal}(m)), \operatorname{tvx\_} \in \operatorname{vec}(m), \operatorname{tvu\_} \in \operatorname{opens}(t). \\ \quad (\operatorname{act}(m)(\operatorname{tvr\_}, \operatorname{tvx\_}) \in \operatorname{tvu\_}) \Rightarrow  \\ \quad \quad \exists \operatorname{tvd\_}. \\ \quad \quad \quad \operatorname{pos\text{-}rr}(\operatorname{tvd\_}) \wedge \\ \quad \quad \quad \exists \operatorname{tvw\_} \in \operatorname{opens}(t). \\ \quad \quad \quad \quad \operatorname{tvx\_} \in \operatorname{tvw\_} \wedge \\ \quad \quad \quad \quad \forall \operatorname{tvs\_} \in \operatorname{carr}(\operatorname{scal}(m)). \\ \quad \quad \quad \quad \quad (\operatorname{fnrm}(\mathit{fld})(\operatorname{add}(\operatorname{scal}(m))(\operatorname{tvs\_}, \operatorname{neg}(\operatorname{scal}(m))(\operatorname{tvr\_}))) < \operatorname{tvd\_}) \Rightarrow  \\ \quad \quad \quad \quad \quad \quad \forall \operatorname{tva\_} \in \operatorname{tvw\_}.\; (\operatorname{act}(m)(\operatorname{tvs\_}, \operatorname{tva\_}) \in \operatorname{tvu\_}) \end{array}
\end{equation*}

\par\medskip\noindent\textbf{Definition} (\texttt{sn-oball}).\;$\operatorname{sn\text{-}oball}(m, \mathit{pf}, \mathit{ff}, x, r) :=$
\begin{equation*}
\{\operatorname{snoy\_} \in \operatorname{vec}(m) : \exists \operatorname{snoe\_} \in \mathbb{R}.\; (\operatorname{pos\text{-}rr}(\operatorname{snoe\_}) \wedge \forall \operatorname{sni\_} \in \mathit{ff}.\; ((\operatorname{pf}(\operatorname{sni\_})(\operatorname{vadd}(m)(\operatorname{snoy\_}, \operatorname{vneg}(m)(x))) + \operatorname{snoe\_}) \leq r))\}
\end{equation*}

\par\medskip\noindent\textbf{Definition} (\texttt{is-extremal-subset}).\;For all $v$, $b$ and $a$, 
\begin{equation*}
\begin{array}{@{}l@{}} \operatorname{is\text{-}extremal\text{-}subset}(v, b, a) \\ \Leftrightarrow\; \operatorname{is\text{-}convex\text{-}subset}(v, b) \wedge \\ \exists \operatorname{exz\_}.\; (\operatorname{exz\_} \in b) \wedge \\ b \subseteq a \wedge \\ \forall \operatorname{exx\_}, \operatorname{exy\_} \in a, \operatorname{ext\_} \in \mathbb{R}. \\ \quad (0 < \operatorname{ext\_}) \Rightarrow  \\ \quad (\operatorname{ext\_} < 1) \Rightarrow  \\ \quad (\operatorname{vadd}(v)(\operatorname{act}(v)(\operatorname{ext\_}, \operatorname{exx\_}), \operatorname{act}(v)(-(1, \operatorname{ext\_}), \operatorname{exy\_})) \in b) \Rightarrow  \\ \quad \quad (\operatorname{exx\_} \in b) \wedge (\operatorname{exy\_} \in b) \end{array}
\end{equation*}


\vnbbookfile{theorem-library/tvs-seminorm-top}{theorem-library/tvs-seminorm-top}
\noindent{\small tb-top.scm -- topological-vector-space.pdf Def 1.3 / Prop 1.4: the topology of a separating family of seminorms is a vector space topology.  Agent TB-01, 2026-10-26. DESTINATION: theorem-library/tvs-seminorm-top.scm.  Library: analysis. Helper prefix: tbt-.}\par
\vnbprop{sn-ball-unfold}{\vnbstmt{\begin{array}{@{}l@{}} \forall m, f, a, x, s. \\ \quad (\operatorname{sn\text{-}ball}(m, f, a, x, s) \simeq \{y \in \operatorname{vec}(m) : \forall i \in a.\; (f(i)(\operatorname{vadd}(m)(y, \operatorname{vneg}(m)(x))) \leq s)\}) \end{array}}}
\vnbprop{sn-ball-membership}{\vnbstmt{\begin{array}{@{}l@{}} \forall m, f, a, x, s, y. \\ \quad y \in \operatorname{sn\text{-}ball}(m, f, a, x, s) \\ \quad \Leftrightarrow\; (y \in \operatorname{vec}(m)) \wedge \forall i \in a.\; (f(i)(\operatorname{vadd}(m)(y, \operatorname{vneg}(m)(x))) \leq s) \end{array}}}
\vnbprop{sntop-pts}{\vnbstmt{\forall m, x, f.\; ((\operatorname{vec}(m) \in \mathbf{Set}) \Rightarrow (\operatorname{pts}(\operatorname{seminorm\text{-}top}(m, x, f)) \simeq \operatorname{vec}(m)))}}
\vnbprop{sntop-opens}{\vnbstmt{\begin{array}{@{}l@{}} \forall m, x, f. \\ \quad (\operatorname{vec}(m) \in \mathbf{Set}) \Rightarrow  \\ \quad \quad (\operatorname{opens}(\operatorname{seminorm\text{-}top}(m, x, f)) \simeq \{u \in \operatorname{power}(\operatorname{vec}(m)) : \operatorname{is\text{-}sn\text{-}open}(m, x, f, u)\}) \end{array}}}
\vnbprop{sn-open-empty}{\vnbstmt{\forall m, x, f.\; (\operatorname{is\text{-}module}(m) \Rightarrow \operatorname{is\text{-}sn\text{-}open}(m, x, f, \emptyset))}}
\vnbprop{sn-open-whole}{\vnbstmt{\forall m, x, f.\; (\operatorname{is\text{-}module}(m) \Rightarrow \operatorname{is\text{-}sn\text{-}open}(m, x, f, \operatorname{vec}(m)))}}
\vnbprop{sn-ball-mono}{\vnbstmt{\begin{array}{@{}l@{}} \forall d, m, a, f, b, g, x, c, h. \\ \quad \operatorname{is\text{-}seminorm\text{-}family\text{-}on}(d, m, a, f) \wedge \\ \quad b \subseteq g \wedge \\ \quad g \subseteq a \wedge \\ \quad x \in \operatorname{vec}(m) \wedge \\ \quad c \in \mathbb{R} \wedge \\ \quad h \in \mathbb{R} \wedge \\ \quad c \leq h \Rightarrow \\ \quad \quad \operatorname{sn\text{-}ball}(m, f, g, x, c) \subseteq \operatorname{sn\text{-}ball}(m, f, b, x, h) \end{array}}}
\vnbprop{sn-open-inter}{\vnbstmt{\begin{array}{@{}l@{}} \forall d, m, x, f, u, v. \\ \quad (\operatorname{is\text{-}seminorm\text{-}family\text{-}on}(d, m, x, f) \wedge (\operatorname{is\text{-}sn\text{-}open}(m, x, f, u) \wedge \operatorname{is\text{-}sn\text{-}open}(m, x, f, v))) \Rightarrow  \\ \quad \quad \operatorname{is\text{-}sn\text{-}open}(m, x, f, (u \cap v)) \end{array}}}
\vnbprop{sn-open-union}{\vnbstmt{\begin{array}{@{}l@{}} \forall m, x, f, a. \\ \quad (\operatorname{is\text{-}module}(m) \wedge \forall w \in a.\; \operatorname{is\text{-}sn\text{-}open}(m, x, f, w)) \Rightarrow  \\ \quad \quad \operatorname{is\text{-}sn\text{-}open}(m, x, f, \operatorname{big\text{-}union}(u, a, u)) \end{array}}}
\vnbprop{sntop-is-top-space}{\vnbstmt{\begin{array}{@{}l@{}} \forall d, m, x, f. \\ \quad \operatorname{is\text{-}seminorm\text{-}family\text{-}on}(d, m, x, f) \Rightarrow \operatorname{is\text{-}top\text{-}space}(\operatorname{seminorm\text{-}top}(m, x, f)) \end{array}}}

\vnbbookfile{theorem-library/tvs-seminorm-tvs}{theorem-library/tvs-seminorm-tvs}
\noindent{\small tb-tvs.scm -- topological-vector-space.pdf Prop 1.4: the topology of a separating family of seminorms is a vector space topology (Hausdorff; addition and scalar multiplication continuous). Agent TB-01, 2026-10-26.  DESTINATION: theorem-library/tvs-seminorm-tvs.scm (after tvs-seminorm-top). Library: analysis.  Helper ...}\par
\vnbprop{sn-oball-unfold}{\vnbstmt{\begin{array}{@{}l@{}} \forall m, f, a, x, r. \\ \quad (\operatorname{sn\text{-}oball}(m, f, a, x, r) \simeq \{y \in \operatorname{vec}(m) : \exists e \in \mathbb{R}.\; (\operatorname{pos\text{-}rr}(e) \wedge \forall i \in a.\; ((f(i)(\operatorname{vadd}(m)(y, \operatorname{vneg}(m)(x))) + e) \leq r))\}) \end{array}}}
\vnbprop{sn-oball-membership}{\vnbstmt{\begin{array}{@{}l@{}} \forall m, f, a, x, r, y. \\ \quad y \in \operatorname{sn\text{-}oball}(m, f, a, x, r) \\ \quad \Leftrightarrow\; y \in \operatorname{vec}(m) \wedge \\ \quad \exists e. \\ \quad \quad \operatorname{pos\text{-}rr}(e) \wedge \forall i \in a.\; ((f(i)(\operatorname{vadd}(m)(y, \operatorname{vneg}(m)(x))) + e) \leq r) \end{array}}}
\vnbprop{sn-oball-open}{\vnbstmt{\begin{array}{@{}l@{}} \forall d, m, a, f, b, x, r. \\ \quad \operatorname{is\text{-}seminorm\text{-}family\text{-}on}(d, m, a, f) \wedge \\ \quad b \in \mathbf{Set} \wedge \\ \quad |b| \in \mathbb{N} \wedge \\ \quad b \subseteq a \wedge \\ \quad x \in \operatorname{vec}(m) \wedge \\ \quad r \in \mathbb{R} \Rightarrow \\ \quad \quad \operatorname{is\text{-}sn\text{-}open}(m, a, f, \operatorname{sn\text{-}oball}(m, f, b, x, r)) \end{array}}}
\vnbprop{sn-oball-center}{\vnbstmt{\begin{array}{@{}l@{}} \forall d, m, a, f, b, x, r. \\ \quad (\operatorname{is\text{-}seminorm\text{-}family\text{-}on}(d, m, a, f) \wedge ((b \subseteq a) \wedge ((x \in \operatorname{vec}(m)) \wedge \operatorname{pos\text{-}rr}(r)))) \Rightarrow  \\ \quad \quad x \in \operatorname{sn\text{-}oball}(m, f, b, x, r) \end{array}}}
\vnbprop{sn-oball-in-ball}{\vnbstmt{\begin{array}{@{}l@{}} \forall d, m, a, f, b, x, r. \\ \quad (\operatorname{is\text{-}seminorm\text{-}family\text{-}on}(d, m, a, f) \wedge ((b \subseteq a) \wedge ((x \in \operatorname{vec}(m)) \wedge (r \in \mathbb{R})))) \Rightarrow  \\ \quad \quad \operatorname{sn\text{-}oball}(m, f, b, x, r) \subseteq \operatorname{sn\text{-}ball}(m, f, b, x, r) \end{array}}}
\vnbprop{sntop-oball-open}{\vnbstmt{\begin{array}{@{}l@{}} \forall d, m, a, f, b, x, r. \\ \quad \operatorname{is\text{-}seminorm\text{-}family\text{-}on}(d, m, a, f) \wedge \\ \quad b \in \mathbf{Set} \wedge \\ \quad |b| \in \mathbb{N} \wedge \\ \quad b \subseteq a \wedge \\ \quad x \in \operatorname{vec}(m) \wedge \\ \quad r \in \mathbb{R} \Rightarrow \\ \quad \quad \operatorname{sn\text{-}oball}(m, f, b, x, r) \in \operatorname{opens}(\operatorname{seminorm\text{-}top}(m, a, f)) \end{array}}}
\vnbprop{sntop-open-basic}{\vnbstmt{\begin{array}{@{}l@{}} \forall d, m, a, f, u, x. \\ \quad (\operatorname{is\text{-}seminorm\text{-}family\text{-}on}(d, m, a, f) \wedge ((u \in \operatorname{opens}(\operatorname{seminorm\text{-}top}(m, a, f))) \wedge (x \in u))) \Rightarrow  \\ \quad \quad \exists \mathit{snf} \in \mathbf{Set}. \\ \quad \quad \quad ((|\mathit{snf}| \in \mathbb{N}) \wedge ((\mathit{snf} \subseteq a) \wedge \exists e.\; (\operatorname{pos\text{-}rr}(e) \wedge (\operatorname{sn\text{-}ball}(m, f, \mathit{snf}, x, e) \subseteq u)))) \end{array}}}
\vnbprop{sntop-hausdorff}{\vnbstmt{\begin{array}{@{}l@{}} \forall d, m, x, f. \\ \quad \operatorname{is\text{-}seminorm\text{-}family\text{-}on}(d, m, x, f) \Rightarrow \operatorname{is\text{-}hausdorff}(\operatorname{seminorm\text{-}top}(m, x, f)) \end{array}}}
\vnbprop{sntop-add-continuous}{\vnbstmt{\begin{array}{@{}l@{}} \forall d, m, x, f, \mathit{tvx}, y, u. \\ \quad \operatorname{is\text{-}seminorm\text{-}family\text{-}on}(d, m, x, f) \wedge \\ \quad \mathit{tvx} \in \operatorname{vec}(m) \wedge \\ \quad y \in \operatorname{vec}(m) \wedge \\ \quad u \in \operatorname{opens}(\operatorname{seminorm\text{-}top}(m, x, f)) \wedge \\ \quad \operatorname{vadd}(m)(\mathit{tvx}, y) \in u \Rightarrow \\ \quad \quad \exists v \in \operatorname{opens}(\operatorname{seminorm\text{-}top}(m, x, f)). \\ \quad \quad \quad \mathit{tvx} \in v \wedge \\ \quad \quad \quad \exists w \in \operatorname{opens}(\operatorname{seminorm\text{-}top}(m, x, f)). \\ \quad \quad \quad \quad (y \in w) \wedge \forall a \in v, b \in w.\; (\operatorname{vadd}(m)(a, b) \in u) \end{array}}}
\vnbprop{sn-finite-bound}{\vnbstmt{\begin{array}{@{}l@{}} \forall d, m, a, f, x.; \forall b \in \mathbf{Set}. \\ \quad (\operatorname{is\text{-}seminorm\text{-}family\text{-}on}(d, m, a, f) \wedge ((x \in \operatorname{vec}(m)) \wedge ((|b| \in \mathbb{N}) \wedge (b \subseteq a)))) \Rightarrow  \\ \quad \quad \exists \mathit{bm} \in \mathbb{R}.\; ((0 \leq \mathit{bm}) \wedge \forall i \in b.\; (f(i)(x) \leq \mathit{bm})) \end{array}}}
\vnbprop{sntop-scal-continuous}{\vnbstmt{\begin{array}{@{}l@{}} \forall d, m, x, f, r, \mathit{tvx}, u. \\ \quad \operatorname{is\text{-}seminorm\text{-}family\text{-}on}(d, m, x, f) \wedge \\ \quad r \in \operatorname{carr}(\operatorname{scal}(m)) \wedge \\ \quad \mathit{tvx} \in \operatorname{vec}(m) \wedge \\ \quad u \in \operatorname{opens}(\operatorname{seminorm\text{-}top}(m, x, f)) \wedge \\ \quad \operatorname{act}(m)(r, \mathit{tvx}) \in u \Rightarrow \\ \quad \quad \exists \mathit{tvd}. \\ \quad \quad \quad \operatorname{pos\text{-}rr}(\mathit{tvd}) \wedge \\ \quad \quad \quad \exists w \in \operatorname{opens}(\operatorname{seminorm\text{-}top}(m, x, f)). \\ \quad \quad \quad \quad \mathit{tvx} \in w \wedge \\ \quad \quad \quad \quad \forall s \in \operatorname{carr}(\operatorname{scal}(m)), a \in w. \\ \quad \quad \quad \quad \quad (\operatorname{fnrm}(d)(\operatorname{add}(\operatorname{scal}(m))(s, \operatorname{neg}(\operatorname{scal}(m))(r))) < \mathit{tvd}) \Rightarrow  \\ \quad \quad \quad \quad \quad \quad \operatorname{act}(m)(s, a) \in u \end{array}}}
\vnbprop{seminorm-top-is-tvs}{\vnbstmt{\begin{array}{@{}l@{}} \forall d, m, x, f. \\ \quad \operatorname{is\text{-}seminorm\text{-}family\text{-}on}(d, m, x, f) \Rightarrow \operatorname{is\text{-}tvs}(d, m, \operatorname{seminorm\text{-}top}(m, x, f)) \end{array}}}

\vnbbookfile{theorem-library/tvs-rn}{theorem-library/tvs-rn}
\noindent{\small tb-rn.scm -- topological-vector-space.pdf Cor 1.5: R\textasciicircum{}n is a topological vector space (and, on the way, every normed vector space is one, with the topology of its norm).  Agent TB-01, 2026-10-26. DESTINATION: theorem-library/tvs-rn.scm, after tvs-seminorm-tvs and the R\textasciicircum{}n files (rn-is-normed-vector-space). Library: ...}\par
\vnbprop{nvs-norm-is-seminorm}{\vnbstmt{\begin{array}{@{}l@{}} \forall m. \\ \quad \operatorname{is\text{-}normed\text{-}vector\text{-}space}(m) \Rightarrow  \\ \quad \quad \operatorname{is\text{-}seminorm}(\operatorname{rr\text{-}normed\text{-}field}, \operatorname{normed\text{-}vector\text{-}space\text{-}as\text{-}module}(m), \operatorname{vnrm}(m)) \end{array}}}
\vnbprop{nvs-norm-family}{\vnbstmt{\begin{array}{@{}l@{}} \forall m. \\ \quad \operatorname{is\text{-}normed\text{-}vector\text{-}space}(m) \Rightarrow  \\ \quad \quad \operatorname{is\text{-}seminorm\text{-}family\text{-}on}(\operatorname{rr\text{-}normed\text{-}field}, \operatorname{normed\text{-}vector\text{-}space\text{-}as\text{-}module}(m), \{0\}, \\ \quad \quad \quad (\lambda k \in \{0\}.\; \operatorname{vnrm}(m))) \end{array}}}
\vnbprop{nvs-norm-top-is-tvs}{\vnbstmt{\begin{array}{@{}l@{}} \forall m. \\ \quad \operatorname{is\text{-}normed\text{-}vector\text{-}space}(m) \Rightarrow  \\ \quad \quad \operatorname{is\text{-}tvs}(\operatorname{rr\text{-}normed\text{-}field}, \operatorname{normed\text{-}vector\text{-}space\text{-}as\text{-}module}(m), \\ \quad \quad \quad \operatorname{seminorm\text{-}top}(\operatorname{normed\text{-}vector\text{-}space\text{-}as\text{-}module}(m), \{0\}, (\lambda k \in \{0\}.\; \operatorname{vnrm}(m)))) \end{array}}}
\vnbprop{rn-is-tvs}{\vnbstmt{\begin{array}{@{}l@{}} \forall n \in \mathbb{N}. \\ \quad \operatorname{is\text{-}tvs}(\operatorname{rr\text{-}normed\text{-}field}, \operatorname{normed\text{-}vector\text{-}space\text{-}as\text{-}module}(\operatorname{rn}(n)), \\ \quad \quad \operatorname{seminorm\text{-}top}(\operatorname{normed\text{-}vector\text{-}space\text{-}as\text{-}module}(\operatorname{rn}(n)), \{0\}, \\ \quad \quad \quad (\lambda k \in \{0\}.\; \operatorname{vnrm}(\operatorname{rn}(n))))) \end{array}}}

\vnbbookfile{theorem-library/tvs-continuity}{theorem-library/tvs-continuity}
\noindent{\small tb-cont.scm -- topological-vector-space.pdf Prop 1.12 / Cor 1.13: a linear map between spaces with separating families of seminorms is continuous iff every q\_j(Tx) is bounded by C max(p\_i1(x), ..., p\_ik(x)). Agent TB-01, 2026-10-26.  DESTINATION: theorem-library/tvs-continuity.scm (after tvs-seminorm-tvs). Library: ...}\par
\vnbprop{sn-bound-combine}{\vnbstmt{\begin{array}{@{}l@{}} \forall d, a, x_{1}, f_{1}, b, x_{2}, f_{2}, t.; \forall g \in \mathbf{Set}. \\ \quad \operatorname{is\text{-}seminorm\text{-}family\text{-}on}(d, a, x_{1}, f_{1}) \wedge \\ \quad \operatorname{is\text{-}seminorm\text{-}family\text{-}on}(d, b, x_{2}, f_{2}) \wedge \\ \quad \operatorname{is\text{-}linear\text{-}map}(a, b, t) \wedge \\ \quad \forall j \in x_{2}. \exists f \in \mathbf{Set}. \\ \quad \quad |f| \in \mathbb{N} \wedge \\ \quad \quad f \subseteq x_{1} \wedge \\ \quad \quad \exists c \in \mathbb{R}. \\ \quad \quad \quad \operatorname{pos\text{-}rr}(c) \wedge \\ \quad \quad \quad \forall x \in \operatorname{vec}(a), \mathit{sbt} \in \mathbb{R}. \\ \quad \quad \quad \quad ((0 \leq \mathit{sbt}) \wedge \forall i \in f.\; (\operatorname{f1}(i)(x) \leq \mathit{sbt})) \Rightarrow  \\ \quad \quad \quad \quad \quad \operatorname{f2}(j)(t(x)) \leq (c \cdot \mathit{sbt}) \wedge \\ \quad |g| \in \mathbb{N} \wedge \\ \quad g \subseteq x_{2} \Rightarrow \\ \quad \quad \exists f \in \mathbf{Set}. \\ \quad \quad \quad |f| \in \mathbb{N} \wedge \\ \quad \quad \quad f \subseteq x_{1} \wedge \\ \quad \quad \quad \exists c \in \mathbb{R}. \\ \quad \quad \quad \quad \operatorname{pos\text{-}rr}(c) \wedge \\ \quad \quad \quad \quad \forall j \in g, x \in \operatorname{vec}(a), \mathit{sbt} \in \mathbb{R}. \\ \quad \quad \quad \quad \quad ((0 \leq \mathit{sbt}) \wedge \forall i \in f.\; (\operatorname{f1}(i)(x) \leq \mathit{sbt})) \Rightarrow  \\ \quad \quad \quad \quad \quad \quad \operatorname{f2}(j)(t(x)) \leq (c \cdot \mathit{sbt}) \end{array}}}
\vnbprop{sn-cont-of-bound}{\vnbstmt{\begin{array}{@{}l@{}} \forall d, a, x_{1}, f_{1}, b, x_{2}, f_{2}, t. \\ \quad \operatorname{is\text{-}seminorm\text{-}family\text{-}on}(d, a, x_{1}, f_{1}) \wedge \\ \quad \operatorname{is\text{-}seminorm\text{-}family\text{-}on}(d, b, x_{2}, f_{2}) \wedge \\ \quad \operatorname{is\text{-}linear\text{-}map}(a, b, t) \wedge \\ \quad \forall j \in x_{2}. \exists f \in \mathbf{Set}. \\ \quad \quad |f| \in \mathbb{N} \wedge \\ \quad \quad f \subseteq x_{1} \wedge \\ \quad \quad \exists c \in \mathbb{R}. \\ \quad \quad \quad \operatorname{pos\text{-}rr}(c) \wedge \\ \quad \quad \quad \forall x \in \operatorname{vec}(a), \mathit{sbt} \in \mathbb{R}. \\ \quad \quad \quad \quad ((0 \leq \mathit{sbt}) \wedge \forall i \in f.\; (\operatorname{f1}(i)(x) \leq \mathit{sbt})) \Rightarrow  \\ \quad \quad \quad \quad \quad \operatorname{f2}(j)(t(x)) \leq (c \cdot \mathit{sbt}) \Rightarrow \\ \quad \quad \operatorname{is\text{-}hom\text{-}top\text{-}space}(\operatorname{seminorm\text{-}top}(a, x_{1}, f_{1}), \operatorname{seminorm\text{-}top}(b, x_{2}, f_{2}), t) \end{array}}}
\vnbprop{sn-bound-of-cont}{\vnbstmt{\begin{array}{@{}l@{}} \forall a, x_{1}, f_{1}, b, x_{2}, f_{2}, t.; \forall j \in x_{2}. \\ \quad \operatorname{is\text{-}seminorm\text{-}family\text{-}on}(\operatorname{rr\text{-}normed\text{-}field}, a, x_{1}, f_{1}) \wedge \\ \quad \operatorname{is\text{-}seminorm\text{-}family\text{-}on}(\operatorname{rr\text{-}normed\text{-}field}, b, x_{2}, f_{2}) \wedge \\ \quad \operatorname{is\text{-}linear\text{-}map}(a, b, t) \wedge \\ \quad \operatorname{is\text{-}hom\text{-}top\text{-}space}(\operatorname{seminorm\text{-}top}(a, x_{1}, f_{1}), \operatorname{seminorm\text{-}top}(b, x_{2}, f_{2}), t) \Rightarrow \\ \quad \quad \exists f \in \mathbf{Set}. \\ \quad \quad \quad |f| \in \mathbb{N} \wedge \\ \quad \quad \quad f \subseteq x_{1} \wedge \\ \quad \quad \quad \exists c \in \mathbb{R}. \\ \quad \quad \quad \quad \operatorname{pos\text{-}rr}(c) \wedge \\ \quad \quad \quad \quad \forall x \in \operatorname{vec}(a), \mathit{sbt} \in \mathbb{R}. \\ \quad \quad \quad \quad \quad ((0 \leq \mathit{sbt}) \wedge \forall i \in f.\; (\operatorname{f1}(i)(x) \leq \mathit{sbt})) \Rightarrow  \\ \quad \quad \quad \quad \quad \quad \operatorname{f2}(j)(t(x)) \leq (c \cdot \mathit{sbt}) \end{array}}}
\vnbprop{sn-cont-iff-bound}{\vnbstmt{\begin{array}{@{}l@{}} \forall a, x_{1}, f_{1}, b, x_{2}, f_{2}, t. \\ \quad \operatorname{is\text{-}seminorm\text{-}family\text{-}on}(\operatorname{rr\text{-}normed\text{-}field}, a, x_{1}, f_{1}) \wedge \\ \quad \operatorname{is\text{-}seminorm\text{-}family\text{-}on}(\operatorname{rr\text{-}normed\text{-}field}, b, x_{2}, f_{2}) \wedge \\ \quad \operatorname{is\text{-}linear\text{-}map}(a, b, t) \Rightarrow \\ \quad \quad \operatorname{is\text{-}hom\text{-}top\text{-}space}(\operatorname{seminorm\text{-}top}(a, x_{1}, f_{1}), \operatorname{seminorm\text{-}top}(b, x_{2}, f_{2}), t) \\ \quad \quad \Leftrightarrow\; \forall j \in x_{2}. \exists f \in \mathbf{Set}. \\ \quad \quad \quad |f| \in \mathbb{N} \wedge \\ \quad \quad \quad f \subseteq x_{1} \wedge \\ \quad \quad \quad \exists c \in \mathbb{R}. \\ \quad \quad \quad \quad \operatorname{pos\text{-}rr}(c) \wedge \\ \quad \quad \quad \quad \forall x \in \operatorname{vec}(a), \mathit{sbt} \in \mathbb{R}. \\ \quad \quad \quad \quad \quad ((0 \leq \mathit{sbt}) \wedge \forall i \in f.\; (\operatorname{f1}(i)(x) \leq \mathit{sbt})) \Rightarrow  \\ \quad \quad \quad \quad \quad \quad \operatorname{f2}(j)(t(x)) \leq (c \cdot \mathit{sbt}) \end{array}}}

\vnbbookfile{theorem-library/tvs-single-bound}{theorem-library/tvs-single-bound}
\noindent{\small tvs-single-bound.scm -- wave of 2026-11-09/10 (agent BA-OT1), wired 2026-11-10 from scratchpad/baot1-1109/tvs-single-bound.scm. Library: analysis.}\par
\vnbprop{sn-cont-of-single-bound}{\vnbstmt{\begin{array}{@{}l@{}} \forall d, a, x_{1}, f_{1}, b, x_{2}, f_{2}, t. \\ \quad \operatorname{is\text{-}seminorm\text{-}family\text{-}on}(d, a, x_{1}, f_{1}) \wedge \\ \quad \operatorname{is\text{-}seminorm\text{-}family\text{-}on}(d, b, x_{2}, f_{2}) \wedge \\ \quad \operatorname{is\text{-}linear\text{-}map}(a, b, t) \wedge \\ \quad \forall j \in x_{2}. \exists i \in x_{1}, c \in \mathbb{R}. \\ \quad \quad (0 \leq c) \wedge \forall x \in \operatorname{vec}(a).\; (\operatorname{f2}(j)(t(x)) \leq (c \cdot \operatorname{f1}(i)(x))) \Rightarrow \\ \quad \quad \operatorname{is\text{-}hom\text{-}top\text{-}space}(\operatorname{seminorm\text{-}top}(a, x_{1}, f_{1}), \operatorname{seminorm\text{-}top}(b, x_{2}, f_{2}), t) \end{array}}}

\vnbbookfile{theorem-library/tvs-functional}{theorem-library/tvs-functional}
\noindent{\small tb-func.scm -- topological-vector-space.pdf Cor 1.13: a linear functional f on a space whose topology is defined by a family of seminorms is continuous iff |f(x)| $<$= C max(p\_i1(x), ..., p\_in(x)). Agent TB-01, 2026-10-26.  DESTINATION: theorem-library/tvs-functional.scm (after tvs-continuity, tvs-rn). Library: analysis. ...}\par
\vnbprop{rr-nvs-norm-apply}{\vnbstmt{\forall v \in \mathbb{R}.\; (\operatorname{vnrm}(\operatorname{rr\text{-}nvs})(v) = \lvert v \rvert)}}
\vnbprop{rr-nvs-vec}{\vnbstmt{(\operatorname{vec}(\operatorname{rr\text{-}nvs}) \simeq \mathbb{R})}}
\vnbprop{functional-cont-iff-bound}{\vnbstmt{\begin{array}{@{}l@{}} \forall m, x, a, f. \\ \quad \operatorname{is\text{-}seminorm\text{-}family\text{-}on}(\operatorname{rr\text{-}normed\text{-}field}, m, x, a) \wedge \\ \quad \operatorname{is\text{-}linear\text{-}map}(m, \operatorname{normed\text{-}vector\text{-}space\text{-}as\text{-}module}(\operatorname{rr\text{-}nvs}), f) \Rightarrow \\ \quad \quad \operatorname{is\text{-}hom\text{-}top\text{-}space}(\operatorname{seminorm\text{-}top}(m, x, a), \\ \quad \quad \quad \operatorname{seminorm\text{-}top}(\operatorname{normed\text{-}vector\text{-}space\text{-}as\text{-}module}(\operatorname{rr\text{-}nvs}), \\ \quad \quad \quad \quad \{0\}, (\lambda k \in \{0\}.\; \operatorname{vnrm}(\operatorname{rr\text{-}nvs}))), \\ \quad \quad \quad f) \\ \quad \quad \Leftrightarrow\; \exists \mathit{sbf} \in \mathbf{Set}. \\ \quad \quad \quad |\mathit{sbf}| \in \mathbb{N} \wedge \\ \quad \quad \quad \mathit{sbf} \subseteq x \wedge \\ \quad \quad \quad \exists c \in \mathbb{R}. \\ \quad \quad \quad \quad \operatorname{pos\text{-}rr}(c) \wedge \\ \quad \quad \quad \quad \forall \mathit{sbx} \in \operatorname{vec}(m), t \in \mathbb{R}. \\ \quad \quad \quad \quad \quad ((0 \leq t) \wedge \forall i \in \mathit{sbf}.\; (a(i)(\mathit{sbx}) \leq t)) \Rightarrow (\lvert f(\mathit{sbx}) \rvert \leq (c \cdot t)) \end{array}}}

\vnbbookfile{theorem-library/tvs-extremal}{theorem-library/tvs-extremal}
\noindent{\small tb-extremal.scm -- topological-vector-space.pdf Prop 2.40: extremal subsets are transitive, closed under intersections (with nonempty intersection) and under inverse images by linear maps. Agent TB-01, 2026-10-26.  DESTINATION: theorem-library/tvs-extremal.scm (after sip-hilbert-vocab, the tvs vocabulary).  Library: ...}\par
\vnbprop{extremal-subset-trans}{\vnbstmt{\begin{array}{@{}l@{}} \forall v, c, b, a. \\ \quad (\operatorname{is\text{-}extremal\text{-}subset}(v, c, b) \wedge \operatorname{is\text{-}extremal\text{-}subset}(v, b, a)) \Rightarrow  \\ \quad \quad \operatorname{is\text{-}extremal\text{-}subset}(v, c, a) \end{array}}}
\vnbprop{extremal-subset-inter}{\vnbstmt{\begin{array}{@{}l@{}} \forall v, x, f, a. \\ \quad \operatorname{is\text{-}module}(v) \wedge \\ \quad \operatorname{scal}(v) = \operatorname{normed\text{-}field\text{-}as\text{-}commutative\text{-}ring}(\operatorname{rr\text{-}normed\text{-}field}) \wedge \\ \quad f \in (x \to \operatorname{power}(\operatorname{vec}(v))) \wedge \\ \quad a \in (x \to \operatorname{power}(\operatorname{vec}(v))) \wedge \\ \quad \forall l \in x.\; \operatorname{is\text{-}extremal\text{-}subset}(v, f(l), a(l)) \wedge \\ \quad \exists z.\; (z \in \{w \in \operatorname{vec}(v) : \forall l \in x.\; (w \in f(l))\}) \Rightarrow \\ \quad \quad \operatorname{is\text{-}extremal\text{-}subset}(v, \{w \in \operatorname{vec}(v) : \forall l \in x.\; (w \in f(l))\}, \\ \quad \quad \quad \{w \in \operatorname{vec}(v) : \forall l \in x.\; (w \in a(l))\}) \end{array}}}
\vnbprop{extremal-subset-preimage}{\vnbstmt{\begin{array}{@{}l@{}} \forall c, d, f, b, a. \\ \quad \operatorname{is\text{-}linear\text{-}map}(c, d, f) \wedge \\ \quad \operatorname{scal}(c) = \operatorname{normed\text{-}field\text{-}as\text{-}commutative\text{-}ring}(\operatorname{rr\text{-}normed\text{-}field}) \wedge \\ \quad \operatorname{is\text{-}extremal\text{-}subset}(d, b, a) \wedge \\ \quad \exists z.\; (z \in \{w \in \operatorname{vec}(c) : (f(w) \in b)\}) \Rightarrow \\ \quad \quad \operatorname{is\text{-}extremal\text{-}subset}(c, \{w \in \operatorname{vec}(c) : (f(w) \in b)\}, \{w \in \operatorname{vec}(c) : (f(w) \in a)\}) \end{array}}}

\vnbbookfile{structure-library/convex-functional}{structure-library/convex-functional}
\noindent{\small hbc-vocab.scm -- convex functionals (topological-vector-space.pdf Def 2.34) and the positively homogeneous minorant used to prove Cor 2.36 from the sublinear Hahn-Banach.  Vocabulary only.  Agent DUAL-FIX, 2026-10-26 (NEW DEFINITIONS, the user's to confirm).  DESTINATION: structure-library/convex-functional.scm. ...}\par
\par\medskip\noindent\textbf{Definition} (\texttt{is-convex-functional}).\;For all $m$ and $p$, 
\begin{equation*}
\begin{array}{@{}l@{}} \operatorname{is\text{-}convex\text{-}functional}(m, p) \\ \Leftrightarrow\; p \in (\operatorname{vec}(m) \to \mathbb{R}) \wedge \\ \forall \operatorname{cfx\_} \in \operatorname{vec}(m).\; (0 \leq p(\operatorname{cfx\_})) \wedge \\ \forall \operatorname{cfx\_}, \operatorname{cfy\_} \in \operatorname{vec}(m), \operatorname{cft\_} \in \mathbb{R}. \\ \quad (0 \leq \operatorname{cft\_}) \Rightarrow  \\ \quad (\operatorname{cft\_} \leq 1) \Rightarrow  \\ \quad \quad p(\operatorname{vadd}(m)(\operatorname{act}(m)(\operatorname{cft\_}, \operatorname{cfx\_}), \operatorname{act}(m)(-(1, \operatorname{cft\_}), \operatorname{cfy\_}))) \\ \quad \quad \leq\; +(\cdot(\operatorname{cft\_}, p(\operatorname{cfx\_})), \cdot(-(1, \operatorname{cft\_}), p(\operatorname{cfy\_}))) \end{array}
\end{equation*}

\par\medskip\noindent\textbf{Definition} (\texttt{homog-set}).\;$\operatorname{homog\text{-}set}(m, p, x) :=$
\begin{equation*}
\{\operatorname{hmr\_} \in \mathbb{R} : \exists \operatorname{hmt\_} \in \mathbb{R}.\; ((0 < \operatorname{hmt\_}) \wedge (p(\operatorname{act}(m)(\operatorname{hmt\_}, x)) \leq (\operatorname{hmt\_} \cdot \operatorname{hmr\_})))\}
\end{equation*}

\par\medskip\noindent\textbf{Definition} (\texttt{homog-minorant}).\;$\operatorname{homog\text{-}minorant}(m, p) :=$ $(\lambda \operatorname{hmx\_} \in \operatorname{vec}(m).\; \operatorname{inf}(\operatorname{homog\text{-}set}(m, p, \operatorname{hmx\_})))$

\vnbbookfile{theorem-library/tvs-hahn-banach-convex}{theorem-library/tvs-hahn-banach-convex}
\noindent{\small hbc.scm -- topological-vector-space.pdf Cor 2.36: Hahn-Banach for a convex dominating functional, by CITING the sublinear Hahn-Banach (hbs-hahn-banach-sublinear) -- no Zorn here.  Agent DUAL-FIX, 2026-10-26. DESTINATION: theorem-library/tvs-hahn-banach-convex.scm.  Library: analysis.  Helper prefix: hbc-. Needs: ...}\par
\vnbprop{homog-set-unfold}{\vnbstmt{\begin{array}{@{}l@{}} \forall m, p, x. \\ \quad (\operatorname{homog\text{-}set}(m, p, x) \simeq \{r \in \mathbb{R} : \exists t \in \mathbb{R}.\; ((0 < t) \wedge (p(\operatorname{act}(m)(t, x)) \leq (t \cdot r)))\}) \end{array}}}
\vnbprop{homog-set-intro}{\vnbstmt{\begin{array}{@{}l@{}} \forall m, p, x, r, t. \\ \quad ((r \in \mathbb{R}) \wedge ((t \in \mathbb{R}) \wedge ((0 < t) \wedge (p(\operatorname{act}(m)(t, x)) \leq (t \cdot r))))) \Rightarrow  \\ \quad \quad r \in \operatorname{homog\text{-}set}(m, p, x) \end{array}}}
\vnbprop{homog-set-elim}{\vnbstmt{\begin{array}{@{}l@{}} \forall m, p, x.; \forall r \in \operatorname{homog\text{-}set}(m, p, x). \\ \quad (r \in \mathbb{R}) \wedge \exists t \in \mathbb{R}.\; ((0 < t) \wedge (p(\operatorname{act}(m)(t, x)) \leq (t \cdot r))) \end{array}}}
\vnbprop{homog-elem-nonneg}{\vnbstmt{\begin{array}{@{}l@{}} \forall m, p, x, r. \\ \quad \operatorname{is\text{-}module}(m) \wedge \\ \quad \operatorname{scal}(m) = \operatorname{normed\text{-}field\text{-}as\text{-}commutative\text{-}ring}(\operatorname{rr\text{-}normed\text{-}field}) \wedge \\ \quad \operatorname{is\text{-}convex\text{-}functional}(m, p) \wedge \\ \quad x \in \operatorname{vec}(m) \wedge \\ \quad r \in \operatorname{homog\text{-}set}(m, p, x) \Rightarrow \\ \quad \quad 0 \leq r \end{array}}}
\vnbprop{homog-p-in-set}{\vnbstmt{\begin{array}{@{}l@{}} \forall m, p, x. \\ \quad \operatorname{is\text{-}module}(m) \wedge \\ \quad \operatorname{scal}(m) = \operatorname{normed\text{-}field\text{-}as\text{-}commutative\text{-}ring}(\operatorname{rr\text{-}normed\text{-}field}) \wedge \\ \quad \operatorname{is\text{-}convex\text{-}functional}(m, p) \wedge \\ \quad x \in \operatorname{vec}(m) \Rightarrow \\ \quad \quad p(x) \in \operatorname{homog\text{-}set}(m, p, x) \end{array}}}
\vnbprop{homog-set-glb}{\vnbstmt{\begin{array}{@{}l@{}} \forall m, p, z. \\ \quad \operatorname{is\text{-}module}(m) \wedge \\ \quad \operatorname{scal}(m) = \operatorname{normed\text{-}field\text{-}as\text{-}commutative\text{-}ring}(\operatorname{rr\text{-}normed\text{-}field}) \wedge \\ \quad \operatorname{is\text{-}convex\text{-}functional}(m, p) \wedge \\ \quad z \in \operatorname{vec}(m) \Rightarrow \\ \quad \quad \operatorname{inf}(\operatorname{homog\text{-}set}(m, p, z)) \in \mathbb{R} \wedge \\ \quad \quad \forall w \in \operatorname{homog\text{-}set}(m, p, z).\; (\operatorname{inf}(\operatorname{homog\text{-}set}(m, p, z)) \leq w) \wedge \\ \quad \quad \forall b \in \mathbb{R}. \\ \quad \quad \quad \forall w \in \operatorname{homog\text{-}set}(m, p, z).\; (b \leq w) \Rightarrow (b \leq \operatorname{inf}(\operatorname{homog\text{-}set}(m, p, z))) \end{array}}}
\vnbprop{homog-value}{\vnbstmt{\begin{array}{@{}l@{}} \forall m, p.; \forall z \in \operatorname{vec}(m). \\ \quad (\operatorname{homog\text{-}minorant}(m, p)(z) \simeq \operatorname{inf}(\operatorname{homog\text{-}set}(m, p, z))) \end{array}}}
\vnbprop{homog-inf-le-p}{\vnbstmt{\begin{array}{@{}l@{}} \forall m, p, x. \\ \quad \operatorname{is\text{-}module}(m) \wedge \\ \quad \operatorname{scal}(m) = \operatorname{normed\text{-}field\text{-}as\text{-}commutative\text{-}ring}(\operatorname{rr\text{-}normed\text{-}field}) \wedge \\ \quad \operatorname{is\text{-}convex\text{-}functional}(m, p) \wedge \\ \quad x \in \operatorname{vec}(m) \Rightarrow \\ \quad \quad \operatorname{inf}(\operatorname{homog\text{-}set}(m, p, x)) \leq p(x) \end{array}}}
\vnbprop{homog-inf-nonneg}{\vnbstmt{\begin{array}{@{}l@{}} \forall m, p, x. \\ \quad \operatorname{is\text{-}module}(m) \wedge \\ \quad \operatorname{scal}(m) = \operatorname{normed\text{-}field\text{-}as\text{-}commutative\text{-}ring}(\operatorname{rr\text{-}normed\text{-}field}) \wedge \\ \quad \operatorname{is\text{-}convex\text{-}functional}(m, p) \wedge \\ \quad x \in \operatorname{vec}(m) \Rightarrow \\ \quad \quad 0 \leq \operatorname{inf}(\operatorname{homog\text{-}set}(m, p, x)) \end{array}}}
\vnbprop{homog-minorant-dominates}{\vnbstmt{\begin{array}{@{}l@{}} \forall m, p, s, f, z. \\ \quad \operatorname{is\text{-}module}(m) \wedge \\ \quad \operatorname{scal}(m) = \operatorname{normed\text{-}field\text{-}as\text{-}commutative\text{-}ring}(\operatorname{rr\text{-}normed\text{-}field}) \wedge \\ \quad \operatorname{is\text{-}convex\text{-}functional}(m, p) \wedge \\ \quad \operatorname{is\text{-}submodule}(m, s) \wedge \\ \quad \operatorname{is\text{-}linear\text{-}functional\text{-}on}(m, s, f) \wedge \\ \quad \forall \mathit{dmz} \in s.\; (f(\mathit{dmz}) \leq p(\mathit{dmz})) \wedge \\ \quad z \in s \Rightarrow \\ \quad \quad f(z) \leq \operatorname{inf}(\operatorname{homog\text{-}set}(m, p, z)) \end{array}}}
\vnbprop{homog-set-zero}{\vnbstmt{\begin{array}{@{}l@{}} \forall m, p, r. \\ \quad \operatorname{is\text{-}module}(m) \wedge \\ \quad \operatorname{scal}(m) = \operatorname{normed\text{-}field\text{-}as\text{-}commutative\text{-}ring}(\operatorname{rr\text{-}normed\text{-}field}) \wedge \\ \quad \operatorname{is\text{-}convex\text{-}functional}(m, p) \wedge \\ \quad r \in \mathbb{R} \wedge \\ \quad 0 < r \Rightarrow \\ \quad \quad r \in \operatorname{homog\text{-}set}(m, p, \operatorname{vzero}(m)) \end{array}}}
\vnbprop{homog-inf-zero}{\vnbstmt{\begin{array}{@{}l@{}} \forall m, p. \\ \quad \operatorname{is\text{-}module}(m) \wedge \\ \quad \operatorname{scal}(m) = \operatorname{normed\text{-}field\text{-}as\text{-}commutative\text{-}ring}(\operatorname{rr\text{-}normed\text{-}field}) \wedge \\ \quad \operatorname{is\text{-}convex\text{-}functional}(m, p) \Rightarrow \\ \quad \quad \operatorname{inf}(\operatorname{homog\text{-}set}(m, p, \operatorname{vzero}(m))) = 0 \end{array}}}
\vnbprop{homog-set-scale}{\vnbstmt{\begin{array}{@{}l@{}} \forall m, p, x, c, r. \\ \quad \operatorname{is\text{-}module}(m) \wedge \\ \quad \operatorname{scal}(m) = \operatorname{normed\text{-}field\text{-}as\text{-}commutative\text{-}ring}(\operatorname{rr\text{-}normed\text{-}field}) \wedge \\ \quad \operatorname{is\text{-}convex\text{-}functional}(m, p) \wedge \\ \quad x \in \operatorname{vec}(m) \wedge \\ \quad c \in \mathbb{R} \wedge \\ \quad 0 < c \wedge \\ \quad r \in \operatorname{homog\text{-}set}(m, p, x) \Rightarrow \\ \quad \quad (c \cdot r) \in \operatorname{homog\text{-}set}(m, p, \operatorname{act}(m)(c, x)) \end{array}}}
\vnbprop{homog-set-unscale}{\vnbstmt{\begin{array}{@{}l@{}} \forall m, p, x, c, r. \\ \quad \operatorname{is\text{-}module}(m) \wedge \\ \quad \operatorname{scal}(m) = \operatorname{normed\text{-}field\text{-}as\text{-}commutative\text{-}ring}(\operatorname{rr\text{-}normed\text{-}field}) \wedge \\ \quad \operatorname{is\text{-}convex\text{-}functional}(m, p) \wedge \\ \quad x \in \operatorname{vec}(m) \wedge \\ \quad c \in \mathbb{R} \wedge \\ \quad 0 < c \wedge \\ \quad r \in \operatorname{homog\text{-}set}(m, p, \operatorname{act}(m)(c, x)) \Rightarrow \\ \quad \quad (\operatorname{recip}(c) \cdot r) \in \operatorname{homog\text{-}set}(m, p, x) \end{array}}}
\vnbprop{homog-inf-scale}{\vnbstmt{\begin{array}{@{}l@{}} \forall m, p, x, c. \\ \quad \operatorname{is\text{-}module}(m) \wedge \\ \quad \operatorname{scal}(m) = \operatorname{normed\text{-}field\text{-}as\text{-}commutative\text{-}ring}(\operatorname{rr\text{-}normed\text{-}field}) \wedge \\ \quad \operatorname{is\text{-}convex\text{-}functional}(m, p) \wedge \\ \quad x \in \operatorname{vec}(m) \wedge \\ \quad c \in \mathbb{R} \wedge \\ \quad 0 < c \Rightarrow \\ \quad \quad \operatorname{inf}(\operatorname{homog\text{-}set}(m, p, \operatorname{act}(m)(c, x))) = (c \cdot \operatorname{inf}(\operatorname{homog\text{-}set}(m, p, x))) \end{array}}}
\vnbprop{homog-set-sum}{\vnbstmt{\begin{array}{@{}l@{}} \forall m, p, x, y, r, a. \\ \quad \operatorname{is\text{-}module}(m) \wedge \\ \quad \operatorname{scal}(m) = \operatorname{normed\text{-}field\text{-}as\text{-}commutative\text{-}ring}(\operatorname{rr\text{-}normed\text{-}field}) \wedge \\ \quad \operatorname{is\text{-}convex\text{-}functional}(m, p) \wedge \\ \quad x \in \operatorname{vec}(m) \wedge \\ \quad y \in \operatorname{vec}(m) \wedge \\ \quad r \in \operatorname{homog\text{-}set}(m, p, x) \wedge \\ \quad a \in \operatorname{homog\text{-}set}(m, p, y) \Rightarrow \\ \quad \quad (r + a) \in \operatorname{homog\text{-}set}(m, p, \operatorname{vadd}(m)(x, y)) \end{array}}}
\vnbprop{homog-inf-subadditive}{\vnbstmt{\begin{array}{@{}l@{}} \forall m, p, x, y. \\ \quad \operatorname{is\text{-}module}(m) \wedge \\ \quad \operatorname{scal}(m) = \operatorname{normed\text{-}field\text{-}as\text{-}commutative\text{-}ring}(\operatorname{rr\text{-}normed\text{-}field}) \wedge \\ \quad \operatorname{is\text{-}convex\text{-}functional}(m, p) \wedge \\ \quad x \in \operatorname{vec}(m) \wedge \\ \quad y \in \operatorname{vec}(m) \Rightarrow \\ \quad \quad \operatorname{inf}(\operatorname{homog\text{-}set}(m, p, \operatorname{vadd}(m)(x, y))) \leq (\operatorname{inf}(\operatorname{homog\text{-}set}(m, p, x)) + \operatorname{inf}(\operatorname{homog\text{-}set}(m, p, y))) \end{array}}}
\vnbprop{homog-minorant-sublinear}{\vnbstmt{\begin{array}{@{}l@{}} \forall m, p. \\ \quad \operatorname{is\text{-}module}(m) \wedge \\ \quad \operatorname{scal}(m) = \operatorname{normed\text{-}field\text{-}as\text{-}commutative\text{-}ring}(\operatorname{rr\text{-}normed\text{-}field}) \wedge \\ \quad \operatorname{is\text{-}convex\text{-}functional}(m, p) \Rightarrow \\ \quad \quad \operatorname{is\text{-}sublinear}(m, \operatorname{homog\text{-}minorant}(m, p)) \end{array}}}
\vnbprop{hahn-banach-convex}{\vnbstmt{\begin{array}{@{}l@{}} \forall m, p, s, f. \\ \quad \operatorname{is\text{-}module}(m) \wedge \\ \quad \operatorname{scal}(m) = \operatorname{normed\text{-}field\text{-}as\text{-}commutative\text{-}ring}(\operatorname{rr\text{-}normed\text{-}field}) \wedge \\ \quad \operatorname{is\text{-}convex\text{-}functional}(m, p) \wedge \\ \quad \operatorname{is\text{-}submodule}(m, s) \wedge \\ \quad \operatorname{is\text{-}linear\text{-}functional\text{-}on}(m, s, f) \wedge \\ \quad \forall z \in s.\; (f(z) \leq p(z)) \Rightarrow \\ \quad \quad \exists a. \\ \quad \quad \quad \operatorname{is\text{-}linear\text{-}functional\text{-}on}(m, \operatorname{vec}(m), a) \wedge \\ \quad \quad \quad \operatorname{extends\text{-}on}(s, a, f) \wedge \\ \quad \quad \quad \forall x \in \operatorname{vec}(m).\; (a(x) \leq p(x)) \end{array}}}

\vnbbookfile{structure-library/top-closed}{structure-library/top-closed}
\noindent{\small tkr-vocab.scm -- closed sets of a topological space and the kernel of a functional (for topological-vector-space.pdf Prop 2.11).  Vocabulary only.  Agent DUAL-FIX, 2026-10-26 (NEW DEFINITIONS, the user's to confirm).  DESTINATION: structure-library/top-closed.scm.  Library: analysis. Load window: after ...}\par
\par\medskip\noindent\textbf{Definition} (\texttt{is-top-closed}).\;For all $t$ and $a$, 
\begin{equation*}
\operatorname{is\text{-}top\text{-}closed}(t, a) \Leftrightarrow ((a \subseteq \operatorname{pts}(t)) \wedge (\operatorname{complement\text{-}in}(\operatorname{pts}(t), a) \in \operatorname{opens}(t)))
\end{equation*}

\par\medskip\noindent\textbf{Definition} (\texttt{fker}).\;$\operatorname{fker}(m, f) :=$ $\{\operatorname{fkz\_} \in \operatorname{vec}(m) : (f(\operatorname{fkz\_}) = 0)\}$

\vnbbookfile{theorem-library/tvs-kernel-closed}{theorem-library/tvs-kernel-closed}
\noindent{\small tkr.scm -- topological-vector-space.pdf Prop 2.11: a nontrivial linear functional on a real topological vector space is continuous iff its kernel f\textasciicircum{}-1(0) is closed.  Agent DUAL-FIX, 2026-10-26. DESTINATION: theorem-library/tvs-kernel-closed.scm.  Library: analysis.  Helper prefix: tkr-. Needs: ...}\par
\vnbprop{rr-nvs-vadd-apply}{\vnbstmt{\forall v, w \in \mathbb{R}.\; (\operatorname{vadd}(\operatorname{rr\text{-}nvs})(v, w) = (v + w))}}
\vnbprop{rr-nvs-vneg-apply}{\vnbstmt{\forall v \in \mathbb{R}.\; (\operatorname{vneg}(\operatorname{rr\text{-}nvs})(v) = -v)}}
\vnbprop{rr-abs-top-facts}{\vnbstmt{\begin{array}{@{}l@{}} \operatorname{is\text{-}top\text{-}space}(\operatorname{seminorm\text{-}top}(\operatorname{normed\text{-}vector\text{-}space\text{-}as\text{-}module}(\operatorname{rr\text{-}nvs}), \\ \quad \quad \{0\}, (\lambda k \in \{0\}.\; \operatorname{vnrm}(\operatorname{rr\text{-}nvs})))) \wedge \\ (\operatorname{pts}(\operatorname{seminorm\text{-}top}(\operatorname{normed\text{-}vector\text{-}space\text{-}as\text{-}module}(\operatorname{rr\text{-}nvs}), \{0\}, (\lambda k \in \{0\}.\; \operatorname{vnrm}(\operatorname{rr\text{-}nvs})))) \simeq \mathbb{R}) \end{array}}}
\vnbprop{rr-abs-open-ball}{\vnbstmt{\begin{array}{@{}l@{}} \forall u \in \operatorname{opens}(\operatorname{seminorm\text{-}top}(\operatorname{normed\text{-}vector\text{-}space\text{-}as\text{-}module}(\operatorname{rr\text{-}nvs}), \{0\}, (\lambda k \in \{0\}.\; \operatorname{vnrm}(\operatorname{rr\text{-}nvs})))), y \in u. \exists e. \\ \quad \operatorname{pos\text{-}rr}(e) \wedge \forall v \in \mathbb{R}.\; ((\lvert (v - y) \rvert \leq e) \Rightarrow (v \in u)) \end{array}}}
\vnbprop{rr-abs-nonzero-open}{\vnbstmt{\begin{array}{@{}l@{}} \{v \in \mathbb{R} : \neg (v = 0)\} \\ \in\; \operatorname{opens}(\operatorname{seminorm\text{-}top}(\operatorname{normed\text{-}vector\text{-}space\text{-}as\text{-}module}(\operatorname{rr\text{-}nvs}), \\ \quad \quad \{0\}, (\lambda k \in \{0\}.\; \operatorname{vnrm}(\operatorname{rr\text{-}nvs})))) \end{array}}}
\vnbprop{tvs-scal-cont-rr}{\vnbstmt{\begin{array}{@{}l@{}} \forall m, t, r, x, u. \\ \quad (\operatorname{is\text{-}tvs}(\operatorname{rr\text{-}normed\text{-}field}, m, t) \wedge ((r \in \mathbb{R}) \wedge ((x \in \operatorname{vec}(m)) \wedge ((u \in \operatorname{opens}(t)) \wedge (\operatorname{act}(m)(r, x) \in u))))) \Rightarrow  \\ \quad \quad \exists d. \\ \quad \quad \quad \operatorname{pos\text{-}rr}(d) \wedge \\ \quad \quad \quad \exists w \in \operatorname{opens}(t). \\ \quad \quad \quad \quad (x \in w) \wedge \forall s \in \mathbb{R}, a \in w.\; ((\lvert (s - r) \rvert < d) \Rightarrow (\operatorname{act}(m)(s, a) \in u)) \end{array}}}
\vnbprop{tvs-translate-open}{\vnbstmt{\begin{array}{@{}l@{}} \forall m, t, w, a. \\ \quad (\operatorname{is\text{-}tvs}(\operatorname{rr\text{-}normed\text{-}field}, m, t) \wedge ((w \in \operatorname{opens}(t)) \wedge (a \in \operatorname{vec}(m)))) \Rightarrow  \\ \quad \quad \{x \in \operatorname{vec}(m) : (\operatorname{vadd}(m)(x, \operatorname{act}(m)(-1, a)) \in w)\} \in \operatorname{opens}(t) \end{array}}}
\vnbprop{tvs-dilate-open}{\vnbstmt{\begin{array}{@{}l@{}} \forall m, t, c, w. \\ \quad (\operatorname{is\text{-}tvs}(\operatorname{rr\text{-}normed\text{-}field}, m, t) \wedge ((c \in \mathbb{R}) \wedge (w \in \operatorname{opens}(t)))) \Rightarrow  \\ \quad \quad \{x \in \operatorname{vec}(m) : (\operatorname{act}(m)(c, x) \in w)\} \in \operatorname{opens}(t) \end{array}}}
\vnbprop{fker-membership}{\vnbstmt{\forall m, f, z.\; ((z \in \operatorname{fker}(m, f)) \Leftrightarrow ((z \in \operatorname{vec}(m)) \wedge (f(z) = 0)))}}
\vnbprop{rr-nonzero-membership}{\vnbstmt{\forall v.\; ((v \in \{\mathit{nzv} \in \mathbb{R} : \neg (\mathit{nzv} = 0)\}) \Leftrightarrow ((v \in \mathbb{R}) \wedge \neg (v = 0)))}}
\vnbprop{tvs-functional-bounded-nbhd}{\vnbstmt{\begin{array}{@{}l@{}} \forall m, t, f, a, s. \\ \quad \operatorname{is\text{-}tvs}(\operatorname{rr\text{-}normed\text{-}field}, m, t) \wedge \\ \quad \operatorname{is\text{-}linear\text{-}functional}(m, f) \wedge \\ \quad \operatorname{complement\text{-}in}(\operatorname{pts}(t), \operatorname{fker}(m, f)) \in \operatorname{opens}(t) \wedge \\ \quad a \in \operatorname{vec}(m) \wedge \\ \quad s \in \mathbb{R} \wedge \\ \quad 0 < s \wedge \\ \quad f(a) = s \Rightarrow \\ \quad \quad \exists w \in \operatorname{opens}(t).\; ((\operatorname{vzero}(m) \in w) \wedge \forall z \in w.\; (\lvert f(z) \rvert \leq s)) \end{array}}}
\vnbprop{tvs-closed-kernel-continuous}{\vnbstmt{\begin{array}{@{}l@{}} \forall m, t, f. \\ \quad \operatorname{is\text{-}tvs}(\operatorname{rr\text{-}normed\text{-}field}, m, t) \wedge \\ \quad \operatorname{is\text{-}linear\text{-}functional}(m, f) \wedge \\ \quad \exists a \in \operatorname{vec}(m).\; \neg (f(a) = 0) \wedge \\ \quad \operatorname{is\text{-}top\text{-}closed}(t, \operatorname{fker}(m, f)) \Rightarrow \\ \quad \quad \operatorname{is\text{-}hom\text{-}top\text{-}space}(t, \\ \quad \quad \quad \operatorname{seminorm\text{-}top}(\operatorname{normed\text{-}vector\text{-}space\text{-}as\text{-}module}(\operatorname{rr\text{-}nvs}), \\ \quad \quad \quad \quad \{0\}, (\lambda k \in \{0\}.\; \operatorname{vnrm}(\operatorname{rr\text{-}nvs}))), \\ \quad \quad \quad f) \end{array}}}
\vnbprop{tvs-continuous-kernel-closed}{\vnbstmt{\begin{array}{@{}l@{}} \forall m, t, f. \\ \quad \operatorname{is\text{-}tvs}(\operatorname{rr\text{-}normed\text{-}field}, m, t) \wedge \\ \quad \operatorname{is\text{-}linear\text{-}functional}(m, f) \wedge \\ \quad \operatorname{is\text{-}hom\text{-}top\text{-}space}(t, \\ \quad \quad \operatorname{seminorm\text{-}top}(\operatorname{normed\text{-}vector\text{-}space\text{-}as\text{-}module}(\operatorname{rr\text{-}nvs}), \\ \quad \quad \quad \{0\}, (\lambda k \in \{0\}.\; \operatorname{vnrm}(\operatorname{rr\text{-}nvs}))), \\ \quad \quad f) \Rightarrow \\ \quad \quad \operatorname{is\text{-}top\text{-}closed}(t, \operatorname{fker}(m, f)) \end{array}}}
\vnbprop{functional-continuous-iff-kernel-closed}{\vnbstmt{\begin{array}{@{}l@{}} \forall m, t, f. \\ \quad \operatorname{is\text{-}tvs}(\operatorname{rr\text{-}normed\text{-}field}, m, t) \wedge \\ \quad \operatorname{is\text{-}linear\text{-}functional}(m, f) \wedge \\ \quad \exists a \in \operatorname{vec}(m).\; \neg (f(a) = 0) \Rightarrow \\ \quad \quad \operatorname{is\text{-}hom\text{-}top\text{-}space}(t, \\ \quad \quad \quad \operatorname{seminorm\text{-}top}(\operatorname{normed\text{-}vector\text{-}space\text{-}as\text{-}module}(\operatorname{rr\text{-}nvs}), \\ \quad \quad \quad \quad \{0\}, (\lambda k \in \{0\}.\; \operatorname{vnrm}(\operatorname{rr\text{-}nvs}))), \\ \quad \quad \quad f) \\ \quad \quad \Leftrightarrow\; \operatorname{is\text{-}top\text{-}closed}(t, \operatorname{fker}(m, f)) \end{array}}}

\vnbbookfile{structure-library/line-coefficient}{structure-library/line-coefficient}
\noindent{\small geo-vocab.scm -- the coefficient of a vector on a line (for the seed functional of the geometric Hahn-Banach, TB-2a).  Vocabulary only.  Agent DUAL-FIX, 2026-10-26 (NEW DEFINITION, the user's to confirm). DESTINATION: structure-library/line-coefficient.scm.  Library: analysis (LINE: linear-functional.scm).}\par
\par\medskip\noindent\textbf{Definition} (\texttt{line-coef}).\;$\operatorname{line\text{-}coef}(m, v, u) :=$ $(\iota\, \operatorname{lnt\_}.\;((\operatorname{lnt\_} \in \mathbb{R}) \wedge (u = \operatorname{act}(m)(\operatorname{lnt\_}, v))))$
\par\medskip\noindent\textbf{Definition} (\texttt{sumset}).\;$\operatorname{sumset}(m, a, p) :=$
\begin{equation*}
\{\operatorname{smx\_} \in \operatorname{vec}(m) : \exists \operatorname{smz\_} \in p.\; (\operatorname{vadd}(m)(\operatorname{smx\_}, \operatorname{act}(m)(-1, \operatorname{smz\_})) \in a)\}
\end{equation*}


\vnbbookfile{theorem-library/tvs-geometric-hb}{theorem-library/tvs-geometric-hb}
\noindent{\small geo.scm -- topological-vector-space.pdf TB-2a (geometric Hahn-Banach): Cor 2.15, Prop 2.13, Prop 2.12, Cor 2.14, by the ANALYTIC route -- the gauge of an open convex set and the sublinear Hahn-Banach -- instead of the notes' (maximal subspaces by Zorn, the closure of H, quotients E/H, connectedness of E \textbackslash{} \{0\}), which ...}\par
\vnbprop{tvs-shift-open}{\vnbstmt{\begin{array}{@{}l@{}} \forall m, t, a, b. \\ \quad (\operatorname{is\text{-}tvs}(\operatorname{rr\text{-}normed\text{-}field}, m, t) \wedge ((a \in \operatorname{opens}(t)) \wedge (b \in \operatorname{vec}(m)))) \Rightarrow  \\ \quad \quad \{y \in \operatorname{vec}(m) : (\operatorname{vadd}(m)(y, b) \in a)\} \in \operatorname{opens}(t) \end{array}}}
\vnbprop{convex-shift}{\vnbstmt{\begin{array}{@{}l@{}} \forall m, a, b. \\ \quad \operatorname{is\text{-}module}(m) \wedge \\ \quad \operatorname{scal}(m) = \operatorname{normed\text{-}field\text{-}as\text{-}commutative\text{-}ring}(\operatorname{rr\text{-}normed\text{-}field}) \wedge \\ \quad \operatorname{is\text{-}convex\text{-}subset}(m, a) \wedge \\ \quad b \in \operatorname{vec}(m) \Rightarrow \\ \quad \quad \operatorname{is\text{-}convex\text{-}subset}(m, \{y \in \operatorname{vec}(m) : (\operatorname{vadd}(m)(y, b) \in a)\}) \end{array}}}
\vnbprop{tvs-open-zero-absorbing}{\vnbstmt{\begin{array}{@{}l@{}} \forall m, t, w. \\ \quad (\operatorname{is\text{-}tvs}(\operatorname{rr\text{-}normed\text{-}field}, m, t) \wedge ((w \in \operatorname{opens}(t)) \wedge (\operatorname{vzero}(m) \in w))) \Rightarrow  \\ \quad \quad \operatorname{is\text{-}absorbing\text{-}subset}(m, w) \end{array}}}
\vnbprop{gauge-lt-one-on-open}{\vnbstmt{\begin{array}{@{}l@{}} \forall m, t, w, x. \\ \quad \operatorname{is\text{-}tvs}(\operatorname{rr\text{-}normed\text{-}field}, m, t) \wedge \\ \quad w \in \operatorname{opens}(t) \wedge \\ \quad \operatorname{is\text{-}convex\text{-}subset}(m, w) \wedge \\ \quad \operatorname{vzero}(m) \in w \wedge \\ \quad x \in w \Rightarrow \\ \quad \quad \operatorname{inf}(\operatorname{mink\text{-}set}(m, w, x)) < 1 \end{array}}}
\vnbprop{gauge-ge-one-off-set}{\vnbstmt{\begin{array}{@{}l@{}} \forall m, t, w, x. \\ \quad \operatorname{is\text{-}tvs}(\operatorname{rr\text{-}normed\text{-}field}, m, t) \wedge \\ \quad w \in \operatorname{opens}(t) \wedge \\ \quad \operatorname{is\text{-}convex\text{-}subset}(m, w) \wedge \\ \quad \operatorname{vzero}(m) \in w \wedge \\ \quad x \in \operatorname{vec}(m) \wedge \\ \quad \neg (x \in w) \Rightarrow \\ \quad \quad 1 \leq \operatorname{inf}(\operatorname{mink\text{-}set}(m, w, x)) \end{array}}}
\vnbprop{gauge-sublinear}{\vnbstmt{\begin{array}{@{}l@{}} \forall m, t, w. \\ \quad (\operatorname{is\text{-}tvs}(\operatorname{rr\text{-}normed\text{-}field}, m, t) \wedge ((w \in \operatorname{opens}(t)) \wedge (\operatorname{is\text{-}convex\text{-}subset}(m, w) \wedge (\operatorname{vzero}(m) \in w)))) \Rightarrow  \\ \quad \quad \operatorname{is\text{-}sublinear}(m, p_{w}) \end{array}}}
\vnbprop{rmod-act-cancel}{\vnbstmt{\begin{array}{@{}l@{}} \forall m, v, a, b. \\ \quad \operatorname{is\text{-}module}(m) \wedge \\ \quad \operatorname{scal}(m) = \operatorname{normed\text{-}field\text{-}as\text{-}commutative\text{-}ring}(\operatorname{rr\text{-}normed\text{-}field}) \wedge \\ \quad v \in \operatorname{vec}(m) \wedge \\ \quad \neg (v = \operatorname{vzero}(m)) \wedge \\ \quad a \in \mathbb{R} \wedge \\ \quad b \in \mathbb{R} \wedge \\ \quad \operatorname{act}(m)(a, v) = \operatorname{act}(m)(b, v) \Rightarrow \\ \quad \quad a = b \end{array}}}
\vnbprop{line-is-submodule-rmod}{\vnbstmt{\begin{array}{@{}l@{}} \forall m, v. \\ \quad (\operatorname{is\text{-}module}(m) \wedge ((\operatorname{scal}(m) = \operatorname{normed\text{-}field\text{-}as\text{-}commutative\text{-}ring}(\operatorname{rr\text{-}normed\text{-}field})) \wedge (v \in \operatorname{vec}(m)))) \Rightarrow  \\ \quad \quad \operatorname{is\text{-}submodule}(m, \operatorname{line}(m, v)) \end{array}}}
\vnbprop{line-coef-spec}{\vnbstmt{\begin{array}{@{}l@{}} \forall m, v, u. \\ \quad \operatorname{is\text{-}module}(m) \wedge \\ \quad \operatorname{scal}(m) = \operatorname{normed\text{-}field\text{-}as\text{-}commutative\text{-}ring}(\operatorname{rr\text{-}normed\text{-}field}) \wedge \\ \quad v \in \operatorname{vec}(m) \wedge \\ \quad \neg (v = \operatorname{vzero}(m)) \wedge \\ \quad u \in \operatorname{line}(m, v) \Rightarrow \\ \quad \quad (\operatorname{line\text{-}coef}(m, v, u) \in \mathbb{R}) \wedge (u = \operatorname{act}(m)(\operatorname{line\text{-}coef}(m, v, u), v)) \end{array}}}
\vnbprop{line-coef-value}{\vnbstmt{\begin{array}{@{}l@{}} \forall m, v, u, c. \\ \quad \operatorname{is\text{-}module}(m) \wedge \\ \quad \operatorname{scal}(m) = \operatorname{normed\text{-}field\text{-}as\text{-}commutative\text{-}ring}(\operatorname{rr\text{-}normed\text{-}field}) \wedge \\ \quad v \in \operatorname{vec}(m) \wedge \\ \quad \neg (v = \operatorname{vzero}(m)) \wedge \\ \quad c \in \mathbb{R} \wedge \\ \quad u = \operatorname{act}(m)(c, v) \Rightarrow \\ \quad \quad \operatorname{line\text{-}coef}(m, v, u) = c \end{array}}}
\vnbprop{line-seed-linear}{\vnbstmt{\begin{array}{@{}l@{}} \forall m, v. \\ \quad \operatorname{is\text{-}module}(m) \wedge \\ \quad \operatorname{scal}(m) = \operatorname{normed\text{-}field\text{-}as\text{-}commutative\text{-}ring}(\operatorname{rr\text{-}normed\text{-}field}) \wedge \\ \quad v \in \operatorname{vec}(m) \wedge \\ \quad \neg (v = \operatorname{vzero}(m)) \Rightarrow \\ \quad \quad \operatorname{is\text{-}linear\text{-}functional\text{-}on}(m, \operatorname{line}(m, v), (\lambda u \in \operatorname{line}(m, v).\; \operatorname{line\text{-}coef}(m, v, u))) \end{array}}}
\vnbprop{line-seed-at-v}{\vnbstmt{\begin{array}{@{}l@{}} \forall m, v. \\ \quad \operatorname{is\text{-}module}(m) \wedge \\ \quad \operatorname{scal}(m) = \operatorname{normed\text{-}field\text{-}as\text{-}commutative\text{-}ring}(\operatorname{rr\text{-}normed\text{-}field}) \wedge \\ \quad v \in \operatorname{vec}(m) \wedge \\ \quad \neg (v = \operatorname{vzero}(m)) \Rightarrow \\ \quad \quad (\lambda u \in \operatorname{line}(m, v).\; \operatorname{line\text{-}coef}(m, v, u))(v) = 1 \end{array}}}
\vnbprop{line-seed-below-gauge}{\vnbstmt{\begin{array}{@{}l@{}} \forall m, t, w, v.; \forall u \in \operatorname{line}(m, v). \\ \quad \operatorname{is\text{-}tvs}(\operatorname{rr\text{-}normed\text{-}field}, m, t) \wedge \\ \quad w \in \operatorname{opens}(t) \wedge \\ \quad \operatorname{is\text{-}convex\text{-}subset}(m, w) \wedge \\ \quad \operatorname{vzero}(m) \in w \wedge \\ \quad v \in \operatorname{vec}(m) \wedge \\ \quad \neg (v = \operatorname{vzero}(m)) \wedge \\ \quad 1 \leq \operatorname{inf}(\operatorname{mink\text{-}set}(m, w, v)) \Rightarrow \\ \quad \quad (\lambda \mathit{lnu} \in \operatorname{line}(m, v).\; \operatorname{line\text{-}coef}(m, v, \mathit{lnu}))(u) \leq p_{w}(u) \end{array}}}
\vnbprop{tvs-local-bound-of-one-sided}{\vnbstmt{\begin{array}{@{}l@{}} \forall m, t, f, n.; \forall s \in \mathbb{R}. \\ \quad \operatorname{is\text{-}tvs}(\operatorname{rr\text{-}normed\text{-}field}, m, t) \wedge \\ \quad \operatorname{is\text{-}linear\text{-}functional}(m, f) \wedge \\ \quad n \in \operatorname{opens}(t) \wedge \\ \quad \operatorname{vzero}(m) \in n \wedge \\ \quad \forall w \in n.\; (f(w) \leq 1) \wedge \\ \quad 0 < s \Rightarrow \\ \quad \quad \exists \mathit{bw} \in \operatorname{opens}(t).\; ((\operatorname{vzero}(m) \in \mathit{bw}) \wedge \forall z \in \mathit{bw}.\; (\lvert f(z) \rvert \leq s)) \end{array}}}
\vnbprop{tvs-continuous-of-local-bound}{\vnbstmt{\begin{array}{@{}l@{}} \forall m, t, f. \\ \quad \operatorname{is\text{-}tvs}(\operatorname{rr\text{-}normed\text{-}field}, m, t) \wedge \\ \quad \operatorname{is\text{-}linear\text{-}functional}(m, f) \wedge \\ \quad \forall s \in \mathbb{R}. \\ \quad \quad (0 < s) \Rightarrow  \\ \quad \quad \quad \exists w \in \operatorname{opens}(t).\; ((\operatorname{vzero}(m) \in w) \wedge \forall z \in w.\; (\lvert f(z) \rvert \leq s)) \Rightarrow \\ \quad \quad \operatorname{is\text{-}hom\text{-}top\text{-}space}(t, \\ \quad \quad \quad \operatorname{seminorm\text{-}top}(\operatorname{normed\text{-}vector\text{-}space\text{-}as\text{-}module}(\operatorname{rr\text{-}nvs}), \\ \quad \quad \quad \quad \{0\}, (\lambda k \in \{0\}.\; \operatorname{vnrm}(\operatorname{rr\text{-}nvs}))), \\ \quad \quad \quad f) \end{array}}}
\vnbprop{open-convex-positive-functional}{\vnbstmt{\begin{array}{@{}l@{}} \forall m, t, a. \\ \quad \operatorname{is\text{-}tvs}(\operatorname{rr\text{-}normed\text{-}field}, m, t) \wedge \\ \quad a \in \operatorname{opens}(t) \wedge \\ \quad \operatorname{is\text{-}convex\text{-}subset}(m, a) \wedge \\ \quad \exists a_{0}.\; (a_{0} \in a) \wedge \\ \quad \neg (\operatorname{vzero}(m) \in a) \Rightarrow \\ \quad \quad \exists f. \\ \quad \quad \quad \operatorname{is\text{-}linear\text{-}functional}(m, f) \wedge \\ \quad \quad \quad \operatorname{is\text{-}hom\text{-}top\text{-}space}(t, \\ \quad \quad \quad \quad \operatorname{seminorm\text{-}top}(\operatorname{normed\text{-}vector\text{-}space\text{-}as\text{-}module}(\operatorname{rr\text{-}nvs}), \\ \quad \quad \quad \quad \quad \{0\}, (\lambda k \in \{0\}.\; \operatorname{vnrm}(\operatorname{rr\text{-}nvs}))), \\ \quad \quad \quad \quad f) \wedge \\ \quad \quad \quad \forall x \in a.\; (0 < f(x)) \end{array}}}
\vnbprop{rmod-add-neg}{\vnbstmt{\begin{array}{@{}l@{}} \forall m, y. \\ \quad (\operatorname{is\text{-}module}(m) \wedge ((\operatorname{scal}(m) = \operatorname{normed\text{-}field\text{-}as\text{-}commutative\text{-}ring}(\operatorname{rr\text{-}normed\text{-}field})) \wedge (y \in \operatorname{vec}(m)))) \Rightarrow  \\ \quad \quad \operatorname{vadd}(m)(y, \operatorname{act}(m)(-1, y)) = \operatorname{vzero}(m) \end{array}}}
\vnbprop{rmod-combination-sub}{\vnbstmt{\begin{array}{@{}l@{}} \forall m, u, t, a, b, c, d. \\ \quad \operatorname{is\text{-}module}(m) \wedge \\ \quad \operatorname{scal}(m) = \operatorname{normed\text{-}field\text{-}as\text{-}commutative\text{-}ring}(\operatorname{rr\text{-}normed\text{-}field}) \wedge \\ \quad u \in \mathbb{R} \wedge \\ \quad t \in \mathbb{R} \wedge \\ \quad a \in \operatorname{vec}(m) \wedge \\ \quad b \in \operatorname{vec}(m) \wedge \\ \quad c \in \operatorname{vec}(m) \wedge \\ \quad d \in \operatorname{vec}(m) \Rightarrow \\ \quad \quad \operatorname{vadd}(m)(\operatorname{vadd}(m)(\operatorname{act}(m)(u, a), \operatorname{act}(m)(t, b)), \\ \quad \quad \quad \operatorname{act}(m)(-1, \operatorname{vadd}(m)(\operatorname{act}(m)(u, c), \operatorname{act}(m)(t, d)))) \\ \quad \quad =\; \operatorname{vadd}(m)(\operatorname{act}(m)(u, \operatorname{vadd}(m)(a, \operatorname{act}(m)(-1, c))), \operatorname{act}(m)(t, \operatorname{vadd}(m)(b, \operatorname{act}(m)(-1, d)))) \end{array}}}
\vnbprop{fker-submodule}{\vnbstmt{\begin{array}{@{}l@{}} \forall m, f. \\ \quad \operatorname{is\text{-}module}(m) \wedge \\ \quad \operatorname{scal}(m) = \operatorname{normed\text{-}field\text{-}as\text{-}commutative\text{-}ring}(\operatorname{rr\text{-}normed\text{-}field}) \wedge \\ \quad \operatorname{is\text{-}linear\text{-}functional}(m, f) \Rightarrow \\ \quad \quad \operatorname{is\text{-}submodule}(m, \operatorname{fker}(m, f)) \end{array}}}
\vnbprop{closed-hyperplane-through-subspace}{\vnbstmt{\begin{array}{@{}l@{}} \forall m, t, a, s. \\ \quad \operatorname{is\text{-}tvs}(\operatorname{rr\text{-}normed\text{-}field}, m, t) \wedge \\ \quad a \in \operatorname{opens}(t) \wedge \\ \quad \operatorname{is\text{-}convex\text{-}subset}(m, a) \wedge \\ \quad \exists a_{0}.\; (a_{0} \in a) \wedge \\ \quad \operatorname{is\text{-}submodule}(m, s) \wedge \\ \quad \forall z \in s.\; \neg (z \in a) \Rightarrow \\ \quad \quad \exists g. \\ \quad \quad \quad \operatorname{is\text{-}linear\text{-}functional}(m, g) \wedge \\ \quad \quad \quad \operatorname{is\text{-}hom\text{-}top\text{-}space}(t, \\ \quad \quad \quad \quad \operatorname{seminorm\text{-}top}(\operatorname{normed\text{-}vector\text{-}space\text{-}as\text{-}module}(\operatorname{rr\text{-}nvs}), \\ \quad \quad \quad \quad \quad \{0\}, (\lambda k \in \{0\}.\; \operatorname{vnrm}(\operatorname{rr\text{-}nvs}))), \\ \quad \quad \quad \quad g) \wedge \\ \quad \quad \quad \exists \mathit{gx} \in \operatorname{vec}(m).\; \neg (g(\mathit{gx}) = 0) \wedge \\ \quad \quad \quad s \subseteq \operatorname{fker}(m, g) \wedge \\ \quad \quad \quad \operatorname{is\text{-}top\text{-}closed}(t, \operatorname{fker}(m, g)) \wedge \\ \quad \quad \quad \forall x \in a.\; \neg (x \in \operatorname{fker}(m, g)) \end{array}}}
\vnbprop{open-cone-misses-subspace}{\vnbstmt{\begin{array}{@{}l@{}} \forall m, t, c. \\ \quad \operatorname{is\text{-}tvs}(\operatorname{rr\text{-}normed\text{-}field}, m, t) \wedge \\ \quad c \in \operatorname{opens}(t) \wedge \\ \quad \operatorname{is\text{-}convex\text{-}subset}(m, c) \wedge \\ \quad \forall x \in c, r \in \mathbb{R}.\; ((0 < r) \Rightarrow (\operatorname{act}(m)(r, x) \in c)) \wedge \\ \quad \exists e \in \operatorname{vec}(m).\; \neg (e \in c) \wedge \\ \quad \exists d_{1}, d_{2} \in \operatorname{vec}(m). \forall a, b \in \mathbb{R}. \\ \quad \quad (\operatorname{vadd}(m)(\operatorname{act}(m)(a, d_{1}), \operatorname{act}(m)(b, d_{2})) = \operatorname{vzero}(m)) \Rightarrow ((a = 0) \wedge (b = 0)) \Rightarrow \\ \quad \quad \exists s. \\ \quad \quad \quad (\operatorname{is\text{-}submodule}(m, s) \wedge (\exists v \in s.\; \neg (v = \operatorname{vzero}(m)) \wedge \forall \mathit{vx} \in s.\; \neg (\mathit{vx} \in c))) \end{array}}}
\vnbprop{sumset-unfold}{\vnbstmt{\begin{array}{@{}l@{}} \forall m, a, p. \\ \quad (\operatorname{sumset}(m, a, p) \simeq \{x \in \operatorname{vec}(m) : \exists z \in p.\; (\operatorname{vadd}(m)(x, \operatorname{act}(m)(-1, z)) \in a)\}) \end{array}}}
\vnbprop{sumset-open}{\vnbstmt{\begin{array}{@{}l@{}} \forall m, t, a, p. \\ \quad (\operatorname{is\text{-}tvs}(\operatorname{rr\text{-}normed\text{-}field}, m, t) \wedge ((a \in \operatorname{opens}(t)) \wedge (p \subseteq \operatorname{vec}(m)))) \Rightarrow  \\ \quad \quad \operatorname{sumset}(m, a, p) \in \operatorname{opens}(t) \end{array}}}
\vnbprop{sumset-convex}{\vnbstmt{\begin{array}{@{}l@{}} \forall m, a, p. \\ \quad \operatorname{is\text{-}module}(m) \wedge \\ \quad \operatorname{scal}(m) = \operatorname{normed\text{-}field\text{-}as\text{-}commutative\text{-}ring}(\operatorname{rr\text{-}normed\text{-}field}) \wedge \\ \quad \operatorname{is\text{-}convex\text{-}subset}(m, a) \wedge \\ \quad \operatorname{is\text{-}convex\text{-}subset}(m, p) \Rightarrow \\ \quad \quad \operatorname{is\text{-}convex\text{-}subset}(m, \operatorname{sumset}(m, a, p)) \end{array}}}
\vnbprop{kernel-on-submodule}{\vnbstmt{\begin{array}{@{}l@{}} \forall m, s, f. \\ \quad \operatorname{is\text{-}module}(m) \wedge \\ \quad \operatorname{scal}(m) = \operatorname{normed\text{-}field\text{-}as\text{-}commutative\text{-}ring}(\operatorname{rr\text{-}normed\text{-}field}) \wedge \\ \quad \operatorname{is\text{-}submodule}(m, s) \wedge \\ \quad \operatorname{is\text{-}linear\text{-}functional\text{-}on}(m, s, f) \Rightarrow \\ \quad \quad \operatorname{is\text{-}submodule}(m, \{z \in s : (f(z) = 0)\}) \end{array}}}
\vnbprop{positive-part-convex}{\vnbstmt{\begin{array}{@{}l@{}} \forall m, s, f. \\ \quad \operatorname{is\text{-}module}(m) \wedge \\ \quad \operatorname{scal}(m) = \operatorname{normed\text{-}field\text{-}as\text{-}commutative\text{-}ring}(\operatorname{rr\text{-}normed\text{-}field}) \wedge \\ \quad \operatorname{is\text{-}submodule}(m, s) \wedge \\ \quad \operatorname{is\text{-}linear\text{-}functional\text{-}on}(m, s, f) \Rightarrow \\ \quad \quad \operatorname{is\text{-}convex\text{-}subset}(m, \{z \in s : (0 < f(z))\}) \end{array}}}
\vnbprop{extend-positive-functional}{\vnbstmt{\begin{array}{@{}l@{}} \forall m, t, a, s, f. \\ \quad \operatorname{is\text{-}tvs}(\operatorname{rr\text{-}normed\text{-}field}, m, t) \wedge \\ \quad a \in \operatorname{opens}(t) \wedge \\ \quad \operatorname{is\text{-}convex\text{-}subset}(m, a) \wedge \\ \quad \operatorname{is\text{-}submodule}(m, s) \wedge \\ \quad \operatorname{is\text{-}linear\text{-}functional\text{-}on}(m, s, f) \wedge \\ \quad \exists \mathit{af} \in a.\; (\mathit{af} \in s) \wedge \\ \quad \forall z \in a.\; ((z \in s) \Rightarrow (0 < f(z))) \Rightarrow \\ \quad \quad \exists g. \\ \quad \quad \quad \operatorname{is\text{-}linear\text{-}functional}(m, g) \wedge \\ \quad \quad \quad \operatorname{is\text{-}hom\text{-}top\text{-}space}(t, \\ \quad \quad \quad \quad \operatorname{seminorm\text{-}top}(\operatorname{normed\text{-}vector\text{-}space\text{-}as\text{-}module}(\operatorname{rr\text{-}nvs}), \\ \quad \quad \quad \quad \quad \{0\}, (\lambda k \in \{0\}.\; \operatorname{vnrm}(\operatorname{rr\text{-}nvs}))), \\ \quad \quad \quad \quad g) \wedge \\ \quad \quad \quad \operatorname{extends\text{-}on}(s, g, f) \wedge \\ \quad \quad \quad \forall x \in a.\; (0 < g(x)) \end{array}}}
\vnbprop{positive-functional-vanishing-on-subspace}{\vnbstmt{\begin{array}{@{}l@{}} \forall m, t, a, s. \\ \quad \operatorname{is\text{-}tvs}(\operatorname{rr\text{-}normed\text{-}field}, m, t) \wedge \\ \quad a \in \operatorname{opens}(t) \wedge \\ \quad \operatorname{is\text{-}convex\text{-}subset}(m, a) \wedge \\ \quad \exists a_{0}.\; (a_{0} \in a) \wedge \\ \quad \operatorname{is\text{-}submodule}(m, s) \wedge \\ \quad \forall z \in s.\; \neg (z \in a) \Rightarrow \\ \quad \quad \exists g. \\ \quad \quad \quad \operatorname{is\text{-}linear\text{-}functional}(m, g) \wedge \\ \quad \quad \quad \operatorname{is\text{-}hom\text{-}top\text{-}space}(t, \\ \quad \quad \quad \quad \operatorname{seminorm\text{-}top}(\operatorname{normed\text{-}vector\text{-}space\text{-}as\text{-}module}(\operatorname{rr\text{-}nvs}), \\ \quad \quad \quad \quad \quad \{0\}, (\lambda k \in \{0\}.\; \operatorname{vnrm}(\operatorname{rr\text{-}nvs}))), \\ \quad \quad \quad \quad g) \wedge \\ \quad \quad \quad s \subseteq \operatorname{fker}(m, g) \wedge \\ \quad \quad \quad \forall x \in a.\; (0 < g(x)) \end{array}}}

\vnbbookfile{structure-library/tvs-product-line}{structure-library/tvs-product-line}
\noindent{\small prod-vocab.scm -- THE PRODUCT E x R OF A REAL MODULE AND THE REALS, AND THE PRODUCT TOPOLOGY (definitions only).  Agent TVS-NEXT (worker-08), 2026-10-26. DESTINATION: structure-library/tvs-product-line.scm.  Library: analysis.}\par
\par\medskip\noindent\textbf{Definition} (\texttt{prod-rmod}).\;$\operatorname{prod\text{-}rmod}(m) :=$
\begin{equation*}
{[}\operatorname{scal}(m), (\operatorname{vec}(m) \times \mathbb{R}), \operatorname{pv\text{-}vadd}(m), {[}\operatorname{vzero}(m), 0], \operatorname{pv\text{-}vneg}(m), \operatorname{pv\text{-}act}(m)]
\end{equation*}

\par\medskip\noindent\textbf{Definition} (\texttt{pt-box}).\;$\operatorname{pt\text{-}box}(t, v, c, e) :=$
\begin{equation*}
\{\operatorname{ptz\_} \in (\operatorname{pts}(t) \times \mathbb{R}) : ((\operatorname{nth}(1, \operatorname{ptz\_}) \in v) \wedge (\lvert (\operatorname{nth}(2, \operatorname{ptz\_}) - c) \rvert < e))\}
\end{equation*}

\par\medskip\noindent\textbf{Definition} (\texttt{is-pt-open}).\;For all $t$ and $u$, 
\begin{equation*}
\begin{array}{@{}l@{}} \operatorname{is\text{-}pt\text{-}open}(t, u) \\ \Leftrightarrow\; u \subseteq (\operatorname{pts}(t) \times \mathbb{R}) \wedge \\ \forall \operatorname{ptw\_} \in u. \exists \operatorname{ptv\_} \in \operatorname{opens}(t). \\ \quad \operatorname{nth}(1, \operatorname{ptw\_}) \in \operatorname{ptv\_} \wedge \\ \quad \exists \operatorname{pte\_}.\; (\operatorname{pos\text{-}rr}(\operatorname{pte\_}) \wedge (\operatorname{pt\text{-}box}(t, \operatorname{ptv\_}, \operatorname{nth}(2, \operatorname{ptw\_}), \operatorname{pte\_}) \subseteq u)) \end{array}
\end{equation*}

\par\medskip\noindent\textbf{Definition} (\texttt{prod-top}).\;$\operatorname{prod\text{-}top}(t) :=$
\begin{equation*}
{[}(\operatorname{pts}(t) \times \mathbb{R}), \{\operatorname{ptu\_} \in \operatorname{power}((\operatorname{pts}(t) \times \mathbb{R})) : \operatorname{is\text{-}pt\text{-}open}(t, \operatorname{ptu\_})\}]
\end{equation*}

\par\medskip\noindent\textbf{Definition} (\texttt{pv-vadd}).\;$\operatorname{pv\text{-}vadd}(m) :=$
\begin{equation*}
\begin{array}{@{}l@{}} (\lambda {[}\operatorname{pva\_}, \operatorname{pvb\_}] \in ((\operatorname{vec}(m) \times \mathbb{R}) \times (\operatorname{vec}(m) \times \mathbb{R})).\; \\ \quad {[}\operatorname{vadd}(m)(\operatorname{nth}(1, \operatorname{pva\_}), \operatorname{nth}(1, \operatorname{pvb\_})), (\operatorname{nth}(2, \operatorname{pva\_}) + \operatorname{nth}(2, \operatorname{pvb\_}))]) \end{array}
\end{equation*}

\par\medskip\noindent\textbf{Definition} (\texttt{pv-vneg}).\;$\operatorname{pv\text{-}vneg}(m) :=$
\begin{equation*}
(\lambda \operatorname{pvc\_} \in (\operatorname{vec}(m) \times \mathbb{R}).\; {[}\operatorname{vneg}(m)(\operatorname{nth}(1, \operatorname{pvc\_})), -\operatorname{nth}(2, \operatorname{pvc\_})])
\end{equation*}

\par\medskip\noindent\textbf{Definition} (\texttt{pv-act}).\;$\operatorname{pv\text{-}act}(m) :=$
\begin{equation*}
\begin{array}{@{}l@{}} (\lambda {[}\operatorname{pvr\_}, \operatorname{pvc\_}] \in (\operatorname{carr}(\operatorname{scal}(m)) \times (\operatorname{vec}(m) \times \mathbb{R})).\; \\ \quad {[}\operatorname{act}(m)(\operatorname{pvr\_}, \operatorname{nth}(1, \operatorname{pvc\_})), (\operatorname{pvr\_} \cdot \operatorname{nth}(2, \operatorname{pvc\_}))]) \end{array}
\end{equation*}

\par\medskip\noindent\textbf{Definition} (\texttt{strict-epigraph}).\;$\operatorname{strict\text{-}epigraph}(m, p) :=$
\begin{equation*}
\{\operatorname{sepw\_} \in (\operatorname{vec}(m) \times \mathbb{R}) : (p(\operatorname{nth}(1, \operatorname{sepw\_})) < \operatorname{nth}(2, \operatorname{sepw\_}))\}
\end{equation*}


\vnbbookfile{theorem-library/tvs-product-line}{theorem-library/tvs-product-line}
\noindent{\small prod.scm -- THE PRODUCT TOPOLOGICAL VECTOR SPACE E x R, and topological-vector-space.pdf Prop 2.35 (the strict epigraph of a continuous convex functional is open and convex). Agent TVS-NEXT (worker-08), 2026-10-26. DESTINATION: theorem-library/tvs-product-line.scm.  Library: analysis.  Helper prefix: ptv-.}\par
\vnbprop{pv-vadd-value}{\vnbstmt{\begin{array}{@{}l@{}} \forall m, x, \mathit{qx}. \\ \quad \operatorname{is\text{-}module}(m) \wedge \\ \quad \operatorname{scal}(m) = \operatorname{normed\text{-}field\text{-}as\text{-}commutative\text{-}ring}(\operatorname{rr\text{-}normed\text{-}field}) \wedge \\ \quad x \in (\operatorname{vec}(m) \times \mathbb{R}) \wedge \\ \quad \mathit{qx} \in (\operatorname{vec}(m) \times \mathbb{R}) \Rightarrow \\ \quad \quad \operatorname{pv\text{-}vadd}(m)(x, \mathit{qx}) = {[}\operatorname{vadd}(m)(\operatorname{nth}(1, x), \operatorname{nth}(1, \mathit{qx})), (\operatorname{nth}(2, x) + \operatorname{nth}(2, \mathit{qx}))] \end{array}}}
\vnbprop{pv-vneg-value}{\vnbstmt{\begin{array}{@{}l@{}} \forall m.; \forall x \in (\operatorname{vec}(m) \times \mathbb{R}). \\ \quad (\operatorname{is\text{-}module}(m) \wedge (\operatorname{scal}(m) = \operatorname{normed\text{-}field\text{-}as\text{-}commutative\text{-}ring}(\operatorname{rr\text{-}normed\text{-}field}))) \Rightarrow  \\ \quad \quad \operatorname{pv\text{-}vneg}(m)(x) = {[}\operatorname{vneg}(m)(\operatorname{nth}(1, x)), -\operatorname{nth}(2, x)] \end{array}}}
\vnbprop{pv-act-value}{\vnbstmt{\begin{array}{@{}l@{}} \forall m, x, \mathit{px}. \\ \quad \operatorname{is\text{-}module}(m) \wedge \\ \quad \operatorname{scal}(m) = \operatorname{normed\text{-}field\text{-}as\text{-}commutative\text{-}ring}(\operatorname{rr\text{-}normed\text{-}field}) \wedge \\ \quad x \in \mathbb{R} \wedge \\ \quad \mathit{px} \in (\operatorname{vec}(m) \times \mathbb{R}) \Rightarrow \\ \quad \quad \operatorname{pv\text{-}act}(m)(x, \mathit{px}) = {[}\operatorname{act}(m)(x, \operatorname{nth}(1, \mathit{px})), (x \cdot \operatorname{nth}(2, \mathit{px}))] \end{array}}}
\vnbprop{prod-rmod-scal}{\vnbstmt{\begin{array}{@{}l@{}} \forall m. \\ \quad (\operatorname{is\text{-}module}(m) \wedge (\operatorname{scal}(m) = \operatorname{normed\text{-}field\text{-}as\text{-}commutative\text{-}ring}(\operatorname{rr\text{-}normed\text{-}field}))) \Rightarrow  \\ \quad \quad (\operatorname{scal}(\operatorname{prod\text{-}rmod}(m)) \simeq \operatorname{scal}(m)) \end{array}}}
\vnbprop{prod-rmod-vec}{\vnbstmt{\begin{array}{@{}l@{}} \forall m. \\ \quad (\operatorname{is\text{-}module}(m) \wedge (\operatorname{scal}(m) = \operatorname{normed\text{-}field\text{-}as\text{-}commutative\text{-}ring}(\operatorname{rr\text{-}normed\text{-}field}))) \Rightarrow  \\ \quad \quad (\operatorname{vec}(\operatorname{prod\text{-}rmod}(m)) \simeq (\operatorname{vec}(m) \times \mathbb{R})) \end{array}}}
\vnbprop{prod-rmod-vadd}{\vnbstmt{\begin{array}{@{}l@{}} \forall m. \\ \quad (\operatorname{is\text{-}module}(m) \wedge (\operatorname{scal}(m) = \operatorname{normed\text{-}field\text{-}as\text{-}commutative\text{-}ring}(\operatorname{rr\text{-}normed\text{-}field}))) \Rightarrow  \\ \quad \quad (\operatorname{vadd}(\operatorname{prod\text{-}rmod}(m)) \simeq \operatorname{pv\text{-}vadd}(m)) \end{array}}}
\vnbprop{prod-rmod-vzero}{\vnbstmt{\begin{array}{@{}l@{}} \forall m. \\ \quad (\operatorname{is\text{-}module}(m) \wedge (\operatorname{scal}(m) = \operatorname{normed\text{-}field\text{-}as\text{-}commutative\text{-}ring}(\operatorname{rr\text{-}normed\text{-}field}))) \Rightarrow  \\ \quad \quad (\operatorname{vzero}(\operatorname{prod\text{-}rmod}(m)) \simeq {[}\operatorname{vzero}(m), 0]) \end{array}}}
\vnbprop{prod-rmod-vneg}{\vnbstmt{\begin{array}{@{}l@{}} \forall m. \\ \quad (\operatorname{is\text{-}module}(m) \wedge (\operatorname{scal}(m) = \operatorname{normed\text{-}field\text{-}as\text{-}commutative\text{-}ring}(\operatorname{rr\text{-}normed\text{-}field}))) \Rightarrow  \\ \quad \quad (\operatorname{vneg}(\operatorname{prod\text{-}rmod}(m)) \simeq \operatorname{pv\text{-}vneg}(m)) \end{array}}}
\vnbprop{prod-rmod-act}{\vnbstmt{\begin{array}{@{}l@{}} \forall m. \\ \quad (\operatorname{is\text{-}module}(m) \wedge (\operatorname{scal}(m) = \operatorname{normed\text{-}field\text{-}as\text{-}commutative\text{-}ring}(\operatorname{rr\text{-}normed\text{-}field}))) \Rightarrow  \\ \quad \quad (\operatorname{act}(\operatorname{prod\text{-}rmod}(m)) \simeq \operatorname{pv\text{-}act}(m)) \end{array}}}
\vnbprop{pv-vadd-fun}{\vnbstmt{\begin{array}{@{}l@{}} \forall m. \\ \quad (\operatorname{is\text{-}module}(m) \wedge (\operatorname{scal}(m) = \operatorname{normed\text{-}field\text{-}as\text{-}commutative\text{-}ring}(\operatorname{rr\text{-}normed\text{-}field}))) \Rightarrow  \\ \quad \quad \operatorname{pv\text{-}vadd}(m) \in (((\operatorname{vec}(m) \times \mathbb{R}) \times (\operatorname{vec}(m) \times \mathbb{R})) \to (\operatorname{vec}(m) \times \mathbb{R})) \end{array}}}
\vnbprop{pv-vneg-fun}{\vnbstmt{\begin{array}{@{}l@{}} \forall m. \\ \quad (\operatorname{is\text{-}module}(m) \wedge (\operatorname{scal}(m) = \operatorname{normed\text{-}field\text{-}as\text{-}commutative\text{-}ring}(\operatorname{rr\text{-}normed\text{-}field}))) \Rightarrow  \\ \quad \quad \operatorname{pv\text{-}vneg}(m) \in ((\operatorname{vec}(m) \times \mathbb{R}) \to (\operatorname{vec}(m) \times \mathbb{R})) \end{array}}}
\vnbprop{pv-act-fun}{\vnbstmt{\begin{array}{@{}l@{}} \forall m. \\ \quad (\operatorname{is\text{-}module}(m) \wedge (\operatorname{scal}(m) = \operatorname{normed\text{-}field\text{-}as\text{-}commutative\text{-}ring}(\operatorname{rr\text{-}normed\text{-}field}))) \Rightarrow  \\ \quad \quad \operatorname{pv\text{-}act}(m) \in ((\operatorname{carr}(\operatorname{scal}(m)) \times (\operatorname{vec}(m) \times \mathbb{R})) \to (\operatorname{vec}(m) \times \mathbb{R})) \end{array}}}
\vnbprop{prod-rmod-assoc}{\vnbstmt{\begin{array}{@{}l@{}} \forall m. \\ \quad (\operatorname{is\text{-}module}(m) \wedge (\operatorname{scal}(m) = \operatorname{normed\text{-}field\text{-}as\text{-}commutative\text{-}ring}(\operatorname{rr\text{-}normed\text{-}field}))) \Rightarrow  \\ \quad \quad \operatorname{is\text{-}associative}(\operatorname{vadd}(\operatorname{prod\text{-}rmod}(m)), \operatorname{vec}(\operatorname{prod\text{-}rmod}(m))) \end{array}}}
\vnbprop{prod-rmod-comm}{\vnbstmt{\begin{array}{@{}l@{}} \forall m. \\ \quad (\operatorname{is\text{-}module}(m) \wedge (\operatorname{scal}(m) = \operatorname{normed\text{-}field\text{-}as\text{-}commutative\text{-}ring}(\operatorname{rr\text{-}normed\text{-}field}))) \Rightarrow  \\ \quad \quad \operatorname{is\text{-}commutative}(\operatorname{vadd}(\operatorname{prod\text{-}rmod}(m)), \operatorname{vec}(\operatorname{prod\text{-}rmod}(m))) \end{array}}}
\vnbprop{prod-rmod-identity}{\vnbstmt{\begin{array}{@{}l@{}} \forall m. \\ \quad (\operatorname{is\text{-}module}(m) \wedge (\operatorname{scal}(m) = \operatorname{normed\text{-}field\text{-}as\text{-}commutative\text{-}ring}(\operatorname{rr\text{-}normed\text{-}field}))) \Rightarrow  \\ \quad \quad \operatorname{is\text{-}identity}(\operatorname{vadd}(\operatorname{prod\text{-}rmod}(m)), \operatorname{vzero}(\operatorname{prod\text{-}rmod}(m)), \operatorname{vec}(\operatorname{prod\text{-}rmod}(m))) \end{array}}}
\vnbprop{prod-rmod-inverses}{\vnbstmt{\begin{array}{@{}l@{}} \forall m. \\ \quad (\operatorname{is\text{-}module}(m) \wedge (\operatorname{scal}(m) = \operatorname{normed\text{-}field\text{-}as\text{-}commutative\text{-}ring}(\operatorname{rr\text{-}normed\text{-}field}))) \Rightarrow  \\ \quad \quad \operatorname{has\text{-}inverses}(\operatorname{vadd}(\operatorname{prod\text{-}rmod}(m)), \operatorname{vzero}(\operatorname{prod\text{-}rmod}(m)), \operatorname{vneg}(\operatorname{prod\text{-}rmod}(m)), \operatorname{vec}(\operatorname{prod\text{-}rmod}(m))) \end{array}}}
\vnbprop{prod-rmod-act-distrib-vec}{\vnbstmt{\begin{array}{@{}l@{}} \forall m.; \forall r \in \operatorname{carr}(\operatorname{scal}(\operatorname{prod\text{-}rmod}(m))), x, y \in \operatorname{vec}(\operatorname{prod\text{-}rmod}(m)). \\ \quad (\operatorname{is\text{-}module}(m) \wedge (\operatorname{scal}(m) = \operatorname{normed\text{-}field\text{-}as\text{-}commutative\text{-}ring}(\operatorname{rr\text{-}normed\text{-}field}))) \Rightarrow  \\ \quad \quad \operatorname{act}(\operatorname{prod\text{-}rmod}(m))(r, \operatorname{vadd}(\operatorname{prod\text{-}rmod}(m))(x, y)) \\ \quad \quad =\; \operatorname{vadd}(\operatorname{prod\text{-}rmod}(m))(\operatorname{act}(\operatorname{prod\text{-}rmod}(m))(r, x), \operatorname{act}(\operatorname{prod\text{-}rmod}(m))(r, y)) \end{array}}}
\vnbprop{prod-rmod-act-distrib-scalar}{\vnbstmt{\begin{array}{@{}l@{}} \forall m.; \forall r, s \in \operatorname{carr}(\operatorname{scal}(\operatorname{prod\text{-}rmod}(m))), x \in \operatorname{vec}(\operatorname{prod\text{-}rmod}(m)). \\ \quad (\operatorname{is\text{-}module}(m) \wedge (\operatorname{scal}(m) = \operatorname{normed\text{-}field\text{-}as\text{-}commutative\text{-}ring}(\operatorname{rr\text{-}normed\text{-}field}))) \Rightarrow  \\ \quad \quad \operatorname{act}(\operatorname{prod\text{-}rmod}(m))(\operatorname{add}(\operatorname{scal}(\operatorname{prod\text{-}rmod}(m)))(r, s), x) \\ \quad \quad =\; \operatorname{vadd}(\operatorname{prod\text{-}rmod}(m))(\operatorname{act}(\operatorname{prod\text{-}rmod}(m))(r, x), \operatorname{act}(\operatorname{prod\text{-}rmod}(m))(s, x)) \end{array}}}
\vnbprop{prod-rmod-act-mul-compat}{\vnbstmt{\begin{array}{@{}l@{}} \forall m.; \forall r, s \in \operatorname{carr}(\operatorname{scal}(\operatorname{prod\text{-}rmod}(m))), x \in \operatorname{vec}(\operatorname{prod\text{-}rmod}(m)). \\ \quad (\operatorname{is\text{-}module}(m) \wedge (\operatorname{scal}(m) = \operatorname{normed\text{-}field\text{-}as\text{-}commutative\text{-}ring}(\operatorname{rr\text{-}normed\text{-}field}))) \Rightarrow  \\ \quad \quad \operatorname{act}(\operatorname{prod\text{-}rmod}(m))(\operatorname{mul}(\operatorname{scal}(\operatorname{prod\text{-}rmod}(m)))(r, s), x) \\ \quad \quad =\; \operatorname{act}(\operatorname{prod\text{-}rmod}(m))(r, \operatorname{act}(\operatorname{prod\text{-}rmod}(m))(s, x)) \end{array}}}
\vnbprop{prod-rmod-act-unital}{\vnbstmt{\begin{array}{@{}l@{}} \forall m.; \forall x \in \operatorname{vec}(\operatorname{prod\text{-}rmod}(m)). \\ \quad (\operatorname{is\text{-}module}(m) \wedge (\operatorname{scal}(m) = \operatorname{normed\text{-}field\text{-}as\text{-}commutative\text{-}ring}(\operatorname{rr\text{-}normed\text{-}field}))) \Rightarrow  \\ \quad \quad \operatorname{act}(\operatorname{prod\text{-}rmod}(m))(\operatorname{one}(\operatorname{scal}(\operatorname{prod\text{-}rmod}(m))), x) = x \end{array}}}
\vnbprop{prod-rmod-is-module}{\vnbstmt{\begin{array}{@{}l@{}} \forall m. \\ \quad (\operatorname{is\text{-}module}(m) \wedge (\operatorname{scal}(m) = \operatorname{normed\text{-}field\text{-}as\text{-}commutative\text{-}ring}(\operatorname{rr\text{-}normed\text{-}field}))) \Rightarrow  \\ \quad \quad \operatorname{is\text{-}module}(\operatorname{prod\text{-}rmod}(m)) \end{array}}}
\vnbprop{pt-box-unfold}{\vnbstmt{\begin{array}{@{}l@{}} \forall t, v, c, e. \\ \quad (\operatorname{pt\text{-}box}(t, v, c, e) \simeq \{z \in (\operatorname{pts}(t) \times \mathbb{R}) : ((\operatorname{nth}(1, z) \in v) \wedge (\lvert (\operatorname{nth}(2, z) - c) \rvert < e))\}) \end{array}}}
\vnbprop{pt-box-membership}{\vnbstmt{\begin{array}{@{}l@{}} \forall t, v, c, e, z. \\ \quad (z \in \operatorname{pt\text{-}box}(t, v, c, e)) \Leftrightarrow ((z \in (\operatorname{pts}(t) \times \mathbb{R})) \wedge ((\operatorname{nth}(1, z) \in v) \wedge (\lvert (\operatorname{nth}(2, z) - c) \rvert < e))) \end{array}}}
\vnbprop{pt-open-empty}{\vnbstmt{\forall t.\; (\operatorname{is\text{-}top\text{-}space}(t) \Rightarrow \operatorname{is\text{-}pt\text{-}open}(t, \emptyset))}}
\vnbprop{pt-open-whole}{\vnbstmt{\forall t.\; (\operatorname{is\text{-}top\text{-}space}(t) \Rightarrow \operatorname{is\text{-}pt\text{-}open}(t, (\operatorname{pts}(t) \times \mathbb{R})))}}
\vnbprop{pt-box-mono}{\vnbstmt{\begin{array}{@{}l@{}} \forall t, a, b, c, d, f. \\ \quad (\operatorname{is\text{-}top\text{-}space}(t) \wedge ((a \subseteq b) \wedge ((c \in \mathbb{R}) \wedge ((d \in \mathbb{R}) \wedge ((f \in \mathbb{R}) \wedge (d \leq f)))))) \Rightarrow  \\ \quad \quad \operatorname{pt\text{-}box}(t, a, c, d) \subseteq \operatorname{pt\text{-}box}(t, b, c, f) \end{array}}}
\vnbprop{pt-open-inter}{\vnbstmt{\begin{array}{@{}l@{}} \forall t, u, v. \\ \quad (\operatorname{is\text{-}top\text{-}space}(t) \wedge (\operatorname{is\text{-}pt\text{-}open}(t, u) \wedge \operatorname{is\text{-}pt\text{-}open}(t, v))) \Rightarrow  \\ \quad \quad \operatorname{is\text{-}pt\text{-}open}(t, (u \cap v)) \end{array}}}
\vnbprop{pt-open-union}{\vnbstmt{\begin{array}{@{}l@{}} \forall t, m. \\ \quad (\operatorname{is\text{-}top\text{-}space}(t) \wedge \forall w \in m.\; \operatorname{is\text{-}pt\text{-}open}(t, w)) \Rightarrow  \\ \quad \quad \operatorname{is\text{-}pt\text{-}open}(t, \operatorname{big\text{-}union}(u, m, u)) \end{array}}}
\vnbprop{prod-top-pts}{\vnbstmt{\forall t.\; (\operatorname{is\text{-}top\text{-}space}(t) \Rightarrow (\operatorname{pts}(\operatorname{prod\text{-}top}(t)) \simeq (\operatorname{pts}(t) \times \mathbb{R})))}}
\vnbprop{prod-top-opens}{\vnbstmt{\begin{array}{@{}l@{}} \forall t. \\ \quad \operatorname{is\text{-}top\text{-}space}(t) \Rightarrow  \\ \quad \quad (\operatorname{opens}(\operatorname{prod\text{-}top}(t)) \simeq \{u \in \operatorname{power}((\operatorname{pts}(t) \times \mathbb{R})) : \operatorname{is\text{-}pt\text{-}open}(t, u)\}) \end{array}}}
\vnbprop{prod-top-is-top-space}{\vnbstmt{\forall t.\; (\operatorname{is\text{-}top\text{-}space}(t) \Rightarrow \operatorname{is\text{-}top\text{-}space}(\operatorname{prod\text{-}top}(t)))}}
\vnbprop{pt-box-open}{\vnbstmt{\begin{array}{@{}l@{}} \forall t, v, c, e. \\ \quad (\operatorname{is\text{-}top\text{-}space}(t) \wedge ((v \in \operatorname{opens}(t)) \wedge ((c \in \mathbb{R}) \wedge \operatorname{pos\text{-}rr}(e)))) \Rightarrow  \\ \quad \quad \operatorname{pt\text{-}box}(t, v, c, e) \in \operatorname{opens}(\operatorname{prod\text{-}top}(t)) \end{array}}}
\vnbprop{prod-top-hausdorff}{\vnbstmt{\forall t.\; (\operatorname{is\text{-}hausdorff}(t) \Rightarrow \operatorname{is\text{-}hausdorff}(\operatorname{prod\text{-}top}(t)))}}
\vnbprop{prod-add-continuous}{\vnbstmt{\begin{array}{@{}l@{}} \forall m, t, x, y, u. \\ \quad \operatorname{is\text{-}tvs}(\operatorname{rr\text{-}normed\text{-}field}, m, t) \wedge \\ \quad x \in (\operatorname{vec}(m) \times \mathbb{R}) \wedge \\ \quad y \in (\operatorname{vec}(m) \times \mathbb{R}) \wedge \\ \quad u \in \operatorname{opens}(\operatorname{prod\text{-}top}(t)) \wedge \\ \quad \operatorname{pv\text{-}vadd}(m)(x, y) \in u \Rightarrow \\ \quad \quad \exists v \in \operatorname{opens}(\operatorname{prod\text{-}top}(t)). \\ \quad \quad \quad x \in v \wedge \\ \quad \quad \quad \exists w \in \operatorname{opens}(\operatorname{prod\text{-}top}(t)). \\ \quad \quad \quad \quad (y \in w) \wedge \forall a \in v, b \in w.\; (\operatorname{pv\text{-}vadd}(m)(a, b) \in u) \end{array}}}
\vnbprop{prod-scal-continuous}{\vnbstmt{\begin{array}{@{}l@{}} \forall m, t, r, x, u. \\ \quad \operatorname{is\text{-}tvs}(\operatorname{rr\text{-}normed\text{-}field}, m, t) \wedge \\ \quad r \in \operatorname{carr}(\operatorname{scal}(m)) \wedge \\ \quad x \in (\operatorname{vec}(m) \times \mathbb{R}) \wedge \\ \quad u \in \operatorname{opens}(\operatorname{prod\text{-}top}(t)) \wedge \\ \quad \operatorname{pv\text{-}act}(m)(r, x) \in u \Rightarrow \\ \quad \quad \exists d. \\ \quad \quad \quad \operatorname{pos\text{-}rr}(d) \wedge \\ \quad \quad \quad \exists w \in \operatorname{opens}(\operatorname{prod\text{-}top}(t)). \\ \quad \quad \quad \quad x \in w \wedge \\ \quad \quad \quad \quad \forall s \in \operatorname{carr}(\operatorname{scal}(m)), a \in w. \\ \quad \quad \quad \quad \quad (\operatorname{fnrm}(\operatorname{rr\text{-}normed\text{-}field})(\operatorname{add}(\operatorname{scal}(m))(s, \operatorname{neg}(\operatorname{scal}(m))(r))) < d) \Rightarrow  \\ \quad \quad \quad \quad \quad \quad \operatorname{pv\text{-}act}(m)(s, a) \in u \end{array}}}
\vnbprop{prod-is-tvs}{\vnbstmt{\begin{array}{@{}l@{}} \forall m, t. \\ \quad \operatorname{is\text{-}tvs}(\operatorname{rr\text{-}normed\text{-}field}, m, t) \Rightarrow  \\ \quad \quad \operatorname{is\text{-}tvs}(\operatorname{rr\text{-}normed\text{-}field}, \operatorname{prod\text{-}rmod}(m), \operatorname{prod\text{-}top}(t)) \end{array}}}
\vnbprop{rr-halfline-open}{\vnbstmt{\begin{array}{@{}l@{}} \forall c \in \mathbb{R}. \\ \quad \{v \in \mathbb{R} : (v < c)\} \\ \quad \in\; \operatorname{opens}(\operatorname{seminorm\text{-}top}(\operatorname{normed\text{-}vector\text{-}space\text{-}as\text{-}module}(\operatorname{rr\text{-}nvs}), \\ \quad \quad \quad \{0\}, (\lambda k \in \{0\}.\; \operatorname{vnrm}(\operatorname{rr\text{-}nvs})))) \end{array}}}
\vnbprop{rr-convex-comb-strict}{\vnbstmt{\begin{array}{@{}l@{}} \forall h, a, b, p, q. \\ \quad ((h \in \mathbb{R}) \wedge ((a \in \mathbb{R}) \wedge ((b \in \mathbb{R}) \wedge ((p \in \mathbb{R}) \wedge ((q \in \mathbb{R}) \wedge ((0 \leq h) \wedge ((h \leq 1) \wedge ((p < a) \wedge (q < b))))))))) \Rightarrow  \\ \quad \quad ((h \cdot q) + ((1 - h) \cdot p)) < ((h \cdot b) + ((1 - h) \cdot a)) \end{array}}}
\vnbprop{strict-epigraph-unfold}{\vnbstmt{\begin{array}{@{}l@{}} \forall m, p. \\ \quad (\operatorname{strict\text{-}epigraph}(m, p) \simeq \{w \in (\operatorname{vec}(m) \times \mathbb{R}) : (p(\operatorname{nth}(1, w)) < \operatorname{nth}(2, w))\}) \end{array}}}
\vnbprop{strict-epigraph-open-convex}{\vnbstmt{\begin{array}{@{}l@{}} \forall m, t, p. \\ \quad \operatorname{is\text{-}tvs}(\operatorname{rr\text{-}normed\text{-}field}, m, t) \wedge \\ \quad \operatorname{is\text{-}convex\text{-}functional}(m, p) \wedge \\ \quad \operatorname{is\text{-}hom\text{-}top\text{-}space}(t, \\ \quad \quad \operatorname{seminorm\text{-}top}(\operatorname{normed\text{-}vector\text{-}space\text{-}as\text{-}module}(\operatorname{rr\text{-}nvs}), \\ \quad \quad \quad \{0\}, (\lambda k \in \{0\}.\; \operatorname{vnrm}(\operatorname{rr\text{-}nvs}))), \\ \quad \quad p) \Rightarrow \\ \quad \quad \operatorname{strict\text{-}epigraph}(m, p) \in \operatorname{opens}(\operatorname{prod\text{-}top}(t)) \wedge \\ \quad \quad \operatorname{is\text{-}convex\text{-}subset}(\operatorname{prod\text{-}rmod}(m), \operatorname{strict\text{-}epigraph}(m, p)) \end{array}}}

\vnbbookfile{structure-library/nets}{structure-library/nets}
\noindent{\small nets-vocab.scm -- DIRECTED SETS, NETS AND SUBNETS in a topological space; Cauchy nets and completeness in a topological vector space (topological-vector-space.pdf Def 4.1).  Definitions only.  Agent NETS (worker-08), 2026-10-26 (NEW DEFINITIONS, the user's to confirm).  DESTINATION: structure-library/nets.scm. Library: ...}\par
\vnbstructdef{directed-set}{A directed set $s$ is a 2-tuple whose components are listed below.\par
\begin{enumerate}[nosep]
\item $\operatorname{pts}(s) = \operatorname{nth}(1, s)$, a carrier (\texttt{carriers}): $\operatorname{pts}(s) \in \mathbf{Set}$.
\item $\operatorname{dleq}(s) = \operatorname{nth}(2, s)$, a relation (\texttt{relation}): $\operatorname{dleq}(s) \in \operatorname{power}((\operatorname{pts}(s) \times \operatorname{pts}(s)))$.
\end{enumerate}
The laws (the \texttt{property} and \texttt{law} clauses):
\begin{itemize}[nosep]
\item $\exists \operatorname{dsn\_}.\; (\operatorname{dsn\_} \in \operatorname{pts}(s))$
\item $\forall \operatorname{dsa\_} \in \operatorname{pts}(s).\; ({[}\operatorname{dsa\_}, \operatorname{dsa\_}] \in \operatorname{dleq}(s))$
\item
\begin{equation*}
\begin{array}{@{}l@{}} \forall \operatorname{dsa\_}, \operatorname{dsb\_}, \operatorname{dsc\_} \in \operatorname{pts}(s). \\ \quad ({[}\operatorname{dsa\_}, \operatorname{dsb\_}] \in \operatorname{dleq}(s)) \Rightarrow  \\ \quad ({[}\operatorname{dsb\_}, \operatorname{dsc\_}] \in \operatorname{dleq}(s)) \Rightarrow ({[}\operatorname{dsa\_}, \operatorname{dsc\_}] \in \operatorname{dleq}(s)) \end{array}
\end{equation*}

\item
\begin{equation*}
\begin{array}{@{}l@{}} \forall \operatorname{dsa\_}, \operatorname{dsb\_} \in \operatorname{pts}(s). \exists \operatorname{dsc\_} \in \operatorname{pts}(s). \\ \quad ({[}\operatorname{dsa\_}, \operatorname{dsc\_}] \in \operatorname{dleq}(s)) \wedge ({[}\operatorname{dsb\_}, \operatorname{dsc\_}] \in \operatorname{dleq}(s)) \end{array}
\end{equation*}

\end{itemize}
The characteristic law (the axiom \texttt{is-directed-set}): $\operatorname{is\text{-}directed\text{-}set}(s)$ holds if and only if $\operatorname{length}(s) = 2$, every component has the typing listed with it, every law holds, and $s \in \mathbf{Set}$ (the structure is a set).
}
\par\medskip\noindent\textbf{Definition} (\texttt{is-eventually-in}).\;For all $d$, $x$ and $u$, 
\begin{equation*}
\begin{array}{@{}l@{}} \operatorname{is\text{-}eventually\text{-}in}(d, x, u) \\ \Leftrightarrow\; \exists \operatorname{evi\_} \in \operatorname{pts}(d). \forall \operatorname{evj\_} \in \operatorname{pts}(d). \\ \quad ({[}\operatorname{evi\_}, \operatorname{evj\_}] \in \operatorname{dleq}(d)) \Rightarrow (x(\operatorname{evj\_}) \in u) \end{array}
\end{equation*}

\par\medskip\noindent\textbf{Definition} (\texttt{net-converges-to}).\;For all $t$, $d$, $x$ and $a$, 
\begin{equation*}
\begin{array}{@{}l@{}} \operatorname{net\text{-}converges\text{-}to}(t, d, x, a) \\ \Leftrightarrow\; \operatorname{is\text{-}top\text{-}space}(t) \wedge \\ \operatorname{is\text{-}directed\text{-}set}(d) \wedge \\ x \in (\operatorname{pts}(d) \to \operatorname{pts}(t)) \wedge \\ a \in \operatorname{pts}(t) \wedge \\ \forall \operatorname{ncu\_} \in \operatorname{opens}(t).\; ((a \in \operatorname{ncu\_}) \Rightarrow \operatorname{is\text{-}eventually\text{-}in}(d, x, \operatorname{ncu\_})) \end{array}
\end{equation*}

\par\medskip\noindent\textbf{Definition} (\texttt{is-cofinal-monotone}).\;For all $e$, $d$ and $h$, 
\begin{equation*}
\begin{array}{@{}l@{}} \operatorname{is\text{-}cofinal\text{-}monotone}(e, d, h) \\ \Leftrightarrow\; h \in (\operatorname{pts}(e) \to \operatorname{pts}(d)) \wedge \\ \forall \operatorname{cmi\_}, \operatorname{cmj\_} \in \operatorname{pts}(e). \\ \quad ({[}\operatorname{cmi\_}, \operatorname{cmj\_}] \in \operatorname{dleq}(e)) \Rightarrow ({[}h(\operatorname{cmi\_}), h(\operatorname{cmj\_})] \in \operatorname{dleq}(d)) \wedge \\ \forall \operatorname{cmk\_} \in \operatorname{pts}(d).\; \exists \operatorname{cml\_} \in \operatorname{pts}(e).\; ({[}\operatorname{cmk\_}, h(\operatorname{cml\_})] \in \operatorname{dleq}(d)) \end{array}
\end{equation*}

\par\medskip\noindent\textbf{Definition} (\texttt{is-subnet}).\;For all $d$, $x$, $e$ and $y$, 
\begin{equation*}
\begin{array}{@{}l@{}} \operatorname{is\text{-}subnet}(d, x, e, y) \\ \Leftrightarrow\; \operatorname{is\text{-}directed\text{-}set}(d) \wedge \\ \operatorname{is\text{-}directed\text{-}set}(e) \wedge \\ \exists \operatorname{snh\_}. \\ \quad (\operatorname{is\text{-}cofinal\text{-}monotone}(e, d, \operatorname{snh\_}) \wedge \forall \operatorname{snj\_} \in \operatorname{pts}(e).\; (y(\operatorname{snj\_}) = x(\operatorname{snh\_}(\operatorname{snj\_})))) \end{array}
\end{equation*}

\par\medskip\noindent\textbf{Definition} (\texttt{nn-directed}).\;
\begin{equation*}
\operatorname{nn\text{-}directed} = {[}\mathbb{N}, \{\operatorname{ndp\_} \in (\mathbb{N} \times \mathbb{N}) : (\operatorname{nth}(1, \operatorname{ndp\_}) \leq \operatorname{nth}(2, \operatorname{ndp\_}))\}]
\end{equation*}

\par\medskip\noindent\textbf{Definition} (\texttt{is-cauchy-net}).\;For all $m$, $t$, $d$ and $x$, 
\begin{equation*}
\begin{array}{@{}l@{}} \operatorname{is\text{-}cauchy\text{-}net}(m, t, d, x) \\ \Leftrightarrow\; \operatorname{is\text{-}directed\text{-}set}(d) \wedge \\ x \in (\operatorname{pts}(d) \to \operatorname{vec}(m)) \wedge \\ \forall \operatorname{cnu\_} \in \operatorname{opens}(t). \\ \quad (\operatorname{vzero}(m) \in \operatorname{cnu\_}) \Rightarrow  \\ \quad \quad \exists \operatorname{cni\_} \in \operatorname{pts}(d). \forall \operatorname{cnj\_} \in \operatorname{pts}(d). \\ \quad \quad \quad ({[}\operatorname{cni\_}, \operatorname{cnj\_}] \in \operatorname{dleq}(d)) \Rightarrow  \\ \quad \quad \quad \quad \forall \operatorname{cnk\_} \in \operatorname{pts}(d). \\ \quad \quad \quad \quad \quad ({[}\operatorname{cni\_}, \operatorname{cnk\_}] \in \operatorname{dleq}(d)) \Rightarrow (\operatorname{vadd}(m)(x(\operatorname{cnj\_}), \operatorname{vneg}(m)(x(\operatorname{cnk\_}))) \in \operatorname{cnu\_}) \end{array}
\end{equation*}

\par\medskip\noindent\textbf{Definition} (\texttt{is-complete-subset}).\;For all $m$, $t$ and $a$, 
\begin{equation*}
\begin{array}{@{}l@{}} \operatorname{is\text{-}complete\text{-}subset}(m, t, a) \\ \Leftrightarrow\; a \subseteq \operatorname{vec}(m) \wedge \\ \forall \operatorname{csd\_}, \operatorname{csx\_}. \\ \quad \operatorname{is\text{-}cauchy\text{-}net}(m, t, \operatorname{csd\_}, \operatorname{csx\_}) \Rightarrow  \\ \quad \forall \operatorname{csi\_} \in \operatorname{pts}(\operatorname{csd\_}).\; (\operatorname{csx\_}(\operatorname{csi\_}) \in a) \Rightarrow  \\ \quad \quad \exists \operatorname{csl\_} \in a.\; \operatorname{net\text{-}converges\text{-}to}(t, \operatorname{csd\_}, \operatorname{csx\_}, \operatorname{csl\_}) \end{array}
\end{equation*}

\par\medskip\noindent\textbf{Definition} (\texttt{ptwise-net-limits}).\;$\operatorname{ptwise\text{-}net\text{-}limits}(t, \mathit{ee}) :=$
\begin{equation*}
\{\operatorname{pwf\_} \in (\operatorname{pts}(t) \to \mathbb{R}) : \exists \operatorname{pwd\_} \in \operatorname{directed\text{-}set}, \operatorname{pwp\_}.\; (\operatorname{is\text{-}directed\text{-}set}(\operatorname{pwd\_}) \wedge ((\operatorname{pwp\_} \in (\operatorname{pts}(\operatorname{pwd\_}) \to \mathit{ee})) \wedge \forall \operatorname{pwx\_} \in \operatorname{pts}(t).\; \operatorname{net\text{-}converges\text{-}to}(\operatorname{metric\text{-}top}(\operatorname{rr\text{-}ms}), \operatorname{pwd\_}, (\lambda \operatorname{pwl\_} \in \operatorname{pts}(\operatorname{pwd\_}).\; \operatorname{pwp\_}(\operatorname{pwl\_})(\operatorname{pwx\_})), \operatorname{pwf\_}(\operatorname{pwx\_}))))\}
\end{equation*}

\par\medskip\noindent\textbf{Definition} (\texttt{is-equicontinuous-t}).\;For all $t$ and $\mathit{ee}$, 
\begin{equation*}
\begin{array}{@{}l@{}} \operatorname{is\text{-}equicontinuous\text{-}t}(t, \mathit{ee}) \\ \Leftrightarrow\; \mathit{ee} \subseteq (\operatorname{pts}(t) \to \mathbb{R}) \wedge \\ \forall \operatorname{eqx\_} \in \operatorname{pts}(t), \operatorname{eqe\_}. \\ \quad \operatorname{pos\text{-}rr}(\operatorname{eqe\_}) \Rightarrow  \\ \quad \quad \exists \operatorname{eqv\_} \in \operatorname{opens}(t). \\ \quad \quad \quad ((\operatorname{eqx\_} \in \operatorname{eqv\_}) \wedge \forall \operatorname{eqf\_} \in \mathit{ee}, \operatorname{eqy\_} \in \operatorname{eqv\_}.\; (\lvert -(\operatorname{eqf\_}(\operatorname{eqx\_}), \operatorname{eqf\_}(\operatorname{eqy\_})) \rvert < \operatorname{eqe\_})) \end{array}
\end{equation*}

\par\medskip\noindent\textbf{Definition} (\texttt{net-unif-conv-on}).\;For all $d$, $p$, $f$ and $k$, 
\begin{equation*}
\begin{array}{@{}l@{}} \operatorname{net\text{-}unif\text{-}conv\text{-}on}(d, p, f, k) \\ \Leftrightarrow\; \forall \operatorname{uce\_}. \\ \quad \operatorname{pos\text{-}rr}(\operatorname{uce\_}) \Rightarrow  \\ \quad \quad \exists \operatorname{uci\_} \in \operatorname{pts}(d). \forall \operatorname{ucj\_} \in \operatorname{pts}(d). \\ \quad \quad \quad ({[}\operatorname{uci\_}, \operatorname{ucj\_}] \in \operatorname{dleq}(d)) \Rightarrow  \\ \quad \quad \quad \quad \forall \operatorname{ucy\_} \in k.\; (\lvert -(p(\operatorname{ucj\_})(\operatorname{ucy\_}), f(\operatorname{ucy\_})) \rvert < \operatorname{uce\_}) \end{array}
\end{equation*}

\par\medskip\noindent\textbf{Definition} (\texttt{hom-directed-set}).\;$\operatorname{hom\text{-}directed\text{-}set}(a, b) :=$
\begin{equation*}
\{\operatorname{homf\_} \in (\operatorname{pts}(a) \to \operatorname{pts}(b)) : \operatorname{is\text{-}hom\text{-}directed\text{-}set}(a, b, \operatorname{homf\_})\}
\end{equation*}

\vnbdef{directed-set-class}{\vnbstmt{\forall s.\; ((s \in \operatorname{directed\text{-}set}) \Leftrightarrow \operatorname{is\text{-}directed\text{-}set}(s))}}
\vnbdef{is-hom-directed-set-def}{\vnbstmt{\begin{array}{@{}l@{}} \forall a, b, f. \\ \quad \operatorname{is\text{-}hom\text{-}directed\text{-}set}(a, b, f) \\ \quad \Leftrightarrow\; \operatorname{is\text{-}directed\text{-}set}(a) \wedge \\ \quad \operatorname{is\text{-}directed\text{-}set}(b) \wedge \\ \quad f \in (\operatorname{pts}(a) \to \operatorname{pts}(b)) \wedge \\ \quad \forall x_{1}, x_{2} \in \operatorname{pts}(a). \\ \quad \quad ({[}x_{1}, x_{2}] \in \operatorname{dleq}(a)) \Rightarrow ({[}f(x_{1}), f(x_{2})] \in \operatorname{dleq}(b)) \end{array}}}
\vnbprop{is-directed-set-is-set}{\vnbstmt{\forall s.\; (\operatorname{is\text{-}directed\text{-}set}(s) \Rightarrow (s \in \mathbf{Set}))}}

\vnbbookfile{theorem-library/nets}{theorem-library/nets}
\noindent{\small nets.scm -- directed sets, nets, subnets, net convergence in a topological space, the metric and sequence bridges, Cauchy nets and completeness in a topological vector space (Def 4.1), and the parts of TB-10 / TB-12 that need nothing but nets.  Agent NETS (worker-08), 2026-10-26. DESTINATION: theorem-library/nets.scm.  ...}\par
\vnbprop{ds-nonempty}{\vnbstmt{\forall d.\; (\operatorname{is\text{-}directed\text{-}set}(d) \Rightarrow \exists n.\; (n \in \operatorname{pts}(d)))}}
\vnbprop{ds-refl}{\vnbstmt{\forall d.; \forall a \in \operatorname{pts}(d).\; (\operatorname{is\text{-}directed\text{-}set}(d) \Rightarrow ({[}a, a] \in \operatorname{dleq}(d)))}}
\vnbprop{ds-trans}{\vnbstmt{\begin{array}{@{}l@{}} \forall d.; \forall a, b, c \in \operatorname{pts}(d). \\ \quad (\operatorname{is\text{-}directed\text{-}set}(d) \wedge (({[}a, b] \in \operatorname{dleq}(d)) \wedge ({[}b, c] \in \operatorname{dleq}(d)))) \Rightarrow  \\ \quad \quad {[}a, c] \in \operatorname{dleq}(d) \end{array}}}
\vnbprop{ds-upper}{\vnbstmt{\begin{array}{@{}l@{}} \forall d.; \forall a, b \in \operatorname{pts}(d). \\ \quad \operatorname{is\text{-}directed\text{-}set}(d) \Rightarrow  \\ \quad \quad \exists c \in \operatorname{pts}(d).\; (({[}a, c] \in \operatorname{dleq}(d)) \wedge ({[}b, c] \in \operatorname{dleq}(d))) \end{array}}}
\vnbprop{eventually-mono}{\vnbstmt{\begin{array}{@{}l@{}} \forall d, x, u, v. \\ \quad (\operatorname{is\text{-}eventually\text{-}in}(d, x, u) \wedge (u \subseteq v)) \Rightarrow \operatorname{is\text{-}eventually\text{-}in}(d, x, v) \end{array}}}
\vnbprop{eventually-and}{\vnbstmt{\begin{array}{@{}l@{}} \forall d, x, u, v. \\ \quad (\operatorname{is\text{-}directed\text{-}set}(d) \wedge (\operatorname{is\text{-}eventually\text{-}in}(d, x, u) \wedge \operatorname{is\text{-}eventually\text{-}in}(d, x, v))) \Rightarrow  \\ \quad \quad \operatorname{is\text{-}eventually\text{-}in}(d, x, (u \cap v)) \end{array}}}
\vnbprop{net-const-converges}{\vnbstmt{\begin{array}{@{}l@{}} \forall t, d, x, a. \\ \quad \operatorname{is\text{-}top\text{-}space}(t) \wedge \\ \quad \operatorname{is\text{-}directed\text{-}set}(d) \wedge \\ \quad x \in (\operatorname{pts}(d) \to \operatorname{pts}(t)) \wedge \\ \quad a \in \operatorname{pts}(t) \wedge \\ \quad \forall j \in \operatorname{pts}(d).\; (x(j) = a) \Rightarrow \\ \quad \quad \operatorname{net\text{-}converges\text{-}to}(t, d, x, a) \end{array}}}
\vnbprop{subnet-converges}{\vnbstmt{\begin{array}{@{}l@{}} \forall t, d, x, a, e, y. \\ \quad (\operatorname{net\text{-}converges\text{-}to}(t, d, x, a) \wedge (\operatorname{is\text{-}subnet}(d, x, e, y) \wedge (y \in (\operatorname{pts}(e) \to \operatorname{pts}(t))))) \Rightarrow  \\ \quad \quad \operatorname{net\text{-}converges\text{-}to}(t, e, y, a) \end{array}}}
\vnbprop{net-continuous-image}{\vnbstmt{\begin{array}{@{}l@{}} \forall t, d, x, a, w, f, y. \\ \quad \operatorname{net\text{-}converges\text{-}to}(t, d, x, a) \wedge \\ \quad \operatorname{is\text{-}hom\text{-}top\text{-}space}(t, w, f) \wedge \\ \quad y \in (\operatorname{pts}(d) \to \operatorname{pts}(w)) \wedge \\ \quad \forall j \in \operatorname{pts}(d).\; (y(j) = f(x(j))) \Rightarrow \\ \quad \quad \operatorname{net\text{-}converges\text{-}to}(w, d, y, f(a)) \end{array}}}
\vnbprop{net-limit-unique}{\vnbstmt{\begin{array}{@{}l@{}} \forall t, d, x, a, b. \\ \quad (\operatorname{is\text{-}hausdorff}(t) \wedge (\operatorname{net\text{-}converges\text{-}to}(t, d, x, a) \wedge \operatorname{net\text{-}converges\text{-}to}(t, d, x, b))) \Rightarrow  \\ \quad \quad a = b \end{array}}}
\vnbprop{metric-top-pts}{\vnbstmt{\forall s.\; (\operatorname{is\text{-}metric\text{-}space}(s) \Rightarrow (\operatorname{pts}(\operatorname{metric\text{-}top}(s)) \simeq \operatorname{pts}(s)))}}
\vnbprop{metric-net-converges-iff}{\vnbstmt{\begin{array}{@{}l@{}} \forall s, d, x, a. \\ \quad (\operatorname{is\text{-}metric\text{-}space}(s) \wedge (\operatorname{is\text{-}directed\text{-}set}(d) \wedge ((x \in (\operatorname{pts}(d) \to \operatorname{pts}(s))) \wedge (a \in \operatorname{pts}(s))))) \Rightarrow  \\ \quad \quad \operatorname{net\text{-}converges\text{-}to}(\operatorname{metric\text{-}top}(s), d, x, a) \\ \quad \quad \Leftrightarrow\; \forall e.\; (\operatorname{pos\text{-}rr}(e) \Rightarrow \operatorname{is\text{-}eventually\text{-}in}(d, x, \operatorname{ball}(s, a, e))) \end{array}}}
\vnbprop{nn-directed-pts}{\vnbstmt{(\operatorname{pts}(\operatorname{nn\text{-}directed}) \simeq \mathbb{N})}}
\vnbprop{nn-directed-dleq}{\vnbstmt{(\operatorname{dleq}(\operatorname{nn\text{-}directed}) \simeq \{p \in (\mathbb{N} \times \mathbb{N}) : (\operatorname{nth}(1, p) \leq \operatorname{nth}(2, p))\})}}
\vnbprop{nn-directed-le-iff}{\vnbstmt{\forall m, n \in \mathbb{N}.\; (({[}m, n] \in \operatorname{dleq}(\operatorname{nn\text{-}directed})) \Leftrightarrow (m \leq n))}}
\vnbprop{nn-directed-is-directed-set}{\vnbstmt{\operatorname{is\text{-}directed\text{-}set}(\operatorname{nn\text{-}directed})}}
\vnbprop{seq-converges-iff-net}{\vnbstmt{\begin{array}{@{}l@{}} \forall s, f, l. \\ \quad (\operatorname{is\text{-}metric\text{-}space}(s) \wedge ((f \in (\mathbb{N} \to \operatorname{pts}(s))) \wedge (l \in \operatorname{pts}(s)))) \Rightarrow  \\ \quad \quad \operatorname{converges\text{-}to}(s, f, l) \Leftrightarrow \operatorname{net\text{-}converges\text{-}to}(\operatorname{metric\text{-}top}(s), \operatorname{nn\text{-}directed}, f, l) \end{array}}}
\vnbprop{rmod-vneg-act}{\vnbstmt{\begin{array}{@{}l@{}} \forall m, x. \\ \quad (\operatorname{is\text{-}module}(m) \wedge ((\operatorname{scal}(m) = \operatorname{normed\text{-}field\text{-}as\text{-}commutative\text{-}ring}(\operatorname{rr\text{-}normed\text{-}field})) \wedge (x \in \operatorname{vec}(m)))) \Rightarrow  \\ \quad \quad \operatorname{vneg}(m)(x) = \operatorname{act}(m)(-1, x) \end{array}}}
\vnbprop{tvs-sub-continuous}{\vnbstmt{\begin{array}{@{}l@{}} \forall m, t, x, y, u. \\ \quad \operatorname{is\text{-}tvs}(\operatorname{rr\text{-}normed\text{-}field}, m, t) \wedge \\ \quad x \in \operatorname{vec}(m) \wedge \\ \quad y \in \operatorname{vec}(m) \wedge \\ \quad u \in \operatorname{opens}(t) \wedge \\ \quad \operatorname{vadd}(m)(x, \operatorname{vneg}(m)(y)) \in u \Rightarrow \\ \quad \quad \exists v \in \operatorname{opens}(t). \\ \quad \quad \quad x \in v \wedge \\ \quad \quad \quad \exists w \in \operatorname{opens}(t). \\ \quad \quad \quad \quad (y \in w) \wedge \forall p \in v, q \in w.\; (\operatorname{vadd}(m)(p, \operatorname{vneg}(m)(q)) \in u) \end{array}}}
\vnbprop{tvs-convergent-net-cauchy}{\vnbstmt{\begin{array}{@{}l@{}} \forall m, t, d, x, a. \\ \quad (\operatorname{is\text{-}tvs}(\operatorname{rr\text{-}normed\text{-}field}, m, t) \wedge \operatorname{net\text{-}converges\text{-}to}(t, d, x, a)) \Rightarrow  \\ \quad \quad \operatorname{is\text{-}cauchy\text{-}net}(m, t, d, x) \end{array}}}
\vnbprop{tvs-cauchy-subnet-converges}{\vnbstmt{\begin{array}{@{}l@{}} \forall m, t, d, x, e, y, a. \\ \quad \operatorname{is\text{-}tvs}(\operatorname{rr\text{-}normed\text{-}field}, m, t) \wedge \\ \quad \operatorname{is\text{-}cauchy\text{-}net}(m, t, d, x) \wedge \\ \quad \operatorname{is\text{-}subnet}(d, x, e, y) \wedge \\ \quad \operatorname{net\text{-}converges\text{-}to}(t, e, y, a) \Rightarrow \\ \quad \quad \operatorname{net\text{-}converges\text{-}to}(t, d, x, a) \end{array}}}
\vnbprop{net-limit-in-closed}{\vnbstmt{\begin{array}{@{}l@{}} \forall t, d, x, l, a. \\ \quad (\operatorname{is\text{-}top\text{-}closed}(t, a) \wedge (\operatorname{net\text{-}converges\text{-}to}(t, d, x, l) \wedge \forall i \in \operatorname{pts}(d).\; (x(i) \in a))) \Rightarrow  \\ \quad \quad l \in a \end{array}}}
\vnbprop{closed-subset-complete}{\vnbstmt{\begin{array}{@{}l@{}} \forall m, t, a. \\ \quad \operatorname{is\text{-}tvs}(\operatorname{rr\text{-}normed\text{-}field}, m, t) \wedge \\ \quad \operatorname{is\text{-}complete\text{-}subset}(m, t, \operatorname{vec}(m)) \wedge \\ \quad a \subseteq \operatorname{vec}(m) \wedge \\ \quad \operatorname{is\text{-}top\text{-}closed}(t, a) \Rightarrow \\ \quad \quad \operatorname{is\text{-}complete\text{-}subset}(m, t, a) \end{array}}}
\vnbprop{adherent-point-net}{\vnbstmt{\begin{array}{@{}l@{}} \forall t, a, l. \\ \quad \operatorname{is\text{-}top\text{-}space}(t) \wedge \\ \quad a \subseteq \operatorname{pts}(t) \wedge \\ \quad l \in \operatorname{pts}(t) \wedge \\ \quad \forall u \in \operatorname{opens}(t).\; ((l \in u) \Rightarrow \exists q \in u.\; (q \in a)) \Rightarrow \\ \quad \quad \exists n, x. \\ \quad \quad \quad \operatorname{is\text{-}directed\text{-}set}(n) \wedge \\ \quad \quad \quad x \in (\operatorname{pts}(n) \to \operatorname{pts}(t)) \wedge \\ \quad \quad \quad \forall i \in \operatorname{pts}(n).\; (x(i) \in a) \wedge \\ \quad \quad \quad \operatorname{net\text{-}converges\text{-}to}(t, n, x, l) \end{array}}}
\vnbprop{complete-subset-closed}{\vnbstmt{\begin{array}{@{}l@{}} \forall m, t, a. \\ \quad (\operatorname{is\text{-}tvs}(\operatorname{rr\text{-}normed\text{-}field}, m, t) \wedge \operatorname{is\text{-}complete\text{-}subset}(m, t, a)) \Rightarrow  \\ \quad \quad \operatorname{is\text{-}top\text{-}closed}(t, a) \end{array}}}
\vnbprop{ptwise-net-limits-unfold}{\vnbstmt{\begin{array}{@{}l@{}} \forall t, e. \\ \quad (\operatorname{ptwise\text{-}net\text{-}limits}(t, e) \simeq \{f \in (\operatorname{pts}(t) \to \mathbb{R}) : \exists d \in \operatorname{directed\text{-}set}, p.\; (\operatorname{is\text{-}directed\text{-}set}(d) \wedge ((p \in (\operatorname{pts}(d) \to e)) \wedge \forall x \in \operatorname{pts}(t).\; \operatorname{net\text{-}converges\text{-}to}(\operatorname{metric\text{-}top}(\operatorname{rr\text{-}ms}), d, (\lambda l \in \operatorname{pts}(d).\; p(l)(x)), f(x))))\}) \end{array}}}
\vnbprop{ptwise-limits-equicontinuous}{\vnbstmt{\begin{array}{@{}l@{}} \forall t, e. \\ \quad (\operatorname{is\text{-}top\text{-}space}(t) \wedge \operatorname{is\text{-}equicontinuous\text{-}t}(t, e)) \Rightarrow  \\ \quad \quad \operatorname{is\text{-}equicontinuous\text{-}t}(t, \operatorname{ptwise\text{-}net\text{-}limits}(t, e)) \end{array}}}
\vnbprop{directed-finite-common-index}{\vnbstmt{\begin{array}{@{}l@{}} \forall d, r.; \forall s \in \mathbf{Set}. \\ \quad \operatorname{is\text{-}directed\text{-}set}(d) \wedge \\ \quad \forall w, i, k. \\ \quad \quad (({[}w, i] \in r) \wedge ((k \in \operatorname{pts}(d)) \wedge ({[}i, k] \in \operatorname{dleq}(d)))) \Rightarrow ({[}w, k] \in r) \wedge \\ \quad |s| \in \mathbb{N} \wedge \\ \quad \forall u \in s.\; \exists j \in \operatorname{pts}(d).\; ({[}u, j] \in r) \Rightarrow \\ \quad \quad \exists m \in \operatorname{pts}(d).\; \forall v \in s.\; ({[}v, m] \in r) \end{array}}}
\vnbprop{equicontinuous-ptwise-implies-compact-unif}{\vnbstmt{\begin{array}{@{}l@{}} \forall t, e, d, p, f, k. \\ \quad \operatorname{is\text{-}top\text{-}space}(t) \wedge \\ \quad \operatorname{is\text{-}equicontinuous\text{-}t}(t, e) \wedge \\ \quad \operatorname{is\text{-}directed\text{-}set}(d) \wedge \\ \quad p \in (\operatorname{pts}(d) \to e) \wedge \\ \quad f \in e \wedge \\ \quad \forall x \in \operatorname{pts}(t). \\ \quad \quad \operatorname{net\text{-}converges\text{-}to}(\operatorname{metric\text{-}top}(\operatorname{rr\text{-}ms}), d, (\lambda l \in \operatorname{pts}(d).\; p(l)(x)), f(x)) \wedge \\ \quad k \subseteq \operatorname{pts}(t) \wedge \\ \quad \operatorname{is\text{-}compact\text{-}t}(\operatorname{subspace\text{-}top}(t, k)) \Rightarrow \\ \quad \quad \operatorname{net\text{-}unif\text{-}conv\text{-}on}(d, p, f, k) \end{array}}}
\vnbprop{singleton-subspace-compact}{\vnbstmt{\begin{array}{@{}l@{}} \forall t, x. \\ \quad (\operatorname{is\text{-}top\text{-}space}(t) \wedge (x \in \operatorname{pts}(t))) \Rightarrow  \\ \quad \quad \operatorname{is\text{-}compact\text{-}t}(\operatorname{subspace\text{-}top}(t, \{x\})) \end{array}}}
\vnbprop{unif-conv-on-pointwise}{\vnbstmt{\begin{array}{@{}l@{}} \forall t, e, d, p, f, k, x. \\ \quad \operatorname{is\text{-}directed\text{-}set}(d) \wedge \\ \quad p \in (\operatorname{pts}(d) \to e) \wedge \\ \quad e \subseteq (\operatorname{pts}(t) \to \mathbb{R}) \wedge \\ \quad f \in (\operatorname{pts}(t) \to \mathbb{R}) \wedge \\ \quad k \subseteq \operatorname{pts}(t) \wedge \\ \quad x \in k \wedge \\ \quad \operatorname{net\text{-}unif\text{-}conv\text{-}on}(d, p, f, k) \Rightarrow \\ \quad \quad \operatorname{net\text{-}converges\text{-}to}(\operatorname{metric\text{-}top}(\operatorname{rr\text{-}ms}), d, (\lambda l \in \operatorname{pts}(d).\; p(l)(x)), f(x)) \end{array}}}

\vnbbookfile{structure-library/weak-star}{structure-library/weak-star}
\noindent{\small tb6-vocab.scm -- THE WEAK* TOPOLOGY AND POLARS (definitions only), for topological-vector-space.pdf Cor 3.27 (Banach-Alaoglu) and Prop 3.29.  Agent TB-6 (worker-08), 2026-10-26 (NEW DEFINITIONS, the user's to confirm).  DESTINATION: structure-library/weak-star.scm.  Library: analysis. Load window: after tb-vocab ...}\par
\par\medskip\noindent\textbf{Definition} (\texttt{top-dual}).\;$\operatorname{top\text{-}dual}(m, t) :=$
\begin{equation*}
\{\operatorname{tdf\_} \in (\operatorname{vec}(m) \to \mathbb{R}) : (\operatorname{is\text{-}linear\text{-}functional}(m, \operatorname{tdf\_}) \wedge \operatorname{is\text{-}hom\text{-}top\text{-}space}(t, \operatorname{seminorm\text{-}top}(\operatorname{normed\text{-}vector\text{-}space\text{-}as\text{-}module}(\operatorname{rr\text{-}nvs}), \{0\}, (\lambda \operatorname{nvk\_} \in \{0\}.\; \operatorname{vnrm}(\operatorname{rr\text{-}nvs}))), \operatorname{tdf\_}))\}
\end{equation*}

\par\medskip\noindent\textbf{Definition} (\texttt{polar}).\;$\operatorname{polar}(m, t, a) :=$
\begin{equation*}
\{\operatorname{plf\_} \in \operatorname{top\text{-}dual}(m, t) : \forall \operatorname{plx\_} \in a.\; (\operatorname{plf\_}(\operatorname{plx\_}) \leq 1)\}
\end{equation*}

\par\medskip\noindent\textbf{Definition} (\texttt{weak-star-top}).\;$\operatorname{weak\text{-}star\text{-}top}(m, t) :=$
\begin{equation*}
\operatorname{subspace\text{-}top}(\operatorname{tproduct\text{-}top}(\operatorname{vec}(m), (\lambda \operatorname{b6n\_} \in \operatorname{vec}(m).\; \operatorname{metric\text{-}top}(\operatorname{rr\text{-}ms}))), \operatorname{top\text{-}dual}(m, t))
\end{equation*}

\par\medskip\noindent\textbf{Definition} (\texttt{dual-weak-star-top}).\;$\operatorname{dual\text{-}weak\text{-}star\text{-}top}(m) :=$
\begin{equation*}
\operatorname{subspace\text{-}top}(\operatorname{tproduct\text{-}top}(\operatorname{vec}(m), (\lambda \operatorname{b6n\_} \in \operatorname{vec}(m).\; \operatorname{metric\text{-}top}(\operatorname{rr\text{-}ms}))), \operatorname{dual\text{-}vec}(m))
\end{equation*}

\par\medskip\noindent\textbf{Definition} (\texttt{dual-unit-ball}).\;$\operatorname{dual\text{-}unit\text{-}ball}(m) :=$ $\{\operatorname{dbf\_} \in \operatorname{dual\text{-}vec}(m) : (\operatorname{dual\text{-}norm}(m, \operatorname{dbf\_}) \leq 1)\}$

\vnbbookfile{theorem-library/tvs-alaoglu}{theorem-library/tvs-alaoglu}
\noindent{\small tb6.scm -- topological-vector-space.pdf TB-6, the parts Tychonoff unblocks: compactness of pointwise-bounded, pointwise-closed sets of real functions (the engine of Prop 3.25), Banach-Alaoglu (Cor 3.27) for the polar of a neighbourhood of 0 in a real TVS, and Prop 3.29 (the dual unit ball of a normed space).  The weak* ...}\par
\vnbprop{subspace-inclusion-continuous}{\vnbstmt{\begin{array}{@{}l@{}} \forall t, a. \\ \quad (\operatorname{is\text{-}top\text{-}space}(t) \wedge (a \subseteq \operatorname{pts}(t))) \Rightarrow  \\ \quad \quad \operatorname{is\text{-}hom\text{-}top\text{-}space}(\operatorname{subspace\text{-}top}(t, a), t, (\lambda q \in a.\; q)) \end{array}}}
\vnbprop{hom-into-subspace}{\vnbstmt{\begin{array}{@{}l@{}} \forall s, t, a, f. \\ \quad (\operatorname{is\text{-}hom\text{-}top\text{-}space}(s, t, f) \wedge ((a \subseteq \operatorname{pts}(t)) \wedge \forall x \in \operatorname{pts}(s).\; (f(x) \in a))) \Rightarrow  \\ \quad \quad \operatorname{is\text{-}hom\text{-}top\text{-}space}(s, \operatorname{subspace\text{-}top}(t, a), f) \end{array}}}
\vnbprop{closed-in-subspace}{\vnbstmt{\begin{array}{@{}l@{}} \forall t, a, b. \\ \quad (\operatorname{is\text{-}top\text{-}space}(t) \wedge ((a \subseteq \operatorname{pts}(t)) \wedge (\operatorname{is\text{-}top\text{-}closed}(t, b) \wedge (b \subseteq a)))) \Rightarrow  \\ \quad \quad \operatorname{is\text{-}top\text{-}closed}(\operatorname{subspace\text{-}top}(t, a), b) \end{array}}}
\vnbprop{closed-subset-compact}{\vnbstmt{\begin{array}{@{}l@{}} \forall s, b. \\ \quad (\operatorname{is\text{-}compact\text{-}t}(s) \wedge \operatorname{is\text{-}top\text{-}closed}(s, b)) \Rightarrow \operatorname{is\text{-}compact\text{-}t}(\operatorname{subspace\text{-}top}(s, b)) \end{array}}}
\vnbprop{subspace-top-unfold}{\vnbstmt{\begin{array}{@{}l@{}} \forall t, a. \\ \quad (\operatorname{subspace\text{-}top}(t, a) \simeq {[}a, \{w \in \operatorname{power}(a) : \exists u \in \operatorname{opens}(t).\; (w = (u \cap a))\}]) \end{array}}}
\vnbprop{subspace-of-subspace}{\vnbstmt{\begin{array}{@{}l@{}} \forall t, a, b. \\ \quad (\operatorname{is\text{-}top\text{-}space}(t) \wedge ((a \subseteq \operatorname{pts}(t)) \wedge (b \subseteq a))) \Rightarrow  \\ \quad \quad (\operatorname{subspace\text{-}top}(\operatorname{subspace\text{-}top}(t, a), b) \simeq \operatorname{subspace\text{-}top}(t, b)) \end{array}}}
\vnbprop{rr-metric-top-pts}{\vnbstmt{(\operatorname{pts}(\operatorname{metric\text{-}top}(\operatorname{rr\text{-}ms})) \simeq \mathbb{R})}}
\vnbprop{interval-product-in-lines}{\vnbstmt{\begin{array}{@{}l@{}} \forall f, b. \\ \quad ((f \in \mathbf{Set}) \wedge (b \in (f \to \mathbb{R}))) \Rightarrow  \\ \quad \quad \operatorname{tprod\text{-}pts}(f, (\lambda n \in f.\; \operatorname{subspace\text{-}top}(\operatorname{metric\text{-}top}(\operatorname{rr\text{-}ms}), \operatorname{ccint}(-b(n), b(n))))) \\ \quad \quad \subseteq\; \operatorname{tprod\text{-}pts}(f, (\lambda n \in f.\; \operatorname{metric\text{-}top}(\operatorname{rr\text{-}ms}))) \end{array}}}
\vnbprop{interval-product-identity}{\vnbstmt{\begin{array}{@{}l@{}} \forall f, b. \\ \quad ((f \in \mathbf{Set}) \wedge (b \in (f \to \mathbb{R}))) \Rightarrow  \\ \quad \quad \operatorname{is\text{-}hom\text{-}top\text{-}space}(\operatorname{tproduct\text{-}top}(f, (\lambda n \in f.\; \operatorname{subspace\text{-}top}(\operatorname{metric\text{-}top}(\operatorname{rr\text{-}ms}), \operatorname{ccint}(-b(n), b(n))))), \\ \quad \quad \quad \operatorname{subspace\text{-}top}(\operatorname{tproduct\text{-}top}(f, (\lambda n \in f.\; \operatorname{metric\text{-}top}(\operatorname{rr\text{-}ms}))), \\ \quad \quad \quad \quad \operatorname{tprod\text{-}pts}(f, (\lambda n \in f.\; \operatorname{subspace\text{-}top}(\operatorname{metric\text{-}top}(\operatorname{rr\text{-}ms}), \operatorname{ccint}(-b(n), b(n)))))), \\ \quad \quad \quad (\lambda p \in \operatorname{tprod\text{-}pts}(f, (\lambda n \in f.\; \operatorname{subspace\text{-}top}(\operatorname{metric\text{-}top}(\operatorname{rr\text{-}ms}), \operatorname{ccint}(-b(n), b(n))))).\; \\ \quad \quad \quad \quad p)) \end{array}}}
\vnbprop{pointwise-bounded-in-intervals}{\vnbstmt{\begin{array}{@{}l@{}} \forall f, b, a. \\ \quad f \in \mathbf{Set} \wedge \\ \quad b \in (f \to \mathbb{R}) \wedge \\ \quad a \subseteq \operatorname{tprod\text{-}pts}(f, (\lambda n \in f.\; \operatorname{metric\text{-}top}(\operatorname{rr\text{-}ms}))) \wedge \\ \quad \forall y \in f, g \in a.\; (\lvert g(y) \rvert \leq b(y)) \Rightarrow \\ \quad \quad a \subseteq \operatorname{tprod\text{-}pts}(f, (\lambda n \in f.\; \operatorname{subspace\text{-}top}(\operatorname{metric\text{-}top}(\operatorname{rr\text{-}ms}), \operatorname{ccint}(-b(n), b(n))))) \end{array}}}
\vnbprop{interval-factors-compact}{\vnbstmt{\begin{array}{@{}l@{}} \forall f, b.; \forall i \in f. \\ \quad ((f \in \mathbf{Set}) \wedge (b \in (f \to \mathbb{R}))) \Rightarrow  \\ \quad \quad \operatorname{is\text{-}compact\text{-}t}((\lambda n \in f.\; \operatorname{subspace\text{-}top}(\operatorname{metric\text{-}top}(\operatorname{rr\text{-}ms}), \operatorname{ccint}(-b(n), b(n))))(i)) \end{array}}}
\vnbprop{pointwise-bounded-closed-compact}{\vnbstmt{\begin{array}{@{}l@{}} \forall f, b, a. \\ \quad f \in \mathbf{Set} \wedge \\ \quad b \in (f \to \mathbb{R}) \wedge \\ \quad a \subseteq \operatorname{tprod\text{-}pts}(f, (\lambda n \in f.\; \operatorname{metric\text{-}top}(\operatorname{rr\text{-}ms}))) \wedge \\ \quad \forall y \in f, g \in a.\; (\lvert g(y) \rvert \leq b(y)) \wedge \\ \quad \operatorname{is\text{-}top\text{-}closed}(\operatorname{tproduct\text{-}top}(f, (\lambda n \in f.\; \operatorname{metric\text{-}top}(\operatorname{rr\text{-}ms}))), a) \Rightarrow \\ \quad \quad \operatorname{is\text{-}compact\text{-}t}(\operatorname{subspace\text{-}top}(\operatorname{tproduct\text{-}top}(f, (\lambda n \in f.\; \operatorname{metric\text{-}top}(\operatorname{rr\text{-}ms}))), a)) \end{array}}}
\vnbprop{line-product-coord-ball-open}{\vnbstmt{\begin{array}{@{}l@{}} \forall f, x, c, e. \\ \quad ((f \in \mathbf{Set}) \wedge ((x \in f) \wedge ((c \in \mathbb{R}) \wedge ((e \in \mathbb{R}) \wedge (0 < e))))) \Rightarrow  \\ \quad \quad \operatorname{preimage}(\operatorname{tproduct\text{-}top}(f, (\lambda n \in f.\; \operatorname{metric\text{-}top}(\operatorname{rr\text{-}ms}))), \\ \quad \quad \quad \operatorname{tproj}(f, (\lambda n \in f.\; \operatorname{metric\text{-}top}(\operatorname{rr\text{-}ms})), x), \\ \quad \quad \quad \operatorname{ball}(\operatorname{rr\text{-}ms}, c, e)) \\ \quad \quad \in\; \operatorname{opens}(\operatorname{tproduct\text{-}top}(f, (\lambda n \in f.\; \operatorname{metric\text{-}top}(\operatorname{rr\text{-}ms})))) \end{array}}}
\vnbprop{line-product-coord-ball-mem}{\vnbstmt{\begin{array}{@{}l@{}} \forall f, x, c, e, h. \\ \quad ((f \in \mathbf{Set}) \wedge ((x \in f) \wedge ((c \in \mathbb{R}) \wedge (h \in \operatorname{tprod\text{-}pts}(f, (\lambda n \in f.\; \operatorname{metric\text{-}top}(\operatorname{rr\text{-}ms}))))))) \Rightarrow  \\ \quad \quad h \\ \quad \quad \in\; \operatorname{preimage}(\operatorname{tproduct\text{-}top}(f, (\lambda n \in f.\; \operatorname{metric\text{-}top}(\operatorname{rr\text{-}ms}))), \\ \quad \quad \quad \operatorname{tproj}(f, (\lambda n \in f.\; \operatorname{metric\text{-}top}(\operatorname{rr\text{-}ms})), x), \\ \quad \quad \quad \operatorname{ball}(\operatorname{rr\text{-}ms}, c, e)) \\ \quad \quad \Leftrightarrow\; \lvert (c - h(x)) \rvert < e \end{array}}}
\vnbprop{line-product-gt-nbhd}{\vnbstmt{\begin{array}{@{}l@{}} \forall f, g, a. \\ \quad ((f \in \mathbf{Set}) \wedge ((g \in \operatorname{tprod\text{-}pts}(f, (\lambda n \in f.\; \operatorname{metric\text{-}top}(\operatorname{rr\text{-}ms})))) \wedge ((a \in f) \wedge (1 < g(a))))) \Rightarrow  \\ \quad \quad \exists w \in \operatorname{opens}(\operatorname{tproduct\text{-}top}(f, (\lambda n \in f.\; \operatorname{metric\text{-}top}(\operatorname{rr\text{-}ms})))). \\ \quad \quad \quad ((g \in w) \wedge \forall h \in w.\; ((h \in \operatorname{tprod\text{-}pts}(f, (\lambda n \in f.\; \operatorname{metric\text{-}top}(\operatorname{rr\text{-}ms})))) \wedge (1 < h(a)))) \end{array}}}
\vnbprop{line-product-additive-nbhd}{\vnbstmt{\begin{array}{@{}l@{}} \forall f, g, a, b, c. \\ \quad f \in \mathbf{Set} \wedge \\ \quad g \in \operatorname{tprod\text{-}pts}(f, (\lambda n \in f.\; \operatorname{metric\text{-}top}(\operatorname{rr\text{-}ms}))) \wedge \\ \quad a \in f \wedge \\ \quad b \in f \wedge \\ \quad c \in f \wedge \\ \quad \neg (g(c) = (g(a) + g(b))) \Rightarrow \\ \quad \quad \exists w \in \operatorname{opens}(\operatorname{tproduct\text{-}top}(f, (\lambda n \in f.\; \operatorname{metric\text{-}top}(\operatorname{rr\text{-}ms})))). \\ \quad \quad \quad g \in w \wedge \\ \quad \quad \quad \forall h \in w. \\ \quad \quad \quad \quad ((h \in \operatorname{tprod\text{-}pts}(f, (\lambda n \in f.\; \operatorname{metric\text{-}top}(\operatorname{rr\text{-}ms})))) \wedge \neg (h(c) = (h(a) + h(b)))) \end{array}}}
\vnbprop{line-product-homogeneous-nbhd}{\vnbstmt{\begin{array}{@{}l@{}} \forall f, g, a, c, r. \\ \quad f \in \mathbf{Set} \wedge \\ \quad g \in \operatorname{tprod\text{-}pts}(f, (\lambda n \in f.\; \operatorname{metric\text{-}top}(\operatorname{rr\text{-}ms}))) \wedge \\ \quad a \in f \wedge \\ \quad c \in f \wedge \\ \quad r \in \mathbb{R} \wedge \\ \quad \neg (g(c) = (r \cdot g(a))) \Rightarrow \\ \quad \quad \exists w \in \operatorname{opens}(\operatorname{tproduct\text{-}top}(f, (\lambda n \in f.\; \operatorname{metric\text{-}top}(\operatorname{rr\text{-}ms})))). \\ \quad \quad \quad g \in w \wedge \\ \quad \quad \quad \forall h \in w. \\ \quad \quad \quad \quad (h \in \operatorname{tprod\text{-}pts}(f, (\lambda n \in f.\; \operatorname{metric\text{-}top}(\operatorname{rr\text{-}ms})))) \wedge \neg (h(c) = (r \cdot h(a))) \end{array}}}
\vnbprop{top-dual-unfold}{\vnbstmt{\begin{array}{@{}l@{}} \forall m, t. \\ \quad (\operatorname{top\text{-}dual}(m, t) \simeq \{f \in (\operatorname{vec}(m) \to \mathbb{R}) : (\operatorname{is\text{-}linear\text{-}functional}(m, f) \wedge \operatorname{is\text{-}hom\text{-}top\text{-}space}(t, \operatorname{seminorm\text{-}top}(\operatorname{normed\text{-}vector\text{-}space\text{-}as\text{-}module}(\operatorname{rr\text{-}nvs}), \{0\}, (\lambda k \in \{0\}.\; \operatorname{vnrm}(\operatorname{rr\text{-}nvs}))), f))\}) \end{array}}}
\vnbprop{polar-unfold}{\vnbstmt{\begin{array}{@{}l@{}} \forall m, t, a. \\ \quad (\operatorname{polar}(m, t, a) \simeq \{f \in \operatorname{top\text{-}dual}(m, t) : \forall x \in a.\; (f(x) \leq 1)\}) \end{array}}}
\vnbprop{weak-star-top-unfold}{\vnbstmt{\begin{array}{@{}l@{}} \forall m, t. \\ \quad (\operatorname{weak\text{-}star\text{-}top}(m, t) \simeq \operatorname{subspace\text{-}top}(\operatorname{tproduct\text{-}top}(\operatorname{vec}(m), (\lambda n \in \operatorname{vec}(m).\; \operatorname{metric\text{-}top}(\operatorname{rr\text{-}ms}))), \operatorname{top\text{-}dual}(m, t))) \end{array}}}
\vnbprop{top-dual-mem}{\vnbstmt{\begin{array}{@{}l@{}} \forall m, t, z. \\ \quad z \in \operatorname{top\text{-}dual}(m, t) \\ \quad \Leftrightarrow\; z \in (\operatorname{vec}(m) \to \mathbb{R}) \wedge \\ \quad \operatorname{is\text{-}linear\text{-}functional}(m, z) \wedge \\ \quad \operatorname{is\text{-}hom\text{-}top\text{-}space}(t, \\ \quad \quad \operatorname{seminorm\text{-}top}(\operatorname{normed\text{-}vector\text{-}space\text{-}as\text{-}module}(\operatorname{rr\text{-}nvs}), \\ \quad \quad \quad \{0\}, (\lambda k \in \{0\}.\; \operatorname{vnrm}(\operatorname{rr\text{-}nvs}))), \\ \quad \quad z) \end{array}}}
\vnbprop{polar-mem}{\vnbstmt{\begin{array}{@{}l@{}} \forall m, t, a, z. \\ \quad (z \in \operatorname{polar}(m, t, a)) \Leftrightarrow ((z \in \operatorname{top\text{-}dual}(m, t)) \wedge \forall x \in a.\; (z(x) \leq 1)) \end{array}}}
\vnbprop{fun-in-line-product}{\vnbstmt{\begin{array}{@{}l@{}} \forall a, f. \\ \quad ((a \in \mathbf{Set}) \wedge (f \in (a \to \mathbb{R}))) \Rightarrow  \\ \quad \quad f \in \operatorname{tprod\text{-}pts}(a, (\lambda n \in a.\; \operatorname{metric\text{-}top}(\operatorname{rr\text{-}ms}))) \end{array}}}
\vnbprop{line-product-fun}{\vnbstmt{\begin{array}{@{}l@{}} \forall a, f. \\ \quad ((a \in \mathbf{Set}) \wedge (f \in \operatorname{tprod\text{-}pts}(a, (\lambda n \in a.\; \operatorname{metric\text{-}top}(\operatorname{rr\text{-}ms}))))) \Rightarrow  \\ \quad \quad f \in (a \to \mathbb{R}) \end{array}}}
\vnbprop{polar-bound-continuous}{\vnbstmt{\begin{array}{@{}l@{}} \forall m, t, v, w, f. \\ \quad \operatorname{is\text{-}tvs}(\operatorname{rr\text{-}normed\text{-}field}, m, t) \wedge \\ \quad v \subseteq \operatorname{vec}(m) \wedge \\ \quad w \in \operatorname{opens}(t) \wedge \\ \quad \operatorname{vzero}(m) \in w \wedge \\ \quad w \subseteq v \wedge \\ \quad \operatorname{is\text{-}linear\text{-}functional}(m, f) \wedge \\ \quad \forall x \in v.\; (f(x) \leq 1) \Rightarrow \\ \quad \quad \operatorname{is\text{-}hom\text{-}top\text{-}space}(t, \\ \quad \quad \quad \operatorname{seminorm\text{-}top}(\operatorname{normed\text{-}vector\text{-}space\text{-}as\text{-}module}(\operatorname{rr\text{-}nvs}), \\ \quad \quad \quad \quad \{0\}, (\lambda k \in \{0\}.\; \operatorname{vnrm}(\operatorname{rr\text{-}nvs}))), \\ \quad \quad \quad f) \end{array}}}
\vnbprop{polar-pointwise-bounded}{\vnbstmt{\begin{array}{@{}l@{}} \forall m, t, v, w. \\ \quad (\operatorname{is\text{-}tvs}(\operatorname{rr\text{-}normed\text{-}field}, m, t) \wedge ((v \subseteq \operatorname{vec}(m)) \wedge ((w \in \operatorname{opens}(t)) \wedge ((\operatorname{vzero}(m) \in w) \wedge (w \subseteq v))))) \Rightarrow  \\ \quad \quad (\lambda x \in \operatorname{vec}(m).\; (\varepsilon(\operatorname{mink\text{-}set}(m, v, x)) + \varepsilon(\operatorname{mink\text{-}set}(m, v, \operatorname{act}(m)(-1, x))))) \\ \quad \quad \in\; (\operatorname{vec}(m) \to \mathbb{R}) \wedge \\ \quad \quad \forall y \in \operatorname{vec}(m), g \in \operatorname{polar}(m, t, v). \\ \quad \quad \quad \lvert g(y) \rvert \\ \quad \quad \quad \leq\; (\lambda x \in \operatorname{vec}(m).\; (\varepsilon(\operatorname{mink\text{-}set}(m, v, x)) + \varepsilon(\operatorname{mink\text{-}set}(m, v, \operatorname{act}(m)(-1, x)))))(y) \end{array}}}
\vnbprop{line-product-facts}{\vnbstmt{\begin{array}{@{}l@{}} \forall f \in \mathbf{Set}. \\ \quad \operatorname{is\text{-}top\text{-}space}(\operatorname{tproduct\text{-}top}(f, (\lambda n \in f.\; \operatorname{metric\text{-}top}(\operatorname{rr\text{-}ms})))) \wedge \\ \quad (\operatorname{pts}(\operatorname{tproduct\text{-}top}(f, (\lambda n \in f.\; \operatorname{metric\text{-}top}(\operatorname{rr\text{-}ms})))) \simeq \operatorname{tprod\text{-}pts}(f, (\lambda n \in f.\; \operatorname{metric\text{-}top}(\operatorname{rr\text{-}ms})))) \wedge \\ \quad \operatorname{tprod\text{-}pts}(f, (\lambda n \in f.\; \operatorname{metric\text{-}top}(\operatorname{rr\text{-}ms}))) \in \mathbf{Set} \end{array}}}
\vnbprop{polar-of-no-failure}{\vnbstmt{\begin{array}{@{}l@{}} \forall m, t, v, w, g. \\ \quad \operatorname{is\text{-}tvs}(\operatorname{rr\text{-}normed\text{-}field}, m, t) \wedge \\ \quad v \subseteq \operatorname{vec}(m) \wedge \\ \quad w \in \operatorname{opens}(t) \wedge \\ \quad \operatorname{vzero}(m) \in w \wedge \\ \quad w \subseteq v \wedge \\ \quad g \in \operatorname{tprod\text{-}pts}(\operatorname{vec}(m), (\lambda n \in \operatorname{vec}(m).\; \operatorname{metric\text{-}top}(\operatorname{rr\text{-}ms}))) \wedge \\ \quad \neg \exists a \in v.\; (1 < g(a)) \wedge \\ \quad \neg \exists a, b \in \operatorname{vec}(m).\; \neg (g(\operatorname{vadd}(m)(a, b)) = (g(a) + g(b))) \wedge \\ \quad \neg \exists r \in \mathbb{R}, a \in \operatorname{vec}(m).\; \neg (g(\operatorname{act}(m)(r, a)) = (r \cdot g(a))) \Rightarrow \\ \quad \quad g \in \operatorname{polar}(m, t, v) \end{array}}}
\vnbprop{polar-closed-in-line-product}{\vnbstmt{\begin{array}{@{}l@{}} \forall m, t, v, w. \\ \quad (\operatorname{is\text{-}tvs}(\operatorname{rr\text{-}normed\text{-}field}, m, t) \wedge ((v \subseteq \operatorname{vec}(m)) \wedge ((w \in \operatorname{opens}(t)) \wedge ((\operatorname{vzero}(m) \in w) \wedge (w \subseteq v))))) \Rightarrow  \\ \quad \quad \operatorname{is\text{-}top\text{-}closed}(\operatorname{tproduct\text{-}top}(\operatorname{vec}(m), (\lambda n \in \operatorname{vec}(m).\; \operatorname{metric\text{-}top}(\operatorname{rr\text{-}ms}))), \operatorname{polar}(m, t, v)) \end{array}}}
\vnbprop{banach-alaoglu}{\vnbstmt{\begin{array}{@{}l@{}} \forall m, t, v. \\ \quad \operatorname{is\text{-}tvs}(\operatorname{rr\text{-}normed\text{-}field}, m, t) \wedge \\ \quad v \subseteq \operatorname{vec}(m) \wedge \\ \quad \exists \mathit{b6v} \in \operatorname{opens}(t).\; ((\operatorname{vzero}(m) \in \mathit{b6v}) \wedge (\mathit{b6v} \subseteq v)) \Rightarrow \\ \quad \quad \operatorname{is\text{-}compact\text{-}t}(\operatorname{subspace\text{-}top}(\operatorname{weak\text{-}star\text{-}top}(m, t), \operatorname{polar}(m, t, v))) \end{array}}}
\vnbprop{line-product-abs-gt-nbhd}{\vnbstmt{\begin{array}{@{}l@{}} \forall f, g, a, k. \\ \quad f \in \mathbf{Set} \wedge \\ \quad g \in \operatorname{tprod\text{-}pts}(f, (\lambda n \in f.\; \operatorname{metric\text{-}top}(\operatorname{rr\text{-}ms}))) \wedge \\ \quad a \in f \wedge \\ \quad k \in \mathbb{R} \wedge \\ \quad k < \lvert g(a) \rvert \Rightarrow \\ \quad \quad \exists w \in \operatorname{opens}(\operatorname{tproduct\text{-}top}(f, (\lambda n \in f.\; \operatorname{metric\text{-}top}(\operatorname{rr\text{-}ms})))). \\ \quad \quad \quad g \in w \wedge \\ \quad \quad \quad \forall h \in w. \\ \quad \quad \quad \quad (h \in \operatorname{tprod\text{-}pts}(f, (\lambda n \in f.\; \operatorname{metric\text{-}top}(\operatorname{rr\text{-}ms})))) \wedge (k < \lvert h(a) \rvert) \end{array}}}
\vnbprop{dual-unit-ball-unfold}{\vnbstmt{\forall m.\; (\operatorname{dual\text{-}unit\text{-}ball}(m) \simeq \{f \in \operatorname{dual\text{-}vec}(m) : (\operatorname{dual\text{-}norm}(m, f) \leq 1)\})}}
\vnbprop{dual-weak-star-top-unfold}{\vnbstmt{\begin{array}{@{}l@{}} \forall m. \\ \quad (\operatorname{dual\text{-}weak\text{-}star\text{-}top}(m) \simeq \operatorname{subspace\text{-}top}(\operatorname{tproduct\text{-}top}(\operatorname{vec}(m), (\lambda n \in \operatorname{vec}(m).\; \operatorname{metric\text{-}top}(\operatorname{rr\text{-}ms}))), \operatorname{dual\text{-}vec}(m))) \end{array}}}
\vnbprop{dual-unit-ball-mem}{\vnbstmt{\begin{array}{@{}l@{}} \forall m, f. \\ \quad \operatorname{is\text{-}normed\text{-}vector\text{-}space}(m) \Rightarrow  \\ \quad \quad f \in \operatorname{dual\text{-}unit\text{-}ball}(m) \\ \quad \quad \Leftrightarrow\; f \in (\operatorname{vec}(m) \to \mathbb{R}) \wedge \\ \quad \quad \operatorname{is\text{-}linear\text{-}functional}(m, f) \wedge \\ \quad \quad \forall x \in \operatorname{vec}(m).\; (\lvert f(x) \rvert \leq \operatorname{vnrm}(m)(x)) \end{array}}}
\vnbprop{dual-unit-ball-of-no-failure}{\vnbstmt{\begin{array}{@{}l@{}} \forall m, g. \\ \quad \operatorname{is\text{-}normed\text{-}vector\text{-}space}(m) \wedge \\ \quad g \in \operatorname{tprod\text{-}pts}(\operatorname{vec}(m), (\lambda n \in \operatorname{vec}(m).\; \operatorname{metric\text{-}top}(\operatorname{rr\text{-}ms}))) \wedge \\ \quad \neg \exists a \in \operatorname{vec}(m).\; (\operatorname{vnrm}(m)(a) < \lvert g(a) \rvert) \wedge \\ \quad \neg \exists a, b \in \operatorname{vec}(m).\; \neg (g(\operatorname{vadd}(m)(a, b)) = (g(a) + g(b))) \wedge \\ \quad \neg \exists r \in \mathbb{R}, a \in \operatorname{vec}(m).\; \neg (g(\operatorname{act}(m)(r, a)) = (r \cdot g(a))) \Rightarrow \\ \quad \quad g \in \operatorname{dual\text{-}unit\text{-}ball}(m) \end{array}}}
\vnbprop{dual-unit-ball-closed-in-line-product}{\vnbstmt{\begin{array}{@{}l@{}} \forall m. \\ \quad \operatorname{is\text{-}normed\text{-}vector\text{-}space}(m) \Rightarrow  \\ \quad \quad \operatorname{is\text{-}top\text{-}closed}(\operatorname{tproduct\text{-}top}(\operatorname{vec}(m), (\lambda n \in \operatorname{vec}(m).\; \operatorname{metric\text{-}top}(\operatorname{rr\text{-}ms}))), \operatorname{dual\text{-}unit\text{-}ball}(m)) \end{array}}}
\vnbprop{dual-unit-ball-weak-star-compact}{\vnbstmt{\begin{array}{@{}l@{}} \forall m. \\ \quad \operatorname{is\text{-}normed\text{-}vector\text{-}space}(m) \Rightarrow  \\ \quad \quad \operatorname{is\text{-}compact\text{-}t}(\operatorname{subspace\text{-}top}(\operatorname{dual\text{-}weak\text{-}star\text{-}top}(m), \operatorname{dual\text{-}unit\text{-}ball}(m))) \end{array}}}

\vnbbookfile{theorem-library/cc-product}{theorem-library/cc-product}
\noindent{\small cc-product.scm -- wave of 2026-11-09/10 (agent BA-1D), wired 2026-11-10 from scratchpad/ba1d-1109/cc-product.scm. Library: analysis.}\par
\vnbprop{cc-metric-top-pts}{\vnbstmt{(\operatorname{pts}(\operatorname{metric\text{-}top}(\operatorname{cc\text{-}ms})) \simeq \mathbb{C})}}
\vnbprop{cc-product-family}{\vnbstmt{\forall f \in \mathbf{Set}.\; \operatorname{is\text{-}top\text{-}family}(f, (\lambda x \in f.\; \operatorname{metric\text{-}top}(\operatorname{cc\text{-}ms})))}}
\vnbprop{fun-in-cc-product}{\vnbstmt{\begin{array}{@{}l@{}} \forall a, f. \\ \quad ((a \in \mathbf{Set}) \wedge (f \in (a \to \mathbb{C}))) \Rightarrow  \\ \quad \quad f \in \operatorname{tprod\text{-}pts}(a, (\lambda x \in a.\; \operatorname{metric\text{-}top}(\operatorname{cc\text{-}ms}))) \end{array}}}
\vnbprop{cc-product-eval-continuous}{\vnbstmt{\begin{array}{@{}l@{}} \forall f, a, x. \\ \quad ((f \in \mathbf{Set}) \wedge ((a \subseteq \operatorname{tprod\text{-}pts}(f, (\lambda \mathit{bgx} \in f.\; \operatorname{metric\text{-}top}(\operatorname{cc\text{-}ms})))) \wedge (x \in f))) \Rightarrow  \\ \quad \quad \operatorname{is\text{-}hom\text{-}top\text{-}space}(\operatorname{subspace\text{-}top}(\operatorname{tproduct\text{-}top}(f, (\lambda \mathit{bgx} \in f.\; \operatorname{metric\text{-}top}(\operatorname{cc\text{-}ms}))), a), \\ \quad \quad \quad \operatorname{metric\text{-}top}(\operatorname{cc\text{-}ms}), (\lambda c \in a.\; c(x))) \end{array}}}

\vnbbookfile{theorem-library/cc-top-continuity}{theorem-library/cc-top-continuity}
\noindent{\small cc-top-continuity.scm -- wave of 2026-11-09/10 (agent BA-1F), wired 2026-11-10 from scratchpad/ba1f-1109/cc-top-continuity.scm. Library: analysis.}\par
\vnbprop{cc-hom-local}{\vnbstmt{\begin{array}{@{}l@{}} \forall s, a.; \forall p \in \operatorname{pts}(s), e \in \mathbb{R}. \\ \quad (\operatorname{is\text{-}hom\text{-}top\text{-}space}(s, \operatorname{metric\text{-}top}(\operatorname{cc\text{-}ms}), a) \wedge (0 < e)) \Rightarrow  \\ \quad \quad \exists w \in \operatorname{opens}(s). \\ \quad \quad \quad (p \in w) \wedge \forall q \in w.\; ((q \in \operatorname{pts}(s)) \wedge (\operatorname{magnitude}((a(p) - a(q))) < e)) \end{array}}}
\vnbprop{cc-hom-of-local}{\vnbstmt{\begin{array}{@{}l@{}} \forall s, a. \\ \quad \operatorname{is\text{-}top\text{-}space}(s) \wedge \\ \quad a \in (\operatorname{pts}(s) \to \mathbb{C}) \wedge \\ \quad \forall p \in \operatorname{pts}(s), e \in \mathbb{R}. \\ \quad \quad (0 < e) \Rightarrow  \\ \quad \quad \quad \exists w \in \operatorname{opens}(s). \\ \quad \quad \quad \quad (p \in w) \wedge \forall q \in w.\; ((q \in \operatorname{pts}(s)) \wedge (\operatorname{magnitude}((a(p) - a(q))) < e)) \Rightarrow \\ \quad \quad \operatorname{is\text{-}hom\text{-}top\text{-}space}(s, \operatorname{metric\text{-}top}(\operatorname{cc\text{-}ms}), a) \end{array}}}
\vnbprop{cc-hom-const}{\vnbstmt{\begin{array}{@{}l@{}} \forall s, c. \\ \quad (\operatorname{is\text{-}top\text{-}space}(s) \wedge (c \in \mathbb{C})) \Rightarrow  \\ \quad \quad \operatorname{is\text{-}hom\text{-}top\text{-}space}(s, \operatorname{metric\text{-}top}(\operatorname{cc\text{-}ms}), (\lambda v \in \operatorname{pts}(s).\; c)) \end{array}}}
\vnbprop{cc-hom-add}{\vnbstmt{\begin{array}{@{}l@{}} \forall s, a, b. \\ \quad (\operatorname{is\text{-}hom\text{-}top\text{-}space}(s, \operatorname{metric\text{-}top}(\operatorname{cc\text{-}ms}), a) \wedge \operatorname{is\text{-}hom\text{-}top\text{-}space}(s, \operatorname{metric\text{-}top}(\operatorname{cc\text{-}ms}), b)) \Rightarrow  \\ \quad \quad \operatorname{is\text{-}hom\text{-}top\text{-}space}(s, \operatorname{metric\text{-}top}(\operatorname{cc\text{-}ms}), (\lambda v \in \operatorname{pts}(s).\; (a(v) + b(v)))) \end{array}}}
\vnbprop{cc-hom-equalizer-closed}{\vnbstmt{\begin{array}{@{}l@{}} \forall s, a, b. \\ \quad (\operatorname{is\text{-}hom\text{-}top\text{-}space}(s, \operatorname{metric\text{-}top}(\operatorname{cc\text{-}ms}), a) \wedge \operatorname{is\text{-}hom\text{-}top\text{-}space}(s, \operatorname{metric\text{-}top}(\operatorname{cc\text{-}ms}), b)) \Rightarrow  \\ \quad \quad \operatorname{is\text{-}top\text{-}closed}(s, \{p \in \operatorname{pts}(s) : (a(p) = b(p))\}) \end{array}}}
\vnbprop{cc-magnitude-prod-bound}{\vnbstmt{\begin{array}{@{}l@{}} \forall u, v, p, q, m, k, t, e. \\ \quad u \in \mathbb{C} \wedge \\ \quad v \in \mathbb{C} \wedge \\ \quad p \in \mathbb{C} \wedge \\ \quad q \in \mathbb{C} \wedge \\ \quad m \in \mathbb{R} \wedge \\ \quad k \in \mathbb{R} \wedge \\ \quad t \in \mathbb{R} \wedge \\ \quad e \in \mathbb{R} \wedge \\ \quad \operatorname{magnitude}(u) \leq m \wedge \\ \quad \operatorname{magnitude}(q) \leq k \wedge \\ \quad \operatorname{magnitude}((u - p)) \leq t \wedge \\ \quad \operatorname{magnitude}((v - q)) \leq t \wedge \\ \quad ((m + k) \cdot t) \leq e \Rightarrow \\ \quad \quad \operatorname{magnitude}(((u \cdot v) - (p \cdot q))) \leq e \end{array}}}
\vnbprop{cc-hom-mul}{\vnbstmt{\begin{array}{@{}l@{}} \forall s, a, b. \\ \quad (\operatorname{is\text{-}hom\text{-}top\text{-}space}(s, \operatorname{metric\text{-}top}(\operatorname{cc\text{-}ms}), a) \wedge \operatorname{is\text{-}hom\text{-}top\text{-}space}(s, \operatorname{metric\text{-}top}(\operatorname{cc\text{-}ms}), b)) \Rightarrow  \\ \quad \quad \operatorname{is\text{-}hom\text{-}top\text{-}space}(s, \operatorname{metric\text{-}top}(\operatorname{cc\text{-}ms}), (\lambda v \in \operatorname{pts}(s).\; (a(v) \cdot b(v)))) \end{array}}}
\vnbprop{cc-hom-separate}{\vnbstmt{\begin{array}{@{}l@{}} \forall s, a, b, h, p. \\ \quad \operatorname{is\text{-}hom\text{-}top\text{-}space}(s, \operatorname{metric\text{-}top}(\operatorname{cc\text{-}ms}), a) \wedge \\ \quad \operatorname{is\text{-}hom\text{-}top\text{-}space}(s, \operatorname{metric\text{-}top}(\operatorname{cc\text{-}ms}), b) \wedge \\ \quad \forall c \in h.\; (a(c) = b(c)) \wedge \\ \quad p \in \operatorname{pts}(s) \wedge \\ \quad \neg (a(p) = b(p)) \Rightarrow \\ \quad \quad \exists w \in \operatorname{opens}(s).\; ((p \in w) \wedge (w \subseteq \operatorname{complement\text{-}in}(\operatorname{pts}(s), h))) \end{array}}}

\vnbbookfile{theorem-library/top-power-compact}{theorem-library/top-power-compact}
\noindent{\small top-power-compact.scm -- wave of 2026-11-09/10 (agent BA-1F), wired 2026-11-10 from scratchpad/ba1f-1109/top-power-compact.scm. Library: analysis.}\par
\vnbprop{confined-product-in-power}{\vnbstmt{\begin{array}{@{}l@{}} \forall f, t, k. \\ \quad f \in \mathbf{Set} \wedge \\ \quad \operatorname{is\text{-}top\text{-}space}(t) \wedge \\ \quad \forall y \in f.\; (k(y) \in \operatorname{power}(\operatorname{pts}(t))) \wedge \\ \quad \forall y \in f.\; \operatorname{is\text{-}compact\text{-}t}(\operatorname{subspace\text{-}top}(t, k(y))) \Rightarrow \\ \quad \quad \operatorname{tprod\text{-}pts}(f, (\lambda n \in f.\; \operatorname{subspace\text{-}top}(t, k(n)))) \subseteq \operatorname{tprod\text{-}pts}(f, (\lambda n \in f.\; t)) \end{array}}}
\vnbprop{confined-product-identity}{\vnbstmt{\begin{array}{@{}l@{}} \forall f, t, k. \\ \quad f \in \mathbf{Set} \wedge \\ \quad \operatorname{is\text{-}top\text{-}space}(t) \wedge \\ \quad \forall y \in f.\; (k(y) \in \operatorname{power}(\operatorname{pts}(t))) \wedge \\ \quad \forall y \in f.\; \operatorname{is\text{-}compact\text{-}t}(\operatorname{subspace\text{-}top}(t, k(y))) \Rightarrow \\ \quad \quad \operatorname{is\text{-}hom\text{-}top\text{-}space}(\operatorname{tproduct\text{-}top}(f, (\lambda n \in f.\; \operatorname{subspace\text{-}top}(t, k(n)))), \\ \quad \quad \quad \operatorname{subspace\text{-}top}(\operatorname{tproduct\text{-}top}(f, (\lambda n \in f.\; t)), \\ \quad \quad \quad \quad \operatorname{tprod\text{-}pts}(f, (\lambda n \in f.\; \operatorname{subspace\text{-}top}(t, k(n))))), \\ \quad \quad \quad (\lambda p \in \operatorname{tprod\text{-}pts}(f, (\lambda n \in f.\; \operatorname{subspace\text{-}top}(t, k(n)))).\; p)) \end{array}}}
\vnbprop{confined-in-product}{\vnbstmt{\begin{array}{@{}l@{}} \forall f, t, k, a. \\ \quad f \in \mathbf{Set} \wedge \\ \quad \operatorname{is\text{-}top\text{-}space}(t) \wedge \\ \quad \forall y \in f.\; (k(y) \in \operatorname{power}(\operatorname{pts}(t))) \wedge \\ \quad \forall y \in f.\; \operatorname{is\text{-}compact\text{-}t}(\operatorname{subspace\text{-}top}(t, k(y))) \wedge \\ \quad a \subseteq \operatorname{tprod\text{-}pts}(f, (\lambda n \in f.\; t)) \wedge \\ \quad \forall y \in f, g \in a.\; (g(y) \in k(y)) \Rightarrow \\ \quad \quad a \subseteq \operatorname{tprod\text{-}pts}(f, (\lambda n \in f.\; \operatorname{subspace\text{-}top}(t, k(n)))) \end{array}}}
\vnbprop{confined-closed-compact}{\vnbstmt{\begin{array}{@{}l@{}} \forall f, t, k, a. \\ \quad f \in \mathbf{Set} \wedge \\ \quad \operatorname{is\text{-}top\text{-}space}(t) \wedge \\ \quad \forall y \in f.\; (k(y) \in \operatorname{power}(\operatorname{pts}(t))) \wedge \\ \quad \forall y \in f.\; \operatorname{is\text{-}compact\text{-}t}(\operatorname{subspace\text{-}top}(t, k(y))) \wedge \\ \quad a \subseteq \operatorname{tprod\text{-}pts}(f, (\lambda n \in f.\; t)) \wedge \\ \quad \forall y \in f, g \in a.\; (g(y) \in k(y)) \wedge \\ \quad \operatorname{is\text{-}top\text{-}closed}(\operatorname{tproduct\text{-}top}(f, (\lambda n \in f.\; t)), a) \Rightarrow \\ \quad \quad \operatorname{is\text{-}compact\text{-}t}(\operatorname{subspace\text{-}top}(\operatorname{tproduct\text{-}top}(f, (\lambda n \in f.\; t)), a)) \end{array}}}

\vnbbookfile{structure-library/weak-star-separable}{structure-library/weak-star-separable}
\noindent{\small tb6s-vocab.scm -- THE SEPARABLE BANACH-ALAOGLU (topological-vector-space.pdf Remark 3.30): definitions only. Agent TB-6 (worker-08), 2026-10-27 (NEW DEFINITIONS, the user's to confirm). DESTINATION: structure-library/weak-star-separable.scm.  Library: analysis. Load window: after product-metric (PRODUCT-METRIC, ...}\par
\par\medskip\noindent\textbf{Definition} (\texttt{pullback-ms}).\;$\operatorname{pullback\text{-}ms}(\mathit{mt}, \mathit{ss}, \mathit{jj}) :=$
\begin{equation*}
{[}\mathit{ss}, (\lambda {[}\operatorname{pbu\_}, \operatorname{pbv\_}] \in (\mathit{ss} \times \mathit{ss}).\; \operatorname{dist}(\mathit{mt})(\operatorname{jj}(\operatorname{pbu\_}), \operatorname{jj}(\operatorname{pbv\_})))]
\end{equation*}

\par\medskip\noindent\textbf{Definition} (\texttt{dbi}).\;$\operatorname{dbi}(m, \mathit{xs}) :=$
\begin{equation*}
(\lambda \operatorname{dbn\_} \in \mathbb{N}.\; \operatorname{subspace\text{-}ms}(\operatorname{rr\text{-}ms}, \operatorname{ccint}(-\operatorname{vnrm}(m)(\operatorname{xs}(\operatorname{dbn\_})), \operatorname{vnrm}(m)(\operatorname{xs}(\operatorname{dbn\_})))))
\end{equation*}

\par\medskip\noindent\textbf{Definition} (\texttt{dbj}).\;$\operatorname{dbj}(m, \mathit{xs}) :=$
\begin{equation*}
(\lambda \operatorname{dbf\_} \in \operatorname{dual\text{-}unit\text{-}ball}(m).\; (\lambda \operatorname{dbn\_} \in \mathbb{N}.\; \operatorname{dbf\_}(\operatorname{xs}(\operatorname{dbn\_}))))
\end{equation*}

\par\medskip\noindent\textbf{Definition} (\texttt{dual-ball-ms}).\;$\operatorname{dual\text{-}ball\text{-}ms}(m, \mathit{xs}) :=$
\begin{equation*}
\operatorname{pullback\text{-}ms}(\operatorname{product\text{-}metric}(\operatorname{dbi}(m, \mathit{xs})), \operatorname{dual\text{-}unit\text{-}ball}(m), \operatorname{dbj}(m, \mathit{xs}))
\end{equation*}


\vnbbookfile{theorem-library/subseq-fun-typing}{theorem-library/subseq-fun-typing}
\noindent{\small subseq-fun-typing.scm -- lemmas MOVED 2026-10-27 (integration of the 10-26/27 wave) out of theorem-library/ca6-hurwitz.scm, so that files of library analysis may cite them (the library gate). The proofs are the originals, verbatim; the helper defines are a copy of ca6-hurwitz.scm's.  Loaded immediately before ...}\par
\vnbprop{subseq-fun-typing}{\vnbstmt{\begin{array}{@{}l@{}} \forall f, a, i. \\ \quad ((f \in (\mathbb{N} \to a)) \wedge \operatorname{strictly\text{-}mono\text{-}nn}(i)) \Rightarrow (\operatorname{subseq}(f, i) \in (\mathbb{N} \to a)) \end{array}}}

\vnbbookfile{theorem-library/tvs-alaoglu-separable}{theorem-library/tvs-alaoglu-separable}
\noindent{\small tb6s.scm -- topological-vector-space.pdf Remark 3.30: the SEPARABLE Banach-Alaoglu theorem by the METRIC route. For a normed space with a dense sequence (x\_n), the dual unit ball B* carries the metric d(f, g) = sum\_n 2\textasciicircum{}-(n+1) rho(f(x\_n), g(x\_n)) (PRODUCT-METRIC of the intervals [-||x\_n||, ||x\_n||], pulled back); its ...}\par
\vnbprop{pullback-pts}{\vnbstmt{\forall t.; \forall s \in \mathbf{Set}, j.\; (\operatorname{pts}(\operatorname{pullback\text{-}ms}(t, s, j)) \simeq s)}}
\vnbprop{pullback-dist-slot}{\vnbstmt{\begin{array}{@{}l@{}} \forall t.; \forall s \in \mathbf{Set}, j. \\ \quad (\operatorname{dist}(\operatorname{pullback\text{-}ms}(t, s, j)) \simeq (\lambda {[}\operatorname{pbu\_}, \operatorname{pbv\_}] \in (s \times s).\; \operatorname{dist}(t)(j(\operatorname{pbu\_}), j(\operatorname{pbv\_})))) \end{array}}}
\vnbprop{pullback-dist}{\vnbstmt{\begin{array}{@{}l@{}} \forall t, s, j, u, v. \\ \quad ((s \in \mathbf{Set}) \wedge ((u \in s) \wedge (v \in s))) \Rightarrow  \\ \quad \quad (\operatorname{dist}(\operatorname{pullback\text{-}ms}(t, s, j))(u, v) \simeq \operatorname{dist}(t)(j(u), j(v))) \end{array}}}
\vnbprop{pullback-is-metric-space}{\vnbstmt{\begin{array}{@{}l@{}} \forall t, s, j. \\ \quad \operatorname{is\text{-}metric\text{-}space}(t) \wedge \\ \quad s \in \mathbf{Set} \wedge \\ \quad j \in (s \to \operatorname{pts}(t)) \wedge \\ \quad \forall x, y \in s.\; ((j(x) = j(y)) \Rightarrow (x = y)) \Rightarrow \\ \quad \quad \operatorname{is\text{-}metric\text{-}space}(\operatorname{pullback\text{-}ms}(t, s, j)) \end{array}}}
\vnbprop{pullback-converges-to}{\vnbstmt{\begin{array}{@{}l@{}} \forall t, s, j, q, a. \\ \quad \operatorname{is\text{-}metric\text{-}space}(t) \wedge \\ \quad s \in \mathbf{Set} \wedge \\ \quad j \in (s \to \operatorname{pts}(t)) \wedge \\ \quad \forall x, y \in s.\; ((j(x) = j(y)) \Rightarrow (x = y)) \wedge \\ \quad q \in (\mathbb{N} \to s) \wedge \\ \quad a \in s \Rightarrow \\ \quad \quad \operatorname{converges\text{-}to}(\operatorname{pullback\text{-}ms}(t, s, j), q, a) \Leftrightarrow \operatorname{converges\text{-}to}(t, \operatorname{compose}(j, q), j(a)) \end{array}}}
\vnbprop{dual-ball-sub}{\vnbstmt{\begin{array}{@{}l@{}} \forall m, f, z, y. \\ \quad (\operatorname{is\text{-}normed\text{-}vector\text{-}space}(m) \wedge ((f \in \operatorname{dual\text{-}unit\text{-}ball}(m)) \wedge ((z \in \operatorname{vec}(m)) \wedge (y \in \operatorname{vec}(m))))) \Rightarrow  \\ \quad \quad f(\operatorname{vadd}(m)(z, \operatorname{vneg}(m)(y))) = (f(z) - f(y)) \end{array}}}
\vnbprop{dual-ball-lipschitz}{\vnbstmt{\begin{array}{@{}l@{}} \forall m, f, z, y. \\ \quad (\operatorname{is\text{-}normed\text{-}vector\text{-}space}(m) \wedge ((f \in \operatorname{dual\text{-}unit\text{-}ball}(m)) \wedge ((z \in \operatorname{vec}(m)) \wedge (y \in \operatorname{vec}(m))))) \Rightarrow  \\ \quad \quad \lvert (f(z) - f(y)) \rvert \leq \operatorname{vnrm}(m)(\operatorname{vadd}(m)(z, \operatorname{vneg}(m)(y))) \end{array}}}
\vnbprop{dbi-at}{\vnbstmt{\begin{array}{@{}l@{}} \forall m, s, n. \\ \quad (\operatorname{is\text{-}normed\text{-}vector\text{-}space}(m) \wedge ((s \in (\mathbb{N} \to \operatorname{vec}(m))) \wedge (n \in \mathbb{N}))) \Rightarrow  \\ \quad \quad (\operatorname{dbi}(m, s)(n) \simeq \operatorname{subspace\text{-}ms}(\operatorname{rr\text{-}ms}, \operatorname{ccint}(-\operatorname{vnrm}(m)(s(n)), \operatorname{vnrm}(m)(s(n))))) \end{array}}}
\vnbprop{dbi-seq-compact}{\vnbstmt{\begin{array}{@{}l@{}} \forall m, s. \\ \quad (\operatorname{is\text{-}normed\text{-}vector\text{-}space}(m) \wedge (s \in (\mathbb{N} \to \operatorname{vec}(m)))) \Rightarrow  \\ \quad \quad \operatorname{is\text{-}ms\text{-}sequence}(\operatorname{dbi}(m, s)) \wedge \forall n \in \mathbb{N}.\; \operatorname{seq\text{-}compact}(\operatorname{dbi}(m, s)(n)) \end{array}}}
\vnbprop{product-carrier-mem}{\vnbstmt{\begin{array}{@{}l@{}} \forall s, x. \\ \quad x \in \operatorname{product\text{-}carrier}(s) \\ \quad \Leftrightarrow\; ((x \in (\mathbb{N} \to \operatorname{big\text{-}union}(n, \mathbb{N}, \operatorname{pts}(s(n))))) \wedge \forall n \in \mathbb{N}.\; (x(n) \in \operatorname{pts}(s(n)))) \end{array}}}
\vnbprop{dual-unit-ball-in-set}{\vnbstmt{\forall m.\; (\operatorname{is\text{-}normed\text{-}vector\text{-}space}(m) \Rightarrow (\operatorname{dual\text{-}unit\text{-}ball}(m) \in \mathbf{Set}))}}
\vnbprop{dbj-value}{\vnbstmt{\begin{array}{@{}l@{}} \forall m, s, f, n. \\ \quad (\operatorname{is\text{-}normed\text{-}vector\text{-}space}(m) \wedge ((s \in (\mathbb{N} \to \operatorname{vec}(m))) \wedge ((f \in \operatorname{dual\text{-}unit\text{-}ball}(m)) \wedge (n \in \mathbb{N})))) \Rightarrow  \\ \quad \quad \operatorname{dbj}(m, s)(f)(n) = f(s(n)) \end{array}}}
\vnbprop{dbj-type}{\vnbstmt{\begin{array}{@{}l@{}} \forall m, s. \\ \quad (\operatorname{is\text{-}normed\text{-}vector\text{-}space}(m) \wedge (s \in (\mathbb{N} \to \operatorname{vec}(m)))) \Rightarrow  \\ \quad \quad \operatorname{dbj}(m, s) \in (\operatorname{dual\text{-}unit\text{-}ball}(m) \to \operatorname{product\text{-}carrier}(\operatorname{dbi}(m, s))) \end{array}}}
\vnbprop{dbj-injective}{\vnbstmt{\begin{array}{@{}l@{}} \forall m, s.; \forall x, y \in \operatorname{dual\text{-}unit\text{-}ball}(m). \\ \quad \operatorname{is\text{-}normed\text{-}vector\text{-}space}(m) \wedge \\ \quad s \in (\mathbb{N} \to \operatorname{vec}(m)) \wedge \\ \quad \forall z \in \operatorname{vec}(m), \mathit{deps}. \\ \quad \quad \operatorname{pos\text{-}rr}(\mathit{deps}) \Rightarrow  \\ \quad \quad \quad \exists k \in \mathbb{N}.\; (\operatorname{vnrm}(m)(\operatorname{vadd}(m)(z, \operatorname{vneg}(m)(s(k)))) < \mathit{deps}) \wedge \\ \quad \operatorname{dbj}(m, s)(x) = \operatorname{dbj}(m, s)(y) \Rightarrow \\ \quad \quad x = y \end{array}}}
\vnbprop{product-metric-unfold}{\vnbstmt{\begin{array}{@{}l@{}} \forall s. \\ \quad (\operatorname{product\text{-}metric}(s) \simeq \operatorname{product\text{-}metric\text{-}w}(s, (\lambda n \in \mathbb{N}.\; /(1, \operatorname{power}(2, (n + 1)))))) \end{array}}}
\vnbprop{dbi-product-facts}{\vnbstmt{\begin{array}{@{}l@{}} \forall m, s. \\ \quad (\operatorname{is\text{-}normed\text{-}vector\text{-}space}(m) \wedge (s \in (\mathbb{N} \to \operatorname{vec}(m)))) \Rightarrow  \\ \quad \quad \operatorname{is\text{-}metric\text{-}space}(\operatorname{product\text{-}metric}(\operatorname{dbi}(m, s))) \wedge \\ \quad \quad (\operatorname{pts}(\operatorname{product\text{-}metric}(\operatorname{dbi}(m, s))) \simeq \operatorname{product\text{-}carrier}(\operatorname{dbi}(m, s))) \end{array}}}
\vnbprop{dual-ball-ms-facts}{\vnbstmt{\begin{array}{@{}l@{}} \forall m, s. \\ \quad \operatorname{is\text{-}normed\text{-}vector\text{-}space}(m) \wedge \\ \quad s \in (\mathbb{N} \to \operatorname{vec}(m)) \wedge \\ \quad \forall z \in \operatorname{vec}(m), \mathit{deps}. \\ \quad \quad \operatorname{pos\text{-}rr}(\mathit{deps}) \Rightarrow  \\ \quad \quad \quad \exists k \in \mathbb{N}.\; (\operatorname{vnrm}(m)(\operatorname{vadd}(m)(z, \operatorname{vneg}(m)(s(k)))) < \mathit{deps}) \Rightarrow \\ \quad \quad \operatorname{is\text{-}metric\text{-}space}(\operatorname{dual\text{-}ball\text{-}ms}(m, s)) \wedge \\ \quad \quad (\operatorname{pts}(\operatorname{dual\text{-}ball\text{-}ms}(m, s)) \simeq \operatorname{dual\text{-}unit\text{-}ball}(m)) \wedge \\ \quad \quad (\operatorname{dual\text{-}ball\text{-}ms}(m, s) \simeq \operatorname{pullback\text{-}ms}(\operatorname{product\text{-}metric}(\operatorname{dbi}(m, s)), \operatorname{dual\text{-}unit\text{-}ball}(m), \operatorname{dbj}(m, s))) \end{array}}}
\vnbprop{dbj-compose-value}{\vnbstmt{\begin{array}{@{}l@{}} \forall m, s, a, k, n. \\ \quad \operatorname{is\text{-}normed\text{-}vector\text{-}space}(m) \wedge \\ \quad s \in (\mathbb{N} \to \operatorname{vec}(m)) \wedge \\ \quad a \in (\mathbb{N} \to \operatorname{dual\text{-}unit\text{-}ball}(m)) \wedge \\ \quad k \in \mathbb{N} \wedge \\ \quad n \in \mathbb{N} \Rightarrow \\ \quad \quad \operatorname{compose}(\operatorname{dbj}(m, s), a)(k)(n) = a(k)(s(n)) \end{array}}}
\vnbprop{dbj-coord-seq}{\vnbstmt{\begin{array}{@{}l@{}} \forall m, s, a, n. \\ \quad (\operatorname{is\text{-}normed\text{-}vector\text{-}space}(m) \wedge ((s \in (\mathbb{N} \to \operatorname{vec}(m))) \wedge ((a \in (\mathbb{N} \to \operatorname{dual\text{-}unit\text{-}ball}(m))) \wedge (n \in \mathbb{N})))) \Rightarrow  \\ \quad \quad (\lambda k \in \mathbb{N}.\; \operatorname{compose}(\operatorname{dbj}(m, s), a)(k)(n)) = (\lambda k \in \mathbb{N}.\; a(k)(s(n))) \end{array}}}
\vnbprop{dbj-coord-converges}{\vnbstmt{\begin{array}{@{}l@{}} \forall m, s, a, f, n. \\ \quad \operatorname{is\text{-}normed\text{-}vector\text{-}space}(m) \wedge \\ \quad s \in (\mathbb{N} \to \operatorname{vec}(m)) \wedge \\ \quad a \in (\mathbb{N} \to \operatorname{dual\text{-}unit\text{-}ball}(m)) \wedge \\ \quad f \in \operatorname{dual\text{-}unit\text{-}ball}(m) \wedge \\ \quad n \in \mathbb{N} \Rightarrow \\ \quad \quad \operatorname{converges\text{-}to}(\operatorname{dbi}(m, s)(n), (\lambda k \in \mathbb{N}.\; \operatorname{compose}(\operatorname{dbj}(m, s), a)(k)(n)), \\ \quad \quad \quad \operatorname{dbj}(m, s)(f)(n)) \\ \quad \quad \Leftrightarrow\; \operatorname{converges\text{-}to}(\operatorname{rr\text{-}ms}, (\lambda k \in \mathbb{N}.\; a(k)(s(n))), f(s(n))) \end{array}}}
\vnbprop{dual-ball-converges-iff}{\vnbstmt{\begin{array}{@{}l@{}} \forall m, s, a, f. \\ \quad \operatorname{is\text{-}normed\text{-}vector\text{-}space}(m) \wedge \\ \quad s \in (\mathbb{N} \to \operatorname{vec}(m)) \wedge \\ \quad \forall z \in \operatorname{vec}(m), \mathit{deps}. \\ \quad \quad \operatorname{pos\text{-}rr}(\mathit{deps}) \Rightarrow  \\ \quad \quad \quad \exists \mathit{dk} \in \mathbb{N}.\; (\operatorname{vnrm}(m)(\operatorname{vadd}(m)(z, \operatorname{vneg}(m)(s(\mathit{dk})))) < \mathit{deps}) \wedge \\ \quad a \in (\mathbb{N} \to \operatorname{dual\text{-}unit\text{-}ball}(m)) \wedge \\ \quad f \in \operatorname{dual\text{-}unit\text{-}ball}(m) \Rightarrow \\ \quad \quad \operatorname{converges\text{-}to}(\operatorname{dual\text{-}ball\text{-}ms}(m, s), a, f) \\ \quad \quad \Leftrightarrow\; \forall n \in \mathbb{N}.\; \operatorname{converges\text{-}to}(\operatorname{rr\text{-}ms}, (\lambda k \in \mathbb{N}.\; a(k)(s(n))), f(s(n))) \end{array}}}
\vnbprop{dual-ball-pointwise-limit}{\vnbstmt{\begin{array}{@{}l@{}} \forall m, s. \\ \quad \operatorname{is\text{-}normed\text{-}vector\text{-}space}(m) \wedge \\ \quad s \in (\mathbb{N} \to \operatorname{dual\text{-}unit\text{-}ball}(m)) \wedge \\ \quad \forall z \in \operatorname{vec}(m). \exists l. \\ \quad \quad \operatorname{converges\text{-}to}(\operatorname{rr\text{-}ms}, (\lambda k \in \mathbb{N}.\; s(k)(z)), l) \Rightarrow \\ \quad \quad \exists f \in \operatorname{dual\text{-}unit\text{-}ball}(m). \forall z \in \operatorname{vec}(m). \\ \quad \quad \quad \operatorname{converges\text{-}to}(\operatorname{rr\text{-}ms}, (\lambda k \in \mathbb{N}.\; s(k)(z)), f(z)) \end{array}}}
\vnbprop{dual-ball-limits-everywhere}{\vnbstmt{\begin{array}{@{}l@{}} \forall m, s, a.; \forall z \in \operatorname{vec}(m). \\ \quad \operatorname{is\text{-}normed\text{-}vector\text{-}space}(m) \wedge \\ \quad s \in (\mathbb{N} \to \operatorname{vec}(m)) \wedge \\ \quad \forall \mathit{dz} \in \operatorname{vec}(m), \mathit{deps}. \\ \quad \quad \operatorname{pos\text{-}rr}(\mathit{deps}) \Rightarrow  \\ \quad \quad \quad \exists \mathit{dk} \in \mathbb{N}.\; (\operatorname{vnrm}(m)(\operatorname{vadd}(m)(\mathit{dz}, \operatorname{vneg}(m)(s(\mathit{dk})))) < \mathit{deps}) \wedge \\ \quad a \in (\mathbb{N} \to \operatorname{dual\text{-}unit\text{-}ball}(m)) \wedge \\ \quad \forall n \in \mathbb{N}. \exists l. \\ \quad \quad \operatorname{converges\text{-}to}(\operatorname{rr\text{-}ms}, (\lambda k \in \mathbb{N}.\; a(k)(s(n))), l) \Rightarrow \\ \quad \quad \exists l.\; \operatorname{converges\text{-}to}(\operatorname{rr\text{-}ms}, (\lambda k \in \mathbb{N}.\; a(k)(z)), l) \end{array}}}
\vnbprop{dbj-diagonal-value}{\vnbstmt{\begin{array}{@{}l@{}} \forall m, s, a, b, k, n. \\ \quad \operatorname{is\text{-}normed\text{-}vector\text{-}space}(m) \wedge \\ \quad s \in (\mathbb{N} \to \operatorname{vec}(m)) \wedge \\ \quad a \in (\mathbb{N} \to \operatorname{dual\text{-}unit\text{-}ball}(m)) \wedge \\ \quad b \in (\mathbb{N} \to \mathbb{N}) \wedge \\ \quad k \in \mathbb{N} \wedge \\ \quad n \in \mathbb{N} \Rightarrow \\ \quad \quad \operatorname{subseq}(\operatorname{compose}(\operatorname{dbj}(m, s), a), b)(k)(n) = \operatorname{subseq}(a, b)(k)(s(n)) \end{array}}}
\vnbprop{dual-ball-seq-coord}{\vnbstmt{\begin{array}{@{}l@{}} \forall m, s, a, k, n. \\ \quad \operatorname{is\text{-}normed\text{-}vector\text{-}space}(m) \wedge \\ \quad s \in (\mathbb{N} \to \operatorname{vec}(m)) \wedge \\ \quad a \in (\mathbb{N} \to \operatorname{dual\text{-}unit\text{-}ball}(m)) \wedge \\ \quad k \in \mathbb{N} \wedge \\ \quad n \in \mathbb{N} \Rightarrow \\ \quad \quad (a(k)(s(n)) \in \operatorname{ccint}(-\operatorname{vnrm}(m)(s(n)), \operatorname{vnrm}(m)(s(n)))) \wedge (a(k)(s(n)) \in \mathbb{R}) \end{array}}}
\vnbprop{dual-ball-subseq-converges}{\vnbstmt{\begin{array}{@{}l@{}} \forall m, s, a. \\ \quad \operatorname{is\text{-}normed\text{-}vector\text{-}space}(m) \wedge \\ \quad s \in (\mathbb{N} \to \operatorname{vec}(m)) \wedge \\ \quad \forall z \in \operatorname{vec}(m), \mathit{deps}. \\ \quad \quad \operatorname{pos\text{-}rr}(\mathit{deps}) \Rightarrow  \\ \quad \quad \quad \exists k \in \mathbb{N}.\; (\operatorname{vnrm}(m)(\operatorname{vadd}(m)(z, \operatorname{vneg}(m)(s(k)))) < \mathit{deps}) \wedge \\ \quad a \in (\mathbb{N} \to \operatorname{dual\text{-}unit\text{-}ball}(m)) \Rightarrow \\ \quad \quad \exists i. \\ \quad \quad \quad \operatorname{strictly\text{-}mono\text{-}nn}(i) \wedge \\ \quad \quad \quad \exists l \in \operatorname{dual\text{-}unit\text{-}ball}(m). \\ \quad \quad \quad \quad \operatorname{converges\text{-}to}(\operatorname{dual\text{-}ball\text{-}ms}(m, s), \operatorname{subseq}(a, i), l) \end{array}}}
\vnbprop{dual-ball-seq-compact}{\vnbstmt{\begin{array}{@{}l@{}} \forall m, s. \\ \quad \operatorname{is\text{-}normed\text{-}vector\text{-}space}(m) \wedge \\ \quad s \in (\mathbb{N} \to \operatorname{vec}(m)) \wedge \\ \quad \forall z \in \operatorname{vec}(m), \mathit{deps}. \\ \quad \quad \operatorname{pos\text{-}rr}(\mathit{deps}) \Rightarrow  \\ \quad \quad \quad \exists k \in \mathbb{N}.\; (\operatorname{vnrm}(m)(\operatorname{vadd}(m)(z, \operatorname{vneg}(m)(s(k)))) < \mathit{deps}) \Rightarrow \\ \quad \quad \operatorname{seq\text{-}compact}(\operatorname{dual\text{-}ball\text{-}ms}(m, s)) \end{array}}}
\vnbprop{dual-ball-compact}{\vnbstmt{\begin{array}{@{}l@{}} \forall m, s. \\ \quad \operatorname{is\text{-}normed\text{-}vector\text{-}space}(m) \wedge \\ \quad s \in (\mathbb{N} \to \operatorname{vec}(m)) \wedge \\ \quad \forall z \in \operatorname{vec}(m), \mathit{deps}. \\ \quad \quad \operatorname{pos\text{-}rr}(\mathit{deps}) \Rightarrow  \\ \quad \quad \quad \exists k \in \mathbb{N}.\; (\operatorname{vnrm}(m)(\operatorname{vadd}(m)(z, \operatorname{vneg}(m)(s(k)))) < \mathit{deps}) \Rightarrow \\ \quad \quad \operatorname{is\text{-}compact}(\operatorname{dual\text{-}ball\text{-}ms}(m, s)) \end{array}}}
\vnbprop{dual-ball-converges-pointwise}{\vnbstmt{\begin{array}{@{}l@{}} \forall m, s, a, f, z. \\ \quad \operatorname{is\text{-}normed\text{-}vector\text{-}space}(m) \wedge \\ \quad s \in (\mathbb{N} \to \operatorname{vec}(m)) \wedge \\ \quad \forall \mathit{dz} \in \operatorname{vec}(m), \mathit{deps}. \\ \quad \quad \operatorname{pos\text{-}rr}(\mathit{deps}) \Rightarrow  \\ \quad \quad \quad \exists \mathit{dk} \in \mathbb{N}.\; (\operatorname{vnrm}(m)(\operatorname{vadd}(m)(\mathit{dz}, \operatorname{vneg}(m)(s(\mathit{dk})))) < \mathit{deps}) \wedge \\ \quad a \in (\mathbb{N} \to \operatorname{dual\text{-}unit\text{-}ball}(m)) \wedge \\ \quad f \in \operatorname{dual\text{-}unit\text{-}ball}(m) \wedge \\ \quad \operatorname{converges\text{-}to}(\operatorname{dual\text{-}ball\text{-}ms}(m, s), a, f) \wedge \\ \quad z \in \operatorname{vec}(m) \Rightarrow \\ \quad \quad \operatorname{converges\text{-}to}(\operatorname{rr\text{-}ms}, (\lambda k \in \mathbb{N}.\; a(k)(z)), f(z)) \end{array}}}
\vnbprop{sqn-unfold-b6s}{\vnbstmt{\forall a.\; (a^{\mathbb{N}} \simeq (\mathbb{N} \to a))}}
\vnbprop{dual-ball-evaluation-continuous}{\vnbstmt{\begin{array}{@{}l@{}} \forall m, s, z. \\ \quad \operatorname{is\text{-}normed\text{-}vector\text{-}space}(m) \wedge \\ \quad s \in (\mathbb{N} \to \operatorname{vec}(m)) \wedge \\ \quad \forall \mathit{dz} \in \operatorname{vec}(m), \mathit{deps}. \\ \quad \quad \operatorname{pos\text{-}rr}(\mathit{deps}) \Rightarrow  \\ \quad \quad \quad \exists k \in \mathbb{N}.\; (\operatorname{vnrm}(m)(\operatorname{vadd}(m)(\mathit{dz}, \operatorname{vneg}(m)(s(k)))) < \mathit{deps}) \wedge \\ \quad z \in \operatorname{vec}(m) \Rightarrow \\ \quad \quad \operatorname{is\text{-}continuous}(\operatorname{dual\text{-}ball\text{-}ms}(m, s), \operatorname{rr\text{-}ms}, (\lambda e \in \operatorname{dual\text{-}unit\text{-}ball}(m).\; e(z))) \end{array}}}
\vnbprop{dual-ball-identity-continuous}{\vnbstmt{\begin{array}{@{}l@{}} \forall m, s. \\ \quad \operatorname{is\text{-}normed\text{-}vector\text{-}space}(m) \wedge \\ \quad s \in (\mathbb{N} \to \operatorname{vec}(m)) \wedge \\ \quad \forall z \in \operatorname{vec}(m), \mathit{deps}. \\ \quad \quad \operatorname{pos\text{-}rr}(\mathit{deps}) \Rightarrow  \\ \quad \quad \quad \exists k \in \mathbb{N}.\; (\operatorname{vnrm}(m)(\operatorname{vadd}(m)(z, \operatorname{vneg}(m)(s(k)))) < \mathit{deps}) \Rightarrow \\ \quad \quad \operatorname{is\text{-}hom\text{-}top\text{-}space}(\operatorname{metric\text{-}top}(\operatorname{dual\text{-}ball\text{-}ms}(m, s)), \\ \quad \quad \quad \operatorname{subspace\text{-}top}(\operatorname{tproduct\text{-}top}(\operatorname{vec}(m), (\lambda n \in \operatorname{vec}(m).\; \operatorname{metric\text{-}top}(\operatorname{rr\text{-}ms}))), \operatorname{dual\text{-}unit\text{-}ball}(m)), \\ \quad \quad \quad (\lambda i \in \operatorname{dual\text{-}unit\text{-}ball}(m).\; i)) \end{array}}}
\vnbprop{line-product-constraint}{\vnbstmt{\begin{array}{@{}l@{}} \forall m, z, c, r. \\ \quad (\operatorname{is\text{-}normed\text{-}vector\text{-}space}(m) \wedge ((z \in \operatorname{vec}(m)) \wedge ((c \in \mathbb{R}) \wedge ((r \in \mathbb{R}) \wedge (0 < r))))) \Rightarrow  \\ \quad \quad {[}z, \operatorname{ball}(\operatorname{rr\text{-}ms}, c, r)] \in \operatorname{tprod\text{-}constr}(\operatorname{vec}(m), (\lambda n \in \operatorname{vec}(m).\; \operatorname{metric\text{-}top}(\operatorname{rr\text{-}ms}))) \end{array}}}
\vnbprop{dual-ball-box-open}{\vnbstmt{\begin{array}{@{}l@{}} \forall m, s, o, b, l. \\ \quad \operatorname{is\text{-}normed\text{-}vector\text{-}space}(m) \wedge \\ \quad s \in (\mathbb{N} \to \operatorname{vec}(m)) \wedge \\ \quad o \in \operatorname{dual\text{-}unit\text{-}ball}(m) \wedge \\ \quad b \in \mathbb{N} \wedge \\ \quad l \in \mathbb{R} \wedge \\ \quad 0 < l \Rightarrow \\ \quad \quad \operatorname{tprod\text{-}box}(\operatorname{vec}(m), \\ \quad \quad \quad (\lambda n \in \operatorname{vec}(m).\; \operatorname{metric\text{-}top}(\operatorname{rr\text{-}ms})), \\ \quad \quad \quad \operatorname{image}((\lambda j \in \mathrm{OS}(b).\; {[}s(j), \operatorname{ball}(\operatorname{rr\text{-}ms}, o(s(j)), l)]), \mathrm{OS}(b))) \\ \quad \quad \in\; \operatorname{opens}(\operatorname{tproduct\text{-}top}(\operatorname{vec}(m), (\lambda n \in \operatorname{vec}(m).\; \operatorname{metric\text{-}top}(\operatorname{rr\text{-}ms})))) \end{array}}}
\vnbprop{dual-ball-box-centre}{\vnbstmt{\begin{array}{@{}l@{}} \forall m, s, o, b, l. \\ \quad \operatorname{is\text{-}normed\text{-}vector\text{-}space}(m) \wedge \\ \quad s \in (\mathbb{N} \to \operatorname{vec}(m)) \wedge \\ \quad o \in \operatorname{dual\text{-}unit\text{-}ball}(m) \wedge \\ \quad b \in \mathbb{N} \wedge \\ \quad l \in \mathbb{R} \wedge \\ \quad 0 < l \Rightarrow \\ \quad \quad o \\ \quad \quad \in\; \operatorname{tprod\text{-}box}(\operatorname{vec}(m), \\ \quad \quad \quad (\lambda n \in \operatorname{vec}(m).\; \operatorname{metric\text{-}top}(\operatorname{rr\text{-}ms})), \\ \quad \quad \quad \operatorname{image}((\lambda j \in \mathrm{OS}(b).\; {[}s(j), \operatorname{ball}(\operatorname{rr\text{-}ms}, o(s(j)), l)]), \mathrm{OS}(b))) \end{array}}}
\vnbprop{dual-ball-box-coord}{\vnbstmt{\begin{array}{@{}l@{}} \forall m, s, o, b, l, g, j. \\ \quad \operatorname{is\text{-}normed\text{-}vector\text{-}space}(m) \wedge \\ \quad s \in (\mathbb{N} \to \operatorname{vec}(m)) \wedge \\ \quad o \in \operatorname{dual\text{-}unit\text{-}ball}(m) \wedge \\ \quad b \in \mathbb{N} \wedge \\ \quad l \in \mathbb{R} \wedge \\ \quad 0 < l \wedge \\ \quad g \\ \quad \in\; \operatorname{tprod\text{-}box}(\operatorname{vec}(m), \\ \quad \quad (\lambda n \in \operatorname{vec}(m).\; \operatorname{metric\text{-}top}(\operatorname{rr\text{-}ms})), \\ \quad \quad \operatorname{image}((\lambda \mathit{b6j} \in \mathrm{OS}(b).\; {[}s(\mathit{b6j}), \operatorname{ball}(\operatorname{rr\text{-}ms}, o(s(\mathit{b6j})), l)]), \mathrm{OS}(b))) \wedge \\ \quad g \in \operatorname{dual\text{-}unit\text{-}ball}(m) \wedge \\ \quad j \in \mathrm{OS}(b) \Rightarrow \\ \quad \quad \lvert (o(s(j)) - g(s(j))) \rvert < l \end{array}}}
\vnbprop{dual-ball-open-has-box}{\vnbstmt{\begin{array}{@{}l@{}} \forall m, s, u, o. \\ \quad \operatorname{is\text{-}normed\text{-}vector\text{-}space}(m) \wedge \\ \quad s \in (\mathbb{N} \to \operatorname{vec}(m)) \wedge \\ \quad \forall z \in \operatorname{vec}(m), \mathit{deps}. \\ \quad \quad \operatorname{pos\text{-}rr}(\mathit{deps}) \Rightarrow  \\ \quad \quad \quad \exists k \in \mathbb{N}.\; (\operatorname{vnrm}(m)(\operatorname{vadd}(m)(z, \operatorname{vneg}(m)(s(k)))) < \mathit{deps}) \wedge \\ \quad \operatorname{is\text{-}open}(\operatorname{dual\text{-}ball\text{-}ms}(m, s), u) \wedge \\ \quad o \in u \Rightarrow \\ \quad \quad \exists b \in \mathbb{N}. \\ \quad \quad \quad (\operatorname{tprod\text{-}box}(\operatorname{vec}(m), (\lambda n \in \operatorname{vec}(m).\; \operatorname{metric\text{-}top}(\operatorname{rr\text{-}ms})), \operatorname{image}((\lambda j \in \mathrm{OS}(b).\; {[}s(j), \operatorname{ball}(\operatorname{rr\text{-}ms}, o(s(j)), \operatorname{recip}((b + 1)))]), \mathrm{OS}(b))) \cap \operatorname{dual\text{-}unit\text{-}ball}(m)) \\ \quad \quad \quad \subseteq\; u \end{array}}}
\vnbprop{dual-ball-topology-agrees}{\vnbstmt{\begin{array}{@{}l@{}} \forall m, s, u. \\ \quad \operatorname{is\text{-}normed\text{-}vector\text{-}space}(m) \wedge \\ \quad s \in (\mathbb{N} \to \operatorname{vec}(m)) \wedge \\ \quad \forall z \in \operatorname{vec}(m), \mathit{deps}. \\ \quad \quad \operatorname{pos\text{-}rr}(\mathit{deps}) \Rightarrow  \\ \quad \quad \quad \exists k \in \mathbb{N}.\; (\operatorname{vnrm}(m)(\operatorname{vadd}(m)(z, \operatorname{vneg}(m)(s(k)))) < \mathit{deps}) \Rightarrow \\ \quad \quad u \in \operatorname{opens}(\operatorname{subspace\text{-}top}(\operatorname{dual\text{-}weak\text{-}star\text{-}top}(m), \operatorname{dual\text{-}unit\text{-}ball}(m))) \\ \quad \quad \Leftrightarrow\; \operatorname{is\text{-}open}(\operatorname{dual\text{-}ball\text{-}ms}(m, s), u) \end{array}}}
\vnbprop{banach-alaoglu-separable}{\vnbstmt{\begin{array}{@{}l@{}} \forall m, s. \\ \quad \operatorname{is\text{-}normed\text{-}vector\text{-}space}(m) \wedge \\ \quad s \in (\mathbb{N} \to \operatorname{vec}(m)) \wedge \\ \quad \forall z \in \operatorname{vec}(m), \mathit{deps}. \\ \quad \quad \operatorname{pos\text{-}rr}(\mathit{deps}) \Rightarrow  \\ \quad \quad \quad \exists k \in \mathbb{N}.\; (\operatorname{vnrm}(m)(\operatorname{vadd}(m)(z, \operatorname{vneg}(m)(s(k)))) < \mathit{deps}) \Rightarrow \\ \quad \quad \operatorname{is\text{-}metric\text{-}space}(\operatorname{dual\text{-}ball\text{-}ms}(m, s)) \wedge \\ \quad \quad (\operatorname{pts}(\operatorname{dual\text{-}ball\text{-}ms}(m, s)) \simeq \operatorname{dual\text{-}unit\text{-}ball}(m)) \wedge \\ \quad \quad \forall u. \\ \quad \quad \quad u \in \operatorname{opens}(\operatorname{subspace\text{-}top}(\operatorname{dual\text{-}weak\text{-}star\text{-}top}(m), \operatorname{dual\text{-}unit\text{-}ball}(m))) \\ \quad \quad \quad \Leftrightarrow\; \operatorname{is\text{-}open}(\operatorname{dual\text{-}ball\text{-}ms}(m, s), u) \wedge \\ \quad \quad \operatorname{is\text{-}compact}(\operatorname{dual\text{-}ball\text{-}ms}(m, s)) \end{array}}}
\vnbprop{banach-alaoglu-separable-metrizable}{\vnbstmt{\begin{array}{@{}l@{}} \forall m. \\ \quad (\operatorname{is\text{-}normed\text{-}vector\text{-}space}(m) \wedge \operatorname{is\text{-}separable}(\operatorname{nvs\text{-}metric\text{-}space}(m))) \Rightarrow  \\ \quad \quad \operatorname{is\text{-}metrizable\text{-}top\text{-}space}(\operatorname{subspace\text{-}top}(\operatorname{dual\text{-}weak\text{-}star\text{-}top}(m), \operatorname{dual\text{-}unit\text{-}ball}(m))) \wedge \\ \quad \quad \operatorname{is\text{-}compact\text{-}t}(\operatorname{subspace\text{-}top}(\operatorname{dual\text{-}weak\text{-}star\text{-}top}(m), \operatorname{dual\text{-}unit\text{-}ball}(m))) \end{array}}}

\vnbbookfile{structure-library/cvx-vocab}{structure-library/cvx-vocab}
\noindent{\small cvx-vocab.scm -- convex combinations and the convex hull in a real normed vector space (compact-groups notes, Lemma 1.12).  DEFINITIONS ONLY.  Agent ISOM, 2026-10-26 (the user's to confirm).}\par
\par\medskip\noindent\textbf{Definition} (\texttt{conv-level}).\;For all $\operatorname{cvm\_}$ and $\operatorname{cve\_}$,  $(\operatorname{conv\text{-}level}(\operatorname{cvm\_}, \operatorname{cve\_}, 0) \simeq \operatorname{cve\_})$\par
Let $\operatorname{cvm\_}$ and $\operatorname{cve\_}$ be arbitrary. Let $\operatorname{cvn\_} \in \mathbb{N}$. Then
\begin{equation*}
(\operatorname{conv\text{-}level}(\operatorname{cvm\_}, \operatorname{cve\_}, \operatorname{succ}(\operatorname{cvn\_})) \simeq \{\operatorname{cvz\_} \in \operatorname{vec}(\operatorname{cvm\_}) : \exists \operatorname{cvx\_} \in \operatorname{conv\text{-}level}(\operatorname{cvm\_}, \operatorname{cve\_}, \operatorname{cvn\_}), \operatorname{cvy\_} \in \operatorname{cve\_}, \operatorname{cvt\_} \in \mathbb{R}.\; ((0 \leq \operatorname{cvt\_}) \wedge ((\operatorname{cvt\_} \leq 1) \wedge (\operatorname{cvz\_} = \operatorname{vadd}(\operatorname{cvm\_})(\operatorname{act}(\operatorname{cvm\_})(-(1, \operatorname{cvt\_}), \operatorname{cvx\_}), \operatorname{act}(\operatorname{cvm\_})(\operatorname{cvt\_}, \operatorname{cvy\_})))))\})
\end{equation*}

\par\medskip\noindent\textbf{Definition} (\texttt{conv-hull}).\;$\operatorname{conv\text{-}hull}(\operatorname{cvm\_}, \operatorname{cve\_}) :=$ $\operatorname{big\text{-}union}(\operatorname{cvn\_}, \mathbb{N}, \operatorname{conv\text{-}level}(\operatorname{cvm\_}, \operatorname{cve\_}, \operatorname{cvn\_}))$
\par\medskip\noindent\textbf{Definition} (\texttt{conv-join}).\;$\operatorname{conv\text{-}join}(\operatorname{cjm\_}, \operatorname{cjc\_}, \operatorname{cja\_}) :=$
\begin{equation*}
\{\operatorname{cjz\_} \in \operatorname{vec}(\operatorname{cjm\_}) : \exists \operatorname{cjy\_} \in \operatorname{cjc\_}, \operatorname{cjt\_} \in \mathbb{R}.\; ((0 \leq \operatorname{cjt\_}) \wedge ((\operatorname{cjt\_} \leq 1) \wedge (\operatorname{cjz\_} = \operatorname{vadd}(\operatorname{cjm\_})(\operatorname{act}(\operatorname{cjm\_})((1 - \operatorname{cjt\_}), \operatorname{cjy\_}), \operatorname{act}(\operatorname{cjm\_})(\operatorname{cjt\_}, \operatorname{cja\_})))))\}
\end{equation*}

\par\medskip\noindent\textbf{Definition} (\texttt{is-totally-bounded-in}).\;For all $\operatorname{tbm\_}$ and $\operatorname{tba\_}$, 
\begin{equation*}
\begin{array}{@{}l@{}} \operatorname{is\text{-}totally\text{-}bounded\text{-}in}(\operatorname{tbm\_}, \operatorname{tba\_}) \\ \Leftrightarrow\; \forall \operatorname{tbr\_}. \\ \quad \operatorname{pos\text{-}rr}(\operatorname{tbr\_}) \Rightarrow  \\ \quad \quad \exists \operatorname{tbf\_} \in \mathbf{Set}. \\ \quad \quad \quad |\operatorname{tbf\_}| \in \mathbb{N} \wedge \\ \quad \quad \quad \operatorname{tbf\_} \subseteq \operatorname{pts}(\operatorname{tbm\_}) \wedge \\ \quad \quad \quad \forall \operatorname{tbp\_} \in \operatorname{tba\_}.\; \exists \operatorname{tbc\_} \in \operatorname{tbf\_}.\; (\operatorname{dist}(\operatorname{tbm\_})(\operatorname{tbc\_}, \operatorname{tbp\_}) < \operatorname{tbr\_}) \end{array}
\end{equation*}


\vnbbookfile{structure-library/hc-vocab}{structure-library/hc-vocab}
\noindent{\small hc-vocab.scm -- wave of 2026-11-09/10 (agent HAAR-C), wired 2026-11-10 from scratchpad/haarc-1109/hc-vocab.scm. Library: analysis.}\par
\par\medskip\noindent\textbf{Definition} (\texttt{is-affine-isometry-group}).\;For all $\operatorname{aim\_}$, $\operatorname{aia\_}$ and $\operatorname{aig\_}$, 
\begin{equation*}
\begin{array}{@{}l@{}} \operatorname{is\text{-}affine\text{-}isometry\text{-}group}(\operatorname{aim\_}, \operatorname{aia\_}, \operatorname{aig\_}) \\ \Leftrightarrow\; \operatorname{aig\_} \subseteq (\operatorname{aia\_} \to \operatorname{aia\_}) \wedge \\ \forall \operatorname{aif\_} \in \operatorname{aig\_}, \operatorname{aiu\_}, \operatorname{aiv\_} \in \operatorname{aia\_}, \operatorname{ait\_} \in \mathbb{R}. \\ \quad (0 \leq \operatorname{ait\_}) \Rightarrow  \\ \quad (\operatorname{ait\_} \leq 1) \Rightarrow  \\ \quad \quad \operatorname{aif\_}(\operatorname{vadd}(\operatorname{aim\_})(\operatorname{act}(\operatorname{aim\_})(-(1, \operatorname{ait\_}), \operatorname{aiu\_}), \operatorname{act}(\operatorname{aim\_})(\operatorname{ait\_}, \operatorname{aiv\_}))) \\ \quad \quad =\; \operatorname{vadd}(\operatorname{aim\_})(\operatorname{act}(\operatorname{aim\_})(-(1, \operatorname{ait\_}), \operatorname{aif\_}(\operatorname{aiu\_})), \operatorname{act}(\operatorname{aim\_})(\operatorname{ait\_}, \operatorname{aif\_}(\operatorname{aiv\_}))) \wedge \\ \forall \operatorname{aif\_} \in \operatorname{aig\_}, \operatorname{aiu\_}, \operatorname{aiv\_} \in \operatorname{aia\_}. \\ \quad \operatorname{dist}(\operatorname{nvs\text{-}metric\text{-}space}(\operatorname{aim\_}))(\operatorname{aif\_}(\operatorname{aiu\_}), \operatorname{aif\_}(\operatorname{aiv\_})) = \operatorname{dist}(\operatorname{nvs\text{-}metric\text{-}space}(\operatorname{aim\_}))(\operatorname{aiu\_}, \operatorname{aiv\_}) \wedge \\ \forall \operatorname{aif\_} \in \operatorname{aig\_}. \exists \operatorname{aih\_} \in \operatorname{aig\_}. \forall \operatorname{aix\_} \in \operatorname{aia\_}. \\ \quad \operatorname{aif\_}(\operatorname{aih\_}(\operatorname{aix\_})) = \operatorname{aix\_} \end{array}
\end{equation*}

\par\medskip\noindent\textbf{Definition} (\texttt{is-invariant-ccs}).\;For all $\operatorname{icm\_}$, $\operatorname{ica\_}$, $\operatorname{icg\_}$ and $\operatorname{icb\_}$, 
\begin{equation*}
\begin{array}{@{}l@{}} \operatorname{is\text{-}invariant\text{-}ccs}(\operatorname{icm\_}, \operatorname{ica\_}, \operatorname{icg\_}, \operatorname{icb\_}) \\ \Leftrightarrow\; \operatorname{icb\_} \subseteq \operatorname{ica\_} \wedge \\ \exists \operatorname{icz\_}.\; (\operatorname{icz\_} \in \operatorname{icb\_}) \wedge \\ \operatorname{is\text{-}convex\text{-}subset}(\operatorname{icm\_}, \operatorname{icb\_}) \wedge \\ \operatorname{is\text{-}closed}(\operatorname{subspace\text{-}ms}(\operatorname{nvs\text{-}metric\text{-}space}(\operatorname{icm\_}), \operatorname{ica\_}), \operatorname{icb\_}) \wedge \\ \forall \operatorname{icf\_} \in \operatorname{icg\_}, \operatorname{icx\_} \in \operatorname{icb\_}.\; (\operatorname{icf\_}(\operatorname{icx\_}) \in \operatorname{icb\_}) \end{array}
\end{equation*}

\par\medskip\noindent\textbf{Definition} (\texttt{is-haar-integral}).\;For all $\operatorname{hig\_}$ and $\operatorname{hiu\_}$, 
\begin{equation*}
\begin{array}{@{}l@{}} \operatorname{is\text{-}haar\text{-}integral}(\operatorname{hig\_}, \operatorname{hiu\_}) \\ \Leftrightarrow\; \operatorname{is\text{-}metric\text{-}group}(\operatorname{hig\_}) \wedge \\ \operatorname{is\text{-}compact}(\operatorname{metric\text{-}group\text{-}ms}(\operatorname{hig\_})) \wedge \\ \operatorname{is\text{-}positive\text{-}functional}(\operatorname{metric\text{-}group\text{-}ms}(\operatorname{hig\_}), \operatorname{hiu\_}) \wedge \\ \operatorname{hiu\_}((\lambda \operatorname{hix\_} \in \operatorname{pts}(\operatorname{metric\text{-}group\text{-}ms}(\operatorname{hig\_})).\; 1)) = 1 \wedge \\ \forall \operatorname{hia\_} \in \operatorname{pts}(\operatorname{metric\text{-}group\text{-}ms}(\operatorname{hig\_})), \operatorname{hif\_} \in \operatorname{cont\text{-}fns}(\operatorname{metric\text{-}group\text{-}ms}(\operatorname{hig\_})). \\ \quad \operatorname{hiu\_}((\lambda \operatorname{hix\_} \in \operatorname{pts}(\operatorname{metric\text{-}group\text{-}ms}(\operatorname{hig\_})).\; \operatorname{hif\_}(\operatorname{opr}(\operatorname{hig\_})(\operatorname{hia\_}, \operatorname{hix\_})))) = \operatorname{hiu\_}(\operatorname{hif\_}) \end{array}
\end{equation*}


\vnbbookfile{structure-library/hc3-vocab}{structure-library/hc3-vocab}
\noindent{\small hc3-vocab.scm -- wave of 2026-11-09/10 (agent HAAR-C3), wired 2026-11-10 from scratchpad/haarc3-1109/hc3-vocab.scm. Library: analysis.}\par
\par\medskip\noindent\textbf{Definition} (\texttt{opp-group}).\;${\operatorname{ogg\_}}^{\mathrm{op}} :=$
\begin{equation*}
{[}\operatorname{carr}(\operatorname{ogg\_}), (\lambda {[}\operatorname{ogp\_}, \operatorname{ogq\_}] \in (\operatorname{carr}(\operatorname{ogg\_}) \times \operatorname{carr}(\operatorname{ogg\_})).\; \operatorname{opr}(\operatorname{ogg\_})(\operatorname{ogq\_}, \operatorname{ogp\_})), \operatorname{iden}(\operatorname{ogg\_}), \operatorname{inv}(\operatorname{ogg\_})]
\end{equation*}

\par\medskip\noindent\textbf{Definition} (\texttt{opp-metric-group}).\;${\operatorname{omg\_}}^{\mathrm{op}} :=$
\begin{equation*}
\operatorname{make\text{-}metric\text{-}group}(\operatorname{metric\text{-}group\text{-}ms}(\operatorname{omg\_}), {\operatorname{metric\text{-}group\text{-}group}(\operatorname{omg\_})}^{\mathrm{op}})
\end{equation*}

\par\medskip\noindent\textbf{Definition} (\texttt{haar-averages}).\;$\operatorname{haar\text{-}averages}(\operatorname{hag\_}) :=$
\begin{equation*}
\{\operatorname{haq\_} \in (\operatorname{cont\text{-}fns}(\operatorname{metric\text{-}group\text{-}ms}(\operatorname{hag\_})) \to \operatorname{cont\text{-}fns}(\operatorname{metric\text{-}group\text{-}ms}(\operatorname{hag\_}))) : (\forall \operatorname{hau\_}, \operatorname{hav\_} \in \operatorname{cont\text{-}fns}(\operatorname{metric\text{-}group\text{-}ms}(\operatorname{hag\_})).\; (\operatorname{haq\_}(\operatorname{vadd}(C({\operatorname{metric\text{-}group\text{-}ms}(\operatorname{hag\_})}))(\operatorname{hau\_}, \operatorname{hav\_})) = \operatorname{vadd}(C({\operatorname{metric\text{-}group\text{-}ms}(\operatorname{hag\_})}))(\operatorname{haq\_}(\operatorname{hau\_}), \operatorname{haq\_}(\operatorname{hav\_}))) \wedge (\forall \operatorname{har\_} \in \mathbb{R}, \operatorname{hau\_} \in \operatorname{cont\text{-}fns}(\operatorname{metric\text{-}group\text{-}ms}(\operatorname{hag\_})).\; (\operatorname{haq\_}(\operatorname{act}(C({\operatorname{metric\text{-}group\text{-}ms}(\operatorname{hag\_})}))(\operatorname{har\_}, \operatorname{hau\_})) = \operatorname{act}(C({\operatorname{metric\text{-}group\text{-}ms}(\operatorname{hag\_})}))(\operatorname{har\_}, \operatorname{haq\_}(\operatorname{hau\_}))) \wedge (\forall \operatorname{hau\_}, \operatorname{hav\_} \in \operatorname{cont\text{-}fns}(\operatorname{metric\text{-}group\text{-}ms}(\operatorname{hag\_})).\; (\operatorname{dist}(\operatorname{nvs\text{-}metric\text{-}space}(C({\operatorname{metric\text{-}group\text{-}ms}(\operatorname{hag\_})})))(\operatorname{haq\_}(\operatorname{hau\_}), \operatorname{haq\_}(\operatorname{hav\_})) \leq \operatorname{dist}(\operatorname{nvs\text{-}metric\text{-}space}(C({\operatorname{metric\text{-}group\text{-}ms}(\operatorname{hag\_})})))(\operatorname{hau\_}, \operatorname{hav\_})) \wedge (\forall \operatorname{hac\_} \in \mathbb{R}.\; (\operatorname{haq\_}((\lambda \operatorname{hix\_} \in \operatorname{pts}(\operatorname{metric\text{-}group\text{-}ms}(\operatorname{hag\_})).\; \operatorname{hac\_})) = (\lambda \operatorname{hix\_} \in \operatorname{pts}(\operatorname{metric\text{-}group\text{-}ms}(\operatorname{hag\_})).\; \operatorname{hac\_})) \wedge (\forall \operatorname{hau\_} \in \operatorname{cont\text{-}fns}(\operatorname{metric\text{-}group\text{-}ms}(\operatorname{hag\_})).\; (\operatorname{haq\_}(\operatorname{hau\_}) \in \operatorname{conv\text{-}hull}(C({\operatorname{metric\text{-}group\text{-}ms}(\operatorname{hag\_})}), \operatorname{image}((\lambda \operatorname{hqa\_} \in \operatorname{pts}(\operatorname{metric\text{-}group\text{-}ms}(\operatorname{hag\_})).\; (\lambda \operatorname{hix\_} \in \operatorname{pts}(\operatorname{metric\text{-}group\text{-}ms}(\operatorname{hag\_})).\; \operatorname{hau\_}(\operatorname{opr}(\operatorname{hag\_})(\operatorname{hqa\_}, \operatorname{hix\_})))), \operatorname{pts}(\operatorname{metric\text{-}group\text{-}ms}(\operatorname{hag\_}))))) \wedge \forall \operatorname{hab\_} \in \operatorname{pts}(\operatorname{metric\text{-}group\text{-}ms}(\operatorname{hag\_})), \operatorname{hau\_} \in \operatorname{cont\text{-}fns}(\operatorname{metric\text{-}group\text{-}ms}(\operatorname{hag\_})).\; (\operatorname{haq\_}((\lambda \operatorname{hix\_} \in \operatorname{pts}(\operatorname{metric\text{-}group\text{-}ms}(\operatorname{hag\_})).\; \operatorname{hau\_}(\operatorname{opr}(\operatorname{hag\_})(\operatorname{hix\_}, \operatorname{hab\_})))) = (\lambda \operatorname{hix\_} \in \operatorname{pts}(\operatorname{metric\text{-}group\text{-}ms}(\operatorname{hag\_})).\; \operatorname{haq\_}(\operatorname{hau\_})(\operatorname{opr}(\operatorname{hag\_})(\operatorname{hix\_}, \operatorname{hab\_})))))))))\}
\end{equation*}

\par\medskip\noindent\textbf{Definition} (\texttt{haar-mean}).\;$\int_{\operatorname{hmg\_}} {\operatorname{hmf\_}} :=$
\begin{equation*}
(\iota\, \operatorname{hmc\_}.\;((\operatorname{hmc\_} \in \mathbb{R}) \wedge ((\lambda \operatorname{hix\_} \in \operatorname{pts}(\operatorname{metric\text{-}group\text{-}ms}(\operatorname{hmg\_})).\; \operatorname{hmc\_}) \in \operatorname{closure}(\operatorname{nvs\text{-}metric\text{-}space}(C({\operatorname{metric\text{-}group\text{-}ms}(\operatorname{hmg\_})})), \operatorname{conv\text{-}hull}(C({\operatorname{metric\text{-}group\text{-}ms}(\operatorname{hmg\_})}), \operatorname{image}((\lambda \operatorname{hqa\_} \in \operatorname{pts}(\operatorname{metric\text{-}group\text{-}ms}(\operatorname{hmg\_})).\; (\lambda \operatorname{hix\_} \in \operatorname{pts}(\operatorname{metric\text{-}group\text{-}ms}(\operatorname{hmg\_})).\; \operatorname{hmf\_}(\operatorname{opr}(\operatorname{hmg\_})(\operatorname{hqa\_}, \operatorname{hix\_})))), \operatorname{pts}(\operatorname{metric\text{-}group\text{-}ms}(\operatorname{hmg\_}))))))))
\end{equation*}

\par\medskip\noindent\textbf{Definition} (\texttt{haar-integral}).\;$\operatorname{haar\text{-}integral}(\operatorname{hng\_}) :=$ $(\iota\, \operatorname{hnu\_}.\;\operatorname{is\text{-}haar\text{-}integral}(\operatorname{hng\_}, \operatorname{hnu\_}))$

\vnbbookfile{theorem-library/cvx-alg}{theorem-library/cvx-alg}
\noindent{\small cvx-alg.scm -- the regrouping identity for convex combinations in a real normed vector space, and the bounds on its coefficients: the algebra under "the join of a convex set to a point is convex" and "the convex hull is convex" (compact-groups notes, Lemma 1.12). Agent ISOM, 2026-10-26.  Helper prefix: cva-.  Library: ...}\par
\vnbprop{convex-combination-regroup}{\vnbstmt{\begin{array}{@{}l@{}} \forall m, b, c, a, s, d, f, g. \\ \quad \operatorname{is\text{-}normed\text{-}vector\text{-}space}(m) \wedge \\ \quad b \in \operatorname{vec}(m) \wedge \\ \quad c \in \operatorname{vec}(m) \wedge \\ \quad a \in \operatorname{vec}(m) \wedge \\ \quad s \in \mathbb{R} \wedge \\ \quad d \in \mathbb{R} \wedge \\ \quad f \in \mathbb{R} \wedge \\ \quad g \in \mathbb{R} \wedge \\ \quad ((1 - (((1 - s) \cdot d) + (s \cdot f))) \cdot (1 - g)) = ((1 - s) \cdot (1 - d)) \wedge \\ \quad ((1 - (((1 - s) \cdot d) + (s \cdot f))) \cdot g) = (s \cdot (1 - f)) \Rightarrow \\ \quad \quad \operatorname{vadd}(m)(\operatorname{act}(m)((1 - s), \operatorname{vadd}(m)(\operatorname{act}(m)((1 - d), b), \operatorname{act}(m)(d, a))), \\ \quad \quad \quad \operatorname{act}(m)(s, \operatorname{vadd}(m)(\operatorname{act}(m)((1 - f), c), \operatorname{act}(m)(f, a)))) \\ \quad \quad =\; \operatorname{vadd}(m)(\operatorname{act}(m)((1 - (((1 - s) \cdot d) + (s \cdot f))), \operatorname{vadd}(m)(\operatorname{act}(m)((1 - g), b), \operatorname{act}(m)(g, c))), \\ \quad \quad \quad \operatorname{act}(m)((((1 - s) \cdot d) + (s \cdot f)), a)) \end{array}}}
\vnbprop{cvx-mul-nonneg}{\vnbstmt{\forall x, y.\; (((x \in \mathbb{R}) \wedge ((y \in \mathbb{R}) \wedge ((0 \leq x) \wedge (0 \leq y)))) \Rightarrow (0 \leq (x \cdot y)))}}
\vnbprop{cvx-mul-le-self}{\vnbstmt{\forall c, x.\; (((c \in \mathbb{R}) \wedge ((x \in \mathbb{R}) \wedge ((0 \leq c) \wedge (x \leq 1)))) \Rightarrow ((c \cdot x) \leq c))}}
\vnbprop{convex-regroup-coefficients}{\vnbstmt{\begin{array}{@{}l@{}} \forall s, a, b. \\ \quad ((s \in \mathbb{R}) \wedge ((0 \leq s) \wedge ((s \leq 1) \wedge ((a \in \mathbb{R}) \wedge ((0 \leq a) \wedge ((a \leq 1) \wedge ((b \in \mathbb{R}) \wedge ((0 \leq b) \wedge (b \leq 1))))))))) \Rightarrow  \\ \quad \quad 0 \leq (((1 - s) \cdot a) + (s \cdot b)) \wedge \\ \quad \quad (((1 - s) \cdot a) + (s \cdot b)) \leq 1 \wedge \\ \quad \quad \exists m \in \mathbb{R}. \\ \quad \quad \quad 0 \leq m \wedge \\ \quad \quad \quad m \leq 1 \wedge \\ \quad \quad \quad ((1 - (((1 - s) \cdot a) + (s \cdot b))) \cdot (1 - m)) = ((1 - s) \cdot (1 - a)) \wedge \\ \quad \quad \quad ((1 - (((1 - s) \cdot a) + (s \cdot b))) \cdot m) = (s \cdot (1 - b)) \end{array}}}

\vnbbookfile{theorem-library/cvx-hull}{theorem-library/cvx-hull}
\noindent{\small cvx-hull.scm -- the join of a convex set to a point is convex; the convex hull (the finite convex combinations, CONV-HULL) contains E, lies in VEC, is convex and lies in every convex set containing E (compact-groups notes, Lemma 1.12, the convexity half). Agent ISOM, 2026-10-26.  Helper prefix: cvh-.  Library: ...}\par
\vnbprop{conv-join-unfold}{\vnbstmt{\begin{array}{@{}l@{}} \forall m, c, a. \\ \quad (\operatorname{conv\text{-}join}(m, c, a) \simeq \{z \in \operatorname{vec}(m) : \exists y \in c, t \in \mathbb{R}.\; ((0 \leq t) \wedge ((t \leq 1) \wedge (z = \operatorname{vadd}(m)(\operatorname{act}(m)((1 - t), y), \operatorname{act}(m)(t, a)))))\}) \end{array}}}
\vnbprop{conv-join-elim}{\vnbstmt{\begin{array}{@{}l@{}} \forall m, c, a.; \forall z \in \operatorname{conv\text{-}join}(m, c, a). \\ \quad z \in \operatorname{vec}(m) \wedge \\ \quad \exists y \in c, t \in \mathbb{R}. \\ \quad \quad (0 \leq t) \wedge ((t \leq 1) \wedge (z = \operatorname{vadd}(m)(\operatorname{act}(m)((1 - t), y), \operatorname{act}(m)(t, a)))) \end{array}}}
\vnbprop{conv-join-intro}{\vnbstmt{\begin{array}{@{}l@{}} \forall m, c, a, y, t. \\ \quad \operatorname{is\text{-}normed\text{-}vector\text{-}space}(m) \wedge \\ \quad y \in \operatorname{vec}(m) \wedge \\ \quad a \in \operatorname{vec}(m) \wedge \\ \quad y \in c \wedge \\ \quad t \in \mathbb{R} \wedge \\ \quad 0 \leq t \wedge \\ \quad t \leq 1 \Rightarrow \\ \quad \quad \operatorname{vadd}(m)(\operatorname{act}(m)((1 - t), y), \operatorname{act}(m)(t, a)) \in \operatorname{conv\text{-}join}(m, c, a) \end{array}}}
\vnbprop{conv-join-convex}{\vnbstmt{\begin{array}{@{}l@{}} \forall m, c, a. \\ \quad (\operatorname{is\text{-}normed\text{-}vector\text{-}space}(m) \wedge (\operatorname{is\text{-}convex\text{-}subset}(m, c) \wedge (a \in \operatorname{vec}(m)))) \Rightarrow  \\ \quad \quad \operatorname{is\text{-}convex\text{-}subset}(m, \operatorname{conv\text{-}join}(m, c, a)) \end{array}}}
\vnbprop{conv-level-is-set}{\vnbstmt{\begin{array}{@{}l@{}} \forall n \in \mathbb{N}, m, e. \\ \quad (\operatorname{is\text{-}normed\text{-}vector\text{-}space}(m) \wedge (e \subseteq \operatorname{vec}(m))) \Rightarrow (\operatorname{conv\text{-}level}(m, e, n) \in \mathbf{Set}) \end{array}}}
\vnbprop{conv-level-succ-elim}{\vnbstmt{\begin{array}{@{}l@{}} \forall m, e, n, z. \\ \quad (\operatorname{is\text{-}normed\text{-}vector\text{-}space}(m) \wedge ((e \subseteq \operatorname{vec}(m)) \wedge ((n \in \mathbb{N}) \wedge (z \in \operatorname{conv\text{-}level}(m, e, (n + 1)))))) \Rightarrow  \\ \quad \quad z \in \operatorname{vec}(m) \wedge \\ \quad \quad \exists x \in \operatorname{conv\text{-}level}(m, e, n), y \in e, t \in \mathbb{R}. \\ \quad \quad \quad (0 \leq t) \wedge ((t \leq 1) \wedge (z = \operatorname{vadd}(m)(\operatorname{act}(m)((1 - t), x), \operatorname{act}(m)(t, y)))) \end{array}}}
\vnbprop{conv-level-succ-intro}{\vnbstmt{\begin{array}{@{}l@{}} \forall m, e, n, x, y, t. \\ \quad \operatorname{is\text{-}normed\text{-}vector\text{-}space}(m) \wedge \\ \quad e \subseteq \operatorname{vec}(m) \wedge \\ \quad n \in \mathbb{N} \wedge \\ \quad x \in \operatorname{conv\text{-}level}(m, e, n) \wedge \\ \quad x \in \operatorname{vec}(m) \wedge \\ \quad y \in e \wedge \\ \quad y \in \operatorname{vec}(m) \wedge \\ \quad t \in \mathbb{R} \wedge \\ \quad 0 \leq t \wedge \\ \quad t \leq 1 \Rightarrow \\ \quad \quad \operatorname{vadd}(m)(\operatorname{act}(m)((1 - t), x), \operatorname{act}(m)(t, y)) \in \operatorname{conv\text{-}level}(m, e, (n + 1)) \end{array}}}
\vnbprop{conv-level-sub-vec}{\vnbstmt{\begin{array}{@{}l@{}} \forall n \in \mathbb{N}, m, e. \\ \quad (\operatorname{is\text{-}normed\text{-}vector\text{-}space}(m) \wedge (e \subseteq \operatorname{vec}(m))) \Rightarrow  \\ \quad \quad \operatorname{conv\text{-}level}(m, e, n) \subseteq \operatorname{vec}(m) \end{array}}}
\vnbprop{conv-hull-unfold}{\vnbstmt{\forall m, e.\; (\operatorname{conv\text{-}hull}(m, e) \simeq \operatorname{big\text{-}union}(n, \mathbb{N}, \operatorname{conv\text{-}level}(m, e, n)))}}
\vnbprop{conv-hull-intro}{\vnbstmt{\begin{array}{@{}l@{}} \forall m, e, k, z. \\ \quad ((k \in \mathbb{N}) \wedge (z \in \operatorname{conv\text{-}level}(m, e, k))) \Rightarrow (z \in \operatorname{conv\text{-}hull}(m, e)) \end{array}}}
\vnbprop{conv-hull-elim}{\vnbstmt{\begin{array}{@{}l@{}} \forall m, e.; \forall z \in \operatorname{conv\text{-}hull}(m, e). \exists k \in \mathbb{N}. \\ \quad z \in \operatorname{conv\text{-}level}(m, e, k) \end{array}}}
\vnbprop{conv-hull-contains}{\vnbstmt{\begin{array}{@{}l@{}} \forall m, e. \\ \quad (\operatorname{is\text{-}normed\text{-}vector\text{-}space}(m) \wedge (e \subseteq \operatorname{vec}(m))) \Rightarrow (e \subseteq \operatorname{conv\text{-}hull}(m, e)) \end{array}}}
\vnbprop{conv-hull-sub-vec}{\vnbstmt{\begin{array}{@{}l@{}} \forall m, e. \\ \quad (\operatorname{is\text{-}normed\text{-}vector\text{-}space}(m) \wedge (e \subseteq \operatorname{vec}(m))) \Rightarrow  \\ \quad \quad \operatorname{conv\text{-}hull}(m, e) \subseteq \operatorname{vec}(m) \end{array}}}
\vnbprop{convex-combination-at-zero}{\vnbstmt{\begin{array}{@{}l@{}} \forall m, u, y. \\ \quad (\operatorname{is\text{-}normed\text{-}vector\text{-}space}(m) \wedge ((u \in \operatorname{vec}(m)) \wedge (y \in \operatorname{vec}(m)))) \Rightarrow  \\ \quad \quad \operatorname{vadd}(m)(\operatorname{act}(m)(1, u), \operatorname{act}(m)(0, y)) = u \end{array}}}
\vnbprop{conv-hull-convex-step}{\vnbstmt{\begin{array}{@{}l@{}} \forall k \in \mathbb{N}, m, e.; \forall u \in \operatorname{conv\text{-}hull}(m, e), w \in \operatorname{conv\text{-}level}(m, e, k), s \in \mathbb{R}. \\ \quad (\operatorname{is\text{-}normed\text{-}vector\text{-}space}(m) \wedge ((e \subseteq \operatorname{vec}(m)) \wedge ((0 \leq s) \wedge (s \leq 1)))) \Rightarrow  \\ \quad \quad \operatorname{vadd}(m)(\operatorname{act}(m)((1 - s), u), \operatorname{act}(m)(s, w)) \in \operatorname{conv\text{-}hull}(m, e) \end{array}}}
\vnbprop{conv-hull-convex}{\vnbstmt{\begin{array}{@{}l@{}} \forall m, e. \\ \quad (\operatorname{is\text{-}normed\text{-}vector\text{-}space}(m) \wedge (e \subseteq \operatorname{vec}(m))) \Rightarrow  \\ \quad \quad \operatorname{is\text{-}convex\text{-}subset}(m, \operatorname{conv\text{-}hull}(m, e)) \end{array}}}
\vnbprop{conv-level-in-convex}{\vnbstmt{\begin{array}{@{}l@{}} \forall k \in \mathbb{N}, m, e, a. \\ \quad (\operatorname{is\text{-}normed\text{-}vector\text{-}space}(m) \wedge ((e \subseteq \operatorname{vec}(m)) \wedge (\operatorname{is\text{-}convex\text{-}subset}(m, a) \wedge (e \subseteq a)))) \Rightarrow  \\ \quad \quad \operatorname{conv\text{-}level}(m, e, k) \subseteq a \end{array}}}
\vnbprop{conv-hull-least}{\vnbstmt{\begin{array}{@{}l@{}} \forall m, e, k. \\ \quad (\operatorname{is\text{-}normed\text{-}vector\text{-}space}(m) \wedge ((e \subseteq \operatorname{vec}(m)) \wedge (\operatorname{is\text{-}convex\text{-}subset}(m, k) \wedge (e \subseteq k)))) \Rightarrow  \\ \quad \quad \operatorname{conv\text{-}hull}(m, e) \subseteq k \end{array}}}

\vnbbookfile{theorem-library/sip-d1}{theorem-library/sip-d1}
\noindent{\small sip-d1.scm -- REAL HILBERT SPACES BY THE COMPLEXIFICATION ROUTE: the norm and metric of a pre-Hilbert space, completeness of the complexification, the orthogonal projection onto a closed subspace, and the complex corollary.  (SIP-D, 2026-10-19; scratchpad/triage/SIP-D.md.)  All PROVEN modulo 0.  Vocabulary:}\par
\vnbprop{sip-rip-neg-self}{\vnbstmt{\begin{array}{@{}l@{}} \forall v.; \forall x \in \operatorname{vec}(v). \\ \quad \operatorname{is\text{-}semi\text{-}inner\text{-}product\text{-}space}(v) \Rightarrow  \\ \quad \quad \langle\cdot\mid\cdot\rangle_{v}(\operatorname{vneg}(v)(x), \operatorname{vneg}(v)(x)) = \langle\cdot\mid\cdot\rangle_{v}(x, x) \end{array}}}
\vnbprop{sip-minkowski}{\vnbstmt{\begin{array}{@{}l@{}} \forall v.; \forall x, y \in \operatorname{vec}(v). \\ \quad \operatorname{is\text{-}semi\text{-}inner\text{-}product\text{-}space}(v) \Rightarrow  \\ \quad \quad \sqrt{\langle\cdot\mid\cdot\rangle_{v}(\operatorname{vadd}(v)(x, y), \operatorname{vadd}(v)(x, y))} \leq (\sqrt{\langle\cdot\mid\cdot\rangle_{v}(x, x)} + \sqrt{\langle\cdot\mid\cdot\rangle_{v}(y, y)}) \end{array}}}
\vnbprop{sip-metric-space-unfold}{\vnbstmt{\forall v.\; (\operatorname{sip\text{-}metric\text{-}space}(v) \simeq \operatorname{nag\text{-}metric\text{-}space}(\operatorname{sip\text{-}normed\text{-}ag}(v)))}}
\vnbprop{sip-normed-ag-carrier}{\vnbstmt{\begin{array}{@{}l@{}} \forall v. \\ \quad \operatorname{is\text{-}semi\text{-}inner\text{-}product\text{-}space}(v) \Rightarrow (\operatorname{carr}(\operatorname{sip\text{-}normed\text{-}ag}(v)) \simeq \operatorname{vec}(v)) \end{array}}}
\vnbprop{sip-normed-ag-opr}{\vnbstmt{\begin{array}{@{}l@{}} \forall v. \\ \quad \operatorname{is\text{-}semi\text{-}inner\text{-}product\text{-}space}(v) \Rightarrow (\operatorname{opr}(\operatorname{sip\text{-}normed\text{-}ag}(v)) \simeq \operatorname{vadd}(v)) \end{array}}}
\vnbprop{sip-normed-ag-inv}{\vnbstmt{\begin{array}{@{}l@{}} \forall v. \\ \quad \operatorname{is\text{-}semi\text{-}inner\text{-}product\text{-}space}(v) \Rightarrow (\operatorname{inv}(\operatorname{sip\text{-}normed\text{-}ag}(v)) \simeq \operatorname{vneg}(v)) \end{array}}}
\vnbprop{sip-normed-ag-norm-value}{\vnbstmt{\begin{array}{@{}l@{}} \forall v.; \forall x \in \operatorname{vec}(v). \\ \quad \operatorname{is\text{-}semi\text{-}inner\text{-}product\text{-}space}(v) \Rightarrow  \\ \quad \quad \operatorname{nrm}(\operatorname{sip\text{-}normed\text{-}ag}(v))(x) = \sqrt{\langle\cdot\mid\cdot\rangle_{v}(x, x)} \end{array}}}
\vnbprop{sip-normed-ag-is-normed-ag}{\vnbstmt{\forall v.\; (\operatorname{is\text{-}pre\text{-}hilbert\text{-}space}(v) \Rightarrow \operatorname{is\text{-}normed\text{-}ag}(\operatorname{sip\text{-}normed\text{-}ag}(v)))}}
\vnbprop{sip-metric-space-is-metric-space}{\vnbstmt{\begin{array}{@{}l@{}} \forall v. \\ \quad \operatorname{is\text{-}pre\text{-}hilbert\text{-}space}(v) \Rightarrow \operatorname{is\text{-}metric\text{-}space}(\operatorname{sip\text{-}metric\text{-}space}(v)) \end{array}}}
\vnbprop{sip-metric-space-pts}{\vnbstmt{\forall v.\; (\operatorname{is\text{-}pre\text{-}hilbert\text{-}space}(v) \Rightarrow (\operatorname{pts}(\operatorname{sip\text{-}metric\text{-}space}(v)) \simeq \operatorname{vec}(v)))}}
\vnbprop{sip-metric-space-dist}{\vnbstmt{\begin{array}{@{}l@{}} \forall v.; \forall x, y \in \operatorname{vec}(v). \\ \quad \operatorname{is\text{-}pre\text{-}hilbert\text{-}space}(v) \Rightarrow  \\ \quad \quad \operatorname{dist}(\operatorname{sip\text{-}metric\text{-}space}(v))(x, y) \\ \quad \quad =\; \sqrt{\langle\cdot\mid\cdot\rangle_{v}(\operatorname{vadd}(v)(x, \operatorname{vneg}(v)(y)), \operatorname{vadd}(v)(x, \operatorname{vneg}(v)(y)))} \end{array}}}
\vnbprop{real-hilbert-space-is-phs}{\vnbstmt{\forall v.\; (\operatorname{is\text{-}real\text{-}hilbert\text{-}space}(v) \Rightarrow \operatorname{is\text{-}pre\text{-}hilbert\text{-}space}(v))}}
\vnbprop{real-hilbert-space-complete}{\vnbstmt{\forall v.\; (\operatorname{is\text{-}real\text{-}hilbert\text{-}space}(v) \Rightarrow \operatorname{is\text{-}complete}(\operatorname{sip\text{-}metric\text{-}space}(v)))}}
\vnbprop{cxw-vec}{\vnbstmt{\begin{array}{@{}l@{}} \forall v. \\ \quad \operatorname{is\text{-}semi\text{-}inner\text{-}product\text{-}space}(v) \Rightarrow  \\ \quad \quad (\operatorname{vec}(\operatorname{csip\text{-}cips}(\operatorname{complexify}(v))) \simeq (\operatorname{vec}(v) \times \operatorname{vec}(v))) \end{array}}}
\vnbprop{cxw-vzero}{\vnbstmt{\begin{array}{@{}l@{}} \forall v. \\ \quad \operatorname{is\text{-}semi\text{-}inner\text{-}product\text{-}space}(v) \Rightarrow  \\ \quad \quad (\operatorname{vzero}(\operatorname{csip\text{-}cips}(\operatorname{complexify}(v))) \simeq {[}\operatorname{vzero}(v), \operatorname{vzero}(v)]) \end{array}}}
\vnbprop{cxw-vadd-value}{\vnbstmt{\begin{array}{@{}l@{}} \forall v, x, \mathit{qx}. \\ \quad \operatorname{is\text{-}semi\text{-}inner\text{-}product\text{-}space}(v) \wedge \\ \quad x \in (\operatorname{vec}(v) \times \operatorname{vec}(v)) \wedge \\ \quad \mathit{qx} \in (\operatorname{vec}(v) \times \operatorname{vec}(v)) \Rightarrow \\ \quad \quad \operatorname{vadd}(\operatorname{csip\text{-}cips}(\operatorname{complexify}(v)))(x, \mathit{qx}) \\ \quad \quad =\; {[}\operatorname{vadd}(v)(\operatorname{nth}(1, x), \operatorname{nth}(1, \mathit{qx})), \operatorname{vadd}(v)(\operatorname{nth}(2, x), \operatorname{nth}(2, \mathit{qx}))] \end{array}}}
\vnbprop{cxw-vneg-value}{\vnbstmt{\begin{array}{@{}l@{}} \forall v.; \forall x \in (\operatorname{vec}(v) \times \operatorname{vec}(v)). \\ \quad \operatorname{is\text{-}semi\text{-}inner\text{-}product\text{-}space}(v) \Rightarrow  \\ \quad \quad \operatorname{vneg}(\operatorname{csip\text{-}cips}(\operatorname{complexify}(v)))(x) = {[}\operatorname{vneg}(v)(\operatorname{nth}(1, x)), \operatorname{vneg}(v)(\operatorname{nth}(2, x))] \end{array}}}
\vnbprop{cxw-ip-value}{\vnbstmt{\begin{array}{@{}l@{}} \forall v, x, \mathit{qx}. \\ \quad \operatorname{is\text{-}semi\text{-}inner\text{-}product\text{-}space}(v) \wedge \\ \quad x \in (\operatorname{vec}(v) \times \operatorname{vec}(v)) \wedge \\ \quad \mathit{qx} \in (\operatorname{vec}(v) \times \operatorname{vec}(v)) \Rightarrow \\ \quad \quad \langle\cdot,\cdot\rangle_{\operatorname{csip\text{-}cips}(\operatorname{complexify}(v))}(x, \mathit{qx}) \\ \quad \quad =\; ((\langle\cdot\mid\cdot\rangle_{v}(\operatorname{nth}(1, x), \operatorname{nth}(1, \mathit{qx})) + \langle\cdot\mid\cdot\rangle_{v}(\operatorname{nth}(2, x), \operatorname{nth}(2, \mathit{qx}))) - (+i \cdot (\langle\cdot\mid\cdot\rangle_{v}(\operatorname{vneg}(v)(\operatorname{nth}(2, x)), \operatorname{nth}(1, \mathit{qx})) + \langle\cdot\mid\cdot\rangle_{v}(\operatorname{nth}(1, x), \operatorname{nth}(2, \mathit{qx}))))) \end{array}}}
\vnbprop{cxw-ip-self}{\vnbstmt{\begin{array}{@{}l@{}} \forall v.; \forall x \in (\operatorname{vec}(v) \times \operatorname{vec}(v)). \\ \quad \operatorname{is\text{-}semi\text{-}inner\text{-}product\text{-}space}(v) \Rightarrow  \\ \quad \quad \langle\cdot,\cdot\rangle_{\operatorname{csip\text{-}cips}(\operatorname{complexify}(v))}(x, x) \\ \quad \quad =\; \langle\cdot\mid\cdot\rangle_{v}(\operatorname{nth}(1, x), \operatorname{nth}(1, x)) + \langle\cdot\mid\cdot\rangle_{v}(\operatorname{nth}(2, x), \operatorname{nth}(2, x)) \end{array}}}
\vnbprop{cxw-act-value}{\vnbstmt{\begin{array}{@{}l@{}} \forall v, z, x. \\ \quad (\operatorname{is\text{-}semi\text{-}inner\text{-}product\text{-}space}(v) \wedge ((z \in \mathbb{C}) \wedge (x \in (\operatorname{vec}(v) \times \operatorname{vec}(v))))) \Rightarrow  \\ \quad \quad \operatorname{act}(\operatorname{csip\text{-}cips}(\operatorname{complexify}(v)))(z, x) \\ \quad \quad =\; {[}\operatorname{vadd}(v)(\operatorname{act}(v)(\operatorname{real\text{-}part}(z), \operatorname{nth}(1, x)), \operatorname{act}(v)(\operatorname{imag\text{-}part}(z), \operatorname{vneg}(v)(\operatorname{nth}(2, x)))), \operatorname{vadd}(v)(\operatorname{act}(v)(\operatorname{real\text{-}part}(z), \operatorname{nth}(2, x)), \operatorname{act}(v)(\operatorname{imag\text{-}part}(z), \operatorname{nth}(1, x)))] \end{array}}}
\vnbprop{complexify-is-cphs}{\vnbstmt{\begin{array}{@{}l@{}} \forall v. \\ \quad \operatorname{is\text{-}pre\text{-}hilbert\text{-}space}(v) \Rightarrow \operatorname{is\text{-}complex\text{-}pre\text{-}hilbert\text{-}space}(\operatorname{complexify}(v)) \end{array}}}
\vnbprop{cxw-is-cips}{\vnbstmt{\begin{array}{@{}l@{}} \forall v. \\ \quad \operatorname{is\text{-}pre\text{-}hilbert\text{-}space}(v) \Rightarrow  \\ \quad \quad \operatorname{is\text{-}complex\text{-}inner\text{-}product\text{-}space}(\operatorname{csip\text{-}cips}(\operatorname{complexify}(v))) \end{array}}}
\vnbprop{cxw-dist}{\vnbstmt{\begin{array}{@{}l@{}} \forall v, x, \mathit{qx}. \\ \quad (\operatorname{is\text{-}pre\text{-}hilbert\text{-}space}(v) \wedge ((x \in (\operatorname{vec}(v) \times \operatorname{vec}(v))) \wedge (\mathit{qx} \in (\operatorname{vec}(v) \times \operatorname{vec}(v))))) \Rightarrow  \\ \quad \quad \operatorname{dist}(\operatorname{ip\text{-}metric\text{-}space}(\operatorname{csip\text{-}cips}(\operatorname{complexify}(v))))(x, \mathit{qx}) \\ \quad \quad =\; \sqrt{(\langle\cdot\mid\cdot\rangle_{v}(\operatorname{vadd}(v)(\operatorname{nth}(1, x), \operatorname{vneg}(v)(\operatorname{nth}(1, \mathit{qx}))), \operatorname{vadd}(v)(\operatorname{nth}(1, x), \operatorname{vneg}(v)(\operatorname{nth}(1, \mathit{qx})))) + \langle\cdot\mid\cdot\rangle_{v}(\operatorname{vadd}(v)(\operatorname{nth}(2, x), \operatorname{vneg}(v)(\operatorname{nth}(2, \mathit{qx}))), \operatorname{vadd}(v)(\operatorname{nth}(2, x), \operatorname{vneg}(v)(\operatorname{nth}(2, \mathit{qx})))))} \end{array}}}
\vnbprop{cxw-dist-comp1}{\vnbstmt{\begin{array}{@{}l@{}} \forall v, x, \mathit{qx}. \\ \quad (\operatorname{is\text{-}pre\text{-}hilbert\text{-}space}(v) \wedge ((x \in (\operatorname{vec}(v) \times \operatorname{vec}(v))) \wedge (\mathit{qx} \in (\operatorname{vec}(v) \times \operatorname{vec}(v))))) \Rightarrow  \\ \quad \quad \operatorname{dist}(\operatorname{sip\text{-}metric\text{-}space}(v))(\operatorname{nth}(1, x), \operatorname{nth}(1, \mathit{qx})) \\ \quad \quad \leq\; \operatorname{dist}(\operatorname{ip\text{-}metric\text{-}space}(\operatorname{csip\text{-}cips}(\operatorname{complexify}(v))))(x, \mathit{qx}) \end{array}}}
\vnbprop{cxw-dist-comp2}{\vnbstmt{\begin{array}{@{}l@{}} \forall v, x, \mathit{qx}. \\ \quad (\operatorname{is\text{-}pre\text{-}hilbert\text{-}space}(v) \wedge ((x \in (\operatorname{vec}(v) \times \operatorname{vec}(v))) \wedge (\mathit{qx} \in (\operatorname{vec}(v) \times \operatorname{vec}(v))))) \Rightarrow  \\ \quad \quad \operatorname{dist}(\operatorname{sip\text{-}metric\text{-}space}(v))(\operatorname{nth}(2, x), \operatorname{nth}(2, \mathit{qx})) \\ \quad \quad \leq\; \operatorname{dist}(\operatorname{ip\text{-}metric\text{-}space}(\operatorname{csip\text{-}cips}(\operatorname{complexify}(v))))(x, \mathit{qx}) \end{array}}}
\vnbprop{cxw-dist-le-sum}{\vnbstmt{\begin{array}{@{}l@{}} \forall v, x, \mathit{qx}. \\ \quad (\operatorname{is\text{-}pre\text{-}hilbert\text{-}space}(v) \wedge ((x \in (\operatorname{vec}(v) \times \operatorname{vec}(v))) \wedge (\mathit{qx} \in (\operatorname{vec}(v) \times \operatorname{vec}(v))))) \Rightarrow  \\ \quad \quad \operatorname{dist}(\operatorname{ip\text{-}metric\text{-}space}(\operatorname{csip\text{-}cips}(\operatorname{complexify}(v))))(x, \mathit{qx}) \\ \quad \quad \leq\; \operatorname{dist}(\operatorname{sip\text{-}metric\text{-}space}(v))(\operatorname{nth}(1, x), \operatorname{nth}(1, \mathit{qx})) + \operatorname{dist}(\operatorname{sip\text{-}metric\text{-}space}(v))(\operatorname{nth}(2, x), \operatorname{nth}(2, \mathit{qx})) \end{array}}}
\vnbprop{ms-cauchy-of-dominated}{\vnbstmt{\begin{array}{@{}l@{}} \forall s, \mathit{tm}, f, g. \\ \quad \operatorname{is\text{-}cauchy\text{-}seq}(s, f) \wedge \\ \quad \operatorname{is\text{-}metric\text{-}space}(\mathit{tm}) \wedge \\ \quad g \in (\mathbb{N} \to \operatorname{pts}(\mathit{tm})) \wedge \\ \quad \forall m, n \in \mathbb{N}.\; (\operatorname{dist}(\mathit{tm})(g(m), g(n)) \leq \operatorname{dist}(s)(f(m), f(n))) \Rightarrow \\ \quad \quad \operatorname{is\text{-}cauchy\text{-}seq}(\mathit{tm}, g) \end{array}}}
\vnbprop{ms-converges-of-dominated-sum}{\vnbstmt{\begin{array}{@{}l@{}} \forall s, t_{1}, t_{2}, f, g_{1}, g_{2}, l, l_{1}, l_{2}. \\ \quad \operatorname{is\text{-}metric\text{-}space}(s) \wedge \\ \quad f \in (\mathbb{N} \to \operatorname{pts}(s)) \wedge \\ \quad l \in \operatorname{pts}(s) \wedge \\ \quad \operatorname{converges\text{-}to}(t_{1}, g_{1}, l_{1}) \wedge \\ \quad \operatorname{converges\text{-}to}(t_{2}, g_{2}, l_{2}) \wedge \\ \quad \forall n \in \mathbb{N}. \\ \quad \quad \operatorname{dist}(s)(f(n), l) \leq (\operatorname{dist}(t_{1})(\operatorname{g1}(n), l_{1}) + \operatorname{dist}(t_{2})(\operatorname{g2}(n), l_{2})) \Rightarrow \\ \quad \quad \operatorname{converges\text{-}to}(s, f, l) \end{array}}}
\vnbprop{complexify-complete}{\vnbstmt{\begin{array}{@{}l@{}} \forall v. \\ \quad \operatorname{is\text{-}real\text{-}hilbert\text{-}space}(v) \Rightarrow  \\ \quad \quad \operatorname{is\text{-}complete}(\operatorname{ip\text{-}metric\text{-}space}(\operatorname{csip\text{-}cips}(\operatorname{complexify}(v)))) \end{array}}}
\vnbprop{complexify-hilbert}{\vnbstmt{\begin{array}{@{}l@{}} \forall v. \\ \quad \operatorname{is\text{-}real\text{-}hilbert\text{-}space}(v) \Rightarrow \operatorname{is\text{-}hilbert\text{-}space}(\operatorname{csip\text{-}cips}(\operatorname{complexify}(v))) \end{array}}}
\vnbprop{sip-perp-unfold}{\vnbstmt{\begin{array}{@{}l@{}} \forall v, s. \\ \quad ({s}^{\perp} \simeq \{p \in \operatorname{vec}(v) : \forall k \in s.\; (\langle\cdot\mid\cdot\rangle_{v}(p, k) = 0)\}) \end{array}}}
\vnbprop{sip-perp-membership}{\vnbstmt{\begin{array}{@{}l@{}} \forall v, s, z. \\ \quad (z \in {s}^{\perp}) \Leftrightarrow ((z \in \operatorname{vec}(v)) \wedge \forall k \in s.\; (\langle\cdot\mid\cdot\rangle_{v}(z, k) = 0)) \end{array}}}
\vnbprop{sip-perp-orth}{\vnbstmt{\forall v, s, z, k.\; (((z \in {s}^{\perp}) \wedge (k \in s)) \Rightarrow (\langle\cdot\mid\cdot\rangle_{v}(z, k) = 0))}}
\vnbprop{sip-perp-subset}{\vnbstmt{\begin{array}{@{}l@{}} \forall v, s. \\ \quad \operatorname{is\text{-}semi\text{-}inner\text{-}product\text{-}space}(v) \Rightarrow ({s}^{\perp} \subseteq \operatorname{vec}(v)) \end{array}}}
\vnbprop{sip-perp-is-submodule}{\vnbstmt{\begin{array}{@{}l@{}} \forall v, s. \\ \quad (\operatorname{is\text{-}semi\text{-}inner\text{-}product\text{-}space}(v) \wedge (s \subseteq \operatorname{vec}(v))) \Rightarrow  \\ \quad \quad \operatorname{is\text{-}submodule}(v, {s}^{\perp}) \end{array}}}
\vnbprop{complexify-submodule}{\vnbstmt{\begin{array}{@{}l@{}} \forall v, k. \\ \quad (\operatorname{is\text{-}pre\text{-}hilbert\text{-}space}(v) \wedge \operatorname{is\text{-}submodule}(v, k)) \Rightarrow  \\ \quad \quad \operatorname{is\text{-}submodule}(\operatorname{csip\text{-}cips}(\operatorname{complexify}(v)), (k \times k)) \end{array}}}
\vnbprop{complexify-closure-comp1}{\vnbstmt{\begin{array}{@{}l@{}} \forall v, k.; \forall z \in \operatorname{closure}(\operatorname{ip\text{-}metric\text{-}space}(\operatorname{csip\text{-}cips}(\operatorname{complexify}(v))), (k \times k)). \\ \quad (\operatorname{is\text{-}pre\text{-}hilbert\text{-}space}(v) \wedge (k \subseteq \operatorname{vec}(v))) \Rightarrow  \\ \quad \quad \operatorname{nth}(1, z) \in \operatorname{closure}(\operatorname{sip\text{-}metric\text{-}space}(v), k) \end{array}}}
\vnbprop{complexify-closure-comp2}{\vnbstmt{\begin{array}{@{}l@{}} \forall v, k.; \forall z \in \operatorname{closure}(\operatorname{ip\text{-}metric\text{-}space}(\operatorname{csip\text{-}cips}(\operatorname{complexify}(v))), (k \times k)). \\ \quad (\operatorname{is\text{-}pre\text{-}hilbert\text{-}space}(v) \wedge (k \subseteq \operatorname{vec}(v))) \Rightarrow  \\ \quad \quad \operatorname{nth}(2, z) \in \operatorname{closure}(\operatorname{sip\text{-}metric\text{-}space}(v), k) \end{array}}}
\vnbprop{complexify-closed}{\vnbstmt{\begin{array}{@{}l@{}} \forall v, k. \\ \quad (\operatorname{is\text{-}pre\text{-}hilbert\text{-}space}(v) \wedge \operatorname{is\text{-}closed}(\operatorname{sip\text{-}metric\text{-}space}(v), k)) \Rightarrow  \\ \quad \quad \operatorname{is\text{-}closed}(\operatorname{ip\text{-}metric\text{-}space}(\operatorname{csip\text{-}cips}(\operatorname{complexify}(v))), (k \times k)) \end{array}}}
\vnbprop{real-projection-point}{\vnbstmt{\begin{array}{@{}l@{}} \forall v, k, x. \\ \quad \operatorname{is\text{-}real\text{-}hilbert\text{-}space}(v) \wedge \\ \quad \operatorname{is\text{-}submodule}(v, k) \wedge \\ \quad \operatorname{is\text{-}closed}(\operatorname{sip\text{-}metric\text{-}space}(v), k) \wedge \\ \quad x \in \operatorname{vec}(v) \Rightarrow \\ \quad \quad \exists a \in k.\; (\operatorname{vadd}(v)(x, \operatorname{vneg}(v)(a)) \in {k}^{\perp}) \end{array}}}
\vnbprop{real-decomposition}{\vnbstmt{\begin{array}{@{}l@{}} \forall v, k, x. \\ \quad \operatorname{is\text{-}real\text{-}hilbert\text{-}space}(v) \wedge \\ \quad \operatorname{is\text{-}submodule}(v, k) \wedge \\ \quad \operatorname{is\text{-}closed}(\operatorname{sip\text{-}metric\text{-}space}(v), k) \wedge \\ \quad x \in \operatorname{vec}(v) \Rightarrow \\ \quad \quad \exists a \in k, b \in {k}^{\perp}.\; (x = \operatorname{vadd}(v)(a, b)) \end{array}}}
\vnbprop{real-decomposition-unique}{\vnbstmt{\begin{array}{@{}l@{}} \forall v, k, a_{1}, a_{2}, b_{1}, b_{2}. \\ \quad \operatorname{is\text{-}pre\text{-}hilbert\text{-}space}(v) \wedge \\ \quad \operatorname{is\text{-}submodule}(v, k) \wedge \\ \quad a_{1} \in k \wedge \\ \quad a_{2} \in k \wedge \\ \quad b_{1} \in {k}^{\perp} \wedge \\ \quad b_{2} \in {k}^{\perp} \wedge \\ \quad \operatorname{vadd}(v)(a_{1}, b_{1}) = \operatorname{vadd}(v)(a_{2}, b_{2}) \Rightarrow \\ \quad \quad (a_{1} = a_{2}) \wedge (b_{1} = b_{2}) \end{array}}}
\vnbprop{sip-pythagoras}{\vnbstmt{\begin{array}{@{}l@{}} \forall v, x, a, m. \\ \quad \operatorname{is\text{-}semi\text{-}inner\text{-}product\text{-}space}(v) \wedge \\ \quad x \in \operatorname{vec}(v) \wedge \\ \quad a \in \operatorname{vec}(v) \wedge \\ \quad m \in \operatorname{vec}(v) \wedge \\ \quad \langle\cdot\mid\cdot\rangle_{v}(\operatorname{vadd}(v)(x, \operatorname{vneg}(v)(a)), \operatorname{vadd}(v)(m, \operatorname{vneg}(v)(a))) = 0 \Rightarrow \\ \quad \quad \langle\cdot\mid\cdot\rangle_{v}(\operatorname{vadd}(v)(m, \operatorname{vneg}(v)(x)), \operatorname{vadd}(v)(m, \operatorname{vneg}(v)(x))) \\ \quad \quad =\; \langle\cdot\mid\cdot\rangle_{v}(\operatorname{vadd}(v)(a, \operatorname{vneg}(v)(x)), \operatorname{vadd}(v)(a, \operatorname{vneg}(v)(x))) + \langle\cdot\mid\cdot\rangle_{v}(\operatorname{vadd}(v)(m, \operatorname{vneg}(v)(a)), \operatorname{vadd}(v)(m, \operatorname{vneg}(v)(a))) \end{array}}}
\vnbprop{real-closest-point}{\vnbstmt{\begin{array}{@{}l@{}} \forall v, k, x. \\ \quad \operatorname{is\text{-}real\text{-}hilbert\text{-}space}(v) \wedge \\ \quad \operatorname{is\text{-}submodule}(v, k) \wedge \\ \quad \operatorname{is\text{-}closed}(\operatorname{sip\text{-}metric\text{-}space}(v), k) \wedge \\ \quad x \in \operatorname{vec}(v) \Rightarrow \\ \quad \quad \exists a \in k. \forall b \in k. \\ \quad \quad \quad \langle\cdot\mid\cdot\rangle_{v}(\operatorname{vadd}(v)(a, \operatorname{vneg}(v)(x)), \operatorname{vadd}(v)(a, \operatorname{vneg}(v)(x))) \\ \quad \quad \quad \leq\; \langle\cdot\mid\cdot\rangle_{v}(\operatorname{vadd}(v)(b, \operatorname{vneg}(v)(x)), \operatorname{vadd}(v)(b, \operatorname{vneg}(v)(x))) \end{array}}}
\vnbprop{real-closest-point-unique}{\vnbstmt{\begin{array}{@{}l@{}} \forall v, k, x, a_{1}, a_{2}. \\ \quad \operatorname{is\text{-}real\text{-}hilbert\text{-}space}(v) \wedge \\ \quad \operatorname{is\text{-}submodule}(v, k) \wedge \\ \quad \operatorname{is\text{-}closed}(\operatorname{sip\text{-}metric\text{-}space}(v), k) \wedge \\ \quad x \in \operatorname{vec}(v) \wedge \\ \quad a_{1} \in k \wedge \\ \quad a_{2} \in k \wedge \\ \quad \forall b \in k. \\ \quad \quad \langle\cdot\mid\cdot\rangle_{v}(\operatorname{vadd}(v)(a_{1}, \operatorname{vneg}(v)(x)), \operatorname{vadd}(v)(a_{1}, \operatorname{vneg}(v)(x))) \\ \quad \quad \leq\; \langle\cdot\mid\cdot\rangle_{v}(\operatorname{vadd}(v)(b, \operatorname{vneg}(v)(x)), \operatorname{vadd}(v)(b, \operatorname{vneg}(v)(x))) \wedge \\ \quad \forall b \in k. \\ \quad \quad \langle\cdot\mid\cdot\rangle_{v}(\operatorname{vadd}(v)(a_{2}, \operatorname{vneg}(v)(x)), \operatorname{vadd}(v)(a_{2}, \operatorname{vneg}(v)(x))) \\ \quad \quad \leq\; \langle\cdot\mid\cdot\rangle_{v}(\operatorname{vadd}(v)(b, \operatorname{vneg}(v)(x)), \operatorname{vadd}(v)(b, \operatorname{vneg}(v)(x))) \Rightarrow \\ \quad \quad a_{1} = a_{2} \end{array}}}
\vnbprop{real-orthogonal-projection-apply}{\vnbstmt{\begin{array}{@{}l@{}} \forall v, k, a, x. \\ \quad (\operatorname{is\text{-}real\text{-}orthogonal\text{-}projection}(v, k, a) \wedge (x \in \operatorname{vec}(v))) \Rightarrow  \\ \quad \quad (a(x) \in k) \wedge (\operatorname{vadd}(v)(x, \operatorname{vneg}(v)(a(x))) \in {k}^{\perp}) \end{array}}}
\vnbprop{real-orthogonal-projection-type}{\vnbstmt{\begin{array}{@{}l@{}} \forall v, k, a. \\ \quad \operatorname{is\text{-}real\text{-}orthogonal\text{-}projection}(v, k, a) \Rightarrow (a \in (\operatorname{vec}(v) \to \operatorname{vec}(v))) \end{array}}}
\vnbprop{real-orthogonal-projection-exists}{\vnbstmt{\begin{array}{@{}l@{}} \forall v, k. \\ \quad (\operatorname{is\text{-}real\text{-}hilbert\text{-}space}(v) \wedge (\operatorname{is\text{-}submodule}(v, k) \wedge \operatorname{is\text{-}closed}(\operatorname{sip\text{-}metric\text{-}space}(v), k))) \Rightarrow  \\ \quad \quad \exists \mathit{srpk}.\; \operatorname{is\text{-}real\text{-}orthogonal\text{-}projection}(v, k, \mathit{srpk}) \end{array}}}
\vnbprop{real-orthogonal-projection-unique}{\vnbstmt{\begin{array}{@{}l@{}} \forall v, k, a, b. \\ \quad \operatorname{is\text{-}pre\text{-}hilbert\text{-}space}(v) \wedge \\ \quad \operatorname{is\text{-}submodule}(v, k) \wedge \\ \quad \operatorname{is\text{-}real\text{-}orthogonal\text{-}projection}(v, k, a) \wedge \\ \quad \operatorname{is\text{-}real\text{-}orthogonal\text{-}projection}(v, k, b) \Rightarrow \\ \quad \quad a = b \end{array}}}
\vnbprop{real-orthogonal-projection-additive}{\vnbstmt{\begin{array}{@{}l@{}} \forall v, k, a, x, y. \\ \quad \operatorname{is\text{-}pre\text{-}hilbert\text{-}space}(v) \wedge \\ \quad \operatorname{is\text{-}submodule}(v, k) \wedge \\ \quad \operatorname{is\text{-}real\text{-}orthogonal\text{-}projection}(v, k, a) \wedge \\ \quad x \in \operatorname{vec}(v) \wedge \\ \quad y \in \operatorname{vec}(v) \Rightarrow \\ \quad \quad a(\operatorname{vadd}(v)(x, y)) = \operatorname{vadd}(v)(a(x), a(y)) \end{array}}}
\vnbprop{real-orthogonal-projection-homogeneous}{\vnbstmt{\begin{array}{@{}l@{}} \forall v, k, a, c, x. \\ \quad \operatorname{is\text{-}pre\text{-}hilbert\text{-}space}(v) \wedge \\ \quad \operatorname{is\text{-}submodule}(v, k) \wedge \\ \quad \operatorname{is\text{-}real\text{-}orthogonal\text{-}projection}(v, k, a) \wedge \\ \quad c \in \mathbb{R} \wedge \\ \quad x \in \operatorname{vec}(v) \Rightarrow \\ \quad \quad a(\operatorname{act}(v)(c, x)) = \operatorname{act}(v)(c, a(x)) \end{array}}}
\vnbprop{real-orthogonal-projection-fixes}{\vnbstmt{\begin{array}{@{}l@{}} \forall v, k, a, y. \\ \quad (\operatorname{is\text{-}pre\text{-}hilbert\text{-}space}(v) \wedge (\operatorname{is\text{-}submodule}(v, k) \wedge (\operatorname{is\text{-}real\text{-}orthogonal\text{-}projection}(v, k, a) \wedge (y \in k)))) \Rightarrow  \\ \quad \quad a(y) = y \end{array}}}
\vnbprop{real-orthogonal-projection-idempotent}{\vnbstmt{\begin{array}{@{}l@{}} \forall v, k, a, x. \\ \quad \operatorname{is\text{-}pre\text{-}hilbert\text{-}space}(v) \wedge \\ \quad \operatorname{is\text{-}submodule}(v, k) \wedge \\ \quad \operatorname{is\text{-}real\text{-}orthogonal\text{-}projection}(v, k, a) \wedge \\ \quad x \in \operatorname{vec}(v) \Rightarrow \\ \quad \quad a(a(x)) = a(x) \end{array}}}
\vnbprop{sip-perp-invariant}{\vnbstmt{\begin{array}{@{}l@{}} \forall v, k, o.; \forall z \in {k}^{\perp}. \\ \quad \operatorname{is\text{-}semi\text{-}inner\text{-}product\text{-}space}(v) \wedge \\ \quad k \subseteq \operatorname{vec}(v) \wedge \\ \quad o \in (\operatorname{vec}(v) \to \operatorname{vec}(v)) \wedge \\ \quad \operatorname{is\text{-}antisymmetric\text{-}op}(v, o) \wedge \\ \quad \forall \mathit{sjk} \in k.\; (o(\mathit{sjk}) \in k) \Rightarrow \\ \quad \quad o(z) \in {k}^{\perp} \end{array}}}
\vnbprop{real-projection-commutes}{\vnbstmt{\begin{array}{@{}l@{}} \forall v, k, a, o.; \forall x \in \operatorname{vec}(v). \\ \quad \operatorname{is\text{-}pre\text{-}hilbert\text{-}space}(v) \wedge \\ \quad \operatorname{is\text{-}submodule}(v, k) \wedge \\ \quad \operatorname{is\text{-}real\text{-}orthogonal\text{-}projection}(v, k, a) \wedge \\ \quad \operatorname{is\text{-}real\text{-}linear\text{-}op}(v, o) \wedge \\ \quad \operatorname{is\text{-}antisymmetric\text{-}op}(v, o) \wedge \\ \quad \forall \mathit{sjk} \in k.\; (o(\mathit{sjk}) \in k) \Rightarrow \\ \quad \quad a(o(x)) = o(a(x)) \end{array}}}
\vnbprop{real-convex-closest-unique}{\vnbstmt{\begin{array}{@{}l@{}} \forall v, s, x, a, b. \\ \quad \operatorname{is\text{-}pre\text{-}hilbert\text{-}space}(v) \wedge \\ \quad \operatorname{is\text{-}convex\text{-}subset}(v, s) \wedge \\ \quad x \in \operatorname{vec}(v) \wedge \\ \quad a \in s \wedge \\ \quad b \in s \wedge \\ \quad \forall c \in s. \\ \quad \quad \langle\cdot\mid\cdot\rangle_{v}(\operatorname{vadd}(v)(a, \operatorname{vneg}(v)(x)), \operatorname{vadd}(v)(a, \operatorname{vneg}(v)(x))) \\ \quad \quad \leq\; \langle\cdot\mid\cdot\rangle_{v}(\operatorname{vadd}(v)(c, \operatorname{vneg}(v)(x)), \operatorname{vadd}(v)(c, \operatorname{vneg}(v)(x))) \wedge \\ \quad \forall c \in s. \\ \quad \quad \langle\cdot\mid\cdot\rangle_{v}(\operatorname{vadd}(v)(b, \operatorname{vneg}(v)(x)), \operatorname{vadd}(v)(b, \operatorname{vneg}(v)(x))) \\ \quad \quad \leq\; \langle\cdot\mid\cdot\rangle_{v}(\operatorname{vadd}(v)(c, \operatorname{vneg}(v)(x)), \operatorname{vadd}(v)(c, \operatorname{vneg}(v)(x))) \Rightarrow \\ \quad \quad a = b \end{array}}}
\vnbprop{real-convex-variational-closest}{\vnbstmt{\begin{array}{@{}l@{}} \forall v, s, x, a.; \forall d \in s. \\ \quad \operatorname{is\text{-}semi\text{-}inner\text{-}product\text{-}space}(v) \wedge \\ \quad s \subseteq \operatorname{vec}(v) \wedge \\ \quad x \in \operatorname{vec}(v) \wedge \\ \quad a \in s \wedge \\ \quad \forall c \in s. \\ \quad \quad \langle\cdot\mid\cdot\rangle_{v}(\operatorname{vadd}(v)(x, \operatorname{vneg}(v)(a)), \operatorname{vadd}(v)(c, \operatorname{vneg}(v)(a))) \leq 0 \Rightarrow \\ \quad \quad \langle\cdot\mid\cdot\rangle_{v}(\operatorname{vadd}(v)(a, \operatorname{vneg}(v)(x)), \operatorname{vadd}(v)(a, \operatorname{vneg}(v)(x))) \\ \quad \quad \leq\; \langle\cdot\mid\cdot\rangle_{v}(\operatorname{vadd}(v)(d, \operatorname{vneg}(v)(x)), \operatorname{vadd}(v)(d, \operatorname{vneg}(v)(x))) \end{array}}}
\vnbprop{cxw-dist-sq}{\vnbstmt{\begin{array}{@{}l@{}} \forall v, x, \mathit{qx}. \\ \quad (\operatorname{is\text{-}pre\text{-}hilbert\text{-}space}(v) \wedge ((x \in (\operatorname{vec}(v) \times \operatorname{vec}(v))) \wedge (\mathit{qx} \in (\operatorname{vec}(v) \times \operatorname{vec}(v))))) \Rightarrow  \\ \quad \quad \langle\cdot,\cdot\rangle_{\operatorname{csip\text{-}cips}(\operatorname{complexify}(v))}(\operatorname{vadd}(\operatorname{csip\text{-}cips}(\operatorname{complexify}(v)))(x, \operatorname{vneg}(\operatorname{csip\text{-}cips}(\operatorname{complexify}(v)))(\mathit{qx})), \\ \quad \quad \quad \operatorname{vadd}(\operatorname{csip\text{-}cips}(\operatorname{complexify}(v)))(x, \operatorname{vneg}(\operatorname{csip\text{-}cips}(\operatorname{complexify}(v)))(\mathit{qx}))) \\ \quad \quad =\; \langle\cdot\mid\cdot\rangle_{v}(\operatorname{vadd}(v)(\operatorname{nth}(1, x), \operatorname{vneg}(v)(\operatorname{nth}(1, \mathit{qx}))), \operatorname{vadd}(v)(\operatorname{nth}(1, x), \operatorname{vneg}(v)(\operatorname{nth}(1, \mathit{qx})))) + \langle\cdot\mid\cdot\rangle_{v}(\operatorname{vadd}(v)(\operatorname{nth}(2, x), \operatorname{vneg}(v)(\operatorname{nth}(2, \mathit{qx}))), \operatorname{vadd}(v)(\operatorname{nth}(2, x), \operatorname{vneg}(v)(\operatorname{nth}(2, \mathit{qx})))) \end{array}}}
\vnbprop{cxw-convex-midpoint-closed}{\vnbstmt{\begin{array}{@{}l@{}} \forall v, s.; \forall m_{1}, m_{2} \in (s \times s). \\ \quad (\operatorname{is\text{-}pre\text{-}hilbert\text{-}space}(v) \wedge \operatorname{is\text{-}convex\text{-}subset}(v, s)) \Rightarrow  \\ \quad \quad \operatorname{act}(\operatorname{csip\text{-}cips}(\operatorname{complexify}(v)))(1/2, \operatorname{vadd}(\operatorname{csip\text{-}cips}(\operatorname{complexify}(v)))(m_{1}, m_{2})) \\ \quad \quad \in\; (s \times s) \end{array}}}
\vnbprop{real-convex-closest-point}{\vnbstmt{\begin{array}{@{}l@{}} \forall v, s, x. \\ \quad \operatorname{is\text{-}real\text{-}hilbert\text{-}space}(v) \wedge \\ \quad \operatorname{is\text{-}convex\text{-}subset}(v, s) \wedge \\ \quad \exists e.\; (e \in s) \wedge \\ \quad \operatorname{is\text{-}closed}(\operatorname{sip\text{-}metric\text{-}space}(v), s) \wedge \\ \quad x \in \operatorname{vec}(v) \Rightarrow \\ \quad \quad \exists a \in s. \forall b \in s. \\ \quad \quad \quad \langle\cdot\mid\cdot\rangle_{v}(\operatorname{vadd}(v)(a, \operatorname{vneg}(v)(x)), \operatorname{vadd}(v)(a, \operatorname{vneg}(v)(x))) \\ \quad \quad \quad \leq\; \langle\cdot\mid\cdot\rangle_{v}(\operatorname{vadd}(v)(b, \operatorname{vneg}(v)(x)), \operatorname{vadd}(v)(b, \operatorname{vneg}(v)(x))) \end{array}}}
\vnbprop{real-convex-variational-necessary}{\vnbstmt{\begin{array}{@{}l@{}} \forall v, s, x, a.; \forall c \in s. \\ \quad \operatorname{is\text{-}semi\text{-}inner\text{-}product\text{-}space}(v) \wedge \\ \quad \operatorname{is\text{-}convex\text{-}subset}(v, s) \wedge \\ \quad x \in \operatorname{vec}(v) \wedge \\ \quad a \in s \wedge \\ \quad \forall d \in s. \\ \quad \quad \langle\cdot\mid\cdot\rangle_{v}(\operatorname{vadd}(v)(a, \operatorname{vneg}(v)(x)), \operatorname{vadd}(v)(a, \operatorname{vneg}(v)(x))) \\ \quad \quad \leq\; \langle\cdot\mid\cdot\rangle_{v}(\operatorname{vadd}(v)(d, \operatorname{vneg}(v)(x)), \operatorname{vadd}(v)(d, \operatorname{vneg}(v)(x))) \Rightarrow \\ \quad \quad \langle\cdot\mid\cdot\rangle_{v}(\operatorname{vadd}(v)(x, \operatorname{vneg}(v)(a)), \operatorname{vadd}(v)(c, \operatorname{vneg}(v)(a))) \leq 0 \end{array}}}
\vnbprop{real-convex-closest-point-unique}{\vnbstmt{\begin{array}{@{}l@{}} \forall v, s, x. \\ \quad \operatorname{is\text{-}real\text{-}hilbert\text{-}space}(v) \wedge \\ \quad \operatorname{is\text{-}convex\text{-}subset}(v, s) \wedge \\ \quad \exists e.\; (e \in s) \wedge \\ \quad \operatorname{is\text{-}closed}(\operatorname{sip\text{-}metric\text{-}space}(v), s) \wedge \\ \quad x \in \operatorname{vec}(v) \Rightarrow \\ \quad \quad \exists a \in s. \\ \quad \quad \quad \forall c \in s. \\ \quad \quad \quad \quad \langle\cdot\mid\cdot\rangle_{v}(\operatorname{vadd}(v)(a, \operatorname{vneg}(v)(x)), \operatorname{vadd}(v)(a, \operatorname{vneg}(v)(x))) \\ \quad \quad \quad \quad \leq\; \langle\cdot\mid\cdot\rangle_{v}(\operatorname{vadd}(v)(c, \operatorname{vneg}(v)(x)), \operatorname{vadd}(v)(c, \operatorname{vneg}(v)(x))) \wedge \\ \quad \quad \quad \forall b \in s. \\ \quad \quad \quad \quad \forall c \in s. \\ \quad \quad \quad \quad \quad \langle\cdot\mid\cdot\rangle_{v}(\operatorname{vadd}(v)(b, \operatorname{vneg}(v)(x)), \operatorname{vadd}(v)(b, \operatorname{vneg}(v)(x))) \\ \quad \quad \quad \quad \quad \leq\; \langle\cdot\mid\cdot\rangle_{v}(\operatorname{vadd}(v)(c, \operatorname{vneg}(v)(x)), \operatorname{vadd}(v)(c, \operatorname{vneg}(v)(x))) \Rightarrow \\ \quad \quad \quad \quad \quad b = a \end{array}}}
\vnbprop{real-convex-closest-iff-variational}{\vnbstmt{\begin{array}{@{}l@{}} \forall v, s, x, a. \\ \quad (\operatorname{is\text{-}semi\text{-}inner\text{-}product\text{-}space}(v) \wedge (\operatorname{is\text{-}convex\text{-}subset}(v, s) \wedge ((x \in \operatorname{vec}(v)) \wedge (a \in s)))) \Rightarrow  \\ \quad \quad \forall d \in s. \\ \quad \quad \quad \langle\cdot\mid\cdot\rangle_{v}(\operatorname{vadd}(v)(a, \operatorname{vneg}(v)(x)), \operatorname{vadd}(v)(a, \operatorname{vneg}(v)(x))) \\ \quad \quad \quad \leq\; \langle\cdot\mid\cdot\rangle_{v}(\operatorname{vadd}(v)(d, \operatorname{vneg}(v)(x)), \operatorname{vadd}(v)(d, \operatorname{vneg}(v)(x))) \\ \quad \quad \Leftrightarrow\; \forall c \in s. \\ \quad \quad \quad \langle\cdot\mid\cdot\rangle_{v}(\operatorname{vadd}(v)(x, \operatorname{vneg}(v)(a)), \operatorname{vadd}(v)(c, \operatorname{vneg}(v)(a))) \leq 0 \end{array}}}

\vnbbookfile{theorem-library/sip-quotient-complete}{theorem-library/sip-quotient-complete}
\noindent{\small sip-quotient-complete.scm -- lemmas MOVED 2026-10-30 (wave-2 integration) out of theorem-library/l2-hilbert.scm into library analysis, for the library gate: cited by real-riesz (analysis); a semi-inner-product-space fact. The proofs are the originals, verbatim; the helper defines are a copy of the source's. ...}\par
\vnbprop{sip-quotient-complete}{\vnbstmt{\begin{array}{@{}l@{}} \forall v. \\ \quad \operatorname{is\text{-}semi\text{-}inner\text{-}product\text{-}space}(v) \wedge \\ \quad \forall s \in (\mathbb{N} \to \operatorname{vec}(v)). \\ \quad \quad \forall b. \\ \quad \quad \quad \operatorname{pos\text{-}rr}(b) \Rightarrow  \\ \quad \quad \quad \quad \exists c \in \mathbb{N}. \forall m, q \in \mathbb{N}. \\ \quad \quad \quad \quad \quad ((c \leq m) \wedge (c \leq q)) \Rightarrow  \\ \quad \quad \quad \quad \quad \quad \sqrt{\langle\cdot\mid\cdot\rangle_{v}(\operatorname{vadd}(v)(s(m), \operatorname{vneg}(v)(s(q))), \operatorname{vadd}(v)(s(m), \operatorname{vneg}(v)(s(q))))} \leq b \Rightarrow \\ \quad \quad \quad \exists a \in \operatorname{vec}(v). \forall b. \\ \quad \quad \quad \quad \operatorname{pos\text{-}rr}(b) \Rightarrow  \\ \quad \quad \quad \quad \quad \exists c \in \mathbb{N}. \forall n \in \mathbb{N}. \\ \quad \quad \quad \quad \quad \quad (c \leq n) \Rightarrow  \\ \quad \quad \quad \quad \quad \quad \quad \sqrt{\langle\cdot\mid\cdot\rangle_{v}(\operatorname{vadd}(v)(s(n), \operatorname{vneg}(v)(a)), \operatorname{vadd}(v)(s(n), \operatorname{vneg}(v)(a)))} \leq b \Rightarrow \\ \quad \quad \operatorname{is\text{-}real\text{-}hilbert\text{-}space}(\operatorname{sip\text{-}quotient}(v)) \end{array}}}

\vnbbookfile{structure-library/cnvs-vocab}{structure-library/cnvs-vocab}
\noindent{\small cnvs-vocab.scm -- wave of 2026-11-09/10 (agent CNVS), wired 2026-11-10 from scratchpad/cnvs-1109/cnvs-vocab.scm. Library: analysis.}\par
\vnbstructdef{complex-normed-vector-space}{A complex normed vector space $s$ is an 8-tuple whose components are listed below.\par
\begin{enumerate}[nosep]
\item $\operatorname{scal}(s) = \operatorname{nth}(1, s)$, a substructure (\texttt{substructure}): $\operatorname{is\text{-}ring}(\operatorname{scal}(s))$.
\item $\operatorname{vec}(s) = \operatorname{nth}(2, s)$, a carrier (\texttt{carriers}): $\operatorname{vec}(s) \in \mathbf{Set}$.
\item $\operatorname{vadd}(s) = \operatorname{nth}(3, s)$, an operation (\texttt{op}): $\operatorname{vadd}(s) \in ((\operatorname{vec}(s) \times \operatorname{vec}(s)) \to \operatorname{vec}(s))$.
\item $\operatorname{vzero}(s) = \operatorname{nth}(4, s)$, a constant (\texttt{constant}): $\operatorname{vzero}(s) \in \operatorname{vec}(s)$.
\item $\operatorname{vneg}(s) = \operatorname{nth}(5, s)$, an operation (\texttt{op}): $\operatorname{vneg}(s) \in (\operatorname{vec}(s) \to \operatorname{vec}(s))$.
\item $\operatorname{act}(s) = \operatorname{nth}(6, s)$, an operation (\texttt{op}): $\operatorname{act}(s) \in ((\operatorname{carr}(\operatorname{scal}(s)) \times \operatorname{vec}(s)) \to \operatorname{vec}(s))$.
\item $\operatorname{vnrm}(s) = \operatorname{nth}(7, s)$, an operation (\texttt{op}): $\operatorname{vnrm}(s) \in (\operatorname{vec}(s) \to \mathbb{R})$.
\item $J_{s} = \operatorname{nth}(8, s)$, an operation (\texttt{op}): $J_{s} \in (\operatorname{vec}(s) \to \operatorname{vec}(s))$.
\end{enumerate}
The laws (the \texttt{property} and \texttt{law} clauses):
\begin{itemize}[nosep]
\item $\operatorname{is\text{-}normed\text{-}vector\text{-}space}({[}\operatorname{scal}(s), \operatorname{vec}(s), \operatorname{vadd}(s), \operatorname{vzero}(s), \operatorname{vneg}(s), \operatorname{act}(s), \operatorname{vnrm}(s)])$
\item $\operatorname{is\text{-}real\text{-}linear\text{-}op}(s, J_{s})$
\item $\operatorname{is\text{-}square\text{-}minus\text{-}id}(s, J_{s})$
\item
\begin{equation*}
\begin{array}{@{}l@{}} \forall \operatorname{cnc\_} \in \mathbb{C}, \operatorname{cnx\_} \in \operatorname{vec}(s). \\ \quad \operatorname{vnrm}(s)(\operatorname{cact}(s)(\operatorname{cnc\_}, \operatorname{cnx\_})) = (\operatorname{magnitude}(\operatorname{cnc\_}) \cdot \operatorname{vnrm}(s)(\operatorname{cnx\_})) \end{array}
\end{equation*}

\end{itemize}
The characteristic law (the axiom \texttt{is-complex-normed-vector-space}): $\operatorname{is\text{-}complex\text{-}normed\text{-}vector\text{-}space}(s)$ holds if and only if $\operatorname{length}(s) = 8$, every component has the typing listed with it, every law holds, and $s \in \mathbf{Set}$ (the structure is a set).
}
\par\medskip\noindent\textbf{Definition} (\texttt{cnvs-metric-space}).\;$\operatorname{cnvs\text{-}metric\text{-}space}(\operatorname{cnv\_}) :=$ $\operatorname{nag\text{-}metric\text{-}space}(\operatorname{cnvs\text{-}as\text{-}normed\text{-}ag}(\operatorname{cnv\_}))$
\par\medskip\noindent\textbf{Definition} (\texttt{cmod-cnvs}).\;$\operatorname{cmod\text{-}cnvs}(m, \mathit{nf}) :=$
\begin{equation*}
{[}\operatorname{normed\text{-}field\text{-}as\text{-}commutative\text{-}ring}(\operatorname{rr\text{-}normed\text{-}field}), \operatorname{vec}(m), \operatorname{vadd}(m), \operatorname{vzero}(m), \operatorname{vneg}(m), \operatorname{cmod\text{-}ract}(m), \mathit{nf}, \operatorname{cmod\text{-}j}(m)]
\end{equation*}

\par\medskip\noindent\textbf{Definition} (\texttt{cips-cnvs}).\;$\operatorname{cips\text{-}cnvs}(w) :=$
\begin{equation*}
\operatorname{cmod\text{-}cnvs}(\operatorname{complex\text{-}inner\text{-}product\text{-}space\text{-}as\text{-}module}(w), \operatorname{cips\text{-}vnrm}(w))
\end{equation*}

\par\medskip\noindent\textbf{Definition} (\texttt{csip-cnvs}).\;$\operatorname{csip\text{-}cnvs}(v) :=$ $\operatorname{cips\text{-}cnvs}(\operatorname{csip\text{-}cips}(v))$
\par\medskip\noindent\textbf{Definition} (\texttt{cnvs-bounded-ops}).\;$\operatorname{cnvs\text{-}bounded\text{-}ops}(\operatorname{cbv\_}, \operatorname{cbw\_}) :=$
\begin{equation*}
\{\operatorname{cbt\_} \in (\operatorname{vec}(\operatorname{cbv\_}) \to \operatorname{vec}(\operatorname{cbw\_})) : \operatorname{is\text{-}cnvs\text{-}bounded\text{-}op}(\operatorname{cbv\_}, \operatorname{cbw\_}, \operatorname{cbt\_})\}
\end{equation*}

\par\medskip\noindent\textbf{Definition} (\texttt{cnvs-opnorm}).\;$\operatorname{cnvs\text{-}opnorm}(\operatorname{cbv\_}, \operatorname{cbw\_}, \operatorname{cbt\_}) :=$
\begin{equation*}
\operatorname{sup}(\{\operatorname{cba\_} \in \mathbb{R} : ((\operatorname{cba\_} = 0) \vee \exists \operatorname{cbx\_} \in \operatorname{vec}(\operatorname{cbv\_}).\; (\neg (\operatorname{vnrm}(\operatorname{cbv\_})(\operatorname{cbx\_}) = 0) \wedge (\operatorname{cba\_} = (\operatorname{vnrm}(\operatorname{cbw\_})(\operatorname{cbt\_}(\operatorname{cbx\_})) \cdot \operatorname{recip}(\operatorname{vnrm}(\operatorname{cbv\_})(\operatorname{cbx\_}))))))\})
\end{equation*}

\par\medskip\noindent\textbf{Definition} (\texttt{cbop-normed-ag}).\;$\operatorname{cbop\text{-}normed\text{-}ag}(\operatorname{cbv\_}, \operatorname{cbw\_}) :=$
\begin{equation*}
{[}\operatorname{cnvs\text{-}bounded\text{-}ops}(\operatorname{cbv\_}, \operatorname{cbw\_}), \operatorname{cbop\text{-}vadd}(\operatorname{cbv\_}, \operatorname{cbw\_}), \operatorname{cbop\text{-}zero}(\operatorname{cbv\_}, \operatorname{cbw\_}), \operatorname{cbop\text{-}vinv}(\operatorname{cbv\_}, \operatorname{cbw\_}), \operatorname{cbop\text{-}vnrm}(\operatorname{cbv\_}, \operatorname{cbw\_})]
\end{equation*}

\par\medskip\noindent\textbf{Definition} (\texttt{cmod-ract}).\;$\operatorname{cmod\text{-}ract}(m) :=$
\begin{equation*}
(\lambda {[}\operatorname{cnr\_}, \operatorname{cnx\_}] \in (\mathbb{R} \times \operatorname{vec}(m)).\; \operatorname{act}(m)(\operatorname{cnr\_}, \operatorname{cnx\_}))
\end{equation*}

\par\medskip\noindent\textbf{Definition} (\texttt{cmod-j}).\;$\operatorname{cmod\text{-}j}(m) :=$ $(\lambda \operatorname{cnx\_} \in \operatorname{vec}(m).\; \operatorname{act}(m)(+i, \operatorname{cnx\_}))$
\par\medskip\noindent\textbf{Definition} (\texttt{is-cmod-norm}).\;For all $\operatorname{cnm\_}$ and $\operatorname{cnf\_}$, 
\begin{equation*}
\begin{array}{@{}l@{}} \operatorname{is\text{-}cmod\text{-}norm}(\operatorname{cnm\_}, \operatorname{cnf\_}) \\ \Leftrightarrow\; \operatorname{cnf\_} \in (\operatorname{vec}(\operatorname{cnm\_}) \to \mathbb{R}) \wedge \\ \forall \operatorname{cnx\_} \in \operatorname{vec}(\operatorname{cnm\_}).\; (0 \leq \operatorname{cnf\_}(\operatorname{cnx\_})) \wedge \\ \forall \operatorname{cnx\_} \in \operatorname{vec}(\operatorname{cnm\_}).\; ((\operatorname{cnf\_}(\operatorname{cnx\_}) = 0) \Leftrightarrow (\operatorname{cnx\_} = \operatorname{vzero}(\operatorname{cnm\_}))) \wedge \\ \forall \operatorname{cnc\_} \in \mathbb{C}, \operatorname{cnx\_} \in \operatorname{vec}(\operatorname{cnm\_}). \\ \quad \operatorname{cnf\_}(\operatorname{act}(\operatorname{cnm\_})(\operatorname{cnc\_}, \operatorname{cnx\_})) = \cdot(\operatorname{magnitude}(\operatorname{cnc\_}), \operatorname{cnf\_}(\operatorname{cnx\_})) \wedge \\ \forall \operatorname{cnx\_}, \operatorname{cny\_} \in \operatorname{vec}(\operatorname{cnm\_}). \\ \quad \operatorname{cnf\_}(\operatorname{vadd}(\operatorname{cnm\_})(\operatorname{cnx\_}, \operatorname{cny\_})) \leq +(\operatorname{cnf\_}(\operatorname{cnx\_}), \operatorname{cnf\_}(\operatorname{cny\_})) \end{array}
\end{equation*}

\par\medskip\noindent\textbf{Definition} (\texttt{cmod-nvs}).\;$\operatorname{cmod\text{-}nvs}(m, \mathit{nf}) :=$
\begin{equation*}
{[}\operatorname{normed\text{-}field\text{-}as\text{-}commutative\text{-}ring}(\operatorname{rr\text{-}normed\text{-}field}), \operatorname{vec}(m), \operatorname{vadd}(m), \operatorname{vzero}(m), \operatorname{vneg}(m), \operatorname{cmod\text{-}ract}(m), \mathit{nf}]
\end{equation*}

\par\medskip\noindent\textbf{Definition} (\texttt{cips-vnrm}).\;$\operatorname{cips\text{-}vnrm}(w) :=$ $(\lambda \operatorname{cnx\_} \in \operatorname{vec}(w).\; \lVert \operatorname{cnx\_} \rVert)$
\par\medskip\noindent\textbf{Definition} (\texttt{is-cnvs-linear-op}).\;For all $\operatorname{cbv\_}$, $\operatorname{cbw\_}$ and $\operatorname{cbt\_}$, 
\begin{equation*}
\begin{array}{@{}l@{}} \operatorname{is\text{-}cnvs\text{-}linear\text{-}op}(\operatorname{cbv\_}, \operatorname{cbw\_}, \operatorname{cbt\_}) \\ \Leftrightarrow\; \operatorname{cbt\_} \in (\operatorname{vec}(\operatorname{cbv\_}) \to \operatorname{vec}(\operatorname{cbw\_})) \wedge \\ \forall \operatorname{cbx\_}, \operatorname{cby\_} \in \operatorname{vec}(\operatorname{cbv\_}). \\ \quad \operatorname{cbt\_}(\operatorname{vadd}(\operatorname{cbv\_})(\operatorname{cbx\_}, \operatorname{cby\_})) = \operatorname{vadd}(\operatorname{cbw\_})(\operatorname{cbt\_}(\operatorname{cbx\_}), \operatorname{cbt\_}(\operatorname{cby\_})) \wedge \\ \forall \operatorname{cbc\_} \in \mathbb{C}, \operatorname{cbx\_} \in \operatorname{vec}(\operatorname{cbv\_}). \\ \quad \operatorname{cbt\_}(\operatorname{cact}(\operatorname{cbv\_})(\operatorname{cbc\_}, \operatorname{cbx\_})) = \operatorname{cact}(\operatorname{cbw\_})(\operatorname{cbc\_}, \operatorname{cbt\_}(\operatorname{cbx\_})) \end{array}
\end{equation*}

\par\medskip\noindent\textbf{Definition} (\texttt{is-cnvs-bounded-op}).\;For all $\operatorname{cbv\_}$, $\operatorname{cbw\_}$ and $\operatorname{cbt\_}$, 
\begin{equation*}
\begin{array}{@{}l@{}} \operatorname{is\text{-}cnvs\text{-}bounded\text{-}op}(\operatorname{cbv\_}, \operatorname{cbw\_}, \operatorname{cbt\_}) \\ \Leftrightarrow\; \operatorname{is\text{-}cnvs\text{-}linear\text{-}op}(\operatorname{cbv\_}, \operatorname{cbw\_}, \operatorname{cbt\_}) \wedge \\ \exists \operatorname{cbm\_} \in \mathbb{R}. \\ \quad ((0 \leq \operatorname{cbm\_}) \wedge \forall \operatorname{cbx\_} \in \operatorname{vec}(\operatorname{cbv\_}).\; (\operatorname{vnrm}(\operatorname{cbw\_})(\operatorname{cbt\_}(\operatorname{cbx\_})) \leq \cdot(\operatorname{cbm\_}, \operatorname{vnrm}(\operatorname{cbv\_})(\operatorname{cbx\_})))) \end{array}
\end{equation*}

\par\medskip\noindent\textbf{Definition} (\texttt{cbop-sum}).\;$\operatorname{cbop\text{-}sum}(\operatorname{cbv\_}, \operatorname{cbw\_}, \operatorname{cbs\_}, \operatorname{cbt\_}) :=$
\begin{equation*}
(\lambda \operatorname{cbx\_} \in \operatorname{vec}(\operatorname{cbv\_}).\; \operatorname{vadd}(\operatorname{cbw\_})(\operatorname{cbs\_}(\operatorname{cbx\_}), \operatorname{cbt\_}(\operatorname{cbx\_})))
\end{equation*}

\par\medskip\noindent\textbf{Definition} (\texttt{cbop-zero}).\;$\operatorname{cbop\text{-}zero}(\operatorname{cbv\_}, \operatorname{cbw\_}) :=$ $(\lambda \operatorname{cbx\_} \in \operatorname{vec}(\operatorname{cbv\_}).\; \operatorname{vzero}(\operatorname{cbw\_}))$
\par\medskip\noindent\textbf{Definition} (\texttt{cbop-neg}).\;$\operatorname{cbop\text{-}neg}(\operatorname{cbv\_}, \operatorname{cbw\_}, \operatorname{cbt\_}) :=$ $(\lambda \operatorname{cbx\_} \in \operatorname{vec}(\operatorname{cbv\_}).\; \operatorname{vneg}(\operatorname{cbw\_})(\operatorname{cbt\_}(\operatorname{cbx\_})))$
\par\medskip\noindent\textbf{Definition} (\texttt{cbop-scale}).\;$\operatorname{cbop\text{-}scale}(\operatorname{cbv\_}, \operatorname{cbw\_}, \operatorname{cbc\_}, \operatorname{cbt\_}) :=$ $(\lambda \operatorname{cbx\_} \in \operatorname{vec}(\operatorname{cbv\_}).\; \operatorname{cact}(\operatorname{cbw\_})(\operatorname{cbc\_}, \operatorname{cbt\_}(\operatorname{cbx\_})))$
\par\medskip\noindent\textbf{Definition} (\texttt{cbop-vadd}).\;$\operatorname{cbop\text{-}vadd}(\operatorname{cbv\_}, \operatorname{cbw\_}) :=$
\begin{equation*}
\begin{array}{@{}l@{}} (\lambda {[}\operatorname{cbs\_}, \operatorname{cbt\_}] \in (\operatorname{cnvs\text{-}bounded\text{-}ops}(\operatorname{cbv\_}, \operatorname{cbw\_}) \times \operatorname{cnvs\text{-}bounded\text{-}ops}(\operatorname{cbv\_}, \operatorname{cbw\_})).\; \\ \quad \operatorname{cbop\text{-}sum}(\operatorname{cbv\_}, \operatorname{cbw\_}, \operatorname{cbs\_}, \operatorname{cbt\_})) \end{array}
\end{equation*}

\par\medskip\noindent\textbf{Definition} (\texttt{cbop-vinv}).\;$\operatorname{cbop\text{-}vinv}(\operatorname{cbv\_}, \operatorname{cbw\_}) :=$
\begin{equation*}
(\lambda \operatorname{cbt\_} \in \operatorname{cnvs\text{-}bounded\text{-}ops}(\operatorname{cbv\_}, \operatorname{cbw\_}).\; \operatorname{cbop\text{-}neg}(\operatorname{cbv\_}, \operatorname{cbw\_}, \operatorname{cbt\_}))
\end{equation*}

\par\medskip\noindent\textbf{Definition} (\texttt{cbop-vnrm}).\;$\operatorname{cbop\text{-}vnrm}(\operatorname{cbv\_}, \operatorname{cbw\_}) :=$
\begin{equation*}
(\lambda \operatorname{cbt\_} \in \operatorname{cnvs\text{-}bounded\text{-}ops}(\operatorname{cbv\_}, \operatorname{cbw\_}).\; \operatorname{cnvs\text{-}opnorm}(\operatorname{cbv\_}, \operatorname{cbw\_}, \operatorname{cbt\_}))
\end{equation*}

\par\medskip\noindent\textbf{Definition} (\texttt{hom-complex-normed-vector-space}).\;$\operatorname{hom\text{-}complex\text{-}normed\text{-}vector\text{-}space}(a, b) :=$
\begin{equation*}
\{\operatorname{homf\_} \in (\operatorname{vec}(a) \to \operatorname{vec}(b)) : \operatorname{is\text{-}hom\text{-}complex\text{-}normed\text{-}vector\text{-}space}(a, b, \operatorname{homf\_})\}
\end{equation*}

\par\medskip\noindent\textbf{Definition} (\texttt{hom-complex-banach-space}).\;$\operatorname{hom\text{-}complex\text{-}banach\text{-}space}(a, b) :=$
\begin{equation*}
\{\operatorname{homf\_} \in (\operatorname{vec}(a) \to \operatorname{vec}(b)) : \operatorname{is\text{-}hom\text{-}complex\text{-}banach\text{-}space}(a, b, \operatorname{homf\_})\}
\end{equation*}

\vnbdef{complex-normed-vector-space-class}{\vnbstmt{\begin{array}{@{}l@{}} \forall s. \\ \quad (s \in \operatorname{complex\text{-}normed\text{-}vector\text{-}space}) \Leftrightarrow \operatorname{is\text{-}complex\text{-}normed\text{-}vector\text{-}space}(s) \end{array}}}
\vnbdef{is-hom-complex-normed-vector-space-def}{\vnbstmt{\begin{array}{@{}l@{}} \forall a, b, f. \\ \quad \operatorname{is\text{-}hom\text{-}complex\text{-}normed\text{-}vector\text{-}space}(a, b, f) \\ \quad \Leftrightarrow\; \operatorname{is\text{-}complex\text{-}normed\text{-}vector\text{-}space}(a) \wedge \\ \quad \operatorname{is\text{-}complex\text{-}normed\text{-}vector\text{-}space}(b) \wedge \\ \quad f \in (\operatorname{vec}(a) \to \operatorname{vec}(b)) \wedge \\ \quad \operatorname{scal}(a) = \operatorname{scal}(b) \wedge \\ \quad \forall x_{1}, x_{2} \in \operatorname{vec}(a). \\ \quad \quad f(\operatorname{vadd}(a)(x_{1}, x_{2})) = \operatorname{vadd}(b)(f(x_{1}), f(x_{2})) \wedge \\ \quad f(\operatorname{vzero}(a)) = \operatorname{vzero}(b) \wedge \\ \quad \forall x_{1} \in \operatorname{vec}(a).\; (f(\operatorname{vneg}(a)(x_{1})) = \operatorname{vneg}(b)(f(x_{1}))) \wedge \\ \quad \forall x_{1} \in \operatorname{carr}(\operatorname{scal}(a)), x_{2} \in \operatorname{vec}(a). \\ \quad \quad f(\operatorname{act}(a)(x_{1}, x_{2})) = \operatorname{act}(b)(x_{1}, f(x_{2})) \wedge \\ \quad \forall x_{1} \in \operatorname{vec}(a).\; (\operatorname{vnrm}(a)(x_{1}) = \operatorname{vnrm}(b)(f(x_{1}))) \wedge \\ \quad \forall x_{1} \in \operatorname{vec}(a).\; (f(J_{a}(x_{1})) = J_{b}(f(x_{1}))) \end{array}}}
\vnbdef{is-hom-complex-banach-space-def}{\vnbstmt{\begin{array}{@{}l@{}} \forall a, b, f. \\ \quad \operatorname{is\text{-}hom\text{-}complex\text{-}banach\text{-}space}(a, b, f) \\ \quad \Leftrightarrow\; \operatorname{is\text{-}complex\text{-}banach\text{-}space}(a) \wedge \\ \quad \operatorname{is\text{-}complex\text{-}banach\text{-}space}(b) \wedge \\ \quad \operatorname{is\text{-}hom\text{-}complex\text{-}normed\text{-}vector\text{-}space}(a, b, f) \end{array}}}
\vnbdef{complex-banach-space-class}{\vnbstmt{\forall s.\; ((s \in \operatorname{complex\text{-}banach\text{-}space}) \Leftrightarrow \operatorname{is\text{-}complex\text{-}banach\text{-}space}(s))}}
\vnbdef{is-complex-banach-space-def}{\vnbstmt{\begin{array}{@{}l@{}} \forall s. \\ \quad \operatorname{is\text{-}complex\text{-}banach\text{-}space}(s) \\ \quad \Leftrightarrow\; \operatorname{is\text{-}complex\text{-}normed\text{-}vector\text{-}space}(s) \wedge \operatorname{is\text{-}complete}(\operatorname{cnvs\text{-}metric\text{-}space}(s)) \end{array}}}
\vnbprop{is-complex-normed-vector-space-is-set}{\vnbstmt{\forall s.\; (\operatorname{is\text{-}complex\text{-}normed\text{-}vector\text{-}space}(s) \Rightarrow (s \in \mathbf{Set}))}}
\vnbprop{is-complex-banach-space-is-set}{\vnbstmt{\forall s.\; (\operatorname{is\text{-}complex\text{-}banach\text{-}space}(s) \Rightarrow (s \in \mathbf{Set}))}}

\vnbbookfile{structure-library/cnvs-comp-vocab}{structure-library/cnvs-comp-vocab}
\noindent{\small cnvs-comp-vocab.scm -- wave of 2026-11-09/10 (agent CNVS2), wired 2026-11-10 from scratchpad/cnvs2-1109/cnvs2-comp-vocab.scm. Library: analysis.}\par
\par\medskip\noindent\textbf{Definition} (\texttt{cbop-comp}).\;${\operatorname{cqs\_}}{\operatorname{cqt\_}} :=$ $(\lambda \operatorname{cqx\_} \in \operatorname{vec}(\operatorname{cqu\_}).\; \operatorname{cqs\_}(\operatorname{cqt\_}(\operatorname{cqx\_})))$
\par\medskip\noindent\textbf{Definition} (\texttt{cbop-id}).\;$1 :=$ $(\lambda \operatorname{cqx\_} \in \operatorname{vec}(\operatorname{cqu\_}).\; \operatorname{cqx\_})$

\vnbbookfile{theorem-library/cnvs-view-typing}{theorem-library/cnvs-view-typing}
\noindent{\small theorem-library/cnvs-view-typing.scm -- THE FLUSH POINT FOR THE COMPLEX NORMED VECTOR SPACE VIEWS (D8c, 2026-11-16).}\par
\vnbprop{cnvs-as-nvs-is-normed-vector-space}{\vnbstmt{\begin{array}{@{}l@{}} \forall r. \\ \quad \operatorname{is\text{-}complex\text{-}normed\text{-}vector\text{-}space}(r) \Rightarrow \operatorname{is\text{-}normed\text{-}vector\text{-}space}(\operatorname{cnvs\text{-}as\text{-}nvs}(r)) \end{array}}}
\vnbprop{cnvs-as-module-is-module}{\vnbstmt{\begin{array}{@{}l@{}} \forall r. \\ \quad \operatorname{is\text{-}complex\text{-}normed\text{-}vector\text{-}space}(r) \Rightarrow \operatorname{is\text{-}module}(\operatorname{cnvs\text{-}as\text{-}module}(r)) \end{array}}}
\vnbprop{cnvs-as-normed-ag-is-normed-ag}{\vnbstmt{\begin{array}{@{}l@{}} \forall r. \\ \quad \operatorname{is\text{-}complex\text{-}normed\text{-}vector\text{-}space}(r) \Rightarrow \operatorname{is\text{-}normed\text{-}ag}(\operatorname{cnvs\text{-}as\text{-}normed\text{-}ag}(r)) \end{array}}}

\vnbbookfile{theorem-library/real-riesz}{theorem-library/real-riesz}
\noindent{\small real-riesz.scm -- WAVE 2 (2026-10-27..30), wired 2026-10-30 from scratchpad/md/me-rz.scm. Library: analysis.}\par
\vnbprop{real-riesz-representation}{\vnbstmt{\begin{array}{@{}l@{}} \forall v, f. \\ \quad \operatorname{is\text{-}real\text{-}hilbert\text{-}space}(v) \wedge \\ \quad f \in (\operatorname{vec}(v) \to \mathbb{R}) \wedge \\ \quad \forall x, y \in \operatorname{vec}(v).\; (f(\operatorname{vadd}(v)(x, y)) = (f(x) + f(y))) \wedge \\ \quad \forall r \in \mathbb{R}, x \in \operatorname{vec}(v).\; (f(\operatorname{act}(v)(r, x)) = (r \cdot f(x))) \wedge \\ \quad \operatorname{is\text{-}continuous}(\operatorname{sip\text{-}metric\text{-}space}(v), \operatorname{rr\text{-}ms}, f) \Rightarrow \\ \quad \quad \exists e \in \operatorname{vec}(v).\; \forall \mathit{rzx} \in \operatorname{vec}(v).\; (f(\mathit{rzx}) = \langle\cdot\mid\cdot\rangle_{v}(\mathit{rzx}, e)) \end{array}}}
\vnbprop{sip-riesz-representation}{\vnbstmt{\begin{array}{@{}l@{}} \forall v, f, b. \\ \quad \operatorname{is\text{-}semi\text{-}inner\text{-}product\text{-}space}(v) \wedge \\ \quad \forall s \in (\mathbb{N} \to \operatorname{vec}(v)). \\ \quad \quad \forall c. \\ \quad \quad \quad \operatorname{pos\text{-}rr}(c) \Rightarrow  \\ \quad \quad \quad \quad \exists d \in \mathbb{N}. \forall m, q \in \mathbb{N}. \\ \quad \quad \quad \quad \quad ((d \leq m) \wedge (d \leq q)) \Rightarrow  \\ \quad \quad \quad \quad \quad \quad \sqrt{\langle\cdot\mid\cdot\rangle_{v}(\operatorname{vadd}(v)(s(m), \operatorname{vneg}(v)(s(q))), \operatorname{vadd}(v)(s(m), \operatorname{vneg}(v)(s(q))))} \leq c \Rightarrow \\ \quad \quad \quad \exists a \in \operatorname{vec}(v). \forall c. \\ \quad \quad \quad \quad \operatorname{pos\text{-}rr}(c) \Rightarrow  \\ \quad \quad \quad \quad \quad \exists d \in \mathbb{N}. \forall n \in \mathbb{N}. \\ \quad \quad \quad \quad \quad \quad (d \leq n) \Rightarrow  \\ \quad \quad \quad \quad \quad \quad \quad \sqrt{\langle\cdot\mid\cdot\rangle_{v}(\operatorname{vadd}(v)(s(n), \operatorname{vneg}(v)(a)), \operatorname{vadd}(v)(s(n), \operatorname{vneg}(v)(a)))} \leq c \wedge \\ \quad f \in (\operatorname{vec}(v) \to \mathbb{R}) \wedge \\ \quad \forall x, y \in \operatorname{vec}(v).\; (f(\operatorname{vadd}(v)(x, y)) = (f(x) + f(y))) \wedge \\ \quad \forall r \in \mathbb{R}, x \in \operatorname{vec}(v).\; (f(\operatorname{act}(v)(r, x)) = (r \cdot f(x))) \wedge \\ \quad b \in \mathbb{R} \wedge \\ \quad 0 \leq b \wedge \\ \quad \forall x \in \operatorname{vec}(v).\; (\lvert f(x) \rvert \leq (b \cdot \sqrt{\langle\cdot\mid\cdot\rangle_{v}(x, x)})) \Rightarrow \\ \quad \quad \exists e \in \operatorname{vec}(v).\; \forall \mathit{rzx} \in \operatorname{vec}(v).\; (f(\mathit{rzx}) = \langle\cdot\mid\cdot\rangle_{v}(\mathit{rzx}, e)) \end{array}}}

\vnbbookfile{theorem-library/series-domination}{theorem-library/series-domination}
\noindent{\small series-domination.scm -- SHIFTING AND EVENTUAL DOMINATION FOR REAL SERIES. Batch 26-A, 2026-09-24.  The two lemmas the batch-25 report named as standing in the way of Proposition 2.10 (theorem-library/cc-power-series-deriv.scm):}\par
\vnbprop{series-shift-partial-sum}{\vnbstmt{\begin{array}{@{}l@{}} \forall k \in \mathbb{N}, f, g \in (\mathbb{N} \to \mathbb{R}). \\ \quad \forall j \in \mathbb{N}.\; (g(j) = f((j + 1))) \Rightarrow  \\ \quad \quad \operatorname{series\text{-}partial\text{-}sum}(f, (k + 1)) = (f(0) + \operatorname{series\text{-}partial\text{-}sum}(g, k)) \end{array}}}
\vnbprop{series-shift-converges-iff}{\vnbstmt{\begin{array}{@{}l@{}} \forall f, g \in (\mathbb{N} \to \mathbb{R}). \\ \quad (\forall j \in \mathbb{N}.\; (0 \leq f(j)) \wedge \forall j \in \mathbb{N}.\; (g(j) = f((j + 1)))) \Rightarrow  \\ \quad \quad \operatorname{series\text{-}converges}(f) \Leftrightarrow \operatorname{series\text{-}converges}(g) \end{array}}}
\vnbprop{series-dominated-partial-bound}{\vnbstmt{\begin{array}{@{}l@{}} \forall k \in \mathbb{N}, f, g \in (\mathbb{N} \to \mathbb{R}), n \in \mathbb{N}. \\ \quad \forall j \in \mathbb{N}.\; (0 \leq f(j)) \wedge \\ \quad \forall j \in \mathbb{N}.\; (0 \leq g(j)) \wedge \\ \quad \forall j \in \mathbb{N}.\; ((n \leq j) \Rightarrow (f(j) \leq g(j))) \Rightarrow \\ \quad \quad \operatorname{series\text{-}partial\text{-}sum}(f, k) \leq (\operatorname{series\text{-}partial\text{-}sum}(f, n) + \operatorname{series\text{-}partial\text{-}sum}(g, k)) \end{array}}}
\vnbprop{series-eventually-dominated-converges}{\vnbstmt{\begin{array}{@{}l@{}} \forall f, g \in (\mathbb{N} \to \mathbb{R}), n \in \mathbb{N}. \\ \quad \forall j \in \mathbb{N}.\; (0 \leq f(j)) \wedge \\ \quad \forall j \in \mathbb{N}.\; (0 \leq g(j)) \wedge \\ \quad \forall j \in \mathbb{N}.\; ((n \leq j) \Rightarrow (f(j) \leq g(j))) \wedge \\ \quad \operatorname{series\text{-}converges}(g) \Rightarrow \\ \quad \quad \operatorname{series\text{-}converges}(f) \end{array}}}

\vnbbookfile{theorem-library/regulated-has-primitive}{theorem-library/regulated-has-primitive}
\noindent{\small regulated-has-primitive.scm -- EVERY REGULATED FUNCTION HAS A PRIMITIVE on [a, b] (batch 22, stage 2: the assembly).}\par
\vnbprop{primitive-shift-const}{\vnbstmt{\begin{array}{@{}l@{}} \forall g, f, a, b.; \forall c \in \mathbb{R}. \\ \quad \operatorname{is\text{-}primitive}(g, f, a, b) \Rightarrow  \\ \quad \quad \operatorname{is\text{-}primitive}((\lambda y \in \operatorname{ccint}(a, b).\; (g(y) - c)), f, a, b) \end{array}}}
\vnbprop{primitive-family-normalised-choice}{\vnbstmt{\begin{array}{@{}l@{}} \forall a, b \in \mathbb{R}, f. \\ \quad \forall k \in \mathbb{N}.\; \exists g.\; \operatorname{is\text{-}primitive}(g, f(k), a, b) \Rightarrow  \\ \quad \quad \exists \mathit{rhpgf} \in (\mathbb{N} \to (\operatorname{ccint}(a, b) \to \mathbb{R})). \forall k \in \mathbb{N}. \\ \quad \quad \quad \operatorname{is\text{-}primitive}(\operatorname{rhpgf}(k), f(k), a, b) \wedge (\operatorname{rhpgf}(k)(a) = 0) \end{array}}}
\vnbprop{regulated-on-has-primitive}{\vnbstmt{\begin{array}{@{}l@{}} \forall a, b \in \mathbb{R}, f. \\ \quad \operatorname{is\text{-}regulated\text{-}on}(f, a, b) \Rightarrow \exists g.\; \operatorname{is\text{-}primitive}(g, f, a, b) \end{array}}}

\vnbbookfile{theorem-library/calint-deriv-limit}{theorem-library/calint-deriv-limit}
\noindent{\small calint-deriv-limit.scm -- WAVE 2 (2026-10-27..30), wired 2026-10-30 from scratchpad/calint/calint-deriv-limit.scm. Library: analysis.}\par
\vnbprop{continuous-at-subspace-abs}{\vnbstmt{\begin{array}{@{}l@{}} \forall w, p, x, s. \\ \quad ((w \subseteq \mathbb{R}) \wedge (\operatorname{is\text{-}continuous\text{-}at}(\operatorname{subspace\text{-}ms}(\operatorname{rr\text{-}ms}, w), \operatorname{rr\text{-}ms}, p, x) \wedge \operatorname{pos\text{-}rr}(s))) \Rightarrow  \\ \quad \quad \exists d. \\ \quad \quad \quad \operatorname{pos\text{-}rr}(d) \wedge \forall y \in w.\; ((\lvert (x - y) \rvert \leq d) \Rightarrow (\lvert (p(x) - p(y)) \rvert \leq s)) \end{array}}}
\vnbprop{continuous-on-abs-fwd}{\vnbstmt{\begin{array}{@{}l@{}} \forall w, p.; \forall x \in w, s. \\ \quad ((w \subseteq \mathbb{R}) \wedge (\operatorname{is\text{-}continuous\text{-}on}(p, w) \wedge \operatorname{pos\text{-}rr}(s))) \Rightarrow  \\ \quad \quad \exists d. \\ \quad \quad \quad \operatorname{pos\text{-}rr}(d) \wedge \forall y \in w.\; ((\lvert (x - y) \rvert \leq d) \Rightarrow (\lvert (p(x) - p(y)) \rvert \leq s)) \end{array}}}
\vnbprop{continuous-on-abs-bwd}{\vnbstmt{\begin{array}{@{}l@{}} \forall w, p. \\ \quad w \subseteq \mathbb{R} \wedge \\ \quad p \in (w \to \mathbb{R}) \wedge \\ \quad \forall x \in w, s. \\ \quad \quad \operatorname{pos\text{-}rr}(s) \Rightarrow  \\ \quad \quad \quad \exists d. \\ \quad \quad \quad \quad \operatorname{pos\text{-}rr}(d) \wedge \forall y \in w.\; ((\lvert (x - y) \rvert \leq d) \Rightarrow (\lvert (p(x) - p(y)) \rvert \leq s)) \Rightarrow \\ \quad \quad \operatorname{is\text{-}continuous\text{-}on}(p, w) \end{array}}}
\vnbprop{deriv-within-local-bound}{\vnbstmt{\begin{array}{@{}l@{}} \forall w, g, x, l, s. \\ \quad ((w \subseteq \mathbb{R}) \wedge (\operatorname{has\text{-}deriv\text{-}at\text{-}within}(g, w, x, l) \wedge \operatorname{pos\text{-}rr}(s))) \Rightarrow  \\ \quad \quad \exists d. \\ \quad \quad \quad \operatorname{pos\text{-}rr}(d) \wedge \forall y \in w.\; ((\lvert (x - y) \rvert \leq d) \Rightarrow (\lvert (g(x) - g(y)) \rvert \leq s)) \end{array}}}
\vnbprop{right-deriv-local-bound}{\vnbstmt{\begin{array}{@{}l@{}} \forall g, x, l, s. \\ \quad (\operatorname{has\text{-}right\text{-}deriv\text{-}at}(g, x, l) \wedge \operatorname{pos\text{-}rr}(s)) \Rightarrow  \\ \quad \quad \exists d. \\ \quad \quad \quad (\operatorname{pos\text{-}rr}(d) \wedge \forall y \in \mathbb{R}.\; (((x \leq y) \wedge (\lvert (x - y) \rvert \leq d)) \Rightarrow (\lvert (g(x) - g(y)) \rvert \leq s))) \end{array}}}
\vnbprop{left-deriv-local-bound}{\vnbstmt{\begin{array}{@{}l@{}} \forall g, x, l, s. \\ \quad (\operatorname{has\text{-}left\text{-}deriv\text{-}at}(g, x, l) \wedge \operatorname{pos\text{-}rr}(s)) \Rightarrow  \\ \quad \quad \exists d. \\ \quad \quad \quad (\operatorname{pos\text{-}rr}(d) \wedge \forall y \in \mathbb{R}.\; (((y \leq x) \wedge (\lvert (x - y) \rvert \leq d)) \Rightarrow (\lvert (g(x) - g(y)) \rvert \leq s))) \end{array}}}
\vnbprop{deriv-on-ccint-one-sided}{\vnbstmt{\begin{array}{@{}l@{}} \forall g, c, a, b.; \forall x \in \operatorname{ccint}(a, b). \\ \quad \operatorname{has\text{-}deriv\text{-}on\text{-}ccint}(g, c, a, b) \Rightarrow  \\ \quad \quad (((x < b) \Rightarrow \operatorname{has\text{-}right\text{-}deriv\text{-}at}(g, x, c(x))) \wedge ((a < x) \Rightarrow \operatorname{has\text{-}left\text{-}deriv\text{-}at}(g, x, c(x)))) \end{array}}}
\vnbprop{deriv-on-ccint-continuous}{\vnbstmt{\begin{array}{@{}l@{}} \forall g, c, a, b. \\ \quad \operatorname{has\text{-}deriv\text{-}on\text{-}ccint}(g, c, a, b) \Rightarrow \operatorname{is\text{-}continuous\text{-}on}(g, \operatorname{ccint}(a, b)) \end{array}}}
\vnbprop{deriv-on-ccint-primitive}{\vnbstmt{\begin{array}{@{}l@{}} \forall g, c, a, b. \\ \quad \operatorname{has\text{-}deriv\text{-}on\text{-}ccint}(g, c, a, b) \Rightarrow \operatorname{is\text{-}primitive}(g, c, a, b) \end{array}}}
\vnbprop{deriv-within-of-eps-bound}{\vnbstmt{\begin{array}{@{}l@{}} \forall g, c, w, l. \\ \quad w \subseteq \mathbb{R} \wedge \\ \quad c \in w \wedge \\ \quad l \in \mathbb{R} \wedge \\ \quad \forall t \in w.\; (g(t) \in \mathbb{R}) \wedge \\ \quad \forall e. \\ \quad \quad \operatorname{pos\text{-}rr}(e) \Rightarrow  \\ \quad \quad \quad \exists d. \\ \quad \quad \quad \quad \operatorname{pos\text{-}rr}(d) \wedge \\ \quad \quad \quad \quad \forall t \in w. \\ \quad \quad \quad \quad \quad (\lvert (t - c) \rvert < d) \Rightarrow (\lvert ((g(t) - g(c)) - ((t - c) \cdot l)) \rvert \leq (e \cdot \lvert (t - c) \rvert)) \Rightarrow \\ \quad \quad \operatorname{has\text{-}deriv\text{-}at\text{-}within}(g, w, c, l) \end{array}}}
\vnbprop{primitive-continuous-eps-bound}{\vnbstmt{\begin{array}{@{}l@{}} \forall g, s, a, b.; \forall t \in \operatorname{ccint}(a, b), e. \\ \quad (\operatorname{is\text{-}primitive}(g, s, a, b) \wedge (\operatorname{is\text{-}continuous\text{-}on}(s, \operatorname{ccint}(a, b)) \wedge \operatorname{pos\text{-}rr}(e))) \Rightarrow  \\ \quad \quad \exists d. \\ \quad \quad \quad \operatorname{pos\text{-}rr}(d) \wedge \\ \quad \quad \quad \forall \mathit{et} \in \operatorname{ccint}(a, b). \\ \quad \quad \quad \quad (\lvert (\mathit{et} - t) \rvert < d) \Rightarrow (\lvert ((g(\mathit{et}) - g(t)) - ((\mathit{et} - t) \cdot s(t))) \rvert \leq (e \cdot \lvert (\mathit{et} - t) \rvert)) \end{array}}}
\vnbprop{primitive-continuous-deriv-on-ccint}{\vnbstmt{\begin{array}{@{}l@{}} \forall g, s, a, b. \\ \quad (\operatorname{is\text{-}primitive}(g, s, a, b) \wedge \operatorname{is\text{-}continuous\text{-}on}(s, \operatorname{ccint}(a, b))) \Rightarrow  \\ \quad \quad \operatorname{has\text{-}deriv\text{-}on\text{-}ccint}(g, s, a, b) \end{array}}}
\vnbprop{uniform-limit-continuous-on-ccint}{\vnbstmt{\begin{array}{@{}l@{}} \forall a, b, m, i. \\ \quad m \in (\mathbb{N} \to (\operatorname{ccint}(a, b) \to \mathbb{R})) \wedge \\ \quad i \in (\operatorname{ccint}(a, b) \to \mathbb{R}) \wedge \\ \quad \forall k \in \mathbb{N}.\; \operatorname{is\text{-}continuous\text{-}on}(m(k), \operatorname{ccint}(a, b)) \wedge \\ \quad \forall s. \\ \quad \quad \operatorname{pos\text{-}rr}(s) \Rightarrow  \\ \quad \quad \quad \exists n \in \mathbb{N}. \forall k \in \mathbb{N}, x \in \operatorname{ccint}(a, b). \\ \quad \quad \quad \quad (n \leq k) \Rightarrow (\lvert (i(x) - m(k)(x)) \rvert \leq s) \Rightarrow \\ \quad \quad \operatorname{is\text{-}continuous\text{-}on}(i, \operatorname{ccint}(a, b)) \end{array}}}
\vnbprop{deriv-uniform-limit-on-ccint}{\vnbstmt{\begin{array}{@{}l@{}} \forall a, b \in \mathbb{R}, m, d \in (\mathbb{N} \to (\operatorname{ccint}(a, b) \to \mathbb{R})), i \in (\operatorname{ccint}(a, b) \to \mathbb{R}). \\ \quad a < b \wedge \\ \quad \forall k \in \mathbb{N}.\; \operatorname{has\text{-}deriv\text{-}on\text{-}ccint}(m(k), d(k), a, b) \wedge \\ \quad \forall k \in \mathbb{N}.\; \operatorname{is\text{-}continuous\text{-}on}(d(k), \operatorname{ccint}(a, b)) \wedge \\ \quad \exists c \in \mathbb{R}.\; \operatorname{converges\text{-}to}(\operatorname{rr\text{-}ms}, (\lambda k \in \mathbb{N}.\; m(k)(a)), c) \wedge \\ \quad \forall s. \\ \quad \quad \operatorname{pos\text{-}rr}(s) \Rightarrow  \\ \quad \quad \quad \exists n \in \mathbb{N}. \forall k \in \mathbb{N}, x \in \operatorname{ccint}(a, b). \\ \quad \quad \quad \quad (n \leq k) \Rightarrow (\lvert (i(x) - d(k)(x)) \rvert \leq s) \Rightarrow \\ \quad \quad \exists f \in (\operatorname{ccint}(a, b) \to \mathbb{R}). \\ \quad \quad \quad \forall x \in \operatorname{ccint}(a, b). \\ \quad \quad \quad \quad \operatorname{converges\text{-}to}(\operatorname{rr\text{-}ms}, (\lambda k \in \mathbb{N}.\; m(k)(x)), f(x)) \wedge \\ \quad \quad \quad \operatorname{has\text{-}deriv\text{-}on\text{-}ccint}(f, i, a, b) \wedge \\ \quad \quad \quad \operatorname{is\text{-}continuous\text{-}on}(i, \operatorname{ccint}(a, b)) \wedge \\ \quad \quad \quad \forall s. \\ \quad \quad \quad \quad \operatorname{pos\text{-}rr}(s) \Rightarrow  \\ \quad \quad \quad \quad \quad \exists n \in \mathbb{N}. \forall k \in \mathbb{N}, x \in \operatorname{ccint}(a, b). \\ \quad \quad \quad \quad \quad \quad (n \leq k) \Rightarrow (\lvert (i(x) - d(k)(x)) \rvert \leq s) \end{array}}}

\vnbbookfile{theorem-library/road-laws-2}{theorem-library/road-laws-2}
\noindent{\small road-laws-2.scm -- THE OPPOSITE ROAD AND DIEUDONNE (9.6.1).}\par
\vnbprop{ccint-reflect-in}{\vnbstmt{\forall a, b \in \mathbb{R}, y \in \operatorname{ccint}(a, b).\; (((a + b) - y) \in \operatorname{ccint}(a, b))}}
\vnbprop{ooint-reflect-in}{\vnbstmt{\forall a, b \in \mathbb{R}, y \in \operatorname{ooint}(a, b).\; (((a + b) - y) \in \operatorname{ooint}(a, b))}}
\vnbprop{reflect-continuous-at-sub}{\vnbstmt{\begin{array}{@{}l@{}} \forall a, b \in \mathbb{R}, f.; \forall c \in \mathbb{R}, q \in (\operatorname{ccint}(a, b) \to \mathbb{R}), t \in \operatorname{ccint}(a, b). \\ \quad \forall y \in \operatorname{ccint}(a, b).\; (q(y) \simeq (c \cdot f(((a + b) - y)))) \wedge \\ \quad \operatorname{is\text{-}continuous\text{-}at}(\operatorname{subspace\text{-}ms}(\operatorname{rr\text{-}ms}, \operatorname{ccint}(a, b)), \operatorname{rr\text{-}ms}, f, ((a + b) - t)) \Rightarrow \\ \quad \quad \operatorname{is\text{-}continuous\text{-}at}(\operatorname{subspace\text{-}ms}(\operatorname{rr\text{-}ms}, \operatorname{ccint}(a, b)), \operatorname{rr\text{-}ms}, q, t) \end{array}}}
\vnbprop{reflect-continuous-on}{\vnbstmt{\begin{array}{@{}l@{}} \forall a, b \in \mathbb{R}, f.; \forall c \in \mathbb{R}, q \in (\operatorname{ccint}(a, b) \to \mathbb{R}). \\ \quad (\operatorname{is\text{-}continuous\text{-}on}(f, \operatorname{ccint}(a, b)) \wedge \forall y \in \operatorname{ccint}(a, b).\; (q(y) \simeq (c \cdot f(((a + b) - y))))) \Rightarrow  \\ \quad \quad \operatorname{is\text{-}continuous\text{-}on}(q, \operatorname{ccint}(a, b)) \end{array}}}
\vnbprop{regulated-on-reflect}{\vnbstmt{\begin{array}{@{}l@{}} \forall a, b \in \mathbb{R}, f.; \forall c \in \mathbb{R}, q \in (\operatorname{ccint}(a, b) \to \mathbb{R}). \\ \quad (\operatorname{is\text{-}regulated\text{-}on}(f, a, b) \wedge \forall y \in \operatorname{ccint}(a, b).\; (q(y) \simeq (c \cdot f(((a + b) - y))))) \Rightarrow  \\ \quad \quad \operatorname{is\text{-}regulated\text{-}on}(q, a, b) \end{array}}}
\vnbprop{reflect-image-finite}{\vnbstmt{\begin{array}{@{}l@{}} \forall a, b.; \forall d \in \mathbf{Set}. \\ \quad ((a \in \mathbb{R}) \wedge ((b \in \mathbb{R}) \wedge ((|d| \in \mathbb{N}) \wedge (d \subseteq \operatorname{ccint}(a, b))))) \Rightarrow  \\ \quad \quad \operatorname{image}((\lambda x \in \operatorname{ccint}(a, b).\; ((a + b) - x)), d) \in \mathbf{Set} \wedge \\ \quad \quad |\operatorname{image}((\lambda x \in \operatorname{ccint}(a, b).\; ((a + b) - x)), d)| \in \mathbb{N} \wedge \\ \quad \quad \operatorname{image}((\lambda x \in \operatorname{ccint}(a, b).\; ((a + b) - x)), d) \subseteq \operatorname{ccint}(a, b) \end{array}}}
\vnbprop{reflect-image-countable}{\vnbstmt{\begin{array}{@{}l@{}} \forall a, b, d. \\ \quad ((a \in \mathbb{R}) \wedge ((b \in \mathbb{R}) \wedge (\operatorname{is\text{-}countable}(d) \wedge (d \subseteq \operatorname{ccint}(a, b))))) \Rightarrow  \\ \quad \quad \operatorname{is\text{-}countable}(\operatorname{image}((\lambda x \in \operatorname{ccint}(a, b).\; ((a + b) - x)), d)) \wedge \\ \quad \quad \operatorname{image}((\lambda x \in \operatorname{ccint}(a, b).\; ((a + b) - x)), d) \subseteq \operatorname{ccint}(a, b) \end{array}}}
\vnbprop{reflect-not-in-image}{\vnbstmt{\begin{array}{@{}l@{}} \forall a, b, d.; \forall t \in \operatorname{ccint}(a, b). \\ \quad a \in \mathbb{R} \wedge \\ \quad b \in \mathbb{R} \wedge \\ \quad d \subseteq \operatorname{ccint}(a, b) \wedge \\ \quad \neg (t \in \operatorname{image}((\lambda x \in \operatorname{ccint}(a, b).\; ((a + b) - x)), d)) \Rightarrow \\ \quad \quad \neg (((a + b) - t) \in d) \end{array}}}
\vnbprop{pw-antiderivative-reflect}{\vnbstmt{\begin{array}{@{}l@{}} \forall a, b \in \mathbb{R}, f, i.; \forall g, k \in (\operatorname{ccint}(a, b) \to \mathbb{R}). \\ \quad \operatorname{is\text{-}primitive}(f, i, a, b) \wedge \\ \quad \forall y \in \operatorname{ccint}(a, b).\; (g(y) \simeq f(((a + b) - y))) \wedge \\ \quad \forall y \in \operatorname{ccint}(a, b).\; (k(y) \simeq (-1 \cdot i(((a + b) - y)))) \Rightarrow \\ \quad \quad \operatorname{is\text{-}primitive}(g, k, a, b) \end{array}}}
\vnbprop{pw-int-reflect}{\vnbstmt{\begin{array}{@{}l@{}} \forall a, b \in \mathbb{R}, f, i.; \forall h, s \in (\operatorname{ccint}(a, b) \to \mathbb{R}). \\ \quad \operatorname{is\text{-}primitive}(f, i, a, b) \wedge \\ \quad \forall y \in \operatorname{ccint}(a, b).\; (h(y) \simeq (-1 \cdot f(((a + b) - y)))) \wedge \\ \quad \forall y \in \operatorname{ccint}(a, b).\; (s(y) \simeq i(((a + b) - y))) \Rightarrow \\ \quad \quad \operatorname{is\text{-}primitive}(h, s, a, b) \wedge (\int_{a}^{b} s = \int_{a}^{b} i) \end{array}}}
\vnbprop{pw-continuous-reflect}{\vnbstmt{\begin{array}{@{}l@{}} \forall a, b \in \mathbb{R}, i.; \forall c \in \mathbb{R}, q \in (\operatorname{ccint}(a, b) \to \mathbb{R}). \\ \quad (\operatorname{is\text{-}pw\text{-}continuous\text{-}on}(i, a, b) \wedge \forall y \in \operatorname{ccint}(a, b).\; (q(y) \simeq (c \cdot i(((a + b) - y))))) \Rightarrow  \\ \quad \quad \operatorname{is\text{-}pw\text{-}continuous\text{-}on}(q, a, b) \end{array}}}
\vnbprop{trace-opposite-subset}{\vnbstmt{\begin{array}{@{}l@{}} \forall a, b \in \mathbb{R}, m, \mathit{prgb}. \\ \quad m \in (\operatorname{ccint}(a, b) \to \mathbb{C}) \wedge \\ \quad \mathit{prgb} \in (\operatorname{ccint}(a, b) \to \mathbb{C}) \wedge \\ \quad \forall y \in \operatorname{ccint}(a, b).\; (\operatorname{prgb}(y) \simeq m(((a + b) - y))) \Rightarrow \\ \quad \quad \operatorname{trace}(\mathit{prgb}, a, b) \subseteq \operatorname{trace}(m, a, b) \end{array}}}
\vnbprop{trace-opposite}{\vnbstmt{\begin{array}{@{}l@{}} \forall a, b \in \mathbb{R}, m, \mathit{prgb}. \\ \quad m \in (\operatorname{ccint}(a, b) \to \mathbb{C}) \wedge \\ \quad \mathit{prgb} \in (\operatorname{ccint}(a, b) \to \mathbb{C}) \wedge \\ \quad \forall y \in \operatorname{ccint}(a, b).\; (\operatorname{prgb}(y) \simeq m(((a + b) - y))) \Rightarrow \\ \quad \quad ((\operatorname{trace}(\mathit{prgb}, a, b) \subseteq \operatorname{trace}(m, a, b)) \wedge (\operatorname{trace}(m, a, b) \subseteq \operatorname{trace}(\mathit{prgb}, a, b))) \end{array}}}
\vnbprop{opposite-is-road}{\vnbstmt{\begin{array}{@{}l@{}} \forall a, b \in \mathbb{R}, m, c, \mathit{prgb}, \mathit{prdb}. \\ \quad \operatorname{is\text{-}road}(m, c, a, b) \wedge \\ \quad \mathit{prgb} \in (\operatorname{ccint}(a, b) \to \mathbb{C}) \wedge \\ \quad \mathit{prdb} \in (\operatorname{ccint}(a, b) \to \mathbb{C}) \wedge \\ \quad \forall y \in \operatorname{ccint}(a, b).\; (\operatorname{prgb}(y) \simeq m(((a + b) - y))) \wedge \\ \quad \forall y \in \operatorname{ccint}(a, b).\; (\operatorname{prdb}(y) \simeq (-1 \cdot c(((a + b) - y)))) \Rightarrow \\ \quad \quad \operatorname{is\text{-}road}(\mathit{prgb}, \mathit{prdb}, a, b) \end{array}}}
\vnbprop{cc-int-reflect}{\vnbstmt{\begin{array}{@{}l@{}} \forall a, b, i, \mathit{prb}, f, w. \\ \quad a \in \mathbb{R} \wedge \\ \quad b \in \mathbb{R} \wedge \\ \quad i \in (\operatorname{ccint}(a, b) \to \mathbb{C}) \wedge \\ \quad \mathit{prb} \in (\operatorname{ccint}(a, b) \to \mathbb{C}) \wedge \\ \quad \forall y \in \operatorname{ccint}(a, b).\; (\operatorname{prb}(y) \simeq i(((a + b) - y))) \wedge \\ \quad \operatorname{is\text{-}primitive}(f, (\lambda t \in \operatorname{ccint}(a, b).\; \operatorname{real\text{-}part}(i(t))), a, b) \wedge \\ \quad \operatorname{is\text{-}primitive}(w, (\lambda t \in \operatorname{ccint}(a, b).\; \operatorname{imag\text{-}part}(i(t))), a, b) \Rightarrow \\ \quad \quad \int_{a}^{b} \mathit{prb} = \int_{a}^{b} i \end{array}}}
\vnbprop{line-int-opposite}{\vnbstmt{\begin{array}{@{}l@{}} \forall a, b, m, c, f, d, \mathit{prgb}, \mathit{prdb}, \mathit{pwf}, w. \\ \quad a \in \mathbb{R} \wedge \\ \quad b \in \mathbb{R} \wedge \\ \quad \operatorname{is\text{-}road}(m, c, a, b) \wedge \\ \quad f \in (d \to \mathbb{C}) \wedge \\ \quad \operatorname{trace}(m, a, b) \subseteq d \wedge \\ \quad \mathit{prgb} \in (\operatorname{ccint}(a, b) \to \mathbb{C}) \wedge \\ \quad \mathit{prdb} \in (\operatorname{ccint}(a, b) \to \mathbb{C}) \wedge \\ \quad \forall y \in \operatorname{ccint}(a, b).\; (\operatorname{prgb}(y) \simeq m(((a + b) - y))) \wedge \\ \quad \forall y \in \operatorname{ccint}(a, b).\; (\operatorname{prdb}(y) \simeq (-1 \cdot c(((a + b) - y)))) \wedge \\ \quad \operatorname{is\text{-}primitive}(\mathit{pwf}, \\ \quad \quad (\lambda t \in \operatorname{ccint}(a, b).\; \operatorname{real\text{-}part}((\lambda t \in \operatorname{ccint}(a, b).\; (f(m(t)) \cdot c(t)))(t))), \\ \quad \quad a, b) \wedge \\ \quad \operatorname{is\text{-}primitive}(w, \\ \quad \quad (\lambda t \in \operatorname{ccint}(a, b).\; \operatorname{imag\text{-}part}((\lambda t \in \operatorname{ccint}(a, b).\; (f(m(t)) \cdot c(t)))(t))), \\ \quad \quad a, b) \Rightarrow \\ \quad \quad \operatorname{line\text{-}int}(f, \mathit{prgb}, \mathit{prdb}, a, b) = (-1 \cdot \operatorname{line\text{-}int}(f, m, c, a, b)) \end{array}}}

\vnbbookfile{theorem-library/nn-power-two-bound}{theorem-library/nn-power-two-bound}
\noindent{\small nn-power-two-bound.scm -- lemmas MOVED 2026-10-27 (integration of the 10-26/27 wave) out of theorem-library/goursat-bricks.scm, so that files of library analysis may cite them (the library gate). The proofs are the originals, verbatim; the helper defines are a copy of goursat-bricks.scm's.  Loaded immediately before ...}\par
\vnbprop{nn-succ-le-power-two}{\vnbstmt{\forall k \in \mathbb{N}.\; ((k + 1) \leq \operatorname{power}(2, k))}}

\vnbbookfile{theorem-library/compact-image}{theorem-library/compact-image}
\noindent{\small compact-image.scm -- THE CONTINUOUS IMAGE OF A COMPACT SPACE IS COMPACT (calculus.pdf, Proposition 3.19), and the COMPACTNESS OF THE TRACE of a path, which follows from it (complex-analysis.pdf 3.1.1: "gamma\_* is sometimes called the trace of gamma.  It is always a compact set").}\par
\vnbprop{preimage-in-power}{\vnbstmt{\forall s, f, u.\; ((\operatorname{pts}(s) \in \mathbf{Set}) \Rightarrow (\operatorname{preimage}(s, f, u) \in \operatorname{power}(\operatorname{pts}(s))))}}
\vnbprop{image-of-preimage-onto}{\vnbstmt{\begin{array}{@{}l@{}} \forall s, f, u. \\ \quad ((\operatorname{pts}(s) \in \mathbf{Set}) \wedge (u \subseteq \operatorname{image}(f, \operatorname{pts}(s)))) \Rightarrow  \\ \quad \quad \operatorname{image}(f, \operatorname{preimage}(s, f, u)) = u \end{array}}}
\vnbprop{preimage-subspace-open}{\vnbstmt{\begin{array}{@{}l@{}} \forall s, t, f, b, u. \\ \quad \operatorname{is\text{-}continuous}(s, t, f) \wedge \\ \quad b \subseteq \operatorname{pts}(t) \wedge \\ \quad \forall p \in \operatorname{pts}(s).\; (f(p) \in b) \wedge \\ \quad \operatorname{is\text{-}open}(\operatorname{subspace\text{-}ms}(t, b), u) \Rightarrow \\ \quad \quad \operatorname{is\text{-}open}(s, \operatorname{preimage}(s, f, u)) \end{array}}}
\vnbprop{compact-image}{\vnbstmt{\begin{array}{@{}l@{}} \forall s, t, f. \\ \quad (\operatorname{is\text{-}continuous}(s, t, f) \wedge \operatorname{is\text{-}compact}(s)) \Rightarrow  \\ \quad \quad \operatorname{is\text{-}compact}(\operatorname{subspace\text{-}ms}(t, \operatorname{image}(f, \operatorname{pts}(s)))) \end{array}}}
\vnbprop{image-restrict}{\vnbstmt{\forall a \in \mathbf{Set}, f.\; (\operatorname{image}(\operatorname{restrict}(f, a), a) = \operatorname{image}(f, a))}}
\vnbprop{continuous-image-of-compact}{\vnbstmt{\begin{array}{@{}l@{}} \forall s, t, f, a. \\ \quad (\operatorname{is\text{-}continuous}(s, t, f) \wedge ((a \subseteq \operatorname{pts}(s)) \wedge \operatorname{is\text{-}compact}(\operatorname{subspace\text{-}ms}(s, a)))) \Rightarrow  \\ \quad \quad \operatorname{is\text{-}compact}(\operatorname{subspace\text{-}ms}(t, \operatorname{image}(f, a))) \end{array}}}
\vnbprop{path-trace-compact}{\vnbstmt{\begin{array}{@{}l@{}} \forall m, v, \mathit{bv}. \\ \quad \operatorname{is\text{-}path}(m, v, \mathit{bv}) \Rightarrow \operatorname{is\text{-}compact}(\operatorname{subspace\text{-}ms}(\operatorname{cc\text{-}ms}, \operatorname{trace}(m, v, \mathit{bv}))) \end{array}}}

\vnbbookfile{structure-library/calc3-vocab}{structure-library/calc3-vocab}
\noindent{\small calc3-vocab.scm -- vocabulary for calculus.pdf chapter 3 (metric topology): the ACCUMULATION POINT of a set (before Prop 3.10) and the HOMEOMORPHISM (Cor 3.20), the DENSE subset and the LIMIT ALONG a subset (Prop 3.26). DEFINITIONS ONLY; the laws are proven in theorem-library/calc3-*.scm.}\par
\par\medskip\noindent\textbf{Definition} (\texttt{is-accumulation-point}).\;For all $\operatorname{apts\_}$, $\operatorname{apta\_}$ and $\operatorname{aptx\_}$, 
\begin{equation*}
\begin{array}{@{}l@{}} \operatorname{is\text{-}accumulation\text{-}point}(\operatorname{apts\_}, \operatorname{apta\_}, \operatorname{aptx\_}) \\ \Leftrightarrow\; \operatorname{is\text{-}metric\text{-}space}(\operatorname{apts\_}) \wedge \\ \operatorname{aptx\_} \in \operatorname{pts}(\operatorname{apts\_}) \wedge \\ \forall \operatorname{aptr\_}. \\ \quad \operatorname{pos\text{-}rr}(\operatorname{aptr\_}) \Rightarrow  \\ \quad \quad \exists \operatorname{apty\_} \in \operatorname{apta\_}. \\ \quad \quad \quad (\operatorname{apty\_} \in \operatorname{ball}(\operatorname{apts\_}, \operatorname{aptx\_}, \operatorname{aptr\_})) \wedge \neg (\operatorname{apty\_} = \operatorname{aptx\_}) \end{array}
\end{equation*}

\par\medskip\noindent\textbf{Definition} (\texttt{is-homeomorphism}).\;For all $\operatorname{hmes\_}$, $\operatorname{hmet\_}$ and $\operatorname{hmef\_}$, 
\begin{equation*}
\begin{array}{@{}l@{}} \operatorname{is\text{-}homeomorphism}(\operatorname{hmes\_}, \operatorname{hmet\_}, \operatorname{hmef\_}) \\ \Leftrightarrow\; \operatorname{hmef\_} \in \mathrm{Bij}(\operatorname{pts}(\operatorname{hmes\_}), \operatorname{pts}(\operatorname{hmet\_})) \wedge \\ \operatorname{is\text{-}continuous}(\operatorname{hmes\_}, \operatorname{hmet\_}, \operatorname{hmef\_}) \wedge \\ \operatorname{is\text{-}continuous}(\operatorname{hmet\_}, \operatorname{hmes\_}, \operatorname{hmef\_}^{-1}) \end{array}
\end{equation*}

\par\medskip\noindent\textbf{Definition} (\texttt{is-dense}).\;For all $\operatorname{dns\_}$ and $\operatorname{dna\_}$, 
\begin{equation*}
\begin{array}{@{}l@{}} \operatorname{is\text{-}dense}(\operatorname{dns\_}, \operatorname{dna\_}) \\ \Leftrightarrow\; (\operatorname{is\text{-}metric\text{-}space}(\operatorname{dns\_}) \wedge ((\operatorname{dna\_} \subseteq \operatorname{pts}(\operatorname{dns\_})) \wedge (\operatorname{pts}(\operatorname{dns\_}) \subseteq \operatorname{closure}(\operatorname{dns\_}, \operatorname{dna\_})))) \end{array}
\end{equation*}

\par\medskip\noindent\textbf{Definition} (\texttt{is-limit-along}).\;For all $\operatorname{lmas\_}$, $\operatorname{lmat\_}$, $\operatorname{lmaa\_}$, $\operatorname{lmaf\_}$, $\operatorname{lmax\_}$ and $\operatorname{lmay\_}$, 
\begin{equation*}
\begin{array}{@{}l@{}} \operatorname{is\text{-}limit\text{-}along}(\operatorname{lmas\_}, \operatorname{lmat\_}, \operatorname{lmaa\_}, \operatorname{lmaf\_}, \operatorname{lmax\_}, \operatorname{lmay\_}) \\ \Leftrightarrow\; \operatorname{is\text{-}metric\text{-}space}(\operatorname{lmas\_}) \wedge \\ \operatorname{is\text{-}metric\text{-}space}(\operatorname{lmat\_}) \wedge \\ \operatorname{lmay\_} \in \operatorname{pts}(\operatorname{lmat\_}) \wedge \\ \forall \operatorname{lmae\_}. \\ \quad \operatorname{pos\text{-}rr}(\operatorname{lmae\_}) \Rightarrow  \\ \quad \quad \exists \operatorname{lmad\_}. \\ \quad \quad \quad \operatorname{pos\text{-}rr}(\operatorname{lmad\_}) \wedge \\ \quad \quad \quad \forall \operatorname{lmaz\_} \in \operatorname{lmaa\_}. \\ \quad \quad \quad \quad (\operatorname{dist}(\operatorname{lmas\_})(\operatorname{lmax\_}, \operatorname{lmaz\_}) \leq \operatorname{lmad\_}) \Rightarrow  \\ \quad \quad \quad \quad \quad \operatorname{dist}(\operatorname{lmat\_})(\operatorname{lmaf\_}(\operatorname{lmaz\_}), \operatorname{lmay\_}) \leq \operatorname{lmae\_} \end{array}
\end{equation*}


\vnbbookfile{theorem-library/calc3-closed}{theorem-library/calc3-closed}
\noindent{\small calc3-closed.scm -- calculus.pdf Prop 3.10: A SET IS CLOSED IFF IT CONTAINS ALL ITS ACCUMULATION POINTS.  Agent CALC-3, 2026-10-19.}\par
\vnbprop{closed-iff-contains-accumulation-points}{\vnbstmt{\begin{array}{@{}l@{}} \forall s, a. \\ \quad (\operatorname{is\text{-}metric\text{-}space}(s) \wedge (a \subseteq \operatorname{pts}(s))) \Rightarrow  \\ \quad \quad \operatorname{is\text{-}closed}(s, a) \Leftrightarrow \forall x.\; (\operatorname{is\text{-}accumulation\text{-}point}(s, a, x) \Rightarrow (x \in a)) \end{array}}}

\vnbbookfile{theorem-library/calc3-compact-closed}{theorem-library/calc3-compact-closed}
\noindent{\small calc3-compact-closed.scm -- calculus.pdf Cor 3.20: A CONTINUOUS BIJECTION FROM A COMPACT METRIC SPACE IS A HOMEOMORPHISM, and the brick it was missing, A COMPACT SUBSET OF A METRIC SPACE IS CLOSED.  Agent CALC-3, 2026-10-19.}\par
\vnbprop{compact-subspace-is-closed}{\vnbstmt{\begin{array}{@{}l@{}} \forall s, a. \\ \quad (\operatorname{is\text{-}metric\text{-}space}(s) \wedge ((a \subseteq \operatorname{pts}(s)) \wedge \operatorname{is\text{-}compact}(\operatorname{subspace\text{-}ms}(s, a)))) \Rightarrow  \\ \quad \quad \operatorname{is\text{-}closed}(s, a) \end{array}}}
\vnbprop{compact-continuous-bijection-is-homeomorphism}{\vnbstmt{\begin{array}{@{}l@{}} \forall s, t, f. \\ \quad (\operatorname{is\text{-}compact}(s) \wedge (\operatorname{is\text{-}continuous}(s, t, f) \wedge (f \in \mathrm{Bij}(\operatorname{pts}(s), \operatorname{pts}(t))))) \Rightarrow  \\ \quad \quad \operatorname{is\text{-}homeomorphism}(s, t, f) \end{array}}}

\vnbbookfile{theorem-library/compact-inter-closed}{theorem-library/compact-inter-closed}
\noindent{\small compact-inter-closed.scm -- lemmas MOVED 2026-11-10 (integration of the 11-09/10 waves) out of theorem-library/mb35-regular.scm into library analysis, for the library gate: cited by ccr4b-base; topology. The proofs are the originals, verbatim; the helper defines are a copy of the source's. compact-inter-closed}\par
\vnbprop{compact-inter-closed}{\vnbstmt{\begin{array}{@{}l@{}} \forall s, k, f. \\ \quad (\operatorname{is\text{-}metric\text{-}space}(s) \wedge ((k \subseteq \operatorname{pts}(s)) \wedge (\operatorname{is\text{-}compact}(\operatorname{subspace\text{-}ms}(s, k)) \wedge \operatorname{is\text{-}closed}(s, f)))) \Rightarrow  \\ \quad \quad \operatorname{is\text{-}compact}(\operatorname{subspace\text{-}ms}(s, (k \cap f))) \end{array}}}

\vnbbookfile{theorem-library/rr-square-facts}{theorem-library/rr-square-facts}
\noindent{\small rr-square-facts.scm -- lemmas MOVED 2026-11-10 (integration of the 11-09/10 waves) out of theorem-library/complex-measure-laws.scm into library analysis, for the library gate: cited by the Hilbert-family files; RR arithmetic. The proofs are the originals, verbatim; the helper defines are a copy of the source's. ...}\par
\vnbprop{sqrt-gt-iff}{\vnbstmt{\begin{array}{@{}l@{}} \forall l, t. \\ \quad ((l \in \mathbb{R}) \wedge ((0 \leq l) \wedge ((t \in \mathbb{R}) \wedge (0 \leq t)))) \Rightarrow ((l < \sqrt{t}) \Leftrightarrow ((l \cdot l) < t)) \end{array}}}
\vnbprop{rr-square-zero}{\vnbstmt{\forall x \in \mathbb{R}.\; (((x \cdot x) = 0) \Rightarrow (x = 0))}}
\vnbprop{rr-square-mono}{\vnbstmt{\forall a, b.\; (((a \in \mathbb{R}) \wedge ((b \in \mathbb{R}) \wedge ((0 \leq a) \wedge (a \leq b)))) \Rightarrow ((a \cdot a) \leq (b \cdot b)))}}
\vnbprop{rr-square-le-imp-le}{\vnbstmt{\begin{array}{@{}l@{}} \forall a, b. \\ \quad ((a \in \mathbb{R}) \wedge ((b \in \mathbb{R}) \wedge ((0 \leq a) \wedge ((0 \leq b) \wedge ((a \cdot a) \leq (b \cdot b)))))) \Rightarrow (a \leq b) \end{array}}}

\vnbbookfile{theorem-library/rr-nonneg-below-all-zero}{theorem-library/rr-nonneg-below-all-zero}
\noindent{\small rr-nonneg-below-all-zero.scm -- lemmas MOVED 2026-11-10 (integration of the 11-09/10 waves) out of theorem-library/mb35-souslin.scm into library analysis, for the library gate: cited by three CCR ch.3 files; RR order. The proofs are the originals, verbatim; the helper defines are a copy of the source's. ...}\par
\vnbprop{rr-nonneg-below-all-zero}{\vnbstmt{\begin{array}{@{}l@{}} \forall x \in \mathbb{R}. \\ \quad ((0 \leq x) \wedge \forall s.\; (\operatorname{pos\text{-}rr}(s) \Rightarrow (x \leq s))) \Rightarrow (x = 0) \end{array}}}

\vnbbookfile{structure-library/ccr2-vocab}{structure-library/ccr2-vocab}
\noindent{\small ccr2-vocab.scm -- NON-ARCHIMEDEAN VALUED FIELDS AND THEIR DOMAINS (CCR.pdf chapters 1-2 and 4: Definitions 1.8, 1.12, 1.14, 2.5, 4.3).  DEFINITIONS ONLY; the laws are in theorem-library/ccr2-nv.scm. Agent CCR-B2, wave MAP-1026, 2026-10-26.  Destination: structure-library/ccr2-vocab.scm.  Library: analysis.  Load ...}\par
\par\medskip\noindent\textbf{Definition} (\texttt{is-fractions-of}).\;For all $\operatorname{cwk\_}$ and $\operatorname{cwa\_}$, 
\begin{equation*}
\begin{array}{@{}l@{}} \operatorname{is\text{-}fractions\text{-}of}(\operatorname{cwk\_}, \operatorname{cwa\_}) \\ \Leftrightarrow\; \forall \operatorname{cwx\_} \in \operatorname{carr}(\operatorname{cwk\_}). \exists \operatorname{cwu\_}, \operatorname{cwv\_} \in \operatorname{cwa\_}. \\ \quad \neg (\operatorname{cwv\_} = \operatorname{zero}(\operatorname{cwk\_})) \wedge (\operatorname{mul}(\operatorname{cwk\_})(\operatorname{cwx\_}, \operatorname{cwv\_}) = \operatorname{cwu\_}) \end{array}
\end{equation*}

\par\medskip\noindent\textbf{Definition} (\texttt{is-nv-ideal}).\;For all $\operatorname{cwk\_}$, $\operatorname{cwa\_}$ and $\operatorname{cwi\_}$, 
\begin{equation*}
\begin{array}{@{}l@{}} \operatorname{is\text{-}nv\text{-}ideal}(\operatorname{cwk\_}, \operatorname{cwa\_}, \operatorname{cwi\_}) \\ \Leftrightarrow\; \forall \operatorname{cwx\_} \in \operatorname{cwi\_}.\; (\operatorname{cwx\_} \in \operatorname{cwa\_}) \wedge \\ \operatorname{zero}(\operatorname{cwk\_}) \in \operatorname{cwi\_} \wedge \\ \forall \operatorname{cwx\_}, \operatorname{cwy\_} \in \operatorname{cwi\_}.\; (\operatorname{add}(\operatorname{cwk\_})(\operatorname{cwx\_}, \operatorname{cwy\_}) \in \operatorname{cwi\_}) \wedge \\ \forall \operatorname{cwx\_} \in \operatorname{cwi\_}.\; (\operatorname{neg}(\operatorname{cwk\_})(\operatorname{cwx\_}) \in \operatorname{cwi\_}) \wedge \\ \forall \operatorname{cwr\_} \in \operatorname{cwa\_}, \operatorname{cwx\_} \in \operatorname{cwi\_}.\; (\operatorname{mul}(\operatorname{cwk\_})(\operatorname{cwr\_}, \operatorname{cwx\_}) \in \operatorname{cwi\_}) \end{array}
\end{equation*}

\par\medskip\noindent\textbf{Definition} (\texttt{is-nv-field}).\;For every $\operatorname{cwk\_}$, 
\begin{equation*}
\begin{array}{@{}l@{}} \operatorname{is\text{-}nv\text{-}field}(\operatorname{cwk\_}) \\ \Leftrightarrow\; \operatorname{is\text{-}normed\text{-}field}(\operatorname{cwk\_}) \wedge \\ \forall \operatorname{cwx\_}, \operatorname{cwy\_} \in \operatorname{carr}(\operatorname{cwk\_}). \\ \quad \operatorname{fnrm}(\operatorname{cwk\_})(\operatorname{add}(\operatorname{cwk\_})(\operatorname{cwx\_}, \operatorname{cwy\_})) \leq \operatorname{max}(\operatorname{fnrm}(\operatorname{cwk\_})(\operatorname{cwx\_}), \operatorname{fnrm}(\operatorname{cwk\_})(\operatorname{cwy\_})) \end{array}
\end{equation*}

\par\medskip\noindent\textbf{Definition} (\texttt{nv-integers}).\;$\operatorname{nv\text{-}integers}(\operatorname{cwk\_}) :=$ $\{\operatorname{cwx\_} \in \operatorname{carr}(\operatorname{cwk\_}) : (\operatorname{fnrm}(\operatorname{cwk\_})(\operatorname{cwx\_}) \leq 1)\}$
\par\medskip\noindent\textbf{Definition} (\texttt{is-nv-domain}).\;For all $\operatorname{cwk\_}$ and $\operatorname{cwa\_}$, 
\begin{equation*}
\begin{array}{@{}l@{}} \operatorname{is\text{-}nv\text{-}domain}(\operatorname{cwk\_}, \operatorname{cwa\_}) \\ \Leftrightarrow\; \operatorname{is\text{-}nv\text{-}field}(\operatorname{cwk\_}) \wedge \\ \forall \operatorname{cwx\_} \in \operatorname{cwa\_}.\; (\operatorname{cwx\_} \in \operatorname{carr}(\operatorname{cwk\_})) \wedge \\ \operatorname{zero}(\operatorname{cwk\_}) \in \operatorname{cwa\_} \wedge \\ \operatorname{one}(\operatorname{cwk\_}) \in \operatorname{cwa\_} \wedge \\ \forall \operatorname{cwx\_}, \operatorname{cwy\_} \in \operatorname{cwa\_}.\; (\operatorname{add}(\operatorname{cwk\_})(\operatorname{cwx\_}, \operatorname{cwy\_}) \in \operatorname{cwa\_}) \wedge \\ \forall \operatorname{cwx\_}, \operatorname{cwy\_} \in \operatorname{cwa\_}.\; (\operatorname{mul}(\operatorname{cwk\_})(\operatorname{cwx\_}, \operatorname{cwy\_}) \in \operatorname{cwa\_}) \wedge \\ \forall \operatorname{cwx\_} \in \operatorname{cwa\_}.\; (\operatorname{neg}(\operatorname{cwk\_})(\operatorname{cwx\_}) \in \operatorname{cwa\_}) \wedge \\ \forall \operatorname{cwx\_} \in \operatorname{cwa\_}.\; (\operatorname{fnrm}(\operatorname{cwk\_})(\operatorname{cwx\_}) \leq 1) \end{array}
\end{equation*}

\par\medskip\noindent\textbf{Definition} (\texttt{is-compact-nv-domain}).\;For all $\operatorname{cwk\_}$ and $\operatorname{cwa\_}$, 
\begin{equation*}
\begin{array}{@{}l@{}} \operatorname{is\text{-}compact\text{-}nv\text{-}domain}(\operatorname{cwk\_}, \operatorname{cwa\_}) \\ \Leftrightarrow\; (\operatorname{is\text{-}nv\text{-}domain}(\operatorname{cwk\_}, \operatorname{cwa\_}) \wedge \operatorname{is\text{-}compact}(\operatorname{subspace\text{-}ms}(\operatorname{nf\text{-}metric\text{-}space}(\operatorname{cwk\_}), \operatorname{cwa\_}))) \end{array}
\end{equation*}

\par\medskip\noindent\textbf{Definition} (\texttt{nv-units}).\;$\operatorname{nv\text{-}units}(\operatorname{cwk\_}, \operatorname{cwa\_}) :=$
\begin{equation*}
\{\operatorname{cwx\_} \in \operatorname{cwa\_} : \exists \operatorname{cwy\_} \in \operatorname{cwa\_}.\; (\operatorname{mul}(\operatorname{cwk\_})(\operatorname{cwx\_}, \operatorname{cwy\_}) = \operatorname{one}(\operatorname{cwk\_}))\}
\end{equation*}

\par\medskip\noindent\textbf{Definition} (\texttt{nv-maximal}).\;$\operatorname{nv\text{-}maximal}(\operatorname{cwk\_}, \operatorname{cwa\_}) :=$ $\{\operatorname{cwx\_} \in \operatorname{cwa\_} : (\operatorname{fnrm}(\operatorname{cwk\_})(\operatorname{cwx\_}) < 1)\}$
\par\medskip\noindent\textbf{Definition} (\texttt{nv-principal}).\;$\operatorname{nv\text{-}principal}(\operatorname{cwk\_}, \operatorname{cwa\_}, \operatorname{cwg\_}) :=$
\begin{equation*}
\{\operatorname{cwx\_} \in \operatorname{cwa\_} : \exists \operatorname{cwm\_} \in \operatorname{cwa\_}.\; (\operatorname{cwx\_} = \operatorname{mul}(\operatorname{cwk\_})(\operatorname{cwm\_}, \operatorname{cwg\_}))\}
\end{equation*}

\par\medskip\noindent\textbf{Definition} (\texttt{nv-ideal2}).\;$\operatorname{nv\text{-}ideal2}(\operatorname{cwk\_}, \operatorname{cwa\_}, \operatorname{cwg\_}, \operatorname{cwh\_}) :=$
\begin{equation*}
\{\operatorname{cwx\_} \in \operatorname{cwa\_} : \exists \operatorname{cwm\_}, \operatorname{cwn\_} \in \operatorname{cwa\_}.\; (\operatorname{cwx\_} = \operatorname{add}(\operatorname{cwk\_})(\operatorname{mul}(\operatorname{cwk\_})(\operatorname{cwm\_}, \operatorname{cwg\_}), \operatorname{mul}(\operatorname{cwk\_})(\operatorname{cwn\_}, \operatorname{cwh\_})))\}
\end{equation*}

\par\medskip\noindent\textbf{Definition} (\texttt{is-nv-pid}).\;For all $\operatorname{cwk\_}$ and $\operatorname{cwa\_}$, 
\begin{equation*}
\begin{array}{@{}l@{}} \operatorname{is\text{-}nv\text{-}pid}(\operatorname{cwk\_}, \operatorname{cwa\_}) \\ \Leftrightarrow\; \operatorname{is\text{-}nv\text{-}domain}(\operatorname{cwk\_}, \operatorname{cwa\_}) \wedge \\ \forall \operatorname{cwi\_}. \\ \quad \operatorname{is\text{-}nv\text{-}ideal}(\operatorname{cwk\_}, \operatorname{cwa\_}, \operatorname{cwi\_}) \Rightarrow  \\ \quad \quad \exists \operatorname{cwg\_} \in \operatorname{cwa\_}.\; (\operatorname{cwi\_} = \operatorname{nv\text{-}principal}(\operatorname{cwk\_}, \operatorname{cwa\_}, \operatorname{cwg\_})) \end{array}
\end{equation*}

\par\medskip\noindent\textbf{Definition} (\texttt{nv-ball}).\;$\operatorname{nv\text{-}ball}(\operatorname{cwk\_}, \operatorname{cwa\_}, \operatorname{cwr\_}) :=$ $\{\operatorname{cwx\_} \in \operatorname{cwa\_} : (\operatorname{fnrm}(\operatorname{cwk\_})(\operatorname{cwx\_}) \leq \operatorname{cwr\_})\}$

\vnbbookfile{structure-library/ccr1109a2-vocab}{structure-library/ccr1109a2-vocab}
\noindent{\small ccr1109a2-vocab.scm -- wave of 2026-11-09/10 (agent CCR-1109-A2), wired 2026-11-10 from scratchpad/ccr1109a2/ccr1109a2-vocab.scm. Library: analysis.}\par
\par\medskip\noindent\textbf{Definition} (\texttt{is-complete-nv-domain}).\;For all $\operatorname{c2k\_}$ and $\operatorname{c2a\_}$, 
\begin{equation*}
\begin{array}{@{}l@{}} \operatorname{is\text{-}complete\text{-}nv\text{-}domain}(\operatorname{c2k\_}, \operatorname{c2a\_}) \\ \Leftrightarrow\; (\operatorname{is\text{-}nv\text{-}domain}(\operatorname{c2k\_}, \operatorname{c2a\_}) \wedge \operatorname{is\text{-}complete}(\operatorname{subspace\text{-}ms}(\operatorname{nf\text{-}metric\text{-}space}(\operatorname{c2k\_}), \operatorname{c2a\_}))) \end{array}
\end{equation*}

\par\medskip\noindent\textbf{Definition} (\texttt{nv-u1}).\;$\operatorname{nv\text{-}u1}(\operatorname{c2k\_}, \operatorname{c2a\_}) :=$
\begin{equation*}
\{\operatorname{c2x\_} \in \operatorname{c2a\_} : (\operatorname{fnrm}(\operatorname{c2k\_})(\operatorname{add}(\operatorname{c2k\_})(\operatorname{c2x\_}, \operatorname{neg}(\operatorname{c2k\_})(\operatorname{one}(\operatorname{c2k\_})))) < 1)\}
\end{equation*}


\vnbbookfile{structure-library/ccr-ch3-vocab}{structure-library/ccr-ch3-vocab}
\noindent{\small ccr-ch3-vocab.scm -- wave of 2026-11-09/10 (agent CCR-1109-B), wired 2026-11-10 from scratchpad/ccr1109b/ccr-ch3-vocab.scm. Library: analysis.}\par
\par\medskip\noindent\textbf{Definition} (\texttt{is-mat-bound}).\;For all $k$, $m$, $n$, $a$ and $c$, 
\begin{equation*}
\begin{array}{@{}l@{}} \operatorname{is\text{-}mat\text{-}bound}(k, m, n, a, c) \\ \Leftrightarrow\; a \in \operatorname{mat}(m, n, \operatorname{carr}(k)) \wedge \\ \forall \operatorname{cbvi\_} \in \operatorname{interval}(1, m), \operatorname{cbvj\_} \in \operatorname{interval}(1, n). \\ \quad \operatorname{fnrm}(k)(\operatorname{entry}(a, \operatorname{cbvi\_}, \operatorname{cbvj\_})) \leq c \end{array}
\end{equation*}

\par\medskip\noindent\textbf{Definition} (\texttt{mat-power}).\;$\operatorname{mat\text{-}power}(k, n, \mathit{tm}, j) :=$
\begin{equation*}
\operatorname{mpow}(\operatorname{ring\text{-}multiplicative\text{-}monoid}(\operatorname{mat\text{-}ring}(\operatorname{normed\text{-}field\text{-}as\text{-}commutative\text{-}ring}(k), n)), \mathit{tm}, j)
\end{equation*}

\par\medskip\noindent\textbf{Definition} (\texttt{is-power-bounded-by}).\;For all $k$, $n$, $\mathit{tm}$ and $c$, 
\begin{equation*}
\begin{array}{@{}l@{}} \operatorname{is\text{-}power\text{-}bounded\text{-}by}(k, n, \mathit{tm}, c) \\ \Leftrightarrow\; \mathit{tm} \in \operatorname{mat}(n, n, \operatorname{carr}(k)) \wedge \\ c \in \mathbb{R} \wedge \\ 0 \leq c \wedge \\ \forall \operatorname{cbvk\_} \in \mathbb{N}.\; \operatorname{is\text{-}mat\text{-}bound}(k, n, n, \operatorname{mat\text{-}power}(k, n, \mathit{tm}, \operatorname{cbvk\_}), c) \end{array}
\end{equation*}

\par\medskip\noindent\textbf{Definition} (\texttt{is-power-bounded}).\;For all $k$, $n$ and $\mathit{tm}$, 
\begin{equation*}
\operatorname{is\text{-}power\text{-}bounded}(k, n, \mathit{tm}) \Leftrightarrow \exists \operatorname{cbvc\_}.\; \operatorname{is\text{-}power\text{-}bounded\text{-}by}(k, n, \mathit{tm}, \operatorname{cbvc\_})
\end{equation*}


\vnbbookfile{structure-library/ccr-ch3b-vocab}{structure-library/ccr-ch3b-vocab}
\noindent{\small ccr-ch3b-vocab.scm -- wave of 2026-11-09/10 (agent CCR-1109-B2), wired 2026-11-10 from scratchpad/ccr1109b2/ccr-ch3b-vocab.scm. Library: analysis.}\par
\par\medskip\noindent\textbf{Definition} (\texttt{col-of}).\;$\operatorname{col\text{-}of}(q, n, j) :=$
\begin{equation*}
\begin{array}{@{}l@{}} \operatorname{matof}(n, 1, \\ \quad (\lambda {[}\operatorname{zkvi\_}, \operatorname{zkvu\_}] \in (\operatorname{interval}(1, n) \times \operatorname{interval}(1, 1)).\; \operatorname{entry}(q, \operatorname{zkvi\_}, j))) \end{array}
\end{equation*}

\par\medskip\noindent\textbf{Definition} (\texttt{krylov-mat}).\;$\operatorname{krylov\text{-}mat}(k, n, \mathit{tm}, \varphi) :=$
\begin{equation*}
\begin{array}{@{}l@{}} \operatorname{matof}(n, n, \\ \quad (\lambda {[}\operatorname{zkvi\_}, \operatorname{zkvj\_}] \in (\operatorname{interval}(1, n) \times \operatorname{interval}(1, n)).\; \\ \quad \quad \operatorname{entry}(\operatorname{matmul}(\operatorname{normed\text{-}field\text{-}as\text{-}commutative\text{-}ring}(k), \operatorname{mat\text{-}power}(k, n, \mathit{tm}, (\operatorname{zkvj\_} - 1)), \varphi), \operatorname{zkvi\_}, 1))) \end{array}
\end{equation*}

\par\medskip\noindent\textbf{Definition} (\texttt{is-irreducible-mat}).\;For all $k$, $n$ and $\mathit{tm}$, 
\begin{equation*}
\begin{array}{@{}l@{}} \operatorname{is\text{-}irreducible\text{-}mat}(k, n, \mathit{tm}) \\ \Leftrightarrow\; \mathit{tm} \in \operatorname{mat}(n, n, \operatorname{carr}(k)) \wedge \\ \forall \operatorname{zkvp\_} \in \operatorname{mat}(n, 1, \operatorname{carr}(k)). \\ \quad \neg (\operatorname{zkvp\_} = \operatorname{zeromat}(\operatorname{normed\text{-}field\text{-}as\text{-}commutative\text{-}ring}(k), n, 1)) \Rightarrow  \\ \quad \quad \forall \operatorname{zkvx\_} \in \operatorname{mat}(n, 1, \operatorname{carr}(k)). \exists \operatorname{zkvc\_} \in \operatorname{mat}(n, 1, \operatorname{carr}(k)). \\ \quad \quad \quad \operatorname{zkvx\_} = \operatorname{matmul}(\operatorname{normed\text{-}field\text{-}as\text{-}commutative\text{-}ring}(k), \operatorname{krylov\text{-}mat}(k, n, \mathit{tm}, \operatorname{zkvp\_}), \operatorname{zkvc\_}) \end{array}
\end{equation*}

\par\medskip\noindent\textbf{Definition} (\texttt{zpoly-alg}).\;$\operatorname{zpoly\text{-}alg}(k, n, \mathit{tm}) :=$
\begin{equation*}
\operatorname{intersection\text{-}of}(\{\operatorname{zkvs\_} \in \operatorname{power}(\operatorname{mat}(n, n, \operatorname{carr}(k))) : ((\operatorname{identmat}(\operatorname{normed\text{-}field\text{-}as\text{-}commutative\text{-}ring}(k), n) \in \operatorname{zkvs\_}) \wedge ((\operatorname{zeromat}(\operatorname{normed\text{-}field\text{-}as\text{-}commutative\text{-}ring}(k), n, n) \in \operatorname{zkvs\_}) \wedge (\forall \operatorname{zkvx\_} \in \operatorname{zkvs\_}.\; (\operatorname{matmul}(\operatorname{normed\text{-}field\text{-}as\text{-}commutative\text{-}ring}(k), \operatorname{zkvx\_}, \mathit{tm}) \in \operatorname{zkvs\_}) \wedge (\forall \operatorname{zkvx\_}, \operatorname{zkvy\_} \in \operatorname{zkvs\_}.\; (\operatorname{matadd}(\operatorname{normed\text{-}field\text{-}as\text{-}commutative\text{-}ring}(k), \operatorname{zkvx\_}, \operatorname{zkvy\_}) \in \operatorname{zkvs\_}) \wedge \forall \operatorname{zkvx\_} \in \operatorname{zkvs\_}, \operatorname{zkvc\_} \in \operatorname{nv\text{-}integers}(k).\; (\operatorname{matscale}(\operatorname{normed\text{-}field\text{-}as\text{-}commutative\text{-}ring}(k), \operatorname{zkvc\_}, \operatorname{zkvx\_}) \in \operatorname{zkvs\_})))))\})
\end{equation*}

\par\medskip\noindent\textbf{Definition} (\texttt{mat-poly-eval}).\;$\operatorname{mat\text{-}poly\text{-}eval}(k, n, p, \mathit{tm}) :=$
\begin{equation*}
\begin{array}{@{}l@{}} \operatorname{sum\text{-}ag}(\operatorname{ring\text{-}additive\text{-}ag}(\operatorname{mat\text{-}ring}(\operatorname{normed\text{-}field\text{-}as\text{-}commutative\text{-}ring}(k), n)), \\ \quad (\lambda \operatorname{zkvk\_} \in \mathbb{N}.\; \\ \quad \quad \operatorname{matscale}(\operatorname{normed\text{-}field\text{-}as\text{-}commutative\text{-}ring}(k), p(\operatorname{zkvk\_}), \operatorname{mat\text{-}power}(k, n, \mathit{tm}, \operatorname{zkvk\_}))), \\ \quad (\operatorname{deg}(\operatorname{normed\text{-}field\text{-}as\text{-}commutative\text{-}ring}(k), p) + 1)) \end{array}
\end{equation*}

\par\medskip\noindent\textbf{Definition} (\texttt{is-norm-on}).\;For all $k$, $\mathit{md}$ and $\nu$, 
\begin{equation*}
\begin{array}{@{}l@{}} \operatorname{is\text{-}norm\text{-}on}(k, \mathit{md}, \nu) \\ \Leftrightarrow\; \operatorname{is\text{-}module}(\mathit{md}) \wedge \\ \operatorname{scal}(\mathit{md}) = \operatorname{normed\text{-}field\text{-}as\text{-}commutative\text{-}ring}(k) \wedge \\ \nu \in (\operatorname{vec}(\mathit{md}) \to \mathbb{R}) \wedge \\ \forall \operatorname{zkvx\_} \in \operatorname{vec}(\mathit{md}).\; (0 \leq \nu(\operatorname{zkvx\_})) \wedge \\ \forall \operatorname{zkvx\_} \in \operatorname{vec}(\mathit{md}).\; ((\nu(\operatorname{zkvx\_}) = 0) \Leftrightarrow (\operatorname{zkvx\_} = \operatorname{vzero}(\mathit{md}))) \wedge \\ \forall \operatorname{zkvx\_}, \operatorname{zkvy\_} \in \operatorname{vec}(\mathit{md}). \\ \quad \nu(\operatorname{vadd}(\mathit{md})(\operatorname{zkvx\_}, \operatorname{zkvy\_})) \leq +(\nu(\operatorname{zkvx\_}), \nu(\operatorname{zkvy\_})) \wedge \\ \forall \operatorname{zkvr\_} \in \operatorname{carr}(k), \operatorname{zkvx\_} \in \operatorname{vec}(\mathit{md}). \\ \quad \nu(\operatorname{act}(\mathit{md})(\operatorname{zkvr\_}, \operatorname{zkvx\_})) = \cdot(\operatorname{fnrm}(k)(\operatorname{zkvr\_}), \nu(\operatorname{zkvx\_})) \end{array}
\end{equation*}


\vnbbookfile{structure-library/ccr-ch3c-vocab}{structure-library/ccr-ch3c-vocab}
\noindent{\small ccr-ch3c-vocab.scm -- wave of 2026-11-09/10 (agent CCR-1109-B3), wired 2026-11-10 from scratchpad/ccr1109b3/ccr-ch3c-vocab.scm. Library: analysis.}\par
\par\medskip\noindent\textbf{Definition} (\texttt{is-mat-close}).\;For all $k$, $m$, $n$, $a$, $b$ and $e$, 
\begin{equation*}
\begin{array}{@{}l@{}} \operatorname{is\text{-}mat\text{-}close}(k, m, n, a, b, e) \\ \Leftrightarrow\; a \in \operatorname{mat}(m, n, \operatorname{carr}(k)) \wedge \\ b \in \operatorname{mat}(m, n, \operatorname{carr}(k)) \wedge \\ \forall \operatorname{cdvi\_} \in \operatorname{interval}(1, m), \operatorname{cdvj\_} \in \operatorname{interval}(1, n). \\ \quad \operatorname{fnrm}(k)(\operatorname{add}(k)(\operatorname{entry}(a, \operatorname{cdvi\_}, \operatorname{cdvj\_}), \operatorname{neg}(k)(\operatorname{entry}(b, \operatorname{cdvi\_}, \operatorname{cdvj\_})))) \leq e \end{array}
\end{equation*}


\vnbbookfile{structure-library/nvlocal-vocab}{structure-library/nvlocal-vocab}
\noindent{\small nvlocal-vocab.scm -- WAVE 2 (2026-10-27..30), wired 2026-10-30 from scratchpad/ccrqp/nvlocal-vocab.scm. Library: analysis.}\par
\par\medskip\noindent\textbf{Definition} (\texttt{is-discrete-ms}).\;For every $\operatorname{dms\_}$, 
\begin{equation*}
\begin{array}{@{}l@{}} \operatorname{is\text{-}discrete\text{-}ms}(\operatorname{dms\_}) \\ \Leftrightarrow\; \forall \operatorname{dmx\_} \in \operatorname{pts}(\operatorname{dms\_}). \exists \operatorname{dmr\_}. \\ \quad \operatorname{pos\text{-}rr}(\operatorname{dmr\_}) \wedge \forall \operatorname{dmy\_} \in \operatorname{ball}(\operatorname{dms\_}, \operatorname{dmx\_}, \operatorname{dmr\_}).\; (\operatorname{dmy\_} = \operatorname{dmx\_}) \end{array}
\end{equation*}

\par\medskip\noindent\textbf{Definition} (\texttt{is-local-nv-field}).\;For every $\operatorname{lck\_}$, 
\begin{equation*}
\begin{array}{@{}l@{}} \operatorname{is\text{-}local\text{-}nv\text{-}field}(\operatorname{lck\_}) \\ \Leftrightarrow\; \operatorname{is\text{-}nv\text{-}field}(\operatorname{lck\_}) \wedge \\ \neg \operatorname{is\text{-}discrete\text{-}ms}(\operatorname{nf\text{-}metric\text{-}space}(\operatorname{lck\_})) \wedge \\ \operatorname{is\text{-}locally\text{-}compact}(\operatorname{nf\text{-}metric\text{-}space}(\operatorname{lck\_})) \end{array}
\end{equation*}

\par\medskip\noindent\textbf{Definition} (\texttt{nv-res-ringoid}).\;$\operatorname{nv\text{-}res\text{-}ringoid}(\operatorname{rsk\_}, \operatorname{rsa\_}) :=$
\begin{equation*}
{[}\operatorname{rsa\_}, (\lambda {[}\operatorname{rsu\_}, \operatorname{rsv\_}] \in (\operatorname{rsa\_} \times \operatorname{rsa\_}).\; \operatorname{add}(\operatorname{rsk\_})(\operatorname{rsu\_}, \operatorname{rsv\_})), (\lambda {[}\operatorname{rsu\_}, \operatorname{rsv\_}] \in (\operatorname{rsa\_} \times \operatorname{rsa\_}).\; \operatorname{mul}(\operatorname{rsk\_})(\operatorname{rsu\_}, \operatorname{rsv\_})), (\lambda \operatorname{rsu\_} \in \operatorname{rsa\_}.\; \operatorname{neg}(\operatorname{rsk\_})(\operatorname{rsu\_})), \operatorname{zero}(\operatorname{rsk\_}), \operatorname{one}(\operatorname{rsk\_}), \operatorname{nv\text{-}maximal}(\operatorname{rsk\_}, \operatorname{rsa\_})]
\end{equation*}

\par\medskip\noindent\textbf{Definition} (\texttt{nv-residue}).\;$\operatorname{nv\text{-}residue}(\operatorname{rsk\_}, \operatorname{rsa\_}) :=$ $\operatorname{ringoid\text{-}quotient\text{-}ring}(\operatorname{nv\text{-}res\text{-}ringoid}(\operatorname{rsk\_}, \operatorname{rsa\_}))$
\par\medskip\noindent\textbf{Definition} (\texttt{is-ring-iso}).\;For all $\operatorname{rib\_}$, $\operatorname{ric\_}$ and $\operatorname{rif\_}$, 
\begin{equation*}
\begin{array}{@{}l@{}} \operatorname{is\text{-}ring\text{-}iso}(\operatorname{rib\_}, \operatorname{ric\_}, \operatorname{rif\_}) \\ \Leftrightarrow\; \operatorname{is\text{-}hom\text{-}ring}(\operatorname{rib\_}, \operatorname{ric\_}, \operatorname{rif\_}) \wedge \\ \forall \operatorname{riu\_}, \operatorname{riv\_} \in \operatorname{carr}(\operatorname{rib\_}). \\ \quad (\operatorname{rif\_}(\operatorname{riu\_}) = \operatorname{rif\_}(\operatorname{riv\_})) \Rightarrow (\operatorname{riu\_} = \operatorname{riv\_}) \wedge \\ \forall \operatorname{riw\_} \in \operatorname{carr}(\operatorname{ric\_}).\; \exists \operatorname{riz\_} \in \operatorname{carr}(\operatorname{rib\_}).\; (\operatorname{rif\_}(\operatorname{riz\_}) = \operatorname{riw\_}) \end{array}
\end{equation*}

\par\medskip\noindent\textbf{Definition} (\texttt{is-add-subgroup}).\;For all $\operatorname{asgk\_}$ and $\operatorname{asgu\_}$, 
\begin{equation*}
\begin{array}{@{}l@{}} \operatorname{is\text{-}add\text{-}subgroup}(\operatorname{asgk\_}, \operatorname{asgu\_}) \\ \Leftrightarrow\; \forall \operatorname{asx\_} \in \operatorname{asgu\_}.\; (\operatorname{asx\_} \in \operatorname{carr}(\operatorname{asgk\_})) \wedge \\ \operatorname{zero}(\operatorname{asgk\_}) \in \operatorname{asgu\_} \wedge \\ \forall \operatorname{asx\_}, \operatorname{asy\_} \in \operatorname{asgu\_}.\; (\operatorname{add}(\operatorname{asgk\_})(\operatorname{asx\_}, \operatorname{asy\_}) \in \operatorname{asgu\_}) \wedge \\ \forall \operatorname{asx\_} \in \operatorname{asgu\_}.\; (\operatorname{neg}(\operatorname{asgk\_})(\operatorname{asx\_}) \in \operatorname{asgu\_}) \end{array}
\end{equation*}


\vnbbookfile{structure-library/ccr4x-ext-vocab}{structure-library/ccr4x-ext-vocab}
\noindent{\small ccr4x-ext-vocab.scm -- wave of 2026-11-09/10 (agent CCR-1109-C2), wired 2026-11-10 from scratchpad/ccr1109c2/ccr4x-ext-vocab.scm. Library: analysis.}\par
\par\medskip\noindent\textbf{Definition} (\texttt{op-mat}).\;$\operatorname{op\text{-}mat}(\operatorname{x4md\_}, \operatorname{x4n\_}, \operatorname{x4u\_}, \operatorname{x4t\_}) :=$
\begin{equation*}
(\iota\, \operatorname{x4a\_}.\;((\operatorname{x4a\_} \in \operatorname{mat}(\operatorname{x4n\_}, \operatorname{x4n\_}, \operatorname{carr}(\operatorname{scal}(\operatorname{x4md\_})))) \wedge (\operatorname{map\text{-}seq}(\operatorname{x4t\_}, \operatorname{x4n\_}, \operatorname{x4u\_}) = \operatorname{matact}(\operatorname{x4md\_}, \operatorname{x4a\_}, \operatorname{x4u\_}))))
\end{equation*}

\par\medskip\noindent\textbf{Definition} (\texttt{is-fin-ext}).\;For all $\operatorname{x4k\_}$, $\operatorname{x4f\_}$, $\operatorname{x4md\_}$, $\operatorname{x4n\_}$ and $\operatorname{x4u\_}$, 
\begin{equation*}
\begin{array}{@{}l@{}} \operatorname{is\text{-}fin\text{-}ext}(\operatorname{x4k\_}, \operatorname{x4f\_}, \operatorname{x4md\_}, \operatorname{x4n\_}, \operatorname{x4u\_}) \\ \Leftrightarrow\; \operatorname{is\text{-}normed\text{-}field}(\operatorname{x4k\_}) \wedge \\ \operatorname{is\text{-}field\text{-}ring}(\operatorname{x4f\_}) \wedge \\ \operatorname{is\text{-}module}(\operatorname{x4md\_}) \wedge \\ \operatorname{scal}(\operatorname{x4md\_}) = \operatorname{normed\text{-}field\text{-}as\text{-}commutative\text{-}ring}(\operatorname{x4k\_}) \wedge \\ \operatorname{vec}(\operatorname{x4md\_}) = \operatorname{carr}(\operatorname{x4f\_}) \wedge \\ \operatorname{vadd}(\operatorname{x4md\_}) = \operatorname{add}(\operatorname{x4f\_}) \wedge \\ \operatorname{vzero}(\operatorname{x4md\_}) = \operatorname{zero}(\operatorname{x4f\_}) \wedge \\ \operatorname{vneg}(\operatorname{x4md\_}) = \operatorname{neg}(\operatorname{x4f\_}) \wedge \\ \forall \operatorname{x4r\_} \in \operatorname{carr}(\operatorname{x4k\_}), \operatorname{x4x\_}, \operatorname{x4y\_} \in \operatorname{carr}(\operatorname{x4f\_}). \\ \quad \operatorname{mul}(\operatorname{x4f\_})(\operatorname{x4x\_}, \operatorname{act}(\operatorname{x4md\_})(\operatorname{x4r\_}, \operatorname{x4y\_})) = \operatorname{act}(\operatorname{x4md\_})(\operatorname{x4r\_}, \operatorname{mul}(\operatorname{x4f\_})(\operatorname{x4x\_}, \operatorname{x4y\_})) \wedge \\ \operatorname{x4n\_} \in \mathbb{N} \wedge \\ 1 \leq \operatorname{x4n\_} \wedge \\ \operatorname{x4u\_} \in \operatorname{mat}(\operatorname{x4n\_}, 1, \operatorname{carr}(\operatorname{x4f\_})) \wedge \\ \operatorname{generates}(\operatorname{x4md\_}, \operatorname{x4n\_}, \operatorname{x4u\_}) \wedge \\ \operatorname{rel\text{-}free}(\operatorname{x4md\_}, \operatorname{x4n\_}, \operatorname{x4u\_}) \end{array}
\end{equation*}

\par\medskip\noindent\textbf{Definition} (\texttt{ext-degree}).\;$\operatorname{ext\text{-}degree}(\operatorname{x4k\_}, \operatorname{x4f\_}, \operatorname{x4md\_}) :=$
\begin{equation*}
(\iota\, \operatorname{x4d\_}.\;\exists \operatorname{x4b\_}.\; \operatorname{is\text{-}fin\text{-}ext}(\operatorname{x4k\_}, \operatorname{x4f\_}, \operatorname{x4md\_}, \operatorname{x4d\_}, \operatorname{x4b\_}))
\end{equation*}

\par\medskip\noindent\textbf{Definition} (\texttt{mult-op}).\;$\operatorname{mult\text{-}op}(\operatorname{x4f\_}, \operatorname{x4xi\_}) :=$ $(\lambda \operatorname{x4p\_} \in \operatorname{carr}(\operatorname{x4f\_}).\; \operatorname{mul}(\operatorname{x4f\_})(\operatorname{x4xi\_}, \operatorname{x4p\_}))$
\par\medskip\noindent\textbf{Definition} (\texttt{mult-mat}).\;$\operatorname{mult\text{-}mat}(\operatorname{x4md\_}, \operatorname{x4n\_}, \operatorname{x4u\_}, \operatorname{x4f\_}, \operatorname{x4xi\_}) :=$ $\operatorname{op\text{-}mat}(\operatorname{x4md\_}, \operatorname{x4n\_}, \operatorname{x4u\_}, \operatorname{mult\text{-}op}(\operatorname{x4f\_}, \operatorname{x4xi\_}))$
\par\medskip\noindent\textbf{Definition} (\texttt{ext-norm}).\;$\operatorname{ext\text{-}norm}(\operatorname{x4k\_}, \operatorname{x4f\_}, \operatorname{x4md\_}, \operatorname{x4n\_}, \operatorname{x4u\_}, \operatorname{x4xi\_}) :=$
\begin{equation*}
\operatorname{det}(\operatorname{normed\text{-}field\text{-}as\text{-}commutative\text{-}ring}(\operatorname{x4k\_}), \operatorname{x4n\_}, \operatorname{mult\text{-}mat}(\operatorname{x4md\_}, \operatorname{x4n\_}, \operatorname{x4u\_}, \operatorname{x4f\_}, \operatorname{x4xi\_}))
\end{equation*}

\par\medskip\noindent\textbf{Definition} (\texttt{ext-abs}).\;$\operatorname{ext\text{-}abs}(\operatorname{x4k\_}, \operatorname{x4f\_}, \operatorname{x4md\_}, \operatorname{x4n\_}, \operatorname{x4u\_}, \operatorname{x4xi\_}) :=$ $\operatorname{fnrm}(\operatorname{x4k\_})(\operatorname{ext\text{-}norm}(\operatorname{x4k\_}, \operatorname{x4f\_}, \operatorname{x4md\_}, \operatorname{x4n\_}, \operatorname{x4u\_}, \operatorname{x4xi\_}))$
\par\medskip\noindent\textbf{Definition} (\texttt{ext-abs-root}).\;$\operatorname{ext\text{-}abs\text{-}root}(\operatorname{x4k\_}, \operatorname{x4f\_}, \operatorname{x4md\_}, \operatorname{x4n\_}, \operatorname{x4u\_}, \operatorname{x4xi\_}) :=$
\begin{equation*}
{\operatorname{ext\text{-}abs}(\operatorname{x4k\_}, \operatorname{x4f\_}, \operatorname{x4md\_}, \operatorname{x4n\_}, \operatorname{x4u\_}, \operatorname{x4xi\_})}^{{\operatorname{x4n\_}}^{\ast\!-1}}
\end{equation*}

\par\medskip\noindent\textbf{Definition} (\texttt{is-nv-valuation}).\;For all $\operatorname{x4f\_}$ and $\operatorname{x4nu\_}$, 
\begin{equation*}
\begin{array}{@{}l@{}} \operatorname{is\text{-}nv\text{-}valuation}(\operatorname{x4f\_}, \operatorname{x4nu\_}) \\ \Leftrightarrow\; \operatorname{x4nu\_} \in (\operatorname{carr}(\operatorname{x4f\_}) \to \mathbb{R}) \wedge \\ \forall \operatorname{x4x\_} \in \operatorname{carr}(\operatorname{x4f\_}).\; (0 \leq \operatorname{x4nu\_}(\operatorname{x4x\_})) \wedge \\ \forall \operatorname{x4x\_} \in \operatorname{carr}(\operatorname{x4f\_}).\; ((\operatorname{x4nu\_}(\operatorname{x4x\_}) = 0) \Leftrightarrow (\operatorname{x4x\_} = \operatorname{zero}(\operatorname{x4f\_}))) \wedge \\ \forall \operatorname{x4x\_}, \operatorname{x4y\_} \in \operatorname{carr}(\operatorname{x4f\_}). \\ \quad \operatorname{x4nu\_}(\operatorname{mul}(\operatorname{x4f\_})(\operatorname{x4x\_}, \operatorname{x4y\_})) = \cdot(\operatorname{x4nu\_}(\operatorname{x4x\_}), \operatorname{x4nu\_}(\operatorname{x4y\_})) \wedge \\ \forall \operatorname{x4x\_}, \operatorname{x4y\_} \in \operatorname{carr}(\operatorname{x4f\_}). \\ \quad \operatorname{x4nu\_}(\operatorname{add}(\operatorname{x4f\_})(\operatorname{x4x\_}, \operatorname{x4y\_})) \leq \operatorname{max}(\operatorname{x4nu\_}(\operatorname{x4x\_}), \operatorname{x4nu\_}(\operatorname{x4y\_})) \end{array}
\end{equation*}


\vnbbookfile{structure-library/ccr4x-vocab}{structure-library/ccr4x-vocab}
\noindent{\small ccr4x-vocab.scm -- wave of 2026-11-09/10 (agent CCR-1109-C), wired 2026-11-10 from scratchpad/ccr1109c/ccr4x-vocab.scm. Library: analysis.}\par
\par\medskip\noindent\textbf{Definition} (\texttt{nv-zpow}).\;$\operatorname{nv\text{-}zpow}(\operatorname{zpk\_}, \operatorname{zpw\_}, \operatorname{zpn\_}) :=$
\begin{equation*}
(\iota\, \operatorname{zpy\_}.\;((\operatorname{zpy\_} \in \operatorname{carr}(\operatorname{zpk\_})) \wedge (\operatorname{mul}(\operatorname{zpk\_})(\operatorname{zpy\_}, \operatorname{ring\text{-}power}(\operatorname{normed\text{-}field\text{-}as\text{-}commutative\text{-}ring}(\operatorname{zpk\_}), \operatorname{zpw\_}, \operatorname{max}(-\operatorname{zpn\_}, 0))) = \operatorname{ring\text{-}power}(\operatorname{normed\text{-}field\text{-}as\text{-}commutative\text{-}ring}(\operatorname{zpk\_}), \operatorname{zpw\_}, \operatorname{max}(\operatorname{zpn\_}, 0)))))
\end{equation*}

\par\medskip\noindent\textbf{Definition} (\texttt{is-nv-rep-set}).\;For all $\operatorname{rpk\_}$, $\operatorname{rpa\_}$ and $\operatorname{rpr\_}$, 
\begin{equation*}
\begin{array}{@{}l@{}} \operatorname{is\text{-}nv\text{-}rep\text{-}set}(\operatorname{rpk\_}, \operatorname{rpa\_}, \operatorname{rpr\_}) \\ \Leftrightarrow\; \forall \operatorname{rpx\_} \in \operatorname{rpr\_}.\; (\operatorname{rpx\_} \in \operatorname{rpa\_}) \wedge \\ \operatorname{zero}(\operatorname{rpk\_}) \in \operatorname{rpr\_} \wedge \\ \operatorname{one}(\operatorname{rpk\_}) \in \operatorname{rpr\_} \wedge \\ \forall \operatorname{rpx\_} \in \operatorname{rpa\_}. \exists \operatorname{rpy\_} \in \operatorname{rpr\_}. \\ \quad \operatorname{fnrm}(\operatorname{rpk\_})(\operatorname{add}(\operatorname{rpk\_})(\operatorname{rpx\_}, \operatorname{neg}(\operatorname{rpk\_})(\operatorname{rpy\_}))) < 1 \wedge \\ \forall \operatorname{rpx\_}, \operatorname{rpy\_} \in \operatorname{rpr\_}. \\ \quad (\operatorname{fnrm}(\operatorname{rpk\_})(\operatorname{add}(\operatorname{rpk\_})(\operatorname{rpx\_}, \operatorname{neg}(\operatorname{rpk\_})(\operatorname{rpy\_}))) < 1) \Rightarrow (\operatorname{rpx\_} = \operatorname{rpy\_}) \end{array}
\end{equation*}

\par\medskip\noindent\textbf{Definition} (\texttt{nv-digit}).\;$\operatorname{nv\text{-}digit}(\operatorname{dgk\_}, \operatorname{dgr\_}, \operatorname{dgz\_}) :=$
\begin{equation*}
(\iota\, \operatorname{dgy\_}.\;((\operatorname{dgy\_} \in \operatorname{dgr\_}) \wedge (\operatorname{fnrm}(\operatorname{dgk\_})(\operatorname{add}(\operatorname{dgk\_})(\operatorname{dgz\_}, \operatorname{neg}(\operatorname{dgk\_})(\operatorname{dgy\_}))) < 1)))
\end{equation*}

\par\medskip\noindent\textbf{Definition} (\texttt{nv-term}).\;$\operatorname{nv\text{-}term}(\operatorname{ntk\_}, \operatorname{ntw\_}, \operatorname{nta\_}, \operatorname{ntn\_}) :=$
\begin{equation*}
(\lambda \operatorname{ntj\_} \in \mathbb{N}.\; \operatorname{mul}(\operatorname{ntk\_})(\operatorname{nta\_}((\operatorname{ntn\_} + \operatorname{ntj\_})), \operatorname{nv\text{-}zpow}(\operatorname{ntk\_}, \operatorname{ntw\_}, (\operatorname{ntn\_} + \operatorname{ntj\_}))))
\end{equation*}

\par\medskip\noindent\textbf{Definition} (\texttt{nv-psum}).\;$\operatorname{nv\text{-}psum}(\operatorname{ntk\_}, \operatorname{ntw\_}, \operatorname{nta\_}, \operatorname{ntn\_}, \operatorname{ntm\_}) :=$
\begin{equation*}
\operatorname{sum\text{-}ag}(\operatorname{normed\text{-}field\text{-}additive\text{-}ag}(\operatorname{ntk\_}), \operatorname{nv\text{-}term}(\operatorname{ntk\_}, \operatorname{ntw\_}, \operatorname{nta\_}, \operatorname{ntn\_}), \operatorname{ntm\_})
\end{equation*}

\par\medskip\noindent\textbf{Definition} (\texttt{is-nv-expansion}).\;For all $\operatorname{nek\_}$, $\operatorname{new\_}$, $\operatorname{ner\_}$, $\operatorname{nea\_}$, $\operatorname{nen\_}$ and $\operatorname{nex\_}$, 
\begin{equation*}
\begin{array}{@{}l@{}} \operatorname{is\text{-}nv\text{-}expansion}(\operatorname{nek\_}, \operatorname{new\_}, \operatorname{ner\_}, \operatorname{nea\_}, \operatorname{nen\_}, \operatorname{nex\_}) \\ \Leftrightarrow\; \operatorname{nea\_} \in (\mathbb{Z} \to \operatorname{ner\_}) \wedge \\ \forall \operatorname{nei\_} \in \mathbb{Z}.\; ((\operatorname{nei\_} < \operatorname{nen\_}) \Rightarrow (\operatorname{nea\_}(\operatorname{nei\_}) = \operatorname{zero}(\operatorname{nek\_}))) \wedge \\ \operatorname{converges\text{-}to}(\operatorname{nf\text{-}metric\text{-}space}(\operatorname{nek\_}), \\ \quad (\lambda \operatorname{nem\_} \in \mathbb{N}.\; \operatorname{nv\text{-}psum}(\operatorname{nek\_}, \operatorname{new\_}, \operatorname{nea\_}, \operatorname{nen\_}, \operatorname{nem\_})), \\ \quad \operatorname{nex\_}) \end{array}
\end{equation*}


\vnbbookfile{structure-library/ccr4b-vocab}{structure-library/ccr4b-vocab}
\noindent{\small ccr4b-vocab.scm -- wave of 2026-11-09/10 (agent CCR-1109-C3), wired 2026-11-10 from scratchpad/ccr1109c3/ccr4b-vocab.scm. Library: analysis.}\par
\par\medskip\noindent\textbf{Definition} (\texttt{is-nf-subfield}).\;For all $\operatorname{bzk\_}$ and $\operatorname{bzs\_}$, 
\begin{equation*}
\begin{array}{@{}l@{}} \operatorname{is\text{-}nf\text{-}subfield}(\operatorname{bzk\_}, \operatorname{bzs\_}) \\ \Leftrightarrow\; \forall \operatorname{bzx\_} \in \operatorname{bzs\_}.\; (\operatorname{bzx\_} \in \operatorname{carr}(\operatorname{bzk\_})) \wedge \\ \operatorname{zero}(\operatorname{bzk\_}) \in \operatorname{bzs\_} \wedge \\ \operatorname{one}(\operatorname{bzk\_}) \in \operatorname{bzs\_} \wedge \\ \forall \operatorname{bzx\_}, \operatorname{bzy\_} \in \operatorname{bzs\_}.\; (\operatorname{add}(\operatorname{bzk\_})(\operatorname{bzx\_}, \operatorname{bzy\_}) \in \operatorname{bzs\_}) \wedge \\ \forall \operatorname{bzx\_}, \operatorname{bzy\_} \in \operatorname{bzs\_}.\; (\operatorname{mul}(\operatorname{bzk\_})(\operatorname{bzx\_}, \operatorname{bzy\_}) \in \operatorname{bzs\_}) \wedge \\ \forall \operatorname{bzx\_} \in \operatorname{bzs\_}.\; (\operatorname{neg}(\operatorname{bzk\_})(\operatorname{bzx\_}) \in \operatorname{bzs\_}) \wedge \\ \forall \operatorname{bzx\_} \in \operatorname{bzs\_}. \\ \quad \neg (\operatorname{bzx\_} = \operatorname{zero}(\operatorname{bzk\_})) \Rightarrow  \\ \quad \quad \exists \operatorname{bzy\_} \in \operatorname{bzs\_}.\; (\operatorname{mul}(\operatorname{bzk\_})(\operatorname{bzx\_}, \operatorname{bzy\_}) = \operatorname{one}(\operatorname{bzk\_})) \end{array}
\end{equation*}

\par\medskip\noindent\textbf{Definition} (\texttt{nf-subfield-gen}).\;$\operatorname{nf\text{-}subfield\text{-}gen}(\operatorname{bzk\_}, \operatorname{bzg\_}) :=$
\begin{equation*}
\{\operatorname{bzz\_} \in \operatorname{carr}(\operatorname{bzk\_}) : \forall \operatorname{bzs\_} \in \operatorname{power}(\operatorname{carr}(\operatorname{bzk\_})).\; (\operatorname{is\text{-}nf\text{-}subfield}(\operatorname{bzk\_}, \operatorname{bzs\_}) \Rightarrow ((\operatorname{bzg\_} \subseteq \operatorname{bzs\_}) \Rightarrow (\operatorname{bzz\_} \in \operatorname{bzs\_})))\}
\end{equation*}

\par\medskip\noindent\textbf{Definition} (\texttt{nv-base-gen}).\;$\operatorname{nv\text{-}base\text{-}gen}(\operatorname{bzk\_}, \operatorname{bzw\_}) :=$
\begin{equation*}
\{\operatorname{bzz\_} \in \operatorname{carr}(\operatorname{bzk\_}) : ((\operatorname{bzz\_} = \operatorname{bzw\_}) \wedge \exists \operatorname{bzp\_}.\; \operatorname{has\text{-}characteristic}(\operatorname{normed\text{-}field\text{-}as\text{-}commutative\text{-}ring}(\operatorname{bzk\_}), \operatorname{bzp\_}))\}
\end{equation*}

\par\medskip\noindent\textbf{Definition} (\texttt{nv-base}).\;$\operatorname{nv\text{-}base}(\operatorname{bzk\_}, \operatorname{bzw\_}) :=$
\begin{equation*}
\operatorname{closure}(\operatorname{nf\text{-}metric\text{-}space}(\operatorname{bzk\_}), \operatorname{nf\text{-}subfield\text{-}gen}(\operatorname{bzk\_}, \operatorname{nv\text{-}base\text{-}gen}(\operatorname{bzk\_}, \operatorname{bzw\_})))
\end{equation*}

\par\medskip\noindent\textbf{Definition} (\texttt{nv-zf}).\;$\operatorname{nv\text{-}zf}(\operatorname{bzk\_}, \operatorname{bzf\_}) :=$
\begin{equation*}
\{\operatorname{bzz\_} \in \operatorname{nv\text{-}integers}(\operatorname{bzk\_}) : (\operatorname{class}(\operatorname{ringoid\text{-}setoid}(\operatorname{nv\text{-}res\text{-}ringoid}(\operatorname{bzk\_}, \operatorname{nv\text{-}integers}(\operatorname{bzk\_}))), \operatorname{bzz\_}) \in \operatorname{bzf\_})\}
\end{equation*}


\vnbbookfile{structure-library/ccr3-vocab}{structure-library/ccr3-vocab}
\noindent{\small ccr3-vocab.scm -- MODULES IN A LOCAL NV FIELD (CCR.pdf Definition 4.11).  DEFINITIONS ONLY; the laws are in theorem-library/ccr3-local.scm.  Agent CCR-B3, wave MAP-1026, 2026-10-26.  Destination: structure-library/ccr3-vocab.scm.  Library: analysis.  Load window: after structure-library/ccr2-vocab.}\par
\par\medskip\noindent\textbf{Definition} (\texttt{is-nv-module}).\;For all $\operatorname{cw3k\_}$, $\operatorname{cw3a\_}$ and $\operatorname{cw3e\_}$, 
\begin{equation*}
\begin{array}{@{}l@{}} \operatorname{is\text{-}nv\text{-}module}(\operatorname{cw3k\_}, \operatorname{cw3a\_}, \operatorname{cw3e\_}) \\ \Leftrightarrow\; \forall \operatorname{cw3x\_} \in \operatorname{cw3e\_}.\; (\operatorname{cw3x\_} \in \operatorname{carr}(\operatorname{cw3k\_})) \wedge \\ \operatorname{zero}(\operatorname{cw3k\_}) \in \operatorname{cw3e\_} \wedge \\ \forall \operatorname{cw3x\_}, \operatorname{cw3y\_} \in \operatorname{cw3e\_}.\; (\operatorname{add}(\operatorname{cw3k\_})(\operatorname{cw3x\_}, \operatorname{cw3y\_}) \in \operatorname{cw3e\_}) \wedge \\ \forall \operatorname{cw3x\_} \in \operatorname{cw3e\_}.\; (\operatorname{neg}(\operatorname{cw3k\_})(\operatorname{cw3x\_}) \in \operatorname{cw3e\_}) \wedge \\ \forall \operatorname{cw3r\_} \in \operatorname{cw3a\_}, \operatorname{cw3x\_} \in \operatorname{cw3e\_}.\; (\operatorname{mul}(\operatorname{cw3k\_})(\operatorname{cw3r\_}, \operatorname{cw3x\_}) \in \operatorname{cw3e\_}) \end{array}
\end{equation*}

\par\medskip\noindent\textbf{Definition} (\texttt{nv-multiples}).\;$\operatorname{nv\text{-}multiples}(\operatorname{cw3k\_}, \operatorname{cw3a\_}, \operatorname{cw3g\_}) :=$
\begin{equation*}
\{\operatorname{cw3x\_} \in \operatorname{carr}(\operatorname{cw3k\_}) : \exists \operatorname{cw3m\_} \in \operatorname{cw3a\_}.\; (\operatorname{cw3x\_} = \operatorname{mul}(\operatorname{cw3k\_})(\operatorname{cw3m\_}, \operatorname{cw3g\_}))\}
\end{equation*}

\par\medskip\noindent\textbf{Definition} (\texttt{is-nv-content-one}).\;For all $\operatorname{cw3k\_}$ and $\operatorname{cw3p\_}$, 
\begin{equation*}
\begin{array}{@{}l@{}} \operatorname{is\text{-}nv\text{-}content\text{-}one}(\operatorname{cw3k\_}, \operatorname{cw3p\_}) \\ \Leftrightarrow\; \forall \operatorname{cw3i\_} \in \mathbb{N}.\; (\operatorname{fnrm}(\operatorname{cw3k\_})(\operatorname{cw3p\_}(\operatorname{cw3i\_})) \leq 1) \wedge \\ \exists \operatorname{cw3i\_} \in \mathbb{N}.\; (\operatorname{fnrm}(\operatorname{cw3k\_})(\operatorname{cw3p\_}(\operatorname{cw3i\_})) = 1) \end{array}
\end{equation*}

\par\medskip\noindent\textbf{Definition} (\texttt{is-nv-content}).\;For all $\operatorname{cw3k\_}$, $\operatorname{cw3p\_}$ and $\operatorname{cw3c\_}$, 
\begin{equation*}
\begin{array}{@{}l@{}} \operatorname{is\text{-}nv\text{-}content}(\operatorname{cw3k\_}, \operatorname{cw3p\_}, \operatorname{cw3c\_}) \\ \Leftrightarrow\; \forall \operatorname{cw3i\_} \in \mathbb{N}.\; (\operatorname{fnrm}(\operatorname{cw3k\_})(\operatorname{cw3p\_}(\operatorname{cw3i\_})) \leq \operatorname{cw3c\_}) \wedge \\ \exists \operatorname{cw3i\_} \in \mathbb{N}.\; (\operatorname{fnrm}(\operatorname{cw3k\_})(\operatorname{cw3p\_}(\operatorname{cw3i\_})) = \operatorname{cw3c\_}) \end{array}
\end{equation*}


\vnbbookfile{theorem-library/ccr2-nv}{theorem-library/ccr2-nv}
\noindent{\small ccr2-nv.scm -- NON-ARCHIMEDEAN VALUED FIELDS AND DOMAINS: CCR.pdf Propositions 1.9, 1.17, 2.4, 2.6, 2.10, 2.11, 2.12 (and what follows when time allows).  Vocabulary: structure-library/ccr2-vocab.scm. Agent CCR-B2, wave MAP-1026, 2026-10-26.  Destination: theorem-library/ccr2-nv.scm.  Library: analysis. Helper prefix ...}\par
\vnbprop{nv-field-normed}{\vnbstmt{\forall k.\; (\operatorname{is\text{-}nv\text{-}field}(k) \Rightarrow \operatorname{is\text{-}normed\text{-}field}(k))}}
\vnbprop{nv-ultra}{\vnbstmt{\begin{array}{@{}l@{}} \forall k, x, y. \\ \quad (\operatorname{is\text{-}nv\text{-}field}(k) \wedge ((x \in \operatorname{carr}(k)) \wedge (y \in \operatorname{carr}(k)))) \Rightarrow  \\ \quad \quad \operatorname{fnrm}(k)(\operatorname{add}(k)(x, y)) \leq \operatorname{max}(\operatorname{fnrm}(k)(x), \operatorname{fnrm}(k)(y)) \end{array}}}
\vnbprop{nf-zero-ne-one}{\vnbstmt{\forall k.\; (\operatorname{is\text{-}normed\text{-}field}(k) \Rightarrow \neg (\operatorname{zero}(k) = \operatorname{one}(k)))}}
\vnbprop{nf-inverse}{\vnbstmt{\begin{array}{@{}l@{}} \forall k.; \forall x \in \operatorname{carr}(k). \\ \quad (\operatorname{is\text{-}normed\text{-}field}(k) \wedge \neg (x = \operatorname{zero}(k))) \Rightarrow  \\ \quad \quad \exists y \in \operatorname{carr}(k).\; (\neg (y = \operatorname{zero}(k)) \wedge (\operatorname{mul}(k)(x, y) = \operatorname{one}(k))) \end{array}}}
\vnbprop{nf-norm-one}{\vnbstmt{\forall k.\; (\operatorname{is\text{-}normed\text{-}field}(k) \Rightarrow (\operatorname{fnrm}(k)(\operatorname{one}(k)) = 1))}}
\vnbprop{rr-unit-interval-prod-one}{\vnbstmt{\begin{array}{@{}l@{}} \forall a, b \in \mathbb{R}. \\ \quad ((0 \leq a) \wedge ((a \leq 1) \wedge ((0 \leq b) \wedge ((b \leq 1) \wedge ((a \cdot b) = 1))))) \Rightarrow (a = 1) \end{array}}}
\vnbprop{nv-domain-field}{\vnbstmt{\forall k, a.\; (\operatorname{is\text{-}nv\text{-}domain}(k, a) \Rightarrow \operatorname{is\text{-}nv\text{-}field}(k))}}
\vnbprop{nv-domain-sub}{\vnbstmt{\forall k, a.; \forall x \in a.\; (\operatorname{is\text{-}nv\text{-}domain}(k, a) \Rightarrow (x \in \operatorname{carr}(k)))}}
\vnbprop{nv-domain-zero}{\vnbstmt{\forall k, a.\; (\operatorname{is\text{-}nv\text{-}domain}(k, a) \Rightarrow (\operatorname{zero}(k) \in a))}}
\vnbprop{nv-domain-one}{\vnbstmt{\forall k, a.\; (\operatorname{is\text{-}nv\text{-}domain}(k, a) \Rightarrow (\operatorname{one}(k) \in a))}}
\vnbprop{nv-domain-add}{\vnbstmt{\forall k, a.; \forall x, y \in a.\; (\operatorname{is\text{-}nv\text{-}domain}(k, a) \Rightarrow (\operatorname{add}(k)(x, y) \in a))}}
\vnbprop{nv-domain-mul}{\vnbstmt{\forall k, a.; \forall x, y \in a.\; (\operatorname{is\text{-}nv\text{-}domain}(k, a) \Rightarrow (\operatorname{mul}(k)(x, y) \in a))}}
\vnbprop{nv-domain-neg}{\vnbstmt{\forall k, a.; \forall x \in a.\; (\operatorname{is\text{-}nv\text{-}domain}(k, a) \Rightarrow (\operatorname{neg}(k)(x) \in a))}}
\vnbprop{nv-domain-norm}{\vnbstmt{\forall k, a.; \forall x \in a.\; (\operatorname{is\text{-}nv\text{-}domain}(k, a) \Rightarrow (\operatorname{fnrm}(k)(x) \leq 1))}}
\vnbprop{nv-domain-normed}{\vnbstmt{\forall k, a.\; (\operatorname{is\text{-}nv\text{-}domain}(k, a) \Rightarrow \operatorname{is\text{-}normed\text{-}field}(k))}}
\vnbprop{nv-integers-unfold}{\vnbstmt{\forall k.\; (\operatorname{nv\text{-}integers}(k) \simeq \{x \in \operatorname{carr}(k) : (\operatorname{fnrm}(k)(x) \leq 1)\})}}
\vnbprop{nv-units-unfold}{\vnbstmt{\begin{array}{@{}l@{}} \forall k, a. \\ \quad (\operatorname{nv\text{-}units}(k, a) \simeq \{x \in a : \exists y \in a.\; (\operatorname{mul}(k)(x, y) = \operatorname{one}(k))\}) \end{array}}}
\vnbprop{nv-maximal-unfold}{\vnbstmt{\forall k, a.\; (\operatorname{nv\text{-}maximal}(k, a) \simeq \{x \in a : (\operatorname{fnrm}(k)(x) < 1)\})}}
\vnbprop{nv-principal-unfold}{\vnbstmt{\begin{array}{@{}l@{}} \forall k, a, g. \\ \quad (\operatorname{nv\text{-}principal}(k, a, g) \simeq \{x \in a : \exists m \in a.\; (x = \operatorname{mul}(k)(m, g))\}) \end{array}}}
\vnbprop{nv-ideal2-unfold}{\vnbstmt{\begin{array}{@{}l@{}} \forall k, a, g, h. \\ \quad (\operatorname{nv\text{-}ideal2}(k, a, g, h) \simeq \{x \in a : \exists m, n \in a.\; (x = \operatorname{add}(k)(\operatorname{mul}(k)(m, g), \operatorname{mul}(k)(n, h)))\}) \end{array}}}
\vnbprop{nv-ball-unfold}{\vnbstmt{\forall k, a, r.\; (\operatorname{nv\text{-}ball}(k, a, r) \simeq \{x \in a : (\operatorname{fnrm}(k)(x) \leq r)\})}}
\vnbprop{nv-domain-integral}{\vnbstmt{\forall k, a.; \forall x \in a.\; (\operatorname{is\text{-}nv\text{-}domain}(k, a) \Rightarrow (x \in \operatorname{nv\text{-}integers}(k)))}}
\vnbprop{nv-integers-iff-divisibility}{\vnbstmt{\begin{array}{@{}l@{}} \forall k, a. \\ \quad (\operatorname{is\text{-}nv\text{-}domain}(k, a) \wedge \operatorname{is\text{-}fractions\text{-}of}(k, a)) \Rightarrow  \\ \quad \quad a = \operatorname{nv\text{-}integers}(k) \\ \quad \quad \Leftrightarrow\; \forall x, y \in a. \\ \quad \quad \quad (\operatorname{fnrm}(k)(x) \leq \operatorname{fnrm}(k)(y)) \Rightarrow \exists l \in a.\; (x = \operatorname{mul}(k)(l, y)) \end{array}}}
\vnbprop{rr-mul-left-one}{\vnbstmt{\forall a, b \in \mathbb{R}.\; ((a = 1) \Rightarrow ((a \cdot b) = b))}}
\vnbprop{nv-unit-norm}{\vnbstmt{\forall k, a.; \forall x \in \operatorname{nv\text{-}units}(k, a).\; (\operatorname{is\text{-}nv\text{-}domain}(k, a) \Rightarrow (\operatorname{fnrm}(k)(x) = 1))}}
\vnbprop{nv-principal-self}{\vnbstmt{\forall k, a.; \forall g \in a.\; (\operatorname{is\text{-}nv\text{-}domain}(k, a) \Rightarrow (g \in \operatorname{nv\text{-}principal}(k, a, g)))}}
\vnbprop{nv-one-unit}{\vnbstmt{\forall k, a.\; (\operatorname{is\text{-}nv\text{-}domain}(k, a) \Rightarrow (\operatorname{one}(k) \in \operatorname{nv\text{-}units}(k, a)))}}
\vnbprop{nv-generators-associate}{\vnbstmt{\begin{array}{@{}l@{}} \forall k, a.; \forall g, h \in a. \\ \quad (\operatorname{is\text{-}nv\text{-}domain}(k, a) \wedge (\operatorname{nv\text{-}principal}(k, a, g) = \operatorname{nv\text{-}principal}(k, a, h))) \Rightarrow  \\ \quad \quad \exists u \in \operatorname{nv\text{-}units}(k, a).\; (g = \operatorname{mul}(k)(u, h)) \end{array}}}
\vnbprop{nv-pow-in-domain-ind}{\vnbstmt{\begin{array}{@{}l@{}} \forall n \in \mathbb{N}, k, a.; \forall x \in a. \\ \quad \operatorname{is\text{-}nv\text{-}domain}(k, a) \Rightarrow  \\ \quad \quad \operatorname{ring\text{-}power}(\operatorname{normed\text{-}field\text{-}as\text{-}commutative\text{-}ring}(k), x, n) \in a \end{array}}}
\vnbprop{nv-pow-in-domain}{\vnbstmt{\begin{array}{@{}l@{}} \forall k, a.; \forall x \in a, n \in \mathbb{N}. \\ \quad \operatorname{is\text{-}nv\text{-}domain}(k, a) \Rightarrow  \\ \quad \quad \operatorname{ring\text{-}power}(\operatorname{normed\text{-}field\text{-}as\text{-}commutative\text{-}ring}(k), x, n) \in a \end{array}}}
\vnbprop{nv-pow-norm-ind}{\vnbstmt{\begin{array}{@{}l@{}} \forall n \in \mathbb{N}, k.; \forall x \in \operatorname{carr}(k). \\ \quad (\operatorname{is\text{-}normed\text{-}field}(k) \wedge (\operatorname{fnrm}(k)(x) = 1)) \Rightarrow  \\ \quad \quad \operatorname{fnrm}(k)(\operatorname{ring\text{-}power}(\operatorname{normed\text{-}field\text{-}as\text{-}commutative\text{-}ring}(k), x, n)) = 1 \end{array}}}
\vnbprop{nv-pow-norm}{\vnbstmt{\begin{array}{@{}l@{}} \forall k.; \forall x \in \operatorname{carr}(k), n \in \mathbb{N}. \\ \quad (\operatorname{is\text{-}normed\text{-}field}(k) \wedge (\operatorname{fnrm}(k)(x) = 1)) \Rightarrow  \\ \quad \quad \operatorname{fnrm}(k)(\operatorname{ring\text{-}power}(\operatorname{normed\text{-}field\text{-}as\text{-}commutative\text{-}ring}(k), x, n)) = 1 \end{array}}}
\vnbprop{nv-compact-domain}{\vnbstmt{\forall k, a.\; (\operatorname{is\text{-}compact\text{-}nv\text{-}domain}(k, a) \Rightarrow \operatorname{is\text{-}nv\text{-}domain}(k, a))}}
\vnbprop{nv-compact-compact}{\vnbstmt{\begin{array}{@{}l@{}} \forall k, a. \\ \quad \operatorname{is\text{-}compact\text{-}nv\text{-}domain}(k, a) \Rightarrow  \\ \quad \quad \operatorname{is\text{-}compact}(\operatorname{subspace\text{-}ms}(\operatorname{nf\text{-}metric\text{-}space}(k), a)) \end{array}}}
\vnbprop{nv-domain-subset-pts}{\vnbstmt{\forall k, a.\; (\operatorname{is\text{-}nv\text{-}domain}(k, a) \Rightarrow (a \subseteq \operatorname{pts}(\operatorname{nf\text{-}metric\text{-}space}(k))))}}
\vnbprop{nv-compact-closed}{\vnbstmt{\begin{array}{@{}l@{}} \forall k, a. \\ \quad \operatorname{is\text{-}compact\text{-}nv\text{-}domain}(k, a) \Rightarrow \operatorname{is\text{-}closed}(\operatorname{nf\text{-}metric\text{-}space}(k), a) \end{array}}}
\vnbprop{nv-compact-near-one}{\vnbstmt{\begin{array}{@{}l@{}} \forall k, a.; \forall x \in a, s. \\ \quad (\operatorname{is\text{-}compact\text{-}nv\text{-}domain}(k, a) \wedge ((\operatorname{fnrm}(k)(x) = 1) \wedge \operatorname{pos\text{-}rr}(s))) \Rightarrow  \\ \quad \quad \exists n, m \in \mathbb{N}. \\ \quad \quad \quad n < m \wedge \\ \quad \quad \quad \operatorname{fnrm}(k)(\operatorname{add}(k)(\operatorname{ring\text{-}power}(\operatorname{normed\text{-}field\text{-}as\text{-}commutative\text{-}ring}(k), x, n), \\ \quad \quad \quad \quad \quad \operatorname{neg}(k)(\operatorname{ring\text{-}power}(\operatorname{normed\text{-}field\text{-}as\text{-}commutative\text{-}ring}(k), x, m)))) \\ \quad \quad \quad \leq\; s \end{array}}}
\vnbprop{nv-compact-norm-one-unit}{\vnbstmt{\begin{array}{@{}l@{}} \forall k, a.; \forall x \in a. \\ \quad (\operatorname{is\text{-}compact\text{-}nv\text{-}domain}(k, a) \wedge (\operatorname{fnrm}(k)(x) = 1)) \Rightarrow (x \in \operatorname{nv\text{-}units}(k, a)) \end{array}}}
\vnbprop{nv-compact-units-norm-one}{\vnbstmt{\begin{array}{@{}l@{}} \forall k, a.; \forall x \in a. \\ \quad \operatorname{is\text{-}compact\text{-}nv\text{-}domain}(k, a) \Rightarrow ((x \in \operatorname{nv\text{-}units}(k, a)) \Leftrightarrow (\operatorname{fnrm}(k)(x) = 1)) \end{array}}}
\vnbprop{nv-ideal-sub}{\vnbstmt{\forall k, a, i.; \forall x \in i.\; (\operatorname{is\text{-}nv\text{-}ideal}(k, a, i) \Rightarrow (x \in a))}}
\vnbprop{nv-ideal-zero}{\vnbstmt{\forall k, a, i.\; (\operatorname{is\text{-}nv\text{-}ideal}(k, a, i) \Rightarrow (\operatorname{zero}(k) \in i))}}
\vnbprop{nv-ideal-add}{\vnbstmt{\begin{array}{@{}l@{}} \forall k, a, i.; \forall x, y \in i. \\ \quad \operatorname{is\text{-}nv\text{-}ideal}(k, a, i) \Rightarrow (\operatorname{add}(k)(x, y) \in i) \end{array}}}
\vnbprop{nv-ideal-neg}{\vnbstmt{\forall k, a, i.; \forall x \in i.\; (\operatorname{is\text{-}nv\text{-}ideal}(k, a, i) \Rightarrow (\operatorname{neg}(k)(x) \in i))}}
\vnbprop{nv-ideal-mul}{\vnbstmt{\begin{array}{@{}l@{}} \forall k, a, i.; \forall r \in a, x \in i. \\ \quad \operatorname{is\text{-}nv\text{-}ideal}(k, a, i) \Rightarrow (\operatorname{mul}(k)(r, x) \in i) \end{array}}}
\vnbprop{nv-maximal-sub}{\vnbstmt{\forall k, a.; \forall x \in \operatorname{nv\text{-}maximal}(k, a).\; (x \in a)}}
\vnbprop{nv-maximal-norm}{\vnbstmt{\forall k, a.; \forall x \in \operatorname{nv\text{-}maximal}(k, a).\; (\operatorname{fnrm}(k)(x) < 1)}}
\vnbprop{nv-maximal-zero}{\vnbstmt{\forall k, a.\; (\operatorname{is\text{-}nv\text{-}domain}(k, a) \Rightarrow (\operatorname{zero}(k) \in \operatorname{nv\text{-}maximal}(k, a)))}}
\vnbprop{nv-maximal-add}{\vnbstmt{\begin{array}{@{}l@{}} \forall k, a.; \forall x, y \in \operatorname{nv\text{-}maximal}(k, a). \\ \quad \operatorname{is\text{-}nv\text{-}domain}(k, a) \Rightarrow (\operatorname{add}(k)(x, y) \in \operatorname{nv\text{-}maximal}(k, a)) \end{array}}}
\vnbprop{nv-maximal-neg}{\vnbstmt{\begin{array}{@{}l@{}} \forall k, a.; \forall x \in \operatorname{nv\text{-}maximal}(k, a). \\ \quad \operatorname{is\text{-}nv\text{-}domain}(k, a) \Rightarrow (\operatorname{neg}(k)(x) \in \operatorname{nv\text{-}maximal}(k, a)) \end{array}}}
\vnbprop{nv-maximal-mul}{\vnbstmt{\begin{array}{@{}l@{}} \forall k, a.; \forall r \in a, x \in \operatorname{nv\text{-}maximal}(k, a). \\ \quad \operatorname{is\text{-}nv\text{-}domain}(k, a) \Rightarrow (\operatorname{mul}(k)(r, x) \in \operatorname{nv\text{-}maximal}(k, a)) \end{array}}}
\vnbprop{nv-maximal-ideal}{\vnbstmt{\forall k, a.\; (\operatorname{is\text{-}nv\text{-}domain}(k, a) \Rightarrow \operatorname{is\text{-}nv\text{-}ideal}(k, a, \operatorname{nv\text{-}maximal}(k, a)))}}
\vnbprop{nv-maximal-proper}{\vnbstmt{\forall k, a.\; (\operatorname{is\text{-}nv\text{-}domain}(k, a) \Rightarrow \neg (\operatorname{one}(k) \in \operatorname{nv\text{-}maximal}(k, a)))}}
\vnbprop{nv-compact-proper-in-maximal}{\vnbstmt{\begin{array}{@{}l@{}} \forall k, a, i.; \forall x \in i. \\ \quad (\operatorname{is\text{-}compact\text{-}nv\text{-}domain}(k, a) \wedge (\operatorname{is\text{-}nv\text{-}ideal}(k, a, i) \wedge \neg (\operatorname{one}(k) \in i))) \Rightarrow  \\ \quad \quad x \in \operatorname{nv\text{-}maximal}(k, a) \end{array}}}
\vnbprop{nv-compact-maximal-unique}{\vnbstmt{\begin{array}{@{}l@{}} \forall k, a. \\ \quad \operatorname{is\text{-}compact\text{-}nv\text{-}domain}(k, a) \Rightarrow  \\ \quad \quad \operatorname{is\text{-}nv\text{-}ideal}(k, a, \operatorname{nv\text{-}maximal}(k, a)) \wedge \\ \quad \quad \neg (\operatorname{one}(k) \in \operatorname{nv\text{-}maximal}(k, a)) \wedge \\ \quad \quad \forall i. \\ \quad \quad \quad (\operatorname{is\text{-}nv\text{-}ideal}(k, a, i) \wedge \neg (\operatorname{one}(k) \in i)) \Rightarrow (i \subseteq \operatorname{nv\text{-}maximal}(k, a)) \end{array}}}
\vnbprop{nv-pid-domain}{\vnbstmt{\forall k, a.\; (\operatorname{is\text{-}nv\text{-}pid}(k, a) \Rightarrow \operatorname{is\text{-}nv\text{-}domain}(k, a))}}
\vnbprop{nv-pid-principal}{\vnbstmt{\begin{array}{@{}l@{}} \forall k, a, i. \\ \quad (\operatorname{is\text{-}nv\text{-}pid}(k, a) \wedge \operatorname{is\text{-}nv\text{-}ideal}(k, a, i)) \Rightarrow  \\ \quad \quad \exists g \in a.\; (i = \operatorname{nv\text{-}principal}(k, a, g)) \end{array}}}
\vnbprop{nv-pid-maximal-generator}{\vnbstmt{\begin{array}{@{}l@{}} \forall k, a. \\ \quad \operatorname{is\text{-}nv\text{-}pid}(k, a) \Rightarrow  \\ \quad \quad \exists p \in a.\; (\operatorname{nv\text{-}maximal}(k, a) = \operatorname{nv\text{-}principal}(k, a, p)) \end{array}}}
\vnbprop{rr-mul-cancel-right-one}{\vnbstmt{\forall a, b \in \mathbb{R}.\; (((0 < b) \wedge ((a \cdot b) = b)) \Rightarrow (a = 1))}}
\vnbprop{nv-compact-units-transitive}{\vnbstmt{\begin{array}{@{}l@{}} \forall k, a.; \forall p, y \in a. \\ \quad \operatorname{is\text{-}compact\text{-}nv\text{-}domain}(k, a) \wedge \\ \quad \operatorname{nv\text{-}maximal}(k, a) = \operatorname{nv\text{-}principal}(k, a, p) \wedge \\ \quad \operatorname{fnrm}(k)(y) = \operatorname{fnrm}(k)(p) \Rightarrow \\ \quad \quad \exists u \in \operatorname{nv\text{-}units}(k, a).\; (y = \operatorname{mul}(k)(u, p)) \end{array}}}
\vnbprop{nv-ideal2-sub}{\vnbstmt{\forall k, a, g, h.; \forall x \in \operatorname{nv\text{-}ideal2}(k, a, g, h).\; (x \in a)}}
\vnbprop{nv-ideal2-zero}{\vnbstmt{\begin{array}{@{}l@{}} \forall k, a.; \forall g, h \in a. \\ \quad \operatorname{is\text{-}nv\text{-}domain}(k, a) \Rightarrow (\operatorname{zero}(k) \in \operatorname{nv\text{-}ideal2}(k, a, g, h)) \end{array}}}
\vnbprop{nv-ideal2-add}{\vnbstmt{\begin{array}{@{}l@{}} \forall k, a.; \forall g, h \in a, x, y \in \operatorname{nv\text{-}ideal2}(k, a, g, h). \\ \quad \operatorname{is\text{-}nv\text{-}domain}(k, a) \Rightarrow (\operatorname{add}(k)(x, y) \in \operatorname{nv\text{-}ideal2}(k, a, g, h)) \end{array}}}
\vnbprop{nv-ideal2-neg}{\vnbstmt{\begin{array}{@{}l@{}} \forall k, a.; \forall g, h \in a, x \in \operatorname{nv\text{-}ideal2}(k, a, g, h). \\ \quad \operatorname{is\text{-}nv\text{-}domain}(k, a) \Rightarrow (\operatorname{neg}(k)(x) \in \operatorname{nv\text{-}ideal2}(k, a, g, h)) \end{array}}}
\vnbprop{nv-ideal2-mul}{\vnbstmt{\begin{array}{@{}l@{}} \forall k, a.; \forall g, h, r \in a, x \in \operatorname{nv\text{-}ideal2}(k, a, g, h). \\ \quad \operatorname{is\text{-}nv\text{-}domain}(k, a) \Rightarrow (\operatorname{mul}(k)(r, x) \in \operatorname{nv\text{-}ideal2}(k, a, g, h)) \end{array}}}
\vnbprop{nv-ideal2-ideal}{\vnbstmt{\begin{array}{@{}l@{}} \forall k, a.; \forall g, h \in a. \\ \quad \operatorname{is\text{-}nv\text{-}domain}(k, a) \Rightarrow \operatorname{is\text{-}nv\text{-}ideal}(k, a, \operatorname{nv\text{-}ideal2}(k, a, g, h)) \end{array}}}
\vnbprop{rr-le-prod-unit}{\vnbstmt{\begin{array}{@{}l@{}} \forall l, b, c \in \mathbb{R}. \\ \quad ((0 \leq l) \wedge ((0 < c) \wedge ((b \leq c) \wedge ((c = (l \cdot b)) \wedge (l \leq 1))))) \Rightarrow (l = 1) \end{array}}}
\vnbprop{nv-compact-pid-principal-ball}{\vnbstmt{\begin{array}{@{}l@{}} \forall k, a.; \forall c \in a. \\ \quad (\operatorname{is\text{-}compact\text{-}nv\text{-}domain}(k, a) \wedge \operatorname{is\text{-}nv\text{-}pid}(k, a)) \Rightarrow  \\ \quad \quad \operatorname{nv\text{-}principal}(k, a, c) = \operatorname{nv\text{-}ball}(k, a, \operatorname{fnrm}(k)(c)) \end{array}}}
\vnbprop{nv-pid-integers}{\vnbstmt{\begin{array}{@{}l@{}} \forall k, a. \\ \quad (\operatorname{is\text{-}compact\text{-}nv\text{-}domain}(k, a) \wedge (\operatorname{is\text{-}nv\text{-}pid}(k, a) \wedge \operatorname{is\text{-}fractions\text{-}of}(k, a))) \Rightarrow  \\ \quad \quad a = \operatorname{nv\text{-}integers}(k) \end{array}}}
\vnbprop{rr-power-decay}{\vnbstmt{\begin{array}{@{}l@{}} \forall r, t \in \mathbb{R}. \\ \quad ((0 \leq r) \wedge ((r < 1) \wedge (0 < t))) \Rightarrow \exists k \in \mathbb{N}.\; (\operatorname{power}(r, k) < t) \end{array}}}
\vnbprop{nv-pow-norm-power-ind}{\vnbstmt{\begin{array}{@{}l@{}} \forall n \in \mathbb{N}, k.; \forall x \in \operatorname{carr}(k). \\ \quad \operatorname{is\text{-}normed\text{-}field}(k) \Rightarrow  \\ \quad \quad \operatorname{fnrm}(k)(\operatorname{ring\text{-}power}(\operatorname{normed\text{-}field\text{-}as\text{-}commutative\text{-}ring}(k), x, n)) = \operatorname{power}(\operatorname{fnrm}(k)(x), n) \end{array}}}
\vnbprop{nv-pow-norm-power}{\vnbstmt{\begin{array}{@{}l@{}} \forall k.; \forall x \in \operatorname{carr}(k), n \in \mathbb{N}. \\ \quad \operatorname{is\text{-}normed\text{-}field}(k) \Rightarrow  \\ \quad \quad \operatorname{fnrm}(k)(\operatorname{ring\text{-}power}(\operatorname{normed\text{-}field\text{-}as\text{-}commutative\text{-}ring}(k), x, n)) = \operatorname{power}(\operatorname{fnrm}(k)(x), n) \end{array}}}
\vnbprop{nv-principal-norm-le}{\vnbstmt{\begin{array}{@{}l@{}} \forall k, a.; \forall g \in a, x \in \operatorname{nv\text{-}principal}(k, a, g). \\ \quad \operatorname{is\text{-}nv\text{-}domain}(k, a) \Rightarrow (\operatorname{fnrm}(k)(x) \leq \operatorname{fnrm}(k)(g)) \end{array}}}
\vnbprop{nv-uniformizer-factorization}{\vnbstmt{\begin{array}{@{}l@{}} \forall k, a, w.; \forall l \in a. \\ \quad \operatorname{is\text{-}compact\text{-}nv\text{-}domain}(k, a) \wedge \\ \quad w \in a \wedge \\ \quad \operatorname{nv\text{-}maximal}(k, a) = \operatorname{nv\text{-}principal}(k, a, w) \wedge \\ \quad \neg (w = \operatorname{zero}(k)) \wedge \\ \quad \neg (l = \operatorname{zero}(k)) \Rightarrow \\ \quad \quad \exists e \in \mathbb{N}, u \in \operatorname{nv\text{-}units}(k, a). \\ \quad \quad \quad l = \operatorname{mul}(k)(\operatorname{ring\text{-}power}(\operatorname{normed\text{-}field\text{-}as\text{-}commutative\text{-}ring}(k), w, e), u) \end{array}}}
\vnbprop{nv-uniformizer-exponent-le}{\vnbstmt{\begin{array}{@{}l@{}} \forall k, a, w.; \forall b, c \in \mathbb{N}, u, v \in \operatorname{nv\text{-}units}(k, a). \\ \quad \operatorname{is\text{-}compact\text{-}nv\text{-}domain}(k, a) \wedge \\ \quad w \in a \wedge \\ \quad \operatorname{nv\text{-}maximal}(k, a) = \operatorname{nv\text{-}principal}(k, a, w) \wedge \\ \quad \neg (w = \operatorname{zero}(k)) \wedge \\ \quad b \leq c \wedge \\ \quad \operatorname{mul}(k)(\operatorname{ring\text{-}power}(\operatorname{normed\text{-}field\text{-}as\text{-}commutative\text{-}ring}(k), w, b), u) \\ \quad =\; \operatorname{mul}(k)(\operatorname{ring\text{-}power}(\operatorname{normed\text{-}field\text{-}as\text{-}commutative\text{-}ring}(k), w, c), v) \Rightarrow \\ \quad \quad b = c \end{array}}}
\vnbprop{nv-uniformizer-factorization-unique}{\vnbstmt{\begin{array}{@{}l@{}} \forall k, a, w.; \forall b, c \in \mathbb{N}, u, v \in \operatorname{nv\text{-}units}(k, a). \\ \quad \operatorname{is\text{-}compact\text{-}nv\text{-}domain}(k, a) \wedge \\ \quad w \in a \wedge \\ \quad \operatorname{nv\text{-}maximal}(k, a) = \operatorname{nv\text{-}principal}(k, a, w) \wedge \\ \quad \neg (w = \operatorname{zero}(k)) \wedge \\ \quad \operatorname{mul}(k)(\operatorname{ring\text{-}power}(\operatorname{normed\text{-}field\text{-}as\text{-}commutative\text{-}ring}(k), w, b), u) \\ \quad =\; \operatorname{mul}(k)(\operatorname{ring\text{-}power}(\operatorname{normed\text{-}field\text{-}as\text{-}commutative\text{-}ring}(k), w, c), v) \Rightarrow \\ \quad \quad (b = c) \wedge (u = v) \end{array}}}
\vnbprop{nv-local-factorization}{\vnbstmt{\begin{array}{@{}l@{}} \forall k, a, w.; \forall x \in \operatorname{carr}(k). \\ \quad \operatorname{is\text{-}compact\text{-}nv\text{-}domain}(k, a) \wedge \\ \quad w \in a \wedge \\ \quad \operatorname{nv\text{-}maximal}(k, a) = \operatorname{nv\text{-}principal}(k, a, w) \wedge \\ \quad \neg (w = \operatorname{zero}(k)) \wedge \\ \quad a = \operatorname{nv\text{-}integers}(k) \wedge \\ \quad \neg (x = \operatorname{zero}(k)) \Rightarrow \\ \quad \quad \exists e \in \mathbb{N}, u \in \operatorname{nv\text{-}units}(k, a). \\ \quad \quad \quad x = \operatorname{mul}(k)(\operatorname{ring\text{-}power}(\operatorname{normed\text{-}field\text{-}as\text{-}commutative\text{-}ring}(k), w, e), u) \vee \\ \quad \quad \quad \operatorname{mul}(k)(x, \operatorname{ring\text{-}power}(\operatorname{normed\text{-}field\text{-}as\text{-}commutative\text{-}ring}(k), w, e)) = u \end{array}}}
\vnbprop{nv-compact-one-minus-unit}{\vnbstmt{\begin{array}{@{}l@{}} \forall k, a.; \forall x \in a. \\ \quad (\operatorname{is\text{-}compact\text{-}nv\text{-}domain}(k, a) \wedge (\operatorname{fnrm}(k)(x) < 1)) \Rightarrow  \\ \quad \quad \operatorname{add}(k)(\operatorname{one}(k), \operatorname{neg}(k)(x)) \in \operatorname{nv\text{-}units}(k, a) \end{array}}}
\vnbprop{nv-multiples-unfold}{\vnbstmt{\begin{array}{@{}l@{}} \forall k, a, g. \\ \quad (\operatorname{nv\text{-}multiples}(k, a, g) \simeq \{x \in \operatorname{carr}(k) : \exists m \in a.\; (x = \operatorname{mul}(k)(m, g))\}) \end{array}}}
\vnbprop{nv-uniformizer-pid}{\vnbstmt{\begin{array}{@{}l@{}} \forall k, a, w. \\ \quad \operatorname{is\text{-}compact\text{-}nv\text{-}domain}(k, a) \wedge \\ \quad w \in a \wedge \\ \quad \operatorname{nv\text{-}maximal}(k, a) = \operatorname{nv\text{-}principal}(k, a, w) \wedge \\ \quad \neg (w = \operatorname{zero}(k)) \Rightarrow \\ \quad \quad \operatorname{is\text{-}nv\text{-}pid}(k, a) \end{array}}}
\vnbprop{nv-module-sub}{\vnbstmt{\forall k, a, e.; \forall x \in e.\; (\operatorname{is\text{-}nv\text{-}module}(k, a, e) \Rightarrow (x \in \operatorname{carr}(k)))}}
\vnbprop{nv-module-zero}{\vnbstmt{\forall k, a, e.\; (\operatorname{is\text{-}nv\text{-}module}(k, a, e) \Rightarrow (\operatorname{zero}(k) \in e))}}
\vnbprop{nv-module-add}{\vnbstmt{\begin{array}{@{}l@{}} \forall k, a, e.; \forall x, y \in e. \\ \quad \operatorname{is\text{-}nv\text{-}module}(k, a, e) \Rightarrow (\operatorname{add}(k)(x, y) \in e) \end{array}}}
\vnbprop{nv-module-neg}{\vnbstmt{\forall k, a, e.; \forall x \in e.\; (\operatorname{is\text{-}nv\text{-}module}(k, a, e) \Rightarrow (\operatorname{neg}(k)(x) \in e))}}
\vnbprop{nv-module-mul}{\vnbstmt{\begin{array}{@{}l@{}} \forall k, a, e.; \forall r \in a, x \in e. \\ \quad \operatorname{is\text{-}nv\text{-}module}(k, a, e) \Rightarrow (\operatorname{mul}(k)(r, x) \in e) \end{array}}}
\vnbprop{nv-local-module-in-multiple}{\vnbstmt{\begin{array}{@{}l@{}} \forall k, a, w, e.; \forall b \in \operatorname{carr}(k), y \in e. \\ \quad \operatorname{is\text{-}compact\text{-}nv\text{-}domain}(k, a) \wedge \\ \quad w \in a \wedge \\ \quad \operatorname{nv\text{-}maximal}(k, a) = \operatorname{nv\text{-}principal}(k, a, w) \wedge \\ \quad \neg (w = \operatorname{zero}(k)) \wedge \\ \quad a = \operatorname{nv\text{-}integers}(k) \wedge \\ \quad \operatorname{is\text{-}nv\text{-}module}(k, a, e) \wedge \\ \quad \neg (b \in e) \Rightarrow \\ \quad \quad \exists t \in a.\; (y = \operatorname{mul}(k)(t, \operatorname{mul}(k)(b, w))) \end{array}}}
\vnbprop{nv-local-module-principal}{\vnbstmt{\begin{array}{@{}l@{}} \forall k, a, w, e. \\ \quad \operatorname{is\text{-}compact\text{-}nv\text{-}domain}(k, a) \wedge \\ \quad w \in a \wedge \\ \quad \operatorname{nv\text{-}maximal}(k, a) = \operatorname{nv\text{-}principal}(k, a, w) \wedge \\ \quad \neg (w = \operatorname{zero}(k)) \wedge \\ \quad a = \operatorname{nv\text{-}integers}(k) \wedge \\ \quad \operatorname{is\text{-}nv\text{-}module}(k, a, e) \wedge \\ \quad \neg (e = \operatorname{carr}(k)) \Rightarrow \\ \quad \quad \exists g \in \operatorname{carr}(k).\; (e = \operatorname{nv\text{-}multiples}(k, a, g)) \end{array}}}
\vnbprop{nv-strict-dominance}{\vnbstmt{\begin{array}{@{}l@{}} \forall k.; \forall x, y \in \operatorname{carr}(k). \\ \quad (\operatorname{is\text{-}nv\text{-}field}(k) \wedge (\operatorname{fnrm}(k)(x) < \operatorname{fnrm}(k)(y))) \Rightarrow  \\ \quad \quad \operatorname{fnrm}(k)(\operatorname{add}(k)(x, y)) = \operatorname{fnrm}(k)(y) \end{array}}}
\vnbprop{nv-finsum-le}{\vnbstmt{\begin{array}{@{}l@{}} \forall k.; \forall q \in \mathbf{Set}, f.; \forall b \in \mathbb{R}. \\ \quad \operatorname{is\text{-}nv\text{-}field}(k) \wedge \\ \quad |q| \in \mathbb{N} \wedge \\ \quad 0 \leq b \wedge \\ \quad \forall z \in q.\; ((f(z) \in \operatorname{carr}(k)) \wedge (\operatorname{fnrm}(k)(f(z)) \leq b)) \Rightarrow \\ \quad \quad \operatorname{fnrm}(k)(\operatorname{finsum}(\operatorname{ring\text{-}additive\text{-}ag}(\operatorname{normed\text{-}field\text{-}as\text{-}commutative\text{-}ring}(k)), f, q)) \leq b \end{array}}}
\vnbprop{nv-finsum-lt}{\vnbstmt{\begin{array}{@{}l@{}} \forall k.; \forall q \in \mathbf{Set}, f.; \forall b \in \mathbb{R}. \\ \quad \operatorname{is\text{-}nv\text{-}field}(k) \wedge \\ \quad |q| \in \mathbb{N} \wedge \\ \quad 0 < b \wedge \\ \quad \forall z \in q.\; ((f(z) \in \operatorname{carr}(k)) \wedge (\operatorname{fnrm}(k)(f(z)) < b)) \Rightarrow \\ \quad \quad \operatorname{fnrm}(k)(\operatorname{finsum}(\operatorname{ring\text{-}additive\text{-}ag}(\operatorname{normed\text{-}field\text{-}as\text{-}commutative\text{-}ring}(k)), f, q)) < b \end{array}}}
\vnbprop{rr-prod-lt-one-left}{\vnbstmt{\begin{array}{@{}l@{}} \forall a, b \in \mathbb{R}. \\ \quad ((0 \leq a) \wedge ((0 \leq b) \wedge ((a < 1) \wedge (b \leq 1)))) \Rightarrow ((a \cdot b) < 1) \end{array}}}
\vnbprop{rr-prod-lt-one-right}{\vnbstmt{\begin{array}{@{}l@{}} \forall a, b \in \mathbb{R}. \\ \quad ((0 \leq a) \wedge ((0 \leq b) \wedge ((a \leq 1) \wedge (b < 1)))) \Rightarrow ((a \cdot b) < 1) \end{array}}}
\vnbprop{rr-prod-le-one}{\vnbstmt{\begin{array}{@{}l@{}} \forall a, b \in \mathbb{R}. \\ \quad ((0 \leq a) \wedge ((0 \leq b) \wedge ((a \leq 1) \wedge (b \leq 1)))) \Rightarrow ((a \cdot b) \leq 1) \end{array}}}
\vnbprop{nv-gauss-content-one}{\vnbstmt{\begin{array}{@{}l@{}} \forall k.; \forall p, q \in \operatorname{finsupp}(\operatorname{normed\text{-}field\text{-}as\text{-}commutative\text{-}ring}(k), \operatorname{nn\text{-}add\text{-}monoid}). \\ \quad (\operatorname{is\text{-}nv\text{-}field}(k) \wedge (\operatorname{is\text{-}nv\text{-}content\text{-}one}(k, p) \wedge \operatorname{is\text{-}nv\text{-}content\text{-}one}(k, q))) \Rightarrow  \\ \quad \quad \operatorname{is\text{-}nv\text{-}content\text{-}one}(k, \operatorname{monalg\text{-}mul}(\operatorname{normed\text{-}field\text{-}as\text{-}commutative\text{-}ring}(k), \operatorname{nn\text{-}add\text{-}monoid}, p, q)) \end{array}}}
\vnbprop{rr-prod-le-gen}{\vnbstmt{\begin{array}{@{}l@{}} \forall a, b, c, d \in \mathbb{R}. \\ \quad ((0 \leq a) \wedge ((0 \leq b) \wedge ((0 \leq c) \wedge ((a \leq c) \wedge (b \leq d))))) \Rightarrow ((a \cdot b) \leq (c \cdot d)) \end{array}}}
\vnbprop{rr-prod-lt-gen-left}{\vnbstmt{\begin{array}{@{}l@{}} \forall a, b, c, d \in \mathbb{R}. \\ \quad ((0 \leq a) \wedge ((0 \leq b) \wedge ((a < c) \wedge ((b \leq d) \wedge (0 < d))))) \Rightarrow ((a \cdot b) < (c \cdot d)) \end{array}}}
\vnbprop{rr-prod-lt-gen-right}{\vnbstmt{\begin{array}{@{}l@{}} \forall a, b, c, d \in \mathbb{R}. \\ \quad ((0 \leq a) \wedge ((0 \leq b) \wedge ((a \leq c) \wedge ((b < d) \wedge (0 < c))))) \Rightarrow ((a \cdot b) < (c \cdot d)) \end{array}}}
\vnbprop{nv-content-multiplicative}{\vnbstmt{\begin{array}{@{}l@{}} \forall k.; \forall p, q \in \operatorname{finsupp}(\operatorname{normed\text{-}field\text{-}as\text{-}commutative\text{-}ring}(k), \operatorname{nn\text{-}add\text{-}monoid}), c, d \in \mathbb{R}. \\ \quad (\operatorname{is\text{-}nv\text{-}field}(k) \wedge ((0 < c) \wedge ((0 < d) \wedge (\operatorname{is\text{-}nv\text{-}content}(k, p, c) \wedge \operatorname{is\text{-}nv\text{-}content}(k, q, d))))) \Rightarrow  \\ \quad \quad \operatorname{is\text{-}nv\text{-}content}(k, \operatorname{monalg\text{-}mul}(\operatorname{normed\text{-}field\text{-}as\text{-}commutative\text{-}ring}(k), \operatorname{nn\text{-}add\text{-}monoid}, p, q), (c \cdot d)) \end{array}}}
\vnbprop{nv-content-ultra}{\vnbstmt{\begin{array}{@{}l@{}} \forall k.; \forall p, q \in \operatorname{finsupp}(\operatorname{normed\text{-}field\text{-}as\text{-}commutative\text{-}ring}(k), \operatorname{nn\text{-}add\text{-}monoid}), c, d \in \mathbb{R}, i \in \mathbb{N}. \\ \quad (\operatorname{is\text{-}nv\text{-}field}(k) \wedge (\operatorname{is\text{-}nv\text{-}content}(k, p, c) \wedge \operatorname{is\text{-}nv\text{-}content}(k, q, d))) \Rightarrow  \\ \quad \quad \operatorname{fnrm}(k)(\operatorname{monalg\text{-}add}(\operatorname{normed\text{-}field\text{-}as\text{-}commutative\text{-}ring}(k), \operatorname{nn\text{-}add\text{-}monoid}, p, q)(i)) \leq \operatorname{max}(c, d) \end{array}}}
\vnbprop{rr-power-decay-scaled}{\vnbstmt{\begin{array}{@{}l@{}} \forall q, r, s \in \mathbb{R}. \\ \quad ((0 \leq q) \wedge ((q < 1) \wedge ((0 < r) \wedge (0 < s)))) \Rightarrow \exists n \in \mathbb{N}.\; ((\operatorname{power}(q, n) \cdot r) < s) \end{array}}}
\vnbprop{nv-scaling-continuous}{\vnbstmt{\begin{array}{@{}l@{}} \forall k.; \forall a \in \mathbf{Set}, t, g \in \operatorname{carr}(k). \\ \quad \operatorname{is\text{-}normed\text{-}field}(k) \wedge \\ \quad a \subseteq \operatorname{pts}(\operatorname{nf\text{-}metric\text{-}space}(k)) \wedge \\ \quad (\operatorname{fnrm}(k)(t) \cdot \operatorname{fnrm}(k)(g)) = 1 \wedge \\ \quad 0 < \operatorname{fnrm}(k)(g) \Rightarrow \\ \quad \quad \operatorname{is\text{-}continuous}(\operatorname{subspace\text{-}ms}(\operatorname{nf\text{-}metric\text{-}space}(k), a), \operatorname{nf\text{-}metric\text{-}space}(k), \\ \quad \quad \quad (\lambda v \in \operatorname{pts}(\operatorname{subspace\text{-}ms}(\operatorname{nf\text{-}metric\text{-}space}(k), a)).\; \operatorname{mul}(k)(t, v))) \end{array}}}
\vnbprop{nv-closed-balls-compact}{\vnbstmt{\begin{array}{@{}l@{}} \forall k.; \forall a \in \mathbb{R}, b \in \operatorname{carr}(k), c \in \mathbb{R}. \\ \quad \operatorname{is\text{-}nv\text{-}field}(k) \wedge \\ \quad 0 < a \wedge \\ \quad \operatorname{is\text{-}compact}(\operatorname{subspace\text{-}ms}(\operatorname{nf\text{-}metric\text{-}space}(k), \operatorname{closed\text{-}ball}(\operatorname{nf\text{-}metric\text{-}space}(k), \operatorname{zero}(k), a))) \wedge \\ \quad 0 < \operatorname{fnrm}(k)(b) \wedge \\ \quad \operatorname{fnrm}(k)(b) < 1 \wedge \\ \quad 0 < c \Rightarrow \\ \quad \quad \operatorname{is\text{-}compact}(\operatorname{subspace\text{-}ms}(\operatorname{nf\text{-}metric\text{-}space}(k), \operatorname{closed\text{-}ball}(\operatorname{nf\text{-}metric\text{-}space}(k), \operatorname{zero}(k), c))) \end{array}}}
\vnbprop{nv-integers-compact}{\vnbstmt{\begin{array}{@{}l@{}} \forall k.; \forall a \in \mathbb{R}, b \in \operatorname{carr}(k). \\ \quad \operatorname{is\text{-}nv\text{-}field}(k) \wedge \\ \quad 0 < a \wedge \\ \quad \operatorname{is\text{-}compact}(\operatorname{subspace\text{-}ms}(\operatorname{nf\text{-}metric\text{-}space}(k), \operatorname{closed\text{-}ball}(\operatorname{nf\text{-}metric\text{-}space}(k), \operatorname{zero}(k), a))) \wedge \\ \quad 0 < \operatorname{fnrm}(k)(b) \wedge \\ \quad \operatorname{fnrm}(k)(b) < 1 \Rightarrow \\ \quad \quad \operatorname{is\text{-}compact}(\operatorname{subspace\text{-}ms}(\operatorname{nf\text{-}metric\text{-}space}(k), \operatorname{nv\text{-}integers}(k))) \end{array}}}
\vnbprop{nv-uniformizer-power-small}{\vnbstmt{\begin{array}{@{}l@{}} \forall k, a, w.; \forall d \in \mathbb{N}. \\ \quad \operatorname{is\text{-}compact\text{-}nv\text{-}domain}(k, a) \wedge \\ \quad w \in a \wedge \\ \quad \operatorname{nv\text{-}maximal}(k, a) = \operatorname{nv\text{-}principal}(k, a, w) \wedge \\ \quad \neg (w = \operatorname{zero}(k)) \wedge \\ \quad \neg (d = 0) \Rightarrow \\ \quad \quad \operatorname{fnrm}(k)(\operatorname{ring\text{-}power}(\operatorname{normed\text{-}field\text{-}as\text{-}commutative\text{-}ring}(k), w, d)) < 1 \end{array}}}
\vnbprop{nv-order-nonneg}{\vnbstmt{\begin{array}{@{}l@{}} \forall k, a, w.; \forall e \in \mathbb{N}, u \in \operatorname{nv\text{-}units}(k, a), j \in \mathbb{N}. \\ \quad \operatorname{is\text{-}compact\text{-}nv\text{-}domain}(k, a) \wedge \\ \quad w \in a \wedge \\ \quad \operatorname{nv\text{-}maximal}(k, a) = \operatorname{nv\text{-}principal}(k, a, w) \wedge \\ \quad \neg (w = \operatorname{zero}(k)) \Rightarrow \\ \quad \quad \operatorname{mul}(k)(\operatorname{ring\text{-}power}(\operatorname{normed\text{-}field\text{-}as\text{-}commutative\text{-}ring}(k), w, e), u) \\ \quad \quad \in\; \operatorname{nv\text{-}principal}(k, a, \operatorname{ring\text{-}power}(\operatorname{normed\text{-}field\text{-}as\text{-}commutative\text{-}ring}(k), w, j)) \\ \quad \quad \Leftrightarrow\; j \leq e \end{array}}}
\vnbprop{nv-order-negative}{\vnbstmt{\begin{array}{@{}l@{}} \forall k, a, w.; \forall x \in \operatorname{carr}(k), e \in \mathbb{N}, u \in \operatorname{nv\text{-}units}(k, a), b \in \mathbb{N}. \\ \quad \operatorname{is\text{-}compact\text{-}nv\text{-}domain}(k, a) \wedge \\ \quad w \in a \wedge \\ \quad \operatorname{nv\text{-}maximal}(k, a) = \operatorname{nv\text{-}principal}(k, a, w) \wedge \\ \quad \neg (w = \operatorname{zero}(k)) \wedge \\ \quad \operatorname{mul}(k)(x, \operatorname{ring\text{-}power}(\operatorname{normed\text{-}field\text{-}as\text{-}commutative\text{-}ring}(k), w, e)) = u \Rightarrow \\ \quad \quad (\operatorname{mul}(k)(\operatorname{ring\text{-}power}(\operatorname{normed\text{-}field\text{-}as\text{-}commutative\text{-}ring}(k), w, b), x) \in a) \Leftrightarrow (e \leq b) \end{array}}}

\vnbbookfile{theorem-library/ccr1109a2-nv}{theorem-library/ccr1109a2-nv}
\noindent{\small ccr1109a2-nv.scm -- wave of 2026-11-09/10 (agent CCR-1109-A2), wired 2026-11-10 from scratchpad/ccr1109a2/ccr1109a2-nv.scm. Library: analysis.}\par
\vnbprop{nv-u1-unfold}{\vnbstmt{\begin{array}{@{}l@{}} \forall k, a. \\ \quad (\operatorname{nv\text{-}u1}(k, a) \simeq \{x \in a : (\operatorname{fnrm}(k)(\operatorname{add}(k)(x, \operatorname{neg}(k)(\operatorname{one}(k)))) < 1)\}) \end{array}}}
\vnbprop{nv-complete-domain}{\vnbstmt{\forall k, a.\; (\operatorname{is\text{-}complete\text{-}nv\text{-}domain}(k, a) \Rightarrow \operatorname{is\text{-}nv\text{-}domain}(k, a))}}
\vnbprop{nv-complete-complete}{\vnbstmt{\begin{array}{@{}l@{}} \forall k, a. \\ \quad \operatorname{is\text{-}complete\text{-}nv\text{-}domain}(k, a) \Rightarrow  \\ \quad \quad \operatorname{is\text{-}complete}(\operatorname{subspace\text{-}ms}(\operatorname{nf\text{-}metric\text{-}space}(k), a)) \end{array}}}
\vnbprop{nv-u1-sub}{\vnbstmt{\forall k, a.; \forall x \in \operatorname{nv\text{-}u1}(k, a).\; (x \in a)}}
\vnbprop{nv-u1-norm}{\vnbstmt{\forall k, a.; \forall x \in \operatorname{nv\text{-}u1}(k, a).\; (\operatorname{fnrm}(k)(\operatorname{add}(k)(x, \operatorname{neg}(k)(\operatorname{one}(k)))) < 1)}}
\vnbprop{nv-u1-intro}{\vnbstmt{\begin{array}{@{}l@{}} \forall k, a.; \forall x \in a. \\ \quad (\operatorname{fnrm}(k)(\operatorname{add}(k)(x, \operatorname{neg}(k)(\operatorname{one}(k)))) < 1) \Rightarrow (x \in \operatorname{nv\text{-}u1}(k, a)) \end{array}}}
\vnbprop{nv-u1-norm-one}{\vnbstmt{\forall k, a.; \forall x \in \operatorname{nv\text{-}u1}(k, a).\; (\operatorname{is\text{-}nv\text{-}domain}(k, a) \Rightarrow (\operatorname{fnrm}(k)(x) = 1))}}
\vnbprop{nv-u1-one}{\vnbstmt{\forall k, a.\; (\operatorname{is\text{-}nv\text{-}domain}(k, a) \Rightarrow (\operatorname{one}(k) \in \operatorname{nv\text{-}u1}(k, a)))}}
\vnbprop{nv-u1-mul}{\vnbstmt{\begin{array}{@{}l@{}} \forall k, a.; \forall x, y \in \operatorname{nv\text{-}u1}(k, a). \\ \quad \operatorname{is\text{-}nv\text{-}domain}(k, a) \Rightarrow (\operatorname{mul}(k)(x, y) \in \operatorname{nv\text{-}u1}(k, a)) \end{array}}}
\vnbprop{nv-u1-inverse-in}{\vnbstmt{\begin{array}{@{}l@{}} \forall k, a.; \forall x \in \operatorname{nv\text{-}u1}(k, a), y \in a. \\ \quad (\operatorname{is\text{-}nv\text{-}domain}(k, a) \wedge (\operatorname{mul}(k)(x, y) = \operatorname{one}(k))) \Rightarrow (y \in \operatorname{nv\text{-}u1}(k, a)) \end{array}}}
\vnbprop{ccr-1-47-neumann}{\vnbstmt{\begin{array}{@{}l@{}} \forall k, a.; \forall x \in \operatorname{nv\text{-}u1}(k, a). \\ \quad \operatorname{is\text{-}complete\text{-}nv\text{-}domain}(k, a) \Rightarrow  \\ \quad \quad \exists y \in a. \\ \quad \quad \quad \operatorname{mul}(k)(x, y) = \operatorname{one}(k) \wedge \\ \quad \quad \quad \operatorname{converges\text{-}to}(\operatorname{nf\text{-}metric\text{-}space}(k), \\ \quad \quad \quad \quad (\lambda n \in \mathbb{N}.\; \\ \quad \quad \quad \quad \quad \operatorname{sum\text{-}ag}(\operatorname{ring\text{-}additive\text{-}ag}(\operatorname{normed\text{-}field\text{-}as\text{-}commutative\text{-}ring}(k)), \\ \quad \quad \quad \quad \quad \quad (\lambda b \in \mathbb{N}.\; \operatorname{ring\text{-}power}(\operatorname{normed\text{-}field\text{-}as\text{-}commutative\text{-}ring}(k), \operatorname{add}(k)(\operatorname{one}(k), \operatorname{neg}(k)(x)), b)), \\ \quad \quad \quad \quad \quad \quad n)), \\ \quad \quad \quad \quad y) \end{array}}}
\vnbprop{ccr-1-47-unit}{\vnbstmt{\begin{array}{@{}l@{}} \forall k, a.; \forall x \in \operatorname{nv\text{-}u1}(k, a). \\ \quad \operatorname{is\text{-}complete\text{-}nv\text{-}domain}(k, a) \Rightarrow (x \in \operatorname{nv\text{-}units}(k, a)) \end{array}}}
\vnbprop{ccr-1-47-subgroup}{\vnbstmt{\begin{array}{@{}l@{}} \forall k, a. \\ \quad \operatorname{is\text{-}complete\text{-}nv\text{-}domain}(k, a) \Rightarrow  \\ \quad \quad \operatorname{nv\text{-}u1}(k, a) \subseteq \operatorname{nv\text{-}units}(k, a) \wedge \\ \quad \quad \operatorname{one}(k) \in \operatorname{nv\text{-}u1}(k, a) \wedge \\ \quad \quad \forall x, y \in \operatorname{nv\text{-}u1}(k, a).\; (\operatorname{mul}(k)(x, y) \in \operatorname{nv\text{-}u1}(k, a)) \wedge \\ \quad \quad \forall x \in \operatorname{nv\text{-}u1}(k, a).\; \exists y \in \operatorname{nv\text{-}u1}(k, a).\; (\operatorname{mul}(k)(x, y) = \operatorname{one}(k)) \end{array}}}

\vnbbookfile{theorem-library/ccr1109a2-laurent}{theorem-library/ccr1109a2-laurent}
\noindent{\small ccr1109a2-laurent.scm -- wave of 2026-11-09/10 (agent CCR-1109-A2), wired 2026-11-10 from scratchpad/ccr1109a2/ccr1109a2-laurent.scm. Library: analysis.}\par
\vnbprop{laurent-set-unfold}{\vnbstmt{\forall k.\; (\operatorname{laurent\text{-}set}(k) \simeq \{f \in (\mathbb{Z} \to \operatorname{carr}(k)) : \operatorname{is\text{-}laurent}(k, f)\})}}
\vnbprop{laurent-set-in}{\vnbstmt{\begin{array}{@{}l@{}} \forall k.; \forall f \in (\mathbb{Z} \to \operatorname{carr}(k)), b \in \mathbb{Z}. \\ \quad \forall j \in \mathbb{Z}.\; ((j < b) \Rightarrow (f(j) = \operatorname{zero}(k))) \Rightarrow (f \in \operatorname{laurent\text{-}set}(k)) \end{array}}}
\vnbprop{laurent-set-out}{\vnbstmt{\begin{array}{@{}l@{}} \forall k.; \forall f \in \operatorname{laurent\text{-}set}(k). \\ \quad ((f \in (\mathbb{Z} \to \operatorname{carr}(k))) \wedge \exists b \in \mathbb{Z}.\; \forall j \in \mathbb{Z}.\; ((j < b) \Rightarrow (f(j) = \operatorname{zero}(k)))) \end{array}}}
\vnbprop{laurent-zero-in}{\vnbstmt{\forall k.\; (\operatorname{is\text{-}ring}(k) \Rightarrow (\operatorname{laurent\text{-}zero}(k) \in \operatorname{laurent\text{-}set}(k)))}}
\vnbprop{laurent-one-in}{\vnbstmt{\forall k.\; (\operatorname{is\text{-}ring}(k) \Rightarrow (\operatorname{laurent\text{-}one}(k) \in \operatorname{laurent\text{-}set}(k)))}}
\vnbprop{laurent-add-in}{\vnbstmt{\begin{array}{@{}l@{}} \forall k.; \forall f, g \in \operatorname{laurent\text{-}set}(k). \\ \quad \operatorname{is\text{-}ring}(k) \Rightarrow (\operatorname{laurent\text{-}add}(k, f, g) \in \operatorname{laurent\text{-}set}(k)) \end{array}}}
\vnbprop{laurent-neg-in}{\vnbstmt{\begin{array}{@{}l@{}} \forall k.; \forall f \in \operatorname{laurent\text{-}set}(k). \\ \quad \operatorname{is\text{-}ring}(k) \Rightarrow (\operatorname{laurent\text{-}neg}(k, f) \in \operatorname{laurent\text{-}set}(k)) \end{array}}}
\vnbprop{laurent-ord-exists}{\vnbstmt{\begin{array}{@{}l@{}} \forall k.; \forall f \in \operatorname{laurent\text{-}set}(k), w \in \mathbb{Z}. \\ \quad (\operatorname{is\text{-}ring}(k) \wedge \neg (f(w) = \operatorname{zero}(k))) \Rightarrow  \\ \quad \quad \exists o \in \mathbb{Z}. \\ \quad \quad \quad \neg (f(o) = \operatorname{zero}(k)) \wedge \forall j \in \mathbb{Z}.\; ((j < o) \Rightarrow (f(j) = \operatorname{zero}(k))) \end{array}}}

\vnbbookfile{theorem-library/ccr1109a3-laurent-ring}{theorem-library/ccr1109a3-laurent-ring}
\noindent{\small ccr1109a3-laurent-ring.scm -- wave of 2026-11-09/10 (agent CCR-1109-A3), wired 2026-11-10 from scratchpad/ccr1109a3/ccr1109a3-laurent-ring.scm. Library: analysis.}\par
\vnbprop{zz-add-monoid-is-comm-monoid}{\vnbstmt{\operatorname{is\text{-}comm\text{-}monoid}(\operatorname{zz\text{-}add\text{-}monoid})}}
\vnbprop{zz-add-monoid-is-monoid}{\vnbstmt{\operatorname{is\text{-}monoid}(\operatorname{zz\text{-}add\text{-}monoid})}}
\vnbprop{zz-add-monoid-carr}{\vnbstmt{(\operatorname{carr}(\operatorname{zz\text{-}add\text{-}monoid}) \simeq \mathbb{Z})}}
\vnbprop{zz-add-monoid-iden}{\vnbstmt{(\operatorname{iden}(\operatorname{zz\text{-}add\text{-}monoid}) \simeq 0)}}
\vnbprop{zz-add-monoid-opr}{\vnbstmt{\forall a, b \in \mathbb{Z}.\; (\operatorname{opr}(\operatorname{zz\text{-}add\text{-}monoid})(a, b) = (a + b))}}
\vnbprop{zz-interval-finite}{\vnbstmt{\begin{array}{@{}l@{}} \forall a, b \in \mathbb{Z}. \\ \quad (\{j \in \mathbb{Z} : ((a \leq j) \wedge (j \leq b))\} \in \mathbf{Set}) \wedge (|\{j \in \mathbb{Z} : ((a \leq j) \wedge (j \leq b))\}| \in \mathbb{N}) \end{array}}}
\vnbprop{laurent-trunc-le}{\vnbstmt{\begin{array}{@{}l@{}} \forall q.; \forall \mathit{fq} \in (\mathbb{Z} \to \operatorname{carr}(q)), \mathit{tq}, \mathit{jq} \in \mathbb{Z}. \\ \quad (\operatorname{is\text{-}ring}(q) \wedge (\mathit{jq} \leq \mathit{tq})) \Rightarrow (\operatorname{laurent\text{-}trunc}(q, \mathit{fq}, \mathit{tq})(\mathit{jq}) = \operatorname{fq}(\mathit{jq})) \end{array}}}
\vnbprop{laurent-trunc-gt}{\vnbstmt{\begin{array}{@{}l@{}} \forall q.; \forall \mathit{fq} \in (\mathbb{Z} \to \operatorname{carr}(q)), \mathit{tq}, \mathit{jq} \in \mathbb{Z}. \\ \quad (\operatorname{is\text{-}ring}(q) \wedge (\mathit{tq} < \mathit{jq})) \Rightarrow (\operatorname{laurent\text{-}trunc}(q, \mathit{fq}, \mathit{tq})(\mathit{jq}) = \operatorname{zero}(q)) \end{array}}}
\vnbprop{laurent-trunc-fun}{\vnbstmt{\begin{array}{@{}l@{}} \forall q.; \forall \mathit{fq} \in (\mathbb{Z} \to \operatorname{carr}(q)), \mathit{tq} \in \mathbb{Z}. \\ \quad \operatorname{is\text{-}ring}(q) \Rightarrow (\operatorname{laurent\text{-}trunc}(q, \mathit{fq}, \mathit{tq}) \in (\mathbb{Z} \to \operatorname{carr}(q))) \end{array}}}
\vnbprop{laurent-trunc-low}{\vnbstmt{\begin{array}{@{}l@{}} \forall q.; \forall \mathit{fq} \in (\mathbb{Z} \to \operatorname{carr}(q)), \mathit{bq}, \mathit{tq}, j \in \mathbb{Z}. \\ \quad (\operatorname{is\text{-}ring}(q) \wedge (\forall j \in \mathbb{Z}.\; ((j < \mathit{bq}) \Rightarrow (\operatorname{fq}(j) = \operatorname{zero}(q))) \wedge (j < \mathit{bq}))) \Rightarrow  \\ \quad \quad \operatorname{laurent\text{-}trunc}(q, \mathit{fq}, \mathit{tq})(j) = \operatorname{zero}(q) \end{array}}}
\vnbprop{laurent-trunc-finsupp}{\vnbstmt{\begin{array}{@{}l@{}} \forall q.; \forall \mathit{fq} \in (\mathbb{Z} \to \operatorname{carr}(q)), \mathit{bq}, \mathit{tq} \in \mathbb{Z}. \\ \quad (\operatorname{is\text{-}ring}(q) \wedge \forall j \in \mathbb{Z}.\; ((j < \mathit{bq}) \Rightarrow (\operatorname{fq}(j) = \operatorname{zero}(q)))) \Rightarrow  \\ \quad \quad \operatorname{laurent\text{-}trunc}(q, \mathit{fq}, \mathit{tq}) \in \operatorname{finsupp}(q, \operatorname{zz\text{-}add\text{-}monoid}) \end{array}}}
\vnbprop{laurent-mul-trunc}{\vnbstmt{\begin{array}{@{}l@{}} \forall q.; \forall \mathit{fq}, \mathit{gq} \in (\mathbb{Z} \to \operatorname{carr}(q)), \mathit{bq}, \mathit{nq}, \mathit{tq} \in \mathbb{Z}. \\ \quad \operatorname{is\text{-}ring}(q) \wedge \\ \quad \forall j \in \mathbb{Z}.\; ((j < \mathit{bq}) \Rightarrow (\operatorname{fq}(j) = \operatorname{zero}(q))) \wedge \\ \quad \forall j \in \mathbb{Z}.\; ((j < \mathit{bq}) \Rightarrow (\operatorname{gq}(j) = \operatorname{zero}(q))) \wedge \\ \quad \mathit{nq} \leq (\mathit{tq} + \mathit{bq}) \Rightarrow \\ \quad \quad \operatorname{monalg\text{-}mul}(q, \operatorname{zz\text{-}add\text{-}monoid}, \mathit{fq}, \mathit{gq})(\mathit{nq}) \\ \quad \quad =\; \operatorname{monalg\text{-}mul}(q, \operatorname{zz\text{-}add\text{-}monoid}, \operatorname{laurent\text{-}trunc}(q, \mathit{fq}, \mathit{tq}), \operatorname{laurent\text{-}trunc}(q, \mathit{gq}, \mathit{tq}))(\mathit{nq}) \end{array}}}
\vnbprop{laurent-trunc-agree}{\vnbstmt{\begin{array}{@{}l@{}} \forall q.; \forall \mathit{fq}, r \in (\mathbb{Z} \to \operatorname{carr}(q)), \mathit{tq} \in \mathbb{Z}. \\ \quad (\operatorname{is\text{-}ring}(q) \wedge \forall j \in \mathbb{Z}.\; ((j \leq \mathit{tq}) \Rightarrow (\operatorname{fq}(j) = r(j)))) \Rightarrow  \\ \quad \quad \operatorname{laurent\text{-}trunc}(q, \mathit{fq}, \mathit{tq}) = \operatorname{laurent\text{-}trunc}(q, r, \mathit{tq}) \end{array}}}
\vnbprop{laurent-mul-local}{\vnbstmt{\begin{array}{@{}l@{}} \forall q.; \forall \mathit{fq}, r, \mathit{gq}, \mathit{gr} \in (\mathbb{Z} \to \operatorname{carr}(q)), \mathit{bq}, \mathit{tq}, \mathit{nq} \in \mathbb{Z}. \\ \quad \operatorname{is\text{-}ring}(q) \wedge \\ \quad \forall j \in \mathbb{Z}.\; ((j < \mathit{bq}) \Rightarrow (\operatorname{fq}(j) = \operatorname{zero}(q))) \wedge \\ \quad \forall j \in \mathbb{Z}.\; ((j < \mathit{bq}) \Rightarrow (r(j) = \operatorname{zero}(q))) \wedge \\ \quad \forall j \in \mathbb{Z}.\; ((j < \mathit{bq}) \Rightarrow (\operatorname{gq}(j) = \operatorname{zero}(q))) \wedge \\ \quad \forall j \in \mathbb{Z}.\; ((j < \mathit{bq}) \Rightarrow (\operatorname{gr}(j) = \operatorname{zero}(q))) \wedge \\ \quad \forall j \in \mathbb{Z}.\; ((j \leq \mathit{tq}) \Rightarrow (\operatorname{fq}(j) = r(j))) \wedge \\ \quad \forall j \in \mathbb{Z}.\; ((j \leq \mathit{tq}) \Rightarrow (\operatorname{gq}(j) = \operatorname{gr}(j))) \wedge \\ \quad \mathit{nq} \leq (\mathit{tq} + \mathit{bq}) \Rightarrow \\ \quad \quad \operatorname{monalg\text{-}mul}(q, \operatorname{zz\text{-}add\text{-}monoid}, \mathit{fq}, \mathit{gq})(\mathit{nq}) = \operatorname{monalg\text{-}mul}(q, \operatorname{zz\text{-}add\text{-}monoid}, r, \mathit{gr})(\mathit{nq}) \end{array}}}
\vnbprop{laurent-mul-vanish}{\vnbstmt{\begin{array}{@{}l@{}} \forall q.; \forall \mathit{fq}, \mathit{gq} \in (\mathbb{Z} \to \operatorname{carr}(q)), \mathit{bq}, j \in \mathbb{Z}. \\ \quad \operatorname{is\text{-}ring}(q) \wedge \\ \quad \forall j \in \mathbb{Z}.\; ((j < \mathit{bq}) \Rightarrow (\operatorname{fq}(j) = \operatorname{zero}(q))) \wedge \\ \quad \forall j \in \mathbb{Z}.\; ((j < \mathit{bq}) \Rightarrow (\operatorname{gq}(j) = \operatorname{zero}(q))) \wedge \\ \quad j < (\mathit{bq} + \mathit{bq}) \Rightarrow \\ \quad \quad \operatorname{monalg\text{-}mul}(q, \operatorname{zz\text{-}add\text{-}monoid}, \mathit{fq}, \mathit{gq})(j) = \operatorname{zero}(q) \end{array}}}
\vnbprop{laurent-mul-fun}{\vnbstmt{\begin{array}{@{}l@{}} \forall q.; \forall \mathit{fq}, \mathit{gq} \in (\mathbb{Z} \to \operatorname{carr}(q)), \mathit{bq} \in \mathbb{Z}. \\ \quad \operatorname{is\text{-}ring}(q) \wedge \\ \quad \forall j \in \mathbb{Z}.\; ((j < \mathit{bq}) \Rightarrow (\operatorname{fq}(j) = \operatorname{zero}(q))) \wedge \\ \quad \forall j \in \mathbb{Z}.\; ((j < \mathit{bq}) \Rightarrow (\operatorname{gq}(j) = \operatorname{zero}(q))) \Rightarrow \\ \quad \quad \operatorname{monalg\text{-}mul}(q, \operatorname{zz\text{-}add\text{-}monoid}, \mathit{fq}, \mathit{gq}) \in (\mathbb{Z} \to \operatorname{carr}(q)) \end{array}}}
\vnbprop{laurent-below-mono}{\vnbstmt{\begin{array}{@{}l@{}} \forall q, \mathit{fq}.; \forall \mathit{bq}, \mathit{cq}, j \in \mathbb{Z}. \\ \quad (\forall j \in \mathbb{Z}.\; ((j < \mathit{bq}) \Rightarrow (\operatorname{fq}(j) = \operatorname{zero}(q))) \wedge ((\mathit{cq} \leq \mathit{bq}) \wedge (j < \mathit{cq}))) \Rightarrow  \\ \quad \quad \operatorname{fq}(j) = \operatorname{zero}(q) \end{array}}}
\vnbprop{laurent-bound-le}{\vnbstmt{\begin{array}{@{}l@{}} \forall q.; \forall \mathit{fq} \in \operatorname{laurent\text{-}set}(q), \mathit{cq} \in \mathbb{Z}. \\ \quad \operatorname{is\text{-}ring}(q) \Rightarrow  \\ \quad \quad \exists \mathit{bq} \in \mathbb{Z}. \\ \quad \quad \quad (\mathit{bq} \leq \mathit{cq}) \wedge \forall j \in \mathbb{Z}.\; ((j < \mathit{bq}) \Rightarrow (\operatorname{fq}(j) = \operatorname{zero}(q))) \end{array}}}
\vnbprop{laurent-add-monalg}{\vnbstmt{\begin{array}{@{}l@{}} \forall q, \mathit{fq}, \mathit{gq}. \\ \quad (\operatorname{laurent\text{-}add}(q, \mathit{fq}, \mathit{gq}) \simeq \operatorname{monalg\text{-}add}(q, \operatorname{zz\text{-}add\text{-}monoid}, \mathit{fq}, \mathit{gq})) \end{array}}}
\vnbprop{laurent-one-monalg}{\vnbstmt{\forall q.\; (\operatorname{laurent\text{-}one}(q) \simeq \operatorname{monalg\text{-}one}(q, \operatorname{zz\text{-}add\text{-}monoid}))}}
\vnbprop{laurent-add-apply}{\vnbstmt{\begin{array}{@{}l@{}} \forall q, \mathit{fq}, \mathit{gq}.; \forall \mathit{jq} \in \mathbb{Z}. \\ \quad (\operatorname{laurent\text{-}add}(q, \mathit{fq}, \mathit{gq})(\mathit{jq}) \simeq \operatorname{add}(q)(\operatorname{fq}(\mathit{jq}), \operatorname{gq}(\mathit{jq}))) \end{array}}}
\vnbprop{laurent-neg-apply}{\vnbstmt{\forall q, \mathit{fq}.; \forall \mathit{jq} \in \mathbb{Z}.\; (\operatorname{laurent\text{-}neg}(q, \mathit{fq})(\mathit{jq}) \simeq \operatorname{neg}(q)(\operatorname{fq}(\mathit{jq})))}}
\vnbprop{laurent-zero-apply}{\vnbstmt{\forall q.; \forall \mathit{jq} \in \mathbb{Z}.\; (\operatorname{laurent\text{-}zero}(q)(\mathit{jq}) \simeq \operatorname{zero}(q))}}
\vnbprop{laurent-one-low}{\vnbstmt{\begin{array}{@{}l@{}} \forall q.; \forall \mathit{bq}, j \in \mathbb{Z}. \\ \quad (\operatorname{is\text{-}ring}(q) \wedge ((\mathit{bq} \leq 0) \wedge (j < \mathit{bq}))) \Rightarrow (\operatorname{laurent\text{-}one}(q)(j) = \operatorname{zero}(q)) \end{array}}}
\vnbprop{laurent-mul-in}{\vnbstmt{\begin{array}{@{}l@{}} \forall q.; \forall \mathit{fq}, \mathit{gq} \in \operatorname{laurent\text{-}set}(q). \\ \quad \operatorname{is\text{-}ring}(q) \Rightarrow (\operatorname{laurent\text{-}mul}(q, \mathit{fq}, \mathit{gq}) \in \operatorname{laurent\text{-}set}(q)) \end{array}}}
\vnbprop{laurent-mul-comm}{\vnbstmt{\begin{array}{@{}l@{}} \forall q.; \forall \mathit{fq}, \mathit{gq} \in \operatorname{laurent\text{-}set}(q). \\ \quad \operatorname{is\text{-}commutative\text{-}ring}(q) \Rightarrow (\operatorname{laurent\text{-}mul}(q, \mathit{fq}, \mathit{gq}) = \operatorname{laurent\text{-}mul}(q, \mathit{gq}, \mathit{fq})) \end{array}}}
\vnbprop{laurent-mul-one}{\vnbstmt{\begin{array}{@{}l@{}} \forall q.; \forall \mathit{fq} \in \operatorname{laurent\text{-}set}(q). \\ \quad \operatorname{is\text{-}ring}(q) \Rightarrow (\operatorname{laurent\text{-}mul}(q, \operatorname{laurent\text{-}one}(q), \mathit{fq}) = \mathit{fq}) \end{array}}}
\vnbprop{laurent-mul-assoc}{\vnbstmt{\begin{array}{@{}l@{}} \forall q.; \forall \mathit{fq}, \mathit{gq}, \mathit{hq} \in \operatorname{laurent\text{-}set}(q). \\ \quad \operatorname{is\text{-}ring}(q) \Rightarrow  \\ \quad \quad \operatorname{laurent\text{-}mul}(q, \operatorname{laurent\text{-}mul}(q, \mathit{fq}, \mathit{gq}), \mathit{hq}) = \operatorname{laurent\text{-}mul}(q, \mathit{fq}, \operatorname{laurent\text{-}mul}(q, \mathit{gq}, \mathit{hq})) \end{array}}}
\vnbprop{laurent-mul-distrib-left}{\vnbstmt{\begin{array}{@{}l@{}} \forall q.; \forall \mathit{fq}, \mathit{gq}, \mathit{hq} \in \operatorname{laurent\text{-}set}(q). \\ \quad \operatorname{is\text{-}ring}(q) \Rightarrow  \\ \quad \quad \operatorname{laurent\text{-}mul}(q, \mathit{fq}, \operatorname{laurent\text{-}add}(q, \mathit{gq}, \mathit{hq})) \\ \quad \quad =\; \operatorname{laurent\text{-}add}(q, \operatorname{laurent\text{-}mul}(q, \mathit{fq}, \mathit{gq}), \operatorname{laurent\text{-}mul}(q, \mathit{fq}, \mathit{hq})) \end{array}}}
\vnbprop{laurent-add-assoc}{\vnbstmt{\begin{array}{@{}l@{}} \forall q.; \forall \mathit{fq}, \mathit{gq}, \mathit{hq} \in \operatorname{laurent\text{-}set}(q). \\ \quad \operatorname{is\text{-}ring}(q) \Rightarrow  \\ \quad \quad \operatorname{laurent\text{-}add}(q, \operatorname{laurent\text{-}add}(q, \mathit{fq}, \mathit{gq}), \mathit{hq}) = \operatorname{laurent\text{-}add}(q, \mathit{fq}, \operatorname{laurent\text{-}add}(q, \mathit{gq}, \mathit{hq})) \end{array}}}
\vnbprop{laurent-add-comm}{\vnbstmt{\begin{array}{@{}l@{}} \forall q.; \forall \mathit{fq}, \mathit{gq} \in \operatorname{laurent\text{-}set}(q). \\ \quad \operatorname{is\text{-}ring}(q) \Rightarrow (\operatorname{laurent\text{-}add}(q, \mathit{fq}, \mathit{gq}) = \operatorname{laurent\text{-}add}(q, \mathit{gq}, \mathit{fq})) \end{array}}}
\vnbprop{laurent-add-zero-left}{\vnbstmt{\begin{array}{@{}l@{}} \forall q.; \forall \mathit{fq} \in \operatorname{laurent\text{-}set}(q). \\ \quad \operatorname{is\text{-}ring}(q) \Rightarrow (\operatorname{laurent\text{-}add}(q, \operatorname{laurent\text{-}zero}(q), \mathit{fq}) = \mathit{fq}) \end{array}}}
\vnbprop{laurent-add-neg-left}{\vnbstmt{\begin{array}{@{}l@{}} \forall q.; \forall \mathit{fq} \in \operatorname{laurent\text{-}set}(q). \\ \quad \operatorname{is\text{-}ring}(q) \Rightarrow (\operatorname{laurent\text{-}add}(q, \operatorname{laurent\text{-}neg}(q, \mathit{fq}), \mathit{fq}) = \operatorname{laurent\text{-}zero}(q)) \end{array}}}
\vnbprop{laurent-ring-carr}{\vnbstmt{\forall q.\; (\operatorname{is\text{-}ring}(q) \Rightarrow (\operatorname{carr}(\operatorname{laurent\text{-}ring}(q)) \simeq \operatorname{laurent\text{-}set}(q)))}}
\vnbprop{laurent-ring-add}{\vnbstmt{\begin{array}{@{}l@{}} \forall q. \\ \quad \operatorname{is\text{-}ring}(q) \Rightarrow  \\ \quad \quad (\operatorname{add}(\operatorname{laurent\text{-}ring}(q)) \simeq (\lambda {[}\operatorname{lrf\_}, \operatorname{lrg\_}] \in (\operatorname{laurent\text{-}set}(q) \times \operatorname{laurent\text{-}set}(q)).\; \operatorname{laurent\text{-}add}(q, \operatorname{lrf\_}, \operatorname{lrg\_}))) \end{array}}}
\vnbprop{laurent-ring-mul}{\vnbstmt{\begin{array}{@{}l@{}} \forall q. \\ \quad \operatorname{is\text{-}ring}(q) \Rightarrow  \\ \quad \quad (\operatorname{mul}(\operatorname{laurent\text{-}ring}(q)) \simeq (\lambda {[}\operatorname{lrf\_}, \operatorname{lrg\_}] \in (\operatorname{laurent\text{-}set}(q) \times \operatorname{laurent\text{-}set}(q)).\; \operatorname{laurent\text{-}mul}(q, \operatorname{lrf\_}, \operatorname{lrg\_}))) \end{array}}}
\vnbprop{laurent-ring-neg}{\vnbstmt{\begin{array}{@{}l@{}} \forall q. \\ \quad \operatorname{is\text{-}ring}(q) \Rightarrow  \\ \quad \quad (\operatorname{neg}(\operatorname{laurent\text{-}ring}(q)) \simeq (\lambda f \in \operatorname{laurent\text{-}set}(q).\; \operatorname{laurent\text{-}neg}(q, f))) \end{array}}}
\vnbprop{laurent-ring-zero}{\vnbstmt{\forall q.\; (\operatorname{is\text{-}ring}(q) \Rightarrow (\operatorname{zero}(\operatorname{laurent\text{-}ring}(q)) \simeq \operatorname{laurent\text{-}zero}(q)))}}
\vnbprop{laurent-ring-one}{\vnbstmt{\forall q.\; (\operatorname{is\text{-}ring}(q) \Rightarrow (\operatorname{one}(\operatorname{laurent\text{-}ring}(q)) \simeq \operatorname{laurent\text{-}one}(q)))}}
\vnbprop{laurent-ring-is-ring}{\vnbstmt{\forall q.\; (\operatorname{is\text{-}commutative\text{-}ring}(q) \Rightarrow \operatorname{is\text{-}ring}(\operatorname{laurent\text{-}ring}(q)))}}
\vnbprop{laurent-ring-is-commutative-ring}{\vnbstmt{\forall q.\; (\operatorname{is\text{-}commutative\text{-}ring}(q) \Rightarrow \operatorname{is\text{-}commutative\text{-}ring}(\operatorname{laurent\text{-}ring}(q)))}}

\vnbbookfile{theorem-library/ccr1109a3-laurent-val}{theorem-library/ccr1109a3-laurent-val}
\noindent{\small ccr1109a3-laurent-val.scm -- wave of 2026-11-09/10 (agent CCR-1109-A3), wired 2026-11-10 from scratchpad/ccr1109a3/ccr1109a3-laurent-val.scm. Library: analysis.}\par
\vnbprop{laurent-mul-vanish2}{\vnbstmt{\begin{array}{@{}l@{}} \forall q.; \forall \mathit{fq}, \mathit{gq} \in (\mathbb{Z} \to \operatorname{carr}(q)), \mathit{pq}, r, j \in \mathbb{Z}. \\ \quad \operatorname{is\text{-}ring}(q) \wedge \\ \quad \forall j \in \mathbb{Z}.\; ((j < \mathit{pq}) \Rightarrow (\operatorname{fq}(j) = \operatorname{zero}(q))) \wedge \\ \quad \forall j \in \mathbb{Z}.\; ((j < r) \Rightarrow (\operatorname{gq}(j) = \operatorname{zero}(q))) \wedge \\ \quad j < (\mathit{pq} + r) \Rightarrow \\ \quad \quad \operatorname{monalg\text{-}mul}(q, \operatorname{zz\text{-}add\text{-}monoid}, \mathit{fq}, \mathit{gq})(j) = \operatorname{zero}(q) \end{array}}}
\vnbprop{laurent-mul-single}{\vnbstmt{\begin{array}{@{}l@{}} \forall q.; \forall \mathit{fq}, \mathit{gq} \in (\mathbb{Z} \to \operatorname{carr}(q)), \mathit{bq}, \mathit{pq}, r, \mathit{nq} \in \mathbb{Z}. \\ \quad \operatorname{is\text{-}ring}(q) \wedge \\ \quad \forall j \in \mathbb{Z}.\; ((j < \mathit{bq}) \Rightarrow (\operatorname{fq}(j) = \operatorname{zero}(q))) \wedge \\ \quad \forall j \in \mathbb{Z}.\; ((j < \mathit{bq}) \Rightarrow (\operatorname{gq}(j) = \operatorname{zero}(q))) \wedge \\ \quad (\mathit{pq} + r) = \mathit{nq} \wedge \\ \quad \forall s, \mathit{qs} \in \mathbb{Z}. \\ \quad \quad (((s + \mathit{qs}) = \mathit{nq}) \wedge \neg (s = \mathit{pq})) \Rightarrow ((\operatorname{fq}(s) = \operatorname{zero}(q)) \vee (\operatorname{gq}(\mathit{qs}) = \operatorname{zero}(q))) \Rightarrow \\ \quad \quad \operatorname{monalg\text{-}mul}(q, \operatorname{zz\text{-}add\text{-}monoid}, \mathit{fq}, \mathit{gq})(\mathit{nq}) = \operatorname{mul}(q)(\operatorname{fq}(\mathit{pq}), \operatorname{gq}(r)) \end{array}}}
\vnbprop{laurent-mul-lowest}{\vnbstmt{\begin{array}{@{}l@{}} \forall q.; \forall \mathit{fq}, \mathit{gq} \in (\mathbb{Z} \to \operatorname{carr}(q)), \mathit{pq}, r \in \mathbb{Z}. \\ \quad \operatorname{is\text{-}ring}(q) \wedge \\ \quad \forall j \in \mathbb{Z}.\; ((j < \mathit{pq}) \Rightarrow (\operatorname{fq}(j) = \operatorname{zero}(q))) \wedge \\ \quad \forall j \in \mathbb{Z}.\; ((j < r) \Rightarrow (\operatorname{gq}(j) = \operatorname{zero}(q))) \Rightarrow \\ \quad \quad \operatorname{monalg\text{-}mul}(q, \operatorname{zz\text{-}add\text{-}monoid}, \mathit{fq}, \mathit{gq})((\mathit{pq} + r)) = \operatorname{mul}(q)(\operatorname{fq}(\mathit{pq}), \operatorname{gq}(r)) \end{array}}}
\vnbprop{laurent-ord-unique}{\vnbstmt{\begin{array}{@{}l@{}} \forall q, \mathit{fq}, \mathit{oq}, r. \\ \quad \operatorname{is\text{-}ring}(q) \wedge \\ \quad \mathit{oq} \in \mathbb{Z} \wedge \\ \quad \neg (\operatorname{fq}(\mathit{oq}) = \operatorname{zero}(q)) \wedge \\ \quad \forall j \in \mathbb{Z}.\; ((j < \mathit{oq}) \Rightarrow (\operatorname{fq}(j) = \operatorname{zero}(q))) \wedge \\ \quad r \in \mathbb{Z} \wedge \\ \quad \neg (\operatorname{fq}(r) = \operatorname{zero}(q)) \wedge \\ \quad \forall j \in \mathbb{Z}.\; ((j < r) \Rightarrow (\operatorname{fq}(j) = \operatorname{zero}(q))) \Rightarrow \\ \quad \quad \mathit{oq} = r \end{array}}}
\vnbprop{laurent-ord-eq}{\vnbstmt{\begin{array}{@{}l@{}} \forall q, \mathit{fq}.; \forall \mathit{oq} \in \mathbb{Z}. \\ \quad (\operatorname{is\text{-}ring}(q) \wedge (\neg (\operatorname{fq}(\mathit{oq}) = \operatorname{zero}(q)) \wedge \forall j \in \mathbb{Z}.\; ((j < \mathit{oq}) \Rightarrow (\operatorname{fq}(j) = \operatorname{zero}(q))))) \Rightarrow  \\ \quad \quad \operatorname{laurent\text{-}ord}(q, \mathit{fq}) = \mathit{oq} \end{array}}}
\vnbprop{laurent-nonzero-ord}{\vnbstmt{\begin{array}{@{}l@{}} \forall q.; \forall \mathit{fq} \in \operatorname{laurent\text{-}set}(q). \\ \quad (\operatorname{is\text{-}ring}(q) \wedge \neg (\mathit{fq} = \operatorname{laurent\text{-}zero}(q))) \Rightarrow  \\ \quad \quad \exists \mathit{oq} \in \mathbb{Z}. \\ \quad \quad \quad \neg (\operatorname{fq}(\mathit{oq}) = \operatorname{zero}(q)) \wedge \forall j \in \mathbb{Z}.\; ((j < \mathit{oq}) \Rightarrow (\operatorname{fq}(j) = \operatorname{zero}(q))) \end{array}}}
\vnbprop{laurent-ord-mul}{\vnbstmt{\begin{array}{@{}l@{}} \forall q.; \forall \mathit{fq}, \mathit{gq} \in \operatorname{laurent\text{-}set}(q), \mathit{oq}, r. \\ \quad \operatorname{is\text{-}integral\text{-}domain}(q) \wedge \\ \quad \mathit{oq} \in \mathbb{Z} \wedge \\ \quad \neg (\operatorname{fq}(\mathit{oq}) = \operatorname{zero}(q)) \wedge \\ \quad \forall j \in \mathbb{Z}.\; ((j < \mathit{oq}) \Rightarrow (\operatorname{fq}(j) = \operatorname{zero}(q))) \wedge \\ \quad r \in \mathbb{Z} \wedge \\ \quad \neg (\operatorname{gq}(r) = \operatorname{zero}(q)) \wedge \\ \quad \forall j \in \mathbb{Z}.\; ((j < r) \Rightarrow (\operatorname{gq}(j) = \operatorname{zero}(q))) \Rightarrow \\ \quad \quad (\mathit{oq} + r) \in \mathbb{Z} \wedge \\ \quad \quad \neg (\operatorname{laurent\text{-}mul}(q, \mathit{fq}, \mathit{gq})((\mathit{oq} + r)) = \operatorname{zero}(q)) \wedge \\ \quad \quad \forall j \in \mathbb{Z}.\; ((j < (\mathit{oq} + r)) \Rightarrow (\operatorname{laurent\text{-}mul}(q, \mathit{fq}, \mathit{gq})(j) = \operatorname{zero}(q))) \end{array}}}
\vnbprop{laurent-ord-add}{\vnbstmt{\begin{array}{@{}l@{}} \forall q.; \forall \mathit{fq}, \mathit{gq} \in \operatorname{laurent\text{-}set}(q), \mathit{oq}, r, s.; \forall \mathit{cq} \in \mathbb{Z}. \\ \quad \operatorname{is\text{-}ring}(q) \wedge \\ \quad \mathit{oq} \in \mathbb{Z} \wedge \\ \quad \neg (\operatorname{fq}(\mathit{oq}) = \operatorname{zero}(q)) \wedge \\ \quad \forall j \in \mathbb{Z}.\; ((j < \mathit{oq}) \Rightarrow (\operatorname{fq}(j) = \operatorname{zero}(q))) \wedge \\ \quad r \in \mathbb{Z} \wedge \\ \quad \neg (\operatorname{gq}(r) = \operatorname{zero}(q)) \wedge \\ \quad \forall j \in \mathbb{Z}.\; ((j < r) \Rightarrow (\operatorname{gq}(j) = \operatorname{zero}(q))) \wedge \\ \quad s \in \mathbb{Z} \wedge \\ \quad \neg (\operatorname{laurent\text{-}add}(q, \mathit{fq}, \mathit{gq})(s) = \operatorname{zero}(q)) \wedge \\ \quad \forall j \in \mathbb{Z}.\; ((j < s) \Rightarrow (\operatorname{laurent\text{-}add}(q, \mathit{fq}, \mathit{gq})(j) = \operatorname{zero}(q))) \wedge \\ \quad \mathit{cq} \leq \mathit{oq} \wedge \\ \quad \mathit{cq} \leq r \Rightarrow \\ \quad \quad \mathit{cq} \leq s \end{array}}}
\vnbprop{laurent-ord-mul-eq}{\vnbstmt{\begin{array}{@{}l@{}} \forall q.; \forall \mathit{fq}, \mathit{gq} \in \operatorname{laurent\text{-}set}(q). \\ \quad (\operatorname{is\text{-}integral\text{-}domain}(q) \wedge (\neg (\mathit{fq} = \operatorname{laurent\text{-}zero}(q)) \wedge \neg (\mathit{gq} = \operatorname{laurent\text{-}zero}(q)))) \Rightarrow  \\ \quad \quad \neg (\operatorname{laurent\text{-}mul}(q, \mathit{fq}, \mathit{gq}) = \operatorname{laurent\text{-}zero}(q)) \wedge \\ \quad \quad \operatorname{laurent\text{-}ord}(q, \operatorname{laurent\text{-}mul}(q, \mathit{fq}, \mathit{gq})) = (\operatorname{laurent\text{-}ord}(q, \mathit{fq}) + \operatorname{laurent\text{-}ord}(q, \mathit{gq})) \end{array}}}
\vnbprop{laurent-mul-zero-left}{\vnbstmt{\begin{array}{@{}l@{}} \forall q.; \forall \mathit{gq} \in \operatorname{laurent\text{-}set}(q). \\ \quad \operatorname{is\text{-}ring}(q) \Rightarrow (\operatorname{laurent\text{-}mul}(q, \operatorname{laurent\text{-}zero}(q), \mathit{gq}) = \operatorname{laurent\text{-}zero}(q)) \end{array}}}
\vnbprop{rpow-star-anti-exp}{\vnbstmt{\begin{array}{@{}l@{}} \forall q, \mathit{sq}, \mathit{tq} \in \mathbb{R}. \\ \quad ((0 < q) \wedge ((q < 1) \wedge (\mathit{sq} \leq \mathit{tq}))) \Rightarrow ({q}^{\mathit{tq}} \leq {q}^{\mathit{sq}}) \end{array}}}
\vnbprop{laurent-abs-zero}{\vnbstmt{\forall q, \mathit{rq}.\; (\operatorname{is\text{-}ring}(q) \Rightarrow (\operatorname{laurent\text{-}abs}(q, \mathit{rq}, \operatorname{laurent\text{-}zero}(q)) = 0))}}
\vnbprop{laurent-abs-nonzero}{\vnbstmt{\begin{array}{@{}l@{}} \forall q.; \forall \mathit{rq} \in \mathbb{R}, \mathit{fq} \in \operatorname{laurent\text{-}set}(q). \\ \quad (\operatorname{is\text{-}ring}(q) \wedge ((0 < \mathit{rq}) \wedge \neg (\mathit{fq} = \operatorname{laurent\text{-}zero}(q)))) \Rightarrow  \\ \quad \quad ((\operatorname{laurent\text{-}abs}(q, \mathit{rq}, \mathit{fq}) = {\mathit{rq}}^{\operatorname{laurent\text{-}ord}(q, \mathit{fq})}) \wedge (0 < \operatorname{laurent\text{-}abs}(q, \mathit{rq}, \mathit{fq}))) \end{array}}}
\vnbprop{laurent-abs-in-rr}{\vnbstmt{\begin{array}{@{}l@{}} \forall q.; \forall \mathit{rq} \in \mathbb{R}, \mathit{fq} \in \operatorname{laurent\text{-}set}(q). \\ \quad (\operatorname{is\text{-}ring}(q) \wedge (0 < \mathit{rq})) \Rightarrow  \\ \quad \quad (\operatorname{laurent\text{-}abs}(q, \mathit{rq}, \mathit{fq}) \in \mathbb{R}) \wedge (0 \leq \operatorname{laurent\text{-}abs}(q, \mathit{rq}, \mathit{fq})) \end{array}}}
\vnbprop{laurent-abs-mul}{\vnbstmt{\begin{array}{@{}l@{}} \forall q.; \forall \mathit{rq} \in \mathbb{R}, \mathit{fq}, \mathit{gq} \in \operatorname{laurent\text{-}set}(q). \\ \quad (\operatorname{is\text{-}integral\text{-}domain}(q) \wedge (0 < \mathit{rq})) \Rightarrow  \\ \quad \quad \operatorname{laurent\text{-}abs}(q, \mathit{rq}, \operatorname{laurent\text{-}mul}(q, \mathit{fq}, \mathit{gq})) = (\operatorname{laurent\text{-}abs}(q, \mathit{rq}, \mathit{fq}) \cdot \operatorname{laurent\text{-}abs}(q, \mathit{rq}, \mathit{gq})) \end{array}}}
\vnbprop{laurent-abs-ultra}{\vnbstmt{\begin{array}{@{}l@{}} \forall q.; \forall \mathit{rq} \in \mathbb{R}, \mathit{fq}, \mathit{gq} \in \operatorname{laurent\text{-}set}(q). \\ \quad (\operatorname{is\text{-}ring}(q) \wedge ((0 < \mathit{rq}) \wedge (\mathit{rq} < 1))) \Rightarrow  \\ \quad \quad \operatorname{laurent\text{-}abs}(q, \mathit{rq}, \operatorname{laurent\text{-}add}(q, \mathit{fq}, \mathit{gq})) \leq \operatorname{max}(\operatorname{laurent\text{-}abs}(q, \mathit{rq}, \mathit{fq}), \operatorname{laurent\text{-}abs}(q, \mathit{rq}, \mathit{gq})) \end{array}}}

\vnbbookfile{theorem-library/ccr1109a3-laurent-field}{theorem-library/ccr1109a3-laurent-field}
\noindent{\small ccr1109a3-laurent-field.scm -- wave of 2026-11-09/10 (agent CCR-1109-A3), wired 2026-11-10 from scratchpad/ccr1109a3/ccr1109a3-laurent-field.scm. Library: analysis.}\par
\vnbprop{laurent-mono-apply}{\vnbstmt{\begin{array}{@{}l@{}} \forall q.; \forall \mathit{cq} \in \operatorname{carr}(q), \mathit{tq}.; \forall \mathit{jq} \in \mathbb{Z}. \\ \quad \operatorname{is\text{-}ring}(q) \Rightarrow (\operatorname{laurent\text{-}mono}(q, \mathit{cq}, \mathit{tq})(\mathit{jq}) \simeq \left(\text{if }(\mathit{jq} = \mathit{tq})\text{ then }\mathit{cq}\text{ else }\operatorname{zero}(q)\right)) \end{array}}}
\vnbprop{laurent-mono-fun}{\vnbstmt{\begin{array}{@{}l@{}} \forall q.; \forall \mathit{cq} \in \operatorname{carr}(q), \mathit{tq} \in \mathbb{Z}. \\ \quad \operatorname{is\text{-}ring}(q) \Rightarrow (\operatorname{laurent\text{-}mono}(q, \mathit{cq}, \mathit{tq}) \in (\mathbb{Z} \to \operatorname{carr}(q))) \end{array}}}
\vnbprop{laurent-mono-at}{\vnbstmt{\begin{array}{@{}l@{}} \forall q.; \forall \mathit{cq} \in \operatorname{carr}(q), \mathit{tq} \in \mathbb{Z}. \\ \quad \operatorname{is\text{-}ring}(q) \Rightarrow (\operatorname{laurent\text{-}mono}(q, \mathit{cq}, \mathit{tq})(\mathit{tq}) = \mathit{cq}) \end{array}}}
\vnbprop{laurent-mono-below}{\vnbstmt{\begin{array}{@{}l@{}} \forall q.; \forall \mathit{cq} \in \operatorname{carr}(q), \mathit{tq}, j \in \mathbb{Z}. \\ \quad (\operatorname{is\text{-}ring}(q) \wedge (j < \mathit{tq})) \Rightarrow (\operatorname{laurent\text{-}mono}(q, \mathit{cq}, \mathit{tq})(j) = \operatorname{zero}(q)) \end{array}}}
\vnbprop{laurent-mono-in}{\vnbstmt{\begin{array}{@{}l@{}} \forall q.; \forall \mathit{cq} \in \operatorname{carr}(q), \mathit{tq} \in \mathbb{Z}. \\ \quad \operatorname{is\text{-}ring}(q) \Rightarrow (\operatorname{laurent\text{-}mono}(q, \mathit{cq}, \mathit{tq}) \in \operatorname{laurent\text{-}set}(q)) \end{array}}}
\vnbprop{laurent-add-below}{\vnbstmt{\begin{array}{@{}l@{}} \forall q.; \forall \mathit{fq}, \mathit{gq} \in (\mathbb{Z} \to \operatorname{carr}(q)), \mathit{bq}, j \in \mathbb{Z}. \\ \quad \operatorname{is\text{-}ring}(q) \wedge \\ \quad \forall j \in \mathbb{Z}.\; ((j < \mathit{bq}) \Rightarrow (\operatorname{fq}(j) = \operatorname{zero}(q))) \wedge \\ \quad \forall j \in \mathbb{Z}.\; ((j < \mathit{bq}) \Rightarrow (\operatorname{gq}(j) = \operatorname{zero}(q))) \wedge \\ \quad j < \mathit{bq} \Rightarrow \\ \quad \quad \operatorname{laurent\text{-}add}(q, \mathit{fq}, \mathit{gq})(j) = \operatorname{zero}(q) \end{array}}}
\vnbprop{laurent-inv-pre-in}{\vnbstmt{\begin{array}{@{}l@{}} \forall q \in \mathbb{N}, \mathit{kq}.; \forall \mathit{fq} \in \operatorname{laurent\text{-}set}(\mathit{kq}), \mathit{oq} \in \mathbb{Z}, \mathit{aq} \in \operatorname{carr}(\mathit{kq}). \\ \quad \operatorname{is\text{-}ring}(\mathit{kq}) \Rightarrow  \\ \quad \quad \operatorname{laurent\text{-}inv\text{-}pre}(\mathit{kq}, \mathit{fq}, \mathit{oq}, \mathit{aq}, q) \in \operatorname{laurent\text{-}set}(\mathit{kq}) \wedge \\ \quad \quad \forall j \in \mathbb{Z}. \\ \quad \quad \quad (j < -\mathit{oq}) \Rightarrow (\operatorname{laurent\text{-}inv\text{-}pre}(\mathit{kq}, \mathit{fq}, \mathit{oq}, \mathit{aq}, q)(j) = \operatorname{zero}(\mathit{kq})) \end{array}}}
\vnbprop{laurent-inv-pre-step}{\vnbstmt{\begin{array}{@{}l@{}} \forall q.; \forall \mathit{fq} \in \operatorname{laurent\text{-}set}(q), \mathit{oq} \in \mathbb{Z}, \mathit{aq} \in \operatorname{carr}(q), \mathit{mq} \in \mathbb{N}, \mathit{jq} \in \mathbb{Z}. \\ \quad (\operatorname{is\text{-}ring}(q) \wedge (\mathit{jq} \leq (\mathit{mq} + -\mathit{oq}))) \Rightarrow  \\ \quad \quad \operatorname{laurent\text{-}inv\text{-}pre}(q, \mathit{fq}, \mathit{oq}, \mathit{aq}, (\mathit{mq} + 1))(\mathit{jq}) = \operatorname{laurent\text{-}inv\text{-}pre}(q, \mathit{fq}, \mathit{oq}, \mathit{aq}, \mathit{mq})(\mathit{jq}) \end{array}}}
\vnbprop{laurent-inv-pre-stable}{\vnbstmt{\begin{array}{@{}l@{}} \forall b \in \mathbb{N}, q.; \forall \mathit{fq} \in \operatorname{laurent\text{-}set}(q), \mathit{oq} \in \mathbb{Z}, \mathit{aq} \in \operatorname{carr}(q), \mathit{mq} \in \mathbb{N}, \mathit{jq} \in \mathbb{Z}. \\ \quad (\operatorname{is\text{-}ring}(q) \wedge ((\mathit{mq} \leq b) \wedge (\mathit{jq} \leq (\mathit{mq} + -\mathit{oq})))) \Rightarrow  \\ \quad \quad \operatorname{laurent\text{-}inv\text{-}pre}(q, \mathit{fq}, \mathit{oq}, \mathit{aq}, b)(\mathit{jq}) = \operatorname{laurent\text{-}inv\text{-}pre}(q, \mathit{fq}, \mathit{oq}, \mathit{aq}, \mathit{mq})(\mathit{jq}) \end{array}}}
\vnbprop{laurent-one-apply}{\vnbstmt{\begin{array}{@{}l@{}} \forall q.; \forall \mathit{jq} \in \mathbb{Z}. \\ \quad \operatorname{is\text{-}ring}(q) \Rightarrow (\operatorname{laurent\text{-}one}(q)(\mathit{jq}) \simeq \left(\text{if }(\mathit{jq} = 0)\text{ then }\operatorname{one}(q)\text{ else }\operatorname{zero}(q)\right)) \end{array}}}
\vnbprop{laurent-one-zero}{\vnbstmt{\forall q.\; (\operatorname{is\text{-}ring}(q) \Rightarrow (\operatorname{laurent\text{-}one}(q)(0) = \operatorname{one}(q)))}}
\vnbprop{laurent-one-off}{\vnbstmt{\begin{array}{@{}l@{}} \forall q.; \forall \mathit{jq} \in \mathbb{Z}. \\ \quad (\operatorname{is\text{-}ring}(q) \wedge \neg (\mathit{jq} = 0)) \Rightarrow (\operatorname{laurent\text{-}one}(q)(\mathit{jq}) = \operatorname{zero}(q)) \end{array}}}
\vnbprop{laurent-mul-mono-top}{\vnbstmt{\begin{array}{@{}l@{}} \forall q.; \forall \mathit{fq} \in (\mathbb{Z} \to \operatorname{carr}(q)), \mathit{oq} \in \mathbb{Z}, \mathit{cq} \in \operatorname{carr}(q), \mathit{tq} \in \mathbb{Z}. \\ \quad (\operatorname{is\text{-}ring}(q) \wedge \forall j \in \mathbb{Z}.\; ((j < \mathit{oq}) \Rightarrow (\operatorname{fq}(j) = \operatorname{zero}(q)))) \Rightarrow  \\ \quad \quad \operatorname{laurent\text{-}mul}(q, \mathit{fq}, \operatorname{laurent\text{-}mono}(q, \mathit{cq}, \mathit{tq}))((\mathit{oq} + \mathit{tq})) = \operatorname{mul}(q)(\operatorname{fq}(\mathit{oq}), \mathit{cq}) \end{array}}}
\vnbprop{laurent-mul-mono-low}{\vnbstmt{\begin{array}{@{}l@{}} \forall q.; \forall \mathit{fq} \in (\mathbb{Z} \to \operatorname{carr}(q)), \mathit{oq} \in \mathbb{Z}, \mathit{cq} \in \operatorname{carr}(q), \mathit{tq}, \mathit{nq} \in \mathbb{Z}. \\ \quad (\operatorname{is\text{-}ring}(q) \wedge (\forall j \in \mathbb{Z}.\; ((j < \mathit{oq}) \Rightarrow (\operatorname{fq}(j) = \operatorname{zero}(q))) \wedge (\mathit{nq} < (\mathit{oq} + \mathit{tq})))) \Rightarrow  \\ \quad \quad \operatorname{laurent\text{-}mul}(q, \mathit{fq}, \operatorname{laurent\text{-}mono}(q, \mathit{cq}, \mathit{tq}))(\mathit{nq}) = \operatorname{zero}(q) \end{array}}}
\vnbprop{laurent-inv-pre-unit}{\vnbstmt{\begin{array}{@{}l@{}} \forall q \in \mathbb{N}, \mathit{kq}.; \forall \mathit{fq} \in \operatorname{laurent\text{-}set}(\mathit{kq}), \mathit{oq} \in \mathbb{Z}, \mathit{aq} \in \operatorname{carr}(\mathit{kq}), \mathit{nq} \in \mathbb{Z}. \\ \quad \operatorname{is\text{-}commutative\text{-}ring}(\mathit{kq}) \wedge \\ \quad \forall j \in \mathbb{Z}.\; ((j < \mathit{oq}) \Rightarrow (\operatorname{fq}(j) = \operatorname{zero}(\mathit{kq}))) \wedge \\ \quad \operatorname{mul}(\mathit{kq})(\operatorname{fq}(\mathit{oq}), \mathit{aq}) = \operatorname{one}(\mathit{kq}) \wedge \\ \quad \mathit{nq} \leq q \Rightarrow \\ \quad \quad \operatorname{laurent\text{-}mul}(\mathit{kq}, \mathit{fq}, \operatorname{laurent\text{-}inv\text{-}pre}(\mathit{kq}, \mathit{fq}, \mathit{oq}, \mathit{aq}, q))(\mathit{nq}) = \operatorname{laurent\text{-}one}(\mathit{kq})(\mathit{nq}) \end{array}}}
\vnbprop{laurent-inv-ser-apply}{\vnbstmt{\begin{array}{@{}l@{}} \forall q.; \forall \mathit{fq} \in \operatorname{laurent\text{-}set}(q), \mathit{oq} \in \mathbb{Z}, \mathit{aq} \in \operatorname{carr}(q), \mathit{jq} \in \mathbb{Z}. \\ \quad \operatorname{is\text{-}ring}(q) \Rightarrow  \\ \quad \quad (\operatorname{laurent\text{-}inv\text{-}ser}(q, \mathit{fq}, \mathit{oq}, \mathit{aq})(\mathit{jq}) \simeq \left(\text{if }(-\mathit{oq} \leq \mathit{jq})\text{ then }\operatorname{laurent\text{-}inv\text{-}pre}(q, \mathit{fq}, \mathit{oq}, \mathit{aq}, (\mathit{jq} + \mathit{oq}))(\mathit{jq})\text{ else }\operatorname{zero}(q)\right)) \end{array}}}
\vnbprop{laurent-inv-ser-fun}{\vnbstmt{\begin{array}{@{}l@{}} \forall q.; \forall \mathit{fq} \in \operatorname{laurent\text{-}set}(q), \mathit{oq} \in \mathbb{Z}, \mathit{aq} \in \operatorname{carr}(q). \\ \quad \operatorname{is\text{-}ring}(q) \Rightarrow (\operatorname{laurent\text{-}inv\text{-}ser}(q, \mathit{fq}, \mathit{oq}, \mathit{aq}) \in (\mathbb{Z} \to \operatorname{carr}(q))) \end{array}}}
\vnbprop{laurent-inv-ser-low}{\vnbstmt{\begin{array}{@{}l@{}} \forall q.; \forall \mathit{fq} \in \operatorname{laurent\text{-}set}(q), \mathit{oq} \in \mathbb{Z}, \mathit{aq} \in \operatorname{carr}(q), j \in \mathbb{Z}. \\ \quad (\operatorname{is\text{-}ring}(q) \wedge (j < -\mathit{oq})) \Rightarrow (\operatorname{laurent\text{-}inv\text{-}ser}(q, \mathit{fq}, \mathit{oq}, \mathit{aq})(j) = \operatorname{zero}(q)) \end{array}}}
\vnbprop{laurent-inv-ser-in}{\vnbstmt{\begin{array}{@{}l@{}} \forall q.; \forall \mathit{fq} \in \operatorname{laurent\text{-}set}(q), \mathit{oq} \in \mathbb{Z}, \mathit{aq} \in \operatorname{carr}(q). \\ \quad \operatorname{is\text{-}ring}(q) \Rightarrow (\operatorname{laurent\text{-}inv\text{-}ser}(q, \mathit{fq}, \mathit{oq}, \mathit{aq}) \in \operatorname{laurent\text{-}set}(q)) \end{array}}}
\vnbprop{laurent-inv-ser-agree}{\vnbstmt{\begin{array}{@{}l@{}} \forall q.; \forall \mathit{fq} \in \operatorname{laurent\text{-}set}(q), \mathit{oq} \in \mathbb{Z}, \mathit{aq} \in \operatorname{carr}(q), b \in \mathbb{N}, j \in \mathbb{Z}. \\ \quad (\operatorname{is\text{-}ring}(q) \wedge (j \leq (b + -\mathit{oq}))) \Rightarrow  \\ \quad \quad \operatorname{laurent\text{-}inv\text{-}ser}(q, \mathit{fq}, \mathit{oq}, \mathit{aq})(j) = \operatorname{laurent\text{-}inv\text{-}pre}(q, \mathit{fq}, \mathit{oq}, \mathit{aq}, b)(j) \end{array}}}
\vnbprop{laurent-mul-unfold}{\vnbstmt{\begin{array}{@{}l@{}} \forall q, \mathit{fq}, \mathit{gq}. \\ \quad (\operatorname{laurent\text{-}mul}(q, \mathit{fq}, \mathit{gq}) \simeq \operatorname{monalg\text{-}mul}(q, \operatorname{zz\text{-}add\text{-}monoid}, \mathit{fq}, \mathit{gq})) \end{array}}}
\vnbprop{laurent-inv-ser-right}{\vnbstmt{\begin{array}{@{}l@{}} \forall q.; \forall \mathit{fq} \in \operatorname{laurent\text{-}set}(q), \mathit{oq} \in \mathbb{Z}, \mathit{aq} \in \operatorname{carr}(q). \\ \quad \operatorname{is\text{-}commutative\text{-}ring}(q) \wedge \\ \quad \forall j \in \mathbb{Z}.\; ((j < \mathit{oq}) \Rightarrow (\operatorname{fq}(j) = \operatorname{zero}(q))) \wedge \\ \quad \operatorname{mul}(q)(\operatorname{fq}(\mathit{oq}), \mathit{aq}) = \operatorname{one}(q) \Rightarrow \\ \quad \quad \operatorname{laurent\text{-}mul}(q, \mathit{fq}, \operatorname{laurent\text{-}inv\text{-}ser}(q, \mathit{fq}, \mathit{oq}, \mathit{aq})) = \operatorname{laurent\text{-}one}(q) \end{array}}}
\vnbprop{laurent-ring-is-field-ring}{\vnbstmt{\forall q.\; (\operatorname{is\text{-}field\text{-}ring}(q) \Rightarrow \operatorname{is\text{-}field\text{-}ring}(\operatorname{laurent\text{-}ring}(q)))}}
\vnbprop{laurent-zero-ne-one}{\vnbstmt{\begin{array}{@{}l@{}} \forall q. \\ \quad (\operatorname{is\text{-}ring}(q) \wedge \neg (\operatorname{one}(q) = \operatorname{zero}(q))) \Rightarrow  \\ \quad \quad \neg (\operatorname{laurent\text{-}zero}(q) = \operatorname{laurent\text{-}one}(q)) \end{array}}}
\vnbprop{laurent-abs-zero-iff}{\vnbstmt{\begin{array}{@{}l@{}} \forall q.; \forall \mathit{rq} \in \mathbb{R}, \mathit{fq} \in \operatorname{laurent\text{-}set}(q). \\ \quad (\operatorname{is\text{-}ring}(q) \wedge (0 < \mathit{rq})) \Rightarrow  \\ \quad \quad (\operatorname{laurent\text{-}abs}(q, \mathit{rq}, \mathit{fq}) = 0) \Leftrightarrow (\mathit{fq} = \operatorname{laurent\text{-}zero}(q)) \end{array}}}
\vnbprop{laurent-nf-carr}{\vnbstmt{\begin{array}{@{}l@{}} \forall q.; \forall \mathit{rq} \in \mathbb{R}. \\ \quad (\operatorname{is\text{-}commutative\text{-}ring}(q) \wedge (0 < \mathit{rq})) \Rightarrow  \\ \quad \quad (\operatorname{carr}(\operatorname{laurent\text{-}nf}(q, \mathit{rq})) \simeq \operatorname{carr}(\operatorname{laurent\text{-}ring}(q))) \end{array}}}
\vnbprop{laurent-nf-add}{\vnbstmt{\begin{array}{@{}l@{}} \forall q.; \forall \mathit{rq} \in \mathbb{R}. \\ \quad (\operatorname{is\text{-}commutative\text{-}ring}(q) \wedge (0 < \mathit{rq})) \Rightarrow  \\ \quad \quad (\operatorname{add}(\operatorname{laurent\text{-}nf}(q, \mathit{rq})) \simeq \operatorname{add}(\operatorname{laurent\text{-}ring}(q))) \end{array}}}
\vnbprop{laurent-nf-mul}{\vnbstmt{\begin{array}{@{}l@{}} \forall q.; \forall \mathit{rq} \in \mathbb{R}. \\ \quad (\operatorname{is\text{-}commutative\text{-}ring}(q) \wedge (0 < \mathit{rq})) \Rightarrow  \\ \quad \quad (\operatorname{mul}(\operatorname{laurent\text{-}nf}(q, \mathit{rq})) \simeq \operatorname{mul}(\operatorname{laurent\text{-}ring}(q))) \end{array}}}
\vnbprop{laurent-nf-neg}{\vnbstmt{\begin{array}{@{}l@{}} \forall q.; \forall \mathit{rq} \in \mathbb{R}. \\ \quad (\operatorname{is\text{-}commutative\text{-}ring}(q) \wedge (0 < \mathit{rq})) \Rightarrow  \\ \quad \quad (\operatorname{neg}(\operatorname{laurent\text{-}nf}(q, \mathit{rq})) \simeq \operatorname{neg}(\operatorname{laurent\text{-}ring}(q))) \end{array}}}
\vnbprop{laurent-nf-zero}{\vnbstmt{\begin{array}{@{}l@{}} \forall q.; \forall \mathit{rq} \in \mathbb{R}. \\ \quad (\operatorname{is\text{-}commutative\text{-}ring}(q) \wedge (0 < \mathit{rq})) \Rightarrow  \\ \quad \quad (\operatorname{zero}(\operatorname{laurent\text{-}nf}(q, \mathit{rq})) \simeq \operatorname{zero}(\operatorname{laurent\text{-}ring}(q))) \end{array}}}
\vnbprop{laurent-nf-one}{\vnbstmt{\begin{array}{@{}l@{}} \forall q.; \forall \mathit{rq} \in \mathbb{R}. \\ \quad (\operatorname{is\text{-}commutative\text{-}ring}(q) \wedge (0 < \mathit{rq})) \Rightarrow  \\ \quad \quad (\operatorname{one}(\operatorname{laurent\text{-}nf}(q, \mathit{rq})) \simeq \operatorname{one}(\operatorname{laurent\text{-}ring}(q))) \end{array}}}
\vnbprop{laurent-nf-fnrm}{\vnbstmt{\begin{array}{@{}l@{}} \forall q.; \forall \mathit{rq} \in \mathbb{R}. \\ \quad (\operatorname{is\text{-}commutative\text{-}ring}(q) \wedge (0 < \mathit{rq})) \Rightarrow  \\ \quad \quad (\operatorname{fnrm}(\operatorname{laurent\text{-}nf}(q, \mathit{rq})) \simeq (\lambda f \in \operatorname{carr}(\operatorname{laurent\text{-}ring}(q)).\; \operatorname{laurent\text{-}abs}(q, \mathit{rq}, f))) \end{array}}}
\vnbprop{laurent-nf-is-normed-field}{\vnbstmt{\begin{array}{@{}l@{}} \forall q.; \forall \mathit{rq} \in \mathbb{R}. \\ \quad (\operatorname{is\text{-}field\text{-}ring}(q) \wedge ((0 < \mathit{rq}) \wedge (\mathit{rq} < 1))) \Rightarrow  \\ \quad \quad \operatorname{is\text{-}normed\text{-}field}(\operatorname{laurent\text{-}nf}(q, \mathit{rq})) \end{array}}}
\vnbprop{laurent-nf-is-nv-field}{\vnbstmt{\begin{array}{@{}l@{}} \forall q.; \forall \mathit{rq} \in \mathbb{R}. \\ \quad (\operatorname{is\text{-}field\text{-}ring}(q) \wedge ((0 < \mathit{rq}) \wedge (\mathit{rq} < 1))) \Rightarrow \operatorname{is\text{-}nv\text{-}field}(\operatorname{laurent\text{-}nf}(q, \mathit{rq})) \end{array}}}

\vnbbookfile{theorem-library/ccr-ch3-ops}{theorem-library/ccr-ch3-ops}
\noindent{\small ccr-ch3-ops.scm -- wave of 2026-11-09/10 (agent CCR-1109-B), wired 2026-11-10 from scratchpad/ccr1109b/ccr-ch3-ops.scm. Library: analysis.}\par
\vnbprop{cb3-mpow-commute-ind}{\vnbstmt{\begin{array}{@{}l@{}} \forall k \in \mathbb{N}, r, x, y. \\ \quad (\operatorname{is\text{-}ring}(r) \wedge ((x \in \operatorname{carr}(r)) \wedge ((y \in \operatorname{carr}(r)) \wedge (\operatorname{mul}(r)(x, y) = \operatorname{mul}(r)(y, x))))) \Rightarrow  \\ \quad \quad \operatorname{mul}(r)(\operatorname{mpow}(\operatorname{ring\text{-}multiplicative\text{-}monoid}(r), x, k), y) \\ \quad \quad =\; \operatorname{mul}(r)(y, \operatorname{mpow}(\operatorname{ring\text{-}multiplicative\text{-}monoid}(r), x, k)) \end{array}}}
\vnbprop{cb3-mpow-commute}{\vnbstmt{\begin{array}{@{}l@{}} \forall r, x, y.; \forall k \in \mathbb{N}. \\ \quad (\operatorname{is\text{-}ring}(r) \wedge ((x \in \operatorname{carr}(r)) \wedge ((y \in \operatorname{carr}(r)) \wedge (\operatorname{mul}(r)(x, y) = \operatorname{mul}(r)(y, x))))) \Rightarrow  \\ \quad \quad \operatorname{mul}(r)(\operatorname{mpow}(\operatorname{ring\text{-}multiplicative\text{-}monoid}(r), x, k), y) \\ \quad \quad =\; \operatorname{mul}(r)(y, \operatorname{mpow}(\operatorname{ring\text{-}multiplicative\text{-}monoid}(r), x, k)) \end{array}}}
\vnbprop{cb3-mpow-mul-commute-ind}{\vnbstmt{\begin{array}{@{}l@{}} \forall k \in \mathbb{N}, r, x, y. \\ \quad (\operatorname{is\text{-}ring}(r) \wedge ((x \in \operatorname{carr}(r)) \wedge ((y \in \operatorname{carr}(r)) \wedge (\operatorname{mul}(r)(x, y) = \operatorname{mul}(r)(y, x))))) \Rightarrow  \\ \quad \quad \operatorname{mpow}(\operatorname{ring\text{-}multiplicative\text{-}monoid}(r), \operatorname{mul}(r)(x, y), k) \\ \quad \quad =\; \operatorname{mul}(r)(\operatorname{mpow}(\operatorname{ring\text{-}multiplicative\text{-}monoid}(r), x, k), \operatorname{mpow}(\operatorname{ring\text{-}multiplicative\text{-}monoid}(r), y, k)) \end{array}}}
\vnbprop{cb3-mpow-mul-commute}{\vnbstmt{\begin{array}{@{}l@{}} \forall r, x, y.; \forall k \in \mathbb{N}. \\ \quad (\operatorname{is\text{-}ring}(r) \wedge ((x \in \operatorname{carr}(r)) \wedge ((y \in \operatorname{carr}(r)) \wedge (\operatorname{mul}(r)(x, y) = \operatorname{mul}(r)(y, x))))) \Rightarrow  \\ \quad \quad \operatorname{mpow}(\operatorname{ring\text{-}multiplicative\text{-}monoid}(r), \operatorname{mul}(r)(x, y), k) \\ \quad \quad =\; \operatorname{mul}(r)(\operatorname{mpow}(\operatorname{ring\text{-}multiplicative\text{-}monoid}(r), x, k), \operatorname{mpow}(\operatorname{ring\text{-}multiplicative\text{-}monoid}(r), y, k)) \end{array}}}
\vnbprop{cb3-sum-powers-in-ind}{\vnbstmt{\begin{array}{@{}l@{}} \forall m \in \mathbb{N}, r, s, x, y.; \forall a, b \in \mathbb{N}. \\ \quad \operatorname{is\text{-}ring}(r) \wedge \\ \quad \forall u, v \in s.\; (\operatorname{add}(r)(u, v) \in s) \wedge \\ \quad x \in \operatorname{carr}(r) \wedge \\ \quad y \in \operatorname{carr}(r) \wedge \\ \quad \operatorname{mul}(r)(x, y) = \operatorname{mul}(r)(y, x) \wedge \\ \quad \forall a, b \in \mathbb{N}. \\ \quad \quad \operatorname{mul}(r)(\operatorname{mpow}(\operatorname{ring\text{-}multiplicative\text{-}monoid}(r), x, a), \operatorname{mpow}(\operatorname{ring\text{-}multiplicative\text{-}monoid}(r), y, b)) \in s \Rightarrow \\ \quad \quad \operatorname{mul}(r)(\operatorname{mul}(r)(\operatorname{mpow}(\operatorname{ring\text{-}multiplicative\text{-}monoid}(r), x, a), \operatorname{mpow}(\operatorname{ring\text{-}multiplicative\text{-}monoid}(r), y, b)), \\ \quad \quad \quad \operatorname{mpow}(\operatorname{ring\text{-}multiplicative\text{-}monoid}(r), \operatorname{add}(r)(x, y), m)) \\ \quad \quad \in\; s \end{array}}}
\vnbprop{cb3-sum-powers-in}{\vnbstmt{\begin{array}{@{}l@{}} \forall r, s, x, y.; \forall m \in \mathbb{N}. \\ \quad \operatorname{is\text{-}ring}(r) \wedge \\ \quad \forall u, v \in s.\; (\operatorname{add}(r)(u, v) \in s) \wedge \\ \quad x \in \operatorname{carr}(r) \wedge \\ \quad y \in \operatorname{carr}(r) \wedge \\ \quad \operatorname{mul}(r)(x, y) = \operatorname{mul}(r)(y, x) \wedge \\ \quad \forall a, b \in \mathbb{N}. \\ \quad \quad \operatorname{mul}(r)(\operatorname{mpow}(\operatorname{ring\text{-}multiplicative\text{-}monoid}(r), x, a), \operatorname{mpow}(\operatorname{ring\text{-}multiplicative\text{-}monoid}(r), y, b)) \in s \Rightarrow \\ \quad \quad \operatorname{mpow}(\operatorname{ring\text{-}multiplicative\text{-}monoid}(r), \operatorname{add}(r)(x, y), m) \in s \end{array}}}
\vnbprop{cb3-mr-ring}{\vnbstmt{\begin{array}{@{}l@{}} \forall k.; \forall n \in \mathbb{N}. \\ \quad \operatorname{is\text{-}normed\text{-}field}(k) \Rightarrow  \\ \quad \quad \operatorname{is\text{-}ring}(\operatorname{mat\text{-}ring}(\operatorname{normed\text{-}field\text{-}as\text{-}commutative\text{-}ring}(k), n)) \end{array}}}
\vnbprop{cb3-mr-carr}{\vnbstmt{\begin{array}{@{}l@{}} \forall k.; \forall n \in \mathbb{N}. \\ \quad \operatorname{is\text{-}normed\text{-}field}(k) \Rightarrow  \\ \quad \quad \operatorname{carr}(\operatorname{mat\text{-}ring}(\operatorname{normed\text{-}field\text{-}as\text{-}commutative\text{-}ring}(k), n)) = \operatorname{mat}(n, n, \operatorname{carr}(k)) \end{array}}}
\vnbprop{cb3-mr-mul-apply}{\vnbstmt{\begin{array}{@{}l@{}} \forall k.; \forall n \in \mathbb{N}, p, q. \\ \quad (\operatorname{is\text{-}normed\text{-}field}(k) \wedge ((p \in \operatorname{mat}(n, n, \operatorname{carr}(k))) \wedge (q \in \operatorname{mat}(n, n, \operatorname{carr}(k))))) \Rightarrow  \\ \quad \quad \operatorname{mul}(\operatorname{mat\text{-}ring}(\operatorname{normed\text{-}field\text{-}as\text{-}commutative\text{-}ring}(k), n))(p, q) \\ \quad \quad =\; \operatorname{matmul}(\operatorname{normed\text{-}field\text{-}as\text{-}commutative\text{-}ring}(k), p, q) \end{array}}}
\vnbprop{cb3-mr-add-apply}{\vnbstmt{\begin{array}{@{}l@{}} \forall k.; \forall n \in \mathbb{N}, p, q. \\ \quad (\operatorname{is\text{-}normed\text{-}field}(k) \wedge ((p \in \operatorname{mat}(n, n, \operatorname{carr}(k))) \wedge (q \in \operatorname{mat}(n, n, \operatorname{carr}(k))))) \Rightarrow  \\ \quad \quad \operatorname{add}(\operatorname{mat\text{-}ring}(\operatorname{normed\text{-}field\text{-}as\text{-}commutative\text{-}ring}(k), n))(p, q) \\ \quad \quad =\; \operatorname{matadd}(\operatorname{normed\text{-}field\text{-}as\text{-}commutative\text{-}ring}(k), p, q) \end{array}}}
\vnbprop{cb3-mat-power-unfold}{\vnbstmt{\begin{array}{@{}l@{}} \forall k, n, m, j. \\ \quad (\operatorname{mat\text{-}power}(k, n, m, j) \simeq \operatorname{mpow}(\operatorname{ring\text{-}multiplicative\text{-}monoid}(\operatorname{mat\text{-}ring}(\operatorname{normed\text{-}field\text{-}as\text{-}commutative\text{-}ring}(k), n)), m, j)) \end{array}}}
\vnbprop{cb3-mat-bound-add}{\vnbstmt{\begin{array}{@{}l@{}} \forall k, m, n, p, q.; \forall c \in \mathbb{R}. \\ \quad (\operatorname{is\text{-}nv\text{-}field}(k) \wedge (\operatorname{is\text{-}mat\text{-}bound}(k, m, n, p, c) \wedge \operatorname{is\text{-}mat\text{-}bound}(k, m, n, q, c))) \Rightarrow  \\ \quad \quad \operatorname{is\text{-}mat\text{-}bound}(k, m, n, \operatorname{matadd}(\operatorname{normed\text{-}field\text{-}as\text{-}commutative\text{-}ring}(k), p, q), c) \end{array}}}
\vnbprop{cb3-mat-bound-mul-entry}{\vnbstmt{\begin{array}{@{}l@{}} \forall k, m, n, l, p, q, c, d.; \forall i \in \operatorname{interval}(1, m), a \in \operatorname{interval}(1, l). \\ \quad \operatorname{is\text{-}nv\text{-}field}(k) \wedge \\ \quad c \in \mathbb{R} \wedge \\ \quad d \in \mathbb{R} \wedge \\ \quad 0 \leq c \wedge \\ \quad 0 \leq d \wedge \\ \quad 1 \leq n \wedge \\ \quad \operatorname{is\text{-}mat\text{-}bound}(k, m, n, p, c) \wedge \\ \quad \operatorname{is\text{-}mat\text{-}bound}(k, n, l, q, d) \Rightarrow \\ \quad \quad \operatorname{fnrm}(k)(\operatorname{entry}(\operatorname{matmul}(\operatorname{normed\text{-}field\text{-}as\text{-}commutative\text{-}ring}(k), p, q), i, a)) \leq (c \cdot d) \end{array}}}
\vnbprop{cb3-mat-bound-mul}{\vnbstmt{\begin{array}{@{}l@{}} \forall k.; \forall n \in \mathbb{N}, p, q, c, d. \\ \quad \operatorname{is\text{-}nv\text{-}field}(k) \wedge \\ \quad c \in \mathbb{R} \wedge \\ \quad d \in \mathbb{R} \wedge \\ \quad 0 \leq c \wedge \\ \quad 0 \leq d \wedge \\ \quad \operatorname{is\text{-}mat\text{-}bound}(k, n, n, p, c) \wedge \\ \quad \operatorname{is\text{-}mat\text{-}bound}(k, n, n, q, d) \Rightarrow \\ \quad \quad \operatorname{is\text{-}mat\text{-}bound}(k, n, n, \operatorname{matmul}(\operatorname{normed\text{-}field\text{-}as\text{-}commutative\text{-}ring}(k), p, q), (c \cdot d)) \end{array}}}
\vnbprop{ccr-3-1-eq19}{\vnbstmt{\begin{array}{@{}l@{}} \forall k.; \forall n \in \mathbb{N}, p, q, c, d. \\ \quad \operatorname{is\text{-}nv\text{-}field}(k) \wedge \\ \quad c \in \mathbb{R} \wedge \\ \quad d \in \mathbb{R} \wedge \\ \quad 0 \leq c \wedge \\ \quad 0 \leq d \wedge \\ \quad \operatorname{is\text{-}mat\text{-}bound}(k, n, n, p, c) \wedge \\ \quad \operatorname{is\text{-}mat\text{-}bound}(k, n, 1, q, d) \Rightarrow \\ \quad \quad \operatorname{is\text{-}mat\text{-}bound}(k, n, 1, \operatorname{matmul}(\operatorname{normed\text{-}field\text{-}as\text{-}commutative\text{-}ring}(k), p, q), (c \cdot d)) \end{array}}}
\vnbprop{cb3-identmat-bound}{\vnbstmt{\begin{array}{@{}l@{}} \forall k.; \forall n \in \mathbb{N}. \\ \quad \operatorname{is\text{-}normed\text{-}field}(k) \Rightarrow  \\ \quad \quad \operatorname{is\text{-}mat\text{-}bound}(k, n, n, \operatorname{identmat}(\operatorname{normed\text{-}field\text{-}as\text{-}commutative\text{-}ring}(k), n), 1) \end{array}}}
\vnbprop{ccr-3-3-product}{\vnbstmt{\begin{array}{@{}l@{}} \forall k.; \forall n \in \mathbb{N}, a, b. \\ \quad \operatorname{is\text{-}nv\text{-}field}(k) \wedge \\ \quad \operatorname{is\text{-}power\text{-}bounded}(k, n, a) \wedge \\ \quad \operatorname{is\text{-}power\text{-}bounded}(k, n, b) \wedge \\ \quad \operatorname{matmul}(\operatorname{normed\text{-}field\text{-}as\text{-}commutative\text{-}ring}(k), a, b) \\ \quad =\; \operatorname{matmul}(\operatorname{normed\text{-}field\text{-}as\text{-}commutative\text{-}ring}(k), b, a) \Rightarrow \\ \quad \quad \operatorname{is\text{-}power\text{-}bounded}(k, n, \operatorname{matmul}(\operatorname{normed\text{-}field\text{-}as\text{-}commutative\text{-}ring}(k), a, b)) \end{array}}}
\vnbprop{ccr-3-3-sum}{\vnbstmt{\begin{array}{@{}l@{}} \forall k.; \forall n \in \mathbb{N}, a, b. \\ \quad \operatorname{is\text{-}nv\text{-}field}(k) \wedge \\ \quad \operatorname{is\text{-}power\text{-}bounded}(k, n, a) \wedge \\ \quad \operatorname{is\text{-}power\text{-}bounded}(k, n, b) \wedge \\ \quad \operatorname{matmul}(\operatorname{normed\text{-}field\text{-}as\text{-}commutative\text{-}ring}(k), a, b) \\ \quad =\; \operatorname{matmul}(\operatorname{normed\text{-}field\text{-}as\text{-}commutative\text{-}ring}(k), b, a) \Rightarrow \\ \quad \quad \operatorname{is\text{-}power\text{-}bounded}(k, n, \operatorname{matadd}(\operatorname{normed\text{-}field\text{-}as\text{-}commutative\text{-}ring}(k), a, b)) \end{array}}}
\vnbprop{cb3-mpow-one-ind}{\vnbstmt{\begin{array}{@{}l@{}} \forall k \in \mathbb{N}, r. \\ \quad \operatorname{is\text{-}ring}(r) \Rightarrow (\operatorname{mpow}(\operatorname{ring\text{-}multiplicative\text{-}monoid}(r), \operatorname{one}(r), k) = \operatorname{one}(r)) \end{array}}}
\vnbprop{cb3-mpow-one}{\vnbstmt{\begin{array}{@{}l@{}} \forall r.; \forall k \in \mathbb{N}. \\ \quad \operatorname{is\text{-}ring}(r) \Rightarrow (\operatorname{mpow}(\operatorname{ring\text{-}multiplicative\text{-}monoid}(r), \operatorname{one}(r), k) = \operatorname{one}(r)) \end{array}}}
\vnbprop{cb3-identmat-power-bounded}{\vnbstmt{\begin{array}{@{}l@{}} \forall k.; \forall n \in \mathbb{N}. \\ \quad \operatorname{is\text{-}normed\text{-}field}(k) \Rightarrow  \\ \quad \quad \operatorname{is\text{-}power\text{-}bounded\text{-}by}(k, n, \operatorname{identmat}(\operatorname{normed\text{-}field\text{-}as\text{-}commutative\text{-}ring}(k), n), 1) \end{array}}}
\vnbprop{cb3-det-mpow-ind}{\vnbstmt{\begin{array}{@{}l@{}} \forall b \in \mathbb{N}, k.; \forall n \in \mathbb{N}, a \in \operatorname{mat}(n, n, \operatorname{carr}(k)). \\ \quad \operatorname{is\text{-}normed\text{-}field}(k) \Rightarrow  \\ \quad \quad \operatorname{det}(\operatorname{normed\text{-}field\text{-}as\text{-}commutative\text{-}ring}(k), n, \\ \quad \quad \quad \operatorname{mpow}(\operatorname{ring\text{-}multiplicative\text{-}monoid}(\operatorname{mat\text{-}ring}(\operatorname{normed\text{-}field\text{-}as\text{-}commutative\text{-}ring}(k), n)), a, b)) \\ \quad \quad =\; \operatorname{ring\text{-}power}(\operatorname{normed\text{-}field\text{-}as\text{-}commutative\text{-}ring}(k), \operatorname{det}(\operatorname{normed\text{-}field\text{-}as\text{-}commutative\text{-}ring}(k), n, a), b) \end{array}}}
\vnbprop{cb3-sign-norm-ind}{\vnbstmt{\begin{array}{@{}l@{}} \forall a \in \mathbb{N}, k. \\ \quad \operatorname{is\text{-}normed\text{-}field}(k) \Rightarrow  \\ \quad \quad \operatorname{fnrm}(k)(\operatorname{mpow}(\operatorname{ring\text{-}multiplicative\text{-}monoid}(\operatorname{normed\text{-}field\text{-}as\text{-}commutative\text{-}ring}(k)), \\ \quad \quad \quad \quad \operatorname{neg}(\operatorname{normed\text{-}field\text{-}as\text{-}commutative\text{-}ring}(k))(\operatorname{one}(\operatorname{normed\text{-}field\text{-}as\text{-}commutative\text{-}ring}(k))), \\ \quad \quad \quad \quad a)) \\ \quad \quad =\; 1 \end{array}}}
\vnbprop{cb3-sign-norm}{\vnbstmt{\begin{array}{@{}l@{}} \forall k.; \forall a \in \mathbb{N}. \\ \quad \operatorname{is\text{-}normed\text{-}field}(k) \Rightarrow  \\ \quad \quad \operatorname{fnrm}(k)(\operatorname{mpow}(\operatorname{ring\text{-}multiplicative\text{-}monoid}(\operatorname{normed\text{-}field\text{-}as\text{-}commutative\text{-}ring}(k)), \\ \quad \quad \quad \quad \operatorname{neg}(\operatorname{normed\text{-}field\text{-}as\text{-}commutative\text{-}ring}(k))(\operatorname{one}(\operatorname{normed\text{-}field\text{-}as\text{-}commutative\text{-}ring}(k))), \\ \quad \quad \quad \quad a)) \\ \quad \quad =\; 1 \end{array}}}
\vnbprop{cb3-minor-bound}{\vnbstmt{\begin{array}{@{}l@{}} \forall k.; \forall n \in \mathbb{N}, a, c.; \forall q \in \operatorname{interval}(1, (n + 1)). \\ \quad (\operatorname{is\text{-}normed\text{-}field}(k) \wedge \operatorname{is\text{-}mat\text{-}bound}(k, (n + 1), (n + 1), a, c)) \Rightarrow  \\ \quad \quad \operatorname{is\text{-}mat\text{-}bound}(k, n, n, \operatorname{minor}(a, 1, q, n), c) \end{array}}}
\vnbprop{cb3-det-bound-ind}{\vnbstmt{\begin{array}{@{}l@{}} \forall n \in \mathbb{N}, k.; \forall a \in \operatorname{mat}(n, n, \operatorname{carr}(k)), c \in \mathbb{R}. \\ \quad (\operatorname{is\text{-}nv\text{-}field}(k) \wedge ((0 \leq c) \wedge \operatorname{is\text{-}mat\text{-}bound}(k, n, n, a, c))) \Rightarrow  \\ \quad \quad \operatorname{fnrm}(k)(\operatorname{det}(\operatorname{normed\text{-}field\text{-}as\text{-}commutative\text{-}ring}(k), n, a)) \leq \operatorname{power}(c, n) \end{array}}}
\vnbprop{ccr-3-1-eq21}{\vnbstmt{\begin{array}{@{}l@{}} \forall k.; \forall n \in \mathbb{N}, a.; \forall c \in \mathbb{R}. \\ \quad (\operatorname{is\text{-}nv\text{-}field}(k) \wedge ((0 \leq c) \wedge \operatorname{is\text{-}mat\text{-}bound}(k, n, n, a, c))) \Rightarrow  \\ \quad \quad \operatorname{fnrm}(k)(\operatorname{det}(\operatorname{normed\text{-}field\text{-}as\text{-}commutative\text{-}ring}(k), n, a)) \leq \operatorname{power}(c, n) \end{array}}}
\vnbprop{cb3-rr-power-bounded-le-one}{\vnbstmt{\begin{array}{@{}l@{}} \forall x, b \in \mathbb{R}. \\ \quad ((0 \leq x) \wedge \forall n \in \mathbb{N}.\; (\operatorname{power}(x, n) \leq b)) \Rightarrow (x \leq 1) \end{array}}}
\vnbprop{cb3-det-shift-bound}{\vnbstmt{\begin{array}{@{}l@{}} \forall k.; \forall n \in \mathbb{N}, a. \\ \quad (\operatorname{is\text{-}nv\text{-}field}(k) \wedge \operatorname{is\text{-}power\text{-}bounded}(k, n, a)) \Rightarrow  \\ \quad \quad \operatorname{fnrm}(k)(\operatorname{det}(\operatorname{normed\text{-}field\text{-}as\text{-}commutative\text{-}ring}(k), n, \\ \quad \quad \quad \quad \operatorname{matadd}(\operatorname{normed\text{-}field\text{-}as\text{-}commutative\text{-}ring}(k), \operatorname{identmat}(\operatorname{normed\text{-}field\text{-}as\text{-}commutative\text{-}ring}(k), n), a))) \\ \quad \quad \leq\; 1 \end{array}}}
\vnbprop{ccr-3-5}{\vnbstmt{\begin{array}{@{}l@{}} \forall k.; \forall n \in \mathbb{N}, a. \\ \quad \operatorname{is\text{-}nv\text{-}field}(k) \wedge \\ \quad \operatorname{is\text{-}power\text{-}bounded}(k, n, a) \wedge \\ \quad \operatorname{fnrm}(k)(\operatorname{det}(\operatorname{normed\text{-}field\text{-}as\text{-}commutative\text{-}ring}(k), n, a)) \leq 1 \Rightarrow \\ \quad \quad \operatorname{fnrm}(k)(\operatorname{det}(\operatorname{normed\text{-}field\text{-}as\text{-}commutative\text{-}ring}(k), n, \\ \quad \quad \quad \quad \operatorname{matadd}(\operatorname{normed\text{-}field\text{-}as\text{-}commutative\text{-}ring}(k), \operatorname{identmat}(\operatorname{normed\text{-}field\text{-}as\text{-}commutative\text{-}ring}(k), n), a))) \\ \quad \quad \leq\; 1 \end{array}}}
\vnbprop{ccr-3-6-multiplicative}{\vnbstmt{\begin{array}{@{}l@{}} \forall k.; \forall n \in \mathbb{N}, a, b. \\ \quad (\operatorname{is\text{-}normed\text{-}field}(k) \wedge ((a \in \operatorname{mat}(n, n, \operatorname{carr}(k))) \wedge (b \in \operatorname{mat}(n, n, \operatorname{carr}(k))))) \Rightarrow  \\ \quad \quad \operatorname{fnrm}(k)(\operatorname{det}(\operatorname{normed\text{-}field\text{-}as\text{-}commutative\text{-}ring}(k), n, \operatorname{matmul}(\operatorname{normed\text{-}field\text{-}as\text{-}commutative\text{-}ring}(k), a, b))) \\ \quad \quad =\; \operatorname{fnrm}(k)(\operatorname{det}(\operatorname{normed\text{-}field\text{-}as\text{-}commutative\text{-}ring}(k), n, a)) \cdot \operatorname{fnrm}(k)(\operatorname{det}(\operatorname{normed\text{-}field\text{-}as\text{-}commutative\text{-}ring}(k), n, b)) \end{array}}}

\vnbbookfile{theorem-library/ccr-ch3b-alg}{theorem-library/ccr-ch3b-alg}
\noindent{\small ccr-ch3b-alg.scm -- wave of 2026-11-09/10 (agent CCR-1109-B2), wired 2026-11-10 from scratchpad/ccr1109b2/ccr-ch3b-alg.scm. Library: analysis.}\par
\vnbprop{cq4-view-field-ring}{\vnbstmt{\begin{array}{@{}l@{}} \forall k. \\ \quad \operatorname{is\text{-}normed\text{-}field}(k) \Rightarrow \operatorname{is\text{-}field\text{-}ring}(\operatorname{normed\text{-}field\text{-}as\text{-}commutative\text{-}ring}(k)) \end{array}}}
\vnbprop{cq4-mat-power-type}{\vnbstmt{\begin{array}{@{}l@{}} \forall k.; \forall n \in \mathbb{N}, m \in \operatorname{mat}(n, n, \operatorname{carr}(k)), j \in \mathbb{N}. \\ \quad \operatorname{is\text{-}normed\text{-}field}(k) \Rightarrow (\operatorname{mat\text{-}power}(k, n, m, j) \in \operatorname{mat}(n, n, \operatorname{carr}(k))) \end{array}}}
\vnbprop{cq4-matscale-mul-left}{\vnbstmt{\begin{array}{@{}l@{}} \forall k.; \forall n \in \mathbb{N}, c \in \operatorname{carr}(k), p, q. \\ \quad (\operatorname{is\text{-}normed\text{-}field}(k) \wedge ((p \in \operatorname{mat}(n, n, \operatorname{carr}(k))) \wedge (q \in \operatorname{mat}(n, n, \operatorname{carr}(k))))) \Rightarrow  \\ \quad \quad \operatorname{matmul}(\operatorname{normed\text{-}field\text{-}as\text{-}commutative\text{-}ring}(k), \operatorname{matscale}(\operatorname{normed\text{-}field\text{-}as\text{-}commutative\text{-}ring}(k), c, p), q) \\ \quad \quad =\; \operatorname{matscale}(\operatorname{normed\text{-}field\text{-}as\text{-}commutative\text{-}ring}(k), c, \operatorname{matmul}(\operatorname{normed\text{-}field\text{-}as\text{-}commutative\text{-}ring}(k), p, q)) \end{array}}}
\vnbprop{cq4-matscale-mul-right}{\vnbstmt{\begin{array}{@{}l@{}} \forall k.; \forall n \in \mathbb{N}, c \in \operatorname{carr}(k), p, q. \\ \quad (\operatorname{is\text{-}normed\text{-}field}(k) \wedge ((p \in \operatorname{mat}(n, n, \operatorname{carr}(k))) \wedge (q \in \operatorname{mat}(n, n, \operatorname{carr}(k))))) \Rightarrow  \\ \quad \quad \operatorname{matmul}(\operatorname{normed\text{-}field\text{-}as\text{-}commutative\text{-}ring}(k), p, \operatorname{matscale}(\operatorname{normed\text{-}field\text{-}as\text{-}commutative\text{-}ring}(k), c, q)) \\ \quad \quad =\; \operatorname{matscale}(\operatorname{normed\text{-}field\text{-}as\text{-}commutative\text{-}ring}(k), c, \operatorname{matmul}(\operatorname{normed\text{-}field\text{-}as\text{-}commutative\text{-}ring}(k), p, q)) \end{array}}}
\vnbprop{cq4-zero-bound}{\vnbstmt{\begin{array}{@{}l@{}} \forall k.; \forall m, n \in \mathbb{N}, c \in \mathbb{R}. \\ \quad (\operatorname{is\text{-}normed\text{-}field}(k) \wedge (0 \leq c)) \Rightarrow  \\ \quad \quad \operatorname{is\text{-}mat\text{-}bound}(k, m, n, \operatorname{zeromat}(\operatorname{normed\text{-}field\text{-}as\text{-}commutative\text{-}ring}(k), m, n), c) \end{array}}}
\vnbprop{cq4-scale-bound}{\vnbstmt{\begin{array}{@{}l@{}} \forall k, m, n.; \forall c \in \operatorname{carr}(k), b \in \mathbb{R}, a. \\ \quad (\operatorname{is\text{-}normed\text{-}field}(k) \wedge ((\operatorname{fnrm}(k)(c) \leq 1) \wedge ((0 \leq b) \wedge \operatorname{is\text{-}mat\text{-}bound}(k, m, n, a, b)))) \Rightarrow  \\ \quad \quad \operatorname{is\text{-}mat\text{-}bound}(k, m, n, \operatorname{matscale}(\operatorname{normed\text{-}field\text{-}as\text{-}commutative\text{-}ring}(k), c, a), b) \end{array}}}
\vnbprop{cq4-nvint-parts}{\vnbstmt{\forall k.; \forall c \in \operatorname{nv\text{-}integers}(k).\; ((c \in \operatorname{carr}(k)) \wedge (\operatorname{fnrm}(k)(c) \leq 1))}}
\vnbprop{cq4-nvint-intro}{\vnbstmt{\forall k.; \forall c \in \operatorname{carr}(k).\; ((\operatorname{fnrm}(k)(c) \leq 1) \Rightarrow (c \in \operatorname{nv\text{-}integers}(k)))}}
\vnbprop{cq4-zalg-unfold}{\vnbstmt{\begin{array}{@{}l@{}} \forall k, n, m. \\ \quad (\operatorname{zpoly\text{-}alg}(k, n, m) \simeq \operatorname{intersection\text{-}of}(\{s \in \operatorname{power}(\operatorname{mat}(n, n, \operatorname{carr}(k))) : ((\operatorname{identmat}(\operatorname{normed\text{-}field\text{-}as\text{-}commutative\text{-}ring}(k), n) \in s) \wedge ((\operatorname{zeromat}(\operatorname{normed\text{-}field\text{-}as\text{-}commutative\text{-}ring}(k), n, n) \in s) \wedge (\forall x \in s.\; (\operatorname{matmul}(\operatorname{normed\text{-}field\text{-}as\text{-}commutative\text{-}ring}(k), x, m) \in s) \wedge (\forall x, y \in s.\; (\operatorname{matadd}(\operatorname{normed\text{-}field\text{-}as\text{-}commutative\text{-}ring}(k), x, y) \in s) \wedge \forall x \in s, c \in \operatorname{nv\text{-}integers}(k).\; (\operatorname{matscale}(\operatorname{normed\text{-}field\text{-}as\text{-}commutative\text{-}ring}(k), c, x) \in s)))))\})) \end{array}}}
\vnbprop{cq4-zalg-fam-has-mat}{\vnbstmt{\begin{array}{@{}l@{}} \forall k.; \forall n \in \mathbb{N}, m \in \operatorname{mat}(n, n, \operatorname{carr}(k)). \\ \quad \operatorname{is\text{-}normed\text{-}field}(k) \Rightarrow  \\ \quad \quad \operatorname{mat}(n, n, \operatorname{carr}(k)) \\ \quad \quad \in\; \{s \in \operatorname{power}(\operatorname{mat}(n, n, \operatorname{carr}(k))) : ((\operatorname{identmat}(\operatorname{normed\text{-}field\text{-}as\text{-}commutative\text{-}ring}(k), n) \in s) \wedge ((\operatorname{zeromat}(\operatorname{normed\text{-}field\text{-}as\text{-}commutative\text{-}ring}(k), n, n) \in s) \wedge (\forall x \in s.\; (\operatorname{matmul}(\operatorname{normed\text{-}field\text{-}as\text{-}commutative\text{-}ring}(k), x, m) \in s) \wedge (\forall x, y \in s.\; (\operatorname{matadd}(\operatorname{normed\text{-}field\text{-}as\text{-}commutative\text{-}ring}(k), x, y) \in s) \wedge \forall x \in s, c \in \operatorname{nv\text{-}integers}(k).\; (\operatorname{matscale}(\operatorname{normed\text{-}field\text{-}as\text{-}commutative\text{-}ring}(k), c, x) \in s)))))\} \end{array}}}
\vnbprop{cq4-zalg-mat}{\vnbstmt{\begin{array}{@{}l@{}} \forall k.; \forall n \in \mathbb{N}, m \in \operatorname{mat}(n, n, \operatorname{carr}(k)), p \in \operatorname{zpoly\text{-}alg}(k, n, m). \\ \quad \operatorname{is\text{-}normed\text{-}field}(k) \Rightarrow (p \in \operatorname{mat}(n, n, \operatorname{carr}(k))) \end{array}}}
\vnbprop{cq4-zalg-id}{\vnbstmt{\begin{array}{@{}l@{}} \forall k.; \forall n \in \mathbb{N}, m \in \operatorname{mat}(n, n, \operatorname{carr}(k)). \\ \quad \operatorname{is\text{-}normed\text{-}field}(k) \Rightarrow  \\ \quad \quad \operatorname{identmat}(\operatorname{normed\text{-}field\text{-}as\text{-}commutative\text{-}ring}(k), n) \in \operatorname{zpoly\text{-}alg}(k, n, m) \end{array}}}
\vnbprop{cq4-zalg-zero}{\vnbstmt{\begin{array}{@{}l@{}} \forall k.; \forall n \in \mathbb{N}, m \in \operatorname{mat}(n, n, \operatorname{carr}(k)). \\ \quad \operatorname{is\text{-}normed\text{-}field}(k) \Rightarrow  \\ \quad \quad \operatorname{zeromat}(\operatorname{normed\text{-}field\text{-}as\text{-}commutative\text{-}ring}(k), n, n) \in \operatorname{zpoly\text{-}alg}(k, n, m) \end{array}}}
\vnbprop{cq4-zalg-mult}{\vnbstmt{\begin{array}{@{}l@{}} \forall k.; \forall n \in \mathbb{N}, m \in \operatorname{mat}(n, n, \operatorname{carr}(k)), p \in \operatorname{zpoly\text{-}alg}(k, n, m). \\ \quad \operatorname{is\text{-}normed\text{-}field}(k) \Rightarrow  \\ \quad \quad \operatorname{matmul}(\operatorname{normed\text{-}field\text{-}as\text{-}commutative\text{-}ring}(k), p, m) \in \operatorname{zpoly\text{-}alg}(k, n, m) \end{array}}}
\vnbprop{cq4-zalg-add}{\vnbstmt{\begin{array}{@{}l@{}} \forall k.; \forall n \in \mathbb{N}, m \in \operatorname{mat}(n, n, \operatorname{carr}(k)), p, q. \\ \quad (\operatorname{is\text{-}normed\text{-}field}(k) \wedge ((p \in \operatorname{zpoly\text{-}alg}(k, n, m)) \wedge (q \in \operatorname{zpoly\text{-}alg}(k, n, m)))) \Rightarrow  \\ \quad \quad \operatorname{matadd}(\operatorname{normed\text{-}field\text{-}as\text{-}commutative\text{-}ring}(k), p, q) \in \operatorname{zpoly\text{-}alg}(k, n, m) \end{array}}}
\vnbprop{cq4-zalg-scale}{\vnbstmt{\begin{array}{@{}l@{}} \forall k.; \forall n \in \mathbb{N}, m \in \operatorname{mat}(n, n, \operatorname{carr}(k)), p, c. \\ \quad (\operatorname{is\text{-}normed\text{-}field}(k) \wedge ((p \in \operatorname{zpoly\text{-}alg}(k, n, m)) \wedge (c \in \operatorname{nv\text{-}integers}(k)))) \Rightarrow  \\ \quad \quad \operatorname{matscale}(\operatorname{normed\text{-}field\text{-}as\text{-}commutative\text{-}ring}(k), c, p) \in \operatorname{zpoly\text{-}alg}(k, n, m) \end{array}}}
\vnbprop{cq4-zalg-induction}{\vnbstmt{\begin{array}{@{}l@{}} \forall k.; \forall n \in \mathbb{N}, m \in \operatorname{mat}(n, n, \operatorname{carr}(k)), s.; \forall p \in \operatorname{zpoly\text{-}alg}(k, n, m). \\ \quad \operatorname{is\text{-}normed\text{-}field}(k) \wedge \\ \quad \operatorname{identmat}(\operatorname{normed\text{-}field\text{-}as\text{-}commutative\text{-}ring}(k), n) \in s \wedge \\ \quad \operatorname{zeromat}(\operatorname{normed\text{-}field\text{-}as\text{-}commutative\text{-}ring}(k), n, n) \in s \wedge \\ \quad \forall x \in s.\; (\operatorname{matmul}(\operatorname{normed\text{-}field\text{-}as\text{-}commutative\text{-}ring}(k), x, m) \in s) \wedge \\ \quad \forall x, y \in s.\; (\operatorname{matadd}(\operatorname{normed\text{-}field\text{-}as\text{-}commutative\text{-}ring}(k), x, y) \in s) \wedge \\ \quad \forall x \in s, c \in \operatorname{nv\text{-}integers}(k). \\ \quad \quad \operatorname{matscale}(\operatorname{normed\text{-}field\text{-}as\text{-}commutative\text{-}ring}(k), c, x) \in s \Rightarrow \\ \quad \quad p \in s \end{array}}}
\vnbprop{cq4-pow-zero}{\vnbstmt{\begin{array}{@{}l@{}} \forall k.; \forall n \in \mathbb{N}, m \in \operatorname{mat}(n, n, \operatorname{carr}(k)). \\ \quad \operatorname{is\text{-}normed\text{-}field}(k) \Rightarrow  \\ \quad \quad \operatorname{mat\text{-}power}(k, n, m, 0) = \operatorname{identmat}(\operatorname{normed\text{-}field\text{-}as\text{-}commutative\text{-}ring}(k), n) \end{array}}}
\vnbprop{cq4-pow-succ}{\vnbstmt{\begin{array}{@{}l@{}} \forall k.; \forall n \in \mathbb{N}, m \in \operatorname{mat}(n, n, \operatorname{carr}(k)), j \in \mathbb{N}. \\ \quad \operatorname{is\text{-}normed\text{-}field}(k) \Rightarrow  \\ \quad \quad \operatorname{mat\text{-}power}(k, n, m, (j + 1)) = \operatorname{matmul}(\operatorname{normed\text{-}field\text{-}as\text{-}commutative\text{-}ring}(k), m, \operatorname{mat\text{-}power}(k, n, m, j)) \end{array}}}
\vnbprop{cq4-zero-mul-left}{\vnbstmt{\begin{array}{@{}l@{}} \forall k.; \forall n \in \mathbb{N}, p \in \operatorname{mat}(n, n, \operatorname{carr}(k)). \\ \quad \operatorname{is\text{-}normed\text{-}field}(k) \Rightarrow  \\ \quad \quad \operatorname{matmul}(\operatorname{normed\text{-}field\text{-}as\text{-}commutative\text{-}ring}(k), \operatorname{zeromat}(\operatorname{normed\text{-}field\text{-}as\text{-}commutative\text{-}ring}(k), n, n), p) \\ \quad \quad =\; \operatorname{zeromat}(\operatorname{normed\text{-}field\text{-}as\text{-}commutative\text{-}ring}(k), n, n) \end{array}}}
\vnbprop{cq4-zalg-bounded}{\vnbstmt{\begin{array}{@{}l@{}} \forall k.; \forall n \in \mathbb{N}, m, c.; \forall p \in \operatorname{zpoly\text{-}alg}(k, n, m), j \in \mathbb{N}. \\ \quad (\operatorname{is\text{-}nv\text{-}field}(k) \wedge \operatorname{is\text{-}power\text{-}bounded\text{-}by}(k, n, m, c)) \Rightarrow  \\ \quad \quad \operatorname{is\text{-}mat\text{-}bound}(k, n, n, \operatorname{matmul}(\operatorname{normed\text{-}field\text{-}as\text{-}commutative\text{-}ring}(k), p, \operatorname{mat\text{-}power}(k, n, m, j)), c) \end{array}}}
\vnbprop{cq4-zero-mul-right}{\vnbstmt{\begin{array}{@{}l@{}} \forall k.; \forall n \in \mathbb{N}, p \in \operatorname{mat}(n, n, \operatorname{carr}(k)). \\ \quad \operatorname{is\text{-}normed\text{-}field}(k) \Rightarrow  \\ \quad \quad \operatorname{matmul}(\operatorname{normed\text{-}field\text{-}as\text{-}commutative\text{-}ring}(k), p, \operatorname{zeromat}(\operatorname{normed\text{-}field\text{-}as\text{-}commutative\text{-}ring}(k), n, n)) \\ \quad \quad =\; \operatorname{zeromat}(\operatorname{normed\text{-}field\text{-}as\text{-}commutative\text{-}ring}(k), n, n) \end{array}}}
\vnbprop{cq4-zalg-mul}{\vnbstmt{\begin{array}{@{}l@{}} \forall k.; \forall n \in \mathbb{N}, m \in \operatorname{mat}(n, n, \operatorname{carr}(k)), p, q. \\ \quad (\operatorname{is\text{-}normed\text{-}field}(k) \wedge ((p \in \operatorname{zpoly\text{-}alg}(k, n, m)) \wedge (q \in \operatorname{zpoly\text{-}alg}(k, n, m)))) \Rightarrow  \\ \quad \quad \operatorname{matmul}(\operatorname{normed\text{-}field\text{-}as\text{-}commutative\text{-}ring}(k), p, q) \in \operatorname{zpoly\text{-}alg}(k, n, m) \end{array}}}
\vnbprop{cq4-zalg-pow}{\vnbstmt{\begin{array}{@{}l@{}} \forall \mathit{zkk} \in \mathbb{N}, k.; \forall n \in \mathbb{N}, m \in \operatorname{mat}(n, n, \operatorname{carr}(k)), p \in \operatorname{zpoly\text{-}alg}(k, n, m). \\ \quad \operatorname{is\text{-}normed\text{-}field}(k) \Rightarrow (\operatorname{mat\text{-}power}(k, n, p, \mathit{zkk}) \in \operatorname{zpoly\text{-}alg}(k, n, m)) \end{array}}}
\vnbprop{cq4-zalg-power-bounded}{\vnbstmt{\begin{array}{@{}l@{}} \forall k.; \forall n \in \mathbb{N}, m, c.; \forall p \in \operatorname{zpoly\text{-}alg}(k, n, m). \\ \quad (\operatorname{is\text{-}nv\text{-}field}(k) \wedge \operatorname{is\text{-}power\text{-}bounded\text{-}by}(k, n, m, c)) \Rightarrow  \\ \quad \quad \operatorname{is\text{-}power\text{-}bounded\text{-}by}(k, n, p, c) \end{array}}}
\vnbprop{cq4-zalg-gen}{\vnbstmt{\begin{array}{@{}l@{}} \forall k.; \forall n \in \mathbb{N}, m \in \operatorname{mat}(n, n, \operatorname{carr}(k)). \\ \quad \operatorname{is\text{-}normed\text{-}field}(k) \Rightarrow (m \in \operatorname{zpoly\text{-}alg}(k, n, m)) \end{array}}}
\vnbprop{cq4-zalg-poly}{\vnbstmt{\begin{array}{@{}l@{}} \forall k.; \forall n \in \mathbb{N}, m \in \operatorname{mat}(n, n, \operatorname{carr}(k)), p \in \operatorname{carr}(\operatorname{poly}(\operatorname{normed\text{-}field\text{-}as\text{-}commutative\text{-}ring}(k))). \\ \quad (\operatorname{is\text{-}normed\text{-}field}(k) \wedge \forall i \in \mathbb{N}.\; (p(i) \in \operatorname{nv\text{-}integers}(k))) \Rightarrow  \\ \quad \quad \operatorname{mat\text{-}poly\text{-}eval}(k, n, p, m) \in \operatorname{zpoly\text{-}alg}(k, n, m) \end{array}}}
\vnbprop{ccr-3-3-poly}{\vnbstmt{\begin{array}{@{}l@{}} \forall k.; \forall n \in \mathbb{N}, m, p. \\ \quad \operatorname{is\text{-}nv\text{-}field}(k) \wedge \\ \quad \operatorname{is\text{-}power\text{-}bounded}(k, n, m) \wedge \\ \quad p \in \operatorname{carr}(\operatorname{poly}(\operatorname{normed\text{-}field\text{-}as\text{-}commutative\text{-}ring}(k))) \wedge \\ \quad \forall i \in \mathbb{N}.\; (p(i) \in \operatorname{nv\text{-}integers}(k)) \Rightarrow \\ \quad \quad \operatorname{is\text{-}power\text{-}bounded}(k, n, \operatorname{mat\text{-}poly\text{-}eval}(k, n, p, m)) \end{array}}}
\vnbprop{ccr-3-3-zalg}{\vnbstmt{\begin{array}{@{}l@{}} \forall k.; \forall n \in \mathbb{N}, m.; \forall p \in \operatorname{zpoly\text{-}alg}(k, n, m). \\ \quad (\operatorname{is\text{-}nv\text{-}field}(k) \wedge \operatorname{is\text{-}power\text{-}bounded}(k, n, m)) \Rightarrow \operatorname{is\text{-}power\text{-}bounded}(k, n, p) \end{array}}}
\vnbprop{cq4-col-type}{\vnbstmt{\begin{array}{@{}l@{}} \forall q, n, l, x, j. \\ \quad ((q \in \operatorname{mat}(n, l, x)) \wedge (j \in \operatorname{interval}(1, l))) \Rightarrow  \\ \quad \quad \operatorname{col\text{-}of}(q, n, j) \in \operatorname{mat}(n, 1, x) \end{array}}}
\vnbprop{cq4-col-entry}{\vnbstmt{\begin{array}{@{}l@{}} \forall q, n, l, x, j, i. \\ \quad ((q \in \operatorname{mat}(n, l, x)) \wedge ((j \in \operatorname{interval}(1, l)) \wedge (i \in \operatorname{interval}(1, n)))) \Rightarrow  \\ \quad \quad \operatorname{entry}(\operatorname{col\text{-}of}(q, n, j), i, 1) = \operatorname{entry}(q, i, j) \end{array}}}
\vnbprop{cq4-col-mul}{\vnbstmt{\begin{array}{@{}l@{}} \forall a, m, n, l, p, q.; \forall i \in \operatorname{interval}(1, m), j \in \operatorname{interval}(1, l). \\ \quad (\operatorname{is\text{-}ring}(a) \wedge ((p \in \operatorname{mat}(m, n, \operatorname{carr}(a))) \wedge ((q \in \operatorname{mat}(n, l, \operatorname{carr}(a))) \wedge (1 \leq n)))) \Rightarrow  \\ \quad \quad \operatorname{entry}(\operatorname{matmul}(a, p, q), i, j) = \operatorname{entry}(\operatorname{matmul}(a, p, \operatorname{col\text{-}of}(q, n, j)), i, 1) \end{array}}}
\vnbprop{cq4-zero-column}{\vnbstmt{\begin{array}{@{}l@{}} \forall r, n, c, o.; \forall k \in \operatorname{interval}(1, n). \\ \quad \operatorname{is\text{-}ring}(r) \wedge \\ \quad n \in \mathbb{N} \wedge \\ \quad c \in \operatorname{mat}(n, n, \operatorname{carr}(r)) \wedge \\ \quad \forall \mathit{zkr} \in \operatorname{interval}(1, n).\; \operatorname{col\text{-}step\text{-}row}(r, n, c, \mathit{zkr}) \wedge \\ \quad o \in \operatorname{interval}(1, n) \wedge \\ \quad \operatorname{entry}(c, o, o) = \operatorname{zero}(r) \wedge \\ \quad \forall k \in \operatorname{interval}(1, n).\; ((o < k) \Rightarrow (\operatorname{entry}(c, k, k) = \operatorname{one}(r))) \Rightarrow \\ \quad \quad \operatorname{entry}(c, k, o) = \operatorname{zero}(r) \end{array}}}
\vnbprop{cq4-diag-zero-or-identity}{\vnbstmt{\begin{array}{@{}l@{}} \forall r, n, c. \\ \quad \operatorname{is\text{-}ring}(r) \wedge \\ \quad n \in \mathbb{N} \wedge \\ \quad c \in \operatorname{mat}(n, n, \operatorname{carr}(r)) \wedge \\ \quad \forall \mathit{zkr} \in \operatorname{interval}(1, n).\; \operatorname{col\text{-}step\text{-}row}(r, n, c, \mathit{zkr}) \Rightarrow \\ \quad \quad \exists e \in \operatorname{interval}(1, n).\; (\operatorname{entry}(c, e, e) = \operatorname{zero}(r)) \vee (c = \operatorname{identmat}(r, n)) \end{array}}}
\vnbprop{cq4-last-zero}{\vnbstmt{\begin{array}{@{}l@{}} \forall r, n, c. \\ \quad \operatorname{is\text{-}ring}(r) \wedge \\ \quad n \in \mathbb{N} \wedge \\ \quad c \in \operatorname{mat}(n, n, \operatorname{carr}(r)) \wedge \\ \quad \forall \mathit{zkr} \in \operatorname{interval}(1, n).\; \operatorname{col\text{-}step\text{-}row}(r, n, c, \mathit{zkr}) \wedge \\ \quad \exists e \in \operatorname{interval}(1, n).\; (\operatorname{entry}(c, e, e) = \operatorname{zero}(r)) \Rightarrow \\ \quad \quad \exists o \in \operatorname{interval}(1, n). \\ \quad \quad \quad ((\operatorname{entry}(c, o, o) = \operatorname{zero}(r)) \wedge \forall k \in \operatorname{interval}(1, n).\; ((o < k) \Rightarrow (\operatorname{entry}(c, k, k) = \operatorname{one}(r)))) \end{array}}}
\vnbprop{cq4-col-of-zero}{\vnbstmt{\begin{array}{@{}l@{}} \forall r.; \forall n \in \mathbb{N}, j \in \operatorname{interval}(1, n). \\ \quad \operatorname{is\text{-}ring}(r) \Rightarrow (\operatorname{col\text{-}of}(\operatorname{zeromat}(r, n, n), n, j) = \operatorname{zeromat}(r, n, 1)) \end{array}}}
\vnbprop{cq4-mm-zero-right}{\vnbstmt{\begin{array}{@{}l@{}} \forall r.; \forall n \in \mathbb{N}, p \in \operatorname{mat}(n, n, \operatorname{carr}(r)). \\ \quad \operatorname{is\text{-}ring}(r) \Rightarrow (\operatorname{matmul}(r, p, \operatorname{zeromat}(r, n, n)) = \operatorname{zeromat}(r, n, n)) \end{array}}}
\vnbprop{cq4-det-zero-kernel}{\vnbstmt{\begin{array}{@{}l@{}} \forall r.; \forall n \in \mathbb{N}, b \in \operatorname{mat}(n, n, \operatorname{carr}(r)). \\ \quad (\operatorname{is\text{-}field\text{-}ring}(r) \wedge (\operatorname{det}(r, n, b) = \operatorname{zero}(r))) \Rightarrow  \\ \quad \quad \exists e \in \operatorname{mat}(n, 1, \operatorname{carr}(r)). \\ \quad \quad \quad \neg (e = \operatorname{zeromat}(r, n, 1)) \wedge (\operatorname{matmul}(r, b, e) = \operatorname{zeromat}(r, n, 1)) \end{array}}}
\vnbprop{ccr-det-zero-kernel}{\vnbstmt{\begin{array}{@{}l@{}} \forall k.; \forall n \in \mathbb{N}, b \in \operatorname{mat}(n, n, \operatorname{carr}(k)). \\ \quad (\operatorname{is\text{-}normed\text{-}field}(k) \wedge (\operatorname{det}(\operatorname{normed\text{-}field\text{-}as\text{-}commutative\text{-}ring}(k), n, b) = \operatorname{zero}(k))) \Rightarrow  \\ \quad \quad \exists e \in \operatorname{mat}(n, 1, \operatorname{carr}(k)). \\ \quad \quad \quad \neg (e = \operatorname{zeromat}(\operatorname{normed\text{-}field\text{-}as\text{-}commutative\text{-}ring}(k), n, 1)) \wedge \\ \quad \quad \quad \operatorname{matmul}(\operatorname{normed\text{-}field\text{-}as\text{-}commutative\text{-}ring}(k), b, e) \\ \quad \quad \quad =\; \operatorname{zeromat}(\operatorname{normed\text{-}field\text{-}as\text{-}commutative\text{-}ring}(k), n, 1) \end{array}}}
\vnbprop{cq4-zero-mul-gen-left}{\vnbstmt{\begin{array}{@{}l@{}} \forall r, m, n, l, p. \\ \quad (\operatorname{is\text{-}ring}(r) \wedge ((m \in \mathbb{N}) \wedge ((l \in \mathbb{N}) \wedge ((1 \leq n) \wedge (p \in \operatorname{mat}(n, l, \operatorname{carr}(r))))))) \Rightarrow  \\ \quad \quad \operatorname{matmul}(r, \operatorname{zeromat}(r, m, n), p) = \operatorname{zeromat}(r, m, l) \end{array}}}
\vnbprop{cq4-zero-mul-gen-right}{\vnbstmt{\begin{array}{@{}l@{}} \forall r, m, n, l, p. \\ \quad (\operatorname{is\text{-}ring}(r) \wedge ((m \in \mathbb{N}) \wedge ((l \in \mathbb{N}) \wedge ((1 \leq n) \wedge (p \in \operatorname{mat}(m, n, \operatorname{carr}(r))))))) \Rightarrow  \\ \quad \quad \operatorname{matmul}(r, p, \operatorname{zeromat}(r, n, l)) = \operatorname{zeromat}(r, m, l) \end{array}}}
\vnbprop{cq4-commute-power}{\vnbstmt{\begin{array}{@{}l@{}} \forall \mathit{zkk} \in \mathbb{N}, k.; \forall n \in \mathbb{N}, m, b. \\ \quad \operatorname{is\text{-}normed\text{-}field}(k) \wedge \\ \quad m \in \operatorname{mat}(n, n, \operatorname{carr}(k)) \wedge \\ \quad b \in \operatorname{mat}(n, n, \operatorname{carr}(k)) \wedge \\ \quad \operatorname{matmul}(\operatorname{normed\text{-}field\text{-}as\text{-}commutative\text{-}ring}(k), b, m) \\ \quad =\; \operatorname{matmul}(\operatorname{normed\text{-}field\text{-}as\text{-}commutative\text{-}ring}(k), m, b) \Rightarrow \\ \quad \quad \operatorname{matmul}(\operatorname{normed\text{-}field\text{-}as\text{-}commutative\text{-}ring}(k), b, \operatorname{mat\text{-}power}(k, n, m, \mathit{zkk})) \\ \quad \quad =\; \operatorname{matmul}(\operatorname{normed\text{-}field\text{-}as\text{-}commutative\text{-}ring}(k), \operatorname{mat\text{-}power}(k, n, m, \mathit{zkk}), b) \end{array}}}
\vnbprop{cq4-krylov-type}{\vnbstmt{\begin{array}{@{}l@{}} \forall k.; \forall n \in \mathbb{N}, m, i. \\ \quad (\operatorname{is\text{-}normed\text{-}field}(k) \wedge ((m \in \operatorname{mat}(n, n, \operatorname{carr}(k))) \wedge (i \in \operatorname{mat}(n, 1, \operatorname{carr}(k))))) \Rightarrow  \\ \quad \quad \operatorname{krylov\text{-}mat}(k, n, m, i) \in \operatorname{mat}(n, n, \operatorname{carr}(k)) \end{array}}}
\vnbprop{cq4-krylov-column}{\vnbstmt{\begin{array}{@{}l@{}} \forall k.; \forall n \in \mathbb{N}, m, i.; \forall j \in \operatorname{interval}(1, n). \\ \quad (\operatorname{is\text{-}normed\text{-}field}(k) \wedge ((m \in \operatorname{mat}(n, n, \operatorname{carr}(k))) \wedge (i \in \operatorname{mat}(n, 1, \operatorname{carr}(k))))) \Rightarrow  \\ \quad \quad \operatorname{col\text{-}of}(\operatorname{krylov\text{-}mat}(k, n, m, i), n, j) \\ \quad \quad =\; \operatorname{matmul}(\operatorname{normed\text{-}field\text{-}as\text{-}commutative\text{-}ring}(k), \operatorname{mat\text{-}power}(k, n, m, (j - 1)), i) \end{array}}}
\vnbprop{cq4-commutant-kernel}{\vnbstmt{\begin{array}{@{}l@{}} \forall k.; \forall n \in \mathbb{N}, m, b, e. \\ \quad \operatorname{is\text{-}normed\text{-}field}(k) \wedge \\ \quad \operatorname{is\text{-}irreducible\text{-}mat}(k, n, m) \wedge \\ \quad b \in \operatorname{mat}(n, n, \operatorname{carr}(k)) \wedge \\ \quad \operatorname{matmul}(\operatorname{normed\text{-}field\text{-}as\text{-}commutative\text{-}ring}(k), b, m) \\ \quad =\; \operatorname{matmul}(\operatorname{normed\text{-}field\text{-}as\text{-}commutative\text{-}ring}(k), m, b) \wedge \\ \quad e \in \operatorname{mat}(n, 1, \operatorname{carr}(k)) \wedge \\ \quad \neg (e = \operatorname{zeromat}(\operatorname{normed\text{-}field\text{-}as\text{-}commutative\text{-}ring}(k), n, 1)) \wedge \\ \quad \operatorname{matmul}(\operatorname{normed\text{-}field\text{-}as\text{-}commutative\text{-}ring}(k), b, e) \\ \quad =\; \operatorname{zeromat}(\operatorname{normed\text{-}field\text{-}as\text{-}commutative\text{-}ring}(k), n, 1) \Rightarrow \\ \quad \quad b = \operatorname{zeromat}(\operatorname{normed\text{-}field\text{-}as\text{-}commutative\text{-}ring}(k), n, n) \end{array}}}
\vnbprop{ccr-3-4-commutant}{\vnbstmt{\begin{array}{@{}l@{}} \forall k.; \forall n \in \mathbb{N}, m, b. \\ \quad \operatorname{is\text{-}normed\text{-}field}(k) \wedge \\ \quad \operatorname{is\text{-}irreducible\text{-}mat}(k, n, m) \wedge \\ \quad b \in \operatorname{mat}(n, n, \operatorname{carr}(k)) \wedge \\ \quad \operatorname{matmul}(\operatorname{normed\text{-}field\text{-}as\text{-}commutative\text{-}ring}(k), b, m) \\ \quad =\; \operatorname{matmul}(\operatorname{normed\text{-}field\text{-}as\text{-}commutative\text{-}ring}(k), m, b) \wedge \\ \quad \operatorname{det}(\operatorname{normed\text{-}field\text{-}as\text{-}commutative\text{-}ring}(k), n, b) = \operatorname{zero}(k) \Rightarrow \\ \quad \quad b = \operatorname{zeromat}(\operatorname{normed\text{-}field\text{-}as\text{-}commutative\text{-}ring}(k), n, n) \end{array}}}

\vnbbookfile{theorem-library/ccr-ch3b-an}{theorem-library/ccr-ch3b-an}
\noindent{\small ccr-ch3b-an.scm -- wave of 2026-11-09/10 (agent CCR-1109-B2), wired 2026-11-10 from scratchpad/ccr1109b2/ccr-ch3b-an.scm. Library: analysis.}\par
\vnbprop{cq5-nondiscrete-small}{\vnbstmt{\begin{array}{@{}l@{}} \forall k. \\ \quad (\operatorname{is\text{-}normed\text{-}field}(k) \wedge \neg \operatorname{is\text{-}discrete\text{-}ms}(\operatorname{nf\text{-}metric\text{-}space}(k))) \Rightarrow  \\ \quad \quad \exists y \in \operatorname{carr}(k).\; ((0 < \operatorname{fnrm}(k)(y)) \wedge (\operatorname{fnrm}(k)(y) < 1)) \end{array}}}
\vnbprop{ccr-local-integers-compact}{\vnbstmt{\begin{array}{@{}l@{}} \forall k. \\ \quad \operatorname{is\text{-}local\text{-}nv\text{-}field}(k) \Rightarrow  \\ \quad \quad \operatorname{is\text{-}compact}(\operatorname{subspace\text{-}ms}(\operatorname{nf\text{-}metric\text{-}space}(k), \operatorname{nv\text{-}integers}(k))) \end{array}}}

\vnbbookfile{theorem-library/ccr-ch3c-det}{theorem-library/ccr-ch3c-det}
\noindent{\small ccr-ch3c-det.scm -- wave of 2026-11-09/10 (agent CCR-1109-B3), wired 2026-11-10 from scratchpad/ccr1109b3/ccr-ch3c-det.scm. Library: analysis.}\par
\vnbprop{cd6-add-sub-cancel}{\vnbstmt{\begin{array}{@{}l@{}} \forall k.; \forall u, w \in \operatorname{carr}(k). \\ \quad \operatorname{is\text{-}normed\text{-}field}(k) \Rightarrow (\operatorname{add}(k)(\operatorname{add}(k)(u, w), \operatorname{neg}(k)(w)) = u) \end{array}}}
\vnbprop{cd6-minor-close}{\vnbstmt{\begin{array}{@{}l@{}} \forall k.; \forall n \in \mathbb{N}, a, b, e.; \forall q \in \operatorname{interval}(1, (n + 1)). \\ \quad (\operatorname{is\text{-}normed\text{-}field}(k) \wedge \operatorname{is\text{-}mat\text{-}close}(k, (n + 1), (n + 1), a, b, e)) \Rightarrow  \\ \quad \quad \operatorname{is\text{-}mat\text{-}close}(k, n, n, \operatorname{minor}(a, 1, q, n), \operatorname{minor}(b, 1, q, n), e) \end{array}}}
\vnbprop{cd6-det-close-ind}{\vnbstmt{\begin{array}{@{}l@{}} \forall n \in \mathbb{N}, k.; \forall a, b \in \operatorname{mat}(n, n, \operatorname{carr}(k)), e \in \mathbb{R}. \\ \quad \operatorname{is\text{-}nv\text{-}field}(k) \wedge \\ \quad 0 \leq e \wedge \\ \quad \operatorname{is\text{-}mat\text{-}bound}(k, n, n, a, 1) \wedge \\ \quad \operatorname{is\text{-}mat\text{-}bound}(k, n, n, b, 1) \wedge \\ \quad \operatorname{is\text{-}mat\text{-}close}(k, n, n, a, b, e) \Rightarrow \\ \quad \quad \operatorname{fnrm}(k)(\operatorname{add}(k)(\operatorname{det}(\operatorname{normed\text{-}field\text{-}as\text{-}commutative\text{-}ring}(k), n, a), \\ \quad \quad \quad \quad \operatorname{neg}(k)(\operatorname{det}(\operatorname{normed\text{-}field\text{-}as\text{-}commutative\text{-}ring}(k), n, b)))) \\ \quad \quad \leq\; e \end{array}}}
\vnbprop{ccr-det-close}{\vnbstmt{\begin{array}{@{}l@{}} \forall k.; \forall n \in \mathbb{N}, a, b.; \forall e \in \mathbb{R}. \\ \quad \operatorname{is\text{-}nv\text{-}field}(k) \wedge \\ \quad 0 \leq e \wedge \\ \quad \operatorname{is\text{-}mat\text{-}bound}(k, n, n, a, 1) \wedge \\ \quad \operatorname{is\text{-}mat\text{-}bound}(k, n, n, b, 1) \wedge \\ \quad \operatorname{is\text{-}mat\text{-}close}(k, n, n, a, b, e) \Rightarrow \\ \quad \quad \operatorname{fnrm}(k)(\operatorname{add}(k)(\operatorname{det}(\operatorname{normed\text{-}field\text{-}as\text{-}commutative\text{-}ring}(k), n, a), \\ \quad \quad \quad \quad \operatorname{neg}(k)(\operatorname{det}(\operatorname{normed\text{-}field\text{-}as\text{-}commutative\text{-}ring}(k), n, b)))) \\ \quad \quad \leq\; e \end{array}}}

\vnbbookfile{theorem-library/ccr-ch3c-ball}{theorem-library/ccr-ch3c-ball}
\noindent{\small ccr-ch3c-ball.scm -- wave of 2026-11-09/10 (agent CCR-1109-B3), wired 2026-11-10 from scratchpad/ccr1109b3/ccr-ch3c-ball.scm. Library: analysis.}\par
\vnbprop{cd7-nvint-set}{\vnbstmt{\begin{array}{@{}l@{}} \forall k. \\ \quad \operatorname{is\text{-}normed\text{-}field}(k) \Rightarrow  \\ \quad \quad (\operatorname{nv\text{-}integers}(k) \in \mathbf{Set}) \wedge (\operatorname{nv\text{-}integers}(k) \subseteq \operatorname{pts}(\operatorname{nf\text{-}metric\text{-}space}(k))) \end{array}}}
\vnbprop{cd7-const-app}{\vnbstmt{\forall x \in \mathbf{Set}, t \in \mathbb{N}.\; ((\lambda \mathit{cd7t} \in \mathbb{N}.\; x)(t) \simeq x)}}
\vnbprop{cd7-const-ms}{\vnbstmt{\forall x.\; (\operatorname{is\text{-}metric\text{-}space}(x) \Rightarrow \operatorname{is\text{-}ms\text{-}sequence}((\lambda t \in \mathbb{N}.\; x)))}}
\vnbprop{cd7-coord-int}{\vnbstmt{\begin{array}{@{}l@{}} \forall k, n, m, s.; \forall a, t \in \mathbb{N}. \\ \quad (\operatorname{is\text{-}normed\text{-}field}(k) \wedge \forall a \in \mathbb{N}.\; \operatorname{is\text{-}mat\text{-}bound}(k, n, m, s(a), 1)) \Rightarrow  \\ \quad \quad \left(\text{if }((\operatorname{nnfst}(t) \in \operatorname{interval}(1, n)) \wedge (\operatorname{nnsnd}(t) \in \operatorname{interval}(1, m)))\text{ then }\operatorname{entry}(s(a), \operatorname{nnfst}(t), \operatorname{nnsnd}(t))\text{ else }\operatorname{zero}(k)\right) \\ \quad \quad \in\; \operatorname{nv\text{-}integers}(k) \end{array}}}
\vnbprop{cd7-sq-type}{\vnbstmt{\begin{array}{@{}l@{}} \forall k, n, m, s. \\ \quad (\operatorname{is\text{-}normed\text{-}field}(k) \wedge \forall a \in \mathbb{N}.\; \operatorname{is\text{-}mat\text{-}bound}(k, n, m, s(a), 1)) \Rightarrow  \\ \quad \quad (\lambda \mathit{cd7k} \in \mathbb{N}.\; \\ \quad \quad \quad (\lambda u \in \mathbb{N}.\; \\ \quad \quad \quad \quad \left(\text{if }((\operatorname{nnfst}(u) \in \operatorname{interval}(1, n)) \wedge (\operatorname{nnsnd}(u) \in \operatorname{interval}(1, m)))\text{ then }\operatorname{entry}(s(\mathit{cd7k}), \operatorname{nnfst}(u), \operatorname{nnsnd}(u))\text{ else }\operatorname{zero}(k)\right))) \\ \quad \quad \in\; (\mathbb{N} \to \operatorname{product\text{-}carrier}((\lambda t \in \mathbb{N}.\; \operatorname{subspace\text{-}ms}(\operatorname{nf\text{-}metric\text{-}space}(k), \operatorname{nv\text{-}integers}(k))))) \end{array}}}
\vnbprop{cd7-b-entry}{\vnbstmt{\begin{array}{@{}l@{}} \forall k.; \forall n, m \in \mathbb{N}, l \in \operatorname{product\text{-}carrier}((\lambda t \in \mathbb{N}.\; \operatorname{subspace\text{-}ms}(\operatorname{nf\text{-}metric\text{-}space}(k), \operatorname{nv\text{-}integers}(k)))), i \in \operatorname{interval}(1, n), j \in \operatorname{interval}(1, m). \\ \quad \operatorname{is\text{-}normed\text{-}field}(k) \Rightarrow  \\ \quad \quad l(\operatorname{nnpair}(i, j)) \in \operatorname{nv\text{-}integers}(k) \wedge \\ \quad \quad \operatorname{entry}(\operatorname{matof}(n, m, \\ \quad \quad \quad \quad (\lambda {[}\operatorname{cd7i\_}, \operatorname{cd7j\_}] \in (\operatorname{interval}(1, n) \times \operatorname{interval}(1, m)).\; l(\operatorname{nnpair}(\operatorname{cd7i\_}, \operatorname{cd7j\_})))), \\ \quad \quad \quad i, j) \\ \quad \quad =\; l(\operatorname{nnpair}(i, j)) \end{array}}}
\vnbprop{cd7-b-bound}{\vnbstmt{\begin{array}{@{}l@{}} \forall k.; \forall n, m \in \mathbb{N}, l \in \operatorname{product\text{-}carrier}((\lambda t \in \mathbb{N}.\; \operatorname{subspace\text{-}ms}(\operatorname{nf\text{-}metric\text{-}space}(k), \operatorname{nv\text{-}integers}(k)))). \\ \quad \operatorname{is\text{-}normed\text{-}field}(k) \Rightarrow  \\ \quad \quad \operatorname{is\text{-}mat\text{-}bound}(k, n, m, \\ \quad \quad \quad \operatorname{matof}(n, m, \\ \quad \quad \quad \quad (\lambda {[}\operatorname{cd7i\_}, \operatorname{cd7j\_}] \in (\operatorname{interval}(1, n) \times \operatorname{interval}(1, m)).\; l(\operatorname{nnpair}(\operatorname{cd7i\_}, \operatorname{cd7j\_})))), \\ \quad \quad \quad 1) \end{array}}}
\vnbprop{cd7-coord-conv}{\vnbstmt{\begin{array}{@{}l@{}} \forall k, n, m, s, a, l, b.; \forall t \in \mathbb{N}. \\ \quad \operatorname{is\text{-}normed\text{-}field}(k) \wedge \\ \quad \forall c \in \mathbb{N}.\; \operatorname{is\text{-}mat\text{-}bound}(k, n, m, s(c), 1) \wedge \\ \quad \operatorname{strictly\text{-}mono\text{-}nn}(a) \wedge \\ \quad \forall t \in \mathbb{N}. \\ \quad \quad \operatorname{converges\text{-}to}((\lambda \mathit{cd7t} \in \mathbb{N}.\; \operatorname{subspace\text{-}ms}(\operatorname{nf\text{-}metric\text{-}space}(k), \operatorname{nv\text{-}integers}(k)))(t), \\ \quad \quad \quad (\lambda v \in \mathbb{N}.\; \\ \quad \quad \quad \quad \operatorname{subseq}((\lambda \mathit{cd7k} \in \mathbb{N}.\; (\lambda u \in \mathbb{N}.\; \left(\text{if }((\operatorname{nnfst}(u) \in \operatorname{interval}(1, n)) \wedge (\operatorname{nnsnd}(u) \in \operatorname{interval}(1, m)))\text{ then }\operatorname{entry}(s(\mathit{cd7k}), \operatorname{nnfst}(u), \operatorname{nnsnd}(u))\text{ else }\operatorname{zero}(k)\right))), a)(v)(t)), \\ \quad \quad \quad l(t)) \wedge \\ \quad \operatorname{pos\text{-}rr}(b) \Rightarrow \\ \quad \quad \exists d \in \mathbb{N}. \forall c \in \mathbb{N}. \\ \quad \quad \quad (d \leq c) \Rightarrow  \\ \quad \quad \quad \quad \operatorname{dist}(\operatorname{subspace\text{-}ms}(\operatorname{nf\text{-}metric\text{-}space}(k), \operatorname{nv\text{-}integers}(k)))(\left(\text{if }((\operatorname{nnfst}(t) \in \operatorname{interval}(1, n)) \wedge (\operatorname{nnsnd}(t) \in \operatorname{interval}(1, m)))\text{ then }\operatorname{entry}(s(a(c)), \operatorname{nnfst}(t), \operatorname{nnsnd}(t))\text{ else }\operatorname{zero}(k)\right), \\ \quad \quad \quad \quad \quad l(t)) \\ \quad \quad \quad \quad \leq\; b \end{array}}}
\vnbprop{cd7-unif-ind}{\vnbstmt{\begin{array}{@{}l@{}} \forall c_{0} \in \mathbb{N}, k, n, m, s, a, l, b. \\ \quad \forall t \in \mathbb{N}. \exists c \in \mathbb{N}. \forall d \in \mathbb{N}. \\ \quad \quad (c \leq d) \Rightarrow  \\ \quad \quad \quad \operatorname{dist}(\operatorname{subspace\text{-}ms}(\operatorname{nf\text{-}metric\text{-}space}(k), \operatorname{nv\text{-}integers}(k)))(\left(\text{if }((\operatorname{nnfst}(t) \in \operatorname{interval}(1, n)) \wedge (\operatorname{nnsnd}(t) \in \operatorname{interval}(1, m)))\text{ then }\operatorname{entry}(s(a(d)), \operatorname{nnfst}(t), \operatorname{nnsnd}(t))\text{ else }\operatorname{zero}(k)\right), \\ \quad \quad \quad \quad l(t)) \\ \quad \quad \quad \leq\; b \Rightarrow \\ \quad \quad \exists c \in \mathbb{N}. \forall t, d \in \mathbb{N}. \\ \quad \quad \quad ((t \leq c_{0}) \wedge (c \leq d)) \Rightarrow  \\ \quad \quad \quad \quad \operatorname{dist}(\operatorname{subspace\text{-}ms}(\operatorname{nf\text{-}metric\text{-}space}(k), \operatorname{nv\text{-}integers}(k)))(\left(\text{if }((\operatorname{nnfst}(t) \in \operatorname{interval}(1, n)) \wedge (\operatorname{nnsnd}(t) \in \operatorname{interval}(1, m)))\text{ then }\operatorname{entry}(s(a(d)), \operatorname{nnfst}(t), \operatorname{nnsnd}(t))\text{ else }\operatorname{zero}(k)\right), \\ \quad \quad \quad \quad \quad l(t)) \\ \quad \quad \quad \quad \leq\; b \end{array}}}
\vnbprop{cd7-code-bound}{\vnbstmt{\begin{array}{@{}l@{}} \forall n, m \in \mathbb{N}, i \in \operatorname{interval}(1, n), j \in \operatorname{interval}(1, m). \\ \quad \operatorname{nnpair}(i, j) \leq (\operatorname{trinum}((n + m)) + m) \end{array}}}
\vnbprop{cd7-read}{\vnbstmt{\begin{array}{@{}l@{}} \forall k.; \forall n, m \in \mathbb{N}, s, a.; \forall l \in \operatorname{product\text{-}carrier}((\lambda t \in \mathbb{N}.\; \operatorname{subspace\text{-}ms}(\operatorname{nf\text{-}metric\text{-}space}(k), \operatorname{nv\text{-}integers}(k)))), b \in \mathbb{N}, i \in \operatorname{interval}(1, n), j \in \operatorname{interval}(1, m). \\ \quad (\operatorname{is\text{-}normed\text{-}field}(k) \wedge (\forall b \in \mathbb{N}.\; \operatorname{is\text{-}mat\text{-}bound}(k, n, m, s(b), 1) \wedge \operatorname{strictly\text{-}mono\text{-}nn}(a))) \Rightarrow  \\ \quad \quad \operatorname{dist}(\operatorname{subspace\text{-}ms}(\operatorname{nf\text{-}metric\text{-}space}(k), \operatorname{nv\text{-}integers}(k)))(\left(\text{if }((\operatorname{nnfst}(\operatorname{nnpair}(i, j)) \in \operatorname{interval}(1, n)) \wedge (\operatorname{nnsnd}(\operatorname{nnpair}(i, j)) \in \operatorname{interval}(1, m)))\text{ then }\operatorname{entry}(s(a(b)), \operatorname{nnfst}(\operatorname{nnpair}(i, j)), \operatorname{nnsnd}(\operatorname{nnpair}(i, j)))\text{ else }\operatorname{zero}(k)\right), \\ \quad \quad \quad l(\operatorname{nnpair}(i, j))) \\ \quad \quad =\; \operatorname{fnrm}(k)(\operatorname{add}(k)(\operatorname{entry}(s(a(b)), i, j), \\ \quad \quad \quad \quad \operatorname{neg}(k)(\operatorname{entry}(\operatorname{matof}(n, m, \\ \quad \quad \quad \quad \quad \quad \quad (\lambda {[}\operatorname{cd7i\_}, \operatorname{cd7j\_}] \in (\operatorname{interval}(1, n) \times \operatorname{interval}(1, m)).\; l(\operatorname{nnpair}(\operatorname{cd7i\_}, \operatorname{cd7j\_})))), \\ \quad \quad \quad \quad \quad \quad i, j)))) \end{array}}}
\vnbprop{ccr-mat-ball-seq-compact}{\vnbstmt{\begin{array}{@{}l@{}} \forall k.; \forall n, m \in \mathbb{N}, s. \\ \quad (\operatorname{is\text{-}local\text{-}nv\text{-}field}(k) \wedge \forall a \in \mathbb{N}.\; \operatorname{is\text{-}mat\text{-}bound}(k, n, m, s(a), 1)) \Rightarrow  \\ \quad \quad \exists i. \\ \quad \quad \quad \operatorname{strictly\text{-}mono\text{-}nn}(i) \wedge \\ \quad \quad \quad \exists b. \\ \quad \quad \quad \quad \operatorname{is\text{-}mat\text{-}bound}(k, n, m, b, 1) \wedge \\ \quad \quad \quad \quad \forall c. \\ \quad \quad \quad \quad \quad \operatorname{pos\text{-}rr}(c) \Rightarrow  \\ \quad \quad \quad \quad \quad \quad \exists d \in \mathbb{N}. \forall a \in \mathbb{N}. \\ \quad \quad \quad \quad \quad \quad \quad (d \leq a) \Rightarrow \operatorname{is\text{-}mat\text{-}close}(k, n, m, s(i(a)), b, c) \end{array}}}

\vnbbookfile{theorem-library/ccr-ch3c-pb}{theorem-library/ccr-ch3c-pb}
\noindent{\small ccr-ch3c-pb.scm -- wave of 2026-11-09/10 (agent CCR-1109-B3), wired 2026-11-10 from scratchpad/ccr1109b3/ccr-ch3c-pb.scm. Library: analysis.}\par
\vnbprop{cd8-scale-bound}{\vnbstmt{\begin{array}{@{}l@{}} \forall k, n, m.; \forall y \in \operatorname{carr}(k), a.; \forall c, d \in \mathbb{R}. \\ \quad (\operatorname{is\text{-}normed\text{-}field}(k) \wedge (\operatorname{is\text{-}mat\text{-}bound}(k, n, m, a, c) \wedge ((\operatorname{fnrm}(k)(y) \cdot c) \leq d))) \Rightarrow  \\ \quad \quad \operatorname{is\text{-}mat\text{-}bound}(k, n, m, \operatorname{matscale}(\operatorname{normed\text{-}field\text{-}as\text{-}commutative\text{-}ring}(k), y, a), d) \end{array}}}
\vnbprop{cd8-argmax-nn}{\vnbstmt{\begin{array}{@{}l@{}} \forall n \in \mathbb{N}, f \in (\mathbb{N} \to \mathbb{R}). \exists t \in \mathbb{N}. \\ \quad (t \leq n) \wedge \forall s \in \mathbb{N}.\; ((s \leq n) \Rightarrow (f(s) \leq f(t))) \end{array}}}
\vnbprop{cd8-power-le-base}{\vnbstmt{\forall n \in \mathbb{N}, r \in \mathbb{R}.\; (((1 \leq n) \wedge ((0 \leq r) \wedge (r \leq 1))) \Rightarrow (\operatorname{power}(r, n) \leq r))}}
\vnbprop{cd8-mat-argmax}{\vnbstmt{\begin{array}{@{}l@{}} \forall k.; \forall n \in \mathbb{N}, a \in \operatorname{mat}(n, n, \operatorname{carr}(k)). \\ \quad (\operatorname{is\text{-}normed\text{-}field}(k) \wedge (1 \leq n)) \Rightarrow  \\ \quad \quad \exists i_{0}, j_{0} \in \operatorname{interval}(1, n). \forall i, j \in \operatorname{interval}(1, n). \\ \quad \quad \quad \operatorname{fnrm}(k)(\operatorname{entry}(a, i, j)) \leq \operatorname{fnrm}(k)(\operatorname{entry}(a, i_{0}, j_{0})) \end{array}}}
\vnbprop{cd8-mat-has-bound}{\vnbstmt{\begin{array}{@{}l@{}} \forall k.; \forall n \in \mathbb{N}, a \in \operatorname{mat}(n, n, \operatorname{carr}(k)). \\ \quad (\operatorname{is\text{-}normed\text{-}field}(k) \wedge (1 \leq n)) \Rightarrow  \\ \quad \quad \exists c \in \mathbb{R}.\; ((0 \leq c) \wedge \operatorname{is\text{-}mat\text{-}bound}(k, n, n, a, c)) \end{array}}}
\vnbprop{cd8-finsum-diff-le}{\vnbstmt{\begin{array}{@{}l@{}} \forall k.; \forall q \in \mathbf{Set}, a, b.; \forall r \in \mathbb{R}. \\ \quad \operatorname{is\text{-}nv\text{-}field}(k) \wedge \\ \quad |q| \in \mathbb{N} \wedge \\ \quad 0 \leq r \wedge \\ \quad \forall z \in q. \\ \quad \quad ((a(z) \in \operatorname{carr}(k)) \wedge ((b(z) \in \operatorname{carr}(k)) \wedge (\operatorname{fnrm}(k)(\operatorname{add}(k)(a(z), \operatorname{neg}(k)(b(z)))) \leq r))) \Rightarrow \\ \quad \quad \operatorname{fnrm}(k)(\operatorname{add}(k)(\operatorname{finsum}(\operatorname{ring\text{-}additive\text{-}ag}(\operatorname{normed\text{-}field\text{-}as\text{-}commutative\text{-}ring}(k)), a, q), \\ \quad \quad \quad \quad \operatorname{neg}(k)(\operatorname{finsum}(\operatorname{ring\text{-}additive\text{-}ag}(\operatorname{normed\text{-}field\text{-}as\text{-}commutative\text{-}ring}(k)), b, q)))) \\ \quad \quad \leq\; r \end{array}}}
\vnbprop{cd8-close-mul-right}{\vnbstmt{\begin{array}{@{}l@{}} \forall k.; \forall n \in \mathbb{N}, a, b, m, e, c. \\ \quad \operatorname{is\text{-}nv\text{-}field}(k) \wedge \\ \quad e \in \mathbb{R} \wedge \\ \quad 0 \leq e \wedge \\ \quad c \in \mathbb{R} \wedge \\ \quad 0 \leq c \wedge \\ \quad \operatorname{is\text{-}mat\text{-}close}(k, n, n, a, b, e) \wedge \\ \quad \operatorname{is\text{-}mat\text{-}bound}(k, n, n, m, c) \Rightarrow \\ \quad \quad \operatorname{is\text{-}mat\text{-}close}(k, n, n, \\ \quad \quad \quad \operatorname{matmul}(\operatorname{normed\text{-}field\text{-}as\text{-}commutative\text{-}ring}(k), a, m), \\ \quad \quad \quad \operatorname{matmul}(\operatorname{normed\text{-}field\text{-}as\text{-}commutative\text{-}ring}(k), b, m), \\ \quad \quad \quad (e \cdot c)) \end{array}}}
\vnbprop{cd8-close-mul-left}{\vnbstmt{\begin{array}{@{}l@{}} \forall k.; \forall n \in \mathbb{N}, a, b, m, e, c. \\ \quad \operatorname{is\text{-}nv\text{-}field}(k) \wedge \\ \quad e \in \mathbb{R} \wedge \\ \quad 0 \leq e \wedge \\ \quad c \in \mathbb{R} \wedge \\ \quad 0 \leq c \wedge \\ \quad \operatorname{is\text{-}mat\text{-}close}(k, n, n, a, b, e) \wedge \\ \quad \operatorname{is\text{-}mat\text{-}bound}(k, n, n, m, c) \Rightarrow \\ \quad \quad \operatorname{is\text{-}mat\text{-}close}(k, n, n, \\ \quad \quad \quad \operatorname{matmul}(\operatorname{normed\text{-}field\text{-}as\text{-}commutative\text{-}ring}(k), m, a), \\ \quad \quad \quad \operatorname{matmul}(\operatorname{normed\text{-}field\text{-}as\text{-}commutative\text{-}ring}(k), m, b), \\ \quad \quad \quad (c \cdot e)) \end{array}}}
\vnbprop{cd8-mpow-type}{\vnbstmt{\begin{array}{@{}l@{}} \forall k.; \forall n \in \mathbb{N}, x \in \operatorname{mat}(n, n, \operatorname{carr}(k)), a \in \mathbb{N}. \\ \quad \operatorname{is\text{-}normed\text{-}field}(k) \Rightarrow (\operatorname{mat\text{-}power}(k, n, x, a) \in \operatorname{mat}(n, n, \operatorname{carr}(k))) \end{array}}}
\vnbprop{cd8-unbounded}{\vnbstmt{\begin{array}{@{}l@{}} \forall k.; \forall n \in \mathbb{N}, x \in \operatorname{mat}(n, n, \operatorname{carr}(k)), c \in \mathbb{R}. \\ \quad (\operatorname{is\text{-}normed\text{-}field}(k) \wedge (\neg \operatorname{is\text{-}power\text{-}bounded}(k, n, x) \wedge (0 \leq c))) \Rightarrow  \\ \quad \quad \exists j \in \mathbb{N}, i, l \in \operatorname{interval}(1, n). \\ \quad \quad \quad c < \operatorname{fnrm}(k)(\operatorname{entry}(\operatorname{mat\text{-}power}(k, n, x, j), i, l)) \end{array}}}
\vnbprop{cd8-scale-commute}{\vnbstmt{\begin{array}{@{}l@{}} \forall k.; \forall n \in \mathbb{N}, p, m, y. \\ \quad \operatorname{is\text{-}normed\text{-}field}(k) \wedge \\ \quad 1 \leq n \wedge \\ \quad p \in \operatorname{mat}(n, n, \operatorname{carr}(k)) \wedge \\ \quad m \in \operatorname{mat}(n, n, \operatorname{carr}(k)) \wedge \\ \quad y \in \operatorname{carr}(k) \wedge \\ \quad \operatorname{matmul}(\operatorname{normed\text{-}field\text{-}as\text{-}commutative\text{-}ring}(k), m, p) \\ \quad =\; \operatorname{matmul}(\operatorname{normed\text{-}field\text{-}as\text{-}commutative\text{-}ring}(k), p, m) \Rightarrow \\ \quad \quad \operatorname{matmul}(\operatorname{normed\text{-}field\text{-}as\text{-}commutative\text{-}ring}(k), \operatorname{matscale}(\operatorname{normed\text{-}field\text{-}as\text{-}commutative\text{-}ring}(k), y, p), m) \\ \quad \quad =\; \operatorname{matmul}(\operatorname{normed\text{-}field\text{-}as\text{-}commutative\text{-}ring}(k), m, \operatorname{matscale}(\operatorname{normed\text{-}field\text{-}as\text{-}commutative\text{-}ring}(k), y, p)) \end{array}}}
\vnbprop{cd8-scale-as-mul}{\vnbstmt{\begin{array}{@{}l@{}} \forall k.; \forall n \in \mathbb{N}, p, y. \\ \quad (\operatorname{is\text{-}normed\text{-}field}(k) \wedge ((1 \leq n) \wedge ((p \in \operatorname{mat}(n, n, \operatorname{carr}(k))) \wedge (y \in \operatorname{carr}(k))))) \Rightarrow  \\ \quad \quad \operatorname{matscale}(\operatorname{normed\text{-}field\text{-}as\text{-}commutative\text{-}ring}(k), y, p) \\ \quad \quad =\; \operatorname{matmul}(\operatorname{normed\text{-}field\text{-}as\text{-}commutative\text{-}ring}(k), \\ \quad \quad \quad \operatorname{matscale}(\operatorname{normed\text{-}field\text{-}as\text{-}commutative\text{-}ring}(k), y, \operatorname{identmat}(\operatorname{normed\text{-}field\text{-}as\text{-}commutative\text{-}ring}(k), n)), \\ \quad \quad \quad p) \end{array}}}
\vnbprop{cd8-step}{\vnbstmt{\begin{array}{@{}l@{}} \forall k.; \forall n \in \mathbb{N}, x, m.; \forall a \in \mathbb{N}. \\ \quad \operatorname{is\text{-}nv\text{-}field}(k) \wedge \\ \quad x \in \operatorname{mat}(n, n, \operatorname{carr}(k)) \wedge \\ \quad m \in \operatorname{mat}(n, n, \operatorname{carr}(k)) \wedge \\ \quad \operatorname{matmul}(\operatorname{normed\text{-}field\text{-}as\text{-}commutative\text{-}ring}(k), x, m) \\ \quad =\; \operatorname{matmul}(\operatorname{normed\text{-}field\text{-}as\text{-}commutative\text{-}ring}(k), m, x) \wedge \\ \quad \operatorname{fnrm}(k)(\operatorname{det}(\operatorname{normed\text{-}field\text{-}as\text{-}commutative\text{-}ring}(k), n, x)) \leq 1 \wedge \\ \quad \neg \operatorname{is\text{-}power\text{-}bounded}(k, n, x) \Rightarrow \\ \quad \quad \exists y \in \operatorname{mat}(n, n, \operatorname{carr}(k)). \\ \quad \quad \quad \operatorname{is\text{-}mat\text{-}bound}(k, n, n, y, 1) \wedge \\ \quad \quad \quad \operatorname{matmul}(\operatorname{normed\text{-}field\text{-}as\text{-}commutative\text{-}ring}(k), y, m) \\ \quad \quad \quad =\; \operatorname{matmul}(\operatorname{normed\text{-}field\text{-}as\text{-}commutative\text{-}ring}(k), m, y) \wedge \\ \quad \quad \quad (\operatorname{fnrm}(k)(\operatorname{det}(\operatorname{normed\text{-}field\text{-}as\text{-}commutative\text{-}ring}(k), n, y)) \cdot (a + 1)) \leq 1 \wedge \\ \quad \quad \quad \exists i, j \in \operatorname{interval}(1, n). \\ \quad \quad \quad \quad \operatorname{fnrm}(k)(\operatorname{entry}(y, i, j)) = 1 \end{array}}}
\vnbprop{ccr-3-4-comm}{\vnbstmt{\begin{array}{@{}l@{}} \forall k.; \forall n \in \mathbb{N}, x, m. \\ \quad \operatorname{is\text{-}local\text{-}nv\text{-}field}(k) \wedge \\ \quad \operatorname{is\text{-}irreducible\text{-}mat}(k, n, m) \wedge \\ \quad x \in \operatorname{mat}(n, n, \operatorname{carr}(k)) \wedge \\ \quad \operatorname{matmul}(\operatorname{normed\text{-}field\text{-}as\text{-}commutative\text{-}ring}(k), x, m) \\ \quad =\; \operatorname{matmul}(\operatorname{normed\text{-}field\text{-}as\text{-}commutative\text{-}ring}(k), m, x) \wedge \\ \quad \operatorname{fnrm}(k)(\operatorname{det}(\operatorname{normed\text{-}field\text{-}as\text{-}commutative\text{-}ring}(k), n, x)) \leq 1 \Rightarrow \\ \quad \quad \operatorname{is\text{-}power\text{-}bounded}(k, n, x) \end{array}}}
\vnbprop{ccr-3-4}{\vnbstmt{\begin{array}{@{}l@{}} \forall k.; \forall n \in \mathbb{N}, m. \\ \quad \operatorname{is\text{-}local\text{-}nv\text{-}field}(k) \wedge \\ \quad \operatorname{is\text{-}irreducible\text{-}mat}(k, n, m) \wedge \\ \quad \operatorname{fnrm}(k)(\operatorname{det}(\operatorname{normed\text{-}field\text{-}as\text{-}commutative\text{-}ring}(k), n, m)) \leq 1 \Rightarrow \\ \quad \quad \operatorname{is\text{-}power\text{-}bounded}(k, n, m) \end{array}}}
\vnbprop{ccr-3-6-ultra-inv}{\vnbstmt{\begin{array}{@{}l@{}} \forall k.; \forall n \in \mathbb{N}, m.; \forall a, b, c \in \operatorname{mat}(n, n, \operatorname{carr}(k)). \\ \quad \operatorname{is\text{-}local\text{-}nv\text{-}field}(k) \wedge \\ \quad \operatorname{is\text{-}irreducible\text{-}mat}(k, n, m) \wedge \\ \quad \operatorname{matmul}(\operatorname{normed\text{-}field\text{-}as\text{-}commutative\text{-}ring}(k), a, m) \\ \quad =\; \operatorname{matmul}(\operatorname{normed\text{-}field\text{-}as\text{-}commutative\text{-}ring}(k), m, a) \wedge \\ \quad \operatorname{matmul}(\operatorname{normed\text{-}field\text{-}as\text{-}commutative\text{-}ring}(k), b, m) \\ \quad =\; \operatorname{matmul}(\operatorname{normed\text{-}field\text{-}as\text{-}commutative\text{-}ring}(k), m, b) \wedge \\ \quad \operatorname{matmul}(\operatorname{normed\text{-}field\text{-}as\text{-}commutative\text{-}ring}(k), c, a) = \operatorname{identmat}(\operatorname{normed\text{-}field\text{-}as\text{-}commutative\text{-}ring}(k), n) \wedge \\ \quad \operatorname{matmul}(\operatorname{normed\text{-}field\text{-}as\text{-}commutative\text{-}ring}(k), a, c) = \operatorname{identmat}(\operatorname{normed\text{-}field\text{-}as\text{-}commutative\text{-}ring}(k), n) \wedge \\ \quad \operatorname{fnrm}(k)(\operatorname{det}(\operatorname{normed\text{-}field\text{-}as\text{-}commutative\text{-}ring}(k), n, b)) \\ \quad \leq\; \operatorname{fnrm}(k)(\operatorname{det}(\operatorname{normed\text{-}field\text{-}as\text{-}commutative\text{-}ring}(k), n, a)) \Rightarrow \\ \quad \quad \operatorname{fnrm}(k)(\operatorname{det}(\operatorname{normed\text{-}field\text{-}as\text{-}commutative\text{-}ring}(k), n, \operatorname{matadd}(\operatorname{normed\text{-}field\text{-}as\text{-}commutative\text{-}ring}(k), a, b))) \\ \quad \quad \leq\; \operatorname{fnrm}(k)(\operatorname{det}(\operatorname{normed\text{-}field\text{-}as\text{-}commutative\text{-}ring}(k), n, a)) \end{array}}}
\vnbprop{ccr-3-6-ultra}{\vnbstmt{\begin{array}{@{}l@{}} \forall k.; \forall n \in \mathbb{N}, m.; \forall a, b \in \operatorname{mat}(n, n, \operatorname{carr}(k)). \\ \quad \operatorname{is\text{-}local\text{-}nv\text{-}field}(k) \wedge \\ \quad \operatorname{is\text{-}irreducible\text{-}mat}(k, n, m) \wedge \\ \quad \operatorname{matmul}(\operatorname{normed\text{-}field\text{-}as\text{-}commutative\text{-}ring}(k), a, m) \\ \quad =\; \operatorname{matmul}(\operatorname{normed\text{-}field\text{-}as\text{-}commutative\text{-}ring}(k), m, a) \wedge \\ \quad \operatorname{matmul}(\operatorname{normed\text{-}field\text{-}as\text{-}commutative\text{-}ring}(k), b, m) \\ \quad =\; \operatorname{matmul}(\operatorname{normed\text{-}field\text{-}as\text{-}commutative\text{-}ring}(k), m, b) \wedge \\ \quad \neg (\operatorname{det}(\operatorname{normed\text{-}field\text{-}as\text{-}commutative\text{-}ring}(k), n, a) = \operatorname{zero}(k)) \Rightarrow  \\ \quad \quad \exists p \in \operatorname{mat}(n, n, \operatorname{carr}(k)). \\ \quad \quad \quad \operatorname{matmul}(\operatorname{normed\text{-}field\text{-}as\text{-}commutative\text{-}ring}(k), p, a) = \operatorname{identmat}(\operatorname{normed\text{-}field\text{-}as\text{-}commutative\text{-}ring}(k), n) \wedge \\ \quad \quad \quad \operatorname{matmul}(\operatorname{normed\text{-}field\text{-}as\text{-}commutative\text{-}ring}(k), a, p) = \operatorname{identmat}(\operatorname{normed\text{-}field\text{-}as\text{-}commutative\text{-}ring}(k), n) \wedge \\ \quad \neg (\operatorname{det}(\operatorname{normed\text{-}field\text{-}as\text{-}commutative\text{-}ring}(k), n, b) = \operatorname{zero}(k)) \Rightarrow  \\ \quad \quad \exists p \in \operatorname{mat}(n, n, \operatorname{carr}(k)). \\ \quad \quad \quad \operatorname{matmul}(\operatorname{normed\text{-}field\text{-}as\text{-}commutative\text{-}ring}(k), p, b) = \operatorname{identmat}(\operatorname{normed\text{-}field\text{-}as\text{-}commutative\text{-}ring}(k), n) \wedge \\ \quad \quad \quad \operatorname{matmul}(\operatorname{normed\text{-}field\text{-}as\text{-}commutative\text{-}ring}(k), b, p) = \operatorname{identmat}(\operatorname{normed\text{-}field\text{-}as\text{-}commutative\text{-}ring}(k), n) \Rightarrow \\ \quad \quad \operatorname{fnrm}(k)(\operatorname{det}(\operatorname{normed\text{-}field\text{-}as\text{-}commutative\text{-}ring}(k), n, \operatorname{matadd}(\operatorname{normed\text{-}field\text{-}as\text{-}commutative\text{-}ring}(k), a, b))) \\ \quad \quad \leq\; \operatorname{max}(\operatorname{fnrm}(k)(\operatorname{det}(\operatorname{normed\text{-}field\text{-}as\text{-}commutative\text{-}ring}(k), n, a)), \\ \quad \quad \quad \operatorname{fnrm}(k)(\operatorname{det}(\operatorname{normed\text{-}field\text{-}as\text{-}commutative\text{-}ring}(k), n, b))) \end{array}}}

\vnbbookfile{theorem-library/ccr4x-ext}{theorem-library/ccr4x-ext}
\noindent{\small ccr4x-ext.scm -- wave of 2026-11-09/10 (agent CCR-1109-C2), wired 2026-11-10 from scratchpad/ccr1109c2/ccr4x-ext.scm. Library: analysis.}\par
\vnbprop{map-seq-compose}{\vnbstmt{\begin{array}{@{}l@{}} \forall s, h, g, a.; \forall n \in \mathbb{N}, u \in \operatorname{mat}(n, 1, s). \\ \quad ((h \in (s \to s)) \wedge ((g \in (s \to s)) \wedge ((a \in (s \to s)) \wedge \forall z \in s.\; (a(z) = g(h(z)))))) \Rightarrow  \\ \quad \quad \operatorname{map\text{-}seq}(a, n, u) = \operatorname{map\text{-}seq}(g, n, \operatorname{map\text{-}seq}(h, n, u)) \end{array}}}
\vnbprop{map-seq-ident}{\vnbstmt{\begin{array}{@{}l@{}} \forall s.; \forall h \in (s \to s), n \in \mathbb{N}, u \in \operatorname{mat}(n, 1, s). \\ \quad \forall z \in s.\; (h(z) = z) \Rightarrow (\operatorname{map\text{-}seq}(h, n, u) = u) \end{array}}}
\vnbprop{op-mat-spec}{\vnbstmt{\begin{array}{@{}l@{}} \forall d, n, u.; \forall t \in (\operatorname{vec}(d) \to \operatorname{vec}(d)). \\ \quad \operatorname{is\text{-}module}(d) \wedge \\ \quad n \in \mathbb{N} \wedge \\ \quad 1 \leq n \wedge \\ \quad u \in \operatorname{mat}(n, 1, \operatorname{vec}(d)) \wedge \\ \quad \operatorname{generates}(d, n, u) \wedge \\ \quad \operatorname{rel\text{-}free}(d, n, u) \Rightarrow \\ \quad \quad \operatorname{op\text{-}mat}(d, n, u, t) \in \operatorname{mat}(n, n, \operatorname{carr}(\operatorname{scal}(d))) \wedge \\ \quad \quad \operatorname{map\text{-}seq}(t, n, u) = \operatorname{matact}(d, \operatorname{op\text{-}mat}(d, n, u, t), u) \end{array}}}
\vnbprop{op-mat-unique}{\vnbstmt{\begin{array}{@{}l@{}} \forall d, n, u.; \forall t \in (\operatorname{vec}(d) \to \operatorname{vec}(d)), m \in \operatorname{mat}(n, n, \operatorname{carr}(\operatorname{scal}(d))). \\ \quad \operatorname{is\text{-}module}(d) \wedge \\ \quad n \in \mathbb{N} \wedge \\ \quad 1 \leq n \wedge \\ \quad u \in \operatorname{mat}(n, 1, \operatorname{vec}(d)) \wedge \\ \quad \operatorname{generates}(d, n, u) \wedge \\ \quad \operatorname{rel\text{-}free}(d, n, u) \wedge \\ \quad \operatorname{map\text{-}seq}(t, n, u) = \operatorname{matact}(d, m, u) \Rightarrow \\ \quad \quad \operatorname{op\text{-}mat}(d, n, u, t) = m \end{array}}}
\vnbprop{op-mat-compose}{\vnbstmt{\begin{array}{@{}l@{}} \forall d, n, u, t, s, r. \\ \quad \operatorname{is\text{-}module}(d) \wedge \\ \quad n \in \mathbb{N} \wedge \\ \quad 1 \leq n \wedge \\ \quad u \in \operatorname{mat}(n, 1, \operatorname{vec}(d)) \wedge \\ \quad \operatorname{generates}(d, n, u) \wedge \\ \quad \operatorname{rel\text{-}free}(d, n, u) \wedge \\ \quad \operatorname{is\text{-}hom\text{-}module}(d, d, t) \wedge \\ \quad \operatorname{is\text{-}hom\text{-}module}(d, d, s) \wedge \\ \quad r \in (\operatorname{vec}(d) \to \operatorname{vec}(d)) \wedge \\ \quad \forall z \in \operatorname{vec}(d).\; (r(z) = s(t(z))) \Rightarrow \\ \quad \quad \operatorname{op\text{-}mat}(d, n, u, r) = \operatorname{matmul}(\operatorname{scal}(d), \operatorname{op\text{-}mat}(d, n, u, t), \operatorname{op\text{-}mat}(d, n, u, s)) \end{array}}}
\vnbprop{op-mat-id}{\vnbstmt{\begin{array}{@{}l@{}} \forall d, n, u.; \forall r \in (\operatorname{vec}(d) \to \operatorname{vec}(d)). \\ \quad \operatorname{is\text{-}module}(d) \wedge \\ \quad n \in \mathbb{N} \wedge \\ \quad 1 \leq n \wedge \\ \quad u \in \operatorname{mat}(n, 1, \operatorname{vec}(d)) \wedge \\ \quad \operatorname{generates}(d, n, u) \wedge \\ \quad \operatorname{rel\text{-}free}(d, n, u) \wedge \\ \quad \forall z \in \operatorname{vec}(d).\; (r(z) = z) \Rightarrow \\ \quad \quad \operatorname{op\text{-}mat}(d, n, u, r) = \operatorname{identmat}(\operatorname{scal}(d), n) \end{array}}}
\vnbprop{mult-op-type}{\vnbstmt{\begin{array}{@{}l@{}} \forall f.; \forall i \in \operatorname{carr}(f). \\ \quad \operatorname{is\text{-}ring}(f) \Rightarrow (\operatorname{mult\text{-}op}(f, i) \in (\operatorname{carr}(f) \to \operatorname{carr}(f))) \end{array}}}
\vnbprop{mult-op-value}{\vnbstmt{\begin{array}{@{}l@{}} \forall f.; \forall i, w \in \operatorname{carr}(f). \\ \quad \operatorname{is\text{-}ring}(f) \Rightarrow (\operatorname{mult\text{-}op}(f, i)(w) = \operatorname{mul}(f)(i, w)) \end{array}}}
\vnbprop{ext-mult-op-hom}{\vnbstmt{\begin{array}{@{}l@{}} \forall k, f, d, n, u.; \forall i \in \operatorname{carr}(f). \\ \quad \operatorname{is\text{-}fin\text{-}ext}(k, f, d, n, u) \Rightarrow \operatorname{is\text{-}hom\text{-}module}(d, d, \operatorname{mult\text{-}op}(f, i)) \end{array}}}
\vnbprop{ext-mult-mat-spec}{\vnbstmt{\begin{array}{@{}l@{}} \forall k, f, d, n, u.; \forall i \in \operatorname{carr}(f). \\ \quad \operatorname{is\text{-}fin\text{-}ext}(k, f, d, n, u) \Rightarrow  \\ \quad \quad \operatorname{mult\text{-}mat}(d, n, u, f, i) \in \operatorname{mat}(n, n, \operatorname{carr}(k)) \wedge \\ \quad \quad \operatorname{map\text{-}seq}(\operatorname{mult\text{-}op}(f, i), n, u) = \operatorname{matact}(d, \operatorname{mult\text{-}mat}(d, n, u, f, i), u) \end{array}}}
\vnbprop{ext-mult-mat-mul}{\vnbstmt{\begin{array}{@{}l@{}} \forall k, f, d, n, u.; \forall i, a \in \operatorname{carr}(f). \\ \quad \operatorname{is\text{-}fin\text{-}ext}(k, f, d, n, u) \Rightarrow  \\ \quad \quad \operatorname{mult\text{-}mat}(d, n, u, f, \operatorname{mul}(f)(i, a)) \\ \quad \quad =\; \operatorname{matmul}(\operatorname{normed\text{-}field\text{-}as\text{-}commutative\text{-}ring}(k), \operatorname{mult\text{-}mat}(d, n, u, f, a), \operatorname{mult\text{-}mat}(d, n, u, f, i)) \end{array}}}
\vnbprop{ext-mult-mat-one}{\vnbstmt{\begin{array}{@{}l@{}} \forall k, f, d, n, u. \\ \quad \operatorname{is\text{-}fin\text{-}ext}(k, f, d, n, u) \Rightarrow  \\ \quad \quad \operatorname{mult\text{-}mat}(d, n, u, f, \operatorname{one}(f)) = \operatorname{identmat}(\operatorname{normed\text{-}field\text{-}as\text{-}commutative\text{-}ring}(k), n) \end{array}}}
\vnbprop{ext-abs-type}{\vnbstmt{\begin{array}{@{}l@{}} \forall k, f, d, n, u.; \forall i \in \operatorname{carr}(f). \\ \quad \operatorname{is\text{-}fin\text{-}ext}(k, f, d, n, u) \Rightarrow  \\ \quad \quad (\operatorname{ext\text{-}abs}(k, f, d, n, u, i) \in \mathbb{R}) \wedge (0 \leq \operatorname{ext\text{-}abs}(k, f, d, n, u, i)) \end{array}}}
\vnbprop{ext-abs-mul}{\vnbstmt{\begin{array}{@{}l@{}} \forall k, f, d, n, u.; \forall i, a \in \operatorname{carr}(f). \\ \quad \operatorname{is\text{-}fin\text{-}ext}(k, f, d, n, u) \Rightarrow  \\ \quad \quad \operatorname{ext\text{-}abs}(k, f, d, n, u, \operatorname{mul}(f)(i, a)) = (\operatorname{ext\text{-}abs}(k, f, d, n, u, i) \cdot \operatorname{ext\text{-}abs}(k, f, d, n, u, a)) \end{array}}}
\vnbprop{ext-abs-one}{\vnbstmt{\begin{array}{@{}l@{}} \forall k, f, d, n, u. \\ \quad \operatorname{is\text{-}fin\text{-}ext}(k, f, d, n, u) \Rightarrow (\operatorname{ext\text{-}abs}(k, f, d, n, u, \operatorname{one}(f)) = 1) \end{array}}}
\vnbprop{ext-abs-nonzero}{\vnbstmt{\begin{array}{@{}l@{}} \forall k, f, d, n, u.; \forall i \in \operatorname{carr}(f). \\ \quad (\operatorname{is\text{-}fin\text{-}ext}(k, f, d, n, u) \wedge \neg (i = \operatorname{zero}(f))) \Rightarrow  \\ \quad \quad \neg (\operatorname{ext\text{-}abs}(k, f, d, n, u, i) = 0) \end{array}}}
\vnbprop{matact-matadd-entry}{\vnbstmt{\begin{array}{@{}l@{}} \forall d, m, n, a, b, u.; \forall i \in \operatorname{interval}(1, m). \\ \quad \operatorname{is\text{-}module}(d) \wedge \\ \quad m \in \mathbb{N} \wedge \\ \quad n \in \mathbb{N} \wedge \\ \quad 1 \leq n \wedge \\ \quad a \in \operatorname{mat}(m, n, \operatorname{carr}(\operatorname{scal}(d))) \wedge \\ \quad b \in \operatorname{mat}(m, n, \operatorname{carr}(\operatorname{scal}(d))) \wedge \\ \quad u \in \operatorname{mat}(n, 1, \operatorname{vec}(d)) \Rightarrow \\ \quad \quad \operatorname{entry}(\operatorname{matact}(d, \operatorname{matadd}(\operatorname{scal}(d), a, b), u), i, 1) \\ \quad \quad =\; \operatorname{vadd}(d)(\operatorname{entry}(\operatorname{matact}(d, a, u), i, 1), \operatorname{entry}(\operatorname{matact}(d, b, u), i, 1)) \end{array}}}
\vnbprop{op-mat-add}{\vnbstmt{\begin{array}{@{}l@{}} \forall d, n, u, t, s, r. \\ \quad \operatorname{is\text{-}module}(d) \wedge \\ \quad n \in \mathbb{N} \wedge \\ \quad 1 \leq n \wedge \\ \quad u \in \operatorname{mat}(n, 1, \operatorname{vec}(d)) \wedge \\ \quad \operatorname{generates}(d, n, u) \wedge \\ \quad \operatorname{rel\text{-}free}(d, n, u) \wedge \\ \quad t \in (\operatorname{vec}(d) \to \operatorname{vec}(d)) \wedge \\ \quad s \in (\operatorname{vec}(d) \to \operatorname{vec}(d)) \wedge \\ \quad r \in (\operatorname{vec}(d) \to \operatorname{vec}(d)) \wedge \\ \quad \forall z \in \operatorname{vec}(d).\; (r(z) = \operatorname{vadd}(d)(t(z), s(z))) \Rightarrow \\ \quad \quad \operatorname{op\text{-}mat}(d, n, u, r) = \operatorname{matadd}(\operatorname{scal}(d), \operatorname{op\text{-}mat}(d, n, u, t), \operatorname{op\text{-}mat}(d, n, u, s)) \end{array}}}
\vnbprop{ext-mult-mat-add}{\vnbstmt{\begin{array}{@{}l@{}} \forall k, f, d, n, u.; \forall i, a \in \operatorname{carr}(f). \\ \quad \operatorname{is\text{-}fin\text{-}ext}(k, f, d, n, u) \Rightarrow  \\ \quad \quad \operatorname{mult\text{-}mat}(d, n, u, f, \operatorname{add}(f)(i, a)) \\ \quad \quad =\; \operatorname{matadd}(\operatorname{normed\text{-}field\text{-}as\text{-}commutative\text{-}ring}(k), \operatorname{mult\text{-}mat}(d, n, u, f, i), \operatorname{mult\text{-}mat}(d, n, u, f, a)) \end{array}}}
\vnbprop{matact-scal-ident-entry}{\vnbstmt{\begin{array}{@{}l@{}} \forall d.; \forall n \in \mathbb{N}, c \in \operatorname{carr}(\operatorname{scal}(d)), u \in \operatorname{mat}(n, 1, \operatorname{vec}(d)), i \in \operatorname{interval}(1, n). \\ \quad (\operatorname{is\text{-}module}(d) \wedge (1 \leq n)) \Rightarrow  \\ \quad \quad \operatorname{entry}(\operatorname{matact}(d, \operatorname{matscale}(\operatorname{scal}(d), c, \operatorname{identmat}(\operatorname{scal}(d), n)), u), i, 1) = \operatorname{act}(d)(c, \operatorname{entry}(u, i, 1)) \end{array}}}
\vnbprop{op-mat-scal}{\vnbstmt{\begin{array}{@{}l@{}} \forall d, n, u.; \forall c \in \operatorname{carr}(\operatorname{scal}(d)), r \in (\operatorname{vec}(d) \to \operatorname{vec}(d)). \\ \quad \operatorname{is\text{-}module}(d) \wedge \\ \quad n \in \mathbb{N} \wedge \\ \quad 1 \leq n \wedge \\ \quad u \in \operatorname{mat}(n, 1, \operatorname{vec}(d)) \wedge \\ \quad \operatorname{generates}(d, n, u) \wedge \\ \quad \operatorname{rel\text{-}free}(d, n, u) \wedge \\ \quad \forall z \in \operatorname{vec}(d).\; (r(z) = \operatorname{act}(d)(c, z)) \Rightarrow \\ \quad \quad \operatorname{op\text{-}mat}(d, n, u, r) = \operatorname{matscale}(\operatorname{scal}(d), c, \operatorname{identmat}(\operatorname{scal}(d), n)) \end{array}}}
\vnbprop{ext-mult-mat-scalar}{\vnbstmt{\begin{array}{@{}l@{}} \forall k, f, d, n, u.; \forall l \in \operatorname{carr}(k). \\ \quad \operatorname{is\text{-}fin\text{-}ext}(k, f, d, n, u) \Rightarrow  \\ \quad \quad \operatorname{mult\text{-}mat}(d, n, u, f, \operatorname{act}(d)(l, \operatorname{one}(f))) \\ \quad \quad =\; \operatorname{matscale}(\operatorname{normed\text{-}field\text{-}as\text{-}commutative\text{-}ring}(k), l, \operatorname{identmat}(\operatorname{normed\text{-}field\text{-}as\text{-}commutative\text{-}ring}(k), n)) \end{array}}}
\vnbprop{prod-ring-const}{\vnbstmt{\begin{array}{@{}l@{}} \forall s \in \mathbb{N}, r.; \forall c \in \operatorname{carr}(r), g. \\ \quad (\operatorname{is\text{-}commutative\text{-}ring}(r) \wedge \forall z \in \operatorname{interval}(1, s).\; (g(z) = c)) \Rightarrow  \\ \quad \quad \operatorname{prod\text{-}ring}(r, g, \operatorname{interval}(1, s)) = \operatorname{ring\text{-}power}(r, c, s) \end{array}}}
\vnbprop{det-scal-ident}{\vnbstmt{\begin{array}{@{}l@{}} \forall r.; \forall n \in \mathbb{N}, c \in \operatorname{carr}(r). \\ \quad \operatorname{is\text{-}commutative\text{-}ring}(r) \Rightarrow  \\ \quad \quad \operatorname{det}(r, n, \operatorname{matscale}(r, c, \operatorname{identmat}(r, n))) = \operatorname{ring\text{-}power}(r, c, n) \end{array}}}
\vnbprop{ext-abs-scalar}{\vnbstmt{\begin{array}{@{}l@{}} \forall k, f, d, n, u.; \forall l \in \operatorname{carr}(k). \\ \quad \operatorname{is\text{-}fin\text{-}ext}(k, f, d, n, u) \Rightarrow  \\ \quad \quad \operatorname{ext\text{-}abs}(k, f, d, n, u, \operatorname{act}(d)(l, \operatorname{one}(f))) = \operatorname{power}(\operatorname{fnrm}(k)(l), n) \end{array}}}
\vnbprop{ext-abs-zero}{\vnbstmt{\begin{array}{@{}l@{}} \forall k, f, d, n, u. \\ \quad \operatorname{is\text{-}fin\text{-}ext}(k, f, d, n, u) \Rightarrow (\operatorname{ext\text{-}abs}(k, f, d, n, u, \operatorname{zero}(f)) = 0) \end{array}}}
\vnbprop{ccr-4-2}{\vnbstmt{\begin{array}{@{}l@{}} \forall k, f, d, n, u. \\ \quad \operatorname{is\text{-}fin\text{-}ext}(k, f, d, n, u) \wedge \\ \quad \forall s, t \in \operatorname{carr}(f). \\ \quad \quad \operatorname{fnrm}(k)(\operatorname{det}(\operatorname{normed\text{-}field\text{-}as\text{-}commutative\text{-}ring}(k), n, \\ \quad \quad \quad \quad \operatorname{matadd}(\operatorname{normed\text{-}field\text{-}as\text{-}commutative\text{-}ring}(k), \operatorname{mult\text{-}mat}(d, n, u, f, s), \operatorname{mult\text{-}mat}(d, n, u, f, t)))) \\ \quad \quad \leq\; \operatorname{max}(\operatorname{fnrm}(k)(\operatorname{det}(\operatorname{normed\text{-}field\text{-}as\text{-}commutative\text{-}ring}(k), n, \operatorname{mult\text{-}mat}(d, n, u, f, s))), \\ \quad \quad \quad \operatorname{fnrm}(k)(\operatorname{det}(\operatorname{normed\text{-}field\text{-}as\text{-}commutative\text{-}ring}(k), n, \operatorname{mult\text{-}mat}(d, n, u, f, t)))) \Rightarrow \\ \quad \quad \operatorname{is\text{-}nv\text{-}valuation}(f, (\lambda v \in \operatorname{carr}(f).\; \operatorname{ext\text{-}abs}(k, f, d, n, u, v))) \wedge \\ \quad \quad \forall l \in \operatorname{carr}(k). \\ \quad \quad \quad \operatorname{ext\text{-}abs}(k, f, d, n, u, \operatorname{act}(d)(l, \operatorname{one}(f))) = \operatorname{power}(\operatorname{fnrm}(k)(l), n) \end{array}}}
\vnbprop{ext-abs-root-scalar}{\vnbstmt{\begin{array}{@{}l@{}} \forall k, f, d, n, u.; \forall l \in \operatorname{carr}(k). \\ \quad \operatorname{is\text{-}fin\text{-}ext}(k, f, d, n, u) \Rightarrow  \\ \quad \quad \operatorname{ext\text{-}abs\text{-}root}(k, f, d, n, u, \operatorname{act}(d)(l, \operatorname{one}(f))) = \operatorname{fnrm}(k)(l) \end{array}}}

\vnbbookfile{theorem-library/ccr-ch3c-ext}{theorem-library/ccr-ch3c-ext}
\noindent{\small ccr-ch3c-ext.scm -- wave of 2026-11-09/10 (agent CCR-1109-B3), wired 2026-11-10 from scratchpad/ccr1109b3/ccr-ch3c-ext.scm. Library: analysis.}\par
\vnbprop{ccr-4-2-full}{\vnbstmt{\begin{array}{@{}l@{}} \forall k, f, d, n, u.; \forall h \in \operatorname{carr}(f). \\ \quad \operatorname{is\text{-}fin\text{-}ext}(k, f, d, n, u) \wedge \\ \quad \operatorname{is\text{-}local\text{-}nv\text{-}field}(k) \wedge \\ \quad \operatorname{is\text{-}irreducible\text{-}mat}(k, n, \operatorname{mult\text{-}mat}(d, n, u, f, h)) \Rightarrow \\ \quad \quad \operatorname{is\text{-}nv\text{-}valuation}(f, (\lambda v \in \operatorname{carr}(f).\; \operatorname{ext\text{-}abs}(k, f, d, n, u, v))) \wedge \\ \quad \quad \forall l \in \operatorname{carr}(k). \\ \quad \quad \quad \operatorname{ext\text{-}abs}(k, f, d, n, u, \operatorname{act}(d)(l, \operatorname{one}(f))) = \operatorname{power}(\operatorname{fnrm}(k)(l), n) \end{array}}}

\vnbbookfile{theorem-library/ccr-ch3d-ultra}{theorem-library/ccr-ch3d-ultra}
\noindent{\small ccr-ch3d-ultra.scm -- wave of 2026-11-09/10 (agent CCR-1109-B4), wired 2026-11-10 from scratchpad/ccr1109b4/ccr-ch3d-ultra.scm. Library: analysis.}\par
\vnbprop{ccr-3-6-ultrametric}{\vnbstmt{\begin{array}{@{}l@{}} \forall k.; \forall n \in \mathbb{N}, m.; \forall a, b \in \operatorname{mat}(n, n, \operatorname{carr}(k)). \\ \quad \operatorname{is\text{-}local\text{-}nv\text{-}field}(k) \wedge \\ \quad \operatorname{is\text{-}irreducible\text{-}mat}(k, n, m) \wedge \\ \quad \operatorname{matmul}(\operatorname{normed\text{-}field\text{-}as\text{-}commutative\text{-}ring}(k), a, m) \\ \quad =\; \operatorname{matmul}(\operatorname{normed\text{-}field\text{-}as\text{-}commutative\text{-}ring}(k), m, a) \wedge \\ \quad \operatorname{matmul}(\operatorname{normed\text{-}field\text{-}as\text{-}commutative\text{-}ring}(k), b, m) \\ \quad =\; \operatorname{matmul}(\operatorname{normed\text{-}field\text{-}as\text{-}commutative\text{-}ring}(k), m, b) \Rightarrow \\ \quad \quad \operatorname{fnrm}(k)(\operatorname{det}(\operatorname{normed\text{-}field\text{-}as\text{-}commutative\text{-}ring}(k), n, \operatorname{matadd}(\operatorname{normed\text{-}field\text{-}as\text{-}commutative\text{-}ring}(k), a, b))) \\ \quad \quad \leq\; \operatorname{max}(\operatorname{fnrm}(k)(\operatorname{det}(\operatorname{normed\text{-}field\text{-}as\text{-}commutative\text{-}ring}(k), n, a)), \\ \quad \quad \quad \operatorname{fnrm}(k)(\operatorname{det}(\operatorname{normed\text{-}field\text{-}as\text{-}commutative\text{-}ring}(k), n, b))) \end{array}}}

\vnbbookfile{theorem-library/ccr-ch3d-fam}{theorem-library/ccr-ch3d-fam}
\noindent{\small ccr-ch3d-fam.scm -- wave of 2026-11-09/10 (agent CCR-1109-B4), wired 2026-11-10 from scratchpad/ccr1109b4/ccr-ch3d-fam.scm. Library: analysis.}\par
\vnbprop{cf4-step}{\vnbstmt{\begin{array}{@{}l@{}} \forall k.; \forall n \in \mathbb{N}, x, m.; \forall a \in \mathbb{N}. \\ \quad \operatorname{is\text{-}nv\text{-}field}(k) \wedge \\ \quad x \in \operatorname{mat}(n, n, \operatorname{carr}(k)) \wedge \\ \quad \forall y \in m.\; (y \in \operatorname{mat}(n, n, \operatorname{carr}(k))) \wedge \\ \quad \forall y \in m. \\ \quad \quad \operatorname{matmul}(\operatorname{normed\text{-}field\text{-}as\text{-}commutative\text{-}ring}(k), x, y) \\ \quad \quad =\; \operatorname{matmul}(\operatorname{normed\text{-}field\text{-}as\text{-}commutative\text{-}ring}(k), y, x) \wedge \\ \quad \operatorname{fnrm}(k)(\operatorname{det}(\operatorname{normed\text{-}field\text{-}as\text{-}commutative\text{-}ring}(k), n, x)) \leq 1 \wedge \\ \quad \neg \operatorname{is\text{-}power\text{-}bounded}(k, n, x) \Rightarrow \\ \quad \quad \exists b \in \operatorname{mat}(n, n, \operatorname{carr}(k)). \\ \quad \quad \quad \operatorname{is\text{-}mat\text{-}bound}(k, n, n, b, 1) \wedge \\ \quad \quad \quad \forall y \in m. \\ \quad \quad \quad \quad \operatorname{matmul}(\operatorname{normed\text{-}field\text{-}as\text{-}commutative\text{-}ring}(k), b, y) \\ \quad \quad \quad \quad =\; \operatorname{matmul}(\operatorname{normed\text{-}field\text{-}as\text{-}commutative\text{-}ring}(k), y, b) \wedge \\ \quad \quad \quad (\operatorname{fnrm}(k)(\operatorname{det}(\operatorname{normed\text{-}field\text{-}as\text{-}commutative\text{-}ring}(k), n, b)) \cdot (a + 1)) \leq 1 \wedge \\ \quad \quad \quad \exists i, j \in \operatorname{interval}(1, n). \\ \quad \quad \quad \quad \operatorname{fnrm}(k)(\operatorname{entry}(b, i, j)) = 1 \end{array}}}
\vnbprop{ccr-3-4-fam}{\vnbstmt{\begin{array}{@{}l@{}} \forall k.; \forall n \in \mathbb{N}, x, m. \\ \quad \operatorname{is\text{-}local\text{-}nv\text{-}field}(k) \wedge \\ \quad \forall y \in m.\; (y \in \operatorname{mat}(n, n, \operatorname{carr}(k))) \wedge \\ \quad \forall b \in \operatorname{mat}(n, n, \operatorname{carr}(k)). \\ \quad \quad \forall y \in m. \\ \quad \quad \quad \operatorname{matmul}(\operatorname{normed\text{-}field\text{-}as\text{-}commutative\text{-}ring}(k), b, y) \\ \quad \quad \quad =\; \operatorname{matmul}(\operatorname{normed\text{-}field\text{-}as\text{-}commutative\text{-}ring}(k), y, b) \wedge \\ \quad \quad \operatorname{det}(\operatorname{normed\text{-}field\text{-}as\text{-}commutative\text{-}ring}(k), n, b) = \operatorname{zero}(k) \Rightarrow \\ \quad \quad \quad b = \operatorname{zeromat}(\operatorname{normed\text{-}field\text{-}as\text{-}commutative\text{-}ring}(k), n, n) \wedge \\ \quad x \in \operatorname{mat}(n, n, \operatorname{carr}(k)) \wedge \\ \quad \forall y \in m. \\ \quad \quad \operatorname{matmul}(\operatorname{normed\text{-}field\text{-}as\text{-}commutative\text{-}ring}(k), x, y) \\ \quad \quad =\; \operatorname{matmul}(\operatorname{normed\text{-}field\text{-}as\text{-}commutative\text{-}ring}(k), y, x) \wedge \\ \quad \operatorname{fnrm}(k)(\operatorname{det}(\operatorname{normed\text{-}field\text{-}as\text{-}commutative\text{-}ring}(k), n, x)) \leq 1 \Rightarrow \\ \quad \quad \operatorname{is\text{-}power\text{-}bounded}(k, n, x) \end{array}}}
\vnbprop{ccr-3-6-fam-inv}{\vnbstmt{\begin{array}{@{}l@{}} \forall k.; \forall n \in \mathbb{N}, m.; \forall a, b, c \in \operatorname{mat}(n, n, \operatorname{carr}(k)). \\ \quad \operatorname{is\text{-}local\text{-}nv\text{-}field}(k) \wedge \\ \quad \forall y \in m.\; (y \in \operatorname{mat}(n, n, \operatorname{carr}(k))) \wedge \\ \quad \forall \mathit{cfb} \in \operatorname{mat}(n, n, \operatorname{carr}(k)). \\ \quad \quad \forall y \in m. \\ \quad \quad \quad \operatorname{matmul}(\operatorname{normed\text{-}field\text{-}as\text{-}commutative\text{-}ring}(k), \mathit{cfb}, y) \\ \quad \quad \quad =\; \operatorname{matmul}(\operatorname{normed\text{-}field\text{-}as\text{-}commutative\text{-}ring}(k), y, \mathit{cfb}) \wedge \\ \quad \quad \operatorname{det}(\operatorname{normed\text{-}field\text{-}as\text{-}commutative\text{-}ring}(k), n, \mathit{cfb}) = \operatorname{zero}(k) \Rightarrow \\ \quad \quad \quad \mathit{cfb} = \operatorname{zeromat}(\operatorname{normed\text{-}field\text{-}as\text{-}commutative\text{-}ring}(k), n, n) \wedge \\ \quad \forall y \in m. \\ \quad \quad \operatorname{matmul}(\operatorname{normed\text{-}field\text{-}as\text{-}commutative\text{-}ring}(k), a, y) \\ \quad \quad =\; \operatorname{matmul}(\operatorname{normed\text{-}field\text{-}as\text{-}commutative\text{-}ring}(k), y, a) \wedge \\ \quad \forall y \in m. \\ \quad \quad \operatorname{matmul}(\operatorname{normed\text{-}field\text{-}as\text{-}commutative\text{-}ring}(k), b, y) \\ \quad \quad =\; \operatorname{matmul}(\operatorname{normed\text{-}field\text{-}as\text{-}commutative\text{-}ring}(k), y, b) \wedge \\ \quad \operatorname{matmul}(\operatorname{normed\text{-}field\text{-}as\text{-}commutative\text{-}ring}(k), c, a) = \operatorname{identmat}(\operatorname{normed\text{-}field\text{-}as\text{-}commutative\text{-}ring}(k), n) \wedge \\ \quad \operatorname{matmul}(\operatorname{normed\text{-}field\text{-}as\text{-}commutative\text{-}ring}(k), a, c) = \operatorname{identmat}(\operatorname{normed\text{-}field\text{-}as\text{-}commutative\text{-}ring}(k), n) \wedge \\ \quad \operatorname{fnrm}(k)(\operatorname{det}(\operatorname{normed\text{-}field\text{-}as\text{-}commutative\text{-}ring}(k), n, b)) \\ \quad \leq\; \operatorname{fnrm}(k)(\operatorname{det}(\operatorname{normed\text{-}field\text{-}as\text{-}commutative\text{-}ring}(k), n, a)) \Rightarrow \\ \quad \quad \operatorname{fnrm}(k)(\operatorname{det}(\operatorname{normed\text{-}field\text{-}as\text{-}commutative\text{-}ring}(k), n, \operatorname{matadd}(\operatorname{normed\text{-}field\text{-}as\text{-}commutative\text{-}ring}(k), a, b))) \\ \quad \quad \leq\; \operatorname{fnrm}(k)(\operatorname{det}(\operatorname{normed\text{-}field\text{-}as\text{-}commutative\text{-}ring}(k), n, a)) \end{array}}}
\vnbprop{ccr-3-6-fam}{\vnbstmt{\begin{array}{@{}l@{}} \forall k.; \forall n \in \mathbb{N}, m.; \forall a, b \in \operatorname{mat}(n, n, \operatorname{carr}(k)). \\ \quad \operatorname{is\text{-}local\text{-}nv\text{-}field}(k) \wedge \\ \quad \forall y \in m.\; (y \in \operatorname{mat}(n, n, \operatorname{carr}(k))) \wedge \\ \quad \forall \mathit{cfb} \in \operatorname{mat}(n, n, \operatorname{carr}(k)). \\ \quad \quad \forall y \in m. \\ \quad \quad \quad \operatorname{matmul}(\operatorname{normed\text{-}field\text{-}as\text{-}commutative\text{-}ring}(k), \mathit{cfb}, y) \\ \quad \quad \quad =\; \operatorname{matmul}(\operatorname{normed\text{-}field\text{-}as\text{-}commutative\text{-}ring}(k), y, \mathit{cfb}) \wedge \\ \quad \quad \operatorname{det}(\operatorname{normed\text{-}field\text{-}as\text{-}commutative\text{-}ring}(k), n, \mathit{cfb}) = \operatorname{zero}(k) \Rightarrow \\ \quad \quad \quad \mathit{cfb} = \operatorname{zeromat}(\operatorname{normed\text{-}field\text{-}as\text{-}commutative\text{-}ring}(k), n, n) \wedge \\ \quad \forall y \in m. \\ \quad \quad \operatorname{matmul}(\operatorname{normed\text{-}field\text{-}as\text{-}commutative\text{-}ring}(k), a, y) \\ \quad \quad =\; \operatorname{matmul}(\operatorname{normed\text{-}field\text{-}as\text{-}commutative\text{-}ring}(k), y, a) \wedge \\ \quad \forall y \in m. \\ \quad \quad \operatorname{matmul}(\operatorname{normed\text{-}field\text{-}as\text{-}commutative\text{-}ring}(k), b, y) \\ \quad \quad =\; \operatorname{matmul}(\operatorname{normed\text{-}field\text{-}as\text{-}commutative\text{-}ring}(k), y, b) \Rightarrow \\ \quad \quad \operatorname{fnrm}(k)(\operatorname{det}(\operatorname{normed\text{-}field\text{-}as\text{-}commutative\text{-}ring}(k), n, \operatorname{matadd}(\operatorname{normed\text{-}field\text{-}as\text{-}commutative\text{-}ring}(k), a, b))) \\ \quad \quad \leq\; \operatorname{max}(\operatorname{fnrm}(k)(\operatorname{det}(\operatorname{normed\text{-}field\text{-}as\text{-}commutative\text{-}ring}(k), n, a)), \\ \quad \quad \quad \operatorname{fnrm}(k)(\operatorname{det}(\operatorname{normed\text{-}field\text{-}as\text{-}commutative\text{-}ring}(k), n, b))) \end{array}}}

\vnbbookfile{theorem-library/ccr-ch3d-ext}{theorem-library/ccr-ch3d-ext}
\noindent{\small ccr-ch3d-ext.scm -- wave of 2026-11-09/10 (agent CCR-1109-B4), wired 2026-11-10 from scratchpad/ccr1109b4/ccr-ch3d-ext.scm. Library: analysis.}\par
\vnbprop{cg5-row-nonzero}{\vnbstmt{\begin{array}{@{}l@{}} \forall k, f, d, n, u.; \forall e \in \operatorname{mat}(n, 1, \operatorname{carr}(k)). \\ \quad (\operatorname{is\text{-}fin\text{-}ext}(k, f, d, n, u) \wedge \neg (e = \operatorname{zeromat}(\operatorname{normed\text{-}field\text{-}as\text{-}commutative\text{-}ring}(k), n, 1))) \Rightarrow  \\ \quad \quad \neg (\operatorname{entry}(\operatorname{matact}(d, \operatorname{transpose}(e, n, 1), u), 1, 1) = \operatorname{zero}(f)) \end{array}}}
\vnbprop{cg5-row-span}{\vnbstmt{\begin{array}{@{}l@{}} \forall k, f, d, n, u.; \forall c, a \in \operatorname{mat}(1, n, \operatorname{carr}(k)). \\ \quad (\operatorname{is\text{-}fin\text{-}ext}(k, f, d, n, u) \wedge \neg (\operatorname{entry}(\operatorname{matact}(d, c, u), 1, 1) = \operatorname{zero}(f))) \Rightarrow  \\ \quad \quad \exists w \in \operatorname{carr}(f). \\ \quad \quad \quad a = \operatorname{matmul}(\operatorname{normed\text{-}field\text{-}as\text{-}commutative\text{-}ring}(k), c, \operatorname{mult\text{-}mat}(d, n, u, f, w)) \end{array}}}
\vnbprop{ext-mult-mat-rigid}{\vnbstmt{\begin{array}{@{}l@{}} \forall k, f, d, n, u.; \forall b \in \operatorname{mat}(n, n, \operatorname{carr}(k)). \\ \quad \operatorname{is\text{-}fin\text{-}ext}(k, f, d, n, u) \wedge \\ \quad \forall x \in \operatorname{carr}(f). \\ \quad \quad \operatorname{matmul}(\operatorname{normed\text{-}field\text{-}as\text{-}commutative\text{-}ring}(k), b, \operatorname{mult\text{-}mat}(d, n, u, f, x)) \\ \quad \quad =\; \operatorname{matmul}(\operatorname{normed\text{-}field\text{-}as\text{-}commutative\text{-}ring}(k), \operatorname{mult\text{-}mat}(d, n, u, f, x), b) \wedge \\ \quad \operatorname{det}(\operatorname{normed\text{-}field\text{-}as\text{-}commutative\text{-}ring}(k), n, b) = \operatorname{zero}(k) \Rightarrow \\ \quad \quad b = \operatorname{zeromat}(\operatorname{normed\text{-}field\text{-}as\text{-}commutative\text{-}ring}(k), n, n) \end{array}}}
\vnbprop{ccr-4-2-local}{\vnbstmt{\begin{array}{@{}l@{}} \forall k, f, d, n, u. \\ \quad (\operatorname{is\text{-}fin\text{-}ext}(k, f, d, n, u) \wedge \operatorname{is\text{-}local\text{-}nv\text{-}field}(k)) \Rightarrow  \\ \quad \quad \operatorname{is\text{-}nv\text{-}valuation}(f, (\lambda v \in \operatorname{carr}(f).\; \operatorname{ext\text{-}abs}(k, f, d, n, u, v))) \wedge \\ \quad \quad \forall l \in \operatorname{carr}(k). \\ \quad \quad \quad \operatorname{ext\text{-}abs}(k, f, d, n, u, \operatorname{act}(d)(l, \operatorname{one}(f))) = \operatorname{power}(\operatorname{fnrm}(k)(l), n) \end{array}}}

\vnbbookfile{theorem-library/ccr-ch3e-norms}{theorem-library/ccr-ch3e-norms}
\noindent{\small ccr-ch3e-norms.scm -- wave of 2026-11-09/10 (agent CCR-1109-B5), wired 2026-11-10 from scratchpad/ccr1109b5/ccr-ch3e-norms.scm. Library: analysis.}\par
\vnbprop{ccr-3-8-upper}{\vnbstmt{\begin{array}{@{}l@{}} \forall k, d, a.; \forall n \in \mathbb{N}, u \in \operatorname{mat}(n, 1, \operatorname{vec}(d)). \\ \quad (\operatorname{is\text{-}normed\text{-}field}(k) \wedge \operatorname{is\text{-}norm\text{-}on}(k, d, a)) \Rightarrow  \\ \quad \quad \exists b \in \mathbb{R}. \\ \quad \quad \quad 0 \leq b \wedge \\ \quad \quad \quad \forall c \in \operatorname{mat}(1, n, \operatorname{carr}(k)), m \in \mathbb{R}. \\ \quad \quad \quad \quad ((0 \leq m) \wedge \operatorname{is\text{-}mat\text{-}bound}(k, 1, n, c, m)) \Rightarrow (a(\operatorname{lincomb}(d, n, c, u)) \leq (b \cdot m)) \end{array}}}
\vnbprop{cn5-row-argmax}{\vnbstmt{\begin{array}{@{}l@{}} \forall k.; \forall n \in \mathbb{N}, a \in \operatorname{mat}(1, n, \operatorname{carr}(k)). \\ \quad (\operatorname{is\text{-}normed\text{-}field}(k) \wedge (1 \leq n)) \Rightarrow  \\ \quad \quad \exists j_{0} \in \operatorname{interval}(1, n). \forall j \in \operatorname{interval}(1, n). \\ \quad \quad \quad \operatorname{fnrm}(k)(\operatorname{entry}(a, 1, j)) \leq \operatorname{fnrm}(k)(\operatorname{entry}(a, 1, j_{0})) \end{array}}}
\vnbprop{cn5-step}{\vnbstmt{\begin{array}{@{}l@{}} \forall k, d, a.; \forall n \in \mathbb{N}, u \in \operatorname{mat}(n, 1, \operatorname{vec}(d)), b \in \mathbb{N}. \\ \quad \operatorname{is\text{-}nv\text{-}field}(k) \wedge \\ \quad \operatorname{is\text{-}norm\text{-}on}(k, d, a) \wedge \\ \quad \forall c \in \mathbb{R}. \\ \quad \quad (0 \leq c) \Rightarrow  \\ \quad \quad \quad \exists f \in \operatorname{mat}(1, n, \operatorname{carr}(k)), j \in \operatorname{interval}(1, n). \\ \quad \quad \quad \quad \neg (\operatorname{fnrm}(k)(\operatorname{entry}(f, 1, j)) \leq (c \cdot a(\operatorname{lincomb}(d, n, f, u)))) \Rightarrow \\ \quad \quad \exists y \in \operatorname{mat}(1, n, \operatorname{carr}(k)). \\ \quad \quad \quad \operatorname{is\text{-}mat\text{-}bound}(k, 1, n, y, 1) \wedge \\ \quad \quad \quad (a(\operatorname{lincomb}(d, n, y, u)) \cdot (b + 1)) \leq 1 \wedge \\ \quad \quad \quad \exists \mathit{cn5j} \in \operatorname{interval}(1, n).\; (\operatorname{fnrm}(k)(\operatorname{entry}(y, 1, \mathit{cn5j})) = 1) \end{array}}}
\vnbprop{ccr-3-8-lower}{\vnbstmt{\begin{array}{@{}l@{}} \forall k, d, a.; \forall n \in \mathbb{N}, u \in \operatorname{mat}(n, 1, \operatorname{vec}(d)). \\ \quad (\operatorname{is\text{-}local\text{-}nv\text{-}field}(k) \wedge (\operatorname{is\text{-}norm\text{-}on}(k, d, a) \wedge \operatorname{rel\text{-}free}(d, n, u))) \Rightarrow  \\ \quad \quad \exists c \in \mathbb{R}. \\ \quad \quad \quad 0 \leq c \wedge \\ \quad \quad \quad \forall b \in \operatorname{mat}(1, n, \operatorname{carr}(k)), j \in \operatorname{interval}(1, n). \\ \quad \quad \quad \quad \operatorname{fnrm}(k)(\operatorname{entry}(b, 1, j)) \leq (c \cdot a(\operatorname{lincomb}(d, n, b, u))) \end{array}}}
\vnbprop{ccr-3-8-dominate}{\vnbstmt{\begin{array}{@{}l@{}} \forall k, d, u_{1}, u_{2}.; \forall n \in \mathbb{N}, u \in \operatorname{mat}(n, 1, \operatorname{vec}(d)). \\ \quad \operatorname{is\text{-}local\text{-}nv\text{-}field}(k) \wedge \\ \quad \operatorname{is\text{-}norm\text{-}on}(k, d, u_{1}) \wedge \\ \quad \operatorname{is\text{-}norm\text{-}on}(k, d, u_{2}) \wedge \\ \quad \operatorname{generates}(d, n, u) \wedge \\ \quad \operatorname{rel\text{-}free}(d, n, u) \Rightarrow \\ \quad \quad \exists a \in \mathbb{R}.\; ((0 < a) \wedge \forall x \in \operatorname{vec}(d).\; (\operatorname{u2}(x) \leq (a \cdot \operatorname{u1}(x)))) \end{array}}}
\vnbprop{ccr-3-8-basis}{\vnbstmt{\begin{array}{@{}l@{}} \forall k, d, u_{1}, u_{2}.; \forall n \in \mathbb{N}, u \in \operatorname{mat}(n, 1, \operatorname{vec}(d)). \\ \quad \operatorname{is\text{-}local\text{-}nv\text{-}field}(k) \wedge \\ \quad \operatorname{is\text{-}norm\text{-}on}(k, d, u_{1}) \wedge \\ \quad \operatorname{is\text{-}norm\text{-}on}(k, d, u_{2}) \wedge \\ \quad \operatorname{generates}(d, n, u) \wedge \\ \quad \operatorname{rel\text{-}free}(d, n, u) \Rightarrow \\ \quad \quad \exists l \in \mathbb{R}. \\ \quad \quad \quad 0 < l \wedge \\ \quad \quad \quad \exists a \in \mathbb{R}. \\ \quad \quad \quad \quad (0 < a) \wedge \forall x \in \operatorname{vec}(d).\; (((l \cdot \operatorname{u1}(x)) \leq \operatorname{u2}(x)) \wedge (\operatorname{u2}(x) \leq (a \cdot \operatorname{u1}(x)))) \end{array}}}
\vnbprop{ccr-3-8}{\vnbstmt{\begin{array}{@{}l@{}} \forall k, d, u_{1}, u_{2}. \\ \quad \operatorname{is\text{-}local\text{-}nv\text{-}field}(k) \wedge \\ \quad \operatorname{is\text{-}norm\text{-}on}(k, d, u_{1}) \wedge \\ \quad \operatorname{is\text{-}norm\text{-}on}(k, d, u_{2}) \wedge \\ \quad \exists n \in \mathbb{N}, u \in \operatorname{mat}(n, 1, \operatorname{vec}(d)). \\ \quad \quad \operatorname{generates}(d, n, u) \wedge \operatorname{rel\text{-}free}(d, n, u) \Rightarrow \\ \quad \quad \exists l \in \mathbb{R}. \\ \quad \quad \quad 0 < l \wedge \\ \quad \quad \quad \exists a \in \mathbb{R}. \\ \quad \quad \quad \quad (0 < a) \wedge \forall x \in \operatorname{vec}(d).\; (((l \cdot \operatorname{u1}(x)) \leq \operatorname{u2}(x)) \wedge (\operatorname{u2}(x) \leq (a \cdot \operatorname{u1}(x)))) \end{array}}}

\vnbbookfile{theorem-library/ccr4x-eisen}{theorem-library/ccr4x-eisen}
\noindent{\small ccr4x-eisen.scm -- wave of 2026-11-09/10 (agent CCR-1109-C), wired 2026-11-10 from scratchpad/ccr1109c/ccr4x-eisen.scm. Library: analysis.}\par
\vnbprop{nv-coef-le}{\vnbstmt{\begin{array}{@{}l@{}} \forall k.; \forall p, q \in \operatorname{finsupp}(\operatorname{normed\text{-}field\text{-}as\text{-}commutative\text{-}ring}(k), \operatorname{nn\text{-}add\text{-}monoid}), a \in \mathbb{N}, c \in \mathbb{R}. \\ \quad \operatorname{is\text{-}nv\text{-}field}(k) \wedge \\ \quad 0 \leq c \wedge \\ \quad \forall i, j \in \mathbb{N}. \\ \quad \quad ((i + j) = a) \Rightarrow ((\operatorname{fnrm}(k)(p(i)) \cdot \operatorname{fnrm}(k)(q(j))) \leq c) \Rightarrow \\ \quad \quad \operatorname{fnrm}(k)(\operatorname{monalg\text{-}mul}(\operatorname{normed\text{-}field\text{-}as\text{-}commutative\text{-}ring}(k), \operatorname{nn\text{-}add\text{-}monoid}, p, q)(a)) \leq c \end{array}}}
\vnbprop{nv-coef-lt}{\vnbstmt{\begin{array}{@{}l@{}} \forall k.; \forall p, q \in \operatorname{finsupp}(\operatorname{normed\text{-}field\text{-}as\text{-}commutative\text{-}ring}(k), \operatorname{nn\text{-}add\text{-}monoid}), a \in \mathbb{N}, c \in \mathbb{R}. \\ \quad \operatorname{is\text{-}nv\text{-}field}(k) \wedge \\ \quad 0 < c \wedge \\ \quad \forall i, j \in \mathbb{N}. \\ \quad \quad ((i + j) = a) \Rightarrow ((\operatorname{fnrm}(k)(p(i)) \cdot \operatorname{fnrm}(k)(q(j))) < c) \Rightarrow \\ \quad \quad \operatorname{fnrm}(k)(\operatorname{monalg\text{-}mul}(\operatorname{normed\text{-}field\text{-}as\text{-}commutative\text{-}ring}(k), \operatorname{nn\text{-}add\text{-}monoid}, p, q)(a)) < c \end{array}}}
\vnbprop{nv-coef-dominant}{\vnbstmt{\begin{array}{@{}l@{}} \forall k.; \forall p, q \in \operatorname{finsupp}(\operatorname{normed\text{-}field\text{-}as\text{-}commutative\text{-}ring}(k), \operatorname{nn\text{-}add\text{-}monoid}), m, n \in \mathbb{N}, c \in \mathbb{R}. \\ \quad \operatorname{is\text{-}nv\text{-}field}(k) \wedge \\ \quad 0 < c \wedge \\ \quad (\operatorname{fnrm}(k)(p(m)) \cdot \operatorname{fnrm}(k)(q(n))) = c \wedge \\ \quad \forall i, j \in \mathbb{N}. \\ \quad \quad (((i + j) = (m + n)) \wedge \neg (i = m)) \Rightarrow ((\operatorname{fnrm}(k)(p(i)) \cdot \operatorname{fnrm}(k)(q(j))) < c) \Rightarrow \\ \quad \quad \operatorname{fnrm}(k)(\operatorname{monalg\text{-}mul}(\operatorname{normed\text{-}field\text{-}as\text{-}commutative\text{-}ring}(k), \operatorname{nn\text{-}add\text{-}monoid}, p, q)((m + n))) = c \end{array}}}
\vnbprop{nv-below-scaled}{\vnbstmt{\begin{array}{@{}l@{}} \forall k.; \forall w, a, x \in \operatorname{carr}(k). \\ \quad \operatorname{is\text{-}nv\text{-}field}(k) \wedge \\ \quad \operatorname{nv\text{-}maximal}(k, \operatorname{nv\text{-}integers}(k)) = \operatorname{nv\text{-}principal}(k, \operatorname{nv\text{-}integers}(k), w) \wedge \\ \quad \operatorname{fnrm}(k)(x) < \operatorname{fnrm}(k)(a) \Rightarrow \\ \quad \quad \operatorname{fnrm}(k)(x) \leq (\operatorname{fnrm}(k)(a) \cdot \operatorname{fnrm}(k)(w)) \end{array}}}
\vnbprop{nv-eisenstein-half}{\vnbstmt{\begin{array}{@{}l@{}} \forall k.; \forall w \in \operatorname{carr}(k), q \in \operatorname{finsupp}(\operatorname{normed\text{-}field\text{-}as\text{-}commutative\text{-}ring}(k), \operatorname{nn\text{-}add\text{-}monoid}), n \in \mathbb{N}, r, s \in \operatorname{finsupp}(\operatorname{normed\text{-}field\text{-}as\text{-}commutative\text{-}ring}(k), \operatorname{nn\text{-}add\text{-}monoid}), a, u \in \operatorname{carr}(k). \\ \quad \operatorname{is\text{-}nv\text{-}field}(k) \wedge \\ \quad \operatorname{nv\text{-}maximal}(k, \operatorname{nv\text{-}integers}(k)) = \operatorname{nv\text{-}principal}(k, \operatorname{nv\text{-}integers}(k), w) \wedge \\ \quad 1 \leq n \wedge \\ \quad \operatorname{deg\text{-}bound}(\operatorname{normed\text{-}field\text{-}as\text{-}commutative\text{-}ring}(k), q, n) \wedge \\ \quad \operatorname{fnrm}(k)(q(n)) = 1 \wedge \\ \quad \forall i \in \mathbb{N}.\; ((i < n) \Rightarrow (\operatorname{fnrm}(k)(q(i)) < 1)) \wedge \\ \quad \neg (\operatorname{fnrm}(k)(q(0)) \leq (\operatorname{fnrm}(k)(w) \cdot \operatorname{fnrm}(k)(w))) \wedge \\ \quad q = \operatorname{monalg\text{-}mul}(\operatorname{normed\text{-}field\text{-}as\text{-}commutative\text{-}ring}(k), \operatorname{nn\text{-}add\text{-}monoid}, r, s) \wedge \\ \quad (\operatorname{fnrm}(k)(a) \cdot \operatorname{fnrm}(k)(u)) = 1 \wedge \\ \quad \forall i \in \mathbb{N}.\; (\operatorname{fnrm}(k)(r(i)) \leq \operatorname{fnrm}(k)(a)) \wedge \\ \quad \forall i \in \mathbb{N}.\; (\operatorname{fnrm}(k)(s(i)) \leq \operatorname{fnrm}(k)(u)) \wedge \\ \quad \operatorname{fnrm}(k)(r(0)) < \operatorname{fnrm}(k)(a) \Rightarrow \\ \quad \quad \operatorname{deg\text{-}bound}(\operatorname{normed\text{-}field\text{-}as\text{-}commutative\text{-}ring}(k), s, 0) \end{array}}}
\vnbprop{nv-eisenstein-scaled}{\vnbstmt{\begin{array}{@{}l@{}} \forall k.; \forall w \in \operatorname{carr}(k), q \in \operatorname{finsupp}(\operatorname{normed\text{-}field\text{-}as\text{-}commutative\text{-}ring}(k), \operatorname{nn\text{-}add\text{-}monoid}), n \in \mathbb{N}, r, s \in \operatorname{finsupp}(\operatorname{normed\text{-}field\text{-}as\text{-}commutative\text{-}ring}(k), \operatorname{nn\text{-}add\text{-}monoid}), a, u \in \operatorname{carr}(k). \\ \quad \operatorname{is\text{-}nv\text{-}field}(k) \wedge \\ \quad \operatorname{nv\text{-}maximal}(k, \operatorname{nv\text{-}integers}(k)) = \operatorname{nv\text{-}principal}(k, \operatorname{nv\text{-}integers}(k), w) \wedge \\ \quad 1 \leq n \wedge \\ \quad \operatorname{deg\text{-}bound}(\operatorname{normed\text{-}field\text{-}as\text{-}commutative\text{-}ring}(k), q, n) \wedge \\ \quad \operatorname{fnrm}(k)(q(n)) = 1 \wedge \\ \quad \forall i \in \mathbb{N}.\; ((i < n) \Rightarrow (\operatorname{fnrm}(k)(q(i)) < 1)) \wedge \\ \quad \neg (\operatorname{fnrm}(k)(q(0)) \leq (\operatorname{fnrm}(k)(w) \cdot \operatorname{fnrm}(k)(w))) \wedge \\ \quad q = \operatorname{monalg\text{-}mul}(\operatorname{normed\text{-}field\text{-}as\text{-}commutative\text{-}ring}(k), \operatorname{nn\text{-}add\text{-}monoid}, r, s) \wedge \\ \quad (\operatorname{fnrm}(k)(a) \cdot \operatorname{fnrm}(k)(u)) = 1 \wedge \\ \quad \forall i \in \mathbb{N}.\; (\operatorname{fnrm}(k)(r(i)) \leq \operatorname{fnrm}(k)(a)) \wedge \\ \quad \forall i \in \mathbb{N}.\; (\operatorname{fnrm}(k)(s(i)) \leq \operatorname{fnrm}(k)(u)) \Rightarrow \\ \quad \quad \operatorname{deg\text{-}bound}(\operatorname{normed\text{-}field\text{-}as\text{-}commutative\text{-}ring}(k), r, 0) \vee \\ \quad \quad \operatorname{deg\text{-}bound}(\operatorname{normed\text{-}field\text{-}as\text{-}commutative\text{-}ring}(k), s, 0) \end{array}}}
\vnbprop{nv-principal-of-norm}{\vnbstmt{\begin{array}{@{}l@{}} \forall k.; \forall g \in \operatorname{nv\text{-}integers}(k), x \in \operatorname{carr}(k). \\ \quad (\operatorname{is\text{-}nv\text{-}field}(k) \wedge (\neg (g = \operatorname{zero}(k)) \wedge (\operatorname{fnrm}(k)(x) \leq \operatorname{fnrm}(k)(g)))) \Rightarrow  \\ \quad \quad x \in \operatorname{nv\text{-}principal}(k, \operatorname{nv\text{-}integers}(k), g) \end{array}}}
\vnbprop{ccr-4-31}{\vnbstmt{\begin{array}{@{}l@{}} \forall k.; \forall w \in \operatorname{nv\text{-}integers}(k), q \in \operatorname{finsupp}(\operatorname{normed\text{-}field\text{-}as\text{-}commutative\text{-}ring}(k), \operatorname{nn\text{-}add\text{-}monoid}), n \in \mathbb{N}, r, s \in \operatorname{finsupp}(\operatorname{normed\text{-}field\text{-}as\text{-}commutative\text{-}ring}(k), \operatorname{nn\text{-}add\text{-}monoid}). \\ \quad \operatorname{is\text{-}nv\text{-}field}(k) \wedge \\ \quad \neg (w = \operatorname{zero}(k)) \wedge \\ \quad \operatorname{nv\text{-}maximal}(k, \operatorname{nv\text{-}integers}(k)) = \operatorname{nv\text{-}principal}(k, \operatorname{nv\text{-}integers}(k), w) \wedge \\ \quad \forall i \in \mathbb{N}.\; (q(i) \in \operatorname{nv\text{-}integers}(k)) \wedge \\ \quad 1 \leq n \wedge \\ \quad \operatorname{deg\text{-}bound}(\operatorname{normed\text{-}field\text{-}as\text{-}commutative\text{-}ring}(k), q, n) \wedge \\ \quad \neg (q(0) \in \operatorname{nv\text{-}principal}(k, \operatorname{nv\text{-}integers}(k), \operatorname{mul}(k)(w, w))) \wedge \\ \quad \forall i \in \mathbb{N}.\; ((i < n) \Rightarrow (q(i) \in \operatorname{nv\text{-}principal}(k, \operatorname{nv\text{-}integers}(k), w))) \wedge \\ \quad \neg (q(n) \in \operatorname{nv\text{-}principal}(k, \operatorname{nv\text{-}integers}(k), w)) \wedge \\ \quad \forall i \in \mathbb{N}.\; (r(i) \in \operatorname{nv\text{-}integers}(k)) \wedge \\ \quad \forall i \in \mathbb{N}.\; (s(i) \in \operatorname{nv\text{-}integers}(k)) \wedge \\ \quad q = \operatorname{monalg\text{-}mul}(\operatorname{normed\text{-}field\text{-}as\text{-}commutative\text{-}ring}(k), \operatorname{nn\text{-}add\text{-}monoid}, r, s) \Rightarrow \\ \quad \quad \operatorname{deg\text{-}bound}(\operatorname{normed\text{-}field\text{-}as\text{-}commutative\text{-}ring}(k), r, 0) \vee \\ \quad \quad \operatorname{deg\text{-}bound}(\operatorname{normed\text{-}field\text{-}as\text{-}commutative\text{-}ring}(k), s, 0) \end{array}}}
\vnbprop{nv-max-coef-ind}{\vnbstmt{\begin{array}{@{}l@{}} \forall m \in \mathbb{N}, k.; \forall p \in (\mathbb{N} \to \operatorname{carr}(k)). \\ \quad \operatorname{is\text{-}nv\text{-}field}(k) \Rightarrow  \\ \quad \quad \exists a \in \mathbb{N}. \\ \quad \quad \quad (a \leq m) \wedge \forall b \in \mathbb{N}.\; ((b \leq m) \Rightarrow (\operatorname{fnrm}(k)(p(b)) \leq \operatorname{fnrm}(k)(p(a)))) \end{array}}}
\vnbprop{nv-content-exists}{\vnbstmt{\begin{array}{@{}l@{}} \forall k.; \forall p \in \operatorname{finsupp}(\operatorname{normed\text{-}field\text{-}as\text{-}commutative\text{-}ring}(k), \operatorname{nn\text{-}add\text{-}monoid}). \\ \quad (\operatorname{is\text{-}nv\text{-}field}(k) \wedge \exists \mathit{e4k} \in \mathbb{N}.\; \neg (p(\mathit{e4k}) = \operatorname{zero}(\operatorname{normed\text{-}field\text{-}as\text{-}commutative\text{-}ring}(k)))) \Rightarrow  \\ \quad \quad \exists c \in \mathbb{N}.\; ((0 < \operatorname{fnrm}(k)(p(c))) \wedge \operatorname{is\text{-}nv\text{-}content}(k, p, \operatorname{fnrm}(k)(p(c)))) \end{array}}}
\vnbprop{nv-content-unique}{\vnbstmt{\begin{array}{@{}l@{}} \forall k, p.; \forall c, d \in \mathbb{R}. \\ \quad (\operatorname{is\text{-}nv\text{-}content}(k, p, c) \wedge \operatorname{is\text{-}nv\text{-}content}(k, p, d)) \Rightarrow (c = d) \end{array}}}
\vnbprop{ccr-4-32}{\vnbstmt{\begin{array}{@{}l@{}} \forall k.; \forall w \in \operatorname{nv\text{-}integers}(k), q \in \operatorname{finsupp}(\operatorname{normed\text{-}field\text{-}as\text{-}commutative\text{-}ring}(k), \operatorname{nn\text{-}add\text{-}monoid}), n \in \mathbb{N}, r, s \in \operatorname{finsupp}(\operatorname{normed\text{-}field\text{-}as\text{-}commutative\text{-}ring}(k), \operatorname{nn\text{-}add\text{-}monoid}). \\ \quad \operatorname{is\text{-}nv\text{-}field}(k) \wedge \\ \quad \neg (w = \operatorname{zero}(k)) \wedge \\ \quad \operatorname{nv\text{-}maximal}(k, \operatorname{nv\text{-}integers}(k)) = \operatorname{nv\text{-}principal}(k, \operatorname{nv\text{-}integers}(k), w) \wedge \\ \quad \forall i \in \mathbb{N}.\; (q(i) \in \operatorname{nv\text{-}integers}(k)) \wedge \\ \quad 1 \leq n \wedge \\ \quad \operatorname{deg\text{-}bound}(\operatorname{normed\text{-}field\text{-}as\text{-}commutative\text{-}ring}(k), q, n) \wedge \\ \quad \neg (q(0) \in \operatorname{nv\text{-}principal}(k, \operatorname{nv\text{-}integers}(k), \operatorname{mul}(k)(w, w))) \wedge \\ \quad \forall i \in \mathbb{N}.\; ((i < n) \Rightarrow (q(i) \in \operatorname{nv\text{-}principal}(k, \operatorname{nv\text{-}integers}(k), w))) \wedge \\ \quad \neg (q(n) \in \operatorname{nv\text{-}principal}(k, \operatorname{nv\text{-}integers}(k), w)) \wedge \\ \quad q = \operatorname{monalg\text{-}mul}(\operatorname{normed\text{-}field\text{-}as\text{-}commutative\text{-}ring}(k), \operatorname{nn\text{-}add\text{-}monoid}, r, s) \Rightarrow \\ \quad \quad \operatorname{deg\text{-}bound}(\operatorname{normed\text{-}field\text{-}as\text{-}commutative\text{-}ring}(k), r, 0) \vee \\ \quad \quad \operatorname{deg\text{-}bound}(\operatorname{normed\text{-}field\text{-}as\text{-}commutative\text{-}ring}(k), s, 0) \end{array}}}

\vnbbookfile{theorem-library/ccr4x-zpow}{theorem-library/ccr4x-zpow}
\noindent{\small ccr4x-zpow.scm -- wave of 2026-11-09/10 (agent CCR-1109-C), wired 2026-11-10 from scratchpad/ccr1109c/ccr4x-zpow.scm. Library: analysis.}\par
\vnbprop{zz-int-halves}{\vnbstmt{\begin{array}{@{}l@{}} \forall k \in \mathbb{Z}. \\ \quad (\operatorname{max}(k, 0) \in \mathbb{N}) \wedge ((\operatorname{max}(-k, 0) \in \mathbb{N}) \wedge ((\operatorname{max}(k, 0) + -\operatorname{max}(-k, 0)) = k)) \end{array}}}
\vnbprop{nv-zpow-def}{\vnbstmt{\begin{array}{@{}l@{}} \forall k.; \forall w \in \operatorname{carr}(k), e \in \mathbb{Z}. \\ \quad (\operatorname{is\text{-}normed\text{-}field}(k) \wedge \neg (w = \operatorname{zero}(k))) \Rightarrow  \\ \quad \quad \operatorname{nv\text{-}zpow}(k, w, e) \in \operatorname{carr}(k) \wedge \\ \quad \quad \operatorname{mul}(k)(\operatorname{nv\text{-}zpow}(k, w, e), \operatorname{ring\text{-}power}(\operatorname{normed\text{-}field\text{-}as\text{-}commutative\text{-}ring}(k), w, \operatorname{max}(-e, 0))) \\ \quad \quad =\; \operatorname{ring\text{-}power}(\operatorname{normed\text{-}field\text{-}as\text{-}commutative\text{-}ring}(k), w, \operatorname{max}(e, 0)) \end{array}}}
\vnbprop{nv-zpow-nat}{\vnbstmt{\begin{array}{@{}l@{}} \forall k.; \forall w \in \operatorname{carr}(k), e \in \mathbb{N}. \\ \quad (\operatorname{is\text{-}normed\text{-}field}(k) \wedge \neg (w = \operatorname{zero}(k))) \Rightarrow  \\ \quad \quad \operatorname{nv\text{-}zpow}(k, w, e) = \operatorname{ring\text{-}power}(\operatorname{normed\text{-}field\text{-}as\text{-}commutative\text{-}ring}(k), w, e) \end{array}}}
\vnbprop{nv-zpow-neg}{\vnbstmt{\begin{array}{@{}l@{}} \forall k.; \forall w \in \operatorname{carr}(k), j \in \mathbb{N}. \\ \quad (\operatorname{is\text{-}normed\text{-}field}(k) \wedge \neg (w = \operatorname{zero}(k))) \Rightarrow  \\ \quad \quad \operatorname{nv\text{-}zpow}(k, w, -j) \in \operatorname{carr}(k) \wedge \\ \quad \quad \operatorname{mul}(k)(\operatorname{nv\text{-}zpow}(k, w, -j), \operatorname{ring\text{-}power}(\operatorname{normed\text{-}field\text{-}as\text{-}commutative\text{-}ring}(k), w, j)) = \operatorname{one}(k) \end{array}}}
\vnbprop{ccr-4-8}{\vnbstmt{\begin{array}{@{}l@{}} \forall k, a, w.; \forall x \in \operatorname{carr}(k). \\ \quad \operatorname{is\text{-}compact\text{-}nv\text{-}domain}(k, a) \wedge \\ \quad w \in a \wedge \\ \quad \operatorname{nv\text{-}maximal}(k, a) = \operatorname{nv\text{-}principal}(k, a, w) \wedge \\ \quad \neg (w = \operatorname{zero}(k)) \wedge \\ \quad a = \operatorname{nv\text{-}integers}(k) \wedge \\ \quad \neg (x = \operatorname{zero}(k)) \Rightarrow \\ \quad \quad \exists e \in \mathbb{Z}, u \in \operatorname{nv\text{-}units}(k, a).\; (x = \operatorname{mul}(k)(\operatorname{nv\text{-}zpow}(k, w, e), u)) \end{array}}}
\vnbprop{ccr-4-8-unique}{\vnbstmt{\begin{array}{@{}l@{}} \forall k, a, w.; \forall b, l \in \mathbb{Z}, u, v \in \operatorname{nv\text{-}units}(k, a). \\ \quad \operatorname{is\text{-}compact\text{-}nv\text{-}domain}(k, a) \wedge \\ \quad w \in a \wedge \\ \quad \operatorname{nv\text{-}maximal}(k, a) = \operatorname{nv\text{-}principal}(k, a, w) \wedge \\ \quad \neg (w = \operatorname{zero}(k)) \wedge \\ \quad a = \operatorname{nv\text{-}integers}(k) \wedge \\ \quad \operatorname{mul}(k)(\operatorname{nv\text{-}zpow}(k, w, b), u) = \operatorname{mul}(k)(\operatorname{nv\text{-}zpow}(k, w, l), v) \Rightarrow \\ \quad \quad (b = l) \wedge (u = v) \end{array}}}
\vnbprop{ccr-remark-4-9}{\vnbstmt{\begin{array}{@{}l@{}} \forall k, a, w.; \forall x \in \operatorname{carr}(k). \\ \quad \operatorname{is\text{-}compact\text{-}nv\text{-}domain}(k, a) \wedge \\ \quad w \in a \wedge \\ \quad \operatorname{nv\text{-}maximal}(k, a) = \operatorname{nv\text{-}principal}(k, a, w) \wedge \\ \quad \neg (w = \operatorname{zero}(k)) \wedge \\ \quad a = \operatorname{nv\text{-}integers}(k) \wedge \\ \quad \neg (x = \operatorname{zero}(k)) \Rightarrow \\ \quad \quad \exists e \in \mathbb{Z}.\; (\operatorname{fnrm}(k)(x) = \operatorname{fnrm}(k)(\operatorname{nv\text{-}zpow}(k, w, e))) \end{array}}}
\vnbprop{nv-fractions-of-maximal-explicit}{\vnbstmt{\begin{array}{@{}l@{}} \forall k.; \forall w, x \in \operatorname{carr}(k). \\ \quad (\operatorname{is\text{-}nv\text{-}field}(k) \wedge (\neg (w = \operatorname{zero}(k)) \wedge (\operatorname{fnrm}(k)(w) < 1))) \Rightarrow  \\ \quad \quad \exists u, v \in \operatorname{nv\text{-}maximal}(k, \operatorname{nv\text{-}integers}(k)). \\ \quad \quad \quad \neg (v = \operatorname{zero}(k)) \wedge (\operatorname{mul}(k)(x, v) = u) \end{array}}}
\vnbprop{nv-fractions-of-maximal}{\vnbstmt{\begin{array}{@{}l@{}} \forall k.; \forall w \in \operatorname{carr}(k). \\ \quad (\operatorname{is\text{-}nv\text{-}field}(k) \wedge (\neg (w = \operatorname{zero}(k)) \wedge (\operatorname{fnrm}(k)(w) < 1))) \Rightarrow  \\ \quad \quad \operatorname{is\text{-}fractions\text{-}of}(k, \operatorname{nv\text{-}maximal}(k, \operatorname{nv\text{-}integers}(k))) \end{array}}}
\vnbprop{nv-rep-set-sub}{\vnbstmt{\forall k, a, p.; \forall x \in p.\; (\operatorname{is\text{-}nv\text{-}rep\text{-}set}(k, a, p) \Rightarrow (x \in a))}}
\vnbprop{nv-rep-set-zero}{\vnbstmt{\forall k, a, p.\; (\operatorname{is\text{-}nv\text{-}rep\text{-}set}(k, a, p) \Rightarrow (\operatorname{zero}(k) \in p))}}
\vnbprop{nv-rep-set-one}{\vnbstmt{\forall k, a, p.\; (\operatorname{is\text{-}nv\text{-}rep\text{-}set}(k, a, p) \Rightarrow (\operatorname{one}(k) \in p))}}
\vnbprop{nv-rep-set-exists}{\vnbstmt{\begin{array}{@{}l@{}} \forall k, a, p.; \forall x \in a. \\ \quad \operatorname{is\text{-}nv\text{-}rep\text{-}set}(k, a, p) \Rightarrow  \\ \quad \quad \exists y \in p.\; (\operatorname{fnrm}(k)(\operatorname{add}(k)(x, \operatorname{neg}(k)(y))) < 1) \end{array}}}
\vnbprop{nv-rep-set-unique}{\vnbstmt{\begin{array}{@{}l@{}} \forall k, a, p.; \forall x, y \in p. \\ \quad (\operatorname{is\text{-}nv\text{-}rep\text{-}set}(k, a, p) \wedge (\operatorname{fnrm}(k)(\operatorname{add}(k)(x, \operatorname{neg}(k)(y))) < 1)) \Rightarrow (x = y) \end{array}}}
\vnbprop{nv-rep-digit-step}{\vnbstmt{\begin{array}{@{}l@{}} \forall k, a, w, p.; \forall z \in a. \\ \quad \operatorname{is\text{-}compact\text{-}nv\text{-}domain}(k, a) \wedge \\ \quad w \in a \wedge \\ \quad \operatorname{nv\text{-}maximal}(k, a) = \operatorname{nv\text{-}principal}(k, a, w) \wedge \\ \quad \neg (w = \operatorname{zero}(k)) \wedge \\ \quad a = \operatorname{nv\text{-}integers}(k) \wedge \\ \quad \operatorname{is\text{-}nv\text{-}rep\text{-}set}(k, a, p) \Rightarrow \\ \quad \quad \exists b \in p, d \in a.\; (z = \operatorname{add}(k)(b, \operatorname{mul}(k)(d, w))) \end{array}}}
\vnbprop{nv-rep-digit-unique}{\vnbstmt{\begin{array}{@{}l@{}} \forall k, a, w, p.; \forall b, c \in p, d, f \in a. \\ \quad \operatorname{is\text{-}compact\text{-}nv\text{-}domain}(k, a) \wedge \\ \quad w \in a \wedge \\ \quad \operatorname{nv\text{-}maximal}(k, a) = \operatorname{nv\text{-}principal}(k, a, w) \wedge \\ \quad \neg (w = \operatorname{zero}(k)) \wedge \\ \quad a = \operatorname{nv\text{-}integers}(k) \wedge \\ \quad \operatorname{is\text{-}nv\text{-}rep\text{-}set}(k, a, p) \wedge \\ \quad \operatorname{add}(k)(b, \operatorname{mul}(k)(d, w)) = \operatorname{add}(k)(c, \operatorname{mul}(k)(f, w)) \Rightarrow \\ \quad \quad (b = c) \wedge (d = f) \end{array}}}
\vnbprop{nv-rep-set-norm-one}{\vnbstmt{\begin{array}{@{}l@{}} \forall k, a, w, p.; \forall b \in p. \\ \quad \operatorname{is\text{-}compact\text{-}nv\text{-}domain}(k, a) \wedge \\ \quad w \in a \wedge \\ \quad \operatorname{nv\text{-}maximal}(k, a) = \operatorname{nv\text{-}principal}(k, a, w) \wedge \\ \quad \neg (w = \operatorname{zero}(k)) \wedge \\ \quad a = \operatorname{nv\text{-}integers}(k) \wedge \\ \quad \operatorname{is\text{-}nv\text{-}rep\text{-}set}(k, a, p) \wedge \\ \quad \neg (b = \operatorname{zero}(k)) \Rightarrow \\ \quad \quad \operatorname{fnrm}(k)(b) = 1 \end{array}}}

\vnbbookfile{theorem-library/ccr4x-series}{theorem-library/ccr4x-series}
\noindent{\small ccr4x-series.scm -- wave of 2026-11-09/10 (agent CCR-1109-C), wired 2026-11-10 from scratchpad/ccr1109c/ccr4x-series.scm. Library: analysis.}\par
\vnbprop{nv-zpow-char}{\vnbstmt{\begin{array}{@{}l@{}} \forall k, w.; \forall e \in \mathbb{Z}, s, t \in \mathbb{N}. \\ \quad (\operatorname{is\text{-}normed\text{-}field}(k) \wedge ((w \in \operatorname{carr}(k)) \wedge (\neg (w = \operatorname{zero}(k)) \wedge ((s + -t) = e)))) \Rightarrow  \\ \quad \quad \operatorname{mul}(k)(\operatorname{nv\text{-}zpow}(k, w, e), \operatorname{ring\text{-}power}(\operatorname{normed\text{-}field\text{-}as\text{-}commutative\text{-}ring}(k), w, t)) \\ \quad \quad =\; \operatorname{ring\text{-}power}(\operatorname{normed\text{-}field\text{-}as\text{-}commutative\text{-}ring}(k), w, s) \end{array}}}
\vnbprop{nv-zpow-succ}{\vnbstmt{\begin{array}{@{}l@{}} \forall k, w.; \forall e \in \mathbb{Z}. \\ \quad (\operatorname{is\text{-}normed\text{-}field}(k) \wedge ((w \in \operatorname{carr}(k)) \wedge \neg (w = \operatorname{zero}(k)))) \Rightarrow  \\ \quad \quad \operatorname{mul}(k)(\operatorname{nv\text{-}zpow}(k, w, e), w) = \operatorname{nv\text{-}zpow}(k, w, (e + 1)) \end{array}}}
\vnbprop{nv-zpow-inv}{\vnbstmt{\begin{array}{@{}l@{}} \forall k, w.; \forall e \in \mathbb{Z}. \\ \quad (\operatorname{is\text{-}normed\text{-}field}(k) \wedge ((w \in \operatorname{carr}(k)) \wedge \neg (w = \operatorname{zero}(k)))) \Rightarrow  \\ \quad \quad \operatorname{mul}(k)(\operatorname{nv\text{-}zpow}(k, w, e), \operatorname{nv\text{-}zpow}(k, w, -e)) = \operatorname{one}(k) \end{array}}}
\vnbprop{nv-zpow-norm-nat}{\vnbstmt{\begin{array}{@{}l@{}} \forall k, w.; \forall e \in \mathbb{N}. \\ \quad (\operatorname{is\text{-}normed\text{-}field}(k) \wedge ((w \in \operatorname{carr}(k)) \wedge \neg (w = \operatorname{zero}(k)))) \Rightarrow  \\ \quad \quad \operatorname{fnrm}(k)(\operatorname{nv\text{-}zpow}(k, w, e)) = \operatorname{power}(\operatorname{fnrm}(k)(w), e) \end{array}}}
\vnbprop{nv-psum-zero}{\vnbstmt{\begin{array}{@{}l@{}} \forall k, w, a, n. \\ \quad \operatorname{is\text{-}normed\text{-}field}(k) \Rightarrow (\operatorname{nv\text{-}psum}(k, w, a, n, 0) = \operatorname{zero}(k)) \end{array}}}
\vnbprop{nv-psum-succ}{\vnbstmt{\begin{array}{@{}l@{}} \forall k, w, a, n.; \forall m \in \mathbb{N}. \\ \quad \operatorname{is\text{-}normed\text{-}field}(k) \Rightarrow  \\ \quad \quad (\operatorname{nv\text{-}psum}(k, w, a, n, (m + 1)) \simeq \operatorname{add}(k)(\operatorname{nv\text{-}psum}(k, w, a, n, m), \operatorname{mul}(k)(a((n + m)), \operatorname{nv\text{-}zpow}(k, w, (n + m))))) \end{array}}}
\vnbprop{nv-psum-type}{\vnbstmt{\begin{array}{@{}l@{}} \forall m \in \mathbb{N}, k, w.; \forall a \in (\mathbb{Z} \to \operatorname{carr}(k)), n \in \mathbb{Z}. \\ \quad (\operatorname{is\text{-}normed\text{-}field}(k) \wedge ((w \in \operatorname{carr}(k)) \wedge \neg (w = \operatorname{zero}(k)))) \Rightarrow  \\ \quad \quad \operatorname{nv\text{-}psum}(k, w, a, n, m) \in \operatorname{carr}(k) \end{array}}}
\vnbprop{nv-digit-def}{\vnbstmt{\begin{array}{@{}l@{}} \forall k, a, w, p.; \forall z \in a. \\ \quad \operatorname{is\text{-}compact\text{-}nv\text{-}domain}(k, a) \wedge \\ \quad w \in a \wedge \\ \quad \operatorname{nv\text{-}maximal}(k, a) = \operatorname{nv\text{-}principal}(k, a, w) \wedge \\ \quad \neg (w = \operatorname{zero}(k)) \wedge \\ \quad a = \operatorname{nv\text{-}integers}(k) \wedge \\ \quad \operatorname{is\text{-}nv\text{-}rep\text{-}set}(k, a, p) \Rightarrow \\ \quad \quad ((\operatorname{nv\text{-}digit}(k, p, z) \in p) \wedge (\operatorname{fnrm}(k)(\operatorname{add}(k)(z, \operatorname{neg}(k)(\operatorname{nv\text{-}digit}(k, p, z)))) < 1)) \end{array}}}
\vnbprop{nv-digit-step}{\vnbstmt{\begin{array}{@{}l@{}} \forall k, a, w, p.; \forall z \in a. \\ \quad \operatorname{is\text{-}compact\text{-}nv\text{-}domain}(k, a) \wedge \\ \quad w \in a \wedge \\ \quad \operatorname{nv\text{-}maximal}(k, a) = \operatorname{nv\text{-}principal}(k, a, w) \wedge \\ \quad \neg (w = \operatorname{zero}(k)) \wedge \\ \quad a = \operatorname{nv\text{-}integers}(k) \wedge \\ \quad \operatorname{is\text{-}nv\text{-}rep\text{-}set}(k, a, p) \Rightarrow \\ \quad \quad \exists d \in a.\; (z = \operatorname{add}(k)(\operatorname{nv\text{-}digit}(k, p, z), \operatorname{mul}(k)(w, d))) \end{array}}}
\vnbprop{nv-digit-sequence}{\vnbstmt{\begin{array}{@{}l@{}} \forall k, a, w, p.; \forall y \in a. \\ \quad \operatorname{is\text{-}compact\text{-}nv\text{-}domain}(k, a) \wedge \\ \quad w \in a \wedge \\ \quad \operatorname{nv\text{-}maximal}(k, a) = \operatorname{nv\text{-}principal}(k, a, w) \wedge \\ \quad \neg (w = \operatorname{zero}(k)) \wedge \\ \quad a = \operatorname{nv\text{-}integers}(k) \wedge \\ \quad \operatorname{is\text{-}nv\text{-}rep\text{-}set}(k, a, p) \Rightarrow \\ \quad \quad \exists f \in (\mathbb{N} \to a). \\ \quad \quad \quad ((f(0) = y) \wedge \forall s \in \mathbb{N}.\; (f(s) = \operatorname{add}(k)(\operatorname{nv\text{-}digit}(k, p, f(s)), \operatorname{mul}(k)(w, f((s + 1)))))) \end{array}}}
\vnbprop{nv-expand-ind}{\vnbstmt{\begin{array}{@{}l@{}} \forall m \in \mathbb{N}, k, a, w, p.; \forall f \in (\mathbb{N} \to a), b \in (\mathbb{Z} \to \operatorname{carr}(k)), n \in \mathbb{Z}. \\ \quad \operatorname{is\text{-}compact\text{-}nv\text{-}domain}(k, a) \wedge \\ \quad w \in a \wedge \\ \quad \operatorname{nv\text{-}maximal}(k, a) = \operatorname{nv\text{-}principal}(k, a, w) \wedge \\ \quad \neg (w = \operatorname{zero}(k)) \wedge \\ \quad a = \operatorname{nv\text{-}integers}(k) \wedge \\ \quad \operatorname{is\text{-}nv\text{-}rep\text{-}set}(k, a, p) \wedge \\ \quad \forall s \in \mathbb{N}. \\ \quad \quad f(s) = \operatorname{add}(k)(\operatorname{nv\text{-}digit}(k, p, f(s)), \operatorname{mul}(k)(w, f((s + 1)))) \wedge \\ \quad \forall j \in \mathbb{N}.\; (b((n + j)) = \operatorname{nv\text{-}digit}(k, p, f(j))) \Rightarrow \\ \quad \quad \operatorname{mul}(k)(\operatorname{nv\text{-}zpow}(k, w, n), f(0)) \\ \quad \quad =\; \operatorname{add}(k)(\operatorname{nv\text{-}psum}(k, w, b, n, m), \operatorname{mul}(k)(\operatorname{nv\text{-}zpow}(k, w, (n + m)), f(m))) \end{array}}}
\vnbprop{ccr-4-17}{\vnbstmt{\begin{array}{@{}l@{}} \forall k, a, w, p.; \forall x \in \operatorname{carr}(k), n \in \mathbb{Z}. \\ \quad \operatorname{is\text{-}compact\text{-}nv\text{-}domain}(k, a) \wedge \\ \quad w \in a \wedge \\ \quad \operatorname{nv\text{-}maximal}(k, a) = \operatorname{nv\text{-}principal}(k, a, w) \wedge \\ \quad \neg (w = \operatorname{zero}(k)) \wedge \\ \quad a = \operatorname{nv\text{-}integers}(k) \wedge \\ \quad \operatorname{is\text{-}nv\text{-}rep\text{-}set}(k, a, p) \wedge \\ \quad \operatorname{mul}(k)(\operatorname{nv\text{-}zpow}(k, w, -n), x) \in a \Rightarrow \\ \quad \quad \exists b.\; \operatorname{is\text{-}nv\text{-}expansion}(k, w, p, b, n, x) \end{array}}}
\vnbprop{ccr-4-17-exists}{\vnbstmt{\begin{array}{@{}l@{}} \forall k, a, w, p.; \forall x \in \operatorname{carr}(k). \\ \quad \operatorname{is\text{-}compact\text{-}nv\text{-}domain}(k, a) \wedge \\ \quad w \in a \wedge \\ \quad \operatorname{nv\text{-}maximal}(k, a) = \operatorname{nv\text{-}principal}(k, a, w) \wedge \\ \quad \neg (w = \operatorname{zero}(k)) \wedge \\ \quad a = \operatorname{nv\text{-}integers}(k) \wedge \\ \quad \operatorname{is\text{-}nv\text{-}rep\text{-}set}(k, a, p) \Rightarrow \\ \quad \quad \exists n \in \mathbb{Z}, b.\; \operatorname{is\text{-}nv\text{-}expansion}(k, w, p, b, n, x) \end{array}}}

\vnbbookfile{theorem-library/ccr4x-unique}{theorem-library/ccr4x-unique}
\noindent{\small ccr4x-unique.scm -- wave of 2026-11-09/10 (agent CCR-1109-C2), wired 2026-11-10 from scratchpad/ccr1109c2/ccr4x-unique.scm. Library: analysis.}\par
\vnbprop{nv-zpow-add}{\vnbstmt{\begin{array}{@{}l@{}} \forall k, w.; \forall e, f \in \mathbb{Z}. \\ \quad (\operatorname{is\text{-}normed\text{-}field}(k) \wedge ((w \in \operatorname{carr}(k)) \wedge \neg (w = \operatorname{zero}(k)))) \Rightarrow  \\ \quad \quad \operatorname{mul}(k)(\operatorname{nv\text{-}zpow}(k, w, e), \operatorname{nv\text{-}zpow}(k, w, f)) = \operatorname{nv\text{-}zpow}(k, w, (e + f)) \end{array}}}
\vnbprop{nv-zpow-nonzero}{\vnbstmt{\begin{array}{@{}l@{}} \forall k, w.; \forall e \in \mathbb{Z}. \\ \quad (\operatorname{is\text{-}normed\text{-}field}(k) \wedge ((w \in \operatorname{carr}(k)) \wedge \neg (w = \operatorname{zero}(k)))) \Rightarrow  \\ \quad \quad \neg (\operatorname{nv\text{-}zpow}(k, w, e) = \operatorname{zero}(k)) \end{array}}}
\vnbprop{nv-zpow-norm-succ}{\vnbstmt{\begin{array}{@{}l@{}} \forall k, w.; \forall e \in \mathbb{Z}. \\ \quad (\operatorname{is\text{-}normed\text{-}field}(k) \wedge ((w \in \operatorname{carr}(k)) \wedge \neg (w = \operatorname{zero}(k)))) \Rightarrow  \\ \quad \quad \operatorname{fnrm}(k)(\operatorname{nv\text{-}zpow}(k, w, (e + 1))) = (\operatorname{fnrm}(k)(\operatorname{nv\text{-}zpow}(k, w, e)) \cdot \operatorname{fnrm}(k)(w)) \end{array}}}
\vnbprop{nv-zpow-norm-mono}{\vnbstmt{\begin{array}{@{}l@{}} \forall j \in \mathbb{N}, k, w.; \forall e \in \mathbb{Z}. \\ \quad (\operatorname{is\text{-}normed\text{-}field}(k) \wedge ((w \in \operatorname{carr}(k)) \wedge (\neg (w = \operatorname{zero}(k)) \wedge (\operatorname{fnrm}(k)(w) \leq 1)))) \Rightarrow  \\ \quad \quad \operatorname{fnrm}(k)(\operatorname{nv\text{-}zpow}(k, w, (e + j))) \leq \operatorname{fnrm}(k)(\operatorname{nv\text{-}zpow}(k, w, e)) \end{array}}}
\vnbprop{nv-ultra-le}{\vnbstmt{\begin{array}{@{}l@{}} \forall k.; \forall u, v, z \in \operatorname{carr}(k), r \in \mathbb{R}. \\ \quad \operatorname{is\text{-}nv\text{-}field}(k) \wedge \\ \quad \operatorname{fnrm}(k)(\operatorname{add}(k)(u, \operatorname{neg}(k)(v))) \leq r \wedge \\ \quad \operatorname{fnrm}(k)(\operatorname{add}(k)(v, \operatorname{neg}(k)(z))) \leq r \Rightarrow \\ \quad \quad \operatorname{fnrm}(k)(\operatorname{add}(k)(u, \operatorname{neg}(k)(z))) \leq r \end{array}}}
\vnbprop{nv-tail-ind}{\vnbstmt{\begin{array}{@{}l@{}} \forall l \in \mathbb{N}, k, a, w, p.; \forall b \in (\mathbb{Z} \to p), n \in \mathbb{Z}, m \in \mathbb{N}. \\ \quad \operatorname{is\text{-}compact\text{-}nv\text{-}domain}(k, a) \wedge \\ \quad w \in a \wedge \\ \quad \operatorname{nv\text{-}maximal}(k, a) = \operatorname{nv\text{-}principal}(k, a, w) \wedge \\ \quad \neg (w = \operatorname{zero}(k)) \wedge \\ \quad a = \operatorname{nv\text{-}integers}(k) \wedge \\ \quad \operatorname{is\text{-}nv\text{-}rep\text{-}set}(k, a, p) \Rightarrow \\ \quad \quad \operatorname{fnrm}(k)(\operatorname{add}(k)(\operatorname{nv\text{-}psum}(k, w, b, n, (m + l)), \operatorname{neg}(k)(\operatorname{nv\text{-}psum}(k, w, b, n, m)))) \\ \quad \quad \leq\; \operatorname{fnrm}(k)(\operatorname{nv\text{-}zpow}(k, w, (n + m))) \end{array}}}
\vnbprop{nv-tail}{\vnbstmt{\begin{array}{@{}l@{}} \forall k, a, w, p.; \forall b \in (\mathbb{Z} \to p), n \in \mathbb{Z}, x.; \forall m \in \mathbb{N}. \\ \quad \operatorname{is\text{-}compact\text{-}nv\text{-}domain}(k, a) \wedge \\ \quad w \in a \wedge \\ \quad \operatorname{nv\text{-}maximal}(k, a) = \operatorname{nv\text{-}principal}(k, a, w) \wedge \\ \quad \neg (w = \operatorname{zero}(k)) \wedge \\ \quad a = \operatorname{nv\text{-}integers}(k) \wedge \\ \quad \operatorname{is\text{-}nv\text{-}rep\text{-}set}(k, a, p) \wedge \\ \quad \operatorname{converges\text{-}to}(\operatorname{nf\text{-}metric\text{-}space}(k), (\lambda \mathit{nem} \in \mathbb{N}.\; \operatorname{nv\text{-}psum}(k, w, b, n, \mathit{nem})), x) \Rightarrow \\ \quad \quad \operatorname{fnrm}(k)(\operatorname{add}(k)(x, \operatorname{neg}(k)(\operatorname{nv\text{-}psum}(k, w, b, n, m)))) \leq \operatorname{fnrm}(k)(\operatorname{nv\text{-}zpow}(k, w, (n + m))) \end{array}}}
\vnbprop{nv-pad-psum}{\vnbstmt{\begin{array}{@{}l@{}} \forall m \in \mathbb{N}, k, a, w, p.; \forall b \in (\mathbb{Z} \to p), l \in \mathbb{Z}, d \in \mathbb{N}, n \in \mathbb{Z}. \\ \quad \operatorname{is\text{-}compact\text{-}nv\text{-}domain}(k, a) \wedge \\ \quad w \in a \wedge \\ \quad \operatorname{nv\text{-}maximal}(k, a) = \operatorname{nv\text{-}principal}(k, a, w) \wedge \\ \quad \neg (w = \operatorname{zero}(k)) \wedge \\ \quad a = \operatorname{nv\text{-}integers}(k) \wedge \\ \quad \operatorname{is\text{-}nv\text{-}rep\text{-}set}(k, a, p) \wedge \\ \quad (l + d) = n \wedge \\ \quad \forall z \in \mathbb{Z}.\; ((z < n) \Rightarrow (b(z) = \operatorname{zero}(k))) \Rightarrow \\ \quad \quad \operatorname{nv\text{-}psum}(k, w, b, l, (d + m)) = \operatorname{nv\text{-}psum}(k, w, b, n, m) \end{array}}}
\vnbprop{nv-pad}{\vnbstmt{\begin{array}{@{}l@{}} \forall k, a, w, p, b.; \forall n \in \mathbb{Z}, x.; \forall l \in \mathbb{Z}. \\ \quad \operatorname{is\text{-}compact\text{-}nv\text{-}domain}(k, a) \wedge \\ \quad w \in a \wedge \\ \quad \operatorname{nv\text{-}maximal}(k, a) = \operatorname{nv\text{-}principal}(k, a, w) \wedge \\ \quad \neg (w = \operatorname{zero}(k)) \wedge \\ \quad a = \operatorname{nv\text{-}integers}(k) \wedge \\ \quad \operatorname{is\text{-}nv\text{-}rep\text{-}set}(k, a, p) \wedge \\ \quad \operatorname{is\text{-}nv\text{-}expansion}(k, w, p, b, n, x) \wedge \\ \quad l \leq n \Rightarrow \\ \quad \quad \operatorname{is\text{-}nv\text{-}expansion}(k, w, p, b, l, x) \end{array}}}
\vnbprop{nv-same-start}{\vnbstmt{\begin{array}{@{}l@{}} \forall b \in \mathbb{N}, k, a, w, p, c, d.; \forall n \in \mathbb{Z}, x.; \forall j \in \mathbb{N}. \\ \quad \operatorname{is\text{-}compact\text{-}nv\text{-}domain}(k, a) \wedge \\ \quad w \in a \wedge \\ \quad \operatorname{nv\text{-}maximal}(k, a) = \operatorname{nv\text{-}principal}(k, a, w) \wedge \\ \quad \neg (w = \operatorname{zero}(k)) \wedge \\ \quad a = \operatorname{nv\text{-}integers}(k) \wedge \\ \quad \operatorname{is\text{-}nv\text{-}rep\text{-}set}(k, a, p) \wedge \\ \quad \operatorname{is\text{-}nv\text{-}expansion}(k, w, p, c, n, x) \wedge \\ \quad \operatorname{is\text{-}nv\text{-}expansion}(k, w, p, d, n, x) \wedge \\ \quad j < b \Rightarrow \\ \quad \quad c((n + j)) = d((n + j)) \end{array}}}
\vnbprop{nv-unique-le}{\vnbstmt{\begin{array}{@{}l@{}} \forall k, a, w, p, b.; \forall n \in \mathbb{Z}, c.; \forall m \in \mathbb{Z}, x.; \forall z \in \mathbb{Z}. \\ \quad \operatorname{is\text{-}compact\text{-}nv\text{-}domain}(k, a) \wedge \\ \quad w \in a \wedge \\ \quad \operatorname{nv\text{-}maximal}(k, a) = \operatorname{nv\text{-}principal}(k, a, w) \wedge \\ \quad \neg (w = \operatorname{zero}(k)) \wedge \\ \quad a = \operatorname{nv\text{-}integers}(k) \wedge \\ \quad \operatorname{is\text{-}nv\text{-}rep\text{-}set}(k, a, p) \wedge \\ \quad \operatorname{is\text{-}nv\text{-}expansion}(k, w, p, b, n, x) \wedge \\ \quad \operatorname{is\text{-}nv\text{-}expansion}(k, w, p, c, m, x) \wedge \\ \quad n \leq m \Rightarrow \\ \quad \quad b(z) = c(z) \end{array}}}
\vnbprop{ccr-4-17-unique}{\vnbstmt{\begin{array}{@{}l@{}} \forall k, a, w, p.; \forall x \in \operatorname{carr}(k), b.; \forall n \in \mathbb{Z}, c.; \forall m, z \in \mathbb{Z}. \\ \quad \operatorname{is\text{-}compact\text{-}nv\text{-}domain}(k, a) \wedge \\ \quad w \in a \wedge \\ \quad \operatorname{nv\text{-}maximal}(k, a) = \operatorname{nv\text{-}principal}(k, a, w) \wedge \\ \quad \neg (w = \operatorname{zero}(k)) \wedge \\ \quad a = \operatorname{nv\text{-}integers}(k) \wedge \\ \quad \operatorname{is\text{-}nv\text{-}rep\text{-}set}(k, a, p) \wedge \\ \quad \operatorname{is\text{-}nv\text{-}expansion}(k, w, p, b, n, x) \wedge \\ \quad \operatorname{is\text{-}nv\text{-}expansion}(k, w, p, c, m, x) \Rightarrow \\ \quad \quad b(z) = c(z) \end{array}}}
\vnbprop{ccr-4-17-full}{\vnbstmt{\begin{array}{@{}l@{}} \forall k, a, w, p.; \forall x \in \operatorname{carr}(k). \\ \quad \operatorname{is\text{-}compact\text{-}nv\text{-}domain}(k, a) \wedge \\ \quad w \in a \wedge \\ \quad \operatorname{nv\text{-}maximal}(k, a) = \operatorname{nv\text{-}principal}(k, a, w) \wedge \\ \quad \neg (w = \operatorname{zero}(k)) \wedge \\ \quad a = \operatorname{nv\text{-}integers}(k) \wedge \\ \quad \operatorname{is\text{-}nv\text{-}rep\text{-}set}(k, a, p) \Rightarrow \\ \quad \quad \exists n \in \mathbb{Z}, b.\; \operatorname{is\text{-}nv\text{-}expansion}(k, w, p, b, n, x) \wedge \\ \quad \quad \forall b.; \forall n \in \mathbb{Z}, c.; \forall m, z \in \mathbb{Z}. \\ \quad \quad \quad (\operatorname{is\text{-}nv\text{-}expansion}(k, w, p, b, n, x) \wedge \operatorname{is\text{-}nv\text{-}expansion}(k, w, p, c, m, x)) \Rightarrow  \\ \quad \quad \quad \quad b(z) = c(z) \end{array}}}
\vnbprop{nv-zpow-norm-gap}{\vnbstmt{\begin{array}{@{}l@{}} \forall k, a, w.; \forall x, y \in \operatorname{carr}(k). \\ \quad \operatorname{is\text{-}compact\text{-}nv\text{-}domain}(k, a) \wedge \\ \quad w \in a \wedge \\ \quad \operatorname{nv\text{-}maximal}(k, a) = \operatorname{nv\text{-}principal}(k, a, w) \wedge \\ \quad \neg (w = \operatorname{zero}(k)) \wedge \\ \quad a = \operatorname{nv\text{-}integers}(k) \wedge \\ \quad \neg (x = \operatorname{zero}(k)) \wedge \\ \quad \neg (y = \operatorname{zero}(k)) \wedge \\ \quad \neg (\operatorname{fnrm}(k)(x) = \operatorname{fnrm}(k)(y)) \Rightarrow \\ \quad \quad ((\operatorname{fnrm}(k)(y) \leq (\operatorname{fnrm}(k)(w) \cdot \operatorname{fnrm}(k)(x))) \vee (\operatorname{fnrm}(k)(x) \leq (\operatorname{fnrm}(k)(w) \cdot \operatorname{fnrm}(k)(y)))) \end{array}}}
\vnbprop{ccr-remark-4-9-accumulation}{\vnbstmt{\begin{array}{@{}l@{}} \forall k, a, w, c. \\ \quad \operatorname{is\text{-}compact\text{-}nv\text{-}domain}(k, a) \wedge \\ \quad w \in a \wedge \\ \quad \operatorname{nv\text{-}maximal}(k, a) = \operatorname{nv\text{-}principal}(k, a, w) \wedge \\ \quad \neg (w = \operatorname{zero}(k)) \wedge \\ \quad a = \operatorname{nv\text{-}integers}(k) \wedge \\ \quad \operatorname{is\text{-}accumulation\text{-}point}(\operatorname{rr\text{-}ms}, \{v \in \mathbb{R} : \exists x \in \operatorname{carr}(k).\; (v = \operatorname{fnrm}(k)(x))\}, c) \Rightarrow \\ \quad \quad c = 0 \end{array}}}

\vnbbookfile{theorem-library/ccr-kn-norms}{theorem-library/ccr-kn-norms}
\noindent{\small ccr-kn-norms.scm -- wave of 2026-11-09/10 (agent CCR-1109-E), wired 2026-11-10 from scratchpad/ccr1109e/ccr-kn-norms.scm. Library: analysis.}\par
\vnbprop{ccr-3-8-kn}{\vnbstmt{\begin{array}{@{}l@{}} \forall k.; \forall n \in \mathbb{N}, u_{1}, u_{2}. \\ \quad \operatorname{is\text{-}local\text{-}nv\text{-}field}(k) \wedge \\ \quad \operatorname{is\text{-}norm\text{-}on}(k, \operatorname{ring\text{-}pow}(\operatorname{normed\text{-}field\text{-}as\text{-}commutative\text{-}ring}(k), n), u_{1}) \wedge \\ \quad \operatorname{is\text{-}norm\text{-}on}(k, \operatorname{ring\text{-}pow}(\operatorname{normed\text{-}field\text{-}as\text{-}commutative\text{-}ring}(k), n), u_{2}) \Rightarrow \\ \quad \quad \exists l \in \mathbb{R}. \\ \quad \quad \quad 0 < l \wedge \\ \quad \quad \quad \exists u \in \mathbb{R}. \\ \quad \quad \quad \quad 0 < u \wedge \\ \quad \quad \quad \quad \forall x \in \operatorname{vec}(\operatorname{ring\text{-}pow}(\operatorname{normed\text{-}field\text{-}as\text{-}commutative\text{-}ring}(k), n)). \\ \quad \quad \quad \quad \quad ((l \cdot \operatorname{u1}(x)) \leq \operatorname{u2}(x)) \wedge (\operatorname{u2}(x) \leq (u \cdot \operatorname{u1}(x))) \end{array}}}

\vnbbookfile{theorem-library/calc3-uniform}{theorem-library/calc3-uniform}
\noindent{\small calc3-uniform.scm -- calculus.pdf Prop 3.22 and Cor 3.23: UNIFORM CONTINUITY OFF A COMPACT SET.  Agent CALC-3, 2026-10-19.}\par
\vnbprop{continuous-uniform-near-compact}{\vnbstmt{\begin{array}{@{}l@{}} \forall s, t, f, k, a. \\ \quad (\operatorname{is\text{-}continuous}(s, t, f) \wedge ((k \subseteq \operatorname{pts}(s)) \wedge (\operatorname{is\text{-}compact}(\operatorname{subspace\text{-}ms}(s, k)) \wedge \operatorname{pos\text{-}rr}(a)))) \Rightarrow  \\ \quad \quad \exists b. \\ \quad \quad \quad (\operatorname{pos\text{-}rr}(b) \wedge \forall x \in k, y \in \operatorname{pts}(s).\; ((\operatorname{dist}(s)(x, y) \leq b) \Rightarrow (\operatorname{dist}(t)(f(x), f(y)) \leq a))) \end{array}}}
\vnbprop{uniformly-continuous-off-compact}{\vnbstmt{\begin{array}{@{}l@{}} \forall s, t, f, k. \\ \quad \operatorname{is\text{-}continuous}(s, t, f) \wedge \\ \quad k \subseteq \operatorname{pts}(s) \wedge \\ \quad \operatorname{is\text{-}compact}(\operatorname{subspace\text{-}ms}(s, k)) \wedge \\ \quad \operatorname{is\text{-}uniformly\text{-}continuous}(\operatorname{subspace\text{-}ms}(s, \operatorname{complement\text{-}in}(\operatorname{pts}(s), k)), t, \\ \quad \quad \operatorname{restrict}(f, \operatorname{complement\text{-}in}(\operatorname{pts}(s), k))) \Rightarrow \\ \quad \quad \operatorname{is\text{-}uniformly\text{-}continuous}(s, t, f) \end{array}}}
\vnbprop{constant-off-compact-uniformly-continuous}{\vnbstmt{\begin{array}{@{}l@{}} \forall s, t, f, k. \\ \quad \operatorname{is\text{-}continuous}(s, t, f) \wedge \\ \quad k \subseteq \operatorname{pts}(s) \wedge \\ \quad \operatorname{is\text{-}compact}(\operatorname{subspace\text{-}ms}(s, k)) \wedge \\ \quad \exists v.\; \forall x \in \operatorname{pts}(s).\; (\neg (x \in k) \Rightarrow (f(x) = v)) \Rightarrow \\ \quad \quad \operatorname{is\text{-}uniformly\text{-}continuous}(s, t, f) \end{array}}}

\vnbbookfile{theorem-library/isom-uc}{theorem-library/isom-uc}
\noindent{\small isom-uc.scm -- calculus.pdf Proposition 3.21: if X is compact, every continuous f : X --$>$ Y is uniformly continuous.  Agent ISOM, 2026-10-26.}\par
\vnbprop{carrier-is-closed}{\vnbstmt{\forall s.\; (\operatorname{is\text{-}metric\text{-}space}(s) \Rightarrow \operatorname{is\text{-}closed}(s, \operatorname{pts}(s)))}}
\vnbprop{compact-continuous-uniformly-continuous}{\vnbstmt{\begin{array}{@{}l@{}} \forall s, t, f. \\ \quad (\operatorname{is\text{-}compact}(s) \wedge \operatorname{is\text{-}continuous}(s, t, f)) \Rightarrow  \\ \quad \quad \operatorname{is\text{-}uniformly\text{-}continuous}(s, t, f) \end{array}}}

\vnbbookfile{theorem-library/isom-equicont}{theorem-library/isom-equicont}
\noindent{\small isom-equicont.scm -- calculus.pdf Proposition 3.33: on a compact metric space, pointwise and uniform convergence agree for an equicontinuous sequence of functions into ANY metric space.  Agent ISOM, 2026-10-26.}\par
\vnbprop{uniform-conv-implies-ptwise}{\vnbstmt{\begin{array}{@{}l@{}} \forall s, t, m, g. \\ \quad \operatorname{converges\text{-}uniformly\text{-}to}(s, t, m, g) \Rightarrow \operatorname{converges\text{-}ptwise\text{-}to}(s, t, m, g) \end{array}}}
\vnbprop{equicont-ptwise-implies-uniform}{\vnbstmt{\begin{array}{@{}l@{}} \forall s, t, m, g. \\ \quad (\operatorname{is\text{-}compact}(s) \wedge (\operatorname{is\text{-}equicontinuous}(s, t, m) \wedge \operatorname{converges\text{-}ptwise\text{-}to}(s, t, m, g))) \Rightarrow  \\ \quad \quad \operatorname{converges\text{-}uniformly\text{-}to}(s, t, m, g) \end{array}}}
\vnbprop{equicont-compact-ptwise-iff-uniform}{\vnbstmt{\begin{array}{@{}l@{}} \forall s, t, m, g. \\ \quad (\operatorname{is\text{-}compact}(s) \wedge \operatorname{is\text{-}equicontinuous}(s, t, m)) \Rightarrow  \\ \quad \quad \operatorname{converges\text{-}ptwise\text{-}to}(s, t, m, g) \Leftrightarrow \operatorname{converges\text{-}uniformly\text{-}to}(s, t, m, g) \end{array}}}
\vnbprop{equicontinuous-set-sequence}{\vnbstmt{\begin{array}{@{}l@{}} \forall s, t, a, m. \\ \quad (\operatorname{is\text{-}equicontinuous\text{-}set}(s, t, a) \wedge (m \in (\mathbb{N} \to a))) \Rightarrow  \\ \quad \quad \operatorname{is\text{-}equicontinuous}(s, t, m) \end{array}}}
\vnbprop{equicont-set-compact-ptwise-iff-uniform}{\vnbstmt{\begin{array}{@{}l@{}} \forall s, t, a, m, g. \\ \quad (\operatorname{is\text{-}compact}(s) \wedge (\operatorname{is\text{-}equicontinuous\text{-}set}(s, t, a) \wedge (m \in (\mathbb{N} \to a)))) \Rightarrow  \\ \quad \quad \operatorname{converges\text{-}ptwise\text{-}to}(s, t, m, g) \Leftrightarrow \operatorname{converges\text{-}uniformly\text{-}to}(s, t, m, g) \end{array}}}

\vnbbookfile{theorem-library/isom-supdist}{theorem-library/isom-supdist}
\noindent{\small isom-supdist.scm -- the sup distance between two functions into a metric space, boundedness of a compact metric space, and the group of isometries of a compact metric space as a metric space (calculus.pdf ch. 3 sec. 8, the metric (54)). Agent ISOM, 2026-10-26.}\par
\vnbprop{sup-dist-upper}{\vnbstmt{\begin{array}{@{}l@{}} \forall s, t, f, g.; \forall x \in \operatorname{pts}(s). \\ \quad (\operatorname{is\text{-}metric\text{-}space}(t) \wedge ((f \in (\operatorname{pts}(s) \to \operatorname{pts}(t))) \wedge (g \in (\operatorname{pts}(s) \to \operatorname{pts}(t))))) \Rightarrow  \\ \quad \quad \operatorname{dist}(t)(f(x), g(x)) \leq \operatorname{sup\text{-}dist}(s, t, f, g) \end{array}}}
\vnbprop{sup-dist-pos-star}{\vnbstmt{\begin{array}{@{}l@{}} \forall s, t, f, g. \\ \quad (\operatorname{is\text{-}metric\text{-}space}(t) \wedge ((f \in (\operatorname{pts}(s) \to \operatorname{pts}(t))) \wedge (g \in (\operatorname{pts}(s) \to \operatorname{pts}(t))))) \Rightarrow  \\ \quad \quad \operatorname{sup\text{-}dist}(s, t, f, g) \in \operatorname{rr\text{-}pos\text{-}star} \end{array}}}
\vnbprop{sup-dist-bound}{\vnbstmt{\begin{array}{@{}l@{}} \forall s, t, f, g, b. \\ \quad \operatorname{is\text{-}metric\text{-}space}(t) \wedge \\ \quad f \in (\operatorname{pts}(s) \to \operatorname{pts}(t)) \wedge \\ \quad g \in (\operatorname{pts}(s) \to \operatorname{pts}(t)) \wedge \\ \quad b \in \mathbb{R} \wedge \\ \quad 0 \leq b \wedge \\ \quad \forall x \in \operatorname{pts}(s).\; (\operatorname{dist}(t)(f(x), g(x)) \leq b) \Rightarrow \\ \quad \quad ((\operatorname{sup\text{-}dist}(s, t, f, g) \in \mathbb{R}) \wedge ((0 \leq \operatorname{sup\text{-}dist}(s, t, f, g)) \wedge (\operatorname{sup\text{-}dist}(s, t, f, g) \leq b))) \end{array}}}
\vnbprop{compact-dist-bound}{\vnbstmt{\begin{array}{@{}l@{}} \forall s. \\ \quad \operatorname{is\text{-}compact}(s) \Rightarrow  \\ \quad \quad \exists b \in \mathbb{R}. \\ \quad \quad \quad ((0 \leq b) \wedge \forall x, y.\; (((x \in \operatorname{pts}(s)) \wedge (y \in \operatorname{pts}(s))) \Rightarrow (\operatorname{dist}(s)(x, y) \leq b))) \end{array}}}
\vnbprop{compact-is-bounded}{\vnbstmt{\forall s.\; (\operatorname{is\text{-}compact}(s) \Rightarrow \operatorname{is\text{-}bounded\text{-}metric\text{-}space}(s))}}
\vnbprop{isometry-group-unfold}{\vnbstmt{\begin{array}{@{}l@{}} \forall s. \\ \quad (\operatorname{isometry\text{-}group}(s) \simeq \{f \in \mathrm{Bij}(\operatorname{pts}(s), \operatorname{pts}(s)) : \operatorname{is\text{-}isometry}(s, s, f)\}) \end{array}}}
\vnbprop{isometry-group-membership}{\vnbstmt{\begin{array}{@{}l@{}} \forall s, i. \\ \quad (i \in \operatorname{isometry\text{-}group}(s)) \Leftrightarrow ((i \in \mathrm{Bij}(\operatorname{pts}(s), \operatorname{pts}(s))) \wedge \operatorname{is\text{-}isometry}(s, s, i)) \end{array}}}
\vnbprop{isometry-group-in-fun}{\vnbstmt{\forall s.; \forall i \in \operatorname{isometry\text{-}group}(s).\; (i \in (\operatorname{pts}(s) \to \operatorname{pts}(s)))}}
\vnbprop{isometry-group-is-set}{\vnbstmt{\forall s.\; (\operatorname{is\text{-}metric\text{-}space}(s) \Rightarrow (\operatorname{isometry\text{-}group}(s) \in \mathbf{Set}))}}
\vnbprop{isometry-group-ms-lam-type}{\vnbstmt{\begin{array}{@{}l@{}} \forall s. \\ \quad \operatorname{is\text{-}metric\text{-}space}(s) \Rightarrow  \\ \quad \quad (\lambda {[}\operatorname{imu\_}, \operatorname{imv\_}] \in (\operatorname{isometry\text{-}group}(s) \times \operatorname{isometry\text{-}group}(s)).\; \\ \quad \quad \quad \operatorname{sup\text{-}dist}(s, s, \operatorname{imu\_}, \operatorname{imv\_})) \\ \quad \quad \in\; ((\operatorname{isometry\text{-}group}(s) \times \operatorname{isometry\text{-}group}(s)) \to \operatorname{rr\text{-}pos\text{-}star}) \end{array}}}
\vnbprop{isometry-group-ms-pts}{\vnbstmt{\begin{array}{@{}l@{}} \forall s. \\ \quad \operatorname{is\text{-}metric\text{-}space}(s) \Rightarrow (\operatorname{pts}(\operatorname{isometry\text{-}group\text{-}ms}(s)) \simeq \operatorname{isometry\text{-}group}(s)) \end{array}}}
\vnbprop{isometry-group-ms-dist-fun}{\vnbstmt{\begin{array}{@{}l@{}} \forall s. \\ \quad \operatorname{is\text{-}metric\text{-}space}(s) \Rightarrow  \\ \quad \quad (\operatorname{dist}(\operatorname{isometry\text{-}group\text{-}ms}(s)) \simeq (\lambda {[}\operatorname{imu\_}, \operatorname{imv\_}] \in (\operatorname{isometry\text{-}group}(s) \times \operatorname{isometry\text{-}group}(s)).\; \operatorname{sup\text{-}dist}(s, s, \operatorname{imu\_}, \operatorname{imv\_}))) \end{array}}}
\vnbprop{isometry-group-ms-dist}{\vnbstmt{\begin{array}{@{}l@{}} \forall s.; \forall u, v \in \operatorname{isometry\text{-}group}(s). \\ \quad \operatorname{is\text{-}metric\text{-}space}(s) \Rightarrow  \\ \quad \quad (\operatorname{dist}(\operatorname{isometry\text{-}group\text{-}ms}(s))(u, v) \simeq \operatorname{sup\text{-}dist}(s, s, u, v)) \end{array}}}
\vnbprop{isometry-group-ms-is-metric-space}{\vnbstmt{\forall s.\; (\operatorname{is\text{-}compact}(s) \Rightarrow \operatorname{is\text{-}metric\text{-}space}(\operatorname{isometry\text{-}group\text{-}ms}(s)))}}

\vnbbookfile{theorem-library/isom-group}{theorem-library/isom-group}
\noindent{\small isom-group.scm -- calculus.pdf Propositions 3.37 and 3.38: the group of isometries of a compact metric space, with the sup metric, is compact, and in it uniform and pointwise convergence agree.  Agent ISOM, 2026-10-26.}\par
\vnbprop{isometry-group-subset-fun}{\vnbstmt{\begin{array}{@{}l@{}} \forall s. \\ \quad \operatorname{is\text{-}metric\text{-}space}(s) \Rightarrow (\operatorname{isometry\text{-}group}(s) \subseteq (\operatorname{pts}(s) \to \operatorname{pts}(s))) \end{array}}}
\vnbprop{isometry-family-equicontinuous}{\vnbstmt{\begin{array}{@{}l@{}} \forall s, m. \\ \quad (\operatorname{is\text{-}metric\text{-}space}(s) \wedge (m \in (\mathbb{N} \to \operatorname{isometry\text{-}group}(s)))) \Rightarrow  \\ \quad \quad \operatorname{is\text{-}equicontinuous}(s, s, m) \end{array}}}
\vnbprop{ptwise-limit-in-isometry-group}{\vnbstmt{\begin{array}{@{}l@{}} \forall s, m, g. \\ \quad (\operatorname{is\text{-}compact}(s) \wedge ((m \in (\mathbb{N} \to \operatorname{isometry\text{-}group}(s))) \wedge \operatorname{converges\text{-}ptwise\text{-}to}(s, s, m, g))) \Rightarrow  \\ \quad \quad g \in \operatorname{isometry\text{-}group}(s) \end{array}}}
\vnbprop{ptwise-conv-iff-converges-to}{\vnbstmt{\begin{array}{@{}l@{}} \forall s, t, m, g. \\ \quad \operatorname{is\text{-}metric\text{-}space}(s) \wedge \\ \quad \operatorname{is\text{-}metric\text{-}space}(t) \wedge \\ \quad m \in (\mathbb{N} \to (\operatorname{pts}(s) \to \operatorname{pts}(t))) \wedge \\ \quad g \in (\operatorname{pts}(s) \to \operatorname{pts}(t)) \Rightarrow \\ \quad \quad \operatorname{converges\text{-}ptwise\text{-}to}(s, t, m, g) \\ \quad \quad \Leftrightarrow\; \forall x \in \operatorname{pts}(s).\; \operatorname{converges\text{-}to}(t, (\lambda k \in \mathbb{N}.\; m(k)(x)), g(x)) \end{array}}}
\vnbprop{isometry-group-conv-iff-uniform}{\vnbstmt{\begin{array}{@{}l@{}} \forall s, m, g. \\ \quad (\operatorname{is\text{-}compact}(s) \wedge ((m \in (\mathbb{N} \to \operatorname{isometry\text{-}group}(s))) \wedge (g \in \operatorname{isometry\text{-}group}(s)))) \Rightarrow  \\ \quad \quad \operatorname{converges\text{-}to}(\operatorname{isometry\text{-}group\text{-}ms}(s), m, g) \Leftrightarrow \operatorname{converges\text{-}uniformly\text{-}to}(s, s, m, g) \end{array}}}
\vnbprop{isometry-group-uniform-iff-ptwise}{\vnbstmt{\begin{array}{@{}l@{}} \forall s, a, g. \\ \quad (\operatorname{is\text{-}compact}(s) \wedge ((a \in (\mathbb{N} \to \operatorname{isometry\text{-}group}(s))) \wedge (g \in \operatorname{isometry\text{-}group}(s)))) \Rightarrow  \\ \quad \quad \operatorname{converges\text{-}to}(\operatorname{isometry\text{-}group\text{-}ms}(s), a, g) \\ \quad \quad \Leftrightarrow\; \forall x \in \operatorname{pts}(s).\; \operatorname{converges\text{-}to}(s, (\lambda k \in \mathbb{N}.\; a(k)(x)), g(x)) \end{array}}}
\vnbprop{isometry-subseq-ptwise}{\vnbstmt{\begin{array}{@{}l@{}} \forall s, f. \\ \quad (\operatorname{is\text{-}compact}(s) \wedge (f \in (\mathbb{N} \to \operatorname{isometry\text{-}group}(s)))) \Rightarrow  \\ \quad \quad \exists a. \\ \quad \quad \quad \operatorname{strictly\text{-}mono\text{-}nn}(a) \wedge \exists g.\; \operatorname{converges\text{-}ptwise\text{-}to}(s, s, \operatorname{subseq}(f, a), g) \end{array}}}
\vnbprop{isometry-group-compact}{\vnbstmt{\forall s.\; (\operatorname{is\text{-}compact}(s) \Rightarrow \operatorname{is\text{-}compact}(\operatorname{isometry\text{-}group\text{-}ms}(s)))}}

\vnbbookfile{theorem-library/isom2-group}{theorem-library/isom2-group}
\noindent{\small isom2-group.scm -- the isometries of a metric space onto itself form a GROUP under composition, and for a compact space the sup metric makes it a compact METRIC GROUP (calculus.pdf ch. 3 sec. 8, after Prop 3.37: "In fact, Gamma is a compact metric group").  Agent ISOM, 2026-10-26.}\par
\vnbprop{isometry-group-compose-in}{\vnbstmt{\begin{array}{@{}l@{}} \forall s.; \forall u, v \in \operatorname{isometry\text{-}group}(s). \\ \quad \operatorname{is\text{-}metric\text{-}space}(s) \Rightarrow ((\lambda x \in \operatorname{pts}(s).\; u(v(x))) \in \operatorname{isometry\text{-}group}(s)) \end{array}}}
\vnbprop{isometry-group-id-in}{\vnbstmt{\begin{array}{@{}l@{}} \forall s. \\ \quad \operatorname{is\text{-}metric\text{-}space}(s) \Rightarrow ((\lambda x \in \operatorname{pts}(s).\; x) \in \operatorname{isometry\text{-}group}(s)) \end{array}}}
\vnbprop{isometry-group-inverse-in}{\vnbstmt{\begin{array}{@{}l@{}} \forall s.; \forall u \in \operatorname{isometry\text{-}group}(s). \\ \quad \operatorname{is\text{-}metric\text{-}space}(s) \Rightarrow (u^{-1} \in \operatorname{isometry\text{-}group}(s)) \end{array}}}
\vnbprop{isometry-group-grp-opr-type}{\vnbstmt{\begin{array}{@{}l@{}} \forall s. \\ \quad \operatorname{is\text{-}metric\text{-}space}(s) \Rightarrow  \\ \quad \quad (\lambda {[}\operatorname{igu\_}, \operatorname{igv\_}] \in (\operatorname{isometry\text{-}group}(s) \times \operatorname{isometry\text{-}group}(s)).\; \\ \quad \quad \quad (\lambda x \in \operatorname{pts}(s).\; \operatorname{igu\_}(\operatorname{igv\_}(x)))) \\ \quad \quad \in\; ((\operatorname{isometry\text{-}group}(s) \times \operatorname{isometry\text{-}group}(s)) \to \operatorname{isometry\text{-}group}(s)) \end{array}}}
\vnbprop{isometry-group-grp-inv-type}{\vnbstmt{\begin{array}{@{}l@{}} \forall s. \\ \quad \operatorname{is\text{-}metric\text{-}space}(s) \Rightarrow  \\ \quad \quad (\lambda i \in \operatorname{isometry\text{-}group}(s).\; i^{-1}) \\ \quad \quad \in\; (\operatorname{isometry\text{-}group}(s) \to \operatorname{isometry\text{-}group}(s)) \end{array}}}
\vnbprop{isometry-group-grp-carr}{\vnbstmt{\begin{array}{@{}l@{}} \forall s. \\ \quad \operatorname{is\text{-}metric\text{-}space}(s) \Rightarrow (\operatorname{carr}(\operatorname{isometry\text{-}group\text{-}grp}(s)) \simeq \operatorname{isometry\text{-}group}(s)) \end{array}}}
\vnbprop{isometry-group-grp-opr}{\vnbstmt{\begin{array}{@{}l@{}} \forall s. \\ \quad \operatorname{is\text{-}metric\text{-}space}(s) \Rightarrow  \\ \quad \quad (\operatorname{opr}(\operatorname{isometry\text{-}group\text{-}grp}(s)) \simeq (\lambda {[}\operatorname{igu\_}, \operatorname{igv\_}] \in (\operatorname{isometry\text{-}group}(s) \times \operatorname{isometry\text{-}group}(s)).\; (\lambda x \in \operatorname{pts}(s).\; \operatorname{igu\_}(\operatorname{igv\_}(x))))) \end{array}}}
\vnbprop{isometry-group-grp-iden}{\vnbstmt{\begin{array}{@{}l@{}} \forall s. \\ \quad \operatorname{is\text{-}metric\text{-}space}(s) \Rightarrow  \\ \quad \quad (\operatorname{iden}(\operatorname{isometry\text{-}group\text{-}grp}(s)) \simeq (\lambda x \in \operatorname{pts}(s).\; x)) \end{array}}}
\vnbprop{isometry-group-grp-inv}{\vnbstmt{\begin{array}{@{}l@{}} \forall s. \\ \quad \operatorname{is\text{-}metric\text{-}space}(s) \Rightarrow  \\ \quad \quad (\operatorname{inv}(\operatorname{isometry\text{-}group\text{-}grp}(s)) \simeq (\lambda i \in \operatorname{isometry\text{-}group}(s).\; i^{-1})) \end{array}}}
\vnbprop{isometry-group-grp-is-group}{\vnbstmt{\forall s.\; (\operatorname{is\text{-}metric\text{-}space}(s) \Rightarrow \operatorname{is\text{-}group}(\operatorname{isometry\text{-}group\text{-}grp}(s)))}}
\vnbprop{isometry-group-compose-dist}{\vnbstmt{\begin{array}{@{}l@{}} \forall s.; \forall u, v, p, q \in \operatorname{isometry\text{-}group}(s). \\ \quad \operatorname{is\text{-}compact}(s) \Rightarrow  \\ \quad \quad \operatorname{sup\text{-}dist}(s, s, (\lambda x \in \operatorname{pts}(s).\; u(v(x))), (\lambda x \in \operatorname{pts}(s).\; p(q(x)))) \\ \quad \quad \leq\; \operatorname{sup\text{-}dist}(s, s, u, p) + \operatorname{sup\text{-}dist}(s, s, v, q) \end{array}}}
\vnbprop{isometry-group-inverse-dist}{\vnbstmt{\begin{array}{@{}l@{}} \forall s.; \forall u, w \in \operatorname{isometry\text{-}group}(s). \\ \quad \operatorname{is\text{-}compact}(s) \Rightarrow  \\ \quad \quad \operatorname{sup\text{-}dist}(s, s, u^{-1}, w^{-1}) \\ \quad \quad \leq\; \operatorname{sup\text{-}dist}(s, s, w, u) \end{array}}}
\vnbprop{isometry-group-is-metric-group-gm}{\vnbstmt{\begin{array}{@{}l@{}} \forall s. \\ \quad \operatorname{is\text{-}compact}(s) \Rightarrow  \\ \quad \quad \operatorname{is\text{-}group\text{-}metric}(\operatorname{isometry\text{-}group\text{-}ms}(s), \operatorname{isometry\text{-}group\text{-}grp}(s)) \end{array}}}
\vnbprop{isometry-group-compact-metric-group-gm}{\vnbstmt{\begin{array}{@{}l@{}} \forall s. \\ \quad \operatorname{is\text{-}compact}(s) \Rightarrow  \\ \quad \quad (\operatorname{is\text{-}compact}(\operatorname{isometry\text{-}group\text{-}ms}(s)) \wedge \operatorname{is\text{-}group\text{-}metric}(\operatorname{isometry\text{-}group\text{-}ms}(s), \operatorname{isometry\text{-}group\text{-}grp}(s))) \end{array}}}

\vnbbookfile{theorem-library/mg-struct}{theorem-library/mg-struct}
\noindent{\small mg-struct.scm -- the structure METRIC-GROUP and the group metrics (in the pattern of tb11/tg-struct.scm): the typing and functoriality obligations of METRIC-GROUP-MS and METRIC-GROUP-GROUP, the bridge IS-METRIC-GROUP(G) =$>$ IS-GROUP-METRIC(space, group), the way back IS-GROUP-METRIC(ms, g) =$>$ ...}\par
\vnbprop{metric-group-ms-is-metric-space}{\vnbstmt{\forall g.\; (\operatorname{is\text{-}metric\text{-}group}(g) \Rightarrow \operatorname{is\text{-}metric\text{-}space}(\operatorname{metric\text{-}group\text{-}ms}(g)))}}
\vnbprop{metric-group-group-is-group}{\vnbstmt{\forall g.\; (\operatorname{is\text{-}metric\text{-}group}(g) \Rightarrow \operatorname{is\text{-}group}(\operatorname{metric\text{-}group\text{-}group}(g)))}}
\vnbprop{metric-group-ms-functorial}{\vnbstmt{\begin{array}{@{}l@{}} \forall a, b, f. \\ \quad \operatorname{is\text{-}hom\text{-}metric\text{-}group}(a, b, f) \Rightarrow  \\ \quad \quad \operatorname{is\text{-}hom\text{-}metric\text{-}space}(\operatorname{metric\text{-}group\text{-}ms}(a), \operatorname{metric\text{-}group\text{-}ms}(b), f) \end{array}}}
\vnbprop{metric-group-group-functorial}{\vnbstmt{\begin{array}{@{}l@{}} \forall a, b, f. \\ \quad \operatorname{is\text{-}hom\text{-}metric\text{-}group}(a, b, f) \Rightarrow  \\ \quad \quad \operatorname{is\text{-}hom\text{-}group}(\operatorname{metric\text{-}group\text{-}group}(a), \operatorname{metric\text{-}group\text{-}group}(b), f) \end{array}}}
\vnbprop{metric-group-group-metric}{\vnbstmt{\begin{array}{@{}l@{}} \forall g. \\ \quad \operatorname{is\text{-}metric\text{-}group}(g) \Rightarrow  \\ \quad \quad \operatorname{is\text{-}group\text{-}metric}(\operatorname{metric\text{-}group\text{-}ms}(g), \operatorname{metric\text{-}group\text{-}group}(g)) \end{array}}}
\vnbprop{make-metric-group-carr}{\vnbstmt{\begin{array}{@{}l@{}} \forall s, g. \\ \quad (\operatorname{is\text{-}metric\text{-}space}(s) \wedge \operatorname{is\text{-}group}(g)) \Rightarrow  \\ \quad \quad (\operatorname{carr}(\operatorname{make\text{-}metric\text{-}group}(s, g)) \simeq \operatorname{carr}(g)) \end{array}}}
\vnbprop{make-metric-group-opr}{\vnbstmt{\begin{array}{@{}l@{}} \forall s, g. \\ \quad (\operatorname{is\text{-}metric\text{-}space}(s) \wedge \operatorname{is\text{-}group}(g)) \Rightarrow  \\ \quad \quad (\operatorname{opr}(\operatorname{make\text{-}metric\text{-}group}(s, g)) \simeq \operatorname{opr}(g)) \end{array}}}
\vnbprop{make-metric-group-iden}{\vnbstmt{\begin{array}{@{}l@{}} \forall s, g. \\ \quad (\operatorname{is\text{-}metric\text{-}space}(s) \wedge \operatorname{is\text{-}group}(g)) \Rightarrow  \\ \quad \quad (\operatorname{iden}(\operatorname{make\text{-}metric\text{-}group}(s, g)) \simeq \operatorname{iden}(g)) \end{array}}}
\vnbprop{make-metric-group-inv}{\vnbstmt{\begin{array}{@{}l@{}} \forall s, g. \\ \quad (\operatorname{is\text{-}metric\text{-}space}(s) \wedge \operatorname{is\text{-}group}(g)) \Rightarrow  \\ \quad \quad (\operatorname{inv}(\operatorname{make\text{-}metric\text{-}group}(s, g)) \simeq \operatorname{inv}(g)) \end{array}}}
\vnbprop{make-metric-group-gdist}{\vnbstmt{\begin{array}{@{}l@{}} \forall s, g. \\ \quad (\operatorname{is\text{-}metric\text{-}space}(s) \wedge \operatorname{is\text{-}group}(g)) \Rightarrow  \\ \quad \quad (\operatorname{gdist}(\operatorname{make\text{-}metric\text{-}group}(s, g)) \simeq \operatorname{dist}(s)) \end{array}}}
\vnbprop{make-metric-group-is-metric-group}{\vnbstmt{\begin{array}{@{}l@{}} \forall s, g. \\ \quad \operatorname{is\text{-}group\text{-}metric}(s, g) \Rightarrow \operatorname{is\text{-}metric\text{-}group}(\operatorname{make\text{-}metric\text{-}group}(s, g)) \end{array}}}
\vnbprop{isometry-group-is-metric-group}{\vnbstmt{\begin{array}{@{}l@{}} \forall s. \\ \quad \operatorname{is\text{-}compact}(s) \Rightarrow  \\ \quad \quad \operatorname{is\text{-}metric\text{-}group}(\operatorname{make\text{-}metric\text{-}group}(\operatorname{isometry\text{-}group\text{-}ms}(s), \operatorname{isometry\text{-}group\text{-}grp}(s))) \end{array}}}
\vnbprop{isometry-group-compact-metric-group}{\vnbstmt{\begin{array}{@{}l@{}} \forall s. \\ \quad \operatorname{is\text{-}compact}(s) \Rightarrow  \\ \quad \quad \operatorname{is\text{-}compact}(\operatorname{isometry\text{-}group\text{-}ms}(s)) \wedge \\ \quad \quad \operatorname{is\text{-}metric\text{-}group}(\operatorname{make\text{-}metric\text{-}group}(\operatorname{isometry\text{-}group\text{-}ms}(s), \operatorname{isometry\text{-}group\text{-}grp}(s))) \end{array}}}

\vnbbookfile{theorem-library/cg1-basic}{theorem-library/cg1-basic}
\noindent{\small cg1-basic.scm -- compact-groups notes, chapter 1 (CG-1): balls are sets in any structure with a set of points; a metric group is a topological group; Isom(X) is a topological group (Proposition 1.6).  Agent ISOM, 2026-10-26.  Helper prefix: cgb-.  Library: analysis.}\par
\vnbprop{ball-is-set-pts}{\vnbstmt{\forall s, x, r.\; ((\operatorname{pts}(s) \in \mathbf{Set}) \Rightarrow (\operatorname{ball}(s, x, r) \in \mathbf{Set}))}}
\vnbprop{metric-group-is-top-group-gt}{\vnbstmt{\begin{array}{@{}l@{}} \forall s, p. \\ \quad \operatorname{is\text{-}group\text{-}metric}(s, p) \Rightarrow \operatorname{is\text{-}group\text{-}topology}(\operatorname{metric\text{-}top}(s), p) \end{array}}}
\vnbprop{isometry-group-is-top-group-gt}{\vnbstmt{\begin{array}{@{}l@{}} \forall s. \\ \quad \operatorname{is\text{-}compact}(s) \Rightarrow  \\ \quad \quad \operatorname{is\text{-}group\text{-}topology}(\operatorname{metric\text{-}top}(\operatorname{isometry\text{-}group\text{-}ms}(s)), \operatorname{isometry\text{-}group\text{-}grp}(s)) \end{array}}}
\vnbprop{metric-group-is-top-group}{\vnbstmt{\begin{array}{@{}l@{}} \forall g. \\ \quad \operatorname{is\text{-}metric\text{-}group}(g) \Rightarrow  \\ \quad \quad \operatorname{is\text{-}top\text{-}group}(\operatorname{make\text{-}top\text{-}group}(\operatorname{metric\text{-}top}(\operatorname{metric\text{-}group\text{-}ms}(g)), \operatorname{metric\text{-}group\text{-}group}(g))) \end{array}}}
\vnbprop{isometry-group-is-top-group}{\vnbstmt{\begin{array}{@{}l@{}} \forall s. \\ \quad \operatorname{is\text{-}compact}(s) \Rightarrow  \\ \quad \quad \operatorname{is\text{-}top\text{-}group}(\operatorname{make\text{-}top\text{-}group}(\operatorname{metric\text{-}top}(\operatorname{isometry\text{-}group\text{-}ms}(s)), \operatorname{isometry\text{-}group\text{-}grp}(s))) \end{array}}}

\vnbbookfile{theorem-library/cg1-conn}{theorem-library/cg1-conn}
\noindent{\small cg1-conn.scm -- connectedness in topological spaces and the connected component of a topological group (compact-groups notes, Lemma 5.7 and Proposition 5.8). Agent ISOM, 2026-10-26.  Helper prefix: cgc-.  Library: analysis.}\par
\vnbprop{local-continuity-preimage-open}{\vnbstmt{\begin{array}{@{}l@{}} \forall s, f.; \forall w \in \operatorname{opens}(s). \\ \quad \operatorname{is\text{-}top\text{-}space}(s) \wedge \\ \quad f \in (\operatorname{pts}(s) \to \operatorname{pts}(s)) \wedge \\ \quad \forall z \in \operatorname{pts}(s), \mathit{lw} \in \operatorname{opens}(s). \\ \quad \quad (f(z) \in \mathit{lw}) \Rightarrow  \\ \quad \quad \quad \exists u \in \operatorname{opens}(s).\; ((z \in u) \wedge \forall q \in u.\; (f(q) \in \mathit{lw})) \Rightarrow \\ \quad \quad \operatorname{preimage}(s, f, w) \in \operatorname{opens}(s) \end{array}}}
\vnbprop{image-connected-t}{\vnbstmt{\begin{array}{@{}l@{}} \forall s, f, a. \\ \quad \operatorname{is\text{-}top\text{-}space}(s) \wedge \\ \quad f \in (\operatorname{pts}(s) \to \operatorname{pts}(s)) \wedge \\ \quad \forall z \in \operatorname{pts}(s), w \in \operatorname{opens}(s). \\ \quad \quad (f(z) \in w) \Rightarrow  \\ \quad \quad \quad \exists u \in \operatorname{opens}(s).\; ((z \in u) \wedge \forall q \in u.\; (f(q) \in w)) \wedge \\ \quad \operatorname{is\text{-}connected\text{-}t}(s, a) \Rightarrow \\ \quad \quad \operatorname{is\text{-}connected\text{-}t}(s, \operatorname{image}(f, a)) \end{array}}}
\vnbprop{union-connected-t}{\vnbstmt{\begin{array}{@{}l@{}} \forall s, a, b, p. \\ \quad (\operatorname{is\text{-}top\text{-}space}(s) \wedge (\operatorname{is\text{-}connected\text{-}t}(s, a) \wedge (\operatorname{is\text{-}connected\text{-}t}(s, b) \wedge ((p \in a) \wedge (p \in b))))) \Rightarrow  \\ \quad \quad \operatorname{is\text{-}connected\text{-}t}(s, (a \cup b)) \end{array}}}
\vnbprop{connected-component-unfold}{\vnbstmt{\begin{array}{@{}l@{}} \forall s, x. \\ \quad (\operatorname{connected\text{-}component}(s, x) \simeq \operatorname{big\text{-}union}(c, \{c \in \operatorname{power}(\operatorname{pts}(s)) : (\operatorname{is\text{-}connected\text{-}t}(s, c) \wedge (x \in c))\}, c)) \end{array}}}
\vnbprop{connected-component-elim}{\vnbstmt{\begin{array}{@{}l@{}} \forall s, x.; \forall y \in \operatorname{connected\text{-}component}(s, x). \exists c. \\ \quad \operatorname{is\text{-}connected\text{-}t}(s, c) \wedge ((x \in c) \wedge (y \in c)) \end{array}}}
\vnbprop{connected-component-intro}{\vnbstmt{\begin{array}{@{}l@{}} \forall s, c, x, y. \\ \quad (\operatorname{is\text{-}top\text{-}space}(s) \wedge (\operatorname{is\text{-}connected\text{-}t}(s, c) \wedge ((x \in c) \wedge (y \in c)))) \Rightarrow  \\ \quad \quad y \in \operatorname{connected\text{-}component}(s, x) \end{array}}}
\vnbprop{singleton-connected-t}{\vnbstmt{\begin{array}{@{}l@{}} \forall s.; \forall x \in \operatorname{pts}(s). \\ \quad \operatorname{is\text{-}top\text{-}space}(s) \Rightarrow \operatorname{is\text{-}connected\text{-}t}(s, \operatorname{singleton}(x)) \end{array}}}
\vnbprop{connected-component-self}{\vnbstmt{\forall s.; \forall x \in \operatorname{pts}(s).\; (\operatorname{is\text{-}top\text{-}space}(s) \Rightarrow (x \in \operatorname{connected\text{-}component}(s, x)))}}
\vnbprop{connected-component-connected}{\vnbstmt{\begin{array}{@{}l@{}} \forall s, x. \\ \quad \operatorname{is\text{-}top\text{-}space}(s) \Rightarrow \operatorname{is\text{-}connected\text{-}t}(s, \operatorname{connected\text{-}component}(s, x)) \end{array}}}
\vnbprop{connected-component-self-map}{\vnbstmt{\begin{array}{@{}l@{}} \forall s, f.; \forall x \in \operatorname{pts}(s), y \in \operatorname{connected\text{-}component}(s, x). \\ \quad \operatorname{is\text{-}top\text{-}space}(s) \wedge \\ \quad f \in (\operatorname{pts}(s) \to \operatorname{pts}(s)) \wedge \\ \quad \forall z \in \operatorname{pts}(s), w \in \operatorname{opens}(s). \\ \quad \quad (f(z) \in w) \Rightarrow  \\ \quad \quad \quad \exists u \in \operatorname{opens}(s).\; ((z \in u) \wedge \forall q \in u.\; (f(q) \in w)) \wedge \\ \quad f(x) \in \operatorname{connected\text{-}component}(s, x) \Rightarrow \\ \quad \quad f(y) \in \operatorname{connected\text{-}component}(s, x) \end{array}}}

\vnbbookfile{theorem-library/cg1-group}{theorem-library/cg1-group}
\noindent{\small cg1-group.scm -- the connected component of the identity of a topological group is a normal subgroup (compact-groups notes, Lemma 5.7, the algebraic half; "closed" needs CLOSURE). Inversion, left translation and conjugation are locally continuous self-maps; each maps the component into itself ...}\par
\vnbprop{top-group-inversion-locally-continuous-gt}{\vnbstmt{\begin{array}{@{}l@{}} \forall s, g.; \forall f \in (\operatorname{pts}(s) \to \operatorname{pts}(s)), z \in \operatorname{pts}(s), w \in \operatorname{opens}(s). \\ \quad (\operatorname{is\text{-}group\text{-}topology}(s, g) \wedge (\forall \mathit{vw} \in \operatorname{pts}(s).\; (f(\mathit{vw}) = \operatorname{inv}(g)(\mathit{vw})) \wedge (f(z) \in w))) \Rightarrow  \\ \quad \quad \exists u \in \operatorname{opens}(s).\; ((z \in u) \wedge \forall q \in u.\; (f(q) \in w)) \end{array}}}
\vnbprop{top-group-inversion-locally-continuous}{\vnbstmt{\begin{array}{@{}l@{}} \forall g.; \forall f \in (\operatorname{pts}(\operatorname{top\text{-}group\text{-}space}(g)) \to \operatorname{pts}(\operatorname{top\text{-}group\text{-}space}(g))), z \in \operatorname{pts}(\operatorname{top\text{-}group\text{-}space}(g)), w \in \operatorname{opens}(\operatorname{top\text{-}group\text{-}space}(g)). \\ \quad \operatorname{is\text{-}top\text{-}group}(g) \wedge \\ \quad \forall \mathit{vw} \in \operatorname{pts}(\operatorname{top\text{-}group\text{-}space}(g)).\; (f(\mathit{vw}) = \operatorname{inv}(\operatorname{top\text{-}group\text{-}group}(g))(\mathit{vw})) \wedge \\ \quad f(z) \in w \Rightarrow \\ \quad \quad \exists u \in \operatorname{opens}(\operatorname{top\text{-}group\text{-}space}(g)). \\ \quad \quad \quad (z \in u) \wedge \forall q \in u.\; (f(q) \in w) \end{array}}}
\vnbprop{top-group-left-translation-locally-continuous-gt}{\vnbstmt{\begin{array}{@{}l@{}} \forall s, g.; \forall a \in \operatorname{carr}(g), f \in (\operatorname{pts}(s) \to \operatorname{pts}(s)), z \in \operatorname{pts}(s), w \in \operatorname{opens}(s). \\ \quad (\operatorname{is\text{-}group\text{-}topology}(s, g) \wedge (\forall \mathit{vw} \in \operatorname{pts}(s).\; (f(\mathit{vw}) = \operatorname{opr}(g)(a, \mathit{vw})) \wedge (f(z) \in w))) \Rightarrow  \\ \quad \quad \exists u \in \operatorname{opens}(s).\; ((z \in u) \wedge \forall q \in u.\; (f(q) \in w)) \end{array}}}
\vnbprop{top-group-left-translation-locally-continuous}{\vnbstmt{\begin{array}{@{}l@{}} \forall g.; \forall a \in \operatorname{carr}(\operatorname{top\text{-}group\text{-}group}(g)), f \in (\operatorname{pts}(\operatorname{top\text{-}group\text{-}space}(g)) \to \operatorname{pts}(\operatorname{top\text{-}group\text{-}space}(g))), z \in \operatorname{pts}(\operatorname{top\text{-}group\text{-}space}(g)), w \in \operatorname{opens}(\operatorname{top\text{-}group\text{-}space}(g)). \\ \quad \operatorname{is\text{-}top\text{-}group}(g) \wedge \\ \quad \forall \mathit{vw} \in \operatorname{pts}(\operatorname{top\text{-}group\text{-}space}(g)). \\ \quad \quad f(\mathit{vw}) = \operatorname{opr}(\operatorname{top\text{-}group\text{-}group}(g))(a, \mathit{vw}) \wedge \\ \quad f(z) \in w \Rightarrow \\ \quad \quad \exists u \in \operatorname{opens}(\operatorname{top\text{-}group\text{-}space}(g)). \\ \quad \quad \quad (z \in u) \wedge \forall q \in u.\; (f(q) \in w) \end{array}}}
\vnbprop{top-group-conjugation-locally-continuous-gt}{\vnbstmt{\begin{array}{@{}l@{}} \forall s, g.; \forall u \in \operatorname{carr}(g), f \in (\operatorname{pts}(s) \to \operatorname{pts}(s)), z \in \operatorname{pts}(s), w \in \operatorname{opens}(s). \\ \quad \operatorname{is\text{-}group\text{-}topology}(s, g) \wedge \\ \quad \forall \mathit{vw} \in \operatorname{pts}(s).\; (f(\mathit{vw}) = \operatorname{opr}(g)(\operatorname{opr}(g)(u, \mathit{vw}), \operatorname{inv}(g)(u))) \wedge \\ \quad f(z) \in w \Rightarrow \\ \quad \quad \exists \mathit{lu} \in \operatorname{opens}(s).\; ((z \in \mathit{lu}) \wedge \forall q \in \mathit{lu}.\; (f(q) \in w)) \end{array}}}
\vnbprop{top-group-conjugation-locally-continuous}{\vnbstmt{\begin{array}{@{}l@{}} \forall g.; \forall u \in \operatorname{carr}(\operatorname{top\text{-}group\text{-}group}(g)), f \in (\operatorname{pts}(\operatorname{top\text{-}group\text{-}space}(g)) \to \operatorname{pts}(\operatorname{top\text{-}group\text{-}space}(g))), z \in \operatorname{pts}(\operatorname{top\text{-}group\text{-}space}(g)), w \in \operatorname{opens}(\operatorname{top\text{-}group\text{-}space}(g)). \\ \quad \operatorname{is\text{-}top\text{-}group}(g) \wedge \\ \quad \forall \mathit{vw} \in \operatorname{pts}(\operatorname{top\text{-}group\text{-}space}(g)). \\ \quad \quad f(\mathit{vw}) = \operatorname{opr}(\operatorname{top\text{-}group\text{-}group}(g))(\operatorname{opr}(\operatorname{top\text{-}group\text{-}group}(g))(u, \mathit{vw}), \operatorname{inv}(\operatorname{top\text{-}group\text{-}group}(g))(u)) \wedge \\ \quad f(z) \in w \Rightarrow \\ \quad \quad \exists \mathit{lu} \in \operatorname{opens}(\operatorname{top\text{-}group\text{-}space}(g)). \\ \quad \quad \quad (z \in \mathit{lu}) \wedge \forall q \in \mathit{lu}.\; (f(q) \in w) \end{array}}}
\vnbprop{top-group-component-normal-subgroup-gt}{\vnbstmt{\begin{array}{@{}l@{}} \forall s, g. \\ \quad \operatorname{is\text{-}group\text{-}topology}(s, g) \Rightarrow  \\ \quad \quad \operatorname{is\text{-}normal\text{-}subgroup}(g, \operatorname{connected\text{-}component}(s, \operatorname{iden}(g))) \end{array}}}
\vnbprop{top-group-component-normal-subgroup}{\vnbstmt{\begin{array}{@{}l@{}} \forall g. \\ \quad \operatorname{is\text{-}top\text{-}group}(g) \Rightarrow  \\ \quad \quad \operatorname{is\text{-}normal\text{-}subgroup}(\operatorname{top\text{-}group\text{-}group}(g), \\ \quad \quad \quad \operatorname{connected\text{-}component}(\operatorname{top\text{-}group\text{-}space}(g), \operatorname{iden}(\operatorname{top\text{-}group\text{-}group}(g)))) \end{array}}}
\vnbprop{top-group-open-subgroup-complement-open-gt}{\vnbstmt{\begin{array}{@{}l@{}} \forall s, g, h. \\ \quad (\operatorname{is\text{-}group\text{-}topology}(s, g) \wedge (\operatorname{is\text{-}subgroup}(g, h) \wedge (h \in \operatorname{opens}(s)))) \Rightarrow  \\ \quad \quad \{z \in \operatorname{pts}(s) : \neg (z \in h)\} \in \operatorname{opens}(s) \end{array}}}
\vnbprop{top-group-open-subgroup-complement-open}{\vnbstmt{\begin{array}{@{}l@{}} \forall g, h. \\ \quad (\operatorname{is\text{-}top\text{-}group}(g) \wedge (\operatorname{is\text{-}subgroup}(\operatorname{top\text{-}group\text{-}group}(g), h) \wedge (h \in \operatorname{opens}(\operatorname{top\text{-}group\text{-}space}(g))))) \Rightarrow  \\ \quad \quad \{z \in \operatorname{pts}(\operatorname{top\text{-}group\text{-}space}(g)) : \neg (z \in h)\} \in \operatorname{opens}(\operatorname{top\text{-}group\text{-}space}(g)) \end{array}}}
\vnbprop{top-group-open-connected-subgroup-is-component-gt}{\vnbstmt{\begin{array}{@{}l@{}} \forall s, g, h. \\ \quad (\operatorname{is\text{-}group\text{-}topology}(s, g) \wedge (\operatorname{is\text{-}subgroup}(g, h) \wedge ((h \in \operatorname{opens}(s)) \wedge \operatorname{is\text{-}connected\text{-}t}(s, h)))) \Rightarrow  \\ \quad \quad h = \operatorname{connected\text{-}component}(s, \operatorname{iden}(g)) \end{array}}}
\vnbprop{top-group-open-connected-subgroup-is-component}{\vnbstmt{\begin{array}{@{}l@{}} \forall g, h. \\ \quad \operatorname{is\text{-}top\text{-}group}(g) \wedge \\ \quad \operatorname{is\text{-}subgroup}(\operatorname{top\text{-}group\text{-}group}(g), h) \wedge \\ \quad h \in \operatorname{opens}(\operatorname{top\text{-}group\text{-}space}(g)) \wedge \\ \quad \operatorname{is\text{-}connected\text{-}t}(\operatorname{top\text{-}group\text{-}space}(g), h) \Rightarrow \\ \quad \quad h = \operatorname{connected\text{-}component}(\operatorname{top\text{-}group\text{-}space}(g), \operatorname{iden}(\operatorname{top\text{-}group\text{-}group}(g))) \end{array}}}

\vnbbookfile{theorem-library/cg2-top}{theorem-library/cg2-top}
\noindent{\small cg2-top.scm -- Proposition 5.8 of the compact-groups notes: for an open V containing e with V = V\textasciicircum{}-1, the union H of the powers V\textasciicircum{}k is an open subgroup; when V is also connected, H is the connected component of e (and so connected).  Each result is proved for a group topology (s, g) under NAME-gt and restated for the ...}\par
\vnbprop{set-prod-open}{\vnbstmt{\begin{array}{@{}l@{}} \forall s, g, a, v. \\ \quad (\operatorname{is\text{-}group\text{-}topology}(s, g) \wedge ((a \subseteq \operatorname{carr}(g)) \wedge (v \in \operatorname{opens}(s)))) \Rightarrow  \\ \quad \quad \operatorname{set\text{-}prod}(g, a, v) \in \operatorname{opens}(s) \end{array}}}
\vnbprop{powers-union-open}{\vnbstmt{\begin{array}{@{}l@{}} \forall s, g.; \forall v \in \operatorname{opens}(s). \\ \quad (\operatorname{is\text{-}group\text{-}topology}(s, g) \wedge (\operatorname{iden}(g) \in v)) \Rightarrow (\operatorname{powers\text{-}union}(g, v) \in \operatorname{opens}(s)) \end{array}}}
\vnbprop{open-subgroup-contains-component-gt}{\vnbstmt{\begin{array}{@{}l@{}} \forall s, g, h. \\ \quad (\operatorname{is\text{-}group\text{-}topology}(s, g) \wedge (\operatorname{is\text{-}subgroup}(g, h) \wedge (h \in \operatorname{opens}(s)))) \Rightarrow  \\ \quad \quad \operatorname{connected\text{-}component}(s, \operatorname{iden}(g)) \subseteq h \end{array}}}
\vnbprop{open-subgroup-contains-component}{\vnbstmt{\begin{array}{@{}l@{}} \forall g, h. \\ \quad (\operatorname{is\text{-}top\text{-}group}(g) \wedge (\operatorname{is\text{-}subgroup}(\operatorname{top\text{-}group\text{-}group}(g), h) \wedge (h \in \operatorname{opens}(\operatorname{top\text{-}group\text{-}space}(g))))) \Rightarrow  \\ \quad \quad \operatorname{connected\text{-}component}(\operatorname{top\text{-}group\text{-}space}(g), \operatorname{iden}(\operatorname{top\text{-}group\text{-}group}(g))) \subseteq h \end{array}}}
\vnbprop{set-powers-in-component}{\vnbstmt{\begin{array}{@{}l@{}} \forall k \in \mathbb{N}, s, g, v. \\ \quad (\operatorname{is\text{-}group\text{-}topology}(s, g) \wedge (v \subseteq \operatorname{connected\text{-}component}(s, \operatorname{iden}(g)))) \Rightarrow  \\ \quad \quad \operatorname{set\text{-}power}(g, v, k) \subseteq \operatorname{connected\text{-}component}(s, \operatorname{iden}(g)) \end{array}}}
\vnbprop{powers-union-in-component}{\vnbstmt{\begin{array}{@{}l@{}} \forall s, g.; \forall v \in \operatorname{opens}(s). \\ \quad (\operatorname{is\text{-}group\text{-}topology}(s, g) \wedge ((\operatorname{iden}(g) \in v) \wedge \operatorname{is\text{-}connected\text{-}t}(s, v))) \Rightarrow  \\ \quad \quad \operatorname{powers\text{-}union}(g, v) \subseteq \operatorname{connected\text{-}component}(s, \operatorname{iden}(g)) \end{array}}}
\vnbprop{top-group-powers-open-subgroup-gt}{\vnbstmt{\begin{array}{@{}l@{}} \forall s, g.; \forall v \in \operatorname{opens}(s). \\ \quad (\operatorname{is\text{-}group\text{-}topology}(s, g) \wedge ((\operatorname{iden}(g) \in v) \wedge \forall p \in v.\; (\operatorname{inv}(g)(p) \in v))) \Rightarrow  \\ \quad \quad \operatorname{is\text{-}subgroup}(g, \operatorname{powers\text{-}union}(g, v)) \wedge (\operatorname{powers\text{-}union}(g, v) \in \operatorname{opens}(s)) \end{array}}}
\vnbprop{top-group-powers-open-subgroup}{\vnbstmt{\begin{array}{@{}l@{}} \forall g.; \forall v \in \operatorname{opens}(\operatorname{top\text{-}group\text{-}space}(g)). \\ \quad \operatorname{is\text{-}top\text{-}group}(g) \wedge \\ \quad \operatorname{iden}(\operatorname{top\text{-}group\text{-}group}(g)) \in v \wedge \\ \quad \forall p \in v.\; (\operatorname{inv}(\operatorname{top\text{-}group\text{-}group}(g))(p) \in v) \Rightarrow \\ \quad \quad \operatorname{is\text{-}subgroup}(\operatorname{top\text{-}group\text{-}group}(g), \operatorname{powers\text{-}union}(\operatorname{top\text{-}group\text{-}group}(g), v)) \wedge \\ \quad \quad \operatorname{powers\text{-}union}(\operatorname{top\text{-}group\text{-}group}(g), v) \in \operatorname{opens}(\operatorname{top\text{-}group\text{-}space}(g)) \end{array}}}
\vnbprop{top-group-powers-component-gt}{\vnbstmt{\begin{array}{@{}l@{}} \forall s, g.; \forall v \in \operatorname{opens}(s). \\ \quad \operatorname{is\text{-}group\text{-}topology}(s, g) \wedge \\ \quad \operatorname{iden}(g) \in v \wedge \\ \quad \forall p \in v.\; (\operatorname{inv}(g)(p) \in v) \wedge \\ \quad \operatorname{is\text{-}connected\text{-}t}(s, v) \Rightarrow \\ \quad \quad ((\operatorname{powers\text{-}union}(g, v) = \operatorname{connected\text{-}component}(s, \operatorname{iden}(g))) \wedge \operatorname{is\text{-}connected\text{-}t}(s, \operatorname{powers\text{-}union}(g, v))) \end{array}}}
\vnbprop{top-group-powers-component}{\vnbstmt{\begin{array}{@{}l@{}} \forall g.; \forall v \in \operatorname{opens}(\operatorname{top\text{-}group\text{-}space}(g)). \\ \quad \operatorname{is\text{-}top\text{-}group}(g) \wedge \\ \quad \operatorname{iden}(\operatorname{top\text{-}group\text{-}group}(g)) \in v \wedge \\ \quad \forall p \in v.\; (\operatorname{inv}(\operatorname{top\text{-}group\text{-}group}(g))(p) \in v) \wedge \\ \quad \operatorname{is\text{-}connected\text{-}t}(\operatorname{top\text{-}group\text{-}space}(g), v) \Rightarrow \\ \quad \quad \operatorname{powers\text{-}union}(\operatorname{top\text{-}group\text{-}group}(g), v) \\ \quad \quad =\; \operatorname{connected\text{-}component}(\operatorname{top\text{-}group\text{-}space}(g), \operatorname{iden}(\operatorname{top\text{-}group\text{-}group}(g))) \wedge \\ \quad \quad \operatorname{is\text{-}connected\text{-}t}(\operatorname{top\text{-}group\text{-}space}(g), \operatorname{powers\text{-}union}(\operatorname{top\text{-}group\text{-}group}(g), v)) \end{array}}}

\vnbbookfile{theorem-library/cg2-compact}{theorem-library/cg2-compact}
\noindent{\small cg2-compact.scm -- Corollary 5.9 of the compact-groups notes: in a COMPACT topological group, for an open V containing e with V = V\textasciicircum{}-1, the powers V\textasciicircum{}k stabilise: H = V\textasciicircum{}m for some m; when V is also connected, V\textasciicircum{}m is the connected component of e.  The notes' proof ("the sequence V\textasciicircum{}k is non-decreasing; by compactness of H ...}\par
\vnbprop{set-power-mono-add}{\vnbstmt{\begin{array}{@{}l@{}} \forall d \in \mathbb{N}, g, v, k. \\ \quad (\operatorname{is\text{-}group}(g) \wedge ((v \subseteq \operatorname{carr}(g)) \wedge ((\operatorname{iden}(g) \in v) \wedge (k \in \mathbb{N})))) \Rightarrow  \\ \quad \quad \operatorname{set\text{-}power}(g, v, k) \subseteq \operatorname{set\text{-}power}(g, v, (k + d)) \end{array}}}
\vnbprop{set-power-succ-open}{\vnbstmt{\begin{array}{@{}l@{}} \forall s, g, v, k. \\ \quad (\operatorname{is\text{-}group\text{-}topology}(s, g) \wedge ((v \in \operatorname{opens}(s)) \wedge (k \in \mathbb{N}))) \Rightarrow  \\ \quad \quad \operatorname{set\text{-}power}(g, v, (k + 1)) \in \operatorname{opens}(s) \end{array}}}
\vnbprop{top-group-compact-powers-stabilize-gt}{\vnbstmt{\begin{array}{@{}l@{}} \forall s, g.; \forall v \in \operatorname{opens}(s). \\ \quad (\operatorname{is\text{-}group\text{-}topology}(s, g) \wedge (\operatorname{is\text{-}compact\text{-}t}(s) \wedge ((\operatorname{iden}(g) \in v) \wedge \forall p \in v.\; (\operatorname{inv}(g)(p) \in v)))) \Rightarrow  \\ \quad \quad \exists m \in \mathbb{N}.\; (\operatorname{set\text{-}power}(g, v, m) = \operatorname{powers\text{-}union}(g, v)) \end{array}}}
\vnbprop{top-group-compact-powers-stabilize}{\vnbstmt{\begin{array}{@{}l@{}} \forall g.; \forall v \in \operatorname{opens}(\operatorname{top\text{-}group\text{-}space}(g)). \\ \quad \operatorname{is\text{-}top\text{-}group}(g) \wedge \\ \quad \operatorname{is\text{-}compact\text{-}t}(\operatorname{top\text{-}group\text{-}space}(g)) \wedge \\ \quad \operatorname{iden}(\operatorname{top\text{-}group\text{-}group}(g)) \in v \wedge \\ \quad \forall p \in v.\; (\operatorname{inv}(\operatorname{top\text{-}group\text{-}group}(g))(p) \in v) \Rightarrow \\ \quad \quad \exists m \in \mathbb{N}. \\ \quad \quad \quad \operatorname{set\text{-}power}(\operatorname{top\text{-}group\text{-}group}(g), v, m) = \operatorname{powers\text{-}union}(\operatorname{top\text{-}group\text{-}group}(g), v) \end{array}}}
\vnbprop{top-group-compact-component-power-gt}{\vnbstmt{\begin{array}{@{}l@{}} \forall s, g.; \forall v \in \operatorname{opens}(s). \\ \quad \operatorname{is\text{-}group\text{-}topology}(s, g) \wedge \\ \quad \operatorname{is\text{-}compact\text{-}t}(s) \wedge \\ \quad \operatorname{iden}(g) \in v \wedge \\ \quad \forall p \in v.\; (\operatorname{inv}(g)(p) \in v) \wedge \\ \quad \operatorname{is\text{-}connected\text{-}t}(s, v) \Rightarrow \\ \quad \quad \exists m \in \mathbb{N}.\; (\operatorname{set\text{-}power}(g, v, m) = \operatorname{connected\text{-}component}(s, \operatorname{iden}(g))) \end{array}}}
\vnbprop{top-group-compact-component-power}{\vnbstmt{\begin{array}{@{}l@{}} \forall g.; \forall v \in \operatorname{opens}(\operatorname{top\text{-}group\text{-}space}(g)). \\ \quad \operatorname{is\text{-}top\text{-}group}(g) \wedge \\ \quad \operatorname{is\text{-}compact\text{-}t}(\operatorname{top\text{-}group\text{-}space}(g)) \wedge \\ \quad \operatorname{iden}(\operatorname{top\text{-}group\text{-}group}(g)) \in v \wedge \\ \quad \forall p \in v.\; (\operatorname{inv}(\operatorname{top\text{-}group\text{-}group}(g))(p) \in v) \wedge \\ \quad \operatorname{is\text{-}connected\text{-}t}(\operatorname{top\text{-}group\text{-}space}(g), v) \Rightarrow \\ \quad \quad \exists m \in \mathbb{N}. \\ \quad \quad \quad \operatorname{set\text{-}power}(\operatorname{top\text{-}group\text{-}group}(g), v, m) \\ \quad \quad \quad =\; \operatorname{connected\text{-}component}(\operatorname{top\text{-}group\text{-}space}(g), \operatorname{iden}(\operatorname{top\text{-}group\text{-}group}(g))) \end{array}}}

\vnbbookfile{theorem-library/cvx-tb}{theorem-library/cvx-tb}
\noindent{\small cvx-tb.scm -- Lemma 1.12 of the compact-groups notes: the convex hull of a totally bounded subset of a real normed vector space is totally bounded.  Route: the join of a totally bounded set to a point is totally bounded (a net of the set times a grid of [0, 1]); the hull of a finite set is totally bounded (induction on ...}\par
\vnbprop{is-totally-bounded-in-classes}{\vnbstmt{\begin{array}{@{}l@{}} \forall m, a. \\ \quad \operatorname{is\text{-}totally\text{-}bounded\text{-}in}(m, a) \\ \quad \Leftrightarrow\; \forall r. \\ \quad \quad \operatorname{pos\text{-}rr}(r) \Rightarrow  \\ \quad \quad \quad \exists f. \\ \quad \quad \quad \quad ((|f| \in \mathbb{N}) \wedge ((f \subseteq \operatorname{pts}(m)) \wedge \forall p \in a.\; \exists c \in f.\; (\operatorname{dist}(m)(c, p) < r))) \end{array}}}
\vnbprop{convex-combination-shift-point}{\vnbstmt{\begin{array}{@{}l@{}} \forall m, x, z, a, t. \\ \quad (\operatorname{is\text{-}normed\text{-}vector\text{-}space}(m) \wedge ((x \in \operatorname{vec}(m)) \wedge ((z \in \operatorname{vec}(m)) \wedge ((a \in \operatorname{vec}(m)) \wedge (t \in \mathbb{R}))))) \Rightarrow  \\ \quad \quad \operatorname{vadd}(m)(\operatorname{act}(m)((1 - t), z), \operatorname{act}(m)(t, a)) \\ \quad \quad =\; \operatorname{vadd}(m)(\operatorname{vadd}(m)(\operatorname{act}(m)((1 - t), x), \operatorname{act}(m)(t, a)), \\ \quad \quad \quad \operatorname{act}(m)((1 - t), \operatorname{vadd}(m)(z, \operatorname{vneg}(m)(x)))) \end{array}}}
\vnbprop{convex-combination-shift-coefficient}{\vnbstmt{\begin{array}{@{}l@{}} \forall m, f, a, t, u. \\ \quad (\operatorname{is\text{-}normed\text{-}vector\text{-}space}(m) \wedge ((f \in \operatorname{vec}(m)) \wedge ((a \in \operatorname{vec}(m)) \wedge ((t \in \mathbb{R}) \wedge (u \in \mathbb{R}))))) \Rightarrow  \\ \quad \quad \operatorname{vadd}(m)(\operatorname{act}(m)((1 - u), f), \operatorname{act}(m)(u, a)) \\ \quad \quad =\; \operatorname{vadd}(m)(\operatorname{vadd}(m)(\operatorname{act}(m)((1 - t), f), \operatorname{act}(m)(t, a)), \\ \quad \quad \quad \operatorname{act}(m)((u - t), \operatorname{vadd}(m)(a, \operatorname{vneg}(m)(f)))) \end{array}}}
\vnbprop{convex-combination-dist-point}{\vnbstmt{\begin{array}{@{}l@{}} \forall m, x, z, a, t. \\ \quad (\operatorname{is\text{-}normed\text{-}vector\text{-}space}(m) \wedge ((x \in \operatorname{vec}(m)) \wedge ((z \in \operatorname{vec}(m)) \wedge ((a \in \operatorname{vec}(m)) \wedge (t \in \mathbb{R}))))) \Rightarrow  \\ \quad \quad \operatorname{dist}(\operatorname{nvs\text{-}metric\text{-}space}(m))(\operatorname{vadd}(m)(\operatorname{act}(m)((1 - t), x), \operatorname{act}(m)(t, a)), \\ \quad \quad \quad \operatorname{vadd}(m)(\operatorname{act}(m)((1 - t), z), \operatorname{act}(m)(t, a))) \\ \quad \quad =\; \lvert (1 - t) \rvert \cdot \operatorname{dist}(\operatorname{nvs\text{-}metric\text{-}space}(m))(z, x) \end{array}}}
\vnbprop{convex-combination-dist-coefficient}{\vnbstmt{\begin{array}{@{}l@{}} \forall m, f, a, t, u. \\ \quad (\operatorname{is\text{-}normed\text{-}vector\text{-}space}(m) \wedge ((f \in \operatorname{vec}(m)) \wedge ((a \in \operatorname{vec}(m)) \wedge ((t \in \mathbb{R}) \wedge (u \in \mathbb{R}))))) \Rightarrow  \\ \quad \quad \operatorname{dist}(\operatorname{nvs\text{-}metric\text{-}space}(m))(\operatorname{vadd}(m)(\operatorname{act}(m)((1 - t), f), \operatorname{act}(m)(t, a)), \\ \quad \quad \quad \operatorname{vadd}(m)(\operatorname{act}(m)((1 - u), f), \operatorname{act}(m)(u, a))) \\ \quad \quad =\; \lvert (u - t) \rvert \cdot \operatorname{dist}(\operatorname{nvs\text{-}metric\text{-}space}(m))(a, f) \end{array}}}
\vnbprop{finite-set-bounded-distance}{\vnbstmt{\begin{array}{@{}l@{}} \forall m, a, f. \\ \quad (\operatorname{is\text{-}normed\text{-}vector\text{-}space}(m) \wedge ((a \in \operatorname{vec}(m)) \wedge ((f \in \mathbf{Set}) \wedge ((|f| \in \mathbb{N}) \wedge (f \subseteq \operatorname{vec}(m)))))) \Rightarrow  \\ \quad \quad \exists d \in \mathbb{R}. \\ \quad \quad \quad (0 \leq d) \wedge \forall \mathit{bdf} \in f.\; (\operatorname{dist}(\operatorname{nvs\text{-}metric\text{-}space}(m))(a, \mathit{bdf}) \leq d) \end{array}}}
\vnbprop{join-net-estimate}{\vnbstmt{\begin{array}{@{}l@{}} \forall m, y, f, a, t, u, o, l, d. \\ \quad \operatorname{is\text{-}normed\text{-}vector\text{-}space}(m) \wedge \\ \quad y \in \operatorname{vec}(m) \wedge \\ \quad f \in \operatorname{vec}(m) \wedge \\ \quad a \in \operatorname{vec}(m) \wedge \\ \quad t \in \mathbb{R} \wedge \\ \quad 0 \leq t \wedge \\ \quad t \leq 1 \wedge \\ \quad u \in \mathbb{R} \wedge \\ \quad u \leq t \wedge \\ \quad t < (u + l) \wedge \\ \quad o \in \mathbb{R} \wedge \\ \quad l \in \mathbb{R} \wedge \\ \quad d \in \mathbb{R} \wedge \\ \quad 0 \leq l \wedge \\ \quad \operatorname{dist}(\operatorname{nvs\text{-}metric\text{-}space}(m))(f, y) < o \wedge \\ \quad \operatorname{dist}(\operatorname{nvs\text{-}metric\text{-}space}(m))(a, f) \leq d \Rightarrow \\ \quad \quad \operatorname{dist}(\operatorname{nvs\text{-}metric\text{-}space}(m))(\operatorname{vadd}(m)(\operatorname{act}(m)((1 - u), f), \operatorname{act}(m)(u, a)), \\ \quad \quad \quad \operatorname{vadd}(m)(\operatorname{act}(m)((1 - t), y), \operatorname{act}(m)(t, a))) \\ \quad \quad <\; o + (l \cdot d) \end{array}}}
\vnbprop{unit-grid-approx}{\vnbstmt{\begin{array}{@{}l@{}} \forall l, k, t. \\ \quad ((l \in \mathbb{R}) \wedge (\operatorname{pos\text{-}rr}(l) \wedge ((k \in \mathbb{N}) \wedge ((1 < ((k \cdot l) + l)) \wedge ((t \in \mathbb{R}) \wedge ((0 \leq t) \wedge (t \leq 1))))))) \Rightarrow  \\ \quad \quad \exists i \in \mathrm{OS}((k + 1)). \\ \quad \quad \quad ((i \cdot l) \in \mathbb{R}) \wedge (((i \cdot l) \leq t) \wedge (t < ((i \cdot l) + l))) \end{array}}}
\vnbprop{conv-join-totally-bounded}{\vnbstmt{\begin{array}{@{}l@{}} \forall m, c, a. \\ \quad \operatorname{is\text{-}normed\text{-}vector\text{-}space}(m) \wedge \\ \quad c \subseteq \operatorname{vec}(m) \wedge \\ \quad a \in \operatorname{vec}(m) \wedge \\ \quad \operatorname{is\text{-}totally\text{-}bounded\text{-}in}(\operatorname{nvs\text{-}metric\text{-}space}(m), c) \Rightarrow \\ \quad \quad \operatorname{is\text{-}totally\text{-}bounded\text{-}in}(\operatorname{nvs\text{-}metric\text{-}space}(m), \operatorname{conv\text{-}join}(m, c, a)) \end{array}}}
\vnbprop{totally-bounded-in-subset}{\vnbstmt{\begin{array}{@{}l@{}} \forall m, a, b. \\ \quad (\operatorname{is\text{-}normed\text{-}vector\text{-}space}(m) \wedge (\operatorname{is\text{-}totally\text{-}bounded\text{-}in}(\operatorname{nvs\text{-}metric\text{-}space}(m), b) \wedge (a \subseteq b))) \Rightarrow  \\ \quad \quad \operatorname{is\text{-}totally\text{-}bounded\text{-}in}(\operatorname{nvs\text{-}metric\text{-}space}(m), a) \end{array}}}
\vnbprop{totally-bounded-in-empty}{\vnbstmt{\begin{array}{@{}l@{}} \forall m. \\ \quad \operatorname{is\text{-}normed\text{-}vector\text{-}space}(m) \Rightarrow  \\ \quad \quad \operatorname{is\text{-}totally\text{-}bounded\text{-}in}(\operatorname{nvs\text{-}metric\text{-}space}(m), \emptyset) \end{array}}}
\vnbprop{totally-bounded-in-add-point}{\vnbstmt{\begin{array}{@{}l@{}} \forall m, s, a. \\ \quad (\operatorname{is\text{-}normed\text{-}vector\text{-}space}(m) \wedge (\operatorname{is\text{-}totally\text{-}bounded\text{-}in}(\operatorname{nvs\text{-}metric\text{-}space}(m), s) \wedge (a \in \operatorname{vec}(m)))) \Rightarrow  \\ \quad \quad \operatorname{is\text{-}totally\text{-}bounded\text{-}in}(\operatorname{nvs\text{-}metric\text{-}space}(m), (s \cup \{a\})) \end{array}}}
\vnbprop{convex-combination-at-one}{\vnbstmt{\begin{array}{@{}l@{}} \forall m, y, a. \\ \quad (\operatorname{is\text{-}normed\text{-}vector\text{-}space}(m) \wedge ((y \in \operatorname{vec}(m)) \wedge (a \in \operatorname{vec}(m)))) \Rightarrow  \\ \quad \quad \operatorname{vadd}(m)(\operatorname{act}(m)(0, y), \operatorname{act}(m)(1, a)) = a \end{array}}}
\vnbprop{convex-combination-same-point}{\vnbstmt{\begin{array}{@{}l@{}} \forall m, a, s. \\ \quad (\operatorname{is\text{-}normed\text{-}vector\text{-}space}(m) \wedge ((a \in \operatorname{vec}(m)) \wedge (s \in \mathbb{R}))) \Rightarrow  \\ \quad \quad \operatorname{vadd}(m)(\operatorname{act}(m)((1 - s), a), \operatorname{act}(m)(s, a)) = a \end{array}}}
\vnbprop{conv-join-with-point-convex}{\vnbstmt{\begin{array}{@{}l@{}} \forall m, c, a. \\ \quad (\operatorname{is\text{-}normed\text{-}vector\text{-}space}(m) \wedge (\operatorname{is\text{-}convex\text{-}subset}(m, c) \wedge (a \in \operatorname{vec}(m)))) \Rightarrow  \\ \quad \quad \operatorname{is\text{-}convex\text{-}subset}(m, (\operatorname{conv\text{-}join}(m, c, a) \cup \{a\})) \end{array}}}
\vnbprop{conv-hull-finite-totally-bounded}{\vnbstmt{\begin{array}{@{}l@{}} \forall m.; \forall a \in \mathbf{Set}. \\ \quad (\operatorname{is\text{-}normed\text{-}vector\text{-}space}(m) \wedge ((|a| \in \mathbb{N}) \wedge (a \subseteq \operatorname{vec}(m)))) \Rightarrow  \\ \quad \quad \operatorname{is\text{-}totally\text{-}bounded\text{-}in}(\operatorname{nvs\text{-}metric\text{-}space}(m), \operatorname{conv\text{-}hull}(m, a)) \end{array}}}
\vnbprop{convex-combination-shift-second-point}{\vnbstmt{\begin{array}{@{}l@{}} \forall m, a, x, z, t. \\ \quad (\operatorname{is\text{-}normed\text{-}vector\text{-}space}(m) \wedge ((a \in \operatorname{vec}(m)) \wedge ((x \in \operatorname{vec}(m)) \wedge ((z \in \operatorname{vec}(m)) \wedge (t \in \mathbb{R}))))) \Rightarrow  \\ \quad \quad \operatorname{vadd}(m)(\operatorname{act}(m)((1 - t), a), \operatorname{act}(m)(t, z)) \\ \quad \quad =\; \operatorname{vadd}(m)(\operatorname{vadd}(m)(\operatorname{act}(m)((1 - t), a), \operatorname{act}(m)(t, x)), \operatorname{act}(m)(t, \operatorname{vadd}(m)(z, \operatorname{vneg}(m)(x)))) \end{array}}}
\vnbprop{convex-combination-dist-second-point}{\vnbstmt{\begin{array}{@{}l@{}} \forall m, a, x, z, t. \\ \quad (\operatorname{is\text{-}normed\text{-}vector\text{-}space}(m) \wedge ((a \in \operatorname{vec}(m)) \wedge ((x \in \operatorname{vec}(m)) \wedge ((z \in \operatorname{vec}(m)) \wedge (t \in \mathbb{R}))))) \Rightarrow  \\ \quad \quad \operatorname{dist}(\operatorname{nvs\text{-}metric\text{-}space}(m))(\operatorname{vadd}(m)(\operatorname{act}(m)((1 - t), a), \operatorname{act}(m)(t, x)), \\ \quad \quad \quad \operatorname{vadd}(m)(\operatorname{act}(m)((1 - t), a), \operatorname{act}(m)(t, z))) \\ \quad \quad =\; \lvert t \rvert \cdot \operatorname{dist}(\operatorname{nvs\text{-}metric\text{-}space}(m))(z, x) \end{array}}}
\vnbprop{conv-hull-totally-bounded}{\vnbstmt{\begin{array}{@{}l@{}} \forall m, e. \\ \quad (\operatorname{is\text{-}normed\text{-}vector\text{-}space}(m) \wedge ((e \subseteq \operatorname{vec}(m)) \wedge \operatorname{is\text{-}totally\text{-}bounded\text{-}in}(\operatorname{nvs\text{-}metric\text{-}space}(m), e))) \Rightarrow  \\ \quad \quad \operatorname{is\text{-}totally\text{-}bounded\text{-}in}(\operatorname{nvs\text{-}metric\text{-}space}(m), \operatorname{conv\text{-}hull}(m, e)) \end{array}}}

\vnbbookfile{structure-library/cg4-vocab}{structure-library/cg4-vocab}
\noindent{\small cg4-vocab.scm -- compact-groups notes, sections 1 and 5 (CG-4 session).  DEFINITIONS ONLY. Agent ISOM, 2026-10-26 (new definitions, the user's to confirm).}\par
\par\medskip\noindent\textbf{Definition} (\texttt{is-extreme-point}).\;For all $\operatorname{epm\_}$, $\operatorname{epk\_}$ and $\operatorname{epx\_}$, 
\begin{equation*}
\begin{array}{@{}l@{}} \operatorname{is\text{-}extreme\text{-}point}(\operatorname{epm\_}, \operatorname{epk\_}, \operatorname{epx\_}) \\ \Leftrightarrow\; \operatorname{epx\_} \in \operatorname{epk\_} \wedge \\ \forall \operatorname{epu\_}, \operatorname{epv\_} \in \operatorname{epk\_}, \operatorname{ept\_} \in \mathbb{R}. \\ \quad (0 < \operatorname{ept\_}) \Rightarrow  \\ \quad (\operatorname{ept\_} < 1) \Rightarrow  \\ \quad (\operatorname{epx\_} = \operatorname{vadd}(\operatorname{epm\_})(\operatorname{act}(\operatorname{epm\_})(-(1, \operatorname{ept\_}), \operatorname{epu\_}), \operatorname{act}(\operatorname{epm\_})(\operatorname{ept\_}, \operatorname{epv\_}))) \Rightarrow  \\ \quad \quad (\operatorname{epu\_} = \operatorname{epx\_}) \wedge (\operatorname{epv\_} = \operatorname{epx\_}) \end{array}
\end{equation*}

\par\medskip\noindent\textbf{Definition} (\texttt{pts-left-cosets}).\;$\operatorname{pts\text{-}left\text{-}cosets}(\operatorname{lcs\_}, \operatorname{lcg\_}, \operatorname{lch\_}) :=$
\begin{equation*}
\{\operatorname{lcu\_} \in \operatorname{power}(\operatorname{pts}(\operatorname{lcs\_})) : \exists \operatorname{lcx\_} \in \operatorname{carr}(\operatorname{lcg\_}).\; (\operatorname{lcu\_} = \{\operatorname{lcz\_} \in \operatorname{pts}(\operatorname{lcs\_}) : (\operatorname{opr}(\operatorname{lcg\_})(\operatorname{inv}(\operatorname{lcg\_})(\operatorname{lcx\_}), \operatorname{lcz\_}) \in \operatorname{lch\_})\})\}
\end{equation*}

\par\medskip\noindent\textbf{Definition} (\texttt{is-one-parameter-group}).\;For all $\operatorname{ops\_}$, $\operatorname{opg\_}$ and $\operatorname{opp\_}$, 
\begin{equation*}
\begin{array}{@{}l@{}} \operatorname{is\text{-}one\text{-}parameter\text{-}group}(\operatorname{ops\_}, \operatorname{opg\_}, \operatorname{opp\_}) \\ \Leftrightarrow\; \operatorname{opp\_} \in (\mathbb{R} \to \operatorname{carr}(\operatorname{opg\_})) \wedge \\ \forall \operatorname{oqr\_}, \operatorname{oqu\_} \in \mathbb{R}. \\ \quad \operatorname{opp\_}(+(\operatorname{oqr\_}, \operatorname{oqu\_})) = \operatorname{opr}(\operatorname{opg\_})(\operatorname{opp\_}(\operatorname{oqr\_}), \operatorname{opp\_}(\operatorname{oqu\_})) \wedge \\ \forall \operatorname{oqt\_} \in \mathbb{R}, \operatorname{oqw\_} \in \operatorname{opens}(\operatorname{ops\_}). \\ \quad (\operatorname{opp\_}(\operatorname{oqt\_}) \in \operatorname{oqw\_}) \Rightarrow  \\ \quad \quad \exists \operatorname{oqd\_}. \\ \quad \quad \quad (\operatorname{pos\text{-}rr}(\operatorname{oqd\_}) \wedge \forall \operatorname{oqv\_} \in \mathbb{R}.\; ((\lvert -(\operatorname{oqv\_}, \operatorname{oqt\_}) \rvert < \operatorname{oqd\_}) \Rightarrow (\operatorname{opp\_}(\operatorname{oqv\_}) \in \operatorname{oqw\_}))) \end{array}
\end{equation*}

\par\medskip\noindent\textbf{Definition} (\texttt{is-group-segment}).\;For all $\operatorname{gss\_}$, $\operatorname{gsg\_}$, $\operatorname{gsa\_}$ and $\operatorname{gsp\_}$, 
\begin{equation*}
\begin{array}{@{}l@{}} \operatorname{is\text{-}group\text{-}segment}(\operatorname{gss\_}, \operatorname{gsg\_}, \operatorname{gsa\_}, \operatorname{gsp\_}) \\ \Leftrightarrow\; \operatorname{pos\text{-}rr}(\operatorname{gsa\_}) \wedge \\ \operatorname{gsp\_} \in (\operatorname{ccint}(0, \operatorname{gsa\_}) \to \operatorname{carr}(\operatorname{gsg\_})) \wedge \\ \forall \operatorname{gsr\_}, \operatorname{gsu\_} \in \operatorname{ccint}(0, \operatorname{gsa\_}). \\ \quad (+(\operatorname{gsr\_}, \operatorname{gsu\_}) \leq \operatorname{gsa\_}) \Rightarrow  \\ \quad \quad \operatorname{gsp\_}(+(\operatorname{gsr\_}, \operatorname{gsu\_})) = \operatorname{opr}(\operatorname{gsg\_})(\operatorname{gsp\_}(\operatorname{gsr\_}), \operatorname{gsp\_}(\operatorname{gsu\_})) \wedge \\ \forall \operatorname{gst\_} \in \operatorname{ccint}(0, \operatorname{gsa\_}), \operatorname{gsw\_} \in \operatorname{opens}(\operatorname{gss\_}). \\ \quad (\operatorname{gsp\_}(\operatorname{gst\_}) \in \operatorname{gsw\_}) \Rightarrow  \\ \quad \quad \exists \operatorname{gsd\_}. \\ \quad \quad \quad \operatorname{pos\text{-}rr}(\operatorname{gsd\_}) \wedge \\ \quad \quad \quad \forall \operatorname{gsv\_} \in \operatorname{ccint}(0, \operatorname{gsa\_}). \\ \quad \quad \quad \quad (\lvert -(\operatorname{gsv\_}, \operatorname{gst\_}) \rvert < \operatorname{gsd\_}) \Rightarrow (\operatorname{gsp\_}(\operatorname{gsv\_}) \in \operatorname{gsw\_}) \end{array}
\end{equation*}

\par\medskip\noindent\textbf{Definition} (\texttt{opg-rep}).\;For all $\operatorname{rpg\_}$, $\operatorname{rpa\_}$, $\operatorname{rpp\_}$, $\operatorname{rpt\_}$ and $\operatorname{rpy\_}$, 
\begin{equation*}
\begin{array}{@{}l@{}} \operatorname{opg\text{-}rep}(\operatorname{rpg\_}, \operatorname{rpa\_}, \operatorname{rpp\_}, \operatorname{rpt\_}, \operatorname{rpy\_}) \\ \Leftrightarrow\; 0 \leq \operatorname{rpt\_} \wedge \\ \exists \operatorname{rpn\_} \in \mathbb{N}, \operatorname{rpr\_} \in \mathbb{R}. \\ \quad 0 \leq \operatorname{rpr\_} \wedge \\ \quad \operatorname{rpr\_} < \operatorname{rpa\_} \wedge \\ \quad \operatorname{rpt\_} = +(\cdot(\operatorname{rpn\_}, \operatorname{rpa\_}), \operatorname{rpr\_}) \wedge \\ \quad \operatorname{rpy\_} = \operatorname{opr}(\operatorname{rpg\_})(\operatorname{gpow}(\operatorname{rpg\_}, \operatorname{rpp\_}(\operatorname{rpa\_}), \operatorname{rpn\_}), \operatorname{rpp\_}(\operatorname{rpr\_})) \vee \\ \operatorname{rpt\_} < 0 \wedge \\ \exists \operatorname{rpn\_} \in \mathbb{N}, \operatorname{rpr\_} \in \mathbb{R}. \\ \quad 0 \leq \operatorname{rpr\_} \wedge \\ \quad \operatorname{rpr\_} < \operatorname{rpa\_} \wedge \\ \quad -(0, \operatorname{rpt\_}) = +(\cdot(\operatorname{rpn\_}, \operatorname{rpa\_}), \operatorname{rpr\_}) \wedge \\ \quad \operatorname{rpy\_} = \operatorname{inv}(\operatorname{rpg\_})(\operatorname{opr}(\operatorname{rpg\_})(\operatorname{gpow}(\operatorname{rpg\_}, \operatorname{rpp\_}(\operatorname{rpa\_}), \operatorname{rpn\_}), \operatorname{rpp\_}(\operatorname{rpr\_}))) \end{array}
\end{equation*}

\par\medskip\noindent\textbf{Definition} (\texttt{opg-val}).\;$\operatorname{opg\text{-}val}(\operatorname{rpg\_}, \operatorname{rpa\_}, \operatorname{rpp\_}, \operatorname{rvt\_}) :=$
\begin{equation*}
(\iota\, \operatorname{rvy\_}.\;((\operatorname{rvy\_} \in \operatorname{carr}(\operatorname{rpg\_})) \wedge \operatorname{opg\text{-}rep}(\operatorname{rpg\_}, \operatorname{rpa\_}, \operatorname{rpp\_}, \operatorname{rvt\_}, \operatorname{rvy\_})))
\end{equation*}

\par\medskip\noindent\textbf{Definition} (\texttt{opg-ext}).\;$\operatorname{opg\text{-}ext}(\operatorname{rpg\_}, \operatorname{rpa\_}, \operatorname{rpp\_}) :=$ $(\lambda \operatorname{rvt\_} \in \mathbb{R}.\; \operatorname{opg\text{-}val}(\operatorname{rpg\_}, \operatorname{rpa\_}, \operatorname{rpp\_}, \operatorname{rvt\_}))$
\par\medskip\noindent\textbf{Definition} (\texttt{is-local-hom}).\;For all $\operatorname{lhs1\_}$, $\operatorname{lhg1\_}$, $\operatorname{lhs2\_}$, $\operatorname{lhg2\_}$, $\operatorname{lhv\_}$ and $\operatorname{lhf\_}$, 
\begin{equation*}
\begin{array}{@{}l@{}} \operatorname{is\text{-}local\text{-}hom}(\operatorname{lhs1\_}, \operatorname{lhg1\_}, \operatorname{lhs2\_}, \operatorname{lhg2\_}, \operatorname{lhv\_}, \operatorname{lhf\_}) \\ \Leftrightarrow\; \operatorname{lhv\_} \in \operatorname{opens}(\operatorname{lhs1\_}) \wedge \\ \operatorname{iden}(\operatorname{lhg1\_}) \in \operatorname{lhv\_} \wedge \\ \operatorname{lhf\_} \in (\operatorname{lhv\_} \to \operatorname{carr}(\operatorname{lhg2\_})) \wedge \\ \forall \operatorname{lhx\_} \in \operatorname{lhv\_}, \operatorname{lho\_} \in \operatorname{opens}(\operatorname{lhs2\_}). \\ \quad (\operatorname{lhf\_}(\operatorname{lhx\_}) \in \operatorname{lho\_}) \Rightarrow  \\ \quad \quad \exists \operatorname{lhu\_} \in \operatorname{opens}(\operatorname{lhs1\_}). \\ \quad \quad \quad (\operatorname{lhx\_} \in \operatorname{lhu\_}) \wedge \forall \operatorname{lhy\_} \in \operatorname{lhu\_}.\; ((\operatorname{lhy\_} \in \operatorname{lhv\_}) \Rightarrow (\operatorname{lhf\_}(\operatorname{lhy\_}) \in \operatorname{lho\_})) \wedge \\ \exists \operatorname{lhw\_} \in \operatorname{opens}(\operatorname{lhs1\_}). \\ \quad \operatorname{iden}(\operatorname{lhg1\_}) \in \operatorname{lhw\_} \wedge \\ \quad \forall \operatorname{lhx\_} \in \operatorname{lhw\_}.\; (\operatorname{inv}(\operatorname{lhg1\_})(\operatorname{lhx\_}) \in \operatorname{lhw\_}) \wedge \\ \quad \forall \operatorname{lhx\_}, \operatorname{lhy\_} \in \operatorname{lhw\_}. \\ \quad \quad \operatorname{opr}(\operatorname{lhg1\_})(\operatorname{lhx\_}, \operatorname{lhy\_}) \in \operatorname{lhv\_} \wedge \\ \quad \quad \operatorname{lhf\_}(\operatorname{opr}(\operatorname{lhg1\_})(\operatorname{lhx\_}, \operatorname{lhy\_})) = \operatorname{opr}(\operatorname{lhg2\_})(\operatorname{lhf\_}(\operatorname{lhx\_}), \operatorname{lhf\_}(\operatorname{lhy\_})) \end{array}
\end{equation*}


\vnbbookfile{structure-library/km-vocab}{structure-library/km-vocab}
\noindent{\small km-vocab.scm -- wave of 2026-11-09/10 (agent BA-KM), wired 2026-11-10 from scratchpad/bakm-1109/km-vocab.scm. Library: analysis.}\par
\par\medskip\noindent\textbf{Definition} (\texttt{is-locally-convex-tvs}).\;For all $m$ and $t$, 
\begin{equation*}
\begin{array}{@{}l@{}} \operatorname{is\text{-}locally\text{-}convex\text{-}tvs}(m, t) \\ \Leftrightarrow\; \operatorname{is\text{-}tvs}(\operatorname{rr\text{-}normed\text{-}field}, m, t) \wedge \\ \forall \operatorname{lcx\_} \in \operatorname{vec}(m), \operatorname{lcu\_} \in \operatorname{opens}(t). \\ \quad (\operatorname{lcx\_} \in \operatorname{lcu\_}) \Rightarrow  \\ \quad \quad \exists \operatorname{lcv\_} \in \operatorname{opens}(t). \\ \quad \quad \quad (\operatorname{lcx\_} \in \operatorname{lcv\_}) \wedge ((\operatorname{lcv\_} \subseteq \operatorname{lcu\_}) \wedge \operatorname{is\text{-}convex\text{-}subset}(m, \operatorname{lcv\_})) \end{array}
\end{equation*}

\par\medskip\noindent\textbf{Definition} (\texttt{extreme-points}).\;$\operatorname{extreme\text{-}points}(m, \mathit{ks}) :=$ $\{\operatorname{kez\_} \in \operatorname{vec}(m) : \operatorname{is\text{-}extreme\text{-}point}(m, \mathit{ks}, \operatorname{kez\_})\}$
\par\medskip\noindent\textbf{Definition} (\texttt{closed-conv-hull}).\;$\operatorname{closed\text{-}conv\text{-}hull}(m, t, e) :=$
\begin{equation*}
\{\operatorname{khx\_} \in \operatorname{vec}(m) : \forall \operatorname{khc\_} \in \mathbf{Set}.\; (\operatorname{is\text{-}convex\text{-}subset}(m, \operatorname{khc\_}) \Rightarrow (\operatorname{is\text{-}top\text{-}closed}(t, \operatorname{khc\_}) \Rightarrow ((e \subseteq \operatorname{khc\_}) \Rightarrow (\operatorname{khx\_} \in \operatorname{khc\_}))))\}
\end{equation*}


\vnbbookfile{theorem-library/cg4-extreme}{theorem-library/cg4-extreme}
\noindent{\small cg4-extreme.scm -- the extreme points of the convex hull of E lie in E (compact-groups notes, Lemma 1.13, its algebraic core: eqs. (10)-(11) of the proof.  The closures of the notes' statement -- "a closed set whose CLOSED convex hull is A" -- are skipped: CLOSURE). Agent ISOM, 2026-10-26.  Helper prefix: ce4-.  ...}\par
\vnbprop{extreme-point-of-level}{\vnbstmt{\begin{array}{@{}l@{}} \forall k \in \mathbb{N}, m, e, x. \\ \quad \operatorname{is\text{-}normed\text{-}vector\text{-}space}(m) \wedge \\ \quad e \subseteq \operatorname{vec}(m) \wedge \\ \quad \operatorname{is\text{-}extreme\text{-}point}(m, \operatorname{conv\text{-}hull}(m, e), x) \wedge \\ \quad x \in \operatorname{conv\text{-}level}(m, e, k) \Rightarrow \\ \quad \quad x \in e \end{array}}}
\vnbprop{extreme-points-of-hull-in-set}{\vnbstmt{\begin{array}{@{}l@{}} \forall m, e, x. \\ \quad (\operatorname{is\text{-}normed\text{-}vector\text{-}space}(m) \wedge ((e \subseteq \operatorname{vec}(m)) \wedge \operatorname{is\text{-}extreme\text{-}point}(m, \operatorname{conv\text{-}hull}(m, e), x))) \Rightarrow  \\ \quad \quad x \in e \end{array}}}

\vnbbookfile{theorem-library/krein-milman}{theorem-library/krein-milman}
\noindent{\small krein-milman.scm -- wave of 2026-11-09/10 (agent BA-KM), wired 2026-11-10 from scratchpad/bakm-1109/krein-milman.scm. Library: analysis.}\par
\vnbprop{extreme-points-unfold}{\vnbstmt{\forall m, s.\; (\operatorname{extreme\text{-}points}(m, s) \simeq \{z \in \operatorname{vec}(m) : \operatorname{is\text{-}extreme\text{-}point}(m, s, z)\})}}
\vnbprop{closed-conv-hull-unfold}{\vnbstmt{\begin{array}{@{}l@{}} \forall m, t, e. \\ \quad (\operatorname{closed\text{-}conv\text{-}hull}(m, t, e) \simeq \{x \in \operatorname{vec}(m) : \forall c \in \mathbf{Set}.\; ((\operatorname{is\text{-}convex\text{-}subset}(m, c) \wedge (\operatorname{is\text{-}top\text{-}closed}(t, c) \wedge (e \subseteq c))) \Rightarrow (x \in c))\}) \end{array}}}
\vnbprop{extreme-points-mem}{\vnbstmt{\begin{array}{@{}l@{}} \forall m, s, z. \\ \quad (z \in \operatorname{extreme\text{-}points}(m, s)) \Leftrightarrow ((z \in \operatorname{vec}(m)) \wedge \operatorname{is\text{-}extreme\text{-}point}(m, s, z)) \end{array}}}
\vnbprop{closed-conv-hull-mem}{\vnbstmt{\begin{array}{@{}l@{}} \forall m, t, e, z. \\ \quad z \in \operatorname{closed\text{-}conv\text{-}hull}(m, t, e) \\ \quad \Leftrightarrow\; z \in \operatorname{vec}(m) \wedge \\ \quad \forall c \in \mathbf{Set}. \\ \quad \quad (\operatorname{is\text{-}convex\text{-}subset}(m, c) \wedge (\operatorname{is\text{-}top\text{-}closed}(t, c) \wedge (e \subseteq c))) \Rightarrow (z \in c) \end{array}}}
\vnbprop{extreme-point-of-extremal-subset}{\vnbstmt{\begin{array}{@{}l@{}} \forall m, b, a, e. \\ \quad (\operatorname{is\text{-}extremal\text{-}subset}(m, b, a) \wedge \operatorname{is\text{-}extreme\text{-}point}(m, b, e)) \Rightarrow  \\ \quad \quad \operatorname{is\text{-}extreme\text{-}point}(m, a, e) \end{array}}}
\vnbprop{extreme-point-in-subset}{\vnbstmt{\begin{array}{@{}l@{}} \forall m, b, a, e. \\ \quad ((b \subseteq a) \wedge (\operatorname{is\text{-}extreme\text{-}point}(m, a, e) \wedge (e \in b))) \Rightarrow  \\ \quad \quad \operatorname{is\text{-}extreme\text{-}point}(m, b, e) \end{array}}}
\vnbprop{extremal-subset-singleton-extreme}{\vnbstmt{\begin{array}{@{}l@{}} \forall m, f, a, e. \\ \quad (\operatorname{is\text{-}extremal\text{-}subset}(m, f, a) \wedge (\forall p, q \in f.\; (p = q) \wedge (e \in f))) \Rightarrow  \\ \quad \quad \operatorname{is\text{-}extreme\text{-}point}(m, a, e) \end{array}}}
\vnbprop{top-closed-intersection}{\vnbstmt{\begin{array}{@{}l@{}} \forall t, a, b. \\ \quad (\operatorname{is\text{-}top\text{-}space}(t) \wedge (\operatorname{is\text{-}top\text{-}closed}(t, a) \wedge \operatorname{is\text{-}top\text{-}closed}(t, b))) \Rightarrow  \\ \quad \quad \operatorname{is\text{-}top\text{-}closed}(t, (a \cap b)) \end{array}}}
\vnbprop{top-union-open}{\vnbstmt{\begin{array}{@{}l@{}} \forall t, u, v. \\ \quad (\operatorname{is\text{-}top\text{-}space}(t) \wedge ((u \in \operatorname{opens}(t)) \wedge (v \in \operatorname{opens}(t)))) \Rightarrow  \\ \quad \quad (u \cup v) \in \operatorname{opens}(t) \end{array}}}
\vnbprop{compact-directed-cover}{\vnbstmt{\begin{array}{@{}l@{}} \forall t, s, d. \\ \quad \operatorname{is\text{-}top\text{-}space}(t) \wedge \\ \quad s \subseteq \operatorname{pts}(t) \wedge \\ \quad \operatorname{is\text{-}compact\text{-}t}(\operatorname{subspace\text{-}top}(t, s)) \wedge \\ \quad \forall u \in d.\; (u \in \operatorname{opens}(t)) \wedge \\ \quad \exists u.\; (u \in d) \wedge \\ \quad \forall a, b \in d.\; \exists c \in d.\; ((a \subseteq c) \wedge (b \subseteq c)) \wedge \\ \quad \forall x \in s.\; \exists u \in d.\; (x \in u) \Rightarrow \\ \quad \quad \exists u \in d.\; (s \subseteq u) \end{array}}}
\vnbprop{compact-hausdorff-closed}{\vnbstmt{\begin{array}{@{}l@{}} \forall t, s. \\ \quad (\operatorname{is\text{-}hausdorff}(t) \wedge ((s \subseteq \operatorname{pts}(t)) \wedge \operatorname{is\text{-}compact\text{-}t}(\operatorname{subspace\text{-}top}(t, s)))) \Rightarrow  \\ \quad \quad \operatorname{is\text{-}top\text{-}closed}(t, s) \end{array}}}
\vnbprop{compact-chain-meet}{\vnbstmt{\begin{array}{@{}l@{}} \forall t, s, h. \\ \quad \operatorname{is\text{-}top\text{-}space}(t) \wedge \\ \quad s \subseteq \operatorname{pts}(t) \wedge \\ \quad \operatorname{is\text{-}compact\text{-}t}(\operatorname{subspace\text{-}top}(t, s)) \wedge \\ \quad \exists a.\; (a \in h) \wedge \\ \quad \forall a \in h.\; \operatorname{is\text{-}top\text{-}closed}(t, a) \wedge \\ \quad \forall a \in h.\; (a \subseteq s) \wedge \\ \quad \forall a \in h.\; \exists z.\; (z \in a) \wedge \\ \quad \forall a, b \in h.\; ((a \subseteq b) \vee (b \subseteq a)) \Rightarrow \\ \quad \quad \exists p \in s.\; \forall a \in h.\; (p \in a) \end{array}}}
\vnbprop{closed-superlevel-set}{\vnbstmt{\begin{array}{@{}l@{}} \forall t, g, s, c. \\ \quad \operatorname{is\text{-}hom\text{-}top\text{-}space}(t, \\ \quad \quad \operatorname{seminorm\text{-}top}(\operatorname{normed\text{-}vector\text{-}space\text{-}as\text{-}module}(\operatorname{rr\text{-}nvs}), \\ \quad \quad \quad \{0\}, (\lambda k \in \{0\}.\; \operatorname{vnrm}(\operatorname{rr\text{-}nvs}))), \\ \quad \quad g) \wedge \\ \quad \operatorname{is\text{-}top\text{-}closed}(t, s) \wedge \\ \quad c \in \mathbb{R} \Rightarrow \\ \quad \quad \operatorname{is\text{-}top\text{-}closed}(t, \{h \in s : (c \leq g(h))\}) \end{array}}}
\vnbprop{compact-max-attained}{\vnbstmt{\begin{array}{@{}l@{}} \forall t, s, g. \\ \quad \operatorname{is\text{-}hausdorff}(t) \wedge \\ \quad s \subseteq \operatorname{pts}(t) \wedge \\ \quad \operatorname{is\text{-}compact\text{-}t}(\operatorname{subspace\text{-}top}(t, s)) \wedge \\ \quad \exists z.\; (z \in s) \wedge \\ \quad \operatorname{is\text{-}hom\text{-}top\text{-}space}(t, \\ \quad \quad \operatorname{seminorm\text{-}top}(\operatorname{normed\text{-}vector\text{-}space\text{-}as\text{-}module}(\operatorname{rr\text{-}nvs}), \\ \quad \quad \quad \{0\}, (\lambda k \in \{0\}.\; \operatorname{vnrm}(\operatorname{rr\text{-}nvs}))), \\ \quad \quad g) \Rightarrow \\ \quad \quad \exists p \in s.\; \forall w \in s.\; (g(w) \leq g(p)) \end{array}}}
\vnbprop{max-face-extremal}{\vnbstmt{\begin{array}{@{}l@{}} \forall m, g, s, p. \\ \quad \operatorname{is\text{-}module}(m) \wedge \\ \quad \operatorname{scal}(m) = \operatorname{normed\text{-}field\text{-}as\text{-}commutative\text{-}ring}(\operatorname{rr\text{-}normed\text{-}field}) \wedge \\ \quad \operatorname{is\text{-}linear\text{-}functional}(m, g) \wedge \\ \quad \operatorname{is\text{-}convex\text{-}subset}(m, s) \wedge \\ \quad p \in s \wedge \\ \quad \forall w \in s.\; (g(w) \leq g(p)) \Rightarrow \\ \quad \quad \operatorname{is\text{-}extremal\text{-}subset}(m, \{h \in s : (g(p) \leq g(h))\}, s) \end{array}}}
\vnbprop{rmod-neg-combination}{\vnbstmt{\begin{array}{@{}l@{}} \forall m, u, t, a, b. \\ \quad \operatorname{is\text{-}module}(m) \wedge \\ \quad \operatorname{scal}(m) = \operatorname{normed\text{-}field\text{-}as\text{-}commutative\text{-}ring}(\operatorname{rr\text{-}normed\text{-}field}) \wedge \\ \quad u \in \mathbb{R} \wedge \\ \quad t \in \mathbb{R} \wedge \\ \quad a \in \operatorname{vec}(m) \wedge \\ \quad b \in \operatorname{vec}(m) \Rightarrow \\ \quad \quad \operatorname{act}(m)(-1, \operatorname{vadd}(m)(\operatorname{act}(m)(u, a), \operatorname{act}(m)(t, b))) \\ \quad \quad =\; \operatorname{vadd}(m)(\operatorname{act}(m)(u, \operatorname{act}(m)(-1, a)), \operatorname{act}(m)(t, \operatorname{act}(m)(-1, b))) \end{array}}}
\vnbprop{rmod-neg-neg}{\vnbstmt{\begin{array}{@{}l@{}} \forall m, x. \\ \quad (\operatorname{is\text{-}module}(m) \wedge ((\operatorname{scal}(m) = \operatorname{normed\text{-}field\text{-}as\text{-}commutative\text{-}ring}(\operatorname{rr\text{-}normed\text{-}field})) \wedge (x \in \operatorname{vec}(m)))) \Rightarrow  \\ \quad \quad \operatorname{act}(m)(-1, \operatorname{act}(m)(-1, x)) = x \end{array}}}
\vnbprop{rmod-add-sub-cancel}{\vnbstmt{\begin{array}{@{}l@{}} \forall m, x, z. \\ \quad \operatorname{is\text{-}module}(m) \wedge \\ \quad \operatorname{scal}(m) = \operatorname{normed\text{-}field\text{-}as\text{-}commutative\text{-}ring}(\operatorname{rr\text{-}normed\text{-}field}) \wedge \\ \quad x \in \operatorname{vec}(m) \wedge \\ \quad z \in \operatorname{vec}(m) \Rightarrow \\ \quad \quad \operatorname{vadd}(m)(\operatorname{vadd}(m)(x, z), \operatorname{act}(m)(-1, z)) = x \end{array}}}
\vnbprop{convex-open-separation}{\vnbstmt{\begin{array}{@{}l@{}} \forall m, t, v, a, c. \\ \quad \operatorname{is\text{-}tvs}(\operatorname{rr\text{-}normed\text{-}field}, m, t) \wedge \\ \quad v \in \operatorname{opens}(t) \wedge \\ \quad \operatorname{is\text{-}convex\text{-}subset}(m, v) \wedge \\ \quad a \in v \wedge \\ \quad \operatorname{is\text{-}convex\text{-}subset}(m, c) \wedge \\ \quad \exists c_{0}.\; (c_{0} \in c) \wedge \\ \quad \forall z \in c.\; \neg (z \in v) \Rightarrow \\ \quad \quad \exists f. \\ \quad \quad \quad \operatorname{is\text{-}linear\text{-}functional}(m, f) \wedge \\ \quad \quad \quad \operatorname{is\text{-}hom\text{-}top\text{-}space}(t, \\ \quad \quad \quad \quad \operatorname{seminorm\text{-}top}(\operatorname{normed\text{-}vector\text{-}space\text{-}as\text{-}module}(\operatorname{rr\text{-}nvs}), \\ \quad \quad \quad \quad \quad \{0\}, (\lambda k \in \{0\}.\; \operatorname{vnrm}(\operatorname{rr\text{-}nvs}))), \\ \quad \quad \quad \quad f) \wedge \\ \quad \quad \quad \forall \mathit{kgz} \in c.\; (f(\mathit{kgz}) < f(a)) \end{array}}}
\vnbprop{locally-convex-separation}{\vnbstmt{\begin{array}{@{}l@{}} \forall m, t, u, a, c. \\ \quad \operatorname{is\text{-}locally\text{-}convex\text{-}tvs}(m, t) \wedge \\ \quad u \in \operatorname{opens}(t) \wedge \\ \quad a \in u \wedge \\ \quad \operatorname{is\text{-}convex\text{-}subset}(m, c) \wedge \\ \quad \exists c_{0}.\; (c_{0} \in c) \wedge \\ \quad \forall z \in c.\; \neg (z \in u) \Rightarrow \\ \quad \quad \exists f. \\ \quad \quad \quad \operatorname{is\text{-}linear\text{-}functional}(m, f) \wedge \\ \quad \quad \quad \operatorname{is\text{-}hom\text{-}top\text{-}space}(t, \\ \quad \quad \quad \quad \operatorname{seminorm\text{-}top}(\operatorname{normed\text{-}vector\text{-}space\text{-}as\text{-}module}(\operatorname{rr\text{-}nvs}), \\ \quad \quad \quad \quad \quad \{0\}, (\lambda k \in \{0\}.\; \operatorname{vnrm}(\operatorname{rr\text{-}nvs}))), \\ \quad \quad \quad \quad f) \wedge \\ \quad \quad \quad \forall \mathit{kgz} \in c.\; (f(\mathit{kgz}) < f(a)) \end{array}}}
\vnbprop{extremal-subset-refl}{\vnbstmt{\begin{array}{@{}l@{}} \forall m, s. \\ \quad (\operatorname{is\text{-}convex\text{-}subset}(m, s) \wedge \exists z_{0}.\; (z_{0} \in s)) \Rightarrow  \\ \quad \quad \operatorname{is\text{-}extremal\text{-}subset}(m, s, s) \end{array}}}
\vnbprop{extremal-meet}{\vnbstmt{\begin{array}{@{}l@{}} \forall m, s, h. \\ \quad \operatorname{is\text{-}module}(m) \wedge \\ \quad \operatorname{is\text{-}convex\text{-}subset}(m, s) \wedge \\ \quad \forall a \in h.\; \operatorname{is\text{-}extremal\text{-}subset}(m, a, s) \wedge \\ \quad \exists z.\; (z \in \{b \in s : \forall a \in h.\; (b \in a)\}) \Rightarrow \\ \quad \quad \operatorname{is\text{-}extremal\text{-}subset}(m, \{b \in s : \forall a \in h.\; (b \in a)\}, s) \end{array}}}
\vnbprop{closed-meet}{\vnbstmt{\begin{array}{@{}l@{}} \forall t, s, h. \\ \quad (\operatorname{is\text{-}top\text{-}space}(t) \wedge (\operatorname{is\text{-}top\text{-}closed}(t, s) \wedge \forall a \in h.\; \operatorname{is\text{-}top\text{-}closed}(t, a))) \Rightarrow  \\ \quad \quad \operatorname{is\text{-}top\text{-}closed}(t, \{b \in s : \forall a \in h.\; (b \in a)\}) \end{array}}}
\vnbprop{locally-convex-separates-points}{\vnbstmt{\begin{array}{@{}l@{}} \forall m, t, x, y. \\ \quad (\operatorname{is\text{-}locally\text{-}convex\text{-}tvs}(m, t) \wedge ((x \in \operatorname{vec}(m)) \wedge ((y \in \operatorname{vec}(m)) \wedge \neg (x = y)))) \Rightarrow  \\ \quad \quad \exists f. \\ \quad \quad \quad \operatorname{is\text{-}linear\text{-}functional}(m, f) \wedge \\ \quad \quad \quad \operatorname{is\text{-}hom\text{-}top\text{-}space}(t, \\ \quad \quad \quad \quad \operatorname{seminorm\text{-}top}(\operatorname{normed\text{-}vector\text{-}space\text{-}as\text{-}module}(\operatorname{rr\text{-}nvs}), \\ \quad \quad \quad \quad \quad \{0\}, (\lambda k \in \{0\}.\; \operatorname{vnrm}(\operatorname{rr\text{-}nvs}))), \\ \quad \quad \quad \quad f) \wedge \\ \quad \quad \quad f(y) < f(x) \end{array}}}
\vnbprop{closed-face-above}{\vnbstmt{\begin{array}{@{}l@{}} \forall m, t, s, f, g, x. \\ \quad \operatorname{is\text{-}tvs}(\operatorname{rr\text{-}normed\text{-}field}, m, t) \wedge \\ \quad s \subseteq \operatorname{vec}(m) \wedge \\ \quad \operatorname{is\text{-}compact\text{-}t}(\operatorname{subspace\text{-}top}(t, s)) \wedge \\ \quad \operatorname{is\text{-}top\text{-}closed}(t, f) \wedge \\ \quad \operatorname{is\text{-}extremal\text{-}subset}(m, f, s) \wedge \\ \quad \operatorname{is\text{-}linear\text{-}functional}(m, g) \wedge \\ \quad \operatorname{is\text{-}hom\text{-}top\text{-}space}(t, \\ \quad \quad \operatorname{seminorm\text{-}top}(\operatorname{normed\text{-}vector\text{-}space\text{-}as\text{-}module}(\operatorname{rr\text{-}nvs}), \\ \quad \quad \quad \{0\}, (\lambda k \in \{0\}.\; \operatorname{vnrm}(\operatorname{rr\text{-}nvs}))), \\ \quad \quad g) \wedge \\ \quad x \in f \Rightarrow \\ \quad \quad \exists \mathit{kfm}. \\ \quad \quad \quad \operatorname{is\text{-}top\text{-}closed}(t, \mathit{kfm}) \wedge \\ \quad \quad \quad \operatorname{is\text{-}extremal\text{-}subset}(m, \mathit{kfm}, s) \wedge \\ \quad \quad \quad \mathit{kfm} \subseteq f \wedge \\ \quad \quad \quad \exists e.\; (e \in \mathit{kfm}) \wedge \\ \quad \quad \quad \forall z \in \mathit{kfm}.\; (g(x) \leq g(z)) \end{array}}}
\vnbprop{extreme-point-exists}{\vnbstmt{\begin{array}{@{}l@{}} \forall m, t, s. \\ \quad \operatorname{is\text{-}locally\text{-}convex\text{-}tvs}(m, t) \wedge \\ \quad s \subseteq \operatorname{vec}(m) \wedge \\ \quad \operatorname{is\text{-}convex\text{-}subset}(m, s) \wedge \\ \quad \operatorname{is\text{-}compact\text{-}t}(\operatorname{subspace\text{-}top}(t, s)) \wedge \\ \quad \exists z_{0}.\; (z_{0} \in s) \Rightarrow \\ \quad \quad \exists p.\; \operatorname{is\text{-}extreme\text{-}point}(m, s, p) \end{array}}}
\vnbprop{closed-conv-hull-least}{\vnbstmt{\begin{array}{@{}l@{}} \forall m, t, e, c. \\ \quad (\operatorname{is\text{-}module}(m) \wedge ((c \in \mathbf{Set}) \wedge (\operatorname{is\text{-}convex\text{-}subset}(m, c) \wedge (\operatorname{is\text{-}top\text{-}closed}(t, c) \wedge (e \subseteq c))))) \Rightarrow  \\ \quad \quad \operatorname{closed\text{-}conv\text{-}hull}(m, t, e) \subseteq c \end{array}}}
\vnbprop{closed-conv-hull-contains}{\vnbstmt{\begin{array}{@{}l@{}} \forall m, t, e. \\ \quad (\operatorname{is\text{-}module}(m) \wedge (e \subseteq \operatorname{vec}(m))) \Rightarrow (e \subseteq \operatorname{closed\text{-}conv\text{-}hull}(m, t, e)) \end{array}}}
\vnbprop{closed-conv-hull-convex}{\vnbstmt{\begin{array}{@{}l@{}} \forall m, t, e. \\ \quad (\operatorname{is\text{-}module}(m) \wedge (\operatorname{scal}(m) = \operatorname{normed\text{-}field\text{-}as\text{-}commutative\text{-}ring}(\operatorname{rr\text{-}normed\text{-}field}))) \Rightarrow  \\ \quad \quad \operatorname{is\text{-}convex\text{-}subset}(m, \operatorname{closed\text{-}conv\text{-}hull}(m, t, e)) \end{array}}}
\vnbprop{closed-conv-hull-closed}{\vnbstmt{\begin{array}{@{}l@{}} \forall m, t, e. \\ \quad \operatorname{is\text{-}tvs}(\operatorname{rr\text{-}normed\text{-}field}, m, t) \Rightarrow  \\ \quad \quad \operatorname{is\text{-}top\text{-}closed}(t, \operatorname{closed\text{-}conv\text{-}hull}(m, t, e)) \end{array}}}
\vnbprop{extreme-points-nonempty}{\vnbstmt{\begin{array}{@{}l@{}} \forall m, t, s. \\ \quad \operatorname{is\text{-}locally\text{-}convex\text{-}tvs}(m, t) \wedge \\ \quad s \subseteq \operatorname{vec}(m) \wedge \\ \quad \operatorname{is\text{-}convex\text{-}subset}(m, s) \wedge \\ \quad \operatorname{is\text{-}compact\text{-}t}(\operatorname{subspace\text{-}top}(t, s)) \wedge \\ \quad \exists z_{0}.\; (z_{0} \in s) \Rightarrow \\ \quad \quad \exists p.\; (p \in \operatorname{extreme\text{-}points}(m, s)) \end{array}}}
\vnbprop{krein-milman}{\vnbstmt{\begin{array}{@{}l@{}} \forall m, t, s. \\ \quad \operatorname{is\text{-}locally\text{-}convex\text{-}tvs}(m, t) \wedge \\ \quad s \subseteq \operatorname{vec}(m) \wedge \\ \quad \operatorname{is\text{-}convex\text{-}subset}(m, s) \wedge \\ \quad \operatorname{is\text{-}compact\text{-}t}(\operatorname{subspace\text{-}top}(t, s)) \Rightarrow \\ \quad \quad s = \operatorname{closed\text{-}conv\text{-}hull}(m, t, \operatorname{extreme\text{-}points}(m, s)) \end{array}}}
\vnbprop{rmod-combination-sub-point}{\vnbstmt{\begin{array}{@{}l@{}} \forall m, t, a, b, x. \\ \quad \operatorname{is\text{-}module}(m) \wedge \\ \quad \operatorname{scal}(m) = \operatorname{normed\text{-}field\text{-}as\text{-}commutative\text{-}ring}(\operatorname{rr\text{-}normed\text{-}field}) \wedge \\ \quad t \in \mathbb{R} \wedge \\ \quad a \in \operatorname{vec}(m) \wedge \\ \quad b \in \operatorname{vec}(m) \wedge \\ \quad x \in \operatorname{vec}(m) \Rightarrow \\ \quad \quad \operatorname{vadd}(m)(\operatorname{vadd}(m)(\operatorname{act}(m)((1 - t), a), \operatorname{act}(m)(t, b)), \operatorname{vneg}(m)(x)) \\ \quad \quad =\; \operatorname{vadd}(m)(\operatorname{act}(m)((1 - t), \operatorname{vadd}(m)(a, \operatorname{vneg}(m)(x))), \operatorname{act}(m)(t, \operatorname{vadd}(m)(b, \operatorname{vneg}(m)(x)))) \end{array}}}
\vnbprop{sn-oball-convex}{\vnbstmt{\begin{array}{@{}l@{}} \forall m, a, f, b, x, r. \\ \quad (\operatorname{is\text{-}seminorm\text{-}family\text{-}on}(\operatorname{rr\text{-}normed\text{-}field}, m, a, f) \wedge ((b \subseteq a) \wedge ((x \in \operatorname{vec}(m)) \wedge (r \in \mathbb{R})))) \Rightarrow  \\ \quad \quad \operatorname{is\text{-}convex\text{-}subset}(m, \operatorname{sn\text{-}oball}(m, f, b, x, r)) \end{array}}}
\vnbprop{seminorm-top-locally-convex}{\vnbstmt{\begin{array}{@{}l@{}} \forall m, x, f. \\ \quad \operatorname{is\text{-}seminorm\text{-}family\text{-}on}(\operatorname{rr\text{-}normed\text{-}field}, m, x, f) \Rightarrow  \\ \quad \quad \operatorname{is\text{-}locally\text{-}convex\text{-}tvs}(m, \operatorname{seminorm\text{-}top}(m, x, f)) \end{array}}}

\vnbbookfile{theorem-library/haar-fixed-point}{theorem-library/haar-fixed-point}
\noindent{\small haar-fixed-point.scm -- wave of 2026-11-09/10 (agent HAAR-C), wired 2026-11-10 from scratchpad/haarc-1109/haar-fixed-point.scm. Library: analysis.}\par
\vnbprop{nvs-dist-convex-combination}{\vnbstmt{\begin{array}{@{}l@{}} \forall m, a, b, y, t. \\ \quad \operatorname{is\text{-}normed\text{-}vector\text{-}space}(m) \wedge \\ \quad a \in \operatorname{vec}(m) \wedge \\ \quad b \in \operatorname{vec}(m) \wedge \\ \quad y \in \operatorname{vec}(m) \wedge \\ \quad t \in \mathbb{R} \wedge \\ \quad 0 \leq t \wedge \\ \quad t \leq 1 \Rightarrow \\ \quad \quad \operatorname{dist}(\operatorname{nvs\text{-}metric\text{-}space}(m))(\operatorname{vadd}(m)(\operatorname{act}(m)((1 - t), a), \operatorname{act}(m)(t, b)), y) \\ \quad \quad \leq\; ((1 - t) \cdot \operatorname{dist}(\operatorname{nvs\text{-}metric\text{-}space}(m))(a, y)) + (t \cdot \operatorname{dist}(\operatorname{nvs\text{-}metric\text{-}space}(m))(b, y)) \end{array}}}
\vnbprop{nvs-dist-convex-ball}{\vnbstmt{\begin{array}{@{}l@{}} \forall m, a, b, y, t, o. \\ \quad \operatorname{is\text{-}normed\text{-}vector\text{-}space}(m) \wedge \\ \quad a \in \operatorname{vec}(m) \wedge \\ \quad b \in \operatorname{vec}(m) \wedge \\ \quad y \in \operatorname{vec}(m) \wedge \\ \quad t \in \mathbb{R} \wedge \\ \quad 0 \leq t \wedge \\ \quad t \leq 1 \wedge \\ \quad o \in \mathbb{R} \wedge \\ \quad \operatorname{dist}(\operatorname{nvs\text{-}metric\text{-}space}(m))(a, y) \leq o \wedge \\ \quad \operatorname{dist}(\operatorname{nvs\text{-}metric\text{-}space}(m))(b, y) \leq o \Rightarrow \\ \quad \quad \operatorname{dist}(\operatorname{nvs\text{-}metric\text{-}space}(m))(\operatorname{vadd}(m)(\operatorname{act}(m)((1 - t), a), \operatorname{act}(m)(t, b)), y) \leq o \end{array}}}
\vnbprop{metric-closed-is-top-closed}{\vnbstmt{\begin{array}{@{}l@{}} \forall s, c. \\ \quad (\operatorname{is\text{-}metric\text{-}space}(s) \wedge \operatorname{is\text{-}closed}(s, c)) \Rightarrow \operatorname{is\text{-}top\text{-}closed}(\operatorname{metric\text{-}top}(s), c) \end{array}}}
\vnbprop{metric-compact-subspace-top-compact}{\vnbstmt{\begin{array}{@{}l@{}} \forall s. \\ \quad \operatorname{is\text{-}compact}(s) \Rightarrow \operatorname{is\text{-}compact\text{-}t}(\operatorname{subspace\text{-}top}(\operatorname{metric\text{-}top}(s), \operatorname{pts}(s))) \end{array}}}
\vnbprop{compact-chain-meet-metric}{\vnbstmt{\begin{array}{@{}l@{}} \forall s, h. \\ \quad \operatorname{is\text{-}compact}(s) \wedge \\ \quad \exists a.\; (a \in h) \wedge \\ \quad \forall a \in h.\; \operatorname{is\text{-}closed}(s, a) \wedge \\ \quad \forall a \in h.\; \exists z.\; (z \in a) \wedge \\ \quad \forall a, b \in h.\; ((a \subseteq b) \vee (b \subseteq a)) \Rightarrow \\ \quad \quad \exists p \in \operatorname{pts}(s).\; \forall a \in h.\; (p \in a) \end{array}}}
\vnbprop{metric-closed-meet}{\vnbstmt{\begin{array}{@{}l@{}} \forall s, h. \\ \quad (\operatorname{is\text{-}metric\text{-}space}(s) \wedge \forall a \in h.\; \operatorname{is\text{-}closed}(s, a)) \Rightarrow  \\ \quad \quad \operatorname{is\text{-}closed}(s, \{u \in \operatorname{pts}(s) : \forall b \in h.\; (u \in b)\}) \end{array}}}
\vnbprop{convex-meet}{\vnbstmt{\begin{array}{@{}l@{}} \forall m, a, h. \\ \quad \operatorname{is\text{-}normed\text{-}vector\text{-}space}(m) \wedge \\ \quad \operatorname{is\text{-}convex\text{-}subset}(m, a) \wedge \\ \quad \forall \mathit{kca} \in h.\; \operatorname{is\text{-}convex\text{-}subset}(m, \mathit{kca}) \Rightarrow \\ \quad \quad \operatorname{is\text{-}convex\text{-}subset}(m, \{u \in a : \forall b \in h.\; (u \in b)\}) \end{array}}}
\vnbprop{invariant-ccs-minimal}{\vnbstmt{\begin{array}{@{}l@{}} \forall m, a, s. \\ \quad \operatorname{is\text{-}normed\text{-}vector\text{-}space}(m) \wedge \\ \quad a \subseteq \operatorname{vec}(m) \wedge \\ \quad \operatorname{is\text{-}compact}(\operatorname{subspace\text{-}ms}(\operatorname{nvs\text{-}metric\text{-}space}(m), a)) \wedge \\ \quad \operatorname{is\text{-}invariant\text{-}ccs}(m, a, s, a) \Rightarrow \\ \quad \quad \exists b. \\ \quad \quad \quad \operatorname{is\text{-}invariant\text{-}ccs}(m, a, s, b) \wedge \\ \quad \quad \quad \forall c.\; ((\operatorname{is\text{-}invariant\text{-}ccs}(m, a, s, c) \wedge (c \subseteq b)) \Rightarrow (b \subseteq c)) \end{array}}}
\vnbprop{nvs-dist-cap-closed}{\vnbstmt{\begin{array}{@{}l@{}} \forall m, a, b, o. \\ \quad \operatorname{is\text{-}normed\text{-}vector\text{-}space}(m) \wedge \\ \quad a \subseteq \operatorname{vec}(m) \wedge \\ \quad b \subseteq a \wedge \\ \quad \operatorname{is\text{-}closed}(\operatorname{subspace\text{-}ms}(\operatorname{nvs\text{-}metric\text{-}space}(m), a), b) \wedge \\ \quad o \in \mathbb{R} \Rightarrow \\ \quad \quad \operatorname{is\text{-}closed}(\operatorname{subspace\text{-}ms}(\operatorname{nvs\text{-}metric\text{-}space}(m), a), \\ \quad \quad \quad \{z \in b : \forall y \in b.\; (\operatorname{dist}(\operatorname{nvs\text{-}metric\text{-}space}(m))(z, y) \leq o)\}) \end{array}}}
\vnbprop{nvs-diameter-sup}{\vnbstmt{\begin{array}{@{}l@{}} \forall m, a, b. \\ \quad \operatorname{is\text{-}normed\text{-}vector\text{-}space}(m) \wedge \\ \quad a \subseteq \operatorname{vec}(m) \wedge \\ \quad \operatorname{is\text{-}compact}(\operatorname{subspace\text{-}ms}(\operatorname{nvs\text{-}metric\text{-}space}(m), a)) \wedge \\ \quad b \subseteq a \Rightarrow \\ \quad \quad \exists d \in \mathbb{R}. \\ \quad \quad \quad \forall x, y \in b.\; (\operatorname{dist}(\operatorname{nvs\text{-}metric\text{-}space}(m))(x, y) \leq d) \wedge \\ \quad \quad \quad \forall e \in \mathbb{R}. \\ \quad \quad \quad \quad ((0 \leq e) \wedge (e < d)) \Rightarrow  \\ \quad \quad \quad \quad \quad \exists x, y \in b.\; (e < \operatorname{dist}(\operatorname{nvs\text{-}metric\text{-}space}(m))(x, y)) \end{array}}}
\vnbprop{nvs-compact-subset-net}{\vnbstmt{\begin{array}{@{}l@{}} \forall m, a, b, d. \\ \quad \operatorname{is\text{-}normed\text{-}vector\text{-}space}(m) \wedge \\ \quad a \subseteq \operatorname{vec}(m) \wedge \\ \quad \operatorname{is\text{-}compact}(\operatorname{subspace\text{-}ms}(\operatorname{nvs\text{-}metric\text{-}space}(m), a)) \wedge \\ \quad b \subseteq a \wedge \\ \quad \operatorname{is\text{-}closed}(\operatorname{subspace\text{-}ms}(\operatorname{nvs\text{-}metric\text{-}space}(m), a), b) \wedge \\ \quad d \in \mathbb{R} \wedge \\ \quad 0 < d \Rightarrow \\ \quad \quad \exists f \in \mathbf{Set}. \\ \quad \quad \quad |f| \in \mathbb{N} \wedge \\ \quad \quad \quad f \subseteq b \wedge \\ \quad \quad \quad \forall p \in b.\; \exists c \in f.\; (\operatorname{dist}(\operatorname{nvs\text{-}metric\text{-}space}(m))(c, p) \leq d) \end{array}}}
\vnbprop{nvs-net-average}{\vnbstmt{\begin{array}{@{}l@{}} \forall m, b, a, d.; \forall s \in \mathbf{Set}. \\ \quad \operatorname{is\text{-}normed\text{-}vector\text{-}space}(m) \wedge \\ \quad b \subseteq \operatorname{vec}(m) \wedge \\ \quad \operatorname{is\text{-}convex\text{-}subset}(m, b) \wedge \\ \quad \exists z.\; (z \in b) \wedge \\ \quad a \in \mathbb{R} \wedge \\ \quad d \in \mathbb{R} \wedge \\ \quad d < a \wedge \\ \quad \forall x, y \in b.\; (\operatorname{dist}(\operatorname{nvs\text{-}metric\text{-}space}(m))(x, y) \leq a) \wedge \\ \quad |s| \in \mathbb{N} \wedge \\ \quad s \subseteq b \Rightarrow \\ \quad \quad \exists u \in b, \mathit{hvd} \in \mathbb{R}. \\ \quad \quad \quad 0 < \mathit{hvd} \wedge \\ \quad \quad \quad \mathit{hvd} \leq (a - d) \wedge \\ \quad \quad \quad \forall c \in s, \mathit{hvy} \in b. \\ \quad \quad \quad \quad (\operatorname{dist}(\operatorname{nvs\text{-}metric\text{-}space}(m))(c, \mathit{hvy}) \leq d) \Rightarrow  \\ \quad \quad \quad \quad \quad \operatorname{dist}(\operatorname{nvs\text{-}metric\text{-}space}(m))(u, \mathit{hvy}) \leq (a - \mathit{hvd}) \end{array}}}
\vnbprop{invariant-ccs-minimal-diameter-zero}{\vnbstmt{\begin{array}{@{}l@{}} \forall m, a, s, b.; \forall x, y \in b. \\ \quad \operatorname{is\text{-}normed\text{-}vector\text{-}space}(m) \wedge \\ \quad a \subseteq \operatorname{vec}(m) \wedge \\ \quad \operatorname{is\text{-}compact}(\operatorname{subspace\text{-}ms}(\operatorname{nvs\text{-}metric\text{-}space}(m), a)) \wedge \\ \quad \operatorname{is\text{-}affine\text{-}isometry\text{-}group}(m, a, s) \wedge \\ \quad \operatorname{is\text{-}invariant\text{-}ccs}(m, a, s, b) \wedge \\ \quad \forall c.\; ((\operatorname{is\text{-}invariant\text{-}ccs}(m, a, s, c) \wedge (c \subseteq b)) \Rightarrow (b \subseteq c)) \Rightarrow \\ \quad \quad \operatorname{dist}(\operatorname{nvs\text{-}metric\text{-}space}(m))(x, y) \leq 0 \end{array}}}
\vnbprop{affine-isometry-fixed-point}{\vnbstmt{\begin{array}{@{}l@{}} \forall m, a, s. \\ \quad \operatorname{is\text{-}normed\text{-}vector\text{-}space}(m) \wedge \\ \quad a \subseteq \operatorname{vec}(m) \wedge \\ \quad \operatorname{is\text{-}convex\text{-}subset}(m, a) \wedge \\ \quad \operatorname{is\text{-}compact}(\operatorname{subspace\text{-}ms}(\operatorname{nvs\text{-}metric\text{-}space}(m), a)) \wedge \\ \quad \exists z.\; (z \in a) \wedge \\ \quad \operatorname{is\text{-}affine\text{-}isometry\text{-}group}(m, a, s) \Rightarrow \\ \quad \quad \exists w \in a.\; \forall g \in s.\; (g(w) = w) \end{array}}}

\vnbbookfile{theorem-library/cg4-cosets}{theorem-library/cg4-cosets}
\noindent{\small cg4-cosets.scm -- a compact topological group is the union of finitely many left cosets of an open subgroup (compact-groups notes, Lemma 5.33: "G/H is finite"; stated without the quotient G/H, and needing neither the invariant metric nor normality of H). Agent ISOM, 2026-10-26.  Helper prefix: cc4-.  Library: analysis.}\par
\vnbprop{pts-left-cosets-unfold}{\vnbstmt{\begin{array}{@{}l@{}} \forall s, g, h. \\ \quad (\operatorname{pts\text{-}left\text{-}cosets}(s, g, h) \simeq \{u \in \operatorname{power}(\operatorname{pts}(s)) : \exists x \in \operatorname{carr}(g).\; (u = \{z \in \operatorname{pts}(s) : (\operatorname{opr}(g)(\operatorname{inv}(g)(x), z) \in h)\})\}) \end{array}}}
\vnbprop{top-group-compact-finitely-many-cosets-gt}{\vnbstmt{\begin{array}{@{}l@{}} \forall s, g, h. \\ \quad (\operatorname{is\text{-}group\text{-}topology}(s, g) \wedge (\operatorname{is\text{-}compact\text{-}t}(s) \wedge (\operatorname{is\text{-}subgroup}(g, h) \wedge (h \in \operatorname{opens}(s))))) \Rightarrow  \\ \quad \quad \exists f. \\ \quad \quad \quad ((f \subseteq \operatorname{pts\text{-}left\text{-}cosets}(s, g, h)) \wedge ((|f| \in \mathbb{N}) \wedge (\operatorname{big\text{-}union}(u, f, u) \simeq \operatorname{pts}(s)))) \end{array}}}
\vnbprop{top-group-compact-finitely-many-cosets}{\vnbstmt{\begin{array}{@{}l@{}} \forall g, h. \\ \quad \operatorname{is\text{-}top\text{-}group}(g) \wedge \\ \quad \operatorname{is\text{-}compact\text{-}t}(\operatorname{top\text{-}group\text{-}space}(g)) \wedge \\ \quad \operatorname{is\text{-}subgroup}(\operatorname{top\text{-}group\text{-}group}(g), h) \wedge \\ \quad h \in \operatorname{opens}(\operatorname{top\text{-}group\text{-}space}(g)) \Rightarrow \\ \quad \quad \exists f. \\ \quad \quad \quad f \subseteq \operatorname{pts\text{-}left\text{-}cosets}(\operatorname{top\text{-}group\text{-}space}(g), \operatorname{top\text{-}group\text{-}group}(g), h) \wedge \\ \quad \quad \quad |f| \in \mathbb{N} \wedge \\ \quad \quad \quad (\operatorname{big\text{-}union}(u, f, u) \simeq \operatorname{pts}(\operatorname{top\text{-}group\text{-}space}(g))) \end{array}}}

\vnbbookfile{theorem-library/cg4-opg}{theorem-library/cg4-opg}
\noindent{\small cg4-opg.scm -- one-parameter groups (compact-groups notes, sec. 5.4, Proposition 5.36, the uniqueness half): two one-parameter groups extending the same group segment are equal. Agent ISOM, 2026-10-26.  Helper prefix: co4-.  Library: analysis.}\par
\vnbprop{group-idempotent-iden}{\vnbstmt{\begin{array}{@{}l@{}} \forall g, x. \\ \quad (\operatorname{is\text{-}group}(g) \wedge ((x \in \operatorname{carr}(g)) \wedge (\operatorname{opr}(g)(x, x) = x))) \Rightarrow (x = \operatorname{iden}(g)) \end{array}}}
\vnbprop{group-right-inverse-unique}{\vnbstmt{\begin{array}{@{}l@{}} \forall g, u, v. \\ \quad (\operatorname{is\text{-}group}(g) \wedge ((u \in \operatorname{carr}(g)) \wedge ((v \in \operatorname{carr}(g)) \wedge (\operatorname{opr}(g)(u, v) = \operatorname{iden}(g))))) \Rightarrow  \\ \quad \quad u = \operatorname{inv}(g)(v) \end{array}}}
\vnbprop{one-parameter-agree-multiples}{\vnbstmt{\begin{array}{@{}l@{}} \forall n \in \mathbb{N}, g, l, a, b.; \forall r \in \mathbb{R}. \\ \quad \operatorname{is\text{-}group}(g) \wedge \\ \quad \operatorname{pos\text{-}rr}(l) \wedge \\ \quad a \in (\mathbb{R} \to \operatorname{carr}(g)) \wedge \\ \quad b \in (\mathbb{R} \to \operatorname{carr}(g)) \wedge \\ \quad \forall \mathit{oqr}, u \in \mathbb{R}.\; (a((\mathit{oqr} + u)) = \operatorname{opr}(g)(a(\mathit{oqr}), a(u))) \wedge \\ \quad \forall \mathit{oqr}, u \in \mathbb{R}.\; (b((\mathit{oqr} + u)) = \operatorname{opr}(g)(b(\mathit{oqr}), b(u))) \wedge \\ \quad \forall t \in \mathbb{R}.\; (((0 \leq t) \wedge (t \leq l)) \Rightarrow (a(t) = b(t))) \wedge \\ \quad 0 \leq r \wedge \\ \quad r \leq l \Rightarrow \\ \quad \quad a(((n \cdot l) + r)) = b(((n \cdot l) + r)) \end{array}}}
\vnbprop{one-parameter-group-extension-unique-gt}{\vnbstmt{\begin{array}{@{}l@{}} \forall s, g, l, i, a, b.; \forall t \in \mathbb{R}. \\ \quad \operatorname{is\text{-}group\text{-}topology}(s, g) \wedge \\ \quad \operatorname{is\text{-}group\text{-}segment}(s, g, l, i) \wedge \\ \quad \operatorname{is\text{-}one\text{-}parameter\text{-}group}(s, g, a) \wedge \\ \quad \operatorname{is\text{-}one\text{-}parameter\text{-}group}(s, g, b) \wedge \\ \quad \forall t \in \operatorname{ccint}(0, l).\; (a(t) = i(t)) \wedge \\ \quad \forall t \in \operatorname{ccint}(0, l).\; (b(t) = i(t)) \Rightarrow \\ \quad \quad a(t) = b(t) \end{array}}}
\vnbprop{one-parameter-group-extension-unique}{\vnbstmt{\begin{array}{@{}l@{}} \forall g, l, i, a, b.; \forall t \in \mathbb{R}. \\ \quad \operatorname{is\text{-}top\text{-}group}(g) \wedge \\ \quad \operatorname{is\text{-}group\text{-}segment}(\operatorname{top\text{-}group\text{-}space}(g), \operatorname{top\text{-}group\text{-}group}(g), l, i) \wedge \\ \quad \operatorname{is\text{-}one\text{-}parameter\text{-}group}(\operatorname{top\text{-}group\text{-}space}(g), \operatorname{top\text{-}group\text{-}group}(g), a) \wedge \\ \quad \operatorname{is\text{-}one\text{-}parameter\text{-}group}(\operatorname{top\text{-}group\text{-}space}(g), \operatorname{top\text{-}group\text{-}group}(g), b) \wedge \\ \quad \forall t \in \operatorname{ccint}(0, l).\; (a(t) = i(t)) \wedge \\ \quad \forall t \in \operatorname{ccint}(0, l).\; (b(t) = i(t)) \Rightarrow \\ \quad \quad a(t) = b(t) \end{array}}}

\vnbbookfile{theorem-library/cg4-seg}{theorem-library/cg4-seg}
\noindent{\small cg4-seg.scm -- the algebra of a one-parameter group segment phi on [0, al] (compact-groups notes, Prop. 5.36): phi(0) = e; phi(r) phi(al - r) = phi(al) = phi(al - r) phi(r); phi(r) commutes with c = phi(al) and its powers; phi(r1) phi(r2) = c phi(r1 + r2 - al) past al. Agent ISOM, 2026-10-26.  Helper prefix: sg4-.  ...}\par
\vnbprop{segment-at-zero}{\vnbstmt{\begin{array}{@{}l@{}} \forall g, l, i. \\ \quad \operatorname{is\text{-}group}(g) \wedge \\ \quad l \in \mathbb{R} \wedge \\ \quad 0 < l \wedge \\ \quad i \in (\operatorname{ccint}(0, l) \to \operatorname{carr}(g)) \wedge \\ \quad \forall r, u \in \operatorname{ccint}(0, l). \\ \quad \quad ((r + u) \leq l) \Rightarrow (i((r + u)) = \operatorname{opr}(g)(i(r), i(u))) \Rightarrow \\ \quad \quad i(0) = \operatorname{iden}(g) \end{array}}}
\vnbprop{segment-complement}{\vnbstmt{\begin{array}{@{}l@{}} \forall g, l, i, r. \\ \quad \operatorname{is\text{-}group}(g) \wedge \\ \quad l \in \mathbb{R} \wedge \\ \quad 0 < l \wedge \\ \quad i \in (\operatorname{ccint}(0, l) \to \operatorname{carr}(g)) \wedge \\ \quad \forall \mathit{gsr}, u \in \operatorname{ccint}(0, l). \\ \quad \quad ((\mathit{gsr} + u) \leq l) \Rightarrow (i((\mathit{gsr} + u)) = \operatorname{opr}(g)(i(\mathit{gsr}), i(u))) \wedge \\ \quad r \in \mathbb{R} \wedge \\ \quad 0 \leq r \wedge \\ \quad r \leq l \Rightarrow \\ \quad \quad (\operatorname{opr}(g)(i(r), i((l - r))) = i(l)) \wedge (\operatorname{opr}(g)(i((l - r)), i(r)) = i(l)) \end{array}}}
\vnbprop{segment-commutes-with-end}{\vnbstmt{\begin{array}{@{}l@{}} \forall g, l, i, r. \\ \quad \operatorname{is\text{-}group}(g) \wedge \\ \quad l \in \mathbb{R} \wedge \\ \quad 0 < l \wedge \\ \quad i \in (\operatorname{ccint}(0, l) \to \operatorname{carr}(g)) \wedge \\ \quad \forall \mathit{gsr}, u \in \operatorname{ccint}(0, l). \\ \quad \quad ((\mathit{gsr} + u) \leq l) \Rightarrow (i((\mathit{gsr} + u)) = \operatorname{opr}(g)(i(\mathit{gsr}), i(u))) \wedge \\ \quad r \in \mathbb{R} \wedge \\ \quad 0 \leq r \wedge \\ \quad r \leq l \Rightarrow \\ \quad \quad \operatorname{opr}(g)(i(r), i(l)) = \operatorname{opr}(g)(i(l), i(r)) \end{array}}}
\vnbprop{segment-overflow}{\vnbstmt{\begin{array}{@{}l@{}} \forall g, l, i, a, b. \\ \quad \operatorname{is\text{-}group}(g) \wedge \\ \quad l \in \mathbb{R} \wedge \\ \quad 0 < l \wedge \\ \quad i \in (\operatorname{ccint}(0, l) \to \operatorname{carr}(g)) \wedge \\ \quad \forall r, u \in \operatorname{ccint}(0, l). \\ \quad \quad ((r + u) \leq l) \Rightarrow (i((r + u)) = \operatorname{opr}(g)(i(r), i(u))) \wedge \\ \quad a \in \mathbb{R} \wedge \\ \quad 0 \leq a \wedge \\ \quad a \leq l \wedge \\ \quad b \in \mathbb{R} \wedge \\ \quad 0 \leq b \wedge \\ \quad b \leq l \wedge \\ \quad l \leq (a + b) \Rightarrow \\ \quad \quad \operatorname{opr}(g)(i(a), i(b)) = \operatorname{opr}(g)(i(l), i(((a + b) - l))) \end{array}}}

\vnbbookfile{theorem-library/cg4-rep}{theorem-library/cg4-rep}
\noindent{\small cg4-rep.scm -- the extension psi of a group segment phi (Prop. 5.36): every real t has exactly one representation t = n al + r (n in NN, 0 $<$= r $<$ al) on its side of 0, so OPG-REP(t, y) holds for exactly one y, and OPG-VAL(t) is that y. Agent ISOM, 2026-10-26.  Helper prefix: rp4-.  Library: analysis.}\par
\vnbprop{representation-index-unique}{\vnbstmt{\begin{array}{@{}l@{}} \forall l, a, b, c, d. \\ \quad l \in \mathbb{R} \wedge \\ \quad 0 < l \wedge \\ \quad a \in \mathbb{N} \wedge \\ \quad c \in \mathbb{N} \wedge \\ \quad b \in \mathbb{R} \wedge \\ \quad 0 \leq b \wedge \\ \quad b < l \wedge \\ \quad d \in \mathbb{R} \wedge \\ \quad 0 \leq d \wedge \\ \quad d < l \wedge \\ \quad ((a \cdot l) + b) = ((c \cdot l) + d) \Rightarrow \\ \quad \quad (a = c) \wedge (b = d) \end{array}}}
\vnbprop{opg-rep-exists}{\vnbstmt{\begin{array}{@{}l@{}} \forall g, l, i, t. \\ \quad \operatorname{is\text{-}group}(g) \wedge \\ \quad l \in \mathbb{R} \wedge \\ \quad 0 < l \wedge \\ \quad i \in (\operatorname{ccint}(0, l) \to \operatorname{carr}(g)) \wedge \\ \quad \forall r, u \in \operatorname{ccint}(0, l). \\ \quad \quad ((r + u) \leq l) \Rightarrow (i((r + u)) = \operatorname{opr}(g)(i(r), i(u))) \wedge \\ \quad t \in \mathbb{R} \Rightarrow \\ \quad \quad \exists y \in \operatorname{carr}(g).\; \operatorname{opg\text{-}rep}(g, l, i, t, y) \end{array}}}
\vnbprop{opg-rep-unique}{\vnbstmt{\begin{array}{@{}l@{}} \forall g, l, i, t, a, b. \\ \quad \operatorname{is\text{-}group}(g) \wedge \\ \quad l \in \mathbb{R} \wedge \\ \quad 0 < l \wedge \\ \quad i \in (\operatorname{ccint}(0, l) \to \operatorname{carr}(g)) \wedge \\ \quad \forall r, u \in \operatorname{ccint}(0, l). \\ \quad \quad ((r + u) \leq l) \Rightarrow (i((r + u)) = \operatorname{opr}(g)(i(r), i(u))) \wedge \\ \quad t \in \mathbb{R} \wedge \\ \quad \operatorname{opg\text{-}rep}(g, l, i, t, a) \wedge \\ \quad \operatorname{opg\text{-}rep}(g, l, i, t, b) \Rightarrow \\ \quad \quad a = b \end{array}}}
\vnbprop{opg-val-value}{\vnbstmt{\begin{array}{@{}l@{}} \forall g, l, i, t, y. \\ \quad \operatorname{is\text{-}group}(g) \wedge \\ \quad l \in \mathbb{R} \wedge \\ \quad 0 < l \wedge \\ \quad i \in (\operatorname{ccint}(0, l) \to \operatorname{carr}(g)) \wedge \\ \quad \forall r, u \in \operatorname{ccint}(0, l). \\ \quad \quad ((r + u) \leq l) \Rightarrow (i((r + u)) = \operatorname{opr}(g)(i(r), i(u))) \wedge \\ \quad t \in \mathbb{R} \wedge \\ \quad y \in \operatorname{carr}(g) \wedge \\ \quad \operatorname{opg\text{-}rep}(g, l, i, t, y) \Rightarrow \\ \quad \quad \operatorname{opg\text{-}val}(g, l, i, t) = y \end{array}}}
\vnbprop{opg-val-rep}{\vnbstmt{\begin{array}{@{}l@{}} \forall g, l, i, t. \\ \quad \operatorname{is\text{-}group}(g) \wedge \\ \quad l \in \mathbb{R} \wedge \\ \quad 0 < l \wedge \\ \quad i \in (\operatorname{ccint}(0, l) \to \operatorname{carr}(g)) \wedge \\ \quad \forall r, u \in \operatorname{ccint}(0, l). \\ \quad \quad ((r + u) \leq l) \Rightarrow (i((r + u)) = \operatorname{opr}(g)(i(r), i(u))) \wedge \\ \quad t \in \mathbb{R} \Rightarrow \\ \quad \quad (\operatorname{opg\text{-}val}(g, l, i, t) \in \operatorname{carr}(g)) \wedge \operatorname{opg\text{-}rep}(g, l, i, t, \operatorname{opg\text{-}val}(g, l, i, t)) \end{array}}}

\vnbbookfile{theorem-library/cg4-hom}{theorem-library/cg4-hom}
\noindent{\small cg4-hom.scm -- the extension psi of a group segment is a homomorphism R -$>$ G (Prop. 5.36): first on t $>$= 0 (psi(n al + r) = c\textasciicircum{}n phi(r), the overflow of phi past al), then on all of R through psi(t) = psi(-t)\textasciicircum{}-1 for t $<$ 0. Agent ISOM, 2026-10-26/27.  Helper prefix: hm4-.  Library: analysis.}\par
\vnbprop{gpow-one}{\vnbstmt{\forall g, x.\; ((\operatorname{is\text{-}group}(g) \wedge (x \in \operatorname{carr}(g))) \Rightarrow (\operatorname{gpow}(g, x, 1) = x))}}
\vnbprop{group-interchange}{\vnbstmt{\begin{array}{@{}l@{}} \forall g, a, b, c, d. \\ \quad \operatorname{is\text{-}group}(g) \wedge \\ \quad a \in \operatorname{carr}(g) \wedge \\ \quad b \in \operatorname{carr}(g) \wedge \\ \quad c \in \operatorname{carr}(g) \wedge \\ \quad d \in \operatorname{carr}(g) \wedge \\ \quad \operatorname{opr}(g)(b, c) = \operatorname{opr}(g)(c, b) \Rightarrow \\ \quad \quad \operatorname{opr}(g)(\operatorname{opr}(g)(a, b), \operatorname{opr}(g)(c, d)) = \operatorname{opr}(g)(\operatorname{opr}(g)(a, c), \operatorname{opr}(g)(b, d)) \end{array}}}
\vnbprop{opg-val-nonneg}{\vnbstmt{\begin{array}{@{}l@{}} \forall g, l, i, n, r, t. \\ \quad \operatorname{is\text{-}group}(g) \wedge \\ \quad l \in \mathbb{R} \wedge \\ \quad 0 < l \wedge \\ \quad i \in (\operatorname{ccint}(0, l) \to \operatorname{carr}(g)) \wedge \\ \quad \forall \mathit{gsr}, u \in \operatorname{ccint}(0, l). \\ \quad \quad ((\mathit{gsr} + u) \leq l) \Rightarrow (i((\mathit{gsr} + u)) = \operatorname{opr}(g)(i(\mathit{gsr}), i(u))) \wedge \\ \quad n \in \mathbb{N} \wedge \\ \quad r \in \mathbb{R} \wedge \\ \quad 0 \leq r \wedge \\ \quad r < l \wedge \\ \quad t \in \mathbb{R} \wedge \\ \quad 0 \leq t \wedge \\ \quad t = ((n \cdot l) + r) \Rightarrow \\ \quad \quad \operatorname{opg\text{-}val}(g, l, i, t) = \operatorname{opr}(g)(\operatorname{gpow}(g, i(l), n), i(r)) \end{array}}}
\vnbprop{opg-val-additive-nonneg}{\vnbstmt{\begin{array}{@{}l@{}} \forall g, l, i, a, b. \\ \quad \operatorname{is\text{-}group}(g) \wedge \\ \quad l \in \mathbb{R} \wedge \\ \quad 0 < l \wedge \\ \quad i \in (\operatorname{ccint}(0, l) \to \operatorname{carr}(g)) \wedge \\ \quad \forall r, u \in \operatorname{ccint}(0, l). \\ \quad \quad ((r + u) \leq l) \Rightarrow (i((r + u)) = \operatorname{opr}(g)(i(r), i(u))) \wedge \\ \quad a \in \mathbb{R} \wedge \\ \quad 0 \leq a \wedge \\ \quad b \in \mathbb{R} \wedge \\ \quad 0 \leq b \Rightarrow \\ \quad \quad \operatorname{opg\text{-}val}(g, l, i, (a + b)) = \operatorname{opr}(g)(\operatorname{opg\text{-}val}(g, l, i, a), \operatorname{opg\text{-}val}(g, l, i, b)) \end{array}}}
\vnbprop{group-solve-factor}{\vnbstmt{\begin{array}{@{}l@{}} \forall g, a, b, c. \\ \quad (\operatorname{is\text{-}group}(g) \wedge ((a \in \operatorname{carr}(g)) \wedge ((b \in \operatorname{carr}(g)) \wedge ((c \in \operatorname{carr}(g)) \wedge (a = \operatorname{opr}(g)(b, c)))))) \Rightarrow  \\ \quad \quad (b = \operatorname{opr}(g)(a, \operatorname{inv}(g)(c))) \wedge (c = \operatorname{opr}(g)(\operatorname{inv}(g)(b), a)) \end{array}}}
\vnbprop{group-inv-quotient}{\vnbstmt{\begin{array}{@{}l@{}} \forall g, a, b. \\ \quad (\operatorname{is\text{-}group}(g) \wedge ((a \in \operatorname{carr}(g)) \wedge (b \in \operatorname{carr}(g)))) \Rightarrow  \\ \quad \quad \operatorname{inv}(g)(\operatorname{opr}(g)(a, \operatorname{inv}(g)(b))) = \operatorname{opr}(g)(b, \operatorname{inv}(g)(a)) \wedge \\ \quad \quad \operatorname{inv}(g)(\operatorname{opr}(g)(\operatorname{inv}(g)(b), a)) = \operatorname{opr}(g)(\operatorname{inv}(g)(a), b) \end{array}}}
\vnbprop{opg-val-negative}{\vnbstmt{\begin{array}{@{}l@{}} \forall g, l, i, t. \\ \quad \operatorname{is\text{-}group}(g) \wedge \\ \quad l \in \mathbb{R} \wedge \\ \quad 0 < l \wedge \\ \quad i \in (\operatorname{ccint}(0, l) \to \operatorname{carr}(g)) \wedge \\ \quad \forall r, u \in \operatorname{ccint}(0, l). \\ \quad \quad ((r + u) \leq l) \Rightarrow (i((r + u)) = \operatorname{opr}(g)(i(r), i(u))) \wedge \\ \quad t \in \mathbb{R} \wedge \\ \quad t < 0 \Rightarrow \\ \quad \quad \operatorname{opg\text{-}val}(g, l, i, t) = \operatorname{inv}(g)(\operatorname{opg\text{-}val}(g, l, i, (0 - t))) \end{array}}}
\vnbprop{opg-val-additive}{\vnbstmt{\begin{array}{@{}l@{}} \forall g, l, i, a, b. \\ \quad \operatorname{is\text{-}group}(g) \wedge \\ \quad l \in \mathbb{R} \wedge \\ \quad 0 < l \wedge \\ \quad i \in (\operatorname{ccint}(0, l) \to \operatorname{carr}(g)) \wedge \\ \quad \forall r, u \in \operatorname{ccint}(0, l). \\ \quad \quad ((r + u) \leq l) \Rightarrow (i((r + u)) = \operatorname{opr}(g)(i(r), i(u))) \wedge \\ \quad a \in \mathbb{R} \wedge \\ \quad b \in \mathbb{R} \Rightarrow \\ \quad \quad \operatorname{opg\text{-}val}(g, l, i, (a + b)) = \operatorname{opr}(g)(\operatorname{opg\text{-}val}(g, l, i, a), \operatorname{opg\text{-}val}(g, l, i, b)) \end{array}}}

\vnbbookfile{theorem-library/cg4-ext}{theorem-library/cg4-ext}
\noindent{\small cg4-ext.scm -- Proposition 5.36, existence: the extension psi = OPG-EXT(g, al, phi) of a group segment phi is a one-parameter group (continuous: at t by continuity at 0, translated; at 0 by phi's continuity at 0 on the right and, through the inverse, on the left) and agrees with phi on [0, al].  Proved for a group ...}\par
\vnbprop{opg-ext-fun}{\vnbstmt{\begin{array}{@{}l@{}} \forall g, l, i. \\ \quad \operatorname{is\text{-}group}(g) \wedge \\ \quad l \in \mathbb{R} \wedge \\ \quad 0 < l \wedge \\ \quad i \in (\operatorname{ccint}(0, l) \to \operatorname{carr}(g)) \wedge \\ \quad \forall r, u \in \operatorname{ccint}(0, l). \\ \quad \quad ((r + u) \leq l) \Rightarrow (i((r + u)) = \operatorname{opr}(g)(i(r), i(u))) \Rightarrow \\ \quad \quad \operatorname{opg\text{-}ext}(g, l, i) \in (\mathbb{R} \to \operatorname{carr}(g)) \end{array}}}
\vnbprop{opg-ext-apply}{\vnbstmt{\begin{array}{@{}l@{}} \forall g, l, i, t. \\ \quad \operatorname{is\text{-}group}(g) \wedge \\ \quad l \in \mathbb{R} \wedge \\ \quad 0 < l \wedge \\ \quad i \in (\operatorname{ccint}(0, l) \to \operatorname{carr}(g)) \wedge \\ \quad \forall r, u \in \operatorname{ccint}(0, l). \\ \quad \quad ((r + u) \leq l) \Rightarrow (i((r + u)) = \operatorname{opr}(g)(i(r), i(u))) \wedge \\ \quad t \in \mathbb{R} \Rightarrow \\ \quad \quad \operatorname{opg\text{-}ext}(g, l, i)(t) = \operatorname{opg\text{-}val}(g, l, i, t) \end{array}}}
\vnbprop{opg-ext-extends}{\vnbstmt{\begin{array}{@{}l@{}} \forall g, l, i, t. \\ \quad \operatorname{is\text{-}group}(g) \wedge \\ \quad l \in \mathbb{R} \wedge \\ \quad 0 < l \wedge \\ \quad i \in (\operatorname{ccint}(0, l) \to \operatorname{carr}(g)) \wedge \\ \quad \forall r, u \in \operatorname{ccint}(0, l). \\ \quad \quad ((r + u) \leq l) \Rightarrow (i((r + u)) = \operatorname{opr}(g)(i(r), i(u))) \wedge \\ \quad t \in \operatorname{ccint}(0, l) \Rightarrow \\ \quad \quad \operatorname{opg\text{-}ext}(g, l, i)(t) = i(t) \end{array}}}
\vnbprop{opg-val-continuous-gt}{\vnbstmt{\begin{array}{@{}l@{}} \forall s, g, l, i.; \forall t \in \mathbb{R}, w \in \operatorname{opens}(s). \\ \quad (\operatorname{is\text{-}group\text{-}topology}(s, g) \wedge (\operatorname{is\text{-}group\text{-}segment}(s, g, l, i) \wedge (\operatorname{opg\text{-}val}(g, l, i, t) \in w))) \Rightarrow  \\ \quad \quad \exists d. \\ \quad \quad \quad \operatorname{pos\text{-}rr}(d) \wedge \forall v \in \mathbb{R}.\; ((\lvert (v - t) \rvert < d) \Rightarrow (\operatorname{opg\text{-}val}(g, l, i, v) \in w)) \end{array}}}
\vnbprop{one-parameter-group-extension-exists-gt}{\vnbstmt{\begin{array}{@{}l@{}} \forall s, g, l, i. \\ \quad (\operatorname{is\text{-}group\text{-}topology}(s, g) \wedge \operatorname{is\text{-}group\text{-}segment}(s, g, l, i)) \Rightarrow  \\ \quad \quad \exists a. \\ \quad \quad \quad \operatorname{is\text{-}one\text{-}parameter\text{-}group}(s, g, a) \wedge \forall t \in \operatorname{ccint}(0, l).\; (a(t) = i(t)) \end{array}}}
\vnbprop{one-parameter-group-extension-exists}{\vnbstmt{\begin{array}{@{}l@{}} \forall g, l, i. \\ \quad (\operatorname{is\text{-}top\text{-}group}(g) \wedge \operatorname{is\text{-}group\text{-}segment}(\operatorname{top\text{-}group\text{-}space}(g), \operatorname{top\text{-}group\text{-}group}(g), l, i)) \Rightarrow  \\ \quad \quad \exists a. \\ \quad \quad \quad \operatorname{is\text{-}one\text{-}parameter\text{-}group}(\operatorname{top\text{-}group\text{-}space}(g), \operatorname{top\text{-}group\text{-}group}(g), a) \wedge \\ \quad \quad \quad \forall t \in \operatorname{ccint}(0, l).\; (a(t) = i(t)) \end{array}}}
\vnbprop{opg-val-continuous}{\vnbstmt{\begin{array}{@{}l@{}} \forall g, l, i.; \forall t \in \mathbb{R}, w \in \operatorname{opens}(\operatorname{top\text{-}group\text{-}space}(g)). \\ \quad \operatorname{is\text{-}top\text{-}group}(g) \wedge \\ \quad \operatorname{is\text{-}group\text{-}segment}(\operatorname{top\text{-}group\text{-}space}(g), \operatorname{top\text{-}group\text{-}group}(g), l, i) \wedge \\ \quad \operatorname{opg\text{-}val}(\operatorname{top\text{-}group\text{-}group}(g), l, i, t) \in w \Rightarrow \\ \quad \quad \exists d. \\ \quad \quad \quad (\operatorname{pos\text{-}rr}(d) \wedge \forall v \in \mathbb{R}.\; ((\lvert (v - t) \rvert < d) \Rightarrow (\operatorname{opg\text{-}val}(\operatorname{top\text{-}group\text{-}group}(g), l, i, v) \in w))) \end{array}}}

\vnbbookfile{theorem-library/cg4-local}{theorem-library/cg4-local}
\noindent{\small cg4-local.scm -- local homomorphisms (compact-groups notes, Definition 5.37, Proposition 5.38): the composite of local homomorphisms f : G -$>$ H (on V) and h : H -$>$ K (on V\_H) is a local homomorphism on V' = \{ x in V : f(x) in V\_H \}, with W' = \{ x in W : f(x) in W\_H \}. Agent ISOM, 2026-10-27.  Helper prefix: lc4-.  ...}\par
\vnbprop{local-homomorphism-composite-gt}{\vnbstmt{\begin{array}{@{}l@{}} \forall a, b, c, d, g, m, v, f, n, h. \\ \quad \operatorname{is\text{-}group\text{-}topology}(a, b) \wedge \\ \quad \operatorname{is\text{-}group\text{-}topology}(c, d) \wedge \\ \quad \operatorname{is\text{-}group\text{-}topology}(g, m) \wedge \\ \quad \operatorname{is\text{-}local\text{-}hom}(a, b, c, d, v, f) \wedge \\ \quad \operatorname{is\text{-}local\text{-}hom}(c, d, g, m, n, h) \Rightarrow \\ \quad \quad \exists p, k. \\ \quad \quad \quad \operatorname{is\text{-}local\text{-}hom}(a, b, g, m, p, k) \wedge \forall x \in p.\; (k(x) = h(f(x))) \end{array}}}
\vnbprop{local-homomorphism-composite}{\vnbstmt{\begin{array}{@{}l@{}} \forall a, b, c, v, f, d, h. \\ \quad \operatorname{is\text{-}top\text{-}group}(a) \wedge \\ \quad \operatorname{is\text{-}top\text{-}group}(b) \wedge \\ \quad \operatorname{is\text{-}top\text{-}group}(c) \wedge \\ \quad \operatorname{is\text{-}local\text{-}hom}(\operatorname{top\text{-}group\text{-}space}(a), \operatorname{top\text{-}group\text{-}group}(a), \operatorname{top\text{-}group\text{-}space}(b), \operatorname{top\text{-}group\text{-}group}(b), v, f) \wedge \\ \quad \operatorname{is\text{-}local\text{-}hom}(\operatorname{top\text{-}group\text{-}space}(b), \operatorname{top\text{-}group\text{-}group}(b), \operatorname{top\text{-}group\text{-}space}(c), \operatorname{top\text{-}group\text{-}group}(c), d, h) \Rightarrow \\ \quad \quad \exists g, k. \\ \quad \quad \quad \operatorname{is\text{-}local\text{-}hom}(\operatorname{top\text{-}group\text{-}space}(a), \operatorname{top\text{-}group\text{-}group}(a), \operatorname{top\text{-}group\text{-}space}(c), \operatorname{top\text{-}group\text{-}group}(c), g, k) \wedge \\ \quad \quad \quad \forall x \in g.\; (k(x) = h(f(x))) \end{array}}}

\vnbbookfile{theorem-library/ccr4c-subnf}{theorem-library/ccr4c-subnf}
\noindent{\small ccr4c-subnf.scm -- wave of 2026-11-09/10 (agent CCR-1109-C4), wired 2026-11-10 from scratchpad/ccr1109c4/ccr4c-subnf.scm. Library: analysis.}\par
\vnbprop{sub-nf-add-val}{\vnbstmt{\begin{array}{@{}l@{}} \forall k, s.; \forall u, v \in s. \\ \quad (\operatorname{is\text{-}normed\text{-}field}(k) \wedge \operatorname{is\text{-}nf\text{-}subfield}(k, s)) \Rightarrow  \\ \quad \quad (\lambda {[}\operatorname{snu\_}, \operatorname{snv\_}] \in (s \times s).\; \operatorname{add}(k)(\operatorname{snu\_}, \operatorname{snv\_}))(u, v) = \operatorname{add}(k)(u, v) \end{array}}}
\vnbprop{sub-nf-mul-val}{\vnbstmt{\begin{array}{@{}l@{}} \forall k, s.; \forall u, v \in s. \\ \quad (\operatorname{is\text{-}normed\text{-}field}(k) \wedge \operatorname{is\text{-}nf\text{-}subfield}(k, s)) \Rightarrow  \\ \quad \quad (\lambda {[}\operatorname{snu\_}, \operatorname{snv\_}] \in (s \times s).\; \operatorname{mul}(k)(\operatorname{snu\_}, \operatorname{snv\_}))(u, v) = \operatorname{mul}(k)(u, v) \end{array}}}
\vnbprop{sub-nf-neg-val}{\vnbstmt{\begin{array}{@{}l@{}} \forall k, s.; \forall u \in s. \\ \quad (\operatorname{is\text{-}normed\text{-}field}(k) \wedge \operatorname{is\text{-}nf\text{-}subfield}(k, s)) \Rightarrow  \\ \quad \quad (\lambda \mathit{snu} \in s.\; \operatorname{neg}(k)(\mathit{snu}))(u) = \operatorname{neg}(k)(u) \end{array}}}
\vnbprop{sub-nf-fnrm-val}{\vnbstmt{\begin{array}{@{}l@{}} \forall k, s.; \forall u \in s. \\ \quad (\operatorname{is\text{-}normed\text{-}field}(k) \wedge \operatorname{is\text{-}nf\text{-}subfield}(k, s)) \Rightarrow  \\ \quad \quad (\lambda \mathit{snu} \in s.\; \operatorname{fnrm}(k)(\mathit{snu}))(u) = \operatorname{fnrm}(k)(u) \end{array}}}
\vnbprop{sub-nf-add-fun}{\vnbstmt{\begin{array}{@{}l@{}} \forall k, s. \\ \quad (\operatorname{is\text{-}normed\text{-}field}(k) \wedge \operatorname{is\text{-}nf\text{-}subfield}(k, s)) \Rightarrow  \\ \quad \quad (\lambda {[}\operatorname{snu\_}, \operatorname{snv\_}] \in (s \times s).\; \operatorname{add}(k)(\operatorname{snu\_}, \operatorname{snv\_})) \in ((s \times s) \to s) \end{array}}}
\vnbprop{sub-nf-mul-fun}{\vnbstmt{\begin{array}{@{}l@{}} \forall k, s. \\ \quad (\operatorname{is\text{-}normed\text{-}field}(k) \wedge \operatorname{is\text{-}nf\text{-}subfield}(k, s)) \Rightarrow  \\ \quad \quad (\lambda {[}\operatorname{snu\_}, \operatorname{snv\_}] \in (s \times s).\; \operatorname{mul}(k)(\operatorname{snu\_}, \operatorname{snv\_})) \in ((s \times s) \to s) \end{array}}}
\vnbprop{sub-nf-neg-fun}{\vnbstmt{\begin{array}{@{}l@{}} \forall k, s. \\ \quad (\operatorname{is\text{-}normed\text{-}field}(k) \wedge \operatorname{is\text{-}nf\text{-}subfield}(k, s)) \Rightarrow  \\ \quad \quad (\lambda u \in s.\; \operatorname{neg}(k)(u)) \in (s \to s) \end{array}}}
\vnbprop{sub-nf-fnrm-fun}{\vnbstmt{\begin{array}{@{}l@{}} \forall k, s. \\ \quad (\operatorname{is\text{-}normed\text{-}field}(k) \wedge \operatorname{is\text{-}nf\text{-}subfield}(k, s)) \Rightarrow  \\ \quad \quad (\lambda u \in s.\; \operatorname{fnrm}(k)(u)) \in (s \to \mathbb{R}) \end{array}}}
\vnbprop{sub-nf-carr}{\vnbstmt{\begin{array}{@{}l@{}} \forall k, s. \\ \quad (\operatorname{is\text{-}normed\text{-}field}(k) \wedge \operatorname{is\text{-}nf\text{-}subfield}(k, s)) \Rightarrow (\operatorname{carr}(\operatorname{sub\text{-}nf}(k, s)) \simeq s) \end{array}}}
\vnbprop{sub-nf-add}{\vnbstmt{\begin{array}{@{}l@{}} \forall k, s. \\ \quad (\operatorname{is\text{-}normed\text{-}field}(k) \wedge \operatorname{is\text{-}nf\text{-}subfield}(k, s)) \Rightarrow  \\ \quad \quad (\operatorname{add}(\operatorname{sub\text{-}nf}(k, s)) \simeq (\lambda {[}\operatorname{snu\_}, \operatorname{snv\_}] \in (s \times s).\; \operatorname{add}(k)(\operatorname{snu\_}, \operatorname{snv\_}))) \end{array}}}
\vnbprop{sub-nf-mul}{\vnbstmt{\begin{array}{@{}l@{}} \forall k, s. \\ \quad (\operatorname{is\text{-}normed\text{-}field}(k) \wedge \operatorname{is\text{-}nf\text{-}subfield}(k, s)) \Rightarrow  \\ \quad \quad (\operatorname{mul}(\operatorname{sub\text{-}nf}(k, s)) \simeq (\lambda {[}\operatorname{snu\_}, \operatorname{snv\_}] \in (s \times s).\; \operatorname{mul}(k)(\operatorname{snu\_}, \operatorname{snv\_}))) \end{array}}}
\vnbprop{sub-nf-neg}{\vnbstmt{\begin{array}{@{}l@{}} \forall k, s. \\ \quad (\operatorname{is\text{-}normed\text{-}field}(k) \wedge \operatorname{is\text{-}nf\text{-}subfield}(k, s)) \Rightarrow  \\ \quad \quad (\operatorname{neg}(\operatorname{sub\text{-}nf}(k, s)) \simeq (\lambda u \in s.\; \operatorname{neg}(k)(u))) \end{array}}}
\vnbprop{sub-nf-zero}{\vnbstmt{\begin{array}{@{}l@{}} \forall k, s. \\ \quad (\operatorname{is\text{-}normed\text{-}field}(k) \wedge \operatorname{is\text{-}nf\text{-}subfield}(k, s)) \Rightarrow  \\ \quad \quad (\operatorname{zero}(\operatorname{sub\text{-}nf}(k, s)) \simeq \operatorname{zero}(k)) \end{array}}}
\vnbprop{sub-nf-one}{\vnbstmt{\begin{array}{@{}l@{}} \forall k, s. \\ \quad (\operatorname{is\text{-}normed\text{-}field}(k) \wedge \operatorname{is\text{-}nf\text{-}subfield}(k, s)) \Rightarrow (\operatorname{one}(\operatorname{sub\text{-}nf}(k, s)) \simeq \operatorname{one}(k)) \end{array}}}
\vnbprop{sub-nf-fnrm}{\vnbstmt{\begin{array}{@{}l@{}} \forall k, s. \\ \quad (\operatorname{is\text{-}normed\text{-}field}(k) \wedge \operatorname{is\text{-}nf\text{-}subfield}(k, s)) \Rightarrow  \\ \quad \quad (\operatorname{fnrm}(\operatorname{sub\text{-}nf}(k, s)) \simeq (\lambda u \in s.\; \operatorname{fnrm}(k)(u))) \end{array}}}
\vnbprop{sub-nf-is-normed-field}{\vnbstmt{\begin{array}{@{}l@{}} \forall k, s. \\ \quad (\operatorname{is\text{-}normed\text{-}field}(k) \wedge \operatorname{is\text{-}nf\text{-}subfield}(k, s)) \Rightarrow  \\ \quad \quad \operatorname{is\text{-}normed\text{-}field}(\operatorname{sub\text{-}nf}(k, s)) \end{array}}}
\vnbprop{sub-nf-is-nv-field}{\vnbstmt{\begin{array}{@{}l@{}} \forall k, s. \\ \quad (\operatorname{is\text{-}nv\text{-}field}(k) \wedge \operatorname{is\text{-}nf\text{-}subfield}(k, s)) \Rightarrow \operatorname{is\text{-}nv\text{-}field}(\operatorname{sub\text{-}nf}(k, s)) \end{array}}}
\vnbprop{nf-metric-space-unfold}{\vnbstmt{\begin{array}{@{}l@{}} \forall f. \\ \quad (\operatorname{nf\text{-}metric\text{-}space}(f) \simeq {[}\operatorname{carr}(f), (\lambda {[}\operatorname{snx\_}, \operatorname{sny\_}] \in (\operatorname{carr}(f) \times \operatorname{carr}(f)).\; \operatorname{fnrm}(f)(\operatorname{add}(f)(\operatorname{snx\_}, \operatorname{neg}(f)(\operatorname{sny\_}))))]) \end{array}}}
\vnbprop{sub-nf-metric}{\vnbstmt{\begin{array}{@{}l@{}} \forall k, s. \\ \quad (\operatorname{is\text{-}normed\text{-}field}(k) \wedge \operatorname{is\text{-}nf\text{-}subfield}(k, s)) \Rightarrow  \\ \quad \quad (\operatorname{nf\text{-}metric\text{-}space}(\operatorname{sub\text{-}nf}(k, s)) \simeq \operatorname{subspace\text{-}ms}(\operatorname{nf\text{-}metric\text{-}space}(k), s)) \end{array}}}
\vnbprop{sub-nf-integers}{\vnbstmt{\begin{array}{@{}l@{}} \forall k, s. \\ \quad (\operatorname{is\text{-}normed\text{-}field}(k) \wedge \operatorname{is\text{-}nf\text{-}subfield}(k, s)) \Rightarrow  \\ \quad \quad \operatorname{nv\text{-}integers}(\operatorname{sub\text{-}nf}(k, s)) = (\operatorname{nv\text{-}integers}(k) \cap s) \end{array}}}
\vnbprop{nf-subfield-ring-power}{\vnbstmt{\begin{array}{@{}l@{}} \forall n \in \mathbb{N}, k, s.; \forall x \in s. \\ \quad (\operatorname{is\text{-}normed\text{-}field}(k) \wedge \operatorname{is\text{-}nf\text{-}subfield}(k, s)) \Rightarrow  \\ \quad \quad \operatorname{ring\text{-}power}(\operatorname{normed\text{-}field\text{-}as\text{-}commutative\text{-}ring}(k), x, n) \in s \end{array}}}
\vnbprop{nf-subfield-small}{\vnbstmt{\begin{array}{@{}l@{}} \forall k, s.; \forall w \in s, a. \\ \quad (\operatorname{is\text{-}normed\text{-}field}(k) \wedge (\operatorname{is\text{-}nf\text{-}subfield}(k, s) \wedge ((0 < \operatorname{fnrm}(k)(w)) \wedge ((\operatorname{fnrm}(k)(w) < 1) \wedge \operatorname{pos\text{-}rr}(a))))) \Rightarrow  \\ \quad \quad \exists b \in s.\; (\neg (b = \operatorname{zero}(k)) \wedge (\operatorname{fnrm}(k)(b) < a)) \end{array}}}
\vnbprop{nf-subfield-zz-act-one}{\vnbstmt{\begin{array}{@{}l@{}} \forall n \in \mathbb{N}, k, s. \\ \quad (\operatorname{is\text{-}normed\text{-}field}(k) \wedge \operatorname{is\text{-}nf\text{-}subfield}(k, s)) \Rightarrow  \\ \quad \quad \operatorname{zz\text{-}act}(\operatorname{ring\text{-}additive\text{-}ag}(\operatorname{normed\text{-}field\text{-}as\text{-}commutative\text{-}ring}(k)), \\ \quad \quad \quad n, \operatorname{one}(\operatorname{normed\text{-}field\text{-}as\text{-}commutative\text{-}ring}(k))) \\ \quad \quad \in\; s \end{array}}}

\vnbbookfile{theorem-library/ccr4d-ext}{theorem-library/ccr4d-ext}
\noindent{\small ccr4d-ext.scm -- wave of 2026-11-09/10 (agent CCR-1109-C5), wired 2026-11-10 from scratchpad/ccr1109c5/ccr4d-ext.scm. Library: analysis.}\par
\vnbprop{nf-view-field-ring}{\vnbstmt{\begin{array}{@{}l@{}} \forall k. \\ \quad \operatorname{is\text{-}normed\text{-}field}(k) \Rightarrow \operatorname{is\text{-}field\text{-}ring}(\operatorname{normed\text{-}field\text{-}as\text{-}commutative\text{-}ring}(k)) \end{array}}}
\vnbprop{nf-subfield-module-act-fun}{\vnbstmt{\begin{array}{@{}l@{}} \forall k, s. \\ \quad (\operatorname{is\text{-}normed\text{-}field}(k) \wedge \operatorname{is\text{-}nf\text{-}subfield}(k, s)) \Rightarrow  \\ \quad \quad (\lambda {[}\operatorname{fmr\_}, \operatorname{fmx\_}] \in (s \times \operatorname{carr}(\operatorname{normed\text{-}field\text{-}as\text{-}commutative\text{-}ring}(k))).\; \\ \quad \quad \quad \operatorname{mul}(\operatorname{normed\text{-}field\text{-}as\text{-}commutative\text{-}ring}(k))(\operatorname{fmr\_}, \operatorname{fmx\_})) \\ \quad \quad \in\; ((s \times \operatorname{carr}(\operatorname{normed\text{-}field\text{-}as\text{-}commutative\text{-}ring}(k))) \to \operatorname{carr}(\operatorname{normed\text{-}field\text{-}as\text{-}commutative\text{-}ring}(k))) \end{array}}}
\vnbprop{nf-subfield-module-scal}{\vnbstmt{\begin{array}{@{}l@{}} \forall k, s. \\ \quad (\operatorname{is\text{-}normed\text{-}field}(k) \wedge \operatorname{is\text{-}nf\text{-}subfield}(k, s)) \Rightarrow  \\ \quad \quad (\operatorname{scal}(\operatorname{nf\text{-}subfield\text{-}module}(k, s)) \simeq \operatorname{normed\text{-}field\text{-}as\text{-}commutative\text{-}ring}(\operatorname{sub\text{-}nf}(k, s))) \end{array}}}
\vnbprop{nf-subfield-module-vec}{\vnbstmt{\begin{array}{@{}l@{}} \forall k, s. \\ \quad (\operatorname{is\text{-}normed\text{-}field}(k) \wedge \operatorname{is\text{-}nf\text{-}subfield}(k, s)) \Rightarrow  \\ \quad \quad (\operatorname{vec}(\operatorname{nf\text{-}subfield\text{-}module}(k, s)) \simeq \operatorname{carr}(\operatorname{normed\text{-}field\text{-}as\text{-}commutative\text{-}ring}(k))) \end{array}}}
\vnbprop{nf-subfield-module-vadd}{\vnbstmt{\begin{array}{@{}l@{}} \forall k, s. \\ \quad (\operatorname{is\text{-}normed\text{-}field}(k) \wedge \operatorname{is\text{-}nf\text{-}subfield}(k, s)) \Rightarrow  \\ \quad \quad (\operatorname{vadd}(\operatorname{nf\text{-}subfield\text{-}module}(k, s)) \simeq \operatorname{add}(\operatorname{normed\text{-}field\text{-}as\text{-}commutative\text{-}ring}(k))) \end{array}}}
\vnbprop{nf-subfield-module-vzero}{\vnbstmt{\begin{array}{@{}l@{}} \forall k, s. \\ \quad (\operatorname{is\text{-}normed\text{-}field}(k) \wedge \operatorname{is\text{-}nf\text{-}subfield}(k, s)) \Rightarrow  \\ \quad \quad (\operatorname{vzero}(\operatorname{nf\text{-}subfield\text{-}module}(k, s)) \simeq \operatorname{zero}(\operatorname{normed\text{-}field\text{-}as\text{-}commutative\text{-}ring}(k))) \end{array}}}
\vnbprop{nf-subfield-module-vneg}{\vnbstmt{\begin{array}{@{}l@{}} \forall k, s. \\ \quad (\operatorname{is\text{-}normed\text{-}field}(k) \wedge \operatorname{is\text{-}nf\text{-}subfield}(k, s)) \Rightarrow  \\ \quad \quad (\operatorname{vneg}(\operatorname{nf\text{-}subfield\text{-}module}(k, s)) \simeq \operatorname{neg}(\operatorname{normed\text{-}field\text{-}as\text{-}commutative\text{-}ring}(k))) \end{array}}}
\vnbprop{nf-subfield-module-act}{\vnbstmt{\begin{array}{@{}l@{}} \forall k, s. \\ \quad (\operatorname{is\text{-}normed\text{-}field}(k) \wedge \operatorname{is\text{-}nf\text{-}subfield}(k, s)) \Rightarrow  \\ \quad \quad (\operatorname{act}(\operatorname{nf\text{-}subfield\text{-}module}(k, s)) \simeq (\lambda {[}\operatorname{fmr\_}, \operatorname{fmx\_}] \in (s \times \operatorname{carr}(\operatorname{normed\text{-}field\text{-}as\text{-}commutative\text{-}ring}(k))).\; \operatorname{mul}(\operatorname{normed\text{-}field\text{-}as\text{-}commutative\text{-}ring}(k))(\operatorname{fmr\_}, \operatorname{fmx\_}))) \end{array}}}
\vnbprop{nf-subfield-module-is-module}{\vnbstmt{\begin{array}{@{}l@{}} \forall k, s. \\ \quad (\operatorname{is\text{-}normed\text{-}field}(k) \wedge \operatorname{is\text{-}nf\text{-}subfield}(k, s)) \Rightarrow  \\ \quad \quad \operatorname{is\text{-}module}(\operatorname{nf\text{-}subfield\text{-}module}(k, s)) \end{array}}}
\vnbprop{nf-subfield-fin-ext}{\vnbstmt{\begin{array}{@{}l@{}} \forall k, s.; \forall n \in \mathbb{N}, u \in \operatorname{mat}(n, 1, \operatorname{carr}(k)). \\ \quad \operatorname{is\text{-}normed\text{-}field}(k) \wedge \\ \quad \operatorname{is\text{-}nf\text{-}subfield}(k, s) \wedge \\ \quad 1 \leq n \wedge \\ \quad \operatorname{generates}(\operatorname{nf\text{-}subfield\text{-}module}(k, s), n, u) \wedge \\ \quad \operatorname{rel\text{-}free}(\operatorname{nf\text{-}subfield\text{-}module}(k, s), n, u) \Rightarrow \\ \quad \quad \operatorname{is\text{-}fin\text{-}ext}(\operatorname{sub\text{-}nf}(k, s), \operatorname{normed\text{-}field\text{-}as\text{-}commutative\text{-}ring}(k), \operatorname{nf\text{-}subfield\text{-}module}(k, s), n, u) \end{array}}}

\vnbbookfile{theorem-library/haar-c2}{theorem-library/haar-c2}
\noindent{\small haar-c2.scm -- wave of 2026-11-09/10 (agent HAAR-C2), wired 2026-11-10 from scratchpad/haarc2-1109/haar-c2.scm. Library: analysis.}\par
\vnbprop{ck-cont-fns-is-set}{\vnbstmt{\forall s.\; (\operatorname{is\text{-}metric\text{-}space}(s) \Rightarrow (\operatorname{cont\text{-}fns}(s) \in \mathbf{Set}))}}
\vnbprop{ck-zero-cont}{\vnbstmt{\forall s.\; (\operatorname{is\text{-}metric\text{-}space}(s) \Rightarrow ((\lambda x \in \operatorname{pts}(s).\; 0) \in \operatorname{cont\text{-}fns}(s)))}}
\vnbprop{ck-neg-cont}{\vnbstmt{\forall s.; \forall f \in \operatorname{cont\text{-}fns}(s).\; ((\lambda x \in \operatorname{pts}(s).\; -f(x)) \in \operatorname{cont\text{-}fns}(s))}}
\vnbprop{ck-norm-upper}{\vnbstmt{\forall s.; \forall h \in (\operatorname{pts}(s) \to \mathbb{R}), x \in \operatorname{pts}(s).\; (\lvert h(x) \rvert \leq \lVert h \rVert_{\infty})}}
\vnbprop{ck-norm-bound}{\vnbstmt{\begin{array}{@{}l@{}} \forall s, h, b. \\ \quad ((h \in (\operatorname{pts}(s) \to \mathbb{R})) \wedge ((b \in \mathbb{R}) \wedge ((0 \leq b) \wedge \forall x \in \operatorname{pts}(s).\; (\lvert h(x) \rvert \leq b)))) \Rightarrow  \\ \quad \quad (\lVert h \rVert_{\infty} \in \mathbb{R}) \wedge ((0 \leq \lVert h \rVert_{\infty}) \wedge (\lVert h \rVert_{\infty} \leq b)) \end{array}}}
\vnbprop{ck-cont-bounded}{\vnbstmt{\begin{array}{@{}l@{}} \forall s, h. \\ \quad (\operatorname{is\text{-}compact}(s) \wedge (h \in \operatorname{cont\text{-}fns}(s))) \Rightarrow  \\ \quad \quad \exists m \in \mathbb{R}.\; ((0 \leq m) \wedge \forall z \in \operatorname{pts}(s).\; (\lvert h(z) \rvert \leq m)) \end{array}}}
\vnbprop{ck-norm-real}{\vnbstmt{\begin{array}{@{}l@{}} \forall s, h. \\ \quad (\operatorname{is\text{-}compact}(s) \wedge (h \in \operatorname{cont\text{-}fns}(s))) \Rightarrow  \\ \quad \quad (\lVert h \rVert_{\infty} \in \mathbb{R}) \wedge (0 \leq \lVert h \rVert_{\infty}) \end{array}}}
\vnbprop{ck-nvs-vec}{\vnbstmt{\forall s.\; (\operatorname{is\text{-}metric\text{-}space}(s) \Rightarrow (\operatorname{vec}(C({s})) \simeq \operatorname{cont\text{-}fns}(s)))}}
\vnbprop{ck-nvs-scal}{\vnbstmt{\begin{array}{@{}l@{}} \forall s. \\ \quad \operatorname{is\text{-}metric\text{-}space}(s) \Rightarrow  \\ \quad \quad (\operatorname{scal}(C({s})) \simeq \operatorname{normed\text{-}field\text{-}as\text{-}commutative\text{-}ring}(\operatorname{rr\text{-}normed\text{-}field})) \end{array}}}
\vnbprop{ck-vadd-readoff}{\vnbstmt{\begin{array}{@{}l@{}} \forall s, f, g. \\ \quad (\operatorname{is\text{-}metric\text{-}space}(s) \wedge ((f \in \operatorname{cont\text{-}fns}(s)) \wedge (g \in \operatorname{cont\text{-}fns}(s)))) \Rightarrow  \\ \quad \quad (\operatorname{vadd}(C({s}))(f, g) \simeq (\lambda x \in \operatorname{pts}(s).\; (f(x) + g(x)))) \end{array}}}
\vnbprop{ck-act-readoff}{\vnbstmt{\begin{array}{@{}l@{}} \forall s, r, f. \\ \quad (\operatorname{is\text{-}metric\text{-}space}(s) \wedge ((r \in \mathbb{R}) \wedge (f \in \operatorname{cont\text{-}fns}(s)))) \Rightarrow  \\ \quad \quad (\operatorname{act}(C({s}))(r, f) \simeq (\lambda x \in \operatorname{pts}(s).\; (r \cdot f(x)))) \end{array}}}
\vnbprop{ck-vneg-readoff}{\vnbstmt{\begin{array}{@{}l@{}} \forall s, f. \\ \quad (\operatorname{is\text{-}metric\text{-}space}(s) \wedge (f \in \operatorname{cont\text{-}fns}(s))) \Rightarrow  \\ \quad \quad (\operatorname{vneg}(C({s}))(f) \simeq (\lambda x \in \operatorname{pts}(s).\; -f(x))) \end{array}}}
\vnbprop{ck-vzero-readoff}{\vnbstmt{\begin{array}{@{}l@{}} \forall s. \\ \quad \operatorname{is\text{-}metric\text{-}space}(s) \Rightarrow (\operatorname{vzero}(C({s})) \simeq (\lambda x \in \operatorname{pts}(s).\; 0)) \end{array}}}
\vnbprop{ck-vnrm-readoff}{\vnbstmt{\begin{array}{@{}l@{}} \forall s, f. \\ \quad (\operatorname{is\text{-}metric\text{-}space}(s) \wedge (f \in \operatorname{cont\text{-}fns}(s))) \Rightarrow  \\ \quad \quad (\operatorname{vnrm}(C({s}))(f) \simeq \lVert f \rVert_{\infty}) \end{array}}}
\vnbprop{ck-vadd-in}{\vnbstmt{\begin{array}{@{}l@{}} \forall s, f, g. \\ \quad (\operatorname{is\text{-}metric\text{-}space}(s) \wedge ((f \in \operatorname{cont\text{-}fns}(s)) \wedge (g \in \operatorname{cont\text{-}fns}(s)))) \Rightarrow  \\ \quad \quad \operatorname{vadd}(C({s}))(f, g) \in \operatorname{cont\text{-}fns}(s) \end{array}}}
\vnbprop{ck-act-in}{\vnbstmt{\begin{array}{@{}l@{}} \forall s, r, f. \\ \quad (\operatorname{is\text{-}metric\text{-}space}(s) \wedge ((r \in \mathbb{R}) \wedge (f \in \operatorname{cont\text{-}fns}(s)))) \Rightarrow  \\ \quad \quad \operatorname{act}(C({s}))(r, f) \in \operatorname{cont\text{-}fns}(s) \end{array}}}
\vnbprop{ck-vneg-in}{\vnbstmt{\begin{array}{@{}l@{}} \forall s, f. \\ \quad (\operatorname{is\text{-}metric\text{-}space}(s) \wedge (f \in \operatorname{cont\text{-}fns}(s))) \Rightarrow  \\ \quad \quad \operatorname{vneg}(C({s}))(f) \in \operatorname{cont\text{-}fns}(s) \end{array}}}
\vnbprop{ck-vzero-in}{\vnbstmt{\forall s.\; (\operatorname{is\text{-}metric\text{-}space}(s) \Rightarrow (\operatorname{vzero}(C({s})) \in \operatorname{cont\text{-}fns}(s)))}}
\vnbprop{ck-vadd-apply}{\vnbstmt{\begin{array}{@{}l@{}} \forall s, f, g, x. \\ \quad (\operatorname{is\text{-}metric\text{-}space}(s) \wedge ((f \in \operatorname{cont\text{-}fns}(s)) \wedge ((g \in \operatorname{cont\text{-}fns}(s)) \wedge (x \in \operatorname{pts}(s))))) \Rightarrow  \\ \quad \quad \operatorname{vadd}(C({s}))(f, g)(x) = (f(x) + g(x)) \end{array}}}
\vnbprop{ck-act-apply}{\vnbstmt{\begin{array}{@{}l@{}} \forall s, r, f, x. \\ \quad (\operatorname{is\text{-}metric\text{-}space}(s) \wedge ((r \in \mathbb{R}) \wedge ((f \in \operatorname{cont\text{-}fns}(s)) \wedge (x \in \operatorname{pts}(s))))) \Rightarrow  \\ \quad \quad \operatorname{act}(C({s}))(r, f)(x) = (r \cdot f(x)) \end{array}}}
\vnbprop{ck-vneg-apply}{\vnbstmt{\begin{array}{@{}l@{}} \forall s, f, x. \\ \quad (\operatorname{is\text{-}metric\text{-}space}(s) \wedge ((f \in \operatorname{cont\text{-}fns}(s)) \wedge (x \in \operatorname{pts}(s)))) \Rightarrow  \\ \quad \quad \operatorname{vneg}(C({s}))(f)(x) = -f(x) \end{array}}}
\vnbprop{ck-vzero-apply}{\vnbstmt{\begin{array}{@{}l@{}} \forall s, x. \\ \quad (\operatorname{is\text{-}metric\text{-}space}(s) \wedge (x \in \operatorname{pts}(s))) \Rightarrow (\operatorname{vzero}(C({s}))(x) = 0) \end{array}}}
\vnbprop{ck-compact-is-ms}{\vnbstmt{\forall s.\; (\operatorname{is\text{-}compact}(s) \Rightarrow \operatorname{is\text{-}metric\text{-}space}(s))}}
\vnbprop{ck-norm-add-le}{\vnbstmt{\begin{array}{@{}l@{}} \forall s, f, g. \\ \quad (\operatorname{is\text{-}compact}(s) \wedge ((f \in \operatorname{cont\text{-}fns}(s)) \wedge (g \in \operatorname{cont\text{-}fns}(s)))) \Rightarrow  \\ \quad \quad \lVert (\lambda x \in \operatorname{pts}(s).\; (f(x) + g(x))) \rVert_{\infty} \leq (\lVert f \rVert_{\infty} + \lVert g \rVert_{\infty}) \end{array}}}
\vnbprop{ck-norm-act-le}{\vnbstmt{\begin{array}{@{}l@{}} \forall s, r, f. \\ \quad (\operatorname{is\text{-}compact}(s) \wedge ((r \in \mathbb{R}) \wedge (f \in \operatorname{cont\text{-}fns}(s)))) \Rightarrow  \\ \quad \quad \lVert (\lambda x \in \operatorname{pts}(s).\; (r \cdot f(x))) \rVert_{\infty} \leq (\lvert r \rvert \cdot \lVert f \rVert_{\infty}) \end{array}}}
\vnbprop{ck-norm-act-ge}{\vnbstmt{\begin{array}{@{}l@{}} \forall s, r, f. \\ \quad (\operatorname{is\text{-}compact}(s) \wedge ((r \in \mathbb{R}) \wedge (f \in \operatorname{cont\text{-}fns}(s)))) \Rightarrow  \\ \quad \quad (\lvert r \rvert \cdot \lVert f \rVert_{\infty}) \leq \lVert (\lambda x \in \operatorname{pts}(s).\; (r \cdot f(x))) \rVert_{\infty} \end{array}}}
\vnbprop{ck-norm-zero-elim}{\vnbstmt{\begin{array}{@{}l@{}} \forall s, f. \\ \quad (\operatorname{is\text{-}metric\text{-}space}(s) \wedge ((f \in \operatorname{cont\text{-}fns}(s)) \wedge (\lVert f \rVert_{\infty} = 0))) \Rightarrow  \\ \quad \quad f = (\lambda x \in \operatorname{pts}(s).\; 0) \end{array}}}
\vnbprop{ck-norm-zero-intro}{\vnbstmt{\forall s.\; (\operatorname{is\text{-}metric\text{-}space}(s) \Rightarrow (\lVert (\lambda x \in \operatorname{pts}(s).\; 0) \rVert_{\infty} = 0))}}
\vnbprop{ck-vnrm-nonneg}{\vnbstmt{\begin{array}{@{}l@{}} \forall s, f. \\ \quad (\operatorname{is\text{-}compact}(s) \wedge (f \in \operatorname{cont\text{-}fns}(s))) \Rightarrow (0 \leq \operatorname{vnrm}(C({s}))(f)) \end{array}}}
\vnbprop{ck-vnrm-triangle}{\vnbstmt{\begin{array}{@{}l@{}} \forall s, f, g. \\ \quad (\operatorname{is\text{-}compact}(s) \wedge ((f \in \operatorname{cont\text{-}fns}(s)) \wedge (g \in \operatorname{cont\text{-}fns}(s)))) \Rightarrow  \\ \quad \quad \operatorname{vnrm}(C({s}))(\operatorname{vadd}(C({s}))(f, g)) \leq (\operatorname{vnrm}(C({s}))(f) + \operatorname{vnrm}(C({s}))(g)) \end{array}}}
\vnbprop{ck-vnrm-homogeneous}{\vnbstmt{\begin{array}{@{}l@{}} \forall s, r, f. \\ \quad (\operatorname{is\text{-}compact}(s) \wedge ((r \in \mathbb{R}) \wedge (f \in \operatorname{cont\text{-}fns}(s)))) \Rightarrow  \\ \quad \quad \operatorname{vnrm}(C({s}))(\operatorname{act}(C({s}))(r, f)) = (\lvert r \rvert \cdot \operatorname{vnrm}(C({s}))(f)) \end{array}}}
\vnbprop{ck-vnrm-definite}{\vnbstmt{\begin{array}{@{}l@{}} \forall s, f. \\ \quad (\operatorname{is\text{-}compact}(s) \wedge (f \in \operatorname{cont\text{-}fns}(s))) \Rightarrow  \\ \quad \quad (\operatorname{vnrm}(C({s}))(f) = 0) \Leftrightarrow (f = \operatorname{vzero}(C({s}))) \end{array}}}
\vnbprop{ck-vadd-typing}{\vnbstmt{\begin{array}{@{}l@{}} \forall s. \\ \quad \operatorname{is\text{-}metric\text{-}space}(s) \Rightarrow  \\ \quad \quad \operatorname{vadd}(C({s})) \in ((\operatorname{cont\text{-}fns}(s) \times \operatorname{cont\text{-}fns}(s)) \to \operatorname{cont\text{-}fns}(s)) \end{array}}}
\vnbprop{ck-act-typing}{\vnbstmt{\begin{array}{@{}l@{}} \forall s. \\ \quad \operatorname{is\text{-}metric\text{-}space}(s) \Rightarrow  \\ \quad \quad \operatorname{act}(C({s})) \in ((\mathbb{R} \times \operatorname{cont\text{-}fns}(s)) \to \operatorname{cont\text{-}fns}(s)) \end{array}}}
\vnbprop{ck-vneg-typing}{\vnbstmt{\begin{array}{@{}l@{}} \forall s. \\ \quad \operatorname{is\text{-}metric\text{-}space}(s) \Rightarrow (\operatorname{vneg}(C({s})) \in (\operatorname{cont\text{-}fns}(s) \to \operatorname{cont\text{-}fns}(s))) \end{array}}}
\vnbprop{ck-vnrm-typing}{\vnbstmt{\forall s.\; (\operatorname{is\text{-}compact}(s) \Rightarrow (\operatorname{vnrm}(C({s})) \in (\operatorname{cont\text{-}fns}(s) \to \mathbb{R})))}}
\vnbprop{ck-nvs-is-normed-vector-space}{\vnbstmt{\forall s.\; (\operatorname{is\text{-}compact}(s) \Rightarrow \operatorname{is\text{-}normed\text{-}vector\text{-}space}(C({s})))}}
\vnbprop{ck-dist-upper}{\vnbstmt{\begin{array}{@{}l@{}} \forall s, f, g, x. \\ \quad (\operatorname{is\text{-}compact}(s) \wedge ((f \in \operatorname{cont\text{-}fns}(s)) \wedge ((g \in \operatorname{cont\text{-}fns}(s)) \wedge (x \in \operatorname{pts}(s))))) \Rightarrow  \\ \quad \quad \lvert (f(x) - g(x)) \rvert \leq \operatorname{dist}(\operatorname{nvs\text{-}metric\text{-}space}(C({s})))(f, g) \end{array}}}
\vnbprop{ck-dist-le}{\vnbstmt{\begin{array}{@{}l@{}} \forall s, f, g, b. \\ \quad \operatorname{is\text{-}compact}(s) \wedge \\ \quad f \in \operatorname{cont\text{-}fns}(s) \wedge \\ \quad g \in \operatorname{cont\text{-}fns}(s) \wedge \\ \quad b \in \mathbb{R} \wedge \\ \quad 0 \leq b \wedge \\ \quad \forall x \in \operatorname{pts}(s).\; (\lvert (f(x) - g(x)) \rvert \leq b) \Rightarrow \\ \quad \quad \operatorname{dist}(\operatorname{nvs\text{-}metric\text{-}space}(C({s})))(f, g) \leq b \end{array}}}
\vnbprop{ck-cont-fns-subset-fun}{\vnbstmt{\forall s.\; (\operatorname{is\text{-}metric\text{-}space}(s) \Rightarrow (\operatorname{cont\text{-}fns}(s) \subseteq (\operatorname{pts}(s) \to \mathbb{R})))}}
\vnbprop{ck-nvs-complete}{\vnbstmt{\forall s.\; (\operatorname{is\text{-}compact}(s) \Rightarrow \operatorname{is\text{-}complete}(\operatorname{nvs\text{-}metric\text{-}space}(C({s}))))}}
\vnbprop{mg-ms-pts}{\vnbstmt{\forall g.\; (\operatorname{is\text{-}metric\text{-}group}(g) \Rightarrow (\operatorname{pts}(\operatorname{metric\text{-}group\text{-}ms}(g)) \simeq \operatorname{carr}(g)))}}
\vnbprop{mg-ms-dist}{\vnbstmt{\forall g.\; (\operatorname{is\text{-}metric\text{-}group}(g) \Rightarrow (\operatorname{dist}(\operatorname{metric\text{-}group\text{-}ms}(g)) \simeq \operatorname{gdist}(g)))}}
\vnbprop{mg-group-carr}{\vnbstmt{\forall g.\; (\operatorname{is\text{-}metric\text{-}group}(g) \Rightarrow (\operatorname{carr}(\operatorname{metric\text{-}group\text{-}group}(g)) \simeq \operatorname{carr}(g)))}}
\vnbprop{mg-group-opr}{\vnbstmt{\forall g.\; (\operatorname{is\text{-}metric\text{-}group}(g) \Rightarrow (\operatorname{opr}(\operatorname{metric\text{-}group\text{-}group}(g)) \simeq \operatorname{opr}(g)))}}
\vnbprop{mg-group-inv}{\vnbstmt{\forall g.\; (\operatorname{is\text{-}metric\text{-}group}(g) \Rightarrow (\operatorname{inv}(\operatorname{metric\text{-}group\text{-}group}(g)) \simeq \operatorname{inv}(g)))}}
\vnbprop{mg-group-iden}{\vnbstmt{\forall g.\; (\operatorname{is\text{-}metric\text{-}group}(g) \Rightarrow (\operatorname{iden}(\operatorname{metric\text{-}group\text{-}group}(g)) \simeq \operatorname{iden}(g)))}}
\vnbprop{mg-opr-in}{\vnbstmt{\begin{array}{@{}l@{}} \forall g, a, b. \\ \quad (\operatorname{is\text{-}metric\text{-}group}(g) \wedge ((a \in \operatorname{pts}(\operatorname{metric\text{-}group\text{-}ms}(g))) \wedge (b \in \operatorname{pts}(\operatorname{metric\text{-}group\text{-}ms}(g))))) \Rightarrow  \\ \quad \quad \operatorname{opr}(g)(a, b) \in \operatorname{pts}(\operatorname{metric\text{-}group\text{-}ms}(g)) \end{array}}}
\vnbprop{mg-inv-in}{\vnbstmt{\begin{array}{@{}l@{}} \forall g, a. \\ \quad (\operatorname{is\text{-}metric\text{-}group}(g) \wedge (a \in \operatorname{pts}(\operatorname{metric\text{-}group\text{-}ms}(g)))) \Rightarrow  \\ \quad \quad \operatorname{inv}(g)(a) \in \operatorname{pts}(\operatorname{metric\text{-}group\text{-}ms}(g)) \end{array}}}
\vnbprop{mg-iden-in}{\vnbstmt{\forall g.\; (\operatorname{is\text{-}metric\text{-}group}(g) \Rightarrow (\operatorname{iden}(g) \in \operatorname{pts}(\operatorname{metric\text{-}group\text{-}ms}(g))))}}
\vnbprop{mg-mul-inv-cancel}{\vnbstmt{\begin{array}{@{}l@{}} \forall g, a, y. \\ \quad (\operatorname{is\text{-}metric\text{-}group}(g) \wedge ((a \in \operatorname{pts}(\operatorname{metric\text{-}group\text{-}ms}(g))) \wedge (y \in \operatorname{pts}(\operatorname{metric\text{-}group\text{-}ms}(g))))) \Rightarrow  \\ \quad \quad \operatorname{opr}(g)(a, \operatorname{opr}(g)(\operatorname{inv}(g)(a), y)) = y \end{array}}}
\vnbprop{mg-inv-mul-cancel}{\vnbstmt{\begin{array}{@{}l@{}} \forall g, a, y. \\ \quad (\operatorname{is\text{-}metric\text{-}group}(g) \wedge ((a \in \operatorname{pts}(\operatorname{metric\text{-}group\text{-}ms}(g))) \wedge (y \in \operatorname{pts}(\operatorname{metric\text{-}group\text{-}ms}(g))))) \Rightarrow  \\ \quad \quad \operatorname{opr}(g)(\operatorname{inv}(g)(a), \operatorname{opr}(g)(a, y)) = y \end{array}}}
\vnbprop{mg-assoc}{\vnbstmt{\begin{array}{@{}l@{}} \forall g, a, b, y. \\ \quad \operatorname{is\text{-}metric\text{-}group}(g) \wedge \\ \quad a \in \operatorname{pts}(\operatorname{metric\text{-}group\text{-}ms}(g)) \wedge \\ \quad b \in \operatorname{pts}(\operatorname{metric\text{-}group\text{-}ms}(g)) \wedge \\ \quad y \in \operatorname{pts}(\operatorname{metric\text{-}group\text{-}ms}(g)) \Rightarrow \\ \quad \quad \operatorname{opr}(g)(\operatorname{opr}(g)(a, b), y) = \operatorname{opr}(g)(a, \operatorname{opr}(g)(b, y)) \end{array}}}
\vnbprop{mg-left-id}{\vnbstmt{\begin{array}{@{}l@{}} \forall g, y. \\ \quad (\operatorname{is\text{-}metric\text{-}group}(g) \wedge (y \in \operatorname{pts}(\operatorname{metric\text{-}group\text{-}ms}(g)))) \Rightarrow  \\ \quad \quad \operatorname{opr}(g)(\operatorname{iden}(g), y) = y \end{array}}}
\vnbprop{mg-right-id}{\vnbstmt{\begin{array}{@{}l@{}} \forall g, y. \\ \quad (\operatorname{is\text{-}metric\text{-}group}(g) \wedge (y \in \operatorname{pts}(\operatorname{metric\text{-}group\text{-}ms}(g)))) \Rightarrow  \\ \quad \quad \operatorname{opr}(g)(y, \operatorname{iden}(g)) = y \end{array}}}
\vnbprop{mg-mul-cont}{\vnbstmt{\begin{array}{@{}l@{}} \forall g, u, v, e. \\ \quad (\operatorname{is\text{-}metric\text{-}group}(g) \wedge ((u \in \operatorname{pts}(\operatorname{metric\text{-}group\text{-}ms}(g))) \wedge ((v \in \operatorname{pts}(\operatorname{metric\text{-}group\text{-}ms}(g))) \wedge \operatorname{pos\text{-}rr}(e)))) \Rightarrow  \\ \quad \quad \exists d. \\ \quad \quad \quad \operatorname{pos\text{-}rr}(d) \wedge \\ \quad \quad \quad \forall p, q \in \operatorname{pts}(\operatorname{metric\text{-}group\text{-}ms}(g)). \\ \quad \quad \quad \quad ((\operatorname{dist}(\operatorname{metric\text{-}group\text{-}ms}(g))(u, p) \leq d) \wedge (\operatorname{dist}(\operatorname{metric\text{-}group\text{-}ms}(g))(v, q) \leq d)) \Rightarrow  \\ \quad \quad \quad \quad \quad \operatorname{dist}(\operatorname{metric\text{-}group\text{-}ms}(g))(\operatorname{opr}(g)(u, v), \operatorname{opr}(g)(p, q)) \leq e \end{array}}}
\vnbprop{mg-ltrans-cont}{\vnbstmt{\begin{array}{@{}l@{}} \forall g, a, h. \\ \quad (\operatorname{is\text{-}metric\text{-}group}(g) \wedge ((a \in \operatorname{pts}(\operatorname{metric\text{-}group\text{-}ms}(g))) \wedge (h \in \operatorname{cont\text{-}fns}(\operatorname{metric\text{-}group\text{-}ms}(g))))) \Rightarrow  \\ \quad \quad (\lambda x \in \operatorname{pts}(\operatorname{metric\text{-}group\text{-}ms}(g)).\; h(\operatorname{opr}(g)(a, x))) \in \operatorname{cont\text{-}fns}(\operatorname{metric\text{-}group\text{-}ms}(g)) \end{array}}}
\vnbprop{ck-precomp-fun}{\vnbstmt{\begin{array}{@{}l@{}} \forall s, i, f. \\ \quad (\operatorname{is\text{-}compact}(s) \wedge ((i \in (\operatorname{pts}(s) \to \operatorname{pts}(s))) \wedge (f \in \operatorname{cont\text{-}fns}(s)))) \Rightarrow  \\ \quad \quad (\lambda x \in \operatorname{pts}(s).\; f(i(x))) \in (\operatorname{pts}(s) \to \mathbb{R}) \end{array}}}
\vnbprop{ck-precomp-affine}{\vnbstmt{\begin{array}{@{}l@{}} \forall s, i, u, v, t. \\ \quad \operatorname{is\text{-}compact}(s) \wedge \\ \quad i \in (\operatorname{pts}(s) \to \operatorname{pts}(s)) \wedge \\ \quad u \in \operatorname{cont\text{-}fns}(s) \wedge \\ \quad v \in \operatorname{cont\text{-}fns}(s) \wedge \\ \quad (\lambda x \in \operatorname{pts}(s).\; u(i(x))) \in \operatorname{cont\text{-}fns}(s) \wedge \\ \quad (\lambda x \in \operatorname{pts}(s).\; v(i(x))) \in \operatorname{cont\text{-}fns}(s) \wedge \\ \quad t \in \mathbb{R} \Rightarrow \\ \quad \quad (\lambda x \in \operatorname{pts}(s).\; \operatorname{vadd}(C({s}))(\operatorname{act}(C({s}))((1 - t), u), \operatorname{act}(C({s}))(t, v))(i(x))) \\ \quad \quad =\; \operatorname{vadd}(C({s}))(\operatorname{act}(C({s}))((1 - t), (\lambda x \in \operatorname{pts}(s).\; u(i(x)))), \\ \quad \quad \quad \operatorname{act}(C({s}))(t, (\lambda x \in \operatorname{pts}(s).\; v(i(x))))) \end{array}}}
\vnbprop{ck-precomp-dist-le}{\vnbstmt{\begin{array}{@{}l@{}} \forall s, i, f, g. \\ \quad \operatorname{is\text{-}compact}(s) \wedge \\ \quad i \in (\operatorname{pts}(s) \to \operatorname{pts}(s)) \wedge \\ \quad f \in \operatorname{cont\text{-}fns}(s) \wedge \\ \quad g \in \operatorname{cont\text{-}fns}(s) \wedge \\ \quad (\lambda x \in \operatorname{pts}(s).\; f(i(x))) \in \operatorname{cont\text{-}fns}(s) \wedge \\ \quad (\lambda x \in \operatorname{pts}(s).\; g(i(x))) \in \operatorname{cont\text{-}fns}(s) \Rightarrow \\ \quad \quad \operatorname{dist}(\operatorname{nvs\text{-}metric\text{-}space}(C({s})))((\lambda x \in \operatorname{pts}(s).\; f(i(x))), (\lambda x \in \operatorname{pts}(s).\; g(i(x)))) \\ \quad \quad \leq\; \operatorname{dist}(\operatorname{nvs\text{-}metric\text{-}space}(C({s})))(f, g) \end{array}}}
\vnbprop{ck-precomp-dist-eq}{\vnbstmt{\begin{array}{@{}l@{}} \forall s, i, f, g. \\ \quad \operatorname{is\text{-}compact}(s) \wedge \\ \quad i \in (\operatorname{pts}(s) \to \operatorname{pts}(s)) \wedge \\ \quad \forall y \in \operatorname{pts}(s).\; \exists x \in \operatorname{pts}(s).\; (i(x) = y) \wedge \\ \quad f \in \operatorname{cont\text{-}fns}(s) \wedge \\ \quad g \in \operatorname{cont\text{-}fns}(s) \wedge \\ \quad (\lambda \mathit{hix} \in \operatorname{pts}(s).\; f(i(\mathit{hix}))) \in \operatorname{cont\text{-}fns}(s) \wedge \\ \quad (\lambda \mathit{hix} \in \operatorname{pts}(s).\; g(i(\mathit{hix}))) \in \operatorname{cont\text{-}fns}(s) \Rightarrow \\ \quad \quad \operatorname{dist}(\operatorname{nvs\text{-}metric\text{-}space}(C({s})))((\lambda \mathit{hix} \in \operatorname{pts}(s).\; f(i(\mathit{hix}))), \\ \quad \quad \quad (\lambda \mathit{hix} \in \operatorname{pts}(s).\; g(i(\mathit{hix})))) \\ \quad \quad =\; \operatorname{dist}(\operatorname{nvs\text{-}metric\text{-}space}(C({s})))(f, g) \end{array}}}
\vnbprop{mg-lmul-fun}{\vnbstmt{\begin{array}{@{}l@{}} \forall g, a. \\ \quad (\operatorname{is\text{-}metric\text{-}group}(g) \wedge (a \in \operatorname{pts}(\operatorname{metric\text{-}group\text{-}ms}(g)))) \Rightarrow  \\ \quad \quad (\lambda x \in \operatorname{pts}(\operatorname{metric\text{-}group\text{-}ms}(g)).\; \operatorname{opr}(g)(a, x)) \\ \quad \quad \in\; (\operatorname{pts}(\operatorname{metric\text{-}group\text{-}ms}(g)) \to \operatorname{pts}(\operatorname{metric\text{-}group\text{-}ms}(g))) \end{array}}}
\vnbprop{mg-lmul-onto}{\vnbstmt{\begin{array}{@{}l@{}} \forall g, a.; \forall y \in \operatorname{pts}(\operatorname{metric\text{-}group\text{-}ms}(g)). \\ \quad (\operatorname{is\text{-}metric\text{-}group}(g) \wedge (a \in \operatorname{pts}(\operatorname{metric\text{-}group\text{-}ms}(g)))) \Rightarrow  \\ \quad \quad \exists x \in \operatorname{pts}(\operatorname{metric\text{-}group\text{-}ms}(g)). \\ \quad \quad \quad (\lambda \mathit{hqx} \in \operatorname{pts}(\operatorname{metric\text{-}group\text{-}ms}(g)).\; \operatorname{opr}(g)(a, \mathit{hqx}))(x) = y \end{array}}}
\vnbprop{mg-ltrans-precomp}{\vnbstmt{\begin{array}{@{}l@{}} \forall g, a, h. \\ \quad (\operatorname{is\text{-}metric\text{-}group}(g) \wedge ((a \in \operatorname{pts}(\operatorname{metric\text{-}group\text{-}ms}(g))) \wedge (h \in (\operatorname{pts}(\operatorname{metric\text{-}group\text{-}ms}(g)) \to \mathbb{R})))) \Rightarrow  \\ \quad \quad (\lambda x \in \operatorname{pts}(\operatorname{metric\text{-}group\text{-}ms}(g)).\; \\ \quad \quad \quad h((\lambda \mathit{hqx} \in \operatorname{pts}(\operatorname{metric\text{-}group\text{-}ms}(g)).\; \operatorname{opr}(g)(a, \mathit{hqx}))(x))) \\ \quad \quad =\; (\lambda x \in \operatorname{pts}(\operatorname{metric\text{-}group\text{-}ms}(g)).\; h(\operatorname{opr}(g)(a, x))) \end{array}}}
\vnbprop{mg-ltrans-dist}{\vnbstmt{\begin{array}{@{}l@{}} \forall g, a, f, b. \\ \quad \operatorname{is\text{-}metric\text{-}group}(g) \wedge \\ \quad \operatorname{is\text{-}compact}(\operatorname{metric\text{-}group\text{-}ms}(g)) \wedge \\ \quad a \in \operatorname{pts}(\operatorname{metric\text{-}group\text{-}ms}(g)) \wedge \\ \quad f \in \operatorname{cont\text{-}fns}(\operatorname{metric\text{-}group\text{-}ms}(g)) \wedge \\ \quad b \in \operatorname{cont\text{-}fns}(\operatorname{metric\text{-}group\text{-}ms}(g)) \Rightarrow \\ \quad \quad \operatorname{dist}(\operatorname{nvs\text{-}metric\text{-}space}(C({\operatorname{metric\text{-}group\text{-}ms}(g)})))((\lambda x \in \operatorname{pts}(\operatorname{metric\text{-}group\text{-}ms}(g)).\; f(\operatorname{opr}(g)(a, x))), \\ \quad \quad \quad (\lambda x \in \operatorname{pts}(\operatorname{metric\text{-}group\text{-}ms}(g)).\; b(\operatorname{opr}(g)(a, x)))) \\ \quad \quad =\; \operatorname{dist}(\operatorname{nvs\text{-}metric\text{-}space}(C({\operatorname{metric\text{-}group\text{-}ms}(g)})))(f, b) \end{array}}}
\vnbprop{mg-ltrans-affine}{\vnbstmt{\begin{array}{@{}l@{}} \forall g, a, u, v, t. \\ \quad \operatorname{is\text{-}metric\text{-}group}(g) \wedge \\ \quad \operatorname{is\text{-}compact}(\operatorname{metric\text{-}group\text{-}ms}(g)) \wedge \\ \quad a \in \operatorname{pts}(\operatorname{metric\text{-}group\text{-}ms}(g)) \wedge \\ \quad u \in \operatorname{cont\text{-}fns}(\operatorname{metric\text{-}group\text{-}ms}(g)) \wedge \\ \quad v \in \operatorname{cont\text{-}fns}(\operatorname{metric\text{-}group\text{-}ms}(g)) \wedge \\ \quad t \in \mathbb{R} \Rightarrow \\ \quad \quad (\lambda x \in \operatorname{pts}(\operatorname{metric\text{-}group\text{-}ms}(g)).\; \\ \quad \quad \quad \operatorname{vadd}(C({\operatorname{metric\text{-}group\text{-}ms}(g)}))(\operatorname{act}(C({\operatorname{metric\text{-}group\text{-}ms}(g)}))((1 - t), u), \operatorname{act}(C({\operatorname{metric\text{-}group\text{-}ms}(g)}))(t, v))(\operatorname{opr}(g)(a, x))) \\ \quad \quad =\; \operatorname{vadd}(C({\operatorname{metric\text{-}group\text{-}ms}(g)}))(\operatorname{act}(C({\operatorname{metric\text{-}group\text{-}ms}(g)}))((1 - t), (\lambda x \in \operatorname{pts}(\operatorname{metric\text{-}group\text{-}ms}(g)).\; u(\operatorname{opr}(g)(a, x)))), \\ \quad \quad \quad \operatorname{act}(C({\operatorname{metric\text{-}group\text{-}ms}(g)}))(t, (\lambda x \in \operatorname{pts}(\operatorname{metric\text{-}group\text{-}ms}(g)).\; v(\operatorname{opr}(g)(a, x))))) \end{array}}}
\vnbprop{mg-orbit-joint-cont}{\vnbstmt{\begin{array}{@{}l@{}} \forall g, f. \\ \quad (\operatorname{is\text{-}metric\text{-}group}(g) \wedge (f \in \operatorname{cont\text{-}fns}(\operatorname{metric\text{-}group\text{-}ms}(g)))) \Rightarrow  \\ \quad \quad \operatorname{is\text{-}continuous}(\operatorname{product2\text{-}ms}(\operatorname{metric\text{-}group\text{-}ms}(g), \operatorname{metric\text{-}group\text{-}ms}(g)), \\ \quad \quad \quad \operatorname{rr\text{-}ms}, \\ \quad \quad \quad (\lambda z \in (\operatorname{pts}(\operatorname{metric\text{-}group\text{-}ms}(g)) \times \operatorname{pts}(\operatorname{metric\text{-}group\text{-}ms}(g))).\; \\ \quad \quad \quad \quad f(\operatorname{opr}(g)(\operatorname{nth}(1, z), \operatorname{nth}(2, z))))) \end{array}}}
\vnbprop{mg-orbit-uc}{\vnbstmt{\begin{array}{@{}l@{}} \forall g, f, e. \\ \quad \operatorname{is\text{-}metric\text{-}group}(g) \wedge \\ \quad \operatorname{is\text{-}compact}(\operatorname{metric\text{-}group\text{-}ms}(g)) \wedge \\ \quad f \in \operatorname{cont\text{-}fns}(\operatorname{metric\text{-}group\text{-}ms}(g)) \wedge \\ \quad \operatorname{pos\text{-}rr}(e) \Rightarrow \\ \quad \quad \exists d. \\ \quad \quad \quad \operatorname{pos\text{-}rr}(d) \wedge \\ \quad \quad \quad \forall a, b \in \operatorname{pts}(\operatorname{metric\text{-}group\text{-}ms}(g)). \\ \quad \quad \quad \quad (\operatorname{dist}(\operatorname{metric\text{-}group\text{-}ms}(g))(a, b) \leq d) \Rightarrow  \\ \quad \quad \quad \quad \quad \operatorname{dist}(\operatorname{nvs\text{-}metric\text{-}space}(C({\operatorname{metric\text{-}group\text{-}ms}(g)})))((\lambda x \in \operatorname{pts}(\operatorname{metric\text{-}group\text{-}ms}(g)).\; f(\operatorname{opr}(g)(a, x))), \\ \quad \quad \quad \quad \quad \quad (\lambda x \in \operatorname{pts}(\operatorname{metric\text{-}group\text{-}ms}(g)).\; f(\operatorname{opr}(g)(b, x)))) \\ \quad \quad \quad \quad \quad \leq\; e \end{array}}}
\vnbprop{mg-orbit-map-fun}{\vnbstmt{\begin{array}{@{}l@{}} \forall g, f. \\ \quad (\operatorname{is\text{-}metric\text{-}group}(g) \wedge (\operatorname{is\text{-}compact}(\operatorname{metric\text{-}group\text{-}ms}(g)) \wedge (f \in \operatorname{cont\text{-}fns}(\operatorname{metric\text{-}group\text{-}ms}(g))))) \Rightarrow  \\ \quad \quad (\lambda a \in \operatorname{pts}(\operatorname{metric\text{-}group\text{-}ms}(g)).\; (\lambda x \in \operatorname{pts}(\operatorname{metric\text{-}group\text{-}ms}(g)).\; f(\operatorname{opr}(g)(a, x)))) \\ \quad \quad \in\; (\operatorname{pts}(\operatorname{metric\text{-}group\text{-}ms}(g)) \to \operatorname{pts}(\operatorname{nvs\text{-}metric\text{-}space}(C({\operatorname{metric\text{-}group\text{-}ms}(g)})))) \end{array}}}
\vnbprop{mg-orbit-map-cont}{\vnbstmt{\begin{array}{@{}l@{}} \forall g, f. \\ \quad (\operatorname{is\text{-}metric\text{-}group}(g) \wedge (\operatorname{is\text{-}compact}(\operatorname{metric\text{-}group\text{-}ms}(g)) \wedge (f \in \operatorname{cont\text{-}fns}(\operatorname{metric\text{-}group\text{-}ms}(g))))) \Rightarrow  \\ \quad \quad \operatorname{is\text{-}continuous}(\operatorname{metric\text{-}group\text{-}ms}(g), \operatorname{nvs\text{-}metric\text{-}space}(C({\operatorname{metric\text{-}group\text{-}ms}(g)})), \\ \quad \quad \quad (\lambda a \in \operatorname{pts}(\operatorname{metric\text{-}group\text{-}ms}(g)).\; (\lambda x \in \operatorname{pts}(\operatorname{metric\text{-}group\text{-}ms}(g)).\; f(\operatorname{opr}(g)(a, x))))) \end{array}}}
\vnbprop{mg-orbit-compact}{\vnbstmt{\begin{array}{@{}l@{}} \forall g, f. \\ \quad (\operatorname{is\text{-}metric\text{-}group}(g) \wedge (\operatorname{is\text{-}compact}(\operatorname{metric\text{-}group\text{-}ms}(g)) \wedge (f \in \operatorname{cont\text{-}fns}(\operatorname{metric\text{-}group\text{-}ms}(g))))) \Rightarrow  \\ \quad \quad \operatorname{is\text{-}compact}(\operatorname{subspace\text{-}ms}(\operatorname{nvs\text{-}metric\text{-}space}(C({\operatorname{metric\text{-}group\text{-}ms}(g)})), \\ \quad \quad \quad \quad \operatorname{image}((\lambda a \in \operatorname{pts}(\operatorname{metric\text{-}group\text{-}ms}(g)).\; (\lambda x \in \operatorname{pts}(\operatorname{metric\text{-}group\text{-}ms}(g)).\; f(\operatorname{opr}(g)(a, x)))), \\ \quad \quad \quad \quad \quad \operatorname{pts}(\operatorname{metric\text{-}group\text{-}ms}(g))))) \end{array}}}
\vnbprop{mg-orbit-sub-vec}{\vnbstmt{\begin{array}{@{}l@{}} \forall g, f. \\ \quad (\operatorname{is\text{-}metric\text{-}group}(g) \wedge (\operatorname{is\text{-}compact}(\operatorname{metric\text{-}group\text{-}ms}(g)) \wedge (f \in \operatorname{cont\text{-}fns}(\operatorname{metric\text{-}group\text{-}ms}(g))))) \Rightarrow  \\ \quad \quad \operatorname{image}((\lambda a \in \operatorname{pts}(\operatorname{metric\text{-}group\text{-}ms}(g)).\; (\lambda x \in \operatorname{pts}(\operatorname{metric\text{-}group\text{-}ms}(g)).\; f(\operatorname{opr}(g)(a, x)))), \\ \quad \quad \quad \operatorname{pts}(\operatorname{metric\text{-}group\text{-}ms}(g))) \\ \quad \quad \subseteq\; \operatorname{vec}(C({\operatorname{metric\text{-}group\text{-}ms}(g)})) \end{array}}}
\vnbprop{mg-orbit-tb}{\vnbstmt{\begin{array}{@{}l@{}} \forall g, f. \\ \quad (\operatorname{is\text{-}metric\text{-}group}(g) \wedge (\operatorname{is\text{-}compact}(\operatorname{metric\text{-}group\text{-}ms}(g)) \wedge (f \in \operatorname{cont\text{-}fns}(\operatorname{metric\text{-}group\text{-}ms}(g))))) \Rightarrow  \\ \quad \quad \operatorname{is\text{-}totally\text{-}bounded\text{-}in}(\operatorname{nvs\text{-}metric\text{-}space}(C({\operatorname{metric\text{-}group\text{-}ms}(g)})), \\ \quad \quad \quad \operatorname{image}((\lambda a \in \operatorname{pts}(\operatorname{metric\text{-}group\text{-}ms}(g)).\; (\lambda x \in \operatorname{pts}(\operatorname{metric\text{-}group\text{-}ms}(g)).\; f(\operatorname{opr}(g)(a, x)))), \\ \quad \quad \quad \quad \operatorname{pts}(\operatorname{metric\text{-}group\text{-}ms}(g)))) \end{array}}}
\vnbprop{ms-tb-in-subspace-tb}{\vnbstmt{\begin{array}{@{}l@{}} \forall s, k. \\ \quad (\operatorname{is\text{-}metric\text{-}space}(s) \wedge ((k \subseteq \operatorname{pts}(s)) \wedge \operatorname{is\text{-}totally\text{-}bounded\text{-}in}(s, k))) \Rightarrow  \\ \quad \quad \operatorname{totally\text{-}bounded}(\operatorname{subspace\text{-}ms}(s, k)) \end{array}}}
\vnbprop{ms-closed-tb-in-compact}{\vnbstmt{\begin{array}{@{}l@{}} \forall s, k. \\ \quad (\operatorname{is\text{-}complete}(s) \wedge ((k \subseteq \operatorname{pts}(s)) \wedge (\operatorname{is\text{-}closed}(s, k) \wedge \operatorname{is\text{-}totally\text{-}bounded\text{-}in}(s, k)))) \Rightarrow  \\ \quad \quad \operatorname{is\text{-}compact}(\operatorname{subspace\text{-}ms}(s, k)) \end{array}}}
\vnbprop{ms-closure-near}{\vnbstmt{\begin{array}{@{}l@{}} \forall s, a.; \forall x \in \operatorname{closure}(s, a), r. \\ \quad (\operatorname{is\text{-}metric\text{-}space}(s) \wedge \operatorname{pos\text{-}rr}(r)) \Rightarrow \exists y \in a.\; (\operatorname{dist}(s)(x, y) \leq r) \end{array}}}
\vnbprop{ms-closure-intro}{\vnbstmt{\begin{array}{@{}l@{}} \forall s, a.; \forall x \in \operatorname{pts}(s). \\ \quad \operatorname{is\text{-}metric\text{-}space}(s) \wedge \\ \quad a \subseteq \operatorname{pts}(s) \wedge \\ \quad \forall r.\; (\operatorname{pos\text{-}rr}(r) \Rightarrow \exists y \in a.\; (\operatorname{dist}(s)(x, y) \leq r)) \Rightarrow \\ \quad \quad x \in \operatorname{closure}(s, a) \end{array}}}
\vnbprop{nvs-closure-convex}{\vnbstmt{\begin{array}{@{}l@{}} \forall m, s. \\ \quad (\operatorname{is\text{-}normed\text{-}vector\text{-}space}(m) \wedge ((s \subseteq \operatorname{vec}(m)) \wedge \operatorname{is\text{-}convex\text{-}subset}(m, s))) \Rightarrow  \\ \quad \quad \operatorname{is\text{-}convex\text{-}subset}(m, \operatorname{closure}(\operatorname{nvs\text{-}metric\text{-}space}(m), s)) \end{array}}}
\vnbprop{ms-closure-tb}{\vnbstmt{\begin{array}{@{}l@{}} \forall s, a. \\ \quad (\operatorname{is\text{-}metric\text{-}space}(s) \wedge ((a \subseteq \operatorname{pts}(s)) \wedge \operatorname{is\text{-}totally\text{-}bounded\text{-}in}(s, a))) \Rightarrow  \\ \quad \quad \operatorname{is\text{-}totally\text{-}bounded\text{-}in}(s, \operatorname{closure}(s, a)) \end{array}}}
\vnbprop{mg-hull-sub-vec}{\vnbstmt{\begin{array}{@{}l@{}} \forall g, f. \\ \quad (\operatorname{is\text{-}metric\text{-}group}(g) \wedge (\operatorname{is\text{-}compact}(\operatorname{metric\text{-}group\text{-}ms}(g)) \wedge (f \in \operatorname{cont\text{-}fns}(\operatorname{metric\text{-}group\text{-}ms}(g))))) \Rightarrow  \\ \quad \quad \operatorname{closure}(\operatorname{nvs\text{-}metric\text{-}space}(C({\operatorname{metric\text{-}group\text{-}ms}(g)})), \\ \quad \quad \quad \operatorname{conv\text{-}hull}(C({\operatorname{metric\text{-}group\text{-}ms}(g)}), \\ \quad \quad \quad \quad \operatorname{image}((\lambda a \in \operatorname{pts}(\operatorname{metric\text{-}group\text{-}ms}(g)).\; (\lambda x \in \operatorname{pts}(\operatorname{metric\text{-}group\text{-}ms}(g)).\; f(\operatorname{opr}(g)(a, x)))), \\ \quad \quad \quad \quad \quad \operatorname{pts}(\operatorname{metric\text{-}group\text{-}ms}(g))))) \\ \quad \quad \subseteq\; \operatorname{vec}(C({\operatorname{metric\text{-}group\text{-}ms}(g)})) \end{array}}}
\vnbprop{mg-hull-convex}{\vnbstmt{\begin{array}{@{}l@{}} \forall g, f. \\ \quad (\operatorname{is\text{-}metric\text{-}group}(g) \wedge (\operatorname{is\text{-}compact}(\operatorname{metric\text{-}group\text{-}ms}(g)) \wedge (f \in \operatorname{cont\text{-}fns}(\operatorname{metric\text{-}group\text{-}ms}(g))))) \Rightarrow  \\ \quad \quad \operatorname{is\text{-}convex\text{-}subset}(C({\operatorname{metric\text{-}group\text{-}ms}(g)}), \\ \quad \quad \quad \operatorname{closure}(\operatorname{nvs\text{-}metric\text{-}space}(C({\operatorname{metric\text{-}group\text{-}ms}(g)})), \\ \quad \quad \quad \quad \operatorname{conv\text{-}hull}(C({\operatorname{metric\text{-}group\text{-}ms}(g)}), \\ \quad \quad \quad \quad \quad \operatorname{image}((\lambda a \in \operatorname{pts}(\operatorname{metric\text{-}group\text{-}ms}(g)).\; (\lambda x \in \operatorname{pts}(\operatorname{metric\text{-}group\text{-}ms}(g)).\; f(\operatorname{opr}(g)(a, x)))), \\ \quad \quad \quad \quad \quad \quad \operatorname{pts}(\operatorname{metric\text{-}group\text{-}ms}(g)))))) \end{array}}}
\vnbprop{mg-hull-compact}{\vnbstmt{\begin{array}{@{}l@{}} \forall g, f. \\ \quad (\operatorname{is\text{-}metric\text{-}group}(g) \wedge (\operatorname{is\text{-}compact}(\operatorname{metric\text{-}group\text{-}ms}(g)) \wedge (f \in \operatorname{cont\text{-}fns}(\operatorname{metric\text{-}group\text{-}ms}(g))))) \Rightarrow  \\ \quad \quad \operatorname{is\text{-}compact}(\operatorname{subspace\text{-}ms}(\operatorname{nvs\text{-}metric\text{-}space}(C({\operatorname{metric\text{-}group\text{-}ms}(g)})), \\ \quad \quad \quad \quad \operatorname{closure}(\operatorname{nvs\text{-}metric\text{-}space}(C({\operatorname{metric\text{-}group\text{-}ms}(g)})), \\ \quad \quad \quad \quad \quad \operatorname{conv\text{-}hull}(C({\operatorname{metric\text{-}group\text{-}ms}(g)}), \\ \quad \quad \quad \quad \quad \quad \operatorname{image}((\lambda a \in \operatorname{pts}(\operatorname{metric\text{-}group\text{-}ms}(g)).\; (\lambda x \in \operatorname{pts}(\operatorname{metric\text{-}group\text{-}ms}(g)).\; f(\operatorname{opr}(g)(a, x)))), \\ \quad \quad \quad \quad \quad \quad \quad \operatorname{pts}(\operatorname{metric\text{-}group\text{-}ms}(g))))))) \end{array}}}
\vnbprop{mg-hull-has-translates}{\vnbstmt{\begin{array}{@{}l@{}} \forall g, f, a. \\ \quad \operatorname{is\text{-}metric\text{-}group}(g) \wedge \\ \quad \operatorname{is\text{-}compact}(\operatorname{metric\text{-}group\text{-}ms}(g)) \wedge \\ \quad f \in \operatorname{cont\text{-}fns}(\operatorname{metric\text{-}group\text{-}ms}(g)) \wedge \\ \quad a \in \operatorname{pts}(\operatorname{metric\text{-}group\text{-}ms}(g)) \Rightarrow \\ \quad \quad (\lambda x \in \operatorname{pts}(\operatorname{metric\text{-}group\text{-}ms}(g)).\; f(\operatorname{opr}(g)(a, x))) \\ \quad \quad \in\; \operatorname{closure}(\operatorname{nvs\text{-}metric\text{-}space}(C({\operatorname{metric\text{-}group\text{-}ms}(g)})), \\ \quad \quad \quad \operatorname{conv\text{-}hull}(C({\operatorname{metric\text{-}group\text{-}ms}(g)}), \\ \quad \quad \quad \quad \operatorname{image}((\lambda \mathit{hqa} \in \operatorname{pts}(\operatorname{metric\text{-}group\text{-}ms}(g)).\; \\ \quad \quad \quad \quad \quad \quad (\lambda x \in \operatorname{pts}(\operatorname{metric\text{-}group\text{-}ms}(g)).\; f(\operatorname{opr}(g)(\mathit{hqa}, x)))), \\ \quad \quad \quad \quad \quad \operatorname{pts}(\operatorname{metric\text{-}group\text{-}ms}(g))))) \end{array}}}
\vnbprop{mg-hull-nonempty}{\vnbstmt{\begin{array}{@{}l@{}} \forall g, f. \\ \quad (\operatorname{is\text{-}metric\text{-}group}(g) \wedge (\operatorname{is\text{-}compact}(\operatorname{metric\text{-}group\text{-}ms}(g)) \wedge (f \in \operatorname{cont\text{-}fns}(\operatorname{metric\text{-}group\text{-}ms}(g))))) \Rightarrow  \\ \quad \quad \exists k. \\ \quad \quad \quad k \\ \quad \quad \quad \in\; \operatorname{closure}(\operatorname{nvs\text{-}metric\text{-}space}(C({\operatorname{metric\text{-}group\text{-}ms}(g)})), \\ \quad \quad \quad \quad \operatorname{conv\text{-}hull}(C({\operatorname{metric\text{-}group\text{-}ms}(g)}), \\ \quad \quad \quad \quad \quad \operatorname{image}((\lambda a \in \operatorname{pts}(\operatorname{metric\text{-}group\text{-}ms}(g)).\; (\lambda x \in \operatorname{pts}(\operatorname{metric\text{-}group\text{-}ms}(g)).\; f(\operatorname{opr}(g)(a, x)))), \\ \quad \quad \quad \quad \quad \quad \operatorname{pts}(\operatorname{metric\text{-}group\text{-}ms}(g))))) \end{array}}}
\vnbprop{mg-ltrans-comp}{\vnbstmt{\begin{array}{@{}l@{}} \forall g, a, b, h. \\ \quad \operatorname{is\text{-}metric\text{-}group}(g) \wedge \\ \quad a \in \operatorname{pts}(\operatorname{metric\text{-}group\text{-}ms}(g)) \wedge \\ \quad b \in \operatorname{pts}(\operatorname{metric\text{-}group\text{-}ms}(g)) \wedge \\ \quad h \in (\operatorname{pts}(\operatorname{metric\text{-}group\text{-}ms}(g)) \to \mathbb{R}) \Rightarrow \\ \quad \quad (\lambda x \in \operatorname{pts}(\operatorname{metric\text{-}group\text{-}ms}(g)).\; \\ \quad \quad \quad (\lambda x \in \operatorname{pts}(\operatorname{metric\text{-}group\text{-}ms}(g)).\; h(\operatorname{opr}(g)(b, x)))(\operatorname{opr}(g)(a, x))) \\ \quad \quad =\; (\lambda x \in \operatorname{pts}(\operatorname{metric\text{-}group\text{-}ms}(g)).\; h(\operatorname{opr}(g)(\operatorname{opr}(g)(b, a), x))) \end{array}}}
\vnbprop{mg-orbit-mem}{\vnbstmt{\begin{array}{@{}l@{}} \forall g, f, a. \\ \quad \operatorname{is\text{-}metric\text{-}group}(g) \wedge \\ \quad \operatorname{is\text{-}compact}(\operatorname{metric\text{-}group\text{-}ms}(g)) \wedge \\ \quad f \in \operatorname{cont\text{-}fns}(\operatorname{metric\text{-}group\text{-}ms}(g)) \wedge \\ \quad a \in \operatorname{pts}(\operatorname{metric\text{-}group\text{-}ms}(g)) \Rightarrow \\ \quad \quad (\lambda x \in \operatorname{pts}(\operatorname{metric\text{-}group\text{-}ms}(g)).\; f(\operatorname{opr}(g)(a, x))) \\ \quad \quad \in\; \operatorname{image}((\lambda \mathit{hqa} \in \operatorname{pts}(\operatorname{metric\text{-}group\text{-}ms}(g)).\; \\ \quad \quad \quad \quad (\lambda x \in \operatorname{pts}(\operatorname{metric\text{-}group\text{-}ms}(g)).\; f(\operatorname{opr}(g)(\mathit{hqa}, x)))), \\ \quad \quad \quad \operatorname{pts}(\operatorname{metric\text{-}group\text{-}ms}(g))) \end{array}}}
\vnbprop{mg-ltrans-orbit}{\vnbstmt{\begin{array}{@{}l@{}} \forall g, f, a, w. \\ \quad \operatorname{is\text{-}metric\text{-}group}(g) \wedge \\ \quad \operatorname{is\text{-}compact}(\operatorname{metric\text{-}group\text{-}ms}(g)) \wedge \\ \quad f \in \operatorname{cont\text{-}fns}(\operatorname{metric\text{-}group\text{-}ms}(g)) \wedge \\ \quad a \in \operatorname{pts}(\operatorname{metric\text{-}group\text{-}ms}(g)) \wedge \\ \quad w \\ \quad \in\; \operatorname{image}((\lambda \mathit{hqa} \in \operatorname{pts}(\operatorname{metric\text{-}group\text{-}ms}(g)).\; \\ \quad \quad \quad (\lambda x \in \operatorname{pts}(\operatorname{metric\text{-}group\text{-}ms}(g)).\; f(\operatorname{opr}(g)(\mathit{hqa}, x)))), \\ \quad \quad \operatorname{pts}(\operatorname{metric\text{-}group\text{-}ms}(g))) \Rightarrow \\ \quad \quad (\lambda x \in \operatorname{pts}(\operatorname{metric\text{-}group\text{-}ms}(g)).\; w(\operatorname{opr}(g)(a, x))) \\ \quad \quad \in\; \operatorname{image}((\lambda \mathit{hqa} \in \operatorname{pts}(\operatorname{metric\text{-}group\text{-}ms}(g)).\; \\ \quad \quad \quad \quad (\lambda x \in \operatorname{pts}(\operatorname{metric\text{-}group\text{-}ms}(g)).\; f(\operatorname{opr}(g)(\mathit{hqa}, x)))), \\ \quad \quad \quad \operatorname{pts}(\operatorname{metric\text{-}group\text{-}ms}(g))) \end{array}}}
\vnbprop{mg-ltrans-hull}{\vnbstmt{\begin{array}{@{}l@{}} \forall g, f, a, w. \\ \quad \operatorname{is\text{-}metric\text{-}group}(g) \wedge \\ \quad \operatorname{is\text{-}compact}(\operatorname{metric\text{-}group\text{-}ms}(g)) \wedge \\ \quad f \in \operatorname{cont\text{-}fns}(\operatorname{metric\text{-}group\text{-}ms}(g)) \wedge \\ \quad a \in \operatorname{pts}(\operatorname{metric\text{-}group\text{-}ms}(g)) \wedge \\ \quad w \\ \quad \in\; \operatorname{conv\text{-}hull}(C({\operatorname{metric\text{-}group\text{-}ms}(g)}), \\ \quad \quad \operatorname{image}((\lambda \mathit{hqa} \in \operatorname{pts}(\operatorname{metric\text{-}group\text{-}ms}(g)).\; \\ \quad \quad \quad \quad (\lambda x \in \operatorname{pts}(\operatorname{metric\text{-}group\text{-}ms}(g)).\; f(\operatorname{opr}(g)(\mathit{hqa}, x)))), \\ \quad \quad \quad \operatorname{pts}(\operatorname{metric\text{-}group\text{-}ms}(g)))) \Rightarrow \\ \quad \quad (\lambda x \in \operatorname{pts}(\operatorname{metric\text{-}group\text{-}ms}(g)).\; w(\operatorname{opr}(g)(a, x))) \\ \quad \quad \in\; \operatorname{conv\text{-}hull}(C({\operatorname{metric\text{-}group\text{-}ms}(g)}), \\ \quad \quad \quad \operatorname{image}((\lambda \mathit{hqa} \in \operatorname{pts}(\operatorname{metric\text{-}group\text{-}ms}(g)).\; \\ \quad \quad \quad \quad \quad (\lambda x \in \operatorname{pts}(\operatorname{metric\text{-}group\text{-}ms}(g)).\; f(\operatorname{opr}(g)(\mathit{hqa}, x)))), \\ \quad \quad \quad \quad \operatorname{pts}(\operatorname{metric\text{-}group\text{-}ms}(g)))) \end{array}}}
\vnbprop{mg-ltrans-k}{\vnbstmt{\begin{array}{@{}l@{}} \forall g, f, a, h. \\ \quad \operatorname{is\text{-}metric\text{-}group}(g) \wedge \\ \quad \operatorname{is\text{-}compact}(\operatorname{metric\text{-}group\text{-}ms}(g)) \wedge \\ \quad f \in \operatorname{cont\text{-}fns}(\operatorname{metric\text{-}group\text{-}ms}(g)) \wedge \\ \quad a \in \operatorname{pts}(\operatorname{metric\text{-}group\text{-}ms}(g)) \wedge \\ \quad h \\ \quad \in\; \operatorname{closure}(\operatorname{nvs\text{-}metric\text{-}space}(C({\operatorname{metric\text{-}group\text{-}ms}(g)})), \\ \quad \quad \operatorname{conv\text{-}hull}(C({\operatorname{metric\text{-}group\text{-}ms}(g)}), \\ \quad \quad \quad \operatorname{image}((\lambda \mathit{hqa} \in \operatorname{pts}(\operatorname{metric\text{-}group\text{-}ms}(g)).\; \\ \quad \quad \quad \quad \quad (\lambda x \in \operatorname{pts}(\operatorname{metric\text{-}group\text{-}ms}(g)).\; f(\operatorname{opr}(g)(\mathit{hqa}, x)))), \\ \quad \quad \quad \quad \operatorname{pts}(\operatorname{metric\text{-}group\text{-}ms}(g))))) \Rightarrow \\ \quad \quad (\lambda x \in \operatorname{pts}(\operatorname{metric\text{-}group\text{-}ms}(g)).\; h(\operatorname{opr}(g)(a, x))) \\ \quad \quad \in\; \operatorname{closure}(\operatorname{nvs\text{-}metric\text{-}space}(C({\operatorname{metric\text{-}group\text{-}ms}(g)})), \\ \quad \quad \quad \operatorname{conv\text{-}hull}(C({\operatorname{metric\text{-}group\text{-}ms}(g)}), \\ \quad \quad \quad \quad \operatorname{image}((\lambda \mathit{hqa} \in \operatorname{pts}(\operatorname{metric\text{-}group\text{-}ms}(g)).\; \\ \quad \quad \quad \quad \quad \quad (\lambda x \in \operatorname{pts}(\operatorname{metric\text{-}group\text{-}ms}(g)).\; f(\operatorname{opr}(g)(\mathit{hqa}, x)))), \\ \quad \quad \quad \quad \quad \operatorname{pts}(\operatorname{metric\text{-}group\text{-}ms}(g))))) \end{array}}}
\vnbprop{mg-left-inv}{\vnbstmt{\begin{array}{@{}l@{}} \forall g, a. \\ \quad (\operatorname{is\text{-}metric\text{-}group}(g) \wedge (a \in \operatorname{pts}(\operatorname{metric\text{-}group\text{-}ms}(g)))) \Rightarrow  \\ \quad \quad \operatorname{opr}(g)(\operatorname{inv}(g)(a), a) = \operatorname{iden}(g) \end{array}}}
\vnbprop{mg-ltrans-iden}{\vnbstmt{\begin{array}{@{}l@{}} \forall g, h. \\ \quad (\operatorname{is\text{-}metric\text{-}group}(g) \wedge (h \in (\operatorname{pts}(\operatorname{metric\text{-}group\text{-}ms}(g)) \to \mathbb{R}))) \Rightarrow  \\ \quad \quad (\lambda x \in \operatorname{pts}(\operatorname{metric\text{-}group\text{-}ms}(g)).\; h(\operatorname{opr}(g)(\operatorname{iden}(g), x))) = h \end{array}}}
\vnbprop{mg-t-fun}{\vnbstmt{\begin{array}{@{}l@{}} \forall g, f, a. \\ \quad \operatorname{is\text{-}metric\text{-}group}(g) \wedge \\ \quad \operatorname{is\text{-}compact}(\operatorname{metric\text{-}group\text{-}ms}(g)) \wedge \\ \quad f \in \operatorname{cont\text{-}fns}(\operatorname{metric\text{-}group\text{-}ms}(g)) \wedge \\ \quad a \in \operatorname{pts}(\operatorname{metric\text{-}group\text{-}ms}(g)) \Rightarrow \\ \quad \quad (\lambda h \in \operatorname{closure}(\operatorname{nvs\text{-}metric\text{-}space}(C({\operatorname{metric\text{-}group\text{-}ms}(g)})), \operatorname{conv\text{-}hull}(C({\operatorname{metric\text{-}group\text{-}ms}(g)}), \operatorname{image}((\lambda \mathit{hqa} \in \operatorname{pts}(\operatorname{metric\text{-}group\text{-}ms}(g)).\; (\lambda x \in \operatorname{pts}(\operatorname{metric\text{-}group\text{-}ms}(g)).\; f(\operatorname{opr}(g)(\mathit{hqa}, x)))), \operatorname{pts}(\operatorname{metric\text{-}group\text{-}ms}(g))))).\; \\ \quad \quad \quad (\lambda x \in \operatorname{pts}(\operatorname{metric\text{-}group\text{-}ms}(g)).\; h(\operatorname{opr}(g)(a, x)))) \\ \quad \quad \in\; (\operatorname{closure}(\operatorname{nvs\text{-}metric\text{-}space}(C({\operatorname{metric\text{-}group\text{-}ms}(g)})), \operatorname{conv\text{-}hull}(C({\operatorname{metric\text{-}group\text{-}ms}(g)}), \operatorname{image}((\lambda \mathit{hqa} \in \operatorname{pts}(\operatorname{metric\text{-}group\text{-}ms}(g)).\; (\lambda x \in \operatorname{pts}(\operatorname{metric\text{-}group\text{-}ms}(g)).\; f(\operatorname{opr}(g)(\mathit{hqa}, x)))), \operatorname{pts}(\operatorname{metric\text{-}group\text{-}ms}(g))))) \to \operatorname{closure}(\operatorname{nvs\text{-}metric\text{-}space}(C({\operatorname{metric\text{-}group\text{-}ms}(g)})), \operatorname{conv\text{-}hull}(C({\operatorname{metric\text{-}group\text{-}ms}(g)}), \operatorname{image}((\lambda \mathit{hqa} \in \operatorname{pts}(\operatorname{metric\text{-}group\text{-}ms}(g)).\; (\lambda x \in \operatorname{pts}(\operatorname{metric\text{-}group\text{-}ms}(g)).\; f(\operatorname{opr}(g)(\mathit{hqa}, x)))), \operatorname{pts}(\operatorname{metric\text{-}group\text{-}ms}(g)))))) \end{array}}}
\vnbprop{mg-t-apply}{\vnbstmt{\begin{array}{@{}l@{}} \forall g, f, a, h. \\ \quad \operatorname{is\text{-}metric\text{-}group}(g) \wedge \\ \quad \operatorname{is\text{-}compact}(\operatorname{metric\text{-}group\text{-}ms}(g)) \wedge \\ \quad f \in \operatorname{cont\text{-}fns}(\operatorname{metric\text{-}group\text{-}ms}(g)) \wedge \\ \quad a \in \operatorname{pts}(\operatorname{metric\text{-}group\text{-}ms}(g)) \wedge \\ \quad h \\ \quad \in\; \operatorname{closure}(\operatorname{nvs\text{-}metric\text{-}space}(C({\operatorname{metric\text{-}group\text{-}ms}(g)})), \\ \quad \quad \operatorname{conv\text{-}hull}(C({\operatorname{metric\text{-}group\text{-}ms}(g)}), \\ \quad \quad \quad \operatorname{image}((\lambda \mathit{hqa} \in \operatorname{pts}(\operatorname{metric\text{-}group\text{-}ms}(g)).\; \\ \quad \quad \quad \quad \quad (\lambda x \in \operatorname{pts}(\operatorname{metric\text{-}group\text{-}ms}(g)).\; f(\operatorname{opr}(g)(\mathit{hqa}, x)))), \\ \quad \quad \quad \quad \operatorname{pts}(\operatorname{metric\text{-}group\text{-}ms}(g))))) \Rightarrow \\ \quad \quad (\lambda \mathit{hqh} \in \operatorname{closure}(\operatorname{nvs\text{-}metric\text{-}space}(C({\operatorname{metric\text{-}group\text{-}ms}(g)})), \operatorname{conv\text{-}hull}(C({\operatorname{metric\text{-}group\text{-}ms}(g)}), \operatorname{image}((\lambda \mathit{hqa} \in \operatorname{pts}(\operatorname{metric\text{-}group\text{-}ms}(g)).\; (\lambda x \in \operatorname{pts}(\operatorname{metric\text{-}group\text{-}ms}(g)).\; f(\operatorname{opr}(g)(\mathit{hqa}, x)))), \operatorname{pts}(\operatorname{metric\text{-}group\text{-}ms}(g))))).\; (\lambda x \in \operatorname{pts}(\operatorname{metric\text{-}group\text{-}ms}(g)).\; \operatorname{hqh}(\operatorname{opr}(g)(a, x))))(h) \\ \quad \quad =\; (\lambda x \in \operatorname{pts}(\operatorname{metric\text{-}group\text{-}ms}(g)).\; h(\operatorname{opr}(g)(a, x))) \end{array}}}
\vnbprop{mg-t-in-gs}{\vnbstmt{\begin{array}{@{}l@{}} \forall g, f, a. \\ \quad \operatorname{is\text{-}metric\text{-}group}(g) \wedge \\ \quad \operatorname{is\text{-}compact}(\operatorname{metric\text{-}group\text{-}ms}(g)) \wedge \\ \quad f \in \operatorname{cont\text{-}fns}(\operatorname{metric\text{-}group\text{-}ms}(g)) \wedge \\ \quad a \in \operatorname{pts}(\operatorname{metric\text{-}group\text{-}ms}(g)) \Rightarrow \\ \quad \quad (\lambda h \in \operatorname{closure}(\operatorname{nvs\text{-}metric\text{-}space}(C({\operatorname{metric\text{-}group\text{-}ms}(g)})), \operatorname{conv\text{-}hull}(C({\operatorname{metric\text{-}group\text{-}ms}(g)}), \operatorname{image}((\lambda \mathit{hqa} \in \operatorname{pts}(\operatorname{metric\text{-}group\text{-}ms}(g)).\; (\lambda x \in \operatorname{pts}(\operatorname{metric\text{-}group\text{-}ms}(g)).\; f(\operatorname{opr}(g)(\mathit{hqa}, x)))), \operatorname{pts}(\operatorname{metric\text{-}group\text{-}ms}(g))))).\; \\ \quad \quad \quad (\lambda x \in \operatorname{pts}(\operatorname{metric\text{-}group\text{-}ms}(g)).\; h(\operatorname{opr}(g)(a, x)))) \\ \quad \quad \in\; \operatorname{image}((\lambda \mathit{hqg} \in \operatorname{pts}(\operatorname{metric\text{-}group\text{-}ms}(g)).\; \\ \quad \quad \quad \quad (\lambda h \in \operatorname{closure}(\operatorname{nvs\text{-}metric\text{-}space}(C({\operatorname{metric\text{-}group\text{-}ms}(g)})), \operatorname{conv\text{-}hull}(C({\operatorname{metric\text{-}group\text{-}ms}(g)}), \operatorname{image}((\lambda \mathit{hqa} \in \operatorname{pts}(\operatorname{metric\text{-}group\text{-}ms}(g)).\; (\lambda x \in \operatorname{pts}(\operatorname{metric\text{-}group\text{-}ms}(g)).\; f(\operatorname{opr}(g)(\mathit{hqa}, x)))), \operatorname{pts}(\operatorname{metric\text{-}group\text{-}ms}(g))))).\; \\ \quad \quad \quad \quad \quad (\lambda x \in \operatorname{pts}(\operatorname{metric\text{-}group\text{-}ms}(g)).\; h(\operatorname{opr}(g)(\mathit{hqg}, x))))), \\ \quad \quad \quad \operatorname{pts}(\operatorname{metric\text{-}group\text{-}ms}(g))) \end{array}}}
\vnbprop{mg-gs-mem}{\vnbstmt{\begin{array}{@{}l@{}} \forall g, f, a. \\ \quad \operatorname{is\text{-}metric\text{-}group}(g) \wedge \\ \quad \operatorname{is\text{-}compact}(\operatorname{metric\text{-}group\text{-}ms}(g)) \wedge \\ \quad f \in \operatorname{cont\text{-}fns}(\operatorname{metric\text{-}group\text{-}ms}(g)) \wedge \\ \quad a \\ \quad \in\; \operatorname{image}((\lambda \mathit{hqg} \in \operatorname{pts}(\operatorname{metric\text{-}group\text{-}ms}(g)).\; \\ \quad \quad \quad (\lambda h \in \operatorname{closure}(\operatorname{nvs\text{-}metric\text{-}space}(C({\operatorname{metric\text{-}group\text{-}ms}(g)})), \operatorname{conv\text{-}hull}(C({\operatorname{metric\text{-}group\text{-}ms}(g)}), \operatorname{image}((\lambda \mathit{hqa} \in \operatorname{pts}(\operatorname{metric\text{-}group\text{-}ms}(g)).\; (\lambda x \in \operatorname{pts}(\operatorname{metric\text{-}group\text{-}ms}(g)).\; f(\operatorname{opr}(g)(\mathit{hqa}, x)))), \operatorname{pts}(\operatorname{metric\text{-}group\text{-}ms}(g))))).\; \\ \quad \quad \quad \quad (\lambda x \in \operatorname{pts}(\operatorname{metric\text{-}group\text{-}ms}(g)).\; h(\operatorname{opr}(g)(\mathit{hqg}, x))))), \\ \quad \quad \operatorname{pts}(\operatorname{metric\text{-}group\text{-}ms}(g))) \Rightarrow \\ \quad \quad \exists b \in \operatorname{pts}(\operatorname{metric\text{-}group\text{-}ms}(g)). \\ \quad \quad \quad a \\ \quad \quad \quad =\; (\lambda h \in \operatorname{closure}(\operatorname{nvs\text{-}metric\text{-}space}(C({\operatorname{metric\text{-}group\text{-}ms}(g)})), \operatorname{conv\text{-}hull}(C({\operatorname{metric\text{-}group\text{-}ms}(g)}), \operatorname{image}((\lambda \mathit{hqa} \in \operatorname{pts}(\operatorname{metric\text{-}group\text{-}ms}(g)).\; (\lambda x \in \operatorname{pts}(\operatorname{metric\text{-}group\text{-}ms}(g)).\; f(\operatorname{opr}(g)(\mathit{hqa}, x)))), \operatorname{pts}(\operatorname{metric\text{-}group\text{-}ms}(g))))).\; \\ \quad \quad \quad \quad (\lambda x \in \operatorname{pts}(\operatorname{metric\text{-}group\text{-}ms}(g)).\; h(\operatorname{opr}(g)(b, x)))) \end{array}}}
\vnbprop{mg-translation-group}{\vnbstmt{\begin{array}{@{}l@{}} \forall g, f. \\ \quad (\operatorname{is\text{-}metric\text{-}group}(g) \wedge (\operatorname{is\text{-}compact}(\operatorname{metric\text{-}group\text{-}ms}(g)) \wedge (f \in \operatorname{cont\text{-}fns}(\operatorname{metric\text{-}group\text{-}ms}(g))))) \Rightarrow  \\ \quad \quad \operatorname{is\text{-}affine\text{-}isometry\text{-}group}(C({\operatorname{metric\text{-}group\text{-}ms}(g)}), \\ \quad \quad \quad \operatorname{closure}(\operatorname{nvs\text{-}metric\text{-}space}(C({\operatorname{metric\text{-}group\text{-}ms}(g)})), \\ \quad \quad \quad \quad \operatorname{conv\text{-}hull}(C({\operatorname{metric\text{-}group\text{-}ms}(g)}), \\ \quad \quad \quad \quad \quad \operatorname{image}((\lambda a \in \operatorname{pts}(\operatorname{metric\text{-}group\text{-}ms}(g)).\; (\lambda x \in \operatorname{pts}(\operatorname{metric\text{-}group\text{-}ms}(g)).\; f(\operatorname{opr}(g)(a, x)))), \\ \quad \quad \quad \quad \quad \quad \operatorname{pts}(\operatorname{metric\text{-}group\text{-}ms}(g))))), \\ \quad \quad \quad \operatorname{image}((\lambda \mathit{hqg} \in \operatorname{pts}(\operatorname{metric\text{-}group\text{-}ms}(g)).\; \\ \quad \quad \quad \quad \quad (\lambda h \in \operatorname{closure}(\operatorname{nvs\text{-}metric\text{-}space}(C({\operatorname{metric\text{-}group\text{-}ms}(g)})), \operatorname{conv\text{-}hull}(C({\operatorname{metric\text{-}group\text{-}ms}(g)}), \operatorname{image}((\lambda a \in \operatorname{pts}(\operatorname{metric\text{-}group\text{-}ms}(g)).\; (\lambda x \in \operatorname{pts}(\operatorname{metric\text{-}group\text{-}ms}(g)).\; f(\operatorname{opr}(g)(a, x)))), \operatorname{pts}(\operatorname{metric\text{-}group\text{-}ms}(g))))).\; \\ \quad \quad \quad \quad \quad \quad (\lambda x \in \operatorname{pts}(\operatorname{metric\text{-}group\text{-}ms}(g)).\; h(\operatorname{opr}(g)(\mathit{hqg}, x))))), \\ \quad \quad \quad \quad \operatorname{pts}(\operatorname{metric\text{-}group\text{-}ms}(g)))) \end{array}}}
\vnbprop{mg-hull-fixed-point}{\vnbstmt{\begin{array}{@{}l@{}} \forall g, f. \\ \quad (\operatorname{is\text{-}metric\text{-}group}(g) \wedge (\operatorname{is\text{-}compact}(\operatorname{metric\text{-}group\text{-}ms}(g)) \wedge (f \in \operatorname{cont\text{-}fns}(\operatorname{metric\text{-}group\text{-}ms}(g))))) \Rightarrow  \\ \quad \quad \exists w \in \operatorname{closure}(\operatorname{nvs\text{-}metric\text{-}space}(C({\operatorname{metric\text{-}group\text{-}ms}(g)})), \operatorname{conv\text{-}hull}(C({\operatorname{metric\text{-}group\text{-}ms}(g)}), \operatorname{image}((\lambda \mathit{hqa} \in \operatorname{pts}(\operatorname{metric\text{-}group\text{-}ms}(g)).\; (\lambda x \in \operatorname{pts}(\operatorname{metric\text{-}group\text{-}ms}(g)).\; f(\operatorname{opr}(g)(\mathit{hqa}, x)))), \operatorname{pts}(\operatorname{metric\text{-}group\text{-}ms}(g))))). \forall a \in \operatorname{pts}(\operatorname{metric\text{-}group\text{-}ms}(g)). \\ \quad \quad \quad (\lambda x \in \operatorname{pts}(\operatorname{metric\text{-}group\text{-}ms}(g)).\; w(\operatorname{opr}(g)(a, x))) = w \end{array}}}
\vnbprop{mg-hull-has-constant}{\vnbstmt{\begin{array}{@{}l@{}} \forall g, f. \\ \quad (\operatorname{is\text{-}metric\text{-}group}(g) \wedge (\operatorname{is\text{-}compact}(\operatorname{metric\text{-}group\text{-}ms}(g)) \wedge (f \in \operatorname{cont\text{-}fns}(\operatorname{metric\text{-}group\text{-}ms}(g))))) \Rightarrow  \\ \quad \quad \exists c \in \mathbb{R}. \\ \quad \quad \quad (\lambda x \in \operatorname{pts}(\operatorname{metric\text{-}group\text{-}ms}(g)).\; c) \\ \quad \quad \quad \in\; \operatorname{closure}(\operatorname{nvs\text{-}metric\text{-}space}(C({\operatorname{metric\text{-}group\text{-}ms}(g)})), \\ \quad \quad \quad \quad \operatorname{conv\text{-}hull}(C({\operatorname{metric\text{-}group\text{-}ms}(g)}), \\ \quad \quad \quad \quad \quad \operatorname{image}((\lambda a \in \operatorname{pts}(\operatorname{metric\text{-}group\text{-}ms}(g)).\; (\lambda x \in \operatorname{pts}(\operatorname{metric\text{-}group\text{-}ms}(g)).\; f(\operatorname{opr}(g)(a, x)))), \\ \quad \quad \quad \quad \quad \quad \operatorname{pts}(\operatorname{metric\text{-}group\text{-}ms}(g))))) \end{array}}}

\vnbbookfile{theorem-library/haar-c3}{theorem-library/haar-c3}
\noindent{\small haar-c3.scm -- wave of 2026-11-09/10 (agent HAAR-C3), wired 2026-11-10 from scratchpad/haarc3-1109/haar-c3.scm. Library: analysis.}\par
\vnbprop{opp-group-carr}{\vnbstmt{\forall g.\; (\operatorname{is\text{-}group}(g) \Rightarrow (\operatorname{carr}({g}^{\mathrm{op}}) \simeq \operatorname{carr}(g)))}}
\vnbprop{opp-group-opr}{\vnbstmt{\begin{array}{@{}l@{}} \forall g. \\ \quad \operatorname{is\text{-}group}(g) \Rightarrow  \\ \quad \quad (\operatorname{opr}({g}^{\mathrm{op}}) \simeq (\lambda {[}\operatorname{ogp\_}, \operatorname{ogq\_}] \in (\operatorname{carr}(g) \times \operatorname{carr}(g)).\; \operatorname{opr}(g)(\operatorname{ogq\_}, \operatorname{ogp\_}))) \end{array}}}
\vnbprop{opp-group-iden}{\vnbstmt{\forall g.\; (\operatorname{is\text{-}group}(g) \Rightarrow (\operatorname{iden}({g}^{\mathrm{op}}) \simeq \operatorname{iden}(g)))}}
\vnbprop{opp-group-inv}{\vnbstmt{\forall g.\; (\operatorname{is\text{-}group}(g) \Rightarrow (\operatorname{inv}({g}^{\mathrm{op}}) \simeq \operatorname{inv}(g)))}}
\vnbprop{opp-group-opr-apply}{\vnbstmt{\begin{array}{@{}l@{}} \forall g, a, b. \\ \quad (\operatorname{is\text{-}group}(g) \wedge ((a \in \operatorname{carr}(g)) \wedge (b \in \operatorname{carr}(g)))) \Rightarrow  \\ \quad \quad \operatorname{opr}({g}^{\mathrm{op}})(a, b) = \operatorname{opr}(g)(b, a) \end{array}}}
\vnbprop{opp-group-is-group}{\vnbstmt{\forall g.\; (\operatorname{is\text{-}group}(g) \Rightarrow \operatorname{is\text{-}group}({g}^{\mathrm{op}}))}}
\vnbprop{opp-group-metric}{\vnbstmt{\begin{array}{@{}l@{}} \forall g. \\ \quad \operatorname{is\text{-}metric\text{-}group}(g) \Rightarrow  \\ \quad \quad \operatorname{is\text{-}group\text{-}metric}(\operatorname{metric\text{-}group\text{-}ms}(g), {\operatorname{metric\text{-}group\text{-}group}(g)}^{\mathrm{op}}) \end{array}}}
\vnbprop{opp-mg-is-metric-group}{\vnbstmt{\forall g.\; (\operatorname{is\text{-}metric\text{-}group}(g) \Rightarrow \operatorname{is\text{-}metric\text{-}group}({g}^{\mathrm{op}}))}}
\vnbprop{opp-mg-carr}{\vnbstmt{\forall g.\; (\operatorname{is\text{-}metric\text{-}group}(g) \Rightarrow (\operatorname{carr}({g}^{\mathrm{op}}) \simeq \operatorname{carr}(g)))}}
\vnbprop{opp-mg-gdist}{\vnbstmt{\forall g.\; (\operatorname{is\text{-}metric\text{-}group}(g) \Rightarrow (\operatorname{gdist}({g}^{\mathrm{op}}) \simeq \operatorname{gdist}(g)))}}
\vnbprop{opp-mg-opr}{\vnbstmt{\begin{array}{@{}l@{}} \forall g. \\ \quad \operatorname{is\text{-}metric\text{-}group}(g) \Rightarrow  \\ \quad \quad (\operatorname{opr}({g}^{\mathrm{op}}) \simeq \operatorname{opr}({\operatorname{metric\text{-}group\text{-}group}(g)}^{\mathrm{op}})) \end{array}}}
\vnbprop{opp-mg-ms}{\vnbstmt{\begin{array}{@{}l@{}} \forall g. \\ \quad \operatorname{is\text{-}metric\text{-}group}(g) \Rightarrow  \\ \quad \quad (\operatorname{metric\text{-}group\text{-}ms}({g}^{\mathrm{op}}) \simeq \operatorname{metric\text{-}group\text{-}ms}(g)) \end{array}}}
\vnbprop{opp-mg-compact}{\vnbstmt{\begin{array}{@{}l@{}} \forall g. \\ \quad (\operatorname{is\text{-}metric\text{-}group}(g) \wedge \operatorname{is\text{-}compact}(\operatorname{metric\text{-}group\text{-}ms}(g))) \Rightarrow  \\ \quad \quad \operatorname{is\text{-}compact}(\operatorname{metric\text{-}group\text{-}ms}({g}^{\mathrm{op}})) \end{array}}}
\vnbprop{opp-mg-opr-apply}{\vnbstmt{\begin{array}{@{}l@{}} \forall g, a, b. \\ \quad (\operatorname{is\text{-}metric\text{-}group}(g) \wedge ((a \in \operatorname{pts}(\operatorname{metric\text{-}group\text{-}ms}(g))) \wedge (b \in \operatorname{pts}(\operatorname{metric\text{-}group\text{-}ms}(g))))) \Rightarrow  \\ \quad \quad \operatorname{opr}({g}^{\mathrm{op}})(a, b) = \operatorname{opr}(g)(b, a) \end{array}}}
\vnbprop{opp-mg-ltrans-eq}{\vnbstmt{\begin{array}{@{}l@{}} \forall g, b, h. \\ \quad (\operatorname{is\text{-}metric\text{-}group}(g) \wedge ((b \in \operatorname{pts}(\operatorname{metric\text{-}group\text{-}ms}(g))) \wedge (h \in (\operatorname{pts}(\operatorname{metric\text{-}group\text{-}ms}(g)) \to \mathbb{R})))) \Rightarrow  \\ \quad \quad (\lambda x \in \operatorname{pts}(\operatorname{metric\text{-}group\text{-}ms}(g)).\; h(\operatorname{opr}({g}^{\mathrm{op}})(b, x))) \\ \quad \quad =\; (\lambda x \in \operatorname{pts}(\operatorname{metric\text{-}group\text{-}ms}(g)).\; h(\operatorname{opr}(g)(x, b))) \end{array}}}
\vnbprop{mg-rtrans-cont}{\vnbstmt{\begin{array}{@{}l@{}} \forall g, b, h. \\ \quad (\operatorname{is\text{-}metric\text{-}group}(g) \wedge ((b \in \operatorname{pts}(\operatorname{metric\text{-}group\text{-}ms}(g))) \wedge (h \in \operatorname{cont\text{-}fns}(\operatorname{metric\text{-}group\text{-}ms}(g))))) \Rightarrow  \\ \quad \quad (\lambda x \in \operatorname{pts}(\operatorname{metric\text{-}group\text{-}ms}(g)).\; h(\operatorname{opr}(g)(x, b))) \in \operatorname{cont\text{-}fns}(\operatorname{metric\text{-}group\text{-}ms}(g)) \end{array}}}
\vnbprop{mg-rtrans-dist}{\vnbstmt{\begin{array}{@{}l@{}} \forall g, b, f, a. \\ \quad \operatorname{is\text{-}metric\text{-}group}(g) \wedge \\ \quad \operatorname{is\text{-}compact}(\operatorname{metric\text{-}group\text{-}ms}(g)) \wedge \\ \quad b \in \operatorname{pts}(\operatorname{metric\text{-}group\text{-}ms}(g)) \wedge \\ \quad f \in \operatorname{cont\text{-}fns}(\operatorname{metric\text{-}group\text{-}ms}(g)) \wedge \\ \quad a \in \operatorname{cont\text{-}fns}(\operatorname{metric\text{-}group\text{-}ms}(g)) \Rightarrow \\ \quad \quad \operatorname{dist}(\operatorname{nvs\text{-}metric\text{-}space}(C({\operatorname{metric\text{-}group\text{-}ms}(g)})))((\lambda x \in \operatorname{pts}(\operatorname{metric\text{-}group\text{-}ms}(g)).\; f(\operatorname{opr}(g)(x, b))), \\ \quad \quad \quad (\lambda x \in \operatorname{pts}(\operatorname{metric\text{-}group\text{-}ms}(g)).\; a(\operatorname{opr}(g)(x, b)))) \\ \quad \quad =\; \operatorname{dist}(\operatorname{nvs\text{-}metric\text{-}space}(C({\operatorname{metric\text{-}group\text{-}ms}(g)})))(f, a) \end{array}}}
\vnbprop{mg-rtrans-affine}{\vnbstmt{\begin{array}{@{}l@{}} \forall g, b, u, v, t. \\ \quad \operatorname{is\text{-}metric\text{-}group}(g) \wedge \\ \quad \operatorname{is\text{-}compact}(\operatorname{metric\text{-}group\text{-}ms}(g)) \wedge \\ \quad b \in \operatorname{pts}(\operatorname{metric\text{-}group\text{-}ms}(g)) \wedge \\ \quad u \in \operatorname{cont\text{-}fns}(\operatorname{metric\text{-}group\text{-}ms}(g)) \wedge \\ \quad v \in \operatorname{cont\text{-}fns}(\operatorname{metric\text{-}group\text{-}ms}(g)) \wedge \\ \quad t \in \mathbb{R} \Rightarrow \\ \quad \quad (\lambda x \in \operatorname{pts}(\operatorname{metric\text{-}group\text{-}ms}(g)).\; \\ \quad \quad \quad \operatorname{vadd}(C({\operatorname{metric\text{-}group\text{-}ms}(g)}))(\operatorname{act}(C({\operatorname{metric\text{-}group\text{-}ms}(g)}))((1 - t), u), \operatorname{act}(C({\operatorname{metric\text{-}group\text{-}ms}(g)}))(t, v))(\operatorname{opr}(g)(x, b))) \\ \quad \quad =\; \operatorname{vadd}(C({\operatorname{metric\text{-}group\text{-}ms}(g)}))(\operatorname{act}(C({\operatorname{metric\text{-}group\text{-}ms}(g)}))((1 - t), (\lambda x \in \operatorname{pts}(\operatorname{metric\text{-}group\text{-}ms}(g)).\; u(\operatorname{opr}(g)(x, b)))), \\ \quad \quad \quad \operatorname{act}(C({\operatorname{metric\text{-}group\text{-}ms}(g)}))(t, (\lambda x \in \operatorname{pts}(\operatorname{metric\text{-}group\text{-}ms}(g)).\; v(\operatorname{opr}(g)(x, b))))) \end{array}}}
\vnbprop{mg-ltrans-add}{\vnbstmt{\begin{array}{@{}l@{}} \forall g, a, u, v. \\ \quad \operatorname{is\text{-}metric\text{-}group}(g) \wedge \\ \quad a \in \operatorname{pts}(\operatorname{metric\text{-}group\text{-}ms}(g)) \wedge \\ \quad u \in \operatorname{cont\text{-}fns}(\operatorname{metric\text{-}group\text{-}ms}(g)) \wedge \\ \quad v \in \operatorname{cont\text{-}fns}(\operatorname{metric\text{-}group\text{-}ms}(g)) \Rightarrow \\ \quad \quad (\lambda x \in \operatorname{pts}(\operatorname{metric\text{-}group\text{-}ms}(g)).\; \operatorname{vadd}(C({\operatorname{metric\text{-}group\text{-}ms}(g)}))(u, v)(\operatorname{opr}(g)(a, x))) \\ \quad \quad =\; \operatorname{vadd}(C({\operatorname{metric\text{-}group\text{-}ms}(g)}))((\lambda x \in \operatorname{pts}(\operatorname{metric\text{-}group\text{-}ms}(g)).\; u(\operatorname{opr}(g)(a, x))), \\ \quad \quad \quad (\lambda x \in \operatorname{pts}(\operatorname{metric\text{-}group\text{-}ms}(g)).\; v(\operatorname{opr}(g)(a, x)))) \end{array}}}
\vnbprop{mg-ltrans-act}{\vnbstmt{\begin{array}{@{}l@{}} \forall g, a, r, u. \\ \quad \operatorname{is\text{-}metric\text{-}group}(g) \wedge \\ \quad a \in \operatorname{pts}(\operatorname{metric\text{-}group\text{-}ms}(g)) \wedge \\ \quad r \in \mathbb{R} \wedge \\ \quad u \in \operatorname{cont\text{-}fns}(\operatorname{metric\text{-}group\text{-}ms}(g)) \Rightarrow \\ \quad \quad (\lambda x \in \operatorname{pts}(\operatorname{metric\text{-}group\text{-}ms}(g)).\; \operatorname{act}(C({\operatorname{metric\text{-}group\text{-}ms}(g)}))(r, u)(\operatorname{opr}(g)(a, x))) \\ \quad \quad =\; \operatorname{act}(C({\operatorname{metric\text{-}group\text{-}ms}(g)}))(r, (\lambda x \in \operatorname{pts}(\operatorname{metric\text{-}group\text{-}ms}(g)).\; u(\operatorname{opr}(g)(a, x)))) \end{array}}}
\vnbprop{mg-ltrans-const}{\vnbstmt{\begin{array}{@{}l@{}} \forall g, a, c. \\ \quad (\operatorname{is\text{-}metric\text{-}group}(g) \wedge ((a \in \operatorname{pts}(\operatorname{metric\text{-}group\text{-}ms}(g))) \wedge (c \in \mathbb{R}))) \Rightarrow  \\ \quad \quad (\lambda x \in \operatorname{pts}(\operatorname{metric\text{-}group\text{-}ms}(g)).\; (\lambda x \in \operatorname{pts}(\operatorname{metric\text{-}group\text{-}ms}(g)).\; c)(\operatorname{opr}(g)(a, x))) \\ \quad \quad =\; (\lambda x \in \operatorname{pts}(\operatorname{metric\text{-}group\text{-}ms}(g)).\; c) \end{array}}}
\vnbprop{mg-rtrans-const}{\vnbstmt{\begin{array}{@{}l@{}} \forall g, b, c. \\ \quad (\operatorname{is\text{-}metric\text{-}group}(g) \wedge ((b \in \operatorname{pts}(\operatorname{metric\text{-}group\text{-}ms}(g))) \wedge (c \in \mathbb{R}))) \Rightarrow  \\ \quad \quad (\lambda x \in \operatorname{pts}(\operatorname{metric\text{-}group\text{-}ms}(g)).\; (\lambda x \in \operatorname{pts}(\operatorname{metric\text{-}group\text{-}ms}(g)).\; c)(\operatorname{opr}(g)(x, b))) \\ \quad \quad =\; (\lambda x \in \operatorname{pts}(\operatorname{metric\text{-}group\text{-}ms}(g)).\; c) \end{array}}}
\vnbprop{mg-lr-commute}{\vnbstmt{\begin{array}{@{}l@{}} \forall g, a, b, h. \\ \quad \operatorname{is\text{-}metric\text{-}group}(g) \wedge \\ \quad a \in \operatorname{pts}(\operatorname{metric\text{-}group\text{-}ms}(g)) \wedge \\ \quad b \in \operatorname{pts}(\operatorname{metric\text{-}group\text{-}ms}(g)) \wedge \\ \quad h \in (\operatorname{pts}(\operatorname{metric\text{-}group\text{-}ms}(g)) \to \mathbb{R}) \Rightarrow \\ \quad \quad (\lambda x \in \operatorname{pts}(\operatorname{metric\text{-}group\text{-}ms}(g)).\; \\ \quad \quad \quad (\lambda x \in \operatorname{pts}(\operatorname{metric\text{-}group\text{-}ms}(g)).\; h(\operatorname{opr}(g)(x, b)))(\operatorname{opr}(g)(a, x))) \\ \quad \quad =\; (\lambda x \in \operatorname{pts}(\operatorname{metric\text{-}group\text{-}ms}(g)).\; \\ \quad \quad \quad (\lambda x \in \operatorname{pts}(\operatorname{metric\text{-}group\text{-}ms}(g)).\; h(\operatorname{opr}(g)(a, x)))(\operatorname{opr}(g)(x, b))) \end{array}}}
\vnbprop{haar-avg-unfold}{\vnbstmt{\begin{array}{@{}l@{}} \forall g. \\ \quad (\operatorname{haar\text{-}averages}(g) \simeq \{q \in (\operatorname{cont\text{-}fns}(\operatorname{metric\text{-}group\text{-}ms}(g)) \to \operatorname{cont\text{-}fns}(\operatorname{metric\text{-}group\text{-}ms}(g))) : (\forall u, v \in \operatorname{cont\text{-}fns}(\operatorname{metric\text{-}group\text{-}ms}(g)).\; (q(\operatorname{vadd}(C({\operatorname{metric\text{-}group\text{-}ms}(g)}))(u, v)) = \operatorname{vadd}(C({\operatorname{metric\text{-}group\text{-}ms}(g)}))(q(u), q(v))) \wedge (\forall r \in \mathbb{R}, u \in \operatorname{cont\text{-}fns}(\operatorname{metric\text{-}group\text{-}ms}(g)).\; (q(\operatorname{act}(C({\operatorname{metric\text{-}group\text{-}ms}(g)}))(r, u)) = \operatorname{act}(C({\operatorname{metric\text{-}group\text{-}ms}(g)}))(r, q(u))) \wedge (\forall u, v \in \operatorname{cont\text{-}fns}(\operatorname{metric\text{-}group\text{-}ms}(g)).\; (\operatorname{dist}(\operatorname{nvs\text{-}metric\text{-}space}(C({\operatorname{metric\text{-}group\text{-}ms}(g)})))(q(u), q(v)) \leq \operatorname{dist}(\operatorname{nvs\text{-}metric\text{-}space}(C({\operatorname{metric\text{-}group\text{-}ms}(g)})))(u, v)) \wedge (\forall c \in \mathbb{R}.\; (q((\lambda x \in \operatorname{pts}(\operatorname{metric\text{-}group\text{-}ms}(g)).\; c)) = (\lambda x \in \operatorname{pts}(\operatorname{metric\text{-}group\text{-}ms}(g)).\; c)) \wedge (\forall u \in \operatorname{cont\text{-}fns}(\operatorname{metric\text{-}group\text{-}ms}(g)).\; (q(u) \in \operatorname{conv\text{-}hull}(C({\operatorname{metric\text{-}group\text{-}ms}(g)}), \operatorname{image}((\lambda a \in \operatorname{pts}(\operatorname{metric\text{-}group\text{-}ms}(g)).\; (\lambda x \in \operatorname{pts}(\operatorname{metric\text{-}group\text{-}ms}(g)).\; u(\operatorname{opr}(g)(a, x)))), \operatorname{pts}(\operatorname{metric\text{-}group\text{-}ms}(g))))) \wedge \forall b \in \operatorname{pts}(\operatorname{metric\text{-}group\text{-}ms}(g)), u \in \operatorname{cont\text{-}fns}(\operatorname{metric\text{-}group\text{-}ms}(g)).\; (q((\lambda x \in \operatorname{pts}(\operatorname{metric\text{-}group\text{-}ms}(g)).\; u(\operatorname{opr}(g)(x, b)))) = (\lambda x \in \operatorname{pts}(\operatorname{metric\text{-}group\text{-}ms}(g)).\; q(u)(\operatorname{opr}(g)(x, b)))))))))\}) \end{array}}}
\vnbprop{haar-avg-elim}{\vnbstmt{\begin{array}{@{}l@{}} \forall g, q. \\ \quad (\operatorname{is\text{-}metric\text{-}group}(g) \wedge (q \in \operatorname{haar\text{-}averages}(g))) \Rightarrow  \\ \quad \quad q \in (\operatorname{cont\text{-}fns}(\operatorname{metric\text{-}group\text{-}ms}(g)) \to \operatorname{cont\text{-}fns}(\operatorname{metric\text{-}group\text{-}ms}(g))) \wedge \\ \quad \quad \forall u, v \in \operatorname{cont\text{-}fns}(\operatorname{metric\text{-}group\text{-}ms}(g)). \\ \quad \quad \quad q(\operatorname{vadd}(C({\operatorname{metric\text{-}group\text{-}ms}(g)}))(u, v)) = \operatorname{vadd}(C({\operatorname{metric\text{-}group\text{-}ms}(g)}))(q(u), q(v)) \wedge \\ \quad \quad \forall r \in \mathbb{R}, u \in \operatorname{cont\text{-}fns}(\operatorname{metric\text{-}group\text{-}ms}(g)). \\ \quad \quad \quad q(\operatorname{act}(C({\operatorname{metric\text{-}group\text{-}ms}(g)}))(r, u)) = \operatorname{act}(C({\operatorname{metric\text{-}group\text{-}ms}(g)}))(r, q(u)) \wedge \\ \quad \quad \forall u, v \in \operatorname{cont\text{-}fns}(\operatorname{metric\text{-}group\text{-}ms}(g)). \\ \quad \quad \quad \operatorname{dist}(\operatorname{nvs\text{-}metric\text{-}space}(C({\operatorname{metric\text{-}group\text{-}ms}(g)})))(q(u), q(v)) \\ \quad \quad \quad \leq\; \operatorname{dist}(\operatorname{nvs\text{-}metric\text{-}space}(C({\operatorname{metric\text{-}group\text{-}ms}(g)})))(u, v) \wedge \\ \quad \quad \forall c \in \mathbb{R}. \\ \quad \quad \quad q((\lambda x \in \operatorname{pts}(\operatorname{metric\text{-}group\text{-}ms}(g)).\; c)) = (\lambda x \in \operatorname{pts}(\operatorname{metric\text{-}group\text{-}ms}(g)).\; c) \wedge \\ \quad \quad \forall u \in \operatorname{cont\text{-}fns}(\operatorname{metric\text{-}group\text{-}ms}(g)). \\ \quad \quad \quad q(u) \\ \quad \quad \quad \in\; \operatorname{conv\text{-}hull}(C({\operatorname{metric\text{-}group\text{-}ms}(g)}), \\ \quad \quad \quad \quad \operatorname{image}((\lambda a \in \operatorname{pts}(\operatorname{metric\text{-}group\text{-}ms}(g)).\; (\lambda x \in \operatorname{pts}(\operatorname{metric\text{-}group\text{-}ms}(g)).\; u(\operatorname{opr}(g)(a, x)))), \\ \quad \quad \quad \quad \quad \operatorname{pts}(\operatorname{metric\text{-}group\text{-}ms}(g)))) \wedge \\ \quad \quad \forall b \in \operatorname{pts}(\operatorname{metric\text{-}group\text{-}ms}(g)), u \in \operatorname{cont\text{-}fns}(\operatorname{metric\text{-}group\text{-}ms}(g)). \\ \quad \quad \quad q((\lambda x \in \operatorname{pts}(\operatorname{metric\text{-}group\text{-}ms}(g)).\; u(\operatorname{opr}(g)(x, b)))) \\ \quad \quad \quad =\; (\lambda x \in \operatorname{pts}(\operatorname{metric\text{-}group\text{-}ms}(g)).\; q(u)(\operatorname{opr}(g)(x, b))) \end{array}}}
\vnbprop{haar-avg-intro}{\vnbstmt{\begin{array}{@{}l@{}} \forall g, q. \\ \quad \operatorname{is\text{-}metric\text{-}group}(g) \wedge \\ \quad q \in (\operatorname{cont\text{-}fns}(\operatorname{metric\text{-}group\text{-}ms}(g)) \to \operatorname{cont\text{-}fns}(\operatorname{metric\text{-}group\text{-}ms}(g))) \wedge \\ \quad \forall u, v \in \operatorname{cont\text{-}fns}(\operatorname{metric\text{-}group\text{-}ms}(g)). \\ \quad \quad q(\operatorname{vadd}(C({\operatorname{metric\text{-}group\text{-}ms}(g)}))(u, v)) = \operatorname{vadd}(C({\operatorname{metric\text{-}group\text{-}ms}(g)}))(q(u), q(v)) \wedge \\ \quad \forall r \in \mathbb{R}, u \in \operatorname{cont\text{-}fns}(\operatorname{metric\text{-}group\text{-}ms}(g)). \\ \quad \quad q(\operatorname{act}(C({\operatorname{metric\text{-}group\text{-}ms}(g)}))(r, u)) = \operatorname{act}(C({\operatorname{metric\text{-}group\text{-}ms}(g)}))(r, q(u)) \wedge \\ \quad \forall u, v \in \operatorname{cont\text{-}fns}(\operatorname{metric\text{-}group\text{-}ms}(g)). \\ \quad \quad \operatorname{dist}(\operatorname{nvs\text{-}metric\text{-}space}(C({\operatorname{metric\text{-}group\text{-}ms}(g)})))(q(u), q(v)) \\ \quad \quad \leq\; \operatorname{dist}(\operatorname{nvs\text{-}metric\text{-}space}(C({\operatorname{metric\text{-}group\text{-}ms}(g)})))(u, v) \wedge \\ \quad \forall c \in \mathbb{R}. \\ \quad \quad q((\lambda x \in \operatorname{pts}(\operatorname{metric\text{-}group\text{-}ms}(g)).\; c)) = (\lambda x \in \operatorname{pts}(\operatorname{metric\text{-}group\text{-}ms}(g)).\; c) \wedge \\ \quad \forall u \in \operatorname{cont\text{-}fns}(\operatorname{metric\text{-}group\text{-}ms}(g)). \\ \quad \quad q(u) \\ \quad \quad \in\; \operatorname{conv\text{-}hull}(C({\operatorname{metric\text{-}group\text{-}ms}(g)}), \\ \quad \quad \quad \operatorname{image}((\lambda a \in \operatorname{pts}(\operatorname{metric\text{-}group\text{-}ms}(g)).\; (\lambda x \in \operatorname{pts}(\operatorname{metric\text{-}group\text{-}ms}(g)).\; u(\operatorname{opr}(g)(a, x)))), \\ \quad \quad \quad \quad \operatorname{pts}(\operatorname{metric\text{-}group\text{-}ms}(g)))) \wedge \\ \quad \forall b \in \operatorname{pts}(\operatorname{metric\text{-}group\text{-}ms}(g)), u \in \operatorname{cont\text{-}fns}(\operatorname{metric\text{-}group\text{-}ms}(g)). \\ \quad \quad q((\lambda x \in \operatorname{pts}(\operatorname{metric\text{-}group\text{-}ms}(g)).\; u(\operatorname{opr}(g)(x, b)))) \\ \quad \quad =\; (\lambda x \in \operatorname{pts}(\operatorname{metric\text{-}group\text{-}ms}(g)).\; q(u)(\operatorname{opr}(g)(x, b))) \Rightarrow \\ \quad \quad q \in \operatorname{haar\text{-}averages}(g) \end{array}}}
\vnbprop{ck-const-cont}{\vnbstmt{\begin{array}{@{}l@{}} \forall s.; \forall c \in \mathbb{R}. \\ \quad \operatorname{is\text{-}metric\text{-}space}(s) \Rightarrow ((\lambda x \in \operatorname{pts}(s).\; c) \in \operatorname{cont\text{-}fns}(s)) \end{array}}}
\vnbprop{haar-ltrans-avg}{\vnbstmt{\begin{array}{@{}l@{}} \forall g, a. \\ \quad (\operatorname{is\text{-}metric\text{-}group}(g) \wedge (\operatorname{is\text{-}compact}(\operatorname{metric\text{-}group\text{-}ms}(g)) \wedge (a \in \operatorname{pts}(\operatorname{metric\text{-}group\text{-}ms}(g))))) \Rightarrow  \\ \quad \quad (\lambda h \in \operatorname{cont\text{-}fns}(\operatorname{metric\text{-}group\text{-}ms}(g)).\; \\ \quad \quad \quad (\lambda x \in \operatorname{pts}(\operatorname{metric\text{-}group\text{-}ms}(g)).\; h(\operatorname{opr}(g)(a, x)))) \\ \quad \quad \in\; \operatorname{haar\text{-}averages}(g) \end{array}}}
\vnbprop{nvs-dist-comb-le}{\vnbstmt{\begin{array}{@{}l@{}} \forall m, a, b, c, d, t. \\ \quad \operatorname{is\text{-}normed\text{-}vector\text{-}space}(m) \wedge \\ \quad a \in \operatorname{vec}(m) \wedge \\ \quad b \in \operatorname{vec}(m) \wedge \\ \quad c \in \operatorname{vec}(m) \wedge \\ \quad d \in \operatorname{vec}(m) \wedge \\ \quad t \in \mathbb{R} \wedge \\ \quad 0 \leq t \wedge \\ \quad t \leq 1 \Rightarrow \\ \quad \quad \operatorname{dist}(\operatorname{nvs\text{-}metric\text{-}space}(m))(\operatorname{vadd}(m)(\operatorname{act}(m)((1 - t), a), \operatorname{act}(m)(t, b)), \\ \quad \quad \quad \operatorname{vadd}(m)(\operatorname{act}(m)((1 - t), c), \operatorname{act}(m)(t, d))) \\ \quad \quad \leq\; ((1 - t) \cdot \operatorname{dist}(\operatorname{nvs\text{-}metric\text{-}space}(m))(a, c)) + (t \cdot \operatorname{dist}(\operatorname{nvs\text{-}metric\text{-}space}(m))(b, d)) \end{array}}}
\vnbprop{rr-convex-le-bound}{\vnbstmt{\begin{array}{@{}l@{}} \forall t, x, y, d. \\ \quad ((t \in \mathbb{R}) \wedge ((0 \leq t) \wedge ((t \leq 1) \wedge ((x \in \mathbb{R}) \wedge ((y \in \mathbb{R}) \wedge ((d \in \mathbb{R}) \wedge ((x \leq d) \wedge (y \leq d)))))))) \Rightarrow  \\ \quad \quad (((1 - t) \cdot x) + (t \cdot y)) \leq d \end{array}}}
\vnbprop{haar-avg-comb}{\vnbstmt{\begin{array}{@{}l@{}} \forall g, a, b, t. \\ \quad \operatorname{is\text{-}metric\text{-}group}(g) \wedge \\ \quad \operatorname{is\text{-}compact}(\operatorname{metric\text{-}group\text{-}ms}(g)) \wedge \\ \quad a \in \operatorname{haar\text{-}averages}(g) \wedge \\ \quad b \in \operatorname{haar\text{-}averages}(g) \wedge \\ \quad t \in \mathbb{R} \wedge \\ \quad 0 \leq t \wedge \\ \quad t \leq 1 \Rightarrow \\ \quad \quad (\lambda h \in \operatorname{cont\text{-}fns}(\operatorname{metric\text{-}group\text{-}ms}(g)).\; \\ \quad \quad \quad \operatorname{vadd}(C({\operatorname{metric\text{-}group\text{-}ms}(g)}))(\operatorname{act}(C({\operatorname{metric\text{-}group\text{-}ms}(g)}))((1 - t), a(h)), \\ \quad \quad \quad \quad \operatorname{act}(C({\operatorname{metric\text{-}group\text{-}ms}(g)}))(t, b(h)))) \\ \quad \quad \in\; \operatorname{haar\text{-}averages}(g) \end{array}}}
\vnbprop{haar-hull-rep}{\vnbstmt{\begin{array}{@{}l@{}} \forall g, f, u. \\ \quad \operatorname{is\text{-}metric\text{-}group}(g) \wedge \\ \quad \operatorname{is\text{-}compact}(\operatorname{metric\text{-}group\text{-}ms}(g)) \wedge \\ \quad f \in \operatorname{cont\text{-}fns}(\operatorname{metric\text{-}group\text{-}ms}(g)) \wedge \\ \quad u \\ \quad \in\; \operatorname{conv\text{-}hull}(C({\operatorname{metric\text{-}group\text{-}ms}(g)}), \\ \quad \quad \operatorname{image}((\lambda a \in \operatorname{pts}(\operatorname{metric\text{-}group\text{-}ms}(g)).\; (\lambda x \in \operatorname{pts}(\operatorname{metric\text{-}group\text{-}ms}(g)).\; f(\operatorname{opr}(g)(a, x)))), \\ \quad \quad \quad \operatorname{pts}(\operatorname{metric\text{-}group\text{-}ms}(g)))) \Rightarrow \\ \quad \quad \exists q \in \operatorname{haar\text{-}averages}(g).\; (u = q(f)) \end{array}}}
\vnbprop{haar-hull-transport}{\vnbstmt{\begin{array}{@{}l@{}} \forall g, q.; \forall f \in \operatorname{cont\text{-}fns}(\operatorname{metric\text{-}group\text{-}ms}(g)), w \in \operatorname{closure}(\operatorname{nvs\text{-}metric\text{-}space}(C({\operatorname{metric\text{-}group\text{-}ms}(g)})), \operatorname{conv\text{-}hull}(C({\operatorname{metric\text{-}group\text{-}ms}(g)}), \operatorname{image}((\lambda a \in \operatorname{pts}(\operatorname{metric\text{-}group\text{-}ms}(g)).\; (\lambda x \in \operatorname{pts}(\operatorname{metric\text{-}group\text{-}ms}(g)).\; f(\operatorname{opr}(g)(a, x)))), \operatorname{pts}(\operatorname{metric\text{-}group\text{-}ms}(g))))). \\ \quad \operatorname{is\text{-}metric\text{-}group}(g) \wedge \\ \quad \operatorname{is\text{-}compact}(\operatorname{metric\text{-}group\text{-}ms}(g)) \wedge \\ \quad q \in (\operatorname{cont\text{-}fns}(\operatorname{metric\text{-}group\text{-}ms}(g)) \to \operatorname{cont\text{-}fns}(\operatorname{metric\text{-}group\text{-}ms}(g))) \wedge \\ \quad \forall u, \mathit{hrw} \in \operatorname{cont\text{-}fns}(\operatorname{metric\text{-}group\text{-}ms}(g)), t \in \mathbb{R}. \\ \quad \quad q(\operatorname{vadd}(C({\operatorname{metric\text{-}group\text{-}ms}(g)}))(\operatorname{act}(C({\operatorname{metric\text{-}group\text{-}ms}(g)}))((1 - t), u), \\ \quad \quad \quad \quad \operatorname{act}(C({\operatorname{metric\text{-}group\text{-}ms}(g)}))(t, \mathit{hrw}))) \\ \quad \quad =\; \operatorname{vadd}(C({\operatorname{metric\text{-}group\text{-}ms}(g)}))(\operatorname{act}(C({\operatorname{metric\text{-}group\text{-}ms}(g)}))((1 - t), q(u)), \\ \quad \quad \quad \operatorname{act}(C({\operatorname{metric\text{-}group\text{-}ms}(g)}))(t, q(\mathit{hrw}))) \wedge \\ \quad \forall u, \mathit{hrw} \in \operatorname{cont\text{-}fns}(\operatorname{metric\text{-}group\text{-}ms}(g)). \\ \quad \quad \operatorname{dist}(\operatorname{nvs\text{-}metric\text{-}space}(C({\operatorname{metric\text{-}group\text{-}ms}(g)})))(q(u), q(\mathit{hrw})) \\ \quad \quad \leq\; \operatorname{dist}(\operatorname{nvs\text{-}metric\text{-}space}(C({\operatorname{metric\text{-}group\text{-}ms}(g)})))(u, \mathit{hrw}) \wedge \\ \quad \forall b \in \operatorname{pts}(\operatorname{metric\text{-}group\text{-}ms}(g)), u \in \operatorname{cont\text{-}fns}(\operatorname{metric\text{-}group\text{-}ms}(g)). \\ \quad \quad q((\lambda x \in \operatorname{pts}(\operatorname{metric\text{-}group\text{-}ms}(g)).\; u(\operatorname{opr}(g)(b, x)))) \\ \quad \quad =\; (\lambda x \in \operatorname{pts}(\operatorname{metric\text{-}group\text{-}ms}(g)).\; q(u)(\operatorname{opr}(g)(b, x))) \Rightarrow \\ \quad \quad q(w) \\ \quad \quad \in\; \operatorname{closure}(\operatorname{nvs\text{-}metric\text{-}space}(C({\operatorname{metric\text{-}group\text{-}ms}(g)})), \\ \quad \quad \quad \operatorname{conv\text{-}hull}(C({\operatorname{metric\text{-}group\text{-}ms}(g)}), \\ \quad \quad \quad \quad \operatorname{image}((\lambda a \in \operatorname{pts}(\operatorname{metric\text{-}group\text{-}ms}(g)).\; \\ \quad \quad \quad \quad \quad \quad (\lambda x \in \operatorname{pts}(\operatorname{metric\text{-}group\text{-}ms}(g)).\; q(f)(\operatorname{opr}(g)(a, x)))), \\ \quad \quad \quad \quad \quad \operatorname{pts}(\operatorname{metric\text{-}group\text{-}ms}(g))))) \end{array}}}
\vnbprop{haar-avg-transport-opp}{\vnbstmt{\begin{array}{@{}l@{}} \forall g, q, f, w. \\ \quad \operatorname{is\text{-}metric\text{-}group}(g) \wedge \\ \quad \operatorname{is\text{-}compact}(\operatorname{metric\text{-}group\text{-}ms}(g)) \wedge \\ \quad q \in \operatorname{haar\text{-}averages}(g) \wedge \\ \quad f \in \operatorname{cont\text{-}fns}(\operatorname{metric\text{-}group\text{-}ms}(g)) \wedge \\ \quad w \\ \quad \in\; \operatorname{closure}(\operatorname{nvs\text{-}metric\text{-}space}(C({\operatorname{metric\text{-}group\text{-}ms}(g)})), \\ \quad \quad \operatorname{conv\text{-}hull}(C({\operatorname{metric\text{-}group\text{-}ms}(g)}), \\ \quad \quad \quad \operatorname{image}((\lambda a \in \operatorname{pts}(\operatorname{metric\text{-}group\text{-}ms}(g)).\; \\ \quad \quad \quad \quad \quad (\lambda x \in \operatorname{pts}(\operatorname{metric\text{-}group\text{-}ms}(g)).\; f(\operatorname{opr}({g}^{\mathrm{op}})(a, x)))), \\ \quad \quad \quad \quad \operatorname{pts}(\operatorname{metric\text{-}group\text{-}ms}(g))))) \Rightarrow \\ \quad \quad q(w) \\ \quad \quad \in\; \operatorname{closure}(\operatorname{nvs\text{-}metric\text{-}space}(C({\operatorname{metric\text{-}group\text{-}ms}(g)})), \\ \quad \quad \quad \operatorname{conv\text{-}hull}(C({\operatorname{metric\text{-}group\text{-}ms}(g)}), \\ \quad \quad \quad \quad \operatorname{image}((\lambda a \in \operatorname{pts}(\operatorname{metric\text{-}group\text{-}ms}(g)).\; \\ \quad \quad \quad \quad \quad \quad (\lambda x \in \operatorname{pts}(\operatorname{metric\text{-}group\text{-}ms}(g)).\; q(f)(\operatorname{opr}({g}^{\mathrm{op}})(a, x)))), \\ \quad \quad \quad \quad \quad \operatorname{pts}(\operatorname{metric\text{-}group\text{-}ms}(g))))) \end{array}}}
\vnbprop{haar-constants-agree}{\vnbstmt{\begin{array}{@{}l@{}} \forall g, f, c, d. \\ \quad \operatorname{is\text{-}metric\text{-}group}(g) \wedge \\ \quad \operatorname{is\text{-}compact}(\operatorname{metric\text{-}group\text{-}ms}(g)) \wedge \\ \quad f \in \operatorname{cont\text{-}fns}(\operatorname{metric\text{-}group\text{-}ms}(g)) \wedge \\ \quad c \in \mathbb{R} \wedge \\ \quad d \in \mathbb{R} \wedge \\ \quad (\lambda x \in \operatorname{pts}(\operatorname{metric\text{-}group\text{-}ms}(g)).\; c) \\ \quad \in\; \operatorname{closure}(\operatorname{nvs\text{-}metric\text{-}space}(C({\operatorname{metric\text{-}group\text{-}ms}(g)})), \\ \quad \quad \operatorname{conv\text{-}hull}(C({\operatorname{metric\text{-}group\text{-}ms}(g)}), \\ \quad \quad \quad \operatorname{image}((\lambda a \in \operatorname{pts}(\operatorname{metric\text{-}group\text{-}ms}(g)).\; (\lambda x \in \operatorname{pts}(\operatorname{metric\text{-}group\text{-}ms}(g)).\; f(\operatorname{opr}(g)(a, x)))), \\ \quad \quad \quad \quad \operatorname{pts}(\operatorname{metric\text{-}group\text{-}ms}(g))))) \wedge \\ \quad (\lambda x \in \operatorname{pts}(\operatorname{metric\text{-}group\text{-}ms}(g)).\; d) \\ \quad \in\; \operatorname{closure}(\operatorname{nvs\text{-}metric\text{-}space}(C({\operatorname{metric\text{-}group\text{-}ms}(g)})), \\ \quad \quad \operatorname{conv\text{-}hull}(C({\operatorname{metric\text{-}group\text{-}ms}(g)}), \\ \quad \quad \quad \operatorname{image}((\lambda a \in \operatorname{pts}(\operatorname{metric\text{-}group\text{-}ms}(g)).\; \\ \quad \quad \quad \quad \quad (\lambda x \in \operatorname{pts}(\operatorname{metric\text{-}group\text{-}ms}(g)).\; f(\operatorname{opr}({g}^{\mathrm{op}})(a, x)))), \\ \quad \quad \quad \quad \operatorname{pts}(\operatorname{metric\text{-}group\text{-}ms}(g))))) \Rightarrow \\ \quad \quad c = d \end{array}}}
\vnbprop{haar-mean-unfold}{\vnbstmt{\begin{array}{@{}l@{}} \forall g, f. \\ \quad (\int_{g} {f} \simeq (\iota\, c.\;((c \in \mathbb{R}) \wedge ((\lambda x \in \operatorname{pts}(\operatorname{metric\text{-}group\text{-}ms}(g)).\; c) \in \operatorname{closure}(\operatorname{nvs\text{-}metric\text{-}space}(C({\operatorname{metric\text{-}group\text{-}ms}(g)})), \operatorname{conv\text{-}hull}(C({\operatorname{metric\text{-}group\text{-}ms}(g)}), \operatorname{image}((\lambda a \in \operatorname{pts}(\operatorname{metric\text{-}group\text{-}ms}(g)).\; (\lambda x \in \operatorname{pts}(\operatorname{metric\text{-}group\text{-}ms}(g)).\; f(\operatorname{opr}(g)(a, x)))), \operatorname{pts}(\operatorname{metric\text{-}group\text{-}ms}(g))))))))) \end{array}}}
\vnbprop{haar-mean-spec}{\vnbstmt{\begin{array}{@{}l@{}} \forall g, f. \\ \quad (\operatorname{is\text{-}metric\text{-}group}(g) \wedge (\operatorname{is\text{-}compact}(\operatorname{metric\text{-}group\text{-}ms}(g)) \wedge (f \in \operatorname{cont\text{-}fns}(\operatorname{metric\text{-}group\text{-}ms}(g))))) \Rightarrow  \\ \quad \quad \int_{g} {f} \in \mathbb{R} \wedge \\ \quad \quad (\lambda x \in \operatorname{pts}(\operatorname{metric\text{-}group\text{-}ms}(g)).\; \int_{g} {f}) \\ \quad \quad \in\; \operatorname{closure}(\operatorname{nvs\text{-}metric\text{-}space}(C({\operatorname{metric\text{-}group\text{-}ms}(g)})), \\ \quad \quad \quad \operatorname{conv\text{-}hull}(C({\operatorname{metric\text{-}group\text{-}ms}(g)}), \\ \quad \quad \quad \quad \operatorname{image}((\lambda a \in \operatorname{pts}(\operatorname{metric\text{-}group\text{-}ms}(g)).\; (\lambda x \in \operatorname{pts}(\operatorname{metric\text{-}group\text{-}ms}(g)).\; f(\operatorname{opr}(g)(a, x)))), \\ \quad \quad \quad \quad \quad \operatorname{pts}(\operatorname{metric\text{-}group\text{-}ms}(g))))) \end{array}}}
\vnbprop{haar-mean-eq}{\vnbstmt{\begin{array}{@{}l@{}} \forall g, f, c. \\ \quad \operatorname{is\text{-}metric\text{-}group}(g) \wedge \\ \quad \operatorname{is\text{-}compact}(\operatorname{metric\text{-}group\text{-}ms}(g)) \wedge \\ \quad f \in \operatorname{cont\text{-}fns}(\operatorname{metric\text{-}group\text{-}ms}(g)) \wedge \\ \quad c \in \mathbb{R} \wedge \\ \quad (\lambda x \in \operatorname{pts}(\operatorname{metric\text{-}group\text{-}ms}(g)).\; c) \\ \quad \in\; \operatorname{closure}(\operatorname{nvs\text{-}metric\text{-}space}(C({\operatorname{metric\text{-}group\text{-}ms}(g)})), \\ \quad \quad \operatorname{conv\text{-}hull}(C({\operatorname{metric\text{-}group\text{-}ms}(g)}), \\ \quad \quad \quad \operatorname{image}((\lambda a \in \operatorname{pts}(\operatorname{metric\text{-}group\text{-}ms}(g)).\; (\lambda x \in \operatorname{pts}(\operatorname{metric\text{-}group\text{-}ms}(g)).\; f(\operatorname{opr}(g)(a, x)))), \\ \quad \quad \quad \quad \operatorname{pts}(\operatorname{metric\text{-}group\text{-}ms}(g))))) \Rightarrow \\ \quad \quad \int_{g} {f} = c \end{array}}}
\vnbprop{haar-mean-opp}{\vnbstmt{\begin{array}{@{}l@{}} \forall g, f. \\ \quad (\operatorname{is\text{-}metric\text{-}group}(g) \wedge (\operatorname{is\text{-}compact}(\operatorname{metric\text{-}group\text{-}ms}(g)) \wedge (f \in \operatorname{cont\text{-}fns}(\operatorname{metric\text{-}group\text{-}ms}(g))))) \Rightarrow  \\ \quad \quad \int_{{g}^{\mathrm{op}}} {f} = \int_{g} {f} \end{array}}}
\vnbprop{haar-avg-near}{\vnbstmt{\begin{array}{@{}l@{}} \forall g, f, e. \\ \quad \operatorname{is\text{-}metric\text{-}group}(g) \wedge \\ \quad \operatorname{is\text{-}compact}(\operatorname{metric\text{-}group\text{-}ms}(g)) \wedge \\ \quad f \in \operatorname{cont\text{-}fns}(\operatorname{metric\text{-}group\text{-}ms}(g)) \wedge \\ \quad \operatorname{pos\text{-}rr}(e) \Rightarrow \\ \quad \quad \exists q \in \operatorname{haar\text{-}averages}(g). \\ \quad \quad \quad \operatorname{dist}(\operatorname{nvs\text{-}metric\text{-}space}(C({\operatorname{metric\text{-}group\text{-}ms}(g)})))(q(f), (\lambda x \in \operatorname{pts}(\operatorname{metric\text{-}group\text{-}ms}(g)).\; \int_{g} {f})) \\ \quad \quad \quad \leq\; e \end{array}}}
\vnbprop{haar-mean-avg}{\vnbstmt{\begin{array}{@{}l@{}} \forall g, q, f. \\ \quad \operatorname{is\text{-}metric\text{-}group}(g) \wedge \\ \quad \operatorname{is\text{-}compact}(\operatorname{metric\text{-}group\text{-}ms}(g)) \wedge \\ \quad q \in \operatorname{haar\text{-}averages}(g) \wedge \\ \quad f \in \operatorname{cont\text{-}fns}(\operatorname{metric\text{-}group\text{-}ms}(g)) \Rightarrow \\ \quad \quad \int_{g} {q(f)} = \int_{g} {f} \end{array}}}
\vnbprop{haar-mean-dist-const}{\vnbstmt{\begin{array}{@{}l@{}} \forall g, f, k. \\ \quad \operatorname{is\text{-}metric\text{-}group}(g) \wedge \\ \quad \operatorname{is\text{-}compact}(\operatorname{metric\text{-}group\text{-}ms}(g)) \wedge \\ \quad f \in \operatorname{cont\text{-}fns}(\operatorname{metric\text{-}group\text{-}ms}(g)) \wedge \\ \quad k \in \mathbb{R} \Rightarrow \\ \quad \quad \lvert (\int_{g} {f} - k) \rvert \\ \quad \quad \leq\; \operatorname{dist}(\operatorname{nvs\text{-}metric\text{-}space}(C({\operatorname{metric\text{-}group\text{-}ms}(g)})))(f, (\lambda x \in \operatorname{pts}(\operatorname{metric\text{-}group\text{-}ms}(g)).\; k)) \end{array}}}
\vnbprop{haar-mean-const}{\vnbstmt{\begin{array}{@{}l@{}} \forall g, k. \\ \quad (\operatorname{is\text{-}metric\text{-}group}(g) \wedge (\operatorname{is\text{-}compact}(\operatorname{metric\text{-}group\text{-}ms}(g)) \wedge (k \in \mathbb{R}))) \Rightarrow  \\ \quad \quad \int_{g} {(\lambda x \in \operatorname{pts}(\operatorname{metric\text{-}group\text{-}ms}(g)).\; k)} = k \end{array}}}
\vnbprop{haar-mean-left-inv}{\vnbstmt{\begin{array}{@{}l@{}} \forall g, a, f. \\ \quad \operatorname{is\text{-}metric\text{-}group}(g) \wedge \\ \quad \operatorname{is\text{-}compact}(\operatorname{metric\text{-}group\text{-}ms}(g)) \wedge \\ \quad a \in \operatorname{pts}(\operatorname{metric\text{-}group\text{-}ms}(g)) \wedge \\ \quad f \in \operatorname{cont\text{-}fns}(\operatorname{metric\text{-}group\text{-}ms}(g)) \Rightarrow \\ \quad \quad \int_{g} {(\lambda x \in \operatorname{pts}(\operatorname{metric\text{-}group\text{-}ms}(g)).\; f(\operatorname{opr}(g)(a, x)))} = \int_{g} {f} \end{array}}}
\vnbprop{haar-mean-right-inv}{\vnbstmt{\begin{array}{@{}l@{}} \forall g, b, f. \\ \quad \operatorname{is\text{-}metric\text{-}group}(g) \wedge \\ \quad \operatorname{is\text{-}compact}(\operatorname{metric\text{-}group\text{-}ms}(g)) \wedge \\ \quad b \in \operatorname{pts}(\operatorname{metric\text{-}group\text{-}ms}(g)) \wedge \\ \quad f \in \operatorname{cont\text{-}fns}(\operatorname{metric\text{-}group\text{-}ms}(g)) \Rightarrow \\ \quad \quad \int_{g} {(\lambda x \in \operatorname{pts}(\operatorname{metric\text{-}group\text{-}ms}(g)).\; f(\operatorname{opr}(g)(x, b)))} = \int_{g} {f} \end{array}}}
\vnbprop{ck-dist-add-le}{\vnbstmt{\begin{array}{@{}l@{}} \forall g, u, v, a, b. \\ \quad \operatorname{is\text{-}metric\text{-}group}(g) \wedge \\ \quad \operatorname{is\text{-}compact}(\operatorname{metric\text{-}group\text{-}ms}(g)) \wedge \\ \quad u \in \operatorname{cont\text{-}fns}(\operatorname{metric\text{-}group\text{-}ms}(g)) \wedge \\ \quad v \in \operatorname{cont\text{-}fns}(\operatorname{metric\text{-}group\text{-}ms}(g)) \wedge \\ \quad a \in \operatorname{cont\text{-}fns}(\operatorname{metric\text{-}group\text{-}ms}(g)) \wedge \\ \quad b \in \operatorname{cont\text{-}fns}(\operatorname{metric\text{-}group\text{-}ms}(g)) \Rightarrow \\ \quad \quad \operatorname{dist}(\operatorname{nvs\text{-}metric\text{-}space}(C({\operatorname{metric\text{-}group\text{-}ms}(g)})))(\operatorname{vadd}(C({\operatorname{metric\text{-}group\text{-}ms}(g)}))(u, v), \\ \quad \quad \quad \operatorname{vadd}(C({\operatorname{metric\text{-}group\text{-}ms}(g)}))(a, b)) \\ \quad \quad \leq\; \operatorname{dist}(\operatorname{nvs\text{-}metric\text{-}space}(C({\operatorname{metric\text{-}group\text{-}ms}(g)})))(u, a) + \operatorname{dist}(\operatorname{nvs\text{-}metric\text{-}space}(C({\operatorname{metric\text{-}group\text{-}ms}(g)})))(v, b) \end{array}}}
\vnbprop{ck-dist-act-le}{\vnbstmt{\begin{array}{@{}l@{}} \forall g, c, u, v. \\ \quad \operatorname{is\text{-}metric\text{-}group}(g) \wedge \\ \quad \operatorname{is\text{-}compact}(\operatorname{metric\text{-}group\text{-}ms}(g)) \wedge \\ \quad c \in \mathbb{R} \wedge \\ \quad u \in \operatorname{cont\text{-}fns}(\operatorname{metric\text{-}group\text{-}ms}(g)) \wedge \\ \quad v \in \operatorname{cont\text{-}fns}(\operatorname{metric\text{-}group\text{-}ms}(g)) \Rightarrow \\ \quad \quad \operatorname{dist}(\operatorname{nvs\text{-}metric\text{-}space}(C({\operatorname{metric\text{-}group\text{-}ms}(g)})))(\operatorname{act}(C({\operatorname{metric\text{-}group\text{-}ms}(g)}))(c, u), \\ \quad \quad \quad \operatorname{act}(C({\operatorname{metric\text{-}group\text{-}ms}(g)}))(c, v)) \\ \quad \quad \leq\; \operatorname{dist}(\operatorname{nvs\text{-}metric\text{-}space}(C({\operatorname{metric\text{-}group\text{-}ms}(g)})))(u, v) \cdot \lvert c \rvert \end{array}}}
\vnbprop{ck-vadd-const}{\vnbstmt{\begin{array}{@{}l@{}} \forall g, j, k. \\ \quad (\operatorname{is\text{-}metric\text{-}group}(g) \wedge (\operatorname{is\text{-}compact}(\operatorname{metric\text{-}group\text{-}ms}(g)) \wedge ((j \in \mathbb{R}) \wedge (k \in \mathbb{R})))) \Rightarrow  \\ \quad \quad \operatorname{vadd}(C({\operatorname{metric\text{-}group\text{-}ms}(g)}))((\lambda x \in \operatorname{pts}(\operatorname{metric\text{-}group\text{-}ms}(g)).\; j), \\ \quad \quad \quad (\lambda x \in \operatorname{pts}(\operatorname{metric\text{-}group\text{-}ms}(g)).\; k)) \\ \quad \quad =\; (\lambda x \in \operatorname{pts}(\operatorname{metric\text{-}group\text{-}ms}(g)).\; (j + k)) \end{array}}}
\vnbprop{ck-act-const}{\vnbstmt{\begin{array}{@{}l@{}} \forall g, c, k. \\ \quad (\operatorname{is\text{-}metric\text{-}group}(g) \wedge (\operatorname{is\text{-}compact}(\operatorname{metric\text{-}group\text{-}ms}(g)) \wedge ((c \in \mathbb{R}) \wedge (k \in \mathbb{R})))) \Rightarrow  \\ \quad \quad \operatorname{act}(C({\operatorname{metric\text{-}group\text{-}ms}(g)}))(c, (\lambda x \in \operatorname{pts}(\operatorname{metric\text{-}group\text{-}ms}(g)).\; k)) \\ \quad \quad =\; (\lambda x \in \operatorname{pts}(\operatorname{metric\text{-}group\text{-}ms}(g)).\; (c \cdot k)) \end{array}}}
\vnbprop{haar-mean-add}{\vnbstmt{\begin{array}{@{}l@{}} \forall g, f, h. \\ \quad \operatorname{is\text{-}metric\text{-}group}(g) \wedge \\ \quad \operatorname{is\text{-}compact}(\operatorname{metric\text{-}group\text{-}ms}(g)) \wedge \\ \quad f \in \operatorname{cont\text{-}fns}(\operatorname{metric\text{-}group\text{-}ms}(g)) \wedge \\ \quad h \in \operatorname{cont\text{-}fns}(\operatorname{metric\text{-}group\text{-}ms}(g)) \Rightarrow \\ \quad \quad \int_{g} {\operatorname{vadd}(C({\operatorname{metric\text{-}group\text{-}ms}(g)}))(f, h)} = (\int_{g} {f} + \int_{g} {h}) \end{array}}}
\vnbprop{haar-mean-act}{\vnbstmt{\begin{array}{@{}l@{}} \forall g, c, f. \\ \quad \operatorname{is\text{-}metric\text{-}group}(g) \wedge \\ \quad \operatorname{is\text{-}compact}(\operatorname{metric\text{-}group\text{-}ms}(g)) \wedge \\ \quad c \in \mathbb{R} \wedge \\ \quad f \in \operatorname{cont\text{-}fns}(\operatorname{metric\text{-}group\text{-}ms}(g)) \Rightarrow \\ \quad \quad \int_{g} {\operatorname{act}(C({\operatorname{metric\text{-}group\text{-}ms}(g)}))(c, f)} = (c \cdot \int_{g} {f}) \end{array}}}
\vnbprop{haar-mean-nonneg}{\vnbstmt{\begin{array}{@{}l@{}} \forall g, f. \\ \quad \operatorname{is\text{-}metric\text{-}group}(g) \wedge \\ \quad \operatorname{is\text{-}compact}(\operatorname{metric\text{-}group\text{-}ms}(g)) \wedge \\ \quad f \in \operatorname{cont\text{-}fns}(\operatorname{metric\text{-}group\text{-}ms}(g)) \wedge \\ \quad \forall x \in \operatorname{pts}(\operatorname{metric\text{-}group\text{-}ms}(g)).\; (0 \leq f(x)) \Rightarrow \\ \quad \quad 0 \leq \int_{g} {f} \end{array}}}
\vnbprop{haar-mean-is-haar-integral}{\vnbstmt{\begin{array}{@{}l@{}} \forall g. \\ \quad (\operatorname{is\text{-}metric\text{-}group}(g) \wedge \operatorname{is\text{-}compact}(\operatorname{metric\text{-}group\text{-}ms}(g))) \Rightarrow  \\ \quad \quad \operatorname{is\text{-}haar\text{-}integral}(g, (\lambda m \in \operatorname{cont\text{-}fns}(\operatorname{metric\text{-}group\text{-}ms}(g)).\; \int_{g} {m})) \end{array}}}
\vnbprop{haar-integral-add}{\vnbstmt{\begin{array}{@{}l@{}} \forall g, a, u, v. \\ \quad \operatorname{is\text{-}metric\text{-}group}(g) \wedge \\ \quad \operatorname{is\text{-}compact}(\operatorname{metric\text{-}group\text{-}ms}(g)) \wedge \\ \quad \operatorname{is\text{-}haar\text{-}integral}(g, a) \wedge \\ \quad u \in \operatorname{cont\text{-}fns}(\operatorname{metric\text{-}group\text{-}ms}(g)) \wedge \\ \quad v \in \operatorname{cont\text{-}fns}(\operatorname{metric\text{-}group\text{-}ms}(g)) \Rightarrow \\ \quad \quad a(\operatorname{vadd}(C({\operatorname{metric\text{-}group\text{-}ms}(g)}))(u, v)) = (a(u) + a(v)) \end{array}}}
\vnbprop{haar-integral-act}{\vnbstmt{\begin{array}{@{}l@{}} \forall g, a, r, u. \\ \quad \operatorname{is\text{-}metric\text{-}group}(g) \wedge \\ \quad \operatorname{is\text{-}compact}(\operatorname{metric\text{-}group\text{-}ms}(g)) \wedge \\ \quad \operatorname{is\text{-}haar\text{-}integral}(g, a) \wedge \\ \quad r \in \mathbb{R} \wedge \\ \quad u \in \operatorname{cont\text{-}fns}(\operatorname{metric\text{-}group\text{-}ms}(g)) \Rightarrow \\ \quad \quad a(\operatorname{act}(C({\operatorname{metric\text{-}group\text{-}ms}(g)}))(r, u)) = (r \cdot a(u)) \end{array}}}
\vnbprop{ck-act-const-one}{\vnbstmt{\begin{array}{@{}l@{}} \forall g, k. \\ \quad (\operatorname{is\text{-}metric\text{-}group}(g) \wedge (\operatorname{is\text{-}compact}(\operatorname{metric\text{-}group\text{-}ms}(g)) \wedge (k \in \mathbb{R}))) \Rightarrow  \\ \quad \quad \operatorname{act}(C({\operatorname{metric\text{-}group\text{-}ms}(g)}))(k, (\lambda x \in \operatorname{pts}(\operatorname{metric\text{-}group\text{-}ms}(g)).\; 1)) \\ \quad \quad =\; (\lambda x \in \operatorname{pts}(\operatorname{metric\text{-}group\text{-}ms}(g)).\; k) \end{array}}}
\vnbprop{haar-integral-const}{\vnbstmt{\begin{array}{@{}l@{}} \forall g, u, k. \\ \quad (\operatorname{is\text{-}metric\text{-}group}(g) \wedge (\operatorname{is\text{-}compact}(\operatorname{metric\text{-}group\text{-}ms}(g)) \wedge (\operatorname{is\text{-}haar\text{-}integral}(g, u) \wedge (k \in \mathbb{R})))) \Rightarrow  \\ \quad \quad u((\lambda x \in \operatorname{pts}(\operatorname{metric\text{-}group\text{-}ms}(g)).\; k)) = k \end{array}}}
\vnbprop{haar-integral-lip}{\vnbstmt{\begin{array}{@{}l@{}} \forall g, a, u, v. \\ \quad \operatorname{is\text{-}metric\text{-}group}(g) \wedge \\ \quad \operatorname{is\text{-}compact}(\operatorname{metric\text{-}group\text{-}ms}(g)) \wedge \\ \quad \operatorname{is\text{-}haar\text{-}integral}(g, a) \wedge \\ \quad u \in \operatorname{cont\text{-}fns}(\operatorname{metric\text{-}group\text{-}ms}(g)) \wedge \\ \quad v \in \operatorname{cont\text{-}fns}(\operatorname{metric\text{-}group\text{-}ms}(g)) \Rightarrow \\ \quad \quad (a(u) - a(v)) \leq \operatorname{dist}(\operatorname{nvs\text{-}metric\text{-}space}(C({\operatorname{metric\text{-}group\text{-}ms}(g)})))(u, v) \end{array}}}
\vnbprop{haar-integral-hull-const}{\vnbstmt{\begin{array}{@{}l@{}} \forall g, u, f.; \forall w \in \operatorname{conv\text{-}hull}(C({\operatorname{metric\text{-}group\text{-}ms}(g)}), \operatorname{image}((\lambda a \in \operatorname{pts}(\operatorname{metric\text{-}group\text{-}ms}(g)).\; (\lambda x \in \operatorname{pts}(\operatorname{metric\text{-}group\text{-}ms}(g)).\; f(\operatorname{opr}(g)(a, x)))), \operatorname{pts}(\operatorname{metric\text{-}group\text{-}ms}(g)))). \\ \quad \operatorname{is\text{-}metric\text{-}group}(g) \wedge \\ \quad \operatorname{is\text{-}compact}(\operatorname{metric\text{-}group\text{-}ms}(g)) \wedge \\ \quad \operatorname{is\text{-}haar\text{-}integral}(g, u) \wedge \\ \quad f \in \operatorname{cont\text{-}fns}(\operatorname{metric\text{-}group\text{-}ms}(g)) \Rightarrow \\ \quad \quad u(w) = u(f) \end{array}}}
\vnbprop{haar-integral-closure-const}{\vnbstmt{\begin{array}{@{}l@{}} \forall g, u, f, w. \\ \quad \operatorname{is\text{-}metric\text{-}group}(g) \wedge \\ \quad \operatorname{is\text{-}compact}(\operatorname{metric\text{-}group\text{-}ms}(g)) \wedge \\ \quad \operatorname{is\text{-}haar\text{-}integral}(g, u) \wedge \\ \quad f \in \operatorname{cont\text{-}fns}(\operatorname{metric\text{-}group\text{-}ms}(g)) \wedge \\ \quad w \\ \quad \in\; \operatorname{closure}(\operatorname{nvs\text{-}metric\text{-}space}(C({\operatorname{metric\text{-}group\text{-}ms}(g)})), \\ \quad \quad \operatorname{conv\text{-}hull}(C({\operatorname{metric\text{-}group\text{-}ms}(g)}), \\ \quad \quad \quad \operatorname{image}((\lambda a \in \operatorname{pts}(\operatorname{metric\text{-}group\text{-}ms}(g)).\; (\lambda x \in \operatorname{pts}(\operatorname{metric\text{-}group\text{-}ms}(g)).\; f(\operatorname{opr}(g)(a, x)))), \\ \quad \quad \quad \quad \operatorname{pts}(\operatorname{metric\text{-}group\text{-}ms}(g))))) \Rightarrow \\ \quad \quad u(w) = u(f) \end{array}}}
\vnbprop{haar-integral-unique}{\vnbstmt{\begin{array}{@{}l@{}} \forall g, u, f. \\ \quad \operatorname{is\text{-}metric\text{-}group}(g) \wedge \\ \quad \operatorname{is\text{-}compact}(\operatorname{metric\text{-}group\text{-}ms}(g)) \wedge \\ \quad \operatorname{is\text{-}haar\text{-}integral}(g, u) \wedge \\ \quad f \in \operatorname{cont\text{-}fns}(\operatorname{metric\text{-}group\text{-}ms}(g)) \Rightarrow \\ \quad \quad u(f) = \int_{g} {f} \end{array}}}
\vnbprop{haar-integral-right-inv}{\vnbstmt{\begin{array}{@{}l@{}} \forall g, u, b, f. \\ \quad \operatorname{is\text{-}metric\text{-}group}(g) \wedge \\ \quad \operatorname{is\text{-}compact}(\operatorname{metric\text{-}group\text{-}ms}(g)) \wedge \\ \quad \operatorname{is\text{-}haar\text{-}integral}(g, u) \wedge \\ \quad b \in \operatorname{pts}(\operatorname{metric\text{-}group\text{-}ms}(g)) \wedge \\ \quad f \in \operatorname{cont\text{-}fns}(\operatorname{metric\text{-}group\text{-}ms}(g)) \Rightarrow \\ \quad \quad u((\lambda x \in \operatorname{pts}(\operatorname{metric\text{-}group\text{-}ms}(g)).\; f(\operatorname{opr}(g)(x, b)))) = u(f) \end{array}}}
\vnbprop{haar-integral-unfold}{\vnbstmt{\forall g.\; (\operatorname{haar\text{-}integral}(g) \simeq (\iota\, u.\;\operatorname{is\text{-}haar\text{-}integral}(g, u)))}}
\vnbprop{haar-integral-is-haar}{\vnbstmt{\begin{array}{@{}l@{}} \forall g. \\ \quad (\operatorname{is\text{-}metric\text{-}group}(g) \wedge \operatorname{is\text{-}compact}(\operatorname{metric\text{-}group\text{-}ms}(g))) \Rightarrow  \\ \quad \quad \operatorname{is\text{-}haar\text{-}integral}(g, \operatorname{haar\text{-}integral}(g)) \end{array}}}
\vnbprop{haar-integral-value}{\vnbstmt{\begin{array}{@{}l@{}} \forall g, f. \\ \quad (\operatorname{is\text{-}metric\text{-}group}(g) \wedge (\operatorname{is\text{-}compact}(\operatorname{metric\text{-}group\text{-}ms}(g)) \wedge (f \in \operatorname{cont\text{-}fns}(\operatorname{metric\text{-}group\text{-}ms}(g))))) \Rightarrow  \\ \quad \quad \operatorname{haar\text{-}integral}(g)(f) = \int_{g} {f} \end{array}}}

\vnbbookfile{theorem-library/group-orders}{theorem-library/group-orders}
\noindent{\small theorem-library/group-orders.scm -- GAUSS'S IDENTITY, THE PARTITION BY ORDERS, LEMMA 5.44 (ALGEBRAIC-NUMBERS notes; ALGNUM, 2026-10-27).  Sums are FINSUMs in C (as inclusion-exclusion). card-disjoint-family       card of the union of a finite disjoint family = sum of the cards zz-gcd-pos-exists, zz-gcd-pos-unique, ...}\par
\vnbprop{card-disjoint-family}{\vnbstmt{\begin{array}{@{}l@{}} \forall i \in \mathbf{Set}, f. \\ \quad |i| \in \mathbb{N} \wedge \\ \quad \forall \mathit{fi} \in i.\; ((f(\mathit{fi}) \in \mathbf{Set}) \wedge (|f(\mathit{fi})| \in \mathbb{N})) \wedge \\ \quad \forall a, b \in i, w \in f(a).\; ((w \in f(b)) \Rightarrow (a = b)) \Rightarrow \\ \quad \quad |\operatorname{big\text{-}union}(u, i, f(u))| \\ \quad \quad =\; \operatorname{finsum}(\operatorname{ring\text{-}additive\text{-}ag}(\operatorname{normed\text{-}field\text{-}as\text{-}commutative\text{-}ring}(\operatorname{cc\text{-}normed\text{-}field})), \\ \quad \quad \quad (\lambda c \in i.\; |f(c)|), i) \end{array}}}
\vnbprop{zz-gcd-pos-exists}{\vnbstmt{\begin{array}{@{}l@{}} \forall a, b. \\ \quad ((a \in \mathbb{Z}) \wedge ((b \in \mathbb{Z}) \wedge \neg (b = 0))) \Rightarrow  \\ \quad \quad \exists g \in \mathbb{N}.\; ((1 \leq g) \wedge \operatorname{zz\text{-}is\text{-}gcd}(a, b, g)) \end{array}}}
\vnbprop{zz-divides-le-pos}{\vnbstmt{\begin{array}{@{}l@{}} \forall g, h. \\ \quad ((g \in \mathbb{Z}) \wedge ((h \in \mathbb{Z}) \wedge ((1 \leq g) \wedge ((1 \leq h) \wedge \operatorname{zz\text{-}divides}(g, h))))) \Rightarrow (g \leq h) \end{array}}}
\vnbprop{zz-gcd-pos-unique}{\vnbstmt{\begin{array}{@{}l@{}} \forall a, b, g, h. \\ \quad ((g \in \mathbb{Z}) \wedge ((h \in \mathbb{Z}) \wedge ((1 \leq g) \wedge ((1 \leq h) \wedge (\operatorname{zz\text{-}is\text{-}gcd}(a, b, g) \wedge \operatorname{zz\text{-}is\text{-}gcd}(a, b, h)))))) \Rightarrow  \\ \quad \quad g = h \end{array}}}
\vnbprop{zz-gcd-scale}{\vnbstmt{\begin{array}{@{}l@{}} \forall c, j, e. \\ \quad ((c \in \mathbb{Z}) \wedge ((j \in \mathbb{Z}) \wedge ((e \in \mathbb{Z}) \wedge (1 \leq c)))) \Rightarrow  \\ \quad \quad \operatorname{zz\text{-}is\text{-}gcd}((c \cdot j), (c \cdot e), c) \Leftrightarrow \operatorname{zz\text{-}coprime}(j, e) \end{array}}}
\vnbprop{gauss-class-card}{\vnbstmt{\begin{array}{@{}l@{}} \forall n, d. \\ \quad ((n \in \mathbb{N}) \wedge ((d \in \mathbb{N}) \wedge ((1 \leq n) \wedge ((1 \leq d) \wedge \operatorname{zz\text{-}divides}(d, n))))) \Rightarrow  \\ \quad \quad |\{k \in \operatorname{interval}(1, n) : \exists g \in \mathbb{N}.\; ((1 \leq g) \wedge (\operatorname{zz\text{-}is\text{-}gcd}(k, n, g) \wedge ((g \cdot d) = n)))\}| \\ \quad \quad =\; \operatorname{euler\text{-}phi}(d) \end{array}}}
\vnbprop{gauss-phi-sum}{\vnbstmt{\begin{array}{@{}l@{}} \forall n \in \mathbb{N}. \\ \quad (1 \leq n) \Rightarrow  \\ \quad \quad \operatorname{finsum}(\operatorname{ring\text{-}additive\text{-}ag}(\operatorname{normed\text{-}field\text{-}as\text{-}commutative\text{-}ring}(\operatorname{cc\text{-}normed\text{-}field})), \\ \quad \quad \quad (\lambda d \in \{\mathit{gd} \in \operatorname{interval}(1, n) : \operatorname{zz\text{-}divides}(\mathit{gd}, n)\}.\; \operatorname{euler\text{-}phi}(d)), \\ \quad \quad \quad \{\mathit{gd} \in \operatorname{interval}(1, n) : \operatorname{zz\text{-}divides}(\mathit{gd}, n)\}) \\ \quad \quad =\; n \end{array}}}
\vnbprop{powers-subgroup}{\vnbstmt{\begin{array}{@{}l@{}} \forall g.; \forall x \in \operatorname{carr}(g). \\ \quad \operatorname{is\text{-}group}(g) \Rightarrow  \\ \quad \quad \operatorname{is\text{-}subgroup}(g, \{y \in \operatorname{carr}(g) : \exists k \in \mathbb{Z}.\; (y = \operatorname{group\text{-}zpow}(g, x, k))\}) \end{array}}}
\vnbprop{powers-card}{\vnbstmt{\begin{array}{@{}l@{}} \forall g.; \forall x \in \operatorname{carr}(g), d. \\ \quad (\operatorname{is\text{-}group}(g) \wedge \operatorname{is\text{-}element\text{-}order}(g, x, d)) \Rightarrow  \\ \quad \quad \{y \in \operatorname{carr}(g) : \exists k \in \mathbb{Z}.\; (y = \operatorname{group\text{-}zpow}(g, x, k))\} \in \mathbf{Set} \wedge \\ \quad \quad |\{y \in \operatorname{carr}(g) : \exists k \in \mathbb{Z}.\; (y = \operatorname{group\text{-}zpow}(g, x, k))\}| = d \end{array}}}
\vnbprop{element-order-divides-card}{\vnbstmt{\begin{array}{@{}l@{}} \forall g.; \forall x \in \operatorname{carr}(g), d. \\ \quad (\operatorname{is\text{-}group}(g) \wedge ((|\operatorname{carr}(g)| \in \mathbb{N}) \wedge \operatorname{is\text{-}element\text{-}order}(g, x, d))) \Rightarrow  \\ \quad \quad \operatorname{zz\text{-}divides}(d, |\operatorname{carr}(g)|) \end{array}}}
\vnbprop{element-order-unique}{\vnbstmt{\begin{array}{@{}l@{}} \forall g.; \forall x \in \operatorname{carr}(g), d, e. \\ \quad (\operatorname{is\text{-}group}(g) \wedge (\operatorname{is\text{-}element\text{-}order}(g, x, d) \wedge \operatorname{is\text{-}element\text{-}order}(g, x, e))) \Rightarrow  \\ \quad \quad d = e \end{array}}}
\vnbprop{order-partition}{\vnbstmt{\begin{array}{@{}l@{}} \forall g. \\ \quad (\operatorname{is\text{-}group}(g) \wedge (|\operatorname{carr}(g)| \in \mathbb{N})) \Rightarrow  \\ \quad \quad |\operatorname{carr}(g)| \\ \quad \quad =\; \operatorname{finsum}(\operatorname{ring\text{-}additive\text{-}ag}(\operatorname{normed\text{-}field\text{-}as\text{-}commutative\text{-}ring}(\operatorname{cc\text{-}normed\text{-}field})), \\ \quad \quad \quad (\lambda d \in \{\mathit{gd} \in \operatorname{interval}(1, |\operatorname{carr}(g)|) : \operatorname{zz\text{-}divides}(\mathit{gd}, |\operatorname{carr}(g)|)\}.\; \\ \quad \quad \quad \quad |\{y \in \operatorname{carr}(g) : \operatorname{is\text{-}element\text{-}order}(g, y, d)\}|), \\ \quad \quad \quad \{\mathit{gd} \in \operatorname{interval}(1, |\operatorname{carr}(g)|) : \operatorname{zz\text{-}divides}(\mathit{gd}, |\operatorname{carr}(g)|)\}) \end{array}}}
\vnbprop{finsum-cc-le-termwise}{\vnbstmt{\begin{array}{@{}l@{}} \forall i \in \mathbf{Set}, a, b. \\ \quad ((|i| \in \mathbb{N}) \wedge \forall \mathit{mi} \in i.\; ((a(\mathit{mi}) \in \mathbb{R}) \wedge ((b(\mathit{mi}) \in \mathbb{R}) \wedge (a(\mathit{mi}) \leq b(\mathit{mi}))))) \Rightarrow  \\ \quad \quad \operatorname{finsum}(\operatorname{ring\text{-}additive\text{-}ag}(\operatorname{normed\text{-}field\text{-}as\text{-}commutative\text{-}ring}(\operatorname{cc\text{-}normed\text{-}field})), a, i) \in \mathbb{R} \wedge \\ \quad \quad \operatorname{finsum}(\operatorname{ring\text{-}additive\text{-}ag}(\operatorname{normed\text{-}field\text{-}as\text{-}commutative\text{-}ring}(\operatorname{cc\text{-}normed\text{-}field})), b, i) \in \mathbb{R} \wedge \\ \quad \quad \operatorname{finsum}(\operatorname{ring\text{-}additive\text{-}ag}(\operatorname{normed\text{-}field\text{-}as\text{-}commutative\text{-}ring}(\operatorname{cc\text{-}normed\text{-}field})), a, i) \\ \quad \quad \leq\; \operatorname{finsum}(\operatorname{ring\text{-}additive\text{-}ag}(\operatorname{normed\text{-}field\text{-}as\text{-}commutative\text{-}ring}(\operatorname{cc\text{-}normed\text{-}field})), b, i) \end{array}}}
\vnbprop{subset-card-eq}{\vnbstmt{\begin{array}{@{}l@{}} \forall a.; \forall b \in \mathbf{Set}. \\ \quad ((|b| \in \mathbb{N}) \wedge ((a \subseteq b) \wedge (|b| \leq |a|))) \Rightarrow (a = b) \end{array}}}
\vnbprop{order-of-power-coprime}{\vnbstmt{\begin{array}{@{}l@{}} \forall g.; \forall x \in \operatorname{carr}(g), d, r. \\ \quad \operatorname{is\text{-}group}(g) \wedge \\ \quad \operatorname{is\text{-}element\text{-}order}(g, x, d) \wedge \\ \quad r \in \mathbb{Z} \wedge \\ \quad 1 \leq r \wedge \\ \quad \operatorname{is\text{-}element\text{-}order}(g, \operatorname{group\text{-}zpow}(g, x, r), d) \Rightarrow \\ \quad \quad \operatorname{zz\text{-}coprime}(r, d) \end{array}}}
\vnbprop{order-count-le-at}{\vnbstmt{\begin{array}{@{}l@{}} \forall g.; \forall x \in \operatorname{carr}(g), d. \\ \quad \operatorname{is\text{-}group}(g) \wedge \\ \quad \operatorname{is\text{-}element\text{-}order}(g, x, d) \wedge \\ \quad |\operatorname{carr}(g)| \in \mathbb{N} \wedge \\ \quad \forall \mathit{hd} \in \mathbb{N}. \\ \quad \quad (1 \leq \mathit{hd}) \Rightarrow (|\{y \in \operatorname{carr}(g) : (\operatorname{gpow}(g, y, \mathit{hd}) = \operatorname{iden}(g))\}| \leq \mathit{hd}) \Rightarrow \\ \quad \quad |\{\mathit{oy} \in \operatorname{carr}(g) : \operatorname{is\text{-}element\text{-}order}(g, \mathit{oy}, d)\}| \leq |\{k \in \operatorname{interval}(1, d) : \operatorname{zz\text{-}coprime}(k, d)\}| \end{array}}}
\vnbprop{order-count-le}{\vnbstmt{\begin{array}{@{}l@{}} \forall g.; \forall d \in \mathbb{N}. \\ \quad \operatorname{is\text{-}group}(g) \wedge \\ \quad |\operatorname{carr}(g)| \in \mathbb{N} \wedge \\ \quad \forall \mathit{hd} \in \mathbb{N}. \\ \quad \quad (1 \leq \mathit{hd}) \Rightarrow (|\{y \in \operatorname{carr}(g) : (\operatorname{gpow}(g, y, \mathit{hd}) = \operatorname{iden}(g))\}| \leq \mathit{hd}) \wedge \\ \quad 1 \leq d \Rightarrow \\ \quad \quad |\{\mathit{oy} \in \operatorname{carr}(g) : \operatorname{is\text{-}element\text{-}order}(g, \mathit{oy}, d)\}| \leq \operatorname{euler\text{-}phi}(d) \end{array}}}
\vnbprop{phiset-card-pos}{\vnbstmt{\begin{array}{@{}l@{}} \forall n \in \mathbb{N}. \\ \quad (1 \leq n) \Rightarrow (1 \leq |\{k \in \operatorname{interval}(1, n) : \operatorname{zz\text{-}coprime}(k, n)\}|) \end{array}}}
\vnbprop{order-card-generator}{\vnbstmt{\begin{array}{@{}l@{}} \forall g.; \forall x \in \operatorname{carr}(g). \\ \quad (\operatorname{is\text{-}group}(g) \wedge ((|\operatorname{carr}(g)| \in \mathbb{N}) \wedge \operatorname{is\text{-}element\text{-}order}(g, x, |\operatorname{carr}(g)|))) \Rightarrow  \\ \quad \quad \operatorname{is\text{-}generator}(g, x) \end{array}}}
\vnbprop{finite-group-cyclic-criterion}{\vnbstmt{\begin{array}{@{}l@{}} \forall g. \\ \quad \operatorname{is\text{-}group}(g) \wedge \\ \quad |\operatorname{carr}(g)| \in \mathbb{N} \wedge \\ \quad \forall d \in \mathbb{N}. \\ \quad \quad (1 \leq d) \Rightarrow (|\{y \in \operatorname{carr}(g) : (\operatorname{gpow}(g, y, d) = \operatorname{iden}(g))\}| \leq d) \Rightarrow \\ \quad \quad \exists x.\; \operatorname{is\text{-}generator}(g, x) \end{array}}}

\vnbbookfile{theorem-library/calc3-extension}{theorem-library/calc3-extension}
\noindent{\small calc3-extension.scm -- calculus.pdf Prop 3.26: A UNIFORMLY CONTINUOUS FUNCTION ON A DENSE SUBSPACE, WITH VALUES IN A COMPLETE SPACE, HAS A UNIQUE UNIFORMLY CONTINUOUS EXTENSION.  Agent CALC-3, 2026-10-19.}\par
\vnbprop{dense-meets-every-ball}{\vnbstmt{\begin{array}{@{}l@{}} \forall s, a.; \forall x \in \operatorname{pts}(s), b. \\ \quad (\operatorname{is\text{-}dense}(s, a) \wedge \operatorname{pos\text{-}rr}(b)) \Rightarrow \exists c \in a.\; (\operatorname{dist}(s)(x, c) < b) \end{array}}}
\vnbprop{limit-along-at-point}{\vnbstmt{\begin{array}{@{}l@{}} \forall s, t, a, f.; \forall x \in a. \\ \quad (\operatorname{is\text{-}metric\text{-}space}(s) \wedge ((a \subseteq \operatorname{pts}(s)) \wedge \operatorname{is\text{-}uniformly\text{-}continuous}(\operatorname{subspace\text{-}ms}(s, a), t, f))) \Rightarrow  \\ \quad \quad \operatorname{is\text{-}limit\text{-}along}(s, t, a, f, x, f(x)) \end{array}}}
\vnbprop{limit-along-of-continuous-at}{\vnbstmt{\begin{array}{@{}l@{}} \forall s, t, h, x, a, f. \\ \quad (\operatorname{is\text{-}continuous\text{-}at}(s, t, h, x) \wedge ((a \subseteq \operatorname{pts}(s)) \wedge \forall z \in a.\; (h(z) = f(z)))) \Rightarrow  \\ \quad \quad \operatorname{is\text{-}limit\text{-}along}(s, t, a, f, x, h(x)) \end{array}}}
\vnbprop{limit-along-unique}{\vnbstmt{\begin{array}{@{}l@{}} \forall s, t, a, f, x, b, c. \\ \quad \operatorname{is\text{-}dense}(s, a) \wedge \\ \quad f \in (a \to \operatorname{pts}(t)) \wedge \\ \quad x \in \operatorname{pts}(s) \wedge \\ \quad \operatorname{is\text{-}limit\text{-}along}(s, t, a, f, x, b) \wedge \\ \quad \operatorname{is\text{-}limit\text{-}along}(s, t, a, f, x, c) \Rightarrow \\ \quad \quad b = c \end{array}}}
\vnbprop{limit-along-exists}{\vnbstmt{\begin{array}{@{}l@{}} \forall s, t, a, f.; \forall x \in \operatorname{pts}(s). \\ \quad (\operatorname{is\text{-}complete}(t) \wedge (\operatorname{is\text{-}dense}(s, a) \wedge \operatorname{is\text{-}uniformly\text{-}continuous}(\operatorname{subspace\text{-}ms}(s, a), t, f))) \Rightarrow  \\ \quad \quad \exists y.\; \operatorname{is\text{-}limit\text{-}along}(s, t, a, f, x, y) \end{array}}}
\vnbprop{uniformly-continuous-extension-from-dense}{\vnbstmt{\begin{array}{@{}l@{}} \forall s, t, a, f. \\ \quad (\operatorname{is\text{-}complete}(t) \wedge (\operatorname{is\text{-}dense}(s, a) \wedge \operatorname{is\text{-}uniformly\text{-}continuous}(\operatorname{subspace\text{-}ms}(s, a), t, f))) \Rightarrow  \\ \quad \quad \exists g. \\ \quad \quad \quad \operatorname{is\text{-}uniformly\text{-}continuous}(s, t, g) \wedge \\ \quad \quad \quad \operatorname{extends\text{-}on}(a, g, f) \wedge \\ \quad \quad \quad \forall h. \\ \quad \quad \quad \quad (\operatorname{is\text{-}uniformly\text{-}continuous}(s, t, h) \wedge \operatorname{extends\text{-}on}(a, h, f)) \Rightarrow (h = g) \end{array}}}

\vnbbookfile{theorem-library/pw-int-order}{theorem-library/pw-int-order}
\noindent{\small pw-int-order.scm -- ORDER PROPERTIES OF THE PIECEWISE INTEGRAL.}\par
\vnbprop{nondecreasing-total}{\vnbstmt{\begin{array}{@{}l@{}} \forall h \in (\mathbb{R} \to \mathbb{R}), u, v \in \mathbb{R}. \\ \quad \forall x \in \mathbb{R}.\; \operatorname{is\text{-}continuous\text{-}at}(\operatorname{rr\text{-}ms}, \operatorname{rr\text{-}ms}, h, x) \wedge \\ \quad u < v \wedge \\ \quad \forall t \in \operatorname{ooint}(u, v).\; \exists l.\; (\operatorname{is\text{-}diff\text{-}at}(h, t, l) \wedge (0 \leq l)) \Rightarrow \\ \quad \quad h(u) \leq h(v) \end{array}}}
\vnbprop{nondecreasing-off-finite-set}{\vnbstmt{\begin{array}{@{}l@{}} \forall s \in \mathbf{Set}, h \in (\mathbb{R} \to \mathbb{R}), u, v \in \mathbb{R}. \\ \quad |s| \in \mathbb{N} \wedge \\ \quad \forall x \in \mathbb{R}.\; \operatorname{is\text{-}continuous\text{-}at}(\operatorname{rr\text{-}ms}, \operatorname{rr\text{-}ms}, h, x) \wedge \\ \quad u < v \wedge \\ \quad \forall t \in \operatorname{ooint}(u, v). \\ \quad \quad \neg (t \in s) \Rightarrow \exists l.\; (\operatorname{is\text{-}diff\text{-}at}(h, t, l) \wedge (0 \leq l)) \Rightarrow \\ \quad \quad h(u) \leq h(v) \end{array}}}
\vnbprop{pw-nondecreasing-on-interval}{\vnbstmt{\begin{array}{@{}l@{}} \forall a, b \in \mathbb{R}, h \in (\operatorname{ccint}(a, b) \to \mathbb{R}), s \in \mathbf{Set}. \\ \quad a < b \wedge \\ \quad \operatorname{is\text{-}continuous\text{-}on}(h, \operatorname{ccint}(a, b)) \wedge \\ \quad |s| \in \mathbb{N} \wedge \\ \quad \forall t \in \operatorname{ooint}(a, b). \\ \quad \quad \neg (t \in s) \Rightarrow \exists l.\; (\operatorname{has\text{-}deriv\text{-}at}(h, t, l) \wedge (0 \leq l)) \Rightarrow \\ \quad \quad h(a) \leq h(b) \end{array}}}
\vnbprop{nondecreasing-off-countable}{\vnbstmt{\begin{array}{@{}l@{}} \forall h \in (\mathbb{R} \to \mathbb{R}), u, v \in \mathbb{R}, s. \\ \quad u < v \wedge \\ \quad \forall x \in \operatorname{ccint}(u, v).\; \operatorname{is\text{-}continuous\text{-}at}(\operatorname{rr\text{-}ms}, \operatorname{rr\text{-}ms}, h, x) \wedge \\ \quad \operatorname{is\text{-}countable}(s) \wedge \\ \quad s \subseteq \mathbb{R} \wedge \\ \quad \forall t \in \operatorname{ooint}(u, v). \\ \quad \quad \neg (t \in s) \Rightarrow \exists l.\; (\operatorname{is\text{-}diff\text{-}at}(h, t, l) \wedge (0 \leq l)) \Rightarrow \\ \quad \quad h(u) \leq h(v) \end{array}}}
\vnbprop{nondecreasing-off-countable-on-interval}{\vnbstmt{\begin{array}{@{}l@{}} \forall a, b \in \mathbb{R}, h \in (\operatorname{ccint}(a, b) \to \mathbb{R}), s. \\ \quad a < b \wedge \\ \quad \operatorname{is\text{-}continuous\text{-}on}(h, \operatorname{ccint}(a, b)) \wedge \\ \quad \operatorname{is\text{-}countable}(s) \wedge \\ \quad s \subseteq \operatorname{ccint}(a, b) \wedge \\ \quad \forall t \in \operatorname{ooint}(a, b). \\ \quad \quad \neg (t \in s) \Rightarrow \exists l.\; (\operatorname{has\text{-}deriv\text{-}at}(h, t, l) \wedge (0 \leq l)) \Rightarrow \\ \quad \quad h(a) \leq h(b) \end{array}}}
\vnbprop{pw-int-nonneg}{\vnbstmt{\begin{array}{@{}l@{}} \forall f, i, a, b. \\ \quad (\operatorname{is\text{-}primitive}(f, i, a, b) \wedge \forall y \in \operatorname{ccint}(a, b).\; (0 \leq i(y))) \Rightarrow  \\ \quad \quad 0 \leq \int_{a}^{b} i \end{array}}}
\vnbprop{pw-int-monotone}{\vnbstmt{\begin{array}{@{}l@{}} \forall f, i, w, \mathit{ppsi}, a, b. \\ \quad \operatorname{is\text{-}primitive}(f, i, a, b) \wedge \\ \quad \operatorname{is\text{-}primitive}(w, \mathit{ppsi}, a, b) \wedge \\ \quad \forall y \in \operatorname{ccint}(a, b).\; (i(y) \leq \operatorname{ppsi}(y)) \Rightarrow \\ \quad \quad \int_{a}^{b} i \leq \int_{a}^{b} \mathit{ppsi} \end{array}}}
\vnbprop{pw-int-abs-bound}{\vnbstmt{\begin{array}{@{}l@{}} \forall f, i, w, s, a, b. \\ \quad \operatorname{is\text{-}primitive}(f, i, a, b) \wedge \\ \quad \operatorname{is\text{-}primitive}(w, s, a, b) \wedge \\ \quad \forall y \in \operatorname{ccint}(a, b).\; (s(y) \simeq \lvert i(y) \rvert) \Rightarrow \\ \quad \quad \lvert \int_{a}^{b} i \rvert \leq \int_{a}^{b} s \end{array}}}

\vnbbookfile{theorem-library/cc-magnitude-sq}{theorem-library/cc-magnitude-sq}
\noindent{\small cc-magnitude-sq.scm -- lemmas MOVED 2026-11-10 (integration of the 11-09/10 waves) out of theorem-library/cc-int-estimate.scm into library analysis, for the library gate: cited by cc-real-bound; CC arithmetic. The proofs are the originals, verbatim; the helper defines are a copy of the source's. cc-magnitude-sq}\par
\vnbprop{cc-magnitude-sq}{\vnbstmt{\begin{array}{@{}l@{}} \forall z \in \mathbb{C}. \\ \quad (\operatorname{magnitude}(z) \cdot \operatorname{magnitude}(z)) = ((\operatorname{real\text{-}part}(z) \cdot \operatorname{real\text{-}part}(z)) + (\operatorname{imag\text{-}part}(z) \cdot \operatorname{imag\text{-}part}(z))) \end{array}}}

\vnbbookfile{theorem-library/regulated-on-glue}{theorem-library/regulated-on-glue}
\noindent{\small regulated-on-glue.scm -- lemmas MOVED 2026-10-30 (wave-2 integration) out of theorem-library/road-juxtaposition-laws.scm into library analysis, for the library gate: cited by calint-regulated-step (analysis); pure real analysis. The proofs are the originals, verbatim; the helper defines are a copy of the source's. ...}\par
\vnbprop{right-limit-within-local}{\vnbstmt{\begin{array}{@{}l@{}} \forall g, v, f, w, x, l, r. \\ \quad \operatorname{is\text{-}right\text{-}limit\text{-}within}(g, v, x, l) \wedge \\ \quad w \subseteq \mathbb{R} \wedge \\ \quad f \in (w \to \mathbb{R}) \wedge \\ \quad \operatorname{pos\text{-}rr}(r) \wedge \\ \quad \forall t \in w.\; (((x < t) \wedge (t < (x + r))) \Rightarrow (t \in v)) \wedge \\ \quad \forall t \in w.\; (((x < t) \wedge (t < (x + r))) \Rightarrow (f(t) \simeq g(t))) \Rightarrow \\ \quad \quad \operatorname{is\text{-}right\text{-}limit\text{-}within}(f, w, x, l) \end{array}}}
\vnbprop{left-limit-within-local}{\vnbstmt{\begin{array}{@{}l@{}} \forall g, v, f, w, x, l, r. \\ \quad \operatorname{is\text{-}left\text{-}limit\text{-}within}(g, v, x, l) \wedge \\ \quad w \subseteq \mathbb{R} \wedge \\ \quad f \in (w \to \mathbb{R}) \wedge \\ \quad \operatorname{pos\text{-}rr}(r) \wedge \\ \quad \forall t \in w.\; ((((x - r) < t) \wedge (t < x)) \Rightarrow (t \in v)) \wedge \\ \quad \forall t \in w.\; ((((x - r) < t) \wedge (t < x)) \Rightarrow (f(t) \simeq g(t))) \Rightarrow \\ \quad \quad \operatorname{is\text{-}left\text{-}limit\text{-}within}(f, w, x, l) \end{array}}}
\vnbprop{regulated-on-glue}{\vnbstmt{\begin{array}{@{}l@{}} \forall a, c, b \in \mathbb{R}, d, g, f. \\ \quad a < c \wedge \\ \quad c < b \wedge \\ \quad \operatorname{is\text{-}regulated\text{-}on}(d, a, c) \wedge \\ \quad \operatorname{is\text{-}regulated\text{-}on}(g, c, b) \wedge \\ \quad f \in (\operatorname{ccint}(a, b) \to \mathbb{R}) \wedge \\ \quad \forall t \in \operatorname{ccint}(a, c).\; ((t < c) \Rightarrow (f(t) \simeq d(t))) \wedge \\ \quad \forall t \in \operatorname{ccint}(c, b).\; ((c < t) \Rightarrow (f(t) \simeq g(t))) \Rightarrow \\ \quad \quad \operatorname{is\text{-}regulated\text{-}on}(f, a, b) \end{array}}}

\vnbbookfile{theorem-library/calint-regulated-step}{theorem-library/calint-regulated-step}
\noindent{\small calint-regulated-step.scm -- WAVE 2 (2026-10-27..30), wired 2026-10-30 from scratchpad/calint/calint-regulated-step.scm. Library: analysis.}\par
\vnbprop{constant-piece-regulated-on}{\vnbstmt{\begin{array}{@{}l@{}} \forall u, v \in \mathbb{R}, f \in (\mathbb{R} \to \mathbb{R}), a \in \mathbb{R}. \\ \quad ((u < v) \wedge \forall t \in \mathbb{R}.\; (((u < t) \wedge (t < v)) \Rightarrow (f(t) = a))) \Rightarrow  \\ \quad \quad \operatorname{is\text{-}regulated\text{-}on}(\operatorname{restrict}(f, \operatorname{ccint}(u, v)), u, v) \end{array}}}
\vnbprop{step-fn-regulated-on-ind}{\vnbstmt{\begin{array}{@{}l@{}} \forall n \in \mathbb{N}, p, a, b.; \forall f \in (\mathbb{R} \to \mathbb{R}). \\ \quad \operatorname{is\text{-}partition}(p, n, a, b) \wedge \\ \quad \forall i \in \operatorname{interval}(0, n). \\ \quad \quad (i < n) \Rightarrow  \\ \quad \quad \quad \exists c \in \mathbb{R}. \forall t \in \mathbb{R}. \\ \quad \quad \quad \quad ((p(i) < t) \wedge (t < p((i + 1)))) \Rightarrow (f(t) = c) \Rightarrow \\ \quad \quad \operatorname{is\text{-}regulated\text{-}on}(\operatorname{restrict}(f, \operatorname{ccint}(a, b)), a, b) \end{array}}}
\vnbprop{step-fn-regulated-on}{\vnbstmt{\begin{array}{@{}l@{}} \forall a, b \in \mathbb{R}, f. \\ \quad \operatorname{is\text{-}step\text{-}fn}(f, a, b) \Rightarrow \operatorname{is\text{-}regulated\text{-}on}(\operatorname{restrict}(f, \operatorname{ccint}(a, b)), a, b) \end{array}}}
\vnbprop{restrict-regulated-on-is-regulated}{\vnbstmt{\begin{array}{@{}l@{}} \forall a, b.; \forall h \in (\mathbb{R} \to \mathbb{R}). \\ \quad \operatorname{is\text{-}regulated\text{-}on}(\operatorname{restrict}(h, \operatorname{ccint}(a, b)), a, b) \Rightarrow \operatorname{is\text{-}regulated}(h, a, b) \end{array}}}
\vnbprop{regulated-extend-const-is-regulated-on}{\vnbstmt{\begin{array}{@{}l@{}} \forall a, b \in \mathbb{R}, f \in (\operatorname{ccint}(a, b) \to \mathbb{R}). \\ \quad \operatorname{is\text{-}regulated}(\operatorname{extend\text{-}const}(f, a, b), a, b) \Rightarrow \operatorname{is\text{-}regulated\text{-}on}(f, a, b) \end{array}}}
\vnbprop{step-fn-is-regulated}{\vnbstmt{\forall a, b \in \mathbb{R}, f.\; (\operatorname{is\text{-}step\text{-}fn}(f, a, b) \Rightarrow \operatorname{is\text{-}regulated}(f, a, b))}}
\vnbprop{regulated-on-iff-step-limit}{\vnbstmt{\begin{array}{@{}l@{}} \forall a, b \in \mathbb{R}, f. \\ \quad ((a < b) \wedge (f \in (\operatorname{ccint}(a, b) \to \mathbb{R}))) \Rightarrow  \\ \quad \quad \operatorname{is\text{-}regulated\text{-}on}(f, a, b) \\ \quad \quad \Leftrightarrow\; \exists m \in (\mathbb{N} \to (\mathbb{R} \to \mathbb{R})). \\ \quad \quad \quad (\forall k \in \mathbb{N}.\; \operatorname{is\text{-}step\text{-}fn}(m(k), a, b) \wedge \operatorname{is\text{-}unif\text{-}limit\text{-}on}(m, \operatorname{extend\text{-}const}(f, a, b), a, b)) \end{array}}}

\vnbbookfile{theorem-library/aif-scale-small}{theorem-library/aif-scale-small}
\noindent{\small aif-scale-small.scm -- lemmas MOVED 2026-11-10 (integration of the 11-09/10 waves) out of theorem-library/analytic-from-integral.scm into library analysis, for the library gate: cited by ccr4d-converge; RR arithmetic. The proofs are the originals, verbatim; the helper defines are a copy of the source's. aif-scale-small}\par
\vnbprop{aif-scale-small}{\vnbstmt{\begin{array}{@{}l@{}} \forall l \in \mathbb{R}, s. \\ \quad ((0 \leq l) \wedge \operatorname{pos\text{-}rr}(s)) \Rightarrow  \\ \quad \quad \operatorname{pos\text{-}rr}((s \cdot \operatorname{recip}((l + 1)))) \wedge ((l \cdot (s \cdot \operatorname{recip}((l + 1)))) \leq s) \end{array}}}

\vnbbookfile{theorem-library/ccr4d-converge}{theorem-library/ccr4d-converge}
\noindent{\small ccr4d-converge.scm -- wave of 2026-11-09/10 (agent CCR-1109-C5), wired 2026-11-10 from scratchpad/ccr1109c5/ccr4d-converge.scm. Library: analysis.}\par
\vnbprop{nf-converges-add}{\vnbstmt{\begin{array}{@{}l@{}} \forall k.; \forall f, g \in (\mathbb{N} \to \operatorname{carr}(k)), a, b \in \operatorname{carr}(k), h \in (\mathbb{N} \to \operatorname{carr}(k)). \\ \quad \operatorname{is\text{-}normed\text{-}field}(k) \wedge \\ \quad \operatorname{converges\text{-}to}(\operatorname{nf\text{-}metric\text{-}space}(k), f, a) \wedge \\ \quad \operatorname{converges\text{-}to}(\operatorname{nf\text{-}metric\text{-}space}(k), g, b) \wedge \\ \quad \forall n \in \mathbb{N}.\; (h(n) = \operatorname{add}(k)(f(n), g(n))) \Rightarrow \\ \quad \quad \operatorname{converges\text{-}to}(\operatorname{nf\text{-}metric\text{-}space}(k), h, \operatorname{add}(k)(a, b)) \end{array}}}
\vnbprop{nf-converges-scale}{\vnbstmt{\begin{array}{@{}l@{}} \forall k.; \forall c \in \operatorname{carr}(k), f \in (\mathbb{N} \to \operatorname{carr}(k)), a \in \operatorname{carr}(k), h \in (\mathbb{N} \to \operatorname{carr}(k)). \\ \quad \operatorname{is\text{-}normed\text{-}field}(k) \wedge \\ \quad \operatorname{converges\text{-}to}(\operatorname{nf\text{-}metric\text{-}space}(k), f, a) \wedge \\ \quad \forall n \in \mathbb{N}.\; (h(n) = \operatorname{mul}(k)(c, f(n))) \Rightarrow \\ \quad \quad \operatorname{converges\text{-}to}(\operatorname{nf\text{-}metric\text{-}space}(k), h, \operatorname{mul}(k)(c, a)) \end{array}}}
\vnbprop{ord-segment-succ-sub}{\vnbstmt{\forall d \in \mathbb{N}, i \in \mathrm{OS}(d).\; (i \in \mathrm{OS}((d + 1)))}}
\vnbprop{ord-segment-succ-top}{\vnbstmt{\forall d \in \mathbb{N}.\; (d \in \mathrm{OS}((d + 1)))}}
\vnbprop{nf-converges-const}{\vnbstmt{\begin{array}{@{}l@{}} \forall k.; \forall c \in \operatorname{carr}(k), h \in (\mathbb{N} \to \operatorname{carr}(k)). \\ \quad (\operatorname{is\text{-}normed\text{-}field}(k) \wedge \forall n \in \mathbb{N}.\; (h(n) = c)) \Rightarrow  \\ \quad \quad \operatorname{converges\text{-}to}(\operatorname{nf\text{-}metric\text{-}space}(k), h, c) \end{array}}}
\vnbprop{nf-converges-finsum}{\vnbstmt{\begin{array}{@{}l@{}} \forall d \in \mathbb{N}, k, f, l, t. \\ \quad \operatorname{is\text{-}normed\text{-}field}(k) \wedge \\ \quad \forall m \in \mathbb{N}.\; (t(m) \in \mathbf{Set}) \wedge \\ \quad \forall i \in \mathrm{OS}(d).\; (f(i) \in (\mathbb{N} \to \operatorname{carr}(k))) \wedge \\ \quad \forall i \in \mathrm{OS}(d).\; (l(i) \in \operatorname{carr}(k)) \wedge \\ \quad \forall i \in \mathrm{OS}(d).\; \operatorname{converges\text{-}to}(\operatorname{nf\text{-}metric\text{-}space}(k), f(i), l(i)) \wedge \\ \quad \forall m \in \mathbb{N}, i \in \mathrm{OS}(d).\; (t(m)(i) \simeq f(i)(m)) \Rightarrow \\ \quad \quad \operatorname{converges\text{-}to}(\operatorname{nf\text{-}metric\text{-}space}(k), \\ \quad \quad \quad (\lambda \mathit{cvm} \in \mathbb{N}.\; \operatorname{sum\text{-}ag}(\operatorname{normed\text{-}field\text{-}additive\text{-}ag}(k), t(\mathit{cvm}), d)), \\ \quad \quad \quad \operatorname{sum\text{-}ag}(\operatorname{normed\text{-}field\text{-}additive\text{-}ag}(k), l, d)) \end{array}}}
\vnbprop{nf-converges-ptwise-eq}{\vnbstmt{\begin{array}{@{}l@{}} \forall k.; \forall f \in (\mathbb{N} \to \operatorname{carr}(k)), a \in \operatorname{carr}(k), h \in (\mathbb{N} \to \operatorname{carr}(k)). \\ \quad (\operatorname{is\text{-}normed\text{-}field}(k) \wedge (\operatorname{converges\text{-}to}(\operatorname{nf\text{-}metric\text{-}space}(k), f, a) \wedge \forall n \in \mathbb{N}.\; (h(n) = f(n)))) \Rightarrow  \\ \quad \quad \operatorname{converges\text{-}to}(\operatorname{nf\text{-}metric\text{-}space}(k), h, a) \end{array}}}
\vnbprop{nf-series-finsum-swap}{\vnbstmt{\begin{array}{@{}l@{}} \forall d \in \mathbb{N}, k, a, c, r, f, t, l. \\ \quad \operatorname{is\text{-}normed\text{-}field}(k) \wedge \\ \quad \forall x \in \mathbb{N}.\; (a(x) \in (\mathrm{OS}(d) \to \operatorname{carr}(k))) \wedge \\ \quad \forall i \in \mathrm{OS}(d).\; (c(i) \in (\mathbb{N} \to \operatorname{carr}(k))) \wedge \\ \quad \forall x \in \mathbb{N}, i \in \mathrm{OS}(d).\; (c(i)(x) \simeq a(x)(i)) \wedge \\ \quad \forall x \in \mathbb{N}.\; (r(x) \simeq \operatorname{sum\text{-}ag}(\operatorname{normed\text{-}field\text{-}additive\text{-}ag}(k), a(x), d)) \wedge \\ \quad \forall i \in \mathrm{OS}(d).\; (f(i) \in (\mathbb{N} \to \operatorname{carr}(k))) \wedge \\ \quad \forall i \in \mathrm{OS}(d), m \in \mathbb{N}. \\ \quad \quad (f(i)(m) \simeq \operatorname{sum\text{-}ag}(\operatorname{normed\text{-}field\text{-}additive\text{-}ag}(k), c(i), m)) \wedge \\ \quad \forall m \in \mathbb{N}.\; (t(m) \in (\mathrm{OS}(d) \to \operatorname{carr}(k))) \wedge \\ \quad \forall m \in \mathbb{N}, i \in \mathrm{OS}(d). \\ \quad \quad (t(m)(i) \simeq \operatorname{sum\text{-}ag}(\operatorname{normed\text{-}field\text{-}additive\text{-}ag}(k), c(i), m)) \wedge \\ \quad \forall i \in \mathrm{OS}(d).\; (l(i) \in \operatorname{carr}(k)) \wedge \\ \quad \forall i \in \mathrm{OS}(d).\; \operatorname{converges\text{-}to}(\operatorname{nf\text{-}metric\text{-}space}(k), f(i), l(i)) \Rightarrow \\ \quad \quad \operatorname{converges\text{-}to}(\operatorname{nf\text{-}metric\text{-}space}(k), \\ \quad \quad \quad (\lambda \mathit{cvm} \in \mathbb{N}.\; \operatorname{sum\text{-}ag}(\operatorname{normed\text{-}field\text{-}additive\text{-}ag}(k), r, \mathit{cvm})), \\ \quad \quad \quad \operatorname{sum\text{-}ag}(\operatorname{normed\text{-}field\text{-}additive\text{-}ag}(k), l, d)) \end{array}}}

\vnbbookfile{theorem-library/cc-magnitude-i}{theorem-library/cc-magnitude-i}
\noindent{\small cc-magnitude-i.scm -- lemmas MOVED 2026-11-10 (integration of the 11-09/10 waves) out of theorem-library/cauchy-estimates.scm into library analysis, for the library gate: cited by cnvs-a; CC arithmetic. The proofs are the originals, verbatim; the helper defines are a copy of the source's. cc-magnitude-i}\par
\vnbprop{cc-magnitude-i}{\vnbstmt{\operatorname{magnitude}(+i) = 1}}

\vnbbookfile{theorem-library/cnvs-a}{theorem-library/cnvs-a}
\noindent{\small cnvs-a.scm -- wave of 2026-11-09/10 (agent CNVS), wired 2026-11-10 from scratchpad/cnvs-1109/cnvs-a.scm. Library: analysis.}\par
\vnbprop{cnvs-vec-set}{\vnbstmt{\forall v.\; (\operatorname{is\text{-}complex\text{-}normed\text{-}vector\text{-}space}(v) \Rightarrow (\operatorname{vec}(v) \in \mathbf{Set}))}}
\vnbprop{cnvs-jop-fun}{\vnbstmt{\begin{array}{@{}l@{}} \forall v. \\ \quad \operatorname{is\text{-}complex\text{-}normed\text{-}vector\text{-}space}(v) \Rightarrow (J_{v} \in (\operatorname{vec}(v) \to \operatorname{vec}(v))) \end{array}}}
\vnbprop{cnvs-vnrm-fun}{\vnbstmt{\begin{array}{@{}l@{}} \forall v. \\ \quad \operatorname{is\text{-}complex\text{-}normed\text{-}vector\text{-}space}(v) \Rightarrow (\operatorname{vnrm}(v) \in (\operatorname{vec}(v) \to \mathbb{R})) \end{array}}}
\vnbprop{cnvs-nvs}{\vnbstmt{\begin{array}{@{}l@{}} \forall v. \\ \quad \operatorname{is\text{-}complex\text{-}normed\text{-}vector\text{-}space}(v) \Rightarrow  \\ \quad \quad \operatorname{is\text{-}normed\text{-}vector\text{-}space}({[}\operatorname{scal}(v), \operatorname{vec}(v), \operatorname{vadd}(v), \operatorname{vzero}(v), \operatorname{vneg}(v), \operatorname{act}(v), \operatorname{vnrm}(v)]) \end{array}}}
\vnbprop{cnvs-real-linear}{\vnbstmt{\begin{array}{@{}l@{}} \forall v. \\ \quad \operatorname{is\text{-}complex\text{-}normed\text{-}vector\text{-}space}(v) \Rightarrow \operatorname{is\text{-}real\text{-}linear\text{-}op}(v, J_{v}) \end{array}}}
\vnbprop{cnvs-square}{\vnbstmt{\begin{array}{@{}l@{}} \forall v. \\ \quad \operatorname{is\text{-}complex\text{-}normed\text{-}vector\text{-}space}(v) \Rightarrow \operatorname{is\text{-}square\text{-}minus\text{-}id}(v, J_{v}) \end{array}}}
\vnbprop{cnvs-vnrm-cact}{\vnbstmt{\begin{array}{@{}l@{}} \forall v.; \forall c \in \mathbb{C}, x \in \operatorname{vec}(v). \\ \quad \operatorname{is\text{-}complex\text{-}normed\text{-}vector\text{-}space}(v) \Rightarrow  \\ \quad \quad \operatorname{vnrm}(v)(\operatorname{cact}(v)(c, x)) = (\operatorname{magnitude}(c) \cdot \operatorname{vnrm}(v)(x)) \end{array}}}
\vnbprop{cbs-is-cnvs}{\vnbstmt{\begin{array}{@{}l@{}} \forall v. \\ \quad \operatorname{is\text{-}complex\text{-}banach\text{-}space}(v) \Rightarrow \operatorname{is\text{-}complex\text{-}normed\text{-}vector\text{-}space}(v) \end{array}}}
\vnbprop{cbs-complete}{\vnbstmt{\begin{array}{@{}l@{}} \forall v. \\ \quad \operatorname{is\text{-}complex\text{-}banach\text{-}space}(v) \Rightarrow \operatorname{is\text{-}complete}(\operatorname{cnvs\text{-}metric\text{-}space}(v)) \end{array}}}
\vnbprop{cnvs-nvs-view-scal}{\vnbstmt{\begin{array}{@{}l@{}} \forall v. \\ \quad \operatorname{is\text{-}complex\text{-}normed\text{-}vector\text{-}space}(v) \Rightarrow (\operatorname{scal}(\operatorname{cnvs\text{-}as\text{-}nvs}(v)) \simeq \operatorname{scal}(v)) \end{array}}}
\vnbprop{cnvs-nvs-view-scal-pin}{\vnbstmt{\begin{array}{@{}l@{}} \forall v. \\ \quad \operatorname{is\text{-}complex\text{-}normed\text{-}vector\text{-}space}(v) \Rightarrow  \\ \quad \quad \operatorname{scal}(\operatorname{cnvs\text{-}as\text{-}nvs}(v)) = \operatorname{normed\text{-}field\text{-}as\text{-}commutative\text{-}ring}(\operatorname{rr\text{-}normed\text{-}field}) \end{array}}}
\vnbprop{cnvs-scal-pin}{\vnbstmt{\begin{array}{@{}l@{}} \forall v. \\ \quad \operatorname{is\text{-}complex\text{-}normed\text{-}vector\text{-}space}(v) \Rightarrow  \\ \quad \quad \operatorname{scal}(v) = \operatorname{normed\text{-}field\text{-}as\text{-}commutative\text{-}ring}(\operatorname{rr\text{-}normed\text{-}field}) \end{array}}}
\vnbprop{cnvs-jop-add}{\vnbstmt{\begin{array}{@{}l@{}} \forall v.; \forall x, y \in \operatorname{vec}(v). \\ \quad \operatorname{is\text{-}complex\text{-}normed\text{-}vector\text{-}space}(v) \Rightarrow  \\ \quad \quad J_{v}(\operatorname{vadd}(v)(x, y)) = \operatorname{vadd}(v)(J_{v}(x), J_{v}(y)) \end{array}}}
\vnbprop{cnvs-jop-act}{\vnbstmt{\begin{array}{@{}l@{}} \forall v.; \forall a \in \mathbb{R}, x \in \operatorname{vec}(v). \\ \quad \operatorname{is\text{-}complex\text{-}normed\text{-}vector\text{-}space}(v) \Rightarrow  \\ \quad \quad J_{v}(\operatorname{act}(v)(a, x)) = \operatorname{act}(v)(a, J_{v}(x)) \end{array}}}
\vnbprop{cnvs-jop-sq}{\vnbstmt{\begin{array}{@{}l@{}} \forall v.; \forall x \in \operatorname{vec}(v). \\ \quad \operatorname{is\text{-}complex\text{-}normed\text{-}vector\text{-}space}(v) \Rightarrow (J_{v}(J_{v}(x)) = \operatorname{vneg}(v)(x)) \end{array}}}
\vnbprop{cnvs-jop-type}{\vnbstmt{\begin{array}{@{}l@{}} \forall v.; \forall x \in \operatorname{vec}(v). \\ \quad \operatorname{is\text{-}complex\text{-}normed\text{-}vector\text{-}space}(v) \Rightarrow (J_{v}(x) \in \operatorname{vec}(v)) \end{array}}}
\vnbprop{cnvs-nf-add}{\vnbstmt{\begin{array}{@{}l@{}} \forall v, p, q, r, s, x. \\ \quad (\operatorname{is\text{-}complex\text{-}normed\text{-}vector\text{-}space}(v) \wedge ((p \in \mathbb{R}) \wedge ((q \in \mathbb{R}) \wedge ((r \in \mathbb{R}) \wedge ((s \in \mathbb{R}) \wedge (x \in \operatorname{vec}(v))))))) \Rightarrow  \\ \quad \quad \operatorname{vadd}(v)(\operatorname{vadd}(v)(\operatorname{act}(v)(p, x), \operatorname{act}(v)(q, J_{v}(x))), \\ \quad \quad \quad \operatorname{vadd}(v)(\operatorname{act}(v)(r, x), \operatorname{act}(v)(s, J_{v}(x)))) \\ \quad \quad =\; \operatorname{vadd}(v)(\operatorname{act}(v)((p + r), x), \operatorname{act}(v)((q + s), J_{v}(x))) \end{array}}}
\vnbprop{cnvs-nf-scale}{\vnbstmt{\begin{array}{@{}l@{}} \forall v, t, p, q, x. \\ \quad (\operatorname{is\text{-}complex\text{-}normed\text{-}vector\text{-}space}(v) \wedge ((t \in \mathbb{R}) \wedge ((p \in \mathbb{R}) \wedge ((q \in \mathbb{R}) \wedge (x \in \operatorname{vec}(v)))))) \Rightarrow  \\ \quad \quad \operatorname{act}(v)(t, \operatorname{vadd}(v)(\operatorname{act}(v)(p, x), \operatorname{act}(v)(q, J_{v}(x)))) \\ \quad \quad =\; \operatorname{vadd}(v)(\operatorname{act}(v)((t \cdot p), x), \operatorname{act}(v)((t \cdot q), J_{v}(x))) \end{array}}}
\vnbprop{cnvs-act-neg-scalar}{\vnbstmt{\begin{array}{@{}l@{}} \forall v, q, x. \\ \quad (\operatorname{is\text{-}complex\text{-}normed\text{-}vector\text{-}space}(v) \wedge ((q \in \mathbb{R}) \wedge (x \in \operatorname{vec}(v)))) \Rightarrow  \\ \quad \quad \operatorname{act}(v)(-q, x) = \operatorname{vneg}(v)(\operatorname{act}(v)(q, x)) \end{array}}}
\vnbprop{cnvs-nf-jop}{\vnbstmt{\begin{array}{@{}l@{}} \forall v, p, q, x. \\ \quad (\operatorname{is\text{-}complex\text{-}normed\text{-}vector\text{-}space}(v) \wedge ((p \in \mathbb{R}) \wedge ((q \in \mathbb{R}) \wedge (x \in \operatorname{vec}(v))))) \Rightarrow  \\ \quad \quad J_{v}(\operatorname{vadd}(v)(\operatorname{act}(v)(p, x), \operatorname{act}(v)(q, J_{v}(x)))) \\ \quad \quad =\; \operatorname{vadd}(v)(\operatorname{act}(v)(-q, x), \operatorname{act}(v)(p, J_{v}(x))) \end{array}}}
\vnbprop{cnvs-cact-value}{\vnbstmt{\begin{array}{@{}l@{}} \forall v.; \forall z \in \mathbb{C}, x \in \operatorname{vec}(v). \\ \quad \operatorname{is\text{-}complex\text{-}normed\text{-}vector\text{-}space}(v) \Rightarrow  \\ \quad \quad \operatorname{cact}(v)(z, x) = \operatorname{vadd}(v)(\operatorname{act}(v)(\operatorname{real\text{-}part}(z), x), \operatorname{act}(v)(\operatorname{imag\text{-}part}(z), J_{v}(x))) \end{array}}}
\vnbprop{cnvs-cact-type-pt}{\vnbstmt{\begin{array}{@{}l@{}} \forall v.; \forall z \in \mathbb{C}, x \in \operatorname{vec}(v). \\ \quad \operatorname{is\text{-}complex\text{-}normed\text{-}vector\text{-}space}(v) \Rightarrow (\operatorname{cact}(v)(z, x) \in \operatorname{vec}(v)) \end{array}}}
\vnbprop{cnvs-cact-fun}{\vnbstmt{\begin{array}{@{}l@{}} \forall v. \\ \quad \operatorname{is\text{-}complex\text{-}normed\text{-}vector\text{-}space}(v) \Rightarrow  \\ \quad \quad \operatorname{cact}(v) \in ((\mathbb{C} \times \operatorname{vec}(v)) \to \operatorname{vec}(v)) \end{array}}}
\vnbprop{cnvs-cact-vadd}{\vnbstmt{\begin{array}{@{}l@{}} \forall v.; \forall z \in \mathbb{C}, x, y \in \operatorname{vec}(v). \\ \quad \operatorname{is\text{-}complex\text{-}normed\text{-}vector\text{-}space}(v) \Rightarrow  \\ \quad \quad \operatorname{cact}(v)(z, \operatorname{vadd}(v)(x, y)) = \operatorname{vadd}(v)(\operatorname{cact}(v)(z, x), \operatorname{cact}(v)(z, y)) \end{array}}}
\vnbprop{cnvs-cact-add}{\vnbstmt{\begin{array}{@{}l@{}} \forall v.; \forall z, w \in \mathbb{C}, x \in \operatorname{vec}(v). \\ \quad \operatorname{is\text{-}complex\text{-}normed\text{-}vector\text{-}space}(v) \Rightarrow  \\ \quad \quad \operatorname{cact}(v)((z + w), x) = \operatorname{vadd}(v)(\operatorname{cact}(v)(z, x), \operatorname{cact}(v)(w, x)) \end{array}}}
\vnbprop{cnvs-cact-mul}{\vnbstmt{\begin{array}{@{}l@{}} \forall v.; \forall z, w \in \mathbb{C}, x \in \operatorname{vec}(v). \\ \quad \operatorname{is\text{-}complex\text{-}normed\text{-}vector\text{-}space}(v) \Rightarrow  \\ \quad \quad \operatorname{cact}(v)((z \cdot w), x) = \operatorname{cact}(v)(z, \operatorname{cact}(v)(w, x)) \end{array}}}
\vnbprop{cnvs-cact-one}{\vnbstmt{\begin{array}{@{}l@{}} \forall v.; \forall x \in \operatorname{vec}(v). \\ \quad \operatorname{is\text{-}complex\text{-}normed\text{-}vector\text{-}space}(v) \Rightarrow (\operatorname{cact}(v)(1, x) = x) \end{array}}}
\vnbprop{cnvs-cact-real}{\vnbstmt{\begin{array}{@{}l@{}} \forall v.; \forall r \in \mathbb{R}, x \in \operatorname{vec}(v). \\ \quad \operatorname{is\text{-}complex\text{-}normed\text{-}vector\text{-}space}(v) \Rightarrow (\operatorname{cact}(v)(r, x) = \operatorname{act}(v)(r, x)) \end{array}}}
\vnbprop{cnvs-cact-i}{\vnbstmt{\begin{array}{@{}l@{}} \forall v.; \forall x \in \operatorname{vec}(v). \\ \quad \operatorname{is\text{-}complex\text{-}normed\text{-}vector\text{-}space}(v) \Rightarrow (\operatorname{cact}(v)(+i, x) = J_{v}(x)) \end{array}}}
\vnbprop{cnvs-jop-isometric}{\vnbstmt{\begin{array}{@{}l@{}} \forall v.; \forall x \in \operatorname{vec}(v). \\ \quad \operatorname{is\text{-}complex\text{-}normed\text{-}vector\text{-}space}(v) \Rightarrow (\operatorname{vnrm}(v)(J_{v}(x)) = \operatorname{vnrm}(v)(x)) \end{array}}}

\vnbbookfile{theorem-library/cnvs-b}{theorem-library/cnvs-b}
\noindent{\small cnvs-b.scm -- wave of 2026-11-09/10 (agent CNVS), wired 2026-11-10 from scratchpad/cnvs-1109/cnvs-b.scm. Library: analysis.}\par
\vnbprop{cmod-norm-fun}{\vnbstmt{\forall m, n.\; (\operatorname{is\text{-}cmod\text{-}norm}(m, n) \Rightarrow (n \in (\operatorname{vec}(m) \to \mathbb{R})))}}
\vnbprop{cmod-norm-nonneg}{\vnbstmt{\forall m, n.; \forall x \in \operatorname{vec}(m).\; (\operatorname{is\text{-}cmod\text{-}norm}(m, n) \Rightarrow (0 \leq n(x)))}}
\vnbprop{cmod-norm-zero-iff}{\vnbstmt{\begin{array}{@{}l@{}} \forall m, n.; \forall x \in \operatorname{vec}(m). \\ \quad \operatorname{is\text{-}cmod\text{-}norm}(m, n) \Rightarrow ((n(x) = 0) \Leftrightarrow (x = \operatorname{vzero}(m))) \end{array}}}
\vnbprop{cmod-norm-homog}{\vnbstmt{\begin{array}{@{}l@{}} \forall m, n.; \forall c \in \mathbb{C}, x \in \operatorname{vec}(m). \\ \quad \operatorname{is\text{-}cmod\text{-}norm}(m, n) \Rightarrow (n(\operatorname{act}(m)(c, x)) = (\operatorname{magnitude}(c) \cdot n(x))) \end{array}}}
\vnbprop{cmod-norm-triangle}{\vnbstmt{\begin{array}{@{}l@{}} \forall m, n.; \forall x, y \in \operatorname{vec}(m). \\ \quad \operatorname{is\text{-}cmod\text{-}norm}(m, n) \Rightarrow (n(\operatorname{vadd}(m)(x, y)) \leq (n(x) + n(y))) \end{array}}}
\vnbprop{cmod-vadd-fun}{\vnbstmt{\begin{array}{@{}l@{}} \forall m. \\ \quad \operatorname{is\text{-}module}(m) \Rightarrow (\operatorname{vadd}(m) \in ((\operatorname{vec}(m) \times \operatorname{vec}(m)) \to \operatorname{vec}(m))) \end{array}}}
\vnbprop{cmod-vneg-fun}{\vnbstmt{\forall m.\; (\operatorname{is\text{-}module}(m) \Rightarrow (\operatorname{vneg}(m) \in (\operatorname{vec}(m) \to \operatorname{vec}(m))))}}
\vnbprop{cmod-assoc-prop}{\vnbstmt{\forall m.\; (\operatorname{is\text{-}module}(m) \Rightarrow \operatorname{is\text{-}associative}(\operatorname{vadd}(m), \operatorname{vec}(m)))}}
\vnbprop{cmod-comm-prop}{\vnbstmt{\forall m.\; (\operatorname{is\text{-}module}(m) \Rightarrow \operatorname{is\text{-}commutative}(\operatorname{vadd}(m), \operatorname{vec}(m)))}}
\vnbprop{cmod-identity-prop}{\vnbstmt{\forall m.\; (\operatorname{is\text{-}module}(m) \Rightarrow \operatorname{is\text{-}identity}(\operatorname{vadd}(m), \operatorname{vzero}(m), \operatorname{vec}(m)))}}
\vnbprop{cmod-inverses-prop}{\vnbstmt{\begin{array}{@{}l@{}} \forall m. \\ \quad \operatorname{is\text{-}module}(m) \Rightarrow \operatorname{has\text{-}inverses}(\operatorname{vadd}(m), \operatorname{vzero}(m), \operatorname{vneg}(m), \operatorname{vec}(m)) \end{array}}}
\vnbprop{cmod-carr}{\vnbstmt{\begin{array}{@{}l@{}} \forall m. \\ \quad (\operatorname{is\text{-}module}(m) \wedge (\operatorname{scal}(m) = \operatorname{normed\text{-}field\text{-}as\text{-}commutative\text{-}ring}(\operatorname{cc\text{-}normed\text{-}field}))) \Rightarrow  \\ \quad \quad \operatorname{carr}(\operatorname{scal}(m)) = \mathbb{C} \end{array}}}
\vnbprop{cmod-act-type-cc}{\vnbstmt{\begin{array}{@{}l@{}} \forall m.; \forall c \in \mathbb{C}, x \in \operatorname{vec}(m). \\ \quad (\operatorname{is\text{-}module}(m) \wedge (\operatorname{scal}(m) = \operatorname{normed\text{-}field\text{-}as\text{-}commutative\text{-}ring}(\operatorname{cc\text{-}normed\text{-}field}))) \Rightarrow  \\ \quad \quad \operatorname{act}(m)(c, x) \in \operatorname{vec}(m) \end{array}}}
\vnbprop{cmod-act-vadd-cc}{\vnbstmt{\begin{array}{@{}l@{}} \forall m.; \forall c \in \mathbb{C}, x, y \in \operatorname{vec}(m). \\ \quad (\operatorname{is\text{-}module}(m) \wedge (\operatorname{scal}(m) = \operatorname{normed\text{-}field\text{-}as\text{-}commutative\text{-}ring}(\operatorname{cc\text{-}normed\text{-}field}))) \Rightarrow  \\ \quad \quad \operatorname{act}(m)(c, \operatorname{vadd}(m)(x, y)) = \operatorname{vadd}(m)(\operatorname{act}(m)(c, x), \operatorname{act}(m)(c, y)) \end{array}}}
\vnbprop{cmod-act-add-cc}{\vnbstmt{\begin{array}{@{}l@{}} \forall m.; \forall c, d \in \mathbb{C}, x \in \operatorname{vec}(m). \\ \quad (\operatorname{is\text{-}module}(m) \wedge (\operatorname{scal}(m) = \operatorname{normed\text{-}field\text{-}as\text{-}commutative\text{-}ring}(\operatorname{cc\text{-}normed\text{-}field}))) \Rightarrow  \\ \quad \quad \operatorname{act}(m)((c + d), x) = \operatorname{vadd}(m)(\operatorname{act}(m)(c, x), \operatorname{act}(m)(d, x)) \end{array}}}
\vnbprop{cmod-act-mul-cc}{\vnbstmt{\begin{array}{@{}l@{}} \forall m.; \forall c, d \in \mathbb{C}, x \in \operatorname{vec}(m). \\ \quad (\operatorname{is\text{-}module}(m) \wedge (\operatorname{scal}(m) = \operatorname{normed\text{-}field\text{-}as\text{-}commutative\text{-}ring}(\operatorname{cc\text{-}normed\text{-}field}))) \Rightarrow  \\ \quad \quad \operatorname{act}(m)((c \cdot d), x) = \operatorname{act}(m)(c, \operatorname{act}(m)(d, x)) \end{array}}}
\vnbprop{cmod-act-one}{\vnbstmt{\begin{array}{@{}l@{}} \forall m.; \forall x \in \operatorname{vec}(m). \\ \quad (\operatorname{is\text{-}module}(m) \wedge (\operatorname{scal}(m) = \operatorname{normed\text{-}field\text{-}as\text{-}commutative\text{-}ring}(\operatorname{cc\text{-}normed\text{-}field}))) \Rightarrow  \\ \quad \quad \operatorname{act}(m)(1, x) = x \end{array}}}
\vnbprop{cmod-act-neg-one}{\vnbstmt{\begin{array}{@{}l@{}} \forall m.; \forall x \in \operatorname{vec}(m). \\ \quad (\operatorname{is\text{-}module}(m) \wedge (\operatorname{scal}(m) = \operatorname{normed\text{-}field\text{-}as\text{-}commutative\text{-}ring}(\operatorname{cc\text{-}normed\text{-}field}))) \Rightarrow  \\ \quad \quad \operatorname{act}(m)(-1, x) = \operatorname{vneg}(m)(x) \end{array}}}
\vnbprop{cmod-j-add}{\vnbstmt{\begin{array}{@{}l@{}} \forall m.; \forall x, y \in \operatorname{vec}(m). \\ \quad (\operatorname{is\text{-}module}(m) \wedge (\operatorname{scal}(m) = \operatorname{normed\text{-}field\text{-}as\text{-}commutative\text{-}ring}(\operatorname{cc\text{-}normed\text{-}field}))) \Rightarrow  \\ \quad \quad \operatorname{act}(m)(+i, \operatorname{vadd}(m)(x, y)) = \operatorname{vadd}(m)(\operatorname{act}(m)(+i, x), \operatorname{act}(m)(+i, y)) \end{array}}}
\vnbprop{cmod-j-act}{\vnbstmt{\begin{array}{@{}l@{}} \forall m.; \forall r \in \mathbb{R}, x \in \operatorname{vec}(m). \\ \quad (\operatorname{is\text{-}module}(m) \wedge (\operatorname{scal}(m) = \operatorname{normed\text{-}field\text{-}as\text{-}commutative\text{-}ring}(\operatorname{cc\text{-}normed\text{-}field}))) \Rightarrow  \\ \quad \quad \operatorname{act}(m)(+i, \operatorname{act}(m)(r, x)) = \operatorname{act}(m)(r, \operatorname{act}(m)(+i, x)) \end{array}}}
\vnbprop{cmod-j-sq}{\vnbstmt{\begin{array}{@{}l@{}} \forall m.; \forall x \in \operatorname{vec}(m). \\ \quad (\operatorname{is\text{-}module}(m) \wedge (\operatorname{scal}(m) = \operatorname{normed\text{-}field\text{-}as\text{-}commutative\text{-}ring}(\operatorname{cc\text{-}normed\text{-}field}))) \Rightarrow  \\ \quad \quad \operatorname{act}(m)(+i, \operatorname{act}(m)(+i, x)) = \operatorname{vneg}(m)(x) \end{array}}}
\vnbprop{cmod-cact-decompose}{\vnbstmt{\begin{array}{@{}l@{}} \forall m.; \forall c \in \mathbb{C}, x \in \operatorname{vec}(m). \\ \quad (\operatorname{is\text{-}module}(m) \wedge (\operatorname{scal}(m) = \operatorname{normed\text{-}field\text{-}as\text{-}commutative\text{-}ring}(\operatorname{cc\text{-}normed\text{-}field}))) \Rightarrow  \\ \quad \quad \operatorname{vadd}(m)(\operatorname{act}(m)(\operatorname{real\text{-}part}(c), x), \operatorname{act}(m)(\operatorname{imag\text{-}part}(c), \operatorname{act}(m)(+i, x))) = \operatorname{act}(m)(c, x) \end{array}}}
\vnbprop{cmod-ract-fun}{\vnbstmt{\begin{array}{@{}l@{}} \forall m. \\ \quad (\operatorname{is\text{-}module}(m) \wedge (\operatorname{scal}(m) = \operatorname{normed\text{-}field\text{-}as\text{-}commutative\text{-}ring}(\operatorname{cc\text{-}normed\text{-}field}))) \Rightarrow  \\ \quad \quad \operatorname{cmod\text{-}ract}(m) \in ((\mathbb{R} \times \operatorname{vec}(m)) \to \operatorname{vec}(m)) \end{array}}}
\vnbprop{cmod-j-fun}{\vnbstmt{\begin{array}{@{}l@{}} \forall m. \\ \quad (\operatorname{is\text{-}module}(m) \wedge (\operatorname{scal}(m) = \operatorname{normed\text{-}field\text{-}as\text{-}commutative\text{-}ring}(\operatorname{cc\text{-}normed\text{-}field}))) \Rightarrow  \\ \quad \quad \operatorname{cmod\text{-}j}(m) \in (\operatorname{vec}(m) \to \operatorname{vec}(m)) \end{array}}}
\vnbprop{cmod-norm-homog-rr}{\vnbstmt{\begin{array}{@{}l@{}} \forall m, n.; \forall r \in \mathbb{R}, x \in \operatorname{vec}(m). \\ \quad (\operatorname{is\text{-}module}(m) \wedge ((\operatorname{scal}(m) = \operatorname{normed\text{-}field\text{-}as\text{-}commutative\text{-}ring}(\operatorname{cc\text{-}normed\text{-}field})) \wedge \operatorname{is\text{-}cmod\text{-}norm}(m, n))) \Rightarrow  \\ \quad \quad n(\operatorname{act}(m)(r, x)) = (\lvert r \rvert \cdot n(x)) \end{array}}}
\vnbprop{cmod-nvs-scal}{\vnbstmt{\begin{array}{@{}l@{}} \forall m, n. \\ \quad (\operatorname{is\text{-}module}(m) \wedge ((\operatorname{scal}(m) = \operatorname{normed\text{-}field\text{-}as\text{-}commutative\text{-}ring}(\operatorname{cc\text{-}normed\text{-}field})) \wedge \operatorname{is\text{-}cmod\text{-}norm}(m, n))) \Rightarrow  \\ \quad \quad (\operatorname{scal}(\operatorname{cmod\text{-}nvs}(m, n)) \simeq \operatorname{normed\text{-}field\text{-}as\text{-}commutative\text{-}ring}(\operatorname{rr\text{-}normed\text{-}field})) \end{array}}}
\vnbprop{cmod-cnvs-scal}{\vnbstmt{\begin{array}{@{}l@{}} \forall m, n. \\ \quad (\operatorname{is\text{-}module}(m) \wedge ((\operatorname{scal}(m) = \operatorname{normed\text{-}field\text{-}as\text{-}commutative\text{-}ring}(\operatorname{cc\text{-}normed\text{-}field})) \wedge \operatorname{is\text{-}cmod\text{-}norm}(m, n))) \Rightarrow  \\ \quad \quad (\operatorname{scal}(\operatorname{cmod\text{-}cnvs}(m, n)) \simeq \operatorname{normed\text{-}field\text{-}as\text{-}commutative\text{-}ring}(\operatorname{rr\text{-}normed\text{-}field})) \end{array}}}
\vnbprop{cmod-nvs-vec}{\vnbstmt{\begin{array}{@{}l@{}} \forall m, n. \\ \quad (\operatorname{is\text{-}module}(m) \wedge ((\operatorname{scal}(m) = \operatorname{normed\text{-}field\text{-}as\text{-}commutative\text{-}ring}(\operatorname{cc\text{-}normed\text{-}field})) \wedge \operatorname{is\text{-}cmod\text{-}norm}(m, n))) \Rightarrow  \\ \quad \quad (\operatorname{vec}(\operatorname{cmod\text{-}nvs}(m, n)) \simeq \operatorname{vec}(m)) \end{array}}}
\vnbprop{cmod-cnvs-vec}{\vnbstmt{\begin{array}{@{}l@{}} \forall m, n. \\ \quad (\operatorname{is\text{-}module}(m) \wedge ((\operatorname{scal}(m) = \operatorname{normed\text{-}field\text{-}as\text{-}commutative\text{-}ring}(\operatorname{cc\text{-}normed\text{-}field})) \wedge \operatorname{is\text{-}cmod\text{-}norm}(m, n))) \Rightarrow  \\ \quad \quad (\operatorname{vec}(\operatorname{cmod\text{-}cnvs}(m, n)) \simeq \operatorname{vec}(m)) \end{array}}}
\vnbprop{cmod-nvs-vadd}{\vnbstmt{\begin{array}{@{}l@{}} \forall m, n. \\ \quad (\operatorname{is\text{-}module}(m) \wedge ((\operatorname{scal}(m) = \operatorname{normed\text{-}field\text{-}as\text{-}commutative\text{-}ring}(\operatorname{cc\text{-}normed\text{-}field})) \wedge \operatorname{is\text{-}cmod\text{-}norm}(m, n))) \Rightarrow  \\ \quad \quad (\operatorname{vadd}(\operatorname{cmod\text{-}nvs}(m, n)) \simeq \operatorname{vadd}(m)) \end{array}}}
\vnbprop{cmod-cnvs-vadd}{\vnbstmt{\begin{array}{@{}l@{}} \forall m, n. \\ \quad (\operatorname{is\text{-}module}(m) \wedge ((\operatorname{scal}(m) = \operatorname{normed\text{-}field\text{-}as\text{-}commutative\text{-}ring}(\operatorname{cc\text{-}normed\text{-}field})) \wedge \operatorname{is\text{-}cmod\text{-}norm}(m, n))) \Rightarrow  \\ \quad \quad (\operatorname{vadd}(\operatorname{cmod\text{-}cnvs}(m, n)) \simeq \operatorname{vadd}(m)) \end{array}}}
\vnbprop{cmod-nvs-vzero}{\vnbstmt{\begin{array}{@{}l@{}} \forall m, n. \\ \quad (\operatorname{is\text{-}module}(m) \wedge ((\operatorname{scal}(m) = \operatorname{normed\text{-}field\text{-}as\text{-}commutative\text{-}ring}(\operatorname{cc\text{-}normed\text{-}field})) \wedge \operatorname{is\text{-}cmod\text{-}norm}(m, n))) \Rightarrow  \\ \quad \quad (\operatorname{vzero}(\operatorname{cmod\text{-}nvs}(m, n)) \simeq \operatorname{vzero}(m)) \end{array}}}
\vnbprop{cmod-cnvs-vzero}{\vnbstmt{\begin{array}{@{}l@{}} \forall m, n. \\ \quad (\operatorname{is\text{-}module}(m) \wedge ((\operatorname{scal}(m) = \operatorname{normed\text{-}field\text{-}as\text{-}commutative\text{-}ring}(\operatorname{cc\text{-}normed\text{-}field})) \wedge \operatorname{is\text{-}cmod\text{-}norm}(m, n))) \Rightarrow  \\ \quad \quad (\operatorname{vzero}(\operatorname{cmod\text{-}cnvs}(m, n)) \simeq \operatorname{vzero}(m)) \end{array}}}
\vnbprop{cmod-nvs-vneg}{\vnbstmt{\begin{array}{@{}l@{}} \forall m, n. \\ \quad (\operatorname{is\text{-}module}(m) \wedge ((\operatorname{scal}(m) = \operatorname{normed\text{-}field\text{-}as\text{-}commutative\text{-}ring}(\operatorname{cc\text{-}normed\text{-}field})) \wedge \operatorname{is\text{-}cmod\text{-}norm}(m, n))) \Rightarrow  \\ \quad \quad (\operatorname{vneg}(\operatorname{cmod\text{-}nvs}(m, n)) \simeq \operatorname{vneg}(m)) \end{array}}}
\vnbprop{cmod-cnvs-vneg}{\vnbstmt{\begin{array}{@{}l@{}} \forall m, n. \\ \quad (\operatorname{is\text{-}module}(m) \wedge ((\operatorname{scal}(m) = \operatorname{normed\text{-}field\text{-}as\text{-}commutative\text{-}ring}(\operatorname{cc\text{-}normed\text{-}field})) \wedge \operatorname{is\text{-}cmod\text{-}norm}(m, n))) \Rightarrow  \\ \quad \quad (\operatorname{vneg}(\operatorname{cmod\text{-}cnvs}(m, n)) \simeq \operatorname{vneg}(m)) \end{array}}}
\vnbprop{cmod-nvs-act}{\vnbstmt{\begin{array}{@{}l@{}} \forall m, n. \\ \quad (\operatorname{is\text{-}module}(m) \wedge ((\operatorname{scal}(m) = \operatorname{normed\text{-}field\text{-}as\text{-}commutative\text{-}ring}(\operatorname{cc\text{-}normed\text{-}field})) \wedge \operatorname{is\text{-}cmod\text{-}norm}(m, n))) \Rightarrow  \\ \quad \quad (\operatorname{act}(\operatorname{cmod\text{-}nvs}(m, n)) \simeq \operatorname{cmod\text{-}ract}(m)) \end{array}}}
\vnbprop{cmod-cnvs-act}{\vnbstmt{\begin{array}{@{}l@{}} \forall m, n. \\ \quad (\operatorname{is\text{-}module}(m) \wedge ((\operatorname{scal}(m) = \operatorname{normed\text{-}field\text{-}as\text{-}commutative\text{-}ring}(\operatorname{cc\text{-}normed\text{-}field})) \wedge \operatorname{is\text{-}cmod\text{-}norm}(m, n))) \Rightarrow  \\ \quad \quad (\operatorname{act}(\operatorname{cmod\text{-}cnvs}(m, n)) \simeq \operatorname{cmod\text{-}ract}(m)) \end{array}}}
\vnbprop{cmod-nvs-vnrm}{\vnbstmt{\begin{array}{@{}l@{}} \forall m, n. \\ \quad (\operatorname{is\text{-}module}(m) \wedge ((\operatorname{scal}(m) = \operatorname{normed\text{-}field\text{-}as\text{-}commutative\text{-}ring}(\operatorname{cc\text{-}normed\text{-}field})) \wedge \operatorname{is\text{-}cmod\text{-}norm}(m, n))) \Rightarrow  \\ \quad \quad (\operatorname{vnrm}(\operatorname{cmod\text{-}nvs}(m, n)) \simeq n) \end{array}}}
\vnbprop{cmod-cnvs-vnrm}{\vnbstmt{\begin{array}{@{}l@{}} \forall m, n. \\ \quad (\operatorname{is\text{-}module}(m) \wedge ((\operatorname{scal}(m) = \operatorname{normed\text{-}field\text{-}as\text{-}commutative\text{-}ring}(\operatorname{cc\text{-}normed\text{-}field})) \wedge \operatorname{is\text{-}cmod\text{-}norm}(m, n))) \Rightarrow  \\ \quad \quad (\operatorname{vnrm}(\operatorname{cmod\text{-}cnvs}(m, n)) \simeq n) \end{array}}}
\vnbprop{cmod-cnvs-jop}{\vnbstmt{\begin{array}{@{}l@{}} \forall m, n. \\ \quad (\operatorname{is\text{-}module}(m) \wedge ((\operatorname{scal}(m) = \operatorname{normed\text{-}field\text{-}as\text{-}commutative\text{-}ring}(\operatorname{cc\text{-}normed\text{-}field})) \wedge \operatorname{is\text{-}cmod\text{-}norm}(m, n))) \Rightarrow  \\ \quad \quad (J_{\operatorname{cmod\text{-}cnvs}(m, n)} \simeq \operatorname{cmod\text{-}j}(m)) \end{array}}}
\vnbprop{cmod-nvs-unfold}{\vnbstmt{\begin{array}{@{}l@{}} \forall m, n. \\ \quad (\operatorname{cmod\text{-}nvs}(m, n) \simeq {[}\operatorname{normed\text{-}field\text{-}as\text{-}commutative\text{-}ring}(\operatorname{rr\text{-}normed\text{-}field}), \operatorname{vec}(m), \operatorname{vadd}(m), \operatorname{vzero}(m), \operatorname{vneg}(m), \operatorname{cmod\text{-}ract}(m), n]) \end{array}}}
\vnbprop{cmod-nvs-length}{\vnbstmt{\begin{array}{@{}l@{}} \forall m, n. \\ \quad (\operatorname{is\text{-}module}(m) \wedge ((\operatorname{scal}(m) = \operatorname{normed\text{-}field\text{-}as\text{-}commutative\text{-}ring}(\operatorname{cc\text{-}normed\text{-}field})) \wedge \operatorname{is\text{-}cmod\text{-}norm}(m, n))) \Rightarrow  \\ \quad \quad \operatorname{length}(\operatorname{cmod\text{-}nvs}(m, n)) = 7 \end{array}}}
\vnbprop{cmod-nvs-scal-ring}{\vnbstmt{\begin{array}{@{}l@{}} \forall m, n. \\ \quad (\operatorname{is\text{-}module}(m) \wedge ((\operatorname{scal}(m) = \operatorname{normed\text{-}field\text{-}as\text{-}commutative\text{-}ring}(\operatorname{cc\text{-}normed\text{-}field})) \wedge \operatorname{is\text{-}cmod\text{-}norm}(m, n))) \Rightarrow  \\ \quad \quad \operatorname{is\text{-}ring}(\operatorname{scal}(\operatorname{cmod\text{-}nvs}(m, n))) \end{array}}}
\vnbprop{cmod-nvs-vec-set}{\vnbstmt{\begin{array}{@{}l@{}} \forall m, n. \\ \quad (\operatorname{is\text{-}module}(m) \wedge ((\operatorname{scal}(m) = \operatorname{normed\text{-}field\text{-}as\text{-}commutative\text{-}ring}(\operatorname{cc\text{-}normed\text{-}field})) \wedge \operatorname{is\text{-}cmod\text{-}norm}(m, n))) \Rightarrow  \\ \quad \quad \operatorname{vec}(\operatorname{cmod\text{-}nvs}(m, n)) \in \mathbf{Set} \end{array}}}
\vnbprop{cmod-nvs-vadd-fun}{\vnbstmt{\begin{array}{@{}l@{}} \forall m, n. \\ \quad (\operatorname{is\text{-}module}(m) \wedge ((\operatorname{scal}(m) = \operatorname{normed\text{-}field\text{-}as\text{-}commutative\text{-}ring}(\operatorname{cc\text{-}normed\text{-}field})) \wedge \operatorname{is\text{-}cmod\text{-}norm}(m, n))) \Rightarrow  \\ \quad \quad \operatorname{vadd}(\operatorname{cmod\text{-}nvs}(m, n)) \\ \quad \quad \in\; ((\operatorname{vec}(\operatorname{cmod\text{-}nvs}(m, n)) \times \operatorname{vec}(\operatorname{cmod\text{-}nvs}(m, n))) \to \operatorname{vec}(\operatorname{cmod\text{-}nvs}(m, n))) \end{array}}}
\vnbprop{cmod-nvs-vzero-in}{\vnbstmt{\begin{array}{@{}l@{}} \forall m, n. \\ \quad (\operatorname{is\text{-}module}(m) \wedge ((\operatorname{scal}(m) = \operatorname{normed\text{-}field\text{-}as\text{-}commutative\text{-}ring}(\operatorname{cc\text{-}normed\text{-}field})) \wedge \operatorname{is\text{-}cmod\text{-}norm}(m, n))) \Rightarrow  \\ \quad \quad \operatorname{vzero}(\operatorname{cmod\text{-}nvs}(m, n)) \in \operatorname{vec}(\operatorname{cmod\text{-}nvs}(m, n)) \end{array}}}
\vnbprop{cmod-nvs-vneg-fun}{\vnbstmt{\begin{array}{@{}l@{}} \forall m, n. \\ \quad (\operatorname{is\text{-}module}(m) \wedge ((\operatorname{scal}(m) = \operatorname{normed\text{-}field\text{-}as\text{-}commutative\text{-}ring}(\operatorname{cc\text{-}normed\text{-}field})) \wedge \operatorname{is\text{-}cmod\text{-}norm}(m, n))) \Rightarrow  \\ \quad \quad \operatorname{vneg}(\operatorname{cmod\text{-}nvs}(m, n)) \in (\operatorname{vec}(\operatorname{cmod\text{-}nvs}(m, n)) \to \operatorname{vec}(\operatorname{cmod\text{-}nvs}(m, n))) \end{array}}}
\vnbprop{cmod-nvs-act-fun}{\vnbstmt{\begin{array}{@{}l@{}} \forall m, n. \\ \quad (\operatorname{is\text{-}module}(m) \wedge ((\operatorname{scal}(m) = \operatorname{normed\text{-}field\text{-}as\text{-}commutative\text{-}ring}(\operatorname{cc\text{-}normed\text{-}field})) \wedge \operatorname{is\text{-}cmod\text{-}norm}(m, n))) \Rightarrow  \\ \quad \quad \operatorname{act}(\operatorname{cmod\text{-}nvs}(m, n)) \\ \quad \quad \in\; ((\operatorname{carr}(\operatorname{scal}(\operatorname{cmod\text{-}nvs}(m, n))) \times \operatorname{vec}(\operatorname{cmod\text{-}nvs}(m, n))) \to \operatorname{vec}(\operatorname{cmod\text{-}nvs}(m, n))) \end{array}}}
\vnbprop{cmod-nvs-vnrm-fun}{\vnbstmt{\begin{array}{@{}l@{}} \forall m, n. \\ \quad (\operatorname{is\text{-}module}(m) \wedge ((\operatorname{scal}(m) = \operatorname{normed\text{-}field\text{-}as\text{-}commutative\text{-}ring}(\operatorname{cc\text{-}normed\text{-}field})) \wedge \operatorname{is\text{-}cmod\text{-}norm}(m, n))) \Rightarrow  \\ \quad \quad \operatorname{vnrm}(\operatorname{cmod\text{-}nvs}(m, n)) \in (\operatorname{vec}(\operatorname{cmod\text{-}nvs}(m, n)) \to \mathbb{R}) \end{array}}}
\vnbprop{cmod-nvs-assoc}{\vnbstmt{\begin{array}{@{}l@{}} \forall m, n. \\ \quad (\operatorname{is\text{-}module}(m) \wedge ((\operatorname{scal}(m) = \operatorname{normed\text{-}field\text{-}as\text{-}commutative\text{-}ring}(\operatorname{cc\text{-}normed\text{-}field})) \wedge \operatorname{is\text{-}cmod\text{-}norm}(m, n))) \Rightarrow  \\ \quad \quad \operatorname{is\text{-}associative}(\operatorname{vadd}(\operatorname{cmod\text{-}nvs}(m, n)), \operatorname{vec}(\operatorname{cmod\text{-}nvs}(m, n))) \end{array}}}
\vnbprop{cmod-nvs-comm}{\vnbstmt{\begin{array}{@{}l@{}} \forall m, n. \\ \quad (\operatorname{is\text{-}module}(m) \wedge ((\operatorname{scal}(m) = \operatorname{normed\text{-}field\text{-}as\text{-}commutative\text{-}ring}(\operatorname{cc\text{-}normed\text{-}field})) \wedge \operatorname{is\text{-}cmod\text{-}norm}(m, n))) \Rightarrow  \\ \quad \quad \operatorname{is\text{-}commutative}(\operatorname{vadd}(\operatorname{cmod\text{-}nvs}(m, n)), \operatorname{vec}(\operatorname{cmod\text{-}nvs}(m, n))) \end{array}}}
\vnbprop{cmod-nvs-identity}{\vnbstmt{\begin{array}{@{}l@{}} \forall m, n. \\ \quad (\operatorname{is\text{-}module}(m) \wedge ((\operatorname{scal}(m) = \operatorname{normed\text{-}field\text{-}as\text{-}commutative\text{-}ring}(\operatorname{cc\text{-}normed\text{-}field})) \wedge \operatorname{is\text{-}cmod\text{-}norm}(m, n))) \Rightarrow  \\ \quad \quad \operatorname{is\text{-}identity}(\operatorname{vadd}(\operatorname{cmod\text{-}nvs}(m, n)), \operatorname{vzero}(\operatorname{cmod\text{-}nvs}(m, n)), \operatorname{vec}(\operatorname{cmod\text{-}nvs}(m, n))) \end{array}}}
\vnbprop{cmod-nvs-inverses}{\vnbstmt{\begin{array}{@{}l@{}} \forall m, n. \\ \quad (\operatorname{is\text{-}module}(m) \wedge ((\operatorname{scal}(m) = \operatorname{normed\text{-}field\text{-}as\text{-}commutative\text{-}ring}(\operatorname{cc\text{-}normed\text{-}field})) \wedge \operatorname{is\text{-}cmod\text{-}norm}(m, n))) \Rightarrow  \\ \quad \quad \operatorname{has\text{-}inverses}(\operatorname{vadd}(\operatorname{cmod\text{-}nvs}(m, n)), \operatorname{vzero}(\operatorname{cmod\text{-}nvs}(m, n)), \operatorname{vneg}(\operatorname{cmod\text{-}nvs}(m, n)), \operatorname{vec}(\operatorname{cmod\text{-}nvs}(m, n))) \end{array}}}
\vnbprop{cmod-nvs-scal-pin}{\vnbstmt{\begin{array}{@{}l@{}} \forall m, n. \\ \quad (\operatorname{is\text{-}module}(m) \wedge ((\operatorname{scal}(m) = \operatorname{normed\text{-}field\text{-}as\text{-}commutative\text{-}ring}(\operatorname{cc\text{-}normed\text{-}field})) \wedge \operatorname{is\text{-}cmod\text{-}norm}(m, n))) \Rightarrow  \\ \quad \quad \operatorname{scal}(\operatorname{cmod\text{-}nvs}(m, n)) = \operatorname{normed\text{-}field\text{-}as\text{-}commutative\text{-}ring}(\operatorname{rr\text{-}normed\text{-}field}) \end{array}}}
\vnbprop{cmod-nvs-act-vadd}{\vnbstmt{\begin{array}{@{}l@{}} \forall m, n.; \forall r \in \operatorname{carr}(\operatorname{scal}(\operatorname{cmod\text{-}nvs}(m, n))), x, y \in \operatorname{vec}(\operatorname{cmod\text{-}nvs}(m, n)). \\ \quad (\operatorname{is\text{-}module}(m) \wedge ((\operatorname{scal}(m) = \operatorname{normed\text{-}field\text{-}as\text{-}commutative\text{-}ring}(\operatorname{cc\text{-}normed\text{-}field})) \wedge \operatorname{is\text{-}cmod\text{-}norm}(m, n))) \Rightarrow  \\ \quad \quad \operatorname{act}(\operatorname{cmod\text{-}nvs}(m, n))(r, \operatorname{vadd}(\operatorname{cmod\text{-}nvs}(m, n))(x, y)) \\ \quad \quad =\; \operatorname{vadd}(\operatorname{cmod\text{-}nvs}(m, n))(\operatorname{act}(\operatorname{cmod\text{-}nvs}(m, n))(r, x), \operatorname{act}(\operatorname{cmod\text{-}nvs}(m, n))(r, y)) \end{array}}}
\vnbprop{cmod-nvs-act-add}{\vnbstmt{\begin{array}{@{}l@{}} \forall m, n.; \forall r, s \in \operatorname{carr}(\operatorname{scal}(\operatorname{cmod\text{-}nvs}(m, n))), x \in \operatorname{vec}(\operatorname{cmod\text{-}nvs}(m, n)). \\ \quad (\operatorname{is\text{-}module}(m) \wedge ((\operatorname{scal}(m) = \operatorname{normed\text{-}field\text{-}as\text{-}commutative\text{-}ring}(\operatorname{cc\text{-}normed\text{-}field})) \wedge \operatorname{is\text{-}cmod\text{-}norm}(m, n))) \Rightarrow  \\ \quad \quad \operatorname{act}(\operatorname{cmod\text{-}nvs}(m, n))(\operatorname{add}(\operatorname{scal}(\operatorname{cmod\text{-}nvs}(m, n)))(r, s), x) \\ \quad \quad =\; \operatorname{vadd}(\operatorname{cmod\text{-}nvs}(m, n))(\operatorname{act}(\operatorname{cmod\text{-}nvs}(m, n))(r, x), \operatorname{act}(\operatorname{cmod\text{-}nvs}(m, n))(s, x)) \end{array}}}
\vnbprop{cmod-nvs-act-mul}{\vnbstmt{\begin{array}{@{}l@{}} \forall m, n.; \forall r, s \in \operatorname{carr}(\operatorname{scal}(\operatorname{cmod\text{-}nvs}(m, n))), x \in \operatorname{vec}(\operatorname{cmod\text{-}nvs}(m, n)). \\ \quad (\operatorname{is\text{-}module}(m) \wedge ((\operatorname{scal}(m) = \operatorname{normed\text{-}field\text{-}as\text{-}commutative\text{-}ring}(\operatorname{cc\text{-}normed\text{-}field})) \wedge \operatorname{is\text{-}cmod\text{-}norm}(m, n))) \Rightarrow  \\ \quad \quad \operatorname{act}(\operatorname{cmod\text{-}nvs}(m, n))(\operatorname{mul}(\operatorname{scal}(\operatorname{cmod\text{-}nvs}(m, n)))(r, s), x) \\ \quad \quad =\; \operatorname{act}(\operatorname{cmod\text{-}nvs}(m, n))(r, \operatorname{act}(\operatorname{cmod\text{-}nvs}(m, n))(s, x)) \end{array}}}
\vnbprop{cmod-nvs-act-one}{\vnbstmt{\begin{array}{@{}l@{}} \forall m, n.; \forall x \in \operatorname{vec}(\operatorname{cmod\text{-}nvs}(m, n)). \\ \quad (\operatorname{is\text{-}module}(m) \wedge ((\operatorname{scal}(m) = \operatorname{normed\text{-}field\text{-}as\text{-}commutative\text{-}ring}(\operatorname{cc\text{-}normed\text{-}field})) \wedge \operatorname{is\text{-}cmod\text{-}norm}(m, n))) \Rightarrow  \\ \quad \quad \operatorname{act}(\operatorname{cmod\text{-}nvs}(m, n))(\operatorname{one}(\operatorname{scal}(\operatorname{cmod\text{-}nvs}(m, n))), x) = x \end{array}}}
\vnbprop{cmod-nvs-nonneg}{\vnbstmt{\begin{array}{@{}l@{}} \forall m, n.; \forall x \in \operatorname{vec}(\operatorname{cmod\text{-}nvs}(m, n)). \\ \quad (\operatorname{is\text{-}module}(m) \wedge ((\operatorname{scal}(m) = \operatorname{normed\text{-}field\text{-}as\text{-}commutative\text{-}ring}(\operatorname{cc\text{-}normed\text{-}field})) \wedge \operatorname{is\text{-}cmod\text{-}norm}(m, n))) \Rightarrow  \\ \quad \quad 0 \leq \operatorname{vnrm}(\operatorname{cmod\text{-}nvs}(m, n))(x) \end{array}}}
\vnbprop{cmod-nvs-definite}{\vnbstmt{\begin{array}{@{}l@{}} \forall m, n.; \forall x \in \operatorname{vec}(\operatorname{cmod\text{-}nvs}(m, n)). \\ \quad (\operatorname{is\text{-}module}(m) \wedge ((\operatorname{scal}(m) = \operatorname{normed\text{-}field\text{-}as\text{-}commutative\text{-}ring}(\operatorname{cc\text{-}normed\text{-}field})) \wedge \operatorname{is\text{-}cmod\text{-}norm}(m, n))) \Rightarrow  \\ \quad \quad (\operatorname{vnrm}(\operatorname{cmod\text{-}nvs}(m, n))(x) = 0) \Leftrightarrow (x = \operatorname{vzero}(\operatorname{cmod\text{-}nvs}(m, n))) \end{array}}}
\vnbprop{cmod-nvs-homog}{\vnbstmt{\begin{array}{@{}l@{}} \forall m, n.; \forall r \in \operatorname{carr}(\operatorname{scal}(\operatorname{cmod\text{-}nvs}(m, n))), x \in \operatorname{vec}(\operatorname{cmod\text{-}nvs}(m, n)). \\ \quad (\operatorname{is\text{-}module}(m) \wedge ((\operatorname{scal}(m) = \operatorname{normed\text{-}field\text{-}as\text{-}commutative\text{-}ring}(\operatorname{cc\text{-}normed\text{-}field})) \wedge \operatorname{is\text{-}cmod\text{-}norm}(m, n))) \Rightarrow  \\ \quad \quad \operatorname{vnrm}(\operatorname{cmod\text{-}nvs}(m, n))(\operatorname{act}(\operatorname{cmod\text{-}nvs}(m, n))(r, x)) = (\lvert r \rvert \cdot \operatorname{vnrm}(\operatorname{cmod\text{-}nvs}(m, n))(x)) \end{array}}}
\vnbprop{cmod-nvs-triangle}{\vnbstmt{\begin{array}{@{}l@{}} \forall m, n.; \forall x, y \in \operatorname{vec}(\operatorname{cmod\text{-}nvs}(m, n)). \\ \quad (\operatorname{is\text{-}module}(m) \wedge ((\operatorname{scal}(m) = \operatorname{normed\text{-}field\text{-}as\text{-}commutative\text{-}ring}(\operatorname{cc\text{-}normed\text{-}field})) \wedge \operatorname{is\text{-}cmod\text{-}norm}(m, n))) \Rightarrow  \\ \quad \quad \operatorname{vnrm}(\operatorname{cmod\text{-}nvs}(m, n))(\operatorname{vadd}(\operatorname{cmod\text{-}nvs}(m, n))(x, y)) \\ \quad \quad \leq\; \operatorname{vnrm}(\operatorname{cmod\text{-}nvs}(m, n))(x) + \operatorname{vnrm}(\operatorname{cmod\text{-}nvs}(m, n))(y) \end{array}}}
\vnbprop{cmod-nvs-is-nvs}{\vnbstmt{\begin{array}{@{}l@{}} \forall m, n. \\ \quad (\operatorname{is\text{-}module}(m) \wedge ((\operatorname{scal}(m) = \operatorname{normed\text{-}field\text{-}as\text{-}commutative\text{-}ring}(\operatorname{cc\text{-}normed\text{-}field})) \wedge \operatorname{is\text{-}cmod\text{-}norm}(m, n))) \Rightarrow  \\ \quad \quad \operatorname{is\text{-}normed\text{-}vector\text{-}space}(\operatorname{cmod\text{-}nvs}(m, n)) \end{array}}}
\vnbprop{cmod-cact-value}{\vnbstmt{\begin{array}{@{}l@{}} \forall m, n.; \forall c \in \mathbb{C}, x \in \operatorname{vec}(m). \\ \quad (\operatorname{is\text{-}module}(m) \wedge ((\operatorname{scal}(m) = \operatorname{normed\text{-}field\text{-}as\text{-}commutative\text{-}ring}(\operatorname{cc\text{-}normed\text{-}field})) \wedge \operatorname{is\text{-}cmod\text{-}norm}(m, n))) \Rightarrow  \\ \quad \quad \operatorname{cact}(\operatorname{cmod\text{-}cnvs}(m, n))(c, x) = \operatorname{act}(m)(c, x) \end{array}}}
\vnbprop{cmod-cnvs-length}{\vnbstmt{\begin{array}{@{}l@{}} \forall m, n. \\ \quad (\operatorname{is\text{-}module}(m) \wedge ((\operatorname{scal}(m) = \operatorname{normed\text{-}field\text{-}as\text{-}commutative\text{-}ring}(\operatorname{cc\text{-}normed\text{-}field})) \wedge \operatorname{is\text{-}cmod\text{-}norm}(m, n))) \Rightarrow  \\ \quad \quad \operatorname{length}(\operatorname{cmod\text{-}cnvs}(m, n)) = 8 \end{array}}}
\vnbprop{cmod-cnvs-scal-ring}{\vnbstmt{\begin{array}{@{}l@{}} \forall m, n. \\ \quad (\operatorname{is\text{-}module}(m) \wedge ((\operatorname{scal}(m) = \operatorname{normed\text{-}field\text{-}as\text{-}commutative\text{-}ring}(\operatorname{cc\text{-}normed\text{-}field})) \wedge \operatorname{is\text{-}cmod\text{-}norm}(m, n))) \Rightarrow  \\ \quad \quad \operatorname{is\text{-}ring}(\operatorname{scal}(\operatorname{cmod\text{-}cnvs}(m, n))) \end{array}}}
\vnbprop{cmod-cnvs-vec-set}{\vnbstmt{\begin{array}{@{}l@{}} \forall m, n. \\ \quad (\operatorname{is\text{-}module}(m) \wedge ((\operatorname{scal}(m) = \operatorname{normed\text{-}field\text{-}as\text{-}commutative\text{-}ring}(\operatorname{cc\text{-}normed\text{-}field})) \wedge \operatorname{is\text{-}cmod\text{-}norm}(m, n))) \Rightarrow  \\ \quad \quad \operatorname{vec}(\operatorname{cmod\text{-}cnvs}(m, n)) \in \mathbf{Set} \end{array}}}
\vnbprop{cmod-cnvs-vadd-fun}{\vnbstmt{\begin{array}{@{}l@{}} \forall m, n. \\ \quad (\operatorname{is\text{-}module}(m) \wedge ((\operatorname{scal}(m) = \operatorname{normed\text{-}field\text{-}as\text{-}commutative\text{-}ring}(\operatorname{cc\text{-}normed\text{-}field})) \wedge \operatorname{is\text{-}cmod\text{-}norm}(m, n))) \Rightarrow  \\ \quad \quad \operatorname{vadd}(\operatorname{cmod\text{-}cnvs}(m, n)) \\ \quad \quad \in\; ((\operatorname{vec}(\operatorname{cmod\text{-}cnvs}(m, n)) \times \operatorname{vec}(\operatorname{cmod\text{-}cnvs}(m, n))) \to \operatorname{vec}(\operatorname{cmod\text{-}cnvs}(m, n))) \end{array}}}
\vnbprop{cmod-cnvs-vzero-in}{\vnbstmt{\begin{array}{@{}l@{}} \forall m, n. \\ \quad (\operatorname{is\text{-}module}(m) \wedge ((\operatorname{scal}(m) = \operatorname{normed\text{-}field\text{-}as\text{-}commutative\text{-}ring}(\operatorname{cc\text{-}normed\text{-}field})) \wedge \operatorname{is\text{-}cmod\text{-}norm}(m, n))) \Rightarrow  \\ \quad \quad \operatorname{vzero}(\operatorname{cmod\text{-}cnvs}(m, n)) \in \operatorname{vec}(\operatorname{cmod\text{-}cnvs}(m, n)) \end{array}}}
\vnbprop{cmod-cnvs-vneg-fun}{\vnbstmt{\begin{array}{@{}l@{}} \forall m, n. \\ \quad (\operatorname{is\text{-}module}(m) \wedge ((\operatorname{scal}(m) = \operatorname{normed\text{-}field\text{-}as\text{-}commutative\text{-}ring}(\operatorname{cc\text{-}normed\text{-}field})) \wedge \operatorname{is\text{-}cmod\text{-}norm}(m, n))) \Rightarrow  \\ \quad \quad \operatorname{vneg}(\operatorname{cmod\text{-}cnvs}(m, n)) \in (\operatorname{vec}(\operatorname{cmod\text{-}cnvs}(m, n)) \to \operatorname{vec}(\operatorname{cmod\text{-}cnvs}(m, n))) \end{array}}}
\vnbprop{cmod-cnvs-act-fun}{\vnbstmt{\begin{array}{@{}l@{}} \forall m, n. \\ \quad (\operatorname{is\text{-}module}(m) \wedge ((\operatorname{scal}(m) = \operatorname{normed\text{-}field\text{-}as\text{-}commutative\text{-}ring}(\operatorname{cc\text{-}normed\text{-}field})) \wedge \operatorname{is\text{-}cmod\text{-}norm}(m, n))) \Rightarrow  \\ \quad \quad \operatorname{act}(\operatorname{cmod\text{-}cnvs}(m, n)) \\ \quad \quad \in\; ((\operatorname{carr}(\operatorname{scal}(\operatorname{cmod\text{-}cnvs}(m, n))) \times \operatorname{vec}(\operatorname{cmod\text{-}cnvs}(m, n))) \to \operatorname{vec}(\operatorname{cmod\text{-}cnvs}(m, n))) \end{array}}}
\vnbprop{cmod-cnvs-vnrm-fun}{\vnbstmt{\begin{array}{@{}l@{}} \forall m, n. \\ \quad (\operatorname{is\text{-}module}(m) \wedge ((\operatorname{scal}(m) = \operatorname{normed\text{-}field\text{-}as\text{-}commutative\text{-}ring}(\operatorname{cc\text{-}normed\text{-}field})) \wedge \operatorname{is\text{-}cmod\text{-}norm}(m, n))) \Rightarrow  \\ \quad \quad \operatorname{vnrm}(\operatorname{cmod\text{-}cnvs}(m, n)) \in (\operatorname{vec}(\operatorname{cmod\text{-}cnvs}(m, n)) \to \mathbb{R}) \end{array}}}
\vnbprop{cmod-cnvs-jop-fun}{\vnbstmt{\begin{array}{@{}l@{}} \forall m, n. \\ \quad (\operatorname{is\text{-}module}(m) \wedge ((\operatorname{scal}(m) = \operatorname{normed\text{-}field\text{-}as\text{-}commutative\text{-}ring}(\operatorname{cc\text{-}normed\text{-}field})) \wedge \operatorname{is\text{-}cmod\text{-}norm}(m, n))) \Rightarrow  \\ \quad \quad J_{\operatorname{cmod\text{-}cnvs}(m, n)} \in (\operatorname{vec}(\operatorname{cmod\text{-}cnvs}(m, n)) \to \operatorname{vec}(\operatorname{cmod\text{-}cnvs}(m, n))) \end{array}}}
\vnbprop{cmod-cnvs-nvs}{\vnbstmt{\begin{array}{@{}l@{}} \forall m, n. \\ \quad (\operatorname{is\text{-}module}(m) \wedge ((\operatorname{scal}(m) = \operatorname{normed\text{-}field\text{-}as\text{-}commutative\text{-}ring}(\operatorname{cc\text{-}normed\text{-}field})) \wedge \operatorname{is\text{-}cmod\text{-}norm}(m, n))) \Rightarrow  \\ \quad \quad \operatorname{is\text{-}normed\text{-}vector\text{-}space}({[}\operatorname{scal}(\operatorname{cmod\text{-}cnvs}(m, n)), \operatorname{vec}(\operatorname{cmod\text{-}cnvs}(m, n)), \operatorname{vadd}(\operatorname{cmod\text{-}cnvs}(m, n)), \operatorname{vzero}(\operatorname{cmod\text{-}cnvs}(m, n)), \operatorname{vneg}(\operatorname{cmod\text{-}cnvs}(m, n)), \operatorname{act}(\operatorname{cmod\text{-}cnvs}(m, n)), \operatorname{vnrm}(\operatorname{cmod\text{-}cnvs}(m, n))]) \end{array}}}
\vnbprop{cmod-cnvs-real-linear}{\vnbstmt{\begin{array}{@{}l@{}} \forall m, n. \\ \quad (\operatorname{is\text{-}module}(m) \wedge ((\operatorname{scal}(m) = \operatorname{normed\text{-}field\text{-}as\text{-}commutative\text{-}ring}(\operatorname{cc\text{-}normed\text{-}field})) \wedge \operatorname{is\text{-}cmod\text{-}norm}(m, n))) \Rightarrow  \\ \quad \quad \operatorname{is\text{-}real\text{-}linear\text{-}op}(\operatorname{cmod\text{-}cnvs}(m, n), J_{\operatorname{cmod\text{-}cnvs}(m, n)}) \end{array}}}
\vnbprop{cmod-cnvs-square}{\vnbstmt{\begin{array}{@{}l@{}} \forall m, n. \\ \quad (\operatorname{is\text{-}module}(m) \wedge ((\operatorname{scal}(m) = \operatorname{normed\text{-}field\text{-}as\text{-}commutative\text{-}ring}(\operatorname{cc\text{-}normed\text{-}field})) \wedge \operatorname{is\text{-}cmod\text{-}norm}(m, n))) \Rightarrow  \\ \quad \quad \operatorname{is\text{-}square\text{-}minus\text{-}id}(\operatorname{cmod\text{-}cnvs}(m, n), J_{\operatorname{cmod\text{-}cnvs}(m, n)}) \end{array}}}
\vnbprop{cmod-cnvs-cact-homog}{\vnbstmt{\begin{array}{@{}l@{}} \forall m, n.; \forall c \in \mathbb{C}, x \in \operatorname{vec}(\operatorname{cmod\text{-}cnvs}(m, n)). \\ \quad (\operatorname{is\text{-}module}(m) \wedge ((\operatorname{scal}(m) = \operatorname{normed\text{-}field\text{-}as\text{-}commutative\text{-}ring}(\operatorname{cc\text{-}normed\text{-}field})) \wedge \operatorname{is\text{-}cmod\text{-}norm}(m, n))) \Rightarrow  \\ \quad \quad \operatorname{vnrm}(\operatorname{cmod\text{-}cnvs}(m, n))(\operatorname{cact}(\operatorname{cmod\text{-}cnvs}(m, n))(c, x)) = (\operatorname{magnitude}(c) \cdot \operatorname{vnrm}(\operatorname{cmod\text{-}cnvs}(m, n))(x)) \end{array}}}
\vnbprop{cmod-cnvs-is-cnvs}{\vnbstmt{\begin{array}{@{}l@{}} \forall m, n. \\ \quad (\operatorname{is\text{-}module}(m) \wedge ((\operatorname{scal}(m) = \operatorname{normed\text{-}field\text{-}as\text{-}commutative\text{-}ring}(\operatorname{cc\text{-}normed\text{-}field})) \wedge \operatorname{is\text{-}cmod\text{-}norm}(m, n))) \Rightarrow  \\ \quad \quad \operatorname{is\text{-}complex\text{-}normed\text{-}vector\text{-}space}(\operatorname{cmod\text{-}cnvs}(m, n)) \end{array}}}

\vnbbookfile{theorem-library/cnvs-c}{theorem-library/cnvs-c}
\noindent{\small cnvs-c.scm -- wave of 2026-11-09/10 (agent CNVS), wired 2026-11-10 from scratchpad/cnvs-1109/cnvs-c.scm. Library: analysis.}\par
\vnbprop{cips-ipn-in-rr}{\vnbstmt{\begin{array}{@{}l@{}} \forall w.; \forall x \in \operatorname{vec}(w). \\ \quad \operatorname{is\text{-}complex\text{-}inner\text{-}product\text{-}space}(w) \Rightarrow (\lVert x \rVert \in \mathbb{R}) \end{array}}}
\vnbprop{cips-ipn-nonneg}{\vnbstmt{\begin{array}{@{}l@{}} \forall w.; \forall x \in \operatorname{vec}(w). \\ \quad \operatorname{is\text{-}complex\text{-}inner\text{-}product\text{-}space}(w) \Rightarrow (0 \leq \lVert x \rVert) \end{array}}}
\vnbprop{cips-ipn-sq}{\vnbstmt{\begin{array}{@{}l@{}} \forall w.; \forall x \in \operatorname{vec}(w). \\ \quad \operatorname{is\text{-}complex\text{-}inner\text{-}product\text{-}space}(w) \Rightarrow  \\ \quad \quad (\lVert x \rVert \cdot \lVert x \rVert) = \langle\cdot,\cdot\rangle_{w}(x, x) \end{array}}}
\vnbprop{cips-ipn-zero-iff}{\vnbstmt{\begin{array}{@{}l@{}} \forall w.; \forall x \in \operatorname{vec}(w). \\ \quad \operatorname{is\text{-}complex\text{-}inner\text{-}product\text{-}space}(w) \Rightarrow ((\lVert x \rVert = 0) \Leftrightarrow (x = \operatorname{vzero}(w))) \end{array}}}
\vnbprop{cips-ipn-homog}{\vnbstmt{\begin{array}{@{}l@{}} \forall w.; \forall c \in \mathbb{C}, x \in \operatorname{vec}(w). \\ \quad \operatorname{is\text{-}complex\text{-}inner\text{-}product\text{-}space}(w) \Rightarrow  \\ \quad \quad \lVert \operatorname{act}(w)(c, x) \rVert = (\operatorname{magnitude}(c) \cdot \lVert x \rVert) \end{array}}}
\vnbprop{cips-vnrm-fun}{\vnbstmt{\begin{array}{@{}l@{}} \forall w. \\ \quad \operatorname{is\text{-}complex\text{-}inner\text{-}product\text{-}space}(w) \Rightarrow (\operatorname{cips\text{-}vnrm}(w) \in (\operatorname{vec}(w) \to \mathbb{R})) \end{array}}}
\vnbprop{cips-vnrm-value}{\vnbstmt{\begin{array}{@{}l@{}} \forall w.; \forall x \in \operatorname{vec}(w). \\ \quad \operatorname{is\text{-}complex\text{-}inner\text{-}product\text{-}space}(w) \Rightarrow (\operatorname{cips\text{-}vnrm}(w)(x) = \lVert x \rVert) \end{array}}}
\vnbprop{cips-module-view-vadd}{\vnbstmt{\begin{array}{@{}l@{}} \forall w. \\ \quad \operatorname{is\text{-}complex\text{-}inner\text{-}product\text{-}space}(w) \Rightarrow  \\ \quad \quad (\operatorname{vadd}(\operatorname{complex\text{-}inner\text{-}product\text{-}space\text{-}as\text{-}module}(w)) \simeq \operatorname{vadd}(w)) \end{array}}}
\vnbprop{cips-module-view-vzero}{\vnbstmt{\begin{array}{@{}l@{}} \forall w. \\ \quad \operatorname{is\text{-}complex\text{-}inner\text{-}product\text{-}space}(w) \Rightarrow  \\ \quad \quad (\operatorname{vzero}(\operatorname{complex\text{-}inner\text{-}product\text{-}space\text{-}as\text{-}module}(w)) \simeq \operatorname{vzero}(w)) \end{array}}}
\vnbprop{cips-cmod-norm}{\vnbstmt{\begin{array}{@{}l@{}} \forall w. \\ \quad \operatorname{is\text{-}complex\text{-}inner\text{-}product\text{-}space}(w) \Rightarrow  \\ \quad \quad \operatorname{is\text{-}cmod\text{-}norm}(\operatorname{complex\text{-}inner\text{-}product\text{-}space\text{-}as\text{-}module}(w), \operatorname{cips\text{-}vnrm}(w)) \end{array}}}
\vnbprop{cips-cnvs-is-cnvs}{\vnbstmt{\begin{array}{@{}l@{}} \forall w. \\ \quad \operatorname{is\text{-}complex\text{-}inner\text{-}product\text{-}space}(w) \Rightarrow  \\ \quad \quad \operatorname{is\text{-}complex\text{-}normed\text{-}vector\text{-}space}(\operatorname{cips\text{-}cnvs}(w)) \end{array}}}
\vnbprop{cips-cnvs-scal}{\vnbstmt{\begin{array}{@{}l@{}} \forall w. \\ \quad \operatorname{is\text{-}complex\text{-}inner\text{-}product\text{-}space}(w) \Rightarrow  \\ \quad \quad (\operatorname{scal}(\operatorname{cips\text{-}cnvs}(w)) \simeq \operatorname{normed\text{-}field\text{-}as\text{-}commutative\text{-}ring}(\operatorname{rr\text{-}normed\text{-}field})) \end{array}}}
\vnbprop{cips-cnvs-vec}{\vnbstmt{\begin{array}{@{}l@{}} \forall w. \\ \quad \operatorname{is\text{-}complex\text{-}inner\text{-}product\text{-}space}(w) \Rightarrow (\operatorname{vec}(\operatorname{cips\text{-}cnvs}(w)) \simeq \operatorname{vec}(w)) \end{array}}}
\vnbprop{cips-cnvs-vadd}{\vnbstmt{\begin{array}{@{}l@{}} \forall w. \\ \quad \operatorname{is\text{-}complex\text{-}inner\text{-}product\text{-}space}(w) \Rightarrow (\operatorname{vadd}(\operatorname{cips\text{-}cnvs}(w)) \simeq \operatorname{vadd}(w)) \end{array}}}
\vnbprop{cips-cnvs-vzero}{\vnbstmt{\begin{array}{@{}l@{}} \forall w. \\ \quad \operatorname{is\text{-}complex\text{-}inner\text{-}product\text{-}space}(w) \Rightarrow (\operatorname{vzero}(\operatorname{cips\text{-}cnvs}(w)) \simeq \operatorname{vzero}(w)) \end{array}}}
\vnbprop{cips-cnvs-vneg}{\vnbstmt{\begin{array}{@{}l@{}} \forall w. \\ \quad \operatorname{is\text{-}complex\text{-}inner\text{-}product\text{-}space}(w) \Rightarrow (\operatorname{vneg}(\operatorname{cips\text{-}cnvs}(w)) \simeq \operatorname{vneg}(w)) \end{array}}}
\vnbprop{cips-cnvs-vnrm}{\vnbstmt{\begin{array}{@{}l@{}} \forall w. \\ \quad \operatorname{is\text{-}complex\text{-}inner\text{-}product\text{-}space}(w) \Rightarrow (\operatorname{vnrm}(\operatorname{cips\text{-}cnvs}(w)) \simeq \operatorname{cips\text{-}vnrm}(w)) \end{array}}}
\vnbprop{cips-cnvs-vnrm-value}{\vnbstmt{\begin{array}{@{}l@{}} \forall w.; \forall x \in \operatorname{vec}(w). \\ \quad \operatorname{is\text{-}complex\text{-}inner\text{-}product\text{-}space}(w) \Rightarrow  \\ \quad \quad \operatorname{vnrm}(\operatorname{cips\text{-}cnvs}(w))(x) = \lVert x \rVert \end{array}}}
\vnbprop{cips-cnvs-cact}{\vnbstmt{\begin{array}{@{}l@{}} \forall w.; \forall c \in \mathbb{C}, x \in \operatorname{vec}(w). \\ \quad \operatorname{is\text{-}complex\text{-}inner\text{-}product\text{-}space}(w) \Rightarrow  \\ \quad \quad \operatorname{cact}(\operatorname{cips\text{-}cnvs}(w))(c, x) = \operatorname{act}(w)(c, x) \end{array}}}
\vnbprop{cips-cnvs-jop-value}{\vnbstmt{\begin{array}{@{}l@{}} \forall w.; \forall x \in \operatorname{vec}(w). \\ \quad \operatorname{is\text{-}complex\text{-}inner\text{-}product\text{-}space}(w) \Rightarrow  \\ \quad \quad J_{\operatorname{cips\text{-}cnvs}(w)}(x) = \operatorname{act}(w)(+i, x) \end{array}}}
\vnbprop{cphs-cnvs-is-cnvs}{\vnbstmt{\begin{array}{@{}l@{}} \forall v. \\ \quad \operatorname{is\text{-}complex\text{-}pre\text{-}hilbert\text{-}space}(v) \Rightarrow  \\ \quad \quad \operatorname{is\text{-}complex\text{-}normed\text{-}vector\text{-}space}(\operatorname{csip\text{-}cnvs}(v)) \end{array}}}
\vnbprop{csip-cnvs-unfold}{\vnbstmt{\begin{array}{@{}l@{}} \forall v. \\ \quad \operatorname{is\text{-}complex\text{-}pre\text{-}hilbert\text{-}space}(v) \Rightarrow (\operatorname{csip\text{-}cnvs}(v) \simeq \operatorname{cips\text{-}cnvs}(\operatorname{csip\text{-}cips}(v))) \end{array}}}
\vnbprop{csip-cnvs-vec}{\vnbstmt{\forall v.\; (\operatorname{is\text{-}complex\text{-}pre\text{-}hilbert\text{-}space}(v) \Rightarrow (\operatorname{vec}(\operatorname{csip\text{-}cnvs}(v)) \simeq \operatorname{vec}(v)))}}
\vnbprop{csip-cnvs-vadd}{\vnbstmt{\begin{array}{@{}l@{}} \forall v. \\ \quad \operatorname{is\text{-}complex\text{-}pre\text{-}hilbert\text{-}space}(v) \Rightarrow (\operatorname{vadd}(\operatorname{csip\text{-}cnvs}(v)) \simeq \operatorname{vadd}(v)) \end{array}}}
\vnbprop{csip-cnvs-vzero}{\vnbstmt{\begin{array}{@{}l@{}} \forall v. \\ \quad \operatorname{is\text{-}complex\text{-}pre\text{-}hilbert\text{-}space}(v) \Rightarrow (\operatorname{vzero}(\operatorname{csip\text{-}cnvs}(v)) \simeq \operatorname{vzero}(v)) \end{array}}}
\vnbprop{csip-cnvs-vneg}{\vnbstmt{\begin{array}{@{}l@{}} \forall v. \\ \quad \operatorname{is\text{-}complex\text{-}pre\text{-}hilbert\text{-}space}(v) \Rightarrow (\operatorname{vneg}(\operatorname{csip\text{-}cnvs}(v)) \simeq \operatorname{vneg}(v)) \end{array}}}
\vnbprop{csip-cnvs-cact}{\vnbstmt{\begin{array}{@{}l@{}} \forall v.; \forall c \in \mathbb{C}, x \in \operatorname{vec}(v). \\ \quad \operatorname{is\text{-}complex\text{-}pre\text{-}hilbert\text{-}space}(v) \Rightarrow  \\ \quad \quad \operatorname{cact}(\operatorname{csip\text{-}cnvs}(v))(c, x) = \operatorname{cact}(v)(c, x) \end{array}}}
\vnbprop{csip-cnvs-vnrm-value}{\vnbstmt{\begin{array}{@{}l@{}} \forall v.; \forall x \in \operatorname{vec}(v). \\ \quad \operatorname{is\text{-}complex\text{-}pre\text{-}hilbert\text{-}space}(v) \Rightarrow  \\ \quad \quad \operatorname{vnrm}(\operatorname{csip\text{-}cnvs}(v))(x) = \lVert x \rVert \end{array}}}
\vnbprop{cips-cnvs-normed-ag}{\vnbstmt{\begin{array}{@{}l@{}} \forall w. \\ \quad \operatorname{is\text{-}complex\text{-}inner\text{-}product\text{-}space}(w) \Rightarrow  \\ \quad \quad (\operatorname{cnvs\text{-}as\text{-}normed\text{-}ag}(\operatorname{cips\text{-}cnvs}(w)) \simeq \operatorname{ip\text{-}normed\text{-}ag}(w)) \end{array}}}
\vnbprop{cips-cnvs-metric-space}{\vnbstmt{\begin{array}{@{}l@{}} \forall w. \\ \quad \operatorname{is\text{-}complex\text{-}inner\text{-}product\text{-}space}(w) \Rightarrow  \\ \quad \quad (\operatorname{cnvs\text{-}metric\text{-}space}(\operatorname{cips\text{-}cnvs}(w)) \simeq \operatorname{ip\text{-}metric\text{-}space}(w)) \end{array}}}
\vnbprop{hilbert-cnvs-is-cbs}{\vnbstmt{\forall w.\; (\operatorname{is\text{-}hilbert\text{-}space}(w) \Rightarrow \operatorname{is\text{-}complex\text{-}banach\text{-}space}(\operatorname{cips\text{-}cnvs}(w)))}}

\vnbbookfile{theorem-library/cnvs-d}{theorem-library/cnvs-d}
\noindent{\small cnvs-d.scm -- wave of 2026-11-09/10 (agent CNVS), wired 2026-11-10 from scratchpad/cnvs-1109/cnvs-d.scm. Library: analysis.}\par
\vnbprop{cnvd-norm-in-rr}{\vnbstmt{\begin{array}{@{}l@{}} \forall v, x. \\ \quad (\operatorname{is\text{-}complex\text{-}normed\text{-}vector\text{-}space}(v) \wedge (x \in \operatorname{vec}(v))) \Rightarrow (\operatorname{vnrm}(v)(x) \in \mathbb{R}) \end{array}}}
\vnbprop{cnvd-norm-nonneg}{\vnbstmt{\begin{array}{@{}l@{}} \forall v, x. \\ \quad (\operatorname{is\text{-}complex\text{-}normed\text{-}vector\text{-}space}(v) \wedge (x \in \operatorname{vec}(v))) \Rightarrow (0 \leq \operatorname{vnrm}(v)(x)) \end{array}}}
\vnbprop{cnvd-norm-zero-vector}{\vnbstmt{\begin{array}{@{}l@{}} \forall v, x. \\ \quad (\operatorname{is\text{-}complex\text{-}normed\text{-}vector\text{-}space}(v) \wedge ((x \in \operatorname{vec}(v)) \wedge (\operatorname{vnrm}(v)(x) = 0))) \Rightarrow  \\ \quad \quad x = \operatorname{vzero}(v) \end{array}}}
\vnbprop{cnvd-norm-vneg}{\vnbstmt{\begin{array}{@{}l@{}} \forall v, x. \\ \quad (\operatorname{is\text{-}complex\text{-}normed\text{-}vector\text{-}space}(v) \wedge (x \in \operatorname{vec}(v))) \Rightarrow  \\ \quad \quad \operatorname{vnrm}(v)(\operatorname{vneg}(v)(x)) = \operatorname{vnrm}(v)(x) \end{array}}}
\vnbprop{cnvd-cact-vneg}{\vnbstmt{\begin{array}{@{}l@{}} \forall v, c, x. \\ \quad (\operatorname{is\text{-}complex\text{-}normed\text{-}vector\text{-}space}(v) \wedge ((c \in \mathbb{C}) \wedge (x \in \operatorname{vec}(v)))) \Rightarrow  \\ \quad \quad \operatorname{cact}(v)(c, \operatorname{vneg}(v)(x)) = \operatorname{vneg}(v)(\operatorname{cact}(v)(c, x)) \end{array}}}
\vnbprop{cnvd-cact-comm}{\vnbstmt{\begin{array}{@{}l@{}} \forall v, c, d, x. \\ \quad (\operatorname{is\text{-}complex\text{-}normed\text{-}vector\text{-}space}(v) \wedge ((c \in \mathbb{C}) \wedge ((d \in \mathbb{C}) \wedge (x \in \operatorname{vec}(v))))) \Rightarrow  \\ \quad \quad \operatorname{cact}(v)(c, \operatorname{cact}(v)(d, x)) = \operatorname{cact}(v)(d, \operatorname{cact}(v)(c, x)) \end{array}}}
\vnbprop{cnvd-cact-vzero}{\vnbstmt{\begin{array}{@{}l@{}} \forall v, c. \\ \quad (\operatorname{is\text{-}complex\text{-}normed\text{-}vector\text{-}space}(v) \wedge (c \in \mathbb{C})) \Rightarrow  \\ \quad \quad \operatorname{cact}(v)(c, \operatorname{vzero}(v)) = \operatorname{vzero}(v) \end{array}}}
\vnbprop{cnvd-ms-is-ms}{\vnbstmt{\begin{array}{@{}l@{}} \forall v. \\ \quad \operatorname{is\text{-}complex\text{-}normed\text{-}vector\text{-}space}(v) \Rightarrow \operatorname{is\text{-}metric\text{-}space}(\operatorname{cnvs\text{-}metric\text{-}space}(v)) \end{array}}}
\vnbprop{cnvd-ms-pts}{\vnbstmt{\begin{array}{@{}l@{}} \forall v. \\ \quad \operatorname{is\text{-}complex\text{-}normed\text{-}vector\text{-}space}(v) \Rightarrow (\operatorname{pts}(\operatorname{cnvs\text{-}metric\text{-}space}(v)) \simeq \operatorname{vec}(v)) \end{array}}}
\vnbprop{cbops-unfold}{\vnbstmt{\begin{array}{@{}l@{}} \forall v, w. \\ \quad (\operatorname{cnvs\text{-}bounded\text{-}ops}(v, w) \simeq \{t \in (\operatorname{vec}(v) \to \operatorname{vec}(w)) : \operatorname{is\text{-}cnvs\text{-}bounded\text{-}op}(v, w, t)\}) \end{array}}}
\vnbprop{cbops-membership}{\vnbstmt{\begin{array}{@{}l@{}} \forall v, w, t. \\ \quad (t \in \operatorname{cnvs\text{-}bounded\text{-}ops}(v, w)) \Leftrightarrow ((t \in (\operatorname{vec}(v) \to \operatorname{vec}(w))) \wedge \operatorname{is\text{-}cnvs\text{-}bounded\text{-}op}(v, w, t)) \end{array}}}
\vnbprop{cbop-fun}{\vnbstmt{\forall v, w.; \forall t \in \operatorname{cnvs\text{-}bounded\text{-}ops}(v, w).\; (t \in (\operatorname{vec}(v) \to \operatorname{vec}(w)))}}
\vnbprop{cbop-linear}{\vnbstmt{\forall v, w.; \forall t \in \operatorname{cnvs\text{-}bounded\text{-}ops}(v, w).\; \operatorname{is\text{-}cnvs\text{-}linear\text{-}op}(v, w, t)}}
\vnbprop{cbop-bound-exists}{\vnbstmt{\begin{array}{@{}l@{}} \forall v, w.; \forall t \in \operatorname{cnvs\text{-}bounded\text{-}ops}(v, w). \exists m \in \mathbb{R}. \\ \quad (0 \leq m) \wedge \forall x \in \operatorname{vec}(v).\; (\operatorname{vnrm}(w)(t(x)) \leq (m \cdot \operatorname{vnrm}(v)(x))) \end{array}}}
\vnbprop{cbop-add}{\vnbstmt{\begin{array}{@{}l@{}} \forall v, w, t, x, y. \\ \quad ((t \in \operatorname{cnvs\text{-}bounded\text{-}ops}(v, w)) \wedge ((x \in \operatorname{vec}(v)) \wedge (y \in \operatorname{vec}(v)))) \Rightarrow  \\ \quad \quad t(\operatorname{vadd}(v)(x, y)) = \operatorname{vadd}(w)(t(x), t(y)) \end{array}}}
\vnbprop{cbop-act}{\vnbstmt{\begin{array}{@{}l@{}} \forall v, w, t, c, x. \\ \quad ((t \in \operatorname{cnvs\text{-}bounded\text{-}ops}(v, w)) \wedge ((c \in \mathbb{C}) \wedge (x \in \operatorname{vec}(v)))) \Rightarrow  \\ \quad \quad t(\operatorname{cact}(v)(c, x)) = \operatorname{cact}(w)(c, t(x)) \end{array}}}
\vnbprop{cbop-intro}{\vnbstmt{\begin{array}{@{}l@{}} \forall v, w, t, m. \\ \quad \operatorname{is\text{-}cnvs\text{-}linear\text{-}op}(v, w, t) \wedge \\ \quad m \in \mathbb{R} \wedge \\ \quad 0 \leq m \wedge \\ \quad \forall x \in \operatorname{vec}(v).\; (\operatorname{vnrm}(w)(t(x)) \leq (m \cdot \operatorname{vnrm}(v)(x))) \Rightarrow \\ \quad \quad t \in \operatorname{cnvs\text{-}bounded\text{-}ops}(v, w) \end{array}}}
\vnbprop{cnvs-opnorm-spec}{\vnbstmt{\begin{array}{@{}l@{}} \forall v, w, t. \\ \quad \operatorname{is\text{-}complex\text{-}normed\text{-}vector\text{-}space}(v) \wedge \\ \quad \operatorname{is\text{-}complex\text{-}normed\text{-}vector\text{-}space}(w) \wedge \\ \quad t \in \operatorname{cnvs\text{-}bounded\text{-}ops}(v, w) \Rightarrow \\ \quad \quad \operatorname{cnvs\text{-}opnorm}(v, w, t) \in \mathbb{R} \wedge \\ \quad \quad 0 \leq \operatorname{cnvs\text{-}opnorm}(v, w, t) \wedge \\ \quad \quad \forall x \in \operatorname{vec}(v).\; (\operatorname{vnrm}(w)(t(x)) \leq (\operatorname{cnvs\text{-}opnorm}(v, w, t) \cdot \operatorname{vnrm}(v)(x))) \wedge \\ \quad \quad \forall d \in \mathbb{R}. \\ \quad \quad \quad ((0 \leq d) \wedge \forall x \in \operatorname{vec}(v).\; (\operatorname{vnrm}(w)(t(x)) \leq (d \cdot \operatorname{vnrm}(v)(x)))) \Rightarrow  \\ \quad \quad \quad \quad \operatorname{cnvs\text{-}opnorm}(v, w, t) \leq d \end{array}}}
\vnbprop{cnvs-opnorm-in-rr}{\vnbstmt{\begin{array}{@{}l@{}} \forall v, w, t. \\ \quad \operatorname{is\text{-}complex\text{-}normed\text{-}vector\text{-}space}(v) \wedge \\ \quad \operatorname{is\text{-}complex\text{-}normed\text{-}vector\text{-}space}(w) \wedge \\ \quad t \in \operatorname{cnvs\text{-}bounded\text{-}ops}(v, w) \Rightarrow \\ \quad \quad \operatorname{cnvs\text{-}opnorm}(v, w, t) \in \mathbb{R} \end{array}}}
\vnbprop{cnvs-opnorm-nonneg}{\vnbstmt{\begin{array}{@{}l@{}} \forall v, w, t. \\ \quad \operatorname{is\text{-}complex\text{-}normed\text{-}vector\text{-}space}(v) \wedge \\ \quad \operatorname{is\text{-}complex\text{-}normed\text{-}vector\text{-}space}(w) \wedge \\ \quad t \in \operatorname{cnvs\text{-}bounded\text{-}ops}(v, w) \Rightarrow \\ \quad \quad 0 \leq \operatorname{cnvs\text{-}opnorm}(v, w, t) \end{array}}}
\vnbprop{cnvs-opnorm-bound}{\vnbstmt{\begin{array}{@{}l@{}} \forall v, w, t.; \forall x \in \operatorname{vec}(v). \\ \quad \operatorname{is\text{-}complex\text{-}normed\text{-}vector\text{-}space}(v) \wedge \\ \quad \operatorname{is\text{-}complex\text{-}normed\text{-}vector\text{-}space}(w) \wedge \\ \quad t \in \operatorname{cnvs\text{-}bounded\text{-}ops}(v, w) \Rightarrow \\ \quad \quad \operatorname{vnrm}(w)(t(x)) \leq (\operatorname{cnvs\text{-}opnorm}(v, w, t) \cdot \operatorname{vnrm}(v)(x)) \end{array}}}
\vnbprop{cbop-sum-apply}{\vnbstmt{\begin{array}{@{}l@{}} \forall v, w, s, t, x. \\ \quad \operatorname{is\text{-}complex\text{-}normed\text{-}vector\text{-}space}(w) \wedge \\ \quad s \in (\operatorname{vec}(v) \to \operatorname{vec}(w)) \wedge \\ \quad t \in (\operatorname{vec}(v) \to \operatorname{vec}(w)) \wedge \\ \quad x \in \operatorname{vec}(v) \Rightarrow \\ \quad \quad \operatorname{cbop\text{-}sum}(v, w, s, t)(x) = \operatorname{vadd}(w)(s(x), t(x)) \end{array}}}
\vnbprop{cbop-neg-apply}{\vnbstmt{\begin{array}{@{}l@{}} \forall v, w, t, x. \\ \quad (\operatorname{is\text{-}complex\text{-}normed\text{-}vector\text{-}space}(w) \wedge ((t \in (\operatorname{vec}(v) \to \operatorname{vec}(w))) \wedge (x \in \operatorname{vec}(v)))) \Rightarrow  \\ \quad \quad \operatorname{cbop\text{-}neg}(v, w, t)(x) = \operatorname{vneg}(w)(t(x)) \end{array}}}
\vnbprop{cbop-scale-apply}{\vnbstmt{\begin{array}{@{}l@{}} \forall v, w, c, t, x. \\ \quad (\operatorname{is\text{-}complex\text{-}normed\text{-}vector\text{-}space}(w) \wedge ((c \in \mathbb{C}) \wedge ((t \in (\operatorname{vec}(v) \to \operatorname{vec}(w))) \wedge (x \in \operatorname{vec}(v))))) \Rightarrow  \\ \quad \quad \operatorname{cbop\text{-}scale}(v, w, c, t)(x) = \operatorname{cact}(w)(c, t(x)) \end{array}}}
\vnbprop{cbop-zero-apply}{\vnbstmt{\begin{array}{@{}l@{}} \forall v, w, x. \\ \quad (\operatorname{is\text{-}complex\text{-}normed\text{-}vector\text{-}space}(w) \wedge (x \in \operatorname{vec}(v))) \Rightarrow  \\ \quad \quad \operatorname{cbop\text{-}zero}(v, w)(x) = \operatorname{vzero}(w) \end{array}}}
\vnbprop{cbop-sum-bounded}{\vnbstmt{\begin{array}{@{}l@{}} \forall v, w, s, t. \\ \quad \operatorname{is\text{-}complex\text{-}normed\text{-}vector\text{-}space}(v) \wedge \\ \quad \operatorname{is\text{-}complex\text{-}normed\text{-}vector\text{-}space}(w) \wedge \\ \quad s \in \operatorname{cnvs\text{-}bounded\text{-}ops}(v, w) \wedge \\ \quad t \in \operatorname{cnvs\text{-}bounded\text{-}ops}(v, w) \Rightarrow \\ \quad \quad \operatorname{cbop\text{-}sum}(v, w, s, t) \in \operatorname{cnvs\text{-}bounded\text{-}ops}(v, w) \end{array}}}
\vnbprop{cbop-neg-bounded}{\vnbstmt{\begin{array}{@{}l@{}} \forall v, w, t. \\ \quad \operatorname{is\text{-}complex\text{-}normed\text{-}vector\text{-}space}(v) \wedge \\ \quad \operatorname{is\text{-}complex\text{-}normed\text{-}vector\text{-}space}(w) \wedge \\ \quad t \in \operatorname{cnvs\text{-}bounded\text{-}ops}(v, w) \Rightarrow \\ \quad \quad \operatorname{cbop\text{-}neg}(v, w, t) \in \operatorname{cnvs\text{-}bounded\text{-}ops}(v, w) \end{array}}}
\vnbprop{cbop-scale-bounded}{\vnbstmt{\begin{array}{@{}l@{}} \forall v, w, c, t. \\ \quad \operatorname{is\text{-}complex\text{-}normed\text{-}vector\text{-}space}(v) \wedge \\ \quad \operatorname{is\text{-}complex\text{-}normed\text{-}vector\text{-}space}(w) \wedge \\ \quad c \in \mathbb{C} \wedge \\ \quad t \in \operatorname{cnvs\text{-}bounded\text{-}ops}(v, w) \Rightarrow \\ \quad \quad \operatorname{cbop\text{-}scale}(v, w, c, t) \in \operatorname{cnvs\text{-}bounded\text{-}ops}(v, w) \end{array}}}
\vnbprop{cbop-zero-bounded}{\vnbstmt{\begin{array}{@{}l@{}} \forall v, w. \\ \quad (\operatorname{is\text{-}complex\text{-}normed\text{-}vector\text{-}space}(v) \wedge \operatorname{is\text{-}complex\text{-}normed\text{-}vector\text{-}space}(w)) \Rightarrow  \\ \quad \quad \operatorname{cbop\text{-}zero}(v, w) \in \operatorname{cnvs\text{-}bounded\text{-}ops}(v, w) \end{array}}}
\vnbprop{cbop-sum-assoc}{\vnbstmt{\begin{array}{@{}l@{}} \forall v, w, r, s, t. \\ \quad \operatorname{is\text{-}complex\text{-}normed\text{-}vector\text{-}space}(v) \wedge \\ \quad \operatorname{is\text{-}complex\text{-}normed\text{-}vector\text{-}space}(w) \wedge \\ \quad r \in \operatorname{cnvs\text{-}bounded\text{-}ops}(v, w) \wedge \\ \quad s \in \operatorname{cnvs\text{-}bounded\text{-}ops}(v, w) \wedge \\ \quad t \in \operatorname{cnvs\text{-}bounded\text{-}ops}(v, w) \Rightarrow \\ \quad \quad \operatorname{cbop\text{-}sum}(v, w, \operatorname{cbop\text{-}sum}(v, w, r, s), t) = \operatorname{cbop\text{-}sum}(v, w, r, \operatorname{cbop\text{-}sum}(v, w, s, t)) \end{array}}}
\vnbprop{cbop-sum-comm}{\vnbstmt{\begin{array}{@{}l@{}} \forall v, w, s, t. \\ \quad \operatorname{is\text{-}complex\text{-}normed\text{-}vector\text{-}space}(v) \wedge \\ \quad \operatorname{is\text{-}complex\text{-}normed\text{-}vector\text{-}space}(w) \wedge \\ \quad s \in \operatorname{cnvs\text{-}bounded\text{-}ops}(v, w) \wedge \\ \quad t \in \operatorname{cnvs\text{-}bounded\text{-}ops}(v, w) \Rightarrow \\ \quad \quad \operatorname{cbop\text{-}sum}(v, w, s, t) = \operatorname{cbop\text{-}sum}(v, w, t, s) \end{array}}}
\vnbprop{cbop-sum-zero-left}{\vnbstmt{\begin{array}{@{}l@{}} \forall v, w, t. \\ \quad \operatorname{is\text{-}complex\text{-}normed\text{-}vector\text{-}space}(v) \wedge \\ \quad \operatorname{is\text{-}complex\text{-}normed\text{-}vector\text{-}space}(w) \wedge \\ \quad t \in \operatorname{cnvs\text{-}bounded\text{-}ops}(v, w) \Rightarrow \\ \quad \quad \operatorname{cbop\text{-}sum}(v, w, \operatorname{cbop\text{-}zero}(v, w), t) = t \end{array}}}
\vnbprop{cbop-sum-zero-right}{\vnbstmt{\begin{array}{@{}l@{}} \forall v, w, t. \\ \quad \operatorname{is\text{-}complex\text{-}normed\text{-}vector\text{-}space}(v) \wedge \\ \quad \operatorname{is\text{-}complex\text{-}normed\text{-}vector\text{-}space}(w) \wedge \\ \quad t \in \operatorname{cnvs\text{-}bounded\text{-}ops}(v, w) \Rightarrow \\ \quad \quad \operatorname{cbop\text{-}sum}(v, w, t, \operatorname{cbop\text{-}zero}(v, w)) = t \end{array}}}
\vnbprop{cbop-sum-neg-left}{\vnbstmt{\begin{array}{@{}l@{}} \forall v, w, t. \\ \quad \operatorname{is\text{-}complex\text{-}normed\text{-}vector\text{-}space}(v) \wedge \\ \quad \operatorname{is\text{-}complex\text{-}normed\text{-}vector\text{-}space}(w) \wedge \\ \quad t \in \operatorname{cnvs\text{-}bounded\text{-}ops}(v, w) \Rightarrow \\ \quad \quad \operatorname{cbop\text{-}sum}(v, w, \operatorname{cbop\text{-}neg}(v, w, t), t) = \operatorname{cbop\text{-}zero}(v, w) \end{array}}}
\vnbprop{cbop-sum-neg-right}{\vnbstmt{\begin{array}{@{}l@{}} \forall v, w, t. \\ \quad \operatorname{is\text{-}complex\text{-}normed\text{-}vector\text{-}space}(v) \wedge \\ \quad \operatorname{is\text{-}complex\text{-}normed\text{-}vector\text{-}space}(w) \wedge \\ \quad t \in \operatorname{cnvs\text{-}bounded\text{-}ops}(v, w) \Rightarrow \\ \quad \quad \operatorname{cbop\text{-}sum}(v, w, t, \operatorname{cbop\text{-}neg}(v, w, t)) = \operatorname{cbop\text{-}zero}(v, w) \end{array}}}
\vnbprop{cbop-scale-sum}{\vnbstmt{\begin{array}{@{}l@{}} \forall v, w, c, s, t. \\ \quad \operatorname{is\text{-}complex\text{-}normed\text{-}vector\text{-}space}(v) \wedge \\ \quad \operatorname{is\text{-}complex\text{-}normed\text{-}vector\text{-}space}(w) \wedge \\ \quad c \in \mathbb{C} \wedge \\ \quad s \in \operatorname{cnvs\text{-}bounded\text{-}ops}(v, w) \wedge \\ \quad t \in \operatorname{cnvs\text{-}bounded\text{-}ops}(v, w) \Rightarrow \\ \quad \quad \operatorname{cbop\text{-}scale}(v, w, c, \operatorname{cbop\text{-}sum}(v, w, s, t)) \\ \quad \quad =\; \operatorname{cbop\text{-}sum}(v, w, \operatorname{cbop\text{-}scale}(v, w, c, s), \operatorname{cbop\text{-}scale}(v, w, c, t)) \end{array}}}
\vnbprop{cbop-scale-add}{\vnbstmt{\begin{array}{@{}l@{}} \forall v, w, c, d, t. \\ \quad \operatorname{is\text{-}complex\text{-}normed\text{-}vector\text{-}space}(v) \wedge \\ \quad \operatorname{is\text{-}complex\text{-}normed\text{-}vector\text{-}space}(w) \wedge \\ \quad c \in \mathbb{C} \wedge \\ \quad d \in \mathbb{C} \wedge \\ \quad t \in \operatorname{cnvs\text{-}bounded\text{-}ops}(v, w) \Rightarrow \\ \quad \quad \operatorname{cbop\text{-}scale}(v, w, (c + d), t) = \operatorname{cbop\text{-}sum}(v, w, \operatorname{cbop\text{-}scale}(v, w, c, t), \operatorname{cbop\text{-}scale}(v, w, d, t)) \end{array}}}
\vnbprop{cbop-scale-mul}{\vnbstmt{\begin{array}{@{}l@{}} \forall v, w, c, d, t. \\ \quad \operatorname{is\text{-}complex\text{-}normed\text{-}vector\text{-}space}(v) \wedge \\ \quad \operatorname{is\text{-}complex\text{-}normed\text{-}vector\text{-}space}(w) \wedge \\ \quad c \in \mathbb{C} \wedge \\ \quad d \in \mathbb{C} \wedge \\ \quad t \in \operatorname{cnvs\text{-}bounded\text{-}ops}(v, w) \Rightarrow \\ \quad \quad \operatorname{cbop\text{-}scale}(v, w, (c \cdot d), t) = \operatorname{cbop\text{-}scale}(v, w, c, \operatorname{cbop\text{-}scale}(v, w, d, t)) \end{array}}}
\vnbprop{cbop-scale-one}{\vnbstmt{\begin{array}{@{}l@{}} \forall v, w, t. \\ \quad \operatorname{is\text{-}complex\text{-}normed\text{-}vector\text{-}space}(v) \wedge \\ \quad \operatorname{is\text{-}complex\text{-}normed\text{-}vector\text{-}space}(w) \wedge \\ \quad t \in \operatorname{cnvs\text{-}bounded\text{-}ops}(v, w) \Rightarrow \\ \quad \quad \operatorname{cbop\text{-}scale}(v, w, 1, t) = t \end{array}}}
\vnbprop{cnvs-opnorm-zero}{\vnbstmt{\begin{array}{@{}l@{}} \forall v, w. \\ \quad (\operatorname{is\text{-}complex\text{-}normed\text{-}vector\text{-}space}(v) \wedge \operatorname{is\text{-}complex\text{-}normed\text{-}vector\text{-}space}(w)) \Rightarrow  \\ \quad \quad \operatorname{cnvs\text{-}opnorm}(v, w, \operatorname{cbop\text{-}zero}(v, w)) = 0 \end{array}}}
\vnbprop{cnvs-opnorm-zero-imp}{\vnbstmt{\begin{array}{@{}l@{}} \forall v, w, t. \\ \quad \operatorname{is\text{-}complex\text{-}normed\text{-}vector\text{-}space}(v) \wedge \\ \quad \operatorname{is\text{-}complex\text{-}normed\text{-}vector\text{-}space}(w) \wedge \\ \quad t \in \operatorname{cnvs\text{-}bounded\text{-}ops}(v, w) \wedge \\ \quad \operatorname{cnvs\text{-}opnorm}(v, w, t) = 0 \Rightarrow \\ \quad \quad t = \operatorname{cbop\text{-}zero}(v, w) \end{array}}}
\vnbprop{cnvs-opnorm-definite}{\vnbstmt{\begin{array}{@{}l@{}} \forall v, w, t. \\ \quad \operatorname{is\text{-}complex\text{-}normed\text{-}vector\text{-}space}(v) \wedge \\ \quad \operatorname{is\text{-}complex\text{-}normed\text{-}vector\text{-}space}(w) \wedge \\ \quad t \in \operatorname{cnvs\text{-}bounded\text{-}ops}(v, w) \Rightarrow \\ \quad \quad (\operatorname{cnvs\text{-}opnorm}(v, w, t) = 0) \Leftrightarrow (t = \operatorname{cbop\text{-}zero}(v, w)) \end{array}}}
\vnbprop{cnvs-opnorm-triangle}{\vnbstmt{\begin{array}{@{}l@{}} \forall v, w, s, t. \\ \quad \operatorname{is\text{-}complex\text{-}normed\text{-}vector\text{-}space}(v) \wedge \\ \quad \operatorname{is\text{-}complex\text{-}normed\text{-}vector\text{-}space}(w) \wedge \\ \quad s \in \operatorname{cnvs\text{-}bounded\text{-}ops}(v, w) \wedge \\ \quad t \in \operatorname{cnvs\text{-}bounded\text{-}ops}(v, w) \Rightarrow \\ \quad \quad \operatorname{cnvs\text{-}opnorm}(v, w, \operatorname{cbop\text{-}sum}(v, w, s, t)) \leq (\operatorname{cnvs\text{-}opnorm}(v, w, s) + \operatorname{cnvs\text{-}opnorm}(v, w, t)) \end{array}}}
\vnbprop{cnvs-opnorm-neg}{\vnbstmt{\begin{array}{@{}l@{}} \forall v, w, t. \\ \quad \operatorname{is\text{-}complex\text{-}normed\text{-}vector\text{-}space}(v) \wedge \\ \quad \operatorname{is\text{-}complex\text{-}normed\text{-}vector\text{-}space}(w) \wedge \\ \quad t \in \operatorname{cnvs\text{-}bounded\text{-}ops}(v, w) \Rightarrow \\ \quad \quad \operatorname{cnvs\text{-}opnorm}(v, w, \operatorname{cbop\text{-}neg}(v, w, t)) = \operatorname{cnvs\text{-}opnorm}(v, w, t) \end{array}}}
\vnbprop{cnvs-opnorm-scale}{\vnbstmt{\begin{array}{@{}l@{}} \forall v, w, c, t. \\ \quad \operatorname{is\text{-}complex\text{-}normed\text{-}vector\text{-}space}(v) \wedge \\ \quad \operatorname{is\text{-}complex\text{-}normed\text{-}vector\text{-}space}(w) \wedge \\ \quad c \in \mathbb{C} \wedge \\ \quad t \in \operatorname{cnvs\text{-}bounded\text{-}ops}(v, w) \Rightarrow \\ \quad \quad \operatorname{cnvs\text{-}opnorm}(v, w, \operatorname{cbop\text{-}scale}(v, w, c, t)) = (\operatorname{magnitude}(c) \cdot \operatorname{cnvs\text{-}opnorm}(v, w, t)) \end{array}}}
\vnbprop{cbops-is-set}{\vnbstmt{\begin{array}{@{}l@{}} \forall v, w. \\ \quad (\operatorname{is\text{-}complex\text{-}normed\text{-}vector\text{-}space}(v) \wedge \operatorname{is\text{-}complex\text{-}normed\text{-}vector\text{-}space}(w)) \Rightarrow  \\ \quad \quad \operatorname{cnvs\text{-}bounded\text{-}ops}(v, w) \in \mathbf{Set} \end{array}}}
\vnbprop{cbop-nag-carr}{\vnbstmt{\begin{array}{@{}l@{}} \forall v, w. \\ \quad (\operatorname{is\text{-}complex\text{-}normed\text{-}vector\text{-}space}(v) \wedge \operatorname{is\text{-}complex\text{-}normed\text{-}vector\text{-}space}(w)) \Rightarrow  \\ \quad \quad (\operatorname{carr}(\operatorname{cbop\text{-}normed\text{-}ag}(v, w)) \simeq \operatorname{cnvs\text{-}bounded\text{-}ops}(v, w)) \end{array}}}
\vnbprop{cbop-nag-opr}{\vnbstmt{\begin{array}{@{}l@{}} \forall v, w. \\ \quad (\operatorname{is\text{-}complex\text{-}normed\text{-}vector\text{-}space}(v) \wedge \operatorname{is\text{-}complex\text{-}normed\text{-}vector\text{-}space}(w)) \Rightarrow  \\ \quad \quad (\operatorname{opr}(\operatorname{cbop\text{-}normed\text{-}ag}(v, w)) \simeq \operatorname{cbop\text{-}vadd}(v, w)) \end{array}}}
\vnbprop{cbop-nag-iden}{\vnbstmt{\begin{array}{@{}l@{}} \forall v, w. \\ \quad (\operatorname{is\text{-}complex\text{-}normed\text{-}vector\text{-}space}(v) \wedge \operatorname{is\text{-}complex\text{-}normed\text{-}vector\text{-}space}(w)) \Rightarrow  \\ \quad \quad (\operatorname{iden}(\operatorname{cbop\text{-}normed\text{-}ag}(v, w)) \simeq \operatorname{cbop\text{-}zero}(v, w)) \end{array}}}
\vnbprop{cbop-nag-inv}{\vnbstmt{\begin{array}{@{}l@{}} \forall v, w. \\ \quad (\operatorname{is\text{-}complex\text{-}normed\text{-}vector\text{-}space}(v) \wedge \operatorname{is\text{-}complex\text{-}normed\text{-}vector\text{-}space}(w)) \Rightarrow  \\ \quad \quad (\operatorname{inv}(\operatorname{cbop\text{-}normed\text{-}ag}(v, w)) \simeq \operatorname{cbop\text{-}vinv}(v, w)) \end{array}}}
\vnbprop{cbop-nag-nrm}{\vnbstmt{\begin{array}{@{}l@{}} \forall v, w. \\ \quad (\operatorname{is\text{-}complex\text{-}normed\text{-}vector\text{-}space}(v) \wedge \operatorname{is\text{-}complex\text{-}normed\text{-}vector\text{-}space}(w)) \Rightarrow  \\ \quad \quad (\operatorname{nrm}(\operatorname{cbop\text{-}normed\text{-}ag}(v, w)) \simeq \operatorname{cbop\text{-}vnrm}(v, w)) \end{array}}}
\vnbprop{cbop-vadd-value}{\vnbstmt{\begin{array}{@{}l@{}} \forall v, w, s, t. \\ \quad \operatorname{is\text{-}complex\text{-}normed\text{-}vector\text{-}space}(v) \wedge \\ \quad \operatorname{is\text{-}complex\text{-}normed\text{-}vector\text{-}space}(w) \wedge \\ \quad s \in \operatorname{cnvs\text{-}bounded\text{-}ops}(v, w) \wedge \\ \quad t \in \operatorname{cnvs\text{-}bounded\text{-}ops}(v, w) \Rightarrow \\ \quad \quad \operatorname{cbop\text{-}vadd}(v, w)(s, t) = \operatorname{cbop\text{-}sum}(v, w, s, t) \end{array}}}
\vnbprop{cbop-vinv-value}{\vnbstmt{\begin{array}{@{}l@{}} \forall v, w, t. \\ \quad \operatorname{is\text{-}complex\text{-}normed\text{-}vector\text{-}space}(v) \wedge \\ \quad \operatorname{is\text{-}complex\text{-}normed\text{-}vector\text{-}space}(w) \wedge \\ \quad t \in \operatorname{cnvs\text{-}bounded\text{-}ops}(v, w) \Rightarrow \\ \quad \quad \operatorname{cbop\text{-}vinv}(v, w)(t) = \operatorname{cbop\text{-}neg}(v, w, t) \end{array}}}
\vnbprop{cbop-vnrm-value}{\vnbstmt{\begin{array}{@{}l@{}} \forall v, w, t. \\ \quad \operatorname{is\text{-}complex\text{-}normed\text{-}vector\text{-}space}(v) \wedge \\ \quad \operatorname{is\text{-}complex\text{-}normed\text{-}vector\text{-}space}(w) \wedge \\ \quad t \in \operatorname{cnvs\text{-}bounded\text{-}ops}(v, w) \Rightarrow \\ \quad \quad \operatorname{cbop\text{-}vnrm}(v, w)(t) = \operatorname{cnvs\text{-}opnorm}(v, w, t) \end{array}}}
\vnbprop{cbop-nag-length}{\vnbstmt{\begin{array}{@{}l@{}} \forall v, w. \\ \quad (\operatorname{is\text{-}complex\text{-}normed\text{-}vector\text{-}space}(v) \wedge \operatorname{is\text{-}complex\text{-}normed\text{-}vector\text{-}space}(w)) \Rightarrow  \\ \quad \quad \operatorname{length}(\operatorname{cbop\text{-}normed\text{-}ag}(v, w)) = 5 \end{array}}}
\vnbprop{cbop-normed-ag-is-normed-ag}{\vnbstmt{\begin{array}{@{}l@{}} \forall v, w. \\ \quad (\operatorname{is\text{-}complex\text{-}normed\text{-}vector\text{-}space}(v) \wedge \operatorname{is\text{-}complex\text{-}normed\text{-}vector\text{-}space}(w)) \Rightarrow  \\ \quad \quad \operatorname{is\text{-}normed\text{-}ag}(\operatorname{cbop\text{-}normed\text{-}ag}(v, w)) \end{array}}}
\vnbprop{cnvd-ms-converges-of-scaled}{\vnbstmt{\begin{array}{@{}l@{}} \forall s, m, f, g, l, \ell, c. \\ \quad \operatorname{is\text{-}metric\text{-}space}(s) \wedge \\ \quad f \in (\mathbb{N} \to \operatorname{pts}(s)) \wedge \\ \quad l \in \operatorname{pts}(s) \wedge \\ \quad \operatorname{converges\text{-}to}(m, g, \ell) \wedge \\ \quad c \in \mathbb{R} \wedge \\ \quad 0 \leq c \wedge \\ \quad \forall n \in \mathbb{N}.\; (\operatorname{dist}(s)(f(n), l) \leq (c \cdot \operatorname{dist}(m)(g(n), \ell))) \Rightarrow \\ \quad \quad \operatorname{converges\text{-}to}(s, f, l) \end{array}}}
\vnbprop{cnvd-ms-cauchy-of-scaled}{\vnbstmt{\begin{array}{@{}l@{}} \forall s, \mathit{tm}, f, g, c. \\ \quad \operatorname{is\text{-}cauchy\text{-}seq}(s, f) \wedge \\ \quad \operatorname{is\text{-}metric\text{-}space}(\mathit{tm}) \wedge \\ \quad g \in (\mathbb{N} \to \operatorname{pts}(\mathit{tm})) \wedge \\ \quad c \in \mathbb{R} \wedge \\ \quad 0 \leq c \wedge \\ \quad \forall m, n \in \mathbb{N}.\; (\operatorname{dist}(\mathit{tm})(g(m), g(n)) \leq (c \cdot \operatorname{dist}(s)(f(m), f(n)))) \Rightarrow \\ \quad \quad \operatorname{is\text{-}cauchy\text{-}seq}(\mathit{tm}, g) \end{array}}}
\vnbprop{cnvd-ms-limit-dist-le}{\vnbstmt{\begin{array}{@{}l@{}} \forall s, g, l, a, b, n_{0}. \\ \quad \operatorname{converges\text{-}to}(s, g, l) \wedge \\ \quad a \in \operatorname{pts}(s) \wedge \\ \quad b \in \mathbb{R} \wedge \\ \quad n_{0} \in \mathbb{N} \wedge \\ \quad \forall n \in \mathbb{N}.\; ((n_{0} \leq n) \Rightarrow (\operatorname{dist}(s)(g(n), a) \leq b)) \Rightarrow \\ \quad \quad \operatorname{dist}(s)(l, a) \leq b \end{array}}}
\vnbprop{cnvdms-dist-norm}{\vnbstmt{\begin{array}{@{}l@{}} \forall v, a, b. \\ \quad (\operatorname{is\text{-}complex\text{-}normed\text{-}vector\text{-}space}(v) \wedge ((a \in \operatorname{vec}(v)) \wedge (b \in \operatorname{vec}(v)))) \Rightarrow  \\ \quad \quad \operatorname{dist}(\operatorname{cnvs\text{-}metric\text{-}space}(v))(a, b) = \operatorname{vnrm}(v)(\operatorname{vadd}(v)(a, \operatorname{vneg}(v)(b))) \end{array}}}
\vnbprop{cbop-nag-dist}{\vnbstmt{\begin{array}{@{}l@{}} \forall v, w, s, t. \\ \quad \operatorname{is\text{-}complex\text{-}normed\text{-}vector\text{-}space}(v) \wedge \\ \quad \operatorname{is\text{-}complex\text{-}normed\text{-}vector\text{-}space}(w) \wedge \\ \quad s \in \operatorname{cnvs\text{-}bounded\text{-}ops}(v, w) \wedge \\ \quad t \in \operatorname{cnvs\text{-}bounded\text{-}ops}(v, w) \Rightarrow \\ \quad \quad \operatorname{dist}(\operatorname{nag\text{-}metric\text{-}space}(\operatorname{cbop\text{-}normed\text{-}ag}(v, w)))(s, t) \\ \quad \quad =\; \operatorname{cnvs\text{-}opnorm}(v, w, \operatorname{cbop\text{-}sum}(v, w, s, \operatorname{cbop\text{-}neg}(v, w, t))) \end{array}}}
\vnbprop{cbop-pt-dist}{\vnbstmt{\begin{array}{@{}l@{}} \forall v, w, s, t, x. \\ \quad \operatorname{is\text{-}complex\text{-}normed\text{-}vector\text{-}space}(v) \wedge \\ \quad \operatorname{is\text{-}complex\text{-}normed\text{-}vector\text{-}space}(w) \wedge \\ \quad s \in \operatorname{cnvs\text{-}bounded\text{-}ops}(v, w) \wedge \\ \quad t \in \operatorname{cnvs\text{-}bounded\text{-}ops}(v, w) \wedge \\ \quad x \in \operatorname{vec}(v) \Rightarrow \\ \quad \quad \operatorname{dist}(\operatorname{cnvs\text{-}metric\text{-}space}(w))(s(x), t(x)) \\ \quad \quad \leq\; \operatorname{vnrm}(v)(x) \cdot \operatorname{dist}(\operatorname{nag\text{-}metric\text{-}space}(\operatorname{cbop\text{-}normed\text{-}ag}(v, w)))(s, t) \end{array}}}
\vnbprop{cnvdms-dist-sum-le}{\vnbstmt{\begin{array}{@{}l@{}} \forall w, a, b, c, e. \\ \quad (\operatorname{is\text{-}complex\text{-}normed\text{-}vector\text{-}space}(w) \wedge ((a \in \operatorname{vec}(w)) \wedge ((b \in \operatorname{vec}(w)) \wedge ((c \in \operatorname{vec}(w)) \wedge (e \in \operatorname{vec}(w)))))) \Rightarrow  \\ \quad \quad \operatorname{dist}(\operatorname{cnvs\text{-}metric\text{-}space}(w))(\operatorname{vadd}(w)(a, b), \operatorname{vadd}(w)(c, e)) \\ \quad \quad \leq\; \operatorname{dist}(\operatorname{cnvs\text{-}metric\text{-}space}(w))(a, c) + \operatorname{dist}(\operatorname{cnvs\text{-}metric\text{-}space}(w))(b, e) \end{array}}}
\vnbprop{cnvdms-dist-act-le}{\vnbstmt{\begin{array}{@{}l@{}} \forall w, c, a, b. \\ \quad (\operatorname{is\text{-}complex\text{-}normed\text{-}vector\text{-}space}(w) \wedge ((c \in \mathbb{C}) \wedge ((a \in \operatorname{vec}(w)) \wedge (b \in \operatorname{vec}(w))))) \Rightarrow  \\ \quad \quad \operatorname{dist}(\operatorname{cnvs\text{-}metric\text{-}space}(w))(\operatorname{cact}(w)(c, a), \operatorname{cact}(w)(c, b)) \\ \quad \quad \leq\; \operatorname{magnitude}(c) \cdot \operatorname{dist}(\operatorname{cnvs\text{-}metric\text{-}space}(w))(a, b) \end{array}}}
\vnbprop{cbop-cauchy-seq-fun}{\vnbstmt{\begin{array}{@{}l@{}} \forall v, w, f. \\ \quad \operatorname{is\text{-}complex\text{-}normed\text{-}vector\text{-}space}(v) \wedge \\ \quad \operatorname{is\text{-}complex\text{-}normed\text{-}vector\text{-}space}(w) \wedge \\ \quad \operatorname{is\text{-}cauchy\text{-}seq}(\operatorname{nag\text{-}metric\text{-}space}(\operatorname{cbop\text{-}normed\text{-}ag}(v, w)), f) \Rightarrow \\ \quad \quad f \in (\mathbb{N} \to \operatorname{cnvs\text{-}bounded\text{-}ops}(v, w)) \end{array}}}
\vnbprop{cbop-ptwise-seq-fun}{\vnbstmt{\begin{array}{@{}l@{}} \forall v, w, f, x. \\ \quad \operatorname{is\text{-}complex\text{-}normed\text{-}vector\text{-}space}(v) \wedge \\ \quad \operatorname{is\text{-}complex\text{-}normed\text{-}vector\text{-}space}(w) \wedge \\ \quad f \in (\mathbb{N} \to \operatorname{cnvs\text{-}bounded\text{-}ops}(v, w)) \wedge \\ \quad x \in \operatorname{vec}(v) \Rightarrow \\ \quad \quad (\lambda q \in \mathbb{N}.\; f(q)(x)) \in (\mathbb{N} \to \operatorname{pts}(\operatorname{cnvs\text{-}metric\text{-}space}(w))) \end{array}}}
\vnbprop{cbop-ptwise-cauchy}{\vnbstmt{\begin{array}{@{}l@{}} \forall v, w, f, x. \\ \quad \operatorname{is\text{-}complex\text{-}normed\text{-}vector\text{-}space}(v) \wedge \\ \quad \operatorname{is\text{-}complex\text{-}normed\text{-}vector\text{-}space}(w) \wedge \\ \quad \operatorname{is\text{-}cauchy\text{-}seq}(\operatorname{nag\text{-}metric\text{-}space}(\operatorname{cbop\text{-}normed\text{-}ag}(v, w)), f) \wedge \\ \quad x \in \operatorname{vec}(v) \Rightarrow \\ \quad \quad \operatorname{is\text{-}cauchy\text{-}seq}(\operatorname{cnvs\text{-}metric\text{-}space}(w), (\lambda q \in \mathbb{N}.\; f(q)(x))) \end{array}}}
\vnbprop{cbop-ptwise-limit}{\vnbstmt{\begin{array}{@{}l@{}} \forall v, w, f, x. \\ \quad \operatorname{is\text{-}complex\text{-}normed\text{-}vector\text{-}space}(v) \wedge \\ \quad \operatorname{is\text{-}complex\text{-}banach\text{-}space}(w) \wedge \\ \quad \operatorname{is\text{-}cauchy\text{-}seq}(\operatorname{nag\text{-}metric\text{-}space}(\operatorname{cbop\text{-}normed\text{-}ag}(v, w)), f) \wedge \\ \quad x \in \operatorname{vec}(v) \Rightarrow \\ \quad \quad \varepsilon(\{q \in \operatorname{vec}(w) : \operatorname{converges\text{-}to}(\operatorname{cnvs\text{-}metric\text{-}space}(w), (\lambda \mathit{kq} \in \mathbb{N}.\; f(\mathit{kq})(x)), q)\}) \\ \quad \quad \in\; \{q \in \operatorname{vec}(w) : \operatorname{converges\text{-}to}(\operatorname{cnvs\text{-}metric\text{-}space}(w), (\lambda \mathit{kq} \in \mathbb{N}.\; f(\mathit{kq})(x)), q)\} \end{array}}}
\vnbprop{cbop-complete}{\vnbstmt{\begin{array}{@{}l@{}} \forall v, w. \\ \quad (\operatorname{is\text{-}complex\text{-}normed\text{-}vector\text{-}space}(v) \wedge \operatorname{is\text{-}complex\text{-}banach\text{-}space}(w)) \Rightarrow  \\ \quad \quad \operatorname{is\text{-}complete}(\operatorname{nag\text{-}metric\text{-}space}(\operatorname{cbop\text{-}normed\text{-}ag}(v, w))) \end{array}}}

\vnbbookfile{theorem-library/cnvs-comp}{theorem-library/cnvs-comp}
\noindent{\small cnvs-comp.scm -- wave of 2026-11-09/10 (agent CNVS2), wired 2026-11-10 from scratchpad/cnvs2-1109/cnvs2-comp.scm. Library: analysis.}\par
\vnbprop{cbop-comp-apply}{\vnbstmt{\begin{array}{@{}l@{}} \forall u, v, w, s, t, x. \\ \quad ((t \in (\operatorname{vec}(u) \to \operatorname{vec}(v))) \wedge ((s \in (\operatorname{vec}(v) \to \operatorname{vec}(w))) \wedge (x \in \operatorname{vec}(u)))) \Rightarrow  \\ \quad \quad {s}{t}(x) = s(t(x)) \end{array}}}
\vnbprop{cbop-id-apply}{\vnbstmt{\forall u.; \forall x \in \operatorname{vec}(u).\; (1(x) = x)}}
\vnbprop{cbop-comp-fun}{\vnbstmt{\begin{array}{@{}l@{}} \forall u, v, w, s, t. \\ \quad (\operatorname{is\text{-}complex\text{-}normed\text{-}vector\text{-}space}(u) \wedge ((t \in (\operatorname{vec}(u) \to \operatorname{vec}(v))) \wedge (s \in (\operatorname{vec}(v) \to \operatorname{vec}(w))))) \Rightarrow  \\ \quad \quad {s}{t} \in (\operatorname{vec}(u) \to \operatorname{vec}(w)) \end{array}}}
\vnbprop{cbop-id-fun}{\vnbstmt{\begin{array}{@{}l@{}} \forall u. \\ \quad \operatorname{is\text{-}complex\text{-}normed\text{-}vector\text{-}space}(u) \Rightarrow (1 \in (\operatorname{vec}(u) \to \operatorname{vec}(u))) \end{array}}}
\vnbprop{cbop-comp-bounded}{\vnbstmt{\begin{array}{@{}l@{}} \forall u, v, w, s, t. \\ \quad \operatorname{is\text{-}complex\text{-}normed\text{-}vector\text{-}space}(u) \wedge \\ \quad \operatorname{is\text{-}complex\text{-}normed\text{-}vector\text{-}space}(v) \wedge \\ \quad \operatorname{is\text{-}complex\text{-}normed\text{-}vector\text{-}space}(w) \wedge \\ \quad t \in \operatorname{cnvs\text{-}bounded\text{-}ops}(u, v) \wedge \\ \quad s \in \operatorname{cnvs\text{-}bounded\text{-}ops}(v, w) \Rightarrow \\ \quad \quad {s}{t} \in \operatorname{cnvs\text{-}bounded\text{-}ops}(u, w) \end{array}}}
\vnbprop{cnvs-opnorm-comp-le}{\vnbstmt{\begin{array}{@{}l@{}} \forall u, v, w, s, t. \\ \quad \operatorname{is\text{-}complex\text{-}normed\text{-}vector\text{-}space}(u) \wedge \\ \quad \operatorname{is\text{-}complex\text{-}normed\text{-}vector\text{-}space}(v) \wedge \\ \quad \operatorname{is\text{-}complex\text{-}normed\text{-}vector\text{-}space}(w) \wedge \\ \quad t \in \operatorname{cnvs\text{-}bounded\text{-}ops}(u, v) \wedge \\ \quad s \in \operatorname{cnvs\text{-}bounded\text{-}ops}(v, w) \Rightarrow \\ \quad \quad \operatorname{cnvs\text{-}opnorm}(u, w, {s}{t}) \leq (\operatorname{cnvs\text{-}opnorm}(v, w, s) \cdot \operatorname{cnvs\text{-}opnorm}(u, v, t)) \end{array}}}
\vnbprop{cbop-id-bounded}{\vnbstmt{\begin{array}{@{}l@{}} \forall u. \\ \quad \operatorname{is\text{-}complex\text{-}normed\text{-}vector\text{-}space}(u) \Rightarrow (1 \in \operatorname{cnvs\text{-}bounded\text{-}ops}(u, u)) \end{array}}}
\vnbprop{cbop-comp-assoc}{\vnbstmt{\begin{array}{@{}l@{}} \forall u, v, w, z, r, s, t. \\ \quad \operatorname{is\text{-}complex\text{-}normed\text{-}vector\text{-}space}(u) \wedge \\ \quad \operatorname{is\text{-}complex\text{-}normed\text{-}vector\text{-}space}(v) \wedge \\ \quad t \in (\operatorname{vec}(u) \to \operatorname{vec}(v)) \wedge \\ \quad s \in (\operatorname{vec}(v) \to \operatorname{vec}(w)) \wedge \\ \quad r \in (\operatorname{vec}(w) \to \operatorname{vec}(z)) \Rightarrow \\ \quad \quad {{r}{s}}{t} = {r}{{s}{t}} \end{array}}}
\vnbprop{cbop-comp-id-left}{\vnbstmt{\begin{array}{@{}l@{}} \forall u, v, t. \\ \quad (\operatorname{is\text{-}complex\text{-}normed\text{-}vector\text{-}space}(u) \wedge (\operatorname{is\text{-}complex\text{-}normed\text{-}vector\text{-}space}(v) \wedge (t \in (\operatorname{vec}(u) \to \operatorname{vec}(v))))) \Rightarrow  \\ \quad \quad {1}{t} = t \end{array}}}
\vnbprop{cbop-comp-id-right}{\vnbstmt{\begin{array}{@{}l@{}} \forall u, v, t. \\ \quad (\operatorname{is\text{-}complex\text{-}normed\text{-}vector\text{-}space}(u) \wedge (t \in (\operatorname{vec}(u) \to \operatorname{vec}(v)))) \Rightarrow  \\ \quad \quad {t}{1} = t \end{array}}}
\vnbprop{cbop-comp-sum-right}{\vnbstmt{\begin{array}{@{}l@{}} \forall u, v, w, r, s, t. \\ \quad \operatorname{is\text{-}complex\text{-}normed\text{-}vector\text{-}space}(u) \wedge \\ \quad \operatorname{is\text{-}complex\text{-}normed\text{-}vector\text{-}space}(v) \wedge \\ \quad \operatorname{is\text{-}complex\text{-}normed\text{-}vector\text{-}space}(w) \wedge \\ \quad t \in \operatorname{cnvs\text{-}bounded\text{-}ops}(u, v) \wedge \\ \quad r \in \operatorname{cnvs\text{-}bounded\text{-}ops}(v, w) \wedge \\ \quad s \in \operatorname{cnvs\text{-}bounded\text{-}ops}(v, w) \Rightarrow \\ \quad \quad {\operatorname{cbop\text{-}sum}(v, w, r, s)}{t} = \operatorname{cbop\text{-}sum}(u, w, {r}{t}, {s}{t}) \end{array}}}
\vnbprop{cbop-comp-sum-left}{\vnbstmt{\begin{array}{@{}l@{}} \forall u, v, w, r, s, t. \\ \quad \operatorname{is\text{-}complex\text{-}normed\text{-}vector\text{-}space}(u) \wedge \\ \quad \operatorname{is\text{-}complex\text{-}normed\text{-}vector\text{-}space}(v) \wedge \\ \quad \operatorname{is\text{-}complex\text{-}normed\text{-}vector\text{-}space}(w) \wedge \\ \quad r \in \operatorname{cnvs\text{-}bounded\text{-}ops}(v, w) \wedge \\ \quad s \in \operatorname{cnvs\text{-}bounded\text{-}ops}(u, v) \wedge \\ \quad t \in \operatorname{cnvs\text{-}bounded\text{-}ops}(u, v) \Rightarrow \\ \quad \quad {r}{\operatorname{cbop\text{-}sum}(u, v, s, t)} = \operatorname{cbop\text{-}sum}(u, w, {r}{s}, {r}{t}) \end{array}}}

\vnbbookfile{theorem-library/cc-exhaust-radius-decrease}{theorem-library/cc-exhaust-radius-decrease}
\noindent{\small cc-exhaust-radius-decrease.scm -- lemmas MOVED 2026-10-30 (wave-2 integration) out of theorem-library/ca5-function-spaces.scm into library analysis, for the library gate: cited by measurable-maps-metric (measure); a fact about 1/(n+1). The proofs are the originals, verbatim; the helper defines are a copy of the ...}\par
\vnbprop{cc-exhaust-radius}{\vnbstmt{\forall n \in \mathbb{N}.\; ((\operatorname{recip}((n + 1)) \in \mathbb{R}) \wedge (0 < \operatorname{recip}((n + 1))))}}
\vnbprop{cc-exhaust-radius-decrease}{\vnbstmt{\forall n \in \mathbb{N}.\; ((\operatorname{recip}(((n + 1) + 1)) < \operatorname{recip}((n + 1))) \wedge (\operatorname{recip}((n + 1)) \leq 1))}}

\vnbbookfile{theorem-library/countable-fin-subsets}{theorem-library/countable-fin-subsets}
\noindent{\small countable-fin-subsets.scm -- wave of 2026-11-09/10 (agent BA-ONB4), wired 2026-11-10 from scratchpad/baonb4-1109/cf-1.scm. Library: analysis.}\par
\vnbprop{countable-fin-subsets}{\vnbstmt{\forall x.\; (\operatorname{is\text{-}countable}(x) \Rightarrow \operatorname{is\text{-}countable}(\operatorname{fin\text{-}subsets}(x)))}}

\vnbbookfile{theorem-library/cc-real-part-facts}{theorem-library/cc-real-part-facts}
\noindent{\small cc-real-part-facts.scm -- lemmas MOVED 2026-11-10 (integration of the 11-09/10 waves) out of theorem-library/ca-lft-cayley.scm into library analysis, for the library gate: cited by cc-real-bound; CC arithmetic. The proofs are the originals, verbatim; the helper defines are a copy of the source's. cc-real-of-im-zero ...}\par
\vnbprop{rr-imag-part-zero}{\vnbstmt{\forall x \in \mathbb{R}.\; (\operatorname{imag\text{-}part}(x) = 0)}}
\vnbprop{cc-real-of-im-zero}{\vnbstmt{\forall z \in \mathbb{C}.\; ((\operatorname{imag\text{-}part}(z) = 0) \Rightarrow (z \in \mathbb{R}))}}

\vnbbookfile{theorem-library/cc-im-i-mul}{theorem-library/cc-im-i-mul}
\noindent{\small cc-im-i-mul.scm -- lemmas MOVED 2026-11-10 (integration of the 11-09/10 waves) out of theorem-library/triangle-orientation.scm into library analysis, for the library gate: cited by cc-real-bound; CC arithmetic. The proofs are the originals, verbatim; the helper defines are a copy of the source's. cc-im-i-mul}\par
\vnbprop{cc-im-i-mul}{\vnbstmt{\forall z \in \mathbb{C}.\; (\operatorname{imag\text{-}part}((+i \cdot z)) = \operatorname{real\text{-}part}(z))}}

\vnbbookfile{theorem-library/cc-real-bound}{theorem-library/cc-real-bound}
\noindent{\small cc-real-bound.scm -- wave of 2026-11-09/10 (agent BA-2), wired 2026-11-10 from scratchpad/ba2-1109/cc-real-bound.scm. Library: analysis.}\par
\vnbprop{crb-linear-bound-zero}{\vnbstmt{\forall y, d \in \mathbb{R}.\; (\forall t \in \mathbb{R}.\; ((2 \cdot (y \cdot t)) \leq d) \Rightarrow (y = 0))}}
\vnbprop{cc-real-of-quadratic-bound}{\vnbstmt{\begin{array}{@{}l@{}} \forall z \in \mathbb{C}, c \in \mathbb{R}. \\ \quad \forall t \in \mathbb{R}.\; ((\operatorname{magnitude}((z + (+i \cdot t))) \cdot \operatorname{magnitude}((z + (+i \cdot t)))) \leq (c + (t \cdot t))) \Rightarrow  \\ \quad \quad z \in \mathbb{R} \end{array}}}

\vnbbookfile{structure-library/functoriality}{structure-library/functoriality}
\noindent{\small A file shared across books; the entries that belong to this one.}\par
\noindent{\small functoriality.scm -- every view-as is a FUNCTOR, and here is the proof.}\par
\vnbprop{cnvs-as-module-functorial}{\vnbstmt{\begin{array}{@{}l@{}} \forall a, b, f. \\ \quad \operatorname{is\text{-}hom\text{-}complex\text{-}normed\text{-}vector\text{-}space}(a, b, f) \Rightarrow  \\ \quad \quad \operatorname{is\text{-}hom\text{-}module}(\operatorname{cnvs\text{-}as\text{-}module}(a), \operatorname{cnvs\text{-}as\text{-}module}(b), f) \end{array}}}
\vnbprop{cnvs-as-normed-ag-functorial}{\vnbstmt{\begin{array}{@{}l@{}} \forall a, b, f. \\ \quad \operatorname{is\text{-}hom\text{-}complex\text{-}normed\text{-}vector\text{-}space}(a, b, f) \Rightarrow  \\ \quad \quad \operatorname{is\text{-}hom\text{-}normed\text{-}ag}(\operatorname{cnvs\text{-}as\text{-}normed\text{-}ag}(a), \operatorname{cnvs\text{-}as\text{-}normed\text{-}ag}(b), f) \end{array}}}
\vnbprop{cnvs-as-nvs-functorial}{\vnbstmt{\begin{array}{@{}l@{}} \forall a, b, f. \\ \quad \operatorname{is\text{-}hom\text{-}complex\text{-}normed\text{-}vector\text{-}space}(a, b, f) \Rightarrow  \\ \quad \quad \operatorname{is\text{-}hom\text{-}normed\text{-}vector\text{-}space}(\operatorname{cnvs\text{-}as\text{-}nvs}(a), \operatorname{cnvs\text{-}as\text{-}nvs}(b), f) \end{array}}}
\vnbprop{complex-inner-product-space-as-module-functorial}{\vnbstmt{\begin{array}{@{}l@{}} \forall a, b, f. \\ \quad \operatorname{is\text{-}hom\text{-}complex\text{-}inner\text{-}product\text{-}space}(a, b, f) \Rightarrow  \\ \quad \quad \operatorname{is\text{-}hom\text{-}module}(\operatorname{complex\text{-}inner\text{-}product\text{-}space\text{-}as\text{-}module}(a), \operatorname{complex\text{-}inner\text{-}product\text{-}space\text{-}as\text{-}module}(b), f) \end{array}}}
\vnbprop{csip-as-sip-functorial}{\vnbstmt{\begin{array}{@{}l@{}} \forall a, b, f. \\ \quad \operatorname{is\text{-}hom\text{-}complex\text{-}semi\text{-}inner\text{-}product\text{-}space}(a, b, f) \Rightarrow  \\ \quad \quad \operatorname{is\text{-}hom\text{-}semi\text{-}inner\text{-}product\text{-}space}(\operatorname{csip\text{-}as\text{-}sip}(a), \operatorname{csip\text{-}as\text{-}sip}(b), f) \end{array}}}
\vnbprop{normed-ag-as-abelian-group-functorial}{\vnbstmt{\begin{array}{@{}l@{}} \forall a, b, f. \\ \quad \operatorname{is\text{-}hom\text{-}normed\text{-}ag}(a, b, f) \Rightarrow  \\ \quad \quad \operatorname{is\text{-}hom\text{-}abelian\text{-}group}(\operatorname{normed\text{-}ag\text{-}as\text{-}abelian\text{-}group}(a), \operatorname{normed\text{-}ag\text{-}as\text{-}abelian\text{-}group}(b), f) \end{array}}}
\vnbprop{normed-field-additive-ag-functorial}{\vnbstmt{\begin{array}{@{}l@{}} \forall a, b, f. \\ \quad \operatorname{is\text{-}hom\text{-}normed\text{-}field}(a, b, f) \Rightarrow  \\ \quad \quad \operatorname{is\text{-}hom\text{-}abelian\text{-}group}(\operatorname{normed\text{-}field\text{-}additive\text{-}ag}(a), \operatorname{normed\text{-}field\text{-}additive\text{-}ag}(b), f) \end{array}}}
\vnbprop{normed-field-as-commutative-ring-functorial}{\vnbstmt{\begin{array}{@{}l@{}} \forall a, b, f. \\ \quad \operatorname{is\text{-}hom\text{-}normed\text{-}field}(a, b, f) \Rightarrow  \\ \quad \quad \operatorname{is\text{-}hom\text{-}commutative\text{-}ring}(\operatorname{normed\text{-}field\text{-}as\text{-}commutative\text{-}ring}(a), \operatorname{normed\text{-}field\text{-}as\text{-}commutative\text{-}ring}(b), f) \end{array}}}
\vnbprop{normed-field-as-integral-domain-functorial}{\vnbstmt{\begin{array}{@{}l@{}} \forall a, b, f. \\ \quad \operatorname{is\text{-}hom\text{-}normed\text{-}field}(a, b, f) \Rightarrow  \\ \quad \quad \operatorname{is\text{-}hom\text{-}integral\text{-}domain}(\operatorname{normed\text{-}field\text{-}as\text{-}integral\text{-}domain}(a), \operatorname{normed\text{-}field\text{-}as\text{-}integral\text{-}domain}(b), f) \end{array}}}
\vnbprop{normed-vector-space-as-module-functorial}{\vnbstmt{\begin{array}{@{}l@{}} \forall a, b, f. \\ \quad \operatorname{is\text{-}hom\text{-}normed\text{-}vector\text{-}space}(a, b, f) \Rightarrow  \\ \quad \quad \operatorname{is\text{-}hom\text{-}module}(\operatorname{normed\text{-}vector\text{-}space\text{-}as\text{-}module}(a), \operatorname{normed\text{-}vector\text{-}space\text{-}as\text{-}module}(b), f) \end{array}}}
\vnbprop{normed-vector-space-as-normed-ag-functorial}{\vnbstmt{\begin{array}{@{}l@{}} \forall a, b, f. \\ \quad \operatorname{is\text{-}hom\text{-}normed\text{-}vector\text{-}space}(a, b, f) \Rightarrow  \\ \quad \quad \operatorname{is\text{-}hom\text{-}normed\text{-}ag}(\operatorname{normed\text{-}vector\text{-}space\text{-}as\text{-}normed\text{-}ag}(a), \operatorname{normed\text{-}vector\text{-}space\text{-}as\text{-}normed\text{-}ag}(b), f) \end{array}}}
\vnbprop{sip-as-module-functorial}{\vnbstmt{\begin{array}{@{}l@{}} \forall a, b, f. \\ \quad \operatorname{is\text{-}hom\text{-}semi\text{-}inner\text{-}product\text{-}space}(a, b, f) \Rightarrow  \\ \quad \quad \operatorname{is\text{-}hom\text{-}module}(\operatorname{sip\text{-}as\text{-}module}(a), \operatorname{sip\text{-}as\text{-}module}(b), f) \end{array}}}

\vnbbookfile{theorem-library/l2-families}{theorem-library/l2-families}
\noindent{\small l2-families.scm -- wave of 2026-11-09/10 (agent BA-ONB2), wired 2026-11-10 from scratchpad/baonb2-1109/bo2-l2a.scm. Library: analysis.}\par
\vnbprop{cc-nag-is-normed-ag}{\vnbstmt{\operatorname{is\text{-}normed\text{-}ag}(\operatorname{cc\text{-}nag})}}
\vnbprop{cc-nag-ag-carr}{\vnbstmt{(\operatorname{carr}(\operatorname{normed\text{-}ag\text{-}as\text{-}abelian\text{-}group}(\operatorname{cc\text{-}nag})) \simeq \mathbb{C})}}
\vnbprop{cc-nag-ag-iden}{\vnbstmt{(\operatorname{iden}(\operatorname{normed\text{-}ag\text{-}as\text{-}abelian\text{-}group}(\operatorname{cc\text{-}nag})) \simeq 0)}}
\vnbprop{cc-nag-ag-opr}{\vnbstmt{\begin{array}{@{}l@{}} \forall a, b \in \mathbb{C}. \\ \quad \operatorname{opr}(\operatorname{normed\text{-}ag\text{-}as\text{-}abelian\text{-}group}(\operatorname{cc\text{-}nag}))(a, b) = (a + b) \end{array}}}
\vnbprop{cc-nag-ag-inv}{\vnbstmt{\forall a \in \mathbb{C}.\; (\operatorname{inv}(\operatorname{normed\text{-}ag\text{-}as\text{-}abelian\text{-}group}(\operatorname{cc\text{-}nag}))(a) = -a)}}
\vnbprop{cc-nag-nrm}{\vnbstmt{\forall a \in \mathbb{C}.\; (\operatorname{nrm}(\operatorname{cc\text{-}nag})(a) = \operatorname{magnitude}(a))}}
\vnbprop{cc-nag-ag-is-abelian-group}{\vnbstmt{\operatorname{is\text{-}abelian\text{-}group}(\operatorname{normed\text{-}ag\text{-}as\text{-}abelian\text{-}group}(\operatorname{cc\text{-}nag}))}}
\vnbprop{finsum-hom-ag}{\vnbstmt{\begin{array}{@{}l@{}} \forall a, b, h.; \forall s \in \mathbf{Set}, f. \\ \quad (\operatorname{is\text{-}hom\text{-}abelian\text{-}group}(a, b, h) \wedge ((|s| \in \mathbb{N}) \wedge \forall z \in s.\; (f(z) \in \operatorname{carr}(a)))) \Rightarrow  \\ \quad \quad h(\operatorname{finsum}(a, f, s)) = \operatorname{finsum}(b, (\lambda z \in s.\; h(f(z))), s) \end{array}}}
\vnbprop{sums-to-uniform-hom}{\vnbstmt{\begin{array}{@{}l@{}} \forall a, b, h.; \forall f \in (\operatorname{dom}(f) \to \operatorname{carr}(a)), r \in \operatorname{carr}(a). \\ \quad \operatorname{is\text{-}normed\text{-}ag}(a) \wedge \\ \quad \operatorname{is\text{-}normed\text{-}ag}(b) \wedge \\ \quad \operatorname{is\text{-}hom\text{-}abelian\text{-}group}(\operatorname{normed\text{-}ag\text{-}as\text{-}abelian\text{-}group}(a), \operatorname{normed\text{-}ag\text{-}as\text{-}abelian\text{-}group}(b), h) \wedge \\ \quad \forall s. \\ \quad \quad \operatorname{pos\text{-}rr}(s) \Rightarrow  \\ \quad \quad \quad \exists l. \\ \quad \quad \quad \quad \operatorname{pos\text{-}rr}(l) \wedge \\ \quad \quad \quad \quad \forall x, y \in \operatorname{carr}(a). \\ \quad \quad \quad \quad \quad (\operatorname{nrm}(a)(\operatorname{opr}(a)(x, \operatorname{inv}(a)(y))) < l) \Rightarrow  \\ \quad \quad \quad \quad \quad \quad \operatorname{nrm}(b)(\operatorname{opr}(b)(h(x), \operatorname{inv}(b)(h(y)))) < s \wedge \\ \quad \operatorname{sums\text{-}to}(a, f, r) \Rightarrow \\ \quad \quad \operatorname{sums\text{-}to}(b, (\lambda z \in \operatorname{dom}(f).\; h(f(z))), h(r)) \end{array}}}
\vnbprop{ip-functional-hom}{\vnbstmt{\begin{array}{@{}l@{}} \forall v.; \forall a \in \operatorname{vec}(v). \\ \quad \operatorname{is\text{-}complex\text{-}inner\text{-}product\text{-}space}(v) \Rightarrow  \\ \quad \quad \operatorname{is\text{-}hom\text{-}abelian\text{-}group}(\operatorname{normed\text{-}ag\text{-}as\text{-}abelian\text{-}group}(\operatorname{ip\text{-}normed\text{-}ag}(v)), \\ \quad \quad \quad \operatorname{normed\text{-}ag\text{-}as\text{-}abelian\text{-}group}(\operatorname{cc\text{-}nag}), \\ \quad \quad \quad (\lambda w \in \operatorname{vec}(v).\; \langle\cdot,\cdot\rangle_{v}(w, a))) \end{array}}}
\vnbprop{ip-functional-ucont}{\vnbstmt{\begin{array}{@{}l@{}} \forall v.; \forall a \in \operatorname{vec}(v), s. \\ \quad (\operatorname{is\text{-}complex\text{-}inner\text{-}product\text{-}space}(v) \wedge \operatorname{pos\text{-}rr}(s)) \Rightarrow  \\ \quad \quad \exists l. \\ \quad \quad \quad \operatorname{pos\text{-}rr}(l) \wedge \\ \quad \quad \quad \forall x, y \in \operatorname{carr}(\operatorname{ip\text{-}normed\text{-}ag}(v)). \\ \quad \quad \quad \quad (\operatorname{nrm}(\operatorname{ip\text{-}normed\text{-}ag}(v))(\operatorname{opr}(\operatorname{ip\text{-}normed\text{-}ag}(v))(x, \operatorname{inv}(\operatorname{ip\text{-}normed\text{-}ag}(v))(y))) < l) \Rightarrow  \\ \quad \quad \quad \quad \quad \operatorname{nrm}(\operatorname{cc\text{-}nag})(\operatorname{opr}(\operatorname{cc\text{-}nag})((\lambda w \in \operatorname{vec}(v).\; \langle\cdot,\cdot\rangle_{v}(w, a))(x), \\ \quad \quad \quad \quad \quad \quad \quad \operatorname{inv}(\operatorname{cc\text{-}nag})((\lambda w \in \operatorname{vec}(v).\; \langle\cdot,\cdot\rangle_{v}(w, a))(y)))) \\ \quad \quad \quad \quad \quad <\; s \end{array}}}
\vnbprop{onb-inner-product-expansion}{\vnbstmt{\begin{array}{@{}l@{}} \forall v, a, e.; \forall x, b \in \operatorname{vec}(v). \\ \quad (\operatorname{is\text{-}hilbert\text{-}space}(v) \wedge \operatorname{is\text{-}orthonormal\text{-}basis}(v, a, e)) \Rightarrow  \\ \quad \quad \operatorname{sums\text{-}to}(\operatorname{cc\text{-}nag}, (\lambda k \in a.\; (\langle\cdot,\cdot\rangle_{v}(x, e(k)) \cdot \langle\cdot,\cdot\rangle_{v}(e(k), b))), \langle\cdot,\cdot\rangle_{v}(x, b)) \end{array}}}
\vnbprop{cc-family-sum-value}{\vnbstmt{\begin{array}{@{}l@{}} \forall f \in (\operatorname{dom}(f) \to \mathbb{C}), r \in \mathbb{C}. \\ \quad \operatorname{sums\text{-}to}(\operatorname{cc\text{-}nag}, f, r) \Rightarrow (\operatorname{cc\text{-}family\text{-}sum}(f) = r) \end{array}}}
\vnbprop{onb-coefficients-in-l2}{\vnbstmt{\begin{array}{@{}l@{}} \forall v, a, e.; \forall x \in \operatorname{vec}(v). \\ \quad (\operatorname{is\text{-}hilbert\text{-}space}(v) \wedge \operatorname{is\text{-}orthonormal\text{-}basis}(v, a, e)) \Rightarrow  \\ \quad \quad (\lambda k \in a.\; \langle\cdot,\cdot\rangle_{v}(x, e(k))) \in \operatorname{l2\text{-}families}(a) \end{array}}}
\vnbprop{onb-coefficients-ip}{\vnbstmt{\begin{array}{@{}l@{}} \forall v, a, e.; \forall x, b \in \operatorname{vec}(v). \\ \quad (\operatorname{is\text{-}hilbert\text{-}space}(v) \wedge \operatorname{is\text{-}orthonormal\text{-}basis}(v, a, e)) \Rightarrow  \\ \quad \quad \operatorname{l2\text{-}ip}(a)((\lambda k \in a.\; \langle\cdot,\cdot\rangle_{v}(x, e(k))), (\lambda k \in a.\; \langle\cdot,\cdot\rangle_{v}(b, e(k)))) = \langle\cdot,\cdot\rangle_{v}(x, b) \end{array}}}

\vnbbookfile{theorem-library/complete-transfer}{theorem-library/complete-transfer}
\noindent{\small complete-transfer.scm -- wave of 2026-11-09/10 (agent BA-ONB3), wired 2026-11-10 from scratchpad/baonb3-1109/b3-complete.scm. Library: analysis.}\par
\vnbprop{complete-of-same-metric}{\vnbstmt{\begin{array}{@{}l@{}} \forall s, t. \\ \quad \operatorname{is\text{-}metric\text{-}space}(t) \wedge \\ \quad (\operatorname{pts}(t) \simeq \operatorname{pts}(s)) \wedge \\ \quad \forall x, y \in \operatorname{pts}(s).\; (\operatorname{dist}(t)(x, y) = \operatorname{dist}(s)(x, y)) \wedge \\ \quad \operatorname{is\text{-}complete}(s) \Rightarrow \\ \quad \quad \operatorname{is\text{-}complete}(t) \end{array}}}
\vnbprop{cc-nag-complete}{\vnbstmt{\operatorname{is\text{-}complete}(\operatorname{nag\text{-}metric\text{-}space}(\operatorname{cc\text{-}nag}))}}

\vnbbookfile{theorem-library/tvs-sub-bound}{theorem-library/tvs-sub-bound}
\noindent{\small tvs-sub-bound.scm -- continuity from seminorm bounds for a map that is only ADDITIVE in the sense tt(a - b) = tt(a) - tt(b), not linear.  The copies of sn-bound-combine / sn-cont-of-bound (tvs-continuity.scm) and sn-cont-of-single-bound (tvs-single-bound.scm) with the hypothesis IS-LINEAR-MAP(m1, m2, tt) replaced by ...}\par
\vnbprop{sn-bound-combine-sub}{\vnbstmt{\begin{array}{@{}l@{}} \forall d, a, x_{1}, f_{1}, b, x_{2}, f_{2}, t.; \forall g \in \mathbf{Set}. \\ \quad \operatorname{is\text{-}seminorm\text{-}family\text{-}on}(d, a, x_{1}, f_{1}) \wedge \\ \quad \operatorname{is\text{-}seminorm\text{-}family\text{-}on}(d, b, x_{2}, f_{2}) \wedge \\ \quad t \in (\operatorname{vec}(a) \to \operatorname{vec}(b)) \wedge \\ \quad \forall \mathit{sba}, \mathit{sbb} \in \operatorname{vec}(a). \\ \quad \quad t(\operatorname{vadd}(a)(\mathit{sba}, \operatorname{vneg}(a)(\mathit{sbb}))) = \operatorname{vadd}(b)(t(\mathit{sba}), \operatorname{vneg}(b)(t(\mathit{sbb}))) \wedge \\ \quad \forall j \in x_{2}. \exists f \in \mathbf{Set}. \\ \quad \quad |f| \in \mathbb{N} \wedge \\ \quad \quad f \subseteq x_{1} \wedge \\ \quad \quad \exists c \in \mathbb{R}. \\ \quad \quad \quad \operatorname{pos\text{-}rr}(c) \wedge \\ \quad \quad \quad \forall x \in \operatorname{vec}(a), \mathit{sbt} \in \mathbb{R}. \\ \quad \quad \quad \quad ((0 \leq \mathit{sbt}) \wedge \forall i \in f.\; (\operatorname{f1}(i)(x) \leq \mathit{sbt})) \Rightarrow  \\ \quad \quad \quad \quad \quad \operatorname{f2}(j)(t(x)) \leq (c \cdot \mathit{sbt}) \wedge \\ \quad |g| \in \mathbb{N} \wedge \\ \quad g \subseteq x_{2} \Rightarrow \\ \quad \quad \exists f \in \mathbf{Set}. \\ \quad \quad \quad |f| \in \mathbb{N} \wedge \\ \quad \quad \quad f \subseteq x_{1} \wedge \\ \quad \quad \quad \exists c \in \mathbb{R}. \\ \quad \quad \quad \quad \operatorname{pos\text{-}rr}(c) \wedge \\ \quad \quad \quad \quad \forall j \in g, x \in \operatorname{vec}(a), \mathit{sbt} \in \mathbb{R}. \\ \quad \quad \quad \quad \quad ((0 \leq \mathit{sbt}) \wedge \forall i \in f.\; (\operatorname{f1}(i)(x) \leq \mathit{sbt})) \Rightarrow  \\ \quad \quad \quad \quad \quad \quad \operatorname{f2}(j)(t(x)) \leq (c \cdot \mathit{sbt}) \end{array}}}
\vnbprop{sn-cont-of-bound-sub}{\vnbstmt{\begin{array}{@{}l@{}} \forall d, a, x_{1}, f_{1}, b, x_{2}, f_{2}, t. \\ \quad \operatorname{is\text{-}seminorm\text{-}family\text{-}on}(d, a, x_{1}, f_{1}) \wedge \\ \quad \operatorname{is\text{-}seminorm\text{-}family\text{-}on}(d, b, x_{2}, f_{2}) \wedge \\ \quad t \in (\operatorname{vec}(a) \to \operatorname{vec}(b)) \wedge \\ \quad \forall \mathit{sba}, \mathit{sbb} \in \operatorname{vec}(a). \\ \quad \quad t(\operatorname{vadd}(a)(\mathit{sba}, \operatorname{vneg}(a)(\mathit{sbb}))) = \operatorname{vadd}(b)(t(\mathit{sba}), \operatorname{vneg}(b)(t(\mathit{sbb}))) \wedge \\ \quad \forall j \in x_{2}. \exists f \in \mathbf{Set}. \\ \quad \quad |f| \in \mathbb{N} \wedge \\ \quad \quad f \subseteq x_{1} \wedge \\ \quad \quad \exists c \in \mathbb{R}. \\ \quad \quad \quad \operatorname{pos\text{-}rr}(c) \wedge \\ \quad \quad \quad \forall x \in \operatorname{vec}(a), \mathit{sbt} \in \mathbb{R}. \\ \quad \quad \quad \quad ((0 \leq \mathit{sbt}) \wedge \forall i \in f.\; (\operatorname{f1}(i)(x) \leq \mathit{sbt})) \Rightarrow  \\ \quad \quad \quad \quad \quad \operatorname{f2}(j)(t(x)) \leq (c \cdot \mathit{sbt}) \Rightarrow \\ \quad \quad \operatorname{is\text{-}hom\text{-}top\text{-}space}(\operatorname{seminorm\text{-}top}(a, x_{1}, f_{1}), \operatorname{seminorm\text{-}top}(b, x_{2}, f_{2}), t) \end{array}}}
\vnbprop{sn-cont-of-single-bound-sub}{\vnbstmt{\begin{array}{@{}l@{}} \forall d, a, x_{1}, f_{1}, b, x_{2}, f_{2}, t. \\ \quad \operatorname{is\text{-}seminorm\text{-}family\text{-}on}(d, a, x_{1}, f_{1}) \wedge \\ \quad \operatorname{is\text{-}seminorm\text{-}family\text{-}on}(d, b, x_{2}, f_{2}) \wedge \\ \quad t \in (\operatorname{vec}(a) \to \operatorname{vec}(b)) \wedge \\ \quad \forall \mathit{sba}, \mathit{sbb} \in \operatorname{vec}(a). \\ \quad \quad t(\operatorname{vadd}(a)(\mathit{sba}, \operatorname{vneg}(a)(\mathit{sbb}))) = \operatorname{vadd}(b)(t(\mathit{sba}), \operatorname{vneg}(b)(t(\mathit{sbb}))) \wedge \\ \quad \forall j \in x_{2}. \exists i \in x_{1}, c \in \mathbb{R}. \\ \quad \quad (0 \leq c) \wedge \forall x \in \operatorname{vec}(a).\; (\operatorname{f2}(j)(t(x)) \leq (c \cdot \operatorname{f1}(i)(x))) \Rightarrow \\ \quad \quad \operatorname{is\text{-}hom\text{-}top\text{-}space}(\operatorname{seminorm\text{-}top}(a, x_{1}, f_{1}), \operatorname{seminorm\text{-}top}(b, x_{2}, f_{2}), t) \end{array}}}
\vnbprop{sn-bound-of-cont-cc}{\vnbstmt{\begin{array}{@{}l@{}} \forall a, x_{1}, f_{1}, b, x_{2}, f_{2}, t.; \forall j \in x_{2}. \\ \quad \operatorname{is\text{-}seminorm\text{-}family\text{-}on}(\operatorname{cc\text{-}normed\text{-}field}, a, x_{1}, f_{1}) \wedge \\ \quad \operatorname{is\text{-}seminorm\text{-}family\text{-}on}(\operatorname{cc\text{-}normed\text{-}field}, b, x_{2}, f_{2}) \wedge \\ \quad \operatorname{is\text{-}linear\text{-}map}(a, b, t) \wedge \\ \quad \operatorname{is\text{-}hom\text{-}top\text{-}space}(\operatorname{seminorm\text{-}top}(a, x_{1}, f_{1}), \operatorname{seminorm\text{-}top}(b, x_{2}, f_{2}), t) \Rightarrow \\ \quad \quad \exists f \in \mathbf{Set}. \\ \quad \quad \quad |f| \in \mathbb{N} \wedge \\ \quad \quad \quad f \subseteq x_{1} \wedge \\ \quad \quad \quad \exists c \in \mathbb{R}. \\ \quad \quad \quad \quad \operatorname{pos\text{-}rr}(c) \wedge \\ \quad \quad \quad \quad \forall x \in \operatorname{vec}(a), \mathit{sbt} \in \mathbb{R}. \\ \quad \quad \quad \quad \quad ((0 \leq \mathit{sbt}) \wedge \forall i \in f.\; (\operatorname{f1}(i)(x) \leq \mathit{sbt})) \Rightarrow  \\ \quad \quad \quad \quad \quad \quad \operatorname{f2}(j)(t(x)) \leq (c \cdot \mathit{sbt}) \end{array}}}
\vnbprop{sn-cont-iff-bound-cc}{\vnbstmt{\begin{array}{@{}l@{}} \forall a, x_{1}, f_{1}, b, x_{2}, f_{2}, t. \\ \quad \operatorname{is\text{-}seminorm\text{-}family\text{-}on}(\operatorname{cc\text{-}normed\text{-}field}, a, x_{1}, f_{1}) \wedge \\ \quad \operatorname{is\text{-}seminorm\text{-}family\text{-}on}(\operatorname{cc\text{-}normed\text{-}field}, b, x_{2}, f_{2}) \wedge \\ \quad \operatorname{is\text{-}linear\text{-}map}(a, b, t) \Rightarrow \\ \quad \quad \operatorname{is\text{-}hom\text{-}top\text{-}space}(\operatorname{seminorm\text{-}top}(a, x_{1}, f_{1}), \operatorname{seminorm\text{-}top}(b, x_{2}, f_{2}), t) \\ \quad \quad \Leftrightarrow\; \forall j \in x_{2}. \exists f \in \mathbf{Set}. \\ \quad \quad \quad |f| \in \mathbb{N} \wedge \\ \quad \quad \quad f \subseteq x_{1} \wedge \\ \quad \quad \quad \exists c \in \mathbb{R}. \\ \quad \quad \quad \quad \operatorname{pos\text{-}rr}(c) \wedge \\ \quad \quad \quad \quad \forall x \in \operatorname{vec}(a), \mathit{sbt} \in \mathbb{R}. \\ \quad \quad \quad \quad \quad ((0 \leq \mathit{sbt}) \wedge \forall i \in f.\; (\operatorname{f1}(i)(x) \leq \mathit{sbt})) \Rightarrow  \\ \quad \quad \quad \quad \quad \quad \operatorname{f2}(j)(t(x)) \leq (c \cdot \mathit{sbt}) \end{array}}}

\part{Complex Analysis}
\noindent Holomorphic functions and power series over CC; the exponential and the logarithm; line integrals, Goursat and the Cauchy integral formula; winding numbers; Cauchy's estimates, Liouville and the fundamental theorem of algebra; Morera; the Gamma function.
\par\noindent\textit{Builds on:} Analysis.
\par\noindent{\small 173 files.}

\vnbbookfile{structure-library/zero-order-vocab}{structure-library/zero-order-vocab}
\noindent{\small zero-order-vocab.scm -- WAVE 2 (2026-10-27..30), wired 2026-10-30 from scratchpad/ch6/zero-order-vocab.scm. Library: complex-analysis.}\par
\par\medskip\noindent\textbf{Definition} (\texttt{has-zero-of-order}).\;For all $\operatorname{zou\_}$, $\operatorname{zof\_}$, $\operatorname{zoa\_}$ and $\operatorname{zom\_}$, 
\begin{equation*}
\begin{array}{@{}l@{}} \operatorname{has\text{-}zero\text{-}of\text{-}order}(\operatorname{zou\_}, \operatorname{zof\_}, \operatorname{zoa\_}, \operatorname{zom\_}) \\ \Leftrightarrow\; \operatorname{holomorphic\text{-}on}(\operatorname{zou\_}, \operatorname{zof\_}) \wedge \\ \operatorname{zoa\_} \in \operatorname{zou\_} \wedge \\ \operatorname{zom\_} \in \mathbb{N} \wedge \\ \exists \operatorname{zoh\_}. \\ \quad \operatorname{holomorphic\text{-}on}(\operatorname{zou\_}, \operatorname{zoh\_}) \wedge \\ \quad \neg (\operatorname{zoh\_}(\operatorname{zoa\_}) = 0) \wedge \\ \quad \forall \operatorname{zox\_} \in \operatorname{zou\_}.\; (\operatorname{zof\_}(\operatorname{zox\_}) = \cdot(\operatorname{power}(-(\operatorname{zox\_}, \operatorname{zoa\_}), \operatorname{zom\_}), \operatorname{zoh\_}(\operatorname{zox\_}))) \end{array}
\end{equation*}

\par\medskip\noindent\textbf{Definition} (\texttt{zero-order}).\;$\operatorname{zero\text{-}order}(\operatorname{zou\_}, \operatorname{zof\_}, \operatorname{zoa\_}) :=$ $(\iota\, \operatorname{zon\_}.\;\operatorname{has\text{-}zero\text{-}of\text{-}order}(\operatorname{zou\_}, \operatorname{zof\_}, \operatorname{zoa\_}, \operatorname{zon\_}))$
\par\medskip\noindent\textbf{Definition} (\texttt{count-hits}).\;For all $\operatorname{chs\_}$ and $\operatorname{chz\_}$,  $(\operatorname{count\text{-}hits}(\operatorname{chs\_}, \operatorname{chz\_}, 0) \simeq 0)$\par
Let $\operatorname{chs\_}$ and $\operatorname{chz\_}$ be arbitrary. Let $\operatorname{chn\_} \in \mathbb{N}$. Then
\begin{equation*}
(\operatorname{count\text{-}hits}(\operatorname{chs\_}, \operatorname{chz\_}, \operatorname{succ}(\operatorname{chn\_})) \simeq \left(\text{if }(\operatorname{chs\_}(\operatorname{chn\_}) = \operatorname{chz\_})\text{ then }\operatorname{succ}(\operatorname{count\text{-}hits}(\operatorname{chs\_}, \operatorname{chz\_}, \operatorname{chn\_}))\text{ else }\operatorname{count\text{-}hits}(\operatorname{chs\_}, \operatorname{chz\_}, \operatorname{chn\_})\right))
\end{equation*}


\vnbbookfile{theorem-library/cc-series}{theorem-library/cc-series}
\noindent{\small cc-series.scm -- THE COMPLEX SERIES LAYER: ordered partial sums in CC, and their convergence in CC-MS.}\par
\par\medskip\noindent\textbf{Definition} (\texttt{cc-series-partial-sum}).\;$\operatorname{cc\text{-}series\text{-}partial\text{-}sum}(f, k) :=$ $\operatorname{sum\text{-}ag}(\operatorname{normed\text{-}field\text{-}additive\text{-}ag}(\operatorname{cc\text{-}normed\text{-}field}), f, k)$
\par\medskip\noindent\textbf{Definition} (\texttt{cc-series-converges-to}).\;For all $f$ and $l$, 
\begin{equation*}
\begin{array}{@{}l@{}} \operatorname{cc\text{-}series\text{-}converges\text{-}to}(f, l) \\ \Leftrightarrow\; \operatorname{converges\text{-}to}(\operatorname{cc\text{-}ms}, (\lambda k \in \mathbb{N}.\; \operatorname{cc\text{-}series\text{-}partial\text{-}sum}(f, k)), l) \end{array}
\end{equation*}

\par\medskip\noindent\textbf{Definition} (\texttt{cc-series-converges}).\;For every $f$, 
\begin{equation*}
\operatorname{cc\text{-}series\text{-}converges}(f) \Leftrightarrow \operatorname{converges}(\operatorname{cc\text{-}ms}, (\lambda k \in \mathbb{N}.\; \operatorname{cc\text{-}series\text{-}partial\text{-}sum}(f, k)))
\end{equation*}

\vnbprop{cc-additive-ag-opr}{\vnbstmt{(\operatorname{opr}(\operatorname{normed\text{-}field\text{-}additive\text{-}ag}(\operatorname{cc\text{-}normed\text{-}field})) \simeq (\lambda {[}\operatorname{x\_}, \operatorname{y\_}] \in (\mathbb{C} \times \mathbb{C}).\; (\operatorname{x\_} + \operatorname{y\_})))}}
\vnbprop{cc-additive-ag-iden}{\vnbstmt{(\operatorname{iden}(\operatorname{normed\text{-}field\text{-}additive\text{-}ag}(\operatorname{cc\text{-}normed\text{-}field})) \simeq 0)}}
\vnbprop{cc-series-partial-sum-zero}{\vnbstmt{\forall f.\; (\operatorname{cc\text{-}series\text{-}partial\text{-}sum}(f, 0) \simeq 0)}}
\vnbprop{sum-ag-cc-in-cc}{\vnbstmt{\begin{array}{@{}l@{}} \forall k \in \mathbb{N}, f \in (\mathbb{N} \to \mathbb{C}). \\ \quad \operatorname{sum\text{-}ag}(\operatorname{normed\text{-}field\text{-}additive\text{-}ag}(\operatorname{cc\text{-}normed\text{-}field}), f, k) \in \mathbb{C} \end{array}}}
\vnbprop{cc-series-partial-sum-in-cc}{\vnbstmt{\forall k \in \mathbb{N}, f \in (\mathbb{N} \to \mathbb{C}).\; (\operatorname{cc\text{-}series\text{-}partial\text{-}sum}(f, k) \in \mathbb{C})}}
\vnbprop{cc-series-partial-sum-unfold}{\vnbstmt{\begin{array}{@{}l@{}} \forall f, k. \\ \quad (\operatorname{cc\text{-}series\text{-}partial\text{-}sum}(f, k) \simeq \operatorname{sum\text{-}ag}(\operatorname{normed\text{-}field\text{-}additive\text{-}ag}(\operatorname{cc\text{-}normed\text{-}field}), f, k)) \end{array}}}
\vnbprop{cc-series-partial-sum-succ}{\vnbstmt{\begin{array}{@{}l@{}} \forall f.; \forall k \in \mathbb{N}. \\ \quad ((\operatorname{cc\text{-}series\text{-}partial\text{-}sum}(f, k) \in \mathbb{C}) \wedge (f(k) \in \mathbb{C})) \Rightarrow  \\ \quad \quad (\operatorname{cc\text{-}series\text{-}partial\text{-}sum}(f, (k + 1)) \simeq (\operatorname{cc\text{-}series\text{-}partial\text{-}sum}(f, k) + f(k))) \end{array}}}
\vnbprop{cc-series-partial-sum-magnitude-le}{\vnbstmt{\begin{array}{@{}l@{}} \forall k \in \mathbb{N}, f \in (\mathbb{N} \to \mathbb{C}), g \in (\mathbb{N} \to \mathbb{R}). \\ \quad \forall j \in \mathbb{N}.\; (g(j) = \operatorname{magnitude}(f(j))) \Rightarrow  \\ \quad \quad \operatorname{magnitude}(\operatorname{cc\text{-}series\text{-}partial\text{-}sum}(f, k)) \leq \operatorname{series\text{-}partial\text{-}sum}(g, k) \end{array}}}

\vnbbookfile{structure-library/poly-fn-vocab}{structure-library/poly-fn-vocab}
\noindent{\small poly-fn-vocab.scm -- POLYNOMIAL FUNCTIONS (CA-MONTEL "POLY", 2026-10-26).  Destination: structure-library/.}\par
\par\medskip\noindent\textbf{Definition} (\texttt{is-cc-subring}).\;For every $\mathit{kk}$, 
\begin{equation*}
\begin{array}{@{}l@{}} \operatorname{is\text{-}cc\text{-}subring}(\mathit{kk}) \\ \Leftrightarrow\; \mathit{kk} \subseteq \mathbb{C} \wedge \\ 0 \in \mathit{kk} \wedge \\ 1 \in \mathit{kk} \wedge \\ \forall \operatorname{psu\_}, \operatorname{psv\_} \in \mathit{kk}.\; (+(\operatorname{psu\_}, \operatorname{psv\_}) \in \mathit{kk}) \wedge \\ \forall \operatorname{psu\_}, \operatorname{psv\_} \in \mathit{kk}.\; (\cdot(\operatorname{psu\_}, \operatorname{psv\_}) \in \mathit{kk}) \end{array}
\end{equation*}

\par\medskip\noindent\textbf{Definition} (\texttt{is-poly-fn-on}).\;For all $\mathit{kk}$, $a$ and $f$, 
\begin{equation*}
\begin{array}{@{}l@{}} \operatorname{is\text{-}poly\text{-}fn\text{-}on}(\mathit{kk}, a, f) \\ \Leftrightarrow\; a \subseteq \mathit{kk} \wedge \\ f \in (a \to \mathit{kk}) \wedge \\ \exists \operatorname{pfn\_} \in \mathbb{N}, \operatorname{pfc\_} \in (\mathbb{N} \to \mathit{kk}). \\ \quad \forall \operatorname{pfk\_} \in \mathbb{N}.\; ((\operatorname{pfn\_} \leq \operatorname{pfk\_}) \Rightarrow (\operatorname{pfc\_}(\operatorname{pfk\_}) = 0)) \wedge \\ \quad \forall \operatorname{pfz\_} \in a. \\ \quad \quad f(\operatorname{pfz\_}) = \operatorname{cc\text{-}series\text{-}partial\text{-}sum}((\lambda \operatorname{pfj\_} \in \mathbb{N}.\; \cdot(\operatorname{pfc\_}(\operatorname{pfj\_}), \operatorname{power}(\operatorname{pfz\_}, \operatorname{pfj\_}))), \operatorname{pfn\_}) \end{array}
\end{equation*}

\par\medskip\noindent\textbf{Definition} (\texttt{cf-cons}).\;$\operatorname{cf\text{-}cons}(b, c) :=$
\begin{equation*}
(\lambda \operatorname{pfk\_} \in \mathbb{N}.\; ((b \cdot \operatorname{power}(0, \operatorname{pfk\_})) + ((1 - \operatorname{power}(0, \operatorname{pfk\_})) \cdot c(\operatorname{nn\text{-}minus}(\operatorname{pfk\_}, 1)))))
\end{equation*}

\par\medskip\noindent\textbf{Definition} (\texttt{is-poly-deg}).\;For all $\mathit{kk}$, $a$, $f$ and $n$, 
\begin{equation*}
\begin{array}{@{}l@{}} \operatorname{is\text{-}poly\text{-}deg}(\mathit{kk}, a, f, n) \\ \Leftrightarrow\; a \subseteq \mathit{kk} \wedge \\ f \in (a \to \mathit{kk}) \wedge \\ \exists \operatorname{pfc\_} \in (\mathbb{N} \to \mathit{kk}). \\ \quad \forall \operatorname{pfk\_} \in \mathbb{N}.\; ((\operatorname{succ}(n) \leq \operatorname{pfk\_}) \Rightarrow (\operatorname{pfc\_}(\operatorname{pfk\_}) = 0)) \wedge \\ \quad \neg (\operatorname{pfc\_}(n) = 0) \wedge \\ \quad \forall \operatorname{pfz\_} \in a. \\ \quad \quad f(\operatorname{pfz\_}) = \operatorname{cc\text{-}series\text{-}partial\text{-}sum}((\lambda \operatorname{pfj\_} \in \mathbb{N}.\; \cdot(\operatorname{pfc\_}(\operatorname{pfj\_}), \operatorname{power}(\operatorname{pfz\_}, \operatorname{pfj\_}))), \operatorname{succ}(n)) \end{array}
\end{equation*}

\par\medskip\noindent\textbf{Definition} (\texttt{poly-fn-deg}).\;$\operatorname{poly\text{-}fn\text{-}deg}(\mathit{kk}, a, f) :=$ $(\iota\, \operatorname{pdgn\_}.\;((\operatorname{pdgn\_} \in \mathbb{N}) \wedge \operatorname{is\text{-}poly\text{-}deg}(\mathit{kk}, a, f, \operatorname{pdgn\_})))$

\vnbbookfile{theorem-library/poly-fn-laws}{theorem-library/poly-fn-laws}
\noindent{\small poly-fn-laws.scm -- laws of polynomial functions (CA-MONTEL "POLY", 2026-10-26).  Vocabulary: poly-fn-vocab.scm.  Library: analysis.  Helper prefix pfl-.}\par
\vnbprop{cf-cons-unfold}{\vnbstmt{\begin{array}{@{}l@{}} \forall b, c. \\ \quad (\operatorname{cf\text{-}cons}(b, c) \simeq (\lambda k \in \mathbb{N}.\; ((b \cdot \operatorname{power}(0, k)) + ((1 - \operatorname{power}(0, k)) \cdot c(\operatorname{nn\text{-}minus}(k, 1)))))) \end{array}}}
\vnbprop{cf-cons-zero}{\vnbstmt{\forall b \in \mathbb{C}, c \in (\mathbb{N} \to \mathbb{C}).\; (\operatorname{cf\text{-}cons}(b, c)(0) = b)}}
\vnbprop{cf-cons-succ}{\vnbstmt{\forall b \in \mathbb{C}, c \in (\mathbb{N} \to \mathbb{C}), j \in \mathbb{N}.\; (\operatorname{cf\text{-}cons}(b, c)((j + 1)) = c(j))}}
\vnbprop{cf-cons-in-fun}{\vnbstmt{\begin{array}{@{}l@{}} \forall k.; \forall b \in k, c \in (\mathbb{N} \to k). \\ \quad \operatorname{is\text{-}cc\text{-}subring}(k) \Rightarrow (\operatorname{cf\text{-}cons}(b, c) \in (\mathbb{N} \to k)) \end{array}}}
\vnbprop{cc-ps-cons}{\vnbstmt{\begin{array}{@{}l@{}} \forall n \in \mathbb{N}, b \in \mathbb{C}, c \in (\mathbb{N} \to \mathbb{C}), z \in \mathbb{C}. \\ \quad \operatorname{cc\text{-}series\text{-}partial\text{-}sum}((\lambda j \in \mathbb{N}.\; (\operatorname{cf\text{-}cons}(b, c)(j) \cdot \operatorname{power}(z, j))), (n + 1)) \\ \quad =\; b + (z \cdot \operatorname{cc\text{-}series\text{-}partial\text{-}sum}((\lambda j \in \mathbb{N}.\; (c(j) \cdot \operatorname{power}(z, j))), n)) \end{array}}}
\vnbprop{cc-ps-stable}{\vnbstmt{\begin{array}{@{}l@{}} \forall m, n \in \mathbb{N}, c \in (\mathbb{N} \to \mathbb{C}), z \in \mathbb{C}. \\ \quad \forall k \in \mathbb{N}.\; ((n \leq k) \Rightarrow (c(k) = 0)) \Rightarrow  \\ \quad \quad \operatorname{cc\text{-}series\text{-}partial\text{-}sum}((\lambda j \in \mathbb{N}.\; (c(j) \cdot \operatorname{power}(z, j))), (n + m)) \\ \quad \quad =\; \operatorname{cc\text{-}series\text{-}partial\text{-}sum}((\lambda j \in \mathbb{N}.\; (c(j) \cdot \operatorname{power}(z, j))), n) \end{array}}}
\vnbprop{cc-subring-power}{\vnbstmt{\forall i \in \mathbb{N}, k.; \forall z \in k.\; (\operatorname{is\text{-}cc\text{-}subring}(k) \Rightarrow (\operatorname{power}(z, i) \in k))}}
\vnbprop{cc-ps-in-subring}{\vnbstmt{\begin{array}{@{}l@{}} \forall n \in \mathbb{N}, k.; \forall c \in (\mathbb{N} \to k), z \in k. \\ \quad \operatorname{is\text{-}cc\text{-}subring}(k) \Rightarrow  \\ \quad \quad \operatorname{cc\text{-}series\text{-}partial\text{-}sum}((\lambda j \in \mathbb{N}.\; (c(j) \cdot \operatorname{power}(z, j))), n) \in k \end{array}}}
\vnbprop{cc-ps-lincomb}{\vnbstmt{\begin{array}{@{}l@{}} \forall n \in \mathbb{N}, c, d, e \in (\mathbb{N} \to \mathbb{C}), l, \mathit{pfbe}, z \in \mathbb{C}. \\ \quad \forall k \in \mathbb{N}.\; (e(k) = ((l \cdot c(k)) + (\mathit{pfbe} \cdot d(k)))) \Rightarrow  \\ \quad \quad \operatorname{cc\text{-}series\text{-}partial\text{-}sum}((\lambda j \in \mathbb{N}.\; (e(j) \cdot \operatorname{power}(z, j))), n) \\ \quad \quad =\; (l \cdot \operatorname{cc\text{-}series\text{-}partial\text{-}sum}((\lambda j \in \mathbb{N}.\; (c(j) \cdot \operatorname{power}(z, j))), n)) + (\mathit{pfbe} \cdot \operatorname{cc\text{-}series\text{-}partial\text{-}sum}((\lambda j \in \mathbb{N}.\; (d(j) \cdot \operatorname{power}(z, j))), n)) \end{array}}}
\vnbprop{poly-fn-ext}{\vnbstmt{\begin{array}{@{}l@{}} \forall k, a, f.; \forall g \in (a \to k). \\ \quad (\operatorname{is\text{-}poly\text{-}fn\text{-}on}(k, a, f) \wedge \forall z \in a.\; (g(z) = f(z))) \Rightarrow  \\ \quad \quad \operatorname{is\text{-}poly\text{-}fn\text{-}on}(k, a, g) \end{array}}}
\vnbprop{poly-fn-const}{\vnbstmt{\begin{array}{@{}l@{}} \forall k, a.; \forall b \in k. \\ \quad (\operatorname{is\text{-}cc\text{-}subring}(k) \wedge (a \subseteq k)) \Rightarrow \operatorname{is\text{-}poly\text{-}fn\text{-}on}(k, a, (\lambda x \in a.\; b)) \end{array}}}
\vnbprop{poly-fn-mulz}{\vnbstmt{\begin{array}{@{}l@{}} \forall k, a, f. \\ \quad (\operatorname{is\text{-}cc\text{-}subring}(k) \wedge \operatorname{is\text{-}poly\text{-}fn\text{-}on}(k, a, f)) \Rightarrow  \\ \quad \quad \operatorname{is\text{-}poly\text{-}fn\text{-}on}(k, a, (\lambda x \in a.\; (x \cdot f(x)))) \end{array}}}
\vnbprop{poly-fn-lin}{\vnbstmt{\begin{array}{@{}l@{}} \forall k, a, f, g.; \forall l, e \in k. \\ \quad (\operatorname{is\text{-}cc\text{-}subring}(k) \wedge (\operatorname{is\text{-}poly\text{-}fn\text{-}on}(k, a, f) \wedge \operatorname{is\text{-}poly\text{-}fn\text{-}on}(k, a, g))) \Rightarrow  \\ \quad \quad \operatorname{is\text{-}poly\text{-}fn\text{-}on}(k, a, (\lambda x \in a.\; ((l \cdot f(x)) + (e \cdot g(x))))) \end{array}}}
\vnbprop{poly-fn-add}{\vnbstmt{\begin{array}{@{}l@{}} \forall k, a, f, g. \\ \quad (\operatorname{is\text{-}cc\text{-}subring}(k) \wedge (\operatorname{is\text{-}poly\text{-}fn\text{-}on}(k, a, f) \wedge \operatorname{is\text{-}poly\text{-}fn\text{-}on}(k, a, g))) \Rightarrow  \\ \quad \quad \operatorname{is\text{-}poly\text{-}fn\text{-}on}(k, a, (\lambda x \in a.\; (f(x) + g(x)))) \end{array}}}
\vnbprop{poly-fn-scale}{\vnbstmt{\begin{array}{@{}l@{}} \forall k, a, f.; \forall l \in k. \\ \quad (\operatorname{is\text{-}cc\text{-}subring}(k) \wedge \operatorname{is\text{-}poly\text{-}fn\text{-}on}(k, a, f)) \Rightarrow  \\ \quad \quad \operatorname{is\text{-}poly\text{-}fn\text{-}on}(k, a, (\lambda x \in a.\; (l \cdot f(x)))) \end{array}}}
\vnbprop{poly-fn-id}{\vnbstmt{\begin{array}{@{}l@{}} \forall k, a. \\ \quad (\operatorname{is\text{-}cc\text{-}subring}(k) \wedge (a \subseteq k)) \Rightarrow \operatorname{is\text{-}poly\text{-}fn\text{-}on}(k, a, (\lambda x \in a.\; x)) \end{array}}}
\vnbprop{cc-ps-congr}{\vnbstmt{\begin{array}{@{}l@{}} \forall n \in \mathbb{N}, c, d \in (\mathbb{N} \to \mathbb{C}), z \in \mathbb{C}. \\ \quad \forall k \in \mathbb{N}.\; (c(k) = d(k)) \Rightarrow  \\ \quad \quad \operatorname{cc\text{-}series\text{-}partial\text{-}sum}((\lambda j \in \mathbb{N}.\; (c(j) \cdot \operatorname{power}(z, j))), n) \\ \quad \quad =\; \operatorname{cc\text{-}series\text{-}partial\text{-}sum}((\lambda j \in \mathbb{N}.\; (d(j) \cdot \operatorname{power}(z, j))), n) \end{array}}}
\vnbprop{poly-fn-mul-ps}{\vnbstmt{\begin{array}{@{}l@{}} \forall n \in \mathbb{N}, k, a.; \forall c \in (\mathbb{N} \to k), g. \\ \quad \operatorname{is\text{-}cc\text{-}subring}(k) \wedge \\ \quad a \subseteq k \wedge \\ \quad \forall \mathit{pfk} \in \mathbb{N}.\; ((n \leq \mathit{pfk}) \Rightarrow (c(\mathit{pfk}) = 0)) \wedge \\ \quad \operatorname{is\text{-}poly\text{-}fn\text{-}on}(k, a, g) \Rightarrow \\ \quad \quad \operatorname{is\text{-}poly\text{-}fn\text{-}on}(k, a, \\ \quad \quad \quad (\lambda x \in a.\; (\operatorname{cc\text{-}series\text{-}partial\text{-}sum}((\lambda j \in \mathbb{N}.\; (c(j) \cdot \operatorname{power}(x, j))), n) \cdot g(x)))) \end{array}}}
\vnbprop{poly-fn-mul}{\vnbstmt{\begin{array}{@{}l@{}} \forall k, a, f, g. \\ \quad (\operatorname{is\text{-}cc\text{-}subring}(k) \wedge (\operatorname{is\text{-}poly\text{-}fn\text{-}on}(k, a, f) \wedge \operatorname{is\text{-}poly\text{-}fn\text{-}on}(k, a, g))) \Rightarrow  \\ \quad \quad \operatorname{is\text{-}poly\text{-}fn\text{-}on}(k, a, (\lambda x \in a.\; (f(x) \cdot g(x)))) \end{array}}}
\vnbprop{poly-fn-affine-ps}{\vnbstmt{\begin{array}{@{}l@{}} \forall n \in \mathbb{N}, k, a.; \forall l, e \in k, c \in (\mathbb{N} \to k). \\ \quad (\operatorname{is\text{-}cc\text{-}subring}(k) \wedge ((a \subseteq k) \wedge \forall \mathit{pfk} \in \mathbb{N}.\; ((n \leq \mathit{pfk}) \Rightarrow (c(\mathit{pfk}) = 0)))) \Rightarrow  \\ \quad \quad \operatorname{is\text{-}poly\text{-}fn\text{-}on}(k, a, \\ \quad \quad \quad (\lambda x \in a.\; \operatorname{cc\text{-}series\text{-}partial\text{-}sum}((\lambda j \in \mathbb{N}.\; (c(j) \cdot \operatorname{power}(((l \cdot x) + e), j))), n))) \end{array}}}
\vnbprop{poly-fn-affine}{\vnbstmt{\begin{array}{@{}l@{}} \forall k, b, f, a.; \forall l, e \in k. \\ \quad (\operatorname{is\text{-}cc\text{-}subring}(k) \wedge (\operatorname{is\text{-}poly\text{-}fn\text{-}on}(k, b, f) \wedge ((a \subseteq k) \wedge \forall x \in a.\; (((l \cdot x) + e) \in b)))) \Rightarrow  \\ \quad \quad \operatorname{is\text{-}poly\text{-}fn\text{-}on}(k, a, (\lambda x \in a.\; f(((l \cdot x) + e)))) \end{array}}}
\vnbprop{poly-fn-comp-ps}{\vnbstmt{\begin{array}{@{}l@{}} \forall n \in \mathbb{N}, k, a, g.; \forall c \in (\mathbb{N} \to k). \\ \quad (\operatorname{is\text{-}cc\text{-}subring}(k) \wedge (\operatorname{is\text{-}poly\text{-}fn\text{-}on}(k, a, g) \wedge \forall \mathit{pfk} \in \mathbb{N}.\; ((n \leq \mathit{pfk}) \Rightarrow (c(\mathit{pfk}) = 0)))) \Rightarrow  \\ \quad \quad \operatorname{is\text{-}poly\text{-}fn\text{-}on}(k, a, (\lambda x \in a.\; \operatorname{cc\text{-}series\text{-}partial\text{-}sum}((\lambda j \in \mathbb{N}.\; (c(j) \cdot \operatorname{power}(g(x), j))), n))) \end{array}}}
\vnbprop{poly-fn-comp}{\vnbstmt{\begin{array}{@{}l@{}} \forall k, b, f, a, g. \\ \quad \operatorname{is\text{-}cc\text{-}subring}(k) \wedge \\ \quad \operatorname{is\text{-}poly\text{-}fn\text{-}on}(k, b, f) \wedge \\ \quad \operatorname{is\text{-}poly\text{-}fn\text{-}on}(k, a, g) \wedge \\ \quad \forall x \in a.\; (g(x) \in b) \Rightarrow \\ \quad \quad \operatorname{is\text{-}poly\text{-}fn\text{-}on}(k, a, (\lambda x \in a.\; f(g(x)))) \end{array}}}
\vnbprop{cc-subring-rr}{\vnbstmt{\operatorname{is\text{-}cc\text{-}subring}(\mathbb{R})}}
\vnbprop{cc-subring-cc}{\vnbstmt{\operatorname{is\text{-}cc\text{-}subring}(\mathbb{C})}}
\vnbprop{poly-fn-restrict}{\vnbstmt{\begin{array}{@{}l@{}} \forall k, a, f, \mathit{pfaa}. \\ \quad (\operatorname{is\text{-}cc\text{-}subring}(k) \wedge (\operatorname{is\text{-}poly\text{-}fn\text{-}on}(k, a, f) \wedge (\mathit{pfaa} \subseteq a))) \Rightarrow  \\ \quad \quad \operatorname{is\text{-}poly\text{-}fn\text{-}on}(k, \mathit{pfaa}, (\lambda x \in \mathit{pfaa}.\; f(x))) \end{array}}}

\vnbbookfile{theorem-library/poly-fn-roots}{theorem-library/poly-fn-roots}
\noindent{\small poly-fn-roots.scm -- division by z - a, the root bound, uniqueness of coefficients, degree, composition (CA-MONTEL "POLY-2", 2026-10-26).  Vocabulary: poly-fn-vocab.scm (+ the degree vocabulary below). Library: analysis.  Helper prefix pfr- in drafting; renamed pfr- with the copied helpers at the head.}\par
\vnbprop{cc-ps-divide}{\vnbstmt{\begin{array}{@{}l@{}} \forall n \in \mathbb{N}, c \in (\mathbb{N} \to \mathbb{C}), a \in \mathbb{C}. \\ \quad \forall k \in \mathbb{N}.\; (((n + 1) \leq k) \Rightarrow (c(k) = 0)) \Rightarrow  \\ \quad \quad \exists d \in (\mathbb{N} \to \mathbb{C}). \\ \quad \quad \quad \forall k \in \mathbb{N}.\; ((n \leq k) \Rightarrow (d(k) = 0)) \wedge \\ \quad \quad \quad \forall k \in \mathbb{N}.\; (c((k + 1)) = (d(k) - (a \cdot d((k + 1))))) \wedge \\ \quad \quad \quad c(0) = (\operatorname{cc\text{-}series\text{-}partial\text{-}sum}((\lambda j \in \mathbb{N}.\; (c(j) \cdot \operatorname{power}(a, j))), (n + 1)) - (a \cdot d(0))) \wedge \\ \quad \quad \quad \forall z \in \mathbb{C}. \\ \quad \quad \quad \quad \operatorname{cc\text{-}series\text{-}partial\text{-}sum}((\lambda j \in \mathbb{N}.\; (c(j) \cdot \operatorname{power}(z, j))), (n + 1)) \\ \quad \quad \quad \quad =\; \operatorname{cc\text{-}series\text{-}partial\text{-}sum}((\lambda j \in \mathbb{N}.\; (c(j) \cdot \operatorname{power}(a, j))), (n + 1)) + ((z - a) \cdot \operatorname{cc\text{-}series\text{-}partial\text{-}sum}((\lambda j \in \mathbb{N}.\; (d(j) \cdot \operatorname{power}(z, j))), n)) \end{array}}}
\vnbprop{cc-ps-root-bound}{\vnbstmt{\begin{array}{@{}l@{}} \forall n \in \mathbb{N}, c, s \in (\mathbb{N} \to \mathbb{C}), k \in \mathbb{N}. \\ \quad \forall k \in \mathbb{N}.\; (((n + 1) \leq k) \Rightarrow (c(k) = 0)) \wedge \\ \quad \forall i, j \in \mathbb{N}.\; (((i \leq n) \wedge ((j \leq n) \wedge (s(i) = s(j)))) \Rightarrow (i = j)) \wedge \\ \quad \forall i \in \mathbb{N}. \\ \quad \quad (i \leq n) \Rightarrow  \\ \quad \quad \quad \operatorname{cc\text{-}series\text{-}partial\text{-}sum}((\lambda \mathit{pfj} \in \mathbb{N}.\; (c(\mathit{pfj}) \cdot \operatorname{power}(s(i), \mathit{pfj}))), (n + 1)) = 0 \Rightarrow \\ \quad \quad c(k) = 0 \end{array}}}
\vnbprop{cc-ps-coef-unique}{\vnbstmt{\begin{array}{@{}l@{}} \forall a.; \forall n \in \mathbb{N}, c \in (\mathbb{N} \to \mathbb{C}), m \in \mathbb{N}, d \in (\mathbb{N} \to \mathbb{C}), k \in \mathbb{N}. \\ \quad a \subseteq \mathbb{C} \wedge \\ \quad \exists s \in (\mathbb{N} \to a). \forall i, j \in \mathbb{N}. \\ \quad \quad (s(i) = s(j)) \Rightarrow (i = j) \wedge \\ \quad \forall k \in \mathbb{N}.\; ((n \leq k) \Rightarrow (c(k) = 0)) \wedge \\ \quad \forall k \in \mathbb{N}.\; ((m \leq k) \Rightarrow (d(k) = 0)) \wedge \\ \quad \forall z \in a. \\ \quad \quad \operatorname{cc\text{-}series\text{-}partial\text{-}sum}((\lambda \mathit{pfj} \in \mathbb{N}.\; (c(\mathit{pfj}) \cdot \operatorname{power}(z, \mathit{pfj}))), n) \\ \quad \quad =\; \operatorname{cc\text{-}series\text{-}partial\text{-}sum}((\lambda \mathit{pfj} \in \mathbb{N}.\; (d(\mathit{pfj}) \cdot \operatorname{power}(z, \mathit{pfj}))), m) \Rightarrow \\ \quad \quad c(k) = d(k) \end{array}}}
\vnbprop{cc-ps-deg-mul}{\vnbstmt{\begin{array}{@{}l@{}} \forall n \in \mathbb{N}, c \in (\mathbb{N} \to \mathbb{C}), m \in \mathbb{N}, d \in (\mathbb{N} \to \mathbb{C}). \\ \quad \forall k \in \mathbb{N}.\; (((n + 1) \leq k) \Rightarrow (c(k) = 0)) \wedge \\ \quad \neg (c(n) = 0) \wedge \\ \quad \forall k \in \mathbb{N}.\; (((m + 1) \leq k) \Rightarrow (d(k) = 0)) \wedge \\ \quad \neg (d(m) = 0) \Rightarrow \\ \quad \quad \exists e \in (\mathbb{N} \to \mathbb{C}). \\ \quad \quad \quad \forall k \in \mathbb{N}.\; ((((n + m) + 1) \leq k) \Rightarrow (e(k) = 0)) \wedge \\ \quad \quad \quad \neg (e((n + m)) = 0) \wedge \\ \quad \quad \quad \forall z \in \mathbb{C}. \\ \quad \quad \quad \quad \operatorname{cc\text{-}series\text{-}partial\text{-}sum}((\lambda j \in \mathbb{N}.\; (c(j) \cdot \operatorname{power}(z, j))), (n + 1)) \cdot \operatorname{cc\text{-}series\text{-}partial\text{-}sum}((\lambda j \in \mathbb{N}.\; (d(j) \cdot \operatorname{power}(z, j))), (m + 1)) \\ \quad \quad \quad \quad =\; \operatorname{cc\text{-}series\text{-}partial\text{-}sum}((\lambda j \in \mathbb{N}.\; (e(j) \cdot \operatorname{power}(z, j))), ((n + m) + 1)) \end{array}}}
\vnbprop{poly-deg-is-poly-fn}{\vnbstmt{\begin{array}{@{}l@{}} \forall k, a, f.; \forall n \in \mathbb{N}. \\ \quad \operatorname{is\text{-}poly\text{-}deg}(k, a, f, n) \Rightarrow \operatorname{is\text{-}poly\text{-}fn\text{-}on}(k, a, f) \end{array}}}
\vnbprop{poly-deg-mul}{\vnbstmt{\begin{array}{@{}l@{}} \forall a, f, g.; \forall n, m \in \mathbb{N}. \\ \quad (\operatorname{is\text{-}poly\text{-}deg}(\mathbb{C}, a, f, n) \wedge \operatorname{is\text{-}poly\text{-}deg}(\mathbb{C}, a, g, m)) \Rightarrow  \\ \quad \quad \operatorname{is\text{-}poly\text{-}deg}(\mathbb{C}, a, (\lambda x \in a.\; (f(x) \cdot g(x))), (n + m)) \end{array}}}
\vnbprop{poly-deg-unique}{\vnbstmt{\begin{array}{@{}l@{}} \forall k, a, f.; \forall n, m \in \mathbb{N}. \\ \quad k \subseteq \mathbb{C} \wedge \\ \quad \exists s \in (\mathbb{N} \to a). \forall i, j \in \mathbb{N}. \\ \quad \quad (s(i) = s(j)) \Rightarrow (i = j) \wedge \\ \quad \operatorname{is\text{-}poly\text{-}deg}(k, a, f, n) \wedge \\ \quad \operatorname{is\text{-}poly\text{-}deg}(k, a, f, m) \Rightarrow \\ \quad \quad n = m \end{array}}}
\vnbprop{deg-value}{\vnbstmt{\begin{array}{@{}l@{}} \forall k, a, f.; \forall n \in \mathbb{N}. \\ \quad k \subseteq \mathbb{C} \wedge \\ \quad \exists s \in (\mathbb{N} \to a). \forall i, j \in \mathbb{N}. \\ \quad \quad (s(i) = s(j)) \Rightarrow (i = j) \wedge \\ \quad \operatorname{is\text{-}poly\text{-}deg}(k, a, f, n) \Rightarrow \\ \quad \quad \operatorname{poly\text{-}fn\text{-}deg}(k, a, f) = n \end{array}}}
\vnbprop{deg-mul}{\vnbstmt{\begin{array}{@{}l@{}} \forall a, f, g.; \forall n, m \in \mathbb{N}. \\ \quad a \subseteq \mathbb{C} \wedge \\ \quad \exists s \in (\mathbb{N} \to a). \forall i, j \in \mathbb{N}. \\ \quad \quad (s(i) = s(j)) \Rightarrow (i = j) \wedge \\ \quad \operatorname{is\text{-}poly\text{-}deg}(\mathbb{C}, a, f, n) \wedge \\ \quad \operatorname{is\text{-}poly\text{-}deg}(\mathbb{C}, a, g, m) \Rightarrow \\ \quad \quad \operatorname{poly\text{-}fn\text{-}deg}(\mathbb{C}, a, (\lambda x \in a.\; (f(x) \cdot g(x)))) = (\operatorname{poly\text{-}fn\text{-}deg}(\mathbb{C}, a, f) + \operatorname{poly\text{-}fn\text{-}deg}(\mathbb{C}, a, g)) \end{array}}}
\vnbprop{poly-deg-roots}{\vnbstmt{\begin{array}{@{}l@{}} \forall a, f.; \forall n \in \mathbb{N}, s \in (\mathbb{N} \to a). \\ \quad \operatorname{is\text{-}poly\text{-}deg}(\mathbb{C}, a, f, n) \wedge \\ \quad \forall i, j \in \mathbb{N}.\; (((i \leq n) \wedge ((j \leq n) \wedge (s(i) = s(j)))) \Rightarrow (i = j)) \Rightarrow \\ \quad \quad \neg \forall i \in \mathbb{N}.\; ((i \leq n) \Rightarrow (f(s(i)) = 0)) \end{array}}}

\vnbbookfile{theorem-library/cc-diff-laws-2}{theorem-library/cc-diff-laws-2}
\noindent{\small cc-diff-laws-2.scm -- the rest of the elementary theory of IS-CC-DIFF-AT, the derivative of a map RR -$>$ CC (structure-library/cc-coords.scm). Continues theorem-library/cc-coords-laws.scm (batch 14-A), which defined the predicate coordinatewise and proved the transfer equivalence}\par
\vnbprop{cc-diff-at-implies-continuous}{\vnbstmt{\begin{array}{@{}l@{}} \forall g \in (\mathbb{R} \to \mathbb{C}), v \in \mathbb{R}, a \in \mathbb{C}. \\ \quad \operatorname{is\text{-}cc\text{-}diff\text{-}at}(g, v, a) \Rightarrow \operatorname{is\text{-}continuous\text{-}at}(\operatorname{rr\text{-}ms}, \operatorname{cc\text{-}ms}, g, v) \end{array}}}
\vnbprop{cc-diff-at-const}{\vnbstmt{\begin{array}{@{}l@{}} \forall v \in \mathbb{C}, g \in (\mathbb{R} \to \mathbb{C}), \mathit{tv} \in \mathbb{R}. \\ \quad \forall x \in \mathbb{R}.\; (g(x) = v) \Rightarrow \operatorname{is\text{-}cc\text{-}diff\text{-}at}(g, \mathit{tv}, 0) \end{array}}}
\vnbprop{cc-diff-at-cc-mul}{\vnbstmt{\begin{array}{@{}l@{}} \forall v \in \mathbb{C}, g, d \in (\mathbb{R} \to \mathbb{C}), \mathit{tv} \in \mathbb{R}, a \in \mathbb{C}. \\ \quad (\forall x \in \mathbb{R}.\; (d(x) = (v \cdot g(x))) \wedge \operatorname{is\text{-}cc\text{-}diff\text{-}at}(g, \mathit{tv}, a)) \Rightarrow  \\ \quad \quad \operatorname{is\text{-}cc\text{-}diff\text{-}at}(d, \mathit{tv}, (v \cdot a)) \end{array}}}
\vnbprop{cc-diff-at-difference}{\vnbstmt{\begin{array}{@{}l@{}} \forall g, h, b \in (\mathbb{R} \to \mathbb{C}), v \in \mathbb{R}, a, c \in \mathbb{C}. \\ \quad (\forall x \in \mathbb{R}.\; (b(x) = (g(x) - h(x))) \wedge (\operatorname{is\text{-}cc\text{-}diff\text{-}at}(g, v, a) \wedge \operatorname{is\text{-}cc\text{-}diff\text{-}at}(h, v, c))) \Rightarrow  \\ \quad \quad \operatorname{is\text{-}cc\text{-}diff\text{-}at}(b, v, (a - c)) \end{array}}}
\vnbprop{cc-diff-at-unique}{\vnbstmt{\begin{array}{@{}l@{}} \forall g \in (\mathbb{R} \to \mathbb{C}), v \in \mathbb{R}, a, b \in \mathbb{C}. \\ \quad (\operatorname{is\text{-}cc\text{-}diff\text{-}at}(g, v, a) \wedge \operatorname{is\text{-}cc\text{-}diff\text{-}at}(g, v, b)) \Rightarrow (a = b) \end{array}}}
\vnbprop{cc-diff-at-chain-real}{\vnbstmt{\begin{array}{@{}l@{}} \forall g \in (\mathbb{R} \to \mathbb{C}), h \in (\mathbb{R} \to \mathbb{R}), m \in (\mathbb{R} \to \mathbb{C}), v, a \in \mathbb{R}, b \in \mathbb{C}. \\ \quad (\forall x \in \mathbb{R}.\; (m(x) = g(h(x))) \wedge (\operatorname{is\text{-}diff\text{-}at}(h, v, a) \wedge \operatorname{is\text{-}cc\text{-}diff\text{-}at}(g, h(v), b))) \Rightarrow  \\ \quad \quad \operatorname{is\text{-}cc\text{-}diff\text{-}at}(m, v, (a \cdot b)) \end{array}}}

\vnbbookfile{theorem-library/poly-fn-weierstrass}{theorem-library/poly-fn-weierstrass}
\noindent{\small poly-fn-weierstrass.scm -- calculus 5.7 (Weierstrass approximation) as the notes state it, and the polynomial-function facts about the Bernstein polynomials it needs (CA-MONTEL "POLY", 2026-10-26). Vocabulary: poly-fn-vocab.scm; laws: poly-fn-laws.scm.  Library: analysis.  Helper prefix pfw-.}\par
\vnbprop{bernstein-basis-poly-fn}{\vnbstmt{\begin{array}{@{}l@{}} \forall n \in \mathbb{N}, k \in \mathbb{Z}. \\ \quad \operatorname{is\text{-}poly\text{-}fn\text{-}on}(\mathbb{R}, \mathbb{R}, (\lambda x \in \mathbb{R}.\; \operatorname{bernstein\text{-}basis}(n, x)(k))) \end{array}}}
\vnbprop{bernstein-partial-poly-fn}{\vnbstmt{\begin{array}{@{}l@{}} \forall m \in \mathbb{N}, g.; \forall n \in \mathbb{N}, a, b \in \mathbb{R}. \\ \quad (\forall y \in \mathbb{R}.\; (g(y) \in \mathbb{R}) \wedge (\neg (n = 0) \wedge (a < b))) \Rightarrow  \\ \quad \quad \operatorname{is\text{-}poly\text{-}fn\text{-}on}(\mathbb{R}, \mathbb{R}, \\ \quad \quad \quad (\lambda x \in \mathbb{R}.\; \\ \quad \quad \quad \quad \operatorname{series\text{-}partial\text{-}sum}((\lambda l \in \mathbb{N}.\; (g(/(l, n)) \cdot \operatorname{bernstein\text{-}basis}(n, ((x - a) \cdot \operatorname{recip}((b - a))))(l))), \\ \quad \quad \quad \quad \quad m))) \end{array}}}
\vnbprop{weierstrass-approximation}{\vnbstmt{\begin{array}{@{}l@{}} \forall a, b \in \mathbb{R}, f. \\ \quad ((a < b) \wedge \operatorname{is\text{-}continuous\text{-}on}(f, \operatorname{ccint}(a, b))) \Rightarrow  \\ \quad \quad \exists p. \\ \quad \quad \quad \forall n \in \mathbb{N}.\; \operatorname{is\text{-}poly\text{-}fn\text{-}on}(\mathbb{R}, \operatorname{ccint}(a, b), p(n)) \wedge \\ \quad \quad \quad \forall s. \\ \quad \quad \quad \quad \operatorname{pos\text{-}rr}(s) \Rightarrow  \\ \quad \quad \quad \quad \quad \exists c \in \mathbb{N}. \forall n \in \mathbb{N}, x \in \operatorname{ccint}(a, b). \\ \quad \quad \quad \quad \quad \quad (c \leq n) \Rightarrow (\lvert (f(x) - p(n)(x)) \rvert \leq s) \end{array}}}

\vnbbookfile{structure-library/cc-power-series}{structure-library/cc-power-series}
\noindent{\small cc-power-series.scm -- THE VOCABULARY OF COMPLEX POWER SERIES (the notes 2.2). DEFINITIONS ONLY; the laws are PROVEN in theorem-library/cc-power-series-laws.scm (and the root / ratio tests they rest on in theorem-library/limsup-tests.scm). Batch 25-A, 2026-09-24.  Every definition below is a DECISION FOR THE USER.}\par
\par\medskip\noindent\textbf{Definition} (\texttt{cps-radius}).\;$\operatorname{cps\text{-}radius}(\mathit{cf}) :=$
\begin{equation*}
\operatorname{esup}(\{\operatorname{cpsr\_} \in \mathbb{R} : ((0 \leq \operatorname{cpsr\_}) \wedge \operatorname{series\text{-}converges}((\lambda \operatorname{cpsk\_} \in \mathbb{N}.\; (\operatorname{magnitude}(\operatorname{cf}(\operatorname{cpsk\_})) \cdot \operatorname{power}(\operatorname{cpsr\_}, \operatorname{cpsk\_})))))\})
\end{equation*}

\par\medskip\noindent\textbf{Definition} (\texttt{cc-sup-norm}).\;$\operatorname{cc\text{-}sup\text{-}norm}(u, \operatorname{cpsxs\_}) :=$
\begin{equation*}
\operatorname{esup}(\operatorname{image}((\lambda \operatorname{cpsx\_} \in \operatorname{cpsxs\_}.\; \operatorname{magnitude}(u(\operatorname{cpsx\_}))), \operatorname{cpsxs\_}))
\end{equation*}

\par\medskip\noindent\textbf{Definition} (\texttt{is-normally-convergent}).\;For all $u$ and $\operatorname{cpsxs\_}$, 
\begin{equation*}
\begin{array}{@{}l@{}} \operatorname{is\text{-}normally\text{-}convergent}(u, \operatorname{cpsxs\_}) \\ \Leftrightarrow\; \forall \operatorname{cpsi\_} \in \mathbb{N}.\; (\operatorname{cc\text{-}sup\text{-}norm}(u(\operatorname{cpsi\_}), \operatorname{cpsxs\_}) \in \mathbb{R}) \wedge \\ \operatorname{series\text{-}converges}((\lambda \operatorname{cpsi\_} \in \mathbb{N}.\; \operatorname{cc\text{-}sup\text{-}norm}(u(\operatorname{cpsi\_}), \operatorname{cpsxs\_}))) \end{array}
\end{equation*}

\par\medskip\noindent\textbf{Definition} (\texttt{cc-series-limit}).\;$\operatorname{cc\text{-}series\text{-}limit}(f) :=$ $(\iota\, \operatorname{cpsl\_}.\;\operatorname{cc\text{-}series\text{-}converges\text{-}to}(f, \operatorname{cpsl\_}))$

\vnbbookfile{theorem-library/cc-power-series-laws}{theorem-library/cc-power-series-laws}
\noindent{\small cc-power-series-laws.scm -- COMPLEX POWER SERIES: THE RADIUS OF CONVERGENCE (Prop 2.7) AND WHAT FOLLOWS FROM IT, as the notes state them. Batch 25-A, 2026-09-24.}\par
\vnbprop{rr-power-mul-base}{\vnbstmt{\forall k \in \mathbb{N}, x, y \in \mathbb{R}.\; (\operatorname{power}((x \cdot y), k) = (\operatorname{power}(x, k) \cdot \operatorname{power}(y, k)))}}
\vnbprop{cps-root-conv}{\vnbstmt{\begin{array}{@{}l@{}} \forall f \in (\mathbb{N} \to \mathbb{C}), s, r \in \mathbb{R}. \\ \quad ((\operatorname{elimsup}((\lambda j \in \mathbb{N}.\; {\operatorname{magnitude}(f(j))}^{{j}^{\ast\!-1}})) < s) \wedge ((0 \leq r) \wedge ((r \cdot s) < 1))) \Rightarrow  \\ \quad \quad \operatorname{series\text{-}converges}((\lambda k \in \mathbb{N}.\; (\operatorname{magnitude}(f(k)) \cdot \operatorname{power}(r, k)))) \end{array}}}
\vnbprop{cps-root-div}{\vnbstmt{\begin{array}{@{}l@{}} \forall f \in (\mathbb{N} \to \mathbb{C}), s, r \in \mathbb{R}. \\ \quad 0 \leq s \wedge \\ \quad s < \operatorname{elimsup}((\lambda j \in \mathbb{N}.\; {\operatorname{magnitude}(f(j))}^{{j}^{\ast\!-1}})) \wedge \\ \quad 0 \leq r \wedge \\ \quad 1 \leq (r \cdot s) \Rightarrow \\ \quad \quad \neg \operatorname{series\text{-}converges}((\lambda k \in \mathbb{N}.\; (\operatorname{magnitude}(f(k)) \cdot \operatorname{power}(r, k)))) \end{array}}}
\vnbprop{cps-small-factor}{\vnbstmt{\forall r \in \mathbb{R}.\; ((0 \leq r) \Rightarrow \exists s \in \mathbb{R}.\; ((0 < s) \wedge ((r \cdot s) < 1)))}}
\vnbprop{cps-radius-infinite}{\vnbstmt{\begin{array}{@{}l@{}} \forall f \in (\mathbb{N} \to \mathbb{C}). \\ \quad (\operatorname{elimsup}((\lambda j \in \mathbb{N}.\; {\operatorname{magnitude}(f(j))}^{{j}^{\ast\!-1}})) = 0) \Rightarrow  \\ \quad \quad \operatorname{cps\text{-}radius}(f) = \operatorname{pos\text{-}inf} \end{array}}}
\vnbprop{cps-radius-zero}{\vnbstmt{\begin{array}{@{}l@{}} \forall f \in (\mathbb{N} \to \mathbb{C}). \\ \quad (\operatorname{elimsup}((\lambda j \in \mathbb{N}.\; {\operatorname{magnitude}(f(j))}^{{j}^{\ast\!-1}})) = \operatorname{pos\text{-}inf}) \Rightarrow  \\ \quad \quad \operatorname{cps\text{-}radius}(f) = 0 \end{array}}}
\vnbprop{cps-radius-recip}{\vnbstmt{\begin{array}{@{}l@{}} \forall f \in (\mathbb{N} \to \mathbb{C}). \\ \quad \operatorname{elimsup}((\lambda j \in \mathbb{N}.\; {\operatorname{magnitude}(f(j))}^{{j}^{\ast\!-1}})) \in \mathbb{R} \wedge \\ \quad 0 < \operatorname{elimsup}((\lambda j \in \mathbb{N}.\; {\operatorname{magnitude}(f(j))}^{{j}^{\ast\!-1}})) \Rightarrow \\ \quad \quad \operatorname{cps\text{-}radius}(f) = \operatorname{recip}(\operatorname{elimsup}((\lambda j \in \mathbb{N}.\; {\operatorname{magnitude}(f(j))}^{{j}^{\ast\!-1}}))) \end{array}}}
\vnbprop{cps-radius-root}{\vnbstmt{\begin{array}{@{}l@{}} \forall f \in (\mathbb{N} \to \mathbb{C}). \\ \quad (\operatorname{elimsup}((\lambda j \in \mathbb{N}.\; {\operatorname{magnitude}(f(j))}^{{j}^{\ast\!-1}})) = 0) \Rightarrow  \\ \quad \quad \operatorname{cps\text{-}radius}(f) = \operatorname{pos\text{-}inf} \wedge \\ \quad (\operatorname{elimsup}((\lambda j \in \mathbb{N}.\; {\operatorname{magnitude}(f(j))}^{{j}^{\ast\!-1}})) = \operatorname{pos\text{-}inf}) \Rightarrow  \\ \quad \quad \operatorname{cps\text{-}radius}(f) = 0 \wedge \\ \quad \operatorname{elimsup}((\lambda j \in \mathbb{N}.\; {\operatorname{magnitude}(f(j))}^{{j}^{\ast\!-1}})) \in \mathbb{R} \wedge \\ \quad 0 < \operatorname{elimsup}((\lambda j \in \mathbb{N}.\; {\operatorname{magnitude}(f(j))}^{{j}^{\ast\!-1}})) \Rightarrow \\ \quad \quad \operatorname{cps\text{-}radius}(f) = \operatorname{recip}(\operatorname{elimsup}((\lambda j \in \mathbb{N}.\; {\operatorname{magnitude}(f(j))}^{{j}^{\ast\!-1}}))) \end{array}}}
\vnbprop{cps-radius-unfold}{\vnbstmt{\begin{array}{@{}l@{}} \forall f. \\ \quad (\operatorname{cps\text{-}radius}(f) \simeq \operatorname{esup}(\{r \in \mathbb{R} : ((0 \leq r) \wedge \operatorname{series\text{-}converges}((\lambda k \in \mathbb{N}.\; (\operatorname{magnitude}(f(k)) \cdot \operatorname{power}(r, k)))))\})) \end{array}}}
\vnbprop{cc-closed-ball-origin}{\vnbstmt{\forall d.; \forall z \in \operatorname{closed\text{-}ball}(\operatorname{cc\text{-}ms}, 0, d).\; ((z \in \mathbb{C}) \wedge (\operatorname{magnitude}(z) \leq d))}}
\vnbprop{cc-magnitude-power}{\vnbstmt{\forall k \in \mathbb{N}, z \in \mathbb{C}.\; (\operatorname{magnitude}(\operatorname{power}(z, k)) = \operatorname{power}(\operatorname{magnitude}(z), k))}}
\vnbprop{cc-sup-norm-bound}{\vnbstmt{\begin{array}{@{}l@{}} \forall u, s.; \forall m \in \mathbb{R}. \\ \quad ((0 \leq m) \wedge (\forall x \in s.\; (u(x) \in \mathbb{C}) \wedge \forall x \in s.\; (\operatorname{magnitude}(u(x)) \leq m))) \Rightarrow  \\ \quad \quad (\operatorname{cc\text{-}sup\text{-}norm}(u, s) \in \mathbb{R}) \wedge ((0 \leq \operatorname{cc\text{-}sup\text{-}norm}(u, s)) \wedge (\operatorname{cc\text{-}sup\text{-}norm}(u, s) \leq m)) \end{array}}}
\vnbprop{cps-radius-approx}{\vnbstmt{\begin{array}{@{}l@{}} \forall f \in (\mathbb{N} \to \mathbb{C}), d \in \mathbb{R}. \\ \quad ((0 \leq d) \wedge (d < \operatorname{cps\text{-}radius}(f))) \Rightarrow  \\ \quad \quad \exists t \in \mathbb{R}. \\ \quad \quad \quad (d < t) \wedge \operatorname{series\text{-}converges}((\lambda k \in \mathbb{N}.\; (\operatorname{magnitude}(f(k)) \cdot \operatorname{power}(t, k)))) \end{array}}}
\vnbprop{cps-normally-convergent}{\vnbstmt{\begin{array}{@{}l@{}} \forall f \in (\mathbb{N} \to \mathbb{C}), d \in \mathbb{R}. \\ \quad ((0 \leq d) \wedge (d < \operatorname{cps\text{-}radius}(f))) \Rightarrow  \\ \quad \quad \operatorname{is\text{-}normally\text{-}convergent}((\lambda k \in \mathbb{N}.\; (\lambda z \in \operatorname{closed\text{-}ball}(\operatorname{cc\text{-}ms}, 0, d).\; (f(k) \cdot \operatorname{power}(z, k)))), \\ \quad \quad \quad \operatorname{closed\text{-}ball}(\operatorname{cc\text{-}ms}, 0, d)) \end{array}}}
\vnbprop{cps-abs-conv-closed-ball}{\vnbstmt{\begin{array}{@{}l@{}} \forall f \in (\mathbb{N} \to \mathbb{C}), d \in \mathbb{R}, z \in \operatorname{closed\text{-}ball}(\operatorname{cc\text{-}ms}, 0, d). \\ \quad ((0 \leq d) \wedge (d < \operatorname{cps\text{-}radius}(f))) \Rightarrow  \\ \quad \quad \operatorname{series\text{-}converges}((\lambda k \in \mathbb{N}.\; \operatorname{magnitude}((f(k) \cdot \operatorname{power}(z, k))))) \end{array}}}
\vnbprop{cc-series-block-le}{\vnbstmt{\begin{array}{@{}l@{}} \forall k, m \in \mathbb{N}, f \in (\mathbb{N} \to \mathbb{C}), g \in (\mathbb{N} \to \mathbb{R}). \\ \quad (\forall j \in \mathbb{N}.\; (g(j) = \operatorname{magnitude}(f(j))) \wedge (m \leq k)) \Rightarrow  \\ \quad \quad \operatorname{magnitude}((\operatorname{cc\text{-}series\text{-}partial\text{-}sum}(f, k) - \operatorname{cc\text{-}series\text{-}partial\text{-}sum}(f, m))) \\ \quad \quad \leq\; (\operatorname{series\text{-}partial\text{-}sum}(g, k) - \operatorname{series\text{-}partial\text{-}sum}(g, m)) \end{array}}}
\vnbprop{cc-series-abs-converges}{\vnbstmt{\begin{array}{@{}l@{}} \forall f \in (\mathbb{N} \to \mathbb{C}), g \in (\mathbb{N} \to \mathbb{R}). \\ \quad (\forall j \in \mathbb{N}.\; (g(j) = \operatorname{magnitude}(f(j))) \wedge \operatorname{series\text{-}converges}(g)) \Rightarrow  \\ \quad \quad \operatorname{cc\text{-}series\text{-}converges}(f) \end{array}}}
\vnbprop{cc-series-limit-converges-to}{\vnbstmt{\begin{array}{@{}l@{}} \forall f \in (\mathbb{N} \to \mathbb{C}). \\ \quad \operatorname{cc\text{-}series\text{-}converges}(f) \Rightarrow \operatorname{cc\text{-}series\text{-}converges\text{-}to}(f, \operatorname{cc\text{-}series\text{-}limit}(f)) \end{array}}}
\vnbprop{cc-series-limit-in-cc}{\vnbstmt{\begin{array}{@{}l@{}} \forall f \in (\mathbb{N} \to \mathbb{C}). \\ \quad \operatorname{cc\text{-}series\text{-}converges}(f) \Rightarrow (\operatorname{cc\text{-}series\text{-}limit}(f) \in \mathbb{C}) \end{array}}}
\vnbprop{cps-converges-closed-ball}{\vnbstmt{\begin{array}{@{}l@{}} \forall f \in (\mathbb{N} \to \mathbb{C}), d \in \mathbb{R}, z \in \operatorname{closed\text{-}ball}(\operatorname{cc\text{-}ms}, 0, d). \\ \quad ((0 \leq d) \wedge (d < \operatorname{cps\text{-}radius}(f))) \Rightarrow  \\ \quad \quad \operatorname{cc\text{-}series\text{-}converges}((\lambda k \in \mathbb{N}.\; (f(k) \cdot \operatorname{power}(z, k)))) \end{array}}}
\vnbprop{cps-uniform-closed-ball}{\vnbstmt{\begin{array}{@{}l@{}} \forall f \in (\mathbb{N} \to \mathbb{C}), d \in \mathbb{R}, e. \\ \quad ((0 \leq d) \wedge ((d < \operatorname{cps\text{-}radius}(f)) \wedge \operatorname{pos\text{-}rr}(e))) \Rightarrow  \\ \quad \quad \exists n \in \mathbb{N}. \forall m \in \mathbb{N}, z \in \operatorname{closed\text{-}ball}(\operatorname{cc\text{-}ms}, 0, d). \\ \quad \quad \quad (n \leq m) \Rightarrow  \\ \quad \quad \quad \quad \operatorname{magnitude}((\operatorname{cc\text{-}series\text{-}limit}((\lambda k \in \mathbb{N}.\; (f(k) \cdot \operatorname{power}(z, k)))) - \operatorname{cc\text{-}series\text{-}partial\text{-}sum}((\lambda k \in \mathbb{N}.\; (f(k) \cdot \operatorname{power}(z, k))), m))) \\ \quad \quad \quad \quad \leq\; e \end{array}}}
\vnbprop{rr-power-power}{\vnbstmt{\forall k, m \in \mathbb{N}, x \in \mathbb{R}.\; (\operatorname{power}(\operatorname{power}(x, k), m) = \operatorname{power}(\operatorname{power}(x, m), k))}}
\vnbprop{cps-poly-geom-bound}{\vnbstmt{\begin{array}{@{}l@{}} \forall m \in \mathbb{N}, q \in \mathbb{R}. \\ \quad ((0 \leq q) \wedge (q < 1)) \Rightarrow  \\ \quad \quad \exists b \in \mathbb{R}.\; ((0 \leq b) \wedge \forall k \in \mathbb{N}.\; ((\operatorname{power}(k, m) \cdot \operatorname{power}(q, k)) \leq b)) \end{array}}}
\vnbprop{series-scale-converges}{\vnbstmt{\begin{array}{@{}l@{}} \forall f \in (\mathbb{N} \to \mathbb{R}), c \in \mathbb{R}, h \in (\mathbb{N} \to \mathbb{R}). \\ \quad \forall k \in \mathbb{N}.\; (0 \leq f(k)) \wedge \\ \quad 0 \leq c \wedge \\ \quad \forall k \in \mathbb{N}.\; (h(k) = (c \cdot f(k))) \wedge \\ \quad \operatorname{series\text{-}converges}(f) \Rightarrow \\ \quad \quad \operatorname{series\text{-}converges}(h) \end{array}}}
\vnbprop{cps-factor-conv}{\vnbstmt{\begin{array}{@{}l@{}} \forall f, g, a \in (\mathbb{N} \to \mathbb{C}), c \in \mathbb{R}, m \in \mathbb{N}, t, d \in \mathbb{R}. \\ \quad \forall k \in \mathbb{N}.\; (a(k) = (f(k) \cdot g(k))) \wedge \\ \quad \forall k \in \mathbb{N}.\; (\operatorname{magnitude}(g(k)) \leq (c \cdot \operatorname{power}(k, m))) \wedge \\ \quad 0 \leq d \wedge \\ \quad d < t \wedge \\ \quad \operatorname{series\text{-}converges}((\lambda k \in \mathbb{N}.\; (\operatorname{magnitude}(f(k)) \cdot \operatorname{power}(t, k)))) \Rightarrow \\ \quad \quad \operatorname{series\text{-}converges}((\lambda k \in \mathbb{N}.\; (\operatorname{magnitude}(a(k)) \cdot \operatorname{power}(d, k)))) \end{array}}}
\vnbprop{cps-radius-factor}{\vnbstmt{\begin{array}{@{}l@{}} \forall f, g, a \in (\mathbb{N} \to \mathbb{C}), c \in \mathbb{R}, m \in \mathbb{N}. \\ \quad (\forall k \in \mathbb{N}.\; (a(k) = (f(k) \cdot g(k))) \wedge \forall k \in \mathbb{N}.\; (\operatorname{magnitude}(g(k)) \leq (c \cdot \operatorname{power}(k, m)))) \Rightarrow  \\ \quad \quad \operatorname{cps\text{-}radius}(f) \leq \operatorname{cps\text{-}radius}(a) \end{array}}}

\vnbbookfile{theorem-library/cc-power-series-deriv}{theorem-library/cc-power-series-deriv}
\noindent{\small cc-power-series-deriv.scm -- PROPOSITION 2.10: THE SUM OF A COMPLEX POWER SERIES IS HOLOMORPHIC INSIDE THE RADIUS, AND ITS DERIVATIVE IS THE TERMWISE DERIVED SERIES, WHICH HAS THE SAME RADIUS. Batch 26-A, 2026-09-24.}\par
\vnbprop{cpd-cc-squeeze}{\vnbstmt{\begin{array}{@{}l@{}} \forall s \in (\mathbb{N} \to \mathbb{C}), v \in \mathbb{C}, m \in (\mathbb{N} \to \mathbb{R}). \\ \quad (\operatorname{converges\text{-}to}(\operatorname{rr\text{-}ms}, m, 0) \wedge \forall j \in \mathbb{N}.\; (\operatorname{magnitude}((s(j) - v)) \leq m(j))) \Rightarrow  \\ \quad \quad \operatorname{converges\text{-}to}(\operatorname{cc\text{-}ms}, s, v) \end{array}}}
\vnbprop{cc-limit-lincomb}{\vnbstmt{\begin{array}{@{}l@{}} \forall s, t, h \in (\mathbb{N} \to \mathbb{C}), v, a, c, d \in \mathbb{C}. \\ \quad \forall j \in \mathbb{N}.\; (h(j) = ((c \cdot s(j)) + (d \cdot t(j)))) \wedge \\ \quad \operatorname{converges\text{-}to}(\operatorname{cc\text{-}ms}, s, v) \wedge \\ \quad \operatorname{converges\text{-}to}(\operatorname{cc\text{-}ms}, t, a) \Rightarrow \\ \quad \quad \operatorname{converges\text{-}to}(\operatorname{cc\text{-}ms}, h, ((c \cdot v) + (d \cdot a))) \end{array}}}
\vnbprop{cc-limit-shift}{\vnbstmt{\begin{array}{@{}l@{}} \forall s, h \in (\mathbb{N} \to \mathbb{C}), v \in \mathbb{C}. \\ \quad (\forall j \in \mathbb{N}.\; (h(j) = s((j + 1))) \wedge \operatorname{converges\text{-}to}(\operatorname{cc\text{-}ms}, s, v)) \Rightarrow  \\ \quad \quad \operatorname{converges\text{-}to}(\operatorname{cc\text{-}ms}, h, v) \end{array}}}
\vnbprop{cc-series-partial-sum-lincomb}{\vnbstmt{\begin{array}{@{}l@{}} \forall k \in \mathbb{N}, f, g, h \in (\mathbb{N} \to \mathbb{C}), c, d \in \mathbb{C}. \\ \quad \forall j \in \mathbb{N}.\; (h(j) = ((c \cdot f(j)) + (d \cdot g(j)))) \Rightarrow  \\ \quad \quad \operatorname{cc\text{-}series\text{-}partial\text{-}sum}(h, k) = ((c \cdot \operatorname{cc\text{-}series\text{-}partial\text{-}sum}(f, k)) + (d \cdot \operatorname{cc\text{-}series\text{-}partial\text{-}sum}(g, k))) \end{array}}}
\vnbprop{cc-series-limit-linear}{\vnbstmt{\begin{array}{@{}l@{}} \forall f, g, h \in (\mathbb{N} \to \mathbb{C}), c, d \in \mathbb{C}. \\ \quad (\forall j \in \mathbb{N}.\; (h(j) = ((c \cdot f(j)) + (d \cdot g(j)))) \wedge (\operatorname{cc\text{-}series\text{-}converges}(f) \wedge \operatorname{cc\text{-}series\text{-}converges}(g))) \Rightarrow  \\ \quad \quad (\operatorname{cc\text{-}series\text{-}converges}(h) \wedge (\operatorname{cc\text{-}series\text{-}limit}(h) = ((c \cdot \operatorname{cc\text{-}series\text{-}limit}(f)) + (d \cdot \operatorname{cc\text{-}series\text{-}limit}(g))))) \end{array}}}
\vnbprop{cpd-power-second-order}{\vnbstmt{\begin{array}{@{}l@{}} \forall k \in \mathbb{N}, z, u \in \mathbb{C}, r \in \mathbb{R}. \\ \quad ((\operatorname{magnitude}(z) \leq r) \wedge (\operatorname{magnitude}(u) \leq r)) \Rightarrow  \\ \quad \quad r \cdot \operatorname{magnitude}(((\operatorname{power}(z, (k + 1)) - \operatorname{power}(u, (k + 1))) - (((k + 1) \cdot \operatorname{power}(u, k)) \cdot (z - u)))) \\ \quad \quad \leq\; (((k + 1) \cdot (k + 1)) \cdot \operatorname{power}(r, k)) \cdot (\operatorname{magnitude}((z - u)) \cdot \operatorname{magnitude}((z - u))) \end{array}}}
\vnbprop{cc-power-diff-bound}{\vnbstmt{\begin{array}{@{}l@{}} \forall k \in \mathbb{N}, z, u \in \mathbb{C}, r \in \mathbb{R}. \\ \quad ((\operatorname{magnitude}(z) \leq r) \wedge (\operatorname{magnitude}(u) \leq r)) \Rightarrow  \\ \quad \quad \operatorname{magnitude}((\operatorname{power}(z, (k + 1)) - \operatorname{power}(u, (k + 1)))) \leq (((k + 1) \cdot \operatorname{power}(r, k)) \cdot \operatorname{magnitude}((z - u))) \end{array}}}
\vnbprop{cps-shift-conv}{\vnbstmt{\begin{array}{@{}l@{}} \forall f, g \in (\mathbb{N} \to \mathbb{C}), t \in \mathbb{R}. \\ \quad (\forall j \in \mathbb{N}.\; (\operatorname{magnitude}(g(j)) = \operatorname{magnitude}(f((j + 1)))) \wedge (0 < t)) \Rightarrow  \\ \quad \quad \operatorname{series\text{-}converges}((\lambda k \in \mathbb{N}.\; (\operatorname{magnitude}(f(k)) \cdot \operatorname{power}(t, k)))) \\ \quad \quad \Leftrightarrow\; \operatorname{series\text{-}converges}((\lambda k \in \mathbb{N}.\; (\operatorname{magnitude}(g(k)) \cdot \operatorname{power}(t, k)))) \end{array}}}
\vnbprop{cps-radius-in}{\vnbstmt{\forall f.\; (\operatorname{cps\text{-}radius}(f) \in \operatorname{rr\text{-}pos\text{-}star})}}
\vnbprop{cps-radius-mono}{\vnbstmt{\begin{array}{@{}l@{}} \forall f, g \in (\mathbb{N} \to \mathbb{C}). \\ \quad \forall t \in \mathbb{R}. \\ \quad \quad ((0 < t) \wedge \operatorname{series\text{-}converges}((\lambda k \in \mathbb{N}.\; (\operatorname{magnitude}(f(k)) \cdot \operatorname{power}(t, k))))) \Rightarrow  \\ \quad \quad \quad \operatorname{series\text{-}converges}((\lambda k \in \mathbb{N}.\; (\operatorname{magnitude}(g(k)) \cdot \operatorname{power}(t, k)))) \Rightarrow \\ \quad \quad \operatorname{cps\text{-}radius}(f) \leq \operatorname{cps\text{-}radius}(g) \end{array}}}
\vnbprop{cps-radius-derived}{\vnbstmt{\begin{array}{@{}l@{}} \forall f, a \in (\mathbb{N} \to \mathbb{C}). \\ \quad \forall j \in \mathbb{N}.\; (a(j) = ((j + 1) \cdot f((j + 1)))) \Rightarrow  \\ \quad \quad \operatorname{cps\text{-}radius}(a) = \operatorname{cps\text{-}radius}(f) \end{array}}}
\vnbprop{cc-diff-on-quadratic}{\vnbstmt{\begin{array}{@{}l@{}} \forall s.; \forall f \in (s \to \mathbb{C}), a \in s, l \in \mathbb{C}, m \in \mathbb{R}. \\ \quad \operatorname{is\text{-}open}(\operatorname{nf\text{-}metric\text{-}space}(\operatorname{cc\text{-}normed\text{-}field}), s) \wedge \\ \quad 0 \leq m \wedge \\ \quad \forall x \in s. \\ \quad \quad \operatorname{magnitude}(((f(x) - f(a)) - ((x - a) \cdot l))) \leq (m \cdot (\operatorname{magnitude}((x - a)) \cdot \operatorname{magnitude}((x - a)))) \Rightarrow \\ \quad \quad \operatorname{is\text{-}diff\text{-}on}(\operatorname{cc\text{-}normed\text{-}field}, s, f, a, l) \end{array}}}
\vnbprop{cps-partial-error-bound}{\vnbstmt{\begin{array}{@{}l@{}} \forall k \in \mathbb{N}, f, a \in (\mathbb{N} \to \mathbb{C}), z, u \in \mathbb{C}, r \in \mathbb{R}, \mathit{cpdtz}, \mathit{cpdtu}, d \in (\mathbb{N} \to \mathbb{C}), v \in (\mathbb{N} \to \mathbb{R}). \\ \quad \forall j \in \mathbb{N}.\; (a(j) = ((j + 1) \cdot f((j + 1)))) \wedge \\ \quad \operatorname{magnitude}(z) \leq r \wedge \\ \quad \operatorname{magnitude}(u) \leq r \wedge \\ \quad \forall j \in \mathbb{N}.\; (\operatorname{cpdtz}(j) = (f(j) \cdot \operatorname{power}(z, j))) \wedge \\ \quad \forall j \in \mathbb{N}.\; (\operatorname{cpdtu}(j) = (f(j) \cdot \operatorname{power}(u, j))) \wedge \\ \quad \forall j \in \mathbb{N}.\; (d(j) = (a(j) \cdot \operatorname{power}(u, j))) \wedge \\ \quad \forall j \in \mathbb{N}. \\ \quad \quad v(j) = (\operatorname{magnitude}((f((j + 1)) \cdot ((j + 1) \cdot (j + 1)))) \cdot \operatorname{power}(r, (j + 1))) \Rightarrow \\ \quad \quad (r \cdot r) \cdot \operatorname{magnitude}(((\operatorname{cc\text{-}series\text{-}partial\text{-}sum}(\mathit{cpdtz}, (k + 1)) - \operatorname{cc\text{-}series\text{-}partial\text{-}sum}(\mathit{cpdtu}, (k + 1))) - ((z - u) \cdot \operatorname{cc\text{-}series\text{-}partial\text{-}sum}(d, k)))) \\ \quad \quad \leq\; (\operatorname{magnitude}((z - u)) \cdot \operatorname{magnitude}((z - u))) \cdot \operatorname{series\text{-}partial\text{-}sum}(v, k) \end{array}}}
\vnbprop{cc-partial-error-limit}{\vnbstmt{\begin{array}{@{}l@{}} \forall f, g, h \in (\mathbb{N} \to \mathbb{C}), w \in \mathbb{C}. \\ \quad (\operatorname{cc\text{-}series\text{-}converges}(f) \wedge (\operatorname{cc\text{-}series\text{-}converges}(g) \wedge \operatorname{cc\text{-}series\text{-}converges}(h))) \Rightarrow  \\ \quad \quad \operatorname{converges\text{-}to}(\operatorname{cc\text{-}ms}, \\ \quad \quad \quad (\lambda j \in \mathbb{N}.\; \\ \quad \quad \quad \quad ((\operatorname{cc\text{-}series\text{-}partial\text{-}sum}(f, (j + 1)) - \operatorname{cc\text{-}series\text{-}partial\text{-}sum}(g, (j + 1))) - (w \cdot \operatorname{cc\text{-}series\text{-}partial\text{-}sum}(h, j)))), \\ \quad \quad \quad ((\operatorname{cc\text{-}series\text{-}limit}(f) - \operatorname{cc\text{-}series\text{-}limit}(g)) - (w \cdot \operatorname{cc\text{-}series\text{-}limit}(h)))) \end{array}}}
\vnbprop{cps-diff-on-ball}{\vnbstmt{\begin{array}{@{}l@{}} \forall f, a \in (\mathbb{N} \to \mathbb{C}), r \in \mathbb{R}, \mathit{cpda} \in \operatorname{ball}(\operatorname{cc\text{-}ms}, 0, r). \\ \quad (\forall j \in \mathbb{N}.\; (a(j) = ((j + 1) \cdot f((j + 1)))) \wedge ((0 < r) \wedge (r < \operatorname{cps\text{-}radius}(f)))) \Rightarrow  \\ \quad \quad \operatorname{is\text{-}diff\text{-}on}(\operatorname{cc\text{-}normed\text{-}field}, \operatorname{ball}(\operatorname{cc\text{-}ms}, 0, r), \\ \quad \quad \quad (\lambda z \in \operatorname{ball}(\operatorname{cc\text{-}ms}, 0, r).\; \operatorname{cc\text{-}series\text{-}limit}((\lambda k \in \mathbb{N}.\; (f(k) \cdot \operatorname{power}(z, k))))), \\ \quad \quad \quad \mathit{cpda}, \\ \quad \quad \quad \operatorname{cc\text{-}series\text{-}limit}((\lambda k \in \mathbb{N}.\; (a(k) \cdot \operatorname{power}(\mathit{cpda}, k))))) \end{array}}}
\vnbprop{cps-holomorphic-ball}{\vnbstmt{\begin{array}{@{}l@{}} \forall f \in (\mathbb{N} \to \mathbb{C}), r \in \mathbb{R}. \\ \quad ((0 < r) \wedge (r < \operatorname{cps\text{-}radius}(f))) \Rightarrow  \\ \quad \quad \operatorname{holomorphic\text{-}on}(\operatorname{ball}(\operatorname{cc\text{-}ms}, 0, r), \\ \quad \quad \quad (\lambda z \in \operatorname{ball}(\operatorname{cc\text{-}ms}, 0, r).\; \operatorname{cc\text{-}series\text{-}limit}((\lambda k \in \mathbb{N}.\; (f(k) \cdot \operatorname{power}(z, k)))))) \wedge \\ \quad \quad \operatorname{cps\text{-}radius}((\lambda n \in \mathbb{N}.\; ((n + 1) \cdot f((n + 1))))) = \operatorname{cps\text{-}radius}(f) \wedge \\ \quad \quad \forall a \in \operatorname{ball}(\operatorname{cc\text{-}ms}, 0, r). \\ \quad \quad \quad \operatorname{is\text{-}diff\text{-}on}(\operatorname{cc\text{-}normed\text{-}field}, \operatorname{ball}(\operatorname{cc\text{-}ms}, 0, r), \\ \quad \quad \quad \quad (\lambda z \in \operatorname{ball}(\operatorname{cc\text{-}ms}, 0, r).\; \operatorname{cc\text{-}series\text{-}limit}((\lambda k \in \mathbb{N}.\; (f(k) \cdot \operatorname{power}(z, k))))), \\ \quad \quad \quad \quad a, \\ \quad \quad \quad \quad \operatorname{cc\text{-}series\text{-}limit}((\lambda k \in \mathbb{N}.\; (((k + 1) \cdot f((k + 1))) \cdot \operatorname{power}(a, k))))) \end{array}}}

\vnbbookfile{theorem-library/cps-taylor-coefficients}{theorem-library/cps-taylor-coefficients}
\noindent{\small cps-taylor-coefficients.scm -- COROLLARY 2.14: THE TAYLOR COEFFICIENTS OF A COMPLEX POWER SERIES.  f\textasciicircum{}(k)(0) = k! a\_k. Batch 27-A, 2026-09-24.}\par
\vnbprop{cps-kth-derived-step}{\vnbstmt{\begin{array}{@{}l@{}} \forall f \in (\mathbb{N} \to \mathbb{C}), k \in \mathbb{N}, g, h \in (\mathbb{N} \to \mathbb{C}), j \in \mathbb{N}. \\ \quad \forall n \in \mathbb{N}.\; (g(n) = ((\operatorname{factorial}((n + k)) \cdot \operatorname{recip}(\operatorname{factorial}(n))) \cdot f((n + k)))) \wedge \\ \quad \forall n \in \mathbb{N}. \\ \quad \quad h(n) = ((\operatorname{factorial}((n + (k + 1))) \cdot \operatorname{recip}(\operatorname{factorial}(n))) \cdot f((n + (k + 1)))) \Rightarrow \\ \quad \quad h(j) = ((j + 1) \cdot g((j + 1))) \end{array}}}
\vnbprop{cps-kth-derived-radius}{\vnbstmt{\begin{array}{@{}l@{}} \forall f \in (\mathbb{N} \to \mathbb{C}), k \in \mathbb{N}, g \in (\mathbb{N} \to \mathbb{C}). \\ \quad \forall n \in \mathbb{N}.\; (g(n) = ((\operatorname{factorial}((n + k)) \cdot \operatorname{recip}(\operatorname{factorial}(n))) \cdot f((n + k)))) \Rightarrow  \\ \quad \quad \operatorname{cps\text{-}radius}(g) = \operatorname{cps\text{-}radius}(f) \end{array}}}
\vnbprop{cps-sum-at-zero}{\vnbstmt{\begin{array}{@{}l@{}} \forall g \in (\mathbb{N} \to \mathbb{C}). \\ \quad \operatorname{cc\text{-}series\text{-}converges}((\lambda k \in \mathbb{N}.\; (g(k) \cdot \operatorname{power}(0, k)))) \wedge \\ \quad \operatorname{cc\text{-}series\text{-}limit}((\lambda k \in \mathbb{N}.\; (g(k) \cdot \operatorname{power}(0, k)))) = g(0) \end{array}}}
\vnbprop{cps-taylor-coefficients}{\vnbstmt{\begin{array}{@{}l@{}} \forall f \in (\mathbb{N} \to \mathbb{C}), k \in \mathbb{N}, g, h \in (\mathbb{N} \to \mathbb{C}), r \in \mathbb{R}. \\ \quad \forall n \in \mathbb{N}.\; (g(n) = ((\operatorname{factorial}((n + k)) \cdot \operatorname{recip}(\operatorname{factorial}(n))) \cdot f((n + k)))) \wedge \\ \quad \forall n \in \mathbb{N}. \\ \quad \quad h(n) = ((\operatorname{factorial}((n + (k + 1))) \cdot \operatorname{recip}(\operatorname{factorial}(n))) \cdot f((n + (k + 1)))) \wedge \\ \quad 0 < r \wedge \\ \quad r < \operatorname{cps\text{-}radius}(f) \Rightarrow \\ \quad \quad \operatorname{holomorphic\text{-}on}(\operatorname{ball}(\operatorname{cc\text{-}ms}, 0, r), \\ \quad \quad \quad (\lambda z \in \operatorname{ball}(\operatorname{cc\text{-}ms}, 0, r).\; \operatorname{cc\text{-}series\text{-}limit}((\lambda \mathit{cpsk} \in \mathbb{N}.\; (g(\mathit{cpsk}) \cdot \operatorname{power}(z, \mathit{cpsk})))))) \wedge \\ \quad \quad (\lambda z \in \operatorname{ball}(\operatorname{cc\text{-}ms}, 0, r).\; \operatorname{cc\text{-}series\text{-}limit}((\lambda \mathit{cpsk} \in \mathbb{N}.\; (g(\mathit{cpsk}) \cdot \operatorname{power}(z, \mathit{cpsk})))))(0) \\ \quad \quad =\; \operatorname{factorial}(k) \cdot f(k) \wedge \\ \quad \quad \forall a \in \operatorname{ball}(\operatorname{cc\text{-}ms}, 0, r). \\ \quad \quad \quad \operatorname{deriv\text{-}on}(\operatorname{cc\text{-}normed\text{-}field}, \operatorname{ball}(\operatorname{cc\text{-}ms}, 0, r), \\ \quad \quad \quad \quad (\lambda z \in \operatorname{ball}(\operatorname{cc\text{-}ms}, 0, r).\; \operatorname{cc\text{-}series\text{-}limit}((\lambda \mathit{cpsk} \in \mathbb{N}.\; (g(\mathit{cpsk}) \cdot \operatorname{power}(z, \mathit{cpsk}))))), \\ \quad \quad \quad \quad a) \\ \quad \quad \quad =\; (\lambda z \in \operatorname{ball}(\operatorname{cc\text{-}ms}, 0, r).\; \operatorname{cc\text{-}series\text{-}limit}((\lambda \mathit{cpsk} \in \mathbb{N}.\; (h(\mathit{cpsk}) \cdot \operatorname{power}(z, \mathit{cpsk})))))(a) \end{array}}}
\vnbprop{cps-coefficient-formula}{\vnbstmt{\begin{array}{@{}l@{}} \forall f \in (\mathbb{N} \to \mathbb{C}), k \in \mathbb{N}, g \in (\mathbb{N} \to \mathbb{C}), r \in \mathbb{R}. \\ \quad (\forall n \in \mathbb{N}.\; (g(n) = ((\operatorname{factorial}((n + k)) \cdot \operatorname{recip}(\operatorname{factorial}(n))) \cdot f((n + k)))) \wedge (0 < r)) \Rightarrow  \\ \quad \quad f(k) \\ \quad \quad =\; (\lambda z \in \operatorname{ball}(\operatorname{cc\text{-}ms}, 0, r).\; \operatorname{cc\text{-}series\text{-}limit}((\lambda \mathit{cpsk} \in \mathbb{N}.\; (g(\mathit{cpsk}) \cdot \operatorname{power}(z, \mathit{cpsk})))))(0) \cdot \operatorname{recip}(\operatorname{factorial}(k)) \end{array}}}

\vnbbookfile{theorem-library/nth-deriv-laws}{theorem-library/nth-deriv-laws}
\noindent{\small nth-deriv-laws.scm -- THE LAWS OF NTH-DERIV-ON, AND COROLLARY 2.14 RESTATED LITERALLY: f\textasciicircum{}(k)(0) = k! a\_k. Batch 36, 2026-09-26.}\par
\vnbprop{nth-deriv-on-succ-apply}{\vnbstmt{\begin{array}{@{}l@{}} \forall f, u, \mathit{ndlf}.; \forall j \in \mathbb{N}, z \in u. \\ \quad (\mathit{ndlf}^{((j + 1))}(z) \simeq \operatorname{deriv\text{-}on}(f, u, \mathit{ndlf}^{(j)}, z)) \end{array}}}
\vnbprop{nth-deriv-on-of-is-diff-on}{\vnbstmt{\begin{array}{@{}l@{}} \forall f, u, \mathit{ndlf}.; \forall j \in \mathbb{N}, g.; \forall w \in u. \\ \quad \operatorname{is\text{-}normed\text{-}field}(f) \wedge \\ \quad \forall e. \\ \quad \quad \operatorname{pos\text{-}rr}(e) \Rightarrow \exists h \in \operatorname{carr}(f).\; (\neg (h = \operatorname{zero}(f)) \wedge (\operatorname{fnrm}(f)(h) < e)) \wedge \\ \quad \forall z \in u.\; \operatorname{is\text{-}diff\text{-}on}(f, u, \mathit{ndlf}^{(j)}, z, g(z)) \Rightarrow \\ \quad \quad \mathit{ndlf}^{((j + 1))}(w) = g(w) \end{array}}}
\vnbprop{nth-deriv-on-succ-eq}{\vnbstmt{\begin{array}{@{}l@{}} \forall f, u, \mathit{ndlf}.; \forall j \in \mathbb{N}, g \in (u \to \operatorname{carr}(f)). \\ \quad \operatorname{is\text{-}normed\text{-}field}(f) \wedge \\ \quad \forall e. \\ \quad \quad \operatorname{pos\text{-}rr}(e) \Rightarrow \exists h \in \operatorname{carr}(f).\; (\neg (h = \operatorname{zero}(f)) \wedge (\operatorname{fnrm}(f)(h) < e)) \wedge \\ \quad \forall z \in u.\; \operatorname{is\text{-}diff\text{-}on}(f, u, \mathit{ndlf}^{(j)}, z, g(z)) \Rightarrow \\ \quad \quad \mathit{ndlf}^{((j + 1))} = g \end{array}}}
\vnbprop{nth-deriv-on-in-fun}{\vnbstmt{\begin{array}{@{}l@{}} \forall f, u.; \forall \mathit{ndlf} \in (u \to \operatorname{carr}(f)), j \in \mathbb{N}. \\ \quad \operatorname{is\text{-}normed\text{-}field}(f) \wedge \\ \quad \forall e. \\ \quad \quad \operatorname{pos\text{-}rr}(e) \Rightarrow \exists h \in \operatorname{carr}(f).\; (\neg (h = \operatorname{zero}(f)) \wedge (\operatorname{fnrm}(f)(h) < e)) \wedge \\ \quad \forall i \in \mathbb{N}, z \in u. \\ \quad \quad (i < j) \Rightarrow \exists l.\; \operatorname{is\text{-}diff\text{-}on}(f, u, \mathit{ndlf}^{(i)}, z, l) \Rightarrow \\ \quad \quad \mathit{ndlf}^{(j)} \in (u \to \operatorname{carr}(f)) \end{array}}}
\vnbprop{cps-nth-deriv}{\vnbstmt{\begin{array}{@{}l@{}} \forall f \in (\mathbb{N} \to \mathbb{C}), r \in \mathbb{R}, k \in \mathbb{N}, g \in (\mathbb{N} \to \mathbb{C}). \\ \quad 0 < r \wedge \\ \quad r < \operatorname{cps\text{-}radius}(f) \wedge \\ \quad \forall n \in \mathbb{N}.\; (g(n) = ((\operatorname{factorial}((n + k)) \cdot \operatorname{recip}(\operatorname{factorial}(n))) \cdot f((n + k)))) \Rightarrow \\ \quad \quad (\lambda z \in \operatorname{ball}(\operatorname{cc\text{-}ms}, 0, r).\; \operatorname{cc\text{-}series\text{-}limit}((\lambda \mathit{cpsk} \in \mathbb{N}.\; (f(\mathit{cpsk}) \cdot \operatorname{power}(z, \mathit{cpsk})))))^{(k)} \\ \quad \quad =\; (\lambda z \in \operatorname{ball}(\operatorname{cc\text{-}ms}, 0, r).\; \operatorname{cc\text{-}series\text{-}limit}((\lambda \mathit{cpsk} \in \mathbb{N}.\; (g(\mathit{cpsk}) \cdot \operatorname{power}(z, \mathit{cpsk}))))) \end{array}}}
\vnbprop{taylor-coefficients}{\vnbstmt{\begin{array}{@{}l@{}} \forall f \in (\mathbb{N} \to \mathbb{C}), r \in \mathbb{R}, k \in \mathbb{N}. \\ \quad ((0 < r) \wedge (r < \operatorname{cps\text{-}radius}(f))) \Rightarrow  \\ \quad \quad (\lambda z \in \operatorname{ball}(\operatorname{cc\text{-}ms}, 0, r).\; \operatorname{cc\text{-}series\text{-}limit}((\lambda \mathit{cpsk} \in \mathbb{N}.\; (f(\mathit{cpsk}) \cdot \operatorname{power}(z, \mathit{cpsk})))))^{(k)}(0) \\ \quad \quad =\; \operatorname{factorial}(k) \cdot f(k) \end{array}}}
\vnbprop{taylor-coefficient-formula}{\vnbstmt{\begin{array}{@{}l@{}} \forall f \in (\mathbb{N} \to \mathbb{C}), r \in \mathbb{R}, k \in \mathbb{N}. \\ \quad ((0 < r) \wedge (r < \operatorname{cps\text{-}radius}(f))) \Rightarrow  \\ \quad \quad f(k) \\ \quad \quad =\; (\lambda z \in \operatorname{ball}(\operatorname{cc\text{-}ms}, 0, r).\; \operatorname{cc\text{-}series\text{-}limit}((\lambda \mathit{cpsk} \in \mathbb{N}.\; (f(\mathit{cpsk}) \cdot \operatorname{power}(z, \mathit{cpsk})))))^{(k)}(0) \cdot \operatorname{recip}(\operatorname{factorial}(k)) \end{array}}}
\vnbprop{nth-deriv-on-holomorphic-ball}{\vnbstmt{\begin{array}{@{}l@{}} \forall f \in (\mathbb{N} \to \mathbb{C}), r \in \mathbb{R}, k \in \mathbb{N}. \\ \quad ((0 < r) \wedge (r < \operatorname{cps\text{-}radius}(f))) \Rightarrow  \\ \quad \quad \operatorname{holomorphic\text{-}on}(\operatorname{ball}(\operatorname{cc\text{-}ms}, 0, r), \\ \quad \quad \quad (\lambda z \in \operatorname{ball}(\operatorname{cc\text{-}ms}, 0, r).\; \operatorname{cc\text{-}series\text{-}limit}((\lambda \mathit{cpsk} \in \mathbb{N}.\; (f(\mathit{cpsk}) \cdot \operatorname{power}(z, \mathit{cpsk})))))^{(k)}) \end{array}}}

\vnbbookfile{theorem-library/line-int-fundamental}{theorem-library/line-int-fundamental}
\noindent{\small line-int-fundamental.scm -- THE FUNDAMENTAL THEOREM FOR INTEGRALS ALONG A ROAD (the user's notes, \textasciitilde{}/docs/complex-analysis.pdf, section 3.1, Proposition 3.2 and equation (50)), and the chain rule it rests on.}\par
\vnbprop{cc-compose-continuous-at}{\vnbstmt{\begin{array}{@{}l@{}} \forall s, u, g, i, x. \\ \quad \operatorname{is\text{-}metric\text{-}space}(s) \wedge \\ \quad u \subseteq \mathbb{C} \wedge \\ \quad \operatorname{is\text{-}continuous\text{-}at}(s, \operatorname{cc\text{-}ms}, g, x) \wedge \\ \quad \forall y \in \operatorname{pts}(s).\; (g(y) \in u) \wedge \\ \quad \operatorname{is\text{-}continuous\text{-}at}(\operatorname{subspace\text{-}ms}(\operatorname{nf\text{-}metric\text{-}space}(\operatorname{cc\text{-}normed\text{-}field}), u), \\ \quad \quad \operatorname{nf\text{-}metric\text{-}space}(\operatorname{cc\text{-}normed\text{-}field}), i, g(x)) \Rightarrow \\ \quad \quad \operatorname{is\text{-}continuous\text{-}at}(s, \operatorname{cc\text{-}ms}, (\lambda z \in \operatorname{pts}(s).\; i(g(z))), x) \end{array}}}
\vnbprop{cc-product-continuous-at}{\vnbstmt{\begin{array}{@{}l@{}} \forall s, p, q, x. \\ \quad (\operatorname{is\text{-}continuous\text{-}at}(s, \operatorname{cc\text{-}ms}, p, x) \wedge \operatorname{is\text{-}continuous\text{-}at}(s, \operatorname{cc\text{-}ms}, q, x)) \Rightarrow  \\ \quad \quad \operatorname{is\text{-}continuous\text{-}at}(s, \operatorname{cc\text{-}ms}, (\lambda z \in \operatorname{pts}(s).\; (p(z) \cdot q(z))), x) \end{array}}}
\vnbprop{lfn-cc-factor-identity}{\vnbstmt{\begin{array}{@{}l@{}} \forall a, b, c, e, h \in \mathbb{C}. \\ \quad (a \cdot ((b \cdot h) + (e \cdot (c \cdot h)))) = (h \cdot (a \cdot (b + (c \cdot e)))) \end{array}}}
\vnbprop{holomorphic-chain-within}{\vnbstmt{\begin{array}{@{}l@{}} \forall m, w, g, x, u, f, l, d. \\ \quad m \subseteq \mathbb{R} \wedge \\ \quad w \subseteq m \wedge \\ \quad g \in (m \to \mathbb{C}) \wedge \\ \quad \operatorname{is\text{-}continuous\text{-}at}(\operatorname{subspace\text{-}ms}(\operatorname{rr\text{-}ms}, m), \operatorname{cc\text{-}ms}, g, x) \wedge \\ \quad \forall y \in m.\; (g(y) \in u) \wedge \\ \quad \operatorname{is\text{-}diff\text{-}on}(\operatorname{cc\text{-}normed\text{-}field}, u, f, g(x), l) \wedge \\ \quad d \in \mathbb{C} \wedge \\ \quad \operatorname{has\text{-}deriv\text{-}at\text{-}within}((\lambda t \in m.\; \operatorname{real\text{-}part}(g(t))), w, x, \operatorname{real\text{-}part}(d)) \wedge \\ \quad \operatorname{has\text{-}deriv\text{-}at\text{-}within}((\lambda t \in m.\; \operatorname{imag\text{-}part}(g(t))), w, x, \operatorname{imag\text{-}part}(d)) \Rightarrow \\ \quad \quad \operatorname{has\text{-}deriv\text{-}at\text{-}within}((\lambda t \in m.\; \operatorname{real\text{-}part}(f(g(t)))), w, x, \operatorname{real\text{-}part}((l \cdot d))) \wedge \\ \quad \quad \operatorname{has\text{-}deriv\text{-}at\text{-}within}((\lambda t \in m.\; \operatorname{imag\text{-}part}(f(g(t)))), w, x, \operatorname{imag\text{-}part}((l \cdot d))) \end{array}}}
\vnbprop{ooint-center-nested}{\vnbstmt{\begin{array}{@{}l@{}} \forall x \in \mathbb{R}, v.; \forall r \in \mathbb{R}. \\ \quad (\operatorname{pos\text{-}rr}(v) \wedge (v \leq r)) \Rightarrow (\operatorname{ooint}((x - v), (x + v)) \subseteq \operatorname{ooint}((x - r), (x + r))) \end{array}}}
\vnbprop{holomorphic-chain}{\vnbstmt{\begin{array}{@{}l@{}} \forall m, g, x, u, f, l, d. \\ \quad m \subseteq \mathbb{R} \wedge \\ \quad g \in (m \to \mathbb{C}) \wedge \\ \quad \operatorname{is\text{-}continuous\text{-}at}(\operatorname{subspace\text{-}ms}(\operatorname{rr\text{-}ms}, m), \operatorname{cc\text{-}ms}, g, x) \wedge \\ \quad \forall y \in m.\; (g(y) \in u) \wedge \\ \quad \operatorname{is\text{-}diff\text{-}on}(\operatorname{cc\text{-}normed\text{-}field}, u, f, g(x), l) \wedge \\ \quad d \in \mathbb{C} \wedge \\ \quad \operatorname{has\text{-}deriv\text{-}at}((\lambda t \in m.\; \operatorname{real\text{-}part}(g(t))), x, \operatorname{real\text{-}part}(d)) \wedge \\ \quad \operatorname{has\text{-}deriv\text{-}at}((\lambda t \in m.\; \operatorname{imag\text{-}part}(g(t))), x, \operatorname{imag\text{-}part}(d)) \Rightarrow \\ \quad \quad \operatorname{has\text{-}deriv\text{-}at}((\lambda t \in m.\; \operatorname{real\text{-}part}(f(g(t)))), x, \operatorname{real\text{-}part}((l \cdot d))) \wedge \\ \quad \quad \operatorname{has\text{-}deriv\text{-}at}((\lambda t \in m.\; \operatorname{imag\text{-}part}(f(g(t)))), x, \operatorname{imag\text{-}part}((l \cdot d))) \end{array}}}
\vnbprop{line-int-of-derivative}{\vnbstmt{\begin{array}{@{}l@{}} \forall u, f, c, m, d, a, b. \\ \quad \operatorname{holomorphic\text{-}on}(u, f) \wedge \\ \quad c \in (u \to \mathbb{C}) \wedge \\ \quad \forall z \in u.\; \operatorname{is\text{-}diff\text{-}on}(\operatorname{cc\text{-}normed\text{-}field}, u, f, z, c(z)) \wedge \\ \quad \operatorname{is\text{-}road}(m, d, a, b) \wedge \\ \quad \operatorname{trace}(m, a, b) \subseteq u \Rightarrow \\ \quad \quad \operatorname{line\text{-}int}(c, m, d, a, b) = (f(m(b)) - f(m(a))) \end{array}}}

\vnbbookfile{theorem-library/goursat-bricks}{theorem-library/goursat-bricks}
\noindent{\small goursat-bricks.scm -- the definition-free bricks of Cauchy-Goursat (complex-analysis.pdf ch. 3, s. 3.2).  Batch 29-B, 2026-09-25.}\par
\vnbprop{quarter-step}{\vnbstmt{\begin{array}{@{}l@{}} \forall j, a, b, c, d \in \mathbb{R}. \\ \quad ((0 \leq a) \wedge ((0 \leq b) \wedge ((0 \leq c) \wedge ((0 \leq d) \wedge (j \leq (((a + b) + c) + d)))))) \Rightarrow  \\ \quad \quad (j \leq (4 \cdot a)) \vee ((j \leq (4 \cdot b)) \vee ((j \leq (4 \cdot c)) \vee (j \leq (4 \cdot d)))) \end{array}}}
\vnbprop{four-power-halving}{\vnbstmt{\forall n \in \mathbb{N}.\; ((\operatorname{power}(4, n) \cdot (\operatorname{recip}(\operatorname{power}(2, n)) \cdot \operatorname{recip}(\operatorname{power}(2, n)))) = 1)}}
\vnbprop{line-int-closed-road-of-derivative}{\vnbstmt{\begin{array}{@{}l@{}} \forall u, f, c, m, d, a, b. \\ \quad \operatorname{holomorphic\text{-}on}(u, f) \wedge \\ \quad c \in (u \to \mathbb{C}) \wedge \\ \quad \forall z \in u.\; \operatorname{is\text{-}diff\text{-}on}(\operatorname{cc\text{-}normed\text{-}field}, u, f, z, c(z)) \wedge \\ \quad \operatorname{is\text{-}road}(m, d, a, b) \wedge \\ \quad \operatorname{trace}(m, a, b) \subseteq u \wedge \\ \quad m(b) = m(a) \Rightarrow \\ \quad \quad \operatorname{line\text{-}int}(c, m, d, a, b) = 0 \end{array}}}
\vnbprop{affine-has-primitive-cc}{\vnbstmt{\begin{array}{@{}l@{}} \forall p, q, w \in \mathbb{C}, f \in (\mathbb{C} \to \mathbb{C}). \\ \quad \forall z \in \mathbb{C}.\; (f(z) \simeq ((p \cdot z) + ((q \cdot ((z - w) \cdot (z - w))) \cdot \operatorname{recip}(2)))) \Rightarrow  \\ \quad \quad (\operatorname{holomorphic\text{-}on}(\mathbb{C}, f) \wedge \forall z \in \mathbb{C}.\; \operatorname{is\text{-}diff\text{-}on}(\operatorname{cc\text{-}normed\text{-}field}, \mathbb{C}, f, z, (p + (q \cdot (z - w))))) \end{array}}}
\vnbprop{affine-closed-road-integral}{\vnbstmt{\begin{array}{@{}l@{}} \forall p, q, w \in \mathbb{C}, g \in (\mathbb{C} \to \mathbb{C}), m, c, a, b. \\ \quad (\forall z \in \mathbb{C}.\; (g(z) \simeq (p + (q \cdot (z - w)))) \wedge (\operatorname{is\text{-}road}(m, c, a, b) \wedge (m(b) = m(a)))) \Rightarrow  \\ \quad \quad \operatorname{line\text{-}int}(g, m, c, a, b) = 0 \end{array}}}
\vnbprop{caratheodory-remainder-small}{\vnbstmt{\begin{array}{@{}l@{}} \forall s, f, a, l, e. \\ \quad (\operatorname{is\text{-}diff\text{-}on}(\operatorname{cc\text{-}normed\text{-}field}, s, f, a, l) \wedge \operatorname{pos\text{-}rr}(e)) \Rightarrow  \\ \quad \quad \exists r. \\ \quad \quad \quad \operatorname{pos\text{-}rr}(r) \wedge \\ \quad \quad \quad \forall z \in s. \\ \quad \quad \quad \quad (\operatorname{magnitude}((z - a)) < r) \Rightarrow  \\ \quad \quad \quad \quad \quad \operatorname{magnitude}((f(z) - (f(a) + (l \cdot (z - a))))) \leq (e \cdot \operatorname{magnitude}((z - a))) \end{array}}}
\vnbprop{cc-telescoping-limit-bound}{\vnbstmt{\begin{array}{@{}l@{}} \forall s \in (\mathbb{N} \to \mathbb{C}), c \in \mathbb{R}. \\ \quad ((0 \leq c) \wedge \forall k \in \mathbb{N}.\; (\operatorname{magnitude}((s((k + 1)) - s(k))) \leq (c \cdot \operatorname{recip}(\operatorname{power}(2, k))))) \Rightarrow  \\ \quad \quad \exists v \in \mathbb{C}. \\ \quad \quad \quad (\operatorname{converges\text{-}to}(\operatorname{cc\text{-}ms}, s, v) \wedge \forall n \in \mathbb{N}.\; (\operatorname{magnitude}((s(n) - v)) \leq ((2 \cdot c) \cdot \operatorname{recip}(\operatorname{power}(2, n))))) \end{array}}}

\vnbbookfile{theorem-library/calint-poly-integral}{theorem-library/calint-poly-integral}
\noindent{\small calint-poly-integral.scm -- WAVE 2 (2026-10-27..30), wired 2026-10-30 from scratchpad/calint/calint-poly-integral.scm. Library: complex-analysis.}\par
\vnbprop{monomial-cc-diff}{\vnbstmt{\begin{array}{@{}l@{}} \forall j \in \mathbb{N}, w \in \mathbb{C}. \\ \quad \operatorname{is\text{-}diff\text{-}on}(\operatorname{cc\text{-}normed\text{-}field}, \mathbb{C}, (\lambda v \in \mathbb{C}.\; \operatorname{power}(v, (j + 1))), w, ((j + 1) \cdot \operatorname{power}(w, j))) \end{array}}}
\vnbprop{cc-ps-antiderivative}{\vnbstmt{\begin{array}{@{}l@{}} \forall n \in \mathbb{N}, c \in (\mathbb{N} \to \mathbb{C}). \exists f \in (\mathbb{C} \to \mathbb{C}). \forall w \in \mathbb{C}. \\ \quad \operatorname{is\text{-}diff\text{-}on}(\operatorname{cc\text{-}normed\text{-}field}, \mathbb{C}, f, w, \operatorname{cc\text{-}series\text{-}partial\text{-}sum}((\lambda j \in \mathbb{N}.\; (c(j) \cdot \operatorname{power}(w, j))), n)) \end{array}}}
\vnbprop{poly-fn-closed-road-integral}{\vnbstmt{\begin{array}{@{}l@{}} \forall f, m, c, a, b. \\ \quad (\operatorname{is\text{-}poly\text{-}fn\text{-}on}(\mathbb{C}, \mathbb{C}, f) \wedge (\operatorname{is\text{-}road}(m, c, a, b) \wedge (m(b) = m(a)))) \Rightarrow  \\ \quad \quad \operatorname{line\text{-}int}(f, m, c, a, b) = 0 \end{array}}}

\vnbbookfile{theorem-library/regulated-algebra}{theorem-library/regulated-algebra}
\noindent{\small regulated-algebra.scm -- THE ALGEBRA OF REGULATED FUNCTIONS ON [a, b], THE INTEGRAND OF A ROAD IS REGULATED, AND THE INTEGRAL ALONG A ROAD EXISTS. Batch 23-A, 2026-09-23.}\par
\vnbprop{right-limit-within-sum}{\vnbstmt{\begin{array}{@{}l@{}} \forall f, g, h, w, x, l, m. \\ \quad \operatorname{is\text{-}right\text{-}limit\text{-}within}(f, w, x, l) \wedge \\ \quad \operatorname{is\text{-}right\text{-}limit\text{-}within}(g, w, x, m) \wedge \\ \quad h \in (w \to \mathbb{R}) \wedge \\ \quad \forall z \in w.\; (h(z) \simeq (f(z) + g(z))) \Rightarrow \\ \quad \quad \operatorname{is\text{-}right\text{-}limit\text{-}within}(h, w, x, (l + m)) \end{array}}}
\vnbprop{left-limit-within-sum}{\vnbstmt{\begin{array}{@{}l@{}} \forall f, g, h, w, x, l, m. \\ \quad \operatorname{is\text{-}left\text{-}limit\text{-}within}(f, w, x, l) \wedge \\ \quad \operatorname{is\text{-}left\text{-}limit\text{-}within}(g, w, x, m) \wedge \\ \quad h \in (w \to \mathbb{R}) \wedge \\ \quad \forall z \in w.\; (h(z) \simeq (f(z) + g(z))) \Rightarrow \\ \quad \quad \operatorname{is\text{-}left\text{-}limit\text{-}within}(h, w, x, (l + m)) \end{array}}}
\vnbprop{right-limit-within-sub}{\vnbstmt{\begin{array}{@{}l@{}} \forall f, g, h, w, x, l, m. \\ \quad \operatorname{is\text{-}right\text{-}limit\text{-}within}(f, w, x, l) \wedge \\ \quad \operatorname{is\text{-}right\text{-}limit\text{-}within}(g, w, x, m) \wedge \\ \quad h \in (w \to \mathbb{R}) \wedge \\ \quad \forall z \in w.\; (h(z) \simeq (f(z) - g(z))) \Rightarrow \\ \quad \quad \operatorname{is\text{-}right\text{-}limit\text{-}within}(h, w, x, (l - m)) \end{array}}}
\vnbprop{left-limit-within-sub}{\vnbstmt{\begin{array}{@{}l@{}} \forall f, g, h, w, x, l, m. \\ \quad \operatorname{is\text{-}left\text{-}limit\text{-}within}(f, w, x, l) \wedge \\ \quad \operatorname{is\text{-}left\text{-}limit\text{-}within}(g, w, x, m) \wedge \\ \quad h \in (w \to \mathbb{R}) \wedge \\ \quad \forall z \in w.\; (h(z) \simeq (f(z) - g(z))) \Rightarrow \\ \quad \quad \operatorname{is\text{-}left\text{-}limit\text{-}within}(h, w, x, (l - m)) \end{array}}}
\vnbprop{right-limit-within-mul}{\vnbstmt{\begin{array}{@{}l@{}} \forall f, g, h, w, x, l, m. \\ \quad \operatorname{is\text{-}right\text{-}limit\text{-}within}(f, w, x, l) \wedge \\ \quad \operatorname{is\text{-}right\text{-}limit\text{-}within}(g, w, x, m) \wedge \\ \quad h \in (w \to \mathbb{R}) \wedge \\ \quad \forall z \in w.\; (h(z) \simeq (f(z) \cdot g(z))) \Rightarrow \\ \quad \quad \operatorname{is\text{-}right\text{-}limit\text{-}within}(h, w, x, (l \cdot m)) \end{array}}}
\vnbprop{left-limit-within-mul}{\vnbstmt{\begin{array}{@{}l@{}} \forall f, g, h, w, x, l, m. \\ \quad \operatorname{is\text{-}left\text{-}limit\text{-}within}(f, w, x, l) \wedge \\ \quad \operatorname{is\text{-}left\text{-}limit\text{-}within}(g, w, x, m) \wedge \\ \quad h \in (w \to \mathbb{R}) \wedge \\ \quad \forall z \in w.\; (h(z) \simeq (f(z) \cdot g(z))) \Rightarrow \\ \quad \quad \operatorname{is\text{-}left\text{-}limit\text{-}within}(h, w, x, (l \cdot m)) \end{array}}}
\vnbprop{right-limit-within-scalar}{\vnbstmt{\begin{array}{@{}l@{}} \forall f, h, w, x, l.; \forall c \in \mathbb{R}. \\ \quad (\operatorname{is\text{-}right\text{-}limit\text{-}within}(f, w, x, l) \wedge ((h \in (w \to \mathbb{R})) \wedge \forall z \in w.\; (h(z) \simeq (c \cdot f(z))))) \Rightarrow  \\ \quad \quad \operatorname{is\text{-}right\text{-}limit\text{-}within}(h, w, x, (c \cdot l)) \end{array}}}
\vnbprop{left-limit-within-scalar}{\vnbstmt{\begin{array}{@{}l@{}} \forall f, h, w, x, l.; \forall c \in \mathbb{R}. \\ \quad (\operatorname{is\text{-}left\text{-}limit\text{-}within}(f, w, x, l) \wedge ((h \in (w \to \mathbb{R})) \wedge \forall z \in w.\; (h(z) \simeq (c \cdot f(z))))) \Rightarrow  \\ \quad \quad \operatorname{is\text{-}left\text{-}limit\text{-}within}(h, w, x, (c \cdot l)) \end{array}}}
\vnbprop{regulated-on-sum}{\vnbstmt{\begin{array}{@{}l@{}} \forall f, g, h, a, b. \\ \quad \operatorname{is\text{-}regulated\text{-}on}(f, a, b) \wedge \\ \quad \operatorname{is\text{-}regulated\text{-}on}(g, a, b) \wedge \\ \quad h \in (\operatorname{ccint}(a, b) \to \mathbb{R}) \wedge \\ \quad \forall z \in \operatorname{ccint}(a, b).\; (h(z) \simeq (f(z) + g(z))) \Rightarrow \\ \quad \quad \operatorname{is\text{-}regulated\text{-}on}(h, a, b) \end{array}}}
\vnbprop{regulated-on-sub}{\vnbstmt{\begin{array}{@{}l@{}} \forall f, g, h, a, b. \\ \quad \operatorname{is\text{-}regulated\text{-}on}(f, a, b) \wedge \\ \quad \operatorname{is\text{-}regulated\text{-}on}(g, a, b) \wedge \\ \quad h \in (\operatorname{ccint}(a, b) \to \mathbb{R}) \wedge \\ \quad \forall z \in \operatorname{ccint}(a, b).\; (h(z) \simeq (f(z) - g(z))) \Rightarrow \\ \quad \quad \operatorname{is\text{-}regulated\text{-}on}(h, a, b) \end{array}}}
\vnbprop{regulated-on-mul}{\vnbstmt{\begin{array}{@{}l@{}} \forall f, g, h, a, b. \\ \quad \operatorname{is\text{-}regulated\text{-}on}(f, a, b) \wedge \\ \quad \operatorname{is\text{-}regulated\text{-}on}(g, a, b) \wedge \\ \quad h \in (\operatorname{ccint}(a, b) \to \mathbb{R}) \wedge \\ \quad \forall z \in \operatorname{ccint}(a, b).\; (h(z) \simeq (f(z) \cdot g(z))) \Rightarrow \\ \quad \quad \operatorname{is\text{-}regulated\text{-}on}(h, a, b) \end{array}}}
\vnbprop{regulated-on-scalar}{\vnbstmt{\begin{array}{@{}l@{}} \forall f, h, a, b.; \forall c \in \mathbb{R}. \\ \quad \operatorname{is\text{-}regulated\text{-}on}(f, a, b) \wedge \\ \quad h \in (\operatorname{ccint}(a, b) \to \mathbb{R}) \wedge \\ \quad \forall z \in \operatorname{ccint}(a, b).\; (h(z) \simeq (c \cdot f(z))) \Rightarrow \\ \quad \quad \operatorname{is\text{-}regulated\text{-}on}(h, a, b) \end{array}}}
\vnbprop{regulated-on-congruence}{\vnbstmt{\begin{array}{@{}l@{}} \forall f, g, a, b. \\ \quad (\operatorname{is\text{-}regulated\text{-}on}(g, a, b) \wedge ((f \in (\operatorname{ccint}(a, b) \to \mathbb{R})) \wedge \forall z \in \operatorname{ccint}(a, b).\; (f(z) \simeq g(z)))) \Rightarrow  \\ \quad \quad \operatorname{is\text{-}regulated\text{-}on}(f, a, b) \end{array}}}
\vnbprop{continuous-on-compose-path}{\vnbstmt{\begin{array}{@{}l@{}} \forall u, f, m, a, b. \\ \quad u \subseteq \mathbb{C} \wedge \\ \quad f \in (u \to \mathbb{C}) \wedge \\ \quad \forall y \in u. \\ \quad \quad \operatorname{is\text{-}continuous\text{-}at}(\operatorname{subspace\text{-}ms}(\operatorname{nf\text{-}metric\text{-}space}(\operatorname{cc\text{-}normed\text{-}field}), u), \\ \quad \quad \quad \operatorname{nf\text{-}metric\text{-}space}(\operatorname{cc\text{-}normed\text{-}field}), f, y) \wedge \\ \quad \operatorname{is\text{-}path}(m, a, b) \wedge \\ \quad \operatorname{trace}(m, a, b) \subseteq u \Rightarrow \\ \quad \quad \operatorname{is\text{-}continuous\text{-}on}((\lambda t \in \operatorname{ccint}(a, b).\; \operatorname{real\text{-}part}(f(m(t)))), \operatorname{ccint}(a, b)) \wedge \\ \quad \quad \operatorname{is\text{-}continuous\text{-}on}((\lambda t \in \operatorname{ccint}(a, b).\; \operatorname{imag\text{-}part}(f(m(t)))), \operatorname{ccint}(a, b)) \end{array}}}
\vnbprop{continuous-on-compose-regulated}{\vnbstmt{\begin{array}{@{}l@{}} \forall u, f, m, a, b. \\ \quad u \subseteq \mathbb{C} \wedge \\ \quad f \in (u \to \mathbb{C}) \wedge \\ \quad \forall y \in u. \\ \quad \quad \operatorname{is\text{-}continuous\text{-}at}(\operatorname{subspace\text{-}ms}(\operatorname{nf\text{-}metric\text{-}space}(\operatorname{cc\text{-}normed\text{-}field}), u), \\ \quad \quad \quad \operatorname{nf\text{-}metric\text{-}space}(\operatorname{cc\text{-}normed\text{-}field}), f, y) \wedge \\ \quad \operatorname{is\text{-}path}(m, a, b) \wedge \\ \quad \operatorname{trace}(m, a, b) \subseteq u \Rightarrow \\ \quad \quad \operatorname{is\text{-}regulated\text{-}on}((\lambda t \in \operatorname{ccint}(a, b).\; \operatorname{real\text{-}part}(f(m(t)))), a, b) \wedge \\ \quad \quad \operatorname{is\text{-}regulated\text{-}on}((\lambda t \in \operatorname{ccint}(a, b).\; \operatorname{imag\text{-}part}(f(m(t)))), a, b) \end{array}}}
\vnbprop{road-integrand-regulated}{\vnbstmt{\begin{array}{@{}l@{}} \forall u, f, m, c, a, b. \\ \quad u \subseteq \mathbb{C} \wedge \\ \quad f \in (u \to \mathbb{C}) \wedge \\ \quad \forall y \in u. \\ \quad \quad \operatorname{is\text{-}continuous\text{-}at}(\operatorname{subspace\text{-}ms}(\operatorname{nf\text{-}metric\text{-}space}(\operatorname{cc\text{-}normed\text{-}field}), u), \\ \quad \quad \quad \operatorname{nf\text{-}metric\text{-}space}(\operatorname{cc\text{-}normed\text{-}field}), f, y) \wedge \\ \quad \operatorname{is\text{-}road}(m, c, a, b) \wedge \\ \quad \operatorname{trace}(m, a, b) \subseteq u \Rightarrow \\ \quad \quad \operatorname{is\text{-}regulated\text{-}on}((\lambda t \in \operatorname{ccint}(a, b).\; \operatorname{real\text{-}part}((f(m(t)) \cdot c(t)))), a, b) \wedge \\ \quad \quad \operatorname{is\text{-}regulated\text{-}on}((\lambda t \in \operatorname{ccint}(a, b).\; \operatorname{imag\text{-}part}((f(m(t)) \cdot c(t)))), a, b) \end{array}}}
\vnbprop{line-int-exists}{\vnbstmt{\begin{array}{@{}l@{}} \forall u, f, m, c, a, b. \\ \quad u \subseteq \mathbb{C} \wedge \\ \quad f \in (u \to \mathbb{C}) \wedge \\ \quad \forall y \in u. \\ \quad \quad \operatorname{is\text{-}continuous\text{-}at}(\operatorname{subspace\text{-}ms}(\operatorname{nf\text{-}metric\text{-}space}(\operatorname{cc\text{-}normed\text{-}field}), u), \\ \quad \quad \quad \operatorname{nf\text{-}metric\text{-}space}(\operatorname{cc\text{-}normed\text{-}field}), f, y) \wedge \\ \quad \operatorname{is\text{-}road}(m, c, a, b) \wedge \\ \quad \operatorname{trace}(m, a, b) \subseteq u \Rightarrow \\ \quad \quad \exists \mathit{pwf}, w. \\ \quad \quad \quad \operatorname{is\text{-}primitive}(\mathit{pwf}, \\ \quad \quad \quad \quad (\lambda t \in \operatorname{ccint}(a, b).\; \operatorname{real\text{-}part}((\lambda t \in \operatorname{ccint}(a, b).\; (f(m(t)) \cdot c(t)))(t))), \\ \quad \quad \quad \quad a, b) \wedge \\ \quad \quad \quad \operatorname{is\text{-}primitive}(w, \\ \quad \quad \quad \quad (\lambda t \in \operatorname{ccint}(a, b).\; \operatorname{imag\text{-}part}((\lambda t \in \operatorname{ccint}(a, b).\; (f(m(t)) \cdot c(t)))(t))), \\ \quad \quad \quad \quad a, b) \end{array}}}
\vnbprop{line-int-in-cc}{\vnbstmt{\begin{array}{@{}l@{}} \forall u, f, m, c, a, b. \\ \quad u \subseteq \mathbb{C} \wedge \\ \quad f \in (u \to \mathbb{C}) \wedge \\ \quad \forall y \in u. \\ \quad \quad \operatorname{is\text{-}continuous\text{-}at}(\operatorname{subspace\text{-}ms}(\operatorname{nf\text{-}metric\text{-}space}(\operatorname{cc\text{-}normed\text{-}field}), u), \\ \quad \quad \quad \operatorname{nf\text{-}metric\text{-}space}(\operatorname{cc\text{-}normed\text{-}field}), f, y) \wedge \\ \quad \operatorname{is\text{-}road}(m, c, a, b) \wedge \\ \quad \operatorname{trace}(m, a, b) \subseteq u \Rightarrow \\ \quad \quad \operatorname{line\text{-}int}(f, m, c, a, b) \in \mathbb{C} \end{array}}}

\vnbbookfile{theorem-library/cc-int-estimate}{theorem-library/cc-int-estimate}
\noindent{\small cc-int-estimate.scm -- THE NOTES' ESTIMATE (45).}\par
\vnbprop{cc-re-inner-bound}{\vnbstmt{\begin{array}{@{}l@{}} \forall w, z \in \mathbb{C}. \\ \quad ((\operatorname{real\text{-}part}(w) \cdot \operatorname{real\text{-}part}(z)) + (\operatorname{imag\text{-}part}(w) \cdot \operatorname{imag\text{-}part}(z))) \leq (\operatorname{magnitude}(w) \cdot \operatorname{magnitude}(z)) \end{array}}}
\vnbprop{cc-int-abs-bound}{\vnbstmt{\begin{array}{@{}l@{}} \forall i, f, w, s, h, a, b. \\ \quad i \in (\operatorname{ccint}(a, b) \to \mathbb{C}) \wedge \\ \quad \operatorname{is\text{-}primitive}(f, (\lambda t \in \operatorname{ccint}(a, b).\; \operatorname{real\text{-}part}(i(t))), a, b) \wedge \\ \quad \operatorname{is\text{-}primitive}(w, (\lambda t \in \operatorname{ccint}(a, b).\; \operatorname{imag\text{-}part}(i(t))), a, b) \wedge \\ \quad \operatorname{is\text{-}primitive}(h, s, a, b) \wedge \\ \quad \forall y \in \operatorname{ccint}(a, b).\; (s(y) \simeq \operatorname{magnitude}(i(y))) \Rightarrow \\ \quad \quad \operatorname{magnitude}(\int_{a}^{b} i) \leq \int_{a}^{b} s \end{array}}}

\vnbbookfile{theorem-library/line-int-laws}{theorem-library/line-int-laws}
\noindent{\small line-int-laws.scm -- THE LAWS OF THE INTEGRAL ALONG A ROAD, for a road and a CONTINUOUS f, with no primitive in the statement. Batch 24-A, 2026-09-23.}\par
\vnbprop{line-int-value}{\vnbstmt{\begin{array}{@{}l@{}} \forall f, m, c, a, b, g_{1}, g_{2}. \\ \quad \operatorname{is\text{-}primitive}(g_{1}, (\lambda t \in \operatorname{ccint}(a, b).\; \operatorname{real\text{-}part}((f(m(t)) \cdot c(t)))), a, b) \wedge \\ \quad \operatorname{is\text{-}primitive}(g_{2}, (\lambda t \in \operatorname{ccint}(a, b).\; \operatorname{imag\text{-}part}((f(m(t)) \cdot c(t)))), a, b) \Rightarrow \\ \quad \quad \operatorname{line\text{-}int}(f, m, c, a, b) = ((\operatorname{g1}(b) - \operatorname{g1}(a)) + ((\operatorname{g2}(b) - \operatorname{g2}(a)) \cdot +i)) \end{array}}}
\vnbprop{line-int-sum}{\vnbstmt{\begin{array}{@{}l@{}} \forall u, f, g, m, c, a, b.; \forall h \in (u \to \mathbb{C}). \\ \quad u \subseteq \mathbb{C} \wedge \\ \quad f \in (u \to \mathbb{C}) \wedge \\ \quad \forall y \in u. \\ \quad \quad \operatorname{is\text{-}continuous\text{-}at}(\operatorname{subspace\text{-}ms}(\operatorname{nf\text{-}metric\text{-}space}(\operatorname{cc\text{-}normed\text{-}field}), u), \\ \quad \quad \quad \operatorname{nf\text{-}metric\text{-}space}(\operatorname{cc\text{-}normed\text{-}field}), f, y) \wedge \\ \quad g \in (u \to \mathbb{C}) \wedge \\ \quad \forall y \in u. \\ \quad \quad \operatorname{is\text{-}continuous\text{-}at}(\operatorname{subspace\text{-}ms}(\operatorname{nf\text{-}metric\text{-}space}(\operatorname{cc\text{-}normed\text{-}field}), u), \\ \quad \quad \quad \operatorname{nf\text{-}metric\text{-}space}(\operatorname{cc\text{-}normed\text{-}field}), g, y) \wedge \\ \quad \operatorname{is\text{-}road}(m, c, a, b) \wedge \\ \quad \operatorname{trace}(m, a, b) \subseteq u \wedge \\ \quad \forall z \in u.\; (h(z) \simeq (f(z) + g(z))) \Rightarrow \\ \quad \quad \operatorname{line\text{-}int}(h, m, c, a, b) = (\operatorname{line\text{-}int}(f, m, c, a, b) + \operatorname{line\text{-}int}(g, m, c, a, b)) \end{array}}}
\vnbprop{line-int-scalar}{\vnbstmt{\begin{array}{@{}l@{}} \forall u, f, m, c, a, b.; \forall \mathit{lic} \in \mathbb{C}, h \in (u \to \mathbb{C}). \\ \quad u \subseteq \mathbb{C} \wedge \\ \quad f \in (u \to \mathbb{C}) \wedge \\ \quad \forall y \in u. \\ \quad \quad \operatorname{is\text{-}continuous\text{-}at}(\operatorname{subspace\text{-}ms}(\operatorname{nf\text{-}metric\text{-}space}(\operatorname{cc\text{-}normed\text{-}field}), u), \\ \quad \quad \quad \operatorname{nf\text{-}metric\text{-}space}(\operatorname{cc\text{-}normed\text{-}field}), f, y) \wedge \\ \quad \operatorname{is\text{-}road}(m, c, a, b) \wedge \\ \quad \operatorname{trace}(m, a, b) \subseteq u \wedge \\ \quad \forall z \in u.\; (h(z) \simeq (\mathit{lic} \cdot f(z))) \Rightarrow \\ \quad \quad \operatorname{line\text{-}int}(h, m, c, a, b) = (\mathit{lic} \cdot \operatorname{line\text{-}int}(f, m, c, a, b)) \end{array}}}
\vnbprop{line-int-opposite-road}{\vnbstmt{\begin{array}{@{}l@{}} \forall u, f, m, c, a, b, g, d. \\ \quad u \subseteq \mathbb{C} \wedge \\ \quad f \in (u \to \mathbb{C}) \wedge \\ \quad \forall y \in u. \\ \quad \quad \operatorname{is\text{-}continuous\text{-}at}(\operatorname{subspace\text{-}ms}(\operatorname{nf\text{-}metric\text{-}space}(\operatorname{cc\text{-}normed\text{-}field}), u), \\ \quad \quad \quad \operatorname{nf\text{-}metric\text{-}space}(\operatorname{cc\text{-}normed\text{-}field}), f, y) \wedge \\ \quad \operatorname{is\text{-}road}(m, c, a, b) \wedge \\ \quad \operatorname{trace}(m, a, b) \subseteq u \wedge \\ \quad g \in (\operatorname{ccint}(a, b) \to \mathbb{C}) \wedge \\ \quad d \in (\operatorname{ccint}(a, b) \to \mathbb{C}) \wedge \\ \quad \forall \mathit{pay} \in \operatorname{ccint}(a, b).\; (g(\mathit{pay}) \simeq m(((a + b) - \mathit{pay}))) \wedge \\ \quad \forall \mathit{pay} \in \operatorname{ccint}(a, b).\; (d(\mathit{pay}) \simeq (-1 \cdot c(((a + b) - \mathit{pay})))) \Rightarrow \\ \quad \quad \operatorname{line\text{-}int}(f, g, d, a, b) = (-1 \cdot \operatorname{line\text{-}int}(f, m, c, a, b)) \end{array}}}
\vnbprop{road-restrict}{\vnbstmt{\begin{array}{@{}l@{}} \forall m, c, a, b.; \forall v, \mathit{lidv} \in \mathbb{R}. \\ \quad (\operatorname{is\text{-}road}(m, c, a, b) \wedge ((a \leq v) \wedge ((v < \mathit{lidv}) \wedge (\mathit{lidv} \leq b)))) \Rightarrow  \\ \quad \quad \operatorname{is\text{-}road}(\operatorname{restrict}(m, \operatorname{ccint}(v, \mathit{lidv})), \operatorname{restrict}(c, \operatorname{ccint}(v, \mathit{lidv})), v, \mathit{lidv}) \end{array}}}
\vnbprop{cc-int-congruence}{\vnbstmt{\begin{array}{@{}l@{}} \forall i, \mathit{ppsi}, a, b, g_{1}, g_{2}. \\ \quad \operatorname{is\text{-}primitive}(g_{1}, (\lambda t \in \operatorname{ccint}(a, b).\; \operatorname{real\text{-}part}(i(t))), a, b) \wedge \\ \quad \operatorname{is\text{-}primitive}(g_{2}, (\lambda t \in \operatorname{ccint}(a, b).\; \operatorname{imag\text{-}part}(i(t))), a, b) \wedge \\ \quad \mathit{ppsi} \in (\operatorname{ccint}(a, b) \to \mathbb{C}) \wedge \\ \quad \forall y \in \operatorname{ccint}(a, b).\; (i(y) \simeq \operatorname{ppsi}(y)) \Rightarrow \\ \quad \quad \int_{a}^{b} i = \int_{a}^{b} \mathit{ppsi} \end{array}}}
\vnbprop{line-int-adjacent}{\vnbstmt{\begin{array}{@{}l@{}} \forall u, f, m, c, a, b.; \forall v \in \operatorname{ooint}(a, b). \\ \quad u \subseteq \mathbb{C} \wedge \\ \quad f \in (u \to \mathbb{C}) \wedge \\ \quad \forall y \in u. \\ \quad \quad \operatorname{is\text{-}continuous\text{-}at}(\operatorname{subspace\text{-}ms}(\operatorname{nf\text{-}metric\text{-}space}(\operatorname{cc\text{-}normed\text{-}field}), u), \\ \quad \quad \quad \operatorname{nf\text{-}metric\text{-}space}(\operatorname{cc\text{-}normed\text{-}field}), f, y) \wedge \\ \quad \operatorname{is\text{-}road}(m, c, a, b) \wedge \\ \quad \operatorname{trace}(m, a, b) \subseteq u \Rightarrow \\ \quad \quad \operatorname{line\text{-}int}(f, m, c, a, b) \\ \quad \quad =\; \operatorname{line\text{-}int}(f, \operatorname{restrict}(m, \operatorname{ccint}(a, v)), \operatorname{restrict}(c, \operatorname{ccint}(a, v)), a, v) + \operatorname{line\text{-}int}(f, \operatorname{restrict}(m, \operatorname{ccint}(v, b)), \operatorname{restrict}(c, \operatorname{ccint}(v, b)), v, b) \end{array}}}
\vnbprop{right-limit-within-magnitude}{\vnbstmt{\begin{array}{@{}l@{}} \forall f, g, h, w, x, l, k. \\ \quad \operatorname{is\text{-}right\text{-}limit\text{-}within}(f, w, x, l) \wedge \\ \quad \operatorname{is\text{-}right\text{-}limit\text{-}within}(g, w, x, k) \wedge \\ \quad h \in (w \to \mathbb{R}) \wedge \\ \quad \forall z \in w.\; (h(z) \simeq \operatorname{magnitude}((f(z) + (g(z) \cdot +i)))) \Rightarrow \\ \quad \quad \operatorname{is\text{-}right\text{-}limit\text{-}within}(h, w, x, \operatorname{magnitude}((l + (k \cdot +i)))) \end{array}}}
\vnbprop{left-limit-within-magnitude}{\vnbstmt{\begin{array}{@{}l@{}} \forall f, g, h, w, x, l, k. \\ \quad \operatorname{is\text{-}left\text{-}limit\text{-}within}(f, w, x, l) \wedge \\ \quad \operatorname{is\text{-}left\text{-}limit\text{-}within}(g, w, x, k) \wedge \\ \quad h \in (w \to \mathbb{R}) \wedge \\ \quad \forall z \in w.\; (h(z) \simeq \operatorname{magnitude}((f(z) + (g(z) \cdot +i)))) \Rightarrow \\ \quad \quad \operatorname{is\text{-}left\text{-}limit\text{-}within}(h, w, x, \operatorname{magnitude}((l + (k \cdot +i)))) \end{array}}}
\vnbprop{regulated-on-magnitude}{\vnbstmt{\begin{array}{@{}l@{}} \forall f, g, h, a, b. \\ \quad \operatorname{is\text{-}regulated\text{-}on}(f, a, b) \wedge \\ \quad \operatorname{is\text{-}regulated\text{-}on}(g, a, b) \wedge \\ \quad h \in (\operatorname{ccint}(a, b) \to \mathbb{R}) \wedge \\ \quad \forall z \in \operatorname{ccint}(a, b).\; (h(z) \simeq \operatorname{magnitude}((f(z) + (g(z) \cdot +i)))) \Rightarrow \\ \quad \quad \operatorname{is\text{-}regulated\text{-}on}(h, a, b) \end{array}}}
\vnbprop{line-int-abs-bound}{\vnbstmt{\begin{array}{@{}l@{}} \forall u, f, m, c, a, b.; \forall \mathit{liab} \in (\operatorname{ccint}(a, b) \to \mathbb{R}). \\ \quad u \subseteq \mathbb{C} \wedge \\ \quad f \in (u \to \mathbb{C}) \wedge \\ \quad \forall y \in u. \\ \quad \quad \operatorname{is\text{-}continuous\text{-}at}(\operatorname{subspace\text{-}ms}(\operatorname{nf\text{-}metric\text{-}space}(\operatorname{cc\text{-}normed\text{-}field}), u), \\ \quad \quad \quad \operatorname{nf\text{-}metric\text{-}space}(\operatorname{cc\text{-}normed\text{-}field}), f, y) \wedge \\ \quad \operatorname{is\text{-}road}(m, c, a, b) \wedge \\ \quad \operatorname{trace}(m, a, b) \subseteq u \wedge \\ \quad \forall \mathit{liy} \in \operatorname{ccint}(a, b). \\ \quad \quad (\operatorname{liab}(\mathit{liy}) \simeq (\operatorname{magnitude}(f(m(\mathit{liy}))) \cdot \operatorname{magnitude}(c(\mathit{liy})))) \Rightarrow \\ \quad \quad \operatorname{magnitude}(\operatorname{line\text{-}int}(f, m, c, a, b)) \leq \int_{a}^{b} \mathit{liab} \end{array}}}

\vnbbookfile{theorem-library/line-int-reparam}{theorem-library/line-int-reparam}
\noindent{\small line-int-reparam.scm -- THE AFFINE CHANGE OF PARAMETER (the notes' (48), affine case) AND THE SEGMENT LEMMAS OF GOURSAT'S DEFINITION-FREE HALF. Batch 29-A, 2026-09-25.}\par
\vnbprop{affine-reparam-ccint-in}{\vnbstmt{\begin{array}{@{}l@{}} \forall a, b, c, d, l, m \in \mathbb{R}, z \in \operatorname{ccint}(c, d). \\ \quad ((0 < l) \wedge (((m + (l \cdot c)) = a) \wedge ((m + (l \cdot d)) = b))) \Rightarrow ((m + (l \cdot z)) \in \operatorname{ccint}(a, b)) \end{array}}}
\vnbprop{affine-reparam-ooint-in}{\vnbstmt{\begin{array}{@{}l@{}} \forall a, b, c, d, l, m \in \mathbb{R}, z \in \operatorname{ooint}(c, d). \\ \quad ((0 < l) \wedge (((m + (l \cdot c)) = a) \wedge ((m + (l \cdot d)) = b))) \Rightarrow ((m + (l \cdot z)) \in \operatorname{ooint}(a, b)) \end{array}}}
\vnbprop{affine-reparam-solve}{\vnbstmt{\begin{array}{@{}l@{}} \forall a, b, c, d, l, m \in \mathbb{R}, z \in \operatorname{ccint}(a, b), q \in \mathbb{R}. \\ \quad ((0 < l) \wedge (((m + (l \cdot c)) = a) \wedge (((m + (l \cdot d)) = b) \wedge ((z - m) = (q \cdot l))))) \Rightarrow  \\ \quad \quad (q \in \operatorname{ccint}(c, d)) \wedge ((m + (l \cdot q)) = z) \end{array}}}
\vnbprop{affine-reparam-inverse}{\vnbstmt{\begin{array}{@{}l@{}} \forall a, b, c, d, l, m \in \mathbb{R}, z \in \operatorname{ccint}(a, b). \\ \quad ((0 < l) \wedge (((m + (l \cdot c)) = a) \wedge ((m + (l \cdot d)) = b))) \Rightarrow  \\ \quad \quad (((z - m) \cdot \operatorname{recip}(l)) \in \operatorname{ccint}(c, d)) \wedge ((m + (l \cdot ((z - m) \cdot \operatorname{recip}(l)))) = z) \end{array}}}
\vnbprop{affine-reparam-inverse-value}{\vnbstmt{\begin{array}{@{}l@{}} \forall l, m, z \in \mathbb{R}. \\ \quad (0 < l) \Rightarrow ((((m + (l \cdot z)) - m) \cdot \operatorname{recip}(l)) = z) \end{array}}}
\vnbprop{affine-reparam-continuous-at-sub}{\vnbstmt{\begin{array}{@{}l@{}} \forall a, b, c, d, l, m \in \mathbb{R}, f.; \forall g \in (\operatorname{ccint}(c, d) \to \mathbb{R}), t \in \operatorname{ccint}(c, d). \\ \quad 0 < l \wedge \\ \quad (m + (l \cdot c)) = a \wedge \\ \quad (m + (l \cdot d)) = b \wedge \\ \quad \forall z \in \operatorname{ccint}(c, d).\; (g(z) \simeq f((m + (l \cdot z)))) \wedge \\ \quad \operatorname{is\text{-}continuous\text{-}at}(\operatorname{subspace\text{-}ms}(\operatorname{rr\text{-}ms}, \operatorname{ccint}(a, b)), \operatorname{rr\text{-}ms}, f, (m + (l \cdot t))) \Rightarrow \\ \quad \quad \operatorname{is\text{-}continuous\text{-}at}(\operatorname{subspace\text{-}ms}(\operatorname{rr\text{-}ms}, \operatorname{ccint}(c, d)), \operatorname{rr\text{-}ms}, g, t) \end{array}}}
\vnbprop{affine-reparam-continuous-on}{\vnbstmt{\begin{array}{@{}l@{}} \forall a, b, c, d, l, m \in \mathbb{R}, f.; \forall g \in (\operatorname{ccint}(c, d) \to \mathbb{R}). \\ \quad 0 < l \wedge \\ \quad (m + (l \cdot c)) = a \wedge \\ \quad (m + (l \cdot d)) = b \wedge \\ \quad \operatorname{is\text{-}continuous\text{-}on}(f, \operatorname{ccint}(a, b)) \wedge \\ \quad \forall z \in \operatorname{ccint}(c, d).\; (g(z) \simeq f((m + (l \cdot z)))) \Rightarrow \\ \quad \quad \operatorname{is\text{-}continuous\text{-}on}(g, \operatorname{ccint}(c, d)) \end{array}}}
\vnbprop{affine-reparam-image-countable}{\vnbstmt{\begin{array}{@{}l@{}} \forall a, b, c, d, l, m \in \mathbb{R}, e. \\ \quad ((0 < l) \wedge (((m + (l \cdot c)) = a) \wedge (((m + (l \cdot d)) = b) \wedge (\operatorname{is\text{-}countable}(e) \wedge (e \subseteq \operatorname{ccint}(a, b)))))) \Rightarrow  \\ \quad \quad \operatorname{is\text{-}countable}(\operatorname{image}((\lambda v \in \operatorname{ccint}(a, b).\; ((v - m) \cdot \operatorname{recip}(l))), e)) \wedge \\ \quad \quad \operatorname{image}((\lambda v \in \operatorname{ccint}(a, b).\; ((v - m) \cdot \operatorname{recip}(l))), e) \subseteq \operatorname{ccint}(c, d) \end{array}}}
\vnbprop{affine-reparam-not-in-image}{\vnbstmt{\begin{array}{@{}l@{}} \forall a, b, c, d, l, m \in \mathbb{R}, e.; \forall t \in \operatorname{ccint}(c, d). \\ \quad 0 < l \wedge \\ \quad (m + (l \cdot c)) = a \wedge \\ \quad (m + (l \cdot d)) = b \wedge \\ \quad e \subseteq \operatorname{ccint}(a, b) \wedge \\ \quad \neg (t \in \operatorname{image}((\lambda v \in \operatorname{ccint}(a, b).\; ((v - m) \cdot \operatorname{recip}(l))), e)) \Rightarrow \\ \quad \quad \neg ((m + (l \cdot t)) \in e) \end{array}}}
\vnbprop{affine-reparam-regulated}{\vnbstmt{\begin{array}{@{}l@{}} \forall a, b, c, d, l, m \in \mathbb{R}, f.; \forall g \in (\operatorname{ccint}(c, d) \to \mathbb{R}). \\ \quad 0 < l \wedge \\ \quad (m + (l \cdot c)) = a \wedge \\ \quad (m + (l \cdot d)) = b \wedge \\ \quad c < d \wedge \\ \quad \operatorname{is\text{-}regulated\text{-}on}(f, a, b) \wedge \\ \quad \forall z \in \operatorname{ccint}(c, d).\; (g(z) \simeq f((m + (l \cdot z)))) \Rightarrow \\ \quad \quad \operatorname{is\text{-}regulated\text{-}on}(g, c, d) \end{array}}}
\vnbprop{pw-int-affine-subst}{\vnbstmt{\begin{array}{@{}l@{}} \forall a, b, c, d, l, m \in \mathbb{R}, f, i.; \forall h, s \in (\operatorname{ccint}(c, d) \to \mathbb{R}). \\ \quad 0 < l \wedge \\ \quad (m + (l \cdot c)) = a \wedge \\ \quad (m + (l \cdot d)) = b \wedge \\ \quad c < d \wedge \\ \quad \operatorname{is\text{-}primitive}(f, i, a, b) \wedge \\ \quad \forall z \in \operatorname{ccint}(c, d).\; (h(z) \simeq f((m + (l \cdot z)))) \wedge \\ \quad \forall z \in \operatorname{ccint}(c, d).\; (s(z) \simeq (l \cdot i((m + (l \cdot z))))) \Rightarrow \\ \quad \quad \operatorname{is\text{-}primitive}(h, s, c, d) \wedge (\int_{c}^{d} s = \int_{a}^{b} i) \end{array}}}
\vnbprop{cc-int-affine-subst}{\vnbstmt{\begin{array}{@{}l@{}} \forall a, b, c, d, l, m \in \mathbb{R}, i, \mathit{lrpb}, f, w. \\ \quad 0 < l \wedge \\ \quad (m + (l \cdot c)) = a \wedge \\ \quad (m + (l \cdot d)) = b \wedge \\ \quad c < d \wedge \\ \quad i \in (\operatorname{ccint}(a, b) \to \mathbb{C}) \wedge \\ \quad \mathit{lrpb} \in (\operatorname{ccint}(c, d) \to \mathbb{C}) \wedge \\ \quad \forall z \in \operatorname{ccint}(c, d).\; (\operatorname{lrpb}(z) \simeq (l \cdot i((m + (l \cdot z))))) \wedge \\ \quad \operatorname{is\text{-}primitive}(f, (\lambda t \in \operatorname{ccint}(a, b).\; \operatorname{real\text{-}part}(i(t))), a, b) \wedge \\ \quad \operatorname{is\text{-}primitive}(w, (\lambda t \in \operatorname{ccint}(a, b).\; \operatorname{imag\text{-}part}(i(t))), a, b) \Rightarrow \\ \quad \quad \int_{c}^{d} \mathit{lrpb} = \int_{a}^{b} i \end{array}}}
\vnbprop{trace-affine-reparam}{\vnbstmt{\begin{array}{@{}l@{}} \forall a, b, c, d, l, m \in \mathbb{R}, f, r. \\ \quad 0 < l \wedge \\ \quad (m + (l \cdot c)) = a \wedge \\ \quad (m + (l \cdot d)) = b \wedge \\ \quad f \in (\operatorname{ccint}(a, b) \to \mathbb{C}) \wedge \\ \quad r \in (\operatorname{ccint}(c, d) \to \mathbb{C}) \wedge \\ \quad \forall z \in \operatorname{ccint}(c, d).\; (r(z) \simeq f((m + (l \cdot z)))) \Rightarrow \\ \quad \quad (\operatorname{trace}(r, c, d) \subseteq \operatorname{trace}(f, a, b)) \wedge (\operatorname{trace}(f, a, b) \subseteq \operatorname{trace}(r, c, d)) \end{array}}}
\vnbprop{road-affine-reparam}{\vnbstmt{\begin{array}{@{}l@{}} \forall a, b, c, d, l, m \in \mathbb{R}, f, g, r, k. \\ \quad 0 < l \wedge \\ \quad (m + (l \cdot c)) = a \wedge \\ \quad (m + (l \cdot d)) = b \wedge \\ \quad c < d \wedge \\ \quad \operatorname{is\text{-}road}(f, g, a, b) \wedge \\ \quad r \in (\operatorname{ccint}(c, d) \to \mathbb{C}) \wedge \\ \quad k \in (\operatorname{ccint}(c, d) \to \mathbb{C}) \wedge \\ \quad \forall z \in \operatorname{ccint}(c, d).\; (r(z) \simeq f((m + (l \cdot z)))) \wedge \\ \quad \forall z \in \operatorname{ccint}(c, d).\; (k(z) \simeq (l \cdot g((m + (l \cdot z))))) \Rightarrow \\ \quad \quad \operatorname{is\text{-}road}(r, k, c, d) \end{array}}}
\vnbprop{line-int-affine-reparam}{\vnbstmt{\begin{array}{@{}l@{}} \forall u, f, m, c, a, b.; \forall \mathit{lrpc}, d, l, \mathit{lrpm} \in \mathbb{R}, r, k. \\ \quad u \subseteq \mathbb{C} \wedge \\ \quad f \in (u \to \mathbb{C}) \wedge \\ \quad \forall y \in u. \\ \quad \quad \operatorname{is\text{-}continuous\text{-}at}(\operatorname{subspace\text{-}ms}(\operatorname{nf\text{-}metric\text{-}space}(\operatorname{cc\text{-}normed\text{-}field}), u), \\ \quad \quad \quad \operatorname{nf\text{-}metric\text{-}space}(\operatorname{cc\text{-}normed\text{-}field}), f, y) \wedge \\ \quad \operatorname{is\text{-}road}(m, c, a, b) \wedge \\ \quad \operatorname{trace}(m, a, b) \subseteq u \wedge \\ \quad 0 < l \wedge \\ \quad (\mathit{lrpm} + (l \cdot \mathit{lrpc})) = a \wedge \\ \quad (\mathit{lrpm} + (l \cdot d)) = b \wedge \\ \quad \mathit{lrpc} < d \wedge \\ \quad r \in (\operatorname{ccint}(\mathit{lrpc}, d) \to \mathbb{C}) \wedge \\ \quad k \in (\operatorname{ccint}(\mathit{lrpc}, d) \to \mathbb{C}) \wedge \\ \quad \forall z \in \operatorname{ccint}(\mathit{lrpc}, d).\; (r(z) \simeq m((\mathit{lrpm} + (l \cdot z)))) \wedge \\ \quad \forall z \in \operatorname{ccint}(\mathit{lrpc}, d).\; (k(z) \simeq (l \cdot c((\mathit{lrpm} + (l \cdot z))))) \Rightarrow \\ \quad \quad \operatorname{line\text{-}int}(f, r, k, \mathit{lrpc}, d) = \operatorname{line\text{-}int}(f, m, c, a, b) \end{array}}}
\vnbprop{trace-restrict-subset}{\vnbstmt{\begin{array}{@{}l@{}} \forall m, a, b.; \forall c, d \in \mathbb{R}. \\ \quad ((m \in (\operatorname{ccint}(a, b) \to \mathbb{C})) \wedge ((a \in \mathbb{R}) \wedge ((b \in \mathbb{R}) \wedge ((a \leq c) \wedge (d \leq b))))) \Rightarrow  \\ \quad \quad \operatorname{trace}(\operatorname{restrict}(m, \operatorname{ccint}(c, d)), c, d) \subseteq \operatorname{trace}(m, a, b) \end{array}}}
\vnbprop{segment-int-split}{\vnbstmt{\begin{array}{@{}l@{}} \forall u, f.; \forall z_{0}, w \in \mathbb{C}, t \in \operatorname{ooint}(0, 1). \\ \quad u \subseteq \mathbb{C} \wedge \\ \quad f \in (u \to \mathbb{C}) \wedge \\ \quad \forall y \in u. \\ \quad \quad \operatorname{is\text{-}continuous\text{-}at}(\operatorname{subspace\text{-}ms}(\operatorname{nf\text{-}metric\text{-}space}(\operatorname{cc\text{-}normed\text{-}field}), u), \\ \quad \quad \quad \operatorname{nf\text{-}metric\text{-}space}(\operatorname{cc\text{-}normed\text{-}field}), f, y) \wedge \\ \quad \operatorname{trace}((\lambda x \in \operatorname{ccint}(0, 1).\; (z_{0} + (x \cdot w))), 0, 1) \subseteq u \Rightarrow \\ \quad \quad \operatorname{line\text{-}int}(f, (\lambda x \in \operatorname{ccint}(0, 1).\; (z_{0} + (x \cdot w))), (\lambda x \in \operatorname{ccint}(0, 1).\; w), 0, 1) \\ \quad \quad =\; \operatorname{line\text{-}int}(f, (\lambda x \in \operatorname{ccint}(0, 1).\; (z_{0} + (x \cdot (t \cdot w)))), (\lambda x \in \operatorname{ccint}(0, 1).\; (t \cdot w)), 0, 1) + \operatorname{line\text{-}int}(f, (\lambda x \in \operatorname{ccint}(0, 1).\; ((z_{0} + (t \cdot w)) + (x \cdot ((1 - t) \cdot w)))), (\lambda x \in \operatorname{ccint}(0, 1).\; ((1 - t) \cdot w)), 0, 1) \end{array}}}
\vnbprop{segment-int-abs-bound}{\vnbstmt{\begin{array}{@{}l@{}} \forall u, f.; \forall z_{0}, w \in \mathbb{C}, n \in \mathbb{R}. \\ \quad u \subseteq \mathbb{C} \wedge \\ \quad f \in (u \to \mathbb{C}) \wedge \\ \quad \forall y \in u. \\ \quad \quad \operatorname{is\text{-}continuous\text{-}at}(\operatorname{subspace\text{-}ms}(\operatorname{nf\text{-}metric\text{-}space}(\operatorname{cc\text{-}normed\text{-}field}), u), \\ \quad \quad \quad \operatorname{nf\text{-}metric\text{-}space}(\operatorname{cc\text{-}normed\text{-}field}), f, y) \wedge \\ \quad \operatorname{trace}((\lambda x \in \operatorname{ccint}(0, 1).\; (z_{0} + (x \cdot w))), 0, 1) \subseteq u \wedge \\ \quad \forall \mathit{lrpy} \in \operatorname{trace}((\lambda x \in \operatorname{ccint}(0, 1).\; (z_{0} + (x \cdot w))), 0, 1). \\ \quad \quad \operatorname{magnitude}(f(\mathit{lrpy})) \leq n \Rightarrow \\ \quad \quad \operatorname{magnitude}(\operatorname{line\text{-}int}(f, (\lambda x \in \operatorname{ccint}(0, 1).\; (z_{0} + (x \cdot w))), (\lambda x \in \operatorname{ccint}(0, 1).\; w), 0, 1)) \\ \quad \quad \leq\; \operatorname{magnitude}(w) \cdot n \end{array}}}
\vnbprop{segment-int-reverse}{\vnbstmt{\begin{array}{@{}l@{}} \forall u, f.; \forall z_{0}, w \in \mathbb{C}. \\ \quad u \subseteq \mathbb{C} \wedge \\ \quad f \in (u \to \mathbb{C}) \wedge \\ \quad \forall y \in u. \\ \quad \quad \operatorname{is\text{-}continuous\text{-}at}(\operatorname{subspace\text{-}ms}(\operatorname{nf\text{-}metric\text{-}space}(\operatorname{cc\text{-}normed\text{-}field}), u), \\ \quad \quad \quad \operatorname{nf\text{-}metric\text{-}space}(\operatorname{cc\text{-}normed\text{-}field}), f, y) \wedge \\ \quad \operatorname{trace}((\lambda x \in \operatorname{ccint}(0, 1).\; (z_{0} + (x \cdot w))), 0, 1) \subseteq u \Rightarrow \\ \quad \quad \operatorname{line\text{-}int}(f, (\lambda x \in \operatorname{ccint}(0, 1).\; ((z_{0} + w) + (x \cdot (-1 \cdot w)))), \\ \quad \quad \quad (\lambda x \in \operatorname{ccint}(0, 1).\; (-1 \cdot w)), 0, 1) \\ \quad \quad =\; -1 \cdot \operatorname{line\text{-}int}(f, (\lambda x \in \operatorname{ccint}(0, 1).\; (z_{0} + (x \cdot w))), (\lambda x \in \operatorname{ccint}(0, 1).\; w), 0, 1) \end{array}}}

\vnbbookfile{theorem-library/goursat}{theorem-library/goursat}
\noindent{\small goursat.scm -- CAUCHY-GOURSAT FOR A TRIANGLE AND CAUCHY ON A CONVEX SET (complex-analysis.pdf ch. 3, s. 3.2: Lemma 3.4, Lemma 3.5, Prop 3.6, Cor 3.7, Prop 3.8, Cor 3.9, all for f holomorphic on the whole open set).  Batch 30, 2026-09-25.}\par
\vnbprop{seg-int-unfold}{\vnbstmt{\begin{array}{@{}l@{}} \forall f, a, b. \\ \quad (\operatorname{seg\text{-}int}(f, a, b) \simeq \operatorname{line\text{-}int}(f, \operatorname{seg\text{-}path}(a, b), \operatorname{seg\text{-}deriv}(a, b), 0, 1)) \end{array}}}
\vnbprop{tri-int-unfold}{\vnbstmt{\begin{array}{@{}l@{}} \forall f, a, b, c. \\ \quad (\operatorname{tri\text{-}int}(f, a, b, c) \simeq ((\operatorname{seg\text{-}int}(f, a, b) + \operatorname{seg\text{-}int}(f, b, c)) + \operatorname{seg\text{-}int}(f, c, a))) \end{array}}}
\vnbprop{cc-mid-unfold}{\vnbstmt{\forall x, y.\; (\operatorname{cc\text{-}mid}(x, y) \simeq ((x + y) \cdot \operatorname{recip}(2)))}}
\vnbprop{seg-path-bridge}{\vnbstmt{\begin{array}{@{}l@{}} \forall a, b, p, w. \\ \quad ((a = p) \wedge ((b - a) = w)) \Rightarrow  \\ \quad \quad (\operatorname{seg\text{-}path}(a, b) \simeq (\lambda x \in \operatorname{ccint}(0, 1).\; (p + (x \cdot w)))) \end{array}}}
\vnbprop{seg-deriv-bridge}{\vnbstmt{\begin{array}{@{}l@{}} \forall a, b, w. \\ \quad ((b - a) = w) \Rightarrow (\operatorname{seg\text{-}deriv}(a, b) \simeq (\lambda x \in \operatorname{ccint}(0, 1).\; w)) \end{array}}}
\vnbprop{seg-int-bridge}{\vnbstmt{\begin{array}{@{}l@{}} \forall f, a, b, p, w. \\ \quad ((a = p) \wedge ((b - a) = w)) \Rightarrow  \\ \quad \quad (\operatorname{seg\text{-}int}(f, a, b) \simeq \operatorname{line\text{-}int}(f, (\lambda x \in \operatorname{ccint}(0, 1).\; (p + (x \cdot w))), (\lambda x \in \operatorname{ccint}(0, 1).\; w), 0, 1)) \end{array}}}
\vnbprop{seg-trace-bridge}{\vnbstmt{\begin{array}{@{}l@{}} \forall a, b, p, w. \\ \quad ((a = p) \wedge ((b - a) = w)) \Rightarrow  \\ \quad \quad (\operatorname{trace}(\operatorname{seg\text{-}path}(a, b), 0, 1) \simeq \operatorname{trace}((\lambda x \in \operatorname{ccint}(0, 1).\; (p + (x \cdot w))), 0, 1)) \end{array}}}
\vnbprop{seg-path-is-road}{\vnbstmt{\forall a, b \in \mathbb{C}.\; \operatorname{is\text{-}road}(\operatorname{seg\text{-}path}(a, b), \operatorname{seg\text{-}deriv}(a, b), 0, 1)}}
\vnbprop{in-transport-eq}{\vnbstmt{\forall x, y, s.\; (((x \in s) \wedge (x = y)) \Rightarrow (y \in s))}}
\vnbprop{seg-path-apply}{\vnbstmt{\begin{array}{@{}l@{}} \forall a, b \in \mathbb{C}, t \in \operatorname{ccint}(0, 1). \\ \quad \operatorname{seg\text{-}path}(a, b)(t) = (a + (t \cdot (b - a))) \end{array}}}
\vnbprop{seg-trace-intro}{\vnbstmt{\begin{array}{@{}l@{}} \forall a, b \in \mathbb{C}, t \in \operatorname{ccint}(0, 1). \\ \quad (a + (t \cdot (b - a))) \in \operatorname{trace}(\operatorname{seg\text{-}path}(a, b), 0, 1) \end{array}}}
\vnbprop{seg-trace-elim}{\vnbstmt{\begin{array}{@{}l@{}} \forall a, b \in \mathbb{C}, z \in \operatorname{trace}(\operatorname{seg\text{-}path}(a, b), 0, 1). \exists t \in \operatorname{ccint}(0, 1). \\ \quad z = (a + (t \cdot (b - a))) \end{array}}}
\vnbprop{conv3-member-iff}{\vnbstmt{\begin{array}{@{}l@{}} \forall a, b, c, z. \\ \quad z \in \operatorname{conv3}(a, b, c) \\ \quad \Leftrightarrow\; z \in \mathbb{C} \wedge \\ \quad \exists l, m, n \in \operatorname{ccint}(0, 1). \\ \quad \quad (((l + m) + n) = 1) \wedge (z = (((l \cdot a) + (m \cdot b)) + (n \cdot c))) \end{array}}}
\vnbprop{conv3-intro}{\vnbstmt{\begin{array}{@{}l@{}} \forall a, b, c \in \mathbb{C}, l, m, n \in \operatorname{ccint}(0, 1), z. \\ \quad ((((l + m) + n) = 1) \wedge (z = (((l \cdot a) + (m \cdot b)) + (n \cdot c)))) \Rightarrow (z \in \operatorname{conv3}(a, b, c)) \end{array}}}
\vnbprop{is-convex-unfold}{\vnbstmt{\begin{array}{@{}l@{}} \forall u.; \forall x, y \in u, t \in \operatorname{ccint}(0, 1). \\ \quad \operatorname{is\text{-}convex\text{-}cc}(u) \Rightarrow ((x + (t \cdot (y - x))) \in u) \end{array}}}
\vnbprop{conv3-is-convex}{\vnbstmt{\forall a, b, c \in \mathbb{C}.\; \operatorname{is\text{-}convex\text{-}cc}(\operatorname{conv3}(a, b, c))}}
\vnbprop{seg-trace-in-convex}{\vnbstmt{\begin{array}{@{}l@{}} \forall u.; \forall a, b \in \mathbb{C}. \\ \quad (\operatorname{is\text{-}convex\text{-}cc}(u) \wedge ((a \in u) \wedge (b \in u))) \Rightarrow  \\ \quad \quad \operatorname{trace}(\operatorname{seg\text{-}path}(a, b), 0, 1) \subseteq u \end{array}}}
\vnbprop{conv3-in-convex}{\vnbstmt{\begin{array}{@{}l@{}} \forall u.; \forall a, b, c \in \mathbb{C}. \\ \quad (\operatorname{is\text{-}convex\text{-}cc}(u) \wedge ((a \in u) \wedge ((b \in u) \wedge (c \in u)))) \Rightarrow (\operatorname{conv3}(a, b, c) \subseteq u) \end{array}}}
\vnbprop{conv3-vertices-in}{\vnbstmt{\begin{array}{@{}l@{}} \forall a, b, c \in \mathbb{C}. \\ \quad (a \in \operatorname{conv3}(a, b, c)) \wedge ((b \in \operatorname{conv3}(a, b, c)) \wedge (c \in \operatorname{conv3}(a, b, c))) \end{array}}}
\vnbprop{conv3-seg-trace}{\vnbstmt{\begin{array}{@{}l@{}} \forall a, b, c \in \mathbb{C}, x, y \in \operatorname{conv3}(a, b, c). \\ \quad \operatorname{trace}(\operatorname{seg\text{-}path}(x, y), 0, 1) \subseteq \operatorname{conv3}(a, b, c) \end{array}}}
\vnbprop{conv3-sub-hull}{\vnbstmt{\begin{array}{@{}l@{}} \forall a, b, c \in \mathbb{C}, x, y, z \in \operatorname{conv3}(a, b, c). \\ \quad \operatorname{conv3}(x, y, z) \subseteq \operatorname{conv3}(a, b, c) \end{array}}}
\vnbprop{cc-mid-affine}{\vnbstmt{\begin{array}{@{}l@{}} \forall x, y \in \mathbb{C}. \\ \quad (\operatorname{cc\text{-}mid}(x, y) = (x + (1/2 \cdot (y - x)))) \wedge (\operatorname{cc\text{-}mid}(x, y) = (y + (1/2 \cdot (x - y)))) \end{array}}}
\vnbprop{cc-mid-in-seg-trace}{\vnbstmt{\begin{array}{@{}l@{}} \forall x, y \in \mathbb{C}. \\ \quad ((\operatorname{cc\text{-}mid}(x, y) \in \operatorname{trace}(\operatorname{seg\text{-}path}(x, y), 0, 1)) \wedge (\operatorname{cc\text{-}mid}(x, y) \in \operatorname{trace}(\operatorname{seg\text{-}path}(y, x), 0, 1))) \end{array}}}
\vnbprop{conv3-mid-in}{\vnbstmt{\begin{array}{@{}l@{}} \forall a, b, c \in \mathbb{C}, x, y \in \operatorname{conv3}(a, b, c). \\ \quad \operatorname{cc\text{-}mid}(x, y) \in \operatorname{conv3}(a, b, c) \end{array}}}
\vnbprop{seg-trace-in-conv3}{\vnbstmt{\begin{array}{@{}l@{}} \forall a, b, c \in \mathbb{C}. \\ \quad \operatorname{trace}(\operatorname{seg\text{-}path}(a, b), 0, 1) \subseteq \operatorname{conv3}(a, b, c) \wedge \\ \quad \operatorname{trace}(\operatorname{seg\text{-}path}(b, c), 0, 1) \subseteq \operatorname{conv3}(a, b, c) \wedge \\ \quad \operatorname{trace}(\operatorname{seg\text{-}path}(c, a), 0, 1) \subseteq \operatorname{conv3}(a, b, c) \end{array}}}
\vnbprop{conv3-sub-triangle}{\vnbstmt{\begin{array}{@{}l@{}} \forall a, b, c \in \mathbb{C}. \\ \quad \operatorname{conv3}(a, \operatorname{cc\text{-}mid}(a, b), \operatorname{cc\text{-}mid}(a, c)) \subseteq \operatorname{conv3}(a, b, c) \wedge \\ \quad \operatorname{conv3}(\operatorname{cc\text{-}mid}(a, b), b, \operatorname{cc\text{-}mid}(b, c)) \subseteq \operatorname{conv3}(a, b, c) \wedge \\ \quad \operatorname{conv3}(\operatorname{cc\text{-}mid}(a, c), \operatorname{cc\text{-}mid}(b, c), c) \subseteq \operatorname{conv3}(a, b, c) \wedge \\ \quad \operatorname{conv3}(\operatorname{cc\text{-}mid}(a, b), \operatorname{cc\text{-}mid}(b, c), \operatorname{cc\text{-}mid}(a, c)) \subseteq \operatorname{conv3}(a, b, c) \end{array}}}
\vnbprop{seg-int-in-cc}{\vnbstmt{\begin{array}{@{}l@{}} \forall u, f.; \forall a, b \in \mathbb{C}. \\ \quad u \subseteq \mathbb{C} \wedge \\ \quad f \in (u \to \mathbb{C}) \wedge \\ \quad \forall y \in u. \\ \quad \quad \operatorname{is\text{-}continuous\text{-}at}(\operatorname{subspace\text{-}ms}(\operatorname{nf\text{-}metric\text{-}space}(\operatorname{cc\text{-}normed\text{-}field}), u), \\ \quad \quad \quad \operatorname{nf\text{-}metric\text{-}space}(\operatorname{cc\text{-}normed\text{-}field}), f, y) \wedge \\ \quad \operatorname{trace}(\operatorname{seg\text{-}path}(a, b), 0, 1) \subseteq u \Rightarrow \\ \quad \quad \operatorname{seg\text{-}int}(f, a, b) \in \mathbb{C} \end{array}}}
\vnbprop{seg-int-degenerate}{\vnbstmt{\begin{array}{@{}l@{}} \forall u, f.; \forall z \in u. \\ \quad u \subseteq \mathbb{C} \wedge \\ \quad f \in (u \to \mathbb{C}) \wedge \\ \quad \forall y \in u. \\ \quad \quad \operatorname{is\text{-}continuous\text{-}at}(\operatorname{subspace\text{-}ms}(\operatorname{nf\text{-}metric\text{-}space}(\operatorname{cc\text{-}normed\text{-}field}), u), \\ \quad \quad \quad \operatorname{nf\text{-}metric\text{-}space}(\operatorname{cc\text{-}normed\text{-}field}), f, y) \Rightarrow \\ \quad \quad \operatorname{seg\text{-}int}(f, z, z) = 0 \end{array}}}
\vnbprop{seg-int-reverse-at}{\vnbstmt{\begin{array}{@{}l@{}} \forall u, f.; \forall a, b \in \mathbb{C}. \\ \quad u \subseteq \mathbb{C} \wedge \\ \quad f \in (u \to \mathbb{C}) \wedge \\ \quad \forall y \in u. \\ \quad \quad \operatorname{is\text{-}continuous\text{-}at}(\operatorname{subspace\text{-}ms}(\operatorname{nf\text{-}metric\text{-}space}(\operatorname{cc\text{-}normed\text{-}field}), u), \\ \quad \quad \quad \operatorname{nf\text{-}metric\text{-}space}(\operatorname{cc\text{-}normed\text{-}field}), f, y) \wedge \\ \quad \operatorname{trace}(\operatorname{seg\text{-}path}(a, b), 0, 1) \subseteq u \Rightarrow \\ \quad \quad \operatorname{seg\text{-}int}(f, b, a) = (-1 \cdot \operatorname{seg\text{-}int}(f, a, b)) \end{array}}}
\vnbprop{seg-int-split-at}{\vnbstmt{\begin{array}{@{}l@{}} \forall u, f.; \forall a, b \in \mathbb{C}, c \in \operatorname{trace}(\operatorname{seg\text{-}path}(a, b), 0, 1). \\ \quad u \subseteq \mathbb{C} \wedge \\ \quad f \in (u \to \mathbb{C}) \wedge \\ \quad \forall y \in u. \\ \quad \quad \operatorname{is\text{-}continuous\text{-}at}(\operatorname{subspace\text{-}ms}(\operatorname{nf\text{-}metric\text{-}space}(\operatorname{cc\text{-}normed\text{-}field}), u), \\ \quad \quad \quad \operatorname{nf\text{-}metric\text{-}space}(\operatorname{cc\text{-}normed\text{-}field}), f, y) \wedge \\ \quad \operatorname{trace}(\operatorname{seg\text{-}path}(a, b), 0, 1) \subseteq u \Rightarrow \\ \quad \quad \operatorname{seg\text{-}int}(f, a, b) = (\operatorname{seg\text{-}int}(f, a, c) + \operatorname{seg\text{-}int}(f, c, b)) \end{array}}}
\vnbprop{tri-int-subdivision}{\vnbstmt{\begin{array}{@{}l@{}} \forall u, f.; \forall a, b, c \in \mathbb{C}. \\ \quad u \subseteq \mathbb{C} \wedge \\ \quad f \in (u \to \mathbb{C}) \wedge \\ \quad \forall y \in u. \\ \quad \quad \operatorname{is\text{-}continuous\text{-}at}(\operatorname{subspace\text{-}ms}(\operatorname{nf\text{-}metric\text{-}space}(\operatorname{cc\text{-}normed\text{-}field}), u), \\ \quad \quad \quad \operatorname{nf\text{-}metric\text{-}space}(\operatorname{cc\text{-}normed\text{-}field}), f, y) \wedge \\ \quad \operatorname{conv3}(a, b, c) \subseteq u \Rightarrow \\ \quad \quad \operatorname{tri\text{-}int}(f, a, b, c) \\ \quad \quad =\; ((\operatorname{tri\text{-}int}(f, a, \operatorname{cc\text{-}mid}(a, b), \operatorname{cc\text{-}mid}(a, c)) + \operatorname{tri\text{-}int}(f, \operatorname{cc\text{-}mid}(a, b), b, \operatorname{cc\text{-}mid}(b, c))) + \operatorname{tri\text{-}int}(f, \operatorname{cc\text{-}mid}(a, c), \operatorname{cc\text{-}mid}(b, c), c)) + \operatorname{tri\text{-}int}(f, \operatorname{cc\text{-}mid}(a, b), \operatorname{cc\text{-}mid}(b, c), \operatorname{cc\text{-}mid}(a, c)) \end{array}}}
\vnbprop{tri-int-abs-bound}{\vnbstmt{\begin{array}{@{}l@{}} \forall u, f.; \forall a, b, c \in \mathbb{C}, m \in \mathbb{R}. \\ \quad u \subseteq \mathbb{C} \wedge \\ \quad f \in (u \to \mathbb{C}) \wedge \\ \quad \forall y \in u. \\ \quad \quad \operatorname{is\text{-}continuous\text{-}at}(\operatorname{subspace\text{-}ms}(\operatorname{nf\text{-}metric\text{-}space}(\operatorname{cc\text{-}normed\text{-}field}), u), \\ \quad \quad \quad \operatorname{nf\text{-}metric\text{-}space}(\operatorname{cc\text{-}normed\text{-}field}), f, y) \wedge \\ \quad \operatorname{conv3}(a, b, c) \subseteq u \wedge \\ \quad \forall z \in \operatorname{conv3}(a, b, c).\; (\operatorname{magnitude}(f(z)) \leq m) \Rightarrow \\ \quad \quad \operatorname{magnitude}(\operatorname{tri\text{-}int}(f, a, b, c)) \leq (((\operatorname{magnitude}((b - a)) + \operatorname{magnitude}((c - b))) + \operatorname{magnitude}((a - c))) \cdot m) \end{array}}}
\vnbprop{tri-int-sum}{\vnbstmt{\begin{array}{@{}l@{}} \forall u, f, g.; \forall h \in (u \to \mathbb{C}), a, b, c \in \mathbb{C}. \\ \quad u \subseteq \mathbb{C} \wedge \\ \quad f \in (u \to \mathbb{C}) \wedge \\ \quad \forall y \in u. \\ \quad \quad \operatorname{is\text{-}continuous\text{-}at}(\operatorname{subspace\text{-}ms}(\operatorname{nf\text{-}metric\text{-}space}(\operatorname{cc\text{-}normed\text{-}field}), u), \\ \quad \quad \quad \operatorname{nf\text{-}metric\text{-}space}(\operatorname{cc\text{-}normed\text{-}field}), f, y) \wedge \\ \quad g \in (u \to \mathbb{C}) \wedge \\ \quad \forall y \in u. \\ \quad \quad \operatorname{is\text{-}continuous\text{-}at}(\operatorname{subspace\text{-}ms}(\operatorname{nf\text{-}metric\text{-}space}(\operatorname{cc\text{-}normed\text{-}field}), u), \\ \quad \quad \quad \operatorname{nf\text{-}metric\text{-}space}(\operatorname{cc\text{-}normed\text{-}field}), g, y) \wedge \\ \quad \forall z \in u.\; (h(z) \simeq (f(z) + g(z))) \wedge \\ \quad \operatorname{conv3}(a, b, c) \subseteq u \Rightarrow \\ \quad \quad \operatorname{tri\text{-}int}(h, a, b, c) = (\operatorname{tri\text{-}int}(f, a, b, c) + \operatorname{tri\text{-}int}(g, a, b, c)) \end{array}}}
\vnbprop{cc-open-subset-cc}{\vnbstmt{\forall u.\; (\operatorname{is\text{-}open}(\operatorname{nf\text{-}metric\text{-}space}(\operatorname{cc\text{-}normed\text{-}field}), u) \Rightarrow (u \subseteq \mathbb{C}))}}
\vnbprop{affine-primitive-on}{\vnbstmt{\begin{array}{@{}l@{}} \forall u.; \forall p, q, w \in \mathbb{C}, f \in (u \to \mathbb{C}). \\ \quad \operatorname{is\text{-}open}(\operatorname{nf\text{-}metric\text{-}space}(\operatorname{cc\text{-}normed\text{-}field}), u) \wedge \\ \quad \forall z \in u.\; (f(z) \simeq ((p \cdot z) + ((q \cdot ((z - w) \cdot (z - w))) \cdot \operatorname{recip}(2)))) \Rightarrow \\ \quad \quad (\operatorname{holomorphic\text{-}on}(u, f) \wedge \forall z \in u.\; \operatorname{is\text{-}diff\text{-}on}(\operatorname{cc\text{-}normed\text{-}field}, u, f, z, (p + (q \cdot (z - w))))) \end{array}}}
\vnbprop{tri-int-affine-zero}{\vnbstmt{\begin{array}{@{}l@{}} \forall u.; \forall p, q, w \in \mathbb{C}, g \in (u \to \mathbb{C}), a, b, c \in \mathbb{C}. \\ \quad \operatorname{is\text{-}open}(\operatorname{nf\text{-}metric\text{-}space}(\operatorname{cc\text{-}normed\text{-}field}), u) \wedge \\ \quad \forall z \in u.\; (g(z) \simeq (p + (q \cdot (z - w)))) \wedge \\ \quad \operatorname{conv3}(a, b, c) \subseteq u \Rightarrow \\ \quad \quad \operatorname{tri\text{-}int}(g, a, b, c) = 0 \end{array}}}
\vnbprop{holomorphic-on-continuous}{\vnbstmt{\begin{array}{@{}l@{}} \forall u, f.; \forall y \in u. \\ \quad \operatorname{holomorphic\text{-}on}(u, f) \Rightarrow  \\ \quad \quad \operatorname{is\text{-}continuous\text{-}at}(\operatorname{subspace\text{-}ms}(\operatorname{nf\text{-}metric\text{-}space}(\operatorname{cc\text{-}normed\text{-}field}), u), \\ \quad \quad \quad \operatorname{nf\text{-}metric\text{-}space}(\operatorname{cc\text{-}normed\text{-}field}), f, y) \end{array}}}
\vnbprop{cc-mid-coords}{\vnbstmt{\begin{array}{@{}l@{}} \forall x, y \in \mathbb{C}. \\ \quad \operatorname{real\text{-}part}(\operatorname{cc\text{-}mid}(x, y)) = (1/2 \cdot (\operatorname{real\text{-}part}(x) + \operatorname{real\text{-}part}(y))) \wedge \\ \quad \operatorname{imag\text{-}part}(\operatorname{cc\text{-}mid}(x, y)) = (1/2 \cdot (\operatorname{imag\text{-}part}(x) + \operatorname{imag\text{-}part}(y))) \end{array}}}
\vnbprop{std-mid}{\vnbstmt{\begin{array}{@{}l@{}} \forall x, y \in \mathbb{C}. \\ \quad 0 \leq \operatorname{real\text{-}part}(x) \wedge \\ \quad 0 \leq \operatorname{imag\text{-}part}(x) \wedge \\ \quad (\operatorname{real\text{-}part}(x) + \operatorname{imag\text{-}part}(x)) \leq 1 \wedge \\ \quad 0 \leq \operatorname{real\text{-}part}(y) \wedge \\ \quad 0 \leq \operatorname{imag\text{-}part}(y) \wedge \\ \quad (\operatorname{real\text{-}part}(y) + \operatorname{imag\text{-}part}(y)) \leq 1 \Rightarrow \\ \quad \quad 0 \leq \operatorname{real\text{-}part}(\operatorname{cc\text{-}mid}(x, y)) \wedge \\ \quad \quad 0 \leq \operatorname{imag\text{-}part}(\operatorname{cc\text{-}mid}(x, y)) \wedge \\ \quad \quad (\operatorname{real\text{-}part}(\operatorname{cc\text{-}mid}(x, y)) + \operatorname{imag\text{-}part}(\operatorname{cc\text{-}mid}(x, y))) \leq 1 \end{array}}}
\vnbprop{std-dist}{\vnbstmt{\begin{array}{@{}l@{}} \forall x, y \in \mathbb{C}. \\ \quad 0 \leq \operatorname{real\text{-}part}(x) \wedge \\ \quad 0 \leq \operatorname{imag\text{-}part}(x) \wedge \\ \quad (\operatorname{real\text{-}part}(x) + \operatorname{imag\text{-}part}(x)) \leq 1 \wedge \\ \quad 0 \leq \operatorname{real\text{-}part}(y) \wedge \\ \quad 0 \leq \operatorname{imag\text{-}part}(y) \wedge \\ \quad (\operatorname{real\text{-}part}(y) + \operatorname{imag\text{-}part}(y)) \leq 1 \Rightarrow \\ \quad \quad \operatorname{magnitude}((x - y)) \leq 2 \end{array}}}
\vnbprop{std-limit}{\vnbstmt{\begin{array}{@{}l@{}} \forall s \in (\mathbb{N} \to \mathbb{C}), w \in \mathbb{C}. \\ \quad \forall n \in \mathbb{N}. \\ \quad \quad ((0 \leq \operatorname{real\text{-}part}(s(n))) \wedge ((0 \leq \operatorname{imag\text{-}part}(s(n))) \wedge ((\operatorname{real\text{-}part}(s(n)) + \operatorname{imag\text{-}part}(s(n))) \leq 1))) \wedge \\ \quad \operatorname{converges\text{-}to}(\operatorname{cc\text{-}ms}, s, w) \Rightarrow \\ \quad \quad (0 \leq \operatorname{real\text{-}part}(w)) \wedge ((0 \leq \operatorname{imag\text{-}part}(w)) \wedge ((\operatorname{real\text{-}part}(w) + \operatorname{imag\text{-}part}(w)) \leq 1)) \end{array}}}
\vnbprop{bary-mid}{\vnbstmt{\begin{array}{@{}l@{}} \forall a, b, c, x, y \in \mathbb{C}. \\ \quad ((\operatorname{real\text{-}part}(\operatorname{cc\text{-}mid}(x, y)) \cdot a) + (\operatorname{imag\text{-}part}(\operatorname{cc\text{-}mid}(x, y)) \cdot b)) + (((1 - \operatorname{real\text{-}part}(\operatorname{cc\text{-}mid}(x, y))) - \operatorname{imag\text{-}part}(\operatorname{cc\text{-}mid}(x, y))) \cdot c) \\ \quad =\; \operatorname{cc\text{-}mid}(((\operatorname{real\text{-}part}(x) \cdot a) + (\operatorname{imag\text{-}part}(x) \cdot b)) + (((1 - \operatorname{real\text{-}part}(x)) - \operatorname{imag\text{-}part}(x)) \cdot c), \\ \quad \quad ((\operatorname{real\text{-}part}(y) \cdot a) + (\operatorname{imag\text{-}part}(y) \cdot b)) + (((1 - \operatorname{real\text{-}part}(y)) - \operatorname{imag\text{-}part}(y)) \cdot c)) \end{array}}}
\vnbprop{bary-in-conv3}{\vnbstmt{\begin{array}{@{}l@{}} \forall a, b, c, x \in \mathbb{C}. \\ \quad ((0 \leq \operatorname{real\text{-}part}(x)) \wedge ((0 \leq \operatorname{imag\text{-}part}(x)) \wedge ((\operatorname{real\text{-}part}(x) + \operatorname{imag\text{-}part}(x)) \leq 1))) \Rightarrow  \\ \quad \quad (((\operatorname{real\text{-}part}(x) \cdot a) + (\operatorname{imag\text{-}part}(x) \cdot b)) + (((1 - \operatorname{real\text{-}part}(x)) - \operatorname{imag\text{-}part}(x)) \cdot c)) \in \operatorname{conv3}(a, b, c) \end{array}}}
\vnbprop{bary-vertices}{\vnbstmt{\begin{array}{@{}l@{}} \forall a, b, c \in \mathbb{C}. \\ \quad (((\operatorname{real\text{-}part}(1) \cdot a) + (\operatorname{imag\text{-}part}(1) \cdot b)) + (((1 - \operatorname{real\text{-}part}(1)) - \operatorname{imag\text{-}part}(1)) \cdot c)) = a \wedge \\ \quad (((\operatorname{real\text{-}part}(+i) \cdot a) + (\operatorname{imag\text{-}part}(+i) \cdot b)) + (((1 - \operatorname{real\text{-}part}(+i)) - \operatorname{imag\text{-}part}(+i)) \cdot c)) = b \wedge \\ \quad (((\operatorname{real\text{-}part}(0) \cdot a) + (\operatorname{imag\text{-}part}(0) \cdot b)) + (((1 - \operatorname{real\text{-}part}(0)) - \operatorname{imag\text{-}part}(0)) \cdot c)) = c \end{array}}}
\vnbprop{bary-lip}{\vnbstmt{\begin{array}{@{}l@{}} \forall a, b, c, x, y \in \mathbb{C}. \\ \quad \operatorname{magnitude}(((((\operatorname{real\text{-}part}(x) \cdot a) + (\operatorname{imag\text{-}part}(x) \cdot b)) + (((1 - \operatorname{real\text{-}part}(x)) - \operatorname{imag\text{-}part}(x)) \cdot c)) - (((\operatorname{real\text{-}part}(y) \cdot a) + (\operatorname{imag\text{-}part}(y) \cdot b)) + (((1 - \operatorname{real\text{-}part}(y)) - \operatorname{imag\text{-}part}(y)) \cdot c)))) \\ \quad \leq\; (\operatorname{magnitude}((a - c)) + \operatorname{magnitude}((b - c))) \cdot \operatorname{magnitude}((x - y)) \end{array}}}
\vnbprop{half-magnitude}{\vnbstmt{\forall x \in \mathbb{C}.\; (\operatorname{magnitude}((1/2 \cdot x)) = (1/2 \cdot \operatorname{magnitude}(x)))}}
\vnbprop{tri-child-geom-1}{\vnbstmt{\begin{array}{@{}l@{}} \forall x, y, z \in \mathbb{C}. \\ \quad (\operatorname{magnitude}((\operatorname{cc\text{-}mid}(x, y) - x)) + \operatorname{magnitude}((\operatorname{cc\text{-}mid}(x, z) - \operatorname{cc\text{-}mid}(x, y)))) + \operatorname{magnitude}((x - \operatorname{cc\text{-}mid}(x, z))) \\ \quad =\; 1/2 \cdot ((\operatorname{magnitude}((y - x)) + \operatorname{magnitude}((z - y))) + \operatorname{magnitude}((x - z))) \wedge \\ \quad \operatorname{magnitude}((x - x)) \leq ((\operatorname{magnitude}((y - x)) + \operatorname{magnitude}((z - y))) + \operatorname{magnitude}((x - z))) \end{array}}}
\vnbprop{tri-child-geom-2}{\vnbstmt{\begin{array}{@{}l@{}} \forall x, y, z \in \mathbb{C}. \\ \quad (\operatorname{magnitude}((y - \operatorname{cc\text{-}mid}(x, y))) + \operatorname{magnitude}((\operatorname{cc\text{-}mid}(y, z) - y))) + \operatorname{magnitude}((\operatorname{cc\text{-}mid}(x, y) - \operatorname{cc\text{-}mid}(y, z))) \\ \quad =\; 1/2 \cdot ((\operatorname{magnitude}((y - x)) + \operatorname{magnitude}((z - y))) + \operatorname{magnitude}((x - z))) \wedge \\ \quad \operatorname{magnitude}((\operatorname{cc\text{-}mid}(x, y) - x)) \leq ((\operatorname{magnitude}((y - x)) + \operatorname{magnitude}((z - y))) + \operatorname{magnitude}((x - z))) \end{array}}}
\vnbprop{tri-child-geom-3}{\vnbstmt{\begin{array}{@{}l@{}} \forall x, y, z \in \mathbb{C}. \\ \quad (\operatorname{magnitude}((\operatorname{cc\text{-}mid}(y, z) - \operatorname{cc\text{-}mid}(x, z))) + \operatorname{magnitude}((z - \operatorname{cc\text{-}mid}(y, z)))) + \operatorname{magnitude}((\operatorname{cc\text{-}mid}(x, z) - z)) \\ \quad =\; 1/2 \cdot ((\operatorname{magnitude}((y - x)) + \operatorname{magnitude}((z - y))) + \operatorname{magnitude}((x - z))) \wedge \\ \quad \operatorname{magnitude}((\operatorname{cc\text{-}mid}(x, z) - x)) \leq ((\operatorname{magnitude}((y - x)) + \operatorname{magnitude}((z - y))) + \operatorname{magnitude}((x - z))) \end{array}}}
\vnbprop{tri-child-geom-4}{\vnbstmt{\begin{array}{@{}l@{}} \forall x, y, z \in \mathbb{C}. \\ \quad (\operatorname{magnitude}((\operatorname{cc\text{-}mid}(y, z) - \operatorname{cc\text{-}mid}(x, y))) + \operatorname{magnitude}((\operatorname{cc\text{-}mid}(x, z) - \operatorname{cc\text{-}mid}(y, z)))) + \operatorname{magnitude}((\operatorname{cc\text{-}mid}(x, y) - \operatorname{cc\text{-}mid}(x, z))) \\ \quad =\; 1/2 \cdot ((\operatorname{magnitude}((y - x)) + \operatorname{magnitude}((z - y))) + \operatorname{magnitude}((x - z))) \wedge \\ \quad \operatorname{magnitude}((\operatorname{cc\text{-}mid}(x, y) - x)) \leq ((\operatorname{magnitude}((y - x)) + \operatorname{magnitude}((z - y))) + \operatorname{magnitude}((x - z))) \end{array}}}
\vnbprop{tri-int-in-cc}{\vnbstmt{\begin{array}{@{}l@{}} \forall u, f.; \forall a, b, c \in \mathbb{C}. \\ \quad u \subseteq \mathbb{C} \wedge \\ \quad f \in (u \to \mathbb{C}) \wedge \\ \quad \forall y \in u. \\ \quad \quad \operatorname{is\text{-}continuous\text{-}at}(\operatorname{subspace\text{-}ms}(\operatorname{nf\text{-}metric\text{-}space}(\operatorname{cc\text{-}normed\text{-}field}), u), \\ \quad \quad \quad \operatorname{nf\text{-}metric\text{-}space}(\operatorname{cc\text{-}normed\text{-}field}), f, y) \wedge \\ \quad \operatorname{conv3}(a, b, c) \subseteq u \Rightarrow \\ \quad \quad \operatorname{tri\text{-}int}(f, a, b, c) \in \mathbb{C} \end{array}}}
\vnbprop{goursat-nest}{\vnbstmt{\begin{array}{@{}l@{}} \forall u, f.; \forall a, b, c \in \mathbb{C}. \\ \quad u \subseteq \mathbb{C} \wedge \\ \quad f \in (u \to \mathbb{C}) \wedge \\ \quad \forall y \in u. \\ \quad \quad \operatorname{is\text{-}continuous\text{-}at}(\operatorname{subspace\text{-}ms}(\operatorname{nf\text{-}metric\text{-}space}(\operatorname{cc\text{-}normed\text{-}field}), u), \\ \quad \quad \quad \operatorname{nf\text{-}metric\text{-}space}(\operatorname{cc\text{-}normed\text{-}field}), f, y) \wedge \\ \quad \operatorname{conv3}(a, b, c) \subseteq u \Rightarrow \\ \quad \quad \exists w \in \mathbb{C}. \\ \quad \quad \quad 0 \leq \operatorname{real\text{-}part}(w) \wedge \\ \quad \quad \quad 0 \leq \operatorname{imag\text{-}part}(w) \wedge \\ \quad \quad \quad (\operatorname{real\text{-}part}(w) + \operatorname{imag\text{-}part}(w)) \leq 1 \wedge \\ \quad \quad \quad \forall n \in \mathbb{N}. \exists x, \mathit{gouy}, z \in \mathbb{C}. \\ \quad \quad \quad \quad 0 \leq \operatorname{real\text{-}part}(x) \wedge \\ \quad \quad \quad \quad 0 \leq \operatorname{imag\text{-}part}(x) \wedge \\ \quad \quad \quad \quad (\operatorname{real\text{-}part}(x) + \operatorname{imag\text{-}part}(x)) \leq 1 \wedge \\ \quad \quad \quad \quad 0 \leq \operatorname{real\text{-}part}(\mathit{gouy}) \wedge \\ \quad \quad \quad \quad 0 \leq \operatorname{imag\text{-}part}(\mathit{gouy}) \wedge \\ \quad \quad \quad \quad (\operatorname{real\text{-}part}(\mathit{gouy}) + \operatorname{imag\text{-}part}(\mathit{gouy})) \leq 1 \wedge \\ \quad \quad \quad \quad 0 \leq \operatorname{real\text{-}part}(z) \wedge \\ \quad \quad \quad \quad 0 \leq \operatorname{imag\text{-}part}(z) \wedge \\ \quad \quad \quad \quad (\operatorname{real\text{-}part}(z) + \operatorname{imag\text{-}part}(z)) \leq 1 \wedge \\ \quad \quad \quad \quad ((\operatorname{magnitude}((\mathit{gouy} - x)) + \operatorname{magnitude}((z - \mathit{gouy}))) + \operatorname{magnitude}((x - z))) \leq (6 \cdot \operatorname{recip}(\operatorname{power}(2, n))) \wedge \\ \quad \quad \quad \quad \operatorname{magnitude}((x - w)) \leq (12 \cdot \operatorname{recip}(\operatorname{power}(2, n))) \wedge \\ \quad \quad \quad \quad \operatorname{magnitude}(\operatorname{tri\text{-}int}(f, a, b, c)) \\ \quad \quad \quad \quad \leq\; \operatorname{power}(4, n) \cdot \operatorname{magnitude}(\operatorname{tri\text{-}int}(f, (((\operatorname{real\text{-}part}(x) \cdot a) + (\operatorname{imag\text{-}part}(x) \cdot b)) + (((1 - \operatorname{real\text{-}part}(x)) - \operatorname{imag\text{-}part}(x)) \cdot c)), (((\operatorname{real\text{-}part}(\mathit{gouy}) \cdot a) + (\operatorname{imag\text{-}part}(\mathit{gouy}) \cdot b)) + (((1 - \operatorname{real\text{-}part}(\mathit{gouy})) - \operatorname{imag\text{-}part}(\mathit{gouy})) \cdot c)), (((\operatorname{real\text{-}part}(z) \cdot a) + (\operatorname{imag\text{-}part}(z) \cdot b)) + (((1 - \operatorname{real\text{-}part}(z)) - \operatorname{imag\text{-}part}(z)) \cdot c)))) \end{array}}}
\vnbprop{affine-holomorphic}{\vnbstmt{\begin{array}{@{}l@{}} \forall u.; \forall p, q, w \in \mathbb{C}, g \in (u \to \mathbb{C}). \\ \quad (\operatorname{is\text{-}open}(\operatorname{nf\text{-}metric\text{-}space}(\operatorname{cc\text{-}normed\text{-}field}), u) \wedge \forall z \in u.\; (g(z) \simeq (p + (q \cdot (z - w))))) \Rightarrow  \\ \quad \quad \operatorname{holomorphic\text{-}on}(u, g) \end{array}}}
\vnbprop{small-dyadic}{\vnbstmt{\begin{array}{@{}l@{}} \forall m \in \mathbb{R}, r. \\ \quad ((0 \leq m) \wedge \operatorname{pos\text{-}rr}(r)) \Rightarrow \exists n \in \mathbb{N}.\; ((m \cdot \operatorname{recip}(\operatorname{power}(2, n))) < r) \end{array}}}
\vnbprop{bary-side-near}{\vnbstmt{\begin{array}{@{}l@{}} \forall a, b, c, p, q, x, w \in \mathbb{C}, r, m \in \mathbb{R}, z \in \operatorname{trace}(\operatorname{seg\text{-}path}((((\operatorname{real\text{-}part}(p) \cdot a) + (\operatorname{imag\text{-}part}(p) \cdot b)) + (((1 - \operatorname{real\text{-}part}(p)) - \operatorname{imag\text{-}part}(p)) \cdot c)), (((\operatorname{real\text{-}part}(q) \cdot a) + (\operatorname{imag\text{-}part}(q) \cdot b)) + (((1 - \operatorname{real\text{-}part}(q)) - \operatorname{imag\text{-}part}(q)) \cdot c))), 0, 1). \\ \quad ((\operatorname{magnitude}((p - x)) \leq r) \wedge ((\operatorname{magnitude}((q - p)) \leq r) \wedge ((r \leq (6 \cdot m)) \wedge (\operatorname{magnitude}((x - w)) \leq (12 \cdot m))))) \Rightarrow  \\ \quad \quad \operatorname{magnitude}((z - (((\operatorname{real\text{-}part}(w) \cdot a) + (\operatorname{imag\text{-}part}(w) \cdot b)) + (((1 - \operatorname{real\text{-}part}(w)) - \operatorname{imag\text{-}part}(w)) \cdot c)))) \\ \quad \quad \leq\; (24 \cdot (\operatorname{magnitude}((a - c)) + \operatorname{magnitude}((b - c)))) \cdot m \end{array}}}
\vnbprop{tri-int-abs-bound-sides}{\vnbstmt{\begin{array}{@{}l@{}} \forall u, f.; \forall a, b, c \in \mathbb{C}, m \in \mathbb{R}. \\ \quad u \subseteq \mathbb{C} \wedge \\ \quad f \in (u \to \mathbb{C}) \wedge \\ \quad \forall y \in u. \\ \quad \quad \operatorname{is\text{-}continuous\text{-}at}(\operatorname{subspace\text{-}ms}(\operatorname{nf\text{-}metric\text{-}space}(\operatorname{cc\text{-}normed\text{-}field}), u), \\ \quad \quad \quad \operatorname{nf\text{-}metric\text{-}space}(\operatorname{cc\text{-}normed\text{-}field}), f, y) \wedge \\ \quad \operatorname{trace}(\operatorname{seg\text{-}path}(a, b), 0, 1) \subseteq u \wedge \\ \quad \operatorname{trace}(\operatorname{seg\text{-}path}(b, c), 0, 1) \subseteq u \wedge \\ \quad \operatorname{trace}(\operatorname{seg\text{-}path}(c, a), 0, 1) \subseteq u \wedge \\ \quad \forall z \in \operatorname{trace}(\operatorname{seg\text{-}path}(a, b), 0, 1).\; (\operatorname{magnitude}(f(z)) \leq m) \wedge \\ \quad \forall z \in \operatorname{trace}(\operatorname{seg\text{-}path}(b, c), 0, 1).\; (\operatorname{magnitude}(f(z)) \leq m) \wedge \\ \quad \forall z \in \operatorname{trace}(\operatorname{seg\text{-}path}(c, a), 0, 1).\; (\operatorname{magnitude}(f(z)) \leq m) \Rightarrow \\ \quad \quad \operatorname{magnitude}(\operatorname{tri\text{-}int}(f, a, b, c)) \leq (((\operatorname{magnitude}((b - a)) + \operatorname{magnitude}((c - b))) + \operatorname{magnitude}((a - c))) \cdot m) \end{array}}}
\vnbprop{goursat-triangle}{\vnbstmt{\begin{array}{@{}l@{}} \forall u, f.; \forall a, b, c \in u. \\ \quad (\operatorname{holomorphic\text{-}on}(u, f) \wedge (\operatorname{conv3}(a, b, c) \subseteq u)) \Rightarrow (\operatorname{tri\text{-}int}(f, a, b, c) = 0) \end{array}}}
\vnbprop{cc-diff-on-local-quadratic}{\vnbstmt{\begin{array}{@{}l@{}} \forall s.; \forall f \in (s \to \mathbb{C}), a \in s, l \in \mathbb{C}, m \in \mathbb{R}, r. \\ \quad \operatorname{is\text{-}open}(\operatorname{nf\text{-}metric\text{-}space}(\operatorname{cc\text{-}normed\text{-}field}), s) \wedge \\ \quad 0 \leq m \wedge \\ \quad \operatorname{pos\text{-}rr}(r) \wedge \\ \quad \forall x \in s. \\ \quad \quad (\operatorname{magnitude}((x - a)) < r) \Rightarrow  \\ \quad \quad \quad \operatorname{magnitude}(((f(x) - f(a)) - ((x - a) \cdot l))) \leq (m \cdot (\operatorname{magnitude}((x - a)) \cdot \operatorname{magnitude}((x - a)))) \Rightarrow \\ \quad \quad \operatorname{is\text{-}diff\text{-}on}(\operatorname{cc\text{-}normed\text{-}field}, s, f, a, l) \end{array}}}
\vnbprop{cauchy-convex-segments}{\vnbstmt{\begin{array}{@{}l@{}} \forall u, f.; \forall a, b, c \in u. \\ \quad (\operatorname{is\text{-}convex\text{-}cc}(u) \wedge \operatorname{holomorphic\text{-}on}(u, f)) \Rightarrow  \\ \quad \quad \operatorname{seg\text{-}int}(f, a, c) = (\operatorname{seg\text{-}int}(f, a, b) + \operatorname{seg\text{-}int}(f, b, c)) \end{array}}}
\vnbprop{seg-int-sum}{\vnbstmt{\begin{array}{@{}l@{}} \forall u, f, g.; \forall h \in (u \to \mathbb{C}), a, b \in \mathbb{C}. \\ \quad u \subseteq \mathbb{C} \wedge \\ \quad f \in (u \to \mathbb{C}) \wedge \\ \quad \forall y \in u. \\ \quad \quad \operatorname{is\text{-}continuous\text{-}at}(\operatorname{subspace\text{-}ms}(\operatorname{nf\text{-}metric\text{-}space}(\operatorname{cc\text{-}normed\text{-}field}), u), \\ \quad \quad \quad \operatorname{nf\text{-}metric\text{-}space}(\operatorname{cc\text{-}normed\text{-}field}), f, y) \wedge \\ \quad g \in (u \to \mathbb{C}) \wedge \\ \quad \forall y \in u. \\ \quad \quad \operatorname{is\text{-}continuous\text{-}at}(\operatorname{subspace\text{-}ms}(\operatorname{nf\text{-}metric\text{-}space}(\operatorname{cc\text{-}normed\text{-}field}), u), \\ \quad \quad \quad \operatorname{nf\text{-}metric\text{-}space}(\operatorname{cc\text{-}normed\text{-}field}), g, y) \wedge \\ \quad \forall z \in u.\; (h(z) \simeq (f(z) + g(z))) \wedge \\ \quad \operatorname{trace}(\operatorname{seg\text{-}path}(a, b), 0, 1) \subseteq u \Rightarrow \\ \quad \quad \operatorname{seg\text{-}int}(h, a, b) = (\operatorname{seg\text{-}int}(f, a, b) + \operatorname{seg\text{-}int}(g, a, b)) \end{array}}}
\vnbprop{seg-int-abs-bound}{\vnbstmt{\begin{array}{@{}l@{}} \forall u, f.; \forall a, b \in \mathbb{C}, m \in \mathbb{R}. \\ \quad u \subseteq \mathbb{C} \wedge \\ \quad f \in (u \to \mathbb{C}) \wedge \\ \quad \forall y \in u. \\ \quad \quad \operatorname{is\text{-}continuous\text{-}at}(\operatorname{subspace\text{-}ms}(\operatorname{nf\text{-}metric\text{-}space}(\operatorname{cc\text{-}normed\text{-}field}), u), \\ \quad \quad \quad \operatorname{nf\text{-}metric\text{-}space}(\operatorname{cc\text{-}normed\text{-}field}), f, y) \wedge \\ \quad \operatorname{trace}(\operatorname{seg\text{-}path}(a, b), 0, 1) \subseteq u \wedge \\ \quad \forall z \in \operatorname{trace}(\operatorname{seg\text{-}path}(a, b), 0, 1).\; (\operatorname{magnitude}(f(z)) \leq m) \Rightarrow \\ \quad \quad \operatorname{magnitude}(\operatorname{seg\text{-}int}(f, a, b)) \leq (\operatorname{magnitude}((b - a)) \cdot m) \end{array}}}
\vnbprop{seg-int-const}{\vnbstmt{\begin{array}{@{}l@{}} \forall u.; \forall p, w \in \mathbb{C}, g \in (u \to \mathbb{C}), a, b \in \mathbb{C}. \\ \quad \operatorname{is\text{-}open}(\operatorname{nf\text{-}metric\text{-}space}(\operatorname{cc\text{-}normed\text{-}field}), u) \wedge \\ \quad \forall z \in u.\; (g(z) \simeq (p + (0 \cdot (z - w)))) \wedge \\ \quad \operatorname{trace}(\operatorname{seg\text{-}path}(a, b), 0, 1) \subseteq u \Rightarrow \\ \quad \quad \operatorname{seg\text{-}int}(g, a, b) = (p \cdot (b - a)) \end{array}}}
\vnbprop{cauchy-convex-primitive}{\vnbstmt{\begin{array}{@{}l@{}} \forall u, f.; \forall o \in u. \\ \quad (\operatorname{is\text{-}convex\text{-}cc}(u) \wedge \operatorname{holomorphic\text{-}on}(u, f)) \Rightarrow  \\ \quad \quad \exists g \in (u \to \mathbb{C}). \\ \quad \quad \quad (\operatorname{holomorphic\text{-}on}(u, g) \wedge \forall z \in u.\; \operatorname{is\text{-}diff\text{-}on}(\operatorname{cc\text{-}normed\text{-}field}, u, g, z, f(z))) \end{array}}}
\vnbprop{cauchy-convex-closed}{\vnbstmt{\begin{array}{@{}l@{}} \forall u, f, m, c, a, b. \\ \quad \operatorname{is\text{-}convex\text{-}cc}(u) \wedge \\ \quad \operatorname{holomorphic\text{-}on}(u, f) \wedge \\ \quad \operatorname{is\text{-}road}(m, c, a, b) \wedge \\ \quad \operatorname{trace}(m, a, b) \subseteq u \wedge \\ \quad m(b) = m(a) \Rightarrow \\ \quad \quad \operatorname{line\text{-}int}(f, m, c, a, b) = 0 \end{array}}}

\vnbbookfile{theorem-library/goursat-exceptional}{theorem-library/goursat-exceptional}
\noindent{\small goursat-exceptional.scm -- CAUCHY-GOURSAT WITH ONE EXCEPTIONAL POINT (complex-analysis.pdf ch. 3, s. 3.2: Prop 3.6, Cor 3.7, Prop 3.8, Cor 3.9 in the notes' form: f CONTINUOUS on the open set U and HOLOMORPHIC on U \textbackslash{} \{x\} for some x in U).  Batch 35, 2026-09-26.}\par
\vnbprop{seg-int-congruence}{\vnbstmt{\begin{array}{@{}l@{}} \forall u, f, w.; \forall g \in (w \to \mathbb{C}), a, b \in \mathbb{C}. \\ \quad u \subseteq \mathbb{C} \wedge \\ \quad f \in (u \to \mathbb{C}) \wedge \\ \quad \forall y \in u. \\ \quad \quad \operatorname{is\text{-}continuous\text{-}at}(\operatorname{subspace\text{-}ms}(\operatorname{nf\text{-}metric\text{-}space}(\operatorname{cc\text{-}normed\text{-}field}), u), \\ \quad \quad \quad \operatorname{nf\text{-}metric\text{-}space}(\operatorname{cc\text{-}normed\text{-}field}), f, y) \wedge \\ \quad \operatorname{trace}(\operatorname{seg\text{-}path}(a, b), 0, 1) \subseteq u \wedge \\ \quad \operatorname{trace}(\operatorname{seg\text{-}path}(a, b), 0, 1) \subseteq w \wedge \\ \quad \forall z \in \operatorname{trace}(\operatorname{seg\text{-}path}(a, b), 0, 1).\; (g(z) \simeq f(z)) \Rightarrow \\ \quad \quad \operatorname{seg\text{-}int}(g, a, b) = \operatorname{seg\text{-}int}(f, a, b) \end{array}}}
\vnbprop{tri-int-restrict}{\vnbstmt{\begin{array}{@{}l@{}} \forall u, f, w.; \forall a, b, c \in \mathbb{C}. \\ \quad u \subseteq \mathbb{C} \wedge \\ \quad f \in (u \to \mathbb{C}) \wedge \\ \quad \forall y \in u. \\ \quad \quad \operatorname{is\text{-}continuous\text{-}at}(\operatorname{subspace\text{-}ms}(\operatorname{nf\text{-}metric\text{-}space}(\operatorname{cc\text{-}normed\text{-}field}), u), \\ \quad \quad \quad \operatorname{nf\text{-}metric\text{-}space}(\operatorname{cc\text{-}normed\text{-}field}), f, y) \wedge \\ \quad w \subseteq u \wedge \\ \quad \operatorname{conv3}(a, b, c) \subseteq w \Rightarrow \\ \quad \quad \operatorname{tri\text{-}int}(\operatorname{restrict}(f, w), a, b, c) = \operatorname{tri\text{-}int}(f, a, b, c) \end{array}}}
\vnbprop{tri-int-fan}{\vnbstmt{\begin{array}{@{}l@{}} \forall u, f.; \forall a, b, c \in \mathbb{C}, x \in \operatorname{conv3}(a, b, c). \\ \quad u \subseteq \mathbb{C} \wedge \\ \quad f \in (u \to \mathbb{C}) \wedge \\ \quad \forall y \in u. \\ \quad \quad \operatorname{is\text{-}continuous\text{-}at}(\operatorname{subspace\text{-}ms}(\operatorname{nf\text{-}metric\text{-}space}(\operatorname{cc\text{-}normed\text{-}field}), u), \\ \quad \quad \quad \operatorname{nf\text{-}metric\text{-}space}(\operatorname{cc\text{-}normed\text{-}field}), f, y) \wedge \\ \quad \operatorname{conv3}(a, b, c) \subseteq u \Rightarrow \\ \quad \quad \operatorname{tri\text{-}int}(f, a, b, c) = ((\operatorname{tri\text{-}int}(f, a, b, x) + \operatorname{tri\text{-}int}(f, b, c, x)) + \operatorname{tri\text{-}int}(f, c, a, x)) \end{array}}}
\vnbprop{tri-int-collinear}{\vnbstmt{\begin{array}{@{}l@{}} \forall u, f.; \forall a, b, c \in \mathbb{C}. \\ \quad u \subseteq \mathbb{C} \wedge \\ \quad f \in (u \to \mathbb{C}) \wedge \\ \quad \forall y \in u. \\ \quad \quad \operatorname{is\text{-}continuous\text{-}at}(\operatorname{subspace\text{-}ms}(\operatorname{nf\text{-}metric\text{-}space}(\operatorname{cc\text{-}normed\text{-}field}), u), \\ \quad \quad \quad \operatorname{nf\text{-}metric\text{-}space}(\operatorname{cc\text{-}normed\text{-}field}), f, y) \wedge \\ \quad \operatorname{conv3}(a, b, c) \subseteq u \wedge \\ \quad c \in \operatorname{trace}(\operatorname{seg\text{-}path}(a, b), 0, 1) \vee \\ \quad a \in \operatorname{trace}(\operatorname{seg\text{-}path}(b, c), 0, 1) \vee \\ \quad b \in \operatorname{trace}(\operatorname{seg\text{-}path}(c, a), 0, 1) \Rightarrow \\ \quad \quad \operatorname{tri\text{-}int}(f, a, b, c) = 0 \end{array}}}
\vnbprop{cc-cross-coords}{\vnbstmt{\begin{array}{@{}l@{}} \forall p, q \in \mathbb{C}. \\ \quad \operatorname{imag\text{-}part}((p \cdot \operatorname{conjugate}(q))) = ((\operatorname{imag\text{-}part}(p) \cdot \operatorname{real\text{-}part}(q)) - (\operatorname{real\text{-}part}(p) \cdot \operatorname{imag\text{-}part}(q))) \end{array}}}
\vnbprop{cc-affinely-independent}{\vnbstmt{\begin{array}{@{}l@{}} \forall a, b, c \in \mathbb{C}, p, q \in \mathbb{R}. \\ \quad (\neg (\operatorname{imag\text{-}part}(((b - a) \cdot \operatorname{conjugate}((c - a)))) = 0) \wedge (((p \cdot (b - a)) + (q \cdot (c - a))) = 0)) \Rightarrow  \\ \quad \quad (p = 0) \wedge (q = 0) \end{array}}}
\vnbprop{conv3-avoids-vertex}{\vnbstmt{\begin{array}{@{}l@{}} \forall a, b, c \in \mathbb{C}, t \in \mathbb{R}. \\ \quad (\neg (\operatorname{imag\text{-}part}(((b - a) \cdot \operatorname{conjugate}((c - a)))) = 0) \wedge ((0 < t) \wedge (t \leq 1))) \Rightarrow  \\ \quad \quad (\neg (a \in \operatorname{conv3}((a + (t \cdot (b - a))), b, c)) \wedge \neg (a \in \operatorname{conv3}((a + (t \cdot (b - a))), c, (a + (t \cdot (c - a)))))) \end{array}}}
\vnbprop{conv3-near-vertex}{\vnbstmt{\begin{array}{@{}l@{}} \forall a, b, c \in \mathbb{C}, z \in \operatorname{conv3}(a, b, c). \\ \quad \operatorname{magnitude}((z - a)) \leq (\operatorname{magnitude}((b - a)) + \operatorname{magnitude}((c - a))) \end{array}}}
\vnbprop{collinear-cases}{\vnbstmt{\begin{array}{@{}l@{}} \forall a, b, c \in \mathbb{C}. \\ \quad (\operatorname{imag\text{-}part}(((b - a) \cdot \operatorname{conjugate}((c - a)))) = 0) \Rightarrow  \\ \quad \quad c \in \operatorname{trace}(\operatorname{seg\text{-}path}(a, b), 0, 1) \vee \\ \quad \quad a \in \operatorname{trace}(\operatorname{seg\text{-}path}(b, c), 0, 1) \vee \\ \quad \quad b \in \operatorname{trace}(\operatorname{seg\text{-}path}(c, a), 0, 1) \end{array}}}
\vnbprop{conv3-vertex-shrink}{\vnbstmt{\begin{array}{@{}l@{}} \forall a, b, c \in \mathbb{C}, t \in \operatorname{ccint}(0, 1). \\ \quad \operatorname{conv3}(a, (a + (t \cdot (b - a))), (a + (t \cdot (c - a)))) \subseteq \operatorname{conv3}(a, b, c) \wedge \\ \quad \operatorname{conv3}((a + (t \cdot (b - a))), b, c) \subseteq \operatorname{conv3}(a, b, c) \wedge \\ \quad \operatorname{conv3}((a + (t \cdot (b - a))), c, (a + (t \cdot (c - a)))) \subseteq \operatorname{conv3}(a, b, c) \end{array}}}
\vnbprop{tri-int-vertex-split}{\vnbstmt{\begin{array}{@{}l@{}} \forall u, f.; \forall a, b, c \in \mathbb{C}, t \in \operatorname{ccint}(0, 1). \\ \quad u \subseteq \mathbb{C} \wedge \\ \quad f \in (u \to \mathbb{C}) \wedge \\ \quad \forall y \in u. \\ \quad \quad \operatorname{is\text{-}continuous\text{-}at}(\operatorname{subspace\text{-}ms}(\operatorname{nf\text{-}metric\text{-}space}(\operatorname{cc\text{-}normed\text{-}field}), u), \\ \quad \quad \quad \operatorname{nf\text{-}metric\text{-}space}(\operatorname{cc\text{-}normed\text{-}field}), f, y) \wedge \\ \quad \operatorname{conv3}(a, b, c) \subseteq u \Rightarrow \\ \quad \quad \operatorname{tri\text{-}int}(f, a, b, c) \\ \quad \quad =\; (\operatorname{tri\text{-}int}(f, a, (a + (t \cdot (b - a))), (a + (t \cdot (c - a)))) + \operatorname{tri\text{-}int}(f, (a + (t \cdot (b - a))), b, c)) + \operatorname{tri\text{-}int}(f, (a + (t \cdot (b - a))), c, (a + (t \cdot (c - a)))) \end{array}}}
\vnbprop{goursat-vertex}{\vnbstmt{\begin{array}{@{}l@{}} \forall u, f, x.; \forall b, c \in u. \\ \quad \operatorname{is\text{-}open}(\operatorname{nf\text{-}metric\text{-}space}(\operatorname{cc\text{-}normed\text{-}field}), u) \wedge \\ \quad f \in (u \to \mathbb{C}) \wedge \\ \quad \forall y \in u. \\ \quad \quad \operatorname{is\text{-}continuous\text{-}at}(\operatorname{subspace\text{-}ms}(\operatorname{nf\text{-}metric\text{-}space}(\operatorname{cc\text{-}normed\text{-}field}), u), \\ \quad \quad \quad \operatorname{nf\text{-}metric\text{-}space}(\operatorname{cc\text{-}normed\text{-}field}), f, y) \wedge \\ \quad x \in u \wedge \\ \quad \operatorname{holomorphic\text{-}on}(\operatorname{difference}(u, \operatorname{singleton}(x)), \operatorname{restrict}(f, \operatorname{difference}(u, \operatorname{singleton}(x)))) \wedge \\ \quad \operatorname{conv3}(x, b, c) \subseteq u \Rightarrow \\ \quad \quad \operatorname{tri\text{-}int}(f, x, b, c) = 0 \end{array}}}
\vnbprop{goursat-triangle-except}{\vnbstmt{\begin{array}{@{}l@{}} \forall u, f, x.; \forall a, b, c \in u. \\ \quad \operatorname{is\text{-}open}(\operatorname{nf\text{-}metric\text{-}space}(\operatorname{cc\text{-}normed\text{-}field}), u) \wedge \\ \quad f \in (u \to \mathbb{C}) \wedge \\ \quad \forall y \in u. \\ \quad \quad \operatorname{is\text{-}continuous\text{-}at}(\operatorname{subspace\text{-}ms}(\operatorname{nf\text{-}metric\text{-}space}(\operatorname{cc\text{-}normed\text{-}field}), u), \\ \quad \quad \quad \operatorname{nf\text{-}metric\text{-}space}(\operatorname{cc\text{-}normed\text{-}field}), f, y) \wedge \\ \quad x \in u \wedge \\ \quad \operatorname{holomorphic\text{-}on}(\operatorname{difference}(u, \operatorname{singleton}(x)), \operatorname{restrict}(f, \operatorname{difference}(u, \operatorname{singleton}(x)))) \wedge \\ \quad \operatorname{conv3}(a, b, c) \subseteq u \Rightarrow \\ \quad \quad \operatorname{tri\text{-}int}(f, a, b, c) = 0 \end{array}}}
\vnbprop{cc-diff-on-of-eps-bound}{\vnbstmt{\begin{array}{@{}l@{}} \forall s.; \forall g \in (s \to \mathbb{C}), a \in s, l \in \mathbb{C}. \\ \quad \operatorname{is\text{-}open}(\operatorname{nf\text{-}metric\text{-}space}(\operatorname{cc\text{-}normed\text{-}field}), s) \wedge \\ \quad \forall e. \\ \quad \quad \operatorname{pos\text{-}rr}(e) \Rightarrow  \\ \quad \quad \quad \exists d. \\ \quad \quad \quad \quad \operatorname{pos\text{-}rr}(d) \wedge \\ \quad \quad \quad \quad \forall y \in s. \\ \quad \quad \quad \quad \quad (\operatorname{magnitude}((y - a)) < d) \Rightarrow  \\ \quad \quad \quad \quad \quad \quad \operatorname{magnitude}(((g(y) - g(a)) - ((y - a) \cdot l))) \leq (e \cdot \operatorname{magnitude}((y - a))) \Rightarrow \\ \quad \quad \operatorname{is\text{-}diff\text{-}on}(\operatorname{cc\text{-}normed\text{-}field}, s, g, a, l) \end{array}}}
\vnbprop{cc-continuous-at-add-const}{\vnbstmt{\begin{array}{@{}l@{}} \forall u.; \forall f, g \in (u \to \mathbb{C}), k \in \mathbb{C}, y \in u. \\ \quad u \subseteq \mathbb{C} \wedge \\ \quad \forall z \in u.\; (g(z) \simeq (f(z) + k)) \wedge \\ \quad \forall y \in u. \\ \quad \quad \operatorname{is\text{-}continuous\text{-}at}(\operatorname{subspace\text{-}ms}(\operatorname{nf\text{-}metric\text{-}space}(\operatorname{cc\text{-}normed\text{-}field}), u), \\ \quad \quad \quad \operatorname{nf\text{-}metric\text{-}space}(\operatorname{cc\text{-}normed\text{-}field}), f, y) \Rightarrow \\ \quad \quad \operatorname{is\text{-}continuous\text{-}at}(\operatorname{subspace\text{-}ms}(\operatorname{nf\text{-}metric\text{-}space}(\operatorname{cc\text{-}normed\text{-}field}), u), \\ \quad \quad \quad \operatorname{nf\text{-}metric\text{-}space}(\operatorname{cc\text{-}normed\text{-}field}), g, y) \end{array}}}
\vnbprop{cauchy-convex-segments-except}{\vnbstmt{\begin{array}{@{}l@{}} \forall u, f, x.; \forall a, b, c \in u. \\ \quad \operatorname{is\text{-}convex\text{-}cc}(u) \wedge \\ \quad \operatorname{is\text{-}open}(\operatorname{nf\text{-}metric\text{-}space}(\operatorname{cc\text{-}normed\text{-}field}), u) \wedge \\ \quad f \in (u \to \mathbb{C}) \wedge \\ \quad \forall y \in u. \\ \quad \quad \operatorname{is\text{-}continuous\text{-}at}(\operatorname{subspace\text{-}ms}(\operatorname{nf\text{-}metric\text{-}space}(\operatorname{cc\text{-}normed\text{-}field}), u), \\ \quad \quad \quad \operatorname{nf\text{-}metric\text{-}space}(\operatorname{cc\text{-}normed\text{-}field}), f, y) \wedge \\ \quad x \in u \wedge \\ \quad \operatorname{holomorphic\text{-}on}(\operatorname{difference}(u, \operatorname{singleton}(x)), \operatorname{restrict}(f, \operatorname{difference}(u, \operatorname{singleton}(x)))) \Rightarrow \\ \quad \quad \operatorname{seg\text{-}int}(f, a, c) = (\operatorname{seg\text{-}int}(f, a, b) + \operatorname{seg\text{-}int}(f, b, c)) \end{array}}}
\vnbprop{cauchy-convex-primitive-except}{\vnbstmt{\begin{array}{@{}l@{}} \forall u, f, x.; \forall o \in u. \\ \quad \operatorname{is\text{-}convex\text{-}cc}(u) \wedge \\ \quad \operatorname{is\text{-}open}(\operatorname{nf\text{-}metric\text{-}space}(\operatorname{cc\text{-}normed\text{-}field}), u) \wedge \\ \quad f \in (u \to \mathbb{C}) \wedge \\ \quad \forall y \in u. \\ \quad \quad \operatorname{is\text{-}continuous\text{-}at}(\operatorname{subspace\text{-}ms}(\operatorname{nf\text{-}metric\text{-}space}(\operatorname{cc\text{-}normed\text{-}field}), u), \\ \quad \quad \quad \operatorname{nf\text{-}metric\text{-}space}(\operatorname{cc\text{-}normed\text{-}field}), f, y) \wedge \\ \quad x \in u \wedge \\ \quad \operatorname{holomorphic\text{-}on}(\operatorname{difference}(u, \operatorname{singleton}(x)), \operatorname{restrict}(f, \operatorname{difference}(u, \operatorname{singleton}(x)))) \Rightarrow \\ \quad \quad \exists g \in (u \to \mathbb{C}). \\ \quad \quad \quad (\operatorname{holomorphic\text{-}on}(u, g) \wedge \forall z \in u.\; \operatorname{is\text{-}diff\text{-}on}(\operatorname{cc\text{-}normed\text{-}field}, u, g, z, f(z))) \end{array}}}
\vnbprop{cauchy-convex-closed-except}{\vnbstmt{\begin{array}{@{}l@{}} \forall u, f, x, m, c, a, b. \\ \quad \operatorname{is\text{-}convex\text{-}cc}(u) \wedge \\ \quad \operatorname{is\text{-}open}(\operatorname{nf\text{-}metric\text{-}space}(\operatorname{cc\text{-}normed\text{-}field}), u) \wedge \\ \quad f \in (u \to \mathbb{C}) \wedge \\ \quad \forall y \in u. \\ \quad \quad \operatorname{is\text{-}continuous\text{-}at}(\operatorname{subspace\text{-}ms}(\operatorname{nf\text{-}metric\text{-}space}(\operatorname{cc\text{-}normed\text{-}field}), u), \\ \quad \quad \quad \operatorname{nf\text{-}metric\text{-}space}(\operatorname{cc\text{-}normed\text{-}field}), f, y) \wedge \\ \quad x \in u \wedge \\ \quad \operatorname{holomorphic\text{-}on}(\operatorname{difference}(u, \operatorname{singleton}(x)), \operatorname{restrict}(f, \operatorname{difference}(u, \operatorname{singleton}(x)))) \wedge \\ \quad \operatorname{is\text{-}road}(m, c, a, b) \wedge \\ \quad \operatorname{trace}(m, a, b) \subseteq u \wedge \\ \quad m(b) = m(a) \Rightarrow \\ \quad \quad \operatorname{line\text{-}int}(f, m, c, a, b) = 0 \end{array}}}

\vnbbookfile{theorem-library/road-juxtaposition-laws}{theorem-library/road-juxtaposition-laws}
\noindent{\small road-juxtaposition-laws.scm -- THE JUXTAPOSITION OF ROADS (Dieudonne 9.6, the notes' Proposition 3.1), THE TRIANGULAR ROAD, AND THE LENGTH OF A ROAD (the notes' (53)). Batch 37, 2026-09-26.  The vocabulary is structure-library/road-juxtaposition.scm (JUXTA, ROAD-LENGTH, TRI-ROAD, TRI-ROAD-DERIV).}\par
\vnbprop{juxta-unfold}{\vnbstmt{\begin{array}{@{}l@{}} \forall p, q, a, b, c. \\ \quad (\operatorname{juxta}(p, q, a, b, c) \simeq (\lambda t \in \operatorname{ccint}(a, c).\; \left(\text{if }(t \leq b)\text{ then }p(t)\text{ else }q(t)\right))) \end{array}}}
\vnbprop{road-length-unfold}{\vnbstmt{\begin{array}{@{}l@{}} \forall d, a, b. \\ \quad (\operatorname{road\text{-}length}(d, a, b) \simeq \int_{a}^{b} (\lambda t \in \operatorname{ccint}(a, b).\; \operatorname{magnitude}(d(t)))) \end{array}}}
\vnbprop{tri-road-unfold}{\vnbstmt{\begin{array}{@{}l@{}} \forall a, b, c. \\ \quad (\operatorname{tri\text{-}road}(a, b, c) \simeq \operatorname{juxta}(\operatorname{juxta}(\operatorname{seg\text{-}path}(a, b), (\lambda s \in \operatorname{ccint}(1, 2).\; \operatorname{seg\text{-}path}(b, c)((s - 1))), 0, 1, 2), (\lambda s \in \operatorname{ccint}(2, 3).\; \operatorname{seg\text{-}path}(c, a)((s - 2))), 0, 2, 3)) \end{array}}}
\vnbprop{tri-road-deriv-unfold}{\vnbstmt{\begin{array}{@{}l@{}} \forall a, b, c. \\ \quad (\operatorname{tri\text{-}road\text{-}deriv}(a, b, c) \simeq \operatorname{juxta}(\operatorname{juxta}(\operatorname{seg\text{-}deriv}(a, b), (\lambda s \in \operatorname{ccint}(1, 2).\; \operatorname{seg\text{-}deriv}(b, c)((s - 1))), 0, 1, 2), (\lambda s \in \operatorname{ccint}(2, 3).\; \operatorname{seg\text{-}deriv}(c, a)((s - 2))), 0, 2, 3)) \end{array}}}
\vnbprop{juxta-apply-left}{\vnbstmt{\begin{array}{@{}l@{}} \forall p, q.; \forall a, b, c \in \mathbb{R}, t \in \operatorname{ccint}(a, b). \\ \quad (b \leq c) \Rightarrow (\operatorname{juxta}(p, q, a, b, c)(t) \simeq p(t)) \end{array}}}
\vnbprop{juxta-apply-right}{\vnbstmt{\begin{array}{@{}l@{}} \forall p, q.; \forall a, b, c \in \mathbb{R}, t \in \operatorname{ccint}(b, c). \\ \quad ((a \leq b) \wedge (b < t)) \Rightarrow (\operatorname{juxta}(p, q, a, b, c)(t) \simeq q(t)) \end{array}}}
\vnbprop{juxta-apply-right-closed}{\vnbstmt{\begin{array}{@{}l@{}} \forall p, q.; \forall a, b, c \in \mathbb{R}, t \in \operatorname{ccint}(b, c). \\ \quad ((a \leq b) \wedge (p(b) = q(b))) \Rightarrow (\operatorname{juxta}(p, q, a, b, c)(t) \simeq q(t)) \end{array}}}
\vnbprop{juxta-in-fun}{\vnbstmt{\begin{array}{@{}l@{}} \forall p, q.; \forall a, b, c \in \mathbb{R}, d. \\ \quad ((a \leq b) \wedge ((b \leq c) \wedge ((p \in (\operatorname{ccint}(a, b) \to d)) \wedge (q \in (\operatorname{ccint}(b, c) \to d))))) \Rightarrow  \\ \quad \quad \operatorname{juxta}(p, q, a, b, c) \in (\operatorname{ccint}(a, c) \to d) \end{array}}}
\vnbprop{juxta-trace-subset}{\vnbstmt{\begin{array}{@{}l@{}} \forall p, q.; \forall a, b, c \in \mathbb{R}. \\ \quad ((a \leq b) \wedge ((b \leq c) \wedge ((p \in (\operatorname{ccint}(a, b) \to \mathbb{C})) \wedge (q \in (\operatorname{ccint}(b, c) \to \mathbb{C}))))) \Rightarrow  \\ \quad \quad \operatorname{trace}(\operatorname{juxta}(p, q, a, b, c), a, c) \subseteq (\operatorname{trace}(p, a, b) \cup \operatorname{trace}(q, b, c)) \end{array}}}
\vnbprop{juxta-trace}{\vnbstmt{\begin{array}{@{}l@{}} \forall p, q.; \forall a, b, c \in \mathbb{R}. \\ \quad ((a \leq b) \wedge ((b \leq c) \wedge ((p \in (\operatorname{ccint}(a, b) \to \mathbb{C})) \wedge ((q \in (\operatorname{ccint}(b, c) \to \mathbb{C})) \wedge (p(b) = q(b)))))) \Rightarrow  \\ \quad \quad \operatorname{trace}(\operatorname{juxta}(p, q, a, b, c), a, c) = (\operatorname{trace}(p, a, b) \cup \operatorname{trace}(q, b, c)) \end{array}}}
\vnbprop{juxta-is-road}{\vnbstmt{\begin{array}{@{}l@{}} \forall m_{1}, d, m_{2}, f, a, b, c. \\ \quad (\operatorname{is\text{-}road}(m_{1}, d, a, b) \wedge (\operatorname{is\text{-}road}(m_{2}, f, b, c) \wedge (\operatorname{m1}(b) = \operatorname{m2}(b)))) \Rightarrow  \\ \quad \quad \operatorname{is\text{-}road}(\operatorname{juxta}(m_{1}, m_{2}, a, b, c), \operatorname{juxta}(d, f, a, b, c), a, c) \end{array}}}
\vnbprop{line-int-interior-congruence}{\vnbstmt{\begin{array}{@{}l@{}} \forall u, f, m, c, a, b.; \forall d, g \in (\operatorname{ccint}(a, b) \to \mathbb{C}). \\ \quad u \subseteq \mathbb{C} \wedge \\ \quad f \in (u \to \mathbb{C}) \wedge \\ \quad \forall y \in u. \\ \quad \quad \operatorname{is\text{-}continuous\text{-}at}(\operatorname{subspace\text{-}ms}(\operatorname{nf\text{-}metric\text{-}space}(\operatorname{cc\text{-}normed\text{-}field}), u), \\ \quad \quad \quad \operatorname{nf\text{-}metric\text{-}space}(\operatorname{cc\text{-}normed\text{-}field}), f, y) \wedge \\ \quad \operatorname{is\text{-}road}(m, c, a, b) \wedge \\ \quad \operatorname{trace}(m, a, b) \subseteq u \wedge \\ \quad \forall t \in \operatorname{ccint}(a, b).\; (d(t) \simeq m(t)) \wedge \\ \quad \forall t \in \operatorname{ooint}(a, b).\; (g(t) \simeq c(t)) \Rightarrow \\ \quad \quad \operatorname{line\text{-}int}(f, d, g, a, b) = \operatorname{line\text{-}int}(f, m, c, a, b) \end{array}}}
\vnbprop{line-int-juxta}{\vnbstmt{\begin{array}{@{}l@{}} \forall u, f, m_{1}, d, m_{2}, g, a, b, c. \\ \quad u \subseteq \mathbb{C} \wedge \\ \quad f \in (u \to \mathbb{C}) \wedge \\ \quad \forall y \in u. \\ \quad \quad \operatorname{is\text{-}continuous\text{-}at}(\operatorname{subspace\text{-}ms}(\operatorname{nf\text{-}metric\text{-}space}(\operatorname{cc\text{-}normed\text{-}field}), u), \\ \quad \quad \quad \operatorname{nf\text{-}metric\text{-}space}(\operatorname{cc\text{-}normed\text{-}field}), f, y) \wedge \\ \quad \operatorname{is\text{-}road}(m_{1}, d, a, b) \wedge \\ \quad \operatorname{is\text{-}road}(m_{2}, g, b, c) \wedge \\ \quad \operatorname{m1}(b) = \operatorname{m2}(b) \wedge \\ \quad \operatorname{trace}(m_{1}, a, b) \subseteq u \wedge \\ \quad \operatorname{trace}(m_{2}, b, c) \subseteq u \Rightarrow \\ \quad \quad \operatorname{line\text{-}int}(f, \operatorname{juxta}(m_{1}, m_{2}, a, b, c), \operatorname{juxta}(d, g, a, b, c), a, c) \\ \quad \quad =\; \operatorname{line\text{-}int}(f, m_{1}, d, a, b) + \operatorname{line\text{-}int}(f, m_{2}, g, b, c) \end{array}}}
\vnbprop{line-int-concat}{\vnbstmt{\begin{array}{@{}l@{}} \forall u, f, m, g, o, h, a, b, c, d. \\ \quad u \subseteq \mathbb{C} \wedge \\ \quad f \in (u \to \mathbb{C}) \wedge \\ \quad \forall y \in u. \\ \quad \quad \operatorname{is\text{-}continuous\text{-}at}(\operatorname{subspace\text{-}ms}(\operatorname{nf\text{-}metric\text{-}space}(\operatorname{cc\text{-}normed\text{-}field}), u), \\ \quad \quad \quad \operatorname{nf\text{-}metric\text{-}space}(\operatorname{cc\text{-}normed\text{-}field}), f, y) \wedge \\ \quad \operatorname{is\text{-}road}(m, g, a, b) \wedge \\ \quad \operatorname{is\text{-}road}(o, h, c, d) \wedge \\ \quad m(b) = o(c) \wedge \\ \quad \operatorname{trace}(m, a, b) \subseteq u \wedge \\ \quad \operatorname{trace}(o, c, d) \subseteq u \Rightarrow \\ \quad \quad \operatorname{line\text{-}int}(f, \\ \quad \quad \quad \operatorname{juxta}(m, (\lambda s \in \operatorname{ccint}(b, (b + (d - c))).\; o((c + (s - b)))), a, b, (b + (d - c))), \\ \quad \quad \quad \operatorname{juxta}(g, (\lambda s \in \operatorname{ccint}(b, (b + (d - c))).\; h((c + (s - b)))), a, b, (b + (d - c))), \\ \quad \quad \quad a, (b + (d - c))) \\ \quad \quad =\; \operatorname{line\text{-}int}(f, m, g, a, b) + \operatorname{line\text{-}int}(f, o, h, c, d) \end{array}}}
\vnbprop{road-length-regulated}{\vnbstmt{\begin{array}{@{}l@{}} \forall m, c, a, b. \\ \quad \operatorname{is\text{-}road}(m, c, a, b) \Rightarrow  \\ \quad \quad \operatorname{is\text{-}regulated\text{-}on}((\lambda t \in \operatorname{ccint}(a, b).\; \operatorname{magnitude}(c(t))), a, b) \end{array}}}
\vnbprop{road-length-primitive}{\vnbstmt{\begin{array}{@{}l@{}} \forall m, c, a, b. \\ \quad \operatorname{is\text{-}road}(m, c, a, b) \Rightarrow  \\ \quad \quad \exists g.\; \operatorname{is\text{-}primitive}(g, (\lambda t \in \operatorname{ccint}(a, b).\; \operatorname{magnitude}(c(t))), a, b) \end{array}}}
\vnbprop{road-length-in-rr}{\vnbstmt{\forall m, c, a, b.\; (\operatorname{is\text{-}road}(m, c, a, b) \Rightarrow (\operatorname{road\text{-}length}(c, a, b) \in \mathbb{R}))}}
\vnbprop{road-length-nonneg}{\vnbstmt{\forall m, c, a, b.\; (\operatorname{is\text{-}road}(m, c, a, b) \Rightarrow (0 \leq \operatorname{road\text{-}length}(c, a, b)))}}
\vnbprop{pw-int-interior-congruence}{\vnbstmt{\begin{array}{@{}l@{}} \forall f, i, \mathit{ppsi}, a, b. \\ \quad \operatorname{is\text{-}primitive}(f, i, a, b) \wedge \\ \quad \mathit{ppsi} \in (\operatorname{ccint}(a, b) \to \mathbb{R}) \wedge \\ \quad \forall y \in \operatorname{ooint}(a, b).\; (i(y) \simeq \operatorname{ppsi}(y)) \Rightarrow \\ \quad \quad \int_{a}^{b} i = \int_{a}^{b} \mathit{ppsi} \end{array}}}
\vnbprop{road-length-juxta}{\vnbstmt{\begin{array}{@{}l@{}} \forall m_{1}, d, m_{2}, f, a, b, c. \\ \quad (\operatorname{is\text{-}road}(m_{1}, d, a, b) \wedge (\operatorname{is\text{-}road}(m_{2}, f, b, c) \wedge (\operatorname{m1}(b) = \operatorname{m2}(b)))) \Rightarrow  \\ \quad \quad \operatorname{road\text{-}length}(\operatorname{juxta}(d, f, a, b, c), a, c) = (\operatorname{road\text{-}length}(d, a, b) + \operatorname{road\text{-}length}(f, b, c)) \end{array}}}
\vnbprop{seg-shift-is-road}{\vnbstmt{\begin{array}{@{}l@{}} \forall p, q \in \mathbb{C}, k, m \in \mathbb{R}. \\ \quad (m = (k + 1)) \Rightarrow  \\ \quad \quad \operatorname{is\text{-}road}((\lambda s \in \operatorname{ccint}(k, m).\; \operatorname{seg\text{-}path}(p, q)((s - k))), \\ \quad \quad \quad (\lambda s \in \operatorname{ccint}(k, m).\; \operatorname{seg\text{-}deriv}(p, q)((s - k))), k, m) \end{array}}}
\vnbprop{seg-shift-trace}{\vnbstmt{\begin{array}{@{}l@{}} \forall p, q \in \mathbb{C}, k, m \in \mathbb{R}. \\ \quad (m = (k + 1)) \Rightarrow  \\ \quad \quad \operatorname{trace}((\lambda s \in \operatorname{ccint}(k, m).\; \operatorname{seg\text{-}path}(p, q)((s - k))), k, m) = \operatorname{trace}(\operatorname{seg\text{-}path}(p, q), 0, 1) \end{array}}}
\vnbprop{road-length-segment}{\vnbstmt{\begin{array}{@{}l@{}} \forall p, q \in \mathbb{C}. \\ \quad \operatorname{road\text{-}length}(\operatorname{seg\text{-}deriv}(p, q), 0, 1) = \operatorname{magnitude}((q - p)) \end{array}}}
\vnbprop{seg-shift-length}{\vnbstmt{\begin{array}{@{}l@{}} \forall p, q \in \mathbb{C}, k, m \in \mathbb{R}. \\ \quad (m = (k + 1)) \Rightarrow  \\ \quad \quad \operatorname{road\text{-}length}((\lambda s \in \operatorname{ccint}(k, m).\; \operatorname{seg\text{-}deriv}(p, q)((s - k))), k, m) = \operatorname{magnitude}((q - p)) \end{array}}}
\vnbprop{seg-shift-line-int}{\vnbstmt{\begin{array}{@{}l@{}} \forall u, f.; \forall p, q \in \mathbb{C}, k, m \in \mathbb{R}. \\ \quad u \subseteq \mathbb{C} \wedge \\ \quad f \in (u \to \mathbb{C}) \wedge \\ \quad \forall y \in u. \\ \quad \quad \operatorname{is\text{-}continuous\text{-}at}(\operatorname{subspace\text{-}ms}(\operatorname{nf\text{-}metric\text{-}space}(\operatorname{cc\text{-}normed\text{-}field}), u), \\ \quad \quad \quad \operatorname{nf\text{-}metric\text{-}space}(\operatorname{cc\text{-}normed\text{-}field}), f, y) \wedge \\ \quad m = (k + 1) \wedge \\ \quad \operatorname{trace}(\operatorname{seg\text{-}path}(p, q), 0, 1) \subseteq u \Rightarrow \\ \quad \quad \operatorname{line\text{-}int}(f, (\lambda s \in \operatorname{ccint}(k, m).\; \operatorname{seg\text{-}path}(p, q)((s - k))), \\ \quad \quad \quad (\lambda s \in \operatorname{ccint}(k, m).\; \operatorname{seg\text{-}deriv}(p, q)((s - k))), k, m) \\ \quad \quad =\; \operatorname{line\text{-}int}(f, \operatorname{seg\text{-}path}(p, q), \operatorname{seg\text{-}deriv}(p, q), 0, 1) \end{array}}}
\vnbprop{tri-road-is-road}{\vnbstmt{\begin{array}{@{}l@{}} \forall a, b, c \in \mathbb{C}. \\ \quad \operatorname{is\text{-}road}(\operatorname{tri\text{-}road}(a, b, c), \operatorname{tri\text{-}road\text{-}deriv}(a, b, c), 0, 3) \end{array}}}
\vnbprop{tri-road-trace}{\vnbstmt{\begin{array}{@{}l@{}} \forall a, b, c \in \mathbb{C}. \\ \quad \operatorname{trace}(\operatorname{tri\text{-}road}(a, b, c), 0, 3) \\ \quad =\; ((\operatorname{trace}(\operatorname{seg\text{-}path}(a, b), 0, 1) \cup \operatorname{trace}(\operatorname{seg\text{-}path}(b, c), 0, 1)) \cup \operatorname{trace}(\operatorname{seg\text{-}path}(c, a), 0, 1)) \end{array}}}
\vnbprop{tri-int-is-line-int}{\vnbstmt{\begin{array}{@{}l@{}} \forall u, f.; \forall a, b, c \in \mathbb{C}. \\ \quad u \subseteq \mathbb{C} \wedge \\ \quad f \in (u \to \mathbb{C}) \wedge \\ \quad \forall y \in u. \\ \quad \quad \operatorname{is\text{-}continuous\text{-}at}(\operatorname{subspace\text{-}ms}(\operatorname{nf\text{-}metric\text{-}space}(\operatorname{cc\text{-}normed\text{-}field}), u), \\ \quad \quad \quad \operatorname{nf\text{-}metric\text{-}space}(\operatorname{cc\text{-}normed\text{-}field}), f, y) \wedge \\ \quad \operatorname{trace}(\operatorname{seg\text{-}path}(a, b), 0, 1) \subseteq u \wedge \\ \quad \operatorname{trace}(\operatorname{seg\text{-}path}(b, c), 0, 1) \subseteq u \wedge \\ \quad \operatorname{trace}(\operatorname{seg\text{-}path}(c, a), 0, 1) \subseteq u \Rightarrow \\ \quad \quad \operatorname{tri\text{-}int}(f, a, b, c) = \operatorname{line\text{-}int}(f, \operatorname{tri\text{-}road}(a, b, c), \operatorname{tri\text{-}road\text{-}deriv}(a, b, c), 0, 3) \end{array}}}
\vnbprop{road-length-tri}{\vnbstmt{\begin{array}{@{}l@{}} \forall a, b, c \in \mathbb{C}. \\ \quad \operatorname{road\text{-}length}(\operatorname{tri\text{-}road\text{-}deriv}(a, b, c), 0, 3) \\ \quad =\; (\operatorname{magnitude}((b - a)) + \operatorname{magnitude}((c - b))) + \operatorname{magnitude}((a - c)) \end{array}}}
\vnbprop{line-int-length-bound}{\vnbstmt{\begin{array}{@{}l@{}} \forall u, f, m, c, a, b.; \forall d \in \mathbb{R}. \\ \quad u \subseteq \mathbb{C} \wedge \\ \quad f \in (u \to \mathbb{C}) \wedge \\ \quad \forall y \in u. \\ \quad \quad \operatorname{is\text{-}continuous\text{-}at}(\operatorname{subspace\text{-}ms}(\operatorname{nf\text{-}metric\text{-}space}(\operatorname{cc\text{-}normed\text{-}field}), u), \\ \quad \quad \quad \operatorname{nf\text{-}metric\text{-}space}(\operatorname{cc\text{-}normed\text{-}field}), f, y) \wedge \\ \quad \operatorname{is\text{-}road}(m, c, a, b) \wedge \\ \quad \operatorname{trace}(m, a, b) \subseteq u \wedge \\ \quad \forall \mathit{bny} \in \operatorname{trace}(m, a, b).\; (\operatorname{magnitude}(f(\mathit{bny})) \leq d) \Rightarrow \\ \quad \quad \operatorname{magnitude}(\operatorname{line\text{-}int}(f, m, c, a, b)) \leq (\operatorname{road\text{-}length}(c, a, b) \cdot d) \end{array}}}
\vnbprop{road-length-affine-reparam}{\vnbstmt{\begin{array}{@{}l@{}} \forall a, b, c, d, l, m \in \mathbb{R}, f, g, k. \\ \quad 0 < l \wedge \\ \quad (m + (l \cdot c)) = a \wedge \\ \quad (m + (l \cdot d)) = b \wedge \\ \quad c < d \wedge \\ \quad \operatorname{is\text{-}road}(f, g, a, b) \wedge \\ \quad k \in (\operatorname{ccint}(c, d) \to \mathbb{C}) \wedge \\ \quad \forall z \in \operatorname{ccint}(c, d).\; (k(z) \simeq (l \cdot g((m + (l \cdot z))))) \Rightarrow \\ \quad \quad \operatorname{road\text{-}length}(k, c, d) = \operatorname{road\text{-}length}(g, a, b) \end{array}}}
\vnbprop{tri-int-is-line-int-conv3}{\vnbstmt{\begin{array}{@{}l@{}} \forall u, f.; \forall a, b, c \in \mathbb{C}. \\ \quad u \subseteq \mathbb{C} \wedge \\ \quad f \in (u \to \mathbb{C}) \wedge \\ \quad \forall y \in u. \\ \quad \quad \operatorname{is\text{-}continuous\text{-}at}(\operatorname{subspace\text{-}ms}(\operatorname{nf\text{-}metric\text{-}space}(\operatorname{cc\text{-}normed\text{-}field}), u), \\ \quad \quad \quad \operatorname{nf\text{-}metric\text{-}space}(\operatorname{cc\text{-}normed\text{-}field}), f, y) \wedge \\ \quad \operatorname{conv3}(a, b, c) \subseteq u \Rightarrow \\ \quad \quad \operatorname{tri\text{-}int}(f, a, b, c) = \operatorname{line\text{-}int}(f, \operatorname{tri\text{-}road}(a, b, c), \operatorname{tri\text{-}road\text{-}deriv}(a, b, c), 0, 3) \end{array}}}

\vnbbookfile{structure-library/cc-elementary}{structure-library/cc-elementary}
\noindent{\small cc-elementary.scm -- THE ELEMENTARY FUNCTIONS exp, cos, sin ON CC, DEFINED BY THEIR POWER SERIES (the notes 2.3.1, 2.3.2). DEFINITIONS ONLY; the laws are PROVEN in theorem-library/cc-exp.scm. Batch 31-A, 2026-09-25.  The user's decision of 2026-09-25: CC-EXP, CC-COS and}\par
\par\medskip\noindent\textbf{Definition} (\texttt{cc-exp}).\;$\exp z :=$
\begin{equation*}
\operatorname{cc\text{-}series\text{-}limit}((\lambda \operatorname{cek\_} \in \mathbb{N}.\; (\operatorname{recip}(\operatorname{factorial}(\operatorname{cek\_})) \cdot \operatorname{power}(z, \operatorname{cek\_}))))
\end{equation*}

\par\medskip\noindent\textbf{Definition} (\texttt{cc-cos}).\;$\cos z :=$
\begin{equation*}
\operatorname{cc\text{-}series\text{-}limit}((\lambda \operatorname{cek\_} \in \mathbb{N}.\; ((\operatorname{power}(-1, \operatorname{cek\_}) \cdot \operatorname{recip}(\operatorname{factorial}((2 \cdot \operatorname{cek\_})))) \cdot \operatorname{power}(z, (2 \cdot \operatorname{cek\_})))))
\end{equation*}

\par\medskip\noindent\textbf{Definition} (\texttt{cc-sin}).\;$\sin z :=$
\begin{equation*}
\operatorname{cc\text{-}series\text{-}limit}((\lambda \operatorname{cek\_} \in \mathbb{N}.\; ((\operatorname{power}(-1, \operatorname{cek\_}) \cdot \operatorname{recip}(\operatorname{factorial}(((2 \cdot \operatorname{cek\_}) + 1)))) \cdot \operatorname{power}(z, ((2 \cdot \operatorname{cek\_}) + 1)))))
\end{equation*}


\vnbbookfile{structure-library/ccr1109a6-vocab}{structure-library/ccr1109a6-vocab}
\noindent{\small ccr1109a6-vocab.scm -- wave of 2026-11-09/10 (agent CCR-1109-A6), wired 2026-11-10 from scratchpad/ccr1109a6/ccr1109a6-vocab.scm. Library: complex-analysis.}\par
\par\medskip\noindent\textbf{Definition} (\texttt{adjoin-char-value}).\;$\operatorname{adjoin\text{-}char\text{-}value}(\operatorname{a6vg\_}, \operatorname{a6vh\_}, \operatorname{a6vc\_}, \operatorname{a6vb\_}, \operatorname{a6vt\_}, \operatorname{a6vy\_}) :=$
\begin{equation*}
(\iota\, \operatorname{a6vw\_}.\;((\operatorname{a6vw\_} \in \mathbb{C}) \wedge \exists \operatorname{a6vq\_} \in \operatorname{a6vh\_}, \operatorname{a6vl\_} \in \mathbb{Z}.\; ((\operatorname{a6vy\_} = \operatorname{opr}(\operatorname{a6vg\_})(\operatorname{a6vq\_}, \operatorname{zz\text{-}act}(\operatorname{a6vg\_}, \operatorname{a6vl\_}, \operatorname{a6vb\_}))) \wedge (\operatorname{a6vw\_} = (\operatorname{a6vc\_}(\operatorname{a6vq\_}) \cdot \exp (\operatorname{a6vl\_} \cdot \operatorname{a6vt\_}))))))
\end{equation*}


\vnbbookfile{structure-library/ca-line-vocab}{structure-library/ca-line-vocab}
\noindent{\small ca-line-vocab.scm -- WAVE 2 (2026-10-27..30), wired 2026-10-30 from scratchpad/ch3e/ch3f-line.scm (vocab), scratchpad/ch3e/ch3f-refl.scm (vocab). Library: complex-analysis. Proof-free: the definitions (with their notation! lines) moved here from the proof file(s) named above.}\par
\par\medskip\noindent\textbf{Definition} (\texttt{cc-line}).\;$\operatorname{cc\text{-}line}(\operatorname{clnp\_}, \operatorname{clnd\_}) :=$
\begin{equation*}
\{\operatorname{clnz\_} \in \mathbb{C} : (\operatorname{imag\text{-}part}(((\operatorname{clnz\_} - \operatorname{clnp\_}) \cdot \operatorname{conjugate}(\operatorname{clnd\_}))) = 0)\}
\end{equation*}

\par\medskip\noindent\textbf{Definition} (\texttt{conj-set}).\;$\operatorname{conj\text{-}set}(\operatorname{cjsa\_}) :=$ $\{\operatorname{cjsz\_} \in \mathbb{C} : (\operatorname{conjugate}(\operatorname{cjsz\_}) \in \operatorname{cjsa\_})\}$

\vnbbookfile{theorem-library/cc-exp}{theorem-library/cc-exp}
\noindent{\small cc-exp.scm -- THE ELEMENTARY FUNCTIONS exp, cos, sin ON CC (the notes 2.3.1, 2.3.2, everything that does not mention pi). Batch 31-A, 2026-09-25.  The definitions are structure-library/cc-elementary.scm (CC-EXP, CC-COS, CC-SIN by their power series, the user's decision of 2026-09-25).}\par
\vnbprop{cc-exp-unfold}{\vnbstmt{\begin{array}{@{}l@{}} \forall z. \\ \quad (\exp z \simeq \operatorname{cc\text{-}series\text{-}limit}((\lambda k \in \mathbb{N}.\; (\operatorname{recip}(\operatorname{factorial}(k)) \cdot \operatorname{power}(z, k))))) \end{array}}}
\vnbprop{cc-cos-unfold}{\vnbstmt{\begin{array}{@{}l@{}} \forall z. \\ \quad (\cos z \simeq \operatorname{cc\text{-}series\text{-}limit}((\lambda k \in \mathbb{N}.\; ((\operatorname{power}(-1, k) \cdot \operatorname{recip}(\operatorname{factorial}((2 \cdot k)))) \cdot \operatorname{power}(z, (2 \cdot k)))))) \end{array}}}
\vnbprop{cc-sin-unfold}{\vnbstmt{\begin{array}{@{}l@{}} \forall z. \\ \quad (\sin z \simeq \operatorname{cc\text{-}series\text{-}limit}((\lambda k \in \mathbb{N}.\; ((\operatorname{power}(-1, k) \cdot \operatorname{recip}(\operatorname{factorial}(((2 \cdot k) + 1)))) \cdot \operatorname{power}(z, ((2 \cdot k) + 1)))))) \end{array}}}
\vnbprop{cps-radius-pos-inf}{\vnbstmt{\begin{array}{@{}l@{}} \forall f \in (\mathbb{N} \to \mathbb{C}). \\ \quad \forall t \in \mathbb{R}.\; ((0 < t) \Rightarrow \operatorname{series\text{-}converges}((\lambda k \in \mathbb{N}.\; (\operatorname{magnitude}(f(k)) \cdot \operatorname{power}(t, k))))) \Rightarrow  \\ \quad \quad \operatorname{cps\text{-}radius}(f) = \operatorname{pos\text{-}inf} \end{array}}}
\vnbprop{cex-exp-series-real-converges}{\vnbstmt{\begin{array}{@{}l@{}} \forall t \in \mathbb{R}. \\ \quad (0 < t) \Rightarrow \operatorname{series\text{-}converges}((\lambda k \in \mathbb{N}.\; (\operatorname{recip}(\operatorname{factorial}(k)) \cdot \operatorname{power}(t, k)))) \end{array}}}
\vnbprop{cc-exp-radius}{\vnbstmt{\begin{array}{@{}l@{}} \forall f \in (\mathbb{N} \to \mathbb{C}). \\ \quad \forall k \in \mathbb{N}.\; (f(k) = \operatorname{recip}(\operatorname{factorial}(k))) \Rightarrow (\operatorname{cps\text{-}radius}(f) = \operatorname{pos\text{-}inf}) \end{array}}}
\vnbprop{cps-entire-converges}{\vnbstmt{\begin{array}{@{}l@{}} \forall f \in (\mathbb{N} \to \mathbb{C}), z \in \mathbb{C}. \\ \quad (\operatorname{cps\text{-}radius}(f) = \operatorname{pos\text{-}inf}) \Rightarrow  \\ \quad \quad \operatorname{cc\text{-}series\text{-}converges}((\lambda k \in \mathbb{N}.\; (f(k) \cdot \operatorname{power}(z, k)))) \end{array}}}
\vnbprop{cc-diff-on-local}{\vnbstmt{\begin{array}{@{}l@{}} \forall u.; \forall f \in (u \to \mathbb{C}), v, g, a, l. \\ \quad \operatorname{is\text{-}open}(\operatorname{nf\text{-}metric\text{-}space}(\operatorname{cc\text{-}normed\text{-}field}), u) \wedge \\ \quad v \subseteq u \wedge \\ \quad \operatorname{is\text{-}diff\text{-}on}(\operatorname{cc\text{-}normed\text{-}field}, v, g, a, l) \wedge \\ \quad \forall x \in v.\; (f(x) = g(x)) \Rightarrow \\ \quad \quad \operatorname{is\text{-}diff\text{-}on}(\operatorname{cc\text{-}normed\text{-}field}, u, f, a, l) \end{array}}}
\vnbprop{cps-entire-diff}{\vnbstmt{\begin{array}{@{}l@{}} \forall f, a \in (\mathbb{N} \to \mathbb{C}), \mathit{cpda} \in \mathbb{C}. \\ \quad (\forall j \in \mathbb{N}.\; (a(j) = ((j + 1) \cdot f((j + 1)))) \wedge (\operatorname{cps\text{-}radius}(f) = \operatorname{pos\text{-}inf})) \Rightarrow  \\ \quad \quad \operatorname{is\text{-}diff\text{-}on}(\operatorname{cc\text{-}normed\text{-}field}, \mathbb{C}, \\ \quad \quad \quad (\lambda z \in \mathbb{C}.\; \operatorname{cc\text{-}series\text{-}limit}((\lambda k \in \mathbb{N}.\; (f(k) \cdot \operatorname{power}(z, k))))), \\ \quad \quad \quad \mathit{cpda}, \\ \quad \quad \quad \operatorname{cc\text{-}series\text{-}limit}((\lambda k \in \mathbb{N}.\; (a(k) \cdot \operatorname{power}(\mathit{cpda}, k))))) \end{array}}}
\vnbprop{cc-exp-as-power-series}{\vnbstmt{\begin{array}{@{}l@{}} \forall f \in (\mathbb{N} \to \mathbb{C}), z \in \mathbb{C}. \\ \quad \forall k \in \mathbb{N}.\; (f(k) = \operatorname{recip}(\operatorname{factorial}(k))) \Rightarrow  \\ \quad \quad \operatorname{cc\text{-}series\text{-}converges}((\lambda k \in \mathbb{N}.\; (\operatorname{recip}(\operatorname{factorial}(k)) \cdot \operatorname{power}(z, k)))) \wedge \\ \quad \quad \exp z = \operatorname{cc\text{-}series\text{-}limit}((\lambda \mathit{cpsk} \in \mathbb{N}.\; (f(\mathit{cpsk}) \cdot \operatorname{power}(z, \mathit{cpsk})))) \end{array}}}
\vnbprop{cc-exp-converges}{\vnbstmt{\begin{array}{@{}l@{}} \forall z \in \mathbb{C}. \\ \quad (\operatorname{cc\text{-}series\text{-}converges}((\lambda k \in \mathbb{N}.\; (\operatorname{recip}(\operatorname{factorial}(k)) \cdot \operatorname{power}(z, k)))) \wedge (\exp z \in \mathbb{C})) \end{array}}}
\vnbprop{cc-exp-deriv}{\vnbstmt{\begin{array}{@{}l@{}} \forall z \in \mathbb{C}. \\ \quad \operatorname{is\text{-}diff\text{-}on}(\operatorname{cc\text{-}normed\text{-}field}, \mathbb{C}, (\lambda y \in \mathbb{C}.\; \exp y), z, \exp z) \end{array}}}
\vnbprop{cc-exp-holomorphic}{\vnbstmt{\operatorname{holomorphic\text{-}on}(\mathbb{C}, (\lambda y \in \mathbb{C}.\; \exp y))}}
\vnbprop{cc-exp-zero}{\vnbstmt{\exp 0 = 1}}
\vnbprop{cex-sub-deriv}{\vnbstmt{\begin{array}{@{}l@{}} \forall w, z \in \mathbb{C}. \\ \quad \operatorname{is\text{-}diff\text{-}on}(\operatorname{cc\text{-}normed\text{-}field}, \mathbb{C}, (\lambda y \in \mathbb{C}.\; (w - y)), z, -1) \end{array}}}
\vnbprop{cex-shift-deriv-zero}{\vnbstmt{\begin{array}{@{}l@{}} \forall w, z \in \mathbb{C}. \\ \quad \operatorname{is\text{-}diff\text{-}on}(\operatorname{cc\text{-}normed\text{-}field}, \mathbb{C}, (\lambda y \in \mathbb{C}.\; (\exp y \cdot \exp (w - y))), z, 0) \end{array}}}
\vnbprop{cc-exp-shift}{\vnbstmt{\forall w, z \in \mathbb{C}.\; ((\exp z \cdot \exp (w - z)) = \exp w)}}
\vnbprop{cc-exp-add}{\vnbstmt{\forall z, w \in \mathbb{C}.\; (\exp (z + w) = (\exp z \cdot \exp w))}}
\vnbprop{cc-exp-nonzero}{\vnbstmt{\forall z \in \mathbb{C}.\; \neg (\exp z = 0)}}
\vnbprop{cc-exp-maps-to-nonzero}{\vnbstmt{\forall z \in \mathbb{C}.\; (\exp z \in \operatorname{difference}(\mathbb{C}, \operatorname{singleton}(0)))}}
\vnbprop{cc-exp-neg}{\vnbstmt{\forall z \in \mathbb{C}.\; (\exp -z = \operatorname{recip}(\exp z))}}
\vnbprop{cc-exp-continuous}{\vnbstmt{\begin{array}{@{}l@{}} \forall z \in \mathbb{C}. \\ \quad \operatorname{is\text{-}continuous\text{-}at}(\operatorname{subspace\text{-}ms}(\operatorname{nf\text{-}metric\text{-}space}(\operatorname{cc\text{-}normed\text{-}field}), \mathbb{C}), \\ \quad \quad \operatorname{nf\text{-}metric\text{-}space}(\operatorname{cc\text{-}normed\text{-}field}), (\lambda y \in \mathbb{C}.\; \exp y), z) \end{array}}}
\vnbprop{cc-conjugate-sub}{\vnbstmt{\forall a, b \in \mathbb{C}.\; (\operatorname{conjugate}((a - b)) = (\operatorname{conjugate}(a) - \operatorname{conjugate}(b)))}}
\vnbprop{cc-conjugate-power}{\vnbstmt{\forall k \in \mathbb{N}, z \in \mathbb{C}.\; (\operatorname{conjugate}(\operatorname{power}(z, k)) = \operatorname{power}(\operatorname{conjugate}(z), k))}}
\vnbprop{cc-series-limit-conjugate}{\vnbstmt{\begin{array}{@{}l@{}} \forall f, g \in (\mathbb{N} \to \mathbb{C}). \\ \quad (\forall j \in \mathbb{N}.\; (g(j) = \operatorname{conjugate}(f(j))) \wedge \operatorname{cc\text{-}series\text{-}converges}(f)) \Rightarrow  \\ \quad \quad \operatorname{cc\text{-}series\text{-}converges}(g) \wedge (\operatorname{cc\text{-}series\text{-}limit}(g) = \operatorname{conjugate}(\operatorname{cc\text{-}series\text{-}limit}(f))) \end{array}}}
\vnbprop{cc-exp-conj}{\vnbstmt{\forall z \in \mathbb{C}.\; (\exp \operatorname{conjugate}(z) = \operatorname{conjugate}(\exp z))}}
\vnbprop{cc-exp-real-valued}{\vnbstmt{\forall t \in \mathbb{R}.\; (\exp t \in \mathbb{R})}}
\vnbprop{cc-exp-real}{\vnbstmt{\forall t \in \mathbb{R}.\; ((\exp t \in \mathbb{R}) \wedge (0 < \exp t))}}
\vnbprop{cc-exp-magnitude-imaginary}{\vnbstmt{\forall t \in \mathbb{R}.\; (\operatorname{magnitude}(\exp (+i \cdot t)) = 1)}}
\vnbprop{cc-power-mul}{\vnbstmt{\forall k, m \in \mathbb{N}, z \in \mathbb{C}.\; (\operatorname{power}(z, (m \cdot k)) = \operatorname{power}(\operatorname{power}(z, m), k))}}
\vnbprop{cc-power-mul-base}{\vnbstmt{\forall k \in \mathbb{N}, a, b \in \mathbb{C}.\; (\operatorname{power}((a \cdot b), k) = (\operatorname{power}(a, k) \cdot \operatorname{power}(b, k)))}}
\vnbprop{factorial-mono}{\vnbstmt{\forall m, k \in \mathbb{N}.\; (\operatorname{factorial}(k) \leq \operatorname{factorial}((k + m)))}}
\vnbprop{cex-sign-magnitude}{\vnbstmt{\forall k \in \mathbb{N}, a \in \mathbb{R}.\; ((0 \leq a) \Rightarrow (\operatorname{magnitude}((\operatorname{power}(-1, k) \cdot a)) = a))}}
\vnbprop{cc-cos-radius}{\vnbstmt{\begin{array}{@{}l@{}} \forall f \in (\mathbb{N} \to \mathbb{C}). \\ \quad \forall k \in \mathbb{N}.\; (f(k) = (\operatorname{power}(-1, k) \cdot \operatorname{recip}(\operatorname{factorial}((2 \cdot k))))) \Rightarrow  \\ \quad \quad \operatorname{cps\text{-}radius}(f) = \operatorname{pos\text{-}inf} \end{array}}}
\vnbprop{cc-sin-radius}{\vnbstmt{\begin{array}{@{}l@{}} \forall f \in (\mathbb{N} \to \mathbb{C}). \\ \quad \forall k \in \mathbb{N}.\; (f(k) = (\operatorname{power}(-1, k) \cdot \operatorname{recip}(\operatorname{factorial}(((2 \cdot k) + 1))))) \Rightarrow  \\ \quad \quad \operatorname{cps\text{-}radius}(f) = \operatorname{pos\text{-}inf} \end{array}}}
\vnbprop{cc-cos-as-power-series}{\vnbstmt{\begin{array}{@{}l@{}} \forall f \in (\mathbb{N} \to \mathbb{C}), z \in \mathbb{C}. \\ \quad \forall k \in \mathbb{N}.\; (f(k) = (\operatorname{power}(-1, k) \cdot \operatorname{recip}(\operatorname{factorial}((2 \cdot k))))) \Rightarrow  \\ \quad \quad \operatorname{cc\text{-}series\text{-}converges}((\lambda k \in \mathbb{N}.\; ((\operatorname{power}(-1, k) \cdot \operatorname{recip}(\operatorname{factorial}((2 \cdot k)))) \cdot \operatorname{power}(z, (2 \cdot k))))) \wedge \\ \quad \quad \cos z = \operatorname{cc\text{-}series\text{-}limit}((\lambda \mathit{cpsk} \in \mathbb{N}.\; (f(\mathit{cpsk}) \cdot \operatorname{power}(\operatorname{power}(z, 2), \mathit{cpsk})))) \end{array}}}
\vnbprop{cc-sin-as-power-series}{\vnbstmt{\begin{array}{@{}l@{}} \forall f \in (\mathbb{N} \to \mathbb{C}), z \in \mathbb{C}. \\ \quad \forall k \in \mathbb{N}.\; (f(k) = (\operatorname{power}(-1, k) \cdot \operatorname{recip}(\operatorname{factorial}(((2 \cdot k) + 1))))) \Rightarrow  \\ \quad \quad \operatorname{cc\text{-}series\text{-}converges}((\lambda k \in \mathbb{N}.\; ((\operatorname{power}(-1, k) \cdot \operatorname{recip}(\operatorname{factorial}(((2 \cdot k) + 1)))) \cdot \operatorname{power}(z, ((2 \cdot k) + 1))))) \wedge \\ \quad \quad \sin z = (z \cdot \operatorname{cc\text{-}series\text{-}limit}((\lambda \mathit{cpsk} \in \mathbb{N}.\; (f(\mathit{cpsk}) \cdot \operatorname{power}(\operatorname{power}(z, 2), \mathit{cpsk}))))) \end{array}}}
\vnbprop{cc-cos-converges}{\vnbstmt{\begin{array}{@{}l@{}} \forall z \in \mathbb{C}. \\ \quad \operatorname{cc\text{-}series\text{-}converges}((\lambda k \in \mathbb{N}.\; ((\operatorname{power}(-1, k) \cdot \operatorname{recip}(\operatorname{factorial}((2 \cdot k)))) \cdot \operatorname{power}(z, (2 \cdot k))))) \wedge \\ \quad \cos z \in \mathbb{C} \end{array}}}
\vnbprop{cc-sin-converges}{\vnbstmt{\begin{array}{@{}l@{}} \forall z \in \mathbb{C}. \\ \quad \operatorname{cc\text{-}series\text{-}converges}((\lambda k \in \mathbb{N}.\; ((\operatorname{power}(-1, k) \cdot \operatorname{recip}(\operatorname{factorial}(((2 \cdot k) + 1)))) \cdot \operatorname{power}(z, ((2 \cdot k) + 1))))) \wedge \\ \quad \sin z \in \mathbb{C} \end{array}}}
\vnbprop{cc-series-even-odd}{\vnbstmt{\begin{array}{@{}l@{}} \forall f, e, o \in (\mathbb{N} \to \mathbb{C}). \\ \quad \forall j \in \mathbb{N}.\; (e(j) = f((2 \cdot j))) \wedge \\ \quad \forall j \in \mathbb{N}.\; (o(j) = f(((2 \cdot j) + 1))) \wedge \\ \quad \operatorname{cc\text{-}series\text{-}converges}(f) \wedge \\ \quad \operatorname{cc\text{-}series\text{-}converges}(e) \wedge \\ \quad \operatorname{cc\text{-}series\text{-}converges}(o) \Rightarrow \\ \quad \quad \operatorname{cc\text{-}series\text{-}limit}(f) = (\operatorname{cc\text{-}series\text{-}limit}(e) + \operatorname{cc\text{-}series\text{-}limit}(o)) \end{array}}}
\vnbprop{cc-exp-i}{\vnbstmt{\forall z \in \mathbb{C}.\; (\exp (+i \cdot z) = (\cos z + (+i \cdot \sin z)))}}
\vnbprop{cc-cos-even}{\vnbstmt{\forall z \in \mathbb{C}.\; (\cos -z = \cos z)}}
\vnbprop{cc-sin-odd}{\vnbstmt{\forall z \in \mathbb{C}.\; (\sin -z = -\sin z)}}
\vnbprop{cc-sin-exp}{\vnbstmt{\begin{array}{@{}l@{}} \forall z \in \mathbb{C}. \\ \quad \sin z = ((\exp (+i \cdot z) - \exp -(+i \cdot z)) \cdot \operatorname{recip}(+2i)) \end{array}}}
\vnbprop{cc-cos-exp}{\vnbstmt{\forall z \in \mathbb{C}.\; (\cos z = ((\exp (+i \cdot z) + \exp -(+i \cdot z)) \cdot \operatorname{recip}(2)))}}
\vnbprop{cex-lin-deriv}{\vnbstmt{\begin{array}{@{}l@{}} \forall c, z \in \mathbb{C}. \\ \quad \operatorname{is\text{-}diff\text{-}on}(\operatorname{cc\text{-}normed\text{-}field}, \mathbb{C}, (\lambda y \in \mathbb{C}.\; (c \cdot y)), z, c) \end{array}}}
\vnbprop{cc-exp-scaled-deriv}{\vnbstmt{\begin{array}{@{}l@{}} \forall c, z \in \mathbb{C}. \\ \quad \operatorname{is\text{-}diff\text{-}on}(\operatorname{cc\text{-}normed\text{-}field}, \mathbb{C}, (\lambda y \in \mathbb{C}.\; \exp (c \cdot y)), z, (c \cdot \exp (c \cdot z))) \end{array}}}
\vnbprop{cc-diff-on-scale}{\vnbstmt{\begin{array}{@{}l@{}} \forall u, f, a, l.; \forall c \in \mathbb{C}. \\ \quad \operatorname{is\text{-}diff\text{-}on}(\operatorname{cc\text{-}normed\text{-}field}, u, f, a, l) \Rightarrow  \\ \quad \quad \operatorname{is\text{-}diff\text{-}on}(\operatorname{cc\text{-}normed\text{-}field}, u, (\lambda y \in u.\; (c \cdot f(y))), a, (c \cdot l)) \end{array}}}
\vnbprop{cex-exp-comb-deriv}{\vnbstmt{\begin{array}{@{}l@{}} \forall a, b, c, d, z \in \mathbb{C}. \\ \quad \operatorname{is\text{-}diff\text{-}on}(\operatorname{cc\text{-}normed\text{-}field}, \mathbb{C}, \\ \quad \quad (\lambda y \in \mathbb{C}.\; ((a \cdot \exp (c \cdot y)) + (b \cdot \exp (d \cdot y)))), \\ \quad \quad z, ((a \cdot (c \cdot \exp (c \cdot z))) + (b \cdot (d \cdot \exp (d \cdot z))))) \end{array}}}
\vnbprop{cc-cos-deriv}{\vnbstmt{\begin{array}{@{}l@{}} \forall z \in \mathbb{C}. \\ \quad \operatorname{is\text{-}diff\text{-}on}(\operatorname{cc\text{-}normed\text{-}field}, \mathbb{C}, (\lambda y \in \mathbb{C}.\; \cos y), z, -\sin z) \end{array}}}
\vnbprop{cc-sin-deriv}{\vnbstmt{\begin{array}{@{}l@{}} \forall z \in \mathbb{C}. \\ \quad \operatorname{is\text{-}diff\text{-}on}(\operatorname{cc\text{-}normed\text{-}field}, \mathbb{C}, (\lambda y \in \mathbb{C}.\; \sin y), z, \cos z) \end{array}}}
\vnbprop{cc-cos-holomorphic}{\vnbstmt{\operatorname{holomorphic\text{-}on}(\mathbb{C}, (\lambda y \in \mathbb{C}.\; \cos y))}}
\vnbprop{cc-sin-holomorphic}{\vnbstmt{\operatorname{holomorphic\text{-}on}(\mathbb{C}, (\lambda y \in \mathbb{C}.\; \sin y))}}
\vnbprop{cc-cos-add}{\vnbstmt{\begin{array}{@{}l@{}} \forall z, w \in \mathbb{C}. \\ \quad \cos (z + w) = ((\cos z \cdot \cos w) - (\sin z \cdot \sin w)) \end{array}}}
\vnbprop{cc-sin-add}{\vnbstmt{\begin{array}{@{}l@{}} \forall z, w \in \mathbb{C}. \\ \quad \sin (z + w) = ((\cos w \cdot \sin z) + (\cos z \cdot \sin w)) \end{array}}}
\vnbprop{cc-cos-sin-one}{\vnbstmt{\forall z \in \mathbb{C}.\; ((\operatorname{power}(\cos z, 2) + \operatorname{power}(\sin z, 2)) = 1)}}

\vnbbookfile{structure-library/cc-pi}{structure-library/cc-pi}
\noindent{\small cc-pi.scm -- THE CONSTANT pi, DEFINED FROM THE KERNEL OF exp (the notes 2.3.3). DEFINITION ONLY; existence, uniqueness and the characterisation `pi-spec' are PROVEN in theorem-library/cc-exp-kernel.scm. Batch 31-B, 2026-09-25.}\par
\par\medskip\noindent\textbf{Definition} (\texttt{pi}).\;
\begin{equation*}
(\pi \simeq (\iota\, \operatorname{cekp\_}.\;((\operatorname{cekp\_} \in \mathbb{R}) \wedge ((0 < \operatorname{cekp\_}) \wedge \forall \operatorname{cekz\_} \in \mathbb{C}.\; ((\exp \operatorname{cekz\_} = 1) \Leftrightarrow \exists \operatorname{cekk\_} \in \mathbb{Z}.\; (\operatorname{cekz\_} = \cdot(\operatorname{cekk\_}, \cdot(\cdot(2, \operatorname{cekp\_}), +i))))))))
\end{equation*}


\vnbbookfile{structure-library/ca-angle-vocab}{structure-library/ca-angle-vocab}
\noindent{\small ca-angle-vocab.scm -- WAVE 2 (2026-10-27..30), wired 2026-10-30 from scratchpad/ch3e/ch3e-angle.scm (vocab). Library: complex-analysis. Proof-free: the definitions (with their notation! lines) moved here from the proof file(s) named above.}\par
\par\medskip\noindent\textbf{Definition} (\texttt{angle}).\;$\operatorname{angle}(\operatorname{ang1\_}, \operatorname{ang2\_}) :=$
\begin{equation*}
\pi + \operatorname{imag\text{-}part}(\varepsilon(\{\operatorname{angz\_} \in \mathbb{C} : ((-\pi \leq \operatorname{imag\text{-}part}(\operatorname{angz\_})) \wedge ((\operatorname{imag\text{-}part}(\operatorname{angz\_}) < \pi) \wedge (\exp \operatorname{angz\_} = -(\operatorname{ang1\_} \cdot \operatorname{conjugate}(\operatorname{ang2\_})))))\}))
\end{equation*}

\par\medskip\noindent\textbf{Definition} (\texttt{positively-oriented}).\;For all $\operatorname{poa\_}$, $\operatorname{pob\_}$ and $\operatorname{poc\_}$, 
\begin{equation*}
\begin{array}{@{}l@{}} \operatorname{positively\text{-}oriented}(\operatorname{poa\_}, \operatorname{pob\_}, \operatorname{poc\_}) \\ \Leftrightarrow\; \forall \operatorname{pou\_} \in \operatorname{interior}(\operatorname{cc\text{-}ms}, \operatorname{conv3}(\operatorname{poa\_}, \operatorname{pob\_}, \operatorname{poc\_})). \\ \quad \operatorname{angle}(-(\operatorname{pob\_}, \operatorname{pou\_}), -(\operatorname{poa\_}, \operatorname{pou\_})) < \pi \wedge \\ \quad \operatorname{angle}(-(\operatorname{poc\_}, \operatorname{pou\_}), -(\operatorname{pob\_}, \operatorname{pou\_})) < \pi \wedge \\ \quad \operatorname{angle}(-(\operatorname{poa\_}, \operatorname{pou\_}), -(\operatorname{poc\_}, \operatorname{pou\_})) < \pi \end{array}
\end{equation*}


\vnbbookfile{structure-library/mvp-vocab}{structure-library/mvp-vocab}
\noindent{\small mvp-vocab.scm -- WAVE 2 (2026-10-27..30), wired 2026-10-30 from scratchpad/calint2/mvp-vocab.scm. Library: complex-analysis.}\par
\par\medskip\noindent\textbf{Definition} (\texttt{has-mean-value-property}).\;For all $\operatorname{mvu\_}$ and $\operatorname{mvf\_}$, 
\begin{equation*}
\begin{array}{@{}l@{}} \operatorname{has\text{-}mean\text{-}value\text{-}property}(\operatorname{mvu\_}, \operatorname{mvf\_}) \\ \Leftrightarrow\; \operatorname{is\text{-}open}(\operatorname{nf\text{-}metric\text{-}space}(\operatorname{cc\text{-}normed\text{-}field}), \operatorname{mvu\_}) \wedge \\ \operatorname{mvf\_} \in (\operatorname{mvu\_} \to \mathbb{C}) \wedge \\ \forall \operatorname{rgy\_} \in \operatorname{mvu\_}. \\ \quad \operatorname{is\text{-}continuous\text{-}at}(\operatorname{subspace\text{-}ms}(\operatorname{nf\text{-}metric\text{-}space}(\operatorname{cc\text{-}normed\text{-}field}), \operatorname{mvu\_}), \\ \quad \quad \operatorname{nf\text{-}metric\text{-}space}(\operatorname{cc\text{-}normed\text{-}field}), \operatorname{mvf\_}, \operatorname{rgy\_}) \wedge \\ \forall \operatorname{mvc\_} \in \mathbb{C}, \operatorname{mvr\_} \in \mathbb{R}. \\ \quad (0 < \operatorname{mvr\_}) \Rightarrow  \\ \quad (\operatorname{closure}(\operatorname{cc\text{-}ms}, \operatorname{ball}(\operatorname{cc\text{-}ms}, \operatorname{mvc\_}, \operatorname{mvr\_})) \subseteq \operatorname{mvu\_}) \Rightarrow  \\ \quad \quad \operatorname{mvf\_}(\operatorname{mvc\_}) \\ \quad \quad =\; \cdot(\operatorname{recip}(\cdot(2, \pi)), \int_{0}^{\cdot(2, \pi)} (\lambda \operatorname{mvt\_} \in \operatorname{ccint}(0, \cdot(2, \pi)).\; \operatorname{mvf\_}(+(\operatorname{mvc\_}, \cdot(\operatorname{mvr\_}, \exp \cdot(+i, \operatorname{mvt\_})))))) \end{array}
\end{equation*}


\vnbbookfile{structure-library/analytic-log}{structure-library/analytic-log}
\noindent{\small analytic-log.scm -- ANALYTICITY AT A POINT (the notes' Def 2.11) AND THE PRINCIPAL LOGARITHM (the notes 2.3.3.1, after Prop 2.19). DEFINITIONS ONLY; the laws are PROVEN in theorem-library/analytic-log-laws.scm. Batch 38, 2026-09-27 (the user's decision of the same day: IS-ANALYTIC-AT}\par
\par\medskip\noindent\textbf{Definition} (\texttt{is-analytic-at}).\;For all $\operatorname{anlf\_}$, $\operatorname{anlu\_}$ and $\operatorname{anla\_}$, 
\begin{equation*}
\begin{array}{@{}l@{}} \operatorname{anlf\_} \text{ is analytic at } \operatorname{anla\_} \\ \Leftrightarrow\; \operatorname{is\text{-}open}(\operatorname{cc\text{-}ms}, \operatorname{anlu\_}) \wedge \\ \operatorname{anlf\_} \in (\operatorname{anlu\_} \to \mathbb{C}) \wedge \\ \operatorname{anla\_} \in \operatorname{anlu\_} \wedge \\ \exists \operatorname{anlcf\_} \in (\mathbb{N} \to \mathbb{C}), \operatorname{anlr\_} \in \mathbb{R}. \\ \quad 0 < \operatorname{anlr\_} \wedge \\ \quad \operatorname{anlr\_} < \operatorname{cps\text{-}radius}(\operatorname{anlcf\_}) \wedge \\ \quad \operatorname{ball}(\operatorname{cc\text{-}ms}, \operatorname{anla\_}, \operatorname{anlr\_}) \subseteq \operatorname{anlu\_} \wedge \\ \quad \forall \operatorname{anlz\_} \in \operatorname{ball}(\operatorname{cc\text{-}ms}, \operatorname{anla\_}, \operatorname{anlr\_}). \\ \quad \quad \operatorname{anlf\_}(\operatorname{anlz\_}) = \operatorname{cc\text{-}series\text{-}limit}((\lambda \operatorname{anlk\_} \in \mathbb{N}.\; \cdot(\operatorname{anlcf\_}(\operatorname{anlk\_}), \operatorname{power}(-(\operatorname{anlz\_}, \operatorname{anla\_}), \operatorname{anlk\_})))) \end{array}
\end{equation*}

\par\medskip\noindent\textbf{Definition} (\texttt{cc-log}).\;$\log\left(\operatorname{clgw\_}\right) :=$
\begin{equation*}
(\iota\, \operatorname{clgz\_}.\;((\operatorname{clgz\_} \in \mathbb{C}) \wedge ((-\pi < \operatorname{imag\text{-}part}(\operatorname{clgz\_})) \wedge ((\operatorname{imag\text{-}part}(\operatorname{clgz\_}) < \pi) \wedge (\exp \operatorname{clgz\_} = \operatorname{clgw\_})))))
\end{equation*}


\vnbbookfile{theorem-library/cc-exp-kernel}{theorem-library/cc-exp-kernel}
\noindent{\small cc-exp-kernel.scm -- THE KERNEL OF exp AND THE CONSTANT pi (the notes 2.3.3: Lemma 2.18, Prop 2.19, and what they rest on). Batch 31-B, 2026-09-25.  exp, cos, sin are theorem-library/cc-exp.scm's (31-A, definitions structure-library/cc-elementary.scm); PI is defined in}\par
\vnbprop{cek-limit-tail-le}{\vnbstmt{\begin{array}{@{}l@{}} \forall s \in (\mathbb{N} \to \mathbb{C}), v, c \in \mathbb{C}, d \in \mathbb{R}, m \in \mathbb{N}. \\ \quad (\operatorname{converges\text{-}to}(\operatorname{cc\text{-}ms}, s, v) \wedge \forall n \in \mathbb{N}.\; ((m \leq n) \Rightarrow (\operatorname{magnitude}((s(n) - c)) \leq d))) \Rightarrow  \\ \quad \quad \operatorname{magnitude}((v - c)) \leq d \end{array}}}
\vnbprop{cek-half-ratio-partial}{\vnbstmt{\begin{array}{@{}l@{}} \forall n \in \mathbb{N}, g \in (\mathbb{N} \to \mathbb{R}), m \in \mathbb{N}. \\ \quad \forall j \in \mathbb{N}.\; ((m \leq j) \Rightarrow ((2 \cdot g((j + 1))) \leq g(j))) \Rightarrow  \\ \quad \quad ((\operatorname{series\text{-}partial\text{-}sum}(g, (m + n)) - \operatorname{series\text{-}partial\text{-}sum}(g, m)) + (2 \cdot g((m + n)))) \leq (2 \cdot g(m)) \end{array}}}
\vnbprop{cc-series-ratio-tail}{\vnbstmt{\begin{array}{@{}l@{}} \forall f \in (\mathbb{N} \to \mathbb{C}), m \in \mathbb{N}. \\ \quad \operatorname{cc\text{-}series\text{-}converges}(f) \wedge \\ \quad \forall j \in \mathbb{N}.\; ((m \leq j) \Rightarrow ((2 \cdot \operatorname{magnitude}(f((j + 1)))) \leq \operatorname{magnitude}(f(j)))) \Rightarrow \\ \quad \quad \operatorname{magnitude}((\operatorname{cc\text{-}series\text{-}limit}(f) - \operatorname{cc\text{-}series\text{-}partial\text{-}sum}(f, m))) \leq (2 \cdot \operatorname{magnitude}(f(m))) \end{array}}}
\vnbprop{cc-exp-tail}{\vnbstmt{\begin{array}{@{}l@{}} \forall z \in \mathbb{C}, m \in \mathbb{N}. \\ \quad ((2 \cdot \operatorname{magnitude}(z)) \leq (m + 1)) \Rightarrow  \\ \quad \quad \operatorname{magnitude}((\exp z - \operatorname{cc\text{-}series\text{-}partial\text{-}sum}((\lambda k \in \mathbb{N}.\; (\operatorname{recip}(\operatorname{factorial}(k)) \cdot \operatorname{power}(z, k))), m))) \\ \quad \quad \leq\; 2 \cdot (\operatorname{recip}(\operatorname{factorial}(m)) \cdot \operatorname{power}(\operatorname{magnitude}(z), m)) \end{array}}}
\vnbprop{cc-exp-near-one}{\vnbstmt{\begin{array}{@{}l@{}} \forall z \in \mathbb{C}. \\ \quad (\operatorname{magnitude}(z) \leq 1) \Rightarrow (\operatorname{magnitude}(((\exp z - 1) - z)) \leq \operatorname{power}(\operatorname{magnitude}(z), 2)) \end{array}}}
\vnbprop{cc-exp-kernel-discrete}{\vnbstmt{\forall z \in \mathbb{C}.\; (((\exp z = 1) \wedge (\operatorname{magnitude}(z) < 1/2)) \Rightarrow (z = 0))}}
\vnbprop{cek-exp-gt-one-small}{\vnbstmt{\forall t \in \mathbb{R}.\; (((0 < t) \wedge (t \leq 1/2)) \Rightarrow (1 < \exp t))}}
\vnbprop{cek-exp-gt-one-n}{\vnbstmt{\forall n \in \mathbb{N}, t \in \mathbb{R}.\; (((0 < t) \wedge (t \leq (n \cdot 1/2))) \Rightarrow (1 < \exp t))}}
\vnbprop{cc-exp-real-gt-one}{\vnbstmt{\forall t \in \mathbb{R}.\; ((0 < t) \Rightarrow (1 < \exp t))}}
\vnbprop{cc-exp-real-one-iff}{\vnbstmt{\forall t \in \mathbb{R}.\; ((\exp t = 1) \Leftrightarrow (t = 0))}}
\vnbprop{cc-exp-kernel-add}{\vnbstmt{\begin{array}{@{}l@{}} \forall z, w \in \mathbb{C}. \\ \quad ((\exp z = 1) \wedge (\exp w = 1)) \Rightarrow (\exp (z + w) = 1) \end{array}}}
\vnbprop{cc-exp-kernel-neg}{\vnbstmt{\forall z \in \mathbb{C}.\; ((\exp z = 1) \Rightarrow (\exp -z = 1))}}
\vnbprop{cc-exp-kernel-sub}{\vnbstmt{\begin{array}{@{}l@{}} \forall z, w \in \mathbb{C}. \\ \quad ((\exp z = 1) \wedge (\exp w = 1)) \Rightarrow (\exp (z - w) = 1) \end{array}}}
\vnbprop{cc-exp-kernel-imaginary}{\vnbstmt{\forall z \in \mathbb{C}.\; ((\exp z = 1) \Rightarrow (\operatorname{real\text{-}part}(z) = 0))}}
\vnbprop{cc-exp-kernel-on-imaginary-axis}{\vnbstmt{\forall z \in \mathbb{C}.\; ((\exp z = 1) \Rightarrow (z = (+i \cdot \operatorname{imag\text{-}part}(z))))}}
\vnbprop{cc-cos-real-part}{\vnbstmt{\forall t \in \mathbb{R}.\; (\cos t = \operatorname{real\text{-}part}(\exp (+i \cdot t)))}}
\vnbprop{cc-sin-real-part}{\vnbstmt{\forall t \in \mathbb{R}.\; (\sin t = \operatorname{imag\text{-}part}(\exp (+i \cdot t)))}}
\vnbprop{cc-cos-sin-real}{\vnbstmt{\forall t \in \mathbb{R}.\; ((\cos t \in \mathbb{R}) \wedge (\sin t \in \mathbb{R}))}}
\vnbprop{cc-cos-one-bounds}{\vnbstmt{(5/12 \leq \cos 1) \wedge (\cos 1 \leq 7/12)}}
\vnbprop{cc-cos-two-neg}{\vnbstmt{\cos 2 < 0}}
\vnbprop{cc-continuous-restrict-rr}{\vnbstmt{\begin{array}{@{}l@{}} \forall f, g, x. \\ \quad \operatorname{is\text{-}continuous\text{-}at}(\operatorname{subspace\text{-}ms}(\operatorname{nf\text{-}metric\text{-}space}(\operatorname{cc\text{-}normed\text{-}field}), \mathbb{C}), \\ \quad \quad \operatorname{nf\text{-}metric\text{-}space}(\operatorname{cc\text{-}normed\text{-}field}), f, x) \wedge \\ \quad g \in (\mathbb{R} \to \mathbb{R}) \wedge \\ \quad x \in \mathbb{R} \wedge \\ \quad \forall v \in \mathbb{R}.\; (g(v) = f(v)) \Rightarrow \\ \quad \quad \operatorname{is\text{-}continuous\text{-}at}(\operatorname{rr\text{-}ms}, \operatorname{rr\text{-}ms}, g, x) \end{array}}}
\vnbprop{cc-neg-cos-continuous-rr}{\vnbstmt{\begin{array}{@{}l@{}} (\lambda v \in \mathbb{R}.\; (-1 \cdot \cos v)) \in (\mathbb{R} \to \mathbb{R}) \wedge \\ \forall x \in \mathbb{R}. \\ \quad \operatorname{is\text{-}continuous\text{-}at}(\operatorname{rr\text{-}ms}, \operatorname{rr\text{-}ms}, (\lambda v \in \mathbb{R}.\; (-1 \cdot \cos v)), x) \end{array}}}
\vnbprop{cc-cos-real-continuous}{\vnbstmt{\begin{array}{@{}l@{}} (\lambda v \in \mathbb{R}.\; \cos v) \in (\mathbb{R} \to \mathbb{R}) \wedge \\ \forall x \in \mathbb{R}. \\ \quad \operatorname{is\text{-}continuous\text{-}at}(\operatorname{rr\text{-}ms}, \operatorname{rr\text{-}ms}, (\lambda v \in \mathbb{R}.\; \cos v), x) \end{array}}}
\vnbprop{cc-cos-zero}{\vnbstmt{\cos 0 = 1}}
\vnbprop{cc-cos-has-zero}{\vnbstmt{\exists t \in \mathbb{R}.\; ((0 < t) \wedge ((t < 2) \wedge (\cos t = 0)))}}
\vnbprop{cc-exp-imaginary-period-exists}{\vnbstmt{\exists t \in \mathbb{R}.\; ((0 < t) \wedge ((t < 8) \wedge (\exp (+i \cdot t) = 1)))}}
\vnbprop{cc-exp-kernel-nontrivial}{\vnbstmt{\exists z \in \mathbb{C}.\; ((\exp z = 1) \wedge \neg (z = 0))}}
\vnbprop{cek-period-ge-half}{\vnbstmt{\forall t \in \mathbb{R}.\; (((0 < t) \wedge (\exp (+i \cdot t) = 1)) \Rightarrow (1/2 \leq t))}}
\vnbprop{cc-exp-least-period}{\vnbstmt{\begin{array}{@{}l@{}} \exists a \in \mathbb{R}. \\ \quad ((0 < a) \wedge ((\exp (+i \cdot a) = 1) \wedge \forall t \in \mathbb{R}.\; (((0 < t) \wedge (\exp (+i \cdot t) = 1)) \Rightarrow (a \leq t)))) \end{array}}}
\vnbprop{cek-kernel-nn-multiple}{\vnbstmt{\forall n \in \mathbb{N}, w \in \mathbb{C}.\; ((\exp w = 1) \Rightarrow (\exp (n \cdot w) = 1))}}
\vnbprop{cc-exp-kernel-zz-multiple}{\vnbstmt{\forall w \in \mathbb{C}, k \in \mathbb{Z}.\; ((\exp w = 1) \Rightarrow (\exp (k \cdot w) = 1))}}
\vnbprop{cek-period-multiple-nn}{\vnbstmt{\begin{array}{@{}l@{}} \forall n \in \mathbb{N}, a, r \in \mathbb{R}. \\ \quad 0 < a \wedge \\ \quad \exp (+i \cdot a) = 1 \wedge \\ \quad \forall t \in \mathbb{R}.\; (((0 < t) \wedge (\exp (+i \cdot t) = 1)) \Rightarrow (a \leq t)) \wedge \\ \quad 0 \leq r \wedge \\ \quad r < (n \cdot a) \wedge \\ \quad \exp (+i \cdot r) = 1 \Rightarrow \\ \quad \quad \exists k \in \mathbb{N}.\; (r = (k \cdot a)) \end{array}}}
\vnbprop{cek-period-multiple-zz}{\vnbstmt{\begin{array}{@{}l@{}} \forall a, r \in \mathbb{R}. \\ \quad 0 < a \wedge \\ \quad \exp (+i \cdot a) = 1 \wedge \\ \quad \forall t \in \mathbb{R}.\; (((0 < t) \wedge (\exp (+i \cdot t) = 1)) \Rightarrow (a \leq t)) \wedge \\ \quad \exp (+i \cdot r) = 1 \Rightarrow \\ \quad \quad \exists k \in \mathbb{Z}.\; (r = (k \cdot a)) \end{array}}}
\vnbprop{cc-pi-exists}{\vnbstmt{\begin{array}{@{}l@{}} \exists p \in \mathbb{R}. \\ \quad ((0 < p) \wedge \forall z \in \mathbb{C}.\; ((\exp z = 1) \Leftrightarrow \exists k \in \mathbb{Z}.\; (z = (k \cdot ((2 \cdot p) \cdot +i))))) \end{array}}}
\vnbprop{cek-pi-le}{\vnbstmt{\begin{array}{@{}l@{}} \forall p, q. \\ \quad p \in \mathbb{R} \wedge \\ \quad 0 < p \wedge \\ \quad \forall z \in \mathbb{C}.\; ((\exp z = 1) \Leftrightarrow \exists k \in \mathbb{Z}.\; (z = (k \cdot ((2 \cdot p) \cdot +i)))) \wedge \\ \quad q \in \mathbb{R} \wedge \\ \quad 0 < q \wedge \\ \quad \forall z \in \mathbb{C}.\; ((\exp z = 1) \Leftrightarrow \exists k \in \mathbb{Z}.\; (z = (k \cdot ((2 \cdot q) \cdot +i)))) \Rightarrow \\ \quad \quad q \leq p \end{array}}}
\vnbprop{cc-pi-unique}{\vnbstmt{\begin{array}{@{}l@{}} \forall p, q. \\ \quad p \in \mathbb{R} \wedge \\ \quad 0 < p \wedge \\ \quad \forall z \in \mathbb{C}.\; ((\exp z = 1) \Leftrightarrow \exists k \in \mathbb{Z}.\; (z = (k \cdot ((2 \cdot p) \cdot +i)))) \wedge \\ \quad q \in \mathbb{R} \wedge \\ \quad 0 < q \wedge \\ \quad \forall z \in \mathbb{C}.\; ((\exp z = 1) \Leftrightarrow \exists k \in \mathbb{Z}.\; (z = (k \cdot ((2 \cdot q) \cdot +i)))) \Rightarrow \\ \quad \quad p = q \end{array}}}
\vnbprop{pi-spec}{\vnbstmt{\begin{array}{@{}l@{}} \pi \in \mathbb{R} \wedge \\ 0 < \pi \wedge \\ \forall z \in \mathbb{C}.\; ((\exp z = 1) \Leftrightarrow \exists k \in \mathbb{Z}.\; (z = (k \cdot ((2 \cdot \pi) \cdot +i)))) \end{array}}}
\vnbprop{pi-pos}{\vnbstmt{(\pi \in \mathbb{R}) \wedge (0 < \pi)}}
\vnbprop{cc-exp-one-iff}{\vnbstmt{\forall z \in \mathbb{C}.\; ((\exp z = 1) \Leftrightarrow \exists k \in \mathbb{Z}.\; (z = (k \cdot ((2 \cdot \pi) \cdot +i))))}}
\vnbprop{cc-exp-two-pi-i}{\vnbstmt{\exp ((2 \cdot \pi) \cdot +i) = 1}}
\vnbprop{cc-cos-sin-two-pi}{\vnbstmt{(\cos (2 \cdot \pi) = 1) \wedge (\sin (2 \cdot \pi) = 0)}}
\vnbprop{cc-exp-i-pi}{\vnbstmt{\exp (+i \cdot \pi) = -1}}
\vnbprop{cc-cos-sin-pi}{\vnbstmt{(\cos \pi = -1) \wedge (\sin \pi = 0)}}
\vnbprop{cc-cos-sin-half-pi}{\vnbstmt{(\cos (1/2 \cdot \pi) = 0) \wedge (\operatorname{power}(\sin (1/2 \cdot \pi), 2) = 1)}}
\vnbprop{cc-exp-injective-strip}{\vnbstmt{\begin{array}{@{}l@{}} \forall z, w \in \mathbb{C}. \\ \quad -\pi \leq \operatorname{imag\text{-}part}(z) \wedge \\ \quad \operatorname{imag\text{-}part}(z) < \pi \wedge \\ \quad -\pi \leq \operatorname{imag\text{-}part}(w) \wedge \\ \quad \operatorname{imag\text{-}part}(w) < \pi \wedge \\ \quad \exp z = \exp w \Rightarrow \\ \quad \quad z = w \end{array}}}
\vnbprop{cc-exp-kernel-generator}{\vnbstmt{\begin{array}{@{}l@{}} \exists g \in \mathbb{C}. \\ \quad ((0 < \operatorname{imag\text{-}part}(g)) \wedge \forall z \in \mathbb{C}.\; ((\exp z = 1) \Leftrightarrow \exists k \in \mathbb{Z}.\; (z = (k \cdot g)))) \end{array}}}
\vnbprop{cc-exp-kernel-generator-unique}{\vnbstmt{\begin{array}{@{}l@{}} \forall g \in \mathbb{C}. \\ \quad ((0 < \operatorname{imag\text{-}part}(g)) \wedge \forall z \in \mathbb{C}.\; ((\exp z = 1) \Leftrightarrow \exists k \in \mathbb{Z}.\; (z = (k \cdot g)))) \Rightarrow  \\ \quad \quad g = ((2 \cdot \pi) \cdot +i) \end{array}}}
\vnbprop{cc-pi-lt-four}{\vnbstmt{\pi < 4}}
\vnbprop{cc-sin-half-pi}{\vnbstmt{\sin (1/2 \cdot \pi) = 1}}
\vnbprop{cc-cos-two-pi}{\vnbstmt{\cos (2 \cdot \pi) = 1}}
\vnbprop{cc-sin-two-pi}{\vnbstmt{\sin (2 \cdot \pi) = 0}}
\vnbprop{cc-cos-pi}{\vnbstmt{\cos \pi = -1}}
\vnbprop{cc-sin-pi}{\vnbstmt{\sin \pi = 0}}
\vnbprop{cc-cos-half-pi}{\vnbstmt{\cos (1/2 \cdot \pi) = 0}}

\vnbbookfile{theorem-library/cc-exp-surjective}{theorem-library/cc-exp-surjective}
\noindent{\small cc-exp-surjective.scm -- PROPOSITION 2.19 COMPLETED (exp is a bijection of each strip onto its image) AND THE BRIDGE cc-exp|RR = r-exp. Batch 34, 2026-09-26.  exp, cos, sin: theorem-library/cc-exp.scm (31-A); PI and the injectivity half of 2.19: theorem-library/cc-exp-kernel.scm (31-B).}\par
\vnbprop{cc-exp-real-continuous-rr}{\vnbstmt{\begin{array}{@{}l@{}} (\lambda v \in \mathbb{R}.\; \exp v) \in (\mathbb{R} \to \mathbb{R}) \wedge \\ \forall x \in \mathbb{R}. \\ \quad \operatorname{is\text{-}continuous\text{-}at}(\operatorname{rr\text{-}ms}, \operatorname{rr\text{-}ms}, (\lambda v \in \mathbb{R}.\; \exp v), x) \end{array}}}
\vnbprop{cc-diff-on-restrict-rr}{\vnbstmt{\begin{array}{@{}l@{}} \forall f, g, a, l. \\ \quad \operatorname{is\text{-}diff\text{-}on}(\operatorname{cc\text{-}normed\text{-}field}, \mathbb{C}, f, a, l) \wedge \\ \quad g \in (\mathbb{R} \to \mathbb{R}) \wedge \\ \quad a \in \mathbb{R} \wedge \\ \quad l \in \mathbb{R} \wedge \\ \quad \forall v \in \mathbb{R}.\; (g(v) = f(v)) \Rightarrow \\ \quad \quad \operatorname{is\text{-}diff\text{-}at}(g, a, l) \end{array}}}
\vnbprop{cc-exp-real-deriv}{\vnbstmt{\forall t \in \mathbb{R}.\; \operatorname{is\text{-}diff\text{-}at}((\lambda v \in \mathbb{R}.\; \exp v), t, \exp t)}}
\vnbprop{cxs-log-exp-deriv}{\vnbstmt{\forall x \in \mathbb{R}.\; \operatorname{is\text{-}diff\text{-}at}((\lambda h \in \mathbb{R}.\; (\log\left(\exp h\right) - h)), x, 0)}}
\vnbprop{cc-exp-real-is-r-exp}{\vnbstmt{\forall t \in \mathbb{R}.\; (\exp t = \exp\left(t\right))}}
\vnbprop{cc-exp-real-surjective-pos}{\vnbstmt{\forall r \in \mathbb{R}.\; ((0 < r) \Rightarrow \exists s \in \mathbb{R}.\; (\exp s = r))}}
\vnbprop{cc-circle-arg}{\vnbstmt{\begin{array}{@{}l@{}} \forall a, b \in \mathbb{R}. \\ \quad ((\operatorname{power}(a, 2) + \operatorname{power}(b, 2)) = 1) \Rightarrow  \\ \quad \quad \exists h \in \mathbb{R}.\; ((-\pi \leq h) \wedge ((h < \pi) \wedge ((\cos h = a) \wedge (\sin h = b)))) \end{array}}}
\vnbprop{cc-polar-form}{\vnbstmt{\begin{array}{@{}l@{}} \forall w \in \mathbb{C}. \\ \quad \neg (w = 0) \Rightarrow  \\ \quad \quad \exists r, h \in \mathbb{R}. \\ \quad \quad \quad (0 < r) \wedge ((-\pi \leq h) \wedge ((h < \pi) \wedge (w = (r \cdot \exp (+i \cdot h))))) \end{array}}}
\vnbprop{cc-exp-surjective-strip}{\vnbstmt{\begin{array}{@{}l@{}} \forall w \in \mathbb{C}. \\ \quad \neg (w = 0) \Rightarrow  \\ \quad \quad \exists z \in \mathbb{C}.\; ((-\pi \leq \operatorname{imag\text{-}part}(z)) \wedge ((\operatorname{imag\text{-}part}(z) < \pi) \wedge (\exp z = w))) \end{array}}}
\vnbprop{cc-exp-bijection-strip}{\vnbstmt{\begin{array}{@{}l@{}} (\lambda z \in \{q \in \mathbb{C} : ((-\pi \leq \operatorname{imag\text{-}part}(q)) \wedge (\operatorname{imag\text{-}part}(q) < \pi))\}.\; \exp z) \\ \in\; \mathrm{Bij}(\{q \in \mathbb{C} : ((-\pi \leq \operatorname{imag\text{-}part}(q)) \wedge (\operatorname{imag\text{-}part}(q) < \pi))\}, \operatorname{difference}(\mathbb{C}, \operatorname{singleton}(0))) \end{array}}}
\vnbprop{cc-exp-open-strip-avoids-negatives}{\vnbstmt{\begin{array}{@{}l@{}} \forall z \in \mathbb{C}. \\ \quad ((-\pi < \operatorname{imag\text{-}part}(z)) \wedge (\operatorname{imag\text{-}part}(z) < \pi)) \Rightarrow  \\ \quad \quad \neg ((\operatorname{imag\text{-}part}(\exp z) = 0) \wedge (\operatorname{real\text{-}part}(\exp z) \leq 0)) \end{array}}}
\vnbprop{cc-exp-open-strip-onto}{\vnbstmt{\begin{array}{@{}l@{}} \forall w \in \mathbb{C}. \\ \quad \neg ((\operatorname{imag\text{-}part}(w) = 0) \wedge (\operatorname{real\text{-}part}(w) \leq 0)) \Rightarrow  \\ \quad \quad \exists z \in \mathbb{C}.\; ((-\pi < \operatorname{imag\text{-}part}(z)) \wedge ((\operatorname{imag\text{-}part}(z) < \pi) \wedge (\exp z = w))) \end{array}}}
\vnbprop{cc-exp-bijection-open-strip}{\vnbstmt{\begin{array}{@{}l@{}} (\lambda z \in \{q \in \mathbb{C} : ((-\pi < \operatorname{imag\text{-}part}(q)) \wedge (\operatorname{imag\text{-}part}(q) < \pi))\}.\; \exp z) \\ \in\; \mathrm{Bij}(\{q \in \mathbb{C} : ((-\pi < \operatorname{imag\text{-}part}(q)) \wedge (\operatorname{imag\text{-}part}(q) < \pi))\}, \{u \in \mathbb{C} : \neg ((\operatorname{imag\text{-}part}(u) = 0) \wedge (\operatorname{real\text{-}part}(u) \leq 0))\}) \end{array}}}

\vnbbookfile{theorem-library/analytic-log-laws}{theorem-library/analytic-log-laws}
\noindent{\small analytic-log-laws.scm -- ANALYTICITY AT A CENTRE (Def 2.11, Cor 2.14 at a), THE PRINCIPAL LOGARITHM, LEMMA 2.20 AND PROPOSITION 2.21. Batch 38, 2026-09-27.  Definitions: structure-library/analytic-log.scm (IS-ANALYTIC-AT, CC-LOG).  The source: scratchpad/triage/ca.txt 1142-1150}\par
\vnbprop{alg-shift-in-ball}{\vnbstmt{\begin{array}{@{}l@{}} \forall a \in \mathbb{C}, r \in \mathbb{R}, z \in \operatorname{ball}(\operatorname{cc\text{-}ms}, a, r). \\ \quad (z \in \mathbb{C}) \wedge ((z - a) \in \operatorname{ball}(\operatorname{cc\text{-}ms}, 0, r)) \end{array}}}
\vnbprop{alg-shift-diff}{\vnbstmt{\begin{array}{@{}l@{}} \forall a \in \mathbb{C}, r \in \mathbb{R}, w \in \operatorname{ball}(\operatorname{cc\text{-}ms}, a, r). \\ \quad (0 < r) \Rightarrow  \\ \quad \quad \operatorname{is\text{-}diff\text{-}on}(\operatorname{cc\text{-}normed\text{-}field}, \operatorname{ball}(\operatorname{cc\text{-}ms}, a, r), (\lambda x \in \operatorname{ball}(\operatorname{cc\text{-}ms}, a, r).\; (x - a)), w, 1) \end{array}}}
\vnbprop{cps-translate-diff}{\vnbstmt{\begin{array}{@{}l@{}} \forall c, d \in (\mathbb{N} \to \mathbb{C}), a \in \mathbb{C}, r \in \mathbb{R}, w \in \operatorname{ball}(\operatorname{cc\text{-}ms}, a, r). \\ \quad (\forall j \in \mathbb{N}.\; (d(j) = ((j + 1) \cdot c((j + 1)))) \wedge ((0 < r) \wedge (r < \operatorname{cps\text{-}radius}(c)))) \Rightarrow  \\ \quad \quad \operatorname{is\text{-}diff\text{-}on}(\operatorname{cc\text{-}normed\text{-}field}, \operatorname{ball}(\operatorname{cc\text{-}ms}, a, r), \\ \quad \quad \quad (\lambda z \in \operatorname{ball}(\operatorname{cc\text{-}ms}, a, r).\; \operatorname{cc\text{-}series\text{-}limit}((\lambda k \in \mathbb{N}.\; (c(k) \cdot \operatorname{power}((z - a), k))))), \\ \quad \quad \quad w, \\ \quad \quad \quad \operatorname{cc\text{-}series\text{-}limit}((\lambda k \in \mathbb{N}.\; (d(k) \cdot \operatorname{power}((w - a), k))))) \end{array}}}
\vnbprop{cps-translate-holomorphic}{\vnbstmt{\begin{array}{@{}l@{}} \forall c \in (\mathbb{N} \to \mathbb{C}), a \in \mathbb{C}, r \in \mathbb{R}. \\ \quad ((0 < r) \wedge (r < \operatorname{cps\text{-}radius}(c))) \Rightarrow  \\ \quad \quad \operatorname{holomorphic\text{-}on}(\operatorname{ball}(\operatorname{cc\text{-}ms}, a, r), \\ \quad \quad \quad (\lambda z \in \operatorname{ball}(\operatorname{cc\text{-}ms}, a, r).\; \operatorname{cc\text{-}series\text{-}limit}((\lambda k \in \mathbb{N}.\; (c(k) \cdot \operatorname{power}((z - a), k)))))) \end{array}}}
\vnbprop{cps-translate-nth-deriv}{\vnbstmt{\begin{array}{@{}l@{}} \forall c \in (\mathbb{N} \to \mathbb{C}), a \in \mathbb{C}, r \in \mathbb{R}, k \in \mathbb{N}, g \in (\mathbb{N} \to \mathbb{C}). \\ \quad 0 < r \wedge \\ \quad r < \operatorname{cps\text{-}radius}(c) \wedge \\ \quad \forall n \in \mathbb{N}.\; (g(n) = ((\operatorname{factorial}((n + k)) \cdot \operatorname{recip}(\operatorname{factorial}(n))) \cdot c((n + k)))) \Rightarrow \\ \quad \quad (\lambda z \in \operatorname{ball}(\operatorname{cc\text{-}ms}, a, r).\; \operatorname{cc\text{-}series\text{-}limit}((\lambda \mathit{cpsk} \in \mathbb{N}.\; (c(\mathit{cpsk}) \cdot \operatorname{power}((z - a), \mathit{cpsk})))))^{(k)} \\ \quad \quad =\; (\lambda z \in \operatorname{ball}(\operatorname{cc\text{-}ms}, a, r).\; \operatorname{cc\text{-}series\text{-}limit}((\lambda \mathit{cpsk} \in \mathbb{N}.\; (g(\mathit{cpsk}) \cdot \operatorname{power}((z - a), \mathit{cpsk}))))) \end{array}}}
\vnbprop{taylor-coefficients-at}{\vnbstmt{\begin{array}{@{}l@{}} \forall f, u.; \forall c \in (\mathbb{N} \to \mathbb{C}), a \in \mathbb{C}, r \in \mathbb{R}, k \in \mathbb{N}. \\ \quad f \in (u \to \mathbb{C}) \wedge \\ \quad 0 < r \wedge \\ \quad r < \operatorname{cps\text{-}radius}(c) \wedge \\ \quad \operatorname{ball}(\operatorname{cc\text{-}ms}, a, r) \subseteq u \wedge \\ \quad \forall z \in \operatorname{ball}(\operatorname{cc\text{-}ms}, a, r). \\ \quad \quad f(z) = \operatorname{cc\text{-}series\text{-}limit}((\lambda \mathit{cpsk} \in \mathbb{N}.\; (c(\mathit{cpsk}) \cdot \operatorname{power}((z - a), \mathit{cpsk})))) \Rightarrow \\ \quad \quad \operatorname{restrict}(f, \operatorname{ball}(\operatorname{cc\text{-}ms}, a, r))^{(k)}(a) \\ \quad \quad =\; \operatorname{factorial}(k) \cdot c(k) \end{array}}}
\vnbprop{analytic-rep-holomorphic}{\vnbstmt{\begin{array}{@{}l@{}} \forall f, u.; \forall c \in (\mathbb{N} \to \mathbb{C}), a \in \mathbb{C}, r \in \mathbb{R}. \\ \quad f \in (u \to \mathbb{C}) \wedge \\ \quad 0 < r \wedge \\ \quad r < \operatorname{cps\text{-}radius}(c) \wedge \\ \quad \operatorname{ball}(\operatorname{cc\text{-}ms}, a, r) \subseteq u \wedge \\ \quad \forall z \in \operatorname{ball}(\operatorname{cc\text{-}ms}, a, r). \\ \quad \quad f(z) = \operatorname{cc\text{-}series\text{-}limit}((\lambda k \in \mathbb{N}.\; (c(k) \cdot \operatorname{power}((z - a), k)))) \Rightarrow \\ \quad \quad \operatorname{holomorphic\text{-}on}(\operatorname{ball}(\operatorname{cc\text{-}ms}, a, r), \operatorname{restrict}(f, \operatorname{ball}(\operatorname{cc\text{-}ms}, a, r))) \end{array}}}
\vnbprop{analytic-deriv}{\vnbstmt{\begin{array}{@{}l@{}} \forall f, u.; \forall c \in (\mathbb{N} \to \mathbb{C}), a \in \mathbb{C}, r \in \mathbb{R}. \\ \quad f \in (u \to \mathbb{C}) \wedge \\ \quad 0 < r \wedge \\ \quad r < \operatorname{cps\text{-}radius}(c) \wedge \\ \quad \operatorname{ball}(\operatorname{cc\text{-}ms}, a, r) \subseteq u \wedge \\ \quad \forall z \in \operatorname{ball}(\operatorname{cc\text{-}ms}, a, r). \\ \quad \quad f(z) = \operatorname{cc\text{-}series\text{-}limit}((\lambda k \in \mathbb{N}.\; (c(k) \cdot \operatorname{power}((z - a), k)))) \wedge \\ \quad \operatorname{is\text{-}open}(\operatorname{cc\text{-}ms}, u) \Rightarrow \\ \quad \quad \operatorname{is\text{-}diff\text{-}on}(\operatorname{cc\text{-}normed\text{-}field}, u, f, a, c(1)) \end{array}}}
\vnbprop{analytic-at-intro}{\vnbstmt{\begin{array}{@{}l@{}} \forall f, u.; \forall c \in (\mathbb{N} \to \mathbb{C}), a \in \mathbb{C}, r \in \mathbb{R}. \\ \quad f \in (u \to \mathbb{C}) \wedge \\ \quad 0 < r \wedge \\ \quad r < \operatorname{cps\text{-}radius}(c) \wedge \\ \quad \operatorname{ball}(\operatorname{cc\text{-}ms}, a, r) \subseteq u \wedge \\ \quad \forall z \in \operatorname{ball}(\operatorname{cc\text{-}ms}, a, r). \\ \quad \quad f(z) = \operatorname{cc\text{-}series\text{-}limit}((\lambda k \in \mathbb{N}.\; (c(k) \cdot \operatorname{power}((z - a), k)))) \wedge \\ \quad \operatorname{is\text{-}open}(\operatorname{cc\text{-}ms}, u) \Rightarrow \\ \quad \quad f \text{ is analytic at } a \end{array}}}
\vnbprop{analytic-holomorphic-ball}{\vnbstmt{\begin{array}{@{}l@{}} \forall f, u, a. \\ \quad f \text{ is analytic at } a \Rightarrow  \\ \quad \quad \exists r \in \mathbb{R}. \\ \quad \quad \quad 0 < r \wedge \\ \quad \quad \quad \operatorname{ball}(\operatorname{cc\text{-}ms}, a, r) \subseteq u \wedge \\ \quad \quad \quad \operatorname{holomorphic\text{-}on}(\operatorname{ball}(\operatorname{cc\text{-}ms}, a, r), \operatorname{restrict}(f, \operatorname{ball}(\operatorname{cc\text{-}ms}, a, r))) \end{array}}}
\vnbprop{cc-log-unique}{\vnbstmt{\begin{array}{@{}l@{}} \forall a, b \in \mathbb{C}. \\ \quad -\pi < \operatorname{imag\text{-}part}(a) \wedge \\ \quad \operatorname{imag\text{-}part}(a) < \pi \wedge \\ \quad -\pi < \operatorname{imag\text{-}part}(b) \wedge \\ \quad \operatorname{imag\text{-}part}(b) < \pi \wedge \\ \quad \exp a = \exp b \Rightarrow \\ \quad \quad a = b \end{array}}}
\vnbprop{cc-log-spec}{\vnbstmt{\begin{array}{@{}l@{}} \forall w \in \mathbb{C}. \\ \quad \neg ((\operatorname{imag\text{-}part}(w) = 0) \wedge (\operatorname{real\text{-}part}(w) \leq 0)) \Rightarrow  \\ \quad \quad \log\left(w\right) \in \mathbb{C} \wedge \\ \quad \quad -\pi < \operatorname{imag\text{-}part}(\log\left(w\right)) \wedge \\ \quad \quad \operatorname{imag\text{-}part}(\log\left(w\right)) < \pi \wedge \\ \quad \quad \exp \log\left(w\right) = w \end{array}}}
\vnbprop{cc-log-in-slit}{\vnbstmt{\begin{array}{@{}l@{}} \forall w \in \mathbb{C}. \\ \quad \neg ((\operatorname{imag\text{-}part}(w) = 0) \wedge (\operatorname{real\text{-}part}(w) \leq 0)) \Rightarrow  \\ \quad \quad (\log\left(w\right) \in \mathbb{C}) \wedge ((-\pi < \operatorname{imag\text{-}part}(\log\left(w\right))) \wedge (\operatorname{imag\text{-}part}(\log\left(w\right)) < \pi)) \end{array}}}
\vnbprop{cc-exp-log}{\vnbstmt{\begin{array}{@{}l@{}} \forall w \in \mathbb{C}. \\ \quad \neg ((\operatorname{imag\text{-}part}(w) = 0) \wedge (\operatorname{real\text{-}part}(w) \leq 0)) \Rightarrow (\exp \log\left(w\right) = w) \end{array}}}
\vnbprop{cc-log-char}{\vnbstmt{\begin{array}{@{}l@{}} \forall w, z \in \mathbb{C}. \\ \quad ((-\pi < \operatorname{imag\text{-}part}(z)) \wedge ((\operatorname{imag\text{-}part}(z) < \pi) \wedge (\exp z = w))) \Rightarrow (\log\left(w\right) = z) \end{array}}}
\vnbprop{cc-log-exp}{\vnbstmt{\begin{array}{@{}l@{}} \forall z \in \mathbb{C}. \\ \quad ((-\pi < \operatorname{imag\text{-}part}(z)) \wedge (\operatorname{imag\text{-}part}(z) < \pi)) \Rightarrow (\log\left(\exp z\right) = z) \end{array}}}
\vnbprop{cc-log-real}{\vnbstmt{\forall t \in \mathbb{R}.\; ((0 < t) \Rightarrow (\log\left(t\right) = \log\left(t\right)))}}
\vnbprop{cc-log-one}{\vnbstmt{\log\left(1\right) = 0}}
\vnbprop{cc-log-polar}{\vnbstmt{\begin{array}{@{}l@{}} \forall r, h \in \mathbb{R}. \\ \quad ((0 < r) \wedge ((-\pi < h) \wedge (h < \pi))) \Rightarrow  \\ \quad \quad \log\left((r \cdot \exp (+i \cdot h))\right) = (\log\left(r\right) + (+i \cdot h)) \end{array}}}
\vnbprop{cc-log-scale}{\vnbstmt{\begin{array}{@{}l@{}} \forall s \in \mathbb{R}, w \in \mathbb{C}. \\ \quad ((0 < s) \wedge \neg ((\operatorname{imag\text{-}part}(w) = 0) \wedge (\operatorname{real\text{-}part}(w) \leq 0))) \Rightarrow  \\ \quad \quad \log\left((s \cdot w)\right) = (\log\left(s\right) + \log\left(w\right)) \end{array}}}
\vnbprop{cc-geometric-alternating}{\vnbstmt{\begin{array}{@{}l@{}} \forall g \in (\mathbb{N} \to \mathbb{C}), z \in \mathbb{C}. \\ \quad (\forall j \in \mathbb{N}.\; (g(j) = \operatorname{power}(-1, j)) \wedge (\operatorname{magnitude}(z) < 1)) \Rightarrow  \\ \quad \quad \operatorname{cc\text{-}series\text{-}converges}((\lambda k \in \mathbb{N}.\; (g(k) \cdot \operatorname{power}(z, k)))) \wedge \\ \quad \quad \operatorname{cc\text{-}series\text{-}limit}((\lambda k \in \mathbb{N}.\; (g(k) \cdot \operatorname{power}(z, k)))) = \operatorname{recip}((1 + z)) \end{array}}}
\vnbprop{cps-log-radius}{\vnbstmt{\begin{array}{@{}l@{}} \forall c \in (\mathbb{N} \to \mathbb{C}), r \in \mathbb{R}. \\ \quad ((c(0) = 0) \wedge (\forall j \in \mathbb{N}.\; (c((j + 1)) = (\operatorname{power}(-1, j) \cdot \operatorname{recip}((j + 1)))) \wedge ((0 \leq r) \wedge (r < 1)))) \Rightarrow  \\ \quad \quad r < \operatorname{cps\text{-}radius}(c) \end{array}}}
\vnbprop{cps-log-series}{\vnbstmt{\begin{array}{@{}l@{}} \forall c \in (\mathbb{N} \to \mathbb{C}), r \in \mathbb{R}, z \in \operatorname{ball}(\operatorname{cc\text{-}ms}, 0, r). \\ \quad ((c(0) = 0) \wedge (\forall j \in \mathbb{N}.\; (c((j + 1)) = (\operatorname{power}(-1, j) \cdot \operatorname{recip}((j + 1)))) \wedge ((0 < r) \wedge (r < 1)))) \Rightarrow  \\ \quad \quad \operatorname{is\text{-}diff\text{-}on}(\operatorname{cc\text{-}normed\text{-}field}, \operatorname{ball}(\operatorname{cc\text{-}ms}, 0, r), \\ \quad \quad \quad (\lambda \mathit{cpdz} \in \operatorname{ball}(\operatorname{cc\text{-}ms}, 0, r).\; \operatorname{cc\text{-}series\text{-}limit}((\lambda k \in \mathbb{N}.\; (c(k) \cdot \operatorname{power}(\mathit{cpdz}, k))))), \\ \quad \quad \quad z, \operatorname{recip}((1 + z))) \end{array}}}
\vnbprop{alg-zero-deriv-ball}{\vnbstmt{\begin{array}{@{}l@{}} \forall a \in \mathbb{C}, r \in \mathbb{R}, h \in (\operatorname{ball}(\operatorname{cc\text{-}ms}, a, r) \to \mathbb{C}), z \in \operatorname{ball}(\operatorname{cc\text{-}ms}, a, r). \\ \quad 0 < r \wedge \\ \quad \forall \mathit{hbz} \in \operatorname{ball}(\operatorname{cc\text{-}ms}, a, r). \\ \quad \quad \operatorname{is\text{-}diff\text{-}on}(\operatorname{cc\text{-}normed\text{-}field}, \operatorname{ball}(\operatorname{cc\text{-}ms}, a, r), h, \mathit{hbz}, 0) \Rightarrow \\ \quad \quad h(z) = h(a) \end{array}}}
\vnbprop{cps-log-deriv-at-one}{\vnbstmt{\begin{array}{@{}l@{}} \forall c \in (\mathbb{N} \to \mathbb{C}), r \in \mathbb{R}, w \in \operatorname{ball}(\operatorname{cc\text{-}ms}, 1, r). \\ \quad ((c(0) = 0) \wedge (\forall j \in \mathbb{N}.\; (c((j + 1)) = (\operatorname{power}(-1, j) \cdot \operatorname{recip}((j + 1)))) \wedge ((0 < r) \wedge (r < 1)))) \Rightarrow  \\ \quad \quad \operatorname{is\text{-}diff\text{-}on}(\operatorname{cc\text{-}normed\text{-}field}, \operatorname{ball}(\operatorname{cc\text{-}ms}, 1, r), \\ \quad \quad \quad (\lambda z \in \operatorname{ball}(\operatorname{cc\text{-}ms}, 1, r).\; \operatorname{cc\text{-}series\text{-}limit}((\lambda k \in \mathbb{N}.\; (c(k) \cdot \operatorname{power}((z - 1), k))))), \\ \quad \quad \quad w, \operatorname{recip}(w)) \end{array}}}
\vnbprop{log0-exp}{\vnbstmt{\begin{array}{@{}l@{}} \forall c \in (\mathbb{N} \to \mathbb{C}), r \in \mathbb{R}, z \in \operatorname{ball}(\operatorname{cc\text{-}ms}, 1, r). \\ \quad ((c(0) = 0) \wedge (\forall j \in \mathbb{N}.\; (c((j + 1)) = (\operatorname{power}(-1, j) \cdot \operatorname{recip}((j + 1)))) \wedge ((0 < r) \wedge (r < 1)))) \Rightarrow  \\ \quad \quad \exp \operatorname{cc\text{-}series\text{-}limit}((\lambda k \in \mathbb{N}.\; (c(k) \cdot \operatorname{power}((z - 1), k)))) = z \end{array}}}
\vnbprop{log0-bound}{\vnbstmt{\begin{array}{@{}l@{}} \forall c \in (\mathbb{N} \to \mathbb{C}), r \in \mathbb{R}, z \in \operatorname{ball}(\operatorname{cc\text{-}ms}, 1, r). \\ \quad ((c(0) = 0) \wedge (\forall j \in \mathbb{N}.\; (c((j + 1)) = (\operatorname{power}(-1, j) \cdot \operatorname{recip}((j + 1)))) \wedge ((0 < r) \wedge (r < 1)))) \Rightarrow  \\ \quad \quad \operatorname{magnitude}(\operatorname{cc\text{-}series\text{-}limit}((\lambda k \in \mathbb{N}.\; (c(k) \cdot \operatorname{power}((z - 1), k))))) \leq (\operatorname{recip}((1 - \operatorname{magnitude}((z - 1)))) - 1) \end{array}}}
\vnbprop{log0-is-log-on-ball}{\vnbstmt{\begin{array}{@{}l@{}} \forall c \in (\mathbb{N} \to \mathbb{C}), r \in \mathbb{R}, z \in \operatorname{ball}(\operatorname{cc\text{-}ms}, 1, r). \\ \quad c(0) = 0 \wedge \\ \quad \forall j \in \mathbb{N}.\; (c((j + 1)) = (\operatorname{power}(-1, j) \cdot \operatorname{recip}((j + 1)))) \wedge \\ \quad 0 < r \wedge \\ \quad (r \cdot (1 + \pi)) \leq \pi \Rightarrow \\ \quad \quad \log\left(z\right) = \operatorname{cc\text{-}series\text{-}limit}((\lambda k \in \mathbb{N}.\; (c(k) \cdot \operatorname{power}((z - 1), k)))) \end{array}}}
\vnbprop{cc-log-holomorphic-near-one}{\vnbstmt{\begin{array}{@{}l@{}} \forall r \in \mathbb{R}. \\ \quad ((0 < r) \wedge ((r \cdot (1 + \pi)) \leq \pi)) \Rightarrow  \\ \quad \quad \operatorname{holomorphic\text{-}on}(\operatorname{ball}(\operatorname{cc\text{-}ms}, 1, r), (\lambda z \in \operatorname{ball}(\operatorname{cc\text{-}ms}, 1, r).\; \log\left(z\right))) \end{array}}}
\vnbprop{cc-log-one-plus}{\vnbstmt{\begin{array}{@{}l@{}} \forall c \in (\mathbb{N} \to \mathbb{C}), r \in \mathbb{R}, z \in \mathbb{C}. \\ \quad c(0) = 0 \wedge \\ \quad \forall j \in \mathbb{N}.\; (c((j + 1)) = (\operatorname{power}(-1, j) \cdot \operatorname{recip}((j + 1)))) \wedge \\ \quad 0 < r \wedge \\ \quad (r \cdot (1 + \pi)) \leq \pi \wedge \\ \quad \operatorname{magnitude}(z) < r \Rightarrow \\ \quad \quad \log\left((1 + z)\right) = \operatorname{cc\text{-}series\text{-}limit}((\lambda k \in \mathbb{N}.\; (c(k) \cdot \operatorname{power}(z, k)))) \end{array}}}
\vnbprop{holomorphic-compose}{\vnbstmt{\begin{array}{@{}l@{}} \forall u, v, f, g. \\ \quad (\operatorname{holomorphic\text{-}on}(u, f) \wedge (\operatorname{holomorphic\text{-}on}(v, g) \wedge \forall y \in u.\; (f(y) \in v))) \Rightarrow  \\ \quad \quad \operatorname{holomorphic\text{-}on}(u, (\lambda x \in u.\; g(f(x)))) \end{array}}}
\vnbprop{cc-log-holomorphic-scale}{\vnbstmt{\begin{array}{@{}l@{}} \forall r, s \in \mathbb{R}. \\ \quad ((0 < r) \wedge (((r \cdot (1 + \pi)) \leq \pi) \wedge (0 < s))) \Rightarrow  \\ \quad \quad \operatorname{holomorphic\text{-}on}(\operatorname{ball}(\operatorname{cc\text{-}ms}, s, (s \cdot r)), (\lambda z \in \operatorname{ball}(\operatorname{cc\text{-}ms}, s, (s \cdot r)).\; \log\left(z\right))) \end{array}}}
\vnbprop{log0-real}{\vnbstmt{\begin{array}{@{}l@{}} \forall c \in (\mathbb{N} \to \mathbb{C}), r \in \mathbb{R}, z \in \operatorname{ball}(\operatorname{cc\text{-}ms}, 1, r). \\ \quad c(0) = 0 \wedge \\ \quad \forall j \in \mathbb{N}.\; (c((j + 1)) = (\operatorname{power}(-1, j) \cdot \operatorname{recip}((j + 1)))) \wedge \\ \quad 0 < r \wedge \\ \quad r < 1 \wedge \\ \quad z \in \mathbb{R} \Rightarrow \\ \quad \quad \operatorname{cc\text{-}series\text{-}limit}((\lambda k \in \mathbb{N}.\; (c(k) \cdot \operatorname{power}((z - 1), k)))) \in \mathbb{C} \wedge \\ \quad \quad \operatorname{imag\text{-}part}(\operatorname{cc\text{-}series\text{-}limit}((\lambda k \in \mathbb{N}.\; (c(k) \cdot \operatorname{power}((z - 1), k))))) = 0 \end{array}}}
\vnbprop{log0-is-ln}{\vnbstmt{\begin{array}{@{}l@{}} \forall c \in (\mathbb{N} \to \mathbb{C}), t \in \mathbb{R}. \\ \quad ((c(0) = 0) \wedge (\forall j \in \mathbb{N}.\; (c((j + 1)) = (\operatorname{power}(-1, j) \cdot \operatorname{recip}((j + 1)))) \wedge ((0 < t) \wedge (t < 2)))) \Rightarrow  \\ \quad \quad \operatorname{cc\text{-}series\text{-}limit}((\lambda k \in \mathbb{N}.\; (c(k) \cdot \operatorname{power}((t - 1), k)))) = \log\left(t\right) \end{array}}}

\vnbbookfile{theorem-library/cc-log-holomorphic}{theorem-library/cc-log-holomorphic}
\noindent{\small cc-log-holomorphic.scm -- THE SIGN OF cos ON [-pi/2, pi], pi $>$ 3, AND THE HOLOMORPHY OF THE PRINCIPAL LOGARITHM ON THE SLIT PLANE C \textbackslash{} (-inf, 0]; the two loose ends of Lemma 2.20 (the radius is exactly 1) and Remark 2.22 (log(1 - z) = - sum z\textasciicircum{}k/k).}\par
\vnbprop{cc-cos-least-zero}{\vnbstmt{\forall t \in \mathbb{R}.\; (((0 < t) \wedge (t < (1/2 \cdot \pi))) \Rightarrow \neg (\cos t = 0))}}
\vnbprop{cc-cos-pos-right}{\vnbstmt{\forall s \in \mathbb{R}.\; (((0 < s) \wedge (s < (1/2 \cdot \pi))) \Rightarrow (0 < \cos s))}}
\vnbprop{cc-cos-pos-half}{\vnbstmt{\begin{array}{@{}l@{}} \forall t \in \mathbb{R}. \\ \quad ((-(1/2 \cdot \pi) < t) \wedge (t < (1/2 \cdot \pi))) \Rightarrow (0 < \cos t) \end{array}}}
\vnbprop{cc-cos-nonpos-far}{\vnbstmt{\forall t \in \mathbb{R}.\; ((((1/2 \cdot \pi) \leq t) \wedge (t \leq \pi)) \Rightarrow (\cos t \leq 0))}}
\vnbprop{cc-cos-pos-small}{\vnbstmt{\forall t \in \mathbb{R}.\; (((0 < t) \wedge (t < 1)) \Rightarrow (0 < \cos t))}}
\vnbprop{cc-pi-gt-two}{\vnbstmt{2 < \pi}}
\vnbprop{cc-cos-three-halves-pos}{\vnbstmt{0 < \cos 3/2}}
\vnbprop{cc-pi-gt-three}{\vnbstmt{3 < \pi}}
\vnbprop{cc-log-right-half}{\vnbstmt{\begin{array}{@{}l@{}} \forall w \in \mathbb{C}. \\ \quad (0 < \operatorname{real\text{-}part}(w)) \Rightarrow  \\ \quad \quad (-(1/2 \cdot \pi) < \operatorname{imag\text{-}part}(\log\left(w\right))) \wedge (\operatorname{imag\text{-}part}(\log\left(w\right)) < (1/2 \cdot \pi)) \end{array}}}
\vnbprop{cc-log-mul}{\vnbstmt{\begin{array}{@{}l@{}} \forall u, v \in \mathbb{C}. \\ \quad \neg ((\operatorname{imag\text{-}part}(u) = 0) \wedge (\operatorname{real\text{-}part}(u) \leq 0)) \wedge \\ \quad \neg ((\operatorname{imag\text{-}part}(v) = 0) \wedge (\operatorname{real\text{-}part}(v) \leq 0)) \wedge \\ \quad -\pi < (\operatorname{imag\text{-}part}(\log\left(u\right)) + \operatorname{imag\text{-}part}(\log\left(v\right))) \wedge \\ \quad (\operatorname{imag\text{-}part}(\log\left(u\right)) + \operatorname{imag\text{-}part}(\log\left(v\right))) < \pi \Rightarrow \\ \quad \quad \log\left((u \cdot v)\right) = (\log\left(u\right) + \log\left(v\right)) \end{array}}}
\vnbprop{cc-log-rotate}{\vnbstmt{\begin{array}{@{}l@{}} \forall w \in \mathbb{C}. \\ \quad (0 < \operatorname{real\text{-}part}(w)) \Rightarrow (\log\left((+i \cdot w)\right) = (\log\left(w\right) + (+i \cdot (1/2 \cdot \pi)))) \end{array}}}
\vnbprop{cc-log-rotate-neg}{\vnbstmt{\begin{array}{@{}l@{}} \forall w \in \mathbb{C}. \\ \quad (0 < \operatorname{real\text{-}part}(w)) \Rightarrow (\log\left((-+i \cdot w)\right) = (\log\left(w\right) - (+i \cdot (1/2 \cdot \pi)))) \end{array}}}
\vnbprop{cc-log-im-near-one}{\vnbstmt{\begin{array}{@{}l@{}} \forall d, r \in \mathbb{R}, z \in \operatorname{ball}(\operatorname{cc\text{-}ms}, 1, r). \\ \quad ((0 < d) \wedge ((0 < r) \wedge (((r \cdot (1 + \pi)) \leq \pi) \wedge ((r \cdot (1 + d)) \leq d)))) \Rightarrow  \\ \quad \quad \neg ((\operatorname{imag\text{-}part}(z) = 0) \wedge (\operatorname{real\text{-}part}(z) \leq 0)) \wedge \\ \quad \quad -d < \operatorname{imag\text{-}part}(\log\left(z\right)) \wedge \\ \quad \quad \operatorname{imag\text{-}part}(\log\left(z\right)) < d \end{array}}}
\vnbprop{cc-log-holomorphic-at}{\vnbstmt{\begin{array}{@{}l@{}} \forall a \in \mathbb{C}. \\ \quad \neg ((\operatorname{imag\text{-}part}(a) = 0) \wedge (\operatorname{real\text{-}part}(a) \leq 0)) \Rightarrow  \\ \quad \quad \exists p \in \mathbb{R}. \\ \quad \quad \quad 0 < p \wedge \\ \quad \quad \quad \operatorname{ball}(\operatorname{cc\text{-}ms}, a, p) \subseteq \{u \in \mathbb{C} : \neg ((\operatorname{imag\text{-}part}(u) = 0) \wedge (\operatorname{real\text{-}part}(u) \leq 0))\} \wedge \\ \quad \quad \quad \operatorname{holomorphic\text{-}on}(\operatorname{ball}(\operatorname{cc\text{-}ms}, a, p), (\lambda x \in \operatorname{ball}(\operatorname{cc\text{-}ms}, a, p).\; \log\left(x\right))) \end{array}}}
\vnbprop{cc-log-holomorphic}{\vnbstmt{\begin{array}{@{}l@{}} \operatorname{holomorphic\text{-}on}(\{u \in \mathbb{C} : \neg ((\operatorname{imag\text{-}part}(u) = 0) \wedge (\operatorname{real\text{-}part}(u) \leq 0))\}, \\ \quad (\lambda w \in \{u \in \mathbb{C} : \neg ((\operatorname{imag\text{-}part}(u) = 0) \wedge (\operatorname{real\text{-}part}(u) \leq 0))\}.\; \log\left(w\right))) \end{array}}}
\vnbprop{cc-log-deriv}{\vnbstmt{\begin{array}{@{}l@{}} \forall b \in \mathbb{C}. \\ \quad \neg ((\operatorname{imag\text{-}part}(b) = 0) \wedge (\operatorname{real\text{-}part}(b) \leq 0)) \Rightarrow  \\ \quad \quad \operatorname{is\text{-}diff\text{-}on}(\operatorname{cc\text{-}normed\text{-}field}, \\ \quad \quad \quad \{u \in \mathbb{C} : \neg ((\operatorname{imag\text{-}part}(u) = 0) \wedge (\operatorname{real\text{-}part}(u) \leq 0))\}, \\ \quad \quad \quad (\lambda w \in \{u \in \mathbb{C} : \neg ((\operatorname{imag\text{-}part}(u) = 0) \wedge (\operatorname{real\text{-}part}(u) \leq 0))\}.\; \log\left(w\right)), \\ \quad \quad \quad b, \operatorname{recip}(b)) \end{array}}}
\vnbprop{cc-log-one-minus}{\vnbstmt{\begin{array}{@{}l@{}} \forall r \in \mathbb{R}, z \in \mathbb{C}. \\ \quad ((0 < r) \wedge (((r \cdot (1 + \pi)) \leq \pi) \wedge (\operatorname{magnitude}(z) < r))) \Rightarrow  \\ \quad \quad \log\left((1 - z)\right) = -\operatorname{cc\text{-}series\text{-}limit}((\lambda k \in \mathbb{N}.\; ({k}^{\ast\!-1} \cdot \operatorname{power}(z, k)))) \end{array}}}
\vnbprop{cc-series-terms-small}{\vnbstmt{\begin{array}{@{}l@{}} \forall f \in (\mathbb{N} \to \mathbb{C}), e. \\ \quad (\operatorname{cc\text{-}series\text{-}converges}(f) \wedge \operatorname{pos\text{-}rr}(e)) \Rightarrow  \\ \quad \quad \exists n \in \mathbb{N}.\; \forall m \in \mathbb{N}.\; ((n \leq m) \Rightarrow (\operatorname{magnitude}(f(m)) \leq e)) \end{array}}}
\vnbprop{cps-log-radius-le-one}{\vnbstmt{\begin{array}{@{}l@{}} \forall c \in (\mathbb{N} \to \mathbb{C}), r \in \mathbb{R}. \\ \quad ((c(0) = 0) \wedge (\forall j \in \mathbb{N}.\; (c((j + 1)) = (\operatorname{power}(-1, j) \cdot \operatorname{recip}((j + 1)))) \wedge (1 < r))) \Rightarrow  \\ \quad \quad \neg (r < \operatorname{cps\text{-}radius}(c)) \end{array}}}
\vnbprop{cc-sin-pos-half}{\vnbstmt{\forall t \in \mathbb{R}.\; (((0 < t) \wedge (t < \pi)) \Rightarrow (0 < \sin t))}}
\vnbprop{cps-log-radius-one}{\vnbstmt{\begin{array}{@{}l@{}} \forall c \in (\mathbb{N} \to \mathbb{C}). \\ \quad ((c(0) = 0) \wedge \forall j \in \mathbb{N}.\; (c((j + 1)) = (\operatorname{power}(-1, j) \cdot \operatorname{recip}((j + 1))))) \Rightarrow  \\ \quad \quad \operatorname{cps\text{-}radius}(c) = 1 \end{array}}}

\vnbbookfile{theorem-library/ca-log-disc}{theorem-library/ca-log-disc}
\noindent{\small ca-log-disc.scm -- Remark/Lemma 2.22 of the complex analysis notes on the WHOLE unit disc (CA-MONTEL, 2026-10-26).}\par
\vnbprop{cc-log-of-deriv-on-ball}{\vnbstmt{\begin{array}{@{}l@{}} \forall r \in \mathbb{R}, g \in (\operatorname{ball}(\operatorname{cc\text{-}ms}, 1, r) \to \mathbb{C}), z \in \operatorname{ball}(\operatorname{cc\text{-}ms}, 1, r). \\ \quad 0 < r \wedge \\ \quad r \leq 1 \wedge \\ \quad \forall w \in \operatorname{ball}(\operatorname{cc\text{-}ms}, 1, r). \\ \quad \quad \operatorname{is\text{-}diff\text{-}on}(\operatorname{cc\text{-}normed\text{-}field}, \operatorname{ball}(\operatorname{cc\text{-}ms}, 1, r), g, w, \operatorname{recip}(w)) \wedge \\ \quad g(1) = 0 \Rightarrow \\ \quad \quad \log\left(z\right) = g(z) \end{array}}}
\vnbprop{log0-is-log-on-unit-ball}{\vnbstmt{\begin{array}{@{}l@{}} \forall c \in (\mathbb{N} \to \mathbb{C}), z \in \operatorname{ball}(\operatorname{cc\text{-}ms}, 1, 1). \\ \quad ((c(0) = 0) \wedge \forall j \in \mathbb{N}.\; (c((j + 1)) = (\operatorname{power}(-1, j) \cdot \operatorname{recip}((j + 1))))) \Rightarrow  \\ \quad \quad \log\left(z\right) = \operatorname{cc\text{-}series\text{-}limit}((\lambda k \in \mathbb{N}.\; (c(k) \cdot \operatorname{power}((z - 1), k)))) \end{array}}}
\vnbprop{cc-log-one-plus-disc}{\vnbstmt{\begin{array}{@{}l@{}} \forall c \in (\mathbb{N} \to \mathbb{C}), z \in \mathbb{C}. \\ \quad ((c(0) = 0) \wedge (\forall j \in \mathbb{N}.\; (c((j + 1)) = (\operatorname{power}(-1, j) \cdot \operatorname{recip}((j + 1)))) \wedge (\operatorname{magnitude}(z) < 1))) \Rightarrow  \\ \quad \quad \log\left((1 + z)\right) = \operatorname{cc\text{-}series\text{-}limit}((\lambda k \in \mathbb{N}.\; (c(k) \cdot \operatorname{power}(z, k)))) \end{array}}}
\vnbprop{cc-log-one-minus-disc}{\vnbstmt{\begin{array}{@{}l@{}} \forall z \in \mathbb{C}. \\ \quad (\operatorname{magnitude}(z) < 1) \Rightarrow  \\ \quad \quad \log\left((1 - z)\right) = -\operatorname{cc\text{-}series\text{-}limit}((\lambda k \in \mathbb{N}.\; ({k}^{\ast\!-1} \cdot \operatorname{power}(z, k)))) \end{array}}}

\vnbbookfile{structure-library/winding-number}{structure-library/winding-number}
\noindent{\small winding-number.scm -- THE WINDING NUMBER AND THE CIRCLE AS A ROAD (complex-analysis.pdf ch. 3: Definition 3.26, equation (73); Example 3.13). DEFINITIONS ONLY; every law is PROVEN in theorem-library/winding-number-laws.scm. Agent CA-1 of the October roadmap, 2026-10-03.}\par
\par\medskip\noindent\textbf{Definition} (\texttt{winding}).\;$\operatorname{winding}(\operatorname{wnp\_}, \operatorname{wnd\_}, \operatorname{wna\_}, \operatorname{wnb\_}, \operatorname{wnz\_}) :=$
\begin{equation*}
\operatorname{recip}(((2 \cdot \pi) \cdot +i)) \cdot \operatorname{line\text{-}int}((\lambda \operatorname{wnw\_} \in \operatorname{difference}(\mathbb{C}, \operatorname{singleton}(\operatorname{wnz\_})).\; \operatorname{recip}((\operatorname{wnw\_} - \operatorname{wnz\_}))), \operatorname{wnp\_}, \operatorname{wnd\_}, \operatorname{wna\_}, \operatorname{wnb\_})
\end{equation*}

\par\medskip\noindent\textbf{Definition} (\texttt{circle-path}).\;$\operatorname{circle\text{-}path}(\operatorname{wnc\_}, \operatorname{wnr\_}) :=$
\begin{equation*}
(\lambda \operatorname{wnt\_} \in \operatorname{ccint}(0, (2 \cdot \pi)).\; (\operatorname{wnc\_} + (\operatorname{wnr\_} \cdot \exp (+i \cdot \operatorname{wnt\_}))))
\end{equation*}

\par\medskip\noindent\textbf{Definition} (\texttt{circle-deriv}).\;$\operatorname{circle\text{-}deriv}(\operatorname{wnc\_}, \operatorname{wnr\_}) :=$
\begin{equation*}
(\lambda \operatorname{wnt\_} \in \operatorname{ccint}(0, (2 \cdot \pi)).\; (+i \cdot (\operatorname{wnr\_} \cdot \exp (+i \cdot \operatorname{wnt\_}))))
\end{equation*}


\vnbbookfile{structure-library/ch4a-vocab}{structure-library/ch4a-vocab}
\noindent{\small ch4a-vocab.scm -- WAVE 2 (2026-10-27..30), wired 2026-10-30 from scratchpad/ch4a/ch4a-vocab.scm. Library: complex-analysis.}\par
\par\medskip\noindent\textbf{Definition} (\texttt{annulus}).\;$\operatorname{annulus}(\operatorname{anla\_}, \operatorname{anlr\_}, \operatorname{anls\_}) :=$
\begin{equation*}
\{\operatorname{anlz\_} \in \mathbb{C} : ((\operatorname{anlr\_} < \operatorname{magnitude}((\operatorname{anlz\_} - \operatorname{anla\_}))) \wedge (\operatorname{magnitude}((\operatorname{anlz\_} - \operatorname{anla\_})) < \operatorname{anls\_}))\}
\end{equation*}

\par\medskip\noindent\textbf{Definition} (\texttt{closed-annulus}).\;$\operatorname{closed\text{-}annulus}(\operatorname{anla\_}, \operatorname{anlr\_}, \operatorname{anls\_}) :=$
\begin{equation*}
\{\operatorname{anlz\_} \in \mathbb{C} : ((\operatorname{anlr\_} \leq \operatorname{magnitude}((\operatorname{anlz\_} - \operatorname{anla\_}))) \wedge (\operatorname{magnitude}((\operatorname{anlz\_} - \operatorname{anla\_})) \leq \operatorname{anls\_}))\}
\end{equation*}

\par\medskip\noindent\textbf{Definition} (\texttt{laurent-sum}).\;$\operatorname{laurent\text{-}sum}(\operatorname{lsc\_}, \operatorname{lsa\_}, \operatorname{lsz\_}) :=$
\begin{equation*}
\operatorname{cc\text{-}series\text{-}limit}((\lambda \operatorname{lsk\_} \in \mathbb{N}.\; (\operatorname{lsc\_}(\operatorname{lsk\_}) \cdot \operatorname{power}((\operatorname{lsz\_} - \operatorname{lsa\_}), \operatorname{lsk\_})))) + (\operatorname{recip}((\operatorname{lsz\_} - \operatorname{lsa\_})) \cdot \operatorname{cc\text{-}series\text{-}limit}((\lambda \operatorname{lsk\_} \in \mathbb{N}.\; (\operatorname{lsc\_}(-(\operatorname{lsk\_} + 1)) \cdot \operatorname{power}(\operatorname{recip}((\operatorname{lsz\_} - \operatorname{lsa\_})), \operatorname{lsk\_})))))
\end{equation*}

\par\medskip\noindent\textbf{Definition} (\texttt{laurent-coefs}).\;$\operatorname{laurent\text{-}coefs}(\operatorname{lcf\_}, \operatorname{lcc\_}, \operatorname{lcr\_}) :=$
\begin{equation*}
\begin{array}{@{}l@{}} (\lambda \operatorname{lck\_} \in \mathbb{Z}.\; \\ \quad \left(\text{if }(0 \leq \operatorname{lck\_})\text{ then }(\operatorname{recip}(((2 \cdot \pi) \cdot +i)) \cdot \operatorname{line\text{-}int}((\lambda \operatorname{lcw\_} \in \operatorname{trace}(\operatorname{circle\text{-}path}(\operatorname{lcc\_}, \operatorname{lcr\_}), 0, (2 \cdot \pi)).\; (\operatorname{lcf\_}(\operatorname{lcw\_}) \cdot \operatorname{power}(\operatorname{recip}((\operatorname{lcw\_} - \operatorname{lcc\_})), (\operatorname{lck\_} + 1)))), \operatorname{circle\text{-}path}(\operatorname{lcc\_}, \operatorname{lcr\_}), \operatorname{circle\text{-}deriv}(\operatorname{lcc\_}, \operatorname{lcr\_}), 0, (2 \cdot \pi)))\text{ else }(\operatorname{recip}(((2 \cdot \pi) \cdot +i)) \cdot \operatorname{line\text{-}int}((\lambda \operatorname{lcw\_} \in \operatorname{trace}(\operatorname{circle\text{-}path}(\operatorname{lcc\_}, \operatorname{lcr\_}), 0, (2 \cdot \pi)).\; (\operatorname{lcf\_}(\operatorname{lcw\_}) \cdot \operatorname{power}((\operatorname{lcw\_} - \operatorname{lcc\_}), -(\operatorname{lck\_} + 1)))), \operatorname{circle\text{-}path}(\operatorname{lcc\_}, \operatorname{lcr\_}), \operatorname{circle\text{-}deriv}(\operatorname{lcc\_}, \operatorname{lcr\_}), 0, (2 \cdot \pi)))\right)) \end{array}
\end{equation*}


\vnbbookfile{structure-library/ch4b-vocab}{structure-library/ch4b-vocab}
\noindent{\small ch4b-vocab.scm -- WAVE 2 (2026-10-27..30), wired 2026-10-30 from scratchpad/ch4b/ch4b-vocab.scm. Library: complex-analysis.}\par
\par\medskip\noindent\textbf{Definition} (\texttt{is-meromorphic-at}).\;For all $\operatorname{mrf\_}$ and $\operatorname{mra\_}$, 
\begin{equation*}
\begin{array}{@{}l@{}} \operatorname{is\text{-}meromorphic\text{-}at}(\operatorname{mrf\_}, \operatorname{mra\_}) \\ \Leftrightarrow\; \operatorname{mra\_} \in \mathbb{C} \wedge \\ \exists \operatorname{mrn\_} \in \mathbb{N}, \operatorname{mrr\_} \in \mathbb{R}. \\ \quad 0 < \operatorname{mrr\_} \wedge \\ \quad \exists \operatorname{mrg\_}. \\ \quad \quad \operatorname{holomorphic\text{-}on}(\operatorname{ball}(\operatorname{cc\text{-}ms}, \operatorname{mra\_}, \operatorname{mrr\_}), \operatorname{mrg\_}) \wedge \\ \quad \quad \forall \operatorname{mrz\_} \in \operatorname{difference}(\operatorname{ball}(\operatorname{cc\text{-}ms}, \operatorname{mra\_}, \operatorname{mrr\_}), \operatorname{singleton}(\operatorname{mra\_})). \\ \quad \quad \quad \operatorname{mrf\_}(\operatorname{mrz\_}) = \cdot(\operatorname{power}(\operatorname{recip}(-(\operatorname{mrz\_}, \operatorname{mra\_})), \operatorname{mrn\_}), \operatorname{mrg\_}(\operatorname{mrz\_})) \end{array}
\end{equation*}

\par\medskip\noindent\textbf{Definition} (\texttt{is-pole}).\;For all $\operatorname{mrf\_}$, $\operatorname{mra\_}$ and $\operatorname{mrn\_}$, 
\begin{equation*}
\begin{array}{@{}l@{}} \operatorname{is\text{-}pole}(\operatorname{mrf\_}, \operatorname{mra\_}, \operatorname{mrn\_}) \\ \Leftrightarrow\; \operatorname{mra\_} \in \mathbb{C} \wedge \\ \operatorname{mrn\_} \in \mathbb{N} \wedge \\ 1 \leq \operatorname{mrn\_} \wedge \\ \exists \operatorname{mrr\_} \in \mathbb{R}. \\ \quad 0 < \operatorname{mrr\_} \wedge \\ \quad \exists \operatorname{mrg\_}. \\ \quad \quad \operatorname{holomorphic\text{-}on}(\operatorname{ball}(\operatorname{cc\text{-}ms}, \operatorname{mra\_}, \operatorname{mrr\_}), \operatorname{mrg\_}) \wedge \\ \quad \quad \neg (\operatorname{mrg\_}(\operatorname{mra\_}) = 0) \wedge \\ \quad \quad \forall \operatorname{mrz\_} \in \operatorname{difference}(\operatorname{ball}(\operatorname{cc\text{-}ms}, \operatorname{mra\_}, \operatorname{mrr\_}), \operatorname{singleton}(\operatorname{mra\_})). \\ \quad \quad \quad \operatorname{mrf\_}(\operatorname{mrz\_}) = \cdot(\operatorname{power}(\operatorname{recip}(-(\operatorname{mrz\_}, \operatorname{mra\_})), \operatorname{mrn\_}), \operatorname{mrg\_}(\operatorname{mrz\_})) \end{array}
\end{equation*}

\par\medskip\noindent\textbf{Definition} (\texttt{is-removable-singularity}).\;For all $\operatorname{mrf\_}$ and $\operatorname{mra\_}$, 
\begin{equation*}
\begin{array}{@{}l@{}} \operatorname{is\text{-}removable\text{-}singularity}(\operatorname{mrf\_}, \operatorname{mra\_}) \\ \Leftrightarrow\; \operatorname{mra\_} \in \mathbb{C} \wedge \\ \exists \operatorname{mrr\_} \in \mathbb{R}. \\ \quad 0 < \operatorname{mrr\_} \wedge \\ \quad \exists \operatorname{mrg\_}. \\ \quad \quad \operatorname{holomorphic\text{-}on}(\operatorname{ball}(\operatorname{cc\text{-}ms}, \operatorname{mra\_}, \operatorname{mrr\_}), \operatorname{mrg\_}) \wedge \\ \quad \quad \forall \operatorname{mrz\_} \in \operatorname{difference}(\operatorname{ball}(\operatorname{cc\text{-}ms}, \operatorname{mra\_}, \operatorname{mrr\_}), \operatorname{singleton}(\operatorname{mra\_})). \\ \quad \quad \quad \operatorname{mrf\_}(\operatorname{mrz\_}) = \operatorname{mrg\_}(\operatorname{mrz\_}) \end{array}
\end{equation*}

\par\medskip\noindent\textbf{Definition} (\texttt{is-essential-singularity}).\;For all $\operatorname{mrf\_}$ and $\operatorname{mra\_}$, 
\begin{equation*}
\begin{array}{@{}l@{}} \operatorname{is\text{-}essential\text{-}singularity}(\operatorname{mrf\_}, \operatorname{mra\_}) \\ \Leftrightarrow\; \operatorname{mra\_} \in \mathbb{C} \wedge \\ \neg \operatorname{is\text{-}removable\text{-}singularity}(\operatorname{mrf\_}, \operatorname{mra\_}) \wedge \\ \neg \exists \operatorname{mrn\_}.\; \operatorname{is\text{-}pole}(\operatorname{mrf\_}, \operatorname{mra\_}, \operatorname{mrn\_}) \end{array}
\end{equation*}

\par\medskip\noindent\textbf{Definition} (\texttt{meromorphic-on}).\;For all $\operatorname{mru\_}$ and $\operatorname{mrf\_}$, 
\begin{equation*}
\begin{array}{@{}l@{}} \operatorname{meromorphic\text{-}on}(\operatorname{mru\_}, \operatorname{mrf\_}) \\ \Leftrightarrow\; \operatorname{is\text{-}open}(\operatorname{nf\text{-}metric\text{-}space}(\operatorname{cc\text{-}normed\text{-}field}), \operatorname{mru\_}) \wedge \\ \exists \operatorname{mrs\_}. \\ \quad \operatorname{mrs\_} \subseteq \operatorname{mru\_} \wedge \\ \quad \operatorname{holomorphic\text{-}on}(\operatorname{difference}(\operatorname{mru\_}, \operatorname{mrs\_}), \operatorname{mrf\_}) \wedge \\ \quad \forall \operatorname{mrx\_} \in \operatorname{mrs\_}. \\ \quad \quad (\operatorname{is\text{-}removable\text{-}singularity}(\operatorname{mrf\_}, \operatorname{mrx\_}) \vee \exists \operatorname{mrn\_}.\; \operatorname{is\text{-}pole}(\operatorname{mrf\_}, \operatorname{mrx\_}, \operatorname{mrn\_})) \end{array}
\end{equation*}

\par\medskip\noindent\textbf{Definition} (\texttt{residue}).\;$\operatorname{residue}(\operatorname{mrf\_}, \operatorname{mra\_}) :=$
\begin{equation*}
(\iota\, \operatorname{mrx\_}.\;((\operatorname{mrx\_} \in \mathbb{C}) \wedge \exists \operatorname{mrr\_} \in \mathbb{R}.\; ((0 < \operatorname{mrr\_}) \wedge \forall \operatorname{mrp\_} \in \mathbb{R}.\; ((0 < \operatorname{mrp\_}) \Rightarrow ((\operatorname{mrp\_} < \operatorname{mrr\_}) \Rightarrow (\operatorname{mrx\_} = (\operatorname{recip}(((2 \cdot \pi) \cdot +i)) \cdot \operatorname{line\text{-}int}((\lambda \operatorname{mrw\_} \in \operatorname{trace}(\operatorname{circle\text{-}path}(\operatorname{mra\_}, \operatorname{mrp\_}), 0, (2 \cdot \pi)).\; \operatorname{mrf\_}(\operatorname{mrw\_})), \operatorname{circle\text{-}path}(\operatorname{mra\_}, \operatorname{mrp\_}), \operatorname{circle\text{-}deriv}(\operatorname{mra\_}, \operatorname{mrp\_}), 0, (2 \cdot \pi)))))))))
\end{equation*}


\vnbbookfile{structure-library/ch4d-vocab}{structure-library/ch4d-vocab}
\noindent{\small ch4d-vocab.scm -- WAVE 2 (2026-10-27..30), wired 2026-10-30 from scratchpad/ch4d/ch4d-vocab.scm. Library: complex-analysis.}\par
\par\medskip\noindent\textbf{Definition} (\texttt{pole-order}).\;$\operatorname{pole\text{-}order}(\operatorname{pof\_}, \operatorname{poa\_}) :=$ $(\iota\, \operatorname{pon\_}.\;\operatorname{is\text{-}pole}(\operatorname{pof\_}, \operatorname{poa\_}, \operatorname{pon\_}))$

\vnbbookfile{theorem-library/winding-number-laws}{theorem-library/winding-number-laws}
\noindent{\small winding-number-laws.scm -- CAUCHY ON A DISC, THE WINDING NUMBER, THE CIRCLE AS A ROAD (complex-analysis.pdf ch. 3: Corollary 3.9 for a disc inside an open set; Definition 3.26 and the disc case of Corollary 3.30; Example 3.13 at the centre). Agent CA-1 of the October roadmap (docs/roadmap-2026-10-03.md, line (a), ...}\par
\vnbprop{cc-ball-is-convex}{\vnbstmt{\forall c \in \mathbb{C}, r \in \mathbb{R}.\; \operatorname{is\text{-}convex\text{-}cc}(\operatorname{ball}(\operatorname{cc\text{-}ms}, c, r))}}
\vnbprop{line-int-domain-congruence}{\vnbstmt{\begin{array}{@{}l@{}} \forall u, f, m, c, a, b, v.; \forall g \in (v \to \mathbb{C}). \\ \quad u \subseteq \mathbb{C} \wedge \\ \quad f \in (u \to \mathbb{C}) \wedge \\ \quad \forall y \in u. \\ \quad \quad \operatorname{is\text{-}continuous\text{-}at}(\operatorname{subspace\text{-}ms}(\operatorname{nf\text{-}metric\text{-}space}(\operatorname{cc\text{-}normed\text{-}field}), u), \\ \quad \quad \quad \operatorname{nf\text{-}metric\text{-}space}(\operatorname{cc\text{-}normed\text{-}field}), f, y) \wedge \\ \quad \operatorname{is\text{-}road}(m, c, a, b) \wedge \\ \quad \operatorname{trace}(m, a, b) \subseteq u \wedge \\ \quad \operatorname{trace}(m, a, b) \subseteq v \wedge \\ \quad \forall x \in \operatorname{trace}(m, a, b).\; (g(x) \simeq f(x)) \Rightarrow \\ \quad \quad \operatorname{line\text{-}int}(g, m, c, a, b) = \operatorname{line\text{-}int}(f, m, c, a, b) \end{array}}}
\vnbprop{holomorphic-on-restrict}{\vnbstmt{\begin{array}{@{}l@{}} \forall u, f, v. \\ \quad (\operatorname{holomorphic\text{-}on}(u, f) \wedge (\operatorname{is\text{-}open}(\operatorname{nf\text{-}metric\text{-}space}(\operatorname{cc\text{-}normed\text{-}field}), v) \wedge (v \subseteq u))) \Rightarrow  \\ \quad \quad \operatorname{holomorphic\text{-}on}(v, \operatorname{restrict}(f, v)) \end{array}}}
\vnbprop{cauchy-convex-sub}{\vnbstmt{\begin{array}{@{}l@{}} \forall u, f, v, m, c, a, b. \\ \quad \operatorname{holomorphic\text{-}on}(u, f) \wedge \\ \quad \operatorname{is\text{-}open}(\operatorname{nf\text{-}metric\text{-}space}(\operatorname{cc\text{-}normed\text{-}field}), v) \wedge \\ \quad \operatorname{is\text{-}convex\text{-}cc}(v) \wedge \\ \quad v \subseteq u \wedge \\ \quad \operatorname{is\text{-}road}(m, c, a, b) \wedge \\ \quad \operatorname{trace}(m, a, b) \subseteq v \wedge \\ \quad m(b) = m(a) \Rightarrow \\ \quad \quad \operatorname{line\text{-}int}(f, m, c, a, b) = 0 \end{array}}}
\vnbprop{cauchy-disc}{\vnbstmt{\begin{array}{@{}l@{}} \forall u, f.; \forall c \in \mathbb{C}, r \in \mathbb{R}, m, d, a, b. \\ \quad \operatorname{holomorphic\text{-}on}(u, f) \wedge \\ \quad 0 < r \wedge \\ \quad \operatorname{ball}(\operatorname{cc\text{-}ms}, c, r) \subseteq u \wedge \\ \quad \operatorname{is\text{-}road}(m, d, a, b) \wedge \\ \quad \operatorname{trace}(m, a, b) \subseteq \operatorname{ball}(\operatorname{cc\text{-}ms}, c, r) \wedge \\ \quad m(b) = m(a) \Rightarrow \\ \quad \quad \operatorname{line\text{-}int}(f, m, d, a, b) = 0 \end{array}}}
\vnbprop{cc-punctured-membership}{\vnbstmt{\begin{array}{@{}l@{}} \forall z \in \mathbb{C}, y. \\ \quad (y \in \operatorname{difference}(\mathbb{C}, \operatorname{singleton}(z))) \Leftrightarrow ((y \in \mathbb{C}) \wedge \neg (y = z)) \end{array}}}
\vnbprop{cc-punctured-open}{\vnbstmt{\forall z \in \mathbb{C}.\; \operatorname{is\text{-}open}(\operatorname{cc\text{-}ms}, \operatorname{difference}(\mathbb{C}, \operatorname{singleton}(z)))}}
\vnbprop{winding-integrand-holomorphic}{\vnbstmt{\begin{array}{@{}l@{}} \forall z \in \mathbb{C}. \\ \quad \operatorname{holomorphic\text{-}on}(\operatorname{difference}(\mathbb{C}, \operatorname{singleton}(z)), \\ \quad \quad (\lambda w \in \operatorname{difference}(\mathbb{C}, \operatorname{singleton}(z)).\; \operatorname{recip}((w - z)))) \end{array}}}
\vnbprop{winding-in-cc}{\vnbstmt{\begin{array}{@{}l@{}} \forall m, c, a, b.; \forall z \in \mathbb{C}. \\ \quad (\operatorname{is\text{-}road}(m, c, a, b) \wedge \neg (z \in \operatorname{trace}(m, a, b))) \Rightarrow  \\ \quad \quad \operatorname{winding}(m, c, a, b, z) \in \mathbb{C} \end{array}}}
\vnbprop{winding-zero-outside-disc}{\vnbstmt{\begin{array}{@{}l@{}} \forall m, d, a, b.; \forall c \in \mathbb{C}, r \in \mathbb{R}, z \in \mathbb{C}. \\ \quad \operatorname{is\text{-}road}(m, d, a, b) \wedge \\ \quad m(b) = m(a) \wedge \\ \quad 0 < r \wedge \\ \quad \operatorname{trace}(m, a, b) \subseteq \operatorname{ball}(\operatorname{cc\text{-}ms}, c, r) \wedge \\ \quad \neg (z \in \operatorname{ball}(\operatorname{cc\text{-}ms}, c, r)) \Rightarrow \\ \quad \quad \operatorname{winding}(m, d, a, b, z) = 0 \end{array}}}
\vnbprop{winding-circle-center}{\vnbstmt{\begin{array}{@{}l@{}} \forall c \in \mathbb{C}, r \in \mathbb{R}. \\ \quad (0 < r) \Rightarrow (\operatorname{winding}(\operatorname{circle\text{-}path}(c, r), \operatorname{circle\text{-}deriv}(c, r), 0, (2 \cdot \pi), c) = 1) \end{array}}}
\vnbprop{road-holomorphic-image}{\vnbstmt{\begin{array}{@{}l@{}} \forall u, f, c, m, d, a, b. \\ \quad \operatorname{holomorphic\text{-}on}(u, f) \wedge \\ \quad c \in (u \to \mathbb{C}) \wedge \\ \quad \forall z \in u.\; \operatorname{is\text{-}diff\text{-}on}(\operatorname{cc\text{-}normed\text{-}field}, u, f, z, c(z)) \wedge \\ \quad \forall y \in u. \\ \quad \quad \operatorname{is\text{-}continuous\text{-}at}(\operatorname{subspace\text{-}ms}(\operatorname{nf\text{-}metric\text{-}space}(\operatorname{cc\text{-}normed\text{-}field}), u), \\ \quad \quad \quad \operatorname{nf\text{-}metric\text{-}space}(\operatorname{cc\text{-}normed\text{-}field}), c, y) \wedge \\ \quad \operatorname{is\text{-}road}(m, d, a, b) \wedge \\ \quad \operatorname{trace}(m, a, b) \subseteq u \Rightarrow \\ \quad \quad \operatorname{is\text{-}road}((\lambda t \in \operatorname{ccint}(a, b).\; f(m(t))), (\lambda t \in \operatorname{ccint}(a, b).\; (c(m(t)) \cdot d(t))), a, b) \end{array}}}
\vnbprop{circle-is-road}{\vnbstmt{\begin{array}{@{}l@{}} \forall c \in \mathbb{C}, r \in \mathbb{R}. \\ \quad \operatorname{is\text{-}road}(\operatorname{circle\text{-}path}(c, r), \operatorname{circle\text{-}deriv}(c, r), 0, (2 \cdot \pi)) \end{array}}}

\vnbbookfile{theorem-library/analytic-from-integral}{theorem-library/analytic-from-integral}
\noindent{\small analytic-from-integral.scm -- THE CAUCHY-TYPE INTEGRAL IS A POWER SERIES (complex-analysis.pdf 3.5: Lemma 3.33 with Remark 3.34, the engine of Theorem 3.35; 3.7: the Cauchy estimate behind Prop 3.47 / Cor 3.48). Agent CA-2 of the October roadmap (docs/roadmap-2026-10-03.md 0(a)), 2026-10-03.}\par
\vnbprop{aif-cont-product}{\vnbstmt{\begin{array}{@{}l@{}} \forall u, f, g, h.; \forall y \in u. \\ \quad u \subseteq \mathbb{C} \wedge \\ \quad f \in (u \to \mathbb{C}) \wedge \\ \quad \forall y \in u. \\ \quad \quad \operatorname{is\text{-}continuous\text{-}at}(\operatorname{subspace\text{-}ms}(\operatorname{nf\text{-}metric\text{-}space}(\operatorname{cc\text{-}normed\text{-}field}), u), \\ \quad \quad \quad \operatorname{nf\text{-}metric\text{-}space}(\operatorname{cc\text{-}normed\text{-}field}), f, y) \wedge \\ \quad g \in (u \to \mathbb{C}) \wedge \\ \quad \forall y \in u. \\ \quad \quad \operatorname{is\text{-}continuous\text{-}at}(\operatorname{subspace\text{-}ms}(\operatorname{nf\text{-}metric\text{-}space}(\operatorname{cc\text{-}normed\text{-}field}), u), \\ \quad \quad \quad \operatorname{nf\text{-}metric\text{-}space}(\operatorname{cc\text{-}normed\text{-}field}), g, y) \wedge \\ \quad h \in (u \to \mathbb{C}) \wedge \\ \quad \forall x \in u.\; (h(x) \simeq (f(x) \cdot g(x))) \Rightarrow \\ \quad \quad \operatorname{is\text{-}continuous\text{-}at}(\operatorname{subspace\text{-}ms}(\operatorname{nf\text{-}metric\text{-}space}(\operatorname{cc\text{-}normed\text{-}field}), u), \\ \quad \quad \quad \operatorname{nf\text{-}metric\text{-}space}(\operatorname{cc\text{-}normed\text{-}field}), h, y) \end{array}}}
\vnbprop{aif-cont-sum}{\vnbstmt{\begin{array}{@{}l@{}} \forall u, f, g, h.; \forall y \in u. \\ \quad u \subseteq \mathbb{C} \wedge \\ \quad f \in (u \to \mathbb{C}) \wedge \\ \quad \forall y \in u. \\ \quad \quad \operatorname{is\text{-}continuous\text{-}at}(\operatorname{subspace\text{-}ms}(\operatorname{nf\text{-}metric\text{-}space}(\operatorname{cc\text{-}normed\text{-}field}), u), \\ \quad \quad \quad \operatorname{nf\text{-}metric\text{-}space}(\operatorname{cc\text{-}normed\text{-}field}), f, y) \wedge \\ \quad g \in (u \to \mathbb{C}) \wedge \\ \quad \forall y \in u. \\ \quad \quad \operatorname{is\text{-}continuous\text{-}at}(\operatorname{subspace\text{-}ms}(\operatorname{nf\text{-}metric\text{-}space}(\operatorname{cc\text{-}normed\text{-}field}), u), \\ \quad \quad \quad \operatorname{nf\text{-}metric\text{-}space}(\operatorname{cc\text{-}normed\text{-}field}), g, y) \wedge \\ \quad h \in (u \to \mathbb{C}) \wedge \\ \quad \forall x \in u.\; (h(x) \simeq (f(x) + g(x))) \Rightarrow \\ \quad \quad \operatorname{is\text{-}continuous\text{-}at}(\operatorname{subspace\text{-}ms}(\operatorname{nf\text{-}metric\text{-}space}(\operatorname{cc\text{-}normed\text{-}field}), u), \\ \quad \quad \quad \operatorname{nf\text{-}metric\text{-}space}(\operatorname{cc\text{-}normed\text{-}field}), h, y) \end{array}}}
\vnbprop{aif-cont-holo-restrict}{\vnbstmt{\begin{array}{@{}l@{}} \forall v, g, u.; \forall h \in (u \to \mathbb{C}), y \in u. \\ \quad (\operatorname{holomorphic\text{-}on}(v, g) \wedge ((u \subseteq v) \wedge \forall x \in u.\; (h(x) \simeq g(x)))) \Rightarrow  \\ \quad \quad \operatorname{is\text{-}continuous\text{-}at}(\operatorname{subspace\text{-}ms}(\operatorname{nf\text{-}metric\text{-}space}(\operatorname{cc\text{-}normed\text{-}field}), u), \\ \quad \quad \quad \operatorname{nf\text{-}metric\text{-}space}(\operatorname{cc\text{-}normed\text{-}field}), h, y) \end{array}}}
\vnbprop{aif-cont-compose}{\vnbstmt{\begin{array}{@{}l@{}} \forall u.; \forall g \in (u \to \mathbb{C}), v, \mathit{aiph}.; \forall h \in (u \to \mathbb{C}), y \in u. \\ \quad u \subseteq \mathbb{C} \wedge \\ \quad \forall y \in u. \\ \quad \quad \operatorname{is\text{-}continuous\text{-}at}(\operatorname{subspace\text{-}ms}(\operatorname{nf\text{-}metric\text{-}space}(\operatorname{cc\text{-}normed\text{-}field}), u), \\ \quad \quad \quad \operatorname{nf\text{-}metric\text{-}space}(\operatorname{cc\text{-}normed\text{-}field}), g, y) \wedge \\ \quad v \subseteq \mathbb{C} \wedge \\ \quad \forall x \in u.\; (g(x) \in v) \wedge \\ \quad \forall y \in v. \\ \quad \quad \operatorname{is\text{-}continuous\text{-}at}(\operatorname{subspace\text{-}ms}(\operatorname{nf\text{-}metric\text{-}space}(\operatorname{cc\text{-}normed\text{-}field}), v), \\ \quad \quad \quad \operatorname{nf\text{-}metric\text{-}space}(\operatorname{cc\text{-}normed\text{-}field}), \mathit{aiph}, y) \wedge \\ \quad \forall x \in u.\; (h(x) \simeq \operatorname{aiph}(g(x))) \Rightarrow \\ \quad \quad \operatorname{is\text{-}continuous\text{-}at}(\operatorname{subspace\text{-}ms}(\operatorname{nf\text{-}metric\text{-}space}(\operatorname{cc\text{-}normed\text{-}field}), u), \\ \quad \quad \quad \operatorname{nf\text{-}metric\text{-}space}(\operatorname{cc\text{-}normed\text{-}field}), h, y) \end{array}}}
\vnbprop{aif-cont-affine}{\vnbstmt{\begin{array}{@{}l@{}} \forall u.; \forall p, q, w \in \mathbb{C}, h \in (u \to \mathbb{C}), y \in u. \\ \quad ((u \subseteq \mathbb{C}) \wedge \forall x \in u.\; (h(x) \simeq (p + (q \cdot (x - w))))) \Rightarrow  \\ \quad \quad \operatorname{is\text{-}continuous\text{-}at}(\operatorname{subspace\text{-}ms}(\operatorname{nf\text{-}metric\text{-}space}(\operatorname{cc\text{-}normed\text{-}field}), u), \\ \quad \quad \quad \operatorname{nf\text{-}metric\text{-}space}(\operatorname{cc\text{-}normed\text{-}field}), h, y) \end{array}}}
\vnbprop{aif-cont-shift}{\vnbstmt{\begin{array}{@{}l@{}} \forall u.; \forall c \in \mathbb{C}, h \in (u \to \mathbb{C}), y \in u. \\ \quad ((u \subseteq \mathbb{C}) \wedge \forall x \in u.\; (h(x) \simeq (x - c))) \Rightarrow  \\ \quad \quad \operatorname{is\text{-}continuous\text{-}at}(\operatorname{subspace\text{-}ms}(\operatorname{nf\text{-}metric\text{-}space}(\operatorname{cc\text{-}normed\text{-}field}), u), \\ \quad \quad \quad \operatorname{nf\text{-}metric\text{-}space}(\operatorname{cc\text{-}normed\text{-}field}), h, y) \end{array}}}
\vnbprop{aif-cont-const}{\vnbstmt{\begin{array}{@{}l@{}} \forall u.; \forall k \in \mathbb{C}, h \in (u \to \mathbb{C}), y \in u. \\ \quad ((u \subseteq \mathbb{C}) \wedge \forall x \in u.\; (h(x) \simeq k)) \Rightarrow  \\ \quad \quad \operatorname{is\text{-}continuous\text{-}at}(\operatorname{subspace\text{-}ms}(\operatorname{nf\text{-}metric\text{-}space}(\operatorname{cc\text{-}normed\text{-}field}), u), \\ \quad \quad \quad \operatorname{nf\text{-}metric\text{-}space}(\operatorname{cc\text{-}normed\text{-}field}), h, y) \end{array}}}
\vnbprop{aif-cont-recip-shift}{\vnbstmt{\begin{array}{@{}l@{}} \forall u.; \forall c \in \mathbb{C}, h \in (u \to \mathbb{C}), y \in u. \\ \quad ((u \subseteq \mathbb{C}) \wedge (\forall x \in u.\; \neg (x = c) \wedge \forall x \in u.\; (h(x) \simeq \operatorname{recip}((x - c))))) \Rightarrow  \\ \quad \quad \operatorname{is\text{-}continuous\text{-}at}(\operatorname{subspace\text{-}ms}(\operatorname{nf\text{-}metric\text{-}space}(\operatorname{cc\text{-}normed\text{-}field}), u), \\ \quad \quad \quad \operatorname{nf\text{-}metric\text{-}space}(\operatorname{cc\text{-}normed\text{-}field}), h, y) \end{array}}}
\vnbprop{aif-cont-power-of}{\vnbstmt{\begin{array}{@{}l@{}} \forall u.; \forall g \in (u \to \mathbb{C}), n \in \mathbb{N}, h \in (u \to \mathbb{C}), y \in u. \\ \quad u \subseteq \mathbb{C} \wedge \\ \quad \forall y \in u. \\ \quad \quad \operatorname{is\text{-}continuous\text{-}at}(\operatorname{subspace\text{-}ms}(\operatorname{nf\text{-}metric\text{-}space}(\operatorname{cc\text{-}normed\text{-}field}), u), \\ \quad \quad \quad \operatorname{nf\text{-}metric\text{-}space}(\operatorname{cc\text{-}normed\text{-}field}), g, y) \wedge \\ \quad \forall x \in u.\; (h(x) \simeq \operatorname{power}(g(x), n)) \Rightarrow \\ \quad \quad \operatorname{is\text{-}continuous\text{-}at}(\operatorname{subspace\text{-}ms}(\operatorname{nf\text{-}metric\text{-}space}(\operatorname{cc\text{-}normed\text{-}field}), u), \\ \quad \quad \quad \operatorname{nf\text{-}metric\text{-}space}(\operatorname{cc\text{-}normed\text{-}field}), h, y) \end{array}}}
\vnbprop{aif-magnitude-recip}{\vnbstmt{\forall d \in \mathbb{C}.\; (\neg (d = 0) \Rightarrow (\operatorname{magnitude}(\operatorname{recip}(d)) = \operatorname{recip}(\operatorname{magnitude}(d))))}}
\vnbprop{aif-geometric-partial}{\vnbstmt{\begin{array}{@{}l@{}} \forall n \in \mathbb{N}, p, v, \mathit{aiqv}, e. \\ \quad ((p \in \mathbb{C}) \wedge ((v \in \mathbb{C}) \wedge ((\mathit{aiqv} \in \mathbb{C}) \wedge ((e \in \mathbb{C}) \wedge ((e \cdot (1 - (\mathit{aiqv} \cdot v))) = v))))) \Rightarrow  \\ \quad \quad \operatorname{cc\text{-}series\text{-}partial\text{-}sum}((\lambda j \in \mathbb{N}.\; ((p \cdot \operatorname{power}(v, (j + 1))) \cdot \operatorname{power}(\mathit{aiqv}, j))), n) \\ \quad \quad =\; ((p \cdot e) - ((p \cdot \operatorname{power}((\mathit{aiqv} \cdot v), n)) \cdot e)) \end{array}}}
\vnbprop{aif-trace-bounded}{\vnbstmt{\begin{array}{@{}l@{}} \forall u.; \forall i \in (u \to \mathbb{C}), m, a, b. \\ \quad u \subseteq \mathbb{C} \wedge \\ \quad \forall y \in u. \\ \quad \quad \operatorname{is\text{-}continuous\text{-}at}(\operatorname{subspace\text{-}ms}(\operatorname{nf\text{-}metric\text{-}space}(\operatorname{cc\text{-}normed\text{-}field}), u), \\ \quad \quad \quad \operatorname{nf\text{-}metric\text{-}space}(\operatorname{cc\text{-}normed\text{-}field}), i, y) \wedge \\ \quad \operatorname{is\text{-}path}(m, a, b) \wedge \\ \quad \operatorname{trace}(m, a, b) \subseteq u \Rightarrow \\ \quad \quad \exists \mathit{aim} \in \mathbb{R}. \\ \quad \quad \quad (0 \leq \mathit{aim}) \wedge \forall \mathit{aiy} \in \operatorname{trace}(m, a, b).\; (\operatorname{magnitude}(i(\mathit{aiy})) \leq \mathit{aim}) \end{array}}}
\vnbprop{aif-cont-scale}{\vnbstmt{\begin{array}{@{}l@{}} \forall u.; \forall k \in \mathbb{C}, g, h \in (u \to \mathbb{C}), y \in u. \\ \quad u \subseteq \mathbb{C} \wedge \\ \quad \forall y \in u. \\ \quad \quad \operatorname{is\text{-}continuous\text{-}at}(\operatorname{subspace\text{-}ms}(\operatorname{nf\text{-}metric\text{-}space}(\operatorname{cc\text{-}normed\text{-}field}), u), \\ \quad \quad \quad \operatorname{nf\text{-}metric\text{-}space}(\operatorname{cc\text{-}normed\text{-}field}), g, y) \wedge \\ \quad \forall x \in u.\; (h(x) \simeq (k \cdot g(x))) \Rightarrow \\ \quad \quad \operatorname{is\text{-}continuous\text{-}at}(\operatorname{subspace\text{-}ms}(\operatorname{nf\text{-}metric\text{-}space}(\operatorname{cc\text{-}normed\text{-}field}), u), \\ \quad \quad \quad \operatorname{nf\text{-}metric\text{-}space}(\operatorname{cc\text{-}normed\text{-}field}), h, y) \end{array}}}
\vnbprop{line-int-of-uniform-limit}{\vnbstmt{\begin{array}{@{}l@{}} \forall u, m, c, a, b.; \forall g \in (u \to \mathbb{C}), d. \\ \quad u \subseteq \mathbb{C} \wedge \\ \quad \operatorname{is\text{-}road}(m, c, a, b) \wedge \\ \quad \operatorname{trace}(m, a, b) \subseteq u \wedge \\ \quad \forall y \in u. \\ \quad \quad \operatorname{is\text{-}continuous\text{-}at}(\operatorname{subspace\text{-}ms}(\operatorname{nf\text{-}metric\text{-}space}(\operatorname{cc\text{-}normed\text{-}field}), u), \\ \quad \quad \quad \operatorname{nf\text{-}metric\text{-}space}(\operatorname{cc\text{-}normed\text{-}field}), g, y) \wedge \\ \quad \forall k \in \mathbb{N}.\; (d(k) \in (u \to \mathbb{C})) \wedge \\ \quad \forall k \in \mathbb{N}, y \in u. \\ \quad \quad \operatorname{is\text{-}continuous\text{-}at}(\operatorname{subspace\text{-}ms}(\operatorname{nf\text{-}metric\text{-}space}(\operatorname{cc\text{-}normed\text{-}field}), u), \\ \quad \quad \quad \operatorname{nf\text{-}metric\text{-}space}(\operatorname{cc\text{-}normed\text{-}field}), d(k), y) \wedge \\ \quad \forall s. \\ \quad \quad \operatorname{pos\text{-}rr}(s) \Rightarrow  \\ \quad \quad \quad \exists p \in \mathbb{N}. \forall k \in \mathbb{N}, \mathit{aiy} \in \operatorname{trace}(m, a, b). \\ \quad \quad \quad \quad (p \leq k) \Rightarrow (\operatorname{magnitude}((d(k)(\mathit{aiy}) - g(\mathit{aiy}))) \leq s) \Rightarrow \\ \quad \quad \operatorname{converges\text{-}to}(\operatorname{cc\text{-}ms}, (\lambda k \in \mathbb{N}.\; \operatorname{line\text{-}int}(d(k), m, c, a, b)), \operatorname{line\text{-}int}(g, m, c, a, b)) \end{array}}}
\vnbprop{aif-geometric-relation}{\vnbstmt{\begin{array}{@{}l@{}} \forall w, c, z \in \mathbb{C}. \\ \quad (\neg (w = c) \wedge \neg (w = z)) \Rightarrow  \\ \quad \quad (\operatorname{recip}((w - z)) \cdot (1 - ((z - c) \cdot \operatorname{recip}((w - c))))) = \operatorname{recip}((w - c)) \end{array}}}
\vnbprop{aif-geometric-remainder}{\vnbstmt{\begin{array}{@{}l@{}} \forall c \in \mathbb{C}, r \in \mathbb{R}, w, z, p \in \mathbb{C}, m \in \mathbb{R}, k \in \mathbb{N}. \\ \quad ((0 < r) \wedge ((r \leq \operatorname{magnitude}((w - c))) \wedge ((\operatorname{magnitude}((z - c)) < r) \wedge (\operatorname{magnitude}(p) \leq m)))) \Rightarrow  \\ \quad \quad \operatorname{magnitude}((\operatorname{cc\text{-}series\text{-}partial\text{-}sum}((\lambda j \in \mathbb{N}.\; ((p \cdot \operatorname{power}(\operatorname{recip}((w - c)), (j + 1))) \cdot \operatorname{power}((z - c), j))), k) - (p \cdot \operatorname{recip}((w - z))))) \\ \quad \quad \in\; \mathbb{R} \wedge \\ \quad \quad \operatorname{magnitude}((\operatorname{cc\text{-}series\text{-}partial\text{-}sum}((\lambda j \in \mathbb{N}.\; ((p \cdot \operatorname{power}(\operatorname{recip}((w - c)), (j + 1))) \cdot \operatorname{power}((z - c), j))), k) - (p \cdot \operatorname{recip}((w - z))))) \\ \quad \quad \leq\; (m \cdot \operatorname{power}((\operatorname{magnitude}((z - c)) \cdot \operatorname{recip}(r)), k)) \cdot \operatorname{recip}((r - \operatorname{magnitude}((z - c)))) \end{array}}}
\vnbprop{geometric-expansion-uniform}{\vnbstmt{\begin{array}{@{}l@{}} \forall u, c, r.; \forall z \in \mathbb{C}, i.; \forall m \in \mathbb{R}, s. \\ \quad u \subseteq \mathbb{C} \wedge \\ \quad c \in \mathbb{C} \wedge \\ \quad r \in \mathbb{R} \wedge \\ \quad 0 < r \wedge \\ \quad \forall x \in u.\; (r \leq \operatorname{magnitude}((x - c))) \wedge \\ \quad \operatorname{magnitude}((z - c)) < r \wedge \\ \quad \forall x \in u.\; (i(x) \in \mathbb{C}) \wedge \\ \quad 0 \leq m \wedge \\ \quad \forall x \in u.\; (\operatorname{magnitude}(i(x)) \leq m) \wedge \\ \quad \operatorname{pos\text{-}rr}(s) \Rightarrow \\ \quad \quad \exists p \in \mathbb{N}. \forall k \in \mathbb{N}, y \in u. \\ \quad \quad \quad (p \leq k) \Rightarrow  \\ \quad \quad \quad \quad \operatorname{magnitude}((\operatorname{cc\text{-}series\text{-}partial\text{-}sum}((\lambda j \in \mathbb{N}.\; ((i(y) \cdot \operatorname{power}(\operatorname{recip}((y - c)), (j + 1))) \cdot \operatorname{power}((z - c), j))), k) - (i(y) \cdot \operatorname{recip}((y - z))))) \\ \quad \quad \quad \quad \leq\; s \end{array}}}
\vnbprop{aif-off-center}{\vnbstmt{\begin{array}{@{}l@{}} \forall c \in \mathbb{C}, r \in \mathbb{R}, w, z \in \mathbb{C}. \\ \quad ((0 < r) \wedge ((r \leq \operatorname{magnitude}((w - c))) \wedge (\operatorname{magnitude}((z - c)) < r))) \Rightarrow  \\ \quad \quad \neg (w = c) \wedge \neg (w = z) \end{array}}}
\vnbprop{aif-partial-integrals}{\vnbstmt{\begin{array}{@{}l@{}} \forall k \in \mathbb{N}, u, c, r, i, m, d, a, b, z. \\ \quad u \subseteq \mathbb{C} \wedge \\ \quad c \in \mathbb{C} \wedge \\ \quad r \in \mathbb{R} \wedge \\ \quad 0 < r \wedge \\ \quad \forall x \in u.\; (r \leq \operatorname{magnitude}((x - c))) \wedge \\ \quad i \in (u \to \mathbb{C}) \wedge \\ \quad \forall y \in u. \\ \quad \quad \operatorname{is\text{-}continuous\text{-}at}(\operatorname{subspace\text{-}ms}(\operatorname{nf\text{-}metric\text{-}space}(\operatorname{cc\text{-}normed\text{-}field}), u), \\ \quad \quad \quad \operatorname{nf\text{-}metric\text{-}space}(\operatorname{cc\text{-}normed\text{-}field}), i, y) \wedge \\ \quad \operatorname{is\text{-}road}(m, d, a, b) \wedge \\ \quad \operatorname{trace}(m, a, b) \subseteq u \wedge \\ \quad z \in \mathbb{C} \wedge \\ \quad \operatorname{magnitude}((z - c)) < r \Rightarrow \\ \quad \quad (\lambda w \in u.\; \\ \quad \quad \quad \operatorname{cc\text{-}series\text{-}partial\text{-}sum}((\lambda j \in \mathbb{N}.\; ((i(w) \cdot \operatorname{power}(\operatorname{recip}((w - c)), (j + 1))) \cdot \operatorname{power}((z - c), j))), k)) \\ \quad \quad \in\; (u \to \mathbb{C}) \wedge \\ \quad \quad \forall y \in u. \\ \quad \quad \quad \operatorname{is\text{-}continuous\text{-}at}(\operatorname{subspace\text{-}ms}(\operatorname{nf\text{-}metric\text{-}space}(\operatorname{cc\text{-}normed\text{-}field}), u), \\ \quad \quad \quad \quad \operatorname{nf\text{-}metric\text{-}space}(\operatorname{cc\text{-}normed\text{-}field}), \\ \quad \quad \quad \quad (\lambda w \in u.\; \\ \quad \quad \quad \quad \quad \operatorname{cc\text{-}series\text{-}partial\text{-}sum}((\lambda j \in \mathbb{N}.\; ((i(w) \cdot \operatorname{power}(\operatorname{recip}((w - c)), (j + 1))) \cdot \operatorname{power}((z - c), j))), k)), \\ \quad \quad \quad \quad y) \wedge \\ \quad \quad \operatorname{line\text{-}int}((\lambda w \in u.\; \\ \quad \quad \quad \quad \operatorname{cc\text{-}series\text{-}partial\text{-}sum}((\lambda j \in \mathbb{N}.\; ((i(w) \cdot \operatorname{power}(\operatorname{recip}((w - c)), (j + 1))) \cdot \operatorname{power}((z - c), j))), k)), \\ \quad \quad \quad m, d, a, b) \\ \quad \quad =\; \operatorname{cc\text{-}series\text{-}partial\text{-}sum}((\lambda \mathit{aitk} \in \mathbb{N}.\; \\ \quad \quad \quad \quad \operatorname{line\text{-}int}((\lambda w \in u.\; (i(w) \cdot \operatorname{power}(\operatorname{recip}((w - c)), (\mathit{aitk} + 1)))), m, d, a, b) \cdot \operatorname{power}((z - c), \mathit{aitk})), \\ \quad \quad \quad k) \end{array}}}
\vnbprop{cauchy-kernel-integral-series}{\vnbstmt{\begin{array}{@{}l@{}} \forall u, c, r, i, m, d, a, b.; \forall z \in \mathbb{C}. \\ \quad u \subseteq \mathbb{C} \wedge \\ \quad c \in \mathbb{C} \wedge \\ \quad r \in \mathbb{R} \wedge \\ \quad 0 < r \wedge \\ \quad \forall x \in u.\; (r \leq \operatorname{magnitude}((x - c))) \wedge \\ \quad i \in (u \to \mathbb{C}) \wedge \\ \quad \forall y \in u. \\ \quad \quad \operatorname{is\text{-}continuous\text{-}at}(\operatorname{subspace\text{-}ms}(\operatorname{nf\text{-}metric\text{-}space}(\operatorname{cc\text{-}normed\text{-}field}), u), \\ \quad \quad \quad \operatorname{nf\text{-}metric\text{-}space}(\operatorname{cc\text{-}normed\text{-}field}), i, y) \wedge \\ \quad \operatorname{is\text{-}road}(m, d, a, b) \wedge \\ \quad \operatorname{trace}(m, a, b) \subseteq u \wedge \\ \quad \operatorname{magnitude}((z - c)) < r \Rightarrow \\ \quad \quad \operatorname{cc\text{-}series\text{-}converges\text{-}to}((\lambda k \in \mathbb{N}.\; \\ \quad \quad \quad \quad \operatorname{line\text{-}int}((\lambda w \in u.\; (i(w) \cdot \operatorname{power}(\operatorname{recip}((w - c)), (k + 1)))), m, d, a, b) \cdot \operatorname{power}((z - c), k)), \\ \quad \quad \quad \operatorname{line\text{-}int}((\lambda w \in u.\; (i(w) \cdot \operatorname{recip}((w - z)))), m, d, a, b)) \end{array}}}
\vnbprop{analytic-from-integral-formula}{\vnbstmt{\begin{array}{@{}l@{}} \forall u, c, r, i, m, d, a, b.; \forall p \in \mathbb{C}, f.; \forall z \in \mathbb{C}. \\ \quad u \subseteq \mathbb{C} \wedge \\ \quad c \in \mathbb{C} \wedge \\ \quad r \in \mathbb{R} \wedge \\ \quad 0 < r \wedge \\ \quad \forall x \in u.\; (r \leq \operatorname{magnitude}((x - c))) \wedge \\ \quad i \in (u \to \mathbb{C}) \wedge \\ \quad \forall y \in u. \\ \quad \quad \operatorname{is\text{-}continuous\text{-}at}(\operatorname{subspace\text{-}ms}(\operatorname{nf\text{-}metric\text{-}space}(\operatorname{cc\text{-}normed\text{-}field}), u), \\ \quad \quad \quad \operatorname{nf\text{-}metric\text{-}space}(\operatorname{cc\text{-}normed\text{-}field}), i, y) \wedge \\ \quad \operatorname{is\text{-}road}(m, d, a, b) \wedge \\ \quad \operatorname{trace}(m, a, b) \subseteq u \wedge \\ \quad \forall z \in \mathbb{C}. \\ \quad \quad (\operatorname{magnitude}((z - c)) < r) \Rightarrow  \\ \quad \quad \quad f(z) = (p \cdot \operatorname{line\text{-}int}((\lambda w \in u.\; (i(w) \cdot \operatorname{recip}((w - z)))), m, d, a, b)) \wedge \\ \quad \operatorname{magnitude}((z - c)) < r \Rightarrow \\ \quad \quad \operatorname{cc\text{-}series\text{-}converges\text{-}to}((\lambda k \in \mathbb{N}.\; \\ \quad \quad \quad \quad (p \cdot \operatorname{line\text{-}int}((\lambda w \in u.\; (i(w) \cdot \operatorname{power}(\operatorname{recip}((w - c)), (k + 1)))), m, d, a, b)) \cdot \operatorname{power}((z - c), k)), \\ \quad \quad \quad f(z)) \end{array}}}
\vnbprop{aif-coefficient-bound}{\vnbstmt{\begin{array}{@{}l@{}} \forall u, c, r, i, m, d, a, b.; \forall \mathit{aim} \in \mathbb{R}, k \in \mathbb{N}. \\ \quad u \subseteq \mathbb{C} \wedge \\ \quad c \in \mathbb{C} \wedge \\ \quad r \in \mathbb{R} \wedge \\ \quad 0 < r \wedge \\ \quad \forall x \in u.\; (r \leq \operatorname{magnitude}((x - c))) \wedge \\ \quad i \in (u \to \mathbb{C}) \wedge \\ \quad \forall y \in u. \\ \quad \quad \operatorname{is\text{-}continuous\text{-}at}(\operatorname{subspace\text{-}ms}(\operatorname{nf\text{-}metric\text{-}space}(\operatorname{cc\text{-}normed\text{-}field}), u), \\ \quad \quad \quad \operatorname{nf\text{-}metric\text{-}space}(\operatorname{cc\text{-}normed\text{-}field}), i, y) \wedge \\ \quad \operatorname{is\text{-}road}(m, d, a, b) \wedge \\ \quad \operatorname{trace}(m, a, b) \subseteq u \wedge \\ \quad \forall \mathit{aiy} \in \operatorname{trace}(m, a, b).\; (\operatorname{magnitude}(i(\mathit{aiy})) \leq \mathit{aim}) \Rightarrow \\ \quad \quad \operatorname{magnitude}(\operatorname{line\text{-}int}((\lambda w \in u.\; (i(w) \cdot \operatorname{power}(\operatorname{recip}((w - c)), (k + 1)))), m, d, a, b)) \\ \quad \quad \leq\; \operatorname{road\text{-}length}(d, a, b) \cdot (\mathit{aim} \cdot \operatorname{power}(\operatorname{recip}(r), (k + 1))) \end{array}}}
\vnbprop{cauchy-estimate-from-formula}{\vnbstmt{\begin{array}{@{}l@{}} \forall u, c, r, i, m, d, a, b.; \forall \mathit{aim} \in \mathbb{R}, p \in \mathbb{C}, k \in \mathbb{N}. \\ \quad u \subseteq \mathbb{C} \wedge \\ \quad c \in \mathbb{C} \wedge \\ \quad r \in \mathbb{R} \wedge \\ \quad 0 < r \wedge \\ \quad \forall x \in u.\; (r \leq \operatorname{magnitude}((x - c))) \wedge \\ \quad i \in (u \to \mathbb{C}) \wedge \\ \quad \forall y \in u. \\ \quad \quad \operatorname{is\text{-}continuous\text{-}at}(\operatorname{subspace\text{-}ms}(\operatorname{nf\text{-}metric\text{-}space}(\operatorname{cc\text{-}normed\text{-}field}), u), \\ \quad \quad \quad \operatorname{nf\text{-}metric\text{-}space}(\operatorname{cc\text{-}normed\text{-}field}), i, y) \wedge \\ \quad \operatorname{is\text{-}road}(m, d, a, b) \wedge \\ \quad \operatorname{trace}(m, a, b) \subseteq u \wedge \\ \quad \forall \mathit{aiy} \in \operatorname{trace}(m, a, b).\; (\operatorname{magnitude}(i(\mathit{aiy})) \leq \mathit{aim}) \Rightarrow \\ \quad \quad \operatorname{magnitude}((p \cdot \operatorname{line\text{-}int}((\lambda w \in u.\; (i(w) \cdot \operatorname{power}(\operatorname{recip}((w - c)), (k + 1)))), m, d, a, b))) \\ \quad \quad \leq\; \operatorname{magnitude}(p) \cdot (\operatorname{road\text{-}length}(d, a, b) \cdot (\mathit{aim} \cdot \operatorname{power}(\operatorname{recip}(r), (k + 1)))) \end{array}}}
\vnbprop{cauchy-kernel-series-absolute}{\vnbstmt{\begin{array}{@{}l@{}} \forall u, c, r, i, m, d, a, b.; \forall z \in \mathbb{C}. \\ \quad u \subseteq \mathbb{C} \wedge \\ \quad c \in \mathbb{C} \wedge \\ \quad r \in \mathbb{R} \wedge \\ \quad 0 < r \wedge \\ \quad \forall x \in u.\; (r \leq \operatorname{magnitude}((x - c))) \wedge \\ \quad i \in (u \to \mathbb{C}) \wedge \\ \quad \forall y \in u. \\ \quad \quad \operatorname{is\text{-}continuous\text{-}at}(\operatorname{subspace\text{-}ms}(\operatorname{nf\text{-}metric\text{-}space}(\operatorname{cc\text{-}normed\text{-}field}), u), \\ \quad \quad \quad \operatorname{nf\text{-}metric\text{-}space}(\operatorname{cc\text{-}normed\text{-}field}), i, y) \wedge \\ \quad \operatorname{is\text{-}road}(m, d, a, b) \wedge \\ \quad \operatorname{trace}(m, a, b) \subseteq u \wedge \\ \quad \operatorname{magnitude}((z - c)) < r \Rightarrow \\ \quad \quad \operatorname{series\text{-}converges}((\lambda k \in \mathbb{N}.\; \\ \quad \quad \quad \quad \operatorname{magnitude}((\operatorname{line\text{-}int}((\lambda w \in u.\; (i(w) \cdot \operatorname{power}(\operatorname{recip}((w - c)), (k + 1)))), m, d, a, b) \cdot \operatorname{power}((z - c), k))))) \end{array}}}
\vnbprop{analytic-at-from-integral-formula}{\vnbstmt{\begin{array}{@{}l@{}} \forall u, c, r, i, m, d, a, b.; \forall p \in \mathbb{C}, \mathit{aiou}.; \forall f \in (\mathit{aiou} \to \mathbb{C}). \\ \quad u \subseteq \mathbb{C} \wedge \\ \quad c \in \mathbb{C} \wedge \\ \quad r \in \mathbb{R} \wedge \\ \quad 0 < r \wedge \\ \quad \forall x \in u.\; (r \leq \operatorname{magnitude}((x - c))) \wedge \\ \quad i \in (u \to \mathbb{C}) \wedge \\ \quad \forall y \in u. \\ \quad \quad \operatorname{is\text{-}continuous\text{-}at}(\operatorname{subspace\text{-}ms}(\operatorname{nf\text{-}metric\text{-}space}(\operatorname{cc\text{-}normed\text{-}field}), u), \\ \quad \quad \quad \operatorname{nf\text{-}metric\text{-}space}(\operatorname{cc\text{-}normed\text{-}field}), i, y) \wedge \\ \quad \operatorname{is\text{-}road}(m, d, a, b) \wedge \\ \quad \operatorname{trace}(m, a, b) \subseteq u \wedge \\ \quad \operatorname{is\text{-}open}(\operatorname{cc\text{-}ms}, \mathit{aiou}) \wedge \\ \quad \operatorname{ball}(\operatorname{cc\text{-}ms}, c, r) \subseteq \mathit{aiou} \wedge \\ \quad \forall z \in \mathbb{C}. \\ \quad \quad (\operatorname{magnitude}((z - c)) < r) \Rightarrow  \\ \quad \quad \quad f(z) = (p \cdot \operatorname{line\text{-}int}((\lambda w \in u.\; (i(w) \cdot \operatorname{recip}((w - z)))), m, d, a, b)) \Rightarrow \\ \quad \quad f \text{ is analytic at } c \end{array}}}
\vnbprop{cauchy-type-integral-deriv}{\vnbstmt{\begin{array}{@{}l@{}} \forall u, c, r, i, m, d, a, b.; \forall p \in \mathbb{C}, \mathit{aiou}.; \forall f \in (\mathit{aiou} \to \mathbb{C}). \\ \quad u \subseteq \mathbb{C} \wedge \\ \quad c \in \mathbb{C} \wedge \\ \quad r \in \mathbb{R} \wedge \\ \quad 0 < r \wedge \\ \quad \forall x \in u.\; (r \leq \operatorname{magnitude}((x - c))) \wedge \\ \quad i \in (u \to \mathbb{C}) \wedge \\ \quad \forall y \in u. \\ \quad \quad \operatorname{is\text{-}continuous\text{-}at}(\operatorname{subspace\text{-}ms}(\operatorname{nf\text{-}metric\text{-}space}(\operatorname{cc\text{-}normed\text{-}field}), u), \\ \quad \quad \quad \operatorname{nf\text{-}metric\text{-}space}(\operatorname{cc\text{-}normed\text{-}field}), i, y) \wedge \\ \quad \operatorname{is\text{-}road}(m, d, a, b) \wedge \\ \quad \operatorname{trace}(m, a, b) \subseteq u \wedge \\ \quad \operatorname{is\text{-}open}(\operatorname{cc\text{-}ms}, \mathit{aiou}) \wedge \\ \quad \operatorname{ball}(\operatorname{cc\text{-}ms}, c, r) \subseteq \mathit{aiou} \wedge \\ \quad \forall z \in \mathbb{C}. \\ \quad \quad (\operatorname{magnitude}((z - c)) < r) \Rightarrow  \\ \quad \quad \quad f(z) = (p \cdot \operatorname{line\text{-}int}((\lambda w \in u.\; (i(w) \cdot \operatorname{recip}((w - z)))), m, d, a, b)) \Rightarrow \\ \quad \quad \operatorname{is\text{-}diff\text{-}on}(\operatorname{cc\text{-}normed\text{-}field}, \mathit{aiou}, f, c, \\ \quad \quad \quad p \cdot \operatorname{line\text{-}int}((\lambda w \in u.\; (i(w) \cdot \operatorname{power}(\operatorname{recip}((w - c)), 2))), m, d, a, b)) \end{array}}}

\vnbbookfile{theorem-library/ch4a-termwise}{theorem-library/ch4a-termwise}
\noindent{\small ch4a-termwise.scm -- WAVE 2 (2026-10-27..30), wired 2026-10-30 from scratchpad/ch4a/ch4a-termwise.scm. Library: complex-analysis.}\par
\vnbprop{termwise-partial-integrals}{\vnbstmt{\begin{array}{@{}l@{}} \forall k \in \mathbb{N}, u, m, c, a, b, f. \\ \quad u \subseteq \mathbb{C} \wedge \\ \quad \operatorname{is\text{-}road}(m, c, a, b) \wedge \\ \quad \operatorname{trace}(m, a, b) \subseteq u \wedge \\ \quad \forall j \in \mathbb{N}.\; (f(j) \in (u \to \mathbb{C})) \wedge \\ \quad \forall j \in \mathbb{N}, y \in u. \\ \quad \quad \operatorname{is\text{-}continuous\text{-}at}(\operatorname{subspace\text{-}ms}(\operatorname{nf\text{-}metric\text{-}space}(\operatorname{cc\text{-}normed\text{-}field}), u), \\ \quad \quad \quad \operatorname{nf\text{-}metric\text{-}space}(\operatorname{cc\text{-}normed\text{-}field}), f(j), y) \Rightarrow \\ \quad \quad (\lambda w \in u.\; \operatorname{cc\text{-}series\text{-}partial\text{-}sum}((\lambda j \in \mathbb{N}.\; f(j)(w)), k)) \in (u \to \mathbb{C}) \wedge \\ \quad \quad \forall y \in u. \\ \quad \quad \quad \operatorname{is\text{-}continuous\text{-}at}(\operatorname{subspace\text{-}ms}(\operatorname{nf\text{-}metric\text{-}space}(\operatorname{cc\text{-}normed\text{-}field}), u), \\ \quad \quad \quad \quad \operatorname{nf\text{-}metric\text{-}space}(\operatorname{cc\text{-}normed\text{-}field}), \\ \quad \quad \quad \quad (\lambda w \in u.\; \operatorname{cc\text{-}series\text{-}partial\text{-}sum}((\lambda j \in \mathbb{N}.\; f(j)(w)), k)), \\ \quad \quad \quad \quad y) \wedge \\ \quad \quad \operatorname{line\text{-}int}((\lambda w \in u.\; \operatorname{cc\text{-}series\text{-}partial\text{-}sum}((\lambda j \in \mathbb{N}.\; f(j)(w)), k)), m, c, a, b) \\ \quad \quad =\; \operatorname{cc\text{-}series\text{-}partial\text{-}sum}((\lambda j \in \mathbb{N}.\; \operatorname{line\text{-}int}(f(j), m, c, a, b)), k) \end{array}}}
\vnbprop{line-int-termwise}{\vnbstmt{\begin{array}{@{}l@{}} \forall u, m, c, a, b, f, g. \\ \quad u \subseteq \mathbb{C} \wedge \\ \quad \operatorname{is\text{-}road}(m, c, a, b) \wedge \\ \quad \operatorname{trace}(m, a, b) \subseteq u \wedge \\ \quad \forall j \in \mathbb{N}.\; (f(j) \in (u \to \mathbb{C})) \wedge \\ \quad \forall j \in \mathbb{N}, y \in u. \\ \quad \quad \operatorname{is\text{-}continuous\text{-}at}(\operatorname{subspace\text{-}ms}(\operatorname{nf\text{-}metric\text{-}space}(\operatorname{cc\text{-}normed\text{-}field}), u), \\ \quad \quad \quad \operatorname{nf\text{-}metric\text{-}space}(\operatorname{cc\text{-}normed\text{-}field}), f(j), y) \wedge \\ \quad g \in (u \to \mathbb{C}) \wedge \\ \quad \forall y \in u. \\ \quad \quad \operatorname{is\text{-}continuous\text{-}at}(\operatorname{subspace\text{-}ms}(\operatorname{nf\text{-}metric\text{-}space}(\operatorname{cc\text{-}normed\text{-}field}), u), \\ \quad \quad \quad \operatorname{nf\text{-}metric\text{-}space}(\operatorname{cc\text{-}normed\text{-}field}), g, y) \wedge \\ \quad \forall s. \\ \quad \quad \operatorname{pos\text{-}rr}(s) \Rightarrow  \\ \quad \quad \quad \exists p \in \mathbb{N}. \forall k \in \mathbb{N}, \mathit{twy} \in \operatorname{trace}(m, a, b). \\ \quad \quad \quad \quad (p \leq k) \Rightarrow  \\ \quad \quad \quad \quad \quad \operatorname{magnitude}((\operatorname{cc\text{-}series\text{-}partial\text{-}sum}((\lambda j \in \mathbb{N}.\; f(j)(\mathit{twy})), k) - g(\mathit{twy}))) \leq s \Rightarrow \\ \quad \quad \operatorname{cc\text{-}series\text{-}converges\text{-}to}((\lambda j \in \mathbb{N}.\; \operatorname{line\text{-}int}(f(j), m, c, a, b)), \operatorname{line\text{-}int}(g, m, c, a, b)) \end{array}}}
\vnbprop{cc-series-limit-of-converges-to}{\vnbstmt{\begin{array}{@{}l@{}} \forall s \in (\mathbb{N} \to \mathbb{C}), x. \\ \quad \operatorname{cc\text{-}series\text{-}converges\text{-}to}(s, x) \Rightarrow  \\ \quad \quad \operatorname{cc\text{-}series\text{-}converges}(s) \wedge (\operatorname{cc\text{-}series\text{-}limit}(s) = x) \end{array}}}

\vnbbookfile{theorem-library/ch4a-outer}{theorem-library/ch4a-outer}
\noindent{\small ch4a-outer.scm -- WAVE 2 (2026-10-27..30), wired 2026-10-30 from scratchpad/ch4a/ch4a-outer.scm. Library: complex-analysis.}\par
\vnbprop{out-cont-recip-rev}{\vnbstmt{\begin{array}{@{}l@{}} \forall u.; \forall z \in \mathbb{C}, h \in (u \to \mathbb{C}), y \in u. \\ \quad ((u \subseteq \mathbb{C}) \wedge (\forall x \in u.\; \neg (x = z) \wedge \forall x \in u.\; (h(x) \simeq \operatorname{recip}((z - x))))) \Rightarrow  \\ \quad \quad \operatorname{is\text{-}continuous\text{-}at}(\operatorname{subspace\text{-}ms}(\operatorname{nf\text{-}metric\text{-}space}(\operatorname{cc\text{-}normed\text{-}field}), u), \\ \quad \quad \quad \operatorname{nf\text{-}metric\text{-}space}(\operatorname{cc\text{-}normed\text{-}field}), h, y) \end{array}}}
\vnbprop{outer-kernel-integral-series}{\vnbstmt{\begin{array}{@{}l@{}} \forall u, c, r, i, m, d, a, b.; \forall z \in \mathbb{C}. \\ \quad u \subseteq \mathbb{C} \wedge \\ \quad c \in \mathbb{C} \wedge \\ \quad r \in \mathbb{R} \wedge \\ \quad 0 \leq r \wedge \\ \quad \forall x \in u.\; (\operatorname{magnitude}((x - c)) \leq r) \wedge \\ \quad i \in (u \to \mathbb{C}) \wedge \\ \quad \forall y \in u. \\ \quad \quad \operatorname{is\text{-}continuous\text{-}at}(\operatorname{subspace\text{-}ms}(\operatorname{nf\text{-}metric\text{-}space}(\operatorname{cc\text{-}normed\text{-}field}), u), \\ \quad \quad \quad \operatorname{nf\text{-}metric\text{-}space}(\operatorname{cc\text{-}normed\text{-}field}), i, y) \wedge \\ \quad \operatorname{is\text{-}road}(m, d, a, b) \wedge \\ \quad \operatorname{trace}(m, a, b) \subseteq u \wedge \\ \quad r < \operatorname{magnitude}((z - c)) \Rightarrow \\ \quad \quad \operatorname{cc\text{-}series\text{-}converges\text{-}to}((\lambda k \in \mathbb{N}.\; \\ \quad \quad \quad \quad \operatorname{line\text{-}int}((\lambda w \in u.\; (i(w) \cdot \operatorname{power}((w - c), k))), m, d, a, b) \cdot \operatorname{power}(\operatorname{recip}((z - c)), (k + 1))), \\ \quad \quad \quad \operatorname{line\text{-}int}((\lambda w \in u.\; (i(w) \cdot \operatorname{recip}((z - w)))), m, d, a, b)) \end{array}}}

\vnbbookfile{theorem-library/cauchy-integral-formula}{theorem-library/cauchy-integral-formula}
\noindent{\small cauchy-integral-formula.scm -- CAUCHY WITH AN EXCEPTIONAL POINT ON A DISC, THE WINDING NUMBER OF A CIRCLE INSIDE IT, THE CAUCHY INTEGRAL FORMULA ON A DISC, THEOREM 3.35 (HOLOMORPHIC =$>$ ANALYTIC) AND THE HOLOMORPHY OF THE DERIVATIVE. (complex-analysis.pdf ch. 3: Cor 3.9 with one exceptional point, Example 3.13 /}\par
\vnbprop{cif-open-minus-point}{\vnbstmt{\begin{array}{@{}l@{}} \forall v.; \forall x \in \mathbb{C}. \\ \quad \operatorname{is\text{-}open}(\operatorname{nf\text{-}metric\text{-}space}(\operatorname{cc\text{-}normed\text{-}field}), v) \Rightarrow  \\ \quad \quad \operatorname{is\text{-}open}(\operatorname{nf\text{-}metric\text{-}space}(\operatorname{cc\text{-}normed\text{-}field}), \operatorname{difference}(v, \operatorname{singleton}(x))) \end{array}}}
\vnbprop{cauchy-convex-sub-except}{\vnbstmt{\begin{array}{@{}l@{}} \forall u, f, x, v, m, c, a, b. \\ \quad \operatorname{is\text{-}open}(\operatorname{nf\text{-}metric\text{-}space}(\operatorname{cc\text{-}normed\text{-}field}), u) \wedge \\ \quad f \in (u \to \mathbb{C}) \wedge \\ \quad \forall y \in u. \\ \quad \quad \operatorname{is\text{-}continuous\text{-}at}(\operatorname{subspace\text{-}ms}(\operatorname{nf\text{-}metric\text{-}space}(\operatorname{cc\text{-}normed\text{-}field}), u), \\ \quad \quad \quad \operatorname{nf\text{-}metric\text{-}space}(\operatorname{cc\text{-}normed\text{-}field}), f, y) \wedge \\ \quad x \in u \wedge \\ \quad \operatorname{holomorphic\text{-}on}(\operatorname{difference}(u, \operatorname{singleton}(x)), \operatorname{restrict}(f, \operatorname{difference}(u, \operatorname{singleton}(x)))) \wedge \\ \quad \operatorname{is\text{-}open}(\operatorname{nf\text{-}metric\text{-}space}(\operatorname{cc\text{-}normed\text{-}field}), v) \wedge \\ \quad \operatorname{is\text{-}convex\text{-}cc}(v) \wedge \\ \quad v \subseteq u \wedge \\ \quad \operatorname{is\text{-}road}(m, c, a, b) \wedge \\ \quad \operatorname{trace}(m, a, b) \subseteq v \wedge \\ \quad m(b) = m(a) \Rightarrow \\ \quad \quad \operatorname{line\text{-}int}(f, m, c, a, b) = 0 \end{array}}}
\vnbprop{cauchy-disc-except}{\vnbstmt{\begin{array}{@{}l@{}} \forall u, f, x.; \forall c \in \mathbb{C}, r \in \mathbb{R}, m, d, a, b. \\ \quad \operatorname{is\text{-}open}(\operatorname{nf\text{-}metric\text{-}space}(\operatorname{cc\text{-}normed\text{-}field}), u) \wedge \\ \quad f \in (u \to \mathbb{C}) \wedge \\ \quad \forall y \in u. \\ \quad \quad \operatorname{is\text{-}continuous\text{-}at}(\operatorname{subspace\text{-}ms}(\operatorname{nf\text{-}metric\text{-}space}(\operatorname{cc\text{-}normed\text{-}field}), u), \\ \quad \quad \quad \operatorname{nf\text{-}metric\text{-}space}(\operatorname{cc\text{-}normed\text{-}field}), f, y) \wedge \\ \quad x \in u \wedge \\ \quad \operatorname{holomorphic\text{-}on}(\operatorname{difference}(u, \operatorname{singleton}(x)), \operatorname{restrict}(f, \operatorname{difference}(u, \operatorname{singleton}(x)))) \wedge \\ \quad 0 < r \wedge \\ \quad \operatorname{ball}(\operatorname{cc\text{-}ms}, c, r) \subseteq u \wedge \\ \quad \operatorname{is\text{-}road}(m, d, a, b) \wedge \\ \quad \operatorname{trace}(m, a, b) \subseteq \operatorname{ball}(\operatorname{cc\text{-}ms}, c, r) \wedge \\ \quad m(b) = m(a) \Rightarrow \\ \quad \quad \operatorname{line\text{-}int}(f, m, d, a, b) = 0 \end{array}}}
\vnbprop{cif-recip-shift-diff}{\vnbstmt{\begin{array}{@{}l@{}} \forall c \in \mathbb{C}, w \in \operatorname{difference}(\mathbb{C}, \operatorname{singleton}(c)). \\ \quad \operatorname{is\text{-}diff\text{-}on}(\operatorname{cc\text{-}normed\text{-}field}, \operatorname{difference}(\mathbb{C}, \operatorname{singleton}(c)), \\ \quad \quad (\lambda v \in \operatorname{difference}(\mathbb{C}, \operatorname{singleton}(c)).\; \operatorname{recip}((v - c))), \\ \quad \quad w, -(\operatorname{recip}((w - c)) \cdot \operatorname{recip}((w - c)))) \end{array}}}
\vnbprop{cif-recip-power-diff}{\vnbstmt{\begin{array}{@{}l@{}} \forall j \in \mathbb{N}, c \in \mathbb{C}, w \in \operatorname{difference}(\mathbb{C}, \operatorname{singleton}(c)). \\ \quad \operatorname{is\text{-}diff\text{-}on}(\operatorname{cc\text{-}normed\text{-}field}, \operatorname{difference}(\mathbb{C}, \operatorname{singleton}(c)), \\ \quad \quad (\lambda v \in \operatorname{difference}(\mathbb{C}, \operatorname{singleton}(c)).\; \operatorname{power}(\operatorname{recip}((v - c)), (j + 1))), \\ \quad \quad w, (-(j + 1) \cdot \operatorname{power}(\operatorname{recip}((w - c)), ((j + 1) + 1)))) \end{array}}}
\vnbprop{cif-recip-power-closed-zero}{\vnbstmt{\begin{array}{@{}l@{}} \forall j \in \mathbb{N}, c \in \mathbb{C}, m, d, a, b. \\ \quad (\operatorname{is\text{-}road}(m, d, a, b) \wedge ((\operatorname{trace}(m, a, b) \subseteq \operatorname{difference}(\mathbb{C}, \operatorname{singleton}(c))) \wedge (m(b) = m(a)))) \Rightarrow  \\ \quad \quad \operatorname{line\text{-}int}((\lambda v \in \operatorname{difference}(\mathbb{C}, \operatorname{singleton}(c)).\; \operatorname{power}(\operatorname{recip}((v - c)), ((j + 1) + 1))), m, d, a, b) = 0 \end{array}}}
\vnbprop{circle-trace-magnitude}{\vnbstmt{\begin{array}{@{}l@{}} \forall c \in \mathbb{C}, r \in \mathbb{R}, y \in \operatorname{trace}(\operatorname{circle\text{-}path}(c, r), 0, (2 \cdot \pi)). \\ \quad (0 \leq r) \Rightarrow (\operatorname{magnitude}((y - c)) = r) \end{array}}}
\vnbprop{circle-path-closed}{\vnbstmt{\forall c \in \mathbb{C}, r \in \mathbb{R}.\; (\operatorname{circle\text{-}path}(c, r)((2 \cdot \pi)) = \operatorname{circle\text{-}path}(c, r)(0))}}
\vnbprop{winding-circle-inside}{\vnbstmt{\begin{array}{@{}l@{}} \forall c \in \mathbb{C}, r \in \mathbb{R}, z \in \mathbb{C}. \\ \quad ((0 < r) \wedge (\operatorname{magnitude}((z - c)) < r)) \Rightarrow  \\ \quad \quad \operatorname{winding}(\operatorname{circle\text{-}path}(c, r), \operatorname{circle\text{-}deriv}(c, r), 0, (2 \cdot \pi), z) = 1 \end{array}}}
\vnbprop{cauchy-integral-formula-disc}{\vnbstmt{\begin{array}{@{}l@{}} \forall u, f.; \forall c \in \mathbb{C}, s, r \in \mathbb{R}, z \in \mathbb{C}. \\ \quad \operatorname{holomorphic\text{-}on}(u, f) \wedge \\ \quad 0 < s \wedge \\ \quad \operatorname{ball}(\operatorname{cc\text{-}ms}, c, s) \subseteq u \wedge \\ \quad 0 < r \wedge \\ \quad r < s \wedge \\ \quad \operatorname{magnitude}((z - c)) < r \Rightarrow \\ \quad \quad f(z) \\ \quad \quad =\; \operatorname{recip}(((2 \cdot \pi) \cdot +i)) \cdot \operatorname{line\text{-}int}((\lambda w \in \operatorname{trace}(\operatorname{circle\text{-}path}(c, r), 0, (2 \cdot \pi)).\; (f(w) \cdot \operatorname{recip}((w - z)))), \operatorname{circle\text{-}path}(c, r), \operatorname{circle\text{-}deriv}(c, r), 0, (2 \cdot \pi)) \end{array}}}
\vnbprop{cauchy-integral-formula-disc-restrict}{\vnbstmt{\begin{array}{@{}l@{}} \forall u, f.; \forall c \in \mathbb{C}, s, r \in \mathbb{R}, z \in \mathbb{C}. \\ \quad \operatorname{holomorphic\text{-}on}(u, f) \wedge \\ \quad 0 < s \wedge \\ \quad \operatorname{ball}(\operatorname{cc\text{-}ms}, c, s) \subseteq u \wedge \\ \quad 0 < r \wedge \\ \quad r < s \wedge \\ \quad \operatorname{magnitude}((z - c)) < r \Rightarrow \\ \quad \quad f(z) \\ \quad \quad =\; \operatorname{recip}(((2 \cdot \pi) \cdot +i)) \cdot \operatorname{line\text{-}int}((\lambda w \in \operatorname{trace}(\operatorname{circle\text{-}path}(c, r), 0, (2 \cdot \pi)).\; (\operatorname{restrict}(f, \operatorname{trace}(\operatorname{circle\text{-}path}(c, r), 0, (2 \cdot \pi)))(w) \cdot \operatorname{recip}((w - z)))), \operatorname{circle\text{-}path}(c, r), \operatorname{circle\text{-}deriv}(c, r), 0, (2 \cdot \pi)) \end{array}}}
\vnbprop{holomorphic-power-series-disc}{\vnbstmt{\begin{array}{@{}l@{}} \forall u, f.; \forall c \in \mathbb{C}, s, r \in \mathbb{R}, z \in \mathbb{C}. \\ \quad \operatorname{holomorphic\text{-}on}(u, f) \wedge \\ \quad 0 < s \wedge \\ \quad \operatorname{ball}(\operatorname{cc\text{-}ms}, c, s) \subseteq u \wedge \\ \quad 0 < r \wedge \\ \quad r < s \wedge \\ \quad \operatorname{magnitude}((z - c)) < r \Rightarrow \\ \quad \quad \operatorname{cc\text{-}series\text{-}converges\text{-}to}((\lambda k \in \mathbb{N}.\; \\ \quad \quad \quad \quad (\operatorname{recip}(((2 \cdot \pi) \cdot +i)) \cdot \operatorname{line\text{-}int}((\lambda w \in \operatorname{trace}(\operatorname{circle\text{-}path}(c, r), 0, (2 \cdot \pi)).\; (\operatorname{restrict}(f, \operatorname{trace}(\operatorname{circle\text{-}path}(c, r), 0, (2 \cdot \pi)))(w) \cdot \operatorname{power}(\operatorname{recip}((w - c)), (k + 1)))), \operatorname{circle\text{-}path}(c, r), \operatorname{circle\text{-}deriv}(c, r), 0, (2 \cdot \pi))) \cdot \operatorname{power}((z - c), k)), \\ \quad \quad \quad f(z)) \end{array}}}
\vnbprop{holomorphic-analytic-at}{\vnbstmt{\forall u, f.; \forall c \in u.\; (\operatorname{holomorphic\text{-}on}(u, f) \Rightarrow f \text{ is analytic at } c)}}
\vnbprop{holomorphic-deriv-holomorphic}{\vnbstmt{\begin{array}{@{}l@{}} \forall u, f. \\ \quad \operatorname{holomorphic\text{-}on}(u, f) \Rightarrow  \\ \quad \quad \operatorname{holomorphic\text{-}on}(u, (\lambda w \in u.\; \operatorname{deriv\text{-}on}(\operatorname{cc\text{-}normed\text{-}field}, u, f, w))) \end{array}}}

\vnbbookfile{theorem-library/cauchy-estimates}{theorem-library/cauchy-estimates}
\noindent{\small cauchy-estimates.scm -- THE CIRCLE'S LENGTH, CAUCHY'S ESTIMATES, THE DERIVATIVE FORMULA, LIOUVILLE AND THE FUNDAMENTAL THEOREM OF ALGEBRA. (complex-analysis.pdf ch. 3: Prop 3.40 / Cor 3.42 through the coefficients, Prop 3.47 and Cor 3.48 (Cauchy's estimates, Liouville); ch. 1, 1.5 (the fundamental theorem of algebra).)}\par
\vnbprop{circle-road-length}{\vnbstmt{\begin{array}{@{}l@{}} \forall c \in \mathbb{C}, r \in \mathbb{R}. \\ \quad (0 \leq r) \Rightarrow (\operatorname{road\text{-}length}(\operatorname{circle\text{-}deriv}(c, r), 0, (2 \cdot \pi)) = ((2 \cdot \pi) \cdot r)) \end{array}}}
\vnbprop{cauchy-kernel-constant-magnitude}{\vnbstmt{\operatorname{magnitude}(\operatorname{recip}(((2 \cdot \pi) \cdot +i))) = \operatorname{recip}((2 \cdot \pi))}}
\vnbprop{cauchy-estimate-circle}{\vnbstmt{\begin{array}{@{}l@{}} \forall u, f.; \forall c \in \mathbb{C}, s, r, d \in \mathbb{R}, k \in \mathbb{N}. \\ \quad \operatorname{holomorphic\text{-}on}(u, f) \wedge \\ \quad 0 < s \wedge \\ \quad \operatorname{ball}(\operatorname{cc\text{-}ms}, c, s) \subseteq u \wedge \\ \quad 0 < r \wedge \\ \quad r < s \wedge \\ \quad \forall y \in \operatorname{trace}(\operatorname{circle\text{-}path}(c, r), 0, (2 \cdot \pi)).\; (\operatorname{magnitude}(f(y)) \leq d) \Rightarrow \\ \quad \quad \operatorname{magnitude}(\operatorname{recip}(((2 \cdot \pi) \cdot +i)) \cdot \operatorname{line\text{-}int}((\lambda w \in \operatorname{trace}(\operatorname{circle\text{-}path}(c, r), 0, (2 \cdot \pi)).\; (\operatorname{restrict}(f, \operatorname{trace}(\operatorname{circle\text{-}path}(c, r), 0, (2 \cdot \pi)))(w) \cdot \operatorname{power}(\operatorname{recip}((w - c)), (k + 1)))), \operatorname{circle\text{-}path}(c, r), \operatorname{circle\text{-}deriv}(c, r), 0, (2 \cdot \pi))) \\ \quad \quad \leq\; d \cdot \operatorname{power}(\operatorname{recip}(r), k) \end{array}}}
\vnbprop{cauchy-estimate-center}{\vnbstmt{\begin{array}{@{}l@{}} \forall u, f.; \forall c \in \mathbb{C}, s, r, d \in \mathbb{R}. \\ \quad \operatorname{holomorphic\text{-}on}(u, f) \wedge \\ \quad 0 < s \wedge \\ \quad \operatorname{ball}(\operatorname{cc\text{-}ms}, c, s) \subseteq u \wedge \\ \quad 0 < r \wedge \\ \quad r < s \wedge \\ \quad \forall y \in \operatorname{trace}(\operatorname{circle\text{-}path}(c, r), 0, (2 \cdot \pi)).\; (\operatorname{magnitude}(f(y)) \leq d) \Rightarrow \\ \quad \quad \operatorname{magnitude}(f(c)) \leq d \end{array}}}
\vnbprop{cauchy-coefficients-radius}{\vnbstmt{\begin{array}{@{}l@{}} \forall u, f.; \forall c \in \mathbb{C}, s, r \in \mathbb{R}, g \in (\mathbb{N} \to \mathbb{C}), o \in \mathbb{R}. \\ \quad \operatorname{holomorphic\text{-}on}(u, f) \wedge \\ \quad 0 < s \wedge \\ \quad \operatorname{ball}(\operatorname{cc\text{-}ms}, c, s) \subseteq u \wedge \\ \quad 0 < r \wedge \\ \quad r < s \wedge \\ \quad \forall k \in \mathbb{N}. \\ \quad \quad g(k) \\ \quad \quad =\; \operatorname{recip}(((2 \cdot \pi) \cdot +i)) \cdot \operatorname{line\text{-}int}((\lambda w \in \operatorname{trace}(\operatorname{circle\text{-}path}(c, r), 0, (2 \cdot \pi)).\; (\operatorname{restrict}(f, \operatorname{trace}(\operatorname{circle\text{-}path}(c, r), 0, (2 \cdot \pi)))(w) \cdot \operatorname{power}(\operatorname{recip}((w - c)), (k + 1)))), \operatorname{circle\text{-}path}(c, r), \operatorname{circle\text{-}deriv}(c, r), 0, (2 \cdot \pi)) \wedge \\ \quad 0 \leq o \wedge \\ \quad o < r \Rightarrow \\ \quad \quad o < \operatorname{cps\text{-}radius}(g) \end{array}}}
\vnbprop{holomorphic-coefficient-series}{\vnbstmt{\begin{array}{@{}l@{}} \forall u, f.; \forall c \in \mathbb{C}, s, r \in \mathbb{R}, g \in (\mathbb{N} \to \mathbb{C}), z \in \mathbb{C}. \\ \quad \operatorname{holomorphic\text{-}on}(u, f) \wedge \\ \quad 0 < s \wedge \\ \quad \operatorname{ball}(\operatorname{cc\text{-}ms}, c, s) \subseteq u \wedge \\ \quad 0 < r \wedge \\ \quad r < s \wedge \\ \quad \forall k \in \mathbb{N}. \\ \quad \quad g(k) \\ \quad \quad =\; \operatorname{recip}(((2 \cdot \pi) \cdot +i)) \cdot \operatorname{line\text{-}int}((\lambda w \in \operatorname{trace}(\operatorname{circle\text{-}path}(c, r), 0, (2 \cdot \pi)).\; (\operatorname{restrict}(f, \operatorname{trace}(\operatorname{circle\text{-}path}(c, r), 0, (2 \cdot \pi)))(w) \cdot \operatorname{power}(\operatorname{recip}((w - c)), (k + 1)))), \operatorname{circle\text{-}path}(c, r), \operatorname{circle\text{-}deriv}(c, r), 0, (2 \cdot \pi)) \wedge \\ \quad \operatorname{magnitude}((z - c)) < r \Rightarrow \\ \quad \quad f(z) = \operatorname{cc\text{-}series\text{-}limit}((\lambda \mathit{cpsk} \in \mathbb{N}.\; (g(\mathit{cpsk}) \cdot \operatorname{power}((z - c), \mathit{cpsk})))) \end{array}}}
\vnbprop{cauchy-derivative-formula}{\vnbstmt{\begin{array}{@{}l@{}} \forall u, f.; \forall c \in \mathbb{C}, s, r, o \in \mathbb{R}, k \in \mathbb{N}. \\ \quad \operatorname{holomorphic\text{-}on}(u, f) \wedge \\ \quad 0 < s \wedge \\ \quad \operatorname{ball}(\operatorname{cc\text{-}ms}, c, s) \subseteq u \wedge \\ \quad 0 < r \wedge \\ \quad r < s \wedge \\ \quad 0 < o \wedge \\ \quad o < r \Rightarrow \\ \quad \quad \operatorname{restrict}(f, \operatorname{ball}(\operatorname{cc\text{-}ms}, c, o))^{(k)}(c) \\ \quad \quad =\; \operatorname{factorial}(k) \cdot (\operatorname{recip}(((2 \cdot \pi) \cdot +i)) \cdot \operatorname{line\text{-}int}((\lambda w \in \operatorname{trace}(\operatorname{circle\text{-}path}(c, r), 0, (2 \cdot \pi)).\; (\operatorname{restrict}(f, \operatorname{trace}(\operatorname{circle\text{-}path}(c, r), 0, (2 \cdot \pi)))(w) \cdot \operatorname{power}(\operatorname{recip}((w - c)), (k + 1)))), \operatorname{circle\text{-}path}(c, r), \operatorname{circle\text{-}deriv}(c, r), 0, (2 \cdot \pi))) \end{array}}}
\vnbprop{cauchy-estimate-derivative}{\vnbstmt{\begin{array}{@{}l@{}} \forall u, f.; \forall c \in \mathbb{C}, s, r, o, d \in \mathbb{R}, k \in \mathbb{N}. \\ \quad \operatorname{holomorphic\text{-}on}(u, f) \wedge \\ \quad 0 < s \wedge \\ \quad \operatorname{ball}(\operatorname{cc\text{-}ms}, c, s) \subseteq u \wedge \\ \quad 0 < r \wedge \\ \quad r < s \wedge \\ \quad 0 < o \wedge \\ \quad o < r \wedge \\ \quad \forall y \in \operatorname{trace}(\operatorname{circle\text{-}path}(c, r), 0, (2 \cdot \pi)).\; (\operatorname{magnitude}(f(y)) \leq d) \Rightarrow \\ \quad \quad \operatorname{magnitude}(\operatorname{restrict}(f, \operatorname{ball}(\operatorname{cc\text{-}ms}, c, o))^{(k)}(c)) \\ \quad \quad \leq\; \operatorname{factorial}(k) \cdot (d \cdot \operatorname{power}(\operatorname{recip}(r), k)) \end{array}}}
\vnbprop{cie-squeeze}{\vnbstmt{\begin{array}{@{}l@{}} \forall x, k, a \in \mathbb{R}. \\ \quad ((0 \leq a) \wedge ((0 \leq k) \wedge \forall q \in \mathbb{R}.\; ((a < q) \Rightarrow (x \leq (k \cdot \operatorname{recip}(q)))))) \Rightarrow  \\ \quad \quad x \leq 0 \end{array}}}
\vnbprop{bounded-entire-deriv-zero}{\vnbstmt{\begin{array}{@{}l@{}} \forall f.; \forall d \in \mathbb{R}, a \in \mathbb{C}. \\ \quad (\operatorname{holomorphic\text{-}on}(\mathbb{C}, f) \wedge \forall y \in \mathbb{C}.\; (\operatorname{magnitude}(f(y)) \leq d)) \Rightarrow  \\ \quad \quad \operatorname{is\text{-}diff\text{-}on}(\operatorname{cc\text{-}normed\text{-}field}, \mathbb{C}, f, a, 0) \end{array}}}
\vnbprop{liouville}{\vnbstmt{\begin{array}{@{}l@{}} \forall f.; \forall d \in \mathbb{R}, z \in \mathbb{C}. \\ \quad (\operatorname{holomorphic\text{-}on}(\mathbb{C}, f) \wedge \forall y \in \mathbb{C}.\; (\operatorname{magnitude}(f(y)) \leq d)) \Rightarrow  \\ \quad \quad f(z) = f(0) \end{array}}}
\vnbprop{polynomial-holomorphic}{\vnbstmt{\begin{array}{@{}l@{}} \forall j \in \mathbb{N}, p \in (\mathbb{N} \to \mathbb{C}). \\ \quad \operatorname{holomorphic\text{-}on}(\mathbb{C}, (\lambda w \in \mathbb{C}.\; \operatorname{cc\text{-}series\text{-}partial\text{-}sum}((\lambda k \in \mathbb{N}.\; (p(k) \cdot \operatorname{power}(w, k))), j))) \end{array}}}
\vnbprop{polynomial-growth-bound}{\vnbstmt{\begin{array}{@{}l@{}} \forall j \in \mathbb{N}, p \in (\mathbb{N} \to \mathbb{C}). \\ \quad \operatorname{cc\text{-}series\text{-}partial\text{-}sum}((\lambda k \in \mathbb{N}.\; \operatorname{magnitude}(p(k))), (j + 1)) \in \mathbb{R} \wedge \\ \quad 0 \leq \operatorname{cc\text{-}series\text{-}partial\text{-}sum}((\lambda k \in \mathbb{N}.\; \operatorname{magnitude}(p(k))), (j + 1)) \wedge \\ \quad \forall w \in \mathbb{C}. \\ \quad \quad (1 \leq \operatorname{magnitude}(w)) \Rightarrow  \\ \quad \quad \quad \operatorname{magnitude}(\operatorname{cc\text{-}series\text{-}partial\text{-}sum}((\lambda k \in \mathbb{N}.\; (p(k) \cdot \operatorname{power}(w, k))), (j + 1))) \\ \quad \quad \quad \leq\; \operatorname{cc\text{-}series\text{-}partial\text{-}sum}((\lambda k \in \mathbb{N}.\; \operatorname{magnitude}(p(k))), (j + 1)) \cdot \operatorname{power}(\operatorname{magnitude}(w), j) \end{array}}}
\vnbprop{fundamental-theorem-of-algebra}{\vnbstmt{\begin{array}{@{}l@{}} \forall p \in (\mathbb{N} \to \mathbb{C}), m \in \mathbb{N}. \\ \quad \neg (p((m + 1)) = 0) \Rightarrow  \\ \quad \quad \exists z \in \mathbb{C}. \\ \quad \quad \quad \operatorname{cc\text{-}series\text{-}partial\text{-}sum}((\lambda k \in \mathbb{N}.\; (p(k) \cdot \operatorname{power}(z, k))), ((m + 1) + 1)) = 0 \end{array}}}

\vnbbookfile{theorem-library/principal-part}{theorem-library/principal-part}
\noindent{\small principal-part.scm -- WAVE 2 (2026-10-27..30), wired 2026-10-30 from scratchpad/ch3e/ch3e-pp.scm. Library: complex-analysis.}\par
\vnbprop{principal-part-integral-step}{\vnbstmt{\begin{array}{@{}l@{}} \forall m \in \mathbb{N}, \mathit{ppa}, c, d, f, a, b. \\ \quad \mathit{ppa} \in (\mathbb{N} \to \mathbb{C}) \wedge \\ \quad c \in \mathbb{C} \wedge \\ \quad \operatorname{is\text{-}road}(d, f, a, b) \wedge \\ \quad \operatorname{trace}(d, a, b) \subseteq \operatorname{difference}(\mathbb{C}, \operatorname{singleton}(c)) \wedge \\ \quad d(b) = d(a) \Rightarrow \\ \quad \quad \operatorname{holomorphic\text{-}on}(\operatorname{difference}(\mathbb{C}, \operatorname{singleton}(c)), \\ \quad \quad \quad (\lambda w \in \operatorname{difference}(\mathbb{C}, \operatorname{singleton}(c)).\; \\ \quad \quad \quad \quad \operatorname{cc\text{-}series\text{-}partial\text{-}sum}((\lambda j \in \mathbb{N}.\; (\operatorname{ppa}((j + 1)) \cdot \operatorname{power}(\operatorname{recip}((w - c)), (j + 1)))), (m + 1)))) \wedge \\ \quad \quad \operatorname{line\text{-}int}((\lambda w \in \operatorname{difference}(\mathbb{C}, \operatorname{singleton}(c)).\; \\ \quad \quad \quad \quad \operatorname{cc\text{-}series\text{-}partial\text{-}sum}((\lambda j \in \mathbb{N}.\; (\operatorname{ppa}((j + 1)) \cdot \operatorname{power}(\operatorname{recip}((w - c)), (j + 1)))), (m + 1))), \\ \quad \quad \quad d, f, a, b) \\ \quad \quad =\; \operatorname{ppa}(1) \cdot \operatorname{line\text{-}int}((\lambda v \in \operatorname{difference}(\mathbb{C}, \operatorname{singleton}(c)).\; \operatorname{power}(\operatorname{recip}((v - c)), 1)), d, f, a, b) \end{array}}}
\vnbprop{principal-part-integral}{\vnbstmt{\begin{array}{@{}l@{}} \forall \mathit{ppa}, c, m, d, a, b.; \forall n \in \mathbb{N}. \\ \quad \mathit{ppa} \in (\mathbb{N} \to \mathbb{C}) \wedge \\ \quad c \in \mathbb{C} \wedge \\ \quad \operatorname{is\text{-}road}(m, d, a, b) \wedge \\ \quad \operatorname{trace}(m, a, b) \subseteq \operatorname{difference}(\mathbb{C}, \operatorname{singleton}(c)) \wedge \\ \quad m(b) = m(a) \wedge \\ \quad 1 \leq n \Rightarrow \\ \quad \quad \operatorname{recip}(((2 \cdot \pi) \cdot +i)) \cdot \operatorname{line\text{-}int}((\lambda w \in \operatorname{difference}(\mathbb{C}, \operatorname{singleton}(c)).\; \operatorname{cc\text{-}series\text{-}partial\text{-}sum}((\lambda j \in \mathbb{N}.\; (\operatorname{ppa}((j + 1)) \cdot \operatorname{power}(\operatorname{recip}((w - c)), (j + 1)))), n)), m, d, a, b) \\ \quad \quad =\; \operatorname{ppa}(1) \cdot \operatorname{winding}(m, d, a, b, c) \end{array}}}

\vnbbookfile{structure-library/ca1-vocab}{structure-library/ca1-vocab}
\noindent{\small ca1-vocab.scm -- THE COMPLEX NUMBERS AS 2-BY-2 REAL MATRICES (complex-analysis.pdf, Proposition 1.3).  DEFINITIONS ONLY; the laws and the proposition are proven in theorem-library/ca1-matrix-rep.scm.  Agent CA-1, 2026-10-19 (NEW DEFINITIONS, the user's to confirm).}\par
\par\medskip\noindent\textbf{Definition} (\texttt{cc-matrix-of}).\;$\operatorname{cc\text{-}matrix\text{-}of}(a, b) :=$ ${[}{[}a, b], {[}-b, a]]$
\par\medskip\noindent\textbf{Definition} (\texttt{cc-matrices}).\;
\begin{equation*}
(\operatorname{cc\text{-}matrices} \simeq \{\operatorname{cmq\_} \in \operatorname{mat}(2, 2, \mathbb{R}) : \exists \operatorname{cma\_}, \operatorname{cmb\_} \in \mathbb{R}.\; (\operatorname{cmq\_} = \operatorname{cc\text{-}matrix\text{-}of}(\operatorname{cma\_}, \operatorname{cmb\_}))\})
\end{equation*}

\par\medskip\noindent\textbf{Definition} (\texttt{cc-matrix-rep}).\;
\begin{equation*}
(\operatorname{cc\text{-}matrix\text{-}rep} \simeq (\lambda \operatorname{cmz\_} \in \mathbb{C}.\; \operatorname{cc\text{-}matrix\text{-}of}(\operatorname{real\text{-}part}(\operatorname{cmz\_}), \operatorname{imag\text{-}part}(\operatorname{cmz\_}))))
\end{equation*}


\vnbbookfile{theorem-library/ca1-matrix-rep}{theorem-library/ca1-matrix-rep}
\noindent{\small ca1-matrix-rep.scm -- THE COMPLEX NUMBERS AS 2-BY-2 REAL MATRICES (complex-analysis.pdf, Proposition 1.3).  Agent CA-1, 2026-10-19.}\par
\vnbprop{mat22-literal-in-mat}{\vnbstmt{\forall x.; \forall a, b, c, d \in x.\; ({[}{[}a, b], {[}c, d]] \in \operatorname{mat}(2, 2, x))}}
\vnbprop{mat22-ext}{\vnbstmt{\begin{array}{@{}l@{}} \forall x.; \forall p, q \in \operatorname{mat}(2, 2, x). \\ \quad \operatorname{entry}(p, 1, 1) = \operatorname{entry}(q, 1, 1) \wedge \\ \quad \operatorname{entry}(p, 1, 2) = \operatorname{entry}(q, 1, 2) \wedge \\ \quad \operatorname{entry}(p, 2, 1) = \operatorname{entry}(q, 2, 1) \wedge \\ \quad \operatorname{entry}(p, 2, 2) = \operatorname{entry}(q, 2, 2) \Rightarrow \\ \quad \quad p = q \end{array}}}
\vnbprop{cc-matrix-of-entries}{\vnbstmt{\begin{array}{@{}l@{}} \forall a, b \in \mathbb{R}. \\ \quad \operatorname{entry}(\operatorname{cc\text{-}matrix\text{-}of}(a, b), 1, 1) = a \wedge \\ \quad \operatorname{entry}(\operatorname{cc\text{-}matrix\text{-}of}(a, b), 1, 2) = b \wedge \\ \quad \operatorname{entry}(\operatorname{cc\text{-}matrix\text{-}of}(a, b), 2, 1) = -b \wedge \\ \quad \operatorname{entry}(\operatorname{cc\text{-}matrix\text{-}of}(a, b), 2, 2) = a \end{array}}}
\vnbprop{cc-matrix-of-in-mat}{\vnbstmt{\forall a, b \in \mathbb{R}.\; (\operatorname{cc\text{-}matrix\text{-}of}(a, b) \in \operatorname{mat}(2, 2, \mathbb{R}))}}
\vnbprop{rr-matmul-2x2-type}{\vnbstmt{\begin{array}{@{}l@{}} \forall p, q \in \operatorname{mat}(2, 2, \mathbb{R}). \\ \quad \operatorname{matmul}(\operatorname{normed\text{-}field\text{-}as\text{-}commutative\text{-}ring}(\operatorname{rr\text{-}normed\text{-}field}), p, q) \in \operatorname{mat}(2, 2, \mathbb{R}) \end{array}}}
\vnbprop{rr-matadd-2x2-type}{\vnbstmt{\begin{array}{@{}l@{}} \forall p, q \in \operatorname{mat}(2, 2, \mathbb{R}). \\ \quad \operatorname{matadd}(\operatorname{normed\text{-}field\text{-}as\text{-}commutative\text{-}ring}(\operatorname{rr\text{-}normed\text{-}field}), p, q) \in \operatorname{mat}(2, 2, \mathbb{R}) \end{array}}}
\vnbprop{rr-matmul-2x2-entry}{\vnbstmt{\begin{array}{@{}l@{}} \forall p, q \in \operatorname{mat}(2, 2, \mathbb{R}), i, k \in \operatorname{interval}(1, 2). \\ \quad \operatorname{entry}(\operatorname{matmul}(\operatorname{normed\text{-}field\text{-}as\text{-}commutative\text{-}ring}(\operatorname{rr\text{-}normed\text{-}field}), p, q), i, k) \\ \quad =\; (\operatorname{entry}(p, i, 1) \cdot \operatorname{entry}(q, 1, k)) + (\operatorname{entry}(p, i, 2) \cdot \operatorname{entry}(q, 2, k)) \end{array}}}
\vnbprop{rr-matadd-2x2-entry}{\vnbstmt{\begin{array}{@{}l@{}} \forall p, q \in \operatorname{mat}(2, 2, \mathbb{R}), i, k \in \operatorname{interval}(1, 2). \\ \quad \operatorname{entry}(\operatorname{matadd}(\operatorname{normed\text{-}field\text{-}as\text{-}commutative\text{-}ring}(\operatorname{rr\text{-}normed\text{-}field}), p, q), i, k) \\ \quad =\; \operatorname{entry}(p, i, k) + \operatorname{entry}(q, i, k) \end{array}}}
\vnbprop{cc-matrix-of-mul}{\vnbstmt{\begin{array}{@{}l@{}} \forall a, b, c, d \in \mathbb{R}. \\ \quad \operatorname{matmul}(\operatorname{normed\text{-}field\text{-}as\text{-}commutative\text{-}ring}(\operatorname{rr\text{-}normed\text{-}field}), \operatorname{cc\text{-}matrix\text{-}of}(a, b), \operatorname{cc\text{-}matrix\text{-}of}(c, d)) \\ \quad =\; \operatorname{cc\text{-}matrix\text{-}of}(((a \cdot c) - (b \cdot d)), ((a \cdot d) + (b \cdot c))) \end{array}}}
\vnbprop{cc-matrix-of-add}{\vnbstmt{\begin{array}{@{}l@{}} \forall a, b, c, d \in \mathbb{R}. \\ \quad \operatorname{matadd}(\operatorname{normed\text{-}field\text{-}as\text{-}commutative\text{-}ring}(\operatorname{rr\text{-}normed\text{-}field}), \operatorname{cc\text{-}matrix\text{-}of}(a, b), \operatorname{cc\text{-}matrix\text{-}of}(c, d)) \\ \quad =\; \operatorname{cc\text{-}matrix\text{-}of}((a + c), (b + d)) \end{array}}}
\vnbprop{cc-matrix-of-scale}{\vnbstmt{\begin{array}{@{}l@{}} \forall c, a, b \in \mathbb{R}. \\ \quad \operatorname{matscale}(\operatorname{normed\text{-}field\text{-}as\text{-}commutative\text{-}ring}(\operatorname{rr\text{-}normed\text{-}field}), c, \operatorname{cc\text{-}matrix\text{-}of}(a, b)) \\ \quad =\; \operatorname{cc\text{-}matrix\text{-}of}((c \cdot a), (c \cdot b)) \end{array}}}
\vnbprop{cc-matrix-of-injective}{\vnbstmt{\begin{array}{@{}l@{}} \forall a, b, c, d \in \mathbb{R}. \\ \quad (\operatorname{cc\text{-}matrix\text{-}of}(a, b) = \operatorname{cc\text{-}matrix\text{-}of}(c, d)) \Rightarrow ((a = c) \wedge (b = d)) \end{array}}}
\vnbprop{cc-matrix-rep-apply}{\vnbstmt{\forall z \in \mathbb{C}.\; (\operatorname{cc\text{-}matrix\text{-}rep}(z) \simeq \operatorname{cc\text{-}matrix\text{-}of}(\operatorname{real\text{-}part}(z), \operatorname{imag\text{-}part}(z)))}}
\vnbprop{cc-matrix-rep-add}{\vnbstmt{\begin{array}{@{}l@{}} \forall z, w \in \mathbb{C}. \\ \quad \operatorname{cc\text{-}matrix\text{-}rep}((z + w)) \\ \quad =\; \operatorname{matadd}(\operatorname{normed\text{-}field\text{-}as\text{-}commutative\text{-}ring}(\operatorname{rr\text{-}normed\text{-}field}), \operatorname{cc\text{-}matrix\text{-}rep}(z), \operatorname{cc\text{-}matrix\text{-}rep}(w)) \end{array}}}
\vnbprop{cc-matrix-rep-mul}{\vnbstmt{\begin{array}{@{}l@{}} \forall z, w \in \mathbb{C}. \\ \quad \operatorname{cc\text{-}matrix\text{-}rep}((z \cdot w)) \\ \quad =\; \operatorname{matmul}(\operatorname{normed\text{-}field\text{-}as\text{-}commutative\text{-}ring}(\operatorname{rr\text{-}normed\text{-}field}), \operatorname{cc\text{-}matrix\text{-}rep}(z), \operatorname{cc\text{-}matrix\text{-}rep}(w)) \end{array}}}
\vnbprop{cc-matrix-rep-real-mul}{\vnbstmt{\begin{array}{@{}l@{}} \forall c \in \mathbb{R}, z \in \mathbb{C}. \\ \quad \operatorname{cc\text{-}matrix\text{-}rep}((c \cdot z)) \\ \quad =\; \operatorname{matscale}(\operatorname{normed\text{-}field\text{-}as\text{-}commutative\text{-}ring}(\operatorname{rr\text{-}normed\text{-}field}), c, \operatorname{cc\text{-}matrix\text{-}rep}(z)) \end{array}}}
\vnbprop{cc-matrix-rep-one}{\vnbstmt{\operatorname{cc\text{-}matrix\text{-}rep}(1) = \operatorname{identmat}(\operatorname{normed\text{-}field\text{-}as\text{-}commutative\text{-}ring}(\operatorname{rr\text{-}normed\text{-}field}), 2)}}
\vnbprop{cc-matrix-rep-in-fun}{\vnbstmt{\operatorname{cc\text{-}matrix\text{-}rep} \in (\mathbb{C} \to \operatorname{cc\text{-}matrices})}}
\vnbprop{cc-matrix-rep-bijection}{\vnbstmt{\operatorname{cc\text{-}matrix\text{-}rep} \in \mathrm{Bij}(\mathbb{C}, \operatorname{cc\text{-}matrices})}}
\vnbprop{cc-matrix-representation}{\vnbstmt{\begin{array}{@{}l@{}} \operatorname{cc\text{-}matrix\text{-}rep} \in \mathrm{Bij}(\mathbb{C}, \operatorname{cc\text{-}matrices}) \wedge \\ \forall z, w \in \mathbb{C}. \\ \quad \operatorname{cc\text{-}matrix\text{-}rep}((z + w)) \\ \quad =\; \operatorname{matadd}(\operatorname{normed\text{-}field\text{-}as\text{-}commutative\text{-}ring}(\operatorname{rr\text{-}normed\text{-}field}), \operatorname{cc\text{-}matrix\text{-}rep}(z), \operatorname{cc\text{-}matrix\text{-}rep}(w)) \wedge \\ \forall z, w \in \mathbb{C}. \\ \quad \operatorname{cc\text{-}matrix\text{-}rep}((z \cdot w)) \\ \quad =\; \operatorname{matmul}(\operatorname{normed\text{-}field\text{-}as\text{-}commutative\text{-}ring}(\operatorname{rr\text{-}normed\text{-}field}), \operatorname{cc\text{-}matrix\text{-}rep}(z), \operatorname{cc\text{-}matrix\text{-}rep}(w)) \wedge \\ \operatorname{cc\text{-}matrix\text{-}rep}(1) = \operatorname{identmat}(\operatorname{normed\text{-}field\text{-}as\text{-}commutative\text{-}ring}(\operatorname{rr\text{-}normed\text{-}field}), 2) \wedge \\ \forall c \in \mathbb{R}, z \in \mathbb{C}. \\ \quad \operatorname{cc\text{-}matrix\text{-}rep}((c \cdot z)) \\ \quad =\; \operatorname{matscale}(\operatorname{normed\text{-}field\text{-}as\text{-}commutative\text{-}ring}(\operatorname{rr\text{-}normed\text{-}field}), c, \operatorname{cc\text{-}matrix\text{-}rep}(z)) \end{array}}}

\vnbbookfile{theorem-library/ca1-alg-closed}{theorem-library/ca1-alg-closed}
\noindent{\small ca1-alg-closed.scm -- C IS ALGEBRAICALLY CLOSED, IN THE NOTES' FORM, AND THE GROWTH ESTIMATE (6)-(7) OF ITS PROOF (complex-analysis.pdf, Theorem 1.15).  Agent CA-1, 2026-10-19.}\par
\vnbprop{cc-algebraically-closed}{\vnbstmt{\begin{array}{@{}l@{}} \forall p \in (\mathbb{N} \to \mathbb{C}), n \in \mathbb{N}. \\ \quad ((1 \leq n) \wedge \neg (p(n) = 0)) \Rightarrow  \\ \quad \quad \exists z \in \mathbb{C}. \\ \quad \quad \quad \operatorname{cc\text{-}series\text{-}partial\text{-}sum}((\lambda k \in \mathbb{N}.\; (p(k) \cdot \operatorname{power}(z, k))), (n + 1)) = 0 \end{array}}}
\vnbprop{cc-polynomial-lower-growth}{\vnbstmt{\begin{array}{@{}l@{}} \forall p \in (\mathbb{N} \to \mathbb{C}), m \in \mathbb{N}, z \in \mathbb{C}. \\ \quad 1 \leq \operatorname{magnitude}(z) \wedge \\ \quad 2 \cdot \operatorname{cc\text{-}series\text{-}partial\text{-}sum}((\lambda k \in \mathbb{N}.\; \operatorname{magnitude}(p(k))), (m + 1)) \\ \quad \leq\; \operatorname{magnitude}(p((m + 1))) \cdot \operatorname{magnitude}(z) \Rightarrow \\ \quad \quad \operatorname{magnitude}(p((m + 1))) \cdot \operatorname{power}(\operatorname{magnitude}(z), (m + 1)) \\ \quad \quad \leq\; 2 \cdot \operatorname{magnitude}(\operatorname{cc\text{-}series\text{-}partial\text{-}sum}((\lambda k \in \mathbb{N}.\; (p(k) \cdot \operatorname{power}(z, k))), ((m + 1) + 1))) \end{array}}}

\vnbbookfile{theorem-library/poly-fn-holomorphic}{theorem-library/poly-fn-holomorphic}
\noindent{\small poly-fn-holomorphic.scm -- complex analysis: a polynomial function on an open set is holomorphic and its derivative is a polynomial function (CA-MONTEL "POLY", 2026-10-26).  Library: complex-analysis. The helpers below are copied from poly-fn-laws.scm (each proof file has its own environment).}\par
\vnbprop{cc-ps-deriv-poly}{\vnbstmt{\begin{array}{@{}l@{}} \forall n \in \mathbb{N}, c \in (\mathbb{N} \to \mathbb{C}). \\ \quad \forall k \in \mathbb{N}.\; ((n \leq k) \Rightarrow (c(k) = 0)) \Rightarrow  \\ \quad \quad \exists d. \\ \quad \quad \quad \operatorname{is\text{-}poly\text{-}fn\text{-}on}(\mathbb{C}, \mathbb{C}, d) \wedge \\ \quad \quad \quad \forall w \in \mathbb{C}. \\ \quad \quad \quad \quad \operatorname{is\text{-}diff\text{-}on}(\operatorname{cc\text{-}normed\text{-}field}, \mathbb{C}, \\ \quad \quad \quad \quad \quad (\lambda x \in \mathbb{C}.\; \operatorname{cc\text{-}series\text{-}partial\text{-}sum}((\lambda j \in \mathbb{N}.\; (c(j) \cdot \operatorname{power}(x, j))), n)), \\ \quad \quad \quad \quad \quad w, d(w)) \end{array}}}
\vnbprop{poly-fn-holomorphic}{\vnbstmt{\begin{array}{@{}l@{}} \forall u, f. \\ \quad (\operatorname{is\text{-}open}(\operatorname{nf\text{-}metric\text{-}space}(\operatorname{cc\text{-}normed\text{-}field}), u) \wedge \operatorname{is\text{-}poly\text{-}fn\text{-}on}(\mathbb{C}, u, f)) \Rightarrow  \\ \quad \quad \operatorname{holomorphic\text{-}on}(u, f) \wedge \\ \quad \quad \exists d. \\ \quad \quad \quad (\operatorname{is\text{-}poly\text{-}fn\text{-}on}(\mathbb{C}, u, d) \wedge \forall w \in u.\; \operatorname{is\text{-}diff\text{-}on}(\operatorname{cc\text{-}normed\text{-}field}, u, f, w, d(w))) \end{array}}}
\vnbprop{poly-fn-fta}{\vnbstmt{\begin{array}{@{}l@{}} \forall f.; \forall n \in \mathbb{N}. \\ \quad ((1 \leq n) \wedge \operatorname{is\text{-}poly\text{-}deg}(\mathbb{C}, \mathbb{C}, f, n)) \Rightarrow \exists z \in \mathbb{C}.\; (f(z) = 0) \end{array}}}

\vnbbookfile{theorem-library/winding-number-integer}{theorem-library/winding-number-integer}
\noindent{\small winding-number-integer.scm -- THE WINDING NUMBER IS AN INTEGER (complex-analysis.pdf ch. 3, Proposition 3.27, first half), with the two pieces of road calculus it rests on. Agent CA-5 of week 3 of the October roadmap, 2026-10-04. Definitions: structure-library/winding-number.scm (WINDING, by the integral (73)).}\par
\vnbprop{road-product-constant}{\vnbstmt{\begin{array}{@{}l@{}} \forall m, c, d, f, a, b. \\ \quad \operatorname{is\text{-}road}(m, c, a, b) \wedge \\ \quad \operatorname{is\text{-}road}(d, f, a, b) \wedge \\ \quad \forall t \in \operatorname{ccint}(a, b).\; (((c(t) \cdot d(t)) + (m(t) \cdot f(t))) = 0) \Rightarrow \\ \quad \quad (m(b) \cdot d(b)) = (m(a) \cdot d(a)) \end{array}}}
\vnbprop{line-int-primitive-road}{\vnbstmt{\begin{array}{@{}l@{}} \forall u, f, m, c, a, b. \\ \quad u \subseteq \mathbb{C} \wedge \\ \quad f \in (u \to \mathbb{C}) \wedge \\ \quad \forall y \in u. \\ \quad \quad \operatorname{is\text{-}continuous\text{-}at}(\operatorname{subspace\text{-}ms}(\operatorname{nf\text{-}metric\text{-}space}(\operatorname{cc\text{-}normed\text{-}field}), u), \\ \quad \quad \quad \operatorname{nf\text{-}metric\text{-}space}(\operatorname{cc\text{-}normed\text{-}field}), f, y) \wedge \\ \quad \operatorname{is\text{-}road}(m, c, a, b) \wedge \\ \quad \operatorname{trace}(m, a, b) \subseteq u \Rightarrow \\ \quad \quad \exists h. \\ \quad \quad \quad \operatorname{is\text{-}road}(h, (\lambda t \in \operatorname{ccint}(a, b).\; (f(m(t)) \cdot c(t))), a, b) \wedge \\ \quad \quad \quad \operatorname{line\text{-}int}(f, m, c, a, b) = (h(b) - h(a)) \end{array}}}
\vnbprop{winding-number-integer}{\vnbstmt{\begin{array}{@{}l@{}} \forall m, c, a, b.; \forall z \in \mathbb{C}. \\ \quad (\operatorname{is\text{-}road}(m, c, a, b) \wedge ((m(b) = m(a)) \wedge \neg (z \in \operatorname{trace}(m, a, b)))) \Rightarrow  \\ \quad \quad \exists k \in \mathbb{Z}.\; (\operatorname{winding}(m, c, a, b, z) = k) \end{array}}}

\vnbbookfile{theorem-library/winding-number-components}{theorem-library/winding-number-components}
\noindent{\small winding-number-components.scm -- THE WINDING NUMBER IS LOCALLY CONSTANT OFF THE TRACE, CONSTANT ON A CONNECTED SET OFF THE TRACE, AND 0 ON THE UNBOUNDED COMPONENT (complex-analysis.pdf ch. 3: Proposition 3.27, second half; Corollary 3.30). Agent CA-5 of week 3 of the October roadmap, 2026-10-04.}\par
\vnbprop{winding-disc-constant}{\vnbstmt{\begin{array}{@{}l@{}} \forall m, d, a, b.; \forall c \in \mathbb{C}, r \in \mathbb{R}, z \in \mathbb{C}. \\ \quad \operatorname{is\text{-}road}(m, d, a, b) \wedge \\ \quad m(b) = m(a) \wedge \\ \quad 0 < r \wedge \\ \quad \forall y \in \operatorname{trace}(m, a, b).\; (r \leq \operatorname{magnitude}((y - c))) \wedge \\ \quad \operatorname{magnitude}((z - c)) < r \Rightarrow \\ \quad \quad \operatorname{winding}(m, d, a, b, z) = \operatorname{winding}(m, d, a, b, c) \end{array}}}
\vnbprop{winding-locally-constant}{\vnbstmt{\begin{array}{@{}l@{}} \forall m, c, a, b.; \forall z \in \mathbb{C}. \\ \quad (\operatorname{is\text{-}road}(m, c, a, b) \wedge ((m(b) = m(a)) \wedge \neg (z \in \operatorname{trace}(m, a, b)))) \Rightarrow  \\ \quad \quad \exists r. \\ \quad \quad \quad \operatorname{pos\text{-}rr}(r) \wedge \\ \quad \quad \quad \forall \mathit{wcz} \in \operatorname{ball}(\operatorname{cc\text{-}ms}, z, r). \\ \quad \quad \quad \quad \neg (\mathit{wcz} \in \operatorname{trace}(m, a, b)) \wedge (\operatorname{winding}(m, c, a, b, \mathit{wcz}) = \operatorname{winding}(m, c, a, b, z)) \end{array}}}
\vnbprop{winding-constant-on-connected}{\vnbstmt{\begin{array}{@{}l@{}} \forall m, c, a, b, \mathit{wca}.; \forall p, q \in \mathit{wca}. \\ \quad \operatorname{is\text{-}road}(m, c, a, b) \wedge \\ \quad m(b) = m(a) \wedge \\ \quad \operatorname{is\text{-}connected}(\operatorname{cc\text{-}ms}, \mathit{wca}) \wedge \\ \quad \forall x \in \mathit{wca}.\; \neg (x \in \operatorname{trace}(m, a, b)) \Rightarrow \\ \quad \quad \operatorname{winding}(m, c, a, b, p) = \operatorname{winding}(m, c, a, b, q) \end{array}}}
\vnbprop{winding-zero-unbounded}{\vnbstmt{\begin{array}{@{}l@{}} \forall m, d, a, b.; \forall c \in \mathbb{C}, r \in \mathbb{R}, \mathit{wca}.; \forall u, p \in \mathit{wca}. \\ \quad \operatorname{is\text{-}road}(m, d, a, b) \wedge \\ \quad m(b) = m(a) \wedge \\ \quad 0 < r \wedge \\ \quad \operatorname{trace}(m, a, b) \subseteq \operatorname{ball}(\operatorname{cc\text{-}ms}, c, r) \wedge \\ \quad \operatorname{is\text{-}connected}(\operatorname{cc\text{-}ms}, \mathit{wca}) \wedge \\ \quad \forall x \in \mathit{wca}.\; \neg (x \in \operatorname{trace}(m, a, b)) \wedge \\ \quad \neg (u \in \operatorname{ball}(\operatorname{cc\text{-}ms}, c, r)) \Rightarrow \\ \quad \quad \operatorname{winding}(m, d, a, b, p) = 0 \end{array}}}
\vnbprop{winding-juxta}{\vnbstmt{\begin{array}{@{}l@{}} \forall m_{1}, d, m_{2}, f, a, b, c.; \forall z \in \mathbb{C}. \\ \quad \operatorname{is\text{-}road}(m_{1}, d, a, b) \wedge \\ \quad \operatorname{is\text{-}road}(m_{2}, f, b, c) \wedge \\ \quad \operatorname{m1}(b) = \operatorname{m2}(b) \wedge \\ \quad \neg (z \in \operatorname{trace}(m_{1}, a, b)) \wedge \\ \quad \neg (z \in \operatorname{trace}(m_{2}, b, c)) \Rightarrow \\ \quad \quad \operatorname{winding}(\operatorname{juxta}(m_{1}, m_{2}, a, b, c), \operatorname{juxta}(d, f, a, b, c), a, c, z) \\ \quad \quad =\; \operatorname{winding}(m_{1}, d, a, b, z) + \operatorname{winding}(m_{2}, f, b, c, z) \end{array}}}
\vnbprop{winding-opposite}{\vnbstmt{\begin{array}{@{}l@{}} \forall m, c, a, b, \mathit{prgb}, \mathit{prdb}.; \forall z \in \mathbb{C}. \\ \quad \operatorname{is\text{-}road}(m, c, a, b) \wedge \\ \quad \mathit{prgb} \in (\operatorname{ccint}(a, b) \to \mathbb{C}) \wedge \\ \quad \mathit{prdb} \in (\operatorname{ccint}(a, b) \to \mathbb{C}) \wedge \\ \quad \forall y \in \operatorname{ccint}(a, b).\; (\operatorname{prgb}(y) \simeq m(((a + b) - y))) \wedge \\ \quad \forall y \in \operatorname{ccint}(a, b).\; (\operatorname{prdb}(y) \simeq (-1 \cdot c(((a + b) - y)))) \wedge \\ \quad \neg (z \in \operatorname{trace}(m, a, b)) \Rightarrow \\ \quad \quad \operatorname{winding}(\mathit{prgb}, \mathit{prdb}, a, b, z) = (-1 \cdot \operatorname{winding}(m, c, a, b, z)) \end{array}}}

\vnbbookfile{theorem-library/cauchy-centre-mean-value}{theorem-library/cauchy-centre-mean-value}
\noindent{\small cauchy-centre-mean-value.scm -- THE INTEGRAL FORMULA AT THE CENTRE, THE MEAN VALUE PROPERTY, THE MEAN VALUE INEQUALITY, THE MAXIMUM MODULUS PRINCIPLE ON A DISC (complex-analysis.pdf 3.7 at z = c; 3.8: equation (92), Definition 3.54, Proposition 3.55 in its disc form).  Agent CA-4 of week 2 of the October roadmap}\par
\vnbprop{cc-continuous-at-local}{\vnbstmt{\begin{array}{@{}l@{}} \forall u.; \forall h \in (u \to \mathbb{C}), v.; \forall w \in v. \\ \quad u \subseteq \mathbb{C} \wedge \\ \quad \operatorname{is\text{-}open}(\operatorname{nf\text{-}metric\text{-}space}(\operatorname{cc\text{-}normed\text{-}field}), v) \wedge \\ \quad v \subseteq u \wedge \\ \quad \operatorname{is\text{-}continuous\text{-}at}(\operatorname{subspace\text{-}ms}(\operatorname{nf\text{-}metric\text{-}space}(\operatorname{cc\text{-}normed\text{-}field}), v), \\ \quad \quad \operatorname{nf\text{-}metric\text{-}space}(\operatorname{cc\text{-}normed\text{-}field}), \operatorname{restrict}(h, v), w) \Rightarrow \\ \quad \quad \operatorname{is\text{-}continuous\text{-}at}(\operatorname{subspace\text{-}ms}(\operatorname{nf\text{-}metric\text{-}space}(\operatorname{cc\text{-}normed\text{-}field}), u), \\ \quad \quad \quad \operatorname{nf\text{-}metric\text{-}space}(\operatorname{cc\text{-}normed\text{-}field}), h, w) \end{array}}}
\vnbprop{holomorphic-difference-quotient}{\vnbstmt{\begin{array}{@{}l@{}} \forall u, f.; \forall c \in u. \\ \quad \operatorname{holomorphic\text{-}on}(u, f) \Rightarrow  \\ \quad \quad \exists q \in (u \to \mathbb{C}). \\ \quad \quad \quad \forall y \in u. \\ \quad \quad \quad \quad \operatorname{is\text{-}continuous\text{-}at}(\operatorname{subspace\text{-}ms}(\operatorname{nf\text{-}metric\text{-}space}(\operatorname{cc\text{-}normed\text{-}field}), u), \\ \quad \quad \quad \quad \quad \operatorname{nf\text{-}metric\text{-}space}(\operatorname{cc\text{-}normed\text{-}field}), q, y) \wedge \\ \quad \quad \quad \operatorname{holomorphic\text{-}on}(\operatorname{difference}(u, \operatorname{singleton}(c)), \operatorname{restrict}(q, \operatorname{difference}(u, \operatorname{singleton}(c)))) \wedge \\ \quad \quad \quad \forall x \in u.\; ((f(x) - f(c)) = (q(x) \cdot (x - c))) \end{array}}}
\vnbprop{circle-trace-in-punctured-ball}{\vnbstmt{\begin{array}{@{}l@{}} \forall c \in \mathbb{C}, r, \mathit{cmvbr} \in \mathbb{R}, y \in \operatorname{trace}(\operatorname{circle\text{-}path}(c, r), 0, (2 \cdot \pi)). \\ \quad ((0 < r) \wedge (r < \mathit{cmvbr})) \Rightarrow ((y \in \operatorname{ball}(\operatorname{cc\text{-}ms}, c, \mathit{cmvbr})) \wedge \neg (y = c)) \end{array}}}
\vnbprop{cauchy-integral-formula-centre}{\vnbstmt{\begin{array}{@{}l@{}} \forall u, f.; \forall c \in \mathbb{C}, r, \mathit{cmvbr} \in \mathbb{R}. \\ \quad (\operatorname{holomorphic\text{-}on}(u, f) \wedge ((0 < r) \wedge ((r < \mathit{cmvbr}) \wedge (\operatorname{ball}(\operatorname{cc\text{-}ms}, c, \mathit{cmvbr}) \subseteq u)))) \Rightarrow  \\ \quad \quad f(c) \\ \quad \quad =\; \operatorname{recip}(((2 \cdot \pi) \cdot +i)) \cdot \operatorname{line\text{-}int}((\lambda w \in \operatorname{difference}(u, \operatorname{singleton}(c)).\; (f(w) \cdot \operatorname{recip}((w - c)))), \operatorname{circle\text{-}path}(c, r), \operatorname{circle\text{-}deriv}(c, r), 0, (2 \cdot \pi)) \end{array}}}
\vnbprop{holomorphic-mean-value}{\vnbstmt{\begin{array}{@{}l@{}} \forall u, f.; \forall c \in \mathbb{C}, r, \mathit{cmvbr} \in \mathbb{R}. \\ \quad (\operatorname{holomorphic\text{-}on}(u, f) \wedge ((0 < r) \wedge ((r < \mathit{cmvbr}) \wedge (\operatorname{ball}(\operatorname{cc\text{-}ms}, c, \mathit{cmvbr}) \subseteq u)))) \Rightarrow  \\ \quad \quad f(c) = (\operatorname{recip}((2 \cdot \pi)) \cdot \int_{0}^{(2 \cdot \pi)} (\lambda t \in \operatorname{ccint}(0, (2 \cdot \pi)).\; f((c + (r \cdot \exp (+i \cdot t)))))) \end{array}}}
\vnbprop{pw-int-nonneg-zero}{\vnbstmt{\begin{array}{@{}l@{}} \forall f, h, a, b.; \forall t \in \operatorname{ccint}(a, b). \\ \quad \operatorname{is\text{-}primitive}(f, h, a, b) \wedge \\ \quad \operatorname{is\text{-}continuous\text{-}on}(h, \operatorname{ccint}(a, b)) \wedge \\ \quad \forall t \in \operatorname{ccint}(a, b).\; (0 \leq h(t)) \wedge \\ \quad \int_{a}^{b} h = 0 \Rightarrow \\ \quad \quad h(t) = 0 \end{array}}}
\vnbprop{holomorphic-circle-max-constant}{\vnbstmt{\begin{array}{@{}l@{}} \forall u, f.; \forall c \in \mathbb{C}, r, \mathit{cmvbr} \in \mathbb{R}, t \in \operatorname{ccint}(0, (2 \cdot \pi)). \\ \quad \operatorname{holomorphic\text{-}on}(u, f) \wedge \\ \quad 0 < r \wedge \\ \quad r < \mathit{cmvbr} \wedge \\ \quad \operatorname{ball}(\operatorname{cc\text{-}ms}, c, \mathit{cmvbr}) \subseteq u \wedge \\ \quad \forall y \in \operatorname{trace}(\operatorname{circle\text{-}path}(c, r), 0, (2 \cdot \pi)). \\ \quad \quad \operatorname{magnitude}(f(y)) \leq \operatorname{magnitude}(f(c)) \Rightarrow \\ \quad \quad f((c + (r \cdot \exp (+i \cdot t)))) = f(c) \end{array}}}
\vnbprop{maximum-modulus-disc}{\vnbstmt{\begin{array}{@{}l@{}} \forall u, f.; \forall c \in \mathbb{C}, r \in \mathbb{R}, z \in \operatorname{ball}(\operatorname{cc\text{-}ms}, c, r). \\ \quad \operatorname{holomorphic\text{-}on}(u, f) \wedge \\ \quad 0 < r \wedge \\ \quad \operatorname{ball}(\operatorname{cc\text{-}ms}, c, r) \subseteq u \wedge \\ \quad \forall z \in \operatorname{ball}(\operatorname{cc\text{-}ms}, c, r).\; (\operatorname{magnitude}(f(z)) \leq \operatorname{magnitude}(f(c))) \Rightarrow \\ \quad \quad f(z) = f(c) \end{array}}}
\vnbprop{holomorphic-mean-value-bound}{\vnbstmt{\begin{array}{@{}l@{}} \forall u, f.; \forall c \in \mathbb{C}, r, \mathit{cmvbr}, m \in \mathbb{R}. \\ \quad \operatorname{holomorphic\text{-}on}(u, f) \wedge \\ \quad 0 < r \wedge \\ \quad r < \mathit{cmvbr} \wedge \\ \quad \operatorname{ball}(\operatorname{cc\text{-}ms}, c, \mathit{cmvbr}) \subseteq u \wedge \\ \quad \forall y \in \operatorname{trace}(\operatorname{circle\text{-}path}(c, r), 0, (2 \cdot \pi)).\; (\operatorname{magnitude}(f(y)) \leq m) \Rightarrow \\ \quad \quad \operatorname{magnitude}(f(c)) \leq m \end{array}}}

\vnbbookfile{theorem-library/mvp-max}{theorem-library/mvp-max}
\noindent{\small mvp-max.scm -- WAVE 2 (2026-10-27..30), wired 2026-10-30 from scratchpad/calint2/mvp-max.scm. Library: complex-analysis.}\par
\vnbprop{mvp-circle-max-constant}{\vnbstmt{\begin{array}{@{}l@{}} \forall u, f.; \forall c \in \mathbb{C}, r, \mathit{cmvbr} \in \mathbb{R}, t \in \operatorname{ccint}(0, (2 \cdot \pi)). \\ \quad \operatorname{has\text{-}mean\text{-}value\text{-}property}(u, f) \wedge \\ \quad 0 < r \wedge \\ \quad r < \mathit{cmvbr} \wedge \\ \quad \operatorname{ball}(\operatorname{cc\text{-}ms}, c, \mathit{cmvbr}) \subseteq u \wedge \\ \quad \forall y \in \operatorname{trace}(\operatorname{circle\text{-}path}(c, r), 0, (2 \cdot \pi)). \\ \quad \quad \operatorname{magnitude}(f(y)) \leq \operatorname{magnitude}(f(c)) \Rightarrow \\ \quad \quad f((c + (r \cdot \exp (+i \cdot t)))) = f(c) \end{array}}}
\vnbprop{mvp-disc-max-constant}{\vnbstmt{\begin{array}{@{}l@{}} \forall u, f.; \forall c \in \mathbb{C}, r \in \mathbb{R}, z \in \operatorname{ball}(\operatorname{cc\text{-}ms}, c, r). \\ \quad \operatorname{has\text{-}mean\text{-}value\text{-}property}(u, f) \wedge \\ \quad 0 < r \wedge \\ \quad \operatorname{ball}(\operatorname{cc\text{-}ms}, c, r) \subseteq u \wedge \\ \quad \forall z \in \operatorname{ball}(\operatorname{cc\text{-}ms}, c, r).\; (\operatorname{magnitude}(f(z)) \leq \operatorname{magnitude}(f(c))) \Rightarrow \\ \quad \quad f(z) = f(c) \end{array}}}

\vnbbookfile{theorem-library/mvp-modulus}{theorem-library/mvp-modulus}
\noindent{\small mvp-modulus.scm -- WAVE 2 (2026-10-27..30), wired 2026-10-30 from scratchpad/calint2/mvp-modulus.scm. Library: complex-analysis.}\par
\vnbprop{connected-clopen}{\vnbstmt{\begin{array}{@{}l@{}} \forall s, a, b.; \forall p \in b, q \in a. \\ \quad \operatorname{is\text{-}connected}(s, a) \wedge \\ \quad b \subseteq a \wedge \\ \quad \forall x \in b. \exists r. \\ \quad \quad \operatorname{pos\text{-}rr}(r) \wedge \forall y \in a.\; ((y \in \operatorname{ball}(s, x, r)) \Rightarrow (y \in b)) \wedge \\ \quad \forall x \in a. \\ \quad \quad \neg (x \in b) \Rightarrow  \\ \quad \quad \quad \exists r. \\ \quad \quad \quad \quad \operatorname{pos\text{-}rr}(r) \wedge \forall y \in a.\; ((y \in \operatorname{ball}(s, x, r)) \Rightarrow \neg (y \in b)) \Rightarrow \\ \quad \quad q \in b \end{array}}}
\vnbprop{mvp-maximum-modulus}{\vnbstmt{\begin{array}{@{}l@{}} \forall u, f.; \forall a, z \in u. \\ \quad \operatorname{has\text{-}mean\text{-}value\text{-}property}(u, f) \wedge \\ \quad \operatorname{is\text{-}connected}(\operatorname{cc\text{-}ms}, u) \wedge \\ \quad \forall z \in u.\; (\operatorname{magnitude}(f(z)) \leq \operatorname{magnitude}(f(a))) \Rightarrow \\ \quad \quad f(z) = f(a) \end{array}}}

\vnbbookfile{theorem-library/morera}{theorem-library/morera}
\noindent{\small morera.scm -- MORERA'S THEOREM AND LOCALLY UNIFORM LIMITS OF HOLOMORPHIC FUNCTIONS (complex-analysis.pdf 3.6: Proposition 3.43, Proposition 3.44), both CONDITIONAL on "the derivative of a holomorphic function is holomorphic" (Corollary 3.37). Agent CA-4 of week 2 of the October roadmap (scratchpad/triage/WEEK2.md), ...}\par
\vnbprop{diff-on-local}{\vnbstmt{\begin{array}{@{}l@{}} \forall u.; \forall f \in (u \to \mathbb{C}), v.; \forall a \in v, l. \\ \quad \operatorname{is\text{-}open}(\operatorname{nf\text{-}metric\text{-}space}(\operatorname{cc\text{-}normed\text{-}field}), u) \wedge \\ \quad \operatorname{is\text{-}open}(\operatorname{nf\text{-}metric\text{-}space}(\operatorname{cc\text{-}normed\text{-}field}), v) \wedge \\ \quad v \subseteq u \wedge \\ \quad \operatorname{is\text{-}diff\text{-}on}(\operatorname{cc\text{-}normed\text{-}field}, v, \operatorname{restrict}(f, v), a, l) \Rightarrow \\ \quad \quad \operatorname{is\text{-}diff\text{-}on}(\operatorname{cc\text{-}normed\text{-}field}, u, f, a, l) \end{array}}}
\vnbprop{cauchy-convex-segments-of-triangles}{\vnbstmt{\begin{array}{@{}l@{}} \forall u, f.; \forall a, b, c \in u. \\ \quad \operatorname{is\text{-}convex\text{-}cc}(u) \wedge \\ \quad \operatorname{is\text{-}open}(\operatorname{nf\text{-}metric\text{-}space}(\operatorname{cc\text{-}normed\text{-}field}), u) \wedge \\ \quad f \in (u \to \mathbb{C}) \wedge \\ \quad \forall y \in u. \\ \quad \quad \operatorname{is\text{-}continuous\text{-}at}(\operatorname{subspace\text{-}ms}(\operatorname{nf\text{-}metric\text{-}space}(\operatorname{cc\text{-}normed\text{-}field}), u), \\ \quad \quad \quad \operatorname{nf\text{-}metric\text{-}space}(\operatorname{cc\text{-}normed\text{-}field}), f, y) \wedge \\ \quad \forall \mathit{mrta}, \mathit{mrtb}, \mathit{mrtc} \in u. \\ \quad \quad (\operatorname{conv3}(\mathit{mrta}, \mathit{mrtb}, \mathit{mrtc}) \subseteq u) \Rightarrow (\operatorname{tri\text{-}int}(f, \mathit{mrta}, \mathit{mrtb}, \mathit{mrtc}) = 0) \Rightarrow \\ \quad \quad \operatorname{seg\text{-}int}(f, a, c) = (\operatorname{seg\text{-}int}(f, a, b) + \operatorname{seg\text{-}int}(f, b, c)) \end{array}}}
\vnbprop{cauchy-convex-primitive-of-triangles}{\vnbstmt{\begin{array}{@{}l@{}} \forall u, f.; \forall o \in u. \\ \quad \operatorname{is\text{-}convex\text{-}cc}(u) \wedge \\ \quad \operatorname{is\text{-}open}(\operatorname{nf\text{-}metric\text{-}space}(\operatorname{cc\text{-}normed\text{-}field}), u) \wedge \\ \quad f \in (u \to \mathbb{C}) \wedge \\ \quad \forall y \in u. \\ \quad \quad \operatorname{is\text{-}continuous\text{-}at}(\operatorname{subspace\text{-}ms}(\operatorname{nf\text{-}metric\text{-}space}(\operatorname{cc\text{-}normed\text{-}field}), u), \\ \quad \quad \quad \operatorname{nf\text{-}metric\text{-}space}(\operatorname{cc\text{-}normed\text{-}field}), f, y) \wedge \\ \quad \forall a, b, c \in u. \\ \quad \quad (\operatorname{conv3}(a, b, c) \subseteq u) \Rightarrow (\operatorname{tri\text{-}int}(f, a, b, c) = 0) \Rightarrow \\ \quad \quad \exists g \in (u \to \mathbb{C}). \\ \quad \quad \quad (\operatorname{holomorphic\text{-}on}(u, g) \wedge \forall z \in u.\; \operatorname{is\text{-}diff\text{-}on}(\operatorname{cc\text{-}normed\text{-}field}, u, g, z, f(z))) \end{array}}}
\vnbprop{morera-conditional}{\vnbstmt{\begin{array}{@{}l@{}} \forall u.; \forall f \in (u \to \mathbb{C}). \\ \quad \forall v, g, d. \\ \quad \quad \operatorname{holomorphic\text{-}on}(v, g) \wedge \\ \quad \quad d \in (v \to \mathbb{C}) \wedge \\ \quad \quad \forall z \in v.\; \operatorname{is\text{-}diff\text{-}on}(\operatorname{cc\text{-}normed\text{-}field}, v, g, z, d(z)) \Rightarrow \\ \quad \quad \quad \operatorname{holomorphic\text{-}on}(v, d) \wedge \\ \quad \operatorname{is\text{-}open}(\operatorname{nf\text{-}metric\text{-}space}(\operatorname{cc\text{-}normed\text{-}field}), u) \wedge \\ \quad \forall y \in u. \\ \quad \quad \operatorname{is\text{-}continuous\text{-}at}(\operatorname{subspace\text{-}ms}(\operatorname{nf\text{-}metric\text{-}space}(\operatorname{cc\text{-}normed\text{-}field}), u), \\ \quad \quad \quad \operatorname{nf\text{-}metric\text{-}space}(\operatorname{cc\text{-}normed\text{-}field}), f, y) \wedge \\ \quad \forall a, b, c \in u. \\ \quad \quad (\operatorname{conv3}(a, b, c) \subseteq u) \Rightarrow (\operatorname{tri\text{-}int}(f, a, b, c) = 0) \Rightarrow \\ \quad \quad \operatorname{holomorphic\text{-}on}(u, f) \end{array}}}
\vnbprop{holomorphic-uniform-limit-conditional}{\vnbstmt{\begin{array}{@{}l@{}} \forall u, m.; \forall f \in (u \to \mathbb{C}). \\ \quad \forall v, g, d. \\ \quad \quad \operatorname{holomorphic\text{-}on}(v, g) \wedge \\ \quad \quad d \in (v \to \mathbb{C}) \wedge \\ \quad \quad \forall z \in v.\; \operatorname{is\text{-}diff\text{-}on}(\operatorname{cc\text{-}normed\text{-}field}, v, g, z, d(z)) \Rightarrow \\ \quad \quad \quad \operatorname{holomorphic\text{-}on}(v, d) \wedge \\ \quad \operatorname{is\text{-}open}(\operatorname{nf\text{-}metric\text{-}space}(\operatorname{cc\text{-}normed\text{-}field}), u) \wedge \\ \quad \forall k \in \mathbb{N}.\; \operatorname{holomorphic\text{-}on}(u, m(k)) \wedge \\ \quad \forall \mathit{mrlz} \in u. \exists r. \\ \quad \quad \operatorname{pos\text{-}rr}(r) \wedge \\ \quad \quad \operatorname{ball}(\operatorname{cc\text{-}ms}, \mathit{mrlz}, r) \subseteq u \wedge \\ \quad \quad \forall s. \\ \quad \quad \quad \operatorname{pos\text{-}rr}(s) \Rightarrow  \\ \quad \quad \quad \quad \exists p \in \mathbb{N}. \forall k \in \mathbb{N}, y \in \operatorname{ball}(\operatorname{cc\text{-}ms}, \mathit{mrlz}, r). \\ \quad \quad \quad \quad \quad (p \leq k) \Rightarrow (\operatorname{magnitude}((m(k)(y) - f(y))) \leq s) \Rightarrow \\ \quad \quad \operatorname{holomorphic\text{-}on}(u, f) \end{array}}}
\vnbprop{deriv-holomorphic-bridge}{\vnbstmt{\begin{array}{@{}l@{}} \forall v, g, d. \\ \quad \forall u, f. \\ \quad \quad \operatorname{holomorphic\text{-}on}(u, f) \Rightarrow  \\ \quad \quad \quad \operatorname{holomorphic\text{-}on}(u, (\lambda w \in u.\; \operatorname{deriv\text{-}on}(\operatorname{cc\text{-}normed\text{-}field}, u, f, w))) \wedge \\ \quad \operatorname{holomorphic\text{-}on}(v, g) \wedge \\ \quad d \in (v \to \mathbb{C}) \wedge \\ \quad \forall z \in v.\; \operatorname{is\text{-}diff\text{-}on}(\operatorname{cc\text{-}normed\text{-}field}, v, g, z, d(z)) \Rightarrow \\ \quad \quad \operatorname{holomorphic\text{-}on}(v, d) \end{array}}}
\vnbprop{morera-of-deriv-holomorphic}{\vnbstmt{\begin{array}{@{}l@{}} \forall u.; \forall f \in (u \to \mathbb{C}). \\ \quad \forall u, f. \\ \quad \quad \operatorname{holomorphic\text{-}on}(u, f) \Rightarrow  \\ \quad \quad \quad \operatorname{holomorphic\text{-}on}(u, (\lambda w \in u.\; \operatorname{deriv\text{-}on}(\operatorname{cc\text{-}normed\text{-}field}, u, f, w))) \wedge \\ \quad \operatorname{is\text{-}open}(\operatorname{nf\text{-}metric\text{-}space}(\operatorname{cc\text{-}normed\text{-}field}), u) \wedge \\ \quad \forall y \in u. \\ \quad \quad \operatorname{is\text{-}continuous\text{-}at}(\operatorname{subspace\text{-}ms}(\operatorname{nf\text{-}metric\text{-}space}(\operatorname{cc\text{-}normed\text{-}field}), u), \\ \quad \quad \quad \operatorname{nf\text{-}metric\text{-}space}(\operatorname{cc\text{-}normed\text{-}field}), f, y) \wedge \\ \quad \forall a, b, c \in u. \\ \quad \quad (\operatorname{conv3}(a, b, c) \subseteq u) \Rightarrow (\operatorname{tri\text{-}int}(f, a, b, c) = 0) \Rightarrow \\ \quad \quad \operatorname{holomorphic\text{-}on}(u, f) \end{array}}}
\vnbprop{holomorphic-uniform-limit-of-deriv-holomorphic}{\vnbstmt{\begin{array}{@{}l@{}} \forall u, m.; \forall f \in (u \to \mathbb{C}). \\ \quad \forall u, f. \\ \quad \quad \operatorname{holomorphic\text{-}on}(u, f) \Rightarrow  \\ \quad \quad \quad \operatorname{holomorphic\text{-}on}(u, (\lambda w \in u.\; \operatorname{deriv\text{-}on}(\operatorname{cc\text{-}normed\text{-}field}, u, f, w))) \wedge \\ \quad \operatorname{is\text{-}open}(\operatorname{nf\text{-}metric\text{-}space}(\operatorname{cc\text{-}normed\text{-}field}), u) \wedge \\ \quad \forall k \in \mathbb{N}.\; \operatorname{holomorphic\text{-}on}(u, m(k)) \wedge \\ \quad \forall z \in u. \exists r. \\ \quad \quad \operatorname{pos\text{-}rr}(r) \wedge \\ \quad \quad \operatorname{ball}(\operatorname{cc\text{-}ms}, z, r) \subseteq u \wedge \\ \quad \quad \forall s. \\ \quad \quad \quad \operatorname{pos\text{-}rr}(s) \Rightarrow  \\ \quad \quad \quad \quad \exists p \in \mathbb{N}. \forall k \in \mathbb{N}, y \in \operatorname{ball}(\operatorname{cc\text{-}ms}, z, r). \\ \quad \quad \quad \quad \quad (p \leq k) \Rightarrow (\operatorname{magnitude}((m(k)(y) - f(y))) \leq s) \Rightarrow \\ \quad \quad \operatorname{holomorphic\text{-}on}(u, f) \end{array}}}
\vnbprop{morera}{\vnbstmt{\begin{array}{@{}l@{}} \forall u.; \forall f \in (u \to \mathbb{C}). \\ \quad \operatorname{is\text{-}open}(\operatorname{nf\text{-}metric\text{-}space}(\operatorname{cc\text{-}normed\text{-}field}), u) \wedge \\ \quad \forall y \in u. \\ \quad \quad \operatorname{is\text{-}continuous\text{-}at}(\operatorname{subspace\text{-}ms}(\operatorname{nf\text{-}metric\text{-}space}(\operatorname{cc\text{-}normed\text{-}field}), u), \\ \quad \quad \quad \operatorname{nf\text{-}metric\text{-}space}(\operatorname{cc\text{-}normed\text{-}field}), f, y) \wedge \\ \quad \forall a, b, c \in u. \\ \quad \quad (\operatorname{conv3}(a, b, c) \subseteq u) \Rightarrow (\operatorname{tri\text{-}int}(f, a, b, c) = 0) \Rightarrow \\ \quad \quad \operatorname{holomorphic\text{-}on}(u, f) \end{array}}}
\vnbprop{holomorphic-uniform-limit}{\vnbstmt{\begin{array}{@{}l@{}} \forall u, m.; \forall f \in (u \to \mathbb{C}). \\ \quad \operatorname{is\text{-}open}(\operatorname{nf\text{-}metric\text{-}space}(\operatorname{cc\text{-}normed\text{-}field}), u) \wedge \\ \quad \forall k \in \mathbb{N}.\; \operatorname{holomorphic\text{-}on}(u, m(k)) \wedge \\ \quad \forall z \in u. \exists r. \\ \quad \quad \operatorname{pos\text{-}rr}(r) \wedge \\ \quad \quad \operatorname{ball}(\operatorname{cc\text{-}ms}, z, r) \subseteq u \wedge \\ \quad \quad \forall s. \\ \quad \quad \quad \operatorname{pos\text{-}rr}(s) \Rightarrow  \\ \quad \quad \quad \quad \exists p \in \mathbb{N}. \forall k \in \mathbb{N}, y \in \operatorname{ball}(\operatorname{cc\text{-}ms}, z, r). \\ \quad \quad \quad \quad \quad (p \leq k) \Rightarrow (\operatorname{magnitude}((m(k)(y) - f(y))) \leq s) \Rightarrow \\ \quad \quad \operatorname{holomorphic\text{-}on}(u, f) \end{array}}}

\vnbbookfile{structure-library/gamma-function}{structure-library/gamma-function}
\noindent{\small gamma-function.scm -- THE IMPROPER INTEGRAL ON (0, +inf) AND THE GAMMA FUNCTION (complex-analysis.pdf ch. 3, section 3.9: equations (96)-(99), Definition 3.58). DEFINITIONS ONLY; every law is PROVEN in theorem-library/gamma-function.scm. Agent CA-6 of week 3 of the October roadmap (scratchpad/triage/WEEK3.md), ...}\par
\par\medskip\noindent\textbf{Definition} (\texttt{improper-int-converges-to}).\;For all $\operatorname{ipf\_}$ and $\operatorname{ipl\_}$, 
\begin{equation*}
\begin{array}{@{}l@{}} \operatorname{improper\text{-}int\text{-}converges\text{-}to}(\operatorname{ipf\_}, \operatorname{ipl\_}) \\ \Leftrightarrow\; \operatorname{ipl\_} \in \mathbb{C} \wedge \\ \forall \operatorname{ipe\_}. \\ \quad \operatorname{pos\text{-}rr}(\operatorname{ipe\_}) \Rightarrow  \\ \quad \quad \exists \operatorname{ipd\_}. \\ \quad \quad \quad \operatorname{pos\text{-}rr}(\operatorname{ipd\_}) \wedge \\ \quad \quad \quad \exists \operatorname{ipm\_} \in \mathbb{R}. \forall \operatorname{ipr\_}, \operatorname{ipq\_} \in \mathbb{R}. \\ \quad \quad \quad \quad (0 < \operatorname{ipr\_}) \Rightarrow  \\ \quad \quad \quad \quad (\operatorname{ipr\_} < \operatorname{ipd\_}) \Rightarrow  \\ \quad \quad \quad \quad (\operatorname{ipm\_} < \operatorname{ipq\_}) \Rightarrow  \\ \quad \quad \quad \quad (\operatorname{ipr\_} < \operatorname{ipq\_}) \Rightarrow (\operatorname{magnitude}(-(\int_{\operatorname{ipr\_}}^{\operatorname{ipq\_}} \operatorname{ipf\_}, \operatorname{ipl\_})) \leq \operatorname{ipe\_}) \end{array}
\end{equation*}

\par\medskip\noindent\textbf{Definition} (\texttt{improper-int}).\;$\int_0^{\infty} \operatorname{ipf\_} :=$ $(\iota\, \operatorname{ipv\_}.\;\operatorname{improper\text{-}int\text{-}converges\text{-}to}(\operatorname{ipf\_}, \operatorname{ipv\_}))$
\par\medskip\noindent\textbf{Definition} (\texttt{gamma-integrand}).\;$\operatorname{gamma\text{-}integrand}(\operatorname{gmz\_}) :=$
\begin{equation*}
\begin{array}{@{}l@{}} (\lambda \operatorname{gmw\_} \in \{\operatorname{gmu\_} \in \mathbb{C} : (0 < \operatorname{real\text{-}part}(\operatorname{gmu\_}))\}.\; \\ \quad \exp ((\operatorname{gmz\_} - 1) \cdot \log\left(\operatorname{gmw\_}\right)) \cdot \exp -\operatorname{gmw\_}) \end{array}
\end{equation*}

\par\medskip\noindent\textbf{Definition} (\texttt{gamma-function}).\;$\Gamma(\operatorname{gmz\_}) :=$ $\int_0^{\infty} \operatorname{gamma\text{-}integrand}(\operatorname{gmz\_})$

\vnbbookfile{structure-library/gamma-continuation-vocab}{structure-library/gamma-continuation-vocab}
\noindent{\small gamma-continuation-vocab.scm -- WAVE 2 (2026-10-27..30), wired 2026-10-30 from scratchpad/ch3g/ch3g-gamma.scm (vocab). Library: complex-analysis. Proof-free: the definitions (with their notation! lines) moved here from the proof file(s) named above.}\par
\par\medskip\noindent\textbf{Definition} (\texttt{poch}).\;For every $\operatorname{pcz\_}$,  $(\operatorname{poch}(\operatorname{pcz\_}, 0) \simeq 1)$\par
Let $\operatorname{pcz\_}$ be arbitrary. Let $\operatorname{pcn\_} \in \mathbb{N}$. Then $(\operatorname{poch}(\operatorname{pcz\_}, \operatorname{succ}(\operatorname{pcn\_})) \simeq \cdot(\operatorname{poch}(\operatorname{pcz\_}, \operatorname{pcn\_}), +(\operatorname{pcz\_}, \operatorname{pcn\_})))$
\par\medskip\noindent\textbf{Definition} (\texttt{gamma-ext}).\;$\operatorname{gamma\text{-}ext}(\operatorname{gez\_}) :=$
\begin{equation*}
\Gamma((\operatorname{gez\_} + \varepsilon(\{\operatorname{gek\_} \in \mathbb{N} : (0 < (\operatorname{real\text{-}part}(\operatorname{gez\_}) + \operatorname{gek\_}))\}))) \cdot \operatorname{recip}(\operatorname{poch}(\operatorname{gez\_}, \varepsilon(\{\operatorname{gek\_} \in \mathbb{N} : (0 < (\operatorname{real\text{-}part}(\operatorname{gez\_}) + \operatorname{gek\_}))\})))
\end{equation*}


\vnbbookfile{theorem-library/gamma-function}{theorem-library/gamma-function}
\noindent{\small gamma-function.scm -- THE IMPROPER INTEGRAL ON (0, +inf), THE GAMMA FUNCTION, ITS CONVERGENCE ON THE RIGHT HALF-PLANE, ITS HOLOMORPHY (PROPOSITION 3.59), GAMMA(1) = 1 AND GAMMA(z + 1) = z GAMMA(z), GAMMA(n + 1) = n!. (complex-analysis.pdf ch. 3, section 3.9: (96)-(101), Definition 3.58, Proposition 3.59.)}\par
\vnbprop{right-half-plane-open}{\vnbstmt{\operatorname{is\text{-}open}(\operatorname{nf\text{-}metric\text{-}space}(\operatorname{cc\text{-}normed\text{-}field}), \{u \in \mathbb{C} : (0 < \operatorname{real\text{-}part}(u))\})}}
\vnbprop{real-interval-is-road}{\vnbstmt{\begin{array}{@{}l@{}} \forall a, b \in \mathbb{R}. \\ \quad (a < b) \Rightarrow  \\ \quad \quad \operatorname{is\text{-}road}((\lambda t \in \operatorname{ccint}(a, b).\; t), (\lambda t \in \operatorname{ccint}(a, b).\; 1), a, b) \end{array}}}
\vnbprop{real-interval-trace}{\vnbstmt{\begin{array}{@{}l@{}} \forall a, b, u. \\ \quad ((a \in \mathbb{R}) \wedge ((b \in \mathbb{R}) \wedge ((a < b) \wedge (\operatorname{ccint}(a, b) \subseteq u)))) \Rightarrow  \\ \quad \quad \operatorname{trace}((\lambda t \in \operatorname{ccint}(a, b).\; t), a, b) \subseteq u \end{array}}}
\vnbprop{cc-int-as-line-int}{\vnbstmt{\begin{array}{@{}l@{}} \forall u, f, a, b. \\ \quad u \subseteq \mathbb{C} \wedge \\ \quad f \in (u \to \mathbb{C}) \wedge \\ \quad \forall y \in u. \\ \quad \quad \operatorname{is\text{-}continuous\text{-}at}(\operatorname{subspace\text{-}ms}(\operatorname{nf\text{-}metric\text{-}space}(\operatorname{cc\text{-}normed\text{-}field}), u), \\ \quad \quad \quad \operatorname{nf\text{-}metric\text{-}space}(\operatorname{cc\text{-}normed\text{-}field}), f, y) \wedge \\ \quad a \in \mathbb{R} \wedge \\ \quad b \in \mathbb{R} \wedge \\ \quad a < b \wedge \\ \quad \operatorname{ccint}(a, b) \subseteq u \Rightarrow \\ \quad \quad \int_{a}^{b} f = \operatorname{line\text{-}int}(f, (\lambda t \in \operatorname{ccint}(a, b).\; t), (\lambda t \in \operatorname{ccint}(a, b).\; 1), a, b) \end{array}}}
\vnbprop{cc-int-in-cc-of-continuous}{\vnbstmt{\begin{array}{@{}l@{}} \forall u, f, a, b. \\ \quad u \subseteq \mathbb{C} \wedge \\ \quad f \in (u \to \mathbb{C}) \wedge \\ \quad \forall y \in u. \\ \quad \quad \operatorname{is\text{-}continuous\text{-}at}(\operatorname{subspace\text{-}ms}(\operatorname{nf\text{-}metric\text{-}space}(\operatorname{cc\text{-}normed\text{-}field}), u), \\ \quad \quad \quad \operatorname{nf\text{-}metric\text{-}space}(\operatorname{cc\text{-}normed\text{-}field}), f, y) \wedge \\ \quad a \in \mathbb{R} \wedge \\ \quad b \in \mathbb{R} \wedge \\ \quad a < b \wedge \\ \quad \operatorname{ccint}(a, b) \subseteq u \Rightarrow \\ \quad \quad \int_{a}^{b} f \in \mathbb{C} \end{array}}}
\vnbprop{cc-int-split}{\vnbstmt{\begin{array}{@{}l@{}} \forall u, f, a, b.; \forall c \in \mathbb{R}. \\ \quad u \subseteq \mathbb{C} \wedge \\ \quad f \in (u \to \mathbb{C}) \wedge \\ \quad \forall y \in u. \\ \quad \quad \operatorname{is\text{-}continuous\text{-}at}(\operatorname{subspace\text{-}ms}(\operatorname{nf\text{-}metric\text{-}space}(\operatorname{cc\text{-}normed\text{-}field}), u), \\ \quad \quad \quad \operatorname{nf\text{-}metric\text{-}space}(\operatorname{cc\text{-}normed\text{-}field}), f, y) \wedge \\ \quad a \in \mathbb{R} \wedge \\ \quad b \in \mathbb{R} \wedge \\ \quad a < b \wedge \\ \quad \operatorname{ccint}(a, b) \subseteq u \wedge \\ \quad a < c \wedge \\ \quad c < b \Rightarrow \\ \quad \quad \int_{a}^{b} f = (\int_{a}^{c} f + \int_{c}^{b} f) \end{array}}}
\vnbprop{cc-int-compare-primitive}{\vnbstmt{\begin{array}{@{}l@{}} \forall u, f, r, c, a, b. \\ \quad u \subseteq \mathbb{C} \wedge \\ \quad f \in (u \to \mathbb{C}) \wedge \\ \quad \forall y \in u. \\ \quad \quad \operatorname{is\text{-}continuous\text{-}at}(\operatorname{subspace\text{-}ms}(\operatorname{nf\text{-}metric\text{-}space}(\operatorname{cc\text{-}normed\text{-}field}), u), \\ \quad \quad \quad \operatorname{nf\text{-}metric\text{-}space}(\operatorname{cc\text{-}normed\text{-}field}), f, y) \wedge \\ \quad a \in \mathbb{R} \wedge \\ \quad b \in \mathbb{R} \wedge \\ \quad a < b \wedge \\ \quad \operatorname{ccint}(a, b) \subseteq u \wedge \\ \quad \operatorname{holomorphic\text{-}on}(u, c) \wedge \\ \quad r \in (u \to \mathbb{C}) \wedge \\ \quad \forall y \in u. \\ \quad \quad \operatorname{is\text{-}continuous\text{-}at}(\operatorname{subspace\text{-}ms}(\operatorname{nf\text{-}metric\text{-}space}(\operatorname{cc\text{-}normed\text{-}field}), u), \\ \quad \quad \quad \operatorname{nf\text{-}metric\text{-}space}(\operatorname{cc\text{-}normed\text{-}field}), r, y) \wedge \\ \quad \forall z \in u.\; \operatorname{is\text{-}diff\text{-}on}(\operatorname{cc\text{-}normed\text{-}field}, u, c, z, r(z)) \wedge \\ \quad \forall \mathit{gmy} \in \operatorname{ccint}(a, b).\; (\operatorname{magnitude}(f(\mathit{gmy})) \leq \operatorname{real\text{-}part}(r(\mathit{gmy}))) \Rightarrow \\ \quad \quad \operatorname{magnitude}(\int_{a}^{b} f) \leq \operatorname{real\text{-}part}((c(b) - c(a))) \end{array}}}
\vnbprop{right-half-plane-in-slit}{\vnbstmt{\begin{array}{@{}l@{}} \forall y \in \{u \in \mathbb{C} : (0 < \operatorname{real\text{-}part}(u))\}. \\ \quad \neg ((\operatorname{imag\text{-}part}(y) = 0) \wedge (\operatorname{real\text{-}part}(y) \leq 0)) \wedge \\ \quad \neg (y = 0) \wedge \\ \quad y \in \{\mathit{cxsu} \in \mathbb{C} : \neg ((\operatorname{imag\text{-}part}(\mathit{cxsu}) = 0) \wedge (\operatorname{real\text{-}part}(\mathit{cxsu}) \leq 0))\} \end{array}}}
\vnbprop{gamma-integrand-holomorphic}{\vnbstmt{\begin{array}{@{}l@{}} \forall z \in \mathbb{C}. \\ \quad \operatorname{holomorphic\text{-}on}(\{u \in \mathbb{C} : (0 < \operatorname{real\text{-}part}(u))\}, \operatorname{gamma\text{-}integrand}(z)) \end{array}}}
\vnbprop{rr-positive-below-two}{\vnbstmt{\begin{array}{@{}l@{}} \forall d, r \in \mathbb{R}. \\ \quad ((0 < d) \wedge (0 < r)) \Rightarrow \exists x \in \mathbb{R}.\; ((0 < x) \wedge ((x < d) \wedge (x < r))) \end{array}}}
\vnbprop{rr-above-two}{\vnbstmt{\forall m, r \in \mathbb{R}.\; \exists x \in \mathbb{R}.\; ((m < x) \wedge (r < x))}}
\vnbprop{improper-int-cauchy}{\vnbstmt{\begin{array}{@{}l@{}} \forall f. \\ \quad \forall r, q \in \mathbb{R}.\; (((0 < r) \wedge (r < q)) \Rightarrow (\int_{r}^{q} f \in \mathbb{C})) \wedge \\ \quad \forall e. \\ \quad \quad \operatorname{pos\text{-}rr}(e) \Rightarrow  \\ \quad \quad \quad \exists d. \\ \quad \quad \quad \quad \operatorname{pos\text{-}rr}(d) \wedge \\ \quad \quad \quad \quad \exists m \in \mathbb{R}. \forall r_{1}, r_{2}, q_{1}, q_{2}. \\ \quad \quad \quad \quad \quad r_{1} \in \mathbb{R} \wedge \\ \quad \quad \quad \quad \quad r_{2} \in \mathbb{R} \wedge \\ \quad \quad \quad \quad \quad q_{1} \in \mathbb{R} \wedge \\ \quad \quad \quad \quad \quad q_{2} \in \mathbb{R} \wedge \\ \quad \quad \quad \quad \quad 0 < r_{1} \wedge \\ \quad \quad \quad \quad \quad r_{1} < d \wedge \\ \quad \quad \quad \quad \quad 0 < r_{2} \wedge \\ \quad \quad \quad \quad \quad r_{2} < d \wedge \\ \quad \quad \quad \quad \quad m < q_{1} \wedge \\ \quad \quad \quad \quad \quad m < q_{2} \wedge \\ \quad \quad \quad \quad \quad r_{1} < q_{1} \wedge \\ \quad \quad \quad \quad \quad r_{2} < q_{2} \Rightarrow \\ \quad \quad \quad \quad \quad \quad \operatorname{magnitude}((\int_{r_{1}}^{q_{1}} f - \int_{r_{2}}^{q_{2}} f)) \leq e \Rightarrow \\ \quad \quad \exists l.\; \operatorname{improper\text{-}int\text{-}converges\text{-}to}(f, l) \end{array}}}
\vnbprop{improper-int-tail-bound}{\vnbstmt{\begin{array}{@{}l@{}} \forall f, l, r, q, t. \\ \quad \forall \mathit{ipr}, \mathit{ipq} \in \mathbb{R}. \\ \quad \quad ((0 < \mathit{ipr}) \wedge (\mathit{ipr} < \mathit{ipq})) \Rightarrow (\int_{\mathit{ipr}}^{\mathit{ipq}} f \in \mathbb{C}) \wedge \\ \quad \operatorname{improper\text{-}int\text{-}converges\text{-}to}(f, l) \wedge \\ \quad r \in \mathbb{R} \wedge \\ \quad q \in \mathbb{R} \wedge \\ \quad t \in \mathbb{R} \wedge \\ \quad 0 < r \wedge \\ \quad r < q \wedge \\ \quad \forall r_{2}, q_{2} \in \mathbb{R}. \\ \quad \quad ((0 < r_{2}) \wedge ((r_{2} < r) \wedge (q < q_{2}))) \Rightarrow  \\ \quad \quad \quad \operatorname{magnitude}((\int_{r_{2}}^{q_{2}} f - \int_{r}^{q} f)) \leq t \Rightarrow \\ \quad \quad \operatorname{magnitude}((l - \int_{r}^{q} f)) \leq t \end{array}}}
\vnbprop{improper-int-unique}{\vnbstmt{\begin{array}{@{}l@{}} \forall f, l_{1}, l_{2}. \\ \quad \forall r, q \in \mathbb{R}.\; (((0 < r) \wedge (r < q)) \Rightarrow (\int_{r}^{q} f \in \mathbb{C})) \wedge \\ \quad \operatorname{improper\text{-}int\text{-}converges\text{-}to}(f, l_{1}) \wedge \\ \quad \operatorname{improper\text{-}int\text{-}converges\text{-}to}(f, l_{2}) \Rightarrow \\ \quad \quad l_{1} = l_{2} \end{array}}}
\vnbprop{improper-int-value}{\vnbstmt{\begin{array}{@{}l@{}} \forall f, l. \\ \quad \forall r, q \in \mathbb{R}.\; (((0 < r) \wedge (r < q)) \Rightarrow (\int_{r}^{q} f \in \mathbb{C})) \wedge \\ \quad \operatorname{improper\text{-}int\text{-}converges\text{-}to}(f, l) \Rightarrow \\ \quad \quad \int_0^{\infty} f = l \end{array}}}
\vnbprop{cc-exp-magnitude}{\vnbstmt{\forall v \in \mathbb{C}.\; (\operatorname{magnitude}(\exp v) = \exp \operatorname{real\text{-}part}(v))}}
\vnbprop{cc-exp-real-mono}{\vnbstmt{\forall u, v \in \mathbb{R}.\; ((u \leq v) \Rightarrow (\exp u \leq \exp v))}}
\vnbprop{cc-exp-log-real}{\vnbstmt{\forall x \in \mathbb{R}.\; ((0 < x) \Rightarrow (\exp \log\left(x\right) = x))}}
\vnbprop{log-cc-exp-real}{\vnbstmt{\forall x \in \mathbb{R}.\; (\log\left(\exp x\right) = x)}}
\vnbprop{gamma-integrand-magnitude}{\vnbstmt{\begin{array}{@{}l@{}} \forall z \in \mathbb{C}, t \in \mathbb{R}. \\ \quad (0 < t) \Rightarrow  \\ \quad \quad \operatorname{magnitude}(\operatorname{gamma\text{-}integrand}(z)(t)) = \exp (((\operatorname{real\text{-}part}(z) - 1) \cdot \log\left(t\right)) - t) \end{array}}}
\vnbprop{log-nonneg-of-ge-one}{\vnbstmt{\forall t \in \mathbb{R}.\; ((1 \leq t) \Rightarrow (0 \leq \log\left(t\right)))}}
\vnbprop{gamma-integrand-bound-small}{\vnbstmt{\begin{array}{@{}l@{}} \forall z, a, t. \\ \quad ((z \in \mathbb{C}) \wedge ((a \in \mathbb{R}) \wedge ((t \in \mathbb{R}) \wedge ((0 < a) \wedge ((a \leq \operatorname{real\text{-}part}(z)) \wedge ((0 < t) \wedge (t \leq 1))))))) \Rightarrow  \\ \quad \quad \operatorname{magnitude}(\operatorname{gamma\text{-}integrand}(z)(t)) \leq \exp ((a - 1) \cdot \log\left(t\right)) \end{array}}}
\vnbprop{gamma-integrand-bound-large}{\vnbstmt{\begin{array}{@{}l@{}} \forall z, b, a, t. \\ \quad z \in \mathbb{C} \wedge \\ \quad b \in \mathbb{R} \wedge \\ \quad a \in \mathbb{R} \wedge \\ \quad t \in \mathbb{R} \wedge \\ \quad \operatorname{real\text{-}part}(z) \leq b \wedge \\ \quad 1 \leq a \wedge \\ \quad (b - 1) \leq a \wedge \\ \quad 1 \leq t \Rightarrow \\ \quad \quad \operatorname{magnitude}(\operatorname{gamma\text{-}integrand}(z)(t)) \leq \exp ((a \cdot (\log\left((2 \cdot a)\right) - 1)) - (1/2 \cdot t)) \end{array}}}
\vnbprop{ccint-in-right-half-plane}{\vnbstmt{\begin{array}{@{}l@{}} \forall a, b \in \mathbb{R}. \\ \quad (0 < a) \Rightarrow (\operatorname{ccint}(a, b) \subseteq \{u \in \mathbb{C} : (0 < \operatorname{real\text{-}part}(u))\}) \end{array}}}
\vnbprop{gamma-integrand-locint}{\vnbstmt{\begin{array}{@{}l@{}} \forall z \in \mathbb{C}, r, q \in \mathbb{R}. \\ \quad ((0 < r) \wedge (r < q)) \Rightarrow (\int_{r}^{q} \operatorname{gamma\text{-}integrand}(z) \in \mathbb{C}) \end{array}}}
\vnbprop{gamma-tail-small}{\vnbstmt{\begin{array}{@{}l@{}} \forall z, a, b, c. \\ \quad z \in \mathbb{C} \wedge \\ \quad a \in \mathbb{R} \wedge \\ \quad b \in \mathbb{R} \wedge \\ \quad c \in \mathbb{R} \wedge \\ \quad 0 < a \wedge \\ \quad a \leq \operatorname{real\text{-}part}(z) \wedge \\ \quad 0 < b \wedge \\ \quad b < c \wedge \\ \quad c \leq 1 \Rightarrow \\ \quad \quad \operatorname{magnitude}(\int_{b}^{c} \operatorname{gamma\text{-}integrand}(z)) \leq (\exp (a \cdot \log\left(c\right)) \cdot \operatorname{recip}(a)) \end{array}}}
\vnbprop{gamma-tail-large}{\vnbstmt{\begin{array}{@{}l@{}} \forall z, b, a, c, d. \\ \quad z \in \mathbb{C} \wedge \\ \quad b \in \mathbb{R} \wedge \\ \quad a \in \mathbb{R} \wedge \\ \quad c \in \mathbb{R} \wedge \\ \quad d \in \mathbb{R} \wedge \\ \quad \operatorname{real\text{-}part}(z) \leq b \wedge \\ \quad 1 \leq a \wedge \\ \quad (b - 1) \leq a \wedge \\ \quad 1 \leq c \wedge \\ \quad c < d \Rightarrow \\ \quad \quad \operatorname{magnitude}(\int_{c}^{d} \operatorname{gamma\text{-}integrand}(z)) \leq (2 \cdot \exp ((a \cdot (\log\left((2 \cdot a)\right) - 1)) - (1/2 \cdot c))) \end{array}}}
\vnbprop{log-monotone}{\vnbstmt{\forall u, v \in \mathbb{R}.\; (((0 < u) \wedge (u \leq v)) \Rightarrow (\log\left(u\right) \leq \log\left(v\right)))}}
\vnbprop{rr-power-small-at-zero}{\vnbstmt{\begin{array}{@{}l@{}} \forall s, e. \\ \quad ((s \in \mathbb{R}) \wedge ((e \in \mathbb{R}) \wedge ((0 < s) \wedge (0 < e)))) \Rightarrow  \\ \quad \quad \exists d \in \mathbb{R}.\; ((0 < d) \wedge ((d \leq 1) \wedge ((\exp (s \cdot \log\left(d\right)) \cdot \operatorname{recip}(s)) \leq e))) \end{array}}}
\vnbprop{rr-power-monotone}{\vnbstmt{\begin{array}{@{}l@{}} \forall s, r, d. \\ \quad ((s \in \mathbb{R}) \wedge ((r \in \mathbb{R}) \wedge ((d \in \mathbb{R}) \wedge ((0 < s) \wedge ((0 < r) \wedge (r \leq d)))))) \Rightarrow  \\ \quad \quad (\exp (s \cdot \log\left(r\right)) \cdot \operatorname{recip}(s)) \leq (\exp (s \cdot \log\left(d\right)) \cdot \operatorname{recip}(s)) \end{array}}}
\vnbprop{exp-tail-small-at-infinity}{\vnbstmt{\begin{array}{@{}l@{}} \forall k, e. \\ \quad ((k \in \mathbb{R}) \wedge ((e \in \mathbb{R}) \wedge (0 < e))) \Rightarrow  \\ \quad \quad \exists m \in \mathbb{R}.\; ((1 \leq m) \wedge ((2 \cdot \exp (k - (1/2 \cdot m))) \leq e)) \end{array}}}
\vnbprop{exp-tail-monotone}{\vnbstmt{\begin{array}{@{}l@{}} \forall k, m, x. \\ \quad ((k \in \mathbb{R}) \wedge ((m \in \mathbb{R}) \wedge ((x \in \mathbb{R}) \wedge (m \leq x)))) \Rightarrow  \\ \quad \quad (2 \cdot \exp (k - (1/2 \cdot x))) \leq (2 \cdot \exp (k - (1/2 \cdot m))) \end{array}}}
\vnbprop{gamma-small-difference}{\vnbstmt{\begin{array}{@{}l@{}} \forall z, a, d, b, c. \\ \quad z \in \mathbb{C} \wedge \\ \quad a \in \mathbb{R} \wedge \\ \quad d \in \mathbb{R} \wedge \\ \quad b \in \mathbb{R} \wedge \\ \quad c \in \mathbb{R} \wedge \\ \quad 0 < a \wedge \\ \quad a \leq \operatorname{real\text{-}part}(z) \wedge \\ \quad 0 < b \wedge \\ \quad b < d \wedge \\ \quad 0 < c \wedge \\ \quad c < d \wedge \\ \quad d \leq 1 \Rightarrow \\ \quad \quad \operatorname{magnitude}((\int_{b}^{1} \operatorname{gamma\text{-}integrand}(z) - \int_{c}^{1} \operatorname{gamma\text{-}integrand}(z))) \\ \quad \quad \leq\; \exp (a \cdot \log\left(d\right)) \cdot \operatorname{recip}(a) \end{array}}}
\vnbprop{gamma-large-difference}{\vnbstmt{\begin{array}{@{}l@{}} \forall z, b, a, m, c, d. \\ \quad z \in \mathbb{C} \wedge \\ \quad b \in \mathbb{R} \wedge \\ \quad a \in \mathbb{R} \wedge \\ \quad m \in \mathbb{R} \wedge \\ \quad c \in \mathbb{R} \wedge \\ \quad d \in \mathbb{R} \wedge \\ \quad \operatorname{real\text{-}part}(z) \leq b \wedge \\ \quad 1 \leq a \wedge \\ \quad (b - 1) \leq a \wedge \\ \quad 1 \leq m \wedge \\ \quad m < c \wedge \\ \quad m < d \Rightarrow \\ \quad \quad \operatorname{magnitude}((\int_{1}^{c} \operatorname{gamma\text{-}integrand}(z) - \int_{1}^{d} \operatorname{gamma\text{-}integrand}(z))) \\ \quad \quad \leq\; 2 \cdot \exp ((a \cdot (\log\left((2 \cdot a)\right) - 1)) - (1/2 \cdot m)) \end{array}}}
\vnbprop{gamma-converges}{\vnbstmt{\begin{array}{@{}l@{}} \forall z \in \mathbb{C}. \\ \quad (0 < \operatorname{real\text{-}part}(z)) \Rightarrow  \\ \quad \quad \operatorname{improper\text{-}int\text{-}converges\text{-}to}(\operatorname{gamma\text{-}integrand}(z), \Gamma(z)) \end{array}}}
\vnbprop{cc-int-linear-combination}{\vnbstmt{\begin{array}{@{}l@{}} \forall u, a, b, p, q.; \forall c \in \mathbb{C}, h \in (u \to \mathbb{C}). \\ \quad u \subseteq \mathbb{C} \wedge \\ \quad a \in (u \to \mathbb{C}) \wedge \\ \quad \forall y \in u. \\ \quad \quad \operatorname{is\text{-}continuous\text{-}at}(\operatorname{subspace\text{-}ms}(\operatorname{nf\text{-}metric\text{-}space}(\operatorname{cc\text{-}normed\text{-}field}), u), \\ \quad \quad \quad \operatorname{nf\text{-}metric\text{-}space}(\operatorname{cc\text{-}normed\text{-}field}), a, y) \wedge \\ \quad p \in \mathbb{R} \wedge \\ \quad q \in \mathbb{R} \wedge \\ \quad p < q \wedge \\ \quad \operatorname{ccint}(p, q) \subseteq u \wedge \\ \quad b \in (u \to \mathbb{C}) \wedge \\ \quad \forall y \in u. \\ \quad \quad \operatorname{is\text{-}continuous\text{-}at}(\operatorname{subspace\text{-}ms}(\operatorname{nf\text{-}metric\text{-}space}(\operatorname{cc\text{-}normed\text{-}field}), u), \\ \quad \quad \quad \operatorname{nf\text{-}metric\text{-}space}(\operatorname{cc\text{-}normed\text{-}field}), b, y) \wedge \\ \quad \forall v \in u.\; (h(v) \simeq (a(v) + (c \cdot b(v)))) \Rightarrow \\ \quad \quad \forall y \in u. \\ \quad \quad \quad \operatorname{is\text{-}continuous\text{-}at}(\operatorname{subspace\text{-}ms}(\operatorname{nf\text{-}metric\text{-}space}(\operatorname{cc\text{-}normed\text{-}field}), u), \\ \quad \quad \quad \quad \operatorname{nf\text{-}metric\text{-}space}(\operatorname{cc\text{-}normed\text{-}field}), h, y) \wedge \\ \quad \quad \int_{p}^{q} h = (\int_{p}^{q} a + (c \cdot \int_{p}^{q} b)) \end{array}}}
\vnbprop{gamma-integrand-expansion}{\vnbstmt{\begin{array}{@{}l@{}} \forall x, a, t. \\ \quad ((x \in \mathbb{C}) \wedge ((a \in \mathbb{C}) \wedge ((t \in \mathbb{R}) \wedge (0 < t)))) \Rightarrow  \\ \quad \quad ((\operatorname{gamma\text{-}integrand}(x)(t) - \operatorname{gamma\text{-}integrand}(a)(t)) - ((x - a) \cdot (\operatorname{gamma\text{-}integrand}(a)(t) \cdot \log\left(t\right)))) \\ \quad \quad =\; \operatorname{gamma\text{-}integrand}(a)(t) \cdot ((\exp ((x - a) \cdot \log\left(t\right)) - 1) - ((x - a) \cdot \log\left(t\right))) \end{array}}}
\vnbprop{log-abs-bound-on-interval}{\vnbstmt{\begin{array}{@{}l@{}} \forall p, q, t. \\ \quad ((p \in \mathbb{R}) \wedge ((q \in \mathbb{R}) \wedge ((t \in \mathbb{R}) \wedge ((0 < p) \wedge ((p \leq t) \wedge (t \leq q)))))) \Rightarrow  \\ \quad \quad \lvert \log\left(t\right) \rvert \leq (\lvert \log\left(p\right) \rvert + \lvert \log\left(q\right) \rvert) \end{array}}}
\vnbprop{gamma-integrand-bound-interval}{\vnbstmt{\begin{array}{@{}l@{}} \forall a, p, q, t. \\ \quad ((a \in \mathbb{C}) \wedge ((p \in \mathbb{R}) \wedge ((q \in \mathbb{R}) \wedge ((t \in \mathbb{R}) \wedge ((0 < p) \wedge ((p \leq t) \wedge (t \leq q))))))) \Rightarrow  \\ \quad \quad \operatorname{magnitude}(\operatorname{gamma\text{-}integrand}(a)(t)) \leq \exp (\lvert (\operatorname{real\text{-}part}(a) - 1) \rvert \cdot (\lvert \log\left(p\right) \rvert + \lvert \log\left(q\right) \rvert)) \end{array}}}
\vnbprop{gamma-integrand-quadratic}{\vnbstmt{\begin{array}{@{}l@{}} \forall a, x, p, q, t. \\ \quad a \in \mathbb{C} \wedge \\ \quad x \in \mathbb{C} \wedge \\ \quad p \in \mathbb{R} \wedge \\ \quad q \in \mathbb{R} \wedge \\ \quad t \in \mathbb{R} \wedge \\ \quad 0 < p \wedge \\ \quad p \leq t \wedge \\ \quad t \leq q \wedge \\ \quad (\operatorname{magnitude}((x - a)) \cdot (\lvert \log\left(p\right) \rvert + \lvert \log\left(q\right) \rvert)) \leq 1 \Rightarrow \\ \quad \quad \operatorname{magnitude}(((\operatorname{gamma\text{-}integrand}(x)(t) - \operatorname{gamma\text{-}integrand}(a)(t)) - ((x - a) \cdot (\operatorname{gamma\text{-}integrand}(a)(t) \cdot \log\left(t\right))))) \\ \quad \quad \leq\; (\exp (\lvert (\operatorname{real\text{-}part}(a) - 1) \rvert \cdot (\lvert \log\left(p\right) \rvert + \lvert \log\left(q\right) \rvert)) \cdot ((\lvert \log\left(p\right) \rvert + \lvert \log\left(q\right) \rvert) \cdot (\lvert \log\left(p\right) \rvert + \lvert \log\left(q\right) \rvert))) \cdot (\operatorname{magnitude}((x - a)) \cdot \operatorname{magnitude}((x - a))) \end{array}}}
\vnbprop{gamma-truncation-quadratic}{\vnbstmt{\begin{array}{@{}l@{}} \forall a, x, p, q. \\ \quad a \in \mathbb{C} \wedge \\ \quad x \in \mathbb{C} \wedge \\ \quad p \in \mathbb{R} \wedge \\ \quad q \in \mathbb{R} \wedge \\ \quad 0 < p \wedge \\ \quad p < q \wedge \\ \quad (\operatorname{magnitude}((x - a)) \cdot (\lvert \log\left(p\right) \rvert + \lvert \log\left(q\right) \rvert)) \leq 1 \Rightarrow \\ \quad \quad \operatorname{magnitude}(((\int_{p}^{q} \operatorname{gamma\text{-}integrand}(x) - \int_{p}^{q} \operatorname{gamma\text{-}integrand}(a)) - ((x - a) \cdot \int_{p}^{q} (\lambda w \in \{u \in \mathbb{C} : (0 < \operatorname{real\text{-}part}(u))\}.\; ((\exp ((a - 1) \cdot \log\left(w\right)) \cdot \exp -w) \cdot \log\left(w\right)))))) \\ \quad \quad \leq\; (q - p) \cdot ((\exp (\lvert (\operatorname{real\text{-}part}(a) - 1) \rvert \cdot (\lvert \log\left(p\right) \rvert + \lvert \log\left(q\right) \rvert)) \cdot ((\lvert \log\left(p\right) \rvert + \lvert \log\left(q\right) \rvert) \cdot (\lvert \log\left(p\right) \rvert + \lvert \log\left(q\right) \rvert))) \cdot (\operatorname{magnitude}((x - a)) \cdot \operatorname{magnitude}((x - a)))) \end{array}}}
\vnbprop{gamma-truncation-holomorphic}{\vnbstmt{\begin{array}{@{}l@{}} \forall p, q. \\ \quad ((p \in \mathbb{R}) \wedge ((q \in \mathbb{R}) \wedge ((0 < p) \wedge (p < q)))) \Rightarrow  \\ \quad \quad \operatorname{holomorphic\text{-}on}(\{u \in \mathbb{C} : (0 < \operatorname{real\text{-}part}(u))\}, \\ \quad \quad \quad (\lambda x \in \{u \in \mathbb{C} : (0 < \operatorname{real\text{-}part}(u))\}.\; \int_{p}^{q} \operatorname{gamma\text{-}integrand}(x))) \end{array}}}
\vnbprop{gamma-truncation-error}{\vnbstmt{\begin{array}{@{}l@{}} \forall z, b, c, a, l, h. \\ \quad z \in \mathbb{C} \wedge \\ \quad b \in \mathbb{R} \wedge \\ \quad c \in \mathbb{R} \wedge \\ \quad a \in \mathbb{R} \wedge \\ \quad l \in \mathbb{R} \wedge \\ \quad h \in \mathbb{R} \wedge \\ \quad 0 < b \wedge \\ \quad b \leq \operatorname{real\text{-}part}(z) \wedge \\ \quad \operatorname{real\text{-}part}(z) \leq c \wedge \\ \quad 1 \leq a \wedge \\ \quad (c - 1) \leq a \wedge \\ \quad 0 < l \wedge \\ \quad l \leq 1 \wedge \\ \quad 1 \leq h \wedge \\ \quad l < h \Rightarrow \\ \quad \quad \operatorname{magnitude}((\Gamma(z) - \int_{l}^{h} \operatorname{gamma\text{-}integrand}(z))) \\ \quad \quad \leq\; (\exp (b \cdot \log\left(l\right)) \cdot \operatorname{recip}(b)) + (2 \cdot \exp ((a \cdot (\log\left((2 \cdot a)\right) - 1)) - (1/2 \cdot h))) \end{array}}}
\vnbprop{gamma-holomorphic}{\vnbstmt{\begin{array}{@{}l@{}} \operatorname{holomorphic\text{-}on}(\{u \in \mathbb{C} : (0 < \operatorname{real\text{-}part}(u))\}, \\ \quad (\lambda x \in \{u \in \mathbb{C} : (0 < \operatorname{real\text{-}part}(u))\}.\; \Gamma(x))) \end{array}}}
\vnbprop{cc-int-of-derivative}{\vnbstmt{\begin{array}{@{}l@{}} \forall u, r, c, a, b. \\ \quad u \subseteq \mathbb{C} \wedge \\ \quad r \in (u \to \mathbb{C}) \wedge \\ \quad \forall y \in u. \\ \quad \quad \operatorname{is\text{-}continuous\text{-}at}(\operatorname{subspace\text{-}ms}(\operatorname{nf\text{-}metric\text{-}space}(\operatorname{cc\text{-}normed\text{-}field}), u), \\ \quad \quad \quad \operatorname{nf\text{-}metric\text{-}space}(\operatorname{cc\text{-}normed\text{-}field}), r, y) \wedge \\ \quad a \in \mathbb{R} \wedge \\ \quad b \in \mathbb{R} \wedge \\ \quad a < b \wedge \\ \quad \operatorname{ccint}(a, b) \subseteq u \wedge \\ \quad \operatorname{holomorphic\text{-}on}(u, c) \wedge \\ \quad \forall z \in u.\; \operatorname{is\text{-}diff\text{-}on}(\operatorname{cc\text{-}normed\text{-}field}, u, c, z, r(z)) \Rightarrow \\ \quad \quad \int_{a}^{b} r = (c(b) - c(a)) \end{array}}}
\vnbprop{gamma-at-one}{\vnbstmt{\Gamma(1) = 1}}
\vnbprop{gamma-product-derivative}{\vnbstmt{\begin{array}{@{}l@{}} \forall z \in \mathbb{C}, \mathit{hbz} \in \{u \in \mathbb{C} : (0 < \operatorname{real\text{-}part}(u))\}. \\ \quad \operatorname{is\text{-}diff\text{-}on}(\operatorname{cc\text{-}normed\text{-}field}, \{u \in \mathbb{C} : (0 < \operatorname{real\text{-}part}(u))\}, \\ \quad \quad (\lambda w \in \{u \in \mathbb{C} : (0 < \operatorname{real\text{-}part}(u))\}.\; -(\exp (z \cdot \log\left(w\right)) \cdot \exp -w)), \\ \quad \quad \mathit{hbz}, \\ \quad \quad (\lambda w \in \{u \in \mathbb{C} : (0 < \operatorname{real\text{-}part}(u))\}.\; (\operatorname{gamma\text{-}integrand}((z + 1))(w) + (-z \cdot \operatorname{gamma\text{-}integrand}(z)(w))))(\mathit{hbz})) \end{array}}}
\vnbprop{rr-power-small-at-zero-plain}{\vnbstmt{\begin{array}{@{}l@{}} \forall s, e. \\ \quad ((s \in \mathbb{R}) \wedge ((e \in \mathbb{R}) \wedge ((0 < s) \wedge (0 < e)))) \Rightarrow  \\ \quad \quad \exists d \in \mathbb{R}. \\ \quad \quad \quad ((0 < d) \wedge ((d \leq 1) \wedge \forall t \in \mathbb{R}.\; (((0 < t) \wedge (t \leq d)) \Rightarrow (\exp (s \cdot \log\left(t\right)) \leq e)))) \end{array}}}
\vnbprop{gamma-recursion}{\vnbstmt{\begin{array}{@{}l@{}} \forall z \in \mathbb{C}. \\ \quad (0 < \operatorname{real\text{-}part}(z)) \Rightarrow (\Gamma((z + 1)) = (z \cdot \Gamma(z))) \end{array}}}
\vnbprop{gamma-factorial}{\vnbstmt{\forall n \in \mathbb{N}.\; (\Gamma((n + 1)) = \operatorname{factorial}(n))}}

\vnbbookfile{structure-library/ca5-function-spaces}{structure-library/ca5-function-spaces}
\noindent{\small ca5-function-spaces.scm -- THE VOCABULARY OF CHAPTER 5 (FUNCTION SPACES) of the user's complex analysis notes (\textasciitilde{}/docs/complex-analysis.pdf, s.5.1), first part. DEFINITIONS ONLY; the laws are PROVEN in theorem-library/ca5-function-spaces.scm. Agent CA-5, 2026-10-14.  Every definition below is a decision for the user,}\par
\par\medskip\noindent\textbf{Definition} (\texttt{cc-exhaust}).\;$\operatorname{cc\text{-}exhaust}(u, n) :=$
\begin{equation*}
\{\operatorname{cxz\_} \in \mathbb{C} : ((\operatorname{magnitude}(\operatorname{cxz\_}) \leq n) \wedge (\operatorname{ball}(\operatorname{cc\text{-}ms}, \operatorname{cxz\_}, \operatorname{recip}((n + 1))) \subseteq u))\}
\end{equation*}
\par\noindent{\small\textit{Comment.} $u$ is an open subset of $\mathbb{C}$ and $n \in \mathbb{N}$; $K_n = \{ z \in \mathbb{C} : |z| \le n,\ B(z, 1/(n+1)) \subseteq u \}$.}\par

\par\medskip\noindent\textbf{Definition} (\texttt{cc-cont-fns}).\;$\operatorname{cc\text{-}cont\text{-}fns}(u) :=$
\begin{equation*}
\{\operatorname{ccf\_} \in (u \to \mathbb{C}) : \forall \operatorname{ccy\_} \in u.\; \operatorname{is\text{-}continuous\text{-}at}(\operatorname{subspace\text{-}ms}(\operatorname{nf\text{-}metric\text{-}space}(\operatorname{cc\text{-}normed\text{-}field}), u), \operatorname{nf\text{-}metric\text{-}space}(\operatorname{cc\text{-}normed\text{-}field}), \operatorname{ccf\_}, \operatorname{ccy\_})\}
\end{equation*}

\par\medskip\noindent\textbf{Definition} (\texttt{cc-hol-fns}).\;$\operatorname{cc\text{-}hol\text{-}fns}(u) :=$ $\{\operatorname{ccf\_} \in (u \to \mathbb{C}) : \operatorname{holomorphic\text{-}on}(u, \operatorname{ccf\_})\}$
\par\medskip\noindent\textbf{Definition} (\texttt{cc-fn-module}).\;$\operatorname{cc\text{-}fn\text{-}module}(u, v) :=$
\begin{equation*}
{[}\operatorname{normed\text{-}field\text{-}as\text{-}commutative\text{-}ring}(\operatorname{cc\text{-}normed\text{-}field}), v, (\lambda {[}\operatorname{cfa\_}, \operatorname{cfb\_}] \in (v \times v).\; (\lambda \operatorname{cfz\_} \in u.\; (\operatorname{cfa\_}(\operatorname{cfz\_}) + \operatorname{cfb\_}(\operatorname{cfz\_})))), (\lambda \operatorname{cfz\_} \in u.\; 0), (\lambda \operatorname{cfa\_} \in v.\; (\lambda \operatorname{cfz\_} \in u.\; -\operatorname{cfa\_}(\operatorname{cfz\_}))), (\lambda {[}\operatorname{cfc\_}, \operatorname{cfa\_}] \in (\mathbb{C} \times v).\; (\lambda \operatorname{cfz\_} \in u.\; (\operatorname{cfc\_} \cdot \operatorname{cfa\_}(\operatorname{cfz\_}))))]
\end{equation*}

\par\medskip\noindent\textbf{Definition} (\texttt{cc-exhaust-seminorms}).\;$\operatorname{cc\text{-}exhaust\text{-}seminorms}(u, v) :=$
\begin{equation*}
(\lambda \operatorname{cfk\_} \in \mathbb{N}.\; (\lambda \operatorname{cfa\_} \in v.\; \operatorname{cc\text{-}sup\text{-}norm}(\operatorname{cfa\_}, \operatorname{cc\text{-}exhaust}(u, \operatorname{cfk\_}))))
\end{equation*}

\par\medskip\noindent\textbf{Definition} (\texttt{cc-hol-bounded}).\;For all $u$ and $a$, 
\begin{equation*}
\begin{array}{@{}l@{}} \operatorname{cc\text{-}hol\text{-}bounded}(u, a) \\ \Leftrightarrow\; a \subseteq \operatorname{cc\text{-}hol\text{-}fns}(u) \wedge \\ \forall \operatorname{chn\_} \in \mathbb{N}. \exists \operatorname{chm\_} \in \mathbb{R}. \forall \operatorname{chf\_} \in a, \operatorname{chz\_} \in \operatorname{cc\text{-}exhaust}(u, \operatorname{chn\_}). \\ \quad \operatorname{magnitude}(\operatorname{chf\_}(\operatorname{chz\_})) \leq \operatorname{chm\_} \end{array}
\end{equation*}


\vnbbookfile{structure-library/infinite-products-vocab}{structure-library/infinite-products-vocab}
\noindent{\small infinite-products-vocab.scm -- WAVE 2 (2026-10-27..30), wired 2026-10-30 from scratchpad/ch6/ch6-prod.scm (vocab). Library: complex-analysis. Proof-free: the definitions (with their notation! lines) moved here from the proof file(s) named above.}\par
\par\medskip\noindent\textbf{Definition} (\texttt{pprod}).\;For every $\operatorname{ppf\_}$,  $(\operatorname{pprod}(\operatorname{ppf\_}, 0) \simeq 1)$\par
Let $\operatorname{ppf\_}$ be arbitrary. Let $\operatorname{ppn\_} \in \mathbb{N}$. Then $(\operatorname{pprod}(\operatorname{ppf\_}, \operatorname{succ}(\operatorname{ppn\_})) \simeq \cdot(\operatorname{pprod}(\operatorname{ppf\_}, \operatorname{ppn\_}), \operatorname{ppf\_}(\operatorname{ppn\_})))$
\par\medskip\noindent\textbf{Definition} (\texttt{prod-converges-to}).\;For all $\operatorname{pcf\_}$ and $\operatorname{pcl\_}$, 
\begin{equation*}
\begin{array}{@{}l@{}} \operatorname{prod\text{-}converges\text{-}to}(\operatorname{pcf\_}, \operatorname{pcl\_}) \\ \Leftrightarrow\; \operatorname{pcf\_} \in (\mathbb{N} \to \mathbb{C}) \wedge \\ \exists \operatorname{pcm\_} \in \mathbb{N}. \\ \quad \forall \operatorname{pck\_} \in \mathbb{N}.\; ((\operatorname{pcm\_} \leq \operatorname{pck\_}) \Rightarrow \neg (\operatorname{pcf\_}(\operatorname{pck\_}) = 0)) \wedge \\ \quad \exists \operatorname{pcz\_} \in \mathbb{C}. \\ \quad \quad \neg (\operatorname{pcz\_} = 0) \wedge \\ \quad \quad \operatorname{converges\text{-}to}(\operatorname{cc\text{-}ms}, \\ \quad \quad \quad (\lambda \operatorname{pcn\_} \in \mathbb{N}.\; \operatorname{pprod}((\lambda \operatorname{pcj\_} \in \mathbb{N}.\; \operatorname{pcf\_}(+(\operatorname{pcm\_}, \operatorname{pcj\_}))), \operatorname{pcn\_})), \\ \quad \quad \quad \operatorname{pcz\_}) \wedge \\ \quad \quad \operatorname{pcl\_} = \cdot(\operatorname{pprod}(\operatorname{pcf\_}, \operatorname{pcm\_}), \operatorname{pcz\_}) \end{array}
\end{equation*}

\par\medskip\noindent\textbf{Definition} (\texttt{prod-converges}).\;For every $\operatorname{pdf\_}$, 
\begin{equation*}
\operatorname{prod\text{-}converges}(\operatorname{pdf\_}) \Leftrightarrow \exists \operatorname{pdl\_}.\; \operatorname{prod\text{-}converges\text{-}to}(\operatorname{pdf\_}, \operatorname{pdl\_})
\end{equation*}

\par\medskip\noindent\textbf{Definition} (\texttt{prod-abs-converges}).\;For every $\operatorname{paa\_}$, 
\begin{equation*}
\begin{array}{@{}l@{}} \operatorname{prod\text{-}abs\text{-}converges}(\operatorname{paa\_}) \\ \Leftrightarrow\; ((\operatorname{paa\_} \in (\mathbb{N} \to \mathbb{C})) \wedge \operatorname{series\text{-}converges}((\lambda \operatorname{pak\_} \in \mathbb{N}.\; \operatorname{magnitude}(\operatorname{paa\_}(\operatorname{pak\_}))))) \end{array}
\end{equation*}

\par\medskip\noindent\textbf{Definition} (\texttt{prod-normal-on-compacta}).\;For all $\operatorname{pnu\_}$ and $\operatorname{pnf\_}$, 
\begin{equation*}
\begin{array}{@{}l@{}} \operatorname{prod\text{-}normal\text{-}on\text{-}compacta}(\operatorname{pnu\_}, \operatorname{pnf\_}) \\ \Leftrightarrow\; \forall \operatorname{pnn\_} \in \mathbb{N}.\; \operatorname{is\text{-}normally\text{-}convergent}(\operatorname{pnf\_}, \operatorname{cc\text{-}exhaust}(\operatorname{pnu\_}, \operatorname{pnn\_})) \end{array}
\end{equation*}

\par\medskip\noindent\textbf{Definition} (\texttt{prod-value}).\;$\operatorname{prod\text{-}value}(\operatorname{pvh\_}) :=$ $(\iota\, \operatorname{pvl\_}.\;((\operatorname{pvl\_} \in \mathbb{C}) \wedge \operatorname{prod\text{-}converges\text{-}to}(\operatorname{pvh\_}, \operatorname{pvl\_})))$
\par\medskip\noindent\textbf{Definition} (\texttt{prod-fun}).\;$\operatorname{prod\text{-}fun}(\operatorname{pfu\_}, \operatorname{pff\_}) :=$
\begin{equation*}
(\lambda \operatorname{pfx\_} \in \operatorname{pfu\_}.\; \operatorname{prod\text{-}value}((\lambda \operatorname{pfk\_} \in \mathbb{N}.\; (1 + \operatorname{pff\_}(\operatorname{pfk\_})(\operatorname{pfx\_})))))
\end{equation*}


\vnbbookfile{structure-library/ca-ac-vocab}{structure-library/ca-ac-vocab}
\noindent{\small ca-ac-vocab.scm -- UNIFORM CONVERGENCE ON COMPACT SUBSETS of an open set of CC. and ZEROS-ACCUMULATE-AT, an accumulation point of the zeros of a function. DEFINITIONS ONLY; the laws and Proposition 3.44 of complex-analysis.pdf are proven in theorem-library/weierstrass-convergence.scm, Proposition 2.12 and Corollary ...}\par
\par\medskip\noindent\textbf{Definition} (\texttt{zeros-accumulate-at}).\;For all $u$, $f$ and $a$, 
\begin{equation*}
\begin{array}{@{}l@{}} \operatorname{zeros\text{-}accumulate\text{-}at}(u, f, a) \\ \Leftrightarrow\; \forall \operatorname{acr\_}. \\ \quad \operatorname{pos\text{-}rr}(\operatorname{acr\_}) \Rightarrow  \\ \quad \quad \exists \operatorname{acx\_} \in u. \\ \quad \quad \quad (\operatorname{acx\_} \in \operatorname{ball}(\operatorname{cc\text{-}ms}, a, \operatorname{acr\_})) \wedge (\neg (\operatorname{acx\_} = a) \wedge (f(\operatorname{acx\_}) = 0)) \end{array}
\end{equation*}

\par\medskip\noindent\textbf{Definition} (\texttt{unif-on-compacta}).\;For all $u$, $\mathit{fam}$ and $f$, 
\begin{equation*}
\begin{array}{@{}l@{}} \operatorname{unif\text{-}on\text{-}compacta}(u, \mathit{fam}, f) \\ \Leftrightarrow\; \forall \operatorname{wcq\_}. \\ \quad (\operatorname{wcq\_} \subseteq u) \Rightarrow  \\ \quad \operatorname{is\text{-}compact}(\operatorname{subspace\text{-}ms}(\operatorname{cc\text{-}ms}, \operatorname{wcq\_})) \Rightarrow  \\ \quad \quad \forall \operatorname{wce\_}. \\ \quad \quad \quad \operatorname{pos\text{-}rr}(\operatorname{wce\_}) \Rightarrow  \\ \quad \quad \quad \quad \exists \operatorname{wcn\_} \in \mathbb{N}. \forall \operatorname{wck\_} \in \mathbb{N}. \\ \quad \quad \quad \quad \quad (\operatorname{wcn\_} \leq \operatorname{wck\_}) \Rightarrow  \\ \quad \quad \quad \quad \quad \quad \forall \operatorname{wcy\_} \in \operatorname{wcq\_}.\; (\operatorname{magnitude}(-(\operatorname{fam}(\operatorname{wck\_})(\operatorname{wcy\_}), f(\operatorname{wcy\_}))) \leq \operatorname{wce\_}) \end{array}
\end{equation*}


\vnbbookfile{theorem-library/ca5-function-spaces}{theorem-library/ca5-function-spaces}
\noindent{\small ca5-function-spaces.scm -- CHAPTER 5 OF THE COMPLEX ANALYSIS NOTES (FUNCTION SPACES), FIRST PART: the exhaustion of an open subset of CC by compact sets (Lemma 5.1), the sup seminorms on it, and the topology of Proposition 5.2. Agent CA-5, 2026-10-14.  Vocabulary: structure-library/ca5-function-spaces.scm}\par
\vnbprop{cc-exhaust-unfold}{\vnbstmt{\begin{array}{@{}l@{}} \forall u, n. \\ \quad (\operatorname{cc\text{-}exhaust}(u, n) \simeq \{z \in \mathbb{C} : ((\operatorname{magnitude}(z) \leq n) \wedge (\operatorname{ball}(\operatorname{cc\text{-}ms}, z, \operatorname{recip}((n + 1))) \subseteq u))\}) \end{array}}}
\vnbprop{cc-exhaust-membership}{\vnbstmt{\begin{array}{@{}l@{}} \forall u, n, z. \\ \quad (z \in \operatorname{cc\text{-}exhaust}(u, n)) \Leftrightarrow ((z \in \mathbb{C}) \wedge ((\operatorname{magnitude}(z) \leq n) \wedge (\operatorname{ball}(\operatorname{cc\text{-}ms}, z, \operatorname{recip}((n + 1))) \subseteq u))) \end{array}}}
\vnbprop{cc-ms-ball-membership}{\vnbstmt{\begin{array}{@{}l@{}} \forall c \in \mathbb{C}, r \in \mathbb{R}, y. \\ \quad (y \in \operatorname{ball}(\operatorname{cc\text{-}ms}, c, r)) \Leftrightarrow ((y \in \mathbb{C}) \wedge (\operatorname{magnitude}((y - c)) < r)) \end{array}}}
\vnbprop{cc-exhaust-subset}{\vnbstmt{\forall u.; \forall n \in \mathbb{N}.\; (\operatorname{cc\text{-}exhaust}(u, n) \subseteq u)}}
\vnbprop{cc-magnitude-sub-chain}{\vnbstmt{\begin{array}{@{}l@{}} \forall a, b, c \in \mathbb{C}. \\ \quad \operatorname{magnitude}((a - c)) \leq (\operatorname{magnitude}((a - b)) + \operatorname{magnitude}((b - c))) \end{array}}}
\vnbprop{cc-ms-ball-centre}{\vnbstmt{\forall c \in \mathbb{C}, r \in \mathbb{R}.\; ((0 < r) \Rightarrow (c \in \operatorname{ball}(\operatorname{cc\text{-}ms}, c, r)))}}
\vnbprop{cc-exhaust-interior}{\vnbstmt{\begin{array}{@{}l@{}} \forall u.; \forall n \in \mathbb{N}. \\ \quad \operatorname{cc\text{-}exhaust}(u, n) \subseteq \operatorname{interior}(\operatorname{cc\text{-}ms}, \operatorname{cc\text{-}exhaust}(u, (n + 1))) \end{array}}}
\vnbprop{cc-ms-ball-mono}{\vnbstmt{\begin{array}{@{}l@{}} \forall c \in \mathbb{C}, r, q \in \mathbb{R}. \\ \quad (r \leq q) \Rightarrow (\operatorname{ball}(\operatorname{cc\text{-}ms}, c, r) \subseteq \operatorname{ball}(\operatorname{cc\text{-}ms}, c, q)) \end{array}}}
\vnbprop{cc-exhaust-mono}{\vnbstmt{\begin{array}{@{}l@{}} \forall u.; \forall a, b \in \mathbb{N}. \\ \quad (a \leq b) \Rightarrow (\operatorname{cc\text{-}exhaust}(u, a) \subseteq \operatorname{cc\text{-}exhaust}(u, b)) \end{array}}}
\vnbprop{cc-exhaust-covers}{\vnbstmt{\begin{array}{@{}l@{}} \forall u.; \forall z \in u. \\ \quad \operatorname{is\text{-}open}(\operatorname{cc\text{-}ms}, u) \Rightarrow \exists n \in \mathbb{N}.\; (z \in \operatorname{cc\text{-}exhaust}(u, n)) \end{array}}}
\vnbprop{not-subset-witness}{\vnbstmt{\forall a, b.\; (\neg (a \subseteq b) \Rightarrow \exists x \in a.\; \neg (x \in b))}}
\vnbprop{cc-exhaust-exhaustive}{\vnbstmt{\begin{array}{@{}l@{}} \forall u, p. \\ \quad (\operatorname{is\text{-}open}(\operatorname{nf\text{-}metric\text{-}space}(\operatorname{cc\text{-}normed\text{-}field}), u) \wedge ((p \subseteq u) \wedge \operatorname{is\text{-}compact}(\operatorname{subspace\text{-}ms}(\operatorname{cc\text{-}ms}, p)))) \Rightarrow  \\ \quad \quad \exists n \in \mathbb{N}.\; (p \subseteq \operatorname{cc\text{-}exhaust}(u, n)) \end{array}}}
\vnbprop{cc-bolzano-weierstrass}{\vnbstmt{\begin{array}{@{}l@{}} \forall h \in (\mathbb{N} \to \mathbb{C}), d \in \mathbb{R}. \\ \quad \forall i \in \mathbb{N}.\; (\operatorname{magnitude}(h(i)) \leq d) \Rightarrow  \\ \quad \quad \exists a. \\ \quad \quad \quad \operatorname{strictly\text{-}mono\text{-}nn}(a) \wedge \exists l \in \mathbb{C}.\; \operatorname{converges\text{-}to}(\operatorname{cc\text{-}ms}, \operatorname{subseq}(h, a), l) \end{array}}}
\vnbprop{cc-exhaust-limit-in}{\vnbstmt{\begin{array}{@{}l@{}} \forall u.; \forall n \in \mathbb{N}, s \in (\mathbb{N} \to \mathbb{C}), l \in \mathbb{C}. \\ \quad (\forall i \in \mathbb{N}.\; (s(i) \in \operatorname{cc\text{-}exhaust}(u, n)) \wedge \operatorname{converges\text{-}to}(\operatorname{cc\text{-}ms}, s, l)) \Rightarrow  \\ \quad \quad l \in \operatorname{cc\text{-}exhaust}(u, n) \end{array}}}
\vnbprop{cc-exhaust-compact}{\vnbstmt{\forall u.; \forall n \in \mathbb{N}.\; \operatorname{is\text{-}compact}(\operatorname{subspace\text{-}ms}(\operatorname{cc\text{-}ms}, \operatorname{cc\text{-}exhaust}(u, n)))}}
\vnbprop{cc-continuous-bounded-on-exhaust}{\vnbstmt{\begin{array}{@{}l@{}} \forall u.; \forall f \in (u \to \mathbb{C}), n \in \mathbb{N}. \\ \quad \operatorname{is\text{-}open}(\operatorname{nf\text{-}metric\text{-}space}(\operatorname{cc\text{-}normed\text{-}field}), u) \wedge \\ \quad \forall y \in u. \\ \quad \quad \operatorname{is\text{-}continuous\text{-}at}(\operatorname{subspace\text{-}ms}(\operatorname{nf\text{-}metric\text{-}space}(\operatorname{cc\text{-}normed\text{-}field}), u), \\ \quad \quad \quad \operatorname{nf\text{-}metric\text{-}space}(\operatorname{cc\text{-}normed\text{-}field}), f, y) \Rightarrow \\ \quad \quad \exists m \in \mathbb{R}. \\ \quad \quad \quad (0 \leq m) \wedge \forall z \in \operatorname{cc\text{-}exhaust}(u, n).\; (\operatorname{magnitude}(f(z)) \leq m) \end{array}}}
\vnbprop{cc-sup-norm-upper}{\vnbstmt{\begin{array}{@{}l@{}} \forall g, s.; \forall x \in s. \\ \quad \forall y \in s.\; (g(y) \in \mathbb{C}) \Rightarrow (\operatorname{magnitude}(g(x)) \leq \operatorname{cc\text{-}sup\text{-}norm}(g, s)) \end{array}}}
\vnbprop{cc-exhaust-seminorm-real}{\vnbstmt{\begin{array}{@{}l@{}} \forall u.; \forall f \in (u \to \mathbb{C}), n \in \mathbb{N}. \\ \quad \operatorname{is\text{-}open}(\operatorname{nf\text{-}metric\text{-}space}(\operatorname{cc\text{-}normed\text{-}field}), u) \wedge \\ \quad \forall y \in u. \\ \quad \quad \operatorname{is\text{-}continuous\text{-}at}(\operatorname{subspace\text{-}ms}(\operatorname{nf\text{-}metric\text{-}space}(\operatorname{cc\text{-}normed\text{-}field}), u), \\ \quad \quad \quad \operatorname{nf\text{-}metric\text{-}space}(\operatorname{cc\text{-}normed\text{-}field}), f, y) \Rightarrow \\ \quad \quad (\operatorname{cc\text{-}sup\text{-}norm}(f, \operatorname{cc\text{-}exhaust}(u, n)) \in \mathbb{R}) \wedge (0 \leq \operatorname{cc\text{-}sup\text{-}norm}(f, \operatorname{cc\text{-}exhaust}(u, n))) \end{array}}}
\vnbprop{cc-exhaust-uniform-iff-compacta}{\vnbstmt{\begin{array}{@{}l@{}} \forall u, m, f. \\ \quad \operatorname{is\text{-}open}(\operatorname{nf\text{-}metric\text{-}space}(\operatorname{cc\text{-}normed\text{-}field}), u) \Rightarrow  \\ \quad \quad \forall n \in \mathbb{N}, s. \\ \quad \quad \quad \operatorname{pos\text{-}rr}(s) \Rightarrow  \\ \quad \quad \quad \quad \exists p \in \mathbb{N}. \forall k \in \mathbb{N}, z \in \operatorname{cc\text{-}exhaust}(u, n). \\ \quad \quad \quad \quad \quad (p \leq k) \Rightarrow (\operatorname{magnitude}((m(k)(z) - f(z))) \leq s) \\ \quad \quad \Leftrightarrow\; \operatorname{unif\text{-}on\text{-}compacta}(u, m, f) \end{array}}}

\vnbbookfile{theorem-library/removable-singularities}{theorem-library/removable-singularities}
\noindent{\small removable-singularities.scm -- WAVE 2 (2026-10-27..30), wired 2026-10-30 from scratchpad/ch3e/ch3f-rem.scm. Library: complex-analysis.}\par
\vnbprop{cc-cont-at-mag}{\vnbstmt{\begin{array}{@{}l@{}} \forall u, g, a, s. \\ \quad \operatorname{is\text{-}open}(\operatorname{nf\text{-}metric\text{-}space}(\operatorname{cc\text{-}normed\text{-}field}), u) \wedge \\ \quad g \in (u \to \mathbb{C}) \wedge \\ \quad a \in u \wedge \\ \quad \operatorname{is\text{-}continuous\text{-}at}(\operatorname{subspace\text{-}ms}(\operatorname{nf\text{-}metric\text{-}space}(\operatorname{cc\text{-}normed\text{-}field}), u), \\ \quad \quad \operatorname{nf\text{-}metric\text{-}space}(\operatorname{cc\text{-}normed\text{-}field}), g, a) \wedge \\ \quad \operatorname{pos\text{-}rr}(s) \Rightarrow \\ \quad \quad \exists l. \\ \quad \quad \quad (\operatorname{pos\text{-}rr}(l) \wedge \forall z \in u.\; ((\operatorname{magnitude}((z - a)) < l) \Rightarrow (\operatorname{magnitude}((g(z) - g(a))) < s))) \end{array}}}
\vnbprop{cc-cont-at-of-mag}{\vnbstmt{\begin{array}{@{}l@{}} \forall u, g, a. \\ \quad \operatorname{is\text{-}open}(\operatorname{nf\text{-}metric\text{-}space}(\operatorname{cc\text{-}normed\text{-}field}), u) \wedge \\ \quad g \in (u \to \mathbb{C}) \wedge \\ \quad a \in u \wedge \\ \quad \forall s. \\ \quad \quad \operatorname{pos\text{-}rr}(s) \Rightarrow  \\ \quad \quad \quad \exists l. \\ \quad \quad \quad \quad (\operatorname{pos\text{-}rr}(l) \wedge \forall z \in u.\; ((\operatorname{magnitude}((z - a)) < l) \Rightarrow (\operatorname{magnitude}((g(z) - g(a))) < s))) \Rightarrow \\ \quad \quad \operatorname{is\text{-}continuous\text{-}at}(\operatorname{subspace\text{-}ms}(\operatorname{nf\text{-}metric\text{-}space}(\operatorname{cc\text{-}normed\text{-}field}), u), \\ \quad \quad \quad \operatorname{nf\text{-}metric\text{-}space}(\operatorname{cc\text{-}normed\text{-}field}), g, a) \end{array}}}
\vnbprop{removable-singularity}{\vnbstmt{\begin{array}{@{}l@{}} \forall u, c, f, g. \\ \quad \operatorname{is\text{-}open}(\operatorname{nf\text{-}metric\text{-}space}(\operatorname{cc\text{-}normed\text{-}field}), u) \wedge \\ \quad c \in u \wedge \\ \quad \operatorname{holomorphic\text{-}on}(\operatorname{difference}(u, \operatorname{singleton}(c)), f) \wedge \\ \quad g \in (u \to \mathbb{C}) \wedge \\ \quad \forall y \in u. \\ \quad \quad \operatorname{is\text{-}continuous\text{-}at}(\operatorname{subspace\text{-}ms}(\operatorname{nf\text{-}metric\text{-}space}(\operatorname{cc\text{-}normed\text{-}field}), u), \\ \quad \quad \quad \operatorname{nf\text{-}metric\text{-}space}(\operatorname{cc\text{-}normed\text{-}field}), g, y) \wedge \\ \quad \forall z \in \operatorname{difference}(u, \operatorname{singleton}(c)).\; (g(z) = f(z)) \Rightarrow \\ \quad \quad \operatorname{holomorphic\text{-}on}(u, g) \end{array}}}
\vnbprop{riemann-removable}{\vnbstmt{\begin{array}{@{}l@{}} \forall u, c, f. \\ \quad \operatorname{is\text{-}open}(\operatorname{nf\text{-}metric\text{-}space}(\operatorname{cc\text{-}normed\text{-}field}), u) \wedge \\ \quad c \in u \wedge \\ \quad \operatorname{holomorphic\text{-}on}(\operatorname{difference}(u, \operatorname{singleton}(c)), f) \Rightarrow \\ \quad \quad \exists g \in (u \to \mathbb{C}). \\ \quad \quad \quad \forall y \in u. \\ \quad \quad \quad \quad \operatorname{is\text{-}continuous\text{-}at}(\operatorname{subspace\text{-}ms}(\operatorname{nf\text{-}metric\text{-}space}(\operatorname{cc\text{-}normed\text{-}field}), u), \\ \quad \quad \quad \quad \quad \operatorname{nf\text{-}metric\text{-}space}(\operatorname{cc\text{-}normed\text{-}field}), g, y) \wedge \\ \quad \quad \quad \forall z \in \operatorname{difference}(u, \operatorname{singleton}(c)).\; (g(z) = f(z)) \\ \quad \quad \Leftrightarrow\; \exists r. \\ \quad \quad \quad \operatorname{pos\text{-}rr}(r) \wedge \\ \quad \quad \quad \exists m \in \mathbb{R}. \forall z \in \operatorname{difference}(u, \operatorname{singleton}(c)). \\ \quad \quad \quad \quad (\operatorname{magnitude}((z - c)) < r) \Rightarrow (\operatorname{magnitude}(f(z)) \leq m) \end{array}}}
\vnbprop{removable-extension-unique}{\vnbstmt{\begin{array}{@{}l@{}} \forall u, c, g, k. \\ \quad \operatorname{is\text{-}open}(\operatorname{nf\text{-}metric\text{-}space}(\operatorname{cc\text{-}normed\text{-}field}), u) \wedge \\ \quad c \in u \wedge \\ \quad g \in (u \to \mathbb{C}) \wedge \\ \quad k \in (u \to \mathbb{C}) \wedge \\ \quad \operatorname{is\text{-}continuous\text{-}at}(\operatorname{subspace\text{-}ms}(\operatorname{nf\text{-}metric\text{-}space}(\operatorname{cc\text{-}normed\text{-}field}), u), \\ \quad \quad \operatorname{nf\text{-}metric\text{-}space}(\operatorname{cc\text{-}normed\text{-}field}), g, c) \wedge \\ \quad \operatorname{is\text{-}continuous\text{-}at}(\operatorname{subspace\text{-}ms}(\operatorname{nf\text{-}metric\text{-}space}(\operatorname{cc\text{-}normed\text{-}field}), u), \\ \quad \quad \operatorname{nf\text{-}metric\text{-}space}(\operatorname{cc\text{-}normed\text{-}field}), k, c) \wedge \\ \quad \forall z \in \operatorname{difference}(u, \operatorname{singleton}(c)).\; (g(z) = k(z)) \Rightarrow \\ \quad \quad g(c) = k(c) \end{array}}}

\vnbbookfile{theorem-library/power-series-coef-bound}{theorem-library/power-series-coef-bound}
\noindent{\small power-series-coef-bound.scm -- WAVE 2 (2026-10-27..30), wired 2026-10-30 from scratchpad/ch3e/ch3e-coef.scm. Library: complex-analysis.}\par
\vnbprop{power-series-coefficient-bound}{\vnbstmt{\begin{array}{@{}l@{}} \forall q, a, r, m.; \forall k \in \mathbb{N}. \\ \quad q \in (\mathbb{N} \to \mathbb{C}) \wedge \\ \quad a \in \mathbb{C} \wedge \\ \quad r \in \mathbb{R} \wedge \\ \quad 0 < r \wedge \\ \quad r < \operatorname{cps\text{-}radius}(q) \wedge \\ \quad m \in \mathbb{R} \wedge \\ \quad \forall y \in \mathbb{C}. \\ \quad \quad (\operatorname{magnitude}((y - a)) = r) \Rightarrow  \\ \quad \quad \quad \operatorname{magnitude}(\operatorname{cc\text{-}series\text{-}limit}((\lambda \mathit{cpsk} \in \mathbb{N}.\; (q(\mathit{cpsk}) \cdot \operatorname{power}((y - a), \mathit{cpsk}))))) \leq m \Rightarrow \\ \quad \quad \operatorname{magnitude}(q(k)) \leq (m \cdot \operatorname{power}(\operatorname{recip}(r), k)) \end{array}}}

\vnbbookfile{theorem-library/ch4a-laurent}{theorem-library/ch4a-laurent}
\noindent{\small ch4a-laurent.scm -- WAVE 2 (2026-10-27..30), wired 2026-10-30 from scratchpad/ch4a/ch4a-laurent.scm. Library: complex-analysis.}\par
\vnbprop{annulus-unfold}{\vnbstmt{\begin{array}{@{}l@{}} \forall a, m, p. \\ \quad (\operatorname{annulus}(a, m, p) \simeq \{z \in \mathbb{C} : ((m < \operatorname{magnitude}((z - a))) \wedge (\operatorname{magnitude}((z - a)) < p))\}) \end{array}}}
\vnbprop{annulus-open}{\vnbstmt{\forall a \in \mathbb{C}, m, p \in \mathbb{R}.\; \operatorname{is\text{-}open}(\operatorname{cc\text{-}ms}, \operatorname{annulus}(a, m, p))}}
\vnbprop{annulus-shift-facts}{\vnbstmt{\begin{array}{@{}l@{}} \forall a \in \mathbb{C}, m, p \in \mathbb{R}, y \in \operatorname{annulus}(a, m, p). \\ \quad (0 < m) \Rightarrow  \\ \quad \quad (\neg ((y - a) = 0) \wedge (((y - a) \in \operatorname{ball}(\operatorname{cc\text{-}ms}, 0, p)) \wedge (\operatorname{recip}((y - a)) \in \operatorname{ball}(\operatorname{cc\text{-}ms}, 0, \operatorname{recip}(m))))) \end{array}}}
\vnbprop{laurent-assembly}{\vnbstmt{\begin{array}{@{}l@{}} \forall a \in \mathbb{C}, m, p \in \mathbb{R}, b, n. \\ \quad 0 < m \wedge \\ \quad m < p \wedge \\ \quad \operatorname{holomorphic\text{-}on}(\operatorname{ball}(\operatorname{cc\text{-}ms}, 0, p), b) \wedge \\ \quad \operatorname{holomorphic\text{-}on}(\operatorname{ball}(\operatorname{cc\text{-}ms}, 0, \operatorname{recip}(m)), n) \Rightarrow \\ \quad \quad \operatorname{holomorphic\text{-}on}(\operatorname{annulus}(a, m, p), \\ \quad \quad \quad (\lambda x \in \operatorname{annulus}(a, m, p).\; (b((x - a)) + (\operatorname{recip}((x - a)) \cdot n(\operatorname{recip}((x - a))))))) \end{array}}}
\vnbprop{laurent-holomorphic}{\vnbstmt{\begin{array}{@{}l@{}} \forall c \in (\mathbb{Z} \to \mathbb{C}), a \in \mathbb{C}, m, p \in \mathbb{R}. \\ \quad 0 < m \wedge \\ \quad m < p \wedge \\ \quad p < \operatorname{cps\text{-}radius}((\lambda k \in \mathbb{N}.\; c(k))) \wedge \\ \quad \operatorname{recip}(m) < \operatorname{cps\text{-}radius}((\lambda k \in \mathbb{N}.\; c(-(k + 1)))) \Rightarrow \\ \quad \quad \operatorname{holomorphic\text{-}on}(\operatorname{annulus}(a, m, p), (\lambda x \in \operatorname{annulus}(a, m, p).\; \operatorname{laurent\text{-}sum}(c, a, x))) \end{array}}}

\vnbbookfile{theorem-library/weierstrass-convergence}{theorem-library/weierstrass-convergence}
\noindent{\small weierstrass-convergence.scm -- UNIFORM CONVERGENCE ON COMPACT SUBSETS AND PROPOSITION 3.44 OF complex-analysis.pdf (with Weierstrass's complement: the derivatives converge too).  Agent CA-AC, 2026-10-14. Vocabulary: structure-library/ca-ac-vocab.scm (UNIF-ON-COMPACTA, new; the user's to confirm).}\par
\vnbprop{cc-closed-ball-mem}{\vnbstmt{\begin{array}{@{}l@{}} \forall c \in \mathbb{C}, r.; \forall z \in \operatorname{closed\text{-}ball}(\operatorname{cc\text{-}ms}, c, r). \\ \quad (z \in \mathbb{C}) \wedge (\operatorname{magnitude}((z - c)) \leq r) \end{array}}}
\vnbprop{cc-closed-ball-mem-intro}{\vnbstmt{\begin{array}{@{}l@{}} \forall c \in \mathbb{C}, r.; \forall z \in \mathbb{C}. \\ \quad (\operatorname{magnitude}((z - c)) \leq r) \Rightarrow (z \in \operatorname{closed\text{-}ball}(\operatorname{cc\text{-}ms}, c, r)) \end{array}}}
\vnbprop{cc-parts-bound-of-dist}{\vnbstmt{\begin{array}{@{}l@{}} \forall c \in \mathbb{C}, r \in \mathbb{R}, z \in \mathbb{C}. \\ \quad (\operatorname{magnitude}((z - c)) \leq r) \Rightarrow  \\ \quad \quad ((\lvert \operatorname{real\text{-}part}(z) \rvert \leq (\lvert \operatorname{real\text{-}part}(c) \rvert + r)) \wedge (\lvert \operatorname{imag\text{-}part}(z) \rvert \leq (\lvert \operatorname{imag\text{-}part}(c) \rvert + r))) \end{array}}}
\vnbprop{rr-subseq-converges-abs}{\vnbstmt{\begin{array}{@{}l@{}} \forall b, g, q, e. \\ \quad (\operatorname{strictly\text{-}mono\text{-}nn}(g) \wedge (\operatorname{converges\text{-}to}(\operatorname{rr\text{-}ms}, \operatorname{subseq}(b, g), q) \wedge \operatorname{pos\text{-}rr}(e))) \Rightarrow  \\ \quad \quad \exists n \in \mathbb{N}.\; \forall j \in \mathbb{N}.\; ((n \leq j) \Rightarrow (\lvert (b(g(j)) - q) \rvert \leq e)) \end{array}}}
\vnbprop{cc-closed-ball-compact}{\vnbstmt{\begin{array}{@{}l@{}} \forall c \in \mathbb{C}, r \in \mathbb{R}. \\ \quad \operatorname{is\text{-}compact}(\operatorname{subspace\text{-}ms}(\operatorname{cc\text{-}ms}, \operatorname{closed\text{-}ball}(\operatorname{cc\text{-}ms}, c, r))) \end{array}}}
\vnbprop{unif-on-compacta-locally}{\vnbstmt{\begin{array}{@{}l@{}} \forall u, m, f.; \forall z \in u. \\ \quad (\operatorname{is\text{-}open}(\operatorname{nf\text{-}metric\text{-}space}(\operatorname{cc\text{-}normed\text{-}field}), u) \wedge \operatorname{unif\text{-}on\text{-}compacta}(u, m, f)) \Rightarrow  \\ \quad \quad \exists r. \\ \quad \quad \quad \operatorname{pos\text{-}rr}(r) \wedge \\ \quad \quad \quad \operatorname{ball}(\operatorname{cc\text{-}ms}, z, r) \subseteq u \wedge \\ \quad \quad \quad \forall s. \\ \quad \quad \quad \quad \operatorname{pos\text{-}rr}(s) \Rightarrow  \\ \quad \quad \quad \quad \quad \exists p \in \mathbb{N}. \forall k \in \mathbb{N}, y \in \operatorname{ball}(\operatorname{cc\text{-}ms}, z, r). \\ \quad \quad \quad \quad \quad \quad (p \leq k) \Rightarrow (\operatorname{magnitude}((m(k)(y) - f(y))) \leq s) \end{array}}}
\vnbprop{holomorphic-limit-on-compacta}{\vnbstmt{\begin{array}{@{}l@{}} \forall u, m.; \forall f \in (u \to \mathbb{C}). \\ \quad \operatorname{is\text{-}open}(\operatorname{nf\text{-}metric\text{-}space}(\operatorname{cc\text{-}normed\text{-}field}), u) \wedge \\ \quad \forall k \in \mathbb{N}.\; \operatorname{holomorphic\text{-}on}(u, m(k)) \wedge \\ \quad \operatorname{unif\text{-}on\text{-}compacta}(u, m, f) \Rightarrow \\ \quad \quad \operatorname{holomorphic\text{-}on}(u, f) \end{array}}}
\vnbprop{compact-cover-by-balls}{\vnbstmt{\begin{array}{@{}l@{}} \forall s.; \forall h \in (s \to \mathbb{R}). \\ \quad ((s \subseteq \mathbb{C}) \wedge (\operatorname{is\text{-}compact}(\operatorname{subspace\text{-}ms}(\operatorname{cc\text{-}ms}, s)) \wedge \forall x \in s.\; (0 < h(x)))) \Rightarrow  \\ \quad \quad \exists p. \\ \quad \quad \quad ((p \subseteq s) \wedge ((|p| \in \mathbb{N}) \wedge \forall y \in s.\; \exists x \in p.\; (y \in \operatorname{ball}(\operatorname{cc\text{-}ms}, x, h(x))))) \end{array}}}
\vnbprop{locally-uniform-on-compacta}{\vnbstmt{\begin{array}{@{}l@{}} \forall u, g, l. \\ \quad u \subseteq \mathbb{C} \wedge \\ \quad \forall z \in u. \exists r. \\ \quad \quad \operatorname{pos\text{-}rr}(r) \wedge \\ \quad \quad \operatorname{ball}(\operatorname{cc\text{-}ms}, z, r) \subseteq u \wedge \\ \quad \quad \forall s. \\ \quad \quad \quad \operatorname{pos\text{-}rr}(s) \Rightarrow  \\ \quad \quad \quad \quad \exists p \in \mathbb{N}. \forall k \in \mathbb{N}, y \in \operatorname{ball}(\operatorname{cc\text{-}ms}, z, r). \\ \quad \quad \quad \quad \quad (p \leq k) \Rightarrow (\operatorname{magnitude}((g(k)(y) - l(y))) \leq s) \Rightarrow \\ \quad \quad \operatorname{unif\text{-}on\text{-}compacta}(u, g, l) \end{array}}}
\vnbprop{holomorphic-limit-derivatives}{\vnbstmt{\begin{array}{@{}l@{}} \forall u, m.; \forall f \in (u \to \mathbb{C}). \\ \quad \operatorname{is\text{-}open}(\operatorname{nf\text{-}metric\text{-}space}(\operatorname{cc\text{-}normed\text{-}field}), u) \wedge \\ \quad \forall k \in \mathbb{N}.\; \operatorname{holomorphic\text{-}on}(u, m(k)) \wedge \\ \quad \operatorname{unif\text{-}on\text{-}compacta}(u, m, f) \Rightarrow \\ \quad \quad \operatorname{unif\text{-}on\text{-}compacta}(u, \\ \quad \quad \quad (\lambda \mathit{wcvk} \in \mathbb{N}.\; (\lambda w \in u.\; \operatorname{deriv\text{-}on}(\operatorname{cc\text{-}normed\text{-}field}, u, m(\mathit{wcvk}), w))), \\ \quad \quad \quad (\lambda w \in u.\; \operatorname{deriv\text{-}on}(\operatorname{cc\text{-}normed\text{-}field}, u, f, w))) \end{array}}}

\vnbbookfile{theorem-library/analytic-continuation}{theorem-library/analytic-continuation}
\noindent{\small analytic-continuation.scm -- PROPOSITION 2.12 (PRINCIPLE OF ANALYTIC CONTINUATION) AND COROLLARY 2.13 (FINITELY MANY ZEROS IN A CLOSED BALL) OF complex-analysis.pdf. Agent CA-AC, 2026-10-14.  Vocabulary: structure-library/ca-ac-vocab.scm (ZEROS-ACCUMULATE-AT, new; the user's to confirm).}\par
\vnbprop{zeros-accumulate-value}{\vnbstmt{\begin{array}{@{}l@{}} \forall v.; \forall h \in (v \to \mathbb{C}), a \in v. \\ \quad \operatorname{is\text{-}open}(\operatorname{nf\text{-}metric\text{-}space}(\operatorname{cc\text{-}normed\text{-}field}), v) \wedge \\ \quad \forall y \in v. \\ \quad \quad \operatorname{is\text{-}continuous\text{-}at}(\operatorname{subspace\text{-}ms}(\operatorname{nf\text{-}metric\text{-}space}(\operatorname{cc\text{-}normed\text{-}field}), v), \\ \quad \quad \quad \operatorname{nf\text{-}metric\text{-}space}(\operatorname{cc\text{-}normed\text{-}field}), h, y) \wedge \\ \quad \operatorname{zeros\text{-}accumulate\text{-}at}(v, h, a) \Rightarrow \\ \quad \quad h(a) = 0 \end{array}}}
\vnbprop{cc-series-of-zeros}{\vnbstmt{\begin{array}{@{}l@{}} \forall t \in (\mathbb{N} \to \mathbb{C}), l. \\ \quad (\forall k \in \mathbb{N}.\; (t(k) = 0) \wedge \operatorname{cc\text{-}series\text{-}converges\text{-}to}(t, l)) \Rightarrow (l = 0) \end{array}}}
\vnbprop{cc-power-nonzero}{\vnbstmt{\forall n \in \mathbb{N}, t \in \mathbb{C}.\; (\neg (t = 0) \Rightarrow \neg (\operatorname{power}(t, n) = 0))}}
\vnbprop{holomorphic-zero-factor}{\vnbstmt{\begin{array}{@{}l@{}} \forall n \in \mathbb{N}, v, h.; \forall a \in v. \\ \quad (\operatorname{holomorphic\text{-}on}(v, h) \wedge \operatorname{zeros\text{-}accumulate\text{-}at}(v, h, a)) \Rightarrow  \\ \quad \quad \exists q. \\ \quad \quad \quad \operatorname{holomorphic\text{-}on}(v, q) \wedge \forall x \in v.\; (h(x) = (\operatorname{power}((x - a), n) \cdot q(x))) \end{array}}}
\vnbprop{holomorphic-zero-coefficients}{\vnbstmt{\begin{array}{@{}l@{}} \forall u, f.; \forall c \in \mathbb{C}, s, r \in \mathbb{R}, k \in \mathbb{N}. \\ \quad \operatorname{holomorphic\text{-}on}(u, f) \wedge \\ \quad 0 < s \wedge \\ \quad \operatorname{ball}(\operatorname{cc\text{-}ms}, c, s) \subseteq u \wedge \\ \quad 0 < r \wedge \\ \quad r < s \wedge \\ \quad \operatorname{zeros\text{-}accumulate\text{-}at}(u, f, c) \Rightarrow \\ \quad \quad \operatorname{recip}(((2 \cdot \pi) \cdot +i)) \cdot \operatorname{line\text{-}int}((\lambda w \in \operatorname{trace}(\operatorname{circle\text{-}path}(c, r), 0, (2 \cdot \pi)).\; (\operatorname{restrict}(f, \operatorname{trace}(\operatorname{circle\text{-}path}(c, r), 0, (2 \cdot \pi)))(w) \cdot \operatorname{power}(\operatorname{recip}((w - c)), (k + 1)))), \operatorname{circle\text{-}path}(c, r), \operatorname{circle\text{-}deriv}(c, r), 0, (2 \cdot \pi)) \\ \quad \quad =\; 0 \end{array}}}
\vnbprop{holomorphic-zeros-local}{\vnbstmt{\begin{array}{@{}l@{}} \forall u, f.; \forall c \in u. \\ \quad (\operatorname{holomorphic\text{-}on}(u, f) \wedge \operatorname{zeros\text{-}accumulate\text{-}at}(u, f, c)) \Rightarrow  \\ \quad \quad \exists r. \\ \quad \quad \quad (\operatorname{pos\text{-}rr}(r) \wedge ((\operatorname{ball}(\operatorname{cc\text{-}ms}, c, r) \subseteq u) \wedge \forall z \in \operatorname{ball}(\operatorname{cc\text{-}ms}, c, r).\; (f(z) = 0))) \end{array}}}
\vnbprop{holomorphic-identity-principle}{\vnbstmt{\begin{array}{@{}l@{}} \forall u, f.; \forall a, z \in u. \\ \quad (\operatorname{is\text{-}connected}(\operatorname{cc\text{-}ms}, u) \wedge (\operatorname{holomorphic\text{-}on}(u, f) \wedge \operatorname{zeros\text{-}accumulate\text{-}at}(u, f, a))) \Rightarrow  \\ \quad \quad f(z) = 0 \end{array}}}
\vnbprop{analytic-on-holomorphic}{\vnbstmt{\begin{array}{@{}l@{}} \forall u.; \forall f \in (u \to \mathbb{C}). \\ \quad (\operatorname{is\text{-}open}(\operatorname{cc\text{-}ms}, u) \wedge \forall b \in u.\; f \text{ is analytic at } b) \Rightarrow  \\ \quad \quad \operatorname{holomorphic\text{-}on}(u, f) \end{array}}}
\vnbprop{analytic-continuation}{\vnbstmt{\begin{array}{@{}l@{}} \forall u.; \forall f \in (u \to \mathbb{C}), z \in u. \\ \quad \operatorname{is\text{-}open}(\operatorname{cc\text{-}ms}, u) \wedge \\ \quad \operatorname{is\text{-}connected}(\operatorname{cc\text{-}ms}, u) \wedge \\ \quad \forall b \in u.\; f \text{ is analytic at } b \wedge \\ \quad \exists a \in u.\; \operatorname{zeros\text{-}accumulate\text{-}at}(u, f, a) \Rightarrow \\ \quad \quad f(z) = 0 \end{array}}}
\vnbprop{analytic-zeros-finite-in-ball}{\vnbstmt{\begin{array}{@{}l@{}} \forall u.; \forall f \in (u \to \mathbb{C}), c \in \mathbb{C}, r \in \mathbb{R}. \\ \quad \operatorname{is\text{-}open}(\operatorname{cc\text{-}ms}, u) \wedge \\ \quad \operatorname{is\text{-}connected}(\operatorname{cc\text{-}ms}, u) \wedge \\ \quad \forall b \in u.\; f \text{ is analytic at } b \wedge \\ \quad \exists e \in u.\; \neg (f(e) = 0) \wedge \\ \quad \operatorname{closed\text{-}ball}(\operatorname{cc\text{-}ms}, c, r) \subseteq u \Rightarrow \\ \quad \quad |\{z \in \operatorname{closed\text{-}ball}(\operatorname{cc\text{-}ms}, c, r) : (f(z) = 0)\}| \in \mathbb{N} \end{array}}}

\vnbbookfile{theorem-library/ca5-frechet}{theorem-library/ca5-frechet}
\noindent{\small ca5-frechet.scm -- CHAPTER 5 OF THE COMPLEX ANALYSIS NOTES, continued: C(U) AND H(U) AS FRECHET SPACES (Proposition 5.2) and H(U) CLOSED IN C(U) (Proposition 5.3). Agent CA-5, 2026-10-14.  Vocabulary: structure-library/ca5-function-spaces.scm (CC-CONT-FNS, CC-HOL-FNS, CC-FN-MODULE, CC-EXHAUST-SEMINORMS); the Frechet ...}\par
\vnbprop{cc-scalar-ring-neg}{\vnbstmt{(\operatorname{neg}(\operatorname{normed\text{-}field\text{-}as\text{-}commutative\text{-}ring}(\operatorname{cc\text{-}normed\text{-}field})) \simeq (\lambda x \in \mathbb{C}.\; -x))}}
\vnbprop{cc-cont-fns-unfold}{\vnbstmt{\begin{array}{@{}l@{}} \forall u. \\ \quad (\operatorname{cc\text{-}cont\text{-}fns}(u) \simeq \{f \in (u \to \mathbb{C}) : \forall y \in u.\; \operatorname{is\text{-}continuous\text{-}at}(\operatorname{subspace\text{-}ms}(\operatorname{nf\text{-}metric\text{-}space}(\operatorname{cc\text{-}normed\text{-}field}), u), \operatorname{nf\text{-}metric\text{-}space}(\operatorname{cc\text{-}normed\text{-}field}), f, y)\}) \end{array}}}
\vnbprop{cc-cont-fns-membership}{\vnbstmt{\begin{array}{@{}l@{}} \forall u, h. \\ \quad h \in \operatorname{cc\text{-}cont\text{-}fns}(u) \\ \quad \Leftrightarrow\; h \in (u \to \mathbb{C}) \wedge \\ \quad \forall y \in u. \\ \quad \quad \operatorname{is\text{-}continuous\text{-}at}(\operatorname{subspace\text{-}ms}(\operatorname{nf\text{-}metric\text{-}space}(\operatorname{cc\text{-}normed\text{-}field}), u), \\ \quad \quad \quad \operatorname{nf\text{-}metric\text{-}space}(\operatorname{cc\text{-}normed\text{-}field}), h, y) \end{array}}}
\vnbprop{cc-fn-module-scal}{\vnbstmt{\begin{array}{@{}l@{}} \forall u, v \in \mathbf{Set}. \\ \quad (\operatorname{scal}(\operatorname{cc\text{-}fn\text{-}module}(u, v)) \simeq \operatorname{normed\text{-}field\text{-}as\text{-}commutative\text{-}ring}(\operatorname{cc\text{-}normed\text{-}field})) \end{array}}}
\vnbprop{cc-fn-module-vec}{\vnbstmt{\forall u, v \in \mathbf{Set}.\; (\operatorname{vec}(\operatorname{cc\text{-}fn\text{-}module}(u, v)) \simeq v)}}
\vnbprop{cc-fn-module-vadd}{\vnbstmt{\begin{array}{@{}l@{}} \forall u, v \in \mathbf{Set}. \\ \quad (\operatorname{vadd}(\operatorname{cc\text{-}fn\text{-}module}(u, v)) \simeq (\lambda {[}\operatorname{cfa\_}, \operatorname{cfb\_}] \in (v \times v).\; (\lambda z \in u.\; (\operatorname{cfa\_}(z) + \operatorname{cfb\_}(z))))) \end{array}}}
\vnbprop{cc-fn-module-vzero}{\vnbstmt{\forall u, v \in \mathbf{Set}.\; (\operatorname{vzero}(\operatorname{cc\text{-}fn\text{-}module}(u, v)) \simeq (\lambda z \in u.\; 0))}}
\vnbprop{cc-fn-module-vneg}{\vnbstmt{\begin{array}{@{}l@{}} \forall u, v \in \mathbf{Set}. \\ \quad (\operatorname{vneg}(\operatorname{cc\text{-}fn\text{-}module}(u, v)) \simeq (\lambda a \in v.\; (\lambda z \in u.\; -a(z)))) \end{array}}}
\vnbprop{cc-fn-module-act}{\vnbstmt{\begin{array}{@{}l@{}} \forall u, v \in \mathbf{Set}. \\ \quad (\operatorname{act}(\operatorname{cc\text{-}fn\text{-}module}(u, v)) \simeq (\lambda {[}\operatorname{cfc\_}, \operatorname{cfa\_}] \in (\mathbb{C} \times v).\; (\lambda z \in u.\; (\operatorname{cfc\_} \cdot \operatorname{cfa\_}(z))))) \end{array}}}
\vnbprop{cc-cont-fns-fun}{\vnbstmt{\forall u.; \forall a \in \operatorname{cc\text{-}cont\text{-}fns}(u).\; ((u \subseteq \mathbb{C}) \Rightarrow (a \in (u \to \mathbb{C})))}}
\vnbprop{cc-cont-fns-add}{\vnbstmt{\begin{array}{@{}l@{}} \forall u.; \forall a, b \in \operatorname{cc\text{-}cont\text{-}fns}(u). \\ \quad (u \subseteq \mathbb{C}) \Rightarrow ((\lambda z \in u.\; (a(z) + b(z))) \in \operatorname{cc\text{-}cont\text{-}fns}(u)) \end{array}}}
\vnbprop{cc-cont-fns-zero}{\vnbstmt{\forall u.\; ((u \subseteq \mathbb{C}) \Rightarrow ((\lambda z \in u.\; 0) \in \operatorname{cc\text{-}cont\text{-}fns}(u)))}}
\vnbprop{cc-cont-fns-scale}{\vnbstmt{\begin{array}{@{}l@{}} \forall u.; \forall c \in \mathbb{C}, a \in \operatorname{cc\text{-}cont\text{-}fns}(u). \\ \quad (u \subseteq \mathbb{C}) \Rightarrow ((\lambda z \in u.\; (c \cdot a(z))) \in \operatorname{cc\text{-}cont\text{-}fns}(u)) \end{array}}}
\vnbprop{cc-cont-fns-neg}{\vnbstmt{\begin{array}{@{}l@{}} \forall u.; \forall a \in \operatorname{cc\text{-}cont\text{-}fns}(u). \\ \quad (u \subseteq \mathbb{C}) \Rightarrow ((\lambda z \in u.\; -a(z)) \in \operatorname{cc\text{-}cont\text{-}fns}(u)) \end{array}}}
\vnbprop{cc-cont-fns-is-set}{\vnbstmt{\forall u.\; ((u \subseteq \mathbb{C}) \Rightarrow (\operatorname{cc\text{-}cont\text{-}fns}(u) \in \mathbf{Set}))}}
\vnbprop{cc-fn-vadd-assoc}{\vnbstmt{\begin{array}{@{}l@{}} \forall u.; \forall a, b, c \in \operatorname{cc\text{-}cont\text{-}fns}(u). \\ \quad (u \subseteq \mathbb{C}) \Rightarrow  \\ \quad \quad (\lambda {[}\operatorname{cfa\_}, \operatorname{cfb\_}] \in (\operatorname{cc\text{-}cont\text{-}fns}(u) \times \operatorname{cc\text{-}cont\text{-}fns}(u)).\; (\lambda z \in u.\; (\operatorname{cfa\_}(z) + \operatorname{cfb\_}(z))))((\lambda {[}\operatorname{cfa\_}, \operatorname{cfb\_}] \in (\operatorname{cc\text{-}cont\text{-}fns}(u) \times \operatorname{cc\text{-}cont\text{-}fns}(u)).\; (\lambda z \in u.\; (\operatorname{cfa\_}(z) + \operatorname{cfb\_}(z))))(a, b), \\ \quad \quad \quad c) \\ \quad \quad =\; (\lambda {[}\operatorname{cfa\_}, \operatorname{cfb\_}] \in (\operatorname{cc\text{-}cont\text{-}fns}(u) \times \operatorname{cc\text{-}cont\text{-}fns}(u)).\; (\lambda z \in u.\; (\operatorname{cfa\_}(z) + \operatorname{cfb\_}(z))))(a, \\ \quad \quad \quad (\lambda {[}\operatorname{cfa\_}, \operatorname{cfb\_}] \in (\operatorname{cc\text{-}cont\text{-}fns}(u) \times \operatorname{cc\text{-}cont\text{-}fns}(u)).\; (\lambda z \in u.\; (\operatorname{cfa\_}(z) + \operatorname{cfb\_}(z))))(b, c)) \end{array}}}
\vnbprop{cc-fn-vadd-comm}{\vnbstmt{\begin{array}{@{}l@{}} \forall u.; \forall a, b \in \operatorname{cc\text{-}cont\text{-}fns}(u). \\ \quad (u \subseteq \mathbb{C}) \Rightarrow  \\ \quad \quad (\lambda {[}\operatorname{cfa\_}, \operatorname{cfb\_}] \in (\operatorname{cc\text{-}cont\text{-}fns}(u) \times \operatorname{cc\text{-}cont\text{-}fns}(u)).\; (\lambda z \in u.\; (\operatorname{cfa\_}(z) + \operatorname{cfb\_}(z))))(a, b) \\ \quad \quad =\; (\lambda {[}\operatorname{cfa\_}, \operatorname{cfb\_}] \in (\operatorname{cc\text{-}cont\text{-}fns}(u) \times \operatorname{cc\text{-}cont\text{-}fns}(u)).\; (\lambda z \in u.\; (\operatorname{cfa\_}(z) + \operatorname{cfb\_}(z))))(b, a) \end{array}}}
\vnbprop{cc-fn-vadd-zero-left}{\vnbstmt{\begin{array}{@{}l@{}} \forall u.; \forall a \in \operatorname{cc\text{-}cont\text{-}fns}(u). \\ \quad (u \subseteq \mathbb{C}) \Rightarrow  \\ \quad \quad (\lambda {[}\operatorname{cfa\_}, \operatorname{cfb\_}] \in (\operatorname{cc\text{-}cont\text{-}fns}(u) \times \operatorname{cc\text{-}cont\text{-}fns}(u)).\; (\lambda z \in u.\; (\operatorname{cfa\_}(z) + \operatorname{cfb\_}(z))))((\lambda z \in u.\; 0), a) \\ \quad \quad =\; a \end{array}}}
\vnbprop{cc-fn-vadd-zero-right}{\vnbstmt{\begin{array}{@{}l@{}} \forall u.; \forall a \in \operatorname{cc\text{-}cont\text{-}fns}(u). \\ \quad (u \subseteq \mathbb{C}) \Rightarrow  \\ \quad \quad (\lambda {[}\operatorname{cfa\_}, \operatorname{cfb\_}] \in (\operatorname{cc\text{-}cont\text{-}fns}(u) \times \operatorname{cc\text{-}cont\text{-}fns}(u)).\; (\lambda z \in u.\; (\operatorname{cfa\_}(z) + \operatorname{cfb\_}(z))))(a, (\lambda z \in u.\; 0)) \\ \quad \quad =\; a \end{array}}}
\vnbprop{cc-fn-vadd-neg-left}{\vnbstmt{\begin{array}{@{}l@{}} \forall u.; \forall a \in \operatorname{cc\text{-}cont\text{-}fns}(u). \\ \quad (u \subseteq \mathbb{C}) \Rightarrow  \\ \quad \quad (\lambda {[}\mathit{cfa}, \operatorname{cfb\_}] \in (\operatorname{cc\text{-}cont\text{-}fns}(u) \times \operatorname{cc\text{-}cont\text{-}fns}(u)).\; (\lambda z \in u.\; (\operatorname{cfa}(z) + \operatorname{cfb\_}(z))))((\lambda \mathit{cfa} \in \operatorname{cc\text{-}cont\text{-}fns}(u).\; (\lambda z \in u.\; -\operatorname{cfa}(z)))(a), \\ \quad \quad \quad a) \\ \quad \quad =\; (\lambda z \in u.\; 0) \end{array}}}
\vnbprop{cc-fn-vadd-neg-right}{\vnbstmt{\begin{array}{@{}l@{}} \forall u.; \forall a \in \operatorname{cc\text{-}cont\text{-}fns}(u). \\ \quad (u \subseteq \mathbb{C}) \Rightarrow  \\ \quad \quad (\lambda {[}\mathit{cfa}, \operatorname{cfb\_}] \in (\operatorname{cc\text{-}cont\text{-}fns}(u) \times \operatorname{cc\text{-}cont\text{-}fns}(u)).\; (\lambda z \in u.\; (\operatorname{cfa}(z) + \operatorname{cfb\_}(z))))(a, \\ \quad \quad \quad (\lambda \mathit{cfa} \in \operatorname{cc\text{-}cont\text{-}fns}(u).\; (\lambda z \in u.\; -\operatorname{cfa}(z)))(a)) \\ \quad \quad =\; (\lambda z \in u.\; 0) \end{array}}}
\vnbprop{cc-fn-act-distrib-vec}{\vnbstmt{\begin{array}{@{}l@{}} \forall u.; \forall a, b \in \operatorname{cc\text{-}cont\text{-}fns}(u), r \in \mathbb{C}. \\ \quad (u \subseteq \mathbb{C}) \Rightarrow  \\ \quad \quad (\lambda {[}\operatorname{cfc\_}, \operatorname{cfa\_}] \in (\mathbb{C} \times \operatorname{cc\text{-}cont\text{-}fns}(u)).\; (\lambda z \in u.\; (\operatorname{cfc\_} \cdot \operatorname{cfa\_}(z))))(r, \\ \quad \quad \quad (\lambda {[}\operatorname{cfa\_}, \operatorname{cfb\_}] \in (\operatorname{cc\text{-}cont\text{-}fns}(u) \times \operatorname{cc\text{-}cont\text{-}fns}(u)).\; (\lambda z \in u.\; (\operatorname{cfa\_}(z) + \operatorname{cfb\_}(z))))(a, b)) \\ \quad \quad =\; (\lambda {[}\operatorname{cfa\_}, \operatorname{cfb\_}] \in (\operatorname{cc\text{-}cont\text{-}fns}(u) \times \operatorname{cc\text{-}cont\text{-}fns}(u)).\; (\lambda z \in u.\; (\operatorname{cfa\_}(z) + \operatorname{cfb\_}(z))))((\lambda {[}\operatorname{cfc\_}, \operatorname{cfa\_}] \in (\mathbb{C} \times \operatorname{cc\text{-}cont\text{-}fns}(u)).\; (\lambda z \in u.\; (\operatorname{cfc\_} \cdot \operatorname{cfa\_}(z))))(r, a), \\ \quad \quad \quad (\lambda {[}\operatorname{cfc\_}, \operatorname{cfa\_}] \in (\mathbb{C} \times \operatorname{cc\text{-}cont\text{-}fns}(u)).\; (\lambda z \in u.\; (\operatorname{cfc\_} \cdot \operatorname{cfa\_}(z))))(r, b)) \end{array}}}
\vnbprop{cc-fn-act-distrib-scalar}{\vnbstmt{\begin{array}{@{}l@{}} \forall u.; \forall a \in \operatorname{cc\text{-}cont\text{-}fns}(u), r, s \in \mathbb{C}. \\ \quad (u \subseteq \mathbb{C}) \Rightarrow  \\ \quad \quad (\lambda {[}\operatorname{cfc\_}, \operatorname{cfa\_}] \in (\mathbb{C} \times \operatorname{cc\text{-}cont\text{-}fns}(u)).\; (\lambda z \in u.\; (\operatorname{cfc\_} \cdot \operatorname{cfa\_}(z))))((\lambda {[}\operatorname{x\_}, \operatorname{y\_}] \in (\mathbb{C} \times \mathbb{C}).\; (\operatorname{x\_} + \operatorname{y\_}))(r, s), \\ \quad \quad \quad a) \\ \quad \quad =\; (\lambda {[}\operatorname{cfa\_}, \operatorname{cfb\_}] \in (\operatorname{cc\text{-}cont\text{-}fns}(u) \times \operatorname{cc\text{-}cont\text{-}fns}(u)).\; (\lambda z \in u.\; (\operatorname{cfa\_}(z) + \operatorname{cfb\_}(z))))((\lambda {[}\operatorname{cfc\_}, \operatorname{cfa\_}] \in (\mathbb{C} \times \operatorname{cc\text{-}cont\text{-}fns}(u)).\; (\lambda z \in u.\; (\operatorname{cfc\_} \cdot \operatorname{cfa\_}(z))))(r, a), \\ \quad \quad \quad (\lambda {[}\operatorname{cfc\_}, \operatorname{cfa\_}] \in (\mathbb{C} \times \operatorname{cc\text{-}cont\text{-}fns}(u)).\; (\lambda z \in u.\; (\operatorname{cfc\_} \cdot \operatorname{cfa\_}(z))))(s, a)) \end{array}}}
\vnbprop{cc-fn-act-mul-compat}{\vnbstmt{\begin{array}{@{}l@{}} \forall u.; \forall a \in \operatorname{cc\text{-}cont\text{-}fns}(u), r, s \in \mathbb{C}. \\ \quad (u \subseteq \mathbb{C}) \Rightarrow  \\ \quad \quad (\lambda {[}\operatorname{cfc\_}, \operatorname{cfa\_}] \in (\mathbb{C} \times \operatorname{cc\text{-}cont\text{-}fns}(u)).\; (\lambda z \in u.\; (\operatorname{cfc\_} \cdot \operatorname{cfa\_}(z))))((\lambda {[}\operatorname{x\_}, \operatorname{y\_}] \in (\mathbb{C} \times \mathbb{C}).\; (\operatorname{x\_} \cdot \operatorname{y\_}))(r, s), \\ \quad \quad \quad a) \\ \quad \quad =\; (\lambda {[}\operatorname{cfc\_}, \operatorname{cfa\_}] \in (\mathbb{C} \times \operatorname{cc\text{-}cont\text{-}fns}(u)).\; (\lambda z \in u.\; (\operatorname{cfc\_} \cdot \operatorname{cfa\_}(z))))(r, \\ \quad \quad \quad (\lambda {[}\operatorname{cfc\_}, \operatorname{cfa\_}] \in (\mathbb{C} \times \operatorname{cc\text{-}cont\text{-}fns}(u)).\; (\lambda z \in u.\; (\operatorname{cfc\_} \cdot \operatorname{cfa\_}(z))))(s, a)) \end{array}}}
\vnbprop{cc-fn-act-unital}{\vnbstmt{\begin{array}{@{}l@{}} \forall u.; \forall a \in \operatorname{cc\text{-}cont\text{-}fns}(u). \\ \quad (u \subseteq \mathbb{C}) \Rightarrow  \\ \quad \quad (\lambda {[}\operatorname{cfc\_}, \operatorname{cfa\_}] \in (\mathbb{C} \times \operatorname{cc\text{-}cont\text{-}fns}(u)).\; (\lambda z \in u.\; (\operatorname{cfc\_} \cdot \operatorname{cfa\_}(z))))(1, a) \\ \quad \quad =\; a \end{array}}}
\vnbprop{cc-fn-module-is-module}{\vnbstmt{\forall u.\; ((u \subseteq \mathbb{C}) \Rightarrow \operatorname{is\text{-}module}(\operatorname{cc\text{-}fn\text{-}module}(u, \operatorname{cc\text{-}cont\text{-}fns}(u))))}}
\vnbprop{cc-exhaust-seminorm-subadditive}{\vnbstmt{\begin{array}{@{}l@{}} \forall u.; \forall a, b \in \operatorname{cc\text{-}cont\text{-}fns}(u), n \in \mathbb{N}. \\ \quad \operatorname{is\text{-}open}(\operatorname{nf\text{-}metric\text{-}space}(\operatorname{cc\text{-}normed\text{-}field}), u) \Rightarrow  \\ \quad \quad \operatorname{cc\text{-}sup\text{-}norm}((\lambda z \in u.\; (a(z) + b(z))), \operatorname{cc\text{-}exhaust}(u, n)) \\ \quad \quad \leq\; \operatorname{cc\text{-}sup\text{-}norm}(a, \operatorname{cc\text{-}exhaust}(u, n)) + \operatorname{cc\text{-}sup\text{-}norm}(b, \operatorname{cc\text{-}exhaust}(u, n)) \end{array}}}
\vnbprop{cc-exhaust-seminorm-homogeneous}{\vnbstmt{\begin{array}{@{}l@{}} \forall u.; \forall c \in \mathbb{C}, a \in \operatorname{cc\text{-}cont\text{-}fns}(u), n \in \mathbb{N}. \\ \quad \operatorname{is\text{-}open}(\operatorname{nf\text{-}metric\text{-}space}(\operatorname{cc\text{-}normed\text{-}field}), u) \Rightarrow  \\ \quad \quad \operatorname{cc\text{-}sup\text{-}norm}((\lambda z \in u.\; (c \cdot a(z))), \operatorname{cc\text{-}exhaust}(u, n)) \\ \quad \quad =\; \operatorname{magnitude}(c) \cdot \operatorname{cc\text{-}sup\text{-}norm}(a, \operatorname{cc\text{-}exhaust}(u, n)) \end{array}}}
\vnbprop{cc-exhaust-seminorms-separate}{\vnbstmt{\begin{array}{@{}l@{}} \forall u.; \forall a \in \operatorname{cc\text{-}cont\text{-}fns}(u). \\ \quad (\operatorname{is\text{-}open}(\operatorname{nf\text{-}metric\text{-}space}(\operatorname{cc\text{-}normed\text{-}field}), u) \wedge \neg (a = (\lambda z \in u.\; 0))) \Rightarrow  \\ \quad \quad \exists k \in \mathbb{N}.\; \neg (\operatorname{cc\text{-}sup\text{-}norm}(a, \operatorname{cc\text{-}exhaust}(u, k)) = 0) \end{array}}}
\vnbprop{cc-exhaust-seminorms-family}{\vnbstmt{\begin{array}{@{}l@{}} \forall u. \\ \quad \operatorname{is\text{-}open}(\operatorname{nf\text{-}metric\text{-}space}(\operatorname{cc\text{-}normed\text{-}field}), u) \Rightarrow  \\ \quad \quad \operatorname{is\text{-}seminorm\text{-}family}(\operatorname{cc\text{-}normed\text{-}field}, \operatorname{cc\text{-}fn\text{-}module}(u, \operatorname{cc\text{-}cont\text{-}fns}(u)), \\ \quad \quad \quad \operatorname{cc\text{-}exhaust\text{-}seminorms}(u, \operatorname{cc\text{-}cont\text{-}fns}(u))) \end{array}}}
\vnbprop{cc-exhaust-seminorms-unfold}{\vnbstmt{\begin{array}{@{}l@{}} \forall u, v. \\ \quad (\operatorname{cc\text{-}exhaust\text{-}seminorms}(u, v) \simeq (\lambda k \in \mathbb{N}.\; (\lambda a \in v.\; \operatorname{cc\text{-}sup\text{-}norm}(a, \operatorname{cc\text{-}exhaust}(u, k))))) \end{array}}}
\vnbprop{cc-cont-fns-diff}{\vnbstmt{\begin{array}{@{}l@{}} \forall u.; \forall a, b \in \operatorname{cc\text{-}cont\text{-}fns}(u). \\ \quad (u \subseteq \mathbb{C}) \Rightarrow ((\lambda z \in u.\; (a(z) - b(z))) \in \operatorname{cc\text{-}cont\text{-}fns}(u)) \end{array}}}
\vnbprop{cc-cont-fns-add-neg}{\vnbstmt{\begin{array}{@{}l@{}} \forall u.; \forall a, b \in \operatorname{cc\text{-}cont\text{-}fns}(u). \\ \quad (u \subseteq \mathbb{C}) \Rightarrow ((\lambda z \in u.\; (a(z) + -b(z))) \in \operatorname{cc\text{-}cont\text{-}fns}(u)) \end{array}}}
\vnbprop{cc-fn-module-seminorm-of-difference}{\vnbstmt{\begin{array}{@{}l@{}} \forall u.; \forall a, b \in \operatorname{cc\text{-}cont\text{-}fns}(u), k \in \mathbb{N}. \\ \quad (u \subseteq \mathbb{C}) \Rightarrow  \\ \quad \quad (\operatorname{cc\text{-}exhaust\text{-}seminorms}(u, \operatorname{cc\text{-}cont\text{-}fns}(u))(k)(\operatorname{vadd}(\operatorname{cc\text{-}fn\text{-}module}(u, \operatorname{cc\text{-}cont\text{-}fns}(u)))(a, \operatorname{vneg}(\operatorname{cc\text{-}fn\text{-}module}(u, \operatorname{cc\text{-}cont\text{-}fns}(u)))(b))) \simeq \operatorname{cc\text{-}sup\text{-}norm}((\lambda z \in u.\; (a(z) - b(z))), \operatorname{cc\text{-}exhaust}(u, k))) \end{array}}}
\vnbprop{cc-fn-module-cauchy-uniform}{\vnbstmt{\begin{array}{@{}l@{}} \forall u.; \forall s \in (\mathbb{N} \to \operatorname{cc\text{-}cont\text{-}fns}(u)), k \in \mathbb{N}, a. \\ \quad \operatorname{is\text{-}open}(\operatorname{nf\text{-}metric\text{-}space}(\operatorname{cc\text{-}normed\text{-}field}), u) \wedge \\ \quad \operatorname{fr\text{-}cauchy}(\operatorname{cc\text{-}fn\text{-}module}(u, \operatorname{cc\text{-}cont\text{-}fns}(u)), \operatorname{cc\text{-}exhaust\text{-}seminorms}(u, \operatorname{cc\text{-}cont\text{-}fns}(u)), s) \wedge \\ \quad \operatorname{pos\text{-}rr}(a) \Rightarrow \\ \quad \quad \exists p \in \mathbb{N}. \forall i, j \in \mathbb{N}, z \in \operatorname{cc\text{-}exhaust}(u, k). \\ \quad \quad \quad ((p \leq i) \wedge (p \leq j)) \Rightarrow (\operatorname{magnitude}((s(i)(z) - s(j)(z))) < a) \end{array}}}
\vnbprop{cc-fn-module-pointwise-limit}{\vnbstmt{\begin{array}{@{}l@{}} \forall u.; \forall s \in (\mathbb{N} \to \operatorname{cc\text{-}cont\text{-}fns}(u)), z \in u. \\ \quad \operatorname{is\text{-}open}(\operatorname{nf\text{-}metric\text{-}space}(\operatorname{cc\text{-}normed\text{-}field}), u) \wedge \\ \quad \forall k \in \mathbb{N}, a. \\ \quad \quad \operatorname{pos\text{-}rr}(a) \Rightarrow  \\ \quad \quad \quad \exists p \in \mathbb{N}. \forall i, j \in \mathbb{N}, z \in \operatorname{cc\text{-}exhaust}(u, k). \\ \quad \quad \quad \quad ((p \leq i) \wedge (p \leq j)) \Rightarrow (\operatorname{magnitude}((s(i)(z) - s(j)(z))) < a) \Rightarrow \\ \quad \quad \exists l \in \mathbb{C}.\; \operatorname{converges\text{-}to}(\operatorname{cc\text{-}ms}, (\lambda \mathit{cfi} \in \mathbb{N}.\; s(\mathit{cfi})(z)), l) \end{array}}}
\vnbprop{cc-fn-module-limit-function}{\vnbstmt{\begin{array}{@{}l@{}} \forall u.; \forall s \in (\mathbb{N} \to \operatorname{cc\text{-}cont\text{-}fns}(u)). \\ \quad \operatorname{is\text{-}open}(\operatorname{nf\text{-}metric\text{-}space}(\operatorname{cc\text{-}normed\text{-}field}), u) \wedge \\ \quad \forall k \in \mathbb{N}, a. \\ \quad \quad \operatorname{pos\text{-}rr}(a) \Rightarrow  \\ \quad \quad \quad \exists p \in \mathbb{N}. \forall i, j \in \mathbb{N}, z \in \operatorname{cc\text{-}exhaust}(u, k). \\ \quad \quad \quad \quad ((p \leq i) \wedge (p \leq j)) \Rightarrow (\operatorname{magnitude}((s(i)(z) - s(j)(z))) < a) \Rightarrow \\ \quad \quad (\lambda \mathit{cfz} \in u.\; \varepsilon(\{l \in \mathbb{C} : \operatorname{converges\text{-}to}(\operatorname{cc\text{-}ms}, (\lambda \mathit{cfi} \in \mathbb{N}.\; s(\mathit{cfi})(\mathit{cfz})), l)\})) \\ \quad \quad \in\; (u \to \mathbb{C}) \wedge \\ \quad \quad \forall z \in u. \\ \quad \quad \quad \operatorname{converges\text{-}to}(\operatorname{cc\text{-}ms}, (\lambda \mathit{cfi} \in \mathbb{N}.\; s(\mathit{cfi})(z)), \\ \quad \quad \quad \quad (\lambda \mathit{cfz} \in u.\; \varepsilon(\{l \in \mathbb{C} : \operatorname{converges\text{-}to}(\operatorname{cc\text{-}ms}, (\lambda \mathit{cfi} \in \mathbb{N}.\; s(\mathit{cfi})(\mathit{cfz})), l)\}))(z)) \end{array}}}
\vnbprop{cc-fn-module-limit-uniform}{\vnbstmt{\begin{array}{@{}l@{}} \forall u.; \forall s \in (\mathbb{N} \to \operatorname{cc\text{-}cont\text{-}fns}(u)), k \in \mathbb{N}, a. \\ \quad \operatorname{is\text{-}open}(\operatorname{nf\text{-}metric\text{-}space}(\operatorname{cc\text{-}normed\text{-}field}), u) \wedge \\ \quad \forall k \in \mathbb{N}, a. \\ \quad \quad \operatorname{pos\text{-}rr}(a) \Rightarrow  \\ \quad \quad \quad \exists p \in \mathbb{N}. \forall i, j \in \mathbb{N}, z \in \operatorname{cc\text{-}exhaust}(u, k). \\ \quad \quad \quad \quad ((p \leq i) \wedge (p \leq j)) \Rightarrow (\operatorname{magnitude}((s(i)(z) - s(j)(z))) < a) \wedge \\ \quad \operatorname{pos\text{-}rr}(a) \Rightarrow \\ \quad \quad \exists p \in \mathbb{N}. \forall i \in \mathbb{N}, z \in \operatorname{cc\text{-}exhaust}(u, k). \\ \quad \quad \quad (p \leq i) \Rightarrow  \\ \quad \quad \quad \quad \operatorname{magnitude}((s(i)(z) - (\lambda \mathit{cfz} \in u.\; \varepsilon(\{l \in \mathbb{C} : \operatorname{converges\text{-}to}(\operatorname{cc\text{-}ms}, (\lambda \mathit{cfi} \in \mathbb{N}.\; s(\mathit{cfi})(\mathit{cfz})), l)\}))(z))) \\ \quad \quad \quad \quad \leq\; a \end{array}}}
\vnbprop{cc-uniform-limit-continuous-at}{\vnbstmt{\begin{array}{@{}l@{}} \forall u.; \forall s \in (\mathbb{N} \to \operatorname{cc\text{-}cont\text{-}fns}(u)), g \in (u \to \mathbb{C}), a \in u, o \in \mathbb{R}. \\ \quad u \subseteq \mathbb{C} \wedge \\ \quad 0 < o \wedge \\ \quad \operatorname{ball}(\operatorname{cc\text{-}ms}, a, o) \subseteq u \wedge \\ \quad \forall b. \\ \quad \quad \operatorname{pos\text{-}rr}(b) \Rightarrow  \\ \quad \quad \quad \exists p \in \mathbb{N}. \forall i \in \mathbb{N}, z \in \operatorname{ball}(\operatorname{cc\text{-}ms}, a, o). \\ \quad \quad \quad \quad (p \leq i) \Rightarrow (\operatorname{magnitude}((s(i)(z) - g(z))) \leq b) \Rightarrow \\ \quad \quad \operatorname{is\text{-}continuous\text{-}at}(\operatorname{subspace\text{-}ms}(\operatorname{nf\text{-}metric\text{-}space}(\operatorname{cc\text{-}normed\text{-}field}), u), \\ \quad \quad \quad \operatorname{nf\text{-}metric\text{-}space}(\operatorname{cc\text{-}normed\text{-}field}), g, a) \end{array}}}
\vnbprop{cc-cont-fns-frechet}{\vnbstmt{\begin{array}{@{}l@{}} \forall u. \\ \quad \operatorname{is\text{-}open}(\operatorname{nf\text{-}metric\text{-}space}(\operatorname{cc\text{-}normed\text{-}field}), u) \Rightarrow  \\ \quad \quad \operatorname{is\text{-}frechet\text{-}structure}(\operatorname{cc\text{-}normed\text{-}field}, \operatorname{cc\text{-}fn\text{-}module}(u, \operatorname{cc\text{-}cont\text{-}fns}(u)), \\ \quad \quad \quad \operatorname{cc\text{-}exhaust\text{-}seminorms}(u, \operatorname{cc\text{-}cont\text{-}fns}(u))) \end{array}}}
\vnbprop{cc-hol-fns-unfold}{\vnbstmt{\forall u.\; (\operatorname{cc\text{-}hol\text{-}fns}(u) \simeq \{f \in (u \to \mathbb{C}) : \operatorname{holomorphic\text{-}on}(u, f)\})}}
\vnbprop{cc-hol-fns-membership}{\vnbstmt{\forall u, h.\; ((h \in \operatorname{cc\text{-}hol\text{-}fns}(u)) \Leftrightarrow ((h \in (u \to \mathbb{C})) \wedge \operatorname{holomorphic\text{-}on}(u, h)))}}
\vnbprop{cc-hol-fns-subset-cont}{\vnbstmt{\forall u.\; (\operatorname{cc\text{-}hol\text{-}fns}(u) \subseteq \operatorname{cc\text{-}cont\text{-}fns}(u))}}
\vnbprop{cc-hol-fns-closed}{\vnbstmt{\begin{array}{@{}l@{}} \forall u.; \forall s \in (\mathbb{N} \to \operatorname{cc\text{-}hol\text{-}fns}(u)), f \in \operatorname{cc\text{-}cont\text{-}fns}(u). \\ \quad \operatorname{is\text{-}open}(\operatorname{nf\text{-}metric\text{-}space}(\operatorname{cc\text{-}normed\text{-}field}), u) \wedge \\ \quad \operatorname{fr\text{-}conv}(\operatorname{cc\text{-}fn\text{-}module}(u, \operatorname{cc\text{-}cont\text{-}fns}(u)), \operatorname{cc\text{-}exhaust\text{-}seminorms}(u, \operatorname{cc\text{-}cont\text{-}fns}(u)), s, f) \Rightarrow \\ \quad \quad f \in \operatorname{cc\text{-}hol\text{-}fns}(u) \end{array}}}

\vnbbookfile{theorem-library/ca5-montel}{theorem-library/ca5-montel}
\noindent{\small ca5-montel.scm -- CHAPTER 5 OF THE COMPLEX ANALYSIS NOTES, section 5.2: BOUNDED SETS OF HOLOMORPHIC FUNCTIONS, Proposition 5.8 (equicontinuity) and MONTEL'S THEOREM (Theorem 5.7 / Corollary 5.19 for sequences).  Agent CA-5, 2026-10-14.}\par
\vnbprop{cc-holomorphic-deriv-estimate}{\vnbstmt{\begin{array}{@{}l@{}} \forall u, f.; \forall y \in \mathbb{C}, q, m \in \mathbb{R}. \\ \quad \operatorname{is\text{-}open}(\operatorname{nf\text{-}metric\text{-}space}(\operatorname{cc\text{-}normed\text{-}field}), u) \wedge \\ \quad \operatorname{holomorphic\text{-}on}(u, f) \wedge \\ \quad 0 < q \wedge \\ \quad \operatorname{ball}(\operatorname{cc\text{-}ms}, y, (2 \cdot q)) \subseteq u \wedge \\ \quad \forall \mathit{cmy} \in \operatorname{ball}(\operatorname{cc\text{-}ms}, y, (2 \cdot q)).\; (\operatorname{magnitude}(f(\mathit{cmy})) \leq m) \Rightarrow \\ \quad \quad \operatorname{magnitude}(\operatorname{deriv\text{-}on}(\operatorname{cc\text{-}normed\text{-}field}, u, f, y)) \leq (m \cdot \operatorname{recip}(q)) \end{array}}}
\vnbprop{cc-holomorphic-lipschitz-at-centre}{\vnbstmt{\begin{array}{@{}l@{}} \forall u, f.; \forall c \in \mathbb{C}, q, m \in \mathbb{R}, z \in \operatorname{ball}(\operatorname{cc\text{-}ms}, c, q). \\ \quad \operatorname{is\text{-}open}(\operatorname{nf\text{-}metric\text{-}space}(\operatorname{cc\text{-}normed\text{-}field}), u) \wedge \\ \quad \operatorname{holomorphic\text{-}on}(u, f) \wedge \\ \quad 0 < q \wedge \\ \quad \operatorname{ball}(\operatorname{cc\text{-}ms}, c, (3 \cdot q)) \subseteq u \wedge \\ \quad \forall y \in \operatorname{ball}(\operatorname{cc\text{-}ms}, c, (3 \cdot q)).\; (\operatorname{magnitude}(f(y)) \leq m) \Rightarrow \\ \quad \quad \operatorname{magnitude}((f(z) - f(c))) \leq ((m \cdot \operatorname{recip}(q)) \cdot \operatorname{magnitude}((z - c))) \end{array}}}
\vnbprop{cc-hol-bounded-equicontinuous}{\vnbstmt{\begin{array}{@{}l@{}} \forall u, a.; \forall b \in u, s. \\ \quad (\operatorname{is\text{-}open}(\operatorname{nf\text{-}metric\text{-}space}(\operatorname{cc\text{-}normed\text{-}field}), u) \wedge (\operatorname{cc\text{-}hol\text{-}bounded}(u, a) \wedge \operatorname{pos\text{-}rr}(s))) \Rightarrow  \\ \quad \quad \exists l. \\ \quad \quad \quad (\operatorname{pos\text{-}rr}(l) \wedge \forall f \in a, z \in u.\; ((\operatorname{magnitude}((z - b)) < l) \Rightarrow (\operatorname{magnitude}((f(z) - f(b))) < s))) \end{array}}}
\vnbprop{cc-re-unif-cauchy-on-compact}{\vnbstmt{\begin{array}{@{}l@{}} \forall a, q, t, s. \\ \quad a \subseteq \mathbb{C} \wedge \\ \quad \operatorname{is\text{-}compact}(\operatorname{subspace\text{-}ms}(\operatorname{cc\text{-}ms}, a)) \wedge \\ \quad \operatorname{is\text{-}dense\text{-}seq}(\operatorname{subspace\text{-}ms}(\operatorname{cc\text{-}ms}, a), q) \wedge \\ \quad \forall k \in \mathbb{N}, z \in a.\; (t(k)(z) \in \mathbb{C}) \wedge \\ \quad \forall x \in a, s. \\ \quad \quad \operatorname{pos\text{-}rr}(s) \Rightarrow  \\ \quad \quad \quad \exists l. \\ \quad \quad \quad \quad \operatorname{pos\text{-}rr}(l) \wedge \\ \quad \quad \quad \quad \forall k \in \mathbb{N}, y \in a. \\ \quad \quad \quad \quad \quad (\operatorname{magnitude}((y - x)) < l) \Rightarrow (\operatorname{magnitude}((t(k)(y) - t(k)(x))) < s) \wedge \\ \quad \forall j \in \mathbb{N}. \exists \ell \in \mathbb{R}. \\ \quad \quad \operatorname{converges\text{-}to}(\operatorname{rr\text{-}ms}, (\lambda \mathit{cmk} \in \mathbb{N}.\; \operatorname{real\text{-}part}(t(\mathit{cmk})(q(j)))), \ell) \wedge \\ \quad \operatorname{pos\text{-}rr}(s) \Rightarrow \\ \quad \quad \exists p \in \mathbb{N}. \forall i, j \in \mathbb{N}, z \in a. \\ \quad \quad \quad ((p \leq i) \wedge (p \leq j)) \Rightarrow (\lvert (\operatorname{real\text{-}part}(t(i)(z)) - \operatorname{real\text{-}part}(t(j)(z))) \rvert < s) \end{array}}}
\vnbprop{cc-im-unif-cauchy-on-compact}{\vnbstmt{\begin{array}{@{}l@{}} \forall a, q, t, s. \\ \quad a \subseteq \mathbb{C} \wedge \\ \quad \operatorname{is\text{-}compact}(\operatorname{subspace\text{-}ms}(\operatorname{cc\text{-}ms}, a)) \wedge \\ \quad \operatorname{is\text{-}dense\text{-}seq}(\operatorname{subspace\text{-}ms}(\operatorname{cc\text{-}ms}, a), q) \wedge \\ \quad \forall k \in \mathbb{N}, z \in a.\; (t(k)(z) \in \mathbb{C}) \wedge \\ \quad \forall x \in a, s. \\ \quad \quad \operatorname{pos\text{-}rr}(s) \Rightarrow  \\ \quad \quad \quad \exists l. \\ \quad \quad \quad \quad \operatorname{pos\text{-}rr}(l) \wedge \\ \quad \quad \quad \quad \forall k \in \mathbb{N}, y \in a. \\ \quad \quad \quad \quad \quad (\operatorname{magnitude}((y - x)) < l) \Rightarrow (\operatorname{magnitude}((t(k)(y) - t(k)(x))) < s) \wedge \\ \quad \forall j \in \mathbb{N}. \exists \ell \in \mathbb{R}. \\ \quad \quad \operatorname{converges\text{-}to}(\operatorname{rr\text{-}ms}, (\lambda \mathit{cmk} \in \mathbb{N}.\; \operatorname{imag\text{-}part}(t(\mathit{cmk})(q(j)))), \ell) \wedge \\ \quad \operatorname{pos\text{-}rr}(s) \Rightarrow \\ \quad \quad \exists p \in \mathbb{N}. \forall i, j \in \mathbb{N}, z \in a. \\ \quad \quad \quad ((p \leq i) \wedge (p \leq j)) \Rightarrow (\lvert (\operatorname{imag\text{-}part}(t(i)(z)) - \operatorname{imag\text{-}part}(t(j)(z))) \rvert < s) \end{array}}}
\vnbprop{cc-unif-cauchy-on-compact}{\vnbstmt{\begin{array}{@{}l@{}} \forall a, q, t, s. \\ \quad a \subseteq \mathbb{C} \wedge \\ \quad \operatorname{is\text{-}compact}(\operatorname{subspace\text{-}ms}(\operatorname{cc\text{-}ms}, a)) \wedge \\ \quad \operatorname{is\text{-}dense\text{-}seq}(\operatorname{subspace\text{-}ms}(\operatorname{cc\text{-}ms}, a), q) \wedge \\ \quad \forall k \in \mathbb{N}, z \in a.\; (t(k)(z) \in \mathbb{C}) \wedge \\ \quad \forall x \in a, s. \\ \quad \quad \operatorname{pos\text{-}rr}(s) \Rightarrow  \\ \quad \quad \quad \exists l. \\ \quad \quad \quad \quad \operatorname{pos\text{-}rr}(l) \wedge \\ \quad \quad \quad \quad \forall k \in \mathbb{N}, y \in a. \\ \quad \quad \quad \quad \quad (\operatorname{magnitude}((y - x)) < l) \Rightarrow (\operatorname{magnitude}((t(k)(y) - t(k)(x))) < s) \wedge \\ \quad \forall j \in \mathbb{N}. \exists \ell \in \mathbb{C}. \\ \quad \quad \operatorname{converges\text{-}to}(\operatorname{cc\text{-}ms}, (\lambda \mathit{cmk} \in \mathbb{N}.\; t(\mathit{cmk})(q(j))), \ell) \wedge \\ \quad \operatorname{pos\text{-}rr}(s) \Rightarrow \\ \quad \quad \exists p \in \mathbb{N}. \forall i, j \in \mathbb{N}, z \in a. \\ \quad \quad \quad ((p \leq i) \wedge (p \leq j)) \Rightarrow (\operatorname{magnitude}((t(i)(z) - t(j)(z))) < s) \end{array}}}
\vnbprop{cc-pointwise-diagonal}{\vnbstmt{\begin{array}{@{}l@{}} \forall u, s, d. \\ \quad u \subseteq \mathbb{C} \wedge \\ \quad \forall k \in \mathbb{N}, z \in u.\; (s(k)(z) \in \mathbb{C}) \wedge \\ \quad \forall z \in u. \exists a \in \mathbb{R}. \forall k \in \mathbb{N}. \\ \quad \quad \operatorname{magnitude}(s(k)(z)) \leq a \wedge \\ \quad \operatorname{is\text{-}dense\text{-}seq}(\operatorname{subspace\text{-}ms}(\operatorname{cc\text{-}ms}, u), d) \Rightarrow \\ \quad \quad \exists i. \\ \quad \quad \quad \operatorname{strictly\text{-}mono\text{-}nn}(i) \wedge \\ \quad \quad \quad \forall j \in \mathbb{N}. \exists \ell \in \mathbb{C}. \\ \quad \quad \quad \quad \operatorname{converges\text{-}to}(\operatorname{cc\text{-}ms}, (\lambda \mathit{cmk} \in \mathbb{N}.\; s(i(\mathit{cmk}))(d(j))), \ell) \end{array}}}
\vnbprop{cc-exhaust-dense-seqs}{\vnbstmt{\begin{array}{@{}l@{}} \forall u.; \forall p \in u. \\ \quad \operatorname{is\text{-}open}(\operatorname{nf\text{-}metric\text{-}space}(\operatorname{cc\text{-}normed\text{-}field}), u) \Rightarrow  \\ \quad \quad \exists a \in \mathbb{N}, q \in (\mathbb{N} \to (\mathbb{N} \to \mathbb{C})). \forall n \in \mathbb{N}. \\ \quad \quad \quad \operatorname{is\text{-}dense\text{-}seq}(\operatorname{subspace\text{-}ms}(\operatorname{cc\text{-}ms}, \operatorname{cc\text{-}exhaust}(u, (n + a))), q(n)) \end{array}}}
\vnbprop{cc-exhaust-dense-in-u}{\vnbstmt{\begin{array}{@{}l@{}} \forall u.; \forall a \in \mathbb{N}, q \in (\mathbb{N} \to (\mathbb{N} \to \mathbb{C})). \\ \quad \operatorname{is\text{-}open}(\operatorname{nf\text{-}metric\text{-}space}(\operatorname{cc\text{-}normed\text{-}field}), u) \wedge \\ \quad \forall n \in \mathbb{N}.\; \operatorname{is\text{-}dense\text{-}seq}(\operatorname{subspace\text{-}ms}(\operatorname{cc\text{-}ms}, \operatorname{cc\text{-}exhaust}(u, (n + a))), q(n)) \Rightarrow \\ \quad \quad \operatorname{is\text{-}dense\text{-}seq}(\operatorname{subspace\text{-}ms}(\operatorname{cc\text{-}ms}, u), (\lambda m \in \mathbb{N}.\; q(\operatorname{nnfst}(m))(\operatorname{nnsnd}(m)))) \end{array}}}
\vnbprop{cc-montel-tail}{\vnbstmt{\begin{array}{@{}l@{}} \forall u.; \forall s \in (\mathbb{N} \to \operatorname{cc\text{-}hol\text{-}fns}(u)). \\ \quad \operatorname{is\text{-}open}(\operatorname{nf\text{-}metric\text{-}space}(\operatorname{cc\text{-}normed\text{-}field}), u) \wedge \\ \quad \forall k \in \mathbb{N}, a. \\ \quad \quad \operatorname{pos\text{-}rr}(a) \Rightarrow  \\ \quad \quad \quad \exists p \in \mathbb{N}. \forall i, j \in \mathbb{N}, z \in \operatorname{cc\text{-}exhaust}(u, k). \\ \quad \quad \quad \quad ((p \leq i) \wedge (p \leq j)) \Rightarrow (\operatorname{magnitude}((s(i)(z) - s(j)(z))) < a) \Rightarrow \\ \quad \quad \exists g \in \operatorname{cc\text{-}hol\text{-}fns}(u).\; \operatorname{unif\text{-}on\text{-}compacta}(u, s, g) \end{array}}}
\vnbprop{cc-montel-subsequence}{\vnbstmt{\begin{array}{@{}l@{}} \forall u.; \forall s \in (\mathbb{N} \to \operatorname{cc\text{-}hol\text{-}fns}(u)). \\ \quad (\operatorname{is\text{-}open}(\operatorname{nf\text{-}metric\text{-}space}(\operatorname{cc\text{-}normed\text{-}field}), u) \wedge \operatorname{cc\text{-}hol\text{-}bounded}(u, \operatorname{image}(s, \mathbb{N}))) \Rightarrow  \\ \quad \quad \exists a. \\ \quad \quad \quad \operatorname{strictly\text{-}mono\text{-}nn}(a) \wedge \\ \quad \quad \quad \forall k \in \mathbb{N}, b. \\ \quad \quad \quad \quad \operatorname{pos\text{-}rr}(b) \Rightarrow  \\ \quad \quad \quad \quad \quad \exists p \in \mathbb{N}. \forall i, j \in \mathbb{N}, z \in \operatorname{cc\text{-}exhaust}(u, k). \\ \quad \quad \quad \quad \quad \quad ((p \leq i) \wedge (p \leq j)) \Rightarrow  \\ \quad \quad \quad \quad \quad \quad \quad \operatorname{magnitude}((\operatorname{subseq}(s, a)(i)(z) - \operatorname{subseq}(s, a)(j)(z))) < b \end{array}}}
\vnbprop{cc-montel}{\vnbstmt{\begin{array}{@{}l@{}} \forall u.; \forall s \in (\mathbb{N} \to \operatorname{cc\text{-}hol\text{-}fns}(u)). \\ \quad (\operatorname{is\text{-}open}(\operatorname{nf\text{-}metric\text{-}space}(\operatorname{cc\text{-}normed\text{-}field}), u) \wedge \operatorname{cc\text{-}hol\text{-}bounded}(u, \operatorname{image}(s, \mathbb{N}))) \Rightarrow  \\ \quad \quad \exists i. \\ \quad \quad \quad (\operatorname{strictly\text{-}mono\text{-}nn}(i) \wedge \exists g \in \operatorname{cc\text{-}hol\text{-}fns}(u).\; \operatorname{unif\text{-}on\text{-}compacta}(u, \operatorname{subseq}(s, i), g)) \end{array}}}
\vnbprop{cc-fr-conv-iff-unif-on-compacta}{\vnbstmt{\begin{array}{@{}l@{}} \forall u.; \forall s \in (\mathbb{N} \to \operatorname{cc\text{-}cont\text{-}fns}(u)), f \in \operatorname{cc\text{-}cont\text{-}fns}(u). \\ \quad \operatorname{is\text{-}open}(\operatorname{nf\text{-}metric\text{-}space}(\operatorname{cc\text{-}normed\text{-}field}), u) \Rightarrow  \\ \quad \quad \operatorname{fr\text{-}conv}(\operatorname{cc\text{-}fn\text{-}module}(u, \operatorname{cc\text{-}cont\text{-}fns}(u)), \operatorname{cc\text{-}exhaust\text{-}seminorms}(u, \operatorname{cc\text{-}cont\text{-}fns}(u)), s, f) \\ \quad \quad \Leftrightarrow\; \operatorname{unif\text{-}on\text{-}compacta}(u, s, f) \end{array}}}
\vnbprop{cc-deriv-continuous-on-hol}{\vnbstmt{\begin{array}{@{}l@{}} \forall u.; \forall s \in (\mathbb{N} \to \operatorname{cc\text{-}hol\text{-}fns}(u)), f \in \operatorname{cc\text{-}hol\text{-}fns}(u). \\ \quad \operatorname{is\text{-}open}(\operatorname{nf\text{-}metric\text{-}space}(\operatorname{cc\text{-}normed\text{-}field}), u) \wedge \\ \quad \operatorname{fr\text{-}conv}(\operatorname{cc\text{-}fn\text{-}module}(u, \operatorname{cc\text{-}cont\text{-}fns}(u)), \operatorname{cc\text{-}exhaust\text{-}seminorms}(u, \operatorname{cc\text{-}cont\text{-}fns}(u)), s, f) \Rightarrow \\ \quad \quad \operatorname{fr\text{-}conv}(\operatorname{cc\text{-}fn\text{-}module}(u, \operatorname{cc\text{-}cont\text{-}fns}(u)), \\ \quad \quad \quad \operatorname{cc\text{-}exhaust\text{-}seminorms}(u, \operatorname{cc\text{-}cont\text{-}fns}(u)), \\ \quad \quad \quad (\lambda k \in \mathbb{N}.\; (\lambda w \in u.\; \operatorname{deriv\text{-}on}(\operatorname{cc\text{-}normed\text{-}field}, u, s(k), w))), \\ \quad \quad \quad (\lambda w \in u.\; \operatorname{deriv\text{-}on}(\operatorname{cc\text{-}normed\text{-}field}, u, f, w))) \end{array}}}

\vnbbookfile{theorem-library/ca5-series}{theorem-library/ca5-series}
\noindent{\small ca5-series.scm -- CHAPTER 5 OF THE COMPLEX ANALYSIS NOTES, section 5.3: DIFFERENTIATION AS A LINEAR OPERATOR, Lemma 5.21, Proposition 5.22 and the linearity half of Proposition 5.20. Agent CA-5, 2026-10-14.}\par
\vnbprop{cc-exhaust-radius-quarter}{\vnbstmt{\forall n \in \mathbb{N}.\; (\operatorname{recip}((((4 \cdot n) + 3) + 1)) = (1/4 \cdot \operatorname{recip}((n + 1))))}}
\vnbprop{cc-exhaust-ball-inside}{\vnbstmt{\begin{array}{@{}l@{}} \forall u.; \forall n \in \mathbb{N}, y \in \operatorname{cc\text{-}exhaust}(u, n). \\ \quad \operatorname{ball}(\operatorname{cc\text{-}ms}, y, (1/2 \cdot \operatorname{recip}((n + 1)))) \subseteq \operatorname{cc\text{-}exhaust}(u, ((4 \cdot n) + 3)) \end{array}}}
\vnbprop{cc-deriv-seminorm-estimate}{\vnbstmt{\begin{array}{@{}l@{}} \forall u, f.; \forall n \in \mathbb{N}, m \in \mathbb{R}, y \in \operatorname{cc\text{-}exhaust}(u, n). \\ \quad \operatorname{is\text{-}open}(\operatorname{nf\text{-}metric\text{-}space}(\operatorname{cc\text{-}normed\text{-}field}), u) \wedge \\ \quad \operatorname{holomorphic\text{-}on}(u, f) \wedge \\ \quad \forall \mathit{cmy} \in \operatorname{cc\text{-}exhaust}(u, ((4 \cdot n) + 3)).\; (\operatorname{magnitude}(f(\mathit{cmy})) \leq m) \Rightarrow \\ \quad \quad \operatorname{magnitude}(\operatorname{deriv\text{-}on}(\operatorname{cc\text{-}normed\text{-}field}, u, f, y)) \leq (m \cdot (4 \cdot (n + 1))) \end{array}}}
\vnbprop{cc-derived-series-normally-convergent}{\vnbstmt{\begin{array}{@{}l@{}} \forall u.; \forall g \in (\mathbb{N} \to \operatorname{cc\text{-}hol\text{-}fns}(u)), n \in \mathbb{N}. \\ \quad \operatorname{is\text{-}open}(\operatorname{nf\text{-}metric\text{-}space}(\operatorname{cc\text{-}normed\text{-}field}), u) \wedge \\ \quad \forall n \in \mathbb{N}.\; \operatorname{is\text{-}normally\text{-}convergent}(g, \operatorname{cc\text{-}exhaust}(u, n)) \Rightarrow \\ \quad \quad \operatorname{is\text{-}normally\text{-}convergent}((\lambda i \in \mathbb{N}.\; (\lambda z \in u.\; \operatorname{deriv\text{-}on}(\operatorname{cc\text{-}normed\text{-}field}, u, g(i), z))), \\ \quad \quad \quad \operatorname{cc\text{-}exhaust}(u, n)) \end{array}}}
\vnbprop{cc-normal-convergence-pointwise}{\vnbstmt{\begin{array}{@{}l@{}} \forall g, s.; \forall z \in s. \\ \quad (\forall i \in \mathbb{N}, y \in s.\; (g(i)(y) \in \mathbb{C}) \wedge \operatorname{is\text{-}normally\text{-}convergent}(g, s)) \Rightarrow  \\ \quad \quad \operatorname{cc\text{-}series\text{-}converges}((\lambda i \in \mathbb{N}.\; g(i)(z))) \end{array}}}
\vnbprop{cc-normal-convergence-uniform}{\vnbstmt{\begin{array}{@{}l@{}} \forall g, s, a. \\ \quad (\forall i \in \mathbb{N}, y \in s.\; (g(i)(y) \in \mathbb{C}) \wedge (\operatorname{is\text{-}normally\text{-}convergent}(g, s) \wedge \operatorname{pos\text{-}rr}(a))) \Rightarrow  \\ \quad \quad \exists p \in \mathbb{N}. \forall n \in \mathbb{N}, z \in s. \\ \quad \quad \quad (p \leq n) \Rightarrow  \\ \quad \quad \quad \quad \operatorname{magnitude}((\operatorname{cc\text{-}series\text{-}limit}((\lambda i \in \mathbb{N}.\; g(i)(z))) - \operatorname{cc\text{-}series\text{-}partial\text{-}sum}((\lambda i \in \mathbb{N}.\; g(i)(z)), n))) \\ \quad \quad \quad \quad \leq\; a \end{array}}}
\vnbprop{cc-partial-sums-diff}{\vnbstmt{\begin{array}{@{}l@{}} \forall n \in \mathbb{N}, u.; \forall g \in (\mathbb{N} \to \operatorname{cc\text{-}hol\text{-}fns}(u)), z \in u. \\ \quad \operatorname{is\text{-}open}(\operatorname{nf\text{-}metric\text{-}space}(\operatorname{cc\text{-}normed\text{-}field}), u) \Rightarrow  \\ \quad \quad \operatorname{is\text{-}diff\text{-}on}(\operatorname{cc\text{-}normed\text{-}field}, u, \\ \quad \quad \quad (\lambda w \in u.\; \operatorname{cc\text{-}series\text{-}partial\text{-}sum}((\lambda i \in \mathbb{N}.\; g(i)(w)), n)), \\ \quad \quad \quad z, \\ \quad \quad \quad \operatorname{cc\text{-}series\text{-}partial\text{-}sum}((\lambda \mathit{cdi} \in \mathbb{N}.\; \operatorname{deriv\text{-}on}(\operatorname{cc\text{-}normed\text{-}field}, u, g(\mathit{cdi}), z)), n)) \end{array}}}
\vnbprop{cc-series-holomorphic}{\vnbstmt{\begin{array}{@{}l@{}} \forall u.; \forall g \in (\mathbb{N} \to \operatorname{cc\text{-}hol\text{-}fns}(u)). \\ \quad \operatorname{is\text{-}open}(\operatorname{nf\text{-}metric\text{-}space}(\operatorname{cc\text{-}normed\text{-}field}), u) \wedge \\ \quad \forall n \in \mathbb{N}.\; \operatorname{is\text{-}normally\text{-}convergent}(g, \operatorname{cc\text{-}exhaust}(u, n)) \Rightarrow \\ \quad \quad \operatorname{holomorphic\text{-}on}(u, (\lambda w \in u.\; \operatorname{cc\text{-}series\text{-}limit}((\lambda i \in \mathbb{N}.\; g(i)(w))))) \wedge \\ \quad \quad \forall z \in u. \\ \quad \quad \quad \operatorname{deriv\text{-}on}(\operatorname{cc\text{-}normed\text{-}field}, u, (\lambda w \in u.\; \operatorname{cc\text{-}series\text{-}limit}((\lambda i \in \mathbb{N}.\; g(i)(w)))), z) \\ \quad \quad \quad =\; \operatorname{cc\text{-}series\text{-}limit}((\lambda \mathit{cdi} \in \mathbb{N}.\; \operatorname{deriv\text{-}on}(\operatorname{cc\text{-}normed\text{-}field}, u, g(\mathit{cdi}), z))) \end{array}}}
\vnbprop{cc-deriv-linear}{\vnbstmt{\begin{array}{@{}l@{}} \forall u, f, g.; \forall a, b \in \mathbb{C}, z \in u. \\ \quad (\operatorname{is\text{-}open}(\operatorname{nf\text{-}metric\text{-}space}(\operatorname{cc\text{-}normed\text{-}field}), u) \wedge (\operatorname{holomorphic\text{-}on}(u, f) \wedge \operatorname{holomorphic\text{-}on}(u, g))) \Rightarrow  \\ \quad \quad \operatorname{deriv\text{-}on}(\operatorname{cc\text{-}normed\text{-}field}, u, (\lambda w \in u.\; ((a \cdot f(w)) + (b \cdot g(w)))), z) \\ \quad \quad =\; (a \cdot \operatorname{deriv\text{-}on}(\operatorname{cc\text{-}normed\text{-}field}, u, f, z)) + (b \cdot \operatorname{deriv\text{-}on}(\operatorname{cc\text{-}normed\text{-}field}, u, g, z)) \end{array}}}

\vnbbookfile{theorem-library/ca6-hurwitz}{theorem-library/ca6-hurwitz}
\noindent{\small ca6-hurwitz.scm -- CHAPTER 5 OF THE COMPLEX ANALYSIS NOTES (\textasciitilde{}/docs/complex-analysis.pdf), the rest of sections 5.2-5.3: COROLLARY 5.19 (pointwise limits of bounded sequences), HURWITZ (Proposition 5.23 and Remark 5.24), INJECTIVE LIMITS (Proposition 5.25), and PROPOSITION 5.20 AS AN OPERATOR (D : f |-$>$ f' a continuous ...}\par
\vnbprop{cc-unif-on-compacta-pointwise}{\vnbstmt{\begin{array}{@{}l@{}} \forall u, s.; \forall h \in (u \to \mathbb{C}), z \in u. \\ \quad \operatorname{is\text{-}open}(\operatorname{nf\text{-}metric\text{-}space}(\operatorname{cc\text{-}normed\text{-}field}), u) \wedge \\ \quad \forall k \in \mathbb{N}, z \in u.\; (s(k)(z) \in \mathbb{C}) \wedge \\ \quad \operatorname{unif\text{-}on\text{-}compacta}(u, s, h) \Rightarrow \\ \quad \quad \operatorname{converges\text{-}to}(\operatorname{cc\text{-}ms}, (\lambda \mathit{cpk} \in \mathbb{N}.\; s(\mathit{cpk})(z)), h(z)) \end{array}}}
\vnbprop{cc-pointwise-subseq}{\vnbstmt{\begin{array}{@{}l@{}} \forall u, s, i.; \forall z \in u, w. \\ \quad \forall k \in \mathbb{N}, z \in u.\; (s(k)(z) \in \mathbb{C}) \wedge \\ \quad \operatorname{strictly\text{-}mono\text{-}nn}(i) \wedge \\ \quad \operatorname{converges\text{-}to}(\operatorname{cc\text{-}ms}, (\lambda \mathit{cpk} \in \mathbb{N}.\; s(\mathit{cpk})(z)), w) \Rightarrow \\ \quad \quad \operatorname{converges\text{-}to}(\operatorname{cc\text{-}ms}, (\lambda \mathit{cpk} \in \mathbb{N}.\; \operatorname{subseq}(s, i)(\mathit{cpk})(z)), w) \end{array}}}
\vnbprop{cc-unif-pointwise-agree}{\vnbstmt{\begin{array}{@{}l@{}} \forall u, s.; \forall h, g \in (u \to \mathbb{C}). \\ \quad \operatorname{is\text{-}open}(\operatorname{nf\text{-}metric\text{-}space}(\operatorname{cc\text{-}normed\text{-}field}), u) \wedge \\ \quad \forall k \in \mathbb{N}, z \in u.\; (s(k)(z) \in \mathbb{C}) \wedge \\ \quad \operatorname{unif\text{-}on\text{-}compacta}(u, s, h) \wedge \\ \quad \forall z \in u.\; \operatorname{converges\text{-}to}(\operatorname{cc\text{-}ms}, (\lambda \mathit{cpk} \in \mathbb{N}.\; s(\mathit{cpk})(z)), g(z)) \Rightarrow \\ \quad \quad h = g \end{array}}}
\vnbprop{cc-hol-bounded-mono}{\vnbstmt{\begin{array}{@{}l@{}} \forall u, b, a. \\ \quad (\operatorname{cc\text{-}hol\text{-}bounded}(u, b) \wedge (a \subseteq b)) \Rightarrow \operatorname{cc\text{-}hol\text{-}bounded}(u, a) \end{array}}}
\vnbprop{cc-hol-bounded-subseq}{\vnbstmt{\begin{array}{@{}l@{}} \forall u, s, i. \\ \quad (\operatorname{cc\text{-}hol\text{-}bounded}(u, \operatorname{image}(s, \mathbb{N})) \wedge \operatorname{strictly\text{-}mono\text{-}nn}(i)) \Rightarrow  \\ \quad \quad \operatorname{cc\text{-}hol\text{-}bounded}(u, \operatorname{image}(\operatorname{subseq}(s, i), \mathbb{N})) \end{array}}}
\vnbprop{cc-montel-pointwise-limit}{\vnbstmt{\begin{array}{@{}l@{}} \forall u.; \forall s \in (\mathbb{N} \to \operatorname{cc\text{-}hol\text{-}fns}(u)), g \in (u \to \mathbb{C}). \\ \quad \operatorname{is\text{-}open}(\operatorname{nf\text{-}metric\text{-}space}(\operatorname{cc\text{-}normed\text{-}field}), u) \wedge \\ \quad \operatorname{cc\text{-}hol\text{-}bounded}(u, \operatorname{image}(s, \mathbb{N})) \wedge \\ \quad \forall z \in u.\; \operatorname{converges\text{-}to}(\operatorname{cc\text{-}ms}, (\lambda k \in \mathbb{N}.\; s(k)(z)), g(z)) \Rightarrow \\ \quad \quad (g \in \operatorname{cc\text{-}hol\text{-}fns}(u)) \wedge \operatorname{unif\text{-}on\text{-}compacta}(u, s, g) \end{array}}}
\vnbprop{cc-hurwitz-disc-bound}{\vnbstmt{\begin{array}{@{}l@{}} \forall u, g.; \forall c \in \mathbb{C}, r, q \in \mathbb{R}, a \in \mathbb{C}, m \in \mathbb{R}. \\ \quad \operatorname{holomorphic\text{-}on}(u, g) \wedge \\ \quad 0 < r \wedge \\ \quad r < q \wedge \\ \quad \operatorname{ball}(\operatorname{cc\text{-}ms}, c, q) \subseteq u \wedge \\ \quad \forall y \in \operatorname{ball}(\operatorname{cc\text{-}ms}, c, q).\; \neg (g(y) = a) \wedge \\ \quad 0 < m \wedge \\ \quad \forall y \in \operatorname{trace}(\operatorname{circle\text{-}path}(c, r), 0, (2 \cdot \pi)).\; (m \leq \operatorname{magnitude}((g(y) - a))) \Rightarrow \\ \quad \quad m \leq \operatorname{magnitude}((g(c) - a)) \end{array}}}
\vnbprop{cc-hurwitz-circle-gap}{\vnbstmt{\begin{array}{@{}l@{}} \forall u, f.; \forall c \in \mathbb{C}, r, q \in \mathbb{R}, a \in \mathbb{C}. \\ \quad \operatorname{holomorphic\text{-}on}(u, f) \wedge \\ \quad 0 < r \wedge \\ \quad r < q \wedge \\ \quad \operatorname{ball}(\operatorname{cc\text{-}ms}, c, q) \subseteq u \wedge \\ \quad \forall y \in \operatorname{ball}(\operatorname{cc\text{-}ms}, c, q).\; (\neg (y = c) \Rightarrow \neg (f(y) = a)) \Rightarrow \\ \quad \quad \exists m \in \mathbb{R}. \\ \quad \quad \quad ((0 < m) \wedge \forall y \in \operatorname{trace}(\operatorname{circle\text{-}path}(c, r), 0, (2 \cdot \pi)).\; (m \leq \operatorname{magnitude}((f(y) - a)))) \end{array}}}
\vnbprop{cc-hurwitz-local}{\vnbstmt{\begin{array}{@{}l@{}} \forall u.; \forall s \in (\mathbb{N} \to \operatorname{cc\text{-}hol\text{-}fns}(u)), f \in (u \to \mathbb{C}), a \in \mathbb{C}, c \in u, p, r \in \mathbb{R}. \\ \quad \operatorname{is\text{-}open}(\operatorname{nf\text{-}metric\text{-}space}(\operatorname{cc\text{-}normed\text{-}field}), u) \wedge \\ \quad \operatorname{unif\text{-}on\text{-}compacta}(u, s, f) \wedge \\ \quad f(c) = a \wedge \\ \quad 0 < p \wedge \\ \quad \forall x \in u.\; (((x \in \operatorname{ball}(\operatorname{cc\text{-}ms}, c, p)) \wedge \neg (x = c)) \Rightarrow \neg (f(x) = a)) \wedge \\ \quad 0 < r \Rightarrow \\ \quad \quad \exists n \in \mathbb{N}. \forall k \in \mathbb{N}. \\ \quad \quad \quad (n \leq k) \Rightarrow \exists z \in u.\; ((z \in \operatorname{ball}(\operatorname{cc\text{-}ms}, c, r)) \wedge (s(k)(z) = a)) \end{array}}}
\vnbprop{cc-hurwitz-eventually}{\vnbstmt{\begin{array}{@{}l@{}} \forall u.; \forall s \in (\mathbb{N} \to \operatorname{cc\text{-}hol\text{-}fns}(u)), f \in (u \to \mathbb{C}), a \in \mathbb{C}, c \in u, r \in \mathbb{R}. \\ \quad \operatorname{is\text{-}open}(\operatorname{nf\text{-}metric\text{-}space}(\operatorname{cc\text{-}normed\text{-}field}), u) \wedge \\ \quad \operatorname{is\text{-}connected}(\operatorname{cc\text{-}ms}, u) \wedge \\ \quad \operatorname{unif\text{-}on\text{-}compacta}(u, s, f) \wedge \\ \quad \neg \exists \mathit{hcc} \in \mathbb{C}.\; \forall z \in u.\; (f(z) = \mathit{hcc}) \wedge \\ \quad f(c) = a \wedge \\ \quad 0 < r \Rightarrow \\ \quad \quad \exists n \in \mathbb{N}. \forall k \in \mathbb{N}. \\ \quad \quad \quad (n \leq k) \Rightarrow \exists \mathit{hzz} \in u.\; ((\mathit{hzz} \in \operatorname{ball}(\operatorname{cc\text{-}ms}, c, r)) \wedge (s(k)(\mathit{hzz}) = a)) \end{array}}}
\vnbprop{cc-hurwitz-value}{\vnbstmt{\begin{array}{@{}l@{}} \forall u.; \forall s \in (\mathbb{N} \to \operatorname{cc\text{-}hol\text{-}fns}(u)), a \in \mathbb{C}, f \in (u \to \mathbb{C}). \\ \quad \operatorname{is\text{-}open}(\operatorname{nf\text{-}metric\text{-}space}(\operatorname{cc\text{-}normed\text{-}field}), u) \wedge \\ \quad \operatorname{is\text{-}connected}(\operatorname{cc\text{-}ms}, u) \wedge \\ \quad \forall k \in \mathbb{N}, z \in u.\; \neg (s(k)(z) = a) \wedge \\ \quad \operatorname{unif\text{-}on\text{-}compacta}(u, s, f) \Rightarrow \\ \quad \quad (\exists c \in \mathbb{C}.\; \forall \mathit{hcz} \in u.\; (f(\mathit{hcz}) = c) \vee \forall z \in u.\; \neg (f(z) = a)) \end{array}}}
\vnbprop{cc-hurwitz}{\vnbstmt{\begin{array}{@{}l@{}} \forall u.; \forall s \in (\mathbb{N} \to \operatorname{cc\text{-}hol\text{-}fns}(u)), f \in (u \to \mathbb{C}). \\ \quad \operatorname{is\text{-}open}(\operatorname{nf\text{-}metric\text{-}space}(\operatorname{cc\text{-}normed\text{-}field}), u) \wedge \\ \quad \operatorname{is\text{-}connected}(\operatorname{cc\text{-}ms}, u) \wedge \\ \quad \forall k \in \mathbb{N}, z \in u.\; \neg (s(k)(z) = 0) \wedge \\ \quad \operatorname{unif\text{-}on\text{-}compacta}(u, s, f) \Rightarrow \\ \quad \quad (\exists c \in \mathbb{C}.\; \forall \mathit{hcz} \in u.\; (f(\mathit{hcz}) = c) \vee \forall z \in u.\; \neg (f(z) = 0)) \end{array}}}
\vnbprop{cc-hurwitz-injective}{\vnbstmt{\begin{array}{@{}l@{}} \forall u.; \forall s \in (\mathbb{N} \to \operatorname{cc\text{-}hol\text{-}fns}(u)), f \in (u \to \mathbb{C}). \\ \quad \operatorname{is\text{-}open}(\operatorname{nf\text{-}metric\text{-}space}(\operatorname{cc\text{-}normed\text{-}field}), u) \wedge \\ \quad \operatorname{is\text{-}connected}(\operatorname{cc\text{-}ms}, u) \wedge \\ \quad \forall k \in \mathbb{N}, z, w \in u.\; ((s(k)(z) = s(k)(w)) \Rightarrow (z = w)) \wedge \\ \quad \operatorname{unif\text{-}on\text{-}compacta}(u, s, f) \Rightarrow \\ \quad \quad \exists c \in \mathbb{C}.\; \forall \mathit{hcz} \in u.\; (f(\mathit{hcz}) = c) \vee \\ \quad \quad \forall z, w \in u.\; ((f(z) = f(w)) \Rightarrow (z = w)) \end{array}}}
\vnbprop{cc-fn-module-is-module-of}{\vnbstmt{\begin{array}{@{}l@{}} \forall u.; \forall v \in \mathbf{Set}. \\ \quad u \subseteq \mathbb{C} \wedge \\ \quad v \subseteq (u \to \mathbb{C}) \wedge \\ \quad \forall a, b \in v.\; ((\lambda z \in u.\; (a(z) + b(z))) \in v) \wedge \\ \quad (\lambda z \in u.\; 0) \in v \wedge \\ \quad \forall a \in v.\; ((\lambda z \in u.\; -a(z)) \in v) \wedge \\ \quad \forall c \in \mathbb{C}, a \in v.\; ((\lambda z \in u.\; (c \cdot a(z))) \in v) \wedge \\ \quad \forall a, b \in v.\; ((\lambda z \in u.\; (a(z) + -b(z))) \in v) \Rightarrow \\ \quad \quad \operatorname{is\text{-}module}(\operatorname{cc\text{-}fn\text{-}module}(u, v)) \end{array}}}
\vnbprop{cc-fn-seminorm-of-difference-of}{\vnbstmt{\begin{array}{@{}l@{}} \forall u.; \forall v \in \mathbf{Set}, a, b \in v, k \in \mathbb{N}. \\ \quad u \subseteq \mathbb{C} \wedge \\ \quad v \subseteq (u \to \mathbb{C}) \wedge \\ \quad \forall a, b \in v.\; ((\lambda z \in u.\; (a(z) + b(z))) \in v) \wedge \\ \quad (\lambda z \in u.\; 0) \in v \wedge \\ \quad \forall a \in v.\; ((\lambda z \in u.\; -a(z)) \in v) \wedge \\ \quad \forall c \in \mathbb{C}, a \in v.\; ((\lambda z \in u.\; (c \cdot a(z))) \in v) \wedge \\ \quad \forall a, b \in v.\; ((\lambda z \in u.\; (a(z) + -b(z))) \in v) \Rightarrow \\ \quad \quad (\operatorname{cc\text{-}exhaust\text{-}seminorms}(u, v)(k)(\operatorname{vadd}(\operatorname{cc\text{-}fn\text{-}module}(u, v))(a, \operatorname{vneg}(\operatorname{cc\text{-}fn\text{-}module}(u, v))(b))) \simeq \operatorname{cc\text{-}sup\text{-}norm}((\lambda z \in u.\; (a(z) - b(z))), \operatorname{cc\text{-}exhaust}(u, k))) \end{array}}}
\vnbprop{cc-fr-conv-transfer}{\vnbstmt{\begin{array}{@{}l@{}} \forall u.; \forall \mathit{fva}, \mathit{fvb} \in \mathbf{Set}, s.; \forall f \in \mathit{fva}. \\ \quad u \subseteq \mathbb{C} \wedge \\ \quad \mathit{fva} \subseteq (u \to \mathbb{C}) \wedge \\ \quad \forall a, b \in \mathit{fva}.\; ((\lambda z \in u.\; (a(z) + b(z))) \in \mathit{fva}) \wedge \\ \quad (\lambda z \in u.\; 0) \in \mathit{fva} \wedge \\ \quad \forall a \in \mathit{fva}.\; ((\lambda z \in u.\; -a(z)) \in \mathit{fva}) \wedge \\ \quad \forall c \in \mathbb{C}, a \in \mathit{fva}.\; ((\lambda z \in u.\; (c \cdot a(z))) \in \mathit{fva}) \wedge \\ \quad \forall a, b \in \mathit{fva}.\; ((\lambda z \in u.\; (a(z) + -b(z))) \in \mathit{fva}) \wedge \\ \quad \mathit{fvb} \subseteq (u \to \mathbb{C}) \wedge \\ \quad \forall a, b \in \mathit{fvb}.\; ((\lambda z \in u.\; (a(z) + b(z))) \in \mathit{fvb}) \wedge \\ \quad (\lambda z \in u.\; 0) \in \mathit{fvb} \wedge \\ \quad \forall a \in \mathit{fvb}.\; ((\lambda z \in u.\; -a(z)) \in \mathit{fvb}) \wedge \\ \quad \forall c \in \mathbb{C}, a \in \mathit{fvb}.\; ((\lambda z \in u.\; (c \cdot a(z))) \in \mathit{fvb}) \wedge \\ \quad \forall a, b \in \mathit{fvb}.\; ((\lambda z \in u.\; (a(z) + -b(z))) \in \mathit{fvb}) \wedge \\ \quad \forall i \in \mathbb{N}.\; (s(i) \in \mathit{fva}) \wedge \\ \quad \forall i \in \mathbb{N}.\; (s(i) \in \mathit{fvb}) \wedge \\ \quad f \in \mathit{fvb} \wedge \\ \quad \operatorname{fr\text{-}conv}(\operatorname{cc\text{-}fn\text{-}module}(u, \mathit{fva}), \operatorname{cc\text{-}exhaust\text{-}seminorms}(u, \mathit{fva}), s, f) \Rightarrow \\ \quad \quad \operatorname{fr\text{-}conv}(\operatorname{cc\text{-}fn\text{-}module}(u, \mathit{fvb}), \operatorname{cc\text{-}exhaust\text{-}seminorms}(u, \mathit{fvb}), s, f) \end{array}}}
\vnbprop{cc-hol-fns-add}{\vnbstmt{\begin{array}{@{}l@{}} \forall u.; \forall a, b \in \operatorname{cc\text{-}hol\text{-}fns}(u). \\ \quad \operatorname{is\text{-}open}(\operatorname{nf\text{-}metric\text{-}space}(\operatorname{cc\text{-}normed\text{-}field}), u) \Rightarrow  \\ \quad \quad (\lambda z \in u.\; (a(z) + b(z))) \in \operatorname{cc\text{-}hol\text{-}fns}(u) \end{array}}}
\vnbprop{cc-hol-fns-zero}{\vnbstmt{\begin{array}{@{}l@{}} \forall u. \\ \quad \operatorname{is\text{-}open}(\operatorname{nf\text{-}metric\text{-}space}(\operatorname{cc\text{-}normed\text{-}field}), u) \Rightarrow  \\ \quad \quad (\lambda z \in u.\; 0) \in \operatorname{cc\text{-}hol\text{-}fns}(u) \end{array}}}
\vnbprop{cc-hol-fns-neg}{\vnbstmt{\begin{array}{@{}l@{}} \forall u.; \forall a \in \operatorname{cc\text{-}hol\text{-}fns}(u). \\ \quad \operatorname{is\text{-}open}(\operatorname{nf\text{-}metric\text{-}space}(\operatorname{cc\text{-}normed\text{-}field}), u) \Rightarrow  \\ \quad \quad (\lambda z \in u.\; -a(z)) \in \operatorname{cc\text{-}hol\text{-}fns}(u) \end{array}}}
\vnbprop{cc-hol-fns-scale}{\vnbstmt{\begin{array}{@{}l@{}} \forall u.; \forall c \in \mathbb{C}, a \in \operatorname{cc\text{-}hol\text{-}fns}(u). \\ \quad \operatorname{is\text{-}open}(\operatorname{nf\text{-}metric\text{-}space}(\operatorname{cc\text{-}normed\text{-}field}), u) \Rightarrow  \\ \quad \quad (\lambda z \in u.\; (c \cdot a(z))) \in \operatorname{cc\text{-}hol\text{-}fns}(u) \end{array}}}
\vnbprop{cc-hol-fns-add-neg}{\vnbstmt{\begin{array}{@{}l@{}} \forall u.; \forall a, b \in \operatorname{cc\text{-}hol\text{-}fns}(u). \\ \quad \operatorname{is\text{-}open}(\operatorname{nf\text{-}metric\text{-}space}(\operatorname{cc\text{-}normed\text{-}field}), u) \Rightarrow  \\ \quad \quad (\lambda z \in u.\; (a(z) + -b(z))) \in \operatorname{cc\text{-}hol\text{-}fns}(u) \end{array}}}
\vnbprop{cc-hol-fn-module-is-module}{\vnbstmt{\begin{array}{@{}l@{}} \forall u. \\ \quad \operatorname{is\text{-}open}(\operatorname{nf\text{-}metric\text{-}space}(\operatorname{cc\text{-}normed\text{-}field}), u) \Rightarrow  \\ \quad \quad \operatorname{is\text{-}module}(\operatorname{cc\text{-}fn\text{-}module}(u, \operatorname{cc\text{-}hol\text{-}fns}(u))) \end{array}}}
\vnbprop{cc-deriv-cont-lin}{\vnbstmt{\begin{array}{@{}l@{}} \forall u. \\ \quad \operatorname{is\text{-}open}(\operatorname{nf\text{-}metric\text{-}space}(\operatorname{cc\text{-}normed\text{-}field}), u) \Rightarrow  \\ \quad \quad \operatorname{is\text{-}cont\text{-}lin}(\operatorname{cc\text{-}fn\text{-}module}(u, \operatorname{cc\text{-}hol\text{-}fns}(u)), \\ \quad \quad \quad \operatorname{cc\text{-}exhaust\text{-}seminorms}(u, \operatorname{cc\text{-}hol\text{-}fns}(u)), \\ \quad \quad \quad \operatorname{cc\text{-}fn\text{-}module}(u, \operatorname{cc\text{-}hol\text{-}fns}(u)), \\ \quad \quad \quad \operatorname{cc\text{-}exhaust\text{-}seminorms}(u, \operatorname{cc\text{-}hol\text{-}fns}(u)), \\ \quad \quad \quad (\lambda f \in \operatorname{cc\text{-}hol\text{-}fns}(u).\; (\lambda w \in u.\; \operatorname{deriv\text{-}on}(\operatorname{cc\text{-}normed\text{-}field}, u, f, w)))) \end{array}}}

\vnbbookfile{theorem-library/ca-connected}{theorem-library/ca-connected}
\noindent{\small ca-connected.scm -- connected sets (CA-MONTEL "CONN", 2026-10-26).  No theorem of the tree concluded IS-CONNECTED of any set before this file.}\par
\vnbprop{unit-interval-clopen}{\vnbstmt{\begin{array}{@{}l@{}} \forall p, q.; \forall t \in \operatorname{ccint}(0, 1). \\ \quad \forall v \in \operatorname{ccint}(0, 1).\; ((v \in p) \vee (v \in q)) \wedge \\ \quad \forall v \in \operatorname{ccint}(0, 1).\; \neg ((v \in p) \wedge (v \in q)) \wedge \\ \quad \forall s \in \operatorname{ccint}(0, 1). \\ \quad \quad (s \in p) \Rightarrow  \\ \quad \quad \quad \exists d. \\ \quad \quad \quad \quad \operatorname{pos\text{-}rr}(d) \wedge \forall u \in \operatorname{ccint}(0, 1).\; ((\lvert (u - s) \rvert < d) \Rightarrow (u \in p)) \wedge \\ \quad \forall s \in \operatorname{ccint}(0, 1). \\ \quad \quad (s \in q) \Rightarrow  \\ \quad \quad \quad \exists d. \\ \quad \quad \quad \quad \operatorname{pos\text{-}rr}(d) \wedge \forall u \in \operatorname{ccint}(0, 1).\; ((\lvert (u - s) \rvert < d) \Rightarrow (u \in q)) \wedge \\ \quad 0 \in p \Rightarrow \\ \quad \quad t \in p \end{array}}}
\vnbprop{cc-segment-clopen}{\vnbstmt{\begin{array}{@{}l@{}} \forall u, v.; \forall x, y \in \mathbb{C}. \\ \quad \operatorname{is\text{-}open}(\operatorname{cc\text{-}ms}, u) \wedge \\ \quad \operatorname{is\text{-}open}(\operatorname{cc\text{-}ms}, v) \wedge \\ \quad \forall t \in \operatorname{ccint}(0, 1).\; (((x + (t \cdot (y - x))) \in u) \vee ((x + (t \cdot (y - x))) \in v)) \wedge \\ \quad \forall t \in \operatorname{ccint}(0, 1).\; \neg (((x + (t \cdot (y - x))) \in u) \wedge ((x + (t \cdot (y - x))) \in v)) \wedge \\ \quad x \in u \Rightarrow \\ \quad \quad y \in u \end{array}}}
\vnbprop{cc-convex-connected}{\vnbstmt{\forall u.\; (\operatorname{is\text{-}convex\text{-}cc}(u) \Rightarrow \operatorname{is\text{-}connected}(\operatorname{cc\text{-}ms}, u))}}
\vnbprop{cc-star-connected}{\vnbstmt{\begin{array}{@{}l@{}} \forall u.; \forall c \in u. \\ \quad ((u \subseteq \mathbb{C}) \wedge \forall y \in u, t \in \operatorname{ccint}(0, 1).\; ((c + (t \cdot (y - c))) \in u)) \Rightarrow  \\ \quad \quad \operatorname{is\text{-}connected}(\operatorname{cc\text{-}ms}, u) \end{array}}}
\vnbprop{unit-interval-connected}{\vnbstmt{\operatorname{is\text{-}connected}(\operatorname{rr\text{-}ms}, \operatorname{ccint}(0, 1))}}
\vnbprop{cc-ball-connected}{\vnbstmt{\forall c \in \mathbb{C}, r \in \mathbb{R}.\; \operatorname{is\text{-}connected}(\operatorname{cc\text{-}ms}, \operatorname{ball}(\operatorname{cc\text{-}ms}, c, r))}}
\vnbprop{cc-plane-connected}{\vnbstmt{\operatorname{is\text{-}connected}(\operatorname{cc\text{-}ms}, \mathbb{C})}}
\vnbprop{identity-principle-ball}{\vnbstmt{\begin{array}{@{}l@{}} \forall c \in \mathbb{C}, r \in \mathbb{R}, f.; \forall a, z \in \operatorname{ball}(\operatorname{cc\text{-}ms}, c, r). \\ \quad ((0 < r) \wedge (\operatorname{holomorphic\text{-}on}(\operatorname{ball}(\operatorname{cc\text{-}ms}, c, r), f) \wedge \operatorname{zeros\text{-}accumulate\text{-}at}(\operatorname{ball}(\operatorname{cc\text{-}ms}, c, r), f, a))) \Rightarrow  \\ \quad \quad f(z) = 0 \end{array}}}
\vnbprop{identity-principle-plane}{\vnbstmt{\begin{array}{@{}l@{}} \forall f.; \forall a, z \in \mathbb{C}. \\ \quad (\operatorname{holomorphic\text{-}on}(\mathbb{C}, f) \wedge \operatorname{zeros\text{-}accumulate\text{-}at}(\mathbb{C}, f, a)) \Rightarrow (f(z) = 0) \end{array}}}
\vnbprop{hurwitz-ball}{\vnbstmt{\begin{array}{@{}l@{}} \forall c \in \mathbb{C}, r \in \mathbb{R}, s \in (\mathbb{N} \to \operatorname{cc\text{-}hol\text{-}fns}(\operatorname{ball}(\operatorname{cc\text{-}ms}, c, r))), f \in (\operatorname{ball}(\operatorname{cc\text{-}ms}, c, r) \to \mathbb{C}). \\ \quad 0 < r \wedge \\ \quad \forall k \in \mathbb{N}, z \in \operatorname{ball}(\operatorname{cc\text{-}ms}, c, r).\; \neg (s(k)(z) = 0) \wedge \\ \quad \operatorname{unif\text{-}on\text{-}compacta}(\operatorname{ball}(\operatorname{cc\text{-}ms}, c, r), s, f) \Rightarrow \\ \quad \quad \exists \mathit{hcc} \in \mathbb{C}.\; \forall \mathit{hcz} \in \operatorname{ball}(\operatorname{cc\text{-}ms}, c, r).\; (f(\mathit{hcz}) = \mathit{hcc}) \vee \\ \quad \quad \forall z \in \operatorname{ball}(\operatorname{cc\text{-}ms}, c, r).\; \neg (f(z) = 0) \end{array}}}
\vnbprop{hurwitz-plane}{\vnbstmt{\begin{array}{@{}l@{}} \forall s \in (\mathbb{N} \to \operatorname{cc\text{-}hol\text{-}fns}(\mathbb{C})), f \in (\mathbb{C} \to \mathbb{C}). \\ \quad (\forall k \in \mathbb{N}, z \in \mathbb{C}.\; \neg (s(k)(z) = 0) \wedge \operatorname{unif\text{-}on\text{-}compacta}(\mathbb{C}, s, f)) \Rightarrow  \\ \quad \quad (\exists c \in \mathbb{C}.\; \forall \mathit{hcz} \in \mathbb{C}.\; (f(\mathit{hcz}) = c) \vee \forall z \in \mathbb{C}.\; \neg (f(z) = 0)) \end{array}}}
\vnbprop{cc-right-half-plane-convex}{\vnbstmt{\operatorname{is\text{-}convex\text{-}cc}(\{z \in \mathbb{C} : (0 < \operatorname{real\text{-}part}(z))\})}}
\vnbprop{cc-right-half-plane-connected}{\vnbstmt{\operatorname{is\text{-}connected}(\operatorname{cc\text{-}ms}, \{z \in \mathbb{C} : (0 < \operatorname{real\text{-}part}(z))\})}}
\vnbprop{cc-slit-plane-connected}{\vnbstmt{\operatorname{is\text{-}connected}(\operatorname{cc\text{-}ms}, \{u \in \mathbb{C} : \neg ((\operatorname{imag\text{-}part}(u) = 0) \wedge (\operatorname{real\text{-}part}(u) \leq 0))\})}}

\vnbbookfile{theorem-library/mvp-bounded}{theorem-library/mvp-bounded}
\noindent{\small mvp-bounded.scm -- WAVE 2 (2026-10-27..30), wired 2026-10-30 from scratchpad/calint2/mvp-bounded.scm. Library: complex-analysis.}\par
\vnbprop{cc-closed-bounded-seq-limit}{\vnbstmt{\begin{array}{@{}l@{}} \forall a, g.; \forall b \in \mathbb{R}, h \in (\mathbb{N} \to \mathbb{C}). \\ \quad a \subseteq \mathbb{C} \wedge \\ \quad \forall x \in \mathbb{C}. \\ \quad \quad \forall r.\; (\operatorname{pos\text{-}rr}(r) \Rightarrow \exists y \in a.\; (\operatorname{magnitude}((y - x)) < r)) \Rightarrow  \\ \quad \quad \quad x \in a \wedge \\ \quad \forall z \in a.\; (\operatorname{magnitude}(z) \leq b) \wedge \\ \quad g \in (a \to \mathbb{R}) \wedge \\ \quad \forall x \in a, s. \\ \quad \quad \operatorname{pos\text{-}rr}(s) \Rightarrow  \\ \quad \quad \quad \exists d. \\ \quad \quad \quad \quad (\operatorname{pos\text{-}rr}(d) \wedge \forall y \in a.\; ((\operatorname{magnitude}((x - y)) \leq d) \Rightarrow (\lvert (g(x) - g(y)) \rvert \leq s))) \wedge \\ \quad \forall n \in \mathbb{N}.\; (h(n) \in a) \Rightarrow \\ \quad \quad \exists i. \\ \quad \quad \quad \operatorname{strictly\text{-}mono\text{-}nn}(i) \wedge \\ \quad \quad \quad \exists l \in a. \forall s. \\ \quad \quad \quad \quad \operatorname{pos\text{-}rr}(s) \Rightarrow  \\ \quad \quad \quad \quad \quad \exists \mathit{mxn} \in \mathbb{N}. \forall k \in \mathbb{N}. \\ \quad \quad \quad \quad \quad \quad (\mathit{mxn} \leq k) \Rightarrow (\lvert (g(h(i(k))) - g(l)) \rvert \leq s) \end{array}}}
\vnbprop{cc-closed-bounded-max}{\vnbstmt{\begin{array}{@{}l@{}} \forall k, g.; \forall b \in \mathbb{R}. \\ \quad k \subseteq \mathbb{C} \wedge \\ \quad \forall x \in \mathbb{C}. \\ \quad \quad \forall r.\; (\operatorname{pos\text{-}rr}(r) \Rightarrow \exists y \in k.\; (\operatorname{magnitude}((y - x)) < r)) \Rightarrow  \\ \quad \quad \quad x \in k \wedge \\ \quad \forall z \in k.\; (\operatorname{magnitude}(z) \leq b) \wedge \\ \quad g \in (k \to \mathbb{R}) \wedge \\ \quad \forall x \in k, s. \\ \quad \quad \operatorname{pos\text{-}rr}(s) \Rightarrow  \\ \quad \quad \quad \exists d. \\ \quad \quad \quad \quad (\operatorname{pos\text{-}rr}(d) \wedge \forall y \in k.\; ((\operatorname{magnitude}((x - y)) \leq d) \Rightarrow (\lvert (g(x) - g(y)) \rvert \leq s))) \wedge \\ \quad \exists c.\; (c \in k) \Rightarrow \\ \quad \quad \exists a \in k.\; \forall z \in k.\; (g(z) \leq g(a)) \end{array}}}
\vnbprop{cc-continuous-at-magnitude-form}{\vnbstmt{\begin{array}{@{}l@{}} \forall w, f, x, s. \\ \quad w \subseteq \mathbb{C} \wedge \\ \quad \operatorname{is\text{-}continuous\text{-}at}(\operatorname{subspace\text{-}ms}(\operatorname{nf\text{-}metric\text{-}space}(\operatorname{cc\text{-}normed\text{-}field}), w), \\ \quad \quad \operatorname{nf\text{-}metric\text{-}space}(\operatorname{cc\text{-}normed\text{-}field}), f, x) \wedge \\ \quad \operatorname{pos\text{-}rr}(s) \Rightarrow \\ \quad \quad \exists d. \\ \quad \quad \quad (\operatorname{pos\text{-}rr}(d) \wedge \forall y \in w.\; ((\operatorname{magnitude}((x - y)) \leq d) \Rightarrow (\operatorname{magnitude}((f(x) - f(y))) \leq s))) \end{array}}}
\vnbprop{closure-cc-sequence-free-closed}{\vnbstmt{\begin{array}{@{}l@{}} \forall u.; \forall x \in \mathbb{C}. \\ \quad u \subseteq \mathbb{C} \wedge \\ \quad \forall r. \\ \quad \quad \operatorname{pos\text{-}rr}(r) \Rightarrow \exists y \in \operatorname{closure}(\operatorname{cc\text{-}ms}, u).\; (\operatorname{magnitude}((y - x)) < r) \Rightarrow \\ \quad \quad x \in \operatorname{closure}(\operatorname{cc\text{-}ms}, u) \end{array}}}
\vnbprop{closure-cc-bounded}{\vnbstmt{\begin{array}{@{}l@{}} \forall u, d.; \forall z \in \operatorname{closure}(\operatorname{cc\text{-}ms}, u). \\ \quad ((u \subseteq \mathbb{C}) \wedge ((d \in \mathbb{R}) \wedge \forall z \in u.\; (\operatorname{magnitude}(z) \leq d))) \Rightarrow  \\ \quad \quad \operatorname{magnitude}(z) \leq d \end{array}}}
\vnbprop{continuous-constant-on-closure}{\vnbstmt{\begin{array}{@{}l@{}} \forall u, f, c.; \forall x \in \operatorname{closure}(\operatorname{cc\text{-}ms}, u). \\ \quad u \subseteq \mathbb{C} \wedge \\ \quad f \in (\operatorname{closure}(\operatorname{cc\text{-}ms}, u) \to \mathbb{C}) \wedge \\ \quad \forall x \in \operatorname{closure}(\operatorname{cc\text{-}ms}, u). \\ \quad \quad \operatorname{is\text{-}continuous\text{-}at}(\operatorname{subspace\text{-}ms}(\operatorname{nf\text{-}metric\text{-}space}(\operatorname{cc\text{-}normed\text{-}field}), \operatorname{closure}(\operatorname{cc\text{-}ms}, u)), \\ \quad \quad \quad \operatorname{nf\text{-}metric\text{-}space}(\operatorname{cc\text{-}normed\text{-}field}), f, x) \wedge \\ \quad c \in \mathbb{C} \wedge \\ \quad \forall z \in u.\; (f(z) = c) \Rightarrow \\ \quad \quad f(x) = c \end{array}}}
\vnbprop{bounded-open-boundary-nonempty}{\vnbstmt{\begin{array}{@{}l@{}} \forall u.; \forall a \in u, d \in \mathbb{R}. \\ \quad (\operatorname{is\text{-}open}(\operatorname{cc\text{-}ms}, u) \wedge \forall z \in u.\; (\operatorname{magnitude}(z) \leq d)) \Rightarrow  \\ \quad \quad \exists b.\; (b \in \operatorname{difference}(\operatorname{closure}(\operatorname{cc\text{-}ms}, u), \operatorname{interior}(\operatorname{cc\text{-}ms}, u))) \end{array}}}
\vnbprop{mvp-maximum-modulus-closure}{\vnbstmt{\begin{array}{@{}l@{}} \forall u, f.; \forall d, m \in \mathbb{R}. \\ \quad \operatorname{is\text{-}connected}(\operatorname{cc\text{-}ms}, u) \wedge \\ \quad \forall z \in u.\; (\operatorname{magnitude}(z) \leq d) \wedge \\ \quad f \in (\operatorname{closure}(\operatorname{cc\text{-}ms}, u) \to \mathbb{C}) \wedge \\ \quad \forall x \in \operatorname{closure}(\operatorname{cc\text{-}ms}, u). \\ \quad \quad \operatorname{is\text{-}continuous\text{-}at}(\operatorname{subspace\text{-}ms}(\operatorname{nf\text{-}metric\text{-}space}(\operatorname{cc\text{-}normed\text{-}field}), \operatorname{closure}(\operatorname{cc\text{-}ms}, u)), \\ \quad \quad \quad \operatorname{nf\text{-}metric\text{-}space}(\operatorname{cc\text{-}normed\text{-}field}), f, x) \wedge \\ \quad \operatorname{has\text{-}mean\text{-}value\text{-}property}(u, \operatorname{restrict}(f, u)) \wedge \\ \quad \forall z \in \operatorname{difference}(\operatorname{closure}(\operatorname{cc\text{-}ms}, u), \operatorname{interior}(\operatorname{cc\text{-}ms}, u)). \\ \quad \quad \operatorname{magnitude}(f(z)) \leq m \Rightarrow \\ \quad \quad \forall w \in u.\; (\operatorname{magnitude}(f(w)) \leq m) \wedge \\ \quad \quad \forall w \in u, z \in \operatorname{closure}(\operatorname{cc\text{-}ms}, u). \\ \quad \quad \quad (\operatorname{magnitude}(f(w)) = m) \Rightarrow (f(z) = f(w)) \end{array}}}

\vnbbookfile{theorem-library/ca-montel-compact}{theorem-library/ca-montel-compact}
\noindent{\small ca-montel-compact.scm -- Theorem 5.7 of the complex analysis notes in SET form (CA-MONTEL, 2026-10-26)}\par
\vnbprop{cc-montel-compact}{\vnbstmt{\begin{array}{@{}l@{}} \forall u, a. \\ \quad \operatorname{is\text{-}open}(\operatorname{nf\text{-}metric\text{-}space}(\operatorname{cc\text{-}normed\text{-}field}), u) \wedge \\ \quad \operatorname{cc\text{-}hol\text{-}bounded}(u, a) \wedge \\ \quad \operatorname{is\text{-}closed}(\operatorname{subspace\text{-}ms}(\operatorname{fr\text{-}metric\text{-}space}(\operatorname{cc\text{-}fn\text{-}module}(u, \operatorname{cc\text{-}cont\text{-}fns}(u)), \operatorname{cc\text{-}exhaust\text{-}seminorms}(u, \operatorname{cc\text{-}cont\text{-}fns}(u))), \\ \quad \quad \quad \operatorname{cc\text{-}hol\text{-}fns}(u)), \\ \quad \quad a) \Rightarrow \\ \quad \quad \operatorname{is\text{-}compact}(\operatorname{subspace\text{-}ms}(\operatorname{fr\text{-}metric\text{-}space}(\operatorname{cc\text{-}fn\text{-}module}(u, \operatorname{cc\text{-}cont\text{-}fns}(u)), \operatorname{cc\text{-}exhaust\text{-}seminorms}(u, \operatorname{cc\text{-}cont\text{-}fns}(u))), \\ \quad \quad \quad \quad a)) \end{array}}}
\vnbprop{cc-montel-compact-in-c}{\vnbstmt{\begin{array}{@{}l@{}} \forall u, a. \\ \quad \operatorname{is\text{-}open}(\operatorname{nf\text{-}metric\text{-}space}(\operatorname{cc\text{-}normed\text{-}field}), u) \wedge \\ \quad \operatorname{cc\text{-}hol\text{-}bounded}(u, a) \wedge \\ \quad \operatorname{is\text{-}closed}(\operatorname{fr\text{-}metric\text{-}space}(\operatorname{cc\text{-}fn\text{-}module}(u, \operatorname{cc\text{-}cont\text{-}fns}(u)), \operatorname{cc\text{-}exhaust\text{-}seminorms}(u, \operatorname{cc\text{-}cont\text{-}fns}(u))), \\ \quad \quad a) \Rightarrow \\ \quad \quad \operatorname{is\text{-}compact}(\operatorname{subspace\text{-}ms}(\operatorname{fr\text{-}metric\text{-}space}(\operatorname{cc\text{-}fn\text{-}module}(u, \operatorname{cc\text{-}cont\text{-}fns}(u)), \operatorname{cc\text{-}exhaust\text{-}seminorms}(u, \operatorname{cc\text{-}cont\text{-}fns}(u))), \\ \quad \quad \quad \quad a)) \end{array}}}

\vnbbookfile{theorem-library/rb2-path}{theorem-library/rb2-path}
\noindent{\small rb2-path.scm -- path-connectedness in a metric space (vocabulary IS-PATH-CONNECTED, rb2-vocab). path-connected-implies-connected  IS-PATH-CONNECTED(s, A) =$>$ IS-CONNECTED(s, A) cc-convex-path-connected          IS-CONVEX-CC(A) =$>$ IS-PATH-CONNECTED(CC-MS, A) rr-convex-path-connected          IS-CONVEX-RR(A) =$>$ ...}\par
\vnbprop{path-connected-implies-connected}{\vnbstmt{\forall s, a.\; (\operatorname{is\text{-}path\text{-}connected}(s, a) \Rightarrow \operatorname{is\text{-}connected}(s, a))}}
\vnbprop{is-convex-rr-unfold}{\vnbstmt{\begin{array}{@{}l@{}} \forall a.; \forall x, y \in a, t \in \operatorname{ccint}(0, 1). \\ \quad \operatorname{is\text{-}convex\text{-}rr}(a) \Rightarrow ((x + (t \cdot (y - x))) \in a) \end{array}}}
\vnbprop{rr-convex-path-connected}{\vnbstmt{\forall a.\; (\operatorname{is\text{-}convex\text{-}rr}(a) \Rightarrow \operatorname{is\text{-}path\text{-}connected}(\operatorname{rr\text{-}ms}, a))}}
\vnbprop{cc-convex-path-connected}{\vnbstmt{\forall a.\; (\operatorname{is\text{-}convex\text{-}cc}(a) \Rightarrow \operatorname{is\text{-}path\text{-}connected}(\operatorname{cc\text{-}ms}, a))}}

\vnbbookfile{theorem-library/ca51-cc}{theorem-library/ca51-cc}
\noindent{\small ca51-cc.scm -- Lemma 5.1 for CC and its open subsets (CA-51, 2026-10-26). Destination: theorem-library/ca51-cc.scm, library complex-analysis.}\par
\vnbprop{cc-ms-locally-compact}{\vnbstmt{\operatorname{is\text{-}locally\text{-}compact}(\operatorname{cc\text{-}ms})}}
\vnbprop{cc-open-locally-compact}{\vnbstmt{\forall u.\; (\operatorname{is\text{-}open}(\operatorname{cc\text{-}ms}, u) \Rightarrow \operatorname{is\text{-}locally\text{-}compact}(\operatorname{subspace\text{-}ms}(\operatorname{cc\text{-}ms}, u)))}}
\vnbprop{cc-open-separable}{\vnbstmt{\begin{array}{@{}l@{}} \forall u.; \forall p \in u. \\ \quad \operatorname{is\text{-}open}(\operatorname{cc\text{-}ms}, u) \Rightarrow \operatorname{is\text{-}separable}(\operatorname{subspace\text{-}ms}(\operatorname{cc\text{-}ms}, u)) \end{array}}}
\vnbprop{cc-ms-separable}{\vnbstmt{\operatorname{is\text{-}separable}(\operatorname{cc\text{-}ms})}}
\vnbprop{cc-has-compact-exhaustion}{\vnbstmt{\exists x.\; \operatorname{is\text{-}compact\text{-}exhaustion}(\operatorname{cc\text{-}ms}, x)}}
\vnbprop{cc-open-has-compact-exhaustion}{\vnbstmt{\begin{array}{@{}l@{}} \forall u.; \forall p \in u. \\ \quad \operatorname{is\text{-}open}(\operatorname{cc\text{-}ms}, u) \Rightarrow  \\ \quad \quad \exists x.\; \operatorname{is\text{-}compact\text{-}exhaustion}(\operatorname{subspace\text{-}ms}(\operatorname{cc\text{-}ms}, u), x) \end{array}}}
\vnbprop{cc-exhaust-is-compact-exhaustion}{\vnbstmt{\begin{array}{@{}l@{}} \forall u. \\ \quad \operatorname{is\text{-}open}(\operatorname{cc\text{-}ms}, u) \Rightarrow  \\ \quad \quad \operatorname{is\text{-}compact\text{-}exhaustion}(\operatorname{subspace\text{-}ms}(\operatorname{cc\text{-}ms}, u), (\lambda n \in \mathbb{N}.\; \operatorname{cc\text{-}exhaust}(u, n))) \end{array}}}

\vnbbookfile{structure-library/ch3a-vocab}{structure-library/ch3a-vocab}
\noindent{\small ch3a-vocab.scm -- automorphisms of an open subset of CC (CH3-A, 2026-10-26). DEFINITIONS ONLY.  Destination: structure-library/ch3a-vocab.scm, library complex-analysis. Load window: after the HOLOMORPHIC-ON vocabulary (structure-library/holomorphic*) and INVERSE-BIJ / BIJECTION; before theorem-library/ch3a-schwarz.}\par
\par\medskip\noindent\textbf{Definition} (\texttt{is-holomorphic-automorphism}).\;For all $\operatorname{hau\_}$ and $\operatorname{hap\_}$, 
\begin{equation*}
\begin{array}{@{}l@{}} \operatorname{is\text{-}holomorphic\text{-}automorphism}(\operatorname{hau\_}, \operatorname{hap\_}) \\ \Leftrightarrow\; (\operatorname{hap\_} \in \mathrm{Bij}(\operatorname{hau\_}, \operatorname{hau\_})) \wedge \operatorname{holomorphic\text{-}on}(\operatorname{hau\_}, \operatorname{hap\_}) \end{array}
\end{equation*}

\par\medskip\noindent\textbf{Definition} (\texttt{cc-partial-prod}).\;For every $\operatorname{cpf\_}$,  $(\operatorname{cc\text{-}partial\text{-}prod}(\operatorname{cpf\_}, 0) \simeq 1)$\par
Let $\operatorname{cpf\_}$ be arbitrary. Let $\operatorname{cpn\_} \in \mathbb{N}$. Then
\begin{equation*}
(\operatorname{cc\text{-}partial\text{-}prod}(\operatorname{cpf\_}, \operatorname{succ}(\operatorname{cpn\_})) \simeq \cdot(\operatorname{cc\text{-}partial\text{-}prod}(\operatorname{cpf\_}, \operatorname{cpn\_}), \operatorname{cpf\_}(\operatorname{cpn\_})))
\end{equation*}


\vnbbookfile{theorem-library/ch3a-maxmod}{theorem-library/ch3a-maxmod}
\noindent{\small ch3a-maxmod.scm -- the maximum modulus principle on a connected open set (Proposition 3.55 of the complex analysis notes, holomorphic form) (CH3-A, 2026-10-26). Destination: theorem-library/ch3a-maxmod.scm, library complex-analysis.}\par
\vnbprop{ball-zeros-accumulate}{\vnbstmt{\begin{array}{@{}l@{}} \forall u, g.; \forall a \in \mathbb{C}, r. \\ \quad (\operatorname{pos\text{-}rr}(r) \wedge ((\operatorname{ball}(\operatorname{cc\text{-}ms}, a, r) \subseteq u) \wedge \forall z \in \operatorname{ball}(\operatorname{cc\text{-}ms}, a, r).\; (g(z) = 0))) \Rightarrow  \\ \quad \quad \operatorname{zeros\text{-}accumulate\text{-}at}(u, g, a) \end{array}}}
\vnbprop{maximum-modulus-connected}{\vnbstmt{\begin{array}{@{}l@{}} \forall u, f.; \forall a, z \in u. \\ \quad \operatorname{is\text{-}connected}(\operatorname{cc\text{-}ms}, u) \wedge \\ \quad \operatorname{holomorphic\text{-}on}(u, f) \wedge \\ \quad \forall z \in u.\; (\operatorname{magnitude}(f(z)) \leq \operatorname{magnitude}(f(a))) \Rightarrow \\ \quad \quad f(z) = f(a) \end{array}}}

\vnbbookfile{theorem-library/ch3a-estimates}{theorem-library/ch3a-estimates}
\noindent{\small ch3a-estimates.scm -- Cauchy's estimates with the full radius (Corollary 3.48 of the complex analysis notes) (CH3-A, 2026-10-26).  Destination: theorem-library/ch3a-estimates.scm, library complex-analysis (the real lemma rr-le-of-power-bounds is analysis vocabulary only).}\par
\vnbprop{holomorphic-nth-deriv}{\vnbstmt{\begin{array}{@{}l@{}} \forall k \in \mathbb{N}, u, f. \\ \quad \operatorname{holomorphic\text{-}on}(u, f) \Rightarrow  \\ \quad \quad \operatorname{holomorphic\text{-}on}(u, f^{(k)}) \end{array}}}
\vnbprop{rr-le-of-power-bounds}{\vnbstmt{\begin{array}{@{}l@{}} \forall k \in \mathbb{N}, v, a, o, i \in \mathbb{R}. \\ \quad ((0 \leq v) \wedge ((0 < o) \wedge ((o < i) \wedge \forall r \in \mathbb{R}.\; (((o < r) \wedge (r < i)) \Rightarrow ((v \cdot \operatorname{power}(r, k)) \leq a))))) \Rightarrow  \\ \quad \quad (v \cdot \operatorname{power}(i, k)) \leq a \end{array}}}
\vnbprop{cauchy-estimates-on-ball}{\vnbstmt{\begin{array}{@{}l@{}} \forall u, f.; \forall c \in \mathbb{C}, r, m, o \in \mathbb{R}, k \in \mathbb{N}. \\ \quad \operatorname{holomorphic\text{-}on}(u, f) \wedge \\ \quad 0 < r \wedge \\ \quad \operatorname{ball}(\operatorname{cc\text{-}ms}, c, r) \subseteq u \wedge \\ \quad \forall z \in \operatorname{ball}(\operatorname{cc\text{-}ms}, c, r).\; (\operatorname{magnitude}(f(z)) \leq m) \wedge \\ \quad 0 < o \wedge \\ \quad o < r \Rightarrow \\ \quad \quad \operatorname{magnitude}(\operatorname{restrict}(f, \operatorname{ball}(\operatorname{cc\text{-}ms}, c, o))^{(k)}(c)) \\ \quad \quad \leq\; \operatorname{factorial}(k) \cdot (m \cdot \operatorname{power}(\operatorname{recip}(r), k)) \end{array}}}

\vnbbookfile{theorem-library/ch3a-schwarz}{theorem-library/ch3a-schwarz}
\noindent{\small ch3a-schwarz.scm -- the maximum of |g| on a compact set, removable points, and Schwarz's lemma (Proposition 3.57 and Lemma 10.11 of the complex analysis notes) (CH3-A, 2026-10-26). Destination: theorem-library/ch3a-schwarz.scm, library complex-analysis. Helper prefix csw-.}\par
\vnbprop{cc-continuous-at-modulus}{\vnbstmt{\begin{array}{@{}l@{}} \forall u, g.; \forall y \in u, s. \\ \quad u \subseteq \mathbb{C} \wedge \\ \quad \operatorname{is\text{-}continuous\text{-}at}(\operatorname{subspace\text{-}ms}(\operatorname{nf\text{-}metric\text{-}space}(\operatorname{cc\text{-}normed\text{-}field}), u), \\ \quad \quad \operatorname{nf\text{-}metric\text{-}space}(\operatorname{cc\text{-}normed\text{-}field}), g, y) \wedge \\ \quad \operatorname{pos\text{-}rr}(s) \Rightarrow \\ \quad \quad \exists l. \\ \quad \quad \quad (\operatorname{pos\text{-}rr}(l) \wedge \forall b \in u.\; ((\operatorname{magnitude}((b - y)) < l) \Rightarrow (\operatorname{magnitude}((g(b) - g(y))) \leq s))) \end{array}}}
\vnbprop{cc-compact-modulus-max}{\vnbstmt{\begin{array}{@{}l@{}} \forall u.; \forall g \in (u \to \mathbb{C}), k.; \forall a \in k. \\ \quad u \subseteq \mathbb{C} \wedge \\ \quad \forall y \in u. \\ \quad \quad \operatorname{is\text{-}continuous\text{-}at}(\operatorname{subspace\text{-}ms}(\operatorname{nf\text{-}metric\text{-}space}(\operatorname{cc\text{-}normed\text{-}field}), u), \\ \quad \quad \quad \operatorname{nf\text{-}metric\text{-}space}(\operatorname{cc\text{-}normed\text{-}field}), g, y) \wedge \\ \quad k \subseteq u \wedge \\ \quad \operatorname{is\text{-}compact}(\operatorname{subspace\text{-}ms}(\operatorname{cc\text{-}ms}, k)) \Rightarrow \\ \quad \quad \exists p \in k.\; \forall z \in k.\; (\operatorname{magnitude}(g(z)) \leq \operatorname{magnitude}(g(p))) \end{array}}}
\vnbprop{holomorphic-removable-point}{\vnbstmt{\begin{array}{@{}l@{}} \forall u.; \forall g \in (u \to \mathbb{C}), a \in u. \\ \quad \operatorname{is\text{-}open}(\operatorname{nf\text{-}metric\text{-}space}(\operatorname{cc\text{-}normed\text{-}field}), u) \wedge \\ \quad \forall y \in u. \\ \quad \quad \operatorname{is\text{-}continuous\text{-}at}(\operatorname{subspace\text{-}ms}(\operatorname{nf\text{-}metric\text{-}space}(\operatorname{cc\text{-}normed\text{-}field}), u), \\ \quad \quad \quad \operatorname{nf\text{-}metric\text{-}space}(\operatorname{cc\text{-}normed\text{-}field}), g, y) \wedge \\ \quad \operatorname{holomorphic\text{-}on}(\operatorname{difference}(u, \operatorname{singleton}(a)), \operatorname{restrict}(g, \operatorname{difference}(u, \operatorname{singleton}(a)))) \Rightarrow \\ \quad \quad \operatorname{holomorphic\text{-}on}(u, g) \end{array}}}
\vnbprop{schwarz-quotient}{\vnbstmt{\begin{array}{@{}l@{}} \forall f. \\ \quad \operatorname{holomorphic\text{-}on}(\operatorname{ball}(\operatorname{cc\text{-}ms}, 0, 1), f) \wedge \\ \quad f(0) = 0 \wedge \\ \quad \forall z \in \operatorname{ball}(\operatorname{cc\text{-}ms}, 0, 1).\; (\operatorname{magnitude}(f(z)) < 1) \Rightarrow \\ \quad \quad \exists q. \\ \quad \quad \quad \operatorname{holomorphic\text{-}on}(\operatorname{ball}(\operatorname{cc\text{-}ms}, 0, 1), q) \wedge \\ \quad \quad \quad \forall w \in \operatorname{ball}(\operatorname{cc\text{-}ms}, 0, 1).\; (f(w) = (q(w) \cdot w)) \wedge \\ \quad \quad \quad \forall w \in \operatorname{ball}(\operatorname{cc\text{-}ms}, 0, 1).\; (\operatorname{magnitude}(q(w)) \leq 1) \end{array}}}
\vnbprop{schwarz-lemma}{\vnbstmt{\begin{array}{@{}l@{}} \forall f. \\ \quad \operatorname{holomorphic\text{-}on}(\operatorname{ball}(\operatorname{cc\text{-}ms}, 0, 1), f) \wedge \\ \quad f(0) = 0 \wedge \\ \quad \forall z \in \operatorname{ball}(\operatorname{cc\text{-}ms}, 0, 1).\; (\operatorname{magnitude}(f(z)) < 1) \Rightarrow \\ \quad \quad \forall z \in \operatorname{ball}(\operatorname{cc\text{-}ms}, 0, 1).\; (\operatorname{magnitude}(f(z)) \leq \operatorname{magnitude}(z)) \wedge \\ \quad \quad \exists a \in \operatorname{ball}(\operatorname{cc\text{-}ms}, 0, 1).\; (\neg (a = 0) \wedge (\operatorname{magnitude}(f(a)) = \operatorname{magnitude}(a))) \Rightarrow  \\ \quad \quad \quad \exists r \in \mathbb{C}. \\ \quad \quad \quad \quad (\operatorname{magnitude}(r) = 1) \wedge \forall z \in \operatorname{ball}(\operatorname{cc\text{-}ms}, 0, 1).\; (f(z) = (r \cdot z)) \end{array}}}
\vnbprop{disc-automorphism-rotation}{\vnbstmt{\begin{array}{@{}l@{}} \forall i. \\ \quad \operatorname{is\text{-}holomorphic\text{-}automorphism}(\operatorname{ball}(\operatorname{cc\text{-}ms}, 0, 1), i) \wedge \\ \quad \operatorname{holomorphic\text{-}on}(\operatorname{ball}(\operatorname{cc\text{-}ms}, 0, 1), i^{-1}) \wedge \\ \quad i(0) = 0 \Rightarrow \\ \quad \quad \exists r \in \mathbb{C}. \\ \quad \quad \quad (\operatorname{magnitude}(r) = 1) \wedge \forall z \in \operatorname{ball}(\operatorname{cc\text{-}ms}, 0, 1).\; (i(z) = (r \cdot z)) \end{array}}}

\vnbbookfile{theorem-library/ch3a-growth}{theorem-library/ch3a-growth}
\noindent{\small ch3a-growth.scm -- entire functions of polynomial growth are polynomials (Proposition 3.53 of the complex analysis notes) (CH3-A, 2026-10-26).  Destination: theorem-library/ch3a-growth.scm, library complex-analysis.}\par
\vnbprop{entire-polynomial-growth}{\vnbstmt{\begin{array}{@{}l@{}} \forall m \in \mathbb{N}, f.; \forall a, b \in \mathbb{R}. \\ \quad (\operatorname{holomorphic\text{-}on}(\mathbb{C}, f) \wedge \forall z \in \mathbb{C}.\; (\operatorname{magnitude}(f(z)) \leq ((a \cdot \operatorname{power}(\operatorname{magnitude}(z), m)) + b))) \Rightarrow  \\ \quad \quad \exists c \in (\mathbb{N} \to \mathbb{C}). \forall z \in \mathbb{C}. \\ \quad \quad \quad f(z) = \operatorname{cc\text{-}series\text{-}partial\text{-}sum}((\lambda j \in \mathbb{N}.\; (c(j) \cdot \operatorname{power}(z, j))), (m + 1)) \end{array}}}

\vnbbookfile{theorem-library/ch3a-factor}{theorem-library/ch3a-factor}
\noindent{\small ch3a-factor.scm -- complete factorization of a complex polynomial (Proposition 3.51 of the complex analysis notes) (CH3-A, 2026-10-26).  Destination: theorem-library/ch3a-factor.scm, library complex-analysis.}\par
\vnbprop{cc-partial-prod-congr}{\vnbstmt{\begin{array}{@{}l@{}} \forall n \in \mathbb{N}, g, h. \\ \quad \forall i \in \mathbb{N}.\; ((i < n) \Rightarrow (g(i) = h(i))) \Rightarrow  \\ \quad \quad (\operatorname{cc\text{-}partial\text{-}prod}(g, n) \simeq \operatorname{cc\text{-}partial\text{-}prod}(h, n)) \end{array}}}
\vnbprop{cc-partial-prod-in-cc}{\vnbstmt{\begin{array}{@{}l@{}} \forall n \in \mathbb{N}, g. \\ \quad \forall i \in \mathbb{N}.\; (g(i) \in \mathbb{C}) \Rightarrow (\operatorname{cc\text{-}partial\text{-}prod}(g, n) \in \mathbb{C}) \end{array}}}
\vnbprop{cc-poly-factorization}{\vnbstmt{\begin{array}{@{}l@{}} \forall n \in \mathbb{N}, c \in (\mathbb{N} \to \mathbb{C}). \\ \quad (\forall k \in \mathbb{N}.\; (((n + 1) \leq k) \Rightarrow (c(k) = 0)) \wedge \neg (c(n) = 0)) \Rightarrow  \\ \quad \quad \exists a \in (\mathbb{N} \to \mathbb{C}). \forall z \in \mathbb{C}. \\ \quad \quad \quad \operatorname{cc\text{-}series\text{-}partial\text{-}sum}((\lambda j \in \mathbb{N}.\; (c(j) \cdot \operatorname{power}(z, j))), (n + 1)) \\ \quad \quad \quad =\; c(n) \cdot \operatorname{cc\text{-}partial\text{-}prod}((\lambda i \in \mathbb{N}.\; (z - a(i))), n) \end{array}}}

\vnbbookfile{structure-library/algnum-vocab}{structure-library/algnum-vocab}
\noindent{\small algnum-vocab.scm -- ALGEBRAIC ELEMENTS OF C (ALGNUM, algebraic-numbers bundle B7, 2026-10-26).  DEFINITIONS ONLY; destination structure-library/.  Source: Thayer, algebraic-numbers notes, ch. 4 s. 3-4, Def. 4.15 and Def. 4.20. Polynomials are written as POLY-2 writes them (poly-fn-vocab.scm): a coefficient sequence c ...}\par
\par\medskip\noindent\textbf{Definition} (\texttt{is-algebraic-over}).\;For all $\mathit{aa}$ and $r$, 
\begin{equation*}
\begin{array}{@{}l@{}} \operatorname{is\text{-}algebraic\text{-}over}(\mathit{aa}, r) \\ \Leftrightarrow\; r \in \mathbb{C} \wedge \\ \exists \operatorname{anan\_} \in \mathbb{N}, \operatorname{anac\_} \in (\mathbb{N} \to \mathit{aa}). \\ \quad \forall \operatorname{anak\_} \in \mathbb{N}.\; ((\operatorname{succ}(\operatorname{anan\_}) \leq \operatorname{anak\_}) \Rightarrow (\operatorname{anac\_}(\operatorname{anak\_}) = 0)) \wedge \\ \quad \operatorname{anac\_}(\operatorname{anan\_}) = 1 \wedge \\ \quad \operatorname{cc\text{-}series\text{-}partial\text{-}sum}((\lambda \operatorname{pfj\_} \in \mathbb{N}.\; \cdot(\operatorname{anac\_}(\operatorname{pfj\_}), \operatorname{power}(r, \operatorname{pfj\_}))), \operatorname{succ}(\operatorname{anan\_})) = 0 \end{array}
\end{equation*}

\par\medskip\noindent\textbf{Definition} (\texttt{is-algebraic-over-field}).\;For all $\mathit{kk}$ and $\mathit{th}$, 
\begin{equation*}
\begin{array}{@{}l@{}} \operatorname{is\text{-}algebraic\text{-}over\text{-}field}(\mathit{kk}, \mathit{th}) \\ \Leftrightarrow\; \mathit{th} \in \mathbb{C} \wedge \\ \exists \operatorname{anan\_} \in \mathbb{N}, \operatorname{anac\_} \in (\mathbb{N} \to \mathit{kk}). \\ \quad \forall \operatorname{anak\_} \in \mathbb{N}.\; ((\operatorname{succ}(\operatorname{anan\_}) \leq \operatorname{anak\_}) \Rightarrow (\operatorname{anac\_}(\operatorname{anak\_}) = 0)) \wedge \\ \quad \neg (\operatorname{anac\_}(\operatorname{anan\_}) = 0) \wedge \\ \quad \operatorname{cc\text{-}series\text{-}partial\text{-}sum}((\lambda \operatorname{pfj\_} \in \mathbb{N}.\; \cdot(\operatorname{anac\_}(\operatorname{pfj\_}), \operatorname{power}(\mathit{th}, \operatorname{pfj\_}))), \operatorname{succ}(\operatorname{anan\_})) = 0 \end{array}
\end{equation*}

\par\medskip\noindent\textbf{Definition} (\texttt{algebraic-numbers-over}).\;$\operatorname{algebraic\text{-}numbers\text{-}over}(\mathit{kk}) :=$ $\{\operatorname{anth\_} \in \mathbb{C} : \operatorname{is\text{-}algebraic\text{-}over\text{-}field}(\mathit{kk}, \operatorname{anth\_})\}$
\par\medskip\noindent\textbf{Definition} (\texttt{is-subring-of-cc}).\;For every $\mathit{aa}$, 
\begin{equation*}
\operatorname{is\text{-}subring\text{-}of\text{-}cc}(\mathit{aa}) \Leftrightarrow (\operatorname{is\text{-}cc\text{-}subring}(\mathit{aa}) \wedge \forall \operatorname{ansn\_} \in \mathit{aa}.\; (-(\operatorname{ansn\_}) \in \mathit{aa}))
\end{equation*}

\par\medskip\noindent\textbf{Definition} (\texttt{cc-lincomb}).\;$\operatorname{cc\text{-}lincomb}(n, b, g) :=$
\begin{equation*}
\begin{array}{@{}l@{}} \operatorname{finsum}(\operatorname{ring\text{-}additive\text{-}ag}(\operatorname{normed\text{-}field\text{-}as\text{-}commutative\text{-}ring}(\operatorname{cc\text{-}normed\text{-}field})), \\ \quad (\lambda \operatorname{almj\_} \in \operatorname{interval}(1, n).\; (b(\operatorname{almj\_}) \cdot g(\operatorname{almj\_}))), \\ \quad \operatorname{interval}(1, n)) \end{array}
\end{equation*}

\par\medskip\noindent\textbf{Definition} (\texttt{is-cc-module-over}).\;For all $\mathit{aa}$ and $\mathit{vv}$, 
\begin{equation*}
\begin{array}{@{}l@{}} \operatorname{is\text{-}cc\text{-}module\text{-}over}(\mathit{aa}, \mathit{vv}) \\ \Leftrightarrow\; \mathit{vv} \subseteq \mathbb{C} \wedge \\ 0 \in \mathit{vv} \wedge \\ \forall \operatorname{amx\_}, \operatorname{amy\_} \in \mathit{vv}.\; (+(\operatorname{amx\_}, \operatorname{amy\_}) \in \mathit{vv}) \wedge \\ \forall \operatorname{amx\_} \in \mathit{vv}.\; (-(\operatorname{amx\_}) \in \mathit{vv}) \wedge \\ \forall \operatorname{ama\_} \in \mathit{aa}, \operatorname{amx\_} \in \mathit{vv}.\; (\cdot(\operatorname{ama\_}, \operatorname{amx\_}) \in \mathit{vv}) \end{array}
\end{equation*}

\par\medskip\noindent\textbf{Definition} (\texttt{is-fg-cc-module-over}).\;For all $\mathit{aa}$ and $\mathit{vv}$, 
\begin{equation*}
\begin{array}{@{}l@{}} \operatorname{is\text{-}fg\text{-}cc\text{-}module\text{-}over}(\mathit{aa}, \mathit{vv}) \\ \Leftrightarrow\; \operatorname{is\text{-}cc\text{-}module\text{-}over}(\mathit{aa}, \mathit{vv}) \wedge \\ \exists \operatorname{amn\_} \in \mathbb{N}, \operatorname{amg\_} \in (\operatorname{interval}(1, \operatorname{amn\_}) \to \mathit{vv}). \forall \operatorname{amv\_} \in \mathit{vv}. \exists \operatorname{amb\_} \in (\operatorname{interval}(1, \operatorname{amn\_}) \to \mathit{aa}). \\ \quad \operatorname{amv\_} = \operatorname{cc\text{-}lincomb}(\operatorname{amn\_}, \operatorname{amb\_}, \operatorname{amg\_}) \end{array}
\end{equation*}

\par\medskip\noindent\textbf{Definition} (\texttt{cc-fraction-field}).\;$\operatorname{cc\text{-}fraction\text{-}field}(\mathit{aa}) :=$
\begin{equation*}
\{\operatorname{anfq\_} \in \mathbb{C} : \exists \operatorname{anfa\_}, \operatorname{anfb\_} \in \mathit{aa}.\; (\neg (\operatorname{anfb\_} = 0) \wedge ((\operatorname{anfq\_} \cdot \operatorname{anfb\_}) = \operatorname{anfa\_}))\}
\end{equation*}

\par\medskip\noindent\textbf{Definition} (\texttt{is-subfield-of-cc}).\;For every $\mathit{kk}$, 
\begin{equation*}
\begin{array}{@{}l@{}} \operatorname{is\text{-}subfield\text{-}of\text{-}cc}(\mathit{kk}) \\ \Leftrightarrow\; (\operatorname{is\text{-}subring\text{-}of\text{-}cc}(\mathit{kk}) \wedge \forall \operatorname{ansf\_} \in \mathit{kk}.\; (\neg (\operatorname{ansf\_} = 0) \Rightarrow (\operatorname{recip}(\operatorname{ansf\_}) \in \mathit{kk}))) \end{array}
\end{equation*}

\par\medskip\noindent\textbf{Definition} (\texttt{is-cc-space-dim-le}).\;For all $\mathit{kk}$, $\mathit{vv}$ and $n$, 
\begin{equation*}
\begin{array}{@{}l@{}} \operatorname{is\text{-}cc\text{-}space\text{-}dim\text{-}le}(\mathit{kk}, \mathit{vv}, n) \\ \Leftrightarrow\; \operatorname{is\text{-}cc\text{-}module\text{-}over}(\mathit{kk}, \mathit{vv}) \wedge \\ \exists \operatorname{amn\_} \in \mathbb{N}. \\ \quad \operatorname{amn\_} \leq n \wedge \\ \quad \exists \operatorname{amg\_} \in (\operatorname{interval}(1, \operatorname{amn\_}) \to \mathit{vv}). \forall \operatorname{amv\_} \in \mathit{vv}. \exists \operatorname{amb\_} \in (\operatorname{interval}(1, \operatorname{amn\_}) \to \mathit{kk}). \\ \quad \quad \operatorname{amv\_} = \operatorname{cc\text{-}lincomb}(\operatorname{amn\_}, \operatorname{amb\_}, \operatorname{amg\_}) \end{array}
\end{equation*}

\par\medskip\noindent\textbf{Definition} (\texttt{is-algebraic-of-degree-le}).\;For all $\mathit{kk}$, $\mathit{th}$ and $n$, 
\begin{equation*}
\begin{array}{@{}l@{}} \operatorname{is\text{-}algebraic\text{-}of\text{-}degree\text{-}le}(\mathit{kk}, \mathit{th}, n) \\ \Leftrightarrow\; \mathit{th} \in \mathbb{C} \wedge \\ \exists \operatorname{anan\_} \in \mathbb{N}. \\ \quad \operatorname{anan\_} \leq n \wedge \\ \quad \exists \operatorname{anac\_} \in (\mathbb{N} \to \mathit{kk}). \\ \quad \quad \forall \operatorname{anak\_} \in \mathbb{N}.\; ((\operatorname{succ}(\operatorname{anan\_}) \leq \operatorname{anak\_}) \Rightarrow (\operatorname{anac\_}(\operatorname{anak\_}) = 0)) \wedge \\ \quad \quad \neg (\operatorname{anac\_}(\operatorname{anan\_}) = 0) \wedge \\ \quad \quad \operatorname{cc\text{-}series\text{-}partial\text{-}sum}((\lambda \operatorname{pfj\_} \in \mathbb{N}.\; \cdot(\operatorname{anac\_}(\operatorname{pfj\_}), \operatorname{power}(\mathit{th}, \operatorname{pfj\_}))), \operatorname{succ}(\operatorname{anan\_})) = 0 \end{array}
\end{equation*}


\vnbbookfile{structure-library/ca-lft-vocab}{structure-library/ca-lft-vocab}
\noindent{\small ca-lft-vocab.scm -- THE EXTENDED PLANE AND LINEAR FRACTIONAL TRANSFORMATIONS, complex-analysis.pdf s. 1.6 (eq. (8), the topology before Lemma 1.16, Def. 1.17, (9), (11)). DEFINITIONS ONLY.  CA-LFT, 2026-10-19.  Library: complex-analysis.}\par
\par\medskip\noindent\textbf{Definition} (\texttt{cc-hat}).\; $(\operatorname{cc\text{-}hat} \simeq (\mathbb{C} \cup \operatorname{singleton}(\operatorname{cc\text{-}inf})))$
\par\medskip\noindent\textbf{Definition} (\texttt{cc-hat-opens}).\;
\begin{equation*}
(\operatorname{cc\text{-}hat\text{-}opens} \simeq \{\operatorname{uh\_} \in \operatorname{power}(\operatorname{cc\text{-}hat}) : (((\operatorname{uh\_} \subseteq \mathbb{C}) \wedge \operatorname{is\text{-}open}(\operatorname{cc\text{-}ms}, \operatorname{uh\_})) \vee \exists \operatorname{kh\_} \in \operatorname{power}(\mathbb{C}).\; ((\operatorname{kh\_} \subseteq \mathbb{C}) \wedge (\operatorname{is\text{-}compact}(\operatorname{subspace\text{-}ms}(\operatorname{cc\text{-}ms}, \operatorname{kh\_})) \wedge (\operatorname{uh\_} = \operatorname{difference}(\operatorname{cc\text{-}hat}, \operatorname{kh\_})))))\})
\end{equation*}

\par\medskip\noindent\textbf{Definition} (\texttt{cc-hat-top}).\; $\operatorname{cc\text{-}hat\text{-}top} = {[}\operatorname{cc\text{-}hat}, \operatorname{cc\text{-}hat\text{-}opens}]$
\par\medskip\noindent\textbf{Definition} (\texttt{top-converges-to}).\;For all $t$, $f$ and $l$, 
\begin{equation*}
\begin{array}{@{}l@{}} \operatorname{top\text{-}converges\text{-}to}(t, f, l) \\ \Leftrightarrow\; f \in (\mathbb{N} \to \operatorname{pts}(t)) \wedge \\ l \in \operatorname{pts}(t) \wedge \\ \forall \operatorname{ut\_} \in \operatorname{opens}(t). \\ \quad (l \in \operatorname{ut\_}) \Rightarrow  \\ \quad \quad \exists \operatorname{nt\_} \in \mathbb{N}.\; \forall \operatorname{mt\_} \in \mathbb{N}.\; ((\operatorname{nt\_} \leq \operatorname{mt\_}) \Rightarrow (f(\operatorname{mt\_}) \in \operatorname{ut\_})) \end{array}
\end{equation*}

\par\medskip\noindent\textbf{Definition} (\texttt{lft}).\;$\operatorname{lft}(a, b, c, d) :=$
\begin{equation*}
\begin{array}{@{}l@{}} (\lambda \operatorname{lfz\_} \in \operatorname{cc\text{-}hat}.\; \\ \quad \left(\text{if }(\operatorname{lfz\_} = \operatorname{cc\text{-}inf})\text{ then }\left(\text{if }(c = 0)\text{ then }\operatorname{cc\text{-}inf}\text{ else }(a \cdot \operatorname{recip}(c))\right)\text{ else }\left(\text{if }(((c \cdot \operatorname{lfz\_}) + d) = 0)\text{ then }\operatorname{cc\text{-}inf}\text{ else }(((a \cdot \operatorname{lfz\_}) + b) \cdot \operatorname{recip}(((c \cdot \operatorname{lfz\_}) + d)))\right)\right)) \end{array}
\end{equation*}

\par\medskip\noindent\textbf{Definition} (\texttt{homog-point}).\;$\operatorname{homog\text{-}point}(u, v) :=$ $\left(\text{if }(v = 0)\text{ then }\operatorname{cc\text{-}inf}\text{ else }(u \cdot \operatorname{recip}(v))\right)$
\par\medskip\noindent\textbf{Definition} (\texttt{lfts}).\;
\begin{equation*}
(\mathit{lfts} \simeq \{\operatorname{lff\_} \in (\operatorname{cc\text{-}hat} \to \operatorname{cc\text{-}hat}) : \exists \operatorname{lfa\_}, \operatorname{lfb\_}, \operatorname{lfc\_}, \operatorname{lfd\_} \in \mathbb{C}.\; (\neg (-(\cdot(\operatorname{lfa\_}, \operatorname{lfd\_}), \cdot(\operatorname{lfb\_}, \operatorname{lfc\_})) = 0) \wedge (\operatorname{lff\_} = \operatorname{lft}(\operatorname{lfa\_}, \operatorname{lfb\_}, \operatorname{lfc\_}, \operatorname{lfd\_})))\})
\end{equation*}

\par\medskip\noindent\textbf{Definition} (\texttt{lft-group}).\;
\begin{equation*}
\begin{array}{@{}l@{}} \operatorname{lft\text{-}group} \\ =\; {[}\mathit{lfts}, (\lambda {[}\operatorname{lfg\_}, \operatorname{lfh\_}] \in (\mathit{lfts} \times \mathit{lfts}).\; \operatorname{compose}(\operatorname{lfg\_}, \operatorname{lfh\_})), \operatorname{id\text{-}fun}(\operatorname{cc\text{-}hat}), (\lambda \operatorname{lfg\_} \in \mathit{lfts}.\; \operatorname{lfg\_}^{-1})] \end{array}
\end{equation*}

\par\medskip\noindent\textbf{Definition} (\texttt{mat-lft}).\;$\operatorname{mat\text{-}lft}(s) :=$
\begin{equation*}
\operatorname{lft}(\operatorname{entry}(s, 1, 1), \operatorname{entry}(s, 1, 2), \operatorname{entry}(s, 2, 1), \operatorname{entry}(s, 2, 2))
\end{equation*}

\par\medskip\noindent\textbf{Definition} (\texttt{cayley}).\; $(\mathit{cayley} \simeq \operatorname{lft}(1, -(+i), 1, +i))$
\par\medskip\noindent\textbf{Definition} (\texttt{upper-half-plane}).\; $(\operatorname{upper\text{-}half\text{-}plane} \simeq \{\operatorname{hpz\_} \in \mathbb{C} : (0 < \operatorname{imag\text{-}part}(\operatorname{hpz\_}))\})$
\par\medskip\noindent\textbf{Definition} (\texttt{unit-circle}).\; $(\operatorname{unit\text{-}circle} \simeq \{\operatorname{ucz\_} \in \mathbb{C} : (\operatorname{magnitude}(\operatorname{ucz\_}) = 1)\})$
\vnbdef{lft-group@carr}{\vnbstmt{(\operatorname{carr}(\operatorname{lft\text{-}group}) \simeq \mathit{lfts})}}
\vnbdef{cc-hat-top@opens}{\vnbstmt{(\operatorname{opens}(\operatorname{cc\text{-}hat\text{-}top}) \simeq \operatorname{cc\text{-}hat\text{-}opens})}}
\vnbdef{lft-group@inv}{\vnbstmt{(\operatorname{inv}(\operatorname{lft\text{-}group}) \simeq (\lambda g \in \mathit{lfts}.\; g^{-1}))}}
\vnbdef{cc-hat-top@pts}{\vnbstmt{(\operatorname{pts}(\operatorname{cc\text{-}hat\text{-}top}) \simeq \operatorname{cc\text{-}hat})}}
\vnbdef{lft-group@opr}{\vnbstmt{(\operatorname{opr}(\operatorname{lft\text{-}group}) \simeq (\lambda {[}\operatorname{lfg\_}, \operatorname{lfh\_}] \in (\mathit{lfts} \times \mathit{lfts}).\; \operatorname{compose}(\operatorname{lfg\_}, \operatorname{lfh\_})))}}
\vnbdef{lft-group@iden}{\vnbstmt{(\operatorname{iden}(\operatorname{lft\text{-}group}) \simeq \operatorname{id\text{-}fun}(\operatorname{cc\text{-}hat}))}}
\vnbaxiom{cc-inf-not-in-cc}{\vnbstmt{\neg (\operatorname{cc\text{-}inf} \in \mathbb{C})}}
\vnbaxiom{cc-inf-in-set}{\vnbstmt{\operatorname{cc\text{-}inf} \in \mathbf{Set}}}

\vnbbookfile{theorem-library/ca-lft-plane}{theorem-library/ca-lft-plane}
\noindent{\small ca-lft-plane.scm -- the point at infinity, the extended plane C u \{oo\}, and the values of LFT(a, b, c, d) and HOMOG-POINT(u, v) (complex-analysis.pdf s. 1.6, eq. (8) and (9)). CA-LFT, 2026-10-19.  Library: complex-analysis.  Vocabulary: structure-library/ca-lft-vocab.scm.}\par
\vnbprop{cc-ne-inf}{\vnbstmt{\forall z \in \mathbb{C}.\; \neg (z = \operatorname{cc\text{-}inf})}}
\vnbprop{cc-hat-membership}{\vnbstmt{\forall x.\; ((x \in \operatorname{cc\text{-}hat}) \Leftrightarrow ((x \in \mathbb{C}) \vee (x = \operatorname{cc\text{-}inf})))}}
\vnbprop{cc-subset-cc-hat}{\vnbstmt{\mathbb{C} \subseteq \operatorname{cc\text{-}hat}}}
\vnbprop{cc-inf-in-cc-hat}{\vnbstmt{\operatorname{cc\text{-}inf} \in \operatorname{cc\text{-}hat}}}
\vnbprop{cc-in-cc-hat}{\vnbstmt{\forall z \in \mathbb{C}.\; (z \in \operatorname{cc\text{-}hat})}}
\vnbprop{cc-hat-in-set}{\vnbstmt{\operatorname{cc\text{-}hat} \in \mathbf{Set}}}
\vnbprop{cc-div-closed}{\vnbstmt{\forall x, y \in \mathbb{C}.\; (\neg (y = 0) \Rightarrow ((x \cdot \operatorname{recip}(y)) \in \mathbb{C}))}}
\vnbprop{lft-at-inf-pole}{\vnbstmt{\begin{array}{@{}l@{}} \forall a, b, c, d \in \mathbb{C}. \\ \quad (c = 0) \Rightarrow (\operatorname{lft}(a, b, c, d)(\operatorname{cc\text{-}inf}) = \operatorname{cc\text{-}inf}) \end{array}}}
\vnbprop{lft-at-inf}{\vnbstmt{\begin{array}{@{}l@{}} \forall a, b, c, d \in \mathbb{C}. \\ \quad \neg (c = 0) \Rightarrow (\operatorname{lft}(a, b, c, d)(\operatorname{cc\text{-}inf}) = (a \cdot \operatorname{recip}(c))) \end{array}}}
\vnbprop{lft-at-pole}{\vnbstmt{\begin{array}{@{}l@{}} \forall a, b, c, d, z \in \mathbb{C}. \\ \quad (((c \cdot z) + d) = 0) \Rightarrow (\operatorname{lft}(a, b, c, d)(z) = \operatorname{cc\text{-}inf}) \end{array}}}
\vnbprop{lft-at-finite}{\vnbstmt{\begin{array}{@{}l@{}} \forall a, b, c, d, z \in \mathbb{C}. \\ \quad \neg (((c \cdot z) + d) = 0) \Rightarrow (\operatorname{lft}(a, b, c, d)(z) = (((a \cdot z) + b) \cdot \operatorname{recip}(((c \cdot z) + d)))) \end{array}}}
\vnbprop{homog-point-at-inf}{\vnbstmt{\forall u \in \mathbb{C}.\; (\operatorname{homog\text{-}point}(u, 0) = \operatorname{cc\text{-}inf})}}
\vnbprop{homog-point-finite}{\vnbstmt{\forall u, v \in \mathbb{C}.\; (\neg (v = 0) \Rightarrow (\operatorname{homog\text{-}point}(u, v) = (u \cdot \operatorname{recip}(v))))}}
\vnbprop{homog-point-in-cc-hat}{\vnbstmt{\forall u, v \in \mathbb{C}.\; (\operatorname{homog\text{-}point}(u, v) \in \operatorname{cc\text{-}hat})}}
\vnbprop{lft-in-fun}{\vnbstmt{\begin{array}{@{}l@{}} \forall a, b, c, d \in \mathbb{C}. \\ \quad \operatorname{lft}(a, b, c, d) \in (\operatorname{cc\text{-}hat} \to \operatorname{cc\text{-}hat}) \end{array}}}

\vnbbookfile{theorem-library/ca-lft-compose}{theorem-library/ca-lft-compose}
\noindent{\small ca-lft-compose.scm -- homogeneous coordinates and the composition of linear fractional transformations (complex-analysis.pdf Prop. 1.18, the computation in its proof). CA-LFT, 2026-10-19.  Library: complex-analysis.  Vocabulary: structure-library/ca-lft-vocab.scm.}\par
\vnbprop{cc-mul-zero-left}{\vnbstmt{\forall x, y \in \mathbb{C}.\; ((((x \cdot y) = 0) \wedge \neg (y = 0)) \Rightarrow (x = 0))}}
\vnbprop{cc-div-mul-back}{\vnbstmt{\forall u, v \in \mathbb{C}.\; (\neg (v = 0) \Rightarrow (((u \cdot \operatorname{recip}(v)) \cdot v) = u))}}
\vnbprop{cc-div-cancel}{\vnbstmt{\begin{array}{@{}l@{}} \forall x, y, w \in \mathbb{C}. \\ \quad (\neg (y = 0) \wedge \neg (w = 0)) \Rightarrow (((x \cdot w) \cdot \operatorname{recip}((y \cdot w))) = (x \cdot \operatorname{recip}(y))) \end{array}}}
\vnbprop{lin-nondegenerate}{\vnbstmt{\begin{array}{@{}l@{}} \forall a, b, c, d, u, v \in \mathbb{C}. \\ \quad (\neg (((a \cdot d) - (b \cdot c)) = 0) \wedge \neg ((u = 0) \wedge (v = 0))) \Rightarrow  \\ \quad \quad \neg ((((a \cdot u) + (b \cdot v)) = 0) \wedge (((c \cdot u) + (d \cdot v)) = 0)) \end{array}}}
\vnbprop{lft-homog}{\vnbstmt{\begin{array}{@{}l@{}} \forall a, b, c, d, u, v \in \mathbb{C}. \\ \quad (\neg (((a \cdot d) - (b \cdot c)) = 0) \wedge \neg ((u = 0) \wedge (v = 0))) \Rightarrow  \\ \quad \quad \operatorname{lft}(a, b, c, d)(\operatorname{homog\text{-}point}(u, v)) = \operatorname{homog\text{-}point}(((a \cdot u) + (b \cdot v)), ((c \cdot u) + (d \cdot v))) \end{array}}}
\vnbprop{homog-point-cover}{\vnbstmt{\begin{array}{@{}l@{}} \forall w \in \operatorname{cc\text{-}hat}. \exists u, v \in \mathbb{C}. \\ \quad \neg ((u = 0) \wedge (v = 0)) \wedge (w = \operatorname{homog\text{-}point}(u, v)) \end{array}}}
\vnbprop{lft-det-mul}{\vnbstmt{\begin{array}{@{}l@{}} \forall a, b, c, d, f, g, h, k \in \mathbb{C}. \\ \quad (\neg (((a \cdot d) - (b \cdot c)) = 0) \wedge \neg (((f \cdot k) - (g \cdot h)) = 0)) \Rightarrow  \\ \quad \quad \neg (((((a \cdot f) + (b \cdot h)) \cdot ((c \cdot g) + (d \cdot k))) - (((a \cdot g) + (b \cdot k)) \cdot ((c \cdot f) + (d \cdot h)))) = 0) \end{array}}}
\vnbprop{lft-compose}{\vnbstmt{\begin{array}{@{}l@{}} \forall a, b, c, d, f, g, h, k \in \mathbb{C}. \\ \quad (\neg (((a \cdot d) - (b \cdot c)) = 0) \wedge \neg (((f \cdot k) - (g \cdot h)) = 0)) \Rightarrow  \\ \quad \quad \operatorname{compose}(\operatorname{lft}(a, b, c, d), \operatorname{lft}(f, g, h, k)) \\ \quad \quad =\; \operatorname{lft}(((a \cdot f) + (b \cdot h)), ((a \cdot g) + (b \cdot k)), ((c \cdot f) + (d \cdot h)), ((c \cdot g) + (d \cdot k))) \end{array}}}
\vnbprop{lft-scalar-identity}{\vnbstmt{\forall r \in \mathbb{C}.\; (\neg (r = 0) \Rightarrow (\operatorname{lft}(r, 0, 0, r) = \operatorname{id\text{-}fun}(\operatorname{cc\text{-}hat})))}}

\vnbbookfile{theorem-library/zero-order-laws}{theorem-library/zero-order-laws}
\noindent{\small zero-order-laws.scm -- WAVE 2 (2026-10-27..30), wired 2026-10-30 from scratchpad/ch6/zero-order-laws.scm. Library: complex-analysis.}\par
\vnbprop{zero-order-zero-iff}{\vnbstmt{\begin{array}{@{}l@{}} \forall u, f.; \forall a \in u. \\ \quad \operatorname{holomorphic\text{-}on}(u, f) \Rightarrow (\operatorname{has\text{-}zero\text{-}of\text{-}order}(u, f, a, 0) \Leftrightarrow \neg (f(a) = 0)) \end{array}}}
\vnbprop{pow-sub-holomorphic}{\vnbstmt{\begin{array}{@{}l@{}} \forall d \in \mathbb{N}, u.; \forall c \in \mathbb{C}. \\ \quad \operatorname{is\text{-}open}(\operatorname{nf\text{-}metric\text{-}space}(\operatorname{cc\text{-}normed\text{-}field}), u) \Rightarrow  \\ \quad \quad \operatorname{holomorphic\text{-}on}(u, (\lambda x \in u.\; \operatorname{power}((x - c), d))) \end{array}}}
\vnbprop{zero-order-linear}{\vnbstmt{\begin{array}{@{}l@{}} \forall u.; \forall c \in u. \\ \quad \operatorname{is\text{-}open}(\operatorname{nf\text{-}metric\text{-}space}(\operatorname{cc\text{-}normed\text{-}field}), u) \Rightarrow  \\ \quad \quad \operatorname{has\text{-}zero\text{-}of\text{-}order}(u, (\lambda x \in u.\; (x - c)), c, 1) \end{array}}}
\vnbprop{zero-order-mul}{\vnbstmt{\begin{array}{@{}l@{}} \forall u, f, g, a, m, n.; \forall k \in (u \to \mathbb{C}). \\ \quad \operatorname{has\text{-}zero\text{-}of\text{-}order}(u, f, a, m) \wedge \\ \quad \operatorname{has\text{-}zero\text{-}of\text{-}order}(u, g, a, n) \wedge \\ \quad \forall y \in u.\; (k(y) = (f(y) \cdot g(y))) \Rightarrow \\ \quad \quad \operatorname{has\text{-}zero\text{-}of\text{-}order}(u, k, a, (m + n)) \end{array}}}
\vnbprop{zero-order-no-higher}{\vnbstmt{\begin{array}{@{}l@{}} \forall u, f, a, m.; \forall n \in \mathbb{N}, k. \\ \quad (\operatorname{has\text{-}zero\text{-}of\text{-}order}(u, f, a, m) \wedge ((m < n) \wedge \operatorname{holomorphic\text{-}on}(u, k))) \Rightarrow  \\ \quad \quad \neg \forall y \in u.\; (f(y) = (\operatorname{power}((y - a), n) \cdot k(y))) \end{array}}}
\vnbprop{zero-order-unique}{\vnbstmt{\begin{array}{@{}l@{}} \forall u, f, a, m, n. \\ \quad (\operatorname{has\text{-}zero\text{-}of\text{-}order}(u, f, a, m) \wedge \operatorname{has\text{-}zero\text{-}of\text{-}order}(u, f, a, n)) \Rightarrow (m = n) \end{array}}}
\vnbprop{zero-order-value}{\vnbstmt{\begin{array}{@{}l@{}} \forall u, f, a, m. \\ \quad \operatorname{has\text{-}zero\text{-}of\text{-}order}(u, f, a, m) \Rightarrow (\operatorname{zero\text{-}order}(u, f, a) = m) \end{array}}}
\vnbprop{count-hits-in-nn}{\vnbstmt{\forall n \in \mathbb{N}, s, z.\; (\operatorname{count\text{-}hits}(s, z, n) \in \mathbb{N})}}
\vnbprop{holomorphic-all-orders-coefficients}{\vnbstmt{\begin{array}{@{}l@{}} \forall u, f.; \forall c \in \mathbb{C}, s, r \in \mathbb{R}, k \in \mathbb{N}. \\ \quad \operatorname{holomorphic\text{-}on}(u, f) \wedge \\ \quad 0 < s \wedge \\ \quad \operatorname{ball}(\operatorname{cc\text{-}ms}, c, s) \subseteq u \wedge \\ \quad 0 < r \wedge \\ \quad r < s \wedge \\ \quad \forall n \in \mathbb{N}. \exists q. \\ \quad \quad \operatorname{holomorphic\text{-}on}(u, q) \wedge \forall x \in u.\; (f(x) = (\operatorname{power}((x - c), n) \cdot q(x))) \Rightarrow \\ \quad \quad \operatorname{recip}(((2 \cdot \pi) \cdot +i)) \cdot \operatorname{line\text{-}int}((\lambda w \in \operatorname{trace}(\operatorname{circle\text{-}path}(c, r), 0, (2 \cdot \pi)).\; (\operatorname{restrict}(f, \operatorname{trace}(\operatorname{circle\text{-}path}(c, r), 0, (2 \cdot \pi)))(w) \cdot \operatorname{power}(\operatorname{recip}((w - c)), (k + 1)))), \operatorname{circle\text{-}path}(c, r), \operatorname{circle\text{-}deriv}(c, r), 0, (2 \cdot \pi)) \\ \quad \quad =\; 0 \end{array}}}
\vnbprop{holomorphic-all-orders-local}{\vnbstmt{\begin{array}{@{}l@{}} \forall u, f.; \forall c \in u. \\ \quad \operatorname{holomorphic\text{-}on}(u, f) \wedge \\ \quad \forall n \in \mathbb{N}. \exists q. \\ \quad \quad \operatorname{holomorphic\text{-}on}(u, q) \wedge \forall x \in u.\; (f(x) = (\operatorname{power}((x - c), n) \cdot q(x))) \Rightarrow \\ \quad \quad \exists r. \\ \quad \quad \quad (\operatorname{pos\text{-}rr}(r) \wedge ((\operatorname{ball}(\operatorname{cc\text{-}ms}, c, r) \subseteq u) \wedge \forall z \in \operatorname{ball}(\operatorname{cc\text{-}ms}, c, r).\; (f(z) = 0))) \end{array}}}
\vnbprop{zero-order-exists}{\vnbstmt{\begin{array}{@{}l@{}} \forall u, f.; \forall a \in u. \\ \quad (\operatorname{is\text{-}connected}(\operatorname{cc\text{-}ms}, u) \wedge (\operatorname{holomorphic\text{-}on}(u, f) \wedge \exists w \in u.\; \neg (f(w) = 0))) \Rightarrow  \\ \quad \quad \exists m \in \mathbb{N}.\; \operatorname{has\text{-}zero\text{-}of\text{-}order}(u, f, a, m) \end{array}}}

\vnbbookfile{theorem-library/ch3b-const}{theorem-library/ch3b-const}
\noindent{\small ch3b-const.scm -- complex-analysis.pdf ch. 3: two holomorphic logarithms of the same function differ by a constant 2 pi i n on a connected open set, and the uniqueness of the principal branch of log (Cor 3.24).  Agent ISOM (brief CH3-B), 2026-10-26.}\par
\vnbprop{holomorphic-continuous-mag}{\vnbstmt{\begin{array}{@{}l@{}} \forall u, f.; \forall a \in u, s. \\ \quad (\operatorname{holomorphic\text{-}on}(u, f) \wedge \operatorname{pos\text{-}rr}(s)) \Rightarrow  \\ \quad \quad \exists l. \\ \quad \quad \quad (\operatorname{pos\text{-}rr}(l) \wedge \forall z \in u.\; ((\operatorname{magnitude}((z - a)) < l) \Rightarrow (\operatorname{magnitude}((f(z) - f(a))) < s))) \end{array}}}
\vnbprop{exp-eq-sub-one}{\vnbstmt{\forall a, b \in \mathbb{C}.\; ((\exp a = \exp b) \Rightarrow (\exp (a - b) = 1))}}
\vnbprop{exp-equal-diff-constant}{\vnbstmt{\begin{array}{@{}l@{}} \forall u, g, h.; \forall p, q \in u. \\ \quad \operatorname{is\text{-}connected}(\operatorname{cc\text{-}ms}, u) \wedge \\ \quad \operatorname{holomorphic\text{-}on}(u, g) \wedge \\ \quad \operatorname{holomorphic\text{-}on}(u, h) \wedge \\ \quad \forall z \in u.\; (\exp g(z) = \exp h(z)) \Rightarrow \\ \quad \quad (g(p) - h(p)) = (g(q) - h(q)) \end{array}}}
\vnbprop{exp-equal-diff-2-pi-i}{\vnbstmt{\begin{array}{@{}l@{}} \forall u, g, h.; \forall a \in u. \\ \quad \operatorname{is\text{-}connected}(\operatorname{cc\text{-}ms}, u) \wedge \\ \quad \operatorname{holomorphic\text{-}on}(u, g) \wedge \\ \quad \operatorname{holomorphic\text{-}on}(u, h) \wedge \\ \quad \forall z \in u.\; (\exp g(z) = \exp h(z)) \Rightarrow \\ \quad \quad \exists n \in \mathbb{Z}.\; \forall z \in u.\; (g(z) = (h(z) + (n \cdot ((2 \cdot \pi) \cdot +i)))) \end{array}}}
\vnbprop{cc-log-unique-branch}{\vnbstmt{\begin{array}{@{}l@{}} \forall g.; \forall z \in \{u \in \mathbb{C} : \neg ((\operatorname{imag\text{-}part}(u) = 0) \wedge (\operatorname{real\text{-}part}(u) \leq 0))\}. \\ \quad \operatorname{holomorphic\text{-}on}(\{u \in \mathbb{C} : \neg ((\operatorname{imag\text{-}part}(u) = 0) \wedge (\operatorname{real\text{-}part}(u) \leq 0))\}, g) \wedge \\ \quad g(1) = 0 \wedge \\ \quad \forall z \in \{u \in \mathbb{C} : \neg ((\operatorname{imag\text{-}part}(u) = 0) \wedge (\operatorname{real\text{-}part}(u) \leq 0))\}. \\ \quad \quad \exp g(z) = z \Rightarrow \\ \quad \quad g(z) = \log\left(z\right) \end{array}}}
\vnbprop{principal-log-exists-unique}{\vnbstmt{\begin{array}{@{}l@{}} \exists g. \\ \quad \operatorname{holomorphic\text{-}on}(\{u \in \mathbb{C} : \neg ((\operatorname{imag\text{-}part}(u) = 0) \wedge (\operatorname{real\text{-}part}(u) \leq 0))\}, g) \wedge \\ \quad g(1) = 0 \wedge \\ \quad \forall z \in \{u \in \mathbb{C} : \neg ((\operatorname{imag\text{-}part}(u) = 0) \wedge (\operatorname{real\text{-}part}(u) \leq 0))\}. \\ \quad \quad \exp g(z) = z \wedge \\ \forall g.; \forall z \in \{u \in \mathbb{C} : \neg ((\operatorname{imag\text{-}part}(u) = 0) \wedge (\operatorname{real\text{-}part}(u) \leq 0))\}. \\ \quad \operatorname{holomorphic\text{-}on}(\{u \in \mathbb{C} : \neg ((\operatorname{imag\text{-}part}(u) = 0) \wedge (\operatorname{real\text{-}part}(u) \leq 0))\}, g) \wedge \\ \quad g(1) = 0 \wedge \\ \quad \forall z \in \{u \in \mathbb{C} : \neg ((\operatorname{imag\text{-}part}(u) = 0) \wedge (\operatorname{real\text{-}part}(u) \leq 0))\}. \\ \quad \quad \exp g(z) = z \Rightarrow \\ \quad \quad g(z) = \log\left(z\right) \end{array}}}

\vnbbookfile{theorem-library/ch3b-log}{theorem-library/ch3b-log}
\noindent{\small ch3b-log.scm -- complex-analysis.pdf: logarithms and n-th roots of a nonvanishing holomorphic function on a CONVEX open set (the convex forms of Prop 3.22 and Cor 3.25), and Lemma 6.19 (a nowhere-zero entire function is an exponential).  Agent ISOM (brief CH3-B), 2026-10-26.}\par
\vnbprop{zero-deriv-connected-constant}{\vnbstmt{\begin{array}{@{}l@{}} \forall u, f.; \forall p, q \in u. \\ \quad \operatorname{is\text{-}connected}(\operatorname{cc\text{-}ms}, u) \wedge \\ \quad f \in (u \to \mathbb{C}) \wedge \\ \quad \forall z \in u.\; \operatorname{is\text{-}diff\text{-}on}(\operatorname{cc\text{-}normed\text{-}field}, u, f, z, 0) \Rightarrow \\ \quad \quad f(p) = f(q) \end{array}}}
\vnbprop{holomorphic-log-convex-exists}{\vnbstmt{\begin{array}{@{}l@{}} \forall u, f. \\ \quad (\operatorname{is\text{-}convex\text{-}cc}(u) \wedge (\operatorname{holomorphic\text{-}on}(u, f) \wedge \forall z \in u.\; \neg (f(z) = 0))) \Rightarrow  \\ \quad \quad \exists g.\; (\operatorname{holomorphic\text{-}on}(u, g) \wedge \forall z \in u.\; (\exp g(z) = f(z))) \end{array}}}
\vnbprop{exp-power-nn}{\vnbstmt{\forall n \in \mathbb{N}, w \in \mathbb{C}.\; (\operatorname{power}(\exp w, n) = \exp (n \cdot w))}}
\vnbprop{holomorphic-log-convex}{\vnbstmt{\begin{array}{@{}l@{}} \forall u, f. \\ \quad (\operatorname{is\text{-}convex\text{-}cc}(u) \wedge (\operatorname{holomorphic\text{-}on}(u, f) \wedge \forall z \in u.\; \neg (f(z) = 0))) \Rightarrow  \\ \quad \quad \exists g.\; (\operatorname{holomorphic\text{-}on}(u, g) \wedge \forall z \in u.\; (\exp g(z) = f(z))) \wedge \\ \quad \quad \forall g, h.; \forall a \in u. \\ \quad \quad \quad \operatorname{holomorphic\text{-}on}(u, g) \wedge \\ \quad \quad \quad \operatorname{holomorphic\text{-}on}(u, h) \wedge \\ \quad \quad \quad \forall z \in u.\; (\exp g(z) = f(z)) \wedge \\ \quad \quad \quad \forall z \in u.\; (\exp h(z) = f(z)) \Rightarrow \\ \quad \quad \quad \quad \exists n \in \mathbb{Z}.\; \forall z \in u.\; (g(z) = (h(z) + (n \cdot ((2 \cdot \pi) \cdot +i)))) \end{array}}}
\vnbprop{cc-is-convex}{\vnbstmt{\operatorname{is\text{-}convex\text{-}cc}(\mathbb{C})}}
\vnbprop{entire-nonvanishing-iff-exp}{\vnbstmt{\begin{array}{@{}l@{}} \forall f. \\ \quad \operatorname{holomorphic\text{-}on}(\mathbb{C}, f) \Rightarrow  \\ \quad \quad \forall z \in \mathbb{C}.\; \neg (f(z) = 0) \\ \quad \quad \Leftrightarrow\; \exists g.\; (\operatorname{holomorphic\text{-}on}(\mathbb{C}, g) \wedge \forall z \in \mathbb{C}.\; (f(z) = \exp g(z))) \end{array}}}
\vnbprop{holomorphic-root-convex-exists}{\vnbstmt{\begin{array}{@{}l@{}} \forall u, f, n. \\ \quad (\operatorname{is\text{-}convex\text{-}cc}(u) \wedge (\operatorname{holomorphic\text{-}on}(u, f) \wedge (\forall z \in u.\; \neg (f(z) = 0) \wedge ((n \in \mathbb{N}) \wedge (1 \leq n))))) \Rightarrow  \\ \quad \quad \exists g.\; (\operatorname{holomorphic\text{-}on}(u, g) \wedge \forall z \in u.\; (\operatorname{power}(g(z), n) = f(z))) \end{array}}}
\vnbprop{holomorphic-root-convex-unique}{\vnbstmt{\begin{array}{@{}l@{}} \forall u, f, n, g, h.; \forall b \in u. \\ \quad \operatorname{is\text{-}convex\text{-}cc}(u) \wedge \\ \quad \operatorname{holomorphic\text{-}on}(u, f) \wedge \\ \quad \forall z \in u.\; \neg (f(z) = 0) \wedge \\ \quad n \in \mathbb{N} \wedge \\ \quad 1 \leq n \wedge \\ \quad \operatorname{holomorphic\text{-}on}(u, g) \wedge \\ \quad \operatorname{holomorphic\text{-}on}(u, h) \wedge \\ \quad \forall z \in u.\; (\operatorname{power}(g(z), n) = f(z)) \wedge \\ \quad \forall z \in u.\; (\operatorname{power}(h(z), n) = f(z)) \Rightarrow \\ \quad \quad \exists a \in \mathbb{C}.\; ((\operatorname{power}(a, n) = 1) \wedge \forall z \in u.\; (h(z) = (a \cdot g(z)))) \end{array}}}

\vnbbookfile{theorem-library/ch4a-circle}{theorem-library/ch4a-circle}
\noindent{\small ch4a-circle.scm -- WAVE 2 (2026-10-27..30), wired 2026-10-30 from scratchpad/ch4a/ch4a-circle.scm. Library: complex-analysis.}\par
\vnbprop{circle-power-integral}{\vnbstmt{\begin{array}{@{}l@{}} \forall c \in \mathbb{C}, r \in \mathbb{R}, k \in \mathbb{N}. \\ \quad \operatorname{line\text{-}int}((\lambda w \in \mathbb{C}.\; \operatorname{power}((w - c), k)), \operatorname{circle\text{-}path}(c, r), \operatorname{circle\text{-}deriv}(c, r), 0, (2 \cdot \pi)) = 0 \end{array}}}
\vnbprop{circle-recip-integral}{\vnbstmt{\begin{array}{@{}l@{}} \forall c \in \mathbb{C}, r \in \mathbb{R}. \\ \quad (0 < r) \Rightarrow  \\ \quad \quad \operatorname{line\text{-}int}((\lambda w \in \operatorname{difference}(\mathbb{C}, \operatorname{singleton}(c)).\; \operatorname{recip}((w - c))), \\ \quad \quad \quad \operatorname{circle\text{-}path}(c, r), \operatorname{circle\text{-}deriv}(c, r), 0, (2 \cdot \pi)) \\ \quad \quad =\; (2 \cdot \pi) \cdot +i \end{array}}}
\vnbprop{circle-recip-power-integral}{\vnbstmt{\begin{array}{@{}l@{}} \forall c \in \mathbb{C}, r \in \mathbb{R}, j \in \mathbb{N}. \\ \quad (0 < r) \Rightarrow  \\ \quad \quad \operatorname{line\text{-}int}((\lambda w \in \operatorname{difference}(\mathbb{C}, \operatorname{singleton}(c)).\; \operatorname{power}(\operatorname{recip}((w - c)), ((j + 1) + 1))), \\ \quad \quad \quad \operatorname{circle\text{-}path}(c, r), \operatorname{circle\text{-}deriv}(c, r), 0, (2 \cdot \pi)) \\ \quad \quad =\; 0 \end{array}}}

\vnbbookfile{theorem-library/ch4a-coef}{theorem-library/ch4a-coef}
\noindent{\small ch4a-coef.scm -- WAVE 2 (2026-10-27..30), wired 2026-10-30 from scratchpad/ch4a/ch4a-coef.scm. Library: complex-analysis.}\par
\vnbprop{lc-partial-sum-before}{\vnbstmt{\begin{array}{@{}l@{}} \forall k \in \mathbb{N}, t \in (\mathbb{N} \to \mathbb{C}), n \in \mathbb{N}. \\ \quad (\forall j \in \mathbb{N}.\; (\neg (j = n) \Rightarrow (t(j) = 0)) \wedge (k \leq n)) \Rightarrow  \\ \quad \quad \operatorname{cc\text{-}series\text{-}partial\text{-}sum}(t, k) = 0 \end{array}}}
\vnbprop{lc-partial-sum-after}{\vnbstmt{\begin{array}{@{}l@{}} \forall m \in \mathbb{N}, t \in (\mathbb{N} \to \mathbb{C}), n \in \mathbb{N}. \\ \quad (\forall j \in \mathbb{N}.\; (\neg (j = n) \Rightarrow (t(j) = 0)) \wedge (n < m)) \Rightarrow  \\ \quad \quad \operatorname{cc\text{-}series\text{-}partial\text{-}sum}(t, m) = t(n) \end{array}}}
\vnbprop{cc-series-single-term}{\vnbstmt{\begin{array}{@{}l@{}} \forall t \in (\mathbb{N} \to \mathbb{C}), n \in \mathbb{N}. \\ \quad \forall j \in \mathbb{N}.\; (\neg (j = n) \Rightarrow (t(j) = 0)) \Rightarrow  \\ \quad \quad \operatorname{cc\text{-}series\text{-}converges\text{-}to}(t, t(n)) \end{array}}}
\vnbprop{lc-power-cancel}{\vnbstmt{\forall x \in \mathbb{C}, j \in \mathbb{N}.\; (\neg (x = 0) \Rightarrow ((\operatorname{power}(x, j) \cdot \operatorname{power}(\operatorname{recip}(x), j)) = 1))}}
\vnbprop{lc-power-one-left}{\vnbstmt{\begin{array}{@{}l@{}} \forall x \in \mathbb{C}, j \in \mathbb{N}. \\ \quad \neg (x = 0) \Rightarrow ((\operatorname{power}(x, j) \cdot \operatorname{power}(\operatorname{recip}(x), (j + 1))) = \operatorname{recip}(x)) \end{array}}}
\vnbprop{lc-power-recip-excess}{\vnbstmt{\begin{array}{@{}l@{}} \forall x \in \mathbb{C}, j, t \in \mathbb{N}. \\ \quad \neg (x = 0) \Rightarrow  \\ \quad \quad (\operatorname{power}(x, j) \cdot \operatorname{power}(\operatorname{recip}(x), (((j + 1) + 1) + t))) = \operatorname{power}(\operatorname{recip}(x), ((t + 1) + 1)) \end{array}}}
\vnbprop{lc-power-x-excess}{\vnbstmt{\begin{array}{@{}l@{}} \forall x \in \mathbb{C}, m, t \in \mathbb{N}. \\ \quad \neg (x = 0) \Rightarrow ((\operatorname{power}(x, (m + t)) \cdot \operatorname{power}(\operatorname{recip}(x), m)) = \operatorname{power}(x, t)) \end{array}}}
\vnbprop{lc-monomial-one}{\vnbstmt{\begin{array}{@{}l@{}} \forall c \in \mathbb{C}, r \in \mathbb{R}, j \in \mathbb{N}. \\ \quad (0 < r) \Rightarrow  \\ \quad \quad \operatorname{line\text{-}int}((\lambda w \in \operatorname{trace}(\operatorname{circle\text{-}path}(c, r), 0, (2 \cdot \pi)).\; (\operatorname{power}((w - c), j) \cdot \operatorname{power}(\operatorname{recip}((w - c)), (j + 1)))), \\ \quad \quad \quad \operatorname{circle\text{-}path}(c, r), \operatorname{circle\text{-}deriv}(c, r), 0, (2 \cdot \pi)) \\ \quad \quad =\; (2 \cdot \pi) \cdot +i \end{array}}}
\vnbprop{lc-monomial-zero}{\vnbstmt{\begin{array}{@{}l@{}} \forall c \in \mathbb{C}, r \in \mathbb{R}, j, m \in \mathbb{N}. \\ \quad ((0 < r) \wedge \neg (m = (j + 1))) \Rightarrow  \\ \quad \quad \operatorname{line\text{-}int}((\lambda w \in \operatorname{trace}(\operatorname{circle\text{-}path}(c, r), 0, (2 \cdot \pi)).\; (\operatorname{power}((w - c), j) \cdot \operatorname{power}(\operatorname{recip}((w - c)), m))), \\ \quad \quad \quad \operatorname{circle\text{-}path}(c, r), \operatorname{circle\text{-}deriv}(c, r), 0, (2 \cdot \pi)) \\ \quad \quad =\; 0 \end{array}}}
\vnbprop{lc-radius-mid}{\vnbstmt{\begin{array}{@{}l@{}} \forall f \in (\mathbb{N} \to \mathbb{C}), o \in \mathbb{R}. \\ \quad ((0 \leq o) \wedge (o < \operatorname{cps\text{-}radius}(f))) \Rightarrow  \\ \quad \quad \exists t \in \mathbb{R}.\; ((o < t) \wedge (t < \operatorname{cps\text{-}radius}(f))) \end{array}}}
\vnbprop{cc-power-series-termwise}{\vnbstmt{\begin{array}{@{}l@{}} \forall u, m, c, a, b.; \forall f \in (\mathbb{N} \to \mathbb{C}), o \in \mathbb{R}, \mathit{lczf}, \mathit{lcwf} \in (u \to \mathbb{C}). \\ \quad u \subseteq \mathbb{C} \wedge \\ \quad \operatorname{is\text{-}road}(m, c, a, b) \wedge \\ \quad \operatorname{trace}(m, a, b) \subseteq u \wedge \\ \quad 0 \leq o \wedge \\ \quad o < \operatorname{cps\text{-}radius}(f) \wedge \\ \quad \forall y \in u. \\ \quad \quad \operatorname{is\text{-}continuous\text{-}at}(\operatorname{subspace\text{-}ms}(\operatorname{nf\text{-}metric\text{-}space}(\operatorname{cc\text{-}normed\text{-}field}), u), \\ \quad \quad \quad \operatorname{nf\text{-}metric\text{-}space}(\operatorname{cc\text{-}normed\text{-}field}), \mathit{lczf}, y) \wedge \\ \quad \forall x \in u.\; (\operatorname{magnitude}(\operatorname{lczf}(x)) \leq o) \wedge \\ \quad \forall y \in u. \\ \quad \quad \operatorname{is\text{-}continuous\text{-}at}(\operatorname{subspace\text{-}ms}(\operatorname{nf\text{-}metric\text{-}space}(\operatorname{cc\text{-}normed\text{-}field}), u), \\ \quad \quad \quad \operatorname{nf\text{-}metric\text{-}space}(\operatorname{cc\text{-}normed\text{-}field}), \mathit{lcwf}, y) \Rightarrow \\ \quad \quad (\lambda w \in u.\; (\operatorname{cc\text{-}series\text{-}limit}((\lambda k \in \mathbb{N}.\; (f(k) \cdot \operatorname{power}(\operatorname{lczf}(w), k)))) \cdot \operatorname{lcwf}(w))) \in (u \to \mathbb{C}) \wedge \\ \quad \quad \forall y \in u. \\ \quad \quad \quad \operatorname{is\text{-}continuous\text{-}at}(\operatorname{subspace\text{-}ms}(\operatorname{nf\text{-}metric\text{-}space}(\operatorname{cc\text{-}normed\text{-}field}), u), \\ \quad \quad \quad \quad \operatorname{nf\text{-}metric\text{-}space}(\operatorname{cc\text{-}normed\text{-}field}), \\ \quad \quad \quad \quad (\lambda w \in u.\; (\operatorname{cc\text{-}series\text{-}limit}((\lambda k \in \mathbb{N}.\; (f(k) \cdot \operatorname{power}(\operatorname{lczf}(w), k)))) \cdot \operatorname{lcwf}(w))), \\ \quad \quad \quad \quad y) \wedge \\ \quad \quad (\lambda \mathit{lck} \in \mathbb{N}.\; (f(\mathit{lck}) \cdot \operatorname{line\text{-}int}((\lambda w \in u.\; (\operatorname{power}(\operatorname{lczf}(w), \mathit{lck}) \cdot \operatorname{lcwf}(w))), m, c, a, b))) \\ \quad \quad \in\; (\mathbb{N} \to \mathbb{C}) \wedge \\ \quad \quad \operatorname{cc\text{-}series\text{-}converges\text{-}to}((\lambda \mathit{lck} \in \mathbb{N}.\; (f(\mathit{lck}) \cdot \operatorname{line\text{-}int}((\lambda w \in u.\; (\operatorname{power}(\operatorname{lczf}(w), \mathit{lck}) \cdot \operatorname{lcwf}(w))), m, c, a, b))), \\ \quad \quad \quad \operatorname{line\text{-}int}((\lambda w \in u.\; (\operatorname{cc\text{-}series\text{-}limit}((\lambda k \in \mathbb{N}.\; (f(k) \cdot \operatorname{power}(\operatorname{lczf}(w), k)))) \cdot \operatorname{lcwf}(w))), \\ \quad \quad \quad \quad m, c, a, b)) \end{array}}}
\vnbprop{laurent-coefficients}{\vnbstmt{\begin{array}{@{}l@{}} \forall \mathit{lcc} \in (\mathbb{Z} \to \mathbb{C}), c \in \mathbb{C}, r \in \mathbb{R}, n \in \mathbb{N}. \\ \quad 0 < r \wedge \\ \quad r < \operatorname{cps\text{-}radius}((\lambda k \in \mathbb{N}.\; \operatorname{lcc}(k))) \wedge \\ \quad \operatorname{recip}(r) < \operatorname{cps\text{-}radius}((\lambda k \in \mathbb{N}.\; \operatorname{lcc}(-(k + 1)))) \Rightarrow \\ \quad \quad \operatorname{lcc}(n) \\ \quad \quad =\; \operatorname{recip}(((2 \cdot \pi) \cdot +i)) \cdot \operatorname{line\text{-}int}((\lambda w \in \operatorname{trace}(\operatorname{circle\text{-}path}(c, r), 0, (2 \cdot \pi)).\; (\operatorname{laurent\text{-}sum}(\mathit{lcc}, c, w) \cdot \operatorname{power}(\operatorname{recip}((w - c)), (n + 1)))), \operatorname{circle\text{-}path}(c, r), \operatorname{circle\text{-}deriv}(c, r), 0, (2 \cdot \pi)) \wedge \\ \quad \quad \operatorname{lcc}(-(n + 1)) \\ \quad \quad =\; \operatorname{recip}(((2 \cdot \pi) \cdot +i)) \cdot \operatorname{line\text{-}int}((\lambda w \in \operatorname{trace}(\operatorname{circle\text{-}path}(c, r), 0, (2 \cdot \pi)).\; (\operatorname{laurent\text{-}sum}(\mathit{lcc}, c, w) \cdot \operatorname{power}((w - c), n))), \operatorname{circle\text{-}path}(c, r), \operatorname{circle\text{-}deriv}(c, r), 0, (2 \cdot \pi)) \end{array}}}
\vnbprop{cps-radius-from-bound}{\vnbstmt{\begin{array}{@{}l@{}} \forall f \in (\mathbb{N} \to \mathbb{C}), c, q, x \in \mathbb{R}. \\ \quad ((0 \leq c) \wedge ((0 < q) \wedge (\forall k \in \mathbb{N}.\; (\operatorname{magnitude}(f(k)) \leq (c \cdot \operatorname{power}(\operatorname{recip}(q), k))) \wedge (x < q)))) \Rightarrow  \\ \quad \quad x < \operatorname{cps\text{-}radius}(f) \end{array}}}

\vnbbookfile{theorem-library/catools-nthlocal}{theorem-library/catools-nthlocal}
\noindent{\small catools-nthlocal.scm -- derivatives are local; Cauchy's estimates as the notes state them (CA-TOOLS, 2026-10-26).  Destination: theorem-library/catools-nthlocal.scm, library complex-analysis.}\par
\vnbprop{nth-deriv-on-local}{\vnbstmt{\begin{array}{@{}l@{}} \forall k \in \mathbb{N}, u, f, v.; \forall z \in v. \\ \quad (\operatorname{holomorphic\text{-}on}(u, f) \wedge (\operatorname{is\text{-}open}(\operatorname{nf\text{-}metric\text{-}space}(\operatorname{cc\text{-}normed\text{-}field}), v) \wedge (v \subseteq u))) \Rightarrow  \\ \quad \quad \operatorname{restrict}(f, v)^{(k)}(z) \\ \quad \quad =\; f^{(k)}(z) \end{array}}}
\vnbprop{cauchy-estimates}{\vnbstmt{\begin{array}{@{}l@{}} \forall u, f.; \forall c \in \mathbb{C}, r, m \in \mathbb{R}, k \in \mathbb{N}. \\ \quad \operatorname{holomorphic\text{-}on}(u, f) \wedge \\ \quad 0 < r \wedge \\ \quad \operatorname{ball}(\operatorname{cc\text{-}ms}, c, r) \subseteq u \wedge \\ \quad \forall z \in \operatorname{ball}(\operatorname{cc\text{-}ms}, c, r).\; (\operatorname{magnitude}(f(z)) \leq m) \Rightarrow \\ \quad \quad \operatorname{magnitude}(f^{(k)}(c)) \leq (\operatorname{factorial}(k) \cdot (m \cdot \operatorname{power}(\operatorname{recip}(r), k))) \end{array}}}

\vnbbookfile{theorem-library/catools-hfrechet}{theorem-library/catools-hfrechet}
\noindent{\small catools-hfrechet.scm -- H(U) is a Frechet space (CA-TOOLS, 2026-10-26). Destination: theorem-library/catools-hfrechet.scm, library complex-analysis.}\par
\vnbprop{cc-hol-seminorms-family}{\vnbstmt{\begin{array}{@{}l@{}} \forall u. \\ \quad \operatorname{is\text{-}open}(\operatorname{nf\text{-}metric\text{-}space}(\operatorname{cc\text{-}normed\text{-}field}), u) \Rightarrow  \\ \quad \quad \operatorname{is\text{-}seminorm\text{-}family}(\operatorname{cc\text{-}normed\text{-}field}, \operatorname{cc\text{-}fn\text{-}module}(u, \operatorname{cc\text{-}hol\text{-}fns}(u)), \\ \quad \quad \quad \operatorname{cc\text{-}exhaust\text{-}seminorms}(u, \operatorname{cc\text{-}hol\text{-}fns}(u))) \end{array}}}
\vnbprop{cc-hol-seminorm-of-difference}{\vnbstmt{\begin{array}{@{}l@{}} \forall u.; \forall a, b \in \operatorname{cc\text{-}hol\text{-}fns}(u), k \in \mathbb{N}. \\ \quad \operatorname{is\text{-}open}(\operatorname{nf\text{-}metric\text{-}space}(\operatorname{cc\text{-}normed\text{-}field}), u) \Rightarrow  \\ \quad \quad (\operatorname{cc\text{-}exhaust\text{-}seminorms}(u, \operatorname{cc\text{-}hol\text{-}fns}(u))(k)(\operatorname{vadd}(\operatorname{cc\text{-}fn\text{-}module}(u, \operatorname{cc\text{-}hol\text{-}fns}(u)))(a, \operatorname{vneg}(\operatorname{cc\text{-}fn\text{-}module}(u, \operatorname{cc\text{-}hol\text{-}fns}(u)))(b))) \simeq \operatorname{cc\text{-}exhaust\text{-}seminorms}(u, \operatorname{cc\text{-}cont\text{-}fns}(u))(k)(\operatorname{vadd}(\operatorname{cc\text{-}fn\text{-}module}(u, \operatorname{cc\text{-}cont\text{-}fns}(u)))(a, \operatorname{vneg}(\operatorname{cc\text{-}fn\text{-}module}(u, \operatorname{cc\text{-}cont\text{-}fns}(u)))(b)))) \end{array}}}
\vnbprop{cc-hol-fns-frechet}{\vnbstmt{\begin{array}{@{}l@{}} \forall u. \\ \quad \operatorname{is\text{-}open}(\operatorname{nf\text{-}metric\text{-}space}(\operatorname{cc\text{-}normed\text{-}field}), u) \Rightarrow  \\ \quad \quad \operatorname{is\text{-}frechet\text{-}structure}(\operatorname{cc\text{-}normed\text{-}field}, \operatorname{cc\text{-}fn\text{-}module}(u, \operatorname{cc\text{-}hol\text{-}fns}(u)), \\ \quad \quad \quad \operatorname{cc\text{-}exhaust\text{-}seminorms}(u, \operatorname{cc\text{-}hol\text{-}fns}(u))) \end{array}}}

\vnbbookfile{theorem-library/compact-margins}{theorem-library/compact-margins}
\noindent{\small ch3c-gen.scm -- general lemmas for the general Cauchy theorem.  Prefix cm-.  Library: analysis. Destination: theorem-library/compact-margins.scm.}\par
\vnbprop{compact-open-margin}{\vnbstmt{\begin{array}{@{}l@{}} \forall k, u. \\ \quad ((k \subseteq u) \wedge (\operatorname{is\text{-}open}(\operatorname{cc\text{-}ms}, u) \wedge \operatorname{is\text{-}compact}(\operatorname{subspace\text{-}ms}(\operatorname{cc\text{-}ms}, k)))) \Rightarrow  \\ \quad \quad \exists e.\; (\operatorname{pos\text{-}rr}(e) \wedge \forall z \in k.\; (\operatorname{ball}(\operatorname{cc\text{-}ms}, z, e) \subseteq u)) \end{array}}}
\vnbprop{pasting-local}{\vnbstmt{\begin{array}{@{}l@{}} \forall s, t, f, a, x, b. \\ \quad \operatorname{is\text{-}metric\text{-}space}(s) \wedge \\ \quad \operatorname{is\text{-}metric\text{-}space}(t) \wedge \\ \quad f \in (\operatorname{pts}(s) \to \operatorname{pts}(t)) \wedge \\ \quad a \subseteq \operatorname{pts}(s) \wedge \\ \quad \operatorname{is\text{-}closed}(s, a) \wedge \\ \quad \forall y \in a.\; \operatorname{is\text{-}continuous\text{-}at}(\operatorname{subspace\text{-}ms}(s, a), t, \operatorname{restrict}(f, a), y) \wedge \\ \quad x \in \operatorname{pts}(s) \wedge \\ \quad \operatorname{pos\text{-}rr}(b) \Rightarrow \\ \quad \quad \exists d. \\ \quad \quad \quad \operatorname{pos\text{-}rr}(d) \wedge \\ \quad \quad \quad \forall \mathit{cmb} \in \operatorname{pts}(s). \\ \quad \quad \quad \quad ((\operatorname{dist}(s)(x, \mathit{cmb}) \leq d) \wedge (\mathit{cmb} \in a)) \Rightarrow (\operatorname{dist}(t)(f(x), f(\mathit{cmb})) \leq b) \end{array}}}
\vnbprop{pasting-closed}{\vnbstmt{\begin{array}{@{}l@{}} \forall s, t, f, a, b.; \forall z \in \operatorname{pts}(s). \\ \quad \operatorname{is\text{-}metric\text{-}space}(s) \wedge \\ \quad \operatorname{is\text{-}metric\text{-}space}(t) \wedge \\ \quad f \in (\operatorname{pts}(s) \to \operatorname{pts}(t)) \wedge \\ \quad a \subseteq \operatorname{pts}(s) \wedge \\ \quad b \subseteq \operatorname{pts}(s) \wedge \\ \quad \operatorname{is\text{-}closed}(s, a) \wedge \\ \quad \operatorname{is\text{-}closed}(s, b) \wedge \\ \quad \forall x \in \operatorname{pts}(s).\; ((x \in a) \vee (x \in b)) \wedge \\ \quad \forall y \in a.\; \operatorname{is\text{-}continuous\text{-}at}(\operatorname{subspace\text{-}ms}(s, a), t, \operatorname{restrict}(f, a), y) \wedge \\ \quad \forall y \in b.\; \operatorname{is\text{-}continuous\text{-}at}(\operatorname{subspace\text{-}ms}(s, b), t, \operatorname{restrict}(f, b), y) \Rightarrow \\ \quad \quad \operatorname{is\text{-}continuous\text{-}at}(s, t, f, z) \end{array}}}

\vnbbookfile{theorem-library/mvp-holo}{theorem-library/mvp-holo}
\noindent{\small mvp-holo.scm -- WAVE 2 (2026-10-27..30), wired 2026-10-30 from scratchpad/calint2/mvp-holo.scm. Library: complex-analysis.}\par
\vnbprop{cc-radial-scale}{\vnbstmt{\begin{array}{@{}l@{}} \forall c, y \in \mathbb{C}, s \in \mathbb{R}. \\ \quad ((0 \leq s) \wedge (s \leq 1)) \Rightarrow  \\ \quad \quad \operatorname{magnitude}(((c + (s \cdot (y - c))) - c)) = (s \cdot \operatorname{magnitude}((y - c))) \wedge \\ \quad \quad \operatorname{magnitude}((y - (c + (s \cdot (y - c))))) = ((1 - s) \cdot \operatorname{magnitude}((y - c))) \end{array}}}
\vnbprop{closed-ball-in-closure-ball}{\vnbstmt{\begin{array}{@{}l@{}} \forall c \in \mathbb{C}, r \in \mathbb{R}. \\ \quad (0 < r) \Rightarrow (\operatorname{closed\text{-}ball}(\operatorname{cc\text{-}ms}, c, r) \subseteq \operatorname{closure}(\operatorname{cc\text{-}ms}, \operatorname{ball}(\operatorname{cc\text{-}ms}, c, r))) \end{array}}}
\vnbprop{closed-disc-margin}{\vnbstmt{\begin{array}{@{}l@{}} \forall u.; \forall c \in \mathbb{C}, r \in \mathbb{R}. \\ \quad (\operatorname{is\text{-}open}(\operatorname{cc\text{-}ms}, u) \wedge ((0 < r) \wedge (\operatorname{closure}(\operatorname{cc\text{-}ms}, \operatorname{ball}(\operatorname{cc\text{-}ms}, c, r)) \subseteq u))) \Rightarrow  \\ \quad \quad \exists \mathit{mvr} \in \mathbb{R}.\; ((r < \mathit{mvr}) \wedge (\operatorname{ball}(\operatorname{cc\text{-}ms}, c, \mathit{mvr}) \subseteq u)) \end{array}}}
\vnbprop{holomorphic-has-mean-value-property}{\vnbstmt{\forall u, f.\; (\operatorname{holomorphic\text{-}on}(u, f) \Rightarrow \operatorname{has\text{-}mean\text{-}value\text{-}property}(u, f))}}
\vnbprop{holomorphic-maximum-modulus}{\vnbstmt{\begin{array}{@{}l@{}} \forall u, f.; \forall a, z \in u. \\ \quad \operatorname{holomorphic\text{-}on}(u, f) \wedge \\ \quad \operatorname{is\text{-}connected}(\operatorname{cc\text{-}ms}, u) \wedge \\ \quad \forall z \in u.\; (\operatorname{magnitude}(f(z)) \leq \operatorname{magnitude}(f(a))) \Rightarrow \\ \quad \quad f(z) = f(a) \end{array}}}

\vnbbookfile{structure-library/ca-homotopy-vocab}{structure-library/ca-homotopy-vocab}
\noindent{\small ch3c-vocab.scm -- vocabulary for the general Cauchy theorem (complex-analysis.pdf 3.3). Destination: structure-library/ca-homotopy-vocab.scm.  Library: complex-analysis. Binders: hm*\_ (hma\_ hmb\_ hmd\_ hmf\_ hmk\_ hmn\_ hms\_ hmt\_ hmw\_ hmx\_ hmz\_), used by nothing else.}\par
\par\medskip\noindent\textbf{Definition} (\texttt{good-approx}).\;For all $\operatorname{hmu\_}$, $\operatorname{hmg\_}$, $\operatorname{hmr\_}$, $\operatorname{hma\_}$ and $\operatorname{hmb\_}$, 
\begin{equation*}
\begin{array}{@{}l@{}} \operatorname{good\text{-}approx}(\operatorname{hmu\_}, \operatorname{hmg\_}, \operatorname{hmr\_}, \operatorname{hma\_}, \operatorname{hmb\_}) \\ \Leftrightarrow\; \operatorname{is\text{-}path}(\operatorname{hmg\_}, \operatorname{hma\_}, \operatorname{hmb\_}) \wedge \\ \operatorname{is\text{-}path}(\operatorname{hmr\_}, \operatorname{hma\_}, \operatorname{hmb\_}) \wedge \\ \operatorname{trace}(\operatorname{hmg\_}, \operatorname{hma\_}, \operatorname{hmb\_}) \subseteq \operatorname{hmu\_} \wedge \\ \operatorname{trace}(\operatorname{hmr\_}, \operatorname{hma\_}, \operatorname{hmb\_}) \subseteq \operatorname{hmu\_} \wedge \\ \operatorname{hmr\_}(\operatorname{hma\_}) = \operatorname{hmg\_}(\operatorname{hma\_}) \wedge \\ \operatorname{hmr\_}(\operatorname{hmb\_}) = \operatorname{hmg\_}(\operatorname{hmb\_}) \wedge \\ \exists \operatorname{hmn\_} \in \mathbb{N}. \\ \quad 1 \leq \operatorname{hmn\_} \wedge \\ \quad \exists \operatorname{hmt\_} \in (\mathbb{N} \to \mathbb{R}). \\ \quad \quad \operatorname{hmt\_}(0) = \operatorname{hma\_} \wedge \\ \quad \quad \operatorname{hmt\_}(\operatorname{hmn\_}) = \operatorname{hmb\_} \wedge \\ \quad \quad \exists \operatorname{hmd\_} \in (\mathbb{N} \to \operatorname{power}(\mathbb{C})). \forall \operatorname{hmk\_} \in \mathbb{N}. \\ \quad \quad \quad (\operatorname{hmk\_} < \operatorname{hmn\_}) \Rightarrow  \\ \quad \quad \quad \quad \operatorname{hmt\_}(\operatorname{hmk\_}) < \operatorname{hmt\_}(\operatorname{succ}(\operatorname{hmk\_})) \wedge \\ \quad \quad \quad \quad \operatorname{is\text{-}open}(\operatorname{nf\text{-}metric\text{-}space}(\operatorname{cc\text{-}normed\text{-}field}), \operatorname{hmd\_}(\operatorname{hmk\_})) \wedge \\ \quad \quad \quad \quad \operatorname{is\text{-}convex\text{-}cc}(\operatorname{hmd\_}(\operatorname{hmk\_})) \wedge \\ \quad \quad \quad \quad \operatorname{hmd\_}(\operatorname{hmk\_}) \subseteq \operatorname{hmu\_} \wedge \\ \quad \quad \quad \quad \forall \operatorname{hmw\_} \in \operatorname{ccint}(\operatorname{hmt\_}(\operatorname{hmk\_}), \operatorname{hmt\_}(\operatorname{succ}(\operatorname{hmk\_}))). \\ \quad \quad \quad \quad \quad (\operatorname{hmg\_}(\operatorname{hmw\_}) \in \operatorname{hmd\_}(\operatorname{hmk\_})) \wedge (\operatorname{hmr\_}(\operatorname{hmw\_}) \in \operatorname{hmd\_}(\operatorname{hmk\_})) \end{array}
\end{equation*}

\par\medskip\noindent\textbf{Definition} (\texttt{homotopic}).\;For all $\operatorname{hmu\_}$, $\operatorname{hmg\_}$, $\operatorname{hmr\_}$, $\operatorname{hma\_}$ and $\operatorname{hmb\_}$, 
\begin{equation*}
\begin{array}{@{}l@{}} \operatorname{homotopic}(\operatorname{hmu\_}, \operatorname{hmg\_}, \operatorname{hmr\_}, \operatorname{hma\_}, \operatorname{hmb\_}) \\ \Leftrightarrow\; \operatorname{is\text{-}path}(\operatorname{hmg\_}, \operatorname{hma\_}, \operatorname{hmb\_}) \wedge \\ \operatorname{is\text{-}path}(\operatorname{hmr\_}, \operatorname{hma\_}, \operatorname{hmb\_}) \wedge \\ \operatorname{hmr\_}(\operatorname{hma\_}) = \operatorname{hmg\_}(\operatorname{hma\_}) \wedge \\ \operatorname{hmr\_}(\operatorname{hmb\_}) = \operatorname{hmg\_}(\operatorname{hmb\_}) \wedge \\ \exists \operatorname{hmf\_} \in ((\operatorname{ccint}(0, 1) \times \operatorname{ccint}(\operatorname{hma\_}, \operatorname{hmb\_})) \to \operatorname{hmu\_}). \\ \quad \forall \operatorname{hmz\_} \in (\operatorname{ccint}(0, 1) \times \operatorname{ccint}(\operatorname{hma\_}, \operatorname{hmb\_})). \\ \quad \quad \operatorname{is\text{-}continuous\text{-}at}(\operatorname{product2\text{-}ms}(\operatorname{subspace\text{-}ms}(\operatorname{rr\text{-}ms}, \operatorname{ccint}(0, 1)), \operatorname{subspace\text{-}ms}(\operatorname{rr\text{-}ms}, \operatorname{ccint}(\operatorname{hma\_}, \operatorname{hmb\_}))), \\ \quad \quad \quad \operatorname{cc\text{-}ms}, \operatorname{hmf\_}, \operatorname{hmz\_}) \wedge \\ \quad \forall \operatorname{hmt\_} \in \operatorname{ccint}(\operatorname{hma\_}, \operatorname{hmb\_}). \\ \quad \quad (\operatorname{hmf\_}({[}0, \operatorname{hmt\_}]) = \operatorname{hmg\_}(\operatorname{hmt\_})) \wedge (\operatorname{hmf\_}({[}1, \operatorname{hmt\_}]) = \operatorname{hmr\_}(\operatorname{hmt\_})) \wedge \\ \quad \forall \operatorname{hms\_} \in \operatorname{ccint}(0, 1). \\ \quad \quad (\operatorname{hmf\_}({[}\operatorname{hms\_}, \operatorname{hma\_}]) = \operatorname{hmg\_}(\operatorname{hma\_})) \wedge (\operatorname{hmf\_}({[}\operatorname{hms\_}, \operatorname{hmb\_}]) = \operatorname{hmg\_}(\operatorname{hmb\_})) \end{array}
\end{equation*}

\par\medskip\noindent\textbf{Definition} (\texttt{null-homotopic}).\;For all $\operatorname{hmu\_}$, $\operatorname{hmg\_}$, $\operatorname{hma\_}$ and $\operatorname{hmb\_}$, 
\begin{equation*}
\begin{array}{@{}l@{}} \operatorname{null\text{-}homotopic}(\operatorname{hmu\_}, \operatorname{hmg\_}, \operatorname{hma\_}, \operatorname{hmb\_}) \\ \Leftrightarrow\; \operatorname{hmg\_}(\operatorname{hmb\_}) = \operatorname{hmg\_}(\operatorname{hma\_}) \wedge \\ \operatorname{homotopic}(\operatorname{hmu\_}, \operatorname{hmg\_}, (\lambda \operatorname{hmx\_} \in \operatorname{ccint}(\operatorname{hma\_}, \operatorname{hmb\_}).\; \operatorname{hmg\_}(\operatorname{hma\_})), \operatorname{hma\_}, \operatorname{hmb\_}) \end{array}
\end{equation*}

\par\medskip\noindent\textbf{Definition} (\texttt{simply-connected}).\;For every $\operatorname{hmu\_}$, 
\begin{equation*}
\begin{array}{@{}l@{}} \operatorname{simply\text{-}connected}(\operatorname{hmu\_}) \\ \Leftrightarrow\; \operatorname{is\text{-}open}(\operatorname{nf\text{-}metric\text{-}space}(\operatorname{cc\text{-}normed\text{-}field}), \operatorname{hmu\_}) \wedge \\ \operatorname{is\text{-}connected}(\operatorname{cc\text{-}ms}, \operatorname{hmu\_}) \wedge \\ \forall \operatorname{hma\_}, \operatorname{hmb\_}, \operatorname{hmg\_}. \\ \quad \operatorname{is\text{-}path}(\operatorname{hmg\_}, \operatorname{hma\_}, \operatorname{hmb\_}) \Rightarrow  \\ \quad (\operatorname{trace}(\operatorname{hmg\_}, \operatorname{hma\_}, \operatorname{hmb\_}) \subseteq \operatorname{hmu\_}) \Rightarrow  \\ \quad (\operatorname{hmg\_}(\operatorname{hmb\_}) = \operatorname{hmg\_}(\operatorname{hma\_})) \Rightarrow \operatorname{null\text{-}homotopic}(\operatorname{hmu\_}, \operatorname{hmg\_}, \operatorname{hma\_}, \operatorname{hmb\_}) \end{array}
\end{equation*}

\par\medskip\noindent\textbf{Definition} (\texttt{aff-seg}).\;$\operatorname{aff\text{-}seg}(\operatorname{hmp\_}, \operatorname{hmq\_}, \operatorname{hmc\_}, \operatorname{hmd\_}) :=$
\begin{equation*}
(\lambda \operatorname{hmz\_} \in \operatorname{ccint}(\operatorname{hmc\_}, \operatorname{hmd\_}).\; (\operatorname{hmp\_} + (((\operatorname{hmz\_} - \operatorname{hmc\_}) \cdot \operatorname{recip}((\operatorname{hmd\_} - \operatorname{hmc\_}))) \cdot (\operatorname{hmq\_} - \operatorname{hmp\_}))))
\end{equation*}

\par\medskip\noindent\textbf{Definition} (\texttt{poly-road}).\;For all $\operatorname{hmt\_}$ and $\operatorname{hmx\_}$, 
\begin{equation*}
(\operatorname{poly\text{-}road}(\operatorname{hmt\_}, \operatorname{hmx\_}, 0) \simeq \operatorname{aff\text{-}seg}(\operatorname{hmx\_}(0), \operatorname{hmx\_}(1), \operatorname{hmt\_}(0), \operatorname{hmt\_}(1)))
\end{equation*}
\par
Let $\operatorname{hmt\_}$ and $\operatorname{hmx\_}$ be arbitrary. Let $\operatorname{hmm\_} \in \mathbb{N}$. Then
\begin{equation*}
(\operatorname{poly\text{-}road}(\operatorname{hmt\_}, \operatorname{hmx\_}, \operatorname{succ}(\operatorname{hmm\_})) \simeq \operatorname{juxta}(\operatorname{poly\text{-}road}(\operatorname{hmt\_}, \operatorname{hmx\_}, \operatorname{hmm\_}), \operatorname{aff\text{-}seg}(\operatorname{hmx\_}(\operatorname{succ}(\operatorname{hmm\_})), \operatorname{hmx\_}(\operatorname{succ}(\operatorname{succ}(\operatorname{hmm\_}))), \operatorname{hmt\_}(\operatorname{succ}(\operatorname{hmm\_})), \operatorname{hmt\_}(\operatorname{succ}(\operatorname{succ}(\operatorname{hmm\_})))), \operatorname{hmt\_}(0), \operatorname{hmt\_}(\operatorname{succ}(\operatorname{hmm\_})), \operatorname{hmt\_}(\operatorname{succ}(\operatorname{succ}(\operatorname{hmm\_})))))
\end{equation*}

\par\medskip\noindent\textbf{Definition} (\texttt{aff-der}).\;$\operatorname{aff\text{-}der}(\operatorname{hmp\_}, \operatorname{hmq\_}, \operatorname{hmc\_}, \operatorname{hmd\_}) :=$
\begin{equation*}
(\lambda \operatorname{hmz\_} \in \operatorname{ccint}(\operatorname{hmc\_}, \operatorname{hmd\_}).\; (\operatorname{recip}((\operatorname{hmd\_} - \operatorname{hmc\_})) \cdot (\operatorname{hmq\_} - \operatorname{hmp\_})))
\end{equation*}

\par\medskip\noindent\textbf{Definition} (\texttt{poly-der}).\;For all $\operatorname{hmt\_}$ and $\operatorname{hmx\_}$, 
\begin{equation*}
(\operatorname{poly\text{-}der}(\operatorname{hmt\_}, \operatorname{hmx\_}, 0) \simeq \operatorname{aff\text{-}der}(\operatorname{hmx\_}(0), \operatorname{hmx\_}(1), \operatorname{hmt\_}(0), \operatorname{hmt\_}(1)))
\end{equation*}
\par
Let $\operatorname{hmt\_}$ and $\operatorname{hmx\_}$ be arbitrary. Let $\operatorname{hmm\_} \in \mathbb{N}$. Then
\begin{equation*}
(\operatorname{poly\text{-}der}(\operatorname{hmt\_}, \operatorname{hmx\_}, \operatorname{succ}(\operatorname{hmm\_})) \simeq \operatorname{juxta}(\operatorname{poly\text{-}der}(\operatorname{hmt\_}, \operatorname{hmx\_}, \operatorname{hmm\_}), \operatorname{aff\text{-}der}(\operatorname{hmx\_}(\operatorname{succ}(\operatorname{hmm\_})), \operatorname{hmx\_}(\operatorname{succ}(\operatorname{succ}(\operatorname{hmm\_}))), \operatorname{hmt\_}(\operatorname{succ}(\operatorname{hmm\_})), \operatorname{hmt\_}(\operatorname{succ}(\operatorname{succ}(\operatorname{hmm\_})))), \operatorname{hmt\_}(0), \operatorname{hmt\_}(\operatorname{succ}(\operatorname{hmm\_})), \operatorname{hmt\_}(\operatorname{succ}(\operatorname{succ}(\operatorname{hmm\_})))))
\end{equation*}


\vnbbookfile{theorem-library/good-approx}{theorem-library/good-approx}
\noindent{\small ch3c-ga.scm -- Proposition 3.12 of complex-analysis.pdf: the integrals of a holomorphic f agree along a good approximation.  Prefix cg-.  Library: complex-analysis. Destination: theorem-library/good-approx.scm, after ch3c-vocab (structure-library/ca-homotopy-vocab).}\par
\vnbprop{trace-restrict-in}{\vnbstmt{\begin{array}{@{}l@{}} \forall p, s, c, e. \\ \quad \forall w \in \operatorname{ccint}(c, e).\; (p(w) \in s) \Rightarrow  \\ \quad \quad \operatorname{trace}(\operatorname{restrict}(p, \operatorname{ccint}(c, e)), c, e) \subseteq s \end{array}}}
\vnbprop{convex-primitive-line-int}{\vnbstmt{\begin{array}{@{}l@{}} \forall u, f, s, z. \\ \quad \operatorname{holomorphic\text{-}on}(u, f) \wedge \\ \quad \operatorname{is\text{-}open}(\operatorname{nf\text{-}metric\text{-}space}(\operatorname{cc\text{-}normed\text{-}field}), s) \wedge \\ \quad \operatorname{is\text{-}convex\text{-}cc}(s) \wedge \\ \quad s \subseteq u \wedge \\ \quad z \in s \Rightarrow \\ \quad \quad \exists g \in (s \to \mathbb{C}). \forall p, d, c, e. \\ \quad \quad \quad (\operatorname{is\text{-}road}(p, d, c, e) \wedge (\operatorname{trace}(p, c, e) \subseteq s)) \Rightarrow  \\ \quad \quad \quad \quad \operatorname{line\text{-}int}(f, p, d, c, e) = (g(p(e)) - g(p(c))) \end{array}}}
\vnbprop{good-approx-piece}{\vnbstmt{\begin{array}{@{}l@{}} \forall u, f, s, m, c, o, d, a, b, \mathit{cgc}, e. \\ \quad \operatorname{holomorphic\text{-}on}(u, f) \wedge \\ \quad \operatorname{is\text{-}open}(\operatorname{nf\text{-}metric\text{-}space}(\operatorname{cc\text{-}normed\text{-}field}), s) \wedge \\ \quad \operatorname{is\text{-}convex\text{-}cc}(s) \wedge \\ \quad s \subseteq u \wedge \\ \quad \operatorname{is\text{-}road}(m, c, a, b) \wedge \\ \quad \operatorname{is\text{-}road}(o, d, a, b) \wedge \\ \quad \mathit{cgc} \in \mathbb{R} \wedge \\ \quad e \in \mathbb{R} \wedge \\ \quad a \leq \mathit{cgc} \wedge \\ \quad \mathit{cgc} < e \wedge \\ \quad e \leq b \wedge \\ \quad \forall w \in \operatorname{ccint}(\mathit{cgc}, e).\; ((m(w) \in s) \wedge (o(w) \in s)) \Rightarrow \\ \quad \quad (\operatorname{line\text{-}int}(f, \operatorname{restrict}(m, \operatorname{ccint}(\mathit{cgc}, e)), \operatorname{restrict}(c, \operatorname{ccint}(\mathit{cgc}, e)), \mathit{cgc}, e) - \operatorname{line\text{-}int}(f, \operatorname{restrict}(o, \operatorname{ccint}(\mathit{cgc}, e)), \operatorname{restrict}(d, \operatorname{ccint}(\mathit{cgc}, e)), \mathit{cgc}, e)) \\ \quad \quad =\; (\operatorname{seg\text{-}int}(f, o(e), m(e)) - \operatorname{seg\text{-}int}(f, o(\mathit{cgc}), m(\mathit{cgc}))) \end{array}}}
\vnbprop{restricted-line-int-split}{\vnbstmt{\begin{array}{@{}l@{}} \forall u, f, m, c, a, b, \mathit{cgc}, e. \\ \quad \operatorname{holomorphic\text{-}on}(u, f) \wedge \\ \quad \operatorname{is\text{-}road}(m, c, a, b) \wedge \\ \quad \operatorname{trace}(m, a, b) \subseteq u \wedge \\ \quad \mathit{cgc} \in \mathbb{R} \wedge \\ \quad e \in \mathbb{R} \wedge \\ \quad a < \mathit{cgc} \wedge \\ \quad \mathit{cgc} < e \wedge \\ \quad e \leq b \Rightarrow \\ \quad \quad \operatorname{line\text{-}int}(f, \operatorname{restrict}(m, \operatorname{ccint}(a, e)), \operatorname{restrict}(c, \operatorname{ccint}(a, e)), a, e) \\ \quad \quad =\; \operatorname{line\text{-}int}(f, \operatorname{restrict}(m, \operatorname{ccint}(a, \mathit{cgc})), \operatorname{restrict}(c, \operatorname{ccint}(a, \mathit{cgc})), a, \mathit{cgc}) + \operatorname{line\text{-}int}(f, \operatorname{restrict}(m, \operatorname{ccint}(\mathit{cgc}, e)), \operatorname{restrict}(c, \operatorname{ccint}(\mathit{cgc}, e)), \mathit{cgc}, e) \end{array}}}
\vnbprop{increasing-seq-from-start}{\vnbstmt{\begin{array}{@{}l@{}} \forall j \in \mathbb{N}, n, q. \\ \quad ((n \in \mathbb{N}) \wedge ((q \in (\mathbb{N} \to \mathbb{R})) \wedge (\forall k \in \mathbb{N}.\; ((k < n) \Rightarrow (q(k) < q((k + 1)))) \wedge (j \leq n)))) \Rightarrow  \\ \quad \quad (q(0) \leq q(j)) \wedge ((1 \leq j) \Rightarrow (q(0) < q(j))) \end{array}}}
\vnbprop{increasing-seq-le}{\vnbstmt{\begin{array}{@{}l@{}} \forall m \in \mathbb{N}, j, n, q. \\ \quad j \in \mathbb{N} \wedge \\ \quad n \in \mathbb{N} \wedge \\ \quad q \in (\mathbb{N} \to \mathbb{R}) \wedge \\ \quad \forall k \in \mathbb{N}.\; ((k < n) \Rightarrow (q(k) < q((k + 1)))) \wedge \\ \quad j \leq m \wedge \\ \quad m \leq n \Rightarrow \\ \quad \quad q(j) \leq q(m) \end{array}}}
\vnbprop{good-approx-invariant}{\vnbstmt{\begin{array}{@{}l@{}} \forall k \in \mathbb{N}, u, f, m, c, o, d, a, b, n, q, s. \\ \quad \operatorname{holomorphic\text{-}on}(u, f) \wedge \\ \quad \operatorname{is\text{-}road}(m, c, a, b) \wedge \\ \quad \operatorname{is\text{-}road}(o, d, a, b) \wedge \\ \quad \operatorname{trace}(m, a, b) \subseteq u \wedge \\ \quad \operatorname{trace}(o, a, b) \subseteq u \wedge \\ \quad o(a) = m(a) \wedge \\ \quad n \in \mathbb{N} \wedge \\ \quad q \in (\mathbb{N} \to \mathbb{R}) \wedge \\ \quad q(0) = a \wedge \\ \quad q(n) = b \wedge \\ \quad s \in (\mathbb{N} \to \operatorname{power}(\mathbb{C})) \wedge \\ \quad \forall \mathit{hmk} \in \mathbb{N}. \\ \quad \quad (\mathit{hmk} < n) \Rightarrow  \\ \quad \quad \quad q(\mathit{hmk}) < q((\mathit{hmk} + 1)) \wedge \\ \quad \quad \quad \operatorname{is\text{-}open}(\operatorname{nf\text{-}metric\text{-}space}(\operatorname{cc\text{-}normed\text{-}field}), s(\mathit{hmk})) \wedge \\ \quad \quad \quad \operatorname{is\text{-}convex\text{-}cc}(s(\mathit{hmk})) \wedge \\ \quad \quad \quad s(\mathit{hmk}) \subseteq u \wedge \\ \quad \quad \quad \forall w \in \operatorname{ccint}(q(\mathit{hmk}), q((\mathit{hmk} + 1))).\; ((m(w) \in s(\mathit{hmk})) \wedge (o(w) \in s(\mathit{hmk}))) \wedge \\ \quad 1 \leq k \wedge \\ \quad k \leq n \Rightarrow \\ \quad \quad (\operatorname{line\text{-}int}(f, \operatorname{restrict}(m, \operatorname{ccint}(a, q(k))), \operatorname{restrict}(c, \operatorname{ccint}(a, q(k))), a, q(k)) - \operatorname{line\text{-}int}(f, \operatorname{restrict}(o, \operatorname{ccint}(a, q(k))), \operatorname{restrict}(d, \operatorname{ccint}(a, q(k))), a, q(k))) \\ \quad \quad =\; \operatorname{seg\text{-}int}(f, o(q(k)), m(q(k))) \end{array}}}
\vnbprop{good-approx-integral}{\vnbstmt{\begin{array}{@{}l@{}} \forall u, f, m, c, o, d, a, b. \\ \quad (\operatorname{holomorphic\text{-}on}(u, f) \wedge (\operatorname{is\text{-}road}(m, c, a, b) \wedge (\operatorname{is\text{-}road}(o, d, a, b) \wedge \operatorname{good\text{-}approx}(u, m, o, a, b)))) \Rightarrow  \\ \quad \quad \operatorname{line\text{-}int}(f, m, c, a, b) = \operatorname{line\text{-}int}(f, o, d, a, b) \end{array}}}
\vnbprop{const-road}{\vnbstmt{\begin{array}{@{}l@{}} \forall c, a, b. \\ \quad ((c \in \mathbb{C}) \wedge ((a \in \mathbb{R}) \wedge ((b \in \mathbb{R}) \wedge (a < b)))) \Rightarrow  \\ \quad \quad \operatorname{is\text{-}road}((\lambda q \in \operatorname{ccint}(a, b).\; c), (\lambda q \in \operatorname{ccint}(a, b).\; 0), a, b) \end{array}}}
\vnbprop{const-road-line-int}{\vnbstmt{\begin{array}{@{}l@{}} \forall u, f, c, a, b. \\ \quad (\operatorname{holomorphic\text{-}on}(u, f) \wedge ((c \in u) \wedge ((a \in \mathbb{R}) \wedge ((b \in \mathbb{R}) \wedge (a < b))))) \Rightarrow  \\ \quad \quad \operatorname{line\text{-}int}(f, (\lambda q \in \operatorname{ccint}(a, b).\; c), (\lambda q \in \operatorname{ccint}(a, b).\; 0), a, b) = 0 \end{array}}}

\vnbbookfile{theorem-library/poly-road}{theorem-library/poly-road}
\noindent{\small ch3c-poly.scm -- polygonal roads on a partition, and Example 3.11 of complex-analysis.pdf. Prefix cp-.  Library: complex-analysis.  Destination: theorem-library/poly-road.scm. Loads after ch3c-vocab (AFF-SEG, AFF-DER, POLY-ROAD, POLY-DER) and ch3c-ga.}\par
\vnbprop{aff-seg-is-road}{\vnbstmt{\begin{array}{@{}l@{}} \forall p, q, c, d. \\ \quad ((p \in \mathbb{C}) \wedge ((q \in \mathbb{C}) \wedge ((c \in \mathbb{R}) \wedge ((d \in \mathbb{R}) \wedge (c < d))))) \Rightarrow  \\ \quad \quad \operatorname{is\text{-}road}(\operatorname{aff\text{-}seg}(p, q, c, d), \operatorname{aff\text{-}der}(p, q, c, d), c, d) \end{array}}}
\vnbprop{aff-seg-ends}{\vnbstmt{\begin{array}{@{}l@{}} \forall p, q, c, d. \\ \quad ((p \in \mathbb{C}) \wedge ((q \in \mathbb{C}) \wedge ((c \in \mathbb{R}) \wedge ((d \in \mathbb{R}) \wedge (c < d))))) \Rightarrow  \\ \quad \quad (\operatorname{aff\text{-}seg}(p, q, c, d)(c) = p) \wedge (\operatorname{aff\text{-}seg}(p, q, c, d)(d) = q) \end{array}}}
\vnbprop{aff-seg-in-convex}{\vnbstmt{\begin{array}{@{}l@{}} \forall p, q, c, d, s, z. \\ \quad p \in \mathbb{C} \wedge \\ \quad q \in \mathbb{C} \wedge \\ \quad c \in \mathbb{R} \wedge \\ \quad d \in \mathbb{R} \wedge \\ \quad c < d \wedge \\ \quad \operatorname{is\text{-}convex\text{-}cc}(s) \wedge \\ \quad p \in s \wedge \\ \quad q \in s \wedge \\ \quad z \in \operatorname{ccint}(c, d) \Rightarrow \\ \quad \quad \operatorname{aff\text{-}seg}(p, q, c, d)(z) \in s \end{array}}}
\vnbprop{poly-road-facts}{\vnbstmt{\begin{array}{@{}l@{}} \forall m \in \mathbb{N}, n, q, a. \\ \quad n \in \mathbb{N} \wedge \\ \quad q \in (\mathbb{N} \to \mathbb{R}) \wedge \\ \quad a \in (\mathbb{N} \to \mathbb{C}) \wedge \\ \quad \forall \mathit{cgk} \in \mathbb{N}.\; ((\mathit{cgk} < n) \Rightarrow (q(\mathit{cgk}) < q((\mathit{cgk} + 1)))) \wedge \\ \quad m < n \Rightarrow \\ \quad \quad \operatorname{is\text{-}road}(\operatorname{poly\text{-}road}(q, a, m), \operatorname{poly\text{-}der}(q, a, m), q(0), q((m + 1))) \wedge \\ \quad \quad \forall k \in \mathbb{N}, h \in \operatorname{ccint}(q(k), q((k + 1))). \\ \quad \quad \quad (k \leq m) \Rightarrow  \\ \quad \quad \quad \quad (\operatorname{poly\text{-}road}(q, a, m)(h) \simeq \operatorname{aff\text{-}seg}(a(k), a((k + 1)), q(k), q((k + 1)))(h)) \end{array}}}
\vnbprop{poly-piece-locate}{\vnbstmt{\begin{array}{@{}l@{}} \forall m \in \mathbb{N}, n, q.; \forall h \in \operatorname{ccint}(q(0), q((m + 1))). \\ \quad n \in \mathbb{N} \wedge \\ \quad q \in (\mathbb{N} \to \mathbb{R}) \wedge \\ \quad \forall \mathit{cgk} \in \mathbb{N}.\; ((\mathit{cgk} < n) \Rightarrow (q(\mathit{cgk}) < q((\mathit{cgk} + 1)))) \wedge \\ \quad m < n \Rightarrow \\ \quad \quad \exists k \in \mathbb{N}.\; ((k \leq m) \wedge (h \in \operatorname{ccint}(q(k), q((k + 1))))) \end{array}}}
\vnbprop{good-approx-polygon}{\vnbstmt{\begin{array}{@{}l@{}} \forall u, c, a, b, m, q, d, s. \\ \quad \operatorname{is\text{-}path}(c, a, b) \wedge \\ \quad \operatorname{trace}(c, a, b) \subseteq u \wedge \\ \quad m \in \mathbb{N} \wedge \\ \quad q \in (\mathbb{N} \to \mathbb{R}) \wedge \\ \quad d \in (\mathbb{N} \to \mathbb{C}) \wedge \\ \quad s \in (\mathbb{N} \to \operatorname{power}(\mathbb{C})) \wedge \\ \quad q(0) = a \wedge \\ \quad q((m + 1)) = b \wedge \\ \quad d(0) = c(a) \wedge \\ \quad d((m + 1)) = c(b) \wedge \\ \quad \forall k \in \mathbb{N}. \\ \quad \quad (k < (m + 1)) \Rightarrow  \\ \quad \quad \quad q(k) < q((k + 1)) \wedge \\ \quad \quad \quad \operatorname{is\text{-}open}(\operatorname{nf\text{-}metric\text{-}space}(\operatorname{cc\text{-}normed\text{-}field}), s(k)) \wedge \\ \quad \quad \quad \operatorname{is\text{-}convex\text{-}cc}(s(k)) \wedge \\ \quad \quad \quad s(k) \subseteq u \wedge \\ \quad \quad \quad d(k) \in s(k) \wedge \\ \quad \quad \quad d((k + 1)) \in s(k) \wedge \\ \quad \quad \quad \forall w \in \operatorname{ccint}(q(k), q((k + 1))).\; (c(w) \in s(k)) \Rightarrow \\ \quad \quad (\operatorname{is\text{-}road}(\operatorname{poly\text{-}road}(q, d, m), \operatorname{poly\text{-}der}(q, d, m), a, b) \wedge \operatorname{good\text{-}approx}(u, c, \operatorname{poly\text{-}road}(q, d, m), a, b)) \end{array}}}
\vnbprop{path-uniformly-continuous}{\vnbstmt{\begin{array}{@{}l@{}} \forall m, a, b, e. \\ \quad (\operatorname{is\text{-}path}(m, a, b) \wedge \operatorname{pos\text{-}rr}(e)) \Rightarrow  \\ \quad \quad \exists d. \\ \quad \quad \quad \operatorname{pos\text{-}rr}(d) \wedge \\ \quad \quad \quad \forall x, y \in \operatorname{ccint}(a, b). \\ \quad \quad \quad \quad (\lvert (x - y) \rvert \leq d) \Rightarrow (\operatorname{magnitude}((m(x) - m(y))) \leq e) \end{array}}}
\vnbprop{uniform-grid}{\vnbstmt{\begin{array}{@{}l@{}} \forall a, b, d. \\ \quad ((a \in \mathbb{R}) \wedge ((b \in \mathbb{R}) \wedge ((d \in \mathbb{R}) \wedge ((a < b) \wedge (0 < d))))) \Rightarrow  \\ \quad \quad \exists m \in \mathbb{N}, h \in \mathbb{R}.\; ((0 < h) \wedge ((h \leq d) \wedge ((a + ((m + 1) \cdot h)) = b))) \end{array}}}
\vnbprop{polygonal-good-approx}{\vnbstmt{\begin{array}{@{}l@{}} \forall u, m, a, b. \\ \quad (\operatorname{is\text{-}open}(\operatorname{nf\text{-}metric\text{-}space}(\operatorname{cc\text{-}normed\text{-}field}), u) \wedge (\operatorname{is\text{-}path}(m, a, b) \wedge (\operatorname{trace}(m, a, b) \subseteq u))) \Rightarrow  \\ \quad \quad \exists r, q.\; (\operatorname{is\text{-}road}(r, q, a, b) \wedge \operatorname{good\text{-}approx}(u, m, r, a, b)) \end{array}}}
\vnbprop{poly-trace-in}{\vnbstmt{\begin{array}{@{}l@{}} \forall u, a, b, m, q, c, s. \\ \quad m \in \mathbb{N} \wedge \\ \quad q \in (\mathbb{N} \to \mathbb{R}) \wedge \\ \quad c \in (\mathbb{N} \to \mathbb{C}) \wedge \\ \quad q(0) = a \wedge \\ \quad q((m + 1)) = b \wedge \\ \quad \operatorname{pos\text{-}rr}(s) \wedge \\ \quad \forall k \in \mathbb{N}. \\ \quad \quad (k < (m + 1)) \Rightarrow  \\ \quad \quad \quad ((q(k) < q((k + 1))) \wedge ((\operatorname{ball}(\operatorname{cc\text{-}ms}, c(k), s) \subseteq u) \wedge (c((k + 1)) \in \operatorname{ball}(\operatorname{cc\text{-}ms}, c(k), s)))) \Rightarrow \\ \quad \quad \operatorname{trace}(\operatorname{poly\text{-}road}(q, c, m), a, b) \subseteq u \end{array}}}

\vnbbookfile{theorem-library/homotopy}{theorem-library/homotopy}
\noindent{\small ch3c-hom.scm -- homotopy (complex-analysis.pdf 3.3.2): 3.16, 3.17, 3.18, 3.20, 3.31 (and 3.15). Prefix ch-.  Library: complex-analysis.  Destination: theorem-library/homotopy.scm. The rectangle [0, 1] x [a, b] is Q(a, b) = PRODUCT2-MS(SUBSPACE-MS(RR-MS, [0, 1]), SUBSPACE-MS(RR-MS, [a, b])).}\par
\vnbprop{rect-dist}{\vnbstmt{\begin{array}{@{}l@{}} \forall a, b, y, z. \\ \quad a \in \mathbb{R} \wedge \\ \quad b \in \mathbb{R} \wedge \\ \quad y \in (\operatorname{ccint}(0, 1) \times \operatorname{ccint}(a, b)) \wedge \\ \quad z \in (\operatorname{ccint}(0, 1) \times \operatorname{ccint}(a, b)) \Rightarrow \\ \quad \quad (\operatorname{dist}(\operatorname{product2\text{-}ms}(\operatorname{subspace\text{-}ms}(\operatorname{rr\text{-}ms}, \operatorname{ccint}(0, 1)), \operatorname{subspace\text{-}ms}(\operatorname{rr\text{-}ms}, \operatorname{ccint}(a, b))))(y, z) \simeq \operatorname{max}(\lvert (\operatorname{nth}(1, y) - \operatorname{nth}(1, z)) \rvert, \lvert (\operatorname{nth}(2, y) - \operatorname{nth}(2, z)) \rvert)) \end{array}}}
\vnbprop{rect-cont-intro}{\vnbstmt{\begin{array}{@{}l@{}} \forall a, b, f, x. \\ \quad a \in \mathbb{R} \wedge \\ \quad b \in \mathbb{R} \wedge \\ \quad a < b \wedge \\ \quad f \in ((\operatorname{ccint}(0, 1) \times \operatorname{ccint}(a, b)) \to \mathbb{C}) \wedge \\ \quad x \in (\operatorname{ccint}(0, 1) \times \operatorname{ccint}(a, b)) \wedge \\ \quad \forall e. \\ \quad \quad \operatorname{pos\text{-}rr}(e) \Rightarrow  \\ \quad \quad \quad \exists d. \\ \quad \quad \quad \quad \operatorname{pos\text{-}rr}(d) \wedge \\ \quad \quad \quad \quad \forall y \in (\operatorname{ccint}(0, 1) \times \operatorname{ccint}(a, b)). \\ \quad \quad \quad \quad \quad ((\lvert (\operatorname{nth}(1, x) - \operatorname{nth}(1, y)) \rvert \leq d) \wedge (\lvert (\operatorname{nth}(2, x) - \operatorname{nth}(2, y)) \rvert \leq d)) \Rightarrow  \\ \quad \quad \quad \quad \quad \quad \operatorname{magnitude}((f(x) - f(y))) \leq e \Rightarrow \\ \quad \quad \operatorname{is\text{-}continuous\text{-}at}(\operatorname{product2\text{-}ms}(\operatorname{subspace\text{-}ms}(\operatorname{rr\text{-}ms}, \operatorname{ccint}(0, 1)), \operatorname{subspace\text{-}ms}(\operatorname{rr\text{-}ms}, \operatorname{ccint}(a, b))), \\ \quad \quad \quad \operatorname{cc\text{-}ms}, f, x) \end{array}}}
\vnbprop{rect-unif}{\vnbstmt{\begin{array}{@{}l@{}} \forall a, b, f, e. \\ \quad a \in \mathbb{R} \wedge \\ \quad b \in \mathbb{R} \wedge \\ \quad a < b \wedge \\ \quad f \in ((\operatorname{ccint}(0, 1) \times \operatorname{ccint}(a, b)) \to \mathbb{C}) \wedge \\ \quad \forall x \in (\operatorname{ccint}(0, 1) \times \operatorname{ccint}(a, b)). \\ \quad \quad \operatorname{is\text{-}continuous\text{-}at}(\operatorname{product2\text{-}ms}(\operatorname{subspace\text{-}ms}(\operatorname{rr\text{-}ms}, \operatorname{ccint}(0, 1)), \operatorname{subspace\text{-}ms}(\operatorname{rr\text{-}ms}, \operatorname{ccint}(a, b))), \\ \quad \quad \quad \operatorname{cc\text{-}ms}, f, x) \wedge \\ \quad \operatorname{pos\text{-}rr}(e) \Rightarrow \\ \quad \quad \exists d. \\ \quad \quad \quad \operatorname{pos\text{-}rr}(d) \wedge \\ \quad \quad \quad \forall y, z \in (\operatorname{ccint}(0, 1) \times \operatorname{ccint}(a, b)). \\ \quad \quad \quad \quad ((\lvert (\operatorname{nth}(1, y) - \operatorname{nth}(1, z)) \rvert \leq d) \wedge (\lvert (\operatorname{nth}(2, y) - \operatorname{nth}(2, z)) \rvert \leq d)) \Rightarrow  \\ \quad \quad \quad \quad \quad \operatorname{magnitude}((f(y) - f(z))) \leq e \end{array}}}
\vnbprop{homotopic-refl}{\vnbstmt{\begin{array}{@{}l@{}} \forall u, m, a, b. \\ \quad (\operatorname{is\text{-}path}(m, a, b) \wedge (\operatorname{trace}(m, a, b) \subseteq u)) \Rightarrow \operatorname{homotopic}(u, m, m, a, b) \end{array}}}
\vnbprop{homotopic-sym}{\vnbstmt{\begin{array}{@{}l@{}} \forall u, m, o, a, b. \\ \quad \operatorname{homotopic}(u, m, o, a, b) \Rightarrow \operatorname{homotopic}(u, o, m, a, b) \end{array}}}
\vnbprop{homotopic-trans}{\vnbstmt{\begin{array}{@{}l@{}} \forall u, m, o, q, a, b. \\ \quad (\operatorname{homotopic}(u, m, o, a, b) \wedge \operatorname{homotopic}(u, o, q, a, b)) \Rightarrow  \\ \quad \quad \operatorname{homotopic}(u, m, q, a, b) \end{array}}}
\vnbprop{convex-comb-diff}{\vnbstmt{\begin{array}{@{}l@{}} \forall a, b, c, d, f, g. \\ \quad ((a \in \mathbb{C}) \wedge ((b \in \mathbb{C}) \wedge ((c \in \mathbb{C}) \wedge ((d \in \mathbb{C}) \wedge ((f \in \mathbb{R}) \wedge (g \in \mathbb{R})))))) \Rightarrow  \\ \quad \quad ((a + (f \cdot (b - a))) - (c + (g \cdot (d - c)))) = ((a - c) + ((g \cdot ((b - a) - (d - c))) + ((f - g) \cdot (b - a)))) \end{array}}}
\vnbprop{cc-unit-scale-le}{\vnbstmt{\begin{array}{@{}l@{}} \forall s, w. \\ \quad ((s \in \operatorname{ccint}(0, 1)) \wedge (w \in \mathbb{C})) \Rightarrow (\operatorname{magnitude}((s \cdot w)) \leq \operatorname{magnitude}(w)) \end{array}}}
\vnbprop{cc-diff-change-le}{\vnbstmt{\begin{array}{@{}l@{}} \forall a, b, c, d. \\ \quad ((a \in \mathbb{C}) \wedge ((b \in \mathbb{C}) \wedge ((c \in \mathbb{C}) \wedge (d \in \mathbb{C})))) \Rightarrow  \\ \quad \quad \operatorname{magnitude}(((a - b) - (c - d))) \leq (\operatorname{magnitude}((a - c)) + \operatorname{magnitude}((b - d))) \end{array}}}
\vnbprop{cc-scale-small}{\vnbstmt{\begin{array}{@{}l@{}} \forall w \in \mathbb{C}, e. \\ \quad \operatorname{pos\text{-}rr}(e) \Rightarrow  \\ \quad \quad \exists d. \\ \quad \quad \quad \operatorname{pos\text{-}rr}(d) \wedge \forall c \in \mathbb{R}.\; ((\lvert c \rvert \leq d) \Rightarrow (\operatorname{magnitude}((c \cdot w)) \leq e)) \end{array}}}
\vnbprop{good-approx-homotopic}{\vnbstmt{\begin{array}{@{}l@{}} \forall u, m, o, a, b. \\ \quad \operatorname{good\text{-}approx}(u, m, o, a, b) \Rightarrow \operatorname{homotopic}(u, m, o, a, b) \end{array}}}
\vnbprop{good-approx-balls}{\vnbstmt{\begin{array}{@{}l@{}} \forall u, c, a, b, m, q, d, s. \\ \quad \operatorname{is\text{-}path}(c, a, b) \wedge \\ \quad \operatorname{trace}(c, a, b) \subseteq u \wedge \\ \quad m \in \mathbb{N} \wedge \\ \quad q \in (\mathbb{N} \to \mathbb{R}) \wedge \\ \quad d \in (\mathbb{N} \to \mathbb{C}) \wedge \\ \quad q(0) = a \wedge \\ \quad q((m + 1)) = b \wedge \\ \quad d(0) = c(a) \wedge \\ \quad d((m + 1)) = c(b) \wedge \\ \quad \operatorname{pos\text{-}rr}(s) \wedge \\ \quad \forall k \in \mathbb{N}. \\ \quad \quad (k < (m + 1)) \Rightarrow  \\ \quad \quad \quad q(k) < q((k + 1)) \wedge \\ \quad \quad \quad \operatorname{ball}(\operatorname{cc\text{-}ms}, d(k), s) \subseteq u \wedge \\ \quad \quad \quad d((k + 1)) \in \operatorname{ball}(\operatorname{cc\text{-}ms}, d(k), s) \wedge \\ \quad \quad \quad \forall w \in \operatorname{ccint}(q(k), q((k + 1))).\; (c(w) \in \operatorname{ball}(\operatorname{cc\text{-}ms}, d(k), s)) \Rightarrow \\ \quad \quad (\operatorname{is\text{-}road}(\operatorname{poly\text{-}road}(q, d, m), \operatorname{poly\text{-}der}(q, d, m), a, b) \wedge \operatorname{good\text{-}approx}(u, c, \operatorname{poly\text{-}road}(q, d, m), a, b)) \end{array}}}
\vnbprop{homotopy-polygon-chain}{\vnbstmt{\begin{array}{@{}l@{}} \forall u, m, o, a, b. \\ \quad (\operatorname{is\text{-}open}(\operatorname{nf\text{-}metric\text{-}space}(\operatorname{cc\text{-}normed\text{-}field}), u) \wedge \operatorname{homotopic}(u, m, o, a, b)) \Rightarrow  \\ \quad \quad \exists \mathit{chm}, n \in \mathbb{N}, q \in (\mathbb{N} \to \mathbb{R}), \mathit{chxq} \in (\mathbb{N} \to (\mathbb{N} \to \mathbb{C})). \\ \quad \quad \quad \forall j \in \mathbb{N}. \\ \quad \quad \quad \quad (j \leq (n + 1)) \Rightarrow  \\ \quad \quad \quad \quad \quad \operatorname{is\text{-}road}(\operatorname{poly\text{-}road}(q, \operatorname{chxq}(j), \mathit{chm}), \operatorname{poly\text{-}der}(q, \operatorname{chxq}(j), \mathit{chm}), a, b) \wedge \\ \quad \quad \quad \operatorname{good\text{-}approx}(u, m, \operatorname{poly\text{-}road}(q, \operatorname{chxq}(0), \mathit{chm}), a, b) \wedge \\ \quad \quad \quad \forall j \in \mathbb{N}. \\ \quad \quad \quad \quad (j \leq n) \Rightarrow  \\ \quad \quad \quad \quad \quad \operatorname{good\text{-}approx}(u, \operatorname{poly\text{-}road}(q, \operatorname{chxq}(j), \mathit{chm}), \operatorname{poly\text{-}road}(q, \operatorname{chxq}((j + 1)), \mathit{chm}), a, b) \wedge \\ \quad \quad \quad \operatorname{good\text{-}approx}(u, o, \operatorname{poly\text{-}road}(q, \operatorname{chxq}((n + 1)), \mathit{chm}), a, b) \end{array}}}
\vnbprop{homotopy-chain-line-int}{\vnbstmt{\begin{array}{@{}l@{}} \forall q \in \mathbb{N}, u, f, m, c, a, b, \mathit{chm}, n, \mathit{chtq}, \mathit{chxq}. \\ \quad \operatorname{holomorphic\text{-}on}(u, f) \wedge \\ \quad \operatorname{is\text{-}road}(m, c, a, b) \wedge \\ \quad n \in \mathbb{N} \wedge \\ \quad \forall j \in \mathbb{N}. \\ \quad \quad (j \leq (n + 1)) \Rightarrow  \\ \quad \quad \quad \operatorname{is\text{-}road}(\operatorname{poly\text{-}road}(\mathit{chtq}, \operatorname{chxq}(j), \mathit{chm}), \operatorname{poly\text{-}der}(\mathit{chtq}, \operatorname{chxq}(j), \mathit{chm}), a, b) \wedge \\ \quad \operatorname{good\text{-}approx}(u, m, \operatorname{poly\text{-}road}(\mathit{chtq}, \operatorname{chxq}(0), \mathit{chm}), a, b) \wedge \\ \quad \forall j \in \mathbb{N}. \\ \quad \quad (j \leq n) \Rightarrow  \\ \quad \quad \quad \operatorname{good\text{-}approx}(u, \operatorname{poly\text{-}road}(\mathit{chtq}, \operatorname{chxq}(j), \mathit{chm}), \operatorname{poly\text{-}road}(\mathit{chtq}, \operatorname{chxq}((j + 1)), \mathit{chm}), a, b) \wedge \\ \quad q \leq (n + 1) \Rightarrow \\ \quad \quad \operatorname{line\text{-}int}(f, \operatorname{poly\text{-}road}(\mathit{chtq}, \operatorname{chxq}(q), \mathit{chm}), \operatorname{poly\text{-}der}(\mathit{chtq}, \operatorname{chxq}(q), \mathit{chm}), a, b) \\ \quad \quad =\; \operatorname{line\text{-}int}(f, m, c, a, b) \end{array}}}
\vnbprop{homotopy-line-int}{\vnbstmt{\begin{array}{@{}l@{}} \forall u, f, m, c, o, d, a, b. \\ \quad \operatorname{is\text{-}open}(\operatorname{nf\text{-}metric\text{-}space}(\operatorname{cc\text{-}normed\text{-}field}), u) \wedge \\ \quad \operatorname{holomorphic\text{-}on}(u, f) \wedge \\ \quad \operatorname{is\text{-}road}(m, c, a, b) \wedge \\ \quad \operatorname{is\text{-}road}(o, d, a, b) \wedge \\ \quad \operatorname{homotopic}(u, m, o, a, b) \Rightarrow \\ \quad \quad \operatorname{line\text{-}int}(f, m, c, a, b) = \operatorname{line\text{-}int}(f, o, d, a, b) \end{array}}}
\vnbprop{null-homotopic-line-int}{\vnbstmt{\begin{array}{@{}l@{}} \forall u, f, m, c, a, b. \\ \quad \operatorname{is\text{-}open}(\operatorname{nf\text{-}metric\text{-}space}(\operatorname{cc\text{-}normed\text{-}field}), u) \wedge \\ \quad \operatorname{holomorphic\text{-}on}(u, f) \wedge \\ \quad \operatorname{is\text{-}road}(m, c, a, b) \wedge \\ \quad \operatorname{null\text{-}homotopic}(u, m, a, b) \Rightarrow \\ \quad \quad \operatorname{line\text{-}int}(f, m, c, a, b) = 0 \end{array}}}
\vnbprop{homotopic-traces}{\vnbstmt{\begin{array}{@{}l@{}} \forall u, m, o, a, b. \\ \quad \operatorname{homotopic}(u, m, o, a, b) \Rightarrow  \\ \quad \quad (\operatorname{trace}(m, a, b) \subseteq u) \wedge (\operatorname{trace}(o, a, b) \subseteq u) \end{array}}}
\vnbprop{homotopy-winding}{\vnbstmt{\begin{array}{@{}l@{}} \forall z, m, c, o, d, a, b. \\ \quad z \in \mathbb{C} \wedge \\ \quad \operatorname{is\text{-}road}(m, c, a, b) \wedge \\ \quad \operatorname{is\text{-}road}(o, d, a, b) \wedge \\ \quad \operatorname{homotopic}(\operatorname{difference}(\mathbb{C}, \operatorname{singleton}(z)), m, o, a, b) \Rightarrow \\ \quad \quad \operatorname{winding}(m, c, a, b, z) = \operatorname{winding}(o, d, a, b, z) \end{array}}}

\vnbbookfile{theorem-library/ch4a-straight}{theorem-library/ch4a-straight}
\noindent{\small ch4a-straight.scm -- WAVE 2 (2026-10-27..30), wired 2026-10-30 from scratchpad/ch4a/ch4a-straight.scm. Library: complex-analysis.}\par
\vnbprop{straight-line-homotopic}{\vnbstmt{\begin{array}{@{}l@{}} \forall u, m, o, a, b. \\ \quad \operatorname{is\text{-}path}(m, a, b) \wedge \\ \quad \operatorname{is\text{-}path}(o, a, b) \wedge \\ \quad o(a) = m(a) \wedge \\ \quad o(b) = m(b) \wedge \\ \quad \forall t \in \operatorname{ccint}(a, b), s \in \operatorname{ccint}(0, 1).\; ((m(t) + (s \cdot (o(t) - m(t)))) \in u) \Rightarrow \\ \quad \quad \operatorname{homotopic}(u, m, o, a, b) \end{array}}}

\vnbbookfile{theorem-library/ch4a-annulus}{theorem-library/ch4a-annulus}
\noindent{\small ch4a-annulus.scm -- WAVE 2 (2026-10-27..30), wired 2026-10-30 from scratchpad/ch4a/ch4a-annulus.scm. Library: complex-analysis.}\par
\vnbprop{trace-in-of-values}{\vnbstmt{\begin{array}{@{}l@{}} \forall p, a, b, s. \\ \quad \forall t \in \operatorname{ccint}(a, b).\; (p(t) \in s) \Rightarrow (\operatorname{trace}(p, a, b) \subseteq s) \end{array}}}
\vnbprop{lerp-between}{\vnbstmt{\begin{array}{@{}l@{}} \forall x, y, m \in \mathbb{R}. \\ \quad ((x \leq y) \wedge ((0 \leq m) \wedge (m \leq 1))) \Rightarrow  \\ \quad \quad (x \leq (y + (m \cdot (x - y)))) \wedge ((y + (m \cdot (x - y))) \leq y) \end{array}}}
\vnbprop{unit-interval-mul}{\vnbstmt{\begin{array}{@{}l@{}} \forall p, q \in \mathbb{R}. \\ \quad ((0 \leq p) \wedge ((p \leq 1) \wedge ((0 \leq q) \wedge (q \leq 1)))) \Rightarrow ((0 \leq (p \cdot q)) \wedge ((p \cdot q) \leq 1)) \end{array}}}
\vnbprop{annulus-cauchy}{\vnbstmt{\begin{array}{@{}l@{}} \forall u, f.; \forall c \in \mathbb{C}, a, b \in \mathbb{R}. \\ \quad \operatorname{is\text{-}open}(\operatorname{nf\text{-}metric\text{-}space}(\operatorname{cc\text{-}normed\text{-}field}), u) \wedge \\ \quad \operatorname{holomorphic\text{-}on}(u, f) \wedge \\ \quad 0 < a \wedge \\ \quad a < b \wedge \\ \quad \forall z \in \mathbb{C}. \\ \quad \quad ((a \leq \operatorname{magnitude}((z - c))) \wedge (\operatorname{magnitude}((z - c)) \leq b)) \Rightarrow (z \in u) \Rightarrow \\ \quad \quad \operatorname{line\text{-}int}(f, \operatorname{circle\text{-}path}(c, a), \operatorname{circle\text{-}deriv}(c, a), 0, (2 \cdot \pi)) \\ \quad \quad =\; \operatorname{line\text{-}int}(f, \operatorname{circle\text{-}path}(c, b), \operatorname{circle\text{-}deriv}(c, b), 0, (2 \cdot \pi)) \end{array}}}

\vnbbookfile{theorem-library/ch4a-expansion}{theorem-library/ch4a-expansion}
\noindent{\small ch4a-expansion.scm -- WAVE 2 (2026-10-27..30), wired 2026-10-30 from scratchpad/ch4a/ch4a-expansion.scm. Library: complex-analysis.}\par
\vnbprop{laurent-circle-split}{\vnbstmt{\begin{array}{@{}l@{}} \forall u, f.; \forall c \in \mathbb{C}, r \in \mathbb{R}, z \in u, q \in (u \to \mathbb{C}). \\ \quad \operatorname{holomorphic\text{-}on}(u, f) \wedge \\ \quad 0 < r \wedge \\ \quad \operatorname{trace}(\operatorname{circle\text{-}path}(c, r), 0, (2 \cdot \pi)) \subseteq u \wedge \\ \quad \neg (\operatorname{magnitude}((z - c)) = r) \wedge \\ \quad \forall y \in u. \\ \quad \quad \operatorname{is\text{-}continuous\text{-}at}(\operatorname{subspace\text{-}ms}(\operatorname{nf\text{-}metric\text{-}space}(\operatorname{cc\text{-}normed\text{-}field}), u), \\ \quad \quad \quad \operatorname{nf\text{-}metric\text{-}space}(\operatorname{cc\text{-}normed\text{-}field}), q, y) \wedge \\ \quad \forall x \in u.\; ((f(x) - f(z)) = (q(x) \cdot (x - z))) \Rightarrow \\ \quad \quad \operatorname{line\text{-}int}((\lambda w \in \operatorname{trace}(\operatorname{circle\text{-}path}(c, r), 0, (2 \cdot \pi)).\; (f(w) \cdot \operatorname{recip}((w - z)))), \\ \quad \quad \quad \operatorname{circle\text{-}path}(c, r), \operatorname{circle\text{-}deriv}(c, r), 0, (2 \cdot \pi)) \\ \quad \quad =\; \operatorname{line\text{-}int}(q, \operatorname{circle\text{-}path}(c, r), \operatorname{circle\text{-}deriv}(c, r), 0, (2 \cdot \pi)) + (f(z) \cdot \operatorname{line\text{-}int}((\lambda \mathit{wnw} \in \operatorname{difference}(\mathbb{C}, \operatorname{singleton}(z)).\; \operatorname{recip}((\mathit{wnw} - z))), \operatorname{circle\text{-}path}(c, r), \operatorname{circle\text{-}deriv}(c, r), 0, (2 \cdot \pi))) \end{array}}}
\vnbprop{laurent-cauchy-difference}{\vnbstmt{\begin{array}{@{}l@{}} \forall u, f.; \forall c \in \mathbb{C}, a, b \in \mathbb{R}, z \in \mathbb{C}. \\ \quad \operatorname{is\text{-}open}(\operatorname{nf\text{-}metric\text{-}space}(\operatorname{cc\text{-}normed\text{-}field}), u) \wedge \\ \quad \operatorname{holomorphic\text{-}on}(u, f) \wedge \\ \quad 0 < a \wedge \\ \quad a < b \wedge \\ \quad \forall \mathit{lxb} \in \mathbb{C}. \\ \quad \quad ((a \leq \operatorname{magnitude}((\mathit{lxb} - c))) \wedge (\operatorname{magnitude}((\mathit{lxb} - c)) \leq b)) \Rightarrow (\mathit{lxb} \in u) \wedge \\ \quad a < \operatorname{magnitude}((z - c)) \wedge \\ \quad \operatorname{magnitude}((z - c)) < b \Rightarrow \\ \quad \quad ((2 \cdot \pi) \cdot +i) \cdot f(z) \\ \quad \quad =\; (\operatorname{line\text{-}int}((\lambda w \in \operatorname{trace}(\operatorname{circle\text{-}path}(c, b), 0, (2 \cdot \pi)).\; (f(w) \cdot \operatorname{recip}((w - z)))), \operatorname{circle\text{-}path}(c, b), \operatorname{circle\text{-}deriv}(c, b), 0, (2 \cdot \pi)) - \operatorname{line\text{-}int}((\lambda w \in \operatorname{trace}(\operatorname{circle\text{-}path}(c, a), 0, (2 \cdot \pi)).\; (f(w) \cdot \operatorname{recip}((w - z)))), \operatorname{circle\text{-}path}(c, a), \operatorname{circle\text{-}deriv}(c, a), 0, (2 \cdot \pi))) \end{array}}}
\vnbprop{laurent-expansion-annulus}{\vnbstmt{\begin{array}{@{}l@{}} \forall u, f.; \forall c \in \mathbb{C}, a, b \in \mathbb{R}, z \in \mathbb{C}. \\ \quad \operatorname{is\text{-}open}(\operatorname{nf\text{-}metric\text{-}space}(\operatorname{cc\text{-}normed\text{-}field}), u) \wedge \\ \quad \operatorname{holomorphic\text{-}on}(u, f) \wedge \\ \quad 0 < a \wedge \\ \quad a < b \wedge \\ \quad \forall \mathit{lxb} \in \mathbb{C}. \\ \quad \quad ((a \leq \operatorname{magnitude}((\mathit{lxb} - c))) \wedge (\operatorname{magnitude}((\mathit{lxb} - c)) \leq b)) \Rightarrow (\mathit{lxb} \in u) \wedge \\ \quad a < \operatorname{magnitude}((z - c)) \wedge \\ \quad \operatorname{magnitude}((z - c)) < b \Rightarrow \\ \quad \quad \operatorname{cc\text{-}series\text{-}converges\text{-}to}((\lambda k \in \mathbb{N}.\; \\ \quad \quad \quad \quad \operatorname{line\text{-}int}((\lambda w \in \operatorname{trace}(\operatorname{circle\text{-}path}(c, b), 0, (2 \cdot \pi)).\; (f(w) \cdot \operatorname{power}(\operatorname{recip}((w - c)), (k + 1)))), \operatorname{circle\text{-}path}(c, b), \operatorname{circle\text{-}deriv}(c, b), 0, (2 \cdot \pi)) \cdot \operatorname{power}((z - c), k)), \\ \quad \quad \quad \operatorname{line\text{-}int}((\lambda w \in \operatorname{trace}(\operatorname{circle\text{-}path}(c, b), 0, (2 \cdot \pi)).\; (f(w) \cdot \operatorname{recip}((w - z)))), \\ \quad \quad \quad \quad \operatorname{circle\text{-}path}(c, b), \operatorname{circle\text{-}deriv}(c, b), 0, (2 \cdot \pi))) \wedge \\ \quad \quad \operatorname{cc\text{-}series\text{-}converges\text{-}to}((\lambda k \in \mathbb{N}.\; \\ \quad \quad \quad \quad \operatorname{line\text{-}int}((\lambda w \in \operatorname{trace}(\operatorname{circle\text{-}path}(c, a), 0, (2 \cdot \pi)).\; (f(w) \cdot \operatorname{power}((w - c), k))), \operatorname{circle\text{-}path}(c, a), \operatorname{circle\text{-}deriv}(c, a), 0, (2 \cdot \pi)) \cdot \operatorname{power}(\operatorname{recip}((z - c)), (k + 1))), \\ \quad \quad \quad \operatorname{line\text{-}int}((\lambda w \in \operatorname{trace}(\operatorname{circle\text{-}path}(c, a), 0, (2 \cdot \pi)).\; (f(w) \cdot \operatorname{recip}((z - w)))), \\ \quad \quad \quad \quad \operatorname{circle\text{-}path}(c, a), \operatorname{circle\text{-}deriv}(c, a), 0, (2 \cdot \pi))) \wedge \\ \quad \quad ((2 \cdot \pi) \cdot +i) \cdot f(z) \\ \quad \quad =\; \operatorname{line\text{-}int}((\lambda w \in \operatorname{trace}(\operatorname{circle\text{-}path}(c, b), 0, (2 \cdot \pi)).\; (f(w) \cdot \operatorname{recip}((w - z)))), \operatorname{circle\text{-}path}(c, b), \operatorname{circle\text{-}deriv}(c, b), 0, (2 \cdot \pi)) + \operatorname{line\text{-}int}((\lambda w \in \operatorname{trace}(\operatorname{circle\text{-}path}(c, a), 0, (2 \cdot \pi)).\; (f(w) \cdot \operatorname{recip}((z - w)))), \operatorname{circle\text{-}path}(c, a), \operatorname{circle\text{-}deriv}(c, a), 0, (2 \cdot \pi)) \end{array}}}
\vnbprop{laurent-coef-radius-free}{\vnbstmt{\begin{array}{@{}l@{}} \forall u, f.; \forall c \in \mathbb{C}, a, b \in \mathbb{R}, k \in \mathbb{N}. \\ \quad \operatorname{is\text{-}open}(\operatorname{nf\text{-}metric\text{-}space}(\operatorname{cc\text{-}normed\text{-}field}), u) \wedge \\ \quad \operatorname{holomorphic\text{-}on}(u, f) \wedge \\ \quad 0 < a \wedge \\ \quad a < b \wedge \\ \quad \forall \mathit{lxb} \in \mathbb{C}. \\ \quad \quad ((a \leq \operatorname{magnitude}((\mathit{lxb} - c))) \wedge (\operatorname{magnitude}((\mathit{lxb} - c)) \leq b)) \Rightarrow (\mathit{lxb} \in u) \Rightarrow \\ \quad \quad \operatorname{line\text{-}int}((\lambda w \in \operatorname{trace}(\operatorname{circle\text{-}path}(c, a), 0, (2 \cdot \pi)).\; (f(w) \cdot \operatorname{power}(\operatorname{recip}((w - c)), (k + 1)))), \\ \quad \quad \quad \operatorname{circle\text{-}path}(c, a), \operatorname{circle\text{-}deriv}(c, a), 0, (2 \cdot \pi)) \\ \quad \quad =\; \operatorname{line\text{-}int}((\lambda w \in \operatorname{trace}(\operatorname{circle\text{-}path}(c, b), 0, (2 \cdot \pi)).\; (f(w) \cdot \operatorname{power}(\operatorname{recip}((w - c)), (k + 1)))), \\ \quad \quad \quad \operatorname{circle\text{-}path}(c, b), \operatorname{circle\text{-}deriv}(c, b), 0, (2 \cdot \pi)) \wedge \\ \quad \quad \operatorname{line\text{-}int}((\lambda w \in \operatorname{trace}(\operatorname{circle\text{-}path}(c, a), 0, (2 \cdot \pi)).\; (f(w) \cdot \operatorname{power}((w - c), k))), \\ \quad \quad \quad \operatorname{circle\text{-}path}(c, a), \operatorname{circle\text{-}deriv}(c, a), 0, (2 \cdot \pi)) \\ \quad \quad =\; \operatorname{line\text{-}int}((\lambda w \in \operatorname{trace}(\operatorname{circle\text{-}path}(c, b), 0, (2 \cdot \pi)).\; (f(w) \cdot \operatorname{power}((w - c), k))), \\ \quad \quad \quad \operatorname{circle\text{-}path}(c, b), \operatorname{circle\text{-}deriv}(c, b), 0, (2 \cdot \pi)) \end{array}}}
\vnbprop{laurent-expansion}{\vnbstmt{\begin{array}{@{}l@{}} \forall u, f.; \forall c \in \mathbb{C}, a, b, o \in \mathbb{R}, z \in \mathbb{C}. \\ \quad \operatorname{is\text{-}open}(\operatorname{nf\text{-}metric\text{-}space}(\operatorname{cc\text{-}normed\text{-}field}), u) \wedge \\ \quad \operatorname{holomorphic\text{-}on}(u, f) \wedge \\ \quad 0 < a \wedge \\ \quad a < o \wedge \\ \quad o < b \wedge \\ \quad \forall \mathit{lxb} \in \mathbb{C}. \\ \quad \quad ((a \leq \operatorname{magnitude}((\mathit{lxb} - c))) \wedge (\operatorname{magnitude}((\mathit{lxb} - c)) \leq b)) \Rightarrow (\mathit{lxb} \in u) \wedge \\ \quad a < \operatorname{magnitude}((z - c)) \wedge \\ \quad \operatorname{magnitude}((z - c)) < b \Rightarrow \\ \quad \quad f(z) = \operatorname{laurent\text{-}sum}(\operatorname{laurent\text{-}coefs}(f, c, o), c, z) \end{array}}}
\vnbprop{laurent-coef-estimate}{\vnbstmt{\begin{array}{@{}l@{}} \forall u, f.; \forall c \in \mathbb{C}, s \in \mathbb{R}. \\ \quad 0 < s \wedge \\ \quad \operatorname{is\text{-}open}(\operatorname{nf\text{-}metric\text{-}space}(\operatorname{cc\text{-}normed\text{-}field}), u) \wedge \\ \quad \operatorname{holomorphic\text{-}on}(u, f) \wedge \\ \quad \operatorname{trace}(\operatorname{circle\text{-}path}(c, s), 0, (2 \cdot \pi)) \subseteq u \Rightarrow \\ \quad \quad \exists m \in \mathbb{R}. \\ \quad \quad \quad 0 \leq m \wedge \\ \quad \quad \quad \forall k \in \mathbb{N}. \\ \quad \quad \quad \quad \operatorname{magnitude}(\operatorname{line\text{-}int}((\lambda w \in \operatorname{trace}(\operatorname{circle\text{-}path}(c, s), 0, (2 \cdot \pi)).\; (f(w) \cdot \operatorname{power}(\operatorname{recip}((w - c)), (k + 1)))), \\ \quad \quad \quad \quad \quad \quad \operatorname{circle\text{-}path}(c, s), \operatorname{circle\text{-}deriv}(c, s), 0, (2 \cdot \pi))) \\ \quad \quad \quad \quad \leq\; ((2 \cdot \pi) \cdot s) \cdot (m \cdot \operatorname{power}(\operatorname{recip}(s), (k + 1))) \wedge \\ \quad \quad \quad \quad \operatorname{magnitude}(\operatorname{line\text{-}int}((\lambda w \in \operatorname{trace}(\operatorname{circle\text{-}path}(c, s), 0, (2 \cdot \pi)).\; (f(w) \cdot \operatorname{power}((w - c), k))), \\ \quad \quad \quad \quad \quad \quad \operatorname{circle\text{-}path}(c, s), \operatorname{circle\text{-}deriv}(c, s), 0, (2 \cdot \pi))) \\ \quad \quad \quad \quad \leq\; ((2 \cdot \pi) \cdot s) \cdot (m \cdot \operatorname{power}(s, k)) \end{array}}}
\vnbprop{laurent-coefs-radii}{\vnbstmt{\begin{array}{@{}l@{}} \forall u, f.; \forall c \in \mathbb{C}, r, o \in \mathbb{R}. \\ \quad 0 < o \wedge \\ \quad o < r \wedge \\ \quad \operatorname{is\text{-}open}(\operatorname{nf\text{-}metric\text{-}space}(\operatorname{cc\text{-}normed\text{-}field}), u) \wedge \\ \quad \operatorname{holomorphic\text{-}on}(u, f) \wedge \\ \quad \forall b \in \mathbb{C}.\; (((0 < \operatorname{magnitude}((b - c))) \wedge (\operatorname{magnitude}((b - c)) < r)) \Rightarrow (b \in u)) \Rightarrow \\ \quad \quad \forall p \in \mathbb{R}. \\ \quad \quad \quad (p < r) \Rightarrow (p < \operatorname{cps\text{-}radius}((\lambda k \in \mathbb{N}.\; \operatorname{laurent\text{-}coefs}(f, c, o)(k)))) \wedge \\ \quad \quad \forall m \in \mathbb{R}. \\ \quad \quad \quad (0 < m) \Rightarrow  \\ \quad \quad \quad \quad \operatorname{recip}(m) < \operatorname{cps\text{-}radius}((\lambda k \in \mathbb{N}.\; \operatorname{laurent\text{-}coefs}(f, c, o)(-(k + 1)))) \end{array}}}
\vnbprop{laurent-coefs-in-fun}{\vnbstmt{\begin{array}{@{}l@{}} \forall u, f.; \forall c \in \mathbb{C}, o \in \mathbb{R}. \\ \quad 0 < o \wedge \\ \quad \operatorname{is\text{-}open}(\operatorname{nf\text{-}metric\text{-}space}(\operatorname{cc\text{-}normed\text{-}field}), u) \wedge \\ \quad \operatorname{holomorphic\text{-}on}(u, f) \wedge \\ \quad \operatorname{trace}(\operatorname{circle\text{-}path}(c, o), 0, (2 \cdot \pi)) \subseteq u \Rightarrow \\ \quad \quad \operatorname{laurent\text{-}coefs}(f, c, o) \in (\mathbb{Z} \to \mathbb{C}) \end{array}}}
\vnbprop{laurent-coef-integrands}{\vnbstmt{\begin{array}{@{}l@{}} \forall u, f.; \forall c \in \mathbb{C}, o \in \mathbb{R}, k \in \mathbb{N}. \\ \quad 0 < o \wedge \\ \quad \operatorname{is\text{-}open}(\operatorname{nf\text{-}metric\text{-}space}(\operatorname{cc\text{-}normed\text{-}field}), u) \wedge \\ \quad \operatorname{holomorphic\text{-}on}(u, f) \wedge \\ \quad \operatorname{trace}(\operatorname{circle\text{-}path}(c, o), 0, (2 \cdot \pi)) \subseteq u \Rightarrow \\ \quad \quad (\lambda w \in \operatorname{trace}(\operatorname{circle\text{-}path}(c, o), 0, (2 \cdot \pi)).\; (f(w) \cdot \operatorname{power}(\operatorname{recip}((w - c)), (k + 1)))) \\ \quad \quad \in\; (\operatorname{trace}(\operatorname{circle\text{-}path}(c, o), 0, (2 \cdot \pi)) \to \mathbb{C}) \wedge \\ \quad \quad \forall y \in \operatorname{trace}(\operatorname{circle\text{-}path}(c, o), 0, (2 \cdot \pi)). \\ \quad \quad \quad \operatorname{is\text{-}continuous\text{-}at}(\operatorname{subspace\text{-}ms}(\operatorname{nf\text{-}metric\text{-}space}(\operatorname{cc\text{-}normed\text{-}field}), \operatorname{trace}(\operatorname{circle\text{-}path}(c, o), 0, (2 \cdot \pi))), \\ \quad \quad \quad \quad \operatorname{nf\text{-}metric\text{-}space}(\operatorname{cc\text{-}normed\text{-}field}), \\ \quad \quad \quad \quad (\lambda w \in \operatorname{trace}(\operatorname{circle\text{-}path}(c, o), 0, (2 \cdot \pi)).\; (f(w) \cdot \operatorname{power}(\operatorname{recip}((w - c)), (k + 1)))), \\ \quad \quad \quad \quad y) \wedge \\ \quad \quad (\lambda w \in \operatorname{trace}(\operatorname{circle\text{-}path}(c, o), 0, (2 \cdot \pi)).\; (f(w) \cdot \operatorname{power}((w - c), k))) \\ \quad \quad \in\; (\operatorname{trace}(\operatorname{circle\text{-}path}(c, o), 0, (2 \cdot \pi)) \to \mathbb{C}) \wedge \\ \quad \quad \forall y \in \operatorname{trace}(\operatorname{circle\text{-}path}(c, o), 0, (2 \cdot \pi)). \\ \quad \quad \quad \operatorname{is\text{-}continuous\text{-}at}(\operatorname{subspace\text{-}ms}(\operatorname{nf\text{-}metric\text{-}space}(\operatorname{cc\text{-}normed\text{-}field}), \operatorname{trace}(\operatorname{circle\text{-}path}(c, o), 0, (2 \cdot \pi))), \\ \quad \quad \quad \quad \operatorname{nf\text{-}metric\text{-}space}(\operatorname{cc\text{-}normed\text{-}field}), \\ \quad \quad \quad \quad (\lambda w \in \operatorname{trace}(\operatorname{circle\text{-}path}(c, o), 0, (2 \cdot \pi)).\; (f(w) \cdot \operatorname{power}((w - c), k))), \\ \quad \quad \quad \quad y) \wedge \\ \quad \quad \operatorname{laurent\text{-}coefs}(f, c, o)(k) \\ \quad \quad =\; \operatorname{recip}(((2 \cdot \pi) \cdot +i)) \cdot \operatorname{line\text{-}int}((\lambda w \in \operatorname{trace}(\operatorname{circle\text{-}path}(c, o), 0, (2 \cdot \pi)).\; (f(w) \cdot \operatorname{power}(\operatorname{recip}((w - c)), (k + 1)))), \operatorname{circle\text{-}path}(c, o), \operatorname{circle\text{-}deriv}(c, o), 0, (2 \cdot \pi)) \wedge \\ \quad \quad \operatorname{laurent\text{-}coefs}(f, c, o)(-(k + 1)) \\ \quad \quad =\; \operatorname{recip}(((2 \cdot \pi) \cdot +i)) \cdot \operatorname{line\text{-}int}((\lambda w \in \operatorname{trace}(\operatorname{circle\text{-}path}(c, o), 0, (2 \cdot \pi)).\; (f(w) \cdot \operatorname{power}((w - c), k))), \operatorname{circle\text{-}path}(c, o), \operatorname{circle\text{-}deriv}(c, o), 0, (2 \cdot \pi)) \end{array}}}
\vnbprop{laurent-coef-sup-bound}{\vnbstmt{\begin{array}{@{}l@{}} \forall u, f.; \forall c \in \mathbb{C}, o, m \in \mathbb{R}, k \in \mathbb{N}. \\ \quad 0 < o \wedge \\ \quad \operatorname{is\text{-}open}(\operatorname{nf\text{-}metric\text{-}space}(\operatorname{cc\text{-}normed\text{-}field}), u) \wedge \\ \quad \operatorname{holomorphic\text{-}on}(u, f) \wedge \\ \quad \operatorname{trace}(\operatorname{circle\text{-}path}(c, o), 0, (2 \cdot \pi)) \subseteq u \wedge \\ \quad 0 \leq m \wedge \\ \quad \forall y \in \operatorname{trace}(\operatorname{circle\text{-}path}(c, o), 0, (2 \cdot \pi)).\; (\operatorname{magnitude}(f(y)) \leq m) \Rightarrow \\ \quad \quad \operatorname{magnitude}(\operatorname{laurent\text{-}coefs}(f, c, o)(k)) \leq ((o \cdot \operatorname{power}(\operatorname{recip}(o), (k + 1))) \cdot m) \wedge \\ \quad \quad \operatorname{magnitude}(\operatorname{laurent\text{-}coefs}(f, c, o)(-(k + 1))) \leq ((o \cdot \operatorname{power}(o, k)) \cdot m) \end{array}}}
\vnbprop{laurent-coef-add}{\vnbstmt{\begin{array}{@{}l@{}} \forall u, f, g, h.; \forall c \in \mathbb{C}, o \in \mathbb{R}, k \in \mathbb{N}. \\ \quad 0 < o \wedge \\ \quad \operatorname{is\text{-}open}(\operatorname{nf\text{-}metric\text{-}space}(\operatorname{cc\text{-}normed\text{-}field}), u) \wedge \\ \quad \operatorname{holomorphic\text{-}on}(u, f) \wedge \\ \quad \operatorname{holomorphic\text{-}on}(u, g) \wedge \\ \quad \operatorname{holomorphic\text{-}on}(u, h) \wedge \\ \quad \operatorname{trace}(\operatorname{circle\text{-}path}(c, o), 0, (2 \cdot \pi)) \subseteq u \wedge \\ \quad \forall z \in u.\; (h(z) \simeq (f(z) + g(z))) \Rightarrow \\ \quad \quad \operatorname{laurent\text{-}coefs}(h, c, o)(k) = (\operatorname{laurent\text{-}coefs}(f, c, o)(k) + \operatorname{laurent\text{-}coefs}(g, c, o)(k)) \wedge \\ \quad \quad \operatorname{laurent\text{-}coefs}(h, c, o)(-(k + 1)) \\ \quad \quad =\; \operatorname{laurent\text{-}coefs}(f, c, o)(-(k + 1)) + \operatorname{laurent\text{-}coefs}(g, c, o)(-(k + 1)) \end{array}}}
\vnbprop{laurent-coef-scale}{\vnbstmt{\begin{array}{@{}l@{}} \forall s \in \mathbb{C}, u, f, h.; \forall c \in \mathbb{C}, o \in \mathbb{R}, k \in \mathbb{N}. \\ \quad 0 < o \wedge \\ \quad \operatorname{is\text{-}open}(\operatorname{nf\text{-}metric\text{-}space}(\operatorname{cc\text{-}normed\text{-}field}), u) \wedge \\ \quad \operatorname{holomorphic\text{-}on}(u, f) \wedge \\ \quad \operatorname{holomorphic\text{-}on}(u, h) \wedge \\ \quad \operatorname{trace}(\operatorname{circle\text{-}path}(c, o), 0, (2 \cdot \pi)) \subseteq u \wedge \\ \quad \forall z \in u.\; (h(z) \simeq (s \cdot f(z))) \Rightarrow \\ \quad \quad \operatorname{laurent\text{-}coefs}(h, c, o)(k) = (s \cdot \operatorname{laurent\text{-}coefs}(f, c, o)(k)) \wedge \\ \quad \quad \operatorname{laurent\text{-}coefs}(h, c, o)(-(k + 1)) = (s \cdot \operatorname{laurent\text{-}coefs}(f, c, o)(-(k + 1))) \end{array}}}

\vnbbookfile{theorem-library/ch4a-punctured}{theorem-library/ch4a-punctured}
\noindent{\small ch4a-punctured.scm -- WAVE 2 (2026-10-27..30), wired 2026-10-30 from scratchpad/ch4a/ch4a-punctured.scm. Library: complex-analysis.}\par
\vnbprop{laurent-punctured-expansion}{\vnbstmt{\begin{array}{@{}l@{}} \forall f.; \forall a \in \mathbb{C}, r, o \in \mathbb{R}, z \in \mathbb{C}. \\ \quad 0 < o \wedge \\ \quad o < r \wedge \\ \quad \operatorname{holomorphic\text{-}on}(\operatorname{annulus}(a, 0, r), f) \wedge \\ \quad 0 < \operatorname{magnitude}((z - a)) \wedge \\ \quad \operatorname{magnitude}((z - a)) < r \Rightarrow \\ \quad \quad f(z) = \operatorname{laurent\text{-}sum}(\operatorname{laurent\text{-}coefs}(f, a, o), a, z) \end{array}}}
\vnbprop{laurent-punctured-holomorphic}{\vnbstmt{\begin{array}{@{}l@{}} \forall c \in (\mathbb{Z} \to \mathbb{C}), a \in \mathbb{C}, r \in \mathbb{R}, f \in (\operatorname{annulus}(a, 0, r) \to \mathbb{C}). \\ \quad 0 < r \wedge \\ \quad \forall x \in \operatorname{annulus}(a, 0, r).\; (f(x) = \operatorname{laurent\text{-}sum}(c, a, x)) \wedge \\ \quad \forall p \in \mathbb{R}.\; ((p < r) \Rightarrow (p < \operatorname{cps\text{-}radius}((\lambda k \in \mathbb{N}.\; c(k))))) \wedge \\ \quad \forall m \in \mathbb{R}. \\ \quad \quad (0 < m) \Rightarrow (\operatorname{recip}(m) < \operatorname{cps\text{-}radius}((\lambda k \in \mathbb{N}.\; c(-(k + 1))))) \Rightarrow \\ \quad \quad \operatorname{holomorphic\text{-}on}(\operatorname{annulus}(a, 0, r), f) \end{array}}}
\vnbprop{laurent-punctured-iff}{\vnbstmt{\begin{array}{@{}l@{}} \forall f.; \forall a \in \mathbb{C}, r \in \mathbb{R}. \\ \quad ((0 < r) \wedge (f \in (\operatorname{annulus}(a, 0, r) \to \mathbb{C}))) \Rightarrow  \\ \quad \quad \operatorname{holomorphic\text{-}on}(\operatorname{annulus}(a, 0, r), f) \\ \quad \quad \Leftrightarrow\; \exists c \in (\mathbb{Z} \to \mathbb{C}). \\ \quad \quad \quad \forall p \in \mathbb{R}.\; ((p < r) \Rightarrow (p < \operatorname{cps\text{-}radius}((\lambda k \in \mathbb{N}.\; c(k))))) \wedge \\ \quad \quad \quad \forall m \in \mathbb{R}. \\ \quad \quad \quad \quad (0 < m) \Rightarrow (\operatorname{recip}(m) < \operatorname{cps\text{-}radius}((\lambda k \in \mathbb{N}.\; c(-(k + 1))))) \wedge \\ \quad \quad \quad \forall x \in \operatorname{annulus}(a, 0, r).\; (f(x) = \operatorname{laurent\text{-}sum}(c, a, x)) \end{array}}}

\vnbbookfile{theorem-library/ch4a-coefcont}{theorem-library/ch4a-coefcont}
\noindent{\small ch4a-coefcont.scm -- WAVE 2 (2026-10-27..30), wired 2026-10-30 from scratchpad/ch4a/ch4a-coefcont.scm. Library: complex-analysis.}\par
\vnbprop{lk-cc-seminorm-family}{\vnbstmt{\begin{array}{@{}l@{}} \operatorname{is\text{-}seminorm\text{-}family}(\operatorname{cc\text{-}normed\text{-}field}, \\ \quad \operatorname{ring\text{-}as\text{-}module}(\operatorname{normed\text{-}field\text{-}as\text{-}commutative\text{-}ring}(\operatorname{cc\text{-}normed\text{-}field})), \\ \quad (\lambda k \in \mathbb{N}.\; (\lambda z \in \mathbb{C}.\; \operatorname{magnitude}(z)))) \end{array}}}
\vnbprop{laurent-coef-cont-lin}{\vnbstmt{\begin{array}{@{}l@{}} \forall u.; \forall a \in \mathbb{C}, o \in \mathbb{R}, q \in \mathbb{Z}. \\ \quad \operatorname{is\text{-}open}(\operatorname{nf\text{-}metric\text{-}space}(\operatorname{cc\text{-}normed\text{-}field}), u) \wedge \\ \quad 0 < o \wedge \\ \quad \operatorname{trace}(\operatorname{circle\text{-}path}(a, o), 0, (2 \cdot \pi)) \subseteq u \Rightarrow \\ \quad \quad \operatorname{is\text{-}cont\text{-}lin}(\operatorname{cc\text{-}fn\text{-}module}(u, \operatorname{cc\text{-}hol\text{-}fns}(u)), \\ \quad \quad \quad \operatorname{cc\text{-}exhaust\text{-}seminorms}(u, \operatorname{cc\text{-}hol\text{-}fns}(u)), \\ \quad \quad \quad \operatorname{ring\text{-}as\text{-}module}(\operatorname{normed\text{-}field\text{-}as\text{-}commutative\text{-}ring}(\operatorname{cc\text{-}normed\text{-}field})), \\ \quad \quad \quad (\lambda k \in \mathbb{N}.\; (\lambda z \in \mathbb{C}.\; \operatorname{magnitude}(z))), \\ \quad \quad \quad (\lambda f \in \operatorname{cc\text{-}hol\text{-}fns}(u).\; \operatorname{laurent\text{-}coefs}(f, a, o)(q))) \end{array}}}

\vnbbookfile{theorem-library/ch4b-residues}{theorem-library/ch4b-residues}
\noindent{\small ch4b-residues.scm -- WAVE 2 (2026-10-27..30), wired 2026-10-30 from scratchpad/ch4b/ch4b-residues.scm. Library: complex-analysis.}\par
\vnbprop{md-recip-power-holomorphic}{\vnbstmt{\begin{array}{@{}l@{}} \forall a \in \mathbb{C}, j \in \mathbb{N}. \\ \quad \operatorname{holomorphic\text{-}on}(\operatorname{difference}(\mathbb{C}, \operatorname{singleton}(a)), \\ \quad \quad (\lambda v \in \operatorname{difference}(\mathbb{C}, \operatorname{singleton}(a)).\; \operatorname{power}(\operatorname{recip}((v - a)), (j + 1)))) \end{array}}}
\vnbprop{md-recip-power-road-zero}{\vnbstmt{\begin{array}{@{}l@{}} \forall a \in \mathbb{C}, j \in \mathbb{N}, p, q, s, t. \\ \quad (\operatorname{is\text{-}road}(p, q, s, t) \wedge ((\operatorname{trace}(p, s, t) \subseteq \operatorname{difference}(\mathbb{C}, \operatorname{singleton}(a))) \wedge (p(t) = p(s)))) \Rightarrow  \\ \quad \quad \operatorname{line\text{-}int}((\lambda v \in \operatorname{difference}(\mathbb{C}, \operatorname{singleton}(a)).\; \operatorname{power}(\operatorname{recip}((v - a)), ((j + 1) + 1))), p, q, s, t) = 0 \end{array}}}
\vnbprop{md-zero-exact}{\vnbstmt{\begin{array}{@{}l@{}} \forall a \in \mathbb{C}. \\ \quad \operatorname{holomorphic\text{-}on}(\operatorname{difference}(\mathbb{C}, \operatorname{singleton}(a)), (\lambda v \in \operatorname{difference}(\mathbb{C}, \operatorname{singleton}(a)).\; 0)) \wedge \\ \quad \forall p, q, s, t. \\ \quad \quad (\operatorname{is\text{-}road}(p, q, s, t) \wedge ((\operatorname{trace}(p, s, t) \subseteq \operatorname{difference}(\mathbb{C}, \operatorname{singleton}(a))) \wedge (p(t) = p(s)))) \Rightarrow  \\ \quad \quad \quad \operatorname{line\text{-}int}((\lambda v \in \operatorname{difference}(\mathbb{C}, \operatorname{singleton}(a)).\; 0), p, q, s, t) = 0 \end{array}}}
\vnbprop{meromorphic-local-decomposition}{\vnbstmt{\begin{array}{@{}l@{}} \forall n \in \mathbb{N}, a \in \mathbb{C}, r \in \mathbb{R}, g. \\ \quad ((0 < r) \wedge \operatorname{holomorphic\text{-}on}(\operatorname{ball}(\operatorname{cc\text{-}ms}, a, r), g)) \Rightarrow  \\ \quad \quad \exists d_{0} \in \mathbb{C}, h. \\ \quad \quad \quad \operatorname{holomorphic\text{-}on}(\operatorname{ball}(\operatorname{cc\text{-}ms}, a, r), h) \wedge \\ \quad \quad \quad \exists f. \\ \quad \quad \quad \quad \operatorname{holomorphic\text{-}on}(\operatorname{difference}(\mathbb{C}, \operatorname{singleton}(a)), f) \wedge \\ \quad \quad \quad \quad \forall p, q, s, t. \\ \quad \quad \quad \quad \quad (\operatorname{is\text{-}road}(p, q, s, t) \wedge ((\operatorname{trace}(p, s, t) \subseteq \operatorname{difference}(\mathbb{C}, \operatorname{singleton}(a))) \wedge (p(t) = p(s)))) \Rightarrow  \\ \quad \quad \quad \quad \quad \quad \operatorname{line\text{-}int}(f, p, q, s, t) = 0 \wedge \\ \quad \quad \quad \quad \forall z \in \operatorname{difference}(\operatorname{ball}(\operatorname{cc\text{-}ms}, a, r), \operatorname{singleton}(a)). \\ \quad \quad \quad \quad \quad (\operatorname{power}(\operatorname{recip}((z - a)), n) \cdot g(z)) = (h(z) + ((d_{0} \cdot \operatorname{recip}((z - a))) + f(z))) \end{array}}}
\vnbprop{residue-circle-integral}{\vnbstmt{\begin{array}{@{}l@{}} \forall a, r, d, h, \mathit{mdf}, f.; \forall o \in \mathbb{R}. \\ \quad a \in \mathbb{C} \wedge \\ \quad r \in \mathbb{R} \wedge \\ \quad 0 < r \wedge \\ \quad d \in \mathbb{C} \wedge \\ \quad \operatorname{holomorphic\text{-}on}(\operatorname{ball}(\operatorname{cc\text{-}ms}, a, r), h) \wedge \\ \quad \operatorname{holomorphic\text{-}on}(\operatorname{difference}(\mathbb{C}, \operatorname{singleton}(a)), \mathit{mdf}) \wedge \\ \quad \forall p, q, s, t. \\ \quad \quad (\operatorname{is\text{-}road}(p, q, s, t) \wedge ((\operatorname{trace}(p, s, t) \subseteq \operatorname{difference}(\mathbb{C}, \operatorname{singleton}(a))) \wedge (p(t) = p(s)))) \Rightarrow  \\ \quad \quad \quad \operatorname{line\text{-}int}(\mathit{mdf}, p, q, s, t) = 0 \wedge \\ \quad \forall z \in \operatorname{difference}(\operatorname{ball}(\operatorname{cc\text{-}ms}, a, r), \operatorname{singleton}(a)). \\ \quad \quad f(z) = (h(z) + ((d \cdot \operatorname{recip}((z - a))) + \operatorname{mdf}(z))) \wedge \\ \quad 0 < o \wedge \\ \quad o < r \Rightarrow \\ \quad \quad \operatorname{recip}(((2 \cdot \pi) \cdot +i)) \cdot \operatorname{line\text{-}int}((\lambda w \in \operatorname{trace}(\operatorname{circle\text{-}path}(a, o), 0, (2 \cdot \pi)).\; f(w)), \operatorname{circle\text{-}path}(a, o), \operatorname{circle\text{-}deriv}(a, o), 0, (2 \cdot \pi)) \\ \quad \quad =\; d \end{array}}}
\vnbprop{residue-property-unique}{\vnbstmt{\begin{array}{@{}l@{}} \forall f, a, x, y. \\ \quad x \in \mathbb{C} \wedge \\ \quad \exists r \in \mathbb{R}. \\ \quad \quad 0 < r \wedge \\ \quad \quad \forall p \in \mathbb{R}. \\ \quad \quad \quad ((0 < p) \wedge (p < r)) \Rightarrow  \\ \quad \quad \quad \quad x \\ \quad \quad \quad \quad =\; \operatorname{recip}(((2 \cdot \pi) \cdot +i)) \cdot \operatorname{line\text{-}int}((\lambda w \in \operatorname{trace}(\operatorname{circle\text{-}path}(a, p), 0, (2 \cdot \pi)).\; f(w)), \operatorname{circle\text{-}path}(a, p), \operatorname{circle\text{-}deriv}(a, p), 0, (2 \cdot \pi)) \wedge \\ \quad y \in \mathbb{C} \wedge \\ \quad \exists r \in \mathbb{R}. \\ \quad \quad 0 < r \wedge \\ \quad \quad \forall p \in \mathbb{R}. \\ \quad \quad \quad ((0 < p) \wedge (p < r)) \Rightarrow  \\ \quad \quad \quad \quad y \\ \quad \quad \quad \quad =\; \operatorname{recip}(((2 \cdot \pi) \cdot +i)) \cdot \operatorname{line\text{-}int}((\lambda w \in \operatorname{trace}(\operatorname{circle\text{-}path}(a, p), 0, (2 \cdot \pi)).\; f(w)), \operatorname{circle\text{-}path}(a, p), \operatorname{circle\text{-}deriv}(a, p), 0, (2 \cdot \pi)) \Rightarrow \\ \quad \quad x = y \end{array}}}
\vnbprop{residue-of-property}{\vnbstmt{\begin{array}{@{}l@{}} \forall f, a.; \forall x \in \mathbb{C}. \\ \quad \exists r \in \mathbb{R}. \\ \quad \quad 0 < r \wedge \\ \quad \quad \forall p \in \mathbb{R}. \\ \quad \quad \quad ((0 < p) \wedge (p < r)) \Rightarrow  \\ \quad \quad \quad \quad x \\ \quad \quad \quad \quad =\; \operatorname{recip}(((2 \cdot \pi) \cdot +i)) \cdot \operatorname{line\text{-}int}((\lambda w \in \operatorname{trace}(\operatorname{circle\text{-}path}(a, p), 0, (2 \cdot \pi)).\; f(w)), \operatorname{circle\text{-}path}(a, p), \operatorname{circle\text{-}deriv}(a, p), 0, (2 \cdot \pi)) \Rightarrow \\ \quad \quad \operatorname{residue}(f, a) = x \end{array}}}
\vnbprop{meromorphic-residue-decomposition}{\vnbstmt{\begin{array}{@{}l@{}} \forall f, a. \\ \quad \operatorname{is\text{-}meromorphic\text{-}at}(f, a) \Rightarrow  \\ \quad \quad \operatorname{residue}(f, a) \in \mathbb{C} \wedge \\ \quad \quad \exists r \in \mathbb{R}. \\ \quad \quad \quad 0 < r \wedge \\ \quad \quad \quad \exists h. \\ \quad \quad \quad \quad \operatorname{holomorphic\text{-}on}(\operatorname{ball}(\operatorname{cc\text{-}ms}, a, r), h) \wedge \\ \quad \quad \quad \quad \exists \mathit{mdf}. \\ \quad \quad \quad \quad \quad \operatorname{holomorphic\text{-}on}(\operatorname{difference}(\mathbb{C}, \operatorname{singleton}(a)), \mathit{mdf}) \wedge \\ \quad \quad \quad \quad \quad \forall p, q, s, t. \\ \quad \quad \quad \quad \quad \quad (\operatorname{is\text{-}road}(p, q, s, t) \wedge ((\operatorname{trace}(p, s, t) \subseteq \operatorname{difference}(\mathbb{C}, \operatorname{singleton}(a))) \wedge (p(t) = p(s)))) \Rightarrow  \\ \quad \quad \quad \quad \quad \quad \quad \operatorname{line\text{-}int}(\mathit{mdf}, p, q, s, t) = 0 \wedge \\ \quad \quad \quad \quad \quad \forall z \in \operatorname{difference}(\operatorname{ball}(\operatorname{cc\text{-}ms}, a, r), \operatorname{singleton}(a)). \\ \quad \quad \quad \quad \quad \quad f(z) = (h(z) + ((\operatorname{residue}(f, a) \cdot \operatorname{recip}((z - a))) + \operatorname{mdf}(z))) \end{array}}}
\vnbprop{fun-value-in-domain}{\vnbstmt{\forall a, b, f, z, t.\; (((f \in (a \to b)) \wedge (f(z) = t)) \Rightarrow (z \in a))}}
\vnbprop{residue-of-decomposition}{\vnbstmt{\begin{array}{@{}l@{}} \forall a, r, d, h, \mathit{mdf}, f. \\ \quad a \in \mathbb{C} \wedge \\ \quad r \in \mathbb{R} \wedge \\ \quad 0 < r \wedge \\ \quad d \in \mathbb{C} \wedge \\ \quad \operatorname{holomorphic\text{-}on}(\operatorname{ball}(\operatorname{cc\text{-}ms}, a, r), h) \wedge \\ \quad \operatorname{holomorphic\text{-}on}(\operatorname{difference}(\mathbb{C}, \operatorname{singleton}(a)), \mathit{mdf}) \wedge \\ \quad \forall p, q, s, t. \\ \quad \quad (\operatorname{is\text{-}road}(p, q, s, t) \wedge ((\operatorname{trace}(p, s, t) \subseteq \operatorname{difference}(\mathbb{C}, \operatorname{singleton}(a))) \wedge (p(t) = p(s)))) \Rightarrow  \\ \quad \quad \quad \operatorname{line\text{-}int}(\mathit{mdf}, p, q, s, t) = 0 \wedge \\ \quad \forall z \in \operatorname{difference}(\operatorname{ball}(\operatorname{cc\text{-}ms}, a, r), \operatorname{singleton}(a)). \\ \quad \quad f(z) = (h(z) + ((d \cdot \operatorname{recip}((z - a))) + \operatorname{mdf}(z))) \Rightarrow \\ \quad \quad \operatorname{residue}(f, a) = d \end{array}}}
\vnbprop{residue-theorem-step}{\vnbstmt{\begin{array}{@{}l@{}} \forall u, s.; \forall b \in u, f. \\ \quad \operatorname{is\text{-}open}(\operatorname{nf\text{-}metric\text{-}space}(\operatorname{cc\text{-}normed\text{-}field}), u) \wedge \\ \quad s \subseteq u \wedge \\ \quad \neg (b \in s) \wedge \\ \quad \operatorname{holomorphic\text{-}on}(\operatorname{difference}(u, (s \cup \{b\})), f) \wedge \\ \quad \exists r \in \mathbb{R}. \\ \quad \quad 0 < r \wedge \\ \quad \quad \exists d \in \mathbb{C}, h. \\ \quad \quad \quad \operatorname{holomorphic\text{-}on}(\operatorname{ball}(\operatorname{cc\text{-}ms}, b, r), h) \wedge \\ \quad \quad \quad \exists \mathit{mqf}. \\ \quad \quad \quad \quad \operatorname{holomorphic\text{-}on}(\operatorname{difference}(\mathbb{C}, \operatorname{singleton}(b)), \mathit{mqf}) \wedge \\ \quad \quad \quad \quad \forall p, q, \mathit{mds}, t. \\ \quad \quad \quad \quad \quad (\operatorname{is\text{-}road}(p, q, \mathit{mds}, t) \wedge ((\operatorname{trace}(p, \mathit{mds}, t) \subseteq \operatorname{difference}(\mathbb{C}, \operatorname{singleton}(b))) \wedge (p(t) = p(\mathit{mds})))) \Rightarrow  \\ \quad \quad \quad \quad \quad \quad \operatorname{line\text{-}int}(\mathit{mqf}, p, q, \mathit{mds}, t) = 0 \wedge \\ \quad \quad \quad \quad \forall z \in \operatorname{difference}(\operatorname{ball}(\operatorname{cc\text{-}ms}, b, r), \operatorname{singleton}(b)). \\ \quad \quad \quad \quad \quad f(z) = (h(z) + ((d \cdot \operatorname{recip}((z - b))) + \operatorname{mqf}(z))) \wedge \\ \quad \forall e \in s. \exists r \in \mathbb{R}. \\ \quad \quad 0 < r \wedge \\ \quad \quad \exists d \in \mathbb{C}, h. \\ \quad \quad \quad \operatorname{holomorphic\text{-}on}(\operatorname{ball}(\operatorname{cc\text{-}ms}, e, r), h) \wedge \\ \quad \quad \quad \exists \mathit{mqf}. \\ \quad \quad \quad \quad \operatorname{holomorphic\text{-}on}(\operatorname{difference}(\mathbb{C}, \operatorname{singleton}(e)), \mathit{mqf}) \wedge \\ \quad \quad \quad \quad \forall p, q, \mathit{mds}, t. \\ \quad \quad \quad \quad \quad (\operatorname{is\text{-}road}(p, q, \mathit{mds}, t) \wedge ((\operatorname{trace}(p, \mathit{mds}, t) \subseteq \operatorname{difference}(\mathbb{C}, \operatorname{singleton}(e))) \wedge (p(t) = p(\mathit{mds})))) \Rightarrow  \\ \quad \quad \quad \quad \quad \quad \operatorname{line\text{-}int}(\mathit{mqf}, p, q, \mathit{mds}, t) = 0 \wedge \\ \quad \quad \quad \quad \forall z \in \operatorname{difference}(\operatorname{ball}(\operatorname{cc\text{-}ms}, e, r), \operatorname{singleton}(e)). \\ \quad \quad \quad \quad \quad f(z) = (h(z) + ((d \cdot \operatorname{recip}((z - e))) + \operatorname{mqf}(z))) \Rightarrow \\ \quad \quad \exists g. \\ \quad \quad \quad \operatorname{holomorphic\text{-}on}(\operatorname{difference}(u, s), g) \wedge \\ \quad \quad \quad \forall e \in s. \\ \quad \quad \quad \quad \exists r \in \mathbb{R}. \\ \quad \quad \quad \quad \quad 0 < r \wedge \\ \quad \quad \quad \quad \quad \exists d \in \mathbb{C}, h. \\ \quad \quad \quad \quad \quad \quad \operatorname{holomorphic\text{-}on}(\operatorname{ball}(\operatorname{cc\text{-}ms}, e, r), h) \wedge \\ \quad \quad \quad \quad \quad \quad \exists \mathit{mqf}. \\ \quad \quad \quad \quad \quad \quad \quad \operatorname{holomorphic\text{-}on}(\operatorname{difference}(\mathbb{C}, \operatorname{singleton}(e)), \mathit{mqf}) \wedge \\ \quad \quad \quad \quad \quad \quad \quad \forall p, q, \mathit{mds}, t. \\ \quad \quad \quad \quad \quad \quad \quad \quad (\operatorname{is\text{-}road}(p, q, \mathit{mds}, t) \wedge ((\operatorname{trace}(p, \mathit{mds}, t) \subseteq \operatorname{difference}(\mathbb{C}, \operatorname{singleton}(e))) \wedge (p(t) = p(\mathit{mds})))) \Rightarrow  \\ \quad \quad \quad \quad \quad \quad \quad \quad \quad \operatorname{line\text{-}int}(\mathit{mqf}, p, q, \mathit{mds}, t) = 0 \wedge \\ \quad \quad \quad \quad \quad \quad \quad \forall z \in \operatorname{difference}(\operatorname{ball}(\operatorname{cc\text{-}ms}, e, r), \operatorname{singleton}(e)). \\ \quad \quad \quad \quad \quad \quad \quad \quad g(z) = (h(z) + ((d \cdot \operatorname{recip}((z - e))) + \operatorname{mqf}(z))) \wedge \\ \quad \quad \quad \quad \operatorname{residue}(g, e) = \operatorname{residue}(f, e) \wedge \\ \quad \quad \quad \forall p, q, \mathit{mds}, t. \\ \quad \quad \quad \quad \operatorname{is\text{-}road}(p, q, \mathit{mds}, t) \wedge \\ \quad \quad \quad \quad \operatorname{trace}(p, \mathit{mds}, t) \subseteq \operatorname{difference}(u, (s \cup \{b\})) \wedge \\ \quad \quad \quad \quad p(t) = p(\mathit{mds}) \Rightarrow \\ \quad \quad \quad \quad \quad \operatorname{line\text{-}int}(f, p, q, \mathit{mds}, t) \\ \quad \quad \quad \quad \quad =\; \operatorname{line\text{-}int}(g, p, q, \mathit{mds}, t) + ((((2 \cdot \pi) \cdot +i) \cdot \operatorname{residue}(f, b)) \cdot \operatorname{winding}(p, q, \mathit{mds}, t, b)) \end{array}}}
\vnbprop{rt-cc-additive-ag-carr}{\vnbstmt{(\operatorname{carr}(\operatorname{normed\text{-}field\text{-}additive\text{-}ag}(\operatorname{cc\text{-}normed\text{-}field})) \simeq \mathbb{C})}}
\vnbprop{residue-theorem-induction}{\vnbstmt{\begin{array}{@{}l@{}} \forall u.; \forall n \in \mathbb{N}, e \in \mathbf{Set}, f, p, q, s, t. \\ \quad \operatorname{is\text{-}open}(\operatorname{nf\text{-}metric\text{-}space}(\operatorname{cc\text{-}normed\text{-}field}), u) \wedge \\ \quad |e| = n \wedge \\ \quad e \subseteq u \wedge \\ \quad \operatorname{holomorphic\text{-}on}(\operatorname{difference}(u, e), f) \wedge \\ \quad \forall \mathit{mqe} \in e. \exists r \in \mathbb{R}. \\ \quad \quad 0 < r \wedge \\ \quad \quad \exists d \in \mathbb{C}, h. \\ \quad \quad \quad \operatorname{holomorphic\text{-}on}(\operatorname{ball}(\operatorname{cc\text{-}ms}, \mathit{mqe}, r), h) \wedge \\ \quad \quad \quad \exists \mathit{mqf}. \\ \quad \quad \quad \quad \operatorname{holomorphic\text{-}on}(\operatorname{difference}(\mathbb{C}, \operatorname{singleton}(\mathit{mqe})), \mathit{mqf}) \wedge \\ \quad \quad \quad \quad \forall p, q, s, t. \\ \quad \quad \quad \quad \quad (\operatorname{is\text{-}road}(p, q, s, t) \wedge ((\operatorname{trace}(p, s, t) \subseteq \operatorname{difference}(\mathbb{C}, \operatorname{singleton}(\mathit{mqe}))) \wedge (p(t) = p(s)))) \Rightarrow  \\ \quad \quad \quad \quad \quad \quad \operatorname{line\text{-}int}(\mathit{mqf}, p, q, s, t) = 0 \wedge \\ \quad \quad \quad \quad \forall z \in \operatorname{difference}(\operatorname{ball}(\operatorname{cc\text{-}ms}, \mathit{mqe}, r), \operatorname{singleton}(\mathit{mqe})). \\ \quad \quad \quad \quad \quad f(z) = (h(z) + ((d \cdot \operatorname{recip}((z - \mathit{mqe}))) + \operatorname{mqf}(z))) \wedge \\ \quad \operatorname{is\text{-}road}(p, q, s, t) \wedge \\ \quad \operatorname{trace}(p, s, t) \subseteq \operatorname{difference}(u, e) \wedge \\ \quad p(t) = p(s) \wedge \\ \quad \operatorname{null\text{-}homotopic}(u, p, s, t) \Rightarrow \\ \quad \quad \operatorname{recip}(((2 \cdot \pi) \cdot +i)) \cdot \operatorname{line\text{-}int}(f, p, q, s, t) \\ \quad \quad =\; \operatorname{finsum}(\operatorname{normed\text{-}field\text{-}additive\text{-}ag}(\operatorname{cc\text{-}normed\text{-}field}), \\ \quad \quad \quad (\lambda x \in \mathbb{C}.\; (\operatorname{residue}(f, x) \cdot \operatorname{winding}(p, q, s, t, x))), \\ \quad \quad \quad e) \end{array}}}
\vnbprop{residue-theorem}{\vnbstmt{\begin{array}{@{}l@{}} \forall u.; \forall e \in \mathbf{Set}, f, p, q, s, t. \\ \quad \operatorname{is\text{-}open}(\operatorname{nf\text{-}metric\text{-}space}(\operatorname{cc\text{-}normed\text{-}field}), u) \wedge \\ \quad |e| \in \mathbb{N} \wedge \\ \quad e \subseteq u \wedge \\ \quad \operatorname{holomorphic\text{-}on}(\operatorname{difference}(u, e), f) \wedge \\ \quad \forall \mathit{mqe} \in e.\; \operatorname{is\text{-}meromorphic\text{-}at}(f, \mathit{mqe}) \wedge \\ \quad \operatorname{is\text{-}road}(p, q, s, t) \wedge \\ \quad \operatorname{trace}(p, s, t) \subseteq \operatorname{difference}(u, e) \wedge \\ \quad p(t) = p(s) \wedge \\ \quad \operatorname{null\text{-}homotopic}(u, p, s, t) \Rightarrow \\ \quad \quad \operatorname{recip}(((2 \cdot \pi) \cdot +i)) \cdot \operatorname{line\text{-}int}(f, p, q, s, t) \\ \quad \quad =\; \operatorname{finsum}(\operatorname{normed\text{-}field\text{-}additive\text{-}ag}(\operatorname{cc\text{-}normed\text{-}field}), \\ \quad \quad \quad (\lambda x \in \mathbb{C}.\; (\operatorname{residue}(f, x) \cdot \operatorname{winding}(p, q, s, t, x))), \\ \quad \quad \quad e) \end{array}}}

\vnbbookfile{theorem-library/ch4b-mero}{theorem-library/ch4b-mero}
\noindent{\small ch4b-mero.scm -- WAVE 2 (2026-10-27..30), wired 2026-10-30 from scratchpad/ch4b/ch4b-mero.scm. Library: complex-analysis.}\par
\vnbprop{cc-continuous-at-of-modulus}{\vnbstmt{\begin{array}{@{}l@{}} \forall u.; \forall g \in (u \to \mathbb{C}), y \in u. \\ \quad u \subseteq \mathbb{C} \wedge \\ \quad \forall s. \\ \quad \quad \operatorname{pos\text{-}rr}(s) \Rightarrow  \\ \quad \quad \quad \exists l. \\ \quad \quad \quad \quad (\operatorname{pos\text{-}rr}(l) \wedge \forall b \in u.\; ((\operatorname{magnitude}((b - y)) < l) \Rightarrow (\operatorname{magnitude}((g(b) - g(y))) \leq s))) \Rightarrow \\ \quad \quad \operatorname{is\text{-}continuous\text{-}at}(\operatorname{subspace\text{-}ms}(\operatorname{nf\text{-}metric\text{-}space}(\operatorname{cc\text{-}normed\text{-}field}), u), \\ \quad \quad \quad \operatorname{nf\text{-}metric\text{-}space}(\operatorname{cc\text{-}normed\text{-}field}), g, y) \end{array}}}
\vnbprop{riemann-removable-singularity}{\vnbstmt{\begin{array}{@{}l@{}} \forall u.; \forall a \in u, f. \\ \quad (\operatorname{is\text{-}open}(\operatorname{nf\text{-}metric\text{-}space}(\operatorname{cc\text{-}normed\text{-}field}), u) \wedge \operatorname{holomorphic\text{-}on}(\operatorname{difference}(u, \operatorname{singleton}(a)), f)) \Rightarrow  \\ \quad \quad \operatorname{is\text{-}removable\text{-}singularity}(f, a) \\ \quad \quad \Leftrightarrow\; \exists r \in \mathbb{R}. \\ \quad \quad \quad 0 < r \wedge \\ \quad \quad \quad \exists m \in \mathbb{R}. \forall z \in \operatorname{difference}(u, \operatorname{singleton}(a)). \\ \quad \quad \quad \quad (\operatorname{magnitude}((z - a)) < r) \Rightarrow (\operatorname{magnitude}(f(z)) \leq m) \end{array}}}
\vnbprop{pole-is-meromorphic-at}{\vnbstmt{\forall f, a, n.\; (\operatorname{is\text{-}pole}(f, a, n) \Rightarrow \operatorname{is\text{-}meromorphic\text{-}at}(f, a))}}
\vnbprop{removable-is-meromorphic-at}{\vnbstmt{\forall f, a.\; (\operatorname{is\text{-}removable\text{-}singularity}(f, a) \Rightarrow \operatorname{is\text{-}meromorphic\text{-}at}(f, a))}}
\vnbprop{meromorphic-point-cases}{\vnbstmt{\begin{array}{@{}l@{}} \forall f, a. \\ \quad \operatorname{is\text{-}meromorphic\text{-}at}(f, a) \Rightarrow  \\ \quad \quad \operatorname{is\text{-}removable\text{-}singularity}(f, a) \vee \exists n.\; \operatorname{is\text{-}pole}(f, a, n) \end{array}}}
\vnbprop{casorati-weierstrass}{\vnbstmt{\begin{array}{@{}l@{}} \forall u.; \forall a \in u, f.; \forall r \in \mathbb{R}, v \in \mathbb{C}, e \in \mathbb{R}. \\ \quad \operatorname{is\text{-}open}(\operatorname{nf\text{-}metric\text{-}space}(\operatorname{cc\text{-}normed\text{-}field}), u) \wedge \\ \quad \operatorname{holomorphic\text{-}on}(\operatorname{difference}(u, \operatorname{singleton}(a)), f) \wedge \\ \quad \operatorname{is\text{-}essential\text{-}singularity}(f, a) \wedge \\ \quad 0 < r \wedge \\ \quad 0 < e \Rightarrow \\ \quad \quad \exists z \in \operatorname{difference}(u, \operatorname{singleton}(a)). \\ \quad \quad \quad (\operatorname{magnitude}((z - a)) < r) \wedge (\operatorname{magnitude}((f(z) - v)) < e) \end{array}}}
\vnbprop{meromorphic-poles-finite-in-compact}{\vnbstmt{\begin{array}{@{}l@{}} \forall u, s, f, k. \\ \quad \operatorname{is\text{-}open}(\operatorname{nf\text{-}metric\text{-}space}(\operatorname{cc\text{-}normed\text{-}field}), u) \wedge \\ \quad s \subseteq u \wedge \\ \quad \operatorname{holomorphic\text{-}on}(\operatorname{difference}(u, s), f) \wedge \\ \quad \forall e \in s.\; \operatorname{is\text{-}meromorphic\text{-}at}(f, e) \wedge \\ \quad k \subseteq u \wedge \\ \quad \operatorname{is\text{-}compact}(\operatorname{subspace\text{-}ms}(\operatorname{cc\text{-}ms}, k)) \Rightarrow \\ \quad \quad |(s \cap k)| \in \mathbb{N} \end{array}}}

\vnbbookfile{theorem-library/sc-primitive}{theorem-library/sc-primitive}
\noindent{\small ch3d-sc.scm -- CH3-D: consequences of simple connectedness (complex analysis, notes ch. 3). Proposition 3.21: a function holomorphic on a simply connected open set has a primitive. The primitive is F(z) = the integral of f along any road in u from a base point to z: the value does not depend on the road (3.20 on the ...}\par
\vnbprop{sc-path-independent}{\vnbstmt{\begin{array}{@{}l@{}} \forall u, f, m, g, a, b, o, h, c, d. \\ \quad \operatorname{simply\text{-}connected}(u) \wedge \\ \quad \operatorname{holomorphic\text{-}on}(u, f) \wedge \\ \quad \operatorname{is\text{-}road}(m, g, a, b) \wedge \\ \quad \operatorname{trace}(m, a, b) \subseteq u \wedge \\ \quad \operatorname{is\text{-}road}(o, h, c, d) \wedge \\ \quad \operatorname{trace}(o, c, d) \subseteq u \wedge \\ \quad o(c) = m(a) \wedge \\ \quad o(d) = m(b) \Rightarrow \\ \quad \quad \operatorname{line\text{-}int}(f, m, g, a, b) = \operatorname{line\text{-}int}(f, o, h, c, d) \end{array}}}
\vnbprop{segment-primitive-diff}{\vnbstmt{\begin{array}{@{}l@{}} \forall u, f, m.; \forall w \in u. \\ \quad \operatorname{holomorphic\text{-}on}(u, f) \wedge \\ \quad m \in (u \to \mathbb{C}) \wedge \\ \quad \forall z \in u. \exists r. \\ \quad \quad \operatorname{pos\text{-}rr}(r) \wedge \\ \quad \quad \operatorname{ball}(\operatorname{cc\text{-}ms}, z, r) \subseteq u \wedge \\ \quad \quad \forall y \in \operatorname{ball}(\operatorname{cc\text{-}ms}, z, r).\; (m(y) = (m(z) + \operatorname{seg\text{-}int}(f, z, y))) \Rightarrow \\ \quad \quad \operatorname{is\text{-}diff\text{-}on}(\operatorname{cc\text{-}normed\text{-}field}, u, m, w, f(w)) \end{array}}}
\vnbprop{road-extend-segment}{\vnbstmt{\begin{array}{@{}l@{}} \forall u, f, c, z, w, m, d, a, b. \\ \quad \operatorname{holomorphic\text{-}on}(u, f) \wedge \\ \quad \operatorname{is\text{-}road}(m, d, a, b) \wedge \\ \quad \operatorname{trace}(m, a, b) \subseteq u \wedge \\ \quad m(a) = c \wedge \\ \quad m(b) = z \wedge \\ \quad z \in \mathbb{C} \wedge \\ \quad w \in \mathbb{C} \wedge \\ \quad \operatorname{trace}(\operatorname{seg\text{-}path}(z, w), 0, 1) \subseteq u \Rightarrow \\ \quad \quad \exists \mathit{sda}, \mathit{sdb} \in \mathbb{R}, g, \mathit{sdd} \in (\operatorname{ccint}(\mathit{sda}, \mathit{sdb}) \to \mathbb{C}). \\ \quad \quad \quad \operatorname{is\text{-}road}(g, \mathit{sdd}, \mathit{sda}, \mathit{sdb}) \wedge \\ \quad \quad \quad \operatorname{trace}(g, \mathit{sda}, \mathit{sdb}) \subseteq u \wedge \\ \quad \quad \quad g(\mathit{sda}) = c \wedge \\ \quad \quad \quad g(\mathit{sdb}) = w \wedge \\ \quad \quad \quad \operatorname{line\text{-}int}(f, g, \mathit{sdd}, \mathit{sda}, \mathit{sdb}) = (\operatorname{line\text{-}int}(f, m, d, a, b) + \operatorname{seg\text{-}int}(f, z, w)) \end{array}}}
\vnbprop{sc-reach}{\vnbstmt{\begin{array}{@{}l@{}} \forall u, f, a.; \forall q \in u. \\ \quad (\operatorname{holomorphic\text{-}on}(u, f) \wedge (\operatorname{is\text{-}connected}(\operatorname{cc\text{-}ms}, u) \wedge (a \in u))) \Rightarrow  \\ \quad \quad \exists v \in \mathbb{C}, \mathit{sda}, b \in \mathbb{R}, g, d \in (\operatorname{ccint}(\mathit{sda}, b) \to \mathbb{C}). \\ \quad \quad \quad \operatorname{is\text{-}road}(g, d, \mathit{sda}, b) \wedge \\ \quad \quad \quad \operatorname{trace}(g, \mathit{sda}, b) \subseteq u \wedge \\ \quad \quad \quad g(\mathit{sda}) = a \wedge \\ \quad \quad \quad g(b) = q \wedge \\ \quad \quad \quad \operatorname{line\text{-}int}(f, g, d, \mathit{sda}, b) = v \end{array}}}
\vnbprop{sc-local-segment-primitive}{\vnbstmt{\begin{array}{@{}l@{}} \forall u, f, a. \\ \quad (\operatorname{simply\text{-}connected}(u) \wedge (\operatorname{holomorphic\text{-}on}(u, f) \wedge (a \in u))) \Rightarrow  \\ \quad \quad \exists r \in (u \to \mathbb{C}). \forall z \in u. \exists \mathit{sdr}. \\ \quad \quad \quad \operatorname{pos\text{-}rr}(\mathit{sdr}) \wedge \\ \quad \quad \quad \operatorname{ball}(\operatorname{cc\text{-}ms}, z, \mathit{sdr}) \subseteq u \wedge \\ \quad \quad \quad \forall y \in \operatorname{ball}(\operatorname{cc\text{-}ms}, z, \mathit{sdr}).\; (r(y) = (r(z) + \operatorname{seg\text{-}int}(f, z, y))) \end{array}}}
\vnbprop{sc-primitive}{\vnbstmt{\begin{array}{@{}l@{}} \forall u, f.; \forall o \in u. \\ \quad (\operatorname{simply\text{-}connected}(u) \wedge \operatorname{holomorphic\text{-}on}(u, f)) \Rightarrow  \\ \quad \quad \exists g \in (u \to \mathbb{C}). \\ \quad \quad \quad (\operatorname{holomorphic\text{-}on}(u, g) \wedge \forall z \in u.\; \operatorname{is\text{-}diff\text{-}on}(\operatorname{cc\text{-}normed\text{-}field}, u, g, z, f(z))) \end{array}}}

\vnbbookfile{theorem-library/sc-log-roots}{theorem-library/sc-log-roots}
\noindent{\small ch3d-log.scm -- CH3-D: Proposition 3.22 and Corollary 3.25 on a SIMPLY CONNECTED open set. The proofs are those of the convex forms (scratchpad/ch3b/ch3b-log.scm, agent ISOM), with the primitive of f'/f taken from 3.21 (sc-primitive) instead of cauchy-convex-primitive, and the connectedness of u read off ...}\par
\vnbprop{holomorphic-log-sc-exists}{\vnbstmt{\begin{array}{@{}l@{}} \forall u, f. \\ \quad (\operatorname{simply\text{-}connected}(u) \wedge (\operatorname{holomorphic\text{-}on}(u, f) \wedge \forall z \in u.\; \neg (f(z) = 0))) \Rightarrow  \\ \quad \quad \exists g.\; (\operatorname{holomorphic\text{-}on}(u, g) \wedge \forall z \in u.\; (\exp g(z) = f(z))) \end{array}}}
\vnbprop{holomorphic-log-sc}{\vnbstmt{\begin{array}{@{}l@{}} \forall u, f. \\ \quad (\operatorname{simply\text{-}connected}(u) \wedge (\operatorname{holomorphic\text{-}on}(u, f) \wedge \forall z \in u.\; \neg (f(z) = 0))) \Rightarrow  \\ \quad \quad \exists g.\; (\operatorname{holomorphic\text{-}on}(u, g) \wedge \forall z \in u.\; (\exp g(z) = f(z))) \wedge \\ \quad \quad \forall g, h.; \forall a \in u. \\ \quad \quad \quad \operatorname{holomorphic\text{-}on}(u, g) \wedge \\ \quad \quad \quad \operatorname{holomorphic\text{-}on}(u, h) \wedge \\ \quad \quad \quad \forall z \in u.\; (\exp g(z) = f(z)) \wedge \\ \quad \quad \quad \forall z \in u.\; (\exp h(z) = f(z)) \Rightarrow \\ \quad \quad \quad \quad \exists n \in \mathbb{Z}.\; \forall z \in u.\; (g(z) = (h(z) + (n \cdot ((2 \cdot \pi) \cdot +i)))) \end{array}}}
\vnbprop{holomorphic-root-sc-exists}{\vnbstmt{\begin{array}{@{}l@{}} \forall u, f, n. \\ \quad \operatorname{simply\text{-}connected}(u) \wedge \\ \quad \operatorname{holomorphic\text{-}on}(u, f) \wedge \\ \quad \forall z \in u.\; \neg (f(z) = 0) \wedge \\ \quad n \in \mathbb{N} \wedge \\ \quad 1 \leq n \Rightarrow \\ \quad \quad \exists g.\; (\operatorname{holomorphic\text{-}on}(u, g) \wedge \forall z \in u.\; (\operatorname{power}(g(z), n) = f(z))) \end{array}}}
\vnbprop{holomorphic-root-sc-unique}{\vnbstmt{\begin{array}{@{}l@{}} \forall u, f, n, g, h.; \forall b \in u. \\ \quad \operatorname{simply\text{-}connected}(u) \wedge \\ \quad \operatorname{holomorphic\text{-}on}(u, f) \wedge \\ \quad \forall z \in u.\; \neg (f(z) = 0) \wedge \\ \quad n \in \mathbb{N} \wedge \\ \quad 1 \leq n \wedge \\ \quad \operatorname{holomorphic\text{-}on}(u, g) \wedge \\ \quad \operatorname{holomorphic\text{-}on}(u, h) \wedge \\ \quad \forall z \in u.\; (\operatorname{power}(g(z), n) = f(z)) \wedge \\ \quad \forall z \in u.\; (\operatorname{power}(h(z), n) = f(z)) \Rightarrow \\ \quad \quad \exists a \in \mathbb{C}.\; ((\operatorname{power}(a, n) = 1) \wedge \forall z \in u.\; (h(z) = (a \cdot g(z)))) \end{array}}}

\vnbbookfile{theorem-library/riemann-reduction}{theorem-library/riemann-reduction}
\noindent{\small ch3d-rm.scm -- CH3-D: Proposition 10.16 and Corollary 10.17 (the reduction step of the Riemann mapping theorem).  An open simply connected u missing a point a carries a bounded simple (injective) holomorphic function, and a simple holomorphic map into the open unit disc whose image contains 0.  The notes take a ball ...}\par
\vnbprop{sc-bounded-simple}{\vnbstmt{\begin{array}{@{}l@{}} \forall u, a. \\ \quad (\operatorname{simply\text{-}connected}(u) \wedge ((a \in \mathbb{C}) \wedge \neg (a \in u))) \Rightarrow  \\ \quad \quad \exists h. \\ \quad \quad \quad \operatorname{holomorphic\text{-}on}(u, h) \wedge \\ \quad \quad \quad \forall z, w \in u.\; ((h(z) = h(w)) \Rightarrow (z = w)) \wedge \\ \quad \quad \quad \exists m \in \mathbb{R}.\; \forall z \in u.\; (\operatorname{magnitude}(h(z)) \leq m) \end{array}}}
\vnbprop{sc-simple-into-disc}{\vnbstmt{\begin{array}{@{}l@{}} \forall u, a, v. \\ \quad (\operatorname{simply\text{-}connected}(u) \wedge ((a \in \mathbb{C}) \wedge (\neg (a \in u) \wedge (v \in u)))) \Rightarrow  \\ \quad \quad \exists k. \\ \quad \quad \quad \operatorname{holomorphic\text{-}on}(u, k) \wedge \\ \quad \quad \quad \forall z, w \in u.\; ((k(z) = k(w)) \Rightarrow (z = w)) \wedge \\ \quad \quad \quad \forall z \in u.\; (\operatorname{magnitude}(k(z)) < 1) \wedge \\ \quad \quad \quad \exists y \in u.\; (k(y) = 0) \end{array}}}

\vnbbookfile{theorem-library/algnum-countable}{theorem-library/algnum-countable}
\noindent{\small theorem-library/algnum-countable.scm -- THE ALGEBRAIC NUMBERS OVER A COUNTABLE SUBRING OF C ARE COUNTABLE (ALGNUM, algebraic-numbers bundle B7, 2026-10-26).  Source: Thayer, algebraic-numbers notes, Cor. 4.23 (the countability half).  Vocabulary: structure-library/algnum-vocab.scm; POLY-2 (poly-fn-vocab.scm, cf-cons-*, ...}\par
\vnbprop{complex-set-no-injective-tuple-countable}{\vnbstmt{\begin{array}{@{}l@{}} \forall n \in \mathbb{N}, s \in \operatorname{power}(\mathbb{C}). \\ \quad \forall q \in (\mathbb{N} \to s). \\ \quad \quad \neg \forall i, j \in \mathbb{N}.\; (((i \leq n) \wedge ((j \leq n) \wedge (q(i) = q(j)))) \Rightarrow (i = j)) \Rightarrow \\ \quad \quad \operatorname{is\text{-}countable}(s) \end{array}}}
\vnbprop{poly-roots-countable}{\vnbstmt{\begin{array}{@{}l@{}} \forall n \in \mathbb{N}, c \in (\mathbb{N} \to \mathbb{C}). \\ \quad \forall k \in \mathbb{N}.\; (((n + 1) \leq k) \Rightarrow (c(k) = 0)) \Rightarrow  \\ \quad \quad \operatorname{is\text{-}countable}(\{z \in \mathbb{C} : (\neg (c(n) = 0) \wedge (\operatorname{cc\text{-}series\text{-}partial\text{-}sum}((\lambda j \in \mathbb{N}.\; (c(j) \cdot \operatorname{power}(z, j))), (n + 1)) = 0))\}) \end{array}}}
\vnbprop{polynomials-below-countable}{\vnbstmt{\begin{array}{@{}l@{}} \forall m \in \mathbb{N}, k. \\ \quad (\operatorname{is\text{-}cc\text{-}subring}(k) \wedge \operatorname{is\text{-}countable}(k)) \Rightarrow  \\ \quad \quad \operatorname{is\text{-}countable}(\{p \in (\mathbb{N} \to k) : \forall \mathit{pfk} \in \mathbb{N}.\; ((m \leq \mathit{pfk}) \Rightarrow (p(\mathit{pfk}) = 0))\}) \end{array}}}
\vnbprop{algebraic-numbers-of-degree-countable}{\vnbstmt{\begin{array}{@{}l@{}} \forall n \in \mathbb{N}, k. \\ \quad (\operatorname{is\text{-}cc\text{-}subring}(k) \wedge \operatorname{is\text{-}countable}(k)) \Rightarrow  \\ \quad \quad \operatorname{is\text{-}countable}(\{z \in \mathbb{C} : \exists c \in \{p \in (\mathbb{N} \to k) : \forall \mathit{pfk} \in \mathbb{N}.\; (((n + 1) \leq \mathit{pfk}) \Rightarrow (p(\mathit{pfk}) = 0))\}.\; (\neg (c(n) = 0) \wedge (\operatorname{cc\text{-}series\text{-}partial\text{-}sum}((\lambda j \in \mathbb{N}.\; (c(j) \cdot \operatorname{power}(z, j))), (n + 1)) = 0))\}) \end{array}}}
\vnbprop{algebraic-numbers-over-unfold}{\vnbstmt{\begin{array}{@{}l@{}} \forall k. \\ \quad (\operatorname{algebraic\text{-}numbers\text{-}over}(k) \simeq \{h \in \mathbb{C} : \operatorname{is\text{-}algebraic\text{-}over\text{-}field}(k, h)\}) \end{array}}}
\vnbprop{algebraic-numbers-countable}{\vnbstmt{\begin{array}{@{}l@{}} \forall k. \\ \quad (\operatorname{is\text{-}cc\text{-}subring}(k) \wedge \operatorname{is\text{-}countable}(k)) \Rightarrow  \\ \quad \quad \operatorname{is\text{-}countable}(\operatorname{algebraic\text{-}numbers\text{-}over}(k)) \end{array}}}
\vnbprop{qq-is-cc-subring}{\vnbstmt{\operatorname{is\text{-}cc\text{-}subring}(\mathbb{Q})}}
\vnbprop{algebraic-numbers-over-qq-countable}{\vnbstmt{\operatorname{is\text{-}countable}(\operatorname{algebraic\text{-}numbers\text{-}over}(\mathbb{Q}))}}

\vnbbookfile{theorem-library/algebraic-numbers}{theorem-library/algebraic-numbers}
\noindent{\small theorem-library/algebraic-numbers.scm -- ALGEBRAIC NUMBERS IN C OVER A SUBRING AND OVER A SUBFIELD (ALGNUM, B7, 2026-10-26): the characteristic polynomial, Props. 4.18, 4.19, 4.22, 4.24, 4.25, Cor. 4.23 (the field half) and Cor. 4.26 of the algebraic-numbers notes. Vocabulary: structure-library/algnum-vocab.scm ...}\par
\vnbprop{cc-sign-in-subring}{\vnbstmt{\begin{array}{@{}l@{}} \forall j \in \mathbb{N}, k. \\ \quad \operatorname{is\text{-}subring\text{-}of\text{-}cc}(k) \Rightarrow  \\ \quad \quad \operatorname{mpow}(\operatorname{ring\text{-}multiplicative\text{-}monoid}(\operatorname{normed\text{-}field\text{-}as\text{-}commutative\text{-}ring}(\operatorname{cc\text{-}normed\text{-}field})), \\ \quad \quad \quad \operatorname{neg}(\operatorname{normed\text{-}field\text{-}as\text{-}commutative\text{-}ring}(\operatorname{cc\text{-}normed\text{-}field}))(\operatorname{one}(\operatorname{normed\text{-}field\text{-}as\text{-}commutative\text{-}ring}(\operatorname{cc\text{-}normed\text{-}field}))), \\ \quad \quad \quad j) \\ \quad \quad \in\; k \wedge \\ \quad \quad \operatorname{mpow}(\operatorname{ring\text{-}multiplicative\text{-}monoid}(\operatorname{normed\text{-}field\text{-}as\text{-}commutative\text{-}ring}(\operatorname{cc\text{-}normed\text{-}field})), \\ \quad \quad \quad \operatorname{neg}(\operatorname{normed\text{-}field\text{-}as\text{-}commutative\text{-}ring}(\operatorname{cc\text{-}normed\text{-}field}))(\operatorname{one}(\operatorname{normed\text{-}field\text{-}as\text{-}commutative\text{-}ring}(\operatorname{cc\text{-}normed\text{-}field}))), \\ \quad \quad \quad j) \\ \quad \quad \in\; \mathbb{C} \end{array}}}
\vnbprop{cp-lin-mat}{\vnbstmt{\begin{array}{@{}l@{}} \forall n, b, e, z. \\ \quad ((n \in \mathbb{N}) \wedge ((b \in \operatorname{mat}(n, n, \mathbb{C})) \wedge ((e \in \operatorname{mat}(n, n, \mathbb{C})) \wedge (z \in \mathbb{C})))) \Rightarrow  \\ \quad \quad \operatorname{matof}(n, n, \\ \quad \quad \quad (\lambda {[}\operatorname{cpi\_}, \operatorname{cpj\_}] \in (\operatorname{interval}(1, n) \times \operatorname{interval}(1, n)).\; \\ \quad \quad \quad \quad \operatorname{entry}(b, \operatorname{cpi\_}, \operatorname{cpj\_}) + (z \cdot \operatorname{entry}(e, \operatorname{cpi\_}, \operatorname{cpj\_})))) \\ \quad \quad \in\; \operatorname{mat}(n, n, \mathbb{C}) \wedge \\ \quad \quad \forall i, j \in \operatorname{interval}(1, n). \\ \quad \quad \quad \operatorname{entry}(\operatorname{matof}(n, n, \\ \quad \quad \quad \quad \quad (\lambda {[}\operatorname{cpi\_}, \operatorname{cpj\_}] \in (\operatorname{interval}(1, n) \times \operatorname{interval}(1, n)).\; \\ \quad \quad \quad \quad \quad \quad \operatorname{entry}(b, \operatorname{cpi\_}, \operatorname{cpj\_}) + (z \cdot \operatorname{entry}(e, \operatorname{cpi\_}, \operatorname{cpj\_})))), \\ \quad \quad \quad \quad i, j) \\ \quad \quad \quad =\; \operatorname{entry}(b, i, j) + (z \cdot \operatorname{entry}(e, i, j)) \end{array}}}
\vnbprop{cp-lin-minor}{\vnbstmt{\begin{array}{@{}l@{}} \forall n, b, e, j, z. \\ \quad n \in \mathbb{N} \wedge \\ \quad b \in \operatorname{mat}((n + 1), (n + 1), \mathbb{C}) \wedge \\ \quad e \in \operatorname{mat}((n + 1), (n + 1), \mathbb{C}) \wedge \\ \quad j \in \operatorname{interval}(1, (n + 1)) \wedge \\ \quad z \in \mathbb{C} \Rightarrow \\ \quad \quad \operatorname{minor}(\operatorname{matof}((n + 1), (n + 1), \\ \quad \quad \quad \quad (\lambda {[}\operatorname{cpi\_}, \operatorname{cpj\_}] \in (\operatorname{interval}(1, (n + 1)) \times \operatorname{interval}(1, (n + 1))).\; \\ \quad \quad \quad \quad \quad \operatorname{entry}(b, \operatorname{cpi\_}, \operatorname{cpj\_}) + (z \cdot \operatorname{entry}(e, \operatorname{cpi\_}, \operatorname{cpj\_})))), \\ \quad \quad \quad 1, j, n) \\ \quad \quad =\; \operatorname{matof}(n, n, \\ \quad \quad \quad (\lambda {[}\operatorname{cpi\_}, \operatorname{cpj\_}] \in (\operatorname{interval}(1, n) \times \operatorname{interval}(1, n)).\; \\ \quad \quad \quad \quad \operatorname{entry}(\operatorname{minor}(b, 1, j, n), \operatorname{cpi\_}, \operatorname{cpj\_}) + (z \cdot \operatorname{entry}(\operatorname{minor}(e, 1, j, n), \operatorname{cpi\_}, \operatorname{cpj\_})))) \end{array}}}
\vnbprop{cp-zero-coefficients}{\vnbstmt{\begin{array}{@{}l@{}} \forall k. \\ \quad \operatorname{is\text{-}subring\text{-}of\text{-}cc}(k) \Rightarrow  \\ \quad \quad (\lambda \mathit{cpk} \in \mathbb{N}.\; 0) \in (\mathbb{N} \to k) \wedge \\ \quad \quad \forall m, \mathit{pfk} \in \mathbb{N}.\; ((m \leq \mathit{pfk}) \Rightarrow ((\lambda \mathit{cpk} \in \mathbb{N}.\; 0)(\mathit{pfk}) = 0)) \wedge \\ \quad \quad \forall a \in \mathbb{N}, z \in \mathbb{C}. \\ \quad \quad \quad \operatorname{cc\text{-}series\text{-}partial\text{-}sum}((\lambda j \in \mathbb{N}.\; ((\lambda \mathit{cpk} \in \mathbb{N}.\; 0)(j) \cdot \operatorname{power}(z, j))), a) = 0 \end{array}}}
\vnbprop{cp-sum-coefficients}{\vnbstmt{\begin{array}{@{}l@{}} \forall n \in \mathbb{N}, a, b, c. \\ \quad \operatorname{is\text{-}subring\text{-}of\text{-}cc}(a) \wedge \\ \quad b \in (\mathbb{N} \to a) \wedge \\ \quad c \in (\mathbb{N} \to a) \wedge \\ \quad \forall \mathit{pfk} \in \mathbb{N}.\; ((n \leq \mathit{pfk}) \Rightarrow (b(\mathit{pfk}) = 0)) \wedge \\ \quad \forall \mathit{pfk} \in \mathbb{N}.\; ((n \leq \mathit{pfk}) \Rightarrow (c(\mathit{pfk}) = 0)) \Rightarrow \\ \quad \quad (\lambda \mathit{cpk} \in \mathbb{N}.\; (b(\mathit{cpk}) + c(\mathit{cpk}))) \in (\mathbb{N} \to a) \wedge \\ \quad \quad \forall \mathit{pfk} \in \mathbb{N}. \\ \quad \quad \quad (n \leq \mathit{pfk}) \Rightarrow ((\lambda \mathit{cpk} \in \mathbb{N}.\; (b(\mathit{cpk}) + c(\mathit{cpk})))(\mathit{pfk}) = 0) \wedge \\ \quad \quad \forall k \in \mathbb{N}.\; ((\lambda \mathit{cpk} \in \mathbb{N}.\; (b(\mathit{cpk}) + c(\mathit{cpk})))(k) = (b(k) + c(k))) \wedge \\ \quad \quad \forall z \in \mathbb{C}. \\ \quad \quad \quad \operatorname{cc\text{-}series\text{-}partial\text{-}sum}((\lambda j \in \mathbb{N}.\; ((\lambda \mathit{cpk} \in \mathbb{N}.\; (b(\mathit{cpk}) + c(\mathit{cpk})))(j) \cdot \operatorname{power}(z, j))), n) \\ \quad \quad \quad =\; \operatorname{cc\text{-}series\text{-}partial\text{-}sum}((\lambda j \in \mathbb{N}.\; (b(j) \cdot \operatorname{power}(z, j))), n) + \operatorname{cc\text{-}series\text{-}partial\text{-}sum}((\lambda j \in \mathbb{N}.\; (c(j) \cdot \operatorname{power}(z, j))), n) \end{array}}}
\vnbprop{cp-shift-coefficients}{\vnbstmt{\begin{array}{@{}l@{}} \forall n \in \mathbb{N}, k, c, l, e. \\ \quad \operatorname{is\text{-}subring\text{-}of\text{-}cc}(k) \wedge \\ \quad c \in (\mathbb{N} \to k) \wedge \\ \quad \forall \mathit{pfk} \in \mathbb{N}.\; (((n + 1) \leq \mathit{pfk}) \Rightarrow (c(\mathit{pfk}) = 0)) \wedge \\ \quad l \in k \wedge \\ \quad e \in k \Rightarrow \\ \quad \quad (\lambda \mathit{cpk} \in \mathbb{N}.\; ((l \cdot c(\mathit{cpk})) + (e \cdot \operatorname{cf\text{-}cons}(0, c)(\mathit{cpk})))) \in (\mathbb{N} \to k) \wedge \\ \quad \quad \forall \mathit{pfk} \in \mathbb{N}. \\ \quad \quad \quad (((n + 1) + 1) \leq \mathit{pfk}) \Rightarrow  \\ \quad \quad \quad \quad (\lambda \mathit{cpk} \in \mathbb{N}.\; ((l \cdot c(\mathit{cpk})) + (e \cdot \operatorname{cf\text{-}cons}(0, c)(\mathit{cpk}))))(\mathit{pfk}) = 0 \wedge \\ \quad \quad (\lambda \mathit{cpk} \in \mathbb{N}.\; ((l \cdot c(\mathit{cpk})) + (e \cdot \operatorname{cf\text{-}cons}(0, c)(\mathit{cpk}))))((n + 1)) = (e \cdot c(n)) \wedge \\ \quad \quad \forall z \in \mathbb{C}. \\ \quad \quad \quad \operatorname{cc\text{-}series\text{-}partial\text{-}sum}((\lambda j \in \mathbb{N}.\; ((\lambda \mathit{cpk} \in \mathbb{N}.\; ((l \cdot c(\mathit{cpk})) + (e \cdot \operatorname{cf\text{-}cons}(0, c)(\mathit{cpk}))))(j) \cdot \operatorname{power}(z, j))), \\ \quad \quad \quad \quad ((n + 1) + 1)) \\ \quad \quad \quad =\; (l + (z \cdot e)) \cdot \operatorname{cc\text{-}series\text{-}partial\text{-}sum}((\lambda j \in \mathbb{N}.\; (c(j) \cdot \operatorname{power}(z, j))), (n + 1)) \end{array}}}
\vnbprop{cp-minor-in-subring}{\vnbstmt{\begin{array}{@{}l@{}} \forall n, k, s, j. \\ \quad ((n \in \mathbb{N}) \wedge ((s \in \operatorname{mat}((n + 1), (n + 1), k)) \wedge (j \in \operatorname{interval}(1, (n + 1))))) \Rightarrow  \\ \quad \quad \operatorname{minor}(s, 1, j, n) \in \operatorname{mat}(n, n, k) \end{array}}}
\vnbprop{cp-summand-value}{\vnbstmt{\begin{array}{@{}l@{}} \forall n, a, x. \\ \quad ((n \in \mathbb{N}) \wedge ((a \in \operatorname{mat}((n + 1), (n + 1), \mathbb{C})) \wedge (x \in \operatorname{interval}(1, (n + 1))))) \Rightarrow  \\ \quad \quad (\lambda j \in \operatorname{interval}(1, (n + 1)).\; \operatorname{mul}(\operatorname{normed\text{-}field\text{-}as\text{-}commutative\text{-}ring}(\operatorname{cc\text{-}normed\text{-}field}))(\operatorname{mpow}(\operatorname{ring\text{-}multiplicative\text{-}monoid}(\operatorname{normed\text{-}field\text{-}as\text{-}commutative\text{-}ring}(\operatorname{cc\text{-}normed\text{-}field})), \operatorname{neg}(\operatorname{normed\text{-}field\text{-}as\text{-}commutative\text{-}ring}(\operatorname{cc\text{-}normed\text{-}field}))(\operatorname{one}(\operatorname{normed\text{-}field\text{-}as\text{-}commutative\text{-}ring}(\operatorname{cc\text{-}normed\text{-}field}))), (j + 1)), \operatorname{mul}(\operatorname{normed\text{-}field\text{-}as\text{-}commutative\text{-}ring}(\operatorname{cc\text{-}normed\text{-}field}))(\operatorname{entry}(a, 1, j), \operatorname{det}(\operatorname{normed\text{-}field\text{-}as\text{-}commutative\text{-}ring}(\operatorname{cc\text{-}normed\text{-}field}), n, \operatorname{minor}(a, 1, j, n)))))(x) \\ \quad \quad =\; \operatorname{mpow}(\operatorname{ring\text{-}multiplicative\text{-}monoid}(\operatorname{normed\text{-}field\text{-}as\text{-}commutative\text{-}ring}(\operatorname{cc\text{-}normed\text{-}field})), \operatorname{neg}(\operatorname{normed\text{-}field\text{-}as\text{-}commutative\text{-}ring}(\operatorname{cc\text{-}normed\text{-}field}))(\operatorname{one}(\operatorname{normed\text{-}field\text{-}as\text{-}commutative\text{-}ring}(\operatorname{cc\text{-}normed\text{-}field}))), (x + 1)) \cdot (\operatorname{entry}(a, 1, x) \cdot \operatorname{det}(\operatorname{normed\text{-}field\text{-}as\text{-}commutative\text{-}ring}(\operatorname{cc\text{-}normed\text{-}field}), n, \operatorname{minor}(a, 1, x, n))) \wedge \\ \quad \quad (\lambda j \in \operatorname{interval}(1, (n + 1)).\; \operatorname{mul}(\operatorname{normed\text{-}field\text{-}as\text{-}commutative\text{-}ring}(\operatorname{cc\text{-}normed\text{-}field}))(\operatorname{mpow}(\operatorname{ring\text{-}multiplicative\text{-}monoid}(\operatorname{normed\text{-}field\text{-}as\text{-}commutative\text{-}ring}(\operatorname{cc\text{-}normed\text{-}field})), \operatorname{neg}(\operatorname{normed\text{-}field\text{-}as\text{-}commutative\text{-}ring}(\operatorname{cc\text{-}normed\text{-}field}))(\operatorname{one}(\operatorname{normed\text{-}field\text{-}as\text{-}commutative\text{-}ring}(\operatorname{cc\text{-}normed\text{-}field}))), (j + 1)), \operatorname{mul}(\operatorname{normed\text{-}field\text{-}as\text{-}commutative\text{-}ring}(\operatorname{cc\text{-}normed\text{-}field}))(\operatorname{entry}(a, 1, j), \operatorname{det}(\operatorname{normed\text{-}field\text{-}as\text{-}commutative\text{-}ring}(\operatorname{cc\text{-}normed\text{-}field}), n, \operatorname{minor}(a, 1, j, n)))))(x) \\ \quad \quad \in\; \mathbb{C} \end{array}}}
\vnbprop{cp-finsum-step}{\vnbstmt{\begin{array}{@{}l@{}} \forall n, a, k. \\ \quad ((n \in \mathbb{N}) \wedge ((a \in \operatorname{mat}((n + 1), (n + 1), \mathbb{C})) \wedge ((k \in \mathbb{N}) \wedge ((k + 1) \leq (n + 1))))) \Rightarrow  \\ \quad \quad \operatorname{finsum}(\operatorname{ring\text{-}additive\text{-}ag}(\operatorname{normed\text{-}field\text{-}as\text{-}commutative\text{-}ring}(\operatorname{cc\text{-}normed\text{-}field})), \\ \quad \quad \quad (\lambda j \in \operatorname{interval}(1, (n + 1)).\; \\ \quad \quad \quad \quad \operatorname{mul}(\operatorname{normed\text{-}field\text{-}as\text{-}commutative\text{-}ring}(\operatorname{cc\text{-}normed\text{-}field}))(\operatorname{mpow}(\operatorname{ring\text{-}multiplicative\text{-}monoid}(\operatorname{normed\text{-}field\text{-}as\text{-}commutative\text{-}ring}(\operatorname{cc\text{-}normed\text{-}field})), \\ \quad \quad \quad \quad \quad \quad \operatorname{neg}(\operatorname{normed\text{-}field\text{-}as\text{-}commutative\text{-}ring}(\operatorname{cc\text{-}normed\text{-}field}))(\operatorname{one}(\operatorname{normed\text{-}field\text{-}as\text{-}commutative\text{-}ring}(\operatorname{cc\text{-}normed\text{-}field}))), \\ \quad \quad \quad \quad \quad \quad (j + 1)), \\ \quad \quad \quad \quad \quad \operatorname{mul}(\operatorname{normed\text{-}field\text{-}as\text{-}commutative\text{-}ring}(\operatorname{cc\text{-}normed\text{-}field}))(\operatorname{entry}(a, 1, j), \\ \quad \quad \quad \quad \quad \quad \operatorname{det}(\operatorname{normed\text{-}field\text{-}as\text{-}commutative\text{-}ring}(\operatorname{cc\text{-}normed\text{-}field}), n, \operatorname{minor}(a, 1, j, n))))), \\ \quad \quad \quad \operatorname{interval}(1, k)) \\ \quad \quad \in\; \mathbb{C} \wedge \\ \quad \quad \operatorname{finsum}(\operatorname{ring\text{-}additive\text{-}ag}(\operatorname{normed\text{-}field\text{-}as\text{-}commutative\text{-}ring}(\operatorname{cc\text{-}normed\text{-}field})), \\ \quad \quad \quad (\lambda j \in \operatorname{interval}(1, (n + 1)).\; \\ \quad \quad \quad \quad \operatorname{mul}(\operatorname{normed\text{-}field\text{-}as\text{-}commutative\text{-}ring}(\operatorname{cc\text{-}normed\text{-}field}))(\operatorname{mpow}(\operatorname{ring\text{-}multiplicative\text{-}monoid}(\operatorname{normed\text{-}field\text{-}as\text{-}commutative\text{-}ring}(\operatorname{cc\text{-}normed\text{-}field})), \\ \quad \quad \quad \quad \quad \quad \operatorname{neg}(\operatorname{normed\text{-}field\text{-}as\text{-}commutative\text{-}ring}(\operatorname{cc\text{-}normed\text{-}field}))(\operatorname{one}(\operatorname{normed\text{-}field\text{-}as\text{-}commutative\text{-}ring}(\operatorname{cc\text{-}normed\text{-}field}))), \\ \quad \quad \quad \quad \quad \quad (j + 1)), \\ \quad \quad \quad \quad \quad \operatorname{mul}(\operatorname{normed\text{-}field\text{-}as\text{-}commutative\text{-}ring}(\operatorname{cc\text{-}normed\text{-}field}))(\operatorname{entry}(a, 1, j), \\ \quad \quad \quad \quad \quad \quad \operatorname{det}(\operatorname{normed\text{-}field\text{-}as\text{-}commutative\text{-}ring}(\operatorname{cc\text{-}normed\text{-}field}), n, \operatorname{minor}(a, 1, j, n))))), \\ \quad \quad \quad \operatorname{interval}(1, (k + 1))) \\ \quad \quad =\; \operatorname{finsum}(\operatorname{ring\text{-}additive\text{-}ag}(\operatorname{normed\text{-}field\text{-}as\text{-}commutative\text{-}ring}(\operatorname{cc\text{-}normed\text{-}field})), (\lambda j \in \operatorname{interval}(1, (n + 1)).\; \operatorname{mul}(\operatorname{normed\text{-}field\text{-}as\text{-}commutative\text{-}ring}(\operatorname{cc\text{-}normed\text{-}field}))(\operatorname{mpow}(\operatorname{ring\text{-}multiplicative\text{-}monoid}(\operatorname{normed\text{-}field\text{-}as\text{-}commutative\text{-}ring}(\operatorname{cc\text{-}normed\text{-}field})), \operatorname{neg}(\operatorname{normed\text{-}field\text{-}as\text{-}commutative\text{-}ring}(\operatorname{cc\text{-}normed\text{-}field}))(\operatorname{one}(\operatorname{normed\text{-}field\text{-}as\text{-}commutative\text{-}ring}(\operatorname{cc\text{-}normed\text{-}field}))), (j + 1)), \operatorname{mul}(\operatorname{normed\text{-}field\text{-}as\text{-}commutative\text{-}ring}(\operatorname{cc\text{-}normed\text{-}field}))(\operatorname{entry}(a, 1, j), \operatorname{det}(\operatorname{normed\text{-}field\text{-}as\text{-}commutative\text{-}ring}(\operatorname{cc\text{-}normed\text{-}field}), n, \operatorname{minor}(a, 1, j, n))))), \operatorname{interval}(1, k)) + (\lambda j \in \operatorname{interval}(1, (n + 1)).\; \operatorname{mul}(\operatorname{normed\text{-}field\text{-}as\text{-}commutative\text{-}ring}(\operatorname{cc\text{-}normed\text{-}field}))(\operatorname{mpow}(\operatorname{ring\text{-}multiplicative\text{-}monoid}(\operatorname{normed\text{-}field\text{-}as\text{-}commutative\text{-}ring}(\operatorname{cc\text{-}normed\text{-}field})), \operatorname{neg}(\operatorname{normed\text{-}field\text{-}as\text{-}commutative\text{-}ring}(\operatorname{cc\text{-}normed\text{-}field}))(\operatorname{one}(\operatorname{normed\text{-}field\text{-}as\text{-}commutative\text{-}ring}(\operatorname{cc\text{-}normed\text{-}field}))), (j + 1)), \operatorname{mul}(\operatorname{normed\text{-}field\text{-}as\text{-}commutative\text{-}ring}(\operatorname{cc\text{-}normed\text{-}field}))(\operatorname{entry}(a, 1, j), \operatorname{det}(\operatorname{normed\text{-}field\text{-}as\text{-}commutative\text{-}ring}(\operatorname{cc\text{-}normed\text{-}field}), n, \operatorname{minor}(a, 1, j, n)))))((k + 1)) \end{array}}}
\vnbprop{cp-cofactor-partial-sums}{\vnbstmt{\begin{array}{@{}l@{}} \forall k, n \in \mathbb{N}, a, b, e. \\ \quad \operatorname{is\text{-}subring\text{-}of\text{-}cc}(a) \wedge \\ \quad b \in \operatorname{mat}((n + 1), (n + 1), a) \wedge \\ \quad e \in \operatorname{mat}((n + 1), (n + 1), a) \wedge \\ \quad \forall \mathit{cpb}, \mathit{cpe}. \\ \quad \quad ((\mathit{cpb} \in \operatorname{mat}(n, n, a)) \wedge (\mathit{cpe} \in \operatorname{mat}(n, n, a))) \Rightarrow  \\ \quad \quad \quad \exists c \in (\mathbb{N} \to a). \\ \quad \quad \quad \quad \forall \mathit{pfk} \in \mathbb{N}.\; (((n + 1) \leq \mathit{pfk}) \Rightarrow (c(\mathit{pfk}) = 0)) \wedge \\ \quad \quad \quad \quad c(n) = \operatorname{det}(\operatorname{normed\text{-}field\text{-}as\text{-}commutative\text{-}ring}(\operatorname{cc\text{-}normed\text{-}field}), n, \mathit{cpe}) \wedge \\ \quad \quad \quad \quad \forall z \in \mathbb{C}. \\ \quad \quad \quad \quad \quad \operatorname{det}(\operatorname{normed\text{-}field\text{-}as\text{-}commutative\text{-}ring}(\operatorname{cc\text{-}normed\text{-}field}), \\ \quad \quad \quad \quad \quad \quad n, \\ \quad \quad \quad \quad \quad \quad \operatorname{matof}(n, n, \\ \quad \quad \quad \quad \quad \quad \quad (\lambda {[}\operatorname{cpi\_}, \operatorname{cpj\_}] \in (\operatorname{interval}(1, n) \times \operatorname{interval}(1, n)).\; \\ \quad \quad \quad \quad \quad \quad \quad \quad \operatorname{entry}(\mathit{cpb}, \operatorname{cpi\_}, \operatorname{cpj\_}) + (z \cdot \operatorname{entry}(\mathit{cpe}, \operatorname{cpi\_}, \operatorname{cpj\_}))))) \\ \quad \quad \quad \quad \quad =\; \operatorname{cc\text{-}series\text{-}partial\text{-}sum}((\lambda \mathit{pfj} \in \mathbb{N}.\; (c(\mathit{pfj}) \cdot \operatorname{power}(z, \mathit{pfj}))), (n + 1)) \wedge \\ \quad k \leq (n + 1) \Rightarrow \\ \quad \quad \exists d \in (\mathbb{N} \to a). \\ \quad \quad \quad \forall \mathit{pfk} \in \mathbb{N}.\; ((((n + 1) + 1) \leq \mathit{pfk}) \Rightarrow (d(\mathit{pfk}) = 0)) \wedge \\ \quad \quad \quad d((n + 1)) \\ \quad \quad \quad =\; \operatorname{finsum}(\operatorname{ring\text{-}additive\text{-}ag}(\operatorname{normed\text{-}field\text{-}as\text{-}commutative\text{-}ring}(\operatorname{cc\text{-}normed\text{-}field})), \\ \quad \quad \quad \quad (\lambda j \in \operatorname{interval}(1, (n + 1)).\; \\ \quad \quad \quad \quad \quad \operatorname{mul}(\operatorname{normed\text{-}field\text{-}as\text{-}commutative\text{-}ring}(\operatorname{cc\text{-}normed\text{-}field}))(\operatorname{mpow}(\operatorname{ring\text{-}multiplicative\text{-}monoid}(\operatorname{normed\text{-}field\text{-}as\text{-}commutative\text{-}ring}(\operatorname{cc\text{-}normed\text{-}field})), \\ \quad \quad \quad \quad \quad \quad \quad \operatorname{neg}(\operatorname{normed\text{-}field\text{-}as\text{-}commutative\text{-}ring}(\operatorname{cc\text{-}normed\text{-}field}))(\operatorname{one}(\operatorname{normed\text{-}field\text{-}as\text{-}commutative\text{-}ring}(\operatorname{cc\text{-}normed\text{-}field}))), \\ \quad \quad \quad \quad \quad \quad \quad (j + 1)), \\ \quad \quad \quad \quad \quad \quad \operatorname{mul}(\operatorname{normed\text{-}field\text{-}as\text{-}commutative\text{-}ring}(\operatorname{cc\text{-}normed\text{-}field}))(\operatorname{entry}(e, 1, j), \\ \quad \quad \quad \quad \quad \quad \quad \operatorname{det}(\operatorname{normed\text{-}field\text{-}as\text{-}commutative\text{-}ring}(\operatorname{cc\text{-}normed\text{-}field}), n, \operatorname{minor}(e, 1, j, n))))), \\ \quad \quad \quad \quad \operatorname{interval}(1, k)) \wedge \\ \quad \quad \quad \forall z \in \mathbb{C}. \\ \quad \quad \quad \quad \operatorname{finsum}(\operatorname{ring\text{-}additive\text{-}ag}(\operatorname{normed\text{-}field\text{-}as\text{-}commutative\text{-}ring}(\operatorname{cc\text{-}normed\text{-}field})), \\ \quad \quad \quad \quad \quad (\lambda j \in \operatorname{interval}(1, (n + 1)).\; \\ \quad \quad \quad \quad \quad \quad \operatorname{mul}(\operatorname{normed\text{-}field\text{-}as\text{-}commutative\text{-}ring}(\operatorname{cc\text{-}normed\text{-}field}))(\operatorname{mpow}(\operatorname{ring\text{-}multiplicative\text{-}monoid}(\operatorname{normed\text{-}field\text{-}as\text{-}commutative\text{-}ring}(\operatorname{cc\text{-}normed\text{-}field})), \\ \quad \quad \quad \quad \quad \quad \quad \quad \operatorname{neg}(\operatorname{normed\text{-}field\text{-}as\text{-}commutative\text{-}ring}(\operatorname{cc\text{-}normed\text{-}field}))(\operatorname{one}(\operatorname{normed\text{-}field\text{-}as\text{-}commutative\text{-}ring}(\operatorname{cc\text{-}normed\text{-}field}))), \\ \quad \quad \quad \quad \quad \quad \quad \quad (j + 1)), \\ \quad \quad \quad \quad \quad \quad \quad \operatorname{mul}(\operatorname{normed\text{-}field\text{-}as\text{-}commutative\text{-}ring}(\operatorname{cc\text{-}normed\text{-}field}))(\operatorname{entry}(\operatorname{matof}((n + 1), (n + 1), \\ \quad \quad \quad \quad \quad \quad \quad \quad \quad \quad (\lambda {[}\operatorname{cpi\_}, \operatorname{cpj\_}] \in (\operatorname{interval}(1, (n + 1)) \times \operatorname{interval}(1, (n + 1))).\; \\ \quad \quad \quad \quad \quad \quad \quad \quad \quad \quad \quad \operatorname{entry}(b, \operatorname{cpi\_}, \operatorname{cpj\_}) + (z \cdot \operatorname{entry}(e, \operatorname{cpi\_}, \operatorname{cpj\_})))), \\ \quad \quad \quad \quad \quad \quad \quad \quad \quad 1, j), \\ \quad \quad \quad \quad \quad \quad \quad \quad \operatorname{det}(\operatorname{normed\text{-}field\text{-}as\text{-}commutative\text{-}ring}(\operatorname{cc\text{-}normed\text{-}field}), \\ \quad \quad \quad \quad \quad \quad \quad \quad \quad n, \\ \quad \quad \quad \quad \quad \quad \quad \quad \quad \operatorname{minor}(\operatorname{matof}((n + 1), (n + 1), \\ \quad \quad \quad \quad \quad \quad \quad \quad \quad \quad \quad (\lambda {[}\operatorname{cpi\_}, \operatorname{cpj\_}] \in (\operatorname{interval}(1, (n + 1)) \times \operatorname{interval}(1, (n + 1))).\; \\ \quad \quad \quad \quad \quad \quad \quad \quad \quad \quad \quad \quad \operatorname{entry}(b, \operatorname{cpi\_}, \operatorname{cpj\_}) + (z \cdot \operatorname{entry}(e, \operatorname{cpi\_}, \operatorname{cpj\_})))), \\ \quad \quad \quad \quad \quad \quad \quad \quad \quad \quad 1, j, n))))), \\ \quad \quad \quad \quad \quad \operatorname{interval}(1, k)) \\ \quad \quad \quad \quad =\; \operatorname{cc\text{-}series\text{-}partial\text{-}sum}((\lambda \mathit{pfj} \in \mathbb{N}.\; (d(\mathit{pfj}) \cdot \operatorname{power}(z, \mathit{pfj}))), ((n + 1) + 1)) \end{array}}}
\vnbprop{det-linear-pencil-polynomial}{\vnbstmt{\begin{array}{@{}l@{}} \forall n \in \mathbb{N}, k, b, e. \\ \quad (\operatorname{is\text{-}subring\text{-}of\text{-}cc}(k) \wedge ((b \in \operatorname{mat}(n, n, k)) \wedge (e \in \operatorname{mat}(n, n, k)))) \Rightarrow  \\ \quad \quad \exists c \in (\mathbb{N} \to k). \\ \quad \quad \quad \forall \mathit{pfk} \in \mathbb{N}.\; (((n + 1) \leq \mathit{pfk}) \Rightarrow (c(\mathit{pfk}) = 0)) \wedge \\ \quad \quad \quad c(n) = \operatorname{det}(\operatorname{normed\text{-}field\text{-}as\text{-}commutative\text{-}ring}(\operatorname{cc\text{-}normed\text{-}field}), n, e) \wedge \\ \quad \quad \quad \forall z \in \mathbb{C}. \\ \quad \quad \quad \quad \operatorname{det}(\operatorname{normed\text{-}field\text{-}as\text{-}commutative\text{-}ring}(\operatorname{cc\text{-}normed\text{-}field}), \\ \quad \quad \quad \quad \quad n, \\ \quad \quad \quad \quad \quad \operatorname{matof}(n, n, \\ \quad \quad \quad \quad \quad \quad (\lambda {[}\operatorname{cpi\_}, \operatorname{cpj\_}] \in (\operatorname{interval}(1, n) \times \operatorname{interval}(1, n)).\; \\ \quad \quad \quad \quad \quad \quad \quad \operatorname{entry}(b, \operatorname{cpi\_}, \operatorname{cpj\_}) + (z \cdot \operatorname{entry}(e, \operatorname{cpi\_}, \operatorname{cpj\_}))))) \\ \quad \quad \quad \quad =\; \operatorname{cc\text{-}series\text{-}partial\text{-}sum}((\lambda j \in \mathbb{N}.\; (c(j) \cdot \operatorname{power}(z, j))), (n + 1)) \end{array}}}
\vnbprop{cp-negi-mat}{\vnbstmt{\begin{array}{@{}l@{}} \forall n \in \mathbb{N}, k. \\ \quad \operatorname{is\text{-}subring\text{-}of\text{-}cc}(k) \Rightarrow  \\ \quad \quad \operatorname{matof}(n, n, \\ \quad \quad \quad (\lambda {[}\operatorname{cqi\_}, \operatorname{cqj\_}] \in (\operatorname{interval}(1, n) \times \operatorname{interval}(1, n)).\; \left(\text{if }(\operatorname{cqi\_} = \operatorname{cqj\_})\text{ then }-1\text{ else }0\right))) \\ \quad \quad \in\; \operatorname{mat}(n, n, k) \wedge \\ \quad \quad \operatorname{matof}(n, n, \\ \quad \quad \quad (\lambda {[}\operatorname{cqi\_}, \operatorname{cqj\_}] \in (\operatorname{interval}(1, n) \times \operatorname{interval}(1, n)).\; \left(\text{if }(\operatorname{cqi\_} = \operatorname{cqj\_})\text{ then }-1\text{ else }0\right))) \\ \quad \quad \in\; \operatorname{mat}(n, n, \mathbb{C}) \wedge \\ \quad \quad \forall i, j \in \operatorname{interval}(1, n). \\ \quad \quad \quad \operatorname{entry}(\operatorname{matof}(n, n, \\ \quad \quad \quad \quad \quad (\lambda {[}\operatorname{cqi\_}, \operatorname{cqj\_}] \in (\operatorname{interval}(1, n) \times \operatorname{interval}(1, n)).\; \left(\text{if }(\operatorname{cqi\_} = \operatorname{cqj\_})\text{ then }-1\text{ else }0\right))), \\ \quad \quad \quad \quad i, j) \\ \quad \quad \quad =\; \left(\text{if }(i = j)\text{ then }-1\text{ else }0\right) \end{array}}}
\vnbprop{cc-is-subring-of-cc}{\vnbstmt{\operatorname{is\text{-}subring\text{-}of\text{-}cc}(\mathbb{C})}}
\vnbprop{cp-negi-det}{\vnbstmt{\begin{array}{@{}l@{}} \forall n \in \mathbb{N}. \\ \quad \operatorname{det}(\operatorname{normed\text{-}field\text{-}as\text{-}commutative\text{-}ring}(\operatorname{cc\text{-}normed\text{-}field}), \\ \quad \quad n, \\ \quad \quad \operatorname{matof}(n, n, \\ \quad \quad \quad (\lambda {[}\operatorname{cqi\_}, \operatorname{cqj\_}] \in (\operatorname{interval}(1, n) \times \operatorname{interval}(1, n)).\; \left(\text{if }(\operatorname{cqi\_} = \operatorname{cqj\_})\text{ then }-1\text{ else }0\right)))) \\ \quad =\; \operatorname{mpow}(\operatorname{ring\text{-}multiplicative\text{-}monoid}(\operatorname{normed\text{-}field\text{-}as\text{-}commutative\text{-}ring}(\operatorname{cc\text{-}normed\text{-}field})), \\ \quad \quad \operatorname{neg}(\operatorname{normed\text{-}field\text{-}as\text{-}commutative\text{-}ring}(\operatorname{cc\text{-}normed\text{-}field}))(\operatorname{one}(\operatorname{normed\text{-}field\text{-}as\text{-}commutative\text{-}ring}(\operatorname{cc\text{-}normed\text{-}field}))), \\ \quad \quad n) \end{array}}}
\vnbprop{characteristic-polynomial}{\vnbstmt{\begin{array}{@{}l@{}} \forall n \in \mathbb{N}, k.; \forall b \in \operatorname{mat}(n, n, k). \\ \quad \operatorname{is\text{-}subring\text{-}of\text{-}cc}(k) \Rightarrow  \\ \quad \quad \exists c \in (\mathbb{N} \to k). \\ \quad \quad \quad \forall \mathit{pfk} \in \mathbb{N}.\; (((n + 1) \leq \mathit{pfk}) \Rightarrow (c(\mathit{pfk}) = 0)) \wedge \\ \quad \quad \quad c(n) \\ \quad \quad \quad =\; \operatorname{mpow}(\operatorname{ring\text{-}multiplicative\text{-}monoid}(\operatorname{normed\text{-}field\text{-}as\text{-}commutative\text{-}ring}(\operatorname{cc\text{-}normed\text{-}field})), \\ \quad \quad \quad \quad \operatorname{neg}(\operatorname{normed\text{-}field\text{-}as\text{-}commutative\text{-}ring}(\operatorname{cc\text{-}normed\text{-}field}))(\operatorname{one}(\operatorname{normed\text{-}field\text{-}as\text{-}commutative\text{-}ring}(\operatorname{cc\text{-}normed\text{-}field}))), \\ \quad \quad \quad \quad n) \wedge \\ \quad \quad \quad \forall z \in \mathbb{C}. \\ \quad \quad \quad \quad \operatorname{det}(\operatorname{normed\text{-}field\text{-}as\text{-}commutative\text{-}ring}(\operatorname{cc\text{-}normed\text{-}field}), \\ \quad \quad \quad \quad \quad n, \\ \quad \quad \quad \quad \quad \operatorname{matof}(n, n, \\ \quad \quad \quad \quad \quad \quad (\lambda {[}\operatorname{cpi\_}, \operatorname{cpj\_}] \in (\operatorname{interval}(1, n) \times \operatorname{interval}(1, n)).\; \\ \quad \quad \quad \quad \quad \quad \quad \operatorname{entry}(b, \operatorname{cpi\_}, \operatorname{cpj\_}) + (z \cdot \operatorname{entry}(\operatorname{matof}(n, n, (\lambda {[}\operatorname{cqi\_}, \operatorname{cqj\_}] \in (\operatorname{interval}(1, n) \times \operatorname{interval}(1, n)).\; \left(\text{if }(\operatorname{cqi\_} = \operatorname{cqj\_})\text{ then }-1\text{ else }0\right))), \operatorname{cpi\_}, \operatorname{cpj\_}))))) \\ \quad \quad \quad \quad =\; \operatorname{cc\text{-}series\text{-}partial\text{-}sum}((\lambda j \in \mathbb{N}.\; (c(j) \cdot \operatorname{power}(z, j))), (n + 1)) \end{array}}}
\vnbprop{am-generator-matrix}{\vnbstmt{\begin{array}{@{}l@{}} \forall n, g. \\ \quad ((n \in \mathbb{N}) \wedge (g \in (\operatorname{interval}(1, n) \to \mathbb{C}))) \Rightarrow  \\ \quad \quad \operatorname{matof}(n, n, (\lambda {[}\operatorname{amp\_}, \operatorname{amq\_}] \in (\operatorname{interval}(1, n) \times \operatorname{interval}(1, n)).\; g(\operatorname{amp\_}))) \\ \quad \quad \in\; \operatorname{mat}(n, n, \mathbb{C}) \wedge \\ \quad \quad \forall i, j \in \operatorname{interval}(1, n). \\ \quad \quad \quad \operatorname{entry}(\operatorname{matof}(n, n, (\lambda {[}\operatorname{amp\_}, \operatorname{amq\_}] \in (\operatorname{interval}(1, n) \times \operatorname{interval}(1, n)).\; g(\operatorname{amp\_}))), \\ \quad \quad \quad \quad i, j) \\ \quad \quad \quad =\; g(i) \end{array}}}
\vnbprop{am-pencil-kills-generators}{\vnbstmt{\begin{array}{@{}l@{}} \forall n, m, g, r.; \forall i, l \in \operatorname{interval}(1, n). \\ \quad n \in \mathbb{N} \wedge \\ \quad m \in \operatorname{mat}(n, n, \mathbb{C}) \wedge \\ \quad g \in (\operatorname{interval}(1, n) \to \mathbb{C}) \wedge \\ \quad r \in \mathbb{C} \wedge \\ \quad \forall \mathit{ami} \in \operatorname{interval}(1, n). \\ \quad \quad r \cdot g(\mathit{ami}) \\ \quad \quad =\; \operatorname{finsum}(\operatorname{ring\text{-}additive\text{-}ag}(\operatorname{normed\text{-}field\text{-}as\text{-}commutative\text{-}ring}(\operatorname{cc\text{-}normed\text{-}field})), \\ \quad \quad \quad (\lambda j \in \operatorname{interval}(1, n).\; (\operatorname{entry}(m, \mathit{ami}, j) \cdot g(j))), \\ \quad \quad \quad \operatorname{interval}(1, n)) \Rightarrow \\ \quad \quad \operatorname{entry}(\operatorname{matmul}(\operatorname{normed\text{-}field\text{-}as\text{-}commutative\text{-}ring}(\operatorname{cc\text{-}normed\text{-}field}), \\ \quad \quad \quad \quad \operatorname{matof}(n, n, \\ \quad \quad \quad \quad \quad (\lambda {[}\operatorname{cpi\_}, \operatorname{cpj\_}] \in (\operatorname{interval}(1, n) \times \operatorname{interval}(1, n)).\; \\ \quad \quad \quad \quad \quad \quad \operatorname{entry}(m, \operatorname{cpi\_}, \operatorname{cpj\_}) + (r \cdot \operatorname{entry}(\operatorname{matof}(n, n, (\lambda {[}\operatorname{cqi\_}, \operatorname{cqj\_}] \in (\operatorname{interval}(1, n) \times \operatorname{interval}(1, n)).\; \left(\text{if }(\operatorname{cqi\_} = \operatorname{cqj\_})\text{ then }-1\text{ else }0\right))), \operatorname{cpi\_}, \operatorname{cpj\_})))), \\ \quad \quad \quad \quad \operatorname{matof}(n, n, (\lambda {[}\operatorname{amp\_}, \operatorname{amq\_}] \in (\operatorname{interval}(1, n) \times \operatorname{interval}(1, n)).\; g(\operatorname{amp\_})))), \\ \quad \quad \quad i, l) \\ \quad \quad =\; 0 \end{array}}}
\vnbprop{am-det-kills-generators}{\vnbstmt{\begin{array}{@{}l@{}} \forall n, m, g, r.; \forall k \in \operatorname{interval}(1, (n + 1)). \\ \quad n \in \mathbb{N} \wedge \\ \quad m \in \operatorname{mat}((n + 1), (n + 1), \mathbb{C}) \wedge \\ \quad g \in (\operatorname{interval}(1, (n + 1)) \to \mathbb{C}) \wedge \\ \quad r \in \mathbb{C} \wedge \\ \quad \forall i \in \operatorname{interval}(1, (n + 1)). \\ \quad \quad r \cdot g(i) \\ \quad \quad =\; \operatorname{finsum}(\operatorname{ring\text{-}additive\text{-}ag}(\operatorname{normed\text{-}field\text{-}as\text{-}commutative\text{-}ring}(\operatorname{cc\text{-}normed\text{-}field})), \\ \quad \quad \quad (\lambda j \in \operatorname{interval}(1, (n + 1)).\; (\operatorname{entry}(m, i, j) \cdot g(j))), \\ \quad \quad \quad \operatorname{interval}(1, (n + 1))) \Rightarrow \\ \quad \quad \operatorname{det}(\operatorname{normed\text{-}field\text{-}as\text{-}commutative\text{-}ring}(\operatorname{cc\text{-}normed\text{-}field}), (n + 1), \operatorname{matof}((n + 1), (n + 1), (\lambda {[}\operatorname{cpi\_}, \operatorname{cpj\_}] \in (\operatorname{interval}(1, (n + 1)) \times \operatorname{interval}(1, (n + 1))).\; (\operatorname{entry}(m, \operatorname{cpi\_}, \operatorname{cpj\_}) + (r \cdot \operatorname{entry}(\operatorname{matof}((n + 1), (n + 1), (\lambda {[}\operatorname{cqi\_}, \operatorname{cqj\_}] \in (\operatorname{interval}(1, (n + 1)) \times \operatorname{interval}(1, (n + 1))).\; \left(\text{if }(\operatorname{cqi\_} = \operatorname{cqj\_})\text{ then }-1\text{ else }0\right))), \operatorname{cpi\_}, \operatorname{cpj\_})))))) \cdot g(k) \\ \quad \quad =\; 0 \end{array}}}
\vnbprop{cc-sign-square}{\vnbstmt{\begin{array}{@{}l@{}} \forall j \in \mathbb{N}. \\ \quad \operatorname{mpow}(\operatorname{ring\text{-}multiplicative\text{-}monoid}(\operatorname{normed\text{-}field\text{-}as\text{-}commutative\text{-}ring}(\operatorname{cc\text{-}normed\text{-}field})), \operatorname{neg}(\operatorname{normed\text{-}field\text{-}as\text{-}commutative\text{-}ring}(\operatorname{cc\text{-}normed\text{-}field}))(\operatorname{one}(\operatorname{normed\text{-}field\text{-}as\text{-}commutative\text{-}ring}(\operatorname{cc\text{-}normed\text{-}field}))), j) \cdot \operatorname{mpow}(\operatorname{ring\text{-}multiplicative\text{-}monoid}(\operatorname{normed\text{-}field\text{-}as\text{-}commutative\text{-}ring}(\operatorname{cc\text{-}normed\text{-}field})), \operatorname{neg}(\operatorname{normed\text{-}field\text{-}as\text{-}commutative\text{-}ring}(\operatorname{cc\text{-}normed\text{-}field}))(\operatorname{one}(\operatorname{normed\text{-}field\text{-}as\text{-}commutative\text{-}ring}(\operatorname{cc\text{-}normed\text{-}field}))), j) \\ \quad =\; 1 \end{array}}}
\vnbprop{cc-lincomb-unfold}{\vnbstmt{\begin{array}{@{}l@{}} \forall n, b, g. \\ \quad (\operatorname{cc\text{-}lincomb}(n, b, g) \simeq \operatorname{finsum}(\operatorname{ring\text{-}additive\text{-}ag}(\operatorname{normed\text{-}field\text{-}as\text{-}commutative\text{-}ring}(\operatorname{cc\text{-}normed\text{-}field})), (\lambda j \in \operatorname{interval}(1, n).\; (b(j) \cdot g(j))), \operatorname{interval}(1, n))) \end{array}}}
\vnbprop{cc-lincomb-zero-generators}{\vnbstmt{\begin{array}{@{}l@{}} \forall n, b, g. \\ \quad ((n \in \mathbb{N}) \wedge (\forall k \in \operatorname{interval}(1, n).\; (b(k) \in \mathbb{C}) \wedge \forall k \in \operatorname{interval}(1, n).\; (g(k) = 0))) \Rightarrow  \\ \quad \quad \operatorname{cc\text{-}lincomb}(n, b, g) = 0 \end{array}}}
\vnbprop{algebraic-monic-of-invariant-space}{\vnbstmt{\begin{array}{@{}l@{}} \forall a, r, v, m. \\ \quad \operatorname{is\text{-}subring\text{-}of\text{-}cc}(a) \wedge \\ \quad r \in \mathbb{C} \wedge \\ \quad m \in \mathbb{N} \wedge \\ \quad \operatorname{is\text{-}cc\text{-}space\text{-}dim\text{-}le}(a, v, m) \wedge \\ \quad \exists z \in v.\; \neg (z = 0) \wedge \\ \quad \forall \mathit{amrv} \in v.\; ((r \cdot \mathit{amrv}) \in v) \Rightarrow \\ \quad \quad \exists n \in \mathbb{N}. \\ \quad \quad \quad n \leq m \wedge \\ \quad \quad \quad \exists c \in (\mathbb{N} \to a). \\ \quad \quad \quad \quad \forall k \in \mathbb{N}.\; (((n + 1) \leq k) \Rightarrow (c(k) = 0)) \wedge \\ \quad \quad \quad \quad c(n) = 1 \wedge \\ \quad \quad \quad \quad \operatorname{cc\text{-}series\text{-}partial\text{-}sum}((\lambda j \in \mathbb{N}.\; (c(j) \cdot \operatorname{power}(r, j))), (n + 1)) = 0 \end{array}}}
\vnbprop{fg-module-dim-le}{\vnbstmt{\begin{array}{@{}l@{}} \forall a, v. \\ \quad \operatorname{is\text{-}fg\text{-}cc\text{-}module\text{-}over}(a, v) \Rightarrow  \\ \quad \quad \exists d \in \mathbb{N}.\; \operatorname{is\text{-}cc\text{-}space\text{-}dim\text{-}le}(a, v, d) \end{array}}}
\vnbprop{dim-le-fg-module}{\vnbstmt{\begin{array}{@{}l@{}} \forall a, v, d. \\ \quad \operatorname{is\text{-}cc\text{-}space\text{-}dim\text{-}le}(a, v, d) \Rightarrow \operatorname{is\text{-}fg\text{-}cc\text{-}module\text{-}over}(a, v) \end{array}}}
\vnbprop{algebraic-of-invariant-module}{\vnbstmt{\begin{array}{@{}l@{}} \forall a, r, v. \\ \quad \operatorname{is\text{-}subring\text{-}of\text{-}cc}(a) \wedge \\ \quad r \in \mathbb{C} \wedge \\ \quad \operatorname{is\text{-}fg\text{-}cc\text{-}module\text{-}over}(a, v) \wedge \\ \quad \exists z \in v.\; \neg (z = 0) \wedge \\ \quad \forall \mathit{amrv} \in v.\; ((r \cdot \mathit{amrv}) \in v) \Rightarrow \\ \quad \quad \operatorname{is\text{-}algebraic\text{-}over}(a, r) \end{array}}}
\vnbprop{cc-lincomb-add}{\vnbstmt{\begin{array}{@{}l@{}} \forall n, a, b, g. \\ \quad n \in \mathbb{N} \wedge \\ \quad a \in (\operatorname{interval}(1, n) \to \mathbb{C}) \wedge \\ \quad b \in (\operatorname{interval}(1, n) \to \mathbb{C}) \wedge \\ \quad g \in (\operatorname{interval}(1, n) \to \mathbb{C}) \Rightarrow \\ \quad \quad \operatorname{cc\text{-}lincomb}(n, a, g) + \operatorname{cc\text{-}lincomb}(n, b, g) \\ \quad \quad =\; \operatorname{cc\text{-}lincomb}(n, (\lambda j \in \operatorname{interval}(1, n).\; (a(j) + b(j))), g) \end{array}}}
\vnbprop{cc-lincomb-scale}{\vnbstmt{\begin{array}{@{}l@{}} \forall n, a, b, g. \\ \quad ((n \in \mathbb{N}) \wedge ((a \in \mathbb{C}) \wedge ((b \in (\operatorname{interval}(1, n) \to \mathbb{C})) \wedge (g \in (\operatorname{interval}(1, n) \to \mathbb{C}))))) \Rightarrow  \\ \quad \quad (a \cdot \operatorname{cc\text{-}lincomb}(n, b, g)) = \operatorname{cc\text{-}lincomb}(n, (\lambda j \in \operatorname{interval}(1, n).\; (a \cdot b(j))), g) \end{array}}}
\vnbprop{cc-lincomb-zero-coefficients}{\vnbstmt{\begin{array}{@{}l@{}} \forall n, g. \\ \quad ((n \in \mathbb{N}) \wedge (g \in (\operatorname{interval}(1, n) \to \mathbb{C}))) \Rightarrow  \\ \quad \quad \operatorname{cc\text{-}lincomb}(n, (\lambda j \in \operatorname{interval}(1, n).\; 0), g) = 0 \end{array}}}
\vnbprop{cc-lincomb-unit}{\vnbstmt{\begin{array}{@{}l@{}} \forall n, g, i. \\ \quad ((n \in \mathbb{N}) \wedge ((g \in (\operatorname{interval}(1, n) \to \mathbb{C})) \wedge (i \in \operatorname{interval}(1, n)))) \Rightarrow  \\ \quad \quad \operatorname{cc\text{-}lincomb}(n, (\lambda j \in \operatorname{interval}(1, n).\; \left(\text{if }(j = i)\text{ then }1\text{ else }0\right)), g) = g(i) \end{array}}}
\vnbprop{cc-power-sum-as-lincomb}{\vnbstmt{\begin{array}{@{}l@{}} \forall k \in \mathbb{N}, b, c, r. \\ \quad ((b \in \mathbb{N}) \wedge ((c \in (\mathbb{N} \to \mathbb{C})) \wedge ((r \in \mathbb{C}) \wedge (k \leq b)))) \Rightarrow  \\ \quad \quad \operatorname{finsum}(\operatorname{ring\text{-}additive\text{-}ag}(\operatorname{normed\text{-}field\text{-}as\text{-}commutative\text{-}ring}(\operatorname{cc\text{-}normed\text{-}field})), \\ \quad \quad \quad (\lambda j \in \operatorname{interval}(1, b).\; \\ \quad \quad \quad \quad (\lambda i \in \operatorname{interval}(1, b).\; c(\operatorname{nn\text{-}minus}(i, 1)))(j) \cdot (\lambda i \in \operatorname{interval}(1, b).\; \operatorname{power}(r, \operatorname{nn\text{-}minus}(i, 1)))(j)), \\ \quad \quad \quad \operatorname{interval}(1, k)) \\ \quad \quad =\; \operatorname{cc\text{-}series\text{-}partial\text{-}sum}((\lambda \mathit{pfj} \in \mathbb{N}.\; (c(\mathit{pfj}) \cdot \operatorname{power}(r, \mathit{pfj}))), k) \end{array}}}
\vnbprop{invariant-space-of-relation}{\vnbstmt{\begin{array}{@{}l@{}} \forall a, r, g, c. \\ \quad \operatorname{is\text{-}subring\text{-}of\text{-}cc}(a) \wedge \\ \quad r \in \mathbb{C} \wedge \\ \quad g \in \mathbb{N} \wedge \\ \quad c \in (\mathbb{N} \to a) \wedge \\ \quad \forall k \in \mathbb{N}.\; (((g + 1) \leq k) \Rightarrow (c(k) = 0)) \wedge \\ \quad c(g) = 1 \wedge \\ \quad \operatorname{cc\text{-}series\text{-}partial\text{-}sum}((\lambda j \in \mathbb{N}.\; (c(j) \cdot \operatorname{power}(r, j))), (g + 1)) = 0 \Rightarrow \\ \quad \quad \exists v. \\ \quad \quad \quad (\operatorname{is\text{-}cc\text{-}space\text{-}dim\text{-}le}(a, v, g) \wedge (\exists z \in v.\; \neg (z = 0) \wedge \forall \mathit{amrv} \in v.\; ((r \cdot \mathit{amrv}) \in v))) \end{array}}}
\vnbprop{invariant-module-of-algebraic}{\vnbstmt{\begin{array}{@{}l@{}} \forall a, r. \\ \quad (\operatorname{is\text{-}subring\text{-}of\text{-}cc}(a) \wedge \operatorname{is\text{-}algebraic\text{-}over}(a, r)) \Rightarrow  \\ \quad \quad \exists v. \\ \quad \quad \quad (\operatorname{is\text{-}fg\text{-}cc\text{-}module\text{-}over}(a, v) \wedge (\exists z \in v.\; \neg (z = 0) \wedge \forall \mathit{amrv} \in v.\; ((r \cdot \mathit{amrv}) \in v))) \end{array}}}
\vnbprop{algebraic-iff-invariant-module}{\vnbstmt{\begin{array}{@{}l@{}} \forall a, r. \\ \quad (\operatorname{is\text{-}subring\text{-}of\text{-}cc}(a) \wedge (r \in \mathbb{C})) \Rightarrow  \\ \quad \quad \operatorname{is\text{-}algebraic\text{-}over}(a, r) \\ \quad \quad \Leftrightarrow\; \exists v. \\ \quad \quad \quad (\operatorname{is\text{-}fg\text{-}cc\text{-}module\text{-}over}(a, v) \wedge (\exists z \in v.\; \neg (z = 0) \wedge \forall \mathit{amrv} \in v.\; ((r \cdot \mathit{amrv}) \in v))) \end{array}}}
\vnbprop{cc-span-module}{\vnbstmt{\begin{array}{@{}l@{}} \forall a, b, u. \\ \quad (\operatorname{is\text{-}subring\text{-}of\text{-}cc}(a) \wedge ((b \in \mathbb{N}) \wedge (u \in (\operatorname{interval}(1, b) \to \mathbb{C})))) \Rightarrow  \\ \quad \quad \operatorname{is\text{-}fg\text{-}cc\text{-}module\text{-}over}(a, \\ \quad \quad \quad \{v \in \mathbb{C} : \exists \mathit{amb} \in (\operatorname{interval}(1, b) \to a).\; (v = \operatorname{cc\text{-}lincomb}(b, \mathit{amb}, u))\}) \wedge \\ \quad \quad u \\ \quad \quad \in\; (\operatorname{interval}(1, b) \to \{v \in \mathbb{C} : \exists \mathit{amb} \in (\operatorname{interval}(1, b) \to a).\; (v = \operatorname{cc\text{-}lincomb}(b, \mathit{amb}, u))\}) \end{array}}}
\vnbprop{cc-span-invariant}{\vnbstmt{\begin{array}{@{}l@{}} \forall a, b, u, o.; \forall v \in \{\mathit{amvv} \in \mathbb{C} : \exists \mathit{amb} \in (\operatorname{interval}(1, b) \to a).\; (\mathit{amvv} = \operatorname{cc\text{-}lincomb}(b, \mathit{amb}, u))\}. \\ \quad \operatorname{is\text{-}subring\text{-}of\text{-}cc}(a) \wedge \\ \quad b \in \mathbb{N} \wedge \\ \quad u \in (\operatorname{interval}(1, b) \to \mathbb{C}) \wedge \\ \quad o \in \mathbb{C} \wedge \\ \quad \forall t \in \operatorname{interval}(1, b). \\ \quad \quad (o \cdot u(t)) \in \{\mathit{amvv} \in \mathbb{C} : \exists \mathit{amb} \in (\operatorname{interval}(1, b) \to a).\; (\mathit{amvv} = \operatorname{cc\text{-}lincomb}(b, \mathit{amb}, u))\} \Rightarrow \\ \quad \quad (o \cdot v) \in \{\mathit{amvv} \in \mathbb{C} : \exists \mathit{amb} \in (\operatorname{interval}(1, b) \to a).\; (\mathit{amvv} = \operatorname{cc\text{-}lincomb}(b, \mathit{amb}, u))\} \end{array}}}
\vnbprop{product-generators-exist}{\vnbstmt{\begin{array}{@{}l@{}} \forall k \in \mathbb{N}, a, n, m, g, h. \\ \quad \operatorname{is\text{-}subring\text{-}of\text{-}cc}(a) \wedge \\ \quad n \in \mathbb{N} \wedge \\ \quad m \in \mathbb{N} \wedge \\ \quad g \in (\operatorname{interval}(1, n) \to \mathbb{C}) \wedge \\ \quad h \in (\operatorname{interval}(1, m) \to \mathbb{C}) \wedge \\ \quad k \leq n \Rightarrow \\ \quad \quad \exists \mathit{apn} \in \mathbb{N}, u \in (\operatorname{interval}(1, \mathit{apn}) \to \mathbb{C}). \\ \quad \quad \quad \forall t \in \operatorname{interval}(1, \mathit{apn}). \exists i \in \operatorname{interval}(1, k), j \in \operatorname{interval}(1, m). \\ \quad \quad \quad \quad u(t) = (g(i) \cdot h(j)) \wedge \\ \quad \quad \quad \forall i \in \operatorname{interval}(1, k), j \in \operatorname{interval}(1, m). \exists b \in (\operatorname{interval}(1, \mathit{apn}) \to a). \\ \quad \quad \quad \quad (g(i) \cdot h(j)) = \operatorname{cc\text{-}lincomb}(\mathit{apn}, b, u) \end{array}}}
\vnbprop{lincomb-in-module}{\vnbstmt{\begin{array}{@{}l@{}} \forall a, v, n, b, w. \\ \quad \operatorname{is\text{-}cc\text{-}module\text{-}over}(a, v) \wedge \\ \quad a \subseteq \mathbb{C} \wedge \\ \quad n \in \mathbb{N} \wedge \\ \quad b \in (\operatorname{interval}(1, n) \to a) \wedge \\ \quad w \in (\operatorname{interval}(1, n) \to v) \Rightarrow \\ \quad \quad \operatorname{cc\text{-}lincomb}(n, b, w) \in v \end{array}}}
\vnbprop{module-preimage-mul}{\vnbstmt{\begin{array}{@{}l@{}} \forall a, v, c. \\ \quad (\operatorname{is\text{-}cc\text{-}module\text{-}over}(a, v) \wedge ((a \subseteq \mathbb{C}) \wedge (c \in \mathbb{C}))) \Rightarrow  \\ \quad \quad \operatorname{is\text{-}cc\text{-}module\text{-}over}(a, \{q \in \mathbb{C} : ((q \cdot c) \in v)\}) \end{array}}}
\vnbprop{fg-module-nonzero-generator}{\vnbstmt{\begin{array}{@{}l@{}} \forall a, v. \\ \quad (\operatorname{is\text{-}subring\text{-}of\text{-}cc}(a) \wedge (\operatorname{is\text{-}fg\text{-}cc\text{-}module\text{-}over}(a, v) \wedge \exists z \in v.\; \neg (z = 0))) \Rightarrow  \\ \quad \quad \exists n \in \mathbb{N}, g \in (\operatorname{interval}(1, n) \to v). \\ \quad \quad \quad \forall \mathit{amv} \in v. \exists b \in (\operatorname{interval}(1, n) \to a). \\ \quad \quad \quad \quad \mathit{amv} = \operatorname{cc\text{-}lincomb}(n, b, g) \wedge \\ \quad \quad \quad \exists k \in \operatorname{interval}(1, n).\; \neg (g(k) = 0) \end{array}}}
\vnbprop{module-product-invariance}{\vnbstmt{\begin{array}{@{}l@{}} \forall a, r, v, n, g, m, h, w.; \forall i \in \operatorname{interval}(1, n), j \in \operatorname{interval}(1, m). \\ \quad \operatorname{is\text{-}subring\text{-}of\text{-}cc}(a) \wedge \\ \quad r \in \mathbb{C} \wedge \\ \quad n \in \mathbb{N} \wedge \\ \quad m \in \mathbb{N} \wedge \\ \quad g \in (\operatorname{interval}(1, n) \to v) \wedge \\ \quad v \subseteq \mathbb{C} \wedge \\ \quad \forall \mathit{amv} \in v. \exists b \in (\operatorname{interval}(1, n) \to a). \\ \quad \quad \mathit{amv} = \operatorname{cc\text{-}lincomb}(n, b, g) \wedge \\ \quad \forall \mathit{amrv} \in v.\; ((r \cdot \mathit{amrv}) \in v) \wedge \\ \quad h \in (\operatorname{interval}(1, m) \to \mathbb{C}) \wedge \\ \quad \operatorname{is\text{-}cc\text{-}module\text{-}over}(a, w) \wedge \\ \quad \forall l \in \operatorname{interval}(1, n), \mathit{apm} \in \operatorname{interval}(1, m).\; ((g(l) \cdot h(\mathit{apm})) \in w) \Rightarrow \\ \quad \quad (r \cdot (g(i) \cdot h(j))) \in w \end{array}}}
\vnbprop{algebraic-add-mul}{\vnbstmt{\begin{array}{@{}l@{}} \forall a, r, s. \\ \quad (\operatorname{is\text{-}subring\text{-}of\text{-}cc}(a) \wedge (\operatorname{is\text{-}algebraic\text{-}over}(a, r) \wedge \operatorname{is\text{-}algebraic\text{-}over}(a, s))) \Rightarrow  \\ \quad \quad \operatorname{is\text{-}algebraic\text{-}over}(a, (r + s)) \wedge \operatorname{is\text{-}algebraic\text{-}over}(a, (r \cdot s)) \end{array}}}
\vnbprop{algebraic-neg}{\vnbstmt{\begin{array}{@{}l@{}} \forall a, r. \\ \quad (\operatorname{is\text{-}subring\text{-}of\text{-}cc}(a) \wedge \operatorname{is\text{-}algebraic\text{-}over}(a, r)) \Rightarrow \operatorname{is\text{-}algebraic\text{-}over}(a, -r) \end{array}}}
\vnbprop{algebraic-of-ring-element}{\vnbstmt{\forall b, a.\; ((\operatorname{is\text{-}subring\text{-}of\text{-}cc}(b) \wedge (a \in b)) \Rightarrow \operatorname{is\text{-}algebraic\text{-}over}(b, a))}}
\vnbprop{algebraic-numbers-subring}{\vnbstmt{\begin{array}{@{}l@{}} \forall a. \\ \quad \operatorname{is\text{-}subring\text{-}of\text{-}cc}(a) \Rightarrow \operatorname{is\text{-}subring\text{-}of\text{-}cc}(\{q \in \mathbb{C} : \operatorname{is\text{-}algebraic\text{-}over}(a, q)\}) \end{array}}}
\vnbprop{cc-fraction-field-sub}{\vnbstmt{\begin{array}{@{}l@{}} \forall a. \\ \quad \operatorname{cc\text{-}fraction\text{-}field}(a) \subseteq \{q \in \mathbb{C} : \exists \mathit{anfa}, b \in a.\; (\neg (b = 0) \wedge ((q \cdot b) = \mathit{anfa}))\} \end{array}}}
\vnbprop{cc-fraction-common-denominator}{\vnbstmt{\begin{array}{@{}l@{}} \forall k \in \mathbb{N}, a, c. \\ \quad \operatorname{is\text{-}subring\text{-}of\text{-}cc}(a) \wedge \\ \quad c \in (\mathbb{N} \to \{q \in \mathbb{C} : \exists \mathit{anfa}, b \in a.\; (\neg (b = 0) \wedge ((q \cdot b) = \mathit{anfa}))\}) \Rightarrow \\ \quad \quad \exists d \in a.\; (\neg (d = 0) \wedge \forall j \in \mathbb{N}.\; ((j < k) \Rightarrow ((d \cdot c(j)) \in a))) \end{array}}}
\vnbprop{algebraic-scale-leading}{\vnbstmt{\begin{array}{@{}l@{}} \forall b, p, a, h, n. \\ \quad \operatorname{is\text{-}subring\text{-}of\text{-}cc}(b) \wedge \\ \quad p \in \mathbb{N} \wedge \\ \quad a \in (\mathbb{N} \to b) \wedge \\ \quad \forall k \in \mathbb{N}.\; ((((p + 1) + 1) \leq k) \Rightarrow (a(k) = 0)) \wedge \\ \quad h \in \mathbb{C} \wedge \\ \quad a((p + 1)) = n \wedge \\ \quad \neg (n = 0) \wedge \\ \quad \operatorname{cc\text{-}series\text{-}partial\text{-}sum}((\lambda j \in \mathbb{N}.\; (a(j) \cdot \operatorname{power}(h, j))), ((p + 1) + 1)) = 0 \Rightarrow \\ \quad \quad \operatorname{is\text{-}algebraic\text{-}over}(b, (n \cdot h)) \end{array}}}
\vnbprop{algebraic-over-fractions-scaled}{\vnbstmt{\begin{array}{@{}l@{}} \forall a, h. \\ \quad (\operatorname{is\text{-}subring\text{-}of\text{-}cc}(a) \wedge \operatorname{is\text{-}algebraic\text{-}over\text{-}field}(\operatorname{cc\text{-}fraction\text{-}field}(a), h)) \Rightarrow  \\ \quad \quad \exists k \in a.\; (\neg (k = 0) \wedge \operatorname{is\text{-}algebraic\text{-}over}(a, (k \cdot h))) \end{array}}}
\vnbprop{zz-is-subring-of-cc}{\vnbstmt{\operatorname{is\text{-}subring\text{-}of\text{-}cc}(\mathbb{Z})}}
\vnbprop{qq-is-subring-of-cc}{\vnbstmt{\operatorname{is\text{-}subring\text{-}of\text{-}cc}(\mathbb{Q})}}
\vnbprop{qq-is-subfield-of-cc}{\vnbstmt{\operatorname{is\text{-}subfield\text{-}of\text{-}cc}(\mathbb{Q})}}
\vnbprop{algebraic-over-mono}{\vnbstmt{\begin{array}{@{}l@{}} \forall a, b, x. \\ \quad ((a \subseteq b) \wedge \operatorname{is\text{-}algebraic\text{-}over}(a, x)) \Rightarrow \operatorname{is\text{-}algebraic\text{-}over}(b, x) \end{array}}}
\vnbprop{algebraic-over-field-mono}{\vnbstmt{\begin{array}{@{}l@{}} \forall a, b, x. \\ \quad ((a \subseteq b) \wedge \operatorname{is\text{-}algebraic\text{-}over\text{-}field}(a, x)) \Rightarrow  \\ \quad \quad \operatorname{is\text{-}algebraic\text{-}over\text{-}field}(b, x) \end{array}}}
\vnbprop{algebraic-ring-to-field}{\vnbstmt{\forall a, x.\; (\operatorname{is\text{-}algebraic\text{-}over}(a, x) \Rightarrow \operatorname{is\text{-}algebraic\text{-}over\text{-}field}(a, x))}}
\vnbprop{qq-in-zz-fractions}{\vnbstmt{\mathbb{Q} \subseteq \{q \in \mathbb{C} : \exists a, b \in \mathbb{Z}.\; (\neg (b = 0) \wedge ((q \cdot b) = a))\}}}
\vnbprop{cc-ps-shift}{\vnbstmt{\begin{array}{@{}l@{}} \forall n, c, z. \\ \quad ((n \in \mathbb{N}) \wedge ((c \in (\mathbb{N} \to \mathbb{C})) \wedge (z \in \mathbb{C}))) \Rightarrow  \\ \quad \quad \operatorname{cc\text{-}series\text{-}partial\text{-}sum}((\lambda j \in \mathbb{N}.\; (c(j) \cdot \operatorname{power}(z, j))), (n + 1)) \\ \quad \quad =\; c(0) + (z \cdot \operatorname{cc\text{-}series\text{-}partial\text{-}sum}((\lambda j \in \mathbb{N}.\; ((\lambda k \in \mathbb{N}.\; c((k + 1)))(j) \cdot \operatorname{power}(z, j))), n)) \end{array}}}
\vnbprop{recip-algebraic}{\vnbstmt{\begin{array}{@{}l@{}} \forall n \in \mathbb{N}, k, c, b. \\ \quad \operatorname{is\text{-}subfield\text{-}of\text{-}cc}(k) \wedge \\ \quad c \in (\mathbb{N} \to k) \wedge \\ \quad \forall \mathit{pfk} \in \mathbb{N}.\; (((n + 1) \leq \mathit{pfk}) \Rightarrow (c(\mathit{pfk}) = 0)) \wedge \\ \quad c(n) = 1 \wedge \\ \quad b \in \mathbb{C} \wedge \\ \quad \neg (b = 0) \wedge \\ \quad \operatorname{cc\text{-}series\text{-}partial\text{-}sum}((\lambda j \in \mathbb{N}.\; (c(j) \cdot \operatorname{power}(b, j))), (n + 1)) = 0 \Rightarrow \\ \quad \quad \operatorname{is\text{-}algebraic\text{-}over}(k, \operatorname{recip}(b)) \end{array}}}
\vnbprop{cc-fraction-field-sup}{\vnbstmt{\begin{array}{@{}l@{}} \forall a. \\ \quad \{q \in \mathbb{C} : \exists \mathit{anfa}, b \in a.\; (\neg (b = 0) \wedge ((q \cdot b) = \mathit{anfa}))\} \subseteq \operatorname{cc\text{-}fraction\text{-}field}(a) \end{array}}}
\vnbprop{algebraic-field-to-monic}{\vnbstmt{\begin{array}{@{}l@{}} \forall k, x. \\ \quad (\operatorname{is\text{-}subfield\text{-}of\text{-}cc}(k) \wedge \operatorname{is\text{-}algebraic\text{-}over\text{-}field}(k, x)) \Rightarrow  \\ \quad \quad \operatorname{is\text{-}algebraic\text{-}over}(k, x) \end{array}}}
\vnbprop{algebraic-numbers-over-subfield-field}{\vnbstmt{\begin{array}{@{}l@{}} \forall k. \\ \quad \operatorname{is\text{-}subfield\text{-}of\text{-}cc}(k) \Rightarrow \operatorname{is\text{-}subfield\text{-}of\text{-}cc}(\operatorname{algebraic\text{-}numbers\text{-}over}(k)) \end{array}}}
\vnbprop{algebraic-over-qq-iff-fraction-of-algebraic-integers}{\vnbstmt{\begin{array}{@{}l@{}} \forall h \in \mathbb{C}. \\ \quad \operatorname{is\text{-}algebraic\text{-}over\text{-}field}(\mathbb{Q}, h) \Leftrightarrow (h \in \operatorname{cc\text{-}fraction\text{-}field}(\{q \in \mathbb{C} : \operatorname{is\text{-}algebraic\text{-}over}(\mathbb{Z}, q)\})) \end{array}}}
\vnbprop{algebraic-degree-of-invariant-space}{\vnbstmt{\begin{array}{@{}l@{}} \forall a, r, v, m. \\ \quad \operatorname{is\text{-}subring\text{-}of\text{-}cc}(a) \wedge \\ \quad r \in \mathbb{C} \wedge \\ \quad m \in \mathbb{N} \wedge \\ \quad \operatorname{is\text{-}cc\text{-}space\text{-}dim\text{-}le}(a, v, m) \wedge \\ \quad \exists z \in v.\; \neg (z = 0) \wedge \\ \quad \forall \mathit{amrv} \in v.\; ((r \cdot \mathit{amrv}) \in v) \Rightarrow \\ \quad \quad \operatorname{is\text{-}algebraic\text{-}of\text{-}degree\text{-}le}(a, r, m) \end{array}}}
\vnbprop{monic-of-leading}{\vnbstmt{\begin{array}{@{}l@{}} \forall k, m, c, x. \\ \quad \operatorname{is\text{-}subfield\text{-}of\text{-}cc}(k) \wedge \\ \quad m \in \mathbb{N} \wedge \\ \quad c \in (\mathbb{N} \to k) \wedge \\ \quad \forall \mathit{pfk} \in \mathbb{N}.\; (((m + 1) \leq \mathit{pfk}) \Rightarrow (c(\mathit{pfk}) = 0)) \wedge \\ \quad \neg (c(m) = 0) \wedge \\ \quad x \in \mathbb{C} \wedge \\ \quad \operatorname{cc\text{-}series\text{-}partial\text{-}sum}((\lambda j \in \mathbb{N}.\; (c(j) \cdot \operatorname{power}(x, j))), (m + 1)) = 0 \Rightarrow \\ \quad \quad \exists \mathit{anmc} \in (\mathbb{N} \to k). \\ \quad \quad \quad \forall \mathit{pfk} \in \mathbb{N}.\; (((m + 1) \leq \mathit{pfk}) \Rightarrow (\operatorname{anmc}(\mathit{pfk}) = 0)) \wedge \\ \quad \quad \quad \operatorname{anmc}(m) = 1 \wedge \\ \quad \quad \quad \operatorname{cc\text{-}series\text{-}partial\text{-}sum}((\lambda j \in \mathbb{N}.\; (\operatorname{anmc}(j) \cdot \operatorname{power}(x, j))), (m + 1)) = 0 \end{array}}}
\vnbprop{cc-space-dim-le-mono}{\vnbstmt{\begin{array}{@{}l@{}} \forall k, v, m, n. \\ \quad ((m \in \mathbb{N}) \wedge ((n \in \mathbb{N}) \wedge ((m \leq n) \wedge \operatorname{is\text{-}cc\text{-}space\text{-}dim\text{-}le}(k, v, m)))) \Rightarrow  \\ \quad \quad \operatorname{is\text{-}cc\text{-}space\text{-}dim\text{-}le}(k, v, n) \end{array}}}
\vnbprop{algebraic-degree-iff-invariant-space}{\vnbstmt{\begin{array}{@{}l@{}} \forall k, h, n. \\ \quad (\operatorname{is\text{-}subfield\text{-}of\text{-}cc}(k) \wedge ((h \in \mathbb{C}) \wedge (n \in \mathbb{N}))) \Rightarrow  \\ \quad \quad \operatorname{is\text{-}algebraic\text{-}of\text{-}degree\text{-}le}(k, h, n) \\ \quad \quad \Leftrightarrow\; \exists v. \\ \quad \quad \quad (\operatorname{is\text{-}cc\text{-}space\text{-}dim\text{-}le}(k, v, n) \wedge (\exists z \in v.\; \neg (z = 0) \wedge \forall \mathit{amrv} \in v.\; ((h \cdot \mathit{amrv}) \in v))) \end{array}}}
\vnbprop{partial-sum-in-module}{\vnbstmt{\begin{array}{@{}l@{}} \forall k \in \mathbb{N}, a, v, f. \\ \quad (\operatorname{is\text{-}cc\text{-}module\text{-}over}(a, v) \wedge ((f \in (\mathbb{N} \to \mathbb{C})) \wedge \forall j \in \mathbb{N}.\; ((j < k) \Rightarrow (f(j) \in v)))) \Rightarrow  \\ \quad \quad \operatorname{cc\text{-}series\text{-}partial\text{-}sum}(f, k) \in v \end{array}}}
\vnbprop{algebraic-tower-module}{\vnbstmt{\begin{array}{@{}l@{}} \forall j \in \mathbb{N}, a, c. \\ \quad \operatorname{is\text{-}subring\text{-}of\text{-}cc}(a) \wedge \\ \quad c \in (\mathbb{N} \to \mathbb{C}) \wedge \\ \quad \forall i \in \mathbb{N}.\; ((i < j) \Rightarrow \operatorname{is\text{-}algebraic\text{-}over}(a, c(i))) \Rightarrow \\ \quad \quad \exists n \in \mathbb{N}, u \in (\operatorname{interval}(1, n) \to \mathbb{C}). \\ \quad \quad \quad \exists z \in \{v \in \mathbb{C} : \exists b \in (\operatorname{interval}(1, n) \to a).\; (v = \operatorname{cc\text{-}lincomb}(n, b, u))\}. \\ \quad \quad \quad \quad \neg (z = 0) \wedge \\ \quad \quad \quad \forall i \in \mathbb{N}, t \in \operatorname{interval}(1, n). \\ \quad \quad \quad \quad (i < j) \Rightarrow  \\ \quad \quad \quad \quad \quad (c(i) \cdot u(t)) \in \{v \in \mathbb{C} : \exists b \in (\operatorname{interval}(1, n) \to a).\; (v = \operatorname{cc\text{-}lincomb}(n, b, u))\} \end{array}}}
\vnbprop{algebraic-over-algebraic}{\vnbstmt{\begin{array}{@{}l@{}} \forall a, n, c, h. \\ \quad \operatorname{is\text{-}subring\text{-}of\text{-}cc}(a) \wedge \\ \quad n \in \mathbb{N} \wedge \\ \quad c \in (\mathbb{N} \to \mathbb{C}) \wedge \\ \quad \forall i \in \mathbb{N}.\; ((i < (n + 1)) \Rightarrow \operatorname{is\text{-}algebraic\text{-}over}(a, c(i))) \wedge \\ \quad c(n) = 1 \wedge \\ \quad h \in \mathbb{C} \wedge \\ \quad \operatorname{cc\text{-}series\text{-}partial\text{-}sum}((\lambda j \in \mathbb{N}.\; (c(j) \cdot \operatorname{power}(h, j))), (n + 1)) = 0 \Rightarrow \\ \quad \quad \operatorname{is\text{-}algebraic\text{-}over}(a, h) \end{array}}}
\vnbprop{algebraic-root-in-closure}{\vnbstmt{\begin{array}{@{}l@{}} \forall k, c, m, z. \\ \quad \operatorname{is\text{-}subfield\text{-}of\text{-}cc}(k) \wedge \\ \quad c \in (\mathbb{N} \to \operatorname{algebraic\text{-}numbers\text{-}over}(k)) \wedge \\ \quad m \in \mathbb{N} \wedge \\ \quad \neg (c(m) = 0) \wedge \\ \quad z \in \mathbb{C} \wedge \\ \quad \operatorname{cc\text{-}series\text{-}partial\text{-}sum}((\lambda j \in \mathbb{N}.\; (c(j) \cdot \operatorname{power}(z, j))), (m + 1)) = 0 \Rightarrow \\ \quad \quad z \in \operatorname{algebraic\text{-}numbers\text{-}over}(k) \end{array}}}
\vnbprop{algebraic-numbers-algebraically-closed}{\vnbstmt{\begin{array}{@{}l@{}} \forall k, c, m. \\ \quad (\operatorname{is\text{-}subfield\text{-}of\text{-}cc}(k) \wedge ((c \in (\mathbb{N} \to \operatorname{algebraic\text{-}numbers\text{-}over}(k))) \wedge ((m \in \mathbb{N}) \wedge \neg (c((m + 1)) = 0)))) \Rightarrow  \\ \quad \quad \exists z \in \operatorname{algebraic\text{-}numbers\text{-}over}(k). \\ \quad \quad \quad \operatorname{cc\text{-}series\text{-}partial\text{-}sum}((\lambda j \in \mathbb{N}.\; (c(j) \cdot \operatorname{power}(z, j))), ((m + 1) + 1)) = 0 \end{array}}}
\vnbprop{cc-partial-prod-zero-factor}{\vnbstmt{\begin{array}{@{}l@{}} \forall n \in \mathbb{N}, g.; \forall q \in \mathbb{N}. \\ \quad (\forall k \in \mathbb{N}.\; (g(k) \in \mathbb{C}) \wedge ((q < n) \wedge (g(q) = 0))) \Rightarrow  \\ \quad \quad \operatorname{cc\text{-}partial\text{-}prod}(g, n) = 0 \end{array}}}
\vnbprop{algebraic-numbers-factorization}{\vnbstmt{\begin{array}{@{}l@{}} \forall k, c, n. \\ \quad \operatorname{is\text{-}subfield\text{-}of\text{-}cc}(k) \wedge \\ \quad c \in (\mathbb{N} \to \operatorname{algebraic\text{-}numbers\text{-}over}(k)) \wedge \\ \quad n \in \mathbb{N} \wedge \\ \quad \forall \mathit{pfk} \in \mathbb{N}.\; (((n + 1) \leq \mathit{pfk}) \Rightarrow (c(\mathit{pfk}) = 0)) \wedge \\ \quad \neg (c(n) = 0) \Rightarrow \\ \quad \quad \exists r \in (\mathbb{N} \to \mathbb{C}). \\ \quad \quad \quad \forall i \in \mathbb{N}.\; ((i < n) \Rightarrow (r(i) \in \operatorname{algebraic\text{-}numbers\text{-}over}(k))) \wedge \\ \quad \quad \quad \forall z \in \mathbb{C}. \\ \quad \quad \quad \quad \operatorname{cc\text{-}series\text{-}partial\text{-}sum}((\lambda j \in \mathbb{N}.\; (c(j) \cdot \operatorname{power}(z, j))), (n + 1)) \\ \quad \quad \quad \quad =\; c(n) \cdot \operatorname{cc\text{-}partial\text{-}prod}((\lambda \mathit{cfi} \in \mathbb{N}.\; (z - r(\mathit{cfi}))), n) \end{array}}}

\vnbbookfile{theorem-library/ca-lft-topology}{theorem-library/ca-lft-topology}
\noindent{\small ca-lft-topology.scm -- Lemma 1.16 of complex-analysis.pdf: convergence in the topology of the extended plane C u \{oo\} is Euclidean convergence at a finite point, and |x\_n| -$>$ +oo at oo. CA-LFT, 2026-10-19.  Library: complex-analysis.  Vocabulary: structure-library/ca-lft-vocab.scm.}\par
\vnbprop{cc-compact-bounded}{\vnbstmt{\begin{array}{@{}l@{}} \forall k. \\ \quad ((k \subseteq \mathbb{C}) \wedge \operatorname{is\text{-}compact}(\operatorname{subspace\text{-}ms}(\operatorname{cc\text{-}ms}, k))) \Rightarrow  \\ \quad \quad \exists n \in \mathbb{N}.\; \forall z \in k.\; (\operatorname{magnitude}(z) \leq n) \end{array}}}
\vnbprop{cc-compact-separated}{\vnbstmt{\begin{array}{@{}l@{}} \forall k.; \forall x \in \mathbb{C}. \\ \quad ((k \subseteq \mathbb{C}) \wedge (\operatorname{is\text{-}compact}(\operatorname{subspace\text{-}ms}(\operatorname{cc\text{-}ms}, k)) \wedge \neg (x \in k))) \Rightarrow  \\ \quad \quad \exists r. \\ \quad \quad \quad \operatorname{pos\text{-}rr}(r) \wedge \forall y \in \mathbb{C}.\; ((\operatorname{magnitude}((y - x)) < r) \Rightarrow \neg (y \in k)) \end{array}}}
\vnbprop{cc-hat-opens-membership}{\vnbstmt{\begin{array}{@{}l@{}} \forall u. \\ \quad u \in \operatorname{cc\text{-}hat\text{-}opens} \\ \quad \Leftrightarrow\; u \in \operatorname{power}(\operatorname{cc\text{-}hat}) \wedge \\ \quad (u \subseteq \mathbb{C}) \wedge \operatorname{is\text{-}open}(\operatorname{cc\text{-}ms}, u) \vee \\ \quad \exists h. \\ \quad \quad (h \subseteq \mathbb{C}) \wedge (\operatorname{is\text{-}compact}(\operatorname{subspace\text{-}ms}(\operatorname{cc\text{-}ms}, h)) \wedge (u = \operatorname{difference}(\operatorname{cc\text{-}hat}, h))) \end{array}}}
\vnbprop{cc-hat-ball-open}{\vnbstmt{\forall x \in \mathbb{C}, r.\; (\operatorname{pos\text{-}rr}(r) \Rightarrow (\operatorname{ball}(\operatorname{cc\text{-}ms}, x, r) \in \operatorname{cc\text{-}hat\text{-}opens}))}}
\vnbprop{cc-hat-ball-complement-open}{\vnbstmt{\forall m \in \mathbb{R}.\; (\operatorname{difference}(\operatorname{cc\text{-}hat}, \operatorname{closed\text{-}ball}(\operatorname{cc\text{-}ms}, 0, m)) \in \operatorname{cc\text{-}hat\text{-}opens})}}
\vnbprop{cc-hat-nbhd-finite}{\vnbstmt{\begin{array}{@{}l@{}} \forall u \in \operatorname{cc\text{-}hat\text{-}opens}, x \in \mathbb{C}. \\ \quad (x \in u) \Rightarrow  \\ \quad \quad \exists r. \\ \quad \quad \quad \operatorname{pos\text{-}rr}(r) \wedge \forall y \in \mathbb{C}.\; ((\operatorname{magnitude}((y - x)) < r) \Rightarrow (y \in u)) \end{array}}}
\vnbprop{cc-hat-nbhd-inf}{\vnbstmt{\begin{array}{@{}l@{}} \forall u \in \operatorname{cc\text{-}hat\text{-}opens}. \\ \quad (\operatorname{cc\text{-}inf} \in u) \Rightarrow  \\ \quad \quad \exists m \in \mathbb{R}.\; \forall y \in \mathbb{C}.\; ((m < \operatorname{magnitude}(y)) \Rightarrow (y \in u)) \end{array}}}
\vnbprop{cc-hat-converges-iff}{\vnbstmt{\begin{array}{@{}l@{}} \forall f \in (\mathbb{N} \to \mathbb{C}), x \in \mathbb{C}. \\ \quad \operatorname{top\text{-}converges\text{-}to}(\operatorname{cc\text{-}hat\text{-}top}, f, x) \Leftrightarrow \operatorname{converges\text{-}to}(\operatorname{cc\text{-}ms}, f, x) \end{array}}}
\vnbprop{cc-hat-converges-inf-iff}{\vnbstmt{\begin{array}{@{}l@{}} \forall f \in (\mathbb{N} \to \mathbb{C}). \\ \quad \operatorname{top\text{-}converges\text{-}to}(\operatorname{cc\text{-}hat\text{-}top}, f, \operatorname{cc\text{-}inf}) \\ \quad \Leftrightarrow\; \forall m \in \mathbb{R}. \exists n \in \mathbb{N}. \forall k \in \mathbb{N}. \\ \quad \quad (n \leq k) \Rightarrow (m \leq \operatorname{magnitude}(f(k))) \end{array}}}

\vnbbookfile{theorem-library/ca-lft-cayley}{theorem-library/ca-lft-cayley}
\noindent{\small ca-lft-cayley.scm -- Prop. 1.21 of complex-analysis.pdf: the Cayley transform C z = (z - i)/(z + i) maps R bijectively onto T \textbackslash{} \{1\} and the upper half plane H+ bijectively onto the disc D = B(0, 1). CA-LFT, 2026-10-19.  Library: complex-analysis.  Vocabulary: structure-library/ca-lft-vocab.scm.}\par
\vnbprop{cayley-mag-sq}{\vnbstmt{\begin{array}{@{}l@{}} \forall w \in \mathbb{C}. \\ \quad (\operatorname{magnitude}((w + +i)) \cdot \operatorname{magnitude}((w + +i))) = ((\operatorname{magnitude}((w - +i)) \cdot \operatorname{magnitude}((w - +i))) + (4 \cdot \operatorname{imag\text{-}part}(w))) \end{array}}}
\vnbprop{cayley-at}{\vnbstmt{\forall w \in \mathbb{C}.\; (\neg ((w + +i) = 0) \Rightarrow (\operatorname{cayley}(w) = ((w - +i) \cdot \operatorname{recip}((w + +i)))))}}
\vnbprop{cayley-mag}{\vnbstmt{\begin{array}{@{}l@{}} \forall w \in \mathbb{C}. \\ \quad \neg ((w + +i) = 0) \Rightarrow  \\ \quad \quad (\operatorname{magnitude}(\operatorname{cayley}(w)) \cdot \operatorname{magnitude}((w + +i))) = \operatorname{magnitude}((w - +i)) \end{array}}}
\vnbprop{cayley-disc-iff}{\vnbstmt{\begin{array}{@{}l@{}} \forall w \in \mathbb{C}. \\ \quad \neg ((w + +i) = 0) \Rightarrow ((\operatorname{magnitude}(\operatorname{cayley}(w)) < 1) \Leftrightarrow (0 < \operatorname{imag\text{-}part}(w))) \end{array}}}
\vnbprop{cayley-circle-iff}{\vnbstmt{\begin{array}{@{}l@{}} \forall w \in \mathbb{C}. \\ \quad \neg ((w + +i) = 0) \Rightarrow ((\operatorname{magnitude}(\operatorname{cayley}(w)) = 1) \Leftrightarrow (\operatorname{imag\text{-}part}(w) = 0)) \end{array}}}
\vnbprop{cc-im-zero}{\vnbstmt{\operatorname{imag\text{-}part}(0) = 0}}
\vnbprop{cc-magnitude-one}{\vnbstmt{\operatorname{magnitude}(1) = 1}}
\vnbprop{cayley-no-pole}{\vnbstmt{\forall w \in \mathbb{C}.\; ((0 \leq \operatorname{imag\text{-}part}(w)) \Rightarrow \neg ((w + +i) = 0))}}
\vnbprop{unit-circle-membership}{\vnbstmt{\forall u.\; ((u \in \operatorname{unit\text{-}circle}) \Leftrightarrow ((u \in \mathbb{C}) \wedge (\operatorname{magnitude}(u) = 1)))}}
\vnbprop{upper-half-plane-membership}{\vnbstmt{\forall z.\; ((z \in \operatorname{upper\text{-}half\text{-}plane}) \Leftrightarrow ((z \in \mathbb{C}) \wedge (0 < \operatorname{imag\text{-}part}(z))))}}
\vnbprop{unit-disc-elim}{\vnbstmt{\forall u \in \operatorname{ball}(\operatorname{cc\text{-}ms}, 0, 1).\; ((u \in \mathbb{C}) \wedge (\operatorname{magnitude}(u) < 1))}}
\vnbprop{unit-disc-intro}{\vnbstmt{\forall u \in \mathbb{C}.\; ((\operatorname{magnitude}(u) < 1) \Rightarrow (u \in \operatorname{ball}(\operatorname{cc\text{-}ms}, 0, 1)))}}
\vnbprop{cayley-in-fun}{\vnbstmt{\mathit{cayley} \in (\operatorname{cc\text{-}hat} \to \operatorname{cc\text{-}hat})}}
\vnbprop{cayley-inv-left}{\vnbstmt{\operatorname{compose}(\operatorname{lft}(+i, +i, -1, 1), \mathit{cayley}) = \operatorname{id\text{-}fun}(\operatorname{cc\text{-}hat})}}
\vnbprop{cayley-inv-right}{\vnbstmt{\operatorname{compose}(\mathit{cayley}, \operatorname{lft}(+i, +i, -1, 1)) = \operatorname{id\text{-}fun}(\operatorname{cc\text{-}hat})}}
\vnbprop{cayley-inv-apply-left}{\vnbstmt{\forall z \in \operatorname{cc\text{-}hat}.\; (\operatorname{lft}(+i, +i, -1, 1)(\operatorname{cayley}(z)) = z)}}
\vnbprop{cayley-inv-apply-right}{\vnbstmt{\forall u \in \operatorname{cc\text{-}hat}.\; (\operatorname{cayley}(\operatorname{lft}(+i, +i, -1, 1)(u)) = u)}}
\vnbprop{cayley-inv-at}{\vnbstmt{\begin{array}{@{}l@{}} \forall u \in \mathbb{C}. \\ \quad \neg (u = 1) \Rightarrow  \\ \quad \quad ((\operatorname{lft}(+i, +i, -1, 1)(u) = (((+i \cdot u) + +i) \cdot \operatorname{recip}(((-1 \cdot u) + 1)))) \wedge (\operatorname{lft}(+i, +i, -1, 1)(u) \in \mathbb{C})) \end{array}}}
\vnbprop{cayley-inv-pole}{\vnbstmt{\operatorname{lft}(+i, +i, -1, 1)(1) = \operatorname{cc\text{-}inf}}}
\vnbprop{cayley-pole}{\vnbstmt{\operatorname{cayley}(-+i) = \operatorname{cc\text{-}inf}}}
\vnbprop{cayley-maps-real}{\vnbstmt{\forall x \in \mathbb{R}.\; (\operatorname{cayley}(x) \in \operatorname{difference}(\operatorname{unit\text{-}circle}, \operatorname{singleton}(1)))}}
\vnbprop{cayley-inv-maps-circle}{\vnbstmt{\begin{array}{@{}l@{}} \forall u \in \operatorname{difference}(\operatorname{unit\text{-}circle}, \operatorname{singleton}(1)). \\ \quad \operatorname{lft}(+i, +i, -1, 1)(u) \in \mathbb{R} \end{array}}}
\vnbprop{cayley-maps-upper}{\vnbstmt{\forall z \in \operatorname{upper\text{-}half\text{-}plane}.\; (\operatorname{cayley}(z) \in \operatorname{ball}(\operatorname{cc\text{-}ms}, 0, 1))}}
\vnbprop{cayley-inv-maps-disc}{\vnbstmt{\forall u \in \operatorname{ball}(\operatorname{cc\text{-}ms}, 0, 1).\; (\operatorname{lft}(+i, +i, -1, 1)(u) \in \operatorname{upper\text{-}half\text{-}plane})}}
\vnbprop{cayley-real-circle}{\vnbstmt{(\lambda r \in \mathbb{R}.\; \operatorname{cayley}(r)) \in \mathrm{Bij}(\mathbb{R}, \operatorname{difference}(\operatorname{unit\text{-}circle}, \operatorname{singleton}(1)))}}
\vnbprop{cayley-upper-disc}{\vnbstmt{(\lambda r \in \operatorname{upper\text{-}half\text{-}plane}.\; \operatorname{cayley}(r)) \in \mathrm{Bij}(\operatorname{upper\text{-}half\text{-}plane}, \operatorname{ball}(\operatorname{cc\text{-}ms}, 0, 1))}}

\vnbbookfile{theorem-library/triangle-orientation}{theorem-library/triangle-orientation}
\noindent{\small triangle-orientation.scm -- WAVE 2 (2026-10-27..30), wired 2026-10-30 from scratchpad/ch3e/ch3e-angle.scm (proofs), scratchpad/ch3e/ch3e-orient.scm. Library: complex-analysis.}\par
\vnbprop{real-scale-re-im}{\vnbstmt{\begin{array}{@{}l@{}} \forall c, w. \\ \quad ((c \in \mathbb{R}) \wedge (w \in \mathbb{C})) \Rightarrow  \\ \quad \quad (\operatorname{imag\text{-}part}((c \cdot w)) = (c \cdot \operatorname{imag\text{-}part}(w))) \wedge (\operatorname{real\text{-}part}((c \cdot w)) = (c \cdot \operatorname{real\text{-}part}(w))) \end{array}}}
\vnbprop{exp-polar-re-im}{\vnbstmt{\begin{array}{@{}l@{}} \forall l \in \mathbb{C}. \\ \quad \operatorname{imag\text{-}part}(\exp l) = (\exp \operatorname{real\text{-}part}(l) \cdot \sin \operatorname{imag\text{-}part}(l)) \wedge \\ \quad \operatorname{real\text{-}part}(\exp l) = (\exp \operatorname{real\text{-}part}(l) \cdot \cos \operatorname{imag\text{-}part}(l)) \end{array}}}
\vnbprop{angle-lt-pi-iff}{\vnbstmt{\begin{array}{@{}l@{}} \forall z, v. \\ \quad ((z \in \mathbb{C}) \wedge ((v \in \mathbb{C}) \wedge (\neg (z = 0) \wedge \neg (v = 0)))) \Rightarrow  \\ \quad \quad \operatorname{angle}(z, v) < \pi \\ \quad \quad \Leftrightarrow\; 0 < \operatorname{imag\text{-}part}((z \cdot \operatorname{conjugate}(v))) \vee \\ \quad \quad (\operatorname{imag\text{-}part}((z \cdot \operatorname{conjugate}(v))) = 0) \wedge (0 < \operatorname{real\text{-}part}((z \cdot \operatorname{conjugate}(v)))) \end{array}}}
\vnbprop{conv3-cross-weights}{\vnbstmt{\begin{array}{@{}l@{}} \forall a, b, c.; \forall z \in \operatorname{conv3}(a, b, c). \\ \quad ((a \in \mathbb{C}) \wedge ((b \in \mathbb{C}) \wedge (c \in \mathbb{C}))) \Rightarrow  \\ \quad \quad \exists t \in \mathbb{R}. \\ \quad \quad \quad ((0 \leq t) \wedge (\operatorname{imag\text{-}part}(((b - z) \cdot \operatorname{conjugate}((a - z)))) = (t \cdot \operatorname{imag\text{-}part}(((b - c) \cdot \operatorname{conjugate}((a - c))))))) \wedge \\ \quad \quad \exists t \in \mathbb{R}. \\ \quad \quad \quad ((0 \leq t) \wedge (\operatorname{imag\text{-}part}(((c - z) \cdot \operatorname{conjugate}((b - z)))) = (t \cdot \operatorname{imag\text{-}part}(((c - a) \cdot \operatorname{conjugate}((b - a))))))) \wedge \\ \quad \quad \exists t \in \mathbb{R}. \\ \quad \quad \quad ((0 \leq t) \wedge (\operatorname{imag\text{-}part}(((a - z) \cdot \operatorname{conjugate}((c - z)))) = (t \cdot \operatorname{imag\text{-}part}(((a - b) \cdot \operatorname{conjugate}((c - b))))))) \end{array}}}
\vnbprop{interior-cross-sign}{\vnbstmt{\begin{array}{@{}l@{}} \forall p, q, k, s, u. \\ \quad p \in \mathbb{C} \wedge \\ \quad q \in \mathbb{C} \wedge \\ \quad k \in \mathbb{R} \wedge \\ \quad \neg (p = q) \wedge \\ \quad u \in \operatorname{interior}(\operatorname{cc\text{-}ms}, s) \wedge \\ \quad \forall z \in s. \exists t \in \mathbb{R}. \\ \quad \quad (0 \leq t) \wedge (\operatorname{imag\text{-}part}(((q - z) \cdot \operatorname{conjugate}((p - z)))) = (t \cdot k)) \Rightarrow \\ \quad \quad 0 < (k \cdot \operatorname{imag\text{-}part}(((q - u) \cdot \operatorname{conjugate}((p - u))))) \end{array}}}
\vnbprop{interior-angle-sign}{\vnbstmt{\begin{array}{@{}l@{}} \forall p, q, k, s.; \forall u \in \operatorname{interior}(\operatorname{cc\text{-}ms}, s). \\ \quad p \in \mathbb{C} \wedge \\ \quad q \in \mathbb{C} \wedge \\ \quad k \in \mathbb{R} \wedge \\ \quad \forall z \in s. \exists t \in \mathbb{R}. \\ \quad \quad (0 \leq t) \wedge (\operatorname{imag\text{-}part}(((q - z) \cdot \operatorname{conjugate}((p - z)))) = (t \cdot k)) \wedge \\ \quad \neg (p = q) \Rightarrow \\ \quad \quad (\operatorname{angle}((q - u), (p - u)) < \pi) \Leftrightarrow (0 < k) \end{array}}}
\vnbprop{positively-oriented-iff}{\vnbstmt{\begin{array}{@{}l@{}} \forall a, b, c. \\ \quad a \in \mathbb{C} \wedge \\ \quad b \in \mathbb{C} \wedge \\ \quad c \in \mathbb{C} \wedge \\ \quad \neg (a = b) \wedge \\ \quad \neg (b = c) \wedge \\ \quad \neg (c = a) \wedge \\ \quad \exists v.\; (v \in \operatorname{interior}(\operatorname{cc\text{-}ms}, \operatorname{conv3}(a, b, c))) \Rightarrow \\ \quad \quad \operatorname{positively\text{-}oriented}(a, b, c) \\ \quad \quad \Leftrightarrow\; \exists u \in \operatorname{interior}(\operatorname{cc\text{-}ms}, \operatorname{conv3}(a, b, c)). \\ \quad \quad \quad ((\operatorname{angle}((b - u), (a - u)) < \pi) \wedge ((\operatorname{angle}((c - u), (b - u)) < \pi) \wedge (\operatorname{angle}((a - u), (c - u)) < \pi))) \end{array}}}

\vnbbookfile{theorem-library/triangle-winding}{theorem-library/triangle-winding}
\noindent{\small triangle-winding.scm -- WAVE 2 (2026-10-27..30), wired 2026-10-30 from scratchpad/ch3e/ch3e-wind.scm, scratchpad/ch3e/ch3e-wext.scm. Library: complex-analysis.}\par
\vnbprop{cc-half-plane-open}{\vnbstmt{\begin{array}{@{}l@{}} \forall u, \mathit{aouu}. \\ \quad ((u \in \mathbb{C}) \wedge (\mathit{aouu} \in \mathbb{C})) \Rightarrow  \\ \quad \quad \operatorname{is\text{-}open}(\operatorname{nf\text{-}metric\text{-}space}(\operatorname{cc\text{-}normed\text{-}field}), \{w \in \mathbb{C} : (0 < \operatorname{real\text{-}part}((u \cdot (w - \mathit{aouu}))))\}) \end{array}}}
\vnbprop{cc-dot-coords}{\vnbstmt{\begin{array}{@{}l@{}} \forall z, w. \\ \quad ((z \in \mathbb{C}) \wedge (w \in \mathbb{C})) \Rightarrow  \\ \quad \quad \operatorname{real\text{-}part}((\operatorname{conjugate}(z) \cdot w)) = ((\operatorname{real\text{-}part}(z) \cdot \operatorname{real\text{-}part}(w)) + (\operatorname{imag\text{-}part}(z) \cdot \operatorname{imag\text{-}part}(w))) \end{array}}}
\vnbprop{cc-line-dot-cross}{\vnbstmt{\begin{array}{@{}l@{}} \forall p, q, u, s. \\ \quad ((p \in \mathbb{C}) \wedge ((q \in \mathbb{C}) \wedge ((u \in \mathbb{C}) \wedge (s \in \mathbb{R})))) \Rightarrow  \\ \quad \quad \operatorname{real\text{-}part}((\operatorname{conjugate}((+i \cdot (p - q))) \cdot ((p + (s \cdot (q - p))) - u))) = \operatorname{imag\text{-}part}(((q - u) \cdot \operatorname{conjugate}((p - u)))) \end{array}}}
\vnbprop{seg-int-log}{\vnbstmt{\begin{array}{@{}l@{}} \forall p, q, u. \\ \quad ((p \in \mathbb{C}) \wedge ((q \in \mathbb{C}) \wedge ((u \in \mathbb{C}) \wedge (0 < \operatorname{imag\text{-}part}(((q - u) \cdot \operatorname{conjugate}((p - u)))))))) \Rightarrow  \\ \quad \quad \operatorname{seg\text{-}int}((\lambda w \in \operatorname{difference}(\mathbb{C}, \operatorname{singleton}(u)).\; \operatorname{recip}((w - u))), p, q) \\ \quad \quad =\; (\log\left((\operatorname{conjugate}((+i \cdot (p - q))) \cdot (q - u))\right) - \log\left((\operatorname{conjugate}((+i \cdot (p - q))) \cdot (p - u))\right)) \end{array}}}
\vnbprop{seg-log-step}{\vnbstmt{\begin{array}{@{}l@{}} \forall p, q, u. \\ \quad ((p \in \mathbb{C}) \wedge ((q \in \mathbb{C}) \wedge ((u \in \mathbb{C}) \wedge (0 < \operatorname{imag\text{-}part}(((q - u) \cdot \operatorname{conjugate}((p - u)))))))) \Rightarrow  \\ \quad \quad \exp (\log\left((\operatorname{conjugate}((+i \cdot (p - q))) \cdot (q - u))\right) - \log\left((\operatorname{conjugate}((+i \cdot (p - q))) \cdot (p - u))\right)) \cdot (p - u) \\ \quad \quad =\; (q - u) \wedge \\ \quad \quad 0 < \operatorname{imag\text{-}part}((\log\left((\operatorname{conjugate}((+i \cdot (p - q))) \cdot (q - u))\right) - \log\left((\operatorname{conjugate}((+i \cdot (p - q))) \cdot (p - u))\right))) \wedge \\ \quad \quad \operatorname{imag\text{-}part}((\log\left((\operatorname{conjugate}((+i \cdot (p - q))) \cdot (q - u))\right) - \log\left((\operatorname{conjugate}((+i \cdot (p - q))) \cdot (p - u))\right))) \\ \quad \quad <\; \pi \end{array}}}
\vnbprop{seg-trace-punctured}{\vnbstmt{\begin{array}{@{}l@{}} \forall p, q, u. \\ \quad ((p \in \mathbb{C}) \wedge ((q \in \mathbb{C}) \wedge ((u \in \mathbb{C}) \wedge (0 < \operatorname{imag\text{-}part}(((q - u) \cdot \operatorname{conjugate}((p - u)))))))) \Rightarrow  \\ \quad \quad \operatorname{trace}(\operatorname{seg\text{-}path}(p, q), 0, 1) \subseteq \operatorname{difference}(\mathbb{C}, \operatorname{singleton}(u)) \end{array}}}
\vnbprop{triangle-crosses-positive}{\vnbstmt{\begin{array}{@{}l@{}} \forall a, b, c.; \forall u \in \operatorname{interior}(\operatorname{cc\text{-}ms}, \operatorname{conv3}(a, b, c)). \\ \quad a \in \mathbb{C} \wedge \\ \quad b \in \mathbb{C} \wedge \\ \quad c \in \mathbb{C} \wedge \\ \quad \neg (a = b) \wedge \\ \quad \neg (b = c) \wedge \\ \quad \neg (c = a) \wedge \\ \quad \operatorname{positively\text{-}oriented}(a, b, c) \Rightarrow \\ \quad \quad 0 < \operatorname{imag\text{-}part}(((b - u) \cdot \operatorname{conjugate}((a - u)))) \wedge \\ \quad \quad 0 < \operatorname{imag\text{-}part}(((c - u) \cdot \operatorname{conjugate}((b - u)))) \wedge \\ \quad \quad 0 < \operatorname{imag\text{-}part}(((a - u) \cdot \operatorname{conjugate}((c - u)))) \end{array}}}
\vnbprop{triangle-tri-int-interior}{\vnbstmt{\begin{array}{@{}l@{}} \forall a, b, c.; \forall u \in \operatorname{interior}(\operatorname{cc\text{-}ms}, \operatorname{conv3}(a, b, c)). \\ \quad a \in \mathbb{C} \wedge \\ \quad b \in \mathbb{C} \wedge \\ \quad c \in \mathbb{C} \wedge \\ \quad \neg (a = b) \wedge \\ \quad \neg (b = c) \wedge \\ \quad \neg (c = a) \wedge \\ \quad \operatorname{positively\text{-}oriented}(a, b, c) \Rightarrow \\ \quad \quad \operatorname{tri\text{-}int}((\lambda w \in \operatorname{difference}(\mathbb{C}, \operatorname{singleton}(u)).\; \operatorname{recip}((w - u))), a, b, c) = ((2 \cdot \pi) \cdot +i) \end{array}}}
\vnbprop{triangle-winding-interior}{\vnbstmt{\begin{array}{@{}l@{}} \forall a, b, c.; \forall u \in \operatorname{interior}(\operatorname{cc\text{-}ms}, \operatorname{conv3}(a, b, c)). \\ \quad a \in \mathbb{C} \wedge \\ \quad b \in \mathbb{C} \wedge \\ \quad c \in \mathbb{C} \wedge \\ \quad \neg (a = b) \wedge \\ \quad \neg (b = c) \wedge \\ \quad \neg (c = a) \wedge \\ \quad \operatorname{positively\text{-}oriented}(a, b, c) \Rightarrow \\ \quad \quad \operatorname{winding}(\operatorname{tri\text{-}road}(a, b, c), \operatorname{tri\text{-}road\text{-}deriv}(a, b, c), 0, 3, u) = 1 \end{array}}}
\vnbprop{cross-half-plane}{\vnbstmt{\begin{array}{@{}l@{}} \forall p, q, u, w. \\ \quad ((p \in \mathbb{C}) \wedge ((q \in \mathbb{C}) \wedge ((u \in \mathbb{C}) \wedge (w \in \mathbb{C})))) \Rightarrow  \\ \quad \quad \operatorname{real\text{-}part}((\operatorname{conjugate}((+i \cdot (q - p))) \cdot (w - (1/2 \cdot (p + u))))) \\ \quad \quad =\; (\operatorname{imag\text{-}part}(((q - w) \cdot \operatorname{conjugate}((p - w)))) - (1/2 \cdot \operatorname{imag\text{-}part}(((q - u) \cdot \operatorname{conjugate}((p - u)))))) \end{array}}}
\vnbprop{cross-barycentric}{\vnbstmt{\begin{array}{@{}l@{}} \forall a, b, c, z. \\ \quad ((a \in \mathbb{C}) \wedge ((b \in \mathbb{C}) \wedge ((c \in \mathbb{C}) \wedge (z \in \mathbb{C})))) \Rightarrow  \\ \quad \quad (\operatorname{imag\text{-}part}(((b - z) \cdot \operatorname{conjugate}((a - z)))) + \operatorname{imag\text{-}part}(((c - z) \cdot \operatorname{conjugate}((b - z))))) + \operatorname{imag\text{-}part}(((a - z) \cdot \operatorname{conjugate}((c - z)))) \\ \quad \quad =\; \operatorname{imag\text{-}part}(((b - c) \cdot \operatorname{conjugate}((a - c)))) \wedge \\ \quad \quad \operatorname{imag\text{-}part}(((b - c) \cdot \operatorname{conjugate}((a - c)))) \cdot z \\ \quad \quad =\; ((\operatorname{imag\text{-}part}(((c - z) \cdot \operatorname{conjugate}((b - z)))) \cdot a) + (\operatorname{imag\text{-}part}(((a - z) \cdot \operatorname{conjugate}((c - z)))) \cdot b)) + (\operatorname{imag\text{-}part}(((b - z) \cdot \operatorname{conjugate}((a - z)))) \cdot c) \end{array}}}
\vnbprop{triangle-outside-side}{\vnbstmt{\begin{array}{@{}l@{}} \forall a, b, c, u, p, q. \\ \quad a \in \mathbb{C} \wedge \\ \quad b \in \mathbb{C} \wedge \\ \quad c \in \mathbb{C} \wedge \\ \quad u \in \mathbb{C} \wedge \\ \quad p \in \mathbb{C} \wedge \\ \quad q \in \mathbb{C} \wedge \\ \quad \forall z \in \operatorname{conv3}(a, b, c).\; (0 \leq \operatorname{imag\text{-}part}(((q - z) \cdot \operatorname{conjugate}((p - z))))) \wedge \\ \quad \operatorname{imag\text{-}part}(((q - u) \cdot \operatorname{conjugate}((p - u)))) < 0 \Rightarrow \\ \quad \quad \operatorname{tri\text{-}int}((\lambda w \in \operatorname{difference}(\mathbb{C}, \operatorname{singleton}(u)).\; \operatorname{recip}((w - u))), a, b, c) = 0 \end{array}}}
\vnbprop{triangle-winding-exterior}{\vnbstmt{\begin{array}{@{}l@{}} \forall a, b, c.; \forall u \in \mathbb{C}. \\ \quad a \in \mathbb{C} \wedge \\ \quad b \in \mathbb{C} \wedge \\ \quad c \in \mathbb{C} \wedge \\ \quad \neg (a = b) \wedge \\ \quad \neg (b = c) \wedge \\ \quad \neg (c = a) \wedge \\ \quad \operatorname{positively\text{-}oriented}(a, b, c) \wedge \\ \quad \exists v.\; (v \in \operatorname{interior}(\operatorname{cc\text{-}ms}, \operatorname{conv3}(a, b, c))) \wedge \\ \quad \neg (u \in \operatorname{conv3}(a, b, c)) \Rightarrow \\ \quad \quad \operatorname{winding}(\operatorname{tri\text{-}road}(a, b, c), \operatorname{tri\text{-}road\text{-}deriv}(a, b, c), 0, 3, u) = 0 \end{array}}}

\vnbbookfile{theorem-library/triangle-cauchy-formula}{theorem-library/triangle-cauchy-formula}
\noindent{\small triangle-cauchy-formula.scm -- WAVE 2 (2026-10-27..30), wired 2026-10-30 from scratchpad/ch3e/ch3e-cif.scm, scratchpad/ch3e/ch3e-cif2.scm. Library: complex-analysis.}\par
\vnbprop{triangle-cauchy-formula}{\vnbstmt{\begin{array}{@{}l@{}} \forall n, f.; \forall a, b, c \in \mathbb{C}, u \in \operatorname{interior}(\operatorname{cc\text{-}ms}, \operatorname{conv3}(a, b, c)). \\ \quad \operatorname{holomorphic\text{-}on}(n, f) \wedge \\ \quad \neg (a = b) \wedge \\ \quad \neg (b = c) \wedge \\ \quad \neg (c = a) \wedge \\ \quad \operatorname{positively\text{-}oriented}(a, b, c) \wedge \\ \quad \operatorname{conv3}(a, b, c) \subseteq n \Rightarrow \\ \quad \quad f(u) \\ \quad \quad =\; \operatorname{recip}(((2 \cdot \pi) \cdot +i)) \cdot \operatorname{tri\text{-}int}((\lambda w \in \operatorname{difference}(n, \operatorname{singleton}(u)).\; (f(w) \cdot \operatorname{recip}((w - u)))), a, b, c) \end{array}}}
\vnbprop{triangle-cauchy-deriv}{\vnbstmt{\begin{array}{@{}l@{}} \forall n, f.; \forall a, b, c \in \mathbb{C}, u \in \operatorname{interior}(\operatorname{cc\text{-}ms}, \operatorname{conv3}(a, b, c)). \\ \quad \operatorname{holomorphic\text{-}on}(n, f) \wedge \\ \quad \neg (a = b) \wedge \\ \quad \neg (b = c) \wedge \\ \quad \neg (c = a) \wedge \\ \quad \operatorname{positively\text{-}oriented}(a, b, c) \wedge \\ \quad \operatorname{conv3}(a, b, c) \subseteq n \Rightarrow \\ \quad \quad \operatorname{deriv\text{-}on}(\operatorname{cc\text{-}normed\text{-}field}, n, f, u) \\ \quad \quad =\; \operatorname{recip}(((2 \cdot \pi) \cdot +i)) \cdot \operatorname{tri\text{-}int}((\lambda w \in \operatorname{difference}(n, \operatorname{singleton}(u)).\; (f(w) \cdot \operatorname{power}(\operatorname{recip}((w - u)), 2))), a, b, c) \end{array}}}

\vnbbookfile{theorem-library/holomorphy-across-line}{theorem-library/holomorphy-across-line}
\noindent{\small holomorphy-across-line.scm -- WAVE 2 (2026-10-27..30), wired 2026-10-30 from scratchpad/ch3e/ch3f-line.scm (proofs). Library: complex-analysis.}\par
\vnbprop{cc-uniform-continuity-ball}{\vnbstmt{\begin{array}{@{}l@{}} \forall u.; \forall f \in (u \to \mathbb{C}), c \in \mathbb{C}, r \in \mathbb{R}, s. \\ \quad \operatorname{is\text{-}open}(\operatorname{nf\text{-}metric\text{-}space}(\operatorname{cc\text{-}normed\text{-}field}), u) \wedge \\ \quad \forall y \in u. \\ \quad \quad \operatorname{is\text{-}continuous\text{-}at}(\operatorname{subspace\text{-}ms}(\operatorname{nf\text{-}metric\text{-}space}(\operatorname{cc\text{-}normed\text{-}field}), u), \\ \quad \quad \quad \operatorname{nf\text{-}metric\text{-}space}(\operatorname{cc\text{-}normed\text{-}field}), f, y) \wedge \\ \quad 0 < r \wedge \\ \quad \operatorname{closed\text{-}ball}(\operatorname{cc\text{-}ms}, c, r) \subseteq u \wedge \\ \quad \operatorname{pos\text{-}rr}(s) \Rightarrow \\ \quad \quad \exists l. \\ \quad \quad \quad \operatorname{pos\text{-}rr}(l) \wedge \\ \quad \quad \quad \forall x, \mathit{lny} \in \operatorname{closed\text{-}ball}(\operatorname{cc\text{-}ms}, c, r). \\ \quad \quad \quad \quad (\operatorname{magnitude}((x - \mathit{lny})) < l) \Rightarrow (\operatorname{magnitude}((f(x) - f(\mathit{lny}))) \leq s) \end{array}}}
\vnbprop{ln-restrict-continuous}{\vnbstmt{\begin{array}{@{}l@{}} \forall u.; \forall f \in (u \to \mathbb{C}), w.; \forall q \in w. \\ \quad \operatorname{is\text{-}open}(\operatorname{nf\text{-}metric\text{-}space}(\operatorname{cc\text{-}normed\text{-}field}), u) \wedge \\ \quad \forall y \in u. \\ \quad \quad \operatorname{is\text{-}continuous\text{-}at}(\operatorname{subspace\text{-}ms}(\operatorname{nf\text{-}metric\text{-}space}(\operatorname{cc\text{-}normed\text{-}field}), u), \\ \quad \quad \quad \operatorname{nf\text{-}metric\text{-}space}(\operatorname{cc\text{-}normed\text{-}field}), f, y) \wedge \\ \quad \operatorname{is\text{-}open}(\operatorname{nf\text{-}metric\text{-}space}(\operatorname{cc\text{-}normed\text{-}field}), w) \wedge \\ \quad w \subseteq u \Rightarrow \\ \quad \quad \operatorname{is\text{-}continuous\text{-}at}(\operatorname{subspace\text{-}ms}(\operatorname{nf\text{-}metric\text{-}space}(\operatorname{cc\text{-}normed\text{-}field}), w), \\ \quad \quad \quad \operatorname{nf\text{-}metric\text{-}space}(\operatorname{cc\text{-}normed\text{-}field}), \operatorname{restrict}(f, w), q) \end{array}}}
\vnbprop{cc-line-unfold}{\vnbstmt{\begin{array}{@{}l@{}} \forall p, d. \\ \quad (\operatorname{cc\text{-}line}(p, d) \simeq \{z \in \mathbb{C} : (\operatorname{imag\text{-}part}(((z - p) \cdot \operatorname{conjugate}(d))) = 0)\}) \end{array}}}
\vnbprop{ln-zero-of-small}{\vnbstmt{\begin{array}{@{}l@{}} \forall z, p. \\ \quad ((z \in \mathbb{C}) \wedge ((p \in \mathbb{R}) \wedge ((0 \leq p) \wedge \forall e.\; (\operatorname{pos\text{-}rr}(e) \Rightarrow (\operatorname{magnitude}(z) \leq (e \cdot p)))))) \Rightarrow  \\ \quad \quad z = 0 \end{array}}}
\vnbprop{ln-psi-conv}{\vnbstmt{\begin{array}{@{}l@{}} \forall a, b, c, p_{0}, d, s.; \forall z \in \operatorname{conv3}(a, b, c). \\ \quad a \in \mathbb{C} \wedge \\ \quad b \in \mathbb{C} \wedge \\ \quad c \in \mathbb{C} \wedge \\ \quad p_{0} \in \mathbb{C} \wedge \\ \quad d \in \mathbb{C} \wedge \\ \quad s \in \mathbb{R} \wedge \\ \quad 0 \leq (s \cdot \operatorname{imag\text{-}part}(((a - p_{0}) \cdot \operatorname{conjugate}(d)))) \wedge \\ \quad 0 \leq (s \cdot \operatorname{imag\text{-}part}(((b - p_{0}) \cdot \operatorname{conjugate}(d)))) \wedge \\ \quad 0 \leq (s \cdot \operatorname{imag\text{-}part}(((c - p_{0}) \cdot \operatorname{conjugate}(d)))) \Rightarrow \\ \quad \quad 0 \leq (s \cdot \operatorname{imag\text{-}part}(((z - p_{0}) \cdot \operatorname{conjugate}(d)))) \end{array}}}
\vnbprop{ln-shift-identities}{\vnbstmt{\begin{array}{@{}l@{}} \forall w, p_{0}, d, s, a. \\ \quad ((w \in \mathbb{C}) \wedge ((p_{0} \in \mathbb{C}) \wedge ((d \in \mathbb{C}) \wedge ((s \in \mathbb{R}) \wedge (a \in \mathbb{R}))))) \Rightarrow  \\ \quad \quad \operatorname{real\text{-}part}((\operatorname{conjugate}((+i \cdot (s \cdot d))) \cdot (w - (p_{0} - ((1/2 \cdot a) \cdot (+i \cdot (s \cdot d))))))) \\ \quad \quad =\; (s \cdot \operatorname{imag\text{-}part}(((w - p_{0}) \cdot \operatorname{conjugate}(d)))) + ((1/2 \cdot a) \cdot ((s \cdot s) \cdot (\operatorname{magnitude}(d) \cdot \operatorname{magnitude}(d)))) \wedge \\ \quad \quad \operatorname{imag\text{-}part}((((w + (a \cdot (+i \cdot (s \cdot d)))) - p_{0}) \cdot \operatorname{conjugate}(d))) \\ \quad \quad =\; \operatorname{imag\text{-}part}(((w - p_{0}) \cdot \operatorname{conjugate}(d))) + (a \cdot (s \cdot (\operatorname{magnitude}(d) \cdot \operatorname{magnitude}(d)))) \end{array}}}
\vnbprop{tri-int-half-plane}{\vnbstmt{\begin{array}{@{}l@{}} \forall u.; \forall f \in (u \to \mathbb{C}), p_{0}, d, c \in \mathbb{C}, r, s \in \mathbb{R}, a, b, \mathit{aoc} \in \operatorname{ball}(\operatorname{cc\text{-}ms}, c, r). \\ \quad \operatorname{is\text{-}open}(\operatorname{nf\text{-}metric\text{-}space}(\operatorname{cc\text{-}normed\text{-}field}), u) \wedge \\ \quad \forall y \in u. \\ \quad \quad \operatorname{is\text{-}continuous\text{-}at}(\operatorname{subspace\text{-}ms}(\operatorname{nf\text{-}metric\text{-}space}(\operatorname{cc\text{-}normed\text{-}field}), u), \\ \quad \quad \quad \operatorname{nf\text{-}metric\text{-}space}(\operatorname{cc\text{-}normed\text{-}field}), f, y) \wedge \\ \quad \neg (d = 0) \wedge \\ \quad \operatorname{holomorphic\text{-}on}(\operatorname{difference}(u, \operatorname{cc\text{-}line}(p_{0}, d)), \operatorname{restrict}(f, \operatorname{difference}(u, \operatorname{cc\text{-}line}(p_{0}, d)))) \wedge \\ \quad 0 < r \wedge \\ \quad \operatorname{closed\text{-}ball}(\operatorname{cc\text{-}ms}, c, (2 \cdot r)) \subseteq u \wedge \\ \quad (s \cdot s) = 1 \wedge \\ \quad \operatorname{magnitude}(s) = 1 \wedge \\ \quad 0 \leq (s \cdot \operatorname{imag\text{-}part}(((a - p_{0}) \cdot \operatorname{conjugate}(d)))) \wedge \\ \quad 0 \leq (s \cdot \operatorname{imag\text{-}part}(((b - p_{0}) \cdot \operatorname{conjugate}(d)))) \wedge \\ \quad 0 \leq (s \cdot \operatorname{imag\text{-}part}(((\mathit{aoc} - p_{0}) \cdot \operatorname{conjugate}(d)))) \Rightarrow \\ \quad \quad \operatorname{tri\text{-}int}(f, a, b, \mathit{aoc}) = 0 \end{array}}}
\vnbprop{ln-psi-seg}{\vnbstmt{\begin{array}{@{}l@{}} \forall p, q, p_{0}, d, t. \\ \quad ((p \in \mathbb{C}) \wedge ((q \in \mathbb{C}) \wedge ((p_{0} \in \mathbb{C}) \wedge ((d \in \mathbb{C}) \wedge (t \in \mathbb{R}))))) \Rightarrow  \\ \quad \quad \operatorname{imag\text{-}part}((((p + (t \cdot (q - p))) - p_{0}) \cdot \operatorname{conjugate}(d))) \\ \quad \quad =\; \operatorname{imag\text{-}part}(((p - p_{0}) \cdot \operatorname{conjugate}(d))) + (t \cdot (\operatorname{imag\text{-}part}(((q - p_{0}) \cdot \operatorname{conjugate}(d))) - \operatorname{imag\text{-}part}(((p - p_{0}) \cdot \operatorname{conjugate}(d))))) \end{array}}}
\vnbprop{ln-cut-point}{\vnbstmt{\begin{array}{@{}l@{}} \forall p, q, p_{0}, d, s. \\ \quad p \in \mathbb{C} \wedge \\ \quad q \in \mathbb{C} \wedge \\ \quad p_{0} \in \mathbb{C} \wedge \\ \quad d \in \mathbb{C} \wedge \\ \quad s \in \mathbb{R} \wedge \\ \quad (s \cdot s) = 1 \wedge \\ \quad 0 < (s \cdot \operatorname{imag\text{-}part}(((p - p_{0}) \cdot \operatorname{conjugate}(d)))) \wedge \\ \quad (s \cdot \operatorname{imag\text{-}part}(((q - p_{0}) \cdot \operatorname{conjugate}(d)))) < 0 \Rightarrow \\ \quad \quad \exists y \in \operatorname{trace}(\operatorname{seg\text{-}path}(p, q), 0, 1). \\ \quad \quad \quad \operatorname{imag\text{-}part}(((y - p_{0}) \cdot \operatorname{conjugate}(d))) = 0 \end{array}}}
\vnbprop{tri-int-line-cross-vertex}{\vnbstmt{\begin{array}{@{}l@{}} \forall u.; \forall f \in (u \to \mathbb{C}), p_{0}, d, c \in \mathbb{C}, r, s \in \mathbb{R}, p, q, x \in \operatorname{ball}(\operatorname{cc\text{-}ms}, c, r). \\ \quad \operatorname{is\text{-}open}(\operatorname{nf\text{-}metric\text{-}space}(\operatorname{cc\text{-}normed\text{-}field}), u) \wedge \\ \quad \forall y \in u. \\ \quad \quad \operatorname{is\text{-}continuous\text{-}at}(\operatorname{subspace\text{-}ms}(\operatorname{nf\text{-}metric\text{-}space}(\operatorname{cc\text{-}normed\text{-}field}), u), \\ \quad \quad \quad \operatorname{nf\text{-}metric\text{-}space}(\operatorname{cc\text{-}normed\text{-}field}), f, y) \wedge \\ \quad \neg (d = 0) \wedge \\ \quad \operatorname{holomorphic\text{-}on}(\operatorname{difference}(u, \operatorname{cc\text{-}line}(p_{0}, d)), \operatorname{restrict}(f, \operatorname{difference}(u, \operatorname{cc\text{-}line}(p_{0}, d)))) \wedge \\ \quad 0 < r \wedge \\ \quad \operatorname{closed\text{-}ball}(\operatorname{cc\text{-}ms}, c, (2 \cdot r)) \subseteq u \wedge \\ \quad (s \cdot s) = 1 \wedge \\ \quad \operatorname{magnitude}(s) = 1 \wedge \\ \quad 0 < (s \cdot \operatorname{imag\text{-}part}(((p - p_{0}) \cdot \operatorname{conjugate}(d)))) \wedge \\ \quad (s \cdot \operatorname{imag\text{-}part}(((q - p_{0}) \cdot \operatorname{conjugate}(d)))) < 0 \wedge \\ \quad \operatorname{imag\text{-}part}(((x - p_{0}) \cdot \operatorname{conjugate}(d))) = 0 \Rightarrow \\ \quad \quad \operatorname{tri\text{-}int}(f, p, q, x) = 0 \end{array}}}
\vnbprop{tri-int-line-vertex}{\vnbstmt{\begin{array}{@{}l@{}} \forall u.; \forall f \in (u \to \mathbb{C}), p_{0}, d, c \in \mathbb{C}, r \in \mathbb{R}, p, q, x \in \operatorname{ball}(\operatorname{cc\text{-}ms}, c, r). \\ \quad \operatorname{is\text{-}open}(\operatorname{nf\text{-}metric\text{-}space}(\operatorname{cc\text{-}normed\text{-}field}), u) \wedge \\ \quad \forall y \in u. \\ \quad \quad \operatorname{is\text{-}continuous\text{-}at}(\operatorname{subspace\text{-}ms}(\operatorname{nf\text{-}metric\text{-}space}(\operatorname{cc\text{-}normed\text{-}field}), u), \\ \quad \quad \quad \operatorname{nf\text{-}metric\text{-}space}(\operatorname{cc\text{-}normed\text{-}field}), f, y) \wedge \\ \quad \neg (d = 0) \wedge \\ \quad \operatorname{holomorphic\text{-}on}(\operatorname{difference}(u, \operatorname{cc\text{-}line}(p_{0}, d)), \operatorname{restrict}(f, \operatorname{difference}(u, \operatorname{cc\text{-}line}(p_{0}, d)))) \wedge \\ \quad 0 < r \wedge \\ \quad \operatorname{closed\text{-}ball}(\operatorname{cc\text{-}ms}, c, (2 \cdot r)) \subseteq u \wedge \\ \quad \operatorname{imag\text{-}part}(((x - p_{0}) \cdot \operatorname{conjugate}(d))) = 0 \Rightarrow \\ \quad \quad \operatorname{tri\text{-}int}(f, p, q, x) = 0 \end{array}}}
\vnbprop{tri-int-line-ball}{\vnbstmt{\begin{array}{@{}l@{}} \forall u.; \forall f \in (u \to \mathbb{C}), p_{0}, d, c \in \mathbb{C}, r \in \mathbb{R}, a, b, \mathit{aoc} \in \operatorname{ball}(\operatorname{cc\text{-}ms}, c, r). \\ \quad \operatorname{is\text{-}open}(\operatorname{nf\text{-}metric\text{-}space}(\operatorname{cc\text{-}normed\text{-}field}), u) \wedge \\ \quad \forall y \in u. \\ \quad \quad \operatorname{is\text{-}continuous\text{-}at}(\operatorname{subspace\text{-}ms}(\operatorname{nf\text{-}metric\text{-}space}(\operatorname{cc\text{-}normed\text{-}field}), u), \\ \quad \quad \quad \operatorname{nf\text{-}metric\text{-}space}(\operatorname{cc\text{-}normed\text{-}field}), f, y) \wedge \\ \quad \neg (d = 0) \wedge \\ \quad \operatorname{holomorphic\text{-}on}(\operatorname{difference}(u, \operatorname{cc\text{-}line}(p_{0}, d)), \operatorname{restrict}(f, \operatorname{difference}(u, \operatorname{cc\text{-}line}(p_{0}, d)))) \wedge \\ \quad 0 < r \wedge \\ \quad \operatorname{closed\text{-}ball}(\operatorname{cc\text{-}ms}, c, (2 \cdot r)) \subseteq u \Rightarrow \\ \quad \quad \operatorname{tri\text{-}int}(f, a, b, \mathit{aoc}) = 0 \end{array}}}
\vnbprop{holomorphic-across-line}{\vnbstmt{\begin{array}{@{}l@{}} \forall u.; \forall f \in (u \to \mathbb{C}), p_{0}, d \in \mathbb{C}. \\ \quad \operatorname{is\text{-}open}(\operatorname{nf\text{-}metric\text{-}space}(\operatorname{cc\text{-}normed\text{-}field}), u) \wedge \\ \quad \forall y \in u. \\ \quad \quad \operatorname{is\text{-}continuous\text{-}at}(\operatorname{subspace\text{-}ms}(\operatorname{nf\text{-}metric\text{-}space}(\operatorname{cc\text{-}normed\text{-}field}), u), \\ \quad \quad \quad \operatorname{nf\text{-}metric\text{-}space}(\operatorname{cc\text{-}normed\text{-}field}), f, y) \wedge \\ \quad \neg (d = 0) \wedge \\ \quad \operatorname{holomorphic\text{-}on}(\operatorname{difference}(u, \operatorname{cc\text{-}line}(p_{0}, d)), \operatorname{restrict}(f, \operatorname{difference}(u, \operatorname{cc\text{-}line}(p_{0}, d)))) \Rightarrow \\ \quad \quad \operatorname{holomorphic\text{-}on}(u, f) \end{array}}}

\vnbbookfile{theorem-library/schwarz-reflection}{theorem-library/schwarz-reflection}
\noindent{\small schwarz-reflection.scm -- WAVE 2 (2026-10-27..30), wired 2026-10-30 from scratchpad/ch3e/ch3f-refl.scm (proofs). Library: complex-analysis.}\par
\vnbprop{cc-cont-at-mag-sub}{\vnbstmt{\begin{array}{@{}l@{}} \forall u, g, a, s. \\ \quad u \subseteq \mathbb{C} \wedge \\ \quad g \in (u \to \mathbb{C}) \wedge \\ \quad a \in u \wedge \\ \quad \operatorname{is\text{-}continuous\text{-}at}(\operatorname{subspace\text{-}ms}(\operatorname{nf\text{-}metric\text{-}space}(\operatorname{cc\text{-}normed\text{-}field}), u), \\ \quad \quad \operatorname{nf\text{-}metric\text{-}space}(\operatorname{cc\text{-}normed\text{-}field}), g, a) \wedge \\ \quad \operatorname{pos\text{-}rr}(s) \Rightarrow \\ \quad \quad \exists l. \\ \quad \quad \quad (\operatorname{pos\text{-}rr}(l) \wedge \forall z \in u.\; ((\operatorname{magnitude}((z - a)) < l) \Rightarrow (\operatorname{magnitude}((g(z) - g(a))) < s))) \end{array}}}
\vnbprop{cc-cont-at-of-mag-sub}{\vnbstmt{\begin{array}{@{}l@{}} \forall u, g, a. \\ \quad u \subseteq \mathbb{C} \wedge \\ \quad g \in (u \to \mathbb{C}) \wedge \\ \quad a \in u \wedge \\ \quad \forall s. \\ \quad \quad \operatorname{pos\text{-}rr}(s) \Rightarrow  \\ \quad \quad \quad \exists l. \\ \quad \quad \quad \quad (\operatorname{pos\text{-}rr}(l) \wedge \forall z \in u.\; ((\operatorname{magnitude}((z - a)) < l) \Rightarrow (\operatorname{magnitude}((g(z) - g(a))) < s))) \Rightarrow \\ \quad \quad \operatorname{is\text{-}continuous\text{-}at}(\operatorname{subspace\text{-}ms}(\operatorname{nf\text{-}metric\text{-}space}(\operatorname{cc\text{-}normed\text{-}field}), u), \\ \quad \quad \quad \operatorname{nf\text{-}metric\text{-}space}(\operatorname{cc\text{-}normed\text{-}field}), g, a) \end{array}}}
\vnbprop{conj-set-unfold}{\vnbstmt{\forall a.\; (\operatorname{conj\text{-}set}(a) \simeq \{z \in \mathbb{C} : (\operatorname{conjugate}(z) \in a)\})}}
\vnbprop{conj-set-open}{\vnbstmt{\begin{array}{@{}l@{}} \forall a. \\ \quad \operatorname{is\text{-}open}(\operatorname{nf\text{-}metric\text{-}space}(\operatorname{cc\text{-}normed\text{-}field}), a) \Rightarrow  \\ \quad \quad \operatorname{is\text{-}open}(\operatorname{nf\text{-}metric\text{-}space}(\operatorname{cc\text{-}normed\text{-}field}), \operatorname{conj\text{-}set}(a)) \end{array}}}
\vnbprop{diff-on-transfer-near}{\vnbstmt{\begin{array}{@{}l@{}} \forall a, h, p, l, b.; \forall g \in (b \to \mathbb{C}). \\ \quad \operatorname{is\text{-}diff\text{-}on}(\operatorname{cc\text{-}normed\text{-}field}, a, h, p, l) \wedge \\ \quad \operatorname{is\text{-}open}(\operatorname{nf\text{-}metric\text{-}space}(\operatorname{cc\text{-}normed\text{-}field}), b) \wedge \\ \quad p \in b \wedge \\ \quad l \in \mathbb{C} \wedge \\ \quad \exists r. \\ \quad \quad (\operatorname{pos\text{-}rr}(r) \wedge \forall z \in b.\; ((\operatorname{magnitude}((z - p)) < r) \Rightarrow ((z \in a) \wedge (g(z) = h(z))))) \Rightarrow \\ \quad \quad \operatorname{is\text{-}diff\text{-}on}(\operatorname{cc\text{-}normed\text{-}field}, b, g, p, l) \end{array}}}
\vnbprop{reflection-holomorphic}{\vnbstmt{\begin{array}{@{}l@{}} \forall a, h. \\ \quad \operatorname{holomorphic\text{-}on}(a, h) \Rightarrow  \\ \quad \quad \operatorname{holomorphic\text{-}on}(\operatorname{conj\text{-}set}(a), (\lambda w \in \operatorname{conj\text{-}set}(a).\; \operatorname{conjugate}(h(\operatorname{conjugate}(w))))) \end{array}}}
\vnbprop{rx-im-conj}{\vnbstmt{\forall z \in \mathbb{C}.\; (\operatorname{imag\text{-}part}(\operatorname{conjugate}(z)) = -\operatorname{imag\text{-}part}(z))}}
\vnbprop{interior-conj-symmetric}{\vnbstmt{\begin{array}{@{}l@{}} \forall s.; \forall z \in \operatorname{interior}(\operatorname{cc\text{-}ms}, s). \\ \quad ((s \subseteq \mathbb{C}) \wedge \forall z \in s.\; (\operatorname{conjugate}(z) \in s)) \Rightarrow  \\ \quad \quad \operatorname{conjugate}(z) \in \operatorname{interior}(\operatorname{cc\text{-}ms}, s) \end{array}}}
\vnbprop{schwarz-reflection}{\vnbstmt{\begin{array}{@{}l@{}} \forall u.; \forall f \in ((u \cup (\operatorname{closure}(\operatorname{cc\text{-}ms}, u) \cap \mathbb{R})) \to \mathbb{C}). \\ \quad \operatorname{is\text{-}open}(\operatorname{nf\text{-}metric\text{-}space}(\operatorname{cc\text{-}normed\text{-}field}), u) \wedge \\ \quad \forall z \in u.\; (0 < \operatorname{imag\text{-}part}(z)) \wedge \\ \quad \forall y \in (u \cup (\operatorname{closure}(\operatorname{cc\text{-}ms}, u) \cap \mathbb{R})). \\ \quad \quad \operatorname{is\text{-}continuous\text{-}at}(\operatorname{subspace\text{-}ms}(\operatorname{nf\text{-}metric\text{-}space}(\operatorname{cc\text{-}normed\text{-}field}), (u \cup (\operatorname{closure}(\operatorname{cc\text{-}ms}, u) \cap \mathbb{R}))), \\ \quad \quad \quad \operatorname{nf\text{-}metric\text{-}space}(\operatorname{cc\text{-}normed\text{-}field}), f, y) \wedge \\ \quad \operatorname{holomorphic\text{-}on}(u, \operatorname{restrict}(f, u)) \wedge \\ \quad \forall x \in (\operatorname{closure}(\operatorname{cc\text{-}ms}, u) \cap \mathbb{R}).\; (f(x) \in \mathbb{R}) \Rightarrow \\ \quad \quad \exists g. \\ \quad \quad \quad \operatorname{holomorphic\text{-}on}(\operatorname{interior}(\operatorname{cc\text{-}ms}, \\ \quad \quad \quad \quad \quad ((u \cup (\operatorname{closure}(\operatorname{cc\text{-}ms}, u) \cap \mathbb{R})) \cup \operatorname{conj\text{-}set}((u \cup (\operatorname{closure}(\operatorname{cc\text{-}ms}, u) \cap \mathbb{R}))))), \\ \quad \quad \quad \quad g) \wedge \\ \quad \quad \quad \forall z \in \operatorname{interior}(\operatorname{cc\text{-}ms}, ((u \cup (\operatorname{closure}(\operatorname{cc\text{-}ms}, u) \cap \mathbb{R})) \cup \operatorname{conj\text{-}set}((u \cup (\operatorname{closure}(\operatorname{cc\text{-}ms}, u) \cap \mathbb{R}))))). \\ \quad \quad \quad \quad (z \in (u \cup (\operatorname{closure}(\operatorname{cc\text{-}ms}, u) \cap \mathbb{R}))) \Rightarrow (g(z) = f(z)) \wedge \\ \quad \quad \quad \forall z \in \operatorname{interior}(\operatorname{cc\text{-}ms}, ((u \cup (\operatorname{closure}(\operatorname{cc\text{-}ms}, u) \cap \mathbb{R})) \cup \operatorname{conj\text{-}set}((u \cup (\operatorname{closure}(\operatorname{cc\text{-}ms}, u) \cap \mathbb{R}))))). \\ \quad \quad \quad \quad \operatorname{conjugate}(z) \\ \quad \quad \quad \quad \in\; \operatorname{interior}(\operatorname{cc\text{-}ms}, \\ \quad \quad \quad \quad \quad ((u \cup (\operatorname{closure}(\operatorname{cc\text{-}ms}, u) \cap \mathbb{R})) \cup \operatorname{conj\text{-}set}((u \cup (\operatorname{closure}(\operatorname{cc\text{-}ms}, u) \cap \mathbb{R}))))) \wedge \\ \quad \quad \quad \quad g(\operatorname{conjugate}(z)) = \operatorname{conjugate}(g(z)) \end{array}}}

\vnbbookfile{theorem-library/gamma-continuation}{theorem-library/gamma-continuation}
\noindent{\small gamma-continuation.scm -- WAVE 2 (2026-10-27..30), wired 2026-10-30 from scratchpad/ch3g/ch3g-gamma.scm (proofs). Library: complex-analysis.}\par
\vnbprop{poch-in-cc}{\vnbstmt{\forall n \in \mathbb{N}, z \in \mathbb{C}.\; (\operatorname{poch}(z, n) \in \mathbb{C})}}
\vnbprop{poch-nonzero}{\vnbstmt{\begin{array}{@{}l@{}} \forall n \in \mathbb{N}, z \in \{\mathit{gdz} \in \mathbb{C} : \forall \mathit{gdn} \in \mathbb{N}.\; \neg ((\mathit{gdz} + \mathit{gdn}) = 0)\}. \\ \quad \neg (\operatorname{poch}(z, n) = 0) \end{array}}}
\vnbprop{poch-shift}{\vnbstmt{\forall n \in \mathbb{N}, z \in \mathbb{C}.\; (\operatorname{poch}(z, (n + 1)) = (z \cdot \operatorname{poch}((z + 1), n)))}}
\vnbprop{gamma-poch-step}{\vnbstmt{\begin{array}{@{}l@{}} \forall n \in \mathbb{N}, z \in \{\mathit{gdz} \in \mathbb{C} : \forall \mathit{gdn} \in \mathbb{N}.\; \neg ((\mathit{gdz} + \mathit{gdn}) = 0)\}. \\ \quad (0 < (\operatorname{real\text{-}part}(z) + n)) \Rightarrow  \\ \quad \quad (\Gamma((z + (n + 1))) \cdot \operatorname{recip}(\operatorname{poch}(z, (n + 1)))) = (\Gamma((z + n)) \cdot \operatorname{recip}(\operatorname{poch}(z, n))) \end{array}}}
\vnbprop{gamma-poch-gap}{\vnbstmt{\begin{array}{@{}l@{}} \forall k, n \in \mathbb{N}, z \in \{\mathit{gdz} \in \mathbb{C} : \forall \mathit{gdn} \in \mathbb{N}.\; \neg ((\mathit{gdz} + \mathit{gdn}) = 0)\}. \\ \quad (0 < (\operatorname{real\text{-}part}(z) + n)) \Rightarrow  \\ \quad \quad (\Gamma((z + (n + k))) \cdot \operatorname{recip}(\operatorname{poch}(z, (n + k)))) = (\Gamma((z + n)) \cdot \operatorname{recip}(\operatorname{poch}(z, n))) \end{array}}}
\vnbprop{gamma-ext-value}{\vnbstmt{\begin{array}{@{}l@{}} \forall n \in \mathbb{N}, z \in \{\mathit{gdz} \in \mathbb{C} : \forall \mathit{gdn} \in \mathbb{N}.\; \neg ((\mathit{gdz} + \mathit{gdn}) = 0)\}. \\ \quad (0 < (\operatorname{real\text{-}part}(z) + n)) \Rightarrow  \\ \quad \quad \operatorname{gamma\text{-}ext}(z) = (\Gamma((z + n)) \cdot \operatorname{recip}(\operatorname{poch}(z, n))) \end{array}}}
\vnbprop{gamma-rhp-in-omega}{\vnbstmt{\begin{array}{@{}l@{}} \forall z \in \mathbb{C}. \\ \quad (0 < \operatorname{real\text{-}part}(z)) \Rightarrow (z \in \{\mathit{gdz} \in \mathbb{C} : \forall n \in \mathbb{N}.\; \neg ((\mathit{gdz} + n) = 0)\}) \end{array}}}
\vnbprop{gamma-ext-extends}{\vnbstmt{\begin{array}{@{}l@{}} \forall z \in \mathbb{C}. \\ \quad (0 < \operatorname{real\text{-}part}(z)) \Rightarrow  \\ \quad \quad ((z \in \{\mathit{gdz} \in \mathbb{C} : \forall n \in \mathbb{N}.\; \neg ((\mathit{gdz} + n) = 0)\}) \wedge (\operatorname{gamma\text{-}ext}(z) = \Gamma(z))) \end{array}}}
\vnbprop{gamma-omega-shift}{\vnbstmt{\begin{array}{@{}l@{}} \forall z \in \{\mathit{gdz} \in \mathbb{C} : \forall n \in \mathbb{N}.\; \neg ((\mathit{gdz} + n) = 0)\}. \\ \quad (z + 1) \in \{\mathit{gdz} \in \mathbb{C} : \forall n \in \mathbb{N}.\; \neg ((\mathit{gdz} + n) = 0)\} \end{array}}}
\vnbprop{gamma-ext-recursion}{\vnbstmt{\begin{array}{@{}l@{}} \forall z \in \{\mathit{gdz} \in \mathbb{C} : \forall n \in \mathbb{N}.\; \neg ((\mathit{gdz} + n) = 0)\}. \\ \quad \operatorname{gamma\text{-}ext}((z + 1)) = (z \cdot \operatorname{gamma\text{-}ext}(z)) \end{array}}}
\vnbprop{gamma-finite-separation}{\vnbstmt{\begin{array}{@{}l@{}} \forall m \in \mathbb{N}, a \in \mathbb{C}. \\ \quad \forall j \in \mathbb{N}.\; ((j < m) \Rightarrow \neg ((a + j) = 0)) \Rightarrow  \\ \quad \quad \exists r \in \mathbb{R}. \\ \quad \quad \quad (0 < r) \wedge \forall j \in \mathbb{N}.\; ((j < m) \Rightarrow (r \leq \operatorname{magnitude}((a + j)))) \end{array}}}
\vnbprop{gamma-omega-ball}{\vnbstmt{\begin{array}{@{}l@{}} \forall a \in \{z \in \mathbb{C} : \forall n \in \mathbb{N}.\; \neg ((z + n) = 0)\}, m \in \mathbb{N}. \\ \quad (0 < (\operatorname{real\text{-}part}(a) + m)) \Rightarrow  \\ \quad \quad \exists r. \\ \quad \quad \quad \operatorname{pos\text{-}rr}(r) \wedge \\ \quad \quad \quad \forall w \in \operatorname{ball}(\operatorname{cc\text{-}ms}, a, r). \\ \quad \quad \quad \quad (w \in \{z \in \mathbb{C} : \forall n \in \mathbb{N}.\; \neg ((z + n) = 0)\}) \wedge (0 < (\operatorname{real\text{-}part}(w) + m)) \end{array}}}
\vnbprop{gamma-omega-open}{\vnbstmt{\operatorname{is\text{-}open}(\operatorname{cc\text{-}ms}, \{z \in \mathbb{C} : \forall n \in \mathbb{N}.\; \neg ((z + n) = 0)\})}}
\vnbprop{poch-holomorphic}{\vnbstmt{\forall m \in \mathbb{N}.\; \operatorname{holomorphic\text{-}on}(\mathbb{C}, (\lambda w \in \mathbb{C}.\; \operatorname{poch}(w, m)))}}
\vnbprop{gamma-ext-holomorphic}{\vnbstmt{\begin{array}{@{}l@{}} \operatorname{holomorphic\text{-}on}(\{z \in \mathbb{C} : \forall n \in \mathbb{N}.\; \neg ((z + n) = 0)\}, \\ \quad (\lambda w \in \{z \in \mathbb{C} : \forall n \in \mathbb{N}.\; \neg ((z + n) = 0)\}.\; \operatorname{gamma\text{-}ext}(w))) \end{array}}}
\vnbprop{gamma-slit-in-omega}{\vnbstmt{\begin{array}{@{}l@{}} \{u \in \mathbb{C} : \neg ((\operatorname{imag\text{-}part}(u) = 0) \wedge (\operatorname{real\text{-}part}(u) \leq 0))\} \\ \subseteq\; \{z \in \mathbb{C} : \forall n \in \mathbb{N}.\; \neg ((z + n) = 0)\} \end{array}}}
\vnbprop{gamma-zero-on-slit}{\vnbstmt{\begin{array}{@{}l@{}} \forall d.; \forall z \in \{u \in \mathbb{C} : \neg ((\operatorname{imag\text{-}part}(u) = 0) \wedge (\operatorname{real\text{-}part}(u) \leq 0))\}. \\ \quad \operatorname{holomorphic\text{-}on}(\{\mathit{gdz} \in \mathbb{C} : \forall n \in \mathbb{N}.\; \neg ((\mathit{gdz} + n) = 0)\}, d) \wedge \\ \quad \forall \mathit{gyz} \in \mathbb{C}.\; ((0 < \operatorname{real\text{-}part}(\mathit{gyz})) \Rightarrow (d(\mathit{gyz}) = 0)) \Rightarrow \\ \quad \quad d(z) = 0 \end{array}}}
\vnbprop{gamma-continuation-zero}{\vnbstmt{\begin{array}{@{}l@{}} \forall d.; \forall z \in \{\mathit{gdz} \in \mathbb{C} : \forall n \in \mathbb{N}.\; \neg ((\mathit{gdz} + n) = 0)\}. \\ \quad \operatorname{holomorphic\text{-}on}(\{\mathit{gdz} \in \mathbb{C} : \forall n \in \mathbb{N}.\; \neg ((\mathit{gdz} + n) = 0)\}, d) \wedge \\ \quad \forall \mathit{gyz} \in \mathbb{C}.\; ((0 < \operatorname{real\text{-}part}(\mathit{gyz})) \Rightarrow (d(\mathit{gyz}) = 0)) \Rightarrow \\ \quad \quad d(z) = 0 \end{array}}}
\vnbprop{gamma-analytic-continuation}{\vnbstmt{\begin{array}{@{}l@{}} \operatorname{holomorphic\text{-}on}(\{z \in \mathbb{C} : \forall n \in \mathbb{N}.\; \neg ((z + n) = 0)\}, \\ \quad (\lambda w \in \{z \in \mathbb{C} : \forall n \in \mathbb{N}.\; \neg ((z + n) = 0)\}.\; \operatorname{gamma\text{-}ext}(w))) \wedge \\ \forall \mathit{gxz} \in \mathbb{C}. \\ \quad (0 < \operatorname{real\text{-}part}(\mathit{gxz})) \Rightarrow (\operatorname{gamma\text{-}ext}(\mathit{gxz}) = \Gamma(\mathit{gxz})) \wedge \\ \forall \mathit{gxz} \in \{z \in \mathbb{C} : \forall n \in \mathbb{N}.\; \neg ((z + n) = 0)\}. \\ \quad (\mathit{gxz} + 1) \in \{z \in \mathbb{C} : \forall n \in \mathbb{N}.\; \neg ((z + n) = 0)\} \wedge \\ \quad \operatorname{gamma\text{-}ext}((\mathit{gxz} + 1)) = (\mathit{gxz} \cdot \operatorname{gamma\text{-}ext}(\mathit{gxz})) \wedge \\ \forall g.; \forall \mathit{gxz} \in \{z \in \mathbb{C} : \forall n \in \mathbb{N}.\; \neg ((z + n) = 0)\}. \\ \quad \operatorname{holomorphic\text{-}on}(\{z \in \mathbb{C} : \forall n \in \mathbb{N}.\; \neg ((z + n) = 0)\}, g) \wedge \\ \quad \forall \mathit{gyz} \in \mathbb{C}.\; ((0 < \operatorname{real\text{-}part}(\mathit{gyz})) \Rightarrow (g(\mathit{gyz}) = \Gamma(\mathit{gyz}))) \Rightarrow \\ \quad \quad g(\mathit{gxz}) = \operatorname{gamma\text{-}ext}(\mathit{gxz}) \end{array}}}

\vnbbookfile{theorem-library/ch4c-inverse}{theorem-library/ch4c-inverse}
\noindent{\small ch4c-inverse.scm -- WAVE 2 (2026-10-27..30), wired 2026-10-30 from scratchpad/ch4c/ch4c-inverse.scm. Library: complex-analysis.}\par
\vnbprop{holomorphic-segment-estimate}{\vnbstmt{\begin{array}{@{}l@{}} \forall u, f.; \forall a \in \mathbb{C}, r \in \mathbb{R}, c \in \mathbb{C}, k \in \mathbb{R}, z_{1}, z_{2} \in \operatorname{ball}(\operatorname{cc\text{-}ms}, a, r). \\ \quad \operatorname{holomorphic\text{-}on}(u, f) \wedge \\ \quad 0 < r \wedge \\ \quad \operatorname{ball}(\operatorname{cc\text{-}ms}, a, r) \subseteq u \wedge \\ \quad \forall w \in \operatorname{ball}(\operatorname{cc\text{-}ms}, a, r). \\ \quad \quad \operatorname{magnitude}((\operatorname{deriv\text{-}on}(\operatorname{cc\text{-}normed\text{-}field}, u, f, w) - c)) \leq k \Rightarrow \\ \quad \quad \operatorname{magnitude}(((f(z_{1}) - f(z_{2})) - (c \cdot (z_{1} - z_{2})))) \leq (\operatorname{magnitude}((z_{1} - z_{2})) \cdot k) \end{array}}}
\vnbprop{holomorphic-locally-bilipschitz}{\vnbstmt{\begin{array}{@{}l@{}} \forall u, f.; \forall a \in u. \\ \quad (\operatorname{holomorphic\text{-}on}(u, f) \wedge \neg (\operatorname{deriv\text{-}on}(\operatorname{cc\text{-}normed\text{-}field}, u, f, a) = 0)) \Rightarrow  \\ \quad \quad \exists r \in \mathbb{R}. \\ \quad \quad \quad 0 < r \wedge \\ \quad \quad \quad \operatorname{ball}(\operatorname{cc\text{-}ms}, a, r) \subseteq u \wedge \\ \quad \quad \quad \forall w \in \operatorname{ball}(\operatorname{cc\text{-}ms}, a, r). \\ \quad \quad \quad \quad \operatorname{magnitude}((\operatorname{deriv\text{-}on}(\operatorname{cc\text{-}normed\text{-}field}, u, f, w) - \operatorname{deriv\text{-}on}(\operatorname{cc\text{-}normed\text{-}field}, u, f, a))) \\ \quad \quad \quad \quad \leq\; 1/2 \cdot \operatorname{magnitude}(\operatorname{deriv\text{-}on}(\operatorname{cc\text{-}normed\text{-}field}, u, f, a)) \wedge \\ \quad \quad \quad \forall z_{1}, z_{2} \in \operatorname{ball}(\operatorname{cc\text{-}ms}, a, r). \\ \quad \quad \quad \quad (1/2 \cdot \operatorname{magnitude}(\operatorname{deriv\text{-}on}(\operatorname{cc\text{-}normed\text{-}field}, u, f, a))) \cdot \operatorname{magnitude}((z_{1} - z_{2})) \\ \quad \quad \quad \quad \leq\; \operatorname{magnitude}((f(z_{1}) - f(z_{2}))) \end{array}}}
\vnbprop{holomorphic-local-onto}{\vnbstmt{\begin{array}{@{}l@{}} \forall u, f.; \forall a \in \mathbb{C}, r, k, p \in \mathbb{R}, v \in \mathbb{C}. \\ \quad \operatorname{holomorphic\text{-}on}(u, f) \wedge \\ \quad 0 < r \wedge \\ \quad \operatorname{ball}(\operatorname{cc\text{-}ms}, a, r) \subseteq u \wedge \\ \quad 0 < k \wedge \\ \quad \forall y \in \operatorname{ball}(\operatorname{cc\text{-}ms}, a, r).\; ((k \cdot \operatorname{magnitude}((y - a))) \leq \operatorname{magnitude}((f(y) - f(a)))) \wedge \\ \quad 0 < p \wedge \\ \quad p < r \wedge \\ \quad \operatorname{magnitude}((v - f(a))) < (1/2 \cdot (k \cdot p)) \Rightarrow \\ \quad \quad \exists z \in \operatorname{ball}(\operatorname{cc\text{-}ms}, a, p).\; (f(z) = v) \end{array}}}
\vnbprop{holomorphic-inverse-holomorphic}{\vnbstmt{\begin{array}{@{}l@{}} \forall u, f, v, w.; \forall g \in (w \to \mathbb{C}). \\ \quad \operatorname{holomorphic\text{-}on}(u, f) \wedge \\ \quad v \subseteq u \wedge \\ \quad \operatorname{is\text{-}open}(\operatorname{nf\text{-}metric\text{-}space}(\operatorname{cc\text{-}normed\text{-}field}), w) \wedge \\ \quad \forall y \in w.\; ((g(y) \in v) \wedge (f(g(y)) = y)) \wedge \\ \quad \forall \mathit{rgy} \in w. \\ \quad \quad \operatorname{is\text{-}continuous\text{-}at}(\operatorname{subspace\text{-}ms}(\operatorname{nf\text{-}metric\text{-}space}(\operatorname{cc\text{-}normed\text{-}field}), w), \\ \quad \quad \quad \operatorname{nf\text{-}metric\text{-}space}(\operatorname{cc\text{-}normed\text{-}field}), g, \mathit{rgy}) \wedge \\ \quad \forall x \in v.\; \neg (\operatorname{deriv\text{-}on}(\operatorname{cc\text{-}normed\text{-}field}, u, f, x) = 0) \Rightarrow \\ \quad \quad \operatorname{holomorphic\text{-}on}(w, g) \end{array}}}
\vnbprop{cc-open-preimage-in-open}{\vnbstmt{\begin{array}{@{}l@{}} \forall u.; \forall f \in (u \to \mathbb{C}), o, w. \\ \quad u \subseteq \mathbb{C} \wedge \\ \quad \forall y \in u. \\ \quad \quad \operatorname{is\text{-}continuous\text{-}at}(\operatorname{subspace\text{-}ms}(\operatorname{nf\text{-}metric\text{-}space}(\operatorname{cc\text{-}normed\text{-}field}), u), \\ \quad \quad \quad \operatorname{nf\text{-}metric\text{-}space}(\operatorname{cc\text{-}normed\text{-}field}), f, y) \wedge \\ \quad \operatorname{is\text{-}open}(\operatorname{cc\text{-}ms}, o) \wedge \\ \quad o \subseteq u \wedge \\ \quad \operatorname{is\text{-}open}(\operatorname{cc\text{-}ms}, w) \Rightarrow \\ \quad \quad \operatorname{is\text{-}open}(\operatorname{cc\text{-}ms}, \{z \in o : (f(z) \in w)\}) \end{array}}}
\vnbprop{holomorphic-inverse-function-theorem}{\vnbstmt{\begin{array}{@{}l@{}} \forall u, f.; \forall a \in u. \\ \quad (\operatorname{holomorphic\text{-}on}(u, f) \wedge \neg (\operatorname{deriv\text{-}on}(\operatorname{cc\text{-}normed\text{-}field}, u, f, a) = 0)) \Rightarrow  \\ \quad \quad \exists v, w, g. \\ \quad \quad \quad \operatorname{is\text{-}open}(\operatorname{nf\text{-}metric\text{-}space}(\operatorname{cc\text{-}normed\text{-}field}), v) \wedge \\ \quad \quad \quad a \in v \wedge \\ \quad \quad \quad v \subseteq u \wedge \\ \quad \quad \quad \operatorname{is\text{-}open}(\operatorname{nf\text{-}metric\text{-}space}(\operatorname{cc\text{-}normed\text{-}field}), w) \wedge \\ \quad \quad \quad g \in (w \to \mathbb{C}) \wedge \\ \quad \quad \quad \forall z \in v.\; ((f(z) \in w) \wedge (g(f(z)) = z)) \wedge \\ \quad \quad \quad \forall y \in w.\; ((g(y) \in v) \wedge (f(g(y)) = y)) \wedge \\ \quad \quad \quad \operatorname{holomorphic\text{-}on}(w, g) \end{array}}}
\vnbprop{holomorphic-local-power-form}{\vnbstmt{\begin{array}{@{}l@{}} \forall u, f, a, m. \\ \quad (\operatorname{has\text{-}zero\text{-}of\text{-}order}(u, f, a, m) \wedge (1 \leq m)) \Rightarrow  \\ \quad \quad \exists r_{7} \in \mathbb{R}, k_{7}. \\ \quad \quad \quad 0 < r_{7} \wedge \\ \quad \quad \quad \operatorname{ball}(\operatorname{cc\text{-}ms}, a, r_{7}) \subseteq u \wedge \\ \quad \quad \quad \operatorname{holomorphic\text{-}on}(\operatorname{ball}(\operatorname{cc\text{-}ms}, a, r_{7}), k_{7}) \wedge \\ \quad \quad \quad \operatorname{k7}(a) = 0 \wedge \\ \quad \quad \quad \neg (\operatorname{deriv\text{-}on}(\operatorname{cc\text{-}normed\text{-}field}, \operatorname{ball}(\operatorname{cc\text{-}ms}, a, r_{7}), k_{7}, a) = 0) \wedge \\ \quad \quad \quad \forall z_{7} \in \operatorname{ball}(\operatorname{cc\text{-}ms}, a, r_{7}).\; (f(z_{7}) = \operatorname{power}(\operatorname{k7}(z_{7}), m)) \end{array}}}
\vnbprop{cc-root-of-unity-exists}{\vnbstmt{\begin{array}{@{}l@{}} \forall m \in \mathbb{N}. \\ \quad (2 \leq m) \Rightarrow  \\ \quad \quad \exists q_{8} \in \mathbb{C}.\; ((\operatorname{power}(q_{8}, m) = 1) \wedge (\neg (q_{8} = 1) \wedge (\operatorname{magnitude}(q_{8}) = 1))) \end{array}}}
\vnbprop{holomorphic-injective-deriv-nonzero}{\vnbstmt{\begin{array}{@{}l@{}} \forall u, f.; \forall a \in u. \\ \quad \operatorname{holomorphic\text{-}on}(u, f) \wedge \\ \quad \operatorname{is\text{-}connected}(\operatorname{cc\text{-}ms}, u) \wedge \\ \quad \forall x_{9}, y_{9} \in u.\; ((f(x_{9}) = f(y_{9})) \Rightarrow (x_{9} = y_{9})) \Rightarrow \\ \quad \quad \neg (\operatorname{deriv\text{-}on}(\operatorname{cc\text{-}normed\text{-}field}, u, f, a) = 0) \end{array}}}
\vnbprop{automorphism-inverse-holomorphic}{\vnbstmt{\begin{array}{@{}l@{}} \forall u, i. \\ \quad (\operatorname{is\text{-}connected}(\operatorname{cc\text{-}ms}, u) \wedge \operatorname{is\text{-}holomorphic\text{-}automorphism}(u, i)) \Rightarrow  \\ \quad \quad \operatorname{holomorphic\text{-}on}(u, i^{-1}) \end{array}}}
\vnbprop{disc-automorphism-is-rotation}{\vnbstmt{\begin{array}{@{}l@{}} \forall i. \\ \quad (\operatorname{is\text{-}holomorphic\text{-}automorphism}(\operatorname{ball}(\operatorname{cc\text{-}ms}, 0, 1), i) \wedge (i(0) = 0)) \Rightarrow  \\ \quad \quad \exists r \in \mathbb{C}. \\ \quad \quad \quad (\operatorname{magnitude}(r) = 1) \wedge \forall z \in \operatorname{ball}(\operatorname{cc\text{-}ms}, 0, 1).\; (i(z) = (r \cdot z)) \end{array}}}
\vnbprop{cc-nth-root-exists}{\vnbstmt{\forall m \in \mathbb{N}, c_{12} \in \mathbb{C}.\; ((1 \leq m) \Rightarrow \exists w_{12} \in \mathbb{C}.\; (\operatorname{power}(w_{12}, m) = c_{12}))}}
\vnbprop{holomorphic-local-openness}{\vnbstmt{\begin{array}{@{}l@{}} \forall u, f, o.; \forall a \in o. \\ \quad \operatorname{holomorphic\text{-}on}(u, f) \wedge \\ \quad \operatorname{is\text{-}connected}(\operatorname{cc\text{-}ms}, u) \wedge \\ \quad \exists x_{13}, y_{13} \in u.\; \neg (f(x_{13}) = f(y_{13})) \wedge \\ \quad \operatorname{is\text{-}open}(\operatorname{nf\text{-}metric\text{-}space}(\operatorname{cc\text{-}normed\text{-}field}), o) \wedge \\ \quad o \subseteq u \Rightarrow \\ \quad \quad \exists e_{13}. \\ \quad \quad \quad \operatorname{pos\text{-}rr}(e_{13}) \wedge \\ \quad \quad \quad \forall v_{13} \in \mathbb{C}. \\ \quad \quad \quad \quad (\operatorname{magnitude}((v_{13} - f(a))) < e_{13}) \Rightarrow \exists z_{13} \in o.\; (f(z_{13}) = v_{13}) \end{array}}}
\vnbprop{open-mapping-theorem-holomorphic}{\vnbstmt{\begin{array}{@{}l@{}} \forall u, f, o. \\ \quad \operatorname{holomorphic\text{-}on}(u, f) \wedge \\ \quad \operatorname{is\text{-}connected}(\operatorname{cc\text{-}ms}, u) \wedge \\ \quad \exists x_{13}, y_{13} \in u.\; \neg (f(x_{13}) = f(y_{13})) \wedge \\ \quad \operatorname{is\text{-}open}(\operatorname{nf\text{-}metric\text{-}space}(\operatorname{cc\text{-}normed\text{-}field}), o) \wedge \\ \quad o \subseteq u \Rightarrow \\ \quad \quad \operatorname{is\text{-}open}(\operatorname{cc\text{-}ms}, \{w_{14} \in \mathbb{C} : \exists z_{14} \in o.\; (f(z_{14}) = w_{14})\}) \end{array}}}
\vnbprop{log-derivative-at-zero}{\vnbstmt{\begin{array}{@{}l@{}} \forall n_{15} \in \mathbb{N}, u.; \forall c \in u, f \in (u \to \mathbb{C}). \\ \quad (\operatorname{is\text{-}open}(\operatorname{nf\text{-}metric\text{-}space}(\operatorname{cc\text{-}normed\text{-}field}), u) \wedge \operatorname{has\text{-}zero\text{-}of\text{-}order}(u, f, c, n_{15})) \Rightarrow  \\ \quad \quad \exists r_{15} \in \mathbb{R}. \\ \quad \quad \quad 0 < r_{15} \wedge \\ \quad \quad \quad \operatorname{ball}(\operatorname{cc\text{-}ms}, c, r_{15}) \subseteq u \wedge \\ \quad \quad \quad \exists h_{15} \in (\operatorname{ball}(\operatorname{cc\text{-}ms}, c, r_{15}) \to \mathbb{C}). \\ \quad \quad \quad \quad \operatorname{holomorphic\text{-}on}(\operatorname{ball}(\operatorname{cc\text{-}ms}, c, r_{15}), h_{15}) \wedge \\ \quad \quad \quad \quad \forall z_{15} \in \operatorname{difference}(\operatorname{ball}(\operatorname{cc\text{-}ms}, c, r_{15}), \operatorname{singleton}(c)). \\ \quad \quad \quad \quad \quad \neg (f(z_{15}) = 0) \wedge \\ \quad \quad \quad \quad \quad (\operatorname{deriv\text{-}on}(\operatorname{cc\text{-}normed\text{-}field}, u, f, z_{15}) \cdot \operatorname{recip}(f(z_{15}))) = (\operatorname{h15}(z_{15}) + (n_{15} \cdot \operatorname{recip}((z_{15} - c)))) \end{array}}}
\vnbprop{argument-principle-holomorphic}{\vnbstmt{\begin{array}{@{}l@{}} \forall u.; \forall e \in \mathbf{Set}, f, p, q, s, t. \\ \quad \operatorname{is\text{-}open}(\operatorname{nf\text{-}metric\text{-}space}(\operatorname{cc\text{-}normed\text{-}field}), u) \wedge \\ \quad |e| \in \mathbb{N} \wedge \\ \quad e \subseteq u \wedge \\ \quad \operatorname{holomorphic\text{-}on}(u, f) \wedge \\ \quad \forall x_{16} \in u.\; ((f(x_{16}) = 0) \Rightarrow (x_{16} \in e)) \wedge \\ \quad \forall x_{16} \in e.\; (f(x_{16}) = 0) \wedge \\ \quad \forall x_{16} \in e.\; \exists m_{16}.\; \operatorname{has\text{-}zero\text{-}of\text{-}order}(u, f, x_{16}, m_{16}) \wedge \\ \quad \operatorname{is\text{-}road}(p, q, s, t) \wedge \\ \quad \operatorname{trace}(p, s, t) \subseteq \operatorname{difference}(u, e) \wedge \\ \quad p(t) = p(s) \wedge \\ \quad \operatorname{null\text{-}homotopic}(u, p, s, t) \Rightarrow \\ \quad \quad \operatorname{recip}(((2 \cdot \pi) \cdot +i)) \cdot \operatorname{line\text{-}int}((\lambda z_{16} \in \operatorname{difference}(u, e).\; (\operatorname{deriv\text{-}on}(\operatorname{cc\text{-}normed\text{-}field}, u, f, z_{16}) \cdot \operatorname{recip}(f(z_{16})))), p, q, s, t) \\ \quad \quad =\; \operatorname{finsum}(\operatorname{normed\text{-}field\text{-}additive\text{-}ag}(\operatorname{cc\text{-}normed\text{-}field}), \\ \quad \quad \quad (\lambda x \in \mathbb{C}.\; (\operatorname{zero\text{-}order}(u, f, x) \cdot \operatorname{winding}(p, q, s, t, x))), \\ \quad \quad \quad e) \end{array}}}

\vnbbookfile{theorem-library/ch4d-meromorphic}{theorem-library/ch4d-meromorphic}
\noindent{\small ch4d-meromorphic.scm -- WAVE 2 (2026-10-27..30), wired 2026-10-30 from scratchpad/ch4d/ch4d-meromorphic.scm. Library: complex-analysis.}\par
\vnbprop{cc-open-minus-finite-induction}{\vnbstmt{\begin{array}{@{}l@{}} \forall u.; \forall n \in \mathbb{N}, s \in \mathbf{Set}. \\ \quad (\operatorname{is\text{-}open}(\operatorname{nf\text{-}metric\text{-}space}(\operatorname{cc\text{-}normed\text{-}field}), u) \wedge ((|s| = n) \wedge (s \subseteq \mathbb{C}))) \Rightarrow  \\ \quad \quad \operatorname{is\text{-}open}(\operatorname{nf\text{-}metric\text{-}space}(\operatorname{cc\text{-}normed\text{-}field}), \operatorname{difference}(u, s)) \end{array}}}
\vnbprop{cc-open-minus-finite}{\vnbstmt{\begin{array}{@{}l@{}} \forall u.; \forall s \in \mathbf{Set}. \\ \quad (\operatorname{is\text{-}open}(\operatorname{nf\text{-}metric\text{-}space}(\operatorname{cc\text{-}normed\text{-}field}), u) \wedge ((|s| \in \mathbb{N}) \wedge (s \subseteq \mathbb{C}))) \Rightarrow  \\ \quad \quad \operatorname{is\text{-}open}(\operatorname{nf\text{-}metric\text{-}space}(\operatorname{cc\text{-}normed\text{-}field}), \operatorname{difference}(u, s)) \end{array}}}
\vnbprop{cc-power-sub-diff}{\vnbstmt{\begin{array}{@{}l@{}} \forall j \in \mathbb{N}, c, w \in \mathbb{C}. \\ \quad \operatorname{is\text{-}diff\text{-}on}(\operatorname{cc\text{-}normed\text{-}field}, \mathbb{C}, (\lambda v \in \mathbb{C}.\; \operatorname{power}((v - c), (j + 1))), w, ((j + 1) \cdot \operatorname{power}((w - c), j))) \end{array}}}
\vnbprop{log-derivative-at-pole}{\vnbstmt{\begin{array}{@{}l@{}} \forall u, f, c, n. \\ \quad (\operatorname{is\text{-}open}(\operatorname{nf\text{-}metric\text{-}space}(\operatorname{cc\text{-}normed\text{-}field}), u) \wedge (\operatorname{holomorphic\text{-}on}(u, f) \wedge \operatorname{is\text{-}pole}(f, c, n))) \Rightarrow  \\ \quad \quad \exists r \in \mathbb{R}. \\ \quad \quad \quad 0 < r \wedge \\ \quad \quad \quad \operatorname{difference}(\operatorname{ball}(\operatorname{cc\text{-}ms}, c, r), \operatorname{singleton}(c)) \subseteq u \wedge \\ \quad \quad \quad \exists h. \\ \quad \quad \quad \quad \operatorname{holomorphic\text{-}on}(\operatorname{ball}(\operatorname{cc\text{-}ms}, c, r), h) \wedge \\ \quad \quad \quad \quad \forall z \in \operatorname{difference}(\operatorname{ball}(\operatorname{cc\text{-}ms}, c, r), \operatorname{singleton}(c)). \\ \quad \quad \quad \quad \quad (\neg (f(z) = 0) \wedge ((\operatorname{deriv\text{-}on}(\operatorname{cc\text{-}normed\text{-}field}, u, f, z) \cdot \operatorname{recip}(f(z))) = (h(z) + (-n \cdot \operatorname{recip}((z - c)))))) \end{array}}}
\vnbprop{simple-pole-of-decomposition}{\vnbstmt{\begin{array}{@{}l@{}} \forall f, c, r, d, h. \\ \quad c \in \mathbb{C} \wedge \\ \quad r \in \mathbb{R} \wedge \\ \quad 0 < r \wedge \\ \quad d \in \mathbb{C} \wedge \\ \quad \neg (d = 0) \wedge \\ \quad \operatorname{holomorphic\text{-}on}(\operatorname{ball}(\operatorname{cc\text{-}ms}, c, r), h) \wedge \\ \quad \forall z \in \operatorname{difference}(\operatorname{ball}(\operatorname{cc\text{-}ms}, c, r), \operatorname{singleton}(c)). \\ \quad \quad f(z) = (h(z) + (d \cdot \operatorname{recip}((z - c)))) \Rightarrow \\ \quad \quad \operatorname{is\text{-}pole}(f, c, 1) \wedge (\operatorname{residue}(f, c) = d) \end{array}}}
\vnbprop{log-derivative-residue-at-pole}{\vnbstmt{\begin{array}{@{}l@{}} \forall u, f, c, n. \\ \quad (\operatorname{is\text{-}open}(\operatorname{nf\text{-}metric\text{-}space}(\operatorname{cc\text{-}normed\text{-}field}), u) \wedge (\operatorname{holomorphic\text{-}on}(u, f) \wedge \operatorname{is\text{-}pole}(f, c, n))) \Rightarrow  \\ \quad \quad \operatorname{is\text{-}pole}((\lambda w \in u.\; (\operatorname{deriv\text{-}on}(\operatorname{cc\text{-}normed\text{-}field}, u, f, w) \cdot \operatorname{recip}(f(w)))), c, 1) \wedge \\ \quad \quad \operatorname{residue}((\lambda w \in u.\; (\operatorname{deriv\text{-}on}(\operatorname{cc\text{-}normed\text{-}field}, u, f, w) \cdot \operatorname{recip}(f(w)))), c) = -n \end{array}}}
\vnbprop{log-derivative-residue-at-zero}{\vnbstmt{\begin{array}{@{}l@{}} \forall u.; \forall c \in u, f, m. \\ \quad (\operatorname{is\text{-}open}(\operatorname{nf\text{-}metric\text{-}space}(\operatorname{cc\text{-}normed\text{-}field}), u) \wedge (\operatorname{has\text{-}zero\text{-}of\text{-}order}(u, f, c, m) \wedge (1 \leq m))) \Rightarrow  \\ \quad \quad \operatorname{is\text{-}pole}((\lambda w \in u.\; (\operatorname{deriv\text{-}on}(\operatorname{cc\text{-}normed\text{-}field}, u, f, w) \cdot \operatorname{recip}(f(w)))), c, 1) \wedge \\ \quad \quad \operatorname{residue}((\lambda w \in u.\; (\operatorname{deriv\text{-}on}(\operatorname{cc\text{-}normed\text{-}field}, u, f, w) \cdot \operatorname{recip}(f(w)))), c) = m \end{array}}}
\vnbprop{pole-order-unique}{\vnbstmt{\forall f, a, n, m.\; ((\operatorname{is\text{-}pole}(f, a, n) \wedge \operatorname{is\text{-}pole}(f, a, m)) \Rightarrow (n = m))}}
\vnbprop{pole-order-value}{\vnbstmt{\forall f, a, n.\; (\operatorname{is\text{-}pole}(f, a, n) \Rightarrow (\operatorname{pole\text{-}order}(f, a) = n))}}
\vnbprop{argument-principle}{\vnbstmt{\begin{array}{@{}l@{}} \forall u.; \forall s \in \mathbf{Set}, f.; \forall \mathit{mzs} \in \mathbf{Set}, p, q, \mathit{mds}, t. \\ \quad \operatorname{is\text{-}open}(\operatorname{nf\text{-}metric\text{-}space}(\operatorname{cc\text{-}normed\text{-}field}), u) \wedge \\ \quad |s| \in \mathbb{N} \wedge \\ \quad s \subseteq u \wedge \\ \quad \operatorname{holomorphic\text{-}on}(\operatorname{difference}(u, s), f) \wedge \\ \quad \forall x \in s.\; \exists n.\; \operatorname{is\text{-}pole}(f, x, n) \wedge \\ \quad |\mathit{mzs}| \in \mathbb{N} \wedge \\ \quad \forall x \in \operatorname{difference}(u, s).\; ((f(x) = 0) \Rightarrow (x \in \mathit{mzs})) \wedge \\ \quad \forall x \in \mathit{mzs}.\; (f(x) = 0) \wedge \\ \quad \forall x \in \mathit{mzs}.\; \exists m.\; \operatorname{has\text{-}zero\text{-}of\text{-}order}(\operatorname{difference}(u, s), f, x, m) \wedge \\ \quad \operatorname{is\text{-}road}(p, q, \mathit{mds}, t) \wedge \\ \quad \operatorname{trace}(p, \mathit{mds}, t) \subseteq \operatorname{difference}(u, (\mathit{mzs} \cup s)) \wedge \\ \quad p(t) = p(\mathit{mds}) \wedge \\ \quad \operatorname{null\text{-}homotopic}(u, p, \mathit{mds}, t) \Rightarrow \\ \quad \quad \operatorname{recip}(((2 \cdot \pi) \cdot +i)) \cdot \operatorname{line\text{-}int}((\lambda z \in \operatorname{difference}(u, (\mathit{mzs} \cup s)).\; (\operatorname{deriv\text{-}on}(\operatorname{cc\text{-}normed\text{-}field}, \operatorname{difference}(u, s), f, z) \cdot \operatorname{recip}(f(z)))), p, q, \mathit{mds}, t) \\ \quad \quad =\; \operatorname{finsum}(\operatorname{normed\text{-}field\text{-}additive\text{-}ag}(\operatorname{cc\text{-}normed\text{-}field}), (\lambda \mathit{mqx} \in \mathbb{C}.\; (\operatorname{zero\text{-}order}(\operatorname{difference}(u, s), f, \mathit{mqx}) \cdot \operatorname{winding}(p, q, \mathit{mds}, t, \mathit{mqx}))), \mathit{mzs}) + \operatorname{finsum}(\operatorname{normed\text{-}field\text{-}additive\text{-}ag}(\operatorname{cc\text{-}normed\text{-}field}), (\lambda \mathit{mqx} \in \mathbb{C}.\; (-\operatorname{pole\text{-}order}(f, \mathit{mqx}) \cdot \operatorname{winding}(p, q, \mathit{mds}, t, \mathit{mqx}))), s) \end{array}}}
\vnbprop{road-holomorphic-composite}{\vnbstmt{\begin{array}{@{}l@{}} \forall u, f, p, q, s, t. \\ \quad (\operatorname{holomorphic\text{-}on}(u, f) \wedge (\operatorname{is\text{-}road}(p, q, s, t) \wedge (\operatorname{trace}(p, s, t) \subseteq u))) \Rightarrow  \\ \quad \quad \operatorname{is\text{-}road}((\lambda \mathit{c4ct} \in \operatorname{ccint}(s, t).\; f(p(\mathit{c4ct}))), \\ \quad \quad \quad (\lambda \mathit{c4ct} \in \operatorname{ccint}(s, t).\; (\operatorname{deriv\text{-}on}(\operatorname{cc\text{-}normed\text{-}field}, u, f, p(\mathit{c4ct})) \cdot q(\mathit{c4ct}))), \\ \quad \quad \quad s, t) \end{array}}}
\vnbprop{winding-unfold-eq}{\vnbstmt{\begin{array}{@{}l@{}} \forall p, q, s, t, z. \\ \quad (\operatorname{winding}(p, q, s, t, z) \simeq (\operatorname{recip}(((2 \cdot \pi) \cdot +i)) \cdot \operatorname{line\text{-}int}((\lambda w \in \operatorname{difference}(\mathbb{C}, \operatorname{singleton}(z)).\; \operatorname{recip}((w - z))), p, q, s, t))) \end{array}}}
\vnbprop{winding-of-composite}{\vnbstmt{\begin{array}{@{}l@{}} \forall u, f, p, q, s, t.; \forall w \in \mathbb{C}, v. \\ \quad \operatorname{holomorphic\text{-}on}(u, f) \wedge \\ \quad \operatorname{is\text{-}road}(p, q, s, t) \wedge \\ \quad \operatorname{trace}(p, s, t) \subseteq u \wedge \\ \quad \operatorname{trace}(p, s, t) \subseteq v \wedge \\ \quad v \subseteq u \wedge \\ \quad \forall y \in v.\; \neg (f(y) = w) \Rightarrow \\ \quad \quad \operatorname{winding}((\lambda \mathit{c4ct} \in \operatorname{ccint}(s, t).\; f(p(\mathit{c4ct}))), \\ \quad \quad \quad (\lambda \mathit{c4ct} \in \operatorname{ccint}(s, t).\; (\operatorname{deriv\text{-}on}(\operatorname{cc\text{-}normed\text{-}field}, u, f, p(\mathit{c4ct})) \cdot q(\mathit{c4ct}))), \\ \quad \quad \quad s, t, w) \\ \quad \quad =\; \operatorname{recip}(((2 \cdot \pi) \cdot +i)) \cdot \operatorname{line\text{-}int}((\lambda z \in v.\; (\operatorname{deriv\text{-}on}(\operatorname{cc\text{-}normed\text{-}field}, u, f, z) \cdot \operatorname{recip}((f(z) - w)))), p, q, s, t) \end{array}}}
\vnbprop{circle-trace-of-magnitude}{\vnbstmt{\begin{array}{@{}l@{}} \forall c \in \mathbb{C}, r \in \mathbb{R}, y \in \mathbb{C}. \\ \quad ((0 < r) \wedge (\operatorname{magnitude}((y - c)) = r)) \Rightarrow (y \in \operatorname{trace}(\operatorname{circle\text{-}path}(c, r), 0, (2 \cdot \pi))) \end{array}}}
\vnbprop{convex-closed-path-null-homotopic}{\vnbstmt{\begin{array}{@{}l@{}} \forall u, p, a, b. \\ \quad \operatorname{is\text{-}open}(\operatorname{nf\text{-}metric\text{-}space}(\operatorname{cc\text{-}normed\text{-}field}), u) \wedge \\ \quad \operatorname{is\text{-}convex\text{-}cc}(u) \wedge \\ \quad \operatorname{is\text{-}path}(p, a, b) \wedge \\ \quad \operatorname{trace}(p, a, b) \subseteq u \wedge \\ \quad p(b) = p(a) \Rightarrow \\ \quad \quad \operatorname{null\text{-}homotopic}(u, p, a, b) \end{array}}}
\vnbprop{zero-order-restrict}{\vnbstmt{\begin{array}{@{}l@{}} \forall u, f, a, m, v. \\ \quad \operatorname{has\text{-}zero\text{-}of\text{-}order}(u, f, a, m) \wedge \\ \quad \operatorname{is\text{-}open}(\operatorname{nf\text{-}metric\text{-}space}(\operatorname{cc\text{-}normed\text{-}field}), v) \wedge \\ \quad v \subseteq u \wedge \\ \quad a \in v \Rightarrow \\ \quad \quad \operatorname{has\text{-}zero\text{-}of\text{-}order}(v, \operatorname{restrict}(f, v), a, m) \end{array}}}
\vnbprop{zeros-poles-finite}{\vnbstmt{\begin{array}{@{}l@{}} \forall u.; \forall s \in \mathbf{Set}, f.; \forall \mathit{mzs} \in \mathbf{Set}. \\ \quad \operatorname{is\text{-}open}(\operatorname{nf\text{-}metric\text{-}space}(\operatorname{cc\text{-}normed\text{-}field}), u) \wedge \\ \quad |s| \in \mathbb{N} \wedge \\ \quad s \subseteq u \wedge \\ \quad \operatorname{holomorphic\text{-}on}(\operatorname{difference}(u, s), f) \wedge \\ \quad |\mathit{mzs}| \in \mathbb{N} \wedge \\ \quad \forall x \in \mathit{mzs}.\; (f(x) = 0) \Rightarrow \\ \quad \quad \mathit{mzs} \subseteq \operatorname{difference}(u, s) \wedge \\ \quad \quad (\mathit{mzs} \cup s) \in \mathbf{Set} \wedge \\ \quad \quad |(\mathit{mzs} \cup s)| \in \mathbb{N} \wedge \\ \quad \quad (\mathit{mzs} \cup s) \subseteq u \end{array}}}
\vnbprop{rouche-theorem}{\vnbstmt{\begin{array}{@{}l@{}} \forall u.; \forall \mathit{mpf} \in \mathbf{Set}, f.; \forall \mathit{mzf}, \mathit{mpg} \in \mathbf{Set}, g.; \forall \mathit{mzg} \in \mathbf{Set}, p, q, s, t. \\ \quad \operatorname{is\text{-}open}(\operatorname{nf\text{-}metric\text{-}space}(\operatorname{cc\text{-}normed\text{-}field}), u) \wedge \\ \quad |\mathit{mpf}| \in \mathbb{N} \wedge \\ \quad \mathit{mpf} \subseteq u \wedge \\ \quad \operatorname{holomorphic\text{-}on}(\operatorname{difference}(u, \mathit{mpf}), f) \wedge \\ \quad \forall x \in \mathit{mpf}.\; \exists n.\; \operatorname{is\text{-}pole}(f, x, n) \wedge \\ \quad |\mathit{mzf}| \in \mathbb{N} \wedge \\ \quad \forall x \in \operatorname{difference}(u, \mathit{mpf}).\; ((f(x) = 0) \Rightarrow (x \in \mathit{mzf})) \wedge \\ \quad \forall x \in \mathit{mzf}.\; (f(x) = 0) \wedge \\ \quad \forall x \in \mathit{mzf}.\; \exists m.\; \operatorname{has\text{-}zero\text{-}of\text{-}order}(\operatorname{difference}(u, \mathit{mpf}), f, x, m) \wedge \\ \quad |\mathit{mpg}| \in \mathbb{N} \wedge \\ \quad \mathit{mpg} \subseteq u \wedge \\ \quad \operatorname{holomorphic\text{-}on}(\operatorname{difference}(u, \mathit{mpg}), g) \wedge \\ \quad \forall x \in \mathit{mpg}.\; \exists n.\; \operatorname{is\text{-}pole}(g, x, n) \wedge \\ \quad |\mathit{mzg}| \in \mathbb{N} \wedge \\ \quad \forall x \in \operatorname{difference}(u, \mathit{mpg}).\; ((g(x) = 0) \Rightarrow (x \in \mathit{mzg})) \wedge \\ \quad \forall x \in \mathit{mzg}.\; (g(x) = 0) \wedge \\ \quad \forall x \in \mathit{mzg}.\; \exists m.\; \operatorname{has\text{-}zero\text{-}of\text{-}order}(\operatorname{difference}(u, \mathit{mpg}), g, x, m) \wedge \\ \quad \operatorname{is\text{-}road}(p, q, s, t) \wedge \\ \quad \operatorname{trace}(p, s, t) \subseteq \operatorname{difference}(u, ((\mathit{mzf} \cup \mathit{mpf}) \cup (\mathit{mzg} \cup \mathit{mpg}))) \wedge \\ \quad p(t) = p(s) \wedge \\ \quad \operatorname{null\text{-}homotopic}(u, p, s, t) \wedge \\ \quad \forall y \in \operatorname{trace}(p, s, t).\; (\operatorname{magnitude}((f(y) - g(y))) < \operatorname{magnitude}(g(y))) \Rightarrow \\ \quad \quad \operatorname{finsum}(\operatorname{normed\text{-}field\text{-}additive\text{-}ag}(\operatorname{cc\text{-}normed\text{-}field}), (\lambda \mathit{mqx} \in \mathbb{C}.\; (\operatorname{zero\text{-}order}(\operatorname{difference}(u, \mathit{mpf}), f, \mathit{mqx}) \cdot \operatorname{winding}(p, q, s, t, \mathit{mqx}))), \mathit{mzf}) + \operatorname{finsum}(\operatorname{normed\text{-}field\text{-}additive\text{-}ag}(\operatorname{cc\text{-}normed\text{-}field}), (\lambda \mathit{mqx} \in \mathbb{C}.\; (-\operatorname{pole\text{-}order}(f, \mathit{mqx}) \cdot \operatorname{winding}(p, q, s, t, \mathit{mqx}))), \mathit{mpf}) \\ \quad \quad =\; \operatorname{finsum}(\operatorname{normed\text{-}field\text{-}additive\text{-}ag}(\operatorname{cc\text{-}normed\text{-}field}), (\lambda \mathit{mqx} \in \mathbb{C}.\; (\operatorname{zero\text{-}order}(\operatorname{difference}(u, \mathit{mpg}), g, \mathit{mqx}) \cdot \operatorname{winding}(p, q, s, t, \mathit{mqx}))), \mathit{mzg}) + \operatorname{finsum}(\operatorname{normed\text{-}field\text{-}additive\text{-}ag}(\operatorname{cc\text{-}normed\text{-}field}), (\lambda \mathit{mqx} \in \mathbb{C}.\; (-\operatorname{pole\text{-}order}(g, \mathit{mqx}) \cdot \operatorname{winding}(p, q, s, t, \mathit{mqx}))), \mathit{mpg}) \end{array}}}
\vnbprop{meromorphic-at-sum}{\vnbstmt{\begin{array}{@{}l@{}} \forall f, g, s, a. \\ \quad \operatorname{is\text{-}meromorphic\text{-}at}(f, a) \wedge \\ \quad \operatorname{is\text{-}meromorphic\text{-}at}(g, a) \wedge \\ \quad \exists r \in \mathbb{R}. \\ \quad \quad ((0 < r) \wedge \forall z \in \operatorname{difference}(\operatorname{ball}(\operatorname{cc\text{-}ms}, a, r), \operatorname{singleton}(a)).\; (s(z) = (f(z) + g(z)))) \Rightarrow \\ \quad \quad \operatorname{is\text{-}meromorphic\text{-}at}(s, a) \end{array}}}
\vnbprop{meromorphic-at-product}{\vnbstmt{\begin{array}{@{}l@{}} \forall f, g, s, a. \\ \quad \operatorname{is\text{-}meromorphic\text{-}at}(f, a) \wedge \\ \quad \operatorname{is\text{-}meromorphic\text{-}at}(g, a) \wedge \\ \quad \exists r \in \mathbb{R}. \\ \quad \quad ((0 < r) \wedge \forall z \in \operatorname{difference}(\operatorname{ball}(\operatorname{cc\text{-}ms}, a, r), \operatorname{singleton}(a)).\; (s(z) = (f(z) \cdot g(z)))) \Rightarrow \\ \quad \quad \operatorname{is\text{-}meromorphic\text{-}at}(s, a) \end{array}}}
\vnbprop{meromorphic-at-recip}{\vnbstmt{\begin{array}{@{}l@{}} \forall f, a. \\ \quad \operatorname{is\text{-}meromorphic\text{-}at}(f, a) \wedge \\ \quad \forall r \in \mathbb{R}. \\ \quad \quad (0 < r) \Rightarrow  \\ \quad \quad \quad \exists z \in \operatorname{difference}(\operatorname{ball}(\operatorname{cc\text{-}ms}, a, r), \operatorname{singleton}(a)).\; \neg (f(z) = 0) \Rightarrow \\ \quad \quad \exists p \in \mathbb{R}. \\ \quad \quad \quad 0 < p \wedge \\ \quad \quad \quad \forall z \in \operatorname{difference}(\operatorname{ball}(\operatorname{cc\text{-}ms}, a, p), \operatorname{singleton}(a)).\; \neg (f(z) = 0) \wedge \\ \quad \quad \quad \forall s. \\ \quad \quad \quad \quad \forall z \in \operatorname{difference}(\operatorname{ball}(\operatorname{cc\text{-}ms}, a, p), \operatorname{singleton}(a)).\; (s(z) = \operatorname{recip}(f(z))) \Rightarrow  \\ \quad \quad \quad \quad \quad \operatorname{is\text{-}meromorphic\text{-}at}(s, a) \end{array}}}
\vnbprop{holomorphic-is-meromorphic-at}{\vnbstmt{\forall u, f.; \forall a \in u.\; (\operatorname{holomorphic\text{-}on}(u, f) \Rightarrow \operatorname{is\text{-}meromorphic\text{-}at}(f, a))}}
\vnbprop{meromorphic-on-sum}{\vnbstmt{\begin{array}{@{}l@{}} \forall u, a, f, b, g. \\ \quad \operatorname{is\text{-}open}(\operatorname{nf\text{-}metric\text{-}space}(\operatorname{cc\text{-}normed\text{-}field}), u) \wedge \\ \quad a \subseteq u \wedge \\ \quad \operatorname{holomorphic\text{-}on}(\operatorname{difference}(u, a), f) \wedge \\ \quad \forall x \in a. \\ \quad \quad \operatorname{is\text{-}removable\text{-}singularity}(f, x) \vee \exists n.\; \operatorname{is\text{-}pole}(f, x, n) \wedge \\ \quad b \subseteq u \wedge \\ \quad \operatorname{holomorphic\text{-}on}(\operatorname{difference}(u, b), g) \wedge \\ \quad \forall x \in b. \\ \quad \quad \operatorname{is\text{-}removable\text{-}singularity}(g, x) \vee \exists n.\; \operatorname{is\text{-}pole}(g, x, n) \Rightarrow \\ \quad \quad \operatorname{meromorphic\text{-}on}(u, (\lambda z \in \operatorname{difference}(u, (a \cup b)).\; (f(z) + g(z)))) \end{array}}}
\vnbprop{meromorphic-on-product}{\vnbstmt{\begin{array}{@{}l@{}} \forall u, a, f, b, g. \\ \quad \operatorname{is\text{-}open}(\operatorname{nf\text{-}metric\text{-}space}(\operatorname{cc\text{-}normed\text{-}field}), u) \wedge \\ \quad a \subseteq u \wedge \\ \quad \operatorname{holomorphic\text{-}on}(\operatorname{difference}(u, a), f) \wedge \\ \quad \forall x \in a. \\ \quad \quad \operatorname{is\text{-}removable\text{-}singularity}(f, x) \vee \exists n.\; \operatorname{is\text{-}pole}(f, x, n) \wedge \\ \quad b \subseteq u \wedge \\ \quad \operatorname{holomorphic\text{-}on}(\operatorname{difference}(u, b), g) \wedge \\ \quad \forall x \in b. \\ \quad \quad \operatorname{is\text{-}removable\text{-}singularity}(g, x) \vee \exists n.\; \operatorname{is\text{-}pole}(g, x, n) \Rightarrow \\ \quad \quad \operatorname{meromorphic\text{-}on}(u, (\lambda z \in \operatorname{difference}(u, (a \cup b)).\; (f(z) \cdot g(z)))) \end{array}}}
\vnbprop{zeros-count-winding}{\vnbstmt{\begin{array}{@{}l@{}} \forall u, f.; \forall c \in \mathbb{C}, r, q \in \mathbb{R}, w \in \mathbb{C}. \\ \quad \operatorname{is\text{-}open}(\operatorname{nf\text{-}metric\text{-}space}(\operatorname{cc\text{-}normed\text{-}field}), u) \wedge \\ \quad \operatorname{is\text{-}connected}(\operatorname{cc\text{-}ms}, u) \wedge \\ \quad \operatorname{holomorphic\text{-}on}(u, f) \wedge \\ \quad 0 < r \wedge \\ \quad r < q \wedge \\ \quad \operatorname{ball}(\operatorname{cc\text{-}ms}, c, q) \subseteq u \wedge \\ \quad \forall y \in \operatorname{trace}(\operatorname{circle\text{-}path}(c, r), 0, (2 \cdot \pi)).\; \neg (f(y) = w) \Rightarrow \\ \quad \quad \{z \in \operatorname{ball}(\operatorname{cc\text{-}ms}, c, r) : (f(z) = w)\} \in \mathbf{Set} \wedge \\ \quad \quad |\{z \in \operatorname{ball}(\operatorname{cc\text{-}ms}, c, r) : (f(z) = w)\}| \in \mathbb{N} \wedge \\ \quad \quad \operatorname{finsum}(\operatorname{normed\text{-}field\text{-}additive\text{-}ag}(\operatorname{cc\text{-}normed\text{-}field}), \\ \quad \quad \quad (\lambda x \in \mathbb{C}.\; \operatorname{zero\text{-}order}(u, (\lambda \mathit{c4nx} \in u.\; (f(\mathit{c4nx}) - w)), x)), \\ \quad \quad \quad \{z \in \operatorname{ball}(\operatorname{cc\text{-}ms}, c, r) : (f(z) = w)\}) \\ \quad \quad =\; \operatorname{winding}((\lambda t \in \operatorname{ccint}(0, (2 \cdot \pi)).\; f(\operatorname{circle\text{-}path}(c, r)(t))), \\ \quad \quad \quad (\lambda t \in \operatorname{ccint}(0, (2 \cdot \pi)).\; \\ \quad \quad \quad \quad \operatorname{deriv\text{-}on}(\operatorname{cc\text{-}normed\text{-}field}, u, f, \operatorname{circle\text{-}path}(c, r)(t)) \cdot \operatorname{circle\text{-}deriv}(c, r)(t)), \\ \quad \quad \quad 0, (2 \cdot \pi), w) \end{array}}}
\vnbprop{finsum-one-card-induction}{\vnbstmt{\begin{array}{@{}l@{}} \forall n \in \mathbb{N}, s \in \mathbf{Set}. \\ \quad ((|s| = n) \wedge (s \subseteq \mathbb{C})) \Rightarrow  \\ \quad \quad \operatorname{finsum}(\operatorname{normed\text{-}field\text{-}additive\text{-}ag}(\operatorname{cc\text{-}normed\text{-}field}), (\lambda x \in \mathbb{C}.\; 1), s) = n \end{array}}}
\vnbprop{simple-zero-order-one}{\vnbstmt{\begin{array}{@{}l@{}} \forall u, f.; \forall b \in u. \\ \quad \operatorname{is\text{-}open}(\operatorname{nf\text{-}metric\text{-}space}(\operatorname{cc\text{-}normed\text{-}field}), u) \wedge \\ \quad \operatorname{is\text{-}connected}(\operatorname{cc\text{-}ms}, u) \wedge \\ \quad \operatorname{holomorphic\text{-}on}(u, f) \wedge \\ \quad \exists w \in u.\; \neg (f(w) = 0) \wedge \\ \quad f(b) = 0 \wedge \\ \quad \neg (\operatorname{deriv\text{-}on}(\operatorname{cc\text{-}normed\text{-}field}, u, f, b) = 0) \Rightarrow \\ \quad \quad \operatorname{has\text{-}zero\text{-}of\text{-}order}(u, f, b, 1) \end{array}}}
\vnbprop{multiplicity-local-simple}{\vnbstmt{\begin{array}{@{}l@{}} \forall u, f, a, n. \\ \quad \operatorname{is\text{-}open}(\operatorname{nf\text{-}metric\text{-}space}(\operatorname{cc\text{-}normed\text{-}field}), u) \wedge \\ \quad \operatorname{holomorphic\text{-}on}(u, f) \wedge \\ \quad \operatorname{has\text{-}zero\text{-}of\text{-}order}(u, (\lambda x \in u.\; (f(x) - f(a))), a, n) \wedge \\ \quad 1 \leq n \Rightarrow \\ \quad \quad \exists r \in \mathbb{R}. \\ \quad \quad \quad 0 < r \wedge \\ \quad \quad \quad \operatorname{ball}(\operatorname{cc\text{-}ms}, a, r) \subseteq u \wedge \\ \quad \quad \quad \forall z \in \operatorname{difference}(\operatorname{ball}(\operatorname{cc\text{-}ms}, a, r), \operatorname{singleton}(a)). \\ \quad \quad \quad \quad \neg (f(z) = f(a)) \wedge \neg (\operatorname{deriv\text{-}on}(\operatorname{cc\text{-}normed\text{-}field}, u, f, z) = 0) \end{array}}}
\vnbprop{local-n-to-one}{\vnbstmt{\begin{array}{@{}l@{}} \forall u, f.; \forall a \in u, n. \\ \quad \operatorname{is\text{-}open}(\operatorname{nf\text{-}metric\text{-}space}(\operatorname{cc\text{-}normed\text{-}field}), u) \wedge \\ \quad \operatorname{is\text{-}connected}(\operatorname{cc\text{-}ms}, u) \wedge \\ \quad \operatorname{holomorphic\text{-}on}(u, f) \wedge \\ \quad \operatorname{has\text{-}zero\text{-}of\text{-}order}(u, (\lambda x \in u.\; (f(x) - f(a))), a, n) \wedge \\ \quad 1 \leq n \Rightarrow \\ \quad \quad \exists v, w. \\ \quad \quad \quad \operatorname{is\text{-}open}(\operatorname{nf\text{-}metric\text{-}space}(\operatorname{cc\text{-}normed\text{-}field}), v) \wedge \\ \quad \quad \quad v \subseteq u \wedge \\ \quad \quad \quad a \in v \wedge \\ \quad \quad \quad \operatorname{is\text{-}open}(\operatorname{nf\text{-}metric\text{-}space}(\operatorname{cc\text{-}normed\text{-}field}), w) \wedge \\ \quad \quad \quad \forall z \in v.\; (f(z) \in w) \wedge \\ \quad \quad \quad \forall y \in w.\; \exists b \in v.\; (f(b) = y) \wedge \\ \quad \quad \quad \forall y \in w.\; (\neg (y = f(a)) \Rightarrow (|\{b \in v : (f(b) = y)\}| = n)) \end{array}}}

\vnbbookfile{theorem-library/ch6-prod}{theorem-library/ch6-prod}
\noindent{\small ch6-prod.scm -- WAVE 2 (2026-10-27..30), wired 2026-10-30 from scratchpad/ch6/ch6-prod.scm (proofs). Library: complex-analysis.}\par
\vnbprop{pprod-in-cc}{\vnbstmt{\forall n \in \mathbb{N}, f \in (\mathbb{N} \to \mathbb{C}).\; (\operatorname{pprod}(f, n) \in \mathbb{C})}}
\vnbprop{pprod-in-rr}{\vnbstmt{\forall n \in \mathbb{N}, f \in (\mathbb{N} \to \mathbb{R}).\; (\operatorname{pprod}(f, n) \in \mathbb{R})}}
\vnbprop{pprod-split}{\vnbstmt{\begin{array}{@{}l@{}} \forall p, n \in \mathbb{N}, f \in (\mathbb{N} \to \mathbb{C}). \\ \quad \operatorname{pprod}(f, (n + p)) = (\operatorname{pprod}(f, n) \cdot \operatorname{pprod}((\lambda k \in \mathbb{N}.\; f((n + k))), p)) \end{array}}}
\vnbprop{pprod-magnitude}{\vnbstmt{\begin{array}{@{}l@{}} \forall n \in \mathbb{N}, f \in (\mathbb{N} \to \mathbb{C}). \\ \quad \operatorname{magnitude}(\operatorname{pprod}(f, n)) = \operatorname{pprod}((\lambda k \in \mathbb{N}.\; \operatorname{magnitude}(f(k))), n) \end{array}}}
\vnbprop{pprod-nonzero}{\vnbstmt{\begin{array}{@{}l@{}} \forall n \in \mathbb{N}, f \in (\mathbb{N} \to \mathbb{C}). \\ \quad \forall j \in \mathbb{N}.\; ((j < n) \Rightarrow \neg (f(j) = 0)) \Rightarrow \neg (\operatorname{pprod}(f, n) = 0) \end{array}}}
\vnbprop{pprod-zero-factor}{\vnbstmt{\begin{array}{@{}l@{}} \forall n \in \mathbb{N}, f \in (\mathbb{N} \to \mathbb{C}), j \in \mathbb{N}. \\ \quad ((j < n) \wedge (f(j) = 0)) \Rightarrow (\operatorname{pprod}(f, n) = 0) \end{array}}}
\vnbprop{pprod-one-plus-le}{\vnbstmt{\begin{array}{@{}l@{}} \forall n \in \mathbb{N}, a \in (\mathbb{N} \to \mathbb{C}). \\ \quad \operatorname{magnitude}((\operatorname{pprod}((\lambda k \in \mathbb{N}.\; (1 + a(k))), n) - 1)) \\ \quad \leq\; (\operatorname{pprod}((\lambda k \in \mathbb{N}.\; (1 + \operatorname{magnitude}(a(k)))), n) - 1) \end{array}}}
\vnbprop{pprod-exp-bounds}{\vnbstmt{\begin{array}{@{}l@{}} \forall n \in \mathbb{N}, r \in (\mathbb{N} \to \mathbb{R}). \\ \quad \forall i \in \mathbb{N}.\; (0 \leq r(i)) \Rightarrow  \\ \quad \quad (1 + \operatorname{series\text{-}partial\text{-}sum}(r, n)) \leq \operatorname{pprod}((\lambda k \in \mathbb{N}.\; (1 + r(k))), n) \wedge \\ \quad \quad \operatorname{pprod}((\lambda k \in \mathbb{N}.\; (1 + r(k))), n) \leq \exp\left(\operatorname{series\text{-}partial\text{-}sum}(r, n)\right) \end{array}}}
\vnbprop{r-exp-neg-two-le}{\vnbstmt{\forall x \in \mathbb{R}.\; (((0 \leq x) \wedge (x \leq 1/2)) \Rightarrow (\exp\left((-2 \cdot x)\right) \leq (1 - x)))}}
\vnbprop{pprod-lower}{\vnbstmt{\begin{array}{@{}l@{}} \forall n \in \mathbb{N}, a \in (\mathbb{N} \to \mathbb{C}). \\ \quad \forall i \in \mathbb{N}.\; (\operatorname{magnitude}(a(i)) \leq 1/2) \Rightarrow  \\ \quad \quad \exp\left((-2 \cdot \operatorname{series\text{-}partial\text{-}sum}((\lambda k \in \mathbb{N}.\; \operatorname{magnitude}(a(k))), n))\right) \\ \quad \quad \leq\; \operatorname{magnitude}(\operatorname{pprod}((\lambda k \in \mathbb{N}.\; (1 + a(k))), n)) \end{array}}}
\vnbprop{pprod-mag-le}{\vnbstmt{\begin{array}{@{}l@{}} \forall n \in \mathbb{N}, a \in (\mathbb{N} \to \mathbb{C}). \\ \quad \operatorname{magnitude}(\operatorname{pprod}((\lambda k \in \mathbb{N}.\; (1 + a(k))), n)) \leq \operatorname{pprod}((\lambda k \in \mathbb{N}.\; (1 + \operatorname{magnitude}(a(k)))), n) \end{array}}}
\vnbprop{pprod-diff-le}{\vnbstmt{\begin{array}{@{}l@{}} \forall p, n \in \mathbb{N}, a \in (\mathbb{N} \to \mathbb{C}). \\ \quad \operatorname{magnitude}((\operatorname{pprod}((\lambda k \in \mathbb{N}.\; (1 + a(k))), (n + p)) - \operatorname{pprod}((\lambda k \in \mathbb{N}.\; (1 + a(k))), n))) \\ \quad \leq\; (\operatorname{pprod}((\lambda k \in \mathbb{N}.\; (1 + \operatorname{magnitude}(a(k)))), (n + p)) - \operatorname{pprod}((\lambda k \in \mathbb{N}.\; (1 + \operatorname{magnitude}(a(k)))), n)) \end{array}}}
\vnbprop{r-exp-sub-one-le}{\vnbstmt{\forall x \in \mathbb{R}.\; ((0 \leq x) \Rightarrow ((\exp\left(x\right) - 1) \leq (x \cdot \exp\left(x\right))))}}
\vnbprop{pprod-ratio-le}{\vnbstmt{\begin{array}{@{}l@{}} \forall p, n \in \mathbb{N}, r \in (\mathbb{N} \to \mathbb{R}). \\ \quad \forall i \in \mathbb{N}.\; (0 \leq r(i)) \Rightarrow  \\ \quad \quad \operatorname{pprod}((\lambda k \in \mathbb{N}.\; (1 + r(k))), (n + p)) \\ \quad \quad \leq\; \operatorname{pprod}((\lambda k \in \mathbb{N}.\; (1 + r(k))), n) \cdot \exp\left((\operatorname{series\text{-}partial\text{-}sum}(r, (n + p)) - \operatorname{series\text{-}partial\text{-}sum}(r, n))\right) \end{array}}}
\vnbprop{pprod-cauchy-bound}{\vnbstmt{\begin{array}{@{}l@{}} \forall p, n \in \mathbb{N}, a \in (\mathbb{N} \to \mathbb{C}), b \in \mathbb{R}. \\ \quad \forall i \in \mathbb{N}.\; (\operatorname{series\text{-}partial\text{-}sum}((\lambda k \in \mathbb{N}.\; \operatorname{magnitude}(a(k))), i) \leq b) \Rightarrow  \\ \quad \quad \operatorname{magnitude}((\operatorname{pprod}((\lambda k \in \mathbb{N}.\; (1 + a(k))), (n + p)) - \operatorname{pprod}((\lambda k \in \mathbb{N}.\; (1 + a(k))), n))) \\ \quad \quad \leq\; (\exp\left(b\right) \cdot \exp\left(b\right)) \cdot (\operatorname{series\text{-}partial\text{-}sum}((\lambda k \in \mathbb{N}.\; \operatorname{magnitude}(a(k))), (n + p)) - \operatorname{series\text{-}partial\text{-}sum}((\lambda k \in \mathbb{N}.\; \operatorname{magnitude}(a(k))), n)) \end{array}}}
\vnbprop{pprod-small-converges}{\vnbstmt{\begin{array}{@{}l@{}} \forall a \in (\mathbb{N} \to \mathbb{C}). \\ \quad \forall i \in \mathbb{N}.\; (\operatorname{magnitude}(a(i)) \leq 1/2) \wedge \\ \quad \operatorname{series\text{-}converges}((\lambda k \in \mathbb{N}.\; \operatorname{magnitude}(a(k)))) \Rightarrow \\ \quad \quad \exists z \in \mathbb{C}. \\ \quad \quad \quad (\neg (z = 0) \wedge \operatorname{converges\text{-}to}(\operatorname{cc\text{-}ms}, (\lambda n \in \mathbb{N}.\; \operatorname{pprod}((\lambda k \in \mathbb{N}.\; (1 + a(k))), n)), z)) \end{array}}}
\vnbprop{sps-shift}{\vnbstmt{\begin{array}{@{}l@{}} \forall n, m \in \mathbb{N}, a \in (\mathbb{N} \to \mathbb{C}). \\ \quad \operatorname{series\text{-}partial\text{-}sum}((\lambda k \in \mathbb{N}.\; \operatorname{magnitude}(a((m + k)))), n) \\ \quad =\; (\operatorname{series\text{-}partial\text{-}sum}((\lambda k \in \mathbb{N}.\; \operatorname{magnitude}(a(k))), (m + n)) - \operatorname{series\text{-}partial\text{-}sum}((\lambda k \in \mathbb{N}.\; \operatorname{magnitude}(a(k))), m)) \end{array}}}
\vnbprop{series-shift-abs-converges}{\vnbstmt{\begin{array}{@{}l@{}} \forall m \in \mathbb{N}, a \in (\mathbb{N} \to \mathbb{C}). \\ \quad \operatorname{series\text{-}converges}((\lambda k \in \mathbb{N}.\; \operatorname{magnitude}(a(k)))) \Rightarrow  \\ \quad \quad \operatorname{series\text{-}converges}((\lambda k \in \mathbb{N}.\; \operatorname{magnitude}(a((m + k))))) \end{array}}}
\vnbprop{abs-series-terms-small}{\vnbstmt{\begin{array}{@{}l@{}} \forall a \in (\mathbb{N} \to \mathbb{C}). \\ \quad \operatorname{series\text{-}converges}((\lambda k \in \mathbb{N}.\; \operatorname{magnitude}(a(k)))) \Rightarrow  \\ \quad \quad \exists m \in \mathbb{N}. \forall q \in \mathbb{N}. \\ \quad \quad \quad (m \leq q) \Rightarrow (\operatorname{magnitude}(a(q)) \leq 1/2) \end{array}}}
\vnbprop{prod-abs-converges-exists}{\vnbstmt{\begin{array}{@{}l@{}} \forall a \in (\mathbb{N} \to \mathbb{C}). \\ \quad \operatorname{series\text{-}converges}((\lambda k \in \mathbb{N}.\; \operatorname{magnitude}(a(k)))) \Rightarrow  \\ \quad \quad \exists l.\; \operatorname{prod\text{-}converges\text{-}to}((\lambda k \in \mathbb{N}.\; (1 + a(k))), l) \end{array}}}
\vnbprop{prod-zero-iff}{\vnbstmt{\begin{array}{@{}l@{}} \forall f \in (\mathbb{N} \to \mathbb{C}), l. \\ \quad \operatorname{prod\text{-}converges\text{-}to}(f, l) \Rightarrow ((l = 0) \Leftrightarrow \exists q \in \mathbb{N}.\; (f(q) = 0)) \end{array}}}
\vnbprop{abs-convergent-product}{\vnbstmt{\begin{array}{@{}l@{}} \forall a \in (\mathbb{N} \to \mathbb{C}). \\ \quad \operatorname{series\text{-}converges}((\lambda k \in \mathbb{N}.\; \operatorname{magnitude}(a(k)))) \Rightarrow  \\ \quad \quad \operatorname{prod\text{-}converges}((\lambda k \in \mathbb{N}.\; (1 + a(k)))) \wedge \\ \quad \quad \forall l. \\ \quad \quad \quad \operatorname{prod\text{-}converges\text{-}to}((\lambda k \in \mathbb{N}.\; (1 + a(k))), l) \Rightarrow  \\ \quad \quad \quad \quad (l = 0) \Leftrightarrow \exists q \in \mathbb{N}.\; ((1 + a(q)) = 0) \end{array}}}
\vnbprop{pprod-pos}{\vnbstmt{\begin{array}{@{}l@{}} \forall n \in \mathbb{N}, f \in (\mathbb{N} \to \mathbb{R}). \\ \quad \forall i \in \mathbb{N}.\; (0 < f(i)) \Rightarrow (0 < \operatorname{pprod}(f, n)) \end{array}}}
\vnbprop{pprod-dec}{\vnbstmt{\begin{array}{@{}l@{}} \forall p, n \in \mathbb{N}, f \in (\mathbb{N} \to \mathbb{R}). \\ \quad (\forall i \in \mathbb{N}.\; (0 < f(i)) \wedge \forall i \in \mathbb{N}.\; (f(i) \leq 1)) \Rightarrow  \\ \quad \quad \operatorname{pprod}(f, (n + p)) \leq \operatorname{pprod}(f, n) \end{array}}}
\vnbprop{log-pprod}{\vnbstmt{\begin{array}{@{}l@{}} \forall n \in \mathbb{N}, r \in (\mathbb{N} \to \mathbb{R}). \\ \quad \forall i \in \mathbb{N}.\; (0 \leq r(i)) \Rightarrow  \\ \quad \quad \log\left(\operatorname{pprod}((\lambda k \in \mathbb{N}.\; (1 + r(k))), n)\right) = \operatorname{series\text{-}partial\text{-}sum}((\lambda k \in \mathbb{N}.\; \log\left((1 + r(k))\right)), n) \end{array}}}
\vnbprop{pprod-omr-ratio}{\vnbstmt{\begin{array}{@{}l@{}} \forall p, n \in \mathbb{N}, r \in (\mathbb{N} \to \mathbb{R}). \\ \quad \forall i \in \mathbb{N}.\; (0 \leq r(i)) \wedge \\ \quad \forall i \in \mathbb{N}.\; (r(i) < 1) \wedge \\ \quad \forall i \in \mathbb{N}.\; ((n \leq i) \Rightarrow (r(i) \leq 1/2)) \Rightarrow \\ \quad \quad \operatorname{pprod}((\lambda k \in \mathbb{N}.\; (1 - r(k))), n) \cdot \exp\left((-2 \cdot (\operatorname{series\text{-}partial\text{-}sum}(r, (n + p)) - \operatorname{series\text{-}partial\text{-}sum}(r, n)))\right) \\ \quad \quad \leq\; \operatorname{pprod}((\lambda k \in \mathbb{N}.\; (1 - r(k))), (n + p)) \end{array}}}
\vnbprop{pprod-opr-omr-le}{\vnbstmt{\begin{array}{@{}l@{}} \forall n \in \mathbb{N}, r \in (\mathbb{N} \to \mathbb{R}). \\ \quad (\forall i \in \mathbb{N}.\; (0 \leq r(i)) \wedge \forall i \in \mathbb{N}.\; (r(i) \leq 1)) \Rightarrow  \\ \quad \quad 0 \leq (\operatorname{pprod}((\lambda k \in \mathbb{N}.\; (1 + r(k))), n) \cdot \operatorname{pprod}((\lambda k \in \mathbb{N}.\; (1 - r(k))), n)) \wedge \\ \quad \quad (\operatorname{pprod}((\lambda k \in \mathbb{N}.\; (1 + r(k))), n) \cdot \operatorname{pprod}((\lambda k \in \mathbb{N}.\; (1 - r(k))), n)) \leq 1 \end{array}}}
\vnbprop{lemma-6-4-4-2}{\vnbstmt{\begin{array}{@{}l@{}} \forall r \in (\mathbb{N} \to \mathbb{R}). \\ \quad (\forall i \in \mathbb{N}.\; (0 < r(i)) \wedge (\forall i \in \mathbb{N}.\; (r(i) < 1) \wedge \operatorname{series\text{-}converges}(r))) \Rightarrow  \\ \quad \quad \exists b \in \mathbb{R}. \forall j \in \mathbb{N}. \\ \quad \quad \quad \operatorname{pprod}((\lambda k \in \mathbb{N}.\; (1 + r(k))), j) \leq b \end{array}}}
\vnbprop{lemma-6-4-2-4}{\vnbstmt{\begin{array}{@{}l@{}} \forall r \in (\mathbb{N} \to \mathbb{R}). \\ \quad \forall i \in \mathbb{N}.\; (0 < r(i)) \wedge \\ \quad \forall i \in \mathbb{N}.\; (r(i) < 1) \wedge \\ \quad \exists b \in \mathbb{R}. \forall j \in \mathbb{N}. \\ \quad \quad \operatorname{pprod}((\lambda k \in \mathbb{N}.\; (1 + r(k))), j) \leq b \Rightarrow \\ \quad \quad \operatorname{series\text{-}converges}(r) \end{array}}}
\vnbprop{lemma-6-4-2-1}{\vnbstmt{\begin{array}{@{}l@{}} \forall r \in (\mathbb{N} \to \mathbb{R}). \\ \quad \forall i \in \mathbb{N}.\; (0 < r(i)) \wedge \\ \quad \forall i \in \mathbb{N}.\; (r(i) < 1) \wedge \\ \quad \exists b \in \mathbb{R}. \forall j \in \mathbb{N}. \\ \quad \quad \operatorname{pprod}((\lambda k \in \mathbb{N}.\; (1 + r(k))), j) \leq b \Rightarrow \\ \quad \quad \operatorname{series\text{-}converges}((\lambda k \in \mathbb{N}.\; \log\left((1 + r(k))\right))) \end{array}}}
\vnbprop{lemma-6-4-1-2}{\vnbstmt{\begin{array}{@{}l@{}} \forall r \in (\mathbb{N} \to \mathbb{R}). \\ \quad \forall i \in \mathbb{N}.\; (0 < r(i)) \wedge \\ \quad \forall i \in \mathbb{N}.\; (r(i) < 1) \wedge \\ \quad \operatorname{series\text{-}converges}((\lambda k \in \mathbb{N}.\; \log\left((1 + r(k))\right))) \Rightarrow \\ \quad \quad \exists b \in \mathbb{R}. \forall j \in \mathbb{N}. \\ \quad \quad \quad \operatorname{pprod}((\lambda k \in \mathbb{N}.\; (1 + r(k))), j) \leq b \end{array}}}
\vnbprop{lemma-6-4-4-3}{\vnbstmt{\begin{array}{@{}l@{}} \forall r \in (\mathbb{N} \to \mathbb{R}). \\ \quad (\forall i \in \mathbb{N}.\; (0 < r(i)) \wedge (\forall i \in \mathbb{N}.\; (r(i) < 1) \wedge \operatorname{series\text{-}converges}(r))) \Rightarrow  \\ \quad \quad \exists c \in \mathbb{R}. \\ \quad \quad \quad (0 < c) \wedge \forall j \in \mathbb{N}.\; (c \leq \operatorname{pprod}((\lambda k \in \mathbb{N}.\; (1 - r(k))), j)) \end{array}}}
\vnbprop{lemma-6-4-3-2}{\vnbstmt{\begin{array}{@{}l@{}} \forall r \in (\mathbb{N} \to \mathbb{R}). \\ \quad \forall i \in \mathbb{N}.\; (0 < r(i)) \wedge \\ \quad \forall i \in \mathbb{N}.\; (r(i) < 1) \wedge \\ \quad \exists c \in \mathbb{R}. \\ \quad \quad (0 < c) \wedge \forall j \in \mathbb{N}.\; (c \leq \operatorname{pprod}((\lambda k \in \mathbb{N}.\; (1 - r(k))), j)) \Rightarrow \\ \quad \quad \exists b \in \mathbb{R}. \forall j \in \mathbb{N}. \\ \quad \quad \quad \operatorname{pprod}((\lambda k \in \mathbb{N}.\; (1 + r(k))), j) \leq b \end{array}}}
\vnbprop{lemma-6-4}{\vnbstmt{\begin{array}{@{}l@{}} \forall r \in (\mathbb{N} \to \mathbb{R}). \\ \quad (\forall i \in \mathbb{N}.\; (0 < r(i)) \wedge \forall i \in \mathbb{N}.\; (r(i) < 1)) \Rightarrow  \\ \quad \quad \operatorname{series\text{-}converges}((\lambda k \in \mathbb{N}.\; \log\left((1 + r(k))\right))) \\ \quad \quad \Leftrightarrow\; \exists b \in \mathbb{R}. \forall j \in \mathbb{N}. \\ \quad \quad \quad \operatorname{pprod}((\lambda k \in \mathbb{N}.\; (1 + r(k))), j) \leq b \wedge \\ \quad \quad \exists b \in \mathbb{R}. \forall j \in \mathbb{N}. \\ \quad \quad \quad \operatorname{pprod}((\lambda k \in \mathbb{N}.\; (1 + r(k))), j) \leq b \\ \quad \quad \Leftrightarrow\; \operatorname{series\text{-}converges}(r) \wedge \\ \quad \quad \exists c \in \mathbb{R}. \\ \quad \quad \quad (0 < c) \wedge \forall j \in \mathbb{N}.\; (c \leq \operatorname{pprod}((\lambda k \in \mathbb{N}.\; (1 - r(k))), j)) \\ \quad \quad \Leftrightarrow\; \operatorname{series\text{-}converges}(r) \end{array}}}
\vnbprop{cc-log-one-plus-near-zero}{\vnbstmt{\begin{array}{@{}l@{}} \forall z \in \mathbb{C}. \\ \quad (\operatorname{magnitude}(z) \leq 1/10) \Rightarrow  \\ \quad \quad \neg ((\operatorname{imag\text{-}part}((1 + z)) = 0) \wedge (\operatorname{real\text{-}part}((1 + z)) \leq 0)) \wedge \\ \quad \quad \operatorname{magnitude}(\log\left((1 + z)\right)) \leq (2 \cdot \operatorname{magnitude}(z)) \wedge \\ \quad \quad \operatorname{magnitude}(z) \leq (2 \cdot \operatorname{magnitude}(\log\left((1 + z)\right))) \end{array}}}
\vnbprop{abs-series-terms-small-eps}{\vnbstmt{\begin{array}{@{}l@{}} \forall a \in (\mathbb{N} \to \mathbb{C}), e. \\ \quad (\operatorname{series\text{-}converges}((\lambda k \in \mathbb{N}.\; \operatorname{magnitude}(a(k)))) \wedge \operatorname{pos\text{-}rr}(e)) \Rightarrow  \\ \quad \quad \exists m \in \mathbb{N}.\; \forall q \in \mathbb{N}.\; ((m \leq q) \Rightarrow (\operatorname{magnitude}(a(q)) \leq e)) \end{array}}}
\vnbprop{series-unshift-abs-converges}{\vnbstmt{\begin{array}{@{}l@{}} \forall m \in \mathbb{N}, a \in (\mathbb{N} \to \mathbb{C}). \\ \quad \operatorname{series\text{-}converges}((\lambda k \in \mathbb{N}.\; \operatorname{magnitude}(a((m + k))))) \Rightarrow  \\ \quad \quad \operatorname{series\text{-}converges}((\lambda k \in \mathbb{N}.\; \operatorname{magnitude}(a(k)))) \end{array}}}
\vnbprop{prop-6-8-forward}{\vnbstmt{\begin{array}{@{}l@{}} \forall m \in \mathbb{N}, a \in (\mathbb{N} \to \mathbb{C}). \\ \quad \operatorname{series\text{-}converges}((\lambda k \in \mathbb{N}.\; \operatorname{magnitude}(a(k)))) \wedge \\ \quad \forall q \in \mathbb{N}.\; ((m \leq q) \Rightarrow (\operatorname{magnitude}(a(q)) \leq 1/10)) \Rightarrow \\ \quad \quad \operatorname{series\text{-}converges}((\lambda k \in \mathbb{N}.\; \operatorname{magnitude}(\log\left((1 + a((m + k)))\right)))) \end{array}}}
\vnbprop{prop-6-8-backward}{\vnbstmt{\begin{array}{@{}l@{}} \forall m \in \mathbb{N}, a \in (\mathbb{N} \to \mathbb{C}). \\ \quad \forall q \in \mathbb{N}. \\ \quad \quad (m \leq q) \Rightarrow \neg ((\operatorname{imag\text{-}part}((1 + a(q))) = 0) \wedge (\operatorname{real\text{-}part}((1 + a(q))) \leq 0)) \wedge \\ \quad \operatorname{series\text{-}converges}((\lambda k \in \mathbb{N}.\; \operatorname{magnitude}(\log\left((1 + a((m + k)))\right)))) \Rightarrow \\ \quad \quad \operatorname{series\text{-}converges}((\lambda k \in \mathbb{N}.\; \operatorname{magnitude}(a(k)))) \end{array}}}
\vnbprop{prop-6-8}{\vnbstmt{\begin{array}{@{}l@{}} \forall a \in (\mathbb{N} \to \mathbb{C}). \\ \quad \operatorname{series\text{-}converges}((\lambda k \in \mathbb{N}.\; \operatorname{magnitude}(a(k)))) \\ \quad \Leftrightarrow\; \exists m \in \mathbb{N}. \\ \quad \quad \forall q \in \mathbb{N}. \\ \quad \quad \quad (m \leq q) \Rightarrow \neg ((\operatorname{imag\text{-}part}((1 + a(q))) = 0) \wedge (\operatorname{real\text{-}part}((1 + a(q))) \leq 0)) \wedge \\ \quad \quad \operatorname{series\text{-}converges}((\lambda k \in \mathbb{N}.\; \operatorname{magnitude}(\log\left((1 + a((m + k)))\right)))) \end{array}}}
\vnbprop{csps-agree-below}{\vnbstmt{\begin{array}{@{}l@{}} \forall n \in \mathbb{N}, f, g \in (\mathbb{N} \to \mathbb{C}). \\ \quad \forall j \in \mathbb{N}.\; ((j < n) \Rightarrow (f(j) = g(j))) \Rightarrow  \\ \quad \quad \operatorname{cc\text{-}series\text{-}partial\text{-}sum}(f, n) = \operatorname{cc\text{-}series\text{-}partial\text{-}sum}(g, n) \end{array}}}
\vnbprop{csps-add}{\vnbstmt{\begin{array}{@{}l@{}} \forall n \in \mathbb{N}, f, g \in (\mathbb{N} \to \mathbb{C}). \\ \quad \operatorname{cc\text{-}series\text{-}partial\text{-}sum}((\lambda k \in \mathbb{N}.\; (f(k) + g(k))), n) \\ \quad =\; \operatorname{cc\text{-}series\text{-}partial\text{-}sum}(f, n) + \operatorname{cc\text{-}series\text{-}partial\text{-}sum}(g, n) \end{array}}}
\vnbprop{csps-delta}{\vnbstmt{\begin{array}{@{}l@{}} \forall n, s \in \mathbb{N}, c \in \mathbb{C}. \\ \quad (s < n) \Rightarrow (\operatorname{cc\text{-}series\text{-}partial\text{-}sum}((\lambda k \in \mathbb{N}.\; \left(\text{if }(k = s)\text{ then }c\text{ else }0\right)), n) = c) \wedge \\ \quad (n \leq s) \Rightarrow (\operatorname{cc\text{-}series\text{-}partial\text{-}sum}((\lambda k \in \mathbb{N}.\; \left(\text{if }(k = s)\text{ then }c\text{ else }0\right)), n) = 0) \end{array}}}
\vnbprop{csps-trunc}{\vnbstmt{\begin{array}{@{}l@{}} \forall p, n \in \mathbb{N}, a \in (\mathbb{N} \to \mathbb{C}). \\ \quad \operatorname{cc\text{-}series\text{-}partial\text{-}sum}((\lambda k \in \mathbb{N}.\; \left(\text{if }(k < n)\text{ then }a(k)\text{ else }0\right)), (n + p)) = \operatorname{cc\text{-}series\text{-}partial\text{-}sum}(a, n) \end{array}}}
\vnbprop{csps-mask}{\vnbstmt{\begin{array}{@{}l@{}} \forall m, n \in \mathbb{N}, a \in (\mathbb{N} \to \mathbb{C}), g \in (\mathbb{N} \to \mathbb{N}). \\ \quad (\forall u, v \in \mathbb{N}.\; ((g(u) = g(v)) \Rightarrow (u = v)) \wedge \forall u \in \mathbb{N}.\; ((u < m) \Rightarrow (g(u) < n))) \Rightarrow  \\ \quad \quad \operatorname{cc\text{-}series\text{-}partial\text{-}sum}((\lambda k \in \mathbb{N}.\; a(g(k))), m) \\ \quad \quad =\; \operatorname{cc\text{-}series\text{-}partial\text{-}sum}((\lambda k \in \mathbb{N}.\; \left(\text{if }\exists t \in \mathbb{N}.\; ((t < m) \wedge (g(t) = k))\text{ then }a(k)\text{ else }0\right)), n) \end{array}}}
\vnbprop{rsps-agree-below}{\vnbstmt{\begin{array}{@{}l@{}} \forall n \in \mathbb{N}, f, g \in (\mathbb{N} \to \mathbb{R}). \\ \quad \forall j \in \mathbb{N}.\; ((j < n) \Rightarrow (f(j) = g(j))) \Rightarrow  \\ \quad \quad \operatorname{series\text{-}partial\text{-}sum}(f, n) = \operatorname{series\text{-}partial\text{-}sum}(g, n) \end{array}}}
\vnbprop{rsps-add}{\vnbstmt{\begin{array}{@{}l@{}} \forall n \in \mathbb{N}, f, g \in (\mathbb{N} \to \mathbb{R}). \\ \quad \operatorname{series\text{-}partial\text{-}sum}((\lambda k \in \mathbb{N}.\; (f(k) + g(k))), n) \\ \quad =\; \operatorname{series\text{-}partial\text{-}sum}(f, n) + \operatorname{series\text{-}partial\text{-}sum}(g, n) \end{array}}}
\vnbprop{rsps-delta}{\vnbstmt{\begin{array}{@{}l@{}} \forall n, s \in \mathbb{N}, c \in \mathbb{R}. \\ \quad (s < n) \Rightarrow (\operatorname{series\text{-}partial\text{-}sum}((\lambda k \in \mathbb{N}.\; \left(\text{if }(k = s)\text{ then }c\text{ else }0\right)), n) = c) \wedge \\ \quad (n \leq s) \Rightarrow (\operatorname{series\text{-}partial\text{-}sum}((\lambda k \in \mathbb{N}.\; \left(\text{if }(k = s)\text{ then }c\text{ else }0\right)), n) = 0) \end{array}}}
\vnbprop{rsps-trunc}{\vnbstmt{\begin{array}{@{}l@{}} \forall p, n \in \mathbb{N}, a \in (\mathbb{N} \to \mathbb{R}). \\ \quad \operatorname{series\text{-}partial\text{-}sum}((\lambda k \in \mathbb{N}.\; \left(\text{if }(k < n)\text{ then }a(k)\text{ else }0\right)), (n + p)) = \operatorname{series\text{-}partial\text{-}sum}(a, n) \end{array}}}
\vnbprop{rsps-mask}{\vnbstmt{\begin{array}{@{}l@{}} \forall m, n \in \mathbb{N}, a \in (\mathbb{N} \to \mathbb{R}), g \in (\mathbb{N} \to \mathbb{N}). \\ \quad (\forall u, v \in \mathbb{N}.\; ((g(u) = g(v)) \Rightarrow (u = v)) \wedge \forall u \in \mathbb{N}.\; ((u < m) \Rightarrow (g(u) < n))) \Rightarrow  \\ \quad \quad \operatorname{series\text{-}partial\text{-}sum}((\lambda k \in \mathbb{N}.\; a(g(k))), m) \\ \quad \quad =\; \operatorname{series\text{-}partial\text{-}sum}((\lambda k \in \mathbb{N}.\; \left(\text{if }\exists t \in \mathbb{N}.\; ((t < m) \wedge (g(t) = k))\text{ then }a(k)\text{ else }0\right)), n) \end{array}}}
\vnbprop{surj-covers}{\vnbstmt{\begin{array}{@{}l@{}} \forall n \in \mathbb{N}, g \in (\mathbb{N} \to \mathbb{N}). \\ \quad \forall w \in \mathbb{N}.\; \exists t \in \mathbb{N}.\; (g(t) = w) \Rightarrow  \\ \quad \quad \exists m \in \mathbb{N}. \forall j \in \mathbb{N}. \\ \quad \quad \quad (j < n) \Rightarrow \exists t \in \mathbb{N}.\; ((t < m) \wedge (g(t) = j)) \end{array}}}
\vnbprop{finite-values-bounded}{\vnbstmt{\begin{array}{@{}l@{}} \forall m \in \mathbb{N}, g \in (\mathbb{N} \to \mathbb{N}). \exists n \in \mathbb{N}. \forall u \in \mathbb{N}. \\ \quad (u < m) \Rightarrow (g(u) < n) \end{array}}}
\vnbprop{rearrangement-abs}{\vnbstmt{\begin{array}{@{}l@{}} \forall a \in (\mathbb{N} \to \mathbb{C}), g \in (\mathbb{N} \to \mathbb{N}). \\ \quad \forall u, v \in \mathbb{N}.\; ((g(u) = g(v)) \Rightarrow (u = v)) \wedge \\ \quad \operatorname{series\text{-}converges}((\lambda k \in \mathbb{N}.\; \operatorname{magnitude}(a(k)))) \Rightarrow \\ \quad \quad \operatorname{series\text{-}converges}((\lambda k \in \mathbb{N}.\; \operatorname{magnitude}(a(g(k))))) \end{array}}}
\vnbprop{rearrangement-value}{\vnbstmt{\begin{array}{@{}l@{}} \forall a \in (\mathbb{N} \to \mathbb{C}), g \in (\mathbb{N} \to \mathbb{N}). \\ \quad \forall u, v \in \mathbb{N}.\; ((g(u) = g(v)) \Rightarrow (u = v)) \wedge \\ \quad \forall w \in \mathbb{N}.\; \exists t \in \mathbb{N}.\; (g(t) = w) \wedge \\ \quad \operatorname{series\text{-}converges}((\lambda k \in \mathbb{N}.\; \operatorname{magnitude}(a(k)))) \Rightarrow \\ \quad \quad \operatorname{cc\text{-}series\text{-}converges\text{-}to}((\lambda k \in \mathbb{N}.\; a(g(k))), \operatorname{cc\text{-}series\text{-}limit}(a)) \end{array}}}
\vnbprop{pps-agree-below}{\vnbstmt{\begin{array}{@{}l@{}} \forall n \in \mathbb{N}, f, g \in (\mathbb{N} \to \mathbb{C}). \\ \quad \forall j \in \mathbb{N}.\; ((j < n) \Rightarrow (f(j) = g(j))) \Rightarrow (\operatorname{pprod}(f, n) = \operatorname{pprod}(g, n)) \end{array}}}
\vnbprop{pps-mul}{\vnbstmt{\begin{array}{@{}l@{}} \forall n \in \mathbb{N}, f, g \in (\mathbb{N} \to \mathbb{C}). \\ \quad \operatorname{pprod}((\lambda k \in \mathbb{N}.\; (f(k) \cdot g(k))), n) = (\operatorname{pprod}(f, n) \cdot \operatorname{pprod}(g, n)) \end{array}}}
\vnbprop{pps-delta}{\vnbstmt{\begin{array}{@{}l@{}} \forall n, s \in \mathbb{N}, c \in \mathbb{C}. \\ \quad (s < n) \Rightarrow (\operatorname{pprod}((\lambda k \in \mathbb{N}.\; \left(\text{if }(k = s)\text{ then }c\text{ else }1\right)), n) = c) \wedge \\ \quad (n \leq s) \Rightarrow (\operatorname{pprod}((\lambda k \in \mathbb{N}.\; \left(\text{if }(k = s)\text{ then }c\text{ else }1\right)), n) = 1) \end{array}}}
\vnbprop{pps-trunc}{\vnbstmt{\begin{array}{@{}l@{}} \forall p, n \in \mathbb{N}, a \in (\mathbb{N} \to \mathbb{C}). \\ \quad \operatorname{pprod}((\lambda k \in \mathbb{N}.\; \left(\text{if }(k < n)\text{ then }a(k)\text{ else }1\right)), (n + p)) = \operatorname{pprod}(a, n) \end{array}}}
\vnbprop{pps-mask}{\vnbstmt{\begin{array}{@{}l@{}} \forall m, n \in \mathbb{N}, a \in (\mathbb{N} \to \mathbb{C}), g \in (\mathbb{N} \to \mathbb{N}). \\ \quad (\forall u, v \in \mathbb{N}.\; ((g(u) = g(v)) \Rightarrow (u = v)) \wedge \forall u \in \mathbb{N}.\; ((u < m) \Rightarrow (g(u) < n))) \Rightarrow  \\ \quad \quad \operatorname{pprod}((\lambda k \in \mathbb{N}.\; a(g(k))), m) \\ \quad \quad =\; \operatorname{pprod}((\lambda k \in \mathbb{N}.\; \left(\text{if }\exists t \in \mathbb{N}.\; ((t < m) \wedge (g(t) = k))\text{ then }a(k)\text{ else }1\right)), n) \end{array}}}
\vnbprop{prod-partial-converges}{\vnbstmt{\begin{array}{@{}l@{}} \forall f \in (\mathbb{N} \to \mathbb{C}), l. \\ \quad \operatorname{prod\text{-}converges\text{-}to}(f, l) \Rightarrow  \\ \quad \quad \operatorname{converges\text{-}to}(\operatorname{cc\text{-}ms}, (\lambda n \in \mathbb{N}.\; \operatorname{pprod}(f, n)), l) \end{array}}}
\vnbprop{prod-rearrangement-limit}{\vnbstmt{\begin{array}{@{}l@{}} \forall a \in (\mathbb{N} \to \mathbb{C}), g \in (\mathbb{N} \to \mathbb{N}), l. \\ \quad \forall u, v \in \mathbb{N}.\; ((g(u) = g(v)) \Rightarrow (u = v)) \wedge \\ \quad \forall w \in \mathbb{N}.\; \exists t \in \mathbb{N}.\; (g(t) = w) \wedge \\ \quad \operatorname{series\text{-}converges}((\lambda k \in \mathbb{N}.\; \operatorname{magnitude}(a(k)))) \wedge \\ \quad \operatorname{prod\text{-}converges\text{-}to}((\lambda k \in \mathbb{N}.\; (1 + a(k))), l) \Rightarrow \\ \quad \quad \operatorname{converges\text{-}to}(\operatorname{cc\text{-}ms}, (\lambda n \in \mathbb{N}.\; \operatorname{pprod}((\lambda k \in \mathbb{N}.\; (1 + a(g(k)))), n)), l) \end{array}}}
\vnbprop{corollary-6-9}{\vnbstmt{\begin{array}{@{}l@{}} \forall a \in (\mathbb{N} \to \mathbb{C}), g \in (\mathbb{N} \to \mathbb{N}). \\ \quad \forall u, v \in \mathbb{N}.\; ((g(u) = g(v)) \Rightarrow (u = v)) \wedge \\ \quad \forall w \in \mathbb{N}.\; \exists t \in \mathbb{N}.\; (g(t) = w) \wedge \\ \quad \operatorname{prod\text{-}abs\text{-}converges}(a) \Rightarrow \\ \quad \quad \operatorname{prod\text{-}abs\text{-}converges}((\lambda k \in \mathbb{N}.\; a(g(k)))) \wedge \\ \quad \quad \forall l. \\ \quad \quad \quad \operatorname{prod\text{-}converges\text{-}to}((\lambda k \in \mathbb{N}.\; (1 + a(k))), l) \Rightarrow  \\ \quad \quad \quad \quad \operatorname{prod\text{-}converges\text{-}to}((\lambda k \in \mathbb{N}.\; (1 + (\lambda k \in \mathbb{N}.\; a(g(k)))(k))), l) \end{array}}}
\vnbprop{prod-value-unique}{\vnbstmt{\begin{array}{@{}l@{}} \forall f \in (\mathbb{N} \to \mathbb{C}), l, m. \\ \quad (\operatorname{prod\text{-}converges\text{-}to}(f, l) \wedge \operatorname{prod\text{-}converges\text{-}to}(f, m)) \Rightarrow (l = m) \end{array}}}
\vnbprop{prod-value-spec}{\vnbstmt{\begin{array}{@{}l@{}} \forall f \in (\mathbb{N} \to \mathbb{C}). \\ \quad \operatorname{prod\text{-}converges}(f) \Rightarrow  \\ \quad \quad (\operatorname{prod\text{-}value}(f) \in \mathbb{C}) \wedge \operatorname{prod\text{-}converges\text{-}to}(f, \operatorname{prod\text{-}value}(f)) \end{array}}}
\vnbprop{normal-on-compacta-pointwise}{\vnbstmt{\begin{array}{@{}l@{}} \forall u.; \forall f \in (\mathbb{N} \to (u \to \mathbb{C})), x \in u. \\ \quad (\operatorname{is\text{-}open}(\operatorname{cc\text{-}ms}, u) \wedge \operatorname{prod\text{-}normal\text{-}on\text{-}compacta}(u, f)) \Rightarrow  \\ \quad \quad \operatorname{series\text{-}converges}((\lambda k \in \mathbb{N}.\; \operatorname{magnitude}((\lambda k \in \mathbb{N}.\; f(k)(x))(k)))) \end{array}}}
\vnbprop{prop-6-11}{\vnbstmt{\begin{array}{@{}l@{}} \forall u.; \forall f \in (\mathbb{N} \to (u \to \mathbb{C})), x \in u. \\ \quad (\operatorname{is\text{-}open}(\operatorname{cc\text{-}ms}, u) \wedge \operatorname{prod\text{-}normal\text{-}on\text{-}compacta}(u, f)) \Rightarrow  \\ \quad \quad \operatorname{prod\text{-}converges}((\lambda k \in \mathbb{N}.\; (1 + f(k)(x)))) \wedge \\ \quad \quad (\operatorname{prod\text{-}value}((\lambda k \in \mathbb{N}.\; (1 + f(k)(x)))) = 0) \Leftrightarrow \exists q \in \mathbb{N}.\; ((1 + f(q)(x)) = 0) \end{array}}}
\vnbprop{sps-tail-dominated}{\vnbstmt{\begin{array}{@{}l@{}} \forall p, n \in \mathbb{N}, f, g \in (\mathbb{N} \to \mathbb{R}). \\ \quad \forall j \in \mathbb{N}.\; (f(j) \leq g(j)) \Rightarrow  \\ \quad \quad (\operatorname{series\text{-}partial\text{-}sum}(f, (n + p)) - \operatorname{series\text{-}partial\text{-}sum}(f, n)) \\ \quad \quad \leq\; (\operatorname{series\text{-}partial\text{-}sum}(g, (n + p)) - \operatorname{series\text{-}partial\text{-}sum}(g, n)) \end{array}}}
\vnbprop{prop-6-12-estimate}{\vnbstmt{\begin{array}{@{}l@{}} \forall u.; \forall f \in (\mathbb{N} \to (u \to \mathbb{C})), n \in \mathbb{N}, e. \\ \quad (\operatorname{is\text{-}open}(\operatorname{cc\text{-}ms}, u) \wedge (\operatorname{prod\text{-}normal\text{-}on\text{-}compacta}(u, f) \wedge \operatorname{pos\text{-}rr}(e))) \Rightarrow  \\ \quad \quad \exists m \in \mathbb{N}. \forall j \in \mathbb{N}, z \in \operatorname{cc\text{-}exhaust}(u, n). \\ \quad \quad \quad (m \leq j) \Rightarrow  \\ \quad \quad \quad \quad \operatorname{magnitude}((\operatorname{pprod}((\lambda k \in \mathbb{N}.\; (1 + f(k)(z))), j) - \operatorname{prod\text{-}value}((\lambda k \in \mathbb{N}.\; (1 + f(k)(z)))))) \\ \quad \quad \quad \quad \leq\; e \end{array}}}
\vnbprop{prop-6-12}{\vnbstmt{\begin{array}{@{}l@{}} \forall u.; \forall f \in (\mathbb{N} \to (u \to \mathbb{C})). \\ \quad (\operatorname{is\text{-}open}(\operatorname{cc\text{-}ms}, u) \wedge \operatorname{prod\text{-}normal\text{-}on\text{-}compacta}(u, f)) \Rightarrow  \\ \quad \quad \operatorname{prod\text{-}fun}(u, f) \in (u \to \mathbb{C}) \wedge \\ \quad \quad \operatorname{unif\text{-}on\text{-}compacta}(u, \\ \quad \quad \quad (\lambda n \in \mathbb{N}.\; (\lambda x \in u.\; \operatorname{pprod}((\lambda k \in \mathbb{N}.\; (1 + f(k)(x))), n))), \\ \quad \quad \quad \operatorname{prod\text{-}fun}(u, f)) \end{array}}}
\vnbprop{ppx-holomorphic}{\vnbstmt{\begin{array}{@{}l@{}} \forall n \in \mathbb{N}, u.; \forall f \in (\mathbb{N} \to (u \to \mathbb{C})). \\ \quad (\operatorname{is\text{-}open}(\operatorname{nf\text{-}metric\text{-}space}(\operatorname{cc\text{-}normed\text{-}field}), u) \wedge \forall i \in \mathbb{N}.\; \operatorname{holomorphic\text{-}on}(u, f(i))) \Rightarrow  \\ \quad \quad \operatorname{holomorphic\text{-}on}(u, (\lambda x \in u.\; \operatorname{pprod}((\lambda k \in \mathbb{N}.\; (1 + f(k)(x))), n))) \end{array}}}
\vnbprop{prod-fun-holomorphic}{\vnbstmt{\begin{array}{@{}l@{}} \forall u.; \forall f \in (\mathbb{N} \to (u \to \mathbb{C})). \\ \quad (\operatorname{is\text{-}open}(\operatorname{cc\text{-}ms}, u) \wedge (\forall i \in \mathbb{N}.\; \operatorname{holomorphic\text{-}on}(u, f(i)) \wedge \operatorname{prod\text{-}normal\text{-}on\text{-}compacta}(u, f))) \Rightarrow  \\ \quad \quad \operatorname{holomorphic\text{-}on}(u, \operatorname{prod\text{-}fun}(u, f)) \end{array}}}
\vnbprop{zeros-not-accumulating}{\vnbstmt{\begin{array}{@{}l@{}} \forall u, g, c. \\ \quad \neg \operatorname{zeros\text{-}accumulate\text{-}at}(u, g, c) \Rightarrow  \\ \quad \quad \exists r. \\ \quad \quad \quad (\operatorname{pos\text{-}rr}(r) \wedge \forall w \in u.\; (((w \in \operatorname{ball}(\operatorname{cc\text{-}ms}, c, r)) \wedge \neg (w = c)) \Rightarrow \neg (g(w) = 0))) \end{array}}}
\vnbprop{finite-zero-cover}{\vnbstmt{\begin{array}{@{}l@{}} \forall n \in \mathbb{N}, u.; \forall f \in (\mathbb{N} \to (u \to \mathbb{C})), c \in u, r. \\ \quad \operatorname{is\text{-}open}(\operatorname{nf\text{-}metric\text{-}space}(\operatorname{cc\text{-}normed\text{-}field}), u) \wedge \\ \quad \operatorname{is\text{-}connected}(\operatorname{cc\text{-}ms}, u) \wedge \\ \quad \forall i \in \mathbb{N}.\; \operatorname{holomorphic\text{-}on}(u, f(i)) \wedge \\ \quad \operatorname{pos\text{-}rr}(r) \wedge \\ \quad \operatorname{ball}(\operatorname{cc\text{-}ms}, c, r) \subseteq u \wedge \\ \quad \forall x \in \operatorname{ball}(\operatorname{cc\text{-}ms}, c, r). \exists q \in \mathbb{N}. \\ \quad \quad (q < n) \wedge ((1 + f(q)(x)) = 0) \Rightarrow \\ \quad \quad \exists q \in \mathbb{N}.\; ((q < n) \wedge \forall z \in u.\; ((1 + f(q)(z)) = 0)) \end{array}}}
\vnbprop{nonneg-series-terms-small}{\vnbstmt{\begin{array}{@{}l@{}} \forall g \in (\mathbb{N} \to \mathbb{R}), e. \\ \quad (\forall i \in \mathbb{N}.\; (0 \leq g(i)) \wedge (\operatorname{series\text{-}converges}(g) \wedge \operatorname{pos\text{-}rr}(e))) \Rightarrow  \\ \quad \quad \exists m \in \mathbb{N}.\; \forall q \in \mathbb{N}.\; ((m \leq q) \Rightarrow (g(q) \leq e)) \end{array}}}
\vnbprop{corollary-6-13}{\vnbstmt{\begin{array}{@{}l@{}} \forall u.; \forall f \in (\mathbb{N} \to (u \to \mathbb{C})). \\ \quad \operatorname{is\text{-}open}(\operatorname{cc\text{-}ms}, u) \wedge \\ \quad \operatorname{is\text{-}connected}(\operatorname{cc\text{-}ms}, u) \wedge \\ \quad \forall i \in \mathbb{N}.\; \operatorname{holomorphic\text{-}on}(u, f(i)) \wedge \\ \quad \forall i \in \mathbb{N}.\; \exists w \in u.\; \neg ((1 + f(i)(w)) = 0) \wedge \\ \quad \operatorname{prod\text{-}normal\text{-}on\text{-}compacta}(u, f) \Rightarrow \\ \quad \quad (\operatorname{holomorphic\text{-}on}(u, \operatorname{prod\text{-}fun}(u, f)) \wedge \exists v \in u.\; \neg (\operatorname{prod\text{-}fun}(u, f)(v) = 0)) \end{array}}}

\vnbbookfile{theorem-library/ch6-fun}{theorem-library/ch6-fun}
\noindent{\small ch6-fun.scm -- WAVE 2 (2026-10-27..30), wired 2026-10-30 from scratchpad/ch6/ch6-fun.scm. Library: complex-analysis.}\par
\vnbprop{ppx-continuous}{\vnbstmt{\begin{array}{@{}l@{}} \forall n \in \mathbb{N}, u.; \forall f \in (\mathbb{N} \to (u \to \mathbb{C})). \\ \quad ((u \subseteq \mathbb{C}) \wedge \forall i \in \mathbb{N}.\; (f(i) \in \operatorname{cc\text{-}cont\text{-}fns}(u))) \Rightarrow  \\ \quad \quad (\lambda x \in u.\; \operatorname{pprod}((\lambda k \in \mathbb{N}.\; (1 + f(k)(x))), n)) \in \operatorname{cc\text{-}cont\text{-}fns}(u) \end{array}}}
\vnbprop{prop-6-12-continuous}{\vnbstmt{\begin{array}{@{}l@{}} \forall u.; \forall f \in (\mathbb{N} \to (u \to \mathbb{C})). \\ \quad (\operatorname{is\text{-}open}(\operatorname{cc\text{-}ms}, u) \wedge (\forall i \in \mathbb{N}.\; (f(i) \in \operatorname{cc\text{-}cont\text{-}fns}(u)) \wedge \operatorname{prod\text{-}normal\text{-}on\text{-}compacta}(u, f))) \Rightarrow  \\ \quad \quad \operatorname{prod\text{-}fun}(u, f) \in \operatorname{cc\text{-}cont\text{-}fns}(u) \end{array}}}
\vnbprop{rseries-tail-converges}{\vnbstmt{\begin{array}{@{}l@{}} \forall m \in \mathbb{N}, g, t \in (\mathbb{N} \to \mathbb{R}). \\ \quad (\forall i \in \mathbb{N}.\; (0 \leq g(i)) \wedge (\forall j \in \mathbb{N}.\; (t(j) = g((m + j))) \wedge \operatorname{series\text{-}converges}(g))) \Rightarrow  \\ \quad \quad \operatorname{series\text{-}converges}(t) \end{array}}}
\vnbprop{rseries-untail-converges}{\vnbstmt{\begin{array}{@{}l@{}} \forall m \in \mathbb{N}, g, t \in (\mathbb{N} \to \mathbb{R}). \\ \quad (\forall i \in \mathbb{N}.\; (0 \leq g(i)) \wedge (\forall j \in \mathbb{N}.\; (t(j) = g((m + j))) \wedge \operatorname{series\text{-}converges}(t))) \Rightarrow  \\ \quad \quad \operatorname{series\text{-}converges}(g) \end{array}}}
\vnbprop{cc-log-one-plus-lower}{\vnbstmt{\begin{array}{@{}l@{}} \forall z \in \mathbb{C}. \\ \quad (\neg ((\operatorname{imag\text{-}part}((1 + z)) = 0) \wedge (\operatorname{real\text{-}part}((1 + z)) \leq 0)) \wedge (\operatorname{magnitude}(\log\left((1 + z)\right)) \leq 1/2)) \Rightarrow  \\ \quad \quad \operatorname{magnitude}(z) \leq (2 \cdot \operatorname{magnitude}(\log\left((1 + z)\right))) \end{array}}}
\vnbprop{prop-6-14-forward}{\vnbstmt{\begin{array}{@{}l@{}} \forall u.; \forall f \in (\mathbb{N} \to (u \to \mathbb{C})), n \in \mathbb{N}. \\ \quad (\operatorname{is\text{-}open}(\operatorname{cc\text{-}ms}, u) \wedge \operatorname{is\text{-}normally\text{-}convergent}(f, \operatorname{cc\text{-}exhaust}(u, n))) \Rightarrow  \\ \quad \quad \exists m \in \mathbb{N}. \\ \quad \quad \quad \forall q \in \mathbb{N}, w \in \operatorname{cc\text{-}exhaust}(u, n). \\ \quad \quad \quad \quad (m \leq q) \Rightarrow  \\ \quad \quad \quad \quad \quad \neg ((\operatorname{imag\text{-}part}((1 + f(q)(w))) = 0) \wedge (\operatorname{real\text{-}part}((1 + f(q)(w))) \leq 0)) \wedge \\ \quad \quad \quad \operatorname{is\text{-}normally\text{-}convergent}((\lambda k \in \mathbb{N}.\; (\lambda x \in \operatorname{cc\text{-}exhaust}(u, n).\; \log\left((1 + f((m + k))(x))\right))), \\ \quad \quad \quad \quad \operatorname{cc\text{-}exhaust}(u, n)) \end{array}}}
\vnbprop{prop-6-14-backward}{\vnbstmt{\begin{array}{@{}l@{}} \forall u.; \forall f \in (\mathbb{N} \to (u \to \mathbb{C})), n, m \in \mathbb{N}. \\ \quad \operatorname{is\text{-}open}(\operatorname{cc\text{-}ms}, u) \wedge \\ \quad \forall i \in \mathbb{N}.\; (f(i) \in \operatorname{cc\text{-}cont\text{-}fns}(u)) \wedge \\ \quad \forall q \in \mathbb{N}, w \in \operatorname{cc\text{-}exhaust}(u, n). \\ \quad \quad (m \leq q) \Rightarrow  \\ \quad \quad \quad \neg ((\operatorname{imag\text{-}part}((1 + f(q)(w))) = 0) \wedge (\operatorname{real\text{-}part}((1 + f(q)(w))) \leq 0)) \wedge \\ \quad \operatorname{is\text{-}normally\text{-}convergent}((\lambda k \in \mathbb{N}.\; (\lambda x \in \operatorname{cc\text{-}exhaust}(u, n).\; \log\left((1 + f((m + k))(x))\right))), \\ \quad \quad \operatorname{cc\text{-}exhaust}(u, n)) \Rightarrow \\ \quad \quad \operatorname{is\text{-}normally\text{-}convergent}(f, \operatorname{cc\text{-}exhaust}(u, n)) \end{array}}}
\vnbprop{prop-6-14}{\vnbstmt{\begin{array}{@{}l@{}} \forall u.; \forall f \in (\mathbb{N} \to (u \to \mathbb{C})). \\ \quad (\operatorname{is\text{-}open}(\operatorname{cc\text{-}ms}, u) \wedge \forall i \in \mathbb{N}.\; (f(i) \in \operatorname{cc\text{-}cont\text{-}fns}(u))) \Rightarrow  \\ \quad \quad \operatorname{prod\text{-}normal\text{-}on\text{-}compacta}(u, f) \\ \quad \quad \Leftrightarrow\; \forall n \in \mathbb{N}. \exists m \in \mathbb{N}. \\ \quad \quad \quad \forall q \in \mathbb{N}, w \in \operatorname{cc\text{-}exhaust}(u, n). \\ \quad \quad \quad \quad (m \leq q) \Rightarrow  \\ \quad \quad \quad \quad \quad \neg ((\operatorname{imag\text{-}part}((1 + f(q)(w))) = 0) \wedge (\operatorname{real\text{-}part}((1 + f(q)(w))) \leq 0)) \wedge \\ \quad \quad \quad \operatorname{is\text{-}normally\text{-}convergent}((\lambda k \in \mathbb{N}.\; (\lambda x \in \operatorname{cc\text{-}exhaust}(u, n).\; \log\left((1 + f((m + k))(x))\right))), \\ \quad \quad \quad \quad \operatorname{cc\text{-}exhaust}(u, n)) \end{array}}}
\vnbprop{corollary-6-16}{\vnbstmt{\begin{array}{@{}l@{}} \forall u.; \forall f \in (\mathbb{N} \to (u \to \mathbb{C})). \\ \quad \operatorname{is\text{-}open}(\operatorname{cc\text{-}ms}, u) \wedge \\ \quad \forall i \in \mathbb{N}.\; (f(i) \in \operatorname{cc\text{-}cont\text{-}fns}(u)) \wedge \\ \quad \forall n \in \mathbb{N}. \exists m \in \mathbb{N}. \\ \quad \quad \forall q \in \mathbb{N}, w \in \operatorname{cc\text{-}exhaust}(u, n). \\ \quad \quad \quad (m \leq q) \Rightarrow  \\ \quad \quad \quad \quad \neg ((\operatorname{imag\text{-}part}((1 + f(q)(w))) = 0) \wedge (\operatorname{real\text{-}part}((1 + f(q)(w))) \leq 0)) \wedge \\ \quad \quad \operatorname{is\text{-}normally\text{-}convergent}((\lambda k \in \mathbb{N}.\; (\lambda x \in \operatorname{cc\text{-}exhaust}(u, n).\; \log\left((1 + f((m + k))(x))\right))), \\ \quad \quad \quad \operatorname{cc\text{-}exhaust}(u, n)) \Rightarrow \\ \quad \quad \operatorname{prod\text{-}normal\text{-}on\text{-}compacta}(u, f) \end{array}}}
\vnbprop{cc-limit-quotient}{\vnbstmt{\begin{array}{@{}l@{}} \forall a, b \in (\mathbb{N} \to \mathbb{C}), p, q \in \mathbb{C}. \\ \quad \neg (q = 0) \wedge \\ \quad \forall i \in \mathbb{N}.\; \neg (b(i) = 0) \wedge \\ \quad \operatorname{converges\text{-}to}(\operatorname{cc\text{-}ms}, a, p) \wedge \\ \quad \operatorname{converges\text{-}to}(\operatorname{cc\text{-}ms}, b, q) \Rightarrow \\ \quad \quad \operatorname{converges\text{-}to}(\operatorname{cc\text{-}ms}, (\lambda k \in \mathbb{N}.\; (a(k) \cdot \operatorname{recip}(b(k)))), (p \cdot \operatorname{recip}(q))) \end{array}}}
\vnbprop{cc-deriv-mul-ptwise}{\vnbstmt{\begin{array}{@{}l@{}} \forall u, p, g.; \forall h \in (u \to \mathbb{C}), z \in u. \\ \quad \operatorname{is\text{-}open}(\operatorname{nf\text{-}metric\text{-}space}(\operatorname{cc\text{-}normed\text{-}field}), u) \wedge \\ \quad \operatorname{holomorphic\text{-}on}(u, p) \wedge \\ \quad \operatorname{holomorphic\text{-}on}(u, g) \wedge \\ \quad \forall y \in u.\; (h(y) \simeq (p(y) \cdot g(y))) \Rightarrow \\ \quad \quad \operatorname{deriv\text{-}on}(\operatorname{cc\text{-}normed\text{-}field}, u, h, z) \\ \quad \quad =\; (\operatorname{deriv\text{-}on}(\operatorname{cc\text{-}normed\text{-}field}, u, p, z) \cdot g(z)) + (p(z) \cdot \operatorname{deriv\text{-}on}(\operatorname{cc\text{-}normed\text{-}field}, u, g, z)) \end{array}}}
\vnbprop{ppx-log-deriv}{\vnbstmt{\begin{array}{@{}l@{}} \forall n \in \mathbb{N}, u.; \forall f \in (\mathbb{N} \to (u \to \mathbb{C})), z \in u. \\ \quad \operatorname{is\text{-}open}(\operatorname{nf\text{-}metric\text{-}space}(\operatorname{cc\text{-}normed\text{-}field}), u) \wedge \\ \quad \forall i \in \mathbb{N}.\; \operatorname{holomorphic\text{-}on}(u, f(i)) \wedge \\ \quad \forall q \in \mathbb{N}.\; \neg ((1 + f(q)(z)) = 0) \Rightarrow \\ \quad \quad \operatorname{deriv\text{-}on}(\operatorname{cc\text{-}normed\text{-}field}, u, (\lambda x \in u.\; \operatorname{pprod}((\lambda k \in \mathbb{N}.\; (1 + f(k)(x))), n)), z) \\ \quad \quad =\; \operatorname{pprod}((\lambda k \in \mathbb{N}.\; (1 + f(k)(z))), n) \cdot \operatorname{cc\text{-}series\text{-}partial\text{-}sum}((\lambda k \in \mathbb{N}.\; (\operatorname{deriv\text{-}on}(\operatorname{cc\text{-}normed\text{-}field}, u, (\lambda x \in u.\; (1 + f(k)(x))), z) \cdot \operatorname{recip}((1 + f(k)(z))))), n) \end{array}}}
\vnbprop{prop-6-17}{\vnbstmt{\begin{array}{@{}l@{}} \forall u.; \forall f \in (\mathbb{N} \to (u \to \mathbb{C})), z \in u. \\ \quad \operatorname{is\text{-}open}(\operatorname{cc\text{-}ms}, u) \wedge \\ \quad \forall i \in \mathbb{N}.\; \operatorname{holomorphic\text{-}on}(u, f(i)) \wedge \\ \quad \operatorname{prod\text{-}normal\text{-}on\text{-}compacta}(u, f) \wedge \\ \quad \neg (\operatorname{prod\text{-}fun}(u, f)(z) = 0) \Rightarrow \\ \quad \quad \operatorname{cc\text{-}series\text{-}converges\text{-}to}((\lambda k \in \mathbb{N}.\; \\ \quad \quad \quad \quad \operatorname{deriv\text{-}on}(\operatorname{cc\text{-}normed\text{-}field}, u, (\lambda x \in u.\; (1 + f(k)(x))), z) \cdot \operatorname{recip}((1 + f(k)(z)))), \\ \quad \quad \quad \operatorname{deriv\text{-}on}(\operatorname{cc\text{-}normed\text{-}field}, u, \operatorname{prod\text{-}fun}(u, f), z) \cdot \operatorname{recip}(\operatorname{prod\text{-}fun}(u, f)(z))) \end{array}}}

\vnbbookfile{theorem-library/ch6-weier}{theorem-library/ch6-weier}
\noindent{\small ch6-weier.scm -- WAVE 2 (2026-10-27..30), wired 2026-10-30 from scratchpad/ch6/ch6-weier.scm. Library: complex-analysis.}\par
\vnbprop{cc-pow-lin-holomorphic}{\vnbstmt{\forall n \in \mathbb{N}, c \in \mathbb{C}.\; \operatorname{holomorphic\text{-}on}(\mathbb{C}, (\lambda x \in \mathbb{C}.\; \operatorname{power}((x \cdot c), n)))}}
\vnbprop{weier-q-holomorphic}{\vnbstmt{\begin{array}{@{}l@{}} \forall n \in \mathbb{N}, c \in \mathbb{C}. \\ \quad \operatorname{holomorphic\text{-}on}(\mathbb{C}, \\ \quad \quad (\lambda x \in \mathbb{C}.\; \operatorname{cc\text{-}series\text{-}partial\text{-}sum}((\lambda k \in \mathbb{N}.\; ({k}^{\ast\!-1} \cdot \operatorname{power}((x \cdot c), k))), n))) \end{array}}}
\vnbprop{weier-e-holomorphic}{\vnbstmt{\begin{array}{@{}l@{}} \forall n \in \mathbb{N}, c \in \mathbb{C}. \\ \quad \operatorname{holomorphic\text{-}on}(\mathbb{C}, \\ \quad \quad (\lambda x \in \mathbb{C}.\; \\ \quad \quad \quad (((1 - (x \cdot c)) \cdot \exp \operatorname{cc\text{-}series\text{-}partial\text{-}sum}((\lambda k \in \mathbb{N}.\; ({k}^{\ast\!-1} \cdot \operatorname{power}((x \cdot c), k))), (n + 1))) - 1))) \end{array}}}
\vnbprop{recip-star-le-one}{\vnbstmt{\forall k \in \mathbb{N}.\; ({k}^{\ast\!-1} \leq 1)}}
\vnbprop{log-series-tail-partial}{\vnbstmt{\begin{array}{@{}l@{}} \forall d, n \in \mathbb{N}, z \in \mathbb{C}. \\ \quad (\operatorname{magnitude}(z) \leq 1/2) \Rightarrow  \\ \quad \quad \operatorname{magnitude}((\operatorname{cc\text{-}series\text{-}partial\text{-}sum}((\lambda k \in \mathbb{N}.\; ({k}^{\ast\!-1} \cdot \operatorname{power}(z, k))), ((n + 1) + d)) - \operatorname{cc\text{-}series\text{-}partial\text{-}sum}((\lambda k \in \mathbb{N}.\; ({k}^{\ast\!-1} \cdot \operatorname{power}(z, k))), (n + 1)))) \\ \quad \quad \leq\; (2 \cdot (\operatorname{power}(\operatorname{magnitude}(z), (n + 1)) \cdot \operatorname{recip}((n + 1)))) \cdot (1 - \operatorname{power}(\operatorname{magnitude}(z), d)) \end{array}}}
\vnbprop{log-series-tail}{\vnbstmt{\begin{array}{@{}l@{}} \forall n \in \mathbb{N}, z \in \mathbb{C}. \\ \quad (\operatorname{magnitude}(z) \leq 1/2) \Rightarrow  \\ \quad \quad \operatorname{magnitude}((\log\left((1 - z)\right) + \operatorname{cc\text{-}series\text{-}partial\text{-}sum}((\lambda k \in \mathbb{N}.\; ({k}^{\ast\!-1} \cdot \operatorname{power}(z, k))), (n + 1)))) \\ \quad \quad \leq\; 2 \cdot (\operatorname{power}(\operatorname{magnitude}(z), (n + 1)) \cdot \operatorname{recip}((n + 1))) \end{array}}}
\vnbprop{cc-log-exp-small}{\vnbstmt{\begin{array}{@{}l@{}} \forall z \in \mathbb{C}. \\ \quad (\operatorname{magnitude}(z) \leq 1/2) \Rightarrow  \\ \quad \quad (\neg ((\operatorname{imag\text{-}part}(\exp z) = 0) \wedge (\operatorname{real\text{-}part}(\exp z) \leq 0)) \wedge (\log\left(\exp z\right) = z)) \end{array}}}
\vnbprop{weier-e-exp}{\vnbstmt{\begin{array}{@{}l@{}} \forall n \in \mathbb{N}, z \in \mathbb{C}. \\ \quad (\operatorname{magnitude}(z) \leq 1/2) \Rightarrow  \\ \quad \quad (1 - z) \cdot \exp \operatorname{cc\text{-}series\text{-}partial\text{-}sum}((\lambda k \in \mathbb{N}.\; ({k}^{\ast\!-1} \cdot \operatorname{power}(z, k))), (n + 1)) \\ \quad \quad =\; \exp (\log\left((1 - z)\right) + \operatorname{cc\text{-}series\text{-}partial\text{-}sum}((\lambda k \in \mathbb{N}.\; ({k}^{\ast\!-1} \cdot \operatorname{power}(z, k))), (n + 1))) \end{array}}}
\vnbprop{weier-e-zero-iff}{\vnbstmt{\begin{array}{@{}l@{}} \forall n \in \mathbb{N}, z \in \mathbb{C}. \\ \quad ((1 - z) \cdot \exp \operatorname{cc\text{-}series\text{-}partial\text{-}sum}((\lambda k \in \mathbb{N}.\; ({k}^{\ast\!-1} \cdot \operatorname{power}(z, k))), (n + 1))) = 0 \\ \quad \Leftrightarrow\; z = 1 \end{array}}}
\vnbprop{weier-factor-log-bound}{\vnbstmt{\begin{array}{@{}l@{}} \forall n \in \mathbb{N}, z \in \mathbb{C}. \\ \quad (\operatorname{magnitude}(z) \leq 1/4) \Rightarrow  \\ \quad \quad \neg ((\operatorname{imag\text{-}part}(((1 - z) \cdot \exp \operatorname{cc\text{-}series\text{-}partial\text{-}sum}((\lambda k \in \mathbb{N}.\; ({k}^{\ast\!-1} \cdot \operatorname{power}(z, k))), (n + 1)))) = 0) \wedge (\operatorname{real\text{-}part}(((1 - z) \cdot \exp \operatorname{cc\text{-}series\text{-}partial\text{-}sum}((\lambda k \in \mathbb{N}.\; ({k}^{\ast\!-1} \cdot \operatorname{power}(z, k))), (n + 1)))) \leq 0)) \wedge \\ \quad \quad \operatorname{magnitude}(\log\left(((1 - z) \cdot \exp \operatorname{cc\text{-}series\text{-}partial\text{-}sum}((\lambda k \in \mathbb{N}.\; ({k}^{\ast\!-1} \cdot \operatorname{power}(z, k))), (n + 1)))\right)) \\ \quad \quad \leq\; 2 \cdot (\operatorname{power}(\operatorname{magnitude}(z), (n + 1)) \cdot \operatorname{recip}((n + 1))) \end{array}}}
\vnbprop{weier-factor-zero-iff}{\vnbstmt{\begin{array}{@{}l@{}} \forall n \in \mathbb{N}, a, z \in \mathbb{C}. \\ \quad \neg (a = 0) \Rightarrow  \\ \quad \quad 1 + (((1 - (z \cdot \operatorname{recip}(a))) \cdot \exp \operatorname{cc\text{-}series\text{-}partial\text{-}sum}((\lambda k \in \mathbb{N}.\; ({k}^{\ast\!-1} \cdot \operatorname{power}((z \cdot \operatorname{recip}(a)), k))), (n + 1))) - 1) \\ \quad \quad =\; 0 \\ \quad \quad \Leftrightarrow\; z = a \end{array}}}
\vnbprop{weier-ratio-bound}{\vnbstmt{\begin{array}{@{}l@{}} \forall a \in \mathbb{C}, r \in \mathbb{R}, z \in \mathbb{C}. \\ \quad (\neg (a = 0) \wedge ((0 \leq r) \wedge (((4 \cdot (r + 1)) \leq \operatorname{magnitude}(a)) \wedge (\operatorname{magnitude}(z) \leq r)))) \Rightarrow  \\ \quad \quad ((\operatorname{magnitude}((z \cdot \operatorname{recip}(a))) \leq 1/4) \wedge (\operatorname{magnitude}((z \cdot \operatorname{recip}(a))) \leq (r \cdot \operatorname{recip}(\operatorname{magnitude}(a))))) \end{array}}}
\vnbprop{theorem-6-20-family}{\vnbstmt{\begin{array}{@{}l@{}} \forall a \in (\mathbb{N} \to \mathbb{C}), u \in (\mathbb{N} \to \mathbb{N}), f \in (\mathbb{N} \to (\mathbb{C} \to \mathbb{C})). \\ \quad \forall q \in \mathbb{N}.\; \neg (a(q) = 0) \wedge \\ \quad \forall r \in \mathbb{R}. \exists m \in \mathbb{N}. \forall q \in \mathbb{N}. \\ \quad \quad (m \leq q) \Rightarrow (r \leq \operatorname{magnitude}(a(q))) \wedge \\ \quad \forall r \in \mathbb{R}. \\ \quad \quad (0 \leq r) \Rightarrow  \\ \quad \quad \quad \operatorname{series\text{-}converges}((\lambda k \in \mathbb{N}.\; (\operatorname{power}((r \cdot \operatorname{recip}(\operatorname{magnitude}(a(k)))), (u(k) + 1)) \cdot \operatorname{recip}((u(k) + 1))))) \wedge \\ \quad \forall k \in \mathbb{N}, x \in \mathbb{C}. \\ \quad \quad f(k)(x) \\ \quad \quad =\; (((1 - (x \cdot \operatorname{recip}(a(k)))) \cdot \exp \operatorname{cc\text{-}series\text{-}partial\text{-}sum}((\lambda \mathit{clhk} \in \mathbb{N}.\; ({\mathit{clhk}}^{\ast\!-1} \cdot \operatorname{power}((x \cdot \operatorname{recip}(a(k))), \mathit{clhk}))), (u(k) + 1))) - 1) \Rightarrow \\ \quad \quad \operatorname{prod\text{-}normal\text{-}on\text{-}compacta}(\mathbb{C}, f) \wedge \\ \quad \quad \operatorname{holomorphic\text{-}on}(\mathbb{C}, \operatorname{prod\text{-}fun}(\mathbb{C}, f)) \wedge \\ \quad \quad \forall z \in \mathbb{C}.\; ((\operatorname{prod\text{-}fun}(\mathbb{C}, f)(z) = 0) \Leftrightarrow \exists q \in \mathbb{N}.\; (z = a(q))) \end{array}}}
\vnbprop{theorem-6-20}{\vnbstmt{\begin{array}{@{}l@{}} \forall a \in (\mathbb{N} \to \mathbb{C}), u \in (\mathbb{N} \to \mathbb{N}). \\ \quad \forall q \in \mathbb{N}.\; \neg (a(q) = 0) \wedge \\ \quad \forall r \in \mathbb{R}. \exists m \in \mathbb{N}. \forall q \in \mathbb{N}. \\ \quad \quad (m \leq q) \Rightarrow (r \leq \operatorname{magnitude}(a(q))) \wedge \\ \quad \forall r \in \mathbb{R}. \\ \quad \quad (0 \leq r) \Rightarrow  \\ \quad \quad \quad \operatorname{series\text{-}converges}((\lambda k \in \mathbb{N}.\; (\operatorname{power}((r \cdot \operatorname{recip}(\operatorname{magnitude}(a(k)))), (u(k) + 1)) \cdot \operatorname{recip}((u(k) + 1))))) \Rightarrow \\ \quad \quad \operatorname{prod\text{-}normal\text{-}on\text{-}compacta}(\mathbb{C}, \\ \quad \quad \quad (\lambda k \in \mathbb{N}.\; \\ \quad \quad \quad \quad (\lambda x \in \mathbb{C}.\; \\ \quad \quad \quad \quad \quad (((1 - (x \cdot \operatorname{recip}(a(k)))) \cdot \exp \operatorname{cc\text{-}series\text{-}partial\text{-}sum}((\lambda \mathit{clhk} \in \mathbb{N}.\; ({\mathit{clhk}}^{\ast\!-1} \cdot \operatorname{power}((x \cdot \operatorname{recip}(a(k))), \mathit{clhk}))), (u(k) + 1))) - 1)))) \wedge \\ \quad \quad \operatorname{holomorphic\text{-}on}(\mathbb{C}, \\ \quad \quad \quad \operatorname{prod\text{-}fun}(\mathbb{C}, \\ \quad \quad \quad \quad (\lambda k \in \mathbb{N}.\; \\ \quad \quad \quad \quad \quad (\lambda x \in \mathbb{C}.\; \\ \quad \quad \quad \quad \quad \quad (((1 - (x \cdot \operatorname{recip}(a(k)))) \cdot \exp \operatorname{cc\text{-}series\text{-}partial\text{-}sum}((\lambda \mathit{clhk} \in \mathbb{N}.\; ({\mathit{clhk}}^{\ast\!-1} \cdot \operatorname{power}((x \cdot \operatorname{recip}(a(k))), \mathit{clhk}))), (u(k) + 1))) - 1))))) \wedge \\ \quad \quad \forall z \in \mathbb{C}. \\ \quad \quad \quad \operatorname{prod\text{-}fun}(\mathbb{C}, (\lambda k \in \mathbb{N}.\; (\lambda x \in \mathbb{C}.\; (((1 - (x \cdot \operatorname{recip}(a(k)))) \cdot \exp \operatorname{cc\text{-}series\text{-}partial\text{-}sum}((\lambda \mathit{clhk} \in \mathbb{N}.\; ({\mathit{clhk}}^{\ast\!-1} \cdot \operatorname{power}((x \cdot \operatorname{recip}(a(k))), \mathit{clhk}))), (u(k) + 1))) - 1))))(z) \\ \quad \quad \quad =\; 0 \\ \quad \quad \quad \Leftrightarrow\; \exists q \in \mathbb{N}.\; (z = a(q)) \end{array}}}

\vnbbookfile{theorem-library/ch6-mult}{theorem-library/ch6-mult}
\noindent{\small ch6-mult.scm -- WAVE 2 (2026-10-27..30), wired 2026-10-30 from scratchpad/ch6/ch6-mult.scm. Library: complex-analysis.}\par
\vnbprop{cc-limit-shift-n}{\vnbstmt{\begin{array}{@{}l@{}} \forall n \in \mathbb{N}, s, h \in (\mathbb{N} \to \mathbb{C}), l \in \mathbb{C}. \\ \quad (\forall j \in \mathbb{N}.\; (h(j) = s((n + j))) \wedge \operatorname{converges\text{-}to}(\operatorname{cc\text{-}ms}, s, l)) \Rightarrow  \\ \quad \quad \operatorname{converges\text{-}to}(\operatorname{cc\text{-}ms}, h, l) \end{array}}}
\vnbprop{prod-value-split}{\vnbstmt{\begin{array}{@{}l@{}} \forall n \in \mathbb{N}, f \in (\mathbb{N} \to \mathbb{C}). \\ \quad (\operatorname{prod\text{-}converges}(f) \wedge \operatorname{prod\text{-}converges}((\lambda k \in \mathbb{N}.\; f((n + k))))) \Rightarrow  \\ \quad \quad \operatorname{prod\text{-}value}(f) = (\operatorname{pprod}(f, n) \cdot \operatorname{prod\text{-}value}((\lambda k \in \mathbb{N}.\; f((n + k))))) \end{array}}}
\vnbprop{prod-normal-tail}{\vnbstmt{\begin{array}{@{}l@{}} \forall n \in \mathbb{N}, u.; \forall f \in (\mathbb{N} \to (u \to \mathbb{C})). \\ \quad (\operatorname{is\text{-}open}(\operatorname{cc\text{-}ms}, u) \wedge \operatorname{prod\text{-}normal\text{-}on\text{-}compacta}(u, f)) \Rightarrow  \\ \quad \quad \operatorname{prod\text{-}normal\text{-}on\text{-}compacta}(u, (\lambda k \in \mathbb{N}.\; f((n + k)))) \end{array}}}
\vnbprop{weier-factor-order-one}{\vnbstmt{\begin{array}{@{}l@{}} \forall a \in (\mathbb{N} \to \mathbb{C}), u \in (\mathbb{N} \to \mathbb{N}), f \in (\mathbb{N} \to (\mathbb{C} \to \mathbb{C})), n \in \mathbb{N}. \\ \quad \forall q \in \mathbb{N}.\; \neg (a(q) = 0) \wedge \\ \quad \forall k \in \mathbb{N}, x \in \mathbb{C}. \\ \quad \quad f(k)(x) \\ \quad \quad =\; (((1 - (x \cdot \operatorname{recip}(a(k)))) \cdot \exp \operatorname{cc\text{-}series\text{-}partial\text{-}sum}((\lambda \mathit{clhk} \in \mathbb{N}.\; ({\mathit{clhk}}^{\ast\!-1} \cdot \operatorname{power}((x \cdot \operatorname{recip}(a(k))), \mathit{clhk}))), (u(k) + 1))) - 1) \Rightarrow \\ \quad \quad \operatorname{has\text{-}zero\text{-}of\text{-}order}(\mathbb{C}, (\lambda x \in \mathbb{C}.\; (1 + f(n)(x))), a(n), 1) \end{array}}}
\vnbprop{weier-pp-order}{\vnbstmt{\begin{array}{@{}l@{}} \forall n \in \mathbb{N}, a \in (\mathbb{N} \to \mathbb{C}), u \in (\mathbb{N} \to \mathbb{N}), f \in (\mathbb{N} \to (\mathbb{C} \to \mathbb{C})), z \in \mathbb{C}. \\ \quad \forall q \in \mathbb{N}.\; \neg (a(q) = 0) \wedge \\ \quad \forall i \in \mathbb{N}.\; \operatorname{holomorphic\text{-}on}(\mathbb{C}, f(i)) \wedge \\ \quad \forall k \in \mathbb{N}, x \in \mathbb{C}. \\ \quad \quad f(k)(x) \\ \quad \quad =\; (((1 - (x \cdot \operatorname{recip}(a(k)))) \cdot \exp \operatorname{cc\text{-}series\text{-}partial\text{-}sum}((\lambda \mathit{clhk} \in \mathbb{N}.\; ({\mathit{clhk}}^{\ast\!-1} \cdot \operatorname{power}((x \cdot \operatorname{recip}(a(k))), \mathit{clhk}))), (u(k) + 1))) - 1) \Rightarrow \\ \quad \quad \operatorname{has\text{-}zero\text{-}of\text{-}order}(\mathbb{C}, (\lambda x \in \mathbb{C}.\; \operatorname{pprod}((\lambda k \in \mathbb{N}.\; (1 + f(k)(x))), n)), \\ \quad \quad \quad z, \operatorname{count\text{-}hits}(a, z, n)) \end{array}}}
\vnbprop{theorem-6-20-multiplicity}{\vnbstmt{\begin{array}{@{}l@{}} \forall a \in (\mathbb{N} \to \mathbb{C}), u \in (\mathbb{N} \to \mathbb{N}), f \in (\mathbb{N} \to (\mathbb{C} \to \mathbb{C})), z \in \mathbb{C}, n \in \mathbb{N}. \\ \quad \forall q \in \mathbb{N}.\; \neg (a(q) = 0) \wedge \\ \quad \forall r \in \mathbb{R}. \exists m \in \mathbb{N}. \forall q \in \mathbb{N}. \\ \quad \quad (m \leq q) \Rightarrow (r \leq \operatorname{magnitude}(a(q))) \wedge \\ \quad \forall r \in \mathbb{R}. \\ \quad \quad (0 \leq r) \Rightarrow  \\ \quad \quad \quad \operatorname{series\text{-}converges}((\lambda k \in \mathbb{N}.\; (\operatorname{power}((r \cdot \operatorname{recip}(\operatorname{magnitude}(a(k)))), (u(k) + 1)) \cdot \operatorname{recip}((u(k) + 1))))) \wedge \\ \quad \forall k \in \mathbb{N}, x \in \mathbb{C}. \\ \quad \quad f(k)(x) \\ \quad \quad =\; (((1 - (x \cdot \operatorname{recip}(a(k)))) \cdot \exp \operatorname{cc\text{-}series\text{-}partial\text{-}sum}((\lambda \mathit{clhk} \in \mathbb{N}.\; ({\mathit{clhk}}^{\ast\!-1} \cdot \operatorname{power}((x \cdot \operatorname{recip}(a(k))), \mathit{clhk}))), (u(k) + 1))) - 1) \wedge \\ \quad \forall q \in \mathbb{N}.\; ((n \leq q) \Rightarrow \neg (a(q) = z)) \Rightarrow \\ \quad \quad \operatorname{has\text{-}zero\text{-}of\text{-}order}(\mathbb{C}, \operatorname{prod\text{-}fun}(\mathbb{C}, f), z, \operatorname{count\text{-}hits}(a, z, n)) \end{array}}}
\vnbprop{zero-factor-value-unique}{\vnbstmt{\begin{array}{@{}l@{}} \forall u, g, h.; \forall a \in u, m \in \mathbb{N}. \\ \quad \operatorname{is\text{-}open}(\operatorname{cc\text{-}ms}, u) \wedge \\ \quad \operatorname{holomorphic\text{-}on}(u, g) \wedge \\ \quad \operatorname{holomorphic\text{-}on}(u, h) \wedge \\ \quad \forall y \in u.\; ((\operatorname{power}((y - a), m) \cdot g(y)) = (\operatorname{power}((y - a), m) \cdot h(y))) \Rightarrow \\ \quad \quad g(a) = h(a) \end{array}}}

\vnbbookfile{theorem-library/ch6-quot}{theorem-library/ch6-quot}
\noindent{\small ch6-quot.scm -- WAVE 2 (2026-10-27..30), wired 2026-10-30 from scratchpad/ch6/ch6-quot.scm. Library: complex-analysis.}\par
\vnbprop{quot-value}{\vnbstmt{\begin{array}{@{}l@{}} \forall f, p.; \forall x \in \mathbb{C}, g, h. \\ \quad \operatorname{holomorphic\text{-}on}(\mathbb{C}, f) \wedge \\ \quad \operatorname{holomorphic\text{-}on}(\mathbb{C}, p) \wedge \\ \quad \operatorname{zero\text{-}order}(\mathbb{C}, p, x) \in \mathbb{N} \wedge \\ \quad \operatorname{holomorphic\text{-}on}(\mathbb{C}, g) \wedge \\ \quad \operatorname{holomorphic\text{-}on}(\mathbb{C}, h) \wedge \\ \quad \neg (h(x) = 0) \wedge \\ \quad \forall y \in \mathbb{C}.\; (f(y) = (\operatorname{power}((y - x), \operatorname{zero\text{-}order}(\mathbb{C}, p, x)) \cdot g(y))) \wedge \\ \quad \forall y \in \mathbb{C}.\; (p(y) = (\operatorname{power}((y - x), \operatorname{zero\text{-}order}(\mathbb{C}, p, x)) \cdot h(y))) \Rightarrow \\ \quad \quad (\iota\, v.\;((v \in \mathbb{C}) \wedge \exists \mathit{zqh}, k \in (\mathbb{C} \to \mathbb{C}).\; (\operatorname{holomorphic\text{-}on}(\mathbb{C}, \mathit{zqh}) \wedge (\operatorname{holomorphic\text{-}on}(\mathbb{C}, k) \wedge (\neg (k(x) = 0) \wedge (\forall y \in \mathbb{C}.\; (f(y) = (\operatorname{power}((y - x), \operatorname{zero\text{-}order}(\mathbb{C}, p, x)) \cdot \operatorname{zqh}(y))) \wedge (\forall y \in \mathbb{C}.\; (p(y) = (\operatorname{power}((y - x), \operatorname{zero\text{-}order}(\mathbb{C}, p, x)) \cdot k(y))) \wedge (v = (\operatorname{zqh}(x) \cdot \operatorname{recip}(k(x))))))))))) \\ \quad \quad =\; g(x) \cdot \operatorname{recip}(h(x)) \end{array}}}
\vnbprop{quot-holomorphic}{\vnbstmt{\begin{array}{@{}l@{}} \forall f, p. \\ \quad \operatorname{holomorphic\text{-}on}(\mathbb{C}, f) \wedge \\ \quad \operatorname{holomorphic\text{-}on}(\mathbb{C}, p) \wedge \\ \quad \exists w \in \mathbb{C}.\; \neg (p(w) = 0) \wedge \\ \quad \forall z \in \mathbb{C}. \exists m \in \mathbb{N}. \\ \quad \quad \operatorname{has\text{-}zero\text{-}of\text{-}order}(\mathbb{C}, f, z, m) \wedge \operatorname{has\text{-}zero\text{-}of\text{-}order}(\mathbb{C}, p, z, m) \Rightarrow \\ \quad \quad \operatorname{holomorphic\text{-}on}(\mathbb{C}, \\ \quad \quad \quad (\lambda x \in \mathbb{C}.\; \\ \quad \quad \quad \quad (\iota\, v.\;((v \in \mathbb{C}) \wedge \exists h, k \in (\mathbb{C} \to \mathbb{C}).\; (\operatorname{holomorphic\text{-}on}(\mathbb{C}, h) \wedge (\operatorname{holomorphic\text{-}on}(\mathbb{C}, k) \wedge (\neg (k(x) = 0) \wedge (\forall y \in \mathbb{C}.\; (f(y) = (\operatorname{power}((y - x), \operatorname{zero\text{-}order}(\mathbb{C}, p, x)) \cdot h(y))) \wedge (\forall y \in \mathbb{C}.\; (p(y) = (\operatorname{power}((y - x), \operatorname{zero\text{-}order}(\mathbb{C}, p, x)) \cdot k(y))) \wedge (v = (h(x) \cdot \operatorname{recip}(k(x))))))))))))) \end{array}}}
\vnbprop{entire-quotient-same-orders}{\vnbstmt{\begin{array}{@{}l@{}} \forall f, p. \\ \quad \operatorname{holomorphic\text{-}on}(\mathbb{C}, f) \wedge \\ \quad \operatorname{holomorphic\text{-}on}(\mathbb{C}, p) \wedge \\ \quad \exists w \in \mathbb{C}.\; \neg (p(w) = 0) \wedge \\ \quad \forall z \in \mathbb{C}. \exists m \in \mathbb{N}. \\ \quad \quad \operatorname{has\text{-}zero\text{-}of\text{-}order}(\mathbb{C}, f, z, m) \wedge \operatorname{has\text{-}zero\text{-}of\text{-}order}(\mathbb{C}, p, z, m) \Rightarrow \\ \quad \quad \exists q. \\ \quad \quad \quad (\operatorname{holomorphic\text{-}on}(\mathbb{C}, q) \wedge (\forall z \in \mathbb{C}.\; \neg (q(z) = 0) \wedge \forall z \in \mathbb{C}.\; (f(z) = (p(z) \cdot q(z))))) \end{array}}}
\vnbprop{weier-h167-mu-id}{\vnbstmt{\begin{array}{@{}l@{}} \forall a \in (\mathbb{N} \to \mathbb{C}), r \in \mathbb{R}. \\ \quad \forall q \in \mathbb{N}.\; \neg (a(q) = 0) \wedge \\ \quad \forall r \in \mathbb{R}. \exists m \in \mathbb{N}. \forall q \in \mathbb{N}. \\ \quad \quad (m \leq q) \Rightarrow (r \leq \operatorname{magnitude}(a(q))) \wedge \\ \quad 0 \leq r \Rightarrow \\ \quad \quad \operatorname{series\text{-}converges}((\lambda k \in \mathbb{N}.\; \\ \quad \quad \quad \quad \operatorname{power}((r \cdot \operatorname{recip}(\operatorname{magnitude}(a(k)))), ((\lambda \mathit{c6mk} \in \mathbb{N}.\; \mathit{c6mk})(k) + 1)) \cdot \operatorname{recip}(((\lambda \mathit{c6mk} \in \mathbb{N}.\; \mathit{c6mk})(k) + 1)))) \end{array}}}
\vnbprop{theorem-6-21}{\vnbstmt{\begin{array}{@{}l@{}} \forall f.; \forall m \in \mathbb{N}, a \in (\mathbb{N} \to \mathbb{C}). \\ \quad \operatorname{holomorphic\text{-}on}(\mathbb{C}, f) \wedge \\ \quad \operatorname{has\text{-}zero\text{-}of\text{-}order}(\mathbb{C}, f, 0, m) \wedge \\ \quad \forall q \in \mathbb{N}.\; \neg (a(q) = 0) \wedge \\ \quad \forall r \in \mathbb{R}. \exists \mathit{c6m} \in \mathbb{N}. \forall q \in \mathbb{N}. \\ \quad \quad (\mathit{c6m} \leq q) \Rightarrow (r \leq \operatorname{magnitude}(a(q))) \wedge \\ \quad \forall z \in \mathbb{C}, n \in \mathbb{N}. \\ \quad \quad (\neg (z = 0) \wedge \forall q \in \mathbb{N}.\; ((n \leq q) \Rightarrow \neg (a(q) = z))) \Rightarrow  \\ \quad \quad \quad \operatorname{has\text{-}zero\text{-}of\text{-}order}(\mathbb{C}, f, z, \operatorname{count\text{-}hits}(a, z, n)) \Rightarrow \\ \quad \quad \exists g. \\ \quad \quad \quad \operatorname{holomorphic\text{-}on}(\mathbb{C}, g) \wedge \\ \quad \quad \quad \exists u \in (\mathbb{N} \to \mathbb{N}). \\ \quad \quad \quad \quad \forall r \in \mathbb{R}. \\ \quad \quad \quad \quad \quad (0 \leq r) \Rightarrow  \\ \quad \quad \quad \quad \quad \quad \operatorname{series\text{-}converges}((\lambda k \in \mathbb{N}.\; (\operatorname{power}((r \cdot \operatorname{recip}(\operatorname{magnitude}(a(k)))), (u(k) + 1)) \cdot \operatorname{recip}((u(k) + 1))))) \wedge \\ \quad \quad \quad \quad \forall z \in \mathbb{C}. \\ \quad \quad \quad \quad \quad f(z) \\ \quad \quad \quad \quad \quad =\; (\operatorname{power}(z, m) \cdot \exp g(z)) \cdot \operatorname{prod\text{-}fun}(\mathbb{C}, (\lambda k \in \mathbb{N}.\; (\lambda x \in \mathbb{C}.\; (((1 - (x \cdot \operatorname{recip}(a(k)))) \cdot \exp \operatorname{cc\text{-}series\text{-}partial\text{-}sum}((\lambda \mathit{clhk} \in \mathbb{N}.\; ({\mathit{clhk}}^{\ast\!-1} \cdot \operatorname{power}((x \cdot \operatorname{recip}(a(k))), \mathit{clhk}))), (u(k) + 1))) - 1))))(z) \end{array}}}
\vnbprop{prop-6-22}{\vnbstmt{\begin{array}{@{}l@{}} \forall a \in (\mathbb{N} \to \mathbb{C}), t \in \mathbb{N}, f.; \forall m \in \mathbb{N}, g, h.; \forall p, r \in \mathbb{C}. \\ \quad \forall q \in \mathbb{N}.\; \neg (a(q) = 0) \wedge \\ \quad \forall r \in \mathbb{R}. \exists m \in \mathbb{N}. \forall q \in \mathbb{N}. \\ \quad \quad (m \leq q) \Rightarrow (r \leq \operatorname{magnitude}(a(q))) \wedge \\ \quad \forall r \in \mathbb{R}. \\ \quad \quad (0 \leq r) \Rightarrow  \\ \quad \quad \quad \operatorname{series\text{-}converges}((\lambda k \in \mathbb{N}.\; \\ \quad \quad \quad \quad \quad \operatorname{power}((r \cdot \operatorname{recip}(\operatorname{magnitude}(a(k)))), ((\lambda \mathit{c6tk} \in \mathbb{N}.\; t)(k) + 1)) \cdot \operatorname{recip}(((\lambda \mathit{c6tk} \in \mathbb{N}.\; t)(k) + 1)))) \wedge \\ \quad \operatorname{holomorphic\text{-}on}(\mathbb{C}, f) \wedge \\ \quad \exists w \in \mathbb{C}.\; \neg (f(w) = 0) \wedge \\ \quad \operatorname{holomorphic\text{-}on}(\mathbb{C}, g) \wedge \\ \quad \operatorname{holomorphic\text{-}on}(\mathbb{C}, h) \wedge \\ \quad \forall z \in \mathbb{C}. \\ \quad \quad f(z) \\ \quad \quad =\; (\operatorname{power}(z, m) \cdot \exp g(z)) \cdot \operatorname{prod\text{-}fun}(\mathbb{C}, (\lambda k \in \mathbb{N}.\; (\lambda x \in \mathbb{C}.\; (((1 - (x \cdot \operatorname{recip}(a(k)))) \cdot \exp \operatorname{cc\text{-}series\text{-}partial\text{-}sum}((\lambda \mathit{clhk} \in \mathbb{N}.\; ({\mathit{clhk}}^{\ast\!-1} \cdot \operatorname{power}((x \cdot \operatorname{recip}(a(k))), \mathit{clhk}))), ((\lambda \mathit{c6tk} \in \mathbb{N}.\; t)(k) + 1))) - 1))))(z) \wedge \\ \quad \forall z \in \mathbb{C}. \\ \quad \quad f(z) \\ \quad \quad =\; (\operatorname{power}(z, m) \cdot \exp h(z)) \cdot \operatorname{prod\text{-}fun}(\mathbb{C}, (\lambda k \in \mathbb{N}.\; (\lambda x \in \mathbb{C}.\; (((1 - (x \cdot \operatorname{recip}(a(k)))) \cdot \exp \operatorname{cc\text{-}series\text{-}partial\text{-}sum}((\lambda \mathit{clhk} \in \mathbb{N}.\; ({\mathit{clhk}}^{\ast\!-1} \cdot \operatorname{power}((x \cdot \operatorname{recip}(a(k))), \mathit{clhk}))), ((\lambda \mathit{c6tk} \in \mathbb{N}.\; t)(k) + 1))) - 1))))(z) \Rightarrow \\ \quad \quad (g(p) - h(p)) = (g(r) - h(r)) \end{array}}}

\vnbbookfile{theorem-library/ch6-gamma}{theorem-library/ch6-gamma}
\noindent{\small ch6-gamma.scm -- WAVE 2 (2026-10-27..30), wired 2026-10-30 from scratchpad/ch6/ch6-gamma.scm. Library: complex-analysis.}\par
\par\medskip\noindent\textbf{Definition} (\texttt{euler-gamma}).\;
\begin{equation*}
(\operatorname{euler\text{-}gamma} \simeq -(1, \operatorname{series\text{-}limit}((\lambda \operatorname{egk\_} \in \mathbb{N}.\; -(-(\log\left(\operatorname{succ}(\operatorname{succ}(\operatorname{egk\_}))\right), \log\left(\operatorname{succ}(\operatorname{egk\_})\right)), \operatorname{recip}(\operatorname{succ}(\operatorname{succ}(\operatorname{egk\_}))))))))
\end{equation*}

\par\medskip\noindent\textbf{Definition} (\texttt{recip-gamma}).\;
\begin{equation*}
(\operatorname{recip\text{-}gamma} \simeq (\lambda \operatorname{egz\_} \in \mathbb{C}.\; \cdot(\cdot(\operatorname{egz\_}, \exp \cdot(\operatorname{euler\text{-}gamma}, \operatorname{egz\_})), \operatorname{prod\text{-}fun}(\mathbb{C}, (\lambda \operatorname{c6k\_} \in \mathbb{N}.\; (\lambda \operatorname{c6x\_} \in \mathbb{C}.\; -(\cdot(-(1, \cdot(\operatorname{c6x\_}, \operatorname{recip}((\lambda \operatorname{c6gk\_} \in \mathbb{N}.\; -(\operatorname{succ}(\operatorname{c6gk\_})))(\operatorname{c6k\_})))), \exp \operatorname{cc\text{-}series\text{-}partial\text{-}sum}((\lambda \operatorname{clhk\_} \in \mathbb{N}.\; \cdot({\operatorname{clhk\_}}^{\ast\!-1}, \operatorname{power}(\cdot(\operatorname{c6x\_}, \operatorname{recip}((\lambda \operatorname{c6gk\_} \in \mathbb{N}.\; -(\operatorname{succ}(\operatorname{c6gk\_})))(\operatorname{c6k\_}))), \operatorname{clhk\_}))), \operatorname{succ}((\lambda \operatorname{c6gk\_} \in \mathbb{N}.\; 1)(\operatorname{c6k\_})))), 1))))(\operatorname{egz\_}))))
\end{equation*}

\vnbprop{log-step-bounds}{\vnbstmt{\begin{array}{@{}l@{}} \forall k \in \mathbb{N}. \\ \quad \operatorname{recip}(((k + 1) + 1)) \leq (\log\left(((k + 1) + 1)\right) - \log\left((k + 1)\right)) \wedge \\ \quad (\log\left(((k + 1) + 1)\right) - \log\left((k + 1)\right)) \leq \operatorname{recip}((k + 1)) \end{array}}}
\vnbprop{log-le-harmonic}{\vnbstmt{\begin{array}{@{}l@{}} \forall n \in \mathbb{N}. \\ \quad \log\left((n + 1)\right) \leq \operatorname{series\text{-}partial\text{-}sum}((\lambda j \in \mathbb{N}.\; \operatorname{recip}((j + 1))), n) \end{array}}}
\vnbprop{harmonic-minus-log-identity}{\vnbstmt{\begin{array}{@{}l@{}} \forall n \in \mathbb{N}. \\ \quad (\operatorname{series\text{-}partial\text{-}sum}((\lambda j \in \mathbb{N}.\; \operatorname{recip}((j + 1))), (n + 1)) - \log\left((n + 1)\right)) \\ \quad =\; (1 - \operatorname{series\text{-}partial\text{-}sum}((\lambda k \in \mathbb{N}.\; ((\log\left(((k + 1) + 1)\right) - \log\left((k + 1)\right)) - \operatorname{recip}(((k + 1) + 1)))), n)) \end{array}}}
\vnbprop{euler-gamma-limit}{\vnbstmt{\begin{array}{@{}l@{}} \operatorname{converges\text{-}to}(\operatorname{rr\text{-}ms}, \\ \quad (\lambda n \in \mathbb{N}.\; (\operatorname{series\text{-}partial\text{-}sum}((\lambda j \in \mathbb{N}.\; \operatorname{recip}((j + 1))), (n + 1)) - \log\left((n + 1)\right))), \\ \quad \operatorname{euler\text{-}gamma}) \end{array}}}
\vnbprop{recip-square-step}{\vnbstmt{\begin{array}{@{}l@{}} \forall n \in \mathbb{N}. \\ \quad ((\operatorname{recip}((n + 1)) \cdot \operatorname{recip}((n + 1))) \cdot 1/2) \leq (\operatorname{recip}((n + 1)) - \operatorname{recip}(((n + 1) + 1))) \wedge \\ \quad 0 \leq (\operatorname{recip}((n + 1)) - \operatorname{recip}(((n + 1) + 1))) \end{array}}}
\vnbprop{recip-square-telescope}{\vnbstmt{\begin{array}{@{}l@{}} \forall n \in \mathbb{N}, r \in \mathbb{R}. \\ \quad \operatorname{series\text{-}partial\text{-}sum}((\lambda q \in \mathbb{N}.\; ((r \cdot r) \cdot (\operatorname{recip}((q + 1)) - \operatorname{recip}(((q + 1) + 1))))), n) + ((r \cdot r) \cdot \operatorname{recip}((n + 1))) \\ \quad =\; r \cdot r \end{array}}}
\vnbprop{gamma-zeros-nonzero}{\vnbstmt{\forall q \in \mathbb{N}.\; \neg ((\lambda k \in \mathbb{N}.\; -(k + 1))(q) = 0)}}
\vnbprop{gamma-zeros-grow}{\vnbstmt{\begin{array}{@{}l@{}} \forall r \in \mathbb{R}. \exists m \in \mathbb{N}. \forall q \in \mathbb{N}. \\ \quad (m \leq q) \Rightarrow (r \leq \operatorname{magnitude}((\lambda k \in \mathbb{N}.\; -(k + 1))(q))) \end{array}}}
\vnbprop{gamma-zeros-h167}{\vnbstmt{\begin{array}{@{}l@{}} \forall r \in \mathbb{R}. \\ \quad (0 \leq r) \Rightarrow  \\ \quad \quad \operatorname{series\text{-}converges}((\lambda k \in \mathbb{N}.\; \\ \quad \quad \quad \quad \operatorname{power}((r \cdot \operatorname{recip}(\operatorname{magnitude}((\lambda \mathit{c6gk} \in \mathbb{N}.\; -(\mathit{c6gk} + 1))(k)))), ((\lambda \mathit{c6gk} \in \mathbb{N}.\; 1)(k) + 1)) \cdot \operatorname{recip}(((\lambda \mathit{c6gk} \in \mathbb{N}.\; 1)(k) + 1)))) \end{array}}}
\vnbprop{theorem-6-20-multiplicity-fam}{\vnbstmt{\begin{array}{@{}l@{}} \forall a \in (\mathbb{N} \to \mathbb{C}), u \in (\mathbb{N} \to \mathbb{N}), z \in \mathbb{C}, n \in \mathbb{N}. \\ \quad \forall q \in \mathbb{N}.\; \neg (a(q) = 0) \wedge \\ \quad \forall r \in \mathbb{R}. \exists m \in \mathbb{N}. \forall q \in \mathbb{N}. \\ \quad \quad (m \leq q) \Rightarrow (r \leq \operatorname{magnitude}(a(q))) \wedge \\ \quad \forall r \in \mathbb{R}. \\ \quad \quad (0 \leq r) \Rightarrow  \\ \quad \quad \quad \operatorname{series\text{-}converges}((\lambda k \in \mathbb{N}.\; (\operatorname{power}((r \cdot \operatorname{recip}(\operatorname{magnitude}(a(k)))), (u(k) + 1)) \cdot \operatorname{recip}((u(k) + 1))))) \wedge \\ \quad \forall q \in \mathbb{N}.\; ((n \leq q) \Rightarrow \neg (a(q) = z)) \Rightarrow \\ \quad \quad \operatorname{has\text{-}zero\text{-}of\text{-}order}(\mathbb{C}, \\ \quad \quad \quad \operatorname{prod\text{-}fun}(\mathbb{C}, \\ \quad \quad \quad \quad (\lambda k \in \mathbb{N}.\; \\ \quad \quad \quad \quad \quad (\lambda x \in \mathbb{C}.\; \\ \quad \quad \quad \quad \quad \quad (((1 - (x \cdot \operatorname{recip}(a(k)))) \cdot \exp \operatorname{cc\text{-}series\text{-}partial\text{-}sum}((\lambda \mathit{clhk} \in \mathbb{N}.\; ({\mathit{clhk}}^{\ast\!-1} \cdot \operatorname{power}((x \cdot \operatorname{recip}(a(k))), \mathit{clhk}))), (u(k) + 1))) - 1)))), \\ \quad \quad \quad z, \operatorname{count\text{-}hits}(a, z, n)) \end{array}}}
\vnbprop{count-hits-none}{\vnbstmt{\begin{array}{@{}l@{}} \forall j \in \mathbb{N}, s, z. \\ \quad \forall q \in \mathbb{N}.\; ((q < j) \Rightarrow \neg (s(q) = z)) \Rightarrow (\operatorname{count\text{-}hits}(s, z, j) = 0) \end{array}}}
\vnbprop{euler-gamma-in-rr}{\vnbstmt{\operatorname{euler\text{-}gamma} \in \mathbb{R}}}
\vnbprop{gamma-g-holomorphic}{\vnbstmt{\begin{array}{@{}l@{}} \forall p. \\ \quad \operatorname{holomorphic\text{-}on}(\mathbb{C}, p) \Rightarrow  \\ \quad \quad \operatorname{holomorphic\text{-}on}(\mathbb{C}, (\lambda z \in \mathbb{C}.\; ((z \cdot \exp (\operatorname{euler\text{-}gamma} \cdot z)) \cdot p(z)))) \end{array}}}
\vnbprop{gamma-g-order-at-zero}{\vnbstmt{\begin{array}{@{}l@{}} \forall p. \\ \quad (\operatorname{holomorphic\text{-}on}(\mathbb{C}, p) \wedge \neg (p(0) = 0)) \Rightarrow  \\ \quad \quad \operatorname{has\text{-}zero\text{-}of\text{-}order}(\mathbb{C}, (\lambda z \in \mathbb{C}.\; ((z \cdot \exp (\operatorname{euler\text{-}gamma} \cdot z)) \cdot p(z))), 0, 1) \end{array}}}
\vnbprop{gamma-g-order-at-p-zero}{\vnbstmt{\begin{array}{@{}l@{}} \forall p.; \forall w \in \mathbb{C}. \\ \quad (\operatorname{holomorphic\text{-}on}(\mathbb{C}, p) \wedge (\neg (w = 0) \wedge \operatorname{has\text{-}zero\text{-}of\text{-}order}(\mathbb{C}, p, w, 1))) \Rightarrow  \\ \quad \quad \operatorname{has\text{-}zero\text{-}of\text{-}order}(\mathbb{C}, (\lambda z \in \mathbb{C}.\; ((z \cdot \exp (\operatorname{euler\text{-}gamma} \cdot z)) \cdot p(z))), w, 1) \end{array}}}
\vnbprop{gamma-g-nonzero}{\vnbstmt{\begin{array}{@{}l@{}} \forall p.; \forall w \in \mathbb{C}. \\ \quad (\operatorname{holomorphic\text{-}on}(\mathbb{C}, p) \wedge (\neg (w = 0) \wedge \neg (p(w) = 0))) \Rightarrow  \\ \quad \quad \neg ((\lambda z \in \mathbb{C}.\; ((z \cdot \exp (\operatorname{euler\text{-}gamma} \cdot z)) \cdot p(z)))(w) = 0) \end{array}}}
\vnbprop{proposition-6-28}{\vnbstmt{\begin{array}{@{}l@{}} \operatorname{holomorphic\text{-}on}(\mathbb{C}, \operatorname{recip\text{-}gamma}) \wedge \\ \forall n \in \mathbb{N}.\; \operatorname{has\text{-}zero\text{-}of\text{-}order}(\mathbb{C}, \operatorname{recip\text{-}gamma}, -n, 1) \wedge \\ \forall w \in \mathbb{C}. \\ \quad \forall n \in \mathbb{N}.\; \neg (w = -n) \Rightarrow \neg (\operatorname{recip\text{-}gamma}(w) = 0) \end{array}}}

\vnbbookfile{theorem-library/ch6-euler}{theorem-library/ch6-euler}
\noindent{\small ch6-euler.scm -- WAVE 2 (2026-10-27..30), wired 2026-10-30 from scratchpad/ch6/ch6-euler.scm. Library: complex-analysis.}\par
\vnbprop{prod-fun-partial-limit}{\vnbstmt{\begin{array}{@{}l@{}} \forall u.; \forall f \in (\mathbb{N} \to (u \to \mathbb{C})), x \in u. \\ \quad (\operatorname{is\text{-}open}(\operatorname{cc\text{-}ms}, u) \wedge \operatorname{prod\text{-}normal\text{-}on\text{-}compacta}(u, f)) \Rightarrow  \\ \quad \quad \operatorname{converges\text{-}to}(\operatorname{cc\text{-}ms}, \\ \quad \quad \quad (\lambda n \in \mathbb{N}.\; \operatorname{pprod}((\lambda k \in \mathbb{N}.\; (1 + f(k)(x))), n)), \\ \quad \quad \quad \operatorname{prod\text{-}fun}(u, f)(x)) \end{array}}}
\vnbprop{weier-q-one}{\vnbstmt{\begin{array}{@{}l@{}} \forall z \in \mathbb{C}. \\ \quad \operatorname{cc\text{-}series\text{-}partial\text{-}sum}((\lambda k \in \mathbb{N}.\; ({k}^{\ast\!-1} \cdot \operatorname{power}(z, k))), 2) = z \end{array}}}
\vnbprop{gamma-factor-value}{\vnbstmt{\begin{array}{@{}l@{}} \forall k \in \mathbb{N}, z \in \mathbb{C}. \\ \quad 1 + (\lambda \mathit{c6k} \in \mathbb{N}.\; (\lambda x \in \mathbb{C}.\; (((1 - (x \cdot \operatorname{recip}((\lambda \mathit{c6gk} \in \mathbb{N}.\; -(\mathit{c6gk} + 1))(\mathit{c6k})))) \cdot \exp \operatorname{cc\text{-}series\text{-}partial\text{-}sum}((\lambda \mathit{clhk} \in \mathbb{N}.\; ({\mathit{clhk}}^{\ast\!-1} \cdot \operatorname{power}((x \cdot \operatorname{recip}((\lambda \mathit{c6gk} \in \mathbb{N}.\; -(\mathit{c6gk} + 1))(\mathit{c6k}))), \mathit{clhk}))), ((\lambda \mathit{c6gk} \in \mathbb{N}.\; 1)(\mathit{c6k}) + 1))) - 1)))(k)(z) \\ \quad =\; (1 + (z \cdot \operatorname{recip}((k + 1)))) \cdot \exp -(z \cdot \operatorname{recip}((k + 1))) \end{array}}}
\vnbprop{pprod-exp-harmonic}{\vnbstmt{\begin{array}{@{}l@{}} \forall m \in \mathbb{N}, z \in \mathbb{C}. \\ \quad \operatorname{pprod}((\lambda k \in \mathbb{N}.\; \exp -(z \cdot \operatorname{recip}((k + 1)))), m) \\ \quad =\; \exp -(z \cdot \operatorname{series\text{-}partial\text{-}sum}((\lambda j \in \mathbb{N}.\; \operatorname{recip}((j + 1))), m)) \end{array}}}
\vnbprop{exp-real-null}{\vnbstmt{\begin{array}{@{}l@{}} \forall z \in \mathbb{C}, r \in (\mathbb{N} \to \mathbb{R}). \\ \quad \operatorname{converges\text{-}to}(\operatorname{rr\text{-}ms}, r, 0) \Rightarrow  \\ \quad \quad \operatorname{converges\text{-}to}(\operatorname{cc\text{-}ms}, (\lambda k \in \mathbb{N}.\; \exp (z \cdot r(k))), 1) \end{array}}}
\vnbprop{euler-d-null}{\vnbstmt{\begin{array}{@{}l@{}} \operatorname{converges\text{-}to}(\operatorname{rr\text{-}ms}, \\ \quad (\lambda d \in \mathbb{N}.\; \\ \quad \quad ((\operatorname{euler\text{-}gamma} + \log\left((d + 1)\right)) - \operatorname{series\text{-}partial\text{-}sum}((\lambda j \in \mathbb{N}.\; \operatorname{recip}((j + 1))), d))), \\ \quad 0) \end{array}}}
\vnbprop{euler-limit-183}{\vnbstmt{\begin{array}{@{}l@{}} \forall z \in \mathbb{C}. \\ \quad \operatorname{converges\text{-}to}(\operatorname{cc\text{-}ms}, \\ \quad \quad (\lambda m \in \mathbb{N}.\; \\ \quad \quad \quad (z \cdot \operatorname{pprod}((\lambda q \in \mathbb{N}.\; (1 + (z \cdot \operatorname{recip}((q + 1))))), m)) \cdot \exp -(z \cdot \log\left((m + 1)\right))), \\ \quad \quad \operatorname{recip\text{-}gamma}(z)) \end{array}}}

\vnbbookfile{theorem-library/ch6-pole}{theorem-library/ch6-pole}
\noindent{\small ch6-pole.scm -- WAVE 2 (2026-10-27..30), wired 2026-10-30 from scratchpad/ch6/ch6-pole.scm. Library: complex-analysis.}\par
\par\medskip\noindent\textbf{Definition} (\texttt{gamma-product}).\;
\begin{equation*}
(\operatorname{gamma\text{-}product} \simeq (\lambda \operatorname{egw\_} \in \{\operatorname{egu\_} \in \mathbb{C} : \forall \operatorname{egn\_} \in \mathbb{N}.\; \neg (\operatorname{egu\_} = -(\operatorname{egn\_}))\}.\; \operatorname{recip}(\operatorname{recip\text{-}gamma}(\operatorname{egw\_}))))
\end{equation*}

\vnbprop{cc-recip-power}{\vnbstmt{\forall m \in \mathbb{N}, x \in \mathbb{C}.\; (\neg (x = 0) \Rightarrow (\operatorname{recip}(\operatorname{power}(x, m)) = \operatorname{power}(\operatorname{recip}(x), m)))}}
\vnbprop{zero-to-pole}{\vnbstmt{\begin{array}{@{}l@{}} \forall f.; \forall a \in \mathbb{C}, m \in \mathbb{N}, g.; \forall r \in \mathbb{R}. \\ \quad \operatorname{holomorphic\text{-}on}(\mathbb{C}, f) \wedge \\ \quad 1 \leq m \wedge \\ \quad \operatorname{has\text{-}zero\text{-}of\text{-}order}(\mathbb{C}, f, a, m) \wedge \\ \quad 0 < r \wedge \\ \quad \forall z \in \operatorname{difference}(\operatorname{ball}(\operatorname{cc\text{-}ms}, a, r), \operatorname{singleton}(a)).\; (g(z) = \operatorname{recip}(f(z))) \Rightarrow \\ \quad \quad \operatorname{is\text{-}pole}(g, a, m) \end{array}}}
\vnbprop{proposition-6-28-gamma}{\vnbstmt{\begin{array}{@{}l@{}} \forall n \in \mathbb{N}.\; \operatorname{is\text{-}pole}(\operatorname{gamma\text{-}product}, -n, 1) \wedge \\ \forall w \in \{u \in \mathbb{C} : \forall n \in \mathbb{N}.\; \neg (u = -n)\}. \\ \quad \neg (\operatorname{gamma\text{-}product}(w) = 0) \end{array}}}

\vnbbookfile{theorem-library/ch6-g629}{theorem-library/ch6-g629}
\noindent{\small ch6-g629.scm -- WAVE 2 (2026-10-27..30), wired 2026-10-30 from scratchpad/ch6/ch6-g629.scm. Library: complex-analysis.}\par
\vnbprop{r-exp-power}{\vnbstmt{\forall k \in \mathbb{N}, a \in \mathbb{R}.\; (\operatorname{power}(\exp\left(a\right), k) = \exp\left((k \cdot a)\right))}}
\vnbprop{bernoulli-inequality}{\vnbstmt{\forall k \in \mathbb{N}, x \in \mathbb{R}.\; (((0 \leq x) \wedge (x \leq 1)) \Rightarrow ((1 - (k \cdot x)) \leq \operatorname{power}((1 - x), k)))}}
\vnbprop{inequality-186}{\vnbstmt{\begin{array}{@{}l@{}} \forall n \in \mathbb{N}, t \in \mathbb{R}. \\ \quad ((0 \leq t) \wedge (t \leq (n + 1))) \Rightarrow  \\ \quad \quad 0 \leq (\exp\left(-t\right) - \operatorname{power}((1 - (t \cdot \operatorname{recip}((n + 1)))), (n + 1))) \wedge \\ \quad \quad (\exp\left(-t\right) - \operatorname{power}((1 - (t \cdot \operatorname{recip}((n + 1)))), (n + 1))) \leq (((t \cdot t) \cdot \exp\left(-t\right)) \cdot \operatorname{recip}((n + 1))) \end{array}}}
\vnbprop{pprod-qq-shift}{\vnbstmt{\begin{array}{@{}l@{}} \forall k \in \mathbb{N}, z \in \mathbb{C}. \\ \quad (k + 1) \cdot \operatorname{pprod}((\lambda q \in \mathbb{N}.\; (1 + (z \cdot \operatorname{recip}((q + 1))))), (k + 1)) \\ \quad =\; (z + 1) \cdot \operatorname{pprod}((\lambda q \in \mathbb{N}.\; (1 + ((z + 1) \cdot \operatorname{recip}((q + 1))))), k) \end{array}}}
\vnbprop{power-lin-diff-on-rhp}{\vnbstmt{\begin{array}{@{}l@{}} \forall k \in \mathbb{N}, c \in \mathbb{C}, a \in \{u \in \mathbb{C} : (0 < \operatorname{real\text{-}part}(u))\}. \\ \quad \operatorname{is\text{-}diff\text{-}on}(\operatorname{cc\text{-}normed\text{-}field}, \{u \in \mathbb{C} : (0 < \operatorname{real\text{-}part}(u))\}, \\ \quad \quad (\lambda w \in \{u \in \mathbb{C} : (0 < \operatorname{real\text{-}part}(u))\}.\; \operatorname{power}((1 - (w \cdot c)), (k + 1))), \\ \quad \quad a, (((k + 1) \cdot \operatorname{power}((1 - (a \cdot c)), k)) \cdot -c)) \end{array}}}
\vnbprop{power-lin-holomorphic-rhp}{\vnbstmt{\begin{array}{@{}l@{}} \forall k \in \mathbb{N}, c \in \mathbb{C}. \\ \quad \operatorname{holomorphic\text{-}on}(\{u \in \mathbb{C} : (0 < \operatorname{real\text{-}part}(u))\}, \\ \quad \quad (\lambda w \in \{u \in \mathbb{C} : (0 < \operatorname{real\text{-}part}(u))\}.\; \operatorname{power}((1 - (w \cdot c)), k))) \end{array}}}
\vnbprop{truncated-integral-base}{\vnbstmt{\begin{array}{@{}l@{}} \forall m \in \mathbb{R}, z \in \mathbb{C}, r \in \mathbb{R}. \\ \quad ((0 < m) \wedge ((0 < \operatorname{real\text{-}part}(z)) \wedge ((0 < r) \wedge (r < m)))) \Rightarrow  \\ \quad \quad \int_{r}^{m} (\lambda w \in \{u \in \mathbb{C} : (0 < \operatorname{real\text{-}part}(u))\}.\; (\operatorname{power}((1 - (w \cdot \operatorname{recip}(m))), 0) \cdot \exp ((z - 1) \cdot \log\left(w\right)))) \\ \quad \quad =\; ((\exp (z \cdot \log\left(m\right)) \cdot \operatorname{recip}(z)) - (\exp (z \cdot \log\left(r\right)) \cdot \operatorname{recip}(z))) \end{array}}}
\vnbprop{truncated-integral-parts}{\vnbstmt{\begin{array}{@{}l@{}} \forall m \in \mathbb{R}, k \in \mathbb{N}, z \in \mathbb{C}, r \in \mathbb{R}. \\ \quad ((0 < m) \wedge ((0 < \operatorname{real\text{-}part}(z)) \wedge ((0 < r) \wedge (r < m)))) \Rightarrow  \\ \quad \quad \int_{r}^{m} (\lambda w \in \{u \in \mathbb{C} : (0 < \operatorname{real\text{-}part}(u))\}.\; (\operatorname{power}((1 - (w \cdot \operatorname{recip}(m))), (k + 1)) \cdot \exp ((z - 1) \cdot \log\left(w\right)))) \\ \quad \quad =\; -(\operatorname{power}((1 - (r \cdot \operatorname{recip}(m))), (k + 1)) \cdot (\exp (z \cdot \log\left(r\right)) \cdot \operatorname{recip}(z))) + (((k + 1) \cdot (\operatorname{recip}(m) \cdot \operatorname{recip}(z))) \cdot \int_{r}^{m} (\lambda w \in \{u \in \mathbb{C} : (0 < \operatorname{real\text{-}part}(u))\}.\; (\operatorname{power}((1 - (w \cdot \operatorname{recip}(m))), k) \cdot \exp (((z + 1) - 1) \cdot \log\left(w\right))))) \end{array}}}
\vnbprop{boundary-power-small}{\vnbstmt{\begin{array}{@{}l@{}} \forall z \in \mathbb{C}, e. \\ \quad ((0 < \operatorname{real\text{-}part}(z)) \wedge \operatorname{pos\text{-}rr}(e)) \Rightarrow  \\ \quad \quad \exists d. \\ \quad \quad \quad \operatorname{pos\text{-}rr}(d) \wedge \\ \quad \quad \quad d \leq 1 \wedge \\ \quad \quad \quad \forall r \in \mathbb{R}. \\ \quad \quad \quad \quad ((0 < r) \wedge (r < d)) \Rightarrow (\operatorname{magnitude}((\exp (z \cdot \log\left(r\right)) \cdot \operatorname{recip}(z))) \leq e) \end{array}}}
\vnbprop{qq-factor-nonzero}{\vnbstmt{\begin{array}{@{}l@{}} \forall z \in \mathbb{C}, j \in \mathbb{N}. \\ \quad (0 < \operatorname{real\text{-}part}(z)) \Rightarrow \neg ((1 + (z \cdot \operatorname{recip}((j + 1)))) = 0) \end{array}}}
\vnbprop{pprod-qq-nonzero}{\vnbstmt{\begin{array}{@{}l@{}} \forall k \in \mathbb{N}, z \in \mathbb{C}. \\ \quad (0 < \operatorname{real\text{-}part}(z)) \Rightarrow  \\ \quad \quad \neg (\operatorname{pprod}((\lambda q \in \mathbb{N}.\; (1 + (z \cdot \operatorname{recip}((q + 1))))), k) = 0) \end{array}}}
\vnbprop{truncated-value-recursion}{\vnbstmt{\begin{array}{@{}l@{}} \forall k \in \mathbb{N}, m \in \mathbb{R}, z \in \mathbb{C}. \\ \quad ((0 < m) \wedge (0 < \operatorname{real\text{-}part}(z))) \Rightarrow  \\ \quad \quad (\exp (z \cdot \log\left(m\right)) \cdot \operatorname{recip}(z)) \cdot \operatorname{recip}(\operatorname{pprod}((\lambda q \in \mathbb{N}.\; (1 + (z \cdot \operatorname{recip}((q + 1))))), (k + 1))) \\ \quad \quad =\; ((k + 1) \cdot (\operatorname{recip}(m) \cdot \operatorname{recip}(z))) \cdot ((\exp ((z + 1) \cdot \log\left(m\right)) \cdot \operatorname{recip}((z + 1))) \cdot \operatorname{recip}(\operatorname{pprod}((\lambda q \in \mathbb{N}.\; (1 + ((z + 1) \cdot \operatorname{recip}((q + 1))))), k))) \end{array}}}
\vnbprop{truncated-limit-at-zero}{\vnbstmt{\begin{array}{@{}l@{}} \forall k \in \mathbb{N}, m \in \mathbb{R}, z \in \mathbb{C}, e. \\ \quad ((0 < m) \wedge ((0 < \operatorname{real\text{-}part}(z)) \wedge \operatorname{pos\text{-}rr}(e))) \Rightarrow  \\ \quad \quad \exists d. \\ \quad \quad \quad \operatorname{pos\text{-}rr}(d) \wedge \\ \quad \quad \quad \forall r \in \mathbb{R}. \\ \quad \quad \quad \quad ((0 < r) \wedge ((r < d) \wedge (r < m))) \Rightarrow  \\ \quad \quad \quad \quad \quad \operatorname{magnitude}((\int_{r}^{m} (\lambda w \in \{u \in \mathbb{C} : (0 < \operatorname{real\text{-}part}(u))\}.\; (\operatorname{power}((1 - (w \cdot \operatorname{recip}(m))), k) \cdot \exp ((z - 1) \cdot \log\left(w\right)))) - ((\exp (z \cdot \log\left(m\right)) \cdot \operatorname{recip}(z)) \cdot \operatorname{recip}(\operatorname{pprod}((\lambda q \in \mathbb{N}.\; (1 + (z \cdot \operatorname{recip}((q + 1))))), k))))) \\ \quad \quad \quad \quad \quad \leq\; e \end{array}}}
\vnbprop{integrand-difference-bound}{\vnbstmt{\begin{array}{@{}l@{}} \forall n \in \mathbb{N}, z \in \mathbb{C}, t \in \mathbb{R}. \\ \quad ((0 < \operatorname{real\text{-}part}(z)) \wedge ((0 < t) \wedge (t \leq (n + 1)))) \Rightarrow  \\ \quad \quad \operatorname{magnitude}((\operatorname{gamma\text{-}integrand}(z)(t) - (\operatorname{power}((1 - (t \cdot \operatorname{recip}((n + 1)))), (n + 1)) \cdot \exp ((z - 1) \cdot \log\left(t\right))))) \\ \quad \quad \leq\; \operatorname{magnitude}(\operatorname{gamma\text{-}integrand}((\operatorname{real\text{-}part}(z) + 2))(t)) \cdot \operatorname{recip}((n + 1)) \end{array}}}
\vnbprop{integral-difference-small}{\vnbstmt{\begin{array}{@{}l@{}} \forall n \in \mathbb{N}, z \in \mathbb{C}, r \in \mathbb{R}. \\ \quad ((0 < \operatorname{real\text{-}part}(z)) \wedge ((0 < r) \wedge (r < 1))) \Rightarrow  \\ \quad \quad \operatorname{magnitude}((\int_{r}^{((n + 1) + 1)} \operatorname{gamma\text{-}integrand}(z) - \int_{r}^{((n + 1) + 1)} (\lambda w \in \{u \in \mathbb{C} : (0 < \operatorname{real\text{-}part}(u))\}.\; (\operatorname{power}((1 - (w \cdot \operatorname{recip}(((n + 1) + 1)))), ((n + 1) + 1)) \cdot \exp ((z - 1) \cdot \log\left(w\right)))))) \\ \quad \quad \leq\; (1 + (2 \cdot \exp (((\operatorname{real\text{-}part}(z) + 2) \cdot (\log\left((2 \cdot (\operatorname{real\text{-}part}(z) + 2))\right) - 1)) - 1/2))) \cdot \operatorname{recip}(((n + 1) + 1)) \end{array}}}
\vnbprop{truncated-values-to-gamma}{\vnbstmt{\begin{array}{@{}l@{}} \forall z \in \mathbb{C}. \\ \quad (0 < \operatorname{real\text{-}part}(z)) \Rightarrow  \\ \quad \quad \operatorname{converges\text{-}to}(\operatorname{cc\text{-}ms}, \\ \quad \quad \quad (\lambda q \in \mathbb{N}.\; \\ \quad \quad \quad \quad (\exp (z \cdot \log\left(((q + 1) + 1)\right)) \cdot \operatorname{recip}(z)) \cdot \operatorname{recip}(\operatorname{pprod}((\lambda \mathit{egq} \in \mathbb{N}.\; (1 + (z \cdot \operatorname{recip}((\mathit{egq} + 1))))), ((q + 1) + 1)))), \\ \quad \quad \quad \Gamma(z)) \end{array}}}
\vnbprop{log-shift-null}{\vnbstmt{\operatorname{converges\text{-}to}(\operatorname{rr\text{-}ms}, (\lambda q \in \mathbb{N}.\; (\log\left((((q + 1) + 1) + 1)\right) - \log\left(((q + 1) + 1)\right))), 0)}}
\vnbprop{proposition-6-29}{\vnbstmt{\begin{array}{@{}l@{}} \forall z \in \mathbb{C}. \\ \quad (0 < \operatorname{real\text{-}part}(z)) \Rightarrow (\Gamma(z) = \operatorname{recip}(\operatorname{recip\text{-}gamma}(z))) \end{array}}}

\vnbbookfile{theorem-library/ch6-g361}{theorem-library/ch6-g361}
\noindent{\small ch6-g361.scm -- WAVE 2 (2026-10-27..30), wired 2026-10-30 from scratchpad/ch6/ch6-g361.scm. Library: complex-analysis.}\par
\vnbprop{recip-gamma-nonzero-rhp}{\vnbstmt{\forall w \in \mathbb{C}.\; ((0 < \operatorname{real\text{-}part}(w)) \Rightarrow \neg (\operatorname{recip\text{-}gamma}(w) = 0))}}
\vnbprop{corollary-3-61}{\vnbstmt{\begin{array}{@{}l@{}} \forall z \in \{\mathit{gdz} \in \mathbb{C} : \forall n \in \mathbb{N}.\; \neg ((\mathit{gdz} + n) = 0)\}. \\ \quad \neg (\operatorname{gamma\text{-}ext}(z) = 0) \end{array}}}

\vnbbookfile{structure-library/ccr4-vocab}{structure-library/ccr4-vocab}
\noindent{\small ccr4-vocab.scm -- ADDITIVE CHARACTERS OF A NORMED FIELD AND THE DUALITY MAP (CCR.pdf chapter 5, Definitions 5.1 and 5.11; chapter 6, Definition 6.1 for E = K).  DEFINITIONS ONLY; the laws are in theorem-library/ccr4-chars.scm.  Agent CCR-B4, 2026-10-26.  Destination: structure-library/ccr4-vocab.scm. Library: analysis. ...}\par
\par\medskip\noindent\textbf{Definition} (\texttt{is-add-character}).\;For all $\operatorname{chk\_}$ and $\operatorname{chc\_}$, 
\begin{equation*}
\begin{array}{@{}l@{}} \operatorname{is\text{-}add\text{-}character}(\operatorname{chk\_}, \operatorname{chc\_}) \\ \Leftrightarrow\; \operatorname{chc\_} \in (\operatorname{carr}(\operatorname{chk\_}) \to \operatorname{unit\text{-}circle}) \wedge \\ \forall \operatorname{chx\_}, \operatorname{chy\_} \in \operatorname{carr}(\operatorname{chk\_}). \\ \quad \operatorname{chc\_}(\operatorname{add}(\operatorname{chk\_})(\operatorname{chx\_}, \operatorname{chy\_})) = \cdot(\operatorname{chc\_}(\operatorname{chx\_}), \operatorname{chc\_}(\operatorname{chy\_})) \wedge \\ \operatorname{is\text{-}continuous}(\operatorname{nf\text{-}metric\text{-}space}(\operatorname{chk\_}), \operatorname{cc\text{-}ms}, \operatorname{chc\_}) \end{array}
\end{equation*}

\par\medskip\noindent\textbf{Definition} (\texttt{is-nontrivial-character}).\;For all $\operatorname{chk\_}$ and $\operatorname{chc\_}$, 
\begin{equation*}
\begin{array}{@{}l@{}} \operatorname{is\text{-}nontrivial\text{-}character}(\operatorname{chk\_}, \operatorname{chc\_}) \\ \Leftrightarrow\; (\operatorname{is\text{-}add\text{-}character}(\operatorname{chk\_}, \operatorname{chc\_}) \wedge \exists \operatorname{chx\_} \in \operatorname{carr}(\operatorname{chk\_}).\; \neg (\operatorname{chc\_}(\operatorname{chx\_}) = 1)) \end{array}
\end{equation*}

\par\medskip\noindent\textbf{Definition} (\texttt{duality-map}).\;$\operatorname{duality\text{-}map}(\operatorname{chk\_}, \operatorname{chc\_}, \operatorname{cha\_}) :=$ $(\lambda \operatorname{chx\_} \in \operatorname{carr}(\operatorname{chk\_}).\; \operatorname{chc\_}(\operatorname{mul}(\operatorname{chk\_})(\operatorname{cha\_}, \operatorname{chx\_})))$
\par\medskip\noindent\textbf{Definition} (\texttt{is-bicharacter}).\;For all $\operatorname{chk\_}$ and $\operatorname{chw\_}$, 
\begin{equation*}
\begin{array}{@{}l@{}} \operatorname{is\text{-}bicharacter}(\operatorname{chk\_}, \operatorname{chw\_}) \\ \Leftrightarrow\; \operatorname{chw\_} \in ((\operatorname{carr}(\operatorname{chk\_}) \times \operatorname{carr}(\operatorname{chk\_})) \to \operatorname{unit\text{-}circle}) \wedge \\ \forall \operatorname{chy\_} \in \operatorname{carr}(\operatorname{chk\_}). \\ \quad \operatorname{is\text{-}add\text{-}character}(\operatorname{chk\_}, (\lambda \operatorname{chx\_} \in \operatorname{carr}(\operatorname{chk\_}).\; \operatorname{chw\_}(\operatorname{chx\_}, \operatorname{chy\_}))) \wedge \\ \forall \operatorname{chy\_} \in \operatorname{carr}(\operatorname{chk\_}). \\ \quad \operatorname{is\text{-}add\text{-}character}(\operatorname{chk\_}, (\lambda \operatorname{chx\_} \in \operatorname{carr}(\operatorname{chk\_}).\; \operatorname{chw\_}(\operatorname{chy\_}, \operatorname{chx\_}))) \end{array}
\end{equation*}

\par\medskip\noindent\textbf{Definition} (\texttt{is-bilinear-bicharacter}).\;For all $\operatorname{chk\_}$ and $\operatorname{chw\_}$, 
\begin{equation*}
\begin{array}{@{}l@{}} \operatorname{is\text{-}bilinear\text{-}bicharacter}(\operatorname{chk\_}, \operatorname{chw\_}) \\ \Leftrightarrow\; \operatorname{is\text{-}bicharacter}(\operatorname{chk\_}, \operatorname{chw\_}) \wedge \\ \forall \operatorname{cha\_}, \operatorname{chx\_}, \operatorname{chy\_} \in \operatorname{carr}(\operatorname{chk\_}). \\ \quad \operatorname{chw\_}(\operatorname{mul}(\operatorname{chk\_})(\operatorname{cha\_}, \operatorname{chx\_}), \operatorname{chy\_}) = \operatorname{chw\_}(\operatorname{chx\_}, \operatorname{mul}(\operatorname{chk\_})(\operatorname{cha\_}, \operatorname{chy\_})) \end{array}
\end{equation*}

\par\medskip\noindent\textbf{Definition} (\texttt{bichar-of}).\;$\operatorname{bichar\text{-}of}(\operatorname{chk\_}, \operatorname{chc\_}, \operatorname{cha\_}) :=$
\begin{equation*}
\begin{array}{@{}l@{}} (\lambda {[}\operatorname{chu\_}, \operatorname{chv\_}] \in (\operatorname{carr}(\operatorname{chk\_}) \times \operatorname{carr}(\operatorname{chk\_})).\; \\ \quad \operatorname{chc\_}(\operatorname{mul}(\operatorname{chk\_})(\operatorname{cha\_}, \operatorname{mul}(\operatorname{chk\_})(\operatorname{chu\_}, \operatorname{chv\_})))) \end{array}
\end{equation*}


\vnbbookfile{theorem-library/ccr1109a5-finite-field}{theorem-library/ccr1109a5-finite-field}
\noindent{\small ccr1109a5-finite-field.scm -- wave of 2026-11-09/10 (agent CCR-1109-A5), wired 2026-11-10 from scratchpad/ccr1109a5/ccr1109a5-finite-field.scm. Library: complex-analysis.}\par
\vnbprop{nf-norm-pos-of-nonzero}{\vnbstmt{\begin{array}{@{}l@{}} \forall f.; \forall x \in \operatorname{carr}(f). \\ \quad (\operatorname{is\text{-}normed\text{-}field}(f) \wedge \neg (x = \operatorname{zero}(f))) \Rightarrow (0 < \operatorname{fnrm}(f)(x)) \end{array}}}
\vnbprop{nf-finite-norm-ge-one}{\vnbstmt{\begin{array}{@{}l@{}} \forall f.; \forall x \in \operatorname{carr}(f). \\ \quad (\operatorname{is\text{-}normed\text{-}field}(f) \wedge ((|\operatorname{carr}(f)| \in \mathbb{N}) \wedge \neg (x = \operatorname{zero}(f)))) \Rightarrow  \\ \quad \quad 1 \leq \operatorname{fnrm}(f)(x) \end{array}}}
\vnbprop{nf-finite-norm-trivial}{\vnbstmt{\begin{array}{@{}l@{}} \forall f.; \forall x \in \operatorname{carr}(f). \\ \quad (\operatorname{is\text{-}normed\text{-}field}(f) \wedge ((|\operatorname{carr}(f)| \in \mathbb{N}) \wedge \neg (x = \operatorname{zero}(f)))) \Rightarrow  \\ \quad \quad \operatorname{fnrm}(f)(x) = 1 \end{array}}}
\vnbprop{ms-locally-constant-continuous}{\vnbstmt{\begin{array}{@{}l@{}} \forall s, t.; \forall f \in (\operatorname{pts}(s) \to \operatorname{pts}(t)). \\ \quad \operatorname{is\text{-}metric\text{-}space}(s) \wedge \\ \quad \operatorname{is\text{-}metric\text{-}space}(t) \wedge \\ \quad \forall a \in \operatorname{pts}(s). \exists d. \\ \quad \quad \operatorname{pos\text{-}rr}(d) \wedge \forall b \in \operatorname{pts}(s).\; ((\operatorname{dist}(s)(a, b) \leq d) \Rightarrow (f(b) = f(a))) \Rightarrow \\ \quad \quad \operatorname{is\text{-}continuous}(s, t, f) \end{array}}}
\vnbprop{nf-finite-continuous}{\vnbstmt{\begin{array}{@{}l@{}} \forall f, t.; \forall a \in (\operatorname{carr}(f) \to \operatorname{pts}(t)). \\ \quad (\operatorname{is\text{-}normed\text{-}field}(f) \wedge ((|\operatorname{carr}(f)| \in \mathbb{N}) \wedge \operatorname{is\text{-}metric\text{-}space}(t))) \Rightarrow  \\ \quad \quad \operatorname{is\text{-}continuous}(\operatorname{nf\text{-}metric\text{-}space}(f), t, a) \end{array}}}
\vnbprop{nf-finite-add-character}{\vnbstmt{\begin{array}{@{}l@{}} \forall f.; \forall c \in (\operatorname{carr}(f) \to \operatorname{unit\text{-}circle}). \\ \quad \operatorname{is\text{-}normed\text{-}field}(f) \wedge \\ \quad |\operatorname{carr}(f)| \in \mathbb{N} \wedge \\ \quad \forall x, y \in \operatorname{carr}(f).\; (c(\operatorname{add}(f)(x, y)) = (c(x) \cdot c(y))) \Rightarrow \\ \quad \quad \operatorname{is\text{-}add\text{-}character}(f, c) \end{array}}}
\vnbprop{triv-abs-zero}{\vnbstmt{\forall q.\; (\operatorname{is\text{-}ring}(q) \Rightarrow (\operatorname{triv\text{-}abs}(q, \operatorname{zero}(q)) = 0))}}
\vnbprop{triv-abs-nonzero}{\vnbstmt{\begin{array}{@{}l@{}} \forall q.; \forall x \in \operatorname{carr}(q). \\ \quad (\operatorname{is\text{-}ring}(q) \wedge \neg (x = \operatorname{zero}(q))) \Rightarrow (\operatorname{triv\text{-}abs}(q, x) = 1) \end{array}}}
\vnbprop{triv-abs-range}{\vnbstmt{\begin{array}{@{}l@{}} \forall q.; \forall x \in \operatorname{carr}(q). \\ \quad \operatorname{is\text{-}ring}(q) \Rightarrow  \\ \quad \quad (\operatorname{triv\text{-}abs}(q, x) \in \mathbb{R}) \wedge ((0 \leq \operatorname{triv\text{-}abs}(q, x)) \wedge (\operatorname{triv\text{-}abs}(q, x) \leq 1)) \end{array}}}
\vnbprop{triv-abs-zero-iff}{\vnbstmt{\forall q.; \forall x \in \operatorname{carr}(q).\; (\operatorname{is\text{-}ring}(q) \Rightarrow ((\operatorname{triv\text{-}abs}(q, x) = 0) \Leftrightarrow (x = \operatorname{zero}(q))))}}
\vnbprop{triv-abs-mul}{\vnbstmt{\begin{array}{@{}l@{}} \forall q.; \forall x, y \in \operatorname{carr}(q). \\ \quad \operatorname{is\text{-}integral\text{-}domain}(q) \Rightarrow  \\ \quad \quad \operatorname{triv\text{-}abs}(q, \operatorname{mul}(q)(x, y)) = (\operatorname{triv\text{-}abs}(q, x) \cdot \operatorname{triv\text{-}abs}(q, y)) \end{array}}}
\vnbprop{triv-abs-add}{\vnbstmt{\begin{array}{@{}l@{}} \forall q.; \forall x, y \in \operatorname{carr}(q). \\ \quad \operatorname{is\text{-}commutative\text{-}ring}(q) \Rightarrow  \\ \quad \quad \operatorname{triv\text{-}abs}(q, \operatorname{add}(q)(x, y)) \leq (\operatorname{triv\text{-}abs}(q, x) + \operatorname{triv\text{-}abs}(q, y)) \end{array}}}
\vnbprop{triv-nf-carr}{\vnbstmt{\forall q.\; (\operatorname{is\text{-}ring}(q) \Rightarrow (\operatorname{carr}(\operatorname{triv\text{-}nf}(q)) \simeq \operatorname{carr}(q)))}}
\vnbprop{triv-nf-add}{\vnbstmt{\forall q.\; (\operatorname{is\text{-}ring}(q) \Rightarrow (\operatorname{add}(\operatorname{triv\text{-}nf}(q)) \simeq \operatorname{add}(q)))}}
\vnbprop{triv-nf-mul}{\vnbstmt{\forall q.\; (\operatorname{is\text{-}ring}(q) \Rightarrow (\operatorname{mul}(\operatorname{triv\text{-}nf}(q)) \simeq \operatorname{mul}(q)))}}
\vnbprop{triv-nf-neg}{\vnbstmt{\forall q.\; (\operatorname{is\text{-}ring}(q) \Rightarrow (\operatorname{neg}(\operatorname{triv\text{-}nf}(q)) \simeq \operatorname{neg}(q)))}}
\vnbprop{triv-nf-zero}{\vnbstmt{\forall q.\; (\operatorname{is\text{-}ring}(q) \Rightarrow (\operatorname{zero}(\operatorname{triv\text{-}nf}(q)) \simeq \operatorname{zero}(q)))}}
\vnbprop{triv-nf-one}{\vnbstmt{\forall q.\; (\operatorname{is\text{-}ring}(q) \Rightarrow (\operatorname{one}(\operatorname{triv\text{-}nf}(q)) \simeq \operatorname{one}(q)))}}
\vnbprop{triv-nf-fnrm}{\vnbstmt{\begin{array}{@{}l@{}} \forall q. \\ \quad \operatorname{is\text{-}ring}(q) \Rightarrow (\operatorname{fnrm}(\operatorname{triv\text{-}nf}(q)) \simeq (\lambda x \in \operatorname{carr}(q).\; \operatorname{triv\text{-}abs}(q, x))) \end{array}}}
\vnbprop{triv-nf-is-normed-field}{\vnbstmt{\forall q.\; (\operatorname{is\text{-}field\text{-}ring}(q) \Rightarrow \operatorname{is\text{-}normed\text{-}field}(\operatorname{triv\text{-}nf}(q)))}}
\vnbprop{triv-nf-add-character}{\vnbstmt{\begin{array}{@{}l@{}} \forall q.; \forall c \in (\operatorname{carr}(q) \to \operatorname{unit\text{-}circle}). \\ \quad \operatorname{is\text{-}field\text{-}ring}(q) \wedge \\ \quad |\operatorname{carr}(q)| \in \mathbb{N} \wedge \\ \quad \forall x, y \in \operatorname{carr}(q).\; (c(\operatorname{add}(q)(x, y)) = (c(x) \cdot c(y))) \Rightarrow \\ \quad \quad \operatorname{is\text{-}add\text{-}character}(\operatorname{triv\text{-}nf}(q), c) \end{array}}}

\vnbbookfile{structure-library/ccr5-vocab}{structure-library/ccr5-vocab}
\noindent{\small ccr5-vocab.scm -- wave of 2026-11-09/10 (agent CCR-1109-D), wired 2026-11-10 from scratchpad/ccr1109d/ccr5-vocab.scm. Library: complex-analysis.}\par
\par\medskip\noindent\textbf{Definition} (\texttt{is-local-uniformizer}).\;For all $\operatorname{c5k\_}$ and $\operatorname{c5u\_}$, 
\begin{equation*}
\begin{array}{@{}l@{}} \operatorname{is\text{-}local\text{-}uniformizer}(\operatorname{c5k\_}, \operatorname{c5u\_}) \\ \Leftrightarrow\; \operatorname{is\text{-}compact\text{-}nv\text{-}domain}(\operatorname{c5k\_}, \operatorname{nv\text{-}integers}(\operatorname{c5k\_})) \wedge \\ \operatorname{c5u\_} \in \operatorname{nv\text{-}integers}(\operatorname{c5k\_}) \wedge \\ \operatorname{nv\text{-}maximal}(\operatorname{c5k\_}, \operatorname{nv\text{-}integers}(\operatorname{c5k\_})) = \operatorname{nv\text{-}principal}(\operatorname{c5k\_}, \operatorname{nv\text{-}integers}(\operatorname{c5k\_}), \operatorname{c5u\_}) \wedge \\ \neg (\operatorname{c5u\_} = \operatorname{zero}(\operatorname{c5k\_})) \end{array}
\end{equation*}

\par\medskip\noindent\textbf{Definition} (\texttt{is-locally-constant-fn}).\;For all $\operatorname{c5k\_}$ and $\operatorname{c5f\_}$, 
\begin{equation*}
\begin{array}{@{}l@{}} \operatorname{is\text{-}locally\text{-}constant\text{-}fn}(\operatorname{c5k\_}, \operatorname{c5f\_}) \\ \Leftrightarrow\; \forall \operatorname{c5x\_} \in \operatorname{carr}(\operatorname{c5k\_}). \exists \operatorname{c5r\_}. \\ \quad \operatorname{pos\text{-}rr}(\operatorname{c5r\_}) \wedge \\ \quad \forall \operatorname{c5y\_} \in \operatorname{carr}(\operatorname{c5k\_}). \\ \quad \quad (\operatorname{fnrm}(\operatorname{c5k\_})(\operatorname{add}(\operatorname{c5k\_})(\operatorname{c5y\_}, \operatorname{neg}(\operatorname{c5k\_})(\operatorname{c5x\_}))) \leq \operatorname{c5r\_}) \Rightarrow  \\ \quad \quad \quad \operatorname{c5f\_}(\operatorname{c5y\_}) = \operatorname{c5f\_}(\operatorname{c5x\_}) \end{array}
\end{equation*}

\par\medskip\noindent\textbf{Definition} (\texttt{char-converges}).\;For all $\operatorname{c5k\_}$, $\operatorname{c5s\_}$ and $\operatorname{c5c\_}$, 
\begin{equation*}
\begin{array}{@{}l@{}} \operatorname{char\text{-}converges}(\operatorname{c5k\_}, \operatorname{c5s\_}, \operatorname{c5c\_}) \\ \Leftrightarrow\; \forall \operatorname{c5e\_}. \\ \quad (\operatorname{c5e\_} \subseteq \operatorname{carr}(\operatorname{c5k\_})) \Rightarrow  \\ \quad \operatorname{is\text{-}compact}(\operatorname{subspace\text{-}ms}(\operatorname{nf\text{-}metric\text{-}space}(\operatorname{c5k\_}), \operatorname{c5e\_})) \Rightarrow  \\ \quad \quad \forall \operatorname{c5p\_}. \\ \quad \quad \quad \operatorname{pos\text{-}rr}(\operatorname{c5p\_}) \Rightarrow  \\ \quad \quad \quad \quad \exists \operatorname{c5n\_} \in \mathbb{N}. \forall \operatorname{c5m\_} \in \mathbb{N}. \\ \quad \quad \quad \quad \quad (\operatorname{c5n\_} \leq \operatorname{c5m\_}) \Rightarrow  \\ \quad \quad \quad \quad \quad \quad \forall \operatorname{c5x\_} \in \operatorname{c5e\_}.\; (\operatorname{magnitude}(-(\operatorname{c5s\_}(\operatorname{c5m\_})(\operatorname{c5x\_}), \operatorname{c5c\_}(\operatorname{c5x\_}))) < \operatorname{c5p\_}) \end{array}
\end{equation*}

\par\medskip\noindent\textbf{Definition} (\texttt{bichar-polar}).\;$\operatorname{bichar\text{-}polar}(\operatorname{c5k\_}, \operatorname{c5w\_}, \operatorname{c5v\_}) :=$
\begin{equation*}
\{\operatorname{c5h\_} \in \operatorname{carr}(\operatorname{c5k\_}) : \forall \operatorname{c5x\_} \in \operatorname{c5v\_}.\; (\operatorname{c5w\_}(\operatorname{c5h\_}, \operatorname{c5x\_}) = 1)\}
\end{equation*}

\par\medskip\noindent\textbf{Definition} (\texttt{is-unipolar}).\;For all $\operatorname{c5k\_}$ and $\operatorname{c5w\_}$, 
\begin{equation*}
\begin{array}{@{}l@{}} \operatorname{is\text{-}unipolar}(\operatorname{c5k\_}, \operatorname{c5w\_}) \\ \Leftrightarrow\; \operatorname{is\text{-}bilinear\text{-}bicharacter}(\operatorname{c5k\_}, \operatorname{c5w\_}) \wedge \\ \operatorname{bichar\text{-}polar}(\operatorname{c5k\_}, \operatorname{c5w\_}, \operatorname{nv\text{-}integers}(\operatorname{c5k\_})) = \operatorname{nv\text{-}integers}(\operatorname{c5k\_}) \end{array}
\end{equation*}

\par\medskip\noindent\textbf{Definition} (\texttt{is-dualizing-character}).\;For all $\operatorname{c5k\_}$ and $\operatorname{c5c\_}$, 
\begin{equation*}
\begin{array}{@{}l@{}} \operatorname{is\text{-}dualizing\text{-}character}(\operatorname{c5k\_}, \operatorname{c5c\_}) \\ \Leftrightarrow\; \operatorname{is\text{-}nontrivial\text{-}character}(\operatorname{c5k\_}, \operatorname{c5c\_}) \wedge \\ \forall \operatorname{c5t\_}. \\ \quad \operatorname{is\text{-}add\text{-}character}(\operatorname{c5k\_}, \operatorname{c5t\_}) \Rightarrow  \\ \quad \quad \exists \operatorname{c5a\_} \in \operatorname{carr}(\operatorname{c5k\_}).\; (\operatorname{c5t\_} = \operatorname{duality\text{-}map}(\operatorname{c5k\_}, \operatorname{c5c\_}, \operatorname{c5a\_})) \end{array}
\end{equation*}

\par\medskip\noindent\textbf{Definition} (\texttt{is-faithful-bicharacter}).\;For all $\operatorname{c5k\_}$ and $\operatorname{c5w\_}$, 
\begin{equation*}
\begin{array}{@{}l@{}} \operatorname{is\text{-}faithful\text{-}bicharacter}(\operatorname{c5k\_}, \operatorname{c5w\_}) \\ \Leftrightarrow\; \operatorname{is\text{-}bicharacter}(\operatorname{c5k\_}, \operatorname{c5w\_}) \wedge \\ \forall \operatorname{c5c\_}. \\ \quad \operatorname{is\text{-}add\text{-}character}(\operatorname{c5k\_}, \operatorname{c5c\_}) \Rightarrow  \\ \quad \quad \exists \operatorname{c5x\_} \in \operatorname{carr}(\operatorname{c5k\_}). \\ \quad \quad \quad \operatorname{c5c\_} = (\lambda \operatorname{c5h\_} \in \operatorname{carr}(\operatorname{c5k\_}).\; \operatorname{c5w\_}(\operatorname{c5h\_}, \operatorname{c5x\_})) \wedge \\ \quad \quad \quad \forall \operatorname{c5y\_} \in \operatorname{carr}(\operatorname{c5k\_}). \\ \quad \quad \quad \quad (\operatorname{c5c\_} = (\lambda \operatorname{c5h\_} \in \operatorname{carr}(\operatorname{c5k\_}).\; \operatorname{c5w\_}(\operatorname{c5h\_}, \operatorname{c5y\_}))) \Rightarrow (\operatorname{c5y\_} = \operatorname{c5x\_}) \wedge \\ \quad \quad \exists \operatorname{c5x\_} \in \operatorname{carr}(\operatorname{c5k\_}). \\ \quad \quad \quad \operatorname{c5c\_} = (\lambda \operatorname{c5h\_} \in \operatorname{carr}(\operatorname{c5k\_}).\; \operatorname{c5w\_}(\operatorname{c5x\_}, \operatorname{c5h\_})) \wedge \\ \quad \quad \quad \forall \operatorname{c5y\_} \in \operatorname{carr}(\operatorname{c5k\_}). \\ \quad \quad \quad \quad (\operatorname{c5c\_} = (\lambda \operatorname{c5h\_} \in \operatorname{carr}(\operatorname{c5k\_}).\; \operatorname{c5w\_}(\operatorname{c5y\_}, \operatorname{c5h\_}))) \Rightarrow (\operatorname{c5y\_} = \operatorname{c5x\_}) \wedge \\ \forall \operatorname{c5x\_} \in \operatorname{carr}(\operatorname{c5k\_}). \\ \quad \neg (\operatorname{c5x\_} = \operatorname{zero}(\operatorname{c5k\_})) \Rightarrow  \\ \quad \quad \exists \operatorname{c5c\_}.\; (\operatorname{is\text{-}add\text{-}character}(\operatorname{c5k\_}, \operatorname{c5c\_}) \wedge \neg (\operatorname{c5c\_}(\operatorname{c5x\_}) = 1)) \end{array}
\end{equation*}


\vnbbookfile{structure-library/ccr6e-vocab}{structure-library/ccr6e-vocab}
\noindent{\small ccr6e-vocab.scm -- wave of 2026-11-09/10 (agent CCR-1109-D3), wired 2026-11-10 from scratchpad/ccr1109d3/ccr6e-vocab.scm. Library: complex-analysis.}\par
\par\medskip\noindent\textbf{Definition} (\texttt{is-k-vector-space}).\;For all $\operatorname{c6vk\_}$ and $\operatorname{c6vm\_}$, 
\begin{equation*}
\begin{array}{@{}l@{}} \operatorname{is\text{-}k\text{-}vector\text{-}space}(\operatorname{c6vk\_}, \operatorname{c6vm\_}) \\ \Leftrightarrow\; \operatorname{is\text{-}module}(\operatorname{c6vm\_}) \wedge (\operatorname{scal}(\operatorname{c6vm\_}) = \operatorname{normed\text{-}field\text{-}as\text{-}commutative\text{-}ring}(\operatorname{c6vk\_})) \end{array}
\end{equation*}

\par\medskip\noindent\textbf{Definition} (\texttt{is-e-character}).\;For all $\operatorname{c6vk\_}$, $\operatorname{c6vm\_}$ and $\operatorname{c6vc\_}$, 
\begin{equation*}
\begin{array}{@{}l@{}} \operatorname{is\text{-}e\text{-}character}(\operatorname{c6vk\_}, \operatorname{c6vm\_}, \operatorname{c6vc\_}) \\ \Leftrightarrow\; \operatorname{c6vc\_} \in (\operatorname{vec}(\operatorname{c6vm\_}) \to \operatorname{unit\text{-}circle}) \wedge \\ \forall \operatorname{c6vx\_}, \operatorname{c6vy\_} \in \operatorname{vec}(\operatorname{c6vm\_}). \\ \quad \operatorname{c6vc\_}(\operatorname{vadd}(\operatorname{c6vm\_})(\operatorname{c6vx\_}, \operatorname{c6vy\_})) = \cdot(\operatorname{c6vc\_}(\operatorname{c6vx\_}), \operatorname{c6vc\_}(\operatorname{c6vy\_})) \wedge \\ \forall \operatorname{c6vx\_} \in \operatorname{vec}(\operatorname{c6vm\_}). \\ \quad \operatorname{is\text{-}add\text{-}character}(\operatorname{c6vk\_}, (\lambda \operatorname{c6va\_} \in \operatorname{carr}(\operatorname{c6vk\_}).\; \operatorname{c6vc\_}(\operatorname{act}(\operatorname{c6vm\_})(\operatorname{c6va\_}, \operatorname{c6vx\_})))) \end{array}
\end{equation*}

\par\medskip\noindent\textbf{Definition} (\texttt{is-e-bicharacter}).\;For all $\operatorname{c6vk\_}$, $\operatorname{c6vm\_}$ and $\operatorname{c6vw\_}$, 
\begin{equation*}
\begin{array}{@{}l@{}} \operatorname{is\text{-}e\text{-}bicharacter}(\operatorname{c6vk\_}, \operatorname{c6vm\_}, \operatorname{c6vw\_}) \\ \Leftrightarrow\; \operatorname{c6vw\_} \in ((\operatorname{vec}(\operatorname{c6vm\_}) \times \operatorname{vec}(\operatorname{c6vm\_})) \to \operatorname{unit\text{-}circle}) \wedge \\ \forall \operatorname{c6vx\_}, \operatorname{c6vy\_}, \operatorname{c6vz\_} \in \operatorname{vec}(\operatorname{c6vm\_}). \\ \quad \operatorname{c6vw\_}(\operatorname{vadd}(\operatorname{c6vm\_})(\operatorname{c6vx\_}, \operatorname{c6vy\_}), \operatorname{c6vz\_}) = \cdot(\operatorname{c6vw\_}(\operatorname{c6vx\_}, \operatorname{c6vz\_}), \operatorname{c6vw\_}(\operatorname{c6vy\_}, \operatorname{c6vz\_})) \wedge \\ \quad \operatorname{c6vw\_}(\operatorname{c6vz\_}, \operatorname{vadd}(\operatorname{c6vm\_})(\operatorname{c6vx\_}, \operatorname{c6vy\_})) = \cdot(\operatorname{c6vw\_}(\operatorname{c6vz\_}, \operatorname{c6vx\_}), \operatorname{c6vw\_}(\operatorname{c6vz\_}, \operatorname{c6vy\_})) \wedge \\ \forall \operatorname{c6vx\_}, \operatorname{c6vy\_} \in \operatorname{vec}(\operatorname{c6vm\_}). \\ \quad \operatorname{is\text{-}add\text{-}character}(\operatorname{c6vk\_}, (\lambda \operatorname{c6va\_} \in \operatorname{carr}(\operatorname{c6vk\_}).\; \operatorname{c6vw\_}(\operatorname{act}(\operatorname{c6vm\_})(\operatorname{c6va\_}, \operatorname{c6vx\_}), \operatorname{c6vy\_}))) \wedge \\ \quad \operatorname{is\text{-}add\text{-}character}(\operatorname{c6vk\_}, (\lambda \operatorname{c6va\_} \in \operatorname{carr}(\operatorname{c6vk\_}).\; \operatorname{c6vw\_}(\operatorname{c6vy\_}, \operatorname{act}(\operatorname{c6vm\_})(\operatorname{c6va\_}, \operatorname{c6vx\_})))) \end{array}
\end{equation*}

\par\medskip\noindent\textbf{Definition} (\texttt{is-e-bilinear-bicharacter}).\;For all $\operatorname{c6vk\_}$, $\operatorname{c6vm\_}$ and $\operatorname{c6vw\_}$, 
\begin{equation*}
\begin{array}{@{}l@{}} \operatorname{is\text{-}e\text{-}bilinear\text{-}bicharacter}(\operatorname{c6vk\_}, \operatorname{c6vm\_}, \operatorname{c6vw\_}) \\ \Leftrightarrow\; \operatorname{is\text{-}e\text{-}bicharacter}(\operatorname{c6vk\_}, \operatorname{c6vm\_}, \operatorname{c6vw\_}) \wedge \\ \forall \operatorname{c6va\_} \in \operatorname{carr}(\operatorname{c6vk\_}), \operatorname{c6vx\_}, \operatorname{c6vy\_} \in \operatorname{vec}(\operatorname{c6vm\_}). \\ \quad \operatorname{c6vw\_}(\operatorname{act}(\operatorname{c6vm\_})(\operatorname{c6va\_}, \operatorname{c6vx\_}), \operatorname{c6vy\_}) = \operatorname{c6vw\_}(\operatorname{c6vx\_}, \operatorname{act}(\operatorname{c6vm\_})(\operatorname{c6va\_}, \operatorname{c6vy\_})) \end{array}
\end{equation*}

\par\medskip\noindent\textbf{Definition} (\texttt{e-polar}).\;$\operatorname{e\text{-}polar}(\operatorname{c6vm\_}, \operatorname{c6vw\_}, \operatorname{c6vv\_}) :=$
\begin{equation*}
\{\operatorname{c6vh\_} \in \operatorname{vec}(\operatorname{c6vm\_}) : \forall \operatorname{c6vx\_} \in \operatorname{c6vv\_}.\; (\operatorname{c6vw\_}(\operatorname{c6vh\_}, \operatorname{c6vx\_}) = 1)\}
\end{equation*}

\par\medskip\noindent\textbf{Definition} (\texttt{norm-ball}).\;$\operatorname{norm\text{-}ball}(\operatorname{c6vm\_}, \operatorname{c6vn\_}, \operatorname{c6vr\_}) :=$ $\{\operatorname{c6vx\_} \in \operatorname{vec}(\operatorname{c6vm\_}) : (\operatorname{c6vn\_}(\operatorname{c6vx\_}) \leq \operatorname{c6vr\_})\}$
\par\medskip\noindent\textbf{Definition} (\texttt{is-e-unipolar}).\;For all $\operatorname{c6vk\_}$, $\operatorname{c6vm\_}$, $\operatorname{c6vn\_}$ and $\operatorname{c6vw\_}$, 
\begin{equation*}
\begin{array}{@{}l@{}} \operatorname{is\text{-}e\text{-}unipolar}(\operatorname{c6vk\_}, \operatorname{c6vm\_}, \operatorname{c6vn\_}, \operatorname{c6vw\_}) \\ \Leftrightarrow\; \operatorname{is\text{-}norm\text{-}on}(\operatorname{c6vk\_}, \operatorname{c6vm\_}, \operatorname{c6vn\_}) \wedge \\ \operatorname{is\text{-}e\text{-}bilinear\text{-}bicharacter}(\operatorname{c6vk\_}, \operatorname{c6vm\_}, \operatorname{c6vw\_}) \wedge \\ \operatorname{e\text{-}polar}(\operatorname{c6vm\_}, \operatorname{c6vw\_}, \operatorname{norm\text{-}ball}(\operatorname{c6vm\_}, \operatorname{c6vn\_}, 1)) = \operatorname{norm\text{-}ball}(\operatorname{c6vm\_}, \operatorname{c6vn\_}, 1) \end{array}
\end{equation*}

\par\medskip\noindent\textbf{Definition} (\texttt{is-k-bilinear-form}).\;For all $\operatorname{c6vk\_}$, $\operatorname{c6vm\_}$ and $\operatorname{c6vl\_}$, 
\begin{equation*}
\begin{array}{@{}l@{}} \operatorname{is\text{-}k\text{-}bilinear\text{-}form}(\operatorname{c6vk\_}, \operatorname{c6vm\_}, \operatorname{c6vl\_}) \\ \Leftrightarrow\; \operatorname{c6vl\_} \in ((\operatorname{vec}(\operatorname{c6vm\_}) \times \operatorname{vec}(\operatorname{c6vm\_})) \to \operatorname{carr}(\operatorname{c6vk\_})) \wedge \\ \forall \operatorname{c6vx\_}, \operatorname{c6vy\_}, \operatorname{c6vz\_} \in \operatorname{vec}(\operatorname{c6vm\_}). \\ \quad \operatorname{c6vl\_}(\operatorname{vadd}(\operatorname{c6vm\_})(\operatorname{c6vx\_}, \operatorname{c6vy\_}), \operatorname{c6vz\_}) = \operatorname{add}(\operatorname{c6vk\_})(\operatorname{c6vl\_}(\operatorname{c6vx\_}, \operatorname{c6vz\_}), \operatorname{c6vl\_}(\operatorname{c6vy\_}, \operatorname{c6vz\_})) \wedge \\ \quad \operatorname{c6vl\_}(\operatorname{c6vz\_}, \operatorname{vadd}(\operatorname{c6vm\_})(\operatorname{c6vx\_}, \operatorname{c6vy\_})) = \operatorname{add}(\operatorname{c6vk\_})(\operatorname{c6vl\_}(\operatorname{c6vz\_}, \operatorname{c6vx\_}), \operatorname{c6vl\_}(\operatorname{c6vz\_}, \operatorname{c6vy\_})) \wedge \\ \forall \operatorname{c6va\_} \in \operatorname{carr}(\operatorname{c6vk\_}), \operatorname{c6vx\_}, \operatorname{c6vy\_} \in \operatorname{vec}(\operatorname{c6vm\_}). \\ \quad \operatorname{c6vl\_}(\operatorname{act}(\operatorname{c6vm\_})(\operatorname{c6va\_}, \operatorname{c6vx\_}), \operatorname{c6vy\_}) = \operatorname{mul}(\operatorname{c6vk\_})(\operatorname{c6va\_}, \operatorname{c6vl\_}(\operatorname{c6vx\_}, \operatorname{c6vy\_})) \wedge \\ \quad \operatorname{c6vl\_}(\operatorname{c6vx\_}, \operatorname{act}(\operatorname{c6vm\_})(\operatorname{c6va\_}, \operatorname{c6vy\_})) = \operatorname{mul}(\operatorname{c6vk\_})(\operatorname{c6va\_}, \operatorname{c6vl\_}(\operatorname{c6vx\_}, \operatorname{c6vy\_})) \end{array}
\end{equation*}


\vnbbookfile{theorem-library/ccr4-chars}{theorem-library/ccr4-chars}
\noindent{\small ccr4-chars.scm -- ADDITIVE CHARACTERS OF A NORMED FIELD (CCR.pdf chapter 5: 5.12, Remark 5.14, Lemma 5.15; chapter 6: the easy half of Lemma 6.3).  Vocabulary: structure-library/ccr4-vocab.scm.  Agent CCR-B4, 2026-10-26.  Destination: theorem-library/ccr4-chars.scm.  Library: analysis.  Load window: after ...}\par
\vnbprop{ch-cc-cancel}{\vnbstmt{\forall z, u, v \in \mathbb{C}.\; ((\neg (z = 0) \wedge ((z \cdot u) = (z \cdot v))) \Rightarrow (u = v))}}
\vnbprop{ch-zero}{\vnbstmt{\begin{array}{@{}l@{}} \forall f, i. \\ \quad (\operatorname{is\text{-}normed\text{-}field}(f) \wedge \operatorname{is\text{-}add\text{-}character}(f, i)) \Rightarrow (i(\operatorname{zero}(f)) = 1) \end{array}}}
\vnbprop{ch-neg}{\vnbstmt{\begin{array}{@{}l@{}} \forall f, i.; \forall x \in \operatorname{carr}(f). \\ \quad (\operatorname{is\text{-}normed\text{-}field}(f) \wedge \operatorname{is\text{-}add\text{-}character}(f, i)) \Rightarrow ((i(\operatorname{neg}(f)(x)) \cdot i(x)) = 1) \end{array}}}
\vnbprop{ch-psi-value}{\vnbstmt{\begin{array}{@{}l@{}} \forall f, i.; \forall a, x \in \operatorname{carr}(f). \\ \quad (\operatorname{is\text{-}normed\text{-}field}(f) \wedge \operatorname{is\text{-}add\text{-}character}(f, i)) \Rightarrow  \\ \quad \quad \operatorname{duality\text{-}map}(f, i, a)(x) = i(\operatorname{mul}(f)(a, x)) \end{array}}}
\vnbprop{ch-unit-conj}{\vnbstmt{\begin{array}{@{}l@{}} \forall z, w \in \mathbb{C}. \\ \quad ((\operatorname{magnitude}(z) = 1) \wedge ((w \cdot z) = 1)) \Rightarrow (\operatorname{conjugate}(z) = w) \end{array}}}
\vnbprop{ch-psi-add}{\vnbstmt{\begin{array}{@{}l@{}} \forall f, i.; \forall a, b, c \in \operatorname{carr}(f). \\ \quad (\operatorname{is\text{-}normed\text{-}field}(f) \wedge \operatorname{is\text{-}add\text{-}character}(f, i)) \Rightarrow  \\ \quad \quad (\operatorname{duality\text{-}map}(f, i, a)(c) \cdot \operatorname{duality\text{-}map}(f, i, b)(c)) = \operatorname{duality\text{-}map}(f, i, \operatorname{add}(f)(a, b))(c) \end{array}}}
\vnbprop{ch-psi-conj}{\vnbstmt{\begin{array}{@{}l@{}} \forall f, i.; \forall a, b \in \operatorname{carr}(f). \\ \quad (\operatorname{is\text{-}normed\text{-}field}(f) \wedge \operatorname{is\text{-}add\text{-}character}(f, i)) \Rightarrow  \\ \quad \quad \operatorname{conjugate}(\operatorname{duality\text{-}map}(f, i, a)(b)) = \operatorname{duality\text{-}map}(f, i, \operatorname{neg}(f)(a))(b) \end{array}}}
\vnbprop{ch-psi-separates}{\vnbstmt{\begin{array}{@{}l@{}} \forall f, i.; \forall a \in \operatorname{carr}(f). \\ \quad (\operatorname{is\text{-}normed\text{-}field}(f) \wedge (\operatorname{is\text{-}nontrivial\text{-}character}(f, i) \wedge \neg (a = \operatorname{zero}(f)))) \Rightarrow  \\ \quad \quad \exists c \in \operatorname{carr}(f). \\ \quad \quad \quad \neg (\operatorname{duality\text{-}map}(f, i, c)(a) = \operatorname{duality\text{-}map}(f, i, c)(\operatorname{zero}(f))) \end{array}}}
\vnbprop{ccr-5-12-injective}{\vnbstmt{\begin{array}{@{}l@{}} \forall f, i.; \forall a, b \in \operatorname{carr}(f). \\ \quad (\operatorname{is\text{-}normed\text{-}field}(f) \wedge (\operatorname{is\text{-}nontrivial\text{-}character}(f, i) \wedge (\operatorname{duality\text{-}map}(f, i, a) = \operatorname{duality\text{-}map}(f, i, b)))) \Rightarrow  \\ \quad \quad a = b \end{array}}}
\vnbprop{ccr-5-15-core}{\vnbstmt{\begin{array}{@{}l@{}} \forall f, i.; \forall b, a \in \operatorname{carr}(f). \\ \quad \operatorname{is\text{-}normed\text{-}field}(f) \wedge \\ \quad \operatorname{is\text{-}add\text{-}character}(f, i) \wedge \\ \quad \neg (b = \operatorname{zero}(f)) \wedge \\ \quad \forall c \in \operatorname{carr}(f).\; ((\operatorname{fnrm}(f)(c) \leq 1) \Rightarrow (i(\operatorname{mul}(f)(b, c)) = 1)) \wedge \\ \quad \operatorname{fnrm}(f)(a) \leq \operatorname{fnrm}(f)(b) \Rightarrow \\ \quad \quad i(a) = 1 \end{array}}}
\vnbprop{ch-scale-continuous}{\vnbstmt{\begin{array}{@{}l@{}} \forall a, f.; \forall c \in \operatorname{carr}(a). \\ \quad (\operatorname{is\text{-}normed\text{-}field}(a) \wedge \operatorname{is\text{-}continuous}(\operatorname{nf\text{-}metric\text{-}space}(a), \operatorname{cc\text{-}ms}, f)) \Rightarrow  \\ \quad \quad \operatorname{is\text{-}continuous}(\operatorname{nf\text{-}metric\text{-}space}(a), \operatorname{cc\text{-}ms}, (\lambda x \in \operatorname{carr}(a).\; f(\operatorname{mul}(a)(c, x)))) \end{array}}}
\vnbprop{ccr-5-12-character}{\vnbstmt{\begin{array}{@{}l@{}} \forall f, i.; \forall a \in \operatorname{carr}(f). \\ \quad (\operatorname{is\text{-}normed\text{-}field}(f) \wedge \operatorname{is\text{-}add\text{-}character}(f, i)) \Rightarrow  \\ \quad \quad \operatorname{is\text{-}add\text{-}character}(f, \operatorname{duality\text{-}map}(f, i, a)) \end{array}}}
\vnbprop{ch-bichar-value}{\vnbstmt{\begin{array}{@{}l@{}} \forall f, i.; \forall a, x, y \in \operatorname{carr}(f). \\ \quad (\operatorname{is\text{-}normed\text{-}field}(f) \wedge \operatorname{is\text{-}add\text{-}character}(f, i)) \Rightarrow  \\ \quad \quad \operatorname{bichar\text{-}of}(f, i, a)(x, y) = i(\operatorname{mul}(f)(a, \operatorname{mul}(f)(x, y))) \end{array}}}
\vnbprop{ch-bichar-type}{\vnbstmt{\begin{array}{@{}l@{}} \forall f, i.; \forall a \in \operatorname{carr}(f). \\ \quad (\operatorname{is\text{-}normed\text{-}field}(f) \wedge \operatorname{is\text{-}add\text{-}character}(f, i)) \Rightarrow  \\ \quad \quad \operatorname{bichar\text{-}of}(f, i, a) \in ((\operatorname{carr}(f) \times \operatorname{carr}(f)) \to \operatorname{unit\text{-}circle}) \end{array}}}
\vnbprop{ccr-6-3-bicharacter}{\vnbstmt{\begin{array}{@{}l@{}} \forall f, i.; \forall a \in \operatorname{carr}(f). \\ \quad (\operatorname{is\text{-}normed\text{-}field}(f) \wedge \operatorname{is\text{-}add\text{-}character}(f, i)) \Rightarrow  \\ \quad \quad \operatorname{is\text{-}bilinear\text{-}bicharacter}(f, \operatorname{bichar\text{-}of}(f, i, a)) \wedge \\ \quad \quad \forall x, y \in \operatorname{carr}(f). \\ \quad \quad \quad \operatorname{bichar\text{-}of}(f, i, a)(x, y) = \operatorname{bichar\text{-}of}(f, i, a)(y, x) \end{array}}}
\vnbprop{ch-unit-product-continuous}{\vnbstmt{\begin{array}{@{}l@{}} \forall s, f, g. \\ \quad \operatorname{is\text{-}metric\text{-}space}(s) \wedge \\ \quad f \in (\operatorname{pts}(s) \to \operatorname{unit\text{-}circle}) \wedge \\ \quad g \in (\operatorname{pts}(s) \to \operatorname{unit\text{-}circle}) \wedge \\ \quad \operatorname{is\text{-}continuous}(s, \operatorname{cc\text{-}ms}, f) \wedge \\ \quad \operatorname{is\text{-}continuous}(s, \operatorname{cc\text{-}ms}, g) \Rightarrow \\ \quad \quad \operatorname{is\text{-}continuous}(s, \operatorname{cc\text{-}ms}, (\lambda x \in \operatorname{pts}(s).\; (f(x) \cdot g(x)))) \end{array}}}
\vnbprop{ch-product-character}{\vnbstmt{\begin{array}{@{}l@{}} \forall f, i, a. \\ \quad (\operatorname{is\text{-}normed\text{-}field}(f) \wedge (\operatorname{is\text{-}add\text{-}character}(f, i) \wedge \operatorname{is\text{-}add\text{-}character}(f, a))) \Rightarrow  \\ \quad \quad \operatorname{is\text{-}add\text{-}character}(f, (\lambda x \in \operatorname{carr}(f).\; (i(x) \cdot a(x)))) \end{array}}}
\vnbprop{ch-inverse-character}{\vnbstmt{\begin{array}{@{}l@{}} \forall f, i. \\ \quad (\operatorname{is\text{-}normed\text{-}field}(f) \wedge \operatorname{is\text{-}add\text{-}character}(f, i)) \Rightarrow  \\ \quad \quad \operatorname{is\text{-}add\text{-}character}(f, \operatorname{duality\text{-}map}(f, i, \operatorname{neg}(f)(\operatorname{one}(f)))) \wedge \\ \quad \quad \forall x \in \operatorname{carr}(f).\; ((i(x) \cdot \operatorname{duality\text{-}map}(f, i, \operatorname{neg}(f)(\operatorname{one}(f)))(x)) = 1) \end{array}}}
\vnbprop{ch-trivial-character}{\vnbstmt{\begin{array}{@{}l@{}} \forall f. \\ \quad \operatorname{is\text{-}normed\text{-}field}(f) \Rightarrow \operatorname{is\text{-}add\text{-}character}(f, (\lambda x \in \operatorname{carr}(f).\; 1)) \end{array}}}

\vnbbookfile{theorem-library/ccr1109a6-chars}{theorem-library/ccr1109a6-chars}
\noindent{\small ccr1109a6-chars.scm -- wave of 2026-11-09/10 (agent CCR-1109-A6), wired 2026-11-10 from scratchpad/ccr1109a6/ccr1109a6-chars.scm. Library: complex-analysis.}\par
\vnbprop{subgroup-character-iden}{\vnbstmt{\begin{array}{@{}l@{}} \forall g, h.; \forall c \in (h \to \operatorname{unit\text{-}circle}). \\ \quad \operatorname{is\text{-}abelian\text{-}group}(g) \wedge \\ \quad \operatorname{is\text{-}subgroup}(g, h) \wedge \\ \quad \forall u, v \in h.\; (c(\operatorname{opr}(g)(u, v)) = (c(u) \cdot c(v))) \Rightarrow \\ \quad \quad c(\operatorname{iden}(g)) = 1 \end{array}}}
\vnbprop{subgroup-character-zz-act-nn}{\vnbstmt{\begin{array}{@{}l@{}} \forall n \in \mathbb{N}, g, h.; \forall c \in (h \to \operatorname{unit\text{-}circle}), a \in h, s \in \mathbb{C}. \\ \quad \operatorname{is\text{-}abelian\text{-}group}(g) \wedge \\ \quad \operatorname{is\text{-}subgroup}(g, h) \wedge \\ \quad \forall u, v \in h.\; (c(\operatorname{opr}(g)(u, v)) = (c(u) \cdot c(v))) \wedge \\ \quad c(a) = \exp s \Rightarrow \\ \quad \quad c(\operatorname{zz\text{-}act}(g, n, a)) = \exp (n \cdot s) \end{array}}}
\vnbprop{subgroup-character-zz-act}{\vnbstmt{\begin{array}{@{}l@{}} \forall n \in \mathbb{Z}, g, h.; \forall c \in (h \to \operatorname{unit\text{-}circle}), a \in h, s \in \mathbb{C}. \\ \quad \operatorname{is\text{-}abelian\text{-}group}(g) \wedge \\ \quad \operatorname{is\text{-}subgroup}(g, h) \wedge \\ \quad \forall u, v \in h.\; (c(\operatorname{opr}(g)(u, v)) = (c(u) \cdot c(v))) \wedge \\ \quad c(a) = \exp s \Rightarrow \\ \quad \quad c(\operatorname{zz\text{-}act}(g, n, a)) = \exp (n \cdot s) \end{array}}}
\vnbprop{ag-subgroup-adjoin}{\vnbstmt{\begin{array}{@{}l@{}} \forall g, h.; \forall b \in \operatorname{carr}(g). \\ \quad (\operatorname{is\text{-}abelian\text{-}group}(g) \wedge \operatorname{is\text{-}subgroup}(g, h)) \Rightarrow  \\ \quad \quad \operatorname{is\text{-}subgroup}(g, \{y \in \operatorname{carr}(g) : \exists p \in h, k \in \mathbb{Z}.\; (y = \operatorname{opr}(g)(p, \operatorname{zz\text{-}act}(g, k, b)))\}) \wedge \\ \quad \quad h \subseteq \{y \in \operatorname{carr}(g) : \exists p \in h, k \in \mathbb{Z}.\; (y = \operatorname{opr}(g)(p, \operatorname{zz\text{-}act}(g, k, b)))\} \wedge \\ \quad \quad b \in \{y \in \operatorname{carr}(g) : \exists p \in h, k \in \mathbb{Z}.\; (y = \operatorname{opr}(g)(p, \operatorname{zz\text{-}act}(g, k, b)))\} \end{array}}}
\vnbprop{subgroup-adjoin-key}{\vnbstmt{\begin{array}{@{}l@{}} \forall g, h.; \forall c \in (h \to \operatorname{unit\text{-}circle}), b \in \operatorname{carr}(g), m \in \mathbb{N}, t \in \mathbb{C}, d \in \mathbb{Z}. \\ \quad \operatorname{is\text{-}abelian\text{-}group}(g) \wedge \\ \quad \operatorname{is\text{-}subgroup}(g, h) \wedge \\ \quad \forall u, v \in h.\; (c(\operatorname{opr}(g)(u, v)) = (c(u) \cdot c(v))) \wedge \\ \quad 1 \leq m \wedge \\ \quad \operatorname{zz\text{-}act}(g, m, b) \in h \wedge \\ \quad \forall j \in \mathbb{N}.\; (((1 \leq j) \wedge (j < m)) \Rightarrow \neg (\operatorname{zz\text{-}act}(g, j, b) \in h)) \wedge \\ \quad c(\operatorname{zz\text{-}act}(g, m, b)) = \exp (m \cdot t) \wedge \\ \quad \operatorname{zz\text{-}act}(g, d, b) \in h \Rightarrow \\ \quad \quad c(\operatorname{zz\text{-}act}(g, d, b)) = \exp (d \cdot t) \end{array}}}
\vnbprop{subgroup-adjoin-rep-unique}{\vnbstmt{\begin{array}{@{}l@{}} \forall g, h.; \forall c \in (h \to \operatorname{unit\text{-}circle}), b \in \operatorname{carr}(g), m \in \mathbb{N}, t \in \mathbb{C}, p_{1} \in h, k_{1} \in \mathbb{Z}, p_{2} \in h, k_{2} \in \mathbb{Z}. \\ \quad \operatorname{is\text{-}abelian\text{-}group}(g) \wedge \\ \quad \operatorname{is\text{-}subgroup}(g, h) \wedge \\ \quad \forall u, v \in h.\; (c(\operatorname{opr}(g)(u, v)) = (c(u) \cdot c(v))) \wedge \\ \quad 1 \leq m \wedge \\ \quad \operatorname{zz\text{-}act}(g, m, b) \in h \wedge \\ \quad \forall j \in \mathbb{N}.\; (((1 \leq j) \wedge (j < m)) \Rightarrow \neg (\operatorname{zz\text{-}act}(g, j, b) \in h)) \wedge \\ \quad c(\operatorname{zz\text{-}act}(g, m, b)) = \exp (m \cdot t) \wedge \\ \quad \operatorname{opr}(g)(p_{1}, \operatorname{zz\text{-}act}(g, k_{1}, b)) = \operatorname{opr}(g)(p_{2}, \operatorname{zz\text{-}act}(g, k_{2}, b)) \Rightarrow \\ \quad \quad (c(p_{1}) \cdot \exp (k_{1} \cdot t)) = (c(p_{2}) \cdot \exp (k_{2} \cdot t)) \end{array}}}
\vnbprop{adjoin-char-value-spec}{\vnbstmt{\begin{array}{@{}l@{}} \forall g, h.; \forall c \in (h \to \operatorname{unit\text{-}circle}), b \in \operatorname{carr}(g), m \in \mathbb{N}, t \in \mathbb{C}, p \in h, k \in \mathbb{Z}. \\ \quad \operatorname{is\text{-}abelian\text{-}group}(g) \wedge \\ \quad \operatorname{is\text{-}subgroup}(g, h) \wedge \\ \quad \forall u, v \in h.\; (c(\operatorname{opr}(g)(u, v)) = (c(u) \cdot c(v))) \wedge \\ \quad 1 \leq m \wedge \\ \quad \operatorname{zz\text{-}act}(g, m, b) \in h \wedge \\ \quad \forall j \in \mathbb{N}.\; (((1 \leq j) \wedge (j < m)) \Rightarrow \neg (\operatorname{zz\text{-}act}(g, j, b) \in h)) \wedge \\ \quad c(\operatorname{zz\text{-}act}(g, m, b)) = \exp (m \cdot t) \Rightarrow \\ \quad \quad \operatorname{adjoin\text{-}char\text{-}value}(g, h, c, b, t, \operatorname{opr}(g)(p, \operatorname{zz\text{-}act}(g, k, b))) = (c(p) \cdot \exp (k \cdot t)) \end{array}}}
\vnbprop{subgroup-character-extend}{\vnbstmt{\begin{array}{@{}l@{}} \forall g, h.; \forall c \in (h \to \operatorname{unit\text{-}circle}), b \in \operatorname{carr}(g), m \in \mathbb{N}, t \in \mathbb{C}. \\ \quad \operatorname{is\text{-}abelian\text{-}group}(g) \wedge \\ \quad \operatorname{is\text{-}subgroup}(g, h) \wedge \\ \quad \forall u, v \in h.\; (c(\operatorname{opr}(g)(u, v)) = (c(u) \cdot c(v))) \wedge \\ \quad 1 \leq m \wedge \\ \quad \operatorname{zz\text{-}act}(g, m, b) \in h \wedge \\ \quad \forall j \in \mathbb{N}.\; (((1 \leq j) \wedge (j < m)) \Rightarrow \neg (\operatorname{zz\text{-}act}(g, j, b) \in h)) \wedge \\ \quad c(\operatorname{zz\text{-}act}(g, m, b)) = \exp (m \cdot t) \wedge \\ \quad \forall r \in \mathbb{Z}.\; (\operatorname{magnitude}(\exp (r \cdot t)) = 1) \Rightarrow \\ \quad \quad (\lambda z \in \{y \in \operatorname{carr}(g) : \exists p \in h, k \in \mathbb{Z}.\; (y = \operatorname{opr}(g)(p, \operatorname{zz\text{-}act}(g, k, b)))\}.\; \\ \quad \quad \quad \operatorname{adjoin\text{-}char\text{-}value}(g, h, c, b, t, z)) \\ \quad \quad \in\; (\{y \in \operatorname{carr}(g) : \exists p \in h, k \in \mathbb{Z}.\; (y = \operatorname{opr}(g)(p, \operatorname{zz\text{-}act}(g, k, b)))\} \to \operatorname{unit\text{-}circle}) \wedge \\ \quad \quad \forall u, v \in \{y \in \operatorname{carr}(g) : \exists p \in h, k \in \mathbb{Z}.\; (y = \operatorname{opr}(g)(p, \operatorname{zz\text{-}act}(g, k, b)))\}. \\ \quad \quad \quad (\lambda z \in \{y \in \operatorname{carr}(g) : \exists p \in h, k \in \mathbb{Z}.\; (y = \operatorname{opr}(g)(p, \operatorname{zz\text{-}act}(g, k, b)))\}.\; \operatorname{adjoin\text{-}char\text{-}value}(g, h, c, b, t, z))(\operatorname{opr}(g)(u, v)) \\ \quad \quad \quad =\; (\lambda z \in \{y \in \operatorname{carr}(g) : \exists p \in h, k \in \mathbb{Z}.\; (y = \operatorname{opr}(g)(p, \operatorname{zz\text{-}act}(g, k, b)))\}.\; \operatorname{adjoin\text{-}char\text{-}value}(g, h, c, b, t, z))(u) \cdot (\lambda z \in \{y \in \operatorname{carr}(g) : \exists p \in h, k \in \mathbb{Z}.\; (y = \operatorname{opr}(g)(p, \operatorname{zz\text{-}act}(g, k, b)))\}.\; \operatorname{adjoin\text{-}char\text{-}value}(g, h, c, b, t, z))(v) \wedge \\ \quad \quad \forall x \in h. \\ \quad \quad \quad (\lambda z \in \{y \in \operatorname{carr}(g) : \exists p \in h, k \in \mathbb{Z}.\; (y = \operatorname{opr}(g)(p, \operatorname{zz\text{-}act}(g, k, b)))\}.\; \operatorname{adjoin\text{-}char\text{-}value}(g, h, c, b, t, z))(x) \\ \quad \quad \quad =\; c(x) \wedge \\ \quad \quad (\lambda z \in \{y \in \operatorname{carr}(g) : \exists p \in h, k \in \mathbb{Z}.\; (y = \operatorname{opr}(g)(p, \operatorname{zz\text{-}act}(g, k, b)))\}.\; \operatorname{adjoin\text{-}char\text{-}value}(g, h, c, b, t, z))(b) \\ \quad \quad =\; \exp t \end{array}}}
\vnbprop{unit-circle-polar}{\vnbstmt{\forall w \in \operatorname{unit\text{-}circle}.\; \exists h \in \mathbb{R}.\; (w = \exp (+i \cdot h))}}
\vnbprop{cc-exp-zz-imaginary-magnitude}{\vnbstmt{\forall h \in \mathbb{R}, r \in \mathbb{Z}.\; (\operatorname{magnitude}(\exp (r \cdot (+i \cdot h))) = 1)}}
\vnbprop{rr-divide-by-positive-nn}{\vnbstmt{\forall m \in \mathbb{N}, h \in \mathbb{R}.\; ((1 \leq m) \Rightarrow \exists s \in \mathbb{R}.\; ((m \cdot s) = h))}}
\vnbprop{subgroup-adjoin-least}{\vnbstmt{\begin{array}{@{}l@{}} \forall g, h.; \forall b \in \operatorname{carr}(g). \\ \quad (\operatorname{is\text{-}abelian\text{-}group}(g) \wedge ((|\operatorname{carr}(g)| \in \mathbb{N}) \wedge \operatorname{is\text{-}subgroup}(g, h))) \Rightarrow  \\ \quad \quad \exists m \in \mathbb{N}. \\ \quad \quad \quad 1 \leq m \wedge \\ \quad \quad \quad \operatorname{zz\text{-}act}(g, m, b) \in h \wedge \\ \quad \quad \quad \forall j \in \mathbb{N}.\; (((1 \leq j) \wedge (j < m)) \Rightarrow \neg (\operatorname{zz\text{-}act}(g, j, b) \in h)) \end{array}}}
\vnbprop{ag-trivial-subgroup-character}{\vnbstmt{\begin{array}{@{}l@{}} \forall g. \\ \quad \operatorname{is\text{-}abelian\text{-}group}(g) \Rightarrow  \\ \quad \quad \operatorname{is\text{-}subgroup}(g, \{y_{0} \in \operatorname{carr}(g) : (y_{0} = \operatorname{iden}(g))\}) \wedge \\ \quad \quad (\lambda z_{0} \in \{y_{0} \in \operatorname{carr}(g) : (y_{0} = \operatorname{iden}(g))\}.\; 1) \in (\{y_{0} \in \operatorname{carr}(g) : (y_{0} = \operatorname{iden}(g))\} \to \operatorname{unit\text{-}circle}) \wedge \\ \quad \quad \forall u, v \in \{y_{0} \in \operatorname{carr}(g) : (y_{0} = \operatorname{iden}(g))\}. \\ \quad \quad \quad (\lambda z_{0} \in \{y_{0} \in \operatorname{carr}(g) : (y_{0} = \operatorname{iden}(g))\}.\; 1)(\operatorname{opr}(g)(u, v)) \\ \quad \quad \quad =\; (\lambda z_{0} \in \{y_{0} \in \operatorname{carr}(g) : (y_{0} = \operatorname{iden}(g))\}.\; 1)(u) \cdot (\lambda z_{0} \in \{y_{0} \in \operatorname{carr}(g) : (y_{0} = \operatorname{iden}(g))\}.\; 1)(v) \end{array}}}
\vnbprop{subgroup-character-step}{\vnbstmt{\begin{array}{@{}l@{}} \forall g, h.; \forall c \in (h \to \operatorname{unit\text{-}circle}), b \in \operatorname{carr}(g). \\ \quad \operatorname{is\text{-}abelian\text{-}group}(g) \wedge \\ \quad \operatorname{is\text{-}subgroup}(g, h) \wedge \\ \quad \forall u, v \in h.\; (c(\operatorname{opr}(g)(u, v)) = (c(u) \cdot c(v))) \wedge \\ \quad |\operatorname{carr}(g)| \in \mathbb{N} \Rightarrow \\ \quad \quad \exists h_{2}. \\ \quad \quad \quad \operatorname{is\text{-}subgroup}(g, h_{2}) \wedge \\ \quad \quad \quad h \subseteq h_{2} \wedge \\ \quad \quad \quad b \in h_{2} \wedge \\ \quad \quad \quad \exists c_{2} \in (h_{2} \to \operatorname{unit\text{-}circle}). \\ \quad \quad \quad \quad (\forall u, v \in h_{2}.\; (\operatorname{c2}(\operatorname{opr}(g)(u, v)) = (\operatorname{c2}(u) \cdot \operatorname{c2}(v))) \wedge \forall x \in h.\; (\operatorname{c2}(x) = c(x))) \end{array}}}
\vnbprop{finite-ag-character-separates}{\vnbstmt{\begin{array}{@{}l@{}} \forall g.; \forall x_{0} \in \operatorname{carr}(g). \\ \quad (\operatorname{is\text{-}abelian\text{-}group}(g) \wedge ((|\operatorname{carr}(g)| \in \mathbb{N}) \wedge \neg (x_{0} = \operatorname{iden}(g)))) \Rightarrow  \\ \quad \quad \exists c \in (\operatorname{carr}(g) \to \operatorname{unit\text{-}circle}). \\ \quad \quad \quad (\forall u, v \in \operatorname{carr}(g).\; (c(\operatorname{opr}(g)(u, v)) = (c(u) \cdot c(v))) \wedge \neg (c(x_{0}) = 1)) \end{array}}}

\vnbbookfile{structure-library/padic-vocab}{structure-library/padic-vocab}
\noindent{\small padic-vocab.scm -- THE p-ADIC ABSOLUTE VALUE ON QQ (CCR.pdf section 1.2).  DEFINITIONS ONLY; the laws are in theorem-library/padic-laws.scm.  Agent PADIC, wave MAP-1026, 2026-10-26.  Destination: structure-library/padic-vocab.scm.  Library: analysis.  Load window: after structure-library/zz-divisibility, normed-field, ...}\par
\par\medskip\noindent\textbf{Definition} (\texttt{zz-prime}).\;For every $\operatorname{cpp\_}$, 
\begin{equation*}
\begin{array}{@{}l@{}} \operatorname{zz\text{-}prime}(\operatorname{cpp\_}) \\ \Leftrightarrow\; \operatorname{cpp\_} \in \mathbb{Z} \wedge \\ 1 < \operatorname{cpp\_} \wedge \\ \forall \operatorname{cpd\_} \in \mathbb{Z}. \\ \quad \operatorname{zz\text{-}divides}(\operatorname{cpd\_}, \operatorname{cpp\_}) \Rightarrow  \\ \quad \quad (\operatorname{cpd\_} = 1) \vee ((\operatorname{cpd\_} = -(1)) \vee ((\operatorname{cpd\_} = \operatorname{cpp\_}) \vee (\operatorname{cpd\_} = -(\operatorname{cpp\_})))) \end{array}
\end{equation*}

\par\medskip\noindent\textbf{Definition} (\texttt{zz-ord-is}).\;For all $\operatorname{cpp\_}$, $\operatorname{cpn\_}$ and $\operatorname{cpk\_}$, 
\begin{equation*}
\begin{array}{@{}l@{}} \operatorname{zz\text{-}ord\text{-}is}(\operatorname{cpp\_}, \operatorname{cpn\_}, \operatorname{cpk\_}) \\ \Leftrightarrow\; \operatorname{zz\text{-}divides}(\operatorname{power}(\operatorname{cpp\_}, \operatorname{cpk\_}), \operatorname{cpn\_}) \wedge \neg \operatorname{zz\text{-}divides}(\operatorname{power}(\operatorname{cpp\_}, \operatorname{succ}(\operatorname{cpk\_})), \operatorname{cpn\_}) \end{array}
\end{equation*}

\par\medskip\noindent\textbf{Definition} (\texttt{padic-rel}).\;For all $\operatorname{cpp\_}$, $\operatorname{cpx\_}$ and $\operatorname{cpv\_}$, 
\begin{equation*}
\begin{array}{@{}l@{}} \operatorname{padic\text{-}rel}(\operatorname{cpp\_}, \operatorname{cpx\_}, \operatorname{cpv\_}) \\ \Leftrightarrow\; \operatorname{cpv\_} \in \mathbb{R} \wedge \\ (\operatorname{cpx\_} = 0) \wedge (\operatorname{cpv\_} = 0) \vee \\ \neg (\operatorname{cpx\_} = 0) \wedge \\ \exists \operatorname{cpa\_}, \operatorname{cpb\_} \in \mathbb{Z}. \\ \quad \neg (\operatorname{cpb\_} = 0) \wedge \\ \quad \cdot(\operatorname{cpx\_}, \operatorname{cpb\_}) = \operatorname{cpa\_} \wedge \\ \quad \exists \operatorname{cpi\_}, \operatorname{cpj\_} \in \mathbb{N}. \\ \quad \quad (\operatorname{zz\text{-}ord\text{-}is}(\operatorname{cpp\_}, \operatorname{cpa\_}, \operatorname{cpi\_}) \wedge (\operatorname{zz\text{-}ord\text{-}is}(\operatorname{cpp\_}, \operatorname{cpb\_}, \operatorname{cpj\_}) \wedge (\cdot(\operatorname{cpv\_}, \operatorname{power}(2, \operatorname{cpi\_})) = \operatorname{power}(2, \operatorname{cpj\_})))) \end{array}
\end{equation*}

\par\medskip\noindent\textbf{Definition} (\texttt{padic-abs}).\;$\operatorname{padic\text{-}abs}(\operatorname{cpp\_}) :=$
\begin{equation*}
(\lambda \operatorname{cpx\_} \in \mathbb{Q}.\; (\iota\, \operatorname{cpv\_}.\;((\operatorname{cpv\_} \in \mathbb{R}) \wedge \operatorname{padic\text{-}rel}(\operatorname{cpp\_}, \operatorname{cpx\_}, \operatorname{cpv\_}))))
\end{equation*}

\par\medskip\noindent\textbf{Definition} (\texttt{qq-padic}).\;$\operatorname{qq\text{-}padic}(\operatorname{cpp\_}) :=$
\begin{equation*}
{[}\mathbb{Q}, (\lambda {[}\operatorname{x\_}, \operatorname{y\_}] \in (\mathbb{Q} \times \mathbb{Q}).\; (\operatorname{x\_} + \operatorname{y\_})), (\lambda {[}\operatorname{x\_}, \operatorname{y\_}] \in (\mathbb{Q} \times \mathbb{Q}).\; (\operatorname{x\_} \cdot \operatorname{y\_})), (\lambda \operatorname{x\_} \in \mathbb{Q}.\; -\operatorname{x\_}), 0, 1, \operatorname{padic\text{-}abs}(\operatorname{cpp\_})]
\end{equation*}

\par\medskip\noindent\textbf{Definition} (\texttt{qq-nv}).\;$\operatorname{qq\text{-}nv}(\operatorname{cpf\_}) :=$
\begin{equation*}
{[}\mathbb{Q}, (\lambda {[}\operatorname{x\_}, \operatorname{y\_}] \in (\mathbb{Q} \times \mathbb{Q}).\; (\operatorname{x\_} + \operatorname{y\_})), (\lambda {[}\operatorname{x\_}, \operatorname{y\_}] \in (\mathbb{Q} \times \mathbb{Q}).\; (\operatorname{x\_} \cdot \operatorname{y\_})), (\lambda \operatorname{x\_} \in \mathbb{Q}.\; -\operatorname{x\_}), 0, 1, \operatorname{cpf\_}]
\end{equation*}

\par\medskip\noindent\textbf{Definition} (\texttt{is-nontrivial-nv}).\;For every $\operatorname{cpk\_}$, 
\begin{equation*}
\begin{array}{@{}l@{}} \operatorname{is\text{-}nontrivial\text{-}nv}(\operatorname{cpk\_}) \\ \Leftrightarrow\; \exists \operatorname{cpa\_} \in \operatorname{carr}(\operatorname{cpk\_}). \\ \quad \neg (\operatorname{fnrm}(\operatorname{cpk\_})(\operatorname{cpa\_}) = 0) \wedge \neg (\operatorname{fnrm}(\operatorname{cpk\_})(\operatorname{cpa\_}) = 1) \end{array}
\end{equation*}

\par\medskip\noindent\textbf{Definition} (\texttt{padic-char}).\;$\operatorname{padic\text{-}char}(\operatorname{cpp\_}) :=$
\begin{equation*}
\begin{array}{@{}l@{}} (\lambda \operatorname{cpx\_} \in \mathbb{Q}.\; \\ \quad (\iota\, \operatorname{cpw\_}.\;((\operatorname{cpw\_} \in \mathbb{C}) \wedge \exists \operatorname{cpr\_} \in \mathbb{Q}.\; (\operatorname{padic\text{-}frac\text{-}rep}(\operatorname{cpp\_}, \operatorname{cpx\_}, \operatorname{cpr\_}) \wedge (\operatorname{cpw\_} = \exp (((2 \cdot \pi) \cdot +i) \cdot \operatorname{cpr\_})))))) \end{array}
\end{equation*}

\par\medskip\noindent\textbf{Definition} (\texttt{padic-frac-rep}).\;For all $\operatorname{cpp\_}$, $\operatorname{cpx\_}$ and $\operatorname{cpr\_}$, 
\begin{equation*}
\begin{array}{@{}l@{}} \operatorname{padic\text{-}frac\text{-}rep}(\operatorname{cpp\_}, \operatorname{cpx\_}, \operatorname{cpr\_}) \\ \Leftrightarrow\; \operatorname{cpr\_} \in \mathbb{Q} \wedge \\ \exists \operatorname{cpz\_} \in \mathbb{Z}, \operatorname{cpm\_} \in \mathbb{N}.\; (\cdot(\operatorname{cpr\_}, \operatorname{power}(\operatorname{cpp\_}, \operatorname{cpm\_})) = \operatorname{cpz\_}) \wedge \\ \operatorname{padic\text{-}abs}(\operatorname{cpp\_})(+(\operatorname{cpx\_}, -(\operatorname{cpr\_}))) \leq 1 \end{array}
\end{equation*}


\vnbbookfile{structure-library/qp-vocab}{structure-library/qp-vocab}
\noindent{\small qp-vocab.scm -- WAVE 2 (2026-10-27..30), wired 2026-10-30 from scratchpad/ccrqp/qp-vocab.scm. Library: complex-analysis.}\par
\par\medskip\noindent\textbf{Definition} (\texttt{qp-ms}).\;$\operatorname{qp\text{-}ms}(\operatorname{qsp\_}) :=$ $\operatorname{nf\text{-}metric\text{-}space}(\operatorname{qq\text{-}padic}(\operatorname{qsp\_}))$
\par\medskip\noindent\textbf{Definition} (\texttt{qp}).\;$\operatorname{qp}(\operatorname{qsp\_}) :=$ $\operatorname{completion}(\operatorname{nf\text{-}metric\text{-}space}(\operatorname{qq\text{-}padic}(\operatorname{qsp\_})))$
\par\medskip\noindent\textbf{Definition} (\texttt{qp-cs}).\;$\operatorname{qp\text{-}cs}(\operatorname{qsp\_}) :=$ $\operatorname{cseq}(\operatorname{nf\text{-}metric\text{-}space}(\operatorname{qq\text{-}padic}(\operatorname{qsp\_})))$
\par\medskip\noindent\textbf{Definition} (\texttt{qp-null}).\;$\operatorname{qp\text{-}null}(\operatorname{qsp\_}) :=$
\begin{equation*}
\{\operatorname{qsu\_} \in \operatorname{qp\text{-}cs}(\operatorname{qsp\_}) : \operatorname{cseq\text{-}equiv}(\operatorname{nf\text{-}metric\text{-}space}(\operatorname{qq\text{-}padic}(\operatorname{qsp\_})), \operatorname{qsu\_}, (\lambda \operatorname{qsn\_} \in \mathbb{N}.\; 0))\}
\end{equation*}

\par\medskip\noindent\textbf{Definition} (\texttt{qp-csr}).\;$\operatorname{qp\text{-}csr}(\operatorname{qsp\_}) :=$
\begin{equation*}
{[}\operatorname{qp\text{-}cs}(\operatorname{qsp\_}), \operatorname{qp\text{-}sadd}(\operatorname{qsp\_}), \operatorname{qp\text{-}smul}(\operatorname{qsp\_}), \operatorname{qp\text{-}sneg}(\operatorname{qsp\_}), (\lambda \operatorname{qsn\_} \in \mathbb{N}.\; 0), (\lambda \operatorname{qsn\_} \in \mathbb{N}.\; 1), \operatorname{qp\text{-}null}(\operatorname{qsp\_})]
\end{equation*}

\par\medskip\noindent\textbf{Definition} (\texttt{qp-ring}).\;$\operatorname{qp\text{-}ring}(\operatorname{qsp\_}) :=$ $\operatorname{ringoid\text{-}quotient\text{-}ring}(\operatorname{qp\text{-}csr}(\operatorname{qsp\_}))$
\par\medskip\noindent\textbf{Definition} (\texttt{qp-abs}).\;$\operatorname{qp\text{-}abs}(\operatorname{qsp\_}) :=$
\begin{equation*}
\begin{array}{@{}l@{}} (\lambda \operatorname{qsx\_} \in \operatorname{quotient}(\operatorname{cauchy\text{-}setoid}(\operatorname{nf\text{-}metric\text{-}space}(\operatorname{qq\text{-}padic}(\operatorname{qsp\_})))).\; \\ \quad \operatorname{dist}(\operatorname{completion}(\operatorname{nf\text{-}metric\text{-}space}(\operatorname{qq\text{-}padic}(\operatorname{qsp\_}))))(\operatorname{qsx\_}, \\ \quad \quad \operatorname{class}(\operatorname{cauchy\text{-}setoid}(\operatorname{nf\text{-}metric\text{-}space}(\operatorname{qq\text{-}padic}(\operatorname{qsp\_}))), (\lambda \operatorname{qsn\_} \in \mathbb{N}.\; 0)))) \end{array}
\end{equation*}

\par\medskip\noindent\textbf{Definition} (\texttt{qp-field}).\;$\operatorname{qp\text{-}field}(\operatorname{qsp\_}) :=$
\begin{equation*}
{[}\operatorname{carr}(\operatorname{qp\text{-}ring}(\operatorname{qsp\_})), \operatorname{add}(\operatorname{qp\text{-}ring}(\operatorname{qsp\_})), \operatorname{mul}(\operatorname{qp\text{-}ring}(\operatorname{qsp\_})), \operatorname{neg}(\operatorname{qp\text{-}ring}(\operatorname{qsp\_})), \operatorname{zero}(\operatorname{qp\text{-}ring}(\operatorname{qsp\_})), \operatorname{one}(\operatorname{qp\text{-}ring}(\operatorname{qsp\_})), \operatorname{qp\text{-}abs}(\operatorname{qsp\_})]
\end{equation*}

\par\medskip\noindent\textbf{Definition} (\texttt{zp}).\;$\operatorname{zp}(\operatorname{qsp\_}) :=$ $\operatorname{nv\text{-}integers}(\operatorname{qp\text{-}field}(\operatorname{qsp\_}))$
\par\medskip\noindent\textbf{Definition} (\texttt{qp-sadd}).\;$\operatorname{qp\text{-}sadd}(\operatorname{qsp\_}) :=$
\begin{equation*}
\begin{array}{@{}l@{}} (\lambda {[}\operatorname{qsu\_}, \operatorname{qsv\_}] \in (\operatorname{qp\text{-}cs}(\operatorname{qsp\_}) \times \operatorname{qp\text{-}cs}(\operatorname{qsp\_})).\; \\ \quad (\lambda \operatorname{qsn\_} \in \mathbb{N}.\; (\operatorname{qsu\_}(\operatorname{qsn\_}) + \operatorname{qsv\_}(\operatorname{qsn\_})))) \end{array}
\end{equation*}

\par\medskip\noindent\textbf{Definition} (\texttt{qp-smul}).\;$\operatorname{qp\text{-}smul}(\operatorname{qsp\_}) :=$
\begin{equation*}
\begin{array}{@{}l@{}} (\lambda {[}\operatorname{qsu\_}, \operatorname{qsv\_}] \in (\operatorname{qp\text{-}cs}(\operatorname{qsp\_}) \times \operatorname{qp\text{-}cs}(\operatorname{qsp\_})).\; \\ \quad (\lambda \operatorname{qsn\_} \in \mathbb{N}.\; (\operatorname{qsu\_}(\operatorname{qsn\_}) \cdot \operatorname{qsv\_}(\operatorname{qsn\_})))) \end{array}
\end{equation*}

\par\medskip\noindent\textbf{Definition} (\texttt{qp-sneg}).\;$\operatorname{qp\text{-}sneg}(\operatorname{qsp\_}) :=$
\begin{equation*}
(\lambda \operatorname{qsu\_} \in \operatorname{qp\text{-}cs}(\operatorname{qsp\_}).\; (\lambda \operatorname{qsn\_} \in \mathbb{N}.\; -\operatorname{qsu\_}(\operatorname{qsn\_})))
\end{equation*}

\par\medskip\noindent\textbf{Definition} (\texttt{qp-ord}).\;$\operatorname{qp\text{-}ord}(\operatorname{qsp\_}) :=$
\begin{equation*}
(\lambda \operatorname{qox\_} \in \operatorname{carr}(\operatorname{qp\text{-}field}(\operatorname{qsp\_})).\; (\log\left(\operatorname{qp\text{-}abs}(\operatorname{qsp\_})(\operatorname{qox\_})\right) \cdot {\log\left({2}^{\ast\!-1}\right)}^{\ast\!-1}))
\end{equation*}

\par\medskip\noindent\textbf{Definition} (\texttt{qp-ppow}).\;$\operatorname{qp\text{-}ppow}(\operatorname{qsp\_}, \operatorname{qsk\_}) :=$
\begin{equation*}
(\iota\, \operatorname{qpy\_}.\;((\operatorname{qpy\_} \in \operatorname{carr}(\operatorname{qp\text{-}field}(\operatorname{qsp\_}))) \wedge (\operatorname{mul}(\operatorname{qp\text{-}field}(\operatorname{qsp\_}))(\operatorname{qpy\_}, \operatorname{embed}(\operatorname{nf\text{-}metric\text{-}space}(\operatorname{qq\text{-}padic}(\operatorname{qsp\_})))(\operatorname{power}(\operatorname{qsp\_}, \operatorname{max}(-\operatorname{qsk\_}, 0)))) = \operatorname{embed}(\operatorname{nf\text{-}metric\text{-}space}(\operatorname{qq\text{-}padic}(\operatorname{qsp\_})))(\operatorname{power}(\operatorname{qsp\_}, \operatorname{max}(\operatorname{qsk\_}, 0))))))
\end{equation*}

\par\medskip\noindent\textbf{Definition} (\texttt{fp-ringoid}).\;$\operatorname{fp\text{-}ringoid}(\operatorname{qsp\_}) :=$
\begin{equation*}
{[}\mathbb{Z}, (\lambda {[}\operatorname{x\_}, \operatorname{y\_}] \in (\mathbb{Z} \times \mathbb{Z}).\; (\operatorname{x\_} + \operatorname{y\_})), (\lambda {[}\operatorname{x\_}, \operatorname{y\_}] \in (\mathbb{Z} \times \mathbb{Z}).\; (\operatorname{x\_} \cdot \operatorname{y\_})), (\lambda \operatorname{x\_} \in \mathbb{Z}.\; -\operatorname{x\_}), 0, 1, \{\operatorname{fpn\_} \in \mathbb{Z} : \operatorname{zz\text{-}divides}(\operatorname{qsp\_}, \operatorname{fpn\_})\}]
\end{equation*}

\par\medskip\noindent\textbf{Definition} (\texttt{fp}).\;$\operatorname{fp}(\operatorname{qsp\_}) :=$ $\operatorname{ringoid\text{-}quotient\text{-}ring}(\operatorname{fp\text{-}ringoid}(\operatorname{qsp\_}))$
\par\medskip\noindent\textbf{Definition} (\texttt{qp-psum}).\;$\operatorname{qp\text{-}psum}(\operatorname{qsp\_}, \operatorname{qta\_}, \operatorname{qtn\_}, \operatorname{qtm\_}) :=$
\begin{equation*}
\operatorname{sum\text{-}ag}(\operatorname{normed\text{-}field\text{-}additive\text{-}ag}(\operatorname{qp\text{-}field}(\operatorname{qsp\_})), \operatorname{qp\text{-}term}(\operatorname{qsp\_}, \operatorname{qta\_}, \operatorname{qtn\_}), \operatorname{qtm\_})
\end{equation*}

\par\medskip\noindent\textbf{Definition} (\texttt{qp-digit}).\;$\operatorname{qp\text{-}digit}(\operatorname{qsp\_}, \operatorname{qdy\_}) :=$
\begin{equation*}
(\iota\, \operatorname{qdd\_}.\;((\operatorname{qdd\_} \in \mathbb{Z}) \wedge ((0 \leq \operatorname{qdd\_}) \wedge ((\operatorname{qdd\_} < \operatorname{qsp\_}) \wedge (\operatorname{qp\text{-}abs}(\operatorname{qsp\_})(\operatorname{add}(\operatorname{qp\text{-}field}(\operatorname{qsp\_}))(\operatorname{qdy\_}, \operatorname{neg}(\operatorname{qp\text{-}field}(\operatorname{qsp\_}))(\operatorname{embed}(\operatorname{nf\text{-}metric\text{-}space}(\operatorname{qq\text{-}padic}(\operatorname{qsp\_})))(\operatorname{qdd\_})))) < 1)))))
\end{equation*}

\par\medskip\noindent\textbf{Definition} (\texttt{qp-term}).\;$\operatorname{qp\text{-}term}(\operatorname{qsp\_}, \operatorname{qta\_}, \operatorname{qtn\_}) :=$
\begin{equation*}
\begin{array}{@{}l@{}} (\lambda \operatorname{qtj\_} \in \mathbb{N}.\; \\ \quad \operatorname{mul}(\operatorname{qp\text{-}field}(\operatorname{qsp\_}))(\operatorname{embed}(\operatorname{nf\text{-}metric\text{-}space}(\operatorname{qq\text{-}padic}(\operatorname{qsp\_})))(\operatorname{qta\_}((\operatorname{qtn\_} + \operatorname{qtj\_}))), \\ \quad \quad \operatorname{qp\text{-}ppow}(\operatorname{qsp\_}, (\operatorname{qtn\_} + \operatorname{qtj\_})))) \end{array}
\end{equation*}

\par\medskip\noindent\textbf{Definition} (\texttt{is-qp-digits}).\;For all $\operatorname{qsp\_}$, $\operatorname{qta\_}$ and $\operatorname{qtn\_}$, 
\begin{equation*}
\begin{array}{@{}l@{}} \operatorname{is\text{-}qp\text{-}digits}(\operatorname{qsp\_}, \operatorname{qta\_}, \operatorname{qtn\_}) \\ \Leftrightarrow\; \operatorname{qta\_} \in (\mathbb{Z} \to \mathbb{Z}) \wedge \\ \forall \operatorname{qtk\_} \in \mathbb{Z}. \\ \quad (0 \leq \operatorname{qta\_}(\operatorname{qtk\_})) \wedge ((\operatorname{qta\_}(\operatorname{qtk\_}) < \operatorname{qsp\_}) \wedge ((\operatorname{qtk\_} < \operatorname{qtn\_}) \Rightarrow (\operatorname{qta\_}(\operatorname{qtk\_}) = 0))) \end{array}
\end{equation*}

\par\medskip\noindent\textbf{Definition} (\texttt{is-qp-expansion}).\;For all $\operatorname{qsp\_}$, $\operatorname{qta\_}$, $\operatorname{qtn\_}$ and $\operatorname{qtx\_}$, 
\begin{equation*}
\begin{array}{@{}l@{}} \operatorname{is\text{-}qp\text{-}expansion}(\operatorname{qsp\_}, \operatorname{qta\_}, \operatorname{qtn\_}, \operatorname{qtx\_}) \\ \Leftrightarrow\; \operatorname{is\text{-}qp\text{-}digits}(\operatorname{qsp\_}, \operatorname{qta\_}, \operatorname{qtn\_}) \wedge \\ \operatorname{converges\text{-}to}(\operatorname{qp}(\operatorname{qsp\_}), (\lambda \operatorname{qtm\_} \in \mathbb{N}.\; \operatorname{qp\text{-}psum}(\operatorname{qsp\_}, \operatorname{qta\_}, \operatorname{qtn\_}, \operatorname{qtm\_})), \operatorname{qtx\_}) \end{array}
\end{equation*}

\par\medskip\noindent\textbf{Definition} (\texttt{qp-lambda-at}).\;$\operatorname{qp\text{-}lambda\text{-}at}(\operatorname{qsp\_}, \operatorname{qlx\_}) :=$
\begin{equation*}
(\iota\, \operatorname{qlw\_}.\;((\operatorname{qlw\_} \in \mathbb{C}) \wedge \exists \operatorname{qlu\_} \in \mathbb{Q}.\; ((\operatorname{qp\text{-}abs}(\operatorname{qsp\_})(\operatorname{add}(\operatorname{qp\text{-}field}(\operatorname{qsp\_}))(\operatorname{qlx\_}, \operatorname{neg}(\operatorname{qp\text{-}field}(\operatorname{qsp\_}))(\operatorname{embed}(\operatorname{nf\text{-}metric\text{-}space}(\operatorname{qq\text{-}padic}(\operatorname{qsp\_})))(\operatorname{qlu\_})))) \leq 1) \wedge (\operatorname{qlw\_} = \operatorname{padic\text{-}char}(\operatorname{qsp\_})(\operatorname{qlu\_})))))
\end{equation*}

\par\medskip\noindent\textbf{Definition} (\texttt{qp-lambda}).\;$\operatorname{qp\text{-}lambda}(\operatorname{qsp\_}) :=$
\begin{equation*}
(\lambda \operatorname{qly\_} \in \operatorname{carr}(\operatorname{qp\text{-}field}(\operatorname{qsp\_})).\; \operatorname{qp\text{-}lambda\text{-}at}(\operatorname{qsp\_}, \operatorname{qly\_}))
\end{equation*}


\vnbbookfile{structure-library/pruefer-vocab}{structure-library/pruefer-vocab}
\noindent{\small pruefer-vocab.scm -- WAVE 2 (2026-10-27..30), wired 2026-10-30 from scratchpad/ccrpr/pruefer-vocab.scm. Library: complex-analysis.}\par
\par\medskip\noindent\textbf{Definition} (\texttt{pruefer}).\;$\operatorname{pruefer}(\operatorname{prp\_}) :=$
\begin{equation*}
\operatorname{quotient\text{-}ag}(\operatorname{normed\text{-}field\text{-}additive\text{-}ag}(\operatorname{qp\text{-}field}(\operatorname{prp\_})), \operatorname{zp}(\operatorname{prp\_}))
\end{equation*}

\par\medskip\noindent\textbf{Definition} (\texttt{is-separating-character}).\;For all $\operatorname{sck\_}$ and $\operatorname{scl\_}$, 
\begin{equation*}
\begin{array}{@{}l@{}} \operatorname{is\text{-}separating\text{-}character}(\operatorname{sck\_}, \operatorname{scl\_}) \\ \Leftrightarrow\; \operatorname{is\text{-}add\text{-}character}(\operatorname{sck\_}, \operatorname{scl\_}) \wedge \\ \forall \operatorname{scx\_} \in \operatorname{carr}(\operatorname{sck\_}).\; ((\operatorname{scl\_}(\operatorname{scx\_}) = 1) \Leftrightarrow (\operatorname{scx\_} \in \operatorname{nv\text{-}integers}(\operatorname{sck\_}))) \end{array}
\end{equation*}

\par\medskip\noindent\textbf{Definition} (\texttt{mu-pinf}).\;$\operatorname{mu\text{-}pinf}(\operatorname{prp\_}) :=$ $\{\operatorname{mpz\_} \in \mathbb{C} : \exists \operatorname{mpk\_} \in \mathbb{N}.\; (\operatorname{power}(\operatorname{mpz\_}, \operatorname{power}(\operatorname{prp\_}, \operatorname{mpk\_})) = 1)\}$
\par\medskip\noindent\textbf{Definition} (\texttt{primitive-roots}).\;$\operatorname{primitive\text{-}roots}(\operatorname{prn\_}) :=$
\begin{equation*}
\{\operatorname{prz\_} \in \mathbb{C} : ((\operatorname{power}(\operatorname{prz\_}, \operatorname{prn\_}) = 1) \wedge \forall \operatorname{prm\_} \in \mathbb{N}.\; ((1 \leq \operatorname{prm\_}) \Rightarrow ((\operatorname{prm\_} < \operatorname{prn\_}) \Rightarrow \neg (\operatorname{power}(\operatorname{prz\_}, \operatorname{prm\_}) = 1))))\}
\end{equation*}


\vnbbookfile{theorem-library/padic-laws}{theorem-library/padic-laws}
\noindent{\small padic-laws.scm -- THE p-ADIC ABSOLUTE VALUE ON QQ AND CCR.pdf 1.19-1.21, 1.24, and Remark 5.21 (the p-adic character of Q; that section cites ch-cc-cancel, so theorem-library/ccr4-chars must load first).  Vocabulary: structure-library/padic-vocab.scm.  Agent PADIC, wave MAP-1026, 2026-10-26.  Destination: ...}\par
\vnbprop{pz-power-in-zz-ind}{\vnbstmt{\forall n \in \mathbb{N}, p \in \mathbb{Z}.\; (\operatorname{power}(p, n) \in \mathbb{Z})}}
\vnbprop{pz-power-in-zz}{\vnbstmt{\forall p \in \mathbb{Z}, n \in \mathbb{N}.\; (\operatorname{power}(p, n) \in \mathbb{Z})}}
\vnbprop{pz-power-ge-ind}{\vnbstmt{\forall n \in \mathbb{N}, p \in \mathbb{Z}.\; ((2 \leq p) \Rightarrow ((n + 1) \leq \operatorname{power}(p, n)))}}
\vnbprop{pz-power-ge}{\vnbstmt{\forall p \in \mathbb{Z}, n \in \mathbb{N}.\; ((2 \leq p) \Rightarrow ((n + 1) \leq \operatorname{power}(p, n)))}}
\vnbprop{pz-dvd-mul}{\vnbstmt{\forall d, n.; \forall c \in \mathbb{Z}.\; (\operatorname{zz\text{-}divides}(d, n) \Rightarrow \operatorname{zz\text{-}divides}(d, (c \cdot n)))}}
\vnbprop{pz-dvd-add}{\vnbstmt{\begin{array}{@{}l@{}} \forall d, a, b. \\ \quad (\operatorname{zz\text{-}divides}(d, a) \wedge \operatorname{zz\text{-}divides}(d, b)) \Rightarrow \operatorname{zz\text{-}divides}(d, (a + b)) \end{array}}}
\vnbprop{pz-dvd-refl}{\vnbstmt{\forall d \in \mathbb{Z}.\; \operatorname{zz\text{-}divides}(d, d)}}
\vnbprop{pz-dvd-zero}{\vnbstmt{\forall d \in \mathbb{Z}.\; \operatorname{zz\text{-}divides}(d, 0)}}
\vnbprop{pz-dvd-neg}{\vnbstmt{\forall d, n.\; (\operatorname{zz\text{-}divides}(d, n) \Rightarrow \operatorname{zz\text{-}divides}(d, -n))}}
\vnbprop{pz-power-dvd}{\vnbstmt{\forall p \in \mathbb{Z}, k, j \in \mathbb{N}.\; ((k \leq j) \Rightarrow \operatorname{zz\text{-}divides}(\operatorname{power}(p, k), \operatorname{power}(p, j)))}}
\vnbprop{pz-dvd-bound}{\vnbstmt{\forall d, n.\; ((\operatorname{zz\text{-}divides}(d, n) \wedge \neg (n = 0)) \Rightarrow (\lvert d \rvert \leq \lvert n \rvert))}}
\vnbprop{pz-euclid}{\vnbstmt{\begin{array}{@{}l@{}} \forall p.; \forall a, b \in \mathbb{Z}. \\ \quad (\operatorname{zz\text{-}prime}(p) \wedge \operatorname{zz\text{-}divides}(p, (a \cdot b))) \Rightarrow  \\ \quad \quad \operatorname{zz\text{-}divides}(p, a) \vee \operatorname{zz\text{-}divides}(p, b) \end{array}}}
\vnbprop{pz-zz-cancel}{\vnbstmt{\forall q, x, y \in \mathbb{Z}.\; ((\neg (q = 0) \wedge ((q \cdot x) = (q \cdot y))) \Rightarrow (x = y))}}
\vnbprop{pz-ord-exists}{\vnbstmt{\begin{array}{@{}l@{}} \forall p.; \forall n \in \mathbb{Z}. \\ \quad (\operatorname{zz\text{-}prime}(p) \wedge \neg (n = 0)) \Rightarrow \exists k \in \mathbb{N}.\; \operatorname{zz\text{-}ord\text{-}is}(p, n, k) \end{array}}}
\vnbprop{pz-ord-unique}{\vnbstmt{\begin{array}{@{}l@{}} \forall p, n.; \forall i, j \in \mathbb{N}. \\ \quad (\operatorname{zz\text{-}prime}(p) \wedge (\operatorname{zz\text{-}ord\text{-}is}(p, n, i) \wedge \operatorname{zz\text{-}ord\text{-}is}(p, n, j))) \Rightarrow (i = j) \end{array}}}
\vnbprop{pz-ord-mult}{\vnbstmt{\begin{array}{@{}l@{}} \forall p.; \forall a, b \in \mathbb{Z}, i, j \in \mathbb{N}. \\ \quad (\operatorname{zz\text{-}prime}(p) \wedge (\operatorname{zz\text{-}ord\text{-}is}(p, a, i) \wedge \operatorname{zz\text{-}ord\text{-}is}(p, b, j))) \Rightarrow  \\ \quad \quad \operatorname{zz\text{-}ord\text{-}is}(p, (a \cdot b), (i + j)) \end{array}}}
\vnbprop{pz-ord-sum}{\vnbstmt{\begin{array}{@{}l@{}} \forall p.; \forall a, b \in \mathbb{Z}, i, j, s, m \in \mathbb{N}. \\ \quad \operatorname{zz\text{-}prime}(p) \wedge \\ \quad \operatorname{zz\text{-}ord\text{-}is}(p, a, i) \wedge \\ \quad \operatorname{zz\text{-}ord\text{-}is}(p, b, j) \wedge \\ \quad \operatorname{zz\text{-}ord\text{-}is}(p, (a + b), s) \wedge \\ \quad m \leq i \wedge \\ \quad m \leq j \Rightarrow \\ \quad \quad m \leq s \end{array}}}
\vnbprop{pz-ord-one}{\vnbstmt{\forall p.\; (\operatorname{zz\text{-}prime}(p) \Rightarrow \operatorname{zz\text{-}ord\text{-}is}(p, 1, 0))}}
\vnbprop{pz-ord-p}{\vnbstmt{\forall p.\; (\operatorname{zz\text{-}prime}(p) \Rightarrow \operatorname{zz\text{-}ord\text{-}is}(p, p, 1))}}
\vnbprop{pz-rr-cancel-pos}{\vnbstmt{\forall a, b, c \in \mathbb{R}.\; (((0 < c) \wedge ((a \cdot c) = (b \cdot c))) \Rightarrow (a = b))}}
\vnbprop{pz-rel-exists}{\vnbstmt{\forall p.; \forall x \in \mathbb{Q}.\; (\operatorname{zz\text{-}prime}(p) \Rightarrow \exists w.\; \operatorname{padic\text{-}rel}(p, x, w))}}
\vnbprop{pz-rel-unique}{\vnbstmt{\begin{array}{@{}l@{}} \forall p.; \forall x \in \mathbb{Q}, v, w. \\ \quad (\operatorname{zz\text{-}prime}(p) \wedge (\operatorname{padic\text{-}rel}(p, x, v) \wedge \operatorname{padic\text{-}rel}(p, x, w))) \Rightarrow (v = w) \end{array}}}
\vnbprop{pz-rel-iota}{\vnbstmt{\begin{array}{@{}l@{}} \forall p.; \forall x \in \mathbb{Q}. \\ \quad \operatorname{zz\text{-}prime}(p) \Rightarrow \operatorname{padic\text{-}rel}(p, x, (\iota\, v.\;((v \in \mathbb{R}) \wedge \operatorname{padic\text{-}rel}(p, x, v)))) \end{array}}}
\vnbprop{pz-abs-type}{\vnbstmt{\forall p.\; (\operatorname{zz\text{-}prime}(p) \Rightarrow (\operatorname{padic\text{-}abs}(p) \in (\mathbb{Q} \to \mathbb{R})))}}
\vnbprop{pz-abs-value}{\vnbstmt{\begin{array}{@{}l@{}} \forall p.; \forall x \in \mathbb{Q}, v. \\ \quad (\operatorname{zz\text{-}prime}(p) \wedge \operatorname{padic\text{-}rel}(p, x, v)) \Rightarrow (\operatorname{padic\text{-}abs}(p)(x) = v) \end{array}}}
\vnbprop{pz-rel-intro}{\vnbstmt{\begin{array}{@{}l@{}} \forall p.; \forall x \in \mathbb{Q}, a, b \in \mathbb{Z}, i, j \in \mathbb{N}, v \in \mathbb{R}. \\ \quad \neg (x = 0) \wedge \\ \quad \neg (b = 0) \wedge \\ \quad (x \cdot b) = a \wedge \\ \quad \operatorname{zz\text{-}ord\text{-}is}(p, a, i) \wedge \\ \quad \operatorname{zz\text{-}ord\text{-}is}(p, b, j) \wedge \\ \quad (v \cdot \operatorname{power}(2, i)) = \operatorname{power}(2, j) \Rightarrow \\ \quad \quad \operatorname{padic\text{-}rel}(p, x, v) \end{array}}}
\vnbprop{pz-two-mono}{\vnbstmt{\forall a, b \in \mathbb{N}.\; ((a \leq b) \Rightarrow (\operatorname{power}(2, a) \leq \operatorname{power}(2, b)))}}
\vnbprop{pz-rel-pos}{\vnbstmt{\begin{array}{@{}l@{}} \forall p.; \forall x \in \mathbb{Q}, v. \\ \quad (\operatorname{zz\text{-}prime}(p) \wedge \operatorname{padic\text{-}rel}(p, x, v)) \Rightarrow ((0 \leq v) \wedge (\neg (x = 0) \Rightarrow (0 < v))) \end{array}}}
\vnbprop{pz-rel-zero}{\vnbstmt{\forall p.\; \operatorname{padic\text{-}rel}(p, 0, 0)}}
\vnbprop{pz-rel-rep}{\vnbstmt{\begin{array}{@{}l@{}} \forall p.; \forall z \in \mathbb{Q}, u.; \forall n, d \in \mathbb{Z}, s, t \in \mathbb{N}. \\ \quad \operatorname{zz\text{-}prime}(p) \wedge \\ \quad \operatorname{padic\text{-}rel}(p, z, u) \wedge \\ \quad \neg (z = 0) \wedge \\ \quad \neg (d = 0) \wedge \\ \quad (z \cdot d) = n \wedge \\ \quad \operatorname{zz\text{-}ord\text{-}is}(p, n, s) \wedge \\ \quad \operatorname{zz\text{-}ord\text{-}is}(p, d, t) \Rightarrow \\ \quad \quad (u \cdot \operatorname{power}(2, s)) = \operatorname{power}(2, t) \end{array}}}
\vnbprop{pz-ultra-num}{\vnbstmt{\begin{array}{@{}l@{}} \forall i, j, l, s \in \mathbb{N}, u, v \in \mathbb{R}. \\ \quad ((0 \leq v) \wedge (((i + l) \leq s) \wedge (((u \cdot \operatorname{power}(2, s)) = \operatorname{power}(2, (j + l))) \wedge ((v \cdot \operatorname{power}(2, i)) = \operatorname{power}(2, j))))) \Rightarrow  \\ \quad \quad u \leq v \end{array}}}
\vnbprop{pz-rel-mul}{\vnbstmt{\begin{array}{@{}l@{}} \forall p.; \forall x, y \in \mathbb{Q}, v, w. \\ \quad (\operatorname{zz\text{-}prime}(p) \wedge (\operatorname{padic\text{-}rel}(p, x, v) \wedge \operatorname{padic\text{-}rel}(p, y, w))) \Rightarrow  \\ \quad \quad \operatorname{padic\text{-}rel}(p, (x \cdot y), (v \cdot w)) \end{array}}}
\vnbprop{pz-rel-ultra}{\vnbstmt{\begin{array}{@{}l@{}} \forall p.; \forall x, y \in \mathbb{Q}, v, w, u. \\ \quad (\operatorname{zz\text{-}prime}(p) \wedge (\operatorname{padic\text{-}rel}(p, x, v) \wedge (\operatorname{padic\text{-}rel}(p, y, w) \wedge \operatorname{padic\text{-}rel}(p, (x + y), u)))) \Rightarrow  \\ \quad \quad u \leq \operatorname{max}(v, w) \end{array}}}
\vnbprop{pz-abs-rel}{\vnbstmt{\forall p.; \forall x \in \mathbb{Q}.\; (\operatorname{zz\text{-}prime}(p) \Rightarrow \operatorname{padic\text{-}rel}(p, x, \operatorname{padic\text{-}abs}(p)(x)))}}
\vnbprop{pz-abs-nonneg}{\vnbstmt{\forall p.; \forall x \in \mathbb{Q}.\; (\operatorname{zz\text{-}prime}(p) \Rightarrow (0 \leq \operatorname{padic\text{-}abs}(p)(x)))}}
\vnbprop{pz-abs-zero}{\vnbstmt{\forall p.; \forall x \in \mathbb{Q}.\; (\operatorname{zz\text{-}prime}(p) \Rightarrow ((\operatorname{padic\text{-}abs}(p)(x) = 0) \Leftrightarrow (x = 0)))}}
\vnbprop{pz-abs-mul}{\vnbstmt{\begin{array}{@{}l@{}} \forall p.; \forall x, y \in \mathbb{Q}. \\ \quad \operatorname{zz\text{-}prime}(p) \Rightarrow  \\ \quad \quad \operatorname{padic\text{-}abs}(p)((x \cdot y)) = (\operatorname{padic\text{-}abs}(p)(x) \cdot \operatorname{padic\text{-}abs}(p)(y)) \end{array}}}
\vnbprop{pz-abs-ultra}{\vnbstmt{\begin{array}{@{}l@{}} \forall p.; \forall x, y \in \mathbb{Q}. \\ \quad \operatorname{zz\text{-}prime}(p) \Rightarrow  \\ \quad \quad \operatorname{padic\text{-}abs}(p)((x + y)) \leq \operatorname{max}(\operatorname{padic\text{-}abs}(p)(x), \operatorname{padic\text{-}abs}(p)(y)) \end{array}}}
\vnbprop{pz-abs-triangle}{\vnbstmt{\begin{array}{@{}l@{}} \forall p.; \forall x, y \in \mathbb{Q}. \\ \quad \operatorname{zz\text{-}prime}(p) \Rightarrow  \\ \quad \quad \operatorname{padic\text{-}abs}(p)((x + y)) \leq (\operatorname{padic\text{-}abs}(p)(x) + \operatorname{padic\text{-}abs}(p)(y)) \end{array}}}
\vnbprop{pq-carr}{\vnbstmt{\forall p.\; (\operatorname{zz\text{-}prime}(p) \Rightarrow (\operatorname{carr}(\operatorname{qq\text{-}padic}(p)) = \mathbb{Q}))}}
\vnbprop{pq-add}{\vnbstmt{\begin{array}{@{}l@{}} \forall p. \\ \quad \operatorname{zz\text{-}prime}(p) \Rightarrow  \\ \quad \quad \operatorname{add}(\operatorname{qq\text{-}padic}(p)) = (\lambda {[}\operatorname{x\_}, \operatorname{y\_}] \in (\mathbb{Q} \times \mathbb{Q}).\; (\operatorname{x\_} + \operatorname{y\_})) \end{array}}}
\vnbprop{pq-mul}{\vnbstmt{\begin{array}{@{}l@{}} \forall p. \\ \quad \operatorname{zz\text{-}prime}(p) \Rightarrow  \\ \quad \quad \operatorname{mul}(\operatorname{qq\text{-}padic}(p)) = (\lambda {[}\operatorname{x\_}, \operatorname{y\_}] \in (\mathbb{Q} \times \mathbb{Q}).\; (\operatorname{x\_} \cdot \operatorname{y\_})) \end{array}}}
\vnbprop{pq-neg}{\vnbstmt{\forall p.\; (\operatorname{zz\text{-}prime}(p) \Rightarrow (\operatorname{neg}(\operatorname{qq\text{-}padic}(p)) = (\lambda x \in \mathbb{Q}.\; -x)))}}
\vnbprop{pq-zero}{\vnbstmt{\forall p.\; (\operatorname{zz\text{-}prime}(p) \Rightarrow (\operatorname{zero}(\operatorname{qq\text{-}padic}(p)) = 0))}}
\vnbprop{pq-one}{\vnbstmt{\forall p.\; (\operatorname{zz\text{-}prime}(p) \Rightarrow (\operatorname{one}(\operatorname{qq\text{-}padic}(p)) = 1))}}
\vnbprop{pq-fnrm}{\vnbstmt{\forall p.\; (\operatorname{zz\text{-}prime}(p) \Rightarrow (\operatorname{fnrm}(\operatorname{qq\text{-}padic}(p)) = \operatorname{padic\text{-}abs}(p)))}}
\vnbprop{qq-padic-is-normed-field}{\vnbstmt{\forall p.\; (\operatorname{zz\text{-}prime}(p) \Rightarrow \operatorname{is\text{-}normed\text{-}field}(\operatorname{qq\text{-}padic}(p)))}}
\vnbprop{qq-padic-is-nv-field}{\vnbstmt{\forall p.\; (\operatorname{zz\text{-}prime}(p) \Rightarrow \operatorname{is\text{-}nv\text{-}field}(\operatorname{qq\text{-}padic}(p)))}}
\vnbprop{pc-bezout-prime}{\vnbstmt{\begin{array}{@{}l@{}} \forall p.; \forall b \in \mathbb{Z}. \\ \quad (\operatorname{zz\text{-}prime}(p) \wedge \neg \operatorname{zz\text{-}divides}(p, b)) \Rightarrow  \\ \quad \quad \exists m, n \in \mathbb{Z}.\; (((m \cdot p) + (n \cdot b)) = 1) \end{array}}}
\vnbprop{pc-bezout-power-ind}{\vnbstmt{\begin{array}{@{}l@{}} \forall k \in \mathbb{N}, p, b \in \mathbb{Z}. \\ \quad \exists m, n \in \mathbb{Z}.\; (((m \cdot p) + (n \cdot b)) = 1) \Rightarrow  \\ \quad \quad \exists m, n \in \mathbb{Z}.\; (((m \cdot \operatorname{power}(p, k)) + (n \cdot b)) = 1) \end{array}}}
\vnbprop{pc-bezout-power}{\vnbstmt{\begin{array}{@{}l@{}} \forall p.; \forall b \in \mathbb{Z}, k \in \mathbb{N}. \\ \quad (\operatorname{zz\text{-}prime}(p) \wedge \neg \operatorname{zz\text{-}divides}(p, b)) \Rightarrow  \\ \quad \quad \exists m, n \in \mathbb{Z}.\; (((m \cdot \operatorname{power}(p, k)) + (n \cdot b)) = 1) \end{array}}}
\vnbprop{pc-abs-one}{\vnbstmt{\forall p.\; (\operatorname{zz\text{-}prime}(p) \Rightarrow (\operatorname{padic\text{-}abs}(p)(1) = 1))}}
\vnbprop{pc-abs-int}{\vnbstmt{\forall p.; \forall n \in \mathbb{Z}.\; (\operatorname{zz\text{-}prime}(p) \Rightarrow (\operatorname{padic\text{-}abs}(p)(n) \leq 1))}}
\vnbprop{pc-abs-power-ind}{\vnbstmt{\begin{array}{@{}l@{}} \forall k \in \mathbb{N}, p. \\ \quad \operatorname{zz\text{-}prime}(p) \Rightarrow (\operatorname{padic\text{-}abs}(p)(\operatorname{power}(p, k)) = \operatorname{power}(\operatorname{padic\text{-}abs}(p)(p), k)) \end{array}}}
\vnbprop{pc-abs-power}{\vnbstmt{\begin{array}{@{}l@{}} \forall p.; \forall k \in \mathbb{N}. \\ \quad \operatorname{zz\text{-}prime}(p) \Rightarrow (\operatorname{padic\text{-}abs}(p)(\operatorname{power}(p, k)) = \operatorname{power}(\operatorname{padic\text{-}abs}(p)(p), k)) \end{array}}}
\vnbprop{pc-two-mono-inv}{\vnbstmt{\forall i, j \in \mathbb{N}.\; ((\operatorname{power}(2, j) \leq \operatorname{power}(2, i)) \Rightarrow (j \leq i))}}
\vnbprop{pc-prime-not-dvd-one}{\vnbstmt{\forall p.\; (\operatorname{zz\text{-}prime}(p) \Rightarrow \neg \operatorname{zz\text{-}divides}(p, 1))}}
\vnbprop{pc-int-denominator}{\vnbstmt{\begin{array}{@{}l@{}} \forall p.; \forall x \in \mathbb{Q}. \\ \quad (\operatorname{zz\text{-}prime}(p) \wedge (\operatorname{padic\text{-}abs}(p)(x) \leq 1)) \Rightarrow  \\ \quad \quad \exists b \in \mathbb{Z}.\; (\neg \operatorname{zz\text{-}divides}(p, b) \wedge ((x \cdot b) \in \mathbb{Z})) \end{array}}}
\vnbprop{pc-approx-identity}{\vnbstmt{\begin{array}{@{}l@{}} \forall x, m, y, n, b \in \mathbb{R}. \\ \quad (((m \cdot y) + (n \cdot b)) = 1) \Rightarrow ((x + -((x \cdot b) \cdot n)) = ((x \cdot m) \cdot y)) \end{array}}}
\vnbprop{pc-nv-integers-eq}{\vnbstmt{\begin{array}{@{}l@{}} \forall p. \\ \quad \operatorname{zz\text{-}prime}(p) \Rightarrow (\operatorname{nv\text{-}integers}(\operatorname{qq\text{-}padic}(p)) = \{x \in \mathbb{Q} : (\operatorname{padic\text{-}abs}(p)(x) \leq 1)\}) \end{array}}}
\vnbprop{pc-approx-raw}{\vnbstmt{\begin{array}{@{}l@{}} \forall p.; \forall x \in \mathbb{Q}, k \in \mathbb{N}. \\ \quad (\operatorname{zz\text{-}prime}(p) \wedge (\operatorname{padic\text{-}abs}(p)(x) \leq 1)) \Rightarrow  \\ \quad \quad \exists h \in \mathbb{Z}.\; (\operatorname{padic\text{-}abs}(p)((x + -h)) \leq \operatorname{padic\text{-}abs}(p)(\operatorname{power}(p, k))) \end{array}}}
\vnbprop{pc-abs-int-mul}{\vnbstmt{\begin{array}{@{}l@{}} \forall p.; \forall n \in \mathbb{Z}, y \in \mathbb{Q}. \\ \quad \operatorname{zz\text{-}prime}(p) \Rightarrow (\operatorname{padic\text{-}abs}(p)((n \cdot y)) \leq \operatorname{padic\text{-}abs}(p)(y)) \end{array}}}
\vnbprop{ccr-1-19}{\vnbstmt{\begin{array}{@{}l@{}} \forall p.; \forall x \in \operatorname{nv\text{-}integers}(\operatorname{qq\text{-}padic}(p)), k \in \mathbb{N}. \\ \quad \operatorname{zz\text{-}prime}(p) \Rightarrow  \\ \quad \quad \exists h \in \mathbb{Z}. \\ \quad \quad \quad (0 \leq h) \wedge ((h < \operatorname{power}(p, k)) \wedge (\operatorname{padic\text{-}abs}(p)((x - h)) \leq \operatorname{power}(\operatorname{padic\text{-}abs}(p)(p), k))) \end{array}}}
\vnbprop{pc-ord-power}{\vnbstmt{\forall p.; \forall k \in \mathbb{N}.\; (\operatorname{zz\text{-}prime}(p) \Rightarrow \operatorname{zz\text{-}ord\text{-}is}(p, \operatorname{power}(p, k), k))}}
\vnbprop{pc-abs-power-two}{\vnbstmt{\forall p.; \forall k \in \mathbb{N}.\; (\operatorname{zz\text{-}prime}(p) \Rightarrow ((\operatorname{padic\text{-}abs}(p)(\operatorname{power}(p, k)) \cdot \operatorname{power}(2, k)) = 1))}}
\vnbprop{pc-power-two-beats}{\vnbstmt{\forall r.\; (\operatorname{pos\text{-}rr}(r) \Rightarrow \exists k \in \mathbb{N}.\; (1 < (r \cdot \operatorname{power}(2, k))))}}
\vnbprop{pc-z-subset-set}{\vnbstmt{\begin{array}{@{}l@{}} \forall p. \\ \quad \operatorname{zz\text{-}prime}(p) \Rightarrow  \\ \quad \quad (\operatorname{nv\text{-}integers}(\operatorname{qq\text{-}padic}(p)) \subseteq \mathbb{Q}) \wedge (\operatorname{nv\text{-}integers}(\operatorname{qq\text{-}padic}(p)) \in \mathbf{Set}) \end{array}}}
\vnbprop{pc-dist}{\vnbstmt{\begin{array}{@{}l@{}} \forall p.; \forall x, y \in \mathbb{Q}. \\ \quad \operatorname{zz\text{-}prime}(p) \Rightarrow  \\ \quad \quad \operatorname{dist}(\operatorname{nf\text{-}metric\text{-}space}(\operatorname{qq\text{-}padic}(p)))(x, y) = \operatorname{padic\text{-}abs}(p)((x + -y)) \end{array}}}
\vnbprop{ccr-1-21}{\vnbstmt{\begin{array}{@{}l@{}} \forall p. \\ \quad \operatorname{zz\text{-}prime}(p) \Rightarrow  \\ \quad \quad \operatorname{totally\text{-}bounded}(\operatorname{subspace\text{-}ms}(\operatorname{nf\text{-}metric\text{-}space}(\operatorname{qq\text{-}padic}(p)), \operatorname{nv\text{-}integers}(\operatorname{qq\text{-}padic}(p)))) \end{array}}}
\vnbprop{qn-carr}{\vnbstmt{\forall f \in (\mathbb{Q} \to \mathbb{R}).\; (\operatorname{carr}(\operatorname{qq\text{-}nv}(f)) = \mathbb{Q})}}
\vnbprop{qn-add}{\vnbstmt{\begin{array}{@{}l@{}} \forall f \in (\mathbb{Q} \to \mathbb{R}). \\ \quad \operatorname{add}(\operatorname{qq\text{-}nv}(f)) = (\lambda {[}\operatorname{x\_}, \operatorname{y\_}] \in (\mathbb{Q} \times \mathbb{Q}).\; (\operatorname{x\_} + \operatorname{y\_})) \end{array}}}
\vnbprop{qn-mul}{\vnbstmt{\begin{array}{@{}l@{}} \forall f \in (\mathbb{Q} \to \mathbb{R}). \\ \quad \operatorname{mul}(\operatorname{qq\text{-}nv}(f)) = (\lambda {[}\operatorname{x\_}, \operatorname{y\_}] \in (\mathbb{Q} \times \mathbb{Q}).\; (\operatorname{x\_} \cdot \operatorname{y\_})) \end{array}}}
\vnbprop{qn-neg}{\vnbstmt{\forall f \in (\mathbb{Q} \to \mathbb{R}).\; (\operatorname{neg}(\operatorname{qq\text{-}nv}(f)) = (\lambda x \in \mathbb{Q}.\; -x))}}
\vnbprop{qn-zero}{\vnbstmt{\forall f \in (\mathbb{Q} \to \mathbb{R}).\; (\operatorname{zero}(\operatorname{qq\text{-}nv}(f)) = 0)}}
\vnbprop{qn-fnrm}{\vnbstmt{\forall f \in (\mathbb{Q} \to \mathbb{R}).\; (\operatorname{fnrm}(\operatorname{qq\text{-}nv}(f)) = f)}}
\vnbprop{qn-in-rr}{\vnbstmt{\forall f \in (\mathbb{Q} \to \mathbb{R}), x \in \mathbb{Q}.\; (\operatorname{is\text{-}nv\text{-}field}(\operatorname{qq\text{-}nv}(f)) \Rightarrow (f(x) \in \mathbb{R}))}}
\vnbprop{qn-nonneg}{\vnbstmt{\forall f \in (\mathbb{Q} \to \mathbb{R}), x \in \mathbb{Q}.\; (\operatorname{is\text{-}nv\text{-}field}(\operatorname{qq\text{-}nv}(f)) \Rightarrow (0 \leq f(x)))}}
\vnbprop{qn-zero-iff}{\vnbstmt{\begin{array}{@{}l@{}} \forall f \in (\mathbb{Q} \to \mathbb{R}), x \in \mathbb{Q}. \\ \quad \operatorname{is\text{-}nv\text{-}field}(\operatorname{qq\text{-}nv}(f)) \Rightarrow ((f(x) = 0) \Leftrightarrow (x = 0)) \end{array}}}
\vnbprop{qn-mul-law}{\vnbstmt{\begin{array}{@{}l@{}} \forall f \in (\mathbb{Q} \to \mathbb{R}), x, y \in \mathbb{Q}. \\ \quad \operatorname{is\text{-}nv\text{-}field}(\operatorname{qq\text{-}nv}(f)) \Rightarrow (f((x \cdot y)) = (f(x) \cdot f(y))) \end{array}}}
\vnbprop{qn-neg-law}{\vnbstmt{\forall f \in (\mathbb{Q} \to \mathbb{R}), x \in \mathbb{Q}.\; (\operatorname{is\text{-}nv\text{-}field}(\operatorname{qq\text{-}nv}(f)) \Rightarrow (f(-x) = f(x)))}}
\vnbprop{qn-ultra}{\vnbstmt{\begin{array}{@{}l@{}} \forall f \in (\mathbb{Q} \to \mathbb{R}), x, y \in \mathbb{Q}. \\ \quad \operatorname{is\text{-}nv\text{-}field}(\operatorname{qq\text{-}nv}(f)) \Rightarrow (f((x + y)) \leq \operatorname{max}(f(x), f(y))) \end{array}}}
\vnbprop{qo-one}{\vnbstmt{\forall f \in (\mathbb{Q} \to \mathbb{R}).\; (\operatorname{is\text{-}nv\text{-}field}(\operatorname{qq\text{-}nv}(f)) \Rightarrow (f(1) = 1))}}
\vnbprop{qo-le-one-nn}{\vnbstmt{\forall n \in \mathbb{N}, f \in (\mathbb{Q} \to \mathbb{R}).\; (\operatorname{is\text{-}nv\text{-}field}(\operatorname{qq\text{-}nv}(f)) \Rightarrow (f(n) \leq 1))}}
\vnbprop{qo-le-one}{\vnbstmt{\forall f \in (\mathbb{Q} \to \mathbb{R}), m \in \mathbb{Z}.\; (\operatorname{is\text{-}nv\text{-}field}(\operatorname{qq\text{-}nv}(f)) \Rightarrow (f(m) \leq 1))}}
\vnbprop{qo-all-units}{\vnbstmt{\begin{array}{@{}l@{}} \forall f \in (\mathbb{Q} \to \mathbb{R}), m \in \mathbb{Z}. \\ \quad (\operatorname{is\text{-}nv\text{-}field}(\operatorname{qq\text{-}nv}(f)) \wedge (\neg \exists n \in \mathbb{N}.\; ((0 < n) \wedge (f(n) < 1)) \wedge \neg (m = 0))) \Rightarrow  \\ \quad \quad f(m) = 1 \end{array}}}
\vnbprop{qo-small-exists}{\vnbstmt{\begin{array}{@{}l@{}} \forall f \in (\mathbb{Q} \to \mathbb{R}). \\ \quad (\operatorname{is\text{-}nv\text{-}field}(\operatorname{qq\text{-}nv}(f)) \wedge \operatorname{is\text{-}nontrivial\text{-}nv}(\operatorname{qq\text{-}nv}(f))) \Rightarrow  \\ \quad \quad \exists n \in \mathbb{N}.\; ((0 < n) \wedge (f(n) < 1)) \end{array}}}
\vnbprop{qo-prime}{\vnbstmt{\begin{array}{@{}l@{}} \forall f \in (\mathbb{Q} \to \mathbb{R}), p \in \mathbb{N}. \\ \quad \operatorname{is\text{-}nv\text{-}field}(\operatorname{qq\text{-}nv}(f)) \wedge \\ \quad 1 < p \wedge \\ \quad f(p) < 1 \wedge \\ \quad \forall e \in \mathbb{N}.\; (((0 < e) \wedge (f(e) < 1)) \Rightarrow (p \leq e)) \Rightarrow \\ \quad \quad \operatorname{zz\text{-}prime}(p) \end{array}}}
\vnbprop{qo-units}{\vnbstmt{\begin{array}{@{}l@{}} \forall f \in (\mathbb{Q} \to \mathbb{R}), p.; \forall m \in \mathbb{Z}. \\ \quad (\operatorname{is\text{-}nv\text{-}field}(\operatorname{qq\text{-}nv}(f)) \wedge (\operatorname{zz\text{-}prime}(p) \wedge ((f(p) < 1) \wedge \neg \operatorname{zz\text{-}divides}(p, m)))) \Rightarrow  \\ \quad \quad f(m) = 1 \end{array}}}
\vnbprop{qo-power}{\vnbstmt{\begin{array}{@{}l@{}} \forall k \in \mathbb{N}, f \in (\mathbb{Q} \to \mathbb{R}), p \in \mathbb{Z}. \\ \quad \operatorname{is\text{-}nv\text{-}field}(\operatorname{qq\text{-}nv}(f)) \Rightarrow (f(\operatorname{power}(p, k)) = \operatorname{power}(f(p), k)) \end{array}}}
\vnbprop{qo-int-value}{\vnbstmt{\begin{array}{@{}l@{}} \forall f \in (\mathbb{Q} \to \mathbb{R}), p.; \forall n \in \mathbb{Z}, i \in \mathbb{N}. \\ \quad (\operatorname{is\text{-}nv\text{-}field}(\operatorname{qq\text{-}nv}(f)) \wedge (\operatorname{zz\text{-}prime}(p) \wedge ((f(p) < 1) \wedge \operatorname{zz\text{-}ord\text{-}is}(p, n, i)))) \Rightarrow  \\ \quad \quad f(n) = \operatorname{power}(f(p), i) \end{array}}}
\vnbprop{ccr-1-24}{\vnbstmt{\begin{array}{@{}l@{}} \forall f \in (\mathbb{Q} \to \mathbb{R}). \\ \quad (\operatorname{is\text{-}nv\text{-}field}(\operatorname{qq\text{-}nv}(f)) \wedge \operatorname{is\text{-}nontrivial\text{-}nv}(\operatorname{qq\text{-}nv}(f))) \Rightarrow  \\ \quad \quad \exists p. \\ \quad \quad \quad \operatorname{zz\text{-}prime}(p) \wedge \\ \quad \quad \quad \exists r \in \mathbb{R}. \\ \quad \quad \quad \quad 0 < r \wedge \\ \quad \quad \quad \quad r < 1 \wedge \\ \quad \quad \quad \quad \forall x \in \mathbb{Q}, a, b \in \mathbb{Z}, i, j \in \mathbb{N}. \\ \quad \quad \quad \quad \quad (\neg (x = 0) \wedge (\neg (b = 0) \wedge (((x \cdot b) = a) \wedge (\operatorname{zz\text{-}ord\text{-}is}(p, a, i) \wedge \operatorname{zz\text{-}ord\text{-}is}(p, b, j))))) \Rightarrow  \\ \quad \quad \quad \quad \quad \quad (f(x) \cdot \operatorname{power}(r, j)) = \operatorname{power}(r, i) \end{array}}}
\vnbprop{ccr-1-20}{\vnbstmt{\begin{array}{@{}l@{}} \forall p.; \forall x \in \operatorname{nv\text{-}integers}(\operatorname{qq\text{-}padic}(p)), s. \\ \quad (\operatorname{zz\text{-}prime}(p) \wedge \operatorname{pos\text{-}rr}(s)) \Rightarrow  \\ \quad \quad \exists h \in \mathbb{Z}.\; (\operatorname{dist}(\operatorname{nf\text{-}metric\text{-}space}(\operatorname{qq\text{-}padic}(p)))(h, x) < s) \end{array}}}
\vnbprop{pc-e-int}{\vnbstmt{\forall e \in \mathbb{Z}.\; (\exp (((2 \cdot \pi) \cdot +i) \cdot e) = 1)}}
\vnbprop{pc-e-add}{\vnbstmt{\begin{array}{@{}l@{}} \forall r, s \in \mathbb{R}. \\ \quad \exp (((2 \cdot \pi) \cdot +i) \cdot (r + s)) = (\exp (((2 \cdot \pi) \cdot +i) \cdot r) \cdot \exp (((2 \cdot \pi) \cdot +i) \cdot s)) \end{array}}}
\vnbprop{pc-e-one-int}{\vnbstmt{\forall r \in \mathbb{R}.\; ((\exp (((2 \cdot \pi) \cdot +i) \cdot r) = 1) \Rightarrow (r \in \mathbb{Z}))}}
\vnbprop{pc-e-unit}{\vnbstmt{\forall r \in \mathbb{R}.\; (\exp (((2 \cdot \pi) \cdot +i) \cdot r) \in \operatorname{unit\text{-}circle})}}
\vnbprop{pc-abs-neg}{\vnbstmt{\forall p.; \forall w \in \mathbb{Q}.\; (\operatorname{zz\text{-}prime}(p) \Rightarrow (\operatorname{padic\text{-}abs}(p)(-w) = \operatorname{padic\text{-}abs}(p)(w)))}}
\vnbprop{pc-abs-unit}{\vnbstmt{\begin{array}{@{}l@{}} \forall p.; \forall m \in \mathbb{Z}. \\ \quad (\operatorname{zz\text{-}prime}(p) \wedge \neg \operatorname{zz\text{-}divides}(p, m)) \Rightarrow (\operatorname{padic\text{-}abs}(p)(m) = 1) \end{array}}}
\vnbprop{pc-frac-int}{\vnbstmt{\begin{array}{@{}l@{}} \forall p.; \forall y \in \mathbb{Q}, m \in \mathbb{N}. \\ \quad (\operatorname{zz\text{-}prime}(p) \wedge ((\operatorname{padic\text{-}abs}(p)(y) \leq 1) \wedge ((y \cdot \operatorname{power}(p, m)) \in \mathbb{Z}))) \Rightarrow (y \in \mathbb{Z}) \end{array}}}
\vnbprop{pc-rep-exists}{\vnbstmt{\forall p.; \forall x \in \mathbb{Q}.\; (\operatorname{zz\text{-}prime}(p) \Rightarrow \exists r.\; \operatorname{padic\text{-}frac\text{-}rep}(p, x, r))}}
\vnbprop{pc-rep-add}{\vnbstmt{\begin{array}{@{}l@{}} \forall p.; \forall x, y \in \mathbb{Q}, r, s. \\ \quad (\operatorname{zz\text{-}prime}(p) \wedge (\operatorname{padic\text{-}frac\text{-}rep}(p, x, r) \wedge \operatorname{padic\text{-}frac\text{-}rep}(p, y, s))) \Rightarrow  \\ \quad \quad \operatorname{padic\text{-}frac\text{-}rep}(p, (x + y), (r + s)) \end{array}}}
\vnbprop{pc-rep-shift}{\vnbstmt{\begin{array}{@{}l@{}} \forall p.; \forall x, y \in \mathbb{Q}, r. \\ \quad (\operatorname{zz\text{-}prime}(p) \wedge (\operatorname{padic\text{-}frac\text{-}rep}(p, y, r) \wedge (\operatorname{padic\text{-}abs}(p)((x + -y)) \leq 1))) \Rightarrow  \\ \quad \quad \operatorname{padic\text{-}frac\text{-}rep}(p, x, r) \end{array}}}
\vnbprop{pc-rep-diff-int}{\vnbstmt{\begin{array}{@{}l@{}} \forall p.; \forall x \in \mathbb{Q}, r, s. \\ \quad (\operatorname{zz\text{-}prime}(p) \wedge (\operatorname{padic\text{-}frac\text{-}rep}(p, x, r) \wedge \operatorname{padic\text{-}frac\text{-}rep}(p, x, s))) \Rightarrow  \\ \quad \quad (r + -s) \in \mathbb{Z} \end{array}}}
\vnbprop{pc-e-rep-unique}{\vnbstmt{\begin{array}{@{}l@{}} \forall p.; \forall x \in \mathbb{Q}, r, s. \\ \quad (\operatorname{zz\text{-}prime}(p) \wedge (\operatorname{padic\text{-}frac\text{-}rep}(p, x, r) \wedge \operatorname{padic\text{-}frac\text{-}rep}(p, x, s))) \Rightarrow  \\ \quad \quad \exp (((2 \cdot \pi) \cdot +i) \cdot r) = \exp (((2 \cdot \pi) \cdot +i) \cdot s) \end{array}}}
\vnbprop{pc-char-iota}{\vnbstmt{\begin{array}{@{}l@{}} \forall p.; \forall x \in \mathbb{Q}. \\ \quad \operatorname{zz\text{-}prime}(p) \Rightarrow  \\ \quad \quad \exists r. \\ \quad \quad \quad \operatorname{padic\text{-}frac\text{-}rep}(p, x, r) \wedge \\ \quad \quad \quad (\iota\, w.\;((w \in \mathbb{C}) \wedge \exists \mathit{cpr} \in \mathbb{Q}.\; (\operatorname{padic\text{-}frac\text{-}rep}(p, x, \mathit{cpr}) \wedge (w = \exp (((2 \cdot \pi) \cdot +i) \cdot \mathit{cpr}))))) \\ \quad \quad \quad =\; \exp (((2 \cdot \pi) \cdot +i) \cdot r) \end{array}}}
\vnbprop{pc-char-value}{\vnbstmt{\begin{array}{@{}l@{}} \forall p.; \forall x \in \mathbb{Q}, r. \\ \quad (\operatorname{zz\text{-}prime}(p) \wedge \operatorname{padic\text{-}frac\text{-}rep}(p, x, r)) \Rightarrow  \\ \quad \quad \operatorname{padic\text{-}char}(p)(x) = \exp (((2 \cdot \pi) \cdot +i) \cdot r) \end{array}}}
\vnbprop{pc-char-add}{\vnbstmt{\begin{array}{@{}l@{}} \forall p.; \forall x, y \in \mathbb{Q}. \\ \quad \operatorname{zz\text{-}prime}(p) \Rightarrow  \\ \quad \quad \operatorname{padic\text{-}char}(p)((x + y)) = (\operatorname{padic\text{-}char}(p)(x) \cdot \operatorname{padic\text{-}char}(p)(y)) \end{array}}}
\vnbprop{pc-char-local}{\vnbstmt{\begin{array}{@{}l@{}} \forall p.; \forall x, y \in \mathbb{Q}. \\ \quad (\operatorname{zz\text{-}prime}(p) \wedge (\operatorname{padic\text{-}abs}(p)((x + -y)) \leq 1)) \Rightarrow  \\ \quad \quad \operatorname{padic\text{-}char}(p)(x) = \operatorname{padic\text{-}char}(p)(y) \end{array}}}
\vnbprop{pc-char-kernel}{\vnbstmt{\begin{array}{@{}l@{}} \forall p.; \forall x \in \mathbb{Q}. \\ \quad \operatorname{zz\text{-}prime}(p) \Rightarrow ((\operatorname{padic\text{-}char}(p)(x) = 1) \Leftrightarrow (\operatorname{padic\text{-}abs}(p)(x) \leq 1)) \end{array}}}
\vnbprop{pc-char-type}{\vnbstmt{\forall p.\; (\operatorname{zz\text{-}prime}(p) \Rightarrow (\operatorname{padic\text{-}char}(p) \in (\mathbb{Q} \to \operatorname{unit\text{-}circle})))}}
\vnbprop{pc-char-type-cc}{\vnbstmt{\forall p.\; (\operatorname{zz\text{-}prime}(p) \Rightarrow (\operatorname{padic\text{-}char}(p) \in (\mathbb{Q} \to \mathbb{C})))}}
\vnbprop{ccr-5-20-character}{\vnbstmt{\forall p.\; (\operatorname{zz\text{-}prime}(p) \Rightarrow \operatorname{is\text{-}add\text{-}character}(\operatorname{qq\text{-}padic}(p), \operatorname{padic\text{-}char}(p)))}}
\vnbprop{ccr-5-21}{\vnbstmt{\begin{array}{@{}l@{}} \forall p.; \forall x \in \mathbb{Q}. \\ \quad \operatorname{zz\text{-}prime}(p) \Rightarrow ((\operatorname{padic\text{-}char}(p)(x) = 1) \Leftrightarrow (x \in \operatorname{nv\text{-}integers}(\operatorname{qq\text{-}padic}(p)))) \end{array}}}

\vnbbookfile{theorem-library/qp-laws}{theorem-library/qp-laws}
\noindent{\small qp-laws.scm -- WAVE 2 (2026-10-27..30), wired 2026-10-30 from scratchpad/ccrqp/qp-laws.scm. Library: complex-analysis.}\par
\vnbprop{qp-isosceles}{\vnbstmt{\begin{array}{@{}l@{}} \forall p.; \forall x, y \in \mathbb{Q}. \\ \quad (\operatorname{zz\text{-}prime}(p) \wedge (\operatorname{padic\text{-}abs}(p)((x + -y)) < \operatorname{padic\text{-}abs}(p)(x))) \Rightarrow  \\ \quad \quad \operatorname{padic\text{-}abs}(p)(y) = \operatorname{padic\text{-}abs}(p)(x) \end{array}}}
\vnbprop{qp-init-bound-ind}{\vnbstmt{\begin{array}{@{}l@{}} \forall k \in \mathbb{N}, p.; \forall f \in (\mathbb{N} \to \mathbb{Q}). \\ \quad \operatorname{zz\text{-}prime}(p) \Rightarrow  \\ \quad \quad \exists b \in \mathbb{R}. \\ \quad \quad \quad (0 \leq b) \wedge \forall n \in \mathbb{N}.\; ((n < k) \Rightarrow (\operatorname{padic\text{-}abs}(p)(f(n)) \leq b)) \end{array}}}
\vnbprop{qp-cauchy-bounded}{\vnbstmt{\begin{array}{@{}l@{}} \forall p, f. \\ \quad (\operatorname{zz\text{-}prime}(p) \wedge \operatorname{is\text{-}cauchy\text{-}seq}(\operatorname{nf\text{-}metric\text{-}space}(\operatorname{qq\text{-}padic}(p)), f)) \Rightarrow  \\ \quad \quad \exists b \in \mathbb{R}.\; ((0 \leq b) \wedge \forall n \in \mathbb{N}.\; (\operatorname{padic\text{-}abs}(p)(f(n)) \leq b)) \end{array}}}
\vnbprop{qp-equiv-intro}{\vnbstmt{\begin{array}{@{}l@{}} \forall p.; \forall f, g \in (\mathbb{N} \to \mathbb{Q}). \\ \quad \operatorname{zz\text{-}prime}(p) \wedge \\ \quad \forall s. \\ \quad \quad \operatorname{pos\text{-}rr}(s) \Rightarrow  \\ \quad \quad \quad \exists b \in \mathbb{N}. \forall a \in \mathbb{N}. \\ \quad \quad \quad \quad (b \leq a) \Rightarrow (\operatorname{padic\text{-}abs}(p)((f(a) + -g(a))) \leq s) \Rightarrow \\ \quad \quad \operatorname{cseq\text{-}equiv}(\operatorname{nf\text{-}metric\text{-}space}(\operatorname{qq\text{-}padic}(p)), f, g) \end{array}}}
\vnbprop{qp-equiv-elim}{\vnbstmt{\begin{array}{@{}l@{}} \forall p.; \forall f, g \in (\mathbb{N} \to \mathbb{Q}), s. \\ \quad (\operatorname{zz\text{-}prime}(p) \wedge (\operatorname{cseq\text{-}equiv}(\operatorname{nf\text{-}metric\text{-}space}(\operatorname{qq\text{-}padic}(p)), f, g) \wedge \operatorname{pos\text{-}rr}(s))) \Rightarrow  \\ \quad \quad \exists b \in \mathbb{N}. \forall a \in \mathbb{N}. \\ \quad \quad \quad (b \leq a) \Rightarrow (\operatorname{padic\text{-}abs}(p)((f(a) + -g(a))) \leq s) \end{array}}}
\vnbprop{qp-cauchy-intro}{\vnbstmt{\begin{array}{@{}l@{}} \forall p.; \forall f \in (\mathbb{N} \to \mathbb{Q}). \\ \quad \operatorname{zz\text{-}prime}(p) \wedge \\ \quad \forall s. \\ \quad \quad \operatorname{pos\text{-}rr}(s) \Rightarrow  \\ \quad \quad \quad \exists b \in \mathbb{N}. \forall a, c \in \mathbb{N}. \\ \quad \quad \quad \quad ((b \leq a) \wedge (b \leq c)) \Rightarrow (\operatorname{padic\text{-}abs}(p)((f(a) + -f(c))) \leq s) \Rightarrow \\ \quad \quad \operatorname{is\text{-}cauchy\text{-}seq}(\operatorname{nf\text{-}metric\text{-}space}(\operatorname{qq\text{-}padic}(p)), f) \end{array}}}
\vnbprop{qp-cauchy-elim}{\vnbstmt{\begin{array}{@{}l@{}} \forall p, f, s. \\ \quad (\operatorname{zz\text{-}prime}(p) \wedge (\operatorname{is\text{-}cauchy\text{-}seq}(\operatorname{nf\text{-}metric\text{-}space}(\operatorname{qq\text{-}padic}(p)), f) \wedge \operatorname{pos\text{-}rr}(s))) \Rightarrow  \\ \quad \quad \exists b \in \mathbb{N}. \forall a, c \in \mathbb{N}. \\ \quad \quad \quad ((b \leq a) \wedge (b \leq c)) \Rightarrow (\operatorname{padic\text{-}abs}(p)((f(a) + -f(c))) \leq s) \end{array}}}
\vnbprop{qp-sum-cauchy}{\vnbstmt{\begin{array}{@{}l@{}} \forall p, f, g. \\ \quad \operatorname{zz\text{-}prime}(p) \wedge \\ \quad \operatorname{is\text{-}cauchy\text{-}seq}(\operatorname{nf\text{-}metric\text{-}space}(\operatorname{qq\text{-}padic}(p)), f) \wedge \\ \quad \operatorname{is\text{-}cauchy\text{-}seq}(\operatorname{nf\text{-}metric\text{-}space}(\operatorname{qq\text{-}padic}(p)), g) \Rightarrow \\ \quad \quad \operatorname{is\text{-}cauchy\text{-}seq}(\operatorname{nf\text{-}metric\text{-}space}(\operatorname{qq\text{-}padic}(p)), (\lambda n \in \mathbb{N}.\; (f(n) + g(n)))) \end{array}}}
\vnbprop{qp-neg-cauchy}{\vnbstmt{\begin{array}{@{}l@{}} \forall p, f. \\ \quad (\operatorname{zz\text{-}prime}(p) \wedge \operatorname{is\text{-}cauchy\text{-}seq}(\operatorname{nf\text{-}metric\text{-}space}(\operatorname{qq\text{-}padic}(p)), f)) \Rightarrow  \\ \quad \quad \operatorname{is\text{-}cauchy\text{-}seq}(\operatorname{nf\text{-}metric\text{-}space}(\operatorname{qq\text{-}padic}(p)), (\lambda n \in \mathbb{N}.\; -f(n))) \end{array}}}
\vnbprop{qp-const-cauchy}{\vnbstmt{\begin{array}{@{}l@{}} \forall p.; \forall c \in \mathbb{Q}. \\ \quad \operatorname{zz\text{-}prime}(p) \Rightarrow  \\ \quad \quad \operatorname{is\text{-}cauchy\text{-}seq}(\operatorname{nf\text{-}metric\text{-}space}(\operatorname{qq\text{-}padic}(p)), (\lambda n \in \mathbb{N}.\; c)) \end{array}}}
\vnbprop{qp-mul-cauchy}{\vnbstmt{\begin{array}{@{}l@{}} \forall p, f, g. \\ \quad \operatorname{zz\text{-}prime}(p) \wedge \\ \quad \operatorname{is\text{-}cauchy\text{-}seq}(\operatorname{nf\text{-}metric\text{-}space}(\operatorname{qq\text{-}padic}(p)), f) \wedge \\ \quad \operatorname{is\text{-}cauchy\text{-}seq}(\operatorname{nf\text{-}metric\text{-}space}(\operatorname{qq\text{-}padic}(p)), g) \Rightarrow \\ \quad \quad \operatorname{is\text{-}cauchy\text{-}seq}(\operatorname{nf\text{-}metric\text{-}space}(\operatorname{qq\text{-}padic}(p)), (\lambda n \in \mathbb{N}.\; (f(n) \cdot g(n)))) \end{array}}}
\vnbprop{qp-cs-iff}{\vnbstmt{\begin{array}{@{}l@{}} \forall p, h. \\ \quad \operatorname{zz\text{-}prime}(p) \Rightarrow  \\ \quad \quad (h \in \operatorname{qp\text{-}cs}(p)) \Leftrightarrow \operatorname{is\text{-}cauchy\text{-}seq}(\operatorname{nf\text{-}metric\text{-}space}(\operatorname{qq\text{-}padic}(p)), h) \end{array}}}
\vnbprop{qp-null-iff}{\vnbstmt{\begin{array}{@{}l@{}} \forall p.; \forall f \in (\mathbb{N} \to \mathbb{Q}). \\ \quad \operatorname{zz\text{-}prime}(p) \Rightarrow  \\ \quad \quad \operatorname{cseq\text{-}equiv}(\operatorname{nf\text{-}metric\text{-}space}(\operatorname{qq\text{-}padic}(p)), f, (\lambda n \in \mathbb{N}.\; 0)) \\ \quad \quad \Leftrightarrow\; \forall s. \\ \quad \quad \quad \operatorname{pos\text{-}rr}(s) \Rightarrow  \\ \quad \quad \quad \quad \exists b \in \mathbb{N}.\; \forall a \in \mathbb{N}.\; ((b \leq a) \Rightarrow (\operatorname{padic\text{-}abs}(p)(f(a)) \leq s)) \end{array}}}
\vnbprop{qp-null-member-iff}{\vnbstmt{\begin{array}{@{}l@{}} \forall p, h. \\ \quad \operatorname{zz\text{-}prime}(p) \Rightarrow  \\ \quad \quad h \in \operatorname{qp\text{-}null}(p) \\ \quad \quad \Leftrightarrow\; \operatorname{is\text{-}cauchy\text{-}seq}(\operatorname{nf\text{-}metric\text{-}space}(\operatorname{qq\text{-}padic}(p)), h) \wedge \\ \quad \quad \forall s. \\ \quad \quad \quad \operatorname{pos\text{-}rr}(s) \Rightarrow  \\ \quad \quad \quad \quad \exists b \in \mathbb{N}.\; \forall a \in \mathbb{N}.\; ((b \leq a) \Rightarrow (\operatorname{padic\text{-}abs}(p)(h(a)) \leq s)) \end{array}}}
\vnbprop{qp-null-add}{\vnbstmt{\begin{array}{@{}l@{}} \forall p.; \forall a, b \in \operatorname{qp\text{-}null}(p). \\ \quad \operatorname{zz\text{-}prime}(p) \Rightarrow ((\lambda n \in \mathbb{N}.\; (a(n) + b(n))) \in \operatorname{qp\text{-}null}(p)) \end{array}}}
\vnbprop{qp-null-neg}{\vnbstmt{\begin{array}{@{}l@{}} \forall p.; \forall a \in \operatorname{qp\text{-}null}(p). \\ \quad \operatorname{zz\text{-}prime}(p) \Rightarrow ((\lambda n \in \mathbb{N}.\; -a(n)) \in \operatorname{qp\text{-}null}(p)) \end{array}}}
\vnbprop{qp-null-zero}{\vnbstmt{\forall p.\; (\operatorname{zz\text{-}prime}(p) \Rightarrow ((\lambda n \in \mathbb{N}.\; 0) \in \operatorname{qp\text{-}null}(p)))}}
\vnbprop{qp-null-mul-left}{\vnbstmt{\begin{array}{@{}l@{}} \forall p.; \forall x \in \operatorname{qp\text{-}cs}(p), a \in \operatorname{qp\text{-}null}(p). \\ \quad \operatorname{zz\text{-}prime}(p) \Rightarrow ((\lambda n \in \mathbb{N}.\; (x(n) \cdot a(n))) \in \operatorname{qp\text{-}null}(p)) \end{array}}}
\vnbprop{qp-null-mul-right}{\vnbstmt{\begin{array}{@{}l@{}} \forall p.; \forall x \in \operatorname{qp\text{-}cs}(p), a \in \operatorname{qp\text{-}null}(p). \\ \quad \operatorname{zz\text{-}prime}(p) \Rightarrow ((\lambda n \in \mathbb{N}.\; (a(n) \cdot x(n))) \in \operatorname{qp\text{-}null}(p)) \end{array}}}
\vnbprop{qp-cs-set}{\vnbstmt{\forall p.\; (\operatorname{zz\text{-}prime}(p) \Rightarrow (\operatorname{qp\text{-}cs}(p) \in \mathbf{Set}))}}
\vnbprop{qp-sadd-at}{\vnbstmt{\begin{array}{@{}l@{}} \forall p.; \forall u, v \in \operatorname{qp\text{-}cs}(p), k \in \mathbb{N}. \\ \quad \operatorname{zz\text{-}prime}(p) \Rightarrow (\operatorname{qp\text{-}sadd}(p)(u, v)(k) = (u(k) + v(k))) \end{array}}}
\vnbprop{qp-smul-at}{\vnbstmt{\begin{array}{@{}l@{}} \forall p.; \forall u, v \in \operatorname{qp\text{-}cs}(p), k \in \mathbb{N}. \\ \quad \operatorname{zz\text{-}prime}(p) \Rightarrow (\operatorname{qp\text{-}smul}(p)(u, v)(k) = (u(k) \cdot v(k))) \end{array}}}
\vnbprop{qp-sneg-at}{\vnbstmt{\begin{array}{@{}l@{}} \forall p.; \forall u \in \operatorname{qp\text{-}cs}(p), k \in \mathbb{N}. \\ \quad \operatorname{zz\text{-}prime}(p) \Rightarrow (\operatorname{qp\text{-}sneg}(p)(u)(k) = -u(k)) \end{array}}}
\vnbprop{qp-sadd-type}{\vnbstmt{\begin{array}{@{}l@{}} \forall p.; \forall u, v \in \operatorname{qp\text{-}cs}(p). \\ \quad \operatorname{zz\text{-}prime}(p) \Rightarrow (\operatorname{qp\text{-}sadd}(p)(u, v) \in \operatorname{qp\text{-}cs}(p)) \end{array}}}
\vnbprop{qp-smul-type}{\vnbstmt{\begin{array}{@{}l@{}} \forall p.; \forall u, v \in \operatorname{qp\text{-}cs}(p). \\ \quad \operatorname{zz\text{-}prime}(p) \Rightarrow (\operatorname{qp\text{-}smul}(p)(u, v) \in \operatorname{qp\text{-}cs}(p)) \end{array}}}
\vnbprop{qp-sneg-type}{\vnbstmt{\forall p.; \forall u \in \operatorname{qp\text{-}cs}(p).\; (\operatorname{zz\text{-}prime}(p) \Rightarrow (\operatorname{qp\text{-}sneg}(p)(u) \in \operatorname{qp\text{-}cs}(p)))}}
\vnbprop{qp-sadd-fun}{\vnbstmt{\begin{array}{@{}l@{}} \forall p. \\ \quad \operatorname{zz\text{-}prime}(p) \Rightarrow (\operatorname{qp\text{-}sadd}(p) \in ((\operatorname{qp\text{-}cs}(p) \times \operatorname{qp\text{-}cs}(p)) \to \operatorname{qp\text{-}cs}(p))) \end{array}}}
\vnbprop{qp-smul-fun}{\vnbstmt{\begin{array}{@{}l@{}} \forall p. \\ \quad \operatorname{zz\text{-}prime}(p) \Rightarrow (\operatorname{qp\text{-}smul}(p) \in ((\operatorname{qp\text{-}cs}(p) \times \operatorname{qp\text{-}cs}(p)) \to \operatorname{qp\text{-}cs}(p))) \end{array}}}
\vnbprop{qp-sneg-fun}{\vnbstmt{\forall p.\; (\operatorname{zz\text{-}prime}(p) \Rightarrow (\operatorname{qp\text{-}sneg}(p) \in (\operatorname{qp\text{-}cs}(p) \to \operatorname{qp\text{-}cs}(p))))}}
\vnbprop{qp-zero-cs}{\vnbstmt{\forall p.\; (\operatorname{zz\text{-}prime}(p) \Rightarrow ((\lambda n \in \mathbb{N}.\; 0) \in \operatorname{qp\text{-}cs}(p)))}}
\vnbprop{qp-one-cs}{\vnbstmt{\forall p.\; (\operatorname{zz\text{-}prime}(p) \Rightarrow ((\lambda n \in \mathbb{N}.\; 1) \in \operatorname{qp\text{-}cs}(p)))}}
\vnbprop{qp-null-cs}{\vnbstmt{\forall p.; \forall z \in \operatorname{qp\text{-}null}(p).\; (\operatorname{zz\text{-}prime}(p) \Rightarrow (z \in \operatorname{qp\text{-}cs}(p)))}}
\vnbprop{qp-cs-fun}{\vnbstmt{\forall p.; \forall h \in \operatorname{qp\text{-}cs}(p).\; (\operatorname{zz\text{-}prime}(p) \Rightarrow (h \in (\mathbb{N} \to \mathbb{Q})))}}
\vnbprop{qp-sadd-value}{\vnbstmt{\begin{array}{@{}l@{}} \forall p.; \forall u, v \in \operatorname{qp\text{-}cs}(p). \\ \quad \operatorname{zz\text{-}prime}(p) \Rightarrow (\operatorname{qp\text{-}sadd}(p)(u, v) = (\lambda n \in \mathbb{N}.\; (u(n) + v(n)))) \end{array}}}
\vnbprop{qp-smul-value}{\vnbstmt{\begin{array}{@{}l@{}} \forall p.; \forall u, v \in \operatorname{qp\text{-}cs}(p). \\ \quad \operatorname{zz\text{-}prime}(p) \Rightarrow (\operatorname{qp\text{-}smul}(p)(u, v) = (\lambda n \in \mathbb{N}.\; (u(n) \cdot v(n)))) \end{array}}}
\vnbprop{qp-sneg-value}{\vnbstmt{\begin{array}{@{}l@{}} \forall p.; \forall u \in \operatorname{qp\text{-}cs}(p). \\ \quad \operatorname{zz\text{-}prime}(p) \Rightarrow (\operatorname{qp\text{-}sneg}(p)(u) = (\lambda n \in \mathbb{N}.\; -u(n))) \end{array}}}
\vnbprop{qp-null-power}{\vnbstmt{\forall p.\; (\operatorname{zz\text{-}prime}(p) \Rightarrow (\operatorname{qp\text{-}null}(p) \in \operatorname{power}(\operatorname{qp\text{-}cs}(p))))}}
\vnbprop{qp-csr-carr}{\vnbstmt{\forall p.\; (\operatorname{zz\text{-}prime}(p) \Rightarrow (\operatorname{carr}(\operatorname{qp\text{-}csr}(p)) \simeq \operatorname{qp\text{-}cs}(p)))}}
\vnbprop{qp-csr-add}{\vnbstmt{\forall p.\; (\operatorname{zz\text{-}prime}(p) \Rightarrow (\operatorname{add}(\operatorname{qp\text{-}csr}(p)) \simeq \operatorname{qp\text{-}sadd}(p)))}}
\vnbprop{qp-csr-mul}{\vnbstmt{\forall p.\; (\operatorname{zz\text{-}prime}(p) \Rightarrow (\operatorname{mul}(\operatorname{qp\text{-}csr}(p)) \simeq \operatorname{qp\text{-}smul}(p)))}}
\vnbprop{qp-csr-neg}{\vnbstmt{\forall p.\; (\operatorname{zz\text{-}prime}(p) \Rightarrow (\operatorname{neg}(\operatorname{qp\text{-}csr}(p)) \simeq \operatorname{qp\text{-}sneg}(p)))}}
\vnbprop{qp-csr-zero}{\vnbstmt{\forall p.\; (\operatorname{zz\text{-}prime}(p) \Rightarrow (\operatorname{zero}(\operatorname{qp\text{-}csr}(p)) \simeq (\lambda n \in \mathbb{N}.\; 0)))}}
\vnbprop{qp-csr-one}{\vnbstmt{\forall p.\; (\operatorname{zz\text{-}prime}(p) \Rightarrow (\operatorname{one}(\operatorname{qp\text{-}csr}(p)) \simeq (\lambda n \in \mathbb{N}.\; 1)))}}
\vnbprop{qp-csr-idl}{\vnbstmt{\forall p.\; (\operatorname{zz\text{-}prime}(p) \Rightarrow (\operatorname{idl}(\operatorname{qp\text{-}csr}(p)) \simeq \operatorname{qp\text{-}null}(p)))}}
\vnbprop{qp-csr-is-ringoid}{\vnbstmt{\forall p.\; (\operatorname{zz\text{-}prime}(p) \Rightarrow \operatorname{is\text{-}ringoid}(\operatorname{qp\text{-}csr}(p)))}}
\vnbprop{qp-ring-is-ring}{\vnbstmt{\forall p.\; (\operatorname{zz\text{-}prime}(p) \Rightarrow \operatorname{is\text{-}ring}(\operatorname{qp\text{-}ring}(p)))}}
\vnbprop{qp-diff-null-iff}{\vnbstmt{\begin{array}{@{}l@{}} \forall p.; \forall a, b \in \operatorname{qp\text{-}cs}(p). \\ \quad \operatorname{zz\text{-}prime}(p) \Rightarrow  \\ \quad \quad (\operatorname{qp\text{-}sadd}(p)(a, \operatorname{qp\text{-}sneg}(p)(b)) \in \operatorname{qp\text{-}null}(p)) \Leftrightarrow \operatorname{cseq\text{-}equiv}(\operatorname{nf\text{-}metric\text{-}space}(\operatorname{qq\text{-}padic}(p)), a, b) \end{array}}}
\vnbprop{qp-cs-cseq}{\vnbstmt{\forall p.\; (\operatorname{qp\text{-}cs}(p) \simeq \operatorname{cseq}(\operatorname{nf\text{-}metric\text{-}space}(\operatorname{qq\text{-}padic}(p))))}}
\vnbprop{qp-rel-eq}{\vnbstmt{\begin{array}{@{}l@{}} \forall p. \\ \quad \operatorname{zz\text{-}prime}(p) \Rightarrow (\operatorname{ringoid\text{-}rel}(\operatorname{qp\text{-}csr}(p)) = \operatorname{crel}(\operatorname{nf\text{-}metric\text{-}space}(\operatorname{qq\text{-}padic}(p)))) \end{array}}}
\vnbprop{qp-setoid-eq}{\vnbstmt{\begin{array}{@{}l@{}} \forall p. \\ \quad \operatorname{zz\text{-}prime}(p) \Rightarrow  \\ \quad \quad (\operatorname{ringoid\text{-}setoid}(\operatorname{qp\text{-}csr}(p)) \simeq \operatorname{cauchy\text{-}setoid}(\operatorname{nf\text{-}metric\text{-}space}(\operatorname{qq\text{-}padic}(p)))) \end{array}}}
\vnbprop{qp-smul-comm}{\vnbstmt{\begin{array}{@{}l@{}} \forall p.; \forall u, v \in \operatorname{qp\text{-}cs}(p). \\ \quad \operatorname{zz\text{-}prime}(p) \Rightarrow (\operatorname{qp\text{-}smul}(p)(u, v) = \operatorname{qp\text{-}smul}(p)(v, u)) \end{array}}}
\vnbprop{qp-ring-is-commutative}{\vnbstmt{\forall p.\; (\operatorname{zz\text{-}prime}(p) \Rightarrow \operatorname{is\text{-}commutative\text{-}ring}(\operatorname{qp\text{-}ring}(p)))}}
\vnbprop{qp-abs-limit}{\vnbstmt{\begin{array}{@{}l@{}} \forall p.; \forall x \in \operatorname{qp\text{-}cs}(p), v \in \mathbb{R}. \\ \quad (\operatorname{zz\text{-}prime}(p) \wedge \operatorname{converges\text{-}to}(\operatorname{rr\text{-}ms}, (\lambda n \in \mathbb{N}.\; \operatorname{padic\text{-}abs}(p)(x(n))), v)) \Rightarrow  \\ \quad \quad \operatorname{qp\text{-}abs}(p)(\operatorname{class}(\operatorname{cauchy\text{-}setoid}(\operatorname{nf\text{-}metric\text{-}space}(\operatorname{qq\text{-}padic}(p))), x)) = v \end{array}}}
\vnbprop{qp-abs-null}{\vnbstmt{\begin{array}{@{}l@{}} \forall p.; \forall x \in \operatorname{qp\text{-}null}(p). \\ \quad \operatorname{zz\text{-}prime}(p) \Rightarrow  \\ \quad \quad \operatorname{qp\text{-}abs}(p)(\operatorname{class}(\operatorname{cauchy\text{-}setoid}(\operatorname{nf\text{-}metric\text{-}space}(\operatorname{qq\text{-}padic}(p))), x)) = 0 \end{array}}}
\vnbprop{qp-abs-const}{\vnbstmt{\begin{array}{@{}l@{}} \forall p.; \forall x \in \operatorname{qp\text{-}cs}(p), a \in \mathbb{N}, c \in \mathbb{R}. \\ \quad (\operatorname{zz\text{-}prime}(p) \wedge \forall k \in \mathbb{N}.\; ((a \leq k) \Rightarrow (\operatorname{padic\text{-}abs}(p)(x(k)) = c))) \Rightarrow  \\ \quad \quad \operatorname{qp\text{-}abs}(p)(\operatorname{class}(\operatorname{cauchy\text{-}setoid}(\operatorname{nf\text{-}metric\text{-}space}(\operatorname{qq\text{-}padic}(p))), x)) = c \end{array}}}
\vnbprop{qp-eventually-const}{\vnbstmt{\begin{array}{@{}l@{}} \forall p.; \forall x \in \operatorname{qp\text{-}cs}(p). \\ \quad (\operatorname{zz\text{-}prime}(p) \wedge \neg (x \in \operatorname{qp\text{-}null}(p))) \Rightarrow  \\ \quad \quad \exists a \in \mathbb{N}, c \in \mathbb{R}. \\ \quad \quad \quad (0 < c) \wedge \forall k \in \mathbb{N}.\; ((a \leq k) \Rightarrow (\operatorname{padic\text{-}abs}(p)(x(k)) = c)) \end{array}}}
\vnbprop{qp-carr}{\vnbstmt{\begin{array}{@{}l@{}} \forall p. \\ \quad \operatorname{zz\text{-}prime}(p) \Rightarrow  \\ \quad \quad (\operatorname{carr}(\operatorname{qp\text{-}ring}(p)) \simeq \operatorname{quotient}(\operatorname{cauchy\text{-}setoid}(\operatorname{nf\text{-}metric\text{-}space}(\operatorname{qq\text{-}padic}(p))))) \end{array}}}
\vnbprop{qp-add-class}{\vnbstmt{\begin{array}{@{}l@{}} \forall p.; \forall x, y \in \operatorname{qp\text{-}cs}(p). \\ \quad \operatorname{zz\text{-}prime}(p) \Rightarrow  \\ \quad \quad \operatorname{add}(\operatorname{qp\text{-}ring}(p))(\operatorname{class}(\operatorname{cauchy\text{-}setoid}(\operatorname{nf\text{-}metric\text{-}space}(\operatorname{qq\text{-}padic}(p))), x), \\ \quad \quad \quad \operatorname{class}(\operatorname{cauchy\text{-}setoid}(\operatorname{nf\text{-}metric\text{-}space}(\operatorname{qq\text{-}padic}(p))), y)) \\ \quad \quad =\; \operatorname{class}(\operatorname{cauchy\text{-}setoid}(\operatorname{nf\text{-}metric\text{-}space}(\operatorname{qq\text{-}padic}(p))), \operatorname{qp\text{-}sadd}(p)(x, y)) \end{array}}}
\vnbprop{qp-mul-class}{\vnbstmt{\begin{array}{@{}l@{}} \forall p.; \forall x, y \in \operatorname{qp\text{-}cs}(p). \\ \quad \operatorname{zz\text{-}prime}(p) \Rightarrow  \\ \quad \quad \operatorname{mul}(\operatorname{qp\text{-}ring}(p))(\operatorname{class}(\operatorname{cauchy\text{-}setoid}(\operatorname{nf\text{-}metric\text{-}space}(\operatorname{qq\text{-}padic}(p))), x), \\ \quad \quad \quad \operatorname{class}(\operatorname{cauchy\text{-}setoid}(\operatorname{nf\text{-}metric\text{-}space}(\operatorname{qq\text{-}padic}(p))), y)) \\ \quad \quad =\; \operatorname{class}(\operatorname{cauchy\text{-}setoid}(\operatorname{nf\text{-}metric\text{-}space}(\operatorname{qq\text{-}padic}(p))), \operatorname{qp\text{-}smul}(p)(x, y)) \end{array}}}
\vnbprop{qp-neg-class}{\vnbstmt{\begin{array}{@{}l@{}} \forall p.; \forall x \in \operatorname{qp\text{-}cs}(p). \\ \quad \operatorname{zz\text{-}prime}(p) \Rightarrow  \\ \quad \quad \operatorname{neg}(\operatorname{qp\text{-}ring}(p))(\operatorname{class}(\operatorname{cauchy\text{-}setoid}(\operatorname{nf\text{-}metric\text{-}space}(\operatorname{qq\text{-}padic}(p))), x)) \\ \quad \quad =\; \operatorname{class}(\operatorname{cauchy\text{-}setoid}(\operatorname{nf\text{-}metric\text{-}space}(\operatorname{qq\text{-}padic}(p))), \operatorname{qp\text{-}sneg}(p)(x)) \end{array}}}
\vnbprop{qp-zero-class}{\vnbstmt{\begin{array}{@{}l@{}} \forall p. \\ \quad \operatorname{zz\text{-}prime}(p) \Rightarrow  \\ \quad \quad \operatorname{zero}(\operatorname{qp\text{-}ring}(p)) = \operatorname{class}(\operatorname{cauchy\text{-}setoid}(\operatorname{nf\text{-}metric\text{-}space}(\operatorname{qq\text{-}padic}(p))), (\lambda n \in \mathbb{N}.\; 0)) \end{array}}}
\vnbprop{qp-one-class}{\vnbstmt{\begin{array}{@{}l@{}} \forall p. \\ \quad \operatorname{zz\text{-}prime}(p) \Rightarrow  \\ \quad \quad \operatorname{one}(\operatorname{qp\text{-}ring}(p)) = \operatorname{class}(\operatorname{cauchy\text{-}setoid}(\operatorname{nf\text{-}metric\text{-}space}(\operatorname{qq\text{-}padic}(p))), (\lambda n \in \mathbb{N}.\; 1)) \end{array}}}
\vnbprop{qp-rep}{\vnbstmt{\begin{array}{@{}l@{}} \forall p.; \forall a \in \operatorname{quotient}(\operatorname{cauchy\text{-}setoid}(\operatorname{nf\text{-}metric\text{-}space}(\operatorname{qq\text{-}padic}(p)))). \\ \quad \operatorname{zz\text{-}prime}(p) \Rightarrow  \\ \quad \quad \exists x \in \operatorname{qp\text{-}cs}(p). \\ \quad \quad \quad a = \operatorname{class}(\operatorname{cauchy\text{-}setoid}(\operatorname{nf\text{-}metric\text{-}space}(\operatorname{qq\text{-}padic}(p))), x) \end{array}}}
\vnbprop{qp-class-zero-iff}{\vnbstmt{\begin{array}{@{}l@{}} \forall p.; \forall x \in \operatorname{qp\text{-}cs}(p). \\ \quad \operatorname{zz\text{-}prime}(p) \Rightarrow  \\ \quad \quad \operatorname{class}(\operatorname{cauchy\text{-}setoid}(\operatorname{nf\text{-}metric\text{-}space}(\operatorname{qq\text{-}padic}(p))), x) \\ \quad \quad =\; \operatorname{class}(\operatorname{cauchy\text{-}setoid}(\operatorname{nf\text{-}metric\text{-}space}(\operatorname{qq\text{-}padic}(p))), (\lambda n \in \mathbb{N}.\; 0)) \\ \quad \quad \Leftrightarrow\; x \in \operatorname{qp\text{-}null}(p) \end{array}}}
\vnbprop{qp-abs-fun}{\vnbstmt{\begin{array}{@{}l@{}} \forall p. \\ \quad \operatorname{zz\text{-}prime}(p) \Rightarrow  \\ \quad \quad \operatorname{qp\text{-}abs}(p) \in (\operatorname{quotient}(\operatorname{cauchy\text{-}setoid}(\operatorname{nf\text{-}metric\text{-}space}(\operatorname{qq\text{-}padic}(p)))) \to \mathbb{R}) \end{array}}}
\vnbprop{qp-abs-nonneg}{\vnbstmt{\begin{array}{@{}l@{}} \forall p.; \forall x \in \operatorname{qp\text{-}cs}(p). \\ \quad \operatorname{zz\text{-}prime}(p) \Rightarrow  \\ \quad \quad 0 \leq \operatorname{qp\text{-}abs}(p)(\operatorname{class}(\operatorname{cauchy\text{-}setoid}(\operatorname{nf\text{-}metric\text{-}space}(\operatorname{qq\text{-}padic}(p))), x)) \end{array}}}
\vnbprop{qp-abs-zero-iff}{\vnbstmt{\begin{array}{@{}l@{}} \forall p.; \forall x \in \operatorname{qp\text{-}cs}(p). \\ \quad \operatorname{zz\text{-}prime}(p) \Rightarrow  \\ \quad \quad (\operatorname{qp\text{-}abs}(p)(\operatorname{class}(\operatorname{cauchy\text{-}setoid}(\operatorname{nf\text{-}metric\text{-}space}(\operatorname{qq\text{-}padic}(p))), x)) = 0) \Leftrightarrow (x \in \operatorname{qp\text{-}null}(p)) \end{array}}}
\vnbprop{qp-abs-mul}{\vnbstmt{\begin{array}{@{}l@{}} \forall p.; \forall x, y \in \operatorname{qp\text{-}cs}(p). \\ \quad \operatorname{zz\text{-}prime}(p) \Rightarrow  \\ \quad \quad \operatorname{qp\text{-}abs}(p)(\operatorname{class}(\operatorname{cauchy\text{-}setoid}(\operatorname{nf\text{-}metric\text{-}space}(\operatorname{qq\text{-}padic}(p))), \operatorname{qp\text{-}smul}(p)(x, y))) \\ \quad \quad =\; \operatorname{qp\text{-}abs}(p)(\operatorname{class}(\operatorname{cauchy\text{-}setoid}(\operatorname{nf\text{-}metric\text{-}space}(\operatorname{qq\text{-}padic}(p))), x)) \cdot \operatorname{qp\text{-}abs}(p)(\operatorname{class}(\operatorname{cauchy\text{-}setoid}(\operatorname{nf\text{-}metric\text{-}space}(\operatorname{qq\text{-}padic}(p))), y)) \end{array}}}
\vnbprop{qp-class-eq-intro}{\vnbstmt{\begin{array}{@{}l@{}} \forall p.; \forall a, b \in \operatorname{qp\text{-}cs}(p). \\ \quad (\operatorname{zz\text{-}prime}(p) \wedge (\operatorname{qp\text{-}sadd}(p)(a, \operatorname{qp\text{-}sneg}(p)(b)) \in \operatorname{qp\text{-}null}(p))) \Rightarrow  \\ \quad \quad \operatorname{class}(\operatorname{cauchy\text{-}setoid}(\operatorname{nf\text{-}metric\text{-}space}(\operatorname{qq\text{-}padic}(p))), a) \\ \quad \quad =\; \operatorname{class}(\operatorname{cauchy\text{-}setoid}(\operatorname{nf\text{-}metric\text{-}space}(\operatorname{qq\text{-}padic}(p))), b) \end{array}}}
\vnbprop{qp-abs-ultra}{\vnbstmt{\begin{array}{@{}l@{}} \forall p.; \forall x, y \in \operatorname{qp\text{-}cs}(p). \\ \quad \operatorname{zz\text{-}prime}(p) \Rightarrow  \\ \quad \quad \operatorname{qp\text{-}abs}(p)(\operatorname{class}(\operatorname{cauchy\text{-}setoid}(\operatorname{nf\text{-}metric\text{-}space}(\operatorname{qq\text{-}padic}(p))), \operatorname{qp\text{-}sadd}(p)(x, y))) \\ \quad \quad \leq\; \operatorname{max}(\operatorname{qp\text{-}abs}(p)(\operatorname{class}(\operatorname{cauchy\text{-}setoid}(\operatorname{nf\text{-}metric\text{-}space}(\operatorname{qq\text{-}padic}(p))), x)), \\ \quad \quad \quad \operatorname{qp\text{-}abs}(p)(\operatorname{class}(\operatorname{cauchy\text{-}setoid}(\operatorname{nf\text{-}metric\text{-}space}(\operatorname{qq\text{-}padic}(p))), y))) \end{array}}}
\vnbprop{qp-zero-ne-one}{\vnbstmt{\begin{array}{@{}l@{}} \forall p. \\ \quad \operatorname{zz\text{-}prime}(p) \Rightarrow  \\ \quad \quad \neg (\operatorname{class}(\operatorname{cauchy\text{-}setoid}(\operatorname{nf\text{-}metric\text{-}space}(\operatorname{qq\text{-}padic}(p))), (\lambda n \in \mathbb{N}.\; 0)) = \operatorname{class}(\operatorname{cauchy\text{-}setoid}(\operatorname{nf\text{-}metric\text{-}space}(\operatorname{qq\text{-}padic}(p))), (\lambda n \in \mathbb{N}.\; 1))) \end{array}}}
\vnbprop{qp-inverse-seq}{\vnbstmt{\begin{array}{@{}l@{}} \forall p.; \forall x \in \operatorname{qp\text{-}cs}(p). \\ \quad (\operatorname{zz\text{-}prime}(p) \wedge \neg (x \in \operatorname{qp\text{-}null}(p))) \Rightarrow  \\ \quad \quad \exists y \in \operatorname{qp\text{-}cs}(p). \\ \quad \quad \quad \operatorname{qp\text{-}sadd}(p)(\operatorname{qp\text{-}smul}(p)(x, y), \operatorname{qp\text{-}sneg}(p)((\lambda n \in \mathbb{N}.\; 1))) \in \operatorname{qp\text{-}null}(p) \end{array}}}
\vnbprop{qp-field-carr}{\vnbstmt{\forall p.\; (\operatorname{zz\text{-}prime}(p) \Rightarrow (\operatorname{carr}(\operatorname{qp\text{-}field}(p)) \simeq \operatorname{carr}(\operatorname{qp\text{-}ring}(p))))}}
\vnbprop{qp-field-add}{\vnbstmt{\forall p.\; (\operatorname{zz\text{-}prime}(p) \Rightarrow (\operatorname{add}(\operatorname{qp\text{-}field}(p)) \simeq \operatorname{add}(\operatorname{qp\text{-}ring}(p))))}}
\vnbprop{qp-field-mul}{\vnbstmt{\forall p.\; (\operatorname{zz\text{-}prime}(p) \Rightarrow (\operatorname{mul}(\operatorname{qp\text{-}field}(p)) \simeq \operatorname{mul}(\operatorname{qp\text{-}ring}(p))))}}
\vnbprop{qp-field-neg}{\vnbstmt{\forall p.\; (\operatorname{zz\text{-}prime}(p) \Rightarrow (\operatorname{neg}(\operatorname{qp\text{-}field}(p)) \simeq \operatorname{neg}(\operatorname{qp\text{-}ring}(p))))}}
\vnbprop{qp-field-zero}{\vnbstmt{\forall p.\; (\operatorname{zz\text{-}prime}(p) \Rightarrow (\operatorname{zero}(\operatorname{qp\text{-}field}(p)) \simeq \operatorname{zero}(\operatorname{qp\text{-}ring}(p))))}}
\vnbprop{qp-field-one}{\vnbstmt{\forall p.\; (\operatorname{zz\text{-}prime}(p) \Rightarrow (\operatorname{one}(\operatorname{qp\text{-}field}(p)) \simeq \operatorname{one}(\operatorname{qp\text{-}ring}(p))))}}
\vnbprop{qp-field-fnrm}{\vnbstmt{\forall p.\; (\operatorname{zz\text{-}prime}(p) \Rightarrow (\operatorname{fnrm}(\operatorname{qp\text{-}field}(p)) \simeq \operatorname{qp\text{-}abs}(p)))}}
\vnbprop{qp-is-normed-field}{\vnbstmt{\forall p.\; (\operatorname{zz\text{-}prime}(p) \Rightarrow \operatorname{is\text{-}normed\text{-}field}(\operatorname{qp\text{-}field}(p)))}}
\vnbprop{qp-is-nv-field}{\vnbstmt{\forall p.\; (\operatorname{zz\text{-}prime}(p) \Rightarrow \operatorname{is\text{-}nv\text{-}field}(\operatorname{qp\text{-}field}(p)))}}
\vnbprop{qp-abs-converges}{\vnbstmt{\begin{array}{@{}l@{}} \forall p.; \forall x \in \operatorname{qp\text{-}cs}(p). \\ \quad \operatorname{zz\text{-}prime}(p) \Rightarrow  \\ \quad \quad \operatorname{converges\text{-}to}(\operatorname{rr\text{-}ms}, (\lambda n \in \mathbb{N}.\; \operatorname{padic\text{-}abs}(p)(x(n))), \\ \quad \quad \quad \operatorname{qp\text{-}abs}(p)(\operatorname{class}(\operatorname{cauchy\text{-}setoid}(\operatorname{nf\text{-}metric\text{-}space}(\operatorname{qq\text{-}padic}(p))), x))) \end{array}}}
\vnbprop{ccr-1-26-complete}{\vnbstmt{\forall p.\; (\operatorname{zz\text{-}prime}(p) \Rightarrow (\operatorname{is\text{-}metric\text{-}space}(\operatorname{qp}(p)) \wedge \operatorname{is\text{-}complete}(\operatorname{qp}(p))))}}
\vnbprop{qp-dist-norm}{\vnbstmt{\begin{array}{@{}l@{}} \forall p.; \forall x, y \in \operatorname{qp\text{-}cs}(p). \\ \quad \operatorname{zz\text{-}prime}(p) \Rightarrow  \\ \quad \quad \operatorname{dist}(\operatorname{qp}(p))(\operatorname{class}(\operatorname{cauchy\text{-}setoid}(\operatorname{nf\text{-}metric\text{-}space}(\operatorname{qq\text{-}padic}(p))), x), \\ \quad \quad \quad \operatorname{class}(\operatorname{cauchy\text{-}setoid}(\operatorname{nf\text{-}metric\text{-}space}(\operatorname{qq\text{-}padic}(p))), y)) \\ \quad \quad =\; \operatorname{qp\text{-}abs}(p)(\operatorname{class}(\operatorname{cauchy\text{-}setoid}(\operatorname{nf\text{-}metric\text{-}space}(\operatorname{qq\text{-}padic}(p))), \operatorname{qp\text{-}sadd}(p)(x, \operatorname{qp\text{-}sneg}(p)(y)))) \end{array}}}
\vnbprop{ccr-1-26-carrier}{\vnbstmt{\forall p.\; (\operatorname{zz\text{-}prime}(p) \Rightarrow (\operatorname{carr}(\operatorname{qp\text{-}field}(p)) \simeq \operatorname{pts}(\operatorname{qp}(p))))}}
\vnbprop{ccr-1-26-metric}{\vnbstmt{\begin{array}{@{}l@{}} \forall p.; \forall a, b \in \operatorname{pts}(\operatorname{qp}(p)). \\ \quad \operatorname{zz\text{-}prime}(p) \Rightarrow  \\ \quad \quad \operatorname{dist}(\operatorname{qp}(p))(a, b) = \operatorname{fnrm}(\operatorname{qp\text{-}field}(p))(\operatorname{add}(\operatorname{qp\text{-}field}(p))(a, \operatorname{neg}(\operatorname{qp\text{-}field}(p))(b))) \end{array}}}
\vnbprop{ccr-1-26-embed}{\vnbstmt{\begin{array}{@{}l@{}} \forall p.; \forall u, v \in \mathbb{Q}. \\ \quad \operatorname{zz\text{-}prime}(p) \Rightarrow  \\ \quad \quad \operatorname{dist}(\operatorname{qp}(p))(\operatorname{embed}(\operatorname{nf\text{-}metric\text{-}space}(\operatorname{qq\text{-}padic}(p)))(u), \operatorname{embed}(\operatorname{nf\text{-}metric\text{-}space}(\operatorname{qq\text{-}padic}(p)))(v)) \\ \quad \quad =\; \operatorname{padic\text{-}abs}(p)((u + -v)) \end{array}}}
\vnbprop{ccr-1-26-zp}{\vnbstmt{\forall p.\; (\operatorname{zz\text{-}prime}(p) \Rightarrow (\operatorname{zp}(p) \simeq \{v \in \operatorname{pts}(\operatorname{qp}(p)) : (\operatorname{qp\text{-}abs}(p)(v) \leq 1)\}))}}
\vnbprop{qp-embed-value}{\vnbstmt{\begin{array}{@{}l@{}} \forall p.; \forall u \in \mathbb{Q}. \\ \quad \operatorname{zz\text{-}prime}(p) \Rightarrow  \\ \quad \quad \operatorname{embed}(\operatorname{nf\text{-}metric\text{-}space}(\operatorname{qq\text{-}padic}(p)))(u) \\ \quad \quad =\; \operatorname{class}(\operatorname{cauchy\text{-}setoid}(\operatorname{nf\text{-}metric\text{-}space}(\operatorname{qq\text{-}padic}(p))), (\lambda n \in \mathbb{N}.\; u)) \end{array}}}
\vnbprop{qp-abs-le}{\vnbstmt{\begin{array}{@{}l@{}} \forall p.; \forall x \in \operatorname{qp\text{-}cs}(p), b \in \mathbb{N}, a \in \mathbb{R}. \\ \quad (\operatorname{zz\text{-}prime}(p) \wedge \forall k \in \mathbb{N}.\; ((b \leq k) \Rightarrow (\operatorname{padic\text{-}abs}(p)(x(k)) \leq a))) \Rightarrow  \\ \quad \quad \operatorname{qp\text{-}abs}(p)(\operatorname{class}(\operatorname{cauchy\text{-}setoid}(\operatorname{nf\text{-}metric\text{-}space}(\operatorname{qq\text{-}padic}(p))), x)) \leq a \end{array}}}
\vnbprop{qp-eventually-le-one}{\vnbstmt{\begin{array}{@{}l@{}} \forall p.; \forall x \in \operatorname{qp\text{-}cs}(p). \\ \quad (\operatorname{zz\text{-}prime}(p) \wedge (\operatorname{qp\text{-}abs}(p)(\operatorname{class}(\operatorname{cauchy\text{-}setoid}(\operatorname{nf\text{-}metric\text{-}space}(\operatorname{qq\text{-}padic}(p))), x)) \leq 1)) \Rightarrow  \\ \quad \quad \exists b \in \mathbb{N}.\; \forall k \in \mathbb{N}.\; ((b \leq k) \Rightarrow (\operatorname{padic\text{-}abs}(p)(x(k)) \leq 1)) \end{array}}}
\vnbprop{ccr-1-28}{\vnbstmt{\begin{array}{@{}l@{}} \forall p. \\ \quad \operatorname{zz\text{-}prime}(p) \Rightarrow  \\ \quad \quad \forall h \in \mathbb{Z}.\; (\operatorname{embed}(\operatorname{nf\text{-}metric\text{-}space}(\operatorname{qq\text{-}padic}(p)))(h) \in \operatorname{zp}(p)) \wedge \\ \quad \quad \forall a \in \operatorname{zp}(p), s. \\ \quad \quad \quad \operatorname{pos\text{-}rr}(s) \Rightarrow  \\ \quad \quad \quad \quad \exists h \in \mathbb{Z}. \\ \quad \quad \quad \quad \quad \operatorname{dist}(\operatorname{qp}(p))(a, \operatorname{embed}(\operatorname{nf\text{-}metric\text{-}space}(\operatorname{qq\text{-}padic}(p)))(h)) < s \end{array}}}
\vnbprop{nv-field-fractions-aux}{\vnbstmt{\begin{array}{@{}l@{}} \forall k.; \forall q \in \operatorname{carr}(k). \\ \quad \operatorname{is\text{-}nv\text{-}field}(k) \Rightarrow  \\ \quad \quad \exists u, v \in \operatorname{nv\text{-}integers}(k). \\ \quad \quad \quad \neg (v = \operatorname{zero}(k)) \wedge (\operatorname{mul}(k)(q, v) = u) \end{array}}}
\vnbprop{nv-field-fractions-of-integers}{\vnbstmt{\forall k.\; (\operatorname{is\text{-}nv\text{-}field}(k) \Rightarrow \operatorname{is\text{-}fractions\text{-}of}(k, \operatorname{nv\text{-}integers}(k)))}}
\vnbprop{ccr-1-29}{\vnbstmt{\forall p.\; (\operatorname{zz\text{-}prime}(p) \Rightarrow \operatorname{is\text{-}fractions\text{-}of}(\operatorname{qp\text{-}field}(p), \operatorname{zp}(p)))}}
\vnbprop{nv-translation-continuous}{\vnbstmt{\begin{array}{@{}l@{}} \forall k.; \forall t \in \operatorname{carr}(k). \\ \quad \operatorname{is\text{-}normed\text{-}field}(k) \Rightarrow  \\ \quad \quad \operatorname{is\text{-}continuous}(\operatorname{nf\text{-}metric\text{-}space}(k), \operatorname{nf\text{-}metric\text{-}space}(k), \\ \quad \quad \quad (\lambda v \in \operatorname{pts}(\operatorname{nf\text{-}metric\text{-}space}(k)).\; \operatorname{add}(k)(t, v))) \end{array}}}
\vnbprop{nv-translate-integers}{\vnbstmt{\begin{array}{@{}l@{}} \forall k.; \forall t \in \operatorname{carr}(k). \\ \quad \operatorname{is\text{-}nv\text{-}field}(k) \Rightarrow  \\ \quad \quad \operatorname{image}((\lambda v \in \operatorname{pts}(\operatorname{nf\text{-}metric\text{-}space}(k)).\; \operatorname{add}(k)(t, v)), \operatorname{nv\text{-}integers}(k)) \\ \quad \quad =\; \operatorname{closed\text{-}ball}(\operatorname{nf\text{-}metric\text{-}space}(k), t, 1) \end{array}}}
\vnbprop{nv-local-of-compact-integers}{\vnbstmt{\begin{array}{@{}l@{}} \forall k. \\ \quad \operatorname{is\text{-}nv\text{-}field}(k) \wedge \\ \quad \operatorname{is\text{-}compact}(\operatorname{subspace\text{-}ms}(\operatorname{nf\text{-}metric\text{-}space}(k), \operatorname{nv\text{-}integers}(k))) \wedge \\ \quad \forall s. \\ \quad \quad \operatorname{pos\text{-}rr}(s) \Rightarrow \exists a \in \operatorname{carr}(k).\; (\neg (a = \operatorname{zero}(k)) \wedge (\operatorname{fnrm}(k)(a) < s)) \Rightarrow \\ \quad \quad \operatorname{is\text{-}local\text{-}nv\text{-}field}(k) \end{array}}}
\vnbprop{qp-embed-abs}{\vnbstmt{\begin{array}{@{}l@{}} \forall p.; \forall u \in \mathbb{Q}. \\ \quad \operatorname{zz\text{-}prime}(p) \Rightarrow  \\ \quad \quad \operatorname{embed}(\operatorname{nf\text{-}metric\text{-}space}(\operatorname{qq\text{-}padic}(p)))(u) \in \operatorname{pts}(\operatorname{qp}(p)) \wedge \\ \quad \quad \operatorname{qp\text{-}abs}(p)(\operatorname{embed}(\operatorname{nf\text{-}metric\text{-}space}(\operatorname{qq\text{-}padic}(p)))(u)) = \operatorname{padic\text{-}abs}(p)(u) \end{array}}}
\vnbprop{qp-abs-dist}{\vnbstmt{\begin{array}{@{}l@{}} \forall p.; \forall z \in \operatorname{pts}(\operatorname{qp}(p)). \\ \quad \operatorname{zz\text{-}prime}(p) \Rightarrow  \\ \quad \quad \operatorname{qp\text{-}abs}(p)(z) \\ \quad \quad =\; \operatorname{dist}(\operatorname{qp}(p))(\operatorname{class}(\operatorname{cauchy\text{-}setoid}(\operatorname{nf\text{-}metric\text{-}space}(\operatorname{qq\text{-}padic}(p))), (\lambda n \in \mathbb{N}.\; 0)), z) \end{array}}}
\vnbprop{qp-zp-closed-ball}{\vnbstmt{\begin{array}{@{}l@{}} \forall p. \\ \quad \operatorname{zz\text{-}prime}(p) \Rightarrow  \\ \quad \quad \operatorname{zp}(p) \\ \quad \quad =\; \operatorname{closed\text{-}ball}(\operatorname{qp}(p), \operatorname{class}(\operatorname{cauchy\text{-}setoid}(\operatorname{nf\text{-}metric\text{-}space}(\operatorname{qq\text{-}padic}(p))), (\lambda n \in \mathbb{N}.\; 0)), 1) \end{array}}}
\vnbprop{qp-unfold}{\vnbstmt{\forall p.\; (\operatorname{qp}(p) \simeq \operatorname{completion}(\operatorname{nf\text{-}metric\text{-}space}(\operatorname{qq\text{-}padic}(p))))}}
\vnbprop{qp-zp-totally-bounded}{\vnbstmt{\forall p.\; (\operatorname{zz\text{-}prime}(p) \Rightarrow \operatorname{totally\text{-}bounded}(\operatorname{subspace\text{-}ms}(\operatorname{qp}(p), \operatorname{zp}(p))))}}
\vnbprop{qp-zp-compact}{\vnbstmt{\forall p.\; (\operatorname{zz\text{-}prime}(p) \Rightarrow \operatorname{is\text{-}compact}(\operatorname{subspace\text{-}ms}(\operatorname{qp}(p), \operatorname{zp}(p))))}}
\vnbprop{qp-nf-metric-eq}{\vnbstmt{\forall p.\; (\operatorname{zz\text{-}prime}(p) \Rightarrow (\operatorname{nf\text{-}metric\text{-}space}(\operatorname{qp\text{-}field}(p)) \simeq \operatorname{qp}(p)))}}
\vnbprop{qp-small-elements}{\vnbstmt{\begin{array}{@{}l@{}} \forall p, s. \\ \quad (\operatorname{zz\text{-}prime}(p) \wedge \operatorname{pos\text{-}rr}(s)) \Rightarrow  \\ \quad \quad \exists a \in \operatorname{carr}(\operatorname{qp\text{-}field}(p)). \\ \quad \quad \quad \neg (a = \operatorname{zero}(\operatorname{qp\text{-}field}(p))) \wedge (\operatorname{fnrm}(\operatorname{qp\text{-}field}(p))(a) < s) \end{array}}}
\vnbprop{ccr-2-18}{\vnbstmt{\forall p.\; (\operatorname{zz\text{-}prime}(p) \Rightarrow \operatorname{is\text{-}local\text{-}nv\text{-}field}(\operatorname{qp\text{-}field}(p)))}}
\vnbprop{qo-rational-abs}{\vnbstmt{\begin{array}{@{}l@{}} \forall p.; \forall q \in \mathbb{Q}. \\ \quad (\operatorname{zz\text{-}prime}(p) \wedge \neg (q = 0)) \Rightarrow  \\ \quad \quad \exists m \in \mathbb{Z}.\; (\operatorname{padic\text{-}abs}(p)(q) = {2}^{m}) \end{array}}}
\vnbprop{qo-abs-power}{\vnbstmt{\begin{array}{@{}l@{}} \forall p.; \forall x \in \operatorname{carr}(\operatorname{qp\text{-}field}(p)). \\ \quad (\operatorname{zz\text{-}prime}(p) \wedge \neg (x = \operatorname{zero}(\operatorname{qp\text{-}field}(p)))) \Rightarrow  \\ \quad \quad \exists m \in \mathbb{Z}.\; (\operatorname{qp\text{-}abs}(p)(x) = {2}^{m}) \end{array}}}
\vnbprop{qo-ord-value}{\vnbstmt{\begin{array}{@{}l@{}} \forall p.; \forall x \in \operatorname{carr}(\operatorname{qp\text{-}field}(p)), m \in \mathbb{R}. \\ \quad (\operatorname{zz\text{-}prime}(p) \wedge (\operatorname{qp\text{-}abs}(p)(x) = {2}^{m})) \Rightarrow (\operatorname{qp\text{-}ord}(p)(x) = -m) \end{array}}}
\vnbprop{qo-ppos-abs}{\vnbstmt{\begin{array}{@{}l@{}} \forall p.; \forall k \in \mathbb{N}. \\ \quad \operatorname{zz\text{-}prime}(p) \Rightarrow  \\ \quad \quad \operatorname{embed}(\operatorname{nf\text{-}metric\text{-}space}(\operatorname{qq\text{-}padic}(p)))(\operatorname{power}(p, k)) \in \operatorname{carr}(\operatorname{qp\text{-}field}(p)) \wedge \\ \quad \quad \neg (\operatorname{embed}(\operatorname{nf\text{-}metric\text{-}space}(\operatorname{qq\text{-}padic}(p)))(\operatorname{power}(p, k)) = \operatorname{zero}(\operatorname{qp\text{-}field}(p))) \wedge \\ \quad \quad \operatorname{qp\text{-}abs}(p)(\operatorname{embed}(\operatorname{nf\text{-}metric\text{-}space}(\operatorname{qq\text{-}padic}(p)))(\operatorname{power}(p, k))) = {2}^{-k} \end{array}}}
\vnbprop{qo-ppos-inverse}{\vnbstmt{\begin{array}{@{}l@{}} \forall p.; \forall k \in \mathbb{N}. \\ \quad \operatorname{zz\text{-}prime}(p) \Rightarrow  \\ \quad \quad \exists x \in \operatorname{carr}(\operatorname{qp\text{-}field}(p)). \\ \quad \quad \quad \neg (x = \operatorname{zero}(\operatorname{qp\text{-}field}(p))) \wedge (\operatorname{qp\text{-}abs}(p)(x) = {2}^{k}) \end{array}}}
\vnbprop{ccr-1-31}{\vnbstmt{\begin{array}{@{}l@{}} \forall p. \\ \quad \operatorname{zz\text{-}prime}(p) \Rightarrow  \\ \quad \quad \forall x \in \operatorname{carr}(\operatorname{qp\text{-}field}(p)). \\ \quad \quad \quad \neg (x = \operatorname{zero}(\operatorname{qp\text{-}field}(p))) \Rightarrow (\operatorname{qp\text{-}ord}(p)(x) \in \mathbb{Z}) \wedge \\ \quad \quad \forall n \in \mathbb{Z}. \exists x \in \operatorname{carr}(\operatorname{qp\text{-}field}(p)). \\ \quad \quad \quad \neg (x = \operatorname{zero}(\operatorname{qp\text{-}field}(p))) \wedge (\operatorname{qp\text{-}ord}(p)(x) = n) \end{array}}}
\vnbprop{qo-int-halves}{\vnbstmt{\begin{array}{@{}l@{}} \forall k \in \mathbb{Z}. \\ \quad (\operatorname{max}(k, 0) \in \mathbb{N}) \wedge ((\operatorname{max}(-k, 0) \in \mathbb{N}) \wedge ((\operatorname{max}(k, 0) + -\operatorname{max}(-k, 0)) = k)) \end{array}}}
\vnbprop{nf-quotient-exists}{\vnbstmt{\begin{array}{@{}l@{}} \forall k.; \forall u, w \in \operatorname{carr}(k). \\ \quad (\operatorname{is\text{-}normed\text{-}field}(k) \wedge \neg (w = \operatorname{zero}(k))) \Rightarrow  \\ \quad \quad \exists v \in \operatorname{carr}(k).\; (\operatorname{mul}(k)(v, w) = u) \end{array}}}
\vnbprop{qo-ppow-def}{\vnbstmt{\begin{array}{@{}l@{}} \forall p.; \forall k \in \mathbb{Z}. \\ \quad \operatorname{zz\text{-}prime}(p) \Rightarrow  \\ \quad \quad \operatorname{qp\text{-}ppow}(p, k) \in \operatorname{carr}(\operatorname{qp\text{-}field}(p)) \wedge \\ \quad \quad \operatorname{mul}(\operatorname{qp\text{-}field}(p))(\operatorname{qp\text{-}ppow}(p, k), \operatorname{embed}(\operatorname{nf\text{-}metric\text{-}space}(\operatorname{qq\text{-}padic}(p)))(\operatorname{power}(p, \operatorname{max}(-k, 0)))) \\ \quad \quad =\; \operatorname{embed}(\operatorname{nf\text{-}metric\text{-}space}(\operatorname{qq\text{-}padic}(p)))(\operatorname{power}(p, \operatorname{max}(k, 0))) \end{array}}}
\vnbprop{qo-ppow-abs}{\vnbstmt{\begin{array}{@{}l@{}} \forall p.; \forall k \in \mathbb{Z}. \\ \quad \operatorname{zz\text{-}prime}(p) \Rightarrow  \\ \quad \quad \operatorname{qp\text{-}ppow}(p, k) \in \operatorname{carr}(\operatorname{qp\text{-}field}(p)) \wedge \\ \quad \quad \neg (\operatorname{qp\text{-}ppow}(p, k) = \operatorname{zero}(\operatorname{qp\text{-}field}(p))) \wedge \\ \quad \quad \operatorname{qp\text{-}abs}(p)(\operatorname{qp\text{-}ppow}(p, k)) = {2}^{-k} \end{array}}}
\vnbprop{ccr-1-32}{\vnbstmt{\begin{array}{@{}l@{}} \forall p.; \forall x \in \operatorname{carr}(\operatorname{qp\text{-}field}(p)). \\ \quad (\operatorname{zz\text{-}prime}(p) \wedge \neg (x = \operatorname{zero}(\operatorname{qp\text{-}field}(p)))) \Rightarrow  \\ \quad \quad \operatorname{mul}(\operatorname{qp\text{-}field}(p))(\operatorname{qp\text{-}ppow}(p, -\operatorname{qp\text{-}ord}(p)(x)), x) \in \operatorname{zp}(p) \wedge \\ \quad \quad \forall q \in \mathbb{Z}. \\ \quad \quad \quad (\operatorname{mul}(\operatorname{qp\text{-}field}(p))(\operatorname{qp\text{-}ppow}(p, q), x) \in \operatorname{zp}(p)) \Rightarrow (-\operatorname{qp\text{-}ord}(p)(x) \leq q) \end{array}}}
\vnbprop{nv-field-ultrametric}{\vnbstmt{\forall k.\; (\operatorname{is\text{-}nv\text{-}field}(k) \Rightarrow \operatorname{is\text{-}ultrametric}(\operatorname{nf\text{-}metric\text{-}space}(k)))}}
\vnbprop{subspace-ultrametric}{\vnbstmt{\begin{array}{@{}l@{}} \forall s.; \forall a \in \mathbf{Set}. \\ \quad (\operatorname{is\text{-}ultrametric}(s) \wedge (a \subseteq \operatorname{pts}(s))) \Rightarrow \operatorname{is\text{-}ultrametric}(\operatorname{subspace\text{-}ms}(s, a)) \end{array}}}
\vnbprop{rs-maximal-power}{\vnbstmt{\forall k, a.\; (\operatorname{is\text{-}nv\text{-}domain}(k, a) \Rightarrow (\operatorname{nv\text{-}maximal}(k, a) \in \operatorname{power}(a)))}}
\vnbprop{rs-add-val}{\vnbstmt{\begin{array}{@{}l@{}} \forall k, a.; \forall u, v \in a. \\ \quad \operatorname{is\text{-}nv\text{-}domain}(k, a) \Rightarrow  \\ \quad \quad (\lambda {[}\operatorname{rsu\_}, \operatorname{rsv\_}] \in (a \times a).\; \operatorname{add}(k)(\operatorname{rsu\_}, \operatorname{rsv\_}))(u, v) = \operatorname{add}(k)(u, v) \end{array}}}
\vnbprop{rs-mul-val}{\vnbstmt{\begin{array}{@{}l@{}} \forall k, a.; \forall u, v \in a. \\ \quad \operatorname{is\text{-}nv\text{-}domain}(k, a) \Rightarrow  \\ \quad \quad (\lambda {[}\operatorname{rsu\_}, \operatorname{rsv\_}] \in (a \times a).\; \operatorname{mul}(k)(\operatorname{rsu\_}, \operatorname{rsv\_}))(u, v) = \operatorname{mul}(k)(u, v) \end{array}}}
\vnbprop{rs-neg-val}{\vnbstmt{\begin{array}{@{}l@{}} \forall k, a.; \forall u \in a. \\ \quad \operatorname{is\text{-}nv\text{-}domain}(k, a) \Rightarrow ((\lambda \mathit{rsu} \in a.\; \operatorname{neg}(k)(\mathit{rsu}))(u) = \operatorname{neg}(k)(u)) \end{array}}}
\vnbprop{rs-add-fun}{\vnbstmt{\begin{array}{@{}l@{}} \forall k, a. \\ \quad \operatorname{is\text{-}nv\text{-}domain}(k, a) \Rightarrow  \\ \quad \quad (\lambda {[}\operatorname{rsu\_}, \operatorname{rsv\_}] \in (a \times a).\; \operatorname{add}(k)(\operatorname{rsu\_}, \operatorname{rsv\_})) \in ((a \times a) \to a) \end{array}}}
\vnbprop{rs-mul-fun}{\vnbstmt{\begin{array}{@{}l@{}} \forall k, a. \\ \quad \operatorname{is\text{-}nv\text{-}domain}(k, a) \Rightarrow  \\ \quad \quad (\lambda {[}\operatorname{rsu\_}, \operatorname{rsv\_}] \in (a \times a).\; \operatorname{mul}(k)(\operatorname{rsu\_}, \operatorname{rsv\_})) \in ((a \times a) \to a) \end{array}}}
\vnbprop{rs-neg-fun}{\vnbstmt{\begin{array}{@{}l@{}} \forall k, a. \\ \quad \operatorname{is\text{-}nv\text{-}domain}(k, a) \Rightarrow ((\lambda u \in a.\; \operatorname{neg}(k)(u)) \in (a \to a)) \end{array}}}
\vnbprop{rs-carr}{\vnbstmt{\forall k, a.\; (\operatorname{is\text{-}nv\text{-}domain}(k, a) \Rightarrow (\operatorname{carr}(\operatorname{nv\text{-}res\text{-}ringoid}(k, a)) \simeq a))}}
\vnbprop{rs-add}{\vnbstmt{\begin{array}{@{}l@{}} \forall k, a. \\ \quad \operatorname{is\text{-}nv\text{-}domain}(k, a) \Rightarrow  \\ \quad \quad (\operatorname{add}(\operatorname{nv\text{-}res\text{-}ringoid}(k, a)) \simeq (\lambda {[}\operatorname{rsu\_}, \operatorname{rsv\_}] \in (a \times a).\; \operatorname{add}(k)(\operatorname{rsu\_}, \operatorname{rsv\_}))) \end{array}}}
\vnbprop{rs-mul}{\vnbstmt{\begin{array}{@{}l@{}} \forall k, a. \\ \quad \operatorname{is\text{-}nv\text{-}domain}(k, a) \Rightarrow  \\ \quad \quad (\operatorname{mul}(\operatorname{nv\text{-}res\text{-}ringoid}(k, a)) \simeq (\lambda {[}\operatorname{rsu\_}, \operatorname{rsv\_}] \in (a \times a).\; \operatorname{mul}(k)(\operatorname{rsu\_}, \operatorname{rsv\_}))) \end{array}}}
\vnbprop{rs-neg}{\vnbstmt{\begin{array}{@{}l@{}} \forall k, a. \\ \quad \operatorname{is\text{-}nv\text{-}domain}(k, a) \Rightarrow  \\ \quad \quad (\operatorname{neg}(\operatorname{nv\text{-}res\text{-}ringoid}(k, a)) \simeq (\lambda u \in a.\; \operatorname{neg}(k)(u))) \end{array}}}
\vnbprop{rs-zero}{\vnbstmt{\forall k, a.\; (\operatorname{is\text{-}nv\text{-}domain}(k, a) \Rightarrow (\operatorname{zero}(\operatorname{nv\text{-}res\text{-}ringoid}(k, a)) \simeq \operatorname{zero}(k)))}}
\vnbprop{rs-one}{\vnbstmt{\forall k, a.\; (\operatorname{is\text{-}nv\text{-}domain}(k, a) \Rightarrow (\operatorname{one}(\operatorname{nv\text{-}res\text{-}ringoid}(k, a)) \simeq \operatorname{one}(k)))}}
\vnbprop{rs-idl}{\vnbstmt{\begin{array}{@{}l@{}} \forall k, a. \\ \quad \operatorname{is\text{-}nv\text{-}domain}(k, a) \Rightarrow (\operatorname{idl}(\operatorname{nv\text{-}res\text{-}ringoid}(k, a)) \simeq \operatorname{nv\text{-}maximal}(k, a)) \end{array}}}
\vnbprop{rs-is-ringoid}{\vnbstmt{\forall k, a.\; (\operatorname{is\text{-}nv\text{-}domain}(k, a) \Rightarrow \operatorname{is\text{-}ringoid}(\operatorname{nv\text{-}res\text{-}ringoid}(k, a)))}}
\vnbprop{rs-class-eq-iff}{\vnbstmt{\begin{array}{@{}l@{}} \forall k, a.; \forall b, c \in a. \\ \quad \operatorname{is\text{-}nv\text{-}domain}(k, a) \Rightarrow  \\ \quad \quad \operatorname{class}(\operatorname{ringoid\text{-}setoid}(\operatorname{nv\text{-}res\text{-}ringoid}(k, a)), b) = \operatorname{class}(\operatorname{ringoid\text{-}setoid}(\operatorname{nv\text{-}res\text{-}ringoid}(k, a)), c) \\ \quad \quad \Leftrightarrow\; \operatorname{add}(k)(b, \operatorname{neg}(k)(c)) \in \operatorname{nv\text{-}maximal}(k, a) \end{array}}}
\vnbprop{rs-setoid-eq}{\vnbstmt{\begin{array}{@{}l@{}} \forall k, a. \\ \quad \operatorname{is\text{-}nv\text{-}domain}(k, a) \Rightarrow  \\ \quad \quad (\operatorname{ringoid\text{-}setoid}(\operatorname{nv\text{-}res\text{-}ringoid}(k, a)) \simeq \operatorname{eps\text{-}setoid}(\operatorname{subspace\text{-}ms}(\operatorname{nf\text{-}metric\text{-}space}(k), a), 1)) \end{array}}}
\vnbprop{ccr-2-7}{\vnbstmt{\begin{array}{@{}l@{}} \forall k, a. \\ \quad \operatorname{is\text{-}compact\text{-}nv\text{-}domain}(k, a) \Rightarrow  \\ \quad \quad \operatorname{is\text{-}field\text{-}ring}(\operatorname{nv\text{-}residue}(k, a)) \wedge (|\operatorname{carr}(\operatorname{nv\text{-}residue}(k, a))| \in \mathbb{N}) \end{array}}}
\vnbprop{nv-unit-iff-norm-one}{\vnbstmt{\begin{array}{@{}l@{}} \forall k.; \forall u \in \operatorname{carr}(k). \\ \quad \operatorname{is\text{-}nv\text{-}field}(k) \Rightarrow ((u \in \operatorname{nv\text{-}units}(k, \operatorname{nv\text{-}integers}(k))) \Leftrightarrow (\operatorname{fnrm}(k)(u) = 1)) \end{array}}}
\vnbprop{nf-pair-inverse}{\vnbstmt{\begin{array}{@{}l@{}} \forall k.; \forall x, y, u, w \in \operatorname{carr}(k). \\ \quad \operatorname{is\text{-}normed\text{-}field}(k) \wedge \\ \quad \neg (u = \operatorname{zero}(k)) \wedge \\ \quad \neg (w = \operatorname{zero}(k)) \wedge \\ \quad \operatorname{mul}(k)(x, w) = u \wedge \\ \quad \operatorname{mul}(k)(y, u) = w \Rightarrow \\ \quad \quad \operatorname{mul}(k)(x, y) = \operatorname{one}(k) \end{array}}}
\vnbprop{qo-ppow-inverse}{\vnbstmt{\begin{array}{@{}l@{}} \forall p.; \forall k \in \mathbb{Z}. \\ \quad \operatorname{zz\text{-}prime}(p) \Rightarrow  \\ \quad \quad \operatorname{mul}(\operatorname{qp\text{-}field}(p))(\operatorname{qp\text{-}ppow}(p, k), \operatorname{qp\text{-}ppow}(p, -k)) = \operatorname{one}(\operatorname{qp\text{-}field}(p)) \end{array}}}
\vnbprop{ccr-4-10}{\vnbstmt{\begin{array}{@{}l@{}} \forall p.; \forall x \in \operatorname{carr}(\operatorname{qp\text{-}field}(p)). \\ \quad (\operatorname{zz\text{-}prime}(p) \wedge \neg (x = \operatorname{zero}(\operatorname{qp\text{-}field}(p)))) \Rightarrow  \\ \quad \quad \exists k \in \mathbb{Z}, u \in \operatorname{nv\text{-}units}(\operatorname{qp\text{-}field}(p), \operatorname{zp}(p)). \\ \quad \quad \quad x = \operatorname{mul}(\operatorname{qp\text{-}field}(p))(u, \operatorname{qp\text{-}ppow}(p, k)) \wedge \\ \quad \quad \forall k, l \in \mathbb{Z}, u, v \in \operatorname{nv\text{-}units}(\operatorname{qp\text{-}field}(p), \operatorname{zp}(p)). \\ \quad \quad \quad (\operatorname{mul}(\operatorname{qp\text{-}field}(p))(u, \operatorname{qp\text{-}ppow}(p, k)) = \operatorname{mul}(\operatorname{qp\text{-}field}(p))(v, \operatorname{qp\text{-}ppow}(p, l))) \Rightarrow  \\ \quad \quad \quad \quad (k = l) \wedge (u = v) \end{array}}}
\vnbprop{fp-zz-carr}{\vnbstmt{(\operatorname{carr}(\operatorname{zz\text{-}ring}) \simeq \mathbb{Z})}}
\vnbprop{fp-zz-zero}{\vnbstmt{(\operatorname{zero}(\operatorname{zz\text{-}ring}) \simeq 0)}}
\vnbprop{fp-zz-one}{\vnbstmt{(\operatorname{one}(\operatorname{zz\text{-}ring}) \simeq 1)}}
\vnbprop{fp-zz-add}{\vnbstmt{(\operatorname{add}(\operatorname{zz\text{-}ring}) \simeq (\lambda {[}\operatorname{x\_}, \operatorname{y\_}] \in (\mathbb{Z} \times \mathbb{Z}).\; (\operatorname{x\_} + \operatorname{y\_})))}}
\vnbprop{fp-zz-mul}{\vnbstmt{(\operatorname{mul}(\operatorname{zz\text{-}ring}) \simeq (\lambda {[}\operatorname{x\_}, \operatorname{y\_}] \in (\mathbb{Z} \times \mathbb{Z}).\; (\operatorname{x\_} \cdot \operatorname{y\_})))}}
\vnbprop{fp-zz-neg}{\vnbstmt{(\operatorname{neg}(\operatorname{zz\text{-}ring}) \simeq (\lambda x \in \mathbb{Z}.\; -x))}}
\vnbprop{fp-add-val}{\vnbstmt{\begin{array}{@{}l@{}} \forall u, v \in \mathbb{Z}. \\ \quad (\lambda {[}\operatorname{x\_}, \operatorname{y\_}] \in (\mathbb{Z} \times \mathbb{Z}).\; (\operatorname{x\_} + \operatorname{y\_}))(u, v) = (u + v) \end{array}}}
\vnbprop{fp-mul-val}{\vnbstmt{\begin{array}{@{}l@{}} \forall u, v \in \mathbb{Z}. \\ \quad (\lambda {[}\operatorname{x\_}, \operatorname{y\_}] \in (\mathbb{Z} \times \mathbb{Z}).\; (\operatorname{x\_} \cdot \operatorname{y\_}))(u, v) = (u \cdot v) \end{array}}}
\vnbprop{fp-neg-val}{\vnbstmt{\forall u \in \mathbb{Z}.\; ((\lambda x \in \mathbb{Z}.\; -x)(u) = -u)}}
\vnbprop{fp-pz-power}{\vnbstmt{\forall p.\; (\operatorname{zz\text{-}prime}(p) \Rightarrow (\{n \in \mathbb{Z} : \operatorname{zz\text{-}divides}(p, n)\} \in \operatorname{power}(\mathbb{Z})))}}
\vnbprop{fp-carr}{\vnbstmt{\forall p.\; (\operatorname{zz\text{-}prime}(p) \Rightarrow (\operatorname{carr}(\operatorname{fp\text{-}ringoid}(p)) \simeq \mathbb{Z}))}}
\vnbprop{fp-add}{\vnbstmt{\begin{array}{@{}l@{}} \forall p. \\ \quad \operatorname{zz\text{-}prime}(p) \Rightarrow  \\ \quad \quad (\operatorname{add}(\operatorname{fp\text{-}ringoid}(p)) \simeq (\lambda {[}\operatorname{x\_}, \operatorname{y\_}] \in (\mathbb{Z} \times \mathbb{Z}).\; (\operatorname{x\_} + \operatorname{y\_}))) \end{array}}}
\vnbprop{fp-mul}{\vnbstmt{\begin{array}{@{}l@{}} \forall p. \\ \quad \operatorname{zz\text{-}prime}(p) \Rightarrow  \\ \quad \quad (\operatorname{mul}(\operatorname{fp\text{-}ringoid}(p)) \simeq (\lambda {[}\operatorname{x\_}, \operatorname{y\_}] \in (\mathbb{Z} \times \mathbb{Z}).\; (\operatorname{x\_} \cdot \operatorname{y\_}))) \end{array}}}
\vnbprop{fp-neg}{\vnbstmt{\forall p.\; (\operatorname{zz\text{-}prime}(p) \Rightarrow (\operatorname{neg}(\operatorname{fp\text{-}ringoid}(p)) \simeq (\lambda x \in \mathbb{Z}.\; -x)))}}
\vnbprop{fp-zero}{\vnbstmt{\forall p.\; (\operatorname{zz\text{-}prime}(p) \Rightarrow (\operatorname{zero}(\operatorname{fp\text{-}ringoid}(p)) \simeq 0))}}
\vnbprop{fp-one}{\vnbstmt{\forall p.\; (\operatorname{zz\text{-}prime}(p) \Rightarrow (\operatorname{one}(\operatorname{fp\text{-}ringoid}(p)) \simeq 1))}}
\vnbprop{fp-idl}{\vnbstmt{\begin{array}{@{}l@{}} \forall p. \\ \quad \operatorname{zz\text{-}prime}(p) \Rightarrow (\operatorname{idl}(\operatorname{fp\text{-}ringoid}(p)) \simeq \{n \in \mathbb{Z} : \operatorname{zz\text{-}divides}(p, n)\}) \end{array}}}
\vnbprop{fp-is-ringoid}{\vnbstmt{\forall p.\; (\operatorname{zz\text{-}prime}(p) \Rightarrow \operatorname{is\text{-}ringoid}(\operatorname{fp\text{-}ringoid}(p)))}}
\vnbprop{fp-class-eq-iff}{\vnbstmt{\begin{array}{@{}l@{}} \forall p.; \forall a, b \in \mathbb{Z}. \\ \quad \operatorname{zz\text{-}prime}(p) \Rightarrow  \\ \quad \quad \operatorname{class}(\operatorname{ringoid\text{-}setoid}(\operatorname{fp\text{-}ringoid}(p)), a) = \operatorname{class}(\operatorname{ringoid\text{-}setoid}(\operatorname{fp\text{-}ringoid}(p)), b) \\ \quad \quad \Leftrightarrow\; \operatorname{zz\text{-}divides}(p, (a + -b)) \end{array}}}
\vnbprop{fp-is-field}{\vnbstmt{\forall p.\; (\operatorname{zz\text{-}prime}(p) \Rightarrow \operatorname{is\text{-}field\text{-}ring}(\operatorname{fp}(p)))}}
\vnbprop{qo-power-one}{\vnbstmt{\forall x \in \mathbb{R}.\; (\operatorname{power}(x, 1) = x)}}
\vnbprop{qo-int-small-iff}{\vnbstmt{\begin{array}{@{}l@{}} \forall p.; \forall n \in \mathbb{Z}. \\ \quad \operatorname{zz\text{-}prime}(p) \Rightarrow ((\operatorname{padic\text{-}abs}(p)(n) < 1) \Leftrightarrow \operatorname{zz\text{-}divides}(p, n)) \end{array}}}
\vnbprop{nv-residue-fp-iso}{\vnbstmt{\begin{array}{@{}l@{}} \forall p, k, a, o. \\ \quad \operatorname{zz\text{-}prime}(p) \wedge \\ \quad \operatorname{is\text{-}nv\text{-}domain}(k, a) \wedge \\ \quad o \in (\mathbb{Z} \to a) \wedge \\ \quad \forall m, n \in \mathbb{Z}. \\ \quad \quad (o((m + n)) = \operatorname{add}(k)(o(m), o(n))) \wedge (o((m \cdot n)) = \operatorname{mul}(k)(o(m), o(n))) \wedge \\ \quad \forall m \in \mathbb{Z}.\; (o(-m) = \operatorname{neg}(k)(o(m))) \wedge \\ \quad o(0) = \operatorname{zero}(k) \wedge \\ \quad o(1) = \operatorname{one}(k) \wedge \\ \quad \forall m \in \mathbb{Z}.\; ((o(m) \in \operatorname{nv\text{-}maximal}(k, a)) \Leftrightarrow \operatorname{zz\text{-}divides}(p, m)) \wedge \\ \quad \forall b \in a. \exists m \in \mathbb{Z}. \\ \quad \quad \operatorname{add}(k)(b, \operatorname{neg}(k)(o(m))) \in \operatorname{nv\text{-}maximal}(k, a) \Rightarrow \\ \quad \quad \exists f.\; \operatorname{is\text{-}ring\text{-}iso}(\operatorname{fp}(p), \operatorname{nv\text{-}residue}(k, a), f) \end{array}}}
\vnbprop{nv-integers-domain}{\vnbstmt{\forall k.\; (\operatorname{is\text{-}nv\text{-}field}(k) \Rightarrow \operatorname{is\text{-}nv\text{-}domain}(k, \operatorname{nv\text{-}integers}(k)))}}
\vnbprop{ccr-1-22}{\vnbstmt{\begin{array}{@{}l@{}} \forall p. \\ \quad \operatorname{zz\text{-}prime}(p) \Rightarrow  \\ \quad \quad \exists f. \\ \quad \quad \quad \operatorname{is\text{-}ring\text{-}iso}(\operatorname{fp}(p), \operatorname{nv\text{-}residue}(\operatorname{qq\text{-}padic}(p), \operatorname{nv\text{-}integers}(\operatorname{qq\text{-}padic}(p))), f) \end{array}}}
\vnbprop{qp-embed-add}{\vnbstmt{\begin{array}{@{}l@{}} \forall p.; \forall u, v \in \mathbb{Q}. \\ \quad \operatorname{zz\text{-}prime}(p) \Rightarrow  \\ \quad \quad \operatorname{embed}(\operatorname{nf\text{-}metric\text{-}space}(\operatorname{qq\text{-}padic}(p)))((u + v)) \\ \quad \quad =\; \operatorname{add}(\operatorname{qp\text{-}field}(p))(\operatorname{embed}(\operatorname{nf\text{-}metric\text{-}space}(\operatorname{qq\text{-}padic}(p)))(u), \\ \quad \quad \quad \operatorname{embed}(\operatorname{nf\text{-}metric\text{-}space}(\operatorname{qq\text{-}padic}(p)))(v)) \end{array}}}
\vnbprop{qp-embed-mul}{\vnbstmt{\begin{array}{@{}l@{}} \forall p.; \forall u, v \in \mathbb{Q}. \\ \quad \operatorname{zz\text{-}prime}(p) \Rightarrow  \\ \quad \quad \operatorname{embed}(\operatorname{nf\text{-}metric\text{-}space}(\operatorname{qq\text{-}padic}(p)))((u \cdot v)) \\ \quad \quad =\; \operatorname{mul}(\operatorname{qp\text{-}field}(p))(\operatorname{embed}(\operatorname{nf\text{-}metric\text{-}space}(\operatorname{qq\text{-}padic}(p)))(u), \\ \quad \quad \quad \operatorname{embed}(\operatorname{nf\text{-}metric\text{-}space}(\operatorname{qq\text{-}padic}(p)))(v)) \end{array}}}
\vnbprop{qp-embed-neg}{\vnbstmt{\begin{array}{@{}l@{}} \forall p.; \forall u \in \mathbb{Q}. \\ \quad \operatorname{zz\text{-}prime}(p) \Rightarrow  \\ \quad \quad \operatorname{embed}(\operatorname{nf\text{-}metric\text{-}space}(\operatorname{qq\text{-}padic}(p)))(-u) \\ \quad \quad =\; \operatorname{neg}(\operatorname{qp\text{-}field}(p))(\operatorname{embed}(\operatorname{nf\text{-}metric\text{-}space}(\operatorname{qq\text{-}padic}(p)))(u)) \end{array}}}
\vnbprop{qp-embed-zero}{\vnbstmt{\begin{array}{@{}l@{}} \forall p. \\ \quad \operatorname{zz\text{-}prime}(p) \Rightarrow (\operatorname{embed}(\operatorname{nf\text{-}metric\text{-}space}(\operatorname{qq\text{-}padic}(p)))(0) = \operatorname{zero}(\operatorname{qp\text{-}field}(p))) \end{array}}}
\vnbprop{qp-embed-one}{\vnbstmt{\begin{array}{@{}l@{}} \forall p. \\ \quad \operatorname{zz\text{-}prime}(p) \Rightarrow (\operatorname{embed}(\operatorname{nf\text{-}metric\text{-}space}(\operatorname{qq\text{-}padic}(p)))(1) = \operatorname{one}(\operatorname{qp\text{-}field}(p))) \end{array}}}
\vnbprop{ccr-2-24}{\vnbstmt{\begin{array}{@{}l@{}} \forall p. \\ \quad \operatorname{zz\text{-}prime}(p) \Rightarrow  \\ \quad \quad \exists f.\; \operatorname{is\text{-}ring\text{-}iso}(\operatorname{fp}(p), \operatorname{nv\text{-}residue}(\operatorname{qp\text{-}field}(p), \operatorname{zp}(p)), f) \end{array}}}
\vnbprop{qp-zp-compact-domain}{\vnbstmt{\forall p.\; (\operatorname{zz\text{-}prime}(p) \Rightarrow \operatorname{is\text{-}compact\text{-}nv\text{-}domain}(\operatorname{qp\text{-}field}(p), \operatorname{zp}(p)))}}
\vnbprop{qo-abs-p}{\vnbstmt{\forall p.\; (\operatorname{zz\text{-}prime}(p) \Rightarrow ((\operatorname{padic\text{-}abs}(p)(p) \cdot 2) = 1))}}
\vnbprop{qo-rpow-half}{\vnbstmt{\forall m \in \mathbb{Z}.\; (({2}^{m} < 1) \Rightarrow (({2}^{m} \cdot 2) \leq 1))}}
\vnbprop{qp-maximal-principal}{\vnbstmt{\begin{array}{@{}l@{}} \forall p. \\ \quad \operatorname{zz\text{-}prime}(p) \Rightarrow  \\ \quad \quad \operatorname{nv\text{-}maximal}(\operatorname{qp\text{-}field}(p), \operatorname{zp}(p)) \\ \quad \quad =\; \operatorname{nv\text{-}principal}(\operatorname{qp\text{-}field}(p), \operatorname{zp}(p), \operatorname{embed}(\operatorname{nf\text{-}metric\text{-}space}(\operatorname{qq\text{-}padic}(p)))(p)) \end{array}}}
\vnbprop{qo-int-mult-nn}{\vnbstmt{\begin{array}{@{}l@{}} \forall n \in \mathbb{N}, p, u.; \forall x \in u. \\ \quad (\operatorname{zz\text{-}prime}(p) \wedge \operatorname{is\text{-}add\text{-}subgroup}(\operatorname{qp\text{-}field}(p), u)) \Rightarrow  \\ \quad \quad \operatorname{mul}(\operatorname{qp\text{-}field}(p))(\operatorname{embed}(\operatorname{nf\text{-}metric\text{-}space}(\operatorname{qq\text{-}padic}(p)))(n), x) \in u \end{array}}}
\vnbprop{qo-int-mult}{\vnbstmt{\begin{array}{@{}l@{}} \forall p, u.; \forall x \in u, n \in \mathbb{Z}. \\ \quad (\operatorname{zz\text{-}prime}(p) \wedge \operatorname{is\text{-}add\text{-}subgroup}(\operatorname{qp\text{-}field}(p), u)) \Rightarrow  \\ \quad \quad \operatorname{mul}(\operatorname{qp\text{-}field}(p))(\operatorname{embed}(\operatorname{nf\text{-}metric\text{-}space}(\operatorname{qq\text{-}padic}(p)))(n), x) \in u \end{array}}}
\vnbprop{qo-open-subgroup-module}{\vnbstmt{\begin{array}{@{}l@{}} \forall p, u. \\ \quad (\operatorname{zz\text{-}prime}(p) \wedge (\operatorname{is\text{-}add\text{-}subgroup}(\operatorname{qp\text{-}field}(p), u) \wedge \operatorname{is\text{-}open}(\operatorname{nf\text{-}metric\text{-}space}(\operatorname{qp\text{-}field}(p)), u))) \Rightarrow  \\ \quad \quad \operatorname{is\text{-}nv\text{-}module}(\operatorname{qp\text{-}field}(p), \operatorname{zp}(p), u) \end{array}}}
\vnbprop{ccr-4-13}{\vnbstmt{\begin{array}{@{}l@{}} \forall p, u. \\ \quad \operatorname{zz\text{-}prime}(p) \wedge \\ \quad \operatorname{is\text{-}add\text{-}subgroup}(\operatorname{qp\text{-}field}(p), u) \wedge \\ \quad \operatorname{is\text{-}open}(\operatorname{nf\text{-}metric\text{-}space}(\operatorname{qp\text{-}field}(p)), u) \wedge \\ \quad \neg (u = \operatorname{carr}(\operatorname{qp\text{-}field}(p))) \Rightarrow \\ \quad \quad \operatorname{is\text{-}nv\text{-}module}(\operatorname{qp\text{-}field}(p), \operatorname{zp}(p), u) \wedge \\ \quad \quad \exists g \in \operatorname{carr}(\operatorname{qp\text{-}field}(p)).\; (u = \operatorname{nv\text{-}multiples}(\operatorname{qp\text{-}field}(p), \operatorname{zp}(p), g)) \end{array}}}
\vnbprop{qe-ppow-char}{\vnbstmt{\begin{array}{@{}l@{}} \forall p.; \forall k \in \mathbb{Z}, s, t \in \mathbb{N}. \\ \quad (\operatorname{zz\text{-}prime}(p) \wedge ((s + -t) = k)) \Rightarrow  \\ \quad \quad \operatorname{mul}(\operatorname{qp\text{-}field}(p))(\operatorname{qp\text{-}ppow}(p, k), \operatorname{embed}(\operatorname{nf\text{-}metric\text{-}space}(\operatorname{qq\text{-}padic}(p)))(\operatorname{power}(p, t))) \\ \quad \quad =\; \operatorname{embed}(\operatorname{nf\text{-}metric\text{-}space}(\operatorname{qq\text{-}padic}(p)))(\operatorname{power}(p, s)) \end{array}}}
\vnbprop{qe-ppow-succ}{\vnbstmt{\begin{array}{@{}l@{}} \forall p.; \forall k \in \mathbb{Z}. \\ \quad \operatorname{zz\text{-}prime}(p) \Rightarrow  \\ \quad \quad \operatorname{mul}(\operatorname{qp\text{-}field}(p))(\operatorname{qp\text{-}ppow}(p, k), \operatorname{embed}(\operatorname{nf\text{-}metric\text{-}space}(\operatorname{qq\text{-}padic}(p)))(p)) = \operatorname{qp\text{-}ppow}(p, (k + 1)) \end{array}}}
\vnbprop{qe-ag-opr}{\vnbstmt{\begin{array}{@{}l@{}} \forall p. \\ \quad \operatorname{zz\text{-}prime}(p) \Rightarrow  \\ \quad \quad (\operatorname{opr}(\operatorname{normed\text{-}field\text{-}additive\text{-}ag}(\operatorname{qp\text{-}field}(p))) \simeq \operatorname{add}(\operatorname{qp\text{-}field}(p))) \end{array}}}
\vnbprop{qe-ag-iden}{\vnbstmt{\begin{array}{@{}l@{}} \forall p. \\ \quad \operatorname{zz\text{-}prime}(p) \Rightarrow  \\ \quad \quad (\operatorname{iden}(\operatorname{normed\text{-}field\text{-}additive\text{-}ag}(\operatorname{qp\text{-}field}(p))) \simeq \operatorname{zero}(\operatorname{qp\text{-}field}(p))) \end{array}}}
\vnbprop{qe-psum-zero}{\vnbstmt{\forall p, a, n.\; (\operatorname{zz\text{-}prime}(p) \Rightarrow (\operatorname{qp\text{-}psum}(p, a, n, 0) = \operatorname{zero}(\operatorname{qp\text{-}field}(p))))}}
\vnbprop{qe-psum-succ}{\vnbstmt{\begin{array}{@{}l@{}} \forall p, a, n.; \forall m \in \mathbb{N}. \\ \quad \operatorname{zz\text{-}prime}(p) \Rightarrow  \\ \quad \quad (\operatorname{qp\text{-}psum}(p, a, n, (m + 1)) \simeq \operatorname{add}(\operatorname{qp\text{-}field}(p))(\operatorname{qp\text{-}psum}(p, a, n, m), \operatorname{mul}(\operatorname{qp\text{-}field}(p))(\operatorname{embed}(\operatorname{nf\text{-}metric\text{-}space}(\operatorname{qq\text{-}padic}(p)))(a((n + m))), \operatorname{qp\text{-}ppow}(p, (n + m))))) \end{array}}}
\vnbprop{qe-psum-type}{\vnbstmt{\begin{array}{@{}l@{}} \forall m \in \mathbb{N}, p.; \forall a \in (\mathbb{Z} \to \mathbb{Z}), n \in \mathbb{Z}. \\ \quad \operatorname{zz\text{-}prime}(p) \Rightarrow (\operatorname{qp\text{-}psum}(p, a, n, m) \in \operatorname{carr}(\operatorname{qp\text{-}field}(p))) \end{array}}}
\vnbprop{qe-abs-swap}{\vnbstmt{\begin{array}{@{}l@{}} \forall p.; \forall u, v \in \operatorname{carr}(\operatorname{qp\text{-}field}(p)). \\ \quad \operatorname{zz\text{-}prime}(p) \Rightarrow  \\ \quad \quad \operatorname{qp\text{-}abs}(p)(\operatorname{add}(\operatorname{qp\text{-}field}(p))(u, \operatorname{neg}(\operatorname{qp\text{-}field}(p))(v))) \\ \quad \quad =\; \operatorname{qp\text{-}abs}(p)(\operatorname{add}(\operatorname{qp\text{-}field}(p))(v, \operatorname{neg}(\operatorname{qp\text{-}field}(p))(u))) \end{array}}}
\vnbprop{qe-ultra-lt}{\vnbstmt{\begin{array}{@{}l@{}} \forall p.; \forall u, v, w \in \operatorname{carr}(\operatorname{qp\text{-}field}(p)), r \in \mathbb{R}. \\ \quad \operatorname{zz\text{-}prime}(p) \wedge \\ \quad \operatorname{qp\text{-}abs}(p)(\operatorname{add}(\operatorname{qp\text{-}field}(p))(u, \operatorname{neg}(\operatorname{qp\text{-}field}(p))(v))) < r \wedge \\ \quad \operatorname{qp\text{-}abs}(p)(\operatorname{add}(\operatorname{qp\text{-}field}(p))(v, \operatorname{neg}(\operatorname{qp\text{-}field}(p))(w))) < r \Rightarrow \\ \quad \quad \operatorname{qp\text{-}abs}(p)(\operatorname{add}(\operatorname{qp\text{-}field}(p))(u, \operatorname{neg}(\operatorname{qp\text{-}field}(p))(w))) < r \end{array}}}
\vnbprop{qe-ultra-le}{\vnbstmt{\begin{array}{@{}l@{}} \forall p.; \forall u, v, w \in \operatorname{carr}(\operatorname{qp\text{-}field}(p)), r \in \mathbb{R}. \\ \quad \operatorname{zz\text{-}prime}(p) \wedge \\ \quad \operatorname{qp\text{-}abs}(p)(\operatorname{add}(\operatorname{qp\text{-}field}(p))(u, \operatorname{neg}(\operatorname{qp\text{-}field}(p))(v))) \leq r \wedge \\ \quad \operatorname{qp\text{-}abs}(p)(\operatorname{add}(\operatorname{qp\text{-}field}(p))(v, \operatorname{neg}(\operatorname{qp\text{-}field}(p))(w))) \leq r \Rightarrow \\ \quad \quad \operatorname{qp\text{-}abs}(p)(\operatorname{add}(\operatorname{qp\text{-}field}(p))(u, \operatorname{neg}(\operatorname{qp\text{-}field}(p))(w))) \leq r \end{array}}}
\vnbprop{qe-embed-diff}{\vnbstmt{\begin{array}{@{}l@{}} \forall p.; \forall u, v \in \mathbb{Q}. \\ \quad \operatorname{zz\text{-}prime}(p) \Rightarrow  \\ \quad \quad \operatorname{qp\text{-}abs}(p)(\operatorname{add}(\operatorname{qp\text{-}field}(p))(\operatorname{embed}(\operatorname{nf\text{-}metric\text{-}space}(\operatorname{qq\text{-}padic}(p)))(u), \\ \quad \quad \quad \quad \operatorname{neg}(\operatorname{qp\text{-}field}(p))(\operatorname{embed}(\operatorname{nf\text{-}metric\text{-}space}(\operatorname{qq\text{-}padic}(p)))(v)))) \\ \quad \quad =\; \operatorname{padic\text{-}abs}(p)((u + -v)) \end{array}}}
\vnbprop{qe-digit-unique}{\vnbstmt{\begin{array}{@{}l@{}} \forall p.; \forall d, e \in \mathbb{Z}. \\ \quad (\operatorname{zz\text{-}prime}(p) \wedge ((0 \leq d) \wedge ((d < p) \wedge ((0 \leq e) \wedge ((e < p) \wedge (\operatorname{padic\text{-}abs}(p)((d + -e)) < 1)))))) \Rightarrow  \\ \quad \quad d = e \end{array}}}
\vnbprop{qe-digit-unique-qp}{\vnbstmt{\begin{array}{@{}l@{}} \forall p.; \forall y \in \operatorname{carr}(\operatorname{qp\text{-}field}(p)), d, e \in \mathbb{Z}. \\ \quad \operatorname{zz\text{-}prime}(p) \wedge \\ \quad 0 \leq d \wedge \\ \quad d < p \wedge \\ \quad 0 \leq e \wedge \\ \quad e < p \wedge \\ \quad \operatorname{qp\text{-}abs}(p)(\operatorname{add}(\operatorname{qp\text{-}field}(p))(y, \operatorname{neg}(\operatorname{qp\text{-}field}(p))(\operatorname{embed}(\operatorname{nf\text{-}metric\text{-}space}(\operatorname{qq\text{-}padic}(p)))(d)))) < 1 \wedge \\ \quad \operatorname{qp\text{-}abs}(p)(\operatorname{add}(\operatorname{qp\text{-}field}(p))(y, \operatorname{neg}(\operatorname{qp\text{-}field}(p))(\operatorname{embed}(\operatorname{nf\text{-}metric\text{-}space}(\operatorname{qq\text{-}padic}(p)))(e)))) < 1 \Rightarrow \\ \quad \quad d = e \end{array}}}
\vnbprop{qe-digit-exists}{\vnbstmt{\begin{array}{@{}l@{}} \forall p.; \forall y \in \operatorname{zp}(p). \\ \quad \operatorname{zz\text{-}prime}(p) \Rightarrow  \\ \quad \quad \exists g \in \mathbb{Z}. \\ \quad \quad \quad 0 \leq g \wedge \\ \quad \quad \quad g < p \wedge \\ \quad \quad \quad \operatorname{qp\text{-}abs}(p)(\operatorname{add}(\operatorname{qp\text{-}field}(p))(y, \operatorname{neg}(\operatorname{qp\text{-}field}(p))(\operatorname{embed}(\operatorname{nf\text{-}metric\text{-}space}(\operatorname{qq\text{-}padic}(p)))(g)))) < 1 \end{array}}}
\vnbprop{qe-digit-def}{\vnbstmt{\begin{array}{@{}l@{}} \forall p.; \forall y \in \operatorname{zp}(p). \\ \quad \operatorname{zz\text{-}prime}(p) \Rightarrow  \\ \quad \quad \operatorname{qp\text{-}digit}(p, y) \in \mathbb{Z} \wedge \\ \quad \quad 0 \leq \operatorname{qp\text{-}digit}(p, y) \wedge \\ \quad \quad \operatorname{qp\text{-}digit}(p, y) < p \wedge \\ \quad \quad \operatorname{qp\text{-}abs}(p)(\operatorname{add}(\operatorname{qp\text{-}field}(p))(y, \operatorname{neg}(\operatorname{qp\text{-}field}(p))(\operatorname{embed}(\operatorname{nf\text{-}metric\text{-}space}(\operatorname{qq\text{-}padic}(p)))(\operatorname{qp\text{-}digit}(p, y))))) \\ \quad \quad <\; 1 \end{array}}}
\vnbprop{qe-digit-step}{\vnbstmt{\begin{array}{@{}l@{}} \forall p.; \forall y \in \operatorname{zp}(p). \\ \quad \operatorname{zz\text{-}prime}(p) \Rightarrow  \\ \quad \quad \exists z \in \operatorname{zp}(p). \\ \quad \quad \quad y \\ \quad \quad \quad =\; \operatorname{add}(\operatorname{qp\text{-}field}(p))(\operatorname{embed}(\operatorname{nf\text{-}metric\text{-}space}(\operatorname{qq\text{-}padic}(p)))(\operatorname{qp\text{-}digit}(p, y)), \\ \quad \quad \quad \quad \operatorname{mul}(\operatorname{qp\text{-}field}(p))(\operatorname{embed}(\operatorname{nf\text{-}metric\text{-}space}(\operatorname{qq\text{-}padic}(p)))(p), z)) \end{array}}}
\vnbprop{qe-digit-sequence}{\vnbstmt{\begin{array}{@{}l@{}} \forall p.; \forall y \in \operatorname{zp}(p). \\ \quad \operatorname{zz\text{-}prime}(p) \Rightarrow  \\ \quad \quad \exists f \in (\mathbb{N} \to \operatorname{zp}(p)). \\ \quad \quad \quad f(0) = y \wedge \\ \quad \quad \quad \forall k \in \mathbb{N}. \\ \quad \quad \quad \quad f(k) \\ \quad \quad \quad \quad =\; \operatorname{add}(\operatorname{qp\text{-}field}(p))(\operatorname{embed}(\operatorname{nf\text{-}metric\text{-}space}(\operatorname{qq\text{-}padic}(p)))(\operatorname{qp\text{-}digit}(p, f(k))), \\ \quad \quad \quad \quad \quad \operatorname{mul}(\operatorname{qp\text{-}field}(p))(\operatorname{embed}(\operatorname{nf\text{-}metric\text{-}space}(\operatorname{qq\text{-}padic}(p)))(p), f((k + 1)))) \end{array}}}
\vnbprop{qe-rp-le}{\vnbstmt{\forall s, t \in \mathbb{R}.\; ((s \leq t) \Rightarrow ({2}^{s} \leq {2}^{t}))}}
\vnbprop{qe-rp-small}{\vnbstmt{\begin{array}{@{}l@{}} \forall s.; \forall n \in \mathbb{Z}. \\ \quad \operatorname{pos\text{-}rr}(s) \Rightarrow  \\ \quad \quad \exists q \in \mathbb{N}.\; \forall m \in \mathbb{N}.\; ((q \leq m) \Rightarrow ({2}^{-(n + m)} < s)) \end{array}}}
\vnbprop{qe-expand-ind}{\vnbstmt{\begin{array}{@{}l@{}} \forall m \in \mathbb{N}, p.; \forall f \in (\mathbb{N} \to \operatorname{zp}(p)), a \in (\mathbb{Z} \to \mathbb{Z}), n \in \mathbb{Z}. \\ \quad \operatorname{zz\text{-}prime}(p) \wedge \\ \quad \forall k \in \mathbb{N}. \\ \quad \quad f(k) \\ \quad \quad =\; \operatorname{add}(\operatorname{qp\text{-}field}(p))(\operatorname{embed}(\operatorname{nf\text{-}metric\text{-}space}(\operatorname{qq\text{-}padic}(p)))(\operatorname{qp\text{-}digit}(p, f(k))), \\ \quad \quad \quad \operatorname{mul}(\operatorname{qp\text{-}field}(p))(\operatorname{embed}(\operatorname{nf\text{-}metric\text{-}space}(\operatorname{qq\text{-}padic}(p)))(p), f((k + 1)))) \wedge \\ \quad \forall j \in \mathbb{N}.\; (a((n + j)) = \operatorname{qp\text{-}digit}(p, f(j))) \Rightarrow \\ \quad \quad \operatorname{mul}(\operatorname{qp\text{-}field}(p))(\operatorname{qp\text{-}ppow}(p, n), f(0)) \\ \quad \quad =\; \operatorname{add}(\operatorname{qp\text{-}field}(p))(\operatorname{qp\text{-}psum}(p, a, n, m), \operatorname{mul}(\operatorname{qp\text{-}field}(p))(\operatorname{qp\text{-}ppow}(p, (n + m)), f(m))) \end{array}}}
\vnbprop{qe-expansion-exists}{\vnbstmt{\begin{array}{@{}l@{}} \forall p.; \forall x \in \operatorname{carr}(\operatorname{qp\text{-}field}(p)), n \in \mathbb{Z}. \\ \quad (\operatorname{zz\text{-}prime}(p) \wedge (\operatorname{mul}(\operatorname{qp\text{-}field}(p))(\operatorname{qp\text{-}ppow}(p, -n), x) \in \operatorname{zp}(p))) \Rightarrow  \\ \quad \quad \exists a.\; \operatorname{is\text{-}qp\text{-}expansion}(p, a, n, x) \end{array}}}
\vnbprop{qe-tail-ind}{\vnbstmt{\begin{array}{@{}l@{}} \forall l \in \mathbb{N}, p.; \forall a \in (\mathbb{Z} \to \mathbb{Z}), n \in \mathbb{Z}, m \in \mathbb{N}. \\ \quad \operatorname{zz\text{-}prime}(p) \Rightarrow  \\ \quad \quad \operatorname{qp\text{-}abs}(p)(\operatorname{add}(\operatorname{qp\text{-}field}(p))(\operatorname{qp\text{-}psum}(p, a, n, (m + l)), \operatorname{neg}(\operatorname{qp\text{-}field}(p))(\operatorname{qp\text{-}psum}(p, a, n, m)))) \\ \quad \quad \leq\; {2}^{-(n + m)} \end{array}}}
\vnbprop{qe-tail}{\vnbstmt{\begin{array}{@{}l@{}} \forall p.; \forall a \in (\mathbb{Z} \to \mathbb{Z}), n \in \mathbb{Z}, x.; \forall m \in \mathbb{N}. \\ \quad (\operatorname{zz\text{-}prime}(p) \wedge \operatorname{converges\text{-}to}(\operatorname{qp}(p), (\lambda \mathit{qtm} \in \mathbb{N}.\; \operatorname{qp\text{-}psum}(p, a, n, \mathit{qtm})), x)) \Rightarrow  \\ \quad \quad \operatorname{qp\text{-}abs}(p)(\operatorname{add}(\operatorname{qp\text{-}field}(p))(x, \operatorname{neg}(\operatorname{qp\text{-}field}(p))(\operatorname{qp\text{-}psum}(p, a, n, m)))) \\ \quad \quad \leq\; {2}^{-(n + m)} \end{array}}}
\vnbprop{qe-pad-psum}{\vnbstmt{\begin{array}{@{}l@{}} \forall m \in \mathbb{N}, p.; \forall a \in (\mathbb{Z} \to \mathbb{Z}), l \in \mathbb{Z}, d \in \mathbb{N}, n \in \mathbb{Z}. \\ \quad (\operatorname{zz\text{-}prime}(p) \wedge (((l + d) = n) \wedge \forall k \in \mathbb{Z}.\; ((k < n) \Rightarrow (a(k) = 0)))) \Rightarrow  \\ \quad \quad \operatorname{qp\text{-}psum}(p, a, l, (d + m)) = \operatorname{qp\text{-}psum}(p, a, n, m) \end{array}}}
\vnbprop{qe-pad}{\vnbstmt{\begin{array}{@{}l@{}} \forall p, a.; \forall n \in \mathbb{Z}, x.; \forall l \in \mathbb{Z}. \\ \quad (\operatorname{zz\text{-}prime}(p) \wedge (\operatorname{is\text{-}qp\text{-}expansion}(p, a, n, x) \wedge (l \leq n))) \Rightarrow  \\ \quad \quad \operatorname{is\text{-}qp\text{-}expansion}(p, a, l, x) \end{array}}}
\vnbprop{qe-same-start}{\vnbstmt{\begin{array}{@{}l@{}} \forall c \in \mathbb{N}, p, a, b.; \forall n \in \mathbb{Z}, x.; \forall j \in \mathbb{N}. \\ \quad (\operatorname{zz\text{-}prime}(p) \wedge (\operatorname{is\text{-}qp\text{-}expansion}(p, a, n, x) \wedge (\operatorname{is\text{-}qp\text{-}expansion}(p, b, n, x) \wedge (j < c)))) \Rightarrow  \\ \quad \quad a((n + j)) = b((n + j)) \end{array}}}
\vnbprop{qe-unique-le}{\vnbstmt{\begin{array}{@{}l@{}} \forall p, a.; \forall n \in \mathbb{Z}, b.; \forall m \in \mathbb{Z}, x.; \forall k \in \mathbb{Z}. \\ \quad (\operatorname{zz\text{-}prime}(p) \wedge (\operatorname{is\text{-}qp\text{-}expansion}(p, a, n, x) \wedge (\operatorname{is\text{-}qp\text{-}expansion}(p, b, m, x) \wedge (n \leq m)))) \Rightarrow  \\ \quad \quad a(k) = b(k) \end{array}}}
\vnbprop{ccr-4-19}{\vnbstmt{\begin{array}{@{}l@{}} \forall p.; \forall x \in \operatorname{carr}(\operatorname{qp\text{-}field}(p)). \\ \quad \operatorname{zz\text{-}prime}(p) \Rightarrow  \\ \quad \quad \exists a.; \exists n \in \mathbb{Z}.\; \operatorname{is\text{-}qp\text{-}expansion}(p, a, n, x) \wedge \\ \quad \quad \forall a.; \forall n \in \mathbb{Z}, b.; \forall m, k \in \mathbb{Z}. \\ \quad \quad \quad (\operatorname{is\text{-}qp\text{-}expansion}(p, a, n, x) \wedge \operatorname{is\text{-}qp\text{-}expansion}(p, b, m, x)) \Rightarrow (a(k) = b(k)) \end{array}}}
\vnbprop{ccr-1-30}{\vnbstmt{\begin{array}{@{}l@{}} \forall p, u. \\ \quad \operatorname{zz\text{-}prime}(p) \wedge \\ \quad \operatorname{is\text{-}add\text{-}subgroup}(\operatorname{qp\text{-}field}(p), u) \wedge \\ \quad u \subseteq \operatorname{zp}(p) \wedge \\ \quad \operatorname{is\text{-}open}(\operatorname{subspace\text{-}ms}(\operatorname{nf\text{-}metric\text{-}space}(\operatorname{qp\text{-}field}(p)), \operatorname{zp}(p)), u) \Rightarrow \\ \quad \quad \operatorname{is\text{-}nv\text{-}module}(\operatorname{qp\text{-}field}(p), \operatorname{zp}(p), u) \wedge \\ \quad \quad \exists g \in \operatorname{zp}(p).\; (u = \operatorname{nv\text{-}multiples}(\operatorname{qp\text{-}field}(p), \operatorname{zp}(p), g)) \end{array}}}
\vnbprop{qe-term-bound}{\vnbstmt{\begin{array}{@{}l@{}} \forall p.; \forall a \in (\mathbb{Z} \to \mathbb{Z}), n \in \mathbb{Z}, j \in \mathbb{N}. \\ \quad \operatorname{zz\text{-}prime}(p) \Rightarrow  \\ \quad \quad \operatorname{qp\text{-}abs}(p)(\operatorname{mul}(\operatorname{qp\text{-}field}(p))(\operatorname{embed}(\operatorname{nf\text{-}metric\text{-}space}(\operatorname{qq\text{-}padic}(p)))(a((n + j))), \operatorname{qp\text{-}ppow}(p, (n + j)))) \\ \quad \quad \in\; \mathbb{R} \wedge \\ \quad \quad 0 \\ \quad \quad \leq\; \operatorname{qp\text{-}abs}(p)(\operatorname{mul}(\operatorname{qp\text{-}field}(p))(\operatorname{embed}(\operatorname{nf\text{-}metric\text{-}space}(\operatorname{qq\text{-}padic}(p)))(a((n + j))), \operatorname{qp\text{-}ppow}(p, (n + j)))) \wedge \\ \quad \quad \operatorname{qp\text{-}abs}(p)(\operatorname{mul}(\operatorname{qp\text{-}field}(p))(\operatorname{embed}(\operatorname{nf\text{-}metric\text{-}space}(\operatorname{qq\text{-}padic}(p)))(a((n + j))), \operatorname{qp\text{-}ppow}(p, (n + j)))) \\ \quad \quad \leq\; {2}^{-(n + j)} \end{array}}}
\vnbprop{qe-abs-convergent}{\vnbstmt{\begin{array}{@{}l@{}} \forall p.; \forall a \in (\mathbb{Z} \to \mathbb{Z}), n \in \mathbb{Z}. \\ \quad \operatorname{zz\text{-}prime}(p) \Rightarrow  \\ \quad \quad \operatorname{series\text{-}converges}((\lambda j \in \mathbb{N}.\; \\ \quad \quad \quad \quad \operatorname{qp\text{-}abs}(p)(\operatorname{mul}(\operatorname{qp\text{-}field}(p))(\operatorname{embed}(\operatorname{nf\text{-}metric\text{-}space}(\operatorname{qq\text{-}padic}(p)))(a((n + j))), \operatorname{qp\text{-}ppow}(p, (n + j)))))) \end{array}}}
\vnbprop{qx-ppow-rat}{\vnbstmt{\begin{array}{@{}l@{}} \forall p.; \forall k \in \mathbb{Z}. \\ \quad \operatorname{zz\text{-}prime}(p) \Rightarrow  \\ \quad \quad (\operatorname{power}(p, \operatorname{max}(k, 0)) \cdot \operatorname{recip}(\operatorname{power}(p, \operatorname{max}(-k, 0)))) \in \mathbb{Q} \wedge \\ \quad \quad \operatorname{qp\text{-}ppow}(p, k) = \operatorname{embed}(\operatorname{nf\text{-}metric\text{-}space}(\operatorname{qq\text{-}padic}(p)))((\operatorname{power}(p, \operatorname{max}(k, 0)) \cdot \operatorname{recip}(\operatorname{power}(p, \operatorname{max}(-k, 0))))) \end{array}}}
\vnbprop{qx-approx}{\vnbstmt{\begin{array}{@{}l@{}} \forall p.; \forall x \in \operatorname{carr}(\operatorname{qp\text{-}field}(p)). \\ \quad \operatorname{zz\text{-}prime}(p) \Rightarrow  \\ \quad \quad \exists k, z \in \mathbb{Z}. \\ \quad \quad \quad \operatorname{qp\text{-}abs}(p)(\operatorname{add}(\operatorname{qp\text{-}field}(p))(x, \\ \quad \quad \quad \quad \quad \operatorname{neg}(\operatorname{qp\text{-}field}(p))(\operatorname{mul}(\operatorname{qp\text{-}field}(p))(\operatorname{qp\text{-}ppow}(p, k), \operatorname{embed}(\operatorname{nf\text{-}metric\text{-}space}(\operatorname{qq\text{-}padic}(p)))(z))))) \\ \quad \quad \quad \leq\; 1 \end{array}}}
\vnbprop{ccr-5-17}{\vnbstmt{\begin{array}{@{}l@{}} \forall p.; \forall x \in \operatorname{carr}(\operatorname{qp\text{-}field}(p)). \\ \quad \operatorname{zz\text{-}prime}(p) \Rightarrow  \\ \quad \quad \exists k, z \in \mathbb{Z}. \\ \quad \quad \quad \operatorname{qp\text{-}abs}(p)(\operatorname{add}(\operatorname{qp\text{-}field}(p))(x, \\ \quad \quad \quad \quad \quad \operatorname{neg}(\operatorname{qp\text{-}field}(p))(\operatorname{mul}(\operatorname{qp\text{-}field}(p))(\operatorname{qp\text{-}ppow}(p, k), \operatorname{embed}(\operatorname{nf\text{-}metric\text{-}space}(\operatorname{qq\text{-}padic}(p)))(z))))) \\ \quad \quad \quad \leq\; 1 \wedge \\ \quad \quad \forall k, z, a, b \in \mathbb{Z}. \\ \quad \quad \quad \operatorname{qp\text{-}abs}(p)(\operatorname{add}(\operatorname{qp\text{-}field}(p))(x, \\ \quad \quad \quad \quad \quad \operatorname{neg}(\operatorname{qp\text{-}field}(p))(\operatorname{mul}(\operatorname{qp\text{-}field}(p))(\operatorname{qp\text{-}ppow}(p, k), \operatorname{embed}(\operatorname{nf\text{-}metric\text{-}space}(\operatorname{qq\text{-}padic}(p)))(z))))) \\ \quad \quad \quad \leq\; 1 \wedge \\ \quad \quad \quad \operatorname{qp\text{-}abs}(p)(\operatorname{add}(\operatorname{qp\text{-}field}(p))(x, \\ \quad \quad \quad \quad \quad \operatorname{neg}(\operatorname{qp\text{-}field}(p))(\operatorname{mul}(\operatorname{qp\text{-}field}(p))(\operatorname{qp\text{-}ppow}(p, a), \operatorname{embed}(\operatorname{nf\text{-}metric\text{-}space}(\operatorname{qq\text{-}padic}(p)))(b))))) \\ \quad \quad \quad \leq\; 1 \Rightarrow \\ \quad \quad \quad \quad \exists q \in \mathbb{Z}. \\ \quad \quad \quad \quad \quad \operatorname{add}(\operatorname{qp\text{-}field}(p))(\operatorname{mul}(\operatorname{qp\text{-}field}(p))(\operatorname{qp\text{-}ppow}(p, k), \operatorname{embed}(\operatorname{nf\text{-}metric\text{-}space}(\operatorname{qq\text{-}padic}(p)))(z)), \\ \quad \quad \quad \quad \quad \quad \operatorname{neg}(\operatorname{qp\text{-}field}(p))(\operatorname{mul}(\operatorname{qp\text{-}field}(p))(\operatorname{qp\text{-}ppow}(p, a), \operatorname{embed}(\operatorname{nf\text{-}metric\text{-}space}(\operatorname{qq\text{-}padic}(p)))(b)))) \\ \quad \quad \quad \quad \quad =\; \operatorname{embed}(\operatorname{nf\text{-}metric\text{-}space}(\operatorname{qq\text{-}padic}(p)))(q) \end{array}}}
\vnbprop{qx-rat-approx}{\vnbstmt{\begin{array}{@{}l@{}} \forall p.; \forall x \in \operatorname{carr}(\operatorname{qp\text{-}field}(p)). \\ \quad \operatorname{zz\text{-}prime}(p) \Rightarrow  \\ \quad \quad \exists u \in \mathbb{Q}. \\ \quad \quad \quad \operatorname{qp\text{-}abs}(p)(\operatorname{add}(\operatorname{qp\text{-}field}(p))(x, \operatorname{neg}(\operatorname{qp\text{-}field}(p))(\operatorname{embed}(\operatorname{nf\text{-}metric\text{-}space}(\operatorname{qq\text{-}padic}(p)))(u)))) \\ \quad \quad \quad \leq\; 1 \end{array}}}
\vnbprop{qx-lam-local}{\vnbstmt{\begin{array}{@{}l@{}} \forall p.; \forall x \in \operatorname{carr}(\operatorname{qp\text{-}field}(p)), u, v \in \mathbb{Q}. \\ \quad \operatorname{zz\text{-}prime}(p) \wedge \\ \quad \operatorname{qp\text{-}abs}(p)(\operatorname{add}(\operatorname{qp\text{-}field}(p))(x, \operatorname{neg}(\operatorname{qp\text{-}field}(p))(\operatorname{embed}(\operatorname{nf\text{-}metric\text{-}space}(\operatorname{qq\text{-}padic}(p)))(u)))) \\ \quad \leq\; 1 \wedge \\ \quad \operatorname{qp\text{-}abs}(p)(\operatorname{add}(\operatorname{qp\text{-}field}(p))(x, \operatorname{neg}(\operatorname{qp\text{-}field}(p))(\operatorname{embed}(\operatorname{nf\text{-}metric\text{-}space}(\operatorname{qq\text{-}padic}(p)))(v)))) \\ \quad \leq\; 1 \Rightarrow \\ \quad \quad \operatorname{padic\text{-}char}(p)(u) = \operatorname{padic\text{-}char}(p)(v) \end{array}}}
\vnbprop{qx-lam-value}{\vnbstmt{\begin{array}{@{}l@{}} \forall p.; \forall x \in \operatorname{carr}(\operatorname{qp\text{-}field}(p)), u \in \mathbb{Q}. \\ \quad \operatorname{zz\text{-}prime}(p) \wedge \\ \quad \operatorname{qp\text{-}abs}(p)(\operatorname{add}(\operatorname{qp\text{-}field}(p))(x, \operatorname{neg}(\operatorname{qp\text{-}field}(p))(\operatorname{embed}(\operatorname{nf\text{-}metric\text{-}space}(\operatorname{qq\text{-}padic}(p)))(u)))) \\ \quad \leq\; 1 \Rightarrow \\ \quad \quad \operatorname{qp\text{-}lambda\text{-}at}(p, x) = \operatorname{padic\text{-}char}(p)(u) \end{array}}}
\vnbprop{qx-lam-app}{\vnbstmt{\begin{array}{@{}l@{}} \forall p.; \forall x \in \operatorname{carr}(\operatorname{qp\text{-}field}(p)). \\ \quad \operatorname{zz\text{-}prime}(p) \Rightarrow (\operatorname{qp\text{-}lambda}(p)(x) \simeq \operatorname{qp\text{-}lambda\text{-}at}(p, x)) \end{array}}}
\vnbprop{qx-lam-type}{\vnbstmt{\begin{array}{@{}l@{}} \forall p. \\ \quad \operatorname{zz\text{-}prime}(p) \Rightarrow (\operatorname{qp\text{-}lambda}(p) \in (\operatorname{carr}(\operatorname{qp\text{-}field}(p)) \to \operatorname{unit\text{-}circle})) \end{array}}}
\vnbprop{qx-lam-type-cc}{\vnbstmt{\forall p.\; (\operatorname{zz\text{-}prime}(p) \Rightarrow (\operatorname{qp\text{-}lambda}(p) \in (\operatorname{carr}(\operatorname{qp\text{-}field}(p)) \to \mathbb{C})))}}
\vnbprop{ccr-5-20}{\vnbstmt{\begin{array}{@{}l@{}} \forall p. \\ \quad \operatorname{zz\text{-}prime}(p) \Rightarrow  \\ \quad \quad \operatorname{is\text{-}add\text{-}character}(\operatorname{qp\text{-}field}(p), \operatorname{qp\text{-}lambda}(p)) \wedge \\ \quad \quad \forall x \in \operatorname{carr}(\operatorname{qp\text{-}field}(p)).\; ((\operatorname{qp\text{-}lambda}(p)(x) = 1) \Leftrightarrow (x \in \operatorname{zp}(p))) \end{array}}}
\vnbprop{ccr-4-18-qp}{\vnbstmt{\begin{array}{@{}l@{}} \forall p, a.; \forall n \in \mathbb{Z}, x. \\ \quad (\operatorname{zz\text{-}prime}(p) \wedge (\operatorname{is\text{-}qp\text{-}expansion}(p, a, n, x) \wedge \neg (a(n) = 0))) \Rightarrow  \\ \quad \quad \operatorname{qp\text{-}abs}(p)(x) = {2}^{-n} \end{array}}}

\vnbbookfile{theorem-library/ccr1109a-late}{theorem-library/ccr1109a-late}
\noindent{\small ccr1109a-late.scm -- wave of 2026-11-09/10 (agent CCR-1109-A), wired 2026-11-10 from scratchpad/ccr1109a/ccr1109a-late.scm. Library: complex-analysis.}\par
\vnbprop{metric-dense-of-near}{\vnbstmt{\begin{array}{@{}l@{}} \forall s, a. \\ \quad \operatorname{is\text{-}metric\text{-}space}(s) \wedge \\ \quad a \subseteq \operatorname{pts}(s) \wedge \\ \quad \forall x \in \operatorname{pts}(s), b. \\ \quad \quad \operatorname{pos\text{-}rr}(b) \Rightarrow \exists c \in a.\; (\operatorname{dist}(s)(x, c) < b) \Rightarrow \\ \quad \quad \operatorname{is\text{-}dense}(s, a) \end{array}}}
\vnbprop{ccr-1-20-dense}{\vnbstmt{\begin{array}{@{}l@{}} \forall p. \\ \quad \operatorname{zz\text{-}prime}(p) \Rightarrow  \\ \quad \quad \operatorname{is\text{-}dense}(\operatorname{subspace\text{-}ms}(\operatorname{nf\text{-}metric\text{-}space}(\operatorname{qq\text{-}padic}(p)), \operatorname{nv\text{-}integers}(\operatorname{qq\text{-}padic}(p))), \mathbb{Z}) \end{array}}}
\vnbprop{ccr-2-23-class}{\vnbstmt{\begin{array}{@{}l@{}} \forall k, a, i.; \forall x \in \operatorname{nv\text{-}integers}(k). \\ \quad \operatorname{is\text{-}nv\text{-}domain}(k, a) \wedge \\ \quad i \in (\operatorname{nv\text{-}integers}(k) \to a) \wedge \\ \quad \forall x \in \operatorname{nv\text{-}integers}(k).\; (\operatorname{fnrm}(k)(\operatorname{add}(k)(x, \operatorname{neg}(k)(i(x)))) < 1) \wedge \\ \quad \forall x \in a.\; (i(x) = x) \Rightarrow \\ \quad \quad \operatorname{class}(\operatorname{ringoid\text{-}setoid}(\operatorname{nv\text{-}res\text{-}ringoid}(k, \operatorname{nv\text{-}integers}(k))), i(x)) \\ \quad \quad =\; \operatorname{class}(\operatorname{ringoid\text{-}setoid}(\operatorname{nv\text{-}res\text{-}ringoid}(k, \operatorname{nv\text{-}integers}(k))), x) \end{array}}}

\vnbbookfile{theorem-library/ccr1109a-nv}{theorem-library/ccr1109a-nv}
\noindent{\small ccr1109a-nv.scm -- wave of 2026-11-09/10 (agent CCR-1109-A), wired 2026-11-10 from scratchpad/ccr1109a/ccr1109a-nv.scm. Library: complex-analysis.}\par
\vnbprop{rr-recip-pair-lt-one-iff}{\vnbstmt{\begin{array}{@{}l@{}} \forall u, w \in \mathbb{R}. \\ \quad ((0 \leq u) \wedge ((0 \leq w) \wedge ((u \cdot w) = 1))) \Rightarrow ((u < 1) \Leftrightarrow (1 < w)) \end{array}}}
\vnbprop{nf-inverse-norm-lt-one-iff}{\vnbstmt{\begin{array}{@{}l@{}} \forall k.; \forall a, b \in \operatorname{carr}(k). \\ \quad (\operatorname{is\text{-}normed\text{-}field}(k) \wedge (\operatorname{mul}(k)(a, b) = \operatorname{one}(k))) \Rightarrow  \\ \quad \quad (\operatorname{fnrm}(k)(b) < 1) \Leftrightarrow (1 < \operatorname{fnrm}(k)(a)) \end{array}}}
\vnbprop{ccr-1-46-ii-iii}{\vnbstmt{\begin{array}{@{}l@{}} \forall a, b. \\ \quad \operatorname{is\text{-}nv\text{-}field}(a) \wedge \\ \quad \operatorname{is\text{-}nv\text{-}field}(b) \wedge \\ \quad (\operatorname{carr}(b) \simeq \operatorname{carr}(a)) \wedge \\ \quad \operatorname{add}(b) = \operatorname{add}(a) \wedge \\ \quad \operatorname{neg}(b) = \operatorname{neg}(a) \wedge \\ \quad \operatorname{mul}(b) = \operatorname{mul}(a) \wedge \\ \quad \operatorname{zero}(b) = \operatorname{zero}(a) \wedge \\ \quad \operatorname{one}(b) = \operatorname{one}(a) \Rightarrow \\ \quad \quad \forall c \in \operatorname{carr}(a).\; ((1 < \operatorname{fnrm}(a)(c)) \Leftrightarrow (1 < \operatorname{fnrm}(b)(c))) \\ \quad \quad \Leftrightarrow\; \forall c \in \operatorname{carr}(a).\; ((\operatorname{fnrm}(a)(c) < 1) \Leftrightarrow (\operatorname{fnrm}(b)(c) < 1)) \end{array}}}
\vnbprop{ccr-1-46-iii-iv}{\vnbstmt{\begin{array}{@{}l@{}} \forall a, b. \\ \quad \operatorname{is\text{-}nv\text{-}field}(a) \wedge \\ \quad \operatorname{is\text{-}nv\text{-}field}(b) \wedge \\ \quad (\operatorname{carr}(b) \simeq \operatorname{carr}(a)) \wedge \\ \quad \operatorname{add}(b) = \operatorname{add}(a) \wedge \\ \quad \operatorname{neg}(b) = \operatorname{neg}(a) \wedge \\ \quad \operatorname{mul}(b) = \operatorname{mul}(a) \wedge \\ \quad \operatorname{zero}(b) = \operatorname{zero}(a) \wedge \\ \quad \operatorname{one}(b) = \operatorname{one}(a) \Rightarrow \\ \quad \quad \forall c \in \operatorname{carr}(a).\; ((\operatorname{fnrm}(a)(c) < 1) \Leftrightarrow (\operatorname{fnrm}(b)(c) < 1)) \\ \quad \quad \Leftrightarrow\; \forall c \in \operatorname{carr}(a).\; ((1 \leq \operatorname{fnrm}(a)(c)) \Leftrightarrow (1 \leq \operatorname{fnrm}(b)(c))) \end{array}}}
\vnbprop{ccr-1-46-ii-v}{\vnbstmt{\begin{array}{@{}l@{}} \forall a, b. \\ \quad \operatorname{is\text{-}nv\text{-}field}(a) \wedge \\ \quad \operatorname{is\text{-}nv\text{-}field}(b) \wedge \\ \quad (\operatorname{carr}(b) \simeq \operatorname{carr}(a)) \wedge \\ \quad \operatorname{add}(b) = \operatorname{add}(a) \wedge \\ \quad \operatorname{neg}(b) = \operatorname{neg}(a) \wedge \\ \quad \operatorname{mul}(b) = \operatorname{mul}(a) \wedge \\ \quad \operatorname{zero}(b) = \operatorname{zero}(a) \wedge \\ \quad \operatorname{one}(b) = \operatorname{one}(a) \Rightarrow \\ \quad \quad \forall c \in \operatorname{carr}(a).\; ((1 < \operatorname{fnrm}(a)(c)) \Leftrightarrow (1 < \operatorname{fnrm}(b)(c))) \\ \quad \quad \Leftrightarrow\; \forall c \in \operatorname{carr}(a).\; ((\operatorname{fnrm}(a)(c) \leq 1) \Leftrightarrow (\operatorname{fnrm}(b)(c) \leq 1)) \end{array}}}
\vnbprop{ccr-1-46-vi-v}{\vnbstmt{\begin{array}{@{}l@{}} \forall a, b.; \forall c \in \operatorname{carr}(a). \\ \quad \operatorname{is\text{-}nv\text{-}field}(a) \wedge \\ \quad \operatorname{is\text{-}nv\text{-}field}(b) \wedge \\ \quad (\operatorname{carr}(b) \simeq \operatorname{carr}(a)) \wedge \\ \quad \operatorname{add}(b) = \operatorname{add}(a) \wedge \\ \quad \operatorname{neg}(b) = \operatorname{neg}(a) \wedge \\ \quad \operatorname{mul}(b) = \operatorname{mul}(a) \wedge \\ \quad \operatorname{zero}(b) = \operatorname{zero}(a) \wedge \\ \quad \operatorname{one}(b) = \operatorname{one}(a) \wedge \\ \quad \exists l \in \mathbb{R}. \\ \quad \quad (0 < l) \wedge \forall c \in \operatorname{carr}(a).\; (\operatorname{fnrm}(a)(c) = {\operatorname{fnrm}(b)(c)}^{l}) \Rightarrow \\ \quad \quad (\operatorname{fnrm}(a)(c) \leq 1) \Leftrightarrow (\operatorname{fnrm}(b)(c) \leq 1) \end{array}}}
\vnbprop{nv-ideal-principal-of}{\vnbstmt{\begin{array}{@{}l@{}} \forall k, a, i, g. \\ \quad \operatorname{is\text{-}nv\text{-}domain}(k, a) \wedge \\ \quad \operatorname{is\text{-}nv\text{-}ideal}(k, a, i) \wedge \\ \quad g \in i \wedge \\ \quad \forall x \in i.\; \exists l \in a.\; (x = \operatorname{mul}(k)(l, g)) \Rightarrow \\ \quad \quad i = \operatorname{nv\text{-}principal}(k, a, g) \end{array}}}
\vnbprop{nv-compact-maximal-max}{\vnbstmt{\begin{array}{@{}l@{}} \forall k, a. \\ \quad \operatorname{is\text{-}compact\text{-}nv\text{-}domain}(k, a) \Rightarrow  \\ \quad \quad \exists w \in \operatorname{nv\text{-}maximal}(k, a). \forall x \in \operatorname{nv\text{-}maximal}(k, a). \\ \quad \quad \quad \operatorname{fnrm}(k)(x) \leq \operatorname{fnrm}(k)(w) \end{array}}}
\vnbprop{ccr-2-13-converse}{\vnbstmt{\begin{array}{@{}l@{}} \forall k, a. \\ \quad (\operatorname{is\text{-}compact\text{-}nv\text{-}domain}(k, a) \wedge (\operatorname{is\text{-}fractions\text{-}of}(k, a) \wedge (a = \operatorname{nv\text{-}integers}(k)))) \Rightarrow  \\ \quad \quad \operatorname{is\text{-}nv\text{-}pid}(k, a) \end{array}}}
\vnbprop{ccr-2-13}{\vnbstmt{\begin{array}{@{}l@{}} \forall k, a. \\ \quad (\operatorname{is\text{-}compact\text{-}nv\text{-}domain}(k, a) \wedge \operatorname{is\text{-}fractions\text{-}of}(k, a)) \Rightarrow  \\ \quad \quad (a = \operatorname{nv\text{-}integers}(k)) \Leftrightarrow \operatorname{is\text{-}nv\text{-}pid}(k, a) \end{array}}}

\vnbbookfile{theorem-library/ccr1109a2-val}{theorem-library/ccr1109a2-val}
\noindent{\small ccr1109a2-val.scm -- wave of 2026-11-09/10 (agent CCR-1109-A2), wired 2026-11-10 from scratchpad/ccr1109a2/ccr1109a2-val.scm. Library: complex-analysis.}\par
\vnbprop{rr-mul-lt-cancel-iff}{\vnbstmt{\forall p, m, l \in \mathbb{R}.\; ((0 < l) \Rightarrow (((p \cdot l) < (m \cdot l)) \Leftrightarrow (p < m)))}}
\vnbprop{rr-prod-lt-one-iff}{\vnbstmt{\begin{array}{@{}l@{}} \forall p, q, r \in \mathbb{R}. \\ \quad ((0 < q) \wedge ((q \cdot r) = 1)) \Rightarrow (((p \cdot q) < 1) \Leftrightarrow (p < r)) \end{array}}}
\vnbprop{rr-log-lt-iff}{\vnbstmt{\forall t, s \in \mathbb{R}.\; (((0 < t) \wedge (0 < s)) \Rightarrow ((t < s) \Leftrightarrow (\log\left(t\right) < \log\left(s\right))))}}
\vnbprop{rr-cut-witness}{\vnbstmt{\begin{array}{@{}l@{}} \forall u, v \in \mathbb{R}. \\ \quad ((0 \leq u) \wedge (u < v)) \Rightarrow \exists m, n \in \mathbb{N}.\; (((n \cdot u) < m) \wedge (m \leq (n \cdot v))) \end{array}}}
\vnbprop{rr-cut-eq}{\vnbstmt{\begin{array}{@{}l@{}} \forall u, v \in \mathbb{R}. \\ \quad ((0 \leq u) \wedge ((0 \leq v) \wedge \forall m, n \in \mathbb{N}.\; (((n \cdot u) < m) \Leftrightarrow ((n \cdot v) < m)))) \Rightarrow  \\ \quad \quad u = v \end{array}}}
\vnbprop{c9v-pow-lt-iff}{\vnbstmt{\begin{array}{@{}l@{}} \forall k.; \forall x, a \in \operatorname{carr}(k), u \in \mathbb{R}, m, n \in \mathbb{N}. \\ \quad \operatorname{is\text{-}normed\text{-}field}(k) \wedge \\ \quad 0 < \operatorname{fnrm}(k)(x) \wedge \\ \quad 1 < \operatorname{fnrm}(k)(a) \wedge \\ \quad (u \cdot \log\left(\operatorname{fnrm}(k)(a)\right)) = \log\left(\operatorname{fnrm}(k)(x)\right) \Rightarrow \\ \quad \quad \operatorname{fnrm}(k)(\operatorname{ring\text{-}power}(\operatorname{normed\text{-}field\text{-}as\text{-}commutative\text{-}ring}(k), x, n)) \\ \quad \quad <\; \operatorname{fnrm}(k)(\operatorname{ring\text{-}power}(\operatorname{normed\text{-}field\text{-}as\text{-}commutative\text{-}ring}(k), a, m)) \\ \quad \quad \Leftrightarrow\; (n \cdot u) < m \end{array}}}
\vnbprop{c9v-pow-same-ind}{\vnbstmt{\begin{array}{@{}l@{}} \forall n \in \mathbb{N}, a, b.; \forall x \in \operatorname{carr}(a). \\ \quad \operatorname{is\text{-}nv\text{-}field}(a) \wedge \\ \quad \operatorname{is\text{-}nv\text{-}field}(b) \wedge \\ \quad (\operatorname{carr}(b) \simeq \operatorname{carr}(a)) \wedge \\ \quad \operatorname{add}(b) = \operatorname{add}(a) \wedge \\ \quad \operatorname{neg}(b) = \operatorname{neg}(a) \wedge \\ \quad \operatorname{mul}(b) = \operatorname{mul}(a) \wedge \\ \quad \operatorname{zero}(b) = \operatorname{zero}(a) \wedge \\ \quad \operatorname{one}(b) = \operatorname{one}(a) \Rightarrow \\ \quad \quad \operatorname{ring\text{-}power}(\operatorname{normed\text{-}field\text{-}as\text{-}commutative\text{-}ring}(a), x, n) \\ \quad \quad =\; \operatorname{ring\text{-}power}(\operatorname{normed\text{-}field\text{-}as\text{-}commutative\text{-}ring}(b), x, n) \end{array}}}
\vnbprop{c9v-ge-one}{\vnbstmt{\begin{array}{@{}l@{}} \forall a, b.; \forall c, x \in \operatorname{carr}(a). \\ \quad \operatorname{is\text{-}nv\text{-}field}(a) \wedge \\ \quad \operatorname{is\text{-}nv\text{-}field}(b) \wedge \\ \quad (\operatorname{carr}(b) \simeq \operatorname{carr}(a)) \wedge \\ \quad \operatorname{add}(b) = \operatorname{add}(a) \wedge \\ \quad \operatorname{neg}(b) = \operatorname{neg}(a) \wedge \\ \quad \operatorname{mul}(b) = \operatorname{mul}(a) \wedge \\ \quad \operatorname{zero}(b) = \operatorname{zero}(a) \wedge \\ \quad \operatorname{one}(b) = \operatorname{one}(a) \wedge \\ \quad \forall c \in \operatorname{carr}(a).\; ((\operatorname{fnrm}(a)(c) < 1) \Leftrightarrow (\operatorname{fnrm}(b)(c) < 1)) \wedge \\ \quad 1 < \operatorname{fnrm}(b)(c) \wedge \\ \quad 1 < \operatorname{fnrm}(a)(c) \wedge \\ \quad 1 \leq \operatorname{fnrm}(b)(x) \wedge \\ \quad 1 \leq \operatorname{fnrm}(a)(x) \Rightarrow \\ \quad \quad \operatorname{fnrm}(a)(x) = {\operatorname{fnrm}(b)(x)}^{(\log\left(\operatorname{fnrm}(a)(c)\right) \cdot \operatorname{recip}(\log\left(\operatorname{fnrm}(b)(c)\right)))} \end{array}}}
\vnbprop{ccr-1-46-v-vi}{\vnbstmt{\begin{array}{@{}l@{}} \forall a, b. \\ \quad \operatorname{is\text{-}nv\text{-}field}(a) \wedge \\ \quad \operatorname{is\text{-}nv\text{-}field}(b) \wedge \\ \quad (\operatorname{carr}(b) \simeq \operatorname{carr}(a)) \wedge \\ \quad \operatorname{add}(b) = \operatorname{add}(a) \wedge \\ \quad \operatorname{neg}(b) = \operatorname{neg}(a) \wedge \\ \quad \operatorname{mul}(b) = \operatorname{mul}(a) \wedge \\ \quad \operatorname{zero}(b) = \operatorname{zero}(a) \wedge \\ \quad \operatorname{one}(b) = \operatorname{one}(a) \wedge \\ \quad \forall c \in \operatorname{carr}(a).\; ((\operatorname{fnrm}(a)(c) \leq 1) \Leftrightarrow (\operatorname{fnrm}(b)(c) \leq 1)) \Rightarrow \\ \quad \quad \exists l \in \mathbb{R}. \\ \quad \quad \quad (0 < l) \wedge \forall c \in \operatorname{carr}(a).\; (\operatorname{fnrm}(a)(c) = {\operatorname{fnrm}(b)(c)}^{l}) \end{array}}}
\vnbprop{metric-open-transfer}{\vnbstmt{\begin{array}{@{}l@{}} \forall s_{1}, s_{2}, o. \\ \quad \operatorname{is\text{-}metric\text{-}space}(s_{1}) \wedge \\ \quad (\operatorname{pts}(s_{2}) \simeq \operatorname{pts}(s_{1})) \wedge \\ \quad \forall c \in \operatorname{pts}(s_{1}), r. \\ \quad \quad \operatorname{pos\text{-}rr}(r) \Rightarrow  \\ \quad \quad \quad \exists t.\; (\operatorname{pos\text{-}rr}(t) \wedge (\operatorname{ball}(s_{1}, c, t) \subseteq \operatorname{ball}(s_{2}, c, r))) \wedge \\ \quad \operatorname{is\text{-}open}(s_{2}, o) \Rightarrow \\ \quad \quad \operatorname{is\text{-}open}(s_{1}, o) \end{array}}}
\vnbprop{nv-ball-transfer-pow}{\vnbstmt{\begin{array}{@{}l@{}} \forall a, b.; \forall l \in \mathbb{R}, c \in \operatorname{pts}(\operatorname{nf\text{-}metric\text{-}space}(a)), r. \\ \quad \operatorname{is\text{-}nv\text{-}field}(a) \wedge \\ \quad \operatorname{is\text{-}nv\text{-}field}(b) \wedge \\ \quad (\operatorname{carr}(b) \simeq \operatorname{carr}(a)) \wedge \\ \quad \operatorname{add}(b) = \operatorname{add}(a) \wedge \\ \quad \operatorname{neg}(b) = \operatorname{neg}(a) \wedge \\ \quad \operatorname{mul}(b) = \operatorname{mul}(a) \wedge \\ \quad \operatorname{zero}(b) = \operatorname{zero}(a) \wedge \\ \quad \operatorname{one}(b) = \operatorname{one}(a) \wedge \\ \quad 0 < l \wedge \\ \quad \forall d \in \operatorname{carr}(a).\; (\operatorname{fnrm}(a)(d) = {\operatorname{fnrm}(b)(d)}^{l}) \wedge \\ \quad \operatorname{pos\text{-}rr}(r) \Rightarrow \\ \quad \quad \exists t. \\ \quad \quad \quad (\operatorname{pos\text{-}rr}(t) \wedge (\operatorname{ball}(\operatorname{nf\text{-}metric\text{-}space}(a), c, t) \subseteq \operatorname{ball}(\operatorname{nf\text{-}metric\text{-}space}(b), c, r))) \end{array}}}
\vnbprop{nv-ball-transfer-root}{\vnbstmt{\begin{array}{@{}l@{}} \forall a, b.; \forall l \in \mathbb{R}, c \in \operatorname{pts}(\operatorname{nf\text{-}metric\text{-}space}(b)), r. \\ \quad \operatorname{is\text{-}nv\text{-}field}(a) \wedge \\ \quad \operatorname{is\text{-}nv\text{-}field}(b) \wedge \\ \quad (\operatorname{carr}(b) \simeq \operatorname{carr}(a)) \wedge \\ \quad \operatorname{add}(b) = \operatorname{add}(a) \wedge \\ \quad \operatorname{neg}(b) = \operatorname{neg}(a) \wedge \\ \quad \operatorname{mul}(b) = \operatorname{mul}(a) \wedge \\ \quad \operatorname{zero}(b) = \operatorname{zero}(a) \wedge \\ \quad \operatorname{one}(b) = \operatorname{one}(a) \wedge \\ \quad 0 < l \wedge \\ \quad \forall d \in \operatorname{carr}(a).\; (\operatorname{fnrm}(a)(d) = {\operatorname{fnrm}(b)(d)}^{l}) \wedge \\ \quad \operatorname{pos\text{-}rr}(r) \Rightarrow \\ \quad \quad \exists t. \\ \quad \quad \quad (\operatorname{pos\text{-}rr}(t) \wedge (\operatorname{ball}(\operatorname{nf\text{-}metric\text{-}space}(b), c, t) \subseteq \operatorname{ball}(\operatorname{nf\text{-}metric\text{-}space}(a), c, r))) \end{array}}}
\vnbprop{ccr-1-46-vi-i}{\vnbstmt{\begin{array}{@{}l@{}} \forall a, b, o. \\ \quad \operatorname{is\text{-}nv\text{-}field}(a) \wedge \\ \quad \operatorname{is\text{-}nv\text{-}field}(b) \wedge \\ \quad (\operatorname{carr}(b) \simeq \operatorname{carr}(a)) \wedge \\ \quad \operatorname{add}(b) = \operatorname{add}(a) \wedge \\ \quad \operatorname{neg}(b) = \operatorname{neg}(a) \wedge \\ \quad \operatorname{mul}(b) = \operatorname{mul}(a) \wedge \\ \quad \operatorname{zero}(b) = \operatorname{zero}(a) \wedge \\ \quad \operatorname{one}(b) = \operatorname{one}(a) \wedge \\ \quad \exists l \in \mathbb{R}. \\ \quad \quad (0 < l) \wedge \forall c \in \operatorname{carr}(a).\; (\operatorname{fnrm}(a)(c) = {\operatorname{fnrm}(b)(c)}^{l}) \Rightarrow \\ \quad \quad \operatorname{is\text{-}open}(\operatorname{nf\text{-}metric\text{-}space}(a), o) \Leftrightarrow \operatorname{is\text{-}open}(\operatorname{nf\text{-}metric\text{-}space}(b), o) \end{array}}}
\vnbprop{nv-open-lt-one}{\vnbstmt{\begin{array}{@{}l@{}} \forall a, b.; \forall c \in \operatorname{carr}(a). \\ \quad \operatorname{is\text{-}nv\text{-}field}(a) \wedge \\ \quad \operatorname{is\text{-}nv\text{-}field}(b) \wedge \\ \quad (\operatorname{carr}(b) \simeq \operatorname{carr}(a)) \wedge \\ \quad \operatorname{add}(b) = \operatorname{add}(a) \wedge \\ \quad \operatorname{neg}(b) = \operatorname{neg}(a) \wedge \\ \quad \operatorname{mul}(b) = \operatorname{mul}(a) \wedge \\ \quad \operatorname{zero}(b) = \operatorname{zero}(a) \wedge \\ \quad \operatorname{one}(b) = \operatorname{one}(a) \wedge \\ \quad \forall o. \\ \quad \quad \operatorname{is\text{-}open}(\operatorname{nf\text{-}metric\text{-}space}(b), o) \Rightarrow \operatorname{is\text{-}open}(\operatorname{nf\text{-}metric\text{-}space}(a), o) \wedge \\ \quad \operatorname{fnrm}(a)(c) < 1 \Rightarrow \\ \quad \quad \operatorname{fnrm}(b)(c) < 1 \end{array}}}
\vnbprop{ccr-1-46-i-iii}{\vnbstmt{\begin{array}{@{}l@{}} \forall a, b.; \forall c \in \operatorname{carr}(a). \\ \quad \operatorname{is\text{-}nv\text{-}field}(a) \wedge \\ \quad \operatorname{is\text{-}nv\text{-}field}(b) \wedge \\ \quad (\operatorname{carr}(b) \simeq \operatorname{carr}(a)) \wedge \\ \quad \operatorname{add}(b) = \operatorname{add}(a) \wedge \\ \quad \operatorname{neg}(b) = \operatorname{neg}(a) \wedge \\ \quad \operatorname{mul}(b) = \operatorname{mul}(a) \wedge \\ \quad \operatorname{zero}(b) = \operatorname{zero}(a) \wedge \\ \quad \operatorname{one}(b) = \operatorname{one}(a) \wedge \\ \quad \forall o.\; (\operatorname{is\text{-}open}(\operatorname{nf\text{-}metric\text{-}space}(a), o) \Leftrightarrow \operatorname{is\text{-}open}(\operatorname{nf\text{-}metric\text{-}space}(b), o)) \Rightarrow \\ \quad \quad (\operatorname{fnrm}(a)(c) < 1) \Leftrightarrow (\operatorname{fnrm}(b)(c) < 1) \end{array}}}

\vnbbookfile{theorem-library/ccr1109a2-res}{theorem-library/ccr1109a2-res}
\noindent{\small ccr1109a2-res.scm -- wave of 2026-11-09/10 (agent CCR-1109-A2), wired 2026-11-10 from scratchpad/ccr1109a2/ccr1109a2-res.scm. Library: complex-analysis.}\par
\vnbprop{ccr-2-23}{\vnbstmt{\begin{array}{@{}l@{}} \forall k, a, i. \\ \quad \operatorname{is\text{-}nv\text{-}domain}(k, a) \wedge \\ \quad i \in (\operatorname{nv\text{-}integers}(k) \to a) \wedge \\ \quad \forall x \in \operatorname{nv\text{-}integers}(k).\; (\operatorname{fnrm}(k)(\operatorname{add}(k)(x, \operatorname{neg}(k)(i(x)))) < 1) \wedge \\ \quad \forall x \in a.\; (i(x) = x) \Rightarrow \\ \quad \quad \exists f.\; \operatorname{is\text{-}ring\text{-}iso}(\operatorname{nv\text{-}residue}(k, a), \operatorname{nv\text{-}residue}(k, \operatorname{nv\text{-}integers}(k)), f) \end{array}}}

\vnbbookfile{theorem-library/ccr1109a4-laurent-complete}{theorem-library/ccr1109a4-laurent-complete}
\noindent{\small ccr1109a4-laurent-complete.scm -- wave of 2026-11-09/10 (agent CCR-1109-A4), wired 2026-11-10 from scratchpad/ccr1109a4/ccr1109a4-laurent-complete.scm. Library: complex-analysis.}\par
\vnbprop{rpow-star-strict-anti}{\vnbstmt{\begin{array}{@{}l@{}} \forall q, \mathit{sq}, \mathit{tq} \in \mathbb{R}. \\ \quad ((0 < q) \wedge ((q < 1) \wedge (\mathit{sq} < \mathit{tq}))) \Rightarrow ({q}^{\mathit{tq}} < {q}^{\mathit{sq}}) \end{array}}}
\vnbprop{rpow-star-small}{\vnbstmt{\begin{array}{@{}l@{}} \forall q \in \mathbb{R}, p. \\ \quad ((0 < q) \wedge ((q < 1) \wedge \operatorname{pos\text{-}rr}(p))) \Rightarrow \exists \mathit{zq} \in \mathbb{N}.\; ({q}^{\mathit{zq}} \leq p) \end{array}}}
\vnbprop{laurent-abs-le-iff}{\vnbstmt{\begin{array}{@{}l@{}} \forall q.; \forall \mathit{rq} \in \mathbb{R}, \mathit{fq} \in \operatorname{laurent\text{-}set}(q), \mathit{nq} \in \mathbb{Z}. \\ \quad (\operatorname{is\text{-}ring}(q) \wedge ((0 < \mathit{rq}) \wedge (\mathit{rq} < 1))) \Rightarrow  \\ \quad \quad (\operatorname{laurent\text{-}abs}(q, \mathit{rq}, \mathit{fq}) \leq {\mathit{rq}}^{\mathit{nq}}) \Leftrightarrow \forall j \in \mathbb{Z}.\; ((j < \mathit{nq}) \Rightarrow (\operatorname{fq}(j) = \operatorname{zero}(q))) \end{array}}}
\vnbprop{laurent-nf-pts}{\vnbstmt{\begin{array}{@{}l@{}} \forall q.; \forall \mathit{rq} \in \mathbb{R}. \\ \quad (\operatorname{is\text{-}field\text{-}ring}(q) \wedge ((0 < \mathit{rq}) \wedge (\mathit{rq} < 1))) \Rightarrow  \\ \quad \quad (\operatorname{pts}(\operatorname{nf\text{-}metric\text{-}space}(\operatorname{laurent\text{-}nf}(q, \mathit{rq}))) \simeq \operatorname{laurent\text{-}set}(q)) \end{array}}}
\vnbprop{laurent-nf-dist}{\vnbstmt{\begin{array}{@{}l@{}} \forall q.; \forall \mathit{rq} \in \mathbb{R}, \mathit{fq}, \mathit{gq} \in \operatorname{laurent\text{-}set}(q). \\ \quad (\operatorname{is\text{-}field\text{-}ring}(q) \wedge ((0 < \mathit{rq}) \wedge (\mathit{rq} < 1))) \Rightarrow  \\ \quad \quad \operatorname{dist}(\operatorname{nf\text{-}metric\text{-}space}(\operatorname{laurent\text{-}nf}(q, \mathit{rq})))(\mathit{fq}, \mathit{gq}) \\ \quad \quad =\; \operatorname{laurent\text{-}abs}(q, \mathit{rq}, \operatorname{laurent\text{-}add}(q, \mathit{fq}, \operatorname{laurent\text{-}neg}(q, \mathit{gq}))) \end{array}}}
\vnbprop{laurent-dist-le-iff}{\vnbstmt{\begin{array}{@{}l@{}} \forall q.; \forall \mathit{rq} \in \mathbb{R}, \mathit{fq}, \mathit{gq} \in \operatorname{laurent\text{-}set}(q), \mathit{nq} \in \mathbb{Z}. \\ \quad (\operatorname{is\text{-}field\text{-}ring}(q) \wedge ((0 < \mathit{rq}) \wedge (\mathit{rq} < 1))) \Rightarrow  \\ \quad \quad \operatorname{dist}(\operatorname{nf\text{-}metric\text{-}space}(\operatorname{laurent\text{-}nf}(q, \mathit{rq})))(\mathit{fq}, \mathit{gq}) \leq {\mathit{rq}}^{\mathit{nq}} \\ \quad \quad \Leftrightarrow\; \forall j \in \mathbb{Z}.\; ((j < \mathit{nq}) \Rightarrow (\operatorname{fq}(j) = \operatorname{gq}(j))) \end{array}}}
\vnbprop{laurent-converges-iff}{\vnbstmt{\begin{array}{@{}l@{}} \forall q.; \forall \mathit{rq} \in \mathbb{R}, \mathit{sq} \in (\mathbb{N} \to \operatorname{laurent\text{-}set}(q)), \mathit{lq} \in \operatorname{laurent\text{-}set}(q). \\ \quad (\operatorname{is\text{-}field\text{-}ring}(q) \wedge ((0 < \mathit{rq}) \wedge (\mathit{rq} < 1))) \Rightarrow  \\ \quad \quad \operatorname{converges\text{-}to}(\operatorname{nf\text{-}metric\text{-}space}(\operatorname{laurent\text{-}nf}(q, \mathit{rq})), \mathit{sq}, \mathit{lq}) \\ \quad \quad \Leftrightarrow\; \forall \mathit{nq} \in \mathbb{Z}. \exists \mathit{mq} \in \mathbb{N}. \forall \mathit{iq} \in \mathbb{N}, j \in \mathbb{Z}. \\ \quad \quad \quad ((\mathit{mq} \leq \mathit{iq}) \wedge (j < \mathit{nq})) \Rightarrow (\operatorname{sq}(\mathit{iq})(j) = \operatorname{lq}(j)) \end{array}}}
\vnbprop{laurent-cauchy-stab}{\vnbstmt{\begin{array}{@{}l@{}} \forall q.; \forall \mathit{rq} \in \mathbb{R}, \mathit{sq}.; \forall \mathit{nq} \in \mathbb{Z}. \\ \quad (\operatorname{is\text{-}field\text{-}ring}(q) \wedge ((0 < \mathit{rq}) \wedge ((\mathit{rq} < 1) \wedge \operatorname{is\text{-}cauchy\text{-}seq}(\operatorname{nf\text{-}metric\text{-}space}(\operatorname{laurent\text{-}nf}(q, \mathit{rq})), \mathit{sq})))) \Rightarrow  \\ \quad \quad \exists \mathit{mq} \in \mathbb{N}. \forall \mathit{iq}, r \in \mathbb{N}, j \in \mathbb{Z}. \\ \quad \quad \quad ((\mathit{mq} \leq \mathit{iq}) \wedge ((\mathit{mq} \leq r) \wedge (j < \mathit{nq}))) \Rightarrow (\operatorname{sq}(\mathit{iq})(j) = \operatorname{sq}(r)(j)) \end{array}}}
\vnbprop{laurent-limc-spec}{\vnbstmt{\begin{array}{@{}l@{}} \forall q.; \forall \mathit{sq} \in (\mathbb{N} \to \operatorname{laurent\text{-}set}(q)), \mathit{jq} \in \mathbb{Z}, \mathit{mq} \in \mathbb{N}. \\ \quad \forall \mathit{iq} \in \mathbb{N}.\; ((\mathit{mq} \leq \mathit{iq}) \Rightarrow (\operatorname{sq}(\mathit{iq})(\mathit{jq}) = \operatorname{sq}(\mathit{mq})(\mathit{jq}))) \Rightarrow  \\ \quad \quad \operatorname{laurent\text{-}limc}(q, \mathit{sq}, \mathit{jq}) \in \operatorname{carr}(q) \wedge \\ \quad \quad \forall \mathit{iq} \in \mathbb{N}.\; ((\mathit{mq} \leq \mathit{iq}) \Rightarrow (\operatorname{sq}(\mathit{iq})(\mathit{jq}) = \operatorname{laurent\text{-}limc}(q, \mathit{sq}, \mathit{jq}))) \end{array}}}
\vnbprop{laurent-lim-apply}{\vnbstmt{\forall q, \mathit{sq}.; \forall \mathit{jq} \in \mathbb{Z}.\; (\operatorname{laurent\text{-}lim}(q, \mathit{sq})(\mathit{jq}) \simeq \operatorname{laurent\text{-}limc}(q, \mathit{sq}, \mathit{jq}))}}
\vnbprop{laurent-lim-near}{\vnbstmt{\begin{array}{@{}l@{}} \forall q.; \forall \mathit{sq} \in (\mathbb{N} \to \operatorname{laurent\text{-}set}(q)), \mathit{nq} \in \mathbb{Z}. \\ \quad \forall \mathit{nq} \in \mathbb{Z}. \exists \mathit{mq} \in \mathbb{N}. \forall \mathit{iq}, r \in \mathbb{N}, j \in \mathbb{Z}. \\ \quad \quad ((\mathit{mq} \leq \mathit{iq}) \wedge ((\mathit{mq} \leq r) \wedge (j < \mathit{nq}))) \Rightarrow (\operatorname{sq}(\mathit{iq})(j) = \operatorname{sq}(r)(j)) \Rightarrow \\ \quad \quad \exists \mathit{mq} \in \mathbb{N}. \forall \mathit{iq} \in \mathbb{N}, j \in \mathbb{Z}. \\ \quad \quad \quad ((\mathit{mq} \leq \mathit{iq}) \wedge (j < \mathit{nq})) \Rightarrow (\operatorname{sq}(\mathit{iq})(j) = \operatorname{laurent\text{-}lim}(q, \mathit{sq})(j)) \end{array}}}
\vnbprop{laurent-lim-in}{\vnbstmt{\begin{array}{@{}l@{}} \forall q.; \forall \mathit{sq} \in (\mathbb{N} \to \operatorname{laurent\text{-}set}(q)). \\ \quad \forall \mathit{nq} \in \mathbb{Z}. \exists \mathit{mq} \in \mathbb{N}. \forall \mathit{iq}, r \in \mathbb{N}, j \in \mathbb{Z}. \\ \quad \quad ((\mathit{mq} \leq \mathit{iq}) \wedge ((\mathit{mq} \leq r) \wedge (j < \mathit{nq}))) \Rightarrow (\operatorname{sq}(\mathit{iq})(j) = \operatorname{sq}(r)(j)) \Rightarrow \\ \quad \quad \operatorname{laurent\text{-}lim}(q, \mathit{sq}) \in \operatorname{laurent\text{-}set}(q) \end{array}}}
\vnbprop{laurent-nf-complete}{\vnbstmt{\begin{array}{@{}l@{}} \forall q.; \forall \mathit{rq} \in \mathbb{R}. \\ \quad (\operatorname{is\text{-}field\text{-}ring}(q) \wedge ((0 < \mathit{rq}) \wedge (\mathit{rq} < 1))) \Rightarrow  \\ \quad \quad \operatorname{is\text{-}complete}(\operatorname{nf\text{-}metric\text{-}space}(\operatorname{laurent\text{-}nf}(q, \mathit{rq}))) \end{array}}}
\vnbprop{laurent-nf-fnrm-apply}{\vnbstmt{\begin{array}{@{}l@{}} \forall q.; \forall \mathit{rq} \in \mathbb{R}, \mathit{fq} \in \operatorname{laurent\text{-}set}(q). \\ \quad (\operatorname{is\text{-}field\text{-}ring}(q) \wedge ((0 < \mathit{rq}) \wedge (\mathit{rq} < 1))) \Rightarrow  \\ \quad \quad \operatorname{fnrm}(\operatorname{laurent\text{-}nf}(q, \mathit{rq}))(\mathit{fq}) = \operatorname{laurent\text{-}abs}(q, \mathit{rq}, \mathit{fq}) \end{array}}}
\vnbprop{laurent-integers-iff}{\vnbstmt{\begin{array}{@{}l@{}} \forall q.; \forall \mathit{rq} \in \mathbb{R}, \mathit{fq}. \\ \quad (\operatorname{is\text{-}field\text{-}ring}(q) \wedge ((0 < \mathit{rq}) \wedge (\mathit{rq} < 1))) \Rightarrow  \\ \quad \quad \mathit{fq} \in \operatorname{nv\text{-}integers}(\operatorname{laurent\text{-}nf}(q, \mathit{rq})) \\ \quad \quad \Leftrightarrow\; (\mathit{fq} \in \operatorname{laurent\text{-}set}(q)) \wedge \forall j \in \mathbb{Z}.\; ((j < 0) \Rightarrow (\operatorname{fq}(j) = \operatorname{zero}(q))) \end{array}}}
\vnbprop{laurent-integers-closed}{\vnbstmt{\begin{array}{@{}l@{}} \forall q.; \forall \mathit{rq} \in \mathbb{R}. \\ \quad (\operatorname{is\text{-}field\text{-}ring}(q) \wedge ((0 < \mathit{rq}) \wedge (\mathit{rq} < 1))) \Rightarrow  \\ \quad \quad \operatorname{is\text{-}closed}(\operatorname{nf\text{-}metric\text{-}space}(\operatorname{laurent\text{-}nf}(q, \mathit{rq})), \operatorname{nv\text{-}integers}(\operatorname{laurent\text{-}nf}(q, \mathit{rq}))) \end{array}}}
\vnbprop{laurent-integers-complete}{\vnbstmt{\begin{array}{@{}l@{}} \forall q.; \forall \mathit{rq} \in \mathbb{R}. \\ \quad (\operatorname{is\text{-}field\text{-}ring}(q) \wedge ((0 < \mathit{rq}) \wedge (\mathit{rq} < 1))) \Rightarrow  \\ \quad \quad \operatorname{is\text{-}complete}(\operatorname{subspace\text{-}ms}(\operatorname{nf\text{-}metric\text{-}space}(\operatorname{laurent\text{-}nf}(q, \mathit{rq})), \operatorname{nv\text{-}integers}(\operatorname{laurent\text{-}nf}(q, \mathit{rq})))) \end{array}}}
\vnbprop{laurent-x-power-abs}{\vnbstmt{\begin{array}{@{}l@{}} \forall q.; \forall \mathit{rq} \in \mathbb{R}, \mathit{zq} \in \mathbb{Z}. \\ \quad (\operatorname{is\text{-}field\text{-}ring}(q) \wedge ((0 < \mathit{rq}) \wedge (\mathit{rq} < 1))) \Rightarrow  \\ \quad \quad \operatorname{laurent\text{-}mono}(q, \operatorname{one}(q), \mathit{zq}) \in \operatorname{laurent\text{-}set}(q) \wedge \\ \quad \quad \neg (\operatorname{laurent\text{-}mono}(q, \operatorname{one}(q), \mathit{zq}) = \operatorname{laurent\text{-}zero}(q)) \wedge \\ \quad \quad \operatorname{laurent\text{-}abs}(q, \mathit{rq}, \operatorname{laurent\text{-}mono}(q, \operatorname{one}(q), \mathit{zq})) = {\mathit{rq}}^{\mathit{zq}} \end{array}}}
\vnbprop{laurent-small-elements}{\vnbstmt{\begin{array}{@{}l@{}} \forall q.; \forall \mathit{rq} \in \mathbb{R}, p. \\ \quad (\operatorname{is\text{-}field\text{-}ring}(q) \wedge ((0 < \mathit{rq}) \wedge ((\mathit{rq} < 1) \wedge \operatorname{pos\text{-}rr}(p)))) \Rightarrow  \\ \quad \quad \exists s \in \operatorname{carr}(\operatorname{laurent\text{-}nf}(q, \mathit{rq})). \\ \quad \quad \quad \neg (s = \operatorname{zero}(\operatorname{laurent\text{-}nf}(q, \mathit{rq}))) \wedge (\operatorname{fnrm}(\operatorname{laurent\text{-}nf}(q, \mathit{rq}))(s) < p) \end{array}}}
\vnbprop{laurent-polys-iff}{\vnbstmt{\begin{array}{@{}l@{}} \forall q, \mathit{mq}, \mathit{gq}. \\ \quad \mathit{gq} \in \operatorname{laurent\text{-}polys}(q, \mathit{mq}) \\ \quad \Leftrightarrow\; \mathit{gq} \in (\mathbb{Z} \to \operatorname{carr}(q)) \wedge \\ \quad \forall j \in \mathbb{Z}.\; ((j < 0) \Rightarrow (\operatorname{gq}(j) = \operatorname{zero}(q))) \wedge \\ \quad \forall j \in \mathbb{Z}.\; ((\mathit{mq} \leq j) \Rightarrow (\operatorname{gq}(j) = \operatorname{zero}(q))) \end{array}}}
\vnbprop{laurent-polys-set}{\vnbstmt{\forall q, \mathit{mq}.\; (\operatorname{is\text{-}ring}(q) \Rightarrow (\operatorname{laurent\text{-}polys}(q, \mathit{mq}) \in \mathbf{Set}))}}
\vnbprop{laurent-mono-off}{\vnbstmt{\begin{array}{@{}l@{}} \forall q.; \forall \mathit{cq} \in \operatorname{carr}(q), \mathit{tq}, \mathit{vq} \in \mathbb{Z}. \\ \quad (\operatorname{is\text{-}ring}(q) \wedge \neg (\mathit{vq} = \mathit{tq})) \Rightarrow (\operatorname{laurent\text{-}mono}(q, \mathit{cq}, \mathit{tq})(\mathit{vq}) = \operatorname{zero}(q)) \end{array}}}
\vnbprop{laurent-polys-zero-finite}{\vnbstmt{\forall q.\; (\operatorname{is\text{-}ring}(q) \Rightarrow (|\operatorname{laurent\text{-}polys}(q, 0)| \in \mathbb{N}))}}
\vnbprop{laurent-polys-step}{\vnbstmt{\begin{array}{@{}l@{}} \forall q.; \forall \mathit{mq} \in \mathbb{Z}. \\ \quad (\operatorname{is\text{-}ring}(q) \wedge ((|\operatorname{carr}(q)| \in \mathbb{N}) \wedge (|\operatorname{laurent\text{-}polys}(q, \mathit{mq})| \in \mathbb{N}))) \Rightarrow  \\ \quad \quad |\operatorname{laurent\text{-}polys}(q, (\mathit{mq} + 1))| \in \mathbb{N} \end{array}}}
\vnbprop{laurent-polys-finite}{\vnbstmt{\begin{array}{@{}l@{}} \forall q \in \mathbb{N}, \mathit{kq}. \\ \quad (\operatorname{is\text{-}ring}(\mathit{kq}) \wedge (|\operatorname{carr}(\mathit{kq})| \in \mathbb{N})) \Rightarrow (|\operatorname{laurent\text{-}polys}(\mathit{kq}, q)| \in \mathbb{N}) \end{array}}}
\vnbprop{laurent-trunc-poly}{\vnbstmt{\begin{array}{@{}l@{}} \forall q.; \forall \mathit{gq} \in (\mathbb{Z} \to \operatorname{carr}(q)), \mathit{zq} \in \mathbb{Z}. \\ \quad (\operatorname{is\text{-}ring}(q) \wedge \forall j \in \mathbb{Z}.\; ((j < 0) \Rightarrow (\operatorname{gq}(j) = \operatorname{zero}(q)))) \Rightarrow  \\ \quad \quad \operatorname{laurent\text{-}trunc}(q, \mathit{gq}, (\mathit{zq} + -1)) \in \operatorname{laurent\text{-}polys}(q, \mathit{zq}) \wedge \\ \quad \quad \forall j \in \mathbb{Z}.\; ((j < \mathit{zq}) \Rightarrow (\operatorname{laurent\text{-}trunc}(q, \mathit{gq}, (\mathit{zq} + -1))(j) = \operatorname{gq}(j))) \end{array}}}
\vnbprop{laurent-polys-sub-integers}{\vnbstmt{\begin{array}{@{}l@{}} \forall q.; \forall \mathit{rq} \in \mathbb{R}, \mathit{zq} \in \mathbb{Z}, \mathit{gq} \in \operatorname{laurent\text{-}polys}(q, \mathit{zq}). \\ \quad (\operatorname{is\text{-}field\text{-}ring}(q) \wedge ((0 < \mathit{rq}) \wedge (\mathit{rq} < 1))) \Rightarrow  \\ \quad \quad \mathit{gq} \in \operatorname{nv\text{-}integers}(\operatorname{laurent\text{-}nf}(q, \mathit{rq})) \end{array}}}
\vnbprop{laurent-integers-set}{\vnbstmt{\begin{array}{@{}l@{}} \forall q.; \forall \mathit{rq} \in \mathbb{R}. \\ \quad (\operatorname{is\text{-}field\text{-}ring}(q) \wedge ((0 < \mathit{rq}) \wedge (\mathit{rq} < 1))) \Rightarrow  \\ \quad \quad \operatorname{nv\text{-}integers}(\operatorname{laurent\text{-}nf}(q, \mathit{rq})) \in \mathbf{Set} \wedge \\ \quad \quad \operatorname{nv\text{-}integers}(\operatorname{laurent\text{-}nf}(q, \mathit{rq})) \subseteq \operatorname{pts}(\operatorname{nf\text{-}metric\text{-}space}(\operatorname{laurent\text{-}nf}(q, \mathit{rq}))) \end{array}}}
\vnbprop{laurent-integers-tb}{\vnbstmt{\begin{array}{@{}l@{}} \forall q.; \forall \mathit{rq} \in \mathbb{R}. \\ \quad (\operatorname{is\text{-}field\text{-}ring}(q) \wedge ((|\operatorname{carr}(q)| \in \mathbb{N}) \wedge ((0 < \mathit{rq}) \wedge (\mathit{rq} < 1)))) \Rightarrow  \\ \quad \quad \operatorname{totally\text{-}bounded}(\operatorname{subspace\text{-}ms}(\operatorname{nf\text{-}metric\text{-}space}(\operatorname{laurent\text{-}nf}(q, \mathit{rq})), \operatorname{nv\text{-}integers}(\operatorname{laurent\text{-}nf}(q, \mathit{rq})))) \end{array}}}
\vnbprop{laurent-integers-compact}{\vnbstmt{\begin{array}{@{}l@{}} \forall q.; \forall \mathit{rq} \in \mathbb{R}. \\ \quad (\operatorname{is\text{-}field\text{-}ring}(q) \wedge ((|\operatorname{carr}(q)| \in \mathbb{N}) \wedge ((0 < \mathit{rq}) \wedge (\mathit{rq} < 1)))) \Rightarrow  \\ \quad \quad \operatorname{is\text{-}compact}(\operatorname{subspace\text{-}ms}(\operatorname{nf\text{-}metric\text{-}space}(\operatorname{laurent\text{-}nf}(q, \mathit{rq})), \operatorname{nv\text{-}integers}(\operatorname{laurent\text{-}nf}(q, \mathit{rq})))) \end{array}}}
\vnbprop{laurent-local-nv-field}{\vnbstmt{\begin{array}{@{}l@{}} \forall q.; \forall \mathit{rq} \in \mathbb{R}. \\ \quad (\operatorname{is\text{-}field\text{-}ring}(q) \wedge ((|\operatorname{carr}(q)| \in \mathbb{N}) \wedge ((0 < \mathit{rq}) \wedge (\mathit{rq} < 1)))) \Rightarrow  \\ \quad \quad \operatorname{is\text{-}local\text{-}nv\text{-}field}(\operatorname{laurent\text{-}nf}(q, \mathit{rq})) \end{array}}}

\vnbbookfile{theorem-library/ccr5-dual}{theorem-library/ccr5-dual}
\noindent{\small ccr5-dual.scm -- wave of 2026-11-09/10 (agent CCR-1109-D), wired 2026-11-10 from scratchpad/ccr1109d/ccr5-dual.scm. Library: complex-analysis.}\par
\vnbprop{c5-e-square}{\vnbstmt{\begin{array}{@{}l@{}} \forall z \in \mathbb{C}. \\ \quad (\operatorname{magnitude}(z) = 1) \Rightarrow  \\ \quad \quad (((z \cdot z) - 1) \cdot \operatorname{conjugate}(((z \cdot z) - 1))) = (((z - 1) \cdot \operatorname{conjugate}((z - 1))) \cdot (4 - ((z - 1) \cdot \operatorname{conjugate}((z - 1))))) \end{array}}}
\vnbprop{c5-nss-ind}{\vnbstmt{\begin{array}{@{}l@{}} \forall n \in \mathbb{N}, f, i.; \forall l \in \mathbb{R}, y \in \operatorname{carr}(f). \\ \quad \operatorname{is\text{-}nv\text{-}field}(f) \wedge \\ \quad \operatorname{is\text{-}add\text{-}character}(f, i) \wedge \\ \quad \forall y \in \operatorname{carr}(f).\; ((\operatorname{fnrm}(f)(y) \leq l) \Rightarrow (\operatorname{magnitude}((i(y) - 1)) < 1)) \wedge \\ \quad \operatorname{fnrm}(f)(y) \leq l \Rightarrow \\ \quad \quad ((n + 1) \cdot ((i(y) - 1) \cdot \operatorname{conjugate}((i(y) - 1)))) < 1 \end{array}}}
\vnbprop{c5-nss}{\vnbstmt{\begin{array}{@{}l@{}} \forall f, i.; \forall l \in \mathbb{R}, y \in \operatorname{carr}(f). \\ \quad \operatorname{is\text{-}nv\text{-}field}(f) \wedge \\ \quad \operatorname{is\text{-}add\text{-}character}(f, i) \wedge \\ \quad \forall y \in \operatorname{carr}(f).\; ((\operatorname{fnrm}(f)(y) \leq l) \Rightarrow (\operatorname{magnitude}((i(y) - 1)) < 1)) \wedge \\ \quad \operatorname{fnrm}(f)(y) \leq l \Rightarrow \\ \quad \quad i(y) = 1 \end{array}}}
\vnbprop{ccr-6-17-ball}{\vnbstmt{\begin{array}{@{}l@{}} \forall f, i. \\ \quad (\operatorname{is\text{-}nv\text{-}field}(f) \wedge \operatorname{is\text{-}add\text{-}character}(f, i)) \Rightarrow  \\ \quad \quad \exists r. \\ \quad \quad \quad \operatorname{pos\text{-}rr}(r) \wedge \forall y \in \operatorname{carr}(f).\; ((\operatorname{fnrm}(f)(y) \leq r) \Rightarrow (i(y) = 1)) \end{array}}}
\vnbprop{ccr-6-17}{\vnbstmt{\begin{array}{@{}l@{}} \forall f, i. \\ \quad (\operatorname{is\text{-}nv\text{-}field}(f) \wedge \operatorname{is\text{-}add\text{-}character}(f, i)) \Rightarrow \operatorname{is\text{-}locally\text{-}constant\text{-}fn}(f, i) \end{array}}}
\vnbprop{ccr-6-16}{\vnbstmt{\begin{array}{@{}l@{}} \forall f, u. \\ \quad \operatorname{is\text{-}nv\text{-}field}(f) \wedge \\ \quad \exists c.\; (\operatorname{is\text{-}nontrivial\text{-}character}(f, c) \wedge \operatorname{is\text{-}locally\text{-}constant\text{-}fn}(f, c)) \wedge \\ \quad \operatorname{is\text{-}add\text{-}character}(f, u) \Rightarrow \\ \quad \quad \operatorname{is\text{-}locally\text{-}constant\text{-}fn}(f, u) \end{array}}}
\vnbprop{c5-compact-norm-bound}{\vnbstmt{\begin{array}{@{}l@{}} \forall f, e. \\ \quad (\operatorname{is\text{-}normed\text{-}field}(f) \wedge ((e \subseteq \operatorname{carr}(f)) \wedge \operatorname{is\text{-}compact}(\operatorname{subspace\text{-}ms}(\operatorname{nf\text{-}metric\text{-}space}(f), e)))) \Rightarrow  \\ \quad \quad \exists b \in \mathbb{R}.\; ((0 < b) \wedge \forall x \in e.\; (\operatorname{fnrm}(f)(x) \leq b)) \end{array}}}
\vnbprop{ccr-5-12-continuous}{\vnbstmt{\begin{array}{@{}l@{}} \forall f, i.; \forall q \in (\mathbb{N} \to \operatorname{carr}(f)), a \in \operatorname{carr}(f). \\ \quad (\operatorname{is\text{-}nv\text{-}field}(f) \wedge (\operatorname{is\text{-}add\text{-}character}(f, i) \wedge \operatorname{converges\text{-}to}(\operatorname{nf\text{-}metric\text{-}space}(f), q, a))) \Rightarrow  \\ \quad \quad \operatorname{char\text{-}converges}(f, (\lambda k \in \mathbb{N}.\; \operatorname{duality\text{-}map}(f, i, q(k))), \operatorname{duality\text{-}map}(f, i, a)) \end{array}}}
\vnbprop{c5-lu-compact}{\vnbstmt{\begin{array}{@{}l@{}} \forall f, w. \\ \quad \operatorname{is\text{-}local\text{-}uniformizer}(f, w) \Rightarrow \operatorname{is\text{-}compact\text{-}nv\text{-}domain}(f, \operatorname{nv\text{-}integers}(f)) \end{array}}}
\vnbprop{c5-lu-in}{\vnbstmt{\forall f, w.\; (\operatorname{is\text{-}local\text{-}uniformizer}(f, w) \Rightarrow (w \in \operatorname{nv\text{-}integers}(f)))}}
\vnbprop{c5-lu-maximal}{\vnbstmt{\begin{array}{@{}l@{}} \forall f, w. \\ \quad \operatorname{is\text{-}local\text{-}uniformizer}(f, w) \Rightarrow  \\ \quad \quad \operatorname{nv\text{-}maximal}(f, \operatorname{nv\text{-}integers}(f)) = \operatorname{nv\text{-}principal}(f, \operatorname{nv\text{-}integers}(f), w) \end{array}}}
\vnbprop{c5-lu-nonzero}{\vnbstmt{\forall f, w.\; (\operatorname{is\text{-}local\text{-}uniformizer}(f, w) \Rightarrow \neg (w = \operatorname{zero}(f)))}}
\vnbprop{c5-lu-nv-field}{\vnbstmt{\forall f, w.\; (\operatorname{is\text{-}local\text{-}uniformizer}(f, w) \Rightarrow \operatorname{is\text{-}nv\text{-}field}(f))}}
\vnbprop{c5-lu-norm}{\vnbstmt{\begin{array}{@{}l@{}} \forall f, w. \\ \quad \operatorname{is\text{-}local\text{-}uniformizer}(f, w) \Rightarrow  \\ \quad \quad (w \in \operatorname{carr}(f)) \wedge ((0 < \operatorname{fnrm}(f)(w)) \wedge (\operatorname{fnrm}(f)(w) < 1)) \end{array}}}
\vnbprop{c5-integers-ball}{\vnbstmt{\begin{array}{@{}l@{}} \forall f. \\ \quad \operatorname{is\text{-}nv\text{-}field}(f) \Rightarrow  \\ \quad \quad \operatorname{nv\text{-}integers}(f) = \operatorname{closed\text{-}ball}(\operatorname{nf\text{-}metric\text{-}space}(f), \operatorname{zero}(f), 1) \end{array}}}
\vnbprop{c5-lu-ball-compact}{\vnbstmt{\begin{array}{@{}l@{}} \forall f, w.; \forall r \in \mathbb{R}. \\ \quad (\operatorname{is\text{-}local\text{-}uniformizer}(f, w) \wedge (0 < r)) \Rightarrow  \\ \quad \quad \operatorname{is\text{-}compact}(\operatorname{subspace\text{-}ms}(\operatorname{nf\text{-}metric\text{-}space}(f), \operatorname{closed\text{-}ball}(\operatorname{nf\text{-}metric\text{-}space}(f), \operatorname{zero}(f), r))) \end{array}}}
\vnbprop{c5-lu-discrete}{\vnbstmt{\begin{array}{@{}l@{}} \forall f, w.; \forall x, y \in \operatorname{carr}(f). \\ \quad (\operatorname{is\text{-}local\text{-}uniformizer}(f, w) \wedge (\operatorname{fnrm}(f)(x) < \operatorname{fnrm}(f)(y))) \Rightarrow  \\ \quad \quad \operatorname{fnrm}(f)(x) \leq (\operatorname{fnrm}(f)(w) \cdot \operatorname{fnrm}(f)(y)) \end{array}}}
\vnbprop{ccr-5-10}{\vnbstmt{\begin{array}{@{}l@{}} \forall f, w, s, i. \\ \quad \operatorname{is\text{-}local\text{-}uniformizer}(f, w) \wedge \\ \quad \forall m \in \mathbb{N}.\; \operatorname{is\text{-}add\text{-}character}(f, s(m)) \wedge \\ \quad \operatorname{is\text{-}add\text{-}character}(f, i) \Rightarrow \\ \quad \quad \operatorname{char\text{-}converges}(f, s, i) \\ \quad \quad \Leftrightarrow\; \forall r. \\ \quad \quad \quad \operatorname{pos\text{-}rr}(r) \Rightarrow  \\ \quad \quad \quad \quad \exists n \in \mathbb{N}. \forall m \in \mathbb{N}, x \in \operatorname{carr}(f). \\ \quad \quad \quad \quad \quad ((n \leq m) \wedge (\operatorname{fnrm}(f)(x) \leq r)) \Rightarrow (s(m)(x) = i(x)) \end{array}}}
\vnbprop{ccr-5-15-bounded}{\vnbstmt{\begin{array}{@{}l@{}} \forall f, i.; \forall q \in (\mathbb{N} \to \operatorname{carr}(f)). \\ \quad \operatorname{is\text{-}nv\text{-}field}(f) \wedge \\ \quad \operatorname{is\text{-}nontrivial\text{-}character}(f, i) \wedge \\ \quad \forall k \in \mathbb{N}, a \in \operatorname{carr}(f). \\ \quad \quad (\operatorname{fnrm}(f)(a) \leq 1) \Rightarrow (i(\operatorname{mul}(f)(q(k), a)) = i(\operatorname{mul}(f)(q(0), a))) \Rightarrow \\ \quad \quad \exists b \in \mathbb{R}.\; ((0 < b) \wedge \forall k \in \mathbb{N}.\; (\operatorname{fnrm}(f)(q(k)) \leq b)) \end{array}}}
\vnbprop{ccr-5-15}{\vnbstmt{\begin{array}{@{}l@{}} \forall f, w, i.; \forall q \in (\mathbb{N} \to \operatorname{carr}(f)). \\ \quad \operatorname{is\text{-}local\text{-}uniformizer}(f, w) \wedge \\ \quad \operatorname{is\text{-}nontrivial\text{-}character}(f, i) \wedge \\ \quad \forall k \in \mathbb{N}, a \in \operatorname{carr}(f). \\ \quad \quad (\operatorname{fnrm}(f)(a) \leq 1) \Rightarrow (i(\operatorname{mul}(f)(q(k), a)) = i(\operatorname{mul}(f)(q(0), a))) \Rightarrow \\ \quad \quad \exists b. \\ \quad \quad \quad (\operatorname{strictly\text{-}mono\text{-}nn}(b) \wedge \exists l \in \operatorname{carr}(f).\; \operatorname{converges\text{-}to}(\operatorname{nf\text{-}metric\text{-}space}(f), \operatorname{subseq}(q, b), l)) \end{array}}}
\vnbprop{c5-chi-min-ind}{\vnbstmt{\begin{array}{@{}l@{}} \forall n \in \mathbb{N}, f, w, i.; \forall a \in \mathbb{R}, b \in \operatorname{carr}(f), d \in \mathbb{R}. \\ \quad \operatorname{is\text{-}local\text{-}uniformizer}(f, w) \wedge \\ \quad \operatorname{is\text{-}add\text{-}character}(f, i) \wedge \\ \quad \forall y \in \operatorname{carr}(f).\; ((\operatorname{fnrm}(f)(y) \leq a) \Rightarrow (i(y) = 1)) \wedge \\ \quad \neg (i(b) = 1) \wedge \\ \quad d = (a - (\operatorname{fnrm}(f)(w) \cdot a)) \Rightarrow \\ \quad \quad \exists c \in \operatorname{carr}(f). \\ \quad \quad \quad (\neg (i(c) = 1) \wedge \forall y \in \operatorname{carr}(f).\; ((\operatorname{fnrm}(f)(y) < \operatorname{fnrm}(f)(c)) \Rightarrow (i(y) = 1))) \vee \\ \quad \quad \exists c \in \operatorname{carr}(f).\; (\neg (i(c) = 1) \wedge ((\operatorname{fnrm}(f)(c) + (n \cdot d)) \leq \operatorname{fnrm}(f)(b))) \end{array}}}
\vnbprop{c5-chi-min}{\vnbstmt{\begin{array}{@{}l@{}} \forall f, w, i, a.; \forall b \in \operatorname{carr}(f). \\ \quad \operatorname{is\text{-}local\text{-}uniformizer}(f, w) \wedge \\ \quad \operatorname{is\text{-}add\text{-}character}(f, i) \wedge \\ \quad \operatorname{pos\text{-}rr}(a) \wedge \\ \quad \forall y \in \operatorname{carr}(f).\; ((\operatorname{fnrm}(f)(y) \leq a) \Rightarrow (i(y) = 1)) \wedge \\ \quad \neg (i(b) = 1) \Rightarrow \\ \quad \quad \exists c \in \operatorname{carr}(f). \\ \quad \quad \quad (\neg (i(c) = 1) \wedge \forall y \in \operatorname{carr}(f).\; ((\operatorname{fnrm}(f)(y) < \operatorname{fnrm}(f)(c)) \Rightarrow (i(y) = 1))) \end{array}}}
\vnbprop{c5-integers-mem}{\vnbstmt{\forall f, x.\; ((x \in \operatorname{nv\text{-}integers}(f)) \Leftrightarrow ((x \in \operatorname{carr}(f)) \wedge (\operatorname{fnrm}(f)(x) \leq 1)))}}
\vnbprop{c5-polar-mem}{\vnbstmt{\begin{array}{@{}l@{}} \forall f, w, v, h. \\ \quad (h \in \operatorname{bichar\text{-}polar}(f, w, v)) \Leftrightarrow ((h \in \operatorname{carr}(f)) \wedge \forall x \in v.\; (w(h, x) = 1)) \end{array}}}
\vnbprop{ccr-6-14}{\vnbstmt{\begin{array}{@{}l@{}} \forall f, w, i, r. \\ \quad \operatorname{is\text{-}local\text{-}uniformizer}(f, w) \wedge \\ \quad \operatorname{is\text{-}nontrivial\text{-}character}(f, i) \wedge \\ \quad \operatorname{pos\text{-}rr}(r) \wedge \\ \quad \forall y \in \operatorname{carr}(f).\; ((\operatorname{fnrm}(f)(y) \leq r) \Rightarrow (i(y) = 1)) \Rightarrow \\ \quad \quad \exists a.\; \operatorname{is\text{-}unipolar}(f, a) \end{array}}}
\vnbprop{ccr-6-3-converse}{\vnbstmt{\begin{array}{@{}l@{}} \forall f, i, w. \\ \quad (\operatorname{is\text{-}normed\text{-}field}(f) \wedge (\operatorname{is\text{-}dualizing\text{-}character}(f, i) \wedge \operatorname{is\text{-}bilinear\text{-}bicharacter}(f, w))) \Rightarrow  \\ \quad \quad \exists a \in \operatorname{carr}(f). \forall x, b \in \operatorname{carr}(f). \\ \quad \quad \quad w(x, b) = \operatorname{bichar\text{-}of}(f, i, a)(x, b) \end{array}}}
\vnbprop{ccr-6-3-unique}{\vnbstmt{\begin{array}{@{}l@{}} \forall f, i.; \forall a, b \in \operatorname{carr}(f). \\ \quad \operatorname{is\text{-}normed\text{-}field}(f) \wedge \\ \quad \operatorname{is\text{-}nontrivial\text{-}character}(f, i) \wedge \\ \quad \forall x, c \in \operatorname{carr}(f). \\ \quad \quad \operatorname{bichar\text{-}of}(f, i, a)(x, c) = \operatorname{bichar\text{-}of}(f, i, b)(x, c) \Rightarrow \\ \quad \quad a = b \end{array}}}
\vnbprop{ccr-6-20}{\vnbstmt{\begin{array}{@{}l@{}} \forall f, i.; \forall c \in \operatorname{carr}(f), w. \\ \quad \operatorname{is\text{-}normed\text{-}field}(f) \wedge \\ \quad \operatorname{is\text{-}dualizing\text{-}character}(f, i) \wedge \\ \quad \neg (c = \operatorname{zero}(f)) \wedge \\ \quad \operatorname{is\text{-}bilinear\text{-}bicharacter}(f, w) \Rightarrow \\ \quad \quad \exists a \in \operatorname{carr}(f). \forall x, b \in \operatorname{carr}(f). \\ \quad \quad \quad w(x, b) = \operatorname{bichar\text{-}of}(f, i, c)(\operatorname{mul}(f)(a, x), b) \wedge \\ \quad \quad \forall a, d \in \operatorname{carr}(f). \\ \quad \quad \quad \forall x, b \in \operatorname{carr}(f). \\ \quad \quad \quad \quad w(x, b) = \operatorname{bichar\text{-}of}(f, i, c)(\operatorname{mul}(f)(a, x), b) \wedge \\ \quad \quad \quad \forall x, b \in \operatorname{carr}(f). \\ \quad \quad \quad \quad w(x, b) = \operatorname{bichar\text{-}of}(f, i, c)(\operatorname{mul}(f)(d, x), b) \Rightarrow \\ \quad \quad \quad \quad a = d \end{array}}}
\vnbprop{ccr-6-18-unipolar}{\vnbstmt{\begin{array}{@{}l@{}} \forall f, w, i. \\ \quad (\operatorname{is\text{-}local\text{-}uniformizer}(f, w) \wedge \operatorname{is\text{-}nontrivial\text{-}character}(f, i)) \Rightarrow  \\ \quad \quad \exists a.\; \operatorname{is\text{-}unipolar}(f, a) \end{array}}}
\vnbprop{ccr-6-21}{\vnbstmt{\begin{array}{@{}l@{}} \forall f, w.; \forall c, a \in \operatorname{carr}(f). \\ \quad (\operatorname{is\text{-}normed\text{-}field}(f) \wedge (\operatorname{is\text{-}unipolar}(f, w) \wedge \neg (c = \operatorname{zero}(f)))) \Rightarrow  \\ \quad \quad a \in \operatorname{bichar\text{-}polar}(f, w, \operatorname{closed\text{-}ball}(\operatorname{nf\text{-}metric\text{-}space}(f), \operatorname{zero}(f), \operatorname{fnrm}(f)(c))) \\ \quad \quad \Leftrightarrow\; \operatorname{fnrm}(f)(\operatorname{mul}(f)(c, a)) \leq 1 \end{array}}}
\vnbprop{c5-unipolar-nonzero}{\vnbstmt{\begin{array}{@{}l@{}} \forall f, a, i, w.; \forall c \in \operatorname{carr}(f). \\ \quad \operatorname{is\text{-}local\text{-}uniformizer}(f, a) \wedge \\ \quad \operatorname{is\text{-}add\text{-}character}(f, i) \wedge \\ \quad \operatorname{is\text{-}unipolar}(f, w) \wedge \\ \quad \forall x, b \in \operatorname{carr}(f).\; (w(x, b) = \operatorname{bichar\text{-}of}(f, i, c)(x, b)) \Rightarrow \\ \quad \quad \neg (c = \operatorname{zero}(f)) \end{array}}}
\vnbprop{ccr-6-23}{\vnbstmt{\begin{array}{@{}l@{}} \forall f, a, i, w, q. \\ \quad \operatorname{is\text{-}local\text{-}uniformizer}(f, a) \wedge \\ \quad \operatorname{is\text{-}dualizing\text{-}character}(f, i) \wedge \\ \quad \operatorname{is\text{-}unipolar}(f, w) \wedge \\ \quad \operatorname{is\text{-}unipolar}(f, q) \Rightarrow \\ \quad \quad \exists b \in \operatorname{carr}(f). \\ \quad \quad \quad ((\operatorname{fnrm}(f)(b) = 1) \wedge \forall x, c \in \operatorname{carr}(f).\; (w(x, c) = q(\operatorname{mul}(f)(b, x), c))) \wedge \\ \quad \quad \forall b, d \in \operatorname{carr}(f). \\ \quad \quad \quad \forall x, c \in \operatorname{carr}(f).\; (w(x, c) = q(\operatorname{mul}(f)(b, x), c)) \wedge \\ \quad \quad \quad \forall x, c \in \operatorname{carr}(f).\; (w(x, c) = q(\operatorname{mul}(f)(d, x), c)) \Rightarrow \\ \quad \quad \quad \quad b = d \end{array}}}
\vnbprop{ccr-6-6}{\vnbstmt{\begin{array}{@{}l@{}} \forall f, i, w. \\ \quad \operatorname{is\text{-}normed\text{-}field}(f) \wedge \\ \quad \operatorname{is\text{-}dualizing\text{-}character}(f, i) \wedge \\ \quad \operatorname{is\text{-}bilinear\text{-}bicharacter}(f, w) \wedge \\ \quad \exists x, y \in \operatorname{carr}(f).\; \neg (w(x, y) = 1) \Rightarrow \\ \quad \quad \operatorname{is\text{-}faithful\text{-}bicharacter}(f, w) \end{array}}}

\vnbbookfile{theorem-library/ccr5-selfdual}{theorem-library/ccr5-selfdual}
\noindent{\small ccr5-selfdual.scm -- wave of 2026-11-09/10 (agent CCR-1109-D2), wired 2026-11-10 from scratchpad/ccr1109d2/ccr5-selfdual.scm. Library: complex-analysis.}\par
\vnbprop{sdl-ball-iff}{\vnbstmt{\begin{array}{@{}l@{}} \forall f.; \forall x \in \operatorname{carr}(f), r, z. \\ \quad \operatorname{is\text{-}normed\text{-}field}(f) \Rightarrow  \\ \quad \quad z \in \operatorname{closed\text{-}ball}(\operatorname{nf\text{-}metric\text{-}space}(f), x, r) \\ \quad \quad \Leftrightarrow\; (z \in \operatorname{carr}(f)) \wedge (\operatorname{fnrm}(f)(\operatorname{add}(f)(x, \operatorname{neg}(f)(z))) \leq r) \end{array}}}
\vnbprop{sdl-ball0-iff}{\vnbstmt{\begin{array}{@{}l@{}} \forall f, r, z. \\ \quad \operatorname{is\text{-}normed\text{-}field}(f) \Rightarrow  \\ \quad \quad (z \in \operatorname{closed\text{-}ball}(\operatorname{nf\text{-}metric\text{-}space}(f), \operatorname{zero}(f), r)) \Leftrightarrow ((z \in \operatorname{carr}(f)) \wedge (\operatorname{fnrm}(f)(z) \leq r)) \end{array}}}
\vnbprop{sdl-ball-eq}{\vnbstmt{\begin{array}{@{}l@{}} \forall f.; \forall x, y \in \operatorname{carr}(f), r \in \mathbb{R}. \\ \quad (\operatorname{is\text{-}nv\text{-}field}(f) \wedge ((0 \leq r) \wedge (\operatorname{fnrm}(f)(\operatorname{add}(f)(x, \operatorname{neg}(f)(y))) \leq r))) \Rightarrow  \\ \quad \quad \operatorname{closed\text{-}ball}(\operatorname{nf\text{-}metric\text{-}space}(f), x, r) = \operatorname{closed\text{-}ball}(\operatorname{nf\text{-}metric\text{-}space}(f), y, r) \end{array}}}
\vnbprop{sdl-cs-in}{\vnbstmt{\begin{array}{@{}l@{}} \forall f, s, b.; \forall x \in \operatorname{carr}(f). \\ \quad (\operatorname{is\text{-}normed\text{-}field}(f) \wedge (\operatorname{fnrm}(f)(x) \leq b)) \Rightarrow  \\ \quad \quad \operatorname{closed\text{-}ball}(\operatorname{nf\text{-}metric\text{-}space}(f), x, s) \\ \quad \quad \in\; \operatorname{image}((\lambda \mathit{sdx} \in \operatorname{closed\text{-}ball}(\operatorname{nf\text{-}metric\text{-}space}(f), \operatorname{zero}(f), b).\; \\ \quad \quad \quad \quad \operatorname{closed\text{-}ball}(\operatorname{nf\text{-}metric\text{-}space}(f), \mathit{sdx}, s)), \\ \quad \quad \quad \operatorname{closed\text{-}ball}(\operatorname{nf\text{-}metric\text{-}space}(f), \operatorname{zero}(f), b)) \end{array}}}
\vnbprop{sdl-cs-rep}{\vnbstmt{\begin{array}{@{}l@{}} \forall f.; \forall s, b \in \mathbb{R}, c \in \operatorname{image}((\lambda x \in \operatorname{closed\text{-}ball}(\operatorname{nf\text{-}metric\text{-}space}(f), \operatorname{zero}(f), b).\; \operatorname{closed\text{-}ball}(\operatorname{nf\text{-}metric\text{-}space}(f), x, s)), \operatorname{closed\text{-}ball}(\operatorname{nf\text{-}metric\text{-}space}(f), \operatorname{zero}(f), b)). \\ \quad (\operatorname{is\text{-}nv\text{-}field}(f) \wedge ((0 < s) \wedge (s \leq b))) \Rightarrow  \\ \quad \quad \varepsilon(c) \in \operatorname{carr}(f) \wedge \\ \quad \quad \operatorname{fnrm}(f)(\varepsilon(c)) \leq b \wedge \\ \quad \quad \varepsilon(c) \in c \wedge \\ \quad \quad c = \operatorname{closed\text{-}ball}(\operatorname{nf\text{-}metric\text{-}space}(f), \varepsilon(c), s) \end{array}}}
\vnbprop{sdl-nv-le}{\vnbstmt{\begin{array}{@{}l@{}} \forall f.; \forall x, y \in \operatorname{carr}(f), r \in \mathbb{R}. \\ \quad (\operatorname{is\text{-}nv\text{-}field}(f) \wedge ((\operatorname{fnrm}(f)(x) \leq r) \wedge (\operatorname{fnrm}(f)(y) \leq r))) \Rightarrow  \\ \quad \quad \operatorname{fnrm}(f)(\operatorname{add}(f)(x, y)) \leq r \end{array}}}
\vnbprop{sdl-shift-in}{\vnbstmt{\begin{array}{@{}l@{}} \forall f.; \forall s, b \in \mathbb{R}, c \in \operatorname{image}((\lambda x \in \operatorname{closed\text{-}ball}(\operatorname{nf\text{-}metric\text{-}space}(f), \operatorname{zero}(f), b).\; \operatorname{closed\text{-}ball}(\operatorname{nf\text{-}metric\text{-}space}(f), x, s)), \operatorname{closed\text{-}ball}(\operatorname{nf\text{-}metric\text{-}space}(f), \operatorname{zero}(f), b)), t \in \operatorname{carr}(f). \\ \quad (\operatorname{is\text{-}nv\text{-}field}(f) \wedge ((0 < s) \wedge ((s \leq b) \wedge (\operatorname{fnrm}(f)(t) \leq b)))) \Rightarrow  \\ \quad \quad \operatorname{closed\text{-}ball}(\operatorname{nf\text{-}metric\text{-}space}(f), \operatorname{add}(f)(\varepsilon(c), t), s) \\ \quad \quad \in\; \operatorname{image}((\lambda x \in \operatorname{closed\text{-}ball}(\operatorname{nf\text{-}metric\text{-}space}(f), \operatorname{zero}(f), b).\; \operatorname{closed\text{-}ball}(\operatorname{nf\text{-}metric\text{-}space}(f), x, s)), \\ \quad \quad \quad \operatorname{closed\text{-}ball}(\operatorname{nf\text{-}metric\text{-}space}(f), \operatorname{zero}(f), b)) \end{array}}}
\vnbprop{sdl-shift-back}{\vnbstmt{\begin{array}{@{}l@{}} \forall f.; \forall s, b \in \mathbb{R}, c \in \operatorname{image}((\lambda x \in \operatorname{closed\text{-}ball}(\operatorname{nf\text{-}metric\text{-}space}(f), \operatorname{zero}(f), b).\; \operatorname{closed\text{-}ball}(\operatorname{nf\text{-}metric\text{-}space}(f), x, s)), \operatorname{closed\text{-}ball}(\operatorname{nf\text{-}metric\text{-}space}(f), \operatorname{zero}(f), b)), t \in \operatorname{carr}(f). \\ \quad (\operatorname{is\text{-}nv\text{-}field}(f) \wedge ((0 < s) \wedge ((s \leq b) \wedge (\operatorname{fnrm}(f)(t) \leq b)))) \Rightarrow  \\ \quad \quad \operatorname{closed\text{-}ball}(\operatorname{nf\text{-}metric\text{-}space}(f), \\ \quad \quad \quad \operatorname{add}(f)(\varepsilon(\operatorname{closed\text{-}ball}(\operatorname{nf\text{-}metric\text{-}space}(f), \operatorname{add}(f)(\varepsilon(c), t), s)), \operatorname{neg}(f)(t)), \\ \quad \quad \quad s) \\ \quad \quad =\; c \end{array}}}
\vnbprop{sdl-cc-fix}{\vnbstmt{\forall z, w \in \mathbb{C}.\; ((\neg (z = 1) \wedge (w = (z \cdot w))) \Rightarrow (w = 0))}}
\vnbprop{sdl-orth}{\vnbstmt{\begin{array}{@{}l@{}} \forall a.; \forall s, b \in \mathbb{R}, g.; \forall f \in (\operatorname{image}((\lambda \mathit{sdx} \in \operatorname{closed\text{-}ball}(\operatorname{nf\text{-}metric\text{-}space}(a), \operatorname{zero}(a), b).\; \operatorname{closed\text{-}ball}(\operatorname{nf\text{-}metric\text{-}space}(a), \mathit{sdx}, s)), \operatorname{closed\text{-}ball}(\operatorname{nf\text{-}metric\text{-}space}(a), \operatorname{zero}(a), b)) \to \mathbb{C}), h \in \operatorname{carr}(a). \\ \quad \operatorname{is\text{-}nv\text{-}field}(a) \wedge \\ \quad 0 < s \wedge \\ \quad s \leq b \wedge \\ \quad |\operatorname{image}((\lambda \mathit{sdx} \in \operatorname{closed\text{-}ball}(\operatorname{nf\text{-}metric\text{-}space}(a), \operatorname{zero}(a), b).\; \operatorname{closed\text{-}ball}(\operatorname{nf\text{-}metric\text{-}space}(a), \mathit{sdx}, s)), \operatorname{closed\text{-}ball}(\operatorname{nf\text{-}metric\text{-}space}(a), \operatorname{zero}(a), b))| \\ \quad \in\; \mathbb{N} \wedge \\ \quad \forall x \in \operatorname{carr}(a).\; (g(x) \in \mathbb{C}) \wedge \\ \quad \forall x, y \in \operatorname{carr}(a).\; (g(\operatorname{add}(a)(x, y)) = (g(x) \cdot g(y))) \wedge \\ \quad \forall y \in \operatorname{carr}(a).\; ((\operatorname{fnrm}(a)(y) \leq s) \Rightarrow (g(y) = 1)) \wedge \\ \quad \forall c \in \operatorname{image}((\lambda \mathit{sdx} \in \operatorname{closed\text{-}ball}(\operatorname{nf\text{-}metric\text{-}space}(a), \operatorname{zero}(a), b).\; \operatorname{closed\text{-}ball}(\operatorname{nf\text{-}metric\text{-}space}(a), \mathit{sdx}, s)), \operatorname{closed\text{-}ball}(\operatorname{nf\text{-}metric\text{-}space}(a), \operatorname{zero}(a), b)). \\ \quad \quad f(c) = g(\varepsilon(c)) \wedge \\ \quad \operatorname{fnrm}(a)(h) \leq b \wedge \\ \quad \neg (g(h) = 1) \Rightarrow \\ \quad \quad \operatorname{finsum}(\operatorname{ring\text{-}additive\text{-}ag}(\operatorname{normed\text{-}field\text{-}as\text{-}commutative\text{-}ring}(\operatorname{cc\text{-}normed\text{-}field})), \\ \quad \quad \quad f, \\ \quad \quad \quad \operatorname{image}((\lambda \mathit{sdx} \in \operatorname{closed\text{-}ball}(\operatorname{nf\text{-}metric\text{-}space}(a), \operatorname{zero}(a), b).\; \\ \quad \quad \quad \quad \quad \operatorname{closed\text{-}ball}(\operatorname{nf\text{-}metric\text{-}space}(a), \mathit{sdx}, s)), \\ \quad \quad \quad \quad \operatorname{closed\text{-}ball}(\operatorname{nf\text{-}metric\text{-}space}(a), \operatorname{zero}(a), b))) \\ \quad \quad =\; 0 \end{array}}}
\vnbprop{sdl-sum-one}{\vnbstmt{\begin{array}{@{}l@{}} \forall n \in \mathbb{N}, f.; \forall s \in \mathbf{Set}. \\ \quad ((|s| = n) \wedge \forall z \in s.\; (f(z) = 1)) \Rightarrow  \\ \quad \quad \operatorname{finsum}(\operatorname{ring\text{-}additive\text{-}ag}(\operatorname{normed\text{-}field\text{-}as\text{-}commutative\text{-}ring}(\operatorname{cc\text{-}normed\text{-}field})), f, s) = n \end{array}}}
\vnbprop{sdl-cs-finite}{\vnbstmt{\begin{array}{@{}l@{}} \forall f, w.; \forall s, b \in \mathbb{R}. \\ \quad (\operatorname{is\text{-}local\text{-}uniformizer}(f, w) \wedge ((0 < s) \wedge (s \leq b))) \Rightarrow  \\ \quad \quad |\operatorname{image}((\lambda x \in \operatorname{closed\text{-}ball}(\operatorname{nf\text{-}metric\text{-}space}(f), \operatorname{zero}(f), b).\; \operatorname{closed\text{-}ball}(\operatorname{nf\text{-}metric\text{-}space}(f), x, s)), \operatorname{closed\text{-}ball}(\operatorname{nf\text{-}metric\text{-}space}(f), \operatorname{zero}(f), b))| \\ \quad \quad \in\; \mathbb{N} \end{array}}}
\vnbprop{sdl-prod1}{\vnbstmt{\begin{array}{@{}l@{}} \forall f.; \forall u, v, p, q \in \operatorname{carr}(f). \\ \quad \operatorname{is\text{-}normed\text{-}field}(f) \wedge \\ \quad \operatorname{mul}(f)(u, v) = \operatorname{one}(f) \wedge \\ \quad \operatorname{fnrm}(f)(p) \leq \operatorname{fnrm}(f)(u) \wedge \\ \quad \operatorname{fnrm}(f)(q) \leq \operatorname{fnrm}(f)(v) \Rightarrow \\ \quad \quad \operatorname{fnrm}(f)(\operatorname{mul}(f)(p, q)) \leq 1 \end{array}}}
\vnbprop{sdl-scale-gt}{\vnbstmt{\begin{array}{@{}l@{}} \forall f.; \forall u, v, x \in \operatorname{carr}(f). \\ \quad (\operatorname{is\text{-}normed\text{-}field}(f) \wedge ((\operatorname{mul}(f)(u, v) = \operatorname{one}(f)) \wedge ((0 < \operatorname{fnrm}(f)(u)) \wedge (\operatorname{fnrm}(f)(v) < \operatorname{fnrm}(f)(x))))) \Rightarrow  \\ \quad \quad 1 < \operatorname{fnrm}(f)(\operatorname{mul}(f)(u, x)) \end{array}}}
\vnbprop{sdl-orth-a}{\vnbstmt{\begin{array}{@{}l@{}} \forall b.; \forall s, d \in \mathbb{R}, g.; \forall a \in \mathbb{C}, f \in (\operatorname{image}((\lambda \mathit{sdx} \in \operatorname{closed\text{-}ball}(\operatorname{nf\text{-}metric\text{-}space}(b), \operatorname{zero}(b), d).\; \operatorname{closed\text{-}ball}(\operatorname{nf\text{-}metric\text{-}space}(b), \mathit{sdx}, s)), \operatorname{closed\text{-}ball}(\operatorname{nf\text{-}metric\text{-}space}(b), \operatorname{zero}(b), d)) \to \mathbb{C}), h \in \operatorname{carr}(b). \\ \quad \operatorname{is\text{-}nv\text{-}field}(b) \wedge \\ \quad 0 < s \wedge \\ \quad s \leq d \wedge \\ \quad |\operatorname{image}((\lambda \mathit{sdx} \in \operatorname{closed\text{-}ball}(\operatorname{nf\text{-}metric\text{-}space}(b), \operatorname{zero}(b), d).\; \operatorname{closed\text{-}ball}(\operatorname{nf\text{-}metric\text{-}space}(b), \mathit{sdx}, s)), \operatorname{closed\text{-}ball}(\operatorname{nf\text{-}metric\text{-}space}(b), \operatorname{zero}(b), d))| \\ \quad \in\; \mathbb{N} \wedge \\ \quad \forall x \in \operatorname{carr}(b).\; (g(x) \in \mathbb{C}) \wedge \\ \quad \forall x, y \in \operatorname{carr}(b).\; (g(\operatorname{add}(b)(x, y)) = (g(x) \cdot g(y))) \wedge \\ \quad \forall y \in \operatorname{carr}(b).\; ((\operatorname{fnrm}(b)(y) \leq s) \Rightarrow (g(y) = 1)) \wedge \\ \quad \forall c \in \operatorname{image}((\lambda \mathit{sdx} \in \operatorname{closed\text{-}ball}(\operatorname{nf\text{-}metric\text{-}space}(b), \operatorname{zero}(b), d).\; \operatorname{closed\text{-}ball}(\operatorname{nf\text{-}metric\text{-}space}(b), \mathit{sdx}, s)), \operatorname{closed\text{-}ball}(\operatorname{nf\text{-}metric\text{-}space}(b), \operatorname{zero}(b), d)). \\ \quad \quad f(c) = (a \cdot g(\varepsilon(c))) \wedge \\ \quad \operatorname{fnrm}(b)(h) \leq d \wedge \\ \quad \neg (g(h) = 1) \Rightarrow \\ \quad \quad \operatorname{finsum}(\operatorname{ring\text{-}additive\text{-}ag}(\operatorname{normed\text{-}field\text{-}as\text{-}commutative\text{-}ring}(\operatorname{cc\text{-}normed\text{-}field})), \\ \quad \quad \quad f, \\ \quad \quad \quad \operatorname{image}((\lambda \mathit{sdx} \in \operatorname{closed\text{-}ball}(\operatorname{nf\text{-}metric\text{-}space}(b), \operatorname{zero}(b), d).\; \\ \quad \quad \quad \quad \quad \operatorname{closed\text{-}ball}(\operatorname{nf\text{-}metric\text{-}space}(b), \mathit{sdx}, s)), \\ \quad \quad \quad \quad \operatorname{closed\text{-}ball}(\operatorname{nf\text{-}metric\text{-}space}(b), \operatorname{zero}(b), d))) \\ \quad \quad =\; 0 \end{array}}}
\vnbprop{sdl-le-mul1}{\vnbstmt{\begin{array}{@{}l@{}} \forall f.; \forall u, x \in \operatorname{carr}(f). \\ \quad (\operatorname{is\text{-}normed\text{-}field}(f) \wedge (\operatorname{fnrm}(f)(x) \leq 1)) \Rightarrow  \\ \quad \quad \operatorname{fnrm}(f)(\operatorname{mul}(f)(u, x)) \leq \operatorname{fnrm}(f)(u) \end{array}}}
\vnbprop{sdl-cc-inv}{\vnbstmt{\forall a, b, v \in \mathbb{C}.\; ((((a \cdot b) = 1) \wedge ((b \cdot v) = 1)) \Rightarrow (a = v))}}
\vnbprop{sdl-dsum}{\vnbstmt{\begin{array}{@{}l@{}} \forall f, w, i, u.; \forall c, a, g, b \in \operatorname{carr}(f). \\ \quad \operatorname{is\text{-}local\text{-}uniformizer}(f, w) \wedge \\ \quad \operatorname{is\text{-}add\text{-}character}(f, i) \wedge \\ \quad \forall m \in \operatorname{carr}(f).\; ((\operatorname{fnrm}(f)(m) \leq 1) \Rightarrow (i(m) = 1)) \wedge \\ \quad \forall h \in \operatorname{carr}(f). \\ \quad \quad (1 < \operatorname{fnrm}(f)(h)) \Rightarrow  \\ \quad \quad \quad \exists x \in \operatorname{carr}(f).\; ((\operatorname{fnrm}(f)(x) \leq 1) \wedge \neg (i(\operatorname{mul}(f)(h, x)) = 1)) \wedge \\ \quad \operatorname{is\text{-}add\text{-}character}(f, u) \wedge \\ \quad \operatorname{mul}(f)(c, a) = \operatorname{one}(f) \wedge \\ \quad 0 < \operatorname{fnrm}(f)(c) \wedge \\ \quad \operatorname{fnrm}(f)(c) \leq 1 \wedge \\ \quad 1 \leq \operatorname{fnrm}(f)(a) \wedge \\ \quad \forall y \in \operatorname{carr}(f).\; ((\operatorname{fnrm}(f)(y) \leq \operatorname{fnrm}(f)(c)) \Rightarrow (u(y) = 1)) \wedge \\ \quad \operatorname{mul}(f)(g, b) = \operatorname{one}(f) \wedge \\ \quad 0 < \operatorname{fnrm}(f)(b) \wedge \\ \quad \operatorname{fnrm}(f)(b) \leq 1 \wedge \\ \quad 1 \leq \operatorname{fnrm}(f)(g) \Rightarrow \\ \quad \quad \exists p \in \operatorname{carr}(f). \forall q \in \operatorname{carr}(f). \\ \quad \quad \quad (\operatorname{fnrm}(f)(q) \leq \operatorname{fnrm}(f)(g)) \Rightarrow (u(q) = i(\operatorname{mul}(f)(p, q))) \end{array}}}
\vnbprop{sdl-inv-ge1}{\vnbstmt{\begin{array}{@{}l@{}} \forall f.; \forall u, v \in \operatorname{carr}(f). \\ \quad (\operatorname{is\text{-}normed\text{-}field}(f) \wedge ((\operatorname{mul}(f)(u, v) = \operatorname{one}(f)) \wedge (\operatorname{fnrm}(f)(u) \leq 1))) \Rightarrow  \\ \quad \quad 1 \leq \operatorname{fnrm}(f)(v) \end{array}}}
\vnbprop{sdl-pow}{\vnbstmt{\begin{array}{@{}l@{}} \forall f.; \forall u, v \in \operatorname{carr}(f), k \in \mathbb{N}. \\ \quad (\operatorname{is\text{-}normed\text{-}field}(f) \wedge ((\operatorname{mul}(f)(u, v) = \operatorname{one}(f)) \wedge ((0 < \operatorname{fnrm}(f)(u)) \wedge (\operatorname{fnrm}(f)(u) \leq 1)))) \Rightarrow  \\ \quad \quad \operatorname{ring\text{-}power}(\operatorname{normed\text{-}field\text{-}as\text{-}commutative\text{-}ring}(f), u, k) \in \operatorname{carr}(f) \wedge \\ \quad \quad \operatorname{ring\text{-}power}(\operatorname{normed\text{-}field\text{-}as\text{-}commutative\text{-}ring}(f), v, k) \in \operatorname{carr}(f) \wedge \\ \quad \quad \operatorname{mul}(f)(\operatorname{ring\text{-}power}(\operatorname{normed\text{-}field\text{-}as\text{-}commutative\text{-}ring}(f), u, k), \\ \quad \quad \quad \operatorname{ring\text{-}power}(\operatorname{normed\text{-}field\text{-}as\text{-}commutative\text{-}ring}(f), v, k)) \\ \quad \quad =\; \operatorname{one}(f) \wedge \\ \quad \quad \operatorname{fnrm}(f)(\operatorname{ring\text{-}power}(\operatorname{normed\text{-}field\text{-}as\text{-}commutative\text{-}ring}(f), u, k)) = \operatorname{power}(\operatorname{fnrm}(f)(u), k) \wedge \\ \quad \quad \operatorname{fnrm}(f)(\operatorname{ring\text{-}power}(\operatorname{normed\text{-}field\text{-}as\text{-}commutative\text{-}ring}(f), v, k)) = \operatorname{power}(\operatorname{fnrm}(f)(v), k) \wedge \\ \quad \quad 0 < \operatorname{fnrm}(f)(\operatorname{ring\text{-}power}(\operatorname{normed\text{-}field\text{-}as\text{-}commutative\text{-}ring}(f), u, k)) \wedge \\ \quad \quad \operatorname{fnrm}(f)(\operatorname{ring\text{-}power}(\operatorname{normed\text{-}field\text{-}as\text{-}commutative\text{-}ring}(f), u, k)) \leq 1 \wedge \\ \quad \quad 1 \leq \operatorname{fnrm}(f)(\operatorname{ring\text{-}power}(\operatorname{normed\text{-}field\text{-}as\text{-}commutative\text{-}ring}(f), v, k)) \end{array}}}
\vnbprop{sdl-core}{\vnbstmt{\begin{array}{@{}l@{}} \forall f, w, i, u. \\ \quad \operatorname{is\text{-}local\text{-}uniformizer}(f, w) \wedge \\ \quad \operatorname{is\text{-}add\text{-}character}(f, i) \wedge \\ \quad \forall m \in \operatorname{carr}(f).\; ((\operatorname{fnrm}(f)(m) \leq 1) \Rightarrow (i(m) = 1)) \wedge \\ \quad \forall h \in \operatorname{carr}(f). \\ \quad \quad (1 < \operatorname{fnrm}(f)(h)) \Rightarrow  \\ \quad \quad \quad \exists x \in \operatorname{carr}(f).\; ((\operatorname{fnrm}(f)(x) \leq 1) \wedge \neg (i(\operatorname{mul}(f)(h, x)) = 1)) \wedge \\ \quad \operatorname{is\text{-}add\text{-}character}(f, u) \Rightarrow \\ \quad \quad \exists a \in \operatorname{carr}(f).\; \forall x \in \operatorname{carr}(f).\; (u(x) = i(\operatorname{mul}(f)(a, x))) \end{array}}}
\vnbprop{sdl-unipolar-phi}{\vnbstmt{\begin{array}{@{}l@{}} \forall f, a, w. \\ \quad (\operatorname{is\text{-}local\text{-}uniformizer}(f, a) \wedge \operatorname{is\text{-}unipolar}(f, w)) \Rightarrow  \\ \quad \quad \operatorname{is\text{-}add\text{-}character}(f, (\lambda \mathit{chx} \in \operatorname{carr}(f).\; w(\operatorname{one}(f), \mathit{chx}))) \wedge \\ \quad \quad \forall m \in \operatorname{carr}(f). \\ \quad \quad \quad (\operatorname{fnrm}(f)(m) \leq 1) \Rightarrow ((\lambda \mathit{chx} \in \operatorname{carr}(f).\; w(\operatorname{one}(f), \mathit{chx}))(m) = 1) \wedge \\ \quad \quad \forall h \in \operatorname{carr}(f). \\ \quad \quad \quad (1 < \operatorname{fnrm}(f)(h)) \Rightarrow  \\ \quad \quad \quad \quad \exists x \in \operatorname{carr}(f). \\ \quad \quad \quad \quad \quad ((\operatorname{fnrm}(f)(x) \leq 1) \wedge \neg ((\lambda \mathit{chx} \in \operatorname{carr}(f).\; w(\operatorname{one}(f), \mathit{chx}))(\operatorname{mul}(f)(h, x)) = 1)) \end{array}}}
\vnbprop{ccr-5-13}{\vnbstmt{\begin{array}{@{}l@{}} \forall f, w, i, u. \\ \quad (\operatorname{is\text{-}local\text{-}uniformizer}(f, w) \wedge (\operatorname{is\text{-}nontrivial\text{-}character}(f, i) \wedge \operatorname{is\text{-}add\text{-}character}(f, u))) \Rightarrow  \\ \quad \quad \exists a \in \operatorname{carr}(f).\; (u = \operatorname{duality\text{-}map}(f, i, a)) \end{array}}}
\vnbprop{ccr-5-13-dualizing}{\vnbstmt{\begin{array}{@{}l@{}} \forall f, w, i. \\ \quad (\operatorname{is\text{-}local\text{-}uniformizer}(f, w) \wedge \operatorname{is\text{-}nontrivial\text{-}character}(f, i)) \Rightarrow  \\ \quad \quad \operatorname{is\text{-}dualizing\text{-}character}(f, i) \end{array}}}
\vnbprop{ccr-6-3-converse-local}{\vnbstmt{\begin{array}{@{}l@{}} \forall f, a, i, w. \\ \quad (\operatorname{is\text{-}local\text{-}uniformizer}(f, a) \wedge (\operatorname{is\text{-}nontrivial\text{-}character}(f, i) \wedge \operatorname{is\text{-}bilinear\text{-}bicharacter}(f, w))) \Rightarrow  \\ \quad \quad \exists b \in \operatorname{carr}(f). \forall x, c \in \operatorname{carr}(f). \\ \quad \quad \quad w(x, c) = \operatorname{bichar\text{-}of}(f, i, b)(x, c) \end{array}}}
\vnbprop{ccr-6-6-local}{\vnbstmt{\begin{array}{@{}l@{}} \forall f, a, i, w. \\ \quad \operatorname{is\text{-}local\text{-}uniformizer}(f, a) \wedge \\ \quad \operatorname{is\text{-}nontrivial\text{-}character}(f, i) \wedge \\ \quad \operatorname{is\text{-}bilinear\text{-}bicharacter}(f, w) \wedge \\ \quad \exists x, y \in \operatorname{carr}(f).\; \neg (w(x, y) = 1) \Rightarrow \\ \quad \quad \operatorname{is\text{-}faithful\text{-}bicharacter}(f, w) \end{array}}}
\vnbprop{ccr-6-20-local}{\vnbstmt{\begin{array}{@{}l@{}} \forall f, a, i.; \forall c \in \operatorname{carr}(f), w. \\ \quad \operatorname{is\text{-}local\text{-}uniformizer}(f, a) \wedge \\ \quad \operatorname{is\text{-}nontrivial\text{-}character}(f, i) \wedge \\ \quad \neg (c = \operatorname{zero}(f)) \wedge \\ \quad \operatorname{is\text{-}bilinear\text{-}bicharacter}(f, w) \Rightarrow \\ \quad \quad \exists b \in \operatorname{carr}(f). \forall x, d \in \operatorname{carr}(f). \\ \quad \quad \quad w(x, d) = \operatorname{bichar\text{-}of}(f, i, c)(\operatorname{mul}(f)(b, x), d) \wedge \\ \quad \quad \forall b, g \in \operatorname{carr}(f). \\ \quad \quad \quad \forall x, d \in \operatorname{carr}(f). \\ \quad \quad \quad \quad w(x, d) = \operatorname{bichar\text{-}of}(f, i, c)(\operatorname{mul}(f)(b, x), d) \wedge \\ \quad \quad \quad \forall x, d \in \operatorname{carr}(f). \\ \quad \quad \quad \quad w(x, d) = \operatorname{bichar\text{-}of}(f, i, c)(\operatorname{mul}(f)(g, x), d) \Rightarrow \\ \quad \quad \quad \quad b = g \end{array}}}
\vnbprop{ccr-6-23-local}{\vnbstmt{\begin{array}{@{}l@{}} \forall f, a, i, w, q. \\ \quad \operatorname{is\text{-}local\text{-}uniformizer}(f, a) \wedge \\ \quad \operatorname{is\text{-}nontrivial\text{-}character}(f, i) \wedge \\ \quad \operatorname{is\text{-}unipolar}(f, w) \wedge \\ \quad \operatorname{is\text{-}unipolar}(f, q) \Rightarrow \\ \quad \quad \exists b \in \operatorname{carr}(f). \\ \quad \quad \quad ((\operatorname{fnrm}(f)(b) = 1) \wedge \forall x, c \in \operatorname{carr}(f).\; (w(x, c) = q(\operatorname{mul}(f)(b, x), c))) \wedge \\ \quad \quad \forall b, d \in \operatorname{carr}(f). \\ \quad \quad \quad \forall x, c \in \operatorname{carr}(f).\; (w(x, c) = q(\operatorname{mul}(f)(b, x), c)) \wedge \\ \quad \quad \quad \forall x, c \in \operatorname{carr}(f).\; (w(x, c) = q(\operatorname{mul}(f)(d, x), c)) \Rightarrow \\ \quad \quad \quad \quad b = d \end{array}}}
\vnbprop{ccr-6-18}{\vnbstmt{\begin{array}{@{}l@{}} \forall f, w, i. \\ \quad (\operatorname{is\text{-}local\text{-}uniformizer}(f, w) \wedge \operatorname{is\text{-}nontrivial\text{-}character}(f, i)) \Rightarrow  \\ \quad \quad \exists \mathit{sdw}.\; (\operatorname{is\text{-}unipolar}(f, \mathit{sdw}) \wedge \operatorname{is\text{-}faithful\text{-}bicharacter}(f, \mathit{sdw})) \end{array}}}
\vnbprop{ccr-6-22-forward}{\vnbstmt{\begin{array}{@{}l@{}} \forall f, a, w.; \forall q \in (\mathbb{N} \to \operatorname{carr}(f)), i, b \in \operatorname{carr}(f). \\ \quad (\operatorname{is\text{-}local\text{-}uniformizer}(f, a) \wedge (\operatorname{is\text{-}unipolar}(f, w) \wedge \operatorname{converges\text{-}to}(\operatorname{nf\text{-}metric\text{-}space}(f), q, i))) \Rightarrow  \\ \quad \quad \operatorname{converges\text{-}to}(\operatorname{cc\text{-}ms}, (\lambda n \in \mathbb{N}.\; w(q(n), b)), w(i, b)) \end{array}}}

\vnbbookfile{theorem-library/ccr1109a5-pairing}{theorem-library/ccr1109a5-pairing}
\noindent{\small ccr1109a5-pairing.scm -- wave of 2026-11-09/10 (agent CCR-1109-A5), wired 2026-11-10 from scratchpad/ccr1109a5/ccr1109a5-pairing.scm. Library: complex-analysis.}\par
\vnbprop{nv-maximal-principal-of-gap}{\vnbstmt{\begin{array}{@{}l@{}} \forall f.; \forall w \in \operatorname{carr}(f). \\ \quad \operatorname{is\text{-}nv\text{-}field}(f) \wedge \\ \quad 0 < \operatorname{fnrm}(f)(w) \wedge \\ \quad \operatorname{fnrm}(f)(w) < 1 \wedge \\ \quad \forall h \in \operatorname{carr}(f).\; ((\operatorname{fnrm}(f)(h) < 1) \Rightarrow (\operatorname{fnrm}(f)(h) \leq \operatorname{fnrm}(f)(w))) \Rightarrow \\ \quad \quad \operatorname{nv\text{-}maximal}(f, \operatorname{nv\text{-}integers}(f)) = \operatorname{nv\text{-}principal}(f, \operatorname{nv\text{-}integers}(f), w) \end{array}}}
\vnbprop{nv-local-uniformizer-of-gap}{\vnbstmt{\begin{array}{@{}l@{}} \forall f.; \forall w \in \operatorname{carr}(f). \\ \quad \operatorname{is\text{-}nv\text{-}field}(f) \wedge \\ \quad \operatorname{is\text{-}compact}(\operatorname{subspace\text{-}ms}(\operatorname{nf\text{-}metric\text{-}space}(f), \operatorname{nv\text{-}integers}(f))) \wedge \\ \quad 0 < \operatorname{fnrm}(f)(w) \wedge \\ \quad \operatorname{fnrm}(f)(w) < 1 \wedge \\ \quad \forall h \in \operatorname{carr}(f).\; ((\operatorname{fnrm}(f)(h) < 1) \Rightarrow (\operatorname{fnrm}(f)(h) \leq \operatorname{fnrm}(f)(w))) \Rightarrow \\ \quad \quad \operatorname{is\text{-}local\text{-}uniformizer}(f, w) \end{array}}}
\vnbprop{laurent-abs-lt-one}{\vnbstmt{\begin{array}{@{}l@{}} \forall q.; \forall \mathit{rq} \in \mathbb{R}, \mathit{fq} \in \operatorname{laurent\text{-}set}(q). \\ \quad (\operatorname{is\text{-}field\text{-}ring}(q) \wedge ((0 < \mathit{rq}) \wedge ((\mathit{rq} < 1) \wedge (\operatorname{laurent\text{-}abs}(q, \mathit{rq}, \mathit{fq}) < 1)))) \Rightarrow  \\ \quad \quad \operatorname{laurent\text{-}abs}(q, \mathit{rq}, \mathit{fq}) \leq \mathit{rq} \end{array}}}
\vnbprop{laurent-local-uniformizer}{\vnbstmt{\begin{array}{@{}l@{}} \forall q.; \forall \mathit{rq} \in \mathbb{R}. \\ \quad (\operatorname{is\text{-}field\text{-}ring}(q) \wedge ((|\operatorname{carr}(q)| \in \mathbb{N}) \wedge ((0 < \mathit{rq}) \wedge (\mathit{rq} < 1)))) \Rightarrow  \\ \quad \quad \operatorname{is\text{-}local\text{-}uniformizer}(\operatorname{laurent\text{-}nf}(q, \mathit{rq}), \operatorname{laurent\text{-}mono}(q, \operatorname{one}(q), 1)) \end{array}}}
\vnbprop{lp-coeff-in}{\vnbstmt{\forall q.; \forall \mathit{fq} \in \operatorname{laurent\text{-}set}(q), \mathit{jq} \in \mathbb{Z}.\; (\operatorname{fq}(\mathit{jq}) \in \operatorname{carr}(q))}}
\vnbprop{lp-nf-add-apply}{\vnbstmt{\begin{array}{@{}l@{}} \forall q.; \forall \mathit{rq} \in \mathbb{R}, \mathit{fq}, \mathit{gq} \in \operatorname{laurent\text{-}set}(q). \\ \quad (\operatorname{is\text{-}field\text{-}ring}(q) \wedge ((0 < \mathit{rq}) \wedge (\mathit{rq} < 1))) \Rightarrow  \\ \quad \quad \operatorname{add}(\operatorname{laurent\text{-}nf}(q, \mathit{rq}))(\mathit{fq}, \mathit{gq}) = \operatorname{laurent\text{-}add}(q, \mathit{fq}, \mathit{gq}) \end{array}}}
\vnbprop{lp-nf-mul-apply}{\vnbstmt{\begin{array}{@{}l@{}} \forall q.; \forall \mathit{rq} \in \mathbb{R}, \mathit{fq}, \mathit{gq} \in \operatorname{laurent\text{-}set}(q). \\ \quad (\operatorname{is\text{-}field\text{-}ring}(q) \wedge ((0 < \mathit{rq}) \wedge (\mathit{rq} < 1))) \Rightarrow  \\ \quad \quad \operatorname{mul}(\operatorname{laurent\text{-}nf}(q, \mathit{rq}))(\mathit{fq}, \mathit{gq}) = \operatorname{laurent\text{-}mul}(q, \mathit{fq}, \mathit{gq}) \end{array}}}
\vnbprop{lp-value}{\vnbstmt{\begin{array}{@{}l@{}} \forall q, c.; \forall \mathit{uq}, \mathit{vq} \in \operatorname{laurent\text{-}set}(q). \\ \quad (\operatorname{laurent\text{-}pairing}(q, c)(\mathit{uq}, \mathit{vq}) \simeq c(\operatorname{laurent\text{-}mul}(q, \mathit{vq}, \mathit{uq})(0))) \end{array}}}
\vnbprop{lp-section-type}{\vnbstmt{\begin{array}{@{}l@{}} \forall q.; \forall \mathit{rq} \in \mathbb{R}, c \in (\operatorname{carr}(q) \to \operatorname{unit\text{-}circle}), \mathit{yq} \in \operatorname{laurent\text{-}set}(q). \\ \quad (\operatorname{is\text{-}field\text{-}ring}(q) \wedge ((0 < \mathit{rq}) \wedge (\mathit{rq} < 1))) \Rightarrow  \\ \quad \quad (\lambda z \in \operatorname{carr}(\operatorname{laurent\text{-}nf}(q, \mathit{rq})).\; c(\operatorname{laurent\text{-}mul}(q, \mathit{yq}, z)(0))) \\ \quad \quad \in\; (\operatorname{carr}(\operatorname{laurent\text{-}nf}(q, \mathit{rq})) \to \operatorname{unit\text{-}circle}) \end{array}}}
\vnbprop{lp-section-character}{\vnbstmt{\begin{array}{@{}l@{}} \forall q.; \forall \mathit{rq} \in \mathbb{R}, c \in (\operatorname{carr}(q) \to \operatorname{unit\text{-}circle}), \mathit{yq} \in \operatorname{laurent\text{-}set}(q). \\ \quad \operatorname{is\text{-}field\text{-}ring}(q) \wedge \\ \quad 0 < \mathit{rq} \wedge \\ \quad \mathit{rq} < 1 \wedge \\ \quad \forall x, y \in \operatorname{carr}(q).\; (c(\operatorname{add}(q)(x, y)) = (c(x) \cdot c(y))) \Rightarrow \\ \quad \quad \operatorname{is\text{-}add\text{-}character}(\operatorname{laurent\text{-}nf}(q, \mathit{rq}), \\ \quad \quad \quad (\lambda z \in \operatorname{carr}(\operatorname{laurent\text{-}nf}(q, \mathit{rq})).\; c(\operatorname{laurent\text{-}mul}(q, \mathit{yq}, z)(0)))) \end{array}}}
\vnbprop{lp-bilinear-symmetric}{\vnbstmt{\begin{array}{@{}l@{}} \forall q.; \forall \mathit{rq} \in \mathbb{R}, c \in (\operatorname{carr}(q) \to \operatorname{unit\text{-}circle}). \\ \quad \operatorname{is\text{-}field\text{-}ring}(q) \wedge \\ \quad 0 < \mathit{rq} \wedge \\ \quad \mathit{rq} < 1 \wedge \\ \quad \forall x, y \in \operatorname{carr}(q).\; (c(\operatorname{add}(q)(x, y)) = (c(x) \cdot c(y))) \Rightarrow \\ \quad \quad \operatorname{is\text{-}bilinear\text{-}bicharacter}(\operatorname{laurent\text{-}nf}(q, \mathit{rq}), \operatorname{laurent\text{-}pairing}(q, c)) \wedge \\ \quad \quad \forall x, y \in \operatorname{carr}(\operatorname{laurent\text{-}nf}(q, \mathit{rq})). \\ \quad \quad \quad \operatorname{laurent\text{-}pairing}(q, c)(x, y) = \operatorname{laurent\text{-}pairing}(q, c)(y, x) \end{array}}}
\vnbprop{lp-triv-char-parts}{\vnbstmt{\begin{array}{@{}l@{}} \forall q, c. \\ \quad (\operatorname{is\text{-}ring}(q) \wedge \operatorname{is\text{-}nontrivial\text{-}character}(\operatorname{triv\text{-}nf}(q), c)) \Rightarrow  \\ \quad \quad c \in (\operatorname{carr}(q) \to \operatorname{unit\text{-}circle}) \wedge \\ \quad \quad \forall x, y \in \operatorname{carr}(q).\; (c(\operatorname{add}(q)(x, y)) = (c(x) \cdot c(y))) \wedge \\ \quad \quad \exists x \in \operatorname{carr}(q).\; \neg (c(x) = 1) \end{array}}}
\vnbprop{lp-nontrivial}{\vnbstmt{\begin{array}{@{}l@{}} \forall q.; \forall \mathit{rq} \in \mathbb{R}, c \in (\operatorname{carr}(q) \to \operatorname{unit\text{-}circle}). \\ \quad \operatorname{is\text{-}field\text{-}ring}(q) \wedge \\ \quad 0 < \mathit{rq} \wedge \\ \quad \mathit{rq} < 1 \wedge \\ \quad \forall x, y \in \operatorname{carr}(q).\; (c(\operatorname{add}(q)(x, y)) = (c(x) \cdot c(y))) \wedge \\ \quad \exists x \in \operatorname{carr}(q).\; \neg (c(x) = 1) \Rightarrow \\ \quad \quad \operatorname{is\text{-}nontrivial\text{-}character}(\operatorname{laurent\text{-}nf}(q, \mathit{rq}), \\ \quad \quad \quad (\lambda z \in \operatorname{carr}(\operatorname{laurent\text{-}nf}(q, \mathit{rq})).\; c(\operatorname{laurent\text{-}mul}(q, \operatorname{laurent\text{-}one}(q), z)(0)))) \wedge \\ \quad \quad \exists a, b \in \operatorname{carr}(\operatorname{laurent\text{-}nf}(q, \mathit{rq})). \\ \quad \quad \quad \neg (\operatorname{laurent\text{-}pairing}(q, c)(a, b) = 1) \end{array}}}
\vnbprop{ccr-6-10}{\vnbstmt{\begin{array}{@{}l@{}} \forall q.; \forall \mathit{rq} \in \mathbb{R}, c. \\ \quad \operatorname{is\text{-}field\text{-}ring}(q) \wedge \\ \quad |\operatorname{carr}(q)| \in \mathbb{N} \wedge \\ \quad 0 < \mathit{rq} \wedge \\ \quad \mathit{rq} < 1 \wedge \\ \quad \operatorname{is\text{-}nontrivial\text{-}character}(\operatorname{triv\text{-}nf}(q), c) \Rightarrow \\ \quad \quad \operatorname{is\text{-}bilinear\text{-}bicharacter}(\operatorname{laurent\text{-}nf}(q, \mathit{rq}), \operatorname{laurent\text{-}pairing}(q, c)) \wedge \\ \quad \quad \forall x, y \in \operatorname{carr}(\operatorname{laurent\text{-}nf}(q, \mathit{rq})). \\ \quad \quad \quad \operatorname{laurent\text{-}pairing}(q, c)(x, y) = \operatorname{laurent\text{-}pairing}(q, c)(y, x) \wedge \\ \quad \quad \operatorname{is\text{-}faithful\text{-}bicharacter}(\operatorname{laurent\text{-}nf}(q, \mathit{rq}), \operatorname{laurent\text{-}pairing}(q, c)) \end{array}}}

\vnbbookfile{theorem-library/ccr1109a6-laurent}{theorem-library/ccr1109a6-laurent}
\noindent{\small ccr1109a6-laurent.scm -- wave of 2026-11-09/10 (agent CCR-1109-A6), wired 2026-11-10 from scratchpad/ccr1109a6/ccr1109a6-laurent.scm. Library: complex-analysis.}\par
\vnbprop{finite-field-nontrivial-character}{\vnbstmt{\begin{array}{@{}l@{}} \forall q. \\ \quad (\operatorname{is\text{-}field\text{-}ring}(q) \wedge (|\operatorname{carr}(q)| \in \mathbb{N})) \Rightarrow  \\ \quad \quad \exists c.\; \operatorname{is\text{-}nontrivial\text{-}character}(\operatorname{triv\text{-}nf}(q), c) \end{array}}}
\vnbprop{ccr-6-10-finite-field}{\vnbstmt{\begin{array}{@{}l@{}} \forall q.; \forall \mathit{rq} \in \mathbb{R}. \\ \quad (\operatorname{is\text{-}field\text{-}ring}(q) \wedge ((|\operatorname{carr}(q)| \in \mathbb{N}) \wedge ((0 < \mathit{rq}) \wedge (\mathit{rq} < 1)))) \Rightarrow  \\ \quad \quad \exists c. \\ \quad \quad \quad \operatorname{is\text{-}nontrivial\text{-}character}(\operatorname{triv\text{-}nf}(q), c) \wedge \\ \quad \quad \quad \operatorname{is\text{-}bilinear\text{-}bicharacter}(\operatorname{laurent\text{-}nf}(q, \mathit{rq}), \operatorname{laurent\text{-}pairing}(q, c)) \wedge \\ \quad \quad \quad \forall x, y \in \operatorname{carr}(\operatorname{laurent\text{-}nf}(q, \mathit{rq})). \\ \quad \quad \quad \quad \operatorname{laurent\text{-}pairing}(q, c)(x, y) = \operatorname{laurent\text{-}pairing}(q, c)(y, x) \wedge \\ \quad \quad \quad \operatorname{is\text{-}faithful\text{-}bicharacter}(\operatorname{laurent\text{-}nf}(q, \mathit{rq}), \operatorname{laurent\text{-}pairing}(q, c)) \end{array}}}
\vnbprop{laurent-coeff-zero-onto}{\vnbstmt{\begin{array}{@{}l@{}} \forall q.; \forall \mathit{cq} \in \operatorname{carr}(q). \\ \quad \operatorname{is\text{-}ring}(q) \Rightarrow \exists \mathit{fq} \in \operatorname{laurent\text{-}set}(q).\; (\operatorname{fq}(0) = \mathit{cq}) \end{array}}}
\vnbprop{laurent-coeff-zero-add}{\vnbstmt{\begin{array}{@{}l@{}} \forall q.; \forall \mathit{fq}, \mathit{gq} \in \operatorname{laurent\text{-}set}(q). \\ \quad \operatorname{is\text{-}ring}(q) \Rightarrow (\operatorname{laurent\text{-}add}(q, \mathit{fq}, \mathit{gq})(0) = \operatorname{add}(q)(\operatorname{fq}(0), \operatorname{gq}(0))) \end{array}}}
\vnbprop{laurent-unit-ball-integers}{\vnbstmt{\begin{array}{@{}l@{}} \forall q.; \forall \mathit{rq} \in \mathbb{R}, \mathit{fq} \in \operatorname{laurent\text{-}set}(q). \\ \quad (\operatorname{is\text{-}field\text{-}ring}(q) \wedge ((0 < \mathit{rq}) \wedge (\mathit{rq} < 1))) \Rightarrow  \\ \quad \quad (\operatorname{laurent\text{-}abs}(q, \mathit{rq}, \mathit{fq}) \leq 1) \Leftrightarrow (\mathit{fq} \in \operatorname{nv\text{-}integers}(\operatorname{laurent\text{-}nf}(q, \mathit{rq}))) \end{array}}}
\vnbprop{laurent-finite-coeff-continuous}{\vnbstmt{\begin{array}{@{}l@{}} \forall q.; \forall \mathit{rq} \in \mathbb{R}, t.; \forall f \in (\operatorname{laurent\text{-}set}(q) \to \operatorname{pts}(t)), n \in \mathbb{Z}. \\ \quad \operatorname{is\text{-}field\text{-}ring}(q) \wedge \\ \quad 0 < \mathit{rq} \wedge \\ \quad \mathit{rq} < 1 \wedge \\ \quad \operatorname{is\text{-}metric\text{-}space}(t) \wedge \\ \quad \forall p, a \in \operatorname{laurent\text{-}set}(q). \\ \quad \quad \forall j \in \mathbb{Z}.\; ((j < n) \Rightarrow (p(j) = a(j))) \Rightarrow (f(p) = f(a)) \Rightarrow \\ \quad \quad \operatorname{is\text{-}continuous}(\operatorname{nf\text{-}metric\text{-}space}(\operatorname{laurent\text{-}nf}(q, \mathit{rq})), t, f) \end{array}}}
\vnbprop{rpow-star-least-exponent}{\vnbstmt{\begin{array}{@{}l@{}} \forall q, o \in \mathbb{R}, z \in \mathbb{N}. \\ \quad ((0 < q) \wedge ((q < 1) \wedge ((0 < o) \wedge \forall t \in \mathbb{Z}.\; (({q}^{t} \leq o) \Rightarrow (-z < t))))) \Rightarrow  \\ \quad \quad \exists n \in \mathbb{Z}. \\ \quad \quad \quad ({q}^{n} \leq o) \wedge \forall a \in \mathbb{Z}.\; (({q}^{a} \leq o) \Rightarrow (n \leq a)) \end{array}}}
\vnbprop{laurent-ball-order}{\vnbstmt{\begin{array}{@{}l@{}} \forall q.; \forall \mathit{rq}, o \in \mathbb{R}. \\ \quad (\operatorname{is\text{-}field\text{-}ring}(q) \wedge ((0 < \mathit{rq}) \wedge ((\mathit{rq} < 1) \wedge (0 < o)))) \Rightarrow  \\ \quad \quad \exists n \in \mathbb{Z}. \\ \quad \quad \quad {\mathit{rq}}^{n} \leq o \wedge \\ \quad \quad \quad \forall m \in \mathbb{Z}.\; ((m < n) \Rightarrow (o < {\mathit{rq}}^{m})) \wedge \\ \quad \quad \quad \forall \mathit{fq} \in \operatorname{laurent\text{-}set}(q). \\ \quad \quad \quad \quad (\operatorname{laurent\text{-}abs}(q, \mathit{rq}, \mathit{fq}) \leq o) \Leftrightarrow \forall j \in \mathbb{Z}.\; ((j < n) \Rightarrow (\operatorname{fq}(j) = \operatorname{zero}(q))) \end{array}}}
\vnbprop{zz-restrict-type}{\vnbstmt{\begin{array}{@{}l@{}} \forall q, s.; \forall \mathit{fq} \in (\mathbb{Z} \to \operatorname{carr}(q)). \\ \quad (\operatorname{is\text{-}ring}(q) \wedge (s \subseteq \mathbb{Z})) \Rightarrow ((\lambda i \in s.\; \operatorname{fq}(i)) \in (s \to \operatorname{carr}(q))) \end{array}}}
\vnbprop{zz-restrict-injective}{\vnbstmt{\begin{array}{@{}l@{}} \forall q, s.; \forall \mathit{fq}, \mathit{gq} \in (\mathbb{Z} \to \operatorname{carr}(q)). \\ \quad \operatorname{is\text{-}ring}(q) \wedge \\ \quad s \subseteq \mathbb{Z} \wedge \\ \quad \forall j \in \mathbb{Z}.\; (\neg (j \in s) \Rightarrow (\operatorname{fq}(j) = \operatorname{zero}(q))) \wedge \\ \quad \forall j \in \mathbb{Z}.\; (\neg (j \in s) \Rightarrow (\operatorname{gq}(j) = \operatorname{zero}(q))) \wedge \\ \quad (\lambda i \in s.\; \operatorname{fq}(i)) = (\lambda i \in s.\; \operatorname{gq}(i)) \Rightarrow \\ \quad \quad \mathit{fq} = \mathit{gq} \end{array}}}
\vnbprop{zz-extend-by-zero}{\vnbstmt{\begin{array}{@{}l@{}} \forall q, s.; \forall h \in (s \to \operatorname{carr}(q)). \\ \quad (\operatorname{is\text{-}ring}(q) \wedge (s \subseteq \mathbb{Z})) \Rightarrow  \\ \quad \quad (\lambda i \in \mathbb{Z}.\; \left(\text{if }(i \in s)\text{ then }h(i)\text{ else }\operatorname{zero}(q)\right)) \in (\mathbb{Z} \to \operatorname{carr}(q)) \wedge \\ \quad \quad \forall j \in \mathbb{Z}. \\ \quad \quad \quad \neg (j \in s) \Rightarrow ((\lambda i \in \mathbb{Z}.\; \left(\text{if }(i \in s)\text{ then }h(i)\text{ else }\operatorname{zero}(q)\right))(j) = \operatorname{zero}(q)) \wedge \\ \quad \quad (\lambda i \in s.\; (\lambda i \in \mathbb{Z}.\; \left(\text{if }(i \in s)\text{ then }h(i)\text{ else }\operatorname{zero}(q)\right))(i)) = h \end{array}}}
\vnbprop{laurent-ball-product}{\vnbstmt{\begin{array}{@{}l@{}} \forall q.; \forall n \in \mathbb{Z}. \\ \quad \operatorname{is\text{-}ring}(q) \Rightarrow  \\ \quad \quad \forall \mathit{fq} \in \operatorname{laurent\text{-}set}(q). \\ \quad \quad \quad \forall j \in \mathbb{Z}.\; ((j < n) \Rightarrow (\operatorname{fq}(j) = \operatorname{zero}(q))) \Rightarrow  \\ \quad \quad \quad \quad (\lambda i \in \{k \in \mathbb{Z} : (n \leq k)\}.\; \operatorname{fq}(i)) \in (\{k \in \mathbb{Z} : (n \leq k)\} \to \operatorname{carr}(q)) \wedge \\ \quad \quad \forall h \in (\{k \in \mathbb{Z} : (n \leq k)\} \to \operatorname{carr}(q)). \exists \mathit{fq} \in \operatorname{laurent\text{-}set}(q). \\ \quad \quad \quad \forall j \in \mathbb{Z}.\; ((j < n) \Rightarrow (\operatorname{fq}(j) = \operatorname{zero}(q))) \wedge \\ \quad \quad \quad (\lambda i \in \{k \in \mathbb{Z} : (n \leq k)\}.\; \operatorname{fq}(i)) = h \wedge \\ \quad \quad \quad \forall \mathit{gq} \in \operatorname{laurent\text{-}set}(q). \\ \quad \quad \quad \quad (\forall j \in \mathbb{Z}.\; ((j < n) \Rightarrow (\operatorname{gq}(j) = \operatorname{zero}(q))) \wedge ((\lambda i \in \{k \in \mathbb{Z} : (n \leq k)\}.\; \operatorname{gq}(i)) = h)) \Rightarrow  \\ \quad \quad \quad \quad \quad \mathit{gq} = \mathit{fq} \end{array}}}
\vnbprop{laurent-ball-product-add}{\vnbstmt{\begin{array}{@{}l@{}} \forall q.; \forall n \in \mathbb{Z}, \mathit{fq}, \mathit{gq} \in \operatorname{laurent\text{-}set}(q), x \in \{k \in \mathbb{Z} : (n \leq k)\}. \\ \quad \operatorname{is\text{-}ring}(q) \Rightarrow  \\ \quad \quad (\lambda i \in \{k \in \mathbb{Z} : (n \leq k)\}.\; \operatorname{laurent\text{-}add}(q, \mathit{fq}, \mathit{gq})(i))(x) \\ \quad \quad =\; \operatorname{add}(q)((\lambda i \in \{k \in \mathbb{Z} : (n \leq k)\}.\; \operatorname{fq}(i))(x), (\lambda i \in \{k \in \mathbb{Z} : (n \leq k)\}.\; \operatorname{gq}(i))(x)) \end{array}}}
\vnbprop{laurent-polar-part-onto}{\vnbstmt{\begin{array}{@{}l@{}} \forall q. \\ \quad \operatorname{is\text{-}ring}(q) \Rightarrow  \\ \quad \quad \forall \mathit{fq} \in \operatorname{laurent\text{-}set}(q). \\ \quad \quad \quad (\lambda i \in \{k \in \mathbb{Z} : (k < 0)\}.\; \operatorname{fq}(i)) \in (\{k \in \mathbb{Z} : (k < 0)\} \to \operatorname{carr}(q)) \wedge \\ \quad \quad \quad \exists b \in \mathbb{Z}. \forall x \in \{k \in \mathbb{Z} : (k < 0)\}. \\ \quad \quad \quad \quad (x < b) \Rightarrow ((\lambda i \in \{k \in \mathbb{Z} : (k < 0)\}.\; \operatorname{fq}(i))(x) = \operatorname{zero}(q)) \wedge \\ \quad \quad \forall h \in (\{k \in \mathbb{Z} : (k < 0)\} \to \operatorname{carr}(q)). \\ \quad \quad \quad \exists b \in \mathbb{Z}.\; \forall x \in \{k \in \mathbb{Z} : (k < 0)\}.\; ((x < b) \Rightarrow (h(x) = \operatorname{zero}(q))) \Rightarrow  \\ \quad \quad \quad \quad \exists \mathit{fq} \in \operatorname{laurent\text{-}set}(q).\; ((\lambda i \in \{k \in \mathbb{Z} : (k < 0)\}.\; \operatorname{fq}(i)) = h) \end{array}}}
\vnbprop{laurent-polar-part-kernel}{\vnbstmt{\begin{array}{@{}l@{}} \forall q.; \forall \mathit{rq} \in \mathbb{R}, \mathit{fq} \in \operatorname{laurent\text{-}set}(q). \\ \quad (\operatorname{is\text{-}field\text{-}ring}(q) \wedge ((0 < \mathit{rq}) \wedge (\mathit{rq} < 1))) \Rightarrow  \\ \quad \quad \mathit{fq} \in \operatorname{nv\text{-}integers}(\operatorname{laurent\text{-}nf}(q, \mathit{rq})) \\ \quad \quad \Leftrightarrow\; \forall x \in \{k \in \mathbb{Z} : (k < 0)\}. \\ \quad \quad \quad (\lambda i \in \{k \in \mathbb{Z} : (k < 0)\}.\; \operatorname{fq}(i))(x) = \operatorname{zero}(q) \end{array}}}
\vnbprop{laurent-polar-part-add}{\vnbstmt{\begin{array}{@{}l@{}} \forall q.; \forall \mathit{fq}, \mathit{gq} \in \operatorname{laurent\text{-}set}(q), x \in \{k \in \mathbb{Z} : (k < 0)\}. \\ \quad \operatorname{is\text{-}ring}(q) \Rightarrow  \\ \quad \quad (\lambda i \in \{k \in \mathbb{Z} : (k < 0)\}.\; \operatorname{laurent\text{-}add}(q, \mathit{fq}, \mathit{gq})(i))(x) \\ \quad \quad =\; \operatorname{add}(q)((\lambda i \in \{k \in \mathbb{Z} : (k < 0)\}.\; \operatorname{fq}(i))(x), (\lambda i \in \{k \in \mathbb{Z} : (k < 0)\}.\; \operatorname{gq}(i))(x)) \end{array}}}

\vnbbookfile{theorem-library/ccr4b-base}{theorem-library/ccr4b-base}
\noindent{\small ccr4b-base.scm -- wave of 2026-11-09/10 (agent CCR-1109-C3), wired 2026-11-10 from scratchpad/ccr1109c3/ccr4b-base.scm. Library: complex-analysis.}\par
\vnbprop{ccr-remark-4-4-norm}{\vnbstmt{\begin{array}{@{}l@{}} \forall k, a.; \forall w \in a. \\ \quad (\operatorname{is\text{-}compact\text{-}nv\text{-}domain}(k, a) \wedge ((\operatorname{nv\text{-}maximal}(k, a) = \operatorname{nv\text{-}principal}(k, a, w)) \wedge \neg (w = \operatorname{zero}(k)))) \Rightarrow  \\ \quad \quad (w \in \operatorname{nv\text{-}maximal}(k, a)) \wedge (\operatorname{fnrm}(k)(w) < 1) \end{array}}}
\vnbprop{ccr-remark-4-4-unit}{\vnbstmt{\begin{array}{@{}l@{}} \forall k, a.; \forall w \in a, u \in \operatorname{nv\text{-}units}(k, a). \\ \quad (\operatorname{is\text{-}compact\text{-}nv\text{-}domain}(k, a) \wedge ((\operatorname{nv\text{-}maximal}(k, a) = \operatorname{nv\text{-}principal}(k, a, w)) \wedge \neg (w = \operatorname{zero}(k)))) \Rightarrow  \\ \quad \quad \operatorname{mul}(k)(u, w) \in a \wedge \\ \quad \quad \operatorname{nv\text{-}maximal}(k, a) = \operatorname{nv\text{-}principal}(k, a, \operatorname{mul}(k)(u, w)) \wedge \\ \quad \quad \neg (\operatorname{mul}(k)(u, w) = \operatorname{zero}(k)) \end{array}}}
\vnbprop{ccr-remark-4-4}{\vnbstmt{\begin{array}{@{}l@{}} \forall k, a.; \forall w, b \in a. \\ \quad \operatorname{is\text{-}compact\text{-}nv\text{-}domain}(k, a) \wedge \\ \quad \operatorname{nv\text{-}maximal}(k, a) = \operatorname{nv\text{-}principal}(k, a, w) \wedge \\ \quad \neg (w = \operatorname{zero}(k)) \wedge \\ \quad \operatorname{nv\text{-}maximal}(k, a) = \operatorname{nv\text{-}principal}(k, a, b) \Rightarrow \\ \quad \quad \exists u \in \operatorname{nv\text{-}units}(k, a).\; (b = \operatorname{mul}(k)(u, w)) \end{array}}}
\vnbprop{ccr-remark-4-5}{\vnbstmt{\begin{array}{@{}l@{}} \forall k, a.; \forall w, b \in a. \\ \quad (\operatorname{is\text{-}compact\text{-}nv\text{-}domain}(k, a) \wedge ((\operatorname{nv\text{-}maximal}(k, a) = \operatorname{nv\text{-}principal}(k, a, w)) \wedge \neg (w = \operatorname{zero}(k)))) \Rightarrow  \\ \quad \quad \operatorname{nv\text{-}maximal}(k, a) = \operatorname{nv\text{-}principal}(k, a, b) \\ \quad \quad \Leftrightarrow\; ((b \in \operatorname{nv\text{-}maximal}(k, a)) \wedge \forall x \in \operatorname{nv\text{-}maximal}(k, a).\; (\operatorname{fnrm}(k)(x) \leq \operatorname{fnrm}(k)(b))) \end{array}}}
\vnbprop{is-fractions-of-mono}{\vnbstmt{\begin{array}{@{}l@{}} \forall k, a, b. \\ \quad ((a \subseteq b) \wedge \operatorname{is\text{-}fractions\text{-}of}(k, a)) \Rightarrow \operatorname{is\text{-}fractions\text{-}of}(k, b) \end{array}}}
\vnbprop{nv-residue-unfold}{\vnbstmt{\forall k, a.\; (\operatorname{nv\text{-}residue}(k, a) \simeq \operatorname{ringoid\text{-}quotient\text{-}ring}(\operatorname{nv\text{-}res\text{-}ringoid}(k, a)))}}
\vnbprop{nv-residue-zero-class}{\vnbstmt{\begin{array}{@{}l@{}} \forall k. \\ \quad \operatorname{is\text{-}nv\text{-}field}(k) \Rightarrow  \\ \quad \quad (\operatorname{zero}(\operatorname{nv\text{-}residue}(k, \operatorname{nv\text{-}integers}(k))) \simeq \operatorname{class}(\operatorname{ringoid\text{-}setoid}(\operatorname{nv\text{-}res\text{-}ringoid}(k, \operatorname{nv\text{-}integers}(k))), \operatorname{zero}(k))) \end{array}}}
\vnbprop{nv-maximal-in-zf}{\vnbstmt{\begin{array}{@{}l@{}} \forall k, f. \\ \quad (\operatorname{is\text{-}nv\text{-}field}(k) \wedge (\operatorname{zero}(\operatorname{nv\text{-}residue}(k, \operatorname{nv\text{-}integers}(k))) \in f)) \Rightarrow  \\ \quad \quad \operatorname{nv\text{-}maximal}(k, \operatorname{nv\text{-}integers}(k)) \subseteq \operatorname{nv\text{-}zf}(k, f) \end{array}}}
\vnbprop{ccr-4-20-collapse}{\vnbstmt{\begin{array}{@{}l@{}} \forall k.; \forall w \in \operatorname{carr}(k), f. \\ \quad \operatorname{is\text{-}nv\text{-}field}(k) \wedge \\ \quad \neg (w = \operatorname{zero}(k)) \wedge \\ \quad \operatorname{fnrm}(k)(w) < 1 \wedge \\ \quad \operatorname{zero}(\operatorname{nv\text{-}residue}(k, \operatorname{nv\text{-}integers}(k))) \in f \Rightarrow \\ \quad \quad \operatorname{is\text{-}fractions\text{-}of}(k, \operatorname{nv\text{-}zf}(k, f)) \end{array}}}
\vnbprop{nf-closure-in-carr}{\vnbstmt{\begin{array}{@{}l@{}} \forall k, a. \\ \quad \operatorname{is\text{-}normed\text{-}field}(k) \Rightarrow (\operatorname{closure}(\operatorname{nf\text{-}metric\text{-}space}(k), a) \subseteq \operatorname{carr}(k)) \end{array}}}
\vnbprop{nf-closure-near}{\vnbstmt{\begin{array}{@{}l@{}} \forall k, a.; \forall x \in \operatorname{closure}(\operatorname{nf\text{-}metric\text{-}space}(k), a), s. \\ \quad (\operatorname{is\text{-}normed\text{-}field}(k) \wedge ((a \subseteq \operatorname{carr}(k)) \wedge \operatorname{pos\text{-}rr}(s))) \Rightarrow  \\ \quad \quad \exists b \in a.\; (\operatorname{fnrm}(k)(\operatorname{add}(k)(x, \operatorname{neg}(k)(b))) < s) \end{array}}}
\vnbprop{nf-closure-approx}{\vnbstmt{\begin{array}{@{}l@{}} \forall k, a.; \forall x \in \operatorname{carr}(k). \\ \quad \operatorname{is\text{-}normed\text{-}field}(k) \wedge \\ \quad a \subseteq \operatorname{carr}(k) \wedge \\ \quad \forall s. \\ \quad \quad \operatorname{pos\text{-}rr}(s) \Rightarrow \exists b \in a.\; (\operatorname{fnrm}(k)(\operatorname{add}(k)(x, \operatorname{neg}(k)(b))) < s) \Rightarrow \\ \quad \quad x \in \operatorname{closure}(\operatorname{nf\text{-}metric\text{-}space}(k), a) \end{array}}}
\vnbprop{nf-closure-near-scaled}{\vnbstmt{\begin{array}{@{}l@{}} \forall k, a.; \forall x \in \operatorname{closure}(\operatorname{nf\text{-}metric\text{-}space}(k), a), c \in \operatorname{carr}(k), b, s. \\ \quad (\operatorname{is\text{-}normed\text{-}field}(k) \wedge ((a \subseteq \operatorname{carr}(k)) \wedge (\operatorname{pos\text{-}rr}(b) \wedge \operatorname{pos\text{-}rr}(s)))) \Rightarrow  \\ \quad \quad \exists d \in a. \\ \quad \quad \quad ((\operatorname{fnrm}(k)(\operatorname{add}(k)(x, \operatorname{neg}(k)(d))) < b) \wedge ((\operatorname{fnrm}(k)(c) \cdot \operatorname{fnrm}(k)(\operatorname{add}(k)(x, \operatorname{neg}(k)(d)))) < s)) \end{array}}}
\vnbprop{rr-recip-unique}{\vnbstmt{\forall p, u, v \in \mathbb{R}.\; ((((p \cdot u) = 1) \wedge ((p \cdot v) = 1)) \Rightarrow (u = v))}}
\vnbprop{nf-closure-neg}{\vnbstmt{\begin{array}{@{}l@{}} \forall k, s.; \forall x \in \operatorname{closure}(\operatorname{nf\text{-}metric\text{-}space}(k), s). \\ \quad (\operatorname{is\text{-}normed\text{-}field}(k) \wedge \operatorname{is\text{-}nf\text{-}subfield}(k, s)) \Rightarrow  \\ \quad \quad \operatorname{neg}(k)(x) \in \operatorname{closure}(\operatorname{nf\text{-}metric\text{-}space}(k), s) \end{array}}}
\vnbprop{nf-closure-add}{\vnbstmt{\begin{array}{@{}l@{}} \forall k, s.; \forall x, y \in \operatorname{closure}(\operatorname{nf\text{-}metric\text{-}space}(k), s). \\ \quad (\operatorname{is\text{-}normed\text{-}field}(k) \wedge \operatorname{is\text{-}nf\text{-}subfield}(k, s)) \Rightarrow  \\ \quad \quad \operatorname{add}(k)(x, y) \in \operatorname{closure}(\operatorname{nf\text{-}metric\text{-}space}(k), s) \end{array}}}
\vnbprop{nf-closure-smul}{\vnbstmt{\begin{array}{@{}l@{}} \forall k, s.; \forall c \in s, y \in \operatorname{closure}(\operatorname{nf\text{-}metric\text{-}space}(k), s). \\ \quad (\operatorname{is\text{-}normed\text{-}field}(k) \wedge \operatorname{is\text{-}nf\text{-}subfield}(k, s)) \Rightarrow  \\ \quad \quad \operatorname{mul}(k)(c, y) \in \operatorname{closure}(\operatorname{nf\text{-}metric\text{-}space}(k), s) \end{array}}}
\vnbprop{nf-closure-mul}{\vnbstmt{\begin{array}{@{}l@{}} \forall k, s.; \forall x, y \in \operatorname{closure}(\operatorname{nf\text{-}metric\text{-}space}(k), s). \\ \quad (\operatorname{is\text{-}normed\text{-}field}(k) \wedge \operatorname{is\text{-}nf\text{-}subfield}(k, s)) \Rightarrow  \\ \quad \quad \operatorname{mul}(k)(x, y) \in \operatorname{closure}(\operatorname{nf\text{-}metric\text{-}space}(k), s) \end{array}}}
\vnbprop{nf-closure-inv}{\vnbstmt{\begin{array}{@{}l@{}} \forall k, s.; \forall x \in \operatorname{closure}(\operatorname{nf\text{-}metric\text{-}space}(k), s). \\ \quad (\operatorname{is\text{-}nv\text{-}field}(k) \wedge (\operatorname{is\text{-}nf\text{-}subfield}(k, s) \wedge \neg (x = \operatorname{zero}(k)))) \Rightarrow  \\ \quad \quad \exists y \in \operatorname{closure}(\operatorname{nf\text{-}metric\text{-}space}(k), s).\; (\operatorname{mul}(k)(x, y) = \operatorname{one}(k)) \end{array}}}
\vnbprop{nf-closure-subfield}{\vnbstmt{\begin{array}{@{}l@{}} \forall k, s. \\ \quad (\operatorname{is\text{-}nv\text{-}field}(k) \wedge \operatorname{is\text{-}nf\text{-}subfield}(k, s)) \Rightarrow  \\ \quad \quad \operatorname{is\text{-}nf\text{-}subfield}(k, \operatorname{closure}(\operatorname{nf\text{-}metric\text{-}space}(k), s)) \end{array}}}
\vnbprop{nf-subfield-gen-unfold}{\vnbstmt{\begin{array}{@{}l@{}} \forall k, x. \\ \quad (\operatorname{nf\text{-}subfield\text{-}gen}(k, x) \simeq \{z \in \operatorname{carr}(k) : \forall s \in \operatorname{power}(\operatorname{carr}(k)).\; ((\operatorname{is\text{-}nf\text{-}subfield}(k, s) \wedge (x \subseteq s)) \Rightarrow (z \in s))\}) \end{array}}}
\vnbprop{nf-subfield-gen-sub}{\vnbstmt{\forall k, x.\; (\operatorname{is\text{-}normed\text{-}field}(k) \Rightarrow (\operatorname{nf\text{-}subfield\text{-}gen}(k, x) \subseteq \operatorname{carr}(k)))}}
\vnbprop{nf-subfield-gen-contains}{\vnbstmt{\begin{array}{@{}l@{}} \forall k, x. \\ \quad (\operatorname{is\text{-}normed\text{-}field}(k) \wedge (x \subseteq \operatorname{carr}(k))) \Rightarrow (x \subseteq \operatorname{nf\text{-}subfield\text{-}gen}(k, x)) \end{array}}}
\vnbprop{nf-subfield-gen-least}{\vnbstmt{\begin{array}{@{}l@{}} \forall k, x.; \forall s \in \operatorname{power}(\operatorname{carr}(k)). \\ \quad (\operatorname{is\text{-}nf\text{-}subfield}(k, s) \wedge (x \subseteq s)) \Rightarrow (\operatorname{nf\text{-}subfield\text{-}gen}(k, x) \subseteq s) \end{array}}}
\vnbprop{nf-subfield-gen-is-subfield}{\vnbstmt{\begin{array}{@{}l@{}} \forall k, x. \\ \quad \operatorname{is\text{-}normed\text{-}field}(k) \Rightarrow \operatorname{is\text{-}nf\text{-}subfield}(k, \operatorname{nf\text{-}subfield\text{-}gen}(k, x)) \end{array}}}
\vnbprop{nv-base-unfold}{\vnbstmt{\begin{array}{@{}l@{}} \forall k, w. \\ \quad (\operatorname{nv\text{-}base}(k, w) \simeq \operatorname{closure}(\operatorname{nf\text{-}metric\text{-}space}(k), \operatorname{nf\text{-}subfield\text{-}gen}(k, \operatorname{nv\text{-}base\text{-}gen}(k, w)))) \end{array}}}
\vnbprop{nv-base-gen-unfold}{\vnbstmt{\begin{array}{@{}l@{}} \forall k, w. \\ \quad (\operatorname{nv\text{-}base\text{-}gen}(k, w) \simeq \{z \in \operatorname{carr}(k) : ((z = w) \wedge \exists p.\; \operatorname{has\text{-}characteristic}(\operatorname{normed\text{-}field\text{-}as\text{-}commutative\text{-}ring}(k), p))\}) \end{array}}}
\vnbprop{nv-base-gen-sub}{\vnbstmt{\forall k, w.\; (\operatorname{is\text{-}normed\text{-}field}(k) \Rightarrow (\operatorname{nv\text{-}base\text{-}gen}(k, w) \subseteq \operatorname{carr}(k)))}}
\vnbprop{nv-base-subfield}{\vnbstmt{\forall k, w.\; (\operatorname{is\text{-}nv\text{-}field}(k) \Rightarrow \operatorname{is\text{-}nf\text{-}subfield}(k, \operatorname{nv\text{-}base}(k, w)))}}
\vnbprop{nv-base-closed}{\vnbstmt{\begin{array}{@{}l@{}} \forall k, w. \\ \quad \operatorname{is\text{-}normed\text{-}field}(k) \Rightarrow \operatorname{is\text{-}closed}(\operatorname{nf\text{-}metric\text{-}space}(k), \operatorname{nv\text{-}base}(k, w)) \end{array}}}
\vnbprop{nv-base-contains-gen}{\vnbstmt{\begin{array}{@{}l@{}} \forall k, w. \\ \quad \operatorname{is\text{-}normed\text{-}field}(k) \Rightarrow  \\ \quad \quad \operatorname{nf\text{-}subfield\text{-}gen}(k, \operatorname{nv\text{-}base\text{-}gen}(k, w)) \subseteq \operatorname{nv\text{-}base}(k, w) \end{array}}}
\vnbprop{nv-base-char-p}{\vnbstmt{\begin{array}{@{}l@{}} \forall k.; \forall w \in \operatorname{carr}(k), p. \\ \quad (\operatorname{is\text{-}normed\text{-}field}(k) \wedge \operatorname{has\text{-}characteristic}(\operatorname{normed\text{-}field\text{-}as\text{-}commutative\text{-}ring}(k), p)) \Rightarrow  \\ \quad \quad w \in \operatorname{nv\text{-}base}(k, w) \end{array}}}
\vnbprop{nv-base-char-0}{\vnbstmt{\begin{array}{@{}l@{}} \forall k, w, z. \\ \quad \forall p.\; \neg \operatorname{has\text{-}characteristic}(\operatorname{normed\text{-}field\text{-}as\text{-}commutative\text{-}ring}(k), p) \Rightarrow  \\ \quad \quad \neg (z \in \operatorname{nv\text{-}base\text{-}gen}(k, w)) \end{array}}}
\vnbprop{ccr-4-21-base}{\vnbstmt{\begin{array}{@{}l@{}} \forall k, w. \\ \quad \operatorname{is\text{-}compact\text{-}nv\text{-}domain}(k, \operatorname{nv\text{-}integers}(k)) \Rightarrow  \\ \quad \quad \operatorname{is\text{-}nf\text{-}subfield}(k, \operatorname{nv\text{-}base}(k, w)) \wedge \\ \quad \quad \operatorname{is\text{-}closed}(\operatorname{nf\text{-}metric\text{-}space}(k), \operatorname{nv\text{-}base}(k, w)) \wedge \\ \quad \quad \operatorname{is\text{-}compact}(\operatorname{subspace\text{-}ms}(\operatorname{nf\text{-}metric\text{-}space}(k), (\operatorname{nv\text{-}integers}(k) \cap \operatorname{nv\text{-}base}(k, w)))) \end{array}}}
\vnbprop{ccr-4-21-base-char-p}{\vnbstmt{\begin{array}{@{}l@{}} \forall k.; \forall w \in \operatorname{nv\text{-}integers}(k), p. \\ \quad \operatorname{is\text{-}compact\text{-}nv\text{-}domain}(k, \operatorname{nv\text{-}integers}(k)) \wedge \\ \quad \operatorname{nv\text{-}maximal}(k, \operatorname{nv\text{-}integers}(k)) = \operatorname{nv\text{-}principal}(k, \operatorname{nv\text{-}integers}(k), w) \wedge \\ \quad \neg (w = \operatorname{zero}(k)) \wedge \\ \quad \operatorname{has\text{-}characteristic}(\operatorname{normed\text{-}field\text{-}as\text{-}commutative\text{-}ring}(k), p) \Rightarrow \\ \quad \quad (w \in \operatorname{nv\text{-}base}(k, w)) \wedge ((0 < \operatorname{fnrm}(k)(w)) \wedge (\operatorname{fnrm}(k)(w) < 1)) \end{array}}}
\vnbprop{ccr-4-23-base-char-p}{\vnbstmt{\begin{array}{@{}l@{}} \forall k.; \forall w \in \operatorname{nv\text{-}integers}(k), p. \\ \quad \operatorname{is\text{-}compact\text{-}nv\text{-}domain}(k, \operatorname{nv\text{-}integers}(k)) \wedge \\ \quad \operatorname{nv\text{-}maximal}(k, \operatorname{nv\text{-}integers}(k)) = \operatorname{nv\text{-}principal}(k, \operatorname{nv\text{-}integers}(k), w) \wedge \\ \quad \neg (w = \operatorname{zero}(k)) \wedge \\ \quad \operatorname{has\text{-}characteristic}(\operatorname{normed\text{-}field\text{-}as\text{-}commutative\text{-}ring}(k), p) \Rightarrow \\ \quad \quad w \in \operatorname{nv\text{-}base}(k, w) \wedge \\ \quad \quad \forall x \in \operatorname{nv\text{-}base}(k, w). \\ \quad \quad \quad (\operatorname{fnrm}(k)(x) < 1) \Rightarrow (\operatorname{fnrm}(k)(x) \leq \operatorname{fnrm}(k)(w)) \end{array}}}

\vnbbookfile{theorem-library/ccr4b-remark418}{theorem-library/ccr4b-remark418}
\noindent{\small ccr4b-remark418.scm -- wave of 2026-11-09/10 (agent CCR-1109-C3), wired 2026-11-10 from scratchpad/ccr1109c3/ccr4b-remark418.scm. Library: complex-analysis.}\par
\vnbprop{ccr-remark-4-18}{\vnbstmt{\begin{array}{@{}l@{}} \forall k, a, w, p.; \forall n \in \mathbb{Z}, b, x. \\ \quad \operatorname{is\text{-}compact\text{-}nv\text{-}domain}(k, a) \wedge \\ \quad w \in a \wedge \\ \quad \operatorname{nv\text{-}maximal}(k, a) = \operatorname{nv\text{-}principal}(k, a, w) \wedge \\ \quad \neg (w = \operatorname{zero}(k)) \wedge \\ \quad a = \operatorname{nv\text{-}integers}(k) \wedge \\ \quad \operatorname{is\text{-}nv\text{-}rep\text{-}set}(k, a, p) \wedge \\ \quad \operatorname{is\text{-}nv\text{-}expansion}(k, w, p, b, n, x) \wedge \\ \quad \neg (b(n) = \operatorname{zero}(k)) \Rightarrow \\ \quad \quad \operatorname{fnrm}(k)(x) = \operatorname{fnrm}(k)(\operatorname{nv\text{-}zpow}(k, w, n)) \end{array}}}

\vnbbookfile{theorem-library/ccr4c-local}{theorem-library/ccr4c-local}
\noindent{\small ccr4c-local.scm -- wave of 2026-11-09/10 (agent CCR-1109-C4), wired 2026-11-10 from scratchpad/ccr1109c4/ccr4c-local.scm. Library: complex-analysis.}\par
\vnbprop{sub-nf-local}{\vnbstmt{\begin{array}{@{}l@{}} \forall k, s.; \forall w \in s. \\ \quad \operatorname{is\text{-}nv\text{-}field}(k) \wedge \\ \quad \operatorname{is\text{-}nf\text{-}subfield}(k, s) \wedge \\ \quad \operatorname{is\text{-}compact}(\operatorname{subspace\text{-}ms}(\operatorname{nf\text{-}metric\text{-}space}(k), (\operatorname{nv\text{-}integers}(k) \cap s))) \wedge \\ \quad 0 < \operatorname{fnrm}(k)(w) \wedge \\ \quad \operatorname{fnrm}(k)(w) < 1 \Rightarrow \\ \quad \quad \operatorname{is\text{-}local\text{-}nv\text{-}field}(\operatorname{sub\text{-}nf}(k, s)) \end{array}}}
\vnbprop{ccr-4-21}{\vnbstmt{\begin{array}{@{}l@{}} \forall k.; \forall w \in \operatorname{nv\text{-}integers}(k), p. \\ \quad \operatorname{is\text{-}compact\text{-}nv\text{-}domain}(k, \operatorname{nv\text{-}integers}(k)) \wedge \\ \quad \operatorname{nv\text{-}maximal}(k, \operatorname{nv\text{-}integers}(k)) = \operatorname{nv\text{-}principal}(k, \operatorname{nv\text{-}integers}(k), w) \wedge \\ \quad \neg (w = \operatorname{zero}(k)) \wedge \\ \quad \operatorname{has\text{-}characteristic}(\operatorname{normed\text{-}field\text{-}as\text{-}commutative\text{-}ring}(k), p) \Rightarrow \\ \quad \quad \operatorname{is\text{-}local\text{-}nv\text{-}field}(\operatorname{sub\text{-}nf}(k, \operatorname{nv\text{-}base}(k, w))) \end{array}}}
\vnbprop{nv-residue-zz-act-one}{\vnbstmt{\begin{array}{@{}l@{}} \forall n \in \mathbb{N}, k. \\ \quad \operatorname{is\text{-}nv\text{-}field}(k) \Rightarrow  \\ \quad \quad \operatorname{zz\text{-}act}(\operatorname{ring\text{-}additive\text{-}ag}(\operatorname{normed\text{-}field\text{-}as\text{-}commutative\text{-}ring}(k)), \\ \quad \quad \quad n, \operatorname{one}(\operatorname{normed\text{-}field\text{-}as\text{-}commutative\text{-}ring}(k))) \\ \quad \quad \in\; \operatorname{nv\text{-}integers}(k) \wedge \\ \quad \quad \operatorname{zz\text{-}act}(\operatorname{ring\text{-}additive\text{-}ag}(\operatorname{ringoid\text{-}quotient\text{-}ring}(\operatorname{nv\text{-}res\text{-}ringoid}(k, \operatorname{nv\text{-}integers}(k)))), \\ \quad \quad \quad n, \\ \quad \quad \quad \operatorname{one}(\operatorname{ringoid\text{-}quotient\text{-}ring}(\operatorname{nv\text{-}res\text{-}ringoid}(k, \operatorname{nv\text{-}integers}(k))))) \\ \quad \quad =\; \operatorname{class}(\operatorname{ringoid\text{-}setoid}(\operatorname{nv\text{-}res\text{-}ringoid}(k, \operatorname{nv\text{-}integers}(k))), \\ \quad \quad \quad \operatorname{zz\text{-}act}(\operatorname{ring\text{-}additive\text{-}ag}(\operatorname{normed\text{-}field\text{-}as\text{-}commutative\text{-}ring}(k)), \\ \quad \quad \quad \quad n, \operatorname{one}(\operatorname{normed\text{-}field\text{-}as\text{-}commutative\text{-}ring}(k)))) \end{array}}}
\vnbprop{ccr-4-21-char-0}{\vnbstmt{\begin{array}{@{}l@{}} \forall k.; \forall w \in \operatorname{nv\text{-}integers}(k). \\ \quad \operatorname{is\text{-}compact\text{-}nv\text{-}domain}(k, \operatorname{nv\text{-}integers}(k)) \wedge \\ \quad \operatorname{nv\text{-}maximal}(k, \operatorname{nv\text{-}integers}(k)) = \operatorname{nv\text{-}principal}(k, \operatorname{nv\text{-}integers}(k), w) \wedge \\ \quad \neg (w = \operatorname{zero}(k)) \wedge \\ \quad \forall p.\; \neg \operatorname{has\text{-}characteristic}(\operatorname{normed\text{-}field\text{-}as\text{-}commutative\text{-}ring}(k), p) \Rightarrow \\ \quad \quad \operatorname{is\text{-}local\text{-}nv\text{-}field}(\operatorname{sub\text{-}nf}(k, \operatorname{nv\text{-}base}(k, w))) \end{array}}}

\vnbbookfile{theorem-library/ccr6-e}{theorem-library/ccr6-e}
\noindent{\small ccr6-e.scm -- wave of 2026-11-09/10 (agent CCR-1109-D3), wired 2026-11-10 from scratchpad/ccr1109d3/ccr6-e.scm. Library: complex-analysis.}\par
\vnbprop{c6-qp-uniformizer}{\vnbstmt{\begin{array}{@{}l@{}} \forall p. \\ \quad \operatorname{zz\text{-}prime}(p) \Rightarrow  \\ \quad \quad \operatorname{is\text{-}local\text{-}uniformizer}(\operatorname{qp\text{-}field}(p), \operatorname{embed}(\operatorname{nf\text{-}metric\text{-}space}(\operatorname{qq\text{-}padic}(p)))(p)) \end{array}}}
\vnbprop{c6-integers-char-nontrivial}{\vnbstmt{\begin{array}{@{}l@{}} \forall f, w, i. \\ \quad \operatorname{is\text{-}local\text{-}uniformizer}(f, w) \wedge \\ \quad \operatorname{is\text{-}add\text{-}character}(f, i) \wedge \\ \quad \forall x \in \operatorname{carr}(f).\; ((i(x) = 1) \Leftrightarrow (x \in \operatorname{nv\text{-}integers}(f))) \Rightarrow \\ \quad \quad \operatorname{is\text{-}nontrivial\text{-}character}(f, i) \end{array}}}
\vnbprop{ccr-5-20-nontrivial}{\vnbstmt{\begin{array}{@{}l@{}} \forall p. \\ \quad \operatorname{zz\text{-}prime}(p) \Rightarrow \operatorname{is\text{-}nontrivial\text{-}character}(\operatorname{qp\text{-}field}(p), \operatorname{qp\text{-}lambda}(p)) \end{array}}}
\vnbprop{ccr-6-18-qp}{\vnbstmt{\begin{array}{@{}l@{}} \forall p. \\ \quad \operatorname{zz\text{-}prime}(p) \Rightarrow  \\ \quad \quad \exists w. \\ \quad \quad \quad \operatorname{is\text{-}unipolar}(\operatorname{qp\text{-}field}(p), w) \wedge \operatorname{is\text{-}faithful\text{-}bicharacter}(\operatorname{qp\text{-}field}(p), w) \end{array}}}
\vnbprop{c6-vs-module}{\vnbstmt{\begin{array}{@{}l@{}} \forall f, d. \\ \quad (\operatorname{is\text{-}normed\text{-}field}(f) \wedge \operatorname{is\text{-}k\text{-}vector\text{-}space}(f, d)) \Rightarrow \operatorname{is\text{-}module}(d) \end{array}}}
\vnbprop{c6-vs-scal}{\vnbstmt{\begin{array}{@{}l@{}} \forall f, d. \\ \quad (\operatorname{is\text{-}normed\text{-}field}(f) \wedge \operatorname{is\text{-}k\text{-}vector\text{-}space}(f, d)) \Rightarrow  \\ \quad \quad \operatorname{scal}(d) = \operatorname{normed\text{-}field\text{-}as\text{-}commutative\text{-}ring}(f) \end{array}}}
\vnbprop{c6-scal-carr}{\vnbstmt{\begin{array}{@{}l@{}} \forall f, d. \\ \quad (\operatorname{is\text{-}normed\text{-}field}(f) \wedge \operatorname{is\text{-}k\text{-}vector\text{-}space}(f, d)) \Rightarrow (\operatorname{carr}(\operatorname{scal}(d)) \simeq \operatorname{carr}(f)) \end{array}}}
\vnbprop{c6-scal-add}{\vnbstmt{\begin{array}{@{}l@{}} \forall f, d. \\ \quad (\operatorname{is\text{-}normed\text{-}field}(f) \wedge \operatorname{is\text{-}k\text{-}vector\text{-}space}(f, d)) \Rightarrow (\operatorname{add}(\operatorname{scal}(d)) \simeq \operatorname{add}(f)) \end{array}}}
\vnbprop{c6-scal-mul}{\vnbstmt{\begin{array}{@{}l@{}} \forall f, d. \\ \quad (\operatorname{is\text{-}normed\text{-}field}(f) \wedge \operatorname{is\text{-}k\text{-}vector\text{-}space}(f, d)) \Rightarrow (\operatorname{mul}(\operatorname{scal}(d)) \simeq \operatorname{mul}(f)) \end{array}}}
\vnbprop{c6-scal-one}{\vnbstmt{\begin{array}{@{}l@{}} \forall f, d. \\ \quad (\operatorname{is\text{-}normed\text{-}field}(f) \wedge \operatorname{is\text{-}k\text{-}vector\text{-}space}(f, d)) \Rightarrow (\operatorname{one}(\operatorname{scal}(d)) \simeq \operatorname{one}(f)) \end{array}}}
\vnbprop{c6-act-in}{\vnbstmt{\begin{array}{@{}l@{}} \forall f, d.; \forall r \in \operatorname{carr}(f), x \in \operatorname{vec}(d). \\ \quad (\operatorname{is\text{-}normed\text{-}field}(f) \wedge \operatorname{is\text{-}k\text{-}vector\text{-}space}(f, d)) \Rightarrow (\operatorname{act}(d)(r, x) \in \operatorname{vec}(d)) \end{array}}}
\vnbprop{c6-act-mul}{\vnbstmt{\begin{array}{@{}l@{}} \forall f, d.; \forall r, s \in \operatorname{carr}(f), x \in \operatorname{vec}(d). \\ \quad (\operatorname{is\text{-}normed\text{-}field}(f) \wedge \operatorname{is\text{-}k\text{-}vector\text{-}space}(f, d)) \Rightarrow  \\ \quad \quad \operatorname{act}(d)(\operatorname{mul}(f)(r, s), x) = \operatorname{act}(d)(r, \operatorname{act}(d)(s, x)) \end{array}}}
\vnbprop{c6-act-add}{\vnbstmt{\begin{array}{@{}l@{}} \forall f, d.; \forall r, s \in \operatorname{carr}(f), x \in \operatorname{vec}(d). \\ \quad (\operatorname{is\text{-}normed\text{-}field}(f) \wedge \operatorname{is\text{-}k\text{-}vector\text{-}space}(f, d)) \Rightarrow  \\ \quad \quad \operatorname{act}(d)(\operatorname{add}(f)(r, s), x) = \operatorname{vadd}(d)(\operatorname{act}(d)(r, x), \operatorname{act}(d)(s, x)) \end{array}}}
\vnbprop{c6-act-one}{\vnbstmt{\begin{array}{@{}l@{}} \forall f, d.; \forall x \in \operatorname{vec}(d). \\ \quad (\operatorname{is\text{-}normed\text{-}field}(f) \wedge \operatorname{is\text{-}k\text{-}vector\text{-}space}(f, d)) \Rightarrow (\operatorname{act}(d)(\operatorname{one}(f), x) = x) \end{array}}}
\vnbprop{c6-act-vadd}{\vnbstmt{\begin{array}{@{}l@{}} \forall f, d.; \forall r \in \operatorname{carr}(f), x, y \in \operatorname{vec}(d). \\ \quad (\operatorname{is\text{-}normed\text{-}field}(f) \wedge \operatorname{is\text{-}k\text{-}vector\text{-}space}(f, d)) \Rightarrow  \\ \quad \quad \operatorname{act}(d)(r, \operatorname{vadd}(d)(x, y)) = \operatorname{vadd}(d)(\operatorname{act}(d)(r, x), \operatorname{act}(d)(r, y)) \end{array}}}
\vnbprop{c6-epolar-mem}{\vnbstmt{\begin{array}{@{}l@{}} \forall d, w, v, h. \\ \quad (h \in \operatorname{e\text{-}polar}(d, w, v)) \Leftrightarrow ((h \in \operatorname{vec}(d)) \wedge \forall x \in v.\; (w(h, x) = 1)) \end{array}}}
\vnbprop{c6-ball-mem}{\vnbstmt{\forall d, u, r, x.\; ((x \in \operatorname{norm\text{-}ball}(d, u, r)) \Leftrightarrow ((x \in \operatorname{vec}(d)) \wedge (u(x) \leq r)))}}
\vnbprop{ccr-6-30}{\vnbstmt{\begin{array}{@{}l@{}} \forall f, d, u, w.; \forall c \in \operatorname{carr}(f), a \in \operatorname{vec}(d). \\ \quad (\operatorname{is\text{-}normed\text{-}field}(f) \wedge (\operatorname{is\text{-}e\text{-}unipolar}(f, d, u, w) \wedge \neg (c = \operatorname{zero}(f)))) \Rightarrow  \\ \quad \quad (a \in \operatorname{e\text{-}polar}(d, w, \operatorname{norm\text{-}ball}(d, u, \operatorname{fnrm}(f)(c)))) \Leftrightarrow (u(\operatorname{act}(d)(c, a)) \leq 1) \end{array}}}
\vnbprop{c6-wk-swap}{\vnbstmt{\begin{array}{@{}l@{}} \forall f, k.; \forall a, l \in \operatorname{carr}(f). \\ \quad (\operatorname{is\text{-}normed\text{-}field}(f) \wedge \operatorname{is\text{-}bilinear\text{-}bicharacter}(f, k)) \Rightarrow  \\ \quad \quad k(\operatorname{mul}(f)(a, l), \operatorname{one}(f)) = k(l, a) \end{array}}}
\vnbprop{ccr-6-32-unique}{\vnbstmt{\begin{array}{@{}l@{}} \forall f, k, d, b, c.; \forall u, v \in \operatorname{vec}(d). \\ \quad \operatorname{is\text{-}normed\text{-}field}(f) \wedge \\ \quad \operatorname{is\text{-}bilinear\text{-}bicharacter}(f, k) \wedge \\ \quad \operatorname{is\text{-}faithful\text{-}bicharacter}(f, k) \wedge \\ \quad \operatorname{is\text{-}k\text{-}vector\text{-}space}(f, d) \wedge \\ \quad \operatorname{is\text{-}k\text{-}bilinear\text{-}form}(f, d, b) \wedge \\ \quad \operatorname{is\text{-}k\text{-}bilinear\text{-}form}(f, d, c) \wedge \\ \quad \forall u, v \in \operatorname{vec}(d).\; (k(b(u, v), \operatorname{one}(f)) = k(c(u, v), \operatorname{one}(f))) \Rightarrow \\ \quad \quad b(u, v) = c(u, v) \end{array}}}
\vnbprop{c6-wk-add-l}{\vnbstmt{\begin{array}{@{}l@{}} \forall f, k.; \forall a, b, y \in \operatorname{carr}(f). \\ \quad (\operatorname{is\text{-}normed\text{-}field}(f) \wedge \operatorname{is\text{-}bilinear\text{-}bicharacter}(f, k)) \Rightarrow  \\ \quad \quad k(\operatorname{add}(f)(a, b), y) = (k(a, y) \cdot k(b, y)) \end{array}}}
\vnbprop{c6-wk-add-r}{\vnbstmt{\begin{array}{@{}l@{}} \forall f, k.; \forall a, b, y \in \operatorname{carr}(f). \\ \quad (\operatorname{is\text{-}normed\text{-}field}(f) \wedge \operatorname{is\text{-}bilinear\text{-}bicharacter}(f, k)) \Rightarrow  \\ \quad \quad k(y, \operatorname{add}(f)(a, b)) = (k(y, a) \cdot k(y, b)) \end{array}}}
\vnbprop{c6-w-value}{\vnbstmt{\begin{array}{@{}l@{}} \forall f, k, d, b.; \forall u, v \in \operatorname{vec}(d). \\ \quad \operatorname{is\text{-}normed\text{-}field}(f) \wedge \\ \quad \operatorname{is\text{-}bilinear\text{-}bicharacter}(f, k) \wedge \\ \quad \operatorname{is\text{-}k\text{-}vector\text{-}space}(f, d) \wedge \\ \quad \operatorname{is\text{-}k\text{-}bilinear\text{-}form}(f, d, b) \Rightarrow \\ \quad \quad (\lambda {[}\operatorname{c6p\_}, \operatorname{c6q\_}] \in (\operatorname{vec}(d) \times \operatorname{vec}(d)).\; k(b(\operatorname{c6p\_}, \operatorname{c6q\_}), \operatorname{one}(f)))(u, v) \\ \quad \quad =\; k(b(u, v), \operatorname{one}(f)) \end{array}}}
\vnbprop{c6-w-type}{\vnbstmt{\begin{array}{@{}l@{}} \forall f, k, d, b. \\ \quad \operatorname{is\text{-}normed\text{-}field}(f) \wedge \\ \quad \operatorname{is\text{-}bilinear\text{-}bicharacter}(f, k) \wedge \\ \quad \operatorname{is\text{-}k\text{-}vector\text{-}space}(f, d) \wedge \\ \quad \operatorname{is\text{-}k\text{-}bilinear\text{-}form}(f, d, b) \Rightarrow \\ \quad \quad (\lambda {[}\operatorname{c6p\_}, \operatorname{c6q\_}] \in (\operatorname{vec}(d) \times \operatorname{vec}(d)).\; k(b(\operatorname{c6p\_}, \operatorname{c6q\_}), \operatorname{one}(f))) \\ \quad \quad \in\; ((\operatorname{vec}(d) \times \operatorname{vec}(d)) \to \operatorname{unit\text{-}circle}) \end{array}}}
\vnbprop{c6-w-69-left}{\vnbstmt{\begin{array}{@{}l@{}} \forall f, k, d, a.; \forall b \in \operatorname{carr}(f), u, v \in \operatorname{vec}(d). \\ \quad \operatorname{is\text{-}normed\text{-}field}(f) \wedge \\ \quad \operatorname{is\text{-}bilinear\text{-}bicharacter}(f, k) \wedge \\ \quad \operatorname{is\text{-}k\text{-}vector\text{-}space}(f, d) \wedge \\ \quad \operatorname{is\text{-}k\text{-}bilinear\text{-}form}(f, d, a) \Rightarrow \\ \quad \quad (\lambda {[}\operatorname{c6p\_}, \operatorname{c6q\_}] \in (\operatorname{vec}(d) \times \operatorname{vec}(d)).\; k(a(\operatorname{c6p\_}, \operatorname{c6q\_}), \operatorname{one}(f)))(\operatorname{act}(d)(b, u), v) \\ \quad \quad =\; k(a(u, v), b) \end{array}}}
\vnbprop{c6-w-69-right}{\vnbstmt{\begin{array}{@{}l@{}} \forall f, k, d, a.; \forall b \in \operatorname{carr}(f), u, v \in \operatorname{vec}(d). \\ \quad \operatorname{is\text{-}normed\text{-}field}(f) \wedge \\ \quad \operatorname{is\text{-}bilinear\text{-}bicharacter}(f, k) \wedge \\ \quad \operatorname{is\text{-}k\text{-}vector\text{-}space}(f, d) \wedge \\ \quad \operatorname{is\text{-}k\text{-}bilinear\text{-}form}(f, d, a) \Rightarrow \\ \quad \quad (\lambda {[}\operatorname{c6p\_}, \operatorname{c6q\_}] \in (\operatorname{vec}(d) \times \operatorname{vec}(d)).\; k(a(\operatorname{c6p\_}, \operatorname{c6q\_}), \operatorname{one}(f)))(u, \operatorname{act}(d)(b, v)) \\ \quad \quad =\; k(a(u, v), b) \end{array}}}
\vnbprop{ccr-6-32-converse}{\vnbstmt{\begin{array}{@{}l@{}} \forall f, k, d, a. \\ \quad \operatorname{is\text{-}normed\text{-}field}(f) \wedge \\ \quad \operatorname{is\text{-}bilinear\text{-}bicharacter}(f, k) \wedge \\ \quad \operatorname{is\text{-}k\text{-}vector\text{-}space}(f, d) \wedge \\ \quad \operatorname{is\text{-}k\text{-}bilinear\text{-}form}(f, d, a) \Rightarrow \\ \quad \quad \exists e. \\ \quad \quad \quad \operatorname{is\text{-}e\text{-}bilinear\text{-}bicharacter}(f, d, e) \wedge \\ \quad \quad \quad \forall b \in \operatorname{carr}(f), u, v \in \operatorname{vec}(d). \\ \quad \quad \quad \quad (e(\operatorname{act}(d)(b, u), v) = k(a(u, v), b)) \wedge (e(u, \operatorname{act}(d)(b, v)) = k(a(u, v), b)) \wedge \\ \quad \quad \quad \forall u, v \in \operatorname{vec}(d).\; (e(u, v) = k(a(u, v), \operatorname{one}(f))) \end{array}}}
\vnbprop{c6-char-finsum}{\vnbstmt{\begin{array}{@{}l@{}} \forall n \in \mathbb{N}, a, i, f, g. \\ \quad \operatorname{is\text{-}abelian\text{-}group}(a) \wedge \\ \quad \forall x \in \operatorname{carr}(a).\; (i(x) \in \mathbb{C}) \wedge \\ \quad \forall x, y \in \operatorname{carr}(a).\; (i(\operatorname{opr}(a)(x, y)) = (i(x) \cdot i(y))) \wedge \\ \quad i(\operatorname{iden}(a)) = 1 \wedge \\ \quad \forall z \in \operatorname{interval}(1, n).\; (f(z) \in \operatorname{carr}(a)) \wedge \\ \quad \forall z \in \operatorname{interval}(1, n).\; (g(z) = i(f(z))) \Rightarrow \\ \quad \quad i(\operatorname{finsum}(a, f, \operatorname{interval}(1, n))) \\ \quad \quad =\; \operatorname{finsum}(\operatorname{commutative\text{-}ring\text{-}multiplicative\text{-}cm}(\operatorname{normed\text{-}field\text{-}as\text{-}commutative\text{-}ring}(\operatorname{cc\text{-}normed\text{-}field})), \\ \quad \quad \quad g, \operatorname{interval}(1, n)) \end{array}}}
\vnbprop{c6-wij-value}{\vnbstmt{\begin{array}{@{}l@{}} \forall f, d, w.; \forall e, a \in \operatorname{vec}(d), y, x \in \operatorname{carr}(f). \\ \quad (\operatorname{is\text{-}normed\text{-}field}(f) \wedge (\operatorname{is\text{-}k\text{-}vector\text{-}space}(f, d) \wedge \operatorname{is\text{-}e\text{-}bilinear\text{-}bicharacter}(f, d, w))) \Rightarrow  \\ \quad \quad (\lambda {[}\operatorname{c6s\_}, \operatorname{c6t\_}] \in (\operatorname{carr}(f) \times \operatorname{carr}(f)).\; w(\operatorname{act}(d)(\operatorname{c6s\_}, e), \operatorname{act}(d)(\operatorname{c6t\_}, a)))(y, x) \\ \quad \quad =\; w(\operatorname{act}(d)(y, e), \operatorname{act}(d)(x, a)) \end{array}}}
\vnbprop{c6-wij-bilinear}{\vnbstmt{\begin{array}{@{}l@{}} \forall f, d, w.; \forall e, a \in \operatorname{vec}(d). \\ \quad (\operatorname{is\text{-}normed\text{-}field}(f) \wedge (\operatorname{is\text{-}k\text{-}vector\text{-}space}(f, d) \wedge \operatorname{is\text{-}e\text{-}bilinear\text{-}bicharacter}(f, d, w))) \Rightarrow  \\ \quad \quad \operatorname{is\text{-}bilinear\text{-}bicharacter}(f, \\ \quad \quad \quad (\lambda {[}\operatorname{c6s\_}, \operatorname{c6t\_}] \in (\operatorname{carr}(f) \times \operatorname{carr}(f)).\; w(\operatorname{act}(d)(\operatorname{c6s\_}, e), \operatorname{act}(d)(\operatorname{c6t\_}, a)))) \end{array}}}
\vnbprop{ccr-6-24-entry}{\vnbstmt{\begin{array}{@{}l@{}} \forall f, a, i.; \forall c \in \operatorname{carr}(f), d, w.; \forall e, b \in \operatorname{vec}(d). \\ \quad \operatorname{is\text{-}local\text{-}uniformizer}(f, a) \wedge \\ \quad \operatorname{is\text{-}nontrivial\text{-}character}(f, i) \wedge \\ \quad \neg (c = \operatorname{zero}(f)) \wedge \\ \quad \operatorname{is\text{-}k\text{-}vector\text{-}space}(f, d) \wedge \\ \quad \operatorname{is\text{-}e\text{-}bilinear\text{-}bicharacter}(f, d, w) \Rightarrow \\ \quad \quad \exists \mathit{c6a} \in \operatorname{carr}(f). \forall y, x \in \operatorname{carr}(f). \\ \quad \quad \quad w(\operatorname{act}(d)(y, e), \operatorname{act}(d)(x, b)) = \operatorname{bichar\text{-}of}(f, i, c)(\operatorname{mul}(f)(\mathit{c6a}, y), x) \wedge \\ \quad \quad \forall \mathit{c6a}, \mathit{c6b} \in \operatorname{carr}(f). \\ \quad \quad \quad \forall y, x \in \operatorname{carr}(f). \\ \quad \quad \quad \quad w(\operatorname{act}(d)(y, e), \operatorname{act}(d)(x, b)) = \operatorname{bichar\text{-}of}(f, i, c)(\operatorname{mul}(f)(\mathit{c6a}, y), x) \wedge \\ \quad \quad \quad \forall y, x \in \operatorname{carr}(f). \\ \quad \quad \quad \quad w(\operatorname{act}(d)(y, e), \operatorname{act}(d)(x, b)) = \operatorname{bichar\text{-}of}(f, i, c)(\operatorname{mul}(f)(\mathit{c6b}, y), x) \Rightarrow \\ \quad \quad \quad \quad \mathit{c6a} = \mathit{c6b} \end{array}}}
\vnbprop{c6-w-zero-l}{\vnbstmt{\begin{array}{@{}l@{}} \forall f, d, w.; \forall y \in \operatorname{vec}(d). \\ \quad (\operatorname{is\text{-}normed\text{-}field}(f) \wedge (\operatorname{is\text{-}k\text{-}vector\text{-}space}(f, d) \wedge \operatorname{is\text{-}e\text{-}bicharacter}(f, d, w))) \Rightarrow  \\ \quad \quad w(\operatorname{vzero}(d), y) = 1 \end{array}}}
\vnbprop{c6-w-zero-r}{\vnbstmt{\begin{array}{@{}l@{}} \forall f, d, w.; \forall y \in \operatorname{vec}(d). \\ \quad (\operatorname{is\text{-}normed\text{-}field}(f) \wedge (\operatorname{is\text{-}k\text{-}vector\text{-}space}(f, d) \wedge \operatorname{is\text{-}e\text{-}bicharacter}(f, d, w))) \Rightarrow  \\ \quad \quad w(y, \operatorname{vzero}(d)) = 1 \end{array}}}
\vnbprop{c6-w-finsum-l}{\vnbstmt{\begin{array}{@{}l@{}} \forall a, d, w.; \forall n \in \mathbb{N}, f.; \forall y \in \operatorname{vec}(d). \\ \quad \operatorname{is\text{-}normed\text{-}field}(a) \wedge \\ \quad \operatorname{is\text{-}k\text{-}vector\text{-}space}(a, d) \wedge \\ \quad \operatorname{is\text{-}e\text{-}bicharacter}(a, d, w) \wedge \\ \quad \forall z \in \operatorname{interval}(1, n).\; (f(z) \in \operatorname{vec}(d)) \Rightarrow \\ \quad \quad w(\operatorname{finsum}(\operatorname{module\text{-}vector\text{-}ag}(d), f, \operatorname{interval}(1, n)), y) \\ \quad \quad =\; \operatorname{finsum}(\operatorname{commutative\text{-}ring\text{-}multiplicative\text{-}cm}(\operatorname{normed\text{-}field\text{-}as\text{-}commutative\text{-}ring}(\operatorname{cc\text{-}normed\text{-}field})), \\ \quad \quad \quad (\lambda i \in \operatorname{interval}(1, n).\; w(f(i), y)), \operatorname{interval}(1, n)) \end{array}}}
\vnbprop{c6-w-finsum-r}{\vnbstmt{\begin{array}{@{}l@{}} \forall a, d, w.; \forall n \in \mathbb{N}, f.; \forall y \in \operatorname{vec}(d). \\ \quad \operatorname{is\text{-}normed\text{-}field}(a) \wedge \\ \quad \operatorname{is\text{-}k\text{-}vector\text{-}space}(a, d) \wedge \\ \quad \operatorname{is\text{-}e\text{-}bicharacter}(a, d, w) \wedge \\ \quad \forall z \in \operatorname{interval}(1, n).\; (f(z) \in \operatorname{vec}(d)) \Rightarrow \\ \quad \quad w(y, \operatorname{finsum}(\operatorname{module\text{-}vector\text{-}ag}(d), f, \operatorname{interval}(1, n))) \\ \quad \quad =\; \operatorname{finsum}(\operatorname{commutative\text{-}ring\text{-}multiplicative\text{-}cm}(\operatorname{normed\text{-}field\text{-}as\text{-}commutative\text{-}ring}(\operatorname{cc\text{-}normed\text{-}field})), \\ \quad \quad \quad (\lambda i \in \operatorname{interval}(1, n).\; w(y, f(i))), \operatorname{interval}(1, n)) \end{array}}}
\vnbprop{ccr-6-24-expand}{\vnbstmt{\begin{array}{@{}l@{}} \forall a, d, w.; \forall n \in \mathbb{N}, f, h. \\ \quad \operatorname{is\text{-}normed\text{-}field}(a) \wedge \\ \quad \operatorname{is\text{-}k\text{-}vector\text{-}space}(a, d) \wedge \\ \quad \operatorname{is\text{-}e\text{-}bicharacter}(a, d, w) \wedge \\ \quad \forall z \in \operatorname{interval}(1, n).\; (f(z) \in \operatorname{vec}(d)) \wedge \\ \quad \forall z \in \operatorname{interval}(1, n).\; (h(z) \in \operatorname{vec}(d)) \Rightarrow \\ \quad \quad w(\operatorname{finsum}(\operatorname{module\text{-}vector\text{-}ag}(d), f, \operatorname{interval}(1, n)), \operatorname{finsum}(\operatorname{module\text{-}vector\text{-}ag}(d), h, \operatorname{interval}(1, n))) \\ \quad \quad =\; \operatorname{finsum}(\operatorname{commutative\text{-}ring\text{-}multiplicative\text{-}cm}(\operatorname{normed\text{-}field\text{-}as\text{-}commutative\text{-}ring}(\operatorname{cc\text{-}normed\text{-}field})), \\ \quad \quad \quad (\lambda i \in \operatorname{interval}(1, n).\; \\ \quad \quad \quad \quad \operatorname{finsum}(\operatorname{commutative\text{-}ring\text{-}multiplicative\text{-}cm}(\operatorname{normed\text{-}field\text{-}as\text{-}commutative\text{-}ring}(\operatorname{cc\text{-}normed\text{-}field})), \\ \quad \quad \quad \quad \quad (\lambda j \in \operatorname{interval}(1, n).\; w(f(i), h(j))), \operatorname{interval}(1, n))), \\ \quad \quad \quad \operatorname{interval}(1, n)) \end{array}}}

\vnbbookfile{theorem-library/ccr6-f}{theorem-library/ccr6-f}
\noindent{\small ccr6-f.scm -- wave of 2026-11-09/10 (agent CCR-1109-D4), wired 2026-11-10 from scratchpad/ccr1109d4/ccr6-f.scm. Library: complex-analysis.}\par
\vnbprop{c6f-wk-inj}{\vnbstmt{\begin{array}{@{}l@{}} \forall f, k.; \forall x, y \in \operatorname{carr}(f). \\ \quad \operatorname{is\text{-}normed\text{-}field}(f) \wedge \\ \quad \operatorname{is\text{-}bilinear\text{-}bicharacter}(f, k) \wedge \\ \quad \operatorname{is\text{-}faithful\text{-}bicharacter}(f, k) \wedge \\ \quad \forall b \in \operatorname{carr}(f).\; (k(x, b) = k(y, b)) \Rightarrow \\ \quad \quad x = y \end{array}}}
\vnbprop{c6f-lambda-raw}{\vnbstmt{\begin{array}{@{}l@{}} \forall f, k, d, e. \\ \quad \operatorname{is\text{-}normed\text{-}field}(f) \wedge \\ \quad \operatorname{is\text{-}bilinear\text{-}bicharacter}(f, k) \wedge \\ \quad \operatorname{is\text{-}faithful\text{-}bicharacter}(f, k) \wedge \\ \quad \operatorname{is\text{-}k\text{-}vector\text{-}space}(f, d) \wedge \\ \quad \operatorname{is\text{-}e\text{-}bilinear\text{-}bicharacter}(f, d, e) \Rightarrow \\ \quad \quad \exists l \in ((\operatorname{vec}(d) \times \operatorname{vec}(d)) \to \operatorname{carr}(f)). \forall u, v \in \operatorname{vec}(d), b \in \operatorname{carr}(f). \\ \quad \quad \quad e(\operatorname{act}(d)(b, u), v) = k(l(u, v), b) \end{array}}}
\vnbprop{ccr-6-32-exists}{\vnbstmt{\begin{array}{@{}l@{}} \forall f, k, d, e. \\ \quad \operatorname{is\text{-}normed\text{-}field}(f) \wedge \\ \quad \operatorname{is\text{-}bilinear\text{-}bicharacter}(f, k) \wedge \\ \quad \operatorname{is\text{-}faithful\text{-}bicharacter}(f, k) \wedge \\ \quad \operatorname{is\text{-}k\text{-}vector\text{-}space}(f, d) \wedge \\ \quad \operatorname{is\text{-}e\text{-}bilinear\text{-}bicharacter}(f, d, e) \Rightarrow \\ \quad \quad \exists l. \\ \quad \quad \quad \operatorname{is\text{-}k\text{-}bilinear\text{-}form}(f, d, l) \wedge \\ \quad \quad \quad \forall b \in \operatorname{carr}(f), u, v \in \operatorname{vec}(d). \\ \quad \quad \quad \quad (e(\operatorname{act}(d)(b, u), v) = k(l(u, v), b)) \wedge (e(u, \operatorname{act}(d)(b, v)) = k(l(u, v), b)) \end{array}}}
\vnbprop{ccr-6-32}{\vnbstmt{\begin{array}{@{}l@{}} \forall f, k, d, e. \\ \quad \operatorname{is\text{-}normed\text{-}field}(f) \wedge \\ \quad \operatorname{is\text{-}bilinear\text{-}bicharacter}(f, k) \wedge \\ \quad \operatorname{is\text{-}faithful\text{-}bicharacter}(f, k) \wedge \\ \quad \operatorname{is\text{-}k\text{-}vector\text{-}space}(f, d) \wedge \\ \quad \operatorname{is\text{-}e\text{-}bilinear\text{-}bicharacter}(f, d, e) \Rightarrow \\ \quad \quad \exists l. \\ \quad \quad \quad \operatorname{is\text{-}k\text{-}bilinear\text{-}form}(f, d, l) \wedge \\ \quad \quad \quad \forall b \in \operatorname{carr}(f), u, v \in \operatorname{vec}(d). \\ \quad \quad \quad \quad (e(\operatorname{act}(d)(b, u), v) = k(l(u, v), b)) \wedge (e(u, \operatorname{act}(d)(b, v)) = k(l(u, v), b)) \wedge \\ \quad \quad \quad \forall m.; \forall u, v \in \operatorname{vec}(d). \\ \quad \quad \quad \quad \operatorname{is\text{-}k\text{-}bilinear\text{-}form}(f, d, m) \wedge \\ \quad \quad \quad \quad \forall b \in \operatorname{carr}(f), u, v \in \operatorname{vec}(d). \\ \quad \quad \quad \quad \quad (e(\operatorname{act}(d)(b, u), v) = k(m(u, v), b)) \wedge (e(u, \operatorname{act}(d)(b, v)) = k(m(u, v), b)) \Rightarrow \\ \quad \quad \quad \quad \quad m(u, v) = l(u, v) \end{array}}}
\vnbprop{ccr-6-24}{\vnbstmt{\begin{array}{@{}l@{}} \forall f, a, b.; \forall c \in \operatorname{carr}(f), d, w.; \forall n \in \mathbb{N}, e. \\ \quad \operatorname{is\text{-}local\text{-}uniformizer}(f, a) \wedge \\ \quad \operatorname{is\text{-}nontrivial\text{-}character}(f, b) \wedge \\ \quad \neg (c = \operatorname{zero}(f)) \wedge \\ \quad \operatorname{is\text{-}k\text{-}vector\text{-}space}(f, d) \wedge \\ \quad \operatorname{is\text{-}e\text{-}bilinear\text{-}bicharacter}(f, d, w) \wedge \\ \quad \forall z \in \operatorname{interval}(1, n).\; (e(z) \in \operatorname{vec}(d)) \Rightarrow \\ \quad \quad \exists m \in \operatorname{mat}(n, n, \operatorname{carr}(f)). \\ \quad \quad \quad \forall i, j \in \operatorname{interval}(1, n), y, x \in \operatorname{carr}(f). \\ \quad \quad \quad \quad w(\operatorname{act}(d)(y, e(i)), \operatorname{act}(d)(x, e(j))) = \operatorname{bichar\text{-}of}(f, b, c)(\operatorname{mul}(f)(\operatorname{entry}(m, i, j), y), x) \wedge \\ \quad \quad \quad \forall p \in \operatorname{mat}(n, n, \operatorname{carr}(f)). \\ \quad \quad \quad \quad \forall i, j \in \operatorname{interval}(1, n), y, x \in \operatorname{carr}(f). \\ \quad \quad \quad \quad \quad w(\operatorname{act}(d)(y, e(i)), \operatorname{act}(d)(x, e(j))) = \operatorname{bichar\text{-}of}(f, b, c)(\operatorname{mul}(f)(\operatorname{entry}(p, i, j), y), x) \Rightarrow \\ \quad \quad \quad \quad \quad p = m \end{array}}}
\vnbprop{ccr-6-24-66}{\vnbstmt{\begin{array}{@{}l@{}} \forall f, d, w, k.; \forall n \in \mathbb{N}, e, m, y, x. \\ \quad \operatorname{is\text{-}normed\text{-}field}(f) \wedge \\ \quad \operatorname{is\text{-}k\text{-}vector\text{-}space}(f, d) \wedge \\ \quad \operatorname{is\text{-}e\text{-}bicharacter}(f, d, w) \wedge \\ \quad \forall z \in \operatorname{interval}(1, n).\; (e(z) \in \operatorname{vec}(d)) \wedge \\ \quad \forall i, j \in \operatorname{interval}(1, n), s, t \in \operatorname{carr}(f). \\ \quad \quad w(\operatorname{act}(d)(s, e(i)), \operatorname{act}(d)(t, e(j))) = k(\operatorname{mul}(f)(\operatorname{entry}(m, i, j), s), t) \wedge \\ \quad \forall z \in \operatorname{interval}(1, n).\; (y(z) \in \operatorname{carr}(f)) \wedge \\ \quad \forall z \in \operatorname{interval}(1, n).\; (x(z) \in \operatorname{carr}(f)) \Rightarrow \\ \quad \quad w(\operatorname{finsum}(\operatorname{module\text{-}vector\text{-}ag}(d), \\ \quad \quad \quad \quad (\lambda \mathit{c6fi} \in \operatorname{interval}(1, n).\; \operatorname{act}(d)(y(\mathit{c6fi}), e(\mathit{c6fi}))), \\ \quad \quad \quad \quad \operatorname{interval}(1, n)), \\ \quad \quad \quad \operatorname{finsum}(\operatorname{module\text{-}vector\text{-}ag}(d), \\ \quad \quad \quad \quad (\lambda \mathit{c6fj} \in \operatorname{interval}(1, n).\; \operatorname{act}(d)(x(\mathit{c6fj}), e(\mathit{c6fj}))), \\ \quad \quad \quad \quad \operatorname{interval}(1, n))) \\ \quad \quad =\; \operatorname{finsum}(\operatorname{commutative\text{-}ring\text{-}multiplicative\text{-}cm}(\operatorname{normed\text{-}field\text{-}as\text{-}commutative\text{-}ring}(\operatorname{cc\text{-}normed\text{-}field})), \\ \quad \quad \quad (\lambda \mathit{c6fi} \in \operatorname{interval}(1, n).\; \\ \quad \quad \quad \quad \operatorname{finsum}(\operatorname{commutative\text{-}ring\text{-}multiplicative\text{-}cm}(\operatorname{normed\text{-}field\text{-}as\text{-}commutative\text{-}ring}(\operatorname{cc\text{-}normed\text{-}field})), \\ \quad \quad \quad \quad \quad (\lambda \mathit{c6fj} \in \operatorname{interval}(1, n).\; k(\operatorname{mul}(f)(\operatorname{entry}(m, \mathit{c6fi}, \mathit{c6fj}), y(\mathit{c6fi})), x(\mathit{c6fj}))), \\ \quad \quad \quad \quad \quad \operatorname{interval}(1, n))), \\ \quad \quad \quad \operatorname{interval}(1, n)) \end{array}}}
\vnbprop{ccr-6-28-ball}{\vnbstmt{\begin{array}{@{}l@{}} \forall f, k.; \forall n \in \mathbb{N}, m, y, x. \\ \quad \operatorname{is\text{-}nv\text{-}field}(f) \wedge \\ \quad \operatorname{is\text{-}unipolar}(f, k) \wedge \\ \quad \forall i, j \in \operatorname{interval}(1, n). \\ \quad \quad \operatorname{entry}(m, i, j) \in \operatorname{nv\text{-}integers}(f) \wedge \\ \quad \forall z \in \operatorname{interval}(1, n).\; (y(z) \in \operatorname{nv\text{-}integers}(f)) \wedge \\ \quad \forall z \in \operatorname{interval}(1, n).\; (x(z) \in \operatorname{nv\text{-}integers}(f)) \Rightarrow \\ \quad \quad \operatorname{finsum}(\operatorname{commutative\text{-}ring\text{-}multiplicative\text{-}cm}(\operatorname{normed\text{-}field\text{-}as\text{-}commutative\text{-}ring}(\operatorname{cc\text{-}normed\text{-}field})), \\ \quad \quad \quad (\lambda \mathit{c6fi} \in \operatorname{interval}(1, n).\; \\ \quad \quad \quad \quad \operatorname{finsum}(\operatorname{commutative\text{-}ring\text{-}multiplicative\text{-}cm}(\operatorname{normed\text{-}field\text{-}as\text{-}commutative\text{-}ring}(\operatorname{cc\text{-}normed\text{-}field})), \\ \quad \quad \quad \quad \quad (\lambda \mathit{c6fj} \in \operatorname{interval}(1, n).\; k(\operatorname{mul}(f)(\operatorname{entry}(m, \mathit{c6fi}, \mathit{c6fj}), y(\mathit{c6fi})), x(\mathit{c6fj}))), \\ \quad \quad \quad \quad \quad \operatorname{interval}(1, n))), \\ \quad \quad \quad \operatorname{interval}(1, n)) \\ \quad \quad =\; 1 \end{array}}}
\vnbprop{c6f-char-finprod}{\vnbstmt{\begin{array}{@{}l@{}} \forall n \in \mathbb{N}, f, s. \\ \quad (\operatorname{is\text{-}normed\text{-}field}(f) \wedge \forall k \in \operatorname{interval}(1, n).\; \operatorname{is\text{-}add\text{-}character}(f, s(k))) \Rightarrow  \\ \quad \quad \operatorname{is\text{-}add\text{-}character}(f, \\ \quad \quad \quad (\lambda h \in \operatorname{carr}(f).\; \\ \quad \quad \quad \quad \operatorname{finsum}(\operatorname{commutative\text{-}ring\text{-}multiplicative\text{-}cm}(\operatorname{normed\text{-}field\text{-}as\text{-}commutative\text{-}ring}(\operatorname{cc\text{-}normed\text{-}field})), \\ \quad \quad \quad \quad \quad (\lambda \mathit{c6fk} \in \operatorname{interval}(1, n).\; s(\mathit{c6fk})(h)), \operatorname{interval}(1, n)))) \end{array}}}

\vnbbookfile{theorem-library/ccr-adj-integral}{theorem-library/ccr-adj-integral}
\noindent{\small ccr-adj-integral.scm -- wave of 2026-11-09/10 (agent CCR-1109-E), wired 2026-11-10 from scratchpad/ccr1109e/ccr-adj-integral.scm. Library: complex-analysis.}\par
\vnbprop{cai-minor-bound}{\vnbstmt{\begin{array}{@{}l@{}} \forall k.; \forall n \in \mathbb{N}, a \in \operatorname{mat}((n + 1), (n + 1), \operatorname{carr}(k)), p, q \in \operatorname{interval}(1, (n + 1)). \\ \quad \operatorname{is\text{-}nv\text{-}field}(k) \wedge \\ \quad \forall i, j \in \operatorname{interval}(1, (n + 1)). \\ \quad \quad \operatorname{entry}(a, i, j) \in \operatorname{nv\text{-}integers}(k) \Rightarrow \\ \quad \quad \operatorname{is\text{-}mat\text{-}bound}(k, n, n, \operatorname{minor}(a, p, q, n), 1) \end{array}}}
\vnbprop{adj-entries-integral}{\vnbstmt{\begin{array}{@{}l@{}} \forall k.; \forall n \in \mathbb{N}, a \in \operatorname{mat}((n + 1), (n + 1), \operatorname{carr}(k)), i, j \in \operatorname{interval}(1, (n + 1)). \\ \quad \operatorname{is\text{-}nv\text{-}field}(k) \wedge \\ \quad \forall i, j \in \operatorname{interval}(1, (n + 1)). \\ \quad \quad \operatorname{entry}(a, i, j) \in \operatorname{nv\text{-}integers}(k) \Rightarrow \\ \quad \quad \operatorname{entry}(\operatorname{adj}(\operatorname{normed\text{-}field\text{-}as\text{-}commutative\text{-}ring}(k), n, a), i, j) \in \operatorname{nv\text{-}integers}(k) \end{array}}}
\vnbprop{ccr-6-27}{\vnbstmt{\begin{array}{@{}l@{}} \forall k.; \forall n \in \mathbb{N}, a \in \operatorname{mat}(n, n, \operatorname{carr}(k)). \\ \quad \operatorname{is\text{-}nv\text{-}field}(k) \wedge \\ \quad \forall i, j \in \operatorname{interval}(1, n). \\ \quad \quad \operatorname{entry}(a, i, j) \in \operatorname{nv\text{-}integers}(k) \wedge \\ \quad \operatorname{det}(\operatorname{normed\text{-}field\text{-}as\text{-}commutative\text{-}ring}(k), n, a) \in \operatorname{nv\text{-}units}(k, \operatorname{nv\text{-}integers}(k)) \Rightarrow \\ \quad \quad \exists b \in \operatorname{mat}(n, n, \operatorname{carr}(k)). \\ \quad \quad \quad \forall i, j \in \operatorname{interval}(1, n). \\ \quad \quad \quad \quad \operatorname{entry}(b, i, j) \in \operatorname{nv\text{-}integers}(k) \wedge \\ \quad \quad \quad \operatorname{matmul}(\operatorname{normed\text{-}field\text{-}as\text{-}commutative\text{-}ring}(k), a, b) = \operatorname{identmat}(\operatorname{normed\text{-}field\text{-}as\text{-}commutative\text{-}ring}(k), n) \wedge \\ \quad \quad \quad \operatorname{matmul}(\operatorname{normed\text{-}field\text{-}as\text{-}commutative\text{-}ring}(k), b, a) = \operatorname{identmat}(\operatorname{normed\text{-}field\text{-}as\text{-}commutative\text{-}ring}(k), n) \end{array}}}

\vnbbookfile{theorem-library/ccr6-g}{theorem-library/ccr6-g}
\noindent{\small ccr6-g.scm -- wave of 2026-11-09/10 (agent CCR-1109-E), wired 2026-11-10 from scratchpad/ccr1109e/ccr6-g.scm. Library: complex-analysis.}\par
\vnbprop{c6g-wk-finsum-r}{\vnbstmt{\begin{array}{@{}l@{}} \forall a, k.; \forall t \in \operatorname{carr}(a), n \in \mathbb{N}, f. \\ \quad (\operatorname{is\text{-}normed\text{-}field}(a) \wedge (\operatorname{is\text{-}bicharacter}(a, k) \wedge \forall z \in \operatorname{interval}(1, n).\; (f(z) \in \operatorname{carr}(a)))) \Rightarrow  \\ \quad \quad k(t, \operatorname{finsum}(\operatorname{ring\text{-}additive\text{-}ag}(\operatorname{normed\text{-}field\text{-}as\text{-}commutative\text{-}ring}(a)), f, \operatorname{interval}(1, n))) \\ \quad \quad =\; \operatorname{finsum}(\operatorname{commutative\text{-}ring\text{-}multiplicative\text{-}cm}(\operatorname{normed\text{-}field\text{-}as\text{-}commutative\text{-}ring}(\operatorname{cc\text{-}normed\text{-}field})), \\ \quad \quad \quad (\lambda j \in \operatorname{interval}(1, n).\; k(t, f(j))), \operatorname{interval}(1, n)) \end{array}}}
\vnbprop{c6g-row-prod}{\vnbstmt{\begin{array}{@{}l@{}} \forall f, a.; \forall n \in \mathbb{N}, m \in \operatorname{mat}(n, n, \operatorname{carr}(f)), x.; \forall k \in \operatorname{interval}(1, n), t \in \operatorname{carr}(f). \\ \quad (\operatorname{is\text{-}normed\text{-}field}(f) \wedge (\operatorname{is\text{-}bicharacter}(f, a) \wedge \forall z \in \operatorname{interval}(1, n).\; (x(z) \in \operatorname{carr}(f)))) \Rightarrow  \\ \quad \quad \operatorname{finsum}(\operatorname{commutative\text{-}ring\text{-}multiplicative\text{-}cm}(\operatorname{normed\text{-}field\text{-}as\text{-}commutative\text{-}ring}(\operatorname{cc\text{-}normed\text{-}field})), \\ \quad \quad \quad (\lambda i \in \operatorname{interval}(1, n).\; \\ \quad \quad \quad \quad \operatorname{finsum}(\operatorname{commutative\text{-}ring\text{-}multiplicative\text{-}cm}(\operatorname{normed\text{-}field\text{-}as\text{-}commutative\text{-}ring}(\operatorname{cc\text{-}normed\text{-}field})), \\ \quad \quad \quad \quad \quad (\lambda j \in \operatorname{interval}(1, n).\; \\ \quad \quad \quad \quad \quad \quad a(\operatorname{mul}(f)(\operatorname{entry}(m, i, j), (\lambda \mathit{c6gi} \in \operatorname{interval}(1, n).\; \left(\text{if }(\mathit{c6gi} = k)\text{ then }t\text{ else }\operatorname{zero}(f)\right))(i)), x(j))), \\ \quad \quad \quad \quad \quad \operatorname{interval}(1, n))), \\ \quad \quad \quad \operatorname{interval}(1, n)) \\ \quad \quad =\; \operatorname{finsum}(\operatorname{commutative\text{-}ring\text{-}multiplicative\text{-}cm}(\operatorname{normed\text{-}field\text{-}as\text{-}commutative\text{-}ring}(\operatorname{cc\text{-}normed\text{-}field})), \\ \quad \quad \quad (\lambda j \in \operatorname{interval}(1, n).\; a(\operatorname{mul}(f)(\operatorname{entry}(m, k, j), t), x(j))), \\ \quad \quad \quad \operatorname{interval}(1, n)) \end{array}}}
\vnbprop{ccr-6-28-polar}{\vnbstmt{\begin{array}{@{}l@{}} \forall f, k.; \forall n \in \mathbb{N}, m \in \operatorname{mat}(n, n, \operatorname{carr}(f)), x.; \forall z \in \operatorname{interval}(1, n). \\ \quad \operatorname{is\text{-}nv\text{-}field}(f) \wedge \\ \quad \operatorname{is\text{-}unipolar}(f, k) \wedge \\ \quad \forall i, j \in \operatorname{interval}(1, n). \\ \quad \quad \operatorname{entry}(m, i, j) \in \operatorname{nv\text{-}integers}(f) \wedge \\ \quad \operatorname{det}(\operatorname{normed\text{-}field\text{-}as\text{-}commutative\text{-}ring}(f), n, m) \in \operatorname{nv\text{-}units}(f, \operatorname{nv\text{-}integers}(f)) \wedge \\ \quad \forall z \in \operatorname{interval}(1, n).\; (x(z) \in \operatorname{carr}(f)) \wedge \\ \quad \forall y. \\ \quad \quad \forall z \in \operatorname{interval}(1, n).\; (y(z) \in \operatorname{nv\text{-}integers}(f)) \Rightarrow  \\ \quad \quad \quad \operatorname{finsum}(\operatorname{commutative\text{-}ring\text{-}multiplicative\text{-}cm}(\operatorname{normed\text{-}field\text{-}as\text{-}commutative\text{-}ring}(\operatorname{cc\text{-}normed\text{-}field})), \\ \quad \quad \quad \quad (\lambda \mathit{c6fi} \in \operatorname{interval}(1, n).\; \\ \quad \quad \quad \quad \quad \operatorname{finsum}(\operatorname{commutative\text{-}ring\text{-}multiplicative\text{-}cm}(\operatorname{normed\text{-}field\text{-}as\text{-}commutative\text{-}ring}(\operatorname{cc\text{-}normed\text{-}field})), \\ \quad \quad \quad \quad \quad \quad (\lambda \mathit{c6fj} \in \operatorname{interval}(1, n).\; k(\operatorname{mul}(f)(\operatorname{entry}(m, \mathit{c6fi}, \mathit{c6fj}), y(\mathit{c6fi})), x(\mathit{c6fj}))), \\ \quad \quad \quad \quad \quad \quad \operatorname{interval}(1, n))), \\ \quad \quad \quad \quad \operatorname{interval}(1, n)) \\ \quad \quad \quad =\; 1 \Rightarrow \\ \quad \quad x(z) \in \operatorname{nv\text{-}integers}(f) \end{array}}}
\vnbprop{ccr-6-28}{\vnbstmt{\begin{array}{@{}l@{}} \forall f, k.; \forall n \in \mathbb{N}, m \in \operatorname{mat}(n, n, \operatorname{carr}(f)), x. \\ \quad \operatorname{is\text{-}nv\text{-}field}(f) \wedge \\ \quad \operatorname{is\text{-}unipolar}(f, k) \wedge \\ \quad \forall i, j \in \operatorname{interval}(1, n). \\ \quad \quad \operatorname{entry}(m, i, j) \in \operatorname{nv\text{-}integers}(f) \wedge \\ \quad \operatorname{det}(\operatorname{normed\text{-}field\text{-}as\text{-}commutative\text{-}ring}(f), n, m) \in \operatorname{nv\text{-}units}(f, \operatorname{nv\text{-}integers}(f)) \wedge \\ \quad \forall z \in \operatorname{interval}(1, n).\; (x(z) \in \operatorname{carr}(f)) \Rightarrow \\ \quad \quad \forall y. \\ \quad \quad \quad \forall z \in \operatorname{interval}(1, n).\; (y(z) \in \operatorname{nv\text{-}integers}(f)) \Rightarrow  \\ \quad \quad \quad \quad \operatorname{finsum}(\operatorname{commutative\text{-}ring\text{-}multiplicative\text{-}cm}(\operatorname{normed\text{-}field\text{-}as\text{-}commutative\text{-}ring}(\operatorname{cc\text{-}normed\text{-}field})), \\ \quad \quad \quad \quad \quad (\lambda \mathit{c6fi} \in \operatorname{interval}(1, n).\; \\ \quad \quad \quad \quad \quad \quad \operatorname{finsum}(\operatorname{commutative\text{-}ring\text{-}multiplicative\text{-}cm}(\operatorname{normed\text{-}field\text{-}as\text{-}commutative\text{-}ring}(\operatorname{cc\text{-}normed\text{-}field})), \\ \quad \quad \quad \quad \quad \quad \quad (\lambda \mathit{c6fj} \in \operatorname{interval}(1, n).\; k(\operatorname{mul}(f)(\operatorname{entry}(m, \mathit{c6fi}, \mathit{c6fj}), y(\mathit{c6fi})), x(\mathit{c6fj}))), \\ \quad \quad \quad \quad \quad \quad \quad \operatorname{interval}(1, n))), \\ \quad \quad \quad \quad \quad \operatorname{interval}(1, n)) \\ \quad \quad \quad \quad =\; 1 \\ \quad \quad \Leftrightarrow\; \forall z \in \operatorname{interval}(1, n).\; (x(z) \in \operatorname{nv\text{-}integers}(f)) \end{array}}}

\vnbbookfile{theorem-library/ccr6-h}{theorem-library/ccr6-h}
\noindent{\small ccr6-h.scm -- wave of 2026-11-09/10 (agent CCR-1109-E2), wired 2026-11-10 from scratchpad/ccr1109e2/ccr6-h.scm. Library: complex-analysis.}\par
\vnbprop{c6h-wk-finsum-l}{\vnbstmt{\begin{array}{@{}l@{}} \forall a, k.; \forall t \in \operatorname{carr}(a), n \in \mathbb{N}, f. \\ \quad (\operatorname{is\text{-}normed\text{-}field}(a) \wedge (\operatorname{is\text{-}bicharacter}(a, k) \wedge \forall z \in \operatorname{interval}(1, n).\; (f(z) \in \operatorname{carr}(a)))) \Rightarrow  \\ \quad \quad k(\operatorname{finsum}(\operatorname{ring\text{-}additive\text{-}ag}(\operatorname{normed\text{-}field\text{-}as\text{-}commutative\text{-}ring}(a)), f, \operatorname{interval}(1, n)), t) \\ \quad \quad =\; \operatorname{finsum}(\operatorname{commutative\text{-}ring\text{-}multiplicative\text{-}cm}(\operatorname{normed\text{-}field\text{-}as\text{-}commutative\text{-}ring}(\operatorname{cc\text{-}normed\text{-}field})), \\ \quad \quad \quad (\lambda j \in \operatorname{interval}(1, n).\; k(f(j), t)), \operatorname{interval}(1, n)) \end{array}}}
\vnbprop{c6h-kn-vector-space}{\vnbstmt{\begin{array}{@{}l@{}} \forall f.; \forall n \in \mathbb{N}. \\ \quad \operatorname{is\text{-}normed\text{-}field}(f) \Rightarrow  \\ \quad \quad \operatorname{is\text{-}k\text{-}vector\text{-}space}(f, \operatorname{ring\text{-}pow}(\operatorname{normed\text{-}field\text{-}as\text{-}commutative\text{-}ring}(f), n)) \end{array}}}
\vnbprop{c6h-form-bilinear}{\vnbstmt{\begin{array}{@{}l@{}} \forall f.; \forall n \in \mathbb{N}, x \in \operatorname{mat}(n, n, \operatorname{carr}(f)). \\ \quad (\operatorname{is\text{-}normed\text{-}field}(f) \wedge (1 \leq n)) \Rightarrow  \\ \quad \quad \operatorname{is\text{-}k\text{-}bilinear\text{-}form}(f, \\ \quad \quad \quad \operatorname{ring\text{-}pow}(\operatorname{normed\text{-}field\text{-}as\text{-}commutative\text{-}ring}(f), n), \\ \quad \quad \quad (\lambda {[}\operatorname{c6hp\_}, \operatorname{c6hq\_}] \in (\operatorname{vec}(\operatorname{ring\text{-}pow}(\operatorname{normed\text{-}field\text{-}as\text{-}commutative\text{-}ring}(f), n)) \times \operatorname{vec}(\operatorname{ring\text{-}pow}(\operatorname{normed\text{-}field\text{-}as\text{-}commutative\text{-}ring}(f), n))).\; \\ \quad \quad \quad \quad \operatorname{entry}(\operatorname{matmul}(\operatorname{normed\text{-}field\text{-}as\text{-}commutative\text{-}ring}(f), \\ \quad \quad \quad \quad \quad \quad \operatorname{ring\text{-}row}(\operatorname{normed\text{-}field\text{-}as\text{-}commutative\text{-}ring}(f), n, \operatorname{c6hp\_}), \\ \quad \quad \quad \quad \quad \quad \operatorname{matmul}(\operatorname{normed\text{-}field\text{-}as\text{-}commutative\text{-}ring}(f), x, \\ \quad \quad \quad \quad \quad \quad \quad \operatorname{transpose}(\operatorname{ring\text{-}row}(\operatorname{normed\text{-}field\text{-}as\text{-}commutative\text{-}ring}(f), n, \operatorname{c6hq\_}), 1, n))), \\ \quad \quad \quad \quad \quad 1, 1))) \end{array}}}
\vnbprop{c6h-row}{\vnbstmt{\begin{array}{@{}l@{}} \forall f, k.; \forall n \in \mathbb{N}, a \in \operatorname{mat}(n, n, \operatorname{carr}(f)), x, i, c. \\ \quad \operatorname{is\text{-}normed\text{-}field}(f) \wedge \\ \quad \operatorname{is\text{-}bilinear\text{-}bicharacter}(f, k) \wedge \\ \quad 1 \leq n \wedge \\ \quad x \in \operatorname{vec}(\operatorname{ring\text{-}pow}(\operatorname{normed\text{-}field\text{-}as\text{-}commutative\text{-}ring}(f), n)) \wedge \\ \quad i \in \operatorname{interval}(1, n) \wedge \\ \quad c \in \operatorname{carr}(f) \Rightarrow \\ \quad \quad k(\operatorname{mul}(\operatorname{normed\text{-}field\text{-}as\text{-}commutative\text{-}ring}(f))(c, \\ \quad \quad \quad \quad \operatorname{entry}(\operatorname{matmul}(\operatorname{normed\text{-}field\text{-}as\text{-}commutative\text{-}ring}(f), a, \\ \quad \quad \quad \quad \quad \quad \operatorname{transpose}(\operatorname{ring\text{-}row}(\operatorname{normed\text{-}field\text{-}as\text{-}commutative\text{-}ring}(f), n, x), 1, n)), \\ \quad \quad \quad \quad \quad i, 1)), \\ \quad \quad \quad \operatorname{one}(f)) \\ \quad \quad =\; \operatorname{finsum}(\operatorname{commutative\text{-}ring\text{-}multiplicative\text{-}cm}(\operatorname{normed\text{-}field\text{-}as\text{-}commutative\text{-}ring}(\operatorname{cc\text{-}normed\text{-}field})), \\ \quad \quad \quad (\lambda j \in \operatorname{interval}(1, n).\; k(\operatorname{mul}(f)(\operatorname{entry}(a, i, j), c), x(j))), \\ \quad \quad \quad \operatorname{interval}(1, n)) \end{array}}}
\vnbprop{ccr-6-24-converse-66}{\vnbstmt{\begin{array}{@{}l@{}} \forall f, k.; \forall n \in \mathbb{N}, a \in \operatorname{mat}(n, n, \operatorname{carr}(f)), y, x. \\ \quad \operatorname{is\text{-}normed\text{-}field}(f) \wedge \\ \quad \operatorname{is\text{-}bilinear\text{-}bicharacter}(f, k) \wedge \\ \quad 1 \leq n \wedge \\ \quad y \in \operatorname{vec}(\operatorname{ring\text{-}pow}(\operatorname{normed\text{-}field\text{-}as\text{-}commutative\text{-}ring}(f), n)) \wedge \\ \quad x \in \operatorname{vec}(\operatorname{ring\text{-}pow}(\operatorname{normed\text{-}field\text{-}as\text{-}commutative\text{-}ring}(f), n)) \Rightarrow \\ \quad \quad k(\operatorname{entry}(\operatorname{matmul}(\operatorname{normed\text{-}field\text{-}as\text{-}commutative\text{-}ring}(f), \\ \quad \quad \quad \quad \quad \operatorname{ring\text{-}row}(\operatorname{normed\text{-}field\text{-}as\text{-}commutative\text{-}ring}(f), n, y), \\ \quad \quad \quad \quad \quad \operatorname{matmul}(\operatorname{normed\text{-}field\text{-}as\text{-}commutative\text{-}ring}(f), a, \\ \quad \quad \quad \quad \quad \quad \operatorname{transpose}(\operatorname{ring\text{-}row}(\operatorname{normed\text{-}field\text{-}as\text{-}commutative\text{-}ring}(f), n, x), 1, n))), \\ \quad \quad \quad \quad 1, 1), \\ \quad \quad \quad \operatorname{one}(f)) \\ \quad \quad =\; \operatorname{finsum}(\operatorname{commutative\text{-}ring\text{-}multiplicative\text{-}cm}(\operatorname{normed\text{-}field\text{-}as\text{-}commutative\text{-}ring}(\operatorname{cc\text{-}normed\text{-}field})), \\ \quad \quad \quad (\lambda i \in \operatorname{interval}(1, n).\; \\ \quad \quad \quad \quad \operatorname{finsum}(\operatorname{commutative\text{-}ring\text{-}multiplicative\text{-}cm}(\operatorname{normed\text{-}field\text{-}as\text{-}commutative\text{-}ring}(\operatorname{cc\text{-}normed\text{-}field})), \\ \quad \quad \quad \quad \quad (\lambda j \in \operatorname{interval}(1, n).\; k(\operatorname{mul}(f)(\operatorname{entry}(a, i, j), y(i)), x(j))), \\ \quad \quad \quad \quad \quad \operatorname{interval}(1, n))), \\ \quad \quad \quad \operatorname{interval}(1, n)) \end{array}}}
\vnbprop{ccr-6-24-converse}{\vnbstmt{\begin{array}{@{}l@{}} \forall f, k.; \forall n \in \mathbb{N}, a \in \operatorname{mat}(n, n, \operatorname{carr}(f)). \\ \quad (\operatorname{is\text{-}normed\text{-}field}(f) \wedge (\operatorname{is\text{-}bilinear\text{-}bicharacter}(f, k) \wedge (1 \leq n))) \Rightarrow  \\ \quad \quad \exists w. \\ \quad \quad \quad \operatorname{is\text{-}e\text{-}bilinear\text{-}bicharacter}(f, \operatorname{ring\text{-}pow}(\operatorname{normed\text{-}field\text{-}as\text{-}commutative\text{-}ring}(f), n), w) \wedge \\ \quad \quad \quad \forall y, x \in \operatorname{vec}(\operatorname{ring\text{-}pow}(\operatorname{normed\text{-}field\text{-}as\text{-}commutative\text{-}ring}(f), n)). \\ \quad \quad \quad \quad w(y, x) \\ \quad \quad \quad \quad =\; \operatorname{finsum}(\operatorname{commutative\text{-}ring\text{-}multiplicative\text{-}cm}(\operatorname{normed\text{-}field\text{-}as\text{-}commutative\text{-}ring}(\operatorname{cc\text{-}normed\text{-}field})), \\ \quad \quad \quad \quad \quad (\lambda i \in \operatorname{interval}(1, n).\; \\ \quad \quad \quad \quad \quad \quad \operatorname{finsum}(\operatorname{commutative\text{-}ring\text{-}multiplicative\text{-}cm}(\operatorname{normed\text{-}field\text{-}as\text{-}commutative\text{-}ring}(\operatorname{cc\text{-}normed\text{-}field})), \\ \quad \quad \quad \quad \quad \quad \quad (\lambda j \in \operatorname{interval}(1, n).\; k(\operatorname{mul}(f)(\operatorname{entry}(a, i, j), y(i)), x(j))), \\ \quad \quad \quad \quad \quad \quad \quad \operatorname{interval}(1, n))), \\ \quad \quad \quad \quad \quad \operatorname{interval}(1, n)) \end{array}}}

\vnbbookfile{theorem-library/pruefer-laws}{theorem-library/pruefer-laws}
\noindent{\small pruefer-laws.scm -- WAVE 2 (2026-10-27..30), wired 2026-10-30 from scratchpad/ccrpr/pruefer-laws.scm. Library: complex-analysis.}\par
\vnbprop{nfag-carr}{\vnbstmt{\begin{array}{@{}l@{}} \forall f. \\ \quad \operatorname{is\text{-}normed\text{-}field}(f) \Rightarrow (\operatorname{carr}(\operatorname{normed\text{-}field\text{-}additive\text{-}ag}(f)) \simeq \operatorname{carr}(f)) \end{array}}}
\vnbprop{nfag-opr}{\vnbstmt{\begin{array}{@{}l@{}} \forall f. \\ \quad \operatorname{is\text{-}normed\text{-}field}(f) \Rightarrow (\operatorname{opr}(\operatorname{normed\text{-}field\text{-}additive\text{-}ag}(f)) \simeq \operatorname{add}(f)) \end{array}}}
\vnbprop{nfag-iden}{\vnbstmt{\begin{array}{@{}l@{}} \forall f. \\ \quad \operatorname{is\text{-}normed\text{-}field}(f) \Rightarrow (\operatorname{iden}(\operatorname{normed\text{-}field\text{-}additive\text{-}ag}(f)) \simeq \operatorname{zero}(f)) \end{array}}}
\vnbprop{nfag-inv}{\vnbstmt{\begin{array}{@{}l@{}} \forall f. \\ \quad \operatorname{is\text{-}normed\text{-}field}(f) \Rightarrow (\operatorname{inv}(\operatorname{normed\text{-}field\text{-}additive\text{-}ag}(f)) \simeq \operatorname{neg}(f)) \end{array}}}
\vnbprop{nv-integers-subgroup}{\vnbstmt{\begin{array}{@{}l@{}} \forall f. \\ \quad \operatorname{is\text{-}nv\text{-}field}(f) \Rightarrow \operatorname{is\text{-}subgroup}(\operatorname{normed\text{-}field\text{-}additive\text{-}ag}(f), \operatorname{nv\text{-}integers}(f)) \end{array}}}
\vnbprop{nv-separating-factor}{\vnbstmt{\begin{array}{@{}l@{}} \forall f, l. \\ \quad (\operatorname{is\text{-}nv\text{-}field}(f) \wedge \operatorname{is\text{-}separating\text{-}character}(f, l)) \Rightarrow  \\ \quad \quad \exists b \in (\operatorname{carr}(\operatorname{quotient\text{-}ag}(\operatorname{normed\text{-}field\text{-}additive\text{-}ag}(f), \operatorname{nv\text{-}integers}(f))) \to \operatorname{unit\text{-}circle}). \\ \quad \quad \quad \forall x \in \operatorname{carr}(f). \\ \quad \quad \quad \quad b(\operatorname{class}(\operatorname{qmod\text{-}setoid}(\operatorname{ag\text{-}zz\text{-}module}(\operatorname{normed\text{-}field\text{-}additive\text{-}ag}(f)), \operatorname{nv\text{-}integers}(f)), x)) = l(x) \wedge \\ \quad \quad \quad \forall c, d \in \operatorname{carr}(\operatorname{quotient\text{-}ag}(\operatorname{normed\text{-}field\text{-}additive\text{-}ag}(f), \operatorname{nv\text{-}integers}(f))). \\ \quad \quad \quad \quad b(\operatorname{opr}(\operatorname{quotient\text{-}ag}(\operatorname{normed\text{-}field\text{-}additive\text{-}ag}(f), \operatorname{nv\text{-}integers}(f)))(c, d)) = (b(c) \cdot b(d)) \wedge \\ \quad \quad \quad \forall c, d \in \operatorname{carr}(\operatorname{quotient\text{-}ag}(\operatorname{normed\text{-}field\text{-}additive\text{-}ag}(f), \operatorname{nv\text{-}integers}(f))). \\ \quad \quad \quad \quad (b(c) = b(d)) \Rightarrow (c = d) \end{array}}}
\vnbprop{ccr-5-22}{\vnbstmt{\forall p.\; (\operatorname{zz\text{-}prime}(p) \Rightarrow \operatorname{is\text{-}abelian\text{-}group}(\operatorname{pruefer}(p)))}}
\vnbprop{qp-lambda-separating}{\vnbstmt{\begin{array}{@{}l@{}} \forall p. \\ \quad \operatorname{zz\text{-}prime}(p) \Rightarrow \operatorname{is\text{-}separating\text{-}character}(\operatorname{qp\text{-}field}(p), \operatorname{qp\text{-}lambda}(p)) \end{array}}}
\vnbprop{ccr-5-23}{\vnbstmt{\begin{array}{@{}l@{}} \forall p, l. \\ \quad (\operatorname{zz\text{-}prime}(p) \wedge \operatorname{is\text{-}separating\text{-}character}(\operatorname{qp\text{-}field}(p), l)) \Rightarrow  \\ \quad \quad \exists b \in (\operatorname{carr}(\operatorname{pruefer}(p)) \to \operatorname{unit\text{-}circle}). \\ \quad \quad \quad \forall x \in \operatorname{carr}(\operatorname{qp\text{-}field}(p)). \\ \quad \quad \quad \quad b(\operatorname{class}(\operatorname{qmod\text{-}setoid}(\operatorname{ag\text{-}zz\text{-}module}(\operatorname{normed\text{-}field\text{-}additive\text{-}ag}(\operatorname{qp\text{-}field}(p))), \operatorname{zp}(p)), x)) = l(x) \wedge \\ \quad \quad \quad \forall c, d \in \operatorname{carr}(\operatorname{pruefer}(p)). \\ \quad \quad \quad \quad b(\operatorname{opr}(\operatorname{pruefer}(p))(c, d)) = (b(c) \cdot b(d)) \wedge \\ \quad \quad \quad \forall c, d \in \operatorname{carr}(\operatorname{pruefer}(p)). \\ \quad \quad \quad \quad (b(c) = b(d)) \Rightarrow (c = d) \end{array}}}
\vnbprop{qp-gpow-embed}{\vnbstmt{\begin{array}{@{}l@{}} \forall n \in \mathbb{N}, p.; \forall x \in \operatorname{carr}(\operatorname{qp\text{-}field}(p)). \\ \quad \operatorname{zz\text{-}prime}(p) \Rightarrow  \\ \quad \quad \operatorname{gpow}(\operatorname{normed\text{-}field\text{-}additive\text{-}ag}(\operatorname{qp\text{-}field}(p)), x, n) \\ \quad \quad =\; \operatorname{mul}(\operatorname{qp\text{-}field}(p))(\operatorname{embed}(\operatorname{nf\text{-}metric\text{-}space}(\operatorname{qq\text{-}padic}(p)))(n), x) \end{array}}}
\vnbprop{qp-ppow-nat}{\vnbstmt{\begin{array}{@{}l@{}} \forall p.; \forall n \in \mathbb{N}. \\ \quad \operatorname{zz\text{-}prime}(p) \Rightarrow (\operatorname{qp\text{-}ppow}(p, n) = \operatorname{embed}(\operatorname{nf\text{-}metric\text{-}space}(\operatorname{qq\text{-}padic}(p)))(\operatorname{power}(p, n))) \end{array}}}
\vnbprop{qp-pruefer-order}{\vnbstmt{\begin{array}{@{}l@{}} \forall p.; \forall x \in \operatorname{carr}(\operatorname{qp\text{-}field}(p)), n \in \mathbb{N}. \\ \quad (\operatorname{zz\text{-}prime}(p) \wedge (\neg (x = \operatorname{zero}(\operatorname{qp\text{-}field}(p))) \wedge (-\operatorname{qp\text{-}ord}(p)(x) = n))) \Rightarrow  \\ \quad \quad \operatorname{is\text{-}element\text{-}order}(\operatorname{pruefer}(p), \\ \quad \quad \quad \operatorname{class}(\operatorname{qmod\text{-}setoid}(\operatorname{ag\text{-}zz\text{-}module}(\operatorname{normed\text{-}field\text{-}additive\text{-}ag}(\operatorname{qp\text{-}field}(p))), \operatorname{zp}(p)), x), \\ \quad \quad \quad \operatorname{power}(p, n)) \end{array}}}
\vnbprop{ccr-1-33}{\vnbstmt{\begin{array}{@{}l@{}} \forall p.; \forall x \in \operatorname{carr}(\operatorname{qp\text{-}field}(p)). \\ \quad (\operatorname{zz\text{-}prime}(p) \wedge (\neg (x = \operatorname{zero}(\operatorname{qp\text{-}field}(p))) \wedge (0 \leq -\operatorname{qp\text{-}ord}(p)(x)))) \Rightarrow  \\ \quad \quad -\operatorname{qp\text{-}ord}(p)(x) \in \mathbb{N} \wedge \\ \quad \quad \operatorname{is\text{-}element\text{-}order}(\operatorname{pruefer}(p), \\ \quad \quad \quad \operatorname{class}(\operatorname{qmod\text{-}setoid}(\operatorname{ag\text{-}zz\text{-}module}(\operatorname{normed\text{-}field\text{-}additive\text{-}ag}(\operatorname{qp\text{-}field}(p))), \operatorname{zp}(p)), x), \\ \quad \quad \quad \operatorname{power}(p, -\operatorname{qp\text{-}ord}(p)(x))) \end{array}}}
\vnbprop{ccr-5-24}{\vnbstmt{\begin{array}{@{}l@{}} \forall p.; \forall x \in \operatorname{carr}(\operatorname{qp\text{-}field}(p)). \\ \quad \operatorname{zz\text{-}prime}(p) \Rightarrow  \\ \quad \quad (x \in \operatorname{zp}(p)) \Rightarrow  \\ \quad \quad \quad \operatorname{is\text{-}element\text{-}order}(\operatorname{pruefer}(p), \\ \quad \quad \quad \quad \operatorname{class}(\operatorname{qmod\text{-}setoid}(\operatorname{ag\text{-}zz\text{-}module}(\operatorname{normed\text{-}field\text{-}additive\text{-}ag}(\operatorname{qp\text{-}field}(p))), \operatorname{zp}(p)), x), \\ \quad \quad \quad \quad 1) \wedge \\ \quad \quad \neg (x \in \operatorname{zp}(p)) \Rightarrow  \\ \quad \quad \quad -\operatorname{qp\text{-}ord}(p)(x) \in \mathbb{N} \wedge \\ \quad \quad \quad \operatorname{is\text{-}element\text{-}order}(\operatorname{pruefer}(p), \\ \quad \quad \quad \quad \operatorname{class}(\operatorname{qmod\text{-}setoid}(\operatorname{ag\text{-}zz\text{-}module}(\operatorname{normed\text{-}field\text{-}additive\text{-}ag}(\operatorname{qp\text{-}field}(p))), \operatorname{zp}(p)), x), \\ \quad \quad \quad \quad \operatorname{power}(p, -\operatorname{qp\text{-}ord}(p)(x))) \end{array}}}
\vnbprop{hom-cc-gpow}{\vnbstmt{\begin{array}{@{}l@{}} \forall n \in \mathbb{N}, q, b.; \forall c \in \operatorname{carr}(q). \\ \quad \operatorname{is\text{-}group}(q) \wedge \\ \quad \forall c \in \operatorname{carr}(q).\; (b(c) \in \mathbb{C}) \wedge \\ \quad b(\operatorname{iden}(q)) = 1 \wedge \\ \quad \forall c, d \in \operatorname{carr}(q).\; (b(\operatorname{opr}(q)(c, d)) = (b(c) \cdot b(d))) \Rightarrow \\ \quad \quad b(\operatorname{gpow}(q, c, n)) = \operatorname{power}(b(c), n) \end{array}}}
\vnbprop{qp-lambda-onto-roots}{\vnbstmt{\begin{array}{@{}l@{}} \forall p.; \forall w \in \operatorname{mu\text{-}pinf}(p). \\ \quad \operatorname{zz\text{-}prime}(p) \Rightarrow  \\ \quad \quad \exists r \in \mathbb{Q}.\; (\operatorname{qp\text{-}lambda}(p)(\operatorname{embed}(\operatorname{nf\text{-}metric\text{-}space}(\operatorname{qq\text{-}padic}(p)))(r)) = w) \end{array}}}
\vnbprop{ccr-5-25}{\vnbstmt{\begin{array}{@{}l@{}} \forall p. \\ \quad \operatorname{zz\text{-}prime}(p) \Rightarrow  \\ \quad \quad \exists f \in \mathrm{Bij}(\operatorname{carr}(\operatorname{pruefer}(p)), \operatorname{mu\text{-}pinf}(p)). \forall c, d \in \operatorname{carr}(\operatorname{pruefer}(p)). \\ \quad \quad \quad f(\operatorname{opr}(\operatorname{pruefer}(p))(c, d)) = (f(c) \cdot f(d)) \end{array}}}
\vnbprop{root-power-iff}{\vnbstmt{\begin{array}{@{}l@{}} \forall n \in \mathbb{N}, j \in \mathbb{Z}, m \in \mathbb{N}. \\ \quad (1 \leq n) \Rightarrow  \\ \quad \quad (\operatorname{power}(\exp (((2 \cdot \pi) \cdot +i) \cdot (j \cdot \operatorname{recip}(n))), m) = 1) \Leftrightarrow \operatorname{zz\text{-}divides}(n, (j \cdot m)) \end{array}}}
\vnbprop{coprime-root-primitive}{\vnbstmt{\begin{array}{@{}l@{}} \forall n \in \mathbb{N}, j \in \mathbb{Z}. \\ \quad ((1 \leq n) \wedge \operatorname{zz\text{-}coprime}(j, n)) \Rightarrow  \\ \quad \quad \exp (((2 \cdot \pi) \cdot +i) \cdot (j \cdot \operatorname{recip}(n))) \in \operatorname{primitive\text{-}roots}(n) \end{array}}}
\vnbprop{primitive-root-coprime}{\vnbstmt{\begin{array}{@{}l@{}} \forall n \in \mathbb{N}, j \in \mathbb{Z}. \\ \quad ((1 \leq n) \wedge (\exp (((2 \cdot \pi) \cdot +i) \cdot (j \cdot \operatorname{recip}(n))) \in \operatorname{primitive\text{-}roots}(n))) \Rightarrow  \\ \quad \quad \operatorname{zz\text{-}coprime}(j, n) \end{array}}}
\vnbprop{root-of-unity-rep}{\vnbstmt{\begin{array}{@{}l@{}} \forall n \in \mathbb{N}, z \in \mathbb{C}. \\ \quad ((1 \leq n) \wedge (\operatorname{power}(z, n) = 1)) \Rightarrow  \\ \quad \quad \exists r \in \operatorname{interval}(1, n).\; (z = \exp (((2 \cdot \pi) \cdot +i) \cdot (r \cdot \operatorname{recip}(n)))) \end{array}}}
\vnbprop{root-rep-injective}{\vnbstmt{\begin{array}{@{}l@{}} \forall n \in \mathbb{N}, a, b \in \operatorname{interval}(1, n). \\ \quad ((1 \leq n) \wedge (\exp (((2 \cdot \pi) \cdot +i) \cdot (a \cdot \operatorname{recip}(n))) = \exp (((2 \cdot \pi) \cdot +i) \cdot (b \cdot \operatorname{recip}(n))))) \Rightarrow  \\ \quad \quad a = b \end{array}}}
\vnbprop{euler-primitive-roots}{\vnbstmt{\forall n \in \mathbb{N}.\; ((1 \leq n) \Rightarrow (|\operatorname{primitive\text{-}roots}(n)| = \operatorname{euler\text{-}phi}(n)))}}
\vnbprop{ccr-5-26}{\vnbstmt{\begin{array}{@{}l@{}} \forall k \in \mathbb{N}, p. \\ \quad \operatorname{is\text{-}prime\text{-}nat}(p) \Rightarrow  \\ \quad \quad |\operatorname{primitive\text{-}roots}(\operatorname{power}(p, (k + 1)))| = (\operatorname{power}(p, (k + 1)) + -\operatorname{power}(p, k)) \end{array}}}

\vnbbookfile{theorem-library/ca-lft-group}{theorem-library/ca-lft-group}
\noindent{\small ca-lft-group.scm -- Cor. 1.19 of complex-analysis.pdf: the linear fractional transformations form a group under composition. CA-LFT, 2026-10-19.  Library: complex-analysis.  Vocabulary: structure-library/ca-lft-vocab.scm.}\par
\vnbprop{lfts-intro}{\vnbstmt{\begin{array}{@{}l@{}} \forall a, b, c, d \in \mathbb{C}. \\ \quad \neg (((a \cdot d) - (b \cdot c)) = 0) \Rightarrow (\operatorname{lft}(a, b, c, d) \in \mathit{lfts}) \end{array}}}
\vnbprop{lfts-in-fun}{\vnbstmt{\forall f \in \mathit{lfts}.\; (f \in (\operatorname{cc\text{-}hat} \to \operatorname{cc\text{-}hat}))}}
\vnbprop{lfts-compose-closed}{\vnbstmt{\forall f, g \in \mathit{lfts}.\; (\operatorname{compose}(f, g) \in \mathit{lfts})}}
\vnbprop{id-in-lfts}{\vnbstmt{\operatorname{id\text{-}fun}(\operatorname{cc\text{-}hat}) \in \mathit{lfts}}}
\vnbprop{lft-adjugate-left}{\vnbstmt{\begin{array}{@{}l@{}} \forall a, b, c, d \in \mathbb{C}. \\ \quad \neg (((a \cdot d) - (b \cdot c)) = 0) \Rightarrow  \\ \quad \quad \operatorname{compose}(\operatorname{lft}(d, -b, -c, a), \operatorname{lft}(a, b, c, d)) = \operatorname{id\text{-}fun}(\operatorname{cc\text{-}hat}) \end{array}}}
\vnbprop{lft-adjugate-right}{\vnbstmt{\begin{array}{@{}l@{}} \forall a, b, c, d \in \mathbb{C}. \\ \quad \neg (((a \cdot d) - (b \cdot c)) = 0) \Rightarrow  \\ \quad \quad \operatorname{compose}(\operatorname{lft}(a, b, c, d), \operatorname{lft}(d, -b, -c, a)) = \operatorname{id\text{-}fun}(\operatorname{cc\text{-}hat}) \end{array}}}
\vnbprop{lft-adjugate-apply-left}{\vnbstmt{\begin{array}{@{}l@{}} \forall a, b, c, d \in \mathbb{C}, u \in \operatorname{cc\text{-}hat}. \\ \quad \neg (((a \cdot d) - (b \cdot c)) = 0) \Rightarrow (\operatorname{lft}(d, -b, -c, a)(\operatorname{lft}(a, b, c, d)(u)) = u) \end{array}}}
\vnbprop{lft-adjugate-apply-right}{\vnbstmt{\begin{array}{@{}l@{}} \forall a, b, c, d \in \mathbb{C}, u \in \operatorname{cc\text{-}hat}. \\ \quad \neg (((a \cdot d) - (b \cdot c)) = 0) \Rightarrow (\operatorname{lft}(a, b, c, d)(\operatorname{lft}(d, -b, -c, a)(u)) = u) \end{array}}}
\vnbprop{lft-bijection}{\vnbstmt{\begin{array}{@{}l@{}} \forall a, b, c, d \in \mathbb{C}. \\ \quad \neg (((a \cdot d) - (b \cdot c)) = 0) \Rightarrow (\operatorname{lft}(a, b, c, d) \in \mathrm{Bij}(\operatorname{cc\text{-}hat}, \operatorname{cc\text{-}hat})) \end{array}}}
\vnbprop{lft-inverse-bij}{\vnbstmt{\begin{array}{@{}l@{}} \forall a, b, c, d \in \mathbb{C}. \\ \quad \neg (((a \cdot d) - (b \cdot c)) = 0) \Rightarrow  \\ \quad \quad \operatorname{lft}(a, b, c, d)^{-1} = \operatorname{lft}(d, -b, -c, a) \end{array}}}
\vnbprop{lfts-in-set}{\vnbstmt{\mathit{lfts} \in \mathbf{Set}}}
\vnbprop{lfts-inverse-in}{\vnbstmt{\forall f \in \mathit{lfts}.\; (f^{-1} \in \mathit{lfts})}}
\vnbprop{lft-group-opr-type}{\vnbstmt{\begin{array}{@{}l@{}} (\lambda {[}\operatorname{lfg\_}, \operatorname{lfh\_}] \in (\mathit{lfts} \times \mathit{lfts}).\; \operatorname{compose}(\operatorname{lfg\_}, \operatorname{lfh\_})) \\ \in\; ((\mathit{lfts} \times \mathit{lfts}) \to \mathit{lfts}) \end{array}}}
\vnbprop{lft-group-inv-type}{\vnbstmt{(\lambda g \in \mathit{lfts}.\; g^{-1}) \in (\mathit{lfts} \to \mathit{lfts})}}
\vnbprop{lft-group-is-group}{\vnbstmt{\operatorname{is\text{-}group}(\operatorname{lft\text{-}group})}}

\vnbbookfile{theorem-library/ca-lft-matrix}{theorem-library/ca-lft-matrix}
\noindent{\small ca-lft-matrix.scm -- Prop. 1.18 of complex-analysis.pdf in matrix form: S |-$>$ L\_S from the invertible 2-by-2 complex matrices onto the linear fractional transformations is a surjective homomorphism whose kernel is the scalar multiples of the identity. CA-LFT, 2026-10-19.  Library: complex-analysis.  Vocabulary: ...}\par
\vnbprop{cc-matmul-2x2-entry}{\vnbstmt{\begin{array}{@{}l@{}} \forall p, q \in \operatorname{mat}(2, 2, \mathbb{C}), i, k \in \operatorname{interval}(1, 2). \\ \quad \operatorname{entry}(\operatorname{matmul}(\operatorname{normed\text{-}field\text{-}as\text{-}commutative\text{-}ring}(\operatorname{cc\text{-}normed\text{-}field}), p, q), i, k) \\ \quad =\; (\operatorname{entry}(p, i, 1) \cdot \operatorname{entry}(q, 1, k)) + (\operatorname{entry}(p, i, 2) \cdot \operatorname{entry}(q, 2, k)) \end{array}}}
\vnbprop{cc-det-2x2}{\vnbstmt{\begin{array}{@{}l@{}} \forall s \in \operatorname{mat}(2, 2, \mathbb{C}). \\ \quad \operatorname{det}(\operatorname{normed\text{-}field\text{-}as\text{-}commutative\text{-}ring}(\operatorname{cc\text{-}normed\text{-}field}), 2, s) \\ \quad =\; ((\operatorname{entry}(s, 1, 1) \cdot \operatorname{entry}(s, 2, 2)) - (\operatorname{entry}(s, 1, 2) \cdot \operatorname{entry}(s, 2, 1))) \end{array}}}
\vnbprop{cc-mat22-entries}{\vnbstmt{\begin{array}{@{}l@{}} \forall a, b, c, d \in \mathbb{C}. \\ \quad \operatorname{entry}({[}{[}a, b], {[}c, d]], 1, 1) = a \wedge \\ \quad \operatorname{entry}({[}{[}a, b], {[}c, d]], 1, 2) = b \wedge \\ \quad \operatorname{entry}({[}{[}a, b], {[}c, d]], 2, 1) = c \wedge \\ \quad \operatorname{entry}({[}{[}a, b], {[}c, d]], 2, 2) = d \end{array}}}
\vnbprop{cc-invertible-mat-is-mat}{\vnbstmt{\begin{array}{@{}l@{}} \forall s. \\ \quad \operatorname{is\text{-}invertible\text{-}mat}(\operatorname{normed\text{-}field\text{-}as\text{-}commutative\text{-}ring}(\operatorname{cc\text{-}normed\text{-}field}), 2, s) \Rightarrow  \\ \quad \quad s \in \operatorname{mat}(2, 2, \mathbb{C}) \end{array}}}
\vnbprop{cc-invertible-det-nonzero}{\vnbstmt{\begin{array}{@{}l@{}} \forall s. \\ \quad \operatorname{is\text{-}invertible\text{-}mat}(\operatorname{normed\text{-}field\text{-}as\text{-}commutative\text{-}ring}(\operatorname{cc\text{-}normed\text{-}field}), 2, s) \Rightarrow  \\ \quad \quad \neg (((\operatorname{entry}(s, 1, 1) \cdot \operatorname{entry}(s, 2, 2)) - (\operatorname{entry}(s, 1, 2) \cdot \operatorname{entry}(s, 2, 1))) = 0) \end{array}}}
\vnbprop{cc-det-nonzero-invertible}{\vnbstmt{\begin{array}{@{}l@{}} \forall s \in \operatorname{mat}(2, 2, \mathbb{C}). \\ \quad \neg (((\operatorname{entry}(s, 1, 1) \cdot \operatorname{entry}(s, 2, 2)) - (\operatorname{entry}(s, 1, 2) \cdot \operatorname{entry}(s, 2, 1))) = 0) \Rightarrow  \\ \quad \quad \operatorname{is\text{-}invertible\text{-}mat}(\operatorname{normed\text{-}field\text{-}as\text{-}commutative\text{-}ring}(\operatorname{cc\text{-}normed\text{-}field}), 2, s) \end{array}}}
\vnbprop{mat-lft-literal}{\vnbstmt{\begin{array}{@{}l@{}} \forall a, b, c, d \in \mathbb{C}. \\ \quad \operatorname{mat\text{-}lft}({[}{[}a, b], {[}c, d]]) = \operatorname{lft}(a, b, c, d) \end{array}}}
\vnbprop{mat-lft-in-lfts}{\vnbstmt{\begin{array}{@{}l@{}} \forall s. \\ \quad \operatorname{is\text{-}invertible\text{-}mat}(\operatorname{normed\text{-}field\text{-}as\text{-}commutative\text{-}ring}(\operatorname{cc\text{-}normed\text{-}field}), 2, s) \Rightarrow  \\ \quad \quad \operatorname{mat\text{-}lft}(s) \in \mathit{lfts} \end{array}}}
\vnbprop{mat-lft-hom}{\vnbstmt{\begin{array}{@{}l@{}} \forall s, t. \\ \quad \operatorname{is\text{-}invertible\text{-}mat}(\operatorname{normed\text{-}field\text{-}as\text{-}commutative\text{-}ring}(\operatorname{cc\text{-}normed\text{-}field}), 2, s) \wedge \\ \quad \operatorname{is\text{-}invertible\text{-}mat}(\operatorname{normed\text{-}field\text{-}as\text{-}commutative\text{-}ring}(\operatorname{cc\text{-}normed\text{-}field}), 2, t) \Rightarrow \\ \quad \quad \operatorname{mat\text{-}lft}(\operatorname{matmul}(\operatorname{normed\text{-}field\text{-}as\text{-}commutative\text{-}ring}(\operatorname{cc\text{-}normed\text{-}field}), s, t)) \\ \quad \quad =\; \operatorname{compose}(\operatorname{mat\text{-}lft}(s), \operatorname{mat\text{-}lft}(t)) \end{array}}}
\vnbprop{mat-lft-onto}{\vnbstmt{\begin{array}{@{}l@{}} \forall f \in \mathit{lfts}. \exists s. \\ \quad (\operatorname{is\text{-}invertible\text{-}mat}(\operatorname{normed\text{-}field\text{-}as\text{-}commutative\text{-}ring}(\operatorname{cc\text{-}normed\text{-}field}), 2, s) \wedge (f = \operatorname{mat\text{-}lft}(s))) \end{array}}}
\vnbprop{mat-lft-kernel}{\vnbstmt{\begin{array}{@{}l@{}} \forall s. \\ \quad \operatorname{is\text{-}invertible\text{-}mat}(\operatorname{normed\text{-}field\text{-}as\text{-}commutative\text{-}ring}(\operatorname{cc\text{-}normed\text{-}field}), 2, s) \Rightarrow  \\ \quad \quad (\operatorname{mat\text{-}lft}(s) = \operatorname{id\text{-}fun}(\operatorname{cc\text{-}hat})) \Leftrightarrow \exists r \in \mathbb{C}.\; (s = {[}{[}r, 0], {[}0, r]]) \end{array}}}

\vnbbookfile{theorem-library/ca-lft-sl2}{theorem-library/ca-lft-sl2}
\noindent{\small ca-lft-sl2.scm -- Cor. 1.20 of complex-analysis.pdf: S |-$>$ L\_S restricted to the matrices of determinant 1 is a surjective homomorphism onto the linear fractional transformations, with kernel \{I, -I\}. CA-LFT, 2026-10-19.  Library: complex-analysis.  Vocabulary: structure-library/ca-lft-vocab.scm.}\par
\vnbprop{cc-sqrt-exists}{\vnbstmt{\forall d \in \mathbb{C}.\; (\neg (d = 0) \Rightarrow \exists r \in \mathbb{C}.\; ((r \cdot r) = d))}}
\vnbprop{homog-point-scale}{\vnbstmt{\begin{array}{@{}l@{}} \forall u, v, k \in \mathbb{C}. \\ \quad \neg (k = 0) \Rightarrow (\operatorname{homog\text{-}point}((k \cdot u), (k \cdot v)) = \operatorname{homog\text{-}point}(u, v)) \end{array}}}
\vnbprop{lft-scale-invariant}{\vnbstmt{\begin{array}{@{}l@{}} \forall a, b, c, d, k \in \mathbb{C}. \\ \quad (\neg (((a \cdot d) - (b \cdot c)) = 0) \wedge \neg (k = 0)) \Rightarrow  \\ \quad \quad \operatorname{lft}((k \cdot a), (k \cdot b), (k \cdot c), (k \cdot d)) = \operatorname{lft}(a, b, c, d) \end{array}}}
\vnbprop{sl2-matmul-closed}{\vnbstmt{\begin{array}{@{}l@{}} \forall s, t \in \operatorname{mat}(2, 2, \mathbb{C}). \\ \quad \operatorname{det}(\operatorname{normed\text{-}field\text{-}as\text{-}commutative\text{-}ring}(\operatorname{cc\text{-}normed\text{-}field}), 2, s) = 1 \wedge \\ \quad \operatorname{det}(\operatorname{normed\text{-}field\text{-}as\text{-}commutative\text{-}ring}(\operatorname{cc\text{-}normed\text{-}field}), 2, t) = 1 \Rightarrow \\ \quad \quad \operatorname{det}(\operatorname{normed\text{-}field\text{-}as\text{-}commutative\text{-}ring}(\operatorname{cc\text{-}normed\text{-}field}), \\ \quad \quad \quad 2, \\ \quad \quad \quad \operatorname{matmul}(\operatorname{normed\text{-}field\text{-}as\text{-}commutative\text{-}ring}(\operatorname{cc\text{-}normed\text{-}field}), s, t)) \\ \quad \quad =\; 1 \end{array}}}
\vnbprop{sl2-lft-hom}{\vnbstmt{\begin{array}{@{}l@{}} \forall s, t \in \operatorname{mat}(2, 2, \mathbb{C}). \\ \quad \operatorname{det}(\operatorname{normed\text{-}field\text{-}as\text{-}commutative\text{-}ring}(\operatorname{cc\text{-}normed\text{-}field}), 2, s) = 1 \wedge \\ \quad \operatorname{det}(\operatorname{normed\text{-}field\text{-}as\text{-}commutative\text{-}ring}(\operatorname{cc\text{-}normed\text{-}field}), 2, t) = 1 \Rightarrow \\ \quad \quad \operatorname{mat\text{-}lft}(\operatorname{matmul}(\operatorname{normed\text{-}field\text{-}as\text{-}commutative\text{-}ring}(\operatorname{cc\text{-}normed\text{-}field}), s, t)) \\ \quad \quad =\; \operatorname{compose}(\operatorname{mat\text{-}lft}(s), \operatorname{mat\text{-}lft}(t)) \end{array}}}
\vnbprop{sl2-lft-onto}{\vnbstmt{\begin{array}{@{}l@{}} \forall f \in \mathit{lfts}. \exists s \in \operatorname{mat}(2, 2, \mathbb{C}). \\ \quad ((\operatorname{det}(\operatorname{normed\text{-}field\text{-}as\text{-}commutative\text{-}ring}(\operatorname{cc\text{-}normed\text{-}field}), 2, s) = 1) \wedge (f = \operatorname{mat\text{-}lft}(s))) \end{array}}}
\vnbprop{sl2-lft-kernel}{\vnbstmt{\begin{array}{@{}l@{}} \forall s \in \operatorname{mat}(2, 2, \mathbb{C}). \\ \quad (\operatorname{det}(\operatorname{normed\text{-}field\text{-}as\text{-}commutative\text{-}ring}(\operatorname{cc\text{-}normed\text{-}field}), 2, s) = 1) \Rightarrow  \\ \quad \quad (\operatorname{mat\text{-}lft}(s) = \operatorname{id\text{-}fun}(\operatorname{cc\text{-}hat})) \Leftrightarrow ((s = {[}{[}1, 0], {[}0, 1]]) \vee (s = {[}{[}-1, 0], {[}0, -1]])) \end{array}}}

\vnbbookfile{structure-library/ca-hopf-vocab}{structure-library/ca-hopf-vocab}
\noindent{\small ca-hopf-vocab.scm -- THE HOPF FIBRATION, complex-analysis.pdf s. 1.6.3 ((13)-(16), Prop. 1.22, Def. 1.23, Cor. 1.24).  DEFINITIONS ONLY.  CA-HOPF, 2026-10-25.  Library: complex-analysis. Destination: structure-library/ca-hopf-vocab.scm.}\par
\par\medskip\noindent\textbf{Definition} (\texttt{cc2-sphere}).\;
\begin{equation*}
(\operatorname{cc2\text{-}sphere} \simeq \{\operatorname{hsw\_} \in (\mathbb{C} \times \mathbb{C}) : (+(\cdot(\operatorname{magnitude}(\operatorname{nth}(1, \operatorname{hsw\_})), \operatorname{magnitude}(\operatorname{nth}(1, \operatorname{hsw\_}))), \cdot(\operatorname{magnitude}(\operatorname{nth}(2, \operatorname{hsw\_})), \operatorname{magnitude}(\operatorname{nth}(2, \operatorname{hsw\_})))) = 1)\})
\end{equation*}

\par\medskip\noindent\textbf{Definition} (\texttt{unit-sphere-cr}).\;
\begin{equation*}
(\operatorname{unit\text{-}sphere\text{-}cr} \simeq \{\operatorname{hsp\_} \in (\mathbb{C} \times \mathbb{R}) : (+(\cdot(\operatorname{magnitude}(\operatorname{nth}(1, \operatorname{hsp\_})), \operatorname{magnitude}(\operatorname{nth}(1, \operatorname{hsp\_}))), \cdot(\operatorname{nth}(2, \operatorname{hsp\_}), \operatorname{nth}(2, \operatorname{hsp\_}))) = 1)\})
\end{equation*}

\par\medskip\noindent\textbf{Definition} (\texttt{hopf-map}).\;$\operatorname{hopf\text{-}map}(w) :=$
\begin{equation*}
{[}(2 \cdot (\operatorname{nth}(1, w) \cdot \operatorname{conjugate}(\operatorname{nth}(2, w)))), ((\operatorname{magnitude}(\operatorname{nth}(1, w)) \cdot \operatorname{magnitude}(\operatorname{nth}(1, w))) - (\operatorname{magnitude}(\operatorname{nth}(2, w)) \cdot \operatorname{magnitude}(\operatorname{nth}(2, w))))]
\end{equation*}

\par\medskip\noindent\textbf{Definition} (\texttt{homog-equiv}).\;For all $u$ and $w$, 
\begin{equation*}
\begin{array}{@{}l@{}} \operatorname{homog\text{-}equiv}(u, w) \\ \Leftrightarrow\; \exists \operatorname{hez\_} \in \mathbb{C}. \\ \quad \neg (\operatorname{hez\_} = 0) \wedge (u = {[}\cdot(\operatorname{hez\_}, \operatorname{nth}(1, w)), \cdot(\operatorname{hez\_}, \operatorname{nth}(2, w))]) \end{array}
\end{equation*}

\par\medskip\noindent\textbf{Definition} (\texttt{hopf-quotient-map}).\;
\begin{equation*}
(\operatorname{hopf\text{-}quotient\text{-}map} \simeq (\lambda \operatorname{hqp\_} \in \operatorname{unit\text{-}sphere\text{-}cr}.\; \operatorname{homog\text{-}point}(\operatorname{nth}(1, \operatorname{hqp\_}), -(1, \operatorname{nth}(2, \operatorname{hqp\_})))))
\end{equation*}


\vnbbookfile{theorem-library/ca-hopf-fibration}{theorem-library/ca-hopf-fibration}
\noindent{\small ca-hopf-fibration.scm -- THE HOPF FIBRATION: complex-analysis.pdf s. 1.6.3, Prop. 1.22, Def. 1.23, Cor. 1.24.  CA-HOPF, 2026-10-25.  Library: complex-analysis. Destination: theorem-library/ca-hopf-fibration.scm.  Vocabulary: structure-library/ca-hopf-vocab.scm.}\par
\vnbprop{cc2-sphere-membership}{\vnbstmt{\begin{array}{@{}l@{}} \forall w. \\ \quad w \in \operatorname{cc2\text{-}sphere} \\ \quad \Leftrightarrow\; w \in (\mathbb{C} \times \mathbb{C}) \wedge \\ \quad ((\operatorname{magnitude}(\operatorname{nth}(1, w)) \cdot \operatorname{magnitude}(\operatorname{nth}(1, w))) + (\operatorname{magnitude}(\operatorname{nth}(2, w)) \cdot \operatorname{magnitude}(\operatorname{nth}(2, w)))) = 1 \end{array}}}
\vnbprop{unit-sphere-cr-membership}{\vnbstmt{\begin{array}{@{}l@{}} \forall p. \\ \quad p \in \operatorname{unit\text{-}sphere\text{-}cr} \\ \quad \Leftrightarrow\; ((p \in (\mathbb{C} \times \mathbb{R})) \wedge (((\operatorname{magnitude}(\operatorname{nth}(1, p)) \cdot \operatorname{magnitude}(\operatorname{nth}(1, p))) + (\operatorname{nth}(2, p) \cdot \operatorname{nth}(2, p))) = 1)) \end{array}}}
\vnbprop{hopf-map-pair}{\vnbstmt{\begin{array}{@{}l@{}} \forall a, b \in \mathbb{C}. \\ \quad \operatorname{hopf\text{-}map}({[}a, b]) \\ \quad =\; {[}(2 \cdot (a \cdot \operatorname{conjugate}(b))), ((\operatorname{magnitude}(a) \cdot \operatorname{magnitude}(a)) - (\operatorname{magnitude}(b) \cdot \operatorname{magnitude}(b)))] \end{array}}}
\vnbprop{cc2-sphere-pair}{\vnbstmt{\begin{array}{@{}l@{}} \forall a, b \in \mathbb{C}. \\ \quad ({[}a, b] \in \operatorname{cc2\text{-}sphere}) \Leftrightarrow (((\operatorname{magnitude}(a) \cdot \operatorname{magnitude}(a)) + (\operatorname{magnitude}(b) \cdot \operatorname{magnitude}(b))) = 1) \end{array}}}
\vnbprop{unit-sphere-cr-pair}{\vnbstmt{\begin{array}{@{}l@{}} \forall z \in \mathbb{C}, r \in \mathbb{R}. \\ \quad ({[}z, r] \in \operatorname{unit\text{-}sphere\text{-}cr}) \Leftrightarrow (((\operatorname{magnitude}(z) \cdot \operatorname{magnitude}(z)) + (r \cdot r)) = 1) \end{array}}}
\vnbprop{hopf-magnitude-first}{\vnbstmt{\begin{array}{@{}l@{}} \forall a, b \in \mathbb{C}. \\ \quad \operatorname{magnitude}((2 \cdot (a \cdot \operatorname{conjugate}(b)))) = (2 \cdot (\operatorname{magnitude}(a) \cdot \operatorname{magnitude}(b))) \end{array}}}
\vnbprop{hopf-norm-identity}{\vnbstmt{\begin{array}{@{}l@{}} \forall a, b \in \mathbb{C}. \\ \quad (\operatorname{magnitude}((2 \cdot (a \cdot \operatorname{conjugate}(b)))) \cdot \operatorname{magnitude}((2 \cdot (a \cdot \operatorname{conjugate}(b))))) + (((\operatorname{magnitude}(a) \cdot \operatorname{magnitude}(a)) - (\operatorname{magnitude}(b) \cdot \operatorname{magnitude}(b))) \cdot ((\operatorname{magnitude}(a) \cdot \operatorname{magnitude}(a)) - (\operatorname{magnitude}(b) \cdot \operatorname{magnitude}(b)))) \\ \quad =\; ((\operatorname{magnitude}(a) \cdot \operatorname{magnitude}(a)) + (\operatorname{magnitude}(b) \cdot \operatorname{magnitude}(b))) \cdot ((\operatorname{magnitude}(a) \cdot \operatorname{magnitude}(a)) + (\operatorname{magnitude}(b) \cdot \operatorname{magnitude}(b))) \end{array}}}
\vnbprop{hopf-pair-in-cr}{\vnbstmt{\begin{array}{@{}l@{}} \forall a, b \in \mathbb{C}. \\ \quad {[}(2 \cdot (a \cdot \operatorname{conjugate}(b))), ((\operatorname{magnitude}(a) \cdot \operatorname{magnitude}(a)) - (\operatorname{magnitude}(b) \cdot \operatorname{magnitude}(b)))] \\ \quad \in\; (\mathbb{C} \times \mathbb{R}) \end{array}}}
\vnbprop{hopf-map-in-cr}{\vnbstmt{\forall w \in (\mathbb{C} \times \mathbb{C}).\; (\operatorname{hopf\text{-}map}(w) \in (\mathbb{C} \times \mathbb{R}))}}
\vnbprop{hopf-maps-sphere}{\vnbstmt{\forall w \in \operatorname{cc2\text{-}sphere}.\; (\operatorname{hopf\text{-}map}(w) \in \operatorname{unit\text{-}sphere\text{-}cr})}}
\vnbprop{rr-magnitude-sq}{\vnbstmt{\forall x \in \mathbb{R}.\; ((\operatorname{magnitude}(x) \cdot \operatorname{magnitude}(x)) = (x \cdot x))}}
\vnbprop{hopf-onto-sphere}{\vnbstmt{\forall p \in \operatorname{unit\text{-}sphere\text{-}cr}.\; \exists w \in \operatorname{cc2\text{-}sphere}.\; (\operatorname{hopf\text{-}map}(w) = p)}}
\vnbprop{cc-conjugate-continuous-at}{\vnbstmt{\begin{array}{@{}l@{}} \forall y \in \mathbb{C}. \\ \quad \operatorname{is\text{-}continuous\text{-}at}(\operatorname{cc\text{-}ms}, \operatorname{cc\text{-}ms}, (\lambda w \in \mathbb{C}.\; \operatorname{conjugate}(w)), y) \end{array}}}
\vnbprop{hc-cont-conj-ptwise}{\vnbstmt{\begin{array}{@{}l@{}} \forall s, g, h, x. \\ \quad \operatorname{is\text{-}continuous\text{-}at}(s, \operatorname{cc\text{-}ms}, g, x) \wedge \\ \quad h \in (\operatorname{pts}(s) \to \mathbb{C}) \wedge \\ \quad \forall q \in \operatorname{pts}(s).\; (h(q) = \operatorname{conjugate}(g(q))) \Rightarrow \\ \quad \quad \operatorname{is\text{-}continuous\text{-}at}(s, \operatorname{cc\text{-}ms}, h, x) \end{array}}}
\vnbprop{hc-cont-mul-ptwise}{\vnbstmt{\begin{array}{@{}l@{}} \forall s, p, q, h, x. \\ \quad \operatorname{is\text{-}continuous\text{-}at}(s, \operatorname{cc\text{-}ms}, p, x) \wedge \\ \quad \operatorname{is\text{-}continuous\text{-}at}(s, \operatorname{cc\text{-}ms}, q, x) \wedge \\ \quad h \in (\operatorname{pts}(s) \to \mathbb{C}) \wedge \\ \quad \forall \mathit{hcq} \in \operatorname{pts}(s).\; (h(\mathit{hcq}) = (p(\mathit{hcq}) \cdot q(\mathit{hcq}))) \Rightarrow \\ \quad \quad \operatorname{is\text{-}continuous\text{-}at}(s, \operatorname{cc\text{-}ms}, h, x) \end{array}}}
\vnbprop{hc-cont-const-ptwise}{\vnbstmt{\begin{array}{@{}l@{}} \forall s, c, h, x. \\ \quad \operatorname{is\text{-}metric\text{-}space}(s) \wedge \\ \quad c \in \mathbb{C} \wedge \\ \quad x \in \operatorname{pts}(s) \wedge \\ \quad h \in (\operatorname{pts}(s) \to \mathbb{C}) \wedge \\ \quad \forall q \in \operatorname{pts}(s).\; (h(q) = c) \Rightarrow \\ \quad \quad \operatorname{is\text{-}continuous\text{-}at}(s, \operatorname{cc\text{-}ms}, h, x) \end{array}}}
\vnbprop{hc-cont-real-ptwise}{\vnbstmt{\begin{array}{@{}l@{}} \forall s, f, r, x. \\ \quad (\operatorname{is\text{-}continuous\text{-}at}(s, \operatorname{cc\text{-}ms}, f, x) \wedge ((r \in (\operatorname{pts}(s) \to \mathbb{R})) \wedge \forall q \in \operatorname{pts}(s).\; (f(q) = r(q)))) \Rightarrow  \\ \quad \quad \operatorname{is\text{-}continuous\text{-}at}(s, \operatorname{rr\text{-}ms}, r, x) \end{array}}}
\vnbprop{hr-cont-sub-ptwise}{\vnbstmt{\begin{array}{@{}l@{}} \forall s, f, g, h, x. \\ \quad \operatorname{is\text{-}continuous\text{-}at}(s, \operatorname{rr\text{-}ms}, f, x) \wedge \\ \quad \operatorname{is\text{-}continuous\text{-}at}(s, \operatorname{rr\text{-}ms}, g, x) \wedge \\ \quad h \in (\operatorname{pts}(s) \to \mathbb{R}) \wedge \\ \quad \forall q \in \operatorname{pts}(s).\; (h(q) = (f(q) - g(q))) \Rightarrow \\ \quad \quad \operatorname{is\text{-}continuous\text{-}at}(s, \operatorname{rr\text{-}ms}, h, x) \end{array}}}
\vnbprop{hopf-coords-continuous}{\vnbstmt{\begin{array}{@{}l@{}} \forall s, a, b, x. \\ \quad (\operatorname{is\text{-}continuous\text{-}at}(s, \operatorname{cc\text{-}ms}, a, x) \wedge \operatorname{is\text{-}continuous\text{-}at}(s, \operatorname{cc\text{-}ms}, b, x)) \Rightarrow  \\ \quad \quad \operatorname{is\text{-}continuous\text{-}at}(s, \operatorname{cc\text{-}ms}, (\lambda z \in \operatorname{pts}(s).\; (2 \cdot (a(z) \cdot \operatorname{conjugate}(b(z))))), x) \wedge \\ \quad \quad \operatorname{is\text{-}continuous\text{-}at}(s, \operatorname{rr\text{-}ms}, \\ \quad \quad \quad (\lambda z \in \operatorname{pts}(s).\; ((\operatorname{magnitude}(a(z)) \cdot \operatorname{magnitude}(a(z))) - (\operatorname{magnitude}(b(z)) \cdot \operatorname{magnitude}(b(z))))), \\ \quad \quad \quad x) \end{array}}}
\vnbprop{hopf-map-continuous}{\vnbstmt{\begin{array}{@{}l@{}} \forall s, a, b, x. \\ \quad (\operatorname{is\text{-}continuous\text{-}at}(s, \operatorname{cc\text{-}ms}, a, x) \wedge \operatorname{is\text{-}continuous\text{-}at}(s, \operatorname{cc\text{-}ms}, b, x)) \Rightarrow  \\ \quad \quad \operatorname{is\text{-}continuous\text{-}at}(s, \operatorname{cc\text{-}ms}, (\lambda z \in \operatorname{pts}(s).\; \operatorname{nth}(1, \operatorname{hopf\text{-}map}({[}a(z), b(z)]))), x) \wedge \\ \quad \quad \operatorname{is\text{-}continuous\text{-}at}(s, \operatorname{rr\text{-}ms}, (\lambda z \in \operatorname{pts}(s).\; \operatorname{nth}(2, \operatorname{hopf\text{-}map}({[}a(z), b(z)]))), x) \end{array}}}
\vnbprop{rr-nonneg-sq-one}{\vnbstmt{\forall x \in \mathbb{R}.\; (((0 \leq x) \wedge ((x \cdot x) = 1)) \Rightarrow (x = 1))}}
\vnbprop{hopf-map-unimodular-scale}{\vnbstmt{\begin{array}{@{}l@{}} \forall r, a, b \in \mathbb{C}. \\ \quad (\operatorname{magnitude}(r) = 1) \Rightarrow (\operatorname{hopf\text{-}map}({[}(r \cdot a), (r \cdot b)]) = \operatorname{hopf\text{-}map}({[}a, b])) \end{array}}}
\vnbprop{hopf-fibre-pair}{\vnbstmt{\begin{array}{@{}l@{}} \forall a, b, c, d \in \mathbb{C}. \\ \quad ((\operatorname{magnitude}(a) \cdot \operatorname{magnitude}(a)) + (\operatorname{magnitude}(b) \cdot \operatorname{magnitude}(b))) = 1 \wedge \\ \quad ((\operatorname{magnitude}(c) \cdot \operatorname{magnitude}(c)) + (\operatorname{magnitude}(d) \cdot \operatorname{magnitude}(d))) = 1 \wedge \\ \quad (2 \cdot (a \cdot \operatorname{conjugate}(b))) = (2 \cdot (c \cdot \operatorname{conjugate}(d))) \wedge \\ \quad ((\operatorname{magnitude}(a) \cdot \operatorname{magnitude}(a)) - (\operatorname{magnitude}(b) \cdot \operatorname{magnitude}(b))) \\ \quad =\; ((\operatorname{magnitude}(c) \cdot \operatorname{magnitude}(c)) - (\operatorname{magnitude}(d) \cdot \operatorname{magnitude}(d))) \Rightarrow \\ \quad \quad \exists r \in \mathbb{C}.\; ((\operatorname{magnitude}(r) = 1) \wedge ((c = (r \cdot a)) \wedge (d = (r \cdot b)))) \end{array}}}
\vnbprop{hopf-fibre-iff}{\vnbstmt{\begin{array}{@{}l@{}} \forall u, v \in \operatorname{cc2\text{-}sphere}. \\ \quad \operatorname{hopf\text{-}map}(u) = \operatorname{hopf\text{-}map}(v) \\ \quad \Leftrightarrow\; \exists r \in \mathbb{C}.\; ((\operatorname{magnitude}(r) = 1) \wedge (v = {[}(r \cdot \operatorname{nth}(1, u)), (r \cdot \operatorname{nth}(2, u))])) \end{array}}}
\vnbprop{cc2-sphere-unimodular}{\vnbstmt{\begin{array}{@{}l@{}} \forall a, b, c, d, r \in \mathbb{C}. \\ \quad ((\operatorname{magnitude}(a) \cdot \operatorname{magnitude}(a)) + (\operatorname{magnitude}(b) \cdot \operatorname{magnitude}(b))) = 1 \wedge \\ \quad ((\operatorname{magnitude}(c) \cdot \operatorname{magnitude}(c)) + (\operatorname{magnitude}(d) \cdot \operatorname{magnitude}(d))) = 1 \wedge \\ \quad c = (r \cdot a) \wedge \\ \quad d = (r \cdot b) \Rightarrow \\ \quad \quad \operatorname{magnitude}(r) = 1 \end{array}}}
\vnbprop{homog-point-sphere-eq}{\vnbstmt{\begin{array}{@{}l@{}} \forall a, b, c, d \in \mathbb{C}. \\ \quad ((\operatorname{magnitude}(a) \cdot \operatorname{magnitude}(a)) + (\operatorname{magnitude}(b) \cdot \operatorname{magnitude}(b))) = 1 \wedge \\ \quad ((\operatorname{magnitude}(c) \cdot \operatorname{magnitude}(c)) + (\operatorname{magnitude}(d) \cdot \operatorname{magnitude}(d))) = 1 \wedge \\ \quad \operatorname{homog\text{-}point}(a, b) = \operatorname{homog\text{-}point}(c, d) \Rightarrow \\ \quad \quad \exists r \in \mathbb{C}.\; ((\operatorname{magnitude}(r) = 1) \wedge ((c = (r \cdot a)) \wedge (d = (r \cdot b)))) \end{array}}}
\vnbprop{unit-sphere-cr-subset}{\vnbstmt{\operatorname{unit\text{-}sphere\text{-}cr} \subseteq (\mathbb{C} \times \mathbb{R})}}
\vnbprop{unit-sphere-cr-in-set}{\vnbstmt{\operatorname{unit\text{-}sphere\text{-}cr} \in \mathbf{Set}}}
\vnbprop{hopf-quotient-map-in-fun}{\vnbstmt{\operatorname{hopf\text{-}quotient\text{-}map} \in (\operatorname{unit\text{-}sphere\text{-}cr} \to \operatorname{cc\text{-}hat})}}
\vnbprop{hopf-quotient-map-apply}{\vnbstmt{\begin{array}{@{}l@{}} \forall p \in \operatorname{unit\text{-}sphere\text{-}cr}. \\ \quad \operatorname{hopf\text{-}quotient\text{-}map}(p) = \operatorname{homog\text{-}point}(\operatorname{nth}(1, p), (1 - \operatorname{nth}(2, p))) \end{array}}}
\vnbprop{hopf-quotient-factor}{\vnbstmt{\begin{array}{@{}l@{}} \forall u \in \operatorname{cc2\text{-}sphere}. \\ \quad \operatorname{hopf\text{-}quotient\text{-}map}(\operatorname{hopf\text{-}map}(u)) = \operatorname{homog\text{-}point}(\operatorname{nth}(1, u), \operatorname{nth}(2, u)) \end{array}}}
\vnbprop{hopf-quotient-map-onto}{\vnbstmt{\forall w \in \operatorname{cc\text{-}hat}.\; \exists z \in \operatorname{unit\text{-}sphere\text{-}cr}.\; (\operatorname{hopf\text{-}quotient\text{-}map}(z) = w)}}
\vnbprop{hopf-quotient-map-injective}{\vnbstmt{\begin{array}{@{}l@{}} \forall a, b \in \operatorname{unit\text{-}sphere\text{-}cr}. \\ \quad (\operatorname{hopf\text{-}quotient\text{-}map}(a) = \operatorname{hopf\text{-}quotient\text{-}map}(b)) \Rightarrow (a = b) \end{array}}}
\vnbprop{hopf-quotient-map-bijection}{\vnbstmt{\operatorname{hopf\text{-}quotient\text{-}map} \in \mathrm{Bij}(\operatorname{unit\text{-}sphere\text{-}cr}, \operatorname{cc\text{-}hat})}}
\vnbprop{hopf-factors-through-bijection}{\vnbstmt{\begin{array}{@{}l@{}} \exists g \in (\operatorname{unit\text{-}sphere\text{-}cr} \to \operatorname{cc\text{-}hat}). \\ \quad g \in \mathrm{Bij}(\operatorname{unit\text{-}sphere\text{-}cr}, \operatorname{cc\text{-}hat}) \wedge \\ \quad \forall u \in \operatorname{cc2\text{-}sphere}.\; (g(\operatorname{hopf\text{-}map}(u)) = \operatorname{homog\text{-}point}(\operatorname{nth}(1, u), \operatorname{nth}(2, u))) \end{array}}}
\vnbprop{cc2-sphere-equiv-iff}{\vnbstmt{\begin{array}{@{}l@{}} \forall u, v \in \operatorname{cc2\text{-}sphere}. \\ \quad \operatorname{homog\text{-}equiv}(u, v) \Leftrightarrow \exists r \in \mathbb{C}.\; ((\operatorname{magnitude}(r) = 1) \wedge (u = {[}(r \cdot \operatorname{nth}(1, v)), (r \cdot \operatorname{nth}(2, v))])) \end{array}}}
\vnbprop{hopf-fibre-equiv}{\vnbstmt{\begin{array}{@{}l@{}} \forall u, v \in \operatorname{cc2\text{-}sphere}. \\ \quad (\operatorname{hopf\text{-}map}(u) = \operatorname{hopf\text{-}map}(v)) \Leftrightarrow \operatorname{homog\text{-}equiv}(u, v) \end{array}}}
\vnbprop{hopf-norm-identity-point}{\vnbstmt{\begin{array}{@{}l@{}} \forall w \in (\mathbb{C} \times \mathbb{C}). \\ \quad (\operatorname{magnitude}(\operatorname{nth}(1, \operatorname{hopf\text{-}map}(w))) \cdot \operatorname{magnitude}(\operatorname{nth}(1, \operatorname{hopf\text{-}map}(w)))) + (\operatorname{nth}(2, \operatorname{hopf\text{-}map}(w)) \cdot \operatorname{nth}(2, \operatorname{hopf\text{-}map}(w))) \\ \quad =\; ((\operatorname{magnitude}(\operatorname{nth}(1, w)) \cdot \operatorname{magnitude}(\operatorname{nth}(1, w))) + (\operatorname{magnitude}(\operatorname{nth}(2, w)) \cdot \operatorname{magnitude}(\operatorname{nth}(2, w)))) \cdot ((\operatorname{magnitude}(\operatorname{nth}(1, w)) \cdot \operatorname{magnitude}(\operatorname{nth}(1, w))) + (\operatorname{magnitude}(\operatorname{nth}(2, w)) \cdot \operatorname{magnitude}(\operatorname{nth}(2, w)))) \end{array}}}

\vnbbookfile{theorem-library/product2-hopf}{theorem-library/product2-hopf}
\noindent{\small product2-hopf.scm -- Prop. 1.22 of the complex-analysis notes in its literal form: "pi is continuous C\textasciicircum{}2 -$>$ C x R", with C\textasciicircum{}2 = PRODUCT2-MS(CC-MS, CC-MS) and C x R = PRODUCT2-MS(CC-MS, RR-MS).  Helper prefix p2h-.}\par
\vnbprop{product2-cc-cc-pts}{\vnbstmt{(\operatorname{pts}(\operatorname{product2\text{-}ms}(\operatorname{cc\text{-}ms}, \operatorname{cc\text{-}ms})) \simeq (\mathbb{C} \times \mathbb{C}))}}
\vnbprop{product2-cc-rr-pts}{\vnbstmt{(\operatorname{pts}(\operatorname{product2\text{-}ms}(\operatorname{cc\text{-}ms}, \operatorname{rr\text{-}ms})) \simeq (\mathbb{C} \times \mathbb{R}))}}
\vnbprop{hopf-map-continuous-at-product}{\vnbstmt{\begin{array}{@{}l@{}} \forall w \in (\mathbb{C} \times \mathbb{C}). \\ \quad \operatorname{is\text{-}continuous\text{-}at}(\operatorname{product2\text{-}ms}(\operatorname{cc\text{-}ms}, \operatorname{cc\text{-}ms}), \operatorname{product2\text{-}ms}(\operatorname{cc\text{-}ms}, \operatorname{rr\text{-}ms}), \\ \quad \quad (\lambda v \in (\mathbb{C} \times \mathbb{C}).\; \operatorname{hopf\text{-}map}(v)), w) \end{array}}}
\vnbprop{hopf-map-continuous-product}{\vnbstmt{\begin{array}{@{}l@{}} \operatorname{is\text{-}continuous}(\operatorname{product2\text{-}ms}(\operatorname{cc\text{-}ms}, \operatorname{cc\text{-}ms}), \operatorname{product2\text{-}ms}(\operatorname{cc\text{-}ms}, \operatorname{rr\text{-}ms}), \\ \quad (\lambda v \in (\mathbb{C} \times \mathbb{C}).\; \operatorname{hopf\text{-}map}(v))) \end{array}}}

\vnbbookfile{theorem-library/ccr4e-unif0}{theorem-library/ccr4e-unif0}
\noindent{\small ccr4e-unif0.scm -- CCR.pdf Remark 4.23 IN CHARACTERISTIC 0, under the user's definition of B(K) (2026-11-09): the residue characteristic p, as the element p.1 of K, is a uniformizing parameter of B(K) -- it lies in B(K), 0 $<$ |p.1| $<$ 1, and every x in B(K) with |x| $<$ 1 has |x| $<$= |p.1|.  Agent CCR-1111, 2026-11-11. ...}\par
\vnbprop{zz-one-act-add}{\vnbstmt{\begin{array}{@{}l@{}} \forall r.; \forall a, b \in \mathbb{Z}. \\ \quad \operatorname{is\text{-}ring}(r) \Rightarrow  \\ \quad \quad \operatorname{zz\text{-}act}(\operatorname{ring\text{-}additive\text{-}ag}(r), (a + b), \operatorname{one}(r)) \\ \quad \quad =\; \operatorname{add}(r)(\operatorname{zz\text{-}act}(\operatorname{ring\text{-}additive\text{-}ag}(r), a, \operatorname{one}(r)), \operatorname{zz\text{-}act}(\operatorname{ring\text{-}additive\text{-}ag}(r), b, \operatorname{one}(r))) \end{array}}}
\vnbprop{zz-one-act-neg}{\vnbstmt{\begin{array}{@{}l@{}} \forall r.; \forall a \in \mathbb{Z}. \\ \quad \operatorname{is\text{-}ring}(r) \Rightarrow  \\ \quad \quad \operatorname{zz\text{-}act}(\operatorname{ring\text{-}additive\text{-}ag}(r), -a, \operatorname{one}(r)) = \operatorname{neg}(r)(\operatorname{zz\text{-}act}(\operatorname{ring\text{-}additive\text{-}ag}(r), a, \operatorname{one}(r))) \end{array}}}
\vnbprop{zz-one-act-mul-zz}{\vnbstmt{\begin{array}{@{}l@{}} \forall r.; \forall j, k \in \mathbb{Z}. \\ \quad \operatorname{is\text{-}commutative\text{-}ring}(r) \Rightarrow  \\ \quad \quad \operatorname{zz\text{-}act}(\operatorname{ring\text{-}additive\text{-}ag}(r), (j \cdot k), \operatorname{one}(r)) \\ \quad \quad =\; \operatorname{mul}(r)(\operatorname{zz\text{-}act}(\operatorname{ring\text{-}additive\text{-}ag}(r), k, \operatorname{one}(r)), \operatorname{zz\text{-}act}(\operatorname{ring\text{-}additive\text{-}ag}(r), j, \operatorname{one}(r))) \end{array}}}
\vnbprop{nv-zz-one-norm-le}{\vnbstmt{\begin{array}{@{}l@{}} \forall k.; \forall m \in \mathbb{Z}. \\ \quad \operatorname{is\text{-}nv\text{-}field}(k) \Rightarrow  \\ \quad \quad \operatorname{fnrm}(k)(\operatorname{zz\text{-}act}(\operatorname{ring\text{-}additive\text{-}ag}(\operatorname{normed\text{-}field\text{-}as\text{-}commutative\text{-}ring}(k)), \\ \quad \quad \quad \quad m, \operatorname{one}(\operatorname{normed\text{-}field\text{-}as\text{-}commutative\text{-}ring}(k)))) \\ \quad \quad \leq\; 1 \end{array}}}
\vnbprop{nv-zz-one-unit}{\vnbstmt{\begin{array}{@{}l@{}} \forall k, q.; \forall m \in \mathbb{Z}. \\ \quad \operatorname{is\text{-}nv\text{-}field}(k) \wedge \\ \quad \operatorname{is\text{-}prime\text{-}nat}(q) \wedge \\ \quad \operatorname{fnrm}(k)(\operatorname{zz\text{-}act}(\operatorname{ring\text{-}additive\text{-}ag}(\operatorname{normed\text{-}field\text{-}as\text{-}commutative\text{-}ring}(k)), \\ \quad \quad \quad q, \operatorname{one}(\operatorname{normed\text{-}field\text{-}as\text{-}commutative\text{-}ring}(k)))) \\ \quad <\; 1 \wedge \\ \quad \neg \operatorname{zz\text{-}divides}(q, m) \Rightarrow \\ \quad \quad \operatorname{fnrm}(k)(\operatorname{zz\text{-}act}(\operatorname{ring\text{-}additive\text{-}ag}(\operatorname{normed\text{-}field\text{-}as\text{-}commutative\text{-}ring}(k)), \\ \quad \quad \quad \quad m, \operatorname{one}(\operatorname{normed\text{-}field\text{-}as\text{-}commutative\text{-}ring}(k)))) \\ \quad \quad =\; 1 \end{array}}}
\vnbprop{nv-zz-one-power}{\vnbstmt{\begin{array}{@{}l@{}} \forall a \in \mathbb{N}, k.; \forall q \in \mathbb{Z}. \\ \quad \operatorname{is\text{-}nv\text{-}field}(k) \Rightarrow  \\ \quad \quad \operatorname{fnrm}(k)(\operatorname{zz\text{-}act}(\operatorname{ring\text{-}additive\text{-}ag}(\operatorname{normed\text{-}field\text{-}as\text{-}commutative\text{-}ring}(k)), \\ \quad \quad \quad \quad \operatorname{power}(q, a), \operatorname{one}(\operatorname{normed\text{-}field\text{-}as\text{-}commutative\text{-}ring}(k)))) \\ \quad \quad =\; \operatorname{power}(\operatorname{fnrm}(k)(\operatorname{zz\text{-}act}(\operatorname{ring\text{-}additive\text{-}ag}(\operatorname{normed\text{-}field\text{-}as\text{-}commutative\text{-}ring}(k)), \\ \quad \quad \quad \quad \quad q, \operatorname{one}(\operatorname{normed\text{-}field\text{-}as\text{-}commutative\text{-}ring}(k)))), \\ \quad \quad \quad a) \end{array}}}
\vnbprop{nv-zz-one-value}{\vnbstmt{\begin{array}{@{}l@{}} \forall k, q.; \forall n \in \mathbb{Z}. \\ \quad \operatorname{is\text{-}nv\text{-}field}(k) \wedge \\ \quad \operatorname{is\text{-}prime\text{-}nat}(q) \wedge \\ \quad \operatorname{fnrm}(k)(\operatorname{zz\text{-}act}(\operatorname{ring\text{-}additive\text{-}ag}(\operatorname{normed\text{-}field\text{-}as\text{-}commutative\text{-}ring}(k)), \\ \quad \quad \quad q, \operatorname{one}(\operatorname{normed\text{-}field\text{-}as\text{-}commutative\text{-}ring}(k)))) \\ \quad <\; 1 \wedge \\ \quad \neg (n = 0) \Rightarrow \\ \quad \quad \exists i \in \mathbb{N}. \\ \quad \quad \quad \operatorname{fnrm}(k)(\operatorname{zz\text{-}act}(\operatorname{ring\text{-}additive\text{-}ag}(\operatorname{normed\text{-}field\text{-}as\text{-}commutative\text{-}ring}(k)), \\ \quad \quad \quad \quad \quad n, \operatorname{one}(\operatorname{normed\text{-}field\text{-}as\text{-}commutative\text{-}ring}(k)))) \\ \quad \quad \quad =\; \operatorname{power}(\operatorname{fnrm}(k)(\operatorname{zz\text{-}act}(\operatorname{ring\text{-}additive\text{-}ag}(\operatorname{normed\text{-}field\text{-}as\text{-}commutative\text{-}ring}(k)), \\ \quad \quad \quad \quad \quad \quad q, \operatorname{one}(\operatorname{normed\text{-}field\text{-}as\text{-}commutative\text{-}ring}(k)))), \\ \quad \quad \quad \quad i) \end{array}}}
\vnbprop{nv-zz-one-nonzero}{\vnbstmt{\begin{array}{@{}l@{}} \forall k, q.; \forall b \in \mathbb{Z}. \\ \quad \operatorname{is\text{-}nv\text{-}field}(k) \wedge \\ \quad \operatorname{is\text{-}prime\text{-}nat}(q) \wedge \\ \quad 0 \\ \quad <\; \operatorname{fnrm}(k)(\operatorname{zz\text{-}act}(\operatorname{ring\text{-}additive\text{-}ag}(\operatorname{normed\text{-}field\text{-}as\text{-}commutative\text{-}ring}(k)), \\ \quad \quad \quad q, \operatorname{one}(\operatorname{normed\text{-}field\text{-}as\text{-}commutative\text{-}ring}(k)))) \wedge \\ \quad \operatorname{fnrm}(k)(\operatorname{zz\text{-}act}(\operatorname{ring\text{-}additive\text{-}ag}(\operatorname{normed\text{-}field\text{-}as\text{-}commutative\text{-}ring}(k)), \\ \quad \quad \quad q, \operatorname{one}(\operatorname{normed\text{-}field\text{-}as\text{-}commutative\text{-}ring}(k)))) \\ \quad <\; 1 \wedge \\ \quad \neg (b = 0) \Rightarrow \\ \quad \quad \neg (\operatorname{zz\text{-}act}(\operatorname{ring\text{-}additive\text{-}ag}(\operatorname{normed\text{-}field\text{-}as\text{-}commutative\text{-}ring}(k)), b, \operatorname{one}(\operatorname{normed\text{-}field\text{-}as\text{-}commutative\text{-}ring}(k))) = \operatorname{zero}(k)) \end{array}}}
\vnbprop{nv-rationals-subfield}{\vnbstmt{\begin{array}{@{}l@{}} \forall k, q. \\ \quad \operatorname{is\text{-}nv\text{-}field}(k) \wedge \\ \quad \operatorname{is\text{-}prime\text{-}nat}(q) \wedge \\ \quad 0 \\ \quad <\; \operatorname{fnrm}(k)(\operatorname{zz\text{-}act}(\operatorname{ring\text{-}additive\text{-}ag}(\operatorname{normed\text{-}field\text{-}as\text{-}commutative\text{-}ring}(k)), \\ \quad \quad \quad q, \operatorname{one}(\operatorname{normed\text{-}field\text{-}as\text{-}commutative\text{-}ring}(k)))) \wedge \\ \quad \operatorname{fnrm}(k)(\operatorname{zz\text{-}act}(\operatorname{ring\text{-}additive\text{-}ag}(\operatorname{normed\text{-}field\text{-}as\text{-}commutative\text{-}ring}(k)), \\ \quad \quad \quad q, \operatorname{one}(\operatorname{normed\text{-}field\text{-}as\text{-}commutative\text{-}ring}(k)))) \\ \quad <\; 1 \Rightarrow \\ \quad \quad \operatorname{is\text{-}nf\text{-}subfield}(k, \\ \quad \quad \quad \{x \in \operatorname{carr}(k) : \exists a, b \in \mathbb{Z}.\; (\neg (b = 0) \wedge (\operatorname{mul}(k)(x, \operatorname{zz\text{-}act}(\operatorname{ring\text{-}additive\text{-}ag}(\operatorname{normed\text{-}field\text{-}as\text{-}commutative\text{-}ring}(k)), b, \operatorname{one}(\operatorname{normed\text{-}field\text{-}as\text{-}commutative\text{-}ring}(k)))) = \operatorname{zz\text{-}act}(\operatorname{ring\text{-}additive\text{-}ag}(\operatorname{normed\text{-}field\text{-}as\text{-}commutative\text{-}ring}(k)), a, \operatorname{one}(\operatorname{normed\text{-}field\text{-}as\text{-}commutative\text{-}ring}(k)))))\}) \end{array}}}
\vnbprop{rr-power-antitone}{\vnbstmt{\begin{array}{@{}l@{}} \forall c \in \mathbb{R}, i, j \in \mathbb{N}. \\ \quad ((0 \leq c) \wedge ((c \leq 1) \wedge (i \leq j))) \Rightarrow (\operatorname{power}(c, j) \leq \operatorname{power}(c, i)) \end{array}}}
\vnbprop{nv-rationals-value}{\vnbstmt{\begin{array}{@{}l@{}} \forall k, q.; \forall x \in \{a \in \operatorname{carr}(k) : \exists \mathit{u3a}, b \in \mathbb{Z}.\; (\neg (b = 0) \wedge (\operatorname{mul}(k)(a, \operatorname{zz\text{-}act}(\operatorname{ring\text{-}additive\text{-}ag}(\operatorname{normed\text{-}field\text{-}as\text{-}commutative\text{-}ring}(k)), b, \operatorname{one}(\operatorname{normed\text{-}field\text{-}as\text{-}commutative\text{-}ring}(k)))) = \operatorname{zz\text{-}act}(\operatorname{ring\text{-}additive\text{-}ag}(\operatorname{normed\text{-}field\text{-}as\text{-}commutative\text{-}ring}(k)), \mathit{u3a}, \operatorname{one}(\operatorname{normed\text{-}field\text{-}as\text{-}commutative\text{-}ring}(k)))))\}. \\ \quad \operatorname{is\text{-}nv\text{-}field}(k) \wedge \\ \quad \operatorname{is\text{-}prime\text{-}nat}(q) \wedge \\ \quad 0 \\ \quad <\; \operatorname{fnrm}(k)(\operatorname{zz\text{-}act}(\operatorname{ring\text{-}additive\text{-}ag}(\operatorname{normed\text{-}field\text{-}as\text{-}commutative\text{-}ring}(k)), \\ \quad \quad \quad q, \operatorname{one}(\operatorname{normed\text{-}field\text{-}as\text{-}commutative\text{-}ring}(k)))) \wedge \\ \quad \operatorname{fnrm}(k)(\operatorname{zz\text{-}act}(\operatorname{ring\text{-}additive\text{-}ag}(\operatorname{normed\text{-}field\text{-}as\text{-}commutative\text{-}ring}(k)), \\ \quad \quad \quad q, \operatorname{one}(\operatorname{normed\text{-}field\text{-}as\text{-}commutative\text{-}ring}(k)))) \\ \quad <\; 1 \wedge \\ \quad \operatorname{fnrm}(k)(x) < 1 \Rightarrow \\ \quad \quad \operatorname{fnrm}(k)(x) \\ \quad \quad \leq\; \operatorname{fnrm}(k)(\operatorname{zz\text{-}act}(\operatorname{ring\text{-}additive\text{-}ag}(\operatorname{normed\text{-}field\text{-}as\text{-}commutative\text{-}ring}(k)), \\ \quad \quad \quad \quad q, \operatorname{one}(\operatorname{normed\text{-}field\text{-}as\text{-}commutative\text{-}ring}(k)))) \end{array}}}
\vnbprop{nv-base-char-0-unif}{\vnbstmt{\begin{array}{@{}l@{}} \forall k, q, w.; \forall x \in \operatorname{nv\text{-}base}(k, w). \\ \quad \operatorname{is\text{-}nv\text{-}field}(k) \wedge \\ \quad \forall p.\; \neg \operatorname{has\text{-}characteristic}(\operatorname{normed\text{-}field\text{-}as\text{-}commutative\text{-}ring}(k), p) \wedge \\ \quad \operatorname{is\text{-}prime\text{-}nat}(q) \wedge \\ \quad 0 \\ \quad <\; \operatorname{fnrm}(k)(\operatorname{zz\text{-}act}(\operatorname{ring\text{-}additive\text{-}ag}(\operatorname{normed\text{-}field\text{-}as\text{-}commutative\text{-}ring}(k)), \\ \quad \quad \quad q, \operatorname{one}(\operatorname{normed\text{-}field\text{-}as\text{-}commutative\text{-}ring}(k)))) \wedge \\ \quad \operatorname{fnrm}(k)(\operatorname{zz\text{-}act}(\operatorname{ring\text{-}additive\text{-}ag}(\operatorname{normed\text{-}field\text{-}as\text{-}commutative\text{-}ring}(k)), \\ \quad \quad \quad q, \operatorname{one}(\operatorname{normed\text{-}field\text{-}as\text{-}commutative\text{-}ring}(k)))) \\ \quad <\; 1 \wedge \\ \quad \operatorname{fnrm}(k)(x) < 1 \Rightarrow \\ \quad \quad \operatorname{fnrm}(k)(x) \\ \quad \quad \leq\; \operatorname{fnrm}(k)(\operatorname{zz\text{-}act}(\operatorname{ring\text{-}additive\text{-}ag}(\operatorname{normed\text{-}field\text{-}as\text{-}commutative\text{-}ring}(k)), \\ \quad \quad \quad \quad q, \operatorname{one}(\operatorname{normed\text{-}field\text{-}as\text{-}commutative\text{-}ring}(k)))) \end{array}}}
\vnbprop{ccr-4-23-char-0}{\vnbstmt{\begin{array}{@{}l@{}} \forall k, a, w. \\ \quad \operatorname{is\text{-}nv\text{-}field}(k) \wedge \\ \quad \forall p.\; \neg \operatorname{has\text{-}characteristic}(\operatorname{normed\text{-}field\text{-}as\text{-}commutative\text{-}ring}(k), p) \wedge \\ \quad \operatorname{has\text{-}characteristic}(\operatorname{nv\text{-}residue}(k, \operatorname{nv\text{-}integers}(k)), a) \Rightarrow \\ \quad \quad \operatorname{zz\text{-}act}(\operatorname{ring\text{-}additive\text{-}ag}(\operatorname{normed\text{-}field\text{-}as\text{-}commutative\text{-}ring}(k)), \\ \quad \quad \quad a, \operatorname{one}(\operatorname{normed\text{-}field\text{-}as\text{-}commutative\text{-}ring}(k))) \\ \quad \quad \in\; \operatorname{nv\text{-}base}(k, w) \wedge \\ \quad \quad 0 \\ \quad \quad <\; \operatorname{fnrm}(k)(\operatorname{zz\text{-}act}(\operatorname{ring\text{-}additive\text{-}ag}(\operatorname{normed\text{-}field\text{-}as\text{-}commutative\text{-}ring}(k)), \\ \quad \quad \quad \quad a, \operatorname{one}(\operatorname{normed\text{-}field\text{-}as\text{-}commutative\text{-}ring}(k)))) \wedge \\ \quad \quad \operatorname{fnrm}(k)(\operatorname{zz\text{-}act}(\operatorname{ring\text{-}additive\text{-}ag}(\operatorname{normed\text{-}field\text{-}as\text{-}commutative\text{-}ring}(k)), \\ \quad \quad \quad \quad a, \operatorname{one}(\operatorname{normed\text{-}field\text{-}as\text{-}commutative\text{-}ring}(k)))) \\ \quad \quad <\; 1 \wedge \\ \quad \quad \forall x \in \operatorname{nv\text{-}base}(k, w). \\ \quad \quad \quad (\operatorname{fnrm}(k)(x) < 1) \Rightarrow  \\ \quad \quad \quad \quad \operatorname{fnrm}(k)(x) \\ \quad \quad \quad \quad \leq\; \operatorname{fnrm}(k)(\operatorname{zz\text{-}act}(\operatorname{ring\text{-}additive\text{-}ag}(\operatorname{normed\text{-}field\text{-}as\text{-}commutative\text{-}ring}(k)), \\ \quad \quad \quad \quad \quad \quad a, \operatorname{one}(\operatorname{normed\text{-}field\text{-}as\text{-}commutative\text{-}ring}(k)))) \end{array}}}
\vnbprop{ccr-4-23-char-0-local}{\vnbstmt{\begin{array}{@{}l@{}} \forall k.; \forall w \in \operatorname{nv\text{-}integers}(k). \\ \quad \operatorname{is\text{-}compact\text{-}nv\text{-}domain}(k, \operatorname{nv\text{-}integers}(k)) \wedge \\ \quad \operatorname{nv\text{-}maximal}(k, \operatorname{nv\text{-}integers}(k)) = \operatorname{nv\text{-}principal}(k, \operatorname{nv\text{-}integers}(k), w) \wedge \\ \quad \neg (w = \operatorname{zero}(k)) \wedge \\ \quad \forall p.\; \neg \operatorname{has\text{-}characteristic}(\operatorname{normed\text{-}field\text{-}as\text{-}commutative\text{-}ring}(k), p) \Rightarrow \\ \quad \quad \exists a. \\ \quad \quad \quad \operatorname{has\text{-}characteristic}(\operatorname{nv\text{-}residue}(k, \operatorname{nv\text{-}integers}(k)), a) \wedge \\ \quad \quad \quad \operatorname{zz\text{-}act}(\operatorname{ring\text{-}additive\text{-}ag}(\operatorname{normed\text{-}field\text{-}as\text{-}commutative\text{-}ring}(k)), \\ \quad \quad \quad \quad a, \operatorname{one}(\operatorname{normed\text{-}field\text{-}as\text{-}commutative\text{-}ring}(k))) \\ \quad \quad \quad \in\; \operatorname{nv\text{-}base}(k, w) \wedge \\ \quad \quad \quad 0 \\ \quad \quad \quad <\; \operatorname{fnrm}(k)(\operatorname{zz\text{-}act}(\operatorname{ring\text{-}additive\text{-}ag}(\operatorname{normed\text{-}field\text{-}as\text{-}commutative\text{-}ring}(k)), \\ \quad \quad \quad \quad \quad a, \operatorname{one}(\operatorname{normed\text{-}field\text{-}as\text{-}commutative\text{-}ring}(k)))) \wedge \\ \quad \quad \quad \operatorname{fnrm}(k)(\operatorname{zz\text{-}act}(\operatorname{ring\text{-}additive\text{-}ag}(\operatorname{normed\text{-}field\text{-}as\text{-}commutative\text{-}ring}(k)), \\ \quad \quad \quad \quad \quad a, \operatorname{one}(\operatorname{normed\text{-}field\text{-}as\text{-}commutative\text{-}ring}(k)))) \\ \quad \quad \quad <\; 1 \wedge \\ \quad \quad \quad \forall x \in \operatorname{nv\text{-}base}(k, w). \\ \quad \quad \quad \quad (\operatorname{fnrm}(k)(x) < 1) \Rightarrow  \\ \quad \quad \quad \quad \quad \operatorname{fnrm}(k)(x) \\ \quad \quad \quad \quad \quad \leq\; \operatorname{fnrm}(k)(\operatorname{zz\text{-}act}(\operatorname{ring\text{-}additive\text{-}ag}(\operatorname{normed\text{-}field\text{-}as\text{-}commutative\text{-}ring}(k)), \\ \quad \quad \quad \quad \quad \quad \quad a, \operatorname{one}(\operatorname{normed\text{-}field\text{-}as\text{-}commutative\text{-}ring}(k)))) \end{array}}}
\vnbprop{nv-zz-one-decomp}{\vnbstmt{\begin{array}{@{}l@{}} \forall k, q.; \forall z \in \mathbb{Z}. \\ \quad \operatorname{is\text{-}nv\text{-}field}(k) \wedge \\ \quad \operatorname{is\text{-}prime\text{-}nat}(q) \wedge \\ \quad \operatorname{fnrm}(k)(\operatorname{zz\text{-}act}(\operatorname{ring\text{-}additive\text{-}ag}(\operatorname{normed\text{-}field\text{-}as\text{-}commutative\text{-}ring}(k)), \\ \quad \quad \quad q, \operatorname{one}(\operatorname{normed\text{-}field\text{-}as\text{-}commutative\text{-}ring}(k)))) \\ \quad <\; 1 \wedge \\ \quad \neg (z = 0) \Rightarrow \\ \quad \quad \exists a \in \mathbb{N}, w \in \mathbb{Z}. \\ \quad \quad \quad z = (\operatorname{power}(q, a) \cdot w) \wedge \\ \quad \quad \quad \neg \operatorname{zz\text{-}divides}(q, w) \wedge \\ \quad \quad \quad \operatorname{fnrm}(k)(\operatorname{zz\text{-}act}(\operatorname{ring\text{-}additive\text{-}ag}(\operatorname{normed\text{-}field\text{-}as\text{-}commutative\text{-}ring}(k)), \\ \quad \quad \quad \quad \quad z, \operatorname{one}(\operatorname{normed\text{-}field\text{-}as\text{-}commutative\text{-}ring}(k)))) \\ \quad \quad \quad =\; \operatorname{power}(\operatorname{fnrm}(k)(\operatorname{zz\text{-}act}(\operatorname{ring\text{-}additive\text{-}ag}(\operatorname{normed\text{-}field\text{-}as\text{-}commutative\text{-}ring}(k)), \\ \quad \quad \quad \quad \quad \quad q, \operatorname{one}(\operatorname{normed\text{-}field\text{-}as\text{-}commutative\text{-}ring}(k)))), \\ \quad \quad \quad \quad a) \end{array}}}
\vnbprop{nv-rationals-residue}{\vnbstmt{\begin{array}{@{}l@{}} \forall k, q.; \forall x \in \{a \in \operatorname{carr}(k) : \exists \mathit{u3a}, b \in \mathbb{Z}.\; (\neg (b = 0) \wedge (\operatorname{mul}(k)(a, \operatorname{zz\text{-}act}(\operatorname{ring\text{-}additive\text{-}ag}(\operatorname{normed\text{-}field\text{-}as\text{-}commutative\text{-}ring}(k)), b, \operatorname{one}(\operatorname{normed\text{-}field\text{-}as\text{-}commutative\text{-}ring}(k)))) = \operatorname{zz\text{-}act}(\operatorname{ring\text{-}additive\text{-}ag}(\operatorname{normed\text{-}field\text{-}as\text{-}commutative\text{-}ring}(k)), \mathit{u3a}, \operatorname{one}(\operatorname{normed\text{-}field\text{-}as\text{-}commutative\text{-}ring}(k)))))\}. \\ \quad \operatorname{is\text{-}nv\text{-}field}(k) \wedge \\ \quad \operatorname{is\text{-}prime\text{-}nat}(q) \wedge \\ \quad 0 \\ \quad <\; \operatorname{fnrm}(k)(\operatorname{zz\text{-}act}(\operatorname{ring\text{-}additive\text{-}ag}(\operatorname{normed\text{-}field\text{-}as\text{-}commutative\text{-}ring}(k)), \\ \quad \quad \quad q, \operatorname{one}(\operatorname{normed\text{-}field\text{-}as\text{-}commutative\text{-}ring}(k)))) \wedge \\ \quad \operatorname{fnrm}(k)(\operatorname{zz\text{-}act}(\operatorname{ring\text{-}additive\text{-}ag}(\operatorname{normed\text{-}field\text{-}as\text{-}commutative\text{-}ring}(k)), \\ \quad \quad \quad q, \operatorname{one}(\operatorname{normed\text{-}field\text{-}as\text{-}commutative\text{-}ring}(k)))) \\ \quad <\; 1 \wedge \\ \quad \operatorname{fnrm}(k)(x) \leq 1 \Rightarrow \\ \quad \quad \exists n \in \mathbb{Z}. \\ \quad \quad \quad \operatorname{fnrm}(k)(\operatorname{add}(k)(x, \\ \quad \quad \quad \quad \quad \operatorname{neg}(k)(\operatorname{zz\text{-}act}(\operatorname{ring\text{-}additive\text{-}ag}(\operatorname{normed\text{-}field\text{-}as\text{-}commutative\text{-}ring}(k)), \\ \quad \quad \quad \quad \quad \quad \quad n, \operatorname{one}(\operatorname{normed\text{-}field\text{-}as\text{-}commutative\text{-}ring}(k)))))) \\ \quad \quad \quad <\; 1 \end{array}}}
\vnbprop{nv-base-char-0-residue}{\vnbstmt{\begin{array}{@{}l@{}} \forall k, q, w.; \forall x \in \operatorname{nv\text{-}base}(k, w). \\ \quad \operatorname{is\text{-}nv\text{-}field}(k) \wedge \\ \quad \forall p.\; \neg \operatorname{has\text{-}characteristic}(\operatorname{normed\text{-}field\text{-}as\text{-}commutative\text{-}ring}(k), p) \wedge \\ \quad \operatorname{is\text{-}prime\text{-}nat}(q) \wedge \\ \quad 0 \\ \quad <\; \operatorname{fnrm}(k)(\operatorname{zz\text{-}act}(\operatorname{ring\text{-}additive\text{-}ag}(\operatorname{normed\text{-}field\text{-}as\text{-}commutative\text{-}ring}(k)), \\ \quad \quad \quad q, \operatorname{one}(\operatorname{normed\text{-}field\text{-}as\text{-}commutative\text{-}ring}(k)))) \wedge \\ \quad \operatorname{fnrm}(k)(\operatorname{zz\text{-}act}(\operatorname{ring\text{-}additive\text{-}ag}(\operatorname{normed\text{-}field\text{-}as\text{-}commutative\text{-}ring}(k)), \\ \quad \quad \quad q, \operatorname{one}(\operatorname{normed\text{-}field\text{-}as\text{-}commutative\text{-}ring}(k)))) \\ \quad <\; 1 \wedge \\ \quad \operatorname{fnrm}(k)(x) \leq 1 \Rightarrow \\ \quad \quad \exists n \in \mathbb{Z}. \\ \quad \quad \quad \operatorname{fnrm}(k)(\operatorname{add}(k)(x, \\ \quad \quad \quad \quad \quad \operatorname{neg}(k)(\operatorname{zz\text{-}act}(\operatorname{ring\text{-}additive\text{-}ag}(\operatorname{normed\text{-}field\text{-}as\text{-}commutative\text{-}ring}(k)), \\ \quad \quad \quad \quad \quad \quad \quad n, \operatorname{one}(\operatorname{normed\text{-}field\text{-}as\text{-}commutative\text{-}ring}(k)))))) \\ \quad \quad \quad <\; 1 \end{array}}}

\vnbbookfile{structure-library/ccr4f-vocab}{structure-library/ccr4f-vocab}
\noindent{\small ccr4f-vocab.scm -- THE DIGIT MAP AND THE COORDINATE MAP AS FUNCTIONS (CCR-1111 finding F1). DEFINITIONS ONLY; the laws are in theorem-library/ccr4f-digits.scm.  Agent CCR-1111B, 2026-11-11. Destination: structure-library/ccr4f-vocab.scm.  Library: analysis (IS-NV-EXPANSION of ccr4x-vocab; LINCOMB / MAT of mod-seq).  ...}\par
\par\medskip\noindent\textbf{Definition} (\texttt{nv-digits}).\;$\operatorname{nv\text{-}digits}(\operatorname{ndsk\_}, \operatorname{ndsw\_}, \operatorname{ndsr\_}, \operatorname{ndsx\_}) :=$
\begin{equation*}
(\iota\, \operatorname{ndsa\_}.\;((\operatorname{ndsa\_} \in (\mathbb{Z} \to \operatorname{ndsr\_})) \wedge \exists \operatorname{ndsn\_} \in \mathbb{Z}.\; \operatorname{is\text{-}nv\text{-}expansion}(\operatorname{ndsk\_}, \operatorname{ndsw\_}, \operatorname{ndsr\_}, \operatorname{ndsa\_}, \operatorname{ndsn\_}, \operatorname{ndsx\_})))
\end{equation*}

\par\medskip\noindent\textbf{Definition} (\texttt{lin-coords}).\;$\operatorname{lin\text{-}coords}(\operatorname{ndsm\_}, \operatorname{ndsd\_}, \operatorname{ndsu\_}, \operatorname{ndsp\_}, \operatorname{ndsv\_}) :=$
\begin{equation*}
(\iota\, \operatorname{ndsc\_}.\;((\operatorname{ndsc\_} \in \operatorname{mat}(1, \operatorname{ndsd\_}, \operatorname{ndsp\_})) \wedge (\operatorname{ndsv\_} = \operatorname{lincomb}(\operatorname{ndsm\_}, \operatorname{ndsd\_}, \operatorname{ndsc\_}, \operatorname{ndsu\_}))))
\end{equation*}

\par\medskip\noindent\textbf{Definition} (\texttt{nf-poly-over}).\;$\operatorname{nf\text{-}poly\text{-}over}(\operatorname{ndsk\_}, \operatorname{ndsf\_}) :=$
\begin{equation*}
\{\operatorname{ndsp\_} \in \operatorname{carr}(\operatorname{poly}(\operatorname{normed\text{-}field\text{-}as\text{-}commutative\text{-}ring}(\operatorname{ndsk\_}))) : \forall \operatorname{ndsi\_} \in \mathbb{N}.\; (\operatorname{ndsp\_}(\operatorname{ndsi\_}) \in \operatorname{ndsf\_})\}
\end{equation*}

\par\medskip\noindent\textbf{Definition} (\texttt{nf-rat-at}).\;$\operatorname{nf\text{-}rat\text{-}at}(\operatorname{ndsk\_}, \operatorname{ndsf\_}, \operatorname{ndsw\_}) :=$
\begin{equation*}
\{\operatorname{ndsx\_} \in \operatorname{carr}(\operatorname{ndsk\_}) : \exists \operatorname{ndsp\_}, \operatorname{ndsq\_} \in \operatorname{nf\text{-}poly\text{-}over}(\operatorname{ndsk\_}, \operatorname{ndsf\_}).\; (\neg (\operatorname{poly\text{-}eval}(\operatorname{normed\text{-}field\text{-}as\text{-}commutative\text{-}ring}(\operatorname{ndsk\_}), \operatorname{ndsq\_}, \operatorname{ndsw\_}) = \operatorname{zero}(\operatorname{ndsk\_})) \wedge (\operatorname{mul}(\operatorname{ndsk\_})(\operatorname{ndsx\_}, \operatorname{poly\text{-}eval}(\operatorname{normed\text{-}field\text{-}as\text{-}commutative\text{-}ring}(\operatorname{ndsk\_}), \operatorname{ndsq\_}, \operatorname{ndsw\_})) = \operatorname{poly\text{-}eval}(\operatorname{normed\text{-}field\text{-}as\text{-}commutative\text{-}ring}(\operatorname{ndsk\_}), \operatorname{ndsp\_}, \operatorname{ndsw\_})))\}
\end{equation*}


\vnbbookfile{theorem-library/ccr4f-digits}{theorem-library/ccr4f-digits}
\noindent{\small ccr4f-digits.scm -- THE DIGIT MAP AND THE COORDINATE MAP ARE FUNCTIONS (CCR-1111 finding F1), and lifts of a matrix of residue classes.  Agent CCR-1111B, 2026-11-11.  Destination: theorem-library/ccr4f-digits.scm. Helper prefix nd6-.  Library: analysis (the 4.17 hypotheses; the setoid and matrix lemmas are core / ...}\par
\vnbprop{nv-digits-spec}{\vnbstmt{\begin{array}{@{}l@{}} \forall k, a, w, p.; \forall x \in \operatorname{carr}(k). \\ \quad \operatorname{is\text{-}compact\text{-}nv\text{-}domain}(k, a) \wedge \\ \quad w \in a \wedge \\ \quad \operatorname{nv\text{-}maximal}(k, a) = \operatorname{nv\text{-}principal}(k, a, w) \wedge \\ \quad \neg (w = \operatorname{zero}(k)) \wedge \\ \quad a = \operatorname{nv\text{-}integers}(k) \wedge \\ \quad \operatorname{is\text{-}nv\text{-}rep\text{-}set}(k, a, p) \Rightarrow \\ \quad \quad \operatorname{nv\text{-}digits}(k, w, p, x) \in (\mathbb{Z} \to p) \wedge \\ \quad \quad \exists n \in \mathbb{Z}.\; \operatorname{is\text{-}nv\text{-}expansion}(k, w, p, \operatorname{nv\text{-}digits}(k, w, p, x), n, x) \end{array}}}
\vnbprop{nv-digits-unique}{\vnbstmt{\begin{array}{@{}l@{}} \forall k, a, w, p.; \forall x \in \operatorname{carr}(k), b.; \forall n \in \mathbb{Z}. \\ \quad \operatorname{is\text{-}compact\text{-}nv\text{-}domain}(k, a) \wedge \\ \quad w \in a \wedge \\ \quad \operatorname{nv\text{-}maximal}(k, a) = \operatorname{nv\text{-}principal}(k, a, w) \wedge \\ \quad \neg (w = \operatorname{zero}(k)) \wedge \\ \quad a = \operatorname{nv\text{-}integers}(k) \wedge \\ \quad \operatorname{is\text{-}nv\text{-}rep\text{-}set}(k, a, p) \wedge \\ \quad \operatorname{is\text{-}nv\text{-}expansion}(k, w, p, b, n, x) \Rightarrow \\ \quad \quad b = \operatorname{nv\text{-}digits}(k, w, p, x) \end{array}}}
\vnbprop{nv-digits-fun}{\vnbstmt{\begin{array}{@{}l@{}} \forall k, a, w, p. \\ \quad \operatorname{is\text{-}compact\text{-}nv\text{-}domain}(k, a) \wedge \\ \quad w \in a \wedge \\ \quad \operatorname{nv\text{-}maximal}(k, a) = \operatorname{nv\text{-}principal}(k, a, w) \wedge \\ \quad \neg (w = \operatorname{zero}(k)) \wedge \\ \quad a = \operatorname{nv\text{-}integers}(k) \wedge \\ \quad \operatorname{is\text{-}nv\text{-}rep\text{-}set}(k, a, p) \Rightarrow \\ \quad \quad (\lambda x \in \operatorname{carr}(k).\; \operatorname{nv\text{-}digits}(k, w, p, x)) \in (\operatorname{carr}(k) \to (\mathbb{Z} \to p)) \end{array}}}
\vnbprop{lin-coords-spec}{\vnbstmt{\begin{array}{@{}l@{}} \forall m, d, u, p.; \forall c \in \operatorname{mat}(1, d, p), r. \\ \quad \forall c_{1}, c_{2} \in \operatorname{mat}(1, d, p). \\ \quad \quad (\operatorname{lincomb}(m, d, c_{1}, u) = \operatorname{lincomb}(m, d, c_{2}, u)) \Rightarrow (c_{1} = c_{2}) \wedge \\ \quad r = \operatorname{lincomb}(m, d, c, u) \Rightarrow \\ \quad \quad \operatorname{lin\text{-}coords}(m, d, u, p, r) = c \end{array}}}
\vnbprop{setoid-lift-matrix}{\vnbstmt{\begin{array}{@{}l@{}} \forall s.; \forall m, q \in \mathbb{N}, e \in \operatorname{mat}(m, q, \operatorname{quotient}(s)). \\ \quad \operatorname{is\text{-}setoid}(s) \Rightarrow  \\ \quad \quad \exists b \in \operatorname{mat}(m, q, \operatorname{pts}(s)). \forall i \in \operatorname{interval}(1, m), j \in \operatorname{interval}(1, q). \\ \quad \quad \quad \operatorname{class}(s, \operatorname{entry}(b, i, j)) = \operatorname{entry}(e, i, j) \end{array}}}

\vnbbookfile{theorem-library/ccr4f-poly}{theorem-library/ccr4f-poly}
\noindent{\small ccr4f-poly.scm -- POLYNOMIALS EVALUATED IN A NORMED FIELD, AND THE SUBFIELD GENERATED BY ONE ELEMENT AS QUOTIENTS OF THEIR VALUES (CCR-1111 finding F2).  Agent CCR-1111B, 2026-11-11. Destination: theorem-library/ccr4f-poly.scm.  Helper prefix fp7-.  Library: analysis (normed fields); the first two lemmas are algebra ...}\par
\vnbprop{monalg-mul-coef-in}{\vnbstmt{\begin{array}{@{}l@{}} \forall a, m, s.; \forall f, g \in \operatorname{finsupp}(a, m), x \in \operatorname{carr}(m). \\ \quad \operatorname{is\text{-}ring}(a) \wedge \\ \quad \operatorname{is\text{-}monoid}(m) \wedge \\ \quad \operatorname{zero}(a) \in s \wedge \\ \quad \forall u, v \in s.\; (\operatorname{add}(a)(u, v) \in s) \wedge \\ \quad \forall u, v \in s.\; (\operatorname{mul}(a)(u, v) \in s) \wedge \\ \quad \forall y \in \operatorname{carr}(m).\; (f(y) \in s) \wedge \\ \quad \forall y \in \operatorname{carr}(m).\; (g(y) \in s) \Rightarrow \\ \quad \quad \operatorname{monalg\text{-}mul}(a, m, f, g)(x) \in s \end{array}}}
\vnbprop{poly-mul-coef-in}{\vnbstmt{\begin{array}{@{}l@{}} \forall a, s.; \forall f, g \in \operatorname{carr}(\operatorname{poly}(a)), x \in \mathbb{N}. \\ \quad \operatorname{is\text{-}ring}(a) \wedge \\ \quad \operatorname{zero}(a) \in s \wedge \\ \quad \forall u, v \in s.\; (\operatorname{add}(a)(u, v) \in s) \wedge \\ \quad \forall u, v \in s.\; (\operatorname{mul}(a)(u, v) \in s) \wedge \\ \quad \forall y \in \mathbb{N}.\; (f(y) \in s) \wedge \\ \quad \forall y \in \mathbb{N}.\; (g(y) \in s) \Rightarrow \\ \quad \quad \operatorname{mul}(\operatorname{poly}(a))(f, g)(x) \in s \end{array}}}
\vnbprop{nf-poly-eval-in-subfield}{\vnbstmt{\begin{array}{@{}l@{}} \forall k, s.; \forall w \in s, p \in \operatorname{carr}(\operatorname{poly}(\operatorname{normed\text{-}field\text{-}as\text{-}commutative\text{-}ring}(k))). \\ \quad (\operatorname{is\text{-}normed\text{-}field}(k) \wedge (\operatorname{is\text{-}nf\text{-}subfield}(k, s) \wedge \forall y \in \mathbb{N}.\; (p(y) \in s))) \Rightarrow  \\ \quad \quad \operatorname{poly\text{-}eval}(\operatorname{normed\text{-}field\text{-}as\text{-}commutative\text{-}ring}(k), p, w) \in s \end{array}}}
\vnbprop{nf-rat-at-subfield}{\vnbstmt{\begin{array}{@{}l@{}} \forall k, b.; \forall w \in \operatorname{carr}(k). \\ \quad (\operatorname{is\text{-}normed\text{-}field}(k) \wedge \operatorname{is\text{-}nf\text{-}subfield}(k, b)) \Rightarrow  \\ \quad \quad \operatorname{is\text{-}nf\text{-}subfield}(k, \operatorname{nf\text{-}rat\text{-}at}(k, b, w)) \end{array}}}
\vnbprop{nf-rat-at-contains}{\vnbstmt{\begin{array}{@{}l@{}} \forall k, b.; \forall w \in \operatorname{carr}(k). \\ \quad (\operatorname{is\text{-}normed\text{-}field}(k) \wedge \operatorname{is\text{-}nf\text{-}subfield}(k, b)) \Rightarrow  \\ \quad \quad \forall z \in b.\; (z \in \operatorname{nf\text{-}rat\text{-}at}(k, b, w)) \wedge (w \in \operatorname{nf\text{-}rat\text{-}at}(k, b, w)) \end{array}}}
\vnbprop{nf-rat-at-least}{\vnbstmt{\begin{array}{@{}l@{}} \forall k, b, s.; \forall w \in s, x \in \operatorname{nf\text{-}rat\text{-}at}(k, b, w). \\ \quad (\operatorname{is\text{-}normed\text{-}field}(k) \wedge (\operatorname{is\text{-}nf\text{-}subfield}(k, b) \wedge (\operatorname{is\text{-}nf\text{-}subfield}(k, s) \wedge \forall z \in b.\; (z \in s)))) \Rightarrow  \\ \quad \quad x \in s \end{array}}}

\vnbbookfile{theorem-library/ccr4g-coords}{theorem-library/ccr4g-coords}
\noindent{\small ccr4g-coords.scm -- LINCOMB AS A SUM-AG, and the REARRANGEMENT of a series of linear combinations by coordinates (CCR 4.24, GENERATES half; brief CCR-1111C target 1).  Agent CCR-1111C, 2026-11-11. Destination: theorem-library/ccr4g-coords.scm.  Library: analysis.  Helper prefix cg1-.}\par
\vnbprop{lincomb-sum-ag-ptwise}{\vnbstmt{\begin{array}{@{}l@{}} \forall d.; \forall n \in \mathbb{N}, c \in \operatorname{mat}(1, n, \operatorname{carr}(\operatorname{scal}(d))), u \in \operatorname{mat}(n, 1, \operatorname{vec}(d)), g. \\ \quad \operatorname{is\text{-}module}(d) \wedge \\ \quad \forall i \in \mathrm{OS}(n). \\ \quad \quad (g(i) \simeq \operatorname{act}(d)(\operatorname{entry}(c, 1, (i + 1)), \operatorname{entry}(u, (i + 1), 1))) \Rightarrow \\ \quad \quad \operatorname{lincomb}(d, n, c, u) = \operatorname{sum\text{-}ag}(\operatorname{module\text{-}vector\text{-}ag}(d), g, n) \end{array}}}
\vnbprop{nf-subfield-module-vector-ag}{\vnbstmt{\begin{array}{@{}l@{}} \forall k, s. \\ \quad (\operatorname{is\text{-}normed\text{-}field}(k) \wedge \operatorname{is\text{-}nf\text{-}subfield}(k, s)) \Rightarrow  \\ \quad \quad (\operatorname{module\text{-}vector\text{-}ag}(\operatorname{nf\text{-}subfield\text{-}module}(k, s)) \simeq \operatorname{normed\text{-}field\text{-}additive\text{-}ag}(k)) \end{array}}}
\vnbprop{nf-lincomb-sum-ag}{\vnbstmt{\begin{array}{@{}l@{}} \forall k, s.; \forall n \in \mathbb{N}, c \in \operatorname{mat}(1, n, s), u \in \operatorname{mat}(n, 1, \operatorname{carr}(k)), g. \\ \quad \operatorname{is\text{-}normed\text{-}field}(k) \wedge \\ \quad \operatorname{is\text{-}nf\text{-}subfield}(k, s) \wedge \\ \quad \forall i \in \mathrm{OS}(n). \\ \quad \quad (g(i) \simeq \operatorname{mul}(k)(\operatorname{entry}(c, 1, (i + 1)), \operatorname{entry}(u, (i + 1), 1))) \Rightarrow \\ \quad \quad \operatorname{lincomb}(\operatorname{nf\text{-}subfield\text{-}module}(k, s), n, c, u) = \operatorname{sum\text{-}ag}(\operatorname{normed\text{-}field\text{-}additive\text{-}ag}(k), g, n) \end{array}}}
\vnbprop{nf-ring-additive-ag-view}{\vnbstmt{\begin{array}{@{}l@{}} \forall k. \\ \quad \operatorname{is\text{-}normed\text{-}field}(k) \Rightarrow  \\ \quad \quad (\operatorname{ring\text{-}additive\text{-}ag}(\operatorname{normed\text{-}field\text{-}as\text{-}commutative\text{-}ring}(k)) \simeq \operatorname{normed\text{-}field\text{-}additive\text{-}ag}(k)) \end{array}}}
\vnbprop{nf-series-lincomb-swap}{\vnbstmt{\begin{array}{@{}l@{}} \forall k, s.; \forall n \in \mathbb{N}, b \in \operatorname{mat}(n, 1, \operatorname{carr}(k)), e \in (\mathbb{N} \to \operatorname{mat}(1, n, s)), r \in \operatorname{mat}(1, n, s). \\ \quad \operatorname{is\text{-}normed\text{-}field}(k) \wedge \\ \quad \operatorname{is\text{-}nf\text{-}subfield}(k, s) \wedge \\ \quad \forall j \in \operatorname{interval}(1, n). \\ \quad \quad \operatorname{converges\text{-}to}(\operatorname{nf\text{-}metric\text{-}space}(k), \\ \quad \quad \quad (\lambda m \in \mathbb{N}.\; \operatorname{sum\text{-}ag}(\operatorname{normed\text{-}field\text{-}additive\text{-}ag}(k), (\lambda \mathit{cgk} \in \mathbb{N}.\; \operatorname{entry}(e(\mathit{cgk}), 1, j)), m)), \\ \quad \quad \quad \operatorname{entry}(r, 1, j)) \Rightarrow \\ \quad \quad \operatorname{converges\text{-}to}(\operatorname{nf\text{-}metric\text{-}space}(k), \\ \quad \quad \quad (\lambda m \in \mathbb{N}.\; \\ \quad \quad \quad \quad \operatorname{sum\text{-}ag}(\operatorname{normed\text{-}field\text{-}additive\text{-}ag}(k), \\ \quad \quad \quad \quad \quad (\lambda \mathit{cgk} \in \mathbb{N}.\; \operatorname{lincomb}(\operatorname{nf\text{-}subfield\text{-}module}(k, s), n, e(\mathit{cgk}), b)), \\ \quad \quad \quad \quad \quad m)), \\ \quad \quad \quad \operatorname{lincomb}(\operatorname{nf\text{-}subfield\text{-}module}(k, s), n, r, b)) \end{array}}}

\vnbbookfile{theorem-library/ccr4g-residue}{theorem-library/ccr4g-residue}
\noindent{\small ccr4g-residue.scm -- THE LEADING TERM OF A POLYNOMIAL VALUE, and THE RESIDUE OF F(w) MEET A IS F (CCR 4.24, REL-FREE half; brief CCR-1111C target 2).  Agent CCR-1111C, 2026-11-11. Destination: theorem-library/ccr4g-residue.scm.  Library: analysis.  Helper prefix cg2-.}\par
\vnbprop{rr-power-strict}{\vnbstmt{\begin{array}{@{}l@{}} \forall x \in \mathbb{R}, s, n \in \mathbb{N}. \\ \quad ((0 < x) \wedge ((x < 1) \wedge (s < n))) \Rightarrow (\operatorname{power}(x, n) < \operatorname{power}(x, s)) \end{array}}}
\vnbprop{nv-sum-leading}{\vnbstmt{\begin{array}{@{}l@{}} \forall n \in \mathbb{N}, k, f, w, a, g. \\ \quad \operatorname{is\text{-}nv\text{-}field}(k) \wedge \\ \quad \forall x \in f.\; (x \in \operatorname{carr}(k)) \wedge \\ \quad \forall x \in f.\; (\operatorname{fnrm}(k)(x) \leq 1) \wedge \\ \quad w \in \operatorname{carr}(k) \wedge \\ \quad 0 < \operatorname{fnrm}(k)(w) \wedge \\ \quad \operatorname{fnrm}(k)(w) < 1 \wedge \\ \quad \forall j \in \mathbb{N}.\; (a(j) \in f) \wedge \\ \quad \forall j \in \mathbb{N}. \\ \quad \quad (g(j) \simeq \operatorname{mul}(k)(a(j), \operatorname{ring\text{-}power}(\operatorname{normed\text{-}field\text{-}as\text{-}commutative\text{-}ring}(k), w, j))) \Rightarrow \\ \quad \quad \operatorname{sum\text{-}ag}(\operatorname{normed\text{-}field\text{-}additive\text{-}ag}(k), g, n) = \operatorname{zero}(k) \vee \\ \quad \quad \exists s \in \mathbb{N}, c \in f. \\ \quad \quad \quad \neg (c = \operatorname{zero}(k)) \wedge \\ \quad \quad \quad s < n \wedge \\ \quad \quad \quad \operatorname{fnrm}(k)(\operatorname{add}(k)(\operatorname{sum\text{-}ag}(\operatorname{normed\text{-}field\text{-}additive\text{-}ag}(k), g, n), \\ \quad \quad \quad \quad \quad \operatorname{neg}(k)(\operatorname{mul}(k)(c, \operatorname{ring\text{-}power}(\operatorname{normed\text{-}field\text{-}as\text{-}commutative\text{-}ring}(k), w, s))))) \\ \quad \quad \quad <\; \operatorname{power}(\operatorname{fnrm}(k)(w), s) \end{array}}}
\vnbprop{nv-rat-at-residue}{\vnbstmt{\begin{array}{@{}l@{}} \forall k, f.; \forall w \in \operatorname{carr}(k), y \in \operatorname{nf\text{-}rat\text{-}at}(k, f, w). \\ \quad \operatorname{is\text{-}nv\text{-}field}(k) \wedge \\ \quad \operatorname{is\text{-}nf\text{-}subfield}(k, f) \wedge \\ \quad \forall x \in f.\; (\operatorname{fnrm}(k)(x) \leq 1) \wedge \\ \quad 0 < \operatorname{fnrm}(k)(w) \wedge \\ \quad \operatorname{fnrm}(k)(w) < 1 \wedge \\ \quad \operatorname{fnrm}(k)(y) \leq 1 \Rightarrow \\ \quad \quad \exists c \in f.\; (\operatorname{fnrm}(k)(\operatorname{add}(k)(y, \operatorname{neg}(k)(c))) < 1) \end{array}}}

\part{Measure}
\noindent Sigma-algebras, outer measures and the Caratheodory construction; measures and measurable maps; the integral of simple and measurable functions; sums over the extended nonnegative reals.
\par\noindent\textit{Builds on:} Analysis.
\par\noindent{\small 125 files.}

\vnbbookfile{structure-library/sigma-algebra}{structure-library/sigma-algebra}
\noindent{\small sigma-algebra.scm -- ALGEBRAS and SIGMA-ALGEBRAS of sets.  Vocabulary only: the two definitions of Chapter 1 of the measure-theory notes, cross-checked against Rudin.  Nothing is asserted here.}\par
\par\medskip\noindent\textbf{Definition} (\texttt{is-sigma-algebra}).\;For all $\omega$ and $\mathit{ca}$, 
\begin{equation*}
\begin{array}{@{}l@{}} \operatorname{is\text{-}sigma\text{-}algebra}(\omega, \mathit{ca}) \\ \Leftrightarrow\; \omega \in \mathbf{Set} \wedge \\ \mathit{ca} \subseteq \operatorname{power}(\omega) \wedge \\ \omega \in \mathit{ca} \wedge \\ \forall \operatorname{a\_} \in \mathit{ca}.\; (\operatorname{complement\text{-}in}(\omega, \operatorname{a\_}) \in \mathit{ca}) \wedge \\ \forall \operatorname{f\_} \in (\mathbb{N} \to \mathit{ca}).\; (\operatorname{big\text{-}union}(\operatorname{n\_}, \mathbb{N}, \operatorname{f\_}(\operatorname{n\_})) \in \mathit{ca}) \end{array}
\end{equation*}

\par\medskip\noindent\textbf{Definition} (\texttt{is-algebra-of-sets}).\;For all $\omega$ and $\mathit{ca}$, 
\begin{equation*}
\begin{array}{@{}l@{}} \operatorname{is\text{-}algebra\text{-}of\text{-}sets}(\omega, \mathit{ca}) \\ \Leftrightarrow\; \omega \in \mathbf{Set} \wedge \\ \mathit{ca} \subseteq \operatorname{power}(\omega) \wedge \\ \exists \operatorname{a\_}.\; (\operatorname{a\_} \in \mathit{ca}) \wedge \\ \forall \operatorname{a\_} \in \mathit{ca}.\; (\operatorname{complement\text{-}in}(\omega, \operatorname{a\_}) \in \mathit{ca}) \wedge \\ \forall \operatorname{a\_}, \operatorname{b\_} \in \mathit{ca}.\; ((\operatorname{a\_} \cup \operatorname{b\_}) \in \mathit{ca}) \wedge \\ \forall \operatorname{a\_}, \operatorname{b\_} \in \mathit{ca}.\; ((\operatorname{a\_} \cap \operatorname{b\_}) \in \mathit{ca}) \end{array}
\end{equation*}


\vnbbookfile{structure-library/measure}{structure-library/measure}
\noindent{\small measure.scm -- OUTER MEASURES, CARATHEODORY MEASURABILITY, MEASURES. Vocabulary.  The Caratheodory construction's five statements were asserted here with `reference' warrants until 2026-10-07; they are PROVEN, statements unchanged, in theorem-library/caratheodory.scm (MEASURE-1), and this file}\par
\par\medskip\noindent\textbf{Definition} (\texttt{is-measure}).\;For all $\omega$, $\mathit{ca}$ and $\mu$, 
\begin{equation*}
\begin{array}{@{}l@{}} \operatorname{is\text{-}measure}(\omega, \mathit{ca}, \mu) \\ \Leftrightarrow\; \operatorname{is\text{-}sigma\text{-}algebra}(\omega, \mathit{ca}) \wedge \\ \mu \in (\mathit{ca} \to \operatorname{rr\text{-}pos\text{-}star}) \wedge \\ \mu(\emptyset) = 0 \wedge \\ \forall \operatorname{f\_} \in (\mathbb{N} \to \mathit{ca}). \\ \quad \forall \operatorname{i\_}, \operatorname{j\_} \in \mathbb{N}. \\ \quad \quad \neg (\operatorname{i\_} = \operatorname{j\_}) \Rightarrow ((\operatorname{f\_}(\operatorname{i\_}) \cap \operatorname{f\_}(\operatorname{j\_})) = \emptyset) \Rightarrow \\ \quad \quad \mu(\operatorname{big\text{-}union}(\operatorname{n\_}, \mathbb{N}, \operatorname{f\_}(\operatorname{n\_}))) = \operatorname{esum}((\lambda \operatorname{n\_} \in \mathbb{N}.\; \mu(\operatorname{f\_}(\operatorname{n\_})))) \end{array}
\end{equation*}

\par\medskip\noindent\textbf{Definition} (\texttt{is-caratheodory-measurable}).\;For all $\omega$, $\mu$ and $\mathit{es}$, 
\begin{equation*}
\begin{array}{@{}l@{}} \operatorname{is\text{-}caratheodory\text{-}measurable}(\omega, \mu, \mathit{es}) \\ \Leftrightarrow\; \mathit{es} \subseteq \omega \wedge \\ \forall \operatorname{a\_} \in \operatorname{power}(\omega). \\ \quad \mu(\operatorname{a\_}) = \operatorname{eplus}(\mu((\operatorname{a\_} \cap \mathit{es})), \mu(\operatorname{difference}(\operatorname{a\_}, \mathit{es}))) \end{array}
\end{equation*}

\par\medskip\noindent\textbf{Definition} (\texttt{is-outer-measure}).\;For all $\omega$ and $\mu$, 
\begin{equation*}
\begin{array}{@{}l@{}} \operatorname{is\text{-}outer\text{-}measure}(\omega, \mu) \\ \Leftrightarrow\; \omega \in \mathbf{Set} \wedge \\ \mu \in (\operatorname{power}(\omega) \to \operatorname{rr\text{-}pos\text{-}star}) \wedge \\ \mu(\emptyset) = 0 \wedge \\ \forall \operatorname{a\_}, \operatorname{b\_} \in \operatorname{power}(\omega). \\ \quad (\operatorname{a\_} \subseteq \operatorname{b\_}) \Rightarrow (\mu(\operatorname{a\_}) \leq \mu(\operatorname{b\_})) \wedge \\ \forall \operatorname{f\_} \in (\mathbb{N} \to \operatorname{power}(\omega)). \\ \quad \mu(\operatorname{big\text{-}union}(\operatorname{n\_}, \mathbb{N}, \operatorname{f\_}(\operatorname{n\_}))) \leq \operatorname{esum}((\lambda \operatorname{n\_} \in \mathbb{N}.\; \mu(\operatorname{f\_}(\operatorname{n\_})))) \end{array}
\end{equation*}

\par\medskip\noindent\textbf{Definition} (\texttt{sigma-generated}).\;$\operatorname{sigma\text{-}generated}(\omega, \mathit{ce}) :=$
\begin{equation*}
\{\operatorname{a\_} \in \operatorname{power}(\omega) : \forall \operatorname{cb\_}.\; (((\operatorname{cb\_} \in \operatorname{power}(\operatorname{power}(\omega))) \wedge (\operatorname{is\text{-}sigma\text{-}algebra}(\omega, \operatorname{cb\_}) \wedge (\mathit{ce} \subseteq \operatorname{cb\_}))) \Rightarrow (\operatorname{a\_} \in \operatorname{cb\_}))\}
\end{equation*}

\par\medskip\noindent\textbf{Definition} (\texttt{caratheodory-sets}).\;$\operatorname{caratheodory\text{-}sets}(\omega, \mu) :=$
\begin{equation*}
\{\operatorname{es\_} \in \operatorname{power}(\omega) : \operatorname{is\text{-}caratheodory\text{-}measurable}(\omega, \mu, \operatorname{es\_})\}
\end{equation*}

\par\medskip\noindent\textbf{Definition} (\texttt{is-measurable-map}).\;For all $\omega$, $\mathit{ca}$, $\omega_{2}$, $\mathit{cb}$ and $f$, 
\begin{equation*}
\begin{array}{@{}l@{}} \operatorname{is\text{-}measurable\text{-}map}(\omega, \mathit{ca}, \omega_{2}, \mathit{cb}, f) \\ \Leftrightarrow\; \operatorname{is\text{-}sigma\text{-}algebra}(\omega, \mathit{ca}) \wedge \\ \operatorname{is\text{-}sigma\text{-}algebra}(\omega_{2}, \mathit{cb}) \wedge \\ f \in (\omega \to \omega_{2}) \wedge \\ \forall \operatorname{b\_} \in \mathit{cb}.\; (\{\operatorname{x\_} \in \omega : (f(\operatorname{x\_}) \in \operatorname{b\_})\} \in \mathit{ca}) \end{array}
\end{equation*}

\vnbdef{sigma-generated-membership}{\vnbstmt{\begin{array}{@{}l@{}} \forall a, e, \alpha. \\ \quad \alpha \in \operatorname{sigma\text{-}generated}(a, e) \\ \quad \Leftrightarrow\; \alpha \in \operatorname{power}(a) \wedge \\ \quad \forall b \in \operatorname{power}(\operatorname{power}(a)). \\ \quad \quad (\operatorname{is\text{-}sigma\text{-}algebra}(a, b) \wedge (e \subseteq b)) \Rightarrow (\alpha \in b) \end{array}}}
\vnbdef{caratheodory-sets-membership}{\vnbstmt{\begin{array}{@{}l@{}} \forall a, u, s. \\ \quad (s \in \operatorname{caratheodory\text{-}sets}(a, u)) \Leftrightarrow ((s \in \operatorname{power}(a)) \wedge \operatorname{is\text{-}caratheodory\text{-}measurable}(a, u, s)) \end{array}}}

\vnbbookfile{structure-library/measurable-space}{structure-library/measurable-space}
\noindent{\small measurable-space.scm -- MEASURABLE-SPACE and MEASURE-SPACE as declared STRUCTURES.}\par
\vnbstructdef{measurable-space}{A measurable space $s$ is a 2-tuple whose components are listed below.\par
\begin{enumerate}[nosep]
\item $\operatorname{pts}(s) = \operatorname{nth}(1, s)$, a carrier (\texttt{carriers}): $\operatorname{pts}(s) \in \mathbf{Set}$.
\item $\sigma(s) = \operatorname{nth}(2, s)$, a family of subsets (\texttt{family}): $\sigma(s) \in \operatorname{power}(\operatorname{power}(\operatorname{pts}(s)))$.
\end{enumerate}
The laws (the \texttt{property} and \texttt{law} clauses):
\begin{itemize}[nosep]
\item $\operatorname{is\text{-}sigma\text{-}algebra}(\operatorname{pts}(s), \sigma(s))$
\end{itemize}
The characteristic law (the axiom \texttt{is-measurable-space}): $\operatorname{is\text{-}measurable\text{-}space}(s)$ holds if and only if $\operatorname{length}(s) = 2$, every component has the typing listed with it, every law holds, and $s \in \mathbf{Set}$ (the structure is a set).
}
\vnbstructdef{measure-space}{A measure space $s$ is a 3-tuple whose components are listed below.\par
\begin{enumerate}[nosep]
\item $\operatorname{pts}(s) = \operatorname{nth}(1, s)$, a carrier (\texttt{carriers}): $\operatorname{pts}(s) \in \mathbf{Set}$.
\item $\sigma(s) = \operatorname{nth}(2, s)$, a family of subsets (\texttt{family}): $\sigma(s) \in \operatorname{power}(\operatorname{power}(\operatorname{pts}(s)))$.
\item $\operatorname{meas}(s) = \operatorname{nth}(3, s)$, a function on a family (\texttt{family-fun}): $\operatorname{meas}(s) \in (\sigma(s) \to \operatorname{rr\text{-}pos\text{-}star})$.
\end{enumerate}
The laws (the \texttt{property} and \texttt{law} clauses):
\begin{itemize}[nosep]
\item $\operatorname{is\text{-}measure}(\operatorname{pts}(s), \sigma(s), \operatorname{meas}(s))$
\end{itemize}
The characteristic law (the axiom \texttt{is-measure-space}): $\operatorname{is\text{-}measure\text{-}space}(s)$ holds if and only if $\operatorname{length}(s) = 3$, every component has the typing listed with it, every law holds, and $s \in \mathbf{Set}$ (the structure is a set).
}
\par\medskip\noindent\textbf{Definition} (\texttt{hom-measurable-space}).\;$\operatorname{hom\text{-}measurable\text{-}space}(a, b) :=$
\begin{equation*}
\{\operatorname{homf\_} \in (\operatorname{pts}(a) \to \operatorname{pts}(b)) : \operatorname{is\text{-}hom\text{-}measurable\text{-}space}(a, b, \operatorname{homf\_})\}
\end{equation*}

\par\medskip\noindent\textbf{Definition} (\texttt{hom-measure-space}).\;$\operatorname{hom\text{-}measure\text{-}space}(a, b) :=$
\begin{equation*}
\{\operatorname{homf\_} \in (\operatorname{pts}(a) \to \operatorname{pts}(b)) : \operatorname{is\text{-}hom\text{-}measure\text{-}space}(a, b, \operatorname{homf\_})\}
\end{equation*}

\vnbdef{measure-space-class}{\vnbstmt{\forall s.\; ((s \in \operatorname{measure\text{-}space}) \Leftrightarrow \operatorname{is\text{-}measure\text{-}space}(s))}}
\vnbdef{measurable-space-class}{\vnbstmt{\forall s.\; ((s \in \operatorname{measurable\text{-}space}) \Leftrightarrow \operatorname{is\text{-}measurable\text{-}space}(s))}}
\vnbdef{is-hom-measurable-space-def}{\vnbstmt{\begin{array}{@{}l@{}} \forall s, t, f. \\ \quad \operatorname{is\text{-}hom\text{-}measurable\text{-}space}(s, t, f) \\ \quad \Leftrightarrow\; \operatorname{is\text{-}measurable\text{-}space}(s) \wedge \\ \quad \operatorname{is\text{-}measurable\text{-}space}(t) \wedge \\ \quad f \in (\operatorname{pts}(s) \to \operatorname{pts}(t)) \wedge \\ \quad \forall u \in \sigma(t).\; (\operatorname{preimage}(s, f, u) \in \sigma(s)) \end{array}}}
\vnbdef{is-hom-measure-space-def}{\vnbstmt{\begin{array}{@{}l@{}} \forall s, t, f. \\ \quad \operatorname{is\text{-}hom\text{-}measure\text{-}space}(s, t, f) \\ \quad \Leftrightarrow\; \operatorname{is\text{-}measure\text{-}space}(s) \wedge \\ \quad \operatorname{is\text{-}measure\text{-}space}(t) \wedge \\ \quad f \in (\operatorname{pts}(s) \to \operatorname{pts}(t)) \wedge \\ \quad \forall u \in \sigma(t).\; (\operatorname{preimage}(s, f, u) \in \sigma(s)) \wedge \\ \quad \forall u \in \sigma(t).\; (\operatorname{meas}(t)(u) = \operatorname{meas}(s)(\operatorname{preimage}(s, f, u))) \end{array}}}

\vnbbookfile{structure-library/integral}{structure-library/integral}
\noindent{\small SINCE 2026-10-04 NOTHING IN THIS FILE IS ASSERTED: the fifteen supports that characterised the integral are THEOREMS -- six in theorem-library/integral-laws.scm (M-2), nine in theorem-library/integral-convergence.scm (M-3) -- and INTEGRAL is DEFINED in structure-library/simple-integral.scm (the supremum of the simple}\par
\par\medskip\noindent\textbf{Definition} (\texttt{is-measurable-fn}).\;For all $\omega$, $\mathit{ca}$ and $f$, 
\begin{equation*}
\begin{array}{@{}l@{}} \operatorname{is\text{-}measurable\text{-}fn}(\omega, \mathit{ca}, f) \\ \Leftrightarrow\; \operatorname{is\text{-}sigma\text{-}algebra}(\omega, \mathit{ca}) \wedge \\ f \in (\omega \to \operatorname{rr\text{-}pos\text{-}star}) \wedge \\ \forall \operatorname{alpha\_} \in \mathbb{R}.\; (\{\operatorname{x\_} \in \omega : (\operatorname{alpha\_} < f(\operatorname{x\_}))\} \in \mathit{ca}) \end{array}
\end{equation*}

\par\medskip\noindent\textbf{Definition} (\texttt{ptwise-eplus}).\;$\operatorname{ptwise\text{-}eplus}(\omega, f, g) :=$ $(\lambda \operatorname{x\_} \in \omega.\; \operatorname{eplus}(f(\operatorname{x\_}), g(\operatorname{x\_})))$
\par\medskip\noindent\textbf{Definition} (\texttt{ptwise-etimes}).\;$\operatorname{ptwise\text{-}etimes}(\omega, f, g) :=$ $(\lambda \operatorname{x\_} \in \omega.\; \operatorname{etimes}(f(\operatorname{x\_}), g(\operatorname{x\_})))$
\par\medskip\noindent\textbf{Definition} (\texttt{ptwise-scale}).\;$\operatorname{ptwise\text{-}scale}(\omega, c, f) :=$ $(\lambda \operatorname{x\_} \in \omega.\; \operatorname{etimes}(c, f(\operatorname{x\_})))$
\par\medskip\noindent\textbf{Definition} (\texttt{ptwise-sup}).\;$\operatorname{ptwise\text{-}sup}(\omega, \mathit{cf}) :=$
\begin{equation*}
(\lambda \operatorname{x\_} \in \omega.\; \operatorname{esup}(\operatorname{image}((\lambda \operatorname{n\_} \in \mathbb{N}.\; \operatorname{cf}(\operatorname{n\_})(\operatorname{x\_})), \mathbb{N})))
\end{equation*}

\par\medskip\noindent\textbf{Definition} (\texttt{ptwise-liminf}).\;$\operatorname{ptwise\text{-}liminf}(\omega, \mathit{cf}) :=$
\begin{equation*}
(\lambda \operatorname{x\_} \in \omega.\; \operatorname{eliminf}((\lambda \operatorname{n\_} \in \mathbb{N}.\; \operatorname{cf}(\operatorname{n\_})(\operatorname{x\_}))))
\end{equation*}

\par\medskip\noindent\textbf{Definition} (\texttt{indicator}).\;$\operatorname{indicator}(\omega, a) :=$ $(\lambda \operatorname{x\_} \in \omega.\; \left(\text{if }(\operatorname{x\_} \in a)\text{ then }1\text{ else }0\right))$
\par\medskip\noindent\textbf{Definition} (\texttt{ptwise-le}).\;For all $\omega$, $f$ and $g$, 
\begin{equation*}
\operatorname{ptwise\text{-}le}(\omega, f, g) \Leftrightarrow \forall \operatorname{x\_} \in \omega.\; (f(\operatorname{x\_}) \leq g(\operatorname{x\_}))
\end{equation*}

\par\medskip\noindent\textbf{Definition} (\texttt{is-simple-fn}).\;For all $\omega$, $\mathit{ca}$ and $f$, 
\begin{equation*}
\begin{array}{@{}l@{}} \operatorname{is\text{-}simple\text{-}fn}(\omega, \mathit{ca}, f) \\ \Leftrightarrow\; \operatorname{is\text{-}measurable\text{-}fn}(\omega, \mathit{ca}, f) \wedge \\ |\operatorname{image}(f, \omega)| \in \mathbb{N} \wedge \\ \forall \operatorname{x\_} \in \omega.\; (f(\operatorname{x\_}) \in \mathbb{R}) \end{array}
\end{equation*}


\vnbbookfile{structure-library/simple-integral}{structure-library/simple-integral}
\noindent{\small simple-integral.scm -- THE INTEGRAL OF A [0,+inf]-VALUED FUNCTION, DEFINED. Definitions only; nothing here is asserted and nothing is proved.}\par
\par\medskip\noindent\textbf{Definition} (\texttt{integral}).\;$\operatorname{integral}(\omega, \mathit{ca}, \mu, f) :=$
\begin{equation*}
\operatorname{esup}(\{\operatorname{iny\_} \in \operatorname{rr\text{-}pos\text{-}star} : \exists \operatorname{ins\_} \in (\omega \to \operatorname{rr\text{-}pos\text{-}star}).\; (\operatorname{is\text{-}simple\text{-}fn}(\omega, \mathit{ca}, \operatorname{ins\_}) \wedge (\operatorname{ptwise\text{-}le}(\omega, \operatorname{ins\_}, f) \wedge (\operatorname{iny\_} = \operatorname{simple\text{-}integral}(\omega, \mu, \operatorname{ins\_}))))\})
\end{equation*}

\par\medskip\noindent\textbf{Definition} (\texttt{simple-integral}).\;$\operatorname{simple\text{-}integral}(\omega, \mu, s) :=$
\begin{equation*}
\begin{array}{@{}l@{}} \operatorname{finsum}(\operatorname{rr\text{-}pos\text{-}star\text{-}add\text{-}monoid}, \\ \quad (\lambda \operatorname{siy\_} \in \operatorname{image}(s, \omega).\; \operatorname{etimes}(\operatorname{siy\_}, \mu(\{\operatorname{six\_} \in \omega : (s(\operatorname{six\_}) = \operatorname{siy\_})\}))), \\ \quad \operatorname{image}(s, \omega)) \end{array}
\end{equation*}


\vnbbookfile{structure-library/signed-measure}{structure-library/signed-measure}
\noindent{\small signed-measure.scm -- WAVE 2 (2026-10-27..30), wired 2026-10-30 from scratchpad/mf/mf-vocab.scm. Library: measure.}\par
\par\medskip\noindent\textbf{Definition} (\texttt{ext-neg}).\;$\operatorname{ext\text{-}neg}(\operatorname{enx\_}) :=$
\begin{equation*}
\left(\text{if }(\operatorname{enx\_} = \operatorname{pos\text{-}inf})\text{ then }\operatorname{neg\text{-}inf}\text{ else }\left(\text{if }(\operatorname{enx\_} = \operatorname{neg\text{-}inf})\text{ then }\operatorname{pos\text{-}inf}\text{ else }-\operatorname{enx\_}\right)\right)
\end{equation*}

\par\medskip\noindent\textbf{Definition} (\texttt{ext-add}).\;$\operatorname{ext\text{-}add}(\operatorname{eax\_}, \operatorname{eay\_}) :=$
\begin{equation*}
\left(\text{if }((\operatorname{eax\_} = \operatorname{pos\text{-}inf}) \vee (\operatorname{eay\_} = \operatorname{pos\text{-}inf}))\text{ then }\operatorname{pos\text{-}inf}\text{ else }\left(\text{if }((\operatorname{eax\_} = \operatorname{neg\text{-}inf}) \vee (\operatorname{eay\_} = \operatorname{neg\text{-}inf}))\text{ then }\operatorname{neg\text{-}inf}\text{ else }(\operatorname{eax\_} + \operatorname{eay\_})\right)\right)
\end{equation*}

\par\medskip\noindent\textbf{Definition} (\texttt{ext-sub}).\;$\operatorname{ext\text{-}sub}(\operatorname{esx\_}, \operatorname{esy\_}) :=$ $\operatorname{ext\text{-}add}(\operatorname{esx\_}, \operatorname{ext\text{-}neg}(\operatorname{esy\_}))$
\par\medskip\noindent\textbf{Definition} (\texttt{ext-partial-sum}).\;$\operatorname{ext\text{-}partial\text{-}sum}(x, n) :=$
\begin{equation*}
\left(\text{if }\exists \operatorname{epk\_} \in \mathbb{N}.\; ((\operatorname{epk\_} < n) \wedge (x(\operatorname{epk\_}) = \operatorname{pos\text{-}inf}))\text{ then }\operatorname{pos\text{-}inf}\text{ else }\left(\text{if }\exists \operatorname{epk\_} \in \mathbb{N}.\; ((\operatorname{epk\_} < n) \wedge (x(\operatorname{epk\_}) = \operatorname{neg\text{-}inf}))\text{ then }\operatorname{neg\text{-}inf}\text{ else }\operatorname{series\text{-}partial\text{-}sum}(x, n)\right)\right)
\end{equation*}

\par\medskip\noindent\textbf{Definition} (\texttt{ext-series-converges-to}).\;For all $x$ and $l$, 
\begin{equation*}
\begin{array}{@{}l@{}} \operatorname{ext\text{-}series\text{-}converges\text{-}to}(x, l) \\ \Leftrightarrow\; l \in \operatorname{rr\text{-}star} \wedge \\ (l \in \mathbb{R}) \Rightarrow  \\ \quad \forall \operatorname{ece\_}. \\ \quad \quad \operatorname{pos\text{-}rr}(\operatorname{ece\_}) \Rightarrow  \\ \quad \quad \quad \exists \operatorname{ecm\_} \in \mathbb{N}. \forall \operatorname{ecn\_} \in \mathbb{N}. \\ \quad \quad \quad \quad (\operatorname{ecm\_} \leq \operatorname{ecn\_}) \Rightarrow  \\ \quad \quad \quad \quad \quad (\operatorname{ext\text{-}partial\text{-}sum}(x, \operatorname{ecn\_}) \in \mathbb{R}) \wedge (\lvert -(\operatorname{ext\text{-}partial\text{-}sum}(x, \operatorname{ecn\_}), l) \rvert < \operatorname{ece\_}) \wedge \\ (l = \operatorname{pos\text{-}inf}) \Rightarrow  \\ \quad \forall \operatorname{ecb\_} \in \mathbb{R}. \exists \operatorname{ecm\_} \in \mathbb{N}. \forall \operatorname{ecn\_} \in \mathbb{N}. \\ \quad \quad (\operatorname{ecm\_} \leq \operatorname{ecn\_}) \Rightarrow (\operatorname{ecb\_} < \operatorname{ext\text{-}partial\text{-}sum}(x, \operatorname{ecn\_})) \wedge \\ (l = \operatorname{neg\text{-}inf}) \Rightarrow  \\ \quad \forall \operatorname{ecb\_} \in \mathbb{R}. \exists \operatorname{ecm\_} \in \mathbb{N}. \forall \operatorname{ecn\_} \in \mathbb{N}. \\ \quad \quad (\operatorname{ecm\_} \leq \operatorname{ecn\_}) \Rightarrow (\operatorname{ext\text{-}partial\text{-}sum}(x, \operatorname{ecn\_}) < \operatorname{ecb\_}) \end{array}
\end{equation*}

\par\medskip\noindent\textbf{Definition} (\texttt{is-signed-measure}).\;For all $\omega$, $\mathit{ca}$ and $\nu$, 
\begin{equation*}
\begin{array}{@{}l@{}} \operatorname{is\text{-}signed\text{-}measure}(\omega, \mathit{ca}, \nu) \\ \Leftrightarrow\; \operatorname{is\text{-}sigma\text{-}algebra}(\omega, \mathit{ca}) \wedge \\ \nu \in (\mathit{ca} \to \operatorname{rr\text{-}star}) \wedge \\ \nu(\emptyset) = 0 \wedge \\ \neg (\exists \operatorname{sma\_} \in \mathit{ca}.\; (\nu(\operatorname{sma\_}) = \operatorname{pos\text{-}inf}) \wedge \exists \operatorname{smb\_} \in \mathit{ca}.\; (\nu(\operatorname{smb\_}) = \operatorname{neg\text{-}inf})) \wedge \\ \forall \operatorname{smf\_} \in (\mathbb{N} \to \mathit{ca}). \\ \quad \forall \operatorname{smi\_}, \operatorname{smj\_} \in \mathbb{N}. \\ \quad \quad \neg (\operatorname{smi\_} = \operatorname{smj\_}) \Rightarrow ((\operatorname{smf\_}(\operatorname{smi\_}) \cap \operatorname{smf\_}(\operatorname{smj\_})) = \emptyset) \Rightarrow \\ \quad \quad \operatorname{ext\text{-}series\text{-}converges\text{-}to}((\lambda \operatorname{smn\_} \in \mathbb{N}.\; \nu(\operatorname{smf\_}(\operatorname{smn\_}))), \nu(\operatorname{big\text{-}union}(\operatorname{smn\_}, \mathbb{N}, \operatorname{smf\_}(\operatorname{smn\_})))) \end{array}
\end{equation*}

\par\medskip\noindent\textbf{Definition} (\texttt{pos-variation}).\;$\operatorname{pos\text{-}variation}(\mathit{ca}, \nu) :=$
\begin{equation*}
\begin{array}{@{}l@{}} (\lambda \operatorname{pva\_} \in \mathit{ca}.\; \\ \quad \operatorname{esup}(\{\operatorname{pvv\_} \in \operatorname{rr\text{-}pos\text{-}star} : \exists \operatorname{pve\_} \in \mathit{ca}.\; ((\operatorname{pve\_} \subseteq \operatorname{pva\_}) \wedge (\nu(\operatorname{pve\_}) = \operatorname{pvv\_}))\})) \end{array}
\end{equation*}

\par\medskip\noindent\textbf{Definition} (\texttt{signed-neg}).\;$\operatorname{signed\text{-}neg}(\mathit{ca}, \nu) :=$ $(\lambda \operatorname{sna\_} \in \mathit{ca}.\; \operatorname{ext\text{-}neg}(\nu(\operatorname{sna\_})))$
\par\medskip\noindent\textbf{Definition} (\texttt{neg-variation}).\;$\operatorname{neg\text{-}variation}(\mathit{ca}, \nu) :=$ $\operatorname{pos\text{-}variation}(\mathit{ca}, \operatorname{signed\text{-}neg}(\mathit{ca}, \nu))$
\par\medskip\noindent\textbf{Definition} (\texttt{total-variation}).\;$\operatorname{total\text{-}variation}(\mathit{ca}, \nu) :=$
\begin{equation*}
(\lambda \operatorname{tva\_} \in \mathit{ca}.\; \operatorname{eplus}(\operatorname{pos\text{-}variation}(\mathit{ca}, \nu)(\operatorname{tva\_}), \operatorname{neg\text{-}variation}(\mathit{ca}, \nu)(\operatorname{tva\_})))
\end{equation*}

\par\medskip\noindent\textbf{Definition} (\texttt{total-variation-norm}).\;$\operatorname{total\text{-}variation\text{-}norm}(\omega, \mathit{ca}, \nu) :=$ $\operatorname{total\text{-}variation}(\mathit{ca}, \nu)(\omega)$

\vnbbookfile{structure-library/is-set-laws}{structure-library/is-set-laws}
\noindent{\small A file shared across books; the entries that belong to this one.}\par
\noindent{\small structure-library/is-set-laws.scm -- THE GENERATED SETHOOD THEOREMS (W1I, 2026-11-06).}\par
\vnbprop{is-measurable-space-is-set}{\vnbstmt{\forall s.\; (\operatorname{is\text{-}measurable\text{-}space}(s) \Rightarrow (s \in \mathbf{Set}))}}
\vnbprop{is-measure-space-is-set}{\vnbstmt{\forall s.\; (\operatorname{is\text{-}measure\text{-}space}(s) \Rightarrow (s \in \mathbf{Set}))}}

\vnbbookfile{theorem-library/view-typing-laws}{theorem-library/view-typing-laws}
\noindent{\small A file shared across books; the entries that belong to this one.}\par
\noindent{\small theorem-library/view-typing-laws.scm -- THE VIEW TYPING THEOREMS, PROVEN (D8, 2026-11-15).}\par
\vnbprop{measure-space-as-measurable-space-is-measurable-space}{\vnbstmt{\begin{array}{@{}l@{}} \forall r. \\ \quad \operatorname{is\text{-}measure\text{-}space}(r) \Rightarrow  \\ \quad \quad \operatorname{is\text{-}measurable\text{-}space}(\operatorname{measure\text{-}space\text{-}as\text{-}measurable\text{-}space}(r)) \end{array}}}

\vnbbookfile{theorem-library/discrete-space}{theorem-library/discrete-space}
\noindent{\small A file shared across books; the entries that belong to this one.}\par
\noindent{\small discrete-space.scm -- THE FIRST WITNESS OF A MEASURABLE SPACE, and the first outright inhabitant of TOP-SPACE.}\par
\vnbprop{discrete-space-is-measurable-space}{\vnbstmt{\forall a \in \mathbf{Set}.\; \operatorname{is\text{-}measurable\text{-}space}(\operatorname{discrete\text{-}space}(a))}}

\vnbbookfile{theorem-library/rake-measure}{theorem-library/rake-measure}
\noindent{\small rake-measure.scm -- rake batch 5, batch W: the light end of measure theory.}\par
\vnbprop{sigma-algebra-ambient-set}{\vnbstmt{\forall a, b.\; (\operatorname{is\text{-}sigma\text{-}algebra}(a, b) \Rightarrow (a \in \mathbf{Set}))}}
\vnbprop{sigma-algebra-whole-in}{\vnbstmt{\forall a, b.\; (\operatorname{is\text{-}sigma\text{-}algebra}(a, b) \Rightarrow (a \in b))}}
\vnbprop{sigma-algebra-complement-closed}{\vnbstmt{\begin{array}{@{}l@{}} \forall a, b, \alpha. \\ \quad (\operatorname{is\text{-}sigma\text{-}algebra}(a, b) \wedge (\alpha \in b)) \Rightarrow (\operatorname{complement\text{-}in}(a, \alpha) \in b) \end{array}}}
\vnbprop{sigma-algebra-union-closed}{\vnbstmt{\begin{array}{@{}l@{}} \forall a, b, f. \\ \quad (\operatorname{is\text{-}sigma\text{-}algebra}(a, b) \wedge (f \in (\mathbb{N} \to b))) \Rightarrow (\operatorname{big\text{-}union}(n, \mathbb{N}, f(n)) \in b) \end{array}}}
\vnbprop{sigma-algebra-in-power-power}{\vnbstmt{\forall a, b.\; (\operatorname{is\text{-}sigma\text{-}algebra}(a, b) \Rightarrow (b \in \operatorname{power}(\operatorname{power}(a))))}}
\vnbprop{sigma-generated-in-power}{\vnbstmt{\forall a, e.; \forall \alpha \in \operatorname{sigma\text{-}generated}(a, e).\; (\alpha \in \operatorname{power}(a))}}
\vnbprop{sigma-generated-in-every}{\vnbstmt{\begin{array}{@{}l@{}} \forall a, e, \alpha, b. \\ \quad ((\alpha \in \operatorname{sigma\text{-}generated}(a, e)) \wedge (\operatorname{is\text{-}sigma\text{-}algebra}(a, b) \wedge (e \subseteq b))) \Rightarrow  \\ \quad \quad \alpha \in b \end{array}}}
\vnbprop{sigma-generated-intro}{\vnbstmt{\begin{array}{@{}l@{}} \forall a, e.; \forall \alpha \in \operatorname{power}(a). \\ \quad \forall b.\; ((\operatorname{is\text{-}sigma\text{-}algebra}(a, b) \wedge (e \subseteq b)) \Rightarrow (\alpha \in b)) \Rightarrow  \\ \quad \quad \alpha \in \operatorname{sigma\text{-}generated}(a, e) \end{array}}}
\vnbprop{sigma-generated-is-sigma-algebra}{\vnbstmt{\begin{array}{@{}l@{}} \forall a, e. \\ \quad ((a \in \mathbf{Set}) \wedge (e \subseteq \operatorname{power}(a))) \Rightarrow  \\ \quad \quad \operatorname{is\text{-}sigma\text{-}algebra}(a, \operatorname{sigma\text{-}generated}(a, e)) \end{array}}}
\vnbprop{measurable-map-source-algebra}{\vnbstmt{\begin{array}{@{}l@{}} \forall a, b, a_{2}, c, f. \\ \quad \operatorname{is\text{-}measurable\text{-}map}(a, b, a_{2}, c, f) \Rightarrow \operatorname{is\text{-}sigma\text{-}algebra}(a, b) \end{array}}}
\vnbprop{measurable-map-target-algebra}{\vnbstmt{\begin{array}{@{}l@{}} \forall a, b, a_{2}, c, f. \\ \quad \operatorname{is\text{-}measurable\text{-}map}(a, b, a_{2}, c, f) \Rightarrow \operatorname{is\text{-}sigma\text{-}algebra}(a_{2}, c) \end{array}}}
\vnbprop{measurable-map-in-fun}{\vnbstmt{\begin{array}{@{}l@{}} \forall a, b, a_{2}, c, f. \\ \quad \operatorname{is\text{-}measurable\text{-}map}(a, b, a_{2}, c, f) \Rightarrow (f \in (a \to a_{2})) \end{array}}}
\vnbprop{measurable-map-preimage}{\vnbstmt{\begin{array}{@{}l@{}} \forall a, b, a_{2}, c, f, \beta. \\ \quad (\operatorname{is\text{-}measurable\text{-}map}(a, b, a_{2}, c, f) \wedge (\beta \in c)) \Rightarrow  \\ \quad \quad \{x \in a : (f(x) \in \beta)\} \in b \end{array}}}
\vnbprop{measurable-map-intro}{\vnbstmt{\begin{array}{@{}l@{}} \forall a, b, a_{2}, c, f. \\ \quad \operatorname{is\text{-}sigma\text{-}algebra}(a, b) \wedge \\ \quad \operatorname{is\text{-}sigma\text{-}algebra}(a_{2}, c) \wedge \\ \quad f \in (a \to a_{2}) \wedge \\ \quad \forall \beta \in c.\; (\{x \in a : (f(x) \in \beta)\} \in b) \Rightarrow \\ \quad \quad \operatorname{is\text{-}measurable\text{-}map}(a, b, a_{2}, c, f) \end{array}}}
\vnbprop{measurable-map-compose}{\vnbstmt{\begin{array}{@{}l@{}} \forall a, b, a_{2}, c, a_{3}, d, f, g. \\ \quad (\operatorname{is\text{-}measurable\text{-}map}(a, b, a_{2}, c, f) \wedge \operatorname{is\text{-}measurable\text{-}map}(a_{2}, c, a_{3}, d, g)) \Rightarrow  \\ \quad \quad \operatorname{is\text{-}measurable\text{-}map}(a, b, a_{3}, d, \operatorname{compose}(g, f)) \end{array}}}

\vnbbookfile{theorem-library/rake-measure2}{theorem-library/rake-measure2}
\noindent{\small rake-measure2.scm -- rake batch 5b-F: the indicator function, and the four SET-LEVEL bricks sigma-algebra.scm's header lists as "NEXT (not stated)".}\par
\vnbprop{sigma-algebra-empty-in}{\vnbstmt{\forall a, b.\; (\operatorname{is\text{-}sigma\text{-}algebra}(a, b) \Rightarrow (\emptyset \in b))}}
\vnbprop{sigma-algebra-union-closed-2}{\vnbstmt{\begin{array}{@{}l@{}} \forall a, b, \alpha, \beta. \\ \quad (\operatorname{is\text{-}sigma\text{-}algebra}(a, b) \wedge ((\alpha \in b) \wedge (\beta \in b))) \Rightarrow  \\ \quad \quad (\alpha \cup \beta) \in b \end{array}}}
\vnbprop{sigma-algebra-inter-closed-2}{\vnbstmt{\begin{array}{@{}l@{}} \forall a, b, \alpha, \beta. \\ \quad (\operatorname{is\text{-}sigma\text{-}algebra}(a, b) \wedge ((\alpha \in b) \wedge (\beta \in b))) \Rightarrow  \\ \quad \quad (\alpha \cap \beta) \in b \end{array}}}
\vnbprop{sigma-algebra-difference-closed}{\vnbstmt{\begin{array}{@{}l@{}} \forall a, b, \alpha, \beta. \\ \quad (\operatorname{is\text{-}sigma\text{-}algebra}(a, b) \wedge ((\alpha \in b) \wedge (\beta \in b))) \Rightarrow  \\ \quad \quad \operatorname{complement\text{-}in}(\alpha, \beta) \in b \end{array}}}
\vnbprop{measurable-fn-indicator}{\vnbstmt{\begin{array}{@{}l@{}} \forall a, b, \alpha. \\ \quad (\operatorname{is\text{-}sigma\text{-}algebra}(a, b) \wedge (\alpha \in b)) \Rightarrow  \\ \quad \quad \operatorname{is\text{-}measurable\text{-}fn}(a, b, \operatorname{indicator}(a, \alpha)) \end{array}}}

\vnbbookfile{theorem-library/hom-kinds}{theorem-library/hom-kinds}
\noindent{\small A file shared across books; the entries that belong to this one.}\par
\noindent{\small hom-kinds.scm -- the category laws for the arrows the generic driver of hom-laws.scm does not reach: the FAMILY slot kinds, and the RELATION kind (SETOID, since batch 40; two carriers in batch 39).}\par
\vnbprop{hom-measurable-space-member-iff}{\vnbstmt{\begin{array}{@{}l@{}} \forall a, b, f. \\ \quad (f \in \operatorname{hom\text{-}measurable\text{-}space}(a, b)) \Leftrightarrow ((f \in (\operatorname{pts}(a) \to \operatorname{pts}(b))) \wedge \operatorname{is\text{-}hom\text{-}measurable\text{-}space}(a, b, f)) \end{array}}}
\vnbprop{hom-measurable-space-id}{\vnbstmt{\begin{array}{@{}l@{}} \forall a. \\ \quad \operatorname{is\text{-}measurable\text{-}space}(a) \Rightarrow \operatorname{is\text{-}hom\text{-}measurable\text{-}space}(a, a, \operatorname{id\text{-}fun}(\operatorname{pts}(a))) \end{array}}}
\vnbprop{hom-measurable-space-in-set}{\vnbstmt{\begin{array}{@{}l@{}} \forall a, b. \\ \quad (\operatorname{is\text{-}measurable\text{-}space}(a) \wedge \operatorname{is\text{-}measurable\text{-}space}(b)) \Rightarrow  \\ \quad \quad \operatorname{hom\text{-}measurable\text{-}space}(a, b) \in \mathbf{Set} \end{array}}}
\vnbprop{hom-measure-space-member-iff}{\vnbstmt{\begin{array}{@{}l@{}} \forall a, b, f. \\ \quad (f \in \operatorname{hom\text{-}measure\text{-}space}(a, b)) \Leftrightarrow ((f \in (\operatorname{pts}(a) \to \operatorname{pts}(b))) \wedge \operatorname{is\text{-}hom\text{-}measure\text{-}space}(a, b, f)) \end{array}}}
\vnbprop{hom-measure-space-id}{\vnbstmt{\begin{array}{@{}l@{}} \forall a. \\ \quad \operatorname{is\text{-}measure\text{-}space}(a) \Rightarrow \operatorname{is\text{-}hom\text{-}measure\text{-}space}(a, a, \operatorname{id\text{-}fun}(\operatorname{pts}(a))) \end{array}}}
\vnbprop{hom-measure-space-in-set}{\vnbstmt{\begin{array}{@{}l@{}} \forall a, b. \\ \quad (\operatorname{is\text{-}measure\text{-}space}(a) \wedge \operatorname{is\text{-}measure\text{-}space}(b)) \Rightarrow  \\ \quad \quad \operatorname{hom\text{-}measure\text{-}space}(a, b) \in \mathbf{Set} \end{array}}}
\vnbprop{measure-space-as-measurable-space-functorial}{\vnbstmt{\begin{array}{@{}l@{}} \forall a, b, f. \\ \quad \operatorname{is\text{-}hom\text{-}measure\text{-}space}(a, b, f) \Rightarrow  \\ \quad \quad \operatorname{is\text{-}hom\text{-}measurable\text{-}space}(\operatorname{measure\text{-}space\text{-}as\text{-}measurable\text{-}space}(a), \\ \quad \quad \quad \operatorname{measure\text{-}space\text{-}as\text{-}measurable\text{-}space}(b), f) \end{array}}}

\vnbbookfile{theorem-library/hom-functors}{theorem-library/hom-functors}
\noindent{\small A file shared across books; the entries that belong to this one.}\par
\noindent{\small hom-functors.scm -- the Hom functors: Hom(a, -) covariant, Hom(-, b) contravariant, as maps on hom-sets.}\par
\vnbprop{hom-measurable-space-post-type}{\vnbstmt{\begin{array}{@{}l@{}} \forall a, b, c, h. \\ \quad (\operatorname{is\text{-}measurable\text{-}space}(a) \wedge (h \in \operatorname{hom\text{-}measurable\text{-}space}(b, c))) \Rightarrow  \\ \quad \quad (\lambda f \in \operatorname{hom\text{-}measurable\text{-}space}(a, b).\; \operatorname{compose}(h, f)) \\ \quad \quad \in\; (\operatorname{hom\text{-}measurable\text{-}space}(a, b) \to \operatorname{hom\text{-}measurable\text{-}space}(a, c)) \end{array}}}
\vnbprop{hom-measurable-space-pre-type}{\vnbstmt{\begin{array}{@{}l@{}} \forall a, b, c, k. \\ \quad (\operatorname{is\text{-}measurable\text{-}space}(b) \wedge (k \in \operatorname{hom\text{-}measurable\text{-}space}(c, a))) \Rightarrow  \\ \quad \quad (\lambda f \in \operatorname{hom\text{-}measurable\text{-}space}(a, b).\; \operatorname{compose}(f, k)) \\ \quad \quad \in\; (\operatorname{hom\text{-}measurable\text{-}space}(a, b) \to \operatorname{hom\text{-}measurable\text{-}space}(c, b)) \end{array}}}
\vnbprop{hom-measure-space-post-type}{\vnbstmt{\begin{array}{@{}l@{}} \forall a, b, c, h. \\ \quad (\operatorname{is\text{-}measure\text{-}space}(a) \wedge (h \in \operatorname{hom\text{-}measure\text{-}space}(b, c))) \Rightarrow  \\ \quad \quad (\lambda f \in \operatorname{hom\text{-}measure\text{-}space}(a, b).\; \operatorname{compose}(h, f)) \\ \quad \quad \in\; (\operatorname{hom\text{-}measure\text{-}space}(a, b) \to \operatorname{hom\text{-}measure\text{-}space}(a, c)) \end{array}}}
\vnbprop{hom-measure-space-pre-type}{\vnbstmt{\begin{array}{@{}l@{}} \forall a, b, c, k. \\ \quad (\operatorname{is\text{-}measure\text{-}space}(b) \wedge (k \in \operatorname{hom\text{-}measure\text{-}space}(c, a))) \Rightarrow  \\ \quad \quad (\lambda f \in \operatorname{hom\text{-}measure\text{-}space}(a, b).\; \operatorname{compose}(f, k)) \\ \quad \quad \in\; (\operatorname{hom\text{-}measure\text{-}space}(a, b) \to \operatorname{hom\text{-}measure\text{-}space}(c, b)) \end{array}}}

\vnbbookfile{theorem-library/measure-laws}{theorem-library/measure-laws}
\noindent{\small measure-laws.scm -- Chapter 1 of the measure notes (Thayer, Construction of Measures, Sec. 3-4): the elementary laws of a countably additive measure. Agent M-1 of the October roadmap (docs/roadmap-2026-10-03.md, line (b), decision 4).  Helper prefix mlw-.  25 theorems, every one `modulo 0'.}\par
\vnbprop{rps-prefix-zero}{\vnbstmt{\forall f.\; (\operatorname{finsum}(\operatorname{rr\text{-}pos\text{-}star\text{-}add\text{-}monoid}, f, \mathrm{OS}(0)) = 0)}}
\vnbprop{rps-prefix-succ}{\vnbstmt{\begin{array}{@{}l@{}} \forall f \in (\mathbb{N} \to \operatorname{rr\text{-}pos\text{-}star}), n \in \mathbb{N}. \\ \quad \operatorname{finsum}(\operatorname{rr\text{-}pos\text{-}star\text{-}add\text{-}monoid}, f, \mathrm{OS}((n + 1))) \\ \quad =\; \operatorname{eplus}(\operatorname{finsum}(\operatorname{rr\text{-}pos\text{-}star\text{-}add\text{-}monoid}, f, \mathrm{OS}(n)), f(n)) \end{array}}}
\vnbprop{esum-nn-prefix-le}{\vnbstmt{\begin{array}{@{}l@{}} \forall f \in (\mathbb{N} \to \operatorname{rr\text{-}pos\text{-}star}), n \in \mathbb{N}. \\ \quad \operatorname{finsum}(\operatorname{rr\text{-}pos\text{-}star\text{-}add\text{-}monoid}, f, \mathrm{OS}(n)) \leq \operatorname{esum}(f) \end{array}}}
\vnbprop{esum-nn-le-of-prefixes}{\vnbstmt{\begin{array}{@{}l@{}} \forall f \in (\mathbb{N} \to \operatorname{rr\text{-}pos\text{-}star}), b \in \operatorname{rr\text{-}pos\text{-}star}. \\ \quad \forall n \in \mathbb{N}.\; (\operatorname{finsum}(\operatorname{rr\text{-}pos\text{-}star\text{-}add\text{-}monoid}, f, \mathrm{OS}(n)) \leq b) \Rightarrow  \\ \quad \quad \operatorname{esum}(f) \leq b \end{array}}}
\vnbprop{rps-prefix-type}{\vnbstmt{\begin{array}{@{}l@{}} \forall f \in (\mathbb{N} \to \operatorname{rr\text{-}pos\text{-}star}), n \in \mathbb{N}. \\ \quad \operatorname{finsum}(\operatorname{rr\text{-}pos\text{-}star\text{-}add\text{-}monoid}, f, \mathrm{OS}(n)) \in \operatorname{rr\text{-}pos\text{-}star} \end{array}}}
\vnbprop{rps-prefix-le-termwise}{\vnbstmt{\begin{array}{@{}l@{}} \forall n \in \mathbb{N}, f, g. \\ \quad ((f \in (\mathbb{N} \to \operatorname{rr\text{-}pos\text{-}star})) \wedge ((g \in (\mathbb{N} \to \operatorname{rr\text{-}pos\text{-}star})) \wedge \forall k \in \mathbb{N}.\; (f(k) \leq g(k)))) \Rightarrow  \\ \quad \quad \operatorname{finsum}(\operatorname{rr\text{-}pos\text{-}star\text{-}add\text{-}monoid}, f, \mathrm{OS}(n)) \\ \quad \quad \leq\; \operatorname{finsum}(\operatorname{rr\text{-}pos\text{-}star\text{-}add\text{-}monoid}, g, \mathrm{OS}(n)) \end{array}}}
\vnbprop{esum-nn-le-termwise}{\vnbstmt{\begin{array}{@{}l@{}} \forall f, g. \\ \quad ((f \in (\mathbb{N} \to \operatorname{rr\text{-}pos\text{-}star})) \wedge ((g \in (\mathbb{N} \to \operatorname{rr\text{-}pos\text{-}star})) \wedge \forall k \in \mathbb{N}.\; (f(k) \leq g(k)))) \Rightarrow  \\ \quad \quad \operatorname{esum}(f) \leq \operatorname{esum}(g) \end{array}}}
\vnbprop{esum-nn-two-support}{\vnbstmt{\begin{array}{@{}l@{}} \forall h \in (\mathbb{N} \to \operatorname{rr\text{-}pos\text{-}star}). \\ \quad \forall q \in \mathbb{N}.\; (h(((q + 1) + 1)) = 0) \Rightarrow  \\ \quad \quad \operatorname{esum}(h) = \operatorname{eplus}(h(0), h(1)) \end{array}}}
\vnbprop{measure-finitely-additive}{\vnbstmt{\begin{array}{@{}l@{}} \forall a, b, u, \alpha, \beta. \\ \quad (\operatorname{is\text{-}measure}(a, b, u) \wedge ((\alpha \in b) \wedge ((\beta \in b) \wedge ((\alpha \cap \beta) = \emptyset)))) \Rightarrow  \\ \quad \quad u((\alpha \cup \beta)) = \operatorname{eplus}(u(\alpha), u(\beta)) \end{array}}}
\vnbprop{measure-monotone}{\vnbstmt{\begin{array}{@{}l@{}} \forall a, b, u, \alpha, \beta. \\ \quad (\operatorname{is\text{-}measure}(a, b, u) \wedge ((\alpha \in b) \wedge ((\beta \in b) \wedge (\alpha \subseteq \beta)))) \Rightarrow  \\ \quad \quad u(\alpha) \leq u(\beta) \end{array}}}
\vnbprop{bu-seg-zero}{\vnbstmt{\forall f.\; (\operatorname{big\text{-}union}(j, \mathrm{OS}(0), f(j)) = \emptyset)}}
\vnbprop{bu-seg-succ}{\vnbstmt{\begin{array}{@{}l@{}} \forall k \in \mathbb{N}, c.; \forall f \in (\mathbb{N} \to c). \\ \quad \operatorname{big\text{-}union}(j, \mathrm{OS}((k + 1)), f(j)) = (\operatorname{big\text{-}union}(j, \mathrm{OS}(k), f(j)) \cup f(k)) \end{array}}}
\vnbprop{sigma-algebra-seg-union-in}{\vnbstmt{\begin{array}{@{}l@{}} \forall k \in \mathbb{N}, a, b, f. \\ \quad (\operatorname{is\text{-}sigma\text{-}algebra}(a, b) \wedge (f \in (\mathbb{N} \to b))) \Rightarrow  \\ \quad \quad \operatorname{big\text{-}union}(j, \mathrm{OS}(k), f(j)) \in b \end{array}}}
\vnbprop{disjointify-type}{\vnbstmt{\begin{array}{@{}l@{}} \forall a, b, f. \\ \quad (\operatorname{is\text{-}sigma\text{-}algebra}(a, b) \wedge (f \in (\mathbb{N} \to b))) \Rightarrow  \\ \quad \quad (\lambda k \in \mathbb{N}.\; \operatorname{complement\text{-}in}(f(k), \operatorname{big\text{-}union}(j, \mathrm{OS}(k), f(j)))) \in (\mathbb{N} \to b) \end{array}}}
\vnbprop{disjointify-disjoint}{\vnbstmt{\begin{array}{@{}l@{}} \forall f, c, i, j. \\ \quad ((f \in (\mathbb{N} \to c)) \wedge ((i \in \mathbb{N}) \wedge ((j \in \mathbb{N}) \wedge \neg (i = j)))) \Rightarrow  \\ \quad \quad ((\lambda k \in \mathbb{N}.\; \operatorname{complement\text{-}in}(f(k), \operatorname{big\text{-}union}(\mathit{uj}, \mathrm{OS}(k), f(\mathit{uj}))))(i) \cap (\lambda k \in \mathbb{N}.\; \operatorname{complement\text{-}in}(f(k), \operatorname{big\text{-}union}(\mathit{uj}, \mathrm{OS}(k), f(\mathit{uj}))))(j)) \\ \quad \quad =\; \emptyset \end{array}}}
\vnbprop{disjointify-union}{\vnbstmt{\begin{array}{@{}l@{}} \forall f, c. \\ \quad (f \in (\mathbb{N} \to c)) \Rightarrow  \\ \quad \quad \operatorname{big\text{-}union}(n, \mathbb{N}, (\lambda k \in \mathbb{N}.\; \operatorname{complement\text{-}in}(f(k), \operatorname{big\text{-}union}(j, \mathrm{OS}(k), f(j))))(n)) \\ \quad \quad =\; \operatorname{big\text{-}union}(n, \mathbb{N}, f(n)) \end{array}}}
\vnbprop{measure-countably-subadditive}{\vnbstmt{\begin{array}{@{}l@{}} \forall a, b, u, f. \\ \quad (\operatorname{is\text{-}measure}(a, b, u) \wedge (f \in (\mathbb{N} \to b))) \Rightarrow  \\ \quad \quad u(\operatorname{big\text{-}union}(n, \mathbb{N}, f(n))) \leq \operatorname{esum}((\lambda n \in \mathbb{N}.\; u(f(n)))) \end{array}}}
\vnbprop{disjointify-prefix}{\vnbstmt{\begin{array}{@{}l@{}} \forall n \in \mathbb{N}, a, b, u, f. \\ \quad (\operatorname{is\text{-}measure}(a, b, u) \wedge (f \in (\mathbb{N} \to b))) \Rightarrow  \\ \quad \quad \operatorname{finsum}(\operatorname{rr\text{-}pos\text{-}star\text{-}add\text{-}monoid}, \\ \quad \quad \quad (\lambda c \in \mathbb{N}.\; \\ \quad \quad \quad \quad u((\lambda k \in \mathbb{N}.\; \operatorname{complement\text{-}in}(f(k), \operatorname{big\text{-}union}(j, \mathrm{OS}(k), f(j))))(c))), \\ \quad \quad \quad \mathrm{OS}(n)) \\ \quad \quad =\; u(\operatorname{big\text{-}union}(j, \mathrm{OS}(n), f(j))) \end{array}}}
\vnbprop{nondecreasing-seg-subset}{\vnbstmt{\begin{array}{@{}l@{}} \forall n \in \mathbb{N}, c.; \forall f \in (\mathbb{N} \to c), j \in \mathrm{OS}(n). \\ \quad \forall a \in \mathbb{N}.\; (f(a) \subseteq f((a + 1))) \Rightarrow (f(j) \subseteq f(n)) \end{array}}}
\vnbprop{measure-continuity-from-below}{\vnbstmt{\begin{array}{@{}l@{}} \forall a, b, u, f. \\ \quad (\operatorname{is\text{-}measure}(a, b, u) \wedge ((f \in (\mathbb{N} \to b)) \wedge \forall n \in \mathbb{N}.\; (f(n) \subseteq f((n + 1))))) \Rightarrow  \\ \quad \quad u(\operatorname{big\text{-}union}(n, \mathbb{N}, f(n))) = \operatorname{esup}(\operatorname{image}((\lambda n \in \mathbb{N}.\; u(f(n))), \mathbb{N})) \end{array}}}
\vnbprop{measure-finite-subset}{\vnbstmt{\begin{array}{@{}l@{}} \forall a, b, u, \alpha, \beta. \\ \quad (\operatorname{is\text{-}measure}(a, b, u) \wedge ((\alpha \in b) \wedge ((\beta \in b) \wedge ((\alpha \subseteq \beta) \wedge (u(\beta) \in \mathbb{R}))))) \Rightarrow  \\ \quad \quad u(\alpha) \in \mathbb{R} \end{array}}}
\vnbprop{measure-subtractive}{\vnbstmt{\begin{array}{@{}l@{}} \forall a, b, u, \alpha, \beta. \\ \quad (\operatorname{is\text{-}measure}(a, b, u) \wedge ((\alpha \in b) \wedge ((\beta \in b) \wedge ((\alpha \subseteq \beta) \wedge (u(\beta) \in \mathbb{R}))))) \Rightarrow  \\ \quad \quad u(\operatorname{complement\text{-}in}(\beta, \alpha)) = (u(\beta) - u(\alpha)) \end{array}}}
\vnbprop{measure-continuity-from-above}{\vnbstmt{\begin{array}{@{}l@{}} \forall a, b, u, f. \\ \quad \operatorname{is\text{-}measure}(a, b, u) \wedge \\ \quad f \in (\mathbb{N} \to b) \wedge \\ \quad \forall n \in \mathbb{N}.\; (f((n + 1)) \subseteq f(n)) \wedge \\ \quad u(f(0)) \in \mathbb{R} \Rightarrow \\ \quad \quad u(\operatorname{complement\text{-}in}(a, \operatorname{big\text{-}union}(n, \mathbb{N}, \operatorname{complement\text{-}in}(a, f(n))))) \\ \quad \quad =\; \operatorname{einf}(\operatorname{image}((\lambda n \in \mathbb{N}.\; u(f(n))), \mathbb{N})) \end{array}}}
\vnbprop{measure-pushforward}{\vnbstmt{\begin{array}{@{}l@{}} \forall a, b, u, a_{2}, c, f. \\ \quad (\operatorname{is\text{-}measure}(a, b, u) \wedge \operatorname{is\text{-}measurable\text{-}map}(a, b, a_{2}, c, f)) \Rightarrow  \\ \quad \quad \operatorname{is\text{-}measure}(a_{2}, c, (\lambda \mathit{pb} \in c.\; u(\{x \in a : (f(x) \in \mathit{pb})\}))) \end{array}}}

\vnbbookfile{theorem-library/integral-laws}{theorem-library/integral-laws}
\noindent{\small integral-laws.scm -- THE INTEGRAL OF A [0,+inf]-VALUED FUNCTION BY SIMPLE FUNCTIONS: the laws, proved from the definition, and the bricks they rest on. Agent M-2 of the October roadmap (week 2, scratchpad/triage/WEEK2.md; decision 4 of docs/roadmap-2026-10-03.md).  Helper prefixes ibk- (the}\par
\vnbprop{etimes-closed}{\vnbstmt{\begin{array}{@{}l@{}} \forall x, y. \\ \quad ((x \in \operatorname{rr\text{-}pos\text{-}star}) \wedge (y \in \operatorname{rr\text{-}pos\text{-}star})) \Rightarrow (\operatorname{etimes}(x, y) \in \operatorname{rr\text{-}pos\text{-}star}) \end{array}}}
\vnbprop{etimes-zero-left}{\vnbstmt{\forall x \in \operatorname{rr\text{-}pos\text{-}star}.\; (\operatorname{etimes}(0, x) = 0)}}
\vnbprop{etimes-zero-right}{\vnbstmt{\forall x \in \operatorname{rr\text{-}pos\text{-}star}.\; (\operatorname{etimes}(x, 0) = 0)}}
\vnbprop{etimes-one-left}{\vnbstmt{\forall x \in \operatorname{rr\text{-}pos\text{-}star}.\; (\operatorname{etimes}(1, x) = x)}}
\vnbprop{etimes-comm}{\vnbstmt{\begin{array}{@{}l@{}} \forall x, y. \\ \quad ((x \in \operatorname{rr\text{-}pos\text{-}star}) \wedge (y \in \operatorname{rr\text{-}pos\text{-}star})) \Rightarrow (\operatorname{etimes}(x, y) = \operatorname{etimes}(y, x)) \end{array}}}
\vnbprop{etimes-distrib-left}{\vnbstmt{\begin{array}{@{}l@{}} \forall c \in \mathbb{R}, x, y. \\ \quad ((0 \leq c) \wedge ((x \in \operatorname{rr\text{-}pos\text{-}star}) \wedge (y \in \operatorname{rr\text{-}pos\text{-}star}))) \Rightarrow  \\ \quad \quad \operatorname{etimes}(c, \operatorname{eplus}(x, y)) = \operatorname{eplus}(\operatorname{etimes}(c, x), \operatorname{etimes}(c, y)) \end{array}}}
\vnbprop{etimes-mono-left}{\vnbstmt{\begin{array}{@{}l@{}} \forall c \in \mathbb{R}, x, y. \\ \quad ((0 \leq c) \wedge (((x \in \operatorname{rr\text{-}pos\text{-}star}) \wedge (y \in \operatorname{rr\text{-}pos\text{-}star})) \wedge (x \leq y))) \Rightarrow  \\ \quad \quad \operatorname{etimes}(c, x) \leq \operatorname{etimes}(c, y) \end{array}}}
\vnbprop{etimes-assoc-real}{\vnbstmt{\begin{array}{@{}l@{}} \forall c, d \in \mathbb{R}, x \in \operatorname{rr\text{-}pos\text{-}star}. \\ \quad ((0 \leq c) \wedge (0 \leq d)) \Rightarrow (\operatorname{etimes}(c, \operatorname{etimes}(d, x)) = \operatorname{etimes}((c \cdot d), x)) \end{array}}}
\vnbprop{etimes-distrib-right-real}{\vnbstmt{\begin{array}{@{}l@{}} \forall c, d \in \mathbb{R}, x \in \operatorname{rr\text{-}pos\text{-}star}. \\ \quad ((0 \leq c) \wedge (0 \leq d)) \Rightarrow  \\ \quad \quad \operatorname{etimes}(\operatorname{eplus}(c, d), x) = \operatorname{eplus}(\operatorname{etimes}(c, x), \operatorname{etimes}(d, x)) \end{array}}}
\vnbprop{etimes-mono-real-left}{\vnbstmt{\begin{array}{@{}l@{}} \forall a, b.; \forall m \in \operatorname{rr\text{-}pos\text{-}star}. \\ \quad ((a \in \mathbb{R}) \wedge ((b \in \mathbb{R}) \wedge ((0 \leq a) \wedge (a \leq b)))) \Rightarrow (\operatorname{etimes}(a, m) \leq \operatorname{etimes}(b, m)) \end{array}}}
\vnbprop{finsum-cm-remove-ptwise}{\vnbstmt{\begin{array}{@{}l@{}} \forall m.; \forall t \in \mathbf{Set}, k \in t, g. \\ \quad (\operatorname{is\text{-}comm\text{-}monoid}(m) \wedge ((|t| \in \mathbb{N}) \wedge \forall z \in t.\; (g(z) \in \operatorname{carr}(m)))) \Rightarrow  \\ \quad \quad \operatorname{finsum}(m, g, t) = \operatorname{opr}(m)(\operatorname{finsum}(m, g, \operatorname{difference}(t, \{k\})), g(k)) \end{array}}}
\vnbprop{sep-fiber-insert-other}{\vnbstmt{\begin{array}{@{}l@{}} \forall s, x, i, z. \\ \quad ((x \in \mathbf{Set}) \wedge \neg (i(x) = z)) \Rightarrow  \\ \quad \quad \{w \in (s \cup \{x\}) : (i(w) = z)\} = \{w \in s : (i(w) = z)\} \end{array}}}
\vnbprop{sep-fiber-insert-same}{\vnbstmt{\begin{array}{@{}l@{}} \forall s, x, i. \\ \quad ((x \in \mathbf{Set}) \wedge (i(x) \in \mathbf{Set})) \Rightarrow  \\ \quad \quad \{w \in (s \cup \{x\}) : (i(w) = i(x))\} = (\{w \in s : (i(w) = i(x))\} \cup \{x\}) \end{array}}}
\vnbprop{finsum-cm-fiber-ptwise}{\vnbstmt{\begin{array}{@{}l@{}} \forall m.; \forall s, t \in \mathbf{Set}, i, f, q. \\ \quad \operatorname{is\text{-}comm\text{-}monoid}(m) \wedge \\ \quad |s| \in \mathbb{N} \wedge \\ \quad |t| \in \mathbb{N} \wedge \\ \quad \forall z \in s.\; (i(z) \in t) \wedge \\ \quad \forall z \in s.\; (f(z) \in \operatorname{carr}(m)) \wedge \\ \quad \forall z \in t.\; (q(z) \simeq \operatorname{finsum}(m, f, \{w \in s : (i(w) = z)\})) \Rightarrow \\ \quad \quad \operatorname{finsum}(m, f, s) = \operatorname{finsum}(m, q, t) \end{array}}}
\vnbprop{rps-finsum-scale-ptwise}{\vnbstmt{\begin{array}{@{}l@{}} \forall c \in \mathbb{R}, s \in \mathbf{Set}, f, h. \\ \quad 0 \leq c \wedge \\ \quad |s| \in \mathbb{N} \wedge \\ \quad \forall z \in s.\; (f(z) \in \operatorname{rr\text{-}pos\text{-}star}) \wedge \\ \quad \forall z \in s.\; (h(z) = \operatorname{etimes}(c, f(z))) \Rightarrow \\ \quad \quad \operatorname{finsum}(\operatorname{rr\text{-}pos\text{-}star\text{-}add\text{-}monoid}, h, s) = \operatorname{etimes}(c, \operatorname{finsum}(\operatorname{rr\text{-}pos\text{-}star\text{-}add\text{-}monoid}, f, s)) \end{array}}}
\vnbprop{sigma-algebra-preimage-finite-in}{\vnbstmt{\begin{array}{@{}l@{}} \forall a, b.; \forall s \in \mathbf{Set}, k. \\ \quad (\operatorname{is\text{-}sigma\text{-}algebra}(a, b) \wedge ((|s| \in \mathbb{N}) \wedge \forall z \in s.\; (\{x \in a : (k(x) = z)\} \in b))) \Rightarrow  \\ \quad \quad \{x \in a : (k(x) \in s)\} \in b \end{array}}}
\vnbprop{measure-preimage-finite-sum}{\vnbstmt{\begin{array}{@{}l@{}} \forall a, b, u.; \forall s \in \mathbf{Set}, k, g. \\ \quad \operatorname{is\text{-}measure}(a, b, u) \wedge \\ \quad |s| \in \mathbb{N} \wedge \\ \quad \forall z \in s.\; (\{x \in a : (k(x) = z)\} \in b) \wedge \\ \quad \forall z \in s.\; (g(z) = u(\{x \in a : (k(x) = z)\})) \Rightarrow \\ \quad \quad u(\{x \in a : (k(x) \in s)\}) = \operatorname{finsum}(\operatorname{rr\text{-}pos\text{-}star\text{-}add\text{-}monoid}, g, s) \end{array}}}
\vnbprop{measurable-fn-ge}{\vnbstmt{\begin{array}{@{}l@{}} \forall b, c, f.; \forall a \in \mathbb{R}. \\ \quad \operatorname{is\text{-}measurable\text{-}fn}(b, c, f) \Rightarrow (\{x \in b : (a \leq f(x))\} \in c) \end{array}}}
\vnbprop{measurable-fn-level-set}{\vnbstmt{\begin{array}{@{}l@{}} \forall b, c, f.; \forall a \in \mathbb{R}. \\ \quad \operatorname{is\text{-}measurable\text{-}fn}(b, c, f) \Rightarrow (\{x \in b : (f(x) = a)\} \in c) \end{array}}}
\vnbprop{simple-integral-unfold}{\vnbstmt{\begin{array}{@{}l@{}} \forall a, u, s. \\ \quad (\operatorname{simple\text{-}integral}(a, u, s) \simeq \operatorname{finsum}(\operatorname{rr\text{-}pos\text{-}star\text{-}add\text{-}monoid}, (\lambda y \in \operatorname{image}(s, a).\; \operatorname{etimes}(y, u(\{x \in a : (s(x) = y)\}))), \operatorname{image}(s, a))) \end{array}}}
\vnbprop{integral-unfold}{\vnbstmt{\begin{array}{@{}l@{}} \forall a, b, u, f. \\ \quad (\operatorname{integral}(a, b, u, f) \simeq \operatorname{esup}(\{y \in \operatorname{rr\text{-}pos\text{-}star} : \exists s \in (a \to \operatorname{rr\text{-}pos\text{-}star}).\; (\operatorname{is\text{-}simple\text{-}fn}(a, b, s) \wedge (\operatorname{ptwise\text{-}le}(a, s, f) \wedge (y = \operatorname{simple\text{-}integral}(a, u, s))))\})) \end{array}}}
\vnbprop{integral-in}{\vnbstmt{\begin{array}{@{}l@{}} \forall a, b, u, f. \\ \quad (\operatorname{is\text{-}measure}(a, b, u) \wedge \operatorname{is\text{-}measurable\text{-}fn}(a, b, f)) \Rightarrow  \\ \quad \quad \operatorname{integral}(a, b, u, f) \in \operatorname{rr\text{-}pos\text{-}star} \end{array}}}
\vnbprop{integral-monotone}{\vnbstmt{\begin{array}{@{}l@{}} \forall a, b, u, f, g. \\ \quad \operatorname{is\text{-}measure}(a, b, u) \wedge \\ \quad \operatorname{is\text{-}measurable\text{-}fn}(a, b, f) \wedge \\ \quad \operatorname{is\text{-}measurable\text{-}fn}(a, b, g) \wedge \\ \quad \operatorname{ptwise\text{-}le}(a, f, g) \Rightarrow \\ \quad \quad \operatorname{integral}(a, b, u, f) \leq \operatorname{integral}(a, b, u, g) \end{array}}}
\vnbprop{simple-integral-partition}{\vnbstmt{\begin{array}{@{}l@{}} \forall a, b, u, s.; \forall c \in \mathbf{Set}, k, h, g. \\ \quad \operatorname{is\text{-}measure}(a, b, u) \wedge \\ \quad |c| \in \mathbb{N} \wedge \\ \quad \forall z \in a.\; (k(z) \in c) \wedge \\ \quad \forall z \in c.\; (\{x \in a : (k(x) = z)\} \in b) \wedge \\ \quad \forall z \in c.\; (h(z) \in \mathbb{R}) \wedge \\ \quad \forall z \in c.\; (0 \leq h(z)) \wedge \\ \quad \forall z \in a.\; (s(z) = h(k(z))) \wedge \\ \quad \forall z \in c.\; (g(z) = \operatorname{etimes}(h(z), u(\{x \in a : (k(x) = z)\}))) \Rightarrow \\ \quad \quad \operatorname{simple\text{-}integral}(a, u, s) = \operatorname{finsum}(\operatorname{rr\text{-}pos\text{-}star\text{-}add\text{-}monoid}, g, c) \end{array}}}
\vnbprop{simple-integral-in}{\vnbstmt{\begin{array}{@{}l@{}} \forall a, b, u, s. \\ \quad (\operatorname{is\text{-}measure}(a, b, u) \wedge \operatorname{is\text{-}simple\text{-}fn}(a, b, s)) \Rightarrow  \\ \quad \quad \operatorname{simple\text{-}integral}(a, u, s) \in \operatorname{rr\text{-}pos\text{-}star} \end{array}}}
\vnbprop{simple-integral-mono}{\vnbstmt{\begin{array}{@{}l@{}} \forall a, b, u, s, t. \\ \quad (\operatorname{is\text{-}measure}(a, b, u) \wedge (\operatorname{is\text{-}simple\text{-}fn}(a, b, s) \wedge (\operatorname{is\text{-}simple\text{-}fn}(a, b, t) \wedge \operatorname{ptwise\text{-}le}(a, s, t)))) \Rightarrow  \\ \quad \quad \operatorname{simple\text{-}integral}(a, u, s) \leq \operatorname{simple\text{-}integral}(a, u, t) \end{array}}}
\vnbprop{integral-simple}{\vnbstmt{\begin{array}{@{}l@{}} \forall a, b, u, s. \\ \quad (\operatorname{is\text{-}measure}(a, b, u) \wedge \operatorname{is\text{-}simple\text{-}fn}(a, b, s)) \Rightarrow  \\ \quad \quad \operatorname{integral}(a, b, u, s) = \operatorname{simple\text{-}integral}(a, u, s) \end{array}}}
\vnbprop{integral-sup-of-simple}{\vnbstmt{\begin{array}{@{}l@{}} \forall a, b, u, f. \\ \quad (\operatorname{is\text{-}measure}(a, b, u) \wedge \operatorname{is\text{-}measurable\text{-}fn}(a, b, f)) \Rightarrow  \\ \quad \quad \operatorname{integral}(a, b, u, f) \\ \quad \quad =\; \operatorname{esup}(\{y \in \operatorname{rr\text{-}pos\text{-}star} : \exists s \in (a \to \operatorname{rr\text{-}pos\text{-}star}).\; (\operatorname{is\text{-}simple\text{-}fn}(a, b, s) \wedge (\operatorname{ptwise\text{-}le}(a, s, f) \wedge (y = \operatorname{integral}(a, b, u, s))))\}) \end{array}}}
\vnbprop{indicator-value-in}{\vnbstmt{\forall a, \alpha.; \forall x \in a.\; ((x \in \alpha) \Rightarrow (\operatorname{indicator}(a, \alpha)(x) = 1))}}
\vnbprop{indicator-value-out}{\vnbstmt{\forall a, \alpha.; \forall x \in a.\; (\neg (x \in \alpha) \Rightarrow (\operatorname{indicator}(a, \alpha)(x) = 0))}}
\vnbprop{indicator-simple}{\vnbstmt{\begin{array}{@{}l@{}} \forall a, b, \alpha. \\ \quad (\operatorname{is\text{-}sigma\text{-}algebra}(a, b) \wedge (\alpha \in b)) \Rightarrow  \\ \quad \quad \operatorname{is\text{-}simple\text{-}fn}(a, b, \operatorname{indicator}(a, \alpha)) \end{array}}}
\vnbprop{simple-integral-indicator}{\vnbstmt{\begin{array}{@{}l@{}} \forall a, b, u, \alpha. \\ \quad (\operatorname{is\text{-}measure}(a, b, u) \wedge (\alpha \in b)) \Rightarrow  \\ \quad \quad \operatorname{simple\text{-}integral}(a, u, \operatorname{indicator}(a, \alpha)) = u(\alpha) \end{array}}}
\vnbprop{integral-indicator}{\vnbstmt{\begin{array}{@{}l@{}} \forall a, b, u, \alpha. \\ \quad (\operatorname{is\text{-}measure}(a, b, u) \wedge (\alpha \in b)) \Rightarrow  \\ \quad \quad \operatorname{integral}(a, b, u, \operatorname{indicator}(a, \alpha)) = u(\alpha) \end{array}}}
\vnbprop{integral-infinite-on-null}{\vnbstmt{\begin{array}{@{}l@{}} \forall a, b, u, \alpha. \\ \quad (\operatorname{is\text{-}measure}(a, b, u) \wedge ((\alpha \in b) \wedge (u(\alpha) = 0))) \Rightarrow  \\ \quad \quad \operatorname{integral}(a, b, u, \operatorname{ptwise\text{-}scale}(a, \operatorname{pos\text{-}inf}, \operatorname{indicator}(a, \alpha))) = 0 \end{array}}}
\vnbprop{simple-integral-additive}{\vnbstmt{\begin{array}{@{}l@{}} \forall a, b, u, s, t. \\ \quad (\operatorname{is\text{-}measure}(a, b, u) \wedge (\operatorname{is\text{-}simple\text{-}fn}(a, b, s) \wedge \operatorname{is\text{-}simple\text{-}fn}(a, b, t))) \Rightarrow  \\ \quad \quad \operatorname{simple\text{-}integral}(a, u, \operatorname{ptwise\text{-}eplus}(a, s, t)) \\ \quad \quad =\; \operatorname{eplus}(\operatorname{simple\text{-}integral}(a, u, s), \operatorname{simple\text{-}integral}(a, u, t)) \end{array}}}
\vnbprop{simple-integral-homogeneous}{\vnbstmt{\begin{array}{@{}l@{}} \forall a, b, u, s, c. \\ \quad (\operatorname{is\text{-}measure}(a, b, u) \wedge (\operatorname{is\text{-}simple\text{-}fn}(a, b, s) \wedge ((c \in \mathbb{R}) \wedge (0 \leq c)))) \Rightarrow  \\ \quad \quad \operatorname{simple\text{-}integral}(a, u, \operatorname{ptwise\text{-}scale}(a, c, s)) = \operatorname{etimes}(c, \operatorname{simple\text{-}integral}(a, u, s)) \end{array}}}
\vnbprop{measurable-fn-finite-range}{\vnbstmt{\begin{array}{@{}l@{}} \forall a, b, u. \\ \quad \operatorname{is\text{-}sigma\text{-}algebra}(a, b) \wedge \\ \quad u \in (a \to \operatorname{rr\text{-}pos\text{-}star}) \wedge \\ \quad |\operatorname{image}(u, a)| \in \mathbb{N} \wedge \\ \quad \forall z \in \operatorname{image}(u, a).\; (\{x \in a : (u(x) = z)\} \in b) \Rightarrow \\ \quad \quad \operatorname{is\text{-}measurable\text{-}fn}(a, b, u) \end{array}}}
\vnbprop{simple-fn-scale}{\vnbstmt{\begin{array}{@{}l@{}} \forall a, b, s, c. \\ \quad (\operatorname{is\text{-}sigma\text{-}algebra}(a, b) \wedge (\operatorname{is\text{-}simple\text{-}fn}(a, b, s) \wedge ((c \in \mathbb{R}) \wedge (0 \leq c)))) \Rightarrow  \\ \quad \quad \operatorname{is\text{-}simple\text{-}fn}(a, b, \operatorname{ptwise\text{-}scale}(a, c, s)) \end{array}}}
\vnbprop{integral-homogeneous}{\vnbstmt{\begin{array}{@{}l@{}} \forall a, b, u, f, c. \\ \quad (\operatorname{is\text{-}measure}(a, b, u) \wedge (\operatorname{is\text{-}measurable\text{-}fn}(a, b, f) \wedge ((c \in \mathbb{R}) \wedge (0 \leq c)))) \Rightarrow  \\ \quad \quad \operatorname{integral}(a, b, u, \operatorname{ptwise\text{-}scale}(a, c, f)) = \operatorname{etimes}(c, \operatorname{integral}(a, b, u, f)) \end{array}}}

\vnbbookfile{theorem-library/esum-scale}{theorem-library/esum-scale}
\noindent{\small ==================================================================== theorem-library/esum-scale.scm -- MEASURE-5-A, 2026-10-10. Homogeneity of ESUM, and a scaled measure is a measure.}\par
\vnbprop{esum-scale-ptwise}{\vnbstmt{\begin{array}{@{}l@{}} \forall c \in \mathbb{R}, d, g, h. \\ \quad 0 \leq c \wedge \\ \quad g \in (d \to \operatorname{rr\text{-}pos\text{-}star}) \wedge \\ \quad h \in (d \to \operatorname{rr\text{-}pos\text{-}star}) \wedge \\ \quad \forall z \in d.\; (h(z) = \operatorname{etimes}(c, g(z))) \Rightarrow \\ \quad \quad \operatorname{esum}(h) = \operatorname{etimes}(c, \operatorname{esum}(g)) \end{array}}}
\vnbprop{esum-scale}{\vnbstmt{\begin{array}{@{}l@{}} \forall c \in \mathbb{R}, d.; \forall f \in (d \to \operatorname{rr\text{-}pos\text{-}star}). \\ \quad (0 \leq c) \Rightarrow (\operatorname{esum}((\lambda z \in d.\; \operatorname{etimes}(c, f(z)))) = \operatorname{etimes}(c, \operatorname{esum}(f))) \end{array}}}
\vnbprop{measure-scale}{\vnbstmt{\begin{array}{@{}l@{}} \forall a, b, u.; \forall c \in \mathbb{R}. \\ \quad (\operatorname{is\text{-}measure}(a, b, u) \wedge (0 \leq c)) \Rightarrow  \\ \quad \quad \operatorname{is\text{-}measure}(a, b, (\lambda x \in b.\; \operatorname{etimes}(c, u(x)))) \end{array}}}

\vnbbookfile{theorem-library/signed-measure-ext}{theorem-library/signed-measure-ext}
\noindent{\small signed-measure-ext.scm -- WAVE 2 (2026-10-27..30), wired 2026-10-30 from scratchpad/mf/mf-ext.scm, scratchpad/mf/mf-psum.scm, scratchpad/mf/mf-conv.scm. Library: measure.}\par
\vnbprop{ext-real-not-pos-inf}{\vnbstmt{\forall x \in \mathbb{R}.\; \neg (x = \operatorname{pos\text{-}inf})}}
\vnbprop{ext-real-not-neg-inf}{\vnbstmt{\forall x \in \mathbb{R}.\; \neg (x = \operatorname{neg\text{-}inf})}}
\vnbprop{ext-neg-real}{\vnbstmt{\forall x \in \mathbb{R}.\; (\operatorname{ext\text{-}neg}(x) = -x)}}
\vnbprop{ext-neg-pos-inf}{\vnbstmt{\operatorname{ext\text{-}neg}(\operatorname{pos\text{-}inf}) = \operatorname{neg\text{-}inf}}}
\vnbprop{ext-neg-neg-inf}{\vnbstmt{\operatorname{ext\text{-}neg}(\operatorname{neg\text{-}inf}) = \operatorname{pos\text{-}inf}}}
\vnbprop{ext-add-real}{\vnbstmt{\forall x, y.\; (((x \in \mathbb{R}) \wedge (y \in \mathbb{R})) \Rightarrow (\operatorname{ext\text{-}add}(x, y) = (x + y)))}}
\vnbprop{ext-add-pos-inf-left}{\vnbstmt{\forall y.\; (\operatorname{ext\text{-}add}(\operatorname{pos\text{-}inf}, y) = \operatorname{pos\text{-}inf})}}
\vnbprop{ext-add-pos-inf-right}{\vnbstmt{\forall y.\; (\operatorname{ext\text{-}add}(y, \operatorname{pos\text{-}inf}) = \operatorname{pos\text{-}inf})}}
\vnbprop{ext-add-neg-inf-left}{\vnbstmt{\forall y.\; (\neg (y = \operatorname{pos\text{-}inf}) \Rightarrow (\operatorname{ext\text{-}add}(\operatorname{neg\text{-}inf}, y) = \operatorname{neg\text{-}inf}))}}
\vnbprop{ext-add-neg-inf-right}{\vnbstmt{\forall y.\; (\neg (y = \operatorname{pos\text{-}inf}) \Rightarrow (\operatorname{ext\text{-}add}(y, \operatorname{neg\text{-}inf}) = \operatorname{neg\text{-}inf}))}}
\vnbprop{ext-add-zero-left}{\vnbstmt{\forall y \in \operatorname{rr\text{-}star}.\; (\operatorname{ext\text{-}add}(0, y) = y)}}
\vnbprop{ext-add-zero-right}{\vnbstmt{\forall y \in \operatorname{rr\text{-}star}.\; (\operatorname{ext\text{-}add}(y, 0) = y)}}
\vnbprop{ext-psum-pos}{\vnbstmt{\begin{array}{@{}l@{}} \forall x, n. \\ \quad \exists k \in \mathbb{N}.\; ((k < n) \wedge (x(k) = \operatorname{pos\text{-}inf})) \Rightarrow  \\ \quad \quad \operatorname{ext\text{-}partial\text{-}sum}(x, n) = \operatorname{pos\text{-}inf} \end{array}}}
\vnbprop{ext-psum-neg}{\vnbstmt{\begin{array}{@{}l@{}} \forall x, n. \\ \quad (\neg \exists k \in \mathbb{N}.\; ((k < n) \wedge (x(k) = \operatorname{pos\text{-}inf})) \wedge \exists k \in \mathbb{N}.\; ((k < n) \wedge (x(k) = \operatorname{neg\text{-}inf}))) \Rightarrow  \\ \quad \quad \operatorname{ext\text{-}partial\text{-}sum}(x, n) = \operatorname{neg\text{-}inf} \end{array}}}
\vnbprop{ext-psum-fin}{\vnbstmt{\begin{array}{@{}l@{}} \forall x, n. \\ \quad \neg \exists k \in \mathbb{N}.\; ((k < n) \wedge (x(k) = \operatorname{pos\text{-}inf})) \wedge \\ \quad \neg \exists k \in \mathbb{N}.\; ((k < n) \wedge (x(k) = \operatorname{neg\text{-}inf})) \Rightarrow \\ \quad \quad (\operatorname{ext\text{-}partial\text{-}sum}(x, n) \simeq \operatorname{series\text{-}partial\text{-}sum}(x, n)) \end{array}}}
\vnbprop{ext-psum-zero}{\vnbstmt{\forall x.\; (\operatorname{ext\text{-}partial\text{-}sum}(x, 0) = 0)}}
\vnbprop{ext-psum-cond-succ}{\vnbstmt{\begin{array}{@{}l@{}} \forall x.; \forall n \in \mathbb{N}, c. \\ \quad \exists k \in \mathbb{N}.\; ((k < (n + 1)) \wedge (x(k) = c)) \Leftrightarrow (\exists k \in \mathbb{N}.\; ((k < n) \wedge (x(k) = c)) \vee (x(n) = c)) \end{array}}}
\vnbprop{ext-psum-real-terms}{\vnbstmt{\begin{array}{@{}l@{}} \forall x.; \forall n, j \in \mathbb{N}. \\ \quad \forall k \in \mathbb{N}.\; (x(k) \in \operatorname{rr\text{-}star}) \wedge \\ \quad \neg \exists k \in \mathbb{N}.\; ((k < n) \wedge (x(k) = \operatorname{pos\text{-}inf})) \wedge \\ \quad \neg \exists k \in \mathbb{N}.\; ((k < n) \wedge (x(k) = \operatorname{neg\text{-}inf})) \wedge \\ \quad j < n \Rightarrow \\ \quad \quad x(j) \in \mathbb{R} \end{array}}}
\vnbprop{ext-psum-in}{\vnbstmt{\begin{array}{@{}l@{}} \forall x.; \forall n \in \mathbb{N}. \\ \quad \forall k \in \mathbb{N}.\; (x(k) \in \operatorname{rr\text{-}star}) \Rightarrow (\operatorname{ext\text{-}partial\text{-}sum}(x, n) \in \operatorname{rr\text{-}star}) \end{array}}}
\vnbprop{ext-psum-succ}{\vnbstmt{\begin{array}{@{}l@{}} \forall x.; \forall n \in \mathbb{N}. \\ \quad \forall k \in \mathbb{N}.\; (x(k) \in \operatorname{rr\text{-}star}) \Rightarrow  \\ \quad \quad \operatorname{ext\text{-}partial\text{-}sum}(x, (n + 1)) = \operatorname{ext\text{-}add}(\operatorname{ext\text{-}partial\text{-}sum}(x, n), x(n)) \end{array}}}
\vnbprop{ext-series-limit-up}{\vnbstmt{\begin{array}{@{}l@{}} \forall x, l. \\ \quad \operatorname{ext\text{-}series\text{-}converges\text{-}to}(x, l) \wedge \\ \quad \forall b \in \mathbb{R}. \exists m \in \mathbb{N}. \forall n \in \mathbb{N}. \\ \quad \quad (m \leq n) \Rightarrow (b \leq \operatorname{ext\text{-}partial\text{-}sum}(x, n)) \Rightarrow \\ \quad \quad l = \operatorname{pos\text{-}inf} \end{array}}}
\vnbprop{ext-series-limit-down}{\vnbstmt{\begin{array}{@{}l@{}} \forall x, l. \\ \quad \operatorname{ext\text{-}series\text{-}converges\text{-}to}(x, l) \wedge \\ \quad \forall b \in \mathbb{R}. \exists m \in \mathbb{N}. \forall n \in \mathbb{N}. \\ \quad \quad (m \leq n) \Rightarrow (\operatorname{ext\text{-}partial\text{-}sum}(x, n) \leq b) \Rightarrow \\ \quad \quad l = \operatorname{neg\text{-}inf} \end{array}}}
\vnbprop{ext-series-limit-const}{\vnbstmt{\begin{array}{@{}l@{}} \forall x, l, c. \\ \quad \operatorname{ext\text{-}series\text{-}converges\text{-}to}(x, l) \wedge \\ \quad c \in \operatorname{rr\text{-}star} \wedge \\ \quad \exists m \in \mathbb{N}.\; \forall n \in \mathbb{N}.\; ((m \leq n) \Rightarrow (\operatorname{ext\text{-}partial\text{-}sum}(x, n) = c)) \Rightarrow \\ \quad \quad l = c \end{array}}}
\vnbprop{ext-series-limit-le}{\vnbstmt{\begin{array}{@{}l@{}} \forall x, l, b. \\ \quad (\operatorname{ext\text{-}series\text{-}converges\text{-}to}(x, l) \wedge ((b \in \operatorname{rr\text{-}pos\text{-}star}) \wedge \forall n \in \mathbb{N}.\; (\operatorname{ext\text{-}partial\text{-}sum}(x, n) \leq b))) \Rightarrow  \\ \quad \quad l \leq b \end{array}}}

\vnbbookfile{theorem-library/integral-convergence}{theorem-library/integral-convergence}
\noindent{\small integral-convergence.scm -- MEASURABLE FUNCTIONS CLOSED UNDER +, *, sup and lim inf; APPROXIMATION BY SIMPLE FUNCTIONS; MONOTONE CONVERGENCE; ADDITIVITY OF THE INTEGRAL; FATOU'S LEMMA; DOMINATED CONVERGENCE.  Agent M-3 of the October roadmap (week 3, scratchpad/triage/WEEK3.md).  Helper prefix icv-.  Every theorem ...}\par
\vnbprop{esup-gt-of-member}{\vnbstmt{\begin{array}{@{}l@{}} \forall s.; \forall y \in s, a \in \operatorname{rr\text{-}pos\text{-}star}. \\ \quad ((s \subseteq \operatorname{rr\text{-}pos\text{-}star}) \wedge (a < y)) \Rightarrow (a < \operatorname{esup}(s)) \end{array}}}
\vnbprop{measurable-fn-of-nonneg-tails}{\vnbstmt{\begin{array}{@{}l@{}} \forall a, b, f. \\ \quad \operatorname{is\text{-}sigma\text{-}algebra}(a, b) \wedge \\ \quad f \in (a \to \operatorname{rr\text{-}pos\text{-}star}) \wedge \\ \quad \forall l \in \mathbb{R}.\; ((0 \leq l) \Rightarrow (\{x \in a : (l < f(x))\} \in b)) \Rightarrow \\ \quad \quad \operatorname{is\text{-}measurable\text{-}fn}(a, b, f) \end{array}}}
\vnbprop{measurable-fn-sup-of-values}{\vnbstmt{\begin{array}{@{}l@{}} \forall a, b, f, g. \\ \quad \operatorname{is\text{-}sigma\text{-}algebra}(a, b) \wedge \\ \quad f \in (\mathbb{N} \to (a \to \operatorname{rr\text{-}pos\text{-}star})) \wedge \\ \quad \forall n \in \mathbb{N}.\; \operatorname{is\text{-}measurable\text{-}fn}(a, b, f(n)) \wedge \\ \quad g \in (a \to \operatorname{rr\text{-}pos\text{-}star}) \wedge \\ \quad \forall x \in a.\; (g(x) = \operatorname{esup}(\operatorname{image}((\lambda n \in \mathbb{N}.\; f(n)(x)), \mathbb{N}))) \Rightarrow \\ \quad \quad \operatorname{is\text{-}measurable\text{-}fn}(a, b, g) \end{array}}}
\vnbprop{measurable-fn-sup}{\vnbstmt{\begin{array}{@{}l@{}} \forall a, b, f. \\ \quad \operatorname{is\text{-}sigma\text{-}algebra}(a, b) \wedge \\ \quad f \in (\mathbb{N} \to (a \to \operatorname{rr\text{-}pos\text{-}star})) \wedge \\ \quad \forall n \in \mathbb{N}.\; \operatorname{is\text{-}measurable\text{-}fn}(a, b, f(n)) \Rightarrow \\ \quad \quad \operatorname{is\text{-}measurable\text{-}fn}(a, b, \operatorname{ptwise\text{-}sup}(a, f)) \end{array}}}
\vnbprop{measurable-fn-guarded-lt}{\vnbstmt{\begin{array}{@{}l@{}} \forall a, c, f, p, s.; \forall b \in \mathbb{R}. \\ \quad \operatorname{is\text{-}measurable\text{-}fn}(a, c, f) \Rightarrow (\{x \in a : ((p \in s) \wedge (f(x) < b))\} \in c) \end{array}}}
\vnbprop{measurable-fn-lower-bound-set}{\vnbstmt{\begin{array}{@{}l@{}} \forall a, c, f, s, b. \\ \quad \operatorname{is\text{-}sigma\text{-}algebra}(a, c) \wedge \\ \quad f \in (\mathbb{N} \to (a \to \operatorname{rr\text{-}pos\text{-}star})) \wedge \\ \quad \forall n \in \mathbb{N}.\; \operatorname{is\text{-}measurable\text{-}fn}(a, c, f(n)) \wedge \\ \quad s \subseteq \mathbb{N} \wedge \\ \quad b \in \mathbb{R} \Rightarrow \\ \quad \quad \{x \in a : \forall p \in s.\; (b \leq f(p)(x))\} \in c \end{array}}}
\vnbprop{measurable-fn-inf-of-values}{\vnbstmt{\begin{array}{@{}l@{}} \forall a, b, f, s, g. \\ \quad \operatorname{is\text{-}sigma\text{-}algebra}(a, b) \wedge \\ \quad f \in (\mathbb{N} \to (a \to \operatorname{rr\text{-}pos\text{-}star})) \wedge \\ \quad \forall n \in \mathbb{N}.\; \operatorname{is\text{-}measurable\text{-}fn}(a, b, f(n)) \wedge \\ \quad s \subseteq \mathbb{N} \wedge \\ \quad g \in (a \to \operatorname{rr\text{-}pos\text{-}star}) \wedge \\ \quad \forall x \in a.\; (g(x) = \operatorname{einf}(\operatorname{image}((\lambda n \in \mathbb{N}.\; f(n)(x)), s))) \Rightarrow \\ \quad \quad \operatorname{is\text{-}measurable\text{-}fn}(a, b, g) \end{array}}}
\vnbprop{measurable-fn-liminf}{\vnbstmt{\begin{array}{@{}l@{}} \forall a, b, f. \\ \quad \operatorname{is\text{-}sigma\text{-}algebra}(a, b) \wedge \\ \quad f \in (\mathbb{N} \to (a \to \operatorname{rr\text{-}pos\text{-}star})) \wedge \\ \quad \forall n \in \mathbb{N}.\; \operatorname{is\text{-}measurable\text{-}fn}(a, b, f(n)) \Rightarrow \\ \quad \quad \operatorname{is\text{-}measurable\text{-}fn}(a, b, \operatorname{ptwise\text{-}liminf}(a, f)) \end{array}}}
\vnbprop{rr-grid-between}{\vnbstmt{\begin{array}{@{}l@{}} \forall a, b \in \mathbb{R}. \\ \quad ((0 \leq a) \wedge (a < b)) \Rightarrow  \\ \quad \quad \exists n, k \in \mathbb{N}.\; ((a < (k \cdot \operatorname{recip}((n + 1)))) \wedge ((k \cdot \operatorname{recip}((n + 1))) < b)) \end{array}}}
\vnbprop{measurable-fn-grid-row}{\vnbstmt{\begin{array}{@{}l@{}} \forall a, b, f, g, l, r. \\ \quad (\operatorname{is\text{-}measurable\text{-}fn}(a, b, f) \wedge (\operatorname{is\text{-}measurable\text{-}fn}(a, b, g) \wedge ((l \in \mathbb{R}) \wedge (r \in \mathbb{R})))) \Rightarrow  \\ \quad \quad \{x \in a : \exists k \in \mathbb{N}.\; (((k \cdot r) < f(x)) \wedge ((l - (k \cdot r)) < g(x)))\} \in b \end{array}}}
\vnbprop{measurable-fn-eplus-of-values}{\vnbstmt{\begin{array}{@{}l@{}} \forall a, b, f, g, h. \\ \quad \operatorname{is\text{-}measurable\text{-}fn}(a, b, f) \wedge \\ \quad \operatorname{is\text{-}measurable\text{-}fn}(a, b, g) \wedge \\ \quad h \in (a \to \operatorname{rr\text{-}pos\text{-}star}) \wedge \\ \quad \forall x \in a.\; (h(x) = \operatorname{eplus}(f(x), g(x))) \Rightarrow \\ \quad \quad \operatorname{is\text{-}measurable\text{-}fn}(a, b, h) \end{array}}}
\vnbprop{measurable-fn-eplus}{\vnbstmt{\begin{array}{@{}l@{}} \forall a, b, f, g. \\ \quad (\operatorname{is\text{-}measurable\text{-}fn}(a, b, f) \wedge \operatorname{is\text{-}measurable\text{-}fn}(a, b, g)) \Rightarrow  \\ \quad \quad \operatorname{is\text{-}measurable\text{-}fn}(a, b, \operatorname{ptwise\text{-}eplus}(a, f, g)) \end{array}}}
\vnbprop{rr-lt-div-of-lt-mul}{\vnbstmt{\begin{array}{@{}l@{}} \forall l, f, g \in \mathbb{R}. \\ \quad ((0 < g) \wedge (l < (f \cdot g))) \Rightarrow ((l \cdot \operatorname{recip}(g)) < f) \end{array}}}
\vnbprop{rr-lt-mul-of-grid}{\vnbstmt{\begin{array}{@{}l@{}} \forall l, q, f, g \in \mathbb{R}. \\ \quad ((0 \leq l) \wedge ((0 < q) \wedge ((q < f) \wedge ((l \cdot \operatorname{recip}(q)) < g)))) \Rightarrow (l < (f \cdot g)) \end{array}}}
\vnbprop{measurable-fn-etimes-cell}{\vnbstmt{\begin{array}{@{}l@{}} \forall a, b, f, g, l, q. \\ \quad (\operatorname{is\text{-}measurable\text{-}fn}(a, b, f) \wedge (\operatorname{is\text{-}measurable\text{-}fn}(a, b, g) \wedge ((l \in \mathbb{R}) \wedge (q \in \mathbb{R})))) \Rightarrow  \\ \quad \quad \{x \in a : ((0 < q) \wedge ((q < f(x)) \wedge ((l \cdot \operatorname{recip}(q)) < g(x))))\} \in b \end{array}}}
\vnbprop{measurable-fn-etimes-row}{\vnbstmt{\begin{array}{@{}l@{}} \forall a, b, f, g, l, r. \\ \quad (\operatorname{is\text{-}measurable\text{-}fn}(a, b, f) \wedge (\operatorname{is\text{-}measurable\text{-}fn}(a, b, g) \wedge ((l \in \mathbb{R}) \wedge (r \in \mathbb{R})))) \Rightarrow  \\ \quad \quad \{x \in a : \exists k \in \mathbb{N}.\; ((0 < (k \cdot r)) \wedge (((k \cdot r) < f(x)) \wedge ((l \cdot \operatorname{recip}((k \cdot r))) < g(x))))\} \in b \end{array}}}
\vnbprop{measurable-fn-etimes-of-values}{\vnbstmt{\begin{array}{@{}l@{}} \forall a, b, f, g, h. \\ \quad \operatorname{is\text{-}measurable\text{-}fn}(a, b, f) \wedge \\ \quad \operatorname{is\text{-}measurable\text{-}fn}(a, b, g) \wedge \\ \quad h \in (a \to \operatorname{rr\text{-}pos\text{-}star}) \wedge \\ \quad \forall x \in a.\; (h(x) = \operatorname{etimes}(f(x), g(x))) \Rightarrow \\ \quad \quad \operatorname{is\text{-}measurable\text{-}fn}(a, b, h) \end{array}}}
\vnbprop{measurable-fn-etimes}{\vnbstmt{\begin{array}{@{}l@{}} \forall a, b, f, g. \\ \quad (\operatorname{is\text{-}measurable\text{-}fn}(a, b, f) \wedge \operatorname{is\text{-}measurable\text{-}fn}(a, b, g)) \Rightarrow  \\ \quad \quad \operatorname{is\text{-}measurable\text{-}fn}(a, b, \operatorname{ptwise\text{-}etimes}(a, f, g)) \end{array}}}
\vnbprop{esup-eplus-le}{\vnbstmt{\begin{array}{@{}l@{}} \forall u, v.; \forall l \in \operatorname{rr\text{-}pos\text{-}star}. \\ \quad u \in (\mathbb{N} \to \operatorname{rr\text{-}pos\text{-}star}) \wedge \\ \quad v \in (\mathbb{N} \to \operatorname{rr\text{-}pos\text{-}star}) \wedge \\ \quad \forall n \in \mathbb{N}.\; (u(n) \leq u((n + 1))) \wedge \\ \quad \forall n \in \mathbb{N}.\; (v(n) \leq v((n + 1))) \wedge \\ \quad \forall n \in \mathbb{N}.\; (\operatorname{eplus}(u(n), v(n)) \leq l) \Rightarrow \\ \quad \quad \operatorname{eplus}(\operatorname{esup}(\operatorname{image}(u, \mathbb{N})), \operatorname{esup}(\operatorname{image}(v, \mathbb{N}))) \leq l \end{array}}}
\vnbprop{rps-finsum-esup-le}{\vnbstmt{\begin{array}{@{}l@{}} \forall s \in \mathbf{Set}, h, a.; \forall l \in \operatorname{rr\text{-}pos\text{-}star}. \\ \quad |s| \in \mathbb{N} \wedge \\ \quad h \in (\mathbb{N} \to (s \to \operatorname{rr\text{-}pos\text{-}star})) \wedge \\ \quad \forall z \in s, n \in \mathbb{N}.\; (h(n)(z) \leq h((n + 1))(z)) \wedge \\ \quad \forall z \in s.\; (a(z) = \operatorname{esup}(\operatorname{image}((\lambda m \in \mathbb{N}.\; h(m)(z)), \mathbb{N}))) \wedge \\ \quad \forall n \in \mathbb{N}.\; (\operatorname{finsum}(\operatorname{rr\text{-}pos\text{-}star\text{-}add\text{-}monoid}, h(n), s) \leq l) \Rightarrow \\ \quad \quad \operatorname{finsum}(\operatorname{rr\text{-}pos\text{-}star\text{-}add\text{-}monoid}, a, s) \leq l \end{array}}}
\vnbprop{simple-fn-of-partition}{\vnbstmt{\begin{array}{@{}l@{}} \forall a, b, s, k, h, t. \\ \quad \operatorname{is\text{-}sigma\text{-}algebra}(a, b) \wedge \\ \quad s \in \mathbf{Set} \wedge \\ \quad |s| \in \mathbb{N} \wedge \\ \quad \forall z \in a.\; (k(z) \in s) \wedge \\ \quad \forall z \in s.\; (\{x \in a : (k(x) = z)\} \in b) \wedge \\ \quad \forall z \in s.\; (h(z) \in \mathbb{R}) \wedge \\ \quad \forall z \in s.\; (0 \leq h(z)) \wedge \\ \quad t \in (a \to \operatorname{rr\text{-}pos\text{-}star}) \wedge \\ \quad \forall z \in a.\; (t(z) = h(k(z))) \Rightarrow \\ \quad \quad \operatorname{is\text{-}simple\text{-}fn}(a, b, t) \end{array}}}
\vnbprop{rps-le-of-scaled}{\vnbstmt{\begin{array}{@{}l@{}} \forall u, l \in \operatorname{rr\text{-}pos\text{-}star}. \\ \quad \forall c \in \mathbb{R}.\; (((0 < c) \wedge (c < 1)) \Rightarrow (\operatorname{etimes}(c, u) \leq l)) \Rightarrow (u \leq l) \end{array}}}
\vnbprop{esup-etimes-ge}{\vnbstmt{\begin{array}{@{}l@{}} \forall w \in \mathbb{R}, m \in (\mathbb{N} \to \operatorname{rr\text{-}pos\text{-}star}). \\ \quad (0 \leq w) \Rightarrow  \\ \quad \quad \operatorname{etimes}(w, \operatorname{esup}(\operatorname{image}(m, \mathbb{N}))) \leq \operatorname{esup}(\operatorname{image}((\lambda n \in \mathbb{N}.\; \operatorname{etimes}(w, m(n))), \mathbb{N})) \end{array}}}
\vnbprop{simple-restrict-integral}{\vnbstmt{\begin{array}{@{}l@{}} \forall a, b, u, s, c, t, f. \\ \quad \operatorname{is\text{-}measure}(a, b, u) \wedge \\ \quad \operatorname{is\text{-}simple\text{-}fn}(a, b, s) \wedge \\ \quad c \in \mathbb{R} \wedge \\ \quad 0 \leq c \wedge \\ \quad \forall y \in \operatorname{image}(s, a).\; ((\{\mathit{six} \in a : (s(\mathit{six}) = y)\} \cap t) \in b) \wedge \\ \quad f \in (a \to \operatorname{rr\text{-}pos\text{-}star}) \wedge \\ \quad \forall x \in a.\; (f(x) = \left(\text{if }(x \in t)\text{ then }(c \cdot s(x))\text{ else }0\right)) \Rightarrow \\ \quad \quad \operatorname{is\text{-}simple\text{-}fn}(a, b, f) \wedge \\ \quad \quad \operatorname{simple\text{-}integral}(a, u, f) \\ \quad \quad =\; \operatorname{finsum}(\operatorname{rr\text{-}pos\text{-}star\text{-}add\text{-}monoid}, \\ \quad \quad \quad (\lambda y \in \operatorname{image}(s, a).\; \operatorname{etimes}((c \cdot y), u((\{\mathit{six} \in a : (s(\mathit{six}) = y)\} \cap t)))), \\ \quad \quad \quad \operatorname{image}(s, a)) \end{array}}}
\vnbprop{mct-simple-below}{\vnbstmt{\begin{array}{@{}l@{}} \forall a, b, u, f, s. \\ \quad \operatorname{is\text{-}measure}(a, b, u) \wedge \\ \quad f \in (\mathbb{N} \to (a \to \operatorname{rr\text{-}pos\text{-}star})) \wedge \\ \quad \forall n \in \mathbb{N}.\; \operatorname{is\text{-}measurable\text{-}fn}(a, b, f(n)) \wedge \\ \quad \forall n \in \mathbb{N}.\; \operatorname{ptwise\text{-}le}(a, f(n), f((n + 1))) \wedge \\ \quad \operatorname{is\text{-}simple\text{-}fn}(a, b, s) \wedge \\ \quad \operatorname{ptwise\text{-}le}(a, s, \operatorname{ptwise\text{-}sup}(a, f)) \Rightarrow \\ \quad \quad \operatorname{simple\text{-}integral}(a, u, s) \leq \operatorname{esup}(\operatorname{image}((\lambda n \in \mathbb{N}.\; \operatorname{integral}(a, b, u, f(n))), \mathbb{N})) \end{array}}}
\vnbprop{monotone-convergence}{\vnbstmt{\begin{array}{@{}l@{}} \forall a, b, u, f. \\ \quad \operatorname{is\text{-}measure}(a, b, u) \wedge \\ \quad f \in (\mathbb{N} \to (a \to \operatorname{rr\text{-}pos\text{-}star})) \wedge \\ \quad \forall n \in \mathbb{N}.\; \operatorname{is\text{-}measurable\text{-}fn}(a, b, f(n)) \wedge \\ \quad \forall n \in \mathbb{N}.\; \operatorname{ptwise\text{-}le}(a, f(n), f((n + 1))) \Rightarrow \\ \quad \quad \operatorname{is\text{-}measurable\text{-}fn}(a, b, \operatorname{ptwise\text{-}sup}(a, f)) \wedge \\ \quad \quad \operatorname{integral}(a, b, u, \operatorname{ptwise\text{-}sup}(a, f)) = \operatorname{esup}(\operatorname{image}((\lambda n \in \mathbb{N}.\; \operatorname{integral}(a, b, u, f(n))), \mathbb{N})) \end{array}}}
\vnbprop{fatou}{\vnbstmt{\begin{array}{@{}l@{}} \forall a, b, u, f. \\ \quad \operatorname{is\text{-}measure}(a, b, u) \wedge \\ \quad f \in (\mathbb{N} \to (a \to \operatorname{rr\text{-}pos\text{-}star})) \wedge \\ \quad \forall n \in \mathbb{N}.\; \operatorname{is\text{-}measurable\text{-}fn}(a, b, f(n)) \Rightarrow \\ \quad \quad \operatorname{is\text{-}measurable\text{-}fn}(a, b, \operatorname{ptwise\text{-}liminf}(a, f)) \wedge \\ \quad \quad \operatorname{integral}(a, b, u, \operatorname{ptwise\text{-}liminf}(a, f)) \leq \operatorname{eliminf}((\lambda n \in \mathbb{N}.\; \operatorname{integral}(a, b, u, f(n)))) \end{array}}}
\vnbprop{esup-grid-max}{\vnbstmt{\begin{array}{@{}l@{}} \forall d, y \in \mathbb{R}, m \in \mathbb{N}, a. \\ \quad ((0 < d) \wedge ((a \subseteq \mathbb{N}) \wedge ((m \in a) \wedge ((y < ((m \cdot d) + d)) \wedge \forall k \in a.\; ((k \cdot d) \leq y))))) \Rightarrow  \\ \quad \quad \operatorname{esup}(\operatorname{image}((\lambda \mathit{sk} \in \mathbb{N}.\; (\mathit{sk} \cdot d)), a)) = (m \cdot d) \end{array}}}
\vnbprop{grid-staircase-value}{\vnbstmt{\begin{array}{@{}l@{}} \forall n \in \mathbb{N}, x \in \operatorname{rr\text{-}pos\text{-}star}. \exists m \in \mathbb{N}. \\ \quad m \in \{q \in \mathbb{N} : (((q \cdot \operatorname{power}(1/2, n)) \leq x) \wedge ((q \cdot \operatorname{power}(1/2, n)) \leq n))\} \wedge \\ \quad \operatorname{esup}(\operatorname{image}((\lambda k \in \mathbb{N}.\; (k \cdot \operatorname{power}(1/2, n))), \\ \quad \quad \quad \{q \in \mathbb{N} : (((q \cdot \operatorname{power}(1/2, n)) \leq x) \wedge ((q \cdot \operatorname{power}(1/2, n)) \leq n))\})) \\ \quad =\; m \cdot \operatorname{power}(1/2, n) \wedge \\ \quad \forall z \in \mathbb{R}.\; (((z \leq x) \wedge (z \leq n)) \Rightarrow (z < ((m \cdot \operatorname{power}(1/2, n)) + \operatorname{power}(1/2, n)))) \end{array}}}
\vnbprop{measurable-fn-guarded-ge}{\vnbstmt{\begin{array}{@{}l@{}} \forall c, d, f, a, b, g. \\ \quad (\operatorname{is\text{-}measurable\text{-}fn}(c, d, f) \wedge (b \in \mathbb{R})) \Rightarrow  \\ \quad \quad \{x \in c : (((a < b) \wedge (b \leq g)) \wedge (b \leq f(x)))\} \in d \end{array}}}
\vnbprop{grid-staircase-simple}{\vnbstmt{\begin{array}{@{}l@{}} \forall a, b, f.; \forall n \in \mathbb{N}. \\ \quad \operatorname{is\text{-}measurable\text{-}fn}(a, b, f) \Rightarrow  \\ \quad \quad \operatorname{is\text{-}simple\text{-}fn}(a, b, \\ \quad \quad \quad (\lambda x \in a.\; \\ \quad \quad \quad \quad \operatorname{esup}(\operatorname{image}((\lambda k \in \mathbb{N}.\; (k \cdot \operatorname{power}(1/2, n))), \\ \quad \quad \quad \quad \quad \quad \{q \in \mathbb{N} : (((q \cdot \operatorname{power}(1/2, n)) \leq f(x)) \wedge ((q \cdot \operatorname{power}(1/2, n)) \leq n))\})))) \end{array}}}
\vnbprop{integral-simple-approx}{\vnbstmt{\begin{array}{@{}l@{}} \forall a, b, f. \\ \quad (\operatorname{is\text{-}sigma\text{-}algebra}(a, b) \wedge \operatorname{is\text{-}measurable\text{-}fn}(a, b, f)) \Rightarrow  \\ \quad \quad \exists c \in (\mathbb{N} \to (a \to \operatorname{rr\text{-}pos\text{-}star})). \\ \quad \quad \quad \forall n \in \mathbb{N}.\; \operatorname{is\text{-}simple\text{-}fn}(a, b, c(n)) \wedge \\ \quad \quad \quad \forall n \in \mathbb{N}.\; \operatorname{ptwise\text{-}le}(a, c(n), c((n + 1))) \wedge \\ \quad \quad \quad \forall x \in a.\; (\operatorname{ptwise\text{-}sup}(a, c)(x) = f(x)) \end{array}}}
\vnbprop{esup-eplus-eq}{\vnbstmt{\begin{array}{@{}l@{}} \forall u, v. \\ \quad u \in (\mathbb{N} \to \operatorname{rr\text{-}pos\text{-}star}) \wedge \\ \quad v \in (\mathbb{N} \to \operatorname{rr\text{-}pos\text{-}star}) \wedge \\ \quad \forall n \in \mathbb{N}.\; (u(n) \leq u((n + 1))) \wedge \\ \quad \forall n \in \mathbb{N}.\; (v(n) \leq v((n + 1))) \Rightarrow \\ \quad \quad \operatorname{esup}(\operatorname{image}((\lambda a \in \mathbb{N}.\; \operatorname{eplus}(u(a), v(a))), \mathbb{N})) \\ \quad \quad =\; \operatorname{eplus}(\operatorname{esup}(\operatorname{image}(u, \mathbb{N})), \operatorname{esup}(\operatorname{image}(v, \mathbb{N}))) \end{array}}}
\vnbprop{simple-fn-eplus}{\vnbstmt{\begin{array}{@{}l@{}} \forall a, b, s, f. \\ \quad (\operatorname{is\text{-}simple\text{-}fn}(a, b, s) \wedge \operatorname{is\text{-}simple\text{-}fn}(a, b, f)) \Rightarrow  \\ \quad \quad \operatorname{is\text{-}simple\text{-}fn}(a, b, \operatorname{ptwise\text{-}eplus}(a, s, f)) \end{array}}}
\vnbprop{integral-additive}{\vnbstmt{\begin{array}{@{}l@{}} \forall a, b, u, f, g. \\ \quad (\operatorname{is\text{-}measure}(a, b, u) \wedge (\operatorname{is\text{-}measurable\text{-}fn}(a, b, f) \wedge \operatorname{is\text{-}measurable\text{-}fn}(a, b, g))) \Rightarrow  \\ \quad \quad \operatorname{integral}(a, b, u, \operatorname{ptwise\text{-}eplus}(a, f, g)) = \operatorname{eplus}(\operatorname{integral}(a, b, u, f), \operatorname{integral}(a, b, u, g)) \end{array}}}
\vnbprop{measurable-fn-lt-row}{\vnbstmt{\begin{array}{@{}l@{}} \forall a, b, u, v, r. \\ \quad (\operatorname{is\text{-}measurable\text{-}fn}(a, b, u) \wedge (\operatorname{is\text{-}measurable\text{-}fn}(a, b, v) \wedge (r \in \mathbb{R}))) \Rightarrow  \\ \quad \quad \{x \in a : \exists k \in \mathbb{N}.\; ((v(x) < (k \cdot r)) \wedge ((k \cdot r) < u(x)))\} \in b \end{array}}}
\vnbprop{measurable-fn-lt-set}{\vnbstmt{\begin{array}{@{}l@{}} \forall a, b, u, v. \\ \quad (\operatorname{is\text{-}measurable\text{-}fn}(a, b, u) \wedge \operatorname{is\text{-}measurable\text{-}fn}(a, b, v)) \Rightarrow  \\ \quad \quad \{x \in a : (v(x) < u(x))\} \in b \end{array}}}
\vnbprop{measurable-fn-const}{\vnbstmt{\begin{array}{@{}l@{}} \forall a, b, c. \\ \quad (\operatorname{is\text{-}sigma\text{-}algebra}(a, b) \wedge (c \in \operatorname{rr\text{-}pos\text{-}star})) \Rightarrow  \\ \quad \quad \operatorname{is\text{-}measurable\text{-}fn}(a, b, (\lambda x \in a.\; c)) \end{array}}}
\vnbprop{measurable-fn-finite-part}{\vnbstmt{\forall a, b, u.\; (\operatorname{is\text{-}measurable\text{-}fn}(a, b, u) \Rightarrow (\{p \in a : (u(p) \in \mathbb{R})\} \in b))}}
\vnbprop{integral-finite-null}{\vnbstmt{\begin{array}{@{}l@{}} \forall a, b, c, u. \\ \quad (\operatorname{is\text{-}measure}(a, b, c) \wedge (\operatorname{is\text{-}measurable\text{-}fn}(a, b, u) \wedge (\operatorname{integral}(a, b, c, u) \in \mathbb{R}))) \Rightarrow  \\ \quad \quad c(\operatorname{complement\text{-}in}(a, \{p \in a : (u(p) \in \mathbb{R})\})) = 0 \end{array}}}
\vnbprop{integral-le-off-null}{\vnbstmt{\begin{array}{@{}l@{}} \forall a, b, u, c, d, f. \\ \quad \operatorname{is\text{-}measure}(a, b, u) \wedge \\ \quad c \in b \wedge \\ \quad u(c) = 0 \wedge \\ \quad \operatorname{is\text{-}measurable\text{-}fn}(a, b, d) \wedge \\ \quad \operatorname{is\text{-}measurable\text{-}fn}(a, b, f) \wedge \\ \quad \forall x \in a.\; (\neg (x \in c) \Rightarrow (d(x) \leq f(x))) \Rightarrow \\ \quad \quad \operatorname{integral}(a, b, u, d) \leq \operatorname{integral}(a, b, u, f) \end{array}}}
\vnbprop{difference-value}{\vnbstmt{\begin{array}{@{}l@{}} \forall a, u, v.; \forall x \in a. \\ \quad ((\lambda \mathit{dx} \in a.\; \left(\text{if }(u(\mathit{dx}) \in \mathbb{R})\text{ then }(u(\mathit{dx}) - v(\mathit{dx}))\text{ else }0\right))(x) \simeq \left(\text{if }(u(x) \in \mathbb{R})\text{ then }(u(x) - v(x))\text{ else }0\right)) \end{array}}}
\vnbprop{measurable-fn-sub}{\vnbstmt{\begin{array}{@{}l@{}} \forall a, b, u, v. \\ \quad (\operatorname{is\text{-}measurable\text{-}fn}(a, b, u) \wedge (\operatorname{is\text{-}measurable\text{-}fn}(a, b, v) \wedge \forall x \in a.\; (v(x) \leq u(x)))) \Rightarrow  \\ \quad \quad \operatorname{is\text{-}measurable\text{-}fn}(a, b, (\lambda \mathit{dx} \in a.\; \left(\text{if }(u(\mathit{dx}) \in \mathbb{R})\text{ then }(u(\mathit{dx}) - v(\mathit{dx}))\text{ else }0\right))) \end{array}}}
\vnbprop{integral-sub}{\vnbstmt{\begin{array}{@{}l@{}} \forall a, b, c, u, v. \\ \quad \operatorname{is\text{-}measure}(a, b, c) \wedge \\ \quad \operatorname{is\text{-}measurable\text{-}fn}(a, b, u) \wedge \\ \quad \operatorname{is\text{-}measurable\text{-}fn}(a, b, v) \wedge \\ \quad \forall x \in a.\; (v(x) \leq u(x)) \wedge \\ \quad \operatorname{integral}(a, b, c, u) \in \mathbb{R} \Rightarrow \\ \quad \quad \operatorname{eplus}(\operatorname{integral}(a, b, c, (\lambda \mathit{dx} \in a.\; \left(\text{if }(u(\mathit{dx}) \in \mathbb{R})\text{ then }(u(\mathit{dx}) - v(\mathit{dx}))\text{ else }0\right))), \operatorname{integral}(a, b, c, v)) \\ \quad \quad =\; \operatorname{integral}(a, b, c, u) \end{array}}}
\vnbprop{elimsup-le-bound}{\vnbstmt{\begin{array}{@{}l@{}} \forall a \in (\mathbb{N} \to \operatorname{rr\text{-}pos\text{-}star}), b \in \operatorname{rr\text{-}pos\text{-}star}. \\ \quad \forall n \in \mathbb{N}.\; (a(n) \leq b) \Rightarrow (\operatorname{elimsup}(a) \leq b) \end{array}}}
\vnbprop{eliminf-le-elimsup}{\vnbstmt{\forall a \in (\mathbb{N} \to \operatorname{rr\text{-}pos\text{-}star}).\; (\operatorname{eliminf}(a) \leq \operatorname{elimsup}(a))}}
\vnbprop{einf-reflection-tail}{\vnbstmt{\begin{array}{@{}l@{}} \forall c \in \mathbb{R}, a, b.; \forall m \in \mathbb{N}. \\ \quad a \in (\mathbb{N} \to \operatorname{rr\text{-}pos\text{-}star}) \wedge \\ \quad b \in (\mathbb{N} \to \operatorname{rr\text{-}pos\text{-}star}) \wedge \\ \quad \forall n \in \mathbb{N}.\; (a(n) \leq c) \wedge \\ \quad \forall n \in \mathbb{N}.\; (b(n) = (c - a(n))) \Rightarrow \\ \quad \quad (\operatorname{esup}(\operatorname{etail}(a, m)) \in \mathbb{R}) \wedge (\operatorname{einf}(\operatorname{etail}(b, m)) = (c - \operatorname{esup}(\operatorname{etail}(a, m)))) \end{array}}}
\vnbprop{eliminf-reflection}{\vnbstmt{\begin{array}{@{}l@{}} \forall c \in \mathbb{R}, a, b. \\ \quad a \in (\mathbb{N} \to \operatorname{rr\text{-}pos\text{-}star}) \wedge \\ \quad b \in (\mathbb{N} \to \operatorname{rr\text{-}pos\text{-}star}) \wedge \\ \quad \forall n \in \mathbb{N}.\; (a(n) \leq c) \wedge \\ \quad \forall n \in \mathbb{N}.\; (b(n) = (c - a(n))) \Rightarrow \\ \quad \quad (\operatorname{elimsup}(a) \in \mathbb{R}) \wedge (\operatorname{eliminf}(b) = (c - \operatorname{elimsup}(a))) \end{array}}}
\vnbprop{dominated-convergence}{\vnbstmt{\begin{array}{@{}l@{}} \forall a, b, u, c, f, g. \\ \quad \operatorname{is\text{-}measure}(a, b, u) \wedge \\ \quad c \in (\mathbb{N} \to (a \to \operatorname{rr\text{-}pos\text{-}star})) \wedge \\ \quad \forall n \in \mathbb{N}.\; \operatorname{is\text{-}measurable\text{-}fn}(a, b, c(n)) \wedge \\ \quad \operatorname{is\text{-}measurable\text{-}fn}(a, b, f) \wedge \\ \quad \operatorname{is\text{-}measurable\text{-}fn}(a, b, g) \wedge \\ \quad \forall n \in \mathbb{N}.\; \operatorname{ptwise\text{-}le}(a, c(n), g) \wedge \\ \quad \operatorname{integral}(a, b, u, g) \in \mathbb{R} \wedge \\ \quad \forall x \in a.\; \operatorname{econverges\text{-}to}((\lambda n \in \mathbb{N}.\; c(n)(x)), f(x)) \Rightarrow \\ \quad \quad \operatorname{econverges\text{-}to}((\lambda n \in \mathbb{N}.\; \operatorname{integral}(a, b, u, c(n))), \operatorname{integral}(a, b, u, f)) \end{array}}}

\vnbbookfile{theorem-library/caratheodory}{theorem-library/caratheodory}
\noindent{\small caratheodory.scm -- THE CARATHEODORY CONSTRUCTION, PROVEN. Agent MEASURE-1, 2026-10-07 (scratchpad/triage/MEASURE-1.md).  Library: measure.  Helper prefix cth-.  16 theorems, every one `modulo 0'.}\par
\vnbprop{difference-unfold}{\vnbstmt{\forall u, b.\; (\operatorname{difference}(u, b) \simeq \operatorname{complement\text{-}in}(u, b))}}
\vnbprop{inter-complement-in-shift}{\vnbstmt{\begin{array}{@{}l@{}} \forall a, u, e. \\ \quad (\operatorname{complement\text{-}in}(a, u) \cap e) = (a \cap \operatorname{complement\text{-}in}(e, u)) \end{array}}}
\vnbprop{complement-in-complement-in-union}{\vnbstmt{\begin{array}{@{}l@{}} \forall a, u, e. \\ \quad \operatorname{complement\text{-}in}(\operatorname{complement\text{-}in}(a, u), e) = \operatorname{complement\text{-}in}(a, (u \cup e)) \end{array}}}
\vnbprop{complement-in-empty-set}{\vnbstmt{\forall a.\; (\operatorname{complement\text{-}in}(a, \emptyset) = a)}}
\vnbprop{union-inter-complement-in-split}{\vnbstmt{\forall a, e.\; (a = ((a \cap e) \cup \operatorname{complement\text{-}in}(a, e)))}}
\vnbprop{inter-complement-in-ambient}{\vnbstmt{\begin{array}{@{}l@{}} \forall t, m, e. \\ \quad (t \subseteq m) \Rightarrow ((t \cap \operatorname{complement\text{-}in}(m, e)) = \operatorname{complement\text{-}in}(t, e)) \end{array}}}
\vnbprop{complement-in-complement-in-ambient}{\vnbstmt{\begin{array}{@{}l@{}} \forall t, m, e. \\ \quad (t \subseteq m) \Rightarrow (\operatorname{complement\text{-}in}(t, \operatorname{complement\text{-}in}(m, e)) = (t \cap e)) \end{array}}}
\vnbprop{inter-of-subset}{\vnbstmt{\forall t, m.\; ((t \subseteq m) \Rightarrow ((t \cap m) = t))}}
\vnbprop{complement-in-of-subset-empty}{\vnbstmt{\forall t, m.\; ((t \subseteq m) \Rightarrow (\operatorname{complement\text{-}in}(t, m) = \emptyset))}}
\vnbprop{outer-measure-binary-subadditive}{\vnbstmt{\begin{array}{@{}l@{}} \forall a, u, \alpha, b. \\ \quad (\operatorname{is\text{-}outer\text{-}measure}(a, u) \wedge ((\alpha \in \operatorname{power}(a)) \wedge (b \in \operatorname{power}(a)))) \Rightarrow  \\ \quad \quad u((\alpha \cup b)) \leq \operatorname{eplus}(u(\alpha), u(b)) \end{array}}}
\vnbprop{outer-measure-split-le}{\vnbstmt{\begin{array}{@{}l@{}} \forall a, u, \alpha, s. \\ \quad (\operatorname{is\text{-}outer\text{-}measure}(a, u) \wedge ((\alpha \in \operatorname{power}(a)) \wedge (s \in \operatorname{power}(a)))) \Rightarrow  \\ \quad \quad u(\alpha) \leq \operatorname{eplus}(u((\alpha \cap s)), u(\operatorname{complement\text{-}in}(\alpha, s))) \end{array}}}
\vnbprop{caratheodory-null-measurable}{\vnbstmt{\begin{array}{@{}l@{}} \forall a, u, s. \\ \quad (\operatorname{is\text{-}outer\text{-}measure}(a, u) \wedge ((s \subseteq a) \wedge (u(s) = 0))) \Rightarrow  \\ \quad \quad \operatorname{is\text{-}caratheodory\text{-}measurable}(a, u, s) \end{array}}}
\vnbprop{caratheodory-superadditive-suffices}{\vnbstmt{\begin{array}{@{}l@{}} \forall a, u, s. \\ \quad \operatorname{is\text{-}outer\text{-}measure}(a, u) \wedge \\ \quad s \subseteq a \wedge \\ \quad \forall \alpha \in \operatorname{power}(a). \\ \quad \quad \operatorname{eplus}(u((\alpha \cap s)), u(\operatorname{difference}(\alpha, s))) \leq u(\alpha) \Rightarrow \\ \quad \quad \operatorname{is\text{-}caratheodory\text{-}measurable}(a, u, s) \end{array}}}
\vnbprop{caratheodory-prefix-split}{\vnbstmt{\begin{array}{@{}l@{}} \forall n \in \mathbb{N}, a, u, f, \alpha. \\ \quad (\operatorname{is\text{-}outer\text{-}measure}(a, u) \wedge ((f \in (\mathbb{N} \to \operatorname{caratheodory\text{-}sets}(a, u))) \wedge (\alpha \in \operatorname{power}(a)))) \Rightarrow  \\ \quad \quad u(\alpha) \\ \quad \quad =\; \operatorname{eplus}(\operatorname{finsum}(\operatorname{rr\text{-}pos\text{-}star\text{-}add\text{-}monoid}, \\ \quad \quad \quad \quad (\lambda k \in \mathbb{N}.\; u((\alpha \cap \operatorname{complement\text{-}in}(f(k), \operatorname{big\text{-}union}(j, \mathrm{OS}(k), f(j)))))), \\ \quad \quad \quad \quad \mathrm{OS}(n)), \\ \quad \quad \quad u(\operatorname{complement\text{-}in}(\alpha, \operatorname{big\text{-}union}(j, \mathrm{OS}(n), f(j))))) \end{array}}}
\vnbprop{caratheodory-measurable-sigma-algebra}{\vnbstmt{\begin{array}{@{}l@{}} \forall a, u. \\ \quad \operatorname{is\text{-}outer\text{-}measure}(a, u) \Rightarrow \operatorname{is\text{-}sigma\text{-}algebra}(a, \operatorname{caratheodory\text{-}sets}(a, u)) \end{array}}}
\vnbprop{caratheodory-measure}{\vnbstmt{\begin{array}{@{}l@{}} \forall a, u. \\ \quad \operatorname{is\text{-}outer\text{-}measure}(a, u) \Rightarrow  \\ \quad \quad \operatorname{is\text{-}measure}(a, \operatorname{caratheodory\text{-}sets}(a, u), \operatorname{res}(u, \operatorname{caratheodory\text{-}sets}(a, u))) \end{array}}}

\vnbbookfile{structure-library/outer-measure-construction}{structure-library/outer-measure-construction}
\noindent{\small outer-measure-construction.scm -- METHOD I OUTER MEASURES; null sets, complete and sigma-finite measures.  Vocabulary only: no proof, no support. Agent MEASURE-2, 2026-10-08 (scratchpad/triage/MEASURE-2.md).  The proofs are theorem-library/outer-measure-construction.scm.}\par
\par\medskip\noindent\textbf{Definition} (\texttt{is-premeasure}).\;For all $\omega$, $\mathit{cls}$ and $\mathit{pm}$, 
\begin{equation*}
\begin{array}{@{}l@{}} \operatorname{is\text{-}premeasure}(\omega, \mathit{cls}, \mathit{pm}) \\ \Leftrightarrow\; \omega \in \mathbf{Set} \wedge \\ \mathit{cls} \subseteq \operatorname{power}(\omega) \wedge \\ \emptyset \in \mathit{cls} \wedge \\ \mathit{pm} \in (\mathit{cls} \to \operatorname{rr\text{-}pos\text{-}star}) \wedge \\ \operatorname{pm}(\emptyset) = 0 \end{array}
\end{equation*}

\par\medskip\noindent\textbf{Definition} (\texttt{cover-sums}).\;$\operatorname{cover\text{-}sums}(\omega, \mathit{cls}, \mathit{pm}, \mathit{es}) :=$
\begin{equation*}
\{\operatorname{cvv\_} \in \operatorname{rr\text{-}pos\text{-}star} : \exists \operatorname{cvs\_} \in (\mathbb{N} \to \mathit{cls}).\; ((\mathit{es} \subseteq \operatorname{big\text{-}union}(\operatorname{cvn\_}, \mathbb{N}, \operatorname{cvs\_}(\operatorname{cvn\_}))) \wedge (\operatorname{cvv\_} = \operatorname{esum}((\lambda \operatorname{cvk\_} \in \mathbb{N}.\; \operatorname{pm}(\operatorname{cvs\_}(\operatorname{cvk\_}))))))\}
\end{equation*}

\par\medskip\noindent\textbf{Definition} (\texttt{method-i-outer}).\;$\operatorname{method\text{-}i\text{-}outer}(\omega, \mathit{cls}, \mathit{pm}) :=$
\begin{equation*}
(\lambda \operatorname{cve\_} \in \operatorname{power}(\omega).\; \operatorname{einf}(\operatorname{cover\text{-}sums}(\omega, \mathit{cls}, \mathit{pm}, \operatorname{cve\_})))
\end{equation*}

\par\medskip\noindent\textbf{Definition} (\texttt{is-null-set}).\;For all $\omega$, $\mathit{ca}$, $\mu$ and $\mathit{ns}$, 
\begin{equation*}
\operatorname{is\text{-}null\text{-}set}(\omega, \mathit{ca}, \mu, \mathit{ns}) \Leftrightarrow \exists \operatorname{nla\_} \in \mathit{ca}.\; ((\mathit{ns} \subseteq \operatorname{nla\_}) \wedge (\mu(\operatorname{nla\_}) = 0))
\end{equation*}

\par\medskip\noindent\textbf{Definition} (\texttt{is-complete-measure}).\;For all $\omega$, $\mathit{ca}$ and $\mu$, 
\begin{equation*}
\begin{array}{@{}l@{}} \operatorname{is\text{-}complete\text{-}measure}(\omega, \mathit{ca}, \mu) \\ \Leftrightarrow\; (\operatorname{is\text{-}measure}(\omega, \mathit{ca}, \mu) \wedge \forall \operatorname{nls\_}.\; (\operatorname{is\text{-}null\text{-}set}(\omega, \mathit{ca}, \mu, \operatorname{nls\_}) \Rightarrow (\operatorname{nls\_} \in \mathit{ca}))) \end{array}
\end{equation*}

\par\medskip\noindent\textbf{Definition} (\texttt{is-sigma-finite}).\;For all $\omega$, $\mathit{ca}$ and $\mu$, 
\begin{equation*}
\begin{array}{@{}l@{}} \operatorname{is\text{-}sigma\text{-}finite}(\omega, \mathit{ca}, \mu) \\ \Leftrightarrow\; \exists \operatorname{sff\_} \in (\mathbb{N} \to \mathit{ca}). \\ \quad (\forall \operatorname{sfn\_} \in \mathbb{N}.\; (\mu(\operatorname{sff\_}(\operatorname{sfn\_})) \in \mathbb{R}) \wedge (\operatorname{big\text{-}union}(\operatorname{sfk\_}, \mathbb{N}, \operatorname{sff\_}(\operatorname{sfk\_})) = \omega)) \end{array}
\end{equation*}

\par\medskip\noindent\textbf{Definition} (\texttt{closed-intervals}).\;
\begin{equation*}
(\operatorname{closed\text{-}intervals} \simeq \{\operatorname{oci\_} \in \operatorname{power}(\mathbb{R}) : ((\operatorname{oci\_} = \emptyset) \vee \exists \operatorname{oca\_}, \operatorname{ocb\_} \in \mathbb{R}.\; ((\operatorname{oca\_} \leq \operatorname{ocb\_}) \wedge (\operatorname{oci\_} = \operatorname{ccint}(\operatorname{oca\_}, \operatorname{ocb\_}))))\})
\end{equation*}

\par\medskip\noindent\textbf{Definition} (\texttt{interval-length}).\;$\operatorname{interval\text{-}length}(s) :=$
\begin{equation*}
\operatorname{esup}(\{\operatorname{odv\_} \in \operatorname{rr\text{-}pos\text{-}star} : ((\operatorname{odv\_} = 0) \vee \exists \operatorname{odx\_}, \operatorname{ody\_} \in s.\; ((\operatorname{odx\_} \leq \operatorname{ody\_}) \wedge (\operatorname{odv\_} = (\operatorname{ody\_} - \operatorname{odx\_}))))\})
\end{equation*}


\vnbbookfile{theorem-library/outer-measure-construction}{theorem-library/outer-measure-construction}
\noindent{\small outer-measure-construction.scm -- METHOD I OUTER MEASURES; the outer measure of a measure, its Caratheodory measure, completeness and the completion; Lebesgue outer measure on the line as an instance. Agent MEASURE-2, 2026-10-08 (scratchpad/triage/MEASURE-2.md).  Library:}\par
\vnbprop{cover-sums-unfold}{\vnbstmt{\begin{array}{@{}l@{}} \forall a, s, m, b. \\ \quad (\operatorname{cover\text{-}sums}(a, s, m, b) \simeq \{v \in \operatorname{rr\text{-}pos\text{-}star} : \exists \mathit{cvs} \in (\mathbb{N} \to s).\; ((b \subseteq \operatorname{big\text{-}union}(n, \mathbb{N}, \operatorname{cvs}(n))) \wedge (v = \operatorname{esum}((\lambda k \in \mathbb{N}.\; m(\operatorname{cvs}(k))))))\}) \end{array}}}
\vnbprop{cover-sums-subset}{\vnbstmt{\forall a, s, m, b.\; (\operatorname{cover\text{-}sums}(a, s, m, b) \subseteq \operatorname{rr\text{-}pos\text{-}star})}}
\vnbprop{cover-sums-elim}{\vnbstmt{\begin{array}{@{}l@{}} \forall a, s, m, b.; \forall v \in \operatorname{cover\text{-}sums}(a, s, m, b). \exists \mathit{cvs} \in (\mathbb{N} \to s). \\ \quad (b \subseteq \operatorname{big\text{-}union}(n, \mathbb{N}, \operatorname{cvs}(n))) \wedge (v = \operatorname{esum}((\lambda k \in \mathbb{N}.\; m(\operatorname{cvs}(k))))) \end{array}}}
\vnbprop{cover-sums-intro}{\vnbstmt{\begin{array}{@{}l@{}} \forall a, s, m, b, q. \\ \quad ((m \in (s \to \operatorname{rr\text{-}pos\text{-}star})) \wedge ((q \in (\mathbb{N} \to s)) \wedge (b \subseteq \operatorname{big\text{-}union}(n, \mathbb{N}, q(n))))) \Rightarrow  \\ \quad \quad \operatorname{esum}((\lambda k \in \mathbb{N}.\; m(q(k)))) \in \operatorname{cover\text{-}sums}(a, s, m, b) \end{array}}}
\vnbprop{method-i-outer-value}{\vnbstmt{\begin{array}{@{}l@{}} \forall a, s, m.; \forall b \in \operatorname{power}(a). \\ \quad \operatorname{method\text{-}i\text{-}outer}(a, s, m)(b) = \operatorname{einf}(\operatorname{cover\text{-}sums}(a, s, m, b)) \end{array}}}
\vnbprop{method-i-outer-type}{\vnbstmt{\begin{array}{@{}l@{}} \forall a \in \mathbf{Set}, s, m. \\ \quad \operatorname{method\text{-}i\text{-}outer}(a, s, m) \in (\operatorname{power}(a) \to \operatorname{rr\text{-}pos\text{-}star}) \end{array}}}
\vnbprop{einf-lt-witness}{\vnbstmt{\begin{array}{@{}l@{}} \forall s.; \forall a \in \operatorname{rr\text{-}pos\text{-}star}. \\ \quad ((s \subseteq \operatorname{rr\text{-}pos\text{-}star}) \wedge (\operatorname{einf}(s) < a)) \Rightarrow \exists y \in s.\; (y < a) \end{array}}}
\vnbprop{rps-le-of-forall-eps}{\vnbstmt{\begin{array}{@{}l@{}} \forall x, y. \\ \quad ((x \in \operatorname{rr\text{-}pos\text{-}star}) \wedge ((y \in \mathbb{R}) \wedge ((0 \leq y) \wedge \forall s.\; (\operatorname{pos\text{-}rr}(s) \Rightarrow (x \leq \operatorname{eplus}(y, s)))))) \Rightarrow  \\ \quad \quad x \leq y \end{array}}}
\vnbprop{eplus-swap4}{\vnbstmt{\begin{array}{@{}l@{}} \forall a, b, c, d. \\ \quad ((a \in \operatorname{rr\text{-}pos\text{-}star}) \wedge ((b \in \operatorname{rr\text{-}pos\text{-}star}) \wedge ((c \in \operatorname{rr\text{-}pos\text{-}star}) \wedge (d \in \operatorname{rr\text{-}pos\text{-}star})))) \Rightarrow  \\ \quad \quad \operatorname{eplus}(\operatorname{eplus}(a, b), \operatorname{eplus}(c, d)) = \operatorname{eplus}(\operatorname{eplus}(a, c), \operatorname{eplus}(b, d)) \end{array}}}
\vnbprop{rps-prefix-add}{\vnbstmt{\begin{array}{@{}l@{}} \forall n \in \mathbb{N}, f, g. \\ \quad ((f \in (\mathbb{N} \to \operatorname{rr\text{-}pos\text{-}star})) \wedge (g \in (\mathbb{N} \to \operatorname{rr\text{-}pos\text{-}star}))) \Rightarrow  \\ \quad \quad \operatorname{finsum}(\operatorname{rr\text{-}pos\text{-}star\text{-}add\text{-}monoid}, (\lambda k \in \mathbb{N}.\; \operatorname{eplus}(f(k), g(k))), \mathrm{OS}(n)) \\ \quad \quad =\; \operatorname{eplus}(\operatorname{finsum}(\operatorname{rr\text{-}pos\text{-}star\text{-}add\text{-}monoid}, f, \mathrm{OS}(n)), \\ \quad \quad \quad \operatorname{finsum}(\operatorname{rr\text{-}pos\text{-}star\text{-}add\text{-}monoid}, g, \mathrm{OS}(n))) \end{array}}}
\vnbprop{esum-nn-add-le}{\vnbstmt{\begin{array}{@{}l@{}} \forall f, g. \\ \quad ((f \in (\mathbb{N} \to \operatorname{rr\text{-}pos\text{-}star})) \wedge (g \in (\mathbb{N} \to \operatorname{rr\text{-}pos\text{-}star}))) \Rightarrow  \\ \quad \quad \operatorname{esum}((\lambda k \in \mathbb{N}.\; \operatorname{eplus}(f(k), g(k)))) \leq \operatorname{eplus}(\operatorname{esum}(f), \operatorname{esum}(g)) \end{array}}}
\vnbprop{dyadic-eps-type}{\vnbstmt{\begin{array}{@{}l@{}} \forall s. \\ \quad \operatorname{pos\text{-}rr}(s) \Rightarrow  \\ \quad \quad (\lambda i \in \mathbb{N}.\; (s \cdot \operatorname{recip}(\operatorname{power}(2, (i + 1))))) \in (\mathbb{N} \to \operatorname{rr\text{-}pos\text{-}star}) \wedge \\ \quad \quad \forall k \in \mathbb{N}.\; (0 < (\lambda i \in \mathbb{N}.\; (s \cdot \operatorname{recip}(\operatorname{power}(2, (i + 1)))))(k)) \end{array}}}
\vnbprop{dyadic-eps-prefix}{\vnbstmt{\begin{array}{@{}l@{}} \forall a \in \mathbb{N}, s. \\ \quad \operatorname{pos\text{-}rr}(s) \Rightarrow  \\ \quad \quad \operatorname{finsum}(\operatorname{rr\text{-}pos\text{-}star\text{-}add\text{-}monoid}, (\lambda i \in \mathbb{N}.\; (s \cdot \operatorname{recip}(\operatorname{power}(2, (i + 1))))), \mathrm{OS}(a)) \\ \quad \quad =\; s \cdot \operatorname{series\text{-}partial\text{-}sum}((\lambda n \in \mathbb{N}.\; /(1, \operatorname{power}(2, (n + 1)))), a) \end{array}}}
\vnbprop{dyadic-eps-sum-le}{\vnbstmt{\begin{array}{@{}l@{}} \forall s. \\ \quad \operatorname{pos\text{-}rr}(s) \Rightarrow (\operatorname{esum}((\lambda i \in \mathbb{N}.\; (s \cdot \operatorname{recip}(\operatorname{power}(2, (i + 1)))))) \leq s) \end{array}}}
\vnbprop{trinum-ge}{\vnbstmt{\forall s \in \mathbb{N}.\; (s \leq \operatorname{trinum}(s))}}
\vnbprop{nn-unpair-le}{\vnbstmt{\forall v \in \mathbb{N}.\; ((\operatorname{nnfst}(v) \leq v) \wedge (\operatorname{nnsnd}(v) \leq v))}}
\vnbprop{rps-finsum-fiber-le}{\vnbstmt{\begin{array}{@{}l@{}} \forall s, t, p, f, g. \\ \quad s \in \mathbf{Set} \wedge \\ \quad |s| \in \mathbb{N} \wedge \\ \quad t \in \mathbf{Set} \wedge \\ \quad |t| \in \mathbb{N} \wedge \\ \quad \forall z \in s.\; (p(z) \in t) \wedge \\ \quad \forall z \in s.\; (f(z) \in \operatorname{rr\text{-}pos\text{-}star}) \wedge \\ \quad \forall z \in t.\; (g(z) \in \operatorname{rr\text{-}pos\text{-}star}) \wedge \\ \quad \forall z \in t.\; (\operatorname{finsum}(\operatorname{rr\text{-}pos\text{-}star\text{-}add\text{-}monoid}, f, \{w \in s : (p(w) = z)\}) \leq g(z)) \Rightarrow \\ \quad \quad \operatorname{finsum}(\operatorname{rr\text{-}pos\text{-}star\text{-}add\text{-}monoid}, f, s) \leq \operatorname{finsum}(\operatorname{rr\text{-}pos\text{-}star\text{-}add\text{-}monoid}, g, t) \end{array}}}
\vnbprop{esum-nn-flatten-le}{\vnbstmt{\begin{array}{@{}l@{}} \forall a, l, b. \\ \quad a \in (\mathbb{N} \to (\mathbb{N} \to \operatorname{rr\text{-}pos\text{-}star})) \wedge \\ \quad l \in (\mathbb{N} \to \operatorname{rr\text{-}pos\text{-}star}) \wedge \\ \quad b \in (\mathbb{N} \to \operatorname{rr\text{-}pos\text{-}star}) \wedge \\ \quad \forall m \in \mathbb{N}.\; (l(m) = a(\operatorname{nnfst}(m))(\operatorname{nnsnd}(m))) \wedge \\ \quad \forall i \in \mathbb{N}.\; (b(i) = \operatorname{esum}(a(i))) \Rightarrow \\ \quad \quad \operatorname{esum}(l) \leq \operatorname{esum}(b) \end{array}}}
\vnbprop{esum-nn-term-le}{\vnbstmt{\forall f \in (\mathbb{N} \to \operatorname{rr\text{-}pos\text{-}star}), k \in \mathbb{N}.\; (f(k) \leq \operatorname{esum}(f))}}
\vnbprop{method-i-outer-empty}{\vnbstmt{\begin{array}{@{}l@{}} \forall a, s, m. \\ \quad \operatorname{is\text{-}premeasure}(a, s, m) \Rightarrow (\operatorname{einf}(\operatorname{cover\text{-}sums}(a, s, m, \emptyset)) = 0) \end{array}}}
\vnbprop{method-i-outer-mono}{\vnbstmt{\begin{array}{@{}l@{}} \forall a, s, m, b, t. \\ \quad ((m \in (s \to \operatorname{rr\text{-}pos\text{-}star})) \wedge (b \subseteq t)) \Rightarrow  \\ \quad \quad \operatorname{einf}(\operatorname{cover\text{-}sums}(a, s, m, b)) \leq \operatorname{einf}(\operatorname{cover\text{-}sums}(a, s, m, t)) \end{array}}}
\vnbprop{method-i-outer-subadditive}{\vnbstmt{\begin{array}{@{}l@{}} \forall a, s, m, f. \\ \quad (\operatorname{is\text{-}premeasure}(a, s, m) \wedge (f \in (\mathbb{N} \to \operatorname{power}(a)))) \Rightarrow  \\ \quad \quad \operatorname{einf}(\operatorname{cover\text{-}sums}(a, s, m, \operatorname{big\text{-}union}(n, \mathbb{N}, f(n)))) \\ \quad \quad \leq\; \operatorname{esum}((\lambda k \in \mathbb{N}.\; \operatorname{einf}(\operatorname{cover\text{-}sums}(a, s, m, f(k))))) \end{array}}}
\vnbprop{method-i-outer-is-outer-measure}{\vnbstmt{\begin{array}{@{}l@{}} \forall a, s, m. \\ \quad \operatorname{is\text{-}premeasure}(a, s, m) \Rightarrow \operatorname{is\text{-}outer\text{-}measure}(a, \operatorname{method\text{-}i\text{-}outer}(a, s, m)) \end{array}}}
\vnbprop{method-i-outer-le-single}{\vnbstmt{\begin{array}{@{}l@{}} \forall a, s, m, b, c. \\ \quad (\operatorname{is\text{-}premeasure}(a, s, m) \wedge ((c \in s) \wedge (b \subseteq c))) \Rightarrow  \\ \quad \quad \operatorname{einf}(\operatorname{cover\text{-}sums}(a, s, m, b)) \leq m(c) \end{array}}}
\vnbprop{measure-is-premeasure}{\vnbstmt{\forall a, b, o.\; (\operatorname{is\text{-}measure}(a, b, o) \Rightarrow \operatorname{is\text{-}premeasure}(a, b, o))}}
\vnbprop{cover-sum-dominates}{\vnbstmt{\begin{array}{@{}l@{}} \forall a, b, o, s, y. \\ \quad (\operatorname{is\text{-}measure}(a, b, o) \wedge (y \in \operatorname{cover\text{-}sums}(a, b, o, s))) \Rightarrow  \\ \quad \quad \exists \mathit{oda} \in b.\; ((s \subseteq \mathit{oda}) \wedge (o(\mathit{oda}) \leq y)) \end{array}}}
\vnbprop{method-i-outer-of-measure}{\vnbstmt{\begin{array}{@{}l@{}} \forall a, b, o, s. \\ \quad (\operatorname{is\text{-}measure}(a, b, o) \wedge (s \in \operatorname{power}(a))) \Rightarrow  \\ \quad \quad \operatorname{method\text{-}i\text{-}outer}(a, b, o)(s) \\ \quad \quad =\; \operatorname{einf}(\{v \in \operatorname{rr\text{-}pos\text{-}star} : \exists \mathit{osa} \in b.\; ((s \subseteq \mathit{osa}) \wedge (v = o(\mathit{osa})))\}) \end{array}}}
\vnbprop{method-i-outer-extends}{\vnbstmt{\begin{array}{@{}l@{}} \forall a, b, o, c. \\ \quad (\operatorname{is\text{-}measure}(a, b, o) \wedge (c \in b)) \Rightarrow (\operatorname{method\text{-}i\text{-}outer}(a, b, o)(c) = o(c)) \end{array}}}
\vnbprop{method-i-outer-measurable}{\vnbstmt{\begin{array}{@{}l@{}} \forall a, b, o, c. \\ \quad (\operatorname{is\text{-}measure}(a, b, o) \wedge (c \in b)) \Rightarrow  \\ \quad \quad \operatorname{is\text{-}caratheodory\text{-}measurable}(a, \operatorname{method\text{-}i\text{-}outer}(a, b, o), c) \end{array}}}
\vnbprop{sigma-algebra-countable-inter}{\vnbstmt{\begin{array}{@{}l@{}} \forall a, b, f. \\ \quad (\operatorname{is\text{-}sigma\text{-}algebra}(a, b) \wedge (f \in (\mathbb{N} \to b))) \Rightarrow  \\ \quad \quad \operatorname{complement\text{-}in}(a, \operatorname{big\text{-}union}(n, \mathbb{N}, \operatorname{complement\text{-}in}(a, f(n)))) \in b \end{array}}}
\vnbprop{method-i-outer-attained}{\vnbstmt{\begin{array}{@{}l@{}} \forall a, b, o, s. \\ \quad (\operatorname{is\text{-}measure}(a, b, o) \wedge ((s \in \operatorname{power}(a)) \wedge (\operatorname{method\text{-}i\text{-}outer}(a, b, o)(s) \in \mathbb{R}))) \Rightarrow  \\ \quad \quad \exists z \in b.\; ((s \subseteq z) \wedge (o(z) = \operatorname{method\text{-}i\text{-}outer}(a, b, o)(s))) \end{array}}}
\vnbprop{method-i-outer-null-iff}{\vnbstmt{\begin{array}{@{}l@{}} \forall a, b, o, s. \\ \quad (\operatorname{is\text{-}measure}(a, b, o) \wedge (s \in \operatorname{power}(a))) \Rightarrow  \\ \quad \quad (\operatorname{method\text{-}i\text{-}outer}(a, b, o)(s) = 0) \Leftrightarrow \operatorname{is\text{-}null\text{-}set}(a, b, o, s) \end{array}}}
\vnbprop{method-i-outer-finite-approx}{\vnbstmt{\begin{array}{@{}l@{}} \forall a, b, o, s. \\ \quad \operatorname{is\text{-}measure}(a, b, o) \wedge \\ \quad \operatorname{is\text{-}caratheodory\text{-}measurable}(a, \operatorname{method\text{-}i\text{-}outer}(a, b, o), s) \wedge \\ \quad \operatorname{method\text{-}i\text{-}outer}(a, b, o)(s) \in \mathbb{R} \Rightarrow \\ \quad \quad \exists y \in b.\; ((s \subseteq y) \wedge \operatorname{is\text{-}null\text{-}set}(a, b, o, \operatorname{complement\text{-}in}(y, s))) \end{array}}}
\vnbprop{caratheodory-measure-complete}{\vnbstmt{\begin{array}{@{}l@{}} \forall a, u. \\ \quad \operatorname{is\text{-}outer\text{-}measure}(a, u) \Rightarrow  \\ \quad \quad \operatorname{is\text{-}complete\text{-}measure}(a, \operatorname{caratheodory\text{-}sets}(a, u), \operatorname{res}(u, \operatorname{caratheodory\text{-}sets}(a, u))) \end{array}}}
\vnbprop{method-i-outer-caratheodory-complete}{\vnbstmt{\begin{array}{@{}l@{}} \forall a, b, o. \\ \quad \operatorname{is\text{-}measure}(a, b, o) \Rightarrow  \\ \quad \quad \operatorname{is\text{-}complete\text{-}measure}(a, \operatorname{caratheodory\text{-}sets}(a, \operatorname{method\text{-}i\text{-}outer}(a, b, o)), \\ \quad \quad \quad \operatorname{res}(\operatorname{method\text{-}i\text{-}outer}(a, b, o), \operatorname{caratheodory\text{-}sets}(a, \operatorname{method\text{-}i\text{-}outer}(a, b, o)))) \end{array}}}
\vnbprop{method-i-outer-extends-completion}{\vnbstmt{\begin{array}{@{}l@{}} \forall a, b, o, c, s. \\ \quad (\operatorname{is\text{-}measure}(a, b, o) \wedge ((c \in b) \wedge \operatorname{is\text{-}null\text{-}set}(a, b, o, s))) \Rightarrow  \\ \quad \quad (c \cup s) \in \operatorname{caratheodory\text{-}sets}(a, \operatorname{method\text{-}i\text{-}outer}(a, b, o)) \wedge \\ \quad \quad \operatorname{res}(\operatorname{method\text{-}i\text{-}outer}(a, b, o), \operatorname{caratheodory\text{-}sets}(a, \operatorname{method\text{-}i\text{-}outer}(a, b, o)))((c \cup s)) = o(c) \end{array}}}
\vnbprop{esum-nn-zero}{\vnbstmt{\begin{array}{@{}l@{}} \forall f \in (\mathbb{N} \to \operatorname{rr\text{-}pos\text{-}star}). \\ \quad \forall n \in \mathbb{N}.\; (f(n) = 0) \Rightarrow (\operatorname{esum}(f) = 0) \end{array}}}
\vnbprop{method-i-outer-sigma-finite-approx}{\vnbstmt{\begin{array}{@{}l@{}} \forall a, b, o, s. \\ \quad \operatorname{is\text{-}measure}(a, b, o) \wedge \\ \quad \operatorname{is\text{-}sigma\text{-}finite}(a, b, o) \wedge \\ \quad \operatorname{is\text{-}caratheodory\text{-}measurable}(a, \operatorname{method\text{-}i\text{-}outer}(a, b, o), s) \Rightarrow \\ \quad \quad \exists y \in b.\; ((s \subseteq y) \wedge \operatorname{is\text{-}null\text{-}set}(a, b, o, \operatorname{complement\text{-}in}(y, s))) \end{array}}}
\vnbprop{method-i-outer-sigma-finite-completion}{\vnbstmt{\begin{array}{@{}l@{}} \forall a, b, o, s. \\ \quad \operatorname{is\text{-}measure}(a, b, o) \wedge \\ \quad \operatorname{is\text{-}sigma\text{-}finite}(a, b, o) \wedge \\ \quad s \in \operatorname{caratheodory\text{-}sets}(a, \operatorname{method\text{-}i\text{-}outer}(a, b, o)) \Rightarrow \\ \quad \quad \exists \mathit{ofa} \in b, \mathit{ofb}.\; (\operatorname{is\text{-}null\text{-}set}(a, b, o, \mathit{ofb}) \wedge (s = (\mathit{ofa} \cup \mathit{ofb}))) \end{array}}}
\vnbprop{lebesgue-premeasure}{\vnbstmt{\operatorname{is\text{-}premeasure}(\mathbb{R}, \operatorname{closed\text{-}intervals}, (\lambda s \in \operatorname{closed\text{-}intervals}.\; \operatorname{interval\text{-}length}(s)))}}
\vnbprop{lebesgue-outer-measure}{\vnbstmt{\begin{array}{@{}l@{}} \operatorname{is\text{-}outer\text{-}measure}(\mathbb{R}, \\ \quad \operatorname{method\text{-}i\text{-}outer}(\mathbb{R}, \operatorname{closed\text{-}intervals}, (\lambda s \in \operatorname{closed\text{-}intervals}.\; \operatorname{interval\text{-}length}(s)))) \end{array}}}
\vnbprop{lebesgue-measure-complete}{\vnbstmt{\begin{array}{@{}l@{}} \operatorname{is\text{-}complete\text{-}measure}(\mathbb{R}, \\ \quad \operatorname{caratheodory\text{-}sets}(\mathbb{R}, \\ \quad \quad \operatorname{method\text{-}i\text{-}outer}(\mathbb{R}, \operatorname{closed\text{-}intervals}, (\lambda s \in \operatorname{closed\text{-}intervals}.\; \operatorname{interval\text{-}length}(s)))), \\ \quad \operatorname{res}(\operatorname{method\text{-}i\text{-}outer}(\mathbb{R}, \operatorname{closed\text{-}intervals}, (\lambda s \in \operatorname{closed\text{-}intervals}.\; \operatorname{interval\text{-}length}(s))), \\ \quad \quad \operatorname{caratheodory\text{-}sets}(\mathbb{R}, \\ \quad \quad \quad \operatorname{method\text{-}i\text{-}outer}(\mathbb{R}, \operatorname{closed\text{-}intervals}, (\lambda s \in \operatorname{closed\text{-}intervals}.\; \operatorname{interval\text{-}length}(s)))))) \end{array}}}
\vnbprop{interval-length-ccint}{\vnbstmt{\begin{array}{@{}l@{}} \forall a, b. \\ \quad ((a \in \mathbb{R}) \wedge ((b \in \mathbb{R}) \wedge (a \leq b))) \Rightarrow  \\ \quad \quad (\lambda s \in \operatorname{closed\text{-}intervals}.\; \operatorname{interval\text{-}length}(s))(\operatorname{ccint}(a, b)) = (b - a) \end{array}}}

\vnbbookfile{structure-library/mb-vocab}{structure-library/mb-vocab}
\noindent{\small mb-vocab.scm -- vocabulary for the M-B wave (construction-of-measures.pdf 2.9, 3.2-3.4, 4.5, 4.13).  DEFINITIONS ONLY.  Agent M-B, 2026-10-26 (new definitions under the rule "missing vocabulary is not a stopping point"). DESTINATION: structure-library/mb-vocab.scm, after structure-library/outer-measure-construction}\par
\par\medskip\noindent\textbf{Definition} (\texttt{is-conull}).\;For all $\omega$, $\mathit{ca}$, $\mu$ and $\mathit{cs}$, 
\begin{equation*}
\operatorname{is\text{-}conull}(\omega, \mathit{ca}, \mu, \mathit{cs}) \Leftrightarrow ((\mathit{cs} \in \mathit{ca}) \wedge (\mu(\operatorname{complement\text{-}in}(\omega, \mathit{cs})) = 0))
\end{equation*}

\par\medskip\noindent\textbf{Definition} (\texttt{sym-difference}).\;$\operatorname{sym\text{-}difference}(\operatorname{sda\_}, \operatorname{sdb\_}) :=$ $(\operatorname{complement\text{-}in}(\operatorname{sda\_}, \operatorname{sdb\_}) \cup \operatorname{complement\text{-}in}(\operatorname{sdb\_}, \operatorname{sda\_}))$
\par\medskip\noindent\textbf{Definition} (\texttt{completion-sigma}).\;$\operatorname{completion\text{-}sigma}(\omega, \mathit{ca}, \mu) :=$
\begin{equation*}
\{\operatorname{cps\_} \in \operatorname{power}(\omega) : \exists \operatorname{cpa\_} \in \mathit{ca}, \operatorname{cpn\_} \in \operatorname{power}(\omega).\; (\operatorname{is\text{-}null\text{-}set}(\omega, \mathit{ca}, \mu, \operatorname{cpn\_}) \wedge (\operatorname{cps\_} = (\operatorname{cpa\_} \cup \operatorname{cpn\_})))\}
\end{equation*}

\par\medskip\noindent\textbf{Definition} (\texttt{completion-measure}).\;$\operatorname{completion\text{-}measure}(\omega, \mathit{ca}, \mu) :=$
\begin{equation*}
\operatorname{res}(\operatorname{method\text{-}i\text{-}outer}(\omega, \mathit{ca}, \mu), \operatorname{completion\text{-}sigma}(\omega, \mathit{ca}, \mu))
\end{equation*}


\vnbbookfile{theorem-library/mb-conull}{theorem-library/mb-conull}
\noindent{\small mb-conull.scm -- construction-of-measures.pdf Prop. 2.9: conull sets are closed under countable intersections, and a (measurable) superset of a conull set is conull. Agent M-B, 2026-10-26.  DESTINATION: theorem-library/mb-conull.scm.  Library: measure.}\par
\vnbprop{conull-countable-inter}{\vnbstmt{\begin{array}{@{}l@{}} \forall a, b, u, f. \\ \quad (\operatorname{is\text{-}measure}(a, b, u) \wedge ((f \in (\mathbb{N} \to b)) \wedge \forall k \in \mathbb{N}.\; \operatorname{is\text{-}conull}(a, b, u, f(k)))) \Rightarrow  \\ \quad \quad \operatorname{is\text{-}conull}(a, b, u, \operatorname{complement\text{-}in}(a, \operatorname{big\text{-}union}(n, \mathbb{N}, \operatorname{complement\text{-}in}(a, f(n))))) \end{array}}}
\vnbprop{conull-superset}{\vnbstmt{\begin{array}{@{}l@{}} \forall a, b, u, \alpha, \beta. \\ \quad (\operatorname{is\text{-}measure}(a, b, u) \wedge (\operatorname{is\text{-}conull}(a, b, u, \alpha) \wedge ((\beta \in b) \wedge (\alpha \subseteq \beta)))) \Rightarrow  \\ \quad \quad \operatorname{is\text{-}conull}(a, b, u, \beta) \end{array}}}

\vnbbookfile{theorem-library/mb-complete}{theorem-library/mb-complete}
\noindent{\small mb-complete.scm -- construction-of-measures.pdf Prop. 3.2: a measure is complete iff "for every A, B subset X, if A (triangle) B is mu-null and A in A then B in A". Agent M-B, 2026-10-26.  DESTINATION: theorem-library/mb-complete.scm.  Library: measure.}\par
\vnbprop{sym-difference-unfold}{\vnbstmt{\begin{array}{@{}l@{}} \forall a, b. \\ \quad (\operatorname{sym\text{-}difference}(a, b) \simeq (\operatorname{complement\text{-}in}(a, b) \cup \operatorname{complement\text{-}in}(b, a))) \end{array}}}
\vnbprop{null-set-subset}{\vnbstmt{\begin{array}{@{}l@{}} \forall a, b, u, \mathit{nsa}, \mathit{nsb}. \\ \quad (\operatorname{is\text{-}null\text{-}set}(a, b, u, \mathit{nsa}) \wedge (\mathit{nsb} \subseteq \mathit{nsa})) \Rightarrow \operatorname{is\text{-}null\text{-}set}(a, b, u, \mathit{nsb}) \end{array}}}
\vnbprop{null-set-in-omega}{\vnbstmt{\begin{array}{@{}l@{}} \forall a, b, u, \mathit{nsa}. \\ \quad (\operatorname{is\text{-}measure}(a, b, u) \wedge \operatorname{is\text{-}null\text{-}set}(a, b, u, \mathit{nsa})) \Rightarrow (\mathit{nsa} \subseteq a) \end{array}}}
\vnbprop{complete-iff-symdiff}{\vnbstmt{\begin{array}{@{}l@{}} \forall a, b, u. \\ \quad \operatorname{is\text{-}measure}(a, b, u) \Rightarrow  \\ \quad \quad \operatorname{is\text{-}complete\text{-}measure}(a, b, u) \\ \quad \quad \Leftrightarrow\; \forall \alpha, \beta. \\ \quad \quad \quad \alpha \subseteq a \wedge \\ \quad \quad \quad \beta \subseteq a \wedge \\ \quad \quad \quad \operatorname{is\text{-}null\text{-}set}(a, b, u, \operatorname{sym\text{-}difference}(\alpha, \beta)) \wedge \\ \quad \quad \quad \alpha \in b \Rightarrow \\ \quad \quad \quad \quad \beta \in b \end{array}}}

\vnbbookfile{theorem-library/mb-outer}{theorem-library/mb-outer}
\noindent{\small mb-outer.scm -- construction-of-measures.pdf Lemma 4.5 (countable form) and Prop. 4.13. Agent M-B, 2026-10-26.  DESTINATION: theorem-library/mb-outer.scm.  Library: measure.}\par
\vnbprop{outer-measure-disjoint-additive}{\vnbstmt{\begin{array}{@{}l@{}} \forall a, u, f, \mathit{ta}. \\ \quad \operatorname{is\text{-}outer\text{-}measure}(a, u) \wedge \\ \quad f \in (\mathbb{N} \to \operatorname{caratheodory\text{-}sets}(a, u)) \wedge \\ \quad \forall i, j. \\ \quad \quad ((i \in \mathbb{N}) \wedge ((j \in \mathbb{N}) \wedge \neg (i = j))) \Rightarrow ((f(i) \cap f(j)) = \emptyset) \wedge \\ \quad \mathit{ta} \in \operatorname{power}(a) \Rightarrow \\ \quad \quad u((\mathit{ta} \cap \operatorname{big\text{-}union}(n, \mathbb{N}, f(n)))) = \operatorname{esum}((\lambda n \in \mathbb{N}.\; u((\mathit{ta} \cap f(n))))) \end{array}}}
\vnbprop{method-i-outer-measurable-iff-finite}{\vnbstmt{\begin{array}{@{}l@{}} \forall a, b, o, s. \\ \quad (\operatorname{is\text{-}measure}(a, b, o) \wedge (s \subseteq a)) \Rightarrow  \\ \quad \quad \operatorname{is\text{-}caratheodory\text{-}measurable}(a, \operatorname{method\text{-}i\text{-}outer}(a, b, o), s) \\ \quad \quad \Leftrightarrow\; \forall \mathit{ta} \in b. \\ \quad \quad \quad (o(\mathit{ta}) \in \mathbb{R}) \Rightarrow  \\ \quad \quad \quad \quad \operatorname{eplus}(\operatorname{method\text{-}i\text{-}outer}(a, b, o)((\mathit{ta} \cap s)), \\ \quad \quad \quad \quad \quad \operatorname{method\text{-}i\text{-}outer}(a, b, o)(\operatorname{complement\text{-}in}(\mathit{ta}, s))) \\ \quad \quad \quad \quad \leq\; o(\mathit{ta}) \end{array}}}

\vnbbookfile{structure-library/real-integral}{structure-library/real-integral}
\noindent{\small real-integral.scm -- THE INTEGRAL OF A REAL-VALUED FUNCTION: positive and negative parts, real measurability, integrability, the real integral. Definitions only; nothing here is asserted and nothing is proved.  The proofs are theorem-library/real-integral.scm (MEASURE-4-C, 2026-10-10).}\par
\par\medskip\noindent\textbf{Definition} (\texttt{pos-part}).\;$\operatorname{pos\text{-}part}(\omega, f) :=$ $(\lambda \operatorname{ppx\_} \in \omega.\; \left(\text{if }(0 \leq f(\operatorname{ppx\_}))\text{ then }f(\operatorname{ppx\_})\text{ else }0\right))$
\par\medskip\noindent\textbf{Definition} (\texttt{neg-part}).\;$\operatorname{neg\text{-}part}(\omega, f) :=$ $(\lambda \operatorname{npx\_} \in \omega.\; \left(\text{if }(0 \leq f(\operatorname{npx\_}))\text{ then }0\text{ else }-f(\operatorname{npx\_})\right))$
\par\medskip\noindent\textbf{Definition} (\texttt{abs-fn}).\;$\operatorname{abs\text{-}fn}(\omega, f) :=$ $(\lambda \operatorname{afx\_} \in \omega.\; \lvert f(\operatorname{afx\_}) \rvert)$
\par\medskip\noindent\textbf{Definition} (\texttt{is-real-measurable-fn}).\;For all $\omega$, $\mathit{ca}$ and $f$, 
\begin{equation*}
\begin{array}{@{}l@{}} \operatorname{is\text{-}real\text{-}measurable\text{-}fn}(\omega, \mathit{ca}, f) \\ \Leftrightarrow\; f \in (\omega \to \mathbb{R}) \wedge \\ \operatorname{is\text{-}measurable\text{-}fn}(\omega, \mathit{ca}, \operatorname{pos\text{-}part}(\omega, f)) \wedge \\ \operatorname{is\text{-}measurable\text{-}fn}(\omega, \mathit{ca}, \operatorname{neg\text{-}part}(\omega, f)) \end{array}
\end{equation*}

\par\medskip\noindent\textbf{Definition} (\texttt{is-integrable}).\;For all $\omega$, $\mathit{ca}$, $\mu$ and $f$, 
\begin{equation*}
\begin{array}{@{}l@{}} \operatorname{is\text{-}integrable}(\omega, \mathit{ca}, \mu, f) \\ \Leftrightarrow\; \operatorname{is\text{-}real\text{-}measurable\text{-}fn}(\omega, \mathit{ca}, f) \wedge \\ \operatorname{integral}(\omega, \mathit{ca}, \mu, \operatorname{pos\text{-}part}(\omega, f)) \in \mathbb{R} \wedge \\ \operatorname{integral}(\omega, \mathit{ca}, \mu, \operatorname{neg\text{-}part}(\omega, f)) \in \mathbb{R} \end{array}
\end{equation*}

\par\medskip\noindent\textbf{Definition} (\texttt{rintegral}).\;$\operatorname{rintegral}(\omega, \mathit{ca}, \mu, f) :=$
\begin{equation*}
(\operatorname{integral}(\omega, \mathit{ca}, \mu, \operatorname{pos\text{-}part}(\omega, f)) - \operatorname{integral}(\omega, \mathit{ca}, \mu, \operatorname{neg\text{-}part}(\omega, f)))
\end{equation*}


\vnbbookfile{structure-library/measure-step}{structure-library/measure-step}
\noindent{\small measure-step.scm -- WAVE 2 (2026-10-27..30), wired 2026-10-30 from scratchpad/md/md-vocab.scm. Library: measure.}\par
\par\medskip\noindent\textbf{Definition} (\texttt{is-measurable-step-fn}).\;For all $\omega$, $\mathit{ca}$ and $s$, 
\begin{equation*}
\begin{array}{@{}l@{}} \operatorname{is\text{-}measurable\text{-}step\text{-}fn}(\omega, \mathit{ca}, s) \\ \Leftrightarrow\; \operatorname{is\text{-}real\text{-}measurable\text{-}fn}(\omega, \mathit{ca}, s) \wedge (|\operatorname{image}(s, \omega)| \in \mathbb{N}) \end{array}
\end{equation*}

\par\medskip\noindent\textbf{Definition} (\texttt{linf-norm}).\;$\operatorname{linf\text{-}norm}(\omega, \mathit{ca}, \mu, f) :=$
\begin{equation*}
\operatorname{einf}(\{\operatorname{lim\_} \in \operatorname{rr\text{-}pos\text{-}star} : (\mu(\{\operatorname{lix\_} \in \omega : (\operatorname{lim\_} < \lvert f(\operatorname{lix\_}) \rvert)\}) = 0)\})
\end{equation*}

\par\medskip\noindent\textbf{Definition} (\texttt{is-measure-isomorphism}).\;For all $\omega$, $\mathit{ca}$, $\mu$, $\omega_{2}$, $\mathit{cb}$, $\nu$ and $\mathit{tt}$, 
\begin{equation*}
\begin{array}{@{}l@{}} \operatorname{is\text{-}measure\text{-}isomorphism}(\omega, \mathit{ca}, \mu, \omega_{2}, \mathit{cb}, \nu, \mathit{tt}) \\ \Leftrightarrow\; \operatorname{is\text{-}measurable\text{-}map}(\omega, \mathit{ca}, \omega_{2}, \mathit{cb}, \mathit{tt}) \wedge \\ \forall \operatorname{mib\_} \in \mathit{cb}.\; (\nu(\operatorname{mib\_}) = \mu(\{\operatorname{mix\_} \in \omega : (\operatorname{tt}(\operatorname{mix\_}) \in \operatorname{mib\_})\})) \wedge \\ \exists \operatorname{mis\_}. \\ \quad \operatorname{is\text{-}measurable\text{-}map}(\omega_{2}, \mathit{cb}, \omega, \mathit{ca}, \operatorname{mis\_}) \wedge \\ \quad \forall \operatorname{mix\_} \in \omega.\; (\operatorname{mis\_}(\operatorname{tt}(\operatorname{mix\_})) = \operatorname{mix\_}) \wedge \\ \quad \forall \operatorname{miy\_} \in \omega_{2}.\; (\operatorname{tt}(\operatorname{mis\_}(\operatorname{miy\_})) = \operatorname{miy\_}) \end{array}
\end{equation*}

\par\medskip\noindent\textbf{Definition} (\texttt{is-absolutely-continuous}).\;For all $\omega$, $\mathit{ca}$, $\mu$ and $\nu$, 
\begin{equation*}
\begin{array}{@{}l@{}} \operatorname{is\text{-}absolutely\text{-}continuous}(\omega, \mathit{ca}, \mu, \nu) \\ \Leftrightarrow\; \forall \operatorname{ace\_} \in \mathit{ca}.\; ((\nu(\operatorname{ace\_}) = 0) \Rightarrow (\mu(\operatorname{ace\_}) = 0)) \end{array}
\end{equation*}


\vnbbookfile{structure-library/l1-dual}{structure-library/l1-dual}
\noindent{\small l1-dual.scm -- WAVE 2 (2026-10-27..30), wired 2026-10-30 from scratchpad/md/me-du-vocab.scm. Library: measure.}\par
\par\medskip\noindent\textbf{Definition} (\texttt{l1-set}).\;$\operatorname{l1\text{-}set}(\omega, \mathit{ca}, \mu) :=$ $\{\operatorname{l1g\_} \in (\omega \to \mathbb{R}) : \operatorname{is\text{-}integrable}(\omega, \mathit{ca}, \mu, \operatorname{l1g\_})\}$
\par\medskip\noindent\textbf{Definition} (\texttt{is-bounded-l1-functional}).\;For all $\omega$, $\mathit{ca}$, $\mu$, $f$ and $\mathit{cb}$, 
\begin{equation*}
\begin{array}{@{}l@{}} \operatorname{is\text{-}bounded\text{-}l1\text{-}functional}(\omega, \mathit{ca}, \mu, f, \mathit{cb}) \\ \Leftrightarrow\; f \in (\operatorname{l1\text{-}set}(\omega, \mathit{ca}, \mu) \to \mathbb{R}) \wedge \\ \mathit{cb} \in \mathbb{R} \wedge \\ 0 \leq \mathit{cb} \wedge \\ \forall \operatorname{bfx\_}, \operatorname{bfy\_} \in \operatorname{l1\text{-}set}(\omega, \mathit{ca}, \mu). \\ \quad f((\lambda \operatorname{sx\_} \in \omega.\; +(\operatorname{bfx\_}(\operatorname{sx\_}), \operatorname{bfy\_}(\operatorname{sx\_})))) = +(f(\operatorname{bfx\_}), f(\operatorname{bfy\_})) \wedge \\ \forall \operatorname{bfr\_} \in \mathbb{R}, \operatorname{bfx\_} \in \operatorname{l1\text{-}set}(\omega, \mathit{ca}, \mu). \\ \quad f((\lambda \operatorname{sx\_} \in \omega.\; \cdot(\operatorname{bfr\_}, \operatorname{bfx\_}(\operatorname{sx\_})))) = \cdot(\operatorname{bfr\_}, f(\operatorname{bfx\_})) \wedge \\ \forall \operatorname{bfx\_} \in \operatorname{l1\text{-}set}(\omega, \mathit{ca}, \mu). \\ \quad \lvert f(\operatorname{bfx\_}) \rvert \leq \cdot(\mathit{cb}, \operatorname{integral}(\omega, \mathit{ca}, \mu, \operatorname{abs\text{-}fn}(\omega, \operatorname{bfx\_}))) \end{array}
\end{equation*}

\par\medskip\noindent\textbf{Definition} (\texttt{l1-pairing}).\;$\operatorname{l1\text{-}pairing}(\omega, \mathit{ca}, \mu, f) :=$
\begin{equation*}
\begin{array}{@{}l@{}} (\lambda \operatorname{lpg\_} \in \operatorname{l1\text{-}set}(\omega, \mathit{ca}, \mu).\; \\ \quad \operatorname{rintegral}(\omega, \mathit{ca}, \mu, (\lambda \operatorname{sx\_} \in \omega.\; (f(\operatorname{sx\_}) \cdot \operatorname{lpg\_}(\operatorname{sx\_}))))) \end{array}
\end{equation*}


\vnbbookfile{structure-library/complex-measure}{structure-library/complex-measure}
\noindent{\small complex-measure.scm -- WAVE 2 (2026-10-27..30), wired 2026-10-30 from scratchpad/mf/mf-cvocab.scm. Library: measure.}\par
\par\medskip\noindent\textbf{Definition} (\texttt{is-finite-measure}).\;For all $\omega$, $\mathit{ca}$ and $\mu$, 
\begin{equation*}
\operatorname{is\text{-}finite\text{-}measure}(\omega, \mathit{ca}, \mu) \Leftrightarrow (\operatorname{is\text{-}measure}(\omega, \mathit{ca}, \mu) \wedge (\mu(\omega) \in \mathbb{R}))
\end{equation*}

\par\medskip\noindent\textbf{Definition} (\texttt{is-complex-measure}).\;For all $\omega$, $\mathit{ca}$ and $\nu$, 
\begin{equation*}
\begin{array}{@{}l@{}} \operatorname{is\text{-}complex\text{-}measure}(\omega, \mathit{ca}, \nu) \\ \Leftrightarrow\; \operatorname{is\text{-}sigma\text{-}algebra}(\omega, \mathit{ca}) \wedge \\ \nu \in (\mathit{ca} \to \mathbb{C}) \wedge \\ \exists \operatorname{cm1\_}, \operatorname{cm2\_}, \operatorname{cm3\_}, \operatorname{cm4\_}. \\ \quad \operatorname{is\text{-}finite\text{-}measure}(\omega, \mathit{ca}, \operatorname{cm1\_}) \wedge \\ \quad \operatorname{is\text{-}finite\text{-}measure}(\omega, \mathit{ca}, \operatorname{cm2\_}) \wedge \\ \quad \operatorname{is\text{-}finite\text{-}measure}(\omega, \mathit{ca}, \operatorname{cm3\_}) \wedge \\ \quad \operatorname{is\text{-}finite\text{-}measure}(\omega, \mathit{ca}, \operatorname{cm4\_}) \wedge \\ \quad \forall \operatorname{cma\_} \in \mathit{ca}. \\ \quad \quad \nu(\operatorname{cma\_}) = +(-(\operatorname{cm1\_}(\operatorname{cma\_}), \operatorname{cm2\_}(\operatorname{cma\_})), \cdot(+i, -(\operatorname{cm3\_}(\operatorname{cma\_}), \operatorname{cm4\_}(\operatorname{cma\_})))) \end{array}
\end{equation*}

\par\medskip\noindent\textbf{Definition} (\texttt{is-complex-measurable-fn}).\;For all $\omega$, $\mathit{ca}$ and $u$, 
\begin{equation*}
\begin{array}{@{}l@{}} \operatorname{is\text{-}complex\text{-}measurable\text{-}fn}(\omega, \mathit{ca}, u) \\ \Leftrightarrow\; u \in (\omega \to \mathbb{C}) \wedge \\ \operatorname{is\text{-}real\text{-}measurable\text{-}fn}(\omega, \mathit{ca}, (\lambda \operatorname{crx\_} \in \omega.\; \operatorname{real\text{-}part}(u(\operatorname{crx\_})))) \wedge \\ \operatorname{is\text{-}real\text{-}measurable\text{-}fn}(\omega, \mathit{ca}, (\lambda \operatorname{crx\_} \in \omega.\; \operatorname{imag\text{-}part}(u(\operatorname{crx\_})))) \end{array}
\end{equation*}

\par\medskip\noindent\textbf{Definition} (\texttt{complex-integral}).\;$\operatorname{complex\text{-}integral}(\omega, \mathit{ca}, \mu, u) :=$
\begin{equation*}
\operatorname{rintegral}(\omega, \mathit{ca}, \mu, (\lambda \operatorname{cix\_} \in \omega.\; \operatorname{real\text{-}part}(u(\operatorname{cix\_})))) + (+i \cdot \operatorname{rintegral}(\omega, \mathit{ca}, \mu, (\lambda \operatorname{cix\_} \in \omega.\; \operatorname{imag\text{-}part}(u(\operatorname{cix\_})))))
\end{equation*}

\par\medskip\noindent\textbf{Definition} (\texttt{is-complex-null-set}).\;For all $\omega$, $\mathit{ca}$, $\nu$ and $a$, 
\begin{equation*}
\begin{array}{@{}l@{}} \operatorname{is\text{-}complex\text{-}null\text{-}set}(\omega, \mathit{ca}, \nu, a) \\ \Leftrightarrow\; (a \in \mathit{ca}) \wedge \forall \operatorname{cnb\_} \in \mathit{ca}.\; ((\operatorname{cnb\_} \subseteq a) \Rightarrow (\nu(\operatorname{cnb\_}) = 0)) \end{array}
\end{equation*}


\vnbbookfile{theorem-library/real-integral}{theorem-library/real-integral}
\noindent{\small real-integral.scm -- THE INTEGRAL OF A REAL-VALUED FUNCTION (MEASURE-4-C, 2026-10-10).  Vocabulary: structure-library/real-integral.scm (POS-PART, NEG-PART, ABS-FN, IS-REAL-MEASURABLE-FN, IS-INTEGRABLE, RINTEGRAL).}\par
\vnbprop{pos-part-unfold}{\vnbstmt{\forall a, f.\; (\operatorname{pos\text{-}part}(a, f) \simeq (\lambda x \in a.\; \left(\text{if }(0 \leq f(x))\text{ then }f(x)\text{ else }0\right)))}}
\vnbprop{neg-part-unfold}{\vnbstmt{\forall a, f.\; (\operatorname{neg\text{-}part}(a, f) \simeq (\lambda x \in a.\; \left(\text{if }(0 \leq f(x))\text{ then }0\text{ else }-f(x)\right)))}}
\vnbprop{abs-fn-unfold}{\vnbstmt{\forall a, f.\; (\operatorname{abs\text{-}fn}(a, f) \simeq (\lambda x \in a.\; \lvert f(x) \rvert))}}
\vnbprop{rintegral-unfold}{\vnbstmt{\begin{array}{@{}l@{}} \forall a, b, u, f. \\ \quad (\operatorname{rintegral}(a, b, u, f) \simeq (\operatorname{integral}(a, b, u, \operatorname{pos\text{-}part}(a, f)) - \operatorname{integral}(a, b, u, \operatorname{neg\text{-}part}(a, f)))) \end{array}}}
\vnbprop{pos-part-value}{\vnbstmt{\forall a, f.; \forall x \in a.\; (\operatorname{pos\text{-}part}(a, f)(x) \simeq \left(\text{if }(0 \leq f(x))\text{ then }f(x)\text{ else }0\right))}}
\vnbprop{neg-part-value}{\vnbstmt{\forall a, f.; \forall x \in a.\; (\operatorname{neg\text{-}part}(a, f)(x) \simeq \left(\text{if }(0 \leq f(x))\text{ then }0\text{ else }-f(x)\right))}}
\vnbprop{abs-fn-value}{\vnbstmt{\forall a, f.; \forall x \in a.\; (\operatorname{abs\text{-}fn}(a, f)(x) \simeq \lvert f(x) \rvert)}}
\vnbprop{ptwise-eplus-value-q}{\vnbstmt{\forall a, b, c.; \forall x \in a.\; (\operatorname{ptwise\text{-}eplus}(a, b, c)(x) \simeq \operatorname{eplus}(b(x), c(x)))}}
\vnbprop{ptwise-etimes-value-q}{\vnbstmt{\forall a, b, c.; \forall x \in a.\; (\operatorname{ptwise\text{-}etimes}(a, b, c)(x) \simeq \operatorname{etimes}(b(x), c(x)))}}
\vnbprop{ptwise-scale-value-q}{\vnbstmt{\forall a, k, g.; \forall x \in a.\; (\operatorname{ptwise\text{-}scale}(a, k, g)(x) \simeq \operatorname{etimes}(k, g(x)))}}
\vnbprop{part-values-at}{\vnbstmt{\begin{array}{@{}l@{}} \forall a.; \forall f \in (a \to \mathbb{R}), x \in a. \\ \quad \operatorname{pos\text{-}part}(a, f)(x) \in \mathbb{R} \wedge \\ \quad 0 \leq \operatorname{pos\text{-}part}(a, f)(x) \wedge \\ \quad \operatorname{neg\text{-}part}(a, f)(x) \in \mathbb{R} \wedge \\ \quad 0 \leq \operatorname{neg\text{-}part}(a, f)(x) \wedge \\ \quad f(x) = (\operatorname{pos\text{-}part}(a, f)(x) - \operatorname{neg\text{-}part}(a, f)(x)) \wedge \\ \quad \lvert f(x) \rvert = (\operatorname{pos\text{-}part}(a, f)(x) + \operatorname{neg\text{-}part}(a, f)(x)) \end{array}}}
\vnbprop{pos-part-fun}{\vnbstmt{\forall a.; \forall f \in (a \to \mathbb{R}).\; (\operatorname{pos\text{-}part}(a, f) \in (a \to \operatorname{rr\text{-}pos\text{-}star}))}}
\vnbprop{neg-part-fun}{\vnbstmt{\forall a.; \forall f \in (a \to \mathbb{R}).\; (\operatorname{neg\text{-}part}(a, f) \in (a \to \operatorname{rr\text{-}pos\text{-}star}))}}
\vnbprop{abs-fn-fun}{\vnbstmt{\forall a.; \forall f \in (a \to \mathbb{R}).\; (\operatorname{abs\text{-}fn}(a, f) \in (a \to \operatorname{rr\text{-}pos\text{-}star}))}}
\vnbprop{part-values-nonneg}{\vnbstmt{\begin{array}{@{}l@{}} \forall a.; \forall f \in (a \to \mathbb{R}), x \in a. \\ \quad (0 \leq f(x)) \Rightarrow ((\operatorname{pos\text{-}part}(a, f)(x) = f(x)) \wedge (\operatorname{neg\text{-}part}(a, f)(x) = 0)) \end{array}}}
\vnbprop{part-values-neg}{\vnbstmt{\begin{array}{@{}l@{}} \forall a.; \forall f \in (a \to \mathbb{R}), x \in a. \\ \quad \neg (0 \leq f(x)) \Rightarrow ((\operatorname{pos\text{-}part}(a, f)(x) = 0) \wedge (\operatorname{neg\text{-}part}(a, f)(x) = -f(x))) \end{array}}}
\vnbprop{pos-part-measurable-of-diff}{\vnbstmt{\begin{array}{@{}l@{}} \forall a, b, f, c, d. \\ \quad f \in (a \to \mathbb{R}) \wedge \\ \quad \operatorname{is\text{-}measurable\text{-}fn}(a, b, c) \wedge \\ \quad \operatorname{is\text{-}measurable\text{-}fn}(a, b, d) \wedge \\ \quad \forall x \in a.\; (c(x) \in \mathbb{R}) \wedge \\ \quad \forall x \in a.\; (d(x) \in \mathbb{R}) \wedge \\ \quad \forall x \in a.\; (f(x) = (c(x) - d(x))) \Rightarrow \\ \quad \quad \operatorname{is\text{-}measurable\text{-}fn}(a, b, \operatorname{pos\text{-}part}(a, f)) \end{array}}}
\vnbprop{neg-part-measurable-of-diff}{\vnbstmt{\begin{array}{@{}l@{}} \forall a, b, f, c, d. \\ \quad f \in (a \to \mathbb{R}) \wedge \\ \quad \operatorname{is\text{-}measurable\text{-}fn}(a, b, c) \wedge \\ \quad \operatorname{is\text{-}measurable\text{-}fn}(a, b, d) \wedge \\ \quad \forall x \in a.\; (c(x) \in \mathbb{R}) \wedge \\ \quad \forall x \in a.\; (d(x) \in \mathbb{R}) \wedge \\ \quad \forall x \in a.\; (f(x) = (c(x) - d(x))) \Rightarrow \\ \quad \quad \operatorname{is\text{-}measurable\text{-}fn}(a, b, \operatorname{neg\text{-}part}(a, f)) \end{array}}}
\vnbprop{rintegral-decomposition}{\vnbstmt{\begin{array}{@{}l@{}} \forall a, b, u, f, c, d. \\ \quad \operatorname{is\text{-}measure}(a, b, u) \wedge \\ \quad f \in (a \to \mathbb{R}) \wedge \\ \quad \operatorname{is\text{-}measurable\text{-}fn}(a, b, c) \wedge \\ \quad \operatorname{is\text{-}measurable\text{-}fn}(a, b, d) \wedge \\ \quad \forall x \in a.\; (c(x) \in \mathbb{R}) \wedge \\ \quad \forall x \in a.\; (d(x) \in \mathbb{R}) \wedge \\ \quad \forall x \in a.\; (f(x) = (c(x) - d(x))) \wedge \\ \quad \operatorname{integral}(a, b, u, c) \in \mathbb{R} \wedge \\ \quad \operatorname{integral}(a, b, u, d) \in \mathbb{R} \Rightarrow \\ \quad \quad (\operatorname{is\text{-}integrable}(a, b, u, f) \wedge (\operatorname{rintegral}(a, b, u, f) = (\operatorname{integral}(a, b, u, c) - \operatorname{integral}(a, b, u, d)))) \end{array}}}
\vnbprop{integral-zero-fn}{\vnbstmt{\begin{array}{@{}l@{}} \forall a, b, u. \\ \quad \operatorname{is\text{-}measure}(a, b, u) \Rightarrow (\operatorname{integral}(a, b, u, (\lambda x \in a.\; 0)) = 0) \end{array}}}
\vnbprop{ptwise-scale-measurable}{\vnbstmt{\begin{array}{@{}l@{}} \forall a, b, g, k. \\ \quad (\operatorname{is\text{-}measurable\text{-}fn}(a, b, g) \wedge ((k \in \mathbb{R}) \wedge (0 \leq k))) \Rightarrow  \\ \quad \quad \operatorname{is\text{-}measurable\text{-}fn}(a, b, \operatorname{ptwise\text{-}scale}(a, k, g)) \end{array}}}
\vnbprop{abs-fn-eq-parts}{\vnbstmt{\begin{array}{@{}l@{}} \forall a.; \forall f \in (a \to \mathbb{R}). \\ \quad \operatorname{abs\text{-}fn}(a, f) = \operatorname{ptwise\text{-}eplus}(a, \operatorname{pos\text{-}part}(a, f), \operatorname{neg\text{-}part}(a, f)) \end{array}}}
\vnbprop{integral-abs-fn}{\vnbstmt{\begin{array}{@{}l@{}} \forall a, b, u, f. \\ \quad (\operatorname{is\text{-}measure}(a, b, u) \wedge \operatorname{is\text{-}real\text{-}measurable\text{-}fn}(a, b, f)) \Rightarrow  \\ \quad \quad \operatorname{is\text{-}measurable\text{-}fn}(a, b, \operatorname{abs\text{-}fn}(a, f)) \wedge \\ \quad \quad \operatorname{integral}(a, b, u, \operatorname{abs\text{-}fn}(a, f)) \\ \quad \quad =\; \operatorname{eplus}(\operatorname{integral}(a, b, u, \operatorname{pos\text{-}part}(a, f)), \operatorname{integral}(a, b, u, \operatorname{neg\text{-}part}(a, f))) \end{array}}}
\vnbprop{integrable-iff-abs-finite}{\vnbstmt{\begin{array}{@{}l@{}} \forall a, b, u, f. \\ \quad (\operatorname{is\text{-}measure}(a, b, u) \wedge \operatorname{is\text{-}real\text{-}measurable\text{-}fn}(a, b, f)) \Rightarrow  \\ \quad \quad \operatorname{is\text{-}integrable}(a, b, u, f) \Leftrightarrow (\operatorname{integral}(a, b, u, \operatorname{abs\text{-}fn}(a, f)) \in \mathbb{R}) \end{array}}}
\vnbprop{rintegral-add}{\vnbstmt{\begin{array}{@{}l@{}} \forall a, b, u, f, g. \\ \quad (\operatorname{is\text{-}measure}(a, b, u) \wedge (\operatorname{is\text{-}integrable}(a, b, u, f) \wedge \operatorname{is\text{-}integrable}(a, b, u, g))) \Rightarrow  \\ \quad \quad \operatorname{is\text{-}integrable}(a, b, u, (\lambda x \in a.\; (f(x) + g(x)))) \wedge \\ \quad \quad \operatorname{rintegral}(a, b, u, (\lambda x \in a.\; (f(x) + g(x)))) = (\operatorname{rintegral}(a, b, u, f) + \operatorname{rintegral}(a, b, u, g)) \end{array}}}
\vnbprop{rintegral-scale}{\vnbstmt{\begin{array}{@{}l@{}} \forall a, b, u, f, c. \\ \quad (\operatorname{is\text{-}measure}(a, b, u) \wedge (\operatorname{is\text{-}integrable}(a, b, u, f) \wedge (c \in \mathbb{R}))) \Rightarrow  \\ \quad \quad \operatorname{is\text{-}integrable}(a, b, u, (\lambda x \in a.\; (c \cdot f(x)))) \wedge \\ \quad \quad \operatorname{rintegral}(a, b, u, (\lambda x \in a.\; (c \cdot f(x)))) = (c \cdot \operatorname{rintegral}(a, b, u, f)) \end{array}}}
\vnbprop{rintegral-mono-off-null}{\vnbstmt{\begin{array}{@{}l@{}} \forall a, b, u, f, g, s. \\ \quad \operatorname{is\text{-}measure}(a, b, u) \wedge \\ \quad \operatorname{is\text{-}integrable}(a, b, u, f) \wedge \\ \quad \operatorname{is\text{-}integrable}(a, b, u, g) \wedge \\ \quad \operatorname{is\text{-}null\text{-}set}(a, b, u, s) \wedge \\ \quad \forall x \in a.\; (\neg (x \in s) \Rightarrow (f(x) \leq g(x))) \Rightarrow \\ \quad \quad \operatorname{rintegral}(a, b, u, f) \leq \operatorname{rintegral}(a, b, u, g) \end{array}}}
\vnbprop{rintegral-monotone}{\vnbstmt{\begin{array}{@{}l@{}} \forall a, b, u, f, g. \\ \quad \operatorname{is\text{-}measure}(a, b, u) \wedge \\ \quad \operatorname{is\text{-}integrable}(a, b, u, f) \wedge \\ \quad \operatorname{is\text{-}integrable}(a, b, u, g) \wedge \\ \quad \forall x \in a.\; (f(x) \leq g(x)) \Rightarrow \\ \quad \quad \operatorname{rintegral}(a, b, u, f) \leq \operatorname{rintegral}(a, b, u, g) \end{array}}}
\vnbprop{rintegral-abs-le}{\vnbstmt{\begin{array}{@{}l@{}} \forall a, b, u, f. \\ \quad (\operatorname{is\text{-}measure}(a, b, u) \wedge \operatorname{is\text{-}integrable}(a, b, u, f)) \Rightarrow  \\ \quad \quad \lvert \operatorname{rintegral}(a, b, u, f) \rvert \leq \operatorname{integral}(a, b, u, \operatorname{abs\text{-}fn}(a, f)) \end{array}}}
\vnbprop{rintegral-of-nonneg}{\vnbstmt{\begin{array}{@{}l@{}} \forall a, b, u, f. \\ \quad (\operatorname{is\text{-}measure}(a, b, u) \wedge (\operatorname{is\text{-}integrable}(a, b, u, f) \wedge \forall x \in a.\; (0 \leq f(x)))) \Rightarrow  \\ \quad \quad \operatorname{is\text{-}measurable\text{-}fn}(a, b, f) \wedge (\operatorname{rintegral}(a, b, u, f) = \operatorname{integral}(a, b, u, f)) \end{array}}}
\vnbprop{rintegral-ae-eq}{\vnbstmt{\begin{array}{@{}l@{}} \forall a, b, u, f, g, s. \\ \quad \operatorname{is\text{-}measure}(a, b, u) \wedge \\ \quad \operatorname{is\text{-}integrable}(a, b, u, f) \wedge \\ \quad \operatorname{is\text{-}real\text{-}measurable\text{-}fn}(a, b, g) \wedge \\ \quad \operatorname{is\text{-}null\text{-}set}(a, b, u, s) \wedge \\ \quad \forall x \in a.\; (\neg (x \in s) \Rightarrow (f(x) = g(x))) \Rightarrow \\ \quad \quad \operatorname{is\text{-}integrable}(a, b, u, g) \wedge (\operatorname{rintegral}(a, b, u, g) = \operatorname{rintegral}(a, b, u, f)) \end{array}}}
\vnbprop{rintegral-restrict}{\vnbstmt{\begin{array}{@{}l@{}} \forall a, b, u, f, c. \\ \quad (\operatorname{is\text{-}measure}(a, b, u) \wedge (\operatorname{is\text{-}integrable}(a, b, u, f) \wedge (c \in b))) \Rightarrow  \\ \quad \quad \operatorname{is\text{-}integrable}(a, b, u, (\lambda x \in a.\; (f(x) \cdot \operatorname{indicator}(a, c)(x)))) \wedge \\ \quad \quad \operatorname{rintegral}(a, b, u, (\lambda x \in a.\; (f(x) \cdot \operatorname{indicator}(a, c)(x)))) \\ \quad \quad =\; (\operatorname{integral}(a, b, u, \operatorname{ptwise\text{-}etimes}(a, \operatorname{pos\text{-}part}(a, f), \operatorname{indicator}(a, c))) - \operatorname{integral}(a, b, u, \operatorname{ptwise\text{-}etimes}(a, \operatorname{neg\text{-}part}(a, f), \operatorname{indicator}(a, c)))) \end{array}}}
\vnbprop{rintegral-nonneg-on-sets-null}{\vnbstmt{\begin{array}{@{}l@{}} \forall a, b, u, f. \\ \quad \operatorname{is\text{-}measure}(a, b, u) \wedge \\ \quad \operatorname{is\text{-}integrable}(a, b, u, f) \wedge \\ \quad \forall c \in b. \\ \quad \quad 0 \leq \operatorname{rintegral}(a, b, u, (\lambda x \in a.\; (f(x) \cdot \operatorname{indicator}(a, c)(x)))) \Rightarrow \\ \quad \quad \operatorname{is\text{-}null\text{-}set}(a, b, u, \{\mathit{nx} \in a : (f(\mathit{nx}) < 0)\}) \end{array}}}
\vnbprop{real-measurable-fn-tails}{\vnbstmt{\begin{array}{@{}l@{}} \forall a, b, f.; \forall l \in \mathbb{R}. \\ \quad \operatorname{is\text{-}real\text{-}measurable\text{-}fn}(a, b, f) \Rightarrow (\{x \in a : (l < f(x))\} \in b) \end{array}}}
\vnbprop{pos-part-measurable-of-tails}{\vnbstmt{\begin{array}{@{}l@{}} \forall a, b, f. \\ \quad (\operatorname{is\text{-}sigma\text{-}algebra}(a, b) \wedge ((f \in (a \to \mathbb{R})) \wedge \forall l \in \mathbb{R}.\; (\{x \in a : (l < f(x))\} \in b))) \Rightarrow  \\ \quad \quad \operatorname{is\text{-}measurable\text{-}fn}(a, b, \operatorname{pos\text{-}part}(a, f)) \end{array}}}
\vnbprop{real-measurable-of-tails}{\vnbstmt{\begin{array}{@{}l@{}} \forall a, b, f. \\ \quad (\operatorname{is\text{-}sigma\text{-}algebra}(a, b) \wedge ((f \in (a \to \mathbb{R})) \wedge \forall l \in \mathbb{R}.\; (\{x \in a : (l < f(x))\} \in b))) \Rightarrow  \\ \quad \quad \operatorname{is\text{-}real\text{-}measurable\text{-}fn}(a, b, f) \end{array}}}
\vnbprop{real-measurable-fn-iff-tails}{\vnbstmt{\begin{array}{@{}l@{}} \forall a, b, f. \\ \quad (\operatorname{is\text{-}sigma\text{-}algebra}(a, b) \wedge (f \in (a \to \mathbb{R}))) \Rightarrow  \\ \quad \quad \operatorname{is\text{-}real\text{-}measurable\text{-}fn}(a, b, f) \Leftrightarrow \forall l \in \mathbb{R}.\; (\{x \in a : (l < f(x))\} \in b) \end{array}}}

\vnbbookfile{structure-library/monotone-class}{structure-library/monotone-class}
\noindent{\small monotone-class.scm -- MONOTONE CLASSES.  Vocabulary only: no proof, no support. Agent MEASURE-2, 2026-10-08.  The proofs (Prop. 1.11, Thm. 1.13) are theorem-library/monotone-class.scm.}\par
\par\medskip\noindent\textbf{Definition} (\texttt{monotone-class-generated}).\;$\operatorname{monotone\text{-}class\text{-}generated}(\omega, \mathit{ce}) :=$
\begin{equation*}
\{\operatorname{mca\_} \in \operatorname{power}(\omega) : \forall \operatorname{mcb\_}.\; (((\operatorname{mcb\_} \in \operatorname{power}(\operatorname{power}(\omega))) \wedge (\operatorname{is\text{-}monotone\text{-}class}(\omega, \operatorname{mcb\_}) \wedge (\mathit{ce} \subseteq \operatorname{mcb\_}))) \Rightarrow (\operatorname{mca\_} \in \operatorname{mcb\_}))\}
\end{equation*}

\par\medskip\noindent\textbf{Definition} (\texttt{is-monotone-class}).\;For all $\omega$ and $\mathit{mc}$, 
\begin{equation*}
\begin{array}{@{}l@{}} \operatorname{is\text{-}monotone\text{-}class}(\omega, \mathit{mc}) \\ \Leftrightarrow\; \omega \in \mathbf{Set} \wedge \\ \mathit{mc} \subseteq \operatorname{power}(\omega) \wedge \\ \forall \operatorname{mcf\_} \in (\mathbb{N} \to \mathit{mc}). \\ \quad \forall \operatorname{mcn\_} \in \mathbb{N}.\; (\operatorname{mcf\_}(\operatorname{mcn\_}) \subseteq \operatorname{mcf\_}(\operatorname{succ}(\operatorname{mcn\_}))) \Rightarrow  \\ \quad \quad \operatorname{big\text{-}union}(\operatorname{mck\_}, \mathbb{N}, \operatorname{mcf\_}(\operatorname{mck\_})) \in \mathit{mc} \wedge \\ \forall \operatorname{mcg\_} \in (\mathbb{N} \to \mathit{mc}). \\ \quad \forall \operatorname{mcm\_} \in \mathbb{N}.\; (\operatorname{mcg\_}(\operatorname{succ}(\operatorname{mcm\_})) \subseteq \operatorname{mcg\_}(\operatorname{mcm\_})) \Rightarrow  \\ \quad \quad \operatorname{complement\text{-}in}(\omega, \operatorname{big\text{-}union}(\operatorname{mcj\_}, \mathbb{N}, \operatorname{complement\text{-}in}(\omega, \operatorname{mcg\_}(\operatorname{mcj\_})))) \in \mathit{mc} \end{array}
\end{equation*}


\vnbbookfile{theorem-library/monotone-class}{theorem-library/monotone-class}
\noindent{\small monotone-class.scm -- MONOTONE CLASSES: Prop. 1.11 and the monotone class theorem (Thm. 1.13).  Agent MEASURE-2, 2026-10-08 (scratchpad/triage/MEASURE-2.md).  Library: measure.  Helper prefix mcl-. 12 theorems, every one `modulo 0', no oracle.  Vocabulary:}\par
\vnbprop{algebra-complement-closed}{\vnbstmt{\begin{array}{@{}l@{}} \forall a, b, \alpha. \\ \quad (\operatorname{is\text{-}algebra\text{-}of\text{-}sets}(a, b) \wedge (\alpha \in b)) \Rightarrow (\operatorname{complement\text{-}in}(a, \alpha) \in b) \end{array}}}
\vnbprop{algebra-union-closed}{\vnbstmt{\begin{array}{@{}l@{}} \forall a, b, \alpha, \beta. \\ \quad (\operatorname{is\text{-}algebra\text{-}of\text{-}sets}(a, b) \wedge ((\alpha \in b) \wedge (\beta \in b))) \Rightarrow  \\ \quad \quad (\alpha \cup \beta) \in b \end{array}}}
\vnbprop{algebra-inter-closed}{\vnbstmt{\begin{array}{@{}l@{}} \forall a, b, \alpha, \beta. \\ \quad (\operatorname{is\text{-}algebra\text{-}of\text{-}sets}(a, b) \wedge ((\alpha \in b) \wedge (\beta \in b))) \Rightarrow  \\ \quad \quad (\alpha \cap \beta) \in b \end{array}}}
\vnbprop{algebra-whole-empty-in}{\vnbstmt{\forall a, b.\; (\operatorname{is\text{-}algebra\text{-}of\text{-}sets}(a, b) \Rightarrow ((a \in b) \wedge (\emptyset \in b)))}}
\vnbprop{algebra-seg-union-in}{\vnbstmt{\begin{array}{@{}l@{}} \forall k \in \mathbb{N}, a, b, f. \\ \quad (\operatorname{is\text{-}algebra\text{-}of\text{-}sets}(a, b) \wedge (f \in (\mathbb{N} \to b))) \Rightarrow  \\ \quad \quad \operatorname{big\text{-}union}(j, \mathrm{OS}(k), f(j)) \in b \end{array}}}
\vnbprop{algebra-monotone-union-sigma-algebra}{\vnbstmt{\begin{array}{@{}l@{}} \forall a, b. \\ \quad \operatorname{is\text{-}algebra\text{-}of\text{-}sets}(a, b) \wedge \\ \quad \forall f \in (\mathbb{N} \to b). \\ \quad \quad \forall n \in \mathbb{N}.\; (f(n) \subseteq f((n + 1))) \Rightarrow (\operatorname{big\text{-}union}(k, \mathbb{N}, f(k)) \in b) \Rightarrow \\ \quad \quad \operatorname{is\text{-}sigma\text{-}algebra}(a, b) \end{array}}}
\vnbprop{monotone-class-generated-unfold}{\vnbstmt{\begin{array}{@{}l@{}} \forall a, e. \\ \quad (\operatorname{monotone\text{-}class\text{-}generated}(a, e) \simeq \{\mathit{mca} \in \operatorname{power}(a) : \forall b \in \operatorname{power}(\operatorname{power}(a)).\; ((\operatorname{is\text{-}monotone\text{-}class}(a, b) \wedge (e \subseteq b)) \Rightarrow (\mathit{mca} \in b))\}) \end{array}}}
\vnbprop{monotone-class-in-power-power}{\vnbstmt{\forall a, c.\; (\operatorname{is\text{-}monotone\text{-}class}(a, c) \Rightarrow (c \in \operatorname{power}(\operatorname{power}(a))))}}
\vnbprop{monotone-class-generated-membership}{\vnbstmt{\begin{array}{@{}l@{}} \forall a, e, \alpha. \\ \quad \alpha \in \operatorname{monotone\text{-}class\text{-}generated}(a, e) \\ \quad \Leftrightarrow\; ((\alpha \in \operatorname{power}(a)) \wedge \forall b.\; ((\operatorname{is\text{-}monotone\text{-}class}(a, b) \wedge (e \subseteq b)) \Rightarrow (\alpha \in b))) \end{array}}}
\vnbprop{monotone-class-generated-is-monotone-class}{\vnbstmt{\begin{array}{@{}l@{}} \forall a, e. \\ \quad ((a \in \mathbf{Set}) \wedge (e \subseteq \operatorname{power}(a))) \Rightarrow  \\ \quad \quad \operatorname{is\text{-}monotone\text{-}class}(a, \operatorname{monotone\text{-}class\text{-}generated}(a, e)) \end{array}}}
\vnbprop{monotone-class-union-section}{\vnbstmt{\begin{array}{@{}l@{}} \forall a, m, d. \\ \quad (\operatorname{is\text{-}monotone\text{-}class}(a, m) \wedge (d \in \operatorname{power}(a))) \Rightarrow  \\ \quad \quad \operatorname{is\text{-}monotone\text{-}class}(a, \{e \in \operatorname{power}(a) : ((d \cup e) \in m)\}) \end{array}}}
\vnbprop{monotone-class-complement-section}{\vnbstmt{\begin{array}{@{}l@{}} \forall a, m. \\ \quad \operatorname{is\text{-}monotone\text{-}class}(a, m) \Rightarrow  \\ \quad \quad \operatorname{is\text{-}monotone\text{-}class}(a, \{c \in \operatorname{power}(a) : (\operatorname{complement\text{-}in}(a, c) \in m)\}) \end{array}}}
\vnbprop{monotone-class-theorem}{\vnbstmt{\begin{array}{@{}l@{}} \forall a, c, e. \\ \quad (\operatorname{is\text{-}monotone\text{-}class}(a, c) \wedge (\operatorname{is\text{-}algebra\text{-}of\text{-}sets}(a, e) \wedge (e \subseteq c))) \Rightarrow  \\ \quad \quad \operatorname{sigma\text{-}generated}(a, e) \subseteq c \end{array}}}

\vnbbookfile{theorem-library/lebesgue-interval-length}{theorem-library/lebesgue-interval-length}
\noindent{\small lebesgue-interval-length.scm -- THE LEBESGUE OUTER MEASURE OF AN INTERVAL IS ITS LENGTH.  Agent MEASURE-3-B, 2026-10-09 (scratchpad/triage/MEASURE-3-B.md). Library: measure.  Helper prefix lil-.  No definition, no support, no stamp.}\par
\vnbprop{interval-length-unfold}{\vnbstmt{\begin{array}{@{}l@{}} \forall s. \\ \quad (\operatorname{interval\text{-}length}(s) \simeq \operatorname{esup}(\{v \in \operatorname{rr\text{-}pos\text{-}star} : ((v = 0) \vee \exists x, y \in s.\; ((x \leq y) \wedge (v = (y - x))))\})) \end{array}}}
\vnbprop{interval-length-type}{\vnbstmt{\forall s.\; (\operatorname{interval\text{-}length}(s) \in \operatorname{rr\text{-}pos\text{-}star})}}
\vnbprop{interval-length-mono}{\vnbstmt{\forall a, b.\; ((a \subseteq b) \Rightarrow (\operatorname{interval\text{-}length}(a) \leq \operatorname{interval\text{-}length}(b)))}}
\vnbprop{interval-length-extend}{\vnbstmt{\begin{array}{@{}l@{}} \forall a, b, w, z. \\ \quad b \subseteq \mathbb{R} \wedge \\ \quad a \subseteq b \wedge \\ \quad w \in \mathbb{R} \wedge \\ \quad z \in \mathbb{R} \wedge \\ \quad w \leq z \wedge \\ \quad w \in b \wedge \\ \quad z \in b \wedge \\ \quad \forall x \in a.\; (x \leq w) \Rightarrow \\ \quad \quad \operatorname{eplus}(\operatorname{interval\text{-}length}(a), (z - w)) \leq \operatorname{interval\text{-}length}(b) \end{array}}}
\vnbprop{interval-length-widen}{\vnbstmt{\begin{array}{@{}l@{}} \forall s, r. \\ \quad ((s \subseteq \mathbb{R}) \wedge ((r \in \mathbb{R}) \wedge (0 < r))) \Rightarrow  \\ \quad \quad \operatorname{interval\text{-}length}(\{z \in \mathbb{R} : \exists u \in s.\; (((z - u) < r) \wedge ((u - z) < r))\}) \\ \quad \quad \leq\; \operatorname{eplus}(\operatorname{interval\text{-}length}(s), \operatorname{eplus}(r, r)) \end{array}}}
\vnbprop{rps-prefix-mono-range}{\vnbstmt{\begin{array}{@{}l@{}} \forall f, m, n. \\ \quad ((f \in (\mathbb{N} \to \operatorname{rr\text{-}pos\text{-}star})) \wedge ((m \in \mathbb{N}) \wedge ((n \in \mathbb{N}) \wedge (m \leq n)))) \Rightarrow  \\ \quad \quad \operatorname{finsum}(\operatorname{rr\text{-}pos\text{-}star\text{-}add\text{-}monoid}, f, \mathrm{OS}(m)) \\ \quad \quad \leq\; \operatorname{finsum}(\operatorname{rr\text{-}pos\text{-}star\text{-}add\text{-}monoid}, f, \mathrm{OS}(n)) \end{array}}}
\vnbprop{rps-prefix-bump}{\vnbstmt{\begin{array}{@{}l@{}} \forall n \in \mathbb{N}, k, f, g, c. \\ \quad k \in \mathbb{N} \wedge \\ \quad f \in (\mathbb{N} \to \operatorname{rr\text{-}pos\text{-}star}) \wedge \\ \quad g \in (\mathbb{N} \to \operatorname{rr\text{-}pos\text{-}star}) \wedge \\ \quad c \in \operatorname{rr\text{-}pos\text{-}star} \wedge \\ \quad \forall q \in \mathbb{N}.\; (f(q) \leq g(q)) \wedge \\ \quad \operatorname{eplus}(f(k), c) \leq g(k) \wedge \\ \quad k < n \Rightarrow \\ \quad \quad \operatorname{eplus}(\operatorname{finsum}(\operatorname{rr\text{-}pos\text{-}star\text{-}add\text{-}monoid}, f, \mathrm{OS}(n)), c) \\ \quad \quad \leq\; \operatorname{finsum}(\operatorname{rr\text{-}pos\text{-}star\text{-}add\text{-}monoid}, g, \mathrm{OS}(n)) \end{array}}}
\vnbprop{rr-near-center-half}{\vnbstmt{\begin{array}{@{}l@{}} \forall t, d, r, w. \\ \quad ((t \in \mathbb{R}) \wedge ((d \in \mathbb{R}) \wedge ((r \in \mathbb{R}) \wedge ((w \in \mathbb{R}) \wedge ((0 < d) \wedge (((d + d) = r) \wedge (((t - d) < w) \wedge (w \leq (t + d))))))))) \Rightarrow  \\ \quad \quad ((w - t) < r) \wedge ((t - w) < r) \end{array}}}
\vnbprop{ccint-cover-widened-sum}{\vnbstmt{\begin{array}{@{}l@{}} \forall a, b, s, e. \\ \quad a \in \mathbb{R} \wedge \\ \quad b \in \mathbb{R} \wedge \\ \quad a \leq b \wedge \\ \quad s \in (\mathbb{N} \to \operatorname{power}(\mathbb{R})) \wedge \\ \quad \operatorname{ccint}(a, b) \subseteq \operatorname{big\text{-}union}(n, \mathbb{N}, s(n)) \wedge \\ \quad \operatorname{pos\text{-}rr}(e) \Rightarrow \\ \quad \quad (b - a) \\ \quad \quad \leq\; \operatorname{esum}((\lambda \mathit{lnn} \in \mathbb{N}.\; \\ \quad \quad \quad \quad \operatorname{interval\text{-}length}(\{z \in \mathbb{R} : \exists u \in s(\mathit{lnn}).\; (((z - u) < (e \cdot \operatorname{recip}(\operatorname{power}(2, (\mathit{lnn} + 1))))) \wedge ((u - z) < (e \cdot \operatorname{recip}(\operatorname{power}(2, (\mathit{lnn} + 1))))))\}))) \end{array}}}
\vnbprop{ccint-cover-sum-ge}{\vnbstmt{\begin{array}{@{}l@{}} \forall a, b, s. \\ \quad a \in \mathbb{R} \wedge \\ \quad b \in \mathbb{R} \wedge \\ \quad a \leq b \wedge \\ \quad s \in (\mathbb{N} \to \operatorname{power}(\mathbb{R})) \wedge \\ \quad \operatorname{ccint}(a, b) \subseteq \operatorname{big\text{-}union}(n, \mathbb{N}, s(n)) \Rightarrow \\ \quad \quad (b - a) \leq \operatorname{esum}((\lambda \mathit{lkn} \in \mathbb{N}.\; \operatorname{interval\text{-}length}(s(\mathit{lkn})))) \end{array}}}
\vnbprop{lebesgue-outer-mono}{\vnbstmt{\begin{array}{@{}l@{}} \forall a, b. \\ \quad ((a \subseteq \mathbb{R}) \wedge ((b \subseteq \mathbb{R}) \wedge (a \subseteq b))) \Rightarrow  \\ \quad \quad \operatorname{method\text{-}i\text{-}outer}(\mathbb{R}, \operatorname{closed\text{-}intervals}, (\lambda s \in \operatorname{closed\text{-}intervals}.\; \operatorname{interval\text{-}length}(s)))(a) \\ \quad \quad \leq\; \operatorname{method\text{-}i\text{-}outer}(\mathbb{R}, \operatorname{closed\text{-}intervals}, (\lambda s \in \operatorname{closed\text{-}intervals}.\; \operatorname{interval\text{-}length}(s)))(b) \end{array}}}
\vnbprop{lebesgue-outer-ccint}{\vnbstmt{\begin{array}{@{}l@{}} \forall a, b. \\ \quad ((a \in \mathbb{R}) \wedge ((b \in \mathbb{R}) \wedge (a \leq b))) \Rightarrow  \\ \quad \quad \operatorname{method\text{-}i\text{-}outer}(\mathbb{R}, \operatorname{closed\text{-}intervals}, (\lambda s \in \operatorname{closed\text{-}intervals}.\; \operatorname{interval\text{-}length}(s)))(\operatorname{ccint}(a, b)) \\ \quad \quad =\; (b - a) \end{array}}}
\vnbprop{lebesgue-outer-point}{\vnbstmt{\begin{array}{@{}l@{}} \forall x \in \mathbb{R}. \\ \quad \operatorname{method\text{-}i\text{-}outer}(\mathbb{R}, \operatorname{closed\text{-}intervals}, (\lambda s \in \operatorname{closed\text{-}intervals}.\; \operatorname{interval\text{-}length}(s)))(\operatorname{singleton}(x)) \\ \quad =\; 0 \end{array}}}
\vnbprop{lebesgue-outer-empty}{\vnbstmt{\begin{array}{@{}l@{}} \operatorname{method\text{-}i\text{-}outer}(\mathbb{R}, \operatorname{closed\text{-}intervals}, (\lambda s \in \operatorname{closed\text{-}intervals}.\; \operatorname{interval\text{-}length}(s)))(\emptyset) \\ =\; 0 \end{array}}}
\vnbprop{lebesgue-outer-between}{\vnbstmt{\begin{array}{@{}l@{}} \forall a, b, s. \\ \quad a \in \mathbb{R} \wedge \\ \quad b \in \mathbb{R} \wedge \\ \quad a \leq b \wedge \\ \quad s \subseteq \operatorname{ccint}(a, b) \wedge \\ \quad \operatorname{ccint}(a, b) \subseteq ((s \cup \operatorname{singleton}(a)) \cup \operatorname{singleton}(b)) \Rightarrow \\ \quad \quad \operatorname{method\text{-}i\text{-}outer}(\mathbb{R}, \operatorname{closed\text{-}intervals}, (\lambda \mathit{ols} \in \operatorname{closed\text{-}intervals}.\; \operatorname{interval\text{-}length}(\mathit{ols})))(s) \\ \quad \quad =\; (b - a) \end{array}}}
\vnbprop{lebesgue-outer-ooint}{\vnbstmt{\begin{array}{@{}l@{}} \forall a, b. \\ \quad ((a \in \mathbb{R}) \wedge ((b \in \mathbb{R}) \wedge (a \leq b))) \Rightarrow  \\ \quad \quad \operatorname{method\text{-}i\text{-}outer}(\mathbb{R}, \operatorname{closed\text{-}intervals}, (\lambda s \in \operatorname{closed\text{-}intervals}.\; \operatorname{interval\text{-}length}(s)))(\operatorname{ooint}(a, b)) \\ \quad \quad =\; (b - a) \end{array}}}
\vnbprop{lebesgue-outer-coint}{\vnbstmt{\begin{array}{@{}l@{}} \forall a, b. \\ \quad ((a \in \mathbb{R}) \wedge ((b \in \mathbb{R}) \wedge (a \leq b))) \Rightarrow  \\ \quad \quad \operatorname{method\text{-}i\text{-}outer}(\mathbb{R}, \operatorname{closed\text{-}intervals}, (\lambda s \in \operatorname{closed\text{-}intervals}.\; \operatorname{interval\text{-}length}(s)))(\operatorname{coint}(a, b)) \\ \quad \quad =\; (b - a) \end{array}}}
\vnbprop{lebesgue-outer-ocint}{\vnbstmt{\begin{array}{@{}l@{}} \forall a, b. \\ \quad ((a \in \mathbb{R}) \wedge ((b \in \mathbb{R}) \wedge (a \leq b))) \Rightarrow  \\ \quad \quad \operatorname{method\text{-}i\text{-}outer}(\mathbb{R}, \operatorname{closed\text{-}intervals}, (\lambda s \in \operatorname{closed\text{-}intervals}.\; \operatorname{interval\text{-}length}(s)))(\operatorname{ocint}(a, b)) \\ \quad \quad =\; (b - a) \end{array}}}
\vnbprop{interval-length-translate-le}{\vnbstmt{\begin{array}{@{}l@{}} \forall s.; \forall c \in \mathbb{R}. \\ \quad \operatorname{interval\text{-}length}(\{z \in \mathbb{R} : ((z - c) \in s)\}) \leq \operatorname{interval\text{-}length}(s) \end{array}}}
\vnbprop{translate-closed-interval}{\vnbstmt{\begin{array}{@{}l@{}} \forall i, c. \\ \quad ((i \in \operatorname{closed\text{-}intervals}) \wedge (c \in \mathbb{R})) \Rightarrow  \\ \quad \quad \{z \in \mathbb{R} : ((z - c) \in i)\} \in \operatorname{closed\text{-}intervals} \end{array}}}
\vnbprop{lebesgue-outer-translate-le}{\vnbstmt{\begin{array}{@{}l@{}} \forall s, c. \\ \quad ((s \subseteq \mathbb{R}) \wedge (c \in \mathbb{R})) \Rightarrow  \\ \quad \quad \operatorname{method\text{-}i\text{-}outer}(\mathbb{R}, \operatorname{closed\text{-}intervals}, (\lambda \mathit{ols} \in \operatorname{closed\text{-}intervals}.\; \operatorname{interval\text{-}length}(\mathit{ols})))(\{z \in \mathbb{R} : ((z - c) \in s)\}) \\ \quad \quad \leq\; \operatorname{method\text{-}i\text{-}outer}(\mathbb{R}, \operatorname{closed\text{-}intervals}, (\lambda \mathit{ols} \in \operatorname{closed\text{-}intervals}.\; \operatorname{interval\text{-}length}(\mathit{ols})))(s) \end{array}}}
\vnbprop{lebesgue-outer-translate}{\vnbstmt{\begin{array}{@{}l@{}} \forall s, c. \\ \quad ((s \subseteq \mathbb{R}) \wedge (c \in \mathbb{R})) \Rightarrow  \\ \quad \quad \operatorname{method\text{-}i\text{-}outer}(\mathbb{R}, \operatorname{closed\text{-}intervals}, (\lambda \mathit{ols} \in \operatorname{closed\text{-}intervals}.\; \operatorname{interval\text{-}length}(\mathit{ols})))(\{z \in \mathbb{R} : ((z - c) \in s)\}) \\ \quad \quad =\; \operatorname{method\text{-}i\text{-}outer}(\mathbb{R}, \operatorname{closed\text{-}intervals}, (\lambda \mathit{ols} \in \operatorname{closed\text{-}intervals}.\; \operatorname{interval\text{-}length}(\mathit{ols})))(s) \end{array}}}

\vnbbookfile{structure-library/metric-outer-measure}{structure-library/metric-outer-measure}
\noindent{\small metric-outer-measure.scm -- the distance between sets, METRIC OUTER MEASURES and the BOREL SETS of a metric space.  Vocabulary only: no proof, no support. Agent MEASURE-3-A, 2026-10-09 (scratchpad/triage/MEASURE-3-A.md).  The proofs are theorem-library/metric-outer-measure.scm.}\par
\par\medskip\noindent\textbf{Definition} (\texttt{pt-set-dist}).\;$\operatorname{pt\text{-}set\text{-}dist}(m, x, \mathit{es}) :=$
\begin{equation*}
\operatorname{einf}(\{\operatorname{mov\_} \in \operatorname{rr\text{-}pos\text{-}star} : \exists \operatorname{moz\_} \in \mathit{es}.\; (\operatorname{mov\_} = \operatorname{dist}(m)(x, \operatorname{moz\_}))\})
\end{equation*}

\par\medskip\noindent\textbf{Definition} (\texttt{set-dist}).\;$\operatorname{set\text{-}dist}(m, \mathit{aa}, \mathit{bb}) :=$
\begin{equation*}
\operatorname{einf}(\{\operatorname{mow\_} \in \operatorname{rr\text{-}pos\text{-}star} : \exists \operatorname{mox\_} \in \mathit{aa}, \operatorname{moy\_} \in \mathit{bb}.\; (\operatorname{mow\_} = \operatorname{dist}(m)(\operatorname{mox\_}, \operatorname{moy\_}))\})
\end{equation*}

\par\medskip\noindent\textbf{Definition} (\texttt{is-metric-outer-measure}).\;For all $m$ and $\mu$, 
\begin{equation*}
\begin{array}{@{}l@{}} \operatorname{is\text{-}metric\text{-}outer\text{-}measure}(m, \mu) \\ \Leftrightarrow\; \operatorname{is\text{-}metric\text{-}space}(m) \wedge \\ \operatorname{is\text{-}outer\text{-}measure}(\operatorname{pts}(m), \mu) \wedge \\ \forall \operatorname{moa\_}, \operatorname{mob\_} \in \operatorname{power}(\operatorname{pts}(m)). \\ \quad \neg (\operatorname{set\text{-}dist}(m, \operatorname{moa\_}, \operatorname{mob\_}) = 0) \Rightarrow  \\ \quad \quad \mu((\operatorname{moa\_} \cup \operatorname{mob\_})) = \operatorname{eplus}(\mu(\operatorname{moa\_}), \mu(\operatorname{mob\_})) \end{array}
\end{equation*}

\par\medskip\noindent\textbf{Definition} (\texttt{borel-sets}).\;$\operatorname{borel\text{-}sets}(m) :=$
\begin{equation*}
\operatorname{sigma\text{-}generated}(\operatorname{pts}(m), \{\operatorname{mou\_} \in \operatorname{power}(\operatorname{pts}(m)) : \operatorname{is\text{-}open}(m, \operatorname{mou\_})\})
\end{equation*}


\vnbbookfile{structure-library/step-maps}{structure-library/step-maps}
\noindent{\small step-maps.scm -- WAVE 2 (2026-10-27..30), wired 2026-10-30 from scratchpad/mf/mf-mcvocab.scm. Library: measure.}\par
\par\medskip\noindent\textbf{Definition} (\texttt{is-step-map}).\;For all $\omega$, $\mathit{ca}$, $s$ and $g$, 
\begin{equation*}
\begin{array}{@{}l@{}} \operatorname{is\text{-}step\text{-}map}(\omega, \mathit{ca}, s, g) \\ \Leftrightarrow\; \operatorname{is\text{-}measurable\text{-}map}(\omega, \mathit{ca}, \operatorname{pts}(s), \operatorname{borel\text{-}sets}(s), g) \wedge \\ \exists \operatorname{stv\_} \in (\mathbb{N} \to \operatorname{pts}(s)). \forall \operatorname{stx\_} \in \omega. \exists \operatorname{stk\_} \in \mathbb{N}. \\ \quad g(\operatorname{stx\_}) = \operatorname{stv\_}(\operatorname{stk\_}) \end{array}
\end{equation*}


\vnbbookfile{theorem-library/metric-outer-measure}{theorem-library/metric-outer-measure}
\noindent{\small metric-outer-measure.scm -- METRIC OUTER MEASURES: separated sets, the layers of a closed set, Borel sets are measurable; Lebesgue outer measure on the line is metric. Agent MEASURE-3-A, 2026-10-09 (scratchpad/triage/MEASURE-3-A.md).  Library: measure. Helper prefix mom-.  58 theorems, every one `modulo 0' (oracles ...}\par
\vnbprop{pt-set-dist-unfold}{\vnbstmt{\begin{array}{@{}l@{}} \forall m, x, s. \\ \quad (\operatorname{pt\text{-}set\text{-}dist}(m, x, s) \simeq \operatorname{einf}(\{v \in \operatorname{rr\text{-}pos\text{-}star} : \exists z \in s.\; (v = \operatorname{dist}(m)(x, z))\})) \end{array}}}
\vnbprop{set-dist-unfold}{\vnbstmt{\begin{array}{@{}l@{}} \forall m, a, b. \\ \quad (\operatorname{set\text{-}dist}(m, a, b) \simeq \operatorname{einf}(\{w \in \operatorname{rr\text{-}pos\text{-}star} : \exists x \in a, y \in b.\; (w = \operatorname{dist}(m)(x, y))\})) \end{array}}}
\vnbprop{pt-set-dist-ge-intro}{\vnbstmt{\begin{array}{@{}l@{}} \forall m, x, s.; \forall r \in \operatorname{rr\text{-}pos\text{-}star}. \\ \quad \forall z \in s.\; (r \leq \operatorname{dist}(m)(x, z)) \Rightarrow (r \leq \operatorname{pt\text{-}set\text{-}dist}(m, x, s)) \end{array}}}
\vnbprop{pt-set-dist-ge-elim}{\vnbstmt{\begin{array}{@{}l@{}} \forall m, x, s, r.; \forall z \in s. \\ \quad \operatorname{is\text{-}metric\text{-}space}(m) \wedge \\ \quad x \in \operatorname{pts}(m) \wedge \\ \quad s \subseteq \operatorname{pts}(m) \wedge \\ \quad r \in \operatorname{rr\text{-}pos\text{-}star} \wedge \\ \quad r \leq \operatorname{pt\text{-}set\text{-}dist}(m, x, s) \Rightarrow \\ \quad \quad r \leq \operatorname{dist}(m)(x, z) \end{array}}}
\vnbprop{set-dist-pos-intro}{\vnbstmt{\begin{array}{@{}l@{}} \forall m, a, b, r. \\ \quad (\operatorname{pos\text{-}rr}(r) \wedge \forall x \in a, y \in b.\; (r \leq \operatorname{dist}(m)(x, y))) \Rightarrow  \\ \quad \quad \neg (\operatorname{set\text{-}dist}(m, a, b) = 0) \end{array}}}
\vnbprop{set-dist-pos-elim}{\vnbstmt{\begin{array}{@{}l@{}} \forall m, a, b. \\ \quad (\operatorname{is\text{-}metric\text{-}space}(m) \wedge ((a \subseteq \operatorname{pts}(m)) \wedge ((b \subseteq \operatorname{pts}(m)) \wedge \neg (\operatorname{set\text{-}dist}(m, a, b) = 0)))) \Rightarrow  \\ \quad \quad \exists r.\; (\operatorname{pos\text{-}rr}(r) \wedge \forall x \in a, y \in b.\; (r \leq \operatorname{dist}(m)(x, y))) \end{array}}}
\vnbprop{set-dist-union-pos}{\vnbstmt{\begin{array}{@{}l@{}} \forall m, a, b, s. \\ \quad \operatorname{is\text{-}metric\text{-}space}(m) \wedge \\ \quad a \subseteq \operatorname{pts}(m) \wedge \\ \quad b \subseteq \operatorname{pts}(m) \wedge \\ \quad s \subseteq \operatorname{pts}(m) \wedge \\ \quad \neg (\operatorname{set\text{-}dist}(m, a, s) = 0) \wedge \\ \quad \neg (\operatorname{set\text{-}dist}(m, b, s) = 0) \Rightarrow \\ \quad \quad \neg (\operatorname{set\text{-}dist}(m, (a \cup b), s) = 0) \end{array}}}
\vnbprop{metric-outer-measure-prefix-separated}{\vnbstmt{\begin{array}{@{}l@{}} \forall n \in \mathbb{N}, m, u, s. \\ \quad \operatorname{is\text{-}metric\text{-}outer\text{-}measure}(m, u) \wedge \\ \quad s \in (\mathbb{N} \to \operatorname{power}(\operatorname{pts}(m))) \wedge \\ \quad \forall k, l \in \mathbb{N}.\; (\neg (k = l) \Rightarrow \neg (\operatorname{set\text{-}dist}(m, s(k), s(l)) = 0)) \Rightarrow \\ \quad \quad u(\operatorname{big\text{-}union}(j, \mathrm{OS}(n), s(j))) \\ \quad \quad =\; \operatorname{finsum}(\operatorname{rr\text{-}pos\text{-}star\text{-}add\text{-}monoid}, (\lambda a \in \mathbb{N}.\; u(s(a))), \mathrm{OS}(n)) \wedge \\ \quad \quad \forall \mathit{moj} \in \mathbb{N}. \\ \quad \quad \quad (n \leq \mathit{moj}) \Rightarrow  \\ \quad \quad \quad \quad \neg (\operatorname{set\text{-}dist}(m, \operatorname{big\text{-}union}(j, \mathrm{OS}(n), s(j)), s(\mathit{moj})) = 0) \end{array}}}
\vnbprop{metric-outer-measure-finite-separated}{\vnbstmt{\begin{array}{@{}l@{}} \forall m, u, s.; \forall n \in \mathbb{N}. \\ \quad \operatorname{is\text{-}metric\text{-}outer\text{-}measure}(m, u) \wedge \\ \quad s \in (\mathbb{N} \to \operatorname{power}(\operatorname{pts}(m))) \wedge \\ \quad \forall k, l \in \mathbb{N}.\; (\neg (k = l) \Rightarrow \neg (\operatorname{set\text{-}dist}(m, s(k), s(l)) = 0)) \Rightarrow \\ \quad \quad u(\operatorname{big\text{-}union}(j, \mathrm{OS}(n), s(j))) \\ \quad \quad =\; \operatorname{finsum}(\operatorname{rr\text{-}pos\text{-}star\text{-}add\text{-}monoid}, (\lambda a \in \mathbb{N}.\; u(s(a))), \mathrm{OS}(n)) \end{array}}}
\vnbprop{metric-outer-measure-separated-sum}{\vnbstmt{\begin{array}{@{}l@{}} \forall m, u, s. \\ \quad \operatorname{is\text{-}metric\text{-}outer\text{-}measure}(m, u) \wedge \\ \quad s \in (\mathbb{N} \to \operatorname{power}(\operatorname{pts}(m))) \wedge \\ \quad \forall k, l \in \mathbb{N}.\; (\neg (k = l) \Rightarrow \neg (\operatorname{set\text{-}dist}(m, s(k), s(l)) = 0)) \Rightarrow \\ \quad \quad u(\operatorname{big\text{-}union}(n, \mathbb{N}, s(n))) = \operatorname{esum}((\lambda n \in \mathbb{N}.\; u(s(n)))) \end{array}}}
\vnbprop{rr-mul-succ-distrib}{\vnbstmt{\forall a, c \in \mathbb{R}.\; ((c \cdot (a + 1)) = ((a \cdot c) + c))}}
\vnbprop{nn-recip-succ-gap}{\vnbstmt{\forall j \in \mathbb{N}.\; \operatorname{pos\text{-}rr}((\operatorname{recip}((j + 1)) - \operatorname{recip}(((j + 1) + 1))))}}
\vnbprop{dist-layer-mem-elim}{\vnbstmt{\begin{array}{@{}l@{}} \forall m, l.; \forall n \in \mathbb{N}, x \in \{\mathit{mex} \in \operatorname{pts}(m) : (\operatorname{recip}((n + 1)) \leq \operatorname{pt\text{-}set\text{-}dist}(m, \mathit{mex}, l))\}. \\ \quad (\operatorname{is\text{-}metric\text{-}space}(m) \wedge (l \subseteq \operatorname{pts}(m))) \Rightarrow  \\ \quad \quad (x \in \operatorname{pts}(m)) \wedge \forall z \in l.\; (\operatorname{recip}((n + 1)) \leq \operatorname{dist}(m)(x, z)) \end{array}}}
\vnbprop{dist-layer-mem-intro}{\vnbstmt{\begin{array}{@{}l@{}} \forall m, l.; \forall n \in \mathbb{N}, x \in \operatorname{pts}(m). \\ \quad \forall z \in l.\; (\operatorname{recip}((n + 1)) \leq \operatorname{dist}(m)(x, z)) \Rightarrow  \\ \quad \quad x \in \{\mathit{mex} \in \operatorname{pts}(m) : (\operatorname{recip}((n + 1)) \leq \operatorname{pt\text{-}set\text{-}dist}(m, \mathit{mex}, l))\} \end{array}}}
\vnbprop{dist-layer-mono}{\vnbstmt{\begin{array}{@{}l@{}} \forall m, l, k, n. \\ \quad (\operatorname{is\text{-}metric\text{-}space}(m) \wedge ((l \subseteq \operatorname{pts}(m)) \wedge ((k \in \mathbb{N}) \wedge ((n \in \mathbb{N}) \wedge (k \leq n))))) \Rightarrow  \\ \quad \quad \{x \in \operatorname{pts}(m) : (\operatorname{recip}((k + 1)) \leq \operatorname{pt\text{-}set\text{-}dist}(m, x, l))\} \\ \quad \quad \subseteq\; \{x \in \operatorname{pts}(m) : (\operatorname{recip}((n + 1)) \leq \operatorname{pt\text{-}set\text{-}dist}(m, x, l))\} \end{array}}}
\vnbprop{dist-layer-disjoint}{\vnbstmt{\begin{array}{@{}l@{}} \forall m, l.; \forall n \in \mathbb{N}, x \in \{\mathit{mex} \in \operatorname{pts}(m) : (\operatorname{recip}((n + 1)) \leq \operatorname{pt\text{-}set\text{-}dist}(m, \mathit{mex}, l))\}. \\ \quad (\operatorname{is\text{-}metric\text{-}space}(m) \wedge (l \subseteq \operatorname{pts}(m))) \Rightarrow \neg (x \in l) \end{array}}}
\vnbprop{dist-layer-cover}{\vnbstmt{\begin{array}{@{}l@{}} \forall m, l.; \forall x \in \operatorname{pts}(m). \\ \quad (\operatorname{is\text{-}closed}(m, l) \wedge \neg (x \in l)) \Rightarrow  \\ \quad \quad \exists k \in \mathbb{N}.\; (x \in \{\mathit{mex} \in \operatorname{pts}(m) : (\operatorname{recip}((k + 1)) \leq \operatorname{pt\text{-}set\text{-}dist}(m, \mathit{mex}, l))\}) \end{array}}}
\vnbprop{dist-layer-gap}{\vnbstmt{\begin{array}{@{}l@{}} \forall m, l.; \forall j \in \mathbb{N}, x, y. \\ \quad \operatorname{is\text{-}metric\text{-}space}(m) \wedge \\ \quad l \subseteq \operatorname{pts}(m) \wedge \\ \quad x \in \{\mathit{mex} \in \operatorname{pts}(m) : (\operatorname{recip}((j + 1)) \leq \operatorname{pt\text{-}set\text{-}dist}(m, \mathit{mex}, l))\} \wedge \\ \quad y \in \operatorname{pts}(m) \wedge \\ \quad \neg (y \in \{\mathit{mex} \in \operatorname{pts}(m) : (\operatorname{recip}(((j + 1) + 1)) \leq \operatorname{pt\text{-}set\text{-}dist}(m, \mathit{mex}, l))\}) \Rightarrow \\ \quad \quad (\operatorname{recip}((j + 1)) - \operatorname{recip}(((j + 1) + 1))) \leq \operatorname{dist}(m)(x, y) \end{array}}}
\vnbprop{rps-prefix-mono}{\vnbstmt{\begin{array}{@{}l@{}} \forall k \in \mathbb{N}, f \in (\mathbb{N} \to \operatorname{rr\text{-}pos\text{-}star}), n \in \mathbb{N}. \\ \quad (n \leq k) \Rightarrow  \\ \quad \quad \operatorname{finsum}(\operatorname{rr\text{-}pos\text{-}star\text{-}add\text{-}monoid}, f, \mathrm{OS}(n)) \\ \quad \quad \leq\; \operatorname{finsum}(\operatorname{rr\text{-}pos\text{-}star\text{-}add\text{-}monoid}, f, \mathrm{OS}(k)) \end{array}}}
\vnbprop{rps-prefix-even-odd}{\vnbstmt{\begin{array}{@{}l@{}} \forall n \in \mathbb{N}, f, g, h. \\ \quad f \in (\mathbb{N} \to \operatorname{rr\text{-}pos\text{-}star}) \wedge \\ \quad g \in (\mathbb{N} \to \operatorname{rr\text{-}pos\text{-}star}) \wedge \\ \quad h \in (\mathbb{N} \to \operatorname{rr\text{-}pos\text{-}star}) \wedge \\ \quad \forall q \in \mathbb{N}.\; (g(q) = f((2 \cdot q))) \wedge \\ \quad \forall q \in \mathbb{N}.\; (h(q) = f(((2 \cdot q) + 1))) \Rightarrow \\ \quad \quad \operatorname{finsum}(\operatorname{rr\text{-}pos\text{-}star\text{-}add\text{-}monoid}, f, \mathrm{OS}((2 \cdot n))) \\ \quad \quad =\; \operatorname{eplus}(\operatorname{finsum}(\operatorname{rr\text{-}pos\text{-}star\text{-}add\text{-}monoid}, g, \mathrm{OS}(n)), \\ \quad \quad \quad \operatorname{finsum}(\operatorname{rr\text{-}pos\text{-}star\text{-}add\text{-}monoid}, h, \mathrm{OS}(n))) \end{array}}}
\vnbprop{esum-nn-even-odd-le}{\vnbstmt{\begin{array}{@{}l@{}} \forall f, g, h. \\ \quad f \in (\mathbb{N} \to \operatorname{rr\text{-}pos\text{-}star}) \wedge \\ \quad g \in (\mathbb{N} \to \operatorname{rr\text{-}pos\text{-}star}) \wedge \\ \quad h \in (\mathbb{N} \to \operatorname{rr\text{-}pos\text{-}star}) \wedge \\ \quad \forall q \in \mathbb{N}.\; (g(q) = f((2 \cdot q))) \wedge \\ \quad \forall q \in \mathbb{N}.\; (h(q) = f(((2 \cdot q) + 1))) \Rightarrow \\ \quad \quad \operatorname{esum}(f) \leq \operatorname{eplus}(\operatorname{esum}(g), \operatorname{esum}(h)) \end{array}}}
\vnbprop{rps-prefix-shift-le}{\vnbstmt{\begin{array}{@{}l@{}} \forall k, n \in \mathbb{N}, f, g. \\ \quad ((f \in (\mathbb{N} \to \operatorname{rr\text{-}pos\text{-}star})) \wedge ((g \in (\mathbb{N} \to \operatorname{rr\text{-}pos\text{-}star})) \wedge \forall j \in \mathbb{N}.\; (g(j) \leq f((n + j))))) \Rightarrow  \\ \quad \quad \operatorname{eplus}(\operatorname{finsum}(\operatorname{rr\text{-}pos\text{-}star\text{-}add\text{-}monoid}, f, \mathrm{OS}(n)), \\ \quad \quad \quad \operatorname{finsum}(\operatorname{rr\text{-}pos\text{-}star\text{-}add\text{-}monoid}, g, \mathrm{OS}(k))) \\ \quad \quad \leq\; \operatorname{finsum}(\operatorname{rr\text{-}pos\text{-}star\text{-}add\text{-}monoid}, f, \mathrm{OS}((n + k))) \end{array}}}
\vnbprop{esum-nn-tail-small}{\vnbstmt{\begin{array}{@{}l@{}} \forall f \in (\mathbb{N} \to \operatorname{rr\text{-}pos\text{-}star}), s. \\ \quad ((\operatorname{esum}(f) \in \mathbb{R}) \wedge \operatorname{pos\text{-}rr}(s)) \Rightarrow  \\ \quad \quad \exists n \in \mathbb{N}. \forall g \in (\mathbb{N} \to \operatorname{rr\text{-}pos\text{-}star}). \\ \quad \quad \quad \forall j \in \mathbb{N}.\; (g(j) \leq f((n + j))) \Rightarrow (\operatorname{esum}(g) \leq s) \end{array}}}
\vnbprop{dist-layer-ring-cover}{\vnbstmt{\begin{array}{@{}l@{}} \forall m, l.; \forall n \in \mathbb{N}, x \in \operatorname{pts}(m). \\ \quad \operatorname{is\text{-}closed}(m, l) \wedge \\ \quad \neg (x \in l) \wedge \\ \quad \neg (x \in \{\mathit{mex} \in \operatorname{pts}(m) : (\operatorname{recip}((n + 1)) \leq \operatorname{pt\text{-}set\text{-}dist}(m, \mathit{mex}, l))\}) \Rightarrow \\ \quad \quad \exists j \in \mathbb{N}. \\ \quad \quad \quad x \\ \quad \quad \quad \in\; \operatorname{complement\text{-}in}(\{\mathit{mex} \in \operatorname{pts}(m) : (\operatorname{recip}((((n + j) + 1) + 1)) \leq \operatorname{pt\text{-}set\text{-}dist}(m, \mathit{mex}, l))\}, \\ \quad \quad \quad \quad \{\mathit{mex} \in \operatorname{pts}(m) : (\operatorname{recip}(((n + j) + 1)) \leq \operatorname{pt\text{-}set\text{-}dist}(m, \mathit{mex}, l))\}) \end{array}}}
\vnbprop{dist-layer-set-sep}{\vnbstmt{\begin{array}{@{}l@{}} \forall m, l.; \forall j \in \mathbb{N}, a, b. \\ \quad \operatorname{is\text{-}metric\text{-}space}(m) \wedge \\ \quad l \subseteq \operatorname{pts}(m) \wedge \\ \quad a \subseteq \{x \in \operatorname{pts}(m) : (\operatorname{recip}((j + 1)) \leq \operatorname{pt\text{-}set\text{-}dist}(m, x, l))\} \wedge \\ \quad \forall y \in b. \\ \quad \quad ((y \in \operatorname{pts}(m)) \wedge \neg (y \in \{x \in \operatorname{pts}(m) : (\operatorname{recip}(((j + 1) + 1)) \leq \operatorname{pt\text{-}set\text{-}dist}(m, x, l))\})) \Rightarrow \\ \quad \quad \neg (\operatorname{set\text{-}dist}(m, a, b) = 0) \wedge \neg (\operatorname{set\text{-}dist}(m, b, a) = 0) \end{array}}}
\vnbprop{dist-ring-even-sum-le}{\vnbstmt{\begin{array}{@{}l@{}} \forall m, u, l.; \forall a \in \operatorname{power}(\operatorname{pts}(m)). \\ \quad (\operatorname{is\text{-}metric\text{-}outer\text{-}measure}(m, u) \wedge \operatorname{is\text{-}closed}(m, l)) \Rightarrow  \\ \quad \quad \operatorname{esum}((\lambda q \in \mathbb{N}.\; \\ \quad \quad \quad \quad u((a \cap \operatorname{complement\text{-}in}(\{x \in \operatorname{pts}(m) : (\operatorname{recip}((((2 \cdot q) + 1) + 1)) \leq \operatorname{pt\text{-}set\text{-}dist}(m, x, l))\}, \{x \in \operatorname{pts}(m) : (\operatorname{recip}(((2 \cdot q) + 1)) \leq \operatorname{pt\text{-}set\text{-}dist}(m, x, l))\}))))) \\ \quad \quad \leq\; \operatorname{esup}(\operatorname{image}((\lambda k \in \mathbb{N}.\; u((a \cap \{x \in \operatorname{pts}(m) : (\operatorname{recip}((k + 1)) \leq \operatorname{pt\text{-}set\text{-}dist}(m, x, l))\}))), \\ \quad \quad \quad \quad \mathbb{N})) \end{array}}}
\vnbprop{dist-ring-odd-sum-le}{\vnbstmt{\begin{array}{@{}l@{}} \forall m, u, l.; \forall a \in \operatorname{power}(\operatorname{pts}(m)). \\ \quad (\operatorname{is\text{-}metric\text{-}outer\text{-}measure}(m, u) \wedge \operatorname{is\text{-}closed}(m, l)) \Rightarrow  \\ \quad \quad \operatorname{esum}((\lambda q \in \mathbb{N}.\; \\ \quad \quad \quad \quad u((a \cap \operatorname{complement\text{-}in}(\{x \in \operatorname{pts}(m) : (\operatorname{recip}(((((2 \cdot q) + 1) + 1) + 1)) \leq \operatorname{pt\text{-}set\text{-}dist}(m, x, l))\}, \{x \in \operatorname{pts}(m) : (\operatorname{recip}((((2 \cdot q) + 1) + 1)) \leq \operatorname{pt\text{-}set\text{-}dist}(m, x, l))\}))))) \\ \quad \quad \leq\; \operatorname{esup}(\operatorname{image}((\lambda k \in \mathbb{N}.\; u((a \cap \{x \in \operatorname{pts}(m) : (\operatorname{recip}((k + 1)) \leq \operatorname{pt\text{-}set\text{-}dist}(m, x, l))\}))), \\ \quad \quad \quad \quad \mathbb{N})) \end{array}}}
\vnbprop{metric-outer-layers-limit}{\vnbstmt{\begin{array}{@{}l@{}} \forall m, u, l. \\ \quad (\operatorname{is\text{-}metric\text{-}outer\text{-}measure}(m, u) \wedge \operatorname{is\text{-}closed}(m, l)) \Rightarrow  \\ \quad \quad \forall n \in \mathbb{N}. \\ \quad \quad \quad \{x \in \operatorname{pts}(m) : (\operatorname{recip}((n + 1)) \leq \operatorname{pt\text{-}set\text{-}dist}(m, x, l))\} \\ \quad \quad \quad \subseteq\; \{x \in \operatorname{pts}(m) : (\operatorname{recip}(((n + 1) + 1)) \leq \operatorname{pt\text{-}set\text{-}dist}(m, x, l))\} \wedge \\ \quad \quad \forall a \in \operatorname{power}(\operatorname{pts}(m)). \\ \quad \quad \quad u((a \cap \operatorname{complement\text{-}in}(\operatorname{pts}(m), l))) \\ \quad \quad \quad =\; \operatorname{esup}(\operatorname{image}((\lambda k \in \mathbb{N}.\; u((a \cap \{x \in \operatorname{pts}(m) : (\operatorname{recip}((k + 1)) \leq \operatorname{pt\text{-}set\text{-}dist}(m, x, l))\}))), \\ \quad \quad \quad \quad \quad \mathbb{N})) \end{array}}}
\vnbprop{metric-outer-closed-measurable}{\vnbstmt{\begin{array}{@{}l@{}} \forall m, u, l. \\ \quad (\operatorname{is\text{-}metric\text{-}outer\text{-}measure}(m, u) \wedge \operatorname{is\text{-}closed}(m, l)) \Rightarrow  \\ \quad \quad \operatorname{is\text{-}caratheodory\text{-}measurable}(\operatorname{pts}(m), u, l) \end{array}}}
\vnbprop{metric-outer-open-measurable}{\vnbstmt{\begin{array}{@{}l@{}} \forall m, u, a. \\ \quad (\operatorname{is\text{-}metric\text{-}outer\text{-}measure}(m, u) \wedge \operatorname{is\text{-}open}(m, a)) \Rightarrow  \\ \quad \quad \operatorname{is\text{-}caratheodory\text{-}measurable}(\operatorname{pts}(m), u, a) \end{array}}}
\vnbprop{borel-sets-unfold}{\vnbstmt{\begin{array}{@{}l@{}} \forall m. \\ \quad (\operatorname{borel\text{-}sets}(m) \simeq \operatorname{sigma\text{-}generated}(\operatorname{pts}(m), \{u \in \operatorname{power}(\operatorname{pts}(m)) : \operatorname{is\text{-}open}(m, u)\})) \end{array}}}
\vnbprop{borel-sets-caratheodory-measurable}{\vnbstmt{\begin{array}{@{}l@{}} \forall m, u.; \forall s \in \operatorname{borel\text{-}sets}(m). \\ \quad \operatorname{is\text{-}metric\text{-}outer\text{-}measure}(m, u) \Rightarrow \operatorname{is\text{-}caratheodory\text{-}measurable}(\operatorname{pts}(m), u, s) \end{array}}}
\vnbprop{ccint-inter-left-ray}{\vnbstmt{\begin{array}{@{}l@{}} \forall a, b, c. \\ \quad ((a \in \mathbb{R}) \wedge ((b \in \mathbb{R}) \wedge ((c \in \mathbb{R}) \wedge (c \leq b)))) \Rightarrow  \\ \quad \quad (\operatorname{ccint}(a, b) \cap \{x \in \mathbb{R} : (x \leq c)\}) = \operatorname{ccint}(a, c) \end{array}}}
\vnbprop{ccint-inter-left-ray-all}{\vnbstmt{\begin{array}{@{}l@{}} \forall a, b, c. \\ \quad ((a \in \mathbb{R}) \wedge ((b \in \mathbb{R}) \wedge ((c \in \mathbb{R}) \wedge (b \leq c)))) \Rightarrow  \\ \quad \quad (\operatorname{ccint}(a, b) \cap \{x \in \mathbb{R} : (x \leq c)\}) = \operatorname{ccint}(a, b) \end{array}}}
\vnbprop{ccint-inter-right-ray}{\vnbstmt{\begin{array}{@{}l@{}} \forall a, b, c. \\ \quad ((a \in \mathbb{R}) \wedge ((b \in \mathbb{R}) \wedge ((c \in \mathbb{R}) \wedge (a \leq c)))) \Rightarrow  \\ \quad \quad (\operatorname{ccint}(a, b) \cap \{x \in \mathbb{R} : (c \leq x)\}) = \operatorname{ccint}(c, b) \end{array}}}
\vnbprop{ccint-inter-right-ray-all}{\vnbstmt{\begin{array}{@{}l@{}} \forall a, b, c. \\ \quad ((a \in \mathbb{R}) \wedge ((b \in \mathbb{R}) \wedge ((c \in \mathbb{R}) \wedge (c \leq a)))) \Rightarrow  \\ \quad \quad (\operatorname{ccint}(a, b) \cap \{x \in \mathbb{R} : (c \leq x)\}) = \operatorname{ccint}(a, b) \end{array}}}
\vnbprop{ccint-empty-of-lt}{\vnbstmt{\begin{array}{@{}l@{}} \forall a, b, c. \\ \quad ((a \in \mathbb{R}) \wedge ((b \in \mathbb{R}) \wedge ((c \in \mathbb{R}) \wedge (b < a)))) \Rightarrow (\operatorname{ccint}(a, b) = \emptyset) \end{array}}}
\vnbprop{empty-inter-left-ray}{\vnbstmt{\begin{array}{@{}l@{}} \forall a, b, c. \\ \quad ((a \in \mathbb{R}) \wedge ((b \in \mathbb{R}) \wedge (c \in \mathbb{R}))) \Rightarrow  \\ \quad \quad (\emptyset \cap \{x \in \mathbb{R} : (x \leq c)\}) = \emptyset \end{array}}}
\vnbprop{empty-inter-right-ray}{\vnbstmt{\begin{array}{@{}l@{}} \forall a, b, c. \\ \quad ((a \in \mathbb{R}) \wedge ((b \in \mathbb{R}) \wedge (c \in \mathbb{R}))) \Rightarrow  \\ \quad \quad (\emptyset \cap \{x \in \mathbb{R} : (c \leq x)\}) = \emptyset \end{array}}}
\vnbprop{ccint-in-closed-intervals}{\vnbstmt{\begin{array}{@{}l@{}} \forall a, b. \\ \quad ((a \in \mathbb{R}) \wedge ((b \in \mathbb{R}) \wedge (a \leq b))) \Rightarrow (\operatorname{ccint}(a, b) \in \operatorname{closed\text{-}intervals}) \end{array}}}
\vnbprop{lebesgue-cut}{\vnbstmt{\begin{array}{@{}l@{}} \forall s, c. \\ \quad ((s \in \operatorname{closed\text{-}intervals}) \wedge (c \in \mathbb{R})) \Rightarrow  \\ \quad \quad (s \cap \{x \in \mathbb{R} : (x \leq c)\}) \in \operatorname{closed\text{-}intervals} \wedge \\ \quad \quad (s \cap \{x \in \mathbb{R} : (c \leq x)\}) \in \operatorname{closed\text{-}intervals} \wedge \\ \quad \quad \operatorname{eplus}((\lambda \mathit{ols} \in \operatorname{closed\text{-}intervals}.\; \operatorname{interval\text{-}length}(\mathit{ols}))((s \cap \{x \in \mathbb{R} : (x \leq c)\})), \\ \quad \quad \quad (\lambda \mathit{ols} \in \operatorname{closed\text{-}intervals}.\; \operatorname{interval\text{-}length}(\mathit{ols}))((s \cap \{x \in \mathbb{R} : (c \leq x)\}))) \\ \quad \quad \leq\; (\lambda \mathit{ols} \in \operatorname{closed\text{-}intervals}.\; \operatorname{interval\text{-}length}(\mathit{ols}))(s) \end{array}}}
\vnbprop{esum-nn-superadditive}{\vnbstmt{\begin{array}{@{}l@{}} \forall f, g, h. \\ \quad f \in (\mathbb{N} \to \operatorname{rr\text{-}pos\text{-}star}) \wedge \\ \quad g \in (\mathbb{N} \to \operatorname{rr\text{-}pos\text{-}star}) \wedge \\ \quad h \in (\mathbb{N} \to \operatorname{rr\text{-}pos\text{-}star}) \wedge \\ \quad \forall k \in \mathbb{N}.\; (\operatorname{eplus}(f(k), g(k)) \leq h(k)) \Rightarrow \\ \quad \quad \operatorname{eplus}(\operatorname{esum}(f), \operatorname{esum}(g)) \leq \operatorname{esum}(h) \end{array}}}
\vnbprop{lebesgue-split-le}{\vnbstmt{\begin{array}{@{}l@{}} \forall s, c, u, v. \\ \quad s \in \operatorname{power}(\mathbb{R}) \wedge \\ \quad c \in \mathbb{R} \wedge \\ \quad u \subseteq (s \cap \{x \in \mathbb{R} : (x \leq c)\}) \wedge \\ \quad v \subseteq (s \cap \{x \in \mathbb{R} : (c \leq x)\}) \Rightarrow \\ \quad \quad \operatorname{eplus}(\operatorname{method\text{-}i\text{-}outer}(\mathbb{R}, \operatorname{closed\text{-}intervals}, (\lambda \mathit{ols} \in \operatorname{closed\text{-}intervals}.\; \operatorname{interval\text{-}length}(\mathit{ols})))(u), \\ \quad \quad \quad \operatorname{method\text{-}i\text{-}outer}(\mathbb{R}, \operatorname{closed\text{-}intervals}, (\lambda \mathit{ols} \in \operatorname{closed\text{-}intervals}.\; \operatorname{interval\text{-}length}(\mathit{ols})))(v)) \\ \quad \quad \leq\; \operatorname{method\text{-}i\text{-}outer}(\mathbb{R}, \operatorname{closed\text{-}intervals}, (\lambda \mathit{ols} \in \operatorname{closed\text{-}intervals}.\; \operatorname{interval\text{-}length}(\mathit{ols})))(s) \end{array}}}
\vnbprop{lebesgue-left-ray-measurable}{\vnbstmt{\begin{array}{@{}l@{}} \forall c \in \mathbb{R}. \\ \quad \operatorname{is\text{-}caratheodory\text{-}measurable}(\mathbb{R}, \\ \quad \quad \operatorname{method\text{-}i\text{-}outer}(\mathbb{R}, \operatorname{closed\text{-}intervals}, (\lambda s \in \operatorname{closed\text{-}intervals}.\; \operatorname{interval\text{-}length}(s))), \\ \quad \quad \{x \in \mathbb{R} : (x \leq c)\}) \end{array}}}
\vnbprop{lebesgue-right-ray-measurable}{\vnbstmt{\begin{array}{@{}l@{}} \forall c \in \mathbb{R}. \\ \quad \operatorname{is\text{-}caratheodory\text{-}measurable}(\mathbb{R}, \\ \quad \quad \operatorname{method\text{-}i\text{-}outer}(\mathbb{R}, \operatorname{closed\text{-}intervals}, (\lambda s \in \operatorname{closed\text{-}intervals}.\; \operatorname{interval\text{-}length}(s))), \\ \quad \quad \{x \in \mathbb{R} : (c \leq x)\}) \end{array}}}
\vnbprop{sigma-algebra-sep-meets}{\vnbstmt{\begin{array}{@{}l@{}} \forall a, b, s, c. \\ \quad (\operatorname{is\text{-}sigma\text{-}algebra}(a, b) \wedge (s \in b)) \Rightarrow  \\ \quad \quad \{x \in s : \exists y \in c.\; (y \in s)\} \in b \end{array}}}
\vnbprop{metric-outer-of-separating-sets}{\vnbstmt{\begin{array}{@{}l@{}} \forall m, u. \\ \quad \operatorname{is\text{-}metric\text{-}space}(m) \wedge \\ \quad \operatorname{is\text{-}outer\text{-}measure}(\operatorname{pts}(m), u) \wedge \\ \quad \forall a, b \in \operatorname{power}(\operatorname{pts}(m)). \\ \quad \quad \neg (\operatorname{set\text{-}dist}(m, a, b) = 0) \Rightarrow  \\ \quad \quad \quad \exists s. \\ \quad \quad \quad \quad (\operatorname{is\text{-}caratheodory\text{-}measurable}(\operatorname{pts}(m), u, s) \wedge ((a \subseteq s) \wedge \forall y \in b.\; \neg (y \in s))) \Rightarrow \\ \quad \quad \operatorname{is\text{-}metric\text{-}outer\text{-}measure}(m, u) \end{array}}}
\vnbprop{grid-cell-cover-nonneg}{\vnbstmt{\begin{array}{@{}l@{}} \forall r.; \forall x \in \mathbb{R}. \\ \quad (\operatorname{pos\text{-}rr}(r) \wedge (0 \leq x)) \Rightarrow  \\ \quad \quad \exists n \in \mathbb{N}. \\ \quad \quad \quad x \in \operatorname{complement\text{-}in}(\{\mathit{mhx} \in \mathbb{R} : ((n \cdot r) \leq \mathit{mhx})\}, \{\mathit{mhx} \in \mathbb{R} : (((n \cdot r) + r) \leq \mathit{mhx})\}) \end{array}}}
\vnbprop{grid-cell-cover-nonpos}{\vnbstmt{\begin{array}{@{}l@{}} \forall r.; \forall x \in \mathbb{R}. \\ \quad (\operatorname{pos\text{-}rr}(r) \wedge (x \leq 0)) \Rightarrow  \\ \quad \quad \exists n \in \mathbb{N}. \\ \quad \quad \quad x \in \operatorname{complement\text{-}in}(\{\mathit{mhx} \in \mathbb{R} : (\mathit{mhx} \leq -(n \cdot r))\}, \{\mathit{mhx} \in \mathbb{R} : (\mathit{mhx} \leq -((n \cdot r) + r))\}) \end{array}}}
\vnbprop{grid-cell-close-nonneg}{\vnbstmt{\begin{array}{@{}l@{}} \forall r, n, x, y. \\ \quad \operatorname{pos\text{-}rr}(r) \wedge \\ \quad n \in \mathbb{N} \wedge \\ \quad x \in \operatorname{complement\text{-}in}(\{\mathit{mhx} \in \mathbb{R} : ((n \cdot r) \leq \mathit{mhx})\}, \{\mathit{mhx} \in \mathbb{R} : (((n \cdot r) + r) \leq \mathit{mhx})\}) \wedge \\ \quad y \in \operatorname{complement\text{-}in}(\{\mathit{mhx} \in \mathbb{R} : ((n \cdot r) \leq \mathit{mhx})\}, \{\mathit{mhx} \in \mathbb{R} : (((n \cdot r) + r) \leq \mathit{mhx})\}) \Rightarrow \\ \quad \quad \lvert (x - y) \rvert < r \end{array}}}
\vnbprop{grid-cell-close-nonpos}{\vnbstmt{\begin{array}{@{}l@{}} \forall r, n, x, y. \\ \quad \operatorname{pos\text{-}rr}(r) \wedge \\ \quad n \in \mathbb{N} \wedge \\ \quad x \in \operatorname{complement\text{-}in}(\{\mathit{mhx} \in \mathbb{R} : (\mathit{mhx} \leq -(n \cdot r))\}, \{\mathit{mhx} \in \mathbb{R} : (\mathit{mhx} \leq -((n \cdot r) + r))\}) \wedge \\ \quad y \in \operatorname{complement\text{-}in}(\{\mathit{mhx} \in \mathbb{R} : (\mathit{mhx} \leq -(n \cdot r))\}, \{\mathit{mhx} \in \mathbb{R} : (\mathit{mhx} \leq -((n \cdot r) + r))\}) \Rightarrow \\ \quad \quad \lvert (x - y) \rvert < r \end{array}}}
\vnbprop{lebesgue-metric-outer-measure}{\vnbstmt{\begin{array}{@{}l@{}} \operatorname{is\text{-}metric\text{-}outer\text{-}measure}(\operatorname{rr\text{-}ms}, \\ \quad \operatorname{method\text{-}i\text{-}outer}(\mathbb{R}, \operatorname{closed\text{-}intervals}, (\lambda s \in \operatorname{closed\text{-}intervals}.\; \operatorname{interval\text{-}length}(s)))) \end{array}}}
\vnbprop{lebesgue-borel-measurable}{\vnbstmt{\begin{array}{@{}l@{}} \forall s \in \operatorname{borel\text{-}sets}(\operatorname{rr\text{-}ms}). \\ \quad \operatorname{is\text{-}caratheodory\text{-}measurable}(\mathbb{R}, \\ \quad \quad \operatorname{method\text{-}i\text{-}outer}(\mathbb{R}, \operatorname{closed\text{-}intervals}, (\lambda \mathit{ols} \in \operatorname{closed\text{-}intervals}.\; \operatorname{interval\text{-}length}(\mathit{ols}))), \\ \quad \quad s) \end{array}}}
\vnbprop{lebesgue-open-left-ray-measurable}{\vnbstmt{\begin{array}{@{}l@{}} \forall c \in \mathbb{R}. \\ \quad \{x \in \mathbb{R} : (x < c)\} \\ \quad \in\; \operatorname{caratheodory\text{-}sets}(\mathbb{R}, \\ \quad \quad \operatorname{method\text{-}i\text{-}outer}(\mathbb{R}, \operatorname{closed\text{-}intervals}, (\lambda s \in \operatorname{closed\text{-}intervals}.\; \operatorname{interval\text{-}length}(s)))) \end{array}}}
\vnbprop{lebesgue-open-right-ray-measurable}{\vnbstmt{\begin{array}{@{}l@{}} \forall c \in \mathbb{R}. \\ \quad \{x \in \mathbb{R} : (c < x)\} \\ \quad \in\; \operatorname{caratheodory\text{-}sets}(\mathbb{R}, \\ \quad \quad \operatorname{method\text{-}i\text{-}outer}(\mathbb{R}, \operatorname{closed\text{-}intervals}, (\lambda s \in \operatorname{closed\text{-}intervals}.\; \operatorname{interval\text{-}length}(s)))) \end{array}}}
\vnbprop{lebesgue-intervals-measurable}{\vnbstmt{\begin{array}{@{}l@{}} \forall a, b. \\ \quad ((a \in \mathbb{R}) \wedge (b \in \mathbb{R})) \Rightarrow  \\ \quad \quad \operatorname{is\text{-}caratheodory\text{-}measurable}(\mathbb{R}, \\ \quad \quad \quad \operatorname{method\text{-}i\text{-}outer}(\mathbb{R}, \operatorname{closed\text{-}intervals}, (\lambda s \in \operatorname{closed\text{-}intervals}.\; \operatorname{interval\text{-}length}(s))), \\ \quad \quad \quad \operatorname{ccint}(a, b)) \wedge \\ \quad \quad \operatorname{is\text{-}caratheodory\text{-}measurable}(\mathbb{R}, \\ \quad \quad \quad \operatorname{method\text{-}i\text{-}outer}(\mathbb{R}, \operatorname{closed\text{-}intervals}, (\lambda s \in \operatorname{closed\text{-}intervals}.\; \operatorname{interval\text{-}length}(s))), \\ \quad \quad \quad \operatorname{ooint}(a, b)) \wedge \\ \quad \quad \operatorname{is\text{-}caratheodory\text{-}measurable}(\mathbb{R}, \\ \quad \quad \quad \operatorname{method\text{-}i\text{-}outer}(\mathbb{R}, \operatorname{closed\text{-}intervals}, (\lambda s \in \operatorname{closed\text{-}intervals}.\; \operatorname{interval\text{-}length}(s))), \\ \quad \quad \quad \operatorname{coint}(a, b)) \wedge \\ \quad \quad \operatorname{is\text{-}caratheodory\text{-}measurable}(\mathbb{R}, \\ \quad \quad \quad \operatorname{method\text{-}i\text{-}outer}(\mathbb{R}, \operatorname{closed\text{-}intervals}, (\lambda s \in \operatorname{closed\text{-}intervals}.\; \operatorname{interval\text{-}length}(s))), \\ \quad \quad \quad \operatorname{ocint}(a, b)) \end{array}}}

\vnbbookfile{theorem-library/riesz-markov-cantor}{theorem-library/riesz-markov-cantor}
\noindent{\small theorem-library/riesz-markov-cantor.scm -- RIESZ-MARKOV, STAGE 2 (RM-2, 2026-10-10): a positive linear functional on the continuous functions of the Cantor model K is integration against a measure.  Vocabulary: structure-library/positive-functional.scm (CONT-FNS, IS-POSITIVE-FUNCTIONAL,}\par
\vnbprop{rmc-indicator-fun}{\vnbstmt{\forall b \in \mathbf{Set}, a.\; (\operatorname{indicator}(b, a) \in (b \to \mathbb{R}))}}
\vnbprop{rmc-indicator-values}{\vnbstmt{\begin{array}{@{}l@{}} \forall b, a.; \forall x \in b. \\ \quad ((\operatorname{indicator}(b, a)(x) \in \mathbb{R}) \wedge ((0 \leq \operatorname{indicator}(b, a)(x)) \wedge (\operatorname{indicator}(b, a)(x) \leq 1))) \end{array}}}
\vnbprop{rmc-indicator-cont}{\vnbstmt{\forall s.; \forall u \in \operatorname{clopen\text{-}sets}(s).\; (\operatorname{indicator}(\operatorname{pts}(s), u) \in \operatorname{cont\text{-}fns}(s))}}
\vnbprop{rmc-cyl-set}{\vnbstmt{\forall s, q, a, b.\; (\operatorname{rm\text{-}cyl}(s, q, a, b) \in \mathbf{Set})}}
\vnbprop{rmc-cyl-clopen}{\vnbstmt{\forall s, q, a.; \forall b \in \mathbb{N}.\; (\operatorname{rm\text{-}cyl}(s, q, a, b) \in \operatorname{clopen\text{-}sets}(\operatorname{rm\text{-}kms}(s, q)))}}
\vnbprop{rmc-cyl-indicator-cont}{\vnbstmt{\begin{array}{@{}l@{}} \forall s, q, a.; \forall b \in \mathbb{N}. \\ \quad \operatorname{indicator}(\operatorname{pts}(\operatorname{rm\text{-}kms}(s, q)), \operatorname{rm\text{-}cyl}(s, q, a, b)) \in \operatorname{cont\text{-}fns}(\operatorname{rm\text{-}kms}(s, q)) \end{array}}}
\vnbprop{rmc-compl-union}{\vnbstmt{\begin{array}{@{}l@{}} \forall a, u, v. \\ \quad \operatorname{complement\text{-}in}(a, (u \cup v)) = (\operatorname{complement\text{-}in}(a, u) \cap \operatorname{complement\text{-}in}(a, v)) \end{array}}}
\vnbprop{rmc-compl-inter}{\vnbstmt{\begin{array}{@{}l@{}} \forall a, u, v. \\ \quad \operatorname{complement\text{-}in}(a, (u \cap v)) = (\operatorname{complement\text{-}in}(a, u) \cup \operatorname{complement\text{-}in}(a, v)) \end{array}}}
\vnbprop{rmc-diff-inter}{\vnbstmt{\begin{array}{@{}l@{}} \forall a, u, v. \\ \quad (u \subseteq a) \Rightarrow (\operatorname{difference}(u, v) = (u \cap \operatorname{complement\text{-}in}(a, v))) \end{array}}}
\vnbprop{rmc-open-union}{\vnbstmt{\begin{array}{@{}l@{}} \forall s, u, v. \\ \quad (\operatorname{is\text{-}open}(s, u) \wedge \operatorname{is\text{-}open}(s, v)) \Rightarrow \operatorname{is\text{-}open}(s, (u \cup v)) \end{array}}}
\vnbprop{rmc-clopen-compl}{\vnbstmt{\forall s.; \forall u \in \operatorname{clopen\text{-}sets}(s).\; (\operatorname{complement\text{-}in}(\operatorname{pts}(s), u) \in \operatorname{clopen\text{-}sets}(s))}}
\vnbprop{rmc-clopen-inter}{\vnbstmt{\begin{array}{@{}l@{}} \forall s, u, v. \\ \quad ((u \in \operatorname{clopen\text{-}sets}(s)) \wedge (v \in \operatorname{clopen\text{-}sets}(s))) \Rightarrow  \\ \quad \quad (u \cap v) \in \operatorname{clopen\text{-}sets}(s) \end{array}}}
\vnbprop{rmc-clopen-union}{\vnbstmt{\begin{array}{@{}l@{}} \forall s, u, v. \\ \quad ((u \in \operatorname{clopen\text{-}sets}(s)) \wedge (v \in \operatorname{clopen\text{-}sets}(s))) \Rightarrow  \\ \quad \quad (u \cup v) \in \operatorname{clopen\text{-}sets}(s) \end{array}}}
\vnbprop{rmc-clopen-diff}{\vnbstmt{\begin{array}{@{}l@{}} \forall s, u, v. \\ \quad ((u \in \operatorname{clopen\text{-}sets}(s)) \wedge (v \in \operatorname{clopen\text{-}sets}(s))) \Rightarrow  \\ \quad \quad \operatorname{difference}(u, v) \in \operatorname{clopen\text{-}sets}(s) \end{array}}}
\vnbprop{rmc-clopen-empty-whole}{\vnbstmt{\begin{array}{@{}l@{}} \forall s. \\ \quad \operatorname{is\text{-}metric\text{-}space}(s) \Rightarrow  \\ \quad \quad (\emptyset \in \operatorname{clopen\text{-}sets}(s)) \wedge (\operatorname{pts}(s) \in \operatorname{clopen\text{-}sets}(s)) \end{array}}}
\vnbprop{rmc-clopen-sub-power}{\vnbstmt{\forall s.\; (\operatorname{is\text{-}metric\text{-}space}(s) \Rightarrow (\operatorname{clopen\text{-}sets}(s) \subseteq \operatorname{power}(\operatorname{pts}(s))))}}
\vnbprop{rmc-clopen-algebra}{\vnbstmt{\begin{array}{@{}l@{}} \forall s. \\ \quad \operatorname{is\text{-}metric\text{-}space}(s) \Rightarrow \operatorname{is\text{-}algebra\text{-}of\text{-}sets}(\operatorname{pts}(s), \operatorname{clopen\text{-}sets}(s)) \end{array}}}
\vnbprop{rmc-pf-add}{\vnbstmt{\begin{array}{@{}l@{}} \forall s, m, f, g. \\ \quad (\operatorname{is\text{-}positive\text{-}functional}(s, m) \wedge ((f \in \operatorname{cont\text{-}fns}(s)) \wedge (g \in \operatorname{cont\text{-}fns}(s)))) \Rightarrow  \\ \quad \quad m((\lambda x \in \operatorname{pts}(s).\; (f(x) + g(x)))) = (m(f) + m(g)) \end{array}}}
\vnbprop{rmc-pf-scale}{\vnbstmt{\begin{array}{@{}l@{}} \forall s, m, c, f. \\ \quad (\operatorname{is\text{-}positive\text{-}functional}(s, m) \wedge ((c \in \mathbb{R}) \wedge (f \in \operatorname{cont\text{-}fns}(s)))) \Rightarrow  \\ \quad \quad m((\lambda x \in \operatorname{pts}(s).\; (c \cdot f(x)))) = (c \cdot m(f)) \end{array}}}
\vnbprop{rmc-pf-mono}{\vnbstmt{\begin{array}{@{}l@{}} \forall s, m, f, g. \\ \quad \operatorname{is\text{-}positive\text{-}functional}(s, m) \wedge \\ \quad f \in \operatorname{cont\text{-}fns}(s) \wedge \\ \quad g \in \operatorname{cont\text{-}fns}(s) \wedge \\ \quad \forall x \in \operatorname{pts}(s).\; (f(x) \leq g(x)) \Rightarrow \\ \quad \quad m(f) \leq m(g) \end{array}}}
\vnbprop{rmc-pf-le-sum}{\vnbstmt{\begin{array}{@{}l@{}} \forall s, m, f, a, b. \\ \quad \operatorname{is\text{-}positive\text{-}functional}(s, m) \wedge \\ \quad f \in \operatorname{cont\text{-}fns}(s) \wedge \\ \quad a \in \operatorname{cont\text{-}fns}(s) \wedge \\ \quad b \in \operatorname{cont\text{-}fns}(s) \wedge \\ \quad \forall x \in \operatorname{pts}(s).\; (f(x) \leq (a(x) + b(x))) \Rightarrow \\ \quad \quad m(f) \leq (m(a) + m(b)) \end{array}}}
\vnbprop{rmc-pf-ge-sum}{\vnbstmt{\begin{array}{@{}l@{}} \forall s, m, f, a, b. \\ \quad \operatorname{is\text{-}positive\text{-}functional}(s, m) \wedge \\ \quad f \in \operatorname{cont\text{-}fns}(s) \wedge \\ \quad a \in \operatorname{cont\text{-}fns}(s) \wedge \\ \quad b \in \operatorname{cont\text{-}fns}(s) \wedge \\ \quad \forall x \in \operatorname{pts}(s).\; ((a(x) + b(x)) \leq f(x)) \Rightarrow \\ \quad \quad (m(a) + m(b)) \leq m(f) \end{array}}}
\vnbprop{rmc-ind-nonneg}{\vnbstmt{\begin{array}{@{}l@{}} \forall s, m, u. \\ \quad (\operatorname{is\text{-}positive\text{-}functional}(s, m) \wedge (u \in \operatorname{clopen\text{-}sets}(s))) \Rightarrow  \\ \quad \quad (m(\operatorname{indicator}(\operatorname{pts}(s), u)) \in \mathbb{R}) \wedge (0 \leq m(\operatorname{indicator}(\operatorname{pts}(s), u))) \end{array}}}
\vnbprop{rmc-ind-split-ge}{\vnbstmt{\begin{array}{@{}l@{}} \forall s, m, c, a. \\ \quad (\operatorname{is\text{-}positive\text{-}functional}(s, m) \wedge ((c \in \operatorname{clopen\text{-}sets}(s)) \wedge (a \in \operatorname{clopen\text{-}sets}(s)))) \Rightarrow  \\ \quad \quad m(\operatorname{indicator}(\operatorname{pts}(s), (c \cap a))) + m(\operatorname{indicator}(\operatorname{pts}(s), \operatorname{difference}(c, a))) \\ \quad \quad \leq\; m(\operatorname{indicator}(\operatorname{pts}(s), c)) \end{array}}}
\vnbprop{rmc-ind-cover-le}{\vnbstmt{\begin{array}{@{}l@{}} \forall s, m, a, c. \\ \quad (\operatorname{is\text{-}positive\text{-}functional}(s, m) \wedge ((a \in \operatorname{clopen\text{-}sets}(s)) \wedge (c \in \operatorname{clopen\text{-}sets}(s)))) \Rightarrow  \\ \quad \quad m(\operatorname{indicator}(\operatorname{pts}(s), a)) \leq (m(\operatorname{indicator}(\operatorname{pts}(s), \operatorname{difference}(a, c))) + m(\operatorname{indicator}(\operatorname{pts}(s), c))) \end{array}}}
\vnbprop{rmc-ind-empty}{\vnbstmt{\begin{array}{@{}l@{}} \forall s, m. \\ \quad (\operatorname{is\text{-}positive\text{-}functional}(s, m) \wedge \operatorname{is\text{-}metric\text{-}space}(s)) \Rightarrow  \\ \quad \quad m(\operatorname{indicator}(\operatorname{pts}(s), \emptyset)) = 0 \end{array}}}
\vnbprop{rmc-ind-additive}{\vnbstmt{\begin{array}{@{}l@{}} \forall s, m, u, v. \\ \quad \operatorname{is\text{-}positive\text{-}functional}(s, m) \wedge \\ \quad u \in \operatorname{clopen\text{-}sets}(s) \wedge \\ \quad v \in \operatorname{clopen\text{-}sets}(s) \wedge \\ \quad (u \cap v) = \emptyset \Rightarrow \\ \quad \quad m(\operatorname{indicator}(\operatorname{pts}(s), (u \cup v))) = (m(\operatorname{indicator}(\operatorname{pts}(s), u)) + m(\operatorname{indicator}(\operatorname{pts}(s), v))) \end{array}}}
\vnbprop{rmc-rps-of-nonneg}{\vnbstmt{\forall x \in \mathbb{R}.\; ((0 \leq x) \Rightarrow (x \in \operatorname{rr\text{-}pos\text{-}star}))}}
\vnbprop{rmc-clopen-set}{\vnbstmt{\forall s.\; (\operatorname{is\text{-}metric\text{-}space}(s) \Rightarrow (\operatorname{clopen\text{-}sets}(s) \in \mathbf{Set}))}}
\vnbprop{rmc-mu0-value}{\vnbstmt{\begin{array}{@{}l@{}} \forall s, m, u. \\ \quad (\operatorname{is\text{-}positive\text{-}functional}(s, m) \wedge (u \in \operatorname{clopen\text{-}sets}(s))) \Rightarrow  \\ \quad \quad (\lambda \mathit{rcu} \in \operatorname{clopen\text{-}sets}(s).\; m(\operatorname{indicator}(\operatorname{pts}(s), \mathit{rcu})))(u) = m(\operatorname{indicator}(\operatorname{pts}(s), u)) \end{array}}}
\vnbprop{rmc-fseq-laws}{\vnbstmt{\begin{array}{@{}l@{}} \forall s, m, q. \\ \quad (\operatorname{is\text{-}positive\text{-}functional}(s, m) \wedge (q \in (\mathbb{N} \to \operatorname{clopen\text{-}sets}(s)))) \Rightarrow  \\ \quad \quad (\lambda \mathit{cvk} \in \mathbb{N}.\; (\lambda u \in \operatorname{clopen\text{-}sets}(s).\; m(\operatorname{indicator}(\operatorname{pts}(s), u)))(q(\mathit{cvk}))) \\ \quad \quad \in\; (\mathbb{N} \to \operatorname{rr\text{-}pos\text{-}star}) \wedge \\ \quad \quad \forall k \in \mathbb{N}. \\ \quad \quad \quad (\lambda \mathit{cvk} \in \mathbb{N}.\; (\lambda u \in \operatorname{clopen\text{-}sets}(s).\; m(\operatorname{indicator}(\operatorname{pts}(s), u)))(q(\mathit{cvk})))(k) \\ \quad \quad \quad =\; m(\operatorname{indicator}(\operatorname{pts}(s), q(k))) \end{array}}}
\vnbprop{rmc-prefix-bound}{\vnbstmt{\begin{array}{@{}l@{}} \forall q \in \mathbb{N}, s, m, b.; \forall a \in \operatorname{clopen\text{-}sets}(s). \\ \quad \operatorname{is\text{-}positive\text{-}functional}(s, m) \wedge \\ \quad \operatorname{is\text{-}metric\text{-}space}(s) \wedge \\ \quad b \in (\mathbb{N} \to \operatorname{clopen\text{-}sets}(s)) \wedge \\ \quad \forall x \in a.\; \exists n \in \mathbb{N}.\; ((n < q) \wedge (x \in b(n))) \Rightarrow \\ \quad \quad m(\operatorname{indicator}(\operatorname{pts}(s), a)) \\ \quad \quad \leq\; \operatorname{finsum}(\operatorname{rr\text{-}pos\text{-}star\text{-}add\text{-}monoid}, \\ \quad \quad \quad (\lambda k \in \mathbb{N}.\; (\lambda u \in \operatorname{clopen\text{-}sets}(s).\; m(\operatorname{indicator}(\operatorname{pts}(s), u)))(b(k))), \\ \quad \quad \quad \mathrm{OS}(q)) \end{array}}}
\vnbprop{rmc-countable-subadditive}{\vnbstmt{\begin{array}{@{}l@{}} \forall s, m, q, a. \\ \quad \operatorname{is\text{-}positive\text{-}functional}(s, m) \wedge \\ \quad \operatorname{is\text{-}compact}(s) \wedge \\ \quad q \in (\mathbb{N} \to \operatorname{clopen\text{-}sets}(s)) \wedge \\ \quad a \in \operatorname{clopen\text{-}sets}(s) \wedge \\ \quad a \subseteq \operatorname{big\text{-}union}(n, \mathbb{N}, q(n)) \Rightarrow \\ \quad \quad m(\operatorname{indicator}(\operatorname{pts}(s), a)) \\ \quad \quad \leq\; \operatorname{esum}((\lambda k \in \mathbb{N}.\; (\lambda u \in \operatorname{clopen\text{-}sets}(s).\; m(\operatorname{indicator}(\operatorname{pts}(s), u)))(q(k)))) \end{array}}}
\vnbprop{rmc-premeasure}{\vnbstmt{\begin{array}{@{}l@{}} \forall s, m. \\ \quad (\operatorname{is\text{-}positive\text{-}functional}(s, m) \wedge \operatorname{is\text{-}metric\text{-}space}(s)) \Rightarrow  \\ \quad \quad \operatorname{is\text{-}premeasure}(\operatorname{pts}(s), \operatorname{clopen\text{-}sets}(s), (\lambda u \in \operatorname{clopen\text{-}sets}(s).\; m(\operatorname{indicator}(\operatorname{pts}(s), u)))) \end{array}}}
\vnbprop{rmc-power-subset}{\vnbstmt{\forall a.; \forall x \in \operatorname{power}(a).\; (x \subseteq a)}}
\vnbprop{rmc-outer-on-clopen}{\vnbstmt{\begin{array}{@{}l@{}} \forall s, m, a. \\ \quad (\operatorname{is\text{-}positive\text{-}functional}(s, m) \wedge (\operatorname{is\text{-}compact}(s) \wedge (a \in \operatorname{clopen\text{-}sets}(s)))) \Rightarrow  \\ \quad \quad \operatorname{method\text{-}i\text{-}outer}(\operatorname{pts}(s), \operatorname{clopen\text{-}sets}(s), (\lambda u \in \operatorname{clopen\text{-}sets}(s).\; m(\operatorname{indicator}(\operatorname{pts}(s), u))))(a) \\ \quad \quad =\; m(\operatorname{indicator}(\operatorname{pts}(s), a)) \end{array}}}
\vnbprop{rmc-clopen-measurable}{\vnbstmt{\begin{array}{@{}l@{}} \forall s, m, a. \\ \quad (\operatorname{is\text{-}positive\text{-}functional}(s, m) \wedge (\operatorname{is\text{-}metric\text{-}space}(s) \wedge (a \in \operatorname{clopen\text{-}sets}(s)))) \Rightarrow  \\ \quad \quad \operatorname{is\text{-}caratheodory\text{-}measurable}(\operatorname{pts}(s), \\ \quad \quad \quad \operatorname{method\text{-}i\text{-}outer}(\operatorname{pts}(s), \operatorname{clopen\text{-}sets}(s), (\lambda u \in \operatorname{clopen\text{-}sets}(s).\; m(\operatorname{indicator}(\operatorname{pts}(s), u)))), \\ \quad \quad \quad a) \end{array}}}
\vnbprop{rmc-measure}{\vnbstmt{\begin{array}{@{}l@{}} \forall s, m. \\ \quad (\operatorname{is\text{-}positive\text{-}functional}(s, m) \wedge \operatorname{is\text{-}compact}(s)) \Rightarrow  \\ \quad \quad \operatorname{is\text{-}measure}(\operatorname{pts}(s), \\ \quad \quad \quad \operatorname{caratheodory\text{-}sets}(\operatorname{pts}(s), \\ \quad \quad \quad \quad \operatorname{method\text{-}i\text{-}outer}(\operatorname{pts}(s), \operatorname{clopen\text{-}sets}(s), (\lambda \mathit{rcu} \in \operatorname{clopen\text{-}sets}(s).\; m(\operatorname{indicator}(\operatorname{pts}(s), \mathit{rcu}))))), \\ \quad \quad \quad \operatorname{res}(\operatorname{method\text{-}i\text{-}outer}(\operatorname{pts}(s), \operatorname{clopen\text{-}sets}(s), (\lambda \mathit{rcu} \in \operatorname{clopen\text{-}sets}(s).\; m(\operatorname{indicator}(\operatorname{pts}(s), \mathit{rcu})))), \\ \quad \quad \quad \quad \operatorname{caratheodory\text{-}sets}(\operatorname{pts}(s), \\ \quad \quad \quad \quad \quad \operatorname{method\text{-}i\text{-}outer}(\operatorname{pts}(s), \operatorname{clopen\text{-}sets}(s), (\lambda \mathit{rcu} \in \operatorname{clopen\text{-}sets}(s).\; m(\operatorname{indicator}(\operatorname{pts}(s), \mathit{rcu}))))))) \wedge \\ \quad \quad \operatorname{clopen\text{-}sets}(s) \\ \quad \quad \subseteq\; \operatorname{caratheodory\text{-}sets}(\operatorname{pts}(s), \\ \quad \quad \quad \operatorname{method\text{-}i\text{-}outer}(\operatorname{pts}(s), \operatorname{clopen\text{-}sets}(s), (\lambda \mathit{rcu} \in \operatorname{clopen\text{-}sets}(s).\; m(\operatorname{indicator}(\operatorname{pts}(s), \mathit{rcu}))))) \wedge \\ \quad \quad \forall u \in \operatorname{clopen\text{-}sets}(s). \\ \quad \quad \quad \operatorname{res}(\operatorname{method\text{-}i\text{-}outer}(\operatorname{pts}(s), \operatorname{clopen\text{-}sets}(s), (\lambda \mathit{rcu} \in \operatorname{clopen\text{-}sets}(s).\; m(\operatorname{indicator}(\operatorname{pts}(s), \mathit{rcu})))), \operatorname{caratheodory\text{-}sets}(\operatorname{pts}(s), \operatorname{method\text{-}i\text{-}outer}(\operatorname{pts}(s), \operatorname{clopen\text{-}sets}(s), (\lambda \mathit{rcu} \in \operatorname{clopen\text{-}sets}(s).\; m(\operatorname{indicator}(\operatorname{pts}(s), \mathit{rcu}))))))(u) \\ \quad \quad \quad =\; m(\operatorname{indicator}(\operatorname{pts}(s), u)) \end{array}}}
\vnbprop{rmc-cyl-same}{\vnbstmt{\begin{array}{@{}l@{}} \forall s, q, e, p, b. \\ \quad ((e \in \operatorname{rm\text{-}k}(s, q)) \wedge (p \in \operatorname{rm\text{-}cyl}(s, q, e, b))) \Rightarrow  \\ \quad \quad \operatorname{rm\text{-}cyl}(s, q, p, b) = \operatorname{rm\text{-}cyl}(s, q, e, b) \end{array}}}
\vnbprop{rmc-saturated-open}{\vnbstmt{\begin{array}{@{}l@{}} \forall s, q, w.; \forall b \in \mathbb{N}. \\ \quad ((w \subseteq \operatorname{rm\text{-}k}(s, q)) \wedge \forall y \in w, z \in \operatorname{rm\text{-}cyl}(s, q, y, b).\; (z \in w)) \Rightarrow  \\ \quad \quad \operatorname{is\text{-}open}(\operatorname{rm\text{-}kms}(s, q), w) \end{array}}}
\vnbprop{rmc-saturated-clopen}{\vnbstmt{\begin{array}{@{}l@{}} \forall s, q, w.; \forall b \in \mathbb{N}. \\ \quad ((w \subseteq \operatorname{rm\text{-}k}(s, q)) \wedge \forall y \in w, z \in \operatorname{rm\text{-}cyl}(s, q, y, b).\; (z \in w)) \Rightarrow  \\ \quad \quad w \in \operatorname{clopen\text{-}sets}(\operatorname{rm\text{-}kms}(s, q)) \end{array}}}
\vnbprop{rmc-w-mem}{\vnbstmt{\begin{array}{@{}l@{}} \forall s, q, u, b, x. \\ \quad x \in \{e \in \operatorname{rm\text{-}k}(s, q) : (\operatorname{rm\text{-}cyl}(s, q, e, b) \subseteq u)\} \\ \quad \Leftrightarrow\; (x \in \operatorname{rm\text{-}k}(s, q)) \wedge (\operatorname{rm\text{-}cyl}(s, q, x, b) \subseteq u) \end{array}}}
\vnbprop{rmc-w-clopen}{\vnbstmt{\begin{array}{@{}l@{}} \forall s, q, u.; \forall b \in \mathbb{N}. \\ \quad \{e \in \operatorname{rm\text{-}k}(s, q) : (\operatorname{rm\text{-}cyl}(s, q, e, b) \subseteq u)\} \in \operatorname{clopen\text{-}sets}(\operatorname{rm\text{-}kms}(s, q)) \end{array}}}
\vnbprop{rmc-open-union-w}{\vnbstmt{\begin{array}{@{}l@{}} \forall s, q, u. \\ \quad \operatorname{is\text{-}open}(\operatorname{rm\text{-}kms}(s, q), u) \Rightarrow  \\ \quad \quad u = \operatorname{big\text{-}union}(n, \mathbb{N}, (\lambda k \in \mathbb{N}.\; \{e \in \operatorname{rm\text{-}k}(s, q) : (\operatorname{rm\text{-}cyl}(s, q, e, k) \subseteq u)\})(n)) \end{array}}}
\vnbprop{rmc-k-open-measurable}{\vnbstmt{\begin{array}{@{}l@{}} \forall s, q, m, u. \\ \quad (\operatorname{is\text{-}positive\text{-}functional}(\operatorname{rm\text{-}kms}(s, q), m) \wedge \operatorname{is\text{-}open}(\operatorname{rm\text{-}kms}(s, q), u)) \Rightarrow  \\ \quad \quad u \\ \quad \quad \in\; \operatorname{caratheodory\text{-}sets}(\operatorname{pts}(\operatorname{rm\text{-}kms}(s, q)), \\ \quad \quad \quad \operatorname{method\text{-}i\text{-}outer}(\operatorname{pts}(\operatorname{rm\text{-}kms}(s, q)), \operatorname{clopen\text{-}sets}(\operatorname{rm\text{-}kms}(s, q)), \\ \quad \quad \quad \quad (\lambda \mathit{rcu} \in \operatorname{clopen\text{-}sets}(\operatorname{rm\text{-}kms}(s, q)).\; m(\operatorname{indicator}(\operatorname{pts}(\operatorname{rm\text{-}kms}(s, q)), \mathit{rcu}))))) \end{array}}}
\vnbprop{rmc-k-measure}{\vnbstmt{\begin{array}{@{}l@{}} \forall s, q, m. \\ \quad \operatorname{is\text{-}positive\text{-}functional}(\operatorname{rm\text{-}kms}(s, q), m) \Rightarrow  \\ \quad \quad \operatorname{is\text{-}measure}(\operatorname{pts}(\operatorname{rm\text{-}kms}(s, q)), \\ \quad \quad \quad \operatorname{caratheodory\text{-}sets}(\operatorname{pts}(\operatorname{rm\text{-}kms}(s, q)), \\ \quad \quad \quad \quad \operatorname{method\text{-}i\text{-}outer}(\operatorname{pts}(\operatorname{rm\text{-}kms}(s, q)), \operatorname{clopen\text{-}sets}(\operatorname{rm\text{-}kms}(s, q)), \\ \quad \quad \quad \quad \quad (\lambda \mathit{rcu} \in \operatorname{clopen\text{-}sets}(\operatorname{rm\text{-}kms}(s, q)).\; m(\operatorname{indicator}(\operatorname{pts}(\operatorname{rm\text{-}kms}(s, q)), \mathit{rcu}))))), \\ \quad \quad \quad \operatorname{res}(\operatorname{method\text{-}i\text{-}outer}(\operatorname{pts}(\operatorname{rm\text{-}kms}(s, q)), \operatorname{clopen\text{-}sets}(\operatorname{rm\text{-}kms}(s, q)), \\ \quad \quad \quad \quad \quad (\lambda \mathit{rcu} \in \operatorname{clopen\text{-}sets}(\operatorname{rm\text{-}kms}(s, q)).\; m(\operatorname{indicator}(\operatorname{pts}(\operatorname{rm\text{-}kms}(s, q)), \mathit{rcu})))), \\ \quad \quad \quad \quad \operatorname{caratheodory\text{-}sets}(\operatorname{pts}(\operatorname{rm\text{-}kms}(s, q)), \\ \quad \quad \quad \quad \quad \operatorname{method\text{-}i\text{-}outer}(\operatorname{pts}(\operatorname{rm\text{-}kms}(s, q)), \operatorname{clopen\text{-}sets}(\operatorname{rm\text{-}kms}(s, q)), \\ \quad \quad \quad \quad \quad \quad (\lambda \mathit{rcu} \in \operatorname{clopen\text{-}sets}(\operatorname{rm\text{-}kms}(s, q)).\; m(\operatorname{indicator}(\operatorname{pts}(\operatorname{rm\text{-}kms}(s, q)), \mathit{rcu}))))))) \wedge \\ \quad \quad \operatorname{borel\text{-}sets}(\operatorname{rm\text{-}kms}(s, q)) \\ \quad \quad \subseteq\; \operatorname{caratheodory\text{-}sets}(\operatorname{pts}(\operatorname{rm\text{-}kms}(s, q)), \\ \quad \quad \quad \operatorname{method\text{-}i\text{-}outer}(\operatorname{pts}(\operatorname{rm\text{-}kms}(s, q)), \operatorname{clopen\text{-}sets}(\operatorname{rm\text{-}kms}(s, q)), \\ \quad \quad \quad \quad (\lambda \mathit{rcu} \in \operatorname{clopen\text{-}sets}(\operatorname{rm\text{-}kms}(s, q)).\; m(\operatorname{indicator}(\operatorname{pts}(\operatorname{rm\text{-}kms}(s, q)), \mathit{rcu}))))) \wedge \\ \quad \quad \operatorname{clopen\text{-}sets}(\operatorname{rm\text{-}kms}(s, q)) \\ \quad \quad \subseteq\; \operatorname{caratheodory\text{-}sets}(\operatorname{pts}(\operatorname{rm\text{-}kms}(s, q)), \\ \quad \quad \quad \operatorname{method\text{-}i\text{-}outer}(\operatorname{pts}(\operatorname{rm\text{-}kms}(s, q)), \operatorname{clopen\text{-}sets}(\operatorname{rm\text{-}kms}(s, q)), \\ \quad \quad \quad \quad (\lambda \mathit{rcu} \in \operatorname{clopen\text{-}sets}(\operatorname{rm\text{-}kms}(s, q)).\; m(\operatorname{indicator}(\operatorname{pts}(\operatorname{rm\text{-}kms}(s, q)), \mathit{rcu}))))) \wedge \\ \quad \quad \forall u \in \operatorname{clopen\text{-}sets}(\operatorname{rm\text{-}kms}(s, q)). \\ \quad \quad \quad \operatorname{res}(\operatorname{method\text{-}i\text{-}outer}(\operatorname{pts}(\operatorname{rm\text{-}kms}(s, q)), \operatorname{clopen\text{-}sets}(\operatorname{rm\text{-}kms}(s, q)), (\lambda \mathit{rcu} \in \operatorname{clopen\text{-}sets}(\operatorname{rm\text{-}kms}(s, q)).\; m(\operatorname{indicator}(\operatorname{pts}(\operatorname{rm\text{-}kms}(s, q)), \mathit{rcu})))), \operatorname{caratheodory\text{-}sets}(\operatorname{pts}(\operatorname{rm\text{-}kms}(s, q)), \operatorname{method\text{-}i\text{-}outer}(\operatorname{pts}(\operatorname{rm\text{-}kms}(s, q)), \operatorname{clopen\text{-}sets}(\operatorname{rm\text{-}kms}(s, q)), (\lambda \mathit{rcu} \in \operatorname{clopen\text{-}sets}(\operatorname{rm\text{-}kms}(s, q)).\; m(\operatorname{indicator}(\operatorname{pts}(\operatorname{rm\text{-}kms}(s, q)), \mathit{rcu}))))))(u) \\ \quad \quad \quad =\; m(\operatorname{indicator}(\operatorname{pts}(\operatorname{rm\text{-}kms}(s, q)), u)) \end{array}}}
\vnbprop{rmc-lt-set-open}{\vnbstmt{\begin{array}{@{}l@{}} \forall m, g, c. \\ \quad ((g \in \operatorname{cont\text{-}fns}(m)) \wedge (c \in \mathbb{R})) \Rightarrow \operatorname{is\text{-}open}(m, \{x \in \operatorname{pts}(m) : (g(x) < c)\}) \end{array}}}
\vnbprop{rmc-gt-set-open}{\vnbstmt{\begin{array}{@{}l@{}} \forall m, g, c. \\ \quad ((g \in \operatorname{cont\text{-}fns}(m)) \wedge (c \in \mathbb{R})) \Rightarrow \operatorname{is\text{-}open}(m, \{x \in \operatorname{pts}(m) : (c < g(x))\}) \end{array}}}
\vnbprop{rmc-cont-bounded}{\vnbstmt{\begin{array}{@{}l@{}} \forall m, g. \\ \quad (\operatorname{is\text{-}compact}(m) \wedge (g \in \operatorname{cont\text{-}fns}(m))) \Rightarrow  \\ \quad \quad \exists d \in \mathbb{R}.\; ((0 \leq d) \wedge \forall x \in \operatorname{pts}(m).\; (\lvert g(x) \rvert \leq d)) \end{array}}}
\vnbprop{rmc-cyl-mono}{\vnbstmt{\begin{array}{@{}l@{}} \forall s, q, e, n, b. \\ \quad ((n \in \mathbb{N}) \wedge ((b \in \mathbb{N}) \wedge (n \leq b))) \Rightarrow  \\ \quad \quad \operatorname{rm\text{-}cyl}(s, q, e, b) \subseteq \operatorname{rm\text{-}cyl}(s, q, e, n) \end{array}}}
\vnbprop{rmc-uniform-level}{\vnbstmt{\begin{array}{@{}l@{}} \forall s, q, g, a. \\ \quad ((g \in \operatorname{cont\text{-}fns}(\operatorname{rm\text{-}kms}(s, q))) \wedge \operatorname{pos\text{-}rr}(a)) \Rightarrow  \\ \quad \quad \exists l \in \mathbb{N}. \forall x \in \operatorname{rm\text{-}k}(s, q), y \in \operatorname{rm\text{-}cyl}(s, q, x, l). \\ \quad \quad \quad \lvert (g(y) - g(x)) \rvert \leq a \end{array}}}
\vnbprop{rmc-bounded-integrable}{\vnbstmt{\begin{array}{@{}l@{}} \forall a, b, u, f, d. \\ \quad \operatorname{is\text{-}measure}(a, b, u) \wedge \\ \quad u(a) \in \mathbb{R} \wedge \\ \quad f \in (a \to \mathbb{R}) \wedge \\ \quad \forall l \in \mathbb{R}.\; (\{c \in a : (l < f(c))\} \in b) \wedge \\ \quad d \in \mathbb{R} \wedge \\ \quad 0 \leq d \wedge \\ \quad \forall x \in a.\; (\lvert f(x) \rvert \leq d) \Rightarrow \\ \quad \quad \operatorname{is\text{-}integrable}(a, b, u, f) \end{array}}}
\vnbprop{rmc-good-ind}{\vnbstmt{\begin{array}{@{}l@{}} \forall s, q, m, u. \\ \quad (\operatorname{is\text{-}positive\text{-}functional}(\operatorname{rm\text{-}kms}(s, q), m) \wedge (u \in \operatorname{clopen\text{-}sets}(\operatorname{rm\text{-}kms}(s, q)))) \Rightarrow  \\ \quad \quad \operatorname{indicator}(\operatorname{pts}(\operatorname{rm\text{-}kms}(s, q)), u) \in \operatorname{cont\text{-}fns}(\operatorname{rm\text{-}kms}(s, q)) \wedge \\ \quad \quad \operatorname{is\text{-}integrable}(\operatorname{pts}(\operatorname{rm\text{-}kms}(s, q)), \\ \quad \quad \quad \operatorname{caratheodory\text{-}sets}(\operatorname{pts}(\operatorname{rm\text{-}kms}(s, q)), \\ \quad \quad \quad \quad \operatorname{method\text{-}i\text{-}outer}(\operatorname{pts}(\operatorname{rm\text{-}kms}(s, q)), \operatorname{clopen\text{-}sets}(\operatorname{rm\text{-}kms}(s, q)), \\ \quad \quad \quad \quad \quad (\lambda \mathit{rcu} \in \operatorname{clopen\text{-}sets}(\operatorname{rm\text{-}kms}(s, q)).\; m(\operatorname{indicator}(\operatorname{pts}(\operatorname{rm\text{-}kms}(s, q)), \mathit{rcu}))))), \\ \quad \quad \quad \operatorname{res}(\operatorname{method\text{-}i\text{-}outer}(\operatorname{pts}(\operatorname{rm\text{-}kms}(s, q)), \operatorname{clopen\text{-}sets}(\operatorname{rm\text{-}kms}(s, q)), \\ \quad \quad \quad \quad \quad (\lambda \mathit{rcu} \in \operatorname{clopen\text{-}sets}(\operatorname{rm\text{-}kms}(s, q)).\; m(\operatorname{indicator}(\operatorname{pts}(\operatorname{rm\text{-}kms}(s, q)), \mathit{rcu})))), \\ \quad \quad \quad \quad \operatorname{caratheodory\text{-}sets}(\operatorname{pts}(\operatorname{rm\text{-}kms}(s, q)), \\ \quad \quad \quad \quad \quad \operatorname{method\text{-}i\text{-}outer}(\operatorname{pts}(\operatorname{rm\text{-}kms}(s, q)), \operatorname{clopen\text{-}sets}(\operatorname{rm\text{-}kms}(s, q)), \\ \quad \quad \quad \quad \quad \quad (\lambda \mathit{rcu} \in \operatorname{clopen\text{-}sets}(\operatorname{rm\text{-}kms}(s, q)).\; m(\operatorname{indicator}(\operatorname{pts}(\operatorname{rm\text{-}kms}(s, q)), \mathit{rcu})))))), \\ \quad \quad \quad \operatorname{indicator}(\operatorname{pts}(\operatorname{rm\text{-}kms}(s, q)), u)) \wedge \\ \quad \quad \operatorname{rintegral}(\operatorname{pts}(\operatorname{rm\text{-}kms}(s, q)), \\ \quad \quad \quad \operatorname{caratheodory\text{-}sets}(\operatorname{pts}(\operatorname{rm\text{-}kms}(s, q)), \\ \quad \quad \quad \quad \operatorname{method\text{-}i\text{-}outer}(\operatorname{pts}(\operatorname{rm\text{-}kms}(s, q)), \operatorname{clopen\text{-}sets}(\operatorname{rm\text{-}kms}(s, q)), \\ \quad \quad \quad \quad \quad (\lambda \mathit{rcu} \in \operatorname{clopen\text{-}sets}(\operatorname{rm\text{-}kms}(s, q)).\; m(\operatorname{indicator}(\operatorname{pts}(\operatorname{rm\text{-}kms}(s, q)), \mathit{rcu}))))), \\ \quad \quad \quad \operatorname{res}(\operatorname{method\text{-}i\text{-}outer}(\operatorname{pts}(\operatorname{rm\text{-}kms}(s, q)), \operatorname{clopen\text{-}sets}(\operatorname{rm\text{-}kms}(s, q)), \\ \quad \quad \quad \quad \quad (\lambda \mathit{rcu} \in \operatorname{clopen\text{-}sets}(\operatorname{rm\text{-}kms}(s, q)).\; m(\operatorname{indicator}(\operatorname{pts}(\operatorname{rm\text{-}kms}(s, q)), \mathit{rcu})))), \\ \quad \quad \quad \quad \operatorname{caratheodory\text{-}sets}(\operatorname{pts}(\operatorname{rm\text{-}kms}(s, q)), \\ \quad \quad \quad \quad \quad \operatorname{method\text{-}i\text{-}outer}(\operatorname{pts}(\operatorname{rm\text{-}kms}(s, q)), \operatorname{clopen\text{-}sets}(\operatorname{rm\text{-}kms}(s, q)), \\ \quad \quad \quad \quad \quad \quad (\lambda \mathit{rcu} \in \operatorname{clopen\text{-}sets}(\operatorname{rm\text{-}kms}(s, q)).\; m(\operatorname{indicator}(\operatorname{pts}(\operatorname{rm\text{-}kms}(s, q)), \mathit{rcu})))))), \\ \quad \quad \quad \operatorname{indicator}(\operatorname{pts}(\operatorname{rm\text{-}kms}(s, q)), u)) \\ \quad \quad =\; m(\operatorname{indicator}(\operatorname{pts}(\operatorname{rm\text{-}kms}(s, q)), u)) \end{array}}}
\vnbprop{rmc-good-scale}{\vnbstmt{\begin{array}{@{}l@{}} \forall s, q, m, c, f. \\ \quad \operatorname{is\text{-}positive\text{-}functional}(\operatorname{rm\text{-}kms}(s, q), m) \wedge \\ \quad c \in \mathbb{R} \wedge \\ \quad f \in \operatorname{cont\text{-}fns}(\operatorname{rm\text{-}kms}(s, q)) \wedge \\ \quad \operatorname{is\text{-}integrable}(\operatorname{pts}(\operatorname{rm\text{-}kms}(s, q)), \\ \quad \quad \operatorname{caratheodory\text{-}sets}(\operatorname{pts}(\operatorname{rm\text{-}kms}(s, q)), \\ \quad \quad \quad \operatorname{method\text{-}i\text{-}outer}(\operatorname{pts}(\operatorname{rm\text{-}kms}(s, q)), \operatorname{clopen\text{-}sets}(\operatorname{rm\text{-}kms}(s, q)), \\ \quad \quad \quad \quad (\lambda u \in \operatorname{clopen\text{-}sets}(\operatorname{rm\text{-}kms}(s, q)).\; m(\operatorname{indicator}(\operatorname{pts}(\operatorname{rm\text{-}kms}(s, q)), u))))), \\ \quad \quad \operatorname{res}(\operatorname{method\text{-}i\text{-}outer}(\operatorname{pts}(\operatorname{rm\text{-}kms}(s, q)), \operatorname{clopen\text{-}sets}(\operatorname{rm\text{-}kms}(s, q)), \\ \quad \quad \quad \quad (\lambda u \in \operatorname{clopen\text{-}sets}(\operatorname{rm\text{-}kms}(s, q)).\; m(\operatorname{indicator}(\operatorname{pts}(\operatorname{rm\text{-}kms}(s, q)), u)))), \\ \quad \quad \quad \operatorname{caratheodory\text{-}sets}(\operatorname{pts}(\operatorname{rm\text{-}kms}(s, q)), \\ \quad \quad \quad \quad \operatorname{method\text{-}i\text{-}outer}(\operatorname{pts}(\operatorname{rm\text{-}kms}(s, q)), \operatorname{clopen\text{-}sets}(\operatorname{rm\text{-}kms}(s, q)), \\ \quad \quad \quad \quad \quad (\lambda u \in \operatorname{clopen\text{-}sets}(\operatorname{rm\text{-}kms}(s, q)).\; m(\operatorname{indicator}(\operatorname{pts}(\operatorname{rm\text{-}kms}(s, q)), u)))))), \\ \quad \quad f) \wedge \\ \quad \operatorname{rintegral}(\operatorname{pts}(\operatorname{rm\text{-}kms}(s, q)), \\ \quad \quad \operatorname{caratheodory\text{-}sets}(\operatorname{pts}(\operatorname{rm\text{-}kms}(s, q)), \\ \quad \quad \quad \operatorname{method\text{-}i\text{-}outer}(\operatorname{pts}(\operatorname{rm\text{-}kms}(s, q)), \operatorname{clopen\text{-}sets}(\operatorname{rm\text{-}kms}(s, q)), \\ \quad \quad \quad \quad (\lambda u \in \operatorname{clopen\text{-}sets}(\operatorname{rm\text{-}kms}(s, q)).\; m(\operatorname{indicator}(\operatorname{pts}(\operatorname{rm\text{-}kms}(s, q)), u))))), \\ \quad \quad \operatorname{res}(\operatorname{method\text{-}i\text{-}outer}(\operatorname{pts}(\operatorname{rm\text{-}kms}(s, q)), \operatorname{clopen\text{-}sets}(\operatorname{rm\text{-}kms}(s, q)), \\ \quad \quad \quad \quad (\lambda u \in \operatorname{clopen\text{-}sets}(\operatorname{rm\text{-}kms}(s, q)).\; m(\operatorname{indicator}(\operatorname{pts}(\operatorname{rm\text{-}kms}(s, q)), u)))), \\ \quad \quad \quad \operatorname{caratheodory\text{-}sets}(\operatorname{pts}(\operatorname{rm\text{-}kms}(s, q)), \\ \quad \quad \quad \quad \operatorname{method\text{-}i\text{-}outer}(\operatorname{pts}(\operatorname{rm\text{-}kms}(s, q)), \operatorname{clopen\text{-}sets}(\operatorname{rm\text{-}kms}(s, q)), \\ \quad \quad \quad \quad \quad (\lambda u \in \operatorname{clopen\text{-}sets}(\operatorname{rm\text{-}kms}(s, q)).\; m(\operatorname{indicator}(\operatorname{pts}(\operatorname{rm\text{-}kms}(s, q)), u)))))), \\ \quad \quad f) \\ \quad =\; m(f) \Rightarrow \\ \quad \quad (\lambda x \in \operatorname{pts}(\operatorname{rm\text{-}kms}(s, q)).\; (c \cdot f(x))) \in \operatorname{cont\text{-}fns}(\operatorname{rm\text{-}kms}(s, q)) \wedge \\ \quad \quad \operatorname{is\text{-}integrable}(\operatorname{pts}(\operatorname{rm\text{-}kms}(s, q)), \\ \quad \quad \quad \operatorname{caratheodory\text{-}sets}(\operatorname{pts}(\operatorname{rm\text{-}kms}(s, q)), \\ \quad \quad \quad \quad \operatorname{method\text{-}i\text{-}outer}(\operatorname{pts}(\operatorname{rm\text{-}kms}(s, q)), \operatorname{clopen\text{-}sets}(\operatorname{rm\text{-}kms}(s, q)), \\ \quad \quad \quad \quad \quad (\lambda u \in \operatorname{clopen\text{-}sets}(\operatorname{rm\text{-}kms}(s, q)).\; m(\operatorname{indicator}(\operatorname{pts}(\operatorname{rm\text{-}kms}(s, q)), u))))), \\ \quad \quad \quad \operatorname{res}(\operatorname{method\text{-}i\text{-}outer}(\operatorname{pts}(\operatorname{rm\text{-}kms}(s, q)), \operatorname{clopen\text{-}sets}(\operatorname{rm\text{-}kms}(s, q)), \\ \quad \quad \quad \quad \quad (\lambda u \in \operatorname{clopen\text{-}sets}(\operatorname{rm\text{-}kms}(s, q)).\; m(\operatorname{indicator}(\operatorname{pts}(\operatorname{rm\text{-}kms}(s, q)), u)))), \\ \quad \quad \quad \quad \operatorname{caratheodory\text{-}sets}(\operatorname{pts}(\operatorname{rm\text{-}kms}(s, q)), \\ \quad \quad \quad \quad \quad \operatorname{method\text{-}i\text{-}outer}(\operatorname{pts}(\operatorname{rm\text{-}kms}(s, q)), \operatorname{clopen\text{-}sets}(\operatorname{rm\text{-}kms}(s, q)), \\ \quad \quad \quad \quad \quad \quad (\lambda u \in \operatorname{clopen\text{-}sets}(\operatorname{rm\text{-}kms}(s, q)).\; m(\operatorname{indicator}(\operatorname{pts}(\operatorname{rm\text{-}kms}(s, q)), u)))))), \\ \quad \quad \quad (\lambda x \in \operatorname{pts}(\operatorname{rm\text{-}kms}(s, q)).\; (c \cdot f(x)))) \wedge \\ \quad \quad \operatorname{rintegral}(\operatorname{pts}(\operatorname{rm\text{-}kms}(s, q)), \\ \quad \quad \quad \operatorname{caratheodory\text{-}sets}(\operatorname{pts}(\operatorname{rm\text{-}kms}(s, q)), \\ \quad \quad \quad \quad \operatorname{method\text{-}i\text{-}outer}(\operatorname{pts}(\operatorname{rm\text{-}kms}(s, q)), \operatorname{clopen\text{-}sets}(\operatorname{rm\text{-}kms}(s, q)), \\ \quad \quad \quad \quad \quad (\lambda u \in \operatorname{clopen\text{-}sets}(\operatorname{rm\text{-}kms}(s, q)).\; m(\operatorname{indicator}(\operatorname{pts}(\operatorname{rm\text{-}kms}(s, q)), u))))), \\ \quad \quad \quad \operatorname{res}(\operatorname{method\text{-}i\text{-}outer}(\operatorname{pts}(\operatorname{rm\text{-}kms}(s, q)), \operatorname{clopen\text{-}sets}(\operatorname{rm\text{-}kms}(s, q)), \\ \quad \quad \quad \quad \quad (\lambda u \in \operatorname{clopen\text{-}sets}(\operatorname{rm\text{-}kms}(s, q)).\; m(\operatorname{indicator}(\operatorname{pts}(\operatorname{rm\text{-}kms}(s, q)), u)))), \\ \quad \quad \quad \quad \operatorname{caratheodory\text{-}sets}(\operatorname{pts}(\operatorname{rm\text{-}kms}(s, q)), \\ \quad \quad \quad \quad \quad \operatorname{method\text{-}i\text{-}outer}(\operatorname{pts}(\operatorname{rm\text{-}kms}(s, q)), \operatorname{clopen\text{-}sets}(\operatorname{rm\text{-}kms}(s, q)), \\ \quad \quad \quad \quad \quad \quad (\lambda u \in \operatorname{clopen\text{-}sets}(\operatorname{rm\text{-}kms}(s, q)).\; m(\operatorname{indicator}(\operatorname{pts}(\operatorname{rm\text{-}kms}(s, q)), u)))))), \\ \quad \quad \quad (\lambda x \in \operatorname{pts}(\operatorname{rm\text{-}kms}(s, q)).\; (c \cdot f(x)))) \\ \quad \quad =\; m((\lambda x \in \operatorname{pts}(\operatorname{rm\text{-}kms}(s, q)).\; (c \cdot f(x)))) \end{array}}}
\vnbprop{rmc-good-add}{\vnbstmt{\begin{array}{@{}l@{}} \forall s, q, m, f, g. \\ \quad \operatorname{is\text{-}positive\text{-}functional}(\operatorname{rm\text{-}kms}(s, q), m) \wedge \\ \quad f \in \operatorname{cont\text{-}fns}(\operatorname{rm\text{-}kms}(s, q)) \wedge \\ \quad \operatorname{is\text{-}integrable}(\operatorname{pts}(\operatorname{rm\text{-}kms}(s, q)), \\ \quad \quad \operatorname{caratheodory\text{-}sets}(\operatorname{pts}(\operatorname{rm\text{-}kms}(s, q)), \\ \quad \quad \quad \operatorname{method\text{-}i\text{-}outer}(\operatorname{pts}(\operatorname{rm\text{-}kms}(s, q)), \operatorname{clopen\text{-}sets}(\operatorname{rm\text{-}kms}(s, q)), \\ \quad \quad \quad \quad (\lambda u \in \operatorname{clopen\text{-}sets}(\operatorname{rm\text{-}kms}(s, q)).\; m(\operatorname{indicator}(\operatorname{pts}(\operatorname{rm\text{-}kms}(s, q)), u))))), \\ \quad \quad \operatorname{res}(\operatorname{method\text{-}i\text{-}outer}(\operatorname{pts}(\operatorname{rm\text{-}kms}(s, q)), \operatorname{clopen\text{-}sets}(\operatorname{rm\text{-}kms}(s, q)), \\ \quad \quad \quad \quad (\lambda u \in \operatorname{clopen\text{-}sets}(\operatorname{rm\text{-}kms}(s, q)).\; m(\operatorname{indicator}(\operatorname{pts}(\operatorname{rm\text{-}kms}(s, q)), u)))), \\ \quad \quad \quad \operatorname{caratheodory\text{-}sets}(\operatorname{pts}(\operatorname{rm\text{-}kms}(s, q)), \\ \quad \quad \quad \quad \operatorname{method\text{-}i\text{-}outer}(\operatorname{pts}(\operatorname{rm\text{-}kms}(s, q)), \operatorname{clopen\text{-}sets}(\operatorname{rm\text{-}kms}(s, q)), \\ \quad \quad \quad \quad \quad (\lambda u \in \operatorname{clopen\text{-}sets}(\operatorname{rm\text{-}kms}(s, q)).\; m(\operatorname{indicator}(\operatorname{pts}(\operatorname{rm\text{-}kms}(s, q)), u)))))), \\ \quad \quad f) \wedge \\ \quad \operatorname{rintegral}(\operatorname{pts}(\operatorname{rm\text{-}kms}(s, q)), \\ \quad \quad \operatorname{caratheodory\text{-}sets}(\operatorname{pts}(\operatorname{rm\text{-}kms}(s, q)), \\ \quad \quad \quad \operatorname{method\text{-}i\text{-}outer}(\operatorname{pts}(\operatorname{rm\text{-}kms}(s, q)), \operatorname{clopen\text{-}sets}(\operatorname{rm\text{-}kms}(s, q)), \\ \quad \quad \quad \quad (\lambda u \in \operatorname{clopen\text{-}sets}(\operatorname{rm\text{-}kms}(s, q)).\; m(\operatorname{indicator}(\operatorname{pts}(\operatorname{rm\text{-}kms}(s, q)), u))))), \\ \quad \quad \operatorname{res}(\operatorname{method\text{-}i\text{-}outer}(\operatorname{pts}(\operatorname{rm\text{-}kms}(s, q)), \operatorname{clopen\text{-}sets}(\operatorname{rm\text{-}kms}(s, q)), \\ \quad \quad \quad \quad (\lambda u \in \operatorname{clopen\text{-}sets}(\operatorname{rm\text{-}kms}(s, q)).\; m(\operatorname{indicator}(\operatorname{pts}(\operatorname{rm\text{-}kms}(s, q)), u)))), \\ \quad \quad \quad \operatorname{caratheodory\text{-}sets}(\operatorname{pts}(\operatorname{rm\text{-}kms}(s, q)), \\ \quad \quad \quad \quad \operatorname{method\text{-}i\text{-}outer}(\operatorname{pts}(\operatorname{rm\text{-}kms}(s, q)), \operatorname{clopen\text{-}sets}(\operatorname{rm\text{-}kms}(s, q)), \\ \quad \quad \quad \quad \quad (\lambda u \in \operatorname{clopen\text{-}sets}(\operatorname{rm\text{-}kms}(s, q)).\; m(\operatorname{indicator}(\operatorname{pts}(\operatorname{rm\text{-}kms}(s, q)), u)))))), \\ \quad \quad f) \\ \quad =\; m(f) \wedge \\ \quad g \in \operatorname{cont\text{-}fns}(\operatorname{rm\text{-}kms}(s, q)) \wedge \\ \quad \operatorname{is\text{-}integrable}(\operatorname{pts}(\operatorname{rm\text{-}kms}(s, q)), \\ \quad \quad \operatorname{caratheodory\text{-}sets}(\operatorname{pts}(\operatorname{rm\text{-}kms}(s, q)), \\ \quad \quad \quad \operatorname{method\text{-}i\text{-}outer}(\operatorname{pts}(\operatorname{rm\text{-}kms}(s, q)), \operatorname{clopen\text{-}sets}(\operatorname{rm\text{-}kms}(s, q)), \\ \quad \quad \quad \quad (\lambda u \in \operatorname{clopen\text{-}sets}(\operatorname{rm\text{-}kms}(s, q)).\; m(\operatorname{indicator}(\operatorname{pts}(\operatorname{rm\text{-}kms}(s, q)), u))))), \\ \quad \quad \operatorname{res}(\operatorname{method\text{-}i\text{-}outer}(\operatorname{pts}(\operatorname{rm\text{-}kms}(s, q)), \operatorname{clopen\text{-}sets}(\operatorname{rm\text{-}kms}(s, q)), \\ \quad \quad \quad \quad (\lambda u \in \operatorname{clopen\text{-}sets}(\operatorname{rm\text{-}kms}(s, q)).\; m(\operatorname{indicator}(\operatorname{pts}(\operatorname{rm\text{-}kms}(s, q)), u)))), \\ \quad \quad \quad \operatorname{caratheodory\text{-}sets}(\operatorname{pts}(\operatorname{rm\text{-}kms}(s, q)), \\ \quad \quad \quad \quad \operatorname{method\text{-}i\text{-}outer}(\operatorname{pts}(\operatorname{rm\text{-}kms}(s, q)), \operatorname{clopen\text{-}sets}(\operatorname{rm\text{-}kms}(s, q)), \\ \quad \quad \quad \quad \quad (\lambda u \in \operatorname{clopen\text{-}sets}(\operatorname{rm\text{-}kms}(s, q)).\; m(\operatorname{indicator}(\operatorname{pts}(\operatorname{rm\text{-}kms}(s, q)), u)))))), \\ \quad \quad g) \wedge \\ \quad \operatorname{rintegral}(\operatorname{pts}(\operatorname{rm\text{-}kms}(s, q)), \\ \quad \quad \operatorname{caratheodory\text{-}sets}(\operatorname{pts}(\operatorname{rm\text{-}kms}(s, q)), \\ \quad \quad \quad \operatorname{method\text{-}i\text{-}outer}(\operatorname{pts}(\operatorname{rm\text{-}kms}(s, q)), \operatorname{clopen\text{-}sets}(\operatorname{rm\text{-}kms}(s, q)), \\ \quad \quad \quad \quad (\lambda u \in \operatorname{clopen\text{-}sets}(\operatorname{rm\text{-}kms}(s, q)).\; m(\operatorname{indicator}(\operatorname{pts}(\operatorname{rm\text{-}kms}(s, q)), u))))), \\ \quad \quad \operatorname{res}(\operatorname{method\text{-}i\text{-}outer}(\operatorname{pts}(\operatorname{rm\text{-}kms}(s, q)), \operatorname{clopen\text{-}sets}(\operatorname{rm\text{-}kms}(s, q)), \\ \quad \quad \quad \quad (\lambda u \in \operatorname{clopen\text{-}sets}(\operatorname{rm\text{-}kms}(s, q)).\; m(\operatorname{indicator}(\operatorname{pts}(\operatorname{rm\text{-}kms}(s, q)), u)))), \\ \quad \quad \quad \operatorname{caratheodory\text{-}sets}(\operatorname{pts}(\operatorname{rm\text{-}kms}(s, q)), \\ \quad \quad \quad \quad \operatorname{method\text{-}i\text{-}outer}(\operatorname{pts}(\operatorname{rm\text{-}kms}(s, q)), \operatorname{clopen\text{-}sets}(\operatorname{rm\text{-}kms}(s, q)), \\ \quad \quad \quad \quad \quad (\lambda u \in \operatorname{clopen\text{-}sets}(\operatorname{rm\text{-}kms}(s, q)).\; m(\operatorname{indicator}(\operatorname{pts}(\operatorname{rm\text{-}kms}(s, q)), u)))))), \\ \quad \quad g) \\ \quad =\; m(g) \Rightarrow \\ \quad \quad (\lambda x \in \operatorname{pts}(\operatorname{rm\text{-}kms}(s, q)).\; (f(x) + g(x))) \in \operatorname{cont\text{-}fns}(\operatorname{rm\text{-}kms}(s, q)) \wedge \\ \quad \quad \operatorname{is\text{-}integrable}(\operatorname{pts}(\operatorname{rm\text{-}kms}(s, q)), \\ \quad \quad \quad \operatorname{caratheodory\text{-}sets}(\operatorname{pts}(\operatorname{rm\text{-}kms}(s, q)), \\ \quad \quad \quad \quad \operatorname{method\text{-}i\text{-}outer}(\operatorname{pts}(\operatorname{rm\text{-}kms}(s, q)), \operatorname{clopen\text{-}sets}(\operatorname{rm\text{-}kms}(s, q)), \\ \quad \quad \quad \quad \quad (\lambda u \in \operatorname{clopen\text{-}sets}(\operatorname{rm\text{-}kms}(s, q)).\; m(\operatorname{indicator}(\operatorname{pts}(\operatorname{rm\text{-}kms}(s, q)), u))))), \\ \quad \quad \quad \operatorname{res}(\operatorname{method\text{-}i\text{-}outer}(\operatorname{pts}(\operatorname{rm\text{-}kms}(s, q)), \operatorname{clopen\text{-}sets}(\operatorname{rm\text{-}kms}(s, q)), \\ \quad \quad \quad \quad \quad (\lambda u \in \operatorname{clopen\text{-}sets}(\operatorname{rm\text{-}kms}(s, q)).\; m(\operatorname{indicator}(\operatorname{pts}(\operatorname{rm\text{-}kms}(s, q)), u)))), \\ \quad \quad \quad \quad \operatorname{caratheodory\text{-}sets}(\operatorname{pts}(\operatorname{rm\text{-}kms}(s, q)), \\ \quad \quad \quad \quad \quad \operatorname{method\text{-}i\text{-}outer}(\operatorname{pts}(\operatorname{rm\text{-}kms}(s, q)), \operatorname{clopen\text{-}sets}(\operatorname{rm\text{-}kms}(s, q)), \\ \quad \quad \quad \quad \quad \quad (\lambda u \in \operatorname{clopen\text{-}sets}(\operatorname{rm\text{-}kms}(s, q)).\; m(\operatorname{indicator}(\operatorname{pts}(\operatorname{rm\text{-}kms}(s, q)), u)))))), \\ \quad \quad \quad (\lambda x \in \operatorname{pts}(\operatorname{rm\text{-}kms}(s, q)).\; (f(x) + g(x)))) \wedge \\ \quad \quad \operatorname{rintegral}(\operatorname{pts}(\operatorname{rm\text{-}kms}(s, q)), \\ \quad \quad \quad \operatorname{caratheodory\text{-}sets}(\operatorname{pts}(\operatorname{rm\text{-}kms}(s, q)), \\ \quad \quad \quad \quad \operatorname{method\text{-}i\text{-}outer}(\operatorname{pts}(\operatorname{rm\text{-}kms}(s, q)), \operatorname{clopen\text{-}sets}(\operatorname{rm\text{-}kms}(s, q)), \\ \quad \quad \quad \quad \quad (\lambda u \in \operatorname{clopen\text{-}sets}(\operatorname{rm\text{-}kms}(s, q)).\; m(\operatorname{indicator}(\operatorname{pts}(\operatorname{rm\text{-}kms}(s, q)), u))))), \\ \quad \quad \quad \operatorname{res}(\operatorname{method\text{-}i\text{-}outer}(\operatorname{pts}(\operatorname{rm\text{-}kms}(s, q)), \operatorname{clopen\text{-}sets}(\operatorname{rm\text{-}kms}(s, q)), \\ \quad \quad \quad \quad \quad (\lambda u \in \operatorname{clopen\text{-}sets}(\operatorname{rm\text{-}kms}(s, q)).\; m(\operatorname{indicator}(\operatorname{pts}(\operatorname{rm\text{-}kms}(s, q)), u)))), \\ \quad \quad \quad \quad \operatorname{caratheodory\text{-}sets}(\operatorname{pts}(\operatorname{rm\text{-}kms}(s, q)), \\ \quad \quad \quad \quad \quad \operatorname{method\text{-}i\text{-}outer}(\operatorname{pts}(\operatorname{rm\text{-}kms}(s, q)), \operatorname{clopen\text{-}sets}(\operatorname{rm\text{-}kms}(s, q)), \\ \quad \quad \quad \quad \quad \quad (\lambda u \in \operatorname{clopen\text{-}sets}(\operatorname{rm\text{-}kms}(s, q)).\; m(\operatorname{indicator}(\operatorname{pts}(\operatorname{rm\text{-}kms}(s, q)), u)))))), \\ \quad \quad \quad (\lambda x \in \operatorname{pts}(\operatorname{rm\text{-}kms}(s, q)).\; (f(x) + g(x)))) \\ \quad \quad =\; m((\lambda x \in \operatorname{pts}(\operatorname{rm\text{-}kms}(s, q)).\; (f(x) + g(x)))) \end{array}}}
\vnbprop{rmc-layer}{\vnbstmt{\begin{array}{@{}l@{}} \forall n \in \mathbb{N}, s, q, m, l, c, a, f. \\ \quad \operatorname{is\text{-}positive\text{-}functional}(\operatorname{rm\text{-}kms}(s, q), m) \wedge \\ \quad l \in \mathbb{N} \wedge \\ \quad c \in \mathbb{R} \wedge \\ \quad \operatorname{pos\text{-}rr}(a) \wedge \\ \quad f \in (\operatorname{pts}(\operatorname{rm\text{-}kms}(s, q)) \to \mathbb{R}) \wedge \\ \quad \forall x \in \operatorname{rm\text{-}k}(s, q), y \in \operatorname{rm\text{-}cyl}(s, q, x, l).\; (f(y) = f(x)) \wedge \\ \quad \forall x \in \operatorname{pts}(\operatorname{rm\text{-}kms}(s, q)).\; ((c \leq f(x)) \wedge (f(x) \leq (c + ((n + 1) \cdot a)))) \Rightarrow \\ \quad \quad \exists h. \\ \quad \quad \quad h \in \operatorname{cont\text{-}fns}(\operatorname{rm\text{-}kms}(s, q)) \wedge \\ \quad \quad \quad \operatorname{is\text{-}integrable}(\operatorname{pts}(\operatorname{rm\text{-}kms}(s, q)), \\ \quad \quad \quad \quad \operatorname{caratheodory\text{-}sets}(\operatorname{pts}(\operatorname{rm\text{-}kms}(s, q)), \\ \quad \quad \quad \quad \quad \operatorname{method\text{-}i\text{-}outer}(\operatorname{pts}(\operatorname{rm\text{-}kms}(s, q)), \operatorname{clopen\text{-}sets}(\operatorname{rm\text{-}kms}(s, q)), \\ \quad \quad \quad \quad \quad \quad (\lambda u \in \operatorname{clopen\text{-}sets}(\operatorname{rm\text{-}kms}(s, q)).\; m(\operatorname{indicator}(\operatorname{pts}(\operatorname{rm\text{-}kms}(s, q)), u))))), \\ \quad \quad \quad \quad \operatorname{res}(\operatorname{method\text{-}i\text{-}outer}(\operatorname{pts}(\operatorname{rm\text{-}kms}(s, q)), \operatorname{clopen\text{-}sets}(\operatorname{rm\text{-}kms}(s, q)), \\ \quad \quad \quad \quad \quad \quad (\lambda u \in \operatorname{clopen\text{-}sets}(\operatorname{rm\text{-}kms}(s, q)).\; m(\operatorname{indicator}(\operatorname{pts}(\operatorname{rm\text{-}kms}(s, q)), u)))), \\ \quad \quad \quad \quad \quad \operatorname{caratheodory\text{-}sets}(\operatorname{pts}(\operatorname{rm\text{-}kms}(s, q)), \\ \quad \quad \quad \quad \quad \quad \operatorname{method\text{-}i\text{-}outer}(\operatorname{pts}(\operatorname{rm\text{-}kms}(s, q)), \operatorname{clopen\text{-}sets}(\operatorname{rm\text{-}kms}(s, q)), \\ \quad \quad \quad \quad \quad \quad \quad (\lambda u \in \operatorname{clopen\text{-}sets}(\operatorname{rm\text{-}kms}(s, q)).\; m(\operatorname{indicator}(\operatorname{pts}(\operatorname{rm\text{-}kms}(s, q)), u)))))), \\ \quad \quad \quad \quad h) \wedge \\ \quad \quad \quad \operatorname{rintegral}(\operatorname{pts}(\operatorname{rm\text{-}kms}(s, q)), \\ \quad \quad \quad \quad \operatorname{caratheodory\text{-}sets}(\operatorname{pts}(\operatorname{rm\text{-}kms}(s, q)), \\ \quad \quad \quad \quad \quad \operatorname{method\text{-}i\text{-}outer}(\operatorname{pts}(\operatorname{rm\text{-}kms}(s, q)), \operatorname{clopen\text{-}sets}(\operatorname{rm\text{-}kms}(s, q)), \\ \quad \quad \quad \quad \quad \quad (\lambda u \in \operatorname{clopen\text{-}sets}(\operatorname{rm\text{-}kms}(s, q)).\; m(\operatorname{indicator}(\operatorname{pts}(\operatorname{rm\text{-}kms}(s, q)), u))))), \\ \quad \quad \quad \quad \operatorname{res}(\operatorname{method\text{-}i\text{-}outer}(\operatorname{pts}(\operatorname{rm\text{-}kms}(s, q)), \operatorname{clopen\text{-}sets}(\operatorname{rm\text{-}kms}(s, q)), \\ \quad \quad \quad \quad \quad \quad (\lambda u \in \operatorname{clopen\text{-}sets}(\operatorname{rm\text{-}kms}(s, q)).\; m(\operatorname{indicator}(\operatorname{pts}(\operatorname{rm\text{-}kms}(s, q)), u)))), \\ \quad \quad \quad \quad \quad \operatorname{caratheodory\text{-}sets}(\operatorname{pts}(\operatorname{rm\text{-}kms}(s, q)), \\ \quad \quad \quad \quad \quad \quad \operatorname{method\text{-}i\text{-}outer}(\operatorname{pts}(\operatorname{rm\text{-}kms}(s, q)), \operatorname{clopen\text{-}sets}(\operatorname{rm\text{-}kms}(s, q)), \\ \quad \quad \quad \quad \quad \quad \quad (\lambda u \in \operatorname{clopen\text{-}sets}(\operatorname{rm\text{-}kms}(s, q)).\; m(\operatorname{indicator}(\operatorname{pts}(\operatorname{rm\text{-}kms}(s, q)), u)))))), \\ \quad \quad \quad \quad h) \\ \quad \quad \quad =\; m(h) \wedge \\ \quad \quad \quad \forall x \in \operatorname{pts}(\operatorname{rm\text{-}kms}(s, q)).\; (\lvert (f(x) - h(x)) \rvert \leq a) \end{array}}}
\vnbprop{rmc-approx}{\vnbstmt{\begin{array}{@{}l@{}} \forall s, q, m, g, a. \\ \quad (\operatorname{is\text{-}positive\text{-}functional}(\operatorname{rm\text{-}kms}(s, q), m) \wedge ((g \in \operatorname{cont\text{-}fns}(\operatorname{rm\text{-}kms}(s, q))) \wedge \operatorname{pos\text{-}rr}(a))) \Rightarrow  \\ \quad \quad \exists h. \\ \quad \quad \quad h \in \operatorname{cont\text{-}fns}(\operatorname{rm\text{-}kms}(s, q)) \wedge \\ \quad \quad \quad \operatorname{is\text{-}integrable}(\operatorname{pts}(\operatorname{rm\text{-}kms}(s, q)), \\ \quad \quad \quad \quad \operatorname{caratheodory\text{-}sets}(\operatorname{pts}(\operatorname{rm\text{-}kms}(s, q)), \\ \quad \quad \quad \quad \quad \operatorname{method\text{-}i\text{-}outer}(\operatorname{pts}(\operatorname{rm\text{-}kms}(s, q)), \operatorname{clopen\text{-}sets}(\operatorname{rm\text{-}kms}(s, q)), \\ \quad \quad \quad \quad \quad \quad (\lambda u \in \operatorname{clopen\text{-}sets}(\operatorname{rm\text{-}kms}(s, q)).\; m(\operatorname{indicator}(\operatorname{pts}(\operatorname{rm\text{-}kms}(s, q)), u))))), \\ \quad \quad \quad \quad \operatorname{res}(\operatorname{method\text{-}i\text{-}outer}(\operatorname{pts}(\operatorname{rm\text{-}kms}(s, q)), \operatorname{clopen\text{-}sets}(\operatorname{rm\text{-}kms}(s, q)), \\ \quad \quad \quad \quad \quad \quad (\lambda u \in \operatorname{clopen\text{-}sets}(\operatorname{rm\text{-}kms}(s, q)).\; m(\operatorname{indicator}(\operatorname{pts}(\operatorname{rm\text{-}kms}(s, q)), u)))), \\ \quad \quad \quad \quad \quad \operatorname{caratheodory\text{-}sets}(\operatorname{pts}(\operatorname{rm\text{-}kms}(s, q)), \\ \quad \quad \quad \quad \quad \quad \operatorname{method\text{-}i\text{-}outer}(\operatorname{pts}(\operatorname{rm\text{-}kms}(s, q)), \operatorname{clopen\text{-}sets}(\operatorname{rm\text{-}kms}(s, q)), \\ \quad \quad \quad \quad \quad \quad \quad (\lambda u \in \operatorname{clopen\text{-}sets}(\operatorname{rm\text{-}kms}(s, q)).\; m(\operatorname{indicator}(\operatorname{pts}(\operatorname{rm\text{-}kms}(s, q)), u)))))), \\ \quad \quad \quad \quad h) \wedge \\ \quad \quad \quad \operatorname{rintegral}(\operatorname{pts}(\operatorname{rm\text{-}kms}(s, q)), \\ \quad \quad \quad \quad \operatorname{caratheodory\text{-}sets}(\operatorname{pts}(\operatorname{rm\text{-}kms}(s, q)), \\ \quad \quad \quad \quad \quad \operatorname{method\text{-}i\text{-}outer}(\operatorname{pts}(\operatorname{rm\text{-}kms}(s, q)), \operatorname{clopen\text{-}sets}(\operatorname{rm\text{-}kms}(s, q)), \\ \quad \quad \quad \quad \quad \quad (\lambda u \in \operatorname{clopen\text{-}sets}(\operatorname{rm\text{-}kms}(s, q)).\; m(\operatorname{indicator}(\operatorname{pts}(\operatorname{rm\text{-}kms}(s, q)), u))))), \\ \quad \quad \quad \quad \operatorname{res}(\operatorname{method\text{-}i\text{-}outer}(\operatorname{pts}(\operatorname{rm\text{-}kms}(s, q)), \operatorname{clopen\text{-}sets}(\operatorname{rm\text{-}kms}(s, q)), \\ \quad \quad \quad \quad \quad \quad (\lambda u \in \operatorname{clopen\text{-}sets}(\operatorname{rm\text{-}kms}(s, q)).\; m(\operatorname{indicator}(\operatorname{pts}(\operatorname{rm\text{-}kms}(s, q)), u)))), \\ \quad \quad \quad \quad \quad \operatorname{caratheodory\text{-}sets}(\operatorname{pts}(\operatorname{rm\text{-}kms}(s, q)), \\ \quad \quad \quad \quad \quad \quad \operatorname{method\text{-}i\text{-}outer}(\operatorname{pts}(\operatorname{rm\text{-}kms}(s, q)), \operatorname{clopen\text{-}sets}(\operatorname{rm\text{-}kms}(s, q)), \\ \quad \quad \quad \quad \quad \quad \quad (\lambda u \in \operatorname{clopen\text{-}sets}(\operatorname{rm\text{-}kms}(s, q)).\; m(\operatorname{indicator}(\operatorname{pts}(\operatorname{rm\text{-}kms}(s, q)), u)))))), \\ \quad \quad \quad \quad h) \\ \quad \quad \quad =\; m(h) \wedge \\ \quad \quad \quad \forall x \in \operatorname{pts}(\operatorname{rm\text{-}kms}(s, q)).\; (\lvert (g(x) - h(x)) \rvert \leq (a + a)) \end{array}}}
\vnbprop{rmc-riesz-markov-k}{\vnbstmt{\begin{array}{@{}l@{}} \forall s, q, m, g. \\ \quad (\operatorname{is\text{-}positive\text{-}functional}(\operatorname{rm\text{-}kms}(s, q), m) \wedge (g \in \operatorname{cont\text{-}fns}(\operatorname{rm\text{-}kms}(s, q)))) \Rightarrow  \\ \quad \quad \operatorname{is\text{-}integrable}(\operatorname{pts}(\operatorname{rm\text{-}kms}(s, q)), \\ \quad \quad \quad \operatorname{caratheodory\text{-}sets}(\operatorname{pts}(\operatorname{rm\text{-}kms}(s, q)), \\ \quad \quad \quad \quad \operatorname{method\text{-}i\text{-}outer}(\operatorname{pts}(\operatorname{rm\text{-}kms}(s, q)), \operatorname{clopen\text{-}sets}(\operatorname{rm\text{-}kms}(s, q)), \\ \quad \quad \quad \quad \quad (\lambda u \in \operatorname{clopen\text{-}sets}(\operatorname{rm\text{-}kms}(s, q)).\; m(\operatorname{indicator}(\operatorname{pts}(\operatorname{rm\text{-}kms}(s, q)), u))))), \\ \quad \quad \quad \operatorname{res}(\operatorname{method\text{-}i\text{-}outer}(\operatorname{pts}(\operatorname{rm\text{-}kms}(s, q)), \operatorname{clopen\text{-}sets}(\operatorname{rm\text{-}kms}(s, q)), \\ \quad \quad \quad \quad \quad (\lambda u \in \operatorname{clopen\text{-}sets}(\operatorname{rm\text{-}kms}(s, q)).\; m(\operatorname{indicator}(\operatorname{pts}(\operatorname{rm\text{-}kms}(s, q)), u)))), \\ \quad \quad \quad \quad \operatorname{caratheodory\text{-}sets}(\operatorname{pts}(\operatorname{rm\text{-}kms}(s, q)), \\ \quad \quad \quad \quad \quad \operatorname{method\text{-}i\text{-}outer}(\operatorname{pts}(\operatorname{rm\text{-}kms}(s, q)), \operatorname{clopen\text{-}sets}(\operatorname{rm\text{-}kms}(s, q)), \\ \quad \quad \quad \quad \quad \quad (\lambda u \in \operatorname{clopen\text{-}sets}(\operatorname{rm\text{-}kms}(s, q)).\; m(\operatorname{indicator}(\operatorname{pts}(\operatorname{rm\text{-}kms}(s, q)), u)))))), \\ \quad \quad \quad g) \wedge \\ \quad \quad \operatorname{rintegral}(\operatorname{pts}(\operatorname{rm\text{-}kms}(s, q)), \\ \quad \quad \quad \operatorname{caratheodory\text{-}sets}(\operatorname{pts}(\operatorname{rm\text{-}kms}(s, q)), \\ \quad \quad \quad \quad \operatorname{method\text{-}i\text{-}outer}(\operatorname{pts}(\operatorname{rm\text{-}kms}(s, q)), \operatorname{clopen\text{-}sets}(\operatorname{rm\text{-}kms}(s, q)), \\ \quad \quad \quad \quad \quad (\lambda u \in \operatorname{clopen\text{-}sets}(\operatorname{rm\text{-}kms}(s, q)).\; m(\operatorname{indicator}(\operatorname{pts}(\operatorname{rm\text{-}kms}(s, q)), u))))), \\ \quad \quad \quad \operatorname{res}(\operatorname{method\text{-}i\text{-}outer}(\operatorname{pts}(\operatorname{rm\text{-}kms}(s, q)), \operatorname{clopen\text{-}sets}(\operatorname{rm\text{-}kms}(s, q)), \\ \quad \quad \quad \quad \quad (\lambda u \in \operatorname{clopen\text{-}sets}(\operatorname{rm\text{-}kms}(s, q)).\; m(\operatorname{indicator}(\operatorname{pts}(\operatorname{rm\text{-}kms}(s, q)), u)))), \\ \quad \quad \quad \quad \operatorname{caratheodory\text{-}sets}(\operatorname{pts}(\operatorname{rm\text{-}kms}(s, q)), \\ \quad \quad \quad \quad \quad \operatorname{method\text{-}i\text{-}outer}(\operatorname{pts}(\operatorname{rm\text{-}kms}(s, q)), \operatorname{clopen\text{-}sets}(\operatorname{rm\text{-}kms}(s, q)), \\ \quad \quad \quad \quad \quad \quad (\lambda u \in \operatorname{clopen\text{-}sets}(\operatorname{rm\text{-}kms}(s, q)).\; m(\operatorname{indicator}(\operatorname{pts}(\operatorname{rm\text{-}kms}(s, q)), u)))))), \\ \quad \quad \quad g) \\ \quad \quad =\; m(g) \end{array}}}

\vnbbookfile{theorem-library/integral-unique}{theorem-library/integral-unique}
\noindent{\small integral-unique.scm -- WAVE 2 (2026-10-27..30), wired 2026-10-30 from scratchpad/md/md-unique.scm. Library: measure.}\par
\vnbprop{integral-unique-step-ptwise}{\vnbstmt{\begin{array}{@{}l@{}} \forall a, s, w, y.; \forall x \in a. \\ \quad ((s(x) \in \mathbb{R}) \wedge ((0 \leq s(x)) \wedge ((y \in \mathbb{R}) \wedge ((0 \leq y) \wedge (\neg (y \in w) \wedge (w \in \mathbf{Set})))))) \Rightarrow  \\ \quad \quad \left(\text{if }(s(x) \in (w \cup \{y\}))\text{ then }s(x)\text{ else }0\right) \\ \quad \quad =\; \operatorname{eplus}(\left(\text{if }(s(x) \in w)\text{ then }s(x)\text{ else }0\right), \operatorname{etimes}(y, \left(\text{if }(x \in \{\mathit{gx} \in a : (s(\mathit{gx}) = y)\})\text{ then }1\text{ else }0\right))) \end{array}}}
\vnbprop{integral-unique-simple}{\vnbstmt{\begin{array}{@{}l@{}} \forall a, b, u, j, s. \\ \quad \operatorname{is\text{-}measure}(a, b, u) \wedge \\ \quad \forall f.\; (\operatorname{is\text{-}measurable\text{-}fn}(a, b, f) \Rightarrow (j(f) \in \operatorname{rr\text{-}pos\text{-}star})) \wedge \\ \quad \forall \mathit{mda} \in b.\; (j(\operatorname{indicator}(a, \mathit{mda})) = u(\mathit{mda})) \wedge \\ \quad \forall f, g. \\ \quad \quad (\operatorname{is\text{-}measurable\text{-}fn}(a, b, f) \wedge \operatorname{is\text{-}measurable\text{-}fn}(a, b, g)) \Rightarrow  \\ \quad \quad \quad j(\operatorname{ptwise\text{-}eplus}(a, f, g)) = \operatorname{eplus}(j(f), j(g)) \wedge \\ \quad \forall f.; \forall c \in \mathbb{R}. \\ \quad \quad (\operatorname{is\text{-}measurable\text{-}fn}(a, b, f) \wedge (0 \leq c)) \Rightarrow  \\ \quad \quad \quad j(\operatorname{ptwise\text{-}scale}(a, c, f)) = \operatorname{etimes}(c, j(f)) \wedge \\ \quad \operatorname{is\text{-}simple\text{-}fn}(a, b, s) \Rightarrow \\ \quad \quad j(s) = \operatorname{integral}(a, b, u, s) \end{array}}}
\vnbprop{integral-unique}{\vnbstmt{\begin{array}{@{}l@{}} \forall a, b, u, j, f. \\ \quad \operatorname{is\text{-}measure}(a, b, u) \wedge \\ \quad \forall \mathit{mdf}.\; (\operatorname{is\text{-}measurable\text{-}fn}(a, b, \mathit{mdf}) \Rightarrow (j(\mathit{mdf}) \in \operatorname{rr\text{-}pos\text{-}star})) \wedge \\ \quad \forall \mathit{mda} \in b.\; (j(\operatorname{indicator}(a, \mathit{mda})) = u(\mathit{mda})) \wedge \\ \quad \forall \mathit{mda} \in b. \\ \quad \quad (u(\mathit{mda}) = 0) \Rightarrow (j(\operatorname{ptwise\text{-}scale}(a, \operatorname{pos\text{-}inf}, \operatorname{indicator}(a, \mathit{mda}))) = 0) \wedge \\ \quad \forall \mathit{mdf}, g. \\ \quad \quad (\operatorname{is\text{-}measurable\text{-}fn}(a, b, \mathit{mdf}) \wedge \operatorname{is\text{-}measurable\text{-}fn}(a, b, g)) \Rightarrow  \\ \quad \quad \quad j(\operatorname{ptwise\text{-}eplus}(a, \mathit{mdf}, g)) = \operatorname{eplus}(j(\mathit{mdf}), j(g)) \wedge \\ \quad \forall \mathit{mdf}.; \forall c \in \mathbb{R}. \\ \quad \quad (\operatorname{is\text{-}measurable\text{-}fn}(a, b, \mathit{mdf}) \wedge (0 \leq c)) \Rightarrow  \\ \quad \quad \quad j(\operatorname{ptwise\text{-}scale}(a, c, \mathit{mdf})) = \operatorname{etimes}(c, j(\mathit{mdf})) \wedge \\ \quad \forall \mathit{mdf}, g. \\ \quad \quad (\operatorname{is\text{-}measurable\text{-}fn}(a, b, \mathit{mdf}) \wedge (\operatorname{is\text{-}measurable\text{-}fn}(a, b, g) \wedge \operatorname{ptwise\text{-}le}(a, \mathit{mdf}, g))) \Rightarrow  \\ \quad \quad \quad j(\mathit{mdf}) \leq j(g) \wedge \\ \quad \forall \mathit{mdcf} \in (\mathbb{N} \to (a \to \operatorname{rr\text{-}pos\text{-}star})). \\ \quad \quad \forall n \in \mathbb{N}.\; \operatorname{is\text{-}measurable\text{-}fn}(a, b, \operatorname{mdcf}(n)) \wedge \\ \quad \quad \forall n \in \mathbb{N}.\; \operatorname{ptwise\text{-}le}(a, \operatorname{mdcf}(n), \operatorname{mdcf}((n + 1))) \Rightarrow \\ \quad \quad \quad j(\operatorname{ptwise\text{-}sup}(a, \mathit{mdcf})) = \operatorname{esup}(\operatorname{image}((\lambda n \in \mathbb{N}.\; j(\operatorname{mdcf}(n))), \mathbb{N})) \wedge \\ \quad \operatorname{is\text{-}measurable\text{-}fn}(a, b, f) \Rightarrow \\ \quad \quad j(f) = \operatorname{integral}(a, b, u, f) \end{array}}}

\vnbbookfile{theorem-library/positive-extension}{theorem-library/positive-extension}
\noindent{\small theorem-library/positive-extension.scm -- THE POSITIVE EXTENSION (M. Riesz) for C(K), K a compact metric space (RM-3A, 2026-10-10).  Prefix pex-.}\par
\vnbprop{pex-c-set}{\vnbstmt{\forall s.\; (\operatorname{is\text{-}metric\text{-}space}(s) \Rightarrow (\operatorname{cont\text{-}fns}(s) \in \mathbf{Set}))}}
\vnbprop{pex-cf-scal}{\vnbstmt{\begin{array}{@{}l@{}} \forall s. \\ \quad \operatorname{is\text{-}metric\text{-}space}(s) \Rightarrow  \\ \quad \quad (\operatorname{scal}({[}\operatorname{normed\text{-}field\text{-}as\text{-}commutative\text{-}ring}(\operatorname{rr\text{-}normed\text{-}field}), \operatorname{cont\text{-}fns}(s), (\lambda {[}a, \operatorname{peb\_}] \in (\operatorname{cont\text{-}fns}(s) \times \operatorname{cont\text{-}fns}(s)).\; (\lambda x \in \operatorname{pts}(s).\; (a(x) + \operatorname{peb\_}(x)))), (\lambda x \in \operatorname{pts}(s).\; 0), (\lambda a \in \operatorname{cont\text{-}fns}(s).\; (\lambda x \in \operatorname{pts}(s).\; (-1 \cdot a(x)))), (\lambda {[}\operatorname{pec\_}, a] \in (\mathbb{R} \times \operatorname{cont\text{-}fns}(s)).\; (\lambda x \in \operatorname{pts}(s).\; (\operatorname{pec\_} \cdot a(x))))]) \simeq \operatorname{normed\text{-}field\text{-}as\text{-}commutative\text{-}ring}(\operatorname{rr\text{-}normed\text{-}field})) \end{array}}}
\vnbprop{pex-cf-vec}{\vnbstmt{\begin{array}{@{}l@{}} \forall s. \\ \quad \operatorname{is\text{-}metric\text{-}space}(s) \Rightarrow  \\ \quad \quad (\operatorname{vec}({[}\operatorname{normed\text{-}field\text{-}as\text{-}commutative\text{-}ring}(\operatorname{rr\text{-}normed\text{-}field}), \operatorname{cont\text{-}fns}(s), (\lambda {[}a, \operatorname{peb\_}] \in (\operatorname{cont\text{-}fns}(s) \times \operatorname{cont\text{-}fns}(s)).\; (\lambda x \in \operatorname{pts}(s).\; (a(x) + \operatorname{peb\_}(x)))), (\lambda x \in \operatorname{pts}(s).\; 0), (\lambda a \in \operatorname{cont\text{-}fns}(s).\; (\lambda x \in \operatorname{pts}(s).\; (-1 \cdot a(x)))), (\lambda {[}\operatorname{pec\_}, a] \in (\mathbb{R} \times \operatorname{cont\text{-}fns}(s)).\; (\lambda x \in \operatorname{pts}(s).\; (\operatorname{pec\_} \cdot a(x))))]) \simeq \operatorname{cont\text{-}fns}(s)) \end{array}}}
\vnbprop{pex-cf-vadd}{\vnbstmt{\begin{array}{@{}l@{}} \forall s. \\ \quad \operatorname{is\text{-}metric\text{-}space}(s) \Rightarrow  \\ \quad \quad (\operatorname{vadd}({[}\operatorname{normed\text{-}field\text{-}as\text{-}commutative\text{-}ring}(\operatorname{rr\text{-}normed\text{-}field}), \operatorname{cont\text{-}fns}(s), (\lambda {[}a, \operatorname{peb\_}] \in (\operatorname{cont\text{-}fns}(s) \times \operatorname{cont\text{-}fns}(s)).\; (\lambda x \in \operatorname{pts}(s).\; (a(x) + \operatorname{peb\_}(x)))), (\lambda x \in \operatorname{pts}(s).\; 0), (\lambda a \in \operatorname{cont\text{-}fns}(s).\; (\lambda x \in \operatorname{pts}(s).\; (-1 \cdot a(x)))), (\lambda {[}\operatorname{pec\_}, a] \in (\mathbb{R} \times \operatorname{cont\text{-}fns}(s)).\; (\lambda x \in \operatorname{pts}(s).\; (\operatorname{pec\_} \cdot a(x))))]) \simeq (\lambda {[}a, \operatorname{peb\_}] \in (\operatorname{cont\text{-}fns}(s) \times \operatorname{cont\text{-}fns}(s)).\; (\lambda x \in \operatorname{pts}(s).\; (a(x) + \operatorname{peb\_}(x))))) \end{array}}}
\vnbprop{pex-cf-vzero}{\vnbstmt{\begin{array}{@{}l@{}} \forall s. \\ \quad \operatorname{is\text{-}metric\text{-}space}(s) \Rightarrow  \\ \quad \quad (\operatorname{vzero}({[}\operatorname{normed\text{-}field\text{-}as\text{-}commutative\text{-}ring}(\operatorname{rr\text{-}normed\text{-}field}), \operatorname{cont\text{-}fns}(s), (\lambda {[}a, \operatorname{peb\_}] \in (\operatorname{cont\text{-}fns}(s) \times \operatorname{cont\text{-}fns}(s)).\; (\lambda x \in \operatorname{pts}(s).\; (a(x) + \operatorname{peb\_}(x)))), (\lambda x \in \operatorname{pts}(s).\; 0), (\lambda a \in \operatorname{cont\text{-}fns}(s).\; (\lambda x \in \operatorname{pts}(s).\; (-1 \cdot a(x)))), (\lambda {[}\operatorname{pec\_}, a] \in (\mathbb{R} \times \operatorname{cont\text{-}fns}(s)).\; (\lambda x \in \operatorname{pts}(s).\; (\operatorname{pec\_} \cdot a(x))))]) \simeq (\lambda x \in \operatorname{pts}(s).\; 0)) \end{array}}}
\vnbprop{pex-cf-vneg}{\vnbstmt{\begin{array}{@{}l@{}} \forall s. \\ \quad \operatorname{is\text{-}metric\text{-}space}(s) \Rightarrow  \\ \quad \quad (\operatorname{vneg}({[}\operatorname{normed\text{-}field\text{-}as\text{-}commutative\text{-}ring}(\operatorname{rr\text{-}normed\text{-}field}), \operatorname{cont\text{-}fns}(s), (\lambda {[}a, \operatorname{peb\_}] \in (\operatorname{cont\text{-}fns}(s) \times \operatorname{cont\text{-}fns}(s)).\; (\lambda x \in \operatorname{pts}(s).\; (a(x) + \operatorname{peb\_}(x)))), (\lambda x \in \operatorname{pts}(s).\; 0), (\lambda a \in \operatorname{cont\text{-}fns}(s).\; (\lambda x \in \operatorname{pts}(s).\; (-1 \cdot a(x)))), (\lambda {[}\operatorname{pec\_}, a] \in (\mathbb{R} \times \operatorname{cont\text{-}fns}(s)).\; (\lambda x \in \operatorname{pts}(s).\; (\operatorname{pec\_} \cdot a(x))))]) \simeq (\lambda a \in \operatorname{cont\text{-}fns}(s).\; (\lambda x \in \operatorname{pts}(s).\; (-1 \cdot a(x))))) \end{array}}}
\vnbprop{pex-cf-act}{\vnbstmt{\begin{array}{@{}l@{}} \forall s. \\ \quad \operatorname{is\text{-}metric\text{-}space}(s) \Rightarrow  \\ \quad \quad (\operatorname{act}({[}\operatorname{normed\text{-}field\text{-}as\text{-}commutative\text{-}ring}(\operatorname{rr\text{-}normed\text{-}field}), \operatorname{cont\text{-}fns}(s), (\lambda {[}a, \operatorname{peb\_}] \in (\operatorname{cont\text{-}fns}(s) \times \operatorname{cont\text{-}fns}(s)).\; (\lambda x \in \operatorname{pts}(s).\; (a(x) + \operatorname{peb\_}(x)))), (\lambda x \in \operatorname{pts}(s).\; 0), (\lambda a \in \operatorname{cont\text{-}fns}(s).\; (\lambda x \in \operatorname{pts}(s).\; (-1 \cdot a(x)))), (\lambda {[}\operatorname{pec\_}, a] \in (\mathbb{R} \times \operatorname{cont\text{-}fns}(s)).\; (\lambda x \in \operatorname{pts}(s).\; (\operatorname{pec\_} \cdot a(x))))]) \simeq (\lambda {[}\operatorname{pec\_}, a] \in (\mathbb{R} \times \operatorname{cont\text{-}fns}(s)).\; (\lambda x \in \operatorname{pts}(s).\; (\operatorname{pec\_} \cdot a(x))))) \end{array}}}
\vnbprop{pex-zero-in}{\vnbstmt{\forall s.\; (\operatorname{is\text{-}metric\text{-}space}(s) \Rightarrow ((\lambda x \in \operatorname{pts}(s).\; 0) \in \operatorname{cont\text{-}fns}(s)))}}
\vnbprop{pex-cf-module}{\vnbstmt{\begin{array}{@{}l@{}} \forall s. \\ \quad \operatorname{is\text{-}metric\text{-}space}(s) \Rightarrow  \\ \quad \quad \operatorname{is\text{-}module}({[}\operatorname{normed\text{-}field\text{-}as\text{-}commutative\text{-}ring}(\operatorname{rr\text{-}normed\text{-}field}), \operatorname{cont\text{-}fns}(s), (\lambda {[}a, \operatorname{peb\_}] \in (\operatorname{cont\text{-}fns}(s) \times \operatorname{cont\text{-}fns}(s)).\; (\lambda x \in \operatorname{pts}(s).\; (a(x) + \operatorname{peb\_}(x)))), (\lambda x \in \operatorname{pts}(s).\; 0), (\lambda a \in \operatorname{cont\text{-}fns}(s).\; (\lambda x \in \operatorname{pts}(s).\; (-1 \cdot a(x)))), (\lambda {[}\operatorname{pec\_}, a] \in (\mathbb{R} \times \operatorname{cont\text{-}fns}(s)).\; (\lambda x \in \operatorname{pts}(s).\; (\operatorname{pec\_} \cdot a(x))))]) \end{array}}}
\vnbprop{pex-vadd-val}{\vnbstmt{\begin{array}{@{}l@{}} \forall s.; \forall f, g \in \operatorname{cont\text{-}fns}(s). \\ \quad \operatorname{is\text{-}metric\text{-}space}(s) \Rightarrow  \\ \quad \quad \operatorname{vadd}({[}\operatorname{normed\text{-}field\text{-}as\text{-}commutative\text{-}ring}(\operatorname{rr\text{-}normed\text{-}field}), \operatorname{cont\text{-}fns}(s), (\lambda {[}a, \operatorname{peb\_}] \in (\operatorname{cont\text{-}fns}(s) \times \operatorname{cont\text{-}fns}(s)).\; (\lambda x \in \operatorname{pts}(s).\; (a(x) + \operatorname{peb\_}(x)))), (\lambda x \in \operatorname{pts}(s).\; 0), (\lambda a \in \operatorname{cont\text{-}fns}(s).\; (\lambda x \in \operatorname{pts}(s).\; (-1 \cdot a(x)))), (\lambda {[}\operatorname{pec\_}, a] \in (\mathbb{R} \times \operatorname{cont\text{-}fns}(s)).\; (\lambda x \in \operatorname{pts}(s).\; (\operatorname{pec\_} \cdot a(x))))])(f, g) \\ \quad \quad =\; (\lambda x \in \operatorname{pts}(s).\; (f(x) + g(x))) \end{array}}}
\vnbprop{pex-act-val}{\vnbstmt{\begin{array}{@{}l@{}} \forall s.; \forall c \in \mathbb{R}, f \in \operatorname{cont\text{-}fns}(s). \\ \quad \operatorname{is\text{-}metric\text{-}space}(s) \Rightarrow  \\ \quad \quad \operatorname{act}({[}\operatorname{normed\text{-}field\text{-}as\text{-}commutative\text{-}ring}(\operatorname{rr\text{-}normed\text{-}field}), \operatorname{cont\text{-}fns}(s), (\lambda {[}a, \operatorname{peb\_}] \in (\operatorname{cont\text{-}fns}(s) \times \operatorname{cont\text{-}fns}(s)).\; (\lambda x \in \operatorname{pts}(s).\; (a(x) + \operatorname{peb\_}(x)))), (\lambda x \in \operatorname{pts}(s).\; 0), (\lambda a \in \operatorname{cont\text{-}fns}(s).\; (\lambda x \in \operatorname{pts}(s).\; (-1 \cdot a(x)))), (\lambda {[}\operatorname{pec\_}, a] \in (\mathbb{R} \times \operatorname{cont\text{-}fns}(s)).\; (\lambda x \in \operatorname{pts}(s).\; (\operatorname{pec\_} \cdot a(x))))])(c, f) \\ \quad \quad =\; (\lambda x \in \operatorname{pts}(s).\; (c \cdot f(x))) \end{array}}}
\vnbprop{pex-cf-pin}{\vnbstmt{\begin{array}{@{}l@{}} \forall s. \\ \quad \operatorname{is\text{-}metric\text{-}space}(s) \Rightarrow  \\ \quad \quad \operatorname{scal}({[}\operatorname{normed\text{-}field\text{-}as\text{-}commutative\text{-}ring}(\operatorname{rr\text{-}normed\text{-}field}), \operatorname{cont\text{-}fns}(s), (\lambda {[}a, \operatorname{peb\_}] \in (\operatorname{cont\text{-}fns}(s) \times \operatorname{cont\text{-}fns}(s)).\; (\lambda x \in \operatorname{pts}(s).\; (a(x) + \operatorname{peb\_}(x)))), (\lambda x \in \operatorname{pts}(s).\; 0), (\lambda a \in \operatorname{cont\text{-}fns}(s).\; (\lambda x \in \operatorname{pts}(s).\; (-1 \cdot a(x)))), (\lambda {[}\operatorname{pec\_}, a] \in (\mathbb{R} \times \operatorname{cont\text{-}fns}(s)).\; (\lambda x \in \operatorname{pts}(s).\; (\operatorname{pec\_} \cdot a(x))))]) \\ \quad \quad =\; \operatorname{normed\text{-}field\text{-}as\text{-}commutative\text{-}ring}(\operatorname{rr\text{-}normed\text{-}field}) \end{array}}}
\vnbprop{pex-vs-zero}{\vnbstmt{\begin{array}{@{}l@{}} \forall s, a.; \forall \ell \in (a \to \mathbb{R}). \\ \quad \operatorname{is\text{-}metric\text{-}space}(s) \wedge \\ \quad \operatorname{is\text{-}compact}(s) \wedge \\ \quad a \subseteq \operatorname{cont\text{-}fns}(s) \wedge \\ \quad \forall f, g \in a.\; ((\lambda x \in \operatorname{pts}(s).\; (f(x) + g(x))) \in a) \wedge \\ \quad \forall c \in \mathbb{R}, f \in a.\; ((\lambda x \in \operatorname{pts}(s).\; (c \cdot f(x))) \in a) \wedge \\ \quad (\lambda x \in \operatorname{pts}(s).\; 1) \in a \wedge \\ \quad \forall f, g \in a. \\ \quad \quad \ell((\lambda x \in \operatorname{pts}(s).\; (f(x) + g(x)))) = (\ell(f) + \ell(g)) \wedge \\ \quad \forall c \in \mathbb{R}, f \in a.\; (\ell((\lambda x \in \operatorname{pts}(s).\; (c \cdot f(x)))) = (c \cdot \ell(f))) \wedge \\ \quad \forall f \in a.\; (\forall b \in \operatorname{pts}(s).\; (0 \leq f(b)) \Rightarrow (0 \leq \ell(f))) \Rightarrow \\ \quad \quad ((\lambda x \in \operatorname{pts}(s).\; 0) \in a) \wedge (\ell((\lambda x \in \operatorname{pts}(s).\; 0)) = 0) \end{array}}}
\vnbprop{pex-lm-mono}{\vnbstmt{\begin{array}{@{}l@{}} \forall s, a.; \forall \ell \in (a \to \mathbb{R}), f, g \in a. \\ \quad \operatorname{is\text{-}metric\text{-}space}(s) \wedge \\ \quad \operatorname{is\text{-}compact}(s) \wedge \\ \quad a \subseteq \operatorname{cont\text{-}fns}(s) \wedge \\ \quad \forall f, g \in a.\; ((\lambda x \in \operatorname{pts}(s).\; (f(x) + g(x))) \in a) \wedge \\ \quad \forall c \in \mathbb{R}, f \in a.\; ((\lambda x \in \operatorname{pts}(s).\; (c \cdot f(x))) \in a) \wedge \\ \quad (\lambda x \in \operatorname{pts}(s).\; 1) \in a \wedge \\ \quad \forall f, g \in a. \\ \quad \quad \ell((\lambda x \in \operatorname{pts}(s).\; (f(x) + g(x)))) = (\ell(f) + \ell(g)) \wedge \\ \quad \forall c \in \mathbb{R}, f \in a.\; (\ell((\lambda x \in \operatorname{pts}(s).\; (c \cdot f(x)))) = (c \cdot \ell(f))) \wedge \\ \quad \forall f \in a.\; (\forall b \in \operatorname{pts}(s).\; (0 \leq f(b)) \Rightarrow (0 \leq \ell(f))) \wedge \\ \quad \forall b \in \operatorname{pts}(s).\; (f(b) \leq g(b)) \Rightarrow \\ \quad \quad \ell(f) \leq \ell(g) \end{array}}}
\vnbprop{pex-vs-submodule}{\vnbstmt{\begin{array}{@{}l@{}} \forall s, a.; \forall \ell \in (a \to \mathbb{R}). \\ \quad \operatorname{is\text{-}metric\text{-}space}(s) \wedge \\ \quad \operatorname{is\text{-}compact}(s) \wedge \\ \quad a \subseteq \operatorname{cont\text{-}fns}(s) \wedge \\ \quad \forall f, g \in a.\; ((\lambda x \in \operatorname{pts}(s).\; (f(x) + g(x))) \in a) \wedge \\ \quad \forall c \in \mathbb{R}, f \in a.\; ((\lambda x \in \operatorname{pts}(s).\; (c \cdot f(x))) \in a) \wedge \\ \quad (\lambda x \in \operatorname{pts}(s).\; 1) \in a \wedge \\ \quad \forall f, g \in a. \\ \quad \quad \ell((\lambda x \in \operatorname{pts}(s).\; (f(x) + g(x)))) = (\ell(f) + \ell(g)) \wedge \\ \quad \forall c \in \mathbb{R}, f \in a.\; (\ell((\lambda x \in \operatorname{pts}(s).\; (c \cdot f(x)))) = (c \cdot \ell(f))) \wedge \\ \quad \forall f \in a.\; (\forall b \in \operatorname{pts}(s).\; (0 \leq f(b)) \Rightarrow (0 \leq \ell(f))) \Rightarrow \\ \quad \quad \operatorname{is\text{-}submodule}({[}\operatorname{normed\text{-}field\text{-}as\text{-}commutative\text{-}ring}(\operatorname{rr\text{-}normed\text{-}field}), \operatorname{cont\text{-}fns}(s), (\lambda {[}\mathit{pea}, \operatorname{peb\_}] \in (\operatorname{cont\text{-}fns}(s) \times \operatorname{cont\text{-}fns}(s)).\; (\lambda x \in \operatorname{pts}(s).\; (\operatorname{pea}(x) + \operatorname{peb\_}(x)))), (\lambda x \in \operatorname{pts}(s).\; 0), (\lambda \mathit{pea} \in \operatorname{cont\text{-}fns}(s).\; (\lambda x \in \operatorname{pts}(s).\; (-1 \cdot \operatorname{pea}(x)))), (\lambda {[}\operatorname{pec\_}, \mathit{pea}] \in (\mathbb{R} \times \operatorname{cont\text{-}fns}(s)).\; (\lambda x \in \operatorname{pts}(s).\; (\operatorname{pec\_} \cdot \operatorname{pea}(x))))], \\ \quad \quad \quad a) \end{array}}}
\vnbprop{pex-lm-linear}{\vnbstmt{\begin{array}{@{}l@{}} \forall s, a.; \forall \ell \in (a \to \mathbb{R}). \\ \quad \operatorname{is\text{-}metric\text{-}space}(s) \wedge \\ \quad \operatorname{is\text{-}compact}(s) \wedge \\ \quad a \subseteq \operatorname{cont\text{-}fns}(s) \wedge \\ \quad \forall f, g \in a.\; ((\lambda x \in \operatorname{pts}(s).\; (f(x) + g(x))) \in a) \wedge \\ \quad \forall c \in \mathbb{R}, f \in a.\; ((\lambda x \in \operatorname{pts}(s).\; (c \cdot f(x))) \in a) \wedge \\ \quad (\lambda x \in \operatorname{pts}(s).\; 1) \in a \wedge \\ \quad \forall f, g \in a. \\ \quad \quad \ell((\lambda x \in \operatorname{pts}(s).\; (f(x) + g(x)))) = (\ell(f) + \ell(g)) \wedge \\ \quad \forall c \in \mathbb{R}, f \in a.\; (\ell((\lambda x \in \operatorname{pts}(s).\; (c \cdot f(x)))) = (c \cdot \ell(f))) \wedge \\ \quad \forall f \in a.\; (\forall b \in \operatorname{pts}(s).\; (0 \leq f(b)) \Rightarrow (0 \leq \ell(f))) \Rightarrow \\ \quad \quad \operatorname{is\text{-}linear\text{-}functional\text{-}on}({[}\operatorname{normed\text{-}field\text{-}as\text{-}commutative\text{-}ring}(\operatorname{rr\text{-}normed\text{-}field}), \operatorname{cont\text{-}fns}(s), (\lambda {[}\mathit{pea}, \operatorname{peb\_}] \in (\operatorname{cont\text{-}fns}(s) \times \operatorname{cont\text{-}fns}(s)).\; (\lambda x \in \operatorname{pts}(s).\; (\operatorname{pea}(x) + \operatorname{peb\_}(x)))), (\lambda x \in \operatorname{pts}(s).\; 0), (\lambda \mathit{pea} \in \operatorname{cont\text{-}fns}(s).\; (\lambda x \in \operatorname{pts}(s).\; (-1 \cdot \operatorname{pea}(x)))), (\lambda {[}\operatorname{pec\_}, \mathit{pea}] \in (\mathbb{R} \times \operatorname{cont\text{-}fns}(s)).\; (\lambda x \in \operatorname{pts}(s).\; (\operatorname{pec\_} \cdot \operatorname{pea}(x))))], \\ \quad \quad \quad a, \ell) \end{array}}}
\vnbprop{pex-iif-mem}{\vnbstmt{\begin{array}{@{}l@{}} \forall s, a.; \forall g \in \operatorname{cont\text{-}fns}(s), v. \\ \quad (\operatorname{is\text{-}metric\text{-}space}(s) \wedge (a \subseteq \operatorname{cont\text{-}fns}(s))) \Rightarrow  \\ \quad \quad v \in (\lambda \mathit{peg} \in \operatorname{cont\text{-}fns}(s).\; \{\mathit{pev} \in a : \forall x \in \operatorname{pts}(s).\; (\operatorname{peg}(x) \leq \operatorname{pev}(x))\})(g) \\ \quad \quad \Leftrightarrow\; v \in \{\mathit{pev} \in a : \forall x \in \operatorname{pts}(s).\; (g(x) \leq \operatorname{pev}(x))\} \end{array}}}
\vnbprop{pex-p-exists}{\vnbstmt{\begin{array}{@{}l@{}} \forall s, a.; \forall \ell \in (a \to \mathbb{R}). \\ \quad \operatorname{is\text{-}metric\text{-}space}(s) \wedge \\ \quad \operatorname{is\text{-}compact}(s) \wedge \\ \quad a \subseteq \operatorname{cont\text{-}fns}(s) \wedge \\ \quad \forall f, g \in a.\; ((\lambda \mathit{sx} \in \operatorname{pts}(s).\; (f(\mathit{sx}) + g(\mathit{sx}))) \in a) \wedge \\ \quad \forall c \in \mathbb{R}, f \in a.\; ((\lambda \mathit{sx} \in \operatorname{pts}(s).\; (c \cdot f(\mathit{sx}))) \in a) \wedge \\ \quad (\lambda \mathit{sx} \in \operatorname{pts}(s).\; 1) \in a \wedge \\ \quad \forall f, g \in a. \\ \quad \quad \ell((\lambda \mathit{sx} \in \operatorname{pts}(s).\; (f(\mathit{sx}) + g(\mathit{sx})))) = (\ell(f) + \ell(g)) \wedge \\ \quad \forall c \in \mathbb{R}, f \in a.\; (\ell((\lambda \mathit{sx} \in \operatorname{pts}(s).\; (c \cdot f(\mathit{sx})))) = (c \cdot \ell(f))) \wedge \\ \quad \forall f \in a.\; (\forall b \in \operatorname{pts}(s).\; (0 \leq f(b)) \Rightarrow (0 \leq \ell(f))) \wedge \\ \quad \forall g \in \operatorname{cont\text{-}fns}(s). \exists d \in \mathbb{R}. \\ \quad \quad (0 \leq d) \wedge \forall x \in \operatorname{pts}(s).\; (\lvert g(x) \rvert \leq d) \Rightarrow \\ \quad \quad \exists p \in (\operatorname{cont\text{-}fns}(s) \to \mathbb{R}). \\ \quad \quad \quad \forall g \in \operatorname{cont\text{-}fns}(s), v \in a. \\ \quad \quad \quad \quad \forall \mathit{pex} \in \operatorname{pts}(s).\; (g(\mathit{pex}) \leq v(\mathit{pex})) \Rightarrow (p(g) \leq \ell(v)) \wedge \\ \quad \quad \quad \forall g \in \operatorname{cont\text{-}fns}(s), \beta \in \mathbb{R}. \\ \quad \quad \quad \quad \forall v \in a.\; (\forall \mathit{pex} \in \operatorname{pts}(s).\; (g(\mathit{pex}) \leq v(\mathit{pex})) \Rightarrow (\beta \leq \ell(v))) \Rightarrow  \\ \quad \quad \quad \quad \quad \beta \leq p(g) \end{array}}}
\vnbprop{pex-p-props}{\vnbstmt{\begin{array}{@{}l@{}} \forall s, a.; \forall \ell \in (a \to \mathbb{R}), p \in (\operatorname{cont\text{-}fns}(s) \to \mathbb{R}). \\ \quad \operatorname{is\text{-}metric\text{-}space}(s) \wedge \\ \quad \operatorname{is\text{-}compact}(s) \wedge \\ \quad a \subseteq \operatorname{cont\text{-}fns}(s) \wedge \\ \quad \forall f, g \in a.\; ((\lambda x \in \operatorname{pts}(s).\; (f(x) + g(x))) \in a) \wedge \\ \quad \forall c \in \mathbb{R}, f \in a.\; ((\lambda x \in \operatorname{pts}(s).\; (c \cdot f(x))) \in a) \wedge \\ \quad (\lambda x \in \operatorname{pts}(s).\; 1) \in a \wedge \\ \quad \forall f, g \in a. \\ \quad \quad \ell((\lambda x \in \operatorname{pts}(s).\; (f(x) + g(x)))) = (\ell(f) + \ell(g)) \wedge \\ \quad \forall c \in \mathbb{R}, f \in a.\; (\ell((\lambda x \in \operatorname{pts}(s).\; (c \cdot f(x)))) = (c \cdot \ell(f))) \wedge \\ \quad \forall f \in a.\; (\forall b \in \operatorname{pts}(s).\; (0 \leq f(b)) \Rightarrow (0 \leq \ell(f))) \wedge \\ \quad \forall g \in \operatorname{cont\text{-}fns}(s), v \in a. \\ \quad \quad \forall \mathit{pex} \in \operatorname{pts}(s).\; (g(\mathit{pex}) \leq v(\mathit{pex})) \Rightarrow (p(g) \leq \ell(v)) \wedge \\ \quad \forall g \in \operatorname{cont\text{-}fns}(s), \beta \in \mathbb{R}. \\ \quad \quad \forall v \in a.\; (\forall \mathit{pex} \in \operatorname{pts}(s).\; (g(\mathit{pex}) \leq v(\mathit{pex})) \Rightarrow (\beta \leq \ell(v))) \Rightarrow  \\ \quad \quad \quad \beta \leq p(g) \Rightarrow \\ \quad \quad \operatorname{is\text{-}sublinear}({[}\operatorname{normed\text{-}field\text{-}as\text{-}commutative\text{-}ring}(\operatorname{rr\text{-}normed\text{-}field}), \operatorname{cont\text{-}fns}(s), (\lambda {[}\mathit{pea}, \operatorname{peb\_}] \in (\operatorname{cont\text{-}fns}(s) \times \operatorname{cont\text{-}fns}(s)).\; (\lambda x \in \operatorname{pts}(s).\; (\operatorname{pea}(x) + \operatorname{peb\_}(x)))), (\lambda x \in \operatorname{pts}(s).\; 0), (\lambda \mathit{pea} \in \operatorname{cont\text{-}fns}(s).\; (\lambda x \in \operatorname{pts}(s).\; (-1 \cdot \operatorname{pea}(x)))), (\lambda {[}\operatorname{pec\_}, \mathit{pea}] \in (\mathbb{R} \times \operatorname{cont\text{-}fns}(s)).\; (\lambda x \in \operatorname{pts}(s).\; (\operatorname{pec\_} \cdot \operatorname{pea}(x))))], \\ \quad \quad \quad p) \wedge \\ \quad \quad \forall v \in a.\; (\ell(v) \leq p(v)) \wedge \\ \quad \quad \forall g \in \operatorname{cont\text{-}fns}(s).\; (\forall b \in \operatorname{pts}(s).\; (g(b) \leq 0) \Rightarrow (p(g) \leq 0)) \end{array}}}
\vnbprop{pex-extension-from-bound}{\vnbstmt{\begin{array}{@{}l@{}} \forall s, a.; \forall \ell \in (a \to \mathbb{R}). \\ \quad \operatorname{is\text{-}metric\text{-}space}(s) \wedge \\ \quad \operatorname{is\text{-}compact}(s) \wedge \\ \quad a \subseteq \operatorname{cont\text{-}fns}(s) \wedge \\ \quad \forall f, g \in a.\; ((\lambda \mathit{sx} \in \operatorname{pts}(s).\; (f(\mathit{sx}) + g(\mathit{sx}))) \in a) \wedge \\ \quad \forall c \in \mathbb{R}, f \in a.\; ((\lambda \mathit{sx} \in \operatorname{pts}(s).\; (c \cdot f(\mathit{sx}))) \in a) \wedge \\ \quad (\lambda \mathit{sx} \in \operatorname{pts}(s).\; 1) \in a \wedge \\ \quad \forall f, g \in a. \\ \quad \quad \ell((\lambda \mathit{sx} \in \operatorname{pts}(s).\; (f(\mathit{sx}) + g(\mathit{sx})))) = (\ell(f) + \ell(g)) \wedge \\ \quad \forall c \in \mathbb{R}, f \in a.\; (\ell((\lambda \mathit{sx} \in \operatorname{pts}(s).\; (c \cdot f(\mathit{sx})))) = (c \cdot \ell(f))) \wedge \\ \quad \forall f \in a.\; (\forall b \in \operatorname{pts}(s).\; (0 \leq f(b)) \Rightarrow (0 \leq \ell(f))) \wedge \\ \quad \forall g \in \operatorname{cont\text{-}fns}(s). \exists d \in \mathbb{R}. \\ \quad \quad (0 \leq d) \wedge \forall x \in \operatorname{pts}(s).\; (\lvert g(x) \rvert \leq d) \Rightarrow \\ \quad \quad \exists m.\; (\operatorname{is\text{-}positive\text{-}functional}(s, m) \wedge \operatorname{extends\text{-}on}(a, m, \ell)) \end{array}}}
\vnbprop{pex-positive-extension}{\vnbstmt{\begin{array}{@{}l@{}} \forall s, a.; \forall \ell \in (a \to \mathbb{R}). \\ \quad \operatorname{is\text{-}metric\text{-}space}(s) \wedge \\ \quad \operatorname{is\text{-}compact}(s) \wedge \\ \quad a \subseteq \operatorname{cont\text{-}fns}(s) \wedge \\ \quad \forall f, g \in a.\; ((\lambda x \in \operatorname{pts}(s).\; (f(x) + g(x))) \in a) \wedge \\ \quad \forall c \in \mathbb{R}, f \in a.\; ((\lambda x \in \operatorname{pts}(s).\; (c \cdot f(x))) \in a) \wedge \\ \quad (\lambda x \in \operatorname{pts}(s).\; 1) \in a \wedge \\ \quad \forall f, g \in a. \\ \quad \quad \ell((\lambda x \in \operatorname{pts}(s).\; (f(x) + g(x)))) = (\ell(f) + \ell(g)) \wedge \\ \quad \forall c \in \mathbb{R}, f \in a.\; (\ell((\lambda x \in \operatorname{pts}(s).\; (c \cdot f(x)))) = (c \cdot \ell(f))) \wedge \\ \quad \forall f \in a.\; (\forall b \in \operatorname{pts}(s).\; (0 \leq f(b)) \Rightarrow (0 \leq \ell(f))) \Rightarrow \\ \quad \quad \exists m.\; (\operatorname{is\text{-}positive\text{-}functional}(s, m) \wedge \operatorname{extends\text{-}on}(a, m, \ell)) \end{array}}}

\vnbbookfile{theorem-library/lebesgue-measure-line}{theorem-library/lebesgue-measure-line}
\noindent{\small ==================================================================== theorem-library/lebesgue-measure-line.scm -- MEASURE-3 (the join), 2026-10-09.}\par
\vnbprop{lebesgue-measure-ccint}{\vnbstmt{\begin{array}{@{}l@{}} \forall a, b. \\ \quad ((a \in \mathbb{R}) \wedge ((b \in \mathbb{R}) \wedge (a \leq b))) \Rightarrow  \\ \quad \quad \operatorname{ccint}(a, b) \\ \quad \quad \in\; \operatorname{caratheodory\text{-}sets}(\mathbb{R}, \\ \quad \quad \quad \operatorname{method\text{-}i\text{-}outer}(\mathbb{R}, \operatorname{closed\text{-}intervals}, (\lambda s \in \operatorname{closed\text{-}intervals}.\; \operatorname{interval\text{-}length}(s)))) \wedge \\ \quad \quad \operatorname{res}(\operatorname{method\text{-}i\text{-}outer}(\mathbb{R}, \operatorname{closed\text{-}intervals}, (\lambda s \in \operatorname{closed\text{-}intervals}.\; \operatorname{interval\text{-}length}(s))), \operatorname{caratheodory\text{-}sets}(\mathbb{R}, \operatorname{method\text{-}i\text{-}outer}(\mathbb{R}, \operatorname{closed\text{-}intervals}, (\lambda s \in \operatorname{closed\text{-}intervals}.\; \operatorname{interval\text{-}length}(s)))))(\operatorname{ccint}(a, b)) \\ \quad \quad =\; (b - a) \end{array}}}
\vnbprop{lebesgue-measure-ooint}{\vnbstmt{\begin{array}{@{}l@{}} \forall a, b. \\ \quad ((a \in \mathbb{R}) \wedge ((b \in \mathbb{R}) \wedge (a \leq b))) \Rightarrow  \\ \quad \quad \operatorname{ooint}(a, b) \\ \quad \quad \in\; \operatorname{caratheodory\text{-}sets}(\mathbb{R}, \\ \quad \quad \quad \operatorname{method\text{-}i\text{-}outer}(\mathbb{R}, \operatorname{closed\text{-}intervals}, (\lambda s \in \operatorname{closed\text{-}intervals}.\; \operatorname{interval\text{-}length}(s)))) \wedge \\ \quad \quad \operatorname{res}(\operatorname{method\text{-}i\text{-}outer}(\mathbb{R}, \operatorname{closed\text{-}intervals}, (\lambda s \in \operatorname{closed\text{-}intervals}.\; \operatorname{interval\text{-}length}(s))), \operatorname{caratheodory\text{-}sets}(\mathbb{R}, \operatorname{method\text{-}i\text{-}outer}(\mathbb{R}, \operatorname{closed\text{-}intervals}, (\lambda s \in \operatorname{closed\text{-}intervals}.\; \operatorname{interval\text{-}length}(s)))))(\operatorname{ooint}(a, b)) \\ \quad \quad =\; (b - a) \end{array}}}
\vnbprop{lebesgue-measure-coint}{\vnbstmt{\begin{array}{@{}l@{}} \forall a, b. \\ \quad ((a \in \mathbb{R}) \wedge ((b \in \mathbb{R}) \wedge (a \leq b))) \Rightarrow  \\ \quad \quad \operatorname{coint}(a, b) \\ \quad \quad \in\; \operatorname{caratheodory\text{-}sets}(\mathbb{R}, \\ \quad \quad \quad \operatorname{method\text{-}i\text{-}outer}(\mathbb{R}, \operatorname{closed\text{-}intervals}, (\lambda s \in \operatorname{closed\text{-}intervals}.\; \operatorname{interval\text{-}length}(s)))) \wedge \\ \quad \quad \operatorname{res}(\operatorname{method\text{-}i\text{-}outer}(\mathbb{R}, \operatorname{closed\text{-}intervals}, (\lambda s \in \operatorname{closed\text{-}intervals}.\; \operatorname{interval\text{-}length}(s))), \operatorname{caratheodory\text{-}sets}(\mathbb{R}, \operatorname{method\text{-}i\text{-}outer}(\mathbb{R}, \operatorname{closed\text{-}intervals}, (\lambda s \in \operatorname{closed\text{-}intervals}.\; \operatorname{interval\text{-}length}(s)))))(\operatorname{coint}(a, b)) \\ \quad \quad =\; (b - a) \end{array}}}
\vnbprop{lebesgue-measure-ocint}{\vnbstmt{\begin{array}{@{}l@{}} \forall a, b. \\ \quad ((a \in \mathbb{R}) \wedge ((b \in \mathbb{R}) \wedge (a \leq b))) \Rightarrow  \\ \quad \quad \operatorname{ocint}(a, b) \\ \quad \quad \in\; \operatorname{caratheodory\text{-}sets}(\mathbb{R}, \\ \quad \quad \quad \operatorname{method\text{-}i\text{-}outer}(\mathbb{R}, \operatorname{closed\text{-}intervals}, (\lambda s \in \operatorname{closed\text{-}intervals}.\; \operatorname{interval\text{-}length}(s)))) \wedge \\ \quad \quad \operatorname{res}(\operatorname{method\text{-}i\text{-}outer}(\mathbb{R}, \operatorname{closed\text{-}intervals}, (\lambda s \in \operatorname{closed\text{-}intervals}.\; \operatorname{interval\text{-}length}(s))), \operatorname{caratheodory\text{-}sets}(\mathbb{R}, \operatorname{method\text{-}i\text{-}outer}(\mathbb{R}, \operatorname{closed\text{-}intervals}, (\lambda s \in \operatorname{closed\text{-}intervals}.\; \operatorname{interval\text{-}length}(s)))))(\operatorname{ocint}(a, b)) \\ \quad \quad =\; (b - a) \end{array}}}
\vnbprop{lebesgue-measure-point}{\vnbstmt{\begin{array}{@{}l@{}} \forall x \in \mathbb{R}. \\ \quad \operatorname{singleton}(x) \\ \quad \in\; \operatorname{caratheodory\text{-}sets}(\mathbb{R}, \\ \quad \quad \operatorname{method\text{-}i\text{-}outer}(\mathbb{R}, \operatorname{closed\text{-}intervals}, (\lambda s \in \operatorname{closed\text{-}intervals}.\; \operatorname{interval\text{-}length}(s)))) \wedge \\ \quad \operatorname{res}(\operatorname{method\text{-}i\text{-}outer}(\mathbb{R}, \operatorname{closed\text{-}intervals}, (\lambda s \in \operatorname{closed\text{-}intervals}.\; \operatorname{interval\text{-}length}(s))), \operatorname{caratheodory\text{-}sets}(\mathbb{R}, \operatorname{method\text{-}i\text{-}outer}(\mathbb{R}, \operatorname{closed\text{-}intervals}, (\lambda s \in \operatorname{closed\text{-}intervals}.\; \operatorname{interval\text{-}length}(s)))))(\operatorname{singleton}(x)) \\ \quad =\; 0 \end{array}}}

\vnbbookfile{structure-library/lebesgue-translate}{structure-library/lebesgue-translate}
\noindent{\small lebesgue-translate.scm -- the TRANSLATE of a set of reals.  Vocabulary only: no proof, no support.  Agent MEASURE-4-A, 2026-10-10 (scratchpad/triage/MEASURE-4-A.md). The proofs that use it are theorem-library/lebesgue-measure-unique.scm.}\par
\par\medskip\noindent\textbf{Definition} (\texttt{translate}).\;$(\mathit{es} + c) :=$ $\{\operatorname{ltz\_} \in \mathbb{R} : ((\operatorname{ltz\_} - c) \in \mathit{es})\}$

\vnbbookfile{theorem-library/lebesgue-measure-unique}{theorem-library/lebesgue-measure-unique}
\noindent{\small ==================================================================== theorem-library/lebesgue-measure-unique.scm -- MEASURE-4-A, 2026-10-10. Lebesgue measure on the line: translation invariance and uniqueness.}\par
\vnbprop{translate-unfold}{\vnbstmt{\forall s, c.\; ((s + c) \simeq \{z \in \mathbb{R} : ((z - c) \in s)\})}}
\vnbprop{translate-membership}{\vnbstmt{\forall s, c, z.\; ((z \in (s + c)) \Leftrightarrow ((z \in \mathbb{R}) \wedge ((z - c) \in s)))}}
\vnbprop{translate-subset-rr}{\vnbstmt{\forall s, c.\; ((s + c) \subseteq \mathbb{R})}}
\vnbprop{lebesgue-outer-translate-invariant}{\vnbstmt{\begin{array}{@{}l@{}} \forall s, c. \\ \quad ((s \subseteq \mathbb{R}) \wedge (c \in \mathbb{R})) \Rightarrow  \\ \quad \quad \operatorname{method\text{-}i\text{-}outer}(\mathbb{R}, \operatorname{closed\text{-}intervals}, (\lambda \mathit{ols} \in \operatorname{closed\text{-}intervals}.\; \operatorname{interval\text{-}length}(\mathit{ols})))((s + c)) \\ \quad \quad =\; \operatorname{method\text{-}i\text{-}outer}(\mathbb{R}, \operatorname{closed\text{-}intervals}, (\lambda \mathit{ols} \in \operatorname{closed\text{-}intervals}.\; \operatorname{interval\text{-}length}(\mathit{ols})))(s) \end{array}}}
\vnbprop{translate-intersection}{\vnbstmt{\begin{array}{@{}l@{}} \forall a, b.; \forall c \in \mathbb{R}. \\ \quad ((a \cap b) + c) = ((a + c) \cap (b + c)) \end{array}}}
\vnbprop{translate-difference}{\vnbstmt{\begin{array}{@{}l@{}} \forall a, b.; \forall c \in \mathbb{R}. \\ \quad (\operatorname{difference}(a, b) + c) = \operatorname{difference}((a + c), (b + c)) \end{array}}}
\vnbprop{translate-neg-back}{\vnbstmt{\begin{array}{@{}l@{}} \forall a, c. \\ \quad ((a \subseteq \mathbb{R}) \wedge (c \in \mathbb{R})) \Rightarrow (((a + (0 - c)) + c) = a) \end{array}}}
\vnbprop{lebesgue-measurable-translate}{\vnbstmt{\begin{array}{@{}l@{}} \forall s \in \operatorname{caratheodory\text{-}sets}(\mathbb{R}, \operatorname{method\text{-}i\text{-}outer}(\mathbb{R}, \operatorname{closed\text{-}intervals}, (\lambda \mathit{ols} \in \operatorname{closed\text{-}intervals}.\; \operatorname{interval\text{-}length}(\mathit{ols})))), c \in \mathbb{R}. \\ \quad (s + c) \\ \quad \in\; \operatorname{caratheodory\text{-}sets}(\mathbb{R}, \\ \quad \quad \operatorname{method\text{-}i\text{-}outer}(\mathbb{R}, \operatorname{closed\text{-}intervals}, (\lambda \mathit{ols} \in \operatorname{closed\text{-}intervals}.\; \operatorname{interval\text{-}length}(\mathit{ols})))) \wedge \\ \quad \operatorname{res}(\operatorname{method\text{-}i\text{-}outer}(\mathbb{R}, \operatorname{closed\text{-}intervals}, (\lambda \mathit{ols} \in \operatorname{closed\text{-}intervals}.\; \operatorname{interval\text{-}length}(\mathit{ols}))), \operatorname{caratheodory\text{-}sets}(\mathbb{R}, \operatorname{method\text{-}i\text{-}outer}(\mathbb{R}, \operatorname{closed\text{-}intervals}, (\lambda \mathit{ols} \in \operatorname{closed\text{-}intervals}.\; \operatorname{interval\text{-}length}(\mathit{ols})))))((s + c)) \\ \quad =\; \operatorname{res}(\operatorname{method\text{-}i\text{-}outer}(\mathbb{R}, \operatorname{closed\text{-}intervals}, (\lambda \mathit{ols} \in \operatorname{closed\text{-}intervals}.\; \operatorname{interval\text{-}length}(\mathit{ols}))), \operatorname{caratheodory\text{-}sets}(\mathbb{R}, \operatorname{method\text{-}i\text{-}outer}(\mathbb{R}, \operatorname{closed\text{-}intervals}, (\lambda \mathit{ols} \in \operatorname{closed\text{-}intervals}.\; \operatorname{interval\text{-}length}(\mathit{ols})))))(s) \end{array}}}
\vnbprop{rr-closed-borel}{\vnbstmt{\forall l.\; (\operatorname{is\text{-}closed}(\operatorname{rr\text{-}ms}, l) \Rightarrow (l \in \operatorname{borel\text{-}sets}(\operatorname{rr\text{-}ms})))}}
\vnbprop{ccint-closed-rr-ms}{\vnbstmt{\forall a, b.\; (((a \in \mathbb{R}) \wedge (b \in \mathbb{R})) \Rightarrow \operatorname{is\text{-}closed}(\operatorname{rr\text{-}ms}, \operatorname{ccint}(a, b)))}}
\vnbprop{closed-intervals-borel}{\vnbstmt{\operatorname{closed\text{-}intervals} \subseteq \operatorname{borel\text{-}sets}(\operatorname{rr\text{-}ms})}}
\vnbprop{borel-measure-le-lebesgue-outer}{\vnbstmt{\begin{array}{@{}l@{}} \forall u, s. \\ \quad \operatorname{is\text{-}measure}(\mathbb{R}, \operatorname{borel\text{-}sets}(\operatorname{rr\text{-}ms}), u) \wedge \\ \quad \forall a, b.\; (((a \in \mathbb{R}) \wedge ((b \in \mathbb{R}) \wedge (a \leq b))) \Rightarrow (u(\operatorname{ccint}(a, b)) = (b - a))) \wedge \\ \quad s \in \operatorname{borel\text{-}sets}(\operatorname{rr\text{-}ms}) \Rightarrow \\ \quad \quad u(s) \\ \quad \quad \leq\; \operatorname{method\text{-}i\text{-}outer}(\mathbb{R}, \operatorname{closed\text{-}intervals}, (\lambda \mathit{ols} \in \operatorname{closed\text{-}intervals}.\; \operatorname{interval\text{-}length}(\mathit{ols})))(s) \end{array}}}
\vnbprop{borel-measure-eq-lebesgue-outer-bounded}{\vnbstmt{\begin{array}{@{}l@{}} \forall u, p, q, s. \\ \quad \operatorname{is\text{-}measure}(\mathbb{R}, \operatorname{borel\text{-}sets}(\operatorname{rr\text{-}ms}), u) \wedge \\ \quad \forall a, b.\; (((a \in \mathbb{R}) \wedge ((b \in \mathbb{R}) \wedge (a \leq b))) \Rightarrow (u(\operatorname{ccint}(a, b)) = (b - a))) \wedge \\ \quad p \in \mathbb{R} \wedge \\ \quad q \in \mathbb{R} \wedge \\ \quad p \leq q \wedge \\ \quad s \in \operatorname{borel\text{-}sets}(\operatorname{rr\text{-}ms}) \wedge \\ \quad s \subseteq \operatorname{ccint}(p, q) \Rightarrow \\ \quad \quad u(s) \\ \quad \quad =\; \operatorname{method\text{-}i\text{-}outer}(\mathbb{R}, \operatorname{closed\text{-}intervals}, (\lambda \mathit{ols} \in \operatorname{closed\text{-}intervals}.\; \operatorname{interval\text{-}length}(\mathit{ols})))(s) \end{array}}}
\vnbprop{measure-le-of-ccint-exhaust}{\vnbstmt{\begin{array}{@{}l@{}} \forall a, u, s, d. \\ \quad \operatorname{is\text{-}measure}(\mathbb{R}, a, u) \wedge \\ \quad \forall n \in \mathbb{N}.\; (\operatorname{ccint}((0 - n), n) \in a) \wedge \\ \quad s \in a \wedge \\ \quad d \in \operatorname{rr\text{-}pos\text{-}star} \wedge \\ \quad \forall n \in \mathbb{N}.\; (u((s \cap \operatorname{ccint}((0 - n), n))) \leq d) \Rightarrow \\ \quad \quad u(s) \leq d \end{array}}}
\vnbprop{lebesgue-measure-unique}{\vnbstmt{\begin{array}{@{}l@{}} \forall u, s. \\ \quad \operatorname{is\text{-}measure}(\mathbb{R}, \operatorname{borel\text{-}sets}(\operatorname{rr\text{-}ms}), u) \wedge \\ \quad \forall a, b.\; (((a \in \mathbb{R}) \wedge ((b \in \mathbb{R}) \wedge (a \leq b))) \Rightarrow (u(\operatorname{ccint}(a, b)) = (b - a))) \wedge \\ \quad s \in \operatorname{borel\text{-}sets}(\operatorname{rr\text{-}ms}) \Rightarrow \\ \quad \quad u(s) \\ \quad \quad =\; \operatorname{res}(\operatorname{method\text{-}i\text{-}outer}(\mathbb{R}, \operatorname{closed\text{-}intervals}, (\lambda \mathit{ols} \in \operatorname{closed\text{-}intervals}.\; \operatorname{interval\text{-}length}(\mathit{ols}))), \operatorname{caratheodory\text{-}sets}(\mathbb{R}, \operatorname{method\text{-}i\text{-}outer}(\mathbb{R}, \operatorname{closed\text{-}intervals}, (\lambda \mathit{ols} \in \operatorname{closed\text{-}intervals}.\; \operatorname{interval\text{-}length}(\mathit{ols})))))(s) \end{array}}}

\vnbbookfile{theorem-library/lebesgue-translation-invariant}{theorem-library/lebesgue-translation-invariant}
\noindent{\small ==================================================================== theorem-library/lebesgue-translation-invariant.scm -- MEASURE-5-A, 2026-10-10. A translation-invariant Borel measure on the line with 0 $<$ mu([0, 1]) = k $<$ +inf is k times Lebesgue measure on the Borel sets.}\par
\vnbprop{translate-ccint}{\vnbstmt{\begin{array}{@{}l@{}} \forall a, b, c. \\ \quad ((a \in \mathbb{R}) \wedge ((b \in \mathbb{R}) \wedge (c \in \mathbb{R}))) \Rightarrow  \\ \quad \quad (\operatorname{ccint}(a, b) + c) = \operatorname{ccint}((a + c), (b + c)) \end{array}}}
\vnbprop{invariant-measure-ccint-shift}{\vnbstmt{\begin{array}{@{}l@{}} \forall u, a, b. \\ \quad \operatorname{is\text{-}measure}(\mathbb{R}, \operatorname{borel\text{-}sets}(\operatorname{rr\text{-}ms}), u) \wedge \\ \quad \forall v, c. \\ \quad \quad ((v \in \operatorname{borel\text{-}sets}(\operatorname{rr\text{-}ms})) \wedge ((c \in \mathbb{R}) \wedge ((v + c) \in \operatorname{borel\text{-}sets}(\operatorname{rr\text{-}ms})))) \Rightarrow  \\ \quad \quad \quad u((v + c)) = u(v) \wedge \\ \quad a \in \mathbb{R} \wedge \\ \quad b \in \mathbb{R} \Rightarrow \\ \quad \quad u(\operatorname{ccint}(a, b)) = u(\operatorname{ccint}(0, (b - a))) \end{array}}}
\vnbprop{invariant-measure-points-null}{\vnbstmt{\begin{array}{@{}l@{}} \forall u.; \forall x \in \mathbb{R}. \\ \quad \operatorname{is\text{-}measure}(\mathbb{R}, \operatorname{borel\text{-}sets}(\operatorname{rr\text{-}ms}), u) \wedge \\ \quad \forall v, c. \\ \quad \quad ((v \in \operatorname{borel\text{-}sets}(\operatorname{rr\text{-}ms})) \wedge ((c \in \mathbb{R}) \wedge ((v + c) \in \operatorname{borel\text{-}sets}(\operatorname{rr\text{-}ms})))) \Rightarrow  \\ \quad \quad \quad u((v + c)) = u(v) \wedge \\ \quad u(\operatorname{ccint}(0, 1)) \in \mathbb{R} \Rightarrow \\ \quad \quad u(\operatorname{ccint}(x, x)) = 0 \end{array}}}
\vnbprop{null-points-ccint-split}{\vnbstmt{\begin{array}{@{}l@{}} \forall u, a, c, b. \\ \quad \operatorname{is\text{-}measure}(\mathbb{R}, \operatorname{borel\text{-}sets}(\operatorname{rr\text{-}ms}), u) \wedge \\ \quad \forall x \in \mathbb{R}.\; (u(\operatorname{ccint}(x, x)) = 0) \wedge \\ \quad a \in \mathbb{R} \wedge \\ \quad c \in \mathbb{R} \wedge \\ \quad b \in \mathbb{R} \wedge \\ \quad a \leq c \wedge \\ \quad c \leq b \Rightarrow \\ \quad \quad u(\operatorname{ccint}(a, b)) = \operatorname{eplus}(u(\operatorname{ccint}(a, c)), u(\operatorname{ccint}(c, b))) \end{array}}}
\vnbprop{invariant-measure-ccint-additive}{\vnbstmt{\begin{array}{@{}l@{}} \forall u, s, t. \\ \quad \operatorname{is\text{-}measure}(\mathbb{R}, \operatorname{borel\text{-}sets}(\operatorname{rr\text{-}ms}), u) \wedge \\ \quad \forall v, c. \\ \quad \quad ((v \in \operatorname{borel\text{-}sets}(\operatorname{rr\text{-}ms})) \wedge ((c \in \mathbb{R}) \wedge ((v + c) \in \operatorname{borel\text{-}sets}(\operatorname{rr\text{-}ms})))) \Rightarrow  \\ \quad \quad \quad u((v + c)) = u(v) \wedge \\ \quad u(\operatorname{ccint}(0, 1)) \in \mathbb{R} \wedge \\ \quad s \in \mathbb{R} \wedge \\ \quad t \in \mathbb{R} \wedge \\ \quad 0 \leq s \wedge \\ \quad 0 \leq t \Rightarrow \\ \quad \quad u(\operatorname{ccint}(0, (s + t))) = \operatorname{eplus}(u(\operatorname{ccint}(0, s)), u(\operatorname{ccint}(0, t))) \end{array}}}
\vnbprop{translation-invariant-measure-nn}{\vnbstmt{\begin{array}{@{}l@{}} \forall u.; \forall m \in \mathbb{N}. \\ \quad \operatorname{is\text{-}measure}(\mathbb{R}, \operatorname{borel\text{-}sets}(\operatorname{rr\text{-}ms}), u) \wedge \\ \quad \forall v, c. \\ \quad \quad ((v \in \operatorname{borel\text{-}sets}(\operatorname{rr\text{-}ms})) \wedge ((c \in \mathbb{R}) \wedge ((v + c) \in \operatorname{borel\text{-}sets}(\operatorname{rr\text{-}ms})))) \Rightarrow  \\ \quad \quad \quad u((v + c)) = u(v) \wedge \\ \quad u(\operatorname{ccint}(0, 1)) = 1 \Rightarrow \\ \quad \quad u(\operatorname{ccint}(0, m)) = m \end{array}}}
\vnbprop{translation-invariant-measure-finite}{\vnbstmt{\begin{array}{@{}l@{}} \forall u, t. \\ \quad \operatorname{is\text{-}measure}(\mathbb{R}, \operatorname{borel\text{-}sets}(\operatorname{rr\text{-}ms}), u) \wedge \\ \quad \forall v, c. \\ \quad \quad ((v \in \operatorname{borel\text{-}sets}(\operatorname{rr\text{-}ms})) \wedge ((c \in \mathbb{R}) \wedge ((v + c) \in \operatorname{borel\text{-}sets}(\operatorname{rr\text{-}ms})))) \Rightarrow  \\ \quad \quad \quad u((v + c)) = u(v) \wedge \\ \quad u(\operatorname{ccint}(0, 1)) = 1 \wedge \\ \quad t \in \mathbb{R} \wedge \\ \quad 0 \leq t \Rightarrow \\ \quad \quad u(\operatorname{ccint}(0, t)) \in \mathbb{R} \end{array}}}
\vnbprop{translation-invariant-measure-nn-mul}{\vnbstmt{\begin{array}{@{}l@{}} \forall j \in \mathbb{N}, u, s. \\ \quad \operatorname{is\text{-}measure}(\mathbb{R}, \operatorname{borel\text{-}sets}(\operatorname{rr\text{-}ms}), u) \wedge \\ \quad \forall v, c. \\ \quad \quad ((v \in \operatorname{borel\text{-}sets}(\operatorname{rr\text{-}ms})) \wedge ((c \in \mathbb{R}) \wedge ((v + c) \in \operatorname{borel\text{-}sets}(\operatorname{rr\text{-}ms})))) \Rightarrow  \\ \quad \quad \quad u((v + c)) = u(v) \wedge \\ \quad u(\operatorname{ccint}(0, 1)) = 1 \wedge \\ \quad s \in \mathbb{R} \wedge \\ \quad 0 \leq s \Rightarrow \\ \quad \quad u(\operatorname{ccint}(0, (j \cdot s))) = (j \cdot u(\operatorname{ccint}(0, s))) \end{array}}}
\vnbprop{translation-invariant-measure-ccint-zero}{\vnbstmt{\begin{array}{@{}l@{}} \forall u, t. \\ \quad \operatorname{is\text{-}measure}(\mathbb{R}, \operatorname{borel\text{-}sets}(\operatorname{rr\text{-}ms}), u) \wedge \\ \quad \forall v, c. \\ \quad \quad ((v \in \operatorname{borel\text{-}sets}(\operatorname{rr\text{-}ms})) \wedge ((c \in \mathbb{R}) \wedge ((v + c) \in \operatorname{borel\text{-}sets}(\operatorname{rr\text{-}ms})))) \Rightarrow  \\ \quad \quad \quad u((v + c)) = u(v) \wedge \\ \quad u(\operatorname{ccint}(0, 1)) = 1 \wedge \\ \quad t \in \mathbb{R} \wedge \\ \quad 0 \leq t \Rightarrow \\ \quad \quad u(\operatorname{ccint}(0, t)) = t \end{array}}}
\vnbprop{translation-invariant-measure-ccint}{\vnbstmt{\begin{array}{@{}l@{}} \forall u, a, b. \\ \quad \operatorname{is\text{-}measure}(\mathbb{R}, \operatorname{borel\text{-}sets}(\operatorname{rr\text{-}ms}), u) \wedge \\ \quad \forall v, c. \\ \quad \quad ((v \in \operatorname{borel\text{-}sets}(\operatorname{rr\text{-}ms})) \wedge ((c \in \mathbb{R}) \wedge ((v + c) \in \operatorname{borel\text{-}sets}(\operatorname{rr\text{-}ms})))) \Rightarrow  \\ \quad \quad \quad u((v + c)) = u(v) \wedge \\ \quad u(\operatorname{ccint}(0, 1)) = 1 \wedge \\ \quad a \in \mathbb{R} \wedge \\ \quad b \in \mathbb{R} \wedge \\ \quad a \leq b \Rightarrow \\ \quad \quad u(\operatorname{ccint}(a, b)) = (b - a) \end{array}}}
\vnbprop{translation-invariant-measure-lebesgue}{\vnbstmt{\begin{array}{@{}l@{}} \forall u, s. \\ \quad \operatorname{is\text{-}measure}(\mathbb{R}, \operatorname{borel\text{-}sets}(\operatorname{rr\text{-}ms}), u) \wedge \\ \quad \forall v, c. \\ \quad \quad ((v \in \operatorname{borel\text{-}sets}(\operatorname{rr\text{-}ms})) \wedge ((c \in \mathbb{R}) \wedge ((v + c) \in \operatorname{borel\text{-}sets}(\operatorname{rr\text{-}ms})))) \Rightarrow  \\ \quad \quad \quad u((v + c)) = u(v) \wedge \\ \quad u(\operatorname{ccint}(0, 1)) = 1 \wedge \\ \quad s \in \operatorname{borel\text{-}sets}(\operatorname{rr\text{-}ms}) \Rightarrow \\ \quad \quad u(s) \\ \quad \quad =\; \operatorname{res}(\operatorname{method\text{-}i\text{-}outer}(\mathbb{R}, \operatorname{closed\text{-}intervals}, (\lambda \mathit{ols} \in \operatorname{closed\text{-}intervals}.\; \operatorname{interval\text{-}length}(\mathit{ols}))), \operatorname{caratheodory\text{-}sets}(\mathbb{R}, \operatorname{method\text{-}i\text{-}outer}(\mathbb{R}, \operatorname{closed\text{-}intervals}, (\lambda \mathit{ols} \in \operatorname{closed\text{-}intervals}.\; \operatorname{interval\text{-}length}(\mathit{ols})))))(s) \end{array}}}
\vnbprop{lebesgue-translation-invariant}{\vnbstmt{\begin{array}{@{}l@{}} \forall u, k, s. \\ \quad \operatorname{is\text{-}measure}(\mathbb{R}, \operatorname{borel\text{-}sets}(\operatorname{rr\text{-}ms}), u) \wedge \\ \quad \forall v, c. \\ \quad \quad ((v \in \operatorname{borel\text{-}sets}(\operatorname{rr\text{-}ms})) \wedge ((c \in \mathbb{R}) \wedge ((v + c) \in \operatorname{borel\text{-}sets}(\operatorname{rr\text{-}ms})))) \Rightarrow  \\ \quad \quad \quad u((v + c)) = u(v) \wedge \\ \quad k \in \mathbb{R} \wedge \\ \quad 0 < k \wedge \\ \quad u(\operatorname{ccint}(0, 1)) = k \wedge \\ \quad s \in \operatorname{borel\text{-}sets}(\operatorname{rr\text{-}ms}) \Rightarrow \\ \quad \quad u(s) \\ \quad \quad =\; \operatorname{etimes}(k, \\ \quad \quad \quad \operatorname{res}(\operatorname{method\text{-}i\text{-}outer}(\mathbb{R}, \operatorname{closed\text{-}intervals}, (\lambda \mathit{ols} \in \operatorname{closed\text{-}intervals}.\; \operatorname{interval\text{-}length}(\mathit{ols}))), \operatorname{caratheodory\text{-}sets}(\mathbb{R}, \operatorname{method\text{-}i\text{-}outer}(\mathbb{R}, \operatorname{closed\text{-}intervals}, (\lambda \mathit{ols} \in \operatorname{closed\text{-}intervals}.\; \operatorname{interval\text{-}length}(\mathit{ols})))))(s)) \end{array}}}

\vnbbookfile{structure-library/lp-inequalities}{structure-library/lp-inequalities}
\noindent{\small lp-inequalities.scm -- POWERS IN [0,+inf] AND THE p-TH POWER INTEGRAL I\_p. Definitions only; nothing here is asserted and nothing is proved.  The laws and the inequalities of Hoelder and Minkowski are proved in theorem-library/lp-inequalities.scm (MEASURE-4-B, 2026-10-10).}\par
\par\medskip\noindent\textbf{Definition} (\texttt{epow}).\;${\operatorname{ept\_}}^{\operatorname{epp\_}} :=$ $\left(\text{if }(\operatorname{ept\_} = \operatorname{pos\text{-}inf})\text{ then }\operatorname{pos\text{-}inf}\text{ else }{\operatorname{ept\_}}^{\operatorname{epp\_}}\right)$
\par\medskip\noindent\textbf{Definition} (\texttt{ptwise-epow}).\;$\operatorname{ptwise\text{-}epow}(\operatorname{epo\_}, \operatorname{epf\_}, \operatorname{epp\_}) :=$ $(\lambda \operatorname{epx\_} \in \operatorname{epo\_}.\; {\operatorname{epf\_}(\operatorname{epx\_})}^{\operatorname{epp\_}})$
\par\medskip\noindent\textbf{Definition} (\texttt{ep-norm}).\;$I_{\operatorname{epp\_}}({\operatorname{epf\_}}) :=$
\begin{equation*}
{\operatorname{integral}(\omega, \mathit{ca}, \mu, \operatorname{ptwise\text{-}epow}(\omega, \operatorname{epf\_}, \operatorname{epp\_}))}^{\operatorname{recip}(\operatorname{epp\_})}
\end{equation*}


\vnbbookfile{structure-library/hausdorff-vocab}{structure-library/hausdorff-vocab}
\noindent{\small mh-vocab.scm -- vocabulary for construction-of-measures.pdf ch. 4, sections 2.4 and 4: method II and Hausdorff measures.  DEFINITIONS ONLY (destination: structure-library/). Agent RA-BASICS (M-H), 2026-10-26.  NEW DEFINITIONS, the user's to confirm.}\par
\par\medskip\noindent\textbf{Definition} (\texttt{hausdorff}).\;$\operatorname{hausdorff}(\operatorname{mhm\_}, \operatorname{mhl\_}) :=$ $\operatorname{method\text{-}ii}(\operatorname{mhm\_}, \operatorname{power}(\operatorname{pts}(\operatorname{mhm\_})), \operatorname{hausdorff\text{-}pm}(\operatorname{mhm\_}, \operatorname{mhl\_}))$
\par\medskip\noindent\textbf{Definition} (\texttt{diam}).\;$\operatorname{diam}(\operatorname{mhm\_}, \operatorname{mhs\_}) :=$
\begin{equation*}
\operatorname{esup}(\{\operatorname{mhw\_} \in \operatorname{rr\text{-}pos\text{-}star} : \exists \operatorname{mhx\_}, \operatorname{mhy\_} \in \operatorname{mhs\_}.\; (\operatorname{mhw\_} = \operatorname{dist}(\operatorname{mhm\_})(\operatorname{mhx\_}, \operatorname{mhy\_}))\})
\end{equation*}

\par\medskip\noindent\textbf{Definition} (\texttt{method-ii-delta}).\;$\operatorname{method\text{-}ii\text{-}delta}(\operatorname{mhm\_}, \operatorname{mhc\_}, \operatorname{mhp\_}, \operatorname{mhd\_}) :=$
\begin{equation*}
\begin{array}{@{}l@{}} \operatorname{method\text{-}i\text{-}outer}(\operatorname{pts}(\operatorname{mhm\_}), \\ \quad \{\operatorname{mha\_} \in \operatorname{mhc\_} : (\operatorname{diam}(\operatorname{mhm\_}, \operatorname{mha\_}) \leq \operatorname{mhd\_})\}, \\ \quad \operatorname{restrict}(\operatorname{mhp\_}, \{\operatorname{mha\_} \in \operatorname{mhc\_} : (\operatorname{diam}(\operatorname{mhm\_}, \operatorname{mha\_}) \leq \operatorname{mhd\_})\})) \end{array}
\end{equation*}

\par\medskip\noindent\textbf{Definition} (\texttt{method-ii}).\;$\operatorname{method\text{-}ii}(\operatorname{mhm\_}, \operatorname{mhc\_}, \operatorname{mhp\_}) :=$
\begin{equation*}
\begin{array}{@{}l@{}} (\lambda \operatorname{mhe\_} \in \operatorname{power}(\operatorname{pts}(\operatorname{mhm\_})).\; \\ \quad \operatorname{esup}(\{\operatorname{mhv\_} \in \operatorname{rr\text{-}pos\text{-}star} : \exists \operatorname{mhd\_} \in \mathbb{R}.\; (\operatorname{pos\text{-}rr}(\operatorname{mhd\_}) \wedge (\operatorname{mhv\_} = \operatorname{method\text{-}ii\text{-}delta}(\operatorname{mhm\_}, \operatorname{mhc\_}, \operatorname{mhp\_}, \operatorname{mhd\_})(\operatorname{mhe\_})))\})) \end{array}
\end{equation*}

\par\medskip\noindent\textbf{Definition} (\texttt{hausdorff-pm}).\;$\operatorname{hausdorff\text{-}pm}(\operatorname{mhm\_}, \operatorname{mhl\_}) :=$
\begin{equation*}
(\lambda \operatorname{mhh\_} \in \operatorname{power}(\operatorname{pts}(\operatorname{mhm\_})).\; {\operatorname{diam}(\operatorname{mhm\_}, \operatorname{mhh\_})}^{\operatorname{mhl\_}})
\end{equation*}

\par\medskip\noindent\textbf{Definition} (\texttt{phi-pm}).\;$\operatorname{phi\text{-}pm}(\operatorname{mhm\_}, \operatorname{mhf\_}) :=$
\begin{equation*}
(\lambda \operatorname{mhh\_} \in \operatorname{power}(\operatorname{pts}(\operatorname{mhm\_})).\; \left(\text{if }(\operatorname{diam}(\operatorname{mhm\_}, \operatorname{mhh\_}) = \operatorname{pos\text{-}inf})\text{ then }\operatorname{pos\text{-}inf}\text{ else }\operatorname{mhf\_}(\operatorname{diam}(\operatorname{mhm\_}, \operatorname{mhh\_}))\right))
\end{equation*}

\par\medskip\noindent\textbf{Definition} (\texttt{hausdorff-phi}).\;$\operatorname{hausdorff\text{-}phi}(\operatorname{mhm\_}, \operatorname{mhf\_}) :=$ $\operatorname{method\text{-}ii}(\operatorname{mhm\_}, \operatorname{power}(\operatorname{pts}(\operatorname{mhm\_})), \operatorname{phi\text{-}pm}(\operatorname{mhm\_}, \operatorname{mhf\_}))$
\par\medskip\noindent\textbf{Definition} (\texttt{nnreal}).\; $(\mathit{nnreal} \simeq \{\operatorname{mhr\_} \in \mathbb{R} : (0 \leq \operatorname{mhr\_})\})$
\par\medskip\noindent\textbf{Definition} (\texttt{is-doubling}).\;For every $\operatorname{mhf\_}$, 
\begin{equation*}
\begin{array}{@{}l@{}} \operatorname{is\text{-}doubling}(\operatorname{mhf\_}) \\ \Leftrightarrow\; \operatorname{mhf\_} \in (\mathit{nnreal} \to \mathit{nnreal}) \wedge \\ \operatorname{mhf\_}(0) = 0 \wedge \\ \forall \operatorname{mhr\_}, \operatorname{mhq\_} \in \mathit{nnreal}. \\ \quad (\operatorname{mhr\_} \leq \operatorname{mhq\_}) \Rightarrow (\operatorname{mhf\_}(\operatorname{mhr\_}) \leq \operatorname{mhf\_}(\operatorname{mhq\_})) \wedge \\ \exists \operatorname{mhp\_}. \\ \quad (\operatorname{pos\text{-}rr}(\operatorname{mhp\_}) \wedge \forall \operatorname{mhr\_} \in \mathbb{R}.\; ((0 < \operatorname{mhr\_}) \Rightarrow (\operatorname{mhf\_}(\cdot(2, \operatorname{mhr\_})) \leq \cdot(\operatorname{mhp\_}, \operatorname{mhf\_}(\operatorname{mhr\_}))))) \end{array}
\end{equation*}

\par\medskip\noindent\textbf{Definition} (\texttt{is-hausdorff-dim}).\;For all $\operatorname{mhm\_}$, $\operatorname{mhs\_}$ and $\operatorname{mhl\_}$, 
\begin{equation*}
\begin{array}{@{}l@{}} \operatorname{is\text{-}hausdorff\text{-}dim}(\operatorname{mhm\_}, \operatorname{mhs\_}, \operatorname{mhl\_}) \\ \Leftrightarrow\; \operatorname{mhl\_} \in \mathbb{R} \wedge \\ 0 \leq \operatorname{mhl\_} \wedge \\ \forall \operatorname{mhr\_} \in \mathbb{R}.\; ((\operatorname{mhl\_} < \operatorname{mhr\_}) \Rightarrow (\operatorname{hausdorff}(\operatorname{mhm\_}, \operatorname{mhr\_})(\operatorname{mhs\_}) = 0)) \wedge \\ \forall \operatorname{mhr\_} \in \mathbb{R}. \\ \quad (0 < \operatorname{mhr\_}) \Rightarrow ((\operatorname{mhr\_} < \operatorname{mhl\_}) \Rightarrow (\operatorname{hausdorff}(\operatorname{mhm\_}, \operatorname{mhr\_})(\operatorname{mhs\_}) = \operatorname{pos\text{-}inf})) \end{array}
\end{equation*}


\vnbbookfile{structure-library/doubling-functions}{structure-library/doubling-functions}
\noindent{\small doubling-functions.scm -- WAVE 2 (2026-10-27..30), wired 2026-10-30 from scratchpad/mf/mf-mivocab.scm. Library: measure.}\par
\par\medskip\noindent\textbf{Definition} (\texttt{is-boundedly-equivalent}).\;For all $\varphi$ and $\psi$, 
\begin{equation*}
\begin{array}{@{}l@{}} \operatorname{is\text{-}boundedly\text{-}equivalent}(\varphi, \psi) \\ \Leftrightarrow\; \varphi \in (\mathit{nnreal} \to \mathit{nnreal}) \wedge \\ \psi \in (\mathit{nnreal} \to \mathit{nnreal}) \wedge \\ \exists \operatorname{bel\_}, \operatorname{beu\_}. \\ \quad \operatorname{pos\text{-}rr}(\operatorname{bel\_}) \wedge \\ \quad \operatorname{pos\text{-}rr}(\operatorname{beu\_}) \wedge \\ \quad \forall \operatorname{ber\_} \in \mathbb{R}. \\ \quad \quad (0 < \operatorname{ber\_}) \Rightarrow ((\cdot(\operatorname{bel\_}, \varphi(\operatorname{ber\_})) \leq \psi(\operatorname{ber\_})) \wedge (\psi(\operatorname{ber\_}) \leq \cdot(\operatorname{beu\_}, \varphi(\operatorname{ber\_})))) \end{array}
\end{equation*}


\vnbbookfile{theorem-library/doubling-functions}{theorem-library/doubling-functions}
\noindent{\small doubling-functions.scm -- WAVE 2 (2026-10-27..30), wired 2026-10-30 from scratchpad/mf/mf-mi1.scm. Library: measure.}\par
\vnbprop{boundedly-equivalent-symmetric}{\vnbstmt{\begin{array}{@{}l@{}} \forall i, a. \\ \quad \operatorname{is\text{-}boundedly\text{-}equivalent}(i, a) \Rightarrow \operatorname{is\text{-}boundedly\text{-}equivalent}(a, i) \end{array}}}
\vnbprop{doubling-of-boundedly-equivalent}{\vnbstmt{\begin{array}{@{}l@{}} \forall i, a. \\ \quad \operatorname{is\text{-}boundedly\text{-}equivalent}(i, a) \wedge \\ \quad a(0) = 0 \wedge \\ \quad \forall r, q \in \mathit{nnreal}.\; ((r \leq q) \Rightarrow (a(r) \leq a(q))) \wedge \\ \quad \operatorname{is\text{-}doubling}(i) \Rightarrow \\ \quad \quad \operatorname{is\text{-}doubling}(a) \end{array}}}
\vnbprop{doubling-iff-boundedly-equivalent}{\vnbstmt{\begin{array}{@{}l@{}} \forall i, a. \\ \quad \operatorname{is\text{-}boundedly\text{-}equivalent}(i, a) \wedge \\ \quad i(0) = 0 \wedge \\ \quad a(0) = 0 \wedge \\ \quad \forall r, q \in \mathit{nnreal}.\; ((r \leq q) \Rightarrow (i(r) \leq i(q))) \wedge \\ \quad \forall r, q \in \mathit{nnreal}.\; ((r \leq q) \Rightarrow (a(r) \leq a(q))) \Rightarrow \\ \quad \quad \operatorname{is\text{-}doubling}(i) \Leftrightarrow \operatorname{is\text{-}doubling}(a) \end{array}}}

\vnbbookfile{theorem-library/method-ii}{theorem-library/method-ii}
\noindent{\small mh-method2.scm -- construction-of-measures.pdf section 2.4: the method II construction. Prop 4.23: METHOD-II(m, C, p) is a metric outer measure.  Helper prefix mh2-.}\par
\vnbprop{diam-in}{\vnbstmt{\forall m, a.\; (\operatorname{diam}(m, a) \in \operatorname{rr\text{-}pos\text{-}star})}}
\vnbprop{diam-empty}{\vnbstmt{\forall m.\; (\operatorname{diam}(m, \emptyset) = 0)}}
\vnbprop{diam-dist-le}{\vnbstmt{\begin{array}{@{}l@{}} \forall m, a.; \forall x, y \in a. \\ \quad (\operatorname{is\text{-}metric\text{-}space}(m) \wedge (a \subseteq \operatorname{pts}(m))) \Rightarrow (\operatorname{dist}(m)(x, y) \leq \operatorname{diam}(m, a)) \end{array}}}
\vnbprop{method-ii-delta-premeasure}{\vnbstmt{\begin{array}{@{}l@{}} \forall m, s, a, l. \\ \quad (\operatorname{is\text{-}premeasure}(\operatorname{pts}(m), s, a) \wedge ((l \in \mathbb{R}) \wedge (0 \leq l))) \Rightarrow  \\ \quad \quad \operatorname{is\text{-}premeasure}(\operatorname{pts}(m), \{\mathit{mha} \in s : (\operatorname{diam}(m, \mathit{mha}) \leq l)\}, \operatorname{restrict}(a, \{\mathit{mha} \in s : (\operatorname{diam}(m, \mathit{mha}) \leq l)\})) \end{array}}}
\vnbprop{method-ii-delta-outer}{\vnbstmt{\begin{array}{@{}l@{}} \forall m, s, a, l. \\ \quad (\operatorname{is\text{-}premeasure}(\operatorname{pts}(m), s, a) \wedge ((l \in \mathbb{R}) \wedge (0 \leq l))) \Rightarrow  \\ \quad \quad \operatorname{is\text{-}outer\text{-}measure}(\operatorname{pts}(m), \operatorname{method\text{-}ii\text{-}delta}(m, s, a, l)) \end{array}}}
\vnbprop{method-ii-delta-value}{\vnbstmt{\begin{array}{@{}l@{}} \forall m, s, a, l.; \forall b \in \operatorname{power}(\operatorname{pts}(m)). \\ \quad \operatorname{method\text{-}ii\text{-}delta}(m, s, a, l)(b) \\ \quad =\; \operatorname{einf}(\operatorname{cover\text{-}sums}(\operatorname{pts}(m), \{\mathit{mha} \in s : (\operatorname{diam}(m, \mathit{mha}) \leq l)\}, \operatorname{restrict}(a, \{\mathit{mha} \in s : (\operatorname{diam}(m, \mathit{mha}) \leq l)\}), b)) \end{array}}}
\vnbprop{method-ii-delta-antitone}{\vnbstmt{\begin{array}{@{}l@{}} \forall m, s, a, b, c, d. \\ \quad (\operatorname{is\text{-}premeasure}(\operatorname{pts}(m), s, a) \wedge ((b \in \mathbb{R}) \wedge ((0 \leq b) \wedge ((c \in \mathbb{R}) \wedge ((b \leq c) \wedge (d \in \operatorname{power}(\operatorname{pts}(m)))))))) \Rightarrow  \\ \quad \quad \operatorname{method\text{-}ii\text{-}delta}(m, s, a, c)(d) \leq \operatorname{method\text{-}ii\text{-}delta}(m, s, a, b)(d) \end{array}}}
\vnbprop{method-ii-value}{\vnbstmt{\begin{array}{@{}l@{}} \forall m, s, a.; \forall b \in \operatorname{power}(\operatorname{pts}(m)). \\ \quad \operatorname{method\text{-}ii}(m, s, a)(b) \\ \quad =\; \operatorname{esup}(\{v \in \operatorname{rr\text{-}pos\text{-}star} : \exists d \in \mathbb{R}.\; (\operatorname{pos\text{-}rr}(d) \wedge (v = \operatorname{method\text{-}ii\text{-}delta}(m, s, a, d)(b)))\}) \end{array}}}
\vnbprop{method-ii-delta-le}{\vnbstmt{\begin{array}{@{}l@{}} \forall m, s, a, b, l. \\ \quad (\operatorname{is\text{-}premeasure}(\operatorname{pts}(m), s, a) \wedge ((b \in \operatorname{power}(\operatorname{pts}(m))) \wedge \operatorname{pos\text{-}rr}(l))) \Rightarrow  \\ \quad \quad \operatorname{method\text{-}ii\text{-}delta}(m, s, a, l)(b) \leq \operatorname{method\text{-}ii}(m, s, a)(b) \end{array}}}
\vnbprop{method-ii-least}{\vnbstmt{\begin{array}{@{}l@{}} \forall m, s, a, c, b. \\ \quad c \in \operatorname{power}(\operatorname{pts}(m)) \wedge \\ \quad b \in \operatorname{rr\text{-}pos\text{-}star} \wedge \\ \quad \forall l.\; (\operatorname{pos\text{-}rr}(l) \Rightarrow (\operatorname{method\text{-}ii\text{-}delta}(m, s, a, l)(c) \leq b)) \Rightarrow \\ \quad \quad \operatorname{method\text{-}ii}(m, s, a)(c) \leq b \end{array}}}
\vnbprop{method-ii-in}{\vnbstmt{\forall m, s, a.; \forall b \in \operatorname{power}(\operatorname{pts}(m)).\; (\operatorname{method\text{-}ii}(m, s, a)(b) \in \operatorname{rr\text{-}pos\text{-}star})}}
\vnbprop{method-ii-outer-measure}{\vnbstmt{\begin{array}{@{}l@{}} \forall m, s, a. \\ \quad \operatorname{is\text{-}premeasure}(\operatorname{pts}(m), s, a) \Rightarrow \operatorname{is\text{-}outer\text{-}measure}(\operatorname{pts}(m), \operatorname{method\text{-}ii}(m, s, a)) \end{array}}}
\vnbprop{union-in-power}{\vnbstmt{\begin{array}{@{}l@{}} \forall x, a, b. \\ \quad ((a \in \operatorname{power}(x)) \wedge (b \in \operatorname{power}(x))) \Rightarrow ((a \cup b) \in \operatorname{power}(x)) \end{array}}}
\vnbprop{method-ii-delta-separated}{\vnbstmt{\begin{array}{@{}l@{}} \forall m, s, a, l, r, b, c. \\ \quad \operatorname{is\text{-}metric\text{-}space}(m) \wedge \\ \quad \operatorname{is\text{-}premeasure}(\operatorname{pts}(m), s, a) \wedge \\ \quad l \in \mathbb{R} \wedge \\ \quad 0 \leq l \wedge \\ \quad \operatorname{pos\text{-}rr}(r) \wedge \\ \quad l < r \wedge \\ \quad b \in \operatorname{power}(\operatorname{pts}(m)) \wedge \\ \quad c \in \operatorname{power}(\operatorname{pts}(m)) \wedge \\ \quad \forall x \in b, y \in c.\; (r \leq \operatorname{dist}(m)(x, y)) \Rightarrow \\ \quad \quad \operatorname{eplus}(\operatorname{method\text{-}ii\text{-}delta}(m, s, a, l)(b), \operatorname{method\text{-}ii\text{-}delta}(m, s, a, l)(c)) \\ \quad \quad \leq\; \operatorname{method\text{-}ii\text{-}delta}(m, s, a, l)((b \cup c)) \end{array}}}
\vnbprop{method-ii-seq}{\vnbstmt{\begin{array}{@{}l@{}} \forall m, s, a, b. \\ \quad (\operatorname{is\text{-}premeasure}(\operatorname{pts}(m), s, a) \wedge (b \in \operatorname{power}(\operatorname{pts}(m)))) \Rightarrow  \\ \quad \quad \operatorname{method\text{-}ii}(m, s, a)(b) \\ \quad \quad =\; \operatorname{esup}(\operatorname{image}((\lambda n \in \mathbb{N}.\; \operatorname{method\text{-}ii\text{-}delta}(m, s, a, \operatorname{recip}((n + 1)))(b)), \mathbb{N})) \end{array}}}
\vnbprop{method-ii-metric-outer-measure}{\vnbstmt{\begin{array}{@{}l@{}} \forall m, s, a. \\ \quad (\operatorname{is\text{-}metric\text{-}space}(m) \wedge \operatorname{is\text{-}premeasure}(\operatorname{pts}(m), s, a)) \Rightarrow  \\ \quad \quad \operatorname{is\text{-}metric\text{-}outer\text{-}measure}(m, \operatorname{method\text{-}ii}(m, s, a)) \end{array}}}

\vnbbookfile{structure-library/mb4-vocab}{structure-library/mb4-vocab}
\noindent{\small mb4-vocab.scm -- vocabulary for construction of measures, chapter 5 (M-B4, 2026-10-26). Destination: structure-library/.  Library: measure (DIAM and IS-POLISH use only metric-space and extended-real vocabulary; they could sit in analysis).}\par
\par\medskip\noindent\textbf{Definition} (\texttt{is-polish}).\;For every $s$,  $\operatorname{is\text{-}polish}(s) \Leftrightarrow (\operatorname{is\text{-}complete}(s) \wedge \operatorname{is\text{-}separable}(s))$
\par\medskip\noindent\textbf{Definition} (\texttt{graph-ms}).\;$\operatorname{graph\text{-}ms}(s, a, g) :=$
\begin{equation*}
{[}a, (\lambda {[}\operatorname{gmu\_}, \operatorname{gmv\_}] \in (a \times a).\; (\operatorname{dist}(s)(\operatorname{gmu\_}, \operatorname{gmv\_}) + \lvert (g(\operatorname{gmu\_}) - g(\operatorname{gmv\_})) \rvert))]
\end{equation*}


\vnbbookfile{theorem-library/mb4-sieve}{theorem-library/mb4-sieve}
\noindent{\small mb4-sieve.scm -- construction of measures, Proposition 5.6 (M-B4, 2026-10-26): in a complete metric space a decreasing sequence of closed nonempty sets whose diameters tend to 0 meets in exactly one point (the notes' proof of 5.6: for each infinite sequence z the sets S\_\{z|n\} of a sieve are such a sequence, so S\_z is a ...}\par
\vnbprop{diam-unfold}{\vnbstmt{\begin{array}{@{}l@{}} \forall s, a. \\ \quad (\operatorname{diam}(s, a) \simeq \operatorname{esup}(\{r \in \operatorname{rr\text{-}pos\text{-}star} : \exists x, y \in a.\; (r = \operatorname{dist}(s)(x, y))\})) \end{array}}}
\vnbprop{nested-closed-intersection}{\vnbstmt{\begin{array}{@{}l@{}} \forall s.; \forall a \in (\mathbb{N} \to \operatorname{power}(\operatorname{pts}(s))). \\ \quad \operatorname{is\text{-}complete}(s) \wedge \\ \quad \forall n \in \mathbb{N}.\; \exists x.\; (x \in a(n)) \wedge \\ \quad \forall n \in \mathbb{N}.\; \operatorname{is\text{-}closed}(s, a(n)) \wedge \\ \quad \forall n \in \mathbb{N}.\; (a((n + 1)) \subseteq a(n)) \wedge \\ \quad \forall b. \\ \quad \quad \operatorname{pos\text{-}rr}(b) \Rightarrow  \\ \quad \quad \quad \exists \mathit{mcbn} \in \mathbb{N}.\; \forall n \in \mathbb{N}.\; ((\mathit{mcbn} \leq n) \Rightarrow (\operatorname{diam}(s, a(n)) \leq b)) \Rightarrow \\ \quad \quad \exists p \in \operatorname{pts}(s). \\ \quad \quad \quad (\forall n \in \mathbb{N}.\; (p \in a(n)) \wedge \forall q \in \operatorname{pts}(s).\; (\forall n \in \mathbb{N}.\; (q \in a(n)) \Rightarrow (q = p))) \end{array}}}

\vnbbookfile{theorem-library/signed-measure-laws}{theorem-library/signed-measure-laws}
\noindent{\small signed-measure-laws.scm -- WAVE 2 (2026-10-27..30), wired 2026-10-30 from scratchpad/mf/mf-fa.scm, scratchpad/mf/mf-pv.scm, scratchpad/mf/mf-pv2.scm, scratchpad/mf/mf-neg.scm, scratchpad/mf/mf-fin.scm, scratchpad/mf/mf-fin2.scm, scratchpad/mf/mf-fin3.scm. Library: measure.}\par
\vnbprop{ext-add-in}{\vnbstmt{\forall x, y.\; (((x \in \operatorname{rr\text{-}star}) \wedge (y \in \operatorname{rr\text{-}star})) \Rightarrow (\operatorname{ext\text{-}add}(x, y) \in \operatorname{rr\text{-}star}))}}
\vnbprop{ext-psum-stable}{\vnbstmt{\begin{array}{@{}l@{}} \forall d \in \mathbb{N}, x.; \forall m \in \mathbb{N}. \\ \quad (\forall k \in \mathbb{N}.\; (x(k) \in \operatorname{rr\text{-}star}) \wedge \forall k \in \mathbb{N}.\; ((m \leq k) \Rightarrow (x(k) = 0))) \Rightarrow  \\ \quad \quad \operatorname{ext\text{-}partial\text{-}sum}(x, (m + d)) = \operatorname{ext\text{-}partial\text{-}sum}(x, m) \end{array}}}
\vnbprop{ext-psum-eventually}{\vnbstmt{\begin{array}{@{}l@{}} \forall x.; \forall m \in \mathbb{N}. \\ \quad (\forall k \in \mathbb{N}.\; (x(k) \in \operatorname{rr\text{-}star}) \wedge \forall k \in \mathbb{N}.\; ((m \leq k) \Rightarrow (x(k) = 0))) \Rightarrow  \\ \quad \quad \exists a \in \mathbb{N}. \forall n \in \mathbb{N}. \\ \quad \quad \quad (a \leq n) \Rightarrow (\operatorname{ext\text{-}partial\text{-}sum}(x, n) = \operatorname{ext\text{-}partial\text{-}sum}(x, m)) \end{array}}}
\vnbprop{ext-series-sum-finite}{\vnbstmt{\begin{array}{@{}l@{}} \forall x, l.; \forall m \in \mathbb{N}. \\ \quad \operatorname{ext\text{-}series\text{-}converges\text{-}to}(x, l) \wedge \\ \quad \forall k \in \mathbb{N}.\; (x(k) \in \operatorname{rr\text{-}star}) \wedge \\ \quad \forall k \in \mathbb{N}.\; ((m \leq k) \Rightarrow (x(k) = 0)) \Rightarrow \\ \quad \quad l = \operatorname{ext\text{-}partial\text{-}sum}(x, m) \end{array}}}
\vnbprop{pad2-family-exists}{\vnbstmt{\begin{array}{@{}l@{}} \forall a, \alpha, b. \\ \quad ((\emptyset \in a) \wedge ((\alpha \in a) \wedge ((b \in a) \wedge ((\alpha \cap b) = \emptyset)))) \Rightarrow  \\ \quad \quad \exists g \in (\mathbb{N} \to a). \\ \quad \quad \quad \forall i, j. \\ \quad \quad \quad \quad ((i \in \mathbb{N}) \wedge ((j \in \mathbb{N}) \wedge \neg (i = j))) \Rightarrow ((g(i) \cap g(j)) = \emptyset) \wedge \\ \quad \quad \quad g(0) = \alpha \wedge \\ \quad \quad \quad g(1) = b \wedge \\ \quad \quad \quad \forall k \in \mathbb{N}.\; ((2 \leq k) \Rightarrow (g(k) = \emptyset)) \wedge \\ \quad \quad \quad \operatorname{big\text{-}union}(n, \mathbb{N}, g(n)) = (\alpha \cup b) \end{array}}}
\vnbprop{signed-sum-two}{\vnbstmt{\begin{array}{@{}l@{}} \forall a, b, u, g. \\ \quad \operatorname{is\text{-}signed\text{-}measure}(a, b, u) \wedge \\ \quad g \in (\mathbb{N} \to b) \wedge \\ \quad \forall i, j. \\ \quad \quad ((i \in \mathbb{N}) \wedge ((j \in \mathbb{N}) \wedge \neg (i = j))) \Rightarrow ((g(i) \cap g(j)) = \emptyset) \wedge \\ \quad \forall k \in \mathbb{N}.\; ((2 \leq k) \Rightarrow (g(k) = \emptyset)) \Rightarrow \\ \quad \quad u(\operatorname{big\text{-}union}(n, \mathbb{N}, g(n))) = \operatorname{ext\text{-}add}(u(g(0)), u(g(1))) \end{array}}}
\vnbprop{signed-measure-fa}{\vnbstmt{\begin{array}{@{}l@{}} \forall a, b, u, \alpha, \beta. \\ \quad (\operatorname{is\text{-}signed\text{-}measure}(a, b, u) \wedge ((\alpha \in b) \wedge ((\beta \in b) \wedge ((\alpha \cap \beta) = \emptyset)))) \Rightarrow  \\ \quad \quad u((\alpha \cup \beta)) = \operatorname{ext\text{-}add}(u(\alpha), u(\beta)) \end{array}}}
\vnbprop{pos-variation-value}{\vnbstmt{\begin{array}{@{}l@{}} \forall b, u.; \forall a \in b. \\ \quad (\operatorname{pos\text{-}variation}(b, u)(a) \simeq \operatorname{esup}(\{v \in \operatorname{rr\text{-}pos\text{-}star} : \exists e \in b.\; ((e \subseteq a) \wedge (u(e) = v))\})) \end{array}}}
\vnbprop{pos-variation-set-sub}{\vnbstmt{\begin{array}{@{}l@{}} \forall b, u, a. \\ \quad \{v \in \operatorname{rr\text{-}pos\text{-}star} : \exists e \in b.\; ((e \subseteq a) \wedge (u(e) = v))\} \subseteq \operatorname{rr\text{-}pos\text{-}star} \end{array}}}
\vnbprop{pos-variation-in}{\vnbstmt{\forall b, u.; \forall a \in b.\; (\operatorname{pos\text{-}variation}(b, u)(a) \in \operatorname{rr\text{-}pos\text{-}star})}}
\vnbprop{pos-variation-fun}{\vnbstmt{\begin{array}{@{}l@{}} \forall a, b, u. \\ \quad \operatorname{is\text{-}sigma\text{-}algebra}(a, b) \Rightarrow (\operatorname{pos\text{-}variation}(b, u) \in (b \to \operatorname{rr\text{-}pos\text{-}star})) \end{array}}}
\vnbprop{pos-variation-upper}{\vnbstmt{\begin{array}{@{}l@{}} \forall b, u, a, e. \\ \quad ((u \in (b \to \operatorname{rr\text{-}star})) \wedge ((a \in b) \wedge ((e \in b) \wedge (e \subseteq a)))) \Rightarrow  \\ \quad \quad u(e) \leq \operatorname{pos\text{-}variation}(b, u)(a) \end{array}}}
\vnbprop{pos-variation-least}{\vnbstmt{\begin{array}{@{}l@{}} \forall c, u, a, b. \\ \quad ((a \in c) \wedge ((b \in \operatorname{rr\text{-}pos\text{-}star}) \wedge \forall e \in c.\; ((e \subseteq a) \Rightarrow (u(e) \leq b)))) \Rightarrow  \\ \quad \quad \operatorname{pos\text{-}variation}(c, u)(a) \leq b \end{array}}}
\vnbprop{pos-variation-empty}{\vnbstmt{\begin{array}{@{}l@{}} \forall a, b, u. \\ \quad \operatorname{is\text{-}signed\text{-}measure}(a, b, u) \Rightarrow (\operatorname{pos\text{-}variation}(b, u)(\emptyset) = 0) \end{array}}}
\vnbprop{pos-variation-mono}{\vnbstmt{\begin{array}{@{}l@{}} \forall c, u, a, b. \\ \quad ((u \in (c \to \operatorname{rr\text{-}star})) \wedge ((a \in c) \wedge ((b \in c) \wedge (a \subseteq b)))) \Rightarrow  \\ \quad \quad \operatorname{pos\text{-}variation}(c, u)(a) \leq \operatorname{pos\text{-}variation}(c, u)(b) \end{array}}}
\vnbprop{pos-variation-ge}{\vnbstmt{\begin{array}{@{}l@{}} \forall b, u, a. \\ \quad ((u \in (b \to \operatorname{rr\text{-}star})) \wedge (a \in b)) \Rightarrow (u(a) \leq \operatorname{pos\text{-}variation}(b, u)(a)) \end{array}}}
\vnbprop{disjoint-of-subsets}{\vnbstmt{\begin{array}{@{}l@{}} \forall a, b, e, f. \\ \quad (((a \cap b) = \emptyset) \wedge ((e \subseteq a) \wedge (f \subseteq b))) \Rightarrow  \\ \quad \quad (e \cap f) = \emptyset \end{array}}}
\vnbprop{union-subset-union}{\vnbstmt{\begin{array}{@{}l@{}} \forall a, b, e, f. \\ \quad ((e \subseteq a) \wedge (f \subseteq b)) \Rightarrow ((e \cup f) \subseteq (a \cup b)) \end{array}}}
\vnbprop{pos-variation-superadditive}{\vnbstmt{\begin{array}{@{}l@{}} \forall c, d, u, a, b. \\ \quad (\operatorname{is\text{-}signed\text{-}measure}(c, d, u) \wedge ((a \in d) \wedge ((b \in d) \wedge ((a \cap b) = \emptyset)))) \Rightarrow  \\ \quad \quad \operatorname{eplus}(\operatorname{pos\text{-}variation}(d, u)(a), \operatorname{pos\text{-}variation}(d, u)(b)) \leq \operatorname{pos\text{-}variation}(d, u)((a \cup b)) \end{array}}}
\vnbprop{esum-nn-type}{\vnbstmt{\forall f \in (\mathbb{N} \to \operatorname{rr\text{-}pos\text{-}star}).\; (\operatorname{esum}(f) \in \operatorname{rr\text{-}pos\text{-}star})}}
\vnbprop{inter-sub-right}{\vnbstmt{\forall a, b.\; ((a \cap b) \subseteq b)}}
\vnbprop{inter-sub-left}{\vnbstmt{\forall a, b.\; ((a \cap b) \subseteq a)}}
\vnbprop{ext-add-le-eplus}{\vnbstmt{\begin{array}{@{}l@{}} \forall a, b, c, d. \\ \quad ((a \in \operatorname{rr\text{-}star}) \wedge ((b \in \operatorname{rr\text{-}star}) \wedge ((c \in \operatorname{rr\text{-}pos\text{-}star}) \wedge ((d \in \operatorname{rr\text{-}pos\text{-}star}) \wedge ((a \leq c) \wedge (b \leq d)))))) \Rightarrow  \\ \quad \quad \operatorname{ext\text{-}add}(a, b) \leq \operatorname{eplus}(c, d) \end{array}}}
\vnbprop{ext-psum-le-prefix}{\vnbstmt{\begin{array}{@{}l@{}} \forall n \in \mathbb{N}, x, y. \\ \quad (\forall k \in \mathbb{N}.\; (x(k) \in \operatorname{rr\text{-}star}) \wedge ((y \in (\mathbb{N} \to \operatorname{rr\text{-}pos\text{-}star})) \wedge \forall k \in \mathbb{N}.\; (x(k) \leq y(k)))) \Rightarrow  \\ \quad \quad \operatorname{ext\text{-}partial\text{-}sum}(x, n) \leq \operatorname{finsum}(\operatorname{rr\text{-}pos\text{-}star\text{-}add\text{-}monoid}, y, \mathrm{OS}(n)) \end{array}}}
\vnbprop{intersect-family-exists}{\vnbstmt{\begin{array}{@{}l@{}} \forall a, b, e, f. \\ \quad \operatorname{is\text{-}sigma\text{-}algebra}(a, b) \wedge \\ \quad e \in b \wedge \\ \quad f \in (\mathbb{N} \to b) \wedge \\ \quad \forall i, j. \\ \quad \quad ((i \in \mathbb{N}) \wedge ((j \in \mathbb{N}) \wedge \neg (i = j))) \Rightarrow ((f(i) \cap f(j)) = \emptyset) \wedge \\ \quad e \subseteq \operatorname{big\text{-}union}(n, \mathbb{N}, f(n)) \Rightarrow \\ \quad \quad \exists h \in (\mathbb{N} \to b). \\ \quad \quad \quad \forall i, j. \\ \quad \quad \quad \quad ((i \in \mathbb{N}) \wedge ((j \in \mathbb{N}) \wedge \neg (i = j))) \Rightarrow ((h(i) \cap h(j)) = \emptyset) \wedge \\ \quad \quad \quad \forall k \in \mathbb{N}.\; (h(k) \subseteq f(k)) \wedge \\ \quad \quad \quad \operatorname{big\text{-}union}(n, \mathbb{N}, h(n)) = e \end{array}}}
\vnbprop{pos-variation-le-esum}{\vnbstmt{\begin{array}{@{}l@{}} \forall a, b, u, f. \\ \quad \operatorname{is\text{-}signed\text{-}measure}(a, b, u) \wedge \\ \quad f \in (\mathbb{N} \to b) \wedge \\ \quad \forall i, j. \\ \quad \quad ((i \in \mathbb{N}) \wedge ((j \in \mathbb{N}) \wedge \neg (i = j))) \Rightarrow ((f(i) \cap f(j)) = \emptyset) \Rightarrow \\ \quad \quad \operatorname{pos\text{-}variation}(b, u)(\operatorname{big\text{-}union}(n, \mathbb{N}, f(n))) \\ \quad \quad \leq\; \operatorname{esum}((\lambda n \in \mathbb{N}.\; \operatorname{pos\text{-}variation}(b, u)(f(n)))) \end{array}}}
\vnbprop{seg-union-disjoint}{\vnbstmt{\begin{array}{@{}l@{}} \forall a, b, n, f. \\ \quad \operatorname{is\text{-}sigma\text{-}algebra}(a, b) \wedge \\ \quad f \in (\mathbb{N} \to b) \wedge \\ \quad n \in \mathbb{N} \wedge \\ \quad \forall i, j. \\ \quad \quad ((i \in \mathbb{N}) \wedge ((j \in \mathbb{N}) \wedge \neg (i = j))) \Rightarrow ((f(i) \cap f(j)) = \emptyset) \Rightarrow \\ \quad \quad (\operatorname{big\text{-}union}(\mathit{uj}, \mathrm{OS}(n), f(\mathit{uj})) \cap f(n)) = \emptyset \end{array}}}
\vnbprop{seg-union-sub}{\vnbstmt{\begin{array}{@{}l@{}} \forall n \in \mathbb{N}, f. \\ \quad \operatorname{big\text{-}union}(j, \mathrm{OS}(n), f(j)) \subseteq \operatorname{big\text{-}union}(a, \mathbb{N}, f(a)) \end{array}}}
\vnbprop{pos-variation-prefix-le}{\vnbstmt{\begin{array}{@{}l@{}} \forall n \in \mathbb{N}, a, b, u, f. \\ \quad \operatorname{is\text{-}signed\text{-}measure}(a, b, u) \wedge \\ \quad f \in (\mathbb{N} \to b) \wedge \\ \quad \forall i, j. \\ \quad \quad ((i \in \mathbb{N}) \wedge ((j \in \mathbb{N}) \wedge \neg (i = j))) \Rightarrow ((f(i) \cap f(j)) = \emptyset) \Rightarrow \\ \quad \quad \operatorname{finsum}(\operatorname{rr\text{-}pos\text{-}star\text{-}add\text{-}monoid}, (\lambda c \in \mathbb{N}.\; \operatorname{pos\text{-}variation}(b, u)(f(c))), \mathrm{OS}(n)) \\ \quad \quad \leq\; \operatorname{pos\text{-}variation}(b, u)(\operatorname{big\text{-}union}(\mathit{uj}, \mathrm{OS}(n), f(\mathit{uj}))) \end{array}}}
\vnbprop{pos-variation-countably-additive}{\vnbstmt{\begin{array}{@{}l@{}} \forall a, b, u, f. \\ \quad \operatorname{is\text{-}signed\text{-}measure}(a, b, u) \wedge \\ \quad f \in (\mathbb{N} \to b) \wedge \\ \quad \forall i, j. \\ \quad \quad ((i \in \mathbb{N}) \wedge ((j \in \mathbb{N}) \wedge \neg (i = j))) \Rightarrow ((f(i) \cap f(j)) = \emptyset) \Rightarrow \\ \quad \quad \operatorname{pos\text{-}variation}(b, u)(\operatorname{big\text{-}union}(n, \mathbb{N}, f(n))) = \operatorname{esum}((\lambda n \in \mathbb{N}.\; \operatorname{pos\text{-}variation}(b, u)(f(n)))) \end{array}}}
\vnbprop{pos-variation-is-measure}{\vnbstmt{\begin{array}{@{}l@{}} \forall a, b, u. \\ \quad \operatorname{is\text{-}signed\text{-}measure}(a, b, u) \Rightarrow \operatorname{is\text{-}measure}(a, b, \operatorname{pos\text{-}variation}(b, u)) \end{array}}}
\vnbprop{ext-neg-in}{\vnbstmt{\forall x \in \operatorname{rr\text{-}star}.\; (\operatorname{ext\text{-}neg}(x) \in \operatorname{rr\text{-}star})}}
\vnbprop{ext-neg-zero}{\vnbstmt{\operatorname{ext\text{-}neg}(0) = 0}}
\vnbprop{ext-neg-to-pos-inf}{\vnbstmt{\forall x \in \operatorname{rr\text{-}star}.\; ((\operatorname{ext\text{-}neg}(x) = \operatorname{pos\text{-}inf}) \Rightarrow (x = \operatorname{neg\text{-}inf}))}}
\vnbprop{ext-neg-to-neg-inf}{\vnbstmt{\forall x \in \operatorname{rr\text{-}star}.\; ((\operatorname{ext\text{-}neg}(x) = \operatorname{neg\text{-}inf}) \Rightarrow (x = \operatorname{pos\text{-}inf}))}}
\vnbprop{ext-neg-real-back}{\vnbstmt{\forall x \in \operatorname{rr\text{-}star}.\; ((\operatorname{ext\text{-}neg}(x) \in \mathbb{R}) \Rightarrow (x \in \mathbb{R}))}}
\vnbprop{ext-neg-add-up}{\vnbstmt{\begin{array}{@{}l@{}} \forall a, b. \\ \quad ((a \in \operatorname{rr\text{-}star}) \wedge ((b \in \operatorname{rr\text{-}star}) \wedge (\neg (a = \operatorname{pos\text{-}inf}) \wedge \neg (b = \operatorname{pos\text{-}inf})))) \Rightarrow  \\ \quad \quad \operatorname{ext\text{-}neg}(\operatorname{ext\text{-}add}(a, b)) = \operatorname{ext\text{-}add}(\operatorname{ext\text{-}neg}(a), \operatorname{ext\text{-}neg}(b)) \end{array}}}
\vnbprop{ext-neg-add-down}{\vnbstmt{\begin{array}{@{}l@{}} \forall a, b. \\ \quad ((a \in \operatorname{rr\text{-}star}) \wedge ((b \in \operatorname{rr\text{-}star}) \wedge (\neg (a = \operatorname{neg\text{-}inf}) \wedge \neg (b = \operatorname{neg\text{-}inf})))) \Rightarrow  \\ \quad \quad \operatorname{ext\text{-}neg}(\operatorname{ext\text{-}add}(a, b)) = \operatorname{ext\text{-}add}(\operatorname{ext\text{-}neg}(a), \operatorname{ext\text{-}neg}(b)) \end{array}}}
\vnbprop{ext-neg-lt-flip}{\vnbstmt{\forall a, b.\; (((a \in \operatorname{rr\text{-}star}) \wedge ((b \in \mathbb{R}) \wedge (a < -b))) \Rightarrow (b < \operatorname{ext\text{-}neg}(a)))}}
\vnbprop{ext-neg-gt-flip}{\vnbstmt{\forall a, b.\; (((a \in \operatorname{rr\text{-}star}) \wedge ((b \in \mathbb{R}) \wedge (-b < a))) \Rightarrow (\operatorname{ext\text{-}neg}(a) < b))}}
\vnbprop{ext-psum-not-pos-inf}{\vnbstmt{\begin{array}{@{}l@{}} \forall x.; \forall n \in \mathbb{N}. \\ \quad (\forall k \in \mathbb{N}.\; (x(k) \in \operatorname{rr\text{-}star}) \wedge \forall k \in \mathbb{N}.\; \neg (x(k) = \operatorname{pos\text{-}inf})) \Rightarrow  \\ \quad \quad \neg (\operatorname{ext\text{-}partial\text{-}sum}(x, n) = \operatorname{pos\text{-}inf}) \end{array}}}
\vnbprop{ext-psum-not-neg-inf}{\vnbstmt{\begin{array}{@{}l@{}} \forall x.; \forall n \in \mathbb{N}. \\ \quad (\forall k \in \mathbb{N}.\; (x(k) \in \operatorname{rr\text{-}star}) \wedge \forall k \in \mathbb{N}.\; \neg (x(k) = \operatorname{neg\text{-}inf})) \Rightarrow  \\ \quad \quad \neg (\operatorname{ext\text{-}partial\text{-}sum}(x, n) = \operatorname{neg\text{-}inf}) \end{array}}}
\vnbprop{ext-psum-negate-up}{\vnbstmt{\begin{array}{@{}l@{}} \forall n \in \mathbb{N}, x, y. \\ \quad \forall k \in \mathbb{N}.\; (x(k) \in \operatorname{rr\text{-}star}) \wedge \\ \quad \forall k \in \mathbb{N}.\; (y(k) = \operatorname{ext\text{-}neg}(x(k))) \wedge \\ \quad \forall k \in \mathbb{N}.\; \neg (x(k) = \operatorname{pos\text{-}inf}) \Rightarrow \\ \quad \quad \operatorname{ext\text{-}partial\text{-}sum}(y, n) = \operatorname{ext\text{-}neg}(\operatorname{ext\text{-}partial\text{-}sum}(x, n)) \end{array}}}
\vnbprop{ext-psum-negate-down}{\vnbstmt{\begin{array}{@{}l@{}} \forall n \in \mathbb{N}, x, y. \\ \quad \forall k \in \mathbb{N}.\; (x(k) \in \operatorname{rr\text{-}star}) \wedge \\ \quad \forall k \in \mathbb{N}.\; (y(k) = \operatorname{ext\text{-}neg}(x(k))) \wedge \\ \quad \forall k \in \mathbb{N}.\; \neg (x(k) = \operatorname{neg\text{-}inf}) \Rightarrow \\ \quad \quad \operatorname{ext\text{-}partial\text{-}sum}(y, n) = \operatorname{ext\text{-}neg}(\operatorname{ext\text{-}partial\text{-}sum}(x, n)) \end{array}}}
\vnbprop{ext-series-negate}{\vnbstmt{\begin{array}{@{}l@{}} \forall x, y, l. \\ \quad \operatorname{ext\text{-}series\text{-}converges\text{-}to}(x, l) \wedge \\ \quad \forall n \in \mathbb{N}.\; (\operatorname{ext\text{-}partial\text{-}sum}(y, n) = \operatorname{ext\text{-}neg}(\operatorname{ext\text{-}partial\text{-}sum}(x, n))) \Rightarrow \\ \quad \quad \operatorname{ext\text{-}series\text{-}converges\text{-}to}(y, \operatorname{ext\text{-}neg}(l)) \end{array}}}
\vnbprop{signed-neg-value}{\vnbstmt{\forall b, u.; \forall a \in b.\; (\operatorname{signed\text{-}neg}(b, u)(a) \simeq \operatorname{ext\text{-}neg}(u(a)))}}
\vnbprop{signed-neg-fun}{\vnbstmt{\begin{array}{@{}l@{}} \forall a, b, u. \\ \quad (\operatorname{is\text{-}sigma\text{-}algebra}(a, b) \wedge (u \in (b \to \operatorname{rr\text{-}star}))) \Rightarrow  \\ \quad \quad \operatorname{signed\text{-}neg}(b, u) \in (b \to \operatorname{rr\text{-}star}) \end{array}}}
\vnbprop{signed-neg-is-signed-measure}{\vnbstmt{\begin{array}{@{}l@{}} \forall a, b, u. \\ \quad \operatorname{is\text{-}signed\text{-}measure}(a, b, u) \Rightarrow \operatorname{is\text{-}signed\text{-}measure}(a, b, \operatorname{signed\text{-}neg}(b, u)) \end{array}}}
\vnbprop{neg-variation-is-measure}{\vnbstmt{\begin{array}{@{}l@{}} \forall a, b, u. \\ \quad \operatorname{is\text{-}signed\text{-}measure}(a, b, u) \Rightarrow \operatorname{is\text{-}measure}(a, b, \operatorname{neg\text{-}variation}(b, u)) \end{array}}}
\vnbprop{nn-above-real}{\vnbstmt{\forall x \in \mathbb{R}.\; \exists n \in \mathbb{N}.\; (x \leq n)}}
\vnbprop{ext-add-ge-left}{\vnbstmt{\begin{array}{@{}l@{}} \forall a, b. \\ \quad ((a \in \operatorname{rr\text{-}star}) \wedge ((b \in \operatorname{rr\text{-}star}) \wedge (0 \leq a))) \Rightarrow (b \leq \operatorname{ext\text{-}add}(b, a)) \end{array}}}
\vnbprop{ext-add-ge-right}{\vnbstmt{\begin{array}{@{}l@{}} \forall a, b. \\ \quad ((a \in \operatorname{rr\text{-}star}) \wedge ((b \in \operatorname{rr\text{-}star}) \wedge (0 \leq a))) \Rightarrow (b \leq \operatorname{ext\text{-}add}(a, b)) \end{array}}}
\vnbprop{ext-psum-mono-nonneg}{\vnbstmt{\begin{array}{@{}l@{}} \forall d \in \mathbb{N}, x.; \forall m \in \mathbb{N}. \\ \quad (\forall k \in \mathbb{N}.\; (x(k) \in \operatorname{rr\text{-}star}) \wedge \forall k \in \mathbb{N}.\; (0 \leq x(k))) \Rightarrow  \\ \quad \quad \operatorname{ext\text{-}partial\text{-}sum}(x, m) \leq \operatorname{ext\text{-}partial\text{-}sum}(x, (m + d)) \end{array}}}
\vnbprop{ext-series-unbounded-nonneg}{\vnbstmt{\begin{array}{@{}l@{}} \forall x, l. \\ \quad \forall k \in \mathbb{N}.\; (x(k) \in \operatorname{rr\text{-}star}) \wedge \\ \quad \forall k \in \mathbb{N}.\; (0 \leq x(k)) \wedge \\ \quad \forall b \in \mathbb{R}.\; \exists k \in \mathbb{N}.\; (b \leq x(k)) \wedge \\ \quad \operatorname{ext\text{-}series\text{-}converges\text{-}to}(x, l) \Rightarrow \\ \quad \quad l = \operatorname{pos\text{-}inf} \end{array}}}
\vnbprop{inter-empty-comm}{\vnbstmt{\begin{array}{@{}l@{}} \forall a, b. \\ \quad ((a \cap b) = \emptyset) \Rightarrow ((b \cap a) = \emptyset) \end{array}}}
\vnbprop{ext-real-of-ge}{\vnbstmt{\begin{array}{@{}l@{}} \forall x, c. \\ \quad ((x \in \operatorname{rr\text{-}star}) \wedge ((c \in \mathbb{R}) \wedge ((c \leq x) \wedge \neg (x = \operatorname{pos\text{-}inf})))) \Rightarrow (x \in \mathbb{R}) \end{array}}}
\vnbprop{signed-union-unbounded}{\vnbstmt{\begin{array}{@{}l@{}} \forall a, b, u, g. \\ \quad \operatorname{is\text{-}signed\text{-}measure}(a, b, u) \wedge \\ \quad g \in (\mathbb{N} \to b) \wedge \\ \quad \forall i, j. \\ \quad \quad ((i \in \mathbb{N}) \wedge ((j \in \mathbb{N}) \wedge \neg (i = j))) \Rightarrow ((g(i) \cap g(j)) = \emptyset) \wedge \\ \quad \forall k \in \mathbb{N}.\; (0 \leq u(g(k))) \wedge \\ \quad \forall \mathit{cb} \in \mathbb{R}.\; \exists k \in \mathbb{N}.\; (\mathit{cb} \leq u(g(k))) \Rightarrow \\ \quad \quad u(\operatorname{big\text{-}union}(n, \mathbb{N}, g(n))) = \operatorname{pos\text{-}inf} \end{array}}}
\vnbprop{diff-family-exists}{\vnbstmt{\begin{array}{@{}l@{}} \forall a, b, q. \\ \quad (\operatorname{is\text{-}sigma\text{-}algebra}(a, b) \wedge ((q \in (\mathbb{N} \to b)) \wedge \forall k \in \mathbb{N}.\; (q((k + 1)) \subseteq q(k)))) \Rightarrow  \\ \quad \quad \exists c \in (\mathbb{N} \to b). \\ \quad \quad \quad \forall i, j. \\ \quad \quad \quad \quad ((i \in \mathbb{N}) \wedge ((j \in \mathbb{N}) \wedge \neg (i = j))) \Rightarrow ((c(i) \cap c(j)) = \emptyset) \wedge \\ \quad \quad \quad \forall k \in \mathbb{N}.\; (c(k) = \operatorname{complement\text{-}in}(q(k), q((k + 1)))) \end{array}}}
\vnbprop{signed-telescope}{\vnbstmt{\begin{array}{@{}l@{}} \forall n \in \mathbb{N}, a, b, u, q, c. \\ \quad \operatorname{is\text{-}signed\text{-}measure}(a, b, u) \wedge \\ \quad \forall \alpha \in b.\; \neg (u(\alpha) = \operatorname{pos\text{-}inf}) \wedge \\ \quad q \in (\mathbb{N} \to b) \wedge \\ \quad \forall k \in \mathbb{N}.\; (q((k + 1)) \subseteq q(k)) \wedge \\ \quad \forall k \in \mathbb{N}.\; (k \leq u(q(k))) \wedge \\ \quad c \in (\mathbb{N} \to b) \wedge \\ \quad \forall k \in \mathbb{N}.\; (c(k) = \operatorname{complement\text{-}in}(q(k), q((k + 1)))) \Rightarrow \\ \quad \quad \operatorname{ext\text{-}partial\text{-}sum}((\lambda \mathit{smn} \in \mathbb{N}.\; u(c(\mathit{smn}))), n) \in \mathbb{R} \wedge \\ \quad \quad \operatorname{ext\text{-}partial\text{-}sum}((\lambda \mathit{smn} \in \mathbb{N}.\; u(c(\mathit{smn}))), n) = (u(q(0)) - u(q(n))) \end{array}}}
\vnbprop{signed-decreasing-neg-inf}{\vnbstmt{\begin{array}{@{}l@{}} \forall a, b, u, q. \\ \quad \operatorname{is\text{-}signed\text{-}measure}(a, b, u) \wedge \\ \quad \forall \alpha \in b.\; \neg (u(\alpha) = \operatorname{pos\text{-}inf}) \wedge \\ \quad q \in (\mathbb{N} \to b) \wedge \\ \quad \forall k \in \mathbb{N}.\; (q((k + 1)) \subseteq q(k)) \wedge \\ \quad \forall k \in \mathbb{N}.\; (k \leq u(q(k))) \Rightarrow \\ \quad \quad u(q(0)) = \operatorname{neg\text{-}inf} \end{array}}}
\vnbprop{rps-real-of-not-pos-inf}{\vnbstmt{\forall x \in \operatorname{rr\text{-}pos\text{-}star}.\; (\neg (x = \operatorname{pos\text{-}inf}) \Rightarrow (x \in \mathbb{R}))}}
\vnbprop{pos-variation-step}{\vnbstmt{\begin{array}{@{}l@{}} \forall a, b, c.; \forall k \in \mathbb{N}, u \in \{\mathit{su} \in b : (\operatorname{pos\text{-}variation}(b, c)(\mathit{su}) = \operatorname{pos\text{-}inf})\}. \\ \quad \operatorname{is\text{-}signed\text{-}measure}(a, b, c) \Rightarrow  \\ \quad \quad \exists y \in \{\mathit{su} \in b : (\operatorname{pos\text{-}variation}(b, c)(\mathit{su}) = \operatorname{pos\text{-}inf})\}. \\ \quad \quad \quad (y \subseteq u) \wedge (((k + 1) \leq c(y)) \vee ((k + 1) \leq c(\operatorname{complement\text{-}in}(u, y)))) \end{array}}}
\vnbprop{pos-variation-chain}{\vnbstmt{\begin{array}{@{}l@{}} \forall a, b, u, c. \\ \quad (\operatorname{is\text{-}signed\text{-}measure}(a, b, u) \wedge ((c \in b) \wedge (\operatorname{pos\text{-}variation}(b, u)(c) = \operatorname{pos\text{-}inf}))) \Rightarrow  \\ \quad \quad \exists f \in (\mathbb{N} \to b). \\ \quad \quad \quad f(0) = c \wedge \\ \quad \quad \quad \forall k \in \mathbb{N}. \\ \quad \quad \quad \quad \operatorname{pos\text{-}variation}(b, u)(f(k)) = \operatorname{pos\text{-}inf} \wedge \\ \quad \quad \quad \quad f((k + 1)) \subseteq f(k) \wedge \\ \quad \quad \quad \quad ((k + 1) \leq u(f((k + 1)))) \vee ((k + 1) \leq u(\operatorname{complement\text{-}in}(f(k), f((k + 1))))) \end{array}}}
\vnbprop{pos-variation-finite}{\vnbstmt{\begin{array}{@{}l@{}} \forall a, b, u.; \forall c \in b. \\ \quad (\operatorname{is\text{-}signed\text{-}measure}(a, b, u) \wedge \forall \alpha \in b.\; \neg (u(\alpha) = \operatorname{pos\text{-}inf})) \Rightarrow  \\ \quad \quad \operatorname{pos\text{-}variation}(b, u)(c) \in \mathbb{R} \end{array}}}

\vnbbookfile{theorem-library/lp-inequalities}{theorem-library/lp-inequalities}
\noindent{\small theorem-library/lp-inequalities.scm -- POWERS IN [0,+inf], HOELDER AND MINKOWSKI, FUNCTIONS EQUAL ALMOST EVERYWHERE.  MEASURE-4-B, 2026-10-10.}\par
\vnbprop{epow-unfold}{\vnbstmt{\forall t, p.\; ({t}^{p} \simeq \left(\text{if }(t = \operatorname{pos\text{-}inf})\text{ then }\operatorname{pos\text{-}inf}\text{ else }{t}^{p}\right))}}
\vnbprop{ep-norm-unfold}{\vnbstmt{\begin{array}{@{}l@{}} \forall a, b, u, p, f. \\ \quad (I_{p}({f}) \simeq {\operatorname{integral}(a, b, u, \operatorname{ptwise\text{-}epow}(a, f, p))}^{\operatorname{recip}(p)}) \end{array}}}
\vnbprop{epow-pos-inf}{\vnbstmt{\forall p.\; ({\operatorname{pos\text{-}inf}}^{p} = \operatorname{pos\text{-}inf})}}
\vnbprop{epow-real}{\vnbstmt{\forall t \in \mathbb{R}, p.\; ({t}^{p} \simeq {t}^{p})}}
\vnbprop{epow-zero}{\vnbstmt{\forall p \in \mathbb{R}.\; ((0 < p) \Rightarrow ({0}^{p} = 0))}}
\vnbprop{epow-one-base}{\vnbstmt{\forall p \in \mathbb{R}.\; ({1}^{p} = 1)}}
\vnbprop{epow-pos-real}{\vnbstmt{\forall t, p \in \mathbb{R}.\; ((0 < t) \Rightarrow (({t}^{p} \in \mathbb{R}) \wedge (0 < {t}^{p})))}}
\vnbprop{epow-real-in}{\vnbstmt{\begin{array}{@{}l@{}} \forall t, p \in \mathbb{R}. \\ \quad ((0 \leq t) \wedge (0 < p)) \Rightarrow (({t}^{p} \in \mathbb{R}) \wedge (0 \leq {t}^{p})) \end{array}}}
\vnbprop{epow-in}{\vnbstmt{\forall t \in \operatorname{rr\text{-}pos\text{-}star}, p \in \mathbb{R}.\; ((0 < p) \Rightarrow ({t}^{p} \in \operatorname{rr\text{-}pos\text{-}star}))}}
\vnbprop{epow-strict-mono}{\vnbstmt{\begin{array}{@{}l@{}} \forall s, t \in \operatorname{rr\text{-}pos\text{-}star}, p \in \mathbb{R}. \\ \quad ((0 < p) \wedge (s < t)) \Rightarrow ({s}^{p} < {t}^{p}) \end{array}}}
\vnbprop{epow-mono}{\vnbstmt{\begin{array}{@{}l@{}} \forall s, t \in \operatorname{rr\text{-}pos\text{-}star}, p \in \mathbb{R}. \\ \quad ((0 < p) \wedge (s \leq t)) \Rightarrow ({s}^{p} \leq {t}^{p}) \end{array}}}
\vnbprop{epow-epow}{\vnbstmt{\begin{array}{@{}l@{}} \forall t \in \operatorname{rr\text{-}pos\text{-}star}, a, b \in \mathbb{R}. \\ \quad ((0 < a) \wedge (0 < b)) \Rightarrow ({{t}^{a}}^{b} = {t}^{(a \cdot b)}) \end{array}}}
\vnbprop{epow-one}{\vnbstmt{\forall t \in \operatorname{rr\text{-}pos\text{-}star}.\; ({t}^{1} = t)}}
\vnbprop{epow-inverse}{\vnbstmt{\begin{array}{@{}l@{}} \forall t \in \operatorname{rr\text{-}pos\text{-}star}, p \in \mathbb{R}. \\ \quad (0 < p) \Rightarrow (({{t}^{p}}^{\operatorname{recip}(p)} = t) \wedge ({{t}^{\operatorname{recip}(p)}}^{p} = t)) \end{array}}}
\vnbprop{epow-eq-zero}{\vnbstmt{\forall t \in \operatorname{rr\text{-}pos\text{-}star}, p \in \mathbb{R}.\; (((0 < p) \wedge ({t}^{p} = 0)) \Rightarrow (t = 0))}}
\vnbprop{epow-eq-pos-inf}{\vnbstmt{\begin{array}{@{}l@{}} \forall t \in \operatorname{rr\text{-}pos\text{-}star}, p \in \mathbb{R}. \\ \quad ((0 < p) \wedge ({t}^{p} = \operatorname{pos\text{-}inf})) \Rightarrow (t = \operatorname{pos\text{-}inf}) \end{array}}}
\vnbprop{epow-etimes}{\vnbstmt{\begin{array}{@{}l@{}} \forall c \in \mathbb{R}, u \in \operatorname{rr\text{-}pos\text{-}star}, p \in \mathbb{R}. \\ \quad ((0 < c) \wedge (0 < p)) \Rightarrow ({\operatorname{etimes}(c, u)}^{p} = \operatorname{etimes}({c}^{p}, {u}^{p})) \end{array}}}
\vnbprop{ptwise-epow-value}{\vnbstmt{\begin{array}{@{}l@{}} \forall a.; \forall f \in (a \to \operatorname{rr\text{-}pos\text{-}star}), p \in \mathbb{R}, x \in a. \\ \quad (0 < p) \Rightarrow (\operatorname{ptwise\text{-}epow}(a, f, p)(x) = {f(x)}^{p}) \end{array}}}
\vnbprop{measurable-fn-epow}{\vnbstmt{\begin{array}{@{}l@{}} \forall a, b, f, p. \\ \quad (\operatorname{is\text{-}measurable\text{-}fn}(a, b, f) \wedge ((p \in \mathbb{R}) \wedge (0 < p))) \Rightarrow  \\ \quad \quad \operatorname{is\text{-}measurable\text{-}fn}(a, b, \operatorname{ptwise\text{-}epow}(a, f, p)) \end{array}}}
\vnbprop{ptwise-eplus-value}{\vnbstmt{\begin{array}{@{}l@{}} \forall a, f, g.; \forall x \in a. \\ \quad ((f(x) \in \operatorname{rr\text{-}pos\text{-}star}) \wedge (g(x) \in \operatorname{rr\text{-}pos\text{-}star})) \Rightarrow  \\ \quad \quad \operatorname{ptwise\text{-}eplus}(a, f, g)(x) = \operatorname{eplus}(f(x), g(x)) \end{array}}}
\vnbprop{ptwise-etimes-value}{\vnbstmt{\begin{array}{@{}l@{}} \forall a, f, g.; \forall x \in a. \\ \quad ((f(x) \in \operatorname{rr\text{-}pos\text{-}star}) \wedge (g(x) \in \operatorname{rr\text{-}pos\text{-}star})) \Rightarrow  \\ \quad \quad \operatorname{ptwise\text{-}etimes}(a, f, g)(x) = \operatorname{etimes}(f(x), g(x)) \end{array}}}
\vnbprop{ptwise-scale-value}{\vnbstmt{\begin{array}{@{}l@{}} \forall a, c, f.; \forall x \in a. \\ \quad ((c \in \operatorname{rr\text{-}pos\text{-}star}) \wedge (f(x) \in \operatorname{rr\text{-}pos\text{-}star})) \Rightarrow  \\ \quad \quad \operatorname{ptwise\text{-}scale}(a, c, f)(x) = \operatorname{etimes}(c, f(x)) \end{array}}}
\vnbprop{young-scaled}{\vnbstmt{\begin{array}{@{}l@{}} \forall c, d \in \mathbb{R}, u, v \in \operatorname{rr\text{-}pos\text{-}star}, p, q \in \mathbb{R}. \\ \quad ((0 < c) \wedge ((0 < d) \wedge ((1 < p) \wedge ((1 < q) \wedge ((/(1, p) + /(1, q)) = 1))))) \Rightarrow  \\ \quad \quad \operatorname{etimes}((c \cdot d), \operatorname{etimes}(u, v)) \\ \quad \quad \leq\; \operatorname{eplus}(\operatorname{etimes}((\operatorname{recip}(p) \cdot {c}^{p}), {u}^{p}), \operatorname{etimes}((\operatorname{recip}(q) \cdot {d}^{q}), {v}^{q})) \end{array}}}
\vnbprop{young-extended}{\vnbstmt{\begin{array}{@{}l@{}} \forall a, b \in \operatorname{rr\text{-}pos\text{-}star}, p, q \in \mathbb{R}. \\ \quad ((1 < p) \wedge ((1 < q) \wedge ((/(1, p) + /(1, q)) = 1))) \Rightarrow  \\ \quad \quad \operatorname{etimes}(a, b) \leq \operatorname{eplus}(\operatorname{etimes}(\operatorname{recip}(p), {a}^{p}), \operatorname{etimes}(\operatorname{recip}(q), {b}^{q})) \end{array}}}
\vnbprop{etimes-eq-zero-pos}{\vnbstmt{\forall c \in \mathbb{R}, m \in \operatorname{rr\text{-}pos\text{-}star}.\; (((0 < c) \wedge (\operatorname{etimes}(c, m) = 0)) \Rightarrow (m = 0))}}
\vnbprop{integral-markov}{\vnbstmt{\begin{array}{@{}l@{}} \forall a, b, u, h, c. \\ \quad (\operatorname{is\text{-}measure}(a, b, u) \wedge (\operatorname{is\text{-}measurable\text{-}fn}(a, b, h) \wedge ((c \in \mathbb{R}) \wedge (0 \leq c)))) \Rightarrow  \\ \quad \quad \operatorname{etimes}(c, u(\{x \in a : (c \leq h(x))\})) \leq \operatorname{integral}(a, b, u, h) \end{array}}}
\vnbprop{integral-zero-null}{\vnbstmt{\begin{array}{@{}l@{}} \forall a, b, u, h. \\ \quad (\operatorname{is\text{-}measure}(a, b, u) \wedge (\operatorname{is\text{-}measurable\text{-}fn}(a, b, h) \wedge (\operatorname{integral}(a, b, u, h) = 0))) \Rightarrow  \\ \quad \quad u(\{x \in a : (0 < h(x))\}) = 0 \end{array}}}
\vnbprop{ep-norm-zero-null}{\vnbstmt{\begin{array}{@{}l@{}} \forall a, b, u, f, p. \\ \quad \operatorname{is\text{-}measure}(a, b, u) \wedge \\ \quad \operatorname{is\text{-}measurable\text{-}fn}(a, b, f) \wedge \\ \quad p \in \mathbb{R} \wedge \\ \quad 0 < p \wedge \\ \quad I_{p}({f}) = 0 \Rightarrow \\ \quad \quad (\{x \in a : (0 < f(x))\} \in b) \wedge (u(\{x \in a : (0 < f(x))\}) = 0) \end{array}}}
\vnbprop{integral-ae-eq}{\vnbstmt{\begin{array}{@{}l@{}} \forall a, b, u, c, f, g. \\ \quad \operatorname{is\text{-}measure}(a, b, u) \wedge \\ \quad c \in b \wedge \\ \quad u(c) = 0 \wedge \\ \quad \operatorname{is\text{-}measurable\text{-}fn}(a, b, f) \wedge \\ \quad \operatorname{is\text{-}measurable\text{-}fn}(a, b, g) \wedge \\ \quad \forall x \in a.\; (\neg (x \in c) \Rightarrow (f(x) = g(x))) \Rightarrow \\ \quad \quad \operatorname{integral}(a, b, u, f) = \operatorname{integral}(a, b, u, g) \end{array}}}
\vnbprop{integral-ptwise-eq}{\vnbstmt{\begin{array}{@{}l@{}} \forall a, b, u, f, g. \\ \quad \operatorname{is\text{-}measure}(a, b, u) \wedge \\ \quad \operatorname{is\text{-}measurable\text{-}fn}(a, b, f) \wedge \\ \quad \operatorname{is\text{-}measurable\text{-}fn}(a, b, g) \wedge \\ \quad \forall x \in a.\; (f(x) = g(x)) \Rightarrow \\ \quad \quad \operatorname{integral}(a, b, u, f) = \operatorname{integral}(a, b, u, g) \end{array}}}
\vnbprop{integral-finite-ae-finite}{\vnbstmt{\begin{array}{@{}l@{}} \forall a, b, u, f. \\ \quad (\operatorname{is\text{-}measure}(a, b, u) \wedge (\operatorname{is\text{-}measurable\text{-}fn}(a, b, f) \wedge (\operatorname{integral}(a, b, u, f) < \operatorname{pos\text{-}inf}))) \Rightarrow  \\ \quad \quad (\{x \in a : (f(x) = \operatorname{pos\text{-}inf})\} \in b) \wedge (u(\{x \in a : (f(x) = \operatorname{pos\text{-}inf})\}) = 0) \end{array}}}
\vnbprop{hoelder}{\vnbstmt{\begin{array}{@{}l@{}} \forall a, b, u, f, g.; \forall p, q \in \mathbb{R}. \\ \quad \operatorname{is\text{-}measure}(a, b, u) \wedge \\ \quad \operatorname{is\text{-}measurable\text{-}fn}(a, b, f) \wedge \\ \quad \operatorname{is\text{-}measurable\text{-}fn}(a, b, g) \wedge \\ \quad 1 < p \wedge \\ \quad 1 < q \wedge \\ \quad (/(1, p) + /(1, q)) = 1 \Rightarrow \\ \quad \quad \operatorname{integral}(a, b, u, \operatorname{ptwise\text{-}etimes}(a, f, g)) \leq \operatorname{etimes}(I_{p}({f}), I_{q}({g})) \end{array}}}
\vnbprop{conjugate-exponent}{\vnbstmt{\begin{array}{@{}l@{}} \forall p \in \mathbb{R}. \\ \quad (1 < p) \Rightarrow  \\ \quad \quad (p \cdot \operatorname{recip}((p - 1))) \in \mathbb{R} \wedge \\ \quad \quad 1 < (p \cdot \operatorname{recip}((p - 1))) \wedge \\ \quad \quad (/(1, p) + /(1, (p \cdot \operatorname{recip}((p - 1))))) = 1 \wedge \\ \quad \quad ((p - 1) \cdot (p \cdot \operatorname{recip}((p - 1)))) = p \end{array}}}
\vnbprop{ep-norm-one}{\vnbstmt{\begin{array}{@{}l@{}} \forall a, b, u, h. \\ \quad (\operatorname{is\text{-}measure}(a, b, u) \wedge \operatorname{is\text{-}measurable\text{-}fn}(a, b, h)) \Rightarrow  \\ \quad \quad I_{1}({h}) = \operatorname{integral}(a, b, u, h) \end{array}}}
\vnbprop{epow-split}{\vnbstmt{\begin{array}{@{}l@{}} \forall t \in \operatorname{rr\text{-}pos\text{-}star}, p \in \mathbb{R}. \\ \quad (1 < p) \Rightarrow ({t}^{p} = \operatorname{etimes}({t}^{(p - 1)}, t)) \end{array}}}
\vnbprop{etimes-distrib-left-star}{\vnbstmt{\begin{array}{@{}l@{}} \forall c \in \operatorname{rr\text{-}pos\text{-}star}, x, y. \\ \quad ((x \in \operatorname{rr\text{-}pos\text{-}star}) \wedge (y \in \operatorname{rr\text{-}pos\text{-}star})) \Rightarrow  \\ \quad \quad \operatorname{etimes}(c, \operatorname{eplus}(x, y)) = \operatorname{eplus}(\operatorname{etimes}(c, x), \operatorname{etimes}(c, y)) \end{array}}}
\vnbprop{epow-eplus-le}{\vnbstmt{\begin{array}{@{}l@{}} \forall u, v \in \operatorname{rr\text{-}pos\text{-}star}, p \in \mathbb{R}. \\ \quad (0 < p) \Rightarrow  \\ \quad \quad {\operatorname{eplus}(u, v)}^{p} \leq \operatorname{etimes}({2}^{p}, \operatorname{eplus}({u}^{p}, {v}^{p})) \end{array}}}
\vnbprop{minkowski}{\vnbstmt{\begin{array}{@{}l@{}} \forall a, b, u, f, g.; \forall p \in \mathbb{R}. \\ \quad (\operatorname{is\text{-}measure}(a, b, u) \wedge (\operatorname{is\text{-}measurable\text{-}fn}(a, b, f) \wedge (\operatorname{is\text{-}measurable\text{-}fn}(a, b, g) \wedge (1 \leq p)))) \Rightarrow  \\ \quad \quad I_{p}({\operatorname{ptwise\text{-}eplus}(a, f, g)}) \leq \operatorname{eplus}(I_{p}({f}), I_{p}({g})) \end{array}}}

\vnbbookfile{structure-library/lp-seminorm}{structure-library/lp-seminorm}
\noindent{\small lp-seminorm.scm -- THE L\textasciicircum{}p SEMINORM OF A REAL-VALUED FUNCTION. Definitions only; nothing here is asserted and nothing is proved.  The laws are proved in theorem-library/lp-seminorm.scm (MEASURE-4-C, 2026-10-10).}\par
\par\medskip\noindent\textbf{Definition} (\texttt{lp-norm}).\;$\|{\operatorname{lpf\_}}\|_{\operatorname{lpp\_}} :=$ $I_{\operatorname{lpp\_}}({\operatorname{abs\text{-}fn}(\omega, \operatorname{lpf\_})})$

\vnbbookfile{theorem-library/lp-seminorm}{theorem-library/lp-seminorm}
\noindent{\small lp-seminorm.scm -- THE L\textasciicircum{}p SEMINORM: the seminorm laws (MEASURE-4-C, 2026-10-10).  Vocabulary: structure-library/lp-seminorm.scm (LP-NORM = EP-NORM of ABS-FN).}\par
\vnbprop{lp-norm-unfold}{\vnbstmt{\begin{array}{@{}l@{}} \forall a, b, u, p, f. \\ \quad (\|{f}\|_{p} \simeq I_{p}({\operatorname{abs\text{-}fn}(a, f)})) \end{array}}}
\vnbprop{real-measurable-add}{\vnbstmt{\begin{array}{@{}l@{}} \forall a, b, f, g. \\ \quad (\operatorname{is\text{-}real\text{-}measurable\text{-}fn}(a, b, f) \wedge \operatorname{is\text{-}real\text{-}measurable\text{-}fn}(a, b, g)) \Rightarrow  \\ \quad \quad \operatorname{is\text{-}real\text{-}measurable\text{-}fn}(a, b, (\lambda x \in a.\; (f(x) + g(x)))) \end{array}}}
\vnbprop{real-measurable-sub}{\vnbstmt{\begin{array}{@{}l@{}} \forall a, b, f, g. \\ \quad (\operatorname{is\text{-}real\text{-}measurable\text{-}fn}(a, b, f) \wedge \operatorname{is\text{-}real\text{-}measurable\text{-}fn}(a, b, g)) \Rightarrow  \\ \quad \quad \operatorname{is\text{-}real\text{-}measurable\text{-}fn}(a, b, (\lambda x \in a.\; (f(x) - g(x)))) \end{array}}}
\vnbprop{ep-norm-mono}{\vnbstmt{\begin{array}{@{}l@{}} \forall a, b, u, p, c, d. \\ \quad \operatorname{is\text{-}measure}(a, b, u) \wedge \\ \quad \operatorname{is\text{-}measurable\text{-}fn}(a, b, c) \wedge \\ \quad \operatorname{is\text{-}measurable\text{-}fn}(a, b, d) \wedge \\ \quad \forall x \in a.\; (c(x) \leq d(x)) \wedge \\ \quad p \in \mathbb{R} \wedge \\ \quad 0 < p \Rightarrow \\ \quad \quad I_{p}({c}) \leq I_{p}({d}) \end{array}}}
\vnbprop{ep-norm-scale}{\vnbstmt{\begin{array}{@{}l@{}} \forall a, b, u, p, k, h. \\ \quad (\operatorname{is\text{-}measure}(a, b, u) \wedge (\operatorname{is\text{-}measurable\text{-}fn}(a, b, h) \wedge ((k \in \mathbb{R}) \wedge ((0 \leq k) \wedge ((p \in \mathbb{R}) \wedge (0 < p)))))) \Rightarrow  \\ \quad \quad I_{p}({\operatorname{ptwise\text{-}scale}(a, k, h)}) = \operatorname{etimes}(k, I_{p}({h})) \end{array}}}
\vnbprop{lp-norm-homogeneous}{\vnbstmt{\begin{array}{@{}l@{}} \forall a, b, u, p, f, c. \\ \quad (\operatorname{is\text{-}measure}(a, b, u) \wedge (\operatorname{is\text{-}real\text{-}measurable\text{-}fn}(a, b, f) \wedge ((c \in \mathbb{R}) \wedge ((p \in \mathbb{R}) \wedge (0 < p))))) \Rightarrow  \\ \quad \quad \|{(\lambda x \in a.\; (c \cdot f(x)))}\|_{p} = \operatorname{etimes}(\lvert c \rvert, \|{f}\|_{p}) \end{array}}}
\vnbprop{lp-norm-triangle}{\vnbstmt{\begin{array}{@{}l@{}} \forall a, b, u, p, f, g. \\ \quad \operatorname{is\text{-}measure}(a, b, u) \wedge \\ \quad \operatorname{is\text{-}real\text{-}measurable\text{-}fn}(a, b, f) \wedge \\ \quad \operatorname{is\text{-}real\text{-}measurable\text{-}fn}(a, b, g) \wedge \\ \quad p \in \mathbb{R} \wedge \\ \quad 1 \leq p \Rightarrow \\ \quad \quad \|{(\lambda x \in a.\; (f(x) + g(x)))}\|_{p} \\ \quad \quad \leq\; \operatorname{eplus}(\|{f}\|_{p}, \|{g}\|_{p}) \end{array}}}
\vnbprop{lp-norm-zero-iff}{\vnbstmt{\begin{array}{@{}l@{}} \forall a, b, u, p, f, g. \\ \quad \operatorname{is\text{-}measure}(a, b, u) \wedge \\ \quad \operatorname{is\text{-}real\text{-}measurable\text{-}fn}(a, b, f) \wedge \\ \quad \operatorname{is\text{-}real\text{-}measurable\text{-}fn}(a, b, g) \wedge \\ \quad p \in \mathbb{R} \wedge \\ \quad 0 < p \Rightarrow \\ \quad \quad \|{(\lambda \mathit{sx} \in a.\; (f(\mathit{sx}) - g(\mathit{sx})))}\|_{p} = 0 \\ \quad \quad \Leftrightarrow\; \exists s. \\ \quad \quad \quad \operatorname{is\text{-}null\text{-}set}(a, b, u, s) \wedge \forall x \in a.\; (\neg (x \in s) \Rightarrow (f(x) = g(x))) \end{array}}}

\vnbbookfile{theorem-library/lp-convergence}{theorem-library/lp-convergence}
\noindent{\small lp-convergence.scm -- WAVE 2 (2026-10-27..30), wired 2026-10-30 from scratchpad/md/md-lp.scm. Library: measure.}\par
\vnbprop{rps-below-all-positive-zero}{\vnbstmt{\forall x \in \operatorname{rr\text{-}pos\text{-}star}.\; (\forall r.\; (\operatorname{pos\text{-}rr}(r) \Rightarrow (x \leq r)) \Rightarrow (x = 0))}}
\vnbprop{econverges-of-eventually-small}{\vnbstmt{\begin{array}{@{}l@{}} \forall a \in (\mathbb{N} \to \mathbb{R}). \\ \quad \forall k \in \mathbb{N}.\; (0 \leq a(k)) \wedge \\ \quad \forall r. \\ \quad \quad \operatorname{pos\text{-}rr}(r) \Rightarrow  \\ \quad \quad \quad \exists b \in \mathbb{N}.\; \forall k \in \mathbb{N}.\; ((b \leq k) \Rightarrow (a(k) \leq r)) \Rightarrow \\ \quad \quad \operatorname{econverges\text{-}to}(a, 0) \end{array}}}
\vnbprop{epow-econverges-zero}{\vnbstmt{\begin{array}{@{}l@{}} \forall a, q. \\ \quad ((a \in (\mathbb{N} \to \mathbb{R})) \wedge (\forall k \in \mathbb{N}.\; (0 \leq a(k)) \wedge (\operatorname{econverges\text{-}to}(a, 0) \wedge ((q \in \mathbb{R}) \wedge (0 < q))))) \Rightarrow  \\ \quad \quad \operatorname{econverges\text{-}to}((\lambda k \in \mathbb{N}.\; {a(k)}^{q}), 0) \end{array}}}
\vnbprop{measure-null-union-2}{\vnbstmt{\begin{array}{@{}l@{}} \forall c, d, u, a, b. \\ \quad (\operatorname{is\text{-}measure}(c, d, u) \wedge ((a \in d) \wedge ((b \in d) \wedge ((u(a) = 0) \wedge (u(b) = 0))))) \Rightarrow  \\ \quad \quad ((a \cup b) \in d) \wedge (u((a \cup b)) = 0) \end{array}}}
\vnbprop{ae-hypotheses-common-null-set}{\vnbstmt{\begin{array}{@{}l@{}} \forall a, b, u, s, f, g. \\ \quad \operatorname{is\text{-}measure}(a, b, u) \wedge \\ \quad \forall k \in \mathbb{N}. \exists c. \\ \quad \quad (\operatorname{is\text{-}null\text{-}set}(a, b, u, c) \wedge \forall x \in a.\; (\neg (x \in c) \Rightarrow (\lvert s(k)(x) \rvert \leq g(x)))) \wedge \\ \quad \exists c. \\ \quad \quad \operatorname{is\text{-}null\text{-}set}(a, b, u, c) \wedge \\ \quad \quad \forall x \in a. \\ \quad \quad \quad \neg (x \in c) \Rightarrow \operatorname{converges\text{-}to}(\operatorname{rr\text{-}ms}, (\lambda k \in \mathbb{N}.\; s(k)(x)), f(x)) \Rightarrow \\ \quad \quad \exists \mathit{mda} \in b. \\ \quad \quad \quad u(\mathit{mda}) = 0 \wedge \\ \quad \quad \quad \forall x \in a. \\ \quad \quad \quad \quad \neg (x \in \mathit{mda}) \Rightarrow  \\ \quad \quad \quad \quad \quad \forall k \in \mathbb{N}.\; (\lvert s(k)(x) \rvert \leq g(x)) \wedge \\ \quad \quad \quad \quad \quad \operatorname{converges\text{-}to}(\operatorname{rr\text{-}ms}, (\lambda k \in \mathbb{N}.\; s(k)(x)), f(x)) \end{array}}}
\vnbprop{rps-below-real-is-real}{\vnbstmt{\forall x, y.\; (((x \in \operatorname{rr\text{-}pos\text{-}star}) \wedge ((y \in \mathbb{R}) \wedge (x \leq y))) \Rightarrow (x \in \mathbb{R}))}}
\vnbprop{epow-abs-diff-le}{\vnbstmt{\begin{array}{@{}l@{}} \forall a, b, c, p. \\ \quad ((a \in \mathbb{R}) \wedge ((b \in \mathbb{R}) \wedge ((c \in \mathbb{R}) \wedge ((\lvert a \rvert \leq c) \wedge ((\lvert b \rvert \leq c) \wedge ((p \in \mathbb{R}) \wedge (0 < p))))))) \Rightarrow  \\ \quad \quad {\lvert (a - b) \rvert}^{p} \leq \operatorname{etimes}({2}^{p}, {\lvert c \rvert}^{p}) \end{array}}}
\vnbprop{econverges-abs-diff-epow}{\vnbstmt{\begin{array}{@{}l@{}} \forall u, l, p. \\ \quad ((u \in (\mathbb{N} \to \mathbb{R})) \wedge ((l \in \mathbb{R}) \wedge (\operatorname{converges\text{-}to}(\operatorname{rr\text{-}ms}, u, l) \wedge ((p \in \mathbb{R}) \wedge (0 < p))))) \Rightarrow  \\ \quad \quad \operatorname{econverges\text{-}to}((\lambda k \in \mathbb{N}.\; {\lvert (u(k) - l) \rvert}^{p}), 0) \end{array}}}
\vnbprop{lp-norm-as-epow}{\vnbstmt{\begin{array}{@{}l@{}} \forall a, b, u, p, h. \\ \quad (\|{h}\|_{p} \simeq {\operatorname{integral}(a, b, u, \operatorname{ptwise\text{-}epow}(a, \operatorname{abs\text{-}fn}(a, h), p))}^{\operatorname{recip}(p)}) \end{array}}}
\vnbprop{lp-norm-finite-iff}{\vnbstmt{\begin{array}{@{}l@{}} \forall a, b, u, p, h. \\ \quad (\operatorname{is\text{-}measure}(a, b, u) \wedge (\operatorname{is\text{-}real\text{-}measurable\text{-}fn}(a, b, h) \wedge ((p \in \mathbb{R}) \wedge (0 < p)))) \Rightarrow  \\ \quad \quad (\|{h}\|_{p} \in \mathbb{R}) \Leftrightarrow (\operatorname{integral}(a, b, u, \operatorname{ptwise\text{-}epow}(a, \operatorname{abs\text{-}fn}(a, h), p)) \in \mathbb{R}) \end{array}}}
\vnbprop{lp-dominated-convergence}{\vnbstmt{\begin{array}{@{}l@{}} \forall a, b, u, p, s, f, g. \\ \quad \operatorname{is\text{-}measure}(a, b, u) \wedge \\ \quad p \in \mathbb{R} \wedge \\ \quad 1 \leq p \wedge \\ \quad s \in (\mathbb{N} \to (a \to \mathbb{R})) \wedge \\ \quad \forall k \in \mathbb{N}.\; \operatorname{is\text{-}real\text{-}measurable\text{-}fn}(a, b, s(k)) \wedge \\ \quad \operatorname{is\text{-}real\text{-}measurable\text{-}fn}(a, b, f) \wedge \\ \quad \operatorname{is\text{-}real\text{-}measurable\text{-}fn}(a, b, g) \wedge \\ \quad \|{g}\|_{p} \in \mathbb{R} \wedge \\ \quad \forall k \in \mathbb{N}. \exists c. \\ \quad \quad (\operatorname{is\text{-}null\text{-}set}(a, b, u, c) \wedge \forall x \in a.\; (\neg (x \in c) \Rightarrow (\lvert s(k)(x) \rvert \leq g(x)))) \wedge \\ \quad \exists c. \\ \quad \quad \operatorname{is\text{-}null\text{-}set}(a, b, u, c) \wedge \\ \quad \quad \forall x \in a. \\ \quad \quad \quad \neg (x \in c) \Rightarrow \operatorname{converges\text{-}to}(\operatorname{rr\text{-}ms}, (\lambda k \in \mathbb{N}.\; s(k)(x)), f(x)) \Rightarrow \\ \quad \quad \|{f}\|_{p} \in \mathbb{R} \wedge \\ \quad \quad \operatorname{econverges\text{-}to}((\lambda k \in \mathbb{N}.\; \|{(\lambda \mathit{sx} \in a.\; (s(k)(\mathit{sx}) - f(\mathit{sx})))}\|_{p}), 0) \end{array}}}

\vnbbookfile{structure-library/product-measure}{structure-library/product-measure}
\noindent{\small product-measure.scm -- PRODUCT SIGMA-ALGEBRAS, SECTIONS, DYNKIN SYSTEMS. Vocabulary only: no proof, no support.  Agent FUBINI-1, 2026-10-10 (scratchpad/triage/FUBINI-1.md).  The proofs are theorem-library/product-measure.scm (prefix pms-) and theorem-library/fubini.scm (prefix fub-).}\par
\par\medskip\noindent\textbf{Definition} (\texttt{rectangles}).\;$\operatorname{rectangles}(\mathit{xs}, \mathit{ca}, \mathit{ys}, \mathit{cb}) :=$
\begin{equation*}
\{\operatorname{rcr\_} \in \operatorname{power}((\mathit{xs} \times \mathit{ys})) : \exists \operatorname{rca\_} \in \mathit{ca}, \operatorname{rcb\_} \in \mathit{cb}.\; (\operatorname{rcr\_} = (\operatorname{rca\_} \times \operatorname{rcb\_}))\}
\end{equation*}

\par\medskip\noindent\textbf{Definition} (\texttt{product-sigma}).\;$\operatorname{product\text{-}sigma}(\mathit{xs}, \mathit{ca}, \mathit{ys}, \mathit{cb}) :=$
\begin{equation*}
\operatorname{sigma\text{-}generated}((\mathit{xs} \times \mathit{ys}), \operatorname{rectangles}(\mathit{xs}, \mathit{ca}, \mathit{ys}, \mathit{cb}))
\end{equation*}

\par\medskip\noindent\textbf{Definition} (\texttt{x-section}).\;$\operatorname{x\text{-}section}(\mathit{ys}, \mathit{cs}, x) :=$ $\{\operatorname{xsy\_} \in \mathit{ys} : ({[}x, \operatorname{xsy\_}] \in \mathit{cs})\}$
\par\medskip\noindent\textbf{Definition} (\texttt{y-section}).\;$\operatorname{y\text{-}section}(\mathit{xs}, \mathit{cs}, y) :=$ $\{\operatorname{ysx\_} \in \mathit{xs} : ({[}\operatorname{ysx\_}, y] \in \mathit{cs})\}$
\par\medskip\noindent\textbf{Definition} (\texttt{fn-x-section}).\;$\operatorname{fn\text{-}x\text{-}section}(\mathit{ys}, f, x) :=$ $(\lambda \operatorname{fsy\_} \in \mathit{ys}.\; f({[}x, \operatorname{fsy\_}]))$
\par\medskip\noindent\textbf{Definition} (\texttt{fn-y-section}).\;$\operatorname{fn\text{-}y\text{-}section}(\mathit{xs}, f, y) :=$ $(\lambda \operatorname{fsx\_} \in \mathit{xs}.\; f({[}\operatorname{fsx\_}, y]))$
\par\medskip\noindent\textbf{Definition} (\texttt{is-dynkin-system}).\;For all $\omega$ and $\mathit{dd}$, 
\begin{equation*}
\begin{array}{@{}l@{}} \operatorname{is\text{-}dynkin\text{-}system}(\omega, \mathit{dd}) \\ \Leftrightarrow\; \omega \in \mathbf{Set} \wedge \\ \mathit{dd} \subseteq \operatorname{power}(\omega) \wedge \\ \omega \in \mathit{dd} \wedge \\ \forall \operatorname{dya\_}, \operatorname{dyb\_} \in \mathit{dd}. \\ \quad (\operatorname{dya\_} \subseteq \operatorname{dyb\_}) \Rightarrow (\operatorname{complement\text{-}in}(\operatorname{dyb\_}, \operatorname{dya\_}) \in \mathit{dd}) \wedge \\ \forall \operatorname{dyf\_} \in (\mathbb{N} \to \mathit{dd}). \\ \quad \forall \operatorname{dyn\_} \in \mathbb{N}.\; (\operatorname{dyf\_}(\operatorname{dyn\_}) \subseteq \operatorname{dyf\_}(\operatorname{succ}(\operatorname{dyn\_}))) \Rightarrow  \\ \quad \quad \operatorname{big\text{-}union}(\operatorname{dyk\_}, \mathbb{N}, \operatorname{dyf\_}(\operatorname{dyk\_})) \in \mathit{dd} \end{array}
\end{equation*}

\par\medskip\noindent\textbf{Definition} (\texttt{is-pi-system}).\;For every $\mathit{pp}$, 
\begin{equation*}
\operatorname{is\text{-}pi\text{-}system}(\mathit{pp}) \Leftrightarrow \forall \operatorname{pya\_}, \operatorname{pyb\_} \in \mathit{pp}.\; ((\operatorname{pya\_} \cap \operatorname{pyb\_}) \in \mathit{pp})
\end{equation*}

\par\medskip\noindent\textbf{Definition} (\texttt{dynkin-generated}).\;$\operatorname{dynkin\text{-}generated}(\omega, \mathit{ce}) :=$
\begin{equation*}
\{\operatorname{dga\_} \in \operatorname{power}(\omega) : \forall \operatorname{dgb\_}.\; (((\operatorname{dgb\_} \in \operatorname{power}(\operatorname{power}(\omega))) \wedge (\operatorname{is\text{-}dynkin\text{-}system}(\omega, \operatorname{dgb\_}) \wedge (\mathit{ce} \subseteq \operatorname{dgb\_}))) \Rightarrow (\operatorname{dga\_} \in \operatorname{dgb\_}))\}
\end{equation*}

\par\medskip\noindent\textbf{Definition} (\texttt{x-section-measure}).\;$\operatorname{x\text{-}section\text{-}measure}(\mathit{xs}, \mathit{ys}, \nu, \mathit{cs}) :=$ $(\lambda \operatorname{smx\_} \in \mathit{xs}.\; \nu(\operatorname{x\text{-}section}(\mathit{ys}, \mathit{cs}, \operatorname{smx\_})))$
\par\medskip\noindent\textbf{Definition} (\texttt{y-section-measure}).\;$\operatorname{y\text{-}section\text{-}measure}(\mathit{xs}, \mathit{ys}, \mu, \mathit{cs}) :=$ $(\lambda \operatorname{smy\_} \in \mathit{ys}.\; \mu(\operatorname{y\text{-}section}(\mathit{xs}, \mathit{cs}, \operatorname{smy\_})))$
\par\medskip\noindent\textbf{Definition} (\texttt{is-fubini-set}).\;For all $\mathit{xs}$, $\mathit{ca}$, $\mu$, $\mathit{ys}$, $\mathit{cb}$, $\nu$ and $\mathit{cs}$, 
\begin{equation*}
\begin{array}{@{}l@{}} \operatorname{is\text{-}fubini\text{-}set}(\mathit{xs}, \mathit{ca}, \mu, \mathit{ys}, \mathit{cb}, \nu, \mathit{cs}) \\ \Leftrightarrow\; \operatorname{is\text{-}measurable\text{-}fn}(\mathit{xs}, \mathit{ca}, \operatorname{x\text{-}section\text{-}measure}(\mathit{xs}, \mathit{ys}, \nu, \mathit{cs})) \wedge \\ \operatorname{is\text{-}measurable\text{-}fn}(\mathit{ys}, \mathit{cb}, \operatorname{y\text{-}section\text{-}measure}(\mathit{xs}, \mathit{ys}, \mu, \mathit{cs})) \wedge \\ \operatorname{integral}(\mathit{xs}, \mathit{ca}, \mu, \operatorname{x\text{-}section\text{-}measure}(\mathit{xs}, \mathit{ys}, \nu, \mathit{cs})) \\ =\; \operatorname{integral}(\mathit{ys}, \mathit{cb}, \nu, \operatorname{y\text{-}section\text{-}measure}(\mathit{xs}, \mathit{ys}, \mu, \mathit{cs})) \end{array}
\end{equation*}

\par\medskip\noindent\textbf{Definition} (\texttt{x-section-integral}).\;$\operatorname{x\text{-}section\text{-}integral}(\mathit{xs}, \mathit{ys}, \mathit{cb}, \nu, f) :=$
\begin{equation*}
(\lambda \operatorname{six\_} \in \mathit{xs}.\; \operatorname{integral}(\mathit{ys}, \mathit{cb}, \nu, \operatorname{fn\text{-}x\text{-}section}(\mathit{ys}, f, \operatorname{six\_})))
\end{equation*}

\par\medskip\noindent\textbf{Definition} (\texttt{y-section-integral}).\;$\operatorname{y\text{-}section\text{-}integral}(\mathit{xs}, \mathit{ca}, \mu, \mathit{ys}, f) :=$
\begin{equation*}
(\lambda \operatorname{siy\_} \in \mathit{ys}.\; \operatorname{integral}(\mathit{xs}, \mathit{ca}, \mu, \operatorname{fn\text{-}y\text{-}section}(\mathit{xs}, f, \operatorname{siy\_})))
\end{equation*}

\par\medskip\noindent\textbf{Definition} (\texttt{product-measure}).\;$\operatorname{product\text{-}measure}(\mathit{xs}, \mathit{ca}, \mu, \mathit{ys}, \mathit{cb}, \nu) :=$
\begin{equation*}
\begin{array}{@{}l@{}} (\lambda \operatorname{pmc\_} \in \operatorname{product\text{-}sigma}(\mathit{xs}, \mathit{ca}, \mathit{ys}, \mathit{cb}).\; \\ \quad \operatorname{integral}(\mathit{xs}, \mathit{ca}, \mu, \operatorname{x\text{-}section\text{-}measure}(\mathit{xs}, \mathit{ys}, \nu, \operatorname{pmc\_}))) \end{array}
\end{equation*}


\vnbbookfile{theorem-library/product-measure}{theorem-library/product-measure}
\noindent{\small ==================================================================== theorem-library/product-measure.scm -- PRODUCT SIGMA-ALGEBRAS, SECTIONS, DYNKIN'S PI-LAMBDA THEOREM, and FUBINI'S THEOREM for non-negative functions. Agent FUBINI-1, 2026-10-10 (scratchpad/triage/FUBINI-1.md).  Library: measure.}\par
\vnbprop{power-complement-in-member}{\vnbstmt{\forall c, a.; \forall b \in \operatorname{power}(c).\; (\operatorname{complement\text{-}in}(b, a) \in \operatorname{power}(c))}}
\vnbprop{dynkin-generated-unfold}{\vnbstmt{\begin{array}{@{}l@{}} \forall a, e. \\ \quad (\operatorname{dynkin\text{-}generated}(a, e) \simeq \{\mathit{dga} \in \operatorname{power}(a) : \forall b \in \operatorname{power}(\operatorname{power}(a)).\; ((\operatorname{is\text{-}dynkin\text{-}system}(a, b) \wedge (e \subseteq b)) \Rightarrow (\mathit{dga} \in b))\}) \end{array}}}
\vnbprop{dynkin-in-power-power}{\vnbstmt{\forall a, d.\; (\operatorname{is\text{-}dynkin\text{-}system}(a, d) \Rightarrow (d \in \operatorname{power}(\operatorname{power}(a))))}}
\vnbprop{dynkin-generated-membership}{\vnbstmt{\begin{array}{@{}l@{}} \forall a, e, \alpha. \\ \quad \alpha \in \operatorname{dynkin\text{-}generated}(a, e) \\ \quad \Leftrightarrow\; ((\alpha \in \operatorname{power}(a)) \wedge \forall b.\; ((\operatorname{is\text{-}dynkin\text{-}system}(a, b) \wedge (e \subseteq b)) \Rightarrow (\alpha \in b))) \end{array}}}
\vnbprop{dynkin-generated-is-dynkin}{\vnbstmt{\begin{array}{@{}l@{}} \forall a, e. \\ \quad ((a \in \mathbf{Set}) \wedge (e \subseteq \operatorname{power}(a))) \Rightarrow  \\ \quad \quad \operatorname{is\text{-}dynkin\text{-}system}(a, \operatorname{dynkin\text{-}generated}(a, e)) \end{array}}}
\vnbprop{dynkin-inter-section}{\vnbstmt{\begin{array}{@{}l@{}} \forall a, d, s. \\ \quad (\operatorname{is\text{-}dynkin\text{-}system}(a, d) \wedge (s \in d)) \Rightarrow  \\ \quad \quad \operatorname{is\text{-}dynkin\text{-}system}(a, \{e \in \operatorname{power}(a) : ((s \cap e) \in d)\}) \end{array}}}
\vnbprop{dynkin-pi-is-sigma-algebra}{\vnbstmt{\begin{array}{@{}l@{}} \forall a, d. \\ \quad (\operatorname{is\text{-}dynkin\text{-}system}(a, d) \wedge \operatorname{is\text{-}pi\text{-}system}(d)) \Rightarrow \operatorname{is\text{-}sigma\text{-}algebra}(a, d) \end{array}}}
\vnbprop{dynkin-pi-lambda}{\vnbstmt{\begin{array}{@{}l@{}} \forall a, d, p. \\ \quad (\operatorname{is\text{-}dynkin\text{-}system}(a, d) \wedge (\operatorname{is\text{-}pi\text{-}system}(p) \wedge (p \subseteq d))) \Rightarrow  \\ \quad \quad \operatorname{sigma\text{-}generated}(a, p) \subseteq d \end{array}}}
\vnbprop{pair-cartesian-iff}{\vnbstmt{\forall a, b, u, v.\; (({[}u, v] \in (a \times b)) \Leftrightarrow ((u \in a) \wedge (v \in b)))}}
\vnbprop{x-section-unfold}{\vnbstmt{\forall s, a, x.\; (\operatorname{x\text{-}section}(s, a, x) \simeq \{y \in s : ({[}x, y] \in a)\})}}
\vnbprop{y-section-unfold}{\vnbstmt{\forall s, a, y.\; (\operatorname{y\text{-}section}(s, a, y) \simeq \{x \in s : ({[}x, y] \in a)\})}}
\vnbprop{x-section-membership}{\vnbstmt{\forall s, a, x, z.\; ((z \in \operatorname{x\text{-}section}(s, a, x)) \Leftrightarrow ((z \in s) \wedge ({[}x, z] \in a)))}}
\vnbprop{y-section-membership}{\vnbstmt{\forall s, a, y, z.\; ((z \in \operatorname{y\text{-}section}(s, a, y)) \Leftrightarrow ((z \in s) \wedge ({[}z, y] \in a)))}}
\vnbprop{product-sigma-unfold}{\vnbstmt{\begin{array}{@{}l@{}} \forall s, a, b, c. \\ \quad (\operatorname{product\text{-}sigma}(s, a, b, c) \simeq \operatorname{sigma\text{-}generated}((s \times b), \operatorname{rectangles}(s, a, b, c))) \end{array}}}
\vnbprop{rectangles-unfold}{\vnbstmt{\begin{array}{@{}l@{}} \forall s, a, b, c. \\ \quad (\operatorname{rectangles}(s, a, b, c) \simeq \{r \in \operatorname{power}((s \times b)) : \exists \mathit{rca} \in a, \mathit{rcb} \in c.\; (r = (\mathit{rca} \times \mathit{rcb}))\}) \end{array}}}
\vnbprop{rectangles-membership}{\vnbstmt{\begin{array}{@{}l@{}} \forall s, a, b, c, r. \\ \quad r \in \operatorname{rectangles}(s, a, b, c) \\ \quad \Leftrightarrow\; ((r \in \operatorname{power}((s \times b))) \wedge \exists \mathit{rca} \in a, \mathit{rcb} \in c.\; (r = (\mathit{rca} \times \mathit{rcb}))) \end{array}}}
\vnbprop{cartesian-in-power-product}{\vnbstmt{\begin{array}{@{}l@{}} \forall s, a, b, c. \\ \quad ((b \in \operatorname{power}(s)) \wedge (c \in \operatorname{power}(a))) \Rightarrow  \\ \quad \quad (b \times c) \in \operatorname{power}((s \times a)) \end{array}}}
\vnbprop{cartesian-intersection}{\vnbstmt{\begin{array}{@{}l@{}} \forall a, b, c, d. \\ \quad ((a \times b) \cap (c \times d)) = ((a \cap c) \times (b \cap d)) \end{array}}}
\vnbprop{rectangle-in-rectangles}{\vnbstmt{\begin{array}{@{}l@{}} \forall s, a, b, c, d, f. \\ \quad (\operatorname{is\text{-}sigma\text{-}algebra}(s, a) \wedge (\operatorname{is\text{-}sigma\text{-}algebra}(b, c) \wedge ((d \in a) \wedge (f \in c)))) \Rightarrow  \\ \quad \quad (d \times f) \in \operatorname{rectangles}(s, a, b, c) \end{array}}}
\vnbprop{rectangles-pi-system}{\vnbstmt{\begin{array}{@{}l@{}} \forall s, a, b, c. \\ \quad (\operatorname{is\text{-}sigma\text{-}algebra}(s, a) \wedge \operatorname{is\text{-}sigma\text{-}algebra}(b, c)) \Rightarrow  \\ \quad \quad \operatorname{is\text{-}pi\text{-}system}(\operatorname{rectangles}(s, a, b, c)) \end{array}}}
\vnbprop{product-sigma-is-sigma-algebra}{\vnbstmt{\begin{array}{@{}l@{}} \forall s, a, b, c. \\ \quad (\operatorname{is\text{-}sigma\text{-}algebra}(s, a) \wedge \operatorname{is\text{-}sigma\text{-}algebra}(b, c)) \Rightarrow  \\ \quad \quad \operatorname{is\text{-}sigma\text{-}algebra}((s \times b), \operatorname{product\text{-}sigma}(s, a, b, c)) \end{array}}}
\vnbprop{rectangle-in-product-sigma}{\vnbstmt{\begin{array}{@{}l@{}} \forall s, a, b, c, d, f. \\ \quad (\operatorname{is\text{-}sigma\text{-}algebra}(s, a) \wedge (\operatorname{is\text{-}sigma\text{-}algebra}(b, c) \wedge ((d \in a) \wedge (f \in c)))) \Rightarrow  \\ \quad \quad (d \times f) \in \operatorname{product\text{-}sigma}(s, a, b, c) \end{array}}}
\vnbprop{x-section-measurable}{\vnbstmt{\begin{array}{@{}l@{}} \forall s, a, b, c, d, x. \\ \quad (\operatorname{is\text{-}sigma\text{-}algebra}(s, a) \wedge (\operatorname{is\text{-}sigma\text{-}algebra}(b, c) \wedge ((d \in \operatorname{product\text{-}sigma}(s, a, b, c)) \wedge (x \in s)))) \Rightarrow  \\ \quad \quad \operatorname{x\text{-}section}(b, d, x) \in c \end{array}}}
\vnbprop{y-section-measurable}{\vnbstmt{\begin{array}{@{}l@{}} \forall s, a, b, c, d, y. \\ \quad (\operatorname{is\text{-}sigma\text{-}algebra}(s, a) \wedge (\operatorname{is\text{-}sigma\text{-}algebra}(b, c) \wedge ((d \in \operatorname{product\text{-}sigma}(s, a, b, c)) \wedge (y \in b)))) \Rightarrow  \\ \quad \quad \operatorname{y\text{-}section}(s, d, y) \in a \end{array}}}
\vnbprop{fn-x-section-apply}{\vnbstmt{\forall s, f, x.; \forall z \in s.\; (\operatorname{fn\text{-}x\text{-}section}(s, f, x)(z) \simeq f({[}x, z]))}}
\vnbprop{fn-y-section-apply}{\vnbstmt{\forall s, f, y.; \forall z \in s.\; (\operatorname{fn\text{-}y\text{-}section}(s, f, y)(z) \simeq f({[}z, y]))}}
\vnbprop{fn-x-section-measurable}{\vnbstmt{\begin{array}{@{}l@{}} \forall s, a, b, c, f, x. \\ \quad \operatorname{is\text{-}sigma\text{-}algebra}(s, a) \wedge \\ \quad \operatorname{is\text{-}sigma\text{-}algebra}(b, c) \wedge \\ \quad \operatorname{is\text{-}measurable\text{-}fn}((s \times b), \operatorname{product\text{-}sigma}(s, a, b, c), f) \wedge \\ \quad x \in s \Rightarrow \\ \quad \quad \operatorname{is\text{-}measurable\text{-}fn}(b, c, \operatorname{fn\text{-}x\text{-}section}(b, f, x)) \end{array}}}
\vnbprop{fn-y-section-measurable}{\vnbstmt{\begin{array}{@{}l@{}} \forall s, a, b, c, f, y. \\ \quad \operatorname{is\text{-}sigma\text{-}algebra}(s, a) \wedge \\ \quad \operatorname{is\text{-}sigma\text{-}algebra}(b, c) \wedge \\ \quad \operatorname{is\text{-}measurable\text{-}fn}((s \times b), \operatorname{product\text{-}sigma}(s, a, b, c), f) \wedge \\ \quad y \in b \Rightarrow \\ \quad \quad \operatorname{is\text{-}measurable\text{-}fn}(s, a, \operatorname{fn\text{-}y\text{-}section}(s, f, y)) \end{array}}}
\vnbprop{x-section-rectangle-in}{\vnbstmt{\begin{array}{@{}l@{}} \forall s, a, b, x. \\ \quad ((b \subseteq s) \wedge (x \in a)) \Rightarrow (\operatorname{x\text{-}section}(s, (a \times b), x) = b) \end{array}}}
\vnbprop{x-section-rectangle-out}{\vnbstmt{\begin{array}{@{}l@{}} \forall s, a, b, x. \\ \quad \neg (x \in a) \Rightarrow (\operatorname{x\text{-}section}(s, (a \times b), x) = \emptyset) \end{array}}}
\vnbprop{y-section-rectangle-in}{\vnbstmt{\begin{array}{@{}l@{}} \forall s, a, b, y. \\ \quad ((a \subseteq s) \wedge (y \in b)) \Rightarrow (\operatorname{y\text{-}section}(s, (a \times b), y) = a) \end{array}}}
\vnbprop{y-section-rectangle-out}{\vnbstmt{\begin{array}{@{}l@{}} \forall s, a, b, y. \\ \quad \neg (y \in b) \Rightarrow (\operatorname{y\text{-}section}(s, (a \times b), y) = \emptyset) \end{array}}}
\vnbprop{integral-congruence-ptwise}{\vnbstmt{\begin{array}{@{}l@{}} \forall a, b, u, f, g. \\ \quad ((f \in (a \to \operatorname{rr\text{-}pos\text{-}star})) \wedge ((g \in (a \to \operatorname{rr\text{-}pos\text{-}star})) \wedge \forall x \in a.\; (f(x) = g(x)))) \Rightarrow  \\ \quad \quad (\operatorname{integral}(a, b, u, f) \simeq \operatorname{integral}(a, b, u, g)) \end{array}}}
\vnbprop{measurable-fn-congruence-ptwise}{\vnbstmt{\begin{array}{@{}l@{}} \forall a, b, f, g. \\ \quad (\operatorname{is\text{-}measurable\text{-}fn}(a, b, f) \wedge ((g \in (a \to \operatorname{rr\text{-}pos\text{-}star})) \wedge \forall x \in a.\; (f(x) = g(x)))) \Rightarrow  \\ \quad \quad \operatorname{is\text{-}measurable\text{-}fn}(a, b, g) \end{array}}}
\vnbprop{x-section-measure-apply}{\vnbstmt{\begin{array}{@{}l@{}} \forall s, a, u, b.; \forall x \in s. \\ \quad (\operatorname{x\text{-}section\text{-}measure}(s, a, u, b)(x) \simeq u(\operatorname{x\text{-}section}(a, b, x))) \end{array}}}
\vnbprop{y-section-measure-apply}{\vnbstmt{\begin{array}{@{}l@{}} \forall s, a, u, b.; \forall y \in a. \\ \quad (\operatorname{y\text{-}section\text{-}measure}(s, a, u, b)(y) \simeq u(\operatorname{y\text{-}section}(s, b, y))) \end{array}}}
\vnbprop{x-section-measure-type}{\vnbstmt{\begin{array}{@{}l@{}} \forall s, a, u, b, c, d, f. \\ \quad (\operatorname{is\text{-}measure}(s, a, u) \wedge (\operatorname{is\text{-}measure}(b, c, d) \wedge (f \in \operatorname{product\text{-}sigma}(s, a, b, c)))) \Rightarrow  \\ \quad \quad \operatorname{x\text{-}section\text{-}measure}(s, b, d, f) \in (s \to \operatorname{rr\text{-}pos\text{-}star}) \end{array}}}
\vnbprop{y-section-measure-type}{\vnbstmt{\begin{array}{@{}l@{}} \forall s, a, u, b, c, d, f. \\ \quad (\operatorname{is\text{-}measure}(s, a, u) \wedge (\operatorname{is\text{-}measure}(b, c, d) \wedge (f \in \operatorname{product\text{-}sigma}(s, a, b, c)))) \Rightarrow  \\ \quad \quad \operatorname{y\text{-}section\text{-}measure}(s, b, u, f) \in (b \to \operatorname{rr\text{-}pos\text{-}star}) \end{array}}}
\vnbprop{x-section-measure-rectangle}{\vnbstmt{\begin{array}{@{}l@{}} \forall s, a, u, b, c, d, f, g. \\ \quad (\operatorname{is\text{-}measure}(s, a, u) \wedge (\operatorname{is\text{-}measure}(b, c, d) \wedge ((f \in a) \wedge ((g \in c) \wedge (d(g) \in \mathbb{R}))))) \Rightarrow  \\ \quad \quad \operatorname{is\text{-}measurable\text{-}fn}(s, a, \operatorname{x\text{-}section\text{-}measure}(s, b, d, (f \times g))) \wedge \\ \quad \quad \operatorname{integral}(s, a, u, \operatorname{x\text{-}section\text{-}measure}(s, b, d, (f \times g))) = \operatorname{etimes}(d(g), u(f)) \end{array}}}
\vnbprop{y-section-measure-rectangle}{\vnbstmt{\begin{array}{@{}l@{}} \forall s, a, u, b, c, d, f, g. \\ \quad (\operatorname{is\text{-}measure}(s, a, u) \wedge (\operatorname{is\text{-}measure}(b, c, d) \wedge ((f \in a) \wedge ((g \in c) \wedge (u(f) \in \mathbb{R}))))) \Rightarrow  \\ \quad \quad \operatorname{is\text{-}measurable\text{-}fn}(b, c, \operatorname{y\text{-}section\text{-}measure}(s, b, u, (f \times g))) \wedge \\ \quad \quad \operatorname{integral}(b, c, d, \operatorname{y\text{-}section\text{-}measure}(s, b, u, (f \times g))) = \operatorname{etimes}(u(f), d(g)) \end{array}}}
\vnbprop{fubini-set-rectangle}{\vnbstmt{\begin{array}{@{}l@{}} \forall s, a, u, b, c, d, f, g. \\ \quad (\operatorname{is\text{-}measure}(s, a, u) \wedge (\operatorname{is\text{-}measure}(b, c, d) \wedge ((f \in a) \wedge ((g \in c) \wedge ((u(f) \in \mathbb{R}) \wedge (d(g) \in \mathbb{R})))))) \Rightarrow  \\ \quad \quad \operatorname{is\text{-}fubini\text{-}set}(s, a, u, b, c, d, (f \times g)) \end{array}}}
\vnbprop{x-section-mono}{\vnbstmt{\begin{array}{@{}l@{}} \forall s, a, b, x. \\ \quad (a \subseteq b) \Rightarrow (\operatorname{x\text{-}section}(s, a, x) \subseteq \operatorname{x\text{-}section}(s, b, x)) \end{array}}}
\vnbprop{y-section-mono}{\vnbstmt{\begin{array}{@{}l@{}} \forall s, a, b, y. \\ \quad (a \subseteq b) \Rightarrow (\operatorname{y\text{-}section}(s, a, y) \subseteq \operatorname{y\text{-}section}(s, b, y)) \end{array}}}
\vnbprop{x-section-difference}{\vnbstmt{\begin{array}{@{}l@{}} \forall s, a, b, x. \\ \quad \operatorname{x\text{-}section}(s, \operatorname{complement\text{-}in}(b, a), x) = \operatorname{complement\text{-}in}(\operatorname{x\text{-}section}(s, b, x), \operatorname{x\text{-}section}(s, a, x)) \end{array}}}
\vnbprop{y-section-difference}{\vnbstmt{\begin{array}{@{}l@{}} \forall s, a, b, y. \\ \quad \operatorname{y\text{-}section}(s, \operatorname{complement\text{-}in}(b, a), y) = \operatorname{complement\text{-}in}(\operatorname{y\text{-}section}(s, b, y), \operatorname{y\text{-}section}(s, a, y)) \end{array}}}
\vnbprop{x-section-measure-difference}{\vnbstmt{\begin{array}{@{}l@{}} \forall s, a, u, b, c, d, f, g. \\ \quad \operatorname{is\text{-}measure}(s, a, u) \wedge \\ \quad \operatorname{is\text{-}measure}(b, c, d) \wedge \\ \quad f \in \operatorname{product\text{-}sigma}(s, a, b, c) \wedge \\ \quad g \in \operatorname{product\text{-}sigma}(s, a, b, c) \wedge \\ \quad f \subseteq g \wedge \\ \quad \operatorname{is\text{-}measurable\text{-}fn}(s, a, \operatorname{x\text{-}section\text{-}measure}(s, b, d, f)) \wedge \\ \quad \operatorname{is\text{-}measurable\text{-}fn}(s, a, \operatorname{x\text{-}section\text{-}measure}(s, b, d, g)) \wedge \\ \quad \forall z \in s.\; (d(\operatorname{x\text{-}section}(b, g, z)) \in \mathbb{R}) \wedge \\ \quad \operatorname{integral}(s, a, u, \operatorname{x\text{-}section\text{-}measure}(s, b, d, g)) \in \mathbb{R} \Rightarrow \\ \quad \quad \operatorname{is\text{-}measurable\text{-}fn}(s, a, \operatorname{x\text{-}section\text{-}measure}(s, b, d, \operatorname{complement\text{-}in}(g, f))) \wedge \\ \quad \quad \operatorname{integral}(s, a, u, \operatorname{x\text{-}section\text{-}measure}(s, b, d, f)) \in \mathbb{R} \wedge \\ \quad \quad \operatorname{integral}(s, a, u, \operatorname{x\text{-}section\text{-}measure}(s, b, d, \operatorname{complement\text{-}in}(g, f))) \\ \quad \quad =\; (\operatorname{integral}(s, a, u, \operatorname{x\text{-}section\text{-}measure}(s, b, d, g)) - \operatorname{integral}(s, a, u, \operatorname{x\text{-}section\text{-}measure}(s, b, d, f))) \end{array}}}
\vnbprop{y-section-measure-difference}{\vnbstmt{\begin{array}{@{}l@{}} \forall s, a, u, b, c, d, f, g. \\ \quad \operatorname{is\text{-}measure}(s, a, u) \wedge \\ \quad \operatorname{is\text{-}measure}(b, c, d) \wedge \\ \quad f \in \operatorname{product\text{-}sigma}(s, a, b, c) \wedge \\ \quad g \in \operatorname{product\text{-}sigma}(s, a, b, c) \wedge \\ \quad f \subseteq g \wedge \\ \quad \operatorname{is\text{-}measurable\text{-}fn}(b, c, \operatorname{y\text{-}section\text{-}measure}(s, b, u, f)) \wedge \\ \quad \operatorname{is\text{-}measurable\text{-}fn}(b, c, \operatorname{y\text{-}section\text{-}measure}(s, b, u, g)) \wedge \\ \quad \forall z \in b.\; (u(\operatorname{y\text{-}section}(s, g, z)) \in \mathbb{R}) \wedge \\ \quad \operatorname{integral}(b, c, d, \operatorname{y\text{-}section\text{-}measure}(s, b, u, g)) \in \mathbb{R} \Rightarrow \\ \quad \quad \operatorname{is\text{-}measurable\text{-}fn}(b, c, \operatorname{y\text{-}section\text{-}measure}(s, b, u, \operatorname{complement\text{-}in}(g, f))) \wedge \\ \quad \quad \operatorname{integral}(b, c, d, \operatorname{y\text{-}section\text{-}measure}(s, b, u, f)) \in \mathbb{R} \wedge \\ \quad \quad \operatorname{integral}(b, c, d, \operatorname{y\text{-}section\text{-}measure}(s, b, u, \operatorname{complement\text{-}in}(g, f))) \\ \quad \quad =\; (\operatorname{integral}(b, c, d, \operatorname{y\text{-}section\text{-}measure}(s, b, u, g)) - \operatorname{integral}(b, c, d, \operatorname{y\text{-}section\text{-}measure}(s, b, u, f))) \end{array}}}
\vnbprop{fubini-set-difference}{\vnbstmt{\begin{array}{@{}l@{}} \forall s, a, u, b, c, d, f, g. \\ \quad \operatorname{is\text{-}measure}(s, a, u) \wedge \\ \quad \operatorname{is\text{-}measure}(b, c, d) \wedge \\ \quad f \in \operatorname{product\text{-}sigma}(s, a, b, c) \wedge \\ \quad g \in \operatorname{product\text{-}sigma}(s, a, b, c) \wedge \\ \quad f \subseteq g \wedge \\ \quad \operatorname{is\text{-}fubini\text{-}set}(s, a, u, b, c, d, f) \wedge \\ \quad \operatorname{is\text{-}fubini\text{-}set}(s, a, u, b, c, d, g) \wedge \\ \quad \forall z \in s.\; (d(\operatorname{x\text{-}section}(b, g, z)) \in \mathbb{R}) \wedge \\ \quad \forall z \in b.\; (u(\operatorname{y\text{-}section}(s, g, z)) \in \mathbb{R}) \wedge \\ \quad \operatorname{integral}(s, a, u, \operatorname{x\text{-}section\text{-}measure}(s, b, d, g)) \in \mathbb{R} \Rightarrow \\ \quad \quad \operatorname{is\text{-}fubini\text{-}set}(s, a, u, b, c, d, \operatorname{complement\text{-}in}(g, f)) \end{array}}}
\vnbprop{x-section-measure-union}{\vnbstmt{\begin{array}{@{}l@{}} \forall s, a, u, b, c, d, f. \\ \quad \operatorname{is\text{-}measure}(s, a, u) \wedge \\ \quad \operatorname{is\text{-}measure}(b, c, d) \wedge \\ \quad f \in (\mathbb{N} \to \operatorname{product\text{-}sigma}(s, a, b, c)) \wedge \\ \quad \forall n \in \mathbb{N}.\; (f(n) \subseteq f((n + 1))) \wedge \\ \quad \forall n \in \mathbb{N}.\; \operatorname{is\text{-}measurable\text{-}fn}(s, a, \operatorname{x\text{-}section\text{-}measure}(s, b, d, f(n))) \Rightarrow \\ \quad \quad \operatorname{is\text{-}measurable\text{-}fn}(s, a, \operatorname{x\text{-}section\text{-}measure}(s, b, d, \operatorname{big\text{-}union}(g, \mathbb{N}, f(g)))) \wedge \\ \quad \quad \operatorname{integral}(s, a, u, \operatorname{x\text{-}section\text{-}measure}(s, b, d, \operatorname{big\text{-}union}(g, \mathbb{N}, f(g)))) \\ \quad \quad =\; \operatorname{esup}(\operatorname{image}((\lambda n \in \mathbb{N}.\; \operatorname{integral}(s, a, u, \operatorname{x\text{-}section\text{-}measure}(s, b, d, f(n)))), \mathbb{N})) \end{array}}}
\vnbprop{y-section-measure-union}{\vnbstmt{\begin{array}{@{}l@{}} \forall s, a, u, b, c, d, f. \\ \quad \operatorname{is\text{-}measure}(s, a, u) \wedge \\ \quad \operatorname{is\text{-}measure}(b, c, d) \wedge \\ \quad f \in (\mathbb{N} \to \operatorname{product\text{-}sigma}(s, a, b, c)) \wedge \\ \quad \forall n \in \mathbb{N}.\; (f(n) \subseteq f((n + 1))) \wedge \\ \quad \forall n \in \mathbb{N}.\; \operatorname{is\text{-}measurable\text{-}fn}(b, c, \operatorname{y\text{-}section\text{-}measure}(s, b, u, f(n))) \Rightarrow \\ \quad \quad \operatorname{is\text{-}measurable\text{-}fn}(b, c, \operatorname{y\text{-}section\text{-}measure}(s, b, u, \operatorname{big\text{-}union}(g, \mathbb{N}, f(g)))) \wedge \\ \quad \quad \operatorname{integral}(b, c, d, \operatorname{y\text{-}section\text{-}measure}(s, b, u, \operatorname{big\text{-}union}(g, \mathbb{N}, f(g)))) \\ \quad \quad =\; \operatorname{esup}(\operatorname{image}((\lambda n \in \mathbb{N}.\; \operatorname{integral}(b, c, d, \operatorname{y\text{-}section\text{-}measure}(s, b, u, f(n)))), \mathbb{N})) \end{array}}}
\vnbprop{fubini-set-union}{\vnbstmt{\begin{array}{@{}l@{}} \forall s, a, u, b, c, d, f. \\ \quad \operatorname{is\text{-}measure}(s, a, u) \wedge \\ \quad \operatorname{is\text{-}measure}(b, c, d) \wedge \\ \quad f \in (\mathbb{N} \to \operatorname{product\text{-}sigma}(s, a, b, c)) \wedge \\ \quad \forall n \in \mathbb{N}.\; (f(n) \subseteq f((n + 1))) \wedge \\ \quad \forall n \in \mathbb{N}.\; \operatorname{is\text{-}fubini\text{-}set}(s, a, u, b, c, d, f(n)) \Rightarrow \\ \quad \quad \operatorname{is\text{-}fubini\text{-}set}(s, a, u, b, c, d, \operatorname{big\text{-}union}(g, \mathbb{N}, f(g))) \end{array}}}
\vnbprop{x-section-in-cut}{\vnbstmt{\begin{array}{@{}l@{}} \forall s, a, b, c, x. \\ \quad \operatorname{x\text{-}section}(s, (a \cap (b \times c)), x) \subseteq c \end{array}}}
\vnbprop{y-section-in-cut}{\vnbstmt{\begin{array}{@{}l@{}} \forall s, a, b, c, y. \\ \quad \operatorname{y\text{-}section}(s, (a \cap (b \times c)), y) \subseteq b \end{array}}}
\vnbprop{fubini-set-rectangle-cut}{\vnbstmt{\begin{array}{@{}l@{}} \forall s, a, u, b, c, d, f, g, h, k. \\ \quad \operatorname{is\text{-}measure}(s, a, u) \wedge \\ \quad \operatorname{is\text{-}measure}(b, c, d) \wedge \\ \quad f \in a \wedge \\ \quad g \in c \wedge \\ \quad u(f) \in \mathbb{R} \wedge \\ \quad d(g) \in \mathbb{R} \wedge \\ \quad h \in a \wedge \\ \quad k \in c \Rightarrow \\ \quad \quad \operatorname{is\text{-}fubini\text{-}set}(s, a, u, b, c, d, ((h \times k) \cap (f \times g))) \end{array}}}
\vnbprop{fubini-set-cut-bounds}{\vnbstmt{\begin{array}{@{}l@{}} \forall s, a, u, b, c, d, f, g, h. \\ \quad \operatorname{is\text{-}measure}(s, a, u) \wedge \\ \quad \operatorname{is\text{-}measure}(b, c, d) \wedge \\ \quad f \in a \wedge \\ \quad g \in c \wedge \\ \quad u(f) \in \mathbb{R} \wedge \\ \quad d(g) \in \mathbb{R} \wedge \\ \quad h \in \operatorname{product\text{-}sigma}(s, a, b, c) \wedge \\ \quad \operatorname{is\text{-}fubini\text{-}set}(s, a, u, b, c, d, (h \cap (f \times g))) \Rightarrow \\ \quad \quad \forall z \in s.\; (d(\operatorname{x\text{-}section}(b, (h \cap (f \times g)), z)) \in \mathbb{R}) \wedge \\ \quad \quad \forall z \in b.\; (u(\operatorname{y\text{-}section}(s, (h \cap (f \times g)), z)) \in \mathbb{R}) \wedge \\ \quad \quad \operatorname{integral}(s, a, u, \operatorname{x\text{-}section\text{-}measure}(s, b, d, (h \cap (f \times g)))) \in \mathbb{R} \end{array}}}
\vnbprop{fubini-set-in-finite-rectangle}{\vnbstmt{\begin{array}{@{}l@{}} \forall s, a, u, b, c, d, f, g, h. \\ \quad \operatorname{is\text{-}measure}(s, a, u) \wedge \\ \quad \operatorname{is\text{-}measure}(b, c, d) \wedge \\ \quad f \in a \wedge \\ \quad g \in c \wedge \\ \quad u(f) \in \mathbb{R} \wedge \\ \quad d(g) \in \mathbb{R} \wedge \\ \quad h \in \operatorname{product\text{-}sigma}(s, a, b, c) \Rightarrow \\ \quad \quad \operatorname{is\text{-}fubini\text{-}set}(s, a, u, b, c, d, (h \cap (f \times g))) \end{array}}}
\vnbprop{measure-union-finite}{\vnbstmt{\begin{array}{@{}l@{}} \forall c, d, u, a, b. \\ \quad (\operatorname{is\text{-}measure}(c, d, u) \wedge ((a \in d) \wedge ((b \in d) \wedge ((u(a) \in \mathbb{R}) \wedge (u(b) \in \mathbb{R}))))) \Rightarrow  \\ \quad \quad u((a \cup b)) \in \mathbb{R} \end{array}}}
\vnbprop{measure-seg-union-finite}{\vnbstmt{\begin{array}{@{}l@{}} \forall k \in \mathbb{N}, a, b, u, f. \\ \quad (\operatorname{is\text{-}measure}(a, b, u) \wedge ((f \in (\mathbb{N} \to b)) \wedge \forall n \in \mathbb{N}.\; (u(f(n)) \in \mathbb{R}))) \Rightarrow  \\ \quad \quad u(\operatorname{big\text{-}union}(j, \mathrm{OS}(k), f(j))) \in \mathbb{R} \end{array}}}
\vnbprop{sigma-finite-increasing}{\vnbstmt{\begin{array}{@{}l@{}} \forall a, b, u. \\ \quad (\operatorname{is\text{-}measure}(a, b, u) \wedge \operatorname{is\text{-}sigma\text{-}finite}(a, b, u)) \Rightarrow  \\ \quad \quad \exists g \in (\mathbb{N} \to b). \\ \quad \quad \quad \forall n \in \mathbb{N}.\; (u(g(n)) \in \mathbb{R}) \wedge \\ \quad \quad \quad \forall n \in \mathbb{N}.\; (g(n) \subseteq g((n + 1))) \wedge \\ \quad \quad \quad \operatorname{big\text{-}union}(k, \mathbb{N}, g(k)) = a \end{array}}}
\vnbprop{cartesian-mono}{\vnbstmt{\begin{array}{@{}l@{}} \forall a, b, c, d. \\ \quad ((a \subseteq c) \wedge (b \subseteq d)) \Rightarrow ((a \times b) \subseteq (c \times d)) \end{array}}}
\vnbprop{fubini-sets}{\vnbstmt{\begin{array}{@{}l@{}} \forall s, a, u, b, c, d, f. \\ \quad \operatorname{is\text{-}measure}(s, a, u) \wedge \\ \quad \operatorname{is\text{-}measure}(b, c, d) \wedge \\ \quad \operatorname{is\text{-}sigma\text{-}finite}(s, a, u) \wedge \\ \quad \operatorname{is\text{-}sigma\text{-}finite}(b, c, d) \wedge \\ \quad f \in \operatorname{product\text{-}sigma}(s, a, b, c) \Rightarrow \\ \quad \quad \operatorname{is\text{-}fubini\text{-}set}(s, a, u, b, c, d, f) \end{array}}}
\vnbprop{integral-finsum}{\vnbstmt{\begin{array}{@{}l@{}} \forall a, b, u, d, s, m. \\ \quad \operatorname{is\text{-}measure}(a, b, u) \wedge \\ \quad s \in \mathbf{Set} \wedge \\ \quad |s| \in \mathbb{N} \wedge \\ \quad s \subseteq d \wedge \\ \quad \forall v \in s.\; \operatorname{is\text{-}measurable\text{-}fn}(a, b, m(v)) \Rightarrow \\ \quad \quad \operatorname{is\text{-}measurable\text{-}fn}(a, b, \\ \quad \quad \quad (\lambda x \in a.\; \operatorname{finsum}(\operatorname{rr\text{-}pos\text{-}star\text{-}add\text{-}monoid}, (\lambda v \in d.\; m(v)(x)), s))) \wedge \\ \quad \quad \operatorname{integral}(a, b, u, (\lambda x \in a.\; \operatorname{finsum}(\operatorname{rr\text{-}pos\text{-}star\text{-}add\text{-}monoid}, (\lambda v \in d.\; m(v)(x)), s))) \\ \quad \quad =\; \operatorname{finsum}(\operatorname{rr\text{-}pos\text{-}star\text{-}add\text{-}monoid}, (\lambda v \in d.\; \operatorname{integral}(a, b, u, m(v))), s) \end{array}}}
\vnbprop{x-section-simple-integral}{\vnbstmt{\begin{array}{@{}l@{}} \forall a, b, c, d, u, s, x. \\ \quad \operatorname{is\text{-}sigma\text{-}algebra}(a, b) \wedge \\ \quad \operatorname{is\text{-}measure}(c, d, u) \wedge \\ \quad \operatorname{is\text{-}simple\text{-}fn}((a \times c), \operatorname{product\text{-}sigma}(a, b, c, d), s) \wedge \\ \quad x \in a \Rightarrow \\ \quad \quad \operatorname{integral}(c, d, u, \operatorname{fn\text{-}x\text{-}section}(c, s, x)) \\ \quad \quad =\; \operatorname{finsum}(\operatorname{rr\text{-}pos\text{-}star\text{-}add\text{-}monoid}, \\ \quad \quad \quad (\lambda v \in \operatorname{image}(s, (a \times c)).\; \\ \quad \quad \quad \quad \operatorname{etimes}(v, u(\operatorname{x\text{-}section}(c, \{\mathit{fx} \in (a \times c) : (s(\mathit{fx}) = v)\}, x)))), \\ \quad \quad \quad \operatorname{image}(s, (a \times c))) \end{array}}}
\vnbprop{y-section-simple-integral}{\vnbstmt{\begin{array}{@{}l@{}} \forall a, b, u, c, d, s, y. \\ \quad \operatorname{is\text{-}measure}(a, b, u) \wedge \\ \quad \operatorname{is\text{-}sigma\text{-}algebra}(c, d) \wedge \\ \quad \operatorname{is\text{-}simple\text{-}fn}((a \times c), \operatorname{product\text{-}sigma}(a, b, c, d), s) \wedge \\ \quad y \in c \Rightarrow \\ \quad \quad \operatorname{integral}(a, b, u, \operatorname{fn\text{-}y\text{-}section}(a, s, y)) \\ \quad \quad =\; \operatorname{finsum}(\operatorname{rr\text{-}pos\text{-}star\text{-}add\text{-}monoid}, \\ \quad \quad \quad (\lambda v \in \operatorname{image}(s, (a \times c)).\; \\ \quad \quad \quad \quad \operatorname{etimes}(v, u(\operatorname{y\text{-}section}(a, \{x \in (a \times c) : (s(x) = v)\}, y)))), \\ \quad \quad \quad \operatorname{image}(s, (a \times c))) \end{array}}}
\vnbprop{x-section-integral-apply}{\vnbstmt{\begin{array}{@{}l@{}} \forall s, a, b, u, f.; \forall x \in s. \\ \quad (\operatorname{x\text{-}section\text{-}integral}(s, a, b, u, f)(x) \simeq \operatorname{integral}(a, b, u, \operatorname{fn\text{-}x\text{-}section}(a, f, x))) \end{array}}}
\vnbprop{y-section-integral-apply}{\vnbstmt{\begin{array}{@{}l@{}} \forall s, a, u, b, f.; \forall y \in b. \\ \quad (\operatorname{y\text{-}section\text{-}integral}(s, a, u, b, f)(y) \simeq \operatorname{integral}(s, a, u, \operatorname{fn\text{-}y\text{-}section}(s, f, y))) \end{array}}}
\vnbprop{x-section-integral-type}{\vnbstmt{\begin{array}{@{}l@{}} \forall s, a, b, c, u, f. \\ \quad \operatorname{is\text{-}sigma\text{-}algebra}(s, a) \wedge \\ \quad \operatorname{is\text{-}sigma\text{-}algebra}(b, c) \wedge \\ \quad \operatorname{is\text{-}measure}(b, c, u) \wedge \\ \quad \operatorname{is\text{-}measurable\text{-}fn}((s \times b), \operatorname{product\text{-}sigma}(s, a, b, c), f) \Rightarrow \\ \quad \quad \operatorname{x\text{-}section\text{-}integral}(s, b, c, u, f) \in (s \to \operatorname{rr\text{-}pos\text{-}star}) \end{array}}}
\vnbprop{y-section-integral-type}{\vnbstmt{\begin{array}{@{}l@{}} \forall s, a, u, b, c, f. \\ \quad \operatorname{is\text{-}sigma\text{-}algebra}(s, a) \wedge \\ \quad \operatorname{is\text{-}sigma\text{-}algebra}(b, c) \wedge \\ \quad \operatorname{is\text{-}measure}(s, a, u) \wedge \\ \quad \operatorname{is\text{-}measurable\text{-}fn}((s \times b), \operatorname{product\text{-}sigma}(s, a, b, c), f) \Rightarrow \\ \quad \quad \operatorname{y\text{-}section\text{-}integral}(s, a, u, b, f) \in (b \to \operatorname{rr\text{-}pos\text{-}star}) \end{array}}}
\vnbprop{fubini-simple}{\vnbstmt{\begin{array}{@{}l@{}} \forall a, b, u, c, d, f, s. \\ \quad \operatorname{is\text{-}measure}(a, b, u) \wedge \\ \quad \operatorname{is\text{-}measure}(c, d, f) \wedge \\ \quad \operatorname{is\text{-}sigma\text{-}finite}(a, b, u) \wedge \\ \quad \operatorname{is\text{-}sigma\text{-}finite}(c, d, f) \wedge \\ \quad \operatorname{is\text{-}simple\text{-}fn}((a \times c), \operatorname{product\text{-}sigma}(a, b, c, d), s) \Rightarrow \\ \quad \quad \operatorname{is\text{-}measurable\text{-}fn}(a, b, \operatorname{x\text{-}section\text{-}integral}(a, c, d, f, s)) \wedge \\ \quad \quad \operatorname{is\text{-}measurable\text{-}fn}(c, d, \operatorname{y\text{-}section\text{-}integral}(a, b, u, c, s)) \wedge \\ \quad \quad \operatorname{integral}(a, b, u, \operatorname{x\text{-}section\text{-}integral}(a, c, d, f, s)) \\ \quad \quad =\; \operatorname{integral}(c, d, f, \operatorname{y\text{-}section\text{-}integral}(a, b, u, c, s)) \end{array}}}
\vnbprop{ptwise-sup-apply}{\vnbstmt{\begin{array}{@{}l@{}} \forall a, f.; \forall x \in a. \\ \quad (\operatorname{ptwise\text{-}sup}(a, f)(x) \simeq \operatorname{esup}(\operatorname{image}((\lambda n \in \mathbb{N}.\; f(n)(x)), \mathbb{N}))) \end{array}}}
\vnbprop{x-section-integral-limit}{\vnbstmt{\begin{array}{@{}l@{}} \forall s, a, u, b, c, d, f, g. \\ \quad \operatorname{is\text{-}measure}(s, a, u) \wedge \\ \quad \operatorname{is\text{-}measure}(b, c, d) \wedge \\ \quad \operatorname{is\text{-}measurable\text{-}fn}((s \times b), \operatorname{product\text{-}sigma}(s, a, b, c), f) \wedge \\ \quad g \in (\mathbb{N} \to ((s \times b) \to \operatorname{rr\text{-}pos\text{-}star})) \wedge \\ \quad \forall n \in \mathbb{N}. \\ \quad \quad \operatorname{is\text{-}measurable\text{-}fn}((s \times b), \operatorname{product\text{-}sigma}(s, a, b, c), g(n)) \wedge \\ \quad \forall n \in \mathbb{N}.\; \operatorname{ptwise\text{-}le}((s \times b), g(n), g((n + 1))) \wedge \\ \quad \forall w \in (s \times b).\; (\operatorname{ptwise\text{-}sup}((s \times b), g)(w) = f(w)) \wedge \\ \quad \forall n \in \mathbb{N}.\; \operatorname{is\text{-}measurable\text{-}fn}(s, a, \operatorname{x\text{-}section\text{-}integral}(s, b, c, d, g(n))) \Rightarrow \\ \quad \quad \operatorname{is\text{-}measurable\text{-}fn}(s, a, \operatorname{x\text{-}section\text{-}integral}(s, b, c, d, f)) \wedge \\ \quad \quad \operatorname{integral}(s, a, u, \operatorname{x\text{-}section\text{-}integral}(s, b, c, d, f)) \\ \quad \quad =\; \operatorname{esup}(\operatorname{image}((\lambda n \in \mathbb{N}.\; \operatorname{integral}(s, a, u, \operatorname{x\text{-}section\text{-}integral}(s, b, c, d, g(n)))), \mathbb{N})) \end{array}}}
\vnbprop{y-section-integral-limit}{\vnbstmt{\begin{array}{@{}l@{}} \forall s, a, u, b, c, d, f, g. \\ \quad \operatorname{is\text{-}measure}(s, a, u) \wedge \\ \quad \operatorname{is\text{-}measure}(b, c, d) \wedge \\ \quad \operatorname{is\text{-}measurable\text{-}fn}((s \times b), \operatorname{product\text{-}sigma}(s, a, b, c), f) \wedge \\ \quad g \in (\mathbb{N} \to ((s \times b) \to \operatorname{rr\text{-}pos\text{-}star})) \wedge \\ \quad \forall n \in \mathbb{N}. \\ \quad \quad \operatorname{is\text{-}measurable\text{-}fn}((s \times b), \operatorname{product\text{-}sigma}(s, a, b, c), g(n)) \wedge \\ \quad \forall n \in \mathbb{N}.\; \operatorname{ptwise\text{-}le}((s \times b), g(n), g((n + 1))) \wedge \\ \quad \forall w \in (s \times b).\; (\operatorname{ptwise\text{-}sup}((s \times b), g)(w) = f(w)) \wedge \\ \quad \forall n \in \mathbb{N}.\; \operatorname{is\text{-}measurable\text{-}fn}(b, c, \operatorname{y\text{-}section\text{-}integral}(s, a, u, b, g(n))) \Rightarrow \\ \quad \quad \operatorname{is\text{-}measurable\text{-}fn}(b, c, \operatorname{y\text{-}section\text{-}integral}(s, a, u, b, f)) \wedge \\ \quad \quad \operatorname{integral}(b, c, d, \operatorname{y\text{-}section\text{-}integral}(s, a, u, b, f)) \\ \quad \quad =\; \operatorname{esup}(\operatorname{image}((\lambda n \in \mathbb{N}.\; \operatorname{integral}(b, c, d, \operatorname{y\text{-}section\text{-}integral}(s, a, u, b, g(n)))), \mathbb{N})) \end{array}}}
\vnbprop{fubini}{\vnbstmt{\begin{array}{@{}l@{}} \forall s, a, u, b, c, d, f. \\ \quad \operatorname{is\text{-}measure}(s, a, u) \wedge \\ \quad \operatorname{is\text{-}measure}(b, c, d) \wedge \\ \quad \operatorname{is\text{-}sigma\text{-}finite}(s, a, u) \wedge \\ \quad \operatorname{is\text{-}sigma\text{-}finite}(b, c, d) \wedge \\ \quad \operatorname{is\text{-}measurable\text{-}fn}((s \times b), \operatorname{product\text{-}sigma}(s, a, b, c), f) \Rightarrow \\ \quad \quad \operatorname{is\text{-}measurable\text{-}fn}(s, a, \operatorname{x\text{-}section\text{-}integral}(s, b, c, d, f)) \wedge \\ \quad \quad \operatorname{is\text{-}measurable\text{-}fn}(b, c, \operatorname{y\text{-}section\text{-}integral}(s, a, u, b, f)) \wedge \\ \quad \quad \operatorname{integral}(s, a, u, \operatorname{x\text{-}section\text{-}integral}(s, b, c, d, f)) \\ \quad \quad =\; \operatorname{integral}(b, c, d, \operatorname{y\text{-}section\text{-}integral}(s, a, u, b, f)) \end{array}}}
\vnbprop{esum-as-esup-prefix}{\vnbstmt{\begin{array}{@{}l@{}} \forall f \in (\mathbb{N} \to \operatorname{rr\text{-}pos\text{-}star}). \\ \quad \operatorname{esum}(f) = \operatorname{esup}(\operatorname{image}((\lambda n \in \mathbb{N}.\; \operatorname{finsum}(\operatorname{rr\text{-}pos\text{-}star\text{-}add\text{-}monoid}, f, \mathrm{OS}(n))), \mathbb{N})) \end{array}}}
\vnbprop{integral-esum}{\vnbstmt{\begin{array}{@{}l@{}} \forall a, b, u, g, h. \\ \quad \operatorname{is\text{-}measure}(a, b, u) \wedge \\ \quad \forall v \in \mathbb{N}.\; \operatorname{is\text{-}measurable\text{-}fn}(a, b, g(v)) \wedge \\ \quad h \in (a \to \operatorname{rr\text{-}pos\text{-}star}) \wedge \\ \quad \forall x \in a.\; (h(x) = \operatorname{esum}((\lambda v \in \mathbb{N}.\; g(v)(x)))) \Rightarrow \\ \quad \quad \operatorname{is\text{-}measurable\text{-}fn}(a, b, h) \wedge \\ \quad \quad \operatorname{integral}(a, b, u, h) = \operatorname{esum}((\lambda v \in \mathbb{N}.\; \operatorname{integral}(a, b, u, g(v)))) \end{array}}}
\vnbprop{product-measure-apply}{\vnbstmt{\begin{array}{@{}l@{}} \forall s, a, u, b, c, d.; \forall f \in \operatorname{product\text{-}sigma}(s, a, b, c). \\ \quad (\operatorname{product\text{-}measure}(s, a, u, b, c, d)(f) \simeq \operatorname{integral}(s, a, u, \operatorname{x\text{-}section\text{-}measure}(s, b, d, f))) \end{array}}}
\vnbprop{x-section-empty}{\vnbstmt{\forall s, x.\; (\operatorname{x\text{-}section}(s, \emptyset, x) = \emptyset)}}
\vnbprop{product-measure-is-measure}{\vnbstmt{\begin{array}{@{}l@{}} \forall s, a, u, b, c, d. \\ \quad \operatorname{is\text{-}measure}(s, a, u) \wedge \\ \quad \operatorname{is\text{-}measure}(b, c, d) \wedge \\ \quad \operatorname{is\text{-}sigma\text{-}finite}(s, a, u) \wedge \\ \quad \operatorname{is\text{-}sigma\text{-}finite}(b, c, d) \Rightarrow \\ \quad \quad \operatorname{is\text{-}measure}((s \times b), \operatorname{product\text{-}sigma}(s, a, b, c), \operatorname{product\text{-}measure}(s, a, u, b, c, d)) \end{array}}}
\vnbprop{product-measure-symmetric}{\vnbstmt{\begin{array}{@{}l@{}} \forall s, a, u, b, c, d, f. \\ \quad \operatorname{is\text{-}measure}(s, a, u) \wedge \\ \quad \operatorname{is\text{-}measure}(b, c, d) \wedge \\ \quad \operatorname{is\text{-}sigma\text{-}finite}(s, a, u) \wedge \\ \quad \operatorname{is\text{-}sigma\text{-}finite}(b, c, d) \wedge \\ \quad f \in \operatorname{product\text{-}sigma}(s, a, b, c) \Rightarrow \\ \quad \quad \operatorname{product\text{-}measure}(s, a, u, b, c, d)(f) = \operatorname{integral}(b, c, d, \operatorname{y\text{-}section\text{-}measure}(s, b, u, f)) \end{array}}}

\vnbbookfile{theorem-library/l2c-b}{theorem-library/l2c-b}
\noindent{\small l2c-b.scm -- wave of 2026-11-09/10 (agent TH3-L2), wired 2026-11-10 from scratchpad/th3-1109/l2/l2c-b.scm. Library: measure.}\par
\vnbprop{measurable-fn-in-fun}{\vnbstmt{\forall a, b, f.\; (\operatorname{is\text{-}measurable\text{-}fn}(a, b, f) \Rightarrow (f \in (a \to \operatorname{rr\text{-}pos\text{-}star})))}}
\vnbprop{simple-fn-is-measurable}{\vnbstmt{\forall a, b, f.\; (\operatorname{is\text{-}simple\text{-}fn}(a, b, f) \Rightarrow \operatorname{is\text{-}measurable\text{-}fn}(a, b, f))}}
\vnbprop{xsi-indicator}{\vnbstmt{\begin{array}{@{}l@{}} \forall s, a, u, b, c, d, f, x. \\ \quad (\operatorname{is\text{-}measure}(s, a, u) \wedge (\operatorname{is\text{-}measure}(b, c, d) \wedge ((f \in \operatorname{product\text{-}sigma}(s, a, b, c)) \wedge (x \in s)))) \Rightarrow  \\ \quad \quad \operatorname{x\text{-}section\text{-}integral}(s, b, c, d, \operatorname{indicator}((s \times b), f))(x) = d(\operatorname{x\text{-}section}(b, f, x)) \end{array}}}
\vnbprop{xsi-eplus}{\vnbstmt{\begin{array}{@{}l@{}} \forall s, a, u, b, c, d, f, g, x. \\ \quad \operatorname{is\text{-}measure}(s, a, u) \wedge \\ \quad \operatorname{is\text{-}measure}(b, c, d) \wedge \\ \quad \operatorname{is\text{-}measurable\text{-}fn}((s \times b), \operatorname{product\text{-}sigma}(s, a, b, c), f) \wedge \\ \quad \operatorname{is\text{-}measurable\text{-}fn}((s \times b), \operatorname{product\text{-}sigma}(s, a, b, c), g) \wedge \\ \quad x \in s \Rightarrow \\ \quad \quad \operatorname{x\text{-}section\text{-}integral}(s, b, c, d, \operatorname{ptwise\text{-}eplus}((s \times b), f, g))(x) \\ \quad \quad =\; \operatorname{eplus}(\operatorname{x\text{-}section\text{-}integral}(s, b, c, d, f)(x), \operatorname{x\text{-}section\text{-}integral}(s, b, c, d, g)(x)) \end{array}}}
\vnbprop{xsi-scale}{\vnbstmt{\begin{array}{@{}l@{}} \forall s, a, u, b, d, g, f, c, x. \\ \quad \operatorname{is\text{-}measure}(s, a, u) \wedge \\ \quad \operatorname{is\text{-}measure}(b, d, g) \wedge \\ \quad \operatorname{is\text{-}measurable\text{-}fn}((s \times b), \operatorname{product\text{-}sigma}(s, a, b, d), f) \wedge \\ \quad c \in \mathbb{R} \wedge \\ \quad 0 \leq c \wedge \\ \quad x \in s \Rightarrow \\ \quad \quad \operatorname{x\text{-}section\text{-}integral}(s, b, d, g, \operatorname{ptwise\text{-}scale}((s \times b), c, f))(x) \\ \quad \quad =\; \operatorname{etimes}(c, \operatorname{x\text{-}section\text{-}integral}(s, b, d, g, f)(x)) \end{array}}}
\vnbprop{product-integral-simple}{\vnbstmt{\begin{array}{@{}l@{}} \forall a, b, u, c, d, f, s. \\ \quad \operatorname{is\text{-}measure}(a, b, u) \wedge \\ \quad \operatorname{is\text{-}measure}(c, d, f) \wedge \\ \quad \operatorname{is\text{-}sigma\text{-}finite}(a, b, u) \wedge \\ \quad \operatorname{is\text{-}sigma\text{-}finite}(c, d, f) \wedge \\ \quad \operatorname{is\text{-}simple\text{-}fn}((a \times c), \operatorname{product\text{-}sigma}(a, b, c, d), s) \Rightarrow \\ \quad \quad \operatorname{integral}((a \times c), \operatorname{product\text{-}sigma}(a, b, c, d), \operatorname{product\text{-}measure}(a, b, u, c, d, f), s) \\ \quad \quad =\; \operatorname{integral}(a, b, u, \operatorname{x\text{-}section\text{-}integral}(a, c, d, f, s)) \end{array}}}
\vnbprop{tonelli-product}{\vnbstmt{\begin{array}{@{}l@{}} \forall s, a, u, b, c, d, f. \\ \quad \operatorname{is\text{-}measure}(s, a, u) \wedge \\ \quad \operatorname{is\text{-}measure}(b, c, d) \wedge \\ \quad \operatorname{is\text{-}sigma\text{-}finite}(s, a, u) \wedge \\ \quad \operatorname{is\text{-}sigma\text{-}finite}(b, c, d) \wedge \\ \quad \operatorname{is\text{-}measurable\text{-}fn}((s \times b), \operatorname{product\text{-}sigma}(s, a, b, c), f) \Rightarrow \\ \quad \quad \operatorname{integral}((s \times b), \operatorname{product\text{-}sigma}(s, a, b, c), \operatorname{product\text{-}measure}(s, a, u, b, c, d), f) \\ \quad \quad =\; \operatorname{integral}(s, a, u, \operatorname{x\text{-}section\text{-}integral}(s, b, c, d, f)) \end{array}}}
\vnbprop{tonelli-product-y}{\vnbstmt{\begin{array}{@{}l@{}} \forall s, a, u, b, c, d, f. \\ \quad \operatorname{is\text{-}measure}(s, a, u) \wedge \\ \quad \operatorname{is\text{-}measure}(b, c, d) \wedge \\ \quad \operatorname{is\text{-}sigma\text{-}finite}(s, a, u) \wedge \\ \quad \operatorname{is\text{-}sigma\text{-}finite}(b, c, d) \wedge \\ \quad \operatorname{is\text{-}measurable\text{-}fn}((s \times b), \operatorname{product\text{-}sigma}(s, a, b, c), f) \Rightarrow \\ \quad \quad \operatorname{integral}((s \times b), \operatorname{product\text{-}sigma}(s, a, b, c), \operatorname{product\text{-}measure}(s, a, u, b, c, d), f) \\ \quad \quad =\; \operatorname{integral}(b, c, d, \operatorname{y\text{-}section\text{-}integral}(s, a, u, b, f)) \end{array}}}

\vnbbookfile{theorem-library/esup-product}{theorem-library/esup-product}
\noindent{\small ==================================================================== theorem-library/esup-product.scm -- THE ESUP PRODUCT LAW in [0, +inf], THE INTEGRAL OF c f FOR c = +inf, and THE PRODUCT MEASURE OF A RECTANGLE WITH ARBITRARY (possibly infinite) SIDES.}\par
\vnbprop{rps-monotone-step-implies-le}{\vnbstmt{\begin{array}{@{}l@{}} \forall f \in (\mathbb{N} \to \operatorname{rr\text{-}pos\text{-}star}), n, m \in \mathbb{N}. \\ \quad (\forall k \in \mathbb{N}.\; (f(k) \leq f((k + 1))) \wedge (m \leq n)) \Rightarrow (f(m) \leq f(n)) \end{array}}}
\vnbprop{etimes-mono-star-left}{\vnbstmt{\begin{array}{@{}l@{}} \forall s, t, w \in \operatorname{rr\text{-}pos\text{-}star}. \\ \quad (s \leq t) \Rightarrow (\operatorname{etimes}(s, w) \leq \operatorname{etimes}(t, w)) \end{array}}}
\vnbprop{etimes-mono-star}{\vnbstmt{\begin{array}{@{}l@{}} \forall p, q, s, t \in \operatorname{rr\text{-}pos\text{-}star}. \\ \quad ((p \leq q) \wedge (s \leq t)) \Rightarrow (\operatorname{etimes}(p, s) \leq \operatorname{etimes}(q, t)) \end{array}}}
\vnbprop{esup-etimes-ge-star}{\vnbstmt{\begin{array}{@{}l@{}} \forall w \in \operatorname{rr\text{-}pos\text{-}star}, m \in (\mathbb{N} \to \operatorname{rr\text{-}pos\text{-}star}). \\ \quad \operatorname{etimes}(w, \operatorname{esup}(\operatorname{image}(m, \mathbb{N}))) \leq \operatorname{esup}(\operatorname{image}((\lambda n \in \mathbb{N}.\; \operatorname{etimes}(w, m(n))), \mathbb{N})) \end{array}}}
\vnbprop{esup-etimes-eq}{\vnbstmt{\begin{array}{@{}l@{}} \forall u, v. \\ \quad u \in (\mathbb{N} \to \operatorname{rr\text{-}pos\text{-}star}) \wedge \\ \quad v \in (\mathbb{N} \to \operatorname{rr\text{-}pos\text{-}star}) \wedge \\ \quad \forall n \in \mathbb{N}.\; (u(n) \leq u((n + 1))) \wedge \\ \quad \forall n \in \mathbb{N}.\; (v(n) \leq v((n + 1))) \Rightarrow \\ \quad \quad \operatorname{esup}(\operatorname{image}((\lambda a \in \mathbb{N}.\; \operatorname{etimes}(u(a), v(a))), \mathbb{N})) \\ \quad \quad =\; \operatorname{etimes}(\operatorname{esup}(\operatorname{image}(u, \mathbb{N})), \operatorname{esup}(\operatorname{image}(v, \mathbb{N}))) \end{array}}}
\vnbprop{esup-nn-identity}{\vnbstmt{\begin{array}{@{}l@{}} \forall g \in (\mathbb{N} \to \operatorname{rr\text{-}pos\text{-}star}). \\ \quad \forall k \in \mathbb{N}.\; (g(k) = k) \Rightarrow (\operatorname{esup}(\operatorname{image}(g, \mathbb{N})) = \operatorname{pos\text{-}inf}) \end{array}}}
\vnbprop{esup-etimes-nn-pos-inf}{\vnbstmt{\begin{array}{@{}l@{}} \forall t \in \operatorname{rr\text{-}pos\text{-}star}, h \in (\mathbb{N} \to \operatorname{rr\text{-}pos\text{-}star}). \\ \quad \forall k \in \mathbb{N}.\; (h(k) = \operatorname{etimes}(k, t)) \Rightarrow  \\ \quad \quad \operatorname{esup}(\operatorname{image}(h, \mathbb{N})) = \operatorname{etimes}(\operatorname{pos\text{-}inf}, t) \end{array}}}
\vnbprop{integral-homogeneous-star}{\vnbstmt{\begin{array}{@{}l@{}} \forall a, b, u, f, c. \\ \quad (\operatorname{is\text{-}measure}(a, b, u) \wedge (\operatorname{is\text{-}measurable\text{-}fn}(a, b, f) \wedge (c \in \operatorname{rr\text{-}pos\text{-}star}))) \Rightarrow  \\ \quad \quad \operatorname{integral}(a, b, u, \operatorname{ptwise\text{-}scale}(a, c, f)) = \operatorname{etimes}(c, \operatorname{integral}(a, b, u, f)) \end{array}}}
\vnbprop{product-measure-rectangle}{\vnbstmt{\begin{array}{@{}l@{}} \forall s, a, u, b, c, d, f, g. \\ \quad (\operatorname{is\text{-}measure}(s, a, u) \wedge (\operatorname{is\text{-}measure}(b, c, d) \wedge ((f \in a) \wedge (g \in c)))) \Rightarrow  \\ \quad \quad \operatorname{product\text{-}measure}(s, a, u, b, c, d)((f \times g)) = \operatorname{etimes}(u(f), d(g)) \end{array}}}

\vnbbookfile{theorem-library/l2c-c}{theorem-library/l2c-c}
\noindent{\small l2c-c.scm -- wave of 2026-11-09/10 (agent TH3-L2), wired 2026-11-10 from scratchpad/th3-1109/l2/l2c-c.scm. Library: measure.}\par
\vnbprop{measurable-fn-tail}{\vnbstmt{\begin{array}{@{}l@{}} \forall a, b, f.; \forall \alpha \in \mathbb{R}. \\ \quad \operatorname{is\text{-}measurable\text{-}fn}(a, b, f) \Rightarrow (\{x \in a : (\alpha < f(x))\} \in b) \end{array}}}
\vnbprop{measurable-fn-pullback-first}{\vnbstmt{\begin{array}{@{}l@{}} \forall s, b, c, d, a. \\ \quad (\operatorname{is\text{-}sigma\text{-}algebra}(s, b) \wedge (\operatorname{is\text{-}sigma\text{-}algebra}(c, d) \wedge \operatorname{is\text{-}measurable\text{-}fn}(s, b, a))) \Rightarrow  \\ \quad \quad \operatorname{is\text{-}measurable\text{-}fn}((s \times c), \operatorname{product\text{-}sigma}(s, b, c, d), \\ \quad \quad \quad (\lambda w \in (s \times c).\; a(\operatorname{nth}(1, w)))) \end{array}}}
\vnbprop{measurable-fn-pullback-second}{\vnbstmt{\begin{array}{@{}l@{}} \forall s, a, c, d, b. \\ \quad (\operatorname{is\text{-}sigma\text{-}algebra}(s, a) \wedge (\operatorname{is\text{-}sigma\text{-}algebra}(c, d) \wedge \operatorname{is\text{-}measurable\text{-}fn}(c, d, b))) \Rightarrow  \\ \quad \quad \operatorname{is\text{-}measurable\text{-}fn}((s \times c), \operatorname{product\text{-}sigma}(s, a, c, d), \\ \quad \quad \quad (\lambda w \in (s \times c).\; b(\operatorname{nth}(2, w)))) \end{array}}}
\vnbprop{measurable-fn-tensor}{\vnbstmt{\begin{array}{@{}l@{}} \forall s, c, d, f, a, b. \\ \quad \operatorname{is\text{-}sigma\text{-}algebra}(s, c) \wedge \\ \quad \operatorname{is\text{-}sigma\text{-}algebra}(d, f) \wedge \\ \quad \operatorname{is\text{-}measurable\text{-}fn}(s, c, a) \wedge \\ \quad \operatorname{is\text{-}measurable\text{-}fn}(d, f, b) \Rightarrow \\ \quad \quad \operatorname{is\text{-}measurable\text{-}fn}((s \times d), \operatorname{product\text{-}sigma}(s, c, d, f), \\ \quad \quad \quad (\lambda w \in (s \times d).\; \operatorname{etimes}(a(\operatorname{nth}(1, w)), b(\operatorname{nth}(2, w))))) \end{array}}}
\vnbprop{xsi-tensor}{\vnbstmt{\begin{array}{@{}l@{}} \forall s, c, d, f, u, a, b, x. \\ \quad \operatorname{is\text{-}sigma\text{-}algebra}(s, c) \wedge \\ \quad \operatorname{is\text{-}sigma\text{-}algebra}(d, f) \wedge \\ \quad \operatorname{is\text{-}measurable\text{-}fn}(s, c, a) \wedge \\ \quad \operatorname{is\text{-}measurable\text{-}fn}(d, f, b) \wedge \\ \quad \operatorname{is\text{-}measure}(d, f, u) \wedge \\ \quad x \in s \Rightarrow \\ \quad \quad \operatorname{x\text{-}section\text{-}integral}(s, d, f, u, (\lambda w \in (s \times d).\; \operatorname{etimes}(a(\operatorname{nth}(1, w)), b(\operatorname{nth}(2, w)))))(x) \\ \quad \quad =\; \operatorname{etimes}(a(x), \operatorname{integral}(d, f, u, b)) \end{array}}}
\vnbprop{tonelli-tensor}{\vnbstmt{\begin{array}{@{}l@{}} \forall s, c, u, d, f, g, a, b. \\ \quad \operatorname{is\text{-}measure}(s, c, u) \wedge \\ \quad \operatorname{is\text{-}measure}(d, f, g) \wedge \\ \quad \operatorname{is\text{-}sigma\text{-}finite}(s, c, u) \wedge \\ \quad \operatorname{is\text{-}sigma\text{-}finite}(d, f, g) \wedge \\ \quad \operatorname{is\text{-}measurable\text{-}fn}(s, c, a) \wedge \\ \quad \operatorname{is\text{-}measurable\text{-}fn}(d, f, b) \Rightarrow \\ \quad \quad \operatorname{integral}((s \times d), \operatorname{product\text{-}sigma}(s, c, d, f), \\ \quad \quad \quad \operatorname{product\text{-}measure}(s, c, u, d, f, g), \\ \quad \quad \quad (\lambda w \in (s \times d).\; \operatorname{etimes}(a(\operatorname{nth}(1, w)), b(\operatorname{nth}(2, w))))) \\ \quad \quad =\; \operatorname{etimes}(\operatorname{integral}(s, c, u, a), \operatorname{integral}(d, f, g, b)) \end{array}}}

\vnbbookfile{theorem-library/hausdorff-measure}{theorem-library/hausdorff-measure}
\noindent{\small mh-hausdorff.scm -- construction-of-measures.pdf ch. 4 section 4: Hausdorff measures. Props 4.32, 4.38 (comparison for comparable metrics), 4.53 (Hausdorff dimension), 4.54 (local dimension), 4.55 (restriction to a subspace).  Needs mh-vocab and mh-method2.  Helper prefix mhh-.}\par
\vnbprop{einf-etimes-bound}{\vnbstmt{\begin{array}{@{}l@{}} \forall s, k, b. \\ \quad s \subseteq \operatorname{rr\text{-}pos\text{-}star} \wedge \\ \quad k \in \mathbb{R} \wedge \\ \quad 0 < k \wedge \\ \quad b \in \operatorname{rr\text{-}pos\text{-}star} \wedge \\ \quad \forall v \in s.\; (b \leq \operatorname{etimes}(k, v)) \Rightarrow \\ \quad \quad b \leq \operatorname{etimes}(k, \operatorname{einf}(s)) \end{array}}}
\vnbprop{method-ii-compare}{\vnbstmt{\begin{array}{@{}l@{}} \forall m, b, p, d, c, k, s. \\ \quad \operatorname{is\text{-}metric\text{-}space}(m) \wedge \\ \quad \operatorname{is\text{-}metric\text{-}space}(b) \wedge \\ \quad \operatorname{pts}(b) = \operatorname{pts}(m) \wedge \\ \quad \operatorname{is\text{-}premeasure}(\operatorname{pts}(m), \operatorname{power}(\operatorname{pts}(m)), p) \wedge \\ \quad \operatorname{is\text{-}premeasure}(\operatorname{pts}(b), \operatorname{power}(\operatorname{pts}(b)), d) \wedge \\ \quad c \in \mathbb{R} \wedge \\ \quad 0 < c \wedge \\ \quad k \in \mathbb{R} \wedge \\ \quad 0 < k \wedge \\ \quad \forall a \in \operatorname{power}(\operatorname{pts}(m)).\; (\operatorname{diam}(m, a) \leq \operatorname{etimes}(c, \operatorname{diam}(b, a))) \wedge \\ \quad \forall a \in \operatorname{power}(\operatorname{pts}(m)).\; (p(a) \leq \operatorname{etimes}(k, d(a))) \wedge \\ \quad s \in \operatorname{power}(\operatorname{pts}(m)) \Rightarrow \\ \quad \quad \operatorname{method\text{-}ii}(m, \operatorname{power}(\operatorname{pts}(m)), p)(s) \leq \operatorname{etimes}(k, \operatorname{method\text{-}ii}(b, \operatorname{power}(\operatorname{pts}(b)), d)(s)) \end{array}}}
\vnbprop{diam-compare}{\vnbstmt{\begin{array}{@{}l@{}} \forall m, b, c, a. \\ \quad \operatorname{is\text{-}metric\text{-}space}(m) \wedge \\ \quad \operatorname{is\text{-}metric\text{-}space}(b) \wedge \\ \quad \operatorname{pts}(b) = \operatorname{pts}(m) \wedge \\ \quad c \in \mathbb{R} \wedge \\ \quad 0 < c \wedge \\ \quad \forall x, y \in \operatorname{pts}(m).\; (\operatorname{dist}(m)(x, y) \leq (c \cdot \operatorname{dist}(b)(x, y))) \wedge \\ \quad a \in \operatorname{power}(\operatorname{pts}(m)) \Rightarrow \\ \quad \quad \operatorname{diam}(m, a) \leq \operatorname{etimes}(c, \operatorname{diam}(b, a)) \end{array}}}
\vnbprop{hausdorff-pm-value}{\vnbstmt{\begin{array}{@{}l@{}} \forall m, b.; \forall a \in \operatorname{power}(\operatorname{pts}(m)). \\ \quad (\operatorname{hausdorff\text{-}pm}(m, b)(a) \simeq {\operatorname{diam}(m, a)}^{b}) \end{array}}}
\vnbprop{hausdorff-pm-premeasure}{\vnbstmt{\begin{array}{@{}l@{}} \forall m, a. \\ \quad (\operatorname{is\text{-}metric\text{-}space}(m) \wedge ((a \in \mathbb{R}) \wedge (0 < a))) \Rightarrow  \\ \quad \quad \operatorname{is\text{-}premeasure}(\operatorname{pts}(m), \operatorname{power}(\operatorname{pts}(m)), \operatorname{hausdorff\text{-}pm}(m, a)) \end{array}}}
\vnbprop{hausdorff-metric-outer-measure}{\vnbstmt{\begin{array}{@{}l@{}} \forall m, a. \\ \quad (\operatorname{is\text{-}metric\text{-}space}(m) \wedge ((a \in \mathbb{R}) \wedge (0 < a))) \Rightarrow  \\ \quad \quad \operatorname{is\text{-}metric\text{-}outer\text{-}measure}(m, \operatorname{hausdorff}(m, a)) \end{array}}}
\vnbprop{hausdorff-compare}{\vnbstmt{\begin{array}{@{}l@{}} \forall m, a, c, b, s. \\ \quad \operatorname{is\text{-}metric\text{-}space}(m) \wedge \\ \quad \operatorname{is\text{-}metric\text{-}space}(a) \wedge \\ \quad \operatorname{pts}(a) = \operatorname{pts}(m) \wedge \\ \quad c \in \mathbb{R} \wedge \\ \quad 0 < c \wedge \\ \quad \forall x, y \in \operatorname{pts}(m).\; (\operatorname{dist}(m)(x, y) \leq (c \cdot \operatorname{dist}(a)(x, y))) \wedge \\ \quad b \in \mathbb{R} \wedge \\ \quad 0 < b \wedge \\ \quad s \in \operatorname{power}(\operatorname{pts}(m)) \Rightarrow \\ \quad \quad \operatorname{hausdorff}(m, b)(s) \leq \operatorname{etimes}({c}^{b}, \operatorname{hausdorff}(a, b)(s)) \end{array}}}
\vnbprop{phi-pm-value}{\vnbstmt{\begin{array}{@{}l@{}} \forall m, i.; \forall a \in \operatorname{power}(\operatorname{pts}(m)). \\ \quad (\operatorname{phi\text{-}pm}(m, i)(a) \simeq \left(\text{if }(\operatorname{diam}(m, a) = \operatorname{pos\text{-}inf})\text{ then }\operatorname{pos\text{-}inf}\text{ else }i(\operatorname{diam}(m, a))\right)) \end{array}}}
\vnbprop{phi-pm-premeasure}{\vnbstmt{\begin{array}{@{}l@{}} \forall m, i. \\ \quad (\operatorname{is\text{-}metric\text{-}space}(m) \wedge \operatorname{is\text{-}doubling}(i)) \Rightarrow  \\ \quad \quad \operatorname{is\text{-}premeasure}(\operatorname{pts}(m), \operatorname{power}(\operatorname{pts}(m)), \operatorname{phi\text{-}pm}(m, i)) \end{array}}}
\vnbprop{doubling-power}{\vnbstmt{\begin{array}{@{}l@{}} \forall n \in \mathbb{N}, i, l.; \forall r \in \mathbb{R}. \\ \quad i \in (\mathit{nnreal} \to \mathit{nnreal}) \wedge \\ \quad \operatorname{pos\text{-}rr}(l) \wedge \\ \quad \forall \mathit{mhr} \in \mathbb{R}.\; ((0 < \mathit{mhr}) \Rightarrow (i((2 \cdot \mathit{mhr})) \leq (l \cdot i(\mathit{mhr})))) \wedge \\ \quad 0 < r \Rightarrow \\ \quad \quad i((\operatorname{power}(2, n) \cdot r)) \leq (\operatorname{power}(l, n) \cdot i(r)) \end{array}}}
\vnbprop{doubling-scale}{\vnbstmt{\begin{array}{@{}l@{}} \forall i.; \forall c \in \mathbb{R}. \\ \quad (\operatorname{is\text{-}doubling}(i) \wedge (0 < c)) \Rightarrow  \\ \quad \quad \exists k. \\ \quad \quad \quad \operatorname{pos\text{-}rr}(k) \wedge \forall r \in \mathbb{R}.\; ((0 \leq r) \Rightarrow (i((c \cdot r)) \leq (k \cdot i(r)))) \end{array}}}
\vnbprop{hausdorff-phi-compare}{\vnbstmt{\begin{array}{@{}l@{}} \forall m, a, c, i. \\ \quad \operatorname{is\text{-}metric\text{-}space}(m) \wedge \\ \quad \operatorname{is\text{-}metric\text{-}space}(a) \wedge \\ \quad \operatorname{pts}(a) = \operatorname{pts}(m) \wedge \\ \quad c \in \mathbb{R} \wedge \\ \quad 0 < c \wedge \\ \quad \forall x, y \in \operatorname{pts}(m).\; (\operatorname{dist}(m)(x, y) \leq (c \cdot \operatorname{dist}(a)(x, y))) \wedge \\ \quad \operatorname{is\text{-}doubling}(i) \Rightarrow \\ \quad \quad \exists k. \\ \quad \quad \quad \operatorname{pos\text{-}rr}(k) \wedge \\ \quad \quad \quad \forall s \in \operatorname{power}(\operatorname{pts}(m)). \\ \quad \quad \quad \quad \operatorname{hausdorff\text{-}phi}(m, i)(s) \leq \operatorname{etimes}(k, \operatorname{hausdorff\text{-}phi}(a, i)(s)) \end{array}}}

\vnbbookfile{theorem-library/hausdorff-dimension}{theorem-library/hausdorff-dimension}
\noindent{\small mh-dim.scm -- Hausdorff dimension (construction-of-measures.pdf, 4.3): Prop 4.53 and its corollary (uniqueness of the dimension), Prop 4.54.  Prefix mhd-.  Library: measure. Loads after mh-vocab.scm, mh-method2.scm, mh-hausdorff.scm.}\par
\vnbprop{mhd-rpow-key}{\vnbstmt{\begin{array}{@{}l@{}} \forall t, l, o, m. \\ \quad ((t \in \mathbb{R}) \wedge ((0 < t) \wedge ((l \in \mathbb{R}) \wedge ((t \leq l) \wedge ((o \in \mathbb{R}) \wedge ((m \in \mathbb{R}) \wedge (o < m))))))) \Rightarrow  \\ \quad \quad ({l}^{(o - m)} \cdot {t}^{m}) \leq {t}^{o} \end{array}}}
\vnbprop{mhd-epow-key}{\vnbstmt{\begin{array}{@{}l@{}} \forall t, l, o, m. \\ \quad ((t \in \operatorname{rr\text{-}pos\text{-}star}) \wedge ((l \in \mathbb{R}) \wedge ((0 < l) \wedge ((t \leq l) \wedge ((o \in \mathbb{R}) \wedge ((0 < o) \wedge ((m \in \mathbb{R}) \wedge (o < m)))))))) \Rightarrow  \\ \quad \quad \operatorname{etimes}({l}^{(o - m)}, {t}^{m}) \leq {t}^{o} \end{array}}}
\vnbprop{method-ii-delta-scale-below}{\vnbstmt{\begin{array}{@{}l@{}} \forall m, s, p, b, d, k, c. \\ \quad \operatorname{is\text{-}premeasure}(\operatorname{pts}(m), s, p) \wedge \\ \quad \operatorname{is\text{-}premeasure}(\operatorname{pts}(m), s, b) \wedge \\ \quad d \in \mathbb{R} \wedge \\ \quad 0 \leq d \wedge \\ \quad k \in \mathbb{R} \wedge \\ \quad 0 < k \wedge \\ \quad \forall a \in s.\; ((\operatorname{diam}(m, a) \leq d) \Rightarrow (\operatorname{etimes}(k, b(a)) \leq p(a))) \wedge \\ \quad c \in \operatorname{power}(\operatorname{pts}(m)) \Rightarrow \\ \quad \quad \operatorname{etimes}(k, \operatorname{method\text{-}ii\text{-}delta}(m, s, b, d)(c)) \leq \operatorname{method\text{-}ii\text{-}delta}(m, s, p, d)(c) \end{array}}}
\vnbprop{hausdorff-unfold}{\vnbstmt{\begin{array}{@{}l@{}} \forall m, a. \\ \quad (\operatorname{hausdorff}(m, a) \simeq \operatorname{method\text{-}ii}(m, \operatorname{power}(\operatorname{pts}(m)), \operatorname{hausdorff\text{-}pm}(m, a))) \end{array}}}
\vnbprop{hausdorff-delta-below}{\vnbstmt{\begin{array}{@{}l@{}} \forall m, d, o, a, s. \\ \quad (\operatorname{is\text{-}metric\text{-}space}(m) \wedge (\operatorname{pos\text{-}rr}(d) \wedge ((o \in \mathbb{R}) \wedge ((0 < o) \wedge ((a \in \mathbb{R}) \wedge ((o < a) \wedge (s \in \operatorname{power}(\operatorname{pts}(m))))))))) \Rightarrow  \\ \quad \quad \operatorname{etimes}({d}^{(o - a)}, \operatorname{method\text{-}ii\text{-}delta}(m, \operatorname{power}(\operatorname{pts}(m)), \operatorname{hausdorff\text{-}pm}(m, a), d)(s)) \\ \quad \quad \leq\; \operatorname{method\text{-}ii\text{-}delta}(m, \operatorname{power}(\operatorname{pts}(m)), \operatorname{hausdorff\text{-}pm}(m, o), d)(s) \end{array}}}
\vnbprop{mhd-small-power}{\vnbstmt{\begin{array}{@{}l@{}} \forall p, q, a. \\ \quad ((p \in \mathbb{R}) \wedge ((0 < p) \wedge ((q \in \mathbb{R}) \wedge ((0 < q) \wedge ((a \in \mathbb{R}) \wedge (0 < a)))))) \Rightarrow  \\ \quad \quad \exists d.\; (\operatorname{pos\text{-}rr}(d) \wedge ((d \leq a) \wedge ({d}^{p} \leq q))) \end{array}}}
\vnbprop{hausdorff-scale-pos}{\vnbstmt{\begin{array}{@{}l@{}} \forall m, a, s. \\ \quad (\operatorname{is\text{-}metric\text{-}space}(m) \wedge ((a \in \mathbb{R}) \wedge ((0 < a) \wedge ((s \in \operatorname{power}(\operatorname{pts}(m))) \wedge (0 < \operatorname{hausdorff}(m, a)(s)))))) \Rightarrow  \\ \quad \quad \exists d. \\ \quad \quad \quad (\operatorname{pos\text{-}rr}(d) \wedge \neg (\operatorname{method\text{-}ii\text{-}delta}(m, \operatorname{power}(\operatorname{pts}(m)), \operatorname{hausdorff\text{-}pm}(m, a), d)(s) = 0)) \end{array}}}
\vnbprop{hausdorff-positive-infinite}{\vnbstmt{\begin{array}{@{}l@{}} \forall m, a, o, s. \\ \quad \operatorname{is\text{-}metric\text{-}space}(m) \wedge \\ \quad a \in \mathbb{R} \wedge \\ \quad o \in \mathbb{R} \wedge \\ \quad 0 < o \wedge \\ \quad o < a \wedge \\ \quad s \in \operatorname{power}(\operatorname{pts}(m)) \wedge \\ \quad 0 < \operatorname{hausdorff}(m, a)(s) \Rightarrow \\ \quad \quad \operatorname{hausdorff}(m, o)(s) = \operatorname{pos\text{-}inf} \end{array}}}
\vnbprop{hausdorff-dim-not-lt}{\vnbstmt{\begin{array}{@{}l@{}} \forall m, s, a, b. \\ \quad (\operatorname{is\text{-}hausdorff\text{-}dim}(m, s, a) \wedge \operatorname{is\text{-}hausdorff\text{-}dim}(m, s, b)) \Rightarrow \neg (a < b) \end{array}}}
\vnbprop{hausdorff-dim-unique}{\vnbstmt{\begin{array}{@{}l@{}} \forall m, s, a, b. \\ \quad (\operatorname{is\text{-}hausdorff\text{-}dim}(m, s, a) \wedge \operatorname{is\text{-}hausdorff\text{-}dim}(m, s, b)) \Rightarrow (a = b) \end{array}}}
\vnbprop{hausdorff-dim-local}{\vnbstmt{\begin{array}{@{}l@{}} \forall m, r, a. \\ \quad \operatorname{is\text{-}separable}(m) \wedge \\ \quad \operatorname{pos\text{-}rr}(r) \wedge \\ \quad \exists \mathit{mhx}.\; (\mathit{mhx} \in \operatorname{pts}(m)) \wedge \\ \quad \forall x \in \operatorname{pts}(m).\; \operatorname{is\text{-}hausdorff\text{-}dim}(m, \operatorname{ball}(m, x, r), a) \Rightarrow \\ \quad \quad \operatorname{is\text{-}hausdorff\text{-}dim}(m, \operatorname{pts}(m), a) \end{array}}}

\vnbbookfile{theorem-library/hausdorff-restrict}{theorem-library/hausdorff-restrict}
\noindent{\small mh-restrict.scm -- restriction of Hausdorff measure to a subspace (construction-of-measures.pdf, 4.4, Prop 4.55).  Prefix mhr-.  Library: measure. Loads after mh-vocab.scm, mh-method2.scm, mh-hausdorff.scm.}\par
\vnbprop{diam-le-of-dist}{\vnbstmt{\begin{array}{@{}l@{}} \forall m, b, a. \\ \quad \operatorname{is\text{-}metric\text{-}space}(m) \wedge \\ \quad \operatorname{is\text{-}metric\text{-}space}(b) \wedge \\ \quad a \subseteq \operatorname{pts}(m) \wedge \\ \quad a \subseteq \operatorname{pts}(b) \wedge \\ \quad \forall x, y \in a.\; (\operatorname{dist}(m)(x, y) \leq \operatorname{dist}(b)(x, y)) \Rightarrow \\ \quad \quad \operatorname{diam}(m, a) \leq \operatorname{diam}(b, a) \end{array}}}
\vnbprop{diam-mono}{\vnbstmt{\begin{array}{@{}l@{}} \forall m, a, b. \\ \quad (\operatorname{is\text{-}metric\text{-}space}(m) \wedge ((a \subseteq b) \wedge (b \subseteq \operatorname{pts}(m)))) \Rightarrow  \\ \quad \quad \operatorname{diam}(m, a) \leq \operatorname{diam}(m, b) \end{array}}}
\vnbprop{subspace-diam}{\vnbstmt{\begin{array}{@{}l@{}} \forall m, s, a. \\ \quad (\operatorname{is\text{-}metric\text{-}space}(m) \wedge ((s \subseteq \operatorname{pts}(m)) \wedge (a \subseteq s))) \Rightarrow  \\ \quad \quad \operatorname{diam}(\operatorname{subspace\text{-}ms}(m, s), a) = \operatorname{diam}(m, a) \end{array}}}
\vnbprop{phi-pm-mono}{\vnbstmt{\begin{array}{@{}l@{}} \forall m, i, a, b. \\ \quad (\operatorname{is\text{-}metric\text{-}space}(m) \wedge (\operatorname{is\text{-}doubling}(i) \wedge ((a \subseteq b) \wedge (b \in \operatorname{power}(\operatorname{pts}(m)))))) \Rightarrow  \\ \quad \quad \operatorname{phi\text{-}pm}(m, i)(a) \leq \operatorname{phi\text{-}pm}(m, i)(b) \end{array}}}
\vnbprop{phi-pm-subspace}{\vnbstmt{\begin{array}{@{}l@{}} \forall m, s, i, a. \\ \quad (\operatorname{is\text{-}metric\text{-}space}(m) \wedge ((s \subseteq \operatorname{pts}(m)) \wedge (a \subseteq s))) \Rightarrow  \\ \quad \quad (\operatorname{phi\text{-}pm}(\operatorname{subspace\text{-}ms}(m, s), i)(a) \simeq \operatorname{phi\text{-}pm}(m, i)(a)) \end{array}}}
\vnbprop{method-ii-delta-subspace-le}{\vnbstmt{\begin{array}{@{}l@{}} \forall m, s, i, d, a. \\ \quad (\operatorname{is\text{-}metric\text{-}space}(m) \wedge (\operatorname{is\text{-}doubling}(i) \wedge ((s \subseteq \operatorname{pts}(m)) \wedge ((d \in \mathbb{R}) \wedge ((0 \leq d) \wedge (a \subseteq s)))))) \Rightarrow  \\ \quad \quad \operatorname{method\text{-}ii\text{-}delta}(m, \operatorname{power}(\operatorname{pts}(m)), \operatorname{phi\text{-}pm}(m, i), d)(a) \\ \quad \quad \leq\; \operatorname{method\text{-}ii\text{-}delta}(\operatorname{subspace\text{-}ms}(m, s), \operatorname{power}(\operatorname{pts}(\operatorname{subspace\text{-}ms}(m, s))), \operatorname{phi\text{-}pm}(\operatorname{subspace\text{-}ms}(m, s), i), d)(a) \end{array}}}
\vnbprop{method-ii-delta-subspace-ge}{\vnbstmt{\begin{array}{@{}l@{}} \forall m, s, i, d, a. \\ \quad (\operatorname{is\text{-}metric\text{-}space}(m) \wedge (\operatorname{is\text{-}doubling}(i) \wedge ((s \subseteq \operatorname{pts}(m)) \wedge ((d \in \mathbb{R}) \wedge ((0 \leq d) \wedge (a \subseteq s)))))) \Rightarrow  \\ \quad \quad \operatorname{method\text{-}ii\text{-}delta}(\operatorname{subspace\text{-}ms}(m, s), \operatorname{power}(\operatorname{pts}(\operatorname{subspace\text{-}ms}(m, s))), \operatorname{phi\text{-}pm}(\operatorname{subspace\text{-}ms}(m, s), i), d)(a) \\ \quad \quad \leq\; \operatorname{method\text{-}ii\text{-}delta}(m, \operatorname{power}(\operatorname{pts}(m)), \operatorname{phi\text{-}pm}(m, i), d)(a) \end{array}}}
\vnbprop{hausdorff-phi-subspace}{\vnbstmt{\begin{array}{@{}l@{}} \forall m, s, i, a. \\ \quad (\operatorname{is\text{-}metric\text{-}space}(m) \wedge (\operatorname{is\text{-}doubling}(i) \wedge ((s \subseteq \operatorname{pts}(m)) \wedge (a \subseteq s)))) \Rightarrow  \\ \quad \quad \operatorname{hausdorff\text{-}phi}(\operatorname{subspace\text{-}ms}(m, s), i)(a) = \operatorname{hausdorff\text{-}phi}(m, i)(a) \end{array}}}

\vnbbookfile{structure-library/sip-l2-vocab}{structure-library/sip-l2-vocab}
\noindent{\small sip-l2-vocab.scm -- EQUALITY ALMOST EVERYWHERE as a setoid (SIP-1, 2026-10-18). Vocabulary only.}\par
\par\medskip\noindent\textbf{Definition} (\texttt{ae-rel}).\;$\operatorname{ae\text{-}rel}(\omega, \mathit{ca}, \mu, \mathit{sf}) :=$
\begin{equation*}
\{\operatorname{aep\_} \in (\mathit{sf} \times \mathit{sf}) : \exists \operatorname{aen\_} \in \mathbf{Set}.\; (\operatorname{is\text{-}null\text{-}set}(\omega, \mathit{ca}, \mu, \operatorname{aen\_}) \wedge \forall \operatorname{aex\_} \in \omega.\; (\neg (\operatorname{aex\_} \in \operatorname{aen\_}) \Rightarrow (\operatorname{nth}(1, \operatorname{aep\_})(\operatorname{aex\_}) = \operatorname{nth}(2, \operatorname{aep\_})(\operatorname{aex\_}))))\}
\end{equation*}

\par\medskip\noindent\textbf{Definition} (\texttt{ae-setoid}).\;$\operatorname{ae\text{-}setoid}(\omega, \mathit{ca}, \mu, \mathit{sf}) :=$ ${[}\mathit{sf}, \operatorname{ae\text{-}rel}(\omega, \mathit{ca}, \mu, \mathit{sf})]$
\par\medskip\noindent\textbf{Definition} (\texttt{l2-set}).\;$\operatorname{l2\text{-}set}(\omega, \mathit{ca}, \mu) :=$
\begin{equation*}
\{\operatorname{lf\_} \in (\omega \to \mathbb{R}) : (\operatorname{is\text{-}real\text{-}measurable\text{-}fn}(\omega, \mathit{ca}, \operatorname{lf\_}) \wedge \operatorname{is\text{-}integrable}(\omega, \mathit{ca}, \mu, (\lambda \operatorname{lx\_} \in \omega.\; (\operatorname{lf\_}(\operatorname{lx\_}) \cdot \operatorname{lf\_}(\operatorname{lx\_})))))\}
\end{equation*}

\par\medskip\noindent\textbf{Definition} (\texttt{l2-vadd}).\;$\operatorname{l2\text{-}vadd}(\omega, \mathit{ca}, \mu) :=$
\begin{equation*}
\begin{array}{@{}l@{}} (\lambda {[}\operatorname{lf\_}, \operatorname{lg\_}] \in (\operatorname{l2\text{-}set}(\omega, \mathit{ca}, \mu) \times \operatorname{l2\text{-}set}(\omega, \mathit{ca}, \mu)).\; \\ \quad (\lambda \operatorname{lx\_} \in \omega.\; (\operatorname{lf\_}(\operatorname{lx\_}) + \operatorname{lg\_}(\operatorname{lx\_})))) \end{array}
\end{equation*}

\par\medskip\noindent\textbf{Definition} (\texttt{l2-vzero}).\;$\operatorname{l2\text{-}vzero}(\omega) :=$ $(\lambda \operatorname{lx\_} \in \omega.\; 0)$
\par\medskip\noindent\textbf{Definition} (\texttt{l2-vneg}).\;$\operatorname{l2\text{-}vneg}(\omega, \mathit{ca}, \mu) :=$
\begin{equation*}
(\lambda \operatorname{lf\_} \in \operatorname{l2\text{-}set}(\omega, \mathit{ca}, \mu).\; (\lambda \operatorname{lx\_} \in \omega.\; -\operatorname{lf\_}(\operatorname{lx\_})))
\end{equation*}

\par\medskip\noindent\textbf{Definition} (\texttt{l2-act}).\;$\operatorname{l2\text{-}act}(\omega, \mathit{ca}, \mu) :=$
\begin{equation*}
(\lambda {[}\operatorname{lr\_}, \operatorname{lf\_}] \in (\mathbb{R} \times \operatorname{l2\text{-}set}(\omega, \mathit{ca}, \mu)).\; (\lambda \operatorname{lx\_} \in \omega.\; (\operatorname{lr\_} \cdot \operatorname{lf\_}(\operatorname{lx\_}))))
\end{equation*}

\par\medskip\noindent\textbf{Definition} (\texttt{l2-rip}).\;$\operatorname{l2\text{-}rip}(\omega, \mathit{ca}, \mu) :=$
\begin{equation*}
\begin{array}{@{}l@{}} (\lambda {[}\operatorname{lf\_}, \operatorname{lg\_}] \in (\operatorname{l2\text{-}set}(\omega, \mathit{ca}, \mu) \times \operatorname{l2\text{-}set}(\omega, \mathit{ca}, \mu)).\; \\ \quad \operatorname{rintegral}(\omega, \mathit{ca}, \mu, (\lambda \operatorname{lx\_} \in \omega.\; (\operatorname{lf\_}(\operatorname{lx\_}) \cdot \operatorname{lg\_}(\operatorname{lx\_}))))) \end{array}
\end{equation*}

\par\medskip\noindent\textbf{Definition} (\texttt{l2-space}).\;$\operatorname{l2\text{-}space}(\omega, \mathit{ca}, \mu) :=$
\begin{equation*}
{[}\operatorname{normed\text{-}field\text{-}as\text{-}commutative\text{-}ring}(\operatorname{rr\text{-}normed\text{-}field}), \operatorname{l2\text{-}set}(\omega, \mathit{ca}, \mu), \operatorname{l2\text{-}vadd}(\omega, \mathit{ca}, \mu), \operatorname{l2\text{-}vzero}(\omega), \operatorname{l2\text{-}vneg}(\omega, \mathit{ca}, \mu), \operatorname{l2\text{-}act}(\omega, \mathit{ca}, \mu), \operatorname{l2\text{-}rip}(\omega, \mathit{ca}, \mu)]
\end{equation*}

\par\medskip\noindent\textbf{Definition} (\texttt{cl2-space}).\;$\operatorname{cl2\text{-}space}(\omega, \mathit{ca}, \mu) :=$ $\operatorname{complexify}(\operatorname{l2\text{-}space}(\omega, \mathit{ca}, \mu))$

\vnbbookfile{theorem-library/sip-c}{theorem-library/sip-c}
\noindent{\small sip-c.scm -- EQUALITY ALMOST EVERYWHERE, REAL AND COMPLEX L\textasciicircum{}2 (SIP-1, 2026-10-18; STOP POINT C of scratchpad/triage/SIP-1.md).  All PROVEN modulo 0. Vocabulary: structure-library/sip-l2-vocab.scm.}\par
\vnbprop{ae-null-empty}{\vnbstmt{\forall a, b, u.\; (\operatorname{is\text{-}measure}(a, b, u) \Rightarrow \operatorname{is\text{-}null\text{-}set}(a, b, u, \emptyset))}}
\vnbprop{ae-null-union}{\vnbstmt{\begin{array}{@{}l@{}} \forall a, b, u, \mathit{na}, \mathit{nb}. \\ \quad (\operatorname{is\text{-}measure}(a, b, u) \wedge (\operatorname{is\text{-}null\text{-}set}(a, b, u, \mathit{na}) \wedge \operatorname{is\text{-}null\text{-}set}(a, b, u, \mathit{nb}))) \Rightarrow  \\ \quad \quad \operatorname{is\text{-}null\text{-}set}(a, b, u, (\mathit{na} \cup \mathit{nb})) \end{array}}}
\vnbprop{ae-rel-unfold}{\vnbstmt{\begin{array}{@{}l@{}} \forall a, b, u, f. \\ \quad (\operatorname{ae\text{-}rel}(a, b, u, f) \simeq \{p \in (f \times f) : \exists n \in \mathbf{Set}.\; (\operatorname{is\text{-}null\text{-}set}(a, b, u, n) \wedge \forall x \in a.\; (\neg (x \in n) \Rightarrow (\operatorname{nth}(1, p)(x) = \operatorname{nth}(2, p)(x))))\}) \end{array}}}
\vnbprop{ae-null-is-set}{\vnbstmt{\forall a, b, u, s.\; (\operatorname{is\text{-}null\text{-}set}(a, b, u, s) \Rightarrow (s \in \mathbf{Set}))}}
\vnbprop{ae-rel-in}{\vnbstmt{\begin{array}{@{}l@{}} \forall a, b, u, f.; \forall \mathit{af}, g \in f, s. \\ \quad (\operatorname{is\text{-}null\text{-}set}(a, b, u, s) \wedge \forall x \in a.\; (\neg (x \in s) \Rightarrow (\operatorname{af}(x) = g(x)))) \Rightarrow  \\ \quad \quad {[}\mathit{af}, g] \in \operatorname{ae\text{-}rel}(a, b, u, f) \end{array}}}
\vnbprop{ae-rel-out}{\vnbstmt{\begin{array}{@{}l@{}} \forall a, b, u, f, \mathit{af}, g. \\ \quad ({[}\mathit{af}, g] \in \operatorname{ae\text{-}rel}(a, b, u, f)) \Rightarrow  \\ \quad \quad \exists n. \\ \quad \quad \quad \operatorname{is\text{-}null\text{-}set}(a, b, u, n) \wedge \forall x \in a.\; (\neg (x \in n) \Rightarrow (\operatorname{af}(x) = g(x))) \end{array}}}
\vnbprop{ae-rel-is-equivalence}{\vnbstmt{\begin{array}{@{}l@{}} \forall a, b, u, f. \\ \quad (\operatorname{is\text{-}measure}(a, b, u) \wedge ((f \in \mathbf{Set}) \wedge (f \subseteq (a \to \mathbb{R})))) \Rightarrow  \\ \quad \quad \operatorname{is\text{-}equivalence}(\operatorname{ae\text{-}rel}(a, b, u, f), f) \end{array}}}
\vnbprop{ae-setoid-is-setoid}{\vnbstmt{\begin{array}{@{}l@{}} \forall a, b, u, f. \\ \quad (\operatorname{is\text{-}measure}(a, b, u) \wedge ((f \in \mathbf{Set}) \wedge (f \subseteq (a \to \mathbb{R})))) \Rightarrow  \\ \quad \quad \operatorname{is\text{-}setoid}(\operatorname{ae\text{-}setoid}(a, b, u, f)) \end{array}}}
\vnbprop{real-measurable-mul}{\vnbstmt{\begin{array}{@{}l@{}} \forall a, b, f, g. \\ \quad (\operatorname{is\text{-}real\text{-}measurable\text{-}fn}(a, b, f) \wedge \operatorname{is\text{-}real\text{-}measurable\text{-}fn}(a, b, g)) \Rightarrow  \\ \quad \quad \operatorname{is\text{-}real\text{-}measurable\text{-}fn}(a, b, (\lambda x \in a.\; (f(x) \cdot g(x)))) \end{array}}}
\vnbprop{integrable-dominated}{\vnbstmt{\begin{array}{@{}l@{}} \forall a, b, u, f, h. \\ \quad \operatorname{is\text{-}measure}(a, b, u) \wedge \\ \quad \operatorname{is\text{-}integrable}(a, b, u, h) \wedge \\ \quad \operatorname{is\text{-}real\text{-}measurable\text{-}fn}(a, b, f) \wedge \\ \quad \forall x \in a.\; (\lvert f(x) \rvert \leq \lvert h(x) \rvert) \Rightarrow \\ \quad \quad \operatorname{is\text{-}integrable}(a, b, u, f) \end{array}}}
\vnbprop{integrable-ptwise-eq}{\vnbstmt{\begin{array}{@{}l@{}} \forall a, b, u, f, h. \\ \quad ((f \in (a \to \mathbb{R})) \wedge ((h \in (a \to \mathbb{R})) \wedge (\forall x \in a.\; (f(x) = h(x)) \wedge \operatorname{is\text{-}integrable}(a, b, u, f)))) \Rightarrow  \\ \quad \quad \operatorname{is\text{-}integrable}(a, b, u, h) \wedge (\operatorname{rintegral}(a, b, u, h) = \operatorname{rintegral}(a, b, u, f)) \end{array}}}
\vnbprop{l2-unfold}{\vnbstmt{\begin{array}{@{}l@{}} \forall a, b, u. \\ \quad (\operatorname{l2\text{-}set}(a, b, u) \simeq \{f \in (a \to \mathbb{R}) : (\operatorname{is\text{-}real\text{-}measurable\text{-}fn}(a, b, f) \wedge \operatorname{is\text{-}integrable}(a, b, u, (\lambda x \in a.\; (f(x) \cdot f(x)))))\}) \end{array}}}
\vnbprop{l2-in}{\vnbstmt{\begin{array}{@{}l@{}} \forall a, b, u.; \forall f \in (a \to \mathbb{R}). \\ \quad (\operatorname{is\text{-}real\text{-}measurable\text{-}fn}(a, b, f) \wedge \operatorname{is\text{-}integrable}(a, b, u, (\lambda x \in a.\; (f(x) \cdot f(x))))) \Rightarrow  \\ \quad \quad f \in \operatorname{l2\text{-}set}(a, b, u) \end{array}}}
\vnbprop{l2-out}{\vnbstmt{\begin{array}{@{}l@{}} \forall a, b, u.; \forall f \in \operatorname{l2\text{-}set}(a, b, u). \\ \quad f \in (a \to \mathbb{R}) \wedge \\ \quad \operatorname{is\text{-}real\text{-}measurable\text{-}fn}(a, b, f) \wedge \\ \quad \operatorname{is\text{-}integrable}(a, b, u, (\lambda x \in a.\; (f(x) \cdot f(x)))) \end{array}}}
\vnbprop{real-measurable-ptwise-eq}{\vnbstmt{\begin{array}{@{}l@{}} \forall a, b, f, h. \\ \quad f \in (a \to \mathbb{R}) \wedge \\ \quad h \in (a \to \mathbb{R}) \wedge \\ \quad \forall x \in a.\; (f(x) = h(x)) \wedge \\ \quad \operatorname{is\text{-}real\text{-}measurable\text{-}fn}(a, b, f) \Rightarrow \\ \quad \quad \operatorname{is\text{-}real\text{-}measurable\text{-}fn}(a, b, h) \end{array}}}
\vnbprop{real-measurable-const}{\vnbstmt{\begin{array}{@{}l@{}} \forall a, b, c. \\ \quad (\operatorname{is\text{-}sigma\text{-}algebra}(a, b) \wedge (c \in \mathbb{R})) \Rightarrow  \\ \quad \quad \operatorname{is\text{-}real\text{-}measurable\text{-}fn}(a, b, (\lambda x \in a.\; c)) \end{array}}}
\vnbprop{integrable-zero}{\vnbstmt{\begin{array}{@{}l@{}} \forall a, b, u. \\ \quad \operatorname{is\text{-}measure}(a, b, u) \Rightarrow \operatorname{is\text{-}integrable}(a, b, u, (\lambda x \in a.\; 0)) \end{array}}}
\vnbprop{l2-prod-integrable}{\vnbstmt{\begin{array}{@{}l@{}} \forall a, b, u, f, g. \\ \quad (\operatorname{is\text{-}measure}(a, b, u) \wedge ((f \in \operatorname{l2\text{-}set}(a, b, u)) \wedge (g \in \operatorname{l2\text{-}set}(a, b, u)))) \Rightarrow  \\ \quad \quad \operatorname{is\text{-}integrable}(a, b, u, (\lambda x \in a.\; (f(x) \cdot g(x)))) \end{array}}}
\vnbprop{l2-add-closed}{\vnbstmt{\begin{array}{@{}l@{}} \forall a, b, u, f, g. \\ \quad (\operatorname{is\text{-}measure}(a, b, u) \wedge ((f \in \operatorname{l2\text{-}set}(a, b, u)) \wedge (g \in \operatorname{l2\text{-}set}(a, b, u)))) \Rightarrow  \\ \quad \quad (\lambda x \in a.\; (f(x) + g(x))) \in \operatorname{l2\text{-}set}(a, b, u) \end{array}}}
\vnbprop{l2-scale-closed}{\vnbstmt{\begin{array}{@{}l@{}} \forall a, b, u, r, f. \\ \quad (\operatorname{is\text{-}measure}(a, b, u) \wedge ((r \in \mathbb{R}) \wedge (f \in \operatorname{l2\text{-}set}(a, b, u)))) \Rightarrow  \\ \quad \quad (\lambda x \in a.\; (r \cdot f(x))) \in \operatorname{l2\text{-}set}(a, b, u) \end{array}}}
\vnbprop{l2-neg-closed}{\vnbstmt{\begin{array}{@{}l@{}} \forall a, b, u, f. \\ \quad (\operatorname{is\text{-}measure}(a, b, u) \wedge (f \in \operatorname{l2\text{-}set}(a, b, u))) \Rightarrow  \\ \quad \quad (\lambda x \in a.\; -f(x)) \in \operatorname{l2\text{-}set}(a, b, u) \end{array}}}
\vnbprop{l2-zero-in}{\vnbstmt{\begin{array}{@{}l@{}} \forall a, b, u. \\ \quad \operatorname{is\text{-}measure}(a, b, u) \Rightarrow ((\lambda x \in a.\; 0) \in \operatorname{l2\text{-}set}(a, b, u)) \end{array}}}
\vnbprop{l2-set-is-set}{\vnbstmt{\forall a, b, u.\; (\operatorname{is\text{-}measure}(a, b, u) \Rightarrow (\operatorname{l2\text{-}set}(a, b, u) \in \mathbf{Set}))}}
\vnbprop{rintegral-in-rr}{\vnbstmt{\begin{array}{@{}l@{}} \forall a, b, u.; \forall f \in (a \to \mathbb{R}). \\ \quad \operatorname{is\text{-}integrable}(a, b, u, f) \Rightarrow (\operatorname{rintegral}(a, b, u, f) \in \mathbb{R}) \end{array}}}
\vnbprop{l2-vadd-value}{\vnbstmt{\begin{array}{@{}l@{}} \forall a, b, u, f, g. \\ \quad (\operatorname{is\text{-}measure}(a, b, u) \wedge ((f \in \operatorname{l2\text{-}set}(a, b, u)) \wedge (g \in \operatorname{l2\text{-}set}(a, b, u)))) \Rightarrow  \\ \quad \quad \operatorname{l2\text{-}vadd}(a, b, u)(f, g) = (\lambda x \in a.\; (f(x) + g(x))) \end{array}}}
\vnbprop{l2-vneg-value}{\vnbstmt{\begin{array}{@{}l@{}} \forall a, b, u.; \forall f \in \operatorname{l2\text{-}set}(a, b, u). \\ \quad \operatorname{is\text{-}measure}(a, b, u) \Rightarrow (\operatorname{l2\text{-}vneg}(a, b, u)(f) = (\lambda x \in a.\; -f(x))) \end{array}}}
\vnbprop{l2-act-value}{\vnbstmt{\begin{array}{@{}l@{}} \forall a, b, u, r, f. \\ \quad (\operatorname{is\text{-}measure}(a, b, u) \wedge ((r \in \mathbb{R}) \wedge (f \in \operatorname{l2\text{-}set}(a, b, u)))) \Rightarrow  \\ \quad \quad \operatorname{l2\text{-}act}(a, b, u)(r, f) = (\lambda x \in a.\; (r \cdot f(x))) \end{array}}}
\vnbprop{l2-rip-value}{\vnbstmt{\begin{array}{@{}l@{}} \forall a, b, u, f, g. \\ \quad (\operatorname{is\text{-}measure}(a, b, u) \wedge ((f \in \operatorname{l2\text{-}set}(a, b, u)) \wedge (g \in \operatorname{l2\text{-}set}(a, b, u)))) \Rightarrow  \\ \quad \quad \operatorname{l2\text{-}rip}(a, b, u)(f, g) = \operatorname{rintegral}(a, b, u, (\lambda x \in a.\; (f(x) \cdot g(x)))) \end{array}}}
\vnbprop{l2-space-scal}{\vnbstmt{\begin{array}{@{}l@{}} \forall a, b, u. \\ \quad \operatorname{is\text{-}measure}(a, b, u) \Rightarrow  \\ \quad \quad (\operatorname{scal}(\operatorname{l2\text{-}space}(a, b, u)) \simeq \operatorname{normed\text{-}field\text{-}as\text{-}commutative\text{-}ring}(\operatorname{rr\text{-}normed\text{-}field})) \end{array}}}
\vnbprop{l2-space-vec}{\vnbstmt{\begin{array}{@{}l@{}} \forall a, b, u. \\ \quad \operatorname{is\text{-}measure}(a, b, u) \Rightarrow (\operatorname{vec}(\operatorname{l2\text{-}space}(a, b, u)) \simeq \operatorname{l2\text{-}set}(a, b, u)) \end{array}}}
\vnbprop{l2-space-vadd}{\vnbstmt{\begin{array}{@{}l@{}} \forall a, b, u. \\ \quad \operatorname{is\text{-}measure}(a, b, u) \Rightarrow (\operatorname{vadd}(\operatorname{l2\text{-}space}(a, b, u)) \simeq \operatorname{l2\text{-}vadd}(a, b, u)) \end{array}}}
\vnbprop{l2-space-vzero}{\vnbstmt{\begin{array}{@{}l@{}} \forall a, b, u. \\ \quad \operatorname{is\text{-}measure}(a, b, u) \Rightarrow (\operatorname{vzero}(\operatorname{l2\text{-}space}(a, b, u)) \simeq \operatorname{l2\text{-}vzero}(a)) \end{array}}}
\vnbprop{l2-space-vneg}{\vnbstmt{\begin{array}{@{}l@{}} \forall a, b, u. \\ \quad \operatorname{is\text{-}measure}(a, b, u) \Rightarrow (\operatorname{vneg}(\operatorname{l2\text{-}space}(a, b, u)) \simeq \operatorname{l2\text{-}vneg}(a, b, u)) \end{array}}}
\vnbprop{l2-space-act}{\vnbstmt{\begin{array}{@{}l@{}} \forall a, b, u. \\ \quad \operatorname{is\text{-}measure}(a, b, u) \Rightarrow (\operatorname{act}(\operatorname{l2\text{-}space}(a, b, u)) \simeq \operatorname{l2\text{-}act}(a, b, u)) \end{array}}}
\vnbprop{l2-space-rip}{\vnbstmt{\begin{array}{@{}l@{}} \forall a, b, u. \\ \quad \operatorname{is\text{-}measure}(a, b, u) \Rightarrow (\langle\cdot\mid\cdot\rangle_{\operatorname{l2\text{-}space}(a, b, u)} \simeq \operatorname{l2\text{-}rip}(a, b, u)) \end{array}}}
\vnbprop{l2-space-length}{\vnbstmt{\forall a, b, u.\; (\operatorname{is\text{-}measure}(a, b, u) \Rightarrow (\operatorname{length}(\operatorname{l2\text{-}space}(a, b, u)) = 7))}}
\vnbprop{l2-space-scal-ring}{\vnbstmt{\forall a, b, u.\; (\operatorname{is\text{-}measure}(a, b, u) \Rightarrow \operatorname{is\text{-}ring}(\operatorname{scal}(\operatorname{l2\text{-}space}(a, b, u))))}}
\vnbprop{l2-space-vec-set}{\vnbstmt{\forall a, b, u.\; (\operatorname{is\text{-}measure}(a, b, u) \Rightarrow (\operatorname{vec}(\operatorname{l2\text{-}space}(a, b, u)) \in \mathbf{Set}))}}
\vnbprop{l2-space-assoc}{\vnbstmt{\begin{array}{@{}l@{}} \forall a, b, u. \\ \quad \operatorname{is\text{-}measure}(a, b, u) \Rightarrow  \\ \quad \quad \operatorname{is\text{-}associative}(\operatorname{vadd}(\operatorname{l2\text{-}space}(a, b, u)), \operatorname{vec}(\operatorname{l2\text{-}space}(a, b, u))) \end{array}}}
\vnbprop{l2-space-comm}{\vnbstmt{\begin{array}{@{}l@{}} \forall a, b, u. \\ \quad \operatorname{is\text{-}measure}(a, b, u) \Rightarrow  \\ \quad \quad \operatorname{is\text{-}commutative}(\operatorname{vadd}(\operatorname{l2\text{-}space}(a, b, u)), \operatorname{vec}(\operatorname{l2\text{-}space}(a, b, u))) \end{array}}}
\vnbprop{l2-space-vadd-fun}{\vnbstmt{\begin{array}{@{}l@{}} \forall a, b, u. \\ \quad \operatorname{is\text{-}measure}(a, b, u) \Rightarrow  \\ \quad \quad \operatorname{vadd}(\operatorname{l2\text{-}space}(a, b, u)) \\ \quad \quad \in\; ((\operatorname{vec}(\operatorname{l2\text{-}space}(a, b, u)) \times \operatorname{vec}(\operatorname{l2\text{-}space}(a, b, u))) \to \operatorname{vec}(\operatorname{l2\text{-}space}(a, b, u))) \end{array}}}
\vnbprop{l2-space-vzero-in}{\vnbstmt{\begin{array}{@{}l@{}} \forall a, b, u. \\ \quad \operatorname{is\text{-}measure}(a, b, u) \Rightarrow (\operatorname{vzero}(\operatorname{l2\text{-}space}(a, b, u)) \in \operatorname{vec}(\operatorname{l2\text{-}space}(a, b, u))) \end{array}}}
\vnbprop{l2-space-vneg-fun}{\vnbstmt{\begin{array}{@{}l@{}} \forall a, b, u. \\ \quad \operatorname{is\text{-}measure}(a, b, u) \Rightarrow  \\ \quad \quad \operatorname{vneg}(\operatorname{l2\text{-}space}(a, b, u)) \in (\operatorname{vec}(\operatorname{l2\text{-}space}(a, b, u)) \to \operatorname{vec}(\operatorname{l2\text{-}space}(a, b, u))) \end{array}}}
\vnbprop{l2-space-act-fun}{\vnbstmt{\begin{array}{@{}l@{}} \forall a, b, u. \\ \quad \operatorname{is\text{-}measure}(a, b, u) \Rightarrow  \\ \quad \quad \operatorname{act}(\operatorname{l2\text{-}space}(a, b, u)) \\ \quad \quad \in\; ((\operatorname{carr}(\operatorname{scal}(\operatorname{l2\text{-}space}(a, b, u))) \times \operatorname{vec}(\operatorname{l2\text{-}space}(a, b, u))) \to \operatorname{vec}(\operatorname{l2\text{-}space}(a, b, u))) \end{array}}}
\vnbprop{l2-space-rip-fun}{\vnbstmt{\begin{array}{@{}l@{}} \forall a, b, u. \\ \quad \operatorname{is\text{-}measure}(a, b, u) \Rightarrow  \\ \quad \quad \langle\cdot\mid\cdot\rangle_{\operatorname{l2\text{-}space}(a, b, u)} \in ((\operatorname{vec}(\operatorname{l2\text{-}space}(a, b, u)) \times \operatorname{vec}(\operatorname{l2\text{-}space}(a, b, u))) \to \mathbb{R}) \end{array}}}
\vnbprop{l2-space-identity}{\vnbstmt{\begin{array}{@{}l@{}} \forall a, b, u. \\ \quad \operatorname{is\text{-}measure}(a, b, u) \Rightarrow  \\ \quad \quad \operatorname{is\text{-}identity}(\operatorname{vadd}(\operatorname{l2\text{-}space}(a, b, u)), \operatorname{vzero}(\operatorname{l2\text{-}space}(a, b, u)), \operatorname{vec}(\operatorname{l2\text{-}space}(a, b, u))) \end{array}}}
\vnbprop{l2-space-inverses}{\vnbstmt{\begin{array}{@{}l@{}} \forall a, b, u. \\ \quad \operatorname{is\text{-}measure}(a, b, u) \Rightarrow  \\ \quad \quad \operatorname{has\text{-}inverses}(\operatorname{vadd}(\operatorname{l2\text{-}space}(a, b, u)), \operatorname{vzero}(\operatorname{l2\text{-}space}(a, b, u)), \\ \quad \quad \quad \operatorname{vneg}(\operatorname{l2\text{-}space}(a, b, u)), \operatorname{vec}(\operatorname{l2\text{-}space}(a, b, u))) \end{array}}}
\vnbprop{l2-space-scal-pin}{\vnbstmt{\begin{array}{@{}l@{}} \forall a, b, u. \\ \quad \operatorname{is\text{-}measure}(a, b, u) \Rightarrow  \\ \quad \quad \operatorname{scal}(\operatorname{l2\text{-}space}(a, b, u)) = \operatorname{normed\text{-}field\text{-}as\text{-}commutative\text{-}ring}(\operatorname{rr\text{-}normed\text{-}field}) \end{array}}}
\vnbprop{l2-space-scal-carr}{\vnbstmt{\begin{array}{@{}l@{}} \forall a, b, u. \\ \quad \operatorname{is\text{-}measure}(a, b, u) \Rightarrow (\operatorname{carr}(\operatorname{scal}(\operatorname{l2\text{-}space}(a, b, u))) = \mathbb{R}) \end{array}}}
\vnbprop{l2-space-act-vadd}{\vnbstmt{\begin{array}{@{}l@{}} \forall a, b, u.; \forall r \in \operatorname{carr}(\operatorname{scal}(\operatorname{l2\text{-}space}(a, b, u))), x, y \in \operatorname{vec}(\operatorname{l2\text{-}space}(a, b, u)). \\ \quad \operatorname{is\text{-}measure}(a, b, u) \Rightarrow  \\ \quad \quad \operatorname{act}(\operatorname{l2\text{-}space}(a, b, u))(r, \operatorname{vadd}(\operatorname{l2\text{-}space}(a, b, u))(x, y)) \\ \quad \quad =\; \operatorname{vadd}(\operatorname{l2\text{-}space}(a, b, u))(\operatorname{act}(\operatorname{l2\text{-}space}(a, b, u))(r, x), \operatorname{act}(\operatorname{l2\text{-}space}(a, b, u))(r, y)) \end{array}}}
\vnbprop{l2-space-act-add}{\vnbstmt{\begin{array}{@{}l@{}} \forall a, b, u.; \forall r, s \in \operatorname{carr}(\operatorname{scal}(\operatorname{l2\text{-}space}(a, b, u))), x \in \operatorname{vec}(\operatorname{l2\text{-}space}(a, b, u)). \\ \quad \operatorname{is\text{-}measure}(a, b, u) \Rightarrow  \\ \quad \quad \operatorname{act}(\operatorname{l2\text{-}space}(a, b, u))(\operatorname{add}(\operatorname{scal}(\operatorname{l2\text{-}space}(a, b, u)))(r, s), x) \\ \quad \quad =\; \operatorname{vadd}(\operatorname{l2\text{-}space}(a, b, u))(\operatorname{act}(\operatorname{l2\text{-}space}(a, b, u))(r, x), \operatorname{act}(\operatorname{l2\text{-}space}(a, b, u))(s, x)) \end{array}}}
\vnbprop{l2-space-act-mul}{\vnbstmt{\begin{array}{@{}l@{}} \forall a, b, u.; \forall r, s \in \operatorname{carr}(\operatorname{scal}(\operatorname{l2\text{-}space}(a, b, u))), x \in \operatorname{vec}(\operatorname{l2\text{-}space}(a, b, u)). \\ \quad \operatorname{is\text{-}measure}(a, b, u) \Rightarrow  \\ \quad \quad \operatorname{act}(\operatorname{l2\text{-}space}(a, b, u))(\operatorname{mul}(\operatorname{scal}(\operatorname{l2\text{-}space}(a, b, u)))(r, s), x) \\ \quad \quad =\; \operatorname{act}(\operatorname{l2\text{-}space}(a, b, u))(r, \operatorname{act}(\operatorname{l2\text{-}space}(a, b, u))(s, x)) \end{array}}}
\vnbprop{l2-space-act-one}{\vnbstmt{\begin{array}{@{}l@{}} \forall a, b, u.; \forall x \in \operatorname{vec}(\operatorname{l2\text{-}space}(a, b, u)). \\ \quad \operatorname{is\text{-}measure}(a, b, u) \Rightarrow  \\ \quad \quad \operatorname{act}(\operatorname{l2\text{-}space}(a, b, u))(\operatorname{one}(\operatorname{scal}(\operatorname{l2\text{-}space}(a, b, u))), x) = x \end{array}}}
\vnbprop{l2-space-rip-sym}{\vnbstmt{\begin{array}{@{}l@{}} \forall a, b, u.; \forall x, y \in \operatorname{vec}(\operatorname{l2\text{-}space}(a, b, u)). \\ \quad \operatorname{is\text{-}measure}(a, b, u) \Rightarrow  \\ \quad \quad \langle\cdot\mid\cdot\rangle_{\operatorname{l2\text{-}space}(a, b, u)}(x, y) = \langle\cdot\mid\cdot\rangle_{\operatorname{l2\text{-}space}(a, b, u)}(y, x) \end{array}}}
\vnbprop{l2-space-rip-add-left}{\vnbstmt{\begin{array}{@{}l@{}} \forall a, b, u.; \forall x, y, z \in \operatorname{vec}(\operatorname{l2\text{-}space}(a, b, u)). \\ \quad \operatorname{is\text{-}measure}(a, b, u) \Rightarrow  \\ \quad \quad \langle\cdot\mid\cdot\rangle_{\operatorname{l2\text{-}space}(a, b, u)}(\operatorname{vadd}(\operatorname{l2\text{-}space}(a, b, u))(x, y), z) \\ \quad \quad =\; \langle\cdot\mid\cdot\rangle_{\operatorname{l2\text{-}space}(a, b, u)}(x, z) + \langle\cdot\mid\cdot\rangle_{\operatorname{l2\text{-}space}(a, b, u)}(y, z) \end{array}}}
\vnbprop{l2-space-rip-homog-left}{\vnbstmt{\begin{array}{@{}l@{}} \forall a, b, u.; \forall \alpha \in \mathbb{R}, x, y \in \operatorname{vec}(\operatorname{l2\text{-}space}(a, b, u)). \\ \quad \operatorname{is\text{-}measure}(a, b, u) \Rightarrow  \\ \quad \quad \langle\cdot\mid\cdot\rangle_{\operatorname{l2\text{-}space}(a, b, u)}(\operatorname{act}(\operatorname{l2\text{-}space}(a, b, u))(\alpha, x), y) \\ \quad \quad =\; \alpha \cdot \langle\cdot\mid\cdot\rangle_{\operatorname{l2\text{-}space}(a, b, u)}(x, y) \end{array}}}
\vnbprop{l2-space-rip-self-nonneg}{\vnbstmt{\begin{array}{@{}l@{}} \forall a, b, u.; \forall x \in \operatorname{vec}(\operatorname{l2\text{-}space}(a, b, u)). \\ \quad \operatorname{is\text{-}measure}(a, b, u) \Rightarrow (0 \leq \langle\cdot\mid\cdot\rangle_{\operatorname{l2\text{-}space}(a, b, u)}(x, x)) \end{array}}}
\vnbprop{l2-space-is-sip}{\vnbstmt{\begin{array}{@{}l@{}} \forall a, b, u. \\ \quad \operatorname{is\text{-}measure}(a, b, u) \Rightarrow \operatorname{is\text{-}semi\text{-}inner\text{-}product\text{-}space}(\operatorname{l2\text{-}space}(a, b, u)) \end{array}}}
\vnbprop{rintegral-zero-fn}{\vnbstmt{\begin{array}{@{}l@{}} \forall a, b, u. \\ \quad \operatorname{is\text{-}measure}(a, b, u) \Rightarrow (\operatorname{rintegral}(a, b, u, (\lambda x \in a.\; 0)) = 0) \end{array}}}
\vnbprop{l2-null-iff-ae}{\vnbstmt{\begin{array}{@{}l@{}} \forall a, b, u.; \forall f, g \in \operatorname{l2\text{-}set}(a, b, u). \\ \quad \operatorname{is\text{-}measure}(a, b, u) \Rightarrow  \\ \quad \quad \langle\cdot\mid\cdot\rangle_{\operatorname{l2\text{-}space}(a, b, u)}(\operatorname{vadd}(\operatorname{l2\text{-}space}(a, b, u))(f, \operatorname{vneg}(\operatorname{l2\text{-}space}(a, b, u))(g)), \\ \quad \quad \quad \operatorname{vadd}(\operatorname{l2\text{-}space}(a, b, u))(f, \operatorname{vneg}(\operatorname{l2\text{-}space}(a, b, u))(g))) \\ \quad \quad =\; 0 \\ \quad \quad \Leftrightarrow\; \exists n. \\ \quad \quad \quad \operatorname{is\text{-}null\text{-}set}(a, b, u, n) \wedge \forall x \in a.\; (\neg (x \in n) \Rightarrow (f(x) = g(x))) \end{array}}}
\vnbprop{l2-null-rel-is-ae-rel}{\vnbstmt{\begin{array}{@{}l@{}} \forall a, b, u. \\ \quad \operatorname{is\text{-}measure}(a, b, u) \Rightarrow  \\ \quad \quad \operatorname{sip\text{-}rel}(\operatorname{l2\text{-}space}(a, b, u)) = \operatorname{ae\text{-}rel}(a, b, u, \operatorname{l2\text{-}set}(a, b, u)) \end{array}}}
\vnbprop{l2-set-subset-fun}{\vnbstmt{\forall a, b, u.\; (\operatorname{is\text{-}measure}(a, b, u) \Rightarrow (\operatorname{l2\text{-}set}(a, b, u) \subseteq (a \to \mathbb{R})))}}
\vnbprop{l2-setoid-is-ae-setoid}{\vnbstmt{\begin{array}{@{}l@{}} \forall a, b, u. \\ \quad \operatorname{is\text{-}measure}(a, b, u) \Rightarrow  \\ \quad \quad \operatorname{sip\text{-}setoid}(\operatorname{l2\text{-}space}(a, b, u)) = \operatorname{ae\text{-}setoid}(a, b, u, \operatorname{l2\text{-}set}(a, b, u)) \end{array}}}
\vnbprop{l2-quotient-is-ae-quotient}{\vnbstmt{\begin{array}{@{}l@{}} \forall a, b, u. \\ \quad \operatorname{is\text{-}measure}(a, b, u) \Rightarrow  \\ \quad \quad \operatorname{quotient}(\operatorname{sip\text{-}setoid}(\operatorname{l2\text{-}space}(a, b, u))) = \operatorname{quotient}(\operatorname{ae\text{-}setoid}(a, b, u, \operatorname{l2\text{-}set}(a, b, u))) \end{array}}}
\vnbprop{l2-quotient-is-pre-hilbert}{\vnbstmt{\begin{array}{@{}l@{}} \forall a, b, u. \\ \quad \operatorname{is\text{-}measure}(a, b, u) \Rightarrow  \\ \quad \quad \operatorname{is\text{-}pre\text{-}hilbert\text{-}space}(\operatorname{sip\text{-}quotient}(\operatorname{l2\text{-}space}(a, b, u))) \end{array}}}
\vnbprop{cl2-space-is-csip}{\vnbstmt{\begin{array}{@{}l@{}} \forall a, b, u. \\ \quad \operatorname{is\text{-}measure}(a, b, u) \Rightarrow  \\ \quad \quad \operatorname{is\text{-}complex\text{-}semi\text{-}inner\text{-}product\text{-}space}(\operatorname{cl2\text{-}space}(a, b, u)) \end{array}}}
\vnbprop{cl2-quotient-is-cphs}{\vnbstmt{\begin{array}{@{}l@{}} \forall a, b, u. \\ \quad \operatorname{is\text{-}measure}(a, b, u) \Rightarrow  \\ \quad \quad \operatorname{is\text{-}complex\text{-}pre\text{-}hilbert\text{-}space}(\operatorname{csip\text{-}quotient}(\operatorname{cl2\text{-}space}(a, b, u))) \end{array}}}
\vnbprop{cl2-quotient-is-cips}{\vnbstmt{\begin{array}{@{}l@{}} \forall a, b, u. \\ \quad \operatorname{is\text{-}measure}(a, b, u) \Rightarrow  \\ \quad \quad \operatorname{is\text{-}complex\text{-}inner\text{-}product\text{-}space}(\operatorname{csip\text{-}cips}(\operatorname{csip\text{-}quotient}(\operatorname{cl2\text{-}space}(a, b, u)))) \end{array}}}

\vnbbookfile{theorem-library/lp-step-dense}{theorem-library/lp-step-dense}
\noindent{\small lp-step-dense.scm -- WAVE 2 (2026-10-27..30), wired 2026-10-30 from scratchpad/md/md-step.scm. Library: measure.}\par
\vnbprop{simple-fn-real-measurable}{\vnbstmt{\forall a, b, u.\; (\operatorname{is\text{-}simple\text{-}fn}(a, b, u) \Rightarrow \operatorname{is\text{-}real\text{-}measurable\text{-}fn}(a, b, u))}}
\vnbprop{real-valued-measurable-fn-real}{\vnbstmt{\begin{array}{@{}l@{}} \forall a, b, u. \\ \quad (\operatorname{is\text{-}measurable\text{-}fn}(a, b, u) \wedge \forall x \in a.\; (u(x) \in \mathbb{R})) \Rightarrow  \\ \quad \quad \operatorname{is\text{-}real\text{-}measurable\text{-}fn}(a, b, u) \end{array}}}
\vnbprop{nondecreasing-converges-to-esup}{\vnbstmt{\begin{array}{@{}l@{}} \forall a, l. \\ \quad a \in (\mathbb{N} \to \mathbb{R}) \wedge \\ \quad \forall k \in \mathbb{N}.\; (0 \leq a(k)) \wedge \\ \quad \forall k \in \mathbb{N}.\; (a(k) \leq a((k + 1))) \wedge \\ \quad l \in \mathbb{R} \wedge \\ \quad \operatorname{esup}(\operatorname{image}(a, \mathbb{N})) = l \Rightarrow \\ \quad \quad \operatorname{converges\text{-}to}(\operatorname{rr\text{-}ms}, a, l) \end{array}}}
\vnbprop{simple-difference-step-fn}{\vnbstmt{\begin{array}{@{}l@{}} \forall a, b, u, v. \\ \quad (\operatorname{is\text{-}simple\text{-}fn}(a, b, u) \wedge \operatorname{is\text{-}simple\text{-}fn}(a, b, v)) \Rightarrow  \\ \quad \quad \operatorname{is\text{-}measurable\text{-}step\text{-}fn}(a, b, (\lambda x \in a.\; (u(x) - v(x)))) \end{array}}}
\vnbprop{pos-neg-parts-ptwise}{\vnbstmt{\begin{array}{@{}l@{}} \forall a, f, x. \\ \quad ((f \in (a \to \mathbb{R})) \wedge (x \in a)) \Rightarrow  \\ \quad \quad \operatorname{pos\text{-}part}(a, f)(x) \in \mathbb{R} \wedge \\ \quad \quad 0 \leq \operatorname{pos\text{-}part}(a, f)(x) \wedge \\ \quad \quad \operatorname{pos\text{-}part}(a, f)(x) \leq \lvert f(x) \rvert \wedge \\ \quad \quad \operatorname{neg\text{-}part}(a, f)(x) \in \mathbb{R} \wedge \\ \quad \quad 0 \leq \operatorname{neg\text{-}part}(a, f)(x) \wedge \\ \quad \quad \operatorname{neg\text{-}part}(a, f)(x) \leq \lvert f(x) \rvert \wedge \\ \quad \quad (\operatorname{pos\text{-}part}(a, f)(x) - \operatorname{neg\text{-}part}(a, f)(x)) = f(x) \end{array}}}
\vnbprop{approximant-below}{\vnbstmt{\begin{array}{@{}l@{}} \forall a, f, h, n, x. \\ \quad f \in (\mathbb{N} \to (a \to \operatorname{rr\text{-}pos\text{-}star})) \wedge \\ \quad \forall b \in a.\; (\operatorname{ptwise\text{-}sup}(a, f)(b) = h(b)) \wedge \\ \quad n \in \mathbb{N} \wedge \\ \quad x \in a \Rightarrow \\ \quad \quad f(n)(x) \leq h(x) \end{array}}}
\vnbprop{approximant-converges}{\vnbstmt{\begin{array}{@{}l@{}} \forall a, f, h, x. \\ \quad f \in (\mathbb{N} \to (a \to \operatorname{rr\text{-}pos\text{-}star})) \wedge \\ \quad \forall n \in \mathbb{N}, b \in a.\; (f(n)(b) \in \mathbb{R}) \wedge \\ \quad \forall n \in \mathbb{N}.\; \operatorname{ptwise\text{-}le}(a, f(n), f((n + 1))) \wedge \\ \quad \forall b \in a.\; (\operatorname{ptwise\text{-}sup}(a, f)(b) = h(b)) \wedge \\ \quad x \in a \wedge \\ \quad h(x) \in \mathbb{R} \Rightarrow \\ \quad \quad \operatorname{converges\text{-}to}(\operatorname{rr\text{-}ms}, (\lambda n \in \mathbb{N}.\; f(n)(x)), h(x)) \end{array}}}
\vnbprop{abs-fn-real-measurable}{\vnbstmt{\begin{array}{@{}l@{}} \forall a, b, u, f. \\ \quad (\operatorname{is\text{-}measure}(a, b, u) \wedge \operatorname{is\text{-}real\text{-}measurable\text{-}fn}(a, b, f)) \Rightarrow  \\ \quad \quad \operatorname{is\text{-}real\text{-}measurable\text{-}fn}(a, b, \operatorname{abs\text{-}fn}(a, f)) \end{array}}}
\vnbprop{step-functions-dense-lp}{\vnbstmt{\begin{array}{@{}l@{}} \forall a, b, u, p, f. \\ \quad \operatorname{is\text{-}measure}(a, b, u) \wedge \\ \quad p \in \mathbb{R} \wedge \\ \quad 1 \leq p \wedge \\ \quad \operatorname{is\text{-}real\text{-}measurable\text{-}fn}(a, b, f) \wedge \\ \quad \|{f}\|_{p} \in \mathbb{R} \Rightarrow \\ \quad \quad \exists s \in (\mathbb{N} \to (a \to \mathbb{R})). \\ \quad \quad \quad \forall k \in \mathbb{N}. \\ \quad \quad \quad \quad \operatorname{is\text{-}measurable\text{-}step\text{-}fn}(a, b, s(k)) \wedge (\|{s(k)}\|_{p} \in \mathbb{R}) \wedge \\ \quad \quad \quad \operatorname{econverges\text{-}to}((\lambda k \in \mathbb{N}.\; \|{(\lambda x \in a.\; (s(k)(x) - f(x)))}\|_{p}), 0) \end{array}}}

\vnbbookfile{theorem-library/lp-complete}{theorem-library/lp-complete}
\noindent{\small lp-complete.scm -- WAVE 2 (2026-10-27..30), wired 2026-10-30 from scratchpad/md/me-rf.scm. Library: measure.}\par
\vnbprop{etail-unfold}{\vnbstmt{\forall f, n.\; (\operatorname{etail}(f, n) \simeq \operatorname{image}(f, \{p \in \mathbb{N} : (n \leq p)\}))}}
\vnbprop{rr-converges-nonneg-eliminf}{\vnbstmt{\begin{array}{@{}l@{}} \forall u, l. \\ \quad ((u \in (\mathbb{N} \to \mathbb{R})) \wedge (\forall k \in \mathbb{N}.\; (0 \leq u(k)) \wedge ((l \in \mathbb{R}) \wedge \operatorname{converges\text{-}to}(\operatorname{rr\text{-}ms}, u, l)))) \Rightarrow  \\ \quad \quad (0 \leq l) \wedge (\operatorname{eliminf}(u) = l) \end{array}}}
\vnbprop{finite-part-real-measurable}{\vnbstmt{\begin{array}{@{}l@{}} \forall a, b, h. \\ \quad \operatorname{is\text{-}measurable\text{-}fn}(a, b, h) \Rightarrow  \\ \quad \quad \operatorname{is\text{-}real\text{-}measurable\text{-}fn}(a, b, \operatorname{ptwise\text{-}etimes}(a, \operatorname{indicator}(a, \{p \in a : (h(p) \in \mathbb{R})\}), h)) \wedge \\ \quad \quad \forall x \in a. \\ \quad \quad \quad (h(x) \in \mathbb{R}) \Rightarrow  \\ \quad \quad \quad \quad \operatorname{ptwise\text{-}etimes}(a, \operatorname{indicator}(a, \{p \in a : (h(p) \in \mathbb{R})\}), h)(x) = h(x) \wedge \\ \quad \quad \forall x \in a. \\ \quad \quad \quad \neg (h(x) \in \mathbb{R}) \Rightarrow  \\ \quad \quad \quad \quad \operatorname{ptwise\text{-}etimes}(a, \operatorname{indicator}(a, \{p \in a : (h(p) \in \mathbb{R})\}), h)(x) = 0 \end{array}}}
\vnbprop{parts-lipschitz}{\vnbstmt{\begin{array}{@{}l@{}} \forall a, f, g, x. \\ \quad ((f \in (a \to \mathbb{R})) \wedge ((g \in (a \to \mathbb{R})) \wedge (x \in a))) \Rightarrow  \\ \quad \quad \lvert (\operatorname{pos\text{-}part}(a, f)(x) - \operatorname{pos\text{-}part}(a, g)(x)) \rvert \leq \lvert (f(x) - g(x)) \rvert \wedge \\ \quad \quad \lvert (\operatorname{neg\text{-}part}(a, f)(x) - \operatorname{neg\text{-}part}(a, g)(x)) \rvert \leq \lvert (f(x) - g(x)) \rvert \end{array}}}
\vnbprop{partial-sum-telescope}{\vnbstmt{\begin{array}{@{}l@{}} \forall a, s, x.; \forall m, n \in \mathbb{N}. \\ \quad ((s \in (\mathbb{N} \to (a \to \mathbb{R}))) \wedge ((x \in a) \wedge (n \leq m))) \Rightarrow  \\ \quad \quad \lvert (s(m)(x) - s(n)(x)) \rvert \\ \quad \quad \leq\; (\operatorname{series\text{-}partial\text{-}sum}((\lambda k \in \mathbb{N}.\; \lvert (s((k + 1))(x) - s(k)(x)) \rvert), m) - \operatorname{series\text{-}partial\text{-}sum}((\lambda k \in \mathbb{N}.\; \lvert (s((k + 1))(x) - s(k)(x)) \rvert), n)) \end{array}}}
\vnbprop{partial-sum-real}{\vnbstmt{\begin{array}{@{}l@{}} \forall a.; \forall s \in (\mathbb{N} \to (a \to \mathbb{R})), x \in a. \\ \quad (\lambda k \in \mathbb{N}.\; \lvert (s((k + 1))(x) - s(k)(x)) \rvert) \in (\mathbb{N} \to \mathbb{R}) \wedge \\ \quad \forall n \in \mathbb{N}. \\ \quad \quad \operatorname{series\text{-}partial\text{-}sum}((\lambda k \in \mathbb{N}.\; \lvert (s((k + 1))(x) - s(k)(x)) \rvert), n) \in \mathbb{R} \wedge \\ \quad \quad 0 \leq \operatorname{series\text{-}partial\text{-}sum}((\lambda k \in \mathbb{N}.\; \lvert (s((k + 1))(x) - s(k)(x)) \rvert), n) \end{array}}}
\vnbprop{partial-sum-fn-measurable}{\vnbstmt{\begin{array}{@{}l@{}} \forall a, b, u, s.; \forall n \in \mathbb{N}. \\ \quad \operatorname{is\text{-}measure}(a, b, u) \wedge \\ \quad s \in (\mathbb{N} \to (a \to \mathbb{R})) \wedge \\ \quad \forall k \in \mathbb{N}.\; \operatorname{is\text{-}real\text{-}measurable\text{-}fn}(a, b, s(k)) \Rightarrow \\ \quad \quad \operatorname{is\text{-}real\text{-}measurable\text{-}fn}(a, b, \\ \quad \quad \quad (\lambda x \in a.\; \operatorname{series\text{-}partial\text{-}sum}((\lambda \mathit{rdk} \in \mathbb{N}.\; \lvert (s((\mathit{rdk} + 1))(x) - s(\mathit{rdk})(x)) \rvert), n))) \end{array}}}
\vnbprop{partial-sum-lp-bound}{\vnbstmt{\begin{array}{@{}l@{}} \forall a, b, u, s, p, d.; \forall n \in \mathbb{N}. \\ \quad \operatorname{is\text{-}measure}(a, b, u) \wedge \\ \quad s \in (\mathbb{N} \to (a \to \mathbb{R})) \wedge \\ \quad \forall k \in \mathbb{N}.\; \operatorname{is\text{-}real\text{-}measurable\text{-}fn}(a, b, s(k)) \wedge \\ \quad p \in \mathbb{R} \wedge \\ \quad 1 \leq p \wedge \\ \quad d \in (\mathbb{N} \to \mathbb{R}) \wedge \\ \quad \forall k \in \mathbb{N}. \\ \quad \quad \|{(\lambda x \in a.\; (s((k + 1))(x) - s(k)(x)))}\|_{p} \leq d(k) \Rightarrow \\ \quad \quad \|{(\lambda \mathit{rdx} \in a.\; \operatorname{series\text{-}partial\text{-}sum}((\lambda \mathit{rdk} \in \mathbb{N}.\; \lvert (s((\mathit{rdk} + 1))(\mathit{rdx}) - s(\mathit{rdk})(\mathit{rdx})) \rvert), n))}\|_{p} \\ \quad \quad \leq\; \operatorname{series\text{-}partial\text{-}sum}(d, n) \end{array}}}
\vnbprop{partial-sum-bounded-ae}{\vnbstmt{\begin{array}{@{}l@{}} \forall a, b, u, s, p, d, c. \\ \quad \operatorname{is\text{-}measure}(a, b, u) \wedge \\ \quad s \in (\mathbb{N} \to (a \to \mathbb{R})) \wedge \\ \quad \forall k \in \mathbb{N}.\; \operatorname{is\text{-}real\text{-}measurable\text{-}fn}(a, b, s(k)) \wedge \\ \quad p \in \mathbb{R} \wedge \\ \quad 1 \leq p \wedge \\ \quad d \in (\mathbb{N} \to \mathbb{R}) \wedge \\ \quad \forall k \in \mathbb{N}. \\ \quad \quad \|{(\lambda \mathit{sx} \in a.\; (s((k + 1))(\mathit{sx}) - s(k)(\mathit{sx})))}\|_{p} \leq d(k) \wedge \\ \quad c \in \mathbb{R} \wedge \\ \quad \forall n \in \mathbb{N}.\; (\operatorname{series\text{-}partial\text{-}sum}(d, n) \leq c) \Rightarrow \\ \quad \quad \operatorname{is\text{-}measurable\text{-}fn}(a, b, \\ \quad \quad \quad \operatorname{ptwise\text{-}sup}(a, \\ \quad \quad \quad \quad (\lambda \mathit{rhn} \in \mathbb{N}.\; \\ \quad \quad \quad \quad \quad \operatorname{ptwise\text{-}epow}(a, \\ \quad \quad \quad \quad \quad \quad \operatorname{abs\text{-}fn}(a, \\ \quad \quad \quad \quad \quad \quad \quad (\lambda \mathit{rdx} \in a.\; \\ \quad \quad \quad \quad \quad \quad \quad \quad \operatorname{series\text{-}partial\text{-}sum}((\lambda \mathit{rdk} \in \mathbb{N}.\; \lvert (s((\mathit{rdk} + 1))(\mathit{rdx}) - s(\mathit{rdk})(\mathit{rdx})) \rvert), \mathit{rhn}))), \\ \quad \quad \quad \quad \quad \quad p)))) \wedge \\ \quad \quad \operatorname{integral}(a, b, u, \\ \quad \quad \quad \operatorname{ptwise\text{-}sup}(a, \\ \quad \quad \quad \quad (\lambda \mathit{rhn} \in \mathbb{N}.\; \\ \quad \quad \quad \quad \quad \operatorname{ptwise\text{-}epow}(a, \\ \quad \quad \quad \quad \quad \quad \operatorname{abs\text{-}fn}(a, \\ \quad \quad \quad \quad \quad \quad \quad (\lambda \mathit{rdx} \in a.\; \\ \quad \quad \quad \quad \quad \quad \quad \quad \operatorname{series\text{-}partial\text{-}sum}((\lambda \mathit{rdk} \in \mathbb{N}.\; \lvert (s((\mathit{rdk} + 1))(\mathit{rdx}) - s(\mathit{rdk})(\mathit{rdx})) \rvert), \mathit{rhn}))), \\ \quad \quad \quad \quad \quad \quad p)))) \\ \quad \quad \in\; \mathbb{R} \wedge \\ \quad \quad \{f \in a : (\operatorname{ptwise\text{-}sup}(a, (\lambda \mathit{rhn} \in \mathbb{N}.\; \operatorname{ptwise\text{-}epow}(a, \operatorname{abs\text{-}fn}(a, (\lambda \mathit{rdx} \in a.\; \operatorname{series\text{-}partial\text{-}sum}((\lambda \mathit{rdk} \in \mathbb{N}.\; \lvert (s((\mathit{rdk} + 1))(\mathit{rdx}) - s(\mathit{rdk})(\mathit{rdx})) \rvert), \mathit{rhn}))), p)))(f) = \operatorname{pos\text{-}inf})\} \\ \quad \quad \in\; b \wedge \\ \quad \quad u(\{f \in a : (\operatorname{ptwise\text{-}sup}(a, (\lambda \mathit{rhn} \in \mathbb{N}.\; \operatorname{ptwise\text{-}epow}(a, \operatorname{abs\text{-}fn}(a, (\lambda \mathit{rdx} \in a.\; \operatorname{series\text{-}partial\text{-}sum}((\lambda \mathit{rdk} \in \mathbb{N}.\; \lvert (s((\mathit{rdk} + 1))(\mathit{rdx}) - s(\mathit{rdk})(\mathit{rdx})) \rvert), \mathit{rhn}))), p)))(f) = \operatorname{pos\text{-}inf})\}) \\ \quad \quad =\; 0 \wedge \\ \quad \quad \forall x \in a. \\ \quad \quad \quad \neg (x \in \{f \in a : (\operatorname{ptwise\text{-}sup}(a, (\lambda \mathit{rhn} \in \mathbb{N}.\; \operatorname{ptwise\text{-}epow}(a, \operatorname{abs\text{-}fn}(a, (\lambda \mathit{rdx} \in a.\; \operatorname{series\text{-}partial\text{-}sum}((\lambda \mathit{rdk} \in \mathbb{N}.\; \lvert (s((\mathit{rdk} + 1))(\mathit{rdx}) - s(\mathit{rdk})(\mathit{rdx})) \rvert), \mathit{rhn}))), p)))(f) = \operatorname{pos\text{-}inf})\}) \Rightarrow \\ \quad \quad \quad \quad \operatorname{ptwise\text{-}sup}(a, (\lambda \mathit{rhn} \in \mathbb{N}.\; \operatorname{ptwise\text{-}epow}(a, \operatorname{abs\text{-}fn}(a, (\lambda \mathit{rdx} \in a.\; \operatorname{series\text{-}partial\text{-}sum}((\lambda \mathit{rdk} \in \mathbb{N}.\; \lvert (s((\mathit{rdk} + 1))(\mathit{rdx}) - s(\mathit{rdk})(\mathit{rdx})) \rvert), \mathit{rhn}))), p)))(x) \\ \quad \quad \quad \quad \in\; \mathbb{R} \wedge \\ \quad \quad \quad \quad \forall n \in \mathbb{N}. \\ \quad \quad \quad \quad \quad \operatorname{series\text{-}partial\text{-}sum}((\lambda \mathit{rdk} \in \mathbb{N}.\; \lvert (s((\mathit{rdk} + 1))(x) - s(\mathit{rdk})(x)) \rvert), n) \\ \quad \quad \quad \quad \quad \leq\; {\operatorname{ptwise\text{-}sup}(a, (\lambda \mathit{rhn} \in \mathbb{N}.\; \operatorname{ptwise\text{-}epow}(a, \operatorname{abs\text{-}fn}(a, (\lambda \mathit{rdx} \in a.\; \operatorname{series\text{-}partial\text{-}sum}((\lambda \mathit{rdk} \in \mathbb{N}.\; \lvert (s((\mathit{rdk} + 1))(\mathit{rdx}) - s(\mathit{rdk})(\mathit{rdx})) \rvert), \mathit{rhn}))), p)))(x)}^{\operatorname{recip}(p)} \end{array}}}
\vnbprop{bounded-nondecreasing-converges}{\vnbstmt{\begin{array}{@{}l@{}} \forall a, b. \\ \quad a \in (\mathbb{N} \to \mathbb{R}) \wedge \\ \quad \forall k \in \mathbb{N}.\; (0 \leq a(k)) \wedge \\ \quad \forall k \in \mathbb{N}.\; (a(k) \leq a((k + 1))) \wedge \\ \quad b \in \mathbb{R} \wedge \\ \quad \forall k \in \mathbb{N}.\; (a(k) \leq b) \Rightarrow \\ \quad \quad (\operatorname{esup}(\operatorname{image}(a, \mathbb{N})) \in \mathbb{R}) \wedge \operatorname{converges\text{-}to}(\operatorname{rr\text{-}ms}, a, \operatorname{esup}(\operatorname{image}(a, \mathbb{N}))) \end{array}}}
\vnbprop{bounded-partial-sums-cauchy}{\vnbstmt{\begin{array}{@{}l@{}} \forall a, s, b, c, d. \\ \quad s \in (\mathbb{N} \to (a \to \mathbb{R})) \wedge \\ \quad b \in a \wedge \\ \quad c \in \mathbb{R} \wedge \\ \quad \forall n \in \mathbb{N}. \\ \quad \quad \operatorname{series\text{-}partial\text{-}sum}((\lambda k \in \mathbb{N}.\; \lvert (s((k + 1))(b) - s(k)(b)) \rvert), n) \leq c \wedge \\ \quad \operatorname{pos\text{-}rr}(d) \Rightarrow \\ \quad \quad \exists f \in \mathbb{N}. \forall m, q \in \mathbb{N}. \\ \quad \quad \quad ((f \leq m) \wedge (f \leq q)) \Rightarrow (\lvert (s(m)(b) - s(q)(b)) \rvert \leq d) \end{array}}}
\vnbprop{ptwise-liminf-apply}{\vnbstmt{\begin{array}{@{}l@{}} \forall a, f.; \forall x \in a. \\ \quad (\operatorname{ptwise\text{-}liminf}(a, f)(x) \simeq \operatorname{eliminf}((\lambda n \in \mathbb{N}.\; f(n)(x)))) \end{array}}}
\vnbprop{cauchy-point-limit}{\vnbstmt{\begin{array}{@{}l@{}} \forall a, s, b. \\ \quad s \in (\mathbb{N} \to (a \to \mathbb{R})) \wedge \\ \quad b \in a \wedge \\ \quad \forall c. \\ \quad \quad \operatorname{pos\text{-}rr}(c) \Rightarrow  \\ \quad \quad \quad \exists d \in \mathbb{N}. \forall m, q \in \mathbb{N}. \\ \quad \quad \quad \quad ((d \leq m) \wedge (d \leq q)) \Rightarrow (\lvert (s(m)(b) - s(q)(b)) \rvert \leq c) \Rightarrow \\ \quad \quad \operatorname{ptwise\text{-}liminf}(a, (\lambda \mathit{rpk} \in \mathbb{N}.\; \operatorname{pos\text{-}part}(a, s(\mathit{rpk}))))(b) \in \mathbb{R} \wedge \\ \quad \quad \operatorname{ptwise\text{-}liminf}(a, (\lambda \mathit{rqk} \in \mathbb{N}.\; \operatorname{neg\text{-}part}(a, s(\mathit{rqk}))))(b) \in \mathbb{R} \wedge \\ \quad \quad \operatorname{converges\text{-}to}(\operatorname{rr\text{-}ms}, (\lambda k \in \mathbb{N}.\; s(k)(b)), \\ \quad \quad \quad (\operatorname{ptwise\text{-}liminf}(a, (\lambda \mathit{rpk} \in \mathbb{N}.\; \operatorname{pos\text{-}part}(a, s(\mathit{rpk}))))(b) - \operatorname{ptwise\text{-}liminf}(a, (\lambda \mathit{rqk} \in \mathbb{N}.\; \operatorname{neg\text{-}part}(a, s(\mathit{rqk}))))(b))) \end{array}}}
\vnbprop{lp-rapid-sequence-converges}{\vnbstmt{\begin{array}{@{}l@{}} \forall a, b, u, s, p, d, c. \\ \quad \operatorname{is\text{-}measure}(a, b, u) \wedge \\ \quad s \in (\mathbb{N} \to (a \to \mathbb{R})) \wedge \\ \quad \forall k \in \mathbb{N}.\; \operatorname{is\text{-}real\text{-}measurable\text{-}fn}(a, b, s(k)) \wedge \\ \quad p \in \mathbb{R} \wedge \\ \quad 1 \leq p \wedge \\ \quad d \in (\mathbb{N} \to \mathbb{R}) \wedge \\ \quad \forall k \in \mathbb{N}. \\ \quad \quad \|{(\lambda x \in a.\; (s((k + 1))(x) - s(k)(x)))}\|_{p} \leq d(k) \wedge \\ \quad c \in \mathbb{R} \wedge \\ \quad \forall n \in \mathbb{N}.\; (\operatorname{series\text{-}partial\text{-}sum}(d, n) \leq c) \wedge \\ \quad \|{s(0)}\|_{p} \in \mathbb{R} \Rightarrow \\ \quad \quad \exists f. \\ \quad \quad \quad \operatorname{is\text{-}real\text{-}measurable\text{-}fn}(a, b, f) \wedge \\ \quad \quad \quad \|{f}\|_{p} \in \mathbb{R} \wedge \\ \quad \quad \quad \operatorname{econverges\text{-}to}((\lambda k \in \mathbb{N}.\; \|{(\lambda x \in a.\; (s(k)(x) - f(x)))}\|_{p}), 0) \end{array}}}
\vnbprop{half-geometric-facts}{\vnbstmt{\begin{array}{@{}l@{}} (\lambda n \in \mathbb{N}.\; \operatorname{power}(1/2, n)) \in (\mathbb{N} \to \mathbb{R}) \wedge \\ \forall k \in \mathbb{N}. \\ \quad ((\lambda n \in \mathbb{N}.\; \operatorname{power}(1/2, n))(k) = \operatorname{power}(1/2, k)) \wedge (0 < \operatorname{power}(1/2, k)) \wedge \\ \forall k \in \mathbb{N}.\; (\operatorname{power}(1/2, (k + 1)) \leq \operatorname{power}(1/2, k)) \wedge \\ \forall q \in \mathbb{N}.\; (\operatorname{series\text{-}partial\text{-}sum}((\lambda n \in \mathbb{N}.\; \operatorname{power}(1/2, n)), q) \leq 2) \end{array}}}
\vnbprop{lp-norm-sub-le}{\vnbstmt{\begin{array}{@{}l@{}} \forall a, b, c, p, u, v. \\ \quad \operatorname{is\text{-}measure}(a, b, c) \wedge \\ \quad \operatorname{is\text{-}real\text{-}measurable\text{-}fn}(a, b, u) \wedge \\ \quad \operatorname{is\text{-}real\text{-}measurable\text{-}fn}(a, b, v) \wedge \\ \quad p \in \mathbb{R} \wedge \\ \quad 1 \leq p \wedge \\ \quad \|{u}\|_{p} \in \mathbb{R} \wedge \\ \quad \|{v}\|_{p} \in \mathbb{R} \Rightarrow \\ \quad \quad \|{(\lambda x \in a.\; (u(x) - v(x)))}\|_{p} \in \mathbb{R} \wedge \\ \quad \quad \|{(\lambda x \in a.\; (u(x) - v(x)))}\|_{p} \\ \quad \quad \leq\; \|{u}\|_{p} + \|{v}\|_{p} \end{array}}}
\vnbprop{lp-complete}{\vnbstmt{\begin{array}{@{}l@{}} \forall a, b, u, p, s. \\ \quad \operatorname{is\text{-}measure}(a, b, u) \wedge \\ \quad p \in \mathbb{R} \wedge \\ \quad 1 \leq p \wedge \\ \quad s \in (\mathbb{N} \to (a \to \mathbb{R})) \wedge \\ \quad \forall k \in \mathbb{N}.\; \operatorname{is\text{-}real\text{-}measurable\text{-}fn}(a, b, s(k)) \wedge \\ \quad \forall k \in \mathbb{N}.\; (\|{s(k)}\|_{p} \in \mathbb{R}) \wedge \\ \quad \forall c. \\ \quad \quad \operatorname{pos\text{-}rr}(c) \Rightarrow  \\ \quad \quad \quad \exists d \in \mathbb{N}. \forall m, q \in \mathbb{N}. \\ \quad \quad \quad \quad ((d \leq m) \wedge (d \leq q)) \Rightarrow  \\ \quad \quad \quad \quad \quad \|{(\lambda x \in a.\; (s(m)(x) - s(q)(x)))}\|_{p} \leq c \Rightarrow \\ \quad \quad \exists f. \\ \quad \quad \quad \operatorname{is\text{-}real\text{-}measurable\text{-}fn}(a, b, f) \wedge \\ \quad \quad \quad \|{f}\|_{p} \in \mathbb{R} \wedge \\ \quad \quad \quad \operatorname{econverges\text{-}to}((\lambda k \in \mathbb{N}.\; \|{(\lambda x \in a.\; (s(k)(x) - f(x)))}\|_{p}), 0) \end{array}}}

\vnbbookfile{theorem-library/mb-completion}{theorem-library/mb-completion}
\noindent{\small mb-completion.scm -- construction-of-measures.pdf Prop. 3.3 (the completion of a measure) and the countable union of null sets it rests on.  Agent M-B, 2026-10-26. DESTINATION: theorem-library/mb-completion.scm.  Library: measure.}\par
\vnbprop{null-countable-union}{\vnbstmt{\begin{array}{@{}l@{}} \forall a, b, u, f. \\ \quad (\operatorname{is\text{-}measure}(a, b, u) \wedge ((f \in (\mathbb{N} \to \operatorname{power}(a))) \wedge \forall k \in \mathbb{N}.\; \operatorname{is\text{-}null\text{-}set}(a, b, u, f(k)))) \Rightarrow  \\ \quad \quad \operatorname{is\text{-}null\text{-}set}(a, b, u, \operatorname{big\text{-}union}(n, \mathbb{N}, f(n))) \end{array}}}
\vnbprop{completion-sigma-unfold}{\vnbstmt{\begin{array}{@{}l@{}} \forall a, b, u. \\ \quad (\operatorname{completion\text{-}sigma}(a, b, u) \simeq \{s \in \operatorname{power}(a) : \exists \mathit{cpa} \in b, n \in \operatorname{power}(a).\; (\operatorname{is\text{-}null\text{-}set}(a, b, u, n) \wedge (s = (\mathit{cpa} \cup n)))\}) \end{array}}}
\vnbprop{completion-measure-unfold}{\vnbstmt{\begin{array}{@{}l@{}} \forall a, b, u. \\ \quad (\operatorname{completion\text{-}measure}(a, b, u) \simeq \operatorname{res}(\operatorname{method\text{-}i\text{-}outer}(a, b, u), \operatorname{completion\text{-}sigma}(a, b, u))) \end{array}}}
\vnbprop{completion-sigma-membership}{\vnbstmt{\begin{array}{@{}l@{}} \forall a, b, u, x. \\ \quad x \in \operatorname{completion\text{-}sigma}(a, b, u) \\ \quad \Leftrightarrow\; ((x \in \operatorname{power}(a)) \wedge \exists \mathit{cpa} \in b, n \in \operatorname{power}(a).\; (\operatorname{is\text{-}null\text{-}set}(a, b, u, n) \wedge (x = (\mathit{cpa} \cup n)))) \end{array}}}
\vnbprop{completion-sigma-subset-power}{\vnbstmt{\forall a, b, u.\; (\operatorname{completion\text{-}sigma}(a, b, u) \subseteq \operatorname{power}(a))}}
\vnbprop{completion-sigma-rep-in}{\vnbstmt{\begin{array}{@{}l@{}} \forall a, b, u, c, s. \\ \quad (\operatorname{is\text{-}measure}(a, b, u) \wedge ((c \in b) \wedge \operatorname{is\text{-}null\text{-}set}(a, b, u, s))) \Rightarrow  \\ \quad \quad (c \cup s) \in \operatorname{completion\text{-}sigma}(a, b, u) \end{array}}}
\vnbprop{completion-measure-value}{\vnbstmt{\begin{array}{@{}l@{}} \forall a, b, u, c, s. \\ \quad (\operatorname{is\text{-}measure}(a, b, u) \wedge ((c \in b) \wedge \operatorname{is\text{-}null\text{-}set}(a, b, u, s))) \Rightarrow  \\ \quad \quad \operatorname{completion\text{-}measure}(a, b, u)((c \cup s)) = u(c) \end{array}}}
\vnbprop{completion-measure-extends}{\vnbstmt{\begin{array}{@{}l@{}} \forall a, b, u, c. \\ \quad (\operatorname{is\text{-}measure}(a, b, u) \wedge (c \in b)) \Rightarrow (\operatorname{completion\text{-}measure}(a, b, u)(c) = u(c)) \end{array}}}
\vnbprop{completion-sigma-algebra}{\vnbstmt{\begin{array}{@{}l@{}} \forall a, b, u. \\ \quad \operatorname{is\text{-}measure}(a, b, u) \Rightarrow \operatorname{is\text{-}sigma\text{-}algebra}(a, \operatorname{completion\text{-}sigma}(a, b, u)) \end{array}}}
\vnbprop{completion-sigma-contains}{\vnbstmt{\begin{array}{@{}l@{}} \forall a, b, u, c. \\ \quad (\operatorname{is\text{-}measure}(a, b, u) \wedge (c \in b)) \Rightarrow (c \in \operatorname{completion\text{-}sigma}(a, b, u)) \end{array}}}
\vnbprop{completion-sigma-subset-cara}{\vnbstmt{\begin{array}{@{}l@{}} \forall a, b, u, x. \\ \quad (\operatorname{is\text{-}measure}(a, b, u) \wedge (x \in \operatorname{completion\text{-}sigma}(a, b, u))) \Rightarrow  \\ \quad \quad x \in \operatorname{caratheodory\text{-}sets}(a, \operatorname{method\text{-}i\text{-}outer}(a, b, u)) \end{array}}}
\vnbprop{completion-measure-is-cara-measure}{\vnbstmt{\begin{array}{@{}l@{}} \forall a, b, u, x. \\ \quad (\operatorname{is\text{-}measure}(a, b, u) \wedge (x \in \operatorname{completion\text{-}sigma}(a, b, u))) \Rightarrow  \\ \quad \quad \operatorname{completion\text{-}measure}(a, b, u)(x) \\ \quad \quad =\; \operatorname{res}(\operatorname{method\text{-}i\text{-}outer}(a, b, u), \operatorname{caratheodory\text{-}sets}(a, \operatorname{method\text{-}i\text{-}outer}(a, b, u)))(x) \end{array}}}
\vnbprop{completion-measure-is-measure}{\vnbstmt{\begin{array}{@{}l@{}} \forall a, b, u. \\ \quad \operatorname{is\text{-}measure}(a, b, u) \Rightarrow  \\ \quad \quad \operatorname{is\text{-}measure}(a, \operatorname{completion\text{-}sigma}(a, b, u), \operatorname{completion\text{-}measure}(a, b, u)) \end{array}}}
\vnbprop{completion-measure-complete}{\vnbstmt{\begin{array}{@{}l@{}} \forall a, b, u. \\ \quad \operatorname{is\text{-}measure}(a, b, u) \Rightarrow  \\ \quad \quad \operatorname{is\text{-}complete\text{-}measure}(a, \operatorname{completion\text{-}sigma}(a, b, u), \operatorname{completion\text{-}measure}(a, b, u)) \end{array}}}
\vnbprop{completion-measure-unique}{\vnbstmt{\begin{array}{@{}l@{}} \forall a, b, u, c.; \forall x \in \operatorname{completion\text{-}sigma}(a, b, u). \\ \quad \operatorname{is\text{-}measure}(a, b, u) \wedge \\ \quad \operatorname{is\text{-}measure}(a, \operatorname{completion\text{-}sigma}(a, b, u), c) \wedge \\ \quad \forall d \in b.\; (c(d) = u(d)) \Rightarrow \\ \quad \quad c(x) = \operatorname{completion\text{-}measure}(a, b, u)(x) \end{array}}}

\vnbbookfile{theorem-library/mb-restrict}{theorem-library/mb-restrict}
\noindent{\small mb-restrict.scm -- construction-of-measures.pdf Prop. 3.4: "If nu is a complete countably additive measure on B, then for any sigma-subalgebra A of B, nu extends the completion of nu restricted to A".  Agent M-B, 2026-10-26. DESTINATION: theorem-library/mb-restrict.scm.  Library: measure.}\par
\vnbprop{measure-restrict}{\vnbstmt{\begin{array}{@{}l@{}} \forall a, b, u, c. \\ \quad (\operatorname{is\text{-}measure}(a, b, u) \wedge (\operatorname{is\text{-}sigma\text{-}algebra}(a, c) \wedge (c \subseteq b))) \Rightarrow  \\ \quad \quad \operatorname{is\text{-}measure}(a, c, \operatorname{res}(u, c)) \end{array}}}
\vnbprop{completion-of-restrict-in}{\vnbstmt{\begin{array}{@{}l@{}} \forall a, b, u, c. \\ \quad (\operatorname{is\text{-}complete\text{-}measure}(a, b, u) \wedge (\operatorname{is\text{-}sigma\text{-}algebra}(a, c) \wedge (c \subseteq b))) \Rightarrow  \\ \quad \quad \operatorname{completion\text{-}sigma}(a, c, \operatorname{res}(u, c)) \subseteq b \end{array}}}
\vnbprop{complete-extends-completion}{\vnbstmt{\begin{array}{@{}l@{}} \forall a, b, u, c.; \forall x \in \operatorname{completion\text{-}sigma}(a, c, \operatorname{res}(u, c)). \\ \quad (\operatorname{is\text{-}complete\text{-}measure}(a, b, u) \wedge (\operatorname{is\text{-}sigma\text{-}algebra}(a, c) \wedge (c \subseteq b))) \Rightarrow  \\ \quad \quad u(x) = \operatorname{completion\text{-}measure}(a, c, \operatorname{res}(u, c))(x) \end{array}}}

\vnbbookfile{theorem-library/l2-hilbert}{theorem-library/l2-hilbert}
\noindent{\small l2-hilbert.scm -- WAVE 2 (2026-10-27..30), wired 2026-10-30 from scratchpad/md/me-l2.scm. Library: measure.}\par
\vnbprop{epow-two-real}{\vnbstmt{\forall t \in \mathbb{R}.\; ((0 \leq t) \Rightarrow ({t}^{2} = (t \cdot t)))}}
\vnbprop{abs-times-abs}{\vnbstmt{\forall a \in \mathbb{R}.\; ((\lvert a \rvert \cdot \lvert a \rvert) = (a \cdot a))}}
\vnbprop{l2-dist-lp-norm}{\vnbstmt{\begin{array}{@{}l@{}} \forall a, b, u, f, g. \\ \quad (\operatorname{is\text{-}measure}(a, b, u) \wedge ((f \in \operatorname{l2\text{-}set}(a, b, u)) \wedge (g \in \operatorname{l2\text{-}set}(a, b, u)))) \Rightarrow  \\ \quad \quad \sqrt{\langle\cdot\mid\cdot\rangle_{\operatorname{l2\text{-}space}(a, b, u)}(\operatorname{vadd}(\operatorname{l2\text{-}space}(a, b, u))(f, \operatorname{vneg}(\operatorname{l2\text{-}space}(a, b, u))(g)), \operatorname{vadd}(\operatorname{l2\text{-}space}(a, b, u))(f, \operatorname{vneg}(\operatorname{l2\text{-}space}(a, b, u))(g)))} \\ \quad \quad =\; \|{(\lambda x \in a.\; (f(x) - g(x)))}\|_{2} \end{array}}}
\vnbprop{l2-iff-lp2-finite}{\vnbstmt{\begin{array}{@{}l@{}} \forall a, b, u, f. \\ \quad (\operatorname{is\text{-}measure}(a, b, u) \wedge \operatorname{is\text{-}real\text{-}measurable\text{-}fn}(a, b, f)) \Rightarrow  \\ \quad \quad (f \in \operatorname{l2\text{-}set}(a, b, u)) \Leftrightarrow (\|{f}\|_{2} \in \mathbb{R}) \end{array}}}
\vnbprop{l2-seminorm-complete}{\vnbstmt{\begin{array}{@{}l@{}} \forall b, c, u.; \forall s \in (\mathbb{N} \to \operatorname{vec}(\operatorname{l2\text{-}space}(b, c, u))). \\ \quad \operatorname{is\text{-}measure}(b, c, u) \wedge \\ \quad \forall d. \\ \quad \quad \operatorname{pos\text{-}rr}(d) \Rightarrow  \\ \quad \quad \quad \exists f \in \mathbb{N}. \forall m, q \in \mathbb{N}. \\ \quad \quad \quad \quad ((f \leq m) \wedge (f \leq q)) \Rightarrow  \\ \quad \quad \quad \quad \quad \sqrt{\langle\cdot\mid\cdot\rangle_{\operatorname{l2\text{-}space}(b, c, u)}(\operatorname{vadd}(\operatorname{l2\text{-}space}(b, c, u))(s(m), \operatorname{vneg}(\operatorname{l2\text{-}space}(b, c, u))(s(q))), \operatorname{vadd}(\operatorname{l2\text{-}space}(b, c, u))(s(m), \operatorname{vneg}(\operatorname{l2\text{-}space}(b, c, u))(s(q))))} \\ \quad \quad \quad \quad \quad \leq\; d \Rightarrow \\ \quad \quad \exists a \in \operatorname{vec}(\operatorname{l2\text{-}space}(b, c, u)). \forall d. \\ \quad \quad \quad \operatorname{pos\text{-}rr}(d) \Rightarrow  \\ \quad \quad \quad \quad \exists f \in \mathbb{N}. \forall n \in \mathbb{N}. \\ \quad \quad \quad \quad \quad (f \leq n) \Rightarrow  \\ \quad \quad \quad \quad \quad \quad \sqrt{\langle\cdot\mid\cdot\rangle_{\operatorname{l2\text{-}space}(b, c, u)}(\operatorname{vadd}(\operatorname{l2\text{-}space}(b, c, u))(s(n), \operatorname{vneg}(\operatorname{l2\text{-}space}(b, c, u))(a)), \operatorname{vadd}(\operatorname{l2\text{-}space}(b, c, u))(s(n), \operatorname{vneg}(\operatorname{l2\text{-}space}(b, c, u))(a)))} \\ \quad \quad \quad \quad \quad \quad \leq\; d \end{array}}}
\vnbprop{l2-quotient-is-real-hilbert}{\vnbstmt{\begin{array}{@{}l@{}} \forall a, b, u. \\ \quad \operatorname{is\text{-}measure}(a, b, u) \Rightarrow  \\ \quad \quad \operatorname{is\text{-}real\text{-}hilbert\text{-}space}(\operatorname{sip\text{-}quotient}(\operatorname{l2\text{-}space}(a, b, u))) \end{array}}}

\vnbbookfile{structure-library/l2c-vocab}{structure-library/l2c-vocab}
\noindent{\small l2c-vocab.scm -- wave of 2026-11-09/10 (agent TH3-L2), wired 2026-11-10 from scratchpad/th3-1109/l2/l2c-vocab.scm. Library: measure.}\par
\par\medskip\noindent\textbf{Definition} (\texttt{cl2-hilbert}).\;$L^2_{\mathbb{C}}({\omega}, {\mu}) :=$ $\operatorname{csip\text{-}cips}(\operatorname{complexify}(\operatorname{sip\text{-}quotient}(\operatorname{l2\text{-}space}(\omega, \mathit{ca}, \mu))))$

\vnbbookfile{theorem-library/l2c-a}{theorem-library/l2c-a}
\noindent{\small l2c-a.scm -- wave of 2026-11-09/10 (agent TH3-L2), wired 2026-11-10 from scratchpad/th3-1109/l2/l2c-a.scm. Library: measure.}\par
\vnbprop{cxq-sip}{\vnbstmt{\begin{array}{@{}l@{}} \forall v. \\ \quad \operatorname{is\text{-}semi\text{-}inner\text{-}product\text{-}space}(v) \Rightarrow  \\ \quad \quad \operatorname{is\text{-}semi\text{-}inner\text{-}product\text{-}space}(\operatorname{sip\text{-}quotient}(v)) \end{array}}}
\vnbprop{cxq-vec}{\vnbstmt{\begin{array}{@{}l@{}} \forall v. \\ \quad \operatorname{is\text{-}semi\text{-}inner\text{-}product\text{-}space}(v) \Rightarrow  \\ \quad \quad (\operatorname{vec}(\operatorname{csip\text{-}cips}(\operatorname{complexify}(\operatorname{sip\text{-}quotient}(v)))) \simeq (\operatorname{quotient}(\operatorname{sip\text{-}setoid}(v)) \times \operatorname{quotient}(\operatorname{sip\text{-}setoid}(v)))) \end{array}}}
\vnbprop{cxq-pair-in}{\vnbstmt{\begin{array}{@{}l@{}} \forall v, a_{1}, a_{2}. \\ \quad (\operatorname{is\text{-}semi\text{-}inner\text{-}product\text{-}space}(v) \wedge ((a_{1} \in \operatorname{vec}(v)) \wedge (a_{2} \in \operatorname{vec}(v)))) \Rightarrow  \\ \quad \quad {[}\operatorname{class}(\operatorname{sip\text{-}setoid}(v), a_{1}), \operatorname{class}(\operatorname{sip\text{-}setoid}(v), a_{2})] \in \operatorname{vec}(\operatorname{csip\text{-}cips}(\operatorname{complexify}(\operatorname{sip\text{-}quotient}(v)))) \end{array}}}
\vnbprop{cxq-vzero}{\vnbstmt{\begin{array}{@{}l@{}} \forall v. \\ \quad \operatorname{is\text{-}semi\text{-}inner\text{-}product\text{-}space}(v) \Rightarrow  \\ \quad \quad (\operatorname{vzero}(\operatorname{csip\text{-}cips}(\operatorname{complexify}(\operatorname{sip\text{-}quotient}(v)))) \simeq {[}\operatorname{class}(\operatorname{sip\text{-}setoid}(v), \operatorname{vzero}(v)), \operatorname{class}(\operatorname{sip\text{-}setoid}(v), \operatorname{vzero}(v))]) \end{array}}}
\vnbprop{cxq-rep}{\vnbstmt{\begin{array}{@{}l@{}} \forall v.; \forall x \in \operatorname{vec}(\operatorname{csip\text{-}cips}(\operatorname{complexify}(\operatorname{sip\text{-}quotient}(v)))). \\ \quad \operatorname{is\text{-}semi\text{-}inner\text{-}product\text{-}space}(v) \Rightarrow  \\ \quad \quad \exists a_{1}, a_{2} \in \operatorname{vec}(v). \\ \quad \quad \quad x = {[}\operatorname{class}(\operatorname{sip\text{-}setoid}(v), a_{1}), \operatorname{class}(\operatorname{sip\text{-}setoid}(v), a_{2})] \end{array}}}
\vnbprop{cxq-vadd-classes}{\vnbstmt{\begin{array}{@{}l@{}} \forall v, a_{1}, a_{2}, b_{1}, b_{2}. \\ \quad (\operatorname{is\text{-}semi\text{-}inner\text{-}product\text{-}space}(v) \wedge ((a_{1} \in \operatorname{vec}(v)) \wedge ((a_{2} \in \operatorname{vec}(v)) \wedge ((b_{1} \in \operatorname{vec}(v)) \wedge (b_{2} \in \operatorname{vec}(v)))))) \Rightarrow  \\ \quad \quad \operatorname{vadd}(\operatorname{csip\text{-}cips}(\operatorname{complexify}(\operatorname{sip\text{-}quotient}(v))))({[}\operatorname{class}(\operatorname{sip\text{-}setoid}(v), a_{1}), \operatorname{class}(\operatorname{sip\text{-}setoid}(v), a_{2})], \\ \quad \quad \quad {[}\operatorname{class}(\operatorname{sip\text{-}setoid}(v), b_{1}), \operatorname{class}(\operatorname{sip\text{-}setoid}(v), b_{2})]) \\ \quad \quad =\; {[}\operatorname{class}(\operatorname{sip\text{-}setoid}(v), \operatorname{vadd}(v)(a_{1}, b_{1})), \operatorname{class}(\operatorname{sip\text{-}setoid}(v), \operatorname{vadd}(v)(a_{2}, b_{2}))] \end{array}}}
\vnbprop{cxq-act-classes}{\vnbstmt{\begin{array}{@{}l@{}} \forall v, z, a_{1}, a_{2}. \\ \quad (\operatorname{is\text{-}semi\text{-}inner\text{-}product\text{-}space}(v) \wedge ((z \in \mathbb{C}) \wedge ((a_{1} \in \operatorname{vec}(v)) \wedge (a_{2} \in \operatorname{vec}(v))))) \Rightarrow  \\ \quad \quad \operatorname{act}(\operatorname{csip\text{-}cips}(\operatorname{complexify}(\operatorname{sip\text{-}quotient}(v))))(z, \\ \quad \quad \quad {[}\operatorname{class}(\operatorname{sip\text{-}setoid}(v), a_{1}), \operatorname{class}(\operatorname{sip\text{-}setoid}(v), a_{2})]) \\ \quad \quad =\; {[}\operatorname{class}(\operatorname{sip\text{-}setoid}(v), \operatorname{vadd}(v)(\operatorname{act}(v)(\operatorname{real\text{-}part}(z), a_{1}), \operatorname{act}(v)(\operatorname{imag\text{-}part}(z), \operatorname{vneg}(v)(a_{2})))), \operatorname{class}(\operatorname{sip\text{-}setoid}(v), \operatorname{vadd}(v)(\operatorname{act}(v)(\operatorname{real\text{-}part}(z), a_{2}), \operatorname{act}(v)(\operatorname{imag\text{-}part}(z), a_{1})))] \end{array}}}
\vnbprop{cxq-ip-classes}{\vnbstmt{\begin{array}{@{}l@{}} \forall v, a_{1}, a_{2}, b_{1}, b_{2}. \\ \quad (\operatorname{is\text{-}semi\text{-}inner\text{-}product\text{-}space}(v) \wedge ((a_{1} \in \operatorname{vec}(v)) \wedge ((a_{2} \in \operatorname{vec}(v)) \wedge ((b_{1} \in \operatorname{vec}(v)) \wedge (b_{2} \in \operatorname{vec}(v)))))) \Rightarrow  \\ \quad \quad \langle\cdot,\cdot\rangle_{\operatorname{csip\text{-}cips}(\operatorname{complexify}(\operatorname{sip\text{-}quotient}(v)))}({[}\operatorname{class}(\operatorname{sip\text{-}setoid}(v), a_{1}), \operatorname{class}(\operatorname{sip\text{-}setoid}(v), a_{2})], \\ \quad \quad \quad {[}\operatorname{class}(\operatorname{sip\text{-}setoid}(v), b_{1}), \operatorname{class}(\operatorname{sip\text{-}setoid}(v), b_{2})]) \\ \quad \quad =\; (\langle\cdot\mid\cdot\rangle_{v}(a_{1}, b_{1}) + \langle\cdot\mid\cdot\rangle_{v}(a_{2}, b_{2})) + (+i \cdot (\langle\cdot\mid\cdot\rangle_{v}(a_{2}, b_{1}) - \langle\cdot\mid\cdot\rangle_{v}(a_{1}, b_{2}))) \end{array}}}
\vnbprop{cl2-hilbert-unfold}{\vnbstmt{\begin{array}{@{}l@{}} \forall a, b, u. \\ \quad (L^2_{\mathbb{C}}({a}, {u}) \simeq \operatorname{csip\text{-}cips}(\operatorname{complexify}(\operatorname{sip\text{-}quotient}(\operatorname{l2\text{-}space}(a, b, u))))) \end{array}}}
\vnbprop{cl2-hilbert-is-hilbert-space}{\vnbstmt{\begin{array}{@{}l@{}} \forall a, b, u. \\ \quad \operatorname{is\text{-}measure}(a, b, u) \Rightarrow \operatorname{is\text{-}hilbert\text{-}space}(L^2_{\mathbb{C}}({a}, {u})) \end{array}}}
\vnbprop{cl2-hilbert-vec}{\vnbstmt{\begin{array}{@{}l@{}} \forall a, b, u. \\ \quad \operatorname{is\text{-}measure}(a, b, u) \Rightarrow  \\ \quad \quad (\operatorname{vec}(L^2_{\mathbb{C}}({a}, {u})) \simeq (\operatorname{quotient}(\operatorname{sip\text{-}setoid}(\operatorname{l2\text{-}space}(a, b, u))) \times \operatorname{quotient}(\operatorname{sip\text{-}setoid}(\operatorname{l2\text{-}space}(a, b, u))))) \end{array}}}
\vnbprop{cl2-hilbert-pair-in}{\vnbstmt{\begin{array}{@{}l@{}} \forall a, b, u, c, d. \\ \quad (\operatorname{is\text{-}measure}(a, b, u) \wedge ((c \in \operatorname{l2\text{-}set}(a, b, u)) \wedge (d \in \operatorname{l2\text{-}set}(a, b, u)))) \Rightarrow  \\ \quad \quad {[}\operatorname{class}(\operatorname{sip\text{-}setoid}(\operatorname{l2\text{-}space}(a, b, u)), c), \operatorname{class}(\operatorname{sip\text{-}setoid}(\operatorname{l2\text{-}space}(a, b, u)), d)] \\ \quad \quad \in\; \operatorname{vec}(L^2_{\mathbb{C}}({a}, {u})) \end{array}}}
\vnbprop{cl2-hilbert-rep}{\vnbstmt{\begin{array}{@{}l@{}} \forall a, b, u.; \forall x \in \operatorname{vec}(L^2_{\mathbb{C}}({a}, {u})). \\ \quad \operatorname{is\text{-}measure}(a, b, u) \Rightarrow  \\ \quad \quad \exists c, d \in \operatorname{l2\text{-}set}(a, b, u). \\ \quad \quad \quad x = {[}\operatorname{class}(\operatorname{sip\text{-}setoid}(\operatorname{l2\text{-}space}(a, b, u)), c), \operatorname{class}(\operatorname{sip\text{-}setoid}(\operatorname{l2\text{-}space}(a, b, u)), d)] \end{array}}}
\vnbprop{cl2-hilbert-ip-classes}{\vnbstmt{\begin{array}{@{}l@{}} \forall a, b, u, c, d, f, g. \\ \quad \operatorname{is\text{-}measure}(a, b, u) \wedge \\ \quad c \in \operatorname{l2\text{-}set}(a, b, u) \wedge \\ \quad d \in \operatorname{l2\text{-}set}(a, b, u) \wedge \\ \quad f \in \operatorname{l2\text{-}set}(a, b, u) \wedge \\ \quad g \in \operatorname{l2\text{-}set}(a, b, u) \Rightarrow \\ \quad \quad \langle\cdot,\cdot\rangle_{L^2_{\mathbb{C}}({a}, {u})}({[}\operatorname{class}(\operatorname{sip\text{-}setoid}(\operatorname{l2\text{-}space}(a, b, u)), c), \operatorname{class}(\operatorname{sip\text{-}setoid}(\operatorname{l2\text{-}space}(a, b, u)), d)], \\ \quad \quad \quad {[}\operatorname{class}(\operatorname{sip\text{-}setoid}(\operatorname{l2\text{-}space}(a, b, u)), f), \operatorname{class}(\operatorname{sip\text{-}setoid}(\operatorname{l2\text{-}space}(a, b, u)), g)]) \\ \quad \quad =\; (\operatorname{rintegral}(a, b, u, (\lambda x \in a.\; (c(x) \cdot f(x)))) + \operatorname{rintegral}(a, b, u, (\lambda x \in a.\; (d(x) \cdot g(x))))) + (+i \cdot (\operatorname{rintegral}(a, b, u, (\lambda x \in a.\; (d(x) \cdot f(x)))) - \operatorname{rintegral}(a, b, u, (\lambda x \in a.\; (c(x) \cdot g(x)))))) \end{array}}}
\vnbprop{cl2-class-eq-iff-ae}{\vnbstmt{\begin{array}{@{}l@{}} \forall a, b, u, f, g. \\ \quad (\operatorname{is\text{-}measure}(a, b, u) \wedge ((f \in \operatorname{l2\text{-}set}(a, b, u)) \wedge (g \in \operatorname{l2\text{-}set}(a, b, u)))) \Rightarrow  \\ \quad \quad \operatorname{class}(\operatorname{sip\text{-}setoid}(\operatorname{l2\text{-}space}(a, b, u)), f) = \operatorname{class}(\operatorname{sip\text{-}setoid}(\operatorname{l2\text{-}space}(a, b, u)), g) \\ \quad \quad \Leftrightarrow\; \exists n. \\ \quad \quad \quad \operatorname{is\text{-}null\text{-}set}(a, b, u, n) \wedge \forall x \in a.\; (\neg (x \in n) \Rightarrow (f(x) = g(x))) \end{array}}}

\vnbbookfile{structure-library/l2-integral-operator}{structure-library/l2-integral-operator}
\noindent{\small l2-integral-operator.scm -- wave of 2026-11-09/10 (agent TH3-L2C), wired 2026-11-10 from scratchpad/th3-1109/l2c/kc-vocab.scm. Library: measure.}\par
\par\medskip\noindent\textbf{Definition} (\texttt{kernel-apply}).\;$\int {\operatorname{kak\_}}(\cdot, y)\, {\operatorname{kaf\_}}(y)\, d{\operatorname{kanu\_}}(y) :=$
\begin{equation*}
\begin{array}{@{}l@{}} (\lambda \operatorname{sgx\_} \in \operatorname{kaxs\_}.\; \\ \quad (\left(\text{if }(\operatorname{x\text{-}section\text{-}integral}(\operatorname{kaxs\_}, \operatorname{kays\_}, \operatorname{kacb\_}, \operatorname{kanu\_}, \operatorname{pos\text{-}part}((\operatorname{kaxs\_} \times \operatorname{kays\_}), (\lambda \operatorname{kpw\_} \in (\operatorname{kaxs\_} \times \operatorname{kays\_}).\; (\operatorname{kak\_}(\operatorname{kpw\_}) \cdot \operatorname{kaf\_}(\operatorname{nth}(2, \operatorname{kpw\_}))))))(\operatorname{sgx\_}) \in \mathbb{R})\text{ then }\operatorname{x\text{-}section\text{-}integral}(\operatorname{kaxs\_}, \operatorname{kays\_}, \operatorname{kacb\_}, \operatorname{kanu\_}, \operatorname{pos\text{-}part}((\operatorname{kaxs\_} \times \operatorname{kays\_}), (\lambda \operatorname{kpw\_} \in (\operatorname{kaxs\_} \times \operatorname{kays\_}).\; (\operatorname{kak\_}(\operatorname{kpw\_}) \cdot \operatorname{kaf\_}(\operatorname{nth}(2, \operatorname{kpw\_}))))))(\operatorname{sgx\_})\text{ else }0\right) - \left(\text{if }(\operatorname{x\text{-}section\text{-}integral}(\operatorname{kaxs\_}, \operatorname{kays\_}, \operatorname{kacb\_}, \operatorname{kanu\_}, \operatorname{neg\text{-}part}((\operatorname{kaxs\_} \times \operatorname{kays\_}), (\lambda \operatorname{kpw\_} \in (\operatorname{kaxs\_} \times \operatorname{kays\_}).\; (\operatorname{kak\_}(\operatorname{kpw\_}) \cdot \operatorname{kaf\_}(\operatorname{nth}(2, \operatorname{kpw\_}))))))(\operatorname{sgx\_}) \in \mathbb{R})\text{ then }\operatorname{x\text{-}section\text{-}integral}(\operatorname{kaxs\_}, \operatorname{kays\_}, \operatorname{kacb\_}, \operatorname{kanu\_}, \operatorname{neg\text{-}part}((\operatorname{kaxs\_} \times \operatorname{kays\_}), (\lambda \operatorname{kpw\_} \in (\operatorname{kaxs\_} \times \operatorname{kays\_}).\; (\operatorname{kak\_}(\operatorname{kpw\_}) \cdot \operatorname{kaf\_}(\operatorname{nth}(2, \operatorname{kpw\_}))))))(\operatorname{sgx\_})\text{ else }0\right))) \end{array}
\end{equation*}

\par\medskip\noindent\textbf{Definition} (\texttt{cl2-int-op}).\;$T_{{\operatorname{kck1\_}} + i {\operatorname{kck2\_}}} :=$
\begin{equation*}
\begin{array}{@{}l@{}} (\lambda \operatorname{kcz\_} \in \operatorname{vec}(L^2_{\mathbb{C}}({\operatorname{kcys\_}}, {\operatorname{kcnu\_}})).\; \\ \quad {[}\operatorname{class}(\operatorname{sip\text{-}setoid}(\operatorname{l2\text{-}space}(\operatorname{kcxs\_}, \operatorname{kcca\_}, \operatorname{kcmu\_})), \operatorname{vadd}(\operatorname{l2\text{-}space}(\operatorname{kcxs\_}, \operatorname{kcca\_}, \operatorname{kcmu\_}))(\int {\operatorname{kck1\_}}(\cdot, y)\, {\varepsilon(\operatorname{nth}(1, \operatorname{kcz\_}))}(y)\, d{\operatorname{kcnu\_}}(y), \operatorname{vneg}(\operatorname{l2\text{-}space}(\operatorname{kcxs\_}, \operatorname{kcca\_}, \operatorname{kcmu\_}))(\int {\operatorname{kck2\_}}(\cdot, y)\, {\varepsilon(\operatorname{nth}(2, \operatorname{kcz\_}))}(y)\, d{\operatorname{kcnu\_}}(y)))), \operatorname{class}(\operatorname{sip\text{-}setoid}(\operatorname{l2\text{-}space}(\operatorname{kcxs\_}, \operatorname{kcca\_}, \operatorname{kcmu\_})), \operatorname{vadd}(\operatorname{l2\text{-}space}(\operatorname{kcxs\_}, \operatorname{kcca\_}, \operatorname{kcmu\_}))(\int {\operatorname{kck1\_}}(\cdot, y)\, {\varepsilon(\operatorname{nth}(2, \operatorname{kcz\_}))}(y)\, d{\operatorname{kcnu\_}}(y), \int {\operatorname{kck2\_}}(\cdot, y)\, {\varepsilon(\operatorname{nth}(1, \operatorname{kcz\_}))}(y)\, d{\operatorname{kcnu\_}}(y)))]) \end{array}
\end{equation*}


\vnbbookfile{structure-library/push-forward}{structure-library/push-forward}
\noindent{\small structure-library/push-forward.scm -- the push-forward of a measure along a map (RM-3B, 2026-10-10).  DEFINITIONS ONLY; every law is PROVEN in theorem-library/push-forward.scm (prefix pfw-).}\par
\par\medskip\noindent\textbf{Definition} (\texttt{push-forward}).\;$\operatorname{push\text{-}forward}(\omega, \mathit{cb}, \mu, h) :=$ $(\lambda \operatorname{pwb\_} \in \mathit{cb}.\; \mu(\{\operatorname{pwx\_} \in \omega : (h(\operatorname{pwx\_}) \in \operatorname{pwb\_})\}))$

\vnbbookfile{theorem-library/push-forward}{theorem-library/push-forward}
\noindent{\small theorem-library/push-forward.scm -- Borel preimages under a continuous map, the push-forward of a measure, and the change of variables formula (RM-3B, 2026-10-10).  Vocabulary: structure-library/push-forward.scm (PUSH-FORWARD). Source: Thayer, Construction of Measures, Prop. 1.19 and Ch. 2 s. 7,}\par
\vnbprop{pfw-push-forward-unfold}{\vnbstmt{\begin{array}{@{}l@{}} \forall a, b, u, h. \\ \quad (\operatorname{push\text{-}forward}(a, b, u, h) \simeq (\lambda \mathit{pwb} \in b.\; u(\{x \in a : (h(x) \in \mathit{pwb})\}))) \end{array}}}
\vnbprop{pfw-push-forward-value}{\vnbstmt{\begin{array}{@{}l@{}} \forall a, c, u, h.; \forall b \in c. \\ \quad (\operatorname{push\text{-}forward}(a, c, u, h)(b) \simeq u(\{x \in a : (h(x) \in b)\})) \end{array}}}
\vnbprop{pfw-preimage-sigma-algebra}{\vnbstmt{\begin{array}{@{}l@{}} \forall a, b, a_{2}, h. \\ \quad (\operatorname{is\text{-}sigma\text{-}algebra}(a, b) \wedge ((a_{2} \in \mathbf{Set}) \wedge (h \in (a \to a_{2})))) \Rightarrow  \\ \quad \quad \operatorname{is\text{-}sigma\text{-}algebra}(a_{2}, \{e \in \operatorname{power}(a_{2}) : (\{x \in a : (h(x) \in e)\} \in b)\}) \end{array}}}
\vnbprop{pfw-borel-sigma-algebra}{\vnbstmt{\forall s.\; (\operatorname{is\text{-}metric\text{-}space}(s) \Rightarrow \operatorname{is\text{-}sigma\text{-}algebra}(\operatorname{pts}(s), \operatorname{borel\text{-}sets}(s)))}}
\vnbprop{pfw-open-borel}{\vnbstmt{\forall s, u.\; (\operatorname{is\text{-}open}(s, u) \Rightarrow (u \in \operatorname{borel\text{-}sets}(s)))}}
\vnbprop{pfw-borel-preimage}{\vnbstmt{\begin{array}{@{}l@{}} \forall s, t, h, b. \\ \quad (\operatorname{is\text{-}continuous}(s, t, h) \wedge (b \in \operatorname{borel\text{-}sets}(t))) \Rightarrow  \\ \quad \quad \{x \in \operatorname{pts}(s) : (h(x) \in b)\} \in \operatorname{borel\text{-}sets}(s) \end{array}}}
\vnbprop{pfw-continuous-borel-measurable}{\vnbstmt{\begin{array}{@{}l@{}} \forall s, t, h. \\ \quad \operatorname{is\text{-}continuous}(s, t, h) \Rightarrow  \\ \quad \quad \operatorname{is\text{-}measurable\text{-}map}(\operatorname{pts}(s), \operatorname{borel\text{-}sets}(s), \operatorname{pts}(t), \operatorname{borel\text{-}sets}(t), h) \end{array}}}
\vnbprop{pfw-measurable-map-enlarge}{\vnbstmt{\begin{array}{@{}l@{}} \forall a, b, a_{2}, c, h, d. \\ \quad (\operatorname{is\text{-}measurable\text{-}map}(a, b, a_{2}, c, h) \wedge (\operatorname{is\text{-}sigma\text{-}algebra}(a, d) \wedge (b \subseteq d))) \Rightarrow  \\ \quad \quad \operatorname{is\text{-}measurable\text{-}map}(a, d, a_{2}, c, h) \end{array}}}
\vnbprop{pfw-push-forward-measure}{\vnbstmt{\begin{array}{@{}l@{}} \forall a, b, u, a_{2}, c, h. \\ \quad (\operatorname{is\text{-}measure}(a, b, u) \wedge \operatorname{is\text{-}measurable\text{-}map}(a, b, a_{2}, c, h)) \Rightarrow  \\ \quad \quad \operatorname{is\text{-}measure}(a_{2}, c, \operatorname{push\text{-}forward}(a, c, u, h)) \end{array}}}
\vnbprop{pfw-comp-apply}{\vnbstmt{\forall a, f, h.; \forall x \in a.\; ((\lambda z \in a.\; f(h(z)))(x) \simeq f(h(x)))}}
\vnbprop{pfw-comp-fun}{\vnbstmt{\begin{array}{@{}l@{}} \forall a, a_{2}, d, f, h. \\ \quad ((a \in \mathbf{Set}) \wedge ((h \in (a \to a_{2})) \wedge (f \in (a_{2} \to d)))) \Rightarrow  \\ \quad \quad (\lambda z \in a.\; f(h(z))) \in (a \to d) \end{array}}}
\vnbprop{pfw-comp-measurable-fn}{\vnbstmt{\begin{array}{@{}l@{}} \forall a, b, a_{2}, c, h, f. \\ \quad (\operatorname{is\text{-}measurable\text{-}map}(a, b, a_{2}, c, h) \wedge \operatorname{is\text{-}measurable\text{-}fn}(a_{2}, c, f)) \Rightarrow  \\ \quad \quad \operatorname{is\text{-}measurable\text{-}fn}(a, b, (\lambda z \in a.\; f(h(z)))) \end{array}}}
\vnbprop{pfw-simple-change-of-variables}{\vnbstmt{\begin{array}{@{}l@{}} \forall a, b, u, a_{2}, c, h, s. \\ \quad (\operatorname{is\text{-}measure}(a, b, u) \wedge (\operatorname{is\text{-}measurable\text{-}map}(a, b, a_{2}, c, h) \wedge \operatorname{is\text{-}simple\text{-}fn}(a_{2}, c, s))) \Rightarrow  \\ \quad \quad \operatorname{is\text{-}simple\text{-}fn}(a, b, (\lambda z \in a.\; s(h(z)))) \wedge \\ \quad \quad \operatorname{integral}(a_{2}, c, \operatorname{push\text{-}forward}(a, c, u, h), s) = \operatorname{integral}(a, b, u, (\lambda z \in a.\; s(h(z)))) \end{array}}}
\vnbprop{pfw-change-of-variables}{\vnbstmt{\begin{array}{@{}l@{}} \forall a, b, u, a_{2}, c, h, f. \\ \quad (\operatorname{is\text{-}measure}(a, b, u) \wedge (\operatorname{is\text{-}measurable\text{-}map}(a, b, a_{2}, c, h) \wedge \operatorname{is\text{-}measurable\text{-}fn}(a_{2}, c, f))) \Rightarrow  \\ \quad \quad \operatorname{is\text{-}measurable\text{-}fn}(a, b, (\lambda z \in a.\; f(h(z)))) \wedge \\ \quad \quad \operatorname{integral}(a_{2}, c, \operatorname{push\text{-}forward}(a, c, u, h), f) = \operatorname{integral}(a, b, u, (\lambda z \in a.\; f(h(z)))) \end{array}}}
\vnbprop{pfw-pos-part-comp}{\vnbstmt{\begin{array}{@{}l@{}} \forall a, a_{2}.; \forall f \in (a_{2} \to \mathbb{R}), h \in (a \to a_{2}), v \in a. \\ \quad \operatorname{pos\text{-}part}(a, (\lambda z \in a.\; f(h(z))))(v) = (\lambda z \in a.\; \operatorname{pos\text{-}part}(a_{2}, f)(h(z)))(v) \end{array}}}
\vnbprop{pfw-neg-part-comp}{\vnbstmt{\begin{array}{@{}l@{}} \forall a, a_{2}.; \forall f \in (a_{2} \to \mathbb{R}), h \in (a \to a_{2}), v \in a. \\ \quad \operatorname{neg\text{-}part}(a, (\lambda z \in a.\; f(h(z))))(v) = (\lambda z \in a.\; \operatorname{neg\text{-}part}(a_{2}, f)(h(z)))(v) \end{array}}}
\vnbprop{pfw-comp-real-measurable}{\vnbstmt{\begin{array}{@{}l@{}} \forall a, b, a_{2}, c, h, f. \\ \quad (\operatorname{is\text{-}measurable\text{-}map}(a, b, a_{2}, c, h) \wedge \operatorname{is\text{-}real\text{-}measurable\text{-}fn}(a_{2}, c, f)) \Rightarrow  \\ \quad \quad \operatorname{is\text{-}real\text{-}measurable\text{-}fn}(a, b, (\lambda z \in a.\; f(h(z)))) \end{array}}}
\vnbprop{pfw-integrable-iff}{\vnbstmt{\begin{array}{@{}l@{}} \forall a, b, u, a_{2}, c, h, f. \\ \quad (\operatorname{is\text{-}measure}(a, b, u) \wedge (\operatorname{is\text{-}measurable\text{-}map}(a, b, a_{2}, c, h) \wedge \operatorname{is\text{-}real\text{-}measurable\text{-}fn}(a_{2}, c, f))) \Rightarrow  \\ \quad \quad \operatorname{is\text{-}integrable}(a_{2}, c, \operatorname{push\text{-}forward}(a, c, u, h), f) \\ \quad \quad \Leftrightarrow\; \operatorname{is\text{-}integrable}(a, b, u, (\lambda z \in a.\; f(h(z)))) \end{array}}}
\vnbprop{pfw-change-of-variables-real}{\vnbstmt{\begin{array}{@{}l@{}} \forall a, b, u, a_{2}, c, h, f. \\ \quad \operatorname{is\text{-}measure}(a, b, u) \wedge \\ \quad \operatorname{is\text{-}measurable\text{-}map}(a, b, a_{2}, c, h) \wedge \\ \quad \operatorname{is\text{-}real\text{-}measurable\text{-}fn}(a_{2}, c, f) \wedge \\ \quad \operatorname{is\text{-}integrable}(a_{2}, c, \operatorname{push\text{-}forward}(a, c, u, h), f) \Rightarrow \\ \quad \quad \operatorname{is\text{-}integrable}(a, b, u, (\lambda z \in a.\; f(h(z)))) \wedge \\ \quad \quad \operatorname{rintegral}(a_{2}, c, \operatorname{push\text{-}forward}(a, c, u, h), f) = \operatorname{rintegral}(a, b, u, (\lambda z \in a.\; f(h(z)))) \end{array}}}

\vnbbookfile{theorem-library/lp-measure-isomorphism}{theorem-library/lp-measure-isomorphism}
\noindent{\small lp-measure-isomorphism.scm -- WAVE 2 (2026-10-27..30), wired 2026-10-30 from scratchpad/md/md-iso.scm. Library: measure.}\par
\vnbprop{measure-preserving-push-forward}{\vnbstmt{\begin{array}{@{}l@{}} \forall a, b, u, a_{2}, c, d, t. \\ \quad \operatorname{is\text{-}measure}(a, b, u) \wedge \\ \quad \operatorname{is\text{-}measure}(a_{2}, c, d) \wedge \\ \quad \operatorname{is\text{-}measurable\text{-}map}(a, b, a_{2}, c, t) \wedge \\ \quad \forall \mathit{mib} \in c.\; (d(\mathit{mib}) = u(\{x \in a : (t(x) \in \mathit{mib})\})) \Rightarrow \\ \quad \quad \operatorname{push\text{-}forward}(a, c, u, t) = d \end{array}}}
\vnbprop{comp-measure-preserving-lp-norm}{\vnbstmt{\begin{array}{@{}l@{}} \forall a, b, u, a_{2}, c, d, t, f, p. \\ \quad \operatorname{is\text{-}measure}(a, b, u) \wedge \\ \quad \operatorname{is\text{-}measure}(a_{2}, c, d) \wedge \\ \quad \operatorname{is\text{-}measurable\text{-}map}(a, b, a_{2}, c, t) \wedge \\ \quad \forall \mathit{mib} \in c.\; (d(\mathit{mib}) = u(\{x \in a : (t(x) \in \mathit{mib})\})) \wedge \\ \quad \operatorname{is\text{-}real\text{-}measurable\text{-}fn}(a_{2}, c, f) \wedge \\ \quad p \in \mathbb{R} \wedge \\ \quad 0 < p \Rightarrow \\ \quad \quad \operatorname{is\text{-}real\text{-}measurable\text{-}fn}(a, b, (\lambda z \in a.\; f(t(z)))) \wedge \\ \quad \quad \|{(\lambda z \in a.\; f(t(z)))}\|_{p} = \|{f}\|_{p} \end{array}}}
\vnbprop{abs-level-set-measurable}{\vnbstmt{\begin{array}{@{}l@{}} \forall a_{2}, b, u, f, m. \\ \quad (\operatorname{is\text{-}measure}(a_{2}, b, u) \wedge (\operatorname{is\text{-}real\text{-}measurable\text{-}fn}(a_{2}, b, f) \wedge (m \in \operatorname{rr\text{-}pos\text{-}star}))) \Rightarrow  \\ \quad \quad \{x \in a_{2} : (m < \lvert f(x) \rvert)\} \in b \end{array}}}
\vnbprop{comp-measure-preserving-level-set}{\vnbstmt{\begin{array}{@{}l@{}} \forall a, b, u, a_{2}, c, d, t, f, m. \\ \quad \operatorname{is\text{-}measure}(a, b, u) \wedge \\ \quad \operatorname{is\text{-}measure}(a_{2}, c, d) \wedge \\ \quad \operatorname{is\text{-}measurable\text{-}map}(a, b, a_{2}, c, t) \wedge \\ \quad \forall \mathit{mib} \in c.\; (d(\mathit{mib}) = u(\{x \in a : (t(x) \in \mathit{mib})\})) \wedge \\ \quad \operatorname{is\text{-}real\text{-}measurable\text{-}fn}(a_{2}, c, f) \wedge \\ \quad m \in \operatorname{rr\text{-}pos\text{-}star} \Rightarrow \\ \quad \quad u(\{\mathit{lix} \in a : (m < \lvert (\lambda z \in a.\; f(t(z)))(\mathit{lix}) \rvert)\}) = d(\{\mathit{lix} \in a_{2} : (m < \lvert f(\mathit{lix}) \rvert)\}) \end{array}}}
\vnbprop{comp-measure-preserving-linf-norm}{\vnbstmt{\begin{array}{@{}l@{}} \forall a, b, u, a_{2}, c, d, t, f. \\ \quad \operatorname{is\text{-}measure}(a, b, u) \wedge \\ \quad \operatorname{is\text{-}measure}(a_{2}, c, d) \wedge \\ \quad \operatorname{is\text{-}measurable\text{-}map}(a, b, a_{2}, c, t) \wedge \\ \quad \forall \mathit{mib} \in c.\; (d(\mathit{mib}) = u(\{x \in a : (t(x) \in \mathit{mib})\})) \wedge \\ \quad \operatorname{is\text{-}real\text{-}measurable\text{-}fn}(a_{2}, c, f) \Rightarrow \\ \quad \quad \operatorname{linf\text{-}norm}(a, b, u, (\lambda z \in a.\; f(t(z)))) = \operatorname{linf\text{-}norm}(a_{2}, c, d, f) \end{array}}}
\vnbprop{measure-isomorphism-lp-isomorphism}{\vnbstmt{\begin{array}{@{}l@{}} \forall a, b, u, a_{2}, d, h, t. \\ \quad (\operatorname{is\text{-}measure}(a, b, u) \wedge (\operatorname{is\text{-}measure}(a_{2}, d, h) \wedge \operatorname{is\text{-}measure\text{-}isomorphism}(a, b, u, a_{2}, d, h, t))) \Rightarrow  \\ \quad \quad \forall f. \\ \quad \quad \quad \operatorname{is\text{-}real\text{-}measurable\text{-}fn}(a_{2}, d, f) \Rightarrow  \\ \quad \quad \quad \quad \operatorname{is\text{-}real\text{-}measurable\text{-}fn}(a, b, (\lambda z \in a.\; f(t(z)))) \wedge \\ \quad \quad \quad \quad \forall p \in \mathbb{R}. \\ \quad \quad \quad \quad \quad (1 \leq p) \Rightarrow  \\ \quad \quad \quad \quad \quad \quad \|{(\lambda z \in a.\; f(t(z)))}\|_{p} = \|{f}\|_{p} \wedge \\ \quad \quad \quad \quad \operatorname{linf\text{-}norm}(a, b, u, (\lambda z \in a.\; f(t(z)))) = \operatorname{linf\text{-}norm}(a_{2}, d, h, f) \wedge \\ \quad \quad \forall f, g \in (a_{2} \to \mathbb{R}), c \in \mathbb{R}, x \in a. \\ \quad \quad \quad (\lambda z \in a.\; (\lambda y \in a_{2}.\; (f(y) + (c \cdot g(y))))(t(z)))(x) \\ \quad \quad \quad =\; (\lambda z \in a.\; f(t(z)))(x) + (c \cdot (\lambda z \in a.\; g(t(z)))(x)) \wedge \\ \quad \quad \forall g. \\ \quad \quad \quad \operatorname{is\text{-}real\text{-}measurable\text{-}fn}(a, b, g) \Rightarrow  \\ \quad \quad \quad \quad \exists f. \\ \quad \quad \quad \quad \quad \operatorname{is\text{-}real\text{-}measurable\text{-}fn}(a_{2}, d, f) \wedge \forall x \in a.\; (g(x) = f(t(x))) \end{array}}}

\vnbbookfile{theorem-library/mb4-borel-subspace}{theorem-library/mb4-borel-subspace}
\noindent{\small mb4-borel-subspace.scm -- construction of measures, Proposition 5.2: the Borel sets of a subspace are the traces of the Borel sets (M-B4, 2026-10-26).  Library: measure.  Helper prefixes mbk- (the kit copied from theorem-library/push-forward.scm, pfw- renamed) and mbs-.}\par
\vnbprop{mbs-trace-borel}{\vnbstmt{\begin{array}{@{}l@{}} \forall s, y.; \forall e \in \operatorname{borel\text{-}sets}(s). \\ \quad (\operatorname{is\text{-}metric\text{-}space}(s) \wedge (y \subseteq \operatorname{pts}(s))) \Rightarrow  \\ \quad \quad (e \cap y) \in \operatorname{borel\text{-}sets}(\operatorname{subspace\text{-}ms}(s, y)) \end{array}}}
\vnbprop{mbs-borel-is-trace}{\vnbstmt{\begin{array}{@{}l@{}} \forall s, y.; \forall b \in \operatorname{borel\text{-}sets}(\operatorname{subspace\text{-}ms}(s, y)). \\ \quad (\operatorname{is\text{-}metric\text{-}space}(s) \wedge (y \subseteq \operatorname{pts}(s))) \Rightarrow  \\ \quad \quad \exists e \in \operatorname{borel\text{-}sets}(s).\; (b = (e \cap y)) \end{array}}}
\vnbprop{borel-sets-subspace}{\vnbstmt{\begin{array}{@{}l@{}} \forall s, y, b. \\ \quad (\operatorname{is\text{-}metric\text{-}space}(s) \wedge (y \subseteq \operatorname{pts}(s))) \Rightarrow  \\ \quad \quad (b \in \operatorname{borel\text{-}sets}(\operatorname{subspace\text{-}ms}(s, y))) \Leftrightarrow \exists e \in \operatorname{borel\text{-}sets}(s).\; (b = (e \cap y)) \end{array}}}

\vnbbookfile{theorem-library/mb4-borel-product}{theorem-library/mb4-borel-product}
\noindent{\small mb4-borel-product.scm -- construction of measures, Proposition 5.3: for separable metric spaces X, Y the Borel sets of X x Y (the max metric, PRODUCT2-MS) are the product sigma-algebra of the Borel sets (M-B4, 2026-10-26).  Library: measure.  Helper prefixes mbk- (kit copied from push-forward.scm) and mbp-.}\par
\vnbprop{mbp-rect-borel}{\vnbstmt{\begin{array}{@{}l@{}} \forall s, t.; \forall a \in \operatorname{borel\text{-}sets}(s), b \in \operatorname{borel\text{-}sets}(t). \\ \quad (\operatorname{is\text{-}metric\text{-}space}(s) \wedge \operatorname{is\text{-}metric\text{-}space}(t)) \Rightarrow  \\ \quad \quad (a \times b) \in \operatorname{borel\text{-}sets}(\operatorname{product2\text{-}ms}(s, t)) \end{array}}}
\vnbprop{mbp-product-in-borel}{\vnbstmt{\begin{array}{@{}l@{}} \forall s, t.; \forall e \in \operatorname{product\text{-}sigma}(\operatorname{pts}(s), \operatorname{borel\text{-}sets}(s), \operatorname{pts}(t), \operatorname{borel\text{-}sets}(t)). \\ \quad (\operatorname{is\text{-}metric\text{-}space}(s) \wedge \operatorname{is\text{-}metric\text{-}space}(t)) \Rightarrow  \\ \quad \quad e \in \operatorname{borel\text{-}sets}(\operatorname{product2\text{-}ms}(s, t)) \end{array}}}
\vnbprop{mbp-gen-in-sigma}{\vnbstmt{\forall a, e.; \forall \mathit{mpa} \in e.\; ((\mathit{mpa} \in \operatorname{power}(a)) \Rightarrow (\mathit{mpa} \in \operatorname{sigma\text{-}generated}(a, e)))}}
\vnbprop{mbp-open-in-product}{\vnbstmt{\begin{array}{@{}l@{}} \forall s, t, u. \\ \quad (\operatorname{is\text{-}separable}(s) \wedge (\operatorname{is\text{-}separable}(t) \wedge \operatorname{is\text{-}open}(\operatorname{product2\text{-}ms}(s, t), u))) \Rightarrow  \\ \quad \quad u \in \operatorname{product\text{-}sigma}(\operatorname{pts}(s), \operatorname{borel\text{-}sets}(s), \operatorname{pts}(t), \operatorname{borel\text{-}sets}(t)) \end{array}}}
\vnbprop{mbp-borel-in-product}{\vnbstmt{\begin{array}{@{}l@{}} \forall s, t.; \forall e \in \operatorname{borel\text{-}sets}(\operatorname{product2\text{-}ms}(s, t)). \\ \quad (\operatorname{is\text{-}separable}(s) \wedge \operatorname{is\text{-}separable}(t)) \Rightarrow  \\ \quad \quad e \in \operatorname{product\text{-}sigma}(\operatorname{pts}(s), \operatorname{borel\text{-}sets}(s), \operatorname{pts}(t), \operatorname{borel\text{-}sets}(t)) \end{array}}}
\vnbprop{borel-sets-product}{\vnbstmt{\begin{array}{@{}l@{}} \forall s, t. \\ \quad (\operatorname{is\text{-}separable}(s) \wedge \operatorname{is\text{-}separable}(t)) \Rightarrow  \\ \quad \quad \operatorname{borel\text{-}sets}(\operatorname{product2\text{-}ms}(s, t)) = \operatorname{product\text{-}sigma}(\operatorname{pts}(s), \operatorname{borel\text{-}sets}(s), \operatorname{pts}(t), \operatorname{borel\text{-}sets}(t)) \end{array}}}

\vnbbookfile{theorem-library/mb4-sigma-finite}{theorem-library/mb4-sigma-finite}
\noindent{\small mb4-sigma-finite.scm -- construction of measures, Proposition 7.7 (M-B4, 2026-10-26): a Borel measure on a separable metric space whose balls of radius r $<$ r0 have finite measure (inequality (192)) is sigma-finite: the balls of a fixed radius r $<$ r0 around a dense sequence cover X.  Library: measure.  Helper prefix ...}\par
\vnbprop{ball-finite-sigma-finite}{\vnbstmt{\begin{array}{@{}l@{}} \forall s, u.; \forall a \in \operatorname{rr\text{-}pos\text{-}star}. \\ \quad \operatorname{is\text{-}separable}(s) \wedge \\ \quad \operatorname{is\text{-}measure}(\operatorname{pts}(s), \operatorname{borel\text{-}sets}(s), u) \wedge \\ \quad 0 < a \wedge \\ \quad \forall x \in \operatorname{pts}(s), r. \\ \quad \quad (\operatorname{pos\text{-}rr}(r) \wedge (r < a)) \Rightarrow ((0 < u(\operatorname{ball}(s, x, r))) \wedge (u(\operatorname{ball}(s, x, r)) \in \mathbb{R})) \Rightarrow \\ \quad \quad \operatorname{is\text{-}sigma\text{-}finite}(\operatorname{pts}(s), \operatorname{borel\text{-}sets}(s), u) \end{array}}}

\vnbbookfile{theorem-library/mb4-separable-measure}{theorem-library/mb4-separable-measure}
\noindent{\small mb4-separable-measure.scm -- construction of measures, Corollary 7.6 (M-B4, 2026-10-26): under the assumptions of Proposition 7.1 with r0 = +inf (every ball has finite positive measure, and for each radius the measures of the balls of that radius are uniformly comparable), X is separable.}\par
\vnbprop{ball-packing-bounded}{\vnbstmt{\begin{array}{@{}l@{}} \forall s, u.; \forall b \in \operatorname{pts}(s), d \in \mathbb{R}, h.; \forall c \in (\mathbb{N} \to \operatorname{pts}(s)), a. \\ \quad \operatorname{is\text{-}metric\text{-}space}(s) \wedge \\ \quad \operatorname{is\text{-}measure}(\operatorname{pts}(s), \operatorname{borel\text{-}sets}(s), u) \wedge \\ \quad \operatorname{pos\text{-}rr}(h) \wedge \\ \quad \forall k \in \mathbb{N}.\; (c(k) \in \operatorname{ball}(s, b, d)) \wedge \\ \quad \forall j, k \in \mathbb{N}.\; ((j < k) \Rightarrow ((h + h) < \operatorname{dist}(s)(c(j), c(k)))) \wedge \\ \quad \operatorname{pos\text{-}rr}(a) \wedge \\ \quad \forall k \in \mathbb{N}.\; (a \leq u(\operatorname{ball}(s, c(k), h))) \Rightarrow \\ \quad \quad \neg (u(\operatorname{ball}(s, b, (d + (h + h)))) \in \mathbb{R}) \end{array}}}
\vnbprop{ball-greedy-net}{\vnbstmt{\begin{array}{@{}l@{}} \forall s, u.; \forall b \in \operatorname{pts}(s), d \in \mathbb{R}, l, h, a. \\ \quad \operatorname{is\text{-}metric\text{-}space}(s) \wedge \\ \quad \operatorname{is\text{-}measure}(\operatorname{pts}(s), \operatorname{borel\text{-}sets}(s), u) \wedge \\ \quad \operatorname{pos\text{-}rr}(l) \wedge \\ \quad \operatorname{pos\text{-}rr}(h) \wedge \\ \quad (h + h) < l \wedge \\ \quad u(\operatorname{ball}(s, b, (d + l))) \in \mathbb{R} \wedge \\ \quad \operatorname{pos\text{-}rr}(a) \wedge \\ \quad \forall x \in \operatorname{pts}(s).\; (a \leq u(\operatorname{ball}(s, x, h))) \Rightarrow \\ \quad \quad \exists f \in (\mathbb{N} \to \operatorname{pts}(s)). \forall p \in \operatorname{ball}(s, b, d). \exists i \in \mathbb{N}. \\ \quad \quad \quad p \in \operatorname{ball}(s, f(i), l) \end{array}}}
\vnbprop{measure-ball-separable}{\vnbstmt{\begin{array}{@{}l@{}} \forall s, u. \\ \quad \operatorname{is\text{-}complete}(s) \wedge \\ \quad \operatorname{is\text{-}measure}(\operatorname{pts}(s), \operatorname{borel\text{-}sets}(s), u) \wedge \\ \quad \forall x \in \operatorname{pts}(s), r. \\ \quad \quad \operatorname{pos\text{-}rr}(r) \Rightarrow ((0 < u(\operatorname{ball}(s, x, r))) \wedge (u(\operatorname{ball}(s, x, r)) \in \mathbb{R})) \wedge \\ \quad \forall r. \\ \quad \quad \operatorname{pos\text{-}rr}(r) \Rightarrow  \\ \quad \quad \quad \exists c. \\ \quad \quad \quad \quad (\operatorname{pos\text{-}rr}(c) \wedge \forall x, y \in \operatorname{pts}(s).\; ((c \cdot u(\operatorname{ball}(s, y, r))) \leq u(\operatorname{ball}(s, x, r)))) \wedge \\ \quad \exists z.\; (z \in \operatorname{pts}(s)) \Rightarrow \\ \quad \quad \operatorname{is\text{-}separable}(s) \end{array}}}

\vnbbookfile{structure-library/measure-ch1-vocab}{structure-library/measure-ch1-vocab}
\noindent{\small measure-ch1-vocab.scm -- ATOMS OF AN ALGEBRA, FINITE UNIONS, RECTANGLES OF R\textasciicircum{}n, THE PULLBACK SIGMA-ALGEBRA AND MEASURE (M-A, wave MAP-1026, 2026-10-26).  DEFINITIONS ONLY; destination structure-library/. Source: Thayer, Construction of Measures, ch. 1: s. 1 (minimal non-empty element, Prop. 1.2; Prop. 1.4; rectangles, ...}\par
\par\medskip\noindent\textbf{Definition} (\texttt{atoms}).\;$\operatorname{atoms}(\mathit{ca}) :=$ $\{\operatorname{mat\_} \in \mathit{ca} : \operatorname{is\text{-}atom}(\mathit{ca}, \operatorname{mat\_})\}$
\par\medskip\noindent\textbf{Definition} (\texttt{is-atom}).\;For all $\mathit{ca}$ and $a$, 
\begin{equation*}
\begin{array}{@{}l@{}} \operatorname{is\text{-}atom}(\mathit{ca}, a) \\ \Leftrightarrow\; a \in \mathit{ca} \wedge \\ \neg (a = \emptyset) \wedge \\ \forall \operatorname{mab\_} \in \mathit{ca}. \\ \quad (\operatorname{mab\_} \subseteq a) \Rightarrow (\neg (\operatorname{mab\_} = \emptyset) \Rightarrow (\operatorname{mab\_} = a)) \end{array}
\end{equation*}

\par\medskip\noindent\textbf{Definition} (\texttt{finite-unions}).\;$\operatorname{finite\text{-}unions}(\omega, \mathit{ci}) :=$
\begin{equation*}
\{\operatorname{mfu\_} \in \operatorname{power}(\omega) : \exists \operatorname{mff\_} \in \mathbf{Set}.\; ((|\operatorname{mff\_}| \in \mathbb{N}) \wedge ((\operatorname{mff\_} \subseteq \mathit{ci}) \wedge (\operatorname{mfu\_} = \operatorname{big\text{-}union}(\operatorname{mfw\_}, \operatorname{mff\_}, \operatorname{mfw\_}))))\}
\end{equation*}

\par\medskip\noindent\textbf{Definition} (\texttt{is-interval}).\;For every $\mathit{es}$, 
\begin{equation*}
\begin{array}{@{}l@{}} \operatorname{is\text{-}interval}(\mathit{es}) \\ \Leftrightarrow\; \mathit{es} \subseteq \mathbb{R} \wedge \\ \forall \operatorname{mia\_}, \operatorname{mib\_} \in \mathit{es}, \operatorname{mix\_} \in \mathbb{R}. \\ \quad (\operatorname{mia\_} \leq \operatorname{mix\_}) \Rightarrow ((\operatorname{mix\_} \leq \operatorname{mib\_}) \Rightarrow (\operatorname{mix\_} \in \mathit{es})) \end{array}
\end{equation*}

\par\medskip\noindent\textbf{Definition} (\texttt{rn-box}).\;$\operatorname{rn\text{-}box}(n, \mathit{jj}) :=$
\begin{equation*}
\{\operatorname{mrx\_} \in (\operatorname{interval}(1, n) \to \mathbb{R}) : \forall \operatorname{mri\_} \in \operatorname{interval}(1, n).\; (\operatorname{mrx\_}(\operatorname{mri\_}) \in \operatorname{jj}(\operatorname{mri\_}))\}
\end{equation*}

\par\medskip\noindent\textbf{Definition} (\texttt{rn-rectangles}).\;$\operatorname{rn\text{-}rectangles}(n) :=$
\begin{equation*}
\{\operatorname{mrr\_} \in \operatorname{power}((\operatorname{interval}(1, n) \to \mathbb{R})) : \exists \operatorname{mrj\_} \in (\operatorname{interval}(1, n) \to \operatorname{power}(\mathbb{R})).\; (\forall \operatorname{mri\_} \in \operatorname{interval}(1, n).\; \operatorname{is\text{-}interval}(\operatorname{mrj\_}(\operatorname{mri\_})) \wedge (\operatorname{mrr\_} = \operatorname{rn\text{-}box}(n, \operatorname{mrj\_})))\}
\end{equation*}

\par\medskip\noindent\textbf{Definition} (\texttt{pullback-sigma}).\;$\operatorname{pullback\text{-}sigma}(\omega, \mathit{cb}, f) :=$
\begin{equation*}
\{\operatorname{mpa\_} \in \operatorname{power}(\omega) : \exists \operatorname{mpb\_} \in \mathit{cb}.\; (\operatorname{mpa\_} = \{\operatorname{mpx\_} \in \omega : (f(\operatorname{mpx\_}) \in \operatorname{mpb\_})\})\}
\end{equation*}

\par\medskip\noindent\textbf{Definition} (\texttt{pullback-measure}).\;$\operatorname{pullback\text{-}measure}(\omega, \mathit{cb}, f, \nu) :=$
\begin{equation*}
(\lambda \operatorname{mpm\_} \in \operatorname{pullback\text{-}sigma}(\omega, \mathit{cb}, f).\; \nu(\operatorname{image}(f, \operatorname{mpm\_})))
\end{equation*}


\vnbbookfile{theorem-library/measure-ch1-algebras}{theorem-library/measure-ch1-algebras}
\noindent{\small theorem-library/measure-ch1-algebras.scm -- ALGEBRAS OF SETS: atoms, finite unions, rectangles of R\textasciicircum{}n (M-A, wave MAP-1026, 2026-10-26).  Vocabulary: structure-library/measure-ch1-vocab.scm. Source: Thayer, Construction of Measures, ch. 1 s. 1.  Helper prefix: ma-.}\par
\vnbprop{algebra-finite-meet-in}{\vnbstmt{\begin{array}{@{}l@{}} \forall a, b.; \forall s \in \mathbf{Set}. \\ \quad (\operatorname{is\text{-}algebra\text{-}of\text{-}sets}(a, b) \wedge ((|s| \in \mathbb{N}) \wedge (s \subseteq b))) \Rightarrow  \\ \quad \quad \{y \in a : \forall \mathit{mmb} \in s.\; (y \in \mathit{mmb})\} \in b \end{array}}}
\vnbprop{atoms-unfold}{\vnbstmt{\forall a.\; (\operatorname{atoms}(a) \simeq \{t \in a : \operatorname{is\text{-}atom}(a, t)\})}}
\vnbprop{atoms-membership}{\vnbstmt{\forall b, a.\; ((a \in \operatorname{atoms}(b)) \Leftrightarrow \operatorname{is\text{-}atom}(b, a))}}
\vnbprop{algebra-point-atom}{\vnbstmt{\begin{array}{@{}l@{}} \forall a, b.; \forall x \in a. \\ \quad (\operatorname{is\text{-}algebra\text{-}of\text{-}sets}(a, b) \wedge (|b| \in \mathbb{N})) \Rightarrow  \\ \quad \quad \{y \in a : \forall \mathit{mmb} \in \{c \in b : (x \in c)\}.\; (y \in \mathit{mmb})\} \in \operatorname{atoms}(b) \wedge \\ \quad \quad x \in \{y \in a : \forall \mathit{mmb} \in \{c \in b : (x \in c)\}.\; (y \in \mathit{mmb})\} \wedge \\ \quad \quad \forall \mathit{mqb} \in b. \\ \quad \quad \quad (x \in \mathit{mqb}) \Rightarrow  \\ \quad \quad \quad \quad \{y \in a : \forall \mathit{mmb} \in \{c \in b : (x \in c)\}.\; (y \in \mathit{mmb})\} \subseteq \mathit{mqb} \end{array}}}
\vnbprop{atoms-pairwise-disjoint}{\vnbstmt{\begin{array}{@{}l@{}} \forall c, d, a, b. \\ \quad (\operatorname{is\text{-}algebra\text{-}of\text{-}sets}(c, d) \wedge ((a \in \operatorname{atoms}(d)) \wedge ((b \in \operatorname{atoms}(d)) \wedge \neg (a = b)))) \Rightarrow  \\ \quad \quad (a \cap b) = \emptyset \end{array}}}
\vnbprop{algebra-atom-decomposition}{\vnbstmt{\begin{array}{@{}l@{}} \forall b, c.; \forall a \in c. \\ \quad (\operatorname{is\text{-}algebra\text{-}of\text{-}sets}(b, c) \wedge (|c| \in \mathbb{N})) \Rightarrow  \\ \quad \quad ((a = \operatorname{big\text{-}union}(u, \{s \in \operatorname{atoms}(c) : (s \subseteq a)\}, u)) \wedge (a = \operatorname{big\text{-}union}(u, \operatorname{atoms}(c), (a \cap u)))) \end{array}}}
\vnbprop{finite-algebra-atoms}{\vnbstmt{\begin{array}{@{}l@{}} \forall c, d. \\ \quad (\operatorname{is\text{-}algebra\text{-}of\text{-}sets}(c, d) \wedge (|d| \in \mathbb{N})) \Rightarrow  \\ \quad \quad \operatorname{atoms}(d) \in \mathbf{Set} \wedge \\ \quad \quad |\operatorname{atoms}(d)| \in \mathbb{N} \wedge \\ \quad \quad \forall a, b. \\ \quad \quad \quad ((a \in \operatorname{atoms}(d)) \wedge ((b \in \operatorname{atoms}(d)) \wedge \neg (a = b))) \Rightarrow  \\ \quad \quad \quad \quad (a \cap b) = \emptyset \wedge \\ \quad \quad \forall a \in d. \\ \quad \quad \quad ((a = \operatorname{big\text{-}union}(u, \{s \in \operatorname{atoms}(d) : (s \subseteq a)\}, u)) \wedge (a = \operatorname{big\text{-}union}(u, \operatorname{atoms}(d), (a \cap u)))) \end{array}}}
\vnbprop{big-union-of-empty-family}{\vnbstmt{\operatorname{big\text{-}union}(w, \emptyset, w) = \emptyset}}
\vnbprop{big-union-of-singleton-family}{\vnbstmt{\forall x \in \mathbf{Set}.\; (\operatorname{big\text{-}union}(w, \{x\}, w) = x)}}
\vnbprop{big-union-of-union-family}{\vnbstmt{\begin{array}{@{}l@{}} \forall f, g. \\ \quad \operatorname{big\text{-}union}(w, (f \cup g), w) = (\operatorname{big\text{-}union}(w, f, w) \cup \operatorname{big\text{-}union}(w, g, w)) \end{array}}}
\vnbprop{finite-unions-unfold}{\vnbstmt{\begin{array}{@{}l@{}} \forall a, i. \\ \quad (\operatorname{finite\text{-}unions}(a, i) \simeq \{u \in \operatorname{power}(a) : \exists f \in \mathbf{Set}.\; ((|f| \in \mathbb{N}) \wedge ((f \subseteq i) \wedge (u = \operatorname{big\text{-}union}(w, f, w))))\}) \end{array}}}
\vnbprop{finite-unions-intro}{\vnbstmt{\begin{array}{@{}l@{}} \forall a, i.; \forall f \in \mathbf{Set}. \\ \quad ((a \in \mathbf{Set}) \wedge ((i \subseteq \operatorname{power}(a)) \wedge ((|f| \in \mathbb{N}) \wedge (f \subseteq i)))) \Rightarrow  \\ \quad \quad \operatorname{big\text{-}union}(w, f, w) \in \operatorname{finite\text{-}unions}(a, i) \end{array}}}
\vnbprop{finite-unions-elim}{\vnbstmt{\begin{array}{@{}l@{}} \forall a, i.; \forall s \in \operatorname{finite\text{-}unions}(a, i). \\ \quad ((s \in \operatorname{power}(a)) \wedge \exists f \in \mathbf{Set}.\; ((|f| \in \mathbb{N}) \wedge ((f \subseteq i) \wedge (s = \operatorname{big\text{-}union}(w, f, w))))) \end{array}}}
\vnbprop{finite-unions-empty-in}{\vnbstmt{\begin{array}{@{}l@{}} \forall a, i. \\ \quad ((a \in \mathbf{Set}) \wedge (i \subseteq \operatorname{power}(a))) \Rightarrow (\emptyset \in \operatorname{finite\text{-}unions}(a, i)) \end{array}}}
\vnbprop{finite-unions-member-in}{\vnbstmt{\begin{array}{@{}l@{}} \forall b, i.; \forall a \in i. \\ \quad ((b \in \mathbf{Set}) \wedge (i \subseteq \operatorname{power}(b))) \Rightarrow (a \in \operatorname{finite\text{-}unions}(b, i)) \end{array}}}
\vnbprop{finite-unions-union-closed}{\vnbstmt{\begin{array}{@{}l@{}} \forall a, i, s, t. \\ \quad ((a \in \mathbf{Set}) \wedge ((i \subseteq \operatorname{power}(a)) \wedge ((s \in \operatorname{finite\text{-}unions}(a, i)) \wedge (t \in \operatorname{finite\text{-}unions}(a, i))))) \Rightarrow  \\ \quad \quad (s \cup t) \in \operatorname{finite\text{-}unions}(a, i) \end{array}}}
\vnbprop{finite-unions-inter-member}{\vnbstmt{\begin{array}{@{}l@{}} \forall b, i, a, t. \\ \quad b \in \mathbf{Set} \wedge \\ \quad i \subseteq \operatorname{power}(b) \wedge \\ \quad \forall \mathit{mia}, \mathit{mib}.\; (((\mathit{mia} \in i) \wedge (\mathit{mib} \in i)) \Rightarrow ((\mathit{mia} \cap \mathit{mib}) \in i)) \wedge \\ \quad a \in i \wedge \\ \quad t \in \operatorname{finite\text{-}unions}(b, i) \Rightarrow \\ \quad \quad (a \cap t) \in \operatorname{finite\text{-}unions}(b, i) \end{array}}}
\vnbprop{finite-unions-inter-closed}{\vnbstmt{\begin{array}{@{}l@{}} \forall a, i, s, t. \\ \quad a \in \mathbf{Set} \wedge \\ \quad i \subseteq \operatorname{power}(a) \wedge \\ \quad \forall \mathit{mia}, b.\; (((\mathit{mia} \in i) \wedge (b \in i)) \Rightarrow ((\mathit{mia} \cap b) \in i)) \wedge \\ \quad s \in \operatorname{finite\text{-}unions}(a, i) \wedge \\ \quad t \in \operatorname{finite\text{-}unions}(a, i) \Rightarrow \\ \quad \quad (s \cap t) \in \operatorname{finite\text{-}unions}(a, i) \end{array}}}
\vnbprop{finite-unions-complement-closed}{\vnbstmt{\begin{array}{@{}l@{}} \forall a, i.; \forall s \in \operatorname{finite\text{-}unions}(a, i). \\ \quad a \in \mathbf{Set} \wedge \\ \quad i \subseteq \operatorname{power}(a) \wedge \\ \quad \exists e.\; (e \in i) \wedge \\ \quad \forall \mathit{mia}, b.\; (((\mathit{mia} \in i) \wedge (b \in i)) \Rightarrow ((\mathit{mia} \cap b) \in i)) \wedge \\ \quad \forall \mathit{mia} \in i.\; (\operatorname{complement\text{-}in}(a, \mathit{mia}) \in \operatorname{finite\text{-}unions}(a, i)) \Rightarrow \\ \quad \quad \operatorname{complement\text{-}in}(a, s) \in \operatorname{finite\text{-}unions}(a, i) \end{array}}}
\vnbprop{finite-unions-algebra}{\vnbstmt{\begin{array}{@{}l@{}} \forall a, i. \\ \quad a \in \mathbf{Set} \wedge \\ \quad i \subseteq \operatorname{power}(a) \wedge \\ \quad \exists e.\; (e \in i) \wedge \\ \quad \forall \mathit{mia}, b.\; (((\mathit{mia} \in i) \wedge (b \in i)) \Rightarrow ((\mathit{mia} \cap b) \in i)) \wedge \\ \quad \forall \mathit{mia} \in i.\; (\operatorname{complement\text{-}in}(a, \mathit{mia}) \in \operatorname{finite\text{-}unions}(a, i)) \Rightarrow \\ \quad \quad \operatorname{is\text{-}algebra\text{-}of\text{-}sets}(a, \operatorname{finite\text{-}unions}(a, i)) \end{array}}}
\vnbprop{is-interval-rr}{\vnbstmt{\operatorname{is\text{-}interval}(\mathbb{R})}}
\vnbprop{is-interval-inter}{\vnbstmt{\begin{array}{@{}l@{}} \forall s, a. \\ \quad (\operatorname{is\text{-}interval}(s) \wedge \operatorname{is\text{-}interval}(a)) \Rightarrow \operatorname{is\text{-}interval}((s \cap a)) \end{array}}}
\vnbprop{rn-ray-below-interval}{\vnbstmt{\begin{array}{@{}l@{}} \forall j, i, k. \\ \quad (j(i) \subseteq \mathbb{R}) \Rightarrow  \\ \quad \quad \operatorname{is\text{-}interval}(\{z \in \mathbb{R} : \forall y \in j(i).\; ((i = k) \Rightarrow (z < y))\}) \end{array}}}
\vnbprop{rn-ray-above-interval}{\vnbstmt{\begin{array}{@{}l@{}} \forall j, i, k. \\ \quad (j(i) \subseteq \mathbb{R}) \Rightarrow  \\ \quad \quad \operatorname{is\text{-}interval}(\{z \in \mathbb{R} : \forall y \in j(i).\; ((i = k) \Rightarrow (y < z))\}) \end{array}}}
\vnbprop{interval-outside-below-or-above}{\vnbstmt{\begin{array}{@{}l@{}} \forall s, x. \\ \quad (\operatorname{is\text{-}interval}(s) \wedge ((x \in \mathbb{R}) \wedge \neg (x \in s))) \Rightarrow  \\ \quad \quad \forall y \in s.\; (x < y) \vee \forall y \in s.\; (y < x) \end{array}}}
\vnbprop{rn-box-unfold}{\vnbstmt{\begin{array}{@{}l@{}} \forall n, j. \\ \quad (\operatorname{rn\text{-}box}(n, j) \simeq \{x \in (\operatorname{interval}(1, n) \to \mathbb{R}) : \forall i \in \operatorname{interval}(1, n).\; (x(i) \in j(i))\}) \end{array}}}
\vnbprop{rn-rectangles-unfold}{\vnbstmt{\begin{array}{@{}l@{}} \forall n. \\ \quad (\operatorname{rn\text{-}rectangles}(n) \simeq \{r \in \operatorname{power}((\operatorname{interval}(1, n) \to \mathbb{R})) : \exists j \in (\operatorname{interval}(1, n) \to \operatorname{power}(\mathbb{R})).\; (\forall i \in \operatorname{interval}(1, n).\; \operatorname{is\text{-}interval}(j(i)) \wedge (r = \operatorname{rn\text{-}box}(n, j)))\}) \end{array}}}
\vnbprop{rn-is-set}{\vnbstmt{\forall n.\; ((\operatorname{interval}(1, n) \in \mathbf{Set}) \wedge ((\operatorname{interval}(1, n) \to \mathbb{R}) \in \mathbf{Set}))}}
\vnbprop{rn-box-membership}{\vnbstmt{\begin{array}{@{}l@{}} \forall n, x, j. \\ \quad (x \in \operatorname{rn\text{-}box}(n, j)) \Leftrightarrow ((x \in (\operatorname{interval}(1, n) \to \mathbb{R})) \wedge \forall i \in \operatorname{interval}(1, n).\; (x(i) \in j(i))) \end{array}}}
\vnbprop{rn-box-in-power}{\vnbstmt{\forall n, j.\; (\operatorname{rn\text{-}box}(n, j) \in \operatorname{power}((\operatorname{interval}(1, n) \to \mathbb{R})))}}
\vnbprop{rn-rectangle-intro}{\vnbstmt{\begin{array}{@{}l@{}} \forall n.; \forall j \in (\operatorname{interval}(1, n) \to \operatorname{power}(\mathbb{R})). \\ \quad \forall i \in \operatorname{interval}(1, n).\; \operatorname{is\text{-}interval}(j(i)) \Rightarrow  \\ \quad \quad \operatorname{rn\text{-}box}(n, j) \in \operatorname{rn\text{-}rectangles}(n) \end{array}}}
\vnbprop{rn-rectangle-elim}{\vnbstmt{\begin{array}{@{}l@{}} \forall n.; \forall r \in \operatorname{rn\text{-}rectangles}(n). \\ \quad r \in \operatorname{power}((\operatorname{interval}(1, n) \to \mathbb{R})) \wedge \\ \quad \exists j \in (\operatorname{interval}(1, n) \to \operatorname{power}(\mathbb{R})). \\ \quad \quad \forall i \in \operatorname{interval}(1, n).\; \operatorname{is\text{-}interval}(j(i)) \wedge (r = \operatorname{rn\text{-}box}(n, j)) \end{array}}}
\vnbprop{rn-rectangles-inter-closed}{\vnbstmt{\begin{array}{@{}l@{}} \forall n, a, b. \\ \quad ((a \in \operatorname{rn\text{-}rectangles}(n)) \wedge (b \in \operatorname{rn\text{-}rectangles}(n))) \Rightarrow  \\ \quad \quad (a \cap b) \in \operatorname{rn\text{-}rectangles}(n) \end{array}}}
\vnbprop{rn-ray-below-rectangle}{\vnbstmt{\begin{array}{@{}l@{}} \forall n, j, k. \\ \quad ((j \in (\operatorname{interval}(1, n) \to \operatorname{power}(\mathbb{R}))) \wedge \forall i \in \operatorname{interval}(1, n).\; \operatorname{is\text{-}interval}(j(i))) \Rightarrow  \\ \quad \quad \operatorname{rn\text{-}box}(n, (\lambda \mathit{mki} \in \operatorname{interval}(1, n).\; \{z \in \mathbb{R} : \forall y \in j(\mathit{mki}).\; ((\mathit{mki} = k) \Rightarrow (z < y))\})) \\ \quad \quad \in\; \operatorname{rn\text{-}rectangles}(n) \end{array}}}
\vnbprop{rn-ray-above-rectangle}{\vnbstmt{\begin{array}{@{}l@{}} \forall n, j, k. \\ \quad ((j \in (\operatorname{interval}(1, n) \to \operatorname{power}(\mathbb{R}))) \wedge \forall i \in \operatorname{interval}(1, n).\; \operatorname{is\text{-}interval}(j(i))) \Rightarrow  \\ \quad \quad \operatorname{rn\text{-}box}(n, (\lambda \mathit{mki} \in \operatorname{interval}(1, n).\; \{z \in \mathbb{R} : \forall y \in j(\mathit{mki}).\; ((\mathit{mki} = k) \Rightarrow (y < z))\})) \\ \quad \quad \in\; \operatorname{rn\text{-}rectangles}(n) \end{array}}}
\vnbprop{rn-rectangles-subset-power}{\vnbstmt{\forall n.\; (\operatorname{rn\text{-}rectangles}(n) \subseteq \operatorname{power}((\operatorname{interval}(1, n) \to \mathbb{R})))}}
\vnbprop{rn-rectangle-complement-finite-union}{\vnbstmt{\begin{array}{@{}l@{}} \forall n \in \mathbb{N}, j \in (\operatorname{interval}(1, n) \to \operatorname{power}(\mathbb{R})). \\ \quad \forall i \in \operatorname{interval}(1, n).\; \operatorname{is\text{-}interval}(j(i)) \Rightarrow  \\ \quad \quad \operatorname{complement\text{-}in}((\operatorname{interval}(1, n) \to \mathbb{R}), \operatorname{rn\text{-}box}(n, j)) \\ \quad \quad \in\; \operatorname{finite\text{-}unions}((\operatorname{interval}(1, n) \to \mathbb{R}), \operatorname{rn\text{-}rectangles}(n)) \end{array}}}
\vnbprop{rn-rectangles-finite-unions-algebra}{\vnbstmt{\begin{array}{@{}l@{}} \forall n \in \mathbb{N}. \\ \quad \operatorname{is\text{-}algebra\text{-}of\text{-}sets}((\operatorname{interval}(1, n) \to \mathbb{R}), \\ \quad \quad \operatorname{finite\text{-}unions}((\operatorname{interval}(1, n) \to \mathbb{R}), \operatorname{rn\text{-}rectangles}(n))) \end{array}}}

\vnbbookfile{theorem-library/measure-ch1-pullback}{theorem-library/measure-ch1-pullback}
\noindent{\small theorem-library/measure-ch1-pullback.scm -- THE PULLBACK SIGMA-ALGEBRA AND THE PULLBACK MEASURE (M-A, wave MAP-1026, 2026-10-26).  Vocabulary: structure-library/measure-ch1-vocab.scm (PULLBACK-SIGMA, PULLBACK-MEASURE).  Source: Thayer, Construction of Measures, ch. 1, Prop. 1.14 and Prop. 1.22. Helper prefix: ma-.}\par
\vnbprop{pullback-sigma-unfold}{\vnbstmt{\begin{array}{@{}l@{}} \forall a, b, f. \\ \quad (\operatorname{pullback\text{-}sigma}(a, b, f) \simeq \{\mathit{mpa} \in \operatorname{power}(a) : \exists \mathit{mpb} \in b.\; (\mathit{mpa} = \{x \in a : (f(x) \in \mathit{mpb})\})\}) \end{array}}}
\vnbprop{pullback-sigma-membership}{\vnbstmt{\begin{array}{@{}l@{}} \forall b, c, f, a. \\ \quad (a \in \operatorname{pullback\text{-}sigma}(b, c, f)) \Leftrightarrow ((a \in \operatorname{power}(b)) \wedge \exists \mathit{mpb} \in c.\; (a = \{x \in b : (f(x) \in \mathit{mpb})\})) \end{array}}}
\vnbprop{pullback-sigma-is-sigma-algebra}{\vnbstmt{\begin{array}{@{}l@{}} \forall a, b, a_{2}, f. \\ \quad (\operatorname{is\text{-}sigma\text{-}algebra}(a_{2}, b) \wedge (f \in (a \to a_{2}))) \Rightarrow  \\ \quad \quad \operatorname{is\text{-}sigma\text{-}algebra}(a, \operatorname{pullback\text{-}sigma}(a, b, f)) \end{array}}}
\vnbprop{pullback-sigma-measurable}{\vnbstmt{\begin{array}{@{}l@{}} \forall a, b, a_{2}, f. \\ \quad (\operatorname{is\text{-}sigma\text{-}algebra}(a_{2}, b) \wedge (f \in (a \to a_{2}))) \Rightarrow  \\ \quad \quad \operatorname{is\text{-}measurable\text{-}map}(a, \operatorname{pullback\text{-}sigma}(a, b, f), a_{2}, b, f) \end{array}}}
\vnbprop{pullback-sigma-smallest}{\vnbstmt{\begin{array}{@{}l@{}} \forall a, b, a_{2}, f, c. \\ \quad \operatorname{is\text{-}sigma\text{-}algebra}(a_{2}, b) \wedge \\ \quad f \in (a \to a_{2}) \wedge \\ \quad \operatorname{is\text{-}sigma\text{-}algebra}(a, c) \wedge \\ \quad \operatorname{is\text{-}measurable\text{-}map}(a, c, a_{2}, b, f) \Rightarrow \\ \quad \quad \operatorname{pullback\text{-}sigma}(a, b, f) \subseteq c \end{array}}}
\vnbprop{pullback-sigma-algebra}{\vnbstmt{\begin{array}{@{}l@{}} \forall a, b, a_{2}, f. \\ \quad (\operatorname{is\text{-}sigma\text{-}algebra}(a_{2}, b) \wedge (f \in (a \to a_{2}))) \Rightarrow  \\ \quad \quad \operatorname{is\text{-}sigma\text{-}algebra}(a, \operatorname{pullback\text{-}sigma}(a, b, f)) \wedge \\ \quad \quad \operatorname{is\text{-}measurable\text{-}map}(a, \operatorname{pullback\text{-}sigma}(a, b, f), a_{2}, b, f) \wedge \\ \quad \quad \forall c. \\ \quad \quad \quad (\operatorname{is\text{-}sigma\text{-}algebra}(a, c) \wedge \operatorname{is\text{-}measurable\text{-}map}(a, c, a_{2}, b, f)) \Rightarrow  \\ \quad \quad \quad \quad \operatorname{pullback\text{-}sigma}(a, b, f) \subseteq c \end{array}}}
\vnbprop{pullback-measure-unfold}{\vnbstmt{\begin{array}{@{}l@{}} \forall a, b, f, u. \\ \quad (\operatorname{pullback\text{-}measure}(a, b, f, u) \simeq (\lambda m \in \operatorname{pullback\text{-}sigma}(a, b, f).\; u(\operatorname{image}(f, m)))) \end{array}}}
\vnbprop{image-of-empty-set}{\vnbstmt{\forall f.\; (\operatorname{image}(f, \emptyset) = \emptyset)}}
\vnbprop{image-preimage-of-surjection}{\vnbstmt{\begin{array}{@{}l@{}} \forall a, a_{2}.; \forall f \in (a \to a_{2}), b. \\ \quad (\forall y \in a_{2}.\; \exists x \in a.\; (f(x) = y) \wedge (b \subseteq a_{2})) \Rightarrow  \\ \quad \quad \operatorname{image}(f, \{\mathit{mpx} \in a : (f(\mathit{mpx}) \in b)\}) = b \end{array}}}
\vnbprop{image-of-countable-union}{\vnbstmt{\begin{array}{@{}l@{}} \forall f, g, h. \\ \quad \forall k \in \mathbb{N}.\; (h(k) \simeq \operatorname{image}(f, g(k))) \Rightarrow  \\ \quad \quad \operatorname{image}(f, \operatorname{big\text{-}union}(n, \mathbb{N}, g(n))) = \operatorname{big\text{-}union}(n, \mathbb{N}, h(n)) \end{array}}}
\vnbprop{pullback-measure-at}{\vnbstmt{\begin{array}{@{}l@{}} \forall b, c, a_{2}, f, u.; \forall a \in \operatorname{pullback\text{-}sigma}(b, c, f). \\ \quad (\operatorname{is\text{-}measure}(a_{2}, c, u) \wedge ((f \in (b \to a_{2})) \wedge \forall y \in a_{2}.\; \exists x \in b.\; (f(x) = y))) \Rightarrow  \\ \quad \quad (\operatorname{image}(f, a) \in c) \wedge (\operatorname{pullback\text{-}measure}(b, c, f, u)(a) = u(\operatorname{image}(f, a))) \end{array}}}
\vnbprop{pullback-measure-value}{\vnbstmt{\begin{array}{@{}l@{}} \forall a, c, a_{2}, f, u.; \forall b \in c. \\ \quad (\operatorname{is\text{-}measure}(a_{2}, c, u) \wedge ((f \in (a \to a_{2})) \wedge \forall y \in a_{2}.\; \exists x \in a.\; (f(x) = y))) \Rightarrow  \\ \quad \quad \operatorname{pullback\text{-}measure}(a, c, f, u)(\{\mathit{mpx} \in a : (f(\mathit{mpx}) \in b)\}) = u(b) \end{array}}}
\vnbprop{pullback-measure-type}{\vnbstmt{\begin{array}{@{}l@{}} \forall a, b, a_{2}, f, u. \\ \quad (\operatorname{is\text{-}measure}(a_{2}, b, u) \wedge ((f \in (a \to a_{2})) \wedge \forall y \in a_{2}.\; \exists x \in a.\; (f(x) = y))) \Rightarrow  \\ \quad \quad \operatorname{pullback\text{-}measure}(a, b, f, u) \in (\operatorname{pullback\text{-}sigma}(a, b, f) \to \operatorname{rr\text{-}pos\text{-}star}) \end{array}}}
\vnbprop{pullback-measure-is-measure}{\vnbstmt{\begin{array}{@{}l@{}} \forall a, b, a_{2}, f, u. \\ \quad (\operatorname{is\text{-}measure}(a_{2}, b, u) \wedge ((f \in (a \to a_{2})) \wedge \forall y \in a_{2}.\; \exists x \in a.\; (f(x) = y))) \Rightarrow  \\ \quad \quad \operatorname{is\text{-}measure}(a, \operatorname{pullback\text{-}sigma}(a, b, f), \operatorname{pullback\text{-}measure}(a, b, f, u)) \end{array}}}
\vnbprop{pullback-measure-unique}{\vnbstmt{\begin{array}{@{}l@{}} \forall a, c, a_{2}, f, u, d. \\ \quad \operatorname{is\text{-}measure}(a_{2}, c, u) \wedge \\ \quad f \in (a \to a_{2}) \wedge \\ \quad \forall y \in a_{2}.\; \exists x \in a.\; (f(x) = y) \wedge \\ \quad \operatorname{is\text{-}measure}(a, \operatorname{pullback\text{-}sigma}(a, c, f), d) \wedge \\ \quad \forall b \in c.\; (d(\{\mathit{mpx} \in a : (f(\mathit{mpx}) \in b)\}) = u(b)) \Rightarrow \\ \quad \quad d = \operatorname{pullback\text{-}measure}(a, c, f, u) \end{array}}}
\vnbprop{pullback-measure-exists-unique}{\vnbstmt{\begin{array}{@{}l@{}} \forall a, c, a_{2}, f, u. \\ \quad (\operatorname{is\text{-}measure}(a_{2}, c, u) \wedge ((f \in (a \to a_{2})) \wedge \forall y \in a_{2}.\; \exists x \in a.\; (f(x) = y))) \Rightarrow  \\ \quad \quad \operatorname{is\text{-}measure}(a, \operatorname{pullback\text{-}sigma}(a, c, f), \operatorname{pullback\text{-}measure}(a, c, f, u)) \wedge \\ \quad \quad \forall b \in c.\; (\operatorname{pullback\text{-}measure}(a, c, f, u)(\{\mathit{mpx} \in a : (f(\mathit{mpx}) \in b)\}) = u(b)) \wedge \\ \quad \quad \forall d. \\ \quad \quad \quad (\operatorname{is\text{-}measure}(a, \operatorname{pullback\text{-}sigma}(a, c, f), d) \wedge \forall b \in c.\; (d(\{\mathit{mpx} \in a : (f(\mathit{mpx}) \in b)\}) = u(b))) \Rightarrow  \\ \quad \quad \quad \quad d = \operatorname{pullback\text{-}measure}(a, c, f, u) \end{array}}}

\vnbbookfile{theorem-library/measure-ch1-vitali}{theorem-library/measure-ch1-vitali}
\noindent{\small theorem-library/measure-ch1-vitali.scm -- NO TRANSLATION-INVARIANT MEASURE ON ALL SUBSETS OF R (VITALI) (M-A, wave MAP-1026, 2026-10-26).  Source: Thayer, Construction of Measures, ch. 1, Prop. 1.23. Vocabulary: TRANSLATE (structure-library/lebesgue-translate.scm).  Helper prefix: ma- (kit), mv- (Vitali).}\par
\vnbprop{vitali-d-refl}{\vnbstmt{\begin{array}{@{}l@{}} \forall x \in \mathbb{R}. \\ \quad (x - x) \in \{d \in \mathbb{R} : \exists a, b, c \in \mathbb{N}.\; ((a - b) = ((c + 1) \cdot d))\} \end{array}}}
\vnbprop{vitali-d-symm}{\vnbstmt{\begin{array}{@{}l@{}} \forall x, y. \\ \quad x \in \mathbb{R} \wedge \\ \quad y \in \mathbb{R} \wedge \\ \quad (y - x) \in \{d \in \mathbb{R} : \exists a, b, c \in \mathbb{N}.\; ((a - b) = ((c + 1) \cdot d))\} \Rightarrow \\ \quad \quad (x - y) \in \{d \in \mathbb{R} : \exists a, b, c \in \mathbb{N}.\; ((a - b) = ((c + 1) \cdot d))\} \end{array}}}
\vnbprop{vitali-d-trans}{\vnbstmt{\begin{array}{@{}l@{}} \forall x, y, z. \\ \quad x \in \mathbb{R} \wedge \\ \quad y \in \mathbb{R} \wedge \\ \quad z \in \mathbb{R} \wedge \\ \quad (x - y) \in \{d \in \mathbb{R} : \exists a, b, c \in \mathbb{N}.\; ((a - b) = ((c + 1) \cdot d))\} \wedge \\ \quad (y - z) \in \{d \in \mathbb{R} : \exists a, b, c \in \mathbb{N}.\; ((a - b) = ((c + 1) \cdot d))\} \Rightarrow \\ \quad \quad (x - z) \in \{d \in \mathbb{R} : \exists a, b, c \in \mathbb{N}.\; ((a - b) = ((c + 1) \cdot d))\} \end{array}}}
\vnbprop{vitali-recip-facts}{\vnbstmt{\forall n \in \mathbb{N}.\; ((\operatorname{recip}((n + 1)) \in \mathbb{R}) \wedge ((0 < \operatorname{recip}((n + 1))) \wedge (\operatorname{recip}((n + 1)) \leq 1)))}}
\vnbprop{vitali-recip-difference-in-d}{\vnbstmt{\begin{array}{@{}l@{}} \forall i, j. \\ \quad ((i \in \mathbb{N}) \wedge (j \in \mathbb{N})) \Rightarrow  \\ \quad \quad (\operatorname{recip}((j + 1)) - \operatorname{recip}((i + 1))) \in \{d \in \mathbb{R} : \exists a, b, c \in \mathbb{N}.\; ((a - b) = ((c + 1) \cdot d))\} \end{array}}}
\vnbprop{vitali-recip-injective}{\vnbstmt{\begin{array}{@{}l@{}} \forall i, j. \\ \quad ((i \in \mathbb{N}) \wedge ((j \in \mathbb{N}) \wedge (\operatorname{recip}((i + 1)) = \operatorname{recip}((j + 1))))) \Rightarrow (i = j) \end{array}}}
\vnbprop{vitali-class-eq}{\vnbstmt{\begin{array}{@{}l@{}} \forall a, b. \\ \quad a \in \mathbb{R} \wedge \\ \quad b \in \mathbb{R} \wedge \\ \quad (a - b) \in \{d \in \mathbb{R} : \exists \mathit{mva}, \mathit{mvb}, c \in \mathbb{N}.\; ((\mathit{mva} - \mathit{mvb}) = ((c + 1) \cdot d))\} \Rightarrow \\ \quad \quad \{y \in \operatorname{ccint}(0, 1) : ((y - a) \in \{d \in \mathbb{R} : \exists \mathit{mva}, \mathit{mvb}, c \in \mathbb{N}.\; ((\mathit{mva} - \mathit{mvb}) = ((c + 1) \cdot d))\})\} \\ \quad \quad =\; \{y \in \operatorname{ccint}(0, 1) : ((y - b) \in \{d \in \mathbb{R} : \exists \mathit{mva}, \mathit{mvb}, c \in \mathbb{N}.\; ((\mathit{mva} - \mathit{mvb}) = ((c + 1) \cdot d))\})\} \end{array}}}
\vnbprop{vitali-choice}{\vnbstmt{\begin{array}{@{}l@{}} \forall x \in \operatorname{ccint}(0, 1). \\ \quad \varepsilon(\{y \in \operatorname{ccint}(0, 1) : ((y - x) \in \{d \in \mathbb{R} : \exists a, b, c \in \mathbb{N}.\; ((a - b) = ((c + 1) \cdot d))\})\}) \\ \quad \in\; \operatorname{ccint}(0, 1) \wedge \\ \quad (\varepsilon(\{y \in \operatorname{ccint}(0, 1) : ((y - x) \in \{d \in \mathbb{R} : \exists a, b, c \in \mathbb{N}.\; ((a - b) = ((c + 1) \cdot d))\})\}) - x) \\ \quad \in\; \{d \in \mathbb{R} : \exists a, b, c \in \mathbb{N}.\; ((a - b) = ((c + 1) \cdot d))\} \wedge \\ \quad \varepsilon(\{y \in \operatorname{ccint}(0, 1) : ((y - x) \in \{d \in \mathbb{R} : \exists a, b, c \in \mathbb{N}.\; ((a - b) = ((c + 1) \cdot d))\})\}) \\ \quad \in\; \{h \in \operatorname{ccint}(0, 1) : \exists \mathit{mvx} \in \operatorname{ccint}(0, 1).\; (h = \varepsilon(\{y \in \operatorname{ccint}(0, 1) : ((y - \mathit{mvx}) \in \{d \in \mathbb{R} : \exists a, b, c \in \mathbb{N}.\; ((a - b) = ((c + 1) \cdot d))\})\}))\} \end{array}}}
\vnbprop{esum-constant-positive-not-finite}{\vnbstmt{\begin{array}{@{}l@{}} \forall f \in (\mathbb{N} \to \operatorname{rr\text{-}pos\text{-}star}), c \in \mathbb{R}. \\ \quad ((0 < c) \wedge \forall k \in \mathbb{N}.\; (f(k) = c)) \Rightarrow \neg (\operatorname{esum}(f) \in \mathbb{R}) \end{array}}}
\vnbprop{translate-in-power-rr}{\vnbstmt{\forall s, c.\; ((s + c) \in \operatorname{power}(\mathbb{R}))}}
\vnbprop{vitali-null-case}{\vnbstmt{\begin{array}{@{}l@{}} \forall u. \\ \quad \operatorname{is\text{-}measure}(\mathbb{R}, \operatorname{power}(\mathbb{R}), u) \wedge \\ \quad \forall s \in \operatorname{power}(\mathbb{R}), t \in \mathbb{R}.\; (u((s + t)) = u(s)) \wedge \\ \quad 0 < u(\operatorname{ccint}(0, 1)) \Rightarrow \\ \quad \quad \neg (u(\{h \in \operatorname{ccint}(0, 1) : \exists x \in \operatorname{ccint}(0, 1).\; (h = \varepsilon(\{y \in \operatorname{ccint}(0, 1) : ((y - x) \in \{d \in \mathbb{R} : \exists a, b, c \in \mathbb{N}.\; ((a - b) = ((c + 1) \cdot d))\})\}))\}) = 0) \end{array}}}
\vnbprop{vitali-positive-case}{\vnbstmt{\begin{array}{@{}l@{}} \forall u. \\ \quad \operatorname{is\text{-}measure}(\mathbb{R}, \operatorname{power}(\mathbb{R}), u) \wedge \\ \quad \forall s \in \operatorname{power}(\mathbb{R}), t \in \mathbb{R}.\; (u((s + t)) = u(s)) \wedge \\ \quad u(\operatorname{ccint}(0, 1)) \in \mathbb{R} \Rightarrow \\ \quad \quad \neg (0 < u(\{h \in \operatorname{ccint}(0, 1) : \exists x \in \operatorname{ccint}(0, 1).\; (h = \varepsilon(\{y \in \operatorname{ccint}(0, 1) : ((y - x) \in \{d \in \mathbb{R} : \exists a, b, c \in \mathbb{N}.\; ((a - b) = ((c + 1) \cdot d))\})\}))\})) \end{array}}}
\vnbprop{vitali-choice-set-finite}{\vnbstmt{\begin{array}{@{}l@{}} \forall u. \\ \quad (\operatorname{is\text{-}measure}(\mathbb{R}, \operatorname{power}(\mathbb{R}), u) \wedge (u(\operatorname{ccint}(0, 1)) \in \mathbb{R})) \Rightarrow  \\ \quad \quad u(\{h \in \operatorname{ccint}(0, 1) : \exists x \in \operatorname{ccint}(0, 1).\; (h = \varepsilon(\{y \in \operatorname{ccint}(0, 1) : ((y - x) \in \{d \in \mathbb{R} : \exists a, b, c \in \mathbb{N}.\; ((a - b) = ((c + 1) \cdot d))\})\}))\}) \\ \quad \quad \in\; \mathbb{R} \wedge \\ \quad \quad 0 \\ \quad \quad \leq\; u(\{h \in \operatorname{ccint}(0, 1) : \exists x \in \operatorname{ccint}(0, 1).\; (h = \varepsilon(\{y \in \operatorname{ccint}(0, 1) : ((y - x) \in \{d \in \mathbb{R} : \exists a, b, c \in \mathbb{N}.\; ((a - b) = ((c + 1) \cdot d))\})\}))\}) \end{array}}}
\vnbprop{vitali-no-invariant-measure-via-d}{\vnbstmt{\neg \exists u.\; (\operatorname{is\text{-}measure}(\mathbb{R}, \operatorname{power}(\mathbb{R}), u) \wedge (\forall s \in \operatorname{power}(\mathbb{R}), t \in \mathbb{R}.\; (u((s + t)) = u(s)) \wedge ((0 < u(\operatorname{ccint}(0, 1))) \wedge (u(\operatorname{ccint}(0, 1)) \in \mathbb{R}))))}}

\vnbbookfile{theorem-library/measure-ch1-vitali-q}{theorem-library/measure-ch1-vitali-q}
\noindent{\small theorem-library/measure-ch1-vitali-q.scm -- VITALI (Prop. 1.23) WITH Q, AS THE NOTES PROVE IT (2026-10-26). The same proof as measure-ch1-vitali.scm, with the equivalence x == y iff x - y in QQ (the notes' (13)), QQ enumerated by qq-enumerated (theorem-library/qq-countable.scm, through the primitive generation axioms). ...}\par
\vnbprop{vitali-q-refl}{\vnbstmt{\forall x \in \mathbb{R}.\; ((x - x) \in \mathbb{Q})}}
\vnbprop{vitali-q-symm}{\vnbstmt{\forall x, y.\; (((x \in \mathbb{R}) \wedge ((y \in \mathbb{R}) \wedge ((y - x) \in \mathbb{Q}))) \Rightarrow ((x - y) \in \mathbb{Q}))}}
\vnbprop{vitali-q-trans}{\vnbstmt{\begin{array}{@{}l@{}} \forall x, y, z. \\ \quad ((x \in \mathbb{R}) \wedge ((y \in \mathbb{R}) \wedge ((z \in \mathbb{R}) \wedge (((x - y) \in \mathbb{Q}) \wedge ((y - z) \in \mathbb{Q}))))) \Rightarrow  \\ \quad \quad (x - z) \in \mathbb{Q} \end{array}}}
\vnbprop{vitali-q-recip-difference}{\vnbstmt{\forall i, j.\; (((i \in \mathbb{N}) \wedge (j \in \mathbb{N})) \Rightarrow ((\operatorname{recip}((j + 1)) - \operatorname{recip}((i + 1))) \in \mathbb{Q}))}}
\vnbprop{vitali-q-class-eq}{\vnbstmt{\begin{array}{@{}l@{}} \forall a, b. \\ \quad ((a \in \mathbb{R}) \wedge ((b \in \mathbb{R}) \wedge ((a - b) \in \mathbb{Q}))) \Rightarrow  \\ \quad \quad \{y \in \operatorname{ccint}(0, 1) : ((y - a) \in \mathbb{Q})\} = \{y \in \operatorname{ccint}(0, 1) : ((y - b) \in \mathbb{Q})\} \end{array}}}
\vnbprop{vitali-q-choice}{\vnbstmt{\begin{array}{@{}l@{}} \forall x \in \operatorname{ccint}(0, 1). \\ \quad \varepsilon(\{y \in \operatorname{ccint}(0, 1) : ((y - x) \in \mathbb{Q})\}) \in \operatorname{ccint}(0, 1) \wedge \\ \quad (\varepsilon(\{y \in \operatorname{ccint}(0, 1) : ((y - x) \in \mathbb{Q})\}) - x) \in \mathbb{Q} \wedge \\ \quad \varepsilon(\{y \in \operatorname{ccint}(0, 1) : ((y - x) \in \mathbb{Q})\}) \\ \quad \in\; \{h \in \operatorname{ccint}(0, 1) : \exists \mathit{mvx} \in \operatorname{ccint}(0, 1).\; (h = \varepsilon(\{y \in \operatorname{ccint}(0, 1) : ((y - \mathit{mvx}) \in \mathbb{Q})\}))\} \end{array}}}
\vnbprop{vitali-q-null-case}{\vnbstmt{\begin{array}{@{}l@{}} \forall u. \\ \quad \operatorname{is\text{-}measure}(\mathbb{R}, \operatorname{power}(\mathbb{R}), u) \wedge \\ \quad \forall s \in \operatorname{power}(\mathbb{R}), t \in \mathbb{R}.\; (u((s + t)) = u(s)) \wedge \\ \quad 0 < u(\operatorname{ccint}(0, 1)) \Rightarrow \\ \quad \quad \neg (u(\{h \in \operatorname{ccint}(0, 1) : \exists x \in \operatorname{ccint}(0, 1).\; (h = \varepsilon(\{y \in \operatorname{ccint}(0, 1) : ((y - x) \in \mathbb{Q})\}))\}) = 0) \end{array}}}
\vnbprop{vitali-q-positive-case}{\vnbstmt{\begin{array}{@{}l@{}} \forall u. \\ \quad \operatorname{is\text{-}measure}(\mathbb{R}, \operatorname{power}(\mathbb{R}), u) \wedge \\ \quad \forall s \in \operatorname{power}(\mathbb{R}), t \in \mathbb{R}.\; (u((s + t)) = u(s)) \wedge \\ \quad u(\operatorname{ccint}(0, 1)) \in \mathbb{R} \Rightarrow \\ \quad \quad \neg (0 < u(\{h \in \operatorname{ccint}(0, 1) : \exists x \in \operatorname{ccint}(0, 1).\; (h = \varepsilon(\{y \in \operatorname{ccint}(0, 1) : ((y - x) \in \mathbb{Q})\}))\})) \end{array}}}
\vnbprop{vitali-q-choice-set-finite}{\vnbstmt{\begin{array}{@{}l@{}} \forall u. \\ \quad (\operatorname{is\text{-}measure}(\mathbb{R}, \operatorname{power}(\mathbb{R}), u) \wedge (u(\operatorname{ccint}(0, 1)) \in \mathbb{R})) \Rightarrow  \\ \quad \quad u(\{h \in \operatorname{ccint}(0, 1) : \exists x \in \operatorname{ccint}(0, 1).\; (h = \varepsilon(\{y \in \operatorname{ccint}(0, 1) : ((y - x) \in \mathbb{Q})\}))\}) \\ \quad \quad \in\; \mathbb{R} \wedge \\ \quad \quad 0 \leq u(\{h \in \operatorname{ccint}(0, 1) : \exists x \in \operatorname{ccint}(0, 1).\; (h = \varepsilon(\{y \in \operatorname{ccint}(0, 1) : ((y - x) \in \mathbb{Q})\}))\}) \end{array}}}
\vnbprop{vitali-no-invariant-measure}{\vnbstmt{\neg \exists u.\; (\operatorname{is\text{-}measure}(\mathbb{R}, \operatorname{power}(\mathbb{R}), u) \wedge (\forall s \in \operatorname{power}(\mathbb{R}), t \in \mathbb{R}.\; (u((s + t)) = u(s)) \wedge ((0 < u(\operatorname{ccint}(0, 1))) \wedge (u(\operatorname{ccint}(0, 1)) \in \mathbb{R}))))}}

\vnbbookfile{theorem-library/riesz-markov}{theorem-library/riesz-markov}
\noindent{\small theorem-library/riesz-markov.scm -- RIESZ-MARKOV, STAGE 3B (RM-3B, 2026-10-10): the assembly on X, conditional on the positive extension.  Helper prefix: rmk-. Source: the user's notes-63, form (b).}\par
\vnbprop{rmk-riesz-markov-from-k}{\vnbstmt{\begin{array}{@{}l@{}} \forall s, q, m, x. \\ \quad \operatorname{is\text{-}compact}(s) \wedge \\ \quad \operatorname{is\text{-}dense\text{-}seq}(s, q) \wedge \\ \quad \operatorname{is\text{-}positive\text{-}functional}(\operatorname{rm\text{-}kms}(s, q), m) \wedge \\ \quad \operatorname{is\text{-}positive\text{-}functional}(s, x) \wedge \\ \quad \forall f \in \operatorname{cont\text{-}fns}(s). \\ \quad \quad m((\lambda z \in \operatorname{pts}(\operatorname{rm\text{-}kms}(s, q)).\; f(\operatorname{rm\text{-}pi}(s, q)(z)))) = x(f) \Rightarrow \\ \quad \quad \operatorname{is\text{-}measure}(\operatorname{pts}(s), \operatorname{borel\text{-}sets}(s), \\ \quad \quad \quad \operatorname{push\text{-}forward}(\operatorname{pts}(\operatorname{rm\text{-}kms}(s, q)), \operatorname{borel\text{-}sets}(s), \\ \quad \quad \quad \quad \operatorname{res}(\operatorname{method\text{-}i\text{-}outer}(\operatorname{pts}(\operatorname{rm\text{-}kms}(s, q)), \operatorname{clopen\text{-}sets}(\operatorname{rm\text{-}kms}(s, q)), \\ \quad \quad \quad \quad \quad \quad (\lambda u \in \operatorname{clopen\text{-}sets}(\operatorname{rm\text{-}kms}(s, q)).\; m(\operatorname{indicator}(\operatorname{pts}(\operatorname{rm\text{-}kms}(s, q)), u)))), \\ \quad \quad \quad \quad \quad \operatorname{caratheodory\text{-}sets}(\operatorname{pts}(\operatorname{rm\text{-}kms}(s, q)), \\ \quad \quad \quad \quad \quad \quad \operatorname{method\text{-}i\text{-}outer}(\operatorname{pts}(\operatorname{rm\text{-}kms}(s, q)), \operatorname{clopen\text{-}sets}(\operatorname{rm\text{-}kms}(s, q)), \\ \quad \quad \quad \quad \quad \quad \quad (\lambda u \in \operatorname{clopen\text{-}sets}(\operatorname{rm\text{-}kms}(s, q)).\; m(\operatorname{indicator}(\operatorname{pts}(\operatorname{rm\text{-}kms}(s, q)), u)))))), \\ \quad \quad \quad \quad \operatorname{rm\text{-}pi}(s, q))) \wedge \\ \quad \quad \forall f \in \operatorname{cont\text{-}fns}(s). \\ \quad \quad \quad \operatorname{is\text{-}integrable}(\operatorname{pts}(s), \operatorname{borel\text{-}sets}(s), \\ \quad \quad \quad \quad \operatorname{push\text{-}forward}(\operatorname{pts}(\operatorname{rm\text{-}kms}(s, q)), \operatorname{borel\text{-}sets}(s), \\ \quad \quad \quad \quad \quad \operatorname{res}(\operatorname{method\text{-}i\text{-}outer}(\operatorname{pts}(\operatorname{rm\text{-}kms}(s, q)), \operatorname{clopen\text{-}sets}(\operatorname{rm\text{-}kms}(s, q)), \\ \quad \quad \quad \quad \quad \quad \quad (\lambda u \in \operatorname{clopen\text{-}sets}(\operatorname{rm\text{-}kms}(s, q)).\; m(\operatorname{indicator}(\operatorname{pts}(\operatorname{rm\text{-}kms}(s, q)), u)))), \\ \quad \quad \quad \quad \quad \quad \operatorname{caratheodory\text{-}sets}(\operatorname{pts}(\operatorname{rm\text{-}kms}(s, q)), \\ \quad \quad \quad \quad \quad \quad \quad \operatorname{method\text{-}i\text{-}outer}(\operatorname{pts}(\operatorname{rm\text{-}kms}(s, q)), \operatorname{clopen\text{-}sets}(\operatorname{rm\text{-}kms}(s, q)), \\ \quad \quad \quad \quad \quad \quad \quad \quad (\lambda u \in \operatorname{clopen\text{-}sets}(\operatorname{rm\text{-}kms}(s, q)).\; m(\operatorname{indicator}(\operatorname{pts}(\operatorname{rm\text{-}kms}(s, q)), u)))))), \\ \quad \quad \quad \quad \quad \operatorname{rm\text{-}pi}(s, q)), \\ \quad \quad \quad \quad f) \wedge \\ \quad \quad \quad \operatorname{rintegral}(\operatorname{pts}(s), \operatorname{borel\text{-}sets}(s), \\ \quad \quad \quad \quad \operatorname{push\text{-}forward}(\operatorname{pts}(\operatorname{rm\text{-}kms}(s, q)), \operatorname{borel\text{-}sets}(s), \\ \quad \quad \quad \quad \quad \operatorname{res}(\operatorname{method\text{-}i\text{-}outer}(\operatorname{pts}(\operatorname{rm\text{-}kms}(s, q)), \operatorname{clopen\text{-}sets}(\operatorname{rm\text{-}kms}(s, q)), \\ \quad \quad \quad \quad \quad \quad \quad (\lambda u \in \operatorname{clopen\text{-}sets}(\operatorname{rm\text{-}kms}(s, q)).\; m(\operatorname{indicator}(\operatorname{pts}(\operatorname{rm\text{-}kms}(s, q)), u)))), \\ \quad \quad \quad \quad \quad \quad \operatorname{caratheodory\text{-}sets}(\operatorname{pts}(\operatorname{rm\text{-}kms}(s, q)), \\ \quad \quad \quad \quad \quad \quad \quad \operatorname{method\text{-}i\text{-}outer}(\operatorname{pts}(\operatorname{rm\text{-}kms}(s, q)), \operatorname{clopen\text{-}sets}(\operatorname{rm\text{-}kms}(s, q)), \\ \quad \quad \quad \quad \quad \quad \quad \quad (\lambda u \in \operatorname{clopen\text{-}sets}(\operatorname{rm\text{-}kms}(s, q)).\; m(\operatorname{indicator}(\operatorname{pts}(\operatorname{rm\text{-}kms}(s, q)), u)))))), \\ \quad \quad \quad \quad \quad \operatorname{rm\text{-}pi}(s, q)), \\ \quad \quad \quad \quad f) \\ \quad \quad \quad =\; x(f) \end{array}}}

\vnbbookfile{theorem-library/riesz-markov-compact}{theorem-library/riesz-markov-compact}
\noindent{\small theorem-library/riesz-markov-compact.scm -- RIESZ-MARKOV on a compact metric space (RM-3A follow-up, 2026-10-10).  Prefix rmz-.  Source: the user's notes-63, form (b).}\par
\vnbprop{rmz-comp-cont}{\vnbstmt{\begin{array}{@{}l@{}} \forall s, q.; \forall f \in \operatorname{cont\text{-}fns}(s). \\ \quad (\operatorname{is\text{-}compact}(s) \wedge \operatorname{is\text{-}dense\text{-}seq}(s, q)) \Rightarrow  \\ \quad \quad (\lambda z \in \operatorname{pts}(\operatorname{rm\text{-}kms}(s, q)).\; f(\operatorname{rm\text{-}pi}(s, q)(z))) \in \operatorname{cont\text{-}fns}(\operatorname{rm\text{-}kms}(s, q)) \end{array}}}
\vnbprop{rmz-one-cont}{\vnbstmt{\forall s.\; (\operatorname{is\text{-}metric\text{-}space}(s) \Rightarrow ((\lambda x \in \operatorname{pts}(s).\; 1) \in \operatorname{cont\text{-}fns}(s)))}}
\vnbprop{rmz-lm-val}{\vnbstmt{\begin{array}{@{}l@{}} \forall s, q, x.; \forall f \in \operatorname{cont\text{-}fns}(s). \\ \quad (\operatorname{is\text{-}compact}(s) \wedge (\operatorname{is\text{-}dense\text{-}seq}(s, q) \wedge \operatorname{is\text{-}positive\text{-}functional}(s, x))) \Rightarrow  \\ \quad \quad (\lambda z \in \operatorname{pts}(\operatorname{rm\text{-}kms}(s, q)).\; f(\operatorname{rm\text{-}pi}(s, q)(z))) \\ \quad \quad \in\; \{g \in \operatorname{cont\text{-}fns}(\operatorname{rm\text{-}kms}(s, q)) : \exists \mathit{rzf} \in \operatorname{cont\text{-}fns}(s).\; (g = (\lambda z \in \operatorname{pts}(\operatorname{rm\text{-}kms}(s, q)).\; \operatorname{rzf}(\operatorname{rm\text{-}pi}(s, q)(z))))\} \wedge \\ \quad \quad (\lambda g \in \{g \in \operatorname{cont\text{-}fns}(\operatorname{rm\text{-}kms}(s, q)) : \exists \mathit{rzf} \in \operatorname{cont\text{-}fns}(s).\; (g = (\lambda z \in \operatorname{pts}(\operatorname{rm\text{-}kms}(s, q)).\; \operatorname{rzf}(\operatorname{rm\text{-}pi}(s, q)(z))))\}.\; x((\lambda \mathit{rzx} \in \operatorname{pts}(s).\; g(\operatorname{rm\text{-}code}(s, q, \mathit{rzx})))))((\lambda z \in \operatorname{pts}(\operatorname{rm\text{-}kms}(s, q)).\; f(\operatorname{rm\text{-}pi}(s, q)(z)))) \\ \quad \quad =\; x(f) \end{array}}}
\vnbprop{rmz-riesz-markov-dense}{\vnbstmt{\begin{array}{@{}l@{}} \forall s, q, x. \\ \quad (\operatorname{is\text{-}compact}(s) \wedge (\operatorname{is\text{-}dense\text{-}seq}(s, q) \wedge \operatorname{is\text{-}positive\text{-}functional}(s, x))) \Rightarrow  \\ \quad \quad \exists u. \\ \quad \quad \quad \operatorname{is\text{-}measure}(\operatorname{pts}(s), \operatorname{borel\text{-}sets}(s), u) \wedge \\ \quad \quad \quad \forall f \in \operatorname{cont\text{-}fns}(s). \\ \quad \quad \quad \quad (\operatorname{is\text{-}integrable}(\operatorname{pts}(s), \operatorname{borel\text{-}sets}(s), u, f) \wedge (\operatorname{rintegral}(\operatorname{pts}(s), \operatorname{borel\text{-}sets}(s), u, f) = x(f))) \end{array}}}
\vnbprop{rmz-riesz-markov}{\vnbstmt{\begin{array}{@{}l@{}} \forall s, a. \\ \quad (\operatorname{is\text{-}compact}(s) \wedge (\exists x.\; (x \in \operatorname{pts}(s)) \wedge \operatorname{is\text{-}positive\text{-}functional}(s, a))) \Rightarrow  \\ \quad \quad \exists u. \\ \quad \quad \quad \operatorname{is\text{-}measure}(\operatorname{pts}(s), \operatorname{borel\text{-}sets}(s), u) \wedge \\ \quad \quad \quad \forall f \in \operatorname{cont\text{-}fns}(s). \\ \quad \quad \quad \quad (\operatorname{is\text{-}integrable}(\operatorname{pts}(s), \operatorname{borel\text{-}sets}(s), u, f) \wedge (\operatorname{rintegral}(\operatorname{pts}(s), \operatorname{borel\text{-}sets}(s), u, f) = a(f))) \end{array}}}

\vnbbookfile{theorem-library/mb4-polish-open}{theorem-library/mb4-polish-open}
\noindent{\small mb4-polish-open.scm -- construction of measures, Lemma 5.18 (M-B4, 2026-10-26): an open set U of a Polish space X is itself Polish -- there is a complete separable metric space homeomorphic to U, namely the points of U with d'(x, y) = d(x, y) + |1/d(x, X\textbackslash{}U) - 1/d(y, X\textbackslash{}U)|, the metric U inherits as the graph Y = \{(x, t) ...}\par
\vnbprop{graph-ms-pts}{\vnbstmt{\forall s.; \forall a \in \mathbf{Set}, g.\; (\operatorname{pts}(\operatorname{graph\text{-}ms}(s, a, g)) \simeq a)}}
\vnbprop{graph-ms-dist-slot}{\vnbstmt{\begin{array}{@{}l@{}} \forall s.; \forall a \in \mathbf{Set}, g. \\ \quad (\operatorname{dist}(\operatorname{graph\text{-}ms}(s, a, g)) \simeq (\lambda {[}\operatorname{gmu\_}, \operatorname{gmv\_}] \in (a \times a).\; (\operatorname{dist}(s)(\operatorname{gmu\_}, \operatorname{gmv\_}) + \lvert (g(\operatorname{gmu\_}) - g(\operatorname{gmv\_})) \rvert))) \end{array}}}
\vnbprop{graph-ms-dist}{\vnbstmt{\begin{array}{@{}l@{}} \forall s, a, g, u, v. \\ \quad ((a \in \mathbf{Set}) \wedge ((u \in a) \wedge (v \in a))) \Rightarrow  \\ \quad \quad (\operatorname{dist}(\operatorname{graph\text{-}ms}(s, a, g))(u, v) \simeq (\operatorname{dist}(s)(u, v) + \lvert (g(u) - g(v)) \rvert)) \end{array}}}
\vnbprop{graph-ms-is-metric-space}{\vnbstmt{\begin{array}{@{}l@{}} \forall s, a.; \forall g \in (a \to \mathbb{R}). \\ \quad (\operatorname{is\text{-}metric\text{-}space}(s) \wedge (a \subseteq \operatorname{pts}(s))) \Rightarrow  \\ \quad \quad \operatorname{is\text{-}metric\text{-}space}(\operatorname{graph\text{-}ms}(s, a, g)) \end{array}}}
\vnbprop{recip-near}{\vnbstmt{\begin{array}{@{}l@{}} \forall c \in \mathbb{R}, s. \\ \quad ((0 < c) \wedge \operatorname{pos\text{-}rr}(s)) \Rightarrow  \\ \quad \quad \exists a. \\ \quad \quad \quad (\operatorname{pos\text{-}rr}(a) \wedge \forall v \in \mathbb{R}.\; ((\lvert (c - v) \rvert \leq a) \Rightarrow ((0 < v) \wedge (\lvert (\operatorname{recip}(c) - \operatorname{recip}(v)) \rvert \leq s)))) \end{array}}}
\vnbprop{pt-set-dist-le}{\vnbstmt{\begin{array}{@{}l@{}} \forall m, s.; \forall x \in \operatorname{pts}(m), z \in s. \\ \quad (\operatorname{is\text{-}metric\text{-}space}(m) \wedge (s \subseteq \operatorname{pts}(m))) \Rightarrow  \\ \quad \quad \operatorname{pt\text{-}set\text{-}dist}(m, x, s) \leq \operatorname{dist}(m)(x, z) \end{array}}}
\vnbprop{pt-set-dist-real}{\vnbstmt{\begin{array}{@{}l@{}} \forall m, s.; \forall z \in s, x \in \operatorname{pts}(m). \\ \quad (\operatorname{is\text{-}metric\text{-}space}(m) \wedge (s \subseteq \operatorname{pts}(m))) \Rightarrow  \\ \quad \quad (\operatorname{pt\text{-}set\text{-}dist}(m, x, s) \in \mathbb{R}) \wedge (0 \leq \operatorname{pt\text{-}set\text{-}dist}(m, x, s)) \end{array}}}
\vnbprop{pt-set-dist-lipschitz}{\vnbstmt{\begin{array}{@{}l@{}} \forall m, s.; \forall z \in s, x, y \in \operatorname{pts}(m). \\ \quad (\operatorname{is\text{-}metric\text{-}space}(m) \wedge (s \subseteq \operatorname{pts}(m))) \Rightarrow  \\ \quad \quad \operatorname{pt\text{-}set\text{-}dist}(m, x, s) \leq (\operatorname{dist}(m)(x, y) + \operatorname{pt\text{-}set\text{-}dist}(m, y, s)) \end{array}}}
\vnbprop{pt-set-dist-open-pos}{\vnbstmt{\begin{array}{@{}l@{}} \forall m, u.; \forall x \in u. \\ \quad \operatorname{is\text{-}open}(m, u) \Rightarrow  \\ \quad \quad \exists r.\; (\operatorname{pos\text{-}rr}(r) \wedge (r \leq \operatorname{pt\text{-}set\text{-}dist}(m, x, \operatorname{complement\text{-}in}(\operatorname{pts}(m), u)))) \end{array}}}
\vnbprop{graph-ms-near}{\vnbstmt{\begin{array}{@{}l@{}} \forall s, a.; \forall g \in (a \to \mathbb{R}), x \in a, b. \\ \quad \operatorname{is\text{-}metric\text{-}space}(s) \wedge \\ \quad a \subseteq \operatorname{pts}(s) \wedge \\ \quad \forall \mathit{gx} \in a, e. \\ \quad \quad \operatorname{pos\text{-}rr}(e) \Rightarrow  \\ \quad \quad \quad \exists d. \\ \quad \quad \quad \quad (\operatorname{pos\text{-}rr}(d) \wedge \forall \mathit{gy} \in a.\; ((\operatorname{dist}(s)(\mathit{gx}, \mathit{gy}) \leq d) \Rightarrow (\lvert (g(\mathit{gx}) - g(\mathit{gy})) \rvert \leq e))) \wedge \\ \quad \operatorname{pos\text{-}rr}(b) \Rightarrow \\ \quad \quad \exists c. \\ \quad \quad \quad (\operatorname{pos\text{-}rr}(c) \wedge \forall y \in a.\; ((\operatorname{dist}(s)(x, y) \leq c) \Rightarrow (\operatorname{dist}(\operatorname{graph\text{-}ms}(s, a, g))(x, y) < b))) \end{array}}}
\vnbprop{coarser-metric-separable}{\vnbstmt{\begin{array}{@{}l@{}} \forall s, a, t. \\ \quad \operatorname{is\text{-}separable}(s) \wedge \\ \quad a \subseteq \operatorname{pts}(s) \wedge \\ \quad \exists z.\; (z \in a) \wedge \\ \quad \operatorname{is\text{-}metric\text{-}space}(t) \wedge \\ \quad (\operatorname{pts}(t) \simeq a) \wedge \\ \quad \forall x \in a, e. \\ \quad \quad \operatorname{pos\text{-}rr}(e) \Rightarrow  \\ \quad \quad \quad \exists d. \\ \quad \quad \quad \quad \operatorname{pos\text{-}rr}(d) \wedge \forall y \in a.\; ((\operatorname{dist}(s)(x, y) \leq d) \Rightarrow (\operatorname{dist}(t)(x, y) < e)) \Rightarrow \\ \quad \quad \operatorname{is\text{-}separable}(t) \end{array}}}
\vnbprop{graph-ms-homeomorphism}{\vnbstmt{\begin{array}{@{}l@{}} \forall s, a.; \forall g \in (a \to \mathbb{R}). \\ \quad \operatorname{is\text{-}metric\text{-}space}(s) \wedge \\ \quad a \subseteq \operatorname{pts}(s) \wedge \\ \quad \forall x \in a, e. \\ \quad \quad \operatorname{pos\text{-}rr}(e) \Rightarrow  \\ \quad \quad \quad \exists d. \\ \quad \quad \quad \quad (\operatorname{pos\text{-}rr}(d) \wedge \forall y \in a.\; ((\operatorname{dist}(s)(x, y) \leq d) \Rightarrow (\lvert (g(x) - g(y)) \rvert \leq e))) \Rightarrow \\ \quad \quad \operatorname{is\text{-}homeomorphism}(\operatorname{graph\text{-}ms}(s, a, g), \operatorname{subspace\text{-}ms}(s, a), (\lambda b \in a.\; b)) \end{array}}}
\vnbprop{graph-ms-complete}{\vnbstmt{\begin{array}{@{}l@{}} \forall s, a.; \forall g \in (a \to \mathbb{R}). \\ \quad \operatorname{is\text{-}complete}(s) \wedge \\ \quad a \subseteq \operatorname{pts}(s) \wedge \\ \quad \forall x \in a, e. \\ \quad \quad \operatorname{pos\text{-}rr}(e) \Rightarrow  \\ \quad \quad \quad \exists d. \\ \quad \quad \quad \quad (\operatorname{pos\text{-}rr}(d) \wedge \forall y \in a.\; ((\operatorname{dist}(s)(x, y) \leq d) \Rightarrow (\lvert (g(x) - g(y)) \rvert \leq e))) \wedge \\ \quad \forall h \in (\mathbb{N} \to a), l \in \operatorname{pts}(s), c \in \mathbb{R}, n \in \mathbb{N}. \\ \quad \quad (\operatorname{converges\text{-}to}(s, h, l) \wedge \forall k \in \mathbb{N}.\; ((n \leq k) \Rightarrow (g(h(k)) \leq c))) \Rightarrow  \\ \quad \quad \quad l \in a \Rightarrow \\ \quad \quad \operatorname{is\text{-}complete}(\operatorname{graph\text{-}ms}(s, a, g)) \end{array}}}
\vnbprop{graph-ms-polish}{\vnbstmt{\begin{array}{@{}l@{}} \forall s, a.; \forall g \in (a \to \mathbb{R}). \\ \quad \operatorname{is\text{-}polish}(s) \wedge \\ \quad a \subseteq \operatorname{pts}(s) \wedge \\ \quad \forall x \in a, e. \\ \quad \quad \operatorname{pos\text{-}rr}(e) \Rightarrow  \\ \quad \quad \quad \exists d. \\ \quad \quad \quad \quad (\operatorname{pos\text{-}rr}(d) \wedge \forall y \in a.\; ((\operatorname{dist}(s)(x, y) \leq d) \Rightarrow (\lvert (g(x) - g(y)) \rvert \leq e))) \wedge \\ \quad \forall h \in (\mathbb{N} \to a), l \in \operatorname{pts}(s), c \in \mathbb{R}, n \in \mathbb{N}. \\ \quad \quad (\operatorname{converges\text{-}to}(s, h, l) \wedge \forall k \in \mathbb{N}.\; ((n \leq k) \Rightarrow (g(h(k)) \leq c))) \Rightarrow  \\ \quad \quad \quad l \in a \wedge \\ \quad \exists z.\; (z \in a) \Rightarrow \\ \quad \quad \operatorname{is\text{-}polish}(\operatorname{graph\text{-}ms}(s, a, g)) \wedge \\ \quad \quad \operatorname{is\text{-}homeomorphism}(\operatorname{graph\text{-}ms}(s, a, g), \operatorname{subspace\text{-}ms}(s, a), (\lambda b \in a.\; b)) \end{array}}}
\vnbprop{polish-open-subspace}{\vnbstmt{\begin{array}{@{}l@{}} \forall s, u. \\ \quad (\operatorname{is\text{-}polish}(s) \wedge (\operatorname{is\text{-}open}(s, u) \wedge \exists z.\; (z \in u))) \Rightarrow  \\ \quad \quad \exists t. \\ \quad \quad \quad \operatorname{is\text{-}polish}(t) \wedge \exists h.\; \operatorname{is\text{-}homeomorphism}(t, \operatorname{subspace\text{-}ms}(s, u), h) \end{array}}}

\vnbbookfile{structure-library/mb35-vocab}{structure-library/mb35-vocab}
\noindent{\small mb35-vocab.scm -- vocabulary for construction of measures, chapter 5 section 2 (M-B35, 2026-10-26). Destination: structure-library/.  Library: measure.}\par
\par\medskip\noindent\textbf{Definition} (\texttt{is-regular-measure}).\;For all $s$ and $\mu$, 
\begin{equation*}
\begin{array}{@{}l@{}} \operatorname{is\text{-}regular\text{-}measure}(s, \mu) \\ \Leftrightarrow\; \operatorname{is\text{-}measure}(\operatorname{pts}(s), \operatorname{borel\text{-}sets}(s), \mu) \wedge \\ \mu(\operatorname{pts}(s)) \in \mathbb{R} \wedge \\ \forall \operatorname{rge\_} \in \operatorname{borel\text{-}sets}(s). \\ \quad \mu(\operatorname{rge\_}) \\ \quad =\; \operatorname{esup}(\{\operatorname{rgv\_} \in \operatorname{rr\text{-}pos\text{-}star} : \exists \operatorname{rgk\_} \in \operatorname{power}(\operatorname{pts}(s)).\; ((\operatorname{rgk\_} \subseteq \operatorname{rge\_}) \wedge (\operatorname{is\text{-}compact}(\operatorname{subspace\text{-}ms}(s, \operatorname{rgk\_})) \wedge (\operatorname{rgv\_} = \mu(\operatorname{rgk\_}))))\}) \end{array}
\end{equation*}


\vnbbookfile{structure-library/mb35-souslin-vocab}{structure-library/mb35-souslin-vocab}
\noindent{\small mb35-souslin-vocab.scm -- vocabulary for construction of measures, chapter 5 section 1.3-1.4: sieves and Souslin sets (M-B35 / B5, 2026-10-26).  Destination: structure-library/.  Library: measure.}\par
\par\medskip\noindent\textbf{Definition} (\texttt{nstar}).\; $(\mathit{nstar} \simeq \operatorname{big\text{-}union}(\operatorname{nsk\_}, \mathbb{N}, (\mathrm{OS}(\operatorname{nsk\_}) \to \mathbb{N})))$
\par\medskip\noindent\textbf{Definition} (\texttt{is-sieve}).\;For all $s$ and $\mathit{sv}$, 
\begin{equation*}
\begin{array}{@{}l@{}} \operatorname{is\text{-}sieve}(s, \mathit{sv}) \\ \Leftrightarrow\; \mathit{sv} \in (\mathit{nstar} \to \operatorname{power}(\operatorname{pts}(s))) \wedge \\ \forall \operatorname{svz\_} \in (\mathbb{N} \to \mathbb{N}). \\ \quad \forall \operatorname{sve\_}. \\ \quad \quad \operatorname{pos\text{-}rr}(\operatorname{sve\_}) \Rightarrow  \\ \quad \quad \quad \exists \operatorname{svb\_} \in \mathbb{N}. \forall \operatorname{svn\_} \in \mathbb{N}. \\ \quad \quad \quad \quad (\operatorname{svb\_} \leq \operatorname{svn\_}) \Rightarrow (\operatorname{diam}(s, \operatorname{sv}(\operatorname{restrict}(\operatorname{svz\_}, \mathrm{OS}(\operatorname{svn\_})))) \leq \operatorname{sve\_}) \wedge \\ \quad \forall \operatorname{svn\_} \in \mathbb{N}. \\ \quad \quad \operatorname{sv}(\operatorname{restrict}(\operatorname{svz\_}, \mathrm{OS}(\operatorname{succ}(\operatorname{svn\_})))) \subseteq \operatorname{sv}(\operatorname{restrict}(\operatorname{svz\_}, \mathrm{OS}(\operatorname{svn\_}))) \wedge \\ \quad \exists \operatorname{svp\_} \in \mathbf{Set}. \forall \operatorname{svn\_} \in \mathbb{N}. \\ \quad \quad \operatorname{svp\_} \in \operatorname{sv}(\operatorname{restrict}(\operatorname{svz\_}, \mathrm{OS}(\operatorname{svn\_}))) \end{array}
\end{equation*}

\par\medskip\noindent\textbf{Definition} (\texttt{souslin-set}).\;$\operatorname{souslin\text{-}set}(s, \mathit{sv}) :=$
\begin{equation*}
\{\operatorname{sxp\_} \in \operatorname{pts}(s) : \exists \operatorname{sxz\_} \in (\mathbb{N} \to \mathbb{N}).\; \forall \operatorname{sxn\_} \in \mathbb{N}.\; (\operatorname{sxp\_} \in \operatorname{sv}(\operatorname{restrict}(\operatorname{sxz\_}, \mathrm{OS}(\operatorname{sxn\_}))))\}
\end{equation*}

\par\medskip\noindent\textbf{Definition} (\texttt{is-souslin}).\;For all $s$ and $e$, 
\begin{equation*}
\operatorname{is\text{-}souslin}(s, e) \Leftrightarrow \exists \operatorname{ssv\_}.\; (\operatorname{is\text{-}sieve}(s, \operatorname{ssv\_}) \wedge (\operatorname{souslin\text{-}set}(s, \operatorname{ssv\_}) = e))
\end{equation*}

\par\medskip\noindent\textbf{Definition} (\texttt{is-baire-continuous}).\;For all $s$ and $f$, 
\begin{equation*}
\begin{array}{@{}l@{}} \operatorname{is\text{-}baire\text{-}continuous}(s, f) \\ \Leftrightarrow\; f \in ((\mathbb{N} \to \mathbb{N}) \to \operatorname{pts}(s)) \wedge \\ \forall \operatorname{bcz\_} \in (\mathbb{N} \to \mathbb{N}), \operatorname{bce\_}. \\ \quad \operatorname{pos\text{-}rr}(\operatorname{bce\_}) \Rightarrow  \\ \quad \quad \exists \operatorname{bcn\_} \in \mathbb{N}. \forall \operatorname{bcw\_} \in (\mathbb{N} \to \mathbb{N}). \\ \quad \quad \quad (\operatorname{restrict}(\operatorname{bcw\_}, \mathrm{OS}(\operatorname{bcn\_})) = \operatorname{restrict}(\operatorname{bcz\_}, \mathrm{OS}(\operatorname{bcn\_}))) \Rightarrow  \\ \quad \quad \quad \quad \operatorname{dist}(s)(f(\operatorname{bcz\_}), f(\operatorname{bcw\_})) \leq \operatorname{bce\_} \end{array}
\end{equation*}


\vnbbookfile{structure-library/mb67-vocab}{structure-library/mb67-vocab}
\noindent{\small mb67-vocab.scm -- vocabulary for construction of measures, chapter 5 sections 1.4-1.6: the Baire space N\textasciicircum{}N as a metric space, Lusin sets, standard Borel spaces (M-B67, 2026-10-26).  Destination: structure-library/. Library: measure.}\par
\par\medskip\noindent\textbf{Definition} (\texttt{baire-ms}).\; $(\operatorname{baire\text{-}ms} \simeq \operatorname{ultra\text{-}seq\text{-}ms}(\mathbb{N}, (\lambda \operatorname{bwk\_} \in \mathbb{N}.\; 2)))$
\par\medskip\noindent\textbf{Definition} (\texttt{snd-iter}).\;For every $\operatorname{sim\_}$,  $(\operatorname{snd\text{-}iter}(\operatorname{sim\_}, 0) \simeq \operatorname{sim\_})$\par
Let $\operatorname{sim\_}$ be arbitrary. Let $\operatorname{sin\_} \in \mathbb{N}$. Then $(\operatorname{snd\text{-}iter}(\operatorname{sim\_}, \operatorname{succ}(\operatorname{sin\_})) \simeq \operatorname{nnsnd}(\operatorname{snd\text{-}iter}(\operatorname{sim\_}, \operatorname{sin\_})))$
\par\medskip\noindent\textbf{Definition} (\texttt{baire-code}).\;$\operatorname{baire\text{-}code}(\operatorname{bcm\_}) :=$ $(\lambda \operatorname{bck\_} \in \mathbb{N}.\; \operatorname{nnfst}(\operatorname{snd\text{-}iter}(\operatorname{bcm\_}, \operatorname{bck\_})))$
\par\medskip\noindent\textbf{Definition} (\texttt{seq-tail}).\;$\operatorname{seq\text{-}tail}(\operatorname{stz\_}) :=$ $(\lambda \operatorname{stj\_} \in \mathbb{N}.\; \operatorname{stz\_}((\operatorname{stj\_} + 1)))$
\par\medskip\noindent\textbf{Definition} (\texttt{seq-cons}).\;$\operatorname{seq\text{-}cons}(\operatorname{sck\_}, \operatorname{scv\_}) :=$
\begin{equation*}
(\lambda \operatorname{scj\_} \in \mathbb{N}.\; \left(\text{if }(\operatorname{scj\_} = 0)\text{ then }\operatorname{sck\_}\text{ else }\operatorname{scv\_}(\operatorname{nn\text{-}minus}(\operatorname{scj\_}, 1))\right))
\end{equation*}

\par\medskip\noindent\textbf{Definition} (\texttt{seq-proj}).\;$\operatorname{seq\text{-}proj}(\operatorname{spk\_}, \operatorname{spz\_}) :=$ $(\lambda \operatorname{spj\_} \in \mathbb{N}.\; \operatorname{spz\_}(\operatorname{nnpair}(\operatorname{spk\_}, \operatorname{spj\_})))$
\par\medskip\noindent\textbf{Definition} (\texttt{sum-pts}).\;$\operatorname{sum\text{-}pts}(\operatorname{smt\_}) :=$
\begin{equation*}
\{\operatorname{smp\_} \in (\mathbb{N} \times \operatorname{big\text{-}union}(\operatorname{smk\_}, \mathbb{N}, \operatorname{pts}(\operatorname{smt\_}(\operatorname{smk\_})))) : (\operatorname{nth}(2, \operatorname{smp\_}) \in \operatorname{pts}(\operatorname{smt\_}(\operatorname{nth}(1, \operatorname{smp\_}))))\}
\end{equation*}

\par\medskip\noindent\textbf{Definition} (\texttt{sum-ms}).\;$\operatorname{sum\text{-}ms}(\operatorname{smt\_}) :=$
\begin{equation*}
{[}\operatorname{sum\text{-}pts}(\operatorname{smt\_}), (\lambda {[}\operatorname{smu\_}, \operatorname{smv\_}] \in (\operatorname{sum\text{-}pts}(\operatorname{smt\_}) \times \operatorname{sum\text{-}pts}(\operatorname{smt\_})).\; \left(\text{if }(\operatorname{nth}(1, \operatorname{smu\_}) = \operatorname{nth}(1, \operatorname{smv\_}))\text{ then }\operatorname{dist}(\operatorname{bdd\text{-}metric}(\operatorname{smt\_}(\operatorname{nth}(1, \operatorname{smu\_}))))(\operatorname{nth}(2, \operatorname{smu\_}), \operatorname{nth}(2, \operatorname{smv\_}))\text{ else }1\right))]
\end{equation*}

\par\medskip\noindent\textbf{Definition} (\texttt{is-countably-separated}).\;For all $\mathit{ys}$ and $\mathit{cb}$, 
\begin{equation*}
\begin{array}{@{}l@{}} \operatorname{is\text{-}countably\text{-}separated}(\mathit{ys}, \mathit{cb}) \\ \Leftrightarrow\; \operatorname{is\text{-}sigma\text{-}algebra}(\mathit{ys}, \mathit{cb}) \wedge \\ \exists \operatorname{csf\_} \in (\mathbb{N} \to \mathit{cb}). \forall \operatorname{csx\_}, \operatorname{csy\_} \in \mathit{ys}. \\ \quad \neg (\operatorname{csx\_} = \operatorname{csy\_}) \Rightarrow  \\ \quad \quad \exists \operatorname{csk\_} \in \mathbb{N}. \\ \quad \quad \quad (((\operatorname{csx\_} \in \operatorname{csf\_}(\operatorname{csk\_})) \wedge \neg (\operatorname{csy\_} \in \operatorname{csf\_}(\operatorname{csk\_}))) \vee ((\operatorname{csy\_} \in \operatorname{csf\_}(\operatorname{csk\_})) \wedge \neg (\operatorname{csx\_} \in \operatorname{csf\_}(\operatorname{csk\_})))) \end{array}
\end{equation*}

\par\medskip\noindent\textbf{Definition} (\texttt{graph-of}).\;$\operatorname{graph\text{-}of}(\operatorname{gfx\_}, \operatorname{gfy\_}, \operatorname{gff\_}) :=$
\begin{equation*}
\{\operatorname{grp\_} \in (\operatorname{gfx\_} \times \operatorname{gfy\_}) : (\operatorname{gff\_}(\operatorname{nth}(1, \operatorname{grp\_})) = \operatorname{nth}(2, \operatorname{grp\_}))\}
\end{equation*}

\par\medskip\noindent\textbf{Definition} (\texttt{is-standard-borel}).\;For all $\mathit{xs}$ and $\mathit{ca}$, 
\begin{equation*}
\begin{array}{@{}l@{}} \operatorname{is\text{-}standard\text{-}borel}(\mathit{xs}, \mathit{ca}) \\ \Leftrightarrow\; \exists \operatorname{sbs\_}. \\ \quad \operatorname{is\text{-}polish}(\operatorname{sbs\_}) \wedge \\ \quad \exists \operatorname{sbf\_} \in \mathrm{Bij}(\mathit{xs}, \operatorname{pts}(\operatorname{sbs\_})). \\ \quad \quad \operatorname{is\text{-}measurable\text{-}map}(\mathit{xs}, \mathit{ca}, \operatorname{pts}(\operatorname{sbs\_}), \operatorname{borel\text{-}sets}(\operatorname{sbs\_}), \operatorname{sbf\_}) \wedge \\ \quad \quad \forall \operatorname{sba\_} \in \mathit{ca}.\; (\operatorname{image}(\operatorname{sbf\_}, \operatorname{sba\_}) \in \operatorname{borel\text{-}sets}(\operatorname{sbs\_})) \end{array}
\end{equation*}

\par\medskip\noindent\textbf{Definition} (\texttt{borel-separated}).\;For all $s$, $a$ and $b$, 
\begin{equation*}
\begin{array}{@{}l@{}} \operatorname{borel\text{-}separated}(s, a, b) \\ \Leftrightarrow\; \exists \operatorname{bsc\_} \in \operatorname{borel\text{-}sets}(s). \\ \quad (a \subseteq \operatorname{bsc\_}) \wedge \forall \operatorname{bsx\_} \in b.\; \neg (\operatorname{bsx\_} \in \operatorname{bsc\_}) \end{array}
\end{equation*}

\par\medskip\noindent\textbf{Definition} (\texttt{cyl}).\;$\operatorname{cyl}(\operatorname{cyw\_}, \operatorname{cyn\_}) :=$
\begin{equation*}
\{\operatorname{cyz\_} \in (\mathbb{N} \to \mathbb{N}) : \forall \operatorname{cyj\_} \in \mathbb{N}.\; ((\operatorname{cyj\_} < \operatorname{cyn\_}) \Rightarrow (\operatorname{cyz\_}(\operatorname{cyj\_}) = \operatorname{cyw\_}(\operatorname{cyj\_})))\}
\end{equation*}

\par\medskip\noindent\textbf{Definition} (\texttt{upd}).\;$\operatorname{upd}(\operatorname{upw\_}, \operatorname{upn\_}, \operatorname{upm\_}) :=$ $(\lambda \operatorname{upj\_} \in \mathbb{N}.\; \left(\text{if }(\operatorname{upj\_} < \operatorname{upn\_})\text{ then }\operatorname{upw\_}(\operatorname{upj\_})\text{ else }\operatorname{upm\_}\right))$
\par\medskip\noindent\textbf{Definition} (\texttt{sum-sigma}).\;$\operatorname{sum\text{-}sigma}(\operatorname{sgx\_}, \operatorname{sgc\_}) :=$
\begin{equation*}
\{\operatorname{sga\_} \in \operatorname{power}(\operatorname{disjoint\text{-}sum\text{-}set}(\operatorname{sgx\_})) : \forall \operatorname{sgk\_} \in \mathbb{N}.\; (\{\operatorname{sgy\_} \in \operatorname{sgx\_}(\operatorname{sgk\_}) : ({[}\operatorname{sgk\_}, \operatorname{sgy\_}] \in \operatorname{sga\_})\} \in \operatorname{sgc\_}(\operatorname{sgk\_}))\}
\end{equation*}

\par\medskip\noindent\textbf{Definition} (\texttt{restrict-sigma}).\;$\operatorname{restrict\text{-}sigma}(\operatorname{rsc\_}, \operatorname{rse\_}) :=$ $\{\operatorname{rsa\_} \in \operatorname{rsc\_} : (\operatorname{rsa\_} \subseteq \operatorname{rse\_})\}$
\par\medskip\noindent\textbf{Definition} (\texttt{is-lusin}).\;For all $s$ and $e$, 
\begin{equation*}
\begin{array}{@{}l@{}} \operatorname{is\text{-}lusin}(s, e) \\ \Leftrightarrow\; \exists \operatorname{lzt\_}. \\ \quad \operatorname{is\text{-}complete}(\operatorname{lzt\_}) \wedge \\ \quad \operatorname{is\text{-}separable}(\operatorname{lzt\_}) \wedge \\ \quad \exists \operatorname{lzh\_}. \\ \quad \quad \operatorname{is\text{-}continuous}(\operatorname{lzt\_}, \operatorname{subspace\text{-}ms}(s, e), \operatorname{lzh\_}) \wedge \\ \quad \quad \operatorname{image}(\operatorname{lzh\_}, \operatorname{pts}(\operatorname{lzt\_})) = e \wedge \\ \quad \quad \forall \operatorname{lzx\_}, \operatorname{lzy\_} \in \operatorname{pts}(\operatorname{lzt\_}). \\ \quad \quad \quad (\operatorname{lzh\_}(\operatorname{lzx\_}) = \operatorname{lzh\_}(\operatorname{lzy\_})) \Rightarrow (\operatorname{lzx\_} = \operatorname{lzy\_}) \end{array}
\end{equation*}

\par\medskip\noindent\textbf{Definition} (\texttt{disjoint-sum-set}).\;$\operatorname{disjoint\text{-}sum\text{-}set}(\operatorname{ssx\_}) :=$
\begin{equation*}
\{\operatorname{ssp\_} \in (\mathbb{N} \times \operatorname{big\text{-}union}(\operatorname{ssk\_}, \mathbb{N}, \operatorname{ssx\_}(\operatorname{ssk\_}))) : (\operatorname{nth}(2, \operatorname{ssp\_}) \in \operatorname{ssx\_}(\operatorname{nth}(1, \operatorname{ssp\_})))\}
\end{equation*}


\vnbbookfile{structure-library/local-measure}{structure-library/local-measure}
\noindent{\small local-measure.scm -- WAVE 2 (2026-10-27..30), wired 2026-10-30 from scratchpad/mf/mf-mgvocab.scm. Library: measure.}\par
\par\medskip\noindent\textbf{Definition} (\texttt{is-hereditary}).\;For all $\mathit{ca}$ and $\mathit{cb}$, 
\begin{equation*}
\begin{array}{@{}l@{}} \operatorname{is\text{-}hereditary}(\mathit{ca}, \mathit{cb}) \\ \Leftrightarrow\; ((\mathit{cb} \subseteq \mathit{ca}) \wedge \forall \operatorname{hde\_} \in \mathit{cb}, \operatorname{hdf\_} \in \mathit{ca}.\; ((\operatorname{hdf\_} \subseteq \operatorname{hde\_}) \Rightarrow (\operatorname{hdf\_} \in \mathit{cb}))) \end{array}
\end{equation*}

\par\medskip\noindent\textbf{Definition} (\texttt{is-local-measure}).\;For all $\omega$, $\mathit{ca}$, $\mathit{cb}$ and $\rho$, 
\begin{equation*}
\begin{array}{@{}l@{}} \operatorname{is\text{-}local\text{-}measure}(\omega, \mathit{ca}, \mathit{cb}, \rho) \\ \Leftrightarrow\; \operatorname{is\text{-}sigma\text{-}algebra}(\omega, \mathit{ca}) \wedge \\ \operatorname{is\text{-}hereditary}(\mathit{ca}, \mathit{cb}) \wedge \\ \rho \in (\mathit{cb} \to \operatorname{rr\text{-}pos\text{-}star}) \wedge \\ \forall \operatorname{lme\_} \in \mathit{cb}. \\ \quad \operatorname{is\text{-}measure}(\operatorname{lme\_}, \operatorname{restrict\text{-}sigma}(\mathit{ca}, \operatorname{lme\_}), \operatorname{restrict}(\rho, \operatorname{restrict\text{-}sigma}(\mathit{ca}, \operatorname{lme\_}))) \end{array}
\end{equation*}

\par\medskip\noindent\textbf{Definition} (\texttt{is-locally-null}).\;For all $\mathit{ca}$, $\mathit{cb}$, $\rho$ and $\mathit{es}$, 
\begin{equation*}
\begin{array}{@{}l@{}} \operatorname{is\text{-}locally\text{-}null}(\mathit{ca}, \mathit{cb}, \rho, \mathit{es}) \\ \Leftrightarrow\; (\mathit{es} \in \mathit{ca}) \wedge \forall \operatorname{lna\_} \in \mathit{cb}.\; (\rho((\operatorname{lna\_} \cap \mathit{es})) = 0) \end{array}
\end{equation*}

\par\medskip\noindent\textbf{Definition} (\texttt{is-sigma-union-of}).\;For all $\mathit{cb}$ and $\mathit{es}$, 
\begin{equation*}
\operatorname{is\text{-}sigma\text{-}union\text{-}of}(\mathit{cb}, \mathit{es}) \Leftrightarrow \exists \operatorname{sug\_} \in (\mathbb{N} \to \mathit{cb}).\; (\mathit{es} = \operatorname{big\text{-}union}(\operatorname{sun\_}, \mathbb{N}, \operatorname{sug\_}(\operatorname{sun\_})))
\end{equation*}


\vnbbookfile{theorem-library/mb35-regular}{theorem-library/mb35-regular}
\noindent{\small mb35-regular.scm -- construction of measures, chapter 5 section 2: approximation of Borel measures on metric spaces by closed, open and compact sets (M-B35 / B3, 2026-10-26).  Library: measure.  Helper prefix mr-. Closed and open sets are IS-CLOSED / IS-OPEN; no CLOSURE, no INTERIOR.}\par
\vnbprop{measure-below-approx}{\vnbstmt{\begin{array}{@{}l@{}} \forall a, b, u, f, s. \\ \quad \operatorname{is\text{-}measure}(a, b, u) \wedge \\ \quad f \in (\mathbb{N} \to b) \wedge \\ \quad \forall n \in \mathbb{N}.\; (f(n) \subseteq f((n + 1))) \wedge \\ \quad u(\operatorname{big\text{-}union}(n, \mathbb{N}, f(n))) \in \mathbb{R} \wedge \\ \quad \operatorname{pos\text{-}rr}(s) \Rightarrow \\ \quad \quad \exists k \in \mathbb{N}.\; (u(\operatorname{big\text{-}union}(n, \mathbb{N}, f(n))) \leq (u(f(k)) + s)) \end{array}}}
\vnbprop{closed-borel}{\vnbstmt{\forall s, a.\; (\operatorname{is\text{-}closed}(s, a) \Rightarrow (a \in \operatorname{borel\text{-}sets}(s)))}}
\vnbprop{dist-layer-closed}{\vnbstmt{\begin{array}{@{}l@{}} \forall m, l.; \forall z \in l, k \in \mathbb{N}. \\ \quad (\operatorname{is\text{-}metric\text{-}space}(m) \wedge (l \subseteq \operatorname{pts}(m))) \Rightarrow  \\ \quad \quad \operatorname{is\text{-}closed}(m, \{x \in \operatorname{pts}(m) : (\operatorname{recip}((k + 1)) \leq \operatorname{pt\text{-}set\text{-}dist}(m, x, l))\}) \end{array}}}
\vnbprop{open-closed-inner-approx}{\vnbstmt{\begin{array}{@{}l@{}} \forall s, a, u, b. \\ \quad \operatorname{is\text{-}metric\text{-}space}(s) \wedge \\ \quad \operatorname{is\text{-}measure}(\operatorname{pts}(s), \operatorname{borel\text{-}sets}(s), a) \wedge \\ \quad a(\operatorname{pts}(s)) \in \mathbb{R} \wedge \\ \quad \operatorname{is\text{-}open}(s, u) \wedge \\ \quad \operatorname{pos\text{-}rr}(b) \Rightarrow \\ \quad \quad \exists f.\; (\operatorname{is\text{-}closed}(s, f) \wedge ((f \subseteq u) \wedge (a(\operatorname{complement\text{-}in}(u, f)) \leq b))) \end{array}}}
\vnbprop{closed-union-2}{\vnbstmt{\begin{array}{@{}l@{}} \forall s, a, b. \\ \quad (\operatorname{is\text{-}closed}(s, a) \wedge \operatorname{is\text{-}closed}(s, b)) \Rightarrow \operatorname{is\text{-}closed}(s, (a \cup b)) \end{array}}}
\vnbprop{closed-seg-union}{\vnbstmt{\begin{array}{@{}l@{}} \forall s, c.; \forall a \in (\mathbb{N} \to c), k \in \mathbb{N}. \\ \quad (\operatorname{is\text{-}metric\text{-}space}(s) \wedge \forall n \in \mathbb{N}.\; \operatorname{is\text{-}closed}(s, a(n))) \Rightarrow  \\ \quad \quad \operatorname{is\text{-}closed}(s, \operatorname{big\text{-}union}(j, \mathrm{OS}(k), a(j))) \end{array}}}
\vnbprop{closed-open-approx-union}{\vnbstmt{\begin{array}{@{}l@{}} \forall s, a.; \forall g \in (\mathbb{N} \to \operatorname{borel\text{-}sets}(s)), b. \\ \quad \operatorname{is\text{-}metric\text{-}space}(s) \wedge \\ \quad \operatorname{is\text{-}measure}(\operatorname{pts}(s), \operatorname{borel\text{-}sets}(s), a) \wedge \\ \quad a(\operatorname{pts}(s)) \in \mathbb{R} \wedge \\ \quad \forall n \in \mathbb{N}, b. \\ \quad \quad \operatorname{pos\text{-}rr}(b) \Rightarrow  \\ \quad \quad \quad \exists f, u. \\ \quad \quad \quad \quad \operatorname{is\text{-}closed}(s, f) \wedge \\ \quad \quad \quad \quad \operatorname{is\text{-}open}(s, u) \wedge \\ \quad \quad \quad \quad f \subseteq g(n) \wedge \\ \quad \quad \quad \quad g(n) \subseteq u \wedge \\ \quad \quad \quad \quad a(\operatorname{complement\text{-}in}(u, f)) \leq b \wedge \\ \quad \operatorname{pos\text{-}rr}(b) \Rightarrow \\ \quad \quad \exists f, u. \\ \quad \quad \quad \operatorname{is\text{-}closed}(s, f) \wedge \\ \quad \quad \quad \operatorname{is\text{-}open}(s, u) \wedge \\ \quad \quad \quad f \subseteq \operatorname{big\text{-}union}(c, \mathbb{N}, g(c)) \wedge \\ \quad \quad \quad \operatorname{big\text{-}union}(c, \mathbb{N}, g(c)) \subseteq u \wedge \\ \quad \quad \quad a(\operatorname{complement\text{-}in}(u, f)) \leq b \end{array}}}
\vnbprop{borel-closed-open-approx-finite}{\vnbstmt{\begin{array}{@{}l@{}} \forall s, a.; \forall e \in \operatorname{borel\text{-}sets}(s), b. \\ \quad (\operatorname{is\text{-}metric\text{-}space}(s) \wedge (\operatorname{is\text{-}measure}(\operatorname{pts}(s), \operatorname{borel\text{-}sets}(s), a) \wedge ((a(\operatorname{pts}(s)) \in \mathbb{R}) \wedge \operatorname{pos\text{-}rr}(b)))) \Rightarrow  \\ \quad \quad \exists f, u. \\ \quad \quad \quad (\operatorname{is\text{-}closed}(s, f) \wedge (\operatorname{is\text{-}open}(s, u) \wedge ((f \subseteq e) \wedge ((e \subseteq u) \wedge (a(\operatorname{complement\text{-}in}(u, f)) \leq b))))) \end{array}}}
\vnbprop{borel-closed-open-approx}{\vnbstmt{\begin{array}{@{}l@{}} \forall s, a.; \forall e \in \operatorname{borel\text{-}sets}(s), b. \\ \quad (\operatorname{is\text{-}metric\text{-}space}(s) \wedge (\operatorname{is\text{-}measure}(\operatorname{pts}(s), \operatorname{borel\text{-}sets}(s), a) \wedge ((a(\operatorname{pts}(s)) = 1) \wedge \operatorname{pos\text{-}rr}(b)))) \Rightarrow  \\ \quad \quad \exists f, u. \\ \quad \quad \quad \operatorname{is\text{-}closed}(s, f) \wedge \\ \quad \quad \quad \operatorname{is\text{-}open}(s, u) \wedge \\ \quad \quad \quad f \subseteq e \wedge \\ \quad \quad \quad e \subseteq u \wedge \\ \quad \quad \quad a(\operatorname{complement\text{-}in}(u, f)) \leq b \wedge \\ \quad \quad \quad (a(u) - b) \leq a(e) \wedge \\ \quad \quad \quad a(e) \leq (a(f) + b) \end{array}}}
\vnbprop{measure-restrict-to-set}{\vnbstmt{\begin{array}{@{}l@{}} \forall a, b, u.; \forall v \in b. \\ \quad \operatorname{is\text{-}measure}(a, b, u) \Rightarrow  \\ \quad \quad \operatorname{is\text{-}measure}(a, b, (\lambda r \in b.\; u((r \cap v)))) \end{array}}}
\vnbprop{borel-outer-approx}{\vnbstmt{\begin{array}{@{}l@{}} \forall s, a.; \forall e \in \operatorname{borel\text{-}sets}(s), b. \\ \quad \operatorname{is\text{-}metric\text{-}space}(s) \wedge \\ \quad \operatorname{is\text{-}measure}(\operatorname{pts}(s), \operatorname{borel\text{-}sets}(s), a) \wedge \\ \quad \exists c. \\ \quad \quad (\forall n \in \mathbb{N}.\; (\operatorname{is\text{-}open}(s, c(n)) \wedge (a(c(n)) \in \mathbb{R})) \wedge (\operatorname{big\text{-}union}(d, \mathbb{N}, c(d)) = \operatorname{pts}(s))) \wedge \\ \quad a(e) \in \mathbb{R} \wedge \\ \quad \operatorname{pos\text{-}rr}(b) \Rightarrow \\ \quad \quad \exists u.\; (\operatorname{is\text{-}open}(s, u) \wedge ((e \subseteq u) \wedge (a(u) \leq (a(e) + b)))) \end{array}}}
\vnbprop{borel-outer-regular}{\vnbstmt{\begin{array}{@{}l@{}} \forall s, u.; \forall e \in \operatorname{borel\text{-}sets}(s). \\ \quad \operatorname{is\text{-}metric\text{-}space}(s) \wedge \\ \quad \operatorname{is\text{-}measure}(\operatorname{pts}(s), \operatorname{borel\text{-}sets}(s), u) \wedge \\ \quad \exists a. \\ \quad \quad (\forall n \in \mathbb{N}.\; (\operatorname{is\text{-}open}(s, a(n)) \wedge (u(a(n)) \in \mathbb{R})) \wedge (\operatorname{big\text{-}union}(b, \mathbb{N}, a(b)) = \operatorname{pts}(s))) \Rightarrow \\ \quad \quad u(e) \\ \quad \quad =\; \operatorname{einf}(\{v \in \operatorname{rr\text{-}pos\text{-}star} : \exists \mathit{mru} \in \operatorname{power}(\operatorname{pts}(s)).\; (\operatorname{is\text{-}open}(s, \mathit{mru}) \wedge ((e \subseteq \mathit{mru}) \wedge (v = u(\mathit{mru}))))\}) \end{array}}}
\vnbprop{borel-measure-compare-open}{\vnbstmt{\begin{array}{@{}l@{}} \forall s, a, b.; \forall e \in \operatorname{borel\text{-}sets}(s). \\ \quad \operatorname{is\text{-}metric\text{-}space}(s) \wedge \\ \quad \operatorname{is\text{-}measure}(\operatorname{pts}(s), \operatorname{borel\text{-}sets}(s), a) \wedge \\ \quad \operatorname{is\text{-}measure}(\operatorname{pts}(s), \operatorname{borel\text{-}sets}(s), b) \wedge \\ \quad \forall u.\; (\operatorname{is\text{-}open}(s, u) \Rightarrow (b(u) \leq a(u))) \wedge \\ \quad \exists c. \\ \quad \quad (\forall n \in \mathbb{N}.\; (\operatorname{is\text{-}open}(s, c(n)) \wedge (a(c(n)) \in \mathbb{R})) \wedge (\operatorname{big\text{-}union}(d, \mathbb{N}, c(d)) = \operatorname{pts}(s))) \Rightarrow \\ \quad \quad b(e) \leq a(e) \end{array}}}
\vnbprop{regular-iff-tight}{\vnbstmt{\begin{array}{@{}l@{}} \forall s, u. \\ \quad (\operatorname{is\text{-}metric\text{-}space}(s) \wedge (\operatorname{is\text{-}measure}(\operatorname{pts}(s), \operatorname{borel\text{-}sets}(s), u) \wedge (u(\operatorname{pts}(s)) \in \mathbb{R}))) \Rightarrow  \\ \quad \quad \operatorname{is\text{-}regular\text{-}measure}(s, u) \\ \quad \quad \Leftrightarrow\; u(\operatorname{pts}(s)) \\ \quad \quad =\; \operatorname{esup}(\{v \in \operatorname{rr\text{-}pos\text{-}star} : \exists k \in \operatorname{power}(\operatorname{pts}(s)).\; ((k \subseteq \operatorname{pts}(s)) \wedge (\operatorname{is\text{-}compact}(\operatorname{subspace\text{-}ms}(s, k)) \wedge (v = u(k))))\}) \end{array}}}
\vnbprop{measure-seg-union-approx}{\vnbstmt{\begin{array}{@{}l@{}} \forall a, b, u, g, s. \\ \quad (\operatorname{is\text{-}measure}(a, b, u) \wedge ((g \in (\mathbb{N} \to b)) \wedge ((u(\operatorname{big\text{-}union}(n, \mathbb{N}, g(n))) \in \mathbb{R}) \wedge \operatorname{pos\text{-}rr}(s)))) \Rightarrow  \\ \quad \quad \exists m \in \mathbb{N}. \\ \quad \quad \quad u(\operatorname{big\text{-}union}(n, \mathbb{N}, g(n))) \leq (u(\operatorname{big\text{-}union}(j, \mathrm{OS}(m), g(j))) + s) \end{array}}}
\vnbprop{tight-of-full-balls}{\vnbstmt{\begin{array}{@{}l@{}} \forall s, u.; \forall x \in (\mathbb{N} \to \operatorname{pts}(s)), a. \\ \quad \operatorname{is\text{-}complete}(s) \wedge \\ \quad \operatorname{is\text{-}measure}(\operatorname{pts}(s), \operatorname{borel\text{-}sets}(s), u) \wedge \\ \quad u(\operatorname{pts}(s)) \in \mathbb{R} \wedge \\ \quad \forall r.\; (\operatorname{pos\text{-}rr}(r) \Rightarrow (u(\operatorname{big\text{-}union}(i, \mathbb{N}, \operatorname{ball}(s, x(i), r))) = u(\operatorname{pts}(s)))) \wedge \\ \quad \operatorname{pos\text{-}rr}(a) \Rightarrow \\ \quad \quad \exists k. \\ \quad \quad \quad (k \subseteq \operatorname{pts}(s)) \wedge (\operatorname{is\text{-}compact}(\operatorname{subspace\text{-}ms}(s, k)) \wedge (u(\operatorname{pts}(s)) \leq (u(k) + a))) \end{array}}}
\vnbprop{tight-implies-regular}{\vnbstmt{\begin{array}{@{}l@{}} \forall s, u. \\ \quad \operatorname{is\text{-}metric\text{-}space}(s) \wedge \\ \quad \operatorname{is\text{-}measure}(\operatorname{pts}(s), \operatorname{borel\text{-}sets}(s), u) \wedge \\ \quad u(\operatorname{pts}(s)) \in \mathbb{R} \wedge \\ \quad \forall a. \\ \quad \quad \operatorname{pos\text{-}rr}(a) \Rightarrow  \\ \quad \quad \quad \exists k. \\ \quad \quad \quad \quad (k \subseteq \operatorname{pts}(s)) \wedge (\operatorname{is\text{-}compact}(\operatorname{subspace\text{-}ms}(s, k)) \wedge (u(\operatorname{pts}(s)) \leq (u(k) + a))) \Rightarrow \\ \quad \quad \operatorname{is\text{-}regular\text{-}measure}(s, u) \end{array}}}
\vnbprop{polish-probability-regular}{\vnbstmt{\begin{array}{@{}l@{}} \forall s, u. \\ \quad (\operatorname{is\text{-}polish}(s) \wedge (\operatorname{is\text{-}measure}(\operatorname{pts}(s), \operatorname{borel\text{-}sets}(s), u) \wedge (u(\operatorname{pts}(s)) = 1))) \Rightarrow  \\ \quad \quad \operatorname{is\text{-}regular\text{-}measure}(s, u) \end{array}}}
\vnbprop{regular-iff-separable-support}{\vnbstmt{\begin{array}{@{}l@{}} \forall s, u. \\ \quad (\operatorname{is\text{-}complete}(s) \wedge (\operatorname{is\text{-}measure}(\operatorname{pts}(s), \operatorname{borel\text{-}sets}(s), u) \wedge (u(\operatorname{pts}(s)) = 1))) \Rightarrow  \\ \quad \quad \operatorname{is\text{-}regular\text{-}measure}(s, u) \\ \quad \quad \Leftrightarrow\; \exists e \in \operatorname{borel\text{-}sets}(s).\; (\operatorname{is\text{-}separable}(\operatorname{subspace\text{-}ms}(s, e)) \wedge (u(e) = 1)) \end{array}}}

\vnbbookfile{theorem-library/radon-nikodym}{theorem-library/radon-nikodym}
\noindent{\small radon-nikodym.scm -- WAVE 2 (2026-10-27..30), wired 2026-10-30 from scratchpad/md/md-rn.scm, scratchpad/md/me-rn.scm, scratchpad/md/me-rn2.scm, scratchpad/md/me-rn3.scm. Library: measure.}\par
\vnbprop{density-implies-absolutely-continuous}{\vnbstmt{\begin{array}{@{}l@{}} \forall a, b, u, c, f. \\ \quad \operatorname{is\text{-}measure}(a, b, u) \wedge \\ \quad \operatorname{is\text{-}measure}(a, b, c) \wedge \\ \quad \operatorname{is\text{-}measurable\text{-}fn}(a, b, f) \wedge \\ \quad \forall \mathit{rna} \in b. \\ \quad \quad u(\mathit{rna}) = \operatorname{integral}(a, b, c, \operatorname{ptwise\text{-}etimes}(a, f, \operatorname{indicator}(a, \mathit{rna}))) \Rightarrow \\ \quad \quad \operatorname{is\text{-}absolutely\text{-}continuous}(a, b, u, c) \end{array}}}
\vnbprop{simple-integral-measure-le}{\vnbstmt{\begin{array}{@{}l@{}} \forall a, b, u, c, s. \\ \quad \operatorname{is\text{-}measure}(a, b, u) \wedge \\ \quad \operatorname{is\text{-}measure}(a, b, c) \wedge \\ \quad \forall \mathit{mma} \in b.\; (u(\mathit{mma}) \leq c(\mathit{mma})) \wedge \\ \quad \operatorname{is\text{-}simple\text{-}fn}(a, b, s) \Rightarrow \\ \quad \quad \operatorname{simple\text{-}integral}(a, u, s) \leq \operatorname{simple\text{-}integral}(a, c, s) \end{array}}}
\vnbprop{integral-measure-le}{\vnbstmt{\begin{array}{@{}l@{}} \forall a, b, u, c, h. \\ \quad \operatorname{is\text{-}measure}(a, b, u) \wedge \\ \quad \operatorname{is\text{-}measure}(a, b, c) \wedge \\ \quad \forall \mathit{mma} \in b.\; (u(\mathit{mma}) \leq c(\mathit{mma})) \wedge \\ \quad \operatorname{is\text{-}measurable\text{-}fn}(a, b, h) \Rightarrow \\ \quad \quad \operatorname{integral}(a, b, u, h) \leq \operatorname{integral}(a, b, c, h) \end{array}}}
\vnbprop{lp-norm-abs-fn}{\vnbstmt{\begin{array}{@{}l@{}} \forall a, b, u, p, g. \\ \quad (\operatorname{is\text{-}measure}(a, b, u) \wedge \operatorname{is\text{-}real\text{-}measurable\text{-}fn}(a, b, g)) \Rightarrow  \\ \quad \quad (\|{\operatorname{abs\text{-}fn}(a, g)}\|_{p} \simeq \|{g}\|_{p}) \end{array}}}
\vnbprop{l2-finite-integrable}{\vnbstmt{\begin{array}{@{}l@{}} \forall a, b, u, g. \\ \quad (\operatorname{is\text{-}measure}(a, b, u) \wedge ((u(a) \in \mathbb{R}) \wedge (g \in \operatorname{l2\text{-}set}(a, b, u)))) \Rightarrow  \\ \quad \quad \operatorname{is\text{-}integrable}(a, b, u, g) \wedge \\ \quad \quad \operatorname{integral}(a, b, u, \operatorname{abs\text{-}fn}(a, g)) \leq (\sqrt{\langle\cdot\mid\cdot\rangle_{\operatorname{l2\text{-}space}(a, b, u)}(g, g)} \cdot \sqrt{u(a)}) \end{array}}}
\vnbprop{rn-finite-dominated}{\vnbstmt{\begin{array}{@{}l@{}} \forall a, b, u, c. \\ \quad (\operatorname{is\text{-}measure}(a, b, u) \wedge (\operatorname{is\text{-}measure}(a, b, c) \wedge ((c(a) \in \mathbb{R}) \wedge \forall \mathit{mma} \in b.\; (u(\mathit{mma}) \leq c(\mathit{mma}))))) \Rightarrow  \\ \quad \quad \exists d. \\ \quad \quad \quad \operatorname{is\text{-}measurable\text{-}fn}(a, b, d) \wedge \\ \quad \quad \quad \forall x \in a.\; (d(x) \in \mathbb{R}) \wedge \\ \quad \quad \quad \forall \mathit{rna} \in b. \\ \quad \quad \quad \quad u(\mathit{rna}) = \operatorname{integral}(a, b, c, \operatorname{ptwise\text{-}etimes}(a, d, \operatorname{indicator}(a, \mathit{rna}))) \end{array}}}
\vnbprop{etimes-left-commute}{\vnbstmt{\begin{array}{@{}l@{}} \forall a, b, c. \\ \quad ((a \in \operatorname{rr\text{-}pos\text{-}star}) \wedge ((b \in \operatorname{rr\text{-}pos\text{-}star}) \wedge ((c \in \mathbb{R}) \wedge (0 \leq c)))) \Rightarrow  \\ \quad \quad \operatorname{etimes}(a, \operatorname{etimes}(c, b)) = \operatorname{etimes}(c, \operatorname{etimes}(a, b)) \end{array}}}
\vnbprop{measure-sum}{\vnbstmt{\begin{array}{@{}l@{}} \forall a, b, u, c. \\ \quad (\operatorname{is\text{-}measure}(a, b, u) \wedge \operatorname{is\text{-}measure}(a, b, c)) \Rightarrow  \\ \quad \quad \operatorname{is\text{-}measure}(a, b, (\lambda \mathit{rsa} \in b.\; \operatorname{eplus}(u(\mathit{rsa}), c(\mathit{rsa})))) \end{array}}}
\vnbprop{scale-inf-indicator-measurable}{\vnbstmt{\begin{array}{@{}l@{}} \forall b, c, a. \\ \quad (\operatorname{is\text{-}sigma\text{-}algebra}(b, c) \wedge (a \in c)) \Rightarrow  \\ \quad \quad \operatorname{is\text{-}measurable\text{-}fn}(b, c, \operatorname{ptwise\text{-}scale}(b, \operatorname{pos\text{-}inf}, \operatorname{indicator}(b, a))) \end{array}}}
\vnbprop{esup-const-image}{\vnbstmt{\forall c \in \operatorname{rr\text{-}pos\text{-}star}.\; (\operatorname{esup}(\operatorname{image}((\lambda u \in \mathbb{N}.\; c), \mathbb{N})) = c)}}
\vnbprop{ptwise-etimes-sup}{\vnbstmt{\begin{array}{@{}l@{}} \forall a, h, f. \\ \quad a \in \mathbf{Set} \wedge \\ \quad h \in (a \to \operatorname{rr\text{-}pos\text{-}star}) \wedge \\ \quad f \in (\mathbb{N} \to (a \to \operatorname{rr\text{-}pos\text{-}star})) \wedge \\ \quad \forall n \in \mathbb{N}.\; \operatorname{ptwise\text{-}le}(a, f(n), f((n + 1))) \Rightarrow \\ \quad \quad \operatorname{ptwise\text{-}etimes}(a, h, \operatorname{ptwise\text{-}sup}(a, f)) \\ \quad \quad =\; \operatorname{ptwise\text{-}sup}(a, (\lambda k \in \mathbb{N}.\; (\lambda x \in a.\; \operatorname{etimes}(h(x), f(k)(x))))) \end{array}}}
\vnbprop{density-change}{\vnbstmt{\begin{array}{@{}l@{}} \forall a, b, o, u, h, g. \\ \quad \operatorname{is\text{-}measure}(a, b, o) \wedge \\ \quad \operatorname{is\text{-}measure}(a, b, u) \wedge \\ \quad \operatorname{is\text{-}measurable\text{-}fn}(a, b, h) \wedge \\ \quad \forall \mathit{rna} \in b. \\ \quad \quad u(\mathit{rna}) = \operatorname{integral}(a, b, o, \operatorname{ptwise\text{-}etimes}(a, h, \operatorname{indicator}(a, \mathit{rna}))) \wedge \\ \quad \operatorname{is\text{-}measurable\text{-}fn}(a, b, g) \Rightarrow \\ \quad \quad \operatorname{integral}(a, b, u, g) = \operatorname{integral}(a, b, o, \operatorname{ptwise\text{-}etimes}(a, h, g)) \end{array}}}
\vnbprop{recip-compare}{\vnbstmt{\forall a, t.\; (((a \in \mathbb{R}) \wedge ((t \in \mathbb{R}) \wedge (0 < t))) \Rightarrow ((a < \operatorname{recip}(t)) \Leftrightarrow ((a \cdot t) < 1)))}}
\vnbprop{measurable-fn-recip}{\vnbstmt{\begin{array}{@{}l@{}} \forall a, b, h. \\ \quad (\operatorname{is\text{-}measurable\text{-}fn}(a, b, h) \wedge \forall x \in a.\; (h(x) \in \mathbb{R})) \Rightarrow  \\ \quad \quad \operatorname{is\text{-}measurable\text{-}fn}(a, b, (\lambda \mathit{rrx} \in a.\; \left(\text{if }(0 < h(\mathit{rrx}))\text{ then }\operatorname{recip}(h(\mathit{rrx}))\text{ else }0\right))) \end{array}}}
\vnbprop{radon-nikodym-finite}{\vnbstmt{\begin{array}{@{}l@{}} \forall a, b, u, c. \\ \quad \operatorname{is\text{-}measure}(a, b, u) \wedge \\ \quad \operatorname{is\text{-}measure}(a, b, c) \wedge \\ \quad u(a) \in \mathbb{R} \wedge \\ \quad c(a) \in \mathbb{R} \wedge \\ \quad \operatorname{is\text{-}absolutely\text{-}continuous}(a, b, u, c) \Rightarrow \\ \quad \quad \exists q. \\ \quad \quad \quad \operatorname{is\text{-}measurable\text{-}fn}(a, b, q) \wedge \\ \quad \quad \quad \forall \mathit{rna} \in b. \\ \quad \quad \quad \quad u(\mathit{rna}) = \operatorname{integral}(a, b, c, \operatorname{ptwise\text{-}etimes}(a, q, \operatorname{indicator}(a, \mathit{rna}))) \end{array}}}
\vnbprop{increasing-sets-chain}{\vnbstmt{\begin{array}{@{}l@{}} \forall m.; \forall j, k \in \mathbb{N}. \\ \quad (\forall k \in \mathbb{N}.\; (m(k) \in \mathbf{Set}) \wedge (\forall k \in \mathbb{N}.\; (m(k) \subseteq m((k + 1))) \wedge (k \leq j))) \Rightarrow  \\ \quad \quad m(k) \subseteq m(j) \end{array}}}
\vnbprop{sigma-finite-common-partition}{\vnbstmt{\begin{array}{@{}l@{}} \forall a, b, u, c. \\ \quad \operatorname{is\text{-}measure}(a, b, u) \wedge \\ \quad \operatorname{is\text{-}measure}(a, b, c) \wedge \\ \quad \operatorname{is\text{-}sigma\text{-}finite}(a, b, u) \wedge \\ \quad \operatorname{is\text{-}sigma\text{-}finite}(a, b, c) \Rightarrow \\ \quad \quad \exists d \in (\mathbb{N} \to b). \\ \quad \quad \quad \forall i, j \in \mathbb{N}. \\ \quad \quad \quad \quad \neg (i = j) \Rightarrow ((d(i) \cap d(j)) = \emptyset) \wedge \\ \quad \quad \quad \operatorname{big\text{-}union}(n, \mathbb{N}, d(n)) = a \wedge \\ \quad \quad \quad \forall \mathit{sfn} \in \mathbb{N}.\; (u(d(\mathit{sfn})) \in \mathbb{R}) \wedge \\ \quad \quad \quad \forall \mathit{sfn} \in \mathbb{N}.\; (c(d(\mathit{sfn})) \in \mathbb{R}) \end{array}}}
\vnbprop{rn-on-finite-piece}{\vnbstmt{\begin{array}{@{}l@{}} \forall a, b, u, c, v. \\ \quad \operatorname{is\text{-}measure}(a, b, u) \wedge \\ \quad \operatorname{is\text{-}measure}(a, b, c) \wedge \\ \quad \operatorname{is\text{-}absolutely\text{-}continuous}(a, b, u, c) \wedge \\ \quad v \in b \wedge \\ \quad u(v) \in \mathbb{R} \wedge \\ \quad c(v) \in \mathbb{R} \Rightarrow \\ \quad \quad \exists q. \\ \quad \quad \quad \operatorname{is\text{-}measurable\text{-}fn}(a, b, q) \wedge \\ \quad \quad \quad \forall \mathit{rna} \in b. \\ \quad \quad \quad \quad u((\mathit{rna} \cap v)) \\ \quad \quad \quad \quad =\; \operatorname{integral}(a, b, c, \operatorname{ptwise\text{-}etimes}(a, \operatorname{indicator}(a, \mathit{rna}), \operatorname{ptwise\text{-}etimes}(a, \operatorname{indicator}(a, v), q))) \end{array}}}
\vnbprop{radon-nikodym-sigma-finite}{\vnbstmt{\begin{array}{@{}l@{}} \forall a, b, u, c. \\ \quad \operatorname{is\text{-}measure}(a, b, u) \wedge \\ \quad \operatorname{is\text{-}measure}(a, b, c) \wedge \\ \quad \operatorname{is\text{-}sigma\text{-}finite}(a, b, u) \wedge \\ \quad \operatorname{is\text{-}sigma\text{-}finite}(a, b, c) \wedge \\ \quad \operatorname{is\text{-}absolutely\text{-}continuous}(a, b, u, c) \Rightarrow \\ \quad \quad \exists q. \\ \quad \quad \quad \operatorname{is\text{-}measurable\text{-}fn}(a, b, q) \wedge \\ \quad \quad \quad \forall \mathit{rna} \in b. \\ \quad \quad \quad \quad u(\mathit{rna}) = \operatorname{integral}(a, b, c, \operatorname{ptwise\text{-}etimes}(a, q, \operatorname{indicator}(a, \mathit{rna}))) \end{array}}}
\vnbprop{radon-nikodym}{\vnbstmt{\begin{array}{@{}l@{}} \forall a, b, u, c. \\ \quad \operatorname{is\text{-}measure}(a, b, u) \wedge \\ \quad \operatorname{is\text{-}measure}(a, b, c) \wedge \\ \quad \operatorname{is\text{-}sigma\text{-}finite}(a, b, u) \wedge \\ \quad \operatorname{is\text{-}sigma\text{-}finite}(a, b, c) \Rightarrow \\ \quad \quad \operatorname{is\text{-}absolutely\text{-}continuous}(a, b, u, c) \\ \quad \quad \Leftrightarrow\; \exists q. \\ \quad \quad \quad \operatorname{is\text{-}measurable\text{-}fn}(a, b, q) \wedge \\ \quad \quad \quad \forall \mathit{rna} \in b. \\ \quad \quad \quad \quad u(\mathit{rna}) = \operatorname{integral}(a, b, c, \operatorname{ptwise\text{-}etimes}(a, q, \operatorname{indicator}(a, \mathit{rna}))) \end{array}}}

\vnbbookfile{theorem-library/l1-dual}{theorem-library/l1-dual}
\noindent{\small l1-dual.scm -- WAVE 2 (2026-10-27..30), wired 2026-10-30 from scratchpad/md/me-du1.scm, scratchpad/md/me-du2.scm, scratchpad/md/me-du3.scm, scratchpad/md/me-du4.scm, scratchpad/md/me-du5.scm. Library: measure.}\par
\vnbprop{linf-norm-essential-bound}{\vnbstmt{\begin{array}{@{}l@{}} \forall a, b, u, f. \\ \quad (\operatorname{is\text{-}measure}(a, b, u) \wedge (\operatorname{is\text{-}real\text{-}measurable\text{-}fn}(a, b, f) \wedge (\operatorname{linf\text{-}norm}(a, b, u, f) \in \mathbb{R}))) \Rightarrow  \\ \quad \quad ((0 \leq \operatorname{linf\text{-}norm}(a, b, u, f)) \wedge (u(\{x \in a : (\operatorname{linf\text{-}norm}(a, b, u, f) < \lvert f(x) \rvert)\}) = 0)) \end{array}}}
\vnbprop{l1-set-membership}{\vnbstmt{\begin{array}{@{}l@{}} \forall a, b, u, g. \\ \quad (g \in \operatorname{l1\text{-}set}(a, b, u)) \Leftrightarrow ((g \in (a \to \mathbb{R})) \wedge \operatorname{is\text{-}integrable}(a, b, u, g)) \end{array}}}
\vnbprop{l1-set-is-set}{\vnbstmt{\forall a \in \mathbf{Set}, b, u.\; (\operatorname{l1\text{-}set}(a, b, u) \in \mathbf{Set})}}
\vnbprop{l1-pairing-value}{\vnbstmt{\begin{array}{@{}l@{}} \forall a, b, u, f.; \forall g \in \operatorname{l1\text{-}set}(a, b, u). \\ \quad (\operatorname{l1\text{-}pairing}(a, b, u, f)(g) \simeq \operatorname{rintegral}(a, b, u, (\lambda x \in a.\; (f(x) \cdot g(x))))) \end{array}}}
\vnbprop{l1-product-bound}{\vnbstmt{\begin{array}{@{}l@{}} \forall a, b, u, f, m, g. \\ \quad \operatorname{is\text{-}measure}(a, b, u) \wedge \\ \quad \operatorname{is\text{-}real\text{-}measurable\text{-}fn}(a, b, f) \wedge \\ \quad m \in \mathbb{R} \wedge \\ \quad 0 \leq m \wedge \\ \quad u(\{x \in a : (m < \lvert f(x) \rvert)\}) = 0 \wedge \\ \quad g \in \operatorname{l1\text{-}set}(a, b, u) \Rightarrow \\ \quad \quad \operatorname{is\text{-}integrable}(a, b, u, (\lambda \mathit{sx} \in a.\; (f(\mathit{sx}) \cdot g(\mathit{sx})))) \wedge \\ \quad \quad \operatorname{integral}(a, b, u, \operatorname{abs\text{-}fn}(a, (\lambda \mathit{sx} \in a.\; (f(\mathit{sx}) \cdot g(\mathit{sx}))))) \\ \quad \quad \leq\; m \cdot \operatorname{integral}(a, b, u, \operatorname{abs\text{-}fn}(a, g)) \end{array}}}
\vnbprop{l1-pairing-bounded}{\vnbstmt{\begin{array}{@{}l@{}} \forall a, b, u, f. \\ \quad (\operatorname{is\text{-}measure}(a, b, u) \wedge (\operatorname{is\text{-}real\text{-}measurable\text{-}fn}(a, b, f) \wedge (\operatorname{linf\text{-}norm}(a, b, u, f) \in \mathbb{R}))) \Rightarrow  \\ \quad \quad \operatorname{is\text{-}bounded\text{-}l1\text{-}functional}(a, b, u, \operatorname{l1\text{-}pairing}(a, b, u, f), \operatorname{linf\text{-}norm}(a, b, u, f)) \end{array}}}
\vnbprop{linf-norm-in}{\vnbstmt{\forall a, b, u, f.\; (\operatorname{linf\text{-}norm}(a, b, u, f) \in \operatorname{rr\text{-}pos\text{-}star})}}
\vnbprop{indicator-l1}{\vnbstmt{\begin{array}{@{}l@{}} \forall a, b, u, p. \\ \quad (\operatorname{is\text{-}measure}(a, b, u) \wedge ((p \in b) \wedge (u(p) \in \mathbb{R}))) \Rightarrow  \\ \quad \quad \operatorname{indicator}(a, p) \in \operatorname{l1\text{-}set}(a, b, u) \wedge \\ \quad \quad \operatorname{integral}(a, b, u, \operatorname{abs\text{-}fn}(a, \operatorname{indicator}(a, p))) = u(p) \wedge \\ \quad \quad \operatorname{rintegral}(a, b, u, \operatorname{indicator}(a, p)) = u(p) \end{array}}}
\vnbprop{positive-excess-null}{\vnbstmt{\begin{array}{@{}l@{}} \forall a, b, d, p, u, c. \\ \quad \operatorname{is\text{-}measure}(a, b, d) \wedge \\ \quad p \in b \wedge \\ \quad d(p) \in \mathbb{R} \wedge \\ \quad c \in \mathbb{R} \wedge \\ \quad 0 \leq c \wedge \\ \quad u \in \operatorname{l1\text{-}set}(a, b, d) \wedge \\ \quad \forall x \in a.\; ((x \in p) \Rightarrow (c < u(x))) \wedge \\ \quad \forall x \in a.\; (\neg (x \in p) \Rightarrow (u(x) = 0)) \wedge \\ \quad \operatorname{rintegral}(a, b, d, u) \leq (c \cdot d(p)) \Rightarrow \\ \quad \quad d(p) = 0 \end{array}}}
\vnbprop{linf-bound-from-indicators}{\vnbstmt{\begin{array}{@{}l@{}} \forall a, b, u, f, c. \\ \quad \operatorname{is\text{-}measure}(a, b, u) \wedge \\ \quad \operatorname{is\text{-}sigma\text{-}finite}(a, b, u) \wedge \\ \quad \operatorname{is\text{-}real\text{-}measurable\text{-}fn}(a, b, f) \wedge \\ \quad c \in \mathbb{R} \wedge \\ \quad 0 \leq c \wedge \\ \quad \forall q \in b. \\ \quad \quad (u(q) \in \mathbb{R}) \Rightarrow  \\ \quad \quad \quad \operatorname{is\text{-}integrable}(a, b, u, (\lambda x \in a.\; (f(x) \cdot \operatorname{indicator}(a, q)(x)))) \wedge \\ \quad \quad \quad \lvert \operatorname{rintegral}(a, b, u, (\lambda x \in a.\; (f(x) \cdot \operatorname{indicator}(a, q)(x)))) \rvert \leq (c \cdot u(q)) \Rightarrow \\ \quad \quad \operatorname{linf\text{-}norm}(a, b, u, f) \leq c \end{array}}}
\vnbprop{l1-pairing-norm-bound}{\vnbstmt{\begin{array}{@{}l@{}} \forall a, b, u, f, c. \\ \quad \operatorname{is\text{-}measure}(a, b, u) \wedge \\ \quad \operatorname{is\text{-}sigma\text{-}finite}(a, b, u) \wedge \\ \quad \operatorname{is\text{-}real\text{-}measurable\text{-}fn}(a, b, f) \wedge \\ \quad \operatorname{is\text{-}bounded\text{-}l1\text{-}functional}(a, b, u, \operatorname{l1\text{-}pairing}(a, b, u, f), c) \wedge \\ \quad \forall g \in \operatorname{l1\text{-}set}(a, b, u). \\ \quad \quad \operatorname{is\text{-}integrable}(a, b, u, (\lambda x \in a.\; (f(x) \cdot g(x)))) \Rightarrow \\ \quad \quad \operatorname{linf\text{-}norm}(a, b, u, f) \leq c \end{array}}}
\vnbprop{l1-functionals-agree}{\vnbstmt{\begin{array}{@{}l@{}} \forall a, b, u, c, d, f, h.; \forall i \in \operatorname{l1\text{-}set}(a, b, u). \\ \quad \operatorname{is\text{-}measure}(a, b, u) \wedge \\ \quad \operatorname{is\text{-}bounded\text{-}l1\text{-}functional}(a, b, u, c, d) \wedge \\ \quad \operatorname{is\text{-}bounded\text{-}l1\text{-}functional}(a, b, u, f, h) \wedge \\ \quad \forall g \in \operatorname{l1\text{-}set}(a, b, u).\; ((g \in \operatorname{l2\text{-}set}(a, b, u)) \Rightarrow (c(g) = f(g))) \Rightarrow \\ \quad \quad c(i) = f(i) \end{array}}}
\vnbprop{finite-measure-sigma-finite}{\vnbstmt{\begin{array}{@{}l@{}} \forall a, b, u. \\ \quad (\operatorname{is\text{-}measure}(a, b, u) \wedge (u(a) \in \mathbb{R})) \Rightarrow \operatorname{is\text{-}sigma\text{-}finite}(a, b, u) \end{array}}}
\vnbprop{l1-dual-onto-finite}{\vnbstmt{\begin{array}{@{}l@{}} \forall a, b, u, f, c. \\ \quad (\operatorname{is\text{-}measure}(a, b, u) \wedge ((u(a) \in \mathbb{R}) \wedge \operatorname{is\text{-}bounded\text{-}l1\text{-}functional}(a, b, u, f, c))) \Rightarrow  \\ \quad \quad \exists \mathit{ldf}. \\ \quad \quad \quad \operatorname{is\text{-}real\text{-}measurable\text{-}fn}(a, b, \mathit{ldf}) \wedge \\ \quad \quad \quad \operatorname{linf\text{-}norm}(a, b, u, \mathit{ldf}) \leq c \wedge \\ \quad \quad \quad \forall g \in \operatorname{l1\text{-}set}(a, b, u).\; (f(g) = \operatorname{l1\text{-}pairing}(a, b, u, \mathit{ldf})(g)) \end{array}}}
\vnbprop{esum-one-term-prefix}{\vnbstmt{\begin{array}{@{}l@{}} \forall h \in (\mathbb{N} \to \operatorname{rr\text{-}pos\text{-}star}), m, n \in \mathbb{N}. \\ \quad \forall k \in \mathbb{N}.\; (\neg (k = m) \Rightarrow (h(k) = 0)) \Rightarrow  \\ \quad \quad \operatorname{finsum}(\operatorname{rr\text{-}pos\text{-}star\text{-}add\text{-}monoid}, h, \mathrm{OS}(n)) \leq h(m) \wedge \\ \quad \quad (n \leq m) \Rightarrow (\operatorname{finsum}(\operatorname{rr\text{-}pos\text{-}star\text{-}add\text{-}monoid}, h, \mathrm{OS}(n)) = 0) \end{array}}}
\vnbprop{esum-one-term}{\vnbstmt{\begin{array}{@{}l@{}} \forall h \in (\mathbb{N} \to \operatorname{rr\text{-}pos\text{-}star}), m \in \mathbb{N}. \\ \quad \forall k \in \mathbb{N}.\; (\neg (k = m) \Rightarrow (h(k) = 0)) \Rightarrow (\operatorname{esum}(h) = h(m)) \end{array}}}
\vnbprop{density-measure}{\vnbstmt{\begin{array}{@{}l@{}} \forall a, b, u, h. \\ \quad (\operatorname{is\text{-}measure}(a, b, u) \wedge \operatorname{is\text{-}measurable\text{-}fn}(a, b, h)) \Rightarrow  \\ \quad \quad \operatorname{is\text{-}measure}(a, b, (\lambda \mathit{dma} \in b.\; \operatorname{integral}(a, b, u, \operatorname{ptwise\text{-}etimes}(a, h, \operatorname{indicator}(a, \mathit{dma}))))) \end{array}}}
\vnbprop{sigma-finite-weight}{\vnbstmt{\begin{array}{@{}l@{}} \forall a, b, u. \\ \quad (\operatorname{is\text{-}measure}(a, b, u) \wedge \operatorname{is\text{-}sigma\text{-}finite}(a, b, u)) \Rightarrow  \\ \quad \quad \exists t. \\ \quad \quad \quad \operatorname{is\text{-}measurable\text{-}fn}(a, b, t) \wedge \\ \quad \quad \quad \forall x \in a.\; (t(x) \in \mathbb{R}) \wedge \\ \quad \quad \quad \forall x \in a.\; (0 < t(x)) \wedge \\ \quad \quad \quad \operatorname{integral}(a, b, u, t) \in \mathbb{R} \end{array}}}
\vnbprop{weight-transfer}{\vnbstmt{\begin{array}{@{}l@{}} \forall a, b, u, c, w, h. \\ \quad \operatorname{is\text{-}measure}(a, b, u) \wedge \\ \quad \operatorname{is\text{-}measure}(a, b, c) \wedge \\ \quad \operatorname{is\text{-}measurable\text{-}fn}(a, b, w) \wedge \\ \quad \forall x \in a.\; (w(x) \in \mathbb{R}) \wedge \\ \quad \forall x \in a.\; (0 < w(x)) \wedge \\ \quad \forall \mathit{rna} \in b. \\ \quad \quad c(\mathit{rna}) = \operatorname{integral}(a, b, u, \operatorname{ptwise\text{-}etimes}(a, w, \operatorname{indicator}(a, \mathit{rna}))) \wedge \\ \quad \operatorname{is\text{-}real\text{-}measurable\text{-}fn}(a, b, h) \Rightarrow \\ \quad \quad \operatorname{integral}(a, b, c, \operatorname{pos\text{-}part}(a, h)) = \operatorname{integral}(a, b, u, \operatorname{pos\text{-}part}(a, (\lambda \mathit{sx} \in a.\; (h(\mathit{sx}) \cdot w(\mathit{sx}))))) \wedge \\ \quad \quad \operatorname{integral}(a, b, c, \operatorname{neg\text{-}part}(a, h)) = \operatorname{integral}(a, b, u, \operatorname{neg\text{-}part}(a, (\lambda \mathit{sx} \in a.\; (h(\mathit{sx}) \cdot w(\mathit{sx}))))) \wedge \\ \quad \quad \operatorname{integral}(a, b, c, \operatorname{abs\text{-}fn}(a, h)) = \operatorname{integral}(a, b, u, \operatorname{abs\text{-}fn}(a, (\lambda \mathit{sx} \in a.\; (h(\mathit{sx}) \cdot w(\mathit{sx}))))) \end{array}}}
\vnbprop{weight-rintegral}{\vnbstmt{\begin{array}{@{}l@{}} \forall a, b, u, c, w, h. \\ \quad \operatorname{is\text{-}measure}(a, b, u) \wedge \\ \quad \operatorname{is\text{-}measure}(a, b, c) \wedge \\ \quad \operatorname{is\text{-}measurable\text{-}fn}(a, b, w) \wedge \\ \quad \forall x \in a.\; (w(x) \in \mathbb{R}) \wedge \\ \quad \forall x \in a.\; (0 < w(x)) \wedge \\ \quad \forall \mathit{rna} \in b. \\ \quad \quad c(\mathit{rna}) = \operatorname{integral}(a, b, u, \operatorname{ptwise\text{-}etimes}(a, w, \operatorname{indicator}(a, \mathit{rna}))) \wedge \\ \quad \operatorname{is\text{-}real\text{-}measurable\text{-}fn}(a, b, h) \Rightarrow \\ \quad \quad \operatorname{is\text{-}integrable}(a, b, c, h) \Leftrightarrow \operatorname{is\text{-}integrable}(a, b, u, (\lambda \mathit{sx} \in a.\; (h(\mathit{sx}) \cdot w(\mathit{sx})))) \wedge \\ \quad \quad \operatorname{is\text{-}integrable}(a, b, c, h) \Rightarrow  \\ \quad \quad \quad \operatorname{rintegral}(a, b, c, h) = \operatorname{rintegral}(a, b, u, (\lambda \mathit{sx} \in a.\; (h(\mathit{sx}) \cdot w(\mathit{sx})))) \end{array}}}
\vnbprop{l1-dual-onto}{\vnbstmt{\begin{array}{@{}l@{}} \forall a, b, u, f, c. \\ \quad (\operatorname{is\text{-}measure}(a, b, u) \wedge (\operatorname{is\text{-}sigma\text{-}finite}(a, b, u) \wedge \operatorname{is\text{-}bounded\text{-}l1\text{-}functional}(a, b, u, f, c))) \Rightarrow  \\ \quad \quad \exists \mathit{ldf}. \\ \quad \quad \quad \operatorname{is\text{-}real\text{-}measurable\text{-}fn}(a, b, \mathit{ldf}) \wedge \\ \quad \quad \quad \operatorname{linf\text{-}norm}(a, b, u, \mathit{ldf}) \leq c \wedge \\ \quad \quad \quad \forall g \in \operatorname{l1\text{-}set}(a, b, u).\; (f(g) = \operatorname{l1\text{-}pairing}(a, b, u, \mathit{ldf})(g)) \end{array}}}
\vnbprop{linf-is-l1-dual}{\vnbstmt{\begin{array}{@{}l@{}} \forall a, b, u. \\ \quad (\operatorname{is\text{-}measure}(a, b, u) \wedge \operatorname{is\text{-}sigma\text{-}finite}(a, b, u)) \Rightarrow  \\ \quad \quad \forall f. \\ \quad \quad \quad (\operatorname{is\text{-}real\text{-}measurable\text{-}fn}(a, b, f) \wedge (\operatorname{linf\text{-}norm}(a, b, u, f) \in \mathbb{R})) \Rightarrow  \\ \quad \quad \quad \quad \operatorname{is\text{-}bounded\text{-}l1\text{-}functional}(a, b, u, \operatorname{l1\text{-}pairing}(a, b, u, f), \operatorname{linf\text{-}norm}(a, b, u, f)) \wedge \\ \quad \quad \forall f, c. \\ \quad \quad \quad \operatorname{is\text{-}real\text{-}measurable\text{-}fn}(a, b, f) \wedge \\ \quad \quad \quad \operatorname{is\text{-}bounded\text{-}l1\text{-}functional}(a, b, u, \operatorname{l1\text{-}pairing}(a, b, u, f), c) \wedge \\ \quad \quad \quad \forall g \in \operatorname{l1\text{-}set}(a, b, u). \\ \quad \quad \quad \quad \operatorname{is\text{-}integrable}(a, b, u, (\lambda x \in a.\; (f(x) \cdot g(x)))) \Rightarrow \\ \quad \quad \quad \quad \operatorname{linf\text{-}norm}(a, b, u, f) \leq c \wedge \\ \quad \quad \forall f, c. \\ \quad \quad \quad \operatorname{is\text{-}bounded\text{-}l1\text{-}functional}(a, b, u, f, c) \Rightarrow  \\ \quad \quad \quad \quad \exists \mathit{ldf}. \\ \quad \quad \quad \quad \quad \operatorname{is\text{-}real\text{-}measurable\text{-}fn}(a, b, \mathit{ldf}) \wedge \\ \quad \quad \quad \quad \quad \operatorname{linf\text{-}norm}(a, b, u, \mathit{ldf}) \leq c \wedge \\ \quad \quad \quad \quad \quad \forall g \in \operatorname{l1\text{-}set}(a, b, u).\; (f(g) = \operatorname{l1\text{-}pairing}(a, b, u, \mathit{ldf})(g)) \end{array}}}

\vnbbookfile{theorem-library/jordan-decomposition}{theorem-library/jordan-decomposition}
\noindent{\small jordan-decomposition.scm -- WAVE 2 (2026-10-27..30), wired 2026-10-30 from scratchpad/mf/mf-jordan.scm, scratchpad/mf/mf-tv.scm. Library: measure.}\par
\vnbprop{neg-variation-unfold}{\vnbstmt{\begin{array}{@{}l@{}} \forall b, u, a. \\ \quad (\operatorname{neg\text{-}variation}(b, u)(a) \simeq \operatorname{pos\text{-}variation}(b, \operatorname{signed\text{-}neg}(b, u))(a)) \end{array}}}
\vnbprop{neg-variation-finite}{\vnbstmt{\begin{array}{@{}l@{}} \forall a, b, u.; \forall \alpha \in b. \\ \quad (\operatorname{is\text{-}signed\text{-}measure}(a, b, u) \wedge \forall \alpha \in b.\; \neg (u(\alpha) = \operatorname{neg\text{-}inf})) \Rightarrow  \\ \quad \quad \operatorname{neg\text{-}variation}(b, u)(\alpha) \in \mathbb{R} \end{array}}}
\vnbprop{signed-measure-one-variation-finite}{\vnbstmt{\begin{array}{@{}l@{}} \forall a, b, u. \\ \quad \operatorname{is\text{-}signed\text{-}measure}(a, b, u) \Rightarrow  \\ \quad \quad \forall \alpha \in b.\; (\operatorname{pos\text{-}variation}(b, u)(\alpha) \in \mathbb{R}) \vee \\ \quad \quad \forall \alpha \in b.\; (\operatorname{neg\text{-}variation}(b, u)(\alpha) \in \mathbb{R}) \end{array}}}
\vnbprop{signed-split-real}{\vnbstmt{\begin{array}{@{}l@{}} \forall b, c, u, a, e. \\ \quad (\operatorname{is\text{-}signed\text{-}measure}(b, c, u) \wedge ((a \in c) \wedge ((e \in c) \wedge ((e \subseteq a) \wedge (u(a) \in \mathbb{R}))))) \Rightarrow  \\ \quad \quad ((u(e) \in \mathbb{R}) \wedge ((u(\operatorname{complement\text{-}in}(a, e)) \in \mathbb{R}) \wedge (u(a) = (u(e) + u(\operatorname{complement\text{-}in}(a, e)))))) \end{array}}}
\vnbprop{jordan-neg-bound}{\vnbstmt{\begin{array}{@{}l@{}} \forall b, c, u, a. \\ \quad (\operatorname{is\text{-}signed\text{-}measure}(b, c, u) \wedge ((a \in c) \wedge ((u(a) \in \mathbb{R}) \wedge (\operatorname{pos\text{-}variation}(c, u)(a) \in \mathbb{R})))) \Rightarrow  \\ \quad \quad \operatorname{neg\text{-}variation}(c, u)(a) \leq (\operatorname{pos\text{-}variation}(c, u)(a) - u(a)) \end{array}}}
\vnbprop{jordan-pos-bound}{\vnbstmt{\begin{array}{@{}l@{}} \forall b, c, u, a. \\ \quad (\operatorname{is\text{-}signed\text{-}measure}(b, c, u) \wedge ((a \in c) \wedge ((u(a) \in \mathbb{R}) \wedge (\operatorname{neg\text{-}variation}(c, u)(a) \in \mathbb{R})))) \Rightarrow  \\ \quad \quad \operatorname{pos\text{-}variation}(c, u)(a) \leq (u(a) + \operatorname{neg\text{-}variation}(c, u)(a)) \end{array}}}
\vnbprop{jordan-decomposition-value}{\vnbstmt{\begin{array}{@{}l@{}} \forall b, c, u.; \forall a \in c. \\ \quad \operatorname{is\text{-}signed\text{-}measure}(b, c, u) \Rightarrow  \\ \quad \quad u(a) = \operatorname{ext\text{-}sub}(\operatorname{pos\text{-}variation}(c, u)(a), \operatorname{neg\text{-}variation}(c, u)(a)) \end{array}}}
\vnbprop{jordan-decomposition}{\vnbstmt{\begin{array}{@{}l@{}} \forall b, c, u. \\ \quad \operatorname{is\text{-}signed\text{-}measure}(b, c, u) \Rightarrow  \\ \quad \quad \operatorname{is\text{-}measure}(b, c, \operatorname{pos\text{-}variation}(c, u)) \wedge \\ \quad \quad \operatorname{is\text{-}measure}(b, c, \operatorname{neg\text{-}variation}(c, u)) \wedge \\ \quad \quad \forall \alpha \in c.\; (\operatorname{pos\text{-}variation}(c, u)(\alpha) \in \mathbb{R}) \vee \\ \quad \quad \forall \alpha \in c.\; (\operatorname{neg\text{-}variation}(c, u)(\alpha) \in \mathbb{R}) \wedge \\ \quad \quad \forall a \in c. \\ \quad \quad \quad u(a) = \operatorname{ext\text{-}sub}(\operatorname{pos\text{-}variation}(c, u)(a), \operatorname{neg\text{-}variation}(c, u)(a)) \end{array}}}
\vnbprop{total-variation-is-measure}{\vnbstmt{\begin{array}{@{}l@{}} \forall a, b, u. \\ \quad \operatorname{is\text{-}signed\text{-}measure}(a, b, u) \Rightarrow \operatorname{is\text{-}measure}(a, b, \operatorname{total\text{-}variation}(b, u)) \end{array}}}
\vnbprop{sigma-algebra-member-subset}{\vnbstmt{\forall b, c, a.\; ((\operatorname{is\text{-}sigma\text{-}algebra}(b, c) \wedge (a \in c)) \Rightarrow (a \subseteq b))}}
\vnbprop{total-variation-value}{\vnbstmt{\begin{array}{@{}l@{}} \forall b, u.; \forall a \in b. \\ \quad (\operatorname{total\text{-}variation}(b, u)(a) \simeq \operatorname{eplus}(\operatorname{pos\text{-}variation}(b, u)(a), \operatorname{neg\text{-}variation}(b, u)(a))) \end{array}}}
\vnbprop{total-variation-norm-finite-iff}{\vnbstmt{\begin{array}{@{}l@{}} \forall a, b, u. \\ \quad \operatorname{is\text{-}signed\text{-}measure}(a, b, u) \Rightarrow  \\ \quad \quad \operatorname{total\text{-}variation\text{-}norm}(a, b, u) \in \mathbb{R} \\ \quad \quad \Leftrightarrow\; \forall \alpha \in b.\; (\operatorname{pos\text{-}variation}(b, u)(\alpha) \in \mathbb{R}) \wedge \\ \quad \quad \forall \alpha \in b.\; (\operatorname{neg\text{-}variation}(b, u)(\alpha) \in \mathbb{R}) \end{array}}}

\vnbbookfile{theorem-library/complex-measure-laws}{theorem-library/complex-measure-laws}
\noindent{\small complex-measure-laws.scm -- WAVE 2 (2026-10-27..30), wired 2026-10-30 from scratchpad/mf/mf-c1.scm, scratchpad/mf/mf-c2.scm, scratchpad/mf/mf-c3.scm, scratchpad/mf/mf-c4.scm. Library: measure.}\par
\vnbprop{real-nonneg-measurable}{\vnbstmt{\begin{array}{@{}l@{}} \forall a, b, f. \\ \quad (\operatorname{is\text{-}real\text{-}measurable\text{-}fn}(a, b, f) \wedge \forall x \in a.\; (0 \leq f(x))) \Rightarrow  \\ \quad \quad \operatorname{is\text{-}measurable\text{-}fn}(a, b, f) \end{array}}}
\vnbprop{sqrt-real-measurable}{\vnbstmt{\begin{array}{@{}l@{}} \forall a, b, g. \\ \quad (\operatorname{is\text{-}sigma\text{-}algebra}(a, b) \wedge (\operatorname{is\text{-}real\text{-}measurable\text{-}fn}(a, b, g) \wedge \forall x \in a.\; (0 \leq g(x)))) \Rightarrow  \\ \quad \quad \operatorname{is\text{-}real\text{-}measurable\text{-}fn}(a, b, (\lambda \mathit{sqx} \in a.\; \sqrt{g(\mathit{sqx})})) \end{array}}}
\vnbprop{alg-polar-unit}{\vnbstmt{\begin{array}{@{}l@{}} \forall a, b, p, q. \\ \quad ((a \in \mathbb{R}) \wedge ((b \in \mathbb{R}) \wedge ((p \in \mathbb{R}) \wedge ((q \in \mathbb{R}) \wedge (((p \cdot p) = ((a \cdot a) + (b \cdot b))) \wedge ((p \cdot q) = 1)))))) \Rightarrow  \\ \quad \quad (((a \cdot q) \cdot (a \cdot q)) + ((b \cdot q) \cdot (b \cdot q))) = 1 \end{array}}}
\vnbprop{polar-pointwise}{\vnbstmt{\begin{array}{@{}l@{}} \forall a, b, p, q, r, i. \\ \quad a \in \mathbb{R} \wedge \\ \quad b \in \mathbb{R} \wedge \\ \quad p = \sqrt{((a \cdot a) + (b \cdot b))} \wedge \\ \quad q = \left(\text{if }(0 < p)\text{ then }\operatorname{recip}(p)\text{ else }0\right) \wedge \\ \quad r = ((a \cdot q) + (1 - (p \cdot q))) \wedge \\ \quad i = (b \cdot q) \Rightarrow \\ \quad \quad p \in \mathbb{R} \wedge \\ \quad \quad 0 \leq p \wedge \\ \quad \quad (p \cdot p) = ((a \cdot a) + (b \cdot b)) \wedge \\ \quad \quad q \in \mathbb{R} \wedge \\ \quad \quad r \in \mathbb{R} \wedge \\ \quad \quad i \in \mathbb{R} \wedge \\ \quad \quad (p \cdot r) = a \wedge \\ \quad \quad (p \cdot i) = b \wedge \\ \quad \quad ((r \cdot r) + (i \cdot i)) = 1 \end{array}}}
\vnbprop{polar-functions-exist}{\vnbstmt{\begin{array}{@{}l@{}} \forall a, b, e, m. \\ \quad (\operatorname{is\text{-}sigma\text{-}algebra}(a, b) \wedge (\operatorname{is\text{-}real\text{-}measurable\text{-}fn}(a, b, e) \wedge \operatorname{is\text{-}real\text{-}measurable\text{-}fn}(a, b, m))) \Rightarrow  \\ \quad \quad \exists p, u, \mathit{iu}. \\ \quad \quad \quad \operatorname{is\text{-}measurable\text{-}fn}(a, b, p) \wedge \\ \quad \quad \quad \operatorname{is\text{-}real\text{-}measurable\text{-}fn}(a, b, p) \wedge \\ \quad \quad \quad \operatorname{is\text{-}real\text{-}measurable\text{-}fn}(a, b, u) \wedge \\ \quad \quad \quad \operatorname{is\text{-}real\text{-}measurable\text{-}fn}(a, b, \mathit{iu}) \wedge \\ \quad \quad \quad \forall x \in a.\; (0 \leq p(x)) \wedge \\ \quad \quad \quad \forall x \in a.\; ((p(x) \cdot p(x)) = ((e(x) \cdot e(x)) + (m(x) \cdot m(x)))) \wedge \\ \quad \quad \quad \forall x \in a.\; ((p(x) \cdot u(x)) = e(x)) \wedge \\ \quad \quad \quad \forall x \in a.\; ((p(x) \cdot \operatorname{iu}(x)) = m(x)) \wedge \\ \quad \quad \quad \forall x \in a.\; (((u(x) \cdot u(x)) + (\operatorname{iu}(x) \cdot \operatorname{iu}(x))) = 1) \end{array}}}
\vnbprop{ptwise-pos-part-scale}{\vnbstmt{\begin{array}{@{}l@{}} \forall p, g. \\ \quad ((p \in \mathbb{R}) \wedge ((g \in \mathbb{R}) \wedge (0 \leq p))) \Rightarrow  \\ \quad \quad \operatorname{etimes}(p, \left(\text{if }(0 \leq g)\text{ then }g\text{ else }0\right)) = \left(\text{if }(0 \leq (p \cdot g))\text{ then }(p \cdot g)\text{ else }0\right) \end{array}}}
\vnbprop{ptwise-neg-part-scale}{\vnbstmt{\begin{array}{@{}l@{}} \forall p, g. \\ \quad ((p \in \mathbb{R}) \wedge ((g \in \mathbb{R}) \wedge (0 \leq p))) \Rightarrow  \\ \quad \quad \operatorname{etimes}(p, \left(\text{if }(0 \leq g)\text{ then }0\text{ else }-g\right)) = \left(\text{if }(0 \leq (p \cdot g))\text{ then }0\text{ else }-(p \cdot g)\right) \end{array}}}
\vnbprop{rintegral-congruence}{\vnbstmt{\begin{array}{@{}l@{}} \forall a, b, u, f, g. \\ \quad ((f \in (a \to \mathbb{R})) \wedge ((g \in (a \to \mathbb{R})) \wedge \forall x \in a.\; (f(x) = g(x)))) \Rightarrow  \\ \quad \quad (\operatorname{rintegral}(a, b, u, f) \simeq \operatorname{rintegral}(a, b, u, g)) \end{array}}}
\vnbprop{polar-bound}{\vnbstmt{\begin{array}{@{}l@{}} \forall a, b, p, c, d, f, g. \\ \quad c \in \mathbb{R} \wedge \\ \quad d \in \mathbb{R} \wedge \\ \quad f \in \mathbb{R} \wedge \\ \quad g \in \mathbb{R} \wedge \\ \quad 0 \leq c \wedge \\ \quad 0 \leq d \wedge \\ \quad 0 \leq f \wedge \\ \quad 0 \leq g \wedge \\ \quad a \in \mathbb{R} \wedge \\ \quad b \in \mathbb{R} \wedge \\ \quad a = (c - d) \wedge \\ \quad b = (f - g) \wedge \\ \quad p \in \mathbb{R} \wedge \\ \quad 0 \leq p \wedge \\ \quad (p \cdot p) = ((a \cdot a) + (b \cdot b)) \Rightarrow \\ \quad \quad p \leq ((c + d) + (f + g)) \end{array}}}
\vnbprop{complex-zero-parts}{\vnbstmt{\forall a, b.\; (((a \in \mathbb{R}) \wedge ((b \in \mathbb{R}) \wedge ((a + (+i \cdot b)) = 0))) \Rightarrow ((a = 0) \wedge (b = 0)))}}
\vnbprop{unit-parts-bound}{\vnbstmt{\begin{array}{@{}l@{}} \forall z \in \mathbb{C}. \\ \quad (\operatorname{magnitude}(z) = 1) \Rightarrow ((\lvert \operatorname{real\text{-}part}(z) \rvert \leq 1) \wedge (\lvert \operatorname{imag\text{-}part}(z) \rvert \leq 1)) \end{array}}}
\vnbprop{abs-mul-le-unit}{\vnbstmt{\forall h, q.\; (((h \in \mathbb{R}) \wedge ((q \in \mathbb{R}) \wedge (\lvert q \rvert \leq 1))) \Rightarrow (\lvert (h \cdot q) \rvert \leq \lvert h \rvert))}}
\vnbprop{complex-integral-unfold}{\vnbstmt{\begin{array}{@{}l@{}} \forall a, b, c, u. \\ \quad (\operatorname{complex\text{-}integral}(a, b, c, u) \simeq (\operatorname{rintegral}(a, b, c, (\lambda x \in a.\; \operatorname{real\text{-}part}(u(x)))) + (+i \cdot \operatorname{rintegral}(a, b, c, (\lambda x \in a.\; \operatorname{imag\text{-}part}(u(x))))))) \end{array}}}
\vnbprop{complex-parts-equal}{\vnbstmt{\begin{array}{@{}l@{}} \forall a, b, c, d. \\ \quad ((a \in \mathbb{R}) \wedge ((b \in \mathbb{R}) \wedge ((c \in \mathbb{R}) \wedge ((d \in \mathbb{R}) \wedge ((a + (+i \cdot b)) = (c + (+i \cdot d))))))) \Rightarrow  \\ \quad \quad (a = c) \wedge (b = d) \end{array}}}

\vnbbookfile{theorem-library/l2-bridge}{theorem-library/l2-bridge}
\noindent{\small l2-bridge.scm -- wave of 2026-11-09/10 (agent TH3-L2B), wired 2026-11-10 from scratchpad/th3-1109/l2b/l2b-a.scm. Library: measure.}\par
\vnbprop{rps-not-pos-zero}{\vnbstmt{\forall y \in \operatorname{rr\text{-}pos\text{-}star}.\; (\neg (0 < y) \Rightarrow (y = 0))}}
\vnbprop{l2-square-integral}{\vnbstmt{\begin{array}{@{}l@{}} \forall a, b, u, f. \\ \quad (\operatorname{is\text{-}measure}(a, b, u) \wedge (f \in \operatorname{l2\text{-}set}(a, b, u))) \Rightarrow  \\ \quad \quad \operatorname{is\text{-}measurable\text{-}fn}(a, b, (\lambda x \in a.\; (f(x) \cdot f(x)))) \wedge \\ \quad \quad \operatorname{integral}(a, b, u, (\lambda x \in a.\; (f(x) \cdot f(x)))) \in \mathbb{R} \wedge \\ \quad \quad \operatorname{rintegral}(a, b, u, (\lambda x \in a.\; (f(x) \cdot f(x)))) = \operatorname{integral}(a, b, u, (\lambda x \in a.\; (f(x) \cdot f(x)))) \end{array}}}
\vnbprop{l2-in-of-square-integral}{\vnbstmt{\begin{array}{@{}l@{}} \forall a, b, u, f. \\ \quad \operatorname{is\text{-}measure}(a, b, u) \wedge \\ \quad \operatorname{is\text{-}real\text{-}measurable\text{-}fn}(a, b, f) \wedge \\ \quad \operatorname{integral}(a, b, u, (\lambda x \in a.\; (f(x) \cdot f(x)))) \in \mathbb{R} \Rightarrow \\ \quad \quad f \in \operatorname{l2\text{-}set}(a, b, u) \end{array}}}
\vnbprop{null-set-sections}{\vnbstmt{\begin{array}{@{}l@{}} \forall s, a, u, b, c, d, f. \\ \quad \operatorname{is\text{-}measure}(s, a, u) \wedge \\ \quad \operatorname{is\text{-}measure}(b, c, d) \wedge \\ \quad \operatorname{is\text{-}sigma\text{-}finite}(s, a, u) \wedge \\ \quad \operatorname{is\text{-}sigma\text{-}finite}(b, c, d) \wedge \\ \quad f \in \operatorname{product\text{-}sigma}(s, a, b, c) \wedge \\ \quad \operatorname{product\text{-}measure}(s, a, u, b, c, d)(f) = 0 \Rightarrow \\ \quad \quad \{x \in s : (0 < \operatorname{x\text{-}section\text{-}measure}(s, b, d, f)(x))\} \in a \wedge \\ \quad \quad u(\{x \in s : (0 < \operatorname{x\text{-}section\text{-}measure}(s, b, d, f)(x))\}) = 0 \end{array}}}
\vnbprop{null-set-sections-ae}{\vnbstmt{\begin{array}{@{}l@{}} \forall s, a, u, b, c, d, f. \\ \quad \operatorname{is\text{-}measure}(s, a, u) \wedge \\ \quad \operatorname{is\text{-}measure}(b, c, d) \wedge \\ \quad \operatorname{is\text{-}sigma\text{-}finite}(s, a, u) \wedge \\ \quad \operatorname{is\text{-}sigma\text{-}finite}(b, c, d) \wedge \\ \quad \operatorname{is\text{-}null\text{-}set}((s \times b), \operatorname{product\text{-}sigma}(s, a, b, c), \operatorname{product\text{-}measure}(s, a, u, b, c, d), f) \Rightarrow \\ \quad \quad \exists n. \\ \quad \quad \quad \operatorname{is\text{-}null\text{-}set}(s, a, u, n) \wedge \\ \quad \quad \quad \forall x \in s.\; (\neg (x \in n) \Rightarrow \operatorname{is\text{-}null\text{-}set}(b, c, d, \operatorname{x\text{-}section}(b, f, x))) \end{array}}}

\vnbbookfile{theorem-library/signed-fubini}{theorem-library/signed-fubini}
\noindent{\small signed-fubini.scm -- wave of 2026-11-09/10 (agent TH3-L2B), wired 2026-11-10 from scratchpad/th3-1109/l2b/l2b-b.scm. Library: measure.}\par
\vnbprop{integral-finite-truncate}{\vnbstmt{\begin{array}{@{}l@{}} \forall a, b, c, u. \\ \quad (\operatorname{is\text{-}measure}(a, b, c) \wedge (\operatorname{is\text{-}measurable\text{-}fn}(a, b, u) \wedge (\operatorname{integral}(a, b, c, u) \in \mathbb{R}))) \Rightarrow  \\ \quad \quad \operatorname{is\text{-}measurable\text{-}fn}(a, b, (\lambda \mathit{tdx} \in a.\; \left(\text{if }(u(\mathit{tdx}) \in \mathbb{R})\text{ then }u(\mathit{tdx})\text{ else }0\right))) \wedge \\ \quad \quad \forall x \in a.\; ((\lambda \mathit{tdx} \in a.\; \left(\text{if }(u(\mathit{tdx}) \in \mathbb{R})\text{ then }u(\mathit{tdx})\text{ else }0\right))(x) \in \mathbb{R}) \wedge \\ \quad \quad \operatorname{integral}(a, b, c, (\lambda \mathit{tdx} \in a.\; \left(\text{if }(u(\mathit{tdx}) \in \mathbb{R})\text{ then }u(\mathit{tdx})\text{ else }0\right))) = \operatorname{integral}(a, b, c, u) \wedge \\ \quad \quad \operatorname{is\text{-}null\text{-}set}(a, b, c, \{d \in a : (u(d) = \operatorname{pos\text{-}inf})\}) \wedge \\ \quad \quad \forall x \in a. \\ \quad \quad \quad \neg (x \in \{d \in a : (u(d) = \operatorname{pos\text{-}inf})\}) \Rightarrow  \\ \quad \quad \quad \quad (u(x) \in \mathbb{R}) \wedge ((\lambda \mathit{tdx} \in a.\; \left(\text{if }(u(\mathit{tdx}) \in \mathbb{R})\text{ then }u(\mathit{tdx})\text{ else }0\right))(x) = u(x)) \end{array}}}
\vnbprop{section-parts}{\vnbstmt{\begin{array}{@{}l@{}} \forall s, a, b, c, u, f, x. \\ \quad \operatorname{is\text{-}sigma\text{-}algebra}(s, a) \wedge \\ \quad \operatorname{is\text{-}measure}(b, c, u) \wedge \\ \quad \operatorname{is\text{-}real\text{-}measurable\text{-}fn}((s \times b), \operatorname{product\text{-}sigma}(s, a, b, c), f) \wedge \\ \quad x \in s \Rightarrow \\ \quad \quad \operatorname{is\text{-}real\text{-}measurable\text{-}fn}(b, c, \operatorname{fn\text{-}x\text{-}section}(b, f, x)) \wedge \\ \quad \quad \operatorname{integral}(b, c, u, \operatorname{pos\text{-}part}(b, \operatorname{fn\text{-}x\text{-}section}(b, f, x))) \\ \quad \quad =\; \operatorname{x\text{-}section\text{-}integral}(s, b, c, u, \operatorname{pos\text{-}part}((s \times b), f))(x) \wedge \\ \quad \quad \operatorname{integral}(b, c, u, \operatorname{neg\text{-}part}(b, \operatorname{fn\text{-}x\text{-}section}(b, f, x))) \\ \quad \quad =\; \operatorname{x\text{-}section\text{-}integral}(s, b, c, u, \operatorname{neg\text{-}part}((s \times b), f))(x) \end{array}}}
\vnbprop{signed-fubini}{\vnbstmt{\begin{array}{@{}l@{}} \forall s, a, u, b, c, d, f. \\ \quad \operatorname{is\text{-}measure}(s, a, u) \wedge \\ \quad \operatorname{is\text{-}measure}(b, c, d) \wedge \\ \quad \operatorname{is\text{-}sigma\text{-}finite}(s, a, u) \wedge \\ \quad \operatorname{is\text{-}sigma\text{-}finite}(b, c, d) \wedge \\ \quad \operatorname{is\text{-}integrable}((s \times b), \operatorname{product\text{-}sigma}(s, a, b, c), \operatorname{product\text{-}measure}(s, a, u, b, c, d), f) \Rightarrow \\ \quad \quad \operatorname{is\text{-}integrable}(s, a, u, \\ \quad \quad \quad (\lambda \mathit{sgx} \in s.\; \\ \quad \quad \quad \quad (\left(\text{if }(\operatorname{x\text{-}section\text{-}integral}(s, b, c, d, \operatorname{pos\text{-}part}((s \times b), f))(\mathit{sgx}) \in \mathbb{R})\text{ then }\operatorname{x\text{-}section\text{-}integral}(s, b, c, d, \operatorname{pos\text{-}part}((s \times b), f))(\mathit{sgx})\text{ else }0\right) - \left(\text{if }(\operatorname{x\text{-}section\text{-}integral}(s, b, c, d, \operatorname{neg\text{-}part}((s \times b), f))(\mathit{sgx}) \in \mathbb{R})\text{ then }\operatorname{x\text{-}section\text{-}integral}(s, b, c, d, \operatorname{neg\text{-}part}((s \times b), f))(\mathit{sgx})\text{ else }0\right)))) \wedge \\ \quad \quad \operatorname{rintegral}(s, a, u, \\ \quad \quad \quad (\lambda \mathit{sgx} \in s.\; \\ \quad \quad \quad \quad (\left(\text{if }(\operatorname{x\text{-}section\text{-}integral}(s, b, c, d, \operatorname{pos\text{-}part}((s \times b), f))(\mathit{sgx}) \in \mathbb{R})\text{ then }\operatorname{x\text{-}section\text{-}integral}(s, b, c, d, \operatorname{pos\text{-}part}((s \times b), f))(\mathit{sgx})\text{ else }0\right) - \left(\text{if }(\operatorname{x\text{-}section\text{-}integral}(s, b, c, d, \operatorname{neg\text{-}part}((s \times b), f))(\mathit{sgx}) \in \mathbb{R})\text{ then }\operatorname{x\text{-}section\text{-}integral}(s, b, c, d, \operatorname{neg\text{-}part}((s \times b), f))(\mathit{sgx})\text{ else }0\right)))) \\ \quad \quad =\; \operatorname{rintegral}((s \times b), \operatorname{product\text{-}sigma}(s, a, b, c), \operatorname{product\text{-}measure}(s, a, u, b, c, d), f) \wedge \\ \quad \quad \exists n. \\ \quad \quad \quad \operatorname{is\text{-}null\text{-}set}(s, a, u, n) \wedge \\ \quad \quad \quad \forall x \in s. \\ \quad \quad \quad \quad \neg (x \in n) \Rightarrow  \\ \quad \quad \quad \quad \quad \operatorname{is\text{-}integrable}(b, c, d, \operatorname{fn\text{-}x\text{-}section}(b, f, x)) \wedge \\ \quad \quad \quad \quad \quad (\lambda \mathit{sgx} \in s.\; (\left(\text{if }(\operatorname{x\text{-}section\text{-}integral}(s, b, c, d, \operatorname{pos\text{-}part}((s \times b), f))(\mathit{sgx}) \in \mathbb{R})\text{ then }\operatorname{x\text{-}section\text{-}integral}(s, b, c, d, \operatorname{pos\text{-}part}((s \times b), f))(\mathit{sgx})\text{ else }0\right) - \left(\text{if }(\operatorname{x\text{-}section\text{-}integral}(s, b, c, d, \operatorname{neg\text{-}part}((s \times b), f))(\mathit{sgx}) \in \mathbb{R})\text{ then }\operatorname{x\text{-}section\text{-}integral}(s, b, c, d, \operatorname{neg\text{-}part}((s \times b), f))(\mathit{sgx})\text{ else }0\right)))(x) \\ \quad \quad \quad \quad \quad =\; \operatorname{rintegral}(b, c, d, \operatorname{fn\text{-}x\text{-}section}(b, f, x)) \end{array}}}

\vnbbookfile{theorem-library/l2-kernel-sections}{theorem-library/l2-kernel-sections}
\noindent{\small l2-kernel-sections.scm -- wave of 2026-11-09/10 (agent TH3-L2B), wired 2026-11-10 from scratchpad/th3-1109/l2b/l2b-c.scm. Library: measure.}\par
\vnbprop{l2-sections-ae}{\vnbstmt{\begin{array}{@{}l@{}} \forall s, a, u, b, c, d, k. \\ \quad \operatorname{is\text{-}measure}(s, a, u) \wedge \\ \quad \operatorname{is\text{-}measure}(b, c, d) \wedge \\ \quad \operatorname{is\text{-}sigma\text{-}finite}(s, a, u) \wedge \\ \quad \operatorname{is\text{-}sigma\text{-}finite}(b, c, d) \wedge \\ \quad k \in \operatorname{l2\text{-}set}((s \times b), \operatorname{product\text{-}sigma}(s, a, b, c), \operatorname{product\text{-}measure}(s, a, u, b, c, d)) \Rightarrow \\ \quad \quad \operatorname{integral}(s, a, u, \operatorname{x\text{-}section\text{-}integral}(s, b, c, d, (\lambda \mathit{lx} \in (s \times b).\; (k(\mathit{lx}) \cdot k(\mathit{lx}))))) \\ \quad \quad =\; \operatorname{integral}((s \times b), \operatorname{product\text{-}sigma}(s, a, b, c), \\ \quad \quad \quad \operatorname{product\text{-}measure}(s, a, u, b, c, d), \\ \quad \quad \quad (\lambda \mathit{lx} \in (s \times b).\; (k(\mathit{lx}) \cdot k(\mathit{lx})))) \wedge \\ \quad \quad \exists n. \\ \quad \quad \quad \operatorname{is\text{-}null\text{-}set}(s, a, u, n) \wedge \\ \quad \quad \quad \forall x \in s.\; (\neg (x \in n) \Rightarrow (\operatorname{fn\text{-}x\text{-}section}(b, k, x) \in \operatorname{l2\text{-}set}(b, c, d))) \end{array}}}
\vnbprop{ae-eq-sections}{\vnbstmt{\begin{array}{@{}l@{}} \forall s, a, u, b, c, d, k, f, g. \\ \quad \operatorname{is\text{-}measure}(s, a, u) \wedge \\ \quad \operatorname{is\text{-}measure}(b, c, d) \wedge \\ \quad \operatorname{is\text{-}sigma\text{-}finite}(s, a, u) \wedge \\ \quad \operatorname{is\text{-}sigma\text{-}finite}(b, c, d) \wedge \\ \quad k \in ((s \times b) \to \mathbb{R}) \wedge \\ \quad f \in ((s \times b) \to \mathbb{R}) \wedge \\ \quad \operatorname{is\text{-}null\text{-}set}((s \times b), \operatorname{product\text{-}sigma}(s, a, b, c), \operatorname{product\text{-}measure}(s, a, u, b, c, d), g) \wedge \\ \quad \forall w \in (s \times b).\; (\neg (w \in g) \Rightarrow (k(w) = f(w))) \Rightarrow \\ \quad \quad \exists n. \\ \quad \quad \quad \operatorname{is\text{-}null\text{-}set}(s, a, u, n) \wedge \\ \quad \quad \quad \forall x \in s. \\ \quad \quad \quad \quad \neg (x \in n) \Rightarrow  \\ \quad \quad \quad \quad \quad \exists \mathit{nb}. \\ \quad \quad \quad \quad \quad \quad \operatorname{is\text{-}null\text{-}set}(b, c, d, \mathit{nb}) \wedge \\ \quad \quad \quad \quad \quad \quad \forall y \in b. \\ \quad \quad \quad \quad \quad \quad \quad \neg (y \in \mathit{nb}) \Rightarrow (\operatorname{fn\text{-}x\text{-}section}(b, k, x)(y) = \operatorname{fn\text{-}x\text{-}section}(b, f, x)(y)) \end{array}}}
\vnbprop{kernel-integral-well-defined}{\vnbstmt{\begin{array}{@{}l@{}} \forall s, a, u, b, c, d, k, f, g, i, h, m. \\ \quad \operatorname{is\text{-}measure}(s, a, u) \wedge \\ \quad \operatorname{is\text{-}measure}(b, c, d) \wedge \\ \quad \operatorname{is\text{-}sigma\text{-}finite}(s, a, u) \wedge \\ \quad \operatorname{is\text{-}sigma\text{-}finite}(b, c, d) \wedge \\ \quad k \in \operatorname{l2\text{-}set}((s \times b), \operatorname{product\text{-}sigma}(s, a, b, c), \operatorname{product\text{-}measure}(s, a, u, b, c, d)) \wedge \\ \quad f \in \operatorname{l2\text{-}set}((s \times b), \operatorname{product\text{-}sigma}(s, a, b, c), \operatorname{product\text{-}measure}(s, a, u, b, c, d)) \wedge \\ \quad \operatorname{is\text{-}null\text{-}set}((s \times b), \operatorname{product\text{-}sigma}(s, a, b, c), \operatorname{product\text{-}measure}(s, a, u, b, c, d), g) \wedge \\ \quad \forall w \in (s \times b).\; (\neg (w \in g) \Rightarrow (k(w) = f(w))) \wedge \\ \quad i \in \operatorname{l2\text{-}set}(b, c, d) \wedge \\ \quad h \in \operatorname{l2\text{-}set}(b, c, d) \wedge \\ \quad \operatorname{is\text{-}null\text{-}set}(b, c, d, m) \wedge \\ \quad \forall y \in b.\; (\neg (y \in m) \Rightarrow (i(y) = h(y))) \Rightarrow \\ \quad \quad \exists n. \\ \quad \quad \quad \operatorname{is\text{-}null\text{-}set}(s, a, u, n) \wedge \\ \quad \quad \quad \forall x \in s. \\ \quad \quad \quad \quad \neg (x \in n) \Rightarrow  \\ \quad \quad \quad \quad \quad \operatorname{is\text{-}integrable}(b, c, d, (\lambda \mathit{lx} \in b.\; (\operatorname{fn\text{-}x\text{-}section}(b, k, x)(\mathit{lx}) \cdot i(\mathit{lx})))) \wedge \\ \quad \quad \quad \quad \quad \operatorname{is\text{-}integrable}(b, c, d, (\lambda \mathit{lx} \in b.\; (\operatorname{fn\text{-}x\text{-}section}(b, f, x)(\mathit{lx}) \cdot h(\mathit{lx})))) \wedge \\ \quad \quad \quad \quad \quad \operatorname{rintegral}(b, c, d, (\lambda \mathit{lx} \in b.\; (\operatorname{fn\text{-}x\text{-}section}(b, k, x)(\mathit{lx}) \cdot i(\mathit{lx})))) \\ \quad \quad \quad \quad \quad =\; \operatorname{rintegral}(b, c, d, (\lambda \mathit{lx} \in b.\; (\operatorname{fn\text{-}x\text{-}section}(b, f, x)(\mathit{lx}) \cdot h(\mathit{lx})))) \end{array}}}

\vnbbookfile{theorem-library/l2-kernel-measurable}{theorem-library/l2-kernel-measurable}
\noindent{\small l2-kernel-measurable.scm -- wave of 2026-11-09/10 (agent TH3-L2B), wired 2026-11-10 from scratchpad/th3-1109/l2b/l2b-d.scm. Library: measure.}\par
\vnbprop{real-measurable-pullback-second}{\vnbstmt{\begin{array}{@{}l@{}} \forall s, a, b, c, i. \\ \quad (\operatorname{is\text{-}sigma\text{-}algebra}(s, a) \wedge (\operatorname{is\text{-}sigma\text{-}algebra}(b, c) \wedge \operatorname{is\text{-}real\text{-}measurable\text{-}fn}(b, c, i))) \Rightarrow  \\ \quad \quad \operatorname{is\text{-}real\text{-}measurable\text{-}fn}((s \times b), \operatorname{product\text{-}sigma}(s, a, b, c), \\ \quad \quad \quad (\lambda w \in (s \times b).\; i(\operatorname{nth}(2, w)))) \end{array}}}
\vnbprop{truncate-measurable}{\vnbstmt{\begin{array}{@{}l@{}} \forall a, b, u. \\ \quad (\operatorname{is\text{-}sigma\text{-}algebra}(a, b) \wedge \operatorname{is\text{-}measurable\text{-}fn}(a, b, u)) \Rightarrow  \\ \quad \quad \operatorname{is\text{-}measurable\text{-}fn}(a, b, (\lambda \mathit{tdx} \in a.\; \left(\text{if }(u(\mathit{tdx}) \in \mathbb{R})\text{ then }u(\mathit{tdx})\text{ else }0\right))) \wedge \\ \quad \quad \forall x \in a.\; ((\lambda \mathit{tdx} \in a.\; \left(\text{if }(u(\mathit{tdx}) \in \mathbb{R})\text{ then }u(\mathit{tdx})\text{ else }0\right))(x) \in \mathbb{R}) \end{array}}}
\vnbprop{section-integral-real-measurable}{\vnbstmt{\begin{array}{@{}l@{}} \forall s, a, u, b, c, d, f. \\ \quad \operatorname{is\text{-}measure}(s, a, u) \wedge \\ \quad \operatorname{is\text{-}measure}(b, c, d) \wedge \\ \quad \operatorname{is\text{-}sigma\text{-}finite}(s, a, u) \wedge \\ \quad \operatorname{is\text{-}sigma\text{-}finite}(b, c, d) \wedge \\ \quad \operatorname{is\text{-}real\text{-}measurable\text{-}fn}((s \times b), \operatorname{product\text{-}sigma}(s, a, b, c), f) \Rightarrow \\ \quad \quad \operatorname{is\text{-}real\text{-}measurable\text{-}fn}(s, a, \\ \quad \quad \quad (\lambda \mathit{sgx} \in s.\; \\ \quad \quad \quad \quad (\left(\text{if }(\operatorname{x\text{-}section\text{-}integral}(s, b, c, d, \operatorname{pos\text{-}part}((s \times b), f))(\mathit{sgx}) \in \mathbb{R})\text{ then }\operatorname{x\text{-}section\text{-}integral}(s, b, c, d, \operatorname{pos\text{-}part}((s \times b), f))(\mathit{sgx})\text{ else }0\right) - \left(\text{if }(\operatorname{x\text{-}section\text{-}integral}(s, b, c, d, \operatorname{neg\text{-}part}((s \times b), f))(\mathit{sgx}) \in \mathbb{R})\text{ then }\operatorname{x\text{-}section\text{-}integral}(s, b, c, d, \operatorname{neg\text{-}part}((s \times b), f))(\mathit{sgx})\text{ else }0\right)))) \wedge \\ \quad \quad \forall x \in s. \\ \quad \quad \quad \operatorname{is\text{-}integrable}(b, c, d, \operatorname{fn\text{-}x\text{-}section}(b, f, x)) \Rightarrow  \\ \quad \quad \quad \quad (\lambda \mathit{sgx} \in s.\; (\left(\text{if }(\operatorname{x\text{-}section\text{-}integral}(s, b, c, d, \operatorname{pos\text{-}part}((s \times b), f))(\mathit{sgx}) \in \mathbb{R})\text{ then }\operatorname{x\text{-}section\text{-}integral}(s, b, c, d, \operatorname{pos\text{-}part}((s \times b), f))(\mathit{sgx})\text{ else }0\right) - \left(\text{if }(\operatorname{x\text{-}section\text{-}integral}(s, b, c, d, \operatorname{neg\text{-}part}((s \times b), f))(\mathit{sgx}) \in \mathbb{R})\text{ then }\operatorname{x\text{-}section\text{-}integral}(s, b, c, d, \operatorname{neg\text{-}part}((s \times b), f))(\mathit{sgx})\text{ else }0\right)))(x) \\ \quad \quad \quad \quad =\; \operatorname{rintegral}(b, c, d, \operatorname{fn\text{-}x\text{-}section}(b, f, x)) \end{array}}}
\vnbprop{l2-schwarz}{\vnbstmt{\begin{array}{@{}l@{}} \forall a, b, u, f, g. \\ \quad (\operatorname{is\text{-}measure}(a, b, u) \wedge ((f \in \operatorname{l2\text{-}set}(a, b, u)) \wedge (g \in \operatorname{l2\text{-}set}(a, b, u)))) \Rightarrow  \\ \quad \quad \operatorname{rintegral}(a, b, u, (\lambda x \in a.\; (f(x) \cdot g(x)))) \cdot \operatorname{rintegral}(a, b, u, (\lambda x \in a.\; (f(x) \cdot g(x)))) \\ \quad \quad \leq\; \operatorname{rintegral}(a, b, u, (\lambda x \in a.\; (f(x) \cdot f(x)))) \cdot \operatorname{rintegral}(a, b, u, (\lambda x \in a.\; (g(x) \cdot g(x)))) \end{array}}}

\vnbbookfile{theorem-library/l2-kernel-operator}{theorem-library/l2-kernel-operator}
\noindent{\small l2-kernel-operator.scm -- wave of 2026-11-09/10 (agent TH3-L2B), wired 2026-11-10 from scratchpad/th3-1109/l2b/l2b-e.scm. Library: measure.}\par
\vnbprop{l2-of-square-bound}{\vnbstmt{\begin{array}{@{}l@{}} \forall a, b, u, g, h, c, d. \\ \quad \operatorname{is\text{-}measure}(a, b, u) \wedge \\ \quad \operatorname{is\text{-}real\text{-}measurable\text{-}fn}(a, b, g) \wedge \\ \quad \operatorname{is\text{-}measurable\text{-}fn}(a, b, h) \wedge \\ \quad c \in \mathbb{R} \wedge \\ \quad 0 \leq c \wedge \\ \quad \operatorname{integral}(a, b, u, h) \in \mathbb{R} \wedge \\ \quad d \in b \wedge \\ \quad u(d) = 0 \wedge \\ \quad \forall x \in a.\; (\neg (x \in d) \Rightarrow (((g(x) \cdot g(x)) \leq (h(x) \cdot c)) \wedge (h(x) \in \mathbb{R}))) \Rightarrow \\ \quad \quad g \in \operatorname{l2\text{-}set}(a, b, u) \wedge \\ \quad \quad \operatorname{integral}(a, b, u, (\lambda \mathit{lx} \in a.\; (g(\mathit{lx}) \cdot g(\mathit{lx})))) \leq (\operatorname{integral}(a, b, u, h) \cdot c) \end{array}}}
\vnbprop{xsi-square-section}{\vnbstmt{\begin{array}{@{}l@{}} \forall s, a, b, c, u, k, x. \\ \quad \operatorname{is\text{-}sigma\text{-}algebra}(s, a) \wedge \\ \quad \operatorname{is\text{-}measure}(b, c, u) \wedge \\ \quad k \in ((s \times b) \to \mathbb{R}) \wedge \\ \quad \operatorname{is\text{-}measurable\text{-}fn}((s \times b), \operatorname{product\text{-}sigma}(s, a, b, c), \\ \quad \quad (\lambda \mathit{lx} \in (s \times b).\; (k(\mathit{lx}) \cdot k(\mathit{lx})))) \wedge \\ \quad x \in s \Rightarrow \\ \quad \quad \operatorname{x\text{-}section\text{-}integral}(s, b, c, u, (\lambda \mathit{lx} \in (s \times b).\; (k(\mathit{lx}) \cdot k(\mathit{lx}))))(x) \\ \quad \quad =\; \operatorname{integral}(b, c, u, (\lambda \mathit{lx} \in b.\; (\operatorname{fn\text{-}x\text{-}section}(b, k, x)(\mathit{lx}) \cdot \operatorname{fn\text{-}x\text{-}section}(b, k, x)(\mathit{lx})))) \end{array}}}
\vnbprop{kernel-operator-l2}{\vnbstmt{\begin{array}{@{}l@{}} \forall s, a, u, b, c, d, k, i. \\ \quad \operatorname{is\text{-}measure}(s, a, u) \wedge \\ \quad \operatorname{is\text{-}measure}(b, c, d) \wedge \\ \quad \operatorname{is\text{-}sigma\text{-}finite}(s, a, u) \wedge \\ \quad \operatorname{is\text{-}sigma\text{-}finite}(b, c, d) \wedge \\ \quad k \in \operatorname{l2\text{-}set}((s \times b), \operatorname{product\text{-}sigma}(s, a, b, c), \operatorname{product\text{-}measure}(s, a, u, b, c, d)) \wedge \\ \quad i \in \operatorname{l2\text{-}set}(b, c, d) \Rightarrow \\ \quad \quad (\lambda \mathit{sgx} \in s.\; \\ \quad \quad \quad (\left(\text{if }(\operatorname{x\text{-}section\text{-}integral}(s, b, c, d, \operatorname{pos\text{-}part}((s \times b), (\lambda w \in (s \times b).\; (k(w) \cdot i(\operatorname{nth}(2, w))))))(\mathit{sgx}) \in \mathbb{R})\text{ then }\operatorname{x\text{-}section\text{-}integral}(s, b, c, d, \operatorname{pos\text{-}part}((s \times b), (\lambda w \in (s \times b).\; (k(w) \cdot i(\operatorname{nth}(2, w))))))(\mathit{sgx})\text{ else }0\right) - \left(\text{if }(\operatorname{x\text{-}section\text{-}integral}(s, b, c, d, \operatorname{neg\text{-}part}((s \times b), (\lambda w \in (s \times b).\; (k(w) \cdot i(\operatorname{nth}(2, w))))))(\mathit{sgx}) \in \mathbb{R})\text{ then }\operatorname{x\text{-}section\text{-}integral}(s, b, c, d, \operatorname{neg\text{-}part}((s \times b), (\lambda w \in (s \times b).\; (k(w) \cdot i(\operatorname{nth}(2, w))))))(\mathit{sgx})\text{ else }0\right))) \\ \quad \quad \in\; \operatorname{l2\text{-}set}(s, a, u) \wedge \\ \quad \quad \operatorname{integral}(s, a, u, \\ \quad \quad \quad (\lambda \mathit{lx} \in s.\; \\ \quad \quad \quad \quad (\lambda \mathit{sgx} \in s.\; (\left(\text{if }(\operatorname{x\text{-}section\text{-}integral}(s, b, c, d, \operatorname{pos\text{-}part}((s \times b), (\lambda w \in (s \times b).\; (k(w) \cdot i(\operatorname{nth}(2, w))))))(\mathit{sgx}) \in \mathbb{R})\text{ then }\operatorname{x\text{-}section\text{-}integral}(s, b, c, d, \operatorname{pos\text{-}part}((s \times b), (\lambda w \in (s \times b).\; (k(w) \cdot i(\operatorname{nth}(2, w))))))(\mathit{sgx})\text{ else }0\right) - \left(\text{if }(\operatorname{x\text{-}section\text{-}integral}(s, b, c, d, \operatorname{neg\text{-}part}((s \times b), (\lambda w \in (s \times b).\; (k(w) \cdot i(\operatorname{nth}(2, w))))))(\mathit{sgx}) \in \mathbb{R})\text{ then }\operatorname{x\text{-}section\text{-}integral}(s, b, c, d, \operatorname{neg\text{-}part}((s \times b), (\lambda w \in (s \times b).\; (k(w) \cdot i(\operatorname{nth}(2, w))))))(\mathit{sgx})\text{ else }0\right)))(\mathit{lx}) \cdot (\lambda \mathit{sgx} \in s.\; (\left(\text{if }(\operatorname{x\text{-}section\text{-}integral}(s, b, c, d, \operatorname{pos\text{-}part}((s \times b), (\lambda w \in (s \times b).\; (k(w) \cdot i(\operatorname{nth}(2, w))))))(\mathit{sgx}) \in \mathbb{R})\text{ then }\operatorname{x\text{-}section\text{-}integral}(s, b, c, d, \operatorname{pos\text{-}part}((s \times b), (\lambda w \in (s \times b).\; (k(w) \cdot i(\operatorname{nth}(2, w))))))(\mathit{sgx})\text{ else }0\right) - \left(\text{if }(\operatorname{x\text{-}section\text{-}integral}(s, b, c, d, \operatorname{neg\text{-}part}((s \times b), (\lambda w \in (s \times b).\; (k(w) \cdot i(\operatorname{nth}(2, w))))))(\mathit{sgx}) \in \mathbb{R})\text{ then }\operatorname{x\text{-}section\text{-}integral}(s, b, c, d, \operatorname{neg\text{-}part}((s \times b), (\lambda w \in (s \times b).\; (k(w) \cdot i(\operatorname{nth}(2, w))))))(\mathit{sgx})\text{ else }0\right)))(\mathit{lx}))) \\ \quad \quad \leq\; \operatorname{integral}((s \times b), \operatorname{product\text{-}sigma}(s, a, b, c), \operatorname{product\text{-}measure}(s, a, u, b, c, d), (\lambda \mathit{lx} \in (s \times b).\; (k(\mathit{lx}) \cdot k(\mathit{lx})))) \cdot \operatorname{integral}(b, c, d, (\lambda \mathit{lx} \in b.\; (i(\mathit{lx}) \cdot i(\mathit{lx})))) \wedge \\ \quad \quad \exists n. \\ \quad \quad \quad \operatorname{is\text{-}null\text{-}set}(s, a, u, n) \wedge \\ \quad \quad \quad \forall x \in s. \\ \quad \quad \quad \quad \neg (x \in n) \Rightarrow  \\ \quad \quad \quad \quad \quad \operatorname{is\text{-}integrable}(b, c, d, (\lambda \mathit{lx} \in b.\; (\operatorname{fn\text{-}x\text{-}section}(b, k, x)(\mathit{lx}) \cdot i(\mathit{lx})))) \wedge \\ \quad \quad \quad \quad \quad (\lambda \mathit{sgx} \in s.\; (\left(\text{if }(\operatorname{x\text{-}section\text{-}integral}(s, b, c, d, \operatorname{pos\text{-}part}((s \times b), (\lambda w \in (s \times b).\; (k(w) \cdot i(\operatorname{nth}(2, w))))))(\mathit{sgx}) \in \mathbb{R})\text{ then }\operatorname{x\text{-}section\text{-}integral}(s, b, c, d, \operatorname{pos\text{-}part}((s \times b), (\lambda w \in (s \times b).\; (k(w) \cdot i(\operatorname{nth}(2, w))))))(\mathit{sgx})\text{ else }0\right) - \left(\text{if }(\operatorname{x\text{-}section\text{-}integral}(s, b, c, d, \operatorname{neg\text{-}part}((s \times b), (\lambda w \in (s \times b).\; (k(w) \cdot i(\operatorname{nth}(2, w))))))(\mathit{sgx}) \in \mathbb{R})\text{ then }\operatorname{x\text{-}section\text{-}integral}(s, b, c, d, \operatorname{neg\text{-}part}((s \times b), (\lambda w \in (s \times b).\; (k(w) \cdot i(\operatorname{nth}(2, w))))))(\mathit{sgx})\text{ else }0\right)))(x) \\ \quad \quad \quad \quad \quad =\; \operatorname{rintegral}(b, c, d, (\lambda \mathit{lx} \in b.\; (\operatorname{fn\text{-}x\text{-}section}(b, k, x)(\mathit{lx}) \cdot i(\mathit{lx})))) \end{array}}}

\vnbbookfile{theorem-library/kernel-apply}{theorem-library/kernel-apply}
\noindent{\small kernel-apply.scm -- wave of 2026-11-09/10 (agent TH3-L2C), wired 2026-11-10 from scratchpad/th3-1109/l2c/kc-a.scm. Library: measure.}\par
\vnbprop{kernel-apply-unfold}{\vnbstmt{\begin{array}{@{}l@{}} \forall s, a, b, u, k, i. \\ \quad (\int {k}(\cdot, y)\, {i}(y)\, d{u}(y) \simeq (\lambda x \in s.\; (\left(\text{if }(\operatorname{x\text{-}section\text{-}integral}(s, a, b, u, \operatorname{pos\text{-}part}((s \times a), (\lambda w \in (s \times a).\; (k(w) \cdot i(\operatorname{nth}(2, w))))))(x) \in \mathbb{R})\text{ then }\operatorname{x\text{-}section\text{-}integral}(s, a, b, u, \operatorname{pos\text{-}part}((s \times a), (\lambda w \in (s \times a).\; (k(w) \cdot i(\operatorname{nth}(2, w))))))(x)\text{ else }0\right) - \left(\text{if }(\operatorname{x\text{-}section\text{-}integral}(s, a, b, u, \operatorname{neg\text{-}part}((s \times a), (\lambda w \in (s \times a).\; (k(w) \cdot i(\operatorname{nth}(2, w))))))(x) \in \mathbb{R})\text{ then }\operatorname{x\text{-}section\text{-}integral}(s, a, b, u, \operatorname{neg\text{-}part}((s \times a), (\lambda w \in (s \times a).\; (k(w) \cdot i(\operatorname{nth}(2, w))))))(x)\text{ else }0\right)))) \end{array}}}
\vnbprop{kernel-apply-spec}{\vnbstmt{\begin{array}{@{}l@{}} \forall s, a, u, b, c, d, k, i. \\ \quad \operatorname{is\text{-}measure}(s, a, u) \wedge \\ \quad \operatorname{is\text{-}measure}(b, c, d) \wedge \\ \quad \operatorname{is\text{-}sigma\text{-}finite}(s, a, u) \wedge \\ \quad \operatorname{is\text{-}sigma\text{-}finite}(b, c, d) \wedge \\ \quad k \in \operatorname{l2\text{-}set}((s \times b), \operatorname{product\text{-}sigma}(s, a, b, c), \operatorname{product\text{-}measure}(s, a, u, b, c, d)) \wedge \\ \quad i \in \operatorname{l2\text{-}set}(b, c, d) \Rightarrow \\ \quad \quad \int {k}(\cdot, y)\, {i}(y)\, d{d}(y) \in \operatorname{l2\text{-}set}(s, a, u) \wedge \\ \quad \quad \operatorname{integral}(s, a, u, \\ \quad \quad \quad (\lambda \mathit{lx} \in s.\; (\int {k}(\cdot, y)\, {i}(y)\, d{d}(y)(\mathit{lx}) \cdot \int {k}(\cdot, y)\, {i}(y)\, d{d}(y)(\mathit{lx})))) \\ \quad \quad \leq\; \operatorname{integral}((s \times b), \operatorname{product\text{-}sigma}(s, a, b, c), \operatorname{product\text{-}measure}(s, a, u, b, c, d), (\lambda \mathit{lx} \in (s \times b).\; (k(\mathit{lx}) \cdot k(\mathit{lx})))) \cdot \operatorname{integral}(b, c, d, (\lambda \mathit{lx} \in b.\; (i(\mathit{lx}) \cdot i(\mathit{lx})))) \wedge \\ \quad \quad \exists n. \\ \quad \quad \quad \operatorname{is\text{-}null\text{-}set}(s, a, u, n) \wedge \\ \quad \quad \quad \forall x \in s. \\ \quad \quad \quad \quad \neg (x \in n) \Rightarrow  \\ \quad \quad \quad \quad \quad \operatorname{is\text{-}integrable}(b, c, d, (\lambda \mathit{lx} \in b.\; (\operatorname{fn\text{-}x\text{-}section}(b, k, x)(\mathit{lx}) \cdot i(\mathit{lx})))) \wedge \\ \quad \quad \quad \quad \quad \int {k}(\cdot, y)\, {i}(y)\, d{d}(y)(x) \\ \quad \quad \quad \quad \quad =\; \operatorname{rintegral}(b, c, d, (\lambda \mathit{lx} \in b.\; (\operatorname{fn\text{-}x\text{-}section}(b, k, x)(\mathit{lx}) \cdot i(\mathit{lx})))) \end{array}}}
\vnbprop{kernel-apply-linear-ae}{\vnbstmt{\begin{array}{@{}l@{}} \forall s, c, d, f, g, m, k, u, w, h, a, b, n. \\ \quad \operatorname{is\text{-}measure}(s, c, d) \wedge \\ \quad \operatorname{is\text{-}measure}(f, g, m) \wedge \\ \quad \operatorname{is\text{-}sigma\text{-}finite}(s, c, d) \wedge \\ \quad \operatorname{is\text{-}sigma\text{-}finite}(f, g, m) \wedge \\ \quad k \in \operatorname{l2\text{-}set}((s \times f), \operatorname{product\text{-}sigma}(s, c, f, g), \operatorname{product\text{-}measure}(s, c, d, f, g, m)) \wedge \\ \quad u \in \operatorname{l2\text{-}set}(f, g, m) \wedge \\ \quad w \in \operatorname{l2\text{-}set}(f, g, m) \wedge \\ \quad h \in \operatorname{l2\text{-}set}(f, g, m) \wedge \\ \quad a \in \mathbb{R} \wedge \\ \quad b \in \mathbb{R} \wedge \\ \quad \operatorname{is\text{-}null\text{-}set}(f, g, m, n) \wedge \\ \quad \forall y \in f.\; (\neg (y \in n) \Rightarrow (h(y) = ((a \cdot u(y)) + (b \cdot w(y))))) \Rightarrow \\ \quad \quad \exists \mathit{en}. \\ \quad \quad \quad \operatorname{is\text{-}null\text{-}set}(s, c, d, \mathit{en}) \wedge \\ \quad \quad \quad \forall x \in s. \\ \quad \quad \quad \quad \neg (x \in \mathit{en}) \Rightarrow  \\ \quad \quad \quad \quad \quad \int {k}(\cdot, y)\, {h}(y)\, d{m}(y)(x) \\ \quad \quad \quad \quad \quad =\; (a \cdot \int {k}(\cdot, y)\, {u}(y)\, d{m}(y)(x)) + (b \cdot \int {k}(\cdot, y)\, {w}(y)\, d{m}(y)(x)) \end{array}}}
\vnbprop{kernel-apply-ae-eq}{\vnbstmt{\begin{array}{@{}l@{}} \forall s, a, b, c, d, f, k, u, h, g. \\ \quad \operatorname{is\text{-}measure}(s, a, b) \wedge \\ \quad \operatorname{is\text{-}measure}(c, d, f) \wedge \\ \quad \operatorname{is\text{-}sigma\text{-}finite}(s, a, b) \wedge \\ \quad \operatorname{is\text{-}sigma\text{-}finite}(c, d, f) \wedge \\ \quad k \in \operatorname{l2\text{-}set}((s \times c), \operatorname{product\text{-}sigma}(s, a, c, d), \operatorname{product\text{-}measure}(s, a, b, c, d, f)) \wedge \\ \quad u \in \operatorname{l2\text{-}set}(c, d, f) \wedge \\ \quad h \in \operatorname{l2\text{-}set}(c, d, f) \wedge \\ \quad \operatorname{is\text{-}null\text{-}set}(c, d, f, g) \wedge \\ \quad \forall y \in c.\; (\neg (y \in g) \Rightarrow (h(y) = u(y))) \Rightarrow \\ \quad \quad \exists n. \\ \quad \quad \quad \operatorname{is\text{-}null\text{-}set}(s, a, b, n) \wedge \\ \quad \quad \quad \forall x \in s. \\ \quad \quad \quad \quad \neg (x \in n) \Rightarrow  \\ \quad \quad \quad \quad \quad \int {k}(\cdot, y)\, {h}(y)\, d{f}(y)(x) = \int {k}(\cdot, y)\, {u}(y)\, d{f}(y)(x) \end{array}}}

\vnbbookfile{theorem-library/cl2-integral-operator}{theorem-library/cl2-integral-operator}
\noindent{\small cl2-integral-operator.scm -- wave of 2026-11-09/10 (agent TH3-L2C), wired 2026-11-10 from scratchpad/th3-1109/l2c/kc-b.scm. Library: measure.}\par
\vnbprop{cl2-hilbert-vadd-classes}{\vnbstmt{\begin{array}{@{}l@{}} \forall a, b, u, c, d, f, g. \\ \quad \operatorname{is\text{-}measure}(a, b, u) \wedge \\ \quad c \in \operatorname{l2\text{-}set}(a, b, u) \wedge \\ \quad d \in \operatorname{l2\text{-}set}(a, b, u) \wedge \\ \quad f \in \operatorname{l2\text{-}set}(a, b, u) \wedge \\ \quad g \in \operatorname{l2\text{-}set}(a, b, u) \Rightarrow \\ \quad \quad \operatorname{vadd}(L^2_{\mathbb{C}}({a}, {u}))({[}\operatorname{class}(\operatorname{sip\text{-}setoid}(\operatorname{l2\text{-}space}(a, b, u)), c), \operatorname{class}(\operatorname{sip\text{-}setoid}(\operatorname{l2\text{-}space}(a, b, u)), d)], \\ \quad \quad \quad {[}\operatorname{class}(\operatorname{sip\text{-}setoid}(\operatorname{l2\text{-}space}(a, b, u)), f), \operatorname{class}(\operatorname{sip\text{-}setoid}(\operatorname{l2\text{-}space}(a, b, u)), g)]) \\ \quad \quad =\; {[}\operatorname{class}(\operatorname{sip\text{-}setoid}(\operatorname{l2\text{-}space}(a, b, u)), \operatorname{vadd}(\operatorname{l2\text{-}space}(a, b, u))(c, f)), \operatorname{class}(\operatorname{sip\text{-}setoid}(\operatorname{l2\text{-}space}(a, b, u)), \operatorname{vadd}(\operatorname{l2\text{-}space}(a, b, u))(d, g))] \end{array}}}
\vnbprop{cl2-hilbert-act-classes}{\vnbstmt{\begin{array}{@{}l@{}} \forall a, b, u, z, c, d. \\ \quad (\operatorname{is\text{-}measure}(a, b, u) \wedge ((z \in \mathbb{C}) \wedge ((c \in \operatorname{l2\text{-}set}(a, b, u)) \wedge (d \in \operatorname{l2\text{-}set}(a, b, u))))) \Rightarrow  \\ \quad \quad \operatorname{act}(L^2_{\mathbb{C}}({a}, {u}))(z, \\ \quad \quad \quad {[}\operatorname{class}(\operatorname{sip\text{-}setoid}(\operatorname{l2\text{-}space}(a, b, u)), c), \operatorname{class}(\operatorname{sip\text{-}setoid}(\operatorname{l2\text{-}space}(a, b, u)), d)]) \\ \quad \quad =\; {[}\operatorname{class}(\operatorname{sip\text{-}setoid}(\operatorname{l2\text{-}space}(a, b, u)), \operatorname{vadd}(\operatorname{l2\text{-}space}(a, b, u))(\operatorname{act}(\operatorname{l2\text{-}space}(a, b, u))(\operatorname{real\text{-}part}(z), c), \operatorname{act}(\operatorname{l2\text{-}space}(a, b, u))(\operatorname{imag\text{-}part}(z), \operatorname{vneg}(\operatorname{l2\text{-}space}(a, b, u))(d)))), \operatorname{class}(\operatorname{sip\text{-}setoid}(\operatorname{l2\text{-}space}(a, b, u)), \operatorname{vadd}(\operatorname{l2\text{-}space}(a, b, u))(\operatorname{act}(\operatorname{l2\text{-}space}(a, b, u))(\operatorname{real\text{-}part}(z), d), \operatorname{act}(\operatorname{l2\text{-}space}(a, b, u))(\operatorname{imag\text{-}part}(z), c)))] \end{array}}}
\vnbprop{cl2-int-op-classes}{\vnbstmt{\begin{array}{@{}l@{}} \forall s, a, u, b, c, d, f, g, h, k. \\ \quad \operatorname{is\text{-}measure}(s, a, u) \wedge \\ \quad \operatorname{is\text{-}measure}(b, c, d) \wedge \\ \quad \operatorname{is\text{-}sigma\text{-}finite}(s, a, u) \wedge \\ \quad \operatorname{is\text{-}sigma\text{-}finite}(b, c, d) \wedge \\ \quad f \in \operatorname{l2\text{-}set}((s \times b), \operatorname{product\text{-}sigma}(s, a, b, c), \operatorname{product\text{-}measure}(s, a, u, b, c, d)) \wedge \\ \quad g \in \operatorname{l2\text{-}set}((s \times b), \operatorname{product\text{-}sigma}(s, a, b, c), \operatorname{product\text{-}measure}(s, a, u, b, c, d)) \wedge \\ \quad h \in \operatorname{l2\text{-}set}(b, c, d) \wedge \\ \quad k \in \operatorname{l2\text{-}set}(b, c, d) \Rightarrow \\ \quad \quad T_{{f} + i {g}}({[}\operatorname{class}(\operatorname{sip\text{-}setoid}(\operatorname{l2\text{-}space}(b, c, d)), h), \operatorname{class}(\operatorname{sip\text{-}setoid}(\operatorname{l2\text{-}space}(b, c, d)), k)]) \\ \quad \quad =\; {[}\operatorname{class}(\operatorname{sip\text{-}setoid}(\operatorname{l2\text{-}space}(s, a, u)), \operatorname{vadd}(\operatorname{l2\text{-}space}(s, a, u))(\int {f}(\cdot, y)\, {h}(y)\, d{d}(y), \operatorname{vneg}(\operatorname{l2\text{-}space}(s, a, u))(\int {g}(\cdot, y)\, {k}(y)\, d{d}(y)))), \operatorname{class}(\operatorname{sip\text{-}setoid}(\operatorname{l2\text{-}space}(s, a, u)), \operatorname{vadd}(\operatorname{l2\text{-}space}(s, a, u))(\int {f}(\cdot, y)\, {k}(y)\, d{d}(y), \int {g}(\cdot, y)\, {h}(y)\, d{d}(y)))] \end{array}}}
\vnbprop{cl2-int-op-fun}{\vnbstmt{\begin{array}{@{}l@{}} \forall s, a, u, b, c, d, f, g. \\ \quad \operatorname{is\text{-}measure}(s, a, u) \wedge \\ \quad \operatorname{is\text{-}measure}(b, c, d) \wedge \\ \quad \operatorname{is\text{-}sigma\text{-}finite}(s, a, u) \wedge \\ \quad \operatorname{is\text{-}sigma\text{-}finite}(b, c, d) \wedge \\ \quad f \in \operatorname{l2\text{-}set}((s \times b), \operatorname{product\text{-}sigma}(s, a, b, c), \operatorname{product\text{-}measure}(s, a, u, b, c, d)) \wedge \\ \quad g \in \operatorname{l2\text{-}set}((s \times b), \operatorname{product\text{-}sigma}(s, a, b, c), \operatorname{product\text{-}measure}(s, a, u, b, c, d)) \Rightarrow \\ \quad \quad T_{{f} + i {g}} \in (\operatorname{vec}(L^2_{\mathbb{C}}({b}, {d})) \to \operatorname{vec}(L^2_{\mathbb{C}}({s}, {u}))) \end{array}}}
\vnbprop{cl2-int-op-add}{\vnbstmt{\begin{array}{@{}l@{}} \forall s, a, u, b, c, d, f, g, z, w. \\ \quad \operatorname{is\text{-}measure}(s, a, u) \wedge \\ \quad \operatorname{is\text{-}measure}(b, c, d) \wedge \\ \quad \operatorname{is\text{-}sigma\text{-}finite}(s, a, u) \wedge \\ \quad \operatorname{is\text{-}sigma\text{-}finite}(b, c, d) \wedge \\ \quad f \in \operatorname{l2\text{-}set}((s \times b), \operatorname{product\text{-}sigma}(s, a, b, c), \operatorname{product\text{-}measure}(s, a, u, b, c, d)) \wedge \\ \quad g \in \operatorname{l2\text{-}set}((s \times b), \operatorname{product\text{-}sigma}(s, a, b, c), \operatorname{product\text{-}measure}(s, a, u, b, c, d)) \wedge \\ \quad z \in \operatorname{vec}(L^2_{\mathbb{C}}({b}, {d})) \wedge \\ \quad w \in \operatorname{vec}(L^2_{\mathbb{C}}({b}, {d})) \Rightarrow \\ \quad \quad T_{{f} + i {g}}(\operatorname{vadd}(L^2_{\mathbb{C}}({b}, {d}))(z, w)) \\ \quad \quad =\; \operatorname{vadd}(L^2_{\mathbb{C}}({s}, {u}))(T_{{f} + i {g}}(z), \\ \quad \quad \quad T_{{f} + i {g}}(w)) \end{array}}}
\vnbprop{cl2-int-op-act}{\vnbstmt{\begin{array}{@{}l@{}} \forall s, a, u, b, c, d, f, g, \gamma, z. \\ \quad \operatorname{is\text{-}measure}(s, a, u) \wedge \\ \quad \operatorname{is\text{-}measure}(b, c, d) \wedge \\ \quad \operatorname{is\text{-}sigma\text{-}finite}(s, a, u) \wedge \\ \quad \operatorname{is\text{-}sigma\text{-}finite}(b, c, d) \wedge \\ \quad f \in \operatorname{l2\text{-}set}((s \times b), \operatorname{product\text{-}sigma}(s, a, b, c), \operatorname{product\text{-}measure}(s, a, u, b, c, d)) \wedge \\ \quad g \in \operatorname{l2\text{-}set}((s \times b), \operatorname{product\text{-}sigma}(s, a, b, c), \operatorname{product\text{-}measure}(s, a, u, b, c, d)) \wedge \\ \quad \gamma \in \mathbb{C} \wedge \\ \quad z \in \operatorname{vec}(L^2_{\mathbb{C}}({b}, {d})) \Rightarrow \\ \quad \quad T_{{f} + i {g}}(\operatorname{act}(L^2_{\mathbb{C}}({b}, {d}))(\gamma, z)) \\ \quad \quad =\; \operatorname{act}(L^2_{\mathbb{C}}({s}, {u}))(\gamma, T_{{f} + i {g}}(z)) \end{array}}}
\vnbprop{cl2-hilbert-ip-self}{\vnbstmt{\begin{array}{@{}l@{}} \forall a, b, u, c, d. \\ \quad (\operatorname{is\text{-}measure}(a, b, u) \wedge ((c \in \operatorname{l2\text{-}set}(a, b, u)) \wedge (d \in \operatorname{l2\text{-}set}(a, b, u)))) \Rightarrow  \\ \quad \quad \langle\cdot,\cdot\rangle_{L^2_{\mathbb{C}}({a}, {u})}({[}\operatorname{class}(\operatorname{sip\text{-}setoid}(\operatorname{l2\text{-}space}(a, b, u)), c), \operatorname{class}(\operatorname{sip\text{-}setoid}(\operatorname{l2\text{-}space}(a, b, u)), d)], \\ \quad \quad \quad {[}\operatorname{class}(\operatorname{sip\text{-}setoid}(\operatorname{l2\text{-}space}(a, b, u)), c), \operatorname{class}(\operatorname{sip\text{-}setoid}(\operatorname{l2\text{-}space}(a, b, u)), d)]) \\ \quad \quad =\; \operatorname{rintegral}(a, b, u, (\lambda x \in a.\; (c(x) \cdot c(x)))) + \operatorname{rintegral}(a, b, u, (\lambda x \in a.\; (d(x) \cdot d(x)))) \end{array}}}
\vnbprop{l2-int-neg-right}{\vnbstmt{\begin{array}{@{}l@{}} \forall a, b, u, p, f. \\ \quad (\operatorname{is\text{-}measure}(a, b, u) \wedge ((p \in \operatorname{l2\text{-}set}(a, b, u)) \wedge (f \in \operatorname{l2\text{-}set}(a, b, u)))) \Rightarrow  \\ \quad \quad \operatorname{rintegral}(a, b, u, (\lambda x \in a.\; (p(x) \cdot \operatorname{vneg}(\operatorname{l2\text{-}space}(a, b, u))(f)(x)))) \\ \quad \quad =\; -\operatorname{rintegral}(a, b, u, (\lambda x \in a.\; (p(x) \cdot f(x)))) \end{array}}}
\vnbprop{l2-int-neg-self}{\vnbstmt{\begin{array}{@{}l@{}} \forall a, b, u.; \forall f \in \operatorname{l2\text{-}set}(a, b, u). \\ \quad \operatorname{is\text{-}measure}(a, b, u) \Rightarrow  \\ \quad \quad \operatorname{rintegral}(a, b, u, \\ \quad \quad \quad (\lambda x \in a.\; (\operatorname{vneg}(\operatorname{l2\text{-}space}(a, b, u))(f)(x) \cdot \operatorname{vneg}(\operatorname{l2\text{-}space}(a, b, u))(f)(x)))) \\ \quad \quad =\; \operatorname{rintegral}(a, b, u, (\lambda x \in a.\; (f(x) \cdot f(x)))) \end{array}}}
\vnbprop{cc-mul-conj-parts}{\vnbstmt{\begin{array}{@{}l@{}} \forall a, b. \\ \quad ((a \in \mathbb{R}) \wedge (b \in \mathbb{R})) \Rightarrow (((a + (+i \cdot b)) \cdot \operatorname{conjugate}((a + (+i \cdot b)))) = ((a \cdot a) + (b \cdot b))) \end{array}}}
\vnbprop{kernel-apply-cschwarz}{\vnbstmt{\begin{array}{@{}l@{}} \forall s, a, u, b, c, d, f, g, h, k. \\ \quad \operatorname{is\text{-}measure}(s, a, u) \wedge \\ \quad \operatorname{is\text{-}measure}(b, c, d) \wedge \\ \quad \operatorname{is\text{-}sigma\text{-}finite}(s, a, u) \wedge \\ \quad \operatorname{is\text{-}sigma\text{-}finite}(b, c, d) \wedge \\ \quad f \in \operatorname{l2\text{-}set}((s \times b), \operatorname{product\text{-}sigma}(s, a, b, c), \operatorname{product\text{-}measure}(s, a, u, b, c, d)) \wedge \\ \quad g \in \operatorname{l2\text{-}set}((s \times b), \operatorname{product\text{-}sigma}(s, a, b, c), \operatorname{product\text{-}measure}(s, a, u, b, c, d)) \wedge \\ \quad h \in \operatorname{l2\text{-}set}(b, c, d) \wedge \\ \quad k \in \operatorname{l2\text{-}set}(b, c, d) \Rightarrow \\ \quad \quad \exists n. \\ \quad \quad \quad \operatorname{is\text{-}null\text{-}set}(s, a, u, n) \wedge \\ \quad \quad \quad \forall x \in s. \\ \quad \quad \quad \quad \neg (x \in n) \Rightarrow  \\ \quad \quad \quad \quad \quad \operatorname{x\text{-}section\text{-}integral}(s, b, c, d, (\lambda \mathit{lx} \in (s \times b).\; (f(\mathit{lx}) \cdot f(\mathit{lx}))))(x) \in \mathbb{R} \wedge \\ \quad \quad \quad \quad \quad \operatorname{x\text{-}section\text{-}integral}(s, b, c, d, (\lambda \mathit{lx} \in (s \times b).\; (g(\mathit{lx}) \cdot g(\mathit{lx}))))(x) \in \mathbb{R} \wedge \\ \quad \quad \quad \quad \quad (\operatorname{vadd}(\operatorname{l2\text{-}space}(s, a, u))(\int {f}(\cdot, y)\, {h}(y)\, d{d}(y), \operatorname{vneg}(\operatorname{l2\text{-}space}(s, a, u))(\int {g}(\cdot, y)\, {k}(y)\, d{d}(y)))(x) \cdot \operatorname{vadd}(\operatorname{l2\text{-}space}(s, a, u))(\int {f}(\cdot, y)\, {h}(y)\, d{d}(y), \operatorname{vneg}(\operatorname{l2\text{-}space}(s, a, u))(\int {g}(\cdot, y)\, {k}(y)\, d{d}(y)))(x)) + (\operatorname{vadd}(\operatorname{l2\text{-}space}(s, a, u))(\int {f}(\cdot, y)\, {k}(y)\, d{d}(y), \int {g}(\cdot, y)\, {h}(y)\, d{d}(y))(x) \cdot \operatorname{vadd}(\operatorname{l2\text{-}space}(s, a, u))(\int {f}(\cdot, y)\, {k}(y)\, d{d}(y), \int {g}(\cdot, y)\, {h}(y)\, d{d}(y))(x)) \\ \quad \quad \quad \quad \quad \leq\; (\operatorname{x\text{-}section\text{-}integral}(s, b, c, d, (\lambda \mathit{lx} \in (s \times b).\; (f(\mathit{lx}) \cdot f(\mathit{lx}))))(x) + \operatorname{x\text{-}section\text{-}integral}(s, b, c, d, (\lambda \mathit{lx} \in (s \times b).\; (g(\mathit{lx}) \cdot g(\mathit{lx}))))(x)) \cdot (\operatorname{integral}(b, c, d, (\lambda \mathit{lx} \in b.\; (h(\mathit{lx}) \cdot h(\mathit{lx})))) + \operatorname{integral}(b, c, d, (\lambda \mathit{lx} \in b.\; (k(\mathit{lx}) \cdot k(\mathit{lx}))))) \end{array}}}
\vnbprop{l2-pair-square-bound}{\vnbstmt{\begin{array}{@{}l@{}} \forall a, b, u, d, f, h, c, g. \\ \quad \operatorname{is\text{-}measure}(a, b, u) \wedge \\ \quad d \in \operatorname{l2\text{-}set}(a, b, u) \wedge \\ \quad f \in \operatorname{l2\text{-}set}(a, b, u) \wedge \\ \quad \operatorname{is\text{-}measurable\text{-}fn}(a, b, h) \wedge \\ \quad c \in \mathbb{R} \wedge \\ \quad 0 \leq c \wedge \\ \quad \operatorname{integral}(a, b, u, h) \in \mathbb{R} \wedge \\ \quad g \in b \wedge \\ \quad u(g) = 0 \wedge \\ \quad \forall x \in a. \\ \quad \quad \neg (x \in g) \Rightarrow ((((d(x) \cdot d(x)) + (f(x) \cdot f(x))) \leq (h(x) \cdot c)) \wedge (h(x) \in \mathbb{R})) \Rightarrow \\ \quad \quad \operatorname{integral}(a, b, u, (\lambda \mathit{lx} \in a.\; (d(\mathit{lx}) \cdot d(\mathit{lx})))) + \operatorname{integral}(a, b, u, (\lambda \mathit{lx} \in a.\; (f(\mathit{lx}) \cdot f(\mathit{lx})))) \\ \quad \quad \leq\; \operatorname{integral}(a, b, u, h) \cdot c \end{array}}}
\vnbprop{cl2-int-op-ip-bound}{\vnbstmt{\begin{array}{@{}l@{}} \forall s, a, u, b, c, d, f, g, z. \\ \quad \operatorname{is\text{-}measure}(s, a, u) \wedge \\ \quad \operatorname{is\text{-}measure}(b, c, d) \wedge \\ \quad \operatorname{is\text{-}sigma\text{-}finite}(s, a, u) \wedge \\ \quad \operatorname{is\text{-}sigma\text{-}finite}(b, c, d) \wedge \\ \quad f \in \operatorname{l2\text{-}set}((s \times b), \operatorname{product\text{-}sigma}(s, a, b, c), \operatorname{product\text{-}measure}(s, a, u, b, c, d)) \wedge \\ \quad g \in \operatorname{l2\text{-}set}((s \times b), \operatorname{product\text{-}sigma}(s, a, b, c), \operatorname{product\text{-}measure}(s, a, u, b, c, d)) \wedge \\ \quad z \in \operatorname{vec}(L^2_{\mathbb{C}}({b}, {d})) \Rightarrow \\ \quad \quad \langle\cdot,\cdot\rangle_{L^2_{\mathbb{C}}({s}, {u})}(T_{{f} + i {g}}(z), \\ \quad \quad \quad T_{{f} + i {g}}(z)) \\ \quad \quad \leq\; (\operatorname{integral}((s \times b), \operatorname{product\text{-}sigma}(s, a, b, c), \operatorname{product\text{-}measure}(s, a, u, b, c, d), (\lambda x \in (s \times b).\; (f(x) \cdot f(x)))) + \operatorname{integral}((s \times b), \operatorname{product\text{-}sigma}(s, a, b, c), \operatorname{product\text{-}measure}(s, a, u, b, c, d), (\lambda x \in (s \times b).\; (g(x) \cdot g(x))))) \cdot \langle\cdot,\cdot\rangle_{L^2_{\mathbb{C}}({b}, {d})}(z, z) \end{array}}}

\vnbbookfile{theorem-library/complex-measure-polar}{theorem-library/complex-measure-polar}
\noindent{\small complex-measure-polar.scm -- WAVE 2 (2026-10-27..30), wired 2026-10-30 from scratchpad/mf/mf-c5.scm, scratchpad/mf/mf-c6.scm, scratchpad/mf/mf-c7.scm, scratchpad/mf/mf-c8.scm. Library: measure.}\par
\vnbprop{rintegral-density}{\vnbstmt{\begin{array}{@{}l@{}} \forall a, b, o, p, g, f. \\ \quad \operatorname{is\text{-}measure}(a, b, o) \wedge \\ \quad \operatorname{is\text{-}measurable\text{-}fn}(a, b, p) \wedge \\ \quad \forall x \in a.\; (p(x) \in \mathbb{R}) \wedge \\ \quad \operatorname{is\text{-}real\text{-}measurable\text{-}fn}(a, b, g) \wedge \\ \quad f \in (a \to \mathbb{R}) \wedge \\ \quad \forall x \in a.\; (f(x) = (p(x) \cdot g(x))) \wedge \\ \quad \operatorname{is\text{-}integrable}(a, b, o, f) \Rightarrow \\ \quad \quad \operatorname{rintegral}(a, b, (\lambda \mathit{dma} \in b.\; \operatorname{integral}(a, b, o, \operatorname{ptwise\text{-}etimes}(a, p, \operatorname{indicator}(a, \mathit{dma})))), g) \\ \quad \quad =\; \operatorname{rintegral}(a, b, o, f) \end{array}}}
\vnbprop{indicator-real-measurable}{\vnbstmt{\begin{array}{@{}l@{}} \forall b, c, a. \\ \quad (\operatorname{is\text{-}sigma\text{-}algebra}(b, c) \wedge (a \in c)) \Rightarrow  \\ \quad \quad \operatorname{is\text{-}real\text{-}measurable\text{-}fn}(b, c, \operatorname{indicator}(b, a)) \end{array}}}
\vnbprop{four-finite-measures-below}{\vnbstmt{\begin{array}{@{}l@{}} \forall a, b, c, d, f, g. \\ \quad \operatorname{is\text{-}finite\text{-}measure}(a, b, c) \wedge \\ \quad \operatorname{is\text{-}finite\text{-}measure}(a, b, d) \wedge \\ \quad \operatorname{is\text{-}finite\text{-}measure}(a, b, f) \wedge \\ \quad \operatorname{is\text{-}finite\text{-}measure}(a, b, g) \Rightarrow \\ \quad \quad \exists o. \\ \quad \quad \quad \operatorname{is\text{-}finite\text{-}measure}(a, b, o) \wedge \\ \quad \quad \quad \forall \alpha \in b.\; (c(\alpha) \leq o(\alpha)) \wedge \\ \quad \quad \quad \forall \alpha \in b.\; (d(\alpha) \leq o(\alpha)) \wedge \\ \quad \quad \quad \forall \alpha \in b.\; (f(\alpha) \leq o(\alpha)) \wedge \\ \quad \quad \quad \forall \alpha \in b.\; (g(\alpha) \leq o(\alpha)) \end{array}}}
\vnbprop{polar-real-part}{\vnbstmt{\begin{array}{@{}l@{}} \forall b, c, o, p, u, d, f, g, h, a. \\ \quad \operatorname{is\text{-}measure}(b, c, o) \wedge \\ \quad \operatorname{is\text{-}measurable\text{-}fn}(b, c, p) \wedge \\ \quad \forall \mathit{rnx} \in b.\; (p(\mathit{rnx}) \in \mathbb{R}) \wedge \\ \quad \operatorname{is\text{-}real\text{-}measurable\text{-}fn}(b, c, u) \wedge \\ \quad \operatorname{is\text{-}measurable\text{-}fn}(b, c, d) \wedge \\ \quad \operatorname{is\text{-}measurable\text{-}fn}(b, c, f) \wedge \\ \quad \forall \mathit{rnx} \in b.\; (d(\mathit{rnx}) \in \mathbb{R}) \wedge \\ \quad \forall \mathit{rnx} \in b.\; (f(\mathit{rnx}) \in \mathbb{R}) \wedge \\ \quad \operatorname{is\text{-}finite\text{-}measure}(b, c, g) \wedge \\ \quad \operatorname{is\text{-}finite\text{-}measure}(b, c, h) \wedge \\ \quad \forall \mathit{rna} \in c. \\ \quad \quad g(\mathit{rna}) = \operatorname{integral}(b, c, o, \operatorname{ptwise\text{-}etimes}(b, d, \operatorname{indicator}(b, \mathit{rna}))) \wedge \\ \quad \forall \mathit{rna} \in c. \\ \quad \quad h(\mathit{rna}) = \operatorname{integral}(b, c, o, \operatorname{ptwise\text{-}etimes}(b, f, \operatorname{indicator}(b, \mathit{rna}))) \wedge \\ \quad \forall x \in b.\; ((p(x) \cdot u(x)) = (d(x) - f(x))) \wedge \\ \quad a \in c \Rightarrow \\ \quad \quad \operatorname{rintegral}(b, c, \\ \quad \quad \quad (\lambda \mathit{dma} \in c.\; \operatorname{integral}(b, c, o, \operatorname{ptwise\text{-}etimes}(b, p, \operatorname{indicator}(b, \mathit{dma})))), \\ \quad \quad \quad (\lambda \mathit{cix} \in b.\; (\operatorname{indicator}(b, a)(\mathit{cix}) \cdot u(\mathit{cix})))) \\ \quad \quad =\; (g(a) - h(a)) \end{array}}}
\vnbprop{rintegral-diff-density}{\vnbstmt{\begin{array}{@{}l@{}} \forall a, b, o, d, f, g, h, c. \\ \quad \operatorname{is\text{-}measure}(a, b, o) \wedge \\ \quad \operatorname{is\text{-}measurable\text{-}fn}(a, b, d) \wedge \\ \quad \operatorname{is\text{-}measurable\text{-}fn}(a, b, f) \wedge \\ \quad \forall x \in a.\; (d(x) \in \mathbb{R}) \wedge \\ \quad \forall x \in a.\; (f(x) \in \mathbb{R}) \wedge \\ \quad \operatorname{is\text{-}finite\text{-}measure}(a, b, g) \wedge \\ \quad \operatorname{is\text{-}finite\text{-}measure}(a, b, h) \wedge \\ \quad \forall \mathit{rna} \in b. \\ \quad \quad g(\mathit{rna}) = \operatorname{integral}(a, b, o, \operatorname{ptwise\text{-}etimes}(a, d, \operatorname{indicator}(a, \mathit{rna}))) \wedge \\ \quad \forall \mathit{rna} \in b. \\ \quad \quad h(\mathit{rna}) = \operatorname{integral}(a, b, o, \operatorname{ptwise\text{-}etimes}(a, f, \operatorname{indicator}(a, \mathit{rna}))) \wedge \\ \quad c \in b \Rightarrow \\ \quad \quad \operatorname{is\text{-}integrable}(a, b, o, (\lambda \mathit{hx} \in a.\; (\operatorname{indicator}(a, c)(\mathit{hx}) \cdot (d(\mathit{hx}) - f(\mathit{hx}))))) \wedge \\ \quad \quad \operatorname{rintegral}(a, b, o, (\lambda \mathit{hx} \in a.\; (\operatorname{indicator}(a, c)(\mathit{hx}) \cdot (d(\mathit{hx}) - f(\mathit{hx}))))) = (g(c) - h(c)) \end{array}}}
\vnbprop{rintegral-null-support}{\vnbstmt{\begin{array}{@{}l@{}} \forall a, c, u, b, g. \\ \quad (\operatorname{is\text{-}measure}(a, c, u) \wedge ((b \in c) \wedge ((u(b) = 0) \wedge \operatorname{is\text{-}real\text{-}measurable\text{-}fn}(a, c, g)))) \Rightarrow  \\ \quad \quad \operatorname{rintegral}(a, c, u, (\lambda x \in a.\; (\operatorname{indicator}(a, b)(x) \cdot g(x)))) = 0 \end{array}}}
\vnbprop{rintegral-zero-on-sets-null}{\vnbstmt{\begin{array}{@{}l@{}} \forall a, b, o, f. \\ \quad \operatorname{is\text{-}measure}(a, b, o) \wedge \\ \quad \operatorname{is\text{-}integrable}(a, b, o, f) \wedge \\ \quad \forall \mathit{bb} \in b. \\ \quad \quad \operatorname{rintegral}(a, b, o, (\lambda x \in a.\; (f(x) \cdot \operatorname{indicator}(a, \mathit{bb})(x)))) = 0 \Rightarrow \\ \quad \quad \operatorname{is\text{-}null\text{-}set}(a, b, o, \{\mathit{nzx} \in a : \neg (f(\mathit{nzx}) = 0)\}) \end{array}}}
\vnbprop{polar-total-finite}{\vnbstmt{\begin{array}{@{}l@{}} \forall a, b, o, p, c, d, f, g, h, k, m, n. \\ \quad \operatorname{is\text{-}measure}(a, b, o) \wedge \\ \quad \operatorname{is\text{-}measurable\text{-}fn}(a, b, p) \wedge \\ \quad \forall \mathit{rnx} \in a.\; (p(\mathit{rnx}) \in \mathbb{R}) \wedge \\ \quad \operatorname{is\text{-}measurable\text{-}fn}(a, b, c) \wedge \\ \quad \operatorname{is\text{-}measurable\text{-}fn}(a, b, d) \wedge \\ \quad \operatorname{is\text{-}measurable\text{-}fn}(a, b, f) \wedge \\ \quad \operatorname{is\text{-}measurable\text{-}fn}(a, b, g) \wedge \\ \quad \forall \mathit{rnx} \in a.\; (c(\mathit{rnx}) \in \mathbb{R}) \wedge \\ \quad \forall \mathit{rnx} \in a.\; (d(\mathit{rnx}) \in \mathbb{R}) \wedge \\ \quad \forall \mathit{rnx} \in a.\; (f(\mathit{rnx}) \in \mathbb{R}) \wedge \\ \quad \forall \mathit{rnx} \in a.\; (g(\mathit{rnx}) \in \mathbb{R}) \wedge \\ \quad \operatorname{is\text{-}finite\text{-}measure}(a, b, h) \wedge \\ \quad \operatorname{is\text{-}finite\text{-}measure}(a, b, k) \wedge \\ \quad \operatorname{is\text{-}finite\text{-}measure}(a, b, m) \wedge \\ \quad \operatorname{is\text{-}finite\text{-}measure}(a, b, n) \wedge \\ \quad \forall \mathit{rna} \in b. \\ \quad \quad h(\mathit{rna}) = \operatorname{integral}(a, b, o, \operatorname{ptwise\text{-}etimes}(a, c, \operatorname{indicator}(a, \mathit{rna}))) \wedge \\ \quad \forall \mathit{rna} \in b. \\ \quad \quad k(\mathit{rna}) = \operatorname{integral}(a, b, o, \operatorname{ptwise\text{-}etimes}(a, d, \operatorname{indicator}(a, \mathit{rna}))) \wedge \\ \quad \forall \mathit{rna} \in b. \\ \quad \quad m(\mathit{rna}) = \operatorname{integral}(a, b, o, \operatorname{ptwise\text{-}etimes}(a, f, \operatorname{indicator}(a, \mathit{rna}))) \wedge \\ \quad \forall \mathit{rna} \in b. \\ \quad \quad n(\mathit{rna}) = \operatorname{integral}(a, b, o, \operatorname{ptwise\text{-}etimes}(a, g, \operatorname{indicator}(a, \mathit{rna}))) \wedge \\ \quad \forall x \in a.\; (p(x) \leq ((c(x) + d(x)) + (f(x) + g(x)))) \Rightarrow \\ \quad \quad \operatorname{integral}(a, b, o, \operatorname{ptwise\text{-}etimes}(a, p, \operatorname{indicator}(a, a))) \in \mathbb{R} \end{array}}}
\vnbprop{polar-decomposition-exists}{\vnbstmt{\begin{array}{@{}l@{}} \forall a, b, u. \\ \quad \operatorname{is\text{-}complex\text{-}measure}(a, b, u) \Rightarrow  \\ \quad \quad \exists v. \\ \quad \quad \quad \operatorname{is\text{-}finite\text{-}measure}(a, b, v) \wedge \\ \quad \quad \quad \forall \alpha \in b.\; ((v(\alpha) = 0) \Leftrightarrow \operatorname{is\text{-}complex\text{-}null\text{-}set}(a, b, u, \alpha)) \wedge \\ \quad \quad \quad \exists c. \\ \quad \quad \quad \quad \operatorname{is\text{-}complex\text{-}measurable\text{-}fn}(a, b, c) \wedge \\ \quad \quad \quad \quad \forall x \in a.\; (\operatorname{magnitude}(c(x)) = 1) \wedge \\ \quad \quad \quad \quad \forall \alpha \in b. \\ \quad \quad \quad \quad \quad u(\alpha) = \operatorname{complex\text{-}integral}(a, b, v, (\lambda \mathit{cux} \in a.\; (\operatorname{indicator}(a, \alpha)(\mathit{cux}) \cdot c(\mathit{cux})))) \end{array}}}
\vnbprop{rintegral-density-hyp}{\vnbstmt{\begin{array}{@{}l@{}} \forall a, b, o, u, h, g, f. \\ \quad \operatorname{is\text{-}measure}(a, b, o) \wedge \\ \quad \operatorname{is\text{-}measure}(a, b, u) \wedge \\ \quad \operatorname{is\text{-}measurable\text{-}fn}(a, b, h) \wedge \\ \quad \forall \mathit{rnx} \in a.\; (h(\mathit{rnx}) \in \mathbb{R}) \wedge \\ \quad \forall \mathit{rna} \in b. \\ \quad \quad u(\mathit{rna}) = \operatorname{integral}(a, b, o, \operatorname{ptwise\text{-}etimes}(a, h, \operatorname{indicator}(a, \mathit{rna}))) \wedge \\ \quad \operatorname{is\text{-}real\text{-}measurable\text{-}fn}(a, b, g) \wedge \\ \quad f \in (a \to \mathbb{R}) \wedge \\ \quad \forall x \in a.\; (f(x) = (h(x) \cdot g(x))) \wedge \\ \quad \operatorname{is\text{-}integrable}(a, b, o, f) \Rightarrow \\ \quad \quad \operatorname{rintegral}(a, b, u, g) = \operatorname{rintegral}(a, b, o, f) \end{array}}}
\vnbprop{density-integrable}{\vnbstmt{\begin{array}{@{}l@{}} \forall a, b, o, u, h. \\ \quad \operatorname{is\text{-}measure}(a, b, o) \wedge \\ \quad \operatorname{is\text{-}measurable\text{-}fn}(a, b, h) \wedge \\ \quad \forall x \in a.\; (h(x) \in \mathbb{R}) \wedge \\ \quad \forall \mathit{rna} \in b. \\ \quad \quad u(\mathit{rna}) = \operatorname{integral}(a, b, o, \operatorname{ptwise\text{-}etimes}(a, h, \operatorname{indicator}(a, \mathit{rna}))) \wedge \\ \quad u(a) \in \mathbb{R} \Rightarrow \\ \quad \quad \operatorname{is\text{-}integrable}(a, b, o, h) \end{array}}}
\vnbprop{polar-part-density-re}{\vnbstmt{\begin{array}{@{}l@{}} \forall b, c, o, d, h, u, a. \\ \quad \operatorname{is\text{-}measure}(b, c, o) \wedge \\ \quad \operatorname{is\text{-}measure}(b, c, d) \wedge \\ \quad \operatorname{is\text{-}measurable\text{-}fn}(b, c, h) \wedge \\ \quad \forall \mathit{rnx} \in b.\; (h(\mathit{rnx}) \in \mathbb{R}) \wedge \\ \quad \forall \mathit{rna} \in c. \\ \quad \quad d(\mathit{rna}) = \operatorname{integral}(b, c, o, \operatorname{ptwise\text{-}etimes}(b, h, \operatorname{indicator}(b, \mathit{rna}))) \wedge \\ \quad d(b) \in \mathbb{R} \wedge \\ \quad \operatorname{is\text{-}complex\text{-}measurable\text{-}fn}(b, c, u) \wedge \\ \quad \forall x \in b.\; (\operatorname{magnitude}(u(x)) = 1) \wedge \\ \quad a \in c \Rightarrow \\ \quad \quad \operatorname{is\text{-}integrable}(b, c, o, (\lambda \mathit{ux} \in b.\; (h(\mathit{ux}) \cdot (\operatorname{indicator}(b, a)(\mathit{ux}) \cdot \operatorname{real\text{-}part}(u(\mathit{ux})))))) \wedge \\ \quad \quad \operatorname{rintegral}(b, c, d, \\ \quad \quad \quad (\lambda \mathit{cix} \in b.\; \operatorname{real\text{-}part}((\lambda \mathit{cux} \in b.\; (\operatorname{indicator}(b, a)(\mathit{cux}) \cdot u(\mathit{cux})))(\mathit{cix})))) \\ \quad \quad =\; \operatorname{rintegral}(b, c, o, (\lambda \mathit{ux} \in b.\; (h(\mathit{ux}) \cdot (\operatorname{indicator}(b, a)(\mathit{ux}) \cdot \operatorname{real\text{-}part}(u(\mathit{ux})))))) \end{array}}}
\vnbprop{polar-part-density-im}{\vnbstmt{\begin{array}{@{}l@{}} \forall b, c, o, d, h, u, a. \\ \quad \operatorname{is\text{-}measure}(b, c, o) \wedge \\ \quad \operatorname{is\text{-}measure}(b, c, d) \wedge \\ \quad \operatorname{is\text{-}measurable\text{-}fn}(b, c, h) \wedge \\ \quad \forall \mathit{rnx} \in b.\; (h(\mathit{rnx}) \in \mathbb{R}) \wedge \\ \quad \forall \mathit{rna} \in c. \\ \quad \quad d(\mathit{rna}) = \operatorname{integral}(b, c, o, \operatorname{ptwise\text{-}etimes}(b, h, \operatorname{indicator}(b, \mathit{rna}))) \wedge \\ \quad d(b) \in \mathbb{R} \wedge \\ \quad \operatorname{is\text{-}complex\text{-}measurable\text{-}fn}(b, c, u) \wedge \\ \quad \forall x \in b.\; (\operatorname{magnitude}(u(x)) = 1) \wedge \\ \quad a \in c \Rightarrow \\ \quad \quad \operatorname{is\text{-}integrable}(b, c, o, (\lambda \mathit{ux} \in b.\; (h(\mathit{ux}) \cdot (\operatorname{indicator}(b, a)(\mathit{ux}) \cdot \operatorname{imag\text{-}part}(u(\mathit{ux})))))) \wedge \\ \quad \quad \operatorname{rintegral}(b, c, d, \\ \quad \quad \quad (\lambda \mathit{cix} \in b.\; \operatorname{imag\text{-}part}((\lambda \mathit{cux} \in b.\; (\operatorname{indicator}(b, a)(\mathit{cux}) \cdot u(\mathit{cux})))(\mathit{cix})))) \\ \quad \quad =\; \operatorname{rintegral}(b, c, o, (\lambda \mathit{ux} \in b.\; (h(\mathit{ux}) \cdot (\operatorname{indicator}(b, a)(\mathit{ux}) \cdot \operatorname{imag\text{-}part}(u(\mathit{ux})))))) \end{array}}}
\vnbprop{polar-decomposition-unique}{\vnbstmt{\begin{array}{@{}l@{}} \forall a, b, u, c, d, f, g. \\ \quad \operatorname{is\text{-}finite\text{-}measure}(a, b, c) \wedge \\ \quad \operatorname{is\text{-}finite\text{-}measure}(a, b, d) \wedge \\ \quad \operatorname{is\text{-}complex\text{-}measurable\text{-}fn}(a, b, f) \wedge \\ \quad \operatorname{is\text{-}complex\text{-}measurable\text{-}fn}(a, b, g) \wedge \\ \quad \forall x \in a.\; (\operatorname{magnitude}(f(x)) = 1) \wedge \\ \quad \forall x \in a.\; (\operatorname{magnitude}(g(x)) = 1) \wedge \\ \quad \forall \alpha \in b. \\ \quad \quad u(\alpha) = \operatorname{complex\text{-}integral}(a, b, c, (\lambda \mathit{cux} \in a.\; (\operatorname{indicator}(a, \alpha)(\mathit{cux}) \cdot f(\mathit{cux})))) \wedge \\ \quad \forall \alpha \in b. \\ \quad \quad u(\alpha) = \operatorname{complex\text{-}integral}(a, b, d, (\lambda \mathit{cux} \in a.\; (\operatorname{indicator}(a, \alpha)(\mathit{cux}) \cdot g(\mathit{cux})))) \Rightarrow \\ \quad \quad (\forall \alpha \in b.\; (c(\alpha) = d(\alpha)) \wedge \operatorname{is\text{-}null\text{-}set}(a, b, c, \{\mathit{ux} \in a : \neg (f(\mathit{ux}) = g(\mathit{ux}))\})) \end{array}}}

\vnbbookfile{theorem-library/mb35-souslin}{theorem-library/mb35-souslin}
\noindent{\small mb35-souslin.scm -- construction of measures, chapter 5 sections 1.3-1.4: Souslin sets (M-B35 / B5, 2026-10-26).  Vocabulary: mb35-souslin-vocab.scm.  Library: measure.  Helper prefixes mr- (copied from mb35-regular.scm) and msn-.}\par
\vnbprop{is-sieve-classes}{\vnbstmt{\begin{array}{@{}l@{}} \forall s, v. \\ \quad \operatorname{is\text{-}sieve}(s, v) \\ \quad \Leftrightarrow\; v \in (\mathit{nstar} \to \operatorname{power}(\operatorname{pts}(s))) \wedge \\ \quad \forall z \in (\mathbb{N} \to \mathbb{N}). \\ \quad \quad \forall e. \\ \quad \quad \quad \operatorname{pos\text{-}rr}(e) \Rightarrow  \\ \quad \quad \quad \quad \exists b \in \mathbb{N}. \forall n \in \mathbb{N}. \\ \quad \quad \quad \quad \quad (b \leq n) \Rightarrow (\operatorname{diam}(s, v(\operatorname{restrict}(z, \mathrm{OS}(n)))) \leq e) \wedge \\ \quad \quad \forall n \in \mathbb{N}. \\ \quad \quad \quad v(\operatorname{restrict}(z, \mathrm{OS}((n + 1)))) \subseteq v(\operatorname{restrict}(z, \mathrm{OS}(n))) \wedge \\ \quad \quad \exists p.\; \forall n \in \mathbb{N}.\; (p \in v(\operatorname{restrict}(z, \mathrm{OS}(n)))) \end{array}}}
\vnbprop{nstar-is-set}{\vnbstmt{\mathit{nstar} \in \mathbf{Set}}}
\vnbprop{prefix-in-nstar}{\vnbstmt{\begin{array}{@{}l@{}} \forall z \in (\mathbb{N} \to \mathbb{N}), n \in \mathbb{N}. \\ \quad ((\operatorname{restrict}(z, \mathrm{OS}(n)) \in (\mathrm{OS}(n) \to \mathbb{N})) \wedge (\operatorname{restrict}(z, \mathrm{OS}(n)) \in \mathit{nstar})) \end{array}}}
\vnbprop{prefix-eq-agree}{\vnbstmt{\begin{array}{@{}l@{}} \forall z, w \in (\mathbb{N} \to \mathbb{N}), n, k \in \mathbb{N}. \\ \quad ((\operatorname{restrict}(w, \mathrm{OS}(n)) = \operatorname{restrict}(z, \mathrm{OS}(n))) \wedge (k < n)) \Rightarrow  \\ \quad \quad w(k) = z(k) \end{array}}}
\vnbprop{prefix-eq-intro}{\vnbstmt{\begin{array}{@{}l@{}} \forall z, w \in (\mathbb{N} \to \mathbb{N}), n \in \mathbb{N}. \\ \quad \forall k \in \mathbb{N}.\; ((k < n) \Rightarrow (w(k) = z(k))) \Rightarrow  \\ \quad \quad \operatorname{restrict}(w, \mathrm{OS}(n)) = \operatorname{restrict}(z, \mathrm{OS}(n)) \end{array}}}
\vnbprop{prefix-eq-le}{\vnbstmt{\begin{array}{@{}l@{}} \forall z, w \in (\mathbb{N} \to \mathbb{N}), n, m \in \mathbb{N}. \\ \quad ((m \leq n) \wedge (\operatorname{restrict}(w, \mathrm{OS}(n)) = \operatorname{restrict}(z, \mathrm{OS}(n)))) \Rightarrow  \\ \quad \quad \operatorname{restrict}(w, \mathrm{OS}(m)) = \operatorname{restrict}(z, \mathrm{OS}(m)) \end{array}}}
\vnbprop{baire-image-souslin}{\vnbstmt{\begin{array}{@{}l@{}} \forall s, f. \\ \quad (\operatorname{is\text{-}metric\text{-}space}(s) \wedge \operatorname{is\text{-}baire\text{-}continuous}(s, f)) \Rightarrow  \\ \quad \quad \operatorname{is\text{-}souslin}(s, \operatorname{image}(f, (\mathbb{N} \to \mathbb{N}))) \end{array}}}
\vnbprop{baire-path-facts}{\vnbstmt{\begin{array}{@{}l@{}} \forall s.; \forall x \in (\mathbb{N} \to \operatorname{pts}(s)), z \in (\mathbb{N} \to \mathbb{N}). \\ \quad \operatorname{is\text{-}metric\text{-}space}(s) \Rightarrow  \\ \quad \quad \operatorname{dc\text{-}iter}(\operatorname{pts}(s), x(z(0)), \\ \quad \quad \quad (\lambda {[}\operatorname{mqk\_}, \operatorname{mqu\_}] \in (\mathbb{N} \times \operatorname{pts}(s)).\; \\ \quad \quad \quad \quad \{y \in \operatorname{pts}(s) : (((\operatorname{dist}(s)(\operatorname{mqu\_}, x(z((\operatorname{mqk\_} + 1)))) < \operatorname{recip}(\operatorname{power}(2, \operatorname{mqk\_}))) \wedge (y = x(z((\operatorname{mqk\_} + 1))))) \vee (\neg (\operatorname{dist}(s)(\operatorname{mqu\_}, x(z((\operatorname{mqk\_} + 1)))) < \operatorname{recip}(\operatorname{power}(2, \operatorname{mqk\_}))) \wedge (y = \operatorname{mqu\_})))\}), \\ \quad \quad \quad 0) \\ \quad \quad =\; x(z(0)) \wedge \\ \quad \quad \forall k \in \mathbb{N}. \\ \quad \quad \quad \operatorname{dc\text{-}iter}(\operatorname{pts}(s), x(z(0)), \\ \quad \quad \quad \quad (\lambda {[}\operatorname{mqk\_}, \operatorname{mqu\_}] \in (\mathbb{N} \times \operatorname{pts}(s)).\; \\ \quad \quad \quad \quad \quad \{y \in \operatorname{pts}(s) : (((\operatorname{dist}(s)(\operatorname{mqu\_}, x(z((\operatorname{mqk\_} + 1)))) < \operatorname{recip}(\operatorname{power}(2, \operatorname{mqk\_}))) \wedge (y = x(z((\operatorname{mqk\_} + 1))))) \vee (\neg (\operatorname{dist}(s)(\operatorname{mqu\_}, x(z((\operatorname{mqk\_} + 1)))) < \operatorname{recip}(\operatorname{power}(2, \operatorname{mqk\_}))) \wedge (y = \operatorname{mqu\_})))\}), \\ \quad \quad \quad \quad k) \\ \quad \quad \quad \in\; \operatorname{pts}(s) \wedge \\ \quad \quad \quad \operatorname{dist}(s)(\operatorname{dc\text{-}iter}(\operatorname{pts}(s), x(z(0)), \\ \quad \quad \quad \quad \quad (\lambda {[}\operatorname{mqk\_}, \operatorname{mqu\_}] \in (\mathbb{N} \times \operatorname{pts}(s)).\; \\ \quad \quad \quad \quad \quad \quad \{y \in \operatorname{pts}(s) : (((\operatorname{dist}(s)(\operatorname{mqu\_}, x(z((\operatorname{mqk\_} + 1)))) < \operatorname{recip}(\operatorname{power}(2, \operatorname{mqk\_}))) \wedge (y = x(z((\operatorname{mqk\_} + 1))))) \vee (\neg (\operatorname{dist}(s)(\operatorname{mqu\_}, x(z((\operatorname{mqk\_} + 1)))) < \operatorname{recip}(\operatorname{power}(2, \operatorname{mqk\_}))) \wedge (y = \operatorname{mqu\_})))\}), \\ \quad \quad \quad \quad \quad k), \\ \quad \quad \quad \quad x(z((k + 1)))) \\ \quad \quad \quad <\; \operatorname{recip}(\operatorname{power}(2, k)) \wedge \\ \quad \quad \quad \operatorname{dc\text{-}iter}(\operatorname{pts}(s), x(z(0)), \\ \quad \quad \quad \quad (\lambda {[}\operatorname{mqk\_}, \operatorname{mqu\_}] \in (\mathbb{N} \times \operatorname{pts}(s)).\; \\ \quad \quad \quad \quad \quad \{y \in \operatorname{pts}(s) : (((\operatorname{dist}(s)(\operatorname{mqu\_}, x(z((\operatorname{mqk\_} + 1)))) < \operatorname{recip}(\operatorname{power}(2, \operatorname{mqk\_}))) \wedge (y = x(z((\operatorname{mqk\_} + 1))))) \vee (\neg (\operatorname{dist}(s)(\operatorname{mqu\_}, x(z((\operatorname{mqk\_} + 1)))) < \operatorname{recip}(\operatorname{power}(2, \operatorname{mqk\_}))) \wedge (y = \operatorname{mqu\_})))\}), \\ \quad \quad \quad \quad (k + 1)) \\ \quad \quad \quad =\; x(z((k + 1))) \vee \\ \quad \quad \quad \neg (\operatorname{dist}(s)(\operatorname{dc\text{-}iter}(\operatorname{pts}(s), x(z(0)), (\lambda {[}\operatorname{mqk\_}, \operatorname{mqu\_}] \in (\mathbb{N} \times \operatorname{pts}(s)).\; \{y \in \operatorname{pts}(s) : (((\operatorname{dist}(s)(\operatorname{mqu\_}, x(z((\operatorname{mqk\_} + 1)))) < \operatorname{recip}(\operatorname{power}(2, \operatorname{mqk\_}))) \wedge (y = x(z((\operatorname{mqk\_} + 1))))) \vee (\neg (\operatorname{dist}(s)(\operatorname{mqu\_}, x(z((\operatorname{mqk\_} + 1)))) < \operatorname{recip}(\operatorname{power}(2, \operatorname{mqk\_}))) \wedge (y = \operatorname{mqu\_})))\}), k), x(z((k + 1)))) < \operatorname{recip}(\operatorname{power}(2, k))) \wedge \\ \quad \quad \quad \operatorname{dc\text{-}iter}(\operatorname{pts}(s), x(z(0)), \\ \quad \quad \quad \quad (\lambda {[}\operatorname{mqk\_}, \operatorname{mqu\_}] \in (\mathbb{N} \times \operatorname{pts}(s)).\; \\ \quad \quad \quad \quad \quad \{y \in \operatorname{pts}(s) : (((\operatorname{dist}(s)(\operatorname{mqu\_}, x(z((\operatorname{mqk\_} + 1)))) < \operatorname{recip}(\operatorname{power}(2, \operatorname{mqk\_}))) \wedge (y = x(z((\operatorname{mqk\_} + 1))))) \vee (\neg (\operatorname{dist}(s)(\operatorname{mqu\_}, x(z((\operatorname{mqk\_} + 1)))) < \operatorname{recip}(\operatorname{power}(2, \operatorname{mqk\_}))) \wedge (y = \operatorname{mqu\_})))\}), \\ \quad \quad \quad \quad (k + 1)) \\ \quad \quad \quad =\; \operatorname{dc\text{-}iter}(\operatorname{pts}(s), x(z(0)), \\ \quad \quad \quad \quad (\lambda {[}\operatorname{mqk\_}, \operatorname{mqu\_}] \in (\mathbb{N} \times \operatorname{pts}(s)).\; \\ \quad \quad \quad \quad \quad \{y \in \operatorname{pts}(s) : (((\operatorname{dist}(s)(\operatorname{mqu\_}, x(z((\operatorname{mqk\_} + 1)))) < \operatorname{recip}(\operatorname{power}(2, \operatorname{mqk\_}))) \wedge (y = x(z((\operatorname{mqk\_} + 1))))) \vee (\neg (\operatorname{dist}(s)(\operatorname{mqu\_}, x(z((\operatorname{mqk\_} + 1)))) < \operatorname{recip}(\operatorname{power}(2, \operatorname{mqk\_}))) \wedge (y = \operatorname{mqu\_})))\}), \\ \quad \quad \quad \quad k) \end{array}}}
\vnbprop{baire-path-close}{\vnbstmt{\begin{array}{@{}l@{}} \forall s.; \forall x \in (\mathbb{N} \to \operatorname{pts}(s)), z \in (\mathbb{N} \to \mathbb{N}), k \in \mathbb{N}. \\ \quad \operatorname{is\text{-}metric\text{-}space}(s) \Rightarrow  \\ \quad \quad \operatorname{dist}(s)(\operatorname{dc\text{-}iter}(\operatorname{pts}(s), x(z(0)), \\ \quad \quad \quad \quad (\lambda {[}\operatorname{mqk\_}, \operatorname{mqu\_}] \in (\mathbb{N} \times \operatorname{pts}(s)).\; \\ \quad \quad \quad \quad \quad \{y \in \operatorname{pts}(s) : (((\operatorname{dist}(s)(\operatorname{mqu\_}, x(z((\operatorname{mqk\_} + 1)))) < \operatorname{recip}(\operatorname{power}(2, \operatorname{mqk\_}))) \wedge (y = x(z((\operatorname{mqk\_} + 1))))) \vee (\neg (\operatorname{dist}(s)(\operatorname{mqu\_}, x(z((\operatorname{mqk\_} + 1)))) < \operatorname{recip}(\operatorname{power}(2, \operatorname{mqk\_}))) \wedge (y = \operatorname{mqu\_})))\}), \\ \quad \quad \quad \quad k), \\ \quad \quad \quad \operatorname{dc\text{-}iter}(\operatorname{pts}(s), x(z(0)), \\ \quad \quad \quad \quad (\lambda {[}\operatorname{mqk\_}, \operatorname{mqu\_}] \in (\mathbb{N} \times \operatorname{pts}(s)).\; \\ \quad \quad \quad \quad \quad \{y \in \operatorname{pts}(s) : (((\operatorname{dist}(s)(\operatorname{mqu\_}, x(z((\operatorname{mqk\_} + 1)))) < \operatorname{recip}(\operatorname{power}(2, \operatorname{mqk\_}))) \wedge (y = x(z((\operatorname{mqk\_} + 1))))) \vee (\neg (\operatorname{dist}(s)(\operatorname{mqu\_}, x(z((\operatorname{mqk\_} + 1)))) < \operatorname{recip}(\operatorname{power}(2, \operatorname{mqk\_}))) \wedge (y = \operatorname{mqu\_})))\}), \\ \quad \quad \quad \quad (k + 1))) \\ \quad \quad \leq\; \operatorname{recip}(\operatorname{power}(2, k)) \end{array}}}
\vnbprop{baire-path-cauchy}{\vnbstmt{\begin{array}{@{}l@{}} \forall s.; \forall x \in (\mathbb{N} \to \operatorname{pts}(s)), z \in (\mathbb{N} \to \mathbb{N}), m, j \in \mathbb{N}. \\ \quad \operatorname{is\text{-}metric\text{-}space}(s) \Rightarrow  \\ \quad \quad \operatorname{dist}(s)(\operatorname{dc\text{-}iter}(\operatorname{pts}(s), x(z(0)), \\ \quad \quad \quad \quad (\lambda {[}\operatorname{mqk\_}, \operatorname{mqu\_}] \in (\mathbb{N} \times \operatorname{pts}(s)).\; \\ \quad \quad \quad \quad \quad \{y \in \operatorname{pts}(s) : (((\operatorname{dist}(s)(\operatorname{mqu\_}, x(z((\operatorname{mqk\_} + 1)))) < \operatorname{recip}(\operatorname{power}(2, \operatorname{mqk\_}))) \wedge (y = x(z((\operatorname{mqk\_} + 1))))) \vee (\neg (\operatorname{dist}(s)(\operatorname{mqu\_}, x(z((\operatorname{mqk\_} + 1)))) < \operatorname{recip}(\operatorname{power}(2, \operatorname{mqk\_}))) \wedge (y = \operatorname{mqu\_})))\}), \\ \quad \quad \quad \quad m), \\ \quad \quad \quad \operatorname{dc\text{-}iter}(\operatorname{pts}(s), x(z(0)), \\ \quad \quad \quad \quad (\lambda {[}\operatorname{mqk\_}, \operatorname{mqu\_}] \in (\mathbb{N} \times \operatorname{pts}(s)).\; \\ \quad \quad \quad \quad \quad \{y \in \operatorname{pts}(s) : (((\operatorname{dist}(s)(\operatorname{mqu\_}, x(z((\operatorname{mqk\_} + 1)))) < \operatorname{recip}(\operatorname{power}(2, \operatorname{mqk\_}))) \wedge (y = x(z((\operatorname{mqk\_} + 1))))) \vee (\neg (\operatorname{dist}(s)(\operatorname{mqu\_}, x(z((\operatorname{mqk\_} + 1)))) < \operatorname{recip}(\operatorname{power}(2, \operatorname{mqk\_}))) \wedge (y = \operatorname{mqu\_})))\}), \\ \quad \quad \quad \quad (m + j))) \\ \quad \quad \leq\; ((2 \cdot \operatorname{recip}(\operatorname{power}(2, m))) - (2 \cdot \operatorname{recip}(\operatorname{power}(2, (m + j))))) \end{array}}}
\vnbprop{baire-path-local}{\vnbstmt{\begin{array}{@{}l@{}} \forall s.; \forall x \in (\mathbb{N} \to \operatorname{pts}(s)), z, w \in (\mathbb{N} \to \mathbb{N}), n, k \in \mathbb{N}. \\ \quad (\operatorname{is\text{-}metric\text{-}space}(s) \wedge ((\operatorname{restrict}(w, \mathrm{OS}(n)) = \operatorname{restrict}(z, \mathrm{OS}(n))) \wedge (k < n))) \Rightarrow  \\ \quad \quad \operatorname{dc\text{-}iter}(\operatorname{pts}(s), x(w(0)), \\ \quad \quad \quad (\lambda {[}\operatorname{mqk\_}, \operatorname{mqu\_}] \in (\mathbb{N} \times \operatorname{pts}(s)).\; \\ \quad \quad \quad \quad \{y \in \operatorname{pts}(s) : (((\operatorname{dist}(s)(\operatorname{mqu\_}, x(w((\operatorname{mqk\_} + 1)))) < \operatorname{recip}(\operatorname{power}(2, \operatorname{mqk\_}))) \wedge (y = x(w((\operatorname{mqk\_} + 1))))) \vee (\neg (\operatorname{dist}(s)(\operatorname{mqu\_}, x(w((\operatorname{mqk\_} + 1)))) < \operatorname{recip}(\operatorname{power}(2, \operatorname{mqk\_}))) \wedge (y = \operatorname{mqu\_})))\}), \\ \quad \quad \quad k) \\ \quad \quad =\; \operatorname{dc\text{-}iter}(\operatorname{pts}(s), x(z(0)), \\ \quad \quad \quad (\lambda {[}\operatorname{mqk\_}, \operatorname{mqu\_}] \in (\mathbb{N} \times \operatorname{pts}(s)).\; \\ \quad \quad \quad \quad \{y \in \operatorname{pts}(s) : (((\operatorname{dist}(s)(\operatorname{mqu\_}, x(z((\operatorname{mqk\_} + 1)))) < \operatorname{recip}(\operatorname{power}(2, \operatorname{mqk\_}))) \wedge (y = x(z((\operatorname{mqk\_} + 1))))) \vee (\neg (\operatorname{dist}(s)(\operatorname{mqu\_}, x(z((\operatorname{mqk\_} + 1)))) < \operatorname{recip}(\operatorname{power}(2, \operatorname{mqk\_}))) \wedge (y = \operatorname{mqu\_})))\}), \\ \quad \quad \quad k) \end{array}}}
\vnbprop{rho-small}{\vnbstmt{\forall s.\; (\operatorname{pos\text{-}rr}(s) \Rightarrow \exists k \in \mathbb{N}.\; (\operatorname{recip}(\operatorname{power}(2, k)) < s))}}
\vnbprop{baire-path-bound}{\vnbstmt{\begin{array}{@{}l@{}} \forall s.; \forall x \in (\mathbb{N} \to \operatorname{pts}(s)), z \in (\mathbb{N} \to \mathbb{N}), k, m \in \mathbb{N}. \\ \quad (\operatorname{is\text{-}metric\text{-}space}(s) \wedge (k \leq m)) \Rightarrow  \\ \quad \quad \operatorname{dist}(s)(\operatorname{dc\text{-}iter}(\operatorname{pts}(s), x(z(0)), \\ \quad \quad \quad \quad (\lambda {[}\operatorname{mqk\_}, \operatorname{mqu\_}] \in (\mathbb{N} \times \operatorname{pts}(s)).\; \\ \quad \quad \quad \quad \quad \{y \in \operatorname{pts}(s) : (((\operatorname{dist}(s)(\operatorname{mqu\_}, x(z((\operatorname{mqk\_} + 1)))) < \operatorname{recip}(\operatorname{power}(2, \operatorname{mqk\_}))) \wedge (y = x(z((\operatorname{mqk\_} + 1))))) \vee (\neg (\operatorname{dist}(s)(\operatorname{mqu\_}, x(z((\operatorname{mqk\_} + 1)))) < \operatorname{recip}(\operatorname{power}(2, \operatorname{mqk\_}))) \wedge (y = \operatorname{mqu\_})))\}), \\ \quad \quad \quad \quad k), \\ \quad \quad \quad \operatorname{dc\text{-}iter}(\operatorname{pts}(s), x(z(0)), \\ \quad \quad \quad \quad (\lambda {[}\operatorname{mqk\_}, \operatorname{mqu\_}] \in (\mathbb{N} \times \operatorname{pts}(s)).\; \\ \quad \quad \quad \quad \quad \{y \in \operatorname{pts}(s) : (((\operatorname{dist}(s)(\operatorname{mqu\_}, x(z((\operatorname{mqk\_} + 1)))) < \operatorname{recip}(\operatorname{power}(2, \operatorname{mqk\_}))) \wedge (y = x(z((\operatorname{mqk\_} + 1))))) \vee (\neg (\operatorname{dist}(s)(\operatorname{mqu\_}, x(z((\operatorname{mqk\_} + 1)))) < \operatorname{recip}(\operatorname{power}(2, \operatorname{mqk\_}))) \wedge (y = \operatorname{mqu\_})))\}), \\ \quad \quad \quad \quad m)) \\ \quad \quad \leq\; 2 \cdot \operatorname{recip}(\operatorname{power}(2, k)) \end{array}}}
\vnbprop{baire-path-limit}{\vnbstmt{\begin{array}{@{}l@{}} \forall s.; \forall x \in (\mathbb{N} \to \operatorname{pts}(s)), z \in (\mathbb{N} \to \mathbb{N}). \\ \quad \operatorname{is\text{-}complete}(s) \Rightarrow  \\ \quad \quad \varepsilon(\{l \in \operatorname{pts}(s) : \operatorname{converges\text{-}to}(s, (\lambda n \in \mathbb{N}.\; \operatorname{dc\text{-}iter}(\operatorname{pts}(s), x(z(0)), (\lambda {[}\operatorname{mqk\_}, \operatorname{mqu\_}] \in (\mathbb{N} \times \operatorname{pts}(s)).\; \{y \in \operatorname{pts}(s) : (((\operatorname{dist}(s)(\operatorname{mqu\_}, x(z((\operatorname{mqk\_} + 1)))) < \operatorname{recip}(\operatorname{power}(2, \operatorname{mqk\_}))) \wedge (y = x(z((\operatorname{mqk\_} + 1))))) \vee (\neg (\operatorname{dist}(s)(\operatorname{mqu\_}, x(z((\operatorname{mqk\_} + 1)))) < \operatorname{recip}(\operatorname{power}(2, \operatorname{mqk\_}))) \wedge (y = \operatorname{mqu\_})))\}), n)), l)\}) \\ \quad \quad \in\; \operatorname{pts}(s) \wedge \\ \quad \quad \operatorname{converges\text{-}to}(s, \\ \quad \quad \quad (\lambda n \in \mathbb{N}.\; \\ \quad \quad \quad \quad \operatorname{dc\text{-}iter}(\operatorname{pts}(s), x(z(0)), \\ \quad \quad \quad \quad \quad (\lambda {[}\operatorname{mqk\_}, \operatorname{mqu\_}] \in (\mathbb{N} \times \operatorname{pts}(s)).\; \\ \quad \quad \quad \quad \quad \quad \{y \in \operatorname{pts}(s) : (((\operatorname{dist}(s)(\operatorname{mqu\_}, x(z((\operatorname{mqk\_} + 1)))) < \operatorname{recip}(\operatorname{power}(2, \operatorname{mqk\_}))) \wedge (y = x(z((\operatorname{mqk\_} + 1))))) \vee (\neg (\operatorname{dist}(s)(\operatorname{mqu\_}, x(z((\operatorname{mqk\_} + 1)))) < \operatorname{recip}(\operatorname{power}(2, \operatorname{mqk\_}))) \wedge (y = \operatorname{mqu\_})))\}), \\ \quad \quad \quad \quad \quad n)), \\ \quad \quad \quad \varepsilon(\{l \in \operatorname{pts}(s) : \operatorname{converges\text{-}to}(s, (\lambda n \in \mathbb{N}.\; \operatorname{dc\text{-}iter}(\operatorname{pts}(s), x(z(0)), (\lambda {[}\operatorname{mqk\_}, \operatorname{mqu\_}] \in (\mathbb{N} \times \operatorname{pts}(s)).\; \{y \in \operatorname{pts}(s) : (((\operatorname{dist}(s)(\operatorname{mqu\_}, x(z((\operatorname{mqk\_} + 1)))) < \operatorname{recip}(\operatorname{power}(2, \operatorname{mqk\_}))) \wedge (y = x(z((\operatorname{mqk\_} + 1))))) \vee (\neg (\operatorname{dist}(s)(\operatorname{mqu\_}, x(z((\operatorname{mqk\_} + 1)))) < \operatorname{recip}(\operatorname{power}(2, \operatorname{mqk\_}))) \wedge (y = \operatorname{mqu\_})))\}), n)), l)\})) \wedge \\ \quad \quad \forall m \in \mathbb{N}. \\ \quad \quad \quad \operatorname{dist}(s)(\operatorname{dc\text{-}iter}(\operatorname{pts}(s), x(z(0)), \\ \quad \quad \quad \quad \quad (\lambda {[}\operatorname{mqk\_}, \operatorname{mqu\_}] \in (\mathbb{N} \times \operatorname{pts}(s)).\; \\ \quad \quad \quad \quad \quad \quad \{y \in \operatorname{pts}(s) : (((\operatorname{dist}(s)(\operatorname{mqu\_}, x(z((\operatorname{mqk\_} + 1)))) < \operatorname{recip}(\operatorname{power}(2, \operatorname{mqk\_}))) \wedge (y = x(z((\operatorname{mqk\_} + 1))))) \vee (\neg (\operatorname{dist}(s)(\operatorname{mqu\_}, x(z((\operatorname{mqk\_} + 1)))) < \operatorname{recip}(\operatorname{power}(2, \operatorname{mqk\_}))) \wedge (y = \operatorname{mqu\_})))\}), \\ \quad \quad \quad \quad \quad m), \\ \quad \quad \quad \quad \varepsilon(\{l \in \operatorname{pts}(s) : \operatorname{converges\text{-}to}(s, (\lambda n \in \mathbb{N}.\; \operatorname{dc\text{-}iter}(\operatorname{pts}(s), x(z(0)), (\lambda {[}\operatorname{mqk\_}, \operatorname{mqu\_}] \in (\mathbb{N} \times \operatorname{pts}(s)).\; \{y \in \operatorname{pts}(s) : (((\operatorname{dist}(s)(\operatorname{mqu\_}, x(z((\operatorname{mqk\_} + 1)))) < \operatorname{recip}(\operatorname{power}(2, \operatorname{mqk\_}))) \wedge (y = x(z((\operatorname{mqk\_} + 1))))) \vee (\neg (\operatorname{dist}(s)(\operatorname{mqu\_}, x(z((\operatorname{mqk\_} + 1)))) < \operatorname{recip}(\operatorname{power}(2, \operatorname{mqk\_}))) \wedge (y = \operatorname{mqu\_})))\}), n)), l)\})) \\ \quad \quad \quad \leq\; 2 \cdot \operatorname{recip}(\operatorname{power}(2, m)) \end{array}}}
\vnbprop{polish-baire-surjection}{\vnbstmt{\begin{array}{@{}l@{}} \forall s. \\ \quad (\operatorname{is\text{-}complete}(s) \wedge \operatorname{is\text{-}separable}(s)) \Rightarrow  \\ \quad \quad \exists f.\; (\operatorname{is\text{-}baire\text{-}continuous}(s, f) \wedge (\operatorname{image}(f, (\mathbb{N} \to \mathbb{N})) = \operatorname{pts}(s))) \end{array}}}
\vnbprop{complete-separable-souslin}{\vnbstmt{\forall s.\; ((\operatorname{is\text{-}complete}(s) \wedge \operatorname{is\text{-}separable}(s)) \Rightarrow \operatorname{is\text{-}souslin}(s, \operatorname{pts}(s)))}}
\vnbprop{continuous-image-polish-souslin}{\vnbstmt{\begin{array}{@{}l@{}} \forall s, e, t, h. \\ \quad \operatorname{is\text{-}metric\text{-}space}(s) \wedge \\ \quad e \subseteq \operatorname{pts}(s) \wedge \\ \quad \operatorname{is\text{-}complete}(t) \wedge \\ \quad \operatorname{is\text{-}separable}(t) \wedge \\ \quad \operatorname{is\text{-}continuous}(t, \operatorname{subspace\text{-}ms}(s, e), h) \wedge \\ \quad \operatorname{image}(h, \operatorname{pts}(t)) = e \Rightarrow \\ \quad \quad \operatorname{is\text{-}souslin}(s, e) \end{array}}}
\vnbprop{polish-open-souslin}{\vnbstmt{\begin{array}{@{}l@{}} \forall s, u. \\ \quad (\operatorname{is\text{-}polish}(s) \wedge (\operatorname{is\text{-}open}(s, u) \wedge \exists z.\; (z \in u))) \Rightarrow \operatorname{is\text{-}souslin}(s, u) \end{array}}}
\vnbprop{souslin-baire-image}{\vnbstmt{\begin{array}{@{}l@{}} \forall s, e. \\ \quad (\operatorname{is\text{-}metric\text{-}space}(s) \wedge \operatorname{is\text{-}souslin}(s, e)) \Rightarrow  \\ \quad \quad \exists f.\; (\operatorname{is\text{-}baire\text{-}continuous}(s, f) \wedge (\operatorname{image}(f, (\mathbb{N} \to \mathbb{N})) = e)) \end{array}}}
\vnbprop{souslin-iff-baire-image}{\vnbstmt{\begin{array}{@{}l@{}} \forall s, e. \\ \quad \operatorname{is\text{-}metric\text{-}space}(s) \Rightarrow  \\ \quad \quad \operatorname{is\text{-}souslin}(s, e) \Leftrightarrow \exists f.\; (\operatorname{is\text{-}baire\text{-}continuous}(s, f) \wedge (\operatorname{image}(f, (\mathbb{N} \to \mathbb{N})) = e)) \end{array}}}

\vnbbookfile{theorem-library/mb67-baire}{theorem-library/mb67-baire}
\noindent{\small mb67-baire.scm -- construction of measures, chapter 5: the Baire space N\textasciicircum{}N as a Polish metric space, and the necessity half of Proposition 5.14 (M-B67, 2026-10-26).  Vocabulary: mb67-vocab.scm (BAIRE-MS, SND-ITER, BAIRE-CODE) over scratchpad/mb12 (ULTRA-SEQ-MS and its laws, M-B1B2).  Library: measure. Helper prefixes ...}\par
\vnbprop{baire-recip-lt}{\vnbstmt{\forall u, v \in \mathbb{R}.\; (((0 < u) \wedge (u < v)) \Rightarrow (\operatorname{recip}(v) < \operatorname{recip}(u)))}}
\vnbprop{baire-wt-facts}{\vnbstmt{(((\lambda k \in \mathbb{N}.\; 2) \in (\mathbb{N} \to \mathbb{R})) \wedge \forall w \in \mathbb{N}.\; ((\lambda k \in \mathbb{N}.\; 2)(w) = 2))}}
\vnbprop{baire-wt-gt-one}{\vnbstmt{\forall w \in \mathbb{N}.\; (1 < (\lambda k \in \mathbb{N}.\; 2)(w))}}
\vnbprop{baire-prod-facts}{\vnbstmt{\begin{array}{@{}l@{}} \forall n \in \mathbb{N}, y, a \in (\mathbb{N} \to \mathbb{N}). \\ \quad \operatorname{prefix\text{-}prod}((\lambda k \in \mathbb{N}.\; 2), y, n) \in \mathbb{R} \wedge \\ \quad (n + 1) \leq \operatorname{prefix\text{-}prod}((\lambda k \in \mathbb{N}.\; 2), y, n) \wedge \\ \quad \operatorname{prefix\text{-}prod}((\lambda k \in \mathbb{N}.\; 2), y, n) = \operatorname{prefix\text{-}prod}((\lambda k \in \mathbb{N}.\; 2), a, n) \end{array}}}
\vnbprop{baire-prod-succ}{\vnbstmt{\begin{array}{@{}l@{}} \forall y \in (\mathbb{N} \to \mathbb{N}), n \in \mathbb{N}. \\ \quad \operatorname{prefix\text{-}prod}((\lambda k \in \mathbb{N}.\; 2), y, (n + 1)) = (\operatorname{prefix\text{-}prod}((\lambda k \in \mathbb{N}.\; 2), y, n) \cdot 2) \end{array}}}
\vnbprop{baire-ms-is-metric-space}{\vnbstmt{\operatorname{is\text{-}metric\text{-}space}(\operatorname{baire\text{-}ms})}}
\vnbprop{baire-ms-pts}{\vnbstmt{(\operatorname{pts}(\operatorname{baire\text{-}ms}) \simeq (\mathbb{N} \to \mathbb{N}))}}
\vnbprop{baire-ms-dist}{\vnbstmt{\begin{array}{@{}l@{}} \forall x, y \in (\mathbb{N} \to \mathbb{N}). \\ \quad \operatorname{dist}(\operatorname{baire\text{-}ms})(x, y) = \operatorname{prefix\text{-}dist}((\lambda k \in \mathbb{N}.\; 2), x, y) \wedge \\ \quad \operatorname{prefix\text{-}dist}((\lambda k \in \mathbb{N}.\; 2), x, y) \in \mathbb{R} \wedge \\ \quad 0 \leq \operatorname{prefix\text{-}dist}((\lambda k \in \mathbb{N}.\; 2), x, y) \end{array}}}
\vnbprop{baire-dist-agree-small}{\vnbstmt{\begin{array}{@{}l@{}} \forall x, y \in (\mathbb{N} \to \mathbb{N}), n \in \mathbb{N}. \\ \quad \forall j \in \mathbb{N}.\; ((j < n) \Rightarrow (x(j) = y(j))) \Rightarrow  \\ \quad \quad \operatorname{dist}(\operatorname{baire\text{-}ms})(x, y) \leq \operatorname{recip}((n + 1)) \end{array}}}
\vnbprop{baire-close-agree}{\vnbstmt{\begin{array}{@{}l@{}} \forall c, x, y \in (\mathbb{N} \to \mathbb{N}), n, j \in \mathbb{N}. \\ \quad ((\operatorname{dist}(\operatorname{baire\text{-}ms})(x, y) < \operatorname{recip}(\operatorname{prefix\text{-}prod}((\lambda k \in \mathbb{N}.\; 2), c, n))) \wedge (j \leq n)) \Rightarrow  \\ \quad \quad x(j) = y(j) \end{array}}}
\vnbprop{baire-delta}{\vnbstmt{\begin{array}{@{}l@{}} \forall c \in (\mathbb{N} \to \mathbb{N}), n \in \mathbb{N}. \\ \quad \operatorname{pos\text{-}rr}(\operatorname{recip}(\operatorname{prefix\text{-}prod}((\lambda k \in \mathbb{N}.\; 2), c, (n + 1)))) \wedge \\ \quad \operatorname{recip}(\operatorname{prefix\text{-}prod}((\lambda k \in \mathbb{N}.\; 2), c, (n + 1))) \\ \quad <\; \operatorname{recip}(\operatorname{prefix\text{-}prod}((\lambda k \in \mathbb{N}.\; 2), c, n)) \end{array}}}
\vnbprop{baire-div}{\vnbstmt{\begin{array}{@{}l@{}} \forall z \in (\mathbb{N} \to \mathbb{N}), m \in \mathbb{R}. \exists n \in \mathbb{N}. \\ \quad m < \operatorname{prefix\text{-}prod}((\lambda k \in \mathbb{N}.\; 2), z, n) \end{array}}}
\vnbprop{baire-cauchy-converges}{\vnbstmt{\forall s.\; (\operatorname{is\text{-}cauchy\text{-}seq}(\operatorname{baire\text{-}ms}, s) \Rightarrow \operatorname{converges}(\operatorname{baire\text{-}ms}, s))}}
\vnbprop{baire-ms-complete}{\vnbstmt{\operatorname{is\text{-}complete}(\operatorname{baire\text{-}ms})}}
\vnbprop{snd-iter-type}{\vnbstmt{\forall k, m \in \mathbb{N}.\; (\operatorname{snd\text{-}iter}(m, k) \in \mathbb{N})}}
\vnbprop{snd-iter-shift}{\vnbstmt{\forall k, m \in \mathbb{N}.\; (\operatorname{snd\text{-}iter}(m, (k + 1)) = \operatorname{snd\text{-}iter}(\operatorname{nnsnd}(m), k))}}
\vnbprop{baire-code-facts}{\vnbstmt{\begin{array}{@{}l@{}} \forall m \in \mathbb{N}. \\ \quad ((\operatorname{baire\text{-}code}(m) \in (\mathbb{N} \to \mathbb{N})) \wedge \forall k \in \mathbb{N}.\; (\operatorname{baire\text{-}code}(m)(k) = \operatorname{nnfst}(\operatorname{snd\text{-}iter}(m, k)))) \end{array}}}
\vnbprop{baire-code-prefix}{\vnbstmt{\begin{array}{@{}l@{}} \forall n \in \mathbb{N}, z \in (\mathbb{N} \to \mathbb{N}). \exists m \in \mathbb{N}. \forall j \in \mathbb{N}. \\ \quad (j < n) \Rightarrow (\operatorname{nnfst}(\operatorname{snd\text{-}iter}(m, j)) = z(j)) \end{array}}}
\vnbprop{baire-ms-separable}{\vnbstmt{\operatorname{is\text{-}separable}(\operatorname{baire\text{-}ms})}}
\vnbprop{baire-ms-polish}{\vnbstmt{\operatorname{is\text{-}polish}(\operatorname{baire\text{-}ms})}}
\vnbprop{seq-tail-facts}{\vnbstmt{\begin{array}{@{}l@{}} \forall z \in (\mathbb{N} \to \mathbb{N}). \\ \quad (\operatorname{seq\text{-}tail}(z) \in (\mathbb{N} \to \mathbb{N})) \wedge \forall j \in \mathbb{N}.\; (\operatorname{seq\text{-}tail}(z)(j) = z((j + 1))) \end{array}}}
\vnbprop{seq-cons-facts}{\vnbstmt{\begin{array}{@{}l@{}} \forall k \in \mathbb{N}, v \in (\mathbb{N} \to \mathbb{N}). \\ \quad ((\operatorname{seq\text{-}cons}(k, v) \in (\mathbb{N} \to \mathbb{N})) \wedge ((\operatorname{seq\text{-}cons}(k, v)(0) = k) \wedge (\operatorname{seq\text{-}tail}(\operatorname{seq\text{-}cons}(k, v)) = v))) \end{array}}}
\vnbprop{seq-proj-facts}{\vnbstmt{\begin{array}{@{}l@{}} \forall k \in \mathbb{N}, z \in (\mathbb{N} \to \mathbb{N}). \\ \quad ((\operatorname{seq\text{-}proj}(k, z) \in (\mathbb{N} \to \mathbb{N})) \wedge \forall j \in \mathbb{N}.\; (\operatorname{seq\text{-}proj}(k, z)(j) = z(\operatorname{nnpair}(k, j)))) \end{array}}}
\vnbprop{nnpair-lt-right}{\vnbstmt{\forall k, j, n \in \mathbb{N}.\; ((j < n) \Rightarrow (\operatorname{nnpair}(k, j) < \operatorname{nnpair}(k, n)))}}
\vnbprop{seq-unproj}{\vnbstmt{\begin{array}{@{}l@{}} \forall v \in (\mathbb{N} \to (\mathbb{N} \to \mathbb{N})). \exists z \in (\mathbb{N} \to \mathbb{N}). \forall k \in \mathbb{N}. \\ \quad \operatorname{seq\text{-}proj}(k, z) = v(k) \end{array}}}

\vnbbookfile{theorem-library/mb67-souslin}{theorem-library/mb67-souslin}
\noindent{\small mb67-souslin.scm -- construction of measures, chapter 5 sections 1.4-1.4.1 (M-B67, 2026-10-26): the necessity half of Proposition 5.14 (so 5.14 as an IFF), Corollary 5.15, Proposition 5.16, Corollary 5.17. Vocabulary: mb35-souslin-vocab.scm (IS-SOUSLIN, IS-BAIRE-CONTINUOUS), mb67-vocab.scm (BAIRE-MS, SEQ-TAIL, ...}\par
\vnbprop{baire-continuous-subspace}{\vnbstmt{\begin{array}{@{}l@{}} \forall s, e, f. \\ \quad (\operatorname{is\text{-}metric\text{-}space}(s) \wedge ((e \subseteq \operatorname{pts}(s)) \wedge (\operatorname{is\text{-}baire\text{-}continuous}(s, f) \wedge (\operatorname{image}(f, (\mathbb{N} \to \mathbb{N})) = e)))) \Rightarrow  \\ \quad \quad \operatorname{is\text{-}continuous}(\operatorname{baire\text{-}ms}, \operatorname{subspace\text{-}ms}(s, e), f) \end{array}}}
\vnbprop{souslin-polish-image}{\vnbstmt{\begin{array}{@{}l@{}} \forall s, e. \\ \quad (\operatorname{is\text{-}metric\text{-}space}(s) \wedge \operatorname{is\text{-}souslin}(s, e)) \Rightarrow  \\ \quad \quad \exists t. \\ \quad \quad \quad \operatorname{is\text{-}complete}(t) \wedge \\ \quad \quad \quad \operatorname{is\text{-}separable}(t) \wedge \\ \quad \quad \quad \exists h.\; (\operatorname{is\text{-}continuous}(t, \operatorname{subspace\text{-}ms}(s, e), h) \wedge (\operatorname{image}(h, \operatorname{pts}(t)) = e)) \end{array}}}
\vnbprop{souslin-iff-polish-image}{\vnbstmt{\begin{array}{@{}l@{}} \forall s, e. \\ \quad (\operatorname{is\text{-}metric\text{-}space}(s) \wedge (e \subseteq \operatorname{pts}(s))) \Rightarrow  \\ \quad \quad \operatorname{is\text{-}souslin}(s, e) \\ \quad \quad \Leftrightarrow\; \exists t. \\ \quad \quad \quad \operatorname{is\text{-}complete}(t) \wedge \\ \quad \quad \quad \operatorname{is\text{-}separable}(t) \wedge \\ \quad \quad \quad \exists h.\; (\operatorname{is\text{-}continuous}(t, \operatorname{subspace\text{-}ms}(s, e), h) \wedge (\operatorname{image}(h, \operatorname{pts}(t)) = e)) \end{array}}}
\vnbprop{souslin-continuous-image}{\vnbstmt{\begin{array}{@{}l@{}} \forall s, t, i, e. \\ \quad (\operatorname{is\text{-}continuous}(s, t, i) \wedge \operatorname{is\text{-}souslin}(s, e)) \Rightarrow \operatorname{is\text{-}souslin}(t, \operatorname{image}(i, e)) \end{array}}}
\vnbprop{souslin-countable-union}{\vnbstmt{\begin{array}{@{}l@{}} \forall s.; \forall e \in (\mathbb{N} \to \operatorname{power}(\operatorname{pts}(s))). \\ \quad (\operatorname{is\text{-}metric\text{-}space}(s) \wedge \forall k \in \mathbb{N}.\; \operatorname{is\text{-}souslin}(s, e(k))) \Rightarrow  \\ \quad \quad \operatorname{is\text{-}souslin}(s, \operatorname{big\text{-}union}(\mathit{mxk}, \mathbb{N}, e(\mathit{mxk}))) \end{array}}}
\vnbprop{souslin-family-choice}{\vnbstmt{\begin{array}{@{}l@{}} \forall s.; \forall e \in (\mathbb{N} \to \operatorname{power}(\operatorname{pts}(s))). \\ \quad (\operatorname{is\text{-}metric\text{-}space}(s) \wedge \forall k \in \mathbb{N}.\; \operatorname{is\text{-}souslin}(s, e(k))) \Rightarrow  \\ \quad \quad \exists g \in (\mathbb{N} \to ((\mathbb{N} \to \mathbb{N}) \to \operatorname{pts}(s))). \forall k \in \mathbb{N}. \\ \quad \quad \quad \operatorname{is\text{-}baire\text{-}continuous}(s, g(k)) \wedge (\operatorname{image}(g(k), (\mathbb{N} \to \mathbb{N})) = e(k)) \end{array}}}
\vnbprop{seq-proj-prefix}{\vnbstmt{\begin{array}{@{}l@{}} \forall k \in \mathbb{N}, z, w \in (\mathbb{N} \to \mathbb{N}), n \in \mathbb{N}. \\ \quad (\operatorname{restrict}(w, \mathrm{OS}(\operatorname{nnpair}(k, n))) = \operatorname{restrict}(z, \mathrm{OS}(\operatorname{nnpair}(k, n)))) \Rightarrow  \\ \quad \quad \operatorname{restrict}(\operatorname{seq\text{-}proj}(k, w), \mathrm{OS}(n)) = \operatorname{restrict}(\operatorname{seq\text{-}proj}(k, z), \mathrm{OS}(n)) \end{array}}}
\vnbprop{souslin-countable-intersection}{\vnbstmt{\begin{array}{@{}l@{}} \forall s.; \forall e \in (\mathbb{N} \to \operatorname{power}(\operatorname{pts}(s))). \\ \quad \operatorname{is\text{-}metric\text{-}space}(s) \wedge \\ \quad \forall k \in \mathbb{N}.\; \operatorname{is\text{-}souslin}(s, e(k)) \wedge \\ \quad \exists x.\; (x \in \{p \in \operatorname{pts}(s) : \forall \mathit{mxk} \in \mathbb{N}.\; (p \in e(\mathit{mxk}))\}) \Rightarrow \\ \quad \quad \operatorname{is\text{-}souslin}(s, \{p \in \operatorname{pts}(s) : \forall \mathit{mxk} \in \mathbb{N}.\; (p \in e(\mathit{mxk}))\}) \end{array}}}
\vnbprop{closed-polish-souslin}{\vnbstmt{\begin{array}{@{}l@{}} \forall s, f. \\ \quad (\operatorname{is\text{-}polish}(s) \wedge (\operatorname{is\text{-}closed}(s, f) \wedge \exists x.\; (x \in f))) \Rightarrow  \\ \quad \quad \operatorname{is\text{-}souslin}(s, f) \end{array}}}
\vnbprop{souslin-or-empty-union}{\vnbstmt{\begin{array}{@{}l@{}} \forall s.; \forall f \in (\mathbb{N} \to \operatorname{power}(\operatorname{pts}(s))). \\ \quad \operatorname{is\text{-}metric\text{-}space}(s) \wedge \\ \quad \forall k \in \mathbb{N}.\; (\exists x.\; (x \in f(k)) \Rightarrow \operatorname{is\text{-}souslin}(s, f(k))) \wedge \\ \quad \exists x.\; (x \in \operatorname{big\text{-}union}(n, \mathbb{N}, f(n))) \Rightarrow \\ \quad \quad \operatorname{is\text{-}souslin}(s, \operatorname{big\text{-}union}(n, \mathbb{N}, f(n))) \end{array}}}
\vnbprop{souslin-or-empty-co-union}{\vnbstmt{\begin{array}{@{}l@{}} \forall s.; \forall f \in (\mathbb{N} \to \operatorname{power}(\operatorname{pts}(s))). \\ \quad \operatorname{is\text{-}metric\text{-}space}(s) \wedge \\ \quad \forall k \in \mathbb{N}. \\ \quad \quad \exists x.\; (x \in \operatorname{complement\text{-}in}(\operatorname{pts}(s), f(k))) \Rightarrow  \\ \quad \quad \quad \operatorname{is\text{-}souslin}(s, \operatorname{complement\text{-}in}(\operatorname{pts}(s), f(k))) \wedge \\ \quad \exists x.\; (x \in \operatorname{complement\text{-}in}(\operatorname{pts}(s), \operatorname{big\text{-}union}(n, \mathbb{N}, f(n)))) \Rightarrow \\ \quad \quad \operatorname{is\text{-}souslin}(s, \operatorname{complement\text{-}in}(\operatorname{pts}(s), \operatorname{big\text{-}union}(n, \mathbb{N}, f(n)))) \end{array}}}
\vnbprop{borel-polish-souslin}{\vnbstmt{\begin{array}{@{}l@{}} \forall s.; \forall b \in \operatorname{borel\text{-}sets}(s). \\ \quad (\operatorname{is\text{-}polish}(s) \wedge \exists x.\; (x \in b)) \Rightarrow \operatorname{is\text{-}souslin}(s, b) \end{array}}}

\vnbbookfile{theorem-library/mb67-lusin}{theorem-library/mb67-lusin}
\noindent{\small mb67-lusin.scm -- construction of measures, chapter 5 section 1.5 (M-B67, 2026-10-26): Lusin sets. lusin-is-souslin           (after Definition 5.19) a Lusin set is Souslin lusin-closed, lusin-open   nonempty closed and nonempty open subsets of a Polish space are Lusin (the base cases of Corollary 5.21)}\par
\vnbprop{lusin-is-souslin}{\vnbstmt{\begin{array}{@{}l@{}} \forall s, e. \\ \quad (\operatorname{is\text{-}metric\text{-}space}(s) \wedge ((e \subseteq \operatorname{pts}(s)) \wedge \operatorname{is\text{-}lusin}(s, e))) \Rightarrow  \\ \quad \quad \operatorname{is\text{-}souslin}(s, e) \end{array}}}
\vnbprop{lusin-closed}{\vnbstmt{\begin{array}{@{}l@{}} \forall s, f. \\ \quad (\operatorname{is\text{-}polish}(s) \wedge (\operatorname{is\text{-}closed}(s, f) \wedge \exists x.\; (x \in f))) \Rightarrow \operatorname{is\text{-}lusin}(s, f) \end{array}}}
\vnbprop{lusin-open}{\vnbstmt{\begin{array}{@{}l@{}} \forall s, u. \\ \quad (\operatorname{is\text{-}polish}(s) \wedge (\operatorname{is\text{-}open}(s, u) \wedge \exists z.\; (z \in u))) \Rightarrow \operatorname{is\text{-}lusin}(s, u) \end{array}}}
\vnbprop{borel-from-inter-disjoint-union}{\vnbstmt{\begin{array}{@{}l@{}} \forall s, b. \\ \quad \operatorname{is\text{-}metric\text{-}space}(s) \wedge \\ \quad b \subseteq \operatorname{power}(\operatorname{pts}(s)) \wedge \\ \quad \forall g \in (\mathbb{N} \to b).\; (\{q \in \operatorname{pts}(s) : \forall n \in \mathbb{N}.\; (q \in g(n))\} \in b) \wedge \\ \quad \forall g \in (\mathbb{N} \to b). \\ \quad \quad \forall i, j \in \mathbb{N}, x \in g(i).\; ((x \in g(j)) \Rightarrow (i = j)) \Rightarrow  \\ \quad \quad \quad \operatorname{big\text{-}union}(a, \mathbb{N}, g(a)) \in b \wedge \\ \quad \forall u.\; (\operatorname{is\text{-}open}(s, u) \Rightarrow (u \in b)) \Rightarrow \\ \quad \quad \operatorname{borel\text{-}sets}(s) \subseteq b \end{array}}}

\vnbbookfile{theorem-library/mb67-standard}{theorem-library/mb67-standard}
\noindent{\small mb67-standard.scm -- construction of measures, chapter 5 sections 1.7-1.8 (M-B67, 2026-10-26): countably separated measurable spaces. graph-measurable                 Lemma 5.33 countably-separated-of-injection Proposition 5.36, sufficiency}\par
\vnbprop{graph-measurable}{\vnbstmt{\begin{array}{@{}l@{}} \forall s, a, b, c, i. \\ \quad (\operatorname{is\text{-}measurable\text{-}map}(s, a, b, c, i) \wedge \operatorname{is\text{-}countably\text{-}separated}(b, c)) \Rightarrow  \\ \quad \quad \operatorname{graph\text{-}of}(s, b, i) \in \operatorname{product\text{-}sigma}(s, a, b, c) \end{array}}}
\vnbprop{countably-separated-of-injection}{\vnbstmt{\begin{array}{@{}l@{}} \forall s, b, z, i. \\ \quad \operatorname{is\text{-}separable}(z) \wedge \\ \quad \operatorname{is\text{-}measurable\text{-}map}(s, b, \operatorname{pts}(z), \operatorname{borel\text{-}sets}(z), i) \wedge \\ \quad \forall x, y \in s.\; ((i(x) = i(y)) \Rightarrow (x = y)) \Rightarrow \\ \quad \quad \operatorname{is\text{-}countably\text{-}separated}(s, b) \end{array}}}
\vnbprop{polish-standard-borel}{\vnbstmt{\forall s.\; (\operatorname{is\text{-}polish}(s) \Rightarrow \operatorname{is\text{-}standard\text{-}borel}(\operatorname{pts}(s), \operatorname{borel\text{-}sets}(s)))}}
\vnbprop{measurable-map-of-generators}{\vnbstmt{\begin{array}{@{}l@{}} \forall s, a, b, c, f. \\ \quad \operatorname{is\text{-}sigma\text{-}algebra}(s, a) \wedge \\ \quad b \in \mathbf{Set} \wedge \\ \quad c \subseteq \operatorname{power}(b) \wedge \\ \quad f \in (s \to b) \wedge \\ \quad \forall e \in c.\; (\{x \in s : (f(x) \in e)\} \in a) \Rightarrow \\ \quad \quad \operatorname{is\text{-}measurable\text{-}map}(s, a, b, \operatorname{sigma\text{-}generated}(b, c), f) \end{array}}}
\vnbprop{countably-separated-embeds-baire}{\vnbstmt{\begin{array}{@{}l@{}} \forall s, b. \\ \quad \operatorname{is\text{-}countably\text{-}separated}(s, b) \Rightarrow  \\ \quad \quad \exists i. \\ \quad \quad \quad \operatorname{is\text{-}measurable\text{-}map}(s, b, \operatorname{pts}(\operatorname{baire\text{-}ms}), \operatorname{borel\text{-}sets}(\operatorname{baire\text{-}ms}), i) \wedge \\ \quad \quad \quad \forall x, y \in s.\; ((i(x) = i(y)) \Rightarrow (x = y)) \end{array}}}
\vnbprop{countably-separated-embeds}{\vnbstmt{\begin{array}{@{}l@{}} \forall s, b. \\ \quad \operatorname{is\text{-}countably\text{-}separated}(s, b) \Rightarrow  \\ \quad \quad \exists z. \\ \quad \quad \quad \operatorname{is\text{-}separable}(z) \wedge \\ \quad \quad \quad \exists i. \\ \quad \quad \quad \quad \operatorname{is\text{-}measurable\text{-}map}(s, b, \operatorname{pts}(z), \operatorname{borel\text{-}sets}(z), i) \wedge \\ \quad \quad \quad \quad \forall x, y \in s.\; ((i(x) = i(y)) \Rightarrow (x = y)) \end{array}}}
\vnbprop{countably-separated-iff}{\vnbstmt{\begin{array}{@{}l@{}} \forall s, b. \\ \quad \operatorname{is\text{-}countably\text{-}separated}(s, b) \\ \quad \Leftrightarrow\; \exists z. \\ \quad \quad \operatorname{is\text{-}separable}(z) \wedge \\ \quad \quad \exists i. \\ \quad \quad \quad \operatorname{is\text{-}measurable\text{-}map}(s, b, \operatorname{pts}(z), \operatorname{borel\text{-}sets}(z), i) \wedge \\ \quad \quad \quad \forall x, y \in s.\; ((i(x) = i(y)) \Rightarrow (x = y)) \end{array}}}

\vnbbookfile{theorem-library/mb67-prod2}{theorem-library/mb67-prod2}
\noindent{\small mb67-prod2.scm -- the product of two Polish spaces is Polish (M-B67, 2026-10-26); X x N\textasciicircum{}N Polish (for 5.38). product2-complete, product2-separable, product2-polish, product2-baire-polish Library: measure.  Helper prefixes mr- (copied) and mpq-. Load window: after mb67-baire (baire-ms-polish), mb4-vocab (IS-POLISH), the ...}\par
\vnbprop{product2-complete}{\vnbstmt{\begin{array}{@{}l@{}} \forall s, t. \\ \quad (\operatorname{is\text{-}complete}(s) \wedge \operatorname{is\text{-}complete}(t)) \Rightarrow \operatorname{is\text{-}complete}(\operatorname{product2\text{-}ms}(s, t)) \end{array}}}
\vnbprop{product2-separable}{\vnbstmt{\begin{array}{@{}l@{}} \forall s, t. \\ \quad (\operatorname{is\text{-}separable}(s) \wedge \operatorname{is\text{-}separable}(t)) \Rightarrow \operatorname{is\text{-}separable}(\operatorname{product2\text{-}ms}(s, t)) \end{array}}}
\vnbprop{product2-polish}{\vnbstmt{\begin{array}{@{}l@{}} \forall s, t. \\ \quad (\operatorname{is\text{-}polish}(s) \wedge \operatorname{is\text{-}polish}(t)) \Rightarrow \operatorname{is\text{-}polish}(\operatorname{product2\text{-}ms}(s, t)) \end{array}}}
\vnbprop{product2-baire-polish}{\vnbstmt{\forall s.\; (\operatorname{is\text{-}polish}(s) \Rightarrow \operatorname{is\text{-}polish}(\operatorname{product2\text{-}ms}(s, \operatorname{baire\text{-}ms})))}}

\vnbbookfile{theorem-library/mb67-prod}{theorem-library/mb67-prod}
\noindent{\small mb67-prod.scm -- countable products of metric spaces (M-B67, 2026-10-26): PRODUCT-METRIC-W of a sequence of metric spaces with a summable weight preserves completeness, separability, hence Polishness. product-coord-le, product-coord-cauchy, product-complete, product-separable, product-polish Library: measure.  Helper ...}\par
\vnbprop{product-coord-le}{\vnbstmt{\begin{array}{@{}l@{}} \forall s, w.; \forall x, y \in \operatorname{product\text{-}carrier}(s), n \in \mathbb{N}. \\ \quad (\operatorname{is\text{-}ms\text{-}sequence}(s) \wedge \operatorname{summable\text{-}weight}(w)) \Rightarrow  \\ \quad \quad (w(n) \cdot \operatorname{dist}(\operatorname{bdd\text{-}metric}(s(n)))(x(n), y(n))) \leq \operatorname{dist}(\operatorname{product\text{-}metric\text{-}w}(s, w))(x, y) \end{array}}}
\vnbprop{product-coord-cauchy}{\vnbstmt{\begin{array}{@{}l@{}} \forall s, w, u.; \forall n \in \mathbb{N}. \\ \quad (\operatorname{is\text{-}ms\text{-}sequence}(s) \wedge (\operatorname{summable\text{-}weight}(w) \wedge \operatorname{is\text{-}cauchy\text{-}seq}(\operatorname{product\text{-}metric\text{-}w}(s, w), u))) \Rightarrow  \\ \quad \quad \operatorname{is\text{-}cauchy\text{-}seq}(s(n), (\lambda j \in \mathbb{N}.\; u(j)(n))) \end{array}}}
\vnbprop{product-complete}{\vnbstmt{\begin{array}{@{}l@{}} \forall s, w. \\ \quad (\operatorname{is\text{-}ms\text{-}sequence}(s) \wedge (\operatorname{summable\text{-}weight}(w) \wedge \forall n \in \mathbb{N}.\; \operatorname{is\text{-}complete}(s(n)))) \Rightarrow  \\ \quad \quad \operatorname{is\text{-}complete}(\operatorname{product\text{-}metric\text{-}w}(s, w)) \end{array}}}
\vnbprop{product-dist-tail-bound}{\vnbstmt{\begin{array}{@{}l@{}} \forall s, w.; \forall x, q \in \operatorname{product\text{-}carrier}(s), t \in \mathbb{N}, l, c \in \mathbb{R}. \\ \quad \operatorname{is\text{-}ms\text{-}sequence}(s) \wedge \\ \quad \operatorname{summable\text{-}weight}(w) \wedge \\ \quad 0 \leq l \wedge \\ \quad l \leq 1 \wedge \\ \quad \forall k \in \mathbb{N}.\; ((k < t) \Rightarrow (\operatorname{dist}(\operatorname{bdd\text{-}metric}(s(k)))(x(k), q(k)) \leq l)) \wedge \\ \quad \forall k \in \mathbb{N}. \\ \quad \quad (t \leq k) \Rightarrow ((\operatorname{series\text{-}partial\text{-}sum}(w, k) - \operatorname{series\text{-}partial\text{-}sum}(w, t)) \leq c) \Rightarrow \\ \quad \quad \operatorname{dist}(\operatorname{product\text{-}metric\text{-}w}(s, w))(x, q) \leq ((l \cdot \operatorname{series\text{-}partial\text{-}sum}(w, t)) + c) \end{array}}}
\vnbprop{product-separable}{\vnbstmt{\begin{array}{@{}l@{}} \forall s, w. \\ \quad (\operatorname{is\text{-}ms\text{-}sequence}(s) \wedge (\operatorname{summable\text{-}weight}(w) \wedge \forall n \in \mathbb{N}.\; \operatorname{is\text{-}separable}(s(n)))) \Rightarrow  \\ \quad \quad \operatorname{is\text{-}separable}(\operatorname{product\text{-}metric\text{-}w}(s, w)) \end{array}}}
\vnbprop{product-polish}{\vnbstmt{\begin{array}{@{}l@{}} \forall s, w. \\ \quad (\operatorname{is\text{-}ms\text{-}sequence}(s) \wedge (\operatorname{summable\text{-}weight}(w) \wedge \forall n \in \mathbb{N}.\; \operatorname{is\text{-}polish}(s(n)))) \Rightarrow  \\ \quad \quad \operatorname{is\text{-}polish}(\operatorname{product\text{-}metric\text{-}w}(s, w)) \end{array}}}

\vnbbookfile{theorem-library/mb67-sum}{theorem-library/mb67-sum}
\noindent{\small mb67-sum.scm -- countable disjoint sums of metric spaces (M-B67, 2026-10-26): SUM-MS(ts) is a metric space, complete when every summand is, separable when every summand is; hence Polish. sum-pts-unfold, sum-pts-parts, sum-pts-intro, sum-ms-pts, sum-ms-dist-slot, sum-dist-value, sum-pts-is-set, sum-ms-is-metric-space, ...}\par
\vnbprop{sum-pts-unfold}{\vnbstmt{\begin{array}{@{}l@{}} \forall s. \\ \quad (\operatorname{sum\text{-}pts}(s) \simeq \{p \in (\mathbb{N} \times \operatorname{big\text{-}union}(k, \mathbb{N}, \operatorname{pts}(s(k)))) : (\operatorname{nth}(2, p) \in \operatorname{pts}(s(\operatorname{nth}(1, p))))\}) \end{array}}}
\vnbprop{sum-pts-parts}{\vnbstmt{\begin{array}{@{}l@{}} \forall s.; \forall p \in \operatorname{sum\text{-}pts}(s). \\ \quad ((\operatorname{nth}(1, p) \in \mathbb{N}) \wedge ((\operatorname{nth}(2, p) \in \operatorname{pts}(s(\operatorname{nth}(1, p)))) \wedge (p = {[}\operatorname{nth}(1, p), \operatorname{nth}(2, p)]))) \end{array}}}
\vnbprop{sum-pts-intro}{\vnbstmt{\forall s.; \forall k \in \mathbb{N}, x \in \operatorname{pts}(s(k)).\; ({[}k, x] \in \operatorname{sum\text{-}pts}(s))}}
\vnbprop{sum-ms-pts}{\vnbstmt{\forall s.\; (\operatorname{pts}(\operatorname{sum\text{-}ms}(s)) \simeq \operatorname{sum\text{-}pts}(s))}}
\vnbprop{sum-ms-dist-slot}{\vnbstmt{\begin{array}{@{}l@{}} \forall s. \\ \quad (\operatorname{dist}(\operatorname{sum\text{-}ms}(s)) \simeq (\lambda {[}\operatorname{smu\_}, \operatorname{smv\_}] \in (\operatorname{sum\text{-}pts}(s) \times \operatorname{sum\text{-}pts}(s)).\; \left(\text{if }(\operatorname{nth}(1, \operatorname{smu\_}) = \operatorname{nth}(1, \operatorname{smv\_}))\text{ then }\operatorname{dist}(\operatorname{bdd\text{-}metric}(s(\operatorname{nth}(1, \operatorname{smu\_}))))(\operatorname{nth}(2, \operatorname{smu\_}), \operatorname{nth}(2, \operatorname{smv\_}))\text{ else }1\right))) \end{array}}}
\vnbprop{sum-dist-value}{\vnbstmt{\begin{array}{@{}l@{}} \forall s.; \forall u, v \in \operatorname{sum\text{-}pts}(s). \\ \quad (\operatorname{dist}(\operatorname{sum\text{-}ms}(s))(u, v) \simeq \left(\text{if }(\operatorname{nth}(1, u) = \operatorname{nth}(1, v))\text{ then }\operatorname{dist}(\operatorname{bdd\text{-}metric}(s(\operatorname{nth}(1, u))))(\operatorname{nth}(2, u), \operatorname{nth}(2, v))\text{ else }1\right)) \end{array}}}
\vnbprop{sum-pts-is-set}{\vnbstmt{\forall s.\; (\operatorname{is\text{-}ms\text{-}sequence}(s) \Rightarrow (\operatorname{sum\text{-}pts}(s) \in \mathbf{Set}))}}
\vnbprop{sum-dist-same}{\vnbstmt{\begin{array}{@{}l@{}} \forall s.; \forall u, v \in \operatorname{sum\text{-}pts}(s). \\ \quad (\operatorname{is\text{-}ms\text{-}sequence}(s) \wedge (\operatorname{nth}(1, u) = \operatorname{nth}(1, v))) \Rightarrow  \\ \quad \quad \operatorname{dist}(\operatorname{sum\text{-}ms}(s))(u, v) = \operatorname{dist}(\operatorname{bdd\text{-}metric}(s(\operatorname{nth}(1, u))))(\operatorname{nth}(2, u), \operatorname{nth}(2, v)) \end{array}}}
\vnbprop{sum-dist-diff}{\vnbstmt{\begin{array}{@{}l@{}} \forall s.; \forall u, v \in \operatorname{sum\text{-}pts}(s). \\ \quad (\operatorname{is\text{-}ms\text{-}sequence}(s) \wedge \neg (\operatorname{nth}(1, u) = \operatorname{nth}(1, v))) \Rightarrow  \\ \quad \quad \operatorname{dist}(\operatorname{sum\text{-}ms}(s))(u, v) = 1 \end{array}}}
\vnbprop{sum-ms-lam-type}{\vnbstmt{\begin{array}{@{}l@{}} \forall s. \\ \quad \operatorname{is\text{-}ms\text{-}sequence}(s) \Rightarrow  \\ \quad \quad (\lambda {[}\operatorname{smu\_}, \operatorname{smv\_}] \in (\operatorname{sum\text{-}pts}(s) \times \operatorname{sum\text{-}pts}(s)).\; \\ \quad \quad \quad \left(\text{if }(\operatorname{nth}(1, \operatorname{smu\_}) = \operatorname{nth}(1, \operatorname{smv\_}))\text{ then }\operatorname{dist}(\operatorname{bdd\text{-}metric}(s(\operatorname{nth}(1, \operatorname{smu\_}))))(\operatorname{nth}(2, \operatorname{smu\_}), \operatorname{nth}(2, \operatorname{smv\_}))\text{ else }1\right)) \\ \quad \quad \in\; ((\operatorname{sum\text{-}pts}(s) \times \operatorname{sum\text{-}pts}(s)) \to \mathbb{R}) \end{array}}}
\vnbprop{sum-ms-is-metric-space}{\vnbstmt{\forall s.\; (\operatorname{is\text{-}ms\text{-}sequence}(s) \Rightarrow \operatorname{is\text{-}metric\text{-}space}(\operatorname{sum\text{-}ms}(s)))}}
\vnbprop{sum-ms-complete}{\vnbstmt{\begin{array}{@{}l@{}} \forall s. \\ \quad (\operatorname{is\text{-}ms\text{-}sequence}(s) \wedge \forall n \in \mathbb{N}.\; \operatorname{is\text{-}complete}(s(n))) \Rightarrow  \\ \quad \quad \operatorname{is\text{-}complete}(\operatorname{sum\text{-}ms}(s)) \end{array}}}
\vnbprop{sum-ms-separable}{\vnbstmt{\begin{array}{@{}l@{}} \forall s. \\ \quad (\operatorname{is\text{-}ms\text{-}sequence}(s) \wedge \forall n \in \mathbb{N}.\; \operatorname{is\text{-}separable}(s(n))) \Rightarrow  \\ \quad \quad \operatorname{is\text{-}separable}(\operatorname{sum\text{-}ms}(s)) \end{array}}}
\vnbprop{sum-ms-polish}{\vnbstmt{\begin{array}{@{}l@{}} \forall s. \\ \quad (\operatorname{is\text{-}ms\text{-}sequence}(s) \wedge \forall n \in \mathbb{N}.\; \operatorname{is\text{-}polish}(s(n))) \Rightarrow  \\ \quad \quad \operatorname{is\text{-}polish}(\operatorname{sum\text{-}ms}(s)) \end{array}}}

\vnbbookfile{theorem-library/mb67-lusin2}{theorem-library/mb67-lusin2}
\noindent{\small mb67-lusin2.scm -- construction of measures, chapter 5 section 1.5 (M-B67, 2026-10-26): Lusin sets are closed under countable disjoint unions and countable intersections (Proposition 5.20); the Borel subsets of a Polish space are Lusin (Corollary 5.21). lusin-disjoint-union      5.20, unions (members Lusin or empty; ...}\par
\vnbprop{lusin-disjoint-union}{\vnbstmt{\begin{array}{@{}l@{}} \forall s.; \forall e \in (\mathbb{N} \to \operatorname{power}(\operatorname{pts}(s))). \\ \quad \operatorname{is\text{-}metric\text{-}space}(s) \wedge \\ \quad \forall k \in \mathbb{N}.\; (\exists x.\; (x \in e(k)) \Rightarrow \operatorname{is\text{-}lusin}(s, e(k))) \wedge \\ \quad \forall i, j \in \mathbb{N}, \mathit{mbx} \in e(i).\; ((\mathit{mbx} \in e(j)) \Rightarrow (i = j)) \wedge \\ \quad \exists x.\; (x \in \operatorname{big\text{-}union}(n, \mathbb{N}, e(n))) \Rightarrow \\ \quad \quad \operatorname{is\text{-}lusin}(s, \operatorname{big\text{-}union}(n, \mathbb{N}, e(n))) \end{array}}}
\vnbprop{lusin-countable-intersection}{\vnbstmt{\begin{array}{@{}l@{}} \forall s.; \forall e \in (\mathbb{N} \to \operatorname{power}(\operatorname{pts}(s))). \\ \quad \operatorname{is\text{-}metric\text{-}space}(s) \wedge \\ \quad \forall k \in \mathbb{N}.\; \operatorname{is\text{-}lusin}(s, e(k)) \wedge \\ \quad \exists x.\; (x \in \{p \in \operatorname{pts}(s) : \forall \mathit{mxk} \in \mathbb{N}.\; (p \in e(\mathit{mxk}))\}) \Rightarrow \\ \quad \quad \operatorname{is\text{-}lusin}(s, \{p \in \operatorname{pts}(s) : \forall \mathit{mxk} \in \mathbb{N}.\; (p \in e(\mathit{mxk}))\}) \end{array}}}
\vnbprop{borel-polish-lusin}{\vnbstmt{\begin{array}{@{}l@{}} \forall s.; \forall b \in \operatorname{borel\text{-}sets}(s). \\ \quad (\operatorname{is\text{-}polish}(s) \wedge \exists x.\; (x \in b)) \Rightarrow \operatorname{is\text{-}lusin}(s, b) \end{array}}}

\vnbbookfile{theorem-library/mb67-sep}{theorem-library/mb67-sep}
\noindent{\small mb67-sep.scm -- construction of measures, chapter 5 section 1.6 (M-B67, 2026-10-26): the separation theorem. cyl-unfold, upd-facts, cyl-agree, cyl-zero, cyl-split, borel-separated-union separation-step, souslin-separation (Proposition 5.25), and its corollaries The proof is the Baire-space form of the notes' argument: ...}\par
\vnbprop{cyl-unfold}{\vnbstmt{\begin{array}{@{}l@{}} \forall w, n. \\ \quad (\operatorname{cyl}(w, n) \simeq \{z \in (\mathbb{N} \to \mathbb{N}) : \forall j \in \mathbb{N}.\; ((j < n) \Rightarrow (z(j) = w(j)))\}) \end{array}}}
\vnbprop{upd-facts}{\vnbstmt{\begin{array}{@{}l@{}} \forall w \in (\mathbb{N} \to \mathbb{N}), n, m \in \mathbb{N}. \\ \quad \operatorname{upd}(w, n, m) \in (\mathbb{N} \to \mathbb{N}) \wedge \\ \quad \forall j \in \mathbb{N}.\; ((j < n) \Rightarrow (\operatorname{upd}(w, n, m)(j) = w(j))) \wedge \\ \quad \forall j \in \mathbb{N}.\; (\neg (j < n) \Rightarrow (\operatorname{upd}(w, n, m)(j) = m)) \end{array}}}
\vnbprop{cyl-agree}{\vnbstmt{\begin{array}{@{}l@{}} \forall w, v \in (\mathbb{N} \to \mathbb{N}), n \in \mathbb{N}. \\ \quad \forall j \in \mathbb{N}.\; ((j < n) \Rightarrow (w(j) = v(j))) \Rightarrow (\operatorname{cyl}(w, n) = \operatorname{cyl}(v, n)) \end{array}}}
\vnbprop{cyl-zero}{\vnbstmt{\forall w \in (\mathbb{N} \to \mathbb{N}).\; (\operatorname{cyl}(w, 0) = (\mathbb{N} \to \mathbb{N}))}}
\vnbprop{cyl-split}{\vnbstmt{\begin{array}{@{}l@{}} \forall w \in (\mathbb{N} \to \mathbb{N}), n \in \mathbb{N}. \\ \quad \operatorname{cyl}(w, n) = \operatorname{big\text{-}union}(m, \mathbb{N}, \operatorname{cyl}(\operatorname{upd}(w, n, m), (n + 1))) \end{array}}}
\vnbprop{borel-separated-union}{\vnbstmt{\begin{array}{@{}l@{}} \forall s.; \forall a, b \in (\mathbb{N} \to \operatorname{power}(\operatorname{pts}(s))). \\ \quad (\operatorname{is\text{-}metric\text{-}space}(s) \wedge \forall k, l \in \mathbb{N}.\; \operatorname{borel\text{-}separated}(s, a(k), b(l))) \Rightarrow  \\ \quad \quad \operatorname{borel\text{-}separated}(s, \operatorname{big\text{-}union}(n, \mathbb{N}, a(n)), \operatorname{big\text{-}union}(n, \mathbb{N}, b(n))) \end{array}}}
\vnbprop{cyl-sub}{\vnbstmt{\forall w \in (\mathbb{N} \to \mathbb{N}), n \in \mathbb{N}.\; (\operatorname{cyl}(w, n) \subseteq (\mathbb{N} \to \mathbb{N}))}}
\vnbprop{image-cyl-power}{\vnbstmt{\begin{array}{@{}l@{}} \forall s.; \forall f \in ((\mathbb{N} \to \mathbb{N}) \to \operatorname{pts}(s)), w \in (\mathbb{N} \to \mathbb{N}), n \in \mathbb{N}. \\ \quad \operatorname{is\text{-}metric\text{-}space}(s) \Rightarrow (\operatorname{image}(f, \operatorname{cyl}(w, n)) \in \operatorname{power}(\operatorname{pts}(s))) \end{array}}}
\vnbprop{image-cyl-split}{\vnbstmt{\begin{array}{@{}l@{}} \forall f.; \forall w \in (\mathbb{N} \to \mathbb{N}), n \in \mathbb{N}. \\ \quad \operatorname{image}(f, \operatorname{cyl}(w, n)) \subseteq \operatorname{big\text{-}union}(m, \mathbb{N}, \operatorname{image}(f, \operatorname{cyl}(\operatorname{upd}(w, n, m), (n + 1)))) \end{array}}}
\vnbprop{borel-separated-mono}{\vnbstmt{\begin{array}{@{}l@{}} \forall s, a, b, c, d. \\ \quad ((c \subseteq a) \wedge ((d \subseteq b) \wedge \operatorname{borel\text{-}separated}(s, a, b))) \Rightarrow  \\ \quad \quad \operatorname{borel\text{-}separated}(s, c, d) \end{array}}}
\vnbprop{separation-step}{\vnbstmt{\begin{array}{@{}l@{}} \forall s.; \forall f, g \in ((\mathbb{N} \to \mathbb{N}) \to \operatorname{pts}(s)), u, v \in (\mathbb{N} \to \mathbb{N}), n \in \mathbb{N}. \\ \quad (\operatorname{is\text{-}metric\text{-}space}(s) \wedge \neg \operatorname{borel\text{-}separated}(s, \operatorname{image}(f, \operatorname{cyl}(u, n)), \operatorname{image}(g, \operatorname{cyl}(v, n)))) \Rightarrow  \\ \quad \quad \exists q_{1}, q_{2} \in \mathbb{N}. \\ \quad \quad \quad \neg \operatorname{borel\text{-}separated}(s, \operatorname{image}(f, \operatorname{cyl}(\operatorname{upd}(u, n, q_{1}), (n + 1))), \operatorname{image}(g, \operatorname{cyl}(\operatorname{upd}(v, n, q_{2}), (n + 1)))) \end{array}}}
\vnbprop{cyl-limit}{\vnbstmt{\begin{array}{@{}l@{}} \forall p \in (\mathbb{N} \to (\mathbb{N} \to \mathbb{N})). \\ \quad \forall \mathit{lqn}, j \in \mathbb{N}.\; ((j < \mathit{lqn}) \Rightarrow (p((\mathit{lqn} + 1))(j) = p(\mathit{lqn})(j))) \Rightarrow  \\ \quad \quad (\lambda \mathit{lyj} \in \mathbb{N}.\; p((\mathit{lyj} + 1))(\mathit{lyj})) \in (\mathbb{N} \to \mathbb{N}) \wedge \\ \quad \quad \forall n \in \mathbb{N}. \\ \quad \quad \quad \operatorname{cyl}((\lambda \mathit{lyj} \in \mathbb{N}.\; p((\mathit{lyj} + 1))(\mathit{lyj})), n) = \operatorname{cyl}(p(n), n) \end{array}}}
\vnbprop{baire-distinct-separated}{\vnbstmt{\begin{array}{@{}l@{}} \forall s, f, g.; \forall y, z \in (\mathbb{N} \to \mathbb{N}). \\ \quad (\operatorname{is\text{-}metric\text{-}space}(s) \wedge (\operatorname{is\text{-}baire\text{-}continuous}(s, f) \wedge (\operatorname{is\text{-}baire\text{-}continuous}(s, g) \wedge \neg (f(y) = g(z))))) \Rightarrow  \\ \quad \quad \exists n \in \mathbb{N}.\; \operatorname{borel\text{-}separated}(s, \operatorname{image}(f, \operatorname{cyl}(y, n)), \operatorname{image}(g, \operatorname{cyl}(z, n))) \end{array}}}
\vnbprop{baire-separation}{\vnbstmt{\begin{array}{@{}l@{}} \forall s, f, g. \\ \quad \operatorname{is\text{-}metric\text{-}space}(s) \wedge \\ \quad \operatorname{is\text{-}baire\text{-}continuous}(s, f) \wedge \\ \quad \operatorname{is\text{-}baire\text{-}continuous}(s, g) \wedge \\ \quad \forall x \in \operatorname{image}(f, (\mathbb{N} \to \mathbb{N})).\; \neg (x \in \operatorname{image}(g, (\mathbb{N} \to \mathbb{N}))) \Rightarrow \\ \quad \quad \operatorname{borel\text{-}separated}(s, \operatorname{image}(f, (\mathbb{N} \to \mathbb{N})), \operatorname{image}(g, (\mathbb{N} \to \mathbb{N}))) \end{array}}}
\vnbprop{souslin-separation}{\vnbstmt{\begin{array}{@{}l@{}} \forall s, a, b. \\ \quad (\operatorname{is\text{-}metric\text{-}space}(s) \wedge (\operatorname{is\text{-}souslin}(s, a) \wedge (\operatorname{is\text{-}souslin}(s, b) \wedge \forall x \in a.\; \neg (x \in b)))) \Rightarrow  \\ \quad \quad \operatorname{borel\text{-}separated}(s, a, b) \end{array}}}
\vnbprop{souslin-or-empty-separation}{\vnbstmt{\begin{array}{@{}l@{}} \forall s, a, b. \\ \quad \operatorname{is\text{-}metric\text{-}space}(s) \wedge \\ \quad a \subseteq \operatorname{pts}(s) \wedge \\ \quad b \subseteq \operatorname{pts}(s) \wedge \\ \quad \exists \mathit{mbx}.\; (\mathit{mbx} \in a) \Rightarrow \operatorname{is\text{-}souslin}(s, a) \wedge \\ \quad \exists \mathit{mbx}.\; (\mathit{mbx} \in b) \Rightarrow \operatorname{is\text{-}souslin}(s, b) \wedge \\ \quad \forall x \in a.\; \neg (x \in b) \Rightarrow \\ \quad \quad \operatorname{borel\text{-}separated}(s, a, b) \end{array}}}
\vnbprop{souslin-complement-borel}{\vnbstmt{\begin{array}{@{}l@{}} \forall s, e. \\ \quad \operatorname{is\text{-}metric\text{-}space}(s) \wedge \\ \quad e \subseteq \operatorname{pts}(s) \wedge \\ \quad \exists x.\; (x \in e) \Rightarrow \operatorname{is\text{-}souslin}(s, e) \wedge \\ \quad \exists x.\; (x \in \operatorname{complement\text{-}in}(\operatorname{pts}(s), e)) \Rightarrow  \\ \quad \quad \operatorname{is\text{-}souslin}(s, \operatorname{complement\text{-}in}(\operatorname{pts}(s), e)) \Rightarrow \\ \quad \quad e \in \operatorname{borel\text{-}sets}(s) \end{array}}}
\vnbprop{souslin-theorem}{\vnbstmt{\begin{array}{@{}l@{}} \forall s, e. \\ \quad (\operatorname{is\text{-}polish}(s) \wedge (e \subseteq \operatorname{pts}(s))) \Rightarrow  \\ \quad \quad e \in \operatorname{borel\text{-}sets}(s) \\ \quad \quad \Leftrightarrow\; \exists x.\; (x \in e) \Rightarrow \operatorname{is\text{-}souslin}(s, e) \wedge \\ \quad \quad \exists x.\; (x \in \operatorname{complement\text{-}in}(\operatorname{pts}(s), e)) \Rightarrow  \\ \quad \quad \quad \operatorname{is\text{-}souslin}(s, \operatorname{complement\text{-}in}(\operatorname{pts}(s), e)) \end{array}}}
\vnbprop{souslin-disjoint-envelopes}{\vnbstmt{\begin{array}{@{}l@{}} \forall s.; \forall a \in (\mathbb{N} \to \operatorname{power}(\operatorname{pts}(s))). \\ \quad \operatorname{is\text{-}metric\text{-}space}(s) \wedge \\ \quad \forall k \in \mathbb{N}.\; (\exists \mathit{mbx}.\; (\mathit{mbx} \in a(k)) \Rightarrow \operatorname{is\text{-}souslin}(s, a(k))) \wedge \\ \quad \forall k, l \in \mathbb{N}, x \in a(k).\; (\neg (k = l) \Rightarrow \neg (x \in a(l))) \Rightarrow \\ \quad \quad \exists e \in (\mathbb{N} \to \operatorname{borel\text{-}sets}(s)). \\ \quad \quad \quad \forall k \in \mathbb{N}.\; (a(k) \subseteq e(k)) \wedge \\ \quad \quad \quad \forall k, l \in \mathbb{N}, x \in e(k).\; (\neg (k = l) \Rightarrow \neg (x \in e(l))) \end{array}}}

\vnbbookfile{theorem-library/mb67-ls}{theorem-library/mb67-ls}
\noindent{\small mb67-ls.scm -- construction of measures, chapter 5 sections 1.6-1.7 (M-B67, 2026-10-26): Borel partitions of a separable space into small pieces (Proposition 5.28) and the Lusin-Souslin theorem (Proposition 5.30: an injective continuous image of a Polish space is Borel), with Corollary 5.31 (in a Polish space the Borel ...}\par
\vnbprop{dense-piece-borel}{\vnbstmt{\begin{array}{@{}l@{}} \forall t.; \forall s \in (\mathbb{N} \to \operatorname{pts}(t)), r.; \forall k \in \mathbb{N}. \\ \quad \operatorname{is\text{-}metric\text{-}space}(t) \wedge \\ \quad \forall x \in \operatorname{pts}(t), a. \\ \quad \quad \operatorname{pos\text{-}rr}(a) \Rightarrow \exists n \in \mathbb{N}.\; (\operatorname{dist}(t)(x, s(n)) < a) \wedge \\ \quad \operatorname{pos\text{-}rr}(r) \Rightarrow \\ \quad \quad \operatorname{complement\text{-}in}(\operatorname{ball}(t, s(k), r), \\ \quad \quad \quad \operatorname{big\text{-}union}(j, \mathrm{OS}(k), (\lambda \mathit{mbj} \in \mathbb{N}.\; \operatorname{ball}(t, s(\mathit{mbj}), r))(j))) \\ \quad \quad \in\; \operatorname{borel\text{-}sets}(t) \end{array}}}
\vnbprop{dense-piece-disjoint}{\vnbstmt{\begin{array}{@{}l@{}} \forall t.; \forall s \in (\mathbb{N} \to \operatorname{pts}(t)), r.; \forall k, a \in \mathbb{N}, x \in \operatorname{complement\text{-}in}(\operatorname{ball}(t, s(k), r), \operatorname{big\text{-}union}(j, \mathrm{OS}(k), (\lambda \mathit{mbj} \in \mathbb{N}.\; \operatorname{ball}(t, s(\mathit{mbj}), r))(j))). \\ \quad \operatorname{is\text{-}metric\text{-}space}(t) \wedge \\ \quad \forall x \in \operatorname{pts}(t), b. \\ \quad \quad \operatorname{pos\text{-}rr}(b) \Rightarrow \exists n \in \mathbb{N}.\; (\operatorname{dist}(t)(x, s(n)) < b) \wedge \\ \quad \operatorname{pos\text{-}rr}(r) \wedge \\ \quad \neg (k = a) \Rightarrow \\ \quad \quad \neg (x \in \operatorname{complement\text{-}in}(\operatorname{ball}(t, s(a), r), \operatorname{big\text{-}union}(j, \mathrm{OS}(a), (\lambda \mathit{mbj} \in \mathbb{N}.\; \operatorname{ball}(t, s(\mathit{mbj}), r))(j)))) \end{array}}}
\vnbprop{dense-piece-cover}{\vnbstmt{\begin{array}{@{}l@{}} \forall t.; \forall s \in (\mathbb{N} \to \operatorname{pts}(t)), r.; \forall x \in \operatorname{pts}(t). \\ \quad \operatorname{is\text{-}metric\text{-}space}(t) \wedge \\ \quad \forall x \in \operatorname{pts}(t), a. \\ \quad \quad \operatorname{pos\text{-}rr}(a) \Rightarrow \exists n \in \mathbb{N}.\; (\operatorname{dist}(t)(x, s(n)) < a) \wedge \\ \quad \operatorname{pos\text{-}rr}(r) \Rightarrow \\ \quad \quad \exists k \in \mathbb{N}. \\ \quad \quad \quad x \\ \quad \quad \quad \in\; \operatorname{complement\text{-}in}(\operatorname{ball}(t, s(k), r), \\ \quad \quad \quad \quad \operatorname{big\text{-}union}(j, \mathrm{OS}(k), (\lambda \mathit{mbj} \in \mathbb{N}.\; \operatorname{ball}(t, s(\mathit{mbj}), r))(j))) \end{array}}}
\vnbprop{dense-piece-small}{\vnbstmt{\begin{array}{@{}l@{}} \forall t.; \forall s \in (\mathbb{N} \to \operatorname{pts}(t)), r.; \forall k \in \mathbb{N}, x, a \in \operatorname{complement\text{-}in}(\operatorname{ball}(t, s(k), r), \operatorname{big\text{-}union}(j, \mathrm{OS}(k), (\lambda \mathit{mbj} \in \mathbb{N}.\; \operatorname{ball}(t, s(\mathit{mbj}), r))(j))). \\ \quad \operatorname{is\text{-}metric\text{-}space}(t) \wedge \\ \quad \forall x \in \operatorname{pts}(t), b. \\ \quad \quad \operatorname{pos\text{-}rr}(b) \Rightarrow \exists n \in \mathbb{N}.\; (\operatorname{dist}(t)(x, s(n)) < b) \wedge \\ \quad \operatorname{pos\text{-}rr}(r) \Rightarrow \\ \quad \quad \operatorname{dist}(t)(x, a) < (r + r) \end{array}}}
\vnbprop{code-least}{\vnbstmt{\begin{array}{@{}l@{}} \forall l \in \mathbb{N}, w \in (\mathbb{N} \to \mathbb{N}). \exists m \in \mathbb{N}. \\ \quad \forall j \in \mathbb{N}.\; ((j < l) \Rightarrow (\operatorname{nnfst}(\operatorname{snd\text{-}iter}(m, j)) = w(j))) \wedge \\ \quad \forall \mathit{mcm} \in \mathbb{N}. \\ \quad \quad (\mathit{mcm} < m) \Rightarrow  \\ \quad \quad \quad \neg \forall \mathit{mcj} \in \mathbb{N}.\; ((\mathit{mcj} < l) \Rightarrow (\operatorname{nnfst}(\operatorname{snd\text{-}iter}(\mathit{mcm}, \mathit{mcj})) = \operatorname{nnfst}(\operatorname{snd\text{-}iter}(m, \mathit{mcj})))) \end{array}}}
\vnbprop{dense-cell-borel}{\vnbstmt{\begin{array}{@{}l@{}} \forall t.; \forall s \in (\mathbb{N} \to \operatorname{pts}(t)), n, m \in \mathbb{N}. \\ \quad \operatorname{is\text{-}metric\text{-}space}(t) \wedge \\ \quad \forall x \in \operatorname{pts}(t), a. \\ \quad \quad \operatorname{pos\text{-}rr}(a) \Rightarrow \exists n \in \mathbb{N}.\; (\operatorname{dist}(t)(x, s(n)) < a) \Rightarrow \\ \quad \quad \{\mathit{mcx} \in \operatorname{pts}(t) : (\forall \mathit{mcm} \in \mathbb{N}.\; ((\mathit{mcm} < m) \Rightarrow \neg \forall j \in \mathbb{N}.\; ((j < (n + 1)) \Rightarrow (\operatorname{nnfst}(\operatorname{snd\text{-}iter}(\mathit{mcm}, j)) = \operatorname{nnfst}(\operatorname{snd\text{-}iter}(m, j))))) \wedge \forall i \in \mathbb{N}.\; ((i < (n + 1)) \Rightarrow (\mathit{mcx} \in \operatorname{complement\text{-}in}(\operatorname{ball}(t, s(\operatorname{nnfst}(\operatorname{snd\text{-}iter}(m, i))), \operatorname{recip}((i + 1))), \operatorname{big\text{-}union}(\mathit{uj}, \mathrm{OS}(\operatorname{nnfst}(\operatorname{snd\text{-}iter}(m, i))), (\lambda \mathit{mbj} \in \mathbb{N}.\; \operatorname{ball}(t, s(\mathit{mbj}), \operatorname{recip}((i + 1))))(\mathit{uj}))))))\} \\ \quad \quad \in\; \operatorname{borel\text{-}sets}(t) \end{array}}}
\vnbprop{dense-cell-disjoint}{\vnbstmt{\begin{array}{@{}l@{}} \forall t.; \forall s \in (\mathbb{N} \to \operatorname{pts}(t)), n, m, a \in \mathbb{N}, x \in \{\mathit{mcx} \in \operatorname{pts}(t) : (\forall \mathit{mcm} \in \mathbb{N}.\; ((\mathit{mcm} < m) \Rightarrow \neg \forall j \in \mathbb{N}.\; ((j < (n + 1)) \Rightarrow (\operatorname{nnfst}(\operatorname{snd\text{-}iter}(\mathit{mcm}, j)) = \operatorname{nnfst}(\operatorname{snd\text{-}iter}(m, j))))) \wedge \forall i \in \mathbb{N}.\; ((i < (n + 1)) \Rightarrow (\mathit{mcx} \in \operatorname{complement\text{-}in}(\operatorname{ball}(t, s(\operatorname{nnfst}(\operatorname{snd\text{-}iter}(m, i))), \operatorname{recip}((i + 1))), \operatorname{big\text{-}union}(\mathit{uj}, \mathrm{OS}(\operatorname{nnfst}(\operatorname{snd\text{-}iter}(m, i))), (\lambda \mathit{mbj} \in \mathbb{N}.\; \operatorname{ball}(t, s(\mathit{mbj}), \operatorname{recip}((i + 1))))(\mathit{uj}))))))\}. \\ \quad \operatorname{is\text{-}metric\text{-}space}(t) \wedge \\ \quad \forall x \in \operatorname{pts}(t), b. \\ \quad \quad \operatorname{pos\text{-}rr}(b) \Rightarrow \exists n \in \mathbb{N}.\; (\operatorname{dist}(t)(x, s(n)) < b) \wedge \\ \quad \neg (m = a) \Rightarrow \\ \quad \quad \neg (x \in \{\mathit{mcx} \in \operatorname{pts}(t) : (\forall \mathit{mcm} \in \mathbb{N}.\; ((\mathit{mcm} < a) \Rightarrow \neg \forall j \in \mathbb{N}.\; ((j < (n + 1)) \Rightarrow (\operatorname{nnfst}(\operatorname{snd\text{-}iter}(\mathit{mcm}, j)) = \operatorname{nnfst}(\operatorname{snd\text{-}iter}(a, j))))) \wedge \forall i \in \mathbb{N}.\; ((i < (n + 1)) \Rightarrow (\mathit{mcx} \in \operatorname{complement\text{-}in}(\operatorname{ball}(t, s(\operatorname{nnfst}(\operatorname{snd\text{-}iter}(a, i))), \operatorname{recip}((i + 1))), \operatorname{big\text{-}union}(\mathit{uj}, \mathrm{OS}(\operatorname{nnfst}(\operatorname{snd\text{-}iter}(a, i))), (\lambda \mathit{mbj} \in \mathbb{N}.\; \operatorname{ball}(t, s(\mathit{mbj}), \operatorname{recip}((i + 1))))(\mathit{uj}))))))\}) \end{array}}}
\vnbprop{dense-cell-cover}{\vnbstmt{\begin{array}{@{}l@{}} \forall t.; \forall s \in (\mathbb{N} \to \operatorname{pts}(t)), x \in \operatorname{pts}(t). \\ \quad \operatorname{is\text{-}metric\text{-}space}(t) \wedge \\ \quad \forall x \in \operatorname{pts}(t), a. \\ \quad \quad \operatorname{pos\text{-}rr}(a) \Rightarrow \exists n \in \mathbb{N}.\; (\operatorname{dist}(t)(x, s(n)) < a) \Rightarrow \\ \quad \quad \exists w \in (\mathbb{N} \to \mathbb{N}). \forall n \in \mathbb{N}. \exists m \in \mathbb{N}. \\ \quad \quad \quad x \\ \quad \quad \quad \in\; \{\mathit{mcx} \in \operatorname{pts}(t) : (\forall \mathit{mcm} \in \mathbb{N}.\; ((\mathit{mcm} < m) \Rightarrow \neg \forall j \in \mathbb{N}.\; ((j < (n + 1)) \Rightarrow (\operatorname{nnfst}(\operatorname{snd\text{-}iter}(\mathit{mcm}, j)) = \operatorname{nnfst}(\operatorname{snd\text{-}iter}(m, j))))) \wedge \forall i \in \mathbb{N}.\; ((i < (n + 1)) \Rightarrow (\mathit{mcx} \in \operatorname{complement\text{-}in}(\operatorname{ball}(t, s(\operatorname{nnfst}(\operatorname{snd\text{-}iter}(m, i))), \operatorname{recip}((i + 1))), \operatorname{big\text{-}union}(\mathit{uj}, \mathrm{OS}(\operatorname{nnfst}(\operatorname{snd\text{-}iter}(m, i))), (\lambda \mathit{mbj} \in \mathbb{N}.\; \operatorname{ball}(t, s(\mathit{mbj}), \operatorname{recip}((i + 1))))(\mathit{uj}))))))\} \wedge \\ \quad \quad \quad \forall \mathit{mpj} \in \mathbb{N}.\; ((\mathit{mpj} < (n + 1)) \Rightarrow (\operatorname{nnfst}(\operatorname{snd\text{-}iter}(m, \mathit{mpj})) = w(\mathit{mpj}))) \end{array}}}
\vnbprop{dense-cell-small}{\vnbstmt{\begin{array}{@{}l@{}} \forall t.; \forall s \in (\mathbb{N} \to \operatorname{pts}(t)), n, m \in \mathbb{N}, x, a \in \{\mathit{mcx} \in \operatorname{pts}(t) : (\forall \mathit{mcm} \in \mathbb{N}.\; ((\mathit{mcm} < m) \Rightarrow \neg \forall j \in \mathbb{N}.\; ((j < (n + 1)) \Rightarrow (\operatorname{nnfst}(\operatorname{snd\text{-}iter}(\mathit{mcm}, j)) = \operatorname{nnfst}(\operatorname{snd\text{-}iter}(m, j))))) \wedge \forall i \in \mathbb{N}.\; ((i < (n + 1)) \Rightarrow (\mathit{mcx} \in \operatorname{complement\text{-}in}(\operatorname{ball}(t, s(\operatorname{nnfst}(\operatorname{snd\text{-}iter}(m, i))), \operatorname{recip}((i + 1))), \operatorname{big\text{-}union}(\mathit{uj}, \mathrm{OS}(\operatorname{nnfst}(\operatorname{snd\text{-}iter}(m, i))), (\lambda \mathit{mbj} \in \mathbb{N}.\; \operatorname{ball}(t, s(\mathit{mbj}), \operatorname{recip}((i + 1))))(\mathit{uj}))))))\}. \\ \quad \operatorname{is\text{-}metric\text{-}space}(t) \wedge \\ \quad \forall x \in \operatorname{pts}(t), b. \\ \quad \quad \operatorname{pos\text{-}rr}(b) \Rightarrow \exists n \in \mathbb{N}.\; (\operatorname{dist}(t)(x, s(n)) < b) \Rightarrow \\ \quad \quad \operatorname{dist}(t)(x, a) < (\operatorname{recip}((n + 1)) + \operatorname{recip}((n + 1))) \end{array}}}
\vnbprop{dense-cell-nested}{\vnbstmt{\begin{array}{@{}l@{}} \forall t.; \forall s \in (\mathbb{N} \to \operatorname{pts}(t)), n, m, i, a \in \mathbb{N}. \\ \quad \operatorname{is\text{-}metric\text{-}space}(t) \wedge \\ \quad \forall x \in \operatorname{pts}(t), b. \\ \quad \quad \operatorname{pos\text{-}rr}(b) \Rightarrow \exists n \in \mathbb{N}.\; (\operatorname{dist}(t)(x, s(n)) < b) \wedge \\ \quad i < n \wedge \\ \quad \forall \mathit{mcm} \in \mathbb{N}. \\ \quad \quad (\mathit{mcm} < a) \Rightarrow  \\ \quad \quad \quad \neg \forall j \in \mathbb{N}.\; ((j < (i + 1)) \Rightarrow (\operatorname{nnfst}(\operatorname{snd\text{-}iter}(\mathit{mcm}, j)) = \operatorname{nnfst}(\operatorname{snd\text{-}iter}(a, j)))) \wedge \\ \quad \forall \mathit{mqj} \in \mathbb{N}. \\ \quad \quad (\mathit{mqj} < (i + 1)) \Rightarrow (\operatorname{nnfst}(\operatorname{snd\text{-}iter}(a, \mathit{mqj})) = \operatorname{nnfst}(\operatorname{snd\text{-}iter}(m, \mathit{mqj}))) \Rightarrow \\ \quad \quad \{\mathit{mcx} \in \operatorname{pts}(t) : (\forall \mathit{mcm} \in \mathbb{N}.\; ((\mathit{mcm} < m) \Rightarrow \neg \forall j \in \mathbb{N}.\; ((j < (n + 1)) \Rightarrow (\operatorname{nnfst}(\operatorname{snd\text{-}iter}(\mathit{mcm}, j)) = \operatorname{nnfst}(\operatorname{snd\text{-}iter}(m, j))))) \wedge \forall \mathit{mci} \in \mathbb{N}.\; ((\mathit{mci} < (n + 1)) \Rightarrow (\mathit{mcx} \in \operatorname{complement\text{-}in}(\operatorname{ball}(t, s(\operatorname{nnfst}(\operatorname{snd\text{-}iter}(m, \mathit{mci}))), \operatorname{recip}((\mathit{mci} + 1))), \operatorname{big\text{-}union}(\mathit{uj}, \mathrm{OS}(\operatorname{nnfst}(\operatorname{snd\text{-}iter}(m, \mathit{mci}))), (\lambda \mathit{mbj} \in \mathbb{N}.\; \operatorname{ball}(t, s(\mathit{mbj}), \operatorname{recip}((\mathit{mci} + 1))))(\mathit{uj}))))))\} \\ \quad \quad \subseteq\; \{\mathit{mcx} \in \operatorname{pts}(t) : (\forall \mathit{mcm} \in \mathbb{N}.\; ((\mathit{mcm} < a) \Rightarrow \neg \forall j \in \mathbb{N}.\; ((j < (i + 1)) \Rightarrow (\operatorname{nnfst}(\operatorname{snd\text{-}iter}(\mathit{mcm}, j)) = \operatorname{nnfst}(\operatorname{snd\text{-}iter}(a, j))))) \wedge \forall \mathit{mci} \in \mathbb{N}.\; ((\mathit{mci} < (i + 1)) \Rightarrow (\mathit{mcx} \in \operatorname{complement\text{-}in}(\operatorname{ball}(t, s(\operatorname{nnfst}(\operatorname{snd\text{-}iter}(a, \mathit{mci}))), \operatorname{recip}((\mathit{mci} + 1))), \operatorname{big\text{-}union}(\mathit{uj}, \mathrm{OS}(\operatorname{nnfst}(\operatorname{snd\text{-}iter}(a, \mathit{mci}))), (\lambda \mathit{mbj} \in \mathbb{N}.\; \operatorname{ball}(t, s(\mathit{mbj}), \operatorname{recip}((\mathit{mci} + 1))))(\mathit{uj}))))))\} \end{array}}}
\vnbprop{nbhd-borel}{\vnbstmt{\begin{array}{@{}l@{}} \forall s, a, o. \\ \quad (\operatorname{is\text{-}metric\text{-}space}(s) \wedge ((a \subseteq \operatorname{pts}(s)) \wedge \operatorname{pos\text{-}rr}(o))) \Rightarrow  \\ \quad \quad \{y \in \operatorname{pts}(s) : \exists \mathit{nba} \in a.\; (\operatorname{dist}(s)(\mathit{nba}, y) < o)\} \in \operatorname{borel\text{-}sets}(s) \end{array}}}
\vnbprop{dense-image-envelopes}{\vnbstmt{\begin{array}{@{}l@{}} \forall t.; \forall a \in (\mathbb{N} \to \operatorname{pts}(t)), s, h.; \forall n \in \mathbb{N}. \\ \quad \operatorname{is\text{-}polish}(t) \wedge \\ \quad \forall x \in \operatorname{pts}(t), b. \\ \quad \quad \operatorname{pos\text{-}rr}(b) \Rightarrow \exists n \in \mathbb{N}.\; (\operatorname{dist}(t)(x, a(n)) < b) \wedge \\ \quad \operatorname{is\text{-}metric\text{-}space}(s) \wedge \\ \quad \operatorname{is\text{-}continuous}(t, s, h) \wedge \\ \quad \forall x, c \in \operatorname{pts}(t).\; ((h(x) = h(c)) \Rightarrow (x = c)) \Rightarrow \\ \quad \quad \exists g \in (\mathbb{N} \to \operatorname{borel\text{-}sets}(s)). \\ \quad \quad \quad \forall m \in \mathbb{N}. \\ \quad \quad \quad \quad \operatorname{image}(h, \\ \quad \quad \quad \quad \quad \{\mathit{mcx} \in \operatorname{pts}(t) : (\forall \mathit{mcm} \in \mathbb{N}.\; ((\mathit{mcm} < m) \Rightarrow \neg \forall j \in \mathbb{N}.\; ((j < (n + 1)) \Rightarrow (\operatorname{nnfst}(\operatorname{snd\text{-}iter}(\mathit{mcm}, j)) = \operatorname{nnfst}(\operatorname{snd\text{-}iter}(m, j))))) \wedge \forall i \in \mathbb{N}.\; ((i < (n + 1)) \Rightarrow (\mathit{mcx} \in \operatorname{complement\text{-}in}(\operatorname{ball}(t, a(\operatorname{nnfst}(\operatorname{snd\text{-}iter}(m, i))), \operatorname{recip}((i + 1))), \operatorname{big\text{-}union}(\mathit{uj}, \mathrm{OS}(\operatorname{nnfst}(\operatorname{snd\text{-}iter}(m, i))), (\lambda \mathit{mbj} \in \mathbb{N}.\; \operatorname{ball}(t, a(\mathit{mbj}), \operatorname{recip}((i + 1))))(\mathit{uj}))))))\}) \\ \quad \quad \quad \quad \subseteq\; g(m) \wedge \\ \quad \quad \quad \forall k, l \in \mathbb{N}, x \in g(k).\; (\neg (k = l) \Rightarrow \neg (x \in g(l))) \end{array}}}
\vnbprop{dense-envelopes-all-levels}{\vnbstmt{\begin{array}{@{}l@{}} \forall t.; \forall a \in (\mathbb{N} \to \operatorname{pts}(t)), s, h. \\ \quad \operatorname{is\text{-}polish}(t) \wedge \\ \quad \forall x \in \operatorname{pts}(t), b. \\ \quad \quad \operatorname{pos\text{-}rr}(b) \Rightarrow \exists n \in \mathbb{N}.\; (\operatorname{dist}(t)(x, a(n)) < b) \wedge \\ \quad \operatorname{is\text{-}metric\text{-}space}(s) \wedge \\ \quad \operatorname{is\text{-}continuous}(t, s, h) \wedge \\ \quad \forall x, c \in \operatorname{pts}(t).\; ((h(x) = h(c)) \Rightarrow (x = c)) \Rightarrow \\ \quad \quad \exists g \in (\mathbb{N} \to (\mathbb{N} \to \operatorname{borel\text{-}sets}(s))). \forall n \in \mathbb{N}. \\ \quad \quad \quad \forall m \in \mathbb{N}. \\ \quad \quad \quad \quad \operatorname{image}(h, \\ \quad \quad \quad \quad \quad \{\mathit{mcx} \in \operatorname{pts}(t) : (\forall \mathit{mcm} \in \mathbb{N}.\; ((\mathit{mcm} < m) \Rightarrow \neg \forall j \in \mathbb{N}.\; ((j < (n + 1)) \Rightarrow (\operatorname{nnfst}(\operatorname{snd\text{-}iter}(\mathit{mcm}, j)) = \operatorname{nnfst}(\operatorname{snd\text{-}iter}(m, j))))) \wedge \forall i \in \mathbb{N}.\; ((i < (n + 1)) \Rightarrow (\mathit{mcx} \in \operatorname{complement\text{-}in}(\operatorname{ball}(t, a(\operatorname{nnfst}(\operatorname{snd\text{-}iter}(m, i))), \operatorname{recip}((i + 1))), \operatorname{big\text{-}union}(\mathit{uj}, \mathrm{OS}(\operatorname{nnfst}(\operatorname{snd\text{-}iter}(m, i))), (\lambda \mathit{mbj} \in \mathbb{N}.\; \operatorname{ball}(t, a(\mathit{mbj}), \operatorname{recip}((i + 1))))(\mathit{uj}))))))\}) \\ \quad \quad \quad \quad \subseteq\; g(n)(m) \wedge \\ \quad \quad \quad \forall k, l \in \mathbb{N}, x \in g(n)(k). \\ \quad \quad \quad \quad \neg (k = l) \Rightarrow \neg (x \in g(n)(l)) \end{array}}}
\vnbprop{dense-final-borel}{\vnbstmt{\begin{array}{@{}l@{}} \forall t.; \forall a \in (\mathbb{N} \to \operatorname{pts}(t)), s, h.; \forall g \in (\mathbb{N} \to (\mathbb{N} \to \operatorname{borel\text{-}sets}(s))). \\ \quad \operatorname{is\text{-}polish}(t) \wedge \\ \quad \forall x \in \operatorname{pts}(t), b. \\ \quad \quad \operatorname{pos\text{-}rr}(b) \Rightarrow \exists n \in \mathbb{N}.\; (\operatorname{dist}(t)(x, a(n)) < b) \wedge \\ \quad \operatorname{is\text{-}metric\text{-}space}(s) \wedge \\ \quad \operatorname{is\text{-}continuous}(t, s, h) \wedge \\ \quad \forall x, c \in \operatorname{pts}(t).\; ((h(x) = h(c)) \Rightarrow (x = c)) \wedge \\ \quad \forall n \in \mathbb{N}. \\ \quad \quad \forall m \in \mathbb{N}. \\ \quad \quad \quad \operatorname{image}(h, \\ \quad \quad \quad \quad \{\mathit{mcx} \in \operatorname{pts}(t) : (\forall \mathit{mcm} \in \mathbb{N}.\; ((\mathit{mcm} < m) \Rightarrow \neg \forall j \in \mathbb{N}.\; ((j < (n + 1)) \Rightarrow (\operatorname{nnfst}(\operatorname{snd\text{-}iter}(\mathit{mcm}, j)) = \operatorname{nnfst}(\operatorname{snd\text{-}iter}(m, j))))) \wedge \forall i \in \mathbb{N}.\; ((i < (n + 1)) \Rightarrow (\mathit{mcx} \in \operatorname{complement\text{-}in}(\operatorname{ball}(t, a(\operatorname{nnfst}(\operatorname{snd\text{-}iter}(m, i))), \operatorname{recip}((i + 1))), \operatorname{big\text{-}union}(\mathit{uj}, \mathrm{OS}(\operatorname{nnfst}(\operatorname{snd\text{-}iter}(m, i))), (\lambda \mathit{mbj} \in \mathbb{N}.\; \operatorname{ball}(t, a(\mathit{mbj}), \operatorname{recip}((i + 1))))(\mathit{uj}))))))\}) \\ \quad \quad \quad \subseteq\; g(n)(m) \wedge \\ \quad \quad \forall k, l \in \mathbb{N}, x \in g(n)(k). \\ \quad \quad \quad \neg (k = l) \Rightarrow \neg (x \in g(n)(l)) \Rightarrow \\ \quad \quad \operatorname{complement\text{-}in}(\operatorname{pts}(s), \\ \quad \quad \quad \operatorname{big\text{-}union}(\mathit{fn}, \mathbb{N}, \\ \quad \quad \quad \quad \operatorname{complement\text{-}in}(\operatorname{pts}(s), \\ \quad \quad \quad \quad \quad \operatorname{big\text{-}union}(\mathit{fm}, \mathbb{N}, \\ \quad \quad \quad \quad \quad \quad (g(\mathit{fn})(\mathit{fm}) \cap (\{y \in \operatorname{pts}(s) : \exists \mathit{nba} \in \operatorname{image}(h, \{\mathit{mcx} \in \operatorname{pts}(t) : (\forall \mathit{mcm} \in \mathbb{N}.\; ((\mathit{mcm} < \mathit{fm}) \Rightarrow \neg \forall j \in \mathbb{N}.\; ((j < (\mathit{fn} + 1)) \Rightarrow (\operatorname{nnfst}(\operatorname{snd\text{-}iter}(\mathit{mcm}, j)) = \operatorname{nnfst}(\operatorname{snd\text{-}iter}(\mathit{fm}, j))))) \wedge \forall i \in \mathbb{N}.\; ((i < (\mathit{fn} + 1)) \Rightarrow (\mathit{mcx} \in \operatorname{complement\text{-}in}(\operatorname{ball}(t, a(\operatorname{nnfst}(\operatorname{snd\text{-}iter}(\mathit{fm}, i))), \operatorname{recip}((i + 1))), \operatorname{big\text{-}union}(\mathit{uj}, \mathrm{OS}(\operatorname{nnfst}(\operatorname{snd\text{-}iter}(\mathit{fm}, i))), (\lambda \mathit{mbj} \in \mathbb{N}.\; \operatorname{ball}(t, a(\mathit{mbj}), \operatorname{recip}((i + 1))))(\mathit{uj}))))))\}).\; (\operatorname{dist}(s)(\mathit{nba}, y) < \operatorname{recip}((\mathit{fn} + 1)))\} \cap \operatorname{complement\text{-}in}(\operatorname{pts}(s), \operatorname{big\text{-}union}(\mathit{hi}, \mathbb{N}, \operatorname{complement\text{-}in}(\operatorname{pts}(s), \{\mathit{mvy} \in \operatorname{pts}(s) : ((\mathit{hi} < \mathit{fn}) \Rightarrow (\mathit{mvy} \in \operatorname{big\text{-}union}(\mathit{vm}, \mathbb{N}, \{\mathit{mwy} \in \operatorname{pts}(s) : ((\forall \mathit{mcm} \in \mathbb{N}.\; ((\mathit{mcm} < \mathit{vm}) \Rightarrow \neg \forall j \in \mathbb{N}.\; ((j < (\mathit{hi} + 1)) \Rightarrow (\operatorname{nnfst}(\operatorname{snd\text{-}iter}(\mathit{mcm}, j)) = \operatorname{nnfst}(\operatorname{snd\text{-}iter}(\mathit{vm}, j))))) \wedge \forall \mathit{mqj} \in \mathbb{N}.\; ((\mathit{mqj} < (\mathit{hi} + 1)) \Rightarrow (\operatorname{nnfst}(\operatorname{snd\text{-}iter}(\mathit{vm}, \mathit{mqj})) = \operatorname{nnfst}(\operatorname{snd\text{-}iter}(\mathit{fm}, \mathit{mqj}))))) \wedge (\mathit{mwy} \in g(\mathit{hi})(\mathit{vm})))\})))\}))))))))) \\ \quad \quad \in\; \operatorname{borel\text{-}sets}(s) \end{array}}}
\vnbprop{dense-image-in-final}{\vnbstmt{\begin{array}{@{}l@{}} \forall t.; \forall a \in (\mathbb{N} \to \operatorname{pts}(t)), s, h.; \forall g \in (\mathbb{N} \to (\mathbb{N} \to \operatorname{borel\text{-}sets}(s))). \\ \quad \operatorname{is\text{-}polish}(t) \wedge \\ \quad \forall x \in \operatorname{pts}(t), b. \\ \quad \quad \operatorname{pos\text{-}rr}(b) \Rightarrow \exists n \in \mathbb{N}.\; (\operatorname{dist}(t)(x, a(n)) < b) \wedge \\ \quad \operatorname{is\text{-}metric\text{-}space}(s) \wedge \\ \quad \operatorname{is\text{-}continuous}(t, s, h) \wedge \\ \quad \forall x, c \in \operatorname{pts}(t).\; ((h(x) = h(c)) \Rightarrow (x = c)) \wedge \\ \quad \forall n \in \mathbb{N}. \\ \quad \quad \forall m \in \mathbb{N}. \\ \quad \quad \quad \operatorname{image}(h, \\ \quad \quad \quad \quad \{\mathit{mcx} \in \operatorname{pts}(t) : (\forall \mathit{mcm} \in \mathbb{N}.\; ((\mathit{mcm} < m) \Rightarrow \neg \forall j \in \mathbb{N}.\; ((j < (n + 1)) \Rightarrow (\operatorname{nnfst}(\operatorname{snd\text{-}iter}(\mathit{mcm}, j)) = \operatorname{nnfst}(\operatorname{snd\text{-}iter}(m, j))))) \wedge \forall i \in \mathbb{N}.\; ((i < (n + 1)) \Rightarrow (\mathit{mcx} \in \operatorname{complement\text{-}in}(\operatorname{ball}(t, a(\operatorname{nnfst}(\operatorname{snd\text{-}iter}(m, i))), \operatorname{recip}((i + 1))), \operatorname{big\text{-}union}(\mathit{uj}, \mathrm{OS}(\operatorname{nnfst}(\operatorname{snd\text{-}iter}(m, i))), (\lambda \mathit{mbj} \in \mathbb{N}.\; \operatorname{ball}(t, a(\mathit{mbj}), \operatorname{recip}((i + 1))))(\mathit{uj}))))))\}) \\ \quad \quad \quad \subseteq\; g(n)(m) \wedge \\ \quad \quad \forall k, l \in \mathbb{N}, x \in g(n)(k). \\ \quad \quad \quad \neg (k = l) \Rightarrow \neg (x \in g(n)(l)) \Rightarrow \\ \quad \quad \operatorname{image}(h, \operatorname{pts}(t)) \\ \quad \quad \subseteq\; \operatorname{complement\text{-}in}(\operatorname{pts}(s), \\ \quad \quad \quad \operatorname{big\text{-}union}(\mathit{fn}, \mathbb{N}, \\ \quad \quad \quad \quad \operatorname{complement\text{-}in}(\operatorname{pts}(s), \\ \quad \quad \quad \quad \quad \operatorname{big\text{-}union}(\mathit{fm}, \mathbb{N}, \\ \quad \quad \quad \quad \quad \quad (g(\mathit{fn})(\mathit{fm}) \cap (\{y \in \operatorname{pts}(s) : \exists \mathit{nba} \in \operatorname{image}(h, \{\mathit{mcx} \in \operatorname{pts}(t) : (\forall \mathit{mcm} \in \mathbb{N}.\; ((\mathit{mcm} < \mathit{fm}) \Rightarrow \neg \forall j \in \mathbb{N}.\; ((j < (\mathit{fn} + 1)) \Rightarrow (\operatorname{nnfst}(\operatorname{snd\text{-}iter}(\mathit{mcm}, j)) = \operatorname{nnfst}(\operatorname{snd\text{-}iter}(\mathit{fm}, j))))) \wedge \forall i \in \mathbb{N}.\; ((i < (\mathit{fn} + 1)) \Rightarrow (\mathit{mcx} \in \operatorname{complement\text{-}in}(\operatorname{ball}(t, a(\operatorname{nnfst}(\operatorname{snd\text{-}iter}(\mathit{fm}, i))), \operatorname{recip}((i + 1))), \operatorname{big\text{-}union}(\mathit{uj}, \mathrm{OS}(\operatorname{nnfst}(\operatorname{snd\text{-}iter}(\mathit{fm}, i))), (\lambda \mathit{mbj} \in \mathbb{N}.\; \operatorname{ball}(t, a(\mathit{mbj}), \operatorname{recip}((i + 1))))(\mathit{uj}))))))\}).\; (\operatorname{dist}(s)(\mathit{nba}, y) < \operatorname{recip}((\mathit{fn} + 1)))\} \cap \operatorname{complement\text{-}in}(\operatorname{pts}(s), \operatorname{big\text{-}union}(\mathit{hi}, \mathbb{N}, \operatorname{complement\text{-}in}(\operatorname{pts}(s), \{\mathit{mvy} \in \operatorname{pts}(s) : ((\mathit{hi} < \mathit{fn}) \Rightarrow (\mathit{mvy} \in \operatorname{big\text{-}union}(\mathit{vm}, \mathbb{N}, \{\mathit{mwy} \in \operatorname{pts}(s) : ((\forall \mathit{mcm} \in \mathbb{N}.\; ((\mathit{mcm} < \mathit{vm}) \Rightarrow \neg \forall j \in \mathbb{N}.\; ((j < (\mathit{hi} + 1)) \Rightarrow (\operatorname{nnfst}(\operatorname{snd\text{-}iter}(\mathit{mcm}, j)) = \operatorname{nnfst}(\operatorname{snd\text{-}iter}(\mathit{vm}, j))))) \wedge \forall \mathit{mqj} \in \mathbb{N}.\; ((\mathit{mqj} < (\mathit{hi} + 1)) \Rightarrow (\operatorname{nnfst}(\operatorname{snd\text{-}iter}(\mathit{vm}, \mathit{mqj})) = \operatorname{nnfst}(\operatorname{snd\text{-}iter}(\mathit{fm}, \mathit{mqj}))))) \wedge (\mathit{mwy} \in g(\mathit{hi})(\mathit{vm})))\})))\}))))))))) \end{array}}}
\vnbprop{dense-final-in-image}{\vnbstmt{\begin{array}{@{}l@{}} \forall t.; \forall a \in (\mathbb{N} \to \operatorname{pts}(t)), s, h.; \forall g \in (\mathbb{N} \to (\mathbb{N} \to \operatorname{borel\text{-}sets}(s))). \\ \quad \operatorname{is\text{-}polish}(t) \wedge \\ \quad \forall x \in \operatorname{pts}(t), b. \\ \quad \quad \operatorname{pos\text{-}rr}(b) \Rightarrow \exists n \in \mathbb{N}.\; (\operatorname{dist}(t)(x, a(n)) < b) \wedge \\ \quad \operatorname{is\text{-}metric\text{-}space}(s) \wedge \\ \quad \operatorname{is\text{-}continuous}(t, s, h) \wedge \\ \quad \forall x, c \in \operatorname{pts}(t).\; ((h(x) = h(c)) \Rightarrow (x = c)) \wedge \\ \quad \forall n \in \mathbb{N}. \\ \quad \quad \forall m \in \mathbb{N}. \\ \quad \quad \quad \operatorname{image}(h, \\ \quad \quad \quad \quad \{\mathit{mcx} \in \operatorname{pts}(t) : (\forall \mathit{mcm} \in \mathbb{N}.\; ((\mathit{mcm} < m) \Rightarrow \neg \forall j \in \mathbb{N}.\; ((j < (n + 1)) \Rightarrow (\operatorname{nnfst}(\operatorname{snd\text{-}iter}(\mathit{mcm}, j)) = \operatorname{nnfst}(\operatorname{snd\text{-}iter}(m, j))))) \wedge \forall i \in \mathbb{N}.\; ((i < (n + 1)) \Rightarrow (\mathit{mcx} \in \operatorname{complement\text{-}in}(\operatorname{ball}(t, a(\operatorname{nnfst}(\operatorname{snd\text{-}iter}(m, i))), \operatorname{recip}((i + 1))), \operatorname{big\text{-}union}(\mathit{uj}, \mathrm{OS}(\operatorname{nnfst}(\operatorname{snd\text{-}iter}(m, i))), (\lambda \mathit{mbj} \in \mathbb{N}.\; \operatorname{ball}(t, a(\mathit{mbj}), \operatorname{recip}((i + 1))))(\mathit{uj}))))))\}) \\ \quad \quad \quad \subseteq\; g(n)(m) \wedge \\ \quad \quad \forall k, l \in \mathbb{N}, x \in g(n)(k). \\ \quad \quad \quad \neg (k = l) \Rightarrow \neg (x \in g(n)(l)) \Rightarrow \\ \quad \quad \operatorname{complement\text{-}in}(\operatorname{pts}(s), \\ \quad \quad \quad \operatorname{big\text{-}union}(\mathit{fn}, \mathbb{N}, \\ \quad \quad \quad \quad \operatorname{complement\text{-}in}(\operatorname{pts}(s), \\ \quad \quad \quad \quad \quad \operatorname{big\text{-}union}(\mathit{fm}, \mathbb{N}, \\ \quad \quad \quad \quad \quad \quad (g(\mathit{fn})(\mathit{fm}) \cap (\{y \in \operatorname{pts}(s) : \exists \mathit{nba} \in \operatorname{image}(h, \{\mathit{mcx} \in \operatorname{pts}(t) : (\forall \mathit{mcm} \in \mathbb{N}.\; ((\mathit{mcm} < \mathit{fm}) \Rightarrow \neg \forall j \in \mathbb{N}.\; ((j < (\mathit{fn} + 1)) \Rightarrow (\operatorname{nnfst}(\operatorname{snd\text{-}iter}(\mathit{mcm}, j)) = \operatorname{nnfst}(\operatorname{snd\text{-}iter}(\mathit{fm}, j))))) \wedge \forall i \in \mathbb{N}.\; ((i < (\mathit{fn} + 1)) \Rightarrow (\mathit{mcx} \in \operatorname{complement\text{-}in}(\operatorname{ball}(t, a(\operatorname{nnfst}(\operatorname{snd\text{-}iter}(\mathit{fm}, i))), \operatorname{recip}((i + 1))), \operatorname{big\text{-}union}(\mathit{uj}, \mathrm{OS}(\operatorname{nnfst}(\operatorname{snd\text{-}iter}(\mathit{fm}, i))), (\lambda \mathit{mbj} \in \mathbb{N}.\; \operatorname{ball}(t, a(\mathit{mbj}), \operatorname{recip}((i + 1))))(\mathit{uj}))))))\}).\; (\operatorname{dist}(s)(\mathit{nba}, y) < \operatorname{recip}((\mathit{fn} + 1)))\} \cap \operatorname{complement\text{-}in}(\operatorname{pts}(s), \operatorname{big\text{-}union}(\mathit{hi}, \mathbb{N}, \operatorname{complement\text{-}in}(\operatorname{pts}(s), \{\mathit{mvy} \in \operatorname{pts}(s) : ((\mathit{hi} < \mathit{fn}) \Rightarrow (\mathit{mvy} \in \operatorname{big\text{-}union}(\mathit{vm}, \mathbb{N}, \{\mathit{mwy} \in \operatorname{pts}(s) : ((\forall \mathit{mcm} \in \mathbb{N}.\; ((\mathit{mcm} < \mathit{vm}) \Rightarrow \neg \forall j \in \mathbb{N}.\; ((j < (\mathit{hi} + 1)) \Rightarrow (\operatorname{nnfst}(\operatorname{snd\text{-}iter}(\mathit{mcm}, j)) = \operatorname{nnfst}(\operatorname{snd\text{-}iter}(\mathit{vm}, j))))) \wedge \forall \mathit{mqj} \in \mathbb{N}.\; ((\mathit{mqj} < (\mathit{hi} + 1)) \Rightarrow (\operatorname{nnfst}(\operatorname{snd\text{-}iter}(\mathit{vm}, \mathit{mqj})) = \operatorname{nnfst}(\operatorname{snd\text{-}iter}(\mathit{fm}, \mathit{mqj}))))) \wedge (\mathit{mwy} \in g(\mathit{hi})(\mathit{vm})))\})))\}))))))))) \\ \quad \quad \subseteq\; \operatorname{image}(h, \operatorname{pts}(t)) \end{array}}}
\vnbprop{lusin-souslin}{\vnbstmt{\begin{array}{@{}l@{}} \forall t, s, h. \\ \quad \operatorname{is\text{-}polish}(t) \wedge \\ \quad \operatorname{is\text{-}metric\text{-}space}(s) \wedge \\ \quad \operatorname{is\text{-}continuous}(t, s, h) \wedge \\ \quad \forall x, a \in \operatorname{pts}(t).\; ((h(x) = h(a)) \Rightarrow (x = a)) \Rightarrow \\ \quad \quad \operatorname{image}(h, \operatorname{pts}(t)) \in \operatorname{borel\text{-}sets}(s) \end{array}}}
\vnbprop{continuous-into-subspace-ambient}{\vnbstmt{\begin{array}{@{}l@{}} \forall s, e, t, h. \\ \quad (\operatorname{is\text{-}metric\text{-}space}(s) \wedge ((e \subseteq \operatorname{pts}(s)) \wedge \operatorname{is\text{-}continuous}(t, \operatorname{subspace\text{-}ms}(s, e), h))) \Rightarrow  \\ \quad \quad \operatorname{is\text{-}continuous}(t, s, h) \end{array}}}
\vnbprop{lusin-borel}{\vnbstmt{\begin{array}{@{}l@{}} \forall s, e. \\ \quad (\operatorname{is\text{-}metric\text{-}space}(s) \wedge ((e \subseteq \operatorname{pts}(s)) \wedge \operatorname{is\text{-}lusin}(s, e))) \Rightarrow  \\ \quad \quad e \in \operatorname{borel\text{-}sets}(s) \end{array}}}
\vnbprop{polish-borel-iff-lusin}{\vnbstmt{\begin{array}{@{}l@{}} \forall s, e. \\ \quad (\operatorname{is\text{-}polish}(s) \wedge (e \subseteq \operatorname{pts}(s))) \Rightarrow  \\ \quad \quad (e \in \operatorname{borel\text{-}sets}(s)) \Leftrightarrow (\exists x.\; (x \in e) \Rightarrow \operatorname{is\text{-}lusin}(s, e)) \end{array}}}
\vnbprop{lusin-souslin-borel-domain}{\vnbstmt{\begin{array}{@{}l@{}} \forall p.; \forall b \in \operatorname{borel\text{-}sets}(p), s, i. \\ \quad \operatorname{is\text{-}polish}(p) \wedge \\ \quad \operatorname{is\text{-}metric\text{-}space}(s) \wedge \\ \quad \operatorname{is\text{-}continuous}(\operatorname{subspace\text{-}ms}(p, b), s, i) \wedge \\ \quad \forall x, a \in b.\; ((i(x) = i(a)) \Rightarrow (x = a)) \Rightarrow \\ \quad \quad \operatorname{image}(i, b) \in \operatorname{borel\text{-}sets}(s) \end{array}}}

\vnbbookfile{theorem-library/mb67-s38}{theorem-library/mb67-s38}
\noindent{\small mb67-s38.scm -- construction of measures, chapter 5 section 1.9 (M-B67, 2026-10-27): Proposition 5.38, Souslin sets of a Polish space are exactly the Souslin sets of families of CLOSED sets (sieves or not). Sufficiency: Z = \{(x, z) : x in S(z|n) for all n\} is closed in X x N\textasciicircum{}N (Polish), and the Souslin set is its first ...}\par
\vnbprop{souslin-set-unfold}{\vnbstmt{\begin{array}{@{}l@{}} \forall s, v. \\ \quad (\operatorname{souslin\text{-}set}(s, v) \simeq \{p \in \operatorname{pts}(s) : \exists z \in (\mathbb{N} \to \mathbb{N}).\; \forall n \in \mathbb{N}.\; (p \in v(\operatorname{restrict}(z, \mathrm{OS}(n))))\}) \end{array}}}
\vnbprop{souslin-closed-graph-closed}{\vnbstmt{\begin{array}{@{}l@{}} \forall s.; \forall v \in (\mathit{nstar} \to \operatorname{power}(\operatorname{pts}(s))). \\ \quad (\operatorname{is\text{-}polish}(s) \wedge \forall p \in \mathit{nstar}.\; \operatorname{is\text{-}closed}(s, v(p))) \Rightarrow  \\ \quad \quad \operatorname{is\text{-}closed}(\operatorname{product2\text{-}ms}(s, \operatorname{baire\text{-}ms}), \\ \quad \quad \quad \{w \in (\operatorname{pts}(s) \times \operatorname{pts}(\operatorname{baire\text{-}ms})) : \forall n \in \mathbb{N}.\; (\operatorname{nth}(1, w) \in v(\operatorname{restrict}(\operatorname{nth}(2, w), \mathrm{OS}(n))))\}) \end{array}}}
\vnbprop{closed-family-souslin}{\vnbstmt{\begin{array}{@{}l@{}} \forall s.; \forall v \in (\mathit{nstar} \to \operatorname{power}(\operatorname{pts}(s))). \\ \quad (\operatorname{is\text{-}polish}(s) \wedge (\forall p \in \mathit{nstar}.\; \operatorname{is\text{-}closed}(s, v(p)) \wedge \exists x.\; (x \in \operatorname{souslin\text{-}set}(s, v)))) \Rightarrow  \\ \quad \quad \operatorname{is\text{-}souslin}(s, \operatorname{souslin\text{-}set}(s, v)) \end{array}}}
\vnbprop{dist0-closed}{\vnbstmt{\begin{array}{@{}l@{}} \forall s, a. \\ \quad (\operatorname{is\text{-}metric\text{-}space}(s) \wedge (a \subseteq \operatorname{pts}(s))) \Rightarrow  \\ \quad \quad \operatorname{is\text{-}closed}(s, \\ \quad \quad \quad \{y \in \operatorname{pts}(s) : \forall e \in \mathbb{R}.\; (\operatorname{pos\text{-}rr}(e) \Rightarrow \exists \mathit{m38a} \in a.\; (\operatorname{dist}(s)(y, \mathit{m38a}) < e))\}) \end{array}}}
\vnbprop{souslin-closed-family}{\vnbstmt{\begin{array}{@{}l@{}} \forall s, e. \\ \quad (\operatorname{is\text{-}metric\text{-}space}(s) \wedge \operatorname{is\text{-}souslin}(s, e)) \Rightarrow  \\ \quad \quad \exists t \in (\mathit{nstar} \to \operatorname{power}(\operatorname{pts}(s))). \\ \quad \quad \quad \forall p \in \mathit{nstar}.\; \operatorname{is\text{-}closed}(s, t(p)) \wedge (\operatorname{souslin\text{-}set}(s, t) = e) \end{array}}}
\vnbprop{souslin-iff-closed-family}{\vnbstmt{\begin{array}{@{}l@{}} \forall s, e. \\ \quad \operatorname{is\text{-}polish}(s) \Rightarrow  \\ \quad \quad \operatorname{is\text{-}souslin}(s, e) \\ \quad \quad \Leftrightarrow\; \exists x.\; (x \in e) \wedge \\ \quad \quad \exists t \in (\mathit{nstar} \to \operatorname{power}(\operatorname{pts}(s))). \\ \quad \quad \quad \forall p \in \mathit{nstar}.\; \operatorname{is\text{-}closed}(s, t(p)) \wedge (\operatorname{souslin\text{-}set}(s, t) = e) \end{array}}}

\vnbbookfile{theorem-library/mb67-b32}{theorem-library/mb67-b32}
\noindent{\small mb67-b32.scm -- construction of measures, chapter 5 sections 1.7-1.8 (M-B67, 2026-10-27): Theorem 5.32 (an injective Borel map of a Polish space into a separable metric space maps Borel sets to Borel sets) and its consequences 5.34 and 5.37.  The core is borel-injection-image: for X, Z Polish and theta: X -$>$ Z ...}\par
\vnbprop{separable-countably-separated}{\vnbstmt{\begin{array}{@{}l@{}} \forall p. \\ \quad \operatorname{is\text{-}separable}(p) \Rightarrow \operatorname{is\text{-}countably\text{-}separated}(\operatorname{pts}(p), \operatorname{borel\text{-}sets}(p)) \end{array}}}
\vnbprop{graph-of-unfold}{\vnbstmt{\begin{array}{@{}l@{}} \forall a, b, f. \\ \quad (\operatorname{graph\text{-}of}(a, b, f) \simeq \{p \in (a \times b) : (f(\operatorname{nth}(1, p)) = \operatorname{nth}(2, p))\}) \end{array}}}
\vnbprop{borel-injection-image}{\vnbstmt{\begin{array}{@{}l@{}} \forall a, z, h.; \forall b \in \operatorname{borel\text{-}sets}(a). \\ \quad \operatorname{is\text{-}polish}(a) \wedge \\ \quad \operatorname{is\text{-}polish}(z) \wedge \\ \quad \operatorname{is\text{-}measurable\text{-}map}(\operatorname{pts}(a), \operatorname{borel\text{-}sets}(a), \operatorname{pts}(z), \operatorname{borel\text{-}sets}(z), h) \wedge \\ \quad \forall x, c \in \operatorname{pts}(a).\; ((h(x) = h(c)) \Rightarrow (x = c)) \Rightarrow \\ \quad \quad \operatorname{image}(h, b) \in \operatorname{borel\text{-}sets}(z) \end{array}}}
\vnbprop{injective-borel-image}{\vnbstmt{\begin{array}{@{}l@{}} \forall a, s, i.; \forall b \in \operatorname{borel\text{-}sets}(a). \\ \quad \operatorname{is\text{-}polish}(a) \wedge \\ \quad \operatorname{is\text{-}separable}(s) \wedge \\ \quad \operatorname{is\text{-}measurable\text{-}map}(\operatorname{pts}(a), \operatorname{borel\text{-}sets}(a), \operatorname{pts}(s), \operatorname{borel\text{-}sets}(s), i) \wedge \\ \quad \forall x, c \in \operatorname{pts}(a).\; ((i(x) = i(c)) \Rightarrow (x = c)) \Rightarrow \\ \quad \quad \operatorname{image}(i, b) \in \operatorname{borel\text{-}sets}(s) \end{array}}}
\vnbprop{standard-injection-images}{\vnbstmt{\begin{array}{@{}l@{}} \forall s, a, b, c, i.; \forall d \in a. \\ \quad \operatorname{is\text{-}standard\text{-}borel}(s, a) \wedge \\ \quad \operatorname{is\text{-}countably\text{-}separated}(b, c) \wedge \\ \quad \operatorname{is\text{-}measurable\text{-}map}(s, a, b, c, i) \wedge \\ \quad \forall x, f \in s.\; ((i(x) = i(f)) \Rightarrow (x = f)) \Rightarrow \\ \quad \quad \operatorname{image}(i, d) \in c \end{array}}}
\vnbprop{restrict-sigma-unfold}{\vnbstmt{\forall a, e.\; (\operatorname{restrict\text{-}sigma}(a, e) \simeq \{\mathit{rsa} \in a : (\mathit{rsa} \subseteq e)\})}}
\vnbprop{restrict-sigma-algebra}{\vnbstmt{\begin{array}{@{}l@{}} \forall s, a.; \forall e \in a. \\ \quad \operatorname{is\text{-}sigma\text{-}algebra}(s, a) \Rightarrow \operatorname{is\text{-}sigma\text{-}algebra}(e, \operatorname{restrict\text{-}sigma}(a, e)) \end{array}}}
\vnbprop{standard-measurable-subset}{\vnbstmt{\begin{array}{@{}l@{}} \forall s, a.; \forall b \in a. \\ \quad (\operatorname{is\text{-}standard\text{-}borel}(s, a) \wedge \exists x.\; (x \in b)) \Rightarrow  \\ \quad \quad \operatorname{is\text{-}standard\text{-}borel}(b, \operatorname{restrict\text{-}sigma}(a, b)) \end{array}}}
\vnbprop{standard-injection-image-standard}{\vnbstmt{\begin{array}{@{}l@{}} \forall s, a, b, c, i. \\ \quad \operatorname{is\text{-}standard\text{-}borel}(s, a) \wedge \\ \quad \operatorname{is\text{-}countably\text{-}separated}(b, c) \wedge \\ \quad \operatorname{is\text{-}measurable\text{-}map}(s, a, b, c, i) \wedge \\ \quad \forall x, d \in s.\; ((i(x) = i(d)) \Rightarrow (x = d)) \Rightarrow \\ \quad \quad \operatorname{is\text{-}standard\text{-}borel}(\operatorname{image}(i, s), \operatorname{restrict\text{-}sigma}(c, \operatorname{image}(i, s))) \end{array}}}

\vnbbookfile{theorem-library/mb67-s35}{theorem-library/mb67-s35}
\noindent{\small mb67-s35.scm -- construction of measures, chapter 5 section 1.8 (M-B67, 2026-10-27): Proposition 5.35, a countable disjoint union of standard Borel spaces is standard.  The witnesses (Polish spaces T(k) and isomorphisms F(k)) are chosen by global choice; the union is carried by G([k, x]) = [k, F(k)(x)] onto SUM-MS(T), ...}\par
\vnbprop{sum-set-unfold}{\vnbstmt{\begin{array}{@{}l@{}} \forall s. \\ \quad (\operatorname{disjoint\text{-}sum\text{-}set}(s) \simeq \{p \in (\mathbb{N} \times \operatorname{big\text{-}union}(k, \mathbb{N}, s(k))) : (\operatorname{nth}(2, p) \in s(\operatorname{nth}(1, p)))\}) \end{array}}}
\vnbprop{sum-sigma-unfold}{\vnbstmt{\begin{array}{@{}l@{}} \forall s, a. \\ \quad (\operatorname{sum\text{-}sigma}(s, a) \simeq \{\mathit{sga} \in \operatorname{power}(\operatorname{disjoint\text{-}sum\text{-}set}(s)) : \forall k \in \mathbb{N}.\; (\{y \in s(k) : ({[}k, y] \in \mathit{sga})\} \in a(k))\}) \end{array}}}
\vnbprop{sum-set-parts}{\vnbstmt{\begin{array}{@{}l@{}} \forall s.; \forall p \in \operatorname{disjoint\text{-}sum\text{-}set}(s). \\ \quad (\operatorname{nth}(1, p) \in \mathbb{N}) \wedge ((\operatorname{nth}(2, p) \in s(\operatorname{nth}(1, p))) \wedge (p = {[}\operatorname{nth}(1, p), \operatorname{nth}(2, p)])) \end{array}}}
\vnbprop{sum-set-intro}{\vnbstmt{\forall s.; \forall k \in \mathbb{N}, x \in s(k).\; ({[}k, x] \in \operatorname{disjoint\text{-}sum\text{-}set}(s))}}
\vnbprop{sum-set-is-set}{\vnbstmt{\forall s.\; (\forall i \in \mathbb{N}.\; (s(i) \in \mathbf{Set}) \Rightarrow (\operatorname{disjoint\text{-}sum\text{-}set}(s) \in \mathbf{Set}))}}
\vnbprop{sum-sigma-algebra}{\vnbstmt{\begin{array}{@{}l@{}} \forall s, a. \\ \quad \forall j \in \mathbb{N}.\; \operatorname{is\text{-}sigma\text{-}algebra}(s(j), a(j)) \Rightarrow  \\ \quad \quad \operatorname{is\text{-}sigma\text{-}algebra}(\operatorname{disjoint\text{-}sum\text{-}set}(s), \operatorname{sum\text{-}sigma}(s, a)) \end{array}}}
\vnbprop{standard-disjoint-union}{\vnbstmt{\begin{array}{@{}l@{}} \forall s, a \in (\mathbb{N} \to \mathbf{Set}). \\ \quad \forall h \in \mathbb{N}.\; \operatorname{is\text{-}standard\text{-}borel}(s(h), a(h)) \Rightarrow  \\ \quad \quad \operatorname{is\text{-}standard\text{-}borel}(\operatorname{disjoint\text{-}sum\text{-}set}(s), \operatorname{sum\text{-}sigma}(s, a)) \end{array}}}

\vnbbookfile{theorem-library/mb12-metric}{theorem-library/mb12-metric}
\noindent{\small mb12-metric.scm -- construction-of-measures.pdf Appendix A.3 (metric odds and ends): A.6 (a compact set and a disjoint closed set are at positive distance), and more below. Agent ISOM (brief M-B1B2), 2026-10-26.}\par
\vnbprop{compact-closed-positive-distance}{\vnbstmt{\begin{array}{@{}l@{}} \forall s, k, f. \\ \quad \operatorname{is\text{-}metric\text{-}space}(s) \wedge \\ \quad k \subseteq \operatorname{pts}(s) \wedge \\ \quad \operatorname{is\text{-}compact}(\operatorname{subspace\text{-}ms}(s, k)) \wedge \\ \quad \operatorname{is\text{-}closed}(s, f) \wedge \\ \quad \forall x \in k.\; \neg (x \in f) \Rightarrow \\ \quad \quad \neg (\operatorname{set\text{-}dist}(s, k, f) = 0) \end{array}}}
\vnbprop{lipschitz-bound-iff-balls}{\vnbstmt{\begin{array}{@{}l@{}} \forall s, t, k. \\ \quad (\operatorname{is\text{-}metric\text{-}space}(s) \wedge (\operatorname{is\text{-}metric\text{-}space}(t) \wedge ((\operatorname{pts}(t) = \operatorname{pts}(s)) \wedge \operatorname{pos\text{-}rr}(k)))) \Rightarrow  \\ \quad \quad \forall x, y \in \operatorname{pts}(s).\; (\operatorname{dist}(s)(x, y) \leq (k \cdot \operatorname{dist}(t)(x, y))) \\ \quad \quad \Leftrightarrow\; \forall x \in \operatorname{pts}(s), r. \\ \quad \quad \quad \operatorname{pos\text{-}rr}(r) \Rightarrow (\operatorname{ball}(t, x, r) \subseteq \operatorname{ball}(s, x, (k \cdot r))) \end{array}}}
\vnbprop{lipschitz-lower-iff-balls}{\vnbstmt{\begin{array}{@{}l@{}} \forall s, t, k. \\ \quad (\operatorname{is\text{-}metric\text{-}space}(s) \wedge (\operatorname{is\text{-}metric\text{-}space}(t) \wedge ((\operatorname{pts}(t) = \operatorname{pts}(s)) \wedge \operatorname{pos\text{-}rr}(k)))) \Rightarrow  \\ \quad \quad \forall x, y \in \operatorname{pts}(s).\; ((k \cdot \operatorname{dist}(t)(x, y)) \leq \operatorname{dist}(s)(x, y)) \\ \quad \quad \Leftrightarrow\; \forall x \in \operatorname{pts}(s), r. \\ \quad \quad \quad \operatorname{pos\text{-}rr}(r) \Rightarrow (\operatorname{ball}(s, x, (k \cdot r)) \subseteq \operatorname{ball}(t, x, r)) \end{array}}}
\vnbprop{lipschitz-equivalent-iff-balls}{\vnbstmt{\begin{array}{@{}l@{}} \forall s, t, a, k. \\ \quad (\operatorname{is\text{-}metric\text{-}space}(s) \wedge (\operatorname{is\text{-}metric\text{-}space}(t) \wedge ((\operatorname{pts}(t) = \operatorname{pts}(s)) \wedge (\operatorname{pos\text{-}rr}(a) \wedge \operatorname{pos\text{-}rr}(k))))) \Rightarrow  \\ \quad \quad \forall x, y \in \operatorname{pts}(s).\; ((a \cdot \operatorname{dist}(t)(x, y)) \leq \operatorname{dist}(s)(x, y)) \wedge \\ \quad \quad \forall x, y \in \operatorname{pts}(s).\; (\operatorname{dist}(s)(x, y) \leq (k \cdot \operatorname{dist}(t)(x, y))) \\ \quad \quad \Leftrightarrow\; \forall x \in \operatorname{pts}(s), r. \\ \quad \quad \quad \operatorname{pos\text{-}rr}(r) \Rightarrow (\operatorname{ball}(s, x, (a \cdot r)) \subseteq \operatorname{ball}(t, x, r)) \wedge \\ \quad \quad \forall x \in \operatorname{pts}(s), r. \\ \quad \quad \quad \operatorname{pos\text{-}rr}(r) \Rightarrow (\operatorname{ball}(t, x, r) \subseteq \operatorname{ball}(s, x, (k \cdot r))) \end{array}}}
\vnbprop{open-complement-closed}{\vnbstmt{\begin{array}{@{}l@{}} \forall s, o. \\ \quad (\operatorname{is\text{-}metric\text{-}space}(s) \wedge \operatorname{is\text{-}open}(s, o)) \Rightarrow  \\ \quad \quad \operatorname{is\text{-}closed}(s, \operatorname{complement\text{-}in}(\operatorname{pts}(s), o)) \end{array}}}
\vnbprop{compact-split-open-cover}{\vnbstmt{\begin{array}{@{}l@{}} \forall s, k, u, v. \\ \quad \operatorname{is\text{-}metric\text{-}space}(s) \wedge \\ \quad k \subseteq \operatorname{pts}(s) \wedge \\ \quad \operatorname{is\text{-}compact}(\operatorname{subspace\text{-}ms}(s, k)) \wedge \\ \quad \operatorname{is\text{-}open}(s, u) \wedge \\ \quad \operatorname{is\text{-}open}(s, v) \wedge \\ \quad k \subseteq (u \cup v) \Rightarrow \\ \quad \quad \exists a, b. \\ \quad \quad \quad \operatorname{is\text{-}compact}(\operatorname{subspace\text{-}ms}(s, a)) \wedge \\ \quad \quad \quad \operatorname{is\text{-}compact}(\operatorname{subspace\text{-}ms}(s, b)) \wedge \\ \quad \quad \quad a \subseteq u \wedge \\ \quad \quad \quad b \subseteq v \wedge \\ \quad \quad \quad k = (a \cup b) \end{array}}}

\vnbbookfile{theorem-library/measurable-maps-metric}{theorem-library/measurable-maps-metric}
\noindent{\small measurable-maps-metric.scm -- WAVE 2 (2026-10-27..30), wired 2026-10-30 from scratchpad/mf/mf-mc1.scm, scratchpad/mf/mf-mc2.scm, scratchpad/mf/mf-mc3.scm. Library: measure.}\par
\vnbprop{sep-exists-family-in-sigma}{\vnbstmt{\begin{array}{@{}l@{}} \forall a, b, g. \\ \quad (\operatorname{is\text{-}sigma\text{-}algebra}(a, b) \wedge (g \in (\mathbb{N} \to b))) \Rightarrow  \\ \quad \quad \{x \in a : \exists k \in \mathbb{N}.\; (x \in g(k))\} \in b \end{array}}}
\vnbprop{sep-forall-family-in-sigma}{\vnbstmt{\begin{array}{@{}l@{}} \forall a, b, g. \\ \quad (\operatorname{is\text{-}sigma\text{-}algebra}(a, b) \wedge (g \in (\mathbb{N} \to b))) \Rightarrow  \\ \quad \quad \{x \in a : \forall k \in \mathbb{N}.\; (x \in g(k))\} \in b \end{array}}}
\vnbprop{metric-nbhd-open}{\vnbstmt{\begin{array}{@{}l@{}} \forall s, a, o. \\ \quad (\operatorname{is\text{-}metric\text{-}space}(s) \wedge ((a \subseteq \operatorname{pts}(s)) \wedge \operatorname{pos\text{-}rr}(o))) \Rightarrow  \\ \quad \quad \operatorname{is\text{-}open}(s, \{y \in \operatorname{pts}(s) : \exists \mathit{nba} \in a.\; (\operatorname{dist}(s)(\mathit{nba}, y) < o)\}) \end{array}}}
\vnbprop{measurable-limit-open-preimage}{\vnbstmt{\begin{array}{@{}l@{}} \forall a, b, s, c, f, u. \\ \quad \operatorname{is\text{-}sigma\text{-}algebra}(a, b) \wedge \\ \quad \operatorname{is\text{-}metric\text{-}space}(s) \wedge \\ \quad c \in (\mathbb{N} \to (a \to \operatorname{pts}(s))) \wedge \\ \quad \forall n \in \mathbb{N}.\; \operatorname{is\text{-}measurable\text{-}map}(a, b, \operatorname{pts}(s), \operatorname{borel\text{-}sets}(s), c(n)) \wedge \\ \quad f \in (a \to \operatorname{pts}(s)) \wedge \\ \quad \forall x \in a.\; \operatorname{converges\text{-}to}(s, (\lambda k \in \mathbb{N}.\; c(k)(x)), f(x)) \wedge \\ \quad \operatorname{is\text{-}open}(s, u) \Rightarrow \\ \quad \quad \{\mathit{mfx} \in a : (f(\mathit{mfx}) \in u)\} \in b \end{array}}}
\vnbprop{measurable-map-ptwise-limit}{\vnbstmt{\begin{array}{@{}l@{}} \forall a, b, s, c, f. \\ \quad \operatorname{is\text{-}sigma\text{-}algebra}(a, b) \wedge \\ \quad \operatorname{is\text{-}metric\text{-}space}(s) \wedge \\ \quad c \in (\mathbb{N} \to (a \to \operatorname{pts}(s))) \wedge \\ \quad \forall n \in \mathbb{N}.\; \operatorname{is\text{-}measurable\text{-}map}(a, b, \operatorname{pts}(s), \operatorname{borel\text{-}sets}(s), c(n)) \wedge \\ \quad f \in (a \to \operatorname{pts}(s)) \wedge \\ \quad \forall x \in a.\; \operatorname{converges\text{-}to}(s, (\lambda k \in \mathbb{N}.\; c(k)(x)), f(x)) \Rightarrow \\ \quad \quad \operatorname{is\text{-}measurable\text{-}map}(a, b, \operatorname{pts}(s), \operatorname{borel\text{-}sets}(s), f) \end{array}}}
\vnbprop{measurable-map-ptwise-eq}{\vnbstmt{\begin{array}{@{}l@{}} \forall a, b, s, c, f, h. \\ \quad (\operatorname{is\text{-}measurable\text{-}map}(a, b, s, c, f) \wedge ((h \in (a \to s)) \wedge \forall x \in a.\; (f(x) = h(x)))) \Rightarrow  \\ \quad \quad \operatorname{is\text{-}measurable\text{-}map}(a, b, s, c, h) \end{array}}}
\vnbprop{measurable-map-pair}{\vnbstmt{\begin{array}{@{}l@{}} \forall a, b, s, x, c, y, f, g. \\ \quad (\operatorname{is\text{-}measurable\text{-}map}(a, b, s, x, f) \wedge \operatorname{is\text{-}measurable\text{-}map}(a, b, c, y, g)) \Rightarrow  \\ \quad \quad \operatorname{is\text{-}measurable\text{-}map}(a, b, (s \times c), \operatorname{product\text{-}sigma}(s, x, c, y), \\ \quad \quad \quad (\lambda \mathit{mfpx} \in a.\; {[}f(\mathit{mfpx}), g(\mathit{mfpx})])) \end{array}}}
\vnbprop{measurable-continuous-combination}{\vnbstmt{\begin{array}{@{}l@{}} \forall a, b, s, i, f, g. \\ \quad \operatorname{is\text{-}separable}(s) \wedge \\ \quad \operatorname{is\text{-}continuous}(\operatorname{product2\text{-}ms}(s, s), s, i) \wedge \\ \quad \operatorname{is\text{-}measurable\text{-}map}(a, b, \operatorname{pts}(s), \operatorname{borel\text{-}sets}(s), f) \wedge \\ \quad \operatorname{is\text{-}measurable\text{-}map}(a, b, \operatorname{pts}(s), \operatorname{borel\text{-}sets}(s), g) \Rightarrow \\ \quad \quad \operatorname{is\text{-}measurable\text{-}map}(a, b, \operatorname{pts}(s), \operatorname{borel\text{-}sets}(s), (\lambda x \in a.\; i({[}f(x), g(x)]))) \end{array}}}
\vnbprop{disjointed-seq-in-sigma}{\vnbstmt{\begin{array}{@{}l@{}} \forall a, b, e, k. \\ \quad (\operatorname{is\text{-}sigma\text{-}algebra}(a, b) \wedge ((e \in (\mathbb{N} \to b)) \wedge (k \in \mathbb{N}))) \Rightarrow  \\ \quad \quad \operatorname{complement\text{-}in}(e(k), \operatorname{big\text{-}union}(j, \mathrm{OS}(k), e(j))) \in b \end{array}}}
\vnbprop{disjointed-seq-unique}{\vnbstmt{\begin{array}{@{}l@{}} \forall a, b, e, x, k, j. \\ \quad \operatorname{is\text{-}sigma\text{-}algebra}(a, b) \wedge \\ \quad e \in (\mathbb{N} \to b) \wedge \\ \quad k \in \mathbb{N} \wedge \\ \quad j \in \mathbb{N} \wedge \\ \quad x \in \operatorname{complement\text{-}in}(e(k), \operatorname{big\text{-}union}(\mathit{uj}, \mathrm{OS}(k), e(\mathit{uj}))) \wedge \\ \quad x \in \operatorname{complement\text{-}in}(e(j), \operatorname{big\text{-}union}(\mathit{uj}, \mathrm{OS}(j), e(\mathit{uj}))) \Rightarrow \\ \quad \quad k = j \end{array}}}
\vnbprop{disjointed-seq-cover}{\vnbstmt{\begin{array}{@{}l@{}} \forall a, b, e, x, k. \\ \quad (\operatorname{is\text{-}sigma\text{-}algebra}(a, b) \wedge ((e \in (\mathbb{N} \to b)) \wedge ((k \in \mathbb{N}) \wedge (x \in e(k))))) \Rightarrow  \\ \quad \quad \exists m \in \mathbb{N}.\; (x \in \operatorname{complement\text{-}in}(e(m), \operatorname{big\text{-}union}(j, \mathrm{OS}(m), e(j)))) \end{array}}}
\vnbprop{step-map-approximation}{\vnbstmt{\begin{array}{@{}l@{}} \forall a, b, s, f, c. \\ \quad (\operatorname{is\text{-}separable}(s) \wedge (\operatorname{is\text{-}measurable\text{-}map}(a, b, \operatorname{pts}(s), \operatorname{borel\text{-}sets}(s), f) \wedge \operatorname{pos\text{-}rr}(c))) \Rightarrow  \\ \quad \quad \exists g. \\ \quad \quad \quad \operatorname{is\text{-}step\text{-}map}(a, b, s, g) \wedge \forall x \in a.\; (\operatorname{dist}(s)(f(x), g(x)) < c) \end{array}}}
\vnbprop{step-map-uniform-limit}{\vnbstmt{\begin{array}{@{}l@{}} \forall a, b, s, f. \\ \quad (\operatorname{is\text{-}separable}(s) \wedge \operatorname{is\text{-}measurable\text{-}map}(a, b, \operatorname{pts}(s), \operatorname{borel\text{-}sets}(s), f)) \Rightarrow  \\ \quad \quad \exists c \in (\mathbb{N} \to (a \to \operatorname{pts}(s))). \forall n \in \mathbb{N}. \\ \quad \quad \quad (\operatorname{is\text{-}step\text{-}map}(a, b, s, c(n)) \wedge \forall x \in a.\; (\operatorname{dist}(s)(f(x), c(n)(x)) < \operatorname{recip}((n + 1)))) \end{array}}}

\vnbbookfile{theorem-library/local-measure}{theorem-library/local-measure}
\noindent{\small local-measure.scm -- WAVE 2 (2026-10-27..30), wired 2026-10-30 from scratchpad/mf/mf-mg1.scm. Library: measure.}\par
\vnbprop{local-measure-monotone}{\vnbstmt{\begin{array}{@{}l@{}} \forall b, c, d, o, a, e. \\ \quad (\operatorname{is\text{-}local\text{-}measure}(b, c, d, o) \wedge ((e \in d) \wedge ((a \in d) \wedge (a \subseteq e)))) \Rightarrow  \\ \quad \quad o(a) \leq o(e) \end{array}}}
\vnbprop{local-measure-cover-le}{\vnbstmt{\begin{array}{@{}l@{}} \forall a, b, c, o, e, d. \\ \quad (\operatorname{is\text{-}local\text{-}measure}(a, b, c, o) \wedge ((e \in c) \wedge ((d \in (\mathbb{N} \to c)) \wedge (e \subseteq \operatorname{big\text{-}union}(n, \mathbb{N}, d(n)))))) \Rightarrow  \\ \quad \quad o(e) \leq \operatorname{esum}((\lambda k \in \mathbb{N}.\; o(d(k)))) \end{array}}}
\vnbprop{local-measure-premeasure}{\vnbstmt{\begin{array}{@{}l@{}} \forall a, b, c, o. \\ \quad (\operatorname{is\text{-}local\text{-}measure}(a, b, c, o) \wedge (\emptyset \in c)) \Rightarrow \operatorname{is\text{-}premeasure}(a, c, o) \end{array}}}
\vnbprop{local-measure-outer-extends}{\vnbstmt{\begin{array}{@{}l@{}} \forall a, b, c, o, e. \\ \quad (\operatorname{is\text{-}local\text{-}measure}(a, b, c, o) \wedge ((\emptyset \in c) \wedge (e \in c))) \Rightarrow  \\ \quad \quad \operatorname{method\text{-}i\text{-}outer}(a, c, o)(e) = o(e) \end{array}}}
\vnbprop{local-measure-outer-measurable}{\vnbstmt{\begin{array}{@{}l@{}} \forall a, b, c, o, s. \\ \quad (\operatorname{is\text{-}local\text{-}measure}(a, b, c, o) \wedge ((\emptyset \in c) \wedge (s \in b))) \Rightarrow  \\ \quad \quad \operatorname{is\text{-}caratheodory\text{-}measurable}(a, \operatorname{method\text{-}i\text{-}outer}(a, c, o), s) \end{array}}}
\vnbprop{local-measure-extension-le-outer}{\vnbstmt{\begin{array}{@{}l@{}} \forall a, b, c, o, u, s. \\ \quad \operatorname{is\text{-}local\text{-}measure}(a, b, c, o) \wedge \\ \quad \emptyset \in c \wedge \\ \quad \operatorname{is\text{-}measure}(a, b, u) \wedge \\ \quad \forall e \in c.\; (u(e) = o(e)) \wedge \\ \quad s \in b \Rightarrow \\ \quad \quad u(s) \leq \operatorname{method\text{-}i\text{-}outer}(a, c, o)(s) \end{array}}}
\vnbprop{local-measure-extensions-agree}{\vnbstmt{\begin{array}{@{}l@{}} \forall a, b, c, o, u, d, g. \\ \quad \operatorname{is\text{-}local\text{-}measure}(a, b, c, o) \wedge \\ \quad \operatorname{is\text{-}measure}(a, b, u) \wedge \\ \quad \operatorname{is\text{-}measure}(a, b, d) \wedge \\ \quad \forall e \in c.\; (u(e) = o(e)) \wedge \\ \quad \forall e \in c.\; (d(e) = o(e)) \wedge \\ \quad g \in (\mathbb{N} \to c) \Rightarrow \\ \quad \quad u(\operatorname{big\text{-}union}(n, \mathbb{N}, g(n))) = d(\operatorname{big\text{-}union}(n, \mathbb{N}, g(n))) \end{array}}}
\vnbprop{local-measure-max-extension}{\vnbstmt{\begin{array}{@{}l@{}} \forall a, b, c, o. \\ \quad (\operatorname{is\text{-}local\text{-}measure}(a, b, c, o) \wedge (\emptyset \in c)) \Rightarrow  \\ \quad \quad \operatorname{is\text{-}measure}(a, b, \operatorname{restrict}(\operatorname{method\text{-}i\text{-}outer}(a, c, o), b)) \wedge \\ \quad \quad \forall e \in c.\; (\operatorname{restrict}(\operatorname{method\text{-}i\text{-}outer}(a, c, o), b)(e) = o(e)) \end{array}}}
\vnbprop{sigma-union-of-in-sigma}{\vnbstmt{\begin{array}{@{}l@{}} \forall b, c, d, o, a. \\ \quad (\operatorname{is\text{-}local\text{-}measure}(b, c, d, o) \wedge \operatorname{is\text{-}sigma\text{-}union\text{-}of}(d, a)) \Rightarrow (a \in c) \end{array}}}
\vnbprop{sigma-union-of-member}{\vnbstmt{\begin{array}{@{}l@{}} \forall a, b, c, o, e. \\ \quad (\operatorname{is\text{-}local\text{-}measure}(a, b, c, o) \wedge (e \in c)) \Rightarrow \operatorname{is\text{-}sigma\text{-}union\text{-}of}(c, e) \end{array}}}
\vnbprop{sigma-union-of-inter}{\vnbstmt{\begin{array}{@{}l@{}} \forall b, c, d, o, a, s. \\ \quad (\operatorname{is\text{-}local\text{-}measure}(b, c, d, o) \wedge (\operatorname{is\text{-}sigma\text{-}union\text{-}of}(d, a) \wedge (s \in c))) \Rightarrow  \\ \quad \quad \operatorname{is\text{-}sigma\text{-}union\text{-}of}(d, (a \cap s)) \end{array}}}
\vnbprop{sigma-union-of-countable}{\vnbstmt{\begin{array}{@{}l@{}} \forall a, b, c, o, h. \\ \quad (\operatorname{is\text{-}local\text{-}measure}(a, b, c, o) \wedge ((h \in (\mathbb{N} \to b)) \wedge \forall k \in \mathbb{N}.\; \operatorname{is\text{-}sigma\text{-}union\text{-}of}(c, h(k)))) \Rightarrow  \\ \quad \quad \operatorname{is\text{-}sigma\text{-}union\text{-}of}(c, \operatorname{big\text{-}union}(n, \mathbb{N}, h(n))) \end{array}}}
\vnbprop{locally-null-subset}{\vnbstmt{\begin{array}{@{}l@{}} \forall a, b, c, o, n, m. \\ \quad (\operatorname{is\text{-}local\text{-}measure}(a, b, c, o) \wedge (\operatorname{is\text{-}locally\text{-}null}(b, c, o, n) \wedge ((m \in b) \wedge (m \subseteq n)))) \Rightarrow  \\ \quad \quad \operatorname{is\text{-}locally\text{-}null}(b, c, o, m) \end{array}}}
\vnbprop{locally-null-empty}{\vnbstmt{\begin{array}{@{}l@{}} \forall a, b, c, o. \\ \quad (\operatorname{is\text{-}local\text{-}measure}(a, b, c, o) \wedge (\emptyset \in c)) \Rightarrow  \\ \quad \quad \operatorname{is\text{-}locally\text{-}null}(b, c, o, \emptyset) \end{array}}}
\vnbprop{locally-null-union}{\vnbstmt{\begin{array}{@{}l@{}} \forall a, b, c, o, f. \\ \quad \operatorname{is\text{-}local\text{-}measure}(a, b, c, o) \wedge \\ \quad f \in (\mathbb{N} \to b) \wedge \\ \quad \forall k \in \mathbb{N}.\; \operatorname{is\text{-}locally\text{-}null}(b, c, o, f(k)) \Rightarrow \\ \quad \quad \operatorname{is\text{-}locally\text{-}null}(b, c, o, \operatorname{big\text{-}union}(n, \mathbb{N}, f(n))) \end{array}}}
\vnbprop{locally-null-sigma-union-outer-null}{\vnbstmt{\begin{array}{@{}l@{}} \forall b, c, d, o, a. \\ \quad \operatorname{is\text{-}local\text{-}measure}(b, c, d, o) \wedge \\ \quad \emptyset \in d \wedge \\ \quad \operatorname{is\text{-}sigma\text{-}union\text{-}of}(d, a) \wedge \\ \quad \operatorname{is\text{-}locally\text{-}null}(c, d, o, a) \Rightarrow \\ \quad \quad \operatorname{method\text{-}i\text{-}outer}(b, d, o)(a) = 0 \end{array}}}
\vnbprop{essential-witness-le}{\vnbstmt{\begin{array}{@{}l@{}} \forall b, c, d, o, s, a, f. \\ \quad \operatorname{is\text{-}local\text{-}measure}(b, c, d, o) \wedge \\ \quad \emptyset \in d \wedge \\ \quad s \in c \wedge \\ \quad \operatorname{is\text{-}sigma\text{-}union\text{-}of}(d, a) \wedge \\ \quad a \subseteq s \wedge \\ \quad \operatorname{is\text{-}locally\text{-}null}(c, d, o, \operatorname{difference}(s, a)) \wedge \\ \quad \operatorname{is\text{-}sigma\text{-}union\text{-}of}(d, f) \wedge \\ \quad f \subseteq s \wedge \\ \quad \operatorname{is\text{-}locally\text{-}null}(c, d, o, \operatorname{difference}(s, f)) \Rightarrow \\ \quad \quad \operatorname{method\text{-}i\text{-}outer}(b, d, o)(a) \leq \operatorname{method\text{-}i\text{-}outer}(b, d, o)(f) \end{array}}}
\vnbprop{essential-witness-agree}{\vnbstmt{\begin{array}{@{}l@{}} \forall b, c, d, o, s, a, f. \\ \quad \operatorname{is\text{-}local\text{-}measure}(b, c, d, o) \wedge \\ \quad \emptyset \in d \wedge \\ \quad s \in c \wedge \\ \quad \operatorname{is\text{-}sigma\text{-}union\text{-}of}(d, a) \wedge \\ \quad a \subseteq s \wedge \\ \quad \operatorname{is\text{-}locally\text{-}null}(c, d, o, \operatorname{difference}(s, a)) \wedge \\ \quad \operatorname{is\text{-}sigma\text{-}union\text{-}of}(d, f) \wedge \\ \quad f \subseteq s \wedge \\ \quad \operatorname{is\text{-}locally\text{-}null}(c, d, o, \operatorname{difference}(s, f)) \Rightarrow \\ \quad \quad \operatorname{method\text{-}i\text{-}outer}(b, d, o)(a) = \operatorname{method\text{-}i\text{-}outer}(b, d, o)(f) \end{array}}}
\vnbprop{sigma-union-of-union}{\vnbstmt{\begin{array}{@{}l@{}} \forall c, d, f, o, a, b. \\ \quad (\operatorname{is\text{-}local\text{-}measure}(c, d, f, o) \wedge (\operatorname{is\text{-}sigma\text{-}union\text{-}of}(f, a) \wedge \operatorname{is\text{-}sigma\text{-}union\text{-}of}(f, b))) \Rightarrow  \\ \quad \quad \operatorname{is\text{-}sigma\text{-}union\text{-}of}(f, (a \cup b)) \end{array}}}
\vnbprop{essential-of-finite-extension}{\vnbstmt{\begin{array}{@{}l@{}} \forall b, c, d, o, u, s. \\ \quad \operatorname{is\text{-}local\text{-}measure}(b, c, d, o) \wedge \\ \quad \emptyset \in d \wedge \\ \quad \operatorname{is\text{-}measure}(b, c, u) \wedge \\ \quad \forall e \in d.\; (u(e) = o(e)) \wedge \\ \quad s \in c \wedge \\ \quad u(s) \in \mathbb{R} \Rightarrow \\ \quad \quad \exists a. \\ \quad \quad \quad (\operatorname{is\text{-}sigma\text{-}union\text{-}of}(d, a) \wedge ((a \subseteq s) \wedge \operatorname{is\text{-}locally\text{-}null}(c, d, o, \operatorname{difference}(s, a)))) \end{array}}}
\vnbprop{local-measure-min-extension}{\vnbstmt{\begin{array}{@{}l@{}} \forall b, c, d, o. \\ \quad (\operatorname{is\text{-}local\text{-}measure}(b, c, d, o) \wedge (\emptyset \in d)) \Rightarrow  \\ \quad \quad \exists m. \\ \quad \quad \quad \operatorname{is\text{-}measure}(b, c, m) \wedge \\ \quad \quad \quad \forall e \in d.\; (m(e) = o(e)) \wedge \\ \quad \quad \quad \forall u.; \forall s \in c. \\ \quad \quad \quad \quad (\operatorname{is\text{-}measure}(b, c, u) \wedge \forall e \in d.\; (u(e) = o(e))) \Rightarrow (m(s) \leq u(s)) \wedge \\ \quad \quad \quad \forall s \in c, a. \\ \quad \quad \quad \quad (\operatorname{is\text{-}sigma\text{-}union\text{-}of}(d, a) \wedge ((a \subseteq s) \wedge \operatorname{is\text{-}locally\text{-}null}(c, d, o, \operatorname{difference}(s, a)))) \Rightarrow  \\ \quad \quad \quad \quad \quad m(s) = \operatorname{method\text{-}i\text{-}outer}(b, d, o)(a) \wedge \\ \quad \quad \quad \forall s \in c. \\ \quad \quad \quad \quad \neg \exists \mathit{mma} \in c.\; (\operatorname{is\text{-}sigma\text{-}union\text{-}of}(d, \mathit{mma}) \wedge ((\mathit{mma} \subseteq s) \wedge \operatorname{is\text{-}locally\text{-}null}(c, d, o, \operatorname{difference}(s, \mathit{mma})))) \Rightarrow \\ \quad \quad \quad \quad \quad m(s) = \operatorname{pos\text{-}inf} \end{array}}}

\vnbbookfile{theorem-library/esum-tonelli}{theorem-library/esum-tonelli}
\noindent{\small esum-tonelli.scm -- wave of 2026-11-09/10 (agent TH3-B), wired 2026-11-10 from scratchpad/th3-1109/b/th3b-esum.scm. Library: measure.}\par
\vnbprop{esum-le-ptwise}{\vnbstmt{\begin{array}{@{}l@{}} \forall d, f, g. \\ \quad ((f \in (d \to \operatorname{rr\text{-}pos\text{-}star})) \wedge ((g \in (d \to \operatorname{rr\text{-}pos\text{-}star})) \wedge \forall z \in d.\; (f(z) \leq g(z)))) \Rightarrow  \\ \quad \quad \operatorname{esum}(f) \leq \operatorname{esum}(g) \end{array}}}
\vnbprop{esum-superadditive}{\vnbstmt{\begin{array}{@{}l@{}} \forall d, f, g, h. \\ \quad f \in (d \to \operatorname{rr\text{-}pos\text{-}star}) \wedge \\ \quad g \in (d \to \operatorname{rr\text{-}pos\text{-}star}) \wedge \\ \quad h \in (d \to \operatorname{rr\text{-}pos\text{-}star}) \wedge \\ \quad \forall z \in d.\; (\operatorname{eplus}(f(z), g(z)) \leq h(z)) \Rightarrow \\ \quad \quad \operatorname{eplus}(\operatorname{esum}(f), \operatorname{esum}(g)) \leq \operatorname{esum}(h) \end{array}}}
\vnbprop{esum-finsum-le}{\vnbstmt{\begin{array}{@{}l@{}} \forall x, b, g, t, a.; \forall s \in \mathbf{Set}. \\ \quad g \in (x \to (b \to \operatorname{rr\text{-}pos\text{-}star})) \wedge \\ \quad t \in (b \to (x \to \operatorname{rr\text{-}pos\text{-}star})) \wedge \\ \quad \forall i \in x, j \in b.\; (t(j)(i) = g(i)(j)) \wedge \\ \quad a \in (x \to \operatorname{rr\text{-}pos\text{-}star}) \wedge \\ \quad \forall i \in x.\; (a(i) = \operatorname{esum}(g(i))) \wedge \\ \quad |s| \in \mathbb{N} \wedge \\ \quad s \subseteq x \Rightarrow \\ \quad \quad \operatorname{finsum}(\operatorname{rr\text{-}pos\text{-}star\text{-}add\text{-}monoid}, a, s) \\ \quad \quad \leq\; \operatorname{esum}((\lambda \mathit{tbj} \in b.\; \operatorname{finsum}(\operatorname{rr\text{-}pos\text{-}star\text{-}add\text{-}monoid}, t(\mathit{tbj}), s))) \end{array}}}
\vnbprop{esum-swap-le}{\vnbstmt{\begin{array}{@{}l@{}} \forall x, c, g, t, a, b. \\ \quad g \in (x \to (c \to \operatorname{rr\text{-}pos\text{-}star})) \wedge \\ \quad t \in (c \to (x \to \operatorname{rr\text{-}pos\text{-}star})) \wedge \\ \quad \forall i \in x, j \in c.\; (t(j)(i) = g(i)(j)) \wedge \\ \quad a \in (x \to \operatorname{rr\text{-}pos\text{-}star}) \wedge \\ \quad \forall i \in x.\; (a(i) = \operatorname{esum}(g(i))) \wedge \\ \quad b \in (c \to \operatorname{rr\text{-}pos\text{-}star}) \wedge \\ \quad \forall j \in c.\; (b(j) = \operatorname{esum}(t(j))) \Rightarrow \\ \quad \quad \operatorname{esum}(a) \leq \operatorname{esum}(b) \end{array}}}
\vnbprop{esum-swap}{\vnbstmt{\begin{array}{@{}l@{}} \forall x, c, g, t, a, b. \\ \quad g \in (x \to (c \to \operatorname{rr\text{-}pos\text{-}star})) \wedge \\ \quad t \in (c \to (x \to \operatorname{rr\text{-}pos\text{-}star})) \wedge \\ \quad \forall i \in x, j \in c.\; (t(j)(i) = g(i)(j)) \wedge \\ \quad a \in (x \to \operatorname{rr\text{-}pos\text{-}star}) \wedge \\ \quad \forall i \in x.\; (a(i) = \operatorname{esum}(g(i))) \wedge \\ \quad b \in (c \to \operatorname{rr\text{-}pos\text{-}star}) \wedge \\ \quad \forall j \in c.\; (b(j) = \operatorname{esum}(t(j))) \Rightarrow \\ \quad \quad \operatorname{esum}(a) = \operatorname{esum}(b) \end{array}}}
\vnbprop{esum-subadditive}{\vnbstmt{\begin{array}{@{}l@{}} \forall d, f, g, h. \\ \quad f \in (d \to \operatorname{rr\text{-}pos\text{-}star}) \wedge \\ \quad g \in (d \to \operatorname{rr\text{-}pos\text{-}star}) \wedge \\ \quad h \in (d \to \operatorname{rr\text{-}pos\text{-}star}) \wedge \\ \quad \forall z \in d.\; (h(z) \leq \operatorname{eplus}(f(z), g(z))) \Rightarrow \\ \quad \quad \operatorname{esum}(h) \leq \operatorname{eplus}(\operatorname{esum}(f), \operatorname{esum}(g)) \end{array}}}

\vnbbookfile{theorem-library/esum-tail}{theorem-library/esum-tail}
\noindent{\small esum-tail.scm -- wave of 2026-11-09/10 (agent TH3-B2), wired 2026-11-10 from scratchpad/th3-1109/b2/tb2-tail.scm. Library: measure.}\par
\vnbprop{esum-tail-small}{\vnbstmt{\begin{array}{@{}l@{}} \forall d.; \forall f \in (d \to \operatorname{rr\text{-}pos\text{-}star}), l. \\ \quad ((\operatorname{esum}(f) \in \mathbb{R}) \wedge \operatorname{pos\text{-}rr}(l)) \Rightarrow  \\ \quad \quad \exists q \in \mathbf{Set}. \\ \quad \quad \quad (|q| \in \mathbb{N}) \wedge ((q \subseteq d) \wedge (\operatorname{esum}((\lambda z \in d.\; \left(\text{if }(z \in q)\text{ then }0\text{ else }f(z)\right))) < l)) \end{array}}}

\vnbbookfile{theorem-library/esum-cartesian}{theorem-library/esum-cartesian}
\noindent{\small esum-cartesian.scm -- wave of 2026-11-09/10 (agent TH3-B3), wired 2026-11-10 from scratchpad/th3-1109/b3/tb3-esum.scm. Library: measure.}\par
\vnbprop{finsum-cm-image-inj}{\vnbstmt{\begin{array}{@{}l@{}} \forall m.; \forall t \in \mathbf{Set}, i, f, k. \\ \quad \operatorname{is\text{-}comm\text{-}monoid}(m) \wedge \\ \quad |t| \in \mathbb{N} \wedge \\ \quad \forall x \in t.\; (i(x) \in \mathbf{Set}) \wedge \\ \quad \forall x, y \in t.\; ((i(x) = i(y)) \Rightarrow (x = y)) \wedge \\ \quad \forall z \in \operatorname{image}(i, t).\; (f(z) \in \operatorname{carr}(m)) \wedge \\ \quad \forall x \in t.\; (k(x) = f(i(x))) \Rightarrow \\ \quad \quad \operatorname{finsum}(m, k, t) = \operatorname{finsum}(m, f, \operatorname{image}(i, t)) \end{array}}}
\vnbprop{esum-cartesian-le}{\vnbstmt{\begin{array}{@{}l@{}} \forall i, j, h, g, a, b, c. \\ \quad h \in ((i \times j) \to \operatorname{rr\text{-}pos\text{-}star}) \wedge \\ \quad g \in (i \to (j \to \operatorname{rr\text{-}pos\text{-}star})) \wedge \\ \quad \forall d \in i, f \in j.\; (g(d)(f) = h({[}d, f])) \wedge \\ \quad a \in (i \to \operatorname{rr\text{-}pos\text{-}star}) \wedge \\ \quad \forall d \in i.\; (a(d) = \operatorname{esum}(g(d))) \wedge \\ \quad b \in ((i \times j) \to i) \wedge \\ \quad c \in ((i \times j) \to j) \wedge \\ \quad \forall w \in (i \times j).\; (w = {[}b(w), c(w)]) \Rightarrow \\ \quad \quad \operatorname{esum}(h) \leq \operatorname{esum}(a) \end{array}}}
\vnbprop{finsum-cm-cartesian}{\vnbstmt{\begin{array}{@{}l@{}} \forall m.; \forall s, u \in \mathbf{Set}, h, t, k, a, b. \\ \quad \operatorname{is\text{-}comm\text{-}monoid}(m) \wedge \\ \quad |s| \in \mathbb{N} \wedge \\ \quad |u| \in \mathbb{N} \wedge \\ \quad \forall z \in (s \times u).\; (h(z) \in \operatorname{carr}(m)) \wedge \\ \quad \forall j \in u, i \in s.\; (t(j)(i) \in \operatorname{carr}(m)) \wedge \\ \quad \forall j \in u, i \in s.\; (t(j)(i) = h({[}i, j])) \wedge \\ \quad \forall j \in u.\; (t(j) \in \mathbf{Set}) \wedge \\ \quad \forall j \in u.\; (k(j) \simeq \operatorname{finsum}(m, t(j), s)) \wedge \\ \quad \forall w \in (s \times u).\; (w = {[}a(w), b(w)]) \Rightarrow \\ \quad \quad \operatorname{finsum}(m, h, (s \times u)) = \operatorname{finsum}(m, k, u) \end{array}}}
\vnbprop{esum-cartesian-ge}{\vnbstmt{\begin{array}{@{}l@{}} \forall i, j, h, g, t, a, b, c. \\ \quad h \in ((i \times j) \to \operatorname{rr\text{-}pos\text{-}star}) \wedge \\ \quad g \in (i \to (j \to \operatorname{rr\text{-}pos\text{-}star})) \wedge \\ \quad \forall d \in i, f \in j.\; (g(d)(f) = h({[}d, f])) \wedge \\ \quad t \in (j \to (i \to \operatorname{rr\text{-}pos\text{-}star})) \wedge \\ \quad \forall d \in i, f \in j.\; (t(f)(d) = g(d)(f)) \wedge \\ \quad a \in (i \to \operatorname{rr\text{-}pos\text{-}star}) \wedge \\ \quad \forall d \in i.\; (a(d) = \operatorname{esum}(g(d))) \wedge \\ \quad b \in ((i \times j) \to i) \wedge \\ \quad c \in ((i \times j) \to j) \wedge \\ \quad \forall w \in (i \times j).\; (w = {[}b(w), c(w)]) \Rightarrow \\ \quad \quad \operatorname{esum}(a) \leq \operatorname{esum}(h) \end{array}}}
\vnbprop{esum-cartesian}{\vnbstmt{\begin{array}{@{}l@{}} \forall i.; \forall j \in \mathbf{Set}, h, g, a. \\ \quad h \in ((i \times j) \to \operatorname{rr\text{-}pos\text{-}star}) \wedge \\ \quad g \in (i \to (j \to \operatorname{rr\text{-}pos\text{-}star})) \wedge \\ \quad \forall b \in i, c \in j.\; (g(b)(c) = h({[}b, c])) \wedge \\ \quad a \in (i \to \operatorname{rr\text{-}pos\text{-}star}) \wedge \\ \quad \forall b \in i.\; (a(b) = \operatorname{esum}(g(b))) \Rightarrow \\ \quad \quad \operatorname{esum}(h) = \operatorname{esum}(a) \end{array}}}

\part{Operator Algebras}
\noindent Banach and C*-algebras, the Gelfand theory, bounded and unbounded operators on Hilbert space, operator topologies, compact and Hilbert-Schmidt operators (the user's BanachAlgebras notes, chapters 1-6; added 2026-11-10).
\par\noindent\textit{Builds on:} Complex Analysis, Measure.
\par\noindent{\small 128 files.}

\vnbbookfile{structure-library/th3a2-vocab}{structure-library/th3a2-vocab}
\noindent{\small th3a2-vocab.scm -- wave of 2026-11-09/10 (agent TH3-A2), wired 2026-11-10 from scratchpad/th3-1109/a2/th3a2-vocab.scm. Library: operator-algebras.}\par
\par\medskip\noindent\textbf{Definition} (\texttt{sub-cips}).\;$\operatorname{sub\text{-}cips}(v, \mathit{kk}) :=$
\begin{equation*}
{[}\operatorname{normed\text{-}field\text{-}as\text{-}commutative\text{-}ring}(\operatorname{cc\text{-}normed\text{-}field}), \mathit{kk}, (\lambda {[}\operatorname{tsbx\_}, \operatorname{tsby\_}] \in (\mathit{kk} \times \mathit{kk}).\; \operatorname{vadd}(v)(\operatorname{tsbx\_}, \operatorname{tsby\_})), \operatorname{vzero}(v), (\lambda \operatorname{tsbx\_} \in \mathit{kk}.\; \operatorname{vneg}(v)(\operatorname{tsbx\_})), (\lambda {[}\operatorname{tsbc\_}, \operatorname{tsbx\_}] \in (\mathbb{C} \times \mathit{kk}).\; \operatorname{act}(v)(\operatorname{tsbc\_}, \operatorname{tsbx\_})), (\lambda {[}\operatorname{tsbx\_}, \operatorname{tsby\_}] \in (\mathit{kk} \times \mathit{kk}).\; \langle\cdot,\cdot\rangle_{v}(\operatorname{tsbx\_}, \operatorname{tsby\_}))]
\end{equation*}


\vnbbookfile{structure-library/ba-ops-vocab}{structure-library/ba-ops-vocab}
\noindent{\small ba-ops-vocab.scm -- wave of 2026-11-09/10 (agent BA-BH), wired 2026-11-10 from scratchpad/babh-1109/ba-ops-vocab.scm. Library: operator-algebras.}\par
\par\medskip\noindent\textbf{Definition} (\texttt{bh-algebra}).\;$L({\operatorname{bov\_}}) :=$
\begin{equation*}
{[}\operatorname{normed\text{-}field\text{-}as\text{-}commutative\text{-}ring}(\operatorname{cc\text{-}normed\text{-}field}), L({\operatorname{bov\_}}, {\operatorname{bov\_}}), \operatorname{bh\text{-}vadd}(\operatorname{bov\_}), 0, \operatorname{bh\text{-}vneg}(\operatorname{bov\_}), \operatorname{bh\text{-}act}(\operatorname{bov\_}), \operatorname{bh\text{-}vnrm}(\operatorname{bov\_}), \operatorname{bh\text{-}vmul}(\operatorname{bov\_}), 1, \operatorname{bh\text{-}vstar}(\operatorname{bov\_})]
\end{equation*}

\par\medskip\noindent\textbf{Definition} (\texttt{bounded-ops}).\;$L({\operatorname{bov\_}}, {\operatorname{bow\_}}) :=$
\begin{equation*}
\{\operatorname{bot\_} \in (\operatorname{vec}(\operatorname{bov\_}) \to \operatorname{vec}(\operatorname{bow\_})) : \operatorname{is\text{-}bounded\text{-}op}(\operatorname{bov\_}, \operatorname{bow\_}, \operatorname{bot\_})\}
\end{equation*}

\par\medskip\noindent\textbf{Definition} (\texttt{opnorm}).\;$\lVert {\operatorname{bot\_}} \rVert :=$
\begin{equation*}
\operatorname{sup}(\{\operatorname{boa\_} \in \mathbb{R} : ((\operatorname{boa\_} = 0) \vee \exists \operatorname{box\_} \in \operatorname{vec}(\operatorname{bov\_}).\; (\neg (\lVert \operatorname{box\_} \rVert = 0) \wedge (\operatorname{boa\_} = (\lVert \operatorname{bot\_}(\operatorname{box\_}) \rVert \cdot \operatorname{recip}(\lVert \operatorname{box\_} \rVert)))))\})
\end{equation*}

\par\medskip\noindent\textbf{Definition} (\texttt{bop-sum}).\;${\operatorname{bos\_}} + {\operatorname{bot\_}} :=$
\begin{equation*}
(\lambda \operatorname{box\_} \in \operatorname{vec}(\operatorname{bov\_}).\; \operatorname{vadd}(\operatorname{bow\_})(\operatorname{bos\_}(\operatorname{box\_}), \operatorname{bot\_}(\operatorname{box\_})))
\end{equation*}

\par\medskip\noindent\textbf{Definition} (\texttt{bop-scale}).\;${\operatorname{boc\_}}\,{\operatorname{bot\_}} :=$ $(\lambda \operatorname{box\_} \in \operatorname{vec}(\operatorname{bov\_}).\; \operatorname{act}(\operatorname{bow\_})(\operatorname{boc\_}, \operatorname{bot\_}(\operatorname{box\_})))$
\par\medskip\noindent\textbf{Definition} (\texttt{hsum}).\;${\operatorname{bov\_}} \oplus {\operatorname{bow\_}} :=$
\begin{equation*}
{[}\operatorname{normed\text{-}field\text{-}as\text{-}commutative\text{-}ring}(\operatorname{cc\text{-}normed\text{-}field}), (\operatorname{vec}(\operatorname{bov\_}) \times \operatorname{vec}(\operatorname{bow\_})), \operatorname{hsum\text{-}vadd}(\operatorname{bov\_}, \operatorname{bow\_}), {[}\operatorname{vzero}(\operatorname{bov\_}), \operatorname{vzero}(\operatorname{bow\_})], \operatorname{hsum\text{-}vneg}(\operatorname{bov\_}, \operatorname{bow\_}), \operatorname{hsum\text{-}act}(\operatorname{bov\_}, \operatorname{bow\_}), \operatorname{hsum\text{-}ip}(\operatorname{bov\_}, \operatorname{bow\_})]
\end{equation*}

\par\medskip\noindent\textbf{Definition} (\texttt{is-linear-op}).\;For all $\operatorname{bov\_}$, $\operatorname{bow\_}$ and $\operatorname{bot\_}$, 
\begin{equation*}
\begin{array}{@{}l@{}} \operatorname{is\text{-}linear\text{-}op}(\operatorname{bov\_}, \operatorname{bow\_}, \operatorname{bot\_}) \\ \Leftrightarrow\; \operatorname{bot\_} \in (\operatorname{vec}(\operatorname{bov\_}) \to \operatorname{vec}(\operatorname{bow\_})) \wedge \\ \forall \operatorname{box\_}, \operatorname{boy\_} \in \operatorname{vec}(\operatorname{bov\_}). \\ \quad \operatorname{bot\_}(\operatorname{vadd}(\operatorname{bov\_})(\operatorname{box\_}, \operatorname{boy\_})) = \operatorname{vadd}(\operatorname{bow\_})(\operatorname{bot\_}(\operatorname{box\_}), \operatorname{bot\_}(\operatorname{boy\_})) \wedge \\ \forall \operatorname{boc\_} \in \mathbb{C}, \operatorname{box\_} \in \operatorname{vec}(\operatorname{bov\_}). \\ \quad \operatorname{bot\_}(\operatorname{act}(\operatorname{bov\_})(\operatorname{boc\_}, \operatorname{box\_})) = \operatorname{act}(\operatorname{bow\_})(\operatorname{boc\_}, \operatorname{bot\_}(\operatorname{box\_})) \end{array}
\end{equation*}

\par\medskip\noindent\textbf{Definition} (\texttt{is-bounded-op}).\;For all $\operatorname{bov\_}$, $\operatorname{bow\_}$ and $\operatorname{bot\_}$, 
\begin{equation*}
\begin{array}{@{}l@{}} \operatorname{is\text{-}bounded\text{-}op}(\operatorname{bov\_}, \operatorname{bow\_}, \operatorname{bot\_}) \\ \Leftrightarrow\; \operatorname{is\text{-}linear\text{-}op}(\operatorname{bov\_}, \operatorname{bow\_}, \operatorname{bot\_}) \wedge \\ \exists \operatorname{bom\_} \in \mathbb{R}. \\ \quad ((0 \leq \operatorname{bom\_}) \wedge \forall \operatorname{box\_} \in \operatorname{vec}(\operatorname{bov\_}).\; (\lVert \operatorname{bot\_}(\operatorname{box\_}) \rVert \leq \cdot(\operatorname{bom\_}, \lVert \operatorname{box\_} \rVert))) \end{array}
\end{equation*}

\par\medskip\noindent\textbf{Definition} (\texttt{bop-zero}).\;$0 :=$ $(\lambda \operatorname{box\_} \in \operatorname{vec}(\operatorname{bov\_}).\; \operatorname{vzero}(\operatorname{bow\_}))$
\par\medskip\noindent\textbf{Definition} (\texttt{bop-neg}).\;$-{\operatorname{bot\_}} :=$ $(\lambda \operatorname{box\_} \in \operatorname{vec}(\operatorname{bov\_}).\; \operatorname{vneg}(\operatorname{bow\_})(\operatorname{bot\_}(\operatorname{box\_})))$
\par\medskip\noindent\textbf{Definition} (\texttt{bop-comp}).\;${\operatorname{bos\_}}{\operatorname{bot\_}} :=$ $(\lambda \operatorname{box\_} \in \operatorname{vec}(\operatorname{bov\_}).\; \operatorname{bos\_}(\operatorname{bot\_}(\operatorname{box\_})))$
\par\medskip\noindent\textbf{Definition} (\texttt{bop-id}).\;$1 :=$ $(\lambda \operatorname{box\_} \in \operatorname{vec}(\operatorname{bov\_}).\; \operatorname{box\_})$
\par\medskip\noindent\textbf{Definition} (\texttt{op-adjoint}).\;${\operatorname{bot\_}}^{*} :=$
\begin{equation*}
\begin{array}{@{}l@{}} (\lambda \operatorname{boy\_} \in \operatorname{vec}(\operatorname{bow\_}).\; \\ \quad (\iota\, \operatorname{boz\_}.\;((\operatorname{boz\_} \in \operatorname{vec}(\operatorname{bov\_})) \wedge \forall \operatorname{box\_} \in \operatorname{vec}(\operatorname{bov\_}).\; (\langle\cdot,\cdot\rangle_{\operatorname{bow\_}}(\operatorname{bot\_}(\operatorname{box\_}), \operatorname{boy\_}) = \langle\cdot,\cdot\rangle_{\operatorname{bov\_}}(\operatorname{box\_}, \operatorname{boz\_}))))) \end{array}
\end{equation*}

\par\medskip\noindent\textbf{Definition} (\texttt{bh-vadd}).\;$\operatorname{bh\text{-}vadd}(\operatorname{bov\_}) :=$
\begin{equation*}
\begin{array}{@{}l@{}} (\lambda {[}\operatorname{bos\_}, \operatorname{bot\_}] \in (L({\operatorname{bov\_}}, {\operatorname{bov\_}}) \times L({\operatorname{bov\_}}, {\operatorname{bov\_}})).\; \\ \quad {\operatorname{bos\_}} + {\operatorname{bot\_}}) \end{array}
\end{equation*}

\par\medskip\noindent\textbf{Definition} (\texttt{bh-vneg}).\;$\operatorname{bh\text{-}vneg}(\operatorname{bov\_}) :=$
\begin{equation*}
(\lambda \operatorname{bot\_} \in L({\operatorname{bov\_}}, {\operatorname{bov\_}}).\; -{\operatorname{bot\_}})
\end{equation*}

\par\medskip\noindent\textbf{Definition} (\texttt{bh-act}).\;$\operatorname{bh\text{-}act}(\operatorname{bov\_}) :=$
\begin{equation*}
(\lambda {[}\operatorname{boc\_}, \operatorname{bot\_}] \in (\mathbb{C} \times L({\operatorname{bov\_}}, {\operatorname{bov\_}})).\; {\operatorname{boc\_}}\,{\operatorname{bot\_}})
\end{equation*}

\par\medskip\noindent\textbf{Definition} (\texttt{bh-vnrm}).\;$\operatorname{bh\text{-}vnrm}(\operatorname{bov\_}) :=$
\begin{equation*}
(\lambda \operatorname{bot\_} \in L({\operatorname{bov\_}}, {\operatorname{bov\_}}).\; \lVert {\operatorname{bot\_}} \rVert)
\end{equation*}

\par\medskip\noindent\textbf{Definition} (\texttt{bh-vmul}).\;$\operatorname{bh\text{-}vmul}(\operatorname{bov\_}) :=$
\begin{equation*}
\begin{array}{@{}l@{}} (\lambda {[}\operatorname{bos\_}, \operatorname{bot\_}] \in (L({\operatorname{bov\_}}, {\operatorname{bov\_}}) \times L({\operatorname{bov\_}}, {\operatorname{bov\_}})).\; \\ \quad {\operatorname{bos\_}}{\operatorname{bot\_}}) \end{array}
\end{equation*}

\par\medskip\noindent\textbf{Definition} (\texttt{bh-vstar}).\;$\operatorname{bh\text{-}vstar}(\operatorname{bov\_}) :=$
\begin{equation*}
(\lambda \operatorname{bot\_} \in L({\operatorname{bov\_}}, {\operatorname{bov\_}}).\; {\operatorname{bot\_}}^{*})
\end{equation*}

\par\medskip\noindent\textbf{Definition} (\texttt{op-range}).\;$\mathrm{ran}\,{\operatorname{bot\_}} :=$
\begin{equation*}
\{\operatorname{boy\_} \in \operatorname{vec}(\operatorname{bow\_}) : \exists \operatorname{box\_} \in \operatorname{vec}(\operatorname{bov\_}).\; (\operatorname{boy\_} = \operatorname{bot\_}(\operatorname{box\_}))\}
\end{equation*}

\par\medskip\noindent\textbf{Definition} (\texttt{op-kernel}).\;$\ker {\operatorname{bot\_}} :=$ $\{\operatorname{box\_} \in \operatorname{vec}(\operatorname{bov\_}) : (\operatorname{bot\_}(\operatorname{box\_}) = \operatorname{vzero}(\operatorname{bow\_}))\}$
\par\medskip\noindent\textbf{Definition} (\texttt{op-support}).\;$\mathrm{supp}\,{\operatorname{bot\_}} :=$ ${\ker {\operatorname{bot\_}}}^{\perp}$
\par\medskip\noindent\textbf{Definition} (\texttt{proj-op}).\;$E_{\operatorname{bok\_}} :=$ $(\iota\, \operatorname{bop\_}.\;\operatorname{is\text{-}orthogonal\text{-}projection}(\operatorname{bov\_}, \operatorname{bok\_}, \operatorname{bop\_}))$
\par\medskip\noindent\textbf{Definition} (\texttt{is-projection-op}).\;For all $\operatorname{bov\_}$ and $\operatorname{bot\_}$, 
\begin{equation*}
\begin{array}{@{}l@{}} \operatorname{is\text{-}projection\text{-}op}(\operatorname{bov\_}, \operatorname{bot\_}) \\ \Leftrightarrow\; \operatorname{bot\_} \in L({\operatorname{bov\_}}, {\operatorname{bov\_}}) \wedge \\ {\operatorname{bot\_}}^{*} = \operatorname{bot\_} \wedge \\ {\operatorname{bot\_}}{\operatorname{bot\_}} = \operatorname{bot\_} \end{array}
\end{equation*}

\par\medskip\noindent\textbf{Definition} (\texttt{is-isometric-on}).\;For all $\operatorname{bov\_}$, $\operatorname{bow\_}$, $\operatorname{bot\_}$ and $\operatorname{bok\_}$, 
\begin{equation*}
\begin{array}{@{}l@{}} \operatorname{is\text{-}isometric\text{-}on}(\operatorname{bov\_}, \operatorname{bow\_}, \operatorname{bot\_}, \operatorname{bok\_}) \\ \Leftrightarrow\; \forall \operatorname{box\_} \in \operatorname{bok\_}.\; (\lVert \operatorname{bot\_}(\operatorname{box\_}) \rVert = \lVert \operatorname{box\_} \rVert) \end{array}
\end{equation*}

\par\medskip\noindent\textbf{Definition} (\texttt{is-partial-isometry}).\;For all $\operatorname{bov\_}$ and $\operatorname{bot\_}$, 
\begin{equation*}
\begin{array}{@{}l@{}} \operatorname{is\text{-}partial\text{-}isometry}(\operatorname{bov\_}, \operatorname{bot\_}) \\ \Leftrightarrow\; \operatorname{bot\_} \in L({\operatorname{bov\_}}, {\operatorname{bov\_}}) \wedge \\ \exists \operatorname{bok\_}. \\ \quad \operatorname{is\text{-}submodule}(\operatorname{bov\_}, \operatorname{bok\_}) \wedge \\ \quad \operatorname{is\text{-}closed}(\operatorname{ip\text{-}metric\text{-}space}(\operatorname{bov\_}), \operatorname{bok\_}) \wedge \\ \quad \operatorname{is\text{-}isometric\text{-}on}(\operatorname{bov\_}, \operatorname{bov\_}, \operatorname{bot\_}, \operatorname{bok\_}) \wedge \\ \quad \forall \operatorname{box\_} \in {\operatorname{bok\_}}^{\perp}.\; (\operatorname{bot\_}(\operatorname{box\_}) = \operatorname{vzero}(\operatorname{bov\_})) \end{array}
\end{equation*}

\par\medskip\noindent\textbf{Definition} (\texttt{hsum-vec}).\;$\operatorname{hsum\text{-}vec}(\operatorname{bov\_}, \operatorname{bow\_}) :=$ $(\operatorname{vec}(\operatorname{bov\_}) \times \operatorname{vec}(\operatorname{bow\_}))$
\par\medskip\noindent\textbf{Definition} (\texttt{hsum-vadd}).\;$\operatorname{hsum\text{-}vadd}(\operatorname{bov\_}, \operatorname{bow\_}) :=$
\begin{equation*}
\begin{array}{@{}l@{}} (\lambda {[}\operatorname{bhp\_}, \operatorname{bhq\_}] \in ((\operatorname{vec}(\operatorname{bov\_}) \times \operatorname{vec}(\operatorname{bow\_})) \times (\operatorname{vec}(\operatorname{bov\_}) \times \operatorname{vec}(\operatorname{bow\_}))).\; \\ \quad {[}\operatorname{vadd}(\operatorname{bov\_})(\operatorname{nth}(1, \operatorname{bhp\_}), \operatorname{nth}(1, \operatorname{bhq\_})), \operatorname{vadd}(\operatorname{bow\_})(\operatorname{nth}(2, \operatorname{bhp\_}), \operatorname{nth}(2, \operatorname{bhq\_}))]) \end{array}
\end{equation*}

\par\medskip\noindent\textbf{Definition} (\texttt{hsum-vneg}).\;$\operatorname{hsum\text{-}vneg}(\operatorname{bov\_}, \operatorname{bow\_}) :=$
\begin{equation*}
\begin{array}{@{}l@{}} (\lambda \operatorname{bhp\_} \in (\operatorname{vec}(\operatorname{bov\_}) \times \operatorname{vec}(\operatorname{bow\_})).\; \\ \quad {[}\operatorname{vneg}(\operatorname{bov\_})(\operatorname{nth}(1, \operatorname{bhp\_})), \operatorname{vneg}(\operatorname{bow\_})(\operatorname{nth}(2, \operatorname{bhp\_}))]) \end{array}
\end{equation*}

\par\medskip\noindent\textbf{Definition} (\texttt{hsum-act}).\;$\operatorname{hsum\text{-}act}(\operatorname{bov\_}, \operatorname{bow\_}) :=$
\begin{equation*}
\begin{array}{@{}l@{}} (\lambda {[}\operatorname{bhc\_}, \operatorname{bhp\_}] \in (\mathbb{C} \times (\operatorname{vec}(\operatorname{bov\_}) \times \operatorname{vec}(\operatorname{bow\_}))).\; \\ \quad {[}\operatorname{act}(\operatorname{bov\_})(\operatorname{bhc\_}, \operatorname{nth}(1, \operatorname{bhp\_})), \operatorname{act}(\operatorname{bow\_})(\operatorname{bhc\_}, \operatorname{nth}(2, \operatorname{bhp\_}))]) \end{array}
\end{equation*}

\par\medskip\noindent\textbf{Definition} (\texttt{hsum-ip}).\;$\operatorname{hsum\text{-}ip}(\operatorname{bov\_}, \operatorname{bow\_}) :=$
\begin{equation*}
\begin{array}{@{}l@{}} (\lambda {[}\operatorname{bhp\_}, \operatorname{bhq\_}] \in ((\operatorname{vec}(\operatorname{bov\_}) \times \operatorname{vec}(\operatorname{bow\_})) \times (\operatorname{vec}(\operatorname{bov\_}) \times \operatorname{vec}(\operatorname{bow\_}))).\; \\ \quad \langle\cdot,\cdot\rangle_{\operatorname{bov\_}}(\operatorname{nth}(1, \operatorname{bhp\_}), \operatorname{nth}(1, \operatorname{bhq\_})) + \langle\cdot,\cdot\rangle_{\operatorname{bow\_}}(\operatorname{nth}(2, \operatorname{bhp\_}), \operatorname{nth}(2, \operatorname{bhq\_}))) \end{array}
\end{equation*}


\vnbbookfile{structure-library/ba-bh2-vocab}{structure-library/ba-bh2-vocab}
\noindent{\small ba-bh2-vocab.scm -- wave of 2026-11-09/10 (agent BA-BH2), wired 2026-11-10 from scratchpad/babh2-1109/ba-bh2-vocab.scm. Library: operator-algebras.}\par
\par\medskip\noindent\textbf{Definition} (\texttt{bop-normed-ag}).\;$\operatorname{bop\text{-}normed\text{-}ag}(\operatorname{bov\_}, \operatorname{bow\_}) :=$
\begin{equation*}
{[}L({\operatorname{bov\_}}, {\operatorname{bow\_}}), \operatorname{bop\text{-}vadd}(\operatorname{bov\_}, \operatorname{bow\_}), 0, \operatorname{bop\text{-}vinv}(\operatorname{bov\_}, \operatorname{bow\_}), \operatorname{bop\text{-}vnrm}(\operatorname{bov\_}, \operatorname{bow\_})]
\end{equation*}

\par\medskip\noindent\textbf{Definition} (\texttt{bop-vadd}).\;$\operatorname{bop\text{-}vadd}(\operatorname{bov\_}, \operatorname{bow\_}) :=$
\begin{equation*}
\begin{array}{@{}l@{}} (\lambda {[}\operatorname{bos\_}, \operatorname{bot\_}] \in (L({\operatorname{bov\_}}, {\operatorname{bow\_}}) \times L({\operatorname{bov\_}}, {\operatorname{bow\_}})).\; \\ \quad {\operatorname{bos\_}} + {\operatorname{bot\_}}) \end{array}
\end{equation*}

\par\medskip\noindent\textbf{Definition} (\texttt{bop-vinv}).\;$\operatorname{bop\text{-}vinv}(\operatorname{bov\_}, \operatorname{bow\_}) :=$
\begin{equation*}
(\lambda \operatorname{bot\_} \in L({\operatorname{bov\_}}, {\operatorname{bow\_}}).\; -{\operatorname{bot\_}})
\end{equation*}

\par\medskip\noindent\textbf{Definition} (\texttt{bop-vnrm}).\;$\operatorname{bop\text{-}vnrm}(\operatorname{bov\_}, \operatorname{bow\_}) :=$
\begin{equation*}
(\lambda \operatorname{bot\_} \in L({\operatorname{bov\_}}, {\operatorname{bow\_}}).\; \lVert {\operatorname{bot\_}} \rVert)
\end{equation*}


\vnbbookfile{structure-library/th3-vocab}{structure-library/th3-vocab}
\noindent{\small th3-vocab.scm -- wave of 2026-11-09/10 (agent TH3-A), wired 2026-11-10 from scratchpad/th3-1109/th3-vocab.scm. Library: operator-algebras.}\par
\par\medskip\noindent\textbf{Definition} (\texttt{compact-ops}).\;$K({\operatorname{kqv\_}}, {\operatorname{kqw\_}}) :=$
\begin{equation*}
\{\operatorname{kqt\_} \in L({\operatorname{kqv\_}}, {\operatorname{kqw\_}}) : \operatorname{is\text{-}compact\text{-}op}(\operatorname{kqv\_}, \operatorname{kqw\_}, \operatorname{kqt\_})\}
\end{equation*}

\par\medskip\noindent\textbf{Definition} (\texttt{is-compact-op}).\;For all $\operatorname{kqv\_}$, $\operatorname{kqw\_}$ and $\operatorname{kqt\_}$, 
\begin{equation*}
\begin{array}{@{}l@{}} \operatorname{is\text{-}compact\text{-}op}(\operatorname{kqv\_}, \operatorname{kqw\_}, \operatorname{kqt\_}) \\ \Leftrightarrow\; \operatorname{kqt\_} \in L({\operatorname{kqv\_}}, {\operatorname{kqw\_}}) \wedge \\ \forall \operatorname{kqf\_} \in (\mathbb{N} \to \operatorname{vec}(\operatorname{kqv\_})). \\ \quad \forall \operatorname{kqn\_} \in \mathbb{N}.\; (\lVert \operatorname{kqf\_}(\operatorname{kqn\_}) \rVert \leq 1) \Rightarrow  \\ \quad \quad \exists \operatorname{kqp\_}. \\ \quad \quad \quad \operatorname{strictly\text{-}mono\text{-}nn}(\operatorname{kqp\_}) \wedge \\ \quad \quad \quad \exists \operatorname{kql\_}. \\ \quad \quad \quad \quad \operatorname{converges\text{-}to}(\operatorname{ip\text{-}metric\text{-}space}(\operatorname{kqw\_}), (\lambda \operatorname{kqn\_} \in \mathbb{N}.\; \operatorname{kqt\_}(\operatorname{kqf\_}(\operatorname{kqp\_}(\operatorname{kqn\_})))), \operatorname{kql\_}) \end{array}
\end{equation*}

\par\medskip\noindent\textbf{Definition} (\texttt{hs-norm}).\;$\lVert {\operatorname{kqt\_}} \rVert_{HS} :=$
\begin{equation*}
\sqrt{\sum_{i \in {\operatorname{canon\text{-}onb}(\operatorname{kqv\_})}} \lVert {\operatorname{kqt\_}} {(\lambda \operatorname{kqe\_} \in \operatorname{canon\text{-}onb}(\operatorname{kqv\_}).\; \operatorname{kqe\_})}_i \rVert^2}
\end{equation*}

\par\medskip\noindent\textbf{Definition} (\texttt{hs-ip}).\;$\langle {\operatorname{kqu\_}} \mid {\operatorname{kqt\_}} \rangle_{HS} :=$
\begin{equation*}
\operatorname{cc\text{-}family\text{-}sum}((\lambda \operatorname{kqe\_} \in \operatorname{canon\text{-}onb}(\operatorname{kqv\_}).\; \langle\cdot,\cdot\rangle_{\operatorname{kqw\_}}(\operatorname{kqu\_}(\operatorname{kqe\_}), \operatorname{kqt\_}(\operatorname{kqe\_}))))
\end{equation*}

\par\medskip\noindent\textbf{Definition} (\texttt{op-trace}).\;$\mathrm{tr}\,{\operatorname{kqt\_}} :=$
\begin{equation*}
\operatorname{cc\text{-}family\text{-}sum}((\lambda \operatorname{kqe\_} \in \operatorname{canon\text{-}onb}(\operatorname{kqv\_}).\; \langle\cdot,\cdot\rangle_{\operatorname{kqv\_}}(\operatorname{kqt\_}(\operatorname{kqe\_}), \operatorname{kqe\_})))
\end{equation*}

\par\medskip\noindent\textbf{Definition} (\texttt{canon-onb}).\;$\operatorname{canon\text{-}onb}(\operatorname{kqv\_}) :=$ $\varepsilon(\operatorname{onb\text{-}sets}(\operatorname{kqv\_}))$
\par\medskip\noindent\textbf{Definition} (\texttt{onb-sets}).\;$\operatorname{onb\text{-}sets}(\operatorname{kqv\_}) :=$
\begin{equation*}
\{\operatorname{kqo\_} \in \operatorname{power}(\operatorname{vec}(\operatorname{kqv\_})) : \operatorname{is\text{-}orthonormal\text{-}basis}(\operatorname{kqv\_}, \operatorname{kqo\_}, (\lambda \operatorname{kqe\_} \in \operatorname{kqo\_}.\; \operatorname{kqe\_}))\}
\end{equation*}

\par\medskip\noindent\textbf{Definition} (\texttt{hs-sum}).\;$\sum_{i \in {\operatorname{kqi\_}}} \lVert {\operatorname{kqt\_}} {\operatorname{kqe\_}}_i \rVert^2 :=$
\begin{equation*}
\operatorname{esum}((\lambda \operatorname{kqk\_} \in \operatorname{kqi\_}.\; (\lVert \operatorname{kqt\_}(\operatorname{kqe\_}(\operatorname{kqk\_})) \rVert \cdot \lVert \operatorname{kqt\_}(\operatorname{kqe\_}(\operatorname{kqk\_})) \rVert)))
\end{equation*}

\par\medskip\noindent\textbf{Definition} (\texttt{is-hilbert-schmidt}).\;For all $\operatorname{kqv\_}$, $\operatorname{kqw\_}$ and $\operatorname{kqt\_}$, 
\begin{equation*}
\begin{array}{@{}l@{}} \operatorname{is\text{-}hilbert\text{-}schmidt}(\operatorname{kqv\_}, \operatorname{kqw\_}, \operatorname{kqt\_}) \\ \Leftrightarrow\; \operatorname{kqt\_} \in L({\operatorname{kqv\_}}, {\operatorname{kqw\_}}) \wedge \\ \exists \operatorname{kqi\_} \in \mathbf{Set}, \operatorname{kqe\_} \in (\operatorname{kqi\_} \to \operatorname{vec}(\operatorname{kqv\_})). \\ \quad (\operatorname{is\text{-}orthonormal\text{-}basis}(\operatorname{kqv\_}, \operatorname{kqi\_}, \operatorname{kqe\_}) \wedge (\sum_{i \in {\operatorname{kqi\_}}} \lVert {\operatorname{kqt\_}} {\operatorname{kqe\_}}_i \rVert^2 \in \mathbb{R})) \end{array}
\end{equation*}

\par\medskip\noindent\textbf{Definition} (\texttt{hs-ops}).\;$HS({\operatorname{kqv\_}}, {\operatorname{kqw\_}}) :=$
\begin{equation*}
\{\operatorname{kqt\_} \in L({\operatorname{kqv\_}}, {\operatorname{kqw\_}}) : \operatorname{is\text{-}hilbert\text{-}schmidt}(\operatorname{kqv\_}, \operatorname{kqw\_}, \operatorname{kqt\_})\}
\end{equation*}

\par\medskip\noindent\textbf{Definition} (\texttt{is-trace-class}).\;For all $\operatorname{kqv\_}$ and $\operatorname{kqt\_}$, 
\begin{equation*}
\begin{array}{@{}l@{}} \operatorname{is\text{-}trace\text{-}class}(\operatorname{kqv\_}, \operatorname{kqt\_}) \\ \Leftrightarrow\; \operatorname{kqt\_} \in L({\operatorname{kqv\_}}, {\operatorname{kqv\_}}) \wedge \\ \exists \operatorname{kqa\_}. \\ \quad \operatorname{is\text{-}hilbert\text{-}schmidt}(\operatorname{kqv\_}, \operatorname{kqv\_}, \operatorname{kqa\_}) \wedge \\ \quad \exists \operatorname{kqb\_}. \\ \quad \quad \operatorname{is\text{-}hilbert\text{-}schmidt}(\operatorname{kqv\_}, \operatorname{kqv\_}, \operatorname{kqb\_}) \wedge (\operatorname{kqt\_} = {\operatorname{kqa\_}}{\operatorname{kqb\_}}) \end{array}
\end{equation*}

\par\medskip\noindent\textbf{Definition} (\texttt{is-self-adjoint-op}).\;For all $\operatorname{kqv\_}$ and $\operatorname{kqt\_}$, 
\begin{equation*}
\begin{array}{@{}l@{}} \operatorname{is\text{-}self\text{-}adjoint\text{-}op}(\operatorname{kqv\_}, \operatorname{kqt\_}) \\ \Leftrightarrow\; \operatorname{kqt\_} \in L({\operatorname{kqv\_}}, {\operatorname{kqv\_}}) \wedge \\ \forall \operatorname{kqx\_}, \operatorname{kqy\_} \in \operatorname{vec}(\operatorname{kqv\_}). \\ \quad \langle\cdot,\cdot\rangle_{\operatorname{kqv\_}}(\operatorname{kqt\_}(\operatorname{kqx\_}), \operatorname{kqy\_}) = \langle\cdot,\cdot\rangle_{\operatorname{kqv\_}}(\operatorname{kqx\_}, \operatorname{kqt\_}(\operatorname{kqy\_})) \end{array}
\end{equation*}

\par\medskip\noindent\textbf{Definition} (\texttt{is-positive-op}).\;For all $\operatorname{kqv\_}$ and $\operatorname{kqt\_}$, 
\begin{equation*}
\begin{array}{@{}l@{}} \operatorname{is\text{-}positive\text{-}op}(\operatorname{kqv\_}, \operatorname{kqt\_}) \\ \Leftrightarrow\; \operatorname{is\text{-}self\text{-}adjoint\text{-}op}(\operatorname{kqv\_}, \operatorname{kqt\_}) \wedge \\ \forall \operatorname{kqx\_} \in \operatorname{vec}(\operatorname{kqv\_}). \\ \quad (\langle\cdot,\cdot\rangle_{\operatorname{kqv\_}}(\operatorname{kqt\_}(\operatorname{kqx\_}), \operatorname{kqx\_}) \in \mathbb{R}) \wedge (0 \leq \langle\cdot,\cdot\rangle_{\operatorname{kqv\_}}(\operatorname{kqt\_}(\operatorname{kqx\_}), \operatorname{kqx\_})) \end{array}
\end{equation*}

\par\medskip\noindent\textbf{Definition} (\texttt{op-sqrt}).\;${\operatorname{kqt\_}}^{1/2} :=$
\begin{equation*}
(\iota\, \operatorname{kqs\_}.\;(\operatorname{is\text{-}positive\text{-}op}(\operatorname{kqv\_}, \operatorname{kqs\_}) \wedge ({\operatorname{kqs\_}}{\operatorname{kqs\_}} = \operatorname{kqt\_})))
\end{equation*}

\par\medskip\noindent\textbf{Definition} (\texttt{abs-op}).\;$\lvert {\operatorname{kqt\_}} \rvert :=$
\begin{equation*}
{{{\operatorname{kqt\_}}^{*}}{\operatorname{kqt\_}}}^{1/2}
\end{equation*}

\par\medskip\noindent\textbf{Definition} (\texttt{eigenspace}).\;$N_{\operatorname{kqc\_}}({\operatorname{kqt\_}}) :=$ $\{\operatorname{kqx\_} \in \operatorname{vec}(\operatorname{kqv\_}) : (\operatorname{kqt\_}(\operatorname{kqx\_}) = \operatorname{act}(\operatorname{kqv\_})(\operatorname{kqc\_}, \operatorname{kqx\_}))\}$
\par\medskip\noindent\textbf{Definition} (\texttt{rank-one-op}).\;$\langle \cdot \mid {\operatorname{kqa\_}} \rangle {\operatorname{kqb\_}} :=$
\begin{equation*}
(\lambda \operatorname{kqx\_} \in \operatorname{vec}(\operatorname{kqv\_}).\; \operatorname{act}(\operatorname{kqw\_})(\langle\cdot,\cdot\rangle_{\operatorname{kqv\_}}(\operatorname{kqx\_}, \operatorname{kqa\_}), \operatorname{kqb\_}))
\end{equation*}

\par\medskip\noindent\textbf{Definition} (\texttt{is-finite-rank-op}).\;For all $\operatorname{kqv\_}$, $\operatorname{kqw\_}$ and $\operatorname{kqt\_}$, 
\begin{equation*}
\begin{array}{@{}l@{}} \operatorname{is\text{-}finite\text{-}rank\text{-}op}(\operatorname{kqv\_}, \operatorname{kqw\_}, \operatorname{kqt\_}) \\ \Leftrightarrow\; \operatorname{kqt\_} \in L({\operatorname{kqv\_}}, {\operatorname{kqw\_}}) \wedge \\ \exists \operatorname{kqi\_} \in \mathbf{Set}. \\ \quad |\operatorname{kqi\_}| \in \mathbb{N} \wedge \\ \quad \exists \operatorname{kqa\_} \in (\operatorname{kqi\_} \to \operatorname{vec}(\operatorname{kqv\_})), \operatorname{kqb\_} \in (\operatorname{kqi\_} \to \operatorname{vec}(\operatorname{kqw\_})). \forall \operatorname{kqx\_} \in \operatorname{vec}(\operatorname{kqv\_}). \\ \quad \quad \operatorname{kqt\_}(\operatorname{kqx\_}) \\ \quad \quad =\; \operatorname{finsum}(\operatorname{normed\text{-}ag\text{-}as\text{-}abelian\text{-}group}(\operatorname{ip\text{-}normed\text{-}ag}(\operatorname{kqw\_})), \\ \quad \quad \quad (\lambda \operatorname{kqk\_} \in \operatorname{kqi\_}.\; \operatorname{act}(\operatorname{kqw\_})(\langle\cdot,\cdot\rangle_{\operatorname{kqv\_}}(\operatorname{kqx\_}, \operatorname{kqa\_}(\operatorname{kqk\_})), \operatorname{kqb\_}(\operatorname{kqk\_}))), \\ \quad \quad \quad \operatorname{kqi\_}) \end{array}
\end{equation*}

\par\medskip\noindent\textbf{Definition} (\texttt{finite-rank-ops}).\;$F({\operatorname{kqv\_}}, {\operatorname{kqw\_}}) :=$
\begin{equation*}
\{\operatorname{kqt\_} \in L({\operatorname{kqv\_}}, {\operatorname{kqw\_}}) : \operatorname{is\text{-}finite\text{-}rank\text{-}op}(\operatorname{kqv\_}, \operatorname{kqw\_}, \operatorname{kqt\_})\}
\end{equation*}

\par\medskip\noindent\textbf{Definition} (\texttt{is-eigenvalue}).\;For all $\operatorname{kqv\_}$, $\operatorname{kqt\_}$ and $\operatorname{kqc\_}$, 
\begin{equation*}
\begin{array}{@{}l@{}} \operatorname{is\text{-}eigenvalue}(\operatorname{kqv\_}, \operatorname{kqt\_}, \operatorname{kqc\_}) \\ \Leftrightarrow\; ((\operatorname{kqc\_} \in \mathbb{C}) \wedge \exists \operatorname{kqx\_} \in N_{\operatorname{kqc\_}}({\operatorname{kqt\_}}).\; \neg (\operatorname{kqx\_} = \operatorname{vzero}(\operatorname{kqv\_}))) \end{array}
\end{equation*}

\par\medskip\noindent\textbf{Definition} (\texttt{trace-class-ops}).\;$T({\operatorname{kqv\_}}) :=$
\begin{equation*}
\{\operatorname{kqt\_} \in L({\operatorname{kqv\_}}, {\operatorname{kqv\_}}) : \operatorname{is\text{-}trace\text{-}class}(\operatorname{kqv\_}, \operatorname{kqt\_})\}
\end{equation*}


\vnbbookfile{theorem-library/th3-vocab-unfold}{theorem-library/th3-vocab-unfold}
\noindent{\small th3-vocab-unfold.scm -- wave of 2026-11-09/10 (agent TH3-A), wired 2026-11-10 from scratchpad/th3-1109/th3-vocab-unfold.scm. Library: operator-algebras.}\par
\vnbprop{compact-ops-unfold}{\vnbstmt{\begin{array}{@{}l@{}} \forall v, w. \\ \quad (K({v}, {w}) \simeq \{t \in L({v}, {w}) : \operatorname{is\text{-}compact\text{-}op}(v, w, t)\}) \end{array}}}
\vnbprop{compact-ops-membership}{\vnbstmt{\begin{array}{@{}l@{}} \forall v, w, z. \\ \quad (z \in K({v}, {w})) \Leftrightarrow ((z \in L({v}, {w})) \wedge \operatorname{is\text{-}compact\text{-}op}(v, w, z)) \end{array}}}
\vnbprop{rank-one-op-unfold}{\vnbstmt{\begin{array}{@{}l@{}} \forall v, w, a, b. \\ \quad (\langle \cdot \mid {a} \rangle {b} \simeq (\lambda x \in \operatorname{vec}(v).\; \operatorname{act}(w)(\langle\cdot,\cdot\rangle_{v}(x, a), b))) \end{array}}}
\vnbprop{finite-rank-ops-unfold}{\vnbstmt{\begin{array}{@{}l@{}} \forall v, w. \\ \quad (F({v}, {w}) \simeq \{t \in L({v}, {w}) : \operatorname{is\text{-}finite\text{-}rank\text{-}op}(v, w, t)\}) \end{array}}}
\vnbprop{finite-rank-ops-membership}{\vnbstmt{\begin{array}{@{}l@{}} \forall v, w, z. \\ \quad (z \in F({v}, {w})) \Leftrightarrow ((z \in L({v}, {w})) \wedge \operatorname{is\text{-}finite\text{-}rank\text{-}op}(v, w, z)) \end{array}}}
\vnbprop{eigenspace-unfold}{\vnbstmt{\forall v, t, c.\; (N_{c}({t}) \simeq \{x \in \operatorname{vec}(v) : (t(x) = \operatorname{act}(v)(c, x))\})}}
\vnbprop{eigenspace-membership}{\vnbstmt{\begin{array}{@{}l@{}} \forall v, t, c, z. \\ \quad (z \in N_{c}({t})) \Leftrightarrow ((z \in \operatorname{vec}(v)) \wedge (t(z) = \operatorname{act}(v)(c, z))) \end{array}}}
\vnbprop{op-sqrt-unfold}{\vnbstmt{\begin{array}{@{}l@{}} \forall v, t. \\ \quad ({t}^{1/2} \simeq (\iota\, s.\;(\operatorname{is\text{-}positive\text{-}op}(v, s) \wedge ({s}{s} = t)))) \end{array}}}
\vnbprop{abs-op-unfold}{\vnbstmt{\forall v, t.\; (\lvert {t} \rvert \simeq {{{t}^{*}}{t}}^{1/2})}}
\vnbprop{onb-sets-unfold}{\vnbstmt{\begin{array}{@{}l@{}} \forall v. \\ \quad (\operatorname{onb\text{-}sets}(v) \simeq \{o \in \operatorname{power}(\operatorname{vec}(v)) : \operatorname{is\text{-}orthonormal\text{-}basis}(v, o, (\lambda e \in o.\; e))\}) \end{array}}}
\vnbprop{onb-sets-membership}{\vnbstmt{\begin{array}{@{}l@{}} \forall v, z. \\ \quad (z \in \operatorname{onb\text{-}sets}(v)) \Leftrightarrow ((z \in \operatorname{power}(\operatorname{vec}(v))) \wedge \operatorname{is\text{-}orthonormal\text{-}basis}(v, z, (\lambda e \in z.\; e))) \end{array}}}
\vnbprop{canon-onb-unfold}{\vnbstmt{\forall v.\; (\operatorname{canon\text{-}onb}(v) \simeq \varepsilon(\operatorname{onb\text{-}sets}(v)))}}
\vnbprop{hs-sum-unfold}{\vnbstmt{\begin{array}{@{}l@{}} \forall v, w, i, e, t. \\ \quad (\sum_{i \in {i}} \lVert {t} {e}_i \rVert^2 \simeq \operatorname{esum}((\lambda k \in i.\; (\lVert t(e(k)) \rVert \cdot \lVert t(e(k)) \rVert)))) \end{array}}}
\vnbprop{hs-ops-unfold}{\vnbstmt{\begin{array}{@{}l@{}} \forall v, w. \\ \quad (HS({v}, {w}) \simeq \{t \in L({v}, {w}) : \operatorname{is\text{-}hilbert\text{-}schmidt}(v, w, t)\}) \end{array}}}
\vnbprop{hs-ops-membership}{\vnbstmt{\begin{array}{@{}l@{}} \forall v, w, z. \\ \quad (z \in HS({v}, {w})) \Leftrightarrow ((z \in L({v}, {w})) \wedge \operatorname{is\text{-}hilbert\text{-}schmidt}(v, w, z)) \end{array}}}
\vnbprop{hs-norm-unfold}{\vnbstmt{\begin{array}{@{}l@{}} \forall v, w, t. \\ \quad (\lVert {t} \rVert_{HS} \simeq \sqrt{\sum_{i \in {\operatorname{canon\text{-}onb}(v)}} \lVert {t} {(\lambda e \in \operatorname{canon\text{-}onb}(v).\; e)}_i \rVert^2}) \end{array}}}
\vnbprop{hs-ip-unfold}{\vnbstmt{\begin{array}{@{}l@{}} \forall v, w, u, t. \\ \quad (\langle {u} \mid {t} \rangle_{HS} \simeq \operatorname{cc\text{-}family\text{-}sum}((\lambda e \in \operatorname{canon\text{-}onb}(v).\; \langle\cdot,\cdot\rangle_{w}(u(e), t(e))))) \end{array}}}
\vnbprop{trace-class-ops-unfold}{\vnbstmt{\begin{array}{@{}l@{}} \forall v. \\ \quad (T({v}) \simeq \{t \in L({v}, {v}) : \operatorname{is\text{-}trace\text{-}class}(v, t)\}) \end{array}}}
\vnbprop{trace-class-ops-membership}{\vnbstmt{\begin{array}{@{}l@{}} \forall v, z. \\ \quad (z \in T({v})) \Leftrightarrow ((z \in L({v}, {v})) \wedge \operatorname{is\text{-}trace\text{-}class}(v, z)) \end{array}}}
\vnbprop{op-trace-unfold}{\vnbstmt{\begin{array}{@{}l@{}} \forall v, t. \\ \quad (\mathrm{tr}\,{t} \simeq \operatorname{cc\text{-}family\text{-}sum}((\lambda e \in \operatorname{canon\text{-}onb}(v).\; \langle\cdot,\cdot\rangle_{v}(t(e), e)))) \end{array}}}

\vnbbookfile{structure-library/ba-vocab}{structure-library/ba-vocab}
\noindent{\small ba-vocab.scm -- wave of 2026-11-09/10 (agent BA-1), wired 2026-11-10 from scratchpad/ba1-1109/ba-vocab.scm. Library: operator-algebras.}\par
\vnbstructdef{complex-normed-algebra}{A complex normed algebra $s$ is an 8-tuple whose components are listed below.\par
\begin{enumerate}[nosep]
\item $\operatorname{scal}(s) = \operatorname{nth}(1, s)$, a substructure (\texttt{substructure}): $\operatorname{is\text{-}ring}(\operatorname{scal}(s))$.
\item $\operatorname{vec}(s) = \operatorname{nth}(2, s)$, a carrier (\texttt{carriers}): $\operatorname{vec}(s) \in \mathbf{Set}$.
\item $\operatorname{vadd}(s) = \operatorname{nth}(3, s)$, an operation (\texttt{op}): $\operatorname{vadd}(s) \in ((\operatorname{vec}(s) \times \operatorname{vec}(s)) \to \operatorname{vec}(s))$.
\item $\operatorname{vzero}(s) = \operatorname{nth}(4, s)$, a constant (\texttt{constant}): $\operatorname{vzero}(s) \in \operatorname{vec}(s)$.
\item $\operatorname{vneg}(s) = \operatorname{nth}(5, s)$, an operation (\texttt{op}): $\operatorname{vneg}(s) \in (\operatorname{vec}(s) \to \operatorname{vec}(s))$.
\item $\operatorname{act}(s) = \operatorname{nth}(6, s)$, an operation (\texttt{op}): $\operatorname{act}(s) \in ((\operatorname{carr}(\operatorname{scal}(s)) \times \operatorname{vec}(s)) \to \operatorname{vec}(s))$.
\item $\operatorname{vnrm}(s) = \operatorname{nth}(7, s)$, an operation (\texttt{op}): $\operatorname{vnrm}(s) \in (\operatorname{vec}(s) \to \mathbb{R})$.
\item $\cdot_{s} = \operatorname{nth}(8, s)$, an operation (\texttt{op}): $\cdot_{s} \in ((\operatorname{vec}(s) \times \operatorname{vec}(s)) \to \operatorname{vec}(s))$.
\end{enumerate}
The laws (the \texttt{property} and \texttt{law} clauses):
\begin{itemize}[nosep]
\item $\operatorname{is\text{-}associative}(\operatorname{vadd}(s), \operatorname{vec}(s))$
\item $\operatorname{is\text{-}commutative}(\operatorname{vadd}(s), \operatorname{vec}(s))$
\item $\operatorname{is\text{-}identity}(\operatorname{vadd}(s), \operatorname{vzero}(s), \operatorname{vec}(s))$
\item $\operatorname{has\text{-}inverses}(\operatorname{vadd}(s), \operatorname{vzero}(s), \operatorname{vneg}(s), \operatorname{vec}(s))$
\item $\operatorname{is\text{-}associative}(\cdot_{s}, \operatorname{vec}(s))$
\item $\operatorname{scal}(s) = \operatorname{normed\text{-}field\text{-}as\text{-}commutative\text{-}ring}(\operatorname{cc\text{-}normed\text{-}field})$
\item $\operatorname{carr}(\operatorname{scal}(s)) = \mathbb{C}$
\item
\begin{equation*}
\begin{array}{@{}l@{}} \forall \operatorname{r\_} \in \operatorname{carr}(\operatorname{scal}(s)), \operatorname{x\_}, \operatorname{y\_} \in \operatorname{vec}(s). \\ \quad \operatorname{act}(s)(\operatorname{r\_}, \operatorname{vadd}(s)(\operatorname{x\_}, \operatorname{y\_})) = \operatorname{vadd}(s)(\operatorname{act}(s)(\operatorname{r\_}, \operatorname{x\_}), \operatorname{act}(s)(\operatorname{r\_}, \operatorname{y\_})) \end{array}
\end{equation*}

\item
\begin{equation*}
\begin{array}{@{}l@{}} \forall \operatorname{r\_}, \operatorname{s\_} \in \operatorname{carr}(\operatorname{scal}(s)), \operatorname{x\_} \in \operatorname{vec}(s). \\ \quad \operatorname{act}(s)(\operatorname{add}(\operatorname{scal}(s))(\operatorname{r\_}, \operatorname{s\_}), \operatorname{x\_}) = \operatorname{vadd}(s)(\operatorname{act}(s)(\operatorname{r\_}, \operatorname{x\_}), \operatorname{act}(s)(\operatorname{s\_}, \operatorname{x\_})) \end{array}
\end{equation*}

\item
\begin{equation*}
\begin{array}{@{}l@{}} \forall \operatorname{r\_}, \operatorname{s\_} \in \operatorname{carr}(\operatorname{scal}(s)), \operatorname{x\_} \in \operatorname{vec}(s). \\ \quad \operatorname{act}(s)(\operatorname{mul}(\operatorname{scal}(s))(\operatorname{r\_}, \operatorname{s\_}), \operatorname{x\_}) = \operatorname{act}(s)(\operatorname{r\_}, \operatorname{act}(s)(\operatorname{s\_}, \operatorname{x\_})) \end{array}
\end{equation*}

\item $\forall \operatorname{x\_} \in \operatorname{vec}(s).\; (\operatorname{act}(s)(\operatorname{one}(\operatorname{scal}(s)), \operatorname{x\_}) = \operatorname{x\_})$
\item $\forall \operatorname{x\_} \in \operatorname{vec}(s).\; (0 \leq \operatorname{vnrm}(s)(\operatorname{x\_}))$
\item $\forall \operatorname{x\_} \in \operatorname{vec}(s).\; ((\operatorname{vnrm}(s)(\operatorname{x\_}) = 0) \Leftrightarrow (\operatorname{x\_} = \operatorname{vzero}(s)))$
\item
\begin{equation*}
\begin{array}{@{}l@{}} \forall \operatorname{r\_} \in \operatorname{carr}(\operatorname{scal}(s)), \operatorname{x\_} \in \operatorname{vec}(s). \\ \quad \operatorname{vnrm}(s)(\operatorname{act}(s)(\operatorname{r\_}, \operatorname{x\_})) = (\operatorname{magnitude}(\operatorname{r\_}) \cdot \operatorname{vnrm}(s)(\operatorname{x\_})) \end{array}
\end{equation*}

\item
\begin{equation*}
\begin{array}{@{}l@{}} \forall \operatorname{x\_}, \operatorname{y\_} \in \operatorname{vec}(s). \\ \quad \operatorname{vnrm}(s)(\operatorname{vadd}(s)(\operatorname{x\_}, \operatorname{y\_})) \leq (\operatorname{vnrm}(s)(\operatorname{x\_}) + \operatorname{vnrm}(s)(\operatorname{y\_})) \end{array}
\end{equation*}

\item
\begin{equation*}
\begin{array}{@{}l@{}} \forall \operatorname{x\_}, \operatorname{y\_}, \operatorname{z\_} \in \operatorname{vec}(s). \\ \quad \cdot_{s}(\operatorname{x\_}, \operatorname{vadd}(s)(\operatorname{y\_}, \operatorname{z\_})) = \operatorname{vadd}(s)(\cdot_{s}(\operatorname{x\_}, \operatorname{y\_}), \cdot_{s}(\operatorname{x\_}, \operatorname{z\_})) \end{array}
\end{equation*}

\item
\begin{equation*}
\begin{array}{@{}l@{}} \forall \operatorname{x\_}, \operatorname{y\_}, \operatorname{z\_} \in \operatorname{vec}(s). \\ \quad \cdot_{s}(\operatorname{vadd}(s)(\operatorname{x\_}, \operatorname{y\_}), \operatorname{z\_}) = \operatorname{vadd}(s)(\cdot_{s}(\operatorname{x\_}, \operatorname{z\_}), \cdot_{s}(\operatorname{y\_}, \operatorname{z\_})) \end{array}
\end{equation*}

\item
\begin{equation*}
\begin{array}{@{}l@{}} \forall \operatorname{r\_} \in \operatorname{carr}(\operatorname{scal}(s)), \operatorname{x\_}, \operatorname{y\_} \in \operatorname{vec}(s). \\ \quad \cdot_{s}(\operatorname{act}(s)(\operatorname{r\_}, \operatorname{x\_}), \operatorname{y\_}) = \operatorname{act}(s)(\operatorname{r\_}, \cdot_{s}(\operatorname{x\_}, \operatorname{y\_})) \end{array}
\end{equation*}

\item
\begin{equation*}
\begin{array}{@{}l@{}} \forall \operatorname{r\_} \in \operatorname{carr}(\operatorname{scal}(s)), \operatorname{x\_}, \operatorname{y\_} \in \operatorname{vec}(s). \\ \quad \cdot_{s}(\operatorname{x\_}, \operatorname{act}(s)(\operatorname{r\_}, \operatorname{y\_})) = \operatorname{act}(s)(\operatorname{r\_}, \cdot_{s}(\operatorname{x\_}, \operatorname{y\_})) \end{array}
\end{equation*}

\item
\begin{equation*}
\begin{array}{@{}l@{}} \forall \operatorname{x\_}, \operatorname{y\_} \in \operatorname{vec}(s). \\ \quad \operatorname{vnrm}(s)(\cdot_{s}(\operatorname{x\_}, \operatorname{y\_})) \leq (\operatorname{vnrm}(s)(\operatorname{x\_}) \cdot \operatorname{vnrm}(s)(\operatorname{y\_})) \end{array}
\end{equation*}

\end{itemize}
The characteristic law (the axiom \texttt{is-complex-normed-algebra}): $\operatorname{is\text{-}complex\text{-}normed\text{-}algebra}(s)$ holds if and only if $\operatorname{length}(s) = 8$, every component has the typing listed with it, every law holds, and $s \in \mathbf{Set}$ (the structure is a set).
}
\vnbstructdef{unital-banach-algebra}{An unital banach algebra $s$ is a 9-tuple whose components are listed below.\par
\begin{enumerate}[nosep]
\item $\operatorname{scal}(s) = \operatorname{nth}(1, s)$, a substructure (\texttt{substructure}): $\operatorname{is\text{-}ring}(\operatorname{scal}(s))$.
\item $\operatorname{vec}(s) = \operatorname{nth}(2, s)$, a carrier (\texttt{carriers}): $\operatorname{vec}(s) \in \mathbf{Set}$.
\item $\operatorname{vadd}(s) = \operatorname{nth}(3, s)$, an operation (\texttt{op}): $\operatorname{vadd}(s) \in ((\operatorname{vec}(s) \times \operatorname{vec}(s)) \to \operatorname{vec}(s))$.
\item $\operatorname{vzero}(s) = \operatorname{nth}(4, s)$, a constant (\texttt{constant}): $\operatorname{vzero}(s) \in \operatorname{vec}(s)$.
\item $\operatorname{vneg}(s) = \operatorname{nth}(5, s)$, an operation (\texttt{op}): $\operatorname{vneg}(s) \in (\operatorname{vec}(s) \to \operatorname{vec}(s))$.
\item $\operatorname{act}(s) = \operatorname{nth}(6, s)$, an operation (\texttt{op}): $\operatorname{act}(s) \in ((\operatorname{carr}(\operatorname{scal}(s)) \times \operatorname{vec}(s)) \to \operatorname{vec}(s))$.
\item $\operatorname{vnrm}(s) = \operatorname{nth}(7, s)$, an operation (\texttt{op}): $\operatorname{vnrm}(s) \in (\operatorname{vec}(s) \to \mathbb{R})$.
\item $\cdot_{s} = \operatorname{nth}(8, s)$, an operation (\texttt{op}): $\cdot_{s} \in ((\operatorname{vec}(s) \times \operatorname{vec}(s)) \to \operatorname{vec}(s))$.
\item $1_{s} = \operatorname{nth}(9, s)$, a constant (\texttt{constant}): $1_{s} \in \operatorname{vec}(s)$.
\end{enumerate}
The laws (the \texttt{property} and \texttt{law} clauses):
\begin{itemize}[nosep]
\item
\begin{equation*}
\operatorname{is\text{-}complex\text{-}banach\text{-}algebra}({[}\operatorname{scal}(s), \operatorname{vec}(s), \operatorname{vadd}(s), \operatorname{vzero}(s), \operatorname{vneg}(s), \operatorname{act}(s), \operatorname{vnrm}(s), \cdot_{s}])
\end{equation*}

\item $\operatorname{carr}(\operatorname{scal}(s)) = \mathbb{C}$
\item $\forall \operatorname{x\_} \in \operatorname{vec}(s).\; (\cdot_{s}(1_{s}, \operatorname{x\_}) = \operatorname{x\_})$
\item $\forall \operatorname{x\_} \in \operatorname{vec}(s).\; (\cdot_{s}(\operatorname{x\_}, 1_{s}) = \operatorname{x\_})$
\item $\operatorname{vnrm}(s)(1_{s}) = 1$
\end{itemize}
The characteristic law (the axiom \texttt{is-unital-banach-algebra}): $\operatorname{is\text{-}unital\text{-}banach\text{-}algebra}(s)$ holds if and only if $\operatorname{length}(s) = 9$, every component has the typing listed with it, every law holds, and $s \in \mathbf{Set}$ (the structure is a set).
}
\vnbstructdef{unital-involutive-banach-algebra}{An unital involutive banach algebra $s$ is a 10-tuple whose components are listed below.\par
\begin{enumerate}[nosep]
\item $\operatorname{scal}(s) = \operatorname{nth}(1, s)$, a substructure (\texttt{substructure}): $\operatorname{is\text{-}ring}(\operatorname{scal}(s))$.
\item $\operatorname{vec}(s) = \operatorname{nth}(2, s)$, a carrier (\texttt{carriers}): $\operatorname{vec}(s) \in \mathbf{Set}$.
\item $\operatorname{vadd}(s) = \operatorname{nth}(3, s)$, an operation (\texttt{op}): $\operatorname{vadd}(s) \in ((\operatorname{vec}(s) \times \operatorname{vec}(s)) \to \operatorname{vec}(s))$.
\item $\operatorname{vzero}(s) = \operatorname{nth}(4, s)$, a constant (\texttt{constant}): $\operatorname{vzero}(s) \in \operatorname{vec}(s)$.
\item $\operatorname{vneg}(s) = \operatorname{nth}(5, s)$, an operation (\texttt{op}): $\operatorname{vneg}(s) \in (\operatorname{vec}(s) \to \operatorname{vec}(s))$.
\item $\operatorname{act}(s) = \operatorname{nth}(6, s)$, an operation (\texttt{op}): $\operatorname{act}(s) \in ((\operatorname{carr}(\operatorname{scal}(s)) \times \operatorname{vec}(s)) \to \operatorname{vec}(s))$.
\item $\operatorname{vnrm}(s) = \operatorname{nth}(7, s)$, an operation (\texttt{op}): $\operatorname{vnrm}(s) \in (\operatorname{vec}(s) \to \mathbb{R})$.
\item $\cdot_{s} = \operatorname{nth}(8, s)$, an operation (\texttt{op}): $\cdot_{s} \in ((\operatorname{vec}(s) \times \operatorname{vec}(s)) \to \operatorname{vec}(s))$.
\item $1_{s} = \operatorname{nth}(9, s)$, a constant (\texttt{constant}): $1_{s} \in \operatorname{vec}(s)$.
\item ${}^{*_{s}} = \operatorname{nth}(10, s)$, an operation (\texttt{op}): ${}^{*_{s}} \in (\operatorname{vec}(s) \to \operatorname{vec}(s))$.
\end{enumerate}
The laws (the \texttt{property} and \texttt{law} clauses):
\begin{itemize}[nosep]
\item
\begin{equation*}
\operatorname{is\text{-}unital\text{-}banach\text{-}algebra}({[}\operatorname{scal}(s), \operatorname{vec}(s), \operatorname{vadd}(s), \operatorname{vzero}(s), \operatorname{vneg}(s), \operatorname{act}(s), \operatorname{vnrm}(s), \cdot_{s}, 1_{s}])
\end{equation*}

\item $\operatorname{carr}(\operatorname{scal}(s)) = \mathbb{C}$
\item
\begin{equation*}
\begin{array}{@{}l@{}} \forall \operatorname{x\_}, \operatorname{y\_} \in \operatorname{vec}(s). \\ \quad {}^{*_{s}}(\operatorname{vadd}(s)(\operatorname{x\_}, \operatorname{y\_})) = \operatorname{vadd}(s)({}^{*_{s}}(\operatorname{x\_}), {}^{*_{s}}(\operatorname{y\_})) \end{array}
\end{equation*}

\item
\begin{equation*}
\begin{array}{@{}l@{}} \forall \operatorname{r\_} \in \mathbb{C}, \operatorname{x\_} \in \operatorname{vec}(s). \\ \quad {}^{*_{s}}(\operatorname{act}(s)(\operatorname{r\_}, \operatorname{x\_})) = \operatorname{act}(s)(\operatorname{conjugate}(\operatorname{r\_}), {}^{*_{s}}(\operatorname{x\_})) \end{array}
\end{equation*}

\item $\forall \operatorname{x\_} \in \operatorname{vec}(s).\; ({}^{*_{s}}({}^{*_{s}}(\operatorname{x\_})) = \operatorname{x\_})$
\item
\begin{equation*}
\begin{array}{@{}l@{}} \forall \operatorname{x\_}, \operatorname{y\_} \in \operatorname{vec}(s). \\ \quad {}^{*_{s}}(\cdot_{s}(\operatorname{x\_}, \operatorname{y\_})) = \cdot_{s}({}^{*_{s}}(\operatorname{y\_}), {}^{*_{s}}(\operatorname{x\_})) \end{array}
\end{equation*}

\item $\forall \operatorname{x\_} \in \operatorname{vec}(s).\; (\operatorname{vnrm}(s)({}^{*_{s}}(\operatorname{x\_})) = \operatorname{vnrm}(s)(\operatorname{x\_}))$
\end{itemize}
The characteristic law (the axiom \texttt{is-unital-involutive-banach-algebra}): $\operatorname{is\text{-}unital\text{-}involutive\text{-}banach\text{-}algebra}(s)$ holds if and only if $\operatorname{length}(s) = 10$, every component has the typing listed with it, every law holds, and $s \in \mathbf{Set}$ (the structure is a set).
}
\vnbstructdef{nu-involutive-banach-algebra}{A nu involutive banach algebra $s$ is a 9-tuple whose components are listed below.\par
\begin{enumerate}[nosep]
\item $\operatorname{scal}(s) = \operatorname{nth}(1, s)$, a substructure (\texttt{substructure}): $\operatorname{is\text{-}ring}(\operatorname{scal}(s))$.
\item $\operatorname{vec}(s) = \operatorname{nth}(2, s)$, a carrier (\texttt{carriers}): $\operatorname{vec}(s) \in \mathbf{Set}$.
\item $\operatorname{vadd}(s) = \operatorname{nth}(3, s)$, an operation (\texttt{op}): $\operatorname{vadd}(s) \in ((\operatorname{vec}(s) \times \operatorname{vec}(s)) \to \operatorname{vec}(s))$.
\item $\operatorname{vzero}(s) = \operatorname{nth}(4, s)$, a constant (\texttt{constant}): $\operatorname{vzero}(s) \in \operatorname{vec}(s)$.
\item $\operatorname{vneg}(s) = \operatorname{nth}(5, s)$, an operation (\texttt{op}): $\operatorname{vneg}(s) \in (\operatorname{vec}(s) \to \operatorname{vec}(s))$.
\item $\operatorname{act}(s) = \operatorname{nth}(6, s)$, an operation (\texttt{op}): $\operatorname{act}(s) \in ((\operatorname{carr}(\operatorname{scal}(s)) \times \operatorname{vec}(s)) \to \operatorname{vec}(s))$.
\item $\operatorname{vnrm}(s) = \operatorname{nth}(7, s)$, an operation (\texttt{op}): $\operatorname{vnrm}(s) \in (\operatorname{vec}(s) \to \mathbb{R})$.
\item $\cdot_{s} = \operatorname{nth}(8, s)$, an operation (\texttt{op}): $\cdot_{s} \in ((\operatorname{vec}(s) \times \operatorname{vec}(s)) \to \operatorname{vec}(s))$.
\item ${}^{*_{s}} = \operatorname{nth}(9, s)$, an operation (\texttt{op}): ${}^{*_{s}} \in (\operatorname{vec}(s) \to \operatorname{vec}(s))$.
\end{enumerate}
The laws (the \texttt{property} and \texttt{law} clauses):
\begin{itemize}[nosep]
\item
\begin{equation*}
\operatorname{is\text{-}complex\text{-}banach\text{-}algebra}({[}\operatorname{scal}(s), \operatorname{vec}(s), \operatorname{vadd}(s), \operatorname{vzero}(s), \operatorname{vneg}(s), \operatorname{act}(s), \operatorname{vnrm}(s), \cdot_{s}])
\end{equation*}

\item $\operatorname{carr}(\operatorname{scal}(s)) = \mathbb{C}$
\item
\begin{equation*}
\begin{array}{@{}l@{}} \forall \operatorname{x\_}, \operatorname{y\_} \in \operatorname{vec}(s). \\ \quad {}^{*_{s}}(\operatorname{vadd}(s)(\operatorname{x\_}, \operatorname{y\_})) = \operatorname{vadd}(s)({}^{*_{s}}(\operatorname{x\_}), {}^{*_{s}}(\operatorname{y\_})) \end{array}
\end{equation*}

\item
\begin{equation*}
\begin{array}{@{}l@{}} \forall \operatorname{r\_} \in \mathbb{C}, \operatorname{x\_} \in \operatorname{vec}(s). \\ \quad {}^{*_{s}}(\operatorname{act}(s)(\operatorname{r\_}, \operatorname{x\_})) = \operatorname{act}(s)(\operatorname{conjugate}(\operatorname{r\_}), {}^{*_{s}}(\operatorname{x\_})) \end{array}
\end{equation*}

\item $\forall \operatorname{x\_} \in \operatorname{vec}(s).\; ({}^{*_{s}}({}^{*_{s}}(\operatorname{x\_})) = \operatorname{x\_})$
\item
\begin{equation*}
\begin{array}{@{}l@{}} \forall \operatorname{x\_}, \operatorname{y\_} \in \operatorname{vec}(s). \\ \quad {}^{*_{s}}(\cdot_{s}(\operatorname{x\_}, \operatorname{y\_})) = \cdot_{s}({}^{*_{s}}(\operatorname{y\_}), {}^{*_{s}}(\operatorname{x\_})) \end{array}
\end{equation*}

\item $\forall \operatorname{x\_} \in \operatorname{vec}(s).\; (\operatorname{vnrm}(s)({}^{*_{s}}(\operatorname{x\_})) = \operatorname{vnrm}(s)(\operatorname{x\_}))$
\end{itemize}
The characteristic law (the axiom \texttt{is-nu-involutive-banach-algebra}): $\operatorname{is\text{-}nu\text{-}involutive\text{-}banach\text{-}algebra}(s)$ holds if and only if $\operatorname{length}(s) = 9$, every component has the typing listed with it, every law holds, and $s \in \mathbf{Set}$ (the structure is a set).
}
\par\medskip\noindent\textbf{Definition} (\texttt{is-alg-hom}).\;For all $\operatorname{baa\_}$, $\operatorname{bab\_}$ and $\operatorname{baf\_}$, 
\begin{equation*}
\begin{array}{@{}l@{}} \operatorname{is\text{-}alg\text{-}hom}(\operatorname{baa\_}, \operatorname{bab\_}, \operatorname{baf\_}) \\ \Leftrightarrow\; \operatorname{baf\_} \in (\operatorname{vec}(\operatorname{baa\_}) \to \operatorname{vec}(\operatorname{bab\_})) \wedge \\ \forall \operatorname{bax\_}, \operatorname{bay\_} \in \operatorname{vec}(\operatorname{baa\_}). \\ \quad \operatorname{baf\_}(\operatorname{vadd}(\operatorname{baa\_})(\operatorname{bax\_}, \operatorname{bay\_})) = \operatorname{vadd}(\operatorname{bab\_})(\operatorname{baf\_}(\operatorname{bax\_}), \operatorname{baf\_}(\operatorname{bay\_})) \wedge \\ \forall \operatorname{baz\_} \in \mathbb{C}, \operatorname{bax\_} \in \operatorname{vec}(\operatorname{baa\_}). \\ \quad \operatorname{baf\_}(\operatorname{act}(\operatorname{baa\_})(\operatorname{baz\_}, \operatorname{bax\_})) = \operatorname{act}(\operatorname{bab\_})(\operatorname{baz\_}, \operatorname{baf\_}(\operatorname{bax\_})) \wedge \\ \forall \operatorname{bax\_}, \operatorname{bay\_} \in \operatorname{vec}(\operatorname{baa\_}). \\ \quad \operatorname{baf\_}(\cdot_{\operatorname{baa\_}}(\operatorname{bax\_}, \operatorname{bay\_})) = \cdot_{\operatorname{bab\_}}(\operatorname{baf\_}(\operatorname{bax\_}), \operatorname{baf\_}(\operatorname{bay\_})) \wedge \\ \operatorname{baf\_}(1_{\operatorname{baa\_}}) = 1_{\operatorname{bab\_}} \end{array}
\end{equation*}

\par\medskip\noindent\textbf{Definition} (\texttt{is-star-hom}).\;For all $\operatorname{baa\_}$, $\operatorname{bab\_}$ and $\operatorname{baf\_}$, 
\begin{equation*}
\begin{array}{@{}l@{}} \operatorname{is\text{-}star\text{-}hom}(\operatorname{baa\_}, \operatorname{bab\_}, \operatorname{baf\_}) \\ \Leftrightarrow\; \operatorname{is\text{-}alg\text{-}hom}(\operatorname{baa\_}, \operatorname{bab\_}, \operatorname{baf\_}) \wedge \\ \forall \operatorname{bax\_} \in \operatorname{vec}(\operatorname{baa\_}). \\ \quad \operatorname{baf\_}({}^{*_{\operatorname{baa\_}}}(\operatorname{bax\_})) = {}^{*_{\operatorname{bab\_}}}(\operatorname{baf\_}(\operatorname{bax\_})) \end{array}
\end{equation*}

\par\medskip\noindent\textbf{Definition} (\texttt{cna-metric-space}).\;$\operatorname{cna\text{-}metric\text{-}space}(\operatorname{baa\_}) :=$ $\operatorname{nag\text{-}metric\text{-}space}(\operatorname{cna\text{-}as\text{-}normed\text{-}ag}(\operatorname{baa\_}))$
\par\medskip\noindent\textbf{Definition} (\texttt{ba-metric-space}).\;$\operatorname{ba\text{-}metric\text{-}space}(\operatorname{baa\_}) :=$ $\operatorname{nag\text{-}metric\text{-}space}(\operatorname{uba\text{-}as\text{-}normed\text{-}ag}(\operatorname{baa\_}))$
\par\medskip\noindent\textbf{Definition} (\texttt{inverse-el}).\;${\operatorname{bax\_}}^{-1} :=$
\begin{equation*}
(\iota\, \operatorname{bau\_}.\;((\operatorname{bau\_} \in \operatorname{vec}(\operatorname{baa\_})) \wedge ((\cdot_{\operatorname{baa\_}}(\operatorname{bau\_}, \operatorname{bax\_}) = 1_{\operatorname{baa\_}}) \wedge (\cdot_{\operatorname{baa\_}}(\operatorname{bax\_}, \operatorname{bau\_}) = 1_{\operatorname{baa\_}}))))
\end{equation*}

\par\medskip\noindent\textbf{Definition} (\texttt{is-invertible-el}).\;For all $\operatorname{baa\_}$ and $\operatorname{bax\_}$, 
\begin{equation*}
\begin{array}{@{}l@{}} \operatorname{is\text{-}invertible\text{-}el}(\operatorname{baa\_}, \operatorname{bax\_}) \\ \Leftrightarrow\; \operatorname{bax\_} \in \operatorname{vec}(\operatorname{baa\_}) \wedge \\ \exists \operatorname{bau\_} \in \operatorname{vec}(\operatorname{baa\_}). \\ \quad (\cdot_{\operatorname{baa\_}}(\operatorname{bau\_}, \operatorname{bax\_}) = 1_{\operatorname{baa\_}}) \wedge (\cdot_{\operatorname{baa\_}}(\operatorname{bax\_}, \operatorname{bau\_}) = 1_{\operatorname{baa\_}}) \end{array}
\end{equation*}

\par\medskip\noindent\textbf{Definition} (\texttt{invertibles}).\;$\operatorname{invertibles}(\operatorname{baa\_}) :=$ $\{\operatorname{bav\_} \in \operatorname{vec}(\operatorname{baa\_}) : \operatorname{is\text{-}invertible\text{-}el}(\operatorname{baa\_}, \operatorname{bav\_})\}$
\par\medskip\noindent\textbf{Definition} (\texttt{el-pow}).\;For all $\operatorname{bpa\_}$ and $\operatorname{bpx\_}$,  $({\operatorname{bpx\_}}^{0} \simeq 1_{\operatorname{bpa\_}})$\par
Let $\operatorname{bpa\_}$ and $\operatorname{bpx\_}$ be arbitrary. Let $\operatorname{bpn\_} \in \mathbb{N}$. Then
\begin{equation*}
({\operatorname{bpx\_}}^{\operatorname{succ}(\operatorname{bpn\_})} \simeq \cdot_{\operatorname{bpa\_}}(\operatorname{bpx\_}, {\operatorname{bpx\_}}^{\operatorname{bpn\_}}))
\end{equation*}

\par\medskip\noindent\textbf{Definition} (\texttt{spectrum}).\;$\sigma({\operatorname{bax\_}}) :=$
\begin{equation*}
\{\operatorname{baz\_} \in \mathbb{C} : \neg \operatorname{is\text{-}invertible\text{-}el}(\operatorname{baa\_}, \operatorname{shift\text{-}el}(\operatorname{baa\_}, \operatorname{baz\_}, \operatorname{bax\_}))\}
\end{equation*}

\par\medskip\noindent\textbf{Definition} (\texttt{shift-el}).\;$\operatorname{shift\text{-}el}(\operatorname{baa\_}, \operatorname{baz\_}, \operatorname{bax\_}) :=$
\begin{equation*}
\operatorname{vadd}(\operatorname{baa\_})(\operatorname{act}(\operatorname{baa\_})(\operatorname{baz\_}, 1_{\operatorname{baa\_}}), \operatorname{vneg}(\operatorname{baa\_})(\operatorname{bax\_}))
\end{equation*}

\par\medskip\noindent\textbf{Definition} (\texttt{resolvent}).\;$\operatorname{resolvent}(\operatorname{baa\_}, \operatorname{bax\_}) :=$
\begin{equation*}
(\lambda \operatorname{brz\_} \in \operatorname{difference}(\mathbb{C}, \sigma({\operatorname{bax\_}})).\; {\operatorname{shift\text{-}el}(\operatorname{baa\_}, \operatorname{brz\_}, \operatorname{bax\_})}^{-1})
\end{equation*}

\par\medskip\noindent\textbf{Definition} (\texttt{spectral-radius}).\;$\rho({\operatorname{bax\_}}) :=$
\begin{equation*}
\operatorname{einf}(\{\operatorname{bsw\_} \in \mathbb{R} : \exists \operatorname{bsn\_} \in \mathbb{N}.\; ((1 \leq \operatorname{bsn\_}) \wedge (\operatorname{bsw\_} = {\operatorname{vnrm}(\operatorname{baa\_})({\operatorname{bax\_}}^{\operatorname{bsn\_}})}^{{\operatorname{bsn\_}}^{\ast\!-1}}))\})
\end{equation*}

\par\medskip\noindent\textbf{Definition} (\texttt{is-self-adjoint-el}).\;For all $\operatorname{baa\_}$ and $\operatorname{bax\_}$, 
\begin{equation*}
\operatorname{is\text{-}self\text{-}adjoint\text{-}el}(\operatorname{baa\_}, \operatorname{bax\_}) \Leftrightarrow ((\operatorname{bax\_} \in \operatorname{vec}(\operatorname{baa\_})) \wedge ({}^{*_{\operatorname{baa\_}}}(\operatorname{bax\_}) = \operatorname{bax\_}))
\end{equation*}

\par\medskip\noindent\textbf{Definition} (\texttt{is-normal-el}).\;For all $\operatorname{baa\_}$ and $\operatorname{bax\_}$, 
\begin{equation*}
\begin{array}{@{}l@{}} \operatorname{is\text{-}normal\text{-}el}(\operatorname{baa\_}, \operatorname{bax\_}) \\ \Leftrightarrow\; \operatorname{bax\_} \in \operatorname{vec}(\operatorname{baa\_}) \wedge \\ \cdot_{\operatorname{baa\_}}({}^{*_{\operatorname{baa\_}}}(\operatorname{bax\_}), \operatorname{bax\_}) = \cdot_{\operatorname{baa\_}}(\operatorname{bax\_}, {}^{*_{\operatorname{baa\_}}}(\operatorname{bax\_})) \end{array}
\end{equation*}

\par\medskip\noindent\textbf{Definition} (\texttt{is-unitary-el}).\;For all $\operatorname{baa\_}$ and $\operatorname{bax\_}$, 
\begin{equation*}
\begin{array}{@{}l@{}} \operatorname{is\text{-}unitary\text{-}el}(\operatorname{baa\_}, \operatorname{bax\_}) \\ \Leftrightarrow\; \operatorname{bax\_} \in \operatorname{vec}(\operatorname{baa\_}) \wedge \\ \cdot_{\operatorname{baa\_}}({}^{*_{\operatorname{baa\_}}}(\operatorname{bax\_}), \operatorname{bax\_}) = 1_{\operatorname{baa\_}} \wedge \\ \cdot_{\operatorname{baa\_}}(\operatorname{bax\_}, {}^{*_{\operatorname{baa\_}}}(\operatorname{bax\_})) = 1_{\operatorname{baa\_}} \end{array}
\end{equation*}

\par\medskip\noindent\textbf{Definition} (\texttt{is-projection-el}).\;For all $\operatorname{baa\_}$ and $\operatorname{bax\_}$, 
\begin{equation*}
\begin{array}{@{}l@{}} \operatorname{is\text{-}projection\text{-}el}(\operatorname{baa\_}, \operatorname{bax\_}) \\ \Leftrightarrow\; ((\operatorname{bax\_} \in \operatorname{vec}(\operatorname{baa\_})) \wedge (({}^{*_{\operatorname{baa\_}}}(\operatorname{bax\_}) = \operatorname{bax\_}) \wedge (\cdot_{\operatorname{baa\_}}(\operatorname{bax\_}, \operatorname{bax\_}) = \operatorname{bax\_}))) \end{array}
\end{equation*}

\par\medskip\noindent\textbf{Definition} (\texttt{sa-part}).\;$\operatorname{sa\text{-}part}(\operatorname{baa\_}) :=$ $\{\operatorname{bav\_} \in \operatorname{vec}(\operatorname{baa\_}) : ({}^{*_{\operatorname{baa\_}}}(\operatorname{bav\_}) = \operatorname{bav\_})\}$
\par\medskip\noindent\textbf{Definition} (\texttt{sa-squares}).\;$\operatorname{sa\text{-}squares}(\operatorname{baa\_}) :=$
\begin{equation*}
\{\operatorname{bav\_} \in \operatorname{vec}(\operatorname{baa\_}) : \exists \operatorname{baw\_} \in \operatorname{sa\text{-}part}(\operatorname{baa\_}).\; (\operatorname{bav\_} = \cdot_{\operatorname{baa\_}}(\operatorname{baw\_}, \operatorname{baw\_}))\}
\end{equation*}

\par\medskip\noindent\textbf{Definition} (\texttt{is-alg-subspace}).\;For all $\operatorname{baa\_}$ and $\operatorname{bai\_}$, 
\begin{equation*}
\begin{array}{@{}l@{}} \operatorname{is\text{-}alg\text{-}subspace}(\operatorname{baa\_}, \operatorname{bai\_}) \\ \Leftrightarrow\; \operatorname{bai\_} \subseteq \operatorname{vec}(\operatorname{baa\_}) \wedge \\ \operatorname{vzero}(\operatorname{baa\_}) \in \operatorname{bai\_} \wedge \\ \forall \operatorname{bax\_}, \operatorname{bay\_} \in \operatorname{bai\_}.\; (\operatorname{vadd}(\operatorname{baa\_})(\operatorname{bax\_}, \operatorname{bay\_}) \in \operatorname{bai\_}) \wedge \\ \forall \operatorname{baz\_} \in \mathbb{C}, \operatorname{bax\_} \in \operatorname{bai\_}.\; (\operatorname{act}(\operatorname{baa\_})(\operatorname{baz\_}, \operatorname{bax\_}) \in \operatorname{bai\_}) \end{array}
\end{equation*}

\par\medskip\noindent\textbf{Definition} (\texttt{is-left-ideal}).\;For all $\operatorname{baa\_}$ and $\operatorname{bai\_}$, 
\begin{equation*}
\begin{array}{@{}l@{}} \operatorname{is\text{-}left\text{-}ideal}(\operatorname{baa\_}, \operatorname{bai\_}) \\ \Leftrightarrow\; \operatorname{is\text{-}alg\text{-}subspace}(\operatorname{baa\_}, \operatorname{bai\_}) \wedge \\ \forall \operatorname{bay\_} \in \operatorname{vec}(\operatorname{baa\_}), \operatorname{bax\_} \in \operatorname{bai\_}.\; (\cdot_{\operatorname{baa\_}}(\operatorname{bay\_}, \operatorname{bax\_}) \in \operatorname{bai\_}) \end{array}
\end{equation*}

\par\medskip\noindent\textbf{Definition} (\texttt{is-right-ideal}).\;For all $\operatorname{baa\_}$ and $\operatorname{bai\_}$, 
\begin{equation*}
\begin{array}{@{}l@{}} \operatorname{is\text{-}right\text{-}ideal}(\operatorname{baa\_}, \operatorname{bai\_}) \\ \Leftrightarrow\; \operatorname{is\text{-}alg\text{-}subspace}(\operatorname{baa\_}, \operatorname{bai\_}) \wedge \\ \forall \operatorname{bay\_} \in \operatorname{vec}(\operatorname{baa\_}), \operatorname{bax\_} \in \operatorname{bai\_}.\; (\cdot_{\operatorname{baa\_}}(\operatorname{bax\_}, \operatorname{bay\_}) \in \operatorname{bai\_}) \end{array}
\end{equation*}

\par\medskip\noindent\textbf{Definition} (\texttt{is-two-sided-ideal}).\;For all $\operatorname{baa\_}$ and $\operatorname{bai\_}$, 
\begin{equation*}
\operatorname{is\text{-}two\text{-}sided\text{-}ideal}(\operatorname{baa\_}, \operatorname{bai\_}) \Leftrightarrow (\operatorname{is\text{-}left\text{-}ideal}(\operatorname{baa\_}, \operatorname{bai\_}) \wedge \operatorname{is\text{-}right\text{-}ideal}(\operatorname{baa\_}, \operatorname{bai\_}))
\end{equation*}

\par\medskip\noindent\textbf{Definition} (\texttt{is-maximal-left-ideal}).\;For all $\operatorname{baa\_}$ and $\operatorname{bai\_}$, 
\begin{equation*}
\begin{array}{@{}l@{}} \operatorname{is\text{-}maximal\text{-}left\text{-}ideal}(\operatorname{baa\_}, \operatorname{bai\_}) \\ \Leftrightarrow\; \operatorname{is\text{-}left\text{-}ideal}(\operatorname{baa\_}, \operatorname{bai\_}) \wedge \\ \neg (\operatorname{bai\_} = \operatorname{vec}(\operatorname{baa\_})) \wedge \\ \forall \operatorname{baj\_}. \\ \quad \operatorname{is\text{-}left\text{-}ideal}(\operatorname{baa\_}, \operatorname{baj\_}) \Rightarrow  \\ \quad (\operatorname{bai\_} \subseteq \operatorname{baj\_}) \Rightarrow (\neg (\operatorname{baj\_} = \operatorname{vec}(\operatorname{baa\_})) \Rightarrow (\operatorname{baj\_} = \operatorname{bai\_})) \end{array}
\end{equation*}

\par\medskip\noindent\textbf{Definition} (\texttt{is-maximal-right-ideal}).\;For all $\operatorname{baa\_}$ and $\operatorname{bai\_}$, 
\begin{equation*}
\begin{array}{@{}l@{}} \operatorname{is\text{-}maximal\text{-}right\text{-}ideal}(\operatorname{baa\_}, \operatorname{bai\_}) \\ \Leftrightarrow\; \operatorname{is\text{-}right\text{-}ideal}(\operatorname{baa\_}, \operatorname{bai\_}) \wedge \\ \neg (\operatorname{bai\_} = \operatorname{vec}(\operatorname{baa\_})) \wedge \\ \forall \operatorname{baj\_}. \\ \quad \operatorname{is\text{-}right\text{-}ideal}(\operatorname{baa\_}, \operatorname{baj\_}) \Rightarrow  \\ \quad (\operatorname{bai\_} \subseteq \operatorname{baj\_}) \Rightarrow (\neg (\operatorname{baj\_} = \operatorname{vec}(\operatorname{baa\_})) \Rightarrow (\operatorname{baj\_} = \operatorname{bai\_})) \end{array}
\end{equation*}

\par\medskip\noindent\textbf{Definition} (\texttt{is-maximal-ideal}).\;For all $\operatorname{baa\_}$ and $\operatorname{bai\_}$, 
\begin{equation*}
\begin{array}{@{}l@{}} \operatorname{is\text{-}maximal\text{-}ideal}(\operatorname{baa\_}, \operatorname{bai\_}) \\ \Leftrightarrow\; \operatorname{is\text{-}two\text{-}sided\text{-}ideal}(\operatorname{baa\_}, \operatorname{bai\_}) \wedge \\ \neg (\operatorname{bai\_} = \operatorname{vec}(\operatorname{baa\_})) \wedge \\ \forall \operatorname{baj\_}. \\ \quad \operatorname{is\text{-}two\text{-}sided\text{-}ideal}(\operatorname{baa\_}, \operatorname{baj\_}) \Rightarrow  \\ \quad (\operatorname{bai\_} \subseteq \operatorname{baj\_}) \Rightarrow (\neg (\operatorname{baj\_} = \operatorname{vec}(\operatorname{baa\_})) \Rightarrow (\operatorname{baj\_} = \operatorname{bai\_})) \end{array}
\end{equation*}

\par\medskip\noindent\textbf{Definition} (\texttt{unitization}).\;$\operatorname{unitization}(\operatorname{baa\_}) :=$
\begin{equation*}
{[}\operatorname{normed\text{-}field\text{-}as\text{-}commutative\text{-}ring}(\operatorname{cc\text{-}normed\text{-}field}), \operatorname{unit\text{-}vec}(\operatorname{baa\_}), \operatorname{unit\text{-}vadd}(\operatorname{baa\_}), {[}0, \operatorname{vzero}(\operatorname{baa\_})], \operatorname{unit\text{-}vneg}(\operatorname{baa\_}), \operatorname{unit\text{-}act}(\operatorname{baa\_}), \operatorname{unit\text{-}vnrm}(\operatorname{baa\_}), \operatorname{unit\text{-}vmul}(\operatorname{baa\_}), {[}1, \operatorname{vzero}(\operatorname{baa\_})]]
\end{equation*}

\par\medskip\noindent\textbf{Definition} (\texttt{unitization-na}).\;$\operatorname{unitization\text{-}na}(\operatorname{baa\_}) :=$
\begin{equation*}
{[}\operatorname{normed\text{-}field\text{-}as\text{-}commutative\text{-}ring}(\operatorname{cc\text{-}normed\text{-}field}), \operatorname{unit\text{-}vec}(\operatorname{baa\_}), \operatorname{unit\text{-}vadd}(\operatorname{baa\_}), {[}0, \operatorname{vzero}(\operatorname{baa\_})], \operatorname{unit\text{-}vneg}(\operatorname{baa\_}), \operatorname{unit\text{-}act}(\operatorname{baa\_}), \operatorname{unit\text{-}vnrm}(\operatorname{baa\_}), \operatorname{unit\text{-}vmul}(\operatorname{baa\_})]
\end{equation*}

\par\medskip\noindent\textbf{Definition} (\texttt{unitization-star}).\;$\operatorname{unitization\text{-}star}(\operatorname{baa\_}) :=$
\begin{equation*}
{[}\operatorname{normed\text{-}field\text{-}as\text{-}commutative\text{-}ring}(\operatorname{cc\text{-}normed\text{-}field}), \operatorname{unit\text{-}vec}(\operatorname{baa\_}), \operatorname{unit\text{-}vadd}(\operatorname{baa\_}), {[}0, \operatorname{vzero}(\operatorname{baa\_})], \operatorname{unit\text{-}vneg}(\operatorname{baa\_}), \operatorname{unit\text{-}act}(\operatorname{baa\_}), \operatorname{unit\text{-}vnrm}(\operatorname{baa\_}), \operatorname{unit\text{-}vmul}(\operatorname{baa\_}), {[}1, \operatorname{vzero}(\operatorname{baa\_})], \operatorname{unit\text{-}istar}(\operatorname{baa\_})]
\end{equation*}

\par\medskip\noindent\textbf{Definition} (\texttt{unit-vec}).\;$\operatorname{unit\text{-}vec}(\operatorname{baa\_}) :=$ $(\mathbb{C} \times \operatorname{vec}(\operatorname{baa\_}))$
\par\medskip\noindent\textbf{Definition} (\texttt{unit-vadd}).\;$\operatorname{unit\text{-}vadd}(\operatorname{baa\_}) :=$
\begin{equation*}
\begin{array}{@{}l@{}} (\lambda {[}\operatorname{bup\_}, \operatorname{buq\_}] \in (\operatorname{unit\text{-}vec}(\operatorname{baa\_}) \times \operatorname{unit\text{-}vec}(\operatorname{baa\_})).\; \\ \quad {[}(\operatorname{nth}(1, \operatorname{bup\_}) + \operatorname{nth}(1, \operatorname{buq\_})), \operatorname{vadd}(\operatorname{baa\_})(\operatorname{nth}(2, \operatorname{bup\_}), \operatorname{nth}(2, \operatorname{buq\_}))]) \end{array}
\end{equation*}

\par\medskip\noindent\textbf{Definition} (\texttt{unit-vneg}).\;$\operatorname{unit\text{-}vneg}(\operatorname{baa\_}) :=$
\begin{equation*}
(\lambda \operatorname{bup\_} \in \operatorname{unit\text{-}vec}(\operatorname{baa\_}).\; {[}-\operatorname{nth}(1, \operatorname{bup\_}), \operatorname{vneg}(\operatorname{baa\_})(\operatorname{nth}(2, \operatorname{bup\_}))])
\end{equation*}

\par\medskip\noindent\textbf{Definition} (\texttt{unit-act}).\;$\operatorname{unit\text{-}act}(\operatorname{baa\_}) :=$
\begin{equation*}
\begin{array}{@{}l@{}} (\lambda {[}\operatorname{buc\_}, \operatorname{bup\_}] \in (\mathbb{C} \times \operatorname{unit\text{-}vec}(\operatorname{baa\_})).\; \\ \quad {[}(\operatorname{buc\_} \cdot \operatorname{nth}(1, \operatorname{bup\_})), \operatorname{act}(\operatorname{baa\_})(\operatorname{buc\_}, \operatorname{nth}(2, \operatorname{bup\_}))]) \end{array}
\end{equation*}

\par\medskip\noindent\textbf{Definition} (\texttt{unit-vnrm}).\;$\operatorname{unit\text{-}vnrm}(\operatorname{baa\_}) :=$
\begin{equation*}
(\lambda \operatorname{bup\_} \in \operatorname{unit\text{-}vec}(\operatorname{baa\_}).\; (\operatorname{magnitude}(\operatorname{nth}(1, \operatorname{bup\_})) + \operatorname{vnrm}(\operatorname{baa\_})(\operatorname{nth}(2, \operatorname{bup\_}))))
\end{equation*}

\par\medskip\noindent\textbf{Definition} (\texttt{unit-vmul}).\;$\operatorname{unit\text{-}vmul}(\operatorname{baa\_}) :=$
\begin{equation*}
\begin{array}{@{}l@{}} (\lambda {[}\operatorname{bup\_}, \operatorname{buq\_}] \in (\operatorname{unit\text{-}vec}(\operatorname{baa\_}) \times \operatorname{unit\text{-}vec}(\operatorname{baa\_})).\; \\ \quad {[}(\operatorname{nth}(1, \operatorname{bup\_}) \cdot \operatorname{nth}(1, \operatorname{buq\_})), \operatorname{vadd}(\operatorname{baa\_})(\operatorname{vadd}(\operatorname{baa\_})(\operatorname{act}(\operatorname{baa\_})(\operatorname{nth}(1, \operatorname{bup\_}), \operatorname{nth}(2, \operatorname{buq\_})), \operatorname{act}(\operatorname{baa\_})(\operatorname{nth}(1, \operatorname{buq\_}), \operatorname{nth}(2, \operatorname{bup\_}))), \cdot_{\operatorname{baa\_}}(\operatorname{nth}(2, \operatorname{bup\_}), \operatorname{nth}(2, \operatorname{buq\_})))]) \end{array}
\end{equation*}

\par\medskip\noindent\textbf{Definition} (\texttt{unit-istar}).\;$\operatorname{unit\text{-}istar}(\operatorname{baa\_}) :=$
\begin{equation*}
(\lambda \operatorname{bup\_} \in \operatorname{unit\text{-}vec}(\operatorname{baa\_}).\; {[}\operatorname{conjugate}(\operatorname{nth}(1, \operatorname{bup\_})), {}^{*_{\operatorname{baa\_}}}(\operatorname{nth}(2, \operatorname{bup\_}))])
\end{equation*}

\par\medskip\noindent\textbf{Definition} (\texttt{characters}).\;$\operatorname{characters}(\operatorname{baa\_}) :=$
\begin{equation*}
\{\operatorname{bac\_} \in (\operatorname{vec}(\operatorname{baa\_}) \to \mathbb{C}) : (\operatorname{is\text{-}cc\text{-}linear\text{-}on}(\operatorname{baa\_}, \operatorname{bac\_}) \wedge (\forall \operatorname{bax\_}, \operatorname{bay\_} \in \operatorname{vec}(\operatorname{baa\_}).\; (\operatorname{bac\_}(\cdot_{\operatorname{baa\_}}(\operatorname{bax\_}, \operatorname{bay\_})) = (\operatorname{bac\_}(\operatorname{bax\_}) \cdot \operatorname{bac\_}(\operatorname{bay\_}))) \wedge (\operatorname{bac\_}(1_{\operatorname{baa\_}}) = 1)))\}
\end{equation*}

\par\medskip\noindent\textbf{Definition} (\texttt{is-cc-linear-on}).\;For all $\operatorname{baa\_}$ and $\operatorname{baf\_}$, 
\begin{equation*}
\begin{array}{@{}l@{}} \operatorname{is\text{-}cc\text{-}linear\text{-}on}(\operatorname{baa\_}, \operatorname{baf\_}) \\ \Leftrightarrow\; \operatorname{baf\_} \in (\operatorname{vec}(\operatorname{baa\_}) \to \mathbb{C}) \wedge \\ \forall \operatorname{bax\_}, \operatorname{bay\_} \in \operatorname{vec}(\operatorname{baa\_}). \\ \quad \operatorname{baf\_}(\operatorname{vadd}(\operatorname{baa\_})(\operatorname{bax\_}, \operatorname{bay\_})) = +(\operatorname{baf\_}(\operatorname{bax\_}), \operatorname{baf\_}(\operatorname{bay\_})) \wedge \\ \forall \operatorname{baz\_} \in \mathbb{C}, \operatorname{bax\_} \in \operatorname{vec}(\operatorname{baa\_}). \\ \quad \operatorname{baf\_}(\operatorname{act}(\operatorname{baa\_})(\operatorname{baz\_}, \operatorname{bax\_})) = \cdot(\operatorname{baz\_}, \operatorname{baf\_}(\operatorname{bax\_})) \end{array}
\end{equation*}

\par\medskip\noindent\textbf{Definition} (\texttt{cdual-vec}).\;$\operatorname{cdual\text{-}vec}(\operatorname{baa\_}) :=$
\begin{equation*}
\{\operatorname{baf\_} \in (\operatorname{vec}(\operatorname{baa\_}) \to \mathbb{C}) : (\operatorname{is\text{-}cc\text{-}linear\text{-}on}(\operatorname{baa\_}, \operatorname{baf\_}) \wedge \exists \operatorname{bam\_} \in \mathbb{R}.\; \forall \operatorname{bax\_} \in \operatorname{vec}(\operatorname{baa\_}).\; (\operatorname{magnitude}(\operatorname{baf\_}(\operatorname{bax\_})) \leq (\operatorname{bam\_} \cdot \operatorname{vnrm}(\operatorname{baa\_})(\operatorname{bax\_}))))\}
\end{equation*}

\par\medskip\noindent\textbf{Definition} (\texttt{cdual-norm}).\;$\operatorname{cdual\text{-}norm}(\operatorname{baa\_}, \operatorname{baf\_}) :=$
\begin{equation*}
\begin{array}{@{}l@{}} \operatorname{esup}(\operatorname{image}((\lambda \operatorname{bdx\_} \in \{\operatorname{bdy\_} \in \operatorname{vec}(\operatorname{baa\_}) : (\operatorname{vnrm}(\operatorname{baa\_})(\operatorname{bdy\_}) \leq 1)\}.\; \operatorname{magnitude}(\operatorname{baf\_}(\operatorname{bdx\_}))), \\ \quad \quad \{\operatorname{bdy\_} \in \operatorname{vec}(\operatorname{baa\_}) : (\operatorname{vnrm}(\operatorname{baa\_})(\operatorname{bdy\_}) \leq 1)\})) \end{array}
\end{equation*}

\par\medskip\noindent\textbf{Definition} (\texttt{gelfand-space}).\;$\widehat{\operatorname{baa\_}} :=$
\begin{equation*}
\begin{array}{@{}l@{}} \operatorname{subspace\text{-}top}(\operatorname{tproduct\text{-}top}(\operatorname{vec}(\operatorname{baa\_}), (\lambda \operatorname{bgx\_} \in \operatorname{vec}(\operatorname{baa\_}).\; \operatorname{metric\text{-}top}(\operatorname{cc\text{-}ms}))), \\ \quad \operatorname{characters}(\operatorname{baa\_})) \end{array}
\end{equation*}

\par\medskip\noindent\textbf{Definition} (\texttt{gelfand-transform}).\;$\operatorname{gelfand\text{-}transform}(\operatorname{baa\_}) :=$
\begin{equation*}
(\lambda \operatorname{bgx\_} \in \operatorname{vec}(\operatorname{baa\_}).\; (\lambda \operatorname{bgc\_} \in \operatorname{characters}(\operatorname{baa\_}).\; \operatorname{bgc\_}(\operatorname{bgx\_})))
\end{equation*}

\par\medskip\noindent\textbf{Definition} (\texttt{gelfand-dual-map}).\;$\operatorname{gelfand\text{-}dual\text{-}map}(\operatorname{baa\_}, \operatorname{bab\_}, \operatorname{baf\_}) :=$
\begin{equation*}
(\lambda \operatorname{bgw\_} \in \operatorname{characters}(\operatorname{bab\_}).\; (\lambda \operatorname{bgx\_} \in \operatorname{vec}(\operatorname{baa\_}).\; \operatorname{bgw\_}(\operatorname{baf\_}(\operatorname{bgx\_}))))
\end{equation*}

\par\medskip\noindent\textbf{Definition} (\texttt{hom-nu-c-star-algebra}).\;$\operatorname{hom\text{-}nu\text{-}c\text{-}star\text{-}algebra}(a, b) :=$
\begin{equation*}
\{\operatorname{homf\_} \in (\operatorname{vec}(a) \to \operatorname{vec}(b)) : \operatorname{is\text{-}hom\text{-}nu\text{-}c\text{-}star\text{-}algebra}(a, b, \operatorname{homf\_})\}
\end{equation*}

\par\medskip\noindent\textbf{Definition} (\texttt{hom-unital-banach-algebra}).\;$\operatorname{hom\text{-}unital\text{-}banach\text{-}algebra}(a, b) :=$
\begin{equation*}
\{\operatorname{homf\_} \in (\operatorname{vec}(a) \to \operatorname{vec}(b)) : \operatorname{is\text{-}hom\text{-}unital\text{-}banach\text{-}algebra}(a, b, \operatorname{homf\_})\}
\end{equation*}

\par\medskip\noindent\textbf{Definition} (\texttt{hom-unital-involutive-banach-algebra}).\;$\operatorname{hom\text{-}unital\text{-}involutive\text{-}banach\text{-}algebra}(a, b) :=$
\begin{equation*}
\{\operatorname{homf\_} \in (\operatorname{vec}(a) \to \operatorname{vec}(b)) : \operatorname{is\text{-}hom\text{-}unital\text{-}involutive\text{-}banach\text{-}algebra}(a, b, \operatorname{homf\_})\}
\end{equation*}

\par\medskip\noindent\textbf{Definition} (\texttt{hom-unital-c-star-algebra}).\;$\operatorname{hom\text{-}unital\text{-}c\text{-}star\text{-}algebra}(a, b) :=$
\begin{equation*}
\{\operatorname{homf\_} \in (\operatorname{vec}(a) \to \operatorname{vec}(b)) : \operatorname{is\text{-}hom\text{-}unital\text{-}c\text{-}star\text{-}algebra}(a, b, \operatorname{homf\_})\}
\end{equation*}

\par\medskip\noindent\textbf{Definition} (\texttt{hom-commutative-c-star-algebra}).\;$\operatorname{hom\text{-}commutative\text{-}c\text{-}star\text{-}algebra}(a, b) :=$
\begin{equation*}
\{\operatorname{homf\_} \in (\operatorname{vec}(a) \to \operatorname{vec}(b)) : \operatorname{is\text{-}hom\text{-}commutative\text{-}c\text{-}star\text{-}algebra}(a, b, \operatorname{homf\_})\}
\end{equation*}

\par\medskip\noindent\textbf{Definition} (\texttt{hom-commutative-banach-algebra}).\;$\operatorname{hom\text{-}commutative\text{-}banach\text{-}algebra}(a, b) :=$
\begin{equation*}
\{\operatorname{homf\_} \in (\operatorname{vec}(a) \to \operatorname{vec}(b)) : \operatorname{is\text{-}hom\text{-}commutative\text{-}banach\text{-}algebra}(a, b, \operatorname{homf\_})\}
\end{equation*}

\par\medskip\noindent\textbf{Definition} (\texttt{hom-nu-involutive-banach-algebra}).\;$\operatorname{hom\text{-}nu\text{-}involutive\text{-}banach\text{-}algebra}(a, b) :=$
\begin{equation*}
\{\operatorname{homf\_} \in (\operatorname{vec}(a) \to \operatorname{vec}(b)) : \operatorname{is\text{-}hom\text{-}nu\text{-}involutive\text{-}banach\text{-}algebra}(a, b, \operatorname{homf\_})\}
\end{equation*}

\par\medskip\noindent\textbf{Definition} (\texttt{hom-complex-normed-algebra}).\;$\operatorname{hom\text{-}complex\text{-}normed\text{-}algebra}(a, b) :=$
\begin{equation*}
\{\operatorname{homf\_} \in (\operatorname{vec}(a) \to \operatorname{vec}(b)) : \operatorname{is\text{-}hom\text{-}complex\text{-}normed\text{-}algebra}(a, b, \operatorname{homf\_})\}
\end{equation*}

\par\medskip\noindent\textbf{Definition} (\texttt{hom-complex-banach-algebra}).\;$\operatorname{hom\text{-}complex\text{-}banach\text{-}algebra}(a, b) :=$
\begin{equation*}
\{\operatorname{homf\_} \in (\operatorname{vec}(a) \to \operatorname{vec}(b)) : \operatorname{is\text{-}hom\text{-}complex\text{-}banach\text{-}algebra}(a, b, \operatorname{homf\_})\}
\end{equation*}

\vnbdef{is-hom-unital-c-star-algebra-def}{\vnbstmt{\begin{array}{@{}l@{}} \forall a, b, f. \\ \quad \operatorname{is\text{-}hom\text{-}unital\text{-}c\text{-}star\text{-}algebra}(a, b, f) \\ \quad \Leftrightarrow\; \operatorname{is\text{-}unital\text{-}c\text{-}star\text{-}algebra}(a) \wedge \\ \quad \operatorname{is\text{-}unital\text{-}c\text{-}star\text{-}algebra}(b) \wedge \\ \quad \operatorname{is\text{-}hom\text{-}unital\text{-}involutive\text{-}banach\text{-}algebra}(a, b, f) \end{array}}}
\vnbdef{complex-normed-algebra-class}{\vnbstmt{\forall s.\; ((s \in \operatorname{complex\text{-}normed\text{-}algebra}) \Leftrightarrow \operatorname{is\text{-}complex\text{-}normed\text{-}algebra}(s))}}
\vnbdef{nu-involutive-banach-algebra-class}{\vnbstmt{\begin{array}{@{}l@{}} \forall s. \\ \quad (s \in \operatorname{nu\text{-}involutive\text{-}banach\text{-}algebra}) \Leftrightarrow \operatorname{is\text{-}nu\text{-}involutive\text{-}banach\text{-}algebra}(s) \end{array}}}
\vnbdef{unital-involutive-banach-algebra-class}{\vnbstmt{\begin{array}{@{}l@{}} \forall s. \\ \quad (s \in \operatorname{unital\text{-}involutive\text{-}banach\text{-}algebra}) \Leftrightarrow \operatorname{is\text{-}unital\text{-}involutive\text{-}banach\text{-}algebra}(s) \end{array}}}
\vnbdef{is-hom-commutative-banach-algebra-def}{\vnbstmt{\begin{array}{@{}l@{}} \forall a, b, f. \\ \quad \operatorname{is\text{-}hom\text{-}commutative\text{-}banach\text{-}algebra}(a, b, f) \\ \quad \Leftrightarrow\; \operatorname{is\text{-}commutative\text{-}banach\text{-}algebra}(a) \wedge \\ \quad \operatorname{is\text{-}commutative\text{-}banach\text{-}algebra}(b) \wedge \\ \quad \operatorname{is\text{-}hom\text{-}unital\text{-}banach\text{-}algebra}(a, b, f) \end{array}}}
\vnbdef{is-hom-unital-involutive-banach-algebra-def}{\vnbstmt{\begin{array}{@{}l@{}} \forall a, b, f. \\ \quad \operatorname{is\text{-}hom\text{-}unital\text{-}involutive\text{-}banach\text{-}algebra}(a, b, f) \\ \quad \Leftrightarrow\; \operatorname{is\text{-}unital\text{-}involutive\text{-}banach\text{-}algebra}(a) \wedge \\ \quad \operatorname{is\text{-}unital\text{-}involutive\text{-}banach\text{-}algebra}(b) \wedge \\ \quad f \in (\operatorname{vec}(a) \to \operatorname{vec}(b)) \wedge \\ \quad \operatorname{scal}(a) = \operatorname{scal}(b) \wedge \\ \quad \forall x_{1}, x_{2} \in \operatorname{vec}(a). \\ \quad \quad f(\operatorname{vadd}(a)(x_{1}, x_{2})) = \operatorname{vadd}(b)(f(x_{1}), f(x_{2})) \wedge \\ \quad f(\operatorname{vzero}(a)) = \operatorname{vzero}(b) \wedge \\ \quad \forall x_{1} \in \operatorname{vec}(a).\; (f(\operatorname{vneg}(a)(x_{1})) = \operatorname{vneg}(b)(f(x_{1}))) \wedge \\ \quad \forall x_{1} \in \operatorname{carr}(\operatorname{scal}(a)), x_{2} \in \operatorname{vec}(a). \\ \quad \quad f(\operatorname{act}(a)(x_{1}, x_{2})) = \operatorname{act}(b)(x_{1}, f(x_{2})) \wedge \\ \quad \forall x_{1} \in \operatorname{vec}(a).\; (\operatorname{vnrm}(a)(x_{1}) = \operatorname{vnrm}(b)(f(x_{1}))) \wedge \\ \quad \forall x_{1}, x_{2} \in \operatorname{vec}(a). \\ \quad \quad f(\cdot_{a}(x_{1}, x_{2})) = \cdot_{b}(f(x_{1}), f(x_{2})) \wedge \\ \quad f(1_{a}) = 1_{b} \wedge \\ \quad \forall x_{1} \in \operatorname{vec}(a).\; (f({}^{*_{a}}(x_{1})) = {}^{*_{b}}(f(x_{1}))) \end{array}}}
\vnbdef{is-hom-nu-c-star-algebra-def}{\vnbstmt{\begin{array}{@{}l@{}} \forall a, b, f. \\ \quad \operatorname{is\text{-}hom\text{-}nu\text{-}c\text{-}star\text{-}algebra}(a, b, f) \\ \quad \Leftrightarrow\; (\operatorname{is\text{-}nu\text{-}c\text{-}star\text{-}algebra}(a) \wedge (\operatorname{is\text{-}nu\text{-}c\text{-}star\text{-}algebra}(b) \wedge \operatorname{is\text{-}hom\text{-}nu\text{-}involutive\text{-}banach\text{-}algebra}(a, b, f))) \end{array}}}
\vnbdef{is-hom-nu-involutive-banach-algebra-def}{\vnbstmt{\begin{array}{@{}l@{}} \forall a, b, f. \\ \quad \operatorname{is\text{-}hom\text{-}nu\text{-}involutive\text{-}banach\text{-}algebra}(a, b, f) \\ \quad \Leftrightarrow\; \operatorname{is\text{-}nu\text{-}involutive\text{-}banach\text{-}algebra}(a) \wedge \\ \quad \operatorname{is\text{-}nu\text{-}involutive\text{-}banach\text{-}algebra}(b) \wedge \\ \quad f \in (\operatorname{vec}(a) \to \operatorname{vec}(b)) \wedge \\ \quad \operatorname{scal}(a) = \operatorname{scal}(b) \wedge \\ \quad \forall x_{1}, x_{2} \in \operatorname{vec}(a). \\ \quad \quad f(\operatorname{vadd}(a)(x_{1}, x_{2})) = \operatorname{vadd}(b)(f(x_{1}), f(x_{2})) \wedge \\ \quad f(\operatorname{vzero}(a)) = \operatorname{vzero}(b) \wedge \\ \quad \forall x_{1} \in \operatorname{vec}(a).\; (f(\operatorname{vneg}(a)(x_{1})) = \operatorname{vneg}(b)(f(x_{1}))) \wedge \\ \quad \forall x_{1} \in \operatorname{carr}(\operatorname{scal}(a)), x_{2} \in \operatorname{vec}(a). \\ \quad \quad f(\operatorname{act}(a)(x_{1}, x_{2})) = \operatorname{act}(b)(x_{1}, f(x_{2})) \wedge \\ \quad \forall x_{1} \in \operatorname{vec}(a).\; (\operatorname{vnrm}(a)(x_{1}) = \operatorname{vnrm}(b)(f(x_{1}))) \wedge \\ \quad \forall x_{1}, x_{2} \in \operatorname{vec}(a). \\ \quad \quad f(\cdot_{a}(x_{1}, x_{2})) = \cdot_{b}(f(x_{1}), f(x_{2})) \wedge \\ \quad \forall x_{1} \in \operatorname{vec}(a).\; (f({}^{*_{a}}(x_{1})) = {}^{*_{b}}(f(x_{1}))) \end{array}}}
\vnbdef{is-hom-complex-banach-algebra-def}{\vnbstmt{\begin{array}{@{}l@{}} \forall a, b, f. \\ \quad \operatorname{is\text{-}hom\text{-}complex\text{-}banach\text{-}algebra}(a, b, f) \\ \quad \Leftrightarrow\; \operatorname{is\text{-}complex\text{-}banach\text{-}algebra}(a) \wedge \\ \quad \operatorname{is\text{-}complex\text{-}banach\text{-}algebra}(b) \wedge \\ \quad \operatorname{is\text{-}hom\text{-}complex\text{-}normed\text{-}algebra}(a, b, f) \end{array}}}
\vnbdef{is-hom-complex-normed-algebra-def}{\vnbstmt{\begin{array}{@{}l@{}} \forall a, b, f. \\ \quad \operatorname{is\text{-}hom\text{-}complex\text{-}normed\text{-}algebra}(a, b, f) \\ \quad \Leftrightarrow\; \operatorname{is\text{-}complex\text{-}normed\text{-}algebra}(a) \wedge \\ \quad \operatorname{is\text{-}complex\text{-}normed\text{-}algebra}(b) \wedge \\ \quad f \in (\operatorname{vec}(a) \to \operatorname{vec}(b)) \wedge \\ \quad \operatorname{scal}(a) = \operatorname{scal}(b) \wedge \\ \quad \forall x_{1}, x_{2} \in \operatorname{vec}(a). \\ \quad \quad f(\operatorname{vadd}(a)(x_{1}, x_{2})) = \operatorname{vadd}(b)(f(x_{1}), f(x_{2})) \wedge \\ \quad f(\operatorname{vzero}(a)) = \operatorname{vzero}(b) \wedge \\ \quad \forall x_{1} \in \operatorname{vec}(a).\; (f(\operatorname{vneg}(a)(x_{1})) = \operatorname{vneg}(b)(f(x_{1}))) \wedge \\ \quad \forall x_{1} \in \operatorname{carr}(\operatorname{scal}(a)), x_{2} \in \operatorname{vec}(a). \\ \quad \quad f(\operatorname{act}(a)(x_{1}, x_{2})) = \operatorname{act}(b)(x_{1}, f(x_{2})) \wedge \\ \quad \forall x_{1} \in \operatorname{vec}(a).\; (\operatorname{vnrm}(a)(x_{1}) = \operatorname{vnrm}(b)(f(x_{1}))) \wedge \\ \quad \forall x_{1}, x_{2} \in \operatorname{vec}(a). \\ \quad \quad f(\cdot_{a}(x_{1}, x_{2})) = \cdot_{b}(f(x_{1}), f(x_{2})) \end{array}}}
\vnbdef{complex-banach-algebra-class}{\vnbstmt{\forall s.\; ((s \in \operatorname{complex\text{-}banach\text{-}algebra}) \Leftrightarrow \operatorname{is\text{-}complex\text{-}banach\text{-}algebra}(s))}}
\vnbdef{commutative-banach-algebra-class}{\vnbstmt{\begin{array}{@{}l@{}} \forall s. \\ \quad (s \in \operatorname{commutative\text{-}banach\text{-}algebra}) \Leftrightarrow \operatorname{is\text{-}commutative\text{-}banach\text{-}algebra}(s) \end{array}}}
\vnbdef{is-unital-c-star-algebra-def}{\vnbstmt{\begin{array}{@{}l@{}} \forall s. \\ \quad \operatorname{is\text{-}unital\text{-}c\text{-}star\text{-}algebra}(s) \\ \quad \Leftrightarrow\; \operatorname{is\text{-}unital\text{-}involutive\text{-}banach\text{-}algebra}(s) \wedge \\ \quad \forall x \in \operatorname{vec}(s). \\ \quad \quad \operatorname{vnrm}(s)(\cdot_{s}({}^{*_{s}}(x), x)) = (\operatorname{vnrm}(s)(x) \cdot \operatorname{vnrm}(s)(x)) \end{array}}}
\vnbdef{is-commutative-c-star-algebra-def}{\vnbstmt{\begin{array}{@{}l@{}} \forall s. \\ \quad \operatorname{is\text{-}commutative\text{-}c\text{-}star\text{-}algebra}(s) \\ \quad \Leftrightarrow\; \operatorname{is\text{-}unital\text{-}c\text{-}star\text{-}algebra}(s) \wedge \\ \quad \forall x, y \in \operatorname{vec}(s).\; (\cdot_{s}(x, y) = \cdot_{s}(y, x)) \end{array}}}
\vnbdef{is-nu-c-star-algebra-def}{\vnbstmt{\begin{array}{@{}l@{}} \forall s. \\ \quad \operatorname{is\text{-}nu\text{-}c\text{-}star\text{-}algebra}(s) \\ \quad \Leftrightarrow\; \operatorname{is\text{-}nu\text{-}involutive\text{-}banach\text{-}algebra}(s) \wedge \\ \quad \forall x \in \operatorname{vec}(s). \\ \quad \quad \operatorname{vnrm}(s)(\cdot_{s}({}^{*_{s}}(x), x)) = (\operatorname{vnrm}(s)(x) \cdot \operatorname{vnrm}(s)(x)) \end{array}}}
\vnbdef{is-commutative-banach-algebra-def}{\vnbstmt{\begin{array}{@{}l@{}} \forall s. \\ \quad \operatorname{is\text{-}commutative\text{-}banach\text{-}algebra}(s) \\ \quad \Leftrightarrow\; \operatorname{is\text{-}unital\text{-}banach\text{-}algebra}(s) \wedge \\ \quad \forall x, y \in \operatorname{vec}(s).\; (\cdot_{s}(x, y) = \cdot_{s}(y, x)) \end{array}}}
\vnbdef{commutative-c-star-algebra-class}{\vnbstmt{\begin{array}{@{}l@{}} \forall s. \\ \quad (s \in \operatorname{commutative\text{-}c\text{-}star\text{-}algebra}) \Leftrightarrow \operatorname{is\text{-}commutative\text{-}c\text{-}star\text{-}algebra}(s) \end{array}}}
\vnbdef{is-hom-unital-banach-algebra-def}{\vnbstmt{\begin{array}{@{}l@{}} \forall a, b, f. \\ \quad \operatorname{is\text{-}hom\text{-}unital\text{-}banach\text{-}algebra}(a, b, f) \\ \quad \Leftrightarrow\; \operatorname{is\text{-}unital\text{-}banach\text{-}algebra}(a) \wedge \\ \quad \operatorname{is\text{-}unital\text{-}banach\text{-}algebra}(b) \wedge \\ \quad f \in (\operatorname{vec}(a) \to \operatorname{vec}(b)) \wedge \\ \quad \operatorname{scal}(a) = \operatorname{scal}(b) \wedge \\ \quad \forall x_{1}, x_{2} \in \operatorname{vec}(a). \\ \quad \quad f(\operatorname{vadd}(a)(x_{1}, x_{2})) = \operatorname{vadd}(b)(f(x_{1}), f(x_{2})) \wedge \\ \quad f(\operatorname{vzero}(a)) = \operatorname{vzero}(b) \wedge \\ \quad \forall x_{1} \in \operatorname{vec}(a).\; (f(\operatorname{vneg}(a)(x_{1})) = \operatorname{vneg}(b)(f(x_{1}))) \wedge \\ \quad \forall x_{1} \in \operatorname{carr}(\operatorname{scal}(a)), x_{2} \in \operatorname{vec}(a). \\ \quad \quad f(\operatorname{act}(a)(x_{1}, x_{2})) = \operatorname{act}(b)(x_{1}, f(x_{2})) \wedge \\ \quad \forall x_{1} \in \operatorname{vec}(a).\; (\operatorname{vnrm}(a)(x_{1}) = \operatorname{vnrm}(b)(f(x_{1}))) \wedge \\ \quad \forall x_{1}, x_{2} \in \operatorname{vec}(a). \\ \quad \quad f(\cdot_{a}(x_{1}, x_{2})) = \cdot_{b}(f(x_{1}), f(x_{2})) \wedge \\ \quad f(1_{a}) = 1_{b} \end{array}}}
\vnbdef{nu-c-star-algebra-class}{\vnbstmt{\forall s.\; ((s \in \operatorname{nu\text{-}c\text{-}star\text{-}algebra}) \Leftrightarrow \operatorname{is\text{-}nu\text{-}c\text{-}star\text{-}algebra}(s))}}
\vnbdef{is-complex-banach-algebra-def}{\vnbstmt{\begin{array}{@{}l@{}} \forall s. \\ \quad \operatorname{is\text{-}complex\text{-}banach\text{-}algebra}(s) \\ \quad \Leftrightarrow\; (\operatorname{is\text{-}complex\text{-}normed\text{-}algebra}(s) \wedge \operatorname{is\text{-}complete}(\operatorname{nag\text{-}metric\text{-}space}(\operatorname{cna\text{-}as\text{-}normed\text{-}ag}(s)))) \end{array}}}
\vnbdef{unital-banach-algebra-class}{\vnbstmt{\forall s.\; ((s \in \operatorname{unital\text{-}banach\text{-}algebra}) \Leftrightarrow \operatorname{is\text{-}unital\text{-}banach\text{-}algebra}(s))}}
\vnbdef{unital-c-star-algebra-class}{\vnbstmt{\forall s.\; ((s \in \operatorname{unital\text{-}c\text{-}star\text{-}algebra}) \Leftrightarrow \operatorname{is\text{-}unital\text{-}c\text{-}star\text{-}algebra}(s))}}
\vnbdef{is-hom-commutative-c-star-algebra-def}{\vnbstmt{\begin{array}{@{}l@{}} \forall a, b, f. \\ \quad \operatorname{is\text{-}hom\text{-}commutative\text{-}c\text{-}star\text{-}algebra}(a, b, f) \\ \quad \Leftrightarrow\; \operatorname{is\text{-}commutative\text{-}c\text{-}star\text{-}algebra}(a) \wedge \\ \quad \operatorname{is\text{-}commutative\text{-}c\text{-}star\text{-}algebra}(b) \wedge \\ \quad \operatorname{is\text{-}hom\text{-}unital\text{-}c\text{-}star\text{-}algebra}(a, b, f) \end{array}}}
\vnbprop{is-complex-normed-algebra-is-set}{\vnbstmt{\forall s.\; (\operatorname{is\text{-}complex\text{-}normed\text{-}algebra}(s) \Rightarrow (s \in \mathbf{Set}))}}
\vnbprop{is-complex-banach-algebra-is-set}{\vnbstmt{\forall s.\; (\operatorname{is\text{-}complex\text{-}banach\text{-}algebra}(s) \Rightarrow (s \in \mathbf{Set}))}}
\vnbprop{is-unital-banach-algebra-is-set}{\vnbstmt{\forall s.\; (\operatorname{is\text{-}unital\text{-}banach\text{-}algebra}(s) \Rightarrow (s \in \mathbf{Set}))}}
\vnbprop{is-commutative-banach-algebra-is-set}{\vnbstmt{\forall s.\; (\operatorname{is\text{-}commutative\text{-}banach\text{-}algebra}(s) \Rightarrow (s \in \mathbf{Set}))}}
\vnbprop{is-unital-involutive-banach-algebra-is-set}{\vnbstmt{\forall s.\; (\operatorname{is\text{-}unital\text{-}involutive\text{-}banach\text{-}algebra}(s) \Rightarrow (s \in \mathbf{Set}))}}
\vnbprop{is-unital-c-star-algebra-is-set}{\vnbstmt{\forall s.\; (\operatorname{is\text{-}unital\text{-}c\text{-}star\text{-}algebra}(s) \Rightarrow (s \in \mathbf{Set}))}}
\vnbprop{is-commutative-c-star-algebra-is-set}{\vnbstmt{\forall s.\; (\operatorname{is\text{-}commutative\text{-}c\text{-}star\text{-}algebra}(s) \Rightarrow (s \in \mathbf{Set}))}}
\vnbprop{is-nu-involutive-banach-algebra-is-set}{\vnbstmt{\forall s.\; (\operatorname{is\text{-}nu\text{-}involutive\text{-}banach\text{-}algebra}(s) \Rightarrow (s \in \mathbf{Set}))}}
\vnbprop{is-nu-c-star-algebra-is-set}{\vnbstmt{\forall s.\; (\operatorname{is\text{-}nu\text{-}c\text{-}star\text{-}algebra}(s) \Rightarrow (s \in \mathbf{Set}))}}

\vnbbookfile{theorem-library/ba-subtype-laws}{theorem-library/ba-subtype-laws}
\noindent{\small ba-subtype-laws.scm -- the refinement step of COMPLEX-BANACH-ALGEBRA, proven. COMPLEX-BANACH-ALGEBRA refines COMPLEX-NORMED-ALGEBRA (same-shape-as, plus completeness). functoriality.scm climbs a refinement chain by citing X-is-PARENT; without this step the functoriality of UBA-AS-CBA and NUIBA-AS-CBA (target ...}\par
\vnbprop{cna-as-module-is-module}{\vnbstmt{\forall r.\; (\operatorname{is\text{-}complex\text{-}normed\text{-}algebra}(r) \Rightarrow \operatorname{is\text{-}module}(\operatorname{cna\text{-}as\text{-}module}(r)))}}
\vnbprop{uba-as-cba-is-complex-banach-algebra}{\vnbstmt{\begin{array}{@{}l@{}} \forall r. \\ \quad \operatorname{is\text{-}unital\text{-}banach\text{-}algebra}(r) \Rightarrow \operatorname{is\text{-}complex\text{-}banach\text{-}algebra}(\operatorname{uba\text{-}as\text{-}cba}(r)) \end{array}}}
\vnbprop{uba-as-ring-is-ring}{\vnbstmt{\forall r.\; (\operatorname{is\text{-}unital\text{-}banach\text{-}algebra}(r) \Rightarrow \operatorname{is\text{-}ring}(\operatorname{uba\text{-}as\text{-}ring}(r)))}}
\vnbprop{uba-as-module-is-module}{\vnbstmt{\forall r.\; (\operatorname{is\text{-}unital\text{-}banach\text{-}algebra}(r) \Rightarrow \operatorname{is\text{-}module}(\operatorname{uba\text{-}as\text{-}module}(r)))}}
\vnbprop{uiba-as-uba-is-unital-banach-algebra}{\vnbstmt{\begin{array}{@{}l@{}} \forall r. \\ \quad \operatorname{is\text{-}unital\text{-}involutive\text{-}banach\text{-}algebra}(r) \Rightarrow  \\ \quad \quad \operatorname{is\text{-}unital\text{-}banach\text{-}algebra}(\operatorname{uiba\text{-}as\text{-}uba}(r)) \end{array}}}
\vnbprop{nuiba-as-cba-is-complex-banach-algebra}{\vnbstmt{\begin{array}{@{}l@{}} \forall r. \\ \quad \operatorname{is\text{-}nu\text{-}involutive\text{-}banach\text{-}algebra}(r) \Rightarrow  \\ \quad \quad \operatorname{is\text{-}complex\text{-}banach\text{-}algebra}(\operatorname{nuiba\text{-}as\text{-}cba}(r)) \end{array}}}
\vnbprop{uiba-as-nuiba-is-nu-involutive-banach-algebra}{\vnbstmt{\begin{array}{@{}l@{}} \forall r. \\ \quad \operatorname{is\text{-}unital\text{-}involutive\text{-}banach\text{-}algebra}(r) \Rightarrow  \\ \quad \quad \operatorname{is\text{-}nu\text{-}involutive\text{-}banach\text{-}algebra}(\operatorname{uiba\text{-}as\text{-}nuiba}(r)) \end{array}}}
\vnbprop{complex-banach-algebra-is-complex-normed-algebra}{\vnbstmt{\forall s.\; (\operatorname{is\text{-}complex\text{-}banach\text{-}algebra}(s) \Rightarrow \operatorname{is\text{-}complex\text{-}normed\text{-}algebra}(s))}}

\vnbbookfile{structure-library/ba-cx-vocab}{structure-library/ba-cx-vocab}
\noindent{\small ba-cx-vocab.scm -- wave of 2026-11-09/10 (agent BA-1), wired 2026-11-10 from scratchpad/ba1-1109/ba-cx-vocab.scm. Library: operator-algebras.}\par
\par\medskip\noindent\textbf{Definition} (\texttt{ccont}).\;$C({\operatorname{kxt\_}}) :=$
\begin{equation*}
\{\operatorname{kxf\_} \in (\operatorname{pts}(\operatorname{kxt\_}) \to \mathbb{C}) : \operatorname{is\text{-}hom\text{-}top\text{-}space}(\operatorname{kxt\_}, \operatorname{metric\text{-}top}(\operatorname{cc\text{-}ms}), \operatorname{kxf\_})\}
\end{equation*}

\par\medskip\noindent\textbf{Definition} (\texttt{sup-norm-t}).\;$\lVert \operatorname{kxf\_} \rVert_{\infty} :=$
\begin{equation*}
\operatorname{esup}(\operatorname{image}((\lambda \operatorname{kxp\_} \in \operatorname{pts}(\operatorname{kxt\_}).\; \operatorname{magnitude}(\operatorname{kxf\_}(\operatorname{kxp\_}))), \operatorname{pts}(\operatorname{kxt\_})))
\end{equation*}

\par\medskip\noindent\textbf{Definition} (\texttt{ccont-algebra}).\;$\operatorname{ccont\text{-}algebra}(\operatorname{kxt\_}) :=$
\begin{equation*}
{[}\operatorname{normed\text{-}field\text{-}as\text{-}commutative\text{-}ring}(\operatorname{cc\text{-}normed\text{-}field}), C({\operatorname{kxt\_}}), \operatorname{ccont\text{-}vadd}(\operatorname{kxt\_}), \operatorname{ccont\text{-}vzero}(\operatorname{kxt\_}), \operatorname{ccont\text{-}vneg}(\operatorname{kxt\_}), \operatorname{ccont\text{-}act}(\operatorname{kxt\_}), \operatorname{ccont\text{-}vnrm}(\operatorname{kxt\_}), \operatorname{ccont\text{-}vmul}(\operatorname{kxt\_}), \operatorname{ccont\text{-}vone}(\operatorname{kxt\_}), \operatorname{ccont\text{-}vstar}(\operatorname{kxt\_})]
\end{equation*}

\par\medskip\noindent\textbf{Definition} (\texttt{ccont-vadd}).\;$\operatorname{ccont\text{-}vadd}(\operatorname{kxt\_}) :=$
\begin{equation*}
\begin{array}{@{}l@{}} (\lambda {[}\operatorname{kxf\_}, \operatorname{kxg\_}] \in (C({\operatorname{kxt\_}}) \times C({\operatorname{kxt\_}})).\; \\ \quad (\lambda \operatorname{kxp\_} \in \operatorname{pts}(\operatorname{kxt\_}).\; (\operatorname{kxf\_}(\operatorname{kxp\_}) + \operatorname{kxg\_}(\operatorname{kxp\_})))) \end{array}
\end{equation*}

\par\medskip\noindent\textbf{Definition} (\texttt{ccont-vzero}).\;$\operatorname{ccont\text{-}vzero}(\operatorname{kxt\_}) :=$ $(\lambda \operatorname{kxp\_} \in \operatorname{pts}(\operatorname{kxt\_}).\; 0)$
\par\medskip\noindent\textbf{Definition} (\texttt{ccont-vneg}).\;$\operatorname{ccont\text{-}vneg}(\operatorname{kxt\_}) :=$
\begin{equation*}
(\lambda \operatorname{kxf\_} \in C({\operatorname{kxt\_}}).\; (\lambda \operatorname{kxp\_} \in \operatorname{pts}(\operatorname{kxt\_}).\; -\operatorname{kxf\_}(\operatorname{kxp\_})))
\end{equation*}

\par\medskip\noindent\textbf{Definition} (\texttt{ccont-act}).\;$\operatorname{ccont\text{-}act}(\operatorname{kxt\_}) :=$
\begin{equation*}
(\lambda {[}\operatorname{kxc\_}, \operatorname{kxf\_}] \in (\mathbb{C} \times C({\operatorname{kxt\_}})).\; (\lambda \operatorname{kxp\_} \in \operatorname{pts}(\operatorname{kxt\_}).\; (\operatorname{kxc\_} \cdot \operatorname{kxf\_}(\operatorname{kxp\_}))))
\end{equation*}

\par\medskip\noindent\textbf{Definition} (\texttt{ccont-vnrm}).\;$\operatorname{ccont\text{-}vnrm}(\operatorname{kxt\_}) :=$ $(\lambda \operatorname{kxf\_} \in C({\operatorname{kxt\_}}).\; \lVert \operatorname{kxf\_} \rVert_{\infty})$
\par\medskip\noindent\textbf{Definition} (\texttt{ccont-vmul}).\;$\operatorname{ccont\text{-}vmul}(\operatorname{kxt\_}) :=$
\begin{equation*}
\begin{array}{@{}l@{}} (\lambda {[}\operatorname{kxf\_}, \operatorname{kxg\_}] \in (C({\operatorname{kxt\_}}) \times C({\operatorname{kxt\_}})).\; \\ \quad (\lambda \operatorname{kxp\_} \in \operatorname{pts}(\operatorname{kxt\_}).\; (\operatorname{kxf\_}(\operatorname{kxp\_}) \cdot \operatorname{kxg\_}(\operatorname{kxp\_})))) \end{array}
\end{equation*}

\par\medskip\noindent\textbf{Definition} (\texttt{ccont-vone}).\;$\operatorname{ccont\text{-}vone}(\operatorname{kxt\_}) :=$ $(\lambda \operatorname{kxp\_} \in \operatorname{pts}(\operatorname{kxt\_}).\; 1)$
\par\medskip\noindent\textbf{Definition} (\texttt{ccont-vstar}).\;$\operatorname{ccont\text{-}vstar}(\operatorname{kxt\_}) :=$
\begin{equation*}
(\lambda \operatorname{kxf\_} \in C({\operatorname{kxt\_}}).\; (\lambda \operatorname{kxp\_} \in \operatorname{pts}(\operatorname{kxt\_}).\; \operatorname{conjugate}(\operatorname{kxf\_}(\operatorname{kxp\_}))))
\end{equation*}


\vnbbookfile{structure-library/ba-neumann-vocab}{structure-library/ba-neumann-vocab}
\noindent{\small ba-neumann-vocab.scm -- wave of 2026-11-09/10 (agent BA-1), wired 2026-11-10 from scratchpad/ba1-1109/ba-neumann-vocab.scm. Library: operator-algebras.}\par
\par\medskip\noindent\textbf{Definition} (\texttt{neumann-sum}).\;For all $\operatorname{bna\_}$ and $\operatorname{bnx\_}$,  $(\sum_{k < 0} {\operatorname{bnx\_}}^{k} \simeq \operatorname{vzero}(\operatorname{bna\_}))$\par
Let $\operatorname{bna\_}$ and $\operatorname{bnx\_}$ be arbitrary. Let $\operatorname{bnk\_} \in \mathbb{N}$. Then
\begin{equation*}
(\sum_{k < \operatorname{succ}(\operatorname{bnk\_})} {\operatorname{bnx\_}}^{k} \simeq \operatorname{vadd}(\operatorname{bna\_})(\sum_{k < \operatorname{bnk\_}} {\operatorname{bnx\_}}^{k}, {\operatorname{bnx\_}}^{\operatorname{bnk\_}}))
\end{equation*}


\vnbbookfile{theorem-library/ba-laws}{theorem-library/ba-laws}
\noindent{\small ba-laws.scm -- wave of 2026-11-09/10 (agent BA-1), wired 2026-11-10 from scratchpad/ba1-1109/ba-laws.scm. Library: operator-algebras.}\par
\vnbprop{bal-cna-scal-pin}{\vnbstmt{\begin{array}{@{}l@{}} \forall a. \\ \quad \operatorname{is\text{-}complex\text{-}normed\text{-}algebra}(a) \Rightarrow  \\ \quad \quad \operatorname{scal}(a) = \operatorname{normed\text{-}field\text{-}as\text{-}commutative\text{-}ring}(\operatorname{cc\text{-}normed\text{-}field}) \end{array}}}
\vnbprop{bal-cna-scal-carr}{\vnbstmt{\forall a.\; (\operatorname{is\text{-}complex\text{-}normed\text{-}algebra}(a) \Rightarrow (\operatorname{carr}(\operatorname{scal}(a)) = \mathbb{C}))}}
\vnbprop{bal-cna-vec-set}{\vnbstmt{\forall a.\; (\operatorname{is\text{-}complex\text{-}normed\text{-}algebra}(a) \Rightarrow (\operatorname{vec}(a) \in \mathbf{Set}))}}
\vnbprop{bal-cna-vadd-fun}{\vnbstmt{\begin{array}{@{}l@{}} \forall a. \\ \quad \operatorname{is\text{-}complex\text{-}normed\text{-}algebra}(a) \Rightarrow  \\ \quad \quad \operatorname{vadd}(a) \in ((\operatorname{vec}(a) \times \operatorname{vec}(a)) \to \operatorname{vec}(a)) \end{array}}}
\vnbprop{bal-cna-vzero-in}{\vnbstmt{\forall a.\; (\operatorname{is\text{-}complex\text{-}normed\text{-}algebra}(a) \Rightarrow (\operatorname{vzero}(a) \in \operatorname{vec}(a)))}}
\vnbprop{bal-cna-vneg-fun}{\vnbstmt{\forall a.\; (\operatorname{is\text{-}complex\text{-}normed\text{-}algebra}(a) \Rightarrow (\operatorname{vneg}(a) \in (\operatorname{vec}(a) \to \operatorname{vec}(a))))}}
\vnbprop{bal-cna-act-fun}{\vnbstmt{\begin{array}{@{}l@{}} \forall a. \\ \quad \operatorname{is\text{-}complex\text{-}normed\text{-}algebra}(a) \Rightarrow  \\ \quad \quad \operatorname{act}(a) \in ((\operatorname{carr}(\operatorname{scal}(a)) \times \operatorname{vec}(a)) \to \operatorname{vec}(a)) \end{array}}}
\vnbprop{bal-cna-vnrm-fun}{\vnbstmt{\forall a.\; (\operatorname{is\text{-}complex\text{-}normed\text{-}algebra}(a) \Rightarrow (\operatorname{vnrm}(a) \in (\operatorname{vec}(a) \to \mathbb{R})))}}
\vnbprop{bal-cna-vmul-fun}{\vnbstmt{\begin{array}{@{}l@{}} \forall a. \\ \quad \operatorname{is\text{-}complex\text{-}normed\text{-}algebra}(a) \Rightarrow  \\ \quad \quad \cdot_{a} \in ((\operatorname{vec}(a) \times \operatorname{vec}(a)) \to \operatorname{vec}(a)) \end{array}}}
\vnbprop{bal-cna-vadd-assoc}{\vnbstmt{\begin{array}{@{}l@{}} \forall a. \\ \quad \operatorname{is\text{-}complex\text{-}normed\text{-}algebra}(a) \Rightarrow \operatorname{is\text{-}associative}(\operatorname{vadd}(a), \operatorname{vec}(a)) \end{array}}}
\vnbprop{bal-cna-vadd-comm}{\vnbstmt{\begin{array}{@{}l@{}} \forall a. \\ \quad \operatorname{is\text{-}complex\text{-}normed\text{-}algebra}(a) \Rightarrow \operatorname{is\text{-}commutative}(\operatorname{vadd}(a), \operatorname{vec}(a)) \end{array}}}
\vnbprop{bal-cna-vadd-ident}{\vnbstmt{\begin{array}{@{}l@{}} \forall a. \\ \quad \operatorname{is\text{-}complex\text{-}normed\text{-}algebra}(a) \Rightarrow \operatorname{is\text{-}identity}(\operatorname{vadd}(a), \operatorname{vzero}(a), \operatorname{vec}(a)) \end{array}}}
\vnbprop{bal-cna-vadd-inv}{\vnbstmt{\begin{array}{@{}l@{}} \forall a. \\ \quad \operatorname{is\text{-}complex\text{-}normed\text{-}algebra}(a) \Rightarrow  \\ \quad \quad \operatorname{has\text{-}inverses}(\operatorname{vadd}(a), \operatorname{vzero}(a), \operatorname{vneg}(a), \operatorname{vec}(a)) \end{array}}}
\vnbprop{bal-cna-vmul-assoc}{\vnbstmt{\begin{array}{@{}l@{}} \forall a. \\ \quad \operatorname{is\text{-}complex\text{-}normed\text{-}algebra}(a) \Rightarrow \operatorname{is\text{-}associative}(\cdot_{a}, \operatorname{vec}(a)) \end{array}}}
\vnbprop{bal-cna-act-vadd}{\vnbstmt{\begin{array}{@{}l@{}} \forall a.; \forall r \in \operatorname{carr}(\operatorname{scal}(a)), x, y \in \operatorname{vec}(a). \\ \quad \operatorname{is\text{-}complex\text{-}normed\text{-}algebra}(a) \Rightarrow  \\ \quad \quad \operatorname{act}(a)(r, \operatorname{vadd}(a)(x, y)) = \operatorname{vadd}(a)(\operatorname{act}(a)(r, x), \operatorname{act}(a)(r, y)) \end{array}}}
\vnbprop{bal-cna-act-add}{\vnbstmt{\begin{array}{@{}l@{}} \forall a.; \forall r, s \in \operatorname{carr}(\operatorname{scal}(a)), x \in \operatorname{vec}(a). \\ \quad \operatorname{is\text{-}complex\text{-}normed\text{-}algebra}(a) \Rightarrow  \\ \quad \quad \operatorname{act}(a)(\operatorname{add}(\operatorname{scal}(a))(r, s), x) = \operatorname{vadd}(a)(\operatorname{act}(a)(r, x), \operatorname{act}(a)(s, x)) \end{array}}}
\vnbprop{bal-cna-act-mul}{\vnbstmt{\begin{array}{@{}l@{}} \forall a.; \forall r, s \in \operatorname{carr}(\operatorname{scal}(a)), x \in \operatorname{vec}(a). \\ \quad \operatorname{is\text{-}complex\text{-}normed\text{-}algebra}(a) \Rightarrow  \\ \quad \quad \operatorname{act}(a)(\operatorname{mul}(\operatorname{scal}(a))(r, s), x) = \operatorname{act}(a)(r, \operatorname{act}(a)(s, x)) \end{array}}}
\vnbprop{bal-cna-act-one}{\vnbstmt{\begin{array}{@{}l@{}} \forall a.; \forall x \in \operatorname{vec}(a). \\ \quad \operatorname{is\text{-}complex\text{-}normed\text{-}algebra}(a) \Rightarrow (\operatorname{act}(a)(\operatorname{one}(\operatorname{scal}(a)), x) = x) \end{array}}}
\vnbprop{bal-cna-vnrm-nonneg}{\vnbstmt{\forall a.; \forall x \in \operatorname{vec}(a).\; (\operatorname{is\text{-}complex\text{-}normed\text{-}algebra}(a) \Rightarrow (0 \leq \operatorname{vnrm}(a)(x)))}}
\vnbprop{bal-cna-vnrm-definite}{\vnbstmt{\begin{array}{@{}l@{}} \forall a.; \forall x \in \operatorname{vec}(a). \\ \quad \operatorname{is\text{-}complex\text{-}normed\text{-}algebra}(a) \Rightarrow ((\operatorname{vnrm}(a)(x) = 0) \Leftrightarrow (x = \operatorname{vzero}(a))) \end{array}}}
\vnbprop{bal-cna-vnrm-homog}{\vnbstmt{\begin{array}{@{}l@{}} \forall a.; \forall r \in \operatorname{carr}(\operatorname{scal}(a)), x \in \operatorname{vec}(a). \\ \quad \operatorname{is\text{-}complex\text{-}normed\text{-}algebra}(a) \Rightarrow  \\ \quad \quad \operatorname{vnrm}(a)(\operatorname{act}(a)(r, x)) = (\operatorname{magnitude}(r) \cdot \operatorname{vnrm}(a)(x)) \end{array}}}
\vnbprop{bal-cna-vnrm-triangle}{\vnbstmt{\begin{array}{@{}l@{}} \forall a.; \forall x, y \in \operatorname{vec}(a). \\ \quad \operatorname{is\text{-}complex\text{-}normed\text{-}algebra}(a) \Rightarrow  \\ \quad \quad \operatorname{vnrm}(a)(\operatorname{vadd}(a)(x, y)) \leq (\operatorname{vnrm}(a)(x) + \operatorname{vnrm}(a)(y)) \end{array}}}
\vnbprop{bal-cna-vmul-dist-left}{\vnbstmt{\begin{array}{@{}l@{}} \forall a.; \forall x, y, z \in \operatorname{vec}(a). \\ \quad \operatorname{is\text{-}complex\text{-}normed\text{-}algebra}(a) \Rightarrow  \\ \quad \quad \cdot_{a}(x, \operatorname{vadd}(a)(y, z)) = \operatorname{vadd}(a)(\cdot_{a}(x, y), \cdot_{a}(x, z)) \end{array}}}
\vnbprop{bal-cna-vmul-dist-right}{\vnbstmt{\begin{array}{@{}l@{}} \forall a.; \forall x, y, z \in \operatorname{vec}(a). \\ \quad \operatorname{is\text{-}complex\text{-}normed\text{-}algebra}(a) \Rightarrow  \\ \quad \quad \cdot_{a}(\operatorname{vadd}(a)(x, y), z) = \operatorname{vadd}(a)(\cdot_{a}(x, z), \cdot_{a}(y, z)) \end{array}}}
\vnbprop{bal-cna-vmul-act-left}{\vnbstmt{\begin{array}{@{}l@{}} \forall a.; \forall r \in \operatorname{carr}(\operatorname{scal}(a)), x, y \in \operatorname{vec}(a). \\ \quad \operatorname{is\text{-}complex\text{-}normed\text{-}algebra}(a) \Rightarrow  \\ \quad \quad \cdot_{a}(\operatorname{act}(a)(r, x), y) = \operatorname{act}(a)(r, \cdot_{a}(x, y)) \end{array}}}
\vnbprop{bal-cna-vmul-act-right}{\vnbstmt{\begin{array}{@{}l@{}} \forall a.; \forall r \in \operatorname{carr}(\operatorname{scal}(a)), x, y \in \operatorname{vec}(a). \\ \quad \operatorname{is\text{-}complex\text{-}normed\text{-}algebra}(a) \Rightarrow  \\ \quad \quad \cdot_{a}(x, \operatorname{act}(a)(r, y)) = \operatorname{act}(a)(r, \cdot_{a}(x, y)) \end{array}}}
\vnbprop{bal-cna-vnrm-submult}{\vnbstmt{\begin{array}{@{}l@{}} \forall a.; \forall x, y \in \operatorname{vec}(a). \\ \quad \operatorname{is\text{-}complex\text{-}normed\text{-}algebra}(a) \Rightarrow  \\ \quad \quad \operatorname{vnrm}(a)(\cdot_{a}(x, y)) \leq (\operatorname{vnrm}(a)(x) \cdot \operatorname{vnrm}(a)(y)) \end{array}}}
\vnbprop{bal-cba-is-cna}{\vnbstmt{\forall a.\; (\operatorname{is\text{-}complex\text{-}banach\text{-}algebra}(a) \Rightarrow \operatorname{is\text{-}complex\text{-}normed\text{-}algebra}(a))}}
\vnbprop{bal-cba-complete}{\vnbstmt{\begin{array}{@{}l@{}} \forall a. \\ \quad \operatorname{is\text{-}complex\text{-}banach\text{-}algebra}(a) \Rightarrow  \\ \quad \quad \operatorname{is\text{-}complete}(\operatorname{nag\text{-}metric\text{-}space}(\operatorname{cna\text{-}as\text{-}normed\text{-}ag}(a))) \end{array}}}
\vnbprop{bal-ucba-scal}{\vnbstmt{\forall a.\; (\operatorname{is\text{-}unital\text{-}banach\text{-}algebra}(a) \Rightarrow (\operatorname{scal}(\operatorname{uba\text{-}as\text{-}cba}(a)) \simeq \operatorname{scal}(a)))}}
\vnbprop{bal-ucba-vec}{\vnbstmt{\forall a.\; (\operatorname{is\text{-}unital\text{-}banach\text{-}algebra}(a) \Rightarrow (\operatorname{vec}(\operatorname{uba\text{-}as\text{-}cba}(a)) \simeq \operatorname{vec}(a)))}}
\vnbprop{bal-ucba-vadd}{\vnbstmt{\forall a.\; (\operatorname{is\text{-}unital\text{-}banach\text{-}algebra}(a) \Rightarrow (\operatorname{vadd}(\operatorname{uba\text{-}as\text{-}cba}(a)) \simeq \operatorname{vadd}(a)))}}
\vnbprop{bal-ucba-vzero}{\vnbstmt{\begin{array}{@{}l@{}} \forall a. \\ \quad \operatorname{is\text{-}unital\text{-}banach\text{-}algebra}(a) \Rightarrow (\operatorname{vzero}(\operatorname{uba\text{-}as\text{-}cba}(a)) \simeq \operatorname{vzero}(a)) \end{array}}}
\vnbprop{bal-ucba-vneg}{\vnbstmt{\forall a.\; (\operatorname{is\text{-}unital\text{-}banach\text{-}algebra}(a) \Rightarrow (\operatorname{vneg}(\operatorname{uba\text{-}as\text{-}cba}(a)) \simeq \operatorname{vneg}(a)))}}
\vnbprop{bal-ucba-act}{\vnbstmt{\forall a.\; (\operatorname{is\text{-}unital\text{-}banach\text{-}algebra}(a) \Rightarrow (\operatorname{act}(\operatorname{uba\text{-}as\text{-}cba}(a)) \simeq \operatorname{act}(a)))}}
\vnbprop{bal-ucba-vnrm}{\vnbstmt{\forall a.\; (\operatorname{is\text{-}unital\text{-}banach\text{-}algebra}(a) \Rightarrow (\operatorname{vnrm}(\operatorname{uba\text{-}as\text{-}cba}(a)) \simeq \operatorname{vnrm}(a)))}}
\vnbprop{bal-ucba-vmul}{\vnbstmt{\forall a.\; (\operatorname{is\text{-}unital\text{-}banach\text{-}algebra}(a) \Rightarrow (\cdot_{\operatorname{uba\text{-}as\text{-}cba}(a)} \simeq \cdot_{a}))}}
\vnbprop{bal-uba-scal-pin}{\vnbstmt{\begin{array}{@{}l@{}} \forall a. \\ \quad \operatorname{is\text{-}unital\text{-}banach\text{-}algebra}(a) \Rightarrow  \\ \quad \quad \operatorname{scal}(a) = \operatorname{normed\text{-}field\text{-}as\text{-}commutative\text{-}ring}(\operatorname{cc\text{-}normed\text{-}field}) \end{array}}}
\vnbprop{bal-uba-scal-carr}{\vnbstmt{\forall a.\; (\operatorname{is\text{-}unital\text{-}banach\text{-}algebra}(a) \Rightarrow (\operatorname{carr}(\operatorname{scal}(a)) = \mathbb{C}))}}
\vnbprop{bal-uba-vec-set}{\vnbstmt{\forall a.\; (\operatorname{is\text{-}unital\text{-}banach\text{-}algebra}(a) \Rightarrow (\operatorname{vec}(a) \in \mathbf{Set}))}}
\vnbprop{bal-uba-vadd-fun}{\vnbstmt{\begin{array}{@{}l@{}} \forall a. \\ \quad \operatorname{is\text{-}unital\text{-}banach\text{-}algebra}(a) \Rightarrow  \\ \quad \quad \operatorname{vadd}(a) \in ((\operatorname{vec}(a) \times \operatorname{vec}(a)) \to \operatorname{vec}(a)) \end{array}}}
\vnbprop{bal-uba-vzero-in}{\vnbstmt{\forall a.\; (\operatorname{is\text{-}unital\text{-}banach\text{-}algebra}(a) \Rightarrow (\operatorname{vzero}(a) \in \operatorname{vec}(a)))}}
\vnbprop{bal-uba-vneg-fun}{\vnbstmt{\forall a.\; (\operatorname{is\text{-}unital\text{-}banach\text{-}algebra}(a) \Rightarrow (\operatorname{vneg}(a) \in (\operatorname{vec}(a) \to \operatorname{vec}(a))))}}
\vnbprop{bal-uba-act-fun}{\vnbstmt{\begin{array}{@{}l@{}} \forall a. \\ \quad \operatorname{is\text{-}unital\text{-}banach\text{-}algebra}(a) \Rightarrow  \\ \quad \quad \operatorname{act}(a) \in ((\operatorname{carr}(\operatorname{scal}(a)) \times \operatorname{vec}(a)) \to \operatorname{vec}(a)) \end{array}}}
\vnbprop{bal-uba-vnrm-fun}{\vnbstmt{\forall a.\; (\operatorname{is\text{-}unital\text{-}banach\text{-}algebra}(a) \Rightarrow (\operatorname{vnrm}(a) \in (\operatorname{vec}(a) \to \mathbb{R})))}}
\vnbprop{bal-uba-vmul-fun}{\vnbstmt{\begin{array}{@{}l@{}} \forall a. \\ \quad \operatorname{is\text{-}unital\text{-}banach\text{-}algebra}(a) \Rightarrow  \\ \quad \quad \cdot_{a} \in ((\operatorname{vec}(a) \times \operatorname{vec}(a)) \to \operatorname{vec}(a)) \end{array}}}
\vnbprop{bal-uba-vadd-assoc}{\vnbstmt{\forall a.\; (\operatorname{is\text{-}unital\text{-}banach\text{-}algebra}(a) \Rightarrow \operatorname{is\text{-}associative}(\operatorname{vadd}(a), \operatorname{vec}(a)))}}
\vnbprop{bal-uba-vadd-comm}{\vnbstmt{\forall a.\; (\operatorname{is\text{-}unital\text{-}banach\text{-}algebra}(a) \Rightarrow \operatorname{is\text{-}commutative}(\operatorname{vadd}(a), \operatorname{vec}(a)))}}
\vnbprop{bal-uba-vadd-ident}{\vnbstmt{\begin{array}{@{}l@{}} \forall a. \\ \quad \operatorname{is\text{-}unital\text{-}banach\text{-}algebra}(a) \Rightarrow \operatorname{is\text{-}identity}(\operatorname{vadd}(a), \operatorname{vzero}(a), \operatorname{vec}(a)) \end{array}}}
\vnbprop{bal-uba-vadd-inv}{\vnbstmt{\begin{array}{@{}l@{}} \forall a. \\ \quad \operatorname{is\text{-}unital\text{-}banach\text{-}algebra}(a) \Rightarrow  \\ \quad \quad \operatorname{has\text{-}inverses}(\operatorname{vadd}(a), \operatorname{vzero}(a), \operatorname{vneg}(a), \operatorname{vec}(a)) \end{array}}}
\vnbprop{bal-uba-vmul-assoc}{\vnbstmt{\forall a.\; (\operatorname{is\text{-}unital\text{-}banach\text{-}algebra}(a) \Rightarrow \operatorname{is\text{-}associative}(\cdot_{a}, \operatorname{vec}(a)))}}
\vnbprop{bal-uba-act-vadd}{\vnbstmt{\begin{array}{@{}l@{}} \forall a.; \forall r \in \operatorname{carr}(\operatorname{scal}(a)), x, y \in \operatorname{vec}(a). \\ \quad \operatorname{is\text{-}unital\text{-}banach\text{-}algebra}(a) \Rightarrow  \\ \quad \quad \operatorname{act}(a)(r, \operatorname{vadd}(a)(x, y)) = \operatorname{vadd}(a)(\operatorname{act}(a)(r, x), \operatorname{act}(a)(r, y)) \end{array}}}
\vnbprop{bal-uba-act-add}{\vnbstmt{\begin{array}{@{}l@{}} \forall a.; \forall r, s \in \operatorname{carr}(\operatorname{scal}(a)), x \in \operatorname{vec}(a). \\ \quad \operatorname{is\text{-}unital\text{-}banach\text{-}algebra}(a) \Rightarrow  \\ \quad \quad \operatorname{act}(a)(\operatorname{add}(\operatorname{scal}(a))(r, s), x) = \operatorname{vadd}(a)(\operatorname{act}(a)(r, x), \operatorname{act}(a)(s, x)) \end{array}}}
\vnbprop{bal-uba-act-mul}{\vnbstmt{\begin{array}{@{}l@{}} \forall a.; \forall r, s \in \operatorname{carr}(\operatorname{scal}(a)), x \in \operatorname{vec}(a). \\ \quad \operatorname{is\text{-}unital\text{-}banach\text{-}algebra}(a) \Rightarrow  \\ \quad \quad \operatorname{act}(a)(\operatorname{mul}(\operatorname{scal}(a))(r, s), x) = \operatorname{act}(a)(r, \operatorname{act}(a)(s, x)) \end{array}}}
\vnbprop{bal-uba-act-one}{\vnbstmt{\begin{array}{@{}l@{}} \forall a.; \forall x \in \operatorname{vec}(a). \\ \quad \operatorname{is\text{-}unital\text{-}banach\text{-}algebra}(a) \Rightarrow (\operatorname{act}(a)(\operatorname{one}(\operatorname{scal}(a)), x) = x) \end{array}}}
\vnbprop{bal-uba-vnrm-nonneg}{\vnbstmt{\forall a.; \forall x \in \operatorname{vec}(a).\; (\operatorname{is\text{-}unital\text{-}banach\text{-}algebra}(a) \Rightarrow (0 \leq \operatorname{vnrm}(a)(x)))}}
\vnbprop{bal-uba-vnrm-definite}{\vnbstmt{\begin{array}{@{}l@{}} \forall a.; \forall x \in \operatorname{vec}(a). \\ \quad \operatorname{is\text{-}unital\text{-}banach\text{-}algebra}(a) \Rightarrow ((\operatorname{vnrm}(a)(x) = 0) \Leftrightarrow (x = \operatorname{vzero}(a))) \end{array}}}
\vnbprop{bal-uba-vnrm-homog}{\vnbstmt{\begin{array}{@{}l@{}} \forall a.; \forall r \in \operatorname{carr}(\operatorname{scal}(a)), x \in \operatorname{vec}(a). \\ \quad \operatorname{is\text{-}unital\text{-}banach\text{-}algebra}(a) \Rightarrow  \\ \quad \quad \operatorname{vnrm}(a)(\operatorname{act}(a)(r, x)) = (\operatorname{magnitude}(r) \cdot \operatorname{vnrm}(a)(x)) \end{array}}}
\vnbprop{bal-uba-vnrm-triangle}{\vnbstmt{\begin{array}{@{}l@{}} \forall a.; \forall x, y \in \operatorname{vec}(a). \\ \quad \operatorname{is\text{-}unital\text{-}banach\text{-}algebra}(a) \Rightarrow  \\ \quad \quad \operatorname{vnrm}(a)(\operatorname{vadd}(a)(x, y)) \leq (\operatorname{vnrm}(a)(x) + \operatorname{vnrm}(a)(y)) \end{array}}}
\vnbprop{bal-uba-vmul-dist-left}{\vnbstmt{\begin{array}{@{}l@{}} \forall a.; \forall x, y, z \in \operatorname{vec}(a). \\ \quad \operatorname{is\text{-}unital\text{-}banach\text{-}algebra}(a) \Rightarrow  \\ \quad \quad \cdot_{a}(x, \operatorname{vadd}(a)(y, z)) = \operatorname{vadd}(a)(\cdot_{a}(x, y), \cdot_{a}(x, z)) \end{array}}}
\vnbprop{bal-uba-vmul-dist-right}{\vnbstmt{\begin{array}{@{}l@{}} \forall a.; \forall x, y, z \in \operatorname{vec}(a). \\ \quad \operatorname{is\text{-}unital\text{-}banach\text{-}algebra}(a) \Rightarrow  \\ \quad \quad \cdot_{a}(\operatorname{vadd}(a)(x, y), z) = \operatorname{vadd}(a)(\cdot_{a}(x, z), \cdot_{a}(y, z)) \end{array}}}
\vnbprop{bal-uba-vmul-act-left}{\vnbstmt{\begin{array}{@{}l@{}} \forall a.; \forall r \in \operatorname{carr}(\operatorname{scal}(a)), x, y \in \operatorname{vec}(a). \\ \quad \operatorname{is\text{-}unital\text{-}banach\text{-}algebra}(a) \Rightarrow  \\ \quad \quad \cdot_{a}(\operatorname{act}(a)(r, x), y) = \operatorname{act}(a)(r, \cdot_{a}(x, y)) \end{array}}}
\vnbprop{bal-uba-vmul-act-right}{\vnbstmt{\begin{array}{@{}l@{}} \forall a.; \forall r \in \operatorname{carr}(\operatorname{scal}(a)), x, y \in \operatorname{vec}(a). \\ \quad \operatorname{is\text{-}unital\text{-}banach\text{-}algebra}(a) \Rightarrow  \\ \quad \quad \cdot_{a}(x, \operatorname{act}(a)(r, y)) = \operatorname{act}(a)(r, \cdot_{a}(x, y)) \end{array}}}
\vnbprop{bal-uba-vnrm-submult}{\vnbstmt{\begin{array}{@{}l@{}} \forall a.; \forall x, y \in \operatorname{vec}(a). \\ \quad \operatorname{is\text{-}unital\text{-}banach\text{-}algebra}(a) \Rightarrow  \\ \quad \quad \operatorname{vnrm}(a)(\cdot_{a}(x, y)) \leq (\operatorname{vnrm}(a)(x) \cdot \operatorname{vnrm}(a)(y)) \end{array}}}
\vnbprop{bal-uba-vone-in}{\vnbstmt{\forall a.\; (\operatorname{is\text{-}unital\text{-}banach\text{-}algebra}(a) \Rightarrow (1_{a} \in \operatorname{vec}(a)))}}
\vnbprop{bal-uba-vone-left}{\vnbstmt{\begin{array}{@{}l@{}} \forall a.; \forall x \in \operatorname{vec}(a). \\ \quad \operatorname{is\text{-}unital\text{-}banach\text{-}algebra}(a) \Rightarrow (\cdot_{a}(1_{a}, x) = x) \end{array}}}
\vnbprop{bal-uba-vone-right}{\vnbstmt{\begin{array}{@{}l@{}} \forall a.; \forall x \in \operatorname{vec}(a). \\ \quad \operatorname{is\text{-}unital\text{-}banach\text{-}algebra}(a) \Rightarrow (\cdot_{a}(x, 1_{a}) = x) \end{array}}}
\vnbprop{bal-uba-vone-norm}{\vnbstmt{\forall a.\; (\operatorname{is\text{-}unital\text{-}banach\text{-}algebra}(a) \Rightarrow (\operatorname{vnrm}(a)(1_{a}) = 1))}}
\vnbprop{bal-uba-nag-view}{\vnbstmt{\begin{array}{@{}l@{}} \forall a. \\ \quad \operatorname{is\text{-}unital\text{-}banach\text{-}algebra}(a) \Rightarrow  \\ \quad \quad (\operatorname{uba\text{-}as\text{-}normed\text{-}ag}(a) \simeq \operatorname{cna\text{-}as\text{-}normed\text{-}ag}(\operatorname{uba\text{-}as\text{-}cba}(a))) \end{array}}}
\vnbprop{bal-uba-complete}{\vnbstmt{\forall a.\; (\operatorname{is\text{-}unital\text{-}banach\text{-}algebra}(a) \Rightarrow \operatorname{is\text{-}complete}(\operatorname{ba\text{-}metric\text{-}space}(a)))}}

\vnbbookfile{theorem-library/ba-elem}{theorem-library/ba-elem}
\noindent{\small ba-elem.scm -- wave of 2026-11-09/10 (agent BA-1), wired 2026-11-10 from scratchpad/ba1-1109/ba-elem.scm. Library: operator-algebras.}\par
\vnbprop{bae-act-type}{\vnbstmt{\begin{array}{@{}l@{}} \forall a.; \forall z \in \mathbb{C}, x \in \operatorname{vec}(a). \\ \quad \operatorname{is\text{-}unital\text{-}banach\text{-}algebra}(a) \Rightarrow (\operatorname{act}(a)(z, x) \in \operatorname{vec}(a)) \end{array}}}
\vnbprop{bae-vnrm-type}{\vnbstmt{\forall a.; \forall x \in \operatorname{vec}(a).\; (\operatorname{is\text{-}unital\text{-}banach\text{-}algebra}(a) \Rightarrow (\operatorname{vnrm}(a)(x) \in \mathbb{R}))}}
\vnbprop{bae-inv-unique}{\vnbstmt{\begin{array}{@{}l@{}} \forall a.; \forall x, u, w \in \operatorname{vec}(a). \\ \quad (\operatorname{is\text{-}unital\text{-}banach\text{-}algebra}(a) \wedge ((\cdot_{a}(u, x) = 1_{a}) \wedge (\cdot_{a}(x, w) = 1_{a}))) \Rightarrow  \\ \quad \quad u = w \end{array}}}
\vnbprop{bae-inverse-prop}{\vnbstmt{\begin{array}{@{}l@{}} \forall a, x. \\ \quad (\operatorname{is\text{-}unital\text{-}banach\text{-}algebra}(a) \wedge \operatorname{is\text{-}invertible\text{-}el}(a, x)) \Rightarrow  \\ \quad \quad {x}^{-1} \in \operatorname{vec}(a) \wedge \\ \quad \quad \cdot_{a}({x}^{-1}, x) = 1_{a} \wedge \\ \quad \quad \cdot_{a}(x, {x}^{-1}) = 1_{a} \end{array}}}
\vnbprop{bae-inverse-in}{\vnbstmt{\begin{array}{@{}l@{}} \forall a, x. \\ \quad (\operatorname{is\text{-}unital\text{-}banach\text{-}algebra}(a) \wedge \operatorname{is\text{-}invertible\text{-}el}(a, x)) \Rightarrow  \\ \quad \quad {x}^{-1} \in \operatorname{vec}(a) \end{array}}}
\vnbprop{bae-inverse-left}{\vnbstmt{\begin{array}{@{}l@{}} \forall a, x. \\ \quad (\operatorname{is\text{-}unital\text{-}banach\text{-}algebra}(a) \wedge \operatorname{is\text{-}invertible\text{-}el}(a, x)) \Rightarrow  \\ \quad \quad \cdot_{a}({x}^{-1}, x) = 1_{a} \end{array}}}
\vnbprop{bae-inverse-right}{\vnbstmt{\begin{array}{@{}l@{}} \forall a, x. \\ \quad (\operatorname{is\text{-}unital\text{-}banach\text{-}algebra}(a) \wedge \operatorname{is\text{-}invertible\text{-}el}(a, x)) \Rightarrow  \\ \quad \quad \cdot_{a}(x, {x}^{-1}) = 1_{a} \end{array}}}
\vnbprop{bae-invertible-intro}{\vnbstmt{\begin{array}{@{}l@{}} \forall a.; \forall x, u \in \operatorname{vec}(a). \\ \quad (\operatorname{is\text{-}unital\text{-}banach\text{-}algebra}(a) \wedge ((\cdot_{a}(u, x) = 1_{a}) \wedge (\cdot_{a}(x, u) = 1_{a}))) \Rightarrow  \\ \quad \quad \operatorname{is\text{-}invertible\text{-}el}(a, x) \end{array}}}
\vnbprop{bae-inverse-eq}{\vnbstmt{\begin{array}{@{}l@{}} \forall a.; \forall x, u \in \operatorname{vec}(a). \\ \quad (\operatorname{is\text{-}unital\text{-}banach\text{-}algebra}(a) \wedge ((\cdot_{a}(u, x) = 1_{a}) \wedge (\cdot_{a}(x, u) = 1_{a}))) \Rightarrow  \\ \quad \quad {x}^{-1} = u \end{array}}}
\vnbprop{bae-one-invertible}{\vnbstmt{\forall a.\; (\operatorname{is\text{-}unital\text{-}banach\text{-}algebra}(a) \Rightarrow \operatorname{is\text{-}invertible\text{-}el}(a, 1_{a}))}}
\vnbprop{bae-invertible-mul}{\vnbstmt{\begin{array}{@{}l@{}} \forall a, x, y. \\ \quad (\operatorname{is\text{-}unital\text{-}banach\text{-}algebra}(a) \wedge (\operatorname{is\text{-}invertible\text{-}el}(a, x) \wedge \operatorname{is\text{-}invertible\text{-}el}(a, y))) \Rightarrow  \\ \quad \quad \operatorname{is\text{-}invertible\text{-}el}(a, \cdot_{a}(x, y)) \end{array}}}
\vnbprop{bae-el-pow-type}{\vnbstmt{\begin{array}{@{}l@{}} \forall n \in \mathbb{N}, a.; \forall x \in \operatorname{vec}(a). \\ \quad \operatorname{is\text{-}unital\text{-}banach\text{-}algebra}(a) \Rightarrow ({x}^{n} \in \operatorname{vec}(a)) \end{array}}}
\vnbprop{bae-el-pow-one}{\vnbstmt{\begin{array}{@{}l@{}} \forall a.; \forall x \in \operatorname{vec}(a). \\ \quad \operatorname{is\text{-}unital\text{-}banach\text{-}algebra}(a) \Rightarrow ({x}^{1} = x) \end{array}}}
\vnbprop{bae-el-pow-comm}{\vnbstmt{\begin{array}{@{}l@{}} \forall n \in \mathbb{N}, a.; \forall x \in \operatorname{vec}(a). \\ \quad \operatorname{is\text{-}unital\text{-}banach\text{-}algebra}(a) \Rightarrow  \\ \quad \quad \cdot_{a}(x, {x}^{n}) = \cdot_{a}({x}^{n}, x) \end{array}}}
\vnbprop{bae-el-pow-succ-right}{\vnbstmt{\begin{array}{@{}l@{}} \forall n \in \mathbb{N}, a.; \forall x \in \operatorname{vec}(a). \\ \quad \operatorname{is\text{-}unital\text{-}banach\text{-}algebra}(a) \Rightarrow  \\ \quad \quad {x}^{(n + 1)} = \cdot_{a}({x}^{n}, x) \end{array}}}
\vnbprop{bae-el-pow-norm}{\vnbstmt{\begin{array}{@{}l@{}} \forall n \in \mathbb{N}, a.; \forall x \in \operatorname{vec}(a). \\ \quad \operatorname{is\text{-}unital\text{-}banach\text{-}algebra}(a) \Rightarrow  \\ \quad \quad \operatorname{vnrm}(a)({x}^{n}) \leq \operatorname{power}(\operatorname{vnrm}(a)(x), n) \end{array}}}

\vnbbookfile{theorem-library/ba-unitization}{theorem-library/ba-unitization}
\noindent{\small ba-unitization.scm -- wave of 2026-11-09/10 (agent BA-1C), wired 2026-11-10 from scratchpad/ba1c-1109/ba-unitization.scm. Library: operator-algebras.}\par
\vnbprop{bun-vadd-apply}{\vnbstmt{\begin{array}{@{}l@{}} \forall a.; \forall l \in \mathbb{C}, x \in \operatorname{vec}(a), r \in \mathbb{C}, y \in \operatorname{vec}(a). \\ \quad \operatorname{is\text{-}complex\text{-}normed\text{-}algebra}(a) \Rightarrow  \\ \quad \quad \operatorname{unit\text{-}vadd}(a)({[}l, x], {[}r, y]) = {[}(l + r), \operatorname{vadd}(a)(x, y)] \end{array}}}
\vnbprop{bun-vneg-apply}{\vnbstmt{\begin{array}{@{}l@{}} \forall a.; \forall l \in \mathbb{C}, x \in \operatorname{vec}(a). \\ \quad \operatorname{is\text{-}complex\text{-}normed\text{-}algebra}(a) \Rightarrow (\operatorname{unit\text{-}vneg}(a)({[}l, x]) = {[}-l, \operatorname{vneg}(a)(x)]) \end{array}}}
\vnbprop{bun-act-apply}{\vnbstmt{\begin{array}{@{}l@{}} \forall a.; \forall c, l \in \mathbb{C}, x \in \operatorname{vec}(a). \\ \quad \operatorname{is\text{-}complex\text{-}normed\text{-}algebra}(a) \Rightarrow  \\ \quad \quad \operatorname{unit\text{-}act}(a)(c, {[}l, x]) = {[}(c \cdot l), \operatorname{act}(a)(c, x)] \end{array}}}
\vnbprop{bun-vnrm-apply}{\vnbstmt{\begin{array}{@{}l@{}} \forall a.; \forall l \in \mathbb{C}, x \in \operatorname{vec}(a). \\ \quad \operatorname{is\text{-}complex\text{-}normed\text{-}algebra}(a) \Rightarrow  \\ \quad \quad \operatorname{unit\text{-}vnrm}(a)({[}l, x]) = (\operatorname{magnitude}(l) + \operatorname{vnrm}(a)(x)) \end{array}}}
\vnbprop{bun-vmul-apply}{\vnbstmt{\begin{array}{@{}l@{}} \forall a.; \forall l \in \mathbb{C}, x \in \operatorname{vec}(a), r \in \mathbb{C}, y \in \operatorname{vec}(a). \\ \quad \operatorname{is\text{-}complex\text{-}normed\text{-}algebra}(a) \Rightarrow  \\ \quad \quad \operatorname{unit\text{-}vmul}(a)({[}l, x], {[}r, y]) \\ \quad \quad =\; {[}(l \cdot r), \operatorname{vadd}(a)(\operatorname{vadd}(a)(\operatorname{act}(a)(l, y), \operatorname{act}(a)(r, x)), \cdot_{a}(x, y))] \end{array}}}
\vnbprop{unitization-norm-unit}{\vnbstmt{\begin{array}{@{}l@{}} \forall a. \\ \quad \operatorname{is\text{-}complex\text{-}normed\text{-}algebra}(a) \Rightarrow (\operatorname{unit\text{-}vnrm}(a)({[}1, \operatorname{vzero}(a)]) = 1) \end{array}}}
\vnbprop{unitization-norm-triangle}{\vnbstmt{\begin{array}{@{}l@{}} \forall a.; \forall l \in \mathbb{C}, x \in \operatorname{vec}(a), r \in \mathbb{C}, y \in \operatorname{vec}(a). \\ \quad \operatorname{is\text{-}complex\text{-}normed\text{-}algebra}(a) \Rightarrow  \\ \quad \quad \operatorname{unit\text{-}vnrm}(a)(\operatorname{unit\text{-}vadd}(a)({[}l, x], {[}r, y])) \leq (\operatorname{unit\text{-}vnrm}(a)({[}l, x]) + \operatorname{unit\text{-}vnrm}(a)({[}r, y])) \end{array}}}
\vnbprop{unitization-norm-submult}{\vnbstmt{\begin{array}{@{}l@{}} \forall a.; \forall l \in \mathbb{C}, x \in \operatorname{vec}(a), r \in \mathbb{C}, y \in \operatorname{vec}(a). \\ \quad \operatorname{is\text{-}complex\text{-}normed\text{-}algebra}(a) \Rightarrow  \\ \quad \quad \operatorname{unit\text{-}vnrm}(a)(\operatorname{unit\text{-}vmul}(a)({[}l, x], {[}r, y])) \leq (\operatorname{unit\text{-}vnrm}(a)({[}l, x]) \cdot \operatorname{unit\text{-}vnrm}(a)({[}r, y])) \end{array}}}
\vnbprop{unitization-norm-homog}{\vnbstmt{\begin{array}{@{}l@{}} \forall a.; \forall c, l \in \mathbb{C}, x \in \operatorname{vec}(a). \\ \quad \operatorname{is\text{-}complex\text{-}normed\text{-}algebra}(a) \Rightarrow  \\ \quad \quad \operatorname{unit\text{-}vnrm}(a)(\operatorname{unit\text{-}act}(a)(c, {[}l, x])) = (\operatorname{magnitude}(c) \cdot \operatorname{unit\text{-}vnrm}(a)({[}l, x])) \end{array}}}
\vnbprop{bun-vnrm-zero}{\vnbstmt{\forall a.\; (\operatorname{is\text{-}complex\text{-}normed\text{-}algebra}(a) \Rightarrow (\operatorname{vnrm}(a)(\operatorname{vzero}(a)) = 0))}}
\vnbprop{bun-vmul-zero-left}{\vnbstmt{\begin{array}{@{}l@{}} \forall a.; \forall y \in \operatorname{vec}(a). \\ \quad \operatorname{is\text{-}complex\text{-}normed\text{-}algebra}(a) \Rightarrow (\cdot_{a}(\operatorname{vzero}(a), y) = \operatorname{vzero}(a)) \end{array}}}
\vnbprop{bun-vmul-zero-right}{\vnbstmt{\begin{array}{@{}l@{}} \forall a.; \forall y \in \operatorname{vec}(a). \\ \quad \operatorname{is\text{-}complex\text{-}normed\text{-}algebra}(a) \Rightarrow (\cdot_{a}(y, \operatorname{vzero}(a)) = \operatorname{vzero}(a)) \end{array}}}
\vnbprop{bun-vadd-zero-right}{\vnbstmt{\begin{array}{@{}l@{}} \forall a.; \forall y \in \operatorname{vec}(a). \\ \quad \operatorname{is\text{-}complex\text{-}normed\text{-}algebra}(a) \Rightarrow (\operatorname{vadd}(a)(y, \operatorname{vzero}(a)) = y) \end{array}}}
\vnbprop{bun-vadd-zero-left}{\vnbstmt{\begin{array}{@{}l@{}} \forall a.; \forall y \in \operatorname{vec}(a). \\ \quad \operatorname{is\text{-}complex\text{-}normed\text{-}algebra}(a) \Rightarrow (\operatorname{vadd}(a)(\operatorname{vzero}(a), y) = y) \end{array}}}
\vnbprop{bun-act-one}{\vnbstmt{\begin{array}{@{}l@{}} \forall a.; \forall y \in \operatorname{vec}(a). \\ \quad \operatorname{is\text{-}complex\text{-}normed\text{-}algebra}(a) \Rightarrow (\operatorname{act}(a)(1, y) = y) \end{array}}}
\vnbprop{bun-act-zero-vec}{\vnbstmt{\begin{array}{@{}l@{}} \forall a.; \forall r \in \mathbb{C}. \\ \quad \operatorname{is\text{-}complex\text{-}normed\text{-}algebra}(a) \Rightarrow (\operatorname{act}(a)(r, \operatorname{vzero}(a)) = \operatorname{vzero}(a)) \end{array}}}
\vnbprop{bun-unit-left-pair}{\vnbstmt{\begin{array}{@{}l@{}} \forall a.; \forall r \in \mathbb{C}, y \in \operatorname{vec}(a). \\ \quad \operatorname{is\text{-}complex\text{-}normed\text{-}algebra}(a) \Rightarrow  \\ \quad \quad \operatorname{unit\text{-}vmul}(a)({[}1, \operatorname{vzero}(a)], {[}r, y]) = {[}r, y] \end{array}}}
\vnbprop{bun-unit-right-pair}{\vnbstmt{\begin{array}{@{}l@{}} \forall a.; \forall r \in \mathbb{C}, y \in \operatorname{vec}(a). \\ \quad \operatorname{is\text{-}complex\text{-}normed\text{-}algebra}(a) \Rightarrow  \\ \quad \quad \operatorname{unit\text{-}vmul}(a)({[}r, y], {[}1, \operatorname{vzero}(a)]) = {[}r, y] \end{array}}}
\vnbprop{unit-vec-unfold}{\vnbstmt{\forall v.\; (\operatorname{unit\text{-}vec}(v) \simeq (\mathbb{C} \times \operatorname{vec}(v)))}}
\vnbprop{unitization-unit-left}{\vnbstmt{\begin{array}{@{}l@{}} \forall a.; \forall u \in \operatorname{unit\text{-}vec}(a). \\ \quad \operatorname{is\text{-}complex\text{-}normed\text{-}algebra}(a) \Rightarrow (\operatorname{unit\text{-}vmul}(a)({[}1, \operatorname{vzero}(a)], u) = u) \end{array}}}
\vnbprop{unitization-unit-right}{\vnbstmt{\begin{array}{@{}l@{}} \forall a.; \forall u \in \operatorname{unit\text{-}vec}(a). \\ \quad \operatorname{is\text{-}complex\text{-}normed\text{-}algebra}(a) \Rightarrow (\operatorname{unit\text{-}vmul}(a)(u, {[}1, \operatorname{vzero}(a)]) = u) \end{array}}}

\vnbbookfile{structure-library/ba-quotient-vocab}{structure-library/ba-quotient-vocab}
\noindent{\small ba-quotient-vocab.scm -- wave of 2026-11-09/10 (agent BA-1C), wired 2026-11-10 from scratchpad/ba1c-1109/ba-quotient-vocab.scm. Library: operator-algebras.}\par
\par\medskip\noindent\textbf{Definition} (\texttt{quot-norm}).\;$\operatorname{quot\text{-}norm}(\operatorname{bqqa\_}, \operatorname{bqqi\_}, \operatorname{bqqx\_}) :=$
\begin{equation*}
\operatorname{einf}(\{\operatorname{bqqw\_} \in \mathbb{R} : \exists \operatorname{bqqy\_} \in \operatorname{bqqi\_}.\; (\operatorname{bqqw\_} = \operatorname{vnrm}(\operatorname{bqqa\_})(\operatorname{vadd}(\operatorname{bqqa\_})(\operatorname{bqqx\_}, \operatorname{vneg}(\operatorname{bqqa\_})(\operatorname{bqqy\_}))))\})
\end{equation*}


\vnbbookfile{theorem-library/ba-maximal-ideal}{theorem-library/ba-maximal-ideal}
\noindent{\small ba-maximal-ideal.scm -- wave of 2026-11-09/10 (agent BA-1D), wired 2026-11-10 from scratchpad/ba1d-1109/ba-maximal-ideal.scm. Library: operator-algebras.}\par
\vnbprop{bmx-subset-antisym}{\vnbstmt{\forall p, q.\; (((p \subseteq q) \wedge (q \subseteq p)) \Rightarrow (p = q))}}
\vnbprop{bmx-ideal-parts}{\vnbstmt{\begin{array}{@{}l@{}} \forall a, q. \\ \quad \operatorname{is\text{-}two\text{-}sided\text{-}ideal}(a, q) \Rightarrow  \\ \quad \quad q \subseteq \operatorname{vec}(a) \wedge \\ \quad \quad \operatorname{vzero}(a) \in q \wedge \\ \quad \quad \forall x, y \in q.\; (\operatorname{vadd}(a)(x, y) \in q) \wedge \\ \quad \quad \forall z \in \mathbb{C}, x \in q.\; (\operatorname{act}(a)(z, x) \in q) \wedge \\ \quad \quad \forall y \in \operatorname{vec}(a), x \in q.\; (\cdot_{a}(y, x) \in q) \wedge \\ \quad \quad \forall y \in \operatorname{vec}(a), x \in q.\; (\cdot_{a}(x, y) \in q) \end{array}}}
\vnbprop{bmx-ideal-intro}{\vnbstmt{\begin{array}{@{}l@{}} \forall a, q. \\ \quad q \subseteq \operatorname{vec}(a) \wedge \\ \quad \operatorname{vzero}(a) \in q \wedge \\ \quad \forall x, y \in q.\; (\operatorname{vadd}(a)(x, y) \in q) \wedge \\ \quad \forall z \in \mathbb{C}, x \in q.\; (\operatorname{act}(a)(z, x) \in q) \wedge \\ \quad \forall y \in \operatorname{vec}(a), x \in q.\; (\cdot_{a}(y, x) \in q) \wedge \\ \quad \forall y \in \operatorname{vec}(a), x \in q.\; (\cdot_{a}(x, y) \in q) \Rightarrow \\ \quad \quad \operatorname{is\text{-}two\text{-}sided\text{-}ideal}(a, q) \end{array}}}
\vnbprop{bmx-one-proper}{\vnbstmt{\begin{array}{@{}l@{}} \forall a, q. \\ \quad (\operatorname{is\text{-}unital\text{-}banach\text{-}algebra}(a) \wedge (\operatorname{is\text{-}two\text{-}sided\text{-}ideal}(a, q) \wedge (1_{a} \in q))) \Rightarrow  \\ \quad \quad q = \operatorname{vec}(a) \end{array}}}
\vnbprop{bmx-grd-in}{\vnbstmt{\begin{array}{@{}l@{}} \forall a, j, q. \\ \quad \operatorname{is\text{-}unital\text{-}banach\text{-}algebra}(a) \wedge \\ \quad \operatorname{is\text{-}two\text{-}sided\text{-}ideal}(a, j) \wedge \\ \quad \operatorname{is\text{-}two\text{-}sided\text{-}ideal}(a, q) \wedge \\ \quad j \subseteq q \wedge \\ \quad \neg (1_{a} \in q) \Rightarrow \\ \quad \quad q \in \{i \in \operatorname{power}(\operatorname{vec}(a)) : (\operatorname{is\text{-}two\text{-}sided\text{-}ideal}(a, i) \wedge ((j \subseteq i) \wedge \neg (1_{a} \in i)))\} \end{array}}}
\vnbprop{bmx-grd-parts}{\vnbstmt{\begin{array}{@{}l@{}} \forall a, j.; \forall q \in \{i \in \operatorname{power}(\operatorname{vec}(a)) : (\operatorname{is\text{-}two\text{-}sided\text{-}ideal}(a, i) \wedge ((j \subseteq i) \wedge \neg (1_{a} \in i)))\}. \\ \quad (\operatorname{is\text{-}unital\text{-}banach\text{-}algebra}(a) \wedge \operatorname{is\text{-}two\text{-}sided\text{-}ideal}(a, j)) \Rightarrow  \\ \quad \quad \operatorname{is\text{-}two\text{-}sided\text{-}ideal}(a, q) \wedge ((j \subseteq q) \wedge \neg (1_{a} \in q)) \end{array}}}
\vnbprop{bmx-porel-in}{\vnbstmt{\begin{array}{@{}l@{}} \forall a, j, p, q. \\ \quad \operatorname{is\text{-}unital\text{-}banach\text{-}algebra}(a) \wedge \\ \quad \operatorname{is\text{-}two\text{-}sided\text{-}ideal}(a, j) \wedge \\ \quad p \in \{i \in \operatorname{power}(\operatorname{vec}(a)) : (\operatorname{is\text{-}two\text{-}sided\text{-}ideal}(a, i) \wedge ((j \subseteq i) \wedge \neg (1_{a} \in i)))\} \wedge \\ \quad q \in \{i \in \operatorname{power}(\operatorname{vec}(a)) : (\operatorname{is\text{-}two\text{-}sided\text{-}ideal}(a, i) \wedge ((j \subseteq i) \wedge \neg (1_{a} \in i)))\} \wedge \\ \quad p \subseteq q \Rightarrow \\ \quad \quad {[}p, q] \\ \quad \quad \in\; \{w \in (\{i \in \operatorname{power}(\operatorname{vec}(a)) : (\operatorname{is\text{-}two\text{-}sided\text{-}ideal}(a, i) \wedge ((j \subseteq i) \wedge \neg (1_{a} \in i)))\} \times \{i \in \operatorname{power}(\operatorname{vec}(a)) : (\operatorname{is\text{-}two\text{-}sided\text{-}ideal}(a, i) \wedge ((j \subseteq i) \wedge \neg (1_{a} \in i)))\}) : (\operatorname{nth}(1, w) \subseteq \operatorname{nth}(2, w))\} \end{array}}}
\vnbprop{bmx-porel-parts}{\vnbstmt{\begin{array}{@{}l@{}} \forall a, j, p, q. \\ \quad \operatorname{is\text{-}unital\text{-}banach\text{-}algebra}(a) \wedge \\ \quad \operatorname{is\text{-}two\text{-}sided\text{-}ideal}(a, j) \wedge \\ \quad {[}p, q] \\ \quad \in\; \{w \in (\{i \in \operatorname{power}(\operatorname{vec}(a)) : (\operatorname{is\text{-}two\text{-}sided\text{-}ideal}(a, i) \wedge ((j \subseteq i) \wedge \neg (1_{a} \in i)))\} \times \{i \in \operatorname{power}(\operatorname{vec}(a)) : (\operatorname{is\text{-}two\text{-}sided\text{-}ideal}(a, i) \wedge ((j \subseteq i) \wedge \neg (1_{a} \in i)))\}) : (\operatorname{nth}(1, w) \subseteq \operatorname{nth}(2, w))\} \Rightarrow \\ \quad \quad p \subseteq q \end{array}}}
\vnbprop{maximal-ideal-exists}{\vnbstmt{\begin{array}{@{}l@{}} \forall a, j. \\ \quad (\operatorname{is\text{-}unital\text{-}banach\text{-}algebra}(a) \wedge (\operatorname{is\text{-}two\text{-}sided\text{-}ideal}(a, j) \wedge \neg (j = \operatorname{vec}(a)))) \Rightarrow  \\ \quad \quad \exists m.\; (\operatorname{is\text{-}maximal\text{-}ideal}(a, m) \wedge (j \subseteq m)) \end{array}}}

\vnbbookfile{structure-library/ba-quotient-na-vocab}{structure-library/ba-quotient-na-vocab}
\noindent{\small ba-quotient-na-vocab.scm -- wave of 2026-11-09/10 (agent BA-1E), wired 2026-11-10 from scratchpad/ba1e-1109/ba-quotient-na-vocab.scm. Library: operator-algebras.}\par
\par\medskip\noindent\textbf{Definition} (\texttt{coset-norm}).\;$\operatorname{coset\text{-}norm}(\operatorname{bqna\_}, \operatorname{bqni\_}, \operatorname{bqnc\_}) :=$
\begin{equation*}
\operatorname{einf}(\{\operatorname{bqnw\_} \in \mathbb{R} : \exists \operatorname{bqnu\_} \in \operatorname{bqnc\_}.\; (\operatorname{bqnw\_} = \operatorname{vnrm}(\operatorname{bqna\_})(\operatorname{bqnu\_}))\})
\end{equation*}

\par\medskip\noindent\textbf{Definition} (\texttt{quotient-ba}).\;${\operatorname{bqna\_}}/{\operatorname{bqni\_}} :=$
\begin{equation*}
{[}\operatorname{scal}(\operatorname{quotient\text{-}module}(\operatorname{uba\text{-}as\text{-}module}(\operatorname{bqna\_}), \operatorname{bqni\_})), \operatorname{vec}(\operatorname{quotient\text{-}module}(\operatorname{uba\text{-}as\text{-}module}(\operatorname{bqna\_}), \operatorname{bqni\_})), \operatorname{vadd}(\operatorname{quotient\text{-}module}(\operatorname{uba\text{-}as\text{-}module}(\operatorname{bqna\_}), \operatorname{bqni\_})), \operatorname{vzero}(\operatorname{quotient\text{-}module}(\operatorname{uba\text{-}as\text{-}module}(\operatorname{bqna\_}), \operatorname{bqni\_})), \operatorname{vneg}(\operatorname{quotient\text{-}module}(\operatorname{uba\text{-}as\text{-}module}(\operatorname{bqna\_}), \operatorname{bqni\_})), \operatorname{act}(\operatorname{quotient\text{-}module}(\operatorname{uba\text{-}as\text{-}module}(\operatorname{bqna\_}), \operatorname{bqni\_})), \operatorname{qna\text{-}vnrm}(\operatorname{bqna\_}, \operatorname{bqni\_}), \operatorname{qna\text{-}vmul}(\operatorname{bqna\_}, \operatorname{bqni\_}), \operatorname{class}(\operatorname{qmod\text{-}setoid}(\operatorname{uba\text{-}as\text{-}module}(\operatorname{bqna\_}), \operatorname{bqni\_}), 1_{\operatorname{bqna\_}})]
\end{equation*}

\par\medskip\noindent\textbf{Definition} (\texttt{qna-vnrm}).\;$\operatorname{qna\text{-}vnrm}(\operatorname{bqna\_}, \operatorname{bqni\_}) :=$
\begin{equation*}
\begin{array}{@{}l@{}} (\lambda \operatorname{bqnc\_} \in \operatorname{quotient}(\operatorname{qmod\text{-}setoid}(\operatorname{uba\text{-}as\text{-}module}(\operatorname{bqna\_}), \operatorname{bqni\_})).\; \\ \quad \operatorname{coset\text{-}norm}(\operatorname{bqna\_}, \operatorname{bqni\_}, \operatorname{bqnc\_})) \end{array}
\end{equation*}

\par\medskip\noindent\textbf{Definition} (\texttt{qna-vmul}).\;$\operatorname{qna\text{-}vmul}(\operatorname{bqna\_}, \operatorname{bqni\_}) :=$
\begin{equation*}
\begin{array}{@{}l@{}} \operatorname{descend2}(\operatorname{qmod\text{-}setoid}(\operatorname{uba\text{-}as\text{-}module}(\operatorname{bqna\_}), \operatorname{bqni\_}), \\ \quad (\lambda {[}\operatorname{bqnx\_}, \operatorname{bqny\_}] \in (\operatorname{vec}(\operatorname{bqna\_}) \times \operatorname{vec}(\operatorname{bqna\_})).\; \\ \quad \quad \operatorname{class}(\operatorname{qmod\text{-}setoid}(\operatorname{uba\text{-}as\text{-}module}(\operatorname{bqna\_}), \operatorname{bqni\_}), \cdot_{\operatorname{bqna\_}}(\operatorname{bqnx\_}, \operatorname{bqny\_})))) \end{array}
\end{equation*}

\par\medskip\noindent\textbf{Definition} (\texttt{quotient-na}).\;${\operatorname{bqna\_}}/{\operatorname{bqni\_}} :=$
\begin{equation*}
{[}\operatorname{scal}(\operatorname{quotient\text{-}module}(\operatorname{uba\text{-}as\text{-}module}(\operatorname{bqna\_}), \operatorname{bqni\_})), \operatorname{vec}(\operatorname{quotient\text{-}module}(\operatorname{uba\text{-}as\text{-}module}(\operatorname{bqna\_}), \operatorname{bqni\_})), \operatorname{vadd}(\operatorname{quotient\text{-}module}(\operatorname{uba\text{-}as\text{-}module}(\operatorname{bqna\_}), \operatorname{bqni\_})), \operatorname{vzero}(\operatorname{quotient\text{-}module}(\operatorname{uba\text{-}as\text{-}module}(\operatorname{bqna\_}), \operatorname{bqni\_})), \operatorname{vneg}(\operatorname{quotient\text{-}module}(\operatorname{uba\text{-}as\text{-}module}(\operatorname{bqna\_}), \operatorname{bqni\_})), \operatorname{act}(\operatorname{quotient\text{-}module}(\operatorname{uba\text{-}as\text{-}module}(\operatorname{bqna\_}), \operatorname{bqni\_})), \operatorname{qna\text{-}vnrm}(\operatorname{bqna\_}, \operatorname{bqni\_}), \operatorname{qna\text{-}vmul}(\operatorname{bqna\_}, \operatorname{bqni\_})]
\end{equation*}


\vnbbookfile{theorem-library/kal-laws}{theorem-library/kal-laws}
\noindent{\small kal-laws.scm -- wave of 2026-11-09/10 (agent KIT-ALG), wired 2026-11-10 from scratchpad/kitalg-1109/kal-laws.scm. Library: operator-algebras.}\par
\vnbprop{kal-cna-vadd-in}{\vnbstmt{\begin{array}{@{}l@{}} \forall a.; \forall x, y \in \operatorname{vec}(a). \\ \quad \operatorname{is\text{-}complex\text{-}normed\text{-}algebra}(a) \Rightarrow (\operatorname{vadd}(a)(x, y) \in \operatorname{vec}(a)) \end{array}}}
\vnbprop{kal-cna-vneg-in}{\vnbstmt{\begin{array}{@{}l@{}} \forall a.; \forall x \in \operatorname{vec}(a). \\ \quad \operatorname{is\text{-}complex\text{-}normed\text{-}algebra}(a) \Rightarrow (\operatorname{vneg}(a)(x) \in \operatorname{vec}(a)) \end{array}}}
\vnbprop{kal-cna-act-in}{\vnbstmt{\begin{array}{@{}l@{}} \forall a.; \forall r \in \mathbb{C}, x \in \operatorname{vec}(a). \\ \quad \operatorname{is\text{-}complex\text{-}normed\text{-}algebra}(a) \Rightarrow (\operatorname{act}(a)(r, x) \in \operatorname{vec}(a)) \end{array}}}
\vnbprop{kal-cna-vzero-in}{\vnbstmt{\forall a.\; (\operatorname{is\text{-}complex\text{-}normed\text{-}algebra}(a) \Rightarrow (\operatorname{vzero}(a) \in \operatorname{vec}(a)))}}
\vnbprop{kal-cna-vmul-in}{\vnbstmt{\begin{array}{@{}l@{}} \forall a.; \forall x, y \in \operatorname{vec}(a). \\ \quad \operatorname{is\text{-}complex\text{-}normed\text{-}algebra}(a) \Rightarrow (\cdot_{a}(x, y) \in \operatorname{vec}(a)) \end{array}}}
\vnbprop{kal-cna-add-assoc}{\vnbstmt{\begin{array}{@{}l@{}} \forall a.; \forall x, y, z \in \operatorname{vec}(a). \\ \quad \operatorname{is\text{-}complex\text{-}normed\text{-}algebra}(a) \Rightarrow  \\ \quad \quad \operatorname{vadd}(a)(\operatorname{vadd}(a)(x, y), z) = \operatorname{vadd}(a)(x, \operatorname{vadd}(a)(y, z)) \end{array}}}
\vnbprop{kal-cna-add-comm}{\vnbstmt{\begin{array}{@{}l@{}} \forall a.; \forall x, y \in \operatorname{vec}(a). \\ \quad \operatorname{is\text{-}complex\text{-}normed\text{-}algebra}(a) \Rightarrow (\operatorname{vadd}(a)(x, y) = \operatorname{vadd}(a)(y, x)) \end{array}}}
\vnbprop{kal-cna-add-zero-l}{\vnbstmt{\begin{array}{@{}l@{}} \forall a.; \forall x \in \operatorname{vec}(a). \\ \quad \operatorname{is\text{-}complex\text{-}normed\text{-}algebra}(a) \Rightarrow (\operatorname{vadd}(a)(\operatorname{vzero}(a), x) = x) \end{array}}}
\vnbprop{kal-cna-add-zero-r}{\vnbstmt{\begin{array}{@{}l@{}} \forall a.; \forall x \in \operatorname{vec}(a). \\ \quad \operatorname{is\text{-}complex\text{-}normed\text{-}algebra}(a) \Rightarrow (\operatorname{vadd}(a)(x, \operatorname{vzero}(a)) = x) \end{array}}}
\vnbprop{kal-cna-act-vadd}{\vnbstmt{\begin{array}{@{}l@{}} \forall a.; \forall r \in \mathbb{C}, x, y \in \operatorname{vec}(a). \\ \quad \operatorname{is\text{-}complex\text{-}normed\text{-}algebra}(a) \Rightarrow  \\ \quad \quad \operatorname{act}(a)(r, \operatorname{vadd}(a)(x, y)) = \operatorname{vadd}(a)(\operatorname{act}(a)(r, x), \operatorname{act}(a)(r, y)) \end{array}}}
\vnbprop{kal-cna-act-add}{\vnbstmt{\begin{array}{@{}l@{}} \forall a.; \forall r, s \in \mathbb{C}, x \in \operatorname{vec}(a). \\ \quad \operatorname{is\text{-}complex\text{-}normed\text{-}algebra}(a) \Rightarrow  \\ \quad \quad \operatorname{act}(a)((r + s), x) = \operatorname{vadd}(a)(\operatorname{act}(a)(r, x), \operatorname{act}(a)(s, x)) \end{array}}}
\vnbprop{kal-cna-act-mul}{\vnbstmt{\begin{array}{@{}l@{}} \forall a.; \forall r, s \in \mathbb{C}, x \in \operatorname{vec}(a). \\ \quad \operatorname{is\text{-}complex\text{-}normed\text{-}algebra}(a) \Rightarrow  \\ \quad \quad \operatorname{act}(a)((r \cdot s), x) = \operatorname{act}(a)(r, \operatorname{act}(a)(s, x)) \end{array}}}
\vnbprop{kal-cna-act-one}{\vnbstmt{\begin{array}{@{}l@{}} \forall a.; \forall x \in \operatorname{vec}(a). \\ \quad \operatorname{is\text{-}complex\text{-}normed\text{-}algebra}(a) \Rightarrow (\operatorname{act}(a)(1, x) = x) \end{array}}}
\vnbprop{kal-cna-act-zero}{\vnbstmt{\begin{array}{@{}l@{}} \forall a.; \forall x \in \operatorname{vec}(a). \\ \quad \operatorname{is\text{-}complex\text{-}normed\text{-}algebra}(a) \Rightarrow (\operatorname{act}(a)(0, x) = \operatorname{vzero}(a)) \end{array}}}
\vnbprop{kal-cna-act-vzero}{\vnbstmt{\begin{array}{@{}l@{}} \forall a.; \forall r \in \mathbb{C}. \\ \quad \operatorname{is\text{-}complex\text{-}normed\text{-}algebra}(a) \Rightarrow (\operatorname{act}(a)(r, \operatorname{vzero}(a)) = \operatorname{vzero}(a)) \end{array}}}
\vnbprop{kal-cna-vneg-act}{\vnbstmt{\begin{array}{@{}l@{}} \forall a.; \forall x \in \operatorname{vec}(a). \\ \quad \operatorname{is\text{-}complex\text{-}normed\text{-}algebra}(a) \Rightarrow (\operatorname{vneg}(a)(x) = \operatorname{act}(a)(-1, x)) \end{array}}}
\vnbprop{kal-cna-mul-assoc}{\vnbstmt{\begin{array}{@{}l@{}} \forall a.; \forall x, y, z \in \operatorname{vec}(a). \\ \quad \operatorname{is\text{-}complex\text{-}normed\text{-}algebra}(a) \Rightarrow  \\ \quad \quad \cdot_{a}(\cdot_{a}(x, y), z) = \cdot_{a}(x, \cdot_{a}(y, z)) \end{array}}}
\vnbprop{kal-cna-mul-dist-l}{\vnbstmt{\begin{array}{@{}l@{}} \forall a.; \forall x, y, z \in \operatorname{vec}(a). \\ \quad \operatorname{is\text{-}complex\text{-}normed\text{-}algebra}(a) \Rightarrow  \\ \quad \quad \cdot_{a}(x, \operatorname{vadd}(a)(y, z)) = \operatorname{vadd}(a)(\cdot_{a}(x, y), \cdot_{a}(x, z)) \end{array}}}
\vnbprop{kal-cna-mul-dist-r}{\vnbstmt{\begin{array}{@{}l@{}} \forall a.; \forall x, y, z \in \operatorname{vec}(a). \\ \quad \operatorname{is\text{-}complex\text{-}normed\text{-}algebra}(a) \Rightarrow  \\ \quad \quad \cdot_{a}(\operatorname{vadd}(a)(x, y), z) = \operatorname{vadd}(a)(\cdot_{a}(x, z), \cdot_{a}(y, z)) \end{array}}}
\vnbprop{kal-cna-mul-act-l}{\vnbstmt{\begin{array}{@{}l@{}} \forall a.; \forall r \in \mathbb{C}, x, y \in \operatorname{vec}(a). \\ \quad \operatorname{is\text{-}complex\text{-}normed\text{-}algebra}(a) \Rightarrow  \\ \quad \quad \cdot_{a}(\operatorname{act}(a)(r, x), y) = \operatorname{act}(a)(r, \cdot_{a}(x, y)) \end{array}}}
\vnbprop{kal-cna-mul-act-r}{\vnbstmt{\begin{array}{@{}l@{}} \forall a.; \forall r \in \mathbb{C}, x, y \in \operatorname{vec}(a). \\ \quad \operatorname{is\text{-}complex\text{-}normed\text{-}algebra}(a) \Rightarrow  \\ \quad \quad \cdot_{a}(x, \operatorname{act}(a)(r, y)) = \operatorname{act}(a)(r, \cdot_{a}(x, y)) \end{array}}}
\vnbprop{kal-cna-mul-zero-l}{\vnbstmt{\begin{array}{@{}l@{}} \forall a.; \forall x \in \operatorname{vec}(a). \\ \quad \operatorname{is\text{-}complex\text{-}normed\text{-}algebra}(a) \Rightarrow (\cdot_{a}(\operatorname{vzero}(a), x) = \operatorname{vzero}(a)) \end{array}}}
\vnbprop{kal-cna-mul-zero-r}{\vnbstmt{\begin{array}{@{}l@{}} \forall a.; \forall x \in \operatorname{vec}(a). \\ \quad \operatorname{is\text{-}complex\text{-}normed\text{-}algebra}(a) \Rightarrow (\cdot_{a}(x, \operatorname{vzero}(a)) = \operatorname{vzero}(a)) \end{array}}}
\vnbprop{kal-uba-vadd-in}{\vnbstmt{\begin{array}{@{}l@{}} \forall a.; \forall x, y \in \operatorname{vec}(a). \\ \quad \operatorname{is\text{-}unital\text{-}banach\text{-}algebra}(a) \Rightarrow (\operatorname{vadd}(a)(x, y) \in \operatorname{vec}(a)) \end{array}}}
\vnbprop{kal-uba-vneg-in}{\vnbstmt{\begin{array}{@{}l@{}} \forall a.; \forall x \in \operatorname{vec}(a). \\ \quad \operatorname{is\text{-}unital\text{-}banach\text{-}algebra}(a) \Rightarrow (\operatorname{vneg}(a)(x) \in \operatorname{vec}(a)) \end{array}}}
\vnbprop{kal-uba-act-in}{\vnbstmt{\begin{array}{@{}l@{}} \forall a.; \forall r \in \mathbb{C}, x \in \operatorname{vec}(a). \\ \quad \operatorname{is\text{-}unital\text{-}banach\text{-}algebra}(a) \Rightarrow (\operatorname{act}(a)(r, x) \in \operatorname{vec}(a)) \end{array}}}
\vnbprop{kal-uba-vmul-in}{\vnbstmt{\begin{array}{@{}l@{}} \forall a.; \forall x, y \in \operatorname{vec}(a). \\ \quad \operatorname{is\text{-}unital\text{-}banach\text{-}algebra}(a) \Rightarrow (\cdot_{a}(x, y) \in \operatorname{vec}(a)) \end{array}}}
\vnbprop{kal-uba-vzero-in}{\vnbstmt{\forall a.\; (\operatorname{is\text{-}unital\text{-}banach\text{-}algebra}(a) \Rightarrow (\operatorname{vzero}(a) \in \operatorname{vec}(a)))}}
\vnbprop{kal-uba-add-assoc}{\vnbstmt{\begin{array}{@{}l@{}} \forall a.; \forall x, y, z \in \operatorname{vec}(a). \\ \quad \operatorname{is\text{-}unital\text{-}banach\text{-}algebra}(a) \Rightarrow  \\ \quad \quad \operatorname{vadd}(a)(\operatorname{vadd}(a)(x, y), z) = \operatorname{vadd}(a)(x, \operatorname{vadd}(a)(y, z)) \end{array}}}
\vnbprop{kal-uba-add-comm}{\vnbstmt{\begin{array}{@{}l@{}} \forall a.; \forall x, y \in \operatorname{vec}(a). \\ \quad \operatorname{is\text{-}unital\text{-}banach\text{-}algebra}(a) \Rightarrow (\operatorname{vadd}(a)(x, y) = \operatorname{vadd}(a)(y, x)) \end{array}}}
\vnbprop{kal-uba-add-zero-l}{\vnbstmt{\begin{array}{@{}l@{}} \forall a.; \forall x \in \operatorname{vec}(a). \\ \quad \operatorname{is\text{-}unital\text{-}banach\text{-}algebra}(a) \Rightarrow (\operatorname{vadd}(a)(\operatorname{vzero}(a), x) = x) \end{array}}}
\vnbprop{kal-uba-add-zero-r}{\vnbstmt{\begin{array}{@{}l@{}} \forall a.; \forall x \in \operatorname{vec}(a). \\ \quad \operatorname{is\text{-}unital\text{-}banach\text{-}algebra}(a) \Rightarrow (\operatorname{vadd}(a)(x, \operatorname{vzero}(a)) = x) \end{array}}}
\vnbprop{kal-uba-vneg-act}{\vnbstmt{\begin{array}{@{}l@{}} \forall a.; \forall x \in \operatorname{vec}(a). \\ \quad \operatorname{is\text{-}unital\text{-}banach\text{-}algebra}(a) \Rightarrow (\operatorname{vneg}(a)(x) = \operatorname{act}(a)(-1, x)) \end{array}}}
\vnbprop{kal-uba-act-vadd}{\vnbstmt{\begin{array}{@{}l@{}} \forall a.; \forall r \in \mathbb{C}, x, y \in \operatorname{vec}(a). \\ \quad \operatorname{is\text{-}unital\text{-}banach\text{-}algebra}(a) \Rightarrow  \\ \quad \quad \operatorname{act}(a)(r, \operatorname{vadd}(a)(x, y)) = \operatorname{vadd}(a)(\operatorname{act}(a)(r, x), \operatorname{act}(a)(r, y)) \end{array}}}
\vnbprop{kal-uba-act-add}{\vnbstmt{\begin{array}{@{}l@{}} \forall a.; \forall r, s \in \mathbb{C}, x \in \operatorname{vec}(a). \\ \quad \operatorname{is\text{-}unital\text{-}banach\text{-}algebra}(a) \Rightarrow  \\ \quad \quad \operatorname{act}(a)((r + s), x) = \operatorname{vadd}(a)(\operatorname{act}(a)(r, x), \operatorname{act}(a)(s, x)) \end{array}}}
\vnbprop{kal-uba-act-mul}{\vnbstmt{\begin{array}{@{}l@{}} \forall a.; \forall r, s \in \mathbb{C}, x \in \operatorname{vec}(a). \\ \quad \operatorname{is\text{-}unital\text{-}banach\text{-}algebra}(a) \Rightarrow  \\ \quad \quad \operatorname{act}(a)((r \cdot s), x) = \operatorname{act}(a)(r, \operatorname{act}(a)(s, x)) \end{array}}}
\vnbprop{kal-uba-act-one}{\vnbstmt{\forall a.; \forall x \in \operatorname{vec}(a).\; (\operatorname{is\text{-}unital\text{-}banach\text{-}algebra}(a) \Rightarrow (\operatorname{act}(a)(1, x) = x))}}
\vnbprop{kal-uba-act-zero}{\vnbstmt{\begin{array}{@{}l@{}} \forall a.; \forall x \in \operatorname{vec}(a). \\ \quad \operatorname{is\text{-}unital\text{-}banach\text{-}algebra}(a) \Rightarrow (\operatorname{act}(a)(0, x) = \operatorname{vzero}(a)) \end{array}}}
\vnbprop{kal-uba-act-vzero}{\vnbstmt{\begin{array}{@{}l@{}} \forall a.; \forall r \in \mathbb{C}. \\ \quad \operatorname{is\text{-}unital\text{-}banach\text{-}algebra}(a) \Rightarrow (\operatorname{act}(a)(r, \operatorname{vzero}(a)) = \operatorname{vzero}(a)) \end{array}}}
\vnbprop{kal-uba-mul-assoc}{\vnbstmt{\begin{array}{@{}l@{}} \forall a.; \forall x, y, z \in \operatorname{vec}(a). \\ \quad \operatorname{is\text{-}unital\text{-}banach\text{-}algebra}(a) \Rightarrow  \\ \quad \quad \cdot_{a}(\cdot_{a}(x, y), z) = \cdot_{a}(x, \cdot_{a}(y, z)) \end{array}}}
\vnbprop{kal-uba-mul-dist-l}{\vnbstmt{\begin{array}{@{}l@{}} \forall a.; \forall x, y, z \in \operatorname{vec}(a). \\ \quad \operatorname{is\text{-}unital\text{-}banach\text{-}algebra}(a) \Rightarrow  \\ \quad \quad \cdot_{a}(x, \operatorname{vadd}(a)(y, z)) = \operatorname{vadd}(a)(\cdot_{a}(x, y), \cdot_{a}(x, z)) \end{array}}}
\vnbprop{kal-uba-mul-dist-r}{\vnbstmt{\begin{array}{@{}l@{}} \forall a.; \forall x, y, z \in \operatorname{vec}(a). \\ \quad \operatorname{is\text{-}unital\text{-}banach\text{-}algebra}(a) \Rightarrow  \\ \quad \quad \cdot_{a}(\operatorname{vadd}(a)(x, y), z) = \operatorname{vadd}(a)(\cdot_{a}(x, z), \cdot_{a}(y, z)) \end{array}}}
\vnbprop{kal-uba-mul-act-l}{\vnbstmt{\begin{array}{@{}l@{}} \forall a.; \forall r \in \mathbb{C}, x, y \in \operatorname{vec}(a). \\ \quad \operatorname{is\text{-}unital\text{-}banach\text{-}algebra}(a) \Rightarrow  \\ \quad \quad \cdot_{a}(\operatorname{act}(a)(r, x), y) = \operatorname{act}(a)(r, \cdot_{a}(x, y)) \end{array}}}
\vnbprop{kal-uba-mul-act-r}{\vnbstmt{\begin{array}{@{}l@{}} \forall a.; \forall r \in \mathbb{C}, x, y \in \operatorname{vec}(a). \\ \quad \operatorname{is\text{-}unital\text{-}banach\text{-}algebra}(a) \Rightarrow  \\ \quad \quad \cdot_{a}(x, \operatorname{act}(a)(r, y)) = \operatorname{act}(a)(r, \cdot_{a}(x, y)) \end{array}}}
\vnbprop{kal-uba-mul-zero-l}{\vnbstmt{\begin{array}{@{}l@{}} \forall a.; \forall x \in \operatorname{vec}(a). \\ \quad \operatorname{is\text{-}unital\text{-}banach\text{-}algebra}(a) \Rightarrow (\cdot_{a}(\operatorname{vzero}(a), x) = \operatorname{vzero}(a)) \end{array}}}
\vnbprop{kal-uba-mul-zero-r}{\vnbstmt{\begin{array}{@{}l@{}} \forall a.; \forall x \in \operatorname{vec}(a). \\ \quad \operatorname{is\text{-}unital\text{-}banach\text{-}algebra}(a) \Rightarrow (\cdot_{a}(x, \operatorname{vzero}(a)) = \operatorname{vzero}(a)) \end{array}}}
\vnbprop{kal-uba-vone-in}{\vnbstmt{\forall a.\; (\operatorname{is\text{-}unital\text{-}banach\text{-}algebra}(a) \Rightarrow (1_{a} \in \operatorname{vec}(a)))}}
\vnbprop{kal-uba-mul-one-l}{\vnbstmt{\begin{array}{@{}l@{}} \forall a.; \forall x \in \operatorname{vec}(a). \\ \quad \operatorname{is\text{-}unital\text{-}banach\text{-}algebra}(a) \Rightarrow (\cdot_{a}(1_{a}, x) = x) \end{array}}}
\vnbprop{kal-uba-mul-one-r}{\vnbstmt{\begin{array}{@{}l@{}} \forall a.; \forall x \in \operatorname{vec}(a). \\ \quad \operatorname{is\text{-}unital\text{-}banach\text{-}algebra}(a) \Rightarrow (\cdot_{a}(x, 1_{a}) = x) \end{array}}}
\vnbprop{kal-uu-scal}{\vnbstmt{\begin{array}{@{}l@{}} \forall a. \\ \quad \operatorname{is\text{-}unital\text{-}involutive\text{-}banach\text{-}algebra}(a) \Rightarrow (\operatorname{scal}(\operatorname{uiba\text{-}as\text{-}uba}(a)) \simeq \operatorname{scal}(a)) \end{array}}}
\vnbprop{kal-uu-vec}{\vnbstmt{\begin{array}{@{}l@{}} \forall a. \\ \quad \operatorname{is\text{-}unital\text{-}involutive\text{-}banach\text{-}algebra}(a) \Rightarrow (\operatorname{vec}(\operatorname{uiba\text{-}as\text{-}uba}(a)) \simeq \operatorname{vec}(a)) \end{array}}}
\vnbprop{kal-uu-vadd}{\vnbstmt{\begin{array}{@{}l@{}} \forall a. \\ \quad \operatorname{is\text{-}unital\text{-}involutive\text{-}banach\text{-}algebra}(a) \Rightarrow (\operatorname{vadd}(\operatorname{uiba\text{-}as\text{-}uba}(a)) \simeq \operatorname{vadd}(a)) \end{array}}}
\vnbprop{kal-uu-vzero}{\vnbstmt{\begin{array}{@{}l@{}} \forall a. \\ \quad \operatorname{is\text{-}unital\text{-}involutive\text{-}banach\text{-}algebra}(a) \Rightarrow (\operatorname{vzero}(\operatorname{uiba\text{-}as\text{-}uba}(a)) \simeq \operatorname{vzero}(a)) \end{array}}}
\vnbprop{kal-uu-vneg}{\vnbstmt{\begin{array}{@{}l@{}} \forall a. \\ \quad \operatorname{is\text{-}unital\text{-}involutive\text{-}banach\text{-}algebra}(a) \Rightarrow (\operatorname{vneg}(\operatorname{uiba\text{-}as\text{-}uba}(a)) \simeq \operatorname{vneg}(a)) \end{array}}}
\vnbprop{kal-uu-act}{\vnbstmt{\begin{array}{@{}l@{}} \forall a. \\ \quad \operatorname{is\text{-}unital\text{-}involutive\text{-}banach\text{-}algebra}(a) \Rightarrow (\operatorname{act}(\operatorname{uiba\text{-}as\text{-}uba}(a)) \simeq \operatorname{act}(a)) \end{array}}}
\vnbprop{kal-uu-vnrm}{\vnbstmt{\begin{array}{@{}l@{}} \forall a. \\ \quad \operatorname{is\text{-}unital\text{-}involutive\text{-}banach\text{-}algebra}(a) \Rightarrow (\operatorname{vnrm}(\operatorname{uiba\text{-}as\text{-}uba}(a)) \simeq \operatorname{vnrm}(a)) \end{array}}}
\vnbprop{kal-uu-vmul}{\vnbstmt{\begin{array}{@{}l@{}} \forall a. \\ \quad \operatorname{is\text{-}unital\text{-}involutive\text{-}banach\text{-}algebra}(a) \Rightarrow (\cdot_{\operatorname{uiba\text{-}as\text{-}uba}(a)} \simeq \cdot_{a}) \end{array}}}
\vnbprop{kal-uu-vone}{\vnbstmt{\begin{array}{@{}l@{}} \forall a. \\ \quad \operatorname{is\text{-}unital\text{-}involutive\text{-}banach\text{-}algebra}(a) \Rightarrow (1_{\operatorname{uiba\text{-}as\text{-}uba}(a)} \simeq 1_{a}) \end{array}}}
\vnbprop{kal-uiba-vadd-in}{\vnbstmt{\begin{array}{@{}l@{}} \forall a.; \forall x, y \in \operatorname{vec}(a). \\ \quad \operatorname{is\text{-}unital\text{-}involutive\text{-}banach\text{-}algebra}(a) \Rightarrow (\operatorname{vadd}(a)(x, y) \in \operatorname{vec}(a)) \end{array}}}
\vnbprop{kal-uiba-vneg-in}{\vnbstmt{\begin{array}{@{}l@{}} \forall a.; \forall x \in \operatorname{vec}(a). \\ \quad \operatorname{is\text{-}unital\text{-}involutive\text{-}banach\text{-}algebra}(a) \Rightarrow (\operatorname{vneg}(a)(x) \in \operatorname{vec}(a)) \end{array}}}
\vnbprop{kal-uiba-act-in}{\vnbstmt{\begin{array}{@{}l@{}} \forall a.; \forall r \in \mathbb{C}, x \in \operatorname{vec}(a). \\ \quad \operatorname{is\text{-}unital\text{-}involutive\text{-}banach\text{-}algebra}(a) \Rightarrow (\operatorname{act}(a)(r, x) \in \operatorname{vec}(a)) \end{array}}}
\vnbprop{kal-uiba-vmul-in}{\vnbstmt{\begin{array}{@{}l@{}} \forall a.; \forall x, y \in \operatorname{vec}(a). \\ \quad \operatorname{is\text{-}unital\text{-}involutive\text{-}banach\text{-}algebra}(a) \Rightarrow (\cdot_{a}(x, y) \in \operatorname{vec}(a)) \end{array}}}
\vnbprop{kal-uiba-vzero-in}{\vnbstmt{\forall a.\; (\operatorname{is\text{-}unital\text{-}involutive\text{-}banach\text{-}algebra}(a) \Rightarrow (\operatorname{vzero}(a) \in \operatorname{vec}(a)))}}
\vnbprop{kal-uiba-add-assoc}{\vnbstmt{\begin{array}{@{}l@{}} \forall a.; \forall x, y, z \in \operatorname{vec}(a). \\ \quad \operatorname{is\text{-}unital\text{-}involutive\text{-}banach\text{-}algebra}(a) \Rightarrow  \\ \quad \quad \operatorname{vadd}(a)(\operatorname{vadd}(a)(x, y), z) = \operatorname{vadd}(a)(x, \operatorname{vadd}(a)(y, z)) \end{array}}}
\vnbprop{kal-uiba-add-comm}{\vnbstmt{\begin{array}{@{}l@{}} \forall a.; \forall x, y \in \operatorname{vec}(a). \\ \quad \operatorname{is\text{-}unital\text{-}involutive\text{-}banach\text{-}algebra}(a) \Rightarrow (\operatorname{vadd}(a)(x, y) = \operatorname{vadd}(a)(y, x)) \end{array}}}
\vnbprop{kal-uiba-add-zero-l}{\vnbstmt{\begin{array}{@{}l@{}} \forall a.; \forall x \in \operatorname{vec}(a). \\ \quad \operatorname{is\text{-}unital\text{-}involutive\text{-}banach\text{-}algebra}(a) \Rightarrow (\operatorname{vadd}(a)(\operatorname{vzero}(a), x) = x) \end{array}}}
\vnbprop{kal-uiba-add-zero-r}{\vnbstmt{\begin{array}{@{}l@{}} \forall a.; \forall x \in \operatorname{vec}(a). \\ \quad \operatorname{is\text{-}unital\text{-}involutive\text{-}banach\text{-}algebra}(a) \Rightarrow (\operatorname{vadd}(a)(x, \operatorname{vzero}(a)) = x) \end{array}}}
\vnbprop{kal-uiba-vneg-act}{\vnbstmt{\begin{array}{@{}l@{}} \forall a.; \forall x \in \operatorname{vec}(a). \\ \quad \operatorname{is\text{-}unital\text{-}involutive\text{-}banach\text{-}algebra}(a) \Rightarrow  \\ \quad \quad \operatorname{vneg}(a)(x) = \operatorname{act}(a)(-1, x) \end{array}}}
\vnbprop{kal-uiba-act-vadd}{\vnbstmt{\begin{array}{@{}l@{}} \forall a.; \forall r \in \mathbb{C}, x, y \in \operatorname{vec}(a). \\ \quad \operatorname{is\text{-}unital\text{-}involutive\text{-}banach\text{-}algebra}(a) \Rightarrow  \\ \quad \quad \operatorname{act}(a)(r, \operatorname{vadd}(a)(x, y)) = \operatorname{vadd}(a)(\operatorname{act}(a)(r, x), \operatorname{act}(a)(r, y)) \end{array}}}
\vnbprop{kal-uiba-act-add}{\vnbstmt{\begin{array}{@{}l@{}} \forall a.; \forall r, s \in \mathbb{C}, x \in \operatorname{vec}(a). \\ \quad \operatorname{is\text{-}unital\text{-}involutive\text{-}banach\text{-}algebra}(a) \Rightarrow  \\ \quad \quad \operatorname{act}(a)((r + s), x) = \operatorname{vadd}(a)(\operatorname{act}(a)(r, x), \operatorname{act}(a)(s, x)) \end{array}}}
\vnbprop{kal-uiba-act-mul}{\vnbstmt{\begin{array}{@{}l@{}} \forall a.; \forall r, s \in \mathbb{C}, x \in \operatorname{vec}(a). \\ \quad \operatorname{is\text{-}unital\text{-}involutive\text{-}banach\text{-}algebra}(a) \Rightarrow  \\ \quad \quad \operatorname{act}(a)((r \cdot s), x) = \operatorname{act}(a)(r, \operatorname{act}(a)(s, x)) \end{array}}}
\vnbprop{kal-uiba-act-one}{\vnbstmt{\begin{array}{@{}l@{}} \forall a.; \forall x \in \operatorname{vec}(a). \\ \quad \operatorname{is\text{-}unital\text{-}involutive\text{-}banach\text{-}algebra}(a) \Rightarrow (\operatorname{act}(a)(1, x) = x) \end{array}}}
\vnbprop{kal-uiba-act-zero}{\vnbstmt{\begin{array}{@{}l@{}} \forall a.; \forall x \in \operatorname{vec}(a). \\ \quad \operatorname{is\text{-}unital\text{-}involutive\text{-}banach\text{-}algebra}(a) \Rightarrow (\operatorname{act}(a)(0, x) = \operatorname{vzero}(a)) \end{array}}}
\vnbprop{kal-uiba-act-vzero}{\vnbstmt{\begin{array}{@{}l@{}} \forall a.; \forall r \in \mathbb{C}. \\ \quad \operatorname{is\text{-}unital\text{-}involutive\text{-}banach\text{-}algebra}(a) \Rightarrow (\operatorname{act}(a)(r, \operatorname{vzero}(a)) = \operatorname{vzero}(a)) \end{array}}}
\vnbprop{kal-uiba-mul-assoc}{\vnbstmt{\begin{array}{@{}l@{}} \forall a.; \forall x, y, z \in \operatorname{vec}(a). \\ \quad \operatorname{is\text{-}unital\text{-}involutive\text{-}banach\text{-}algebra}(a) \Rightarrow  \\ \quad \quad \cdot_{a}(\cdot_{a}(x, y), z) = \cdot_{a}(x, \cdot_{a}(y, z)) \end{array}}}
\vnbprop{kal-uiba-mul-dist-l}{\vnbstmt{\begin{array}{@{}l@{}} \forall a.; \forall x, y, z \in \operatorname{vec}(a). \\ \quad \operatorname{is\text{-}unital\text{-}involutive\text{-}banach\text{-}algebra}(a) \Rightarrow  \\ \quad \quad \cdot_{a}(x, \operatorname{vadd}(a)(y, z)) = \operatorname{vadd}(a)(\cdot_{a}(x, y), \cdot_{a}(x, z)) \end{array}}}
\vnbprop{kal-uiba-mul-dist-r}{\vnbstmt{\begin{array}{@{}l@{}} \forall a.; \forall x, y, z \in \operatorname{vec}(a). \\ \quad \operatorname{is\text{-}unital\text{-}involutive\text{-}banach\text{-}algebra}(a) \Rightarrow  \\ \quad \quad \cdot_{a}(\operatorname{vadd}(a)(x, y), z) = \operatorname{vadd}(a)(\cdot_{a}(x, z), \cdot_{a}(y, z)) \end{array}}}
\vnbprop{kal-uiba-mul-act-l}{\vnbstmt{\begin{array}{@{}l@{}} \forall a.; \forall r \in \mathbb{C}, x, y \in \operatorname{vec}(a). \\ \quad \operatorname{is\text{-}unital\text{-}involutive\text{-}banach\text{-}algebra}(a) \Rightarrow  \\ \quad \quad \cdot_{a}(\operatorname{act}(a)(r, x), y) = \operatorname{act}(a)(r, \cdot_{a}(x, y)) \end{array}}}
\vnbprop{kal-uiba-mul-act-r}{\vnbstmt{\begin{array}{@{}l@{}} \forall a.; \forall r \in \mathbb{C}, x, y \in \operatorname{vec}(a). \\ \quad \operatorname{is\text{-}unital\text{-}involutive\text{-}banach\text{-}algebra}(a) \Rightarrow  \\ \quad \quad \cdot_{a}(x, \operatorname{act}(a)(r, y)) = \operatorname{act}(a)(r, \cdot_{a}(x, y)) \end{array}}}
\vnbprop{kal-uiba-mul-zero-l}{\vnbstmt{\begin{array}{@{}l@{}} \forall a.; \forall x \in \operatorname{vec}(a). \\ \quad \operatorname{is\text{-}unital\text{-}involutive\text{-}banach\text{-}algebra}(a) \Rightarrow (\cdot_{a}(\operatorname{vzero}(a), x) = \operatorname{vzero}(a)) \end{array}}}
\vnbprop{kal-uiba-mul-zero-r}{\vnbstmt{\begin{array}{@{}l@{}} \forall a.; \forall x \in \operatorname{vec}(a). \\ \quad \operatorname{is\text{-}unital\text{-}involutive\text{-}banach\text{-}algebra}(a) \Rightarrow (\cdot_{a}(x, \operatorname{vzero}(a)) = \operatorname{vzero}(a)) \end{array}}}
\vnbprop{kal-uiba-vone-in}{\vnbstmt{\forall a.\; (\operatorname{is\text{-}unital\text{-}involutive\text{-}banach\text{-}algebra}(a) \Rightarrow (1_{a} \in \operatorname{vec}(a)))}}
\vnbprop{kal-uiba-mul-one-l}{\vnbstmt{\begin{array}{@{}l@{}} \forall a.; \forall x \in \operatorname{vec}(a). \\ \quad \operatorname{is\text{-}unital\text{-}involutive\text{-}banach\text{-}algebra}(a) \Rightarrow (\cdot_{a}(1_{a}, x) = x) \end{array}}}
\vnbprop{kal-uiba-mul-one-r}{\vnbstmt{\begin{array}{@{}l@{}} \forall a.; \forall x \in \operatorname{vec}(a). \\ \quad \operatorname{is\text{-}unital\text{-}involutive\text{-}banach\text{-}algebra}(a) \Rightarrow (\cdot_{a}(x, 1_{a}) = x) \end{array}}}
\vnbprop{kal-uiba-vstar-fun}{\vnbstmt{\begin{array}{@{}l@{}} \forall a. \\ \quad \operatorname{is\text{-}unital\text{-}involutive\text{-}banach\text{-}algebra}(a) \Rightarrow ({}^{*_{a}} \in (\operatorname{vec}(a) \to \operatorname{vec}(a))) \end{array}}}
\vnbprop{kal-uiba-star-vadd}{\vnbstmt{\begin{array}{@{}l@{}} \forall a.; \forall x, y \in \operatorname{vec}(a). \\ \quad \operatorname{is\text{-}unital\text{-}involutive\text{-}banach\text{-}algebra}(a) \Rightarrow  \\ \quad \quad {}^{*_{a}}(\operatorname{vadd}(a)(x, y)) = \operatorname{vadd}(a)({}^{*_{a}}(x), {}^{*_{a}}(y)) \end{array}}}
\vnbprop{kal-uiba-star-act}{\vnbstmt{\begin{array}{@{}l@{}} \forall a.; \forall r \in \mathbb{C}, x \in \operatorname{vec}(a). \\ \quad \operatorname{is\text{-}unital\text{-}involutive\text{-}banach\text{-}algebra}(a) \Rightarrow  \\ \quad \quad {}^{*_{a}}(\operatorname{act}(a)(r, x)) = \operatorname{act}(a)(\operatorname{conjugate}(r), {}^{*_{a}}(x)) \end{array}}}
\vnbprop{kal-uiba-star-star}{\vnbstmt{\begin{array}{@{}l@{}} \forall a.; \forall x \in \operatorname{vec}(a). \\ \quad \operatorname{is\text{-}unital\text{-}involutive\text{-}banach\text{-}algebra}(a) \Rightarrow ({}^{*_{a}}({}^{*_{a}}(x)) = x) \end{array}}}
\vnbprop{kal-uiba-star-mul}{\vnbstmt{\begin{array}{@{}l@{}} \forall a.; \forall x, y \in \operatorname{vec}(a). \\ \quad \operatorname{is\text{-}unital\text{-}involutive\text{-}banach\text{-}algebra}(a) \Rightarrow  \\ \quad \quad {}^{*_{a}}(\cdot_{a}(x, y)) = \cdot_{a}({}^{*_{a}}(y), {}^{*_{a}}(x)) \end{array}}}
\vnbprop{kal-uiba-vstar-in}{\vnbstmt{\begin{array}{@{}l@{}} \forall a.; \forall x \in \operatorname{vec}(a). \\ \quad \operatorname{is\text{-}unital\text{-}involutive\text{-}banach\text{-}algebra}(a) \Rightarrow ({}^{*_{a}}(x) \in \operatorname{vec}(a)) \end{array}}}
\vnbprop{kal-uiba-star-zero}{\vnbstmt{\begin{array}{@{}l@{}} \forall a. \\ \quad \operatorname{is\text{-}unital\text{-}involutive\text{-}banach\text{-}algebra}(a) \Rightarrow ({}^{*_{a}}(\operatorname{vzero}(a)) = \operatorname{vzero}(a)) \end{array}}}
\vnbprop{kal-uiba-star-one}{\vnbstmt{\begin{array}{@{}l@{}} \forall a. \\ \quad \operatorname{is\text{-}unital\text{-}involutive\text{-}banach\text{-}algebra}(a) \Rightarrow ({}^{*_{a}}(1_{a}) = 1_{a}) \end{array}}}

\vnbbookfile{theorem-library/kal-unit}{theorem-library/kal-unit}
\noindent{\small kal-unit.scm -- wave of 2026-11-09/10 (agent KIT-ALG), wired 2026-11-10 from scratchpad/kitalg-1109/kal-unit.scm. Library: operator-algebras.}\par
\vnbprop{kun-vadd-assoc-pair}{\vnbstmt{\begin{array}{@{}l@{}} \forall a.; \forall l \in \mathbb{C}, x \in \operatorname{vec}(a), m \in \mathbb{C}, y \in \operatorname{vec}(a), n \in \mathbb{C}, z \in \operatorname{vec}(a). \\ \quad \operatorname{is\text{-}complex\text{-}normed\text{-}algebra}(a) \Rightarrow  \\ \quad \quad \operatorname{unit\text{-}vadd}(a)(\operatorname{unit\text{-}vadd}(a)({[}l, x], {[}m, y]), {[}n, z]) \\ \quad \quad =\; \operatorname{unit\text{-}vadd}(a)({[}l, x], \operatorname{unit\text{-}vadd}(a)({[}m, y], {[}n, z])) \end{array}}}
\vnbprop{unitization-vadd-assoc}{\vnbstmt{\begin{array}{@{}l@{}} \forall a.; \forall u, v, w \in \operatorname{unit\text{-}vec}(a). \\ \quad \operatorname{is\text{-}complex\text{-}normed\text{-}algebra}(a) \Rightarrow  \\ \quad \quad \operatorname{unit\text{-}vadd}(a)(\operatorname{unit\text{-}vadd}(a)(u, v), w) = \operatorname{unit\text{-}vadd}(a)(u, \operatorname{unit\text{-}vadd}(a)(v, w)) \end{array}}}
\vnbprop{kun-vadd-comm-pair}{\vnbstmt{\begin{array}{@{}l@{}} \forall a.; \forall l \in \mathbb{C}, x \in \operatorname{vec}(a), m \in \mathbb{C}, y \in \operatorname{vec}(a). \\ \quad \operatorname{is\text{-}complex\text{-}normed\text{-}algebra}(a) \Rightarrow  \\ \quad \quad \operatorname{unit\text{-}vadd}(a)({[}l, x], {[}m, y]) = \operatorname{unit\text{-}vadd}(a)({[}m, y], {[}l, x]) \end{array}}}
\vnbprop{unitization-vadd-comm}{\vnbstmt{\begin{array}{@{}l@{}} \forall a.; \forall u, v \in \operatorname{unit\text{-}vec}(a). \\ \quad \operatorname{is\text{-}complex\text{-}normed\text{-}algebra}(a) \Rightarrow (\operatorname{unit\text{-}vadd}(a)(u, v) = \operatorname{unit\text{-}vadd}(a)(v, u)) \end{array}}}
\vnbprop{kun-vadd-zero-left-pair}{\vnbstmt{\begin{array}{@{}l@{}} \forall a.; \forall l \in \mathbb{C}, x \in \operatorname{vec}(a). \\ \quad \operatorname{is\text{-}complex\text{-}normed\text{-}algebra}(a) \Rightarrow  \\ \quad \quad \operatorname{unit\text{-}vadd}(a)({[}0, \operatorname{vzero}(a)], {[}l, x]) = {[}l, x] \end{array}}}
\vnbprop{unitization-vadd-zero-left}{\vnbstmt{\begin{array}{@{}l@{}} \forall a.; \forall u \in \operatorname{unit\text{-}vec}(a). \\ \quad \operatorname{is\text{-}complex\text{-}normed\text{-}algebra}(a) \Rightarrow (\operatorname{unit\text{-}vadd}(a)({[}0, \operatorname{vzero}(a)], u) = u) \end{array}}}
\vnbprop{kun-vadd-zero-right-pair}{\vnbstmt{\begin{array}{@{}l@{}} \forall a.; \forall l \in \mathbb{C}, x \in \operatorname{vec}(a). \\ \quad \operatorname{is\text{-}complex\text{-}normed\text{-}algebra}(a) \Rightarrow  \\ \quad \quad \operatorname{unit\text{-}vadd}(a)({[}l, x], {[}0, \operatorname{vzero}(a)]) = {[}l, x] \end{array}}}
\vnbprop{unitization-vadd-zero-right}{\vnbstmt{\begin{array}{@{}l@{}} \forall a.; \forall u \in \operatorname{unit\text{-}vec}(a). \\ \quad \operatorname{is\text{-}complex\text{-}normed\text{-}algebra}(a) \Rightarrow (\operatorname{unit\text{-}vadd}(a)(u, {[}0, \operatorname{vzero}(a)]) = u) \end{array}}}
\vnbprop{kun-vadd-neg-right-pair}{\vnbstmt{\begin{array}{@{}l@{}} \forall a.; \forall l \in \mathbb{C}, x \in \operatorname{vec}(a). \\ \quad \operatorname{is\text{-}complex\text{-}normed\text{-}algebra}(a) \Rightarrow  \\ \quad \quad \operatorname{unit\text{-}vadd}(a)({[}l, x], \operatorname{unit\text{-}vneg}(a)({[}l, x])) = {[}0, \operatorname{vzero}(a)] \end{array}}}
\vnbprop{unitization-vadd-neg-right}{\vnbstmt{\begin{array}{@{}l@{}} \forall a.; \forall u \in \operatorname{unit\text{-}vec}(a). \\ \quad \operatorname{is\text{-}complex\text{-}normed\text{-}algebra}(a) \Rightarrow  \\ \quad \quad \operatorname{unit\text{-}vadd}(a)(u, \operatorname{unit\text{-}vneg}(a)(u)) = {[}0, \operatorname{vzero}(a)] \end{array}}}
\vnbprop{kun-vadd-neg-left-pair}{\vnbstmt{\begin{array}{@{}l@{}} \forall a.; \forall l \in \mathbb{C}, x \in \operatorname{vec}(a). \\ \quad \operatorname{is\text{-}complex\text{-}normed\text{-}algebra}(a) \Rightarrow  \\ \quad \quad \operatorname{unit\text{-}vadd}(a)(\operatorname{unit\text{-}vneg}(a)({[}l, x]), {[}l, x]) = {[}0, \operatorname{vzero}(a)] \end{array}}}
\vnbprop{unitization-vadd-neg-left}{\vnbstmt{\begin{array}{@{}l@{}} \forall a.; \forall u \in \operatorname{unit\text{-}vec}(a). \\ \quad \operatorname{is\text{-}complex\text{-}normed\text{-}algebra}(a) \Rightarrow  \\ \quad \quad \operatorname{unit\text{-}vadd}(a)(\operatorname{unit\text{-}vneg}(a)(u), u) = {[}0, \operatorname{vzero}(a)] \end{array}}}
\vnbprop{kun-act-vadd-pair}{\vnbstmt{\begin{array}{@{}l@{}} \forall a.; \forall c, l \in \mathbb{C}, x \in \operatorname{vec}(a), m \in \mathbb{C}, y \in \operatorname{vec}(a). \\ \quad \operatorname{is\text{-}complex\text{-}normed\text{-}algebra}(a) \Rightarrow  \\ \quad \quad \operatorname{unit\text{-}act}(a)(c, \operatorname{unit\text{-}vadd}(a)({[}l, x], {[}m, y])) \\ \quad \quad =\; \operatorname{unit\text{-}vadd}(a)(\operatorname{unit\text{-}act}(a)(c, {[}l, x]), \operatorname{unit\text{-}act}(a)(c, {[}m, y])) \end{array}}}
\vnbprop{unitization-act-vadd}{\vnbstmt{\begin{array}{@{}l@{}} \forall a.; \forall c \in \mathbb{C}, u, v \in \operatorname{unit\text{-}vec}(a). \\ \quad \operatorname{is\text{-}complex\text{-}normed\text{-}algebra}(a) \Rightarrow  \\ \quad \quad \operatorname{unit\text{-}act}(a)(c, \operatorname{unit\text{-}vadd}(a)(u, v)) = \operatorname{unit\text{-}vadd}(a)(\operatorname{unit\text{-}act}(a)(c, u), \operatorname{unit\text{-}act}(a)(c, v)) \end{array}}}
\vnbprop{kun-act-add-pair}{\vnbstmt{\begin{array}{@{}l@{}} \forall a.; \forall c, d, l \in \mathbb{C}, x \in \operatorname{vec}(a). \\ \quad \operatorname{is\text{-}complex\text{-}normed\text{-}algebra}(a) \Rightarrow  \\ \quad \quad \operatorname{unit\text{-}act}(a)((c + d), {[}l, x]) = \operatorname{unit\text{-}vadd}(a)(\operatorname{unit\text{-}act}(a)(c, {[}l, x]), \operatorname{unit\text{-}act}(a)(d, {[}l, x])) \end{array}}}
\vnbprop{unitization-act-add}{\vnbstmt{\begin{array}{@{}l@{}} \forall a.; \forall c, d \in \mathbb{C}, u \in \operatorname{unit\text{-}vec}(a). \\ \quad \operatorname{is\text{-}complex\text{-}normed\text{-}algebra}(a) \Rightarrow  \\ \quad \quad \operatorname{unit\text{-}act}(a)((c + d), u) = \operatorname{unit\text{-}vadd}(a)(\operatorname{unit\text{-}act}(a)(c, u), \operatorname{unit\text{-}act}(a)(d, u)) \end{array}}}
\vnbprop{kun-act-mul-pair}{\vnbstmt{\begin{array}{@{}l@{}} \forall a.; \forall c, d, l \in \mathbb{C}, x \in \operatorname{vec}(a). \\ \quad \operatorname{is\text{-}complex\text{-}normed\text{-}algebra}(a) \Rightarrow  \\ \quad \quad \operatorname{unit\text{-}act}(a)((c \cdot d), {[}l, x]) = \operatorname{unit\text{-}act}(a)(c, \operatorname{unit\text{-}act}(a)(d, {[}l, x])) \end{array}}}
\vnbprop{unitization-act-mul}{\vnbstmt{\begin{array}{@{}l@{}} \forall a.; \forall c, d \in \mathbb{C}, u \in \operatorname{unit\text{-}vec}(a). \\ \quad \operatorname{is\text{-}complex\text{-}normed\text{-}algebra}(a) \Rightarrow  \\ \quad \quad \operatorname{unit\text{-}act}(a)((c \cdot d), u) = \operatorname{unit\text{-}act}(a)(c, \operatorname{unit\text{-}act}(a)(d, u)) \end{array}}}
\vnbprop{kun-act-one-pair}{\vnbstmt{\begin{array}{@{}l@{}} \forall a.; \forall l \in \mathbb{C}, x \in \operatorname{vec}(a). \\ \quad \operatorname{is\text{-}complex\text{-}normed\text{-}algebra}(a) \Rightarrow (\operatorname{unit\text{-}act}(a)(1, {[}l, x]) = {[}l, x]) \end{array}}}
\vnbprop{unitization-act-one}{\vnbstmt{\begin{array}{@{}l@{}} \forall a.; \forall u \in \operatorname{unit\text{-}vec}(a). \\ \quad \operatorname{is\text{-}complex\text{-}normed\text{-}algebra}(a) \Rightarrow (\operatorname{unit\text{-}act}(a)(1, u) = u) \end{array}}}
\vnbprop{kun-vmul-assoc-pair}{\vnbstmt{\begin{array}{@{}l@{}} \forall a.; \forall l \in \mathbb{C}, x \in \operatorname{vec}(a), m \in \mathbb{C}, y \in \operatorname{vec}(a), n \in \mathbb{C}, z \in \operatorname{vec}(a). \\ \quad \operatorname{is\text{-}complex\text{-}normed\text{-}algebra}(a) \Rightarrow  \\ \quad \quad \operatorname{unit\text{-}vmul}(a)(\operatorname{unit\text{-}vmul}(a)({[}l, x], {[}m, y]), {[}n, z]) \\ \quad \quad =\; \operatorname{unit\text{-}vmul}(a)({[}l, x], \operatorname{unit\text{-}vmul}(a)({[}m, y], {[}n, z])) \end{array}}}
\vnbprop{unitization-vmul-assoc}{\vnbstmt{\begin{array}{@{}l@{}} \forall a.; \forall u, v, w \in \operatorname{unit\text{-}vec}(a). \\ \quad \operatorname{is\text{-}complex\text{-}normed\text{-}algebra}(a) \Rightarrow  \\ \quad \quad \operatorname{unit\text{-}vmul}(a)(\operatorname{unit\text{-}vmul}(a)(u, v), w) = \operatorname{unit\text{-}vmul}(a)(u, \operatorname{unit\text{-}vmul}(a)(v, w)) \end{array}}}
\vnbprop{kun-vmul-dist-left-pair}{\vnbstmt{\begin{array}{@{}l@{}} \forall a.; \forall l \in \mathbb{C}, x \in \operatorname{vec}(a), m \in \mathbb{C}, y \in \operatorname{vec}(a), n \in \mathbb{C}, z \in \operatorname{vec}(a). \\ \quad \operatorname{is\text{-}complex\text{-}normed\text{-}algebra}(a) \Rightarrow  \\ \quad \quad \operatorname{unit\text{-}vmul}(a)({[}l, x], \operatorname{unit\text{-}vadd}(a)({[}m, y], {[}n, z])) \\ \quad \quad =\; \operatorname{unit\text{-}vadd}(a)(\operatorname{unit\text{-}vmul}(a)({[}l, x], {[}m, y]), \operatorname{unit\text{-}vmul}(a)({[}l, x], {[}n, z])) \end{array}}}
\vnbprop{unitization-vmul-dist-left}{\vnbstmt{\begin{array}{@{}l@{}} \forall a.; \forall u, v, w \in \operatorname{unit\text{-}vec}(a). \\ \quad \operatorname{is\text{-}complex\text{-}normed\text{-}algebra}(a) \Rightarrow  \\ \quad \quad \operatorname{unit\text{-}vmul}(a)(u, \operatorname{unit\text{-}vadd}(a)(v, w)) = \operatorname{unit\text{-}vadd}(a)(\operatorname{unit\text{-}vmul}(a)(u, v), \operatorname{unit\text{-}vmul}(a)(u, w)) \end{array}}}
\vnbprop{kun-vmul-dist-right-pair}{\vnbstmt{\begin{array}{@{}l@{}} \forall a.; \forall l \in \mathbb{C}, x \in \operatorname{vec}(a), m \in \mathbb{C}, y \in \operatorname{vec}(a), n \in \mathbb{C}, z \in \operatorname{vec}(a). \\ \quad \operatorname{is\text{-}complex\text{-}normed\text{-}algebra}(a) \Rightarrow  \\ \quad \quad \operatorname{unit\text{-}vmul}(a)(\operatorname{unit\text{-}vadd}(a)({[}l, x], {[}m, y]), {[}n, z]) \\ \quad \quad =\; \operatorname{unit\text{-}vadd}(a)(\operatorname{unit\text{-}vmul}(a)({[}l, x], {[}n, z]), \operatorname{unit\text{-}vmul}(a)({[}m, y], {[}n, z])) \end{array}}}
\vnbprop{unitization-vmul-dist-right}{\vnbstmt{\begin{array}{@{}l@{}} \forall a.; \forall u, v, w \in \operatorname{unit\text{-}vec}(a). \\ \quad \operatorname{is\text{-}complex\text{-}normed\text{-}algebra}(a) \Rightarrow  \\ \quad \quad \operatorname{unit\text{-}vmul}(a)(\operatorname{unit\text{-}vadd}(a)(u, v), w) = \operatorname{unit\text{-}vadd}(a)(\operatorname{unit\text{-}vmul}(a)(u, w), \operatorname{unit\text{-}vmul}(a)(v, w)) \end{array}}}
\vnbprop{kun-vmul-act-left-pair}{\vnbstmt{\begin{array}{@{}l@{}} \forall a.; \forall c, l \in \mathbb{C}, x \in \operatorname{vec}(a), m \in \mathbb{C}, y \in \operatorname{vec}(a). \\ \quad \operatorname{is\text{-}complex\text{-}normed\text{-}algebra}(a) \Rightarrow  \\ \quad \quad \operatorname{unit\text{-}vmul}(a)(\operatorname{unit\text{-}act}(a)(c, {[}l, x]), {[}m, y]) = \operatorname{unit\text{-}act}(a)(c, \operatorname{unit\text{-}vmul}(a)({[}l, x], {[}m, y])) \end{array}}}
\vnbprop{unitization-vmul-act-left}{\vnbstmt{\begin{array}{@{}l@{}} \forall a.; \forall c \in \mathbb{C}, u, v \in \operatorname{unit\text{-}vec}(a). \\ \quad \operatorname{is\text{-}complex\text{-}normed\text{-}algebra}(a) \Rightarrow  \\ \quad \quad \operatorname{unit\text{-}vmul}(a)(\operatorname{unit\text{-}act}(a)(c, u), v) = \operatorname{unit\text{-}act}(a)(c, \operatorname{unit\text{-}vmul}(a)(u, v)) \end{array}}}
\vnbprop{kun-vmul-act-right-pair}{\vnbstmt{\begin{array}{@{}l@{}} \forall a.; \forall c, l \in \mathbb{C}, x \in \operatorname{vec}(a), m \in \mathbb{C}, y \in \operatorname{vec}(a). \\ \quad \operatorname{is\text{-}complex\text{-}normed\text{-}algebra}(a) \Rightarrow  \\ \quad \quad \operatorname{unit\text{-}vmul}(a)({[}l, x], \operatorname{unit\text{-}act}(a)(c, {[}m, y])) = \operatorname{unit\text{-}act}(a)(c, \operatorname{unit\text{-}vmul}(a)({[}l, x], {[}m, y])) \end{array}}}
\vnbprop{unitization-vmul-act-right}{\vnbstmt{\begin{array}{@{}l@{}} \forall a.; \forall c \in \mathbb{C}, u, v \in \operatorname{unit\text{-}vec}(a). \\ \quad \operatorname{is\text{-}complex\text{-}normed\text{-}algebra}(a) \Rightarrow  \\ \quad \quad \operatorname{unit\text{-}vmul}(a)(u, \operatorname{unit\text{-}act}(a)(c, v)) = \operatorname{unit\text{-}act}(a)(c, \operatorname{unit\text{-}vmul}(a)(u, v)) \end{array}}}

\vnbbookfile{theorem-library/kal-unit-cna}{theorem-library/kal-unit-cna}
\noindent{\small kal-unit-cna.scm -- wave of 2026-11-09/10 (agent KIT-ALG), wired 2026-11-10 from scratchpad/kitalg-1109/kal-unit-cna.scm. Library: operator-algebras.}\par
\vnbprop{kuc-vnrm-nonneg-pair}{\vnbstmt{\begin{array}{@{}l@{}} \forall a.; \forall l \in \mathbb{C}, x \in \operatorname{vec}(a). \\ \quad \operatorname{is\text{-}complex\text{-}normed\text{-}algebra}(a) \Rightarrow (0 \leq \operatorname{unit\text{-}vnrm}(a)({[}l, x])) \end{array}}}
\vnbprop{kuc-vnrm-definite-pair}{\vnbstmt{\begin{array}{@{}l@{}} \forall a.; \forall l \in \mathbb{C}, x \in \operatorname{vec}(a). \\ \quad \operatorname{is\text{-}complex\text{-}normed\text{-}algebra}(a) \Rightarrow  \\ \quad \quad (\operatorname{unit\text{-}vnrm}(a)({[}l, x]) = 0) \Leftrightarrow ({[}l, x] = {[}0, \operatorname{vzero}(a)]) \end{array}}}
\vnbprop{unitization-vnrm-nonneg}{\vnbstmt{\begin{array}{@{}l@{}} \forall a.; \forall u \in \operatorname{unit\text{-}vec}(a). \\ \quad \operatorname{is\text{-}complex\text{-}normed\text{-}algebra}(a) \Rightarrow (0 \leq \operatorname{unit\text{-}vnrm}(a)(u)) \end{array}}}
\vnbprop{unitization-vnrm-definite}{\vnbstmt{\begin{array}{@{}l@{}} \forall a.; \forall u \in \operatorname{unit\text{-}vec}(a). \\ \quad \operatorname{is\text{-}complex\text{-}normed\text{-}algebra}(a) \Rightarrow ((\operatorname{unit\text{-}vnrm}(a)(u) = 0) \Leftrightarrow (u = {[}0, \operatorname{vzero}(a)])) \end{array}}}
\vnbprop{unitization-vnrm-homog}{\vnbstmt{\begin{array}{@{}l@{}} \forall a.; \forall c \in \mathbb{C}, u \in \operatorname{unit\text{-}vec}(a). \\ \quad \operatorname{is\text{-}complex\text{-}normed\text{-}algebra}(a) \Rightarrow  \\ \quad \quad \operatorname{unit\text{-}vnrm}(a)(\operatorname{unit\text{-}act}(a)(c, u)) = (\operatorname{magnitude}(c) \cdot \operatorname{unit\text{-}vnrm}(a)(u)) \end{array}}}
\vnbprop{unitization-vnrm-triangle}{\vnbstmt{\begin{array}{@{}l@{}} \forall a.; \forall u, v \in \operatorname{unit\text{-}vec}(a). \\ \quad \operatorname{is\text{-}complex\text{-}normed\text{-}algebra}(a) \Rightarrow  \\ \quad \quad \operatorname{unit\text{-}vnrm}(a)(\operatorname{unit\text{-}vadd}(a)(u, v)) \leq (\operatorname{unit\text{-}vnrm}(a)(u) + \operatorname{unit\text{-}vnrm}(a)(v)) \end{array}}}
\vnbprop{unitization-vnrm-submult}{\vnbstmt{\begin{array}{@{}l@{}} \forall a.; \forall u, v \in \operatorname{unit\text{-}vec}(a). \\ \quad \operatorname{is\text{-}complex\text{-}normed\text{-}algebra}(a) \Rightarrow  \\ \quad \quad \operatorname{unit\text{-}vnrm}(a)(\operatorname{unit\text{-}vmul}(a)(u, v)) \leq (\operatorname{unit\text{-}vnrm}(a)(u) \cdot \operatorname{unit\text{-}vnrm}(a)(v)) \end{array}}}
\vnbprop{kuc-uv-set}{\vnbstmt{\forall a.\; (\operatorname{is\text{-}complex\text{-}normed\text{-}algebra}(a) \Rightarrow (\operatorname{unit\text{-}vec}(a) \in \mathbf{Set}))}}
\vnbprop{kuc-zero-in}{\vnbstmt{\forall a.\; (\operatorname{is\text{-}complex\text{-}normed\text{-}algebra}(a) \Rightarrow ({[}0, \operatorname{vzero}(a)] \in \operatorname{unit\text{-}vec}(a)))}}
\vnbprop{kuc-typ-vadd}{\vnbstmt{\begin{array}{@{}l@{}} \forall a. \\ \quad \operatorname{is\text{-}complex\text{-}normed\text{-}algebra}(a) \Rightarrow  \\ \quad \quad \operatorname{unit\text{-}vadd}(a) \in ((\operatorname{unit\text{-}vec}(a) \times \operatorname{unit\text{-}vec}(a)) \to \operatorname{unit\text{-}vec}(a)) \end{array}}}
\vnbprop{kuc-typ-vneg}{\vnbstmt{\begin{array}{@{}l@{}} \forall a. \\ \quad \operatorname{is\text{-}complex\text{-}normed\text{-}algebra}(a) \Rightarrow  \\ \quad \quad \operatorname{unit\text{-}vneg}(a) \in (\operatorname{unit\text{-}vec}(a) \to \operatorname{unit\text{-}vec}(a)) \end{array}}}
\vnbprop{kuc-typ-act}{\vnbstmt{\begin{array}{@{}l@{}} \forall a. \\ \quad \operatorname{is\text{-}complex\text{-}normed\text{-}algebra}(a) \Rightarrow  \\ \quad \quad \operatorname{unit\text{-}act}(a) \in ((\mathbb{C} \times \operatorname{unit\text{-}vec}(a)) \to \operatorname{unit\text{-}vec}(a)) \end{array}}}
\vnbprop{kuc-typ-vnrm}{\vnbstmt{\begin{array}{@{}l@{}} \forall a. \\ \quad \operatorname{is\text{-}complex\text{-}normed\text{-}algebra}(a) \Rightarrow (\operatorname{unit\text{-}vnrm}(a) \in (\operatorname{unit\text{-}vec}(a) \to \mathbb{R})) \end{array}}}
\vnbprop{kuc-typ-vmul}{\vnbstmt{\begin{array}{@{}l@{}} \forall a. \\ \quad \operatorname{is\text{-}complex\text{-}normed\text{-}algebra}(a) \Rightarrow  \\ \quad \quad \operatorname{unit\text{-}vmul}(a) \in ((\operatorname{unit\text{-}vec}(a) \times \operatorname{unit\text{-}vec}(a)) \to \operatorname{unit\text{-}vec}(a)) \end{array}}}
\vnbprop{kuc-ro-scal}{\vnbstmt{\begin{array}{@{}l@{}} \forall a. \\ \quad \operatorname{is\text{-}complex\text{-}normed\text{-}algebra}(a) \Rightarrow  \\ \quad \quad (\operatorname{scal}(\operatorname{unitization\text{-}na}(a)) \simeq \operatorname{normed\text{-}field\text{-}as\text{-}commutative\text{-}ring}(\operatorname{cc\text{-}normed\text{-}field})) \end{array}}}
\vnbprop{kuc-ro-vec}{\vnbstmt{\begin{array}{@{}l@{}} \forall a. \\ \quad \operatorname{is\text{-}complex\text{-}normed\text{-}algebra}(a) \Rightarrow (\operatorname{vec}(\operatorname{unitization\text{-}na}(a)) \simeq \operatorname{unit\text{-}vec}(a)) \end{array}}}
\vnbprop{kuc-ro-vadd}{\vnbstmt{\begin{array}{@{}l@{}} \forall a. \\ \quad \operatorname{is\text{-}complex\text{-}normed\text{-}algebra}(a) \Rightarrow (\operatorname{vadd}(\operatorname{unitization\text{-}na}(a)) \simeq \operatorname{unit\text{-}vadd}(a)) \end{array}}}
\vnbprop{kuc-ro-vzero}{\vnbstmt{\begin{array}{@{}l@{}} \forall a. \\ \quad \operatorname{is\text{-}complex\text{-}normed\text{-}algebra}(a) \Rightarrow (\operatorname{vzero}(\operatorname{unitization\text{-}na}(a)) \simeq {[}0, \operatorname{vzero}(a)]) \end{array}}}
\vnbprop{kuc-ro-vneg}{\vnbstmt{\begin{array}{@{}l@{}} \forall a. \\ \quad \operatorname{is\text{-}complex\text{-}normed\text{-}algebra}(a) \Rightarrow (\operatorname{vneg}(\operatorname{unitization\text{-}na}(a)) \simeq \operatorname{unit\text{-}vneg}(a)) \end{array}}}
\vnbprop{kuc-ro-act}{\vnbstmt{\begin{array}{@{}l@{}} \forall a. \\ \quad \operatorname{is\text{-}complex\text{-}normed\text{-}algebra}(a) \Rightarrow (\operatorname{act}(\operatorname{unitization\text{-}na}(a)) \simeq \operatorname{unit\text{-}act}(a)) \end{array}}}
\vnbprop{kuc-ro-vnrm}{\vnbstmt{\begin{array}{@{}l@{}} \forall a. \\ \quad \operatorname{is\text{-}complex\text{-}normed\text{-}algebra}(a) \Rightarrow (\operatorname{vnrm}(\operatorname{unitization\text{-}na}(a)) \simeq \operatorname{unit\text{-}vnrm}(a)) \end{array}}}
\vnbprop{kuc-ro-vmul}{\vnbstmt{\begin{array}{@{}l@{}} \forall a. \\ \quad \operatorname{is\text{-}complex\text{-}normed\text{-}algebra}(a) \Rightarrow (\cdot_{\operatorname{unitization\text{-}na}(a)} \simeq \operatorname{unit\text{-}vmul}(a)) \end{array}}}
\vnbprop{unitization-is-complex-normed-algebra}{\vnbstmt{\begin{array}{@{}l@{}} \forall a. \\ \quad \operatorname{is\text{-}complex\text{-}normed\text{-}algebra}(a) \Rightarrow  \\ \quad \quad \operatorname{is\text{-}complex\text{-}normed\text{-}algebra}(\operatorname{unitization\text{-}na}(a)) \end{array}}}

\vnbbookfile{structure-library/ba2-vocab}{structure-library/ba2-vocab}
\noindent{\small ba2-vocab.scm -- wave of 2026-11-09/10 (agent BA-2), wired 2026-11-10 from scratchpad/ba2-1109/ba2-vocab.scm. Library: operator-algebras.}\par
\par\medskip\noindent\textbf{Definition} (\texttt{is-unital-c-star-subalgebra}).\;For all $\operatorname{c2a\_}$ and $\operatorname{c2b\_}$, 
\begin{equation*}
\begin{array}{@{}l@{}} \operatorname{is\text{-}unital\text{-}c\text{-}star\text{-}subalgebra}(\operatorname{c2a\_}, \operatorname{c2b\_}) \\ \Leftrightarrow\; \operatorname{is\text{-}unital\text{-}c\text{-}star\text{-}algebra}(\operatorname{c2a\_}) \wedge \\ \operatorname{is\text{-}unital\text{-}c\text{-}star\text{-}algebra}(\operatorname{c2b\_}) \wedge \\ \forall \operatorname{c2x\_} \in \operatorname{vec}(\operatorname{c2a\_}).\; (\operatorname{c2x\_} \in \operatorname{vec}(\operatorname{c2b\_})) \wedge \\ 1_{\operatorname{c2a\_}} = 1_{\operatorname{c2b\_}} \wedge \\ \forall \operatorname{c2x\_}, \operatorname{c2y\_} \in \operatorname{vec}(\operatorname{c2a\_}). \\ \quad \operatorname{vadd}(\operatorname{c2a\_})(\operatorname{c2x\_}, \operatorname{c2y\_}) = \operatorname{vadd}(\operatorname{c2b\_})(\operatorname{c2x\_}, \operatorname{c2y\_}) \wedge \\ \forall \operatorname{c2x\_}, \operatorname{c2y\_} \in \operatorname{vec}(\operatorname{c2a\_}). \\ \quad \cdot_{\operatorname{c2a\_}}(\operatorname{c2x\_}, \operatorname{c2y\_}) = \cdot_{\operatorname{c2b\_}}(\operatorname{c2x\_}, \operatorname{c2y\_}) \wedge \\ \forall \operatorname{c2z\_} \in \mathbb{C}, \operatorname{c2x\_} \in \operatorname{vec}(\operatorname{c2a\_}). \\ \quad \operatorname{act}(\operatorname{c2a\_})(\operatorname{c2z\_}, \operatorname{c2x\_}) = \operatorname{act}(\operatorname{c2b\_})(\operatorname{c2z\_}, \operatorname{c2x\_}) \wedge \\ \forall \operatorname{c2x\_} \in \operatorname{vec}(\operatorname{c2a\_}).\; (\operatorname{vneg}(\operatorname{c2a\_})(\operatorname{c2x\_}) = \operatorname{vneg}(\operatorname{c2b\_})(\operatorname{c2x\_})) \wedge \\ \forall \operatorname{c2x\_} \in \operatorname{vec}(\operatorname{c2a\_}).\; ({}^{*_{\operatorname{c2a\_}}}(\operatorname{c2x\_}) = {}^{*_{\operatorname{c2b\_}}}(\operatorname{c2x\_})) \wedge \\ \forall \operatorname{c2x\_} \in \operatorname{vec}(\operatorname{c2a\_}).\; (\operatorname{vnrm}(\operatorname{c2a\_})(\operatorname{c2x\_}) = \operatorname{vnrm}(\operatorname{c2b\_})(\operatorname{c2x\_})) \end{array}
\end{equation*}


\vnbbookfile{structure-library/ba2b-gen-vocab}{structure-library/ba2b-gen-vocab}
\noindent{\small ba2b-gen-vocab.scm -- wave of 2026-11-09/10 (agent BA-2B), wired 2026-11-10 from scratchpad/ba2b-1109/ba2b-gen-vocab.scm. Library: operator-algebras.}\par
\par\medskip\noindent\textbf{Definition} (\texttt{is-closed-star-subalgebra-set}).\;For all $\operatorname{b2va\_}$ and $\operatorname{b2vs\_}$, 
\begin{equation*}
\begin{array}{@{}l@{}} \operatorname{is\text{-}closed\text{-}star\text{-}subalgebra\text{-}set}(\operatorname{b2va\_}, \operatorname{b2vs\_}) \\ \Leftrightarrow\; \operatorname{b2vs\_} \subseteq \operatorname{vec}(\operatorname{b2va\_}) \wedge \\ 1_{\operatorname{b2va\_}} \in \operatorname{b2vs\_} \wedge \\ \forall \operatorname{b2vx\_}, \operatorname{b2vy\_} \in \operatorname{b2vs\_}.\; (\operatorname{vadd}(\operatorname{b2va\_})(\operatorname{b2vx\_}, \operatorname{b2vy\_}) \in \operatorname{b2vs\_}) \wedge \\ \forall \operatorname{b2vz\_} \in \mathbb{C}, \operatorname{b2vx\_} \in \operatorname{b2vs\_}.\; (\operatorname{act}(\operatorname{b2va\_})(\operatorname{b2vz\_}, \operatorname{b2vx\_}) \in \operatorname{b2vs\_}) \wedge \\ \forall \operatorname{b2vx\_}, \operatorname{b2vy\_} \in \operatorname{b2vs\_}.\; (\cdot_{\operatorname{b2va\_}}(\operatorname{b2vx\_}, \operatorname{b2vy\_}) \in \operatorname{b2vs\_}) \wedge \\ \forall \operatorname{b2vx\_} \in \operatorname{b2vs\_}.\; ({}^{*_{\operatorname{b2va\_}}}(\operatorname{b2vx\_}) \in \operatorname{b2vs\_}) \wedge \\ \forall \operatorname{b2vy\_} \in \operatorname{vec}(\operatorname{b2va\_}). \\ \quad \forall \operatorname{b2ve\_}. \\ \quad \quad \operatorname{pos\text{-}rr}(\operatorname{b2ve\_}) \Rightarrow  \\ \quad \quad \quad \exists \operatorname{b2vx\_} \in \operatorname{b2vs\_}. \\ \quad \quad \quad \quad \operatorname{vnrm}(\operatorname{b2va\_})(\operatorname{vadd}(\operatorname{b2va\_})(\operatorname{b2vy\_}, \operatorname{vneg}(\operatorname{b2va\_})(\operatorname{b2vx\_}))) < \operatorname{b2ve\_} \Rightarrow \\ \quad \quad \operatorname{b2vy\_} \in \operatorname{b2vs\_} \end{array}
\end{equation*}

\par\medskip\noindent\textbf{Definition} (\texttt{cstar-gen}).\;$\operatorname{cstar\text{-}gen}(\operatorname{b2va\_}, \operatorname{b2vx\_}) :=$
\begin{equation*}
\operatorname{intersection\text{-}of}(\{\operatorname{b2vs\_} \in \operatorname{power}(\operatorname{vec}(\operatorname{b2va\_})) : (\operatorname{is\text{-}closed\text{-}star\text{-}subalgebra\text{-}set}(\operatorname{b2va\_}, \operatorname{b2vs\_}) \wedge (\operatorname{b2vx\_} \in \operatorname{b2vs\_}))\})
\end{equation*}


\vnbbookfile{structure-library/basw2-vocab}{structure-library/basw2-vocab}
\noindent{\small basw2-vocab.scm -- wave of 2026-11-09/10 (agent BA-SW2), wired 2026-11-10 from scratchpad/basw2-1109/basw2-vocab.scm. Library: operator-algebras.}\par
\par\medskip\noindent\textbf{Definition} (\texttt{pos-cone}).\;${\operatorname{pca\_}}_{+} :=$
\begin{equation*}
\{\operatorname{pcx\_} \in \operatorname{vec}(\operatorname{pca\_}) : \forall \operatorname{pcc\_} \in \operatorname{power}(\operatorname{vec}(\operatorname{pca\_})).\; (\forall \operatorname{pcy\_} \in \operatorname{vec}(\operatorname{pca\_}).\; (\cdot_{\operatorname{pca\_}}({}^{*_{\operatorname{pca\_}}}(\operatorname{pcy\_}), \operatorname{pcy\_}) \in \operatorname{pcc\_}) \Rightarrow (\forall \operatorname{pcu\_}, \operatorname{pcv\_} \in \operatorname{pcc\_}.\; (\operatorname{vadd}(\operatorname{pca\_})(\operatorname{pcu\_}, \operatorname{pcv\_}) \in \operatorname{pcc\_}) \Rightarrow (\operatorname{pcx\_} \in \operatorname{pcc\_})))\}
\end{equation*}


\vnbbookfile{theorem-library/ba-fekete}{theorem-library/ba-fekete}
\noindent{\small ba-fekete.scm -- wave of 2026-11-09/10 (agent BA-1C), wired 2026-11-10 from scratchpad/ba1c-1109/ba-fekete.scm. Library: operator-algebras.}\par
\vnbprop{baf-el-pow-add}{\vnbstmt{\begin{array}{@{}l@{}} \forall n \in \mathbb{N}, a.; \forall x \in \operatorname{vec}(a), m \in \mathbb{N}. \\ \quad \operatorname{is\text{-}unital\text{-}banach\text{-}algebra}(a) \Rightarrow  \\ \quad \quad {x}^{(m + n)} = \cdot_{a}({x}^{m}, {x}^{n}) \end{array}}}
\vnbprop{baf-norm-pow-submult}{\vnbstmt{\begin{array}{@{}l@{}} \forall a.; \forall x \in \operatorname{vec}(a), m, n \in \mathbb{N}. \\ \quad \operatorname{is\text{-}unital\text{-}banach\text{-}algebra}(a) \Rightarrow  \\ \quad \quad \operatorname{vnrm}(a)({x}^{(m + n)}) \leq (\operatorname{vnrm}(a)({x}^{m}) \cdot \operatorname{vnrm}(a)({x}^{n})) \end{array}}}
\vnbprop{spectral-radius-limit}{\vnbstmt{\begin{array}{@{}l@{}} \forall a.; \forall x \in \operatorname{vec}(a). \\ \quad \operatorname{is\text{-}unital\text{-}banach\text{-}algebra}(a) \Rightarrow  \\ \quad \quad \operatorname{converges\text{-}to}(\operatorname{rr\text{-}ms}, \\ \quad \quad \quad (\lambda n \in \mathbb{N}.\; {\operatorname{vnrm}(a)({x}^{(n + 1)})}^{{(n + 1)}^{\ast\!-1}}), \\ \quad \quad \quad \rho({x})) \end{array}}}

\vnbbookfile{structure-library/basw3-vocab}{structure-library/basw3-vocab}
\noindent{\small basw3-vocab.scm -- wave of 2026-11-09/10 (agent BA-SW3), wired 2026-11-10 from scratchpad/basw3-1109/basw3-vocab.scm. Library: operator-algebras.}\par
\par\medskip\noindent\textbf{Definition} (\texttt{point-eval}).\;$\operatorname{point\text{-}eval}(\operatorname{uet\_}, \operatorname{uex\_}) :=$ $(\lambda \operatorname{uef\_} \in C({\operatorname{uet\_}}).\; \operatorname{uef\_}(\operatorname{uex\_}))$
\par\medskip\noindent\textbf{Definition} (\texttt{ury-l}).\;$\operatorname{ury\text{-}l}(\operatorname{uat\_}, \operatorname{uac\_}) :=$ ${[}\operatorname{nth}(1, \operatorname{uac\_}), \operatorname{ury\text{-}v}(\operatorname{uat\_}, \operatorname{uac\_})]$
\par\medskip\noindent\textbf{Definition} (\texttt{ury-r}).\;$\operatorname{ury\text{-}r}(\operatorname{uat\_}, \operatorname{uac\_}) :=$ ${[}\operatorname{ury\text{-}k}(\operatorname{uat\_}, \operatorname{uac\_}), \operatorname{nth}(2, \operatorname{uac\_})]$
\par\medskip\noindent\textbf{Definition} (\texttt{cu-pairs}).\;$\operatorname{cu\text{-}pairs}(\operatorname{uat\_}) :=$
\begin{equation*}
\{\operatorname{ucw\_} \in (\operatorname{power}(\operatorname{pts}(\operatorname{uat\_})) \times \operatorname{opens}(\operatorname{uat\_})) : (\operatorname{is\text{-}top\text{-}closed}(\operatorname{uat\_}, \operatorname{nth}(1, \operatorname{ucw\_})) \wedge (\operatorname{nth}(1, \operatorname{ucw\_}) \subseteq \operatorname{nth}(2, \operatorname{ucw\_})))\}
\end{equation*}

\par\medskip\noindent\textbf{Definition} (\texttt{ury-v}).\;$\operatorname{ury\text{-}v}(\operatorname{uat\_}, \operatorname{uac\_}) :=$
\begin{equation*}
\varepsilon(\{\operatorname{uvv\_} \in \operatorname{opens}(\operatorname{uat\_}) : \exists \operatorname{uvk\_} \in \operatorname{power}(\operatorname{pts}(\operatorname{uat\_})).\; (\operatorname{is\text{-}top\text{-}closed}(\operatorname{uat\_}, \operatorname{uvk\_}) \wedge ((\operatorname{nth}(1, \operatorname{uac\_}) \subseteq \operatorname{uvv\_}) \wedge ((\operatorname{uvv\_} \subseteq \operatorname{uvk\_}) \wedge (\operatorname{uvk\_} \subseteq \operatorname{nth}(2, \operatorname{uac\_})))))\})
\end{equation*}

\par\medskip\noindent\textbf{Definition} (\texttt{ury-k}).\;$\operatorname{ury\text{-}k}(\operatorname{uat\_}, \operatorname{uac\_}) :=$
\begin{equation*}
\varepsilon(\{\operatorname{uvk\_} \in \operatorname{power}(\operatorname{pts}(\operatorname{uat\_})) : (\operatorname{is\text{-}top\text{-}closed}(\operatorname{uat\_}, \operatorname{uvk\_}) \wedge ((\operatorname{ury\text{-}v}(\operatorname{uat\_}, \operatorname{uac\_}) \subseteq \operatorname{uvk\_}) \wedge (\operatorname{uvk\_} \subseteq \operatorname{nth}(2, \operatorname{uac\_}))))\})
\end{equation*}

\par\medskip\noindent\textbf{Definition} (\texttt{ury-apx}).\;For all $\operatorname{uat\_}$ and $\operatorname{uax\_}$, 
\begin{equation*}
(\operatorname{ury\text{-}apx}(\operatorname{uat\_}, \operatorname{uax\_}, 0) \simeq (\lambda \operatorname{uac\_} \in \operatorname{cu\text{-}pairs}(\operatorname{uat\_}).\; \left(\text{if }(\operatorname{uax\_} \in \operatorname{nth}(2, \operatorname{uac\_}))\text{ then }0\text{ else }1\right)))
\end{equation*}
\par
Let $\operatorname{uat\_}$ and $\operatorname{uax\_}$ be arbitrary. Let $\operatorname{uan\_} \in \mathbb{N}$. Then
\begin{equation*}
(\operatorname{ury\text{-}apx}(\operatorname{uat\_}, \operatorname{uax\_}, \operatorname{succ}(\operatorname{uan\_})) \simeq (\lambda \operatorname{uac\_} \in \operatorname{cu\text{-}pairs}(\operatorname{uat\_}).\; \cdot(1/2, +(\operatorname{ury\text{-}apx}(\operatorname{uat\_}, \operatorname{uax\_}, \operatorname{uan\_})(\operatorname{ury\text{-}l}(\operatorname{uat\_}, \operatorname{uac\_})), \operatorname{ury\text{-}apx}(\operatorname{uat\_}, \operatorname{uax\_}, \operatorname{uan\_})(\operatorname{ury\text{-}r}(\operatorname{uat\_}, \operatorname{uac\_}))))))
\end{equation*}

\par\medskip\noindent\textbf{Definition} (\texttt{ury-lim}).\;$\operatorname{ury\text{-}lim}(\operatorname{uat\_}, \operatorname{uac\_}, \operatorname{uax\_}) :=$
\begin{equation*}
\operatorname{esup}(\{\operatorname{ulr\_} \in \mathbb{R} : \exists \operatorname{uln\_} \in \mathbb{N}.\; (\operatorname{ulr\_} = \operatorname{ury\text{-}apx}(\operatorname{uat\_}, \operatorname{uax\_}, \operatorname{uln\_})(\operatorname{uac\_}))\})
\end{equation*}


\vnbbookfile{structure-library/cnvs-ba-vocab}{structure-library/cnvs-ba-vocab}
\noindent{\small cnvs-ba-vocab.scm -- wave of 2026-11-09/10 (agent CNVS), wired 2026-11-10 from scratchpad/cnvs-1109/cnvs-ba-vocab.scm. Library: operator-algebras.}\par
\par\medskip\noindent\textbf{Definition} (\texttt{cna-cnvs}).\;$\operatorname{cna\text{-}cnvs}(\operatorname{cna\_}) :=$ $\operatorname{cmod\text{-}cnvs}(\operatorname{cna\text{-}as\text{-}module}(\operatorname{cna\_}), \operatorname{vnrm}(\operatorname{cna\_}))$

\vnbbookfile{structure-library/ba-optop-vocab}{structure-library/ba-optop-vocab}
\noindent{\small ba-optop-vocab.scm -- wave of 2026-11-09/10 (agent BA-OT1), wired 2026-11-10 from scratchpad/baot1-1109/ba-optop-vocab.scm. Library: operator-algebras.}\par
\par\medskip\noindent\textbf{Definition} (\texttt{bh-module}).\;$L({\operatorname{otv\_}}) :=$
\begin{equation*}
{[}\operatorname{normed\text{-}field\text{-}as\text{-}commutative\text{-}ring}(\operatorname{cc\text{-}normed\text{-}field}), L({\operatorname{otv\_}}, {\operatorname{otv\_}}), \operatorname{bh\text{-}vadd}(\operatorname{otv\_}), 0, \operatorname{bh\text{-}vneg}(\operatorname{otv\_}), \operatorname{bh\text{-}act}(\operatorname{otv\_})]
\end{equation*}

\par\medskip\noindent\textbf{Definition} (\texttt{sot}).\;$L({\operatorname{otv\_}})_{s} :=$ $\operatorname{seminorm\text{-}top}(L({\operatorname{otv\_}}), \operatorname{vec}(\operatorname{otv\_}), \operatorname{sot\text{-}pf}(\operatorname{otv\_}))$
\par\medskip\noindent\textbf{Definition} (\texttt{wot}).\;$L({\operatorname{otv\_}})_{w} :=$
\begin{equation*}
\operatorname{seminorm\text{-}top}(L({\operatorname{otv\_}}), (\operatorname{vec}(\operatorname{otv\_}) \times \operatorname{vec}(\operatorname{otv\_})), \operatorname{wot\text{-}pf}(\operatorname{otv\_}))
\end{equation*}

\par\medskip\noindent\textbf{Definition} (\texttt{norm-optop}).\;$L({\operatorname{otv\_}}) :=$
\begin{equation*}
\operatorname{seminorm\text{-}top}(L({\operatorname{otv\_}}), \{0\}, \operatorname{norm\text{-}optop\text{-}pf}(\operatorname{otv\_}))
\end{equation*}

\par\medskip\noindent\textbf{Definition} (\texttt{ust}).\;$L({\operatorname{otv\_}})_{us} :=$
\begin{equation*}
\operatorname{seminorm\text{-}top}(L({\operatorname{otv\_}}), \operatorname{sq\text{-}summable\text{-}seqs}(\operatorname{otv\_}), \operatorname{ust\text{-}pf}(\operatorname{otv\_}))
\end{equation*}

\par\medskip\noindent\textbf{Definition} (\texttt{sq-summable-seqs}).\;$\operatorname{sq\text{-}summable\text{-}seqs}(\operatorname{otv\_}) :=$
\begin{equation*}
\{\operatorname{ots\_} \in (\mathbb{N} \to \operatorname{vec}(\operatorname{otv\_})) : (\operatorname{esum}((\lambda \operatorname{otk\_} \in \mathbb{N}.\; (\lVert \operatorname{ots\_}(\operatorname{otk\_}) \rVert \cdot \lVert \operatorname{ots\_}(\operatorname{otk\_}) \rVert))) \in \mathbb{R})\}
\end{equation*}

\par\medskip\noindent\textbf{Definition} (\texttt{uwt}).\;$L({\operatorname{otv\_}})_{uw} :=$
\begin{equation*}
\begin{array}{@{}l@{}} \operatorname{seminorm\text{-}top}(L({\operatorname{otv\_}}), \\ \quad (\operatorname{sq\text{-}summable\text{-}seqs}(\operatorname{otv\_}) \times \operatorname{sq\text{-}summable\text{-}seqs}(\operatorname{otv\_})), \\ \quad \operatorname{uwt\text{-}pf}(\operatorname{otv\_})) \end{array}
\end{equation*}

\par\medskip\noindent\textbf{Definition} (\texttt{cc-sn-top}).\;
\begin{equation*}
(\operatorname{cc\text{-}sn\text{-}top} \simeq \operatorname{seminorm\text{-}top}(\operatorname{cc\text{-}sn\text{-}module}, \{0\}, \operatorname{cc\text{-}sn\text{-}pf}))
\end{equation*}

\par\medskip\noindent\textbf{Definition} (\texttt{cc-sn-module}).\;
\begin{equation*}
(\operatorname{cc\text{-}sn\text{-}module} \simeq \operatorname{ring\text{-}as\text{-}module}(\operatorname{normed\text{-}field\text{-}as\text{-}commutative\text{-}ring}(\operatorname{cc\text{-}normed\text{-}field})))
\end{equation*}

\par\medskip\noindent\textbf{Definition} (\texttt{strong-dual}).\;$L_{\bullet}({\operatorname{otv\_}}) :=$ $\operatorname{op\text{-}dual}(\operatorname{otv\_}, L({\operatorname{otv\_}})_{s})$
\par\medskip\noindent\textbf{Definition} (\texttt{weak-dual}).\;$\operatorname{weak\text{-}dual}(\operatorname{otv\_}) :=$ $\operatorname{op\text{-}dual}(\operatorname{otv\_}, L({\operatorname{otv\_}})_{w})$
\par\medskip\noindent\textbf{Definition} (\texttt{ultra-dual}).\;$L_{*}({\operatorname{otv\_}}) :=$ $\operatorname{op\text{-}dual}(\operatorname{otv\_}, L({\operatorname{otv\_}})_{us})$
\par\medskip\noindent\textbf{Definition} (\texttt{vec-functional}).\;$\omega_{{\operatorname{otx\_}},{\operatorname{oty\_}}} :=$
\begin{equation*}
(\lambda \operatorname{ott\_} \in L({\operatorname{otv\_}}, {\operatorname{otv\_}}).\; \langle\cdot,\cdot\rangle_{\operatorname{otv\_}}(\operatorname{ott\_}(\operatorname{otx\_}), \operatorname{oty\_}))
\end{equation*}

\par\medskip\noindent\textbf{Definition} (\texttt{weak-pairing-top}).\;$w\langle {\operatorname{otq\_}}, {\operatorname{ota\_}}\rangle :=$ $\operatorname{seminorm\text{-}top}(\operatorname{otq\_}, \operatorname{ota\_}, \operatorname{weak\text{-}pairing\text{-}pf}(\operatorname{otq\_}, \operatorname{ota\_}))$
\par\medskip\noindent\textbf{Definition} (\texttt{top-closure}).\;$\overline{\operatorname{ota\_}} :=$
\begin{equation*}
\{\operatorname{otx\_} \in \operatorname{pts}(\operatorname{otq\_}) : \forall \operatorname{otu\_} \in \operatorname{opens}(\operatorname{otq\_}).\; ((\operatorname{otx\_} \in \operatorname{otu\_}) \Rightarrow \exists \operatorname{oty\_} \in \operatorname{ota\_}.\; (\operatorname{oty\_} \in \operatorname{otu\_}))\}
\end{equation*}

\par\medskip\noindent\textbf{Definition} (\texttt{cc-sn-pf}).\;
\begin{equation*}
(\operatorname{cc\text{-}sn\text{-}pf} \simeq (\lambda \operatorname{otk\_} \in \{0\}.\; (\lambda \operatorname{otz\_} \in \mathbb{C}.\; \operatorname{magnitude}(\operatorname{otz\_}))))
\end{equation*}

\par\medskip\noindent\textbf{Definition} (\texttt{sot-pf}).\;$\operatorname{sot\text{-}pf}(\operatorname{otv\_}) :=$
\begin{equation*}
(\lambda \operatorname{otx\_} \in \operatorname{vec}(\operatorname{otv\_}).\; (\lambda \operatorname{ott\_} \in L({\operatorname{otv\_}}, {\operatorname{otv\_}}).\; \lVert \operatorname{ott\_}(\operatorname{otx\_}) \rVert))
\end{equation*}

\par\medskip\noindent\textbf{Definition} (\texttt{wot-pf}).\;$\operatorname{wot\text{-}pf}(\operatorname{otv\_}) :=$
\begin{equation*}
\begin{array}{@{}l@{}} (\lambda \operatorname{otp\_} \in (\operatorname{vec}(\operatorname{otv\_}) \times \operatorname{vec}(\operatorname{otv\_})).\; \\ \quad (\lambda \operatorname{ott\_} \in L({\operatorname{otv\_}}, {\operatorname{otv\_}}).\; \operatorname{magnitude}(\langle\cdot,\cdot\rangle_{\operatorname{otv\_}}(\operatorname{ott\_}(\operatorname{nth}(1, \operatorname{otp\_})), \operatorname{nth}(2, \operatorname{otp\_}))))) \end{array}
\end{equation*}

\par\medskip\noindent\textbf{Definition} (\texttt{norm-optop-pf}).\;$\operatorname{norm\text{-}optop\text{-}pf}(\operatorname{otv\_}) :=$
\begin{equation*}
(\lambda \operatorname{otk\_} \in \{0\}.\; (\lambda \operatorname{ott\_} \in L({\operatorname{otv\_}}, {\operatorname{otv\_}}).\; \lVert {\operatorname{ott\_}} \rVert))
\end{equation*}

\par\medskip\noindent\textbf{Definition} (\texttt{ust-pf}).\;$\operatorname{ust\text{-}pf}(\operatorname{otv\_}) :=$
\begin{equation*}
\begin{array}{@{}l@{}} (\lambda \operatorname{ots\_} \in \operatorname{sq\text{-}summable\text{-}seqs}(\operatorname{otv\_}).\; \\ \quad (\lambda \operatorname{ott\_} \in L({\operatorname{otv\_}}, {\operatorname{otv\_}}).\; \\ \quad \quad \sqrt{\operatorname{esum}((\lambda \operatorname{otk\_} \in \mathbb{N}.\; (\lVert \operatorname{ott\_}(\operatorname{ots\_}(\operatorname{otk\_})) \rVert \cdot \lVert \operatorname{ott\_}(\operatorname{ots\_}(\operatorname{otk\_})) \rVert)))})) \end{array}
\end{equation*}

\par\medskip\noindent\textbf{Definition} (\texttt{uwt-pf}).\;$\operatorname{uwt\text{-}pf}(\operatorname{otv\_}) :=$
\begin{equation*}
\begin{array}{@{}l@{}} (\lambda \operatorname{otp\_} \in (\operatorname{sq\text{-}summable\text{-}seqs}(\operatorname{otv\_}) \times \operatorname{sq\text{-}summable\text{-}seqs}(\operatorname{otv\_})).\; \\ \quad (\lambda \operatorname{ott\_} \in L({\operatorname{otv\_}}, {\operatorname{otv\_}}).\; \\ \quad \quad \operatorname{magnitude}(\operatorname{cc\text{-}series\text{-}limit}((\lambda \operatorname{otk\_} \in \mathbb{N}.\; \langle\cdot,\cdot\rangle_{\operatorname{otv\_}}(\operatorname{ott\_}(\operatorname{nth}(1, \operatorname{otp\_})(\operatorname{otk\_})), \operatorname{nth}(2, \operatorname{otp\_})(\operatorname{otk\_}))))))) \end{array}
\end{equation*}

\par\medskip\noindent\textbf{Definition} (\texttt{op-dual}).\;$\operatorname{op\text{-}dual}(\operatorname{otv\_}, \operatorname{otu\_}) :=$
\begin{equation*}
\{\operatorname{otf\_} \in (L({\operatorname{otv\_}}, {\operatorname{otv\_}}) \to \mathbb{C}) : (\operatorname{is\text{-}linear\text{-}map}(L({\operatorname{otv\_}}), \operatorname{cc\text{-}sn\text{-}module}, \operatorname{otf\_}) \wedge \operatorname{is\text{-}hom\text{-}top\text{-}space}(\operatorname{otu\_}, \operatorname{cc\text{-}sn\text{-}top}, \operatorname{otf\_}))\}
\end{equation*}

\par\medskip\noindent\textbf{Definition} (\texttt{ultraweak-dual}).\;$\operatorname{ultraweak\text{-}dual}(\operatorname{otv\_}) :=$ $\operatorname{op\text{-}dual}(\operatorname{otv\_}, L({\operatorname{otv\_}})_{uw})$
\par\medskip\noindent\textbf{Definition} (\texttt{weak-pairing-pf}).\;$\operatorname{weak\text{-}pairing\text{-}pf}(\operatorname{otq\_}, \operatorname{ota\_}) :=$
\begin{equation*}
(\lambda \operatorname{otf\_} \in \operatorname{ota\_}.\; (\lambda \operatorname{otx\_} \in \operatorname{vec}(\operatorname{otq\_}).\; \operatorname{magnitude}(\operatorname{otf\_}(\operatorname{otx\_}))))
\end{equation*}


\vnbbookfile{structure-library/ba-unbdd-vocab}{structure-library/ba-unbdd-vocab}
\noindent{\small ba-unbdd-vocab.scm -- wave of 2026-11-09/10 (agent BA-G1), wired 2026-11-10 from scratchpad/bag1-1109/ba-unbdd-vocab.scm. Library: operator-algebras.}\par
\par\medskip\noindent\textbf{Definition} (\texttt{graph-j}).\;$J{\operatorname{ubg\_}} :=$
\begin{equation*}
\{\operatorname{ubp\_} \in (\operatorname{vec}(\operatorname{ubv\_}) \times \operatorname{vec}(\operatorname{ubv\_})) : ({[}\operatorname{nth}(2, \operatorname{ubp\_}), \operatorname{vneg}(\operatorname{ubv\_})(\operatorname{nth}(1, \operatorname{ubp\_}))] \in \operatorname{ubg\_})\}
\end{equation*}

\par\medskip\noindent\textbf{Definition} (\texttt{form-plus-space}).\;$H_{{\operatorname{ubf\_}}+} :=$
\begin{equation*}
{[}\operatorname{normed\text{-}field\text{-}as\text{-}commutative\text{-}ring}(\operatorname{cc\text{-}normed\text{-}field}), \operatorname{ubd\_}, (\lambda {[}\operatorname{ubx\_}, \operatorname{uby\_}] \in (\operatorname{ubd\_} \times \operatorname{ubd\_}).\; \operatorname{vadd}(\operatorname{ubv\_})(\operatorname{ubx\_}, \operatorname{uby\_})), \operatorname{vzero}(\operatorname{ubv\_}), (\lambda \operatorname{ubx\_} \in \operatorname{ubd\_}.\; \operatorname{vneg}(\operatorname{ubv\_})(\operatorname{ubx\_})), (\lambda {[}\operatorname{ubk\_}, \operatorname{ubx\_}] \in (\mathbb{C} \times \operatorname{ubd\_}).\; \operatorname{act}(\operatorname{ubv\_})(\operatorname{ubk\_}, \operatorname{ubx\_})), {\operatorname{ubf\_}}_{+}]
\end{equation*}

\par\medskip\noindent\textbf{Definition} (\texttt{form-operator}).\;$T_{\operatorname{ubf\_}} :=$
\begin{equation*}
\{\operatorname{ubp\_} \in (\operatorname{vec}(\operatorname{ubv\_}) \times \operatorname{vec}(\operatorname{ubv\_})) : ((\operatorname{nth}(1, \operatorname{ubp\_}) \in \operatorname{ubd\_}) \wedge \forall \operatorname{ubx\_} \in \operatorname{ubd\_}.\; (\langle\cdot,\cdot\rangle_{\operatorname{ubv\_}}(\operatorname{ubx\_}, \operatorname{nth}(2, \operatorname{ubp\_})) = \operatorname{ubf\_}(\operatorname{ubx\_}, \operatorname{nth}(1, \operatorname{ubp\_}))))\}
\end{equation*}

\par\medskip\noindent\textbf{Definition} (\texttt{is-linear-graph}).\;For all $\operatorname{ubv\_}$ and $\operatorname{ubg\_}$, 
\begin{equation*}
\operatorname{is\text{-}linear\text{-}graph}(\operatorname{ubv\_}, \operatorname{ubg\_}) \Leftrightarrow \operatorname{is\text{-}submodule}({\operatorname{ubv\_}} \oplus {\operatorname{ubv\_}}, \operatorname{ubg\_})
\end{equation*}

\par\medskip\noindent\textbf{Definition} (\texttt{graph-dom}).\;$\mathrm{dom}\,{\operatorname{ubg\_}} :=$
\begin{equation*}
\{\operatorname{ubx\_} \in \operatorname{vec}(\operatorname{ubv\_}) : \exists \operatorname{uby\_} \in \operatorname{vec}(\operatorname{ubv\_}).\; ({[}\operatorname{ubx\_}, \operatorname{uby\_}] \in \operatorname{ubg\_})\}
\end{equation*}

\par\medskip\noindent\textbf{Definition} (\texttt{graph-ran}).\;$\mathrm{ran}\,{\operatorname{ubg\_}} :=$
\begin{equation*}
\{\operatorname{uby\_} \in \operatorname{vec}(\operatorname{ubv\_}) : \exists \operatorname{ubx\_} \in \operatorname{vec}(\operatorname{ubv\_}).\; ({[}\operatorname{ubx\_}, \operatorname{uby\_}] \in \operatorname{ubg\_})\}
\end{equation*}

\par\medskip\noindent\textbf{Definition} (\texttt{is-functional-graph}).\;For all $\operatorname{ubv\_}$ and $\operatorname{ubg\_}$, 
\begin{equation*}
\begin{array}{@{}l@{}} \operatorname{is\text{-}functional\text{-}graph}(\operatorname{ubv\_}, \operatorname{ubg\_}) \\ \Leftrightarrow\; \forall \operatorname{ubx\_}, \operatorname{uby\_}, \operatorname{ubz\_} \in \operatorname{vec}(\operatorname{ubv\_}). \\ \quad ({[}\operatorname{ubx\_}, \operatorname{uby\_}] \in \operatorname{ubg\_}) \Rightarrow (({[}\operatorname{ubx\_}, \operatorname{ubz\_}] \in \operatorname{ubg\_}) \Rightarrow (\operatorname{uby\_} = \operatorname{ubz\_})) \end{array}
\end{equation*}

\par\medskip\noindent\textbf{Definition} (\texttt{graph-adjoint}).\;${\operatorname{ubg\_}}^{*} :=$
\begin{equation*}
\{\operatorname{ubp\_} \in (\operatorname{vec}(\operatorname{ubv\_}) \times \operatorname{vec}(\operatorname{ubv\_})) : \forall \operatorname{ubq\_} \in \operatorname{ubg\_}.\; (\langle\cdot,\cdot\rangle_{\operatorname{ubv\_}}(\operatorname{nth}(1, \operatorname{ubq\_}), \operatorname{nth}(2, \operatorname{ubp\_})) = \langle\cdot,\cdot\rangle_{\operatorname{ubv\_}}(\operatorname{nth}(2, \operatorname{ubq\_}), \operatorname{nth}(1, \operatorname{ubp\_})))\}
\end{equation*}

\par\medskip\noindent\textbf{Definition} (\texttt{graph-inverse}).\;${\operatorname{ubg\_}}^{-1} :=$
\begin{equation*}
\{\operatorname{ubp\_} \in (\operatorname{vec}(\operatorname{ubv\_}) \times \operatorname{vec}(\operatorname{ubv\_})) : ({[}\operatorname{nth}(2, \operatorname{ubp\_}), \operatorname{nth}(1, \operatorname{ubp\_})] \in \operatorname{ubg\_})\}
\end{equation*}

\par\medskip\noindent\textbf{Definition} (\texttt{graph-closure}).\;$\overline{\operatorname{ubg\_}} :=$ $\operatorname{closure}(\operatorname{ip\text{-}metric\text{-}space}({\operatorname{ubv\_}} \oplus {\operatorname{ubv\_}}), \operatorname{ubg\_})$
\par\medskip\noindent\textbf{Definition} (\texttt{is-closed-graph}).\;For all $\operatorname{ubv\_}$ and $\operatorname{ubg\_}$, 
\begin{equation*}
\operatorname{is\text{-}closed\text{-}graph}(\operatorname{ubv\_}, \operatorname{ubg\_}) \Leftrightarrow \operatorname{is\text{-}closed}(\operatorname{ip\text{-}metric\text{-}space}({\operatorname{ubv\_}} \oplus {\operatorname{ubv\_}}), \operatorname{ubg\_})
\end{equation*}

\par\medskip\noindent\textbf{Definition} (\texttt{is-closable-graph}).\;For all $\operatorname{ubv\_}$ and $\operatorname{ubg\_}$, 
\begin{equation*}
\operatorname{is\text{-}closable\text{-}graph}(\operatorname{ubv\_}, \operatorname{ubg\_}) \Leftrightarrow \operatorname{is\text{-}functional\text{-}graph}(\operatorname{ubv\_}, \overline{\operatorname{ubg\_}})
\end{equation*}

\par\medskip\noindent\textbf{Definition} (\texttt{is-symmetric-graph}).\;For all $\operatorname{ubv\_}$ and $\operatorname{ubg\_}$, 
\begin{equation*}
\operatorname{is\text{-}symmetric\text{-}graph}(\operatorname{ubv\_}, \operatorname{ubg\_}) \Leftrightarrow (\operatorname{ubg\_} \subseteq {\operatorname{ubg\_}}^{*})
\end{equation*}

\par\medskip\noindent\textbf{Definition} (\texttt{is-self-adjoint-graph}).\;For all $\operatorname{ubv\_}$ and $\operatorname{ubg\_}$, 
\begin{equation*}
\operatorname{is\text{-}self\text{-}adjoint\text{-}graph}(\operatorname{ubv\_}, \operatorname{ubg\_}) \Leftrightarrow ({\operatorname{ubg\_}}^{*} = \operatorname{ubg\_})
\end{equation*}

\par\medskip\noindent\textbf{Definition} (\texttt{is-ess-self-adjoint-graph}).\;For all $\operatorname{ubv\_}$ and $\operatorname{ubg\_}$, 
\begin{equation*}
\operatorname{is\text{-}ess\text{-}self\text{-}adjoint\text{-}graph}(\operatorname{ubv\_}, \operatorname{ubg\_}) \Leftrightarrow \operatorname{is\text{-}self\text{-}adjoint\text{-}graph}(\operatorname{ubv\_}, \overline{\operatorname{ubg\_}})
\end{equation*}

\par\medskip\noindent\textbf{Definition} (\texttt{is-isometric-graph}).\;For all $\operatorname{ubv\_}$ and $\operatorname{ubg\_}$, 
\begin{equation*}
\begin{array}{@{}l@{}} \operatorname{is\text{-}isometric\text{-}graph}(\operatorname{ubv\_}, \operatorname{ubg\_}) \\ \Leftrightarrow\; \forall \operatorname{ubp\_} \in \operatorname{ubg\_}. \\ \quad \lVert \operatorname{nth}(1, \operatorname{ubp\_}) \rVert = \lVert \operatorname{nth}(2, \operatorname{ubp\_}) \rVert \end{array}
\end{equation*}

\par\medskip\noindent\textbf{Definition} (\texttt{is-unitary-graph}).\;For all $\operatorname{ubv\_}$ and $\operatorname{ubg\_}$, 
\begin{equation*}
\begin{array}{@{}l@{}} \operatorname{is\text{-}unitary\text{-}graph}(\operatorname{ubv\_}, \operatorname{ubg\_}) \\ \Leftrightarrow\; \operatorname{is\text{-}linear\text{-}graph}(\operatorname{ubv\_}, \operatorname{ubg\_}) \wedge \\ \operatorname{is\text{-}isometric\text{-}graph}(\operatorname{ubv\_}, \operatorname{ubg\_}) \wedge \\ \mathrm{dom}\,{\operatorname{ubg\_}} = \operatorname{vec}(\operatorname{ubv\_}) \wedge \\ \mathrm{ran}\,{\operatorname{ubg\_}} = \operatorname{vec}(\operatorname{ubv\_}) \end{array}
\end{equation*}

\par\medskip\noindent\textbf{Definition} (\texttt{cayley-graph}).\;$C({\operatorname{ubg\_}}) :=$
\begin{equation*}
\{\operatorname{ubp\_} \in (\operatorname{vec}(\operatorname{ubv\_}) \times \operatorname{vec}(\operatorname{ubv\_})) : \exists \operatorname{ubq\_} \in \operatorname{ubg\_}.\; (\operatorname{ubp\_} = {[}\operatorname{vadd}(\operatorname{ubv\_})(\operatorname{act}(\operatorname{ubv\_})(+i, \operatorname{nth}(1, \operatorname{ubq\_})), \operatorname{nth}(2, \operatorname{ubq\_})), \operatorname{vadd}(\operatorname{ubv\_})(\operatorname{act}(\operatorname{ubv\_})(-+i, \operatorname{nth}(1, \operatorname{ubq\_})), \operatorname{nth}(2, \operatorname{ubq\_}))])\}
\end{equation*}

\par\medskip\noindent\textbf{Definition} (\texttt{graph-compose}).\;${\operatorname{ubg\_}}{\operatorname{ubh\_}} :=$
\begin{equation*}
\{\operatorname{ubp\_} \in (\operatorname{vec}(\operatorname{ubv\_}) \times \operatorname{vec}(\operatorname{ubv\_})) : \exists \operatorname{uby\_} \in \operatorname{vec}(\operatorname{ubv\_}).\; (({[}\operatorname{nth}(1, \operatorname{ubp\_}), \operatorname{uby\_}] \in \operatorname{ubh\_}) \wedge ({[}\operatorname{uby\_}, \operatorname{nth}(2, \operatorname{ubp\_})] \in \operatorname{ubg\_}))\}
\end{equation*}

\par\medskip\noindent\textbf{Definition} (\texttt{is-quadratic-form}).\;For all $\operatorname{ubv\_}$, $\operatorname{ubd\_}$ and $\operatorname{ubf\_}$, 
\begin{equation*}
\begin{array}{@{}l@{}} \operatorname{is\text{-}quadratic\text{-}form}(\operatorname{ubv\_}, \operatorname{ubd\_}, \operatorname{ubf\_}) \\ \Leftrightarrow\; \operatorname{is\text{-}submodule}(\operatorname{ubv\_}, \operatorname{ubd\_}) \wedge \\ \operatorname{is\text{-}dense}(\operatorname{ip\text{-}metric\text{-}space}(\operatorname{ubv\_}), \operatorname{ubd\_}) \wedge \\ \operatorname{ubf\_} \in ((\operatorname{ubd\_} \times \operatorname{ubd\_}) \to \mathbb{C}) \wedge \\ \forall \operatorname{ubx\_}, \operatorname{uby\_}, \operatorname{ubz\_} \in \operatorname{ubd\_}. \\ \quad \operatorname{ubf\_}(\operatorname{vadd}(\operatorname{ubv\_})(\operatorname{ubx\_}, \operatorname{uby\_}), \operatorname{ubz\_}) = +(\operatorname{ubf\_}(\operatorname{ubx\_}, \operatorname{ubz\_}), \operatorname{ubf\_}(\operatorname{uby\_}, \operatorname{ubz\_})) \wedge \\ \forall \operatorname{ubk\_} \in \mathbb{C}, \operatorname{ubx\_}, \operatorname{uby\_} \in \operatorname{ubd\_}. \\ \quad \operatorname{ubf\_}(\operatorname{act}(\operatorname{ubv\_})(\operatorname{ubk\_}, \operatorname{ubx\_}), \operatorname{uby\_}) = \cdot(\operatorname{ubk\_}, \operatorname{ubf\_}(\operatorname{ubx\_}, \operatorname{uby\_})) \wedge \\ \forall \operatorname{ubx\_}, \operatorname{uby\_} \in \operatorname{ubd\_}. \\ \quad \operatorname{ubf\_}(\operatorname{uby\_}, \operatorname{ubx\_}) = \operatorname{conjugate}(\operatorname{ubf\_}(\operatorname{ubx\_}, \operatorname{uby\_})) \end{array}
\end{equation*}

\par\medskip\noindent\textbf{Definition} (\texttt{is-positive-form}).\;For all $\operatorname{ubv\_}$, $\operatorname{ubd\_}$ and $\operatorname{ubf\_}$, 
\begin{equation*}
\begin{array}{@{}l@{}} \operatorname{is\text{-}positive\text{-}form}(\operatorname{ubv\_}, \operatorname{ubd\_}, \operatorname{ubf\_}) \\ \Leftrightarrow\; \operatorname{is\text{-}quadratic\text{-}form}(\operatorname{ubv\_}, \operatorname{ubd\_}, \operatorname{ubf\_}) \wedge \\ \forall \operatorname{ubx\_} \in \operatorname{ubd\_}.\; ((\operatorname{ubf\_}(\operatorname{ubx\_}, \operatorname{ubx\_}) \in \mathbb{R}) \wedge (0 \leq \operatorname{ubf\_}(\operatorname{ubx\_}, \operatorname{ubx\_}))) \end{array}
\end{equation*}

\par\medskip\noindent\textbf{Definition} (\texttt{form-plus}).\;${\operatorname{ubf\_}}_{+} :=$
\begin{equation*}
(\lambda {[}\operatorname{ubx\_}, \operatorname{uby\_}] \in (\operatorname{ubd\_} \times \operatorname{ubd\_}).\; (\operatorname{ubf\_}(\operatorname{ubx\_}, \operatorname{uby\_}) + \langle\cdot,\cdot\rangle_{\operatorname{ubv\_}}(\operatorname{ubx\_}, \operatorname{uby\_})))
\end{equation*}

\par\medskip\noindent\textbf{Definition} (\texttt{is-closed-form}).\;For all $\operatorname{ubv\_}$, $\operatorname{ubd\_}$ and $\operatorname{ubf\_}$, 
\begin{equation*}
\begin{array}{@{}l@{}} \operatorname{is\text{-}closed\text{-}form}(\operatorname{ubv\_}, \operatorname{ubd\_}, \operatorname{ubf\_}) \\ \Leftrightarrow\; \operatorname{is\text{-}positive\text{-}form}(\operatorname{ubv\_}, \operatorname{ubd\_}, \operatorname{ubf\_}) \wedge \\ \forall \operatorname{ubq\_} \in (\mathbb{N} \to \operatorname{ubd\_}), \operatorname{ubx\_}. \\ \quad \operatorname{converges\text{-}to}(\operatorname{ip\text{-}metric\text{-}space}(\operatorname{ubv\_}), \operatorname{ubq\_}, \operatorname{ubx\_}) \Rightarrow  \\ \quad \forall \operatorname{ube\_}. \\ \quad \quad \operatorname{pos\text{-}rr}(\operatorname{ube\_}) \Rightarrow  \\ \quad \quad \quad \exists \operatorname{ubn\_} \in \mathbb{N}. \forall \operatorname{ubm\_}, \operatorname{ubk\_} \in \mathbb{N}. \\ \quad \quad \quad \quad (\operatorname{ubn\_} \leq \operatorname{ubm\_}) \Rightarrow  \\ \quad \quad \quad \quad (\operatorname{ubn\_} \leq \operatorname{ubk\_}) \Rightarrow  \\ \quad \quad \quad \quad \quad \operatorname{magnitude}(\operatorname{ubf\_}(\operatorname{vadd}(\operatorname{ubv\_})(\operatorname{ubq\_}(\operatorname{ubm\_}), \operatorname{vneg}(\operatorname{ubv\_})(\operatorname{ubq\_}(\operatorname{ubk\_}))), \\ \quad \quad \quad \quad \quad \quad \quad \operatorname{vadd}(\operatorname{ubv\_})(\operatorname{ubq\_}(\operatorname{ubm\_}), \operatorname{vneg}(\operatorname{ubv\_})(\operatorname{ubq\_}(\operatorname{ubk\_}))))) \\ \quad \quad \quad \quad \quad \leq\; \operatorname{ube\_} \Rightarrow \\ \quad \quad \operatorname{ubx\_} \in \operatorname{ubd\_} \wedge \\ \quad \quad \forall \operatorname{ube\_}. \\ \quad \quad \quad \operatorname{pos\text{-}rr}(\operatorname{ube\_}) \Rightarrow  \\ \quad \quad \quad \quad \exists \operatorname{ubn\_} \in \mathbb{N}. \forall \operatorname{ubm\_} \in \mathbb{N}. \\ \quad \quad \quad \quad \quad (\operatorname{ubn\_} \leq \operatorname{ubm\_}) \Rightarrow  \\ \quad \quad \quad \quad \quad \quad \operatorname{magnitude}(\operatorname{ubf\_}(\operatorname{vadd}(\operatorname{ubv\_})(\operatorname{ubq\_}(\operatorname{ubm\_}), \operatorname{vneg}(\operatorname{ubv\_})(\operatorname{ubx\_})), \operatorname{vadd}(\operatorname{ubv\_})(\operatorname{ubq\_}(\operatorname{ubm\_}), \operatorname{vneg}(\operatorname{ubv\_})(\operatorname{ubx\_})))) \\ \quad \quad \quad \quad \quad \quad \leq\; \operatorname{ube\_} \end{array}
\end{equation*}


\vnbbookfile{theorem-library/ba-bh}{theorem-library/ba-bh}
\noindent{\small ba-bh.scm -- wave of 2026-11-09/10 (agent BA-BH), wired 2026-11-10 from scratchpad/babh-1109/ba-bh.scm. Library: operator-algebras.}\par
\vnbprop{proj-op-is-orthogonal-projection}{\vnbstmt{\begin{array}{@{}l@{}} \forall v, k. \\ \quad (\operatorname{is\text{-}hilbert\text{-}space}(v) \wedge (\operatorname{is\text{-}submodule}(v, k) \wedge \operatorname{is\text{-}closed}(\operatorname{ip\text{-}metric\text{-}space}(v), k))) \Rightarrow  \\ \quad \quad \operatorname{is\text{-}orthogonal\text{-}projection}(v, k, E_{k}) \end{array}}}
\vnbprop{proj-op-unique}{\vnbstmt{\begin{array}{@{}l@{}} \forall v, k, a. \\ \quad \operatorname{is\text{-}hilbert\text{-}space}(v) \wedge \\ \quad \operatorname{is\text{-}submodule}(v, k) \wedge \\ \quad \operatorname{is\text{-}closed}(\operatorname{ip\text{-}metric\text{-}space}(v), k) \wedge \\ \quad \operatorname{is\text{-}orthogonal\text{-}projection}(v, k, a) \Rightarrow \\ \quad \quad E_{k} = a \end{array}}}
\vnbprop{proj-op-linear}{\vnbstmt{\begin{array}{@{}l@{}} \forall v, k. \\ \quad (\operatorname{is\text{-}hilbert\text{-}space}(v) \wedge (\operatorname{is\text{-}submodule}(v, k) \wedge \operatorname{is\text{-}closed}(\operatorname{ip\text{-}metric\text{-}space}(v), k))) \Rightarrow  \\ \quad \quad \operatorname{is\text{-}linear\text{-}op}(v, v, E_{k}) \end{array}}}
\vnbprop{minimizer-orthogonal}{\vnbstmt{\begin{array}{@{}l@{}} \forall v, k, i, a. \\ \quad \operatorname{is\text{-}complex\text{-}inner\text{-}product\text{-}space}(v) \wedge \\ \quad \operatorname{is\text{-}submodule}(v, k) \wedge \\ \quad i \in \operatorname{vec}(v) \wedge \\ \quad a \in k \wedge \\ \quad \forall p \in k. \\ \quad \quad \langle\cdot,\cdot\rangle_{v}(\operatorname{vadd}(v)(a, \operatorname{vneg}(v)(i)), \operatorname{vadd}(v)(a, \operatorname{vneg}(v)(i))) \\ \quad \quad \leq\; \langle\cdot,\cdot\rangle_{v}(\operatorname{vadd}(v)(p, \operatorname{vneg}(v)(i)), \operatorname{vadd}(v)(p, \operatorname{vneg}(v)(i))) \Rightarrow \\ \quad \quad \operatorname{vadd}(v)(i, \operatorname{vneg}(v)(a)) \in {k}^{\perp} \end{array}}}
\vnbprop{orthogonal-point-unique}{\vnbstmt{\begin{array}{@{}l@{}} \forall v, k, i, a_{1}, a_{2}. \\ \quad \operatorname{is\text{-}complex\text{-}inner\text{-}product\text{-}space}(v) \wedge \\ \quad \operatorname{is\text{-}submodule}(v, k) \wedge \\ \quad i \in \operatorname{vec}(v) \wedge \\ \quad a_{1} \in k \wedge \\ \quad a_{2} \in k \wedge \\ \quad \operatorname{vadd}(v)(i, \operatorname{vneg}(v)(a_{1})) \in {k}^{\perp} \wedge \\ \quad \operatorname{vadd}(v)(i, \operatorname{vneg}(v)(a_{2})) \in {k}^{\perp} \Rightarrow \\ \quad \quad a_{1} = a_{2} \end{array}}}
\vnbprop{op-range-unfold}{\vnbstmt{\begin{array}{@{}l@{}} \forall v, w, t. \\ \quad (\mathrm{ran}\,{t} \simeq \{y \in \operatorname{vec}(w) : \exists x \in \operatorname{vec}(v).\; (y = t(x))\}) \end{array}}}
\vnbprop{op-range-membership}{\vnbstmt{\begin{array}{@{}l@{}} \forall v, w, t, y. \\ \quad (y \in \mathrm{ran}\,{t}) \Leftrightarrow ((y \in \operatorname{vec}(w)) \wedge \exists x \in \operatorname{vec}(v).\; (y = t(x))) \end{array}}}
\vnbprop{proj-op-range}{\vnbstmt{\begin{array}{@{}l@{}} \forall v, k. \\ \quad (\operatorname{is\text{-}hilbert\text{-}space}(v) \wedge (\operatorname{is\text{-}submodule}(v, k) \wedge \operatorname{is\text{-}closed}(\operatorname{ip\text{-}metric\text{-}space}(v), k))) \Rightarrow  \\ \quad \quad \mathrm{ran}\,{E_{k}} = k \end{array}}}
\vnbprop{perp-perp-closed}{\vnbstmt{\begin{array}{@{}l@{}} \forall v, k. \\ \quad (\operatorname{is\text{-}hilbert\text{-}space}(v) \wedge (\operatorname{is\text{-}submodule}(v, k) \wedge \operatorname{is\text{-}closed}(\operatorname{ip\text{-}metric\text{-}space}(v), k))) \Rightarrow  \\ \quad \quad {{k}^{\perp}}^{\perp} = k \end{array}}}
\vnbprop{ip-norm-in-rr}{\vnbstmt{\begin{array}{@{}l@{}} \forall v, x. \\ \quad (\operatorname{is\text{-}complex\text{-}inner\text{-}product\text{-}space}(v) \wedge (x \in \operatorname{vec}(v))) \Rightarrow (\lVert x \rVert \in \mathbb{R}) \end{array}}}
\vnbprop{ip-norm-nonneg}{\vnbstmt{\begin{array}{@{}l@{}} \forall v, x. \\ \quad (\operatorname{is\text{-}complex\text{-}inner\text{-}product\text{-}space}(v) \wedge (x \in \operatorname{vec}(v))) \Rightarrow (0 \leq \lVert x \rVert) \end{array}}}
\vnbprop{ip-norm-sq}{\vnbstmt{\begin{array}{@{}l@{}} \forall v, x. \\ \quad (\operatorname{is\text{-}complex\text{-}inner\text{-}product\text{-}space}(v) \wedge (x \in \operatorname{vec}(v))) \Rightarrow  \\ \quad \quad (\lVert x \rVert \cdot \lVert x \rVert) = \langle\cdot,\cdot\rangle_{v}(x, x) \end{array}}}
\vnbprop{ip-norm-zero-vector}{\vnbstmt{\begin{array}{@{}l@{}} \forall v, x. \\ \quad (\operatorname{is\text{-}complex\text{-}inner\text{-}product\text{-}space}(v) \wedge ((x \in \operatorname{vec}(v)) \wedge (\lVert x \rVert = 0))) \Rightarrow  \\ \quad \quad x = \operatorname{vzero}(v) \end{array}}}
\vnbprop{ip-norm-zero}{\vnbstmt{\forall v.\; (\operatorname{is\text{-}complex\text{-}inner\text{-}product\text{-}space}(v) \Rightarrow (\lVert \operatorname{vzero}(v) \rVert = 0))}}
\vnbprop{ip-norm-act}{\vnbstmt{\begin{array}{@{}l@{}} \forall v, c, x. \\ \quad (\operatorname{is\text{-}complex\text{-}inner\text{-}product\text{-}space}(v) \wedge ((c \in \mathbb{C}) \wedge (x \in \operatorname{vec}(v)))) \Rightarrow  \\ \quad \quad \lVert \operatorname{act}(v)(c, x) \rVert = (\operatorname{magnitude}(c) \cdot \lVert x \rVert) \end{array}}}
\vnbprop{bounded-ops-unfold}{\vnbstmt{\begin{array}{@{}l@{}} \forall v, w. \\ \quad (L({v}, {w}) \simeq \{t \in (\operatorname{vec}(v) \to \operatorname{vec}(w)) : \operatorname{is\text{-}bounded\text{-}op}(v, w, t)\}) \end{array}}}
\vnbprop{bounded-ops-membership}{\vnbstmt{\begin{array}{@{}l@{}} \forall v, w, t. \\ \quad (t \in L({v}, {w})) \Leftrightarrow ((t \in (\operatorname{vec}(v) \to \operatorname{vec}(w))) \wedge \operatorname{is\text{-}bounded\text{-}op}(v, w, t)) \end{array}}}
\vnbprop{bop-fun}{\vnbstmt{\forall v, w.; \forall t \in L({v}, {w}).\; (t \in (\operatorname{vec}(v) \to \operatorname{vec}(w)))}}
\vnbprop{bop-linear}{\vnbstmt{\forall v, w.; \forall t \in L({v}, {w}).\; \operatorname{is\text{-}linear\text{-}op}(v, w, t)}}
\vnbprop{bop-bound-exists}{\vnbstmt{\begin{array}{@{}l@{}} \forall v, w.; \forall t \in L({v}, {w}). \exists m \in \mathbb{R}. \\ \quad (0 \leq m) \wedge \forall x \in \operatorname{vec}(v).\; (\lVert t(x) \rVert \leq (m \cdot \lVert x \rVert)) \end{array}}}
\vnbprop{bop-add}{\vnbstmt{\begin{array}{@{}l@{}} \forall v, w, t, x, y. \\ \quad ((t \in L({v}, {w})) \wedge ((x \in \operatorname{vec}(v)) \wedge (y \in \operatorname{vec}(v)))) \Rightarrow  \\ \quad \quad t(\operatorname{vadd}(v)(x, y)) = \operatorname{vadd}(w)(t(x), t(y)) \end{array}}}
\vnbprop{bop-act}{\vnbstmt{\begin{array}{@{}l@{}} \forall v, w, t, c, x. \\ \quad ((t \in L({v}, {w})) \wedge ((c \in \mathbb{C}) \wedge (x \in \operatorname{vec}(v)))) \Rightarrow  \\ \quad \quad t(\operatorname{act}(v)(c, x)) = \operatorname{act}(w)(c, t(x)) \end{array}}}
\vnbprop{bop-intro}{\vnbstmt{\begin{array}{@{}l@{}} \forall v, w, t, m. \\ \quad \operatorname{is\text{-}linear\text{-}op}(v, w, t) \wedge \\ \quad m \in \mathbb{R} \wedge \\ \quad 0 \leq m \wedge \\ \quad \forall x \in \operatorname{vec}(v).\; (\lVert t(x) \rVert \leq (m \cdot \lVert x \rVert)) \Rightarrow \\ \quad \quad t \in L({v}, {w}) \end{array}}}
\vnbprop{opnorm-spec}{\vnbstmt{\begin{array}{@{}l@{}} \forall v, w, t. \\ \quad (\operatorname{is\text{-}complex\text{-}inner\text{-}product\text{-}space}(v) \wedge (\operatorname{is\text{-}complex\text{-}inner\text{-}product\text{-}space}(w) \wedge (t \in L({v}, {w})))) \Rightarrow  \\ \quad \quad \lVert {t} \rVert \in \mathbb{R} \wedge \\ \quad \quad 0 \leq \lVert {t} \rVert \wedge \\ \quad \quad \forall x \in \operatorname{vec}(v).\; (\lVert t(x) \rVert \leq (\lVert {t} \rVert \cdot \lVert x \rVert)) \wedge \\ \quad \quad \forall d \in \mathbb{R}. \\ \quad \quad \quad ((0 \leq d) \wedge \forall x \in \operatorname{vec}(v).\; (\lVert t(x) \rVert \leq (d \cdot \lVert x \rVert))) \Rightarrow  \\ \quad \quad \quad \quad \lVert {t} \rVert \leq d \end{array}}}
\vnbprop{opnorm-in-rr}{\vnbstmt{\begin{array}{@{}l@{}} \forall v, w, t. \\ \quad (\operatorname{is\text{-}complex\text{-}inner\text{-}product\text{-}space}(v) \wedge (\operatorname{is\text{-}complex\text{-}inner\text{-}product\text{-}space}(w) \wedge (t \in L({v}, {w})))) \Rightarrow  \\ \quad \quad \lVert {t} \rVert \in \mathbb{R} \end{array}}}
\vnbprop{opnorm-nonneg}{\vnbstmt{\begin{array}{@{}l@{}} \forall v, w, t. \\ \quad (\operatorname{is\text{-}complex\text{-}inner\text{-}product\text{-}space}(v) \wedge (\operatorname{is\text{-}complex\text{-}inner\text{-}product\text{-}space}(w) \wedge (t \in L({v}, {w})))) \Rightarrow  \\ \quad \quad 0 \leq \lVert {t} \rVert \end{array}}}
\vnbprop{opnorm-bound}{\vnbstmt{\begin{array}{@{}l@{}} \forall v, w, t.; \forall x \in \operatorname{vec}(v). \\ \quad (\operatorname{is\text{-}complex\text{-}inner\text{-}product\text{-}space}(v) \wedge (\operatorname{is\text{-}complex\text{-}inner\text{-}product\text{-}space}(w) \wedge (t \in L({v}, {w})))) \Rightarrow  \\ \quad \quad \lVert t(x) \rVert \leq (\lVert {t} \rVert \cdot \lVert x \rVert) \end{array}}}
\vnbprop{cc-functional-continuous-of-bound}{\vnbstmt{\begin{array}{@{}l@{}} \forall v, f, m. \\ \quad \operatorname{is\text{-}complex\text{-}inner\text{-}product\text{-}space}(v) \wedge \\ \quad \operatorname{is\text{-}cc\text{-}linear\text{-}functional}(v, f) \wedge \\ \quad m \in \mathbb{R} \wedge \\ \quad 0 \leq m \wedge \\ \quad \forall x \in \operatorname{vec}(v).\; (\operatorname{magnitude}(f(x)) \leq (m \cdot \lVert x \rVert)) \Rightarrow \\ \quad \quad \operatorname{is\text{-}continuous}(\operatorname{ip\text{-}metric\text{-}space}(v), \operatorname{cc\text{-}ms}, f) \end{array}}}
\vnbprop{bop-ip-bound}{\vnbstmt{\begin{array}{@{}l@{}} \forall v, w, t, m, y.; \forall x \in \operatorname{vec}(v). \\ \quad \operatorname{is\text{-}complex\text{-}inner\text{-}product\text{-}space}(v) \wedge \\ \quad \operatorname{is\text{-}complex\text{-}inner\text{-}product\text{-}space}(w) \wedge \\ \quad t \in (\operatorname{vec}(v) \to \operatorname{vec}(w)) \wedge \\ \quad m \in \mathbb{R} \wedge \\ \quad 0 \leq m \wedge \\ \quad \forall \mathit{box} \in \operatorname{vec}(v).\; (\lVert t(\mathit{box}) \rVert \leq (m \cdot \lVert \mathit{box} \rVert)) \wedge \\ \quad y \in \operatorname{vec}(w) \Rightarrow \\ \quad \quad \operatorname{magnitude}(\langle\cdot,\cdot\rangle_{w}(t(x), y)) \leq ((m \cdot \lVert y \rVert) \cdot \lVert x \rVert) \end{array}}}
\vnbprop{bop-ip-functional-linear}{\vnbstmt{\begin{array}{@{}l@{}} \forall v, w, t, y. \\ \quad \operatorname{is\text{-}complex\text{-}inner\text{-}product\text{-}space}(v) \wedge \\ \quad \operatorname{is\text{-}complex\text{-}inner\text{-}product\text{-}space}(w) \wedge \\ \quad t \in L({v}, {w}) \wedge \\ \quad y \in \operatorname{vec}(w) \Rightarrow \\ \quad \quad \operatorname{is\text{-}cc\text{-}linear\text{-}functional}(v, (\lambda x \in \operatorname{vec}(v).\; \langle\cdot,\cdot\rangle_{w}(t(x), y))) \end{array}}}
\vnbprop{bop-ip-functional-continuous}{\vnbstmt{\begin{array}{@{}l@{}} \forall v, w, t, y. \\ \quad \operatorname{is\text{-}complex\text{-}inner\text{-}product\text{-}space}(v) \wedge \\ \quad \operatorname{is\text{-}complex\text{-}inner\text{-}product\text{-}space}(w) \wedge \\ \quad t \in L({v}, {w}) \wedge \\ \quad y \in \operatorname{vec}(w) \Rightarrow \\ \quad \quad \operatorname{is\text{-}continuous}(\operatorname{ip\text{-}metric\text{-}space}(v), \operatorname{cc\text{-}ms}, (\lambda x \in \operatorname{vec}(v).\; \langle\cdot,\cdot\rangle_{w}(t(x), y))) \end{array}}}
\vnbprop{op-adjoint-unfold}{\vnbstmt{\begin{array}{@{}l@{}} \forall v, w, t. \\ \quad ({t}^{*} \simeq (\lambda y \in \operatorname{vec}(w).\; (\iota\, z.\;((z \in \operatorname{vec}(v)) \wedge \forall x \in \operatorname{vec}(v).\; (\langle\cdot,\cdot\rangle_{w}(t(x), y) = \langle\cdot,\cdot\rangle_{v}(x, z)))))) \end{array}}}
\vnbprop{op-adjoint-spec}{\vnbstmt{\begin{array}{@{}l@{}} \forall v, w, t, y. \\ \quad (\operatorname{is\text{-}hilbert\text{-}space}(v) \wedge (\operatorname{is\text{-}complex\text{-}inner\text{-}product\text{-}space}(w) \wedge ((t \in L({v}, {w})) \wedge (y \in \operatorname{vec}(w))))) \Rightarrow  \\ \quad \quad {t}^{*}(y) \in \operatorname{vec}(v) \wedge \\ \quad \quad \forall x \in \operatorname{vec}(v).\; (\langle\cdot,\cdot\rangle_{w}(t(x), y) = \langle\cdot,\cdot\rangle_{v}(x, {t}^{*}(y))) \end{array}}}
\vnbprop{op-adjoint-type}{\vnbstmt{\begin{array}{@{}l@{}} \forall v, w, t, y. \\ \quad (\operatorname{is\text{-}hilbert\text{-}space}(v) \wedge (\operatorname{is\text{-}complex\text{-}inner\text{-}product\text{-}space}(w) \wedge ((t \in L({v}, {w})) \wedge (y \in \operatorname{vec}(w))))) \Rightarrow  \\ \quad \quad {t}^{*}(y) \in \operatorname{vec}(v) \end{array}}}
\vnbprop{op-adjoint-ip}{\vnbstmt{\begin{array}{@{}l@{}} \forall v, w, t, x, y. \\ \quad \operatorname{is\text{-}hilbert\text{-}space}(v) \wedge \\ \quad \operatorname{is\text{-}complex\text{-}inner\text{-}product\text{-}space}(w) \wedge \\ \quad t \in L({v}, {w}) \wedge \\ \quad x \in \operatorname{vec}(v) \wedge \\ \quad y \in \operatorname{vec}(w) \Rightarrow \\ \quad \quad \langle\cdot,\cdot\rangle_{w}(t(x), y) = \langle\cdot,\cdot\rangle_{v}(x, {t}^{*}(y)) \end{array}}}
\vnbprop{op-adjoint-fun}{\vnbstmt{\begin{array}{@{}l@{}} \forall v, w, t. \\ \quad (\operatorname{is\text{-}hilbert\text{-}space}(v) \wedge (\operatorname{is\text{-}complex\text{-}inner\text{-}product\text{-}space}(w) \wedge (t \in L({v}, {w})))) \Rightarrow  \\ \quad \quad {t}^{*} \in (\operatorname{vec}(w) \to \operatorname{vec}(v)) \end{array}}}
\vnbprop{op-adjoint-unique}{\vnbstmt{\begin{array}{@{}l@{}} \forall v, w, t, g. \\ \quad \operatorname{is\text{-}hilbert\text{-}space}(v) \wedge \\ \quad \operatorname{is\text{-}complex\text{-}inner\text{-}product\text{-}space}(w) \wedge \\ \quad t \in L({v}, {w}) \wedge \\ \quad g \in (\operatorname{vec}(w) \to \operatorname{vec}(v)) \wedge \\ \quad \forall x \in \operatorname{vec}(v), y \in \operatorname{vec}(w).\; (\langle\cdot,\cdot\rangle_{w}(t(x), y) = \langle\cdot,\cdot\rangle_{v}(x, g(y))) \Rightarrow \\ \quad \quad {t}^{*} = g \end{array}}}
\vnbprop{op-adjoint-bound}{\vnbstmt{\begin{array}{@{}l@{}} \forall v, w, t, m, y. \\ \quad \operatorname{is\text{-}hilbert\text{-}space}(v) \wedge \\ \quad \operatorname{is\text{-}complex\text{-}inner\text{-}product\text{-}space}(w) \wedge \\ \quad t \in L({v}, {w}) \wedge \\ \quad m \in \mathbb{R} \wedge \\ \quad 0 \leq m \wedge \\ \quad \forall x \in \operatorname{vec}(v).\; (\lVert t(x) \rVert \leq (m \cdot \lVert x \rVert)) \wedge \\ \quad y \in \operatorname{vec}(w) \Rightarrow \\ \quad \quad \lVert {t}^{*}(y) \rVert \leq (m \cdot \lVert y \rVert) \end{array}}}
\vnbprop{op-adjoint-linear}{\vnbstmt{\begin{array}{@{}l@{}} \forall v, w, t. \\ \quad (\operatorname{is\text{-}hilbert\text{-}space}(v) \wedge (\operatorname{is\text{-}complex\text{-}inner\text{-}product\text{-}space}(w) \wedge (t \in L({v}, {w})))) \Rightarrow  \\ \quad \quad \operatorname{is\text{-}linear\text{-}op}(w, v, {t}^{*}) \end{array}}}
\vnbprop{op-adjoint-bounded}{\vnbstmt{\begin{array}{@{}l@{}} \forall v, w, t. \\ \quad (\operatorname{is\text{-}hilbert\text{-}space}(v) \wedge (\operatorname{is\text{-}complex\text{-}inner\text{-}product\text{-}space}(w) \wedge (t \in L({v}, {w})))) \Rightarrow  \\ \quad \quad {t}^{*} \in L({w}, {v}) \end{array}}}
\vnbprop{op-adjoint-opnorm-le}{\vnbstmt{\begin{array}{@{}l@{}} \forall v, w, t. \\ \quad (\operatorname{is\text{-}hilbert\text{-}space}(v) \wedge (\operatorname{is\text{-}complex\text{-}inner\text{-}product\text{-}space}(w) \wedge (t \in L({v}, {w})))) \Rightarrow  \\ \quad \quad \lVert {{t}^{*}} \rVert \leq \lVert {t} \rVert \end{array}}}
\vnbprop{op-adjoint-involutive}{\vnbstmt{\begin{array}{@{}l@{}} \forall v, w, t. \\ \quad (\operatorname{is\text{-}hilbert\text{-}space}(v) \wedge (\operatorname{is\text{-}hilbert\text{-}space}(w) \wedge (t \in L({v}, {w})))) \Rightarrow  \\ \quad \quad {{t}^{*}}^{*} = t \end{array}}}
\vnbprop{bop-comp-apply}{\vnbstmt{\begin{array}{@{}l@{}} \forall v, s, t, x. \\ \quad ((s \in (\operatorname{vec}(v) \to \operatorname{vec}(v))) \wedge ((t \in (\operatorname{vec}(v) \to \operatorname{vec}(v))) \wedge (x \in \operatorname{vec}(v)))) \Rightarrow  \\ \quad \quad {s}{t}(x) = s(t(x)) \end{array}}}
\vnbprop{bop-id-apply}{\vnbstmt{\forall v.; \forall x \in \operatorname{vec}(v).\; (1(x) = x)}}
\vnbprop{bop-comp-bounded}{\vnbstmt{\begin{array}{@{}l@{}} \forall v, s, t. \\ \quad (\operatorname{is\text{-}complex\text{-}inner\text{-}product\text{-}space}(v) \wedge ((s \in L({v}, {v})) \wedge (t \in L({v}, {v})))) \Rightarrow  \\ \quad \quad {s}{t} \in L({v}, {v}) \end{array}}}
\vnbprop{bop-id-bounded}{\vnbstmt{\begin{array}{@{}l@{}} \forall v. \\ \quad \operatorname{is\text{-}complex\text{-}inner\text{-}product\text{-}space}(v) \Rightarrow (1 \in L({v}, {v})) \end{array}}}
\vnbprop{op-adjoint-comp}{\vnbstmt{\begin{array}{@{}l@{}} \forall v, s, t. \\ \quad (\operatorname{is\text{-}hilbert\text{-}space}(v) \wedge ((s \in L({v}, {v})) \wedge (t \in L({v}, {v})))) \Rightarrow  \\ \quad \quad {{s}{t}}^{*} = {{t}^{*}}{{s}^{*}} \end{array}}}
\vnbprop{op-adjoint-id}{\vnbstmt{\forall v.\; (\operatorname{is\text{-}hilbert\text{-}space}(v) \Rightarrow ({1}^{*} = 1))}}
\vnbprop{opnorm-comp-le}{\vnbstmt{\begin{array}{@{}l@{}} \forall v, s, t. \\ \quad (\operatorname{is\text{-}complex\text{-}inner\text{-}product\text{-}space}(v) \wedge ((s \in L({v}, {v})) \wedge (t \in L({v}, {v})))) \Rightarrow  \\ \quad \quad \lVert {{s}{t}} \rVert \leq (\lVert {s} \rVert \cdot \lVert {t} \rVert) \end{array}}}
\vnbprop{bh-c-star-identity}{\vnbstmt{\begin{array}{@{}l@{}} \forall v, t. \\ \quad (\operatorname{is\text{-}hilbert\text{-}space}(v) \wedge (t \in L({v}, {v}))) \Rightarrow  \\ \quad \quad \lVert {{{t}^{*}}{t}} \rVert = (\lVert {t} \rVert \cdot \lVert {t} \rVert) \end{array}}}
\vnbprop{bop-vneg}{\vnbstmt{\begin{array}{@{}l@{}} \forall v, w, t, x. \\ \quad \operatorname{is\text{-}complex\text{-}inner\text{-}product\text{-}space}(v) \wedge \\ \quad \operatorname{is\text{-}complex\text{-}inner\text{-}product\text{-}space}(w) \wedge \\ \quad t \in L({v}, {w}) \wedge \\ \quad x \in \operatorname{vec}(v) \Rightarrow \\ \quad \quad t(\operatorname{vneg}(v)(x)) = \operatorname{vneg}(w)(t(x)) \end{array}}}
\vnbprop{op-kernel-unfold}{\vnbstmt{\forall v, w, t.\; (\ker {t} \simeq \{x \in \operatorname{vec}(v) : (t(x) = \operatorname{vzero}(w))\})}}
\vnbprop{op-kernel-membership}{\vnbstmt{\begin{array}{@{}l@{}} \forall v, w, t, y. \\ \quad (y \in \ker {t}) \Leftrightarrow ((y \in \operatorname{vec}(v)) \wedge (t(y) = \operatorname{vzero}(w))) \end{array}}}
\vnbprop{op-kernel-subset}{\vnbstmt{\begin{array}{@{}l@{}} \forall v, w, t. \\ \quad \operatorname{is\text{-}complex\text{-}inner\text{-}product\text{-}space}(v) \Rightarrow (\ker {t} \subseteq \operatorname{vec}(v)) \end{array}}}
\vnbprop{projection-op-symmetric}{\vnbstmt{\begin{array}{@{}l@{}} \forall v, e, x, y. \\ \quad (\operatorname{is\text{-}hilbert\text{-}space}(v) \wedge (\operatorname{is\text{-}projection\text{-}op}(v, e) \wedge ((x \in \operatorname{vec}(v)) \wedge (y \in \operatorname{vec}(v))))) \Rightarrow  \\ \quad \quad \langle\cdot,\cdot\rangle_{v}(e(x), y) = \langle\cdot,\cdot\rangle_{v}(x, e(y)) \end{array}}}
\vnbprop{projection-op-idempotent}{\vnbstmt{\begin{array}{@{}l@{}} \forall v, e, x. \\ \quad (\operatorname{is\text{-}hilbert\text{-}space}(v) \wedge (\operatorname{is\text{-}projection\text{-}op}(v, e) \wedge (x \in \operatorname{vec}(v)))) \Rightarrow  \\ \quad \quad e(e(x)) = e(x) \end{array}}}
\vnbprop{projection-op-kills-complement}{\vnbstmt{\begin{array}{@{}l@{}} \forall v, e, x. \\ \quad (\operatorname{is\text{-}hilbert\text{-}space}(v) \wedge (\operatorname{is\text{-}projection\text{-}op}(v, e) \wedge (x \in \operatorname{vec}(v)))) \Rightarrow  \\ \quad \quad e(\operatorname{vadd}(v)(x, \operatorname{vneg}(v)(e(x)))) = \operatorname{vzero}(v) \end{array}}}
\vnbprop{projection-op-fixes-range}{\vnbstmt{\begin{array}{@{}l@{}} \forall v, e, k. \\ \quad (\operatorname{is\text{-}hilbert\text{-}space}(v) \wedge (\operatorname{is\text{-}projection\text{-}op}(v, e) \wedge (k \in \mathrm{ran}\,{e}))) \Rightarrow  \\ \quad \quad e(k) = k \end{array}}}
\vnbprop{projection-op-range-perp-kernel}{\vnbstmt{\begin{array}{@{}l@{}} \forall v, e. \\ \quad (\operatorname{is\text{-}hilbert\text{-}space}(v) \wedge \operatorname{is\text{-}projection\text{-}op}(v, e)) \Rightarrow  \\ \quad \quad \mathrm{ran}\,{e} = {\ker {e}}^{\perp} \end{array}}}
\vnbprop{projection-op-range-closed}{\vnbstmt{\begin{array}{@{}l@{}} \forall v, e. \\ \quad (\operatorname{is\text{-}hilbert\text{-}space}(v) \wedge \operatorname{is\text{-}projection\text{-}op}(v, e)) \Rightarrow  \\ \quad \quad (\operatorname{is\text{-}submodule}(v, \mathrm{ran}\,{e}) \wedge \operatorname{is\text{-}closed}(\operatorname{ip\text{-}metric\text{-}space}(v), \mathrm{ran}\,{e})) \end{array}}}
\vnbprop{projection-op-is-proj-op-of-range}{\vnbstmt{\begin{array}{@{}l@{}} \forall v, e. \\ \quad (\operatorname{is\text{-}hilbert\text{-}space}(v) \wedge \operatorname{is\text{-}projection\text{-}op}(v, e)) \Rightarrow  \\ \quad \quad E_{\mathrm{ran}\,{e}} = e \end{array}}}
\vnbprop{bop-zero-value}{\vnbstmt{\begin{array}{@{}l@{}} \forall v, w, t. \\ \quad (\operatorname{is\text{-}complex\text{-}inner\text{-}product\text{-}space}(v) \wedge (\operatorname{is\text{-}complex\text{-}inner\text{-}product\text{-}space}(w) \wedge (t \in L({v}, {w})))) \Rightarrow  \\ \quad \quad t(\operatorname{vzero}(v)) = \operatorname{vzero}(w) \end{array}}}
\vnbprop{adjoint-square-ip}{\vnbstmt{\begin{array}{@{}l@{}} \forall v, u, x, y. \\ \quad (\operatorname{is\text{-}hilbert\text{-}space}(v) \wedge ((u \in L({v}, {v})) \wedge ((x \in \operatorname{vec}(v)) \wedge (y \in \operatorname{vec}(v))))) \Rightarrow  \\ \quad \quad \langle\cdot,\cdot\rangle_{v}(x, {{u}^{*}}{u}(y)) = \langle\cdot,\cdot\rangle_{v}(u(x), u(y)) \end{array}}}
\vnbprop{adjoint-square-kernel}{\vnbstmt{\begin{array}{@{}l@{}} \forall v, u, x. \\ \quad \operatorname{is\text{-}hilbert\text{-}space}(v) \wedge \\ \quad u \in L({v}, {v}) \wedge \\ \quad x \in \operatorname{vec}(v) \wedge \\ \quad {{u}^{*}}{u}(x) = \operatorname{vzero}(v) \Rightarrow \\ \quad \quad u(x) = \operatorname{vzero}(v) \end{array}}}
\vnbprop{adjoint-square-projection-absorbs}{\vnbstmt{\begin{array}{@{}l@{}} \forall v, u, x. \\ \quad \operatorname{is\text{-}hilbert\text{-}space}(v) \wedge \\ \quad u \in L({v}, {v}) \wedge \\ \quad \operatorname{is\text{-}projection\text{-}op}(v, {{u}^{*}}{u}) \wedge \\ \quad x \in \operatorname{vec}(v) \Rightarrow \\ \quad \quad u({{u}^{*}}{u}(x)) = u(x) \end{array}}}
\vnbprop{partial-isometry-of-projection}{\vnbstmt{\begin{array}{@{}l@{}} \forall v, u. \\ \quad \operatorname{is\text{-}hilbert\text{-}space}(v) \wedge \\ \quad u \in L({v}, {v}) \wedge \\ \quad \operatorname{is\text{-}projection\text{-}op}(v, {{u}^{*}}{u}) \Rightarrow \\ \quad \quad \operatorname{is\text{-}partial\text{-}isometry}(v, u) \end{array}}}
\vnbprop{adjoint-square-projection-swap}{\vnbstmt{\begin{array}{@{}l@{}} \forall v, u. \\ \quad \operatorname{is\text{-}hilbert\text{-}space}(v) \wedge \\ \quad u \in L({v}, {v}) \wedge \\ \quad \operatorname{is\text{-}projection\text{-}op}(v, {{u}^{*}}{u}) \Rightarrow \\ \quad \quad \operatorname{is\text{-}projection\text{-}op}(v, {u}{{u}^{*}}) \end{array}}}
\vnbprop{adjoint-square-projection-swap-back}{\vnbstmt{\begin{array}{@{}l@{}} \forall v, u. \\ \quad \operatorname{is\text{-}hilbert\text{-}space}(v) \wedge \\ \quad u \in L({v}, {v}) \wedge \\ \quad \operatorname{is\text{-}projection\text{-}op}(v, {u}{{u}^{*}}) \Rightarrow \\ \quad \quad \operatorname{is\text{-}projection\text{-}op}(v, {{u}^{*}}{u}) \end{array}}}
\vnbprop{support-projection-of-partial-isometry}{\vnbstmt{\begin{array}{@{}l@{}} \forall v, u. \\ \quad \operatorname{is\text{-}hilbert\text{-}space}(v) \wedge \\ \quad u \in L({v}, {v}) \wedge \\ \quad \operatorname{is\text{-}projection\text{-}op}(v, {{u}^{*}}{u}) \Rightarrow \\ \quad \quad E_{\mathrm{supp}\,{u}} = {{u}^{*}}{u} \end{array}}}
\vnbprop{range-projection-of-partial-isometry}{\vnbstmt{\begin{array}{@{}l@{}} \forall v, u. \\ \quad \operatorname{is\text{-}hilbert\text{-}space}(v) \wedge \\ \quad u \in L({v}, {v}) \wedge \\ \quad \operatorname{is\text{-}projection\text{-}op}(v, {{u}^{*}}{u}) \Rightarrow \\ \quad \quad E_{\mathrm{ran}\,{u}} = {u}{{u}^{*}} \end{array}}}
\vnbprop{cips-expand-scaled}{\vnbstmt{\begin{array}{@{}l@{}} \forall v, c, p, q. \\ \quad (\operatorname{is\text{-}complex\text{-}inner\text{-}product\text{-}space}(v) \wedge ((c \in \mathbb{C}) \wedge ((p \in \operatorname{vec}(v)) \wedge (q \in \operatorname{vec}(v))))) \Rightarrow  \\ \quad \quad \langle\cdot,\cdot\rangle_{v}(\operatorname{vadd}(v)(p, \operatorname{act}(v)(c, q)), \operatorname{vadd}(v)(p, \operatorname{act}(v)(c, q))) \\ \quad \quad =\; (\langle\cdot,\cdot\rangle_{v}(p, p) + (\operatorname{conjugate}(c) \cdot \langle\cdot,\cdot\rangle_{v}(p, q))) + ((c \cdot \langle\cdot,\cdot\rangle_{v}(q, p)) + ((c \cdot \operatorname{conjugate}(c)) \cdot \langle\cdot,\cdot\rangle_{v}(q, q))) \end{array}}}
\vnbprop{isometric-on-preserves-ip}{\vnbstmt{\begin{array}{@{}l@{}} \forall v, k, u, a, b. \\ \quad \operatorname{is\text{-}complex\text{-}inner\text{-}product\text{-}space}(v) \wedge \\ \quad \operatorname{is\text{-}submodule}(v, k) \wedge \\ \quad u \in L({v}, {v}) \wedge \\ \quad \operatorname{is\text{-}isometric\text{-}on}(v, v, u, k) \wedge \\ \quad a \in k \wedge \\ \quad b \in k \Rightarrow \\ \quad \quad \langle\cdot,\cdot\rangle_{v}(u(a), u(b)) = \langle\cdot,\cdot\rangle_{v}(a, b) \end{array}}}
\vnbprop{projection-of-partial-isometry}{\vnbstmt{\begin{array}{@{}l@{}} \forall v, u. \\ \quad (\operatorname{is\text{-}hilbert\text{-}space}(v) \wedge ((u \in L({v}, {v})) \wedge \operatorname{is\text{-}partial\text{-}isometry}(v, u))) \Rightarrow  \\ \quad \quad \operatorname{is\text{-}projection\text{-}op}(v, {{u}^{*}}{u}) \end{array}}}
\vnbprop{partial-isometry-iff-adjoint-square-projection}{\vnbstmt{\begin{array}{@{}l@{}} \forall v, u. \\ \quad (\operatorname{is\text{-}hilbert\text{-}space}(v) \wedge (u \in L({v}, {v}))) \Rightarrow  \\ \quad \quad \operatorname{is\text{-}projection\text{-}op}(v, {{u}^{*}}{u}) \Leftrightarrow \operatorname{is\text{-}partial\text{-}isometry}(v, u) \end{array}}}
\vnbprop{adjoint-square-projection-iff}{\vnbstmt{\begin{array}{@{}l@{}} \forall v, u. \\ \quad (\operatorname{is\text{-}hilbert\text{-}space}(v) \wedge (u \in L({v}, {v}))) \Rightarrow  \\ \quad \quad \operatorname{is\text{-}projection\text{-}op}(v, {{u}^{*}}{u}) \\ \quad \quad \Leftrightarrow\; \operatorname{is\text{-}projection\text{-}op}(v, {u}{{u}^{*}}) \end{array}}}
\vnbprop{adjoint-partial-isometry-iff}{\vnbstmt{\begin{array}{@{}l@{}} \forall v, u. \\ \quad (\operatorname{is\text{-}hilbert\text{-}space}(v) \wedge (u \in L({v}, {v}))) \Rightarrow  \\ \quad \quad \operatorname{is\text{-}projection\text{-}op}(v, {u}{{u}^{*}}) \\ \quad \quad \Leftrightarrow\; \operatorname{is\text{-}partial\text{-}isometry}(v, {u}^{*}) \end{array}}}
\vnbprop{partial-isometry-support-projection}{\vnbstmt{\begin{array}{@{}l@{}} \forall v, u. \\ \quad (\operatorname{is\text{-}hilbert\text{-}space}(v) \wedge ((u \in L({v}, {v})) \wedge \operatorname{is\text{-}partial\text{-}isometry}(v, u))) \Rightarrow  \\ \quad \quad E_{\mathrm{supp}\,{u}} = {{u}^{*}}{u} \end{array}}}
\vnbprop{partial-isometry-range-projection}{\vnbstmt{\begin{array}{@{}l@{}} \forall v, u. \\ \quad (\operatorname{is\text{-}hilbert\text{-}space}(v) \wedge ((u \in L({v}, {v})) \wedge \operatorname{is\text{-}partial\text{-}isometry}(v, u))) \Rightarrow  \\ \quad \quad E_{\mathrm{ran}\,{u}} = {u}{{u}^{*}} \end{array}}}
\vnbprop{op-adjoint-opnorm-eq}{\vnbstmt{\begin{array}{@{}l@{}} \forall v, w, t. \\ \quad (\operatorname{is\text{-}hilbert\text{-}space}(v) \wedge (\operatorname{is\text{-}hilbert\text{-}space}(w) \wedge (t \in L({v}, {w})))) \Rightarrow  \\ \quad \quad \lVert {{t}^{*}} \rVert = \lVert {t} \rVert \end{array}}}

\vnbbookfile{theorem-library/ba-bh2}{theorem-library/ba-bh2}
\noindent{\small ba-bh2.scm -- wave of 2026-11-09/10 (agent BA-BH2), wired 2026-11-10 from scratchpad/babh2-1109/ba-bh2.scm. Library: operator-algebras.}\par
\vnbprop{bop-sum-apply}{\vnbstmt{\begin{array}{@{}l@{}} \forall v, w, s, t, x. \\ \quad \operatorname{is\text{-}complex\text{-}inner\text{-}product\text{-}space}(w) \wedge \\ \quad s \in (\operatorname{vec}(v) \to \operatorname{vec}(w)) \wedge \\ \quad t \in (\operatorname{vec}(v) \to \operatorname{vec}(w)) \wedge \\ \quad x \in \operatorname{vec}(v) \Rightarrow \\ \quad \quad {s} + {t}(x) = \operatorname{vadd}(w)(s(x), t(x)) \end{array}}}
\vnbprop{bop-neg-apply}{\vnbstmt{\begin{array}{@{}l@{}} \forall v, w, t, x. \\ \quad (\operatorname{is\text{-}complex\text{-}inner\text{-}product\text{-}space}(w) \wedge ((t \in (\operatorname{vec}(v) \to \operatorname{vec}(w))) \wedge (x \in \operatorname{vec}(v)))) \Rightarrow  \\ \quad \quad -{t}(x) = \operatorname{vneg}(w)(t(x)) \end{array}}}
\vnbprop{bop-scale-apply}{\vnbstmt{\begin{array}{@{}l@{}} \forall v, w, c, t, x. \\ \quad (\operatorname{is\text{-}complex\text{-}inner\text{-}product\text{-}space}(w) \wedge ((c \in \mathbb{C}) \wedge ((t \in (\operatorname{vec}(v) \to \operatorname{vec}(w))) \wedge (x \in \operatorname{vec}(v))))) \Rightarrow  \\ \quad \quad {c}\,{t}(x) = \operatorname{act}(w)(c, t(x)) \end{array}}}
\vnbprop{bop-zero-apply}{\vnbstmt{\begin{array}{@{}l@{}} \forall v, w, x. \\ \quad (\operatorname{is\text{-}complex\text{-}inner\text{-}product\text{-}space}(w) \wedge (x \in \operatorname{vec}(v))) \Rightarrow  \\ \quad \quad 0(x) = \operatorname{vzero}(w) \end{array}}}
\vnbprop{ip-norm-vneg}{\vnbstmt{\begin{array}{@{}l@{}} \forall v, x. \\ \quad (\operatorname{is\text{-}complex\text{-}inner\text{-}product\text{-}space}(v) \wedge (x \in \operatorname{vec}(v))) \Rightarrow  \\ \quad \quad \lVert \operatorname{vneg}(v)(x) \rVert = \lVert x \rVert \end{array}}}
\vnbprop{bop-sum-bounded}{\vnbstmt{\begin{array}{@{}l@{}} \forall v, w, s, t. \\ \quad \operatorname{is\text{-}complex\text{-}inner\text{-}product\text{-}space}(v) \wedge \\ \quad \operatorname{is\text{-}complex\text{-}inner\text{-}product\text{-}space}(w) \wedge \\ \quad s \in L({v}, {w}) \wedge \\ \quad t \in L({v}, {w}) \Rightarrow \\ \quad \quad {s} + {t} \in L({v}, {w}) \end{array}}}
\vnbprop{bop-neg-bounded}{\vnbstmt{\begin{array}{@{}l@{}} \forall v, w, t. \\ \quad (\operatorname{is\text{-}complex\text{-}inner\text{-}product\text{-}space}(v) \wedge (\operatorname{is\text{-}complex\text{-}inner\text{-}product\text{-}space}(w) \wedge (t \in L({v}, {w})))) \Rightarrow  \\ \quad \quad -{t} \in L({v}, {w}) \end{array}}}
\vnbprop{bop-scale-bounded}{\vnbstmt{\begin{array}{@{}l@{}} \forall v, w, c, t. \\ \quad \operatorname{is\text{-}complex\text{-}inner\text{-}product\text{-}space}(v) \wedge \\ \quad \operatorname{is\text{-}complex\text{-}inner\text{-}product\text{-}space}(w) \wedge \\ \quad c \in \mathbb{C} \wedge \\ \quad t \in L({v}, {w}) \Rightarrow \\ \quad \quad {c}\,{t} \in L({v}, {w}) \end{array}}}
\vnbprop{bop-zero-bounded}{\vnbstmt{\begin{array}{@{}l@{}} \forall v, w. \\ \quad (\operatorname{is\text{-}complex\text{-}inner\text{-}product\text{-}space}(v) \wedge \operatorname{is\text{-}complex\text{-}inner\text{-}product\text{-}space}(w)) \Rightarrow  \\ \quad \quad 0 \in L({v}, {w}) \end{array}}}
\vnbprop{bop-sum-assoc}{\vnbstmt{\begin{array}{@{}l@{}} \forall v, w, r, s, t. \\ \quad \operatorname{is\text{-}complex\text{-}inner\text{-}product\text{-}space}(v) \wedge \\ \quad \operatorname{is\text{-}complex\text{-}inner\text{-}product\text{-}space}(w) \wedge \\ \quad r \in L({v}, {w}) \wedge \\ \quad s \in L({v}, {w}) \wedge \\ \quad t \in L({v}, {w}) \Rightarrow \\ \quad \quad {{r} + {s}} + {t} = {r} + {{s} + {t}} \end{array}}}
\vnbprop{bop-sum-comm}{\vnbstmt{\begin{array}{@{}l@{}} \forall v, w, s, t. \\ \quad \operatorname{is\text{-}complex\text{-}inner\text{-}product\text{-}space}(v) \wedge \\ \quad \operatorname{is\text{-}complex\text{-}inner\text{-}product\text{-}space}(w) \wedge \\ \quad s \in L({v}, {w}) \wedge \\ \quad t \in L({v}, {w}) \Rightarrow \\ \quad \quad {s} + {t} = {t} + {s} \end{array}}}
\vnbprop{bop-sum-zero-left}{\vnbstmt{\begin{array}{@{}l@{}} \forall v, w, t. \\ \quad (\operatorname{is\text{-}complex\text{-}inner\text{-}product\text{-}space}(v) \wedge (\operatorname{is\text{-}complex\text{-}inner\text{-}product\text{-}space}(w) \wedge (t \in L({v}, {w})))) \Rightarrow  \\ \quad \quad {0} + {t} = t \end{array}}}
\vnbprop{bop-sum-zero-right}{\vnbstmt{\begin{array}{@{}l@{}} \forall v, w, t. \\ \quad (\operatorname{is\text{-}complex\text{-}inner\text{-}product\text{-}space}(v) \wedge (\operatorname{is\text{-}complex\text{-}inner\text{-}product\text{-}space}(w) \wedge (t \in L({v}, {w})))) \Rightarrow  \\ \quad \quad {t} + {0} = t \end{array}}}
\vnbprop{bop-sum-neg-left}{\vnbstmt{\begin{array}{@{}l@{}} \forall v, w, t. \\ \quad (\operatorname{is\text{-}complex\text{-}inner\text{-}product\text{-}space}(v) \wedge (\operatorname{is\text{-}complex\text{-}inner\text{-}product\text{-}space}(w) \wedge (t \in L({v}, {w})))) \Rightarrow  \\ \quad \quad {-{t}} + {t} = 0 \end{array}}}
\vnbprop{bop-sum-neg-right}{\vnbstmt{\begin{array}{@{}l@{}} \forall v, w, t. \\ \quad (\operatorname{is\text{-}complex\text{-}inner\text{-}product\text{-}space}(v) \wedge (\operatorname{is\text{-}complex\text{-}inner\text{-}product\text{-}space}(w) \wedge (t \in L({v}, {w})))) \Rightarrow  \\ \quad \quad {t} + {-{t}} = 0 \end{array}}}
\vnbprop{bop-scale-sum}{\vnbstmt{\begin{array}{@{}l@{}} \forall v, w, c, s, t. \\ \quad \operatorname{is\text{-}complex\text{-}inner\text{-}product\text{-}space}(v) \wedge \\ \quad \operatorname{is\text{-}complex\text{-}inner\text{-}product\text{-}space}(w) \wedge \\ \quad c \in \mathbb{C} \wedge \\ \quad s \in L({v}, {w}) \wedge \\ \quad t \in L({v}, {w}) \Rightarrow \\ \quad \quad {c}\,{{s} + {t}} \\ \quad \quad =\; {{c}\,{s}} + {{c}\,{t}} \end{array}}}
\vnbprop{bop-scale-add}{\vnbstmt{\begin{array}{@{}l@{}} \forall v, w, c, d, t. \\ \quad \operatorname{is\text{-}complex\text{-}inner\text{-}product\text{-}space}(v) \wedge \\ \quad \operatorname{is\text{-}complex\text{-}inner\text{-}product\text{-}space}(w) \wedge \\ \quad c \in \mathbb{C} \wedge \\ \quad d \in \mathbb{C} \wedge \\ \quad t \in L({v}, {w}) \Rightarrow \\ \quad \quad {(c + d)}\,{t} = {{c}\,{t}} + {{d}\,{t}} \end{array}}}
\vnbprop{bop-scale-mul}{\vnbstmt{\begin{array}{@{}l@{}} \forall v, w, c, d, t. \\ \quad \operatorname{is\text{-}complex\text{-}inner\text{-}product\text{-}space}(v) \wedge \\ \quad \operatorname{is\text{-}complex\text{-}inner\text{-}product\text{-}space}(w) \wedge \\ \quad c \in \mathbb{C} \wedge \\ \quad d \in \mathbb{C} \wedge \\ \quad t \in L({v}, {w}) \Rightarrow \\ \quad \quad {(c \cdot d)}\,{t} = {c}\,{{d}\,{t}} \end{array}}}
\vnbprop{bop-scale-one}{\vnbstmt{\begin{array}{@{}l@{}} \forall v, w, t. \\ \quad (\operatorname{is\text{-}complex\text{-}inner\text{-}product\text{-}space}(v) \wedge (\operatorname{is\text{-}complex\text{-}inner\text{-}product\text{-}space}(w) \wedge (t \in L({v}, {w})))) \Rightarrow  \\ \quad \quad {1}\,{t} = t \end{array}}}
\vnbprop{bop-comp-assoc}{\vnbstmt{\begin{array}{@{}l@{}} \forall v, r, s, t. \\ \quad \operatorname{is\text{-}complex\text{-}inner\text{-}product\text{-}space}(v) \wedge \\ \quad r \in L({v}, {v}) \wedge \\ \quad s \in L({v}, {v}) \wedge \\ \quad t \in L({v}, {v}) \Rightarrow \\ \quad \quad {{r}{s}}{t} = {r}{{s}{t}} \end{array}}}
\vnbprop{bop-comp-sum-left}{\vnbstmt{\begin{array}{@{}l@{}} \forall v, r, s, t. \\ \quad \operatorname{is\text{-}complex\text{-}inner\text{-}product\text{-}space}(v) \wedge \\ \quad r \in L({v}, {v}) \wedge \\ \quad s \in L({v}, {v}) \wedge \\ \quad t \in L({v}, {v}) \Rightarrow \\ \quad \quad {{r} + {s}}{t} = {{r}{t}} + {{s}{t}} \end{array}}}
\vnbprop{bop-comp-sum-right}{\vnbstmt{\begin{array}{@{}l@{}} \forall v, r, s, t. \\ \quad \operatorname{is\text{-}complex\text{-}inner\text{-}product\text{-}space}(v) \wedge \\ \quad r \in L({v}, {v}) \wedge \\ \quad s \in L({v}, {v}) \wedge \\ \quad t \in L({v}, {v}) \Rightarrow \\ \quad \quad {r}{{s} + {t}} = {{r}{s}} + {{r}{t}} \end{array}}}
\vnbprop{bop-comp-scale-left}{\vnbstmt{\begin{array}{@{}l@{}} \forall v, c, s, t. \\ \quad (\operatorname{is\text{-}complex\text{-}inner\text{-}product\text{-}space}(v) \wedge ((c \in \mathbb{C}) \wedge ((s \in L({v}, {v})) \wedge (t \in L({v}, {v}))))) \Rightarrow  \\ \quad \quad {{c}\,{s}}{t} = {c}\,{{s}{t}} \end{array}}}
\vnbprop{bop-comp-scale-right}{\vnbstmt{\begin{array}{@{}l@{}} \forall v, c, s, t. \\ \quad (\operatorname{is\text{-}complex\text{-}inner\text{-}product\text{-}space}(v) \wedge ((c \in \mathbb{C}) \wedge ((s \in L({v}, {v})) \wedge (t \in L({v}, {v}))))) \Rightarrow  \\ \quad \quad {s}{{c}\,{t}} = {c}\,{{s}{t}} \end{array}}}
\vnbprop{bop-comp-id-left}{\vnbstmt{\begin{array}{@{}l@{}} \forall v, t. \\ \quad (\operatorname{is\text{-}complex\text{-}inner\text{-}product\text{-}space}(v) \wedge (t \in L({v}, {v}))) \Rightarrow  \\ \quad \quad {1}{t} = t \end{array}}}
\vnbprop{bop-comp-id-right}{\vnbstmt{\begin{array}{@{}l@{}} \forall v, t. \\ \quad (\operatorname{is\text{-}complex\text{-}inner\text{-}product\text{-}space}(v) \wedge (t \in L({v}, {v}))) \Rightarrow  \\ \quad \quad {t}{1} = t \end{array}}}
\vnbprop{opnorm-zero}{\vnbstmt{\begin{array}{@{}l@{}} \forall v, w. \\ \quad (\operatorname{is\text{-}complex\text{-}inner\text{-}product\text{-}space}(v) \wedge \operatorname{is\text{-}complex\text{-}inner\text{-}product\text{-}space}(w)) \Rightarrow  \\ \quad \quad \lVert {0} \rVert = 0 \end{array}}}
\vnbprop{opnorm-zero-imp}{\vnbstmt{\begin{array}{@{}l@{}} \forall v, w, t. \\ \quad \operatorname{is\text{-}complex\text{-}inner\text{-}product\text{-}space}(v) \wedge \\ \quad \operatorname{is\text{-}complex\text{-}inner\text{-}product\text{-}space}(w) \wedge \\ \quad t \in L({v}, {w}) \wedge \\ \quad \lVert {t} \rVert = 0 \Rightarrow \\ \quad \quad t = 0 \end{array}}}
\vnbprop{opnorm-definite}{\vnbstmt{\begin{array}{@{}l@{}} \forall v, w, t. \\ \quad (\operatorname{is\text{-}complex\text{-}inner\text{-}product\text{-}space}(v) \wedge (\operatorname{is\text{-}complex\text{-}inner\text{-}product\text{-}space}(w) \wedge (t \in L({v}, {w})))) \Rightarrow  \\ \quad \quad (\lVert {t} \rVert = 0) \Leftrightarrow (t = 0) \end{array}}}
\vnbprop{opnorm-triangle}{\vnbstmt{\begin{array}{@{}l@{}} \forall v, w, s, t. \\ \quad \operatorname{is\text{-}complex\text{-}inner\text{-}product\text{-}space}(v) \wedge \\ \quad \operatorname{is\text{-}complex\text{-}inner\text{-}product\text{-}space}(w) \wedge \\ \quad s \in L({v}, {w}) \wedge \\ \quad t \in L({v}, {w}) \Rightarrow \\ \quad \quad \lVert {{s} + {t}} \rVert \leq (\lVert {s} \rVert + \lVert {t} \rVert) \end{array}}}
\vnbprop{opnorm-neg}{\vnbstmt{\begin{array}{@{}l@{}} \forall v, w, t. \\ \quad (\operatorname{is\text{-}complex\text{-}inner\text{-}product\text{-}space}(v) \wedge (\operatorname{is\text{-}complex\text{-}inner\text{-}product\text{-}space}(w) \wedge (t \in L({v}, {w})))) \Rightarrow  \\ \quad \quad \lVert {-{t}} \rVert = \lVert {t} \rVert \end{array}}}
\vnbprop{opnorm-scale}{\vnbstmt{\begin{array}{@{}l@{}} \forall v, w, c, t. \\ \quad \operatorname{is\text{-}complex\text{-}inner\text{-}product\text{-}space}(v) \wedge \\ \quad \operatorname{is\text{-}complex\text{-}inner\text{-}product\text{-}space}(w) \wedge \\ \quad c \in \mathbb{C} \wedge \\ \quad t \in L({v}, {w}) \Rightarrow \\ \quad \quad \lVert {{c}\,{t}} \rVert = (\operatorname{magnitude}(c) \cdot \lVert {t} \rVert) \end{array}}}
\vnbprop{opnorm-id}{\vnbstmt{\begin{array}{@{}l@{}} \forall v. \\ \quad (\operatorname{is\text{-}complex\text{-}inner\text{-}product\text{-}space}(v) \wedge \exists x \in \operatorname{vec}(v).\; \neg (x = \operatorname{vzero}(v))) \Rightarrow  \\ \quad \quad \lVert {1} \rVert = 1 \end{array}}}
\vnbprop{op-adjoint-sum}{\vnbstmt{\begin{array}{@{}l@{}} \forall v, w, s, t. \\ \quad \operatorname{is\text{-}hilbert\text{-}space}(v) \wedge \\ \quad \operatorname{is\text{-}complex\text{-}inner\text{-}product\text{-}space}(w) \wedge \\ \quad s \in L({v}, {w}) \wedge \\ \quad t \in L({v}, {w}) \Rightarrow \\ \quad \quad {{s} + {t}}^{*} = {{s}^{*}} + {{t}^{*}} \end{array}}}
\vnbprop{op-adjoint-scale}{\vnbstmt{\begin{array}{@{}l@{}} \forall v, w, c, t. \\ \quad (\operatorname{is\text{-}hilbert\text{-}space}(v) \wedge (\operatorname{is\text{-}complex\text{-}inner\text{-}product\text{-}space}(w) \wedge ((c \in \mathbb{C}) \wedge (t \in L({v}, {w}))))) \Rightarrow  \\ \quad \quad {{c}\,{t}}^{*} = {\operatorname{conjugate}(c)}\,{{t}^{*}} \end{array}}}
\vnbprop{bounded-ops-is-set}{\vnbstmt{\begin{array}{@{}l@{}} \forall v, w. \\ \quad (\operatorname{is\text{-}complex\text{-}inner\text{-}product\text{-}space}(v) \wedge \operatorname{is\text{-}complex\text{-}inner\text{-}product\text{-}space}(w)) \Rightarrow  \\ \quad \quad L({v}, {w}) \in \mathbf{Set} \end{array}}}
\vnbprop{bop-nag-carr}{\vnbstmt{\begin{array}{@{}l@{}} \forall v, w. \\ \quad (\operatorname{is\text{-}complex\text{-}inner\text{-}product\text{-}space}(v) \wedge \operatorname{is\text{-}complex\text{-}inner\text{-}product\text{-}space}(w)) \Rightarrow  \\ \quad \quad (\operatorname{carr}(\operatorname{bop\text{-}normed\text{-}ag}(v, w)) \simeq L({v}, {w})) \end{array}}}
\vnbprop{bop-nag-opr}{\vnbstmt{\begin{array}{@{}l@{}} \forall v, w. \\ \quad (\operatorname{is\text{-}complex\text{-}inner\text{-}product\text{-}space}(v) \wedge \operatorname{is\text{-}complex\text{-}inner\text{-}product\text{-}space}(w)) \Rightarrow  \\ \quad \quad (\operatorname{opr}(\operatorname{bop\text{-}normed\text{-}ag}(v, w)) \simeq \operatorname{bop\text{-}vadd}(v, w)) \end{array}}}
\vnbprop{bop-nag-iden}{\vnbstmt{\begin{array}{@{}l@{}} \forall v, w. \\ \quad (\operatorname{is\text{-}complex\text{-}inner\text{-}product\text{-}space}(v) \wedge \operatorname{is\text{-}complex\text{-}inner\text{-}product\text{-}space}(w)) \Rightarrow  \\ \quad \quad (\operatorname{iden}(\operatorname{bop\text{-}normed\text{-}ag}(v, w)) \simeq 0) \end{array}}}
\vnbprop{bop-nag-inv}{\vnbstmt{\begin{array}{@{}l@{}} \forall v, w. \\ \quad (\operatorname{is\text{-}complex\text{-}inner\text{-}product\text{-}space}(v) \wedge \operatorname{is\text{-}complex\text{-}inner\text{-}product\text{-}space}(w)) \Rightarrow  \\ \quad \quad (\operatorname{inv}(\operatorname{bop\text{-}normed\text{-}ag}(v, w)) \simeq \operatorname{bop\text{-}vinv}(v, w)) \end{array}}}
\vnbprop{bop-nag-nrm}{\vnbstmt{\begin{array}{@{}l@{}} \forall v, w. \\ \quad (\operatorname{is\text{-}complex\text{-}inner\text{-}product\text{-}space}(v) \wedge \operatorname{is\text{-}complex\text{-}inner\text{-}product\text{-}space}(w)) \Rightarrow  \\ \quad \quad (\operatorname{nrm}(\operatorname{bop\text{-}normed\text{-}ag}(v, w)) \simeq \operatorname{bop\text{-}vnrm}(v, w)) \end{array}}}
\vnbprop{bop-vadd-value}{\vnbstmt{\begin{array}{@{}l@{}} \forall v, w, s, t. \\ \quad \operatorname{is\text{-}complex\text{-}inner\text{-}product\text{-}space}(v) \wedge \\ \quad \operatorname{is\text{-}complex\text{-}inner\text{-}product\text{-}space}(w) \wedge \\ \quad s \in L({v}, {w}) \wedge \\ \quad t \in L({v}, {w}) \Rightarrow \\ \quad \quad \operatorname{bop\text{-}vadd}(v, w)(s, t) = {s} + {t} \end{array}}}
\vnbprop{bop-vinv-value}{\vnbstmt{\begin{array}{@{}l@{}} \forall v, w, t. \\ \quad (\operatorname{is\text{-}complex\text{-}inner\text{-}product\text{-}space}(v) \wedge (\operatorname{is\text{-}complex\text{-}inner\text{-}product\text{-}space}(w) \wedge (t \in L({v}, {w})))) \Rightarrow  \\ \quad \quad \operatorname{bop\text{-}vinv}(v, w)(t) = -{t} \end{array}}}
\vnbprop{bop-vnrm-value}{\vnbstmt{\begin{array}{@{}l@{}} \forall v, w, t. \\ \quad (\operatorname{is\text{-}complex\text{-}inner\text{-}product\text{-}space}(v) \wedge (\operatorname{is\text{-}complex\text{-}inner\text{-}product\text{-}space}(w) \wedge (t \in L({v}, {w})))) \Rightarrow  \\ \quad \quad \operatorname{bop\text{-}vnrm}(v, w)(t) = \lVert {t} \rVert \end{array}}}
\vnbprop{bop-nag-length}{\vnbstmt{\begin{array}{@{}l@{}} \forall v, w. \\ \quad (\operatorname{is\text{-}complex\text{-}inner\text{-}product\text{-}space}(v) \wedge \operatorname{is\text{-}complex\text{-}inner\text{-}product\text{-}space}(w)) \Rightarrow  \\ \quad \quad \operatorname{length}(\operatorname{bop\text{-}normed\text{-}ag}(v, w)) = 5 \end{array}}}
\vnbprop{bop-normed-ag-is-normed-ag}{\vnbstmt{\begin{array}{@{}l@{}} \forall v, w. \\ \quad (\operatorname{is\text{-}complex\text{-}inner\text{-}product\text{-}space}(v) \wedge \operatorname{is\text{-}complex\text{-}inner\text{-}product\text{-}space}(w)) \Rightarrow  \\ \quad \quad \operatorname{is\text{-}normed\text{-}ag}(\operatorname{bop\text{-}normed\text{-}ag}(v, w)) \end{array}}}
\vnbprop{ms-converges-of-scaled}{\vnbstmt{\begin{array}{@{}l@{}} \forall s, m, f, g, l, \ell, c. \\ \quad \operatorname{is\text{-}metric\text{-}space}(s) \wedge \\ \quad f \in (\mathbb{N} \to \operatorname{pts}(s)) \wedge \\ \quad l \in \operatorname{pts}(s) \wedge \\ \quad \operatorname{converges\text{-}to}(m, g, \ell) \wedge \\ \quad c \in \mathbb{R} \wedge \\ \quad 0 \leq c \wedge \\ \quad \forall n \in \mathbb{N}.\; (\operatorname{dist}(s)(f(n), l) \leq (c \cdot \operatorname{dist}(m)(g(n), \ell))) \Rightarrow \\ \quad \quad \operatorname{converges\text{-}to}(s, f, l) \end{array}}}
\vnbprop{ms-cauchy-of-scaled}{\vnbstmt{\begin{array}{@{}l@{}} \forall s, \mathit{tm}, f, g, c. \\ \quad \operatorname{is\text{-}cauchy\text{-}seq}(s, f) \wedge \\ \quad \operatorname{is\text{-}metric\text{-}space}(\mathit{tm}) \wedge \\ \quad g \in (\mathbb{N} \to \operatorname{pts}(\mathit{tm})) \wedge \\ \quad c \in \mathbb{R} \wedge \\ \quad 0 \leq c \wedge \\ \quad \forall m, n \in \mathbb{N}.\; (\operatorname{dist}(\mathit{tm})(g(m), g(n)) \leq (c \cdot \operatorname{dist}(s)(f(m), f(n)))) \Rightarrow \\ \quad \quad \operatorname{is\text{-}cauchy\text{-}seq}(\mathit{tm}, g) \end{array}}}
\vnbprop{ms-limit-dist-le}{\vnbstmt{\begin{array}{@{}l@{}} \forall s, g, l, a, b, n_{0}. \\ \quad \operatorname{converges\text{-}to}(s, g, l) \wedge \\ \quad a \in \operatorname{pts}(s) \wedge \\ \quad b \in \mathbb{R} \wedge \\ \quad n_{0} \in \mathbb{N} \wedge \\ \quad \forall n \in \mathbb{N}.\; ((n_{0} \leq n) \Rightarrow (\operatorname{dist}(s)(g(n), a) \leq b)) \Rightarrow \\ \quad \quad \operatorname{dist}(s)(l, a) \leq b \end{array}}}
\vnbprop{ipms-dist-norm}{\vnbstmt{\begin{array}{@{}l@{}} \forall v, a, b. \\ \quad (\operatorname{is\text{-}complex\text{-}inner\text{-}product\text{-}space}(v) \wedge ((a \in \operatorname{vec}(v)) \wedge (b \in \operatorname{vec}(v)))) \Rightarrow  \\ \quad \quad \operatorname{dist}(\operatorname{ip\text{-}metric\text{-}space}(v))(a, b) = \lVert \operatorname{vadd}(v)(a, \operatorname{vneg}(v)(b)) \rVert \end{array}}}
\vnbprop{bop-nag-dist}{\vnbstmt{\begin{array}{@{}l@{}} \forall v, w, s, t. \\ \quad \operatorname{is\text{-}complex\text{-}inner\text{-}product\text{-}space}(v) \wedge \\ \quad \operatorname{is\text{-}complex\text{-}inner\text{-}product\text{-}space}(w) \wedge \\ \quad s \in L({v}, {w}) \wedge \\ \quad t \in L({v}, {w}) \Rightarrow \\ \quad \quad \operatorname{dist}(\operatorname{nag\text{-}metric\text{-}space}(\operatorname{bop\text{-}normed\text{-}ag}(v, w)))(s, t) \\ \quad \quad =\; \lVert {{s} + {-{t}}} \rVert \end{array}}}
\vnbprop{bop-pt-dist}{\vnbstmt{\begin{array}{@{}l@{}} \forall v, w, s, t, x. \\ \quad \operatorname{is\text{-}complex\text{-}inner\text{-}product\text{-}space}(v) \wedge \\ \quad \operatorname{is\text{-}complex\text{-}inner\text{-}product\text{-}space}(w) \wedge \\ \quad s \in L({v}, {w}) \wedge \\ \quad t \in L({v}, {w}) \wedge \\ \quad x \in \operatorname{vec}(v) \Rightarrow \\ \quad \quad \operatorname{dist}(\operatorname{ip\text{-}metric\text{-}space}(w))(s(x), t(x)) \\ \quad \quad \leq\; \lVert x \rVert \cdot \operatorname{dist}(\operatorname{nag\text{-}metric\text{-}space}(\operatorname{bop\text{-}normed\text{-}ag}(v, w)))(s, t) \end{array}}}
\vnbprop{ipms-dist-sum-le}{\vnbstmt{\begin{array}{@{}l@{}} \forall w, a, b, c, e. \\ \quad (\operatorname{is\text{-}complex\text{-}inner\text{-}product\text{-}space}(w) \wedge ((a \in \operatorname{vec}(w)) \wedge ((b \in \operatorname{vec}(w)) \wedge ((c \in \operatorname{vec}(w)) \wedge (e \in \operatorname{vec}(w)))))) \Rightarrow  \\ \quad \quad \operatorname{dist}(\operatorname{ip\text{-}metric\text{-}space}(w))(\operatorname{vadd}(w)(a, b), \operatorname{vadd}(w)(c, e)) \\ \quad \quad \leq\; \operatorname{dist}(\operatorname{ip\text{-}metric\text{-}space}(w))(a, c) + \operatorname{dist}(\operatorname{ip\text{-}metric\text{-}space}(w))(b, e) \end{array}}}
\vnbprop{ipms-dist-act-le}{\vnbstmt{\begin{array}{@{}l@{}} \forall w, c, a, b. \\ \quad (\operatorname{is\text{-}complex\text{-}inner\text{-}product\text{-}space}(w) \wedge ((c \in \mathbb{C}) \wedge ((a \in \operatorname{vec}(w)) \wedge (b \in \operatorname{vec}(w))))) \Rightarrow  \\ \quad \quad \operatorname{dist}(\operatorname{ip\text{-}metric\text{-}space}(w))(\operatorname{act}(w)(c, a), \operatorname{act}(w)(c, b)) \\ \quad \quad \leq\; \operatorname{magnitude}(c) \cdot \operatorname{dist}(\operatorname{ip\text{-}metric\text{-}space}(w))(a, b) \end{array}}}
\vnbprop{bop-cauchy-seq-fun}{\vnbstmt{\begin{array}{@{}l@{}} \forall v, w, f. \\ \quad \operatorname{is\text{-}complex\text{-}inner\text{-}product\text{-}space}(v) \wedge \\ \quad \operatorname{is\text{-}complex\text{-}inner\text{-}product\text{-}space}(w) \wedge \\ \quad \operatorname{is\text{-}cauchy\text{-}seq}(\operatorname{nag\text{-}metric\text{-}space}(\operatorname{bop\text{-}normed\text{-}ag}(v, w)), f) \Rightarrow \\ \quad \quad f \in (\mathbb{N} \to L({v}, {w})) \end{array}}}
\vnbprop{bop-ptwise-seq-fun}{\vnbstmt{\begin{array}{@{}l@{}} \forall v, w, f, x. \\ \quad \operatorname{is\text{-}complex\text{-}inner\text{-}product\text{-}space}(v) \wedge \\ \quad \operatorname{is\text{-}complex\text{-}inner\text{-}product\text{-}space}(w) \wedge \\ \quad f \in (\mathbb{N} \to L({v}, {w})) \wedge \\ \quad x \in \operatorname{vec}(v) \Rightarrow \\ \quad \quad (\lambda q \in \mathbb{N}.\; f(q)(x)) \in (\mathbb{N} \to \operatorname{pts}(\operatorname{ip\text{-}metric\text{-}space}(w))) \end{array}}}
\vnbprop{bop-ptwise-cauchy}{\vnbstmt{\begin{array}{@{}l@{}} \forall v, w, f, x. \\ \quad \operatorname{is\text{-}complex\text{-}inner\text{-}product\text{-}space}(v) \wedge \\ \quad \operatorname{is\text{-}complex\text{-}inner\text{-}product\text{-}space}(w) \wedge \\ \quad \operatorname{is\text{-}cauchy\text{-}seq}(\operatorname{nag\text{-}metric\text{-}space}(\operatorname{bop\text{-}normed\text{-}ag}(v, w)), f) \wedge \\ \quad x \in \operatorname{vec}(v) \Rightarrow \\ \quad \quad \operatorname{is\text{-}cauchy\text{-}seq}(\operatorname{ip\text{-}metric\text{-}space}(w), (\lambda q \in \mathbb{N}.\; f(q)(x))) \end{array}}}
\vnbprop{bop-ptwise-limit}{\vnbstmt{\begin{array}{@{}l@{}} \forall v, w, f, x. \\ \quad \operatorname{is\text{-}complex\text{-}inner\text{-}product\text{-}space}(v) \wedge \\ \quad \operatorname{is\text{-}hilbert\text{-}space}(w) \wedge \\ \quad \operatorname{is\text{-}cauchy\text{-}seq}(\operatorname{nag\text{-}metric\text{-}space}(\operatorname{bop\text{-}normed\text{-}ag}(v, w)), f) \wedge \\ \quad x \in \operatorname{vec}(v) \Rightarrow \\ \quad \quad \varepsilon(\{q \in \operatorname{vec}(w) : \operatorname{converges\text{-}to}(\operatorname{ip\text{-}metric\text{-}space}(w), (\lambda \mathit{kq} \in \mathbb{N}.\; f(\mathit{kq})(x)), q)\}) \\ \quad \quad \in\; \{q \in \operatorname{vec}(w) : \operatorname{converges\text{-}to}(\operatorname{ip\text{-}metric\text{-}space}(w), (\lambda \mathit{kq} \in \mathbb{N}.\; f(\mathit{kq})(x)), q)\} \end{array}}}
\vnbprop{bop-complete}{\vnbstmt{\begin{array}{@{}l@{}} \forall v, w. \\ \quad (\operatorname{is\text{-}complex\text{-}inner\text{-}product\text{-}space}(v) \wedge \operatorname{is\text{-}hilbert\text{-}space}(w)) \Rightarrow  \\ \quad \quad \operatorname{is\text{-}complete}(\operatorname{nag\text{-}metric\text{-}space}(\operatorname{bop\text{-}normed\text{-}ag}(v, w))) \end{array}}}
\vnbprop{bh-vadd-type}{\vnbstmt{\begin{array}{@{}l@{}} \forall v. \\ \quad \operatorname{is\text{-}complex\text{-}inner\text{-}product\text{-}space}(v) \Rightarrow  \\ \quad \quad \operatorname{bh\text{-}vadd}(v) \in ((L({v}, {v}) \times L({v}, {v})) \to L({v}, {v})) \end{array}}}
\vnbprop{bh-vneg-type}{\vnbstmt{\begin{array}{@{}l@{}} \forall v. \\ \quad \operatorname{is\text{-}complex\text{-}inner\text{-}product\text{-}space}(v) \Rightarrow  \\ \quad \quad \operatorname{bh\text{-}vneg}(v) \in (L({v}, {v}) \to L({v}, {v})) \end{array}}}
\vnbprop{bh-act-type}{\vnbstmt{\begin{array}{@{}l@{}} \forall v. \\ \quad \operatorname{is\text{-}complex\text{-}inner\text{-}product\text{-}space}(v) \Rightarrow  \\ \quad \quad \operatorname{bh\text{-}act}(v) \in ((\mathbb{C} \times L({v}, {v})) \to L({v}, {v})) \end{array}}}
\vnbprop{bh-vnrm-type}{\vnbstmt{\begin{array}{@{}l@{}} \forall v. \\ \quad \operatorname{is\text{-}complex\text{-}inner\text{-}product\text{-}space}(v) \Rightarrow  \\ \quad \quad \operatorname{bh\text{-}vnrm}(v) \in (L({v}, {v}) \to \mathbb{R}) \end{array}}}
\vnbprop{bh-vmul-type}{\vnbstmt{\begin{array}{@{}l@{}} \forall v. \\ \quad \operatorname{is\text{-}complex\text{-}inner\text{-}product\text{-}space}(v) \Rightarrow  \\ \quad \quad \operatorname{bh\text{-}vmul}(v) \in ((L({v}, {v}) \times L({v}, {v})) \to L({v}, {v})) \end{array}}}
\vnbprop{bh-vstar-type}{\vnbstmt{\begin{array}{@{}l@{}} \forall v. \\ \quad \operatorname{is\text{-}hilbert\text{-}space}(v) \Rightarrow  \\ \quad \quad \operatorname{bh\text{-}vstar}(v) \in (L({v}, {v}) \to L({v}, {v})) \end{array}}}
\vnbprop{bh-vadd-value}{\vnbstmt{\begin{array}{@{}l@{}} \forall v, s, t. \\ \quad (\operatorname{is\text{-}complex\text{-}inner\text{-}product\text{-}space}(v) \wedge ((s \in L({v}, {v})) \wedge (t \in L({v}, {v})))) \Rightarrow  \\ \quad \quad \operatorname{bh\text{-}vadd}(v)(s, t) = {s} + {t} \end{array}}}
\vnbprop{bh-vneg-value}{\vnbstmt{\begin{array}{@{}l@{}} \forall v, t. \\ \quad (\operatorname{is\text{-}complex\text{-}inner\text{-}product\text{-}space}(v) \wedge (t \in L({v}, {v}))) \Rightarrow  \\ \quad \quad \operatorname{bh\text{-}vneg}(v)(t) = -{t} \end{array}}}
\vnbprop{bh-act-value}{\vnbstmt{\begin{array}{@{}l@{}} \forall v, c, t. \\ \quad (\operatorname{is\text{-}complex\text{-}inner\text{-}product\text{-}space}(v) \wedge ((c \in \mathbb{C}) \wedge (t \in L({v}, {v})))) \Rightarrow  \\ \quad \quad \operatorname{bh\text{-}act}(v)(c, t) = {c}\,{t} \end{array}}}
\vnbprop{bh-vnrm-value}{\vnbstmt{\begin{array}{@{}l@{}} \forall v, t. \\ \quad (\operatorname{is\text{-}complex\text{-}inner\text{-}product\text{-}space}(v) \wedge (t \in L({v}, {v}))) \Rightarrow  \\ \quad \quad \operatorname{bh\text{-}vnrm}(v)(t) = \lVert {t} \rVert \end{array}}}
\vnbprop{bh-vmul-value}{\vnbstmt{\begin{array}{@{}l@{}} \forall v, s, t. \\ \quad (\operatorname{is\text{-}complex\text{-}inner\text{-}product\text{-}space}(v) \wedge ((s \in L({v}, {v})) \wedge (t \in L({v}, {v})))) \Rightarrow  \\ \quad \quad \operatorname{bh\text{-}vmul}(v)(s, t) = {s}{t} \end{array}}}
\vnbprop{bh-vstar-value}{\vnbstmt{\begin{array}{@{}l@{}} \forall v, t. \\ \quad (\operatorname{is\text{-}hilbert\text{-}space}(v) \wedge (t \in L({v}, {v}))) \Rightarrow  \\ \quad \quad \operatorname{bh\text{-}vstar}(v)(t) = {t}^{*} \end{array}}}
\vnbprop{bh-vadd-assoc}{\vnbstmt{\begin{array}{@{}l@{}} \forall v. \\ \quad \operatorname{is\text{-}complex\text{-}inner\text{-}product\text{-}space}(v) \Rightarrow  \\ \quad \quad \operatorname{is\text{-}associative}(\operatorname{bh\text{-}vadd}(v), L({v}, {v})) \end{array}}}
\vnbprop{bh-vadd-comm}{\vnbstmt{\begin{array}{@{}l@{}} \forall v. \\ \quad \operatorname{is\text{-}complex\text{-}inner\text{-}product\text{-}space}(v) \Rightarrow  \\ \quad \quad \operatorname{is\text{-}commutative}(\operatorname{bh\text{-}vadd}(v), L({v}, {v})) \end{array}}}
\vnbprop{bh-vadd-ident}{\vnbstmt{\begin{array}{@{}l@{}} \forall v. \\ \quad \operatorname{is\text{-}complex\text{-}inner\text{-}product\text{-}space}(v) \Rightarrow  \\ \quad \quad \operatorname{is\text{-}identity}(\operatorname{bh\text{-}vadd}(v), 0, L({v}, {v})) \end{array}}}
\vnbprop{bh-vadd-inv}{\vnbstmt{\begin{array}{@{}l@{}} \forall v. \\ \quad \operatorname{is\text{-}complex\text{-}inner\text{-}product\text{-}space}(v) \Rightarrow  \\ \quad \quad \operatorname{has\text{-}inverses}(\operatorname{bh\text{-}vadd}(v), 0, \operatorname{bh\text{-}vneg}(v), L({v}, {v})) \end{array}}}
\vnbprop{bh-vmul-assoc}{\vnbstmt{\begin{array}{@{}l@{}} \forall v. \\ \quad \operatorname{is\text{-}complex\text{-}inner\text{-}product\text{-}space}(v) \Rightarrow  \\ \quad \quad \operatorname{is\text{-}associative}(\operatorname{bh\text{-}vmul}(v), L({v}, {v})) \end{array}}}
\vnbprop{bh-act-vadd}{\vnbstmt{\begin{array}{@{}l@{}} \forall v.; \forall r \in \operatorname{carr}(\operatorname{normed\text{-}field\text{-}as\text{-}commutative\text{-}ring}(\operatorname{cc\text{-}normed\text{-}field})), x, y \in L({v}, {v}). \\ \quad \operatorname{is\text{-}complex\text{-}inner\text{-}product\text{-}space}(v) \Rightarrow  \\ \quad \quad \operatorname{bh\text{-}act}(v)(r, \operatorname{bh\text{-}vadd}(v)(x, y)) = \operatorname{bh\text{-}vadd}(v)(\operatorname{bh\text{-}act}(v)(r, x), \operatorname{bh\text{-}act}(v)(r, y)) \end{array}}}
\vnbprop{bh-act-add}{\vnbstmt{\begin{array}{@{}l@{}} \forall v.; \forall r, s \in \operatorname{carr}(\operatorname{normed\text{-}field\text{-}as\text{-}commutative\text{-}ring}(\operatorname{cc\text{-}normed\text{-}field})), x \in L({v}, {v}). \\ \quad \operatorname{is\text{-}complex\text{-}inner\text{-}product\text{-}space}(v) \Rightarrow  \\ \quad \quad \operatorname{bh\text{-}act}(v)(\operatorname{add}(\operatorname{normed\text{-}field\text{-}as\text{-}commutative\text{-}ring}(\operatorname{cc\text{-}normed\text{-}field}))(r, s), x) \\ \quad \quad =\; \operatorname{bh\text{-}vadd}(v)(\operatorname{bh\text{-}act}(v)(r, x), \operatorname{bh\text{-}act}(v)(s, x)) \end{array}}}
\vnbprop{bh-act-mul}{\vnbstmt{\begin{array}{@{}l@{}} \forall v.; \forall r, s \in \operatorname{carr}(\operatorname{normed\text{-}field\text{-}as\text{-}commutative\text{-}ring}(\operatorname{cc\text{-}normed\text{-}field})), x \in L({v}, {v}). \\ \quad \operatorname{is\text{-}complex\text{-}inner\text{-}product\text{-}space}(v) \Rightarrow  \\ \quad \quad \operatorname{bh\text{-}act}(v)(\operatorname{mul}(\operatorname{normed\text{-}field\text{-}as\text{-}commutative\text{-}ring}(\operatorname{cc\text{-}normed\text{-}field}))(r, s), x) \\ \quad \quad =\; \operatorname{bh\text{-}act}(v)(r, \operatorname{bh\text{-}act}(v)(s, x)) \end{array}}}
\vnbprop{bh-act-one}{\vnbstmt{\begin{array}{@{}l@{}} \forall v.; \forall x \in L({v}, {v}). \\ \quad \operatorname{is\text{-}complex\text{-}inner\text{-}product\text{-}space}(v) \Rightarrow  \\ \quad \quad \operatorname{bh\text{-}act}(v)(\operatorname{one}(\operatorname{normed\text{-}field\text{-}as\text{-}commutative\text{-}ring}(\operatorname{cc\text{-}normed\text{-}field})), x) = x \end{array}}}
\vnbprop{bh-vnrm-nonneg}{\vnbstmt{\begin{array}{@{}l@{}} \forall v.; \forall x \in L({v}, {v}). \\ \quad \operatorname{is\text{-}complex\text{-}inner\text{-}product\text{-}space}(v) \Rightarrow (0 \leq \operatorname{bh\text{-}vnrm}(v)(x)) \end{array}}}
\vnbprop{bh-vnrm-zero-iff}{\vnbstmt{\begin{array}{@{}l@{}} \forall v.; \forall x \in L({v}, {v}). \\ \quad \operatorname{is\text{-}complex\text{-}inner\text{-}product\text{-}space}(v) \Rightarrow  \\ \quad \quad (\operatorname{bh\text{-}vnrm}(v)(x) = 0) \Leftrightarrow (x = 0) \end{array}}}
\vnbprop{bh-vnrm-act}{\vnbstmt{\begin{array}{@{}l@{}} \forall v.; \forall r \in \operatorname{carr}(\operatorname{normed\text{-}field\text{-}as\text{-}commutative\text{-}ring}(\operatorname{cc\text{-}normed\text{-}field})), x \in L({v}, {v}). \\ \quad \operatorname{is\text{-}complex\text{-}inner\text{-}product\text{-}space}(v) \Rightarrow  \\ \quad \quad \operatorname{bh\text{-}vnrm}(v)(\operatorname{bh\text{-}act}(v)(r, x)) = (\operatorname{magnitude}(r) \cdot \operatorname{bh\text{-}vnrm}(v)(x)) \end{array}}}
\vnbprop{bh-vnrm-triangle}{\vnbstmt{\begin{array}{@{}l@{}} \forall v.; \forall x, y \in L({v}, {v}). \\ \quad \operatorname{is\text{-}complex\text{-}inner\text{-}product\text{-}space}(v) \Rightarrow  \\ \quad \quad \operatorname{bh\text{-}vnrm}(v)(\operatorname{bh\text{-}vadd}(v)(x, y)) \leq (\operatorname{bh\text{-}vnrm}(v)(x) + \operatorname{bh\text{-}vnrm}(v)(y)) \end{array}}}
\vnbprop{bh-vmul-vadd-right}{\vnbstmt{\begin{array}{@{}l@{}} \forall v.; \forall x, y, z \in L({v}, {v}). \\ \quad \operatorname{is\text{-}complex\text{-}inner\text{-}product\text{-}space}(v) \Rightarrow  \\ \quad \quad \operatorname{bh\text{-}vmul}(v)(x, \operatorname{bh\text{-}vadd}(v)(y, z)) = \operatorname{bh\text{-}vadd}(v)(\operatorname{bh\text{-}vmul}(v)(x, y), \operatorname{bh\text{-}vmul}(v)(x, z)) \end{array}}}
\vnbprop{bh-vmul-vadd-left}{\vnbstmt{\begin{array}{@{}l@{}} \forall v.; \forall x, y, z \in L({v}, {v}). \\ \quad \operatorname{is\text{-}complex\text{-}inner\text{-}product\text{-}space}(v) \Rightarrow  \\ \quad \quad \operatorname{bh\text{-}vmul}(v)(\operatorname{bh\text{-}vadd}(v)(x, y), z) = \operatorname{bh\text{-}vadd}(v)(\operatorname{bh\text{-}vmul}(v)(x, z), \operatorname{bh\text{-}vmul}(v)(y, z)) \end{array}}}
\vnbprop{bh-vmul-act-left}{\vnbstmt{\begin{array}{@{}l@{}} \forall v.; \forall r \in \operatorname{carr}(\operatorname{normed\text{-}field\text{-}as\text{-}commutative\text{-}ring}(\operatorname{cc\text{-}normed\text{-}field})), x, y \in L({v}, {v}). \\ \quad \operatorname{is\text{-}complex\text{-}inner\text{-}product\text{-}space}(v) \Rightarrow  \\ \quad \quad \operatorname{bh\text{-}vmul}(v)(\operatorname{bh\text{-}act}(v)(r, x), y) = \operatorname{bh\text{-}act}(v)(r, \operatorname{bh\text{-}vmul}(v)(x, y)) \end{array}}}
\vnbprop{bh-vmul-act-right}{\vnbstmt{\begin{array}{@{}l@{}} \forall v.; \forall r \in \operatorname{carr}(\operatorname{normed\text{-}field\text{-}as\text{-}commutative\text{-}ring}(\operatorname{cc\text{-}normed\text{-}field})), x, y \in L({v}, {v}). \\ \quad \operatorname{is\text{-}complex\text{-}inner\text{-}product\text{-}space}(v) \Rightarrow  \\ \quad \quad \operatorname{bh\text{-}vmul}(v)(x, \operatorname{bh\text{-}act}(v)(r, y)) = \operatorname{bh\text{-}act}(v)(r, \operatorname{bh\text{-}vmul}(v)(x, y)) \end{array}}}
\vnbprop{bh-vnrm-submult}{\vnbstmt{\begin{array}{@{}l@{}} \forall v.; \forall x, y \in L({v}, {v}). \\ \quad \operatorname{is\text{-}complex\text{-}inner\text{-}product\text{-}space}(v) \Rightarrow  \\ \quad \quad \operatorname{bh\text{-}vnrm}(v)(\operatorname{bh\text{-}vmul}(v)(x, y)) \leq (\operatorname{bh\text{-}vnrm}(v)(x) \cdot \operatorname{bh\text{-}vnrm}(v)(y)) \end{array}}}
\vnbprop{bh-vmul-unit-left}{\vnbstmt{\begin{array}{@{}l@{}} \forall v.; \forall x \in L({v}, {v}). \\ \quad \operatorname{is\text{-}complex\text{-}inner\text{-}product\text{-}space}(v) \Rightarrow (\operatorname{bh\text{-}vmul}(v)(1, x) = x) \end{array}}}
\vnbprop{bh-vmul-unit-right}{\vnbstmt{\begin{array}{@{}l@{}} \forall v.; \forall x \in L({v}, {v}). \\ \quad \operatorname{is\text{-}complex\text{-}inner\text{-}product\text{-}space}(v) \Rightarrow (\operatorname{bh\text{-}vmul}(v)(x, 1) = x) \end{array}}}
\vnbprop{bh-vnrm-vone}{\vnbstmt{\begin{array}{@{}l@{}} \forall v. \\ \quad (\operatorname{is\text{-}complex\text{-}inner\text{-}product\text{-}space}(v) \wedge \exists x \in \operatorname{vec}(v).\; \neg (x = \operatorname{vzero}(v))) \Rightarrow  \\ \quad \quad \operatorname{bh\text{-}vnrm}(v)(1) = 1 \end{array}}}
\vnbprop{bh-vstar-vadd}{\vnbstmt{\begin{array}{@{}l@{}} \forall v.; \forall x, y \in L({v}, {v}). \\ \quad \operatorname{is\text{-}hilbert\text{-}space}(v) \Rightarrow  \\ \quad \quad \operatorname{bh\text{-}vstar}(v)(\operatorname{bh\text{-}vadd}(v)(x, y)) = \operatorname{bh\text{-}vadd}(v)(\operatorname{bh\text{-}vstar}(v)(x), \operatorname{bh\text{-}vstar}(v)(y)) \end{array}}}
\vnbprop{bh-vstar-act}{\vnbstmt{\begin{array}{@{}l@{}} \forall v.; \forall r \in \mathbb{C}, x \in L({v}, {v}). \\ \quad \operatorname{is\text{-}hilbert\text{-}space}(v) \Rightarrow  \\ \quad \quad \operatorname{bh\text{-}vstar}(v)(\operatorname{bh\text{-}act}(v)(r, x)) = \operatorname{bh\text{-}act}(v)(\operatorname{conjugate}(r), \operatorname{bh\text{-}vstar}(v)(x)) \end{array}}}
\vnbprop{bh-vstar-vstar}{\vnbstmt{\begin{array}{@{}l@{}} \forall v.; \forall x \in L({v}, {v}). \\ \quad \operatorname{is\text{-}hilbert\text{-}space}(v) \Rightarrow (\operatorname{bh\text{-}vstar}(v)(\operatorname{bh\text{-}vstar}(v)(x)) = x) \end{array}}}
\vnbprop{bh-vstar-vmul}{\vnbstmt{\begin{array}{@{}l@{}} \forall v.; \forall x, y \in L({v}, {v}). \\ \quad \operatorname{is\text{-}hilbert\text{-}space}(v) \Rightarrow  \\ \quad \quad \operatorname{bh\text{-}vstar}(v)(\operatorname{bh\text{-}vmul}(v)(x, y)) = \operatorname{bh\text{-}vmul}(v)(\operatorname{bh\text{-}vstar}(v)(y), \operatorname{bh\text{-}vstar}(v)(x)) \end{array}}}
\vnbprop{bh-vnrm-vstar}{\vnbstmt{\begin{array}{@{}l@{}} \forall v.; \forall x \in L({v}, {v}). \\ \quad \operatorname{is\text{-}hilbert\text{-}space}(v) \Rightarrow (\operatorname{bh\text{-}vnrm}(v)(\operatorname{bh\text{-}vstar}(v)(x)) = \operatorname{bh\text{-}vnrm}(v)(x)) \end{array}}}
\vnbprop{bh-vnrm-cstar}{\vnbstmt{\begin{array}{@{}l@{}} \forall v.; \forall x \in L({v}, {v}). \\ \quad \operatorname{is\text{-}hilbert\text{-}space}(v) \Rightarrow  \\ \quad \quad \operatorname{bh\text{-}vnrm}(v)(\operatorname{bh\text{-}vmul}(v)(\operatorname{bh\text{-}vstar}(v)(x), x)) = (\operatorname{bh\text{-}vnrm}(v)(x) \cdot \operatorname{bh\text{-}vnrm}(v)(x)) \end{array}}}
\vnbprop{bh-algebra-is-cna}{\vnbstmt{\begin{array}{@{}l@{}} \forall v. \\ \quad \operatorname{is\text{-}complex\text{-}inner\text{-}product\text{-}space}(v) \Rightarrow  \\ \quad \quad \operatorname{is\text{-}complex\text{-}normed\text{-}algebra}({[}\operatorname{normed\text{-}field\text{-}as\text{-}commutative\text{-}ring}(\operatorname{cc\text{-}normed\text{-}field}), L({v}, {v}), \operatorname{bh\text{-}vadd}(v), 0, \operatorname{bh\text{-}vneg}(v), \operatorname{bh\text{-}act}(v), \operatorname{bh\text{-}vnrm}(v), \operatorname{bh\text{-}vmul}(v)]) \end{array}}}
\vnbprop{bh-nag-view}{\vnbstmt{\begin{array}{@{}l@{}} \forall v. \\ \quad \operatorname{is\text{-}complex\text{-}inner\text{-}product\text{-}space}(v) \Rightarrow  \\ \quad \quad (\operatorname{cna\text{-}as\text{-}normed\text{-}ag}({[}\operatorname{normed\text{-}field\text{-}as\text{-}commutative\text{-}ring}(\operatorname{cc\text{-}normed\text{-}field}), L({v}, {v}), \operatorname{bh\text{-}vadd}(v), 0, \operatorname{bh\text{-}vneg}(v), \operatorname{bh\text{-}act}(v), \operatorname{bh\text{-}vnrm}(v), \operatorname{bh\text{-}vmul}(v)]) \simeq \operatorname{bop\text{-}normed\text{-}ag}(v, v)) \end{array}}}
\vnbprop{bh-algebra-complete}{\vnbstmt{\begin{array}{@{}l@{}} \forall v. \\ \quad \operatorname{is\text{-}hilbert\text{-}space}(v) \Rightarrow  \\ \quad \quad \operatorname{is\text{-}complete}(\operatorname{nag\text{-}metric\text{-}space}(\operatorname{cna\text{-}as\text{-}normed\text{-}ag}({[}\operatorname{normed\text{-}field\text{-}as\text{-}commutative\text{-}ring}(\operatorname{cc\text{-}normed\text{-}field}), L({v}, {v}), \operatorname{bh\text{-}vadd}(v), 0, \operatorname{bh\text{-}vneg}(v), \operatorname{bh\text{-}act}(v), \operatorname{bh\text{-}vnrm}(v), \operatorname{bh\text{-}vmul}(v)]))) \end{array}}}
\vnbprop{bh-algebra-is-cba}{\vnbstmt{\begin{array}{@{}l@{}} \forall v. \\ \quad \operatorname{is\text{-}hilbert\text{-}space}(v) \Rightarrow  \\ \quad \quad \operatorname{is\text{-}complex\text{-}banach\text{-}algebra}({[}\operatorname{normed\text{-}field\text{-}as\text{-}commutative\text{-}ring}(\operatorname{cc\text{-}normed\text{-}field}), L({v}, {v}), \operatorname{bh\text{-}vadd}(v), 0, \operatorname{bh\text{-}vneg}(v), \operatorname{bh\text{-}act}(v), \operatorname{bh\text{-}vnrm}(v), \operatorname{bh\text{-}vmul}(v)]) \end{array}}}
\vnbprop{bh-algebra-is-uba}{\vnbstmt{\begin{array}{@{}l@{}} \forall v. \\ \quad (\operatorname{is\text{-}hilbert\text{-}space}(v) \wedge \exists x \in \operatorname{vec}(v).\; \neg (x = \operatorname{vzero}(v))) \Rightarrow  \\ \quad \quad \operatorname{is\text{-}unital\text{-}banach\text{-}algebra}({[}\operatorname{normed\text{-}field\text{-}as\text{-}commutative\text{-}ring}(\operatorname{cc\text{-}normed\text{-}field}), L({v}, {v}), \operatorname{bh\text{-}vadd}(v), 0, \operatorname{bh\text{-}vneg}(v), \operatorname{bh\text{-}act}(v), \operatorname{bh\text{-}vnrm}(v), \operatorname{bh\text{-}vmul}(v), 1]) \end{array}}}
\vnbprop{bh-algebra-is-uiba}{\vnbstmt{\begin{array}{@{}l@{}} \forall v. \\ \quad (\operatorname{is\text{-}hilbert\text{-}space}(v) \wedge \exists x \in \operatorname{vec}(v).\; \neg (x = \operatorname{vzero}(v))) \Rightarrow  \\ \quad \quad \operatorname{is\text{-}unital\text{-}involutive\text{-}banach\text{-}algebra}({[}\operatorname{normed\text{-}field\text{-}as\text{-}commutative\text{-}ring}(\operatorname{cc\text{-}normed\text{-}field}), L({v}, {v}), \operatorname{bh\text{-}vadd}(v), 0, \operatorname{bh\text{-}vneg}(v), \operatorname{bh\text{-}act}(v), \operatorname{bh\text{-}vnrm}(v), \operatorname{bh\text{-}vmul}(v), 1, \operatorname{bh\text{-}vstar}(v)]) \end{array}}}
\vnbprop{bh-algebra-is-ucsa-tuple}{\vnbstmt{\begin{array}{@{}l@{}} \forall v. \\ \quad (\operatorname{is\text{-}hilbert\text{-}space}(v) \wedge \exists x \in \operatorname{vec}(v).\; \neg (x = \operatorname{vzero}(v))) \Rightarrow  \\ \quad \quad \operatorname{is\text{-}unital\text{-}c\text{-}star\text{-}algebra}({[}\operatorname{normed\text{-}field\text{-}as\text{-}commutative\text{-}ring}(\operatorname{cc\text{-}normed\text{-}field}), L({v}, {v}), \operatorname{bh\text{-}vadd}(v), 0, \operatorname{bh\text{-}vneg}(v), \operatorname{bh\text{-}act}(v), \operatorname{bh\text{-}vnrm}(v), \operatorname{bh\text{-}vmul}(v), 1, \operatorname{bh\text{-}vstar}(v)]) \end{array}}}
\vnbprop{bh-algebra-unital-c-star}{\vnbstmt{\begin{array}{@{}l@{}} \forall v. \\ \quad (\operatorname{is\text{-}hilbert\text{-}space}(v) \wedge \exists x \in \operatorname{vec}(v).\; \neg (x = \operatorname{vzero}(v))) \Rightarrow  \\ \quad \quad \operatorname{is\text{-}unital\text{-}c\text{-}star\text{-}algebra}(L({v})) \end{array}}}

\vnbbookfile{theorem-library/hilbert-separable}{theorem-library/hilbert-separable}
\noindent{\small hilbert-separable.scm -- wave of 2026-11-09/10 (agent BA-ONB2), wired 2026-11-10 from scratchpad/baonb2-1109/bo2-sep.scm. Library: operator-algebras.}\par
\vnbprop{bo2-onf-dist-gt-one}{\vnbstmt{\begin{array}{@{}l@{}} \forall v, x, e.; \forall i, j \in x. \\ \quad (\operatorname{is\text{-}complex\text{-}inner\text{-}product\text{-}space}(v) \wedge (\operatorname{is\text{-}orthonormal\text{-}family}(v, x, e) \wedge \neg (i = j))) \Rightarrow  \\ \quad \quad 1 < \operatorname{dist}(\operatorname{ip\text{-}metric\text{-}space}(v))(e(i), e(j)) \end{array}}}
\vnbprop{bo2-nn-countable}{\vnbstmt{\operatorname{is\text{-}countable}(\mathbb{N})}}
\vnbprop{orthonormal-family-countable-of-separable}{\vnbstmt{\begin{array}{@{}l@{}} \forall v, x, e. \\ \quad \operatorname{is\text{-}complex\text{-}inner\text{-}product\text{-}space}(v) \wedge \\ \quad \operatorname{is\text{-}separable}(\operatorname{ip\text{-}metric\text{-}space}(v)) \wedge \\ \quad \operatorname{is\text{-}orthonormal\text{-}family}(v, x, e) \Rightarrow \\ \quad \quad \operatorname{is\text{-}countable}(x) \end{array}}}
\vnbprop{separable-hilbert-countable-basis}{\vnbstmt{\begin{array}{@{}l@{}} \forall v. \\ \quad (\operatorname{is\text{-}complex\text{-}inner\text{-}product\text{-}space}(v) \wedge \operatorname{is\text{-}separable}(\operatorname{ip\text{-}metric\text{-}space}(v))) \Rightarrow  \\ \quad \quad \exists x, e.\; (\operatorname{is\text{-}orthonormal\text{-}basis}(v, x, e) \wedge \operatorname{is\text{-}countable}(x)) \end{array}}}

\vnbbookfile{theorem-library/cl2-integral-operator-bounded}{theorem-library/cl2-integral-operator-bounded}
\noindent{\small cl2-integral-operator-bounded.scm -- wave of 2026-11-09/10 (agent TH3-L2C), wired 2026-11-10 from scratchpad/th3-1109/l2c/kc-d.scm. Library: operator-algebras.}\par
\vnbprop{cl2-int-op-linear}{\vnbstmt{\begin{array}{@{}l@{}} \forall s, a, u, b, c, d, f, g. \\ \quad \operatorname{is\text{-}measure}(s, a, u) \wedge \\ \quad \operatorname{is\text{-}measure}(b, c, d) \wedge \\ \quad \operatorname{is\text{-}sigma\text{-}finite}(s, a, u) \wedge \\ \quad \operatorname{is\text{-}sigma\text{-}finite}(b, c, d) \wedge \\ \quad f \in \operatorname{l2\text{-}set}((s \times b), \operatorname{product\text{-}sigma}(s, a, b, c), \operatorname{product\text{-}measure}(s, a, u, b, c, d)) \wedge \\ \quad g \in \operatorname{l2\text{-}set}((s \times b), \operatorname{product\text{-}sigma}(s, a, b, c), \operatorname{product\text{-}measure}(s, a, u, b, c, d)) \Rightarrow \\ \quad \quad \operatorname{is\text{-}linear\text{-}op}(L^2_{\mathbb{C}}({b}, {d}), L^2_{\mathbb{C}}({s}, {u}), T_{{f} + i {g}}) \end{array}}}
\vnbprop{cl2-int-op-norm-bound}{\vnbstmt{\begin{array}{@{}l@{}} \forall s, a, u, b, c, d, f, g, z. \\ \quad \operatorname{is\text{-}measure}(s, a, u) \wedge \\ \quad \operatorname{is\text{-}measure}(b, c, d) \wedge \\ \quad \operatorname{is\text{-}sigma\text{-}finite}(s, a, u) \wedge \\ \quad \operatorname{is\text{-}sigma\text{-}finite}(b, c, d) \wedge \\ \quad f \in \operatorname{l2\text{-}set}((s \times b), \operatorname{product\text{-}sigma}(s, a, b, c), \operatorname{product\text{-}measure}(s, a, u, b, c, d)) \wedge \\ \quad g \in \operatorname{l2\text{-}set}((s \times b), \operatorname{product\text{-}sigma}(s, a, b, c), \operatorname{product\text{-}measure}(s, a, u, b, c, d)) \wedge \\ \quad z \in \operatorname{vec}(L^2_{\mathbb{C}}({b}, {d})) \Rightarrow \\ \quad \quad \lVert T_{{f} + i {g}}(z) \rVert \\ \quad \quad \leq\; \sqrt{(\operatorname{integral}((s \times b), \operatorname{product\text{-}sigma}(s, a, b, c), \operatorname{product\text{-}measure}(s, a, u, b, c, d), (\lambda x \in (s \times b).\; (f(x) \cdot f(x)))) + \operatorname{integral}((s \times b), \operatorname{product\text{-}sigma}(s, a, b, c), \operatorname{product\text{-}measure}(s, a, u, b, c, d), (\lambda x \in (s \times b).\; (g(x) \cdot g(x)))))} \cdot \lVert z \rVert \end{array}}}

\vnbbookfile{theorem-library/ba1i-unit-banach}{theorem-library/ba1i-unit-banach}
\noindent{\small ba1i-unit-banach.scm -- wave of 2026-11-09/10 (agent BA-1I), wired 2026-11-10 from scratchpad/ba1i-1109/ba1i-unit-banach.scm. Library: operator-algebras.}\par
\vnbprop{b1u-dist-pair}{\vnbstmt{\begin{array}{@{}l@{}} \forall a.; \forall l \in \mathbb{C}, x \in \operatorname{vec}(a), r \in \mathbb{C}, y \in \operatorname{vec}(a). \\ \quad \operatorname{is\text{-}complex\text{-}normed\text{-}algebra}(a) \Rightarrow  \\ \quad \quad \operatorname{dist}(\operatorname{nag\text{-}metric\text{-}space}(\operatorname{cna\text{-}as\text{-}normed\text{-}ag}(\operatorname{unitization\text{-}na}(a))))({[}l, x], {[}r, y]) \\ \quad \quad =\; \operatorname{magnitude}((l - r)) + \operatorname{vnrm}(a)(\operatorname{vadd}(a)(x, \operatorname{vneg}(a)(y))) \end{array}}}
\vnbprop{b1u-dom1-pair}{\vnbstmt{\begin{array}{@{}l@{}} \forall a.; \forall l \in \mathbb{C}, x \in \operatorname{vec}(a), r \in \mathbb{C}, y \in \operatorname{vec}(a). \\ \quad \operatorname{is\text{-}complex\text{-}normed\text{-}algebra}(a) \Rightarrow  \\ \quad \quad \operatorname{dist}(\operatorname{product2\text{-}ms}(\operatorname{cc\text{-}ms}, \operatorname{nag\text{-}metric\text{-}space}(\operatorname{cna\text{-}as\text{-}normed\text{-}ag}(a))))({[}l, x], {[}r, y]) \\ \quad \quad \leq\; \operatorname{dist}(\operatorname{nag\text{-}metric\text{-}space}(\operatorname{cna\text{-}as\text{-}normed\text{-}ag}(\operatorname{unitization\text{-}na}(a))))({[}l, x], {[}r, y]) \end{array}}}
\vnbprop{b1u-dom2-pair}{\vnbstmt{\begin{array}{@{}l@{}} \forall a.; \forall l \in \mathbb{C}, x \in \operatorname{vec}(a), r \in \mathbb{C}, y \in \operatorname{vec}(a). \\ \quad \operatorname{is\text{-}complex\text{-}normed\text{-}algebra}(a) \Rightarrow  \\ \quad \quad \operatorname{dist}(\operatorname{nag\text{-}metric\text{-}space}(\operatorname{cna\text{-}as\text{-}normed\text{-}ag}(\operatorname{unitization\text{-}na}(a))))({[}l, x], {[}r, y]) \\ \quad \quad \leq\; 2 \cdot \operatorname{dist}(\operatorname{product2\text{-}ms}(\operatorname{cc\text{-}ms}, \operatorname{nag\text{-}metric\text{-}space}(\operatorname{cna\text{-}as\text{-}normed\text{-}ag}(a))))({[}l, x], {[}r, y]) \end{array}}}
\vnbprop{b1u-dom1}{\vnbstmt{\begin{array}{@{}l@{}} \forall a.; \forall p, q \in \operatorname{pts}(\operatorname{nag\text{-}metric\text{-}space}(\operatorname{cna\text{-}as\text{-}normed\text{-}ag}(\operatorname{unitization\text{-}na}(a)))). \\ \quad \operatorname{is\text{-}complex\text{-}normed\text{-}algebra}(a) \Rightarrow  \\ \quad \quad \operatorname{dist}(\operatorname{product2\text{-}ms}(\operatorname{cc\text{-}ms}, \operatorname{nag\text{-}metric\text{-}space}(\operatorname{cna\text{-}as\text{-}normed\text{-}ag}(a))))(p, q) \\ \quad \quad \leq\; \operatorname{dist}(\operatorname{nag\text{-}metric\text{-}space}(\operatorname{cna\text{-}as\text{-}normed\text{-}ag}(\operatorname{unitization\text{-}na}(a))))(p, q) \end{array}}}
\vnbprop{b1u-dom2}{\vnbstmt{\begin{array}{@{}l@{}} \forall a.; \forall p, q \in \operatorname{pts}(\operatorname{nag\text{-}metric\text{-}space}(\operatorname{cna\text{-}as\text{-}normed\text{-}ag}(\operatorname{unitization\text{-}na}(a)))). \\ \quad \operatorname{is\text{-}complex\text{-}normed\text{-}algebra}(a) \Rightarrow  \\ \quad \quad \operatorname{dist}(\operatorname{nag\text{-}metric\text{-}space}(\operatorname{cna\text{-}as\text{-}normed\text{-}ag}(\operatorname{unitization\text{-}na}(a))))(p, q) \\ \quad \quad \leq\; 2 \cdot \operatorname{dist}(\operatorname{product2\text{-}ms}(\operatorname{cc\text{-}ms}, \operatorname{nag\text{-}metric\text{-}space}(\operatorname{cna\text{-}as\text{-}normed\text{-}ag}(a))))(p, q) \end{array}}}
\vnbprop{unitization-complete}{\vnbstmt{\begin{array}{@{}l@{}} \forall a. \\ \quad \operatorname{is\text{-}complex\text{-}banach\text{-}algebra}(a) \Rightarrow  \\ \quad \quad \operatorname{is\text{-}complete}(\operatorname{nag\text{-}metric\text{-}space}(\operatorname{cna\text{-}as\text{-}normed\text{-}ag}(\operatorname{unitization\text{-}na}(a)))) \end{array}}}
\vnbprop{unitization-is-complex-banach-algebra}{\vnbstmt{\begin{array}{@{}l@{}} \forall a. \\ \quad \operatorname{is\text{-}complex\text{-}banach\text{-}algebra}(a) \Rightarrow  \\ \quad \quad \operatorname{is\text{-}complex\text{-}banach\text{-}algebra}(\operatorname{unitization\text{-}na}(a)) \end{array}}}
\vnbprop{b1u-one-in}{\vnbstmt{\forall a.\; (\operatorname{is\text{-}complex\text{-}normed\text{-}algebra}(a) \Rightarrow ({[}1, \operatorname{vzero}(a)] \in \operatorname{unit\text{-}vec}(a)))}}
\vnbprop{b1u-ro-scal}{\vnbstmt{\begin{array}{@{}l@{}} \forall a. \\ \quad \operatorname{is\text{-}complex\text{-}normed\text{-}algebra}(a) \Rightarrow  \\ \quad \quad (\operatorname{scal}(\operatorname{unitization}(a)) \simeq \operatorname{normed\text{-}field\text{-}as\text{-}commutative\text{-}ring}(\operatorname{cc\text{-}normed\text{-}field})) \end{array}}}
\vnbprop{b1u-ro-vec}{\vnbstmt{\begin{array}{@{}l@{}} \forall a. \\ \quad \operatorname{is\text{-}complex\text{-}normed\text{-}algebra}(a) \Rightarrow (\operatorname{vec}(\operatorname{unitization}(a)) \simeq \operatorname{unit\text{-}vec}(a)) \end{array}}}
\vnbprop{b1u-ro-vadd}{\vnbstmt{\begin{array}{@{}l@{}} \forall a. \\ \quad \operatorname{is\text{-}complex\text{-}normed\text{-}algebra}(a) \Rightarrow (\operatorname{vadd}(\operatorname{unitization}(a)) \simeq \operatorname{unit\text{-}vadd}(a)) \end{array}}}
\vnbprop{b1u-ro-vzero}{\vnbstmt{\begin{array}{@{}l@{}} \forall a. \\ \quad \operatorname{is\text{-}complex\text{-}normed\text{-}algebra}(a) \Rightarrow (\operatorname{vzero}(\operatorname{unitization}(a)) \simeq {[}0, \operatorname{vzero}(a)]) \end{array}}}
\vnbprop{b1u-ro-vneg}{\vnbstmt{\begin{array}{@{}l@{}} \forall a. \\ \quad \operatorname{is\text{-}complex\text{-}normed\text{-}algebra}(a) \Rightarrow (\operatorname{vneg}(\operatorname{unitization}(a)) \simeq \operatorname{unit\text{-}vneg}(a)) \end{array}}}
\vnbprop{b1u-ro-act}{\vnbstmt{\begin{array}{@{}l@{}} \forall a. \\ \quad \operatorname{is\text{-}complex\text{-}normed\text{-}algebra}(a) \Rightarrow (\operatorname{act}(\operatorname{unitization}(a)) \simeq \operatorname{unit\text{-}act}(a)) \end{array}}}
\vnbprop{b1u-ro-vnrm}{\vnbstmt{\begin{array}{@{}l@{}} \forall a. \\ \quad \operatorname{is\text{-}complex\text{-}normed\text{-}algebra}(a) \Rightarrow (\operatorname{vnrm}(\operatorname{unitization}(a)) \simeq \operatorname{unit\text{-}vnrm}(a)) \end{array}}}
\vnbprop{b1u-ro-vmul}{\vnbstmt{\begin{array}{@{}l@{}} \forall a. \\ \quad \operatorname{is\text{-}complex\text{-}normed\text{-}algebra}(a) \Rightarrow (\cdot_{\operatorname{unitization}(a)} \simeq \operatorname{unit\text{-}vmul}(a)) \end{array}}}
\vnbprop{b1u-ro-vone}{\vnbstmt{\begin{array}{@{}l@{}} \forall a. \\ \quad \operatorname{is\text{-}complex\text{-}normed\text{-}algebra}(a) \Rightarrow (1_{\operatorname{unitization}(a)} \simeq {[}1, \operatorname{vzero}(a)]) \end{array}}}
\vnbprop{unitization-na-unfold}{\vnbstmt{\begin{array}{@{}l@{}} \forall a. \\ \quad (\operatorname{unitization\text{-}na}(a) \simeq {[}\operatorname{normed\text{-}field\text{-}as\text{-}commutative\text{-}ring}(\operatorname{cc\text{-}normed\text{-}field}), \operatorname{unit\text{-}vec}(a), \operatorname{unit\text{-}vadd}(a), {[}0, \operatorname{vzero}(a)], \operatorname{unit\text{-}vneg}(a), \operatorname{unit\text{-}act}(a), \operatorname{unit\text{-}vnrm}(a), \operatorname{unit\text{-}vmul}(a)]) \end{array}}}
\vnbprop{unitization-is-unital-banach-algebra}{\vnbstmt{\begin{array}{@{}l@{}} \forall a. \\ \quad \operatorname{is\text{-}complex\text{-}banach\text{-}algebra}(a) \Rightarrow \operatorname{is\text{-}unital\text{-}banach\text{-}algebra}(\operatorname{unitization}(a)) \end{array}}}

\vnbbookfile{theorem-library/ba1i-unit-star}{theorem-library/ba1i-unit-star}
\noindent{\small ba1i-unit-star.scm -- wave of 2026-11-09/10 (agent BA-1I), wired 2026-11-10 from scratchpad/ba1i-1109/ba1i-unit-star.scm. Library: operator-algebras.}\par
\vnbprop{b1s-nv-scal}{\vnbstmt{\begin{array}{@{}l@{}} \forall a. \\ \quad \operatorname{is\text{-}nu\text{-}involutive\text{-}banach\text{-}algebra}(a) \Rightarrow (\operatorname{scal}(\operatorname{nuiba\text{-}as\text{-}cba}(a)) \simeq \operatorname{scal}(a)) \end{array}}}
\vnbprop{b1s-nv-vec}{\vnbstmt{\begin{array}{@{}l@{}} \forall a. \\ \quad \operatorname{is\text{-}nu\text{-}involutive\text{-}banach\text{-}algebra}(a) \Rightarrow (\operatorname{vec}(\operatorname{nuiba\text{-}as\text{-}cba}(a)) \simeq \operatorname{vec}(a)) \end{array}}}
\vnbprop{b1s-nv-vadd}{\vnbstmt{\begin{array}{@{}l@{}} \forall a. \\ \quad \operatorname{is\text{-}nu\text{-}involutive\text{-}banach\text{-}algebra}(a) \Rightarrow (\operatorname{vadd}(\operatorname{nuiba\text{-}as\text{-}cba}(a)) \simeq \operatorname{vadd}(a)) \end{array}}}
\vnbprop{b1s-nv-vzero}{\vnbstmt{\begin{array}{@{}l@{}} \forall a. \\ \quad \operatorname{is\text{-}nu\text{-}involutive\text{-}banach\text{-}algebra}(a) \Rightarrow (\operatorname{vzero}(\operatorname{nuiba\text{-}as\text{-}cba}(a)) \simeq \operatorname{vzero}(a)) \end{array}}}
\vnbprop{b1s-nv-vneg}{\vnbstmt{\begin{array}{@{}l@{}} \forall a. \\ \quad \operatorname{is\text{-}nu\text{-}involutive\text{-}banach\text{-}algebra}(a) \Rightarrow (\operatorname{vneg}(\operatorname{nuiba\text{-}as\text{-}cba}(a)) \simeq \operatorname{vneg}(a)) \end{array}}}
\vnbprop{b1s-nv-act}{\vnbstmt{\begin{array}{@{}l@{}} \forall a. \\ \quad \operatorname{is\text{-}nu\text{-}involutive\text{-}banach\text{-}algebra}(a) \Rightarrow (\operatorname{act}(\operatorname{nuiba\text{-}as\text{-}cba}(a)) \simeq \operatorname{act}(a)) \end{array}}}
\vnbprop{b1s-nv-vnrm}{\vnbstmt{\begin{array}{@{}l@{}} \forall a. \\ \quad \operatorname{is\text{-}nu\text{-}involutive\text{-}banach\text{-}algebra}(a) \Rightarrow (\operatorname{vnrm}(\operatorname{nuiba\text{-}as\text{-}cba}(a)) \simeq \operatorname{vnrm}(a)) \end{array}}}
\vnbprop{b1s-nv-vmul}{\vnbstmt{\begin{array}{@{}l@{}} \forall a. \\ \quad \operatorname{is\text{-}nu\text{-}involutive\text{-}banach\text{-}algebra}(a) \Rightarrow (\cdot_{\operatorname{nuiba\text{-}as\text{-}cba}(a)} \simeq \cdot_{a}) \end{array}}}
\vnbprop{b1s-uf-unit-vec}{\vnbstmt{\begin{array}{@{}l@{}} \forall a. \\ \quad \operatorname{is\text{-}nu\text{-}involutive\text{-}banach\text{-}algebra}(a) \Rightarrow  \\ \quad \quad (\operatorname{unit\text{-}vec}(\operatorname{nuiba\text{-}as\text{-}cba}(a)) \simeq \operatorname{unit\text{-}vec}(a)) \end{array}}}
\vnbprop{b1s-uf-unit-vadd}{\vnbstmt{\begin{array}{@{}l@{}} \forall a. \\ \quad \operatorname{is\text{-}nu\text{-}involutive\text{-}banach\text{-}algebra}(a) \Rightarrow  \\ \quad \quad (\operatorname{unit\text{-}vadd}(\operatorname{nuiba\text{-}as\text{-}cba}(a)) \simeq \operatorname{unit\text{-}vadd}(a)) \end{array}}}
\vnbprop{b1s-uf-unit-vneg}{\vnbstmt{\begin{array}{@{}l@{}} \forall a. \\ \quad \operatorname{is\text{-}nu\text{-}involutive\text{-}banach\text{-}algebra}(a) \Rightarrow  \\ \quad \quad (\operatorname{unit\text{-}vneg}(\operatorname{nuiba\text{-}as\text{-}cba}(a)) \simeq \operatorname{unit\text{-}vneg}(a)) \end{array}}}
\vnbprop{b1s-uf-unit-act}{\vnbstmt{\begin{array}{@{}l@{}} \forall a. \\ \quad \operatorname{is\text{-}nu\text{-}involutive\text{-}banach\text{-}algebra}(a) \Rightarrow  \\ \quad \quad (\operatorname{unit\text{-}act}(\operatorname{nuiba\text{-}as\text{-}cba}(a)) \simeq \operatorname{unit\text{-}act}(a)) \end{array}}}
\vnbprop{b1s-uf-unit-vnrm}{\vnbstmt{\begin{array}{@{}l@{}} \forall a. \\ \quad \operatorname{is\text{-}nu\text{-}involutive\text{-}banach\text{-}algebra}(a) \Rightarrow  \\ \quad \quad (\operatorname{unit\text{-}vnrm}(\operatorname{nuiba\text{-}as\text{-}cba}(a)) \simeq \operatorname{unit\text{-}vnrm}(a)) \end{array}}}
\vnbprop{b1s-uf-unit-vmul}{\vnbstmt{\begin{array}{@{}l@{}} \forall a. \\ \quad \operatorname{is\text{-}nu\text{-}involutive\text{-}banach\text{-}algebra}(a) \Rightarrow  \\ \quad \quad (\operatorname{unit\text{-}vmul}(\operatorname{nuiba\text{-}as\text{-}cba}(a)) \simeq \operatorname{unit\text{-}vmul}(a)) \end{array}}}
\vnbprop{b1s-nuiba-vadd-in}{\vnbstmt{\begin{array}{@{}l@{}} \forall a.; \forall x, y \in \operatorname{vec}(a). \\ \quad \operatorname{is\text{-}nu\text{-}involutive\text{-}banach\text{-}algebra}(a) \Rightarrow (\operatorname{vadd}(a)(x, y) \in \operatorname{vec}(a)) \end{array}}}
\vnbprop{b1s-nuiba-vneg-in}{\vnbstmt{\begin{array}{@{}l@{}} \forall a.; \forall x \in \operatorname{vec}(a). \\ \quad \operatorname{is\text{-}nu\text{-}involutive\text{-}banach\text{-}algebra}(a) \Rightarrow (\operatorname{vneg}(a)(x) \in \operatorname{vec}(a)) \end{array}}}
\vnbprop{b1s-nuiba-act-in}{\vnbstmt{\begin{array}{@{}l@{}} \forall a.; \forall r \in \mathbb{C}, x \in \operatorname{vec}(a). \\ \quad \operatorname{is\text{-}nu\text{-}involutive\text{-}banach\text{-}algebra}(a) \Rightarrow (\operatorname{act}(a)(r, x) \in \operatorname{vec}(a)) \end{array}}}
\vnbprop{b1s-nuiba-vmul-in}{\vnbstmt{\begin{array}{@{}l@{}} \forall a.; \forall x, y \in \operatorname{vec}(a). \\ \quad \operatorname{is\text{-}nu\text{-}involutive\text{-}banach\text{-}algebra}(a) \Rightarrow (\cdot_{a}(x, y) \in \operatorname{vec}(a)) \end{array}}}
\vnbprop{b1s-nuiba-vzero-in}{\vnbstmt{\forall a.\; (\operatorname{is\text{-}nu\text{-}involutive\text{-}banach\text{-}algebra}(a) \Rightarrow (\operatorname{vzero}(a) \in \operatorname{vec}(a)))}}
\vnbprop{b1s-nuiba-add-assoc}{\vnbstmt{\begin{array}{@{}l@{}} \forall a.; \forall x, y, z \in \operatorname{vec}(a). \\ \quad \operatorname{is\text{-}nu\text{-}involutive\text{-}banach\text{-}algebra}(a) \Rightarrow  \\ \quad \quad \operatorname{vadd}(a)(\operatorname{vadd}(a)(x, y), z) = \operatorname{vadd}(a)(x, \operatorname{vadd}(a)(y, z)) \end{array}}}
\vnbprop{b1s-nuiba-add-comm}{\vnbstmt{\begin{array}{@{}l@{}} \forall a.; \forall x, y \in \operatorname{vec}(a). \\ \quad \operatorname{is\text{-}nu\text{-}involutive\text{-}banach\text{-}algebra}(a) \Rightarrow (\operatorname{vadd}(a)(x, y) = \operatorname{vadd}(a)(y, x)) \end{array}}}
\vnbprop{b1s-nuiba-add-zero-l}{\vnbstmt{\begin{array}{@{}l@{}} \forall a.; \forall x \in \operatorname{vec}(a). \\ \quad \operatorname{is\text{-}nu\text{-}involutive\text{-}banach\text{-}algebra}(a) \Rightarrow (\operatorname{vadd}(a)(\operatorname{vzero}(a), x) = x) \end{array}}}
\vnbprop{b1s-nuiba-add-zero-r}{\vnbstmt{\begin{array}{@{}l@{}} \forall a.; \forall x \in \operatorname{vec}(a). \\ \quad \operatorname{is\text{-}nu\text{-}involutive\text{-}banach\text{-}algebra}(a) \Rightarrow (\operatorname{vadd}(a)(x, \operatorname{vzero}(a)) = x) \end{array}}}
\vnbprop{b1s-nuiba-vneg-act}{\vnbstmt{\begin{array}{@{}l@{}} \forall a.; \forall x \in \operatorname{vec}(a). \\ \quad \operatorname{is\text{-}nu\text{-}involutive\text{-}banach\text{-}algebra}(a) \Rightarrow (\operatorname{vneg}(a)(x) = \operatorname{act}(a)(-1, x)) \end{array}}}
\vnbprop{b1s-nuiba-act-vadd}{\vnbstmt{\begin{array}{@{}l@{}} \forall a.; \forall r \in \mathbb{C}, x, y \in \operatorname{vec}(a). \\ \quad \operatorname{is\text{-}nu\text{-}involutive\text{-}banach\text{-}algebra}(a) \Rightarrow  \\ \quad \quad \operatorname{act}(a)(r, \operatorname{vadd}(a)(x, y)) = \operatorname{vadd}(a)(\operatorname{act}(a)(r, x), \operatorname{act}(a)(r, y)) \end{array}}}
\vnbprop{b1s-nuiba-act-add}{\vnbstmt{\begin{array}{@{}l@{}} \forall a.; \forall r, s \in \mathbb{C}, x \in \operatorname{vec}(a). \\ \quad \operatorname{is\text{-}nu\text{-}involutive\text{-}banach\text{-}algebra}(a) \Rightarrow  \\ \quad \quad \operatorname{act}(a)((r + s), x) = \operatorname{vadd}(a)(\operatorname{act}(a)(r, x), \operatorname{act}(a)(s, x)) \end{array}}}
\vnbprop{b1s-nuiba-act-mul}{\vnbstmt{\begin{array}{@{}l@{}} \forall a.; \forall r, s \in \mathbb{C}, x \in \operatorname{vec}(a). \\ \quad \operatorname{is\text{-}nu\text{-}involutive\text{-}banach\text{-}algebra}(a) \Rightarrow  \\ \quad \quad \operatorname{act}(a)((r \cdot s), x) = \operatorname{act}(a)(r, \operatorname{act}(a)(s, x)) \end{array}}}
\vnbprop{b1s-nuiba-act-one}{\vnbstmt{\begin{array}{@{}l@{}} \forall a.; \forall x \in \operatorname{vec}(a). \\ \quad \operatorname{is\text{-}nu\text{-}involutive\text{-}banach\text{-}algebra}(a) \Rightarrow (\operatorname{act}(a)(1, x) = x) \end{array}}}
\vnbprop{b1s-nuiba-act-zero}{\vnbstmt{\begin{array}{@{}l@{}} \forall a.; \forall x \in \operatorname{vec}(a). \\ \quad \operatorname{is\text{-}nu\text{-}involutive\text{-}banach\text{-}algebra}(a) \Rightarrow (\operatorname{act}(a)(0, x) = \operatorname{vzero}(a)) \end{array}}}
\vnbprop{b1s-nuiba-act-vzero}{\vnbstmt{\begin{array}{@{}l@{}} \forall a.; \forall r \in \mathbb{C}. \\ \quad \operatorname{is\text{-}nu\text{-}involutive\text{-}banach\text{-}algebra}(a) \Rightarrow (\operatorname{act}(a)(r, \operatorname{vzero}(a)) = \operatorname{vzero}(a)) \end{array}}}
\vnbprop{b1s-nuiba-mul-assoc}{\vnbstmt{\begin{array}{@{}l@{}} \forall a.; \forall x, y, z \in \operatorname{vec}(a). \\ \quad \operatorname{is\text{-}nu\text{-}involutive\text{-}banach\text{-}algebra}(a) \Rightarrow  \\ \quad \quad \cdot_{a}(\cdot_{a}(x, y), z) = \cdot_{a}(x, \cdot_{a}(y, z)) \end{array}}}
\vnbprop{b1s-nuiba-mul-dist-l}{\vnbstmt{\begin{array}{@{}l@{}} \forall a.; \forall x, y, z \in \operatorname{vec}(a). \\ \quad \operatorname{is\text{-}nu\text{-}involutive\text{-}banach\text{-}algebra}(a) \Rightarrow  \\ \quad \quad \cdot_{a}(x, \operatorname{vadd}(a)(y, z)) = \operatorname{vadd}(a)(\cdot_{a}(x, y), \cdot_{a}(x, z)) \end{array}}}
\vnbprop{b1s-nuiba-mul-dist-r}{\vnbstmt{\begin{array}{@{}l@{}} \forall a.; \forall x, y, z \in \operatorname{vec}(a). \\ \quad \operatorname{is\text{-}nu\text{-}involutive\text{-}banach\text{-}algebra}(a) \Rightarrow  \\ \quad \quad \cdot_{a}(\operatorname{vadd}(a)(x, y), z) = \operatorname{vadd}(a)(\cdot_{a}(x, z), \cdot_{a}(y, z)) \end{array}}}
\vnbprop{b1s-nuiba-mul-act-l}{\vnbstmt{\begin{array}{@{}l@{}} \forall a.; \forall r \in \mathbb{C}, x, y \in \operatorname{vec}(a). \\ \quad \operatorname{is\text{-}nu\text{-}involutive\text{-}banach\text{-}algebra}(a) \Rightarrow  \\ \quad \quad \cdot_{a}(\operatorname{act}(a)(r, x), y) = \operatorname{act}(a)(r, \cdot_{a}(x, y)) \end{array}}}
\vnbprop{b1s-nuiba-mul-act-r}{\vnbstmt{\begin{array}{@{}l@{}} \forall a.; \forall r \in \mathbb{C}, x, y \in \operatorname{vec}(a). \\ \quad \operatorname{is\text{-}nu\text{-}involutive\text{-}banach\text{-}algebra}(a) \Rightarrow  \\ \quad \quad \cdot_{a}(x, \operatorname{act}(a)(r, y)) = \operatorname{act}(a)(r, \cdot_{a}(x, y)) \end{array}}}
\vnbprop{b1s-nuiba-mul-zero-l}{\vnbstmt{\begin{array}{@{}l@{}} \forall a.; \forall x \in \operatorname{vec}(a). \\ \quad \operatorname{is\text{-}nu\text{-}involutive\text{-}banach\text{-}algebra}(a) \Rightarrow (\cdot_{a}(\operatorname{vzero}(a), x) = \operatorname{vzero}(a)) \end{array}}}
\vnbprop{b1s-nuiba-mul-zero-r}{\vnbstmt{\begin{array}{@{}l@{}} \forall a.; \forall x \in \operatorname{vec}(a). \\ \quad \operatorname{is\text{-}nu\text{-}involutive\text{-}banach\text{-}algebra}(a) \Rightarrow (\cdot_{a}(x, \operatorname{vzero}(a)) = \operatorname{vzero}(a)) \end{array}}}
\vnbprop{b1s-istar-fun}{\vnbstmt{\begin{array}{@{}l@{}} \forall a. \\ \quad \operatorname{is\text{-}nu\text{-}involutive\text{-}banach\text{-}algebra}(a) \Rightarrow ({}^{*_{a}} \in (\operatorname{vec}(a) \to \operatorname{vec}(a))) \end{array}}}
\vnbprop{b1s-istar-isom}{\vnbstmt{\begin{array}{@{}l@{}} \forall a.; \forall x \in \operatorname{vec}(a). \\ \quad \operatorname{is\text{-}nu\text{-}involutive\text{-}banach\text{-}algebra}(a) \Rightarrow  \\ \quad \quad \operatorname{vnrm}(a)({}^{*_{a}}(x)) = \operatorname{vnrm}(a)(x) \end{array}}}
\vnbprop{b1s-nuiba-star-vadd}{\vnbstmt{\begin{array}{@{}l@{}} \forall a.; \forall x, y \in \operatorname{vec}(a). \\ \quad \operatorname{is\text{-}nu\text{-}involutive\text{-}banach\text{-}algebra}(a) \Rightarrow  \\ \quad \quad {}^{*_{a}}(\operatorname{vadd}(a)(x, y)) = \operatorname{vadd}(a)({}^{*_{a}}(x), {}^{*_{a}}(y)) \end{array}}}
\vnbprop{b1s-nuiba-star-act}{\vnbstmt{\begin{array}{@{}l@{}} \forall a.; \forall r \in \mathbb{C}, x \in \operatorname{vec}(a). \\ \quad \operatorname{is\text{-}nu\text{-}involutive\text{-}banach\text{-}algebra}(a) \Rightarrow  \\ \quad \quad {}^{*_{a}}(\operatorname{act}(a)(r, x)) = \operatorname{act}(a)(\operatorname{conjugate}(r), {}^{*_{a}}(x)) \end{array}}}
\vnbprop{b1s-nuiba-star-star}{\vnbstmt{\begin{array}{@{}l@{}} \forall a.; \forall x \in \operatorname{vec}(a). \\ \quad \operatorname{is\text{-}nu\text{-}involutive\text{-}banach\text{-}algebra}(a) \Rightarrow ({}^{*_{a}}({}^{*_{a}}(x)) = x) \end{array}}}
\vnbprop{b1s-nuiba-star-mul}{\vnbstmt{\begin{array}{@{}l@{}} \forall a.; \forall x, y \in \operatorname{vec}(a). \\ \quad \operatorname{is\text{-}nu\text{-}involutive\text{-}banach\text{-}algebra}(a) \Rightarrow  \\ \quad \quad {}^{*_{a}}(\cdot_{a}(x, y)) = \cdot_{a}({}^{*_{a}}(y), {}^{*_{a}}(x)) \end{array}}}
\vnbprop{b1s-nuiba-vstar-in}{\vnbstmt{\begin{array}{@{}l@{}} \forall a.; \forall x \in \operatorname{vec}(a). \\ \quad \operatorname{is\text{-}nu\text{-}involutive\text{-}banach\text{-}algebra}(a) \Rightarrow ({}^{*_{a}}(x) \in \operatorname{vec}(a)) \end{array}}}
\vnbprop{b1s-nuiba-star-zero}{\vnbstmt{\begin{array}{@{}l@{}} \forall a. \\ \quad \operatorname{is\text{-}nu\text{-}involutive\text{-}banach\text{-}algebra}(a) \Rightarrow ({}^{*_{a}}(\operatorname{vzero}(a)) = \operatorname{vzero}(a)) \end{array}}}
\vnbprop{b1s-istar-apply}{\vnbstmt{\begin{array}{@{}l@{}} \forall a.; \forall l \in \mathbb{C}, x \in \operatorname{vec}(a). \\ \quad \operatorname{is\text{-}nu\text{-}involutive\text{-}banach\text{-}algebra}(a) \Rightarrow  \\ \quad \quad \operatorname{unit\text{-}istar}(a)({[}l, x]) = {[}\operatorname{conjugate}(l), {}^{*_{a}}(x)] \end{array}}}
\vnbprop{b1s-typ-istar}{\vnbstmt{\begin{array}{@{}l@{}} \forall a. \\ \quad \operatorname{is\text{-}nu\text{-}involutive\text{-}banach\text{-}algebra}(a) \Rightarrow  \\ \quad \quad \operatorname{unit\text{-}istar}(a) \in (\operatorname{unit\text{-}vec}(a) \to \operatorname{unit\text{-}vec}(a)) \end{array}}}
\vnbprop{b1s-star-vadd-pair}{\vnbstmt{\begin{array}{@{}l@{}} \forall a.; \forall l \in \mathbb{C}, x \in \operatorname{vec}(a), m \in \mathbb{C}, y \in \operatorname{vec}(a). \\ \quad \operatorname{is\text{-}nu\text{-}involutive\text{-}banach\text{-}algebra}(a) \Rightarrow  \\ \quad \quad \operatorname{unit\text{-}istar}(a)(\operatorname{unit\text{-}vadd}(a)({[}l, x], {[}m, y])) \\ \quad \quad =\; \operatorname{unit\text{-}vadd}(a)(\operatorname{unit\text{-}istar}(a)({[}l, x]), \operatorname{unit\text{-}istar}(a)({[}m, y])) \end{array}}}
\vnbprop{unitization-star-vadd}{\vnbstmt{\begin{array}{@{}l@{}} \forall a.; \forall u, v \in \operatorname{unit\text{-}vec}(a). \\ \quad \operatorname{is\text{-}nu\text{-}involutive\text{-}banach\text{-}algebra}(a) \Rightarrow  \\ \quad \quad \operatorname{unit\text{-}istar}(a)(\operatorname{unit\text{-}vadd}(a)(u, v)) = \operatorname{unit\text{-}vadd}(a)(\operatorname{unit\text{-}istar}(a)(u), \operatorname{unit\text{-}istar}(a)(v)) \end{array}}}
\vnbprop{b1s-star-act-pair}{\vnbstmt{\begin{array}{@{}l@{}} \forall a.; \forall c, l \in \mathbb{C}, x \in \operatorname{vec}(a). \\ \quad \operatorname{is\text{-}nu\text{-}involutive\text{-}banach\text{-}algebra}(a) \Rightarrow  \\ \quad \quad \operatorname{unit\text{-}istar}(a)(\operatorname{unit\text{-}act}(a)(c, {[}l, x])) = \operatorname{unit\text{-}act}(a)(\operatorname{conjugate}(c), \operatorname{unit\text{-}istar}(a)({[}l, x])) \end{array}}}
\vnbprop{unitization-star-act}{\vnbstmt{\begin{array}{@{}l@{}} \forall a.; \forall c \in \mathbb{C}, u \in \operatorname{unit\text{-}vec}(a). \\ \quad \operatorname{is\text{-}nu\text{-}involutive\text{-}banach\text{-}algebra}(a) \Rightarrow  \\ \quad \quad \operatorname{unit\text{-}istar}(a)(\operatorname{unit\text{-}act}(a)(c, u)) = \operatorname{unit\text{-}act}(a)(\operatorname{conjugate}(c), \operatorname{unit\text{-}istar}(a)(u)) \end{array}}}
\vnbprop{b1s-star-star-pair}{\vnbstmt{\begin{array}{@{}l@{}} \forall a.; \forall l \in \mathbb{C}, x \in \operatorname{vec}(a). \\ \quad \operatorname{is\text{-}nu\text{-}involutive\text{-}banach\text{-}algebra}(a) \Rightarrow  \\ \quad \quad \operatorname{unit\text{-}istar}(a)(\operatorname{unit\text{-}istar}(a)({[}l, x])) = {[}l, x] \end{array}}}
\vnbprop{unitization-star-star}{\vnbstmt{\begin{array}{@{}l@{}} \forall a.; \forall u \in \operatorname{unit\text{-}vec}(a). \\ \quad \operatorname{is\text{-}nu\text{-}involutive\text{-}banach\text{-}algebra}(a) \Rightarrow  \\ \quad \quad \operatorname{unit\text{-}istar}(a)(\operatorname{unit\text{-}istar}(a)(u)) = u \end{array}}}
\vnbprop{b1s-star-mul-pair}{\vnbstmt{\begin{array}{@{}l@{}} \forall a.; \forall l \in \mathbb{C}, x \in \operatorname{vec}(a), m \in \mathbb{C}, y \in \operatorname{vec}(a). \\ \quad \operatorname{is\text{-}nu\text{-}involutive\text{-}banach\text{-}algebra}(a) \Rightarrow  \\ \quad \quad \operatorname{unit\text{-}istar}(a)(\operatorname{unit\text{-}vmul}(a)({[}l, x], {[}m, y])) \\ \quad \quad =\; \operatorname{unit\text{-}vmul}(a)(\operatorname{unit\text{-}istar}(a)({[}m, y]), \operatorname{unit\text{-}istar}(a)({[}l, x])) \end{array}}}
\vnbprop{unitization-star-mul}{\vnbstmt{\begin{array}{@{}l@{}} \forall a.; \forall u, v \in \operatorname{unit\text{-}vec}(a). \\ \quad \operatorname{is\text{-}nu\text{-}involutive\text{-}banach\text{-}algebra}(a) \Rightarrow  \\ \quad \quad \operatorname{unit\text{-}istar}(a)(\operatorname{unit\text{-}vmul}(a)(u, v)) = \operatorname{unit\text{-}vmul}(a)(\operatorname{unit\text{-}istar}(a)(v), \operatorname{unit\text{-}istar}(a)(u)) \end{array}}}
\vnbprop{b1s-star-isom-pair}{\vnbstmt{\begin{array}{@{}l@{}} \forall a.; \forall l \in \mathbb{C}, x \in \operatorname{vec}(a). \\ \quad \operatorname{is\text{-}nu\text{-}involutive\text{-}banach\text{-}algebra}(a) \Rightarrow  \\ \quad \quad \operatorname{unit\text{-}vnrm}(a)(\operatorname{unit\text{-}istar}(a)({[}l, x])) = \operatorname{unit\text{-}vnrm}(a)({[}l, x]) \end{array}}}
\vnbprop{unitization-star-isometric}{\vnbstmt{\begin{array}{@{}l@{}} \forall a.; \forall u \in \operatorname{unit\text{-}vec}(a). \\ \quad \operatorname{is\text{-}nu\text{-}involutive\text{-}banach\text{-}algebra}(a) \Rightarrow  \\ \quad \quad \operatorname{unit\text{-}vnrm}(a)(\operatorname{unit\text{-}istar}(a)(u)) = \operatorname{unit\text{-}vnrm}(a)(u) \end{array}}}
\vnbprop{unitization-unfold}{\vnbstmt{\begin{array}{@{}l@{}} \forall a. \\ \quad (\operatorname{unitization}(a) \simeq {[}\operatorname{normed\text{-}field\text{-}as\text{-}commutative\text{-}ring}(\operatorname{cc\text{-}normed\text{-}field}), \operatorname{unit\text{-}vec}(a), \operatorname{unit\text{-}vadd}(a), {[}0, \operatorname{vzero}(a)], \operatorname{unit\text{-}vneg}(a), \operatorname{unit\text{-}act}(a), \operatorname{unit\text{-}vnrm}(a), \operatorname{unit\text{-}vmul}(a), {[}1, \operatorname{vzero}(a)]]) \end{array}}}
\vnbprop{b1s-ro-scal}{\vnbstmt{\begin{array}{@{}l@{}} \forall a. \\ \quad \operatorname{is\text{-}nu\text{-}involutive\text{-}banach\text{-}algebra}(a) \Rightarrow  \\ \quad \quad (\operatorname{scal}(\operatorname{unitization\text{-}star}(a)) \simeq \operatorname{normed\text{-}field\text{-}as\text{-}commutative\text{-}ring}(\operatorname{cc\text{-}normed\text{-}field})) \end{array}}}
\vnbprop{b1s-ro-vec}{\vnbstmt{\begin{array}{@{}l@{}} \forall a. \\ \quad \operatorname{is\text{-}nu\text{-}involutive\text{-}banach\text{-}algebra}(a) \Rightarrow  \\ \quad \quad (\operatorname{vec}(\operatorname{unitization\text{-}star}(a)) \simeq \operatorname{unit\text{-}vec}(a)) \end{array}}}
\vnbprop{b1s-ro-vadd}{\vnbstmt{\begin{array}{@{}l@{}} \forall a. \\ \quad \operatorname{is\text{-}nu\text{-}involutive\text{-}banach\text{-}algebra}(a) \Rightarrow  \\ \quad \quad (\operatorname{vadd}(\operatorname{unitization\text{-}star}(a)) \simeq \operatorname{unit\text{-}vadd}(a)) \end{array}}}
\vnbprop{b1s-ro-vzero}{\vnbstmt{\begin{array}{@{}l@{}} \forall a. \\ \quad \operatorname{is\text{-}nu\text{-}involutive\text{-}banach\text{-}algebra}(a) \Rightarrow  \\ \quad \quad (\operatorname{vzero}(\operatorname{unitization\text{-}star}(a)) \simeq {[}0, \operatorname{vzero}(a)]) \end{array}}}
\vnbprop{b1s-ro-vneg}{\vnbstmt{\begin{array}{@{}l@{}} \forall a. \\ \quad \operatorname{is\text{-}nu\text{-}involutive\text{-}banach\text{-}algebra}(a) \Rightarrow  \\ \quad \quad (\operatorname{vneg}(\operatorname{unitization\text{-}star}(a)) \simeq \operatorname{unit\text{-}vneg}(a)) \end{array}}}
\vnbprop{b1s-ro-act}{\vnbstmt{\begin{array}{@{}l@{}} \forall a. \\ \quad \operatorname{is\text{-}nu\text{-}involutive\text{-}banach\text{-}algebra}(a) \Rightarrow  \\ \quad \quad (\operatorname{act}(\operatorname{unitization\text{-}star}(a)) \simeq \operatorname{unit\text{-}act}(a)) \end{array}}}
\vnbprop{b1s-ro-vnrm}{\vnbstmt{\begin{array}{@{}l@{}} \forall a. \\ \quad \operatorname{is\text{-}nu\text{-}involutive\text{-}banach\text{-}algebra}(a) \Rightarrow  \\ \quad \quad (\operatorname{vnrm}(\operatorname{unitization\text{-}star}(a)) \simeq \operatorname{unit\text{-}vnrm}(a)) \end{array}}}
\vnbprop{b1s-ro-vmul}{\vnbstmt{\begin{array}{@{}l@{}} \forall a. \\ \quad \operatorname{is\text{-}nu\text{-}involutive\text{-}banach\text{-}algebra}(a) \Rightarrow  \\ \quad \quad (\cdot_{\operatorname{unitization\text{-}star}(a)} \simeq \operatorname{unit\text{-}vmul}(a)) \end{array}}}
\vnbprop{b1s-ro-vone}{\vnbstmt{\begin{array}{@{}l@{}} \forall a. \\ \quad \operatorname{is\text{-}nu\text{-}involutive\text{-}banach\text{-}algebra}(a) \Rightarrow  \\ \quad \quad (1_{\operatorname{unitization\text{-}star}(a)} \simeq {[}1, \operatorname{vzero}(a)]) \end{array}}}
\vnbprop{b1s-ro-vstar}{\vnbstmt{\begin{array}{@{}l@{}} \forall a. \\ \quad \operatorname{is\text{-}nu\text{-}involutive\text{-}banach\text{-}algebra}(a) \Rightarrow  \\ \quad \quad ({}^{*_{\operatorname{unitization\text{-}star}(a)}} \simeq \operatorname{unit\text{-}istar}(a)) \end{array}}}
\vnbprop{unitization-is-involutive-banach-algebra}{\vnbstmt{\begin{array}{@{}l@{}} \forall a. \\ \quad \operatorname{is\text{-}nu\text{-}involutive\text{-}banach\text{-}algebra}(a) \Rightarrow  \\ \quad \quad \operatorname{is\text{-}unital\text{-}involutive\text{-}banach\text{-}algebra}(\operatorname{unitization\text{-}star}(a)) \end{array}}}

\vnbbookfile{theorem-library/ba-neumann}{theorem-library/ba-neumann}
\noindent{\small ba-neumann.scm -- wave of 2026-11-09/10 (agent BA-1), wired 2026-11-10 from scratchpad/ba1-1109/ba-neumann.scm. Library: operator-algebras.}\par
\vnbprop{ban-ms}{\vnbstmt{\begin{array}{@{}l@{}} \forall a. \\ \quad \operatorname{is\text{-}unital\text{-}banach\text{-}algebra}(a) \Rightarrow \operatorname{is\text{-}metric\text{-}space}(\operatorname{ba\text{-}metric\text{-}space}(a)) \end{array}}}
\vnbprop{ban-pts}{\vnbstmt{\begin{array}{@{}l@{}} \forall a. \\ \quad \operatorname{is\text{-}unital\text{-}banach\text{-}algebra}(a) \Rightarrow (\operatorname{pts}(\operatorname{ba\text{-}metric\text{-}space}(a)) \simeq \operatorname{vec}(a)) \end{array}}}
\vnbprop{ban-dist}{\vnbstmt{\begin{array}{@{}l@{}} \forall a.; \forall u, v \in \operatorname{vec}(a). \\ \quad \operatorname{is\text{-}unital\text{-}banach\text{-}algebra}(a) \Rightarrow  \\ \quad \quad \operatorname{dist}(\operatorname{ba\text{-}metric\text{-}space}(a))(u, v) = \operatorname{vnrm}(a)(\operatorname{vadd}(a)(u, \operatorname{vneg}(a)(v))) \end{array}}}
\vnbprop{ban-dist-sym}{\vnbstmt{\begin{array}{@{}l@{}} \forall a.; \forall u, v \in \operatorname{vec}(a). \\ \quad \operatorname{is\text{-}unital\text{-}banach\text{-}algebra}(a) \Rightarrow  \\ \quad \quad \operatorname{vnrm}(a)(\operatorname{vadd}(a)(u, \operatorname{vneg}(a)(v))) = \operatorname{vnrm}(a)(\operatorname{vadd}(a)(v, \operatorname{vneg}(a)(u))) \end{array}}}
\vnbprop{ban-neg-norm}{\vnbstmt{\begin{array}{@{}l@{}} \forall a.; \forall q \in \operatorname{vec}(a). \\ \quad \operatorname{is\text{-}unital\text{-}banach\text{-}algebra}(a) \Rightarrow (\operatorname{vnrm}(a)(\operatorname{vneg}(a)(q)) = \operatorname{vnrm}(a)(q)) \end{array}}}
\vnbprop{ban-neg-mul-right}{\vnbstmt{\begin{array}{@{}l@{}} \forall a.; \forall y, q \in \operatorname{vec}(a). \\ \quad \operatorname{is\text{-}unital\text{-}banach\text{-}algebra}(a) \Rightarrow  \\ \quad \quad \cdot_{a}(y, \operatorname{vneg}(a)(q)) = \operatorname{vneg}(a)(\cdot_{a}(y, q)) \end{array}}}
\vnbprop{ban-mul-sub-left}{\vnbstmt{\begin{array}{@{}l@{}} \forall a.; \forall y, p, q \in \operatorname{vec}(a). \\ \quad \operatorname{is\text{-}unital\text{-}banach\text{-}algebra}(a) \Rightarrow  \\ \quad \quad \cdot_{a}(y, \operatorname{vadd}(a)(p, \operatorname{vneg}(a)(q))) = \operatorname{vadd}(a)(\cdot_{a}(y, p), \operatorname{vneg}(a)(\cdot_{a}(y, q))) \end{array}}}
\vnbprop{ban-mul-sub-right}{\vnbstmt{\begin{array}{@{}l@{}} \forall a.; \forall y, p, q \in \operatorname{vec}(a). \\ \quad \operatorname{is\text{-}unital\text{-}banach\text{-}algebra}(a) \Rightarrow  \\ \quad \quad \cdot_{a}(\operatorname{vadd}(a)(p, \operatorname{vneg}(a)(q)), y) = \operatorname{vadd}(a)(\cdot_{a}(p, y), \operatorname{vneg}(a)(\cdot_{a}(q, y))) \end{array}}}
\vnbprop{ban-telescope}{\vnbstmt{\begin{array}{@{}l@{}} \forall a.; \forall p, q, r \in \operatorname{vec}(a). \\ \quad \operatorname{is\text{-}unital\text{-}banach\text{-}algebra}(a) \Rightarrow  \\ \quad \quad \operatorname{vadd}(a)(\operatorname{vadd}(a)(p, \operatorname{vneg}(a)(q)), \operatorname{vadd}(a)(q, \operatorname{vneg}(a)(r))) = \operatorname{vadd}(a)(p, \operatorname{vneg}(a)(r)) \end{array}}}
\vnbprop{ban-nsum-type}{\vnbstmt{\begin{array}{@{}l@{}} \forall n \in \mathbb{N}, a.; \forall x \in \operatorname{vec}(a). \\ \quad \operatorname{is\text{-}unital\text{-}banach\text{-}algebra}(a) \Rightarrow (\sum_{k < n} {x}^{k} \in \operatorname{vec}(a)) \end{array}}}
\vnbprop{ban-one-minus-left}{\vnbstmt{\begin{array}{@{}l@{}} \forall a.; \forall x, y \in \operatorname{vec}(a). \\ \quad \operatorname{is\text{-}unital\text{-}banach\text{-}algebra}(a) \Rightarrow  \\ \quad \quad \cdot_{a}(\operatorname{vadd}(a)(1_{a}, \operatorname{vneg}(a)(x)), y) = \operatorname{vadd}(a)(y, \operatorname{vneg}(a)(\cdot_{a}(x, y))) \end{array}}}
\vnbprop{ban-one-minus-right}{\vnbstmt{\begin{array}{@{}l@{}} \forall a.; \forall x, y \in \operatorname{vec}(a). \\ \quad \operatorname{is\text{-}unital\text{-}banach\text{-}algebra}(a) \Rightarrow  \\ \quad \quad \cdot_{a}(y, \operatorname{vadd}(a)(1_{a}, \operatorname{vneg}(a)(x))) = \operatorname{vadd}(a)(y, \operatorname{vneg}(a)(\cdot_{a}(y, x))) \end{array}}}
\vnbprop{ban-nsum-left}{\vnbstmt{\begin{array}{@{}l@{}} \forall n \in \mathbb{N}, a.; \forall x \in \operatorname{vec}(a). \\ \quad \operatorname{is\text{-}unital\text{-}banach\text{-}algebra}(a) \Rightarrow  \\ \quad \quad \cdot_{a}(\operatorname{vadd}(a)(1_{a}, \operatorname{vneg}(a)(x)), \sum_{k < n} {x}^{k}) \\ \quad \quad =\; \operatorname{vadd}(a)(1_{a}, \operatorname{vneg}(a)({x}^{n})) \end{array}}}
\vnbprop{ban-nsum-right}{\vnbstmt{\begin{array}{@{}l@{}} \forall n \in \mathbb{N}, a.; \forall x \in \operatorname{vec}(a). \\ \quad \operatorname{is\text{-}unital\text{-}banach\text{-}algebra}(a) \Rightarrow  \\ \quad \quad \cdot_{a}(\sum_{k < n} {x}^{k}, \operatorname{vadd}(a)(1_{a}, \operatorname{vneg}(a)(x))) \\ \quad \quad =\; \operatorname{vadd}(a)(1_{a}, \operatorname{vneg}(a)({x}^{n})) \end{array}}}
\vnbprop{ban-nsum-step}{\vnbstmt{\begin{array}{@{}l@{}} \forall n \in \mathbb{N}, a.; \forall x \in \operatorname{vec}(a). \\ \quad \operatorname{is\text{-}unital\text{-}banach\text{-}algebra}(a) \Rightarrow  \\ \quad \quad \operatorname{dist}(\operatorname{ba\text{-}metric\text{-}space}(a))(\sum_{k < n} {x}^{k}, \sum_{k < (n + 1)} {x}^{k}) \leq \operatorname{power}(\operatorname{vnrm}(a)(x), n) \end{array}}}
\vnbprop{ban-lmul-conv}{\vnbstmt{\begin{array}{@{}l@{}} \forall a.; \forall y, \ell \in \operatorname{vec}(a), f, g \in (\mathbb{N} \to \operatorname{vec}(a)). \\ \quad \operatorname{is\text{-}unital\text{-}banach\text{-}algebra}(a) \wedge \\ \quad \operatorname{converges\text{-}to}(\operatorname{ba\text{-}metric\text{-}space}(a), f, \ell) \wedge \\ \quad \forall j \in \mathbb{N}.\; (g(j) = \cdot_{a}(y, f(j))) \Rightarrow \\ \quad \quad \operatorname{converges\text{-}to}(\operatorname{ba\text{-}metric\text{-}space}(a), g, \cdot_{a}(y, \ell)) \end{array}}}
\vnbprop{ban-rmul-conv}{\vnbstmt{\begin{array}{@{}l@{}} \forall a.; \forall y, \ell \in \operatorname{vec}(a), f, g \in (\mathbb{N} \to \operatorname{vec}(a)). \\ \quad \operatorname{is\text{-}unital\text{-}banach\text{-}algebra}(a) \wedge \\ \quad \operatorname{converges\text{-}to}(\operatorname{ba\text{-}metric\text{-}space}(a), f, \ell) \wedge \\ \quad \forall j \in \mathbb{N}.\; (g(j) = \cdot_{a}(f(j), y)) \Rightarrow \\ \quad \quad \operatorname{converges\text{-}to}(\operatorname{ba\text{-}metric\text{-}space}(a), g, \cdot_{a}(\ell, y)) \end{array}}}
\vnbprop{ban-conv-one-minus-pow}{\vnbstmt{\begin{array}{@{}l@{}} \forall a.; \forall x \in \operatorname{vec}(a), g \in (\mathbb{N} \to \operatorname{vec}(a)). \\ \quad \operatorname{is\text{-}unital\text{-}banach\text{-}algebra}(a) \wedge \\ \quad \operatorname{vnrm}(a)(x) < 1 \wedge \\ \quad \forall j \in \mathbb{N}.\; (g(j) = \operatorname{vadd}(a)(1_{a}, \operatorname{vneg}(a)({x}^{j}))) \Rightarrow \\ \quad \quad \operatorname{converges\text{-}to}(\operatorname{ba\text{-}metric\text{-}space}(a), g, 1_{a}) \end{array}}}
\vnbprop{ban-nsum-converges}{\vnbstmt{\begin{array}{@{}l@{}} \forall a.; \forall x \in \operatorname{vec}(a). \\ \quad (\operatorname{is\text{-}unital\text{-}banach\text{-}algebra}(a) \wedge (\operatorname{vnrm}(a)(x) < 1)) \Rightarrow  \\ \quad \quad \operatorname{converges}(\operatorname{ba\text{-}metric\text{-}space}(a), (\lambda j \in \mathbb{N}.\; \sum_{k < j} {x}^{k})) \end{array}}}
\vnbprop{neumann-series}{\vnbstmt{\begin{array}{@{}l@{}} \forall a.; \forall x \in \operatorname{vec}(a). \\ \quad (\operatorname{is\text{-}unital\text{-}banach\text{-}algebra}(a) \wedge (\operatorname{vnrm}(a)(x) < 1)) \Rightarrow  \\ \quad \quad \operatorname{is\text{-}invertible\text{-}el}(a, \operatorname{vadd}(a)(1_{a}, \operatorname{vneg}(a)(x))) \wedge \\ \quad \quad \operatorname{converges\text{-}to}(\operatorname{ba\text{-}metric\text{-}space}(a), (\lambda j \in \mathbb{N}.\; \sum_{k < j} {x}^{k}), \\ \quad \quad \quad {\operatorname{vadd}(a)(1_{a}, \operatorname{vneg}(a)(x))}^{-1}) \end{array}}}
\vnbprop{ban-sub-sub}{\vnbstmt{\begin{array}{@{}l@{}} \forall a.; \forall p, w \in \operatorname{vec}(a). \\ \quad \operatorname{is\text{-}unital\text{-}banach\text{-}algebra}(a) \Rightarrow  \\ \quad \quad \operatorname{vadd}(a)(p, \operatorname{vneg}(a)(\operatorname{vadd}(a)(p, \operatorname{vneg}(a)(w)))) = w \end{array}}}
\vnbprop{ban-inverse-norm-pos}{\vnbstmt{\begin{array}{@{}l@{}} \forall a, b. \\ \quad (\operatorname{is\text{-}unital\text{-}banach\text{-}algebra}(a) \wedge \operatorname{is\text{-}invertible\text{-}el}(a, b)) \Rightarrow  \\ \quad \quad 0 < \operatorname{vnrm}(a)({b}^{-1}) \end{array}}}
\vnbprop{invertible-near-invertible}{\vnbstmt{\begin{array}{@{}l@{}} \forall a, b.; \forall x \in \operatorname{vec}(a). \\ \quad \operatorname{is\text{-}unital\text{-}banach\text{-}algebra}(a) \wedge \\ \quad \operatorname{is\text{-}invertible\text{-}el}(a, b) \wedge \\ \quad \operatorname{vnrm}(a)(\operatorname{vadd}(a)(b, \operatorname{vneg}(a)(x))) < \operatorname{recip}(\operatorname{vnrm}(a)({b}^{-1})) \Rightarrow \\ \quad \quad \operatorname{is\text{-}invertible\text{-}el}(a, x) \end{array}}}
\vnbprop{invertibles-open}{\vnbstmt{\begin{array}{@{}l@{}} \forall a. \\ \quad \operatorname{is\text{-}unital\text{-}banach\text{-}algebra}(a) \Rightarrow \operatorname{is\text{-}open}(\operatorname{ba\text{-}metric\text{-}space}(a), \operatorname{invertibles}(a)) \end{array}}}

\vnbbookfile{theorem-library/sub-riesz}{theorem-library/sub-riesz}
\noindent{\small sub-riesz.scm -- agent BA-1117-TRACE (round 2), 2026-11-17. Finding F-5.1 of BA-1111-C2-REPORT.md: a bounded CC-linear functional on a submodule M of a Hilbert space H is $<$m | Y$>$ on M for some Y in H.  Proof: extend it by the complex Hahn-Banach theorem (hahn-banach-seminorm, at the seminorm x |-$>$ ||x|| c) to a bounded ...}\par
\vnbprop{submodule-functional-riesz}{\vnbstmt{\begin{array}{@{}l@{}} \forall v, s, f, c. \\ \quad \operatorname{is\text{-}hilbert\text{-}space}(v) \wedge \\ \quad \operatorname{is\text{-}submodule}(v, s) \wedge \\ \quad f \in (s \to \mathbb{C}) \wedge \\ \quad \forall a, b \in s.\; (f(\operatorname{vadd}(v)(a, b)) = (f(a) + f(b))) \wedge \\ \quad \forall k \in \mathbb{C}, a \in s.\; (f(\operatorname{act}(v)(k, a)) = (k \cdot f(a))) \wedge \\ \quad c \in \mathbb{R} \wedge \\ \quad 0 \leq c \wedge \\ \quad \forall a \in s.\; (\operatorname{magnitude}(f(a)) \leq (c \cdot \lVert a \rVert)) \Rightarrow \\ \quad \quad \exists y \in \operatorname{vec}(v).\; \forall a \in s.\; (f(a) = \langle\cdot,\cdot\rangle_{v}(a, y)) \end{array}}}

\vnbbookfile{theorem-library/cnvs-ba}{theorem-library/cnvs-ba}
\noindent{\small cnvs-ba.scm -- wave of 2026-11-09/10 (agent CNVS), wired 2026-11-10 from scratchpad/cnvs-1109/cnvs-ba.scm. Library: operator-algebras.}\par
\vnbprop{cna-vnrm-fun}{\vnbstmt{\forall a.\; (\operatorname{is\text{-}complex\text{-}normed\text{-}algebra}(a) \Rightarrow (\operatorname{vnrm}(a) \in (\operatorname{vec}(a) \to \mathbb{R})))}}
\vnbprop{cna-scal-pin}{\vnbstmt{\begin{array}{@{}l@{}} \forall a. \\ \quad \operatorname{is\text{-}complex\text{-}normed\text{-}algebra}(a) \Rightarrow  \\ \quad \quad \operatorname{scal}(a) = \operatorname{normed\text{-}field\text{-}as\text{-}commutative\text{-}ring}(\operatorname{cc\text{-}normed\text{-}field}) \end{array}}}
\vnbprop{cna-scal-carr}{\vnbstmt{\forall a.\; (\operatorname{is\text{-}complex\text{-}normed\text{-}algebra}(a) \Rightarrow (\operatorname{carr}(\operatorname{scal}(a)) = \mathbb{C}))}}
\vnbprop{cna-vnrm-nonneg}{\vnbstmt{\forall a.; \forall x \in \operatorname{vec}(a).\; (\operatorname{is\text{-}complex\text{-}normed\text{-}algebra}(a) \Rightarrow (0 \leq \operatorname{vnrm}(a)(x)))}}
\vnbprop{cna-vnrm-definite}{\vnbstmt{\begin{array}{@{}l@{}} \forall a.; \forall x \in \operatorname{vec}(a). \\ \quad \operatorname{is\text{-}complex\text{-}normed\text{-}algebra}(a) \Rightarrow ((\operatorname{vnrm}(a)(x) = 0) \Leftrightarrow (x = \operatorname{vzero}(a))) \end{array}}}
\vnbprop{cna-vnrm-homog-scal}{\vnbstmt{\begin{array}{@{}l@{}} \forall a.; \forall r \in \operatorname{carr}(\operatorname{scal}(a)), x \in \operatorname{vec}(a). \\ \quad \operatorname{is\text{-}complex\text{-}normed\text{-}algebra}(a) \Rightarrow  \\ \quad \quad \operatorname{vnrm}(a)(\operatorname{act}(a)(r, x)) = (\operatorname{magnitude}(r) \cdot \operatorname{vnrm}(a)(x)) \end{array}}}
\vnbprop{cna-vnrm-triangle}{\vnbstmt{\begin{array}{@{}l@{}} \forall a.; \forall x, y \in \operatorname{vec}(a). \\ \quad \operatorname{is\text{-}complex\text{-}normed\text{-}algebra}(a) \Rightarrow  \\ \quad \quad \operatorname{vnrm}(a)(\operatorname{vadd}(a)(x, y)) \leq (\operatorname{vnrm}(a)(x) + \operatorname{vnrm}(a)(y)) \end{array}}}
\vnbprop{cna-vnrm-homog}{\vnbstmt{\begin{array}{@{}l@{}} \forall a.; \forall c \in \mathbb{C}, x \in \operatorname{vec}(a). \\ \quad \operatorname{is\text{-}complex\text{-}normed\text{-}algebra}(a) \Rightarrow  \\ \quad \quad \operatorname{vnrm}(a)(\operatorname{act}(a)(c, x)) = (\operatorname{magnitude}(c) \cdot \operatorname{vnrm}(a)(x)) \end{array}}}
\vnbprop{cna-module-view-scal}{\vnbstmt{\begin{array}{@{}l@{}} \forall a. \\ \quad \operatorname{is\text{-}complex\text{-}normed\text{-}algebra}(a) \Rightarrow (\operatorname{scal}(\operatorname{cna\text{-}as\text{-}module}(a)) \simeq \operatorname{scal}(a)) \end{array}}}
\vnbprop{cna-module-view-vec}{\vnbstmt{\begin{array}{@{}l@{}} \forall a. \\ \quad \operatorname{is\text{-}complex\text{-}normed\text{-}algebra}(a) \Rightarrow (\operatorname{vec}(\operatorname{cna\text{-}as\text{-}module}(a)) \simeq \operatorname{vec}(a)) \end{array}}}
\vnbprop{cna-module-view-vadd}{\vnbstmt{\begin{array}{@{}l@{}} \forall a. \\ \quad \operatorname{is\text{-}complex\text{-}normed\text{-}algebra}(a) \Rightarrow (\operatorname{vadd}(\operatorname{cna\text{-}as\text{-}module}(a)) \simeq \operatorname{vadd}(a)) \end{array}}}
\vnbprop{cna-module-view-vzero}{\vnbstmt{\begin{array}{@{}l@{}} \forall a. \\ \quad \operatorname{is\text{-}complex\text{-}normed\text{-}algebra}(a) \Rightarrow (\operatorname{vzero}(\operatorname{cna\text{-}as\text{-}module}(a)) \simeq \operatorname{vzero}(a)) \end{array}}}
\vnbprop{cna-module-view-vneg}{\vnbstmt{\begin{array}{@{}l@{}} \forall a. \\ \quad \operatorname{is\text{-}complex\text{-}normed\text{-}algebra}(a) \Rightarrow (\operatorname{vneg}(\operatorname{cna\text{-}as\text{-}module}(a)) \simeq \operatorname{vneg}(a)) \end{array}}}
\vnbprop{cna-module-view-act}{\vnbstmt{\begin{array}{@{}l@{}} \forall a. \\ \quad \operatorname{is\text{-}complex\text{-}normed\text{-}algebra}(a) \Rightarrow (\operatorname{act}(\operatorname{cna\text{-}as\text{-}module}(a)) \simeq \operatorname{act}(a)) \end{array}}}
\vnbprop{cna-cmod-norm}{\vnbstmt{\begin{array}{@{}l@{}} \forall a. \\ \quad \operatorname{is\text{-}complex\text{-}normed\text{-}algebra}(a) \Rightarrow \operatorname{is\text{-}cmod\text{-}norm}(\operatorname{cna\text{-}as\text{-}module}(a), \operatorname{vnrm}(a)) \end{array}}}
\vnbprop{cna-cnvs-is-cnvs}{\vnbstmt{\begin{array}{@{}l@{}} \forall a. \\ \quad \operatorname{is\text{-}complex\text{-}normed\text{-}algebra}(a) \Rightarrow \operatorname{is\text{-}complex\text{-}normed\text{-}vector\text{-}space}(\operatorname{cna\text{-}cnvs}(a)) \end{array}}}
\vnbprop{cna-cnvs-vec}{\vnbstmt{\forall a.\; (\operatorname{is\text{-}complex\text{-}normed\text{-}algebra}(a) \Rightarrow (\operatorname{vec}(\operatorname{cna\text{-}cnvs}(a)) \simeq \operatorname{vec}(a)))}}
\vnbprop{cna-cnvs-vadd}{\vnbstmt{\forall a.\; (\operatorname{is\text{-}complex\text{-}normed\text{-}algebra}(a) \Rightarrow (\operatorname{vadd}(\operatorname{cna\text{-}cnvs}(a)) \simeq \operatorname{vadd}(a)))}}
\vnbprop{cna-cnvs-vzero}{\vnbstmt{\forall a.\; (\operatorname{is\text{-}complex\text{-}normed\text{-}algebra}(a) \Rightarrow (\operatorname{vzero}(\operatorname{cna\text{-}cnvs}(a)) \simeq \operatorname{vzero}(a)))}}
\vnbprop{cna-cnvs-vneg}{\vnbstmt{\forall a.\; (\operatorname{is\text{-}complex\text{-}normed\text{-}algebra}(a) \Rightarrow (\operatorname{vneg}(\operatorname{cna\text{-}cnvs}(a)) \simeq \operatorname{vneg}(a)))}}
\vnbprop{cna-cnvs-vnrm}{\vnbstmt{\forall a.\; (\operatorname{is\text{-}complex\text{-}normed\text{-}algebra}(a) \Rightarrow (\operatorname{vnrm}(\operatorname{cna\text{-}cnvs}(a)) \simeq \operatorname{vnrm}(a)))}}
\vnbprop{cna-cnvs-cact}{\vnbstmt{\begin{array}{@{}l@{}} \forall a.; \forall c \in \mathbb{C}, x \in \operatorname{vec}(a). \\ \quad \operatorname{is\text{-}complex\text{-}normed\text{-}algebra}(a) \Rightarrow (\operatorname{cact}(\operatorname{cna\text{-}cnvs}(a))(c, x) = \operatorname{act}(a)(c, x)) \end{array}}}
\vnbprop{cna-cnvs-normed-ag}{\vnbstmt{\begin{array}{@{}l@{}} \forall a. \\ \quad \operatorname{is\text{-}complex\text{-}normed\text{-}algebra}(a) \Rightarrow  \\ \quad \quad (\operatorname{cnvs\text{-}as\text{-}normed\text{-}ag}(\operatorname{cna\text{-}cnvs}(a)) \simeq \operatorname{cna\text{-}as\text{-}normed\text{-}ag}(a)) \end{array}}}
\vnbprop{cna-cnvs-metric-space}{\vnbstmt{\begin{array}{@{}l@{}} \forall a. \\ \quad \operatorname{is\text{-}complex\text{-}normed\text{-}algebra}(a) \Rightarrow  \\ \quad \quad (\operatorname{cnvs\text{-}metric\text{-}space}(\operatorname{cna\text{-}cnvs}(a)) \simeq \operatorname{cna\text{-}metric\text{-}space}(a)) \end{array}}}
\vnbprop{cba-cnvs-is-cbs}{\vnbstmt{\begin{array}{@{}l@{}} \forall a. \\ \quad \operatorname{is\text{-}complex\text{-}banach\text{-}algebra}(a) \Rightarrow \operatorname{is\text{-}complex\text{-}banach\text{-}space}(\operatorname{cna\text{-}cnvs}(a)) \end{array}}}

\vnbbookfile{theorem-library/cnvs-bop-bridge}{theorem-library/cnvs-bop-bridge}
\noindent{\small cnvs-bop-bridge.scm -- wave of 2026-11-09/10 (agent CNVS2), wired 2026-11-10 from scratchpad/cnvs2-1109/cnvs2-bop.scm. Library: operator-algebras.}\par
\vnbprop{bops-cnvs-bops-eq}{\vnbstmt{\begin{array}{@{}l@{}} \forall v, w. \\ \quad (\operatorname{is\text{-}complex\text{-}inner\text{-}product\text{-}space}(v) \wedge \operatorname{is\text{-}complex\text{-}inner\text{-}product\text{-}space}(w)) \Rightarrow  \\ \quad \quad L({v}, {w}) = \operatorname{cnvs\text{-}bounded\text{-}ops}(\operatorname{cips\text{-}cnvs}(v), \operatorname{cips\text{-}cnvs}(w)) \end{array}}}
\vnbprop{opnorm-body-cnvs-iff}{\vnbstmt{\begin{array}{@{}l@{}} \forall v, w, t, e. \\ \quad (\operatorname{is\text{-}complex\text{-}inner\text{-}product\text{-}space}(v) \wedge (\operatorname{is\text{-}complex\text{-}inner\text{-}product\text{-}space}(w) \wedge (t \in (\operatorname{vec}(v) \to \operatorname{vec}(w))))) \Rightarrow  \\ \quad \quad e = 0 \vee \\ \quad \quad \exists x \in \operatorname{vec}(v). \\ \quad \quad \quad \neg (\lVert x \rVert = 0) \wedge (e = (\lVert t(x) \rVert \cdot \operatorname{recip}(\lVert x \rVert))) \\ \quad \quad \Leftrightarrow\; e = 0 \vee \\ \quad \quad \exists \mathit{cbx} \in \operatorname{vec}(v). \\ \quad \quad \quad \neg (\operatorname{vnrm}(\operatorname{cips\text{-}cnvs}(v))(\mathit{cbx}) = 0) \wedge \\ \quad \quad \quad e = (\operatorname{vnrm}(\operatorname{cips\text{-}cnvs}(w))(t(\mathit{cbx})) \cdot \operatorname{recip}(\operatorname{vnrm}(\operatorname{cips\text{-}cnvs}(v))(\mathit{cbx}))) \end{array}}}
\vnbprop{opnorm-cnvs-opnorm}{\vnbstmt{\begin{array}{@{}l@{}} \forall v, w, t. \\ \quad (\operatorname{is\text{-}complex\text{-}inner\text{-}product\text{-}space}(v) \wedge (\operatorname{is\text{-}complex\text{-}inner\text{-}product\text{-}space}(w) \wedge (t \in (\operatorname{vec}(v) \to \operatorname{vec}(w))))) \Rightarrow  \\ \quad \quad (\lVert {t} \rVert \simeq \operatorname{cnvs\text{-}opnorm}(\operatorname{cips\text{-}cnvs}(v), \operatorname{cips\text{-}cnvs}(w), t)) \end{array}}}
\vnbprop{bop-sum-unfold}{\vnbstmt{\begin{array}{@{}l@{}} \forall a, b, s, t. \\ \quad ({s} + {t} \simeq (\lambda x \in \operatorname{vec}(a).\; \operatorname{vadd}(b)(s(x), t(x)))) \end{array}}}
\vnbprop{cbop-sum-unfold}{\vnbstmt{\begin{array}{@{}l@{}} \forall a, b, s, t. \\ \quad (\operatorname{cbop\text{-}sum}(a, b, s, t) \simeq (\lambda x \in \operatorname{vec}(a).\; \operatorname{vadd}(b)(s(x), t(x)))) \end{array}}}
\vnbprop{bop-neg-unfold}{\vnbstmt{\forall a, b, t.\; (-{t} \simeq (\lambda x \in \operatorname{vec}(a).\; \operatorname{vneg}(b)(t(x))))}}
\vnbprop{cbop-neg-unfold}{\vnbstmt{\forall a, b, t.\; (\operatorname{cbop\text{-}neg}(a, b, t) \simeq (\lambda x \in \operatorname{vec}(a).\; \operatorname{vneg}(b)(t(x))))}}
\vnbprop{bop-zero-cnvs}{\vnbstmt{\begin{array}{@{}l@{}} \forall v, w. \\ \quad (\operatorname{is\text{-}complex\text{-}inner\text{-}product\text{-}space}(v) \wedge \operatorname{is\text{-}complex\text{-}inner\text{-}product\text{-}space}(w)) \Rightarrow  \\ \quad \quad (0 \simeq \operatorname{cbop\text{-}zero}(\operatorname{cips\text{-}cnvs}(v), \operatorname{cips\text{-}cnvs}(w))) \end{array}}}
\vnbprop{bop-vinv-cnvs}{\vnbstmt{\begin{array}{@{}l@{}} \forall v, w. \\ \quad (\operatorname{is\text{-}complex\text{-}inner\text{-}product\text{-}space}(v) \wedge \operatorname{is\text{-}complex\text{-}inner\text{-}product\text{-}space}(w)) \Rightarrow  \\ \quad \quad (\operatorname{bop\text{-}vinv}(v, w) \simeq \operatorname{cbop\text{-}vinv}(\operatorname{cips\text{-}cnvs}(v), \operatorname{cips\text{-}cnvs}(w))) \end{array}}}
\vnbprop{bop-vadd-cnvs}{\vnbstmt{\begin{array}{@{}l@{}} \forall v, w. \\ \quad (\operatorname{is\text{-}complex\text{-}inner\text{-}product\text{-}space}(v) \wedge \operatorname{is\text{-}complex\text{-}inner\text{-}product\text{-}space}(w)) \Rightarrow  \\ \quad \quad (\operatorname{bop\text{-}vadd}(v, w) \simeq \operatorname{cbop\text{-}vadd}(\operatorname{cips\text{-}cnvs}(v), \operatorname{cips\text{-}cnvs}(w))) \end{array}}}
\vnbprop{bop-vnrm-cnvs}{\vnbstmt{\begin{array}{@{}l@{}} \forall v, w. \\ \quad (\operatorname{is\text{-}complex\text{-}inner\text{-}product\text{-}space}(v) \wedge \operatorname{is\text{-}complex\text{-}inner\text{-}product\text{-}space}(w)) \Rightarrow  \\ \quad \quad \operatorname{bop\text{-}vnrm}(v, w) = \operatorname{cbop\text{-}vnrm}(\operatorname{cips\text{-}cnvs}(v), \operatorname{cips\text{-}cnvs}(w)) \end{array}}}
\vnbprop{bop-normed-ag-cnvs}{\vnbstmt{\begin{array}{@{}l@{}} \forall v, w. \\ \quad (\operatorname{is\text{-}complex\text{-}inner\text{-}product\text{-}space}(v) \wedge \operatorname{is\text{-}complex\text{-}inner\text{-}product\text{-}space}(w)) \Rightarrow  \\ \quad \quad (\operatorname{bop\text{-}normed\text{-}ag}(v, w) \simeq \operatorname{cbop\text{-}normed\text{-}ag}(\operatorname{cips\text{-}cnvs}(v), \operatorname{cips\text{-}cnvs}(w))) \end{array}}}
\vnbprop{bop-complete-via-cnvs}{\vnbstmt{\begin{array}{@{}l@{}} \forall v, w. \\ \quad (\operatorname{is\text{-}complex\text{-}inner\text{-}product\text{-}space}(v) \wedge \operatorname{is\text{-}hilbert\text{-}space}(w)) \Rightarrow  \\ \quad \quad \operatorname{is\text{-}complete}(\operatorname{nag\text{-}metric\text{-}space}(\operatorname{bop\text{-}normed\text{-}ag}(v, w))) \end{array}}}

\vnbbookfile{theorem-library/ba-spectrum}{theorem-library/ba-spectrum}
\noindent{\small ba-spectrum.scm -- wave of 2026-11-09/10 (agent BA-1), wired 2026-11-10 from scratchpad/ba1-1109/ba-spectrum.scm. Library: operator-algebras.}\par
\vnbprop{bas-scal-neg-eq}{\vnbstmt{\forall a.; \forall z \in \mathbb{C}.\; (\operatorname{is\text{-}unital\text{-}banach\text{-}algebra}(a) \Rightarrow (\operatorname{neg}(\operatorname{scal}(a))(z) = -z))}}
\vnbprop{bas-act-sub}{\vnbstmt{\begin{array}{@{}l@{}} \forall a.; \forall u, z \in \mathbb{C}, v \in \operatorname{vec}(a). \\ \quad \operatorname{is\text{-}unital\text{-}banach\text{-}algebra}(a) \Rightarrow  \\ \quad \quad \operatorname{act}(a)((u + -z), v) = \operatorname{vadd}(a)(\operatorname{act}(a)(u, v), \operatorname{vneg}(a)(\operatorname{act}(a)(z, v))) \end{array}}}
\vnbprop{bas-shift-unfold}{\vnbstmt{\begin{array}{@{}l@{}} \forall a, z, x. \\ \quad (\operatorname{shift\text{-}el}(a, z, x) \simeq \operatorname{vadd}(a)(\operatorname{act}(a)(z, 1_{a}), \operatorname{vneg}(a)(x))) \end{array}}}
\vnbprop{bas-shift-type}{\vnbstmt{\begin{array}{@{}l@{}} \forall a.; \forall z \in \mathbb{C}, x \in \operatorname{vec}(a). \\ \quad \operatorname{is\text{-}unital\text{-}banach\text{-}algebra}(a) \Rightarrow (\operatorname{shift\text{-}el}(a, z, x) \in \operatorname{vec}(a)) \end{array}}}
\vnbprop{bas-neg-sub}{\vnbstmt{\begin{array}{@{}l@{}} \forall a.; \forall r, q \in \operatorname{vec}(a). \\ \quad \operatorname{is\text{-}unital\text{-}banach\text{-}algebra}(a) \Rightarrow  \\ \quad \quad \operatorname{vneg}(a)(\operatorname{vadd}(a)(r, \operatorname{vneg}(a)(q))) = \operatorname{vadd}(a)(q, \operatorname{vneg}(a)(r)) \end{array}}}
\vnbprop{bas-sub-sub-common}{\vnbstmt{\begin{array}{@{}l@{}} \forall a.; \forall p, q, r \in \operatorname{vec}(a). \\ \quad \operatorname{is\text{-}unital\text{-}banach\text{-}algebra}(a) \Rightarrow  \\ \quad \quad \operatorname{vadd}(a)(\operatorname{vadd}(a)(p, \operatorname{vneg}(a)(q)), \operatorname{vneg}(a)(\operatorname{vadd}(a)(r, \operatorname{vneg}(a)(q)))) \\ \quad \quad =\; \operatorname{vadd}(a)(p, \operatorname{vneg}(a)(r)) \end{array}}}
\vnbprop{bas-shift-diff}{\vnbstmt{\begin{array}{@{}l@{}} \forall a.; \forall u, z \in \mathbb{C}, x \in \operatorname{vec}(a). \\ \quad \operatorname{is\text{-}unital\text{-}banach\text{-}algebra}(a) \Rightarrow  \\ \quad \quad \operatorname{vadd}(a)(\operatorname{shift\text{-}el}(a, u, x), \operatorname{vneg}(a)(\operatorname{shift\text{-}el}(a, z, x))) = \operatorname{act}(a)((u + -z), 1_{a}) \end{array}}}
\vnbprop{resolvent-identity}{\vnbstmt{\begin{array}{@{}l@{}} \forall a.; \forall x \in \operatorname{vec}(a), z, u \in \mathbb{C}. \\ \quad \operatorname{is\text{-}unital\text{-}banach\text{-}algebra}(a) \wedge \\ \quad \operatorname{is\text{-}invertible\text{-}el}(a, \operatorname{shift\text{-}el}(a, z, x)) \wedge \\ \quad \operatorname{is\text{-}invertible\text{-}el}(a, \operatorname{shift\text{-}el}(a, u, x)) \Rightarrow \\ \quad \quad \operatorname{vadd}(a)({\operatorname{shift\text{-}el}(a, z, x)}^{-1}, \operatorname{vneg}(a)({\operatorname{shift\text{-}el}(a, u, x)}^{-1})) \\ \quad \quad =\; \operatorname{act}(a)((u + -z), \cdot_{a}({\operatorname{shift\text{-}el}(a, z, x)}^{-1}, {\operatorname{shift\text{-}el}(a, u, x)}^{-1})) \end{array}}}
\vnbprop{bas-act-one}{\vnbstmt{\forall a.; \forall v \in \operatorname{vec}(a).\; (\operatorname{is\text{-}unital\text{-}banach\text{-}algebra}(a) \Rightarrow (\operatorname{act}(a)(1, v) = v))}}
\vnbprop{bas-act-invertible}{\vnbstmt{\begin{array}{@{}l@{}} \forall a.; \forall z \in \mathbb{C}, u. \\ \quad (\operatorname{is\text{-}unital\text{-}banach\text{-}algebra}(a) \wedge (\neg (z = 0) \wedge \operatorname{is\text{-}invertible\text{-}el}(a, u))) \Rightarrow  \\ \quad \quad \operatorname{is\text{-}invertible\text{-}el}(a, \operatorname{act}(a)(z, u)) \end{array}}}
\vnbprop{shift-invertible-outside-ball}{\vnbstmt{\begin{array}{@{}l@{}} \forall a.; \forall x \in \operatorname{vec}(a), z \in \mathbb{C}. \\ \quad (\operatorname{is\text{-}unital\text{-}banach\text{-}algebra}(a) \wedge (\operatorname{vnrm}(a)(x) < \operatorname{magnitude}(z))) \Rightarrow  \\ \quad \quad \operatorname{is\text{-}invertible\text{-}el}(a, \operatorname{shift\text{-}el}(a, z, x)) \end{array}}}
\vnbprop{bas-spectrum-unfold}{\vnbstmt{\begin{array}{@{}l@{}} \forall a, x. \\ \quad (\sigma({x}) \simeq \{z \in \mathbb{C} : \neg \operatorname{is\text{-}invertible\text{-}el}(a, \operatorname{shift\text{-}el}(a, z, x))\}) \end{array}}}
\vnbprop{spectrum-in-closed-ball}{\vnbstmt{\begin{array}{@{}l@{}} \forall a.; \forall x \in \operatorname{vec}(a), z \in \sigma({x}). \\ \quad \operatorname{is\text{-}unital\text{-}banach\text{-}algebra}(a) \Rightarrow (\operatorname{magnitude}(z) \leq \operatorname{vnrm}(a)(x)) \end{array}}}
\vnbprop{bas-act-zero}{\vnbstmt{\begin{array}{@{}l@{}} \forall a.; \forall v \in \operatorname{vec}(a). \\ \quad \operatorname{is\text{-}unital\text{-}banach\text{-}algebra}(a) \Rightarrow (\operatorname{act}(a)(0, v) = \operatorname{vzero}(a)) \end{array}}}
\vnbprop{bas-characters-unfold}{\vnbstmt{\begin{array}{@{}l@{}} \forall a. \\ \quad (\operatorname{characters}(a) \simeq \{c \in (\operatorname{vec}(a) \to \mathbb{C}) : (\operatorname{is\text{-}cc\text{-}linear\text{-}on}(a, c) \wedge (\forall x, y \in \operatorname{vec}(a).\; (c(\cdot_{a}(x, y)) = (c(x) \cdot c(y))) \wedge (c(1_{a}) = 1)))\}) \end{array}}}
\vnbprop{character-nonzero-at-invertible}{\vnbstmt{\begin{array}{@{}l@{}} \forall a.; \forall c \in \operatorname{characters}(a), u. \\ \quad (\operatorname{is\text{-}unital\text{-}banach\text{-}algebra}(a) \wedge \operatorname{is\text{-}invertible\text{-}el}(a, u)) \Rightarrow \neg (c(u) = 0) \end{array}}}
\vnbprop{character-bounded}{\vnbstmt{\begin{array}{@{}l@{}} \forall a.; \forall c \in \operatorname{characters}(a), x \in \operatorname{vec}(a). \\ \quad \operatorname{is\text{-}unital\text{-}banach\text{-}algebra}(a) \Rightarrow (\operatorname{magnitude}(c(x)) \leq \operatorname{vnrm}(a)(x)) \end{array}}}
\vnbprop{left-ideal-invertible-one}{\vnbstmt{\begin{array}{@{}l@{}} \forall a, i.; \forall y \in i. \\ \quad (\operatorname{is\text{-}unital\text{-}banach\text{-}algebra}(a) \wedge (\operatorname{is\text{-}left\text{-}ideal}(a, i) \wedge \operatorname{is\text{-}invertible\text{-}el}(a, y))) \Rightarrow  \\ \quad \quad 1_{a} \in i \end{array}}}
\vnbprop{left-ideal-near-one}{\vnbstmt{\begin{array}{@{}l@{}} \forall a, i.; \forall y \in i. \\ \quad \operatorname{is\text{-}unital\text{-}banach\text{-}algebra}(a) \wedge \\ \quad \operatorname{is\text{-}left\text{-}ideal}(a, i) \wedge \\ \quad \operatorname{vnrm}(a)(\operatorname{vadd}(a)(1_{a}, \operatorname{vneg}(a)(y))) < 1 \Rightarrow \\ \quad \quad 1_{a} \in i \end{array}}}

\vnbbookfile{theorem-library/ba-inverse}{theorem-library/ba-inverse}
\noindent{\small ba-inverse.scm -- wave of 2026-11-09/10 (agent BA-1B), wired 2026-11-10 from scratchpad/ba1b-1109/ba-inverse.scm. Library: operator-algebras.}\par
\vnbprop{bbi-inverse-fix}{\vnbstmt{\begin{array}{@{}l@{}} \forall a.; \forall x \in \operatorname{vec}(a). \\ \quad (\operatorname{is\text{-}unital\text{-}banach\text{-}algebra}(a) \wedge \operatorname{is\text{-}invertible\text{-}el}(a, \operatorname{vadd}(a)(1_{a}, \operatorname{vneg}(a)(x)))) \Rightarrow  \\ \quad \quad {\operatorname{vadd}(a)(1_{a}, \operatorname{vneg}(a)(x))}^{-1} \\ \quad \quad =\; \operatorname{vadd}(a)(1_{a}, \cdot_{a}(x, {\operatorname{vadd}(a)(1_{a}, \operatorname{vneg}(a)(x))}^{-1})) \end{array}}}
\vnbprop{bbi-fix-expand}{\vnbstmt{\begin{array}{@{}l@{}} \forall a.; \forall x, v \in \operatorname{vec}(a). \\ \quad (\operatorname{is\text{-}unital\text{-}banach\text{-}algebra}(a) \wedge (v = \operatorname{vadd}(a)(1_{a}, \cdot_{a}(x, v)))) \Rightarrow  \\ \quad \quad \operatorname{vadd}(a)(v, \operatorname{vneg}(a)(\operatorname{vadd}(a)(1_{a}, x))) = \cdot_{a}(x, \cdot_{a}(x, v)) \end{array}}}
\vnbprop{neumann-inverse-norm-bound}{\vnbstmt{\begin{array}{@{}l@{}} \forall a.; \forall x \in \operatorname{vec}(a). \\ \quad (\operatorname{is\text{-}unital\text{-}banach\text{-}algebra}(a) \wedge (\operatorname{vnrm}(a)(x) < 1)) \Rightarrow  \\ \quad \quad \operatorname{vnrm}(a)({\operatorname{vadd}(a)(1_{a}, \operatorname{vneg}(a)(x))}^{-1}) \leq \operatorname{recip}((1 - \operatorname{vnrm}(a)(x))) \end{array}}}
\vnbprop{neumann-series-bound}{\vnbstmt{\begin{array}{@{}l@{}} \forall a.; \forall x \in \operatorname{vec}(a). \\ \quad (\operatorname{is\text{-}unital\text{-}banach\text{-}algebra}(a) \wedge (\operatorname{vnrm}(a)(x) < 1)) \Rightarrow  \\ \quad \quad \operatorname{vnrm}(a)(\operatorname{vadd}(a)({\operatorname{vadd}(a)(1_{a}, \operatorname{vneg}(a)(x))}^{-1}, \operatorname{vneg}(a)(\operatorname{vadd}(a)(1_{a}, x)))) \\ \quad \quad \leq\; (\operatorname{vnrm}(a)(x) \cdot \operatorname{vnrm}(a)(x)) \cdot \operatorname{recip}((1 - \operatorname{vnrm}(a)(x))) \end{array}}}
\vnbprop{bbi-inverse-diff}{\vnbstmt{\begin{array}{@{}l@{}} \forall a, b, v. \\ \quad (\operatorname{is\text{-}unital\text{-}banach\text{-}algebra}(a) \wedge (\operatorname{is\text{-}invertible\text{-}el}(a, b) \wedge \operatorname{is\text{-}invertible\text{-}el}(a, v))) \Rightarrow  \\ \quad \quad \operatorname{vadd}(a)({b}^{-1}, \operatorname{vneg}(a)({v}^{-1})) \\ \quad \quad =\; \cdot_{a}({v}^{-1}, \cdot_{a}(\operatorname{vadd}(a)(v, \operatorname{vneg}(a)(b)), {b}^{-1})) \end{array}}}
\vnbprop{inverse-diff-norm}{\vnbstmt{\begin{array}{@{}l@{}} \forall a, b, v. \\ \quad (\operatorname{is\text{-}unital\text{-}banach\text{-}algebra}(a) \wedge (\operatorname{is\text{-}invertible\text{-}el}(a, b) \wedge \operatorname{is\text{-}invertible\text{-}el}(a, v))) \Rightarrow  \\ \quad \quad \operatorname{vnrm}(a)(\operatorname{vadd}(a)({b}^{-1}, \operatorname{vneg}(a)({v}^{-1}))) \\ \quad \quad \leq\; (\operatorname{vnrm}(a)(\operatorname{vadd}(a)(b, \operatorname{vneg}(a)(v))) \cdot \operatorname{vnrm}(a)({b}^{-1})) \cdot \operatorname{vnrm}(a)({v}^{-1}) \end{array}}}
\vnbprop{inverse-near-bound}{\vnbstmt{\begin{array}{@{}l@{}} \forall a, b.; \forall v \in \operatorname{vec}(a). \\ \quad \operatorname{is\text{-}unital\text{-}banach\text{-}algebra}(a) \wedge \\ \quad \operatorname{is\text{-}invertible\text{-}el}(a, b) \wedge \\ \quad (\operatorname{vnrm}(a)(\operatorname{vadd}(a)(b, \operatorname{vneg}(a)(v))) \cdot \operatorname{vnrm}(a)({b}^{-1})) \leq 1/2 \Rightarrow \\ \quad \quad (\operatorname{is\text{-}invertible\text{-}el}(a, v) \wedge (\operatorname{vnrm}(a)({v}^{-1}) \leq (2 \cdot \operatorname{vnrm}(a)({b}^{-1})))) \end{array}}}
\vnbprop{inverse-lipschitz}{\vnbstmt{\begin{array}{@{}l@{}} \forall a, b.; \forall v \in \operatorname{vec}(a). \\ \quad \operatorname{is\text{-}unital\text{-}banach\text{-}algebra}(a) \wedge \\ \quad \operatorname{is\text{-}invertible\text{-}el}(a, b) \wedge \\ \quad (\operatorname{vnrm}(a)(\operatorname{vadd}(a)(b, \operatorname{vneg}(a)(v))) \cdot \operatorname{vnrm}(a)({b}^{-1})) \leq 1/2 \Rightarrow \\ \quad \quad \operatorname{vnrm}(a)(\operatorname{vadd}(a)({b}^{-1}, \operatorname{vneg}(a)({v}^{-1}))) \\ \quad \quad \leq\; \operatorname{vnrm}(a)(\operatorname{vadd}(a)(b, \operatorname{vneg}(a)(v))) \cdot (2 \cdot (\operatorname{vnrm}(a)({b}^{-1}) \cdot \operatorname{vnrm}(a)({b}^{-1}))) \end{array}}}
\vnbprop{inverse-continuous}{\vnbstmt{\begin{array}{@{}l@{}} \forall a, b, s. \\ \quad (\operatorname{is\text{-}unital\text{-}banach\text{-}algebra}(a) \wedge (\operatorname{is\text{-}invertible\text{-}el}(a, b) \wedge \operatorname{pos\text{-}rr}(s))) \Rightarrow  \\ \quad \quad \exists l. \\ \quad \quad \quad \operatorname{pos\text{-}rr}(l) \wedge \\ \quad \quad \quad \forall v \in \operatorname{vec}(a). \\ \quad \quad \quad \quad (\operatorname{vnrm}(a)(\operatorname{vadd}(a)(b, \operatorname{vneg}(a)(v))) \leq l) \Rightarrow  \\ \quad \quad \quad \quad \quad (\operatorname{is\text{-}invertible\text{-}el}(a, v) \wedge (\operatorname{vnrm}(a)(\operatorname{vadd}(a)({b}^{-1}, \operatorname{vneg}(a)({v}^{-1}))) \leq s)) \end{array}}}

\vnbbookfile{theorem-library/ba-ideal-closure}{theorem-library/ba-ideal-closure}
\noindent{\small ba-ideal-closure.scm -- wave of 2026-11-09/10 (agent BA-1B), wired 2026-11-10 from scratchpad/ba1b-1109/ba-ideal-closure.scm. Library: operator-algebras.}\par
\vnbprop{bic-closure-add-one}{\vnbstmt{\begin{array}{@{}l@{}} \forall a, i.; \forall q \in i, x \in \operatorname{closure}(\operatorname{ba\text{-}metric\text{-}space}(a), i). \\ \quad (\operatorname{is\text{-}unital\text{-}banach\text{-}algebra}(a) \wedge \operatorname{is\text{-}left\text{-}ideal}(a, i)) \Rightarrow  \\ \quad \quad \operatorname{vadd}(a)(x, q) \in \operatorname{closure}(\operatorname{ba\text{-}metric\text{-}space}(a), i) \end{array}}}
\vnbprop{bic-closure-add}{\vnbstmt{\begin{array}{@{}l@{}} \forall a, i.; \forall x, y \in \operatorname{closure}(\operatorname{ba\text{-}metric\text{-}space}(a), i). \\ \quad (\operatorname{is\text{-}unital\text{-}banach\text{-}algebra}(a) \wedge \operatorname{is\text{-}left\text{-}ideal}(a, i)) \Rightarrow  \\ \quad \quad \operatorname{vadd}(a)(x, y) \in \operatorname{closure}(\operatorname{ba\text{-}metric\text{-}space}(a), i) \end{array}}}
\vnbprop{bic-closure-act}{\vnbstmt{\begin{array}{@{}l@{}} \forall a, i.; \forall z \in \mathbb{C}, x \in \operatorname{closure}(\operatorname{ba\text{-}metric\text{-}space}(a), i). \\ \quad (\operatorname{is\text{-}unital\text{-}banach\text{-}algebra}(a) \wedge \operatorname{is\text{-}left\text{-}ideal}(a, i)) \Rightarrow  \\ \quad \quad \operatorname{act}(a)(z, x) \in \operatorname{closure}(\operatorname{ba\text{-}metric\text{-}space}(a), i) \end{array}}}
\vnbprop{bic-closure-mul}{\vnbstmt{\begin{array}{@{}l@{}} \forall a, i.; \forall y \in \operatorname{vec}(a), x \in \operatorname{closure}(\operatorname{ba\text{-}metric\text{-}space}(a), i). \\ \quad (\operatorname{is\text{-}unital\text{-}banach\text{-}algebra}(a) \wedge \operatorname{is\text{-}left\text{-}ideal}(a, i)) \Rightarrow  \\ \quad \quad \cdot_{a}(y, x) \in \operatorname{closure}(\operatorname{ba\text{-}metric\text{-}space}(a), i) \end{array}}}
\vnbprop{left-ideal-closure}{\vnbstmt{\begin{array}{@{}l@{}} \forall a, i. \\ \quad (\operatorname{is\text{-}unital\text{-}banach\text{-}algebra}(a) \wedge \operatorname{is\text{-}left\text{-}ideal}(a, i)) \Rightarrow  \\ \quad \quad \operatorname{is\text{-}left\text{-}ideal}(a, \operatorname{closure}(\operatorname{ba\text{-}metric\text{-}space}(a), i)) \end{array}}}
\vnbprop{left-ideal-one-all}{\vnbstmt{\begin{array}{@{}l@{}} \forall a, j. \\ \quad (\operatorname{is\text{-}unital\text{-}banach\text{-}algebra}(a) \wedge (\operatorname{is\text{-}left\text{-}ideal}(a, j) \wedge (1_{a} \in j))) \Rightarrow  \\ \quad \quad j = \operatorname{vec}(a) \end{array}}}
\vnbprop{maximal-left-ideal-closed}{\vnbstmt{\begin{array}{@{}l@{}} \forall a, i. \\ \quad (\operatorname{is\text{-}unital\text{-}banach\text{-}algebra}(a) \wedge \operatorname{is\text{-}maximal\text{-}left\text{-}ideal}(a, i)) \Rightarrow  \\ \quad \quad \operatorname{is\text{-}closed}(\operatorname{ba\text{-}metric\text{-}space}(a), i) \end{array}}}

\vnbbookfile{theorem-library/ba-resolvent}{theorem-library/ba-resolvent}
\noindent{\small ba-resolvent.scm -- wave of 2026-11-09/10 (agent BA-1B), wired 2026-11-10 from scratchpad/ba1b-1109/ba-resolvent.scm. Library: operator-algebras.}\par
\vnbprop{brv-inverse-one}{\vnbstmt{\begin{array}{@{}l@{}} \forall a. \\ \quad \operatorname{is\text{-}unital\text{-}banach\text{-}algebra}(a) \Rightarrow ({1_{a}}^{-1} = 1_{a}) \end{array}}}
\vnbprop{brv-act-inverse}{\vnbstmt{\begin{array}{@{}l@{}} \forall a.; \forall z \in \mathbb{C}, u. \\ \quad (\operatorname{is\text{-}unital\text{-}banach\text{-}algebra}(a) \wedge (\neg (z = 0) \wedge \operatorname{is\text{-}invertible\text{-}el}(a, u))) \Rightarrow  \\ \quad \quad {\operatorname{act}(a)(z, u)}^{-1} = \operatorname{act}(a)(\operatorname{recip}(z), {u}^{-1}) \end{array}}}
\vnbprop{brv-shift-factor}{\vnbstmt{\begin{array}{@{}l@{}} \forall a.; \forall x \in \operatorname{vec}(a), z \in \mathbb{C}. \\ \quad (\operatorname{is\text{-}unital\text{-}banach\text{-}algebra}(a) \wedge \neg (z = 0)) \Rightarrow  \\ \quad \quad \operatorname{shift\text{-}el}(a, z, x) = \operatorname{act}(a)(z, \operatorname{vadd}(a)(1_{a}, \operatorname{vneg}(a)(\operatorname{act}(a)(\operatorname{recip}(z), x)))) \end{array}}}
\vnbprop{brv-w-norm}{\vnbstmt{\begin{array}{@{}l@{}} \forall a.; \forall x \in \operatorname{vec}(a), z \in \mathbb{C}. \\ \quad (\operatorname{is\text{-}unital\text{-}banach\text{-}algebra}(a) \wedge \neg (z = 0)) \Rightarrow  \\ \quad \quad \operatorname{vnrm}(a)(\operatorname{act}(a)(\operatorname{recip}(z), x)) = (\operatorname{recip}(\operatorname{magnitude}(z)) \cdot \operatorname{vnrm}(a)(x)) \end{array}}}
\vnbprop{resolvent-norm-bound}{\vnbstmt{\begin{array}{@{}l@{}} \forall a.; \forall x \in \operatorname{vec}(a), z \in \mathbb{C}. \\ \quad (\operatorname{is\text{-}unital\text{-}banach\text{-}algebra}(a) \wedge ((0 < \operatorname{magnitude}(z)) \wedge ((2 \cdot \operatorname{vnrm}(a)(x)) \leq \operatorname{magnitude}(z)))) \Rightarrow  \\ \quad \quad \operatorname{is\text{-}invertible\text{-}el}(a, \operatorname{shift\text{-}el}(a, z, x)) \wedge \\ \quad \quad \operatorname{vnrm}(a)({\operatorname{shift\text{-}el}(a, z, x)}^{-1}) \leq (2 \cdot \operatorname{recip}(\operatorname{magnitude}(z))) \end{array}}}
\vnbprop{resolvent-continuous}{\vnbstmt{\begin{array}{@{}l@{}} \forall a.; \forall x \in \operatorname{vec}(a), u \in \mathbb{C}, s. \\ \quad (\operatorname{is\text{-}unital\text{-}banach\text{-}algebra}(a) \wedge (\operatorname{is\text{-}invertible\text{-}el}(a, \operatorname{shift\text{-}el}(a, u, x)) \wedge \operatorname{pos\text{-}rr}(s))) \Rightarrow  \\ \quad \quad \exists l. \\ \quad \quad \quad \operatorname{pos\text{-}rr}(l) \wedge \\ \quad \quad \quad \forall z \in \mathbb{C}. \\ \quad \quad \quad \quad (\operatorname{magnitude}((z + -u)) \leq l) \Rightarrow  \\ \quad \quad \quad \quad \quad \operatorname{is\text{-}invertible\text{-}el}(a, \operatorname{shift\text{-}el}(a, z, x)) \wedge \\ \quad \quad \quad \quad \quad \operatorname{vnrm}(a)(\operatorname{vadd}(a)({\operatorname{shift\text{-}el}(a, z, x)}^{-1}, \operatorname{vneg}(a)({\operatorname{shift\text{-}el}(a, u, x)}^{-1}))) \\ \quad \quad \quad \quad \quad \leq\; s \end{array}}}
\vnbprop{resolvent-derivative}{\vnbstmt{\begin{array}{@{}l@{}} \forall a.; \forall x \in \operatorname{vec}(a), u \in \mathbb{C}, s. \\ \quad (\operatorname{is\text{-}unital\text{-}banach\text{-}algebra}(a) \wedge (\operatorname{is\text{-}invertible\text{-}el}(a, \operatorname{shift\text{-}el}(a, u, x)) \wedge \operatorname{pos\text{-}rr}(s))) \Rightarrow  \\ \quad \quad \exists l. \\ \quad \quad \quad \operatorname{pos\text{-}rr}(l) \wedge \\ \quad \quad \quad \forall z \in \mathbb{C}. \\ \quad \quad \quad \quad (\operatorname{magnitude}((z + -u)) \leq l) \Rightarrow  \\ \quad \quad \quad \quad \quad \operatorname{is\text{-}invertible\text{-}el}(a, \operatorname{shift\text{-}el}(a, z, x)) \wedge \\ \quad \quad \quad \quad \quad \operatorname{vnrm}(a)(\operatorname{vadd}(a)(\operatorname{vadd}(a)({\operatorname{shift\text{-}el}(a, z, x)}^{-1}, \operatorname{vneg}(a)({\operatorname{shift\text{-}el}(a, u, x)}^{-1})), \\ \quad \quad \quad \quad \quad \quad \quad \operatorname{vneg}(a)(\operatorname{act}(a)((u + -z), \cdot_{a}({\operatorname{shift\text{-}el}(a, u, x)}^{-1}, {\operatorname{shift\text{-}el}(a, u, x)}^{-1}))))) \\ \quad \quad \quad \quad \quad \leq\; s \cdot \operatorname{magnitude}((z + -u)) \end{array}}}
\vnbprop{resolvent-laurent}{\vnbstmt{\begin{array}{@{}l@{}} \forall a.; \forall x \in \operatorname{vec}(a), z \in \mathbb{C}. \\ \quad (\operatorname{is\text{-}unital\text{-}banach\text{-}algebra}(a) \wedge (\operatorname{vnrm}(a)(x) < \operatorname{magnitude}(z))) \Rightarrow  \\ \quad \quad \operatorname{converges\text{-}to}(\operatorname{ba\text{-}metric\text{-}space}(a), \\ \quad \quad \quad (\lambda n \in \mathbb{N}.\; \operatorname{act}(a)(\operatorname{recip}(z), \sum_{k < n} {\operatorname{act}(a)(\operatorname{recip}(z), x)}^{k})), \\ \quad \quad \quad {\operatorname{shift\text{-}el}(a, z, x)}^{-1}) \end{array}}}

\vnbbookfile{theorem-library/ba-gelfand-mazur}{theorem-library/ba-gelfand-mazur}
\noindent{\small ba-gelfand-mazur.scm -- wave of 2026-11-09/10 (agent BA-1B), wired 2026-11-10 from scratchpad/ba1b-1109/ba-gelfand-mazur.scm. Library: operator-algebras.}\par
\vnbprop{gelfand-mazur-spectrum-point}{\vnbstmt{\begin{array}{@{}l@{}} \forall a.; \forall x \in \operatorname{vec}(a), z \in \sigma({x}). \\ \quad \operatorname{is\text{-}unital\text{-}banach\text{-}algebra}(a) \wedge \\ \quad \forall y \in \operatorname{vec}(a).\; (\neg (y = \operatorname{vzero}(a)) \Rightarrow \operatorname{is\text{-}invertible\text{-}el}(a, y)) \Rightarrow \\ \quad \quad x = \operatorname{act}(a)(z, 1_{a}) \end{array}}}
\vnbprop{scalar-embedding-injective}{\vnbstmt{\begin{array}{@{}l@{}} \forall a.; \forall z, w \in \mathbb{C}. \\ \quad (\operatorname{is\text{-}unital\text{-}banach\text{-}algebra}(a) \wedge (\operatorname{act}(a)(z, 1_{a}) = \operatorname{act}(a)(w, 1_{a}))) \Rightarrow  \\ \quad \quad z = w \end{array}}}
\vnbprop{gelfand-mazur-from-nonempty-spectrum}{\vnbstmt{\begin{array}{@{}l@{}} \forall a.; \forall x \in \operatorname{vec}(a). \\ \quad \operatorname{is\text{-}unital\text{-}banach\text{-}algebra}(a) \wedge \\ \quad \forall y \in \operatorname{vec}(a).\; (\neg (y = \operatorname{vzero}(a)) \Rightarrow \operatorname{is\text{-}invertible\text{-}el}(a, y)) \wedge \\ \quad \forall x \in \operatorname{vec}(a).\; \exists z.\; (z \in \sigma({x})) \Rightarrow \\ \quad \quad \exists z \in \mathbb{C}.\; (x = \operatorname{act}(a)(z, 1_{a})) \end{array}}}

\vnbbookfile{theorem-library/ba-gelfand}{theorem-library/ba-gelfand}
\noindent{\small ba-gelfand.scm -- wave of 2026-11-09/10 (agent BA-1B), wired 2026-11-10 from scratchpad/ba1b-1109/ba-gelfand.scm. Library: operator-algebras.}\par
\vnbprop{gelfand-transform-value}{\vnbstmt{\begin{array}{@{}l@{}} \forall a.; \forall x \in \operatorname{vec}(a), c \in \operatorname{characters}(a). \\ \quad \operatorname{is\text{-}unital\text{-}banach\text{-}algebra}(a) \Rightarrow (\operatorname{gelfand\text{-}transform}(a)(x)(c) = c(x)) \end{array}}}
\vnbprop{gelfand-transform-mul}{\vnbstmt{\begin{array}{@{}l@{}} \forall a.; \forall x, y \in \operatorname{vec}(a), c \in \operatorname{characters}(a). \\ \quad \operatorname{is\text{-}unital\text{-}banach\text{-}algebra}(a) \Rightarrow  \\ \quad \quad \operatorname{gelfand\text{-}transform}(a)(\cdot_{a}(x, y))(c) \\ \quad \quad =\; \operatorname{gelfand\text{-}transform}(a)(x)(c) \cdot \operatorname{gelfand\text{-}transform}(a)(y)(c) \end{array}}}
\vnbprop{gelfand-transform-add}{\vnbstmt{\begin{array}{@{}l@{}} \forall a.; \forall x, y \in \operatorname{vec}(a), c \in \operatorname{characters}(a). \\ \quad \operatorname{is\text{-}unital\text{-}banach\text{-}algebra}(a) \Rightarrow  \\ \quad \quad \operatorname{gelfand\text{-}transform}(a)(\operatorname{vadd}(a)(x, y))(c) \\ \quad \quad =\; \operatorname{gelfand\text{-}transform}(a)(x)(c) + \operatorname{gelfand\text{-}transform}(a)(y)(c) \end{array}}}
\vnbprop{gelfand-transform-act}{\vnbstmt{\begin{array}{@{}l@{}} \forall a.; \forall z \in \mathbb{C}, x \in \operatorname{vec}(a), c \in \operatorname{characters}(a). \\ \quad \operatorname{is\text{-}unital\text{-}banach\text{-}algebra}(a) \Rightarrow  \\ \quad \quad \operatorname{gelfand\text{-}transform}(a)(\operatorname{act}(a)(z, x))(c) = (z \cdot \operatorname{gelfand\text{-}transform}(a)(x)(c)) \end{array}}}
\vnbprop{gelfand-transform-one}{\vnbstmt{\begin{array}{@{}l@{}} \forall a.; \forall c \in \operatorname{characters}(a). \\ \quad \operatorname{is\text{-}unital\text{-}banach\text{-}algebra}(a) \Rightarrow (\operatorname{gelfand\text{-}transform}(a)(1_{a})(c) = 1) \end{array}}}
\vnbprop{gelfand-transform-bounded}{\vnbstmt{\begin{array}{@{}l@{}} \forall a.; \forall x \in \operatorname{vec}(a), c \in \operatorname{characters}(a). \\ \quad \operatorname{is\text{-}unital\text{-}banach\text{-}algebra}(a) \Rightarrow  \\ \quad \quad \operatorname{magnitude}(\operatorname{gelfand\text{-}transform}(a)(x)(c)) \leq \operatorname{vnrm}(a)(x) \end{array}}}
\vnbprop{character-value-in-spectrum}{\vnbstmt{\begin{array}{@{}l@{}} \forall a.; \forall x \in \operatorname{vec}(a), c \in \operatorname{characters}(a). \\ \quad \operatorname{is\text{-}unital\text{-}banach\text{-}algebra}(a) \Rightarrow (c(x) \in \sigma({x})) \end{array}}}

\vnbbookfile{theorem-library/ba-quotient}{theorem-library/ba-quotient}
\noindent{\small ba-quotient.scm -- wave of 2026-11-09/10 (agent BA-1C), wired 2026-11-10 from scratchpad/ba1c-1109/ba-quotient.scm. Library: operator-algebras.}\par
\vnbprop{quot-norm-unfold}{\vnbstmt{\begin{array}{@{}l@{}} \forall a, i, x. \\ \quad (\operatorname{quot\text{-}norm}(a, i, x) \simeq \operatorname{einf}(\{w \in \mathbb{R} : \exists y \in i.\; (w = \operatorname{vnrm}(a)(\operatorname{vadd}(a)(x, \operatorname{vneg}(a)(y))))\})) \end{array}}}
\vnbprop{ba-left-ideal-dist-one-ge}{\vnbstmt{\begin{array}{@{}l@{}} \forall a, i.; \forall y \in i. \\ \quad (\operatorname{is\text{-}unital\text{-}banach\text{-}algebra}(a) \wedge (\operatorname{is\text{-}left\text{-}ideal}(a, i) \wedge \neg (1_{a} \in i))) \Rightarrow  \\ \quad \quad 1 \leq \operatorname{vnrm}(a)(\operatorname{vadd}(a)(1_{a}, \operatorname{vneg}(a)(y))) \end{array}}}
\vnbprop{quot-norm-le-norm}{\vnbstmt{\begin{array}{@{}l@{}} \forall a, i.; \forall x \in \operatorname{vec}(a). \\ \quad (\operatorname{is\text{-}unital\text{-}banach\text{-}algebra}(a) \wedge \operatorname{is\text{-}left\text{-}ideal}(a, i)) \Rightarrow  \\ \quad \quad \operatorname{quot\text{-}norm}(a, i, x) \leq \operatorname{vnrm}(a)(x) \end{array}}}
\vnbprop{quot-norm-one}{\vnbstmt{\begin{array}{@{}l@{}} \forall a, i. \\ \quad (\operatorname{is\text{-}unital\text{-}banach\text{-}algebra}(a) \wedge (\operatorname{is\text{-}left\text{-}ideal}(a, i) \wedge \neg (1_{a} \in i))) \Rightarrow  \\ \quad \quad \operatorname{quot\text{-}norm}(a, i, 1_{a}) = 1 \end{array}}}

\vnbbookfile{theorem-library/ba-gelfand-continuous}{theorem-library/ba-gelfand-continuous}
\noindent{\small ba-gelfand-continuous.scm -- wave of 2026-11-09/10 (agent BA-1D), wired 2026-11-10 from scratchpad/ba1d-1109/ba-gelfand-continuous.scm. Library: operator-algebras.}\par
\vnbprop{character-in-fun}{\vnbstmt{\begin{array}{@{}l@{}} \forall a.; \forall c \in \operatorname{characters}(a). \\ \quad \operatorname{is\text{-}unital\text{-}banach\text{-}algebra}(a) \Rightarrow (c \in (\operatorname{vec}(a) \to \mathbb{C})) \end{array}}}
\vnbprop{gelfand-transform-continuous}{\vnbstmt{\begin{array}{@{}l@{}} \forall a.; \forall x \in \operatorname{vec}(a). \\ \quad \operatorname{is\text{-}unital\text{-}banach\text{-}algebra}(a) \Rightarrow  \\ \quad \quad \operatorname{is\text{-}hom\text{-}top\text{-}space}(\widehat{a}, \operatorname{metric\text{-}top}(\operatorname{cc\text{-}ms}), \operatorname{gelfand\text{-}transform}(a)(x)) \end{array}}}

\vnbbookfile{theorem-library/ba-quotient-na}{theorem-library/ba-quotient-na}
\noindent{\small ba-quotient-na.scm -- wave of 2026-11-09/10 (agent BA-1E), wired 2026-11-10 from scratchpad/ba1e-1109/ba-quotient-na.scm. Library: operator-algebras.}\par
\vnbprop{ba-ideal-proper-iff}{\vnbstmt{\begin{array}{@{}l@{}} \forall a, i. \\ \quad (\operatorname{is\text{-}unital\text{-}banach\text{-}algebra}(a) \wedge \operatorname{is\text{-}left\text{-}ideal}(a, i)) \Rightarrow  \\ \quad \quad \neg (1_{a} \in i) \Leftrightarrow \neg (i = \operatorname{vec}(a)) \end{array}}}
\vnbprop{bqn-uam-scal}{\vnbstmt{\begin{array}{@{}l@{}} \forall a. \\ \quad \operatorname{is\text{-}unital\text{-}banach\text{-}algebra}(a) \Rightarrow (\operatorname{scal}(\operatorname{uba\text{-}as\text{-}module}(a)) \simeq \operatorname{scal}(a)) \end{array}}}
\vnbprop{bqn-uam-vec}{\vnbstmt{\forall a.\; (\operatorname{is\text{-}unital\text{-}banach\text{-}algebra}(a) \Rightarrow (\operatorname{vec}(\operatorname{uba\text{-}as\text{-}module}(a)) \simeq \operatorname{vec}(a)))}}
\vnbprop{bqn-uam-vadd}{\vnbstmt{\begin{array}{@{}l@{}} \forall a. \\ \quad \operatorname{is\text{-}unital\text{-}banach\text{-}algebra}(a) \Rightarrow (\operatorname{vadd}(\operatorname{uba\text{-}as\text{-}module}(a)) \simeq \operatorname{vadd}(a)) \end{array}}}
\vnbprop{bqn-uam-vzero}{\vnbstmt{\begin{array}{@{}l@{}} \forall a. \\ \quad \operatorname{is\text{-}unital\text{-}banach\text{-}algebra}(a) \Rightarrow (\operatorname{vzero}(\operatorname{uba\text{-}as\text{-}module}(a)) \simeq \operatorname{vzero}(a)) \end{array}}}
\vnbprop{bqn-uam-vneg}{\vnbstmt{\begin{array}{@{}l@{}} \forall a. \\ \quad \operatorname{is\text{-}unital\text{-}banach\text{-}algebra}(a) \Rightarrow (\operatorname{vneg}(\operatorname{uba\text{-}as\text{-}module}(a)) \simeq \operatorname{vneg}(a)) \end{array}}}
\vnbprop{bqn-uam-act}{\vnbstmt{\forall a.\; (\operatorname{is\text{-}unital\text{-}banach\text{-}algebra}(a) \Rightarrow (\operatorname{act}(\operatorname{uba\text{-}as\text{-}module}(a)) \simeq \operatorname{act}(a)))}}
\vnbprop{ba-left-ideal-submodule}{\vnbstmt{\begin{array}{@{}l@{}} \forall a, i. \\ \quad (\operatorname{is\text{-}unital\text{-}banach\text{-}algebra}(a) \wedge \operatorname{is\text{-}left\text{-}ideal}(a, i)) \Rightarrow  \\ \quad \quad \operatorname{is\text{-}submodule}(\operatorname{uba\text{-}as\text{-}module}(a), i) \end{array}}}
\vnbprop{bqn-mul-neg-right}{\vnbstmt{\begin{array}{@{}l@{}} \forall a.; \forall x, y \in \operatorname{vec}(a). \\ \quad \operatorname{is\text{-}unital\text{-}banach\text{-}algebra}(a) \Rightarrow  \\ \quad \quad \cdot_{a}(x, \operatorname{vneg}(a)(y)) = \operatorname{vneg}(a)(\cdot_{a}(x, y)) \end{array}}}
\vnbprop{qna-lvm-fun}{\vnbstmt{\begin{array}{@{}l@{}} \forall a, i. \\ \quad (\operatorname{is\text{-}unital\text{-}banach\text{-}algebra}(a) \wedge \operatorname{is\text{-}two\text{-}sided\text{-}ideal}(a, i)) \Rightarrow  \\ \quad \quad (\lambda {[}\operatorname{bqnx\_}, \operatorname{bqny\_}] \in (\operatorname{vec}(a) \times \operatorname{vec}(a)).\; \\ \quad \quad \quad \operatorname{class}(\operatorname{qmod\text{-}setoid}(\operatorname{uba\text{-}as\text{-}module}(a), i), \cdot_{a}(\operatorname{bqnx\_}, \operatorname{bqny\_}))) \\ \quad \quad \in\; ((\operatorname{pts}(\operatorname{qmod\text{-}setoid}(\operatorname{uba\text{-}as\text{-}module}(a), i)) \times \operatorname{pts}(\operatorname{qmod\text{-}setoid}(\operatorname{uba\text{-}as\text{-}module}(a), i))) \to \operatorname{quotient}(\operatorname{qmod\text{-}setoid}(\operatorname{uba\text{-}as\text{-}module}(a), i))) \end{array}}}
\vnbprop{qna-lvm-respects}{\vnbstmt{\begin{array}{@{}l@{}} \forall a, i. \\ \quad (\operatorname{is\text{-}unital\text{-}banach\text{-}algebra}(a) \wedge \operatorname{is\text{-}two\text{-}sided\text{-}ideal}(a, i)) \Rightarrow  \\ \quad \quad \operatorname{respects2}(\operatorname{qmod\text{-}setoid}(\operatorname{uba\text{-}as\text{-}module}(a), i), \\ \quad \quad \quad (\lambda {[}\operatorname{bqnx\_}, \operatorname{bqny\_}] \in (\operatorname{vec}(a) \times \operatorname{vec}(a)).\; \\ \quad \quad \quad \quad \operatorname{class}(\operatorname{qmod\text{-}setoid}(\operatorname{uba\text{-}as\text{-}module}(a), i), \cdot_{a}(\operatorname{bqnx\_}, \operatorname{bqny\_})))) \end{array}}}
\vnbprop{qna-vmul-fun}{\vnbstmt{\begin{array}{@{}l@{}} \forall a, i. \\ \quad (\operatorname{is\text{-}unital\text{-}banach\text{-}algebra}(a) \wedge \operatorname{is\text{-}two\text{-}sided\text{-}ideal}(a, i)) \Rightarrow  \\ \quad \quad \operatorname{qna\text{-}vmul}(a, i) \\ \quad \quad \in\; ((\operatorname{quotient}(\operatorname{qmod\text{-}setoid}(\operatorname{uba\text{-}as\text{-}module}(a), i)) \times \operatorname{quotient}(\operatorname{qmod\text{-}setoid}(\operatorname{uba\text{-}as\text{-}module}(a), i))) \to \operatorname{quotient}(\operatorname{qmod\text{-}setoid}(\operatorname{uba\text{-}as\text{-}module}(a), i))) \end{array}}}
\vnbprop{qna-vmul-computes}{\vnbstmt{\begin{array}{@{}l@{}} \forall a, i.; \forall x, y \in \operatorname{vec}(a). \\ \quad (\operatorname{is\text{-}unital\text{-}banach\text{-}algebra}(a) \wedge \operatorname{is\text{-}two\text{-}sided\text{-}ideal}(a, i)) \Rightarrow  \\ \quad \quad \operatorname{qna\text{-}vmul}(a, i)(\operatorname{class}(\operatorname{qmod\text{-}setoid}(\operatorname{uba\text{-}as\text{-}module}(a), i), x), \\ \quad \quad \quad \operatorname{class}(\operatorname{qmod\text{-}setoid}(\operatorname{uba\text{-}as\text{-}module}(a), i), y)) \\ \quad \quad =\; \operatorname{class}(\operatorname{qmod\text{-}setoid}(\operatorname{uba\text{-}as\text{-}module}(a), i), \cdot_{a}(x, y)) \end{array}}}
\vnbprop{coset-norm-le}{\vnbstmt{\begin{array}{@{}l@{}} \forall a, i.; \forall c \in \operatorname{quotient}(\operatorname{qmod\text{-}setoid}(\operatorname{uba\text{-}as\text{-}module}(a), i)), u \in c. \\ \quad (\operatorname{is\text{-}unital\text{-}banach\text{-}algebra}(a) \wedge \operatorname{is\text{-}two\text{-}sided\text{-}ideal}(a, i)) \Rightarrow  \\ \quad \quad \operatorname{coset\text{-}norm}(a, i, c) \leq \operatorname{vnrm}(a)(u) \end{array}}}
\vnbprop{coset-norm-real}{\vnbstmt{\begin{array}{@{}l@{}} \forall a, i.; \forall c \in \operatorname{quotient}(\operatorname{qmod\text{-}setoid}(\operatorname{uba\text{-}as\text{-}module}(a), i)). \\ \quad (\operatorname{is\text{-}unital\text{-}banach\text{-}algebra}(a) \wedge \operatorname{is\text{-}two\text{-}sided\text{-}ideal}(a, i)) \Rightarrow  \\ \quad \quad (\operatorname{coset\text{-}norm}(a, i, c) \in \mathbb{R}) \wedge (0 \leq \operatorname{coset\text{-}norm}(a, i, c)) \end{array}}}
\vnbprop{coset-norm-approx}{\vnbstmt{\begin{array}{@{}l@{}} \forall a, i.; \forall c \in \operatorname{quotient}(\operatorname{qmod\text{-}setoid}(\operatorname{uba\text{-}as\text{-}module}(a), i)), s. \\ \quad (\operatorname{is\text{-}unital\text{-}banach\text{-}algebra}(a) \wedge (\operatorname{is\text{-}two\text{-}sided\text{-}ideal}(a, i) \wedge \operatorname{pos\text{-}rr}(s))) \Rightarrow  \\ \quad \quad \exists u \in c.\; (\operatorname{vnrm}(a)(u) < (\operatorname{coset\text{-}norm}(a, i, c) + s)) \end{array}}}
\vnbprop{coset-norm-scale-greatest}{\vnbstmt{\begin{array}{@{}l@{}} \forall a, i.; \forall c \in \operatorname{quotient}(\operatorname{qmod\text{-}setoid}(\operatorname{uba\text{-}as\text{-}module}(a), i)), t, b \in \mathbb{R}. \\ \quad \operatorname{is\text{-}unital\text{-}banach\text{-}algebra}(a) \wedge \\ \quad \operatorname{is\text{-}two\text{-}sided\text{-}ideal}(a, i) \wedge \\ \quad 0 < t \wedge \\ \quad 0 \leq b \wedge \\ \quad \forall u \in c.\; (b \leq (t \cdot \operatorname{vnrm}(a)(u))) \Rightarrow \\ \quad \quad b \leq (t \cdot \operatorname{coset\text{-}norm}(a, i, c)) \end{array}}}
\vnbprop{coset-norm-triangle}{\vnbstmt{\begin{array}{@{}l@{}} \forall a, i.; \forall c, d \in \operatorname{quotient}(\operatorname{qmod\text{-}setoid}(\operatorname{uba\text{-}as\text{-}module}(a), i)). \\ \quad (\operatorname{is\text{-}unital\text{-}banach\text{-}algebra}(a) \wedge \operatorname{is\text{-}two\text{-}sided\text{-}ideal}(a, i)) \Rightarrow  \\ \quad \quad \operatorname{coset\text{-}norm}(a, i, \operatorname{vadd}(\operatorname{quotient\text{-}module}(\operatorname{uba\text{-}as\text{-}module}(a), i))(c, d)) \\ \quad \quad \leq\; \operatorname{coset\text{-}norm}(a, i, c) + \operatorname{coset\text{-}norm}(a, i, d) \end{array}}}
\vnbprop{coset-norm-submult}{\vnbstmt{\begin{array}{@{}l@{}} \forall a, i.; \forall c, d \in \operatorname{quotient}(\operatorname{qmod\text{-}setoid}(\operatorname{uba\text{-}as\text{-}module}(a), i)). \\ \quad (\operatorname{is\text{-}unital\text{-}banach\text{-}algebra}(a) \wedge \operatorname{is\text{-}two\text{-}sided\text{-}ideal}(a, i)) \Rightarrow  \\ \quad \quad \operatorname{coset\text{-}norm}(a, i, \operatorname{qna\text{-}vmul}(a, i)(c, d)) \leq (\operatorname{coset\text{-}norm}(a, i, c) \cdot \operatorname{coset\text{-}norm}(a, i, d)) \end{array}}}
\vnbprop{coset-norm-act-le}{\vnbstmt{\begin{array}{@{}l@{}} \forall a, i.; \forall c \in \operatorname{quotient}(\operatorname{qmod\text{-}setoid}(\operatorname{uba\text{-}as\text{-}module}(a), i)), m \in \mathbb{C}. \\ \quad (\operatorname{is\text{-}unital\text{-}banach\text{-}algebra}(a) \wedge \operatorname{is\text{-}two\text{-}sided\text{-}ideal}(a, i)) \Rightarrow  \\ \quad \quad \operatorname{coset\text{-}norm}(a, i, \operatorname{act}(\operatorname{quotient\text{-}module}(\operatorname{uba\text{-}as\text{-}module}(a), i))(m, c)) \\ \quad \quad \leq\; \operatorname{magnitude}(m) \cdot \operatorname{coset\text{-}norm}(a, i, c) \end{array}}}
\vnbprop{coset-norm-homog}{\vnbstmt{\begin{array}{@{}l@{}} \forall a, i.; \forall c \in \operatorname{quotient}(\operatorname{qmod\text{-}setoid}(\operatorname{uba\text{-}as\text{-}module}(a), i)), m \in \mathbb{C}. \\ \quad (\operatorname{is\text{-}unital\text{-}banach\text{-}algebra}(a) \wedge \operatorname{is\text{-}two\text{-}sided\text{-}ideal}(a, i)) \Rightarrow  \\ \quad \quad \operatorname{coset\text{-}norm}(a, i, \operatorname{act}(\operatorname{quotient\text{-}module}(\operatorname{uba\text{-}as\text{-}module}(a), i))(m, c)) \\ \quad \quad =\; \operatorname{magnitude}(m) \cdot \operatorname{coset\text{-}norm}(a, i, c) \end{array}}}
\vnbprop{coset-norm-definite}{\vnbstmt{\begin{array}{@{}l@{}} \forall a, i.; \forall c \in \operatorname{quotient}(\operatorname{qmod\text{-}setoid}(\operatorname{uba\text{-}as\text{-}module}(a), i)). \\ \quad (\operatorname{is\text{-}unital\text{-}banach\text{-}algebra}(a) \wedge (\operatorname{is\text{-}two\text{-}sided\text{-}ideal}(a, i) \wedge \operatorname{is\text{-}closed}(\operatorname{ba\text{-}metric\text{-}space}(a), i))) \Rightarrow  \\ \quad \quad (\operatorname{coset\text{-}norm}(a, i, c) = 0) \Leftrightarrow (c = \operatorname{class}(\operatorname{qmod\text{-}setoid}(\operatorname{uba\text{-}as\text{-}module}(a), i), \operatorname{vzero}(a))) \end{array}}}
\vnbprop{qna-vmul-assoc}{\vnbstmt{\begin{array}{@{}l@{}} \forall a, i.; \forall c, d, e \in \operatorname{quotient}(\operatorname{qmod\text{-}setoid}(\operatorname{uba\text{-}as\text{-}module}(a), i)). \\ \quad (\operatorname{is\text{-}unital\text{-}banach\text{-}algebra}(a) \wedge \operatorname{is\text{-}two\text{-}sided\text{-}ideal}(a, i)) \Rightarrow  \\ \quad \quad \operatorname{qna\text{-}vmul}(a, i)(\operatorname{qna\text{-}vmul}(a, i)(c, d), e) = \operatorname{qna\text{-}vmul}(a, i)(c, \operatorname{qna\text{-}vmul}(a, i)(d, e)) \end{array}}}
\vnbprop{qna-vmul-dist-left}{\vnbstmt{\begin{array}{@{}l@{}} \forall a, i.; \forall c, d, e \in \operatorname{quotient}(\operatorname{qmod\text{-}setoid}(\operatorname{uba\text{-}as\text{-}module}(a), i)). \\ \quad (\operatorname{is\text{-}unital\text{-}banach\text{-}algebra}(a) \wedge \operatorname{is\text{-}two\text{-}sided\text{-}ideal}(a, i)) \Rightarrow  \\ \quad \quad \operatorname{qna\text{-}vmul}(a, i)(c, \operatorname{vadd}(\operatorname{quotient\text{-}module}(\operatorname{uba\text{-}as\text{-}module}(a), i))(d, e)) \\ \quad \quad =\; \operatorname{vadd}(\operatorname{quotient\text{-}module}(\operatorname{uba\text{-}as\text{-}module}(a), i))(\operatorname{qna\text{-}vmul}(a, i)(c, d), \operatorname{qna\text{-}vmul}(a, i)(c, e)) \end{array}}}
\vnbprop{qna-vmul-dist-right}{\vnbstmt{\begin{array}{@{}l@{}} \forall a, i.; \forall c, d, e \in \operatorname{quotient}(\operatorname{qmod\text{-}setoid}(\operatorname{uba\text{-}as\text{-}module}(a), i)). \\ \quad (\operatorname{is\text{-}unital\text{-}banach\text{-}algebra}(a) \wedge \operatorname{is\text{-}two\text{-}sided\text{-}ideal}(a, i)) \Rightarrow  \\ \quad \quad \operatorname{qna\text{-}vmul}(a, i)(\operatorname{vadd}(\operatorname{quotient\text{-}module}(\operatorname{uba\text{-}as\text{-}module}(a), i))(c, d), e) \\ \quad \quad =\; \operatorname{vadd}(\operatorname{quotient\text{-}module}(\operatorname{uba\text{-}as\text{-}module}(a), i))(\operatorname{qna\text{-}vmul}(a, i)(c, e), \operatorname{qna\text{-}vmul}(a, i)(d, e)) \end{array}}}
\vnbprop{qna-vmul-act-left}{\vnbstmt{\begin{array}{@{}l@{}} \forall a, i.; \forall m \in \mathbb{C}, c, d \in \operatorname{quotient}(\operatorname{qmod\text{-}setoid}(\operatorname{uba\text{-}as\text{-}module}(a), i)). \\ \quad (\operatorname{is\text{-}unital\text{-}banach\text{-}algebra}(a) \wedge \operatorname{is\text{-}two\text{-}sided\text{-}ideal}(a, i)) \Rightarrow  \\ \quad \quad \operatorname{qna\text{-}vmul}(a, i)(\operatorname{act}(\operatorname{quotient\text{-}module}(\operatorname{uba\text{-}as\text{-}module}(a), i))(m, c), d) \\ \quad \quad =\; \operatorname{act}(\operatorname{quotient\text{-}module}(\operatorname{uba\text{-}as\text{-}module}(a), i))(m, \operatorname{qna\text{-}vmul}(a, i)(c, d)) \end{array}}}
\vnbprop{qna-vmul-act-right}{\vnbstmt{\begin{array}{@{}l@{}} \forall a, i.; \forall m \in \mathbb{C}, c, d \in \operatorname{quotient}(\operatorname{qmod\text{-}setoid}(\operatorname{uba\text{-}as\text{-}module}(a), i)). \\ \quad (\operatorname{is\text{-}unital\text{-}banach\text{-}algebra}(a) \wedge \operatorname{is\text{-}two\text{-}sided\text{-}ideal}(a, i)) \Rightarrow  \\ \quad \quad \operatorname{qna\text{-}vmul}(a, i)(c, \operatorname{act}(\operatorname{quotient\text{-}module}(\operatorname{uba\text{-}as\text{-}module}(a), i))(m, d)) \\ \quad \quad =\; \operatorname{act}(\operatorname{quotient\text{-}module}(\operatorname{uba\text{-}as\text{-}module}(a), i))(m, \operatorname{qna\text{-}vmul}(a, i)(c, d)) \end{array}}}
\vnbprop{qna-vnrm-apply}{\vnbstmt{\begin{array}{@{}l@{}} \forall a, i.; \forall c \in \operatorname{quotient}(\operatorname{qmod\text{-}setoid}(\operatorname{uba\text{-}as\text{-}module}(a), i)). \\ \quad (\operatorname{is\text{-}unital\text{-}banach\text{-}algebra}(a) \wedge \operatorname{is\text{-}two\text{-}sided\text{-}ideal}(a, i)) \Rightarrow  \\ \quad \quad \operatorname{qna\text{-}vnrm}(a, i)(c) = \operatorname{coset\text{-}norm}(a, i, c) \end{array}}}
\vnbprop{qna-vnrm-fun}{\vnbstmt{\begin{array}{@{}l@{}} \forall a, i. \\ \quad (\operatorname{is\text{-}unital\text{-}banach\text{-}algebra}(a) \wedge \operatorname{is\text{-}two\text{-}sided\text{-}ideal}(a, i)) \Rightarrow  \\ \quad \quad \operatorname{qna\text{-}vnrm}(a, i) \in (\operatorname{quotient}(\operatorname{qmod\text{-}setoid}(\operatorname{uba\text{-}as\text{-}module}(a), i)) \to \mathbb{R}) \end{array}}}
\vnbprop{qna-ro-scal}{\vnbstmt{\begin{array}{@{}l@{}} \forall a, i. \\ \quad (\operatorname{is\text{-}unital\text{-}banach\text{-}algebra}(a) \wedge \operatorname{is\text{-}two\text{-}sided\text{-}ideal}(a, i)) \Rightarrow  \\ \quad \quad (\operatorname{scal}({a}/{i}) \simeq \operatorname{scal}(\operatorname{quotient\text{-}module}(\operatorname{uba\text{-}as\text{-}module}(a), i))) \end{array}}}
\vnbprop{qba-ro-scal}{\vnbstmt{\begin{array}{@{}l@{}} \forall a, i. \\ \quad (\operatorname{is\text{-}unital\text{-}banach\text{-}algebra}(a) \wedge \operatorname{is\text{-}two\text{-}sided\text{-}ideal}(a, i)) \Rightarrow  \\ \quad \quad (\operatorname{scal}({a}/{i}) \simeq \operatorname{scal}(\operatorname{quotient\text{-}module}(\operatorname{uba\text{-}as\text{-}module}(a), i))) \end{array}}}
\vnbprop{qna-ro-vec}{\vnbstmt{\begin{array}{@{}l@{}} \forall a, i. \\ \quad (\operatorname{is\text{-}unital\text{-}banach\text{-}algebra}(a) \wedge \operatorname{is\text{-}two\text{-}sided\text{-}ideal}(a, i)) \Rightarrow  \\ \quad \quad (\operatorname{vec}({a}/{i}) \simeq \operatorname{vec}(\operatorname{quotient\text{-}module}(\operatorname{uba\text{-}as\text{-}module}(a), i))) \end{array}}}
\vnbprop{qba-ro-vec}{\vnbstmt{\begin{array}{@{}l@{}} \forall a, i. \\ \quad (\operatorname{is\text{-}unital\text{-}banach\text{-}algebra}(a) \wedge \operatorname{is\text{-}two\text{-}sided\text{-}ideal}(a, i)) \Rightarrow  \\ \quad \quad (\operatorname{vec}({a}/{i}) \simeq \operatorname{vec}(\operatorname{quotient\text{-}module}(\operatorname{uba\text{-}as\text{-}module}(a), i))) \end{array}}}
\vnbprop{qna-ro-vadd}{\vnbstmt{\begin{array}{@{}l@{}} \forall a, i. \\ \quad (\operatorname{is\text{-}unital\text{-}banach\text{-}algebra}(a) \wedge \operatorname{is\text{-}two\text{-}sided\text{-}ideal}(a, i)) \Rightarrow  \\ \quad \quad (\operatorname{vadd}({a}/{i}) \simeq \operatorname{vadd}(\operatorname{quotient\text{-}module}(\operatorname{uba\text{-}as\text{-}module}(a), i))) \end{array}}}
\vnbprop{qba-ro-vadd}{\vnbstmt{\begin{array}{@{}l@{}} \forall a, i. \\ \quad (\operatorname{is\text{-}unital\text{-}banach\text{-}algebra}(a) \wedge \operatorname{is\text{-}two\text{-}sided\text{-}ideal}(a, i)) \Rightarrow  \\ \quad \quad (\operatorname{vadd}({a}/{i}) \simeq \operatorname{vadd}(\operatorname{quotient\text{-}module}(\operatorname{uba\text{-}as\text{-}module}(a), i))) \end{array}}}
\vnbprop{qna-ro-vzero}{\vnbstmt{\begin{array}{@{}l@{}} \forall a, i. \\ \quad (\operatorname{is\text{-}unital\text{-}banach\text{-}algebra}(a) \wedge \operatorname{is\text{-}two\text{-}sided\text{-}ideal}(a, i)) \Rightarrow  \\ \quad \quad (\operatorname{vzero}({a}/{i}) \simeq \operatorname{vzero}(\operatorname{quotient\text{-}module}(\operatorname{uba\text{-}as\text{-}module}(a), i))) \end{array}}}
\vnbprop{qba-ro-vzero}{\vnbstmt{\begin{array}{@{}l@{}} \forall a, i. \\ \quad (\operatorname{is\text{-}unital\text{-}banach\text{-}algebra}(a) \wedge \operatorname{is\text{-}two\text{-}sided\text{-}ideal}(a, i)) \Rightarrow  \\ \quad \quad (\operatorname{vzero}({a}/{i}) \simeq \operatorname{vzero}(\operatorname{quotient\text{-}module}(\operatorname{uba\text{-}as\text{-}module}(a), i))) \end{array}}}
\vnbprop{qna-ro-vneg}{\vnbstmt{\begin{array}{@{}l@{}} \forall a, i. \\ \quad (\operatorname{is\text{-}unital\text{-}banach\text{-}algebra}(a) \wedge \operatorname{is\text{-}two\text{-}sided\text{-}ideal}(a, i)) \Rightarrow  \\ \quad \quad (\operatorname{vneg}({a}/{i}) \simeq \operatorname{vneg}(\operatorname{quotient\text{-}module}(\operatorname{uba\text{-}as\text{-}module}(a), i))) \end{array}}}
\vnbprop{qba-ro-vneg}{\vnbstmt{\begin{array}{@{}l@{}} \forall a, i. \\ \quad (\operatorname{is\text{-}unital\text{-}banach\text{-}algebra}(a) \wedge \operatorname{is\text{-}two\text{-}sided\text{-}ideal}(a, i)) \Rightarrow  \\ \quad \quad (\operatorname{vneg}({a}/{i}) \simeq \operatorname{vneg}(\operatorname{quotient\text{-}module}(\operatorname{uba\text{-}as\text{-}module}(a), i))) \end{array}}}
\vnbprop{qna-ro-act}{\vnbstmt{\begin{array}{@{}l@{}} \forall a, i. \\ \quad (\operatorname{is\text{-}unital\text{-}banach\text{-}algebra}(a) \wedge \operatorname{is\text{-}two\text{-}sided\text{-}ideal}(a, i)) \Rightarrow  \\ \quad \quad (\operatorname{act}({a}/{i}) \simeq \operatorname{act}(\operatorname{quotient\text{-}module}(\operatorname{uba\text{-}as\text{-}module}(a), i))) \end{array}}}
\vnbprop{qba-ro-act}{\vnbstmt{\begin{array}{@{}l@{}} \forall a, i. \\ \quad (\operatorname{is\text{-}unital\text{-}banach\text{-}algebra}(a) \wedge \operatorname{is\text{-}two\text{-}sided\text{-}ideal}(a, i)) \Rightarrow  \\ \quad \quad (\operatorname{act}({a}/{i}) \simeq \operatorname{act}(\operatorname{quotient\text{-}module}(\operatorname{uba\text{-}as\text{-}module}(a), i))) \end{array}}}
\vnbprop{qna-ro-vnrm}{\vnbstmt{\begin{array}{@{}l@{}} \forall a, i. \\ \quad (\operatorname{is\text{-}unital\text{-}banach\text{-}algebra}(a) \wedge \operatorname{is\text{-}two\text{-}sided\text{-}ideal}(a, i)) \Rightarrow  \\ \quad \quad (\operatorname{vnrm}({a}/{i}) \simeq \operatorname{qna\text{-}vnrm}(a, i)) \end{array}}}
\vnbprop{qba-ro-vnrm}{\vnbstmt{\begin{array}{@{}l@{}} \forall a, i. \\ \quad (\operatorname{is\text{-}unital\text{-}banach\text{-}algebra}(a) \wedge \operatorname{is\text{-}two\text{-}sided\text{-}ideal}(a, i)) \Rightarrow  \\ \quad \quad (\operatorname{vnrm}({a}/{i}) \simeq \operatorname{qna\text{-}vnrm}(a, i)) \end{array}}}
\vnbprop{qna-ro-vmul}{\vnbstmt{\begin{array}{@{}l@{}} \forall a, i. \\ \quad (\operatorname{is\text{-}unital\text{-}banach\text{-}algebra}(a) \wedge \operatorname{is\text{-}two\text{-}sided\text{-}ideal}(a, i)) \Rightarrow  \\ \quad \quad (\cdot_{{a}/{i}} \simeq \operatorname{qna\text{-}vmul}(a, i)) \end{array}}}
\vnbprop{qba-ro-vmul}{\vnbstmt{\begin{array}{@{}l@{}} \forall a, i. \\ \quad (\operatorname{is\text{-}unital\text{-}banach\text{-}algebra}(a) \wedge \operatorname{is\text{-}two\text{-}sided\text{-}ideal}(a, i)) \Rightarrow  \\ \quad \quad (\cdot_{{a}/{i}} \simeq \operatorname{qna\text{-}vmul}(a, i)) \end{array}}}
\vnbprop{qba-ro-vone}{\vnbstmt{\begin{array}{@{}l@{}} \forall a, i. \\ \quad (\operatorname{is\text{-}unital\text{-}banach\text{-}algebra}(a) \wedge \operatorname{is\text{-}two\text{-}sided\text{-}ideal}(a, i)) \Rightarrow  \\ \quad \quad (1_{{a}/{i}} \simeq \operatorname{class}(\operatorname{qmod\text{-}setoid}(\operatorname{uba\text{-}as\text{-}module}(a), i), 1_{a})) \end{array}}}
\vnbprop{qna-vnrm-nonneg}{\vnbstmt{\begin{array}{@{}l@{}} \forall a, i.; \forall x \in \operatorname{vec}(\operatorname{quotient\text{-}module}(\operatorname{uba\text{-}as\text{-}module}(a), i)). \\ \quad (\operatorname{is\text{-}unital\text{-}banach\text{-}algebra}(a) \wedge \operatorname{is\text{-}two\text{-}sided\text{-}ideal}(a, i)) \Rightarrow  \\ \quad \quad 0 \leq \operatorname{qna\text{-}vnrm}(a, i)(x) \end{array}}}
\vnbprop{qna-vnrm-homog}{\vnbstmt{\begin{array}{@{}l@{}} \forall a, i.; \forall r \in \operatorname{carr}(\operatorname{scal}(\operatorname{quotient\text{-}module}(\operatorname{uba\text{-}as\text{-}module}(a), i))), x \in \operatorname{vec}(\operatorname{quotient\text{-}module}(\operatorname{uba\text{-}as\text{-}module}(a), i)). \\ \quad (\operatorname{is\text{-}unital\text{-}banach\text{-}algebra}(a) \wedge \operatorname{is\text{-}two\text{-}sided\text{-}ideal}(a, i)) \Rightarrow  \\ \quad \quad \operatorname{qna\text{-}vnrm}(a, i)(\operatorname{act}(\operatorname{quotient\text{-}module}(\operatorname{uba\text{-}as\text{-}module}(a), i))(r, x)) \\ \quad \quad =\; \operatorname{magnitude}(r) \cdot \operatorname{qna\text{-}vnrm}(a, i)(x) \end{array}}}
\vnbprop{qna-vnrm-triangle}{\vnbstmt{\begin{array}{@{}l@{}} \forall a, i.; \forall x, y \in \operatorname{vec}(\operatorname{quotient\text{-}module}(\operatorname{uba\text{-}as\text{-}module}(a), i)). \\ \quad (\operatorname{is\text{-}unital\text{-}banach\text{-}algebra}(a) \wedge \operatorname{is\text{-}two\text{-}sided\text{-}ideal}(a, i)) \Rightarrow  \\ \quad \quad \operatorname{qna\text{-}vnrm}(a, i)(\operatorname{vadd}(\operatorname{quotient\text{-}module}(\operatorname{uba\text{-}as\text{-}module}(a), i))(x, y)) \\ \quad \quad \leq\; \operatorname{qna\text{-}vnrm}(a, i)(x) + \operatorname{qna\text{-}vnrm}(a, i)(y) \end{array}}}
\vnbprop{qna-vnrm-submult}{\vnbstmt{\begin{array}{@{}l@{}} \forall a, i.; \forall x, y \in \operatorname{vec}(\operatorname{quotient\text{-}module}(\operatorname{uba\text{-}as\text{-}module}(a), i)). \\ \quad (\operatorname{is\text{-}unital\text{-}banach\text{-}algebra}(a) \wedge \operatorname{is\text{-}two\text{-}sided\text{-}ideal}(a, i)) \Rightarrow  \\ \quad \quad \operatorname{qna\text{-}vnrm}(a, i)(\operatorname{qna\text{-}vmul}(a, i)(x, y)) \leq (\operatorname{qna\text{-}vnrm}(a, i)(x) \cdot \operatorname{qna\text{-}vnrm}(a, i)(y)) \end{array}}}
\vnbprop{qna-vmul-dist-left-na}{\vnbstmt{\begin{array}{@{}l@{}} \forall a, i.; \forall x, y, z \in \operatorname{vec}(\operatorname{quotient\text{-}module}(\operatorname{uba\text{-}as\text{-}module}(a), i)). \\ \quad (\operatorname{is\text{-}unital\text{-}banach\text{-}algebra}(a) \wedge \operatorname{is\text{-}two\text{-}sided\text{-}ideal}(a, i)) \Rightarrow  \\ \quad \quad \operatorname{qna\text{-}vmul}(a, i)(x, \operatorname{vadd}(\operatorname{quotient\text{-}module}(\operatorname{uba\text{-}as\text{-}module}(a), i))(y, z)) \\ \quad \quad =\; \operatorname{vadd}(\operatorname{quotient\text{-}module}(\operatorname{uba\text{-}as\text{-}module}(a), i))(\operatorname{qna\text{-}vmul}(a, i)(x, y), \operatorname{qna\text{-}vmul}(a, i)(x, z)) \end{array}}}
\vnbprop{qna-vmul-dist-right-na}{\vnbstmt{\begin{array}{@{}l@{}} \forall a, i.; \forall x, y, z \in \operatorname{vec}(\operatorname{quotient\text{-}module}(\operatorname{uba\text{-}as\text{-}module}(a), i)). \\ \quad (\operatorname{is\text{-}unital\text{-}banach\text{-}algebra}(a) \wedge \operatorname{is\text{-}two\text{-}sided\text{-}ideal}(a, i)) \Rightarrow  \\ \quad \quad \operatorname{qna\text{-}vmul}(a, i)(\operatorname{vadd}(\operatorname{quotient\text{-}module}(\operatorname{uba\text{-}as\text{-}module}(a), i))(x, y), z) \\ \quad \quad =\; \operatorname{vadd}(\operatorname{quotient\text{-}module}(\operatorname{uba\text{-}as\text{-}module}(a), i))(\operatorname{qna\text{-}vmul}(a, i)(x, z), \operatorname{qna\text{-}vmul}(a, i)(y, z)) \end{array}}}
\vnbprop{qna-vmul-act-left-na}{\vnbstmt{\begin{array}{@{}l@{}} \forall a, i.; \forall r \in \operatorname{carr}(\operatorname{scal}(\operatorname{quotient\text{-}module}(\operatorname{uba\text{-}as\text{-}module}(a), i))), x, y \in \operatorname{vec}(\operatorname{quotient\text{-}module}(\operatorname{uba\text{-}as\text{-}module}(a), i)). \\ \quad (\operatorname{is\text{-}unital\text{-}banach\text{-}algebra}(a) \wedge \operatorname{is\text{-}two\text{-}sided\text{-}ideal}(a, i)) \Rightarrow  \\ \quad \quad \operatorname{qna\text{-}vmul}(a, i)(\operatorname{act}(\operatorname{quotient\text{-}module}(\operatorname{uba\text{-}as\text{-}module}(a), i))(r, x), y) \\ \quad \quad =\; \operatorname{act}(\operatorname{quotient\text{-}module}(\operatorname{uba\text{-}as\text{-}module}(a), i))(r, \operatorname{qna\text{-}vmul}(a, i)(x, y)) \end{array}}}
\vnbprop{qna-vmul-act-right-na}{\vnbstmt{\begin{array}{@{}l@{}} \forall a, i.; \forall r \in \operatorname{carr}(\operatorname{scal}(\operatorname{quotient\text{-}module}(\operatorname{uba\text{-}as\text{-}module}(a), i))), x, y \in \operatorname{vec}(\operatorname{quotient\text{-}module}(\operatorname{uba\text{-}as\text{-}module}(a), i)). \\ \quad (\operatorname{is\text{-}unital\text{-}banach\text{-}algebra}(a) \wedge \operatorname{is\text{-}two\text{-}sided\text{-}ideal}(a, i)) \Rightarrow  \\ \quad \quad \operatorname{qna\text{-}vmul}(a, i)(x, \operatorname{act}(\operatorname{quotient\text{-}module}(\operatorname{uba\text{-}as\text{-}module}(a), i))(r, y)) \\ \quad \quad =\; \operatorname{act}(\operatorname{quotient\text{-}module}(\operatorname{uba\text{-}as\text{-}module}(a), i))(r, \operatorname{qna\text{-}vmul}(a, i)(x, y)) \end{array}}}
\vnbprop{qna-vnrm-definite}{\vnbstmt{\begin{array}{@{}l@{}} \forall a, i.; \forall x \in \operatorname{vec}(\operatorname{quotient\text{-}module}(\operatorname{uba\text{-}as\text{-}module}(a), i)). \\ \quad (\operatorname{is\text{-}unital\text{-}banach\text{-}algebra}(a) \wedge (\operatorname{is\text{-}two\text{-}sided\text{-}ideal}(a, i) \wedge \operatorname{is\text{-}closed}(\operatorname{ba\text{-}metric\text{-}space}(a), i))) \Rightarrow  \\ \quad \quad (\operatorname{qna\text{-}vnrm}(a, i)(x) = 0) \Leftrightarrow (x = \operatorname{vzero}(\operatorname{quotient\text{-}module}(\operatorname{uba\text{-}as\text{-}module}(a), i))) \end{array}}}
\vnbprop{qna-vmul-assoc-na}{\vnbstmt{\begin{array}{@{}l@{}} \forall a, i. \\ \quad (\operatorname{is\text{-}unital\text{-}banach\text{-}algebra}(a) \wedge \operatorname{is\text{-}two\text{-}sided\text{-}ideal}(a, i)) \Rightarrow  \\ \quad \quad \operatorname{is\text{-}associative}(\operatorname{qna\text{-}vmul}(a, i), \operatorname{vec}(\operatorname{quotient\text{-}module}(\operatorname{uba\text{-}as\text{-}module}(a), i))) \end{array}}}
\vnbprop{quotient-na-is-cna}{\vnbstmt{\begin{array}{@{}l@{}} \forall a, i. \\ \quad (\operatorname{is\text{-}unital\text{-}banach\text{-}algebra}(a) \wedge (\operatorname{is\text{-}two\text{-}sided\text{-}ideal}(a, i) \wedge \operatorname{is\text{-}closed}(\operatorname{ba\text{-}metric\text{-}space}(a), i))) \Rightarrow  \\ \quad \quad \operatorname{is\text{-}complex\text{-}normed\text{-}algebra}({a}/{i}) \end{array}}}
\vnbprop{bqn-sub-swap}{\vnbstmt{\begin{array}{@{}l@{}} \forall a.; \forall u, v, w \in \operatorname{vec}(a). \\ \quad \operatorname{is\text{-}unital\text{-}banach\text{-}algebra}(a) \Rightarrow  \\ \quad \quad \operatorname{vadd}(a)(\operatorname{vadd}(a)(u, \operatorname{vneg}(a)(v)), \operatorname{vneg}(a)(w)) \\ \quad \quad =\; \operatorname{vadd}(a)(\operatorname{vadd}(a)(u, \operatorname{vneg}(a)(w)), \operatorname{vneg}(a)(v)) \end{array}}}
\vnbprop{quotient-na-complete}{\vnbstmt{\begin{array}{@{}l@{}} \forall a, i. \\ \quad (\operatorname{is\text{-}unital\text{-}banach\text{-}algebra}(a) \wedge (\operatorname{is\text{-}two\text{-}sided\text{-}ideal}(a, i) \wedge \operatorname{is\text{-}closed}(\operatorname{ba\text{-}metric\text{-}space}(a), i))) \Rightarrow  \\ \quad \quad \operatorname{is\text{-}complete}(\operatorname{nag\text{-}metric\text{-}space}(\operatorname{cna\text{-}as\text{-}normed\text{-}ag}({a}/{i}))) \end{array}}}
\vnbprop{quotient-na-is-cba}{\vnbstmt{\begin{array}{@{}l@{}} \forall a, i. \\ \quad (\operatorname{is\text{-}unital\text{-}banach\text{-}algebra}(a) \wedge (\operatorname{is\text{-}two\text{-}sided\text{-}ideal}(a, i) \wedge \operatorname{is\text{-}closed}(\operatorname{ba\text{-}metric\text{-}space}(a), i))) \Rightarrow  \\ \quad \quad \operatorname{is\text{-}complex\text{-}banach\text{-}algebra}({a}/{i}) \end{array}}}
\vnbprop{coset-norm-class-quot-norm}{\vnbstmt{\begin{array}{@{}l@{}} \forall a, i.; \forall x \in \operatorname{vec}(a). \\ \quad (\operatorname{is\text{-}unital\text{-}banach\text{-}algebra}(a) \wedge \operatorname{is\text{-}two\text{-}sided\text{-}ideal}(a, i)) \Rightarrow  \\ \quad \quad \operatorname{coset\text{-}norm}(a, i, \operatorname{class}(\operatorname{qmod\text{-}setoid}(\operatorname{uba\text{-}as\text{-}module}(a), i), x)) = \operatorname{quot\text{-}norm}(a, i, x) \end{array}}}
\vnbprop{qba-vone-left}{\vnbstmt{\begin{array}{@{}l@{}} \forall a, i.; \forall c \in \operatorname{quotient}(\operatorname{qmod\text{-}setoid}(\operatorname{uba\text{-}as\text{-}module}(a), i)). \\ \quad (\operatorname{is\text{-}unital\text{-}banach\text{-}algebra}(a) \wedge \operatorname{is\text{-}two\text{-}sided\text{-}ideal}(a, i)) \Rightarrow  \\ \quad \quad \operatorname{qna\text{-}vmul}(a, i)(\operatorname{class}(\operatorname{qmod\text{-}setoid}(\operatorname{uba\text{-}as\text{-}module}(a), i), 1_{a}), c) = c \end{array}}}
\vnbprop{qba-vone-right}{\vnbstmt{\begin{array}{@{}l@{}} \forall a, i.; \forall c \in \operatorname{quotient}(\operatorname{qmod\text{-}setoid}(\operatorname{uba\text{-}as\text{-}module}(a), i)). \\ \quad (\operatorname{is\text{-}unital\text{-}banach\text{-}algebra}(a) \wedge \operatorname{is\text{-}two\text{-}sided\text{-}ideal}(a, i)) \Rightarrow  \\ \quad \quad \operatorname{qna\text{-}vmul}(a, i)(c, \operatorname{class}(\operatorname{qmod\text{-}setoid}(\operatorname{uba\text{-}as\text{-}module}(a), i), 1_{a})) = c \end{array}}}
\vnbprop{qba-vone-left-na}{\vnbstmt{\begin{array}{@{}l@{}} \forall a, i.; \forall x \in \operatorname{vec}(\operatorname{quotient\text{-}module}(\operatorname{uba\text{-}as\text{-}module}(a), i)). \\ \quad (\operatorname{is\text{-}unital\text{-}banach\text{-}algebra}(a) \wedge \operatorname{is\text{-}two\text{-}sided\text{-}ideal}(a, i)) \Rightarrow  \\ \quad \quad \operatorname{qna\text{-}vmul}(a, i)(\operatorname{class}(\operatorname{qmod\text{-}setoid}(\operatorname{uba\text{-}as\text{-}module}(a), i), 1_{a}), x) = x \end{array}}}
\vnbprop{qba-vone-right-na}{\vnbstmt{\begin{array}{@{}l@{}} \forall a, i.; \forall x \in \operatorname{vec}(\operatorname{quotient\text{-}module}(\operatorname{uba\text{-}as\text{-}module}(a), i)). \\ \quad (\operatorname{is\text{-}unital\text{-}banach\text{-}algebra}(a) \wedge \operatorname{is\text{-}two\text{-}sided\text{-}ideal}(a, i)) \Rightarrow  \\ \quad \quad \operatorname{qna\text{-}vmul}(a, i)(x, \operatorname{class}(\operatorname{qmod\text{-}setoid}(\operatorname{uba\text{-}as\text{-}module}(a), i), 1_{a})) = x \end{array}}}
\vnbprop{qba-vone-norm}{\vnbstmt{\begin{array}{@{}l@{}} \forall a, i. \\ \quad (\operatorname{is\text{-}unital\text{-}banach\text{-}algebra}(a) \wedge (\operatorname{is\text{-}two\text{-}sided\text{-}ideal}(a, i) \wedge \neg (1_{a} \in i))) \Rightarrow  \\ \quad \quad \operatorname{qna\text{-}vnrm}(a, i)(\operatorname{class}(\operatorname{qmod\text{-}setoid}(\operatorname{uba\text{-}as\text{-}module}(a), i), 1_{a})) = 1 \end{array}}}
\vnbprop{quotient-na-unfold}{\vnbstmt{\begin{array}{@{}l@{}} \forall a, i. \\ \quad (\operatorname{is\text{-}unital\text{-}banach\text{-}algebra}(a) \wedge \operatorname{is\text{-}two\text{-}sided\text{-}ideal}(a, i)) \Rightarrow  \\ \quad \quad ({a}/{i} \simeq {[}\operatorname{scal}(\operatorname{quotient\text{-}module}(\operatorname{uba\text{-}as\text{-}module}(a), i)), \operatorname{vec}(\operatorname{quotient\text{-}module}(\operatorname{uba\text{-}as\text{-}module}(a), i)), \operatorname{vadd}(\operatorname{quotient\text{-}module}(\operatorname{uba\text{-}as\text{-}module}(a), i)), \operatorname{vzero}(\operatorname{quotient\text{-}module}(\operatorname{uba\text{-}as\text{-}module}(a), i)), \operatorname{vneg}(\operatorname{quotient\text{-}module}(\operatorname{uba\text{-}as\text{-}module}(a), i)), \operatorname{act}(\operatorname{quotient\text{-}module}(\operatorname{uba\text{-}as\text{-}module}(a), i)), \operatorname{qna\text{-}vnrm}(a, i), \operatorname{qna\text{-}vmul}(a, i)]) \end{array}}}
\vnbprop{quotient-ba-is-uba}{\vnbstmt{\begin{array}{@{}l@{}} \forall a, i. \\ \quad \operatorname{is\text{-}unital\text{-}banach\text{-}algebra}(a) \wedge \\ \quad \operatorname{is\text{-}two\text{-}sided\text{-}ideal}(a, i) \wedge \\ \quad \operatorname{is\text{-}closed}(\operatorname{ba\text{-}metric\text{-}space}(a), i) \wedge \\ \quad \neg (1_{a} \in i) \Rightarrow \\ \quad \quad \operatorname{is\text{-}unital\text{-}banach\text{-}algebra}({a}/{i}) \end{array}}}
\vnbprop{quotient-ba-is-uba-proper}{\vnbstmt{\begin{array}{@{}l@{}} \forall a, i. \\ \quad \operatorname{is\text{-}unital\text{-}banach\text{-}algebra}(a) \wedge \\ \quad \operatorname{is\text{-}two\text{-}sided\text{-}ideal}(a, i) \wedge \\ \quad \operatorname{is\text{-}closed}(\operatorname{ba\text{-}metric\text{-}space}(a), i) \wedge \\ \quad \neg (i = \operatorname{vec}(a)) \Rightarrow \\ \quad \quad \operatorname{is\text{-}unital\text{-}banach\text{-}algebra}({a}/{i}) \end{array}}}

\vnbbookfile{theorem-library/ba-spectral-radius}{theorem-library/ba-spectral-radius}
\noindent{\small ba-spectral-radius.scm -- wave of 2026-11-09/10 (agent BA-1F), wired 2026-11-10 from scratchpad/ba1f-1109/ba-spectral-radius.scm. Library: operator-algebras.}\par
\vnbprop{spectral-radius-unfold}{\vnbstmt{\begin{array}{@{}l@{}} \forall a, x. \\ \quad (\rho({x}) \simeq \operatorname{einf}(\{w \in \mathbb{R} : \exists n \in \mathbb{N}.\; ((1 \leq n) \wedge (w = {\operatorname{vnrm}(a)({x}^{n})}^{{n}^{\ast\!-1}}))\})) \end{array}}}
\vnbprop{spectral-radius-real}{\vnbstmt{\begin{array}{@{}l@{}} \forall a.; \forall x \in \operatorname{vec}(a). \\ \quad \operatorname{is\text{-}unital\text{-}banach\text{-}algebra}(a) \Rightarrow  \\ \quad \quad (\rho({x}) \in \mathbb{R}) \wedge (0 \leq \rho({x})) \end{array}}}
\vnbprop{bsr-el-pow-act}{\vnbstmt{\begin{array}{@{}l@{}} \forall n \in \mathbb{N}, a.; \forall c \in \mathbb{C}, x \in \operatorname{vec}(a). \\ \quad \operatorname{is\text{-}unital\text{-}banach\text{-}algebra}(a) \Rightarrow  \\ \quad \quad {\operatorname{act}(a)(c, x)}^{n} = \operatorname{act}(a)(\operatorname{power}(c, n), {x}^{n}) \end{array}}}
\vnbprop{bsr-one-minus-pow-invertible}{\vnbstmt{\begin{array}{@{}l@{}} \forall n \in \mathbb{N}, a.; \forall x \in \operatorname{vec}(a). \\ \quad (\operatorname{is\text{-}unital\text{-}banach\text{-}algebra}(a) \wedge \operatorname{is\text{-}invertible\text{-}el}(a, \operatorname{vadd}(a)(1_{a}, \operatorname{vneg}(a)({x}^{n})))) \Rightarrow  \\ \quad \quad \operatorname{is\text{-}invertible\text{-}el}(a, \operatorname{vadd}(a)(1_{a}, \operatorname{vneg}(a)(x))) \end{array}}}
\vnbprop{shift-invertible-beyond-spectral-radius}{\vnbstmt{\begin{array}{@{}l@{}} \forall a.; \forall x \in \operatorname{vec}(a), z \in \mathbb{C}. \\ \quad (\operatorname{is\text{-}unital\text{-}banach\text{-}algebra}(a) \wedge (\rho({x}) < \operatorname{magnitude}(z))) \Rightarrow  \\ \quad \quad \operatorname{is\text{-}invertible\text{-}el}(a, \operatorname{shift\text{-}el}(a, z, x)) \end{array}}}
\vnbprop{spectrum-in-spectral-radius-disc}{\vnbstmt{\begin{array}{@{}l@{}} \forall a.; \forall x \in \operatorname{vec}(a), z \in \sigma({x}). \\ \quad \operatorname{is\text{-}unital\text{-}banach\text{-}algebra}(a) \Rightarrow (\operatorname{magnitude}(z) \leq \rho({x})) \end{array}}}
\vnbprop{resolvent-set-open}{\vnbstmt{\begin{array}{@{}l@{}} \forall a.; \forall x \in \operatorname{vec}(a). \\ \quad \operatorname{is\text{-}unital\text{-}banach\text{-}algebra}(a) \Rightarrow  \\ \quad \quad \operatorname{is\text{-}open}(\operatorname{nf\text{-}metric\text{-}space}(\operatorname{cc\text{-}normed\text{-}field}), \operatorname{difference}(\mathbb{C}, \sigma({x}))) \end{array}}}

\vnbbookfile{theorem-library/ba-gelfand-compact}{theorem-library/ba-gelfand-compact}
\noindent{\small ba-gelfand-compact.scm -- wave of 2026-11-09/10 (agent BA-1F), wired 2026-11-10 from scratchpad/ba1f-1109/ba-gelfand-compact.scm. Library: operator-algebras.}\par
\vnbprop{cc-product-fun}{\vnbstmt{\begin{array}{@{}l@{}} \forall a, f. \\ \quad ((a \in \mathbf{Set}) \wedge (f \in \operatorname{tprod\text{-}pts}(a, (\lambda x \in a.\; \operatorname{metric\text{-}top}(\operatorname{cc\text{-}ms}))))) \Rightarrow  \\ \quad \quad f \in (a \to \mathbb{C}) \end{array}}}
\vnbprop{characters-closed-in-product}{\vnbstmt{\begin{array}{@{}l@{}} \forall a. \\ \quad \operatorname{is\text{-}unital\text{-}banach\text{-}algebra}(a) \Rightarrow  \\ \quad \quad \operatorname{is\text{-}top\text{-}closed}(\operatorname{tproduct\text{-}top}(\operatorname{vec}(a), (\lambda x \in \operatorname{vec}(a).\; \operatorname{metric\text{-}top}(\operatorname{cc\text{-}ms}))), \operatorname{characters}(a)) \end{array}}}
\vnbprop{gelfand-space-compact}{\vnbstmt{\forall a.\; (\operatorname{is\text{-}unital\text{-}banach\text{-}algebra}(a) \Rightarrow \operatorname{is\text{-}compact\text{-}t}(\widehat{a}))}}

\vnbbookfile{theorem-library/ba-cstar-laws}{theorem-library/ba-cstar-laws}
\noindent{\small ba-cstar-laws.scm -- wave of 2026-11-09/10 (agent BA-1H), wired 2026-11-10 from scratchpad/ba1h-1109/ba-cstar-laws.scm. Library: operator-algebras.}\par
\vnbprop{bcs-ccsa-ucsa}{\vnbstmt{\begin{array}{@{}l@{}} \forall a. \\ \quad \operatorname{is\text{-}commutative\text{-}c\text{-}star\text{-}algebra}(a) \Rightarrow \operatorname{is\text{-}unital\text{-}c\text{-}star\text{-}algebra}(a) \end{array}}}
\vnbprop{bcs-ccsa-comm}{\vnbstmt{\begin{array}{@{}l@{}} \forall a.; \forall x, y \in \operatorname{vec}(a). \\ \quad \operatorname{is\text{-}commutative\text{-}c\text{-}star\text{-}algebra}(a) \Rightarrow (\cdot_{a}(x, y) = \cdot_{a}(y, x)) \end{array}}}
\vnbprop{bcs-ucsa-uiba}{\vnbstmt{\begin{array}{@{}l@{}} \forall a. \\ \quad \operatorname{is\text{-}unital\text{-}c\text{-}star\text{-}algebra}(a) \Rightarrow \operatorname{is\text{-}unital\text{-}involutive\text{-}banach\text{-}algebra}(a) \end{array}}}
\vnbprop{bcs-ucsa-cstar}{\vnbstmt{\begin{array}{@{}l@{}} \forall a.; \forall x \in \operatorname{vec}(a). \\ \quad \operatorname{is\text{-}unital\text{-}c\text{-}star\text{-}algebra}(a) \Rightarrow  \\ \quad \quad \operatorname{vnrm}(a)(\cdot_{a}({}^{*_{a}}(x), x)) = (\operatorname{vnrm}(a)(x) \cdot \operatorname{vnrm}(a)(x)) \end{array}}}
\vnbprop{bcs-uiba-uba-tuple}{\vnbstmt{\begin{array}{@{}l@{}} \forall a. \\ \quad \operatorname{is\text{-}unital\text{-}involutive\text{-}banach\text{-}algebra}(a) \Rightarrow  \\ \quad \quad \operatorname{is\text{-}unital\text{-}banach\text{-}algebra}({[}\operatorname{scal}(a), \operatorname{vec}(a), \operatorname{vadd}(a), \operatorname{vzero}(a), \operatorname{vneg}(a), \operatorname{act}(a), \operatorname{vnrm}(a), \cdot_{a}, 1_{a}]) \end{array}}}
\vnbprop{bcs-uiba-scal-carr}{\vnbstmt{\forall a.\; (\operatorname{is\text{-}unital\text{-}involutive\text{-}banach\text{-}algebra}(a) \Rightarrow (\operatorname{carr}(\operatorname{scal}(a)) = \mathbb{C}))}}
\vnbprop{bcs-uiba-vstar-fun}{\vnbstmt{\begin{array}{@{}l@{}} \forall a. \\ \quad \operatorname{is\text{-}unital\text{-}involutive\text{-}banach\text{-}algebra}(a) \Rightarrow ({}^{*_{a}} \in (\operatorname{vec}(a) \to \operatorname{vec}(a))) \end{array}}}
\vnbprop{bcs-uiba-star-add}{\vnbstmt{\begin{array}{@{}l@{}} \forall a.; \forall x, y \in \operatorname{vec}(a). \\ \quad \operatorname{is\text{-}unital\text{-}involutive\text{-}banach\text{-}algebra}(a) \Rightarrow  \\ \quad \quad {}^{*_{a}}(\operatorname{vadd}(a)(x, y)) = \operatorname{vadd}(a)({}^{*_{a}}(x), {}^{*_{a}}(y)) \end{array}}}
\vnbprop{bcs-uiba-star-act}{\vnbstmt{\begin{array}{@{}l@{}} \forall a.; \forall r \in \mathbb{C}, x \in \operatorname{vec}(a). \\ \quad \operatorname{is\text{-}unital\text{-}involutive\text{-}banach\text{-}algebra}(a) \Rightarrow  \\ \quad \quad {}^{*_{a}}(\operatorname{act}(a)(r, x)) = \operatorname{act}(a)(\operatorname{conjugate}(r), {}^{*_{a}}(x)) \end{array}}}
\vnbprop{bcs-uiba-star-invol}{\vnbstmt{\begin{array}{@{}l@{}} \forall a.; \forall x \in \operatorname{vec}(a). \\ \quad \operatorname{is\text{-}unital\text{-}involutive\text{-}banach\text{-}algebra}(a) \Rightarrow ({}^{*_{a}}({}^{*_{a}}(x)) = x) \end{array}}}
\vnbprop{bcs-uiba-star-mul}{\vnbstmt{\begin{array}{@{}l@{}} \forall a.; \forall x, y \in \operatorname{vec}(a). \\ \quad \operatorname{is\text{-}unital\text{-}involutive\text{-}banach\text{-}algebra}(a) \Rightarrow  \\ \quad \quad {}^{*_{a}}(\cdot_{a}(x, y)) = \cdot_{a}({}^{*_{a}}(y), {}^{*_{a}}(x)) \end{array}}}
\vnbprop{bcs-uiba-star-norm}{\vnbstmt{\begin{array}{@{}l@{}} \forall a.; \forall x \in \operatorname{vec}(a). \\ \quad \operatorname{is\text{-}unital\text{-}involutive\text{-}banach\text{-}algebra}(a) \Rightarrow  \\ \quad \quad \operatorname{vnrm}(a)({}^{*_{a}}(x)) = \operatorname{vnrm}(a)(x) \end{array}}}
\vnbprop{bcs-uiba-view-unfold}{\vnbstmt{\begin{array}{@{}l@{}} \forall a. \\ \quad (\operatorname{uiba\text{-}as\text{-}uba}(a) \simeq {[}\operatorname{scal}(a), \operatorname{vec}(a), \operatorname{vadd}(a), \operatorname{vzero}(a), \operatorname{vneg}(a), \operatorname{act}(a), \operatorname{vnrm}(a), \cdot_{a}, 1_{a}]) \end{array}}}
\vnbprop{bcs-uiba-is-uba}{\vnbstmt{\begin{array}{@{}l@{}} \forall a. \\ \quad \operatorname{is\text{-}unital\text{-}involutive\text{-}banach\text{-}algebra}(a) \Rightarrow  \\ \quad \quad \operatorname{is\text{-}unital\text{-}banach\text{-}algebra}(\operatorname{uiba\text{-}as\text{-}uba}(a)) \end{array}}}
\vnbprop{bcs-uu-scal}{\vnbstmt{\begin{array}{@{}l@{}} \forall a. \\ \quad \operatorname{is\text{-}unital\text{-}involutive\text{-}banach\text{-}algebra}(a) \Rightarrow (\operatorname{scal}(\operatorname{uiba\text{-}as\text{-}uba}(a)) \simeq \operatorname{scal}(a)) \end{array}}}
\vnbprop{bcs-uu-vec}{\vnbstmt{\begin{array}{@{}l@{}} \forall a. \\ \quad \operatorname{is\text{-}unital\text{-}involutive\text{-}banach\text{-}algebra}(a) \Rightarrow (\operatorname{vec}(\operatorname{uiba\text{-}as\text{-}uba}(a)) \simeq \operatorname{vec}(a)) \end{array}}}
\vnbprop{bcs-uu-vadd}{\vnbstmt{\begin{array}{@{}l@{}} \forall a. \\ \quad \operatorname{is\text{-}unital\text{-}involutive\text{-}banach\text{-}algebra}(a) \Rightarrow (\operatorname{vadd}(\operatorname{uiba\text{-}as\text{-}uba}(a)) \simeq \operatorname{vadd}(a)) \end{array}}}
\vnbprop{bcs-uu-vzero}{\vnbstmt{\begin{array}{@{}l@{}} \forall a. \\ \quad \operatorname{is\text{-}unital\text{-}involutive\text{-}banach\text{-}algebra}(a) \Rightarrow (\operatorname{vzero}(\operatorname{uiba\text{-}as\text{-}uba}(a)) \simeq \operatorname{vzero}(a)) \end{array}}}
\vnbprop{bcs-uu-vneg}{\vnbstmt{\begin{array}{@{}l@{}} \forall a. \\ \quad \operatorname{is\text{-}unital\text{-}involutive\text{-}banach\text{-}algebra}(a) \Rightarrow (\operatorname{vneg}(\operatorname{uiba\text{-}as\text{-}uba}(a)) \simeq \operatorname{vneg}(a)) \end{array}}}
\vnbprop{bcs-uu-act}{\vnbstmt{\begin{array}{@{}l@{}} \forall a. \\ \quad \operatorname{is\text{-}unital\text{-}involutive\text{-}banach\text{-}algebra}(a) \Rightarrow (\operatorname{act}(\operatorname{uiba\text{-}as\text{-}uba}(a)) \simeq \operatorname{act}(a)) \end{array}}}
\vnbprop{bcs-uu-vnrm}{\vnbstmt{\begin{array}{@{}l@{}} \forall a. \\ \quad \operatorname{is\text{-}unital\text{-}involutive\text{-}banach\text{-}algebra}(a) \Rightarrow (\operatorname{vnrm}(\operatorname{uiba\text{-}as\text{-}uba}(a)) \simeq \operatorname{vnrm}(a)) \end{array}}}
\vnbprop{bcs-uu-vmul}{\vnbstmt{\begin{array}{@{}l@{}} \forall a. \\ \quad \operatorname{is\text{-}unital\text{-}involutive\text{-}banach\text{-}algebra}(a) \Rightarrow (\cdot_{\operatorname{uiba\text{-}as\text{-}uba}(a)} \simeq \cdot_{a}) \end{array}}}
\vnbprop{bcs-uu-vone}{\vnbstmt{\begin{array}{@{}l@{}} \forall a. \\ \quad \operatorname{is\text{-}unital\text{-}involutive\text{-}banach\text{-}algebra}(a) \Rightarrow (1_{\operatorname{uiba\text{-}as\text{-}uba}(a)} \simeq 1_{a}) \end{array}}}
\vnbprop{bcs-star-one}{\vnbstmt{\begin{array}{@{}l@{}} \forall a. \\ \quad \operatorname{is\text{-}unital\text{-}involutive\text{-}banach\text{-}algebra}(a) \Rightarrow ({}^{*_{a}}(1_{a}) = 1_{a}) \end{array}}}
\vnbprop{bcs-characters-view}{\vnbstmt{\begin{array}{@{}l@{}} \forall a. \\ \quad \operatorname{is\text{-}unital\text{-}involutive\text{-}banach\text{-}algebra}(a) \Rightarrow  \\ \quad \quad (\operatorname{characters}(\operatorname{uiba\text{-}as\text{-}uba}(a)) \simeq \operatorname{characters}(a)) \end{array}}}
\vnbprop{bcs-character-bound}{\vnbstmt{\begin{array}{@{}l@{}} \forall a.; \forall c \in \operatorname{characters}(a), x \in \operatorname{vec}(a). \\ \quad \operatorname{is\text{-}unital\text{-}involutive\text{-}banach\text{-}algebra}(a) \Rightarrow (\operatorname{magnitude}(c(x)) \leq \operatorname{vnrm}(a)(x)) \end{array}}}
\vnbprop{bcs-linear-bound-zero}{\vnbstmt{\begin{array}{@{}l@{}} \forall l, e, c \in \mathbb{R}. \\ \quad \forall t \in \mathbb{R}.\; (((l \cdot l) + ((e + t) \cdot (e + t))) \leq (c + (t \cdot t))) \Rightarrow (e = 0) \end{array}}}

\vnbbookfile{theorem-library/ba-spectrum-nonempty}{theorem-library/ba-spectrum-nonempty}
\noindent{\small ba-spectrum-nonempty.scm -- wave of 2026-11-09/10 (agent BA-1D), wired 2026-11-10 from scratchpad/ba1d-1109/ba-spectrum-nonempty.scm. Library: operator-algebras.}\par
\vnbprop{bsn-uam-vec}{\vnbstmt{\forall a.\; (\operatorname{is\text{-}unital\text{-}banach\text{-}algebra}(a) \Rightarrow (\operatorname{vec}(\operatorname{uba\text{-}as\text{-}module}(a)) \simeq \operatorname{vec}(a)))}}
\vnbprop{bsn-uam-vadd}{\vnbstmt{\begin{array}{@{}l@{}} \forall a. \\ \quad \operatorname{is\text{-}unital\text{-}banach\text{-}algebra}(a) \Rightarrow (\operatorname{vadd}(\operatorname{uba\text{-}as\text{-}module}(a)) \simeq \operatorname{vadd}(a)) \end{array}}}
\vnbprop{bsn-uam-act}{\vnbstmt{\forall a.\; (\operatorname{is\text{-}unital\text{-}banach\text{-}algebra}(a) \Rightarrow (\operatorname{act}(\operatorname{uba\text{-}as\text{-}module}(a)) \simeq \operatorname{act}(a)))}}
\vnbprop{bsn-uam-scal}{\vnbstmt{\begin{array}{@{}l@{}} \forall a. \\ \quad \operatorname{is\text{-}unital\text{-}banach\text{-}algebra}(a) \Rightarrow (\operatorname{scal}(\operatorname{uba\text{-}as\text{-}module}(a)) \simeq \operatorname{scal}(a)) \end{array}}}
\vnbprop{uba-norm-seminorm}{\vnbstmt{\begin{array}{@{}l@{}} \forall a. \\ \quad \operatorname{is\text{-}unital\text{-}banach\text{-}algebra}(a) \Rightarrow  \\ \quad \quad \operatorname{is\text{-}seminorm}(\operatorname{cc\text{-}normed\text{-}field}, \operatorname{uba\text{-}as\text{-}module}(a), \operatorname{vnrm}(a)) \end{array}}}
\vnbprop{bsn-cl-fun}{\vnbstmt{\begin{array}{@{}l@{}} \forall a, f. \\ \quad (\operatorname{is\text{-}unital\text{-}banach\text{-}algebra}(a) \wedge \operatorname{is\text{-}cc\text{-}linear\text{-}on}(a, f)) \Rightarrow  \\ \quad \quad f \in (\operatorname{vec}(a) \to \mathbb{C}) \end{array}}}
\vnbprop{bsn-cl-add}{\vnbstmt{\begin{array}{@{}l@{}} \forall a, f.; \forall x, y \in \operatorname{vec}(a). \\ \quad (\operatorname{is\text{-}unital\text{-}banach\text{-}algebra}(a) \wedge \operatorname{is\text{-}cc\text{-}linear\text{-}on}(a, f)) \Rightarrow  \\ \quad \quad f(\operatorname{vadd}(a)(x, y)) = (f(x) + f(y)) \end{array}}}
\vnbprop{bsn-cl-act}{\vnbstmt{\begin{array}{@{}l@{}} \forall a, f.; \forall z \in \mathbb{C}, x \in \operatorname{vec}(a). \\ \quad (\operatorname{is\text{-}unital\text{-}banach\text{-}algebra}(a) \wedge \operatorname{is\text{-}cc\text{-}linear\text{-}on}(a, f)) \Rightarrow  \\ \quad \quad f(\operatorname{act}(a)(z, x)) = (z \cdot f(x)) \end{array}}}
\vnbprop{bsn-psi-neg}{\vnbstmt{\begin{array}{@{}l@{}} \forall a, f.; \forall w \in \operatorname{vec}(a). \\ \quad (\operatorname{is\text{-}unital\text{-}banach\text{-}algebra}(a) \wedge \operatorname{is\text{-}cc\text{-}linear\text{-}on}(a, f)) \Rightarrow  \\ \quad \quad f(\operatorname{vneg}(a)(w)) = -f(w) \end{array}}}
\vnbprop{bsn-klin}{\vnbstmt{\begin{array}{@{}l@{}} \forall a, f. \\ \quad (\operatorname{is\text{-}unital\text{-}banach\text{-}algebra}(a) \wedge \operatorname{is\text{-}k\text{-}linear}(\operatorname{uba\text{-}as\text{-}module}(a), f)) \Rightarrow  \\ \quad \quad \operatorname{is\text{-}cc\text{-}linear\text{-}on}(a, f) \end{array}}}
\vnbprop{uba-norming-functional}{\vnbstmt{\begin{array}{@{}l@{}} \forall a.; \forall y \in \operatorname{vec}(a). \\ \quad \operatorname{is\text{-}unital\text{-}banach\text{-}algebra}(a) \Rightarrow  \\ \quad \quad \exists f. \\ \quad \quad \quad \operatorname{is\text{-}cc\text{-}linear\text{-}on}(a, f) \wedge \\ \quad \quad \quad \forall v \in \operatorname{vec}(a).\; (\operatorname{magnitude}(f(v)) \leq \operatorname{vnrm}(a)(v)) \wedge \\ \quad \quad \quad f(y) = \operatorname{vnrm}(a)(y) \end{array}}}
\vnbprop{resolvent-functional-holomorphic}{\vnbstmt{\begin{array}{@{}l@{}} \forall a.; \forall x \in \operatorname{vec}(a), f, u. \\ \quad \operatorname{is\text{-}unital\text{-}banach\text{-}algebra}(a) \wedge \\ \quad \operatorname{is\text{-}cc\text{-}linear\text{-}on}(a, f) \wedge \\ \quad \forall v \in \operatorname{vec}(a).\; (\operatorname{magnitude}(f(v)) \leq \operatorname{vnrm}(a)(v)) \wedge \\ \quad \operatorname{is\text{-}open}(\operatorname{nf\text{-}metric\text{-}space}(\operatorname{cc\text{-}normed\text{-}field}), u) \wedge \\ \quad \forall z \in u.\; \operatorname{is\text{-}invertible\text{-}el}(a, \operatorname{shift\text{-}el}(a, z, x)) \Rightarrow \\ \quad \quad \operatorname{holomorphic\text{-}on}(u, (\lambda w \in u.\; f({\operatorname{shift\text{-}el}(a, w, x)}^{-1}))) \end{array}}}
\vnbprop{bsn-no-spectrum}{\vnbstmt{\begin{array}{@{}l@{}} \forall a.; \forall x \in \operatorname{vec}(a), z \in \mathbb{C}. \\ \quad (\operatorname{is\text{-}unital\text{-}banach\text{-}algebra}(a) \wedge \neg \exists q.\; (q \in \sigma({x}))) \Rightarrow  \\ \quad \quad \operatorname{is\text{-}invertible\text{-}el}(a, \operatorname{shift\text{-}el}(a, z, x)) \end{array}}}
\vnbprop{cc-plane-open}{\vnbstmt{\operatorname{is\text{-}open}(\operatorname{nf\text{-}metric\text{-}space}(\operatorname{cc\text{-}normed\text{-}field}), \mathbb{C})}}
\vnbprop{spectrum-nonempty}{\vnbstmt{\begin{array}{@{}l@{}} \forall a.; \forall x \in \operatorname{vec}(a). \\ \quad \operatorname{is\text{-}unital\text{-}banach\text{-}algebra}(a) \Rightarrow \exists q.\; (q \in \sigma({x})) \end{array}}}
\vnbprop{gelfand-mazur}{\vnbstmt{\begin{array}{@{}l@{}} \forall a.; \forall x \in \operatorname{vec}(a). \\ \quad \operatorname{is\text{-}unital\text{-}banach\text{-}algebra}(a) \wedge \\ \quad \forall y \in \operatorname{vec}(a).\; (\neg (y = \operatorname{vzero}(a)) \Rightarrow \operatorname{is\text{-}invertible\text{-}el}(a, y)) \Rightarrow \\ \quad \quad \exists z \in \mathbb{C}.\; (x = \operatorname{act}(a)(z, 1_{a})) \end{array}}}

\vnbbookfile{theorem-library/th3-compact}{theorem-library/th3-compact}
\noindent{\small th3-compact.scm -- wave of 2026-11-09/10 (agent TH3-A), wired 2026-11-10 from scratchpad/th3-1109/th3-compact.scm. Library: operator-algebras.}\par
\vnbprop{strictly-mono-compose}{\vnbstmt{\begin{array}{@{}l@{}} \forall i, a. \\ \quad (\operatorname{strictly\text{-}mono\text{-}nn}(i) \wedge \operatorname{strictly\text{-}mono\text{-}nn}(a)) \Rightarrow  \\ \quad \quad \operatorname{strictly\text{-}mono\text{-}nn}((\lambda k \in \mathbb{N}.\; i(a(k)))) \end{array}}}
\vnbprop{converges-reindex-ptwise}{\vnbstmt{\begin{array}{@{}l@{}} \forall s, a, b, i, l. \\ \quad \operatorname{converges\text{-}to}(s, a, l) \wedge \\ \quad \operatorname{strictly\text{-}mono\text{-}nn}(i) \wedge \\ \quad b \in (\mathbb{N} \to \operatorname{pts}(s)) \wedge \\ \quad \forall n \in \mathbb{N}.\; (b(n) = a(i(n))) \Rightarrow \\ \quad \quad \operatorname{converges\text{-}to}(s, b, l) \end{array}}}
\vnbprop{converges-ptwise-eq}{\vnbstmt{\begin{array}{@{}l@{}} \forall s, a, b, l. \\ \quad (\operatorname{converges\text{-}to}(s, a, l) \wedge ((b \in (\mathbb{N} \to \operatorname{pts}(s))) \wedge \forall n \in \mathbb{N}.\; (b(n) = a(n)))) \Rightarrow  \\ \quad \quad \operatorname{converges\text{-}to}(s, b, l) \end{array}}}
\vnbprop{ipc-op}{\vnbstmt{\begin{array}{@{}l@{}} \forall v, w, t, a, b, l. \\ \quad \operatorname{is\text{-}complex\text{-}inner\text{-}product\text{-}space}(v) \wedge \\ \quad \operatorname{is\text{-}complex\text{-}inner\text{-}product\text{-}space}(w) \wedge \\ \quad t \in L({v}, {w}) \wedge \\ \quad \operatorname{converges\text{-}to}(\operatorname{ip\text{-}metric\text{-}space}(v), a, l) \wedge \\ \quad b \in (\mathbb{N} \to \operatorname{vec}(w)) \wedge \\ \quad \forall n \in \mathbb{N}.\; (b(n) = t(a(n))) \Rightarrow \\ \quad \quad \operatorname{converges\text{-}to}(\operatorname{ip\text{-}metric\text{-}space}(w), b, t(l)) \end{array}}}
\vnbprop{ipc-sum}{\vnbstmt{\begin{array}{@{}l@{}} \forall w, a, b, c, \mathit{la}, \mathit{lb}. \\ \quad \operatorname{is\text{-}complex\text{-}inner\text{-}product\text{-}space}(w) \wedge \\ \quad \operatorname{converges\text{-}to}(\operatorname{ip\text{-}metric\text{-}space}(w), a, \mathit{la}) \wedge \\ \quad \operatorname{converges\text{-}to}(\operatorname{ip\text{-}metric\text{-}space}(w), b, \mathit{lb}) \wedge \\ \quad c \in (\mathbb{N} \to \operatorname{vec}(w)) \wedge \\ \quad \forall n \in \mathbb{N}.\; (c(n) = \operatorname{vadd}(w)(a(n), b(n))) \Rightarrow \\ \quad \quad \operatorname{converges\text{-}to}(\operatorname{ip\text{-}metric\text{-}space}(w), c, \operatorname{vadd}(w)(\mathit{la}, \mathit{lb})) \end{array}}}
\vnbprop{ipc-act}{\vnbstmt{\begin{array}{@{}l@{}} \forall w, c, a, b, \mathit{la}. \\ \quad \operatorname{is\text{-}complex\text{-}inner\text{-}product\text{-}space}(w) \wedge \\ \quad c \in \mathbb{C} \wedge \\ \quad \operatorname{converges\text{-}to}(\operatorname{ip\text{-}metric\text{-}space}(w), a, \mathit{la}) \wedge \\ \quad b \in (\mathbb{N} \to \operatorname{vec}(w)) \wedge \\ \quad \forall n \in \mathbb{N}.\; (b(n) = \operatorname{act}(w)(c, a(n))) \Rightarrow \\ \quad \quad \operatorname{converges\text{-}to}(\operatorname{ip\text{-}metric\text{-}space}(w), b, \operatorname{act}(w)(c, \mathit{la})) \end{array}}}
\vnbprop{compact-op-seq}{\vnbstmt{\begin{array}{@{}l@{}} \forall v, w, t, f. \\ \quad \operatorname{is\text{-}complex\text{-}inner\text{-}product\text{-}space}(w) \wedge \\ \quad \operatorname{is\text{-}compact\text{-}op}(v, w, t) \wedge \\ \quad f \in (\mathbb{N} \to \operatorname{vec}(v)) \wedge \\ \quad \forall q \in \mathbb{N}.\; (\lVert f(q) \rVert \leq 1) \Rightarrow \\ \quad \quad \exists \mathit{phq}. \\ \quad \quad \quad \operatorname{strictly\text{-}mono\text{-}nn}(\mathit{phq}) \wedge \\ \quad \quad \quad \exists \mathit{lq}. \forall \mathit{bq} \in (\mathbb{N} \to \operatorname{vec}(w)). \\ \quad \quad \quad \quad \forall q \in \mathbb{N}.\; (\operatorname{bq}(q) = t(f(\operatorname{phq}(q)))) \Rightarrow  \\ \quad \quad \quad \quad \quad \operatorname{converges\text{-}to}(\operatorname{ip\text{-}metric\text{-}space}(w), \mathit{bq}, \mathit{lq}) \end{array}}}
\vnbprop{compact-op-intro}{\vnbstmt{\begin{array}{@{}l@{}} \forall v, w, t. \\ \quad \operatorname{is\text{-}complex\text{-}inner\text{-}product\text{-}space}(w) \wedge \\ \quad t \in L({v}, {w}) \wedge \\ \quad \forall f \in (\mathbb{N} \to \operatorname{vec}(v)). \\ \quad \quad \forall q \in \mathbb{N}.\; (\lVert f(q) \rVert \leq 1) \Rightarrow  \\ \quad \quad \quad \exists \mathit{phq}. \\ \quad \quad \quad \quad \operatorname{strictly\text{-}mono\text{-}nn}(\mathit{phq}) \wedge \\ \quad \quad \quad \quad \exists \mathit{lq}. \forall \mathit{bq} \in (\mathbb{N} \to \operatorname{vec}(w)). \\ \quad \quad \quad \quad \quad \forall q \in \mathbb{N}.\; (\operatorname{bq}(q) = t(f(\operatorname{phq}(q)))) \Rightarrow  \\ \quad \quad \quad \quad \quad \quad \operatorname{converges\text{-}to}(\operatorname{ip\text{-}metric\text{-}space}(w), \mathit{bq}, \mathit{lq}) \Rightarrow \\ \quad \quad \operatorname{is\text{-}compact\text{-}op}(v, w, t) \end{array}}}
\vnbprop{compact-op-bounded}{\vnbstmt{\forall v, w, t.\; (\operatorname{is\text{-}compact\text{-}op}(v, w, t) \Rightarrow (t \in L({v}, {w})))}}
\vnbprop{compact-op-sum}{\vnbstmt{\begin{array}{@{}l@{}} \forall v, w, s, t. \\ \quad \operatorname{is\text{-}complex\text{-}inner\text{-}product\text{-}space}(v) \wedge \\ \quad \operatorname{is\text{-}complex\text{-}inner\text{-}product\text{-}space}(w) \wedge \\ \quad \operatorname{is\text{-}compact\text{-}op}(v, w, s) \wedge \\ \quad \operatorname{is\text{-}compact\text{-}op}(v, w, t) \Rightarrow \\ \quad \quad \operatorname{is\text{-}compact\text{-}op}(v, w, {s} + {t}) \end{array}}}
\vnbprop{compact-op-scale}{\vnbstmt{\begin{array}{@{}l@{}} \forall v, w, c, t. \\ \quad \operatorname{is\text{-}complex\text{-}inner\text{-}product\text{-}space}(v) \wedge \\ \quad \operatorname{is\text{-}complex\text{-}inner\text{-}product\text{-}space}(w) \wedge \\ \quad c \in \mathbb{C} \wedge \\ \quad \operatorname{is\text{-}compact\text{-}op}(v, w, t) \Rightarrow \\ \quad \quad \operatorname{is\text{-}compact\text{-}op}(v, w, {c}\,{t}) \end{array}}}
\vnbprop{compact-op-comp-left}{\vnbstmt{\begin{array}{@{}l@{}} \forall v, s, t. \\ \quad (\operatorname{is\text{-}complex\text{-}inner\text{-}product\text{-}space}(v) \wedge ((s \in L({v}, {v})) \wedge \operatorname{is\text{-}compact\text{-}op}(v, v, t))) \Rightarrow  \\ \quad \quad \operatorname{is\text{-}compact\text{-}op}(v, v, {s}{t}) \end{array}}}
\vnbprop{compact-op-seq-r}{\vnbstmt{\begin{array}{@{}l@{}} \forall v, w, t, f, r. \\ \quad \operatorname{is\text{-}complex\text{-}inner\text{-}product\text{-}space}(v) \wedge \\ \quad \operatorname{is\text{-}complex\text{-}inner\text{-}product\text{-}space}(w) \wedge \\ \quad \operatorname{is\text{-}compact\text{-}op}(v, w, t) \wedge \\ \quad r \in \mathbb{R} \wedge \\ \quad 0 \leq r \wedge \\ \quad f \in (\mathbb{N} \to \operatorname{vec}(v)) \wedge \\ \quad \forall q \in \mathbb{N}.\; (\lVert f(q) \rVert \leq r) \Rightarrow \\ \quad \quad \exists \mathit{phq}. \\ \quad \quad \quad \operatorname{strictly\text{-}mono\text{-}nn}(\mathit{phq}) \wedge \\ \quad \quad \quad \exists \mathit{lq}. \forall \mathit{bq} \in (\mathbb{N} \to \operatorname{vec}(w)). \\ \quad \quad \quad \quad \forall q \in \mathbb{N}.\; (\operatorname{bq}(q) = t(f(\operatorname{phq}(q)))) \Rightarrow  \\ \quad \quad \quad \quad \quad \operatorname{converges\text{-}to}(\operatorname{ip\text{-}metric\text{-}space}(w), \mathit{bq}, \mathit{lq}) \end{array}}}
\vnbprop{compact-op-comp-right}{\vnbstmt{\begin{array}{@{}l@{}} \forall v, s, t. \\ \quad (\operatorname{is\text{-}complex\text{-}inner\text{-}product\text{-}space}(v) \wedge (\operatorname{is\text{-}compact\text{-}op}(v, v, t) \wedge (s \in L({v}, {v})))) \Rightarrow  \\ \quad \quad \operatorname{is\text{-}compact\text{-}op}(v, v, {t}{s}) \end{array}}}
\vnbprop{sa-eigenvalue-real}{\vnbstmt{\begin{array}{@{}l@{}} \forall v, t, c, x. \\ \quad \operatorname{is\text{-}complex\text{-}inner\text{-}product\text{-}space}(v) \wedge \\ \quad \operatorname{is\text{-}self\text{-}adjoint\text{-}op}(v, t) \wedge \\ \quad c \in \mathbb{C} \wedge \\ \quad x \in \operatorname{vec}(v) \wedge \\ \quad \neg (x = \operatorname{vzero}(v)) \wedge \\ \quad t(x) = \operatorname{act}(v)(c, x) \Rightarrow \\ \quad \quad c = \operatorname{conjugate}(c) \end{array}}}
\vnbprop{sa-eigenvectors-orthogonal}{\vnbstmt{\begin{array}{@{}l@{}} \forall v, t, a, b, x, y. \\ \quad \operatorname{is\text{-}complex\text{-}inner\text{-}product\text{-}space}(v) \wedge \\ \quad \operatorname{is\text{-}self\text{-}adjoint\text{-}op}(v, t) \wedge \\ \quad a \in \mathbb{C} \wedge \\ \quad b \in \mathbb{C} \wedge \\ \quad b = \operatorname{conjugate}(b) \wedge \\ \quad \neg (a = b) \wedge \\ \quad x \in \operatorname{vec}(v) \wedge \\ \quad y \in \operatorname{vec}(v) \wedge \\ \quad t(x) = \operatorname{act}(v)(a, x) \wedge \\ \quad t(y) = \operatorname{act}(v)(b, y) \Rightarrow \\ \quad \quad \langle\cdot,\cdot\rangle_{v}(x, y) = 0 \end{array}}}
\vnbprop{sa-eigenvalue-real-ev}{\vnbstmt{\begin{array}{@{}l@{}} \forall v, t, c. \\ \quad (\operatorname{is\text{-}complex\text{-}inner\text{-}product\text{-}space}(v) \wedge (\operatorname{is\text{-}self\text{-}adjoint\text{-}op}(v, t) \wedge \operatorname{is\text{-}eigenvalue}(v, t, c))) \Rightarrow  \\ \quad \quad c = \operatorname{conjugate}(c) \end{array}}}
\vnbprop{sa-eigenspaces-orthogonal}{\vnbstmt{\begin{array}{@{}l@{}} \forall v, t, a, b, x, y. \\ \quad \operatorname{is\text{-}complex\text{-}inner\text{-}product\text{-}space}(v) \wedge \\ \quad \operatorname{is\text{-}self\text{-}adjoint\text{-}op}(v, t) \wedge \\ \quad \operatorname{is\text{-}eigenvalue}(v, t, b) \wedge \\ \quad a \in \mathbb{C} \wedge \\ \quad \neg (a = b) \wedge \\ \quad x \in N_{a}({t}) \wedge \\ \quad y \in N_{b}({t}) \Rightarrow \\ \quad \quad \langle\cdot,\cdot\rangle_{v}(x, y) = 0 \end{array}}}
\vnbprop{compact-op-no-separated}{\vnbstmt{\begin{array}{@{}l@{}} \forall v, w, t, f, d. \\ \quad \operatorname{is\text{-}complex\text{-}inner\text{-}product\text{-}space}(v) \wedge \\ \quad \operatorname{is\text{-}complex\text{-}inner\text{-}product\text{-}space}(w) \wedge \\ \quad \operatorname{is\text{-}compact\text{-}op}(v, w, t) \wedge \\ \quad f \in (\mathbb{N} \to \operatorname{vec}(v)) \wedge \\ \quad \forall q \in \mathbb{N}.\; (\lVert f(q) \rVert \leq 1) \wedge \\ \quad \operatorname{pos\text{-}rr}(d) \wedge \\ \quad \forall \mathit{iq}, \mathit{jq} \in \mathbb{N}. \\ \quad \quad (\mathit{iq} < \mathit{jq}) \Rightarrow (d \leq \lVert \operatorname{vadd}(w)(t(f(\mathit{iq})), \operatorname{vneg}(w)(t(f(\mathit{jq})))) \rVert) \Rightarrow \\ \quad \quad \mathit{falsity} \end{array}}}
\vnbprop{compact-op-eigen-not-separated}{\vnbstmt{\begin{array}{@{}l@{}} \forall v, t, e, \ell, d. \\ \quad \operatorname{is\text{-}complex\text{-}inner\text{-}product\text{-}space}(v) \wedge \\ \quad \operatorname{is\text{-}compact\text{-}op}(v, v, t) \wedge \\ \quad e \in (\mathbb{N} \to \operatorname{vec}(v)) \wedge \\ \quad \ell \in (\mathbb{N} \to \mathbb{C}) \wedge \\ \quad \forall n \in \mathbb{N}.\; (\lVert e(n) \rVert = 1) \wedge \\ \quad \forall i, j \in \mathbb{N}.\; ((i < j) \Rightarrow (\langle\cdot,\cdot\rangle_{v}(e(i), e(j)) = 0)) \wedge \\ \quad \forall \mathit{en} \in \mathbb{N}.\; (t(e(\mathit{en})) = \operatorname{act}(v)(\ell(\mathit{en}), e(\mathit{en}))) \wedge \\ \quad \operatorname{pos\text{-}rr}(d) \wedge \\ \quad \forall \mathit{ln} \in \mathbb{N}.\; (d \leq \operatorname{magnitude}(\ell(\mathit{ln}))) \Rightarrow \\ \quad \quad \mathit{falsity} \end{array}}}
\vnbprop{sa-approx-eigen-estimate}{\vnbstmt{\begin{array}{@{}l@{}} \forall v, t, x. \\ \quad \operatorname{is\text{-}complex\text{-}inner\text{-}product\text{-}space}(v) \wedge \\ \quad \operatorname{is\text{-}self\text{-}adjoint\text{-}op}(v, t) \wedge \\ \quad x \in \operatorname{vec}(v) \wedge \\ \quad \lVert x \rVert \leq 1 \Rightarrow \\ \quad \quad \langle\cdot,\cdot\rangle_{v}(\operatorname{vadd}(v)(t(t(x)), \operatorname{vneg}(v)(\operatorname{act}(v)((\lVert {t} \rVert \cdot \lVert {t} \rVert), x))), \\ \quad \quad \quad \operatorname{vadd}(v)(t(t(x)), \operatorname{vneg}(v)(\operatorname{act}(v)((\lVert {t} \rVert \cdot \lVert {t} \rVert), x)))) \\ \quad \quad \leq\; (((\lVert {t} \rVert \cdot \lVert {t} \rVert) \cdot (\lVert {t} \rVert \cdot \lVert {t} \rVert)) - ((\lVert {t} \rVert \cdot \lVert {t} \rVert) \cdot \langle\cdot,\cdot\rangle_{v}(t(x), t(x)))) \end{array}}}

\vnbbookfile{theorem-library/th3a2-compact}{theorem-library/th3a2-compact}
\noindent{\small th3a2-compact.scm -- wave of 2026-11-09/10 (agent TH3-A2), wired 2026-11-10 from scratchpad/th3-1109/a2/th3a2-compact.scm. Library: operator-algebras.}\par
\vnbprop{subseq-ptwise-converges}{\vnbstmt{\begin{array}{@{}l@{}} \forall s, g, i, p, b. \\ \quad (\operatorname{converges\text{-}to}(s, \operatorname{subseq}(g, i), p) \wedge ((b \in (\mathbb{N} \to \operatorname{pts}(s))) \wedge \forall n \in \mathbb{N}.\; (b(n) = g(i(n))))) \Rightarrow  \\ \quad \quad \operatorname{converges\text{-}to}(s, b, p) \end{array}}}
\vnbprop{ptwise-subseq-converges}{\vnbstmt{\begin{array}{@{}l@{}} \forall s, g, i, p, b. \\ \quad \operatorname{converges\text{-}to}(s, b, p) \wedge \\ \quad \operatorname{strictly\text{-}mono\text{-}nn}(i) \wedge \\ \quad g \in (\mathbb{N} \to \operatorname{pts}(s)) \wedge \\ \quad \forall n \in \mathbb{N}.\; (b(n) = g(i(n))) \Rightarrow \\ \quad \quad \operatorname{converges\text{-}to}(s, \operatorname{subseq}(g, i), p) \end{array}}}
\vnbprop{cc-bolzano-weierstrass-ptwise}{\vnbstmt{\begin{array}{@{}l@{}} \forall h \in (\mathbb{N} \to \mathbb{C}), d \in \mathbb{R}. \\ \quad \forall i \in \mathbb{N}.\; (\operatorname{magnitude}(h(i)) \leq d) \Rightarrow  \\ \quad \quad \exists a. \\ \quad \quad \quad \operatorname{strictly\text{-}mono\text{-}nn}(a) \wedge \\ \quad \quad \quad \exists l \in \mathbb{C}. \forall q \in (\mathbb{N} \to \mathbb{C}). \\ \quad \quad \quad \quad \forall \mathit{nq} \in \mathbb{N}.\; (q(\mathit{nq}) = h(a(\mathit{nq}))) \Rightarrow \operatorname{converges\text{-}to}(\operatorname{cc\text{-}ms}, q, l) \end{array}}}
\vnbprop{seq-block-step}{\vnbstmt{\begin{array}{@{}l@{}} \forall s.; \forall h \in (\mathbb{N} \to \operatorname{pts}(s)), b \in \mathrm{Inf}(\mathbb{N}). \\ \quad \forall i. \\ \quad \quad \operatorname{strictly\text{-}mono\text{-}nn}(i) \Rightarrow  \\ \quad \quad \quad \exists a. \\ \quad \quad \quad \quad (\operatorname{strictly\text{-}mono\text{-}nn}(a) \wedge \exists p \in \operatorname{pts}(s).\; \operatorname{converges\text{-}to}(s, \operatorname{subseq}(\operatorname{subseq}(h, i), a), p)) \Rightarrow \\ \quad \quad \exists \beta \in \mathrm{Inf}(\mathbb{N}). \\ \quad \quad \quad (\beta \subseteq b) \wedge \exists p \in \operatorname{pts}(s).\; \operatorname{converges\text{-}along}(s, h, \beta, p) \end{array}}}
\vnbprop{diagonal-subseq}{\vnbstmt{\begin{array}{@{}l@{}} \forall s.; \forall f \in (\mathbb{N} \to (\mathbb{N} \to \operatorname{pts}(s))). \\ \quad \forall k \in \mathbb{N}, i. \\ \quad \quad \operatorname{strictly\text{-}mono\text{-}nn}(i) \Rightarrow  \\ \quad \quad \quad \exists a. \\ \quad \quad \quad \quad (\operatorname{strictly\text{-}mono\text{-}nn}(a) \wedge \exists p \in \operatorname{pts}(s).\; \operatorname{converges\text{-}to}(s, \operatorname{subseq}(\operatorname{subseq}(f(k), i), a), p)) \Rightarrow \\ \quad \quad \exists b. \\ \quad \quad \quad \operatorname{strictly\text{-}mono\text{-}nn}(b) \wedge \\ \quad \quad \quad \forall k \in \mathbb{N}.\; \exists p \in \operatorname{pts}(s).\; \operatorname{converges\text{-}to}(s, \operatorname{subseq}(f(k), b), p) \end{array}}}
\vnbprop{compact-op-subseq-hyp}{\vnbstmt{\begin{array}{@{}l@{}} \forall v, w, t, f, g, i. \\ \quad \operatorname{is\text{-}complex\text{-}inner\text{-}product\text{-}space}(v) \wedge \\ \quad \operatorname{is\text{-}complex\text{-}inner\text{-}product\text{-}space}(w) \wedge \\ \quad \operatorname{is\text{-}compact\text{-}op}(v, w, t) \wedge \\ \quad f \in (\mathbb{N} \to \operatorname{vec}(v)) \wedge \\ \quad \forall q \in \mathbb{N}.\; (\lVert f(q) \rVert \leq 1) \wedge \\ \quad g \in (\mathbb{N} \to \operatorname{vec}(w)) \wedge \\ \quad \forall n \in \mathbb{N}.\; (g(n) = t(f(n))) \wedge \\ \quad \operatorname{strictly\text{-}mono\text{-}nn}(i) \Rightarrow \\ \quad \quad \exists a. \\ \quad \quad \quad \operatorname{strictly\text{-}mono\text{-}nn}(a) \wedge \\ \quad \quad \quad \exists p \in \operatorname{pts}(\operatorname{ip\text{-}metric\text{-}space}(w)). \\ \quad \quad \quad \quad \operatorname{converges\text{-}to}(\operatorname{ip\text{-}metric\text{-}space}(w), \operatorname{subseq}(\operatorname{subseq}(g, i), a), p) \end{array}}}
\vnbprop{compact-op-norm-limit}{\vnbstmt{\begin{array}{@{}l@{}} \forall v, w, t, a. \\ \quad \operatorname{is\text{-}complex\text{-}inner\text{-}product\text{-}space}(v) \wedge \\ \quad \operatorname{is\text{-}hilbert\text{-}space}(w) \wedge \\ \quad t \in L({v}, {w}) \wedge \\ \quad a \in (\mathbb{N} \to L({v}, {w})) \wedge \\ \quad \forall k \in \mathbb{N}.\; \operatorname{is\text{-}compact\text{-}op}(v, w, a(k)) \wedge \\ \quad \forall s. \\ \quad \quad \operatorname{pos\text{-}rr}(s) \Rightarrow  \\ \quad \quad \quad \exists q \in \mathbb{N}. \forall \mathit{xq} \in \operatorname{vec}(v). \\ \quad \quad \quad \quad (\lVert \mathit{xq} \rVert \leq 1) \Rightarrow  \\ \quad \quad \quad \quad \quad \lVert \operatorname{vadd}(w)(t(\mathit{xq}), \operatorname{vneg}(w)(a(q)(\mathit{xq}))) \rVert \leq s \Rightarrow \\ \quad \quad \operatorname{is\text{-}compact\text{-}op}(v, w, t) \end{array}}}
\vnbprop{compact-op-adjoint}{\vnbstmt{\begin{array}{@{}l@{}} \forall v, t. \\ \quad (\operatorname{is\text{-}hilbert\text{-}space}(v) \wedge \operatorname{is\text{-}compact\text{-}op}(v, v, t)) \Rightarrow  \\ \quad \quad \operatorname{is\text{-}compact\text{-}op}(v, v, {t}^{*}) \end{array}}}
\vnbprop{rank-one-op-apply}{\vnbstmt{\begin{array}{@{}l@{}} \forall v, w, a, b, x. \\ \quad \operatorname{is\text{-}complex\text{-}inner\text{-}product\text{-}space}(v) \wedge \\ \quad \operatorname{is\text{-}complex\text{-}inner\text{-}product\text{-}space}(w) \wedge \\ \quad a \in \operatorname{vec}(v) \wedge \\ \quad b \in \operatorname{vec}(w) \wedge \\ \quad x \in \operatorname{vec}(v) \Rightarrow \\ \quad \quad \langle \cdot \mid {a} \rangle {b}(x) = \operatorname{act}(w)(\langle\cdot,\cdot\rangle_{v}(x, a), b) \end{array}}}
\vnbprop{rank-one-op-fun}{\vnbstmt{\begin{array}{@{}l@{}} \forall v, w, a, b. \\ \quad \operatorname{is\text{-}complex\text{-}inner\text{-}product\text{-}space}(v) \wedge \\ \quad \operatorname{is\text{-}complex\text{-}inner\text{-}product\text{-}space}(w) \wedge \\ \quad a \in \operatorname{vec}(v) \wedge \\ \quad b \in \operatorname{vec}(w) \Rightarrow \\ \quad \quad \langle \cdot \mid {a} \rangle {b} \in (\operatorname{vec}(v) \to \operatorname{vec}(w)) \end{array}}}
\vnbprop{rank-one-op-bounded}{\vnbstmt{\begin{array}{@{}l@{}} \forall v, w, a, b. \\ \quad \operatorname{is\text{-}complex\text{-}inner\text{-}product\text{-}space}(v) \wedge \\ \quad \operatorname{is\text{-}complex\text{-}inner\text{-}product\text{-}space}(w) \wedge \\ \quad a \in \operatorname{vec}(v) \wedge \\ \quad b \in \operatorname{vec}(w) \Rightarrow \\ \quad \quad \langle \cdot \mid {a} \rangle {b} \in L({v}, {w}) \end{array}}}
\vnbprop{ipms-dist-act-diff}{\vnbstmt{\begin{array}{@{}l@{}} \forall w, c, d, x. \\ \quad (\operatorname{is\text{-}complex\text{-}inner\text{-}product\text{-}space}(w) \wedge ((c \in \mathbb{C}) \wedge ((d \in \mathbb{C}) \wedge (x \in \operatorname{vec}(w))))) \Rightarrow  \\ \quad \quad \operatorname{dist}(\operatorname{ip\text{-}metric\text{-}space}(w))(\operatorname{act}(w)(c, x), \operatorname{act}(w)(d, x)) = (\lVert x \rVert \cdot \operatorname{magnitude}((c - d))) \end{array}}}
\vnbprop{rank-one-op-compact}{\vnbstmt{\begin{array}{@{}l@{}} \forall v, w, a, b. \\ \quad \operatorname{is\text{-}complex\text{-}inner\text{-}product\text{-}space}(v) \wedge \\ \quad \operatorname{is\text{-}complex\text{-}inner\text{-}product\text{-}space}(w) \wedge \\ \quad a \in \operatorname{vec}(v) \wedge \\ \quad b \in \operatorname{vec}(w) \Rightarrow \\ \quad \quad \operatorname{is\text{-}compact\text{-}op}(v, w, \langle \cdot \mid {a} \rangle {b}) \end{array}}}
\vnbprop{sub-cips-scal}{\vnbstmt{\begin{array}{@{}l@{}} \forall v, k. \\ \quad (\operatorname{is\text{-}complex\text{-}inner\text{-}product\text{-}space}(v) \wedge \operatorname{is\text{-}submodule}(v, k)) \Rightarrow  \\ \quad \quad (\operatorname{scal}(\operatorname{sub\text{-}cips}(v, k)) \simeq \operatorname{normed\text{-}field\text{-}as\text{-}commutative\text{-}ring}(\operatorname{cc\text{-}normed\text{-}field})) \end{array}}}
\vnbprop{sub-cips-vec}{\vnbstmt{\begin{array}{@{}l@{}} \forall v, k. \\ \quad (\operatorname{is\text{-}complex\text{-}inner\text{-}product\text{-}space}(v) \wedge \operatorname{is\text{-}submodule}(v, k)) \Rightarrow  \\ \quad \quad (\operatorname{vec}(\operatorname{sub\text{-}cips}(v, k)) \simeq k) \end{array}}}
\vnbprop{sub-cips-vadd}{\vnbstmt{\begin{array}{@{}l@{}} \forall v, k. \\ \quad (\operatorname{is\text{-}complex\text{-}inner\text{-}product\text{-}space}(v) \wedge \operatorname{is\text{-}submodule}(v, k)) \Rightarrow  \\ \quad \quad (\operatorname{vadd}(\operatorname{sub\text{-}cips}(v, k)) \simeq (\lambda {[}\operatorname{tsbx\_}, \operatorname{tsby\_}] \in (k \times k).\; \operatorname{vadd}(v)(\operatorname{tsbx\_}, \operatorname{tsby\_}))) \end{array}}}
\vnbprop{sub-cips-vzero}{\vnbstmt{\begin{array}{@{}l@{}} \forall v, k. \\ \quad (\operatorname{is\text{-}complex\text{-}inner\text{-}product\text{-}space}(v) \wedge \operatorname{is\text{-}submodule}(v, k)) \Rightarrow  \\ \quad \quad (\operatorname{vzero}(\operatorname{sub\text{-}cips}(v, k)) \simeq \operatorname{vzero}(v)) \end{array}}}
\vnbprop{sub-cips-vneg}{\vnbstmt{\begin{array}{@{}l@{}} \forall v, k. \\ \quad (\operatorname{is\text{-}complex\text{-}inner\text{-}product\text{-}space}(v) \wedge \operatorname{is\text{-}submodule}(v, k)) \Rightarrow  \\ \quad \quad (\operatorname{vneg}(\operatorname{sub\text{-}cips}(v, k)) \simeq (\lambda x \in k.\; \operatorname{vneg}(v)(x))) \end{array}}}
\vnbprop{sub-cips-act}{\vnbstmt{\begin{array}{@{}l@{}} \forall v, k. \\ \quad (\operatorname{is\text{-}complex\text{-}inner\text{-}product\text{-}space}(v) \wedge \operatorname{is\text{-}submodule}(v, k)) \Rightarrow  \\ \quad \quad (\operatorname{act}(\operatorname{sub\text{-}cips}(v, k)) \simeq (\lambda {[}\operatorname{tsbc\_}, \operatorname{tsbx\_}] \in (\mathbb{C} \times k).\; \operatorname{act}(v)(\operatorname{tsbc\_}, \operatorname{tsbx\_}))) \end{array}}}
\vnbprop{sub-cips-ip}{\vnbstmt{\begin{array}{@{}l@{}} \forall v, k. \\ \quad (\operatorname{is\text{-}complex\text{-}inner\text{-}product\text{-}space}(v) \wedge \operatorname{is\text{-}submodule}(v, k)) \Rightarrow  \\ \quad \quad (\langle\cdot,\cdot\rangle_{\operatorname{sub\text{-}cips}(v, k)} \simeq (\lambda {[}\operatorname{tsbx\_}, \operatorname{tsby\_}] \in (k \times k).\; \langle\cdot,\cdot\rangle_{v}(\operatorname{tsbx\_}, \operatorname{tsby\_}))) \end{array}}}
\vnbprop{sub-cips-vadd-apply}{\vnbstmt{\begin{array}{@{}l@{}} \forall v, k, a, b. \\ \quad (\operatorname{is\text{-}complex\text{-}inner\text{-}product\text{-}space}(v) \wedge (\operatorname{is\text{-}submodule}(v, k) \wedge ((a \in k) \wedge (b \in k)))) \Rightarrow  \\ \quad \quad (\operatorname{vadd}(\operatorname{sub\text{-}cips}(v, k))(a, b) \simeq \operatorname{vadd}(v)(a, b)) \end{array}}}
\vnbprop{sub-cips-vneg-apply}{\vnbstmt{\begin{array}{@{}l@{}} \forall v, k, a. \\ \quad (\operatorname{is\text{-}complex\text{-}inner\text{-}product\text{-}space}(v) \wedge (\operatorname{is\text{-}submodule}(v, k) \wedge (a \in k))) \Rightarrow  \\ \quad \quad (\operatorname{vneg}(\operatorname{sub\text{-}cips}(v, k))(a) \simeq \operatorname{vneg}(v)(a)) \end{array}}}
\vnbprop{sub-cips-act-apply}{\vnbstmt{\begin{array}{@{}l@{}} \forall v, k, c, a. \\ \quad (\operatorname{is\text{-}complex\text{-}inner\text{-}product\text{-}space}(v) \wedge (\operatorname{is\text{-}submodule}(v, k) \wedge ((c \in \mathbb{C}) \wedge (a \in k)))) \Rightarrow  \\ \quad \quad (\operatorname{act}(\operatorname{sub\text{-}cips}(v, k))(c, a) \simeq \operatorname{act}(v)(c, a)) \end{array}}}
\vnbprop{sub-cips-ip-apply}{\vnbstmt{\begin{array}{@{}l@{}} \forall v, k, a, b. \\ \quad (\operatorname{is\text{-}complex\text{-}inner\text{-}product\text{-}space}(v) \wedge (\operatorname{is\text{-}submodule}(v, k) \wedge ((a \in k) \wedge (b \in k)))) \Rightarrow  \\ \quad \quad (\langle\cdot,\cdot\rangle_{\operatorname{sub\text{-}cips}(v, k)}(a, b) \simeq \langle\cdot,\cdot\rangle_{v}(a, b)) \end{array}}}
\vnbprop{sub-cips-is-cips}{\vnbstmt{\begin{array}{@{}l@{}} \forall v, k. \\ \quad (\operatorname{is\text{-}complex\text{-}inner\text{-}product\text{-}space}(v) \wedge \operatorname{is\text{-}submodule}(v, k)) \Rightarrow  \\ \quad \quad \operatorname{is\text{-}complex\text{-}inner\text{-}product\text{-}space}(\operatorname{sub\text{-}cips}(v, k)) \end{array}}}
\vnbprop{sub-cips-norm}{\vnbstmt{\begin{array}{@{}l@{}} \forall v, k, a. \\ \quad (\operatorname{is\text{-}complex\text{-}inner\text{-}product\text{-}space}(v) \wedge (\operatorname{is\text{-}submodule}(v, k) \wedge (a \in k))) \Rightarrow  \\ \quad \quad (\lVert a \rVert \simeq \lVert a \rVert) \end{array}}}
\vnbprop{sub-cips-dist}{\vnbstmt{\begin{array}{@{}l@{}} \forall v, k, a, b. \\ \quad (\operatorname{is\text{-}complex\text{-}inner\text{-}product\text{-}space}(v) \wedge (\operatorname{is\text{-}submodule}(v, k) \wedge ((a \in k) \wedge (b \in k)))) \Rightarrow  \\ \quad \quad \operatorname{dist}(\operatorname{ip\text{-}metric\text{-}space}(\operatorname{sub\text{-}cips}(v, k)))(a, b) = \operatorname{dist}(\operatorname{ip\text{-}metric\text{-}space}(v))(a, b) \end{array}}}
\vnbprop{sub-cips-hilbert}{\vnbstmt{\begin{array}{@{}l@{}} \forall v, k. \\ \quad (\operatorname{is\text{-}hilbert\text{-}space}(v) \wedge (\operatorname{is\text{-}submodule}(v, k) \wedge \operatorname{is\text{-}closed}(\operatorname{ip\text{-}metric\text{-}space}(v), k))) \Rightarrow  \\ \quad \quad \operatorname{is\text{-}hilbert\text{-}space}(\operatorname{sub\text{-}cips}(v, k)) \end{array}}}
\vnbprop{sub-cips-restrict-bounded}{\vnbstmt{\begin{array}{@{}l@{}} \forall v, k, t. \\ \quad \operatorname{is\text{-}complex\text{-}inner\text{-}product\text{-}space}(v) \wedge \\ \quad \operatorname{is\text{-}submodule}(v, k) \wedge \\ \quad t \in L({v}, {v}) \wedge \\ \quad \forall x \in k.\; (t(x) \in k) \Rightarrow \\ \quad \quad \operatorname{restrict}(t, k) \in L({\operatorname{sub\text{-}cips}(v, k)}, {\operatorname{sub\text{-}cips}(v, k)}) \end{array}}}
\vnbprop{sub-cips-restrict-self-adjoint}{\vnbstmt{\begin{array}{@{}l@{}} \forall v, k, t. \\ \quad \operatorname{is\text{-}complex\text{-}inner\text{-}product\text{-}space}(v) \wedge \\ \quad \operatorname{is\text{-}submodule}(v, k) \wedge \\ \quad \operatorname{is\text{-}self\text{-}adjoint\text{-}op}(v, t) \wedge \\ \quad \forall x \in k.\; (t(x) \in k) \Rightarrow \\ \quad \quad \operatorname{is\text{-}self\text{-}adjoint\text{-}op}(\operatorname{sub\text{-}cips}(v, k), \operatorname{restrict}(t, k)) \end{array}}}
\vnbprop{sub-cips-restrict-compact}{\vnbstmt{\begin{array}{@{}l@{}} \forall v, k, t. \\ \quad \operatorname{is\text{-}complex\text{-}inner\text{-}product\text{-}space}(v) \wedge \\ \quad \operatorname{is\text{-}submodule}(v, k) \wedge \\ \quad \operatorname{is\text{-}closed}(\operatorname{ip\text{-}metric\text{-}space}(v), k) \wedge \\ \quad \operatorname{is\text{-}compact\text{-}op}(v, v, t) \wedge \\ \quad \forall x \in k.\; (t(x) \in k) \Rightarrow \\ \quad \quad \operatorname{is\text{-}compact\text{-}op}(\operatorname{sub\text{-}cips}(v, k), \operatorname{sub\text{-}cips}(v, k), \operatorname{restrict}(t, k)) \end{array}}}
\vnbprop{sa-perp-invariant}{\vnbstmt{\begin{array}{@{}l@{}} \forall v, k, t.; \forall x \in {k}^{\perp}. \\ \quad \operatorname{is\text{-}complex\text{-}inner\text{-}product\text{-}space}(v) \wedge \\ \quad k \subseteq \operatorname{vec}(v) \wedge \\ \quad \operatorname{is\text{-}self\text{-}adjoint\text{-}op}(v, t) \wedge \\ \quad \forall x \in k.\; (t(x) \in k) \Rightarrow \\ \quad \quad t(x) \in {k}^{\perp} \end{array}}}
\vnbprop{sub-cips-perp-self-adjoint}{\vnbstmt{\begin{array}{@{}l@{}} \forall v, k, t. \\ \quad \operatorname{is\text{-}complex\text{-}inner\text{-}product\text{-}space}(v) \wedge \\ \quad k \subseteq \operatorname{vec}(v) \wedge \\ \quad \operatorname{is\text{-}self\text{-}adjoint\text{-}op}(v, t) \wedge \\ \quad \forall x \in k.\; (t(x) \in k) \Rightarrow \\ \quad \quad \operatorname{is\text{-}self\text{-}adjoint\text{-}op}(\operatorname{sub\text{-}cips}(v, {k}^{\perp}), \operatorname{restrict}(t, {k}^{\perp})) \end{array}}}
\vnbprop{sub-cips-perp-hilbert}{\vnbstmt{\begin{array}{@{}l@{}} \forall v, k. \\ \quad (\operatorname{is\text{-}hilbert\text{-}space}(v) \wedge (k \subseteq \operatorname{vec}(v))) \Rightarrow  \\ \quad \quad \operatorname{is\text{-}hilbert\text{-}space}(\operatorname{sub\text{-}cips}(v, {k}^{\perp})) \end{array}}}
\vnbprop{sq-eigen-split}{\vnbstmt{\begin{array}{@{}l@{}} \forall v, t, z, m. \\ \quad \operatorname{is\text{-}complex\text{-}inner\text{-}product\text{-}space}(v) \wedge \\ \quad t \in L({v}, {v}) \wedge \\ \quad z \in \operatorname{vec}(v) \wedge \\ \quad \neg (z = \operatorname{vzero}(v)) \wedge \\ \quad m \in \mathbb{R} \wedge \\ \quad t(t(z)) = \operatorname{act}(v)((m \cdot m), z) \Rightarrow \\ \quad \quad \exists l.\; (\operatorname{is\text{-}eigenvalue}(v, t, l) \wedge ((l = m) \vee (l = -m))) \end{array}}}

\vnbbookfile{theorem-library/th3a3-finrank}{theorem-library/th3a3-finrank}
\noindent{\small th3a3-finrank.scm -- wave of 2026-11-09/10 (agent TH3-A3), wired 2026-11-10 from scratchpad/th3-1109/a3/th3a3-finrank.scm. Library: operator-algebras.}\par
\vnbprop{hom-ag-finsum}{\vnbstmt{\begin{array}{@{}l@{}} \forall a, b, h.; \forall s \in \mathbf{Set}, f, g. \\ \quad \operatorname{is\text{-}hom\text{-}abelian\text{-}group}(a, b, h) \wedge \\ \quad |s| \in \mathbb{N} \wedge \\ \quad \forall z \in s.\; (f(z) \in \operatorname{carr}(a)) \wedge \\ \quad \forall z \in s.\; (g(z) = h(f(z))) \Rightarrow \\ \quad \quad h(\operatorname{finsum}(a, f, s)) = \operatorname{finsum}(b, g, s) \end{array}}}
\vnbprop{eval-op-hom}{\vnbstmt{\begin{array}{@{}l@{}} \forall v, w, x. \\ \quad (\operatorname{is\text{-}complex\text{-}inner\text{-}product\text{-}space}(v) \wedge (\operatorname{is\text{-}complex\text{-}inner\text{-}product\text{-}space}(w) \wedge (x \in \operatorname{vec}(v)))) \Rightarrow  \\ \quad \quad \operatorname{is\text{-}hom\text{-}abelian\text{-}group}(\operatorname{normed\text{-}ag\text{-}as\text{-}abelian\text{-}group}(\operatorname{bop\text{-}normed\text{-}ag}(v, w)), \\ \quad \quad \quad \operatorname{normed\text{-}ag\text{-}as\text{-}abelian\text{-}group}(\operatorname{ip\text{-}normed\text{-}ag}(w)), \\ \quad \quad \quad (\lambda t \in L({v}, {w}).\; t(x))) \end{array}}}
\vnbprop{op-finsum-apply}{\vnbstmt{\begin{array}{@{}l@{}} \forall v, w, s, p, x. \\ \quad \operatorname{is\text{-}complex\text{-}inner\text{-}product\text{-}space}(v) \wedge \\ \quad \operatorname{is\text{-}complex\text{-}inner\text{-}product\text{-}space}(w) \wedge \\ \quad s \in \mathbf{Set} \wedge \\ \quad |s| \in \mathbb{N} \wedge \\ \quad \forall y \in s.\; (p(y) \in L({v}, {w})) \wedge \\ \quad x \in \operatorname{vec}(v) \Rightarrow \\ \quad \quad \operatorname{finsum}(\operatorname{normed\text{-}ag\text{-}as\text{-}abelian\text{-}group}(\operatorname{bop\text{-}normed\text{-}ag}(v, w)), p, s)(x) \\ \quad \quad =\; \operatorname{finsum}(\operatorname{normed\text{-}ag\text{-}as\text{-}abelian\text{-}group}(\operatorname{ip\text{-}normed\text{-}ag}(w)), (\lambda k \in s.\; p(k)(x)), s) \end{array}}}
\vnbprop{bop-zero-compact}{\vnbstmt{\begin{array}{@{}l@{}} \forall v, w. \\ \quad (\operatorname{is\text{-}complex\text{-}inner\text{-}product\text{-}space}(v) \wedge \operatorname{is\text{-}complex\text{-}inner\text{-}product\text{-}space}(w)) \Rightarrow  \\ \quad \quad \operatorname{is\text{-}compact\text{-}op}(v, w, 0) \end{array}}}
\vnbprop{finite-rank-op-compact}{\vnbstmt{\begin{array}{@{}l@{}} \forall v, w, t. \\ \quad \operatorname{is\text{-}complex\text{-}inner\text{-}product\text{-}space}(v) \wedge \\ \quad \operatorname{is\text{-}complex\text{-}inner\text{-}product\text{-}space}(w) \wedge \\ \quad \operatorname{is\text{-}finite\text{-}rank\text{-}op}(v, w, t) \Rightarrow \\ \quad \quad \operatorname{is\text{-}compact\text{-}op}(v, w, t) \end{array}}}

\vnbbookfile{theorem-library/th3a3-spectral}{theorem-library/th3a3-spectral}
\noindent{\small th3a3-spectral.scm -- wave of 2026-11-09/10 (agent TH3-A3), wired 2026-11-10 from scratchpad/th3-1109/a3/th3a3-spectral.scm. Library: operator-algebras.}\par
\vnbprop{opnorm-approx}{\vnbstmt{\begin{array}{@{}l@{}} \forall v, w, t, e. \\ \quad \operatorname{is\text{-}complex\text{-}inner\text{-}product\text{-}space}(v) \wedge \\ \quad \operatorname{is\text{-}complex\text{-}inner\text{-}product\text{-}space}(w) \wedge \\ \quad t \in L({v}, {w}) \wedge \\ \quad e \in \mathbb{R} \wedge \\ \quad 0 < e \Rightarrow \\ \quad \quad \exists x \in \operatorname{vec}(v). \\ \quad \quad \quad (\lVert x \rVert \leq 1) \wedge ((\lVert {t} \rVert - e) < \lVert t(x) \rVert) \end{array}}}
\vnbprop{sa-max-defect-bound}{\vnbstmt{\begin{array}{@{}l@{}} \forall v, t, x, r, k. \\ \quad \operatorname{is\text{-}complex\text{-}inner\text{-}product\text{-}space}(v) \wedge \\ \quad \operatorname{is\text{-}self\text{-}adjoint\text{-}op}(v, t) \wedge \\ \quad x \in \operatorname{vec}(v) \wedge \\ \quad \lVert x \rVert \leq 1 \wedge \\ \quad r \in \mathbb{R} \wedge \\ \quad 0 \leq r \wedge \\ \quad k \in \mathbb{R} \wedge \\ \quad 0 \leq k \wedge \\ \quad (((2 \cdot \lVert {t} \rVert) \cdot (\lVert {t} \rVert \cdot \lVert {t} \rVert)) \cdot k) \leq 1 \wedge \\ \quad (\lVert {t} \rVert - ((r \cdot r) \cdot k)) < \lVert t(x) \rVert \Rightarrow \\ \quad \quad \lVert \operatorname{vadd}(v)(t(t(x)), \operatorname{vneg}(v)(\operatorname{act}(v)((\lVert {t} \rVert \cdot \lVert {t} \rVert), x))) \rVert \leq r \end{array}}}
\vnbprop{opnorm-approx-seq}{\vnbstmt{\begin{array}{@{}l@{}} \forall v, t, k. \\ \quad \operatorname{is\text{-}complex\text{-}inner\text{-}product\text{-}space}(v) \wedge \\ \quad t \in L({v}, {v}) \wedge \\ \quad 0 < \lVert {t} \rVert \wedge \\ \quad k \in \mathbb{R} \wedge \\ \quad 0 < k \Rightarrow \\ \quad \quad \exists g \in (\mathbb{N} \to \operatorname{vec}(v)). \forall n \in \mathbb{N}. \\ \quad \quad \quad \lVert g(n) \rVert \leq 1 \wedge \\ \quad \quad \quad (\lVert {t} \rVert - ((\operatorname{recip}((n + 1)) \cdot \operatorname{recip}((n + 1))) \cdot k)) < \lVert t(g(n)) \rVert \wedge \\ \quad \quad \quad \lVert {t} \rVert < (2 \cdot \lVert t(g(n)) \rVert) \end{array}}}
\vnbprop{sa-norm-eigenvalue}{\vnbstmt{\begin{array}{@{}l@{}} \forall v, t. \\ \quad \operatorname{is\text{-}complex\text{-}inner\text{-}product\text{-}space}(v) \wedge \\ \quad \operatorname{is\text{-}compact\text{-}op}(v, v, t) \wedge \\ \quad \operatorname{is\text{-}self\text{-}adjoint\text{-}op}(v, t) \wedge \\ \quad \neg (t = 0) \Rightarrow \\ \quad \quad \exists l. \\ \quad \quad \quad \operatorname{is\text{-}eigenvalue}(v, t, l) \wedge ((l = \lVert {t} \rVert) \vee (l = -\lVert {t} \rVert)) \end{array}}}
\vnbprop{sa-has-eigenvector}{\vnbstmt{\begin{array}{@{}l@{}} \forall v, t, x. \\ \quad \operatorname{is\text{-}complex\text{-}inner\text{-}product\text{-}space}(v) \wedge \\ \quad \operatorname{is\text{-}compact\text{-}op}(v, v, t) \wedge \\ \quad \operatorname{is\text{-}self\text{-}adjoint\text{-}op}(v, t) \wedge \\ \quad x \in \operatorname{vec}(v) \wedge \\ \quad \neg (x = \operatorname{vzero}(v)) \Rightarrow \\ \quad \quad \exists l.\; \operatorname{is\text{-}eigenvalue}(v, t, l) \end{array}}}
\vnbprop{sa-eigen-perp-zero}{\vnbstmt{\begin{array}{@{}l@{}} \forall v, t.; \forall x \in {\{z \in \operatorname{vec}(v) : \exists c \in \mathbb{C}.\; (z \in N_{c}({t}))\}}^{\perp}. \\ \quad (\operatorname{is\text{-}hilbert\text{-}space}(v) \wedge (\operatorname{is\text{-}compact\text{-}op}(v, v, t) \wedge \operatorname{is\text{-}self\text{-}adjoint\text{-}op}(v, t))) \Rightarrow  \\ \quad \quad x = \operatorname{vzero}(v) \end{array}}}

\vnbbookfile{theorem-library/ba-g1}{theorem-library/ba-g1}
\noindent{\small ba-g1.scm -- wave of 2026-11-09/10 (agent BA-G1), wired 2026-11-10 from scratchpad/bag1-1109/ba-g1.scm. Library: operator-algebras.}\par
\vnbprop{hsum-scal-acc}{\vnbstmt{\begin{array}{@{}l@{}} \forall v, w. \\ \quad (\operatorname{is\text{-}complex\text{-}inner\text{-}product\text{-}space}(v) \wedge \operatorname{is\text{-}complex\text{-}inner\text{-}product\text{-}space}(w)) \Rightarrow  \\ \quad \quad (\operatorname{scal}({v} \oplus {w}) \simeq \operatorname{normed\text{-}field\text{-}as\text{-}commutative\text{-}ring}(\operatorname{cc\text{-}normed\text{-}field})) \end{array}}}
\vnbprop{hsum-vec-acc}{\vnbstmt{\begin{array}{@{}l@{}} \forall v, w. \\ \quad (\operatorname{is\text{-}complex\text{-}inner\text{-}product\text{-}space}(v) \wedge \operatorname{is\text{-}complex\text{-}inner\text{-}product\text{-}space}(w)) \Rightarrow  \\ \quad \quad (\operatorname{vec}({v} \oplus {w}) \simeq (\operatorname{vec}(v) \times \operatorname{vec}(w))) \end{array}}}
\vnbprop{hsum-vadd-acc}{\vnbstmt{\begin{array}{@{}l@{}} \forall v, w. \\ \quad (\operatorname{is\text{-}complex\text{-}inner\text{-}product\text{-}space}(v) \wedge \operatorname{is\text{-}complex\text{-}inner\text{-}product\text{-}space}(w)) \Rightarrow  \\ \quad \quad (\operatorname{vadd}({v} \oplus {w}) \simeq \operatorname{hsum\text{-}vadd}(v, w)) \end{array}}}
\vnbprop{hsum-vzero-acc}{\vnbstmt{\begin{array}{@{}l@{}} \forall v, w. \\ \quad (\operatorname{is\text{-}complex\text{-}inner\text{-}product\text{-}space}(v) \wedge \operatorname{is\text{-}complex\text{-}inner\text{-}product\text{-}space}(w)) \Rightarrow  \\ \quad \quad (\operatorname{vzero}({v} \oplus {w}) \simeq {[}\operatorname{vzero}(v), \operatorname{vzero}(w)]) \end{array}}}
\vnbprop{hsum-vneg-acc}{\vnbstmt{\begin{array}{@{}l@{}} \forall v, w. \\ \quad (\operatorname{is\text{-}complex\text{-}inner\text{-}product\text{-}space}(v) \wedge \operatorname{is\text{-}complex\text{-}inner\text{-}product\text{-}space}(w)) \Rightarrow  \\ \quad \quad (\operatorname{vneg}({v} \oplus {w}) \simeq \operatorname{hsum\text{-}vneg}(v, w)) \end{array}}}
\vnbprop{hsum-act-acc}{\vnbstmt{\begin{array}{@{}l@{}} \forall v, w. \\ \quad (\operatorname{is\text{-}complex\text{-}inner\text{-}product\text{-}space}(v) \wedge \operatorname{is\text{-}complex\text{-}inner\text{-}product\text{-}space}(w)) \Rightarrow  \\ \quad \quad (\operatorname{act}({v} \oplus {w}) \simeq \operatorname{hsum\text{-}act}(v, w)) \end{array}}}
\vnbprop{hsum-ip-acc}{\vnbstmt{\begin{array}{@{}l@{}} \forall v, w. \\ \quad (\operatorname{is\text{-}complex\text{-}inner\text{-}product\text{-}space}(v) \wedge \operatorname{is\text{-}complex\text{-}inner\text{-}product\text{-}space}(w)) \Rightarrow  \\ \quad \quad (\langle\cdot,\cdot\rangle_{{v} \oplus {w}} \simeq \operatorname{hsum\text{-}ip}(v, w)) \end{array}}}
\vnbprop{hsum-ip-pair}{\vnbstmt{\begin{array}{@{}l@{}} \forall v, w, a, b, c, d. \\ \quad \operatorname{is\text{-}complex\text{-}inner\text{-}product\text{-}space}(v) \wedge \\ \quad \operatorname{is\text{-}complex\text{-}inner\text{-}product\text{-}space}(w) \wedge \\ \quad a \in \operatorname{vec}(v) \wedge \\ \quad b \in \operatorname{vec}(w) \wedge \\ \quad c \in \operatorname{vec}(v) \wedge \\ \quad d \in \operatorname{vec}(w) \Rightarrow \\ \quad \quad \langle\cdot,\cdot\rangle_{{v} \oplus {w}}({[}a, b], {[}c, d]) = (\langle\cdot,\cdot\rangle_{v}(a, c) + \langle\cdot,\cdot\rangle_{w}(b, d)) \end{array}}}
\vnbprop{hsum-vadd-pair}{\vnbstmt{\begin{array}{@{}l@{}} \forall v, w, a, b, c, d. \\ \quad \operatorname{is\text{-}complex\text{-}inner\text{-}product\text{-}space}(v) \wedge \\ \quad \operatorname{is\text{-}complex\text{-}inner\text{-}product\text{-}space}(w) \wedge \\ \quad a \in \operatorname{vec}(v) \wedge \\ \quad b \in \operatorname{vec}(w) \wedge \\ \quad c \in \operatorname{vec}(v) \wedge \\ \quad d \in \operatorname{vec}(w) \Rightarrow \\ \quad \quad \operatorname{vadd}({v} \oplus {w})({[}a, b], {[}c, d]) = {[}\operatorname{vadd}(v)(a, c), \operatorname{vadd}(w)(b, d)] \end{array}}}
\vnbprop{hsum-vneg-pair}{\vnbstmt{\begin{array}{@{}l@{}} \forall v, w, a, b. \\ \quad \operatorname{is\text{-}complex\text{-}inner\text{-}product\text{-}space}(v) \wedge \\ \quad \operatorname{is\text{-}complex\text{-}inner\text{-}product\text{-}space}(w) \wedge \\ \quad a \in \operatorname{vec}(v) \wedge \\ \quad b \in \operatorname{vec}(w) \Rightarrow \\ \quad \quad \operatorname{vneg}({v} \oplus {w})({[}a, b]) = {[}\operatorname{vneg}(v)(a), \operatorname{vneg}(w)(b)] \end{array}}}
\vnbprop{hsum-act-pair}{\vnbstmt{\begin{array}{@{}l@{}} \forall v, w, k, a, b. \\ \quad \operatorname{is\text{-}complex\text{-}inner\text{-}product\text{-}space}(v) \wedge \\ \quad \operatorname{is\text{-}complex\text{-}inner\text{-}product\text{-}space}(w) \wedge \\ \quad k \in \mathbb{C} \wedge \\ \quad a \in \operatorname{vec}(v) \wedge \\ \quad b \in \operatorname{vec}(w) \Rightarrow \\ \quad \quad \operatorname{act}({v} \oplus {w})(k, {[}a, b]) = {[}\operatorname{act}(v)(k, a), \operatorname{act}(w)(k, b)] \end{array}}}
\vnbprop{hsum-vadd-type}{\vnbstmt{\begin{array}{@{}l@{}} \forall v, w. \\ \quad (\operatorname{is\text{-}complex\text{-}inner\text{-}product\text{-}space}(v) \wedge \operatorname{is\text{-}complex\text{-}inner\text{-}product\text{-}space}(w)) \Rightarrow  \\ \quad \quad \operatorname{hsum\text{-}vadd}(v, w) \\ \quad \quad \in\; (((\operatorname{vec}(v) \times \operatorname{vec}(w)) \times (\operatorname{vec}(v) \times \operatorname{vec}(w))) \to (\operatorname{vec}(v) \times \operatorname{vec}(w))) \end{array}}}
\vnbprop{hsum-vneg-type}{\vnbstmt{\begin{array}{@{}l@{}} \forall v, w. \\ \quad (\operatorname{is\text{-}complex\text{-}inner\text{-}product\text{-}space}(v) \wedge \operatorname{is\text{-}complex\text{-}inner\text{-}product\text{-}space}(w)) \Rightarrow  \\ \quad \quad \operatorname{hsum\text{-}vneg}(v, w) \in ((\operatorname{vec}(v) \times \operatorname{vec}(w)) \to (\operatorname{vec}(v) \times \operatorname{vec}(w))) \end{array}}}
\vnbprop{hsum-act-type}{\vnbstmt{\begin{array}{@{}l@{}} \forall v, w. \\ \quad (\operatorname{is\text{-}complex\text{-}inner\text{-}product\text{-}space}(v) \wedge \operatorname{is\text{-}complex\text{-}inner\text{-}product\text{-}space}(w)) \Rightarrow  \\ \quad \quad \operatorname{hsum\text{-}act}(v, w) \in ((\mathbb{C} \times (\operatorname{vec}(v) \times \operatorname{vec}(w))) \to (\operatorname{vec}(v) \times \operatorname{vec}(w))) \end{array}}}
\vnbprop{hsum-ip-type}{\vnbstmt{\begin{array}{@{}l@{}} \forall v, w. \\ \quad (\operatorname{is\text{-}complex\text{-}inner\text{-}product\text{-}space}(v) \wedge \operatorname{is\text{-}complex\text{-}inner\text{-}product\text{-}space}(w)) \Rightarrow  \\ \quad \quad \operatorname{hsum\text{-}ip}(v, w) \in (((\operatorname{vec}(v) \times \operatorname{vec}(w)) \times (\operatorname{vec}(v) \times \operatorname{vec}(w))) \to \mathbb{C}) \end{array}}}
\vnbprop{hsum-is-set}{\vnbstmt{\begin{array}{@{}l@{}} \forall v, w. \\ \quad (\operatorname{is\text{-}complex\text{-}inner\text{-}product\text{-}space}(v) \wedge \operatorname{is\text{-}complex\text{-}inner\text{-}product\text{-}space}(w)) \Rightarrow  \\ \quad \quad {v} \oplus {w} \in \mathbf{Set} \end{array}}}
\vnbprop{relation-pair-ext}{\vnbstmt{\begin{array}{@{}l@{}} \forall s, a, b, c. \\ \quad b \subseteq (s \times a) \wedge \\ \quad c \subseteq (s \times a) \wedge \\ \quad \forall \mathit{ugc} \in s, d \in a.\; (({[}\mathit{ugc}, d] \in b) \Leftrightarrow ({[}\mathit{ugc}, d] \in c)) \Rightarrow \\ \quad \quad b = c \end{array}}}
\vnbprop{relation-pair-subset}{\vnbstmt{\begin{array}{@{}l@{}} \forall s, a, b, c. \\ \quad ((b \subseteq (s \times a)) \wedge \forall \mathit{ugc} \in s, d \in a.\; (({[}\mathit{ugc}, d] \in b) \Rightarrow ({[}\mathit{ugc}, d] \in c))) \Rightarrow  \\ \quad \quad b \subseteq c \end{array}}}
\vnbprop{graph-adjoint-subset}{\vnbstmt{\begin{array}{@{}l@{}} \forall v, g. \\ \quad \operatorname{is\text{-}complex\text{-}inner\text{-}product\text{-}space}(v) \Rightarrow  \\ \quad \quad {g}^{*} \subseteq (\operatorname{vec}(v) \times \operatorname{vec}(v)) \end{array}}}
\vnbprop{graph-j-subset}{\vnbstmt{\begin{array}{@{}l@{}} \forall v, g. \\ \quad \operatorname{is\text{-}complex\text{-}inner\text{-}product\text{-}space}(v) \Rightarrow  \\ \quad \quad J{g} \subseteq (\operatorname{vec}(v) \times \operatorname{vec}(v)) \end{array}}}
\vnbprop{graph-inverse-subset}{\vnbstmt{\begin{array}{@{}l@{}} \forall v, g. \\ \quad \operatorname{is\text{-}complex\text{-}inner\text{-}product\text{-}space}(v) \Rightarrow  \\ \quad \quad {g}^{-1} \subseteq (\operatorname{vec}(v) \times \operatorname{vec}(v)) \end{array}}}
\vnbprop{cayley-graph-subset}{\vnbstmt{\begin{array}{@{}l@{}} \forall v, g. \\ \quad \operatorname{is\text{-}complex\text{-}inner\text{-}product\text{-}space}(v) \Rightarrow  \\ \quad \quad C({g}) \subseteq (\operatorname{vec}(v) \times \operatorname{vec}(v)) \end{array}}}
\vnbprop{hsum-perp-subset}{\vnbstmt{\begin{array}{@{}l@{}} \forall v, s. \\ \quad \operatorname{is\text{-}complex\text{-}inner\text{-}product\text{-}space}(v) \Rightarrow  \\ \quad \quad {s}^{\perp} \subseteq (\operatorname{vec}(v) \times \operatorname{vec}(v)) \end{array}}}
\vnbprop{graph-adjoint-is-set}{\vnbstmt{\begin{array}{@{}l@{}} \forall v, g. \\ \quad \operatorname{is\text{-}complex\text{-}inner\text{-}product\text{-}space}(v) \Rightarrow ({g}^{*} \in \mathbf{Set}) \end{array}}}
\vnbprop{graph-j-is-set}{\vnbstmt{\forall v, g.\; (\operatorname{is\text{-}complex\text{-}inner\text{-}product\text{-}space}(v) \Rightarrow (J{g} \in \mathbf{Set}))}}
\vnbprop{graph-inverse-is-set}{\vnbstmt{\begin{array}{@{}l@{}} \forall v, g. \\ \quad \operatorname{is\text{-}complex\text{-}inner\text{-}product\text{-}space}(v) \Rightarrow ({g}^{-1} \in \mathbf{Set}) \end{array}}}
\vnbprop{cayley-graph-is-set}{\vnbstmt{\begin{array}{@{}l@{}} \forall v, g. \\ \quad \operatorname{is\text{-}complex\text{-}inner\text{-}product\text{-}space}(v) \Rightarrow (C({g}) \in \mathbf{Set}) \end{array}}}
\vnbprop{hsum-perp-is-set}{\vnbstmt{\begin{array}{@{}l@{}} \forall v, s. \\ \quad \operatorname{is\text{-}complex\text{-}inner\text{-}product\text{-}space}(v) \Rightarrow ({s}^{\perp} \in \mathbf{Set}) \end{array}}}
\vnbprop{graph-j-pair-iff}{\vnbstmt{\begin{array}{@{}l@{}} \forall v, g, a, b. \\ \quad (\operatorname{is\text{-}complex\text{-}inner\text{-}product\text{-}space}(v) \wedge ((a \in \operatorname{vec}(v)) \wedge (b \in \operatorname{vec}(v)))) \Rightarrow  \\ \quad \quad ({[}a, b] \in J{g}) \Leftrightarrow ({[}b, \operatorname{vneg}(v)(a)] \in g) \end{array}}}
\vnbprop{graph-inverse-pair-iff}{\vnbstmt{\begin{array}{@{}l@{}} \forall v, g, a, b. \\ \quad (\operatorname{is\text{-}complex\text{-}inner\text{-}product\text{-}space}(v) \wedge ((a \in \operatorname{vec}(v)) \wedge (b \in \operatorname{vec}(v)))) \Rightarrow  \\ \quad \quad ({[}a, b] \in {g}^{-1}) \Leftrightarrow ({[}b, a] \in g) \end{array}}}
\vnbprop{graph-adjoint-pair-iff}{\vnbstmt{\begin{array}{@{}l@{}} \forall v, g, a, b. \\ \quad \operatorname{is\text{-}complex\text{-}inner\text{-}product\text{-}space}(v) \wedge \\ \quad g \subseteq (\operatorname{vec}(v) \times \operatorname{vec}(v)) \wedge \\ \quad a \in \operatorname{vec}(v) \wedge \\ \quad b \in \operatorname{vec}(v) \Rightarrow \\ \quad \quad {[}a, b] \in {g}^{*} \\ \quad \quad \Leftrightarrow\; \forall c, d \in \operatorname{vec}(v). \\ \quad \quad \quad ({[}c, d] \in g) \Rightarrow (\langle\cdot,\cdot\rangle_{v}(c, b) = \langle\cdot,\cdot\rangle_{v}(d, a)) \end{array}}}
\vnbprop{hsum-perp-pair-iff}{\vnbstmt{\begin{array}{@{}l@{}} \forall v, s, a, b. \\ \quad \operatorname{is\text{-}complex\text{-}inner\text{-}product\text{-}space}(v) \wedge \\ \quad s \subseteq (\operatorname{vec}(v) \times \operatorname{vec}(v)) \wedge \\ \quad a \in \operatorname{vec}(v) \wedge \\ \quad b \in \operatorname{vec}(v) \Rightarrow \\ \quad \quad {[}a, b] \in {s}^{\perp} \\ \quad \quad \Leftrightarrow\; \forall c, d \in \operatorname{vec}(v). \\ \quad \quad \quad ({[}c, d] \in s) \Rightarrow ((\langle\cdot,\cdot\rangle_{v}(a, c) + \langle\cdot,\cdot\rangle_{v}(b, d)) = 0) \end{array}}}
\vnbprop{cc-eq-of-add-neg-zero}{\vnbstmt{\forall x, y.\; (((x \in \mathbb{C}) \wedge ((y \in \mathbb{C}) \wedge ((x + -y) = 0))) \Rightarrow (x = y))}}
\vnbprop{graph-adjoint-j-perp}{\vnbstmt{\begin{array}{@{}l@{}} \forall v, g. \\ \quad (\operatorname{is\text{-}complex\text{-}inner\text{-}product\text{-}space}(v) \wedge (g \subseteq (\operatorname{vec}(v) \times \operatorname{vec}(v)))) \Rightarrow  \\ \quad \quad {g}^{*} = J{{g}^{\perp}} \end{array}}}
\vnbprop{graph-adjoint-perp-j}{\vnbstmt{\begin{array}{@{}l@{}} \forall v, g. \\ \quad (\operatorname{is\text{-}complex\text{-}inner\text{-}product\text{-}space}(v) \wedge (g \subseteq (\operatorname{vec}(v) \times \operatorname{vec}(v)))) \Rightarrow  \\ \quad \quad {g}^{*} = {J{g}}^{\perp} \end{array}}}
\vnbprop{graph-adjoint-via-j}{\vnbstmt{\begin{array}{@{}l@{}} \forall v, g. \\ \quad (\operatorname{is\text{-}complex\text{-}inner\text{-}product\text{-}space}(v) \wedge (g \subseteq (\operatorname{vec}(v) \times \operatorname{vec}(v)))) \Rightarrow  \\ \quad \quad {g}^{*} = {J{g}}^{\perp} \wedge \\ \quad \quad {g}^{*} = J{{g}^{\perp}} \end{array}}}
\vnbprop{perp-of-closure}{\vnbstmt{\begin{array}{@{}l@{}} \forall v, k. \\ \quad (\operatorname{is\text{-}complex\text{-}inner\text{-}product\text{-}space}(v) \wedge (k \subseteq \operatorname{vec}(v))) \Rightarrow  \\ \quad \quad {\operatorname{closure}(\operatorname{ip\text{-}metric\text{-}space}(v), k)}^{\perp} = {k}^{\perp} \end{array}}}
\vnbprop{perp-perp-closure}{\vnbstmt{\begin{array}{@{}l@{}} \forall v, k. \\ \quad (\operatorname{is\text{-}hilbert\text{-}space}(v) \wedge \operatorname{is\text{-}submodule}(v, k)) \Rightarrow  \\ \quad \quad {{k}^{\perp}}^{\perp} = \operatorname{closure}(\operatorname{ip\text{-}metric\text{-}space}(v), k) \end{array}}}
\vnbprop{cayley-pair-in}{\vnbstmt{\begin{array}{@{}l@{}} \forall v, g, c, d. \\ \quad (\operatorname{is\text{-}complex\text{-}inner\text{-}product\text{-}space}(v) \wedge ((c \in \operatorname{vec}(v)) \wedge ((d \in \operatorname{vec}(v)) \wedge ({[}c, d] \in g)))) \Rightarrow  \\ \quad \quad {[}\operatorname{vadd}(v)(\operatorname{act}(v)(+i, c), d), \operatorname{vadd}(v)(\operatorname{act}(v)(-+i, c), d)] \in C({g}) \end{array}}}
\vnbprop{cayley-norm-iff}{\vnbstmt{\begin{array}{@{}l@{}} \forall v, c, d. \\ \quad (\operatorname{is\text{-}complex\text{-}inner\text{-}product\text{-}space}(v) \wedge ((c \in \operatorname{vec}(v)) \wedge (d \in \operatorname{vec}(v)))) \Rightarrow  \\ \quad \quad \lVert \operatorname{vadd}(v)(\operatorname{act}(v)(+i, c), d) \rVert = \lVert \operatorname{vadd}(v)(\operatorname{act}(v)(-+i, c), d) \rVert \\ \quad \quad \Leftrightarrow\; \langle\cdot,\cdot\rangle_{v}(c, d) = \langle\cdot,\cdot\rangle_{v}(d, c) \end{array}}}
\vnbprop{cayley-isometric-of-symmetric}{\vnbstmt{\begin{array}{@{}l@{}} \forall v, g. \\ \quad \operatorname{is\text{-}complex\text{-}inner\text{-}product\text{-}space}(v) \wedge \\ \quad g \subseteq (\operatorname{vec}(v) \times \operatorname{vec}(v)) \wedge \\ \quad \operatorname{is\text{-}symmetric\text{-}graph}(v, g) \Rightarrow \\ \quad \quad \operatorname{is\text{-}isometric\text{-}graph}(v, C({g})) \end{array}}}
\vnbprop{cayley-symmetric-of-isometric}{\vnbstmt{\begin{array}{@{}l@{}} \forall v, g. \\ \quad \operatorname{is\text{-}complex\text{-}inner\text{-}product\text{-}space}(v) \wedge \\ \quad \operatorname{is\text{-}linear\text{-}graph}(v, g) \wedge \\ \quad \operatorname{is\text{-}isometric\text{-}graph}(v, C({g})) \Rightarrow \\ \quad \quad \operatorname{is\text{-}symmetric\text{-}graph}(v, g) \end{array}}}
\vnbprop{cayley-symmetric-iff-isometric}{\vnbstmt{\begin{array}{@{}l@{}} \forall v, g. \\ \quad (\operatorname{is\text{-}complex\text{-}inner\text{-}product\text{-}space}(v) \wedge \operatorname{is\text{-}linear\text{-}graph}(v, g)) \Rightarrow  \\ \quad \quad \operatorname{is\text{-}symmetric\text{-}graph}(v, g) \Leftrightarrow \operatorname{is\text{-}isometric\text{-}graph}(v, C({g})) \end{array}}}
\vnbprop{cayley-pair-iff}{\vnbstmt{\begin{array}{@{}l@{}} \forall v, h, a, b. \\ \quad \operatorname{is\text{-}complex\text{-}inner\text{-}product\text{-}space}(v) \wedge \\ \quad h \subseteq (\operatorname{vec}(v) \times \operatorname{vec}(v)) \wedge \\ \quad a \in \operatorname{vec}(v) \wedge \\ \quad b \in \operatorname{vec}(v) \Rightarrow \\ \quad \quad {[}a, b] \in C({h}) \\ \quad \quad \Leftrightarrow\; {[}\operatorname{act}(v)(-+i, \operatorname{act}(v)(1/2, \operatorname{vadd}(v)(a, \operatorname{vneg}(v)(b)))), \operatorname{act}(v)(1/2, \operatorname{vadd}(v)(a, b))] \\ \quad \quad \in\; h \end{array}}}
\vnbprop{cayley-of-adjoint}{\vnbstmt{\begin{array}{@{}l@{}} \forall v, g. \\ \quad (\operatorname{is\text{-}complex\text{-}inner\text{-}product\text{-}space}(v) \wedge (g \subseteq (\operatorname{vec}(v) \times \operatorname{vec}(v)))) \Rightarrow  \\ \quad \quad C({{g}^{*}}) = {{C({g})}^{*}}^{-1} \end{array}}}
\vnbprop{unitary-graph-preserves-ip}{\vnbstmt{\begin{array}{@{}l@{}} \forall v, u, c, d, e, f. \\ \quad \operatorname{is\text{-}complex\text{-}inner\text{-}product\text{-}space}(v) \wedge \\ \quad \operatorname{is\text{-}unitary\text{-}graph}(v, u) \wedge \\ \quad c \in \operatorname{vec}(v) \wedge \\ \quad d \in \operatorname{vec}(v) \wedge \\ \quad e \in \operatorname{vec}(v) \wedge \\ \quad f \in \operatorname{vec}(v) \wedge \\ \quad {[}c, d] \in u \wedge \\ \quad {[}e, f] \in u \Rightarrow \\ \quad \quad \langle\cdot,\cdot\rangle_{v}(c, e) = \langle\cdot,\cdot\rangle_{v}(d, f) \end{array}}}
\vnbprop{graph-dom-unfold}{\vnbstmt{\begin{array}{@{}l@{}} \forall v, u. \\ \quad (\mathrm{dom}\,{u} \simeq \{x \in \operatorname{vec}(v) : \exists y \in \operatorname{vec}(v).\; ({[}x, y] \in u)\}) \end{array}}}
\vnbprop{graph-ran-unfold}{\vnbstmt{\begin{array}{@{}l@{}} \forall v, u. \\ \quad (\mathrm{ran}\,{u} \simeq \{y \in \operatorname{vec}(v) : \exists x \in \operatorname{vec}(v).\; ({[}x, y] \in u)\}) \end{array}}}
\vnbprop{unitary-graph-adjoint-is-inverse}{\vnbstmt{\begin{array}{@{}l@{}} \forall v, u. \\ \quad (\operatorname{is\text{-}complex\text{-}inner\text{-}product\text{-}space}(v) \wedge \operatorname{is\text{-}unitary\text{-}graph}(v, u)) \Rightarrow  \\ \quad \quad {u}^{*} = {u}^{-1} \end{array}}}
\vnbprop{graph-inverse-involutive}{\vnbstmt{\begin{array}{@{}l@{}} \forall v, g. \\ \quad (\operatorname{is\text{-}complex\text{-}inner\text{-}product\text{-}space}(v) \wedge (g \subseteq (\operatorname{vec}(v) \times \operatorname{vec}(v)))) \Rightarrow  \\ \quad \quad {{g}^{-1}}^{-1} = g \end{array}}}
\vnbprop{cayley-graph-injective}{\vnbstmt{\begin{array}{@{}l@{}} \forall v, h, k. \\ \quad \operatorname{is\text{-}complex\text{-}inner\text{-}product\text{-}space}(v) \wedge \\ \quad h \subseteq (\operatorname{vec}(v) \times \operatorname{vec}(v)) \wedge \\ \quad k \subseteq (\operatorname{vec}(v) \times \operatorname{vec}(v)) \wedge \\ \quad C({h}) = C({k}) \Rightarrow \\ \quad \quad h = k \end{array}}}
\vnbprop{self-adjoint-of-cayley-unitary}{\vnbstmt{\begin{array}{@{}l@{}} \forall v, g. \\ \quad \operatorname{is\text{-}complex\text{-}inner\text{-}product\text{-}space}(v) \wedge \\ \quad g \subseteq (\operatorname{vec}(v) \times \operatorname{vec}(v)) \wedge \\ \quad \operatorname{is\text{-}unitary\text{-}graph}(v, C({g})) \Rightarrow \\ \quad \quad \operatorname{is\text{-}self\text{-}adjoint\text{-}graph}(v, g) \end{array}}}
\vnbprop{graph-j-of-perp}{\vnbstmt{\begin{array}{@{}l@{}} \forall v, s. \\ \quad (\operatorname{is\text{-}complex\text{-}inner\text{-}product\text{-}space}(v) \wedge (s \subseteq (\operatorname{vec}(v) \times \operatorname{vec}(v)))) \Rightarrow  \\ \quad \quad J{{s}^{\perp}} = {J{s}}^{\perp} \end{array}}}
\vnbprop{graph-j-j-linear}{\vnbstmt{\begin{array}{@{}l@{}} \forall v, g. \\ \quad (\operatorname{is\text{-}complex\text{-}inner\text{-}product\text{-}space}(v) \wedge \operatorname{is\text{-}linear\text{-}graph}(v, g)) \Rightarrow  \\ \quad \quad J{J{g}} = g \end{array}}}
\vnbprop{graph-closure-is-double-adjoint}{\vnbstmt{\begin{array}{@{}l@{}} \forall v, g. \\ \quad (\operatorname{is\text{-}complex\text{-}inner\text{-}product\text{-}space}(v) \wedge (\operatorname{is\text{-}hilbert\text{-}space}({v} \oplus {v}) \wedge \operatorname{is\text{-}linear\text{-}graph}(v, g))) \Rightarrow  \\ \quad \quad \overline{g} = {{g}^{*}}^{*} \end{array}}}

\vnbbookfile{theorem-library/ba-g2}{theorem-library/ba-g2}
\noindent{\small ba-g2.scm -- wave of 2026-11-09/10 (agent BA-G2), wired 2026-11-10 from scratchpad/bag2-1109/ba-g2.scm. Library: operator-algebras.}\par
\vnbprop{hsum-is-cips}{\vnbstmt{\begin{array}{@{}l@{}} \forall v, w. \\ \quad (\operatorname{is\text{-}complex\text{-}inner\text{-}product\text{-}space}(v) \wedge \operatorname{is\text{-}complex\text{-}inner\text{-}product\text{-}space}(w)) \Rightarrow  \\ \quad \quad \operatorname{is\text{-}complex\text{-}inner\text{-}product\text{-}space}({v} \oplus {w}) \end{array}}}
\vnbprop{rr-sqrt-add-le}{\vnbstmt{\begin{array}{@{}l@{}} \forall x, y. \\ \quad ((x \in \mathbb{R}) \wedge ((y \in \mathbb{R}) \wedge ((0 \leq x) \wedge (0 \leq y)))) \Rightarrow  \\ \quad \quad \sqrt{(x + y)} \leq (\sqrt{x} + \sqrt{y}) \end{array}}}
\vnbprop{hsum-dist-ge-1}{\vnbstmt{\begin{array}{@{}l@{}} \forall v, w, p, q. \\ \quad \operatorname{is\text{-}complex\text{-}inner\text{-}product\text{-}space}(v) \wedge \\ \quad \operatorname{is\text{-}complex\text{-}inner\text{-}product\text{-}space}(w) \wedge \\ \quad p \in (\operatorname{vec}(v) \times \operatorname{vec}(w)) \wedge \\ \quad q \in (\operatorname{vec}(v) \times \operatorname{vec}(w)) \Rightarrow \\ \quad \quad \operatorname{dist}(\operatorname{ip\text{-}metric\text{-}space}(v))(\operatorname{nth}(1, p), \operatorname{nth}(1, q)) \leq \operatorname{dist}(\operatorname{ip\text{-}metric\text{-}space}({v} \oplus {w}))(p, q) \end{array}}}
\vnbprop{hsum-dist-ge-2}{\vnbstmt{\begin{array}{@{}l@{}} \forall v, w, p, q. \\ \quad \operatorname{is\text{-}complex\text{-}inner\text{-}product\text{-}space}(v) \wedge \\ \quad \operatorname{is\text{-}complex\text{-}inner\text{-}product\text{-}space}(w) \wedge \\ \quad p \in (\operatorname{vec}(v) \times \operatorname{vec}(w)) \wedge \\ \quad q \in (\operatorname{vec}(v) \times \operatorname{vec}(w)) \Rightarrow \\ \quad \quad \operatorname{dist}(\operatorname{ip\text{-}metric\text{-}space}(w))(\operatorname{nth}(2, p), \operatorname{nth}(2, q)) \leq \operatorname{dist}(\operatorname{ip\text{-}metric\text{-}space}({v} \oplus {w}))(p, q) \end{array}}}
\vnbprop{hsum-dist-le-sum}{\vnbstmt{\begin{array}{@{}l@{}} \forall v, w, p, c, d. \\ \quad \operatorname{is\text{-}complex\text{-}inner\text{-}product\text{-}space}(v) \wedge \\ \quad \operatorname{is\text{-}complex\text{-}inner\text{-}product\text{-}space}(w) \wedge \\ \quad p \in (\operatorname{vec}(v) \times \operatorname{vec}(w)) \wedge \\ \quad c \in \operatorname{vec}(v) \wedge \\ \quad d \in \operatorname{vec}(w) \Rightarrow \\ \quad \quad \operatorname{dist}(\operatorname{ip\text{-}metric\text{-}space}({v} \oplus {w}))(p, {[}c, d]) \\ \quad \quad \leq\; \operatorname{dist}(\operatorname{ip\text{-}metric\text{-}space}(v))(\operatorname{nth}(1, p), c) + \operatorname{dist}(\operatorname{ip\text{-}metric\text{-}space}(w))(\operatorname{nth}(2, p), d) \end{array}}}
\vnbprop{hsum-is-hilbert}{\vnbstmt{\begin{array}{@{}l@{}} \forall v, w. \\ \quad (\operatorname{is\text{-}hilbert\text{-}space}(v) \wedge \operatorname{is\text{-}hilbert\text{-}space}(w)) \Rightarrow \operatorname{is\text{-}hilbert\text{-}space}({v} \oplus {w}) \end{array}}}
\vnbprop{graph-closure-double-adjoint}{\vnbstmt{\begin{array}{@{}l@{}} \forall v, g. \\ \quad (\operatorname{is\text{-}hilbert\text{-}space}(v) \wedge \operatorname{is\text{-}linear\text{-}graph}(v, g)) \Rightarrow  \\ \quad \quad \overline{g} = {{g}^{*}}^{*} \end{array}}}
\vnbprop{graph-adjoint-linear}{\vnbstmt{\begin{array}{@{}l@{}} \forall v, g. \\ \quad (\operatorname{is\text{-}complex\text{-}inner\text{-}product\text{-}space}(v) \wedge (g \subseteq (\operatorname{vec}(v) \times \operatorname{vec}(v)))) \Rightarrow  \\ \quad \quad \operatorname{is\text{-}linear\text{-}graph}(v, {g}^{*}) \end{array}}}
\vnbprop{graph-adjoint-closed}{\vnbstmt{\begin{array}{@{}l@{}} \forall v, g. \\ \quad (\operatorname{is\text{-}complex\text{-}inner\text{-}product\text{-}space}(v) \wedge (g \subseteq (\operatorname{vec}(v) \times \operatorname{vec}(v)))) \Rightarrow  \\ \quad \quad \operatorname{is\text{-}closed\text{-}graph}(v, {g}^{*}) \end{array}}}
\vnbprop{linear-graph-functional-iff}{\vnbstmt{\begin{array}{@{}l@{}} \forall v, g. \\ \quad (\operatorname{is\text{-}complex\text{-}inner\text{-}product\text{-}space}(v) \wedge \operatorname{is\text{-}linear\text{-}graph}(v, g)) \Rightarrow  \\ \quad \quad \operatorname{is\text{-}functional\text{-}graph}(v, g) \Leftrightarrow \forall y \in \operatorname{vec}(v).\; (({[}\operatorname{vzero}(v), y] \in g) \Rightarrow (y = \operatorname{vzero}(v))) \end{array}}}
\vnbprop{graph-adjoint-zero-pair-iff}{\vnbstmt{\begin{array}{@{}l@{}} \forall v, g, y. \\ \quad (\operatorname{is\text{-}complex\text{-}inner\text{-}product\text{-}space}(v) \wedge ((g \subseteq (\operatorname{vec}(v) \times \operatorname{vec}(v))) \wedge (y \in \operatorname{vec}(v)))) \Rightarrow  \\ \quad \quad ({[}\operatorname{vzero}(v), y] \in {g}^{*}) \Leftrightarrow (y \in {\mathrm{dom}\,{g}}^{\perp}) \end{array}}}
\vnbprop{hilbert-dense-iff-perp-trivial}{\vnbstmt{\begin{array}{@{}l@{}} \forall v, k. \\ \quad (\operatorname{is\text{-}hilbert\text{-}space}(v) \wedge \operatorname{is\text{-}submodule}(v, k)) \Rightarrow  \\ \quad \quad \operatorname{is\text{-}dense}(\operatorname{ip\text{-}metric\text{-}space}(v), k) \Leftrightarrow \forall z \in {k}^{\perp}.\; (z = \operatorname{vzero}(v)) \end{array}}}
\vnbprop{graph-dom-submodule}{\vnbstmt{\begin{array}{@{}l@{}} \forall v, g. \\ \quad (\operatorname{is\text{-}complex\text{-}inner\text{-}product\text{-}space}(v) \wedge \operatorname{is\text{-}linear\text{-}graph}(v, g)) \Rightarrow  \\ \quad \quad \operatorname{is\text{-}submodule}(v, \mathrm{dom}\,{g}) \end{array}}}
\vnbprop{graph-ran-submodule}{\vnbstmt{\begin{array}{@{}l@{}} \forall v, g. \\ \quad (\operatorname{is\text{-}complex\text{-}inner\text{-}product\text{-}space}(v) \wedge \operatorname{is\text{-}linear\text{-}graph}(v, g)) \Rightarrow  \\ \quad \quad \operatorname{is\text{-}submodule}(v, \mathrm{ran}\,{g}) \end{array}}}
\vnbprop{graph-adjoint-functional-iff-dense}{\vnbstmt{\begin{array}{@{}l@{}} \forall v, g. \\ \quad (\operatorname{is\text{-}hilbert\text{-}space}(v) \wedge \operatorname{is\text{-}linear\text{-}graph}(v, g)) \Rightarrow  \\ \quad \quad \operatorname{is\text{-}functional\text{-}graph}(v, {g}^{*}) \Leftrightarrow \operatorname{is\text{-}dense}(\operatorname{ip\text{-}metric\text{-}space}(v), \mathrm{dom}\,{g}) \end{array}}}
\vnbprop{closable-iff-adjoint-dense}{\vnbstmt{\begin{array}{@{}l@{}} \forall v, g. \\ \quad (\operatorname{is\text{-}hilbert\text{-}space}(v) \wedge \operatorname{is\text{-}linear\text{-}graph}(v, g)) \Rightarrow  \\ \quad \quad \operatorname{is\text{-}closable\text{-}graph}(v, g) \Leftrightarrow \operatorname{is\text{-}dense}(\operatorname{ip\text{-}metric\text{-}space}(v), \mathrm{dom}\,{{g}^{*}}) \end{array}}}
\vnbprop{graph-dom-of-inverse}{\vnbstmt{\begin{array}{@{}l@{}} \forall v, g. \\ \quad (\operatorname{is\text{-}complex\text{-}inner\text{-}product\text{-}space}(v) \wedge (g \subseteq (\operatorname{vec}(v) \times \operatorname{vec}(v)))) \Rightarrow  \\ \quad \quad \mathrm{dom}\,{{g}^{-1}} = \mathrm{ran}\,{g} \end{array}}}
\vnbprop{graph-adjoint-of-inverse}{\vnbstmt{\begin{array}{@{}l@{}} \forall v, g. \\ \quad (\operatorname{is\text{-}complex\text{-}inner\text{-}product\text{-}space}(v) \wedge (g \subseteq (\operatorname{vec}(v) \times \operatorname{vec}(v)))) \Rightarrow  \\ \quad \quad {{g}^{-1}}^{*} = {{g}^{*}}^{-1} \end{array}}}
\vnbprop{graph-inverse-linear}{\vnbstmt{\begin{array}{@{}l@{}} \forall v, g. \\ \quad (\operatorname{is\text{-}complex\text{-}inner\text{-}product\text{-}space}(v) \wedge \operatorname{is\text{-}linear\text{-}graph}(v, g)) \Rightarrow  \\ \quad \quad \operatorname{is\text{-}linear\text{-}graph}(v, {g}^{-1}) \end{array}}}
\vnbprop{inverse-adjoint-functional-iff-dense-range}{\vnbstmt{\begin{array}{@{}l@{}} \forall v, g. \\ \quad (\operatorname{is\text{-}hilbert\text{-}space}(v) \wedge \operatorname{is\text{-}linear\text{-}graph}(v, g)) \Rightarrow  \\ \quad \quad \operatorname{is\text{-}functional\text{-}graph}(v, {{g}^{-1}}^{*}) \\ \quad \quad \Leftrightarrow\; \operatorname{is\text{-}dense}(\operatorname{ip\text{-}metric\text{-}space}(v), \mathrm{ran}\,{g}) \end{array}}}
\vnbprop{inverse-closable-iff-adjoint-dense-range}{\vnbstmt{\begin{array}{@{}l@{}} \forall v, g. \\ \quad (\operatorname{is\text{-}hilbert\text{-}space}(v) \wedge \operatorname{is\text{-}linear\text{-}graph}(v, g)) \Rightarrow  \\ \quad \quad \operatorname{is\text{-}closable\text{-}graph}(v, {g}^{-1}) \\ \quad \quad \Leftrightarrow\; \operatorname{is\text{-}dense}(\operatorname{ip\text{-}metric\text{-}space}(v), \mathrm{ran}\,{{g}^{*}}) \end{array}}}
\vnbprop{cayley-graph-linear}{\vnbstmt{\begin{array}{@{}l@{}} \forall v, g. \\ \quad (\operatorname{is\text{-}complex\text{-}inner\text{-}product\text{-}space}(v) \wedge \operatorname{is\text{-}linear\text{-}graph}(v, g)) \Rightarrow  \\ \quad \quad \operatorname{is\text{-}linear\text{-}graph}(v, C({g})) \end{array}}}
\vnbprop{graph-adjoint-zero-pair-ran-iff}{\vnbstmt{\begin{array}{@{}l@{}} \forall v, g, y. \\ \quad (\operatorname{is\text{-}complex\text{-}inner\text{-}product\text{-}space}(v) \wedge ((g \subseteq (\operatorname{vec}(v) \times \operatorname{vec}(v))) \wedge (y \in \operatorname{vec}(v)))) \Rightarrow  \\ \quad \quad ({[}y, \operatorname{vzero}(v)] \in {g}^{*}) \Leftrightarrow (y \in {\mathrm{ran}\,{g}}^{\perp}) \end{array}}}
\vnbprop{cayley-adjoint-of-self-adjoint}{\vnbstmt{\begin{array}{@{}l@{}} \forall v, g. \\ \quad \operatorname{is\text{-}complex\text{-}inner\text{-}product\text{-}space}(v) \wedge \\ \quad g \subseteq (\operatorname{vec}(v) \times \operatorname{vec}(v)) \wedge \\ \quad \operatorname{is\text{-}self\text{-}adjoint\text{-}graph}(v, g) \Rightarrow \\ \quad \quad {C({g})}^{*} = {C({g})}^{-1} \end{array}}}
\vnbprop{cayley-dom-dense-of-self-adjoint}{\vnbstmt{\begin{array}{@{}l@{}} \forall v, g. \\ \quad (\operatorname{is\text{-}hilbert\text{-}space}(v) \wedge ((g \subseteq (\operatorname{vec}(v) \times \operatorname{vec}(v))) \wedge \operatorname{is\text{-}self\text{-}adjoint\text{-}graph}(v, g))) \Rightarrow  \\ \quad \quad \operatorname{is\text{-}dense}(\operatorname{ip\text{-}metric\text{-}space}(v), \mathrm{dom}\,{C({g})}) \end{array}}}
\vnbprop{cayley-ran-dense-of-self-adjoint}{\vnbstmt{\begin{array}{@{}l@{}} \forall v, g. \\ \quad (\operatorname{is\text{-}hilbert\text{-}space}(v) \wedge ((g \subseteq (\operatorname{vec}(v) \times \operatorname{vec}(v))) \wedge \operatorname{is\text{-}self\text{-}adjoint\text{-}graph}(v, g))) \Rightarrow  \\ \quad \quad \operatorname{is\text{-}dense}(\operatorname{ip\text{-}metric\text{-}space}(v), \mathrm{ran}\,{C({g})}) \end{array}}}
\vnbprop{cayley-dom-of-self-adjoint}{\vnbstmt{\begin{array}{@{}l@{}} \forall v, g. \\ \quad (\operatorname{is\text{-}hilbert\text{-}space}(v) \wedge ((g \subseteq (\operatorname{vec}(v) \times \operatorname{vec}(v))) \wedge \operatorname{is\text{-}self\text{-}adjoint\text{-}graph}(v, g))) \Rightarrow  \\ \quad \quad \mathrm{dom}\,{C({g})} = \operatorname{vec}(v) \end{array}}}
\vnbprop{cayley-ran-of-self-adjoint}{\vnbstmt{\begin{array}{@{}l@{}} \forall v, g. \\ \quad (\operatorname{is\text{-}hilbert\text{-}space}(v) \wedge ((g \subseteq (\operatorname{vec}(v) \times \operatorname{vec}(v))) \wedge \operatorname{is\text{-}self\text{-}adjoint\text{-}graph}(v, g))) \Rightarrow  \\ \quad \quad \mathrm{ran}\,{C({g})} = \operatorname{vec}(v) \end{array}}}
\vnbprop{cayley-unitary-of-self-adjoint}{\vnbstmt{\begin{array}{@{}l@{}} \forall v, g. \\ \quad (\operatorname{is\text{-}hilbert\text{-}space}(v) \wedge ((g \subseteq (\operatorname{vec}(v) \times \operatorname{vec}(v))) \wedge \operatorname{is\text{-}self\text{-}adjoint\text{-}graph}(v, g))) \Rightarrow  \\ \quad \quad \operatorname{is\text{-}unitary\text{-}graph}(v, C({g})) \end{array}}}
\vnbprop{self-adjoint-iff-cayley-unitary}{\vnbstmt{\begin{array}{@{}l@{}} \forall v, g. \\ \quad (\operatorname{is\text{-}hilbert\text{-}space}(v) \wedge (g \subseteq (\operatorname{vec}(v) \times \operatorname{vec}(v)))) \Rightarrow  \\ \quad \quad \operatorname{is\text{-}self\text{-}adjoint\text{-}graph}(v, g) \Leftrightarrow \operatorname{is\text{-}unitary\text{-}graph}(v, C({g})) \end{array}}}

\vnbbookfile{theorem-library/ba-optop}{theorem-library/ba-optop}
\noindent{\small ba-optop.scm -- wave of 2026-11-09/10 (agent BA-OT1), wired 2026-11-10 from scratchpad/baot1-1109/ba-optop.scm. Library: operator-algebras.}\par
\vnbprop{bh-module-scal}{\vnbstmt{\begin{array}{@{}l@{}} \forall v. \\ \quad \operatorname{is\text{-}complex\text{-}inner\text{-}product\text{-}space}(v) \Rightarrow  \\ \quad \quad (\operatorname{scal}(L({v})) \simeq \operatorname{normed\text{-}field\text{-}as\text{-}commutative\text{-}ring}(\operatorname{cc\text{-}normed\text{-}field})) \end{array}}}
\vnbprop{bh-module-vec}{\vnbstmt{\begin{array}{@{}l@{}} \forall v. \\ \quad \operatorname{is\text{-}complex\text{-}inner\text{-}product\text{-}space}(v) \Rightarrow (\operatorname{vec}(L({v})) \simeq L({v}, {v})) \end{array}}}
\vnbprop{bh-module-vadd}{\vnbstmt{\begin{array}{@{}l@{}} \forall v. \\ \quad \operatorname{is\text{-}complex\text{-}inner\text{-}product\text{-}space}(v) \Rightarrow (\operatorname{vadd}(L({v})) \simeq \operatorname{bh\text{-}vadd}(v)) \end{array}}}
\vnbprop{bh-module-vzero}{\vnbstmt{\begin{array}{@{}l@{}} \forall v. \\ \quad \operatorname{is\text{-}complex\text{-}inner\text{-}product\text{-}space}(v) \Rightarrow (\operatorname{vzero}(L({v})) \simeq 0) \end{array}}}
\vnbprop{bh-module-vneg}{\vnbstmt{\begin{array}{@{}l@{}} \forall v. \\ \quad \operatorname{is\text{-}complex\text{-}inner\text{-}product\text{-}space}(v) \Rightarrow (\operatorname{vneg}(L({v})) \simeq \operatorname{bh\text{-}vneg}(v)) \end{array}}}
\vnbprop{bh-module-act}{\vnbstmt{\begin{array}{@{}l@{}} \forall v. \\ \quad \operatorname{is\text{-}complex\text{-}inner\text{-}product\text{-}space}(v) \Rightarrow (\operatorname{act}(L({v})) \simeq \operatorname{bh\text{-}act}(v)) \end{array}}}
\vnbprop{bh-module-is-module}{\vnbstmt{\forall v.\; (\operatorname{is\text{-}complex\text{-}inner\text{-}product\text{-}space}(v) \Rightarrow \operatorname{is\text{-}module}(L({v})))}}
\vnbprop{cc-sn-module-scal}{\vnbstmt{(\operatorname{scal}(\operatorname{cc\text{-}sn\text{-}module}) \simeq \operatorname{normed\text{-}field\text{-}as\text{-}commutative\text{-}ring}(\operatorname{cc\text{-}normed\text{-}field}))}}
\vnbprop{cc-sn-module-vec}{\vnbstmt{(\operatorname{vec}(\operatorname{cc\text{-}sn\text{-}module}) \simeq \operatorname{carr}(\operatorname{normed\text{-}field\text{-}as\text{-}commutative\text{-}ring}(\operatorname{cc\text{-}normed\text{-}field})))}}
\vnbprop{cc-sn-module-vadd}{\vnbstmt{(\operatorname{vadd}(\operatorname{cc\text{-}sn\text{-}module}) \simeq \operatorname{add}(\operatorname{normed\text{-}field\text{-}as\text{-}commutative\text{-}ring}(\operatorname{cc\text{-}normed\text{-}field})))}}
\vnbprop{cc-sn-module-vzero}{\vnbstmt{(\operatorname{vzero}(\operatorname{cc\text{-}sn\text{-}module}) \simeq \operatorname{zero}(\operatorname{normed\text{-}field\text{-}as\text{-}commutative\text{-}ring}(\operatorname{cc\text{-}normed\text{-}field})))}}
\vnbprop{cc-sn-module-vneg}{\vnbstmt{(\operatorname{vneg}(\operatorname{cc\text{-}sn\text{-}module}) \simeq \operatorname{neg}(\operatorname{normed\text{-}field\text{-}as\text{-}commutative\text{-}ring}(\operatorname{cc\text{-}normed\text{-}field})))}}
\vnbprop{cc-sn-module-act}{\vnbstmt{(\operatorname{act}(\operatorname{cc\text{-}sn\text{-}module}) \simeq \operatorname{mul}(\operatorname{normed\text{-}field\text{-}as\text{-}commutative\text{-}ring}(\operatorname{cc\text{-}normed\text{-}field})))}}
\vnbprop{cc-sn-module-is-module}{\vnbstmt{\operatorname{is\text{-}module}(\operatorname{cc\text{-}sn\text{-}module})}}
\vnbprop{cc-sn-family}{\vnbstmt{\operatorname{is\text{-}seminorm\text{-}family\text{-}on}(\operatorname{cc\text{-}normed\text{-}field}, \operatorname{cc\text{-}sn\text{-}module}, \{0\}, \operatorname{cc\text{-}sn\text{-}pf})}}
\vnbprop{cc-sn-top-is-tvs}{\vnbstmt{\operatorname{is\text{-}tvs}(\operatorname{cc\text{-}normed\text{-}field}, \operatorname{cc\text{-}sn\text{-}module}, \operatorname{cc\text{-}sn\text{-}top})}}
\vnbprop{sot-pf-type}{\vnbstmt{\begin{array}{@{}l@{}} \forall v. \\ \quad \operatorname{is\text{-}complex\text{-}inner\text{-}product\text{-}space}(v) \Rightarrow  \\ \quad \quad \operatorname{sot\text{-}pf}(v) \in (\operatorname{vec}(v) \to (L({v}, {v}) \to \mathbb{R})) \end{array}}}
\vnbprop{sot-pf-apply}{\vnbstmt{\begin{array}{@{}l@{}} \forall v, x, t. \\ \quad (\operatorname{is\text{-}complex\text{-}inner\text{-}product\text{-}space}(v) \wedge ((x \in \operatorname{vec}(v)) \wedge (t \in L({v}, {v})))) \Rightarrow  \\ \quad \quad \operatorname{sot\text{-}pf}(v)(x)(t) = \lVert t(x) \rVert \end{array}}}
\vnbprop{sot-seminorm}{\vnbstmt{\begin{array}{@{}l@{}} \forall v, x. \\ \quad (\operatorname{is\text{-}complex\text{-}inner\text{-}product\text{-}space}(v) \wedge (x \in \operatorname{vec}(v))) \Rightarrow  \\ \quad \quad \operatorname{is\text{-}seminorm}(\operatorname{cc\text{-}normed\text{-}field}, L({v}), \operatorname{sot\text{-}pf}(v)(x)) \end{array}}}
\vnbprop{sot-separating}{\vnbstmt{\begin{array}{@{}l@{}} \forall v, t. \\ \quad \operatorname{is\text{-}complex\text{-}inner\text{-}product\text{-}space}(v) \wedge \\ \quad t \in L({v}, {v}) \wedge \\ \quad \forall x \in \operatorname{vec}(v).\; (\operatorname{sot\text{-}pf}(v)(x)(t) = 0) \Rightarrow \\ \quad \quad t = 0 \end{array}}}
\vnbprop{sot-seminorm-family}{\vnbstmt{\begin{array}{@{}l@{}} \forall v. \\ \quad \operatorname{is\text{-}complex\text{-}inner\text{-}product\text{-}space}(v) \Rightarrow  \\ \quad \quad \operatorname{is\text{-}seminorm\text{-}family\text{-}on}(\operatorname{cc\text{-}normed\text{-}field}, L({v}), \operatorname{vec}(v), \operatorname{sot\text{-}pf}(v)) \end{array}}}
\vnbprop{sot-is-tvs}{\vnbstmt{\begin{array}{@{}l@{}} \forall v. \\ \quad \operatorname{is\text{-}complex\text{-}inner\text{-}product\text{-}space}(v) \Rightarrow  \\ \quad \quad \operatorname{is\text{-}tvs}(\operatorname{cc\text{-}normed\text{-}field}, L({v}), L({v})_{s}) \end{array}}}
\vnbprop{wot-pf-type}{\vnbstmt{\begin{array}{@{}l@{}} \forall v. \\ \quad \operatorname{is\text{-}complex\text{-}inner\text{-}product\text{-}space}(v) \Rightarrow  \\ \quad \quad \operatorname{wot\text{-}pf}(v) \in ((\operatorname{vec}(v) \times \operatorname{vec}(v)) \to (L({v}, {v}) \to \mathbb{R})) \end{array}}}
\vnbprop{wot-pf-apply}{\vnbstmt{\begin{array}{@{}l@{}} \forall v, p, t. \\ \quad (\operatorname{is\text{-}complex\text{-}inner\text{-}product\text{-}space}(v) \wedge ((p \in (\operatorname{vec}(v) \times \operatorname{vec}(v))) \wedge (t \in L({v}, {v})))) \Rightarrow  \\ \quad \quad \operatorname{wot\text{-}pf}(v)(p)(t) = \operatorname{magnitude}(\langle\cdot,\cdot\rangle_{v}(t(\operatorname{nth}(1, p)), \operatorname{nth}(2, p))) \end{array}}}
\vnbprop{wot-pf-apply-pair}{\vnbstmt{\begin{array}{@{}l@{}} \forall v, x, y, t. \\ \quad (\operatorname{is\text{-}complex\text{-}inner\text{-}product\text{-}space}(v) \wedge ((x \in \operatorname{vec}(v)) \wedge ((y \in \operatorname{vec}(v)) \wedge (t \in L({v}, {v}))))) \Rightarrow  \\ \quad \quad \operatorname{wot\text{-}pf}(v)({[}x, y])(t) = \operatorname{magnitude}(\langle\cdot,\cdot\rangle_{v}(t(x), y)) \end{array}}}
\vnbprop{wot-seminorm}{\vnbstmt{\begin{array}{@{}l@{}} \forall v, p. \\ \quad (\operatorname{is\text{-}complex\text{-}inner\text{-}product\text{-}space}(v) \wedge (p \in (\operatorname{vec}(v) \times \operatorname{vec}(v)))) \Rightarrow  \\ \quad \quad \operatorname{is\text{-}seminorm}(\operatorname{cc\text{-}normed\text{-}field}, L({v}), \operatorname{wot\text{-}pf}(v)(p)) \end{array}}}
\vnbprop{wot-separating}{\vnbstmt{\begin{array}{@{}l@{}} \forall v, t. \\ \quad \operatorname{is\text{-}complex\text{-}inner\text{-}product\text{-}space}(v) \wedge \\ \quad t \in L({v}, {v}) \wedge \\ \quad \forall p \in (\operatorname{vec}(v) \times \operatorname{vec}(v)).\; (\operatorname{wot\text{-}pf}(v)(p)(t) = 0) \Rightarrow \\ \quad \quad t = 0 \end{array}}}
\vnbprop{wot-seminorm-family}{\vnbstmt{\begin{array}{@{}l@{}} \forall v. \\ \quad \operatorname{is\text{-}complex\text{-}inner\text{-}product\text{-}space}(v) \Rightarrow  \\ \quad \quad \operatorname{is\text{-}seminorm\text{-}family\text{-}on}(\operatorname{cc\text{-}normed\text{-}field}, L({v}), (\operatorname{vec}(v) \times \operatorname{vec}(v)), \operatorname{wot\text{-}pf}(v)) \end{array}}}
\vnbprop{wot-is-tvs}{\vnbstmt{\begin{array}{@{}l@{}} \forall v. \\ \quad \operatorname{is\text{-}complex\text{-}inner\text{-}product\text{-}space}(v) \Rightarrow  \\ \quad \quad \operatorname{is\text{-}tvs}(\operatorname{cc\text{-}normed\text{-}field}, L({v}), L({v})_{w}) \end{array}}}
\vnbprop{norm-optop-pf-type}{\vnbstmt{\begin{array}{@{}l@{}} \forall v. \\ \quad \operatorname{is\text{-}complex\text{-}inner\text{-}product\text{-}space}(v) \Rightarrow  \\ \quad \quad \operatorname{norm\text{-}optop\text{-}pf}(v) \in (\{0\} \to (L({v}, {v}) \to \mathbb{R})) \end{array}}}
\vnbprop{norm-optop-pf-apply}{\vnbstmt{\begin{array}{@{}l@{}} \forall v, k, t. \\ \quad (\operatorname{is\text{-}complex\text{-}inner\text{-}product\text{-}space}(v) \wedge ((k \in \{0\}) \wedge (t \in L({v}, {v})))) \Rightarrow  \\ \quad \quad \operatorname{norm\text{-}optop\text{-}pf}(v)(k)(t) = \lVert {t} \rVert \end{array}}}
\vnbprop{norm-optop-seminorm}{\vnbstmt{\begin{array}{@{}l@{}} \forall v, k. \\ \quad (\operatorname{is\text{-}complex\text{-}inner\text{-}product\text{-}space}(v) \wedge (k \in \{0\})) \Rightarrow  \\ \quad \quad \operatorname{is\text{-}seminorm}(\operatorname{cc\text{-}normed\text{-}field}, L({v}), \operatorname{norm\text{-}optop\text{-}pf}(v)(k)) \end{array}}}
\vnbprop{norm-optop-separating}{\vnbstmt{\begin{array}{@{}l@{}} \forall v, t. \\ \quad \operatorname{is\text{-}complex\text{-}inner\text{-}product\text{-}space}(v) \wedge \\ \quad t \in L({v}, {v}) \wedge \\ \quad \forall k \in \{0\}.\; (\operatorname{norm\text{-}optop\text{-}pf}(v)(k)(t) = 0) \Rightarrow \\ \quad \quad t = 0 \end{array}}}
\vnbprop{norm-optop-seminorm-family}{\vnbstmt{\begin{array}{@{}l@{}} \forall v. \\ \quad \operatorname{is\text{-}complex\text{-}inner\text{-}product\text{-}space}(v) \Rightarrow  \\ \quad \quad \operatorname{is\text{-}seminorm\text{-}family\text{-}on}(\operatorname{cc\text{-}normed\text{-}field}, L({v}), \{0\}, \operatorname{norm\text{-}optop\text{-}pf}(v)) \end{array}}}
\vnbprop{norm-optop-is-tvs}{\vnbstmt{\begin{array}{@{}l@{}} \forall v. \\ \quad \operatorname{is\text{-}complex\text{-}inner\text{-}product\text{-}space}(v) \Rightarrow  \\ \quad \quad \operatorname{is\text{-}tvs}(\operatorname{cc\text{-}normed\text{-}field}, L({v}), L({v})) \end{array}}}
\vnbprop{bh-id-linear}{\vnbstmt{\begin{array}{@{}l@{}} \forall v. \\ \quad \operatorname{is\text{-}complex\text{-}inner\text{-}product\text{-}space}(v) \Rightarrow  \\ \quad \quad \operatorname{is\text{-}linear\text{-}map}(L({v}), L({v}), \operatorname{id\text{-}fun}(L({v}, {v}))) \end{array}}}
\vnbprop{bop-right-mult-linear}{\vnbstmt{\begin{array}{@{}l@{}} \forall v, t. \\ \quad (\operatorname{is\text{-}complex\text{-}inner\text{-}product\text{-}space}(v) \wedge (t \in L({v}, {v}))) \Rightarrow  \\ \quad \quad \operatorname{is\text{-}linear\text{-}map}(L({v}), L({v}), (\lambda s \in L({v}, {v}).\; {s}{t})) \end{array}}}
\vnbprop{bop-left-mult-linear}{\vnbstmt{\begin{array}{@{}l@{}} \forall v, t. \\ \quad (\operatorname{is\text{-}complex\text{-}inner\text{-}product\text{-}space}(v) \wedge (t \in L({v}, {v}))) \Rightarrow  \\ \quad \quad \operatorname{is\text{-}linear\text{-}map}(L({v}), L({v}), (\lambda s \in L({v}, {v}).\; {t}{s})) \end{array}}}
\vnbprop{norm-to-sot-continuous}{\vnbstmt{\begin{array}{@{}l@{}} \forall v. \\ \quad \operatorname{is\text{-}complex\text{-}inner\text{-}product\text{-}space}(v) \Rightarrow  \\ \quad \quad \operatorname{is\text{-}hom\text{-}top\text{-}space}(L({v}), L({v})_{s}, \operatorname{id\text{-}fun}(L({v}, {v}))) \end{array}}}
\vnbprop{sot-to-wot-continuous}{\vnbstmt{\begin{array}{@{}l@{}} \forall v. \\ \quad \operatorname{is\text{-}complex\text{-}inner\text{-}product\text{-}space}(v) \Rightarrow  \\ \quad \quad \operatorname{is\text{-}hom\text{-}top\text{-}space}(L({v})_{s}, L({v})_{w}, \operatorname{id\text{-}fun}(L({v}, {v}))) \end{array}}}
\vnbprop{bop-right-mult-sot-continuous}{\vnbstmt{\begin{array}{@{}l@{}} \forall v, t. \\ \quad (\operatorname{is\text{-}complex\text{-}inner\text{-}product\text{-}space}(v) \wedge (t \in L({v}, {v}))) \Rightarrow  \\ \quad \quad \operatorname{is\text{-}hom\text{-}top\text{-}space}(L({v})_{s}, L({v})_{s}, (\lambda s \in L({v}, {v}).\; {s}{t})) \end{array}}}
\vnbprop{bop-right-mult-wot-continuous}{\vnbstmt{\begin{array}{@{}l@{}} \forall v, t. \\ \quad (\operatorname{is\text{-}complex\text{-}inner\text{-}product\text{-}space}(v) \wedge (t \in L({v}, {v}))) \Rightarrow  \\ \quad \quad \operatorname{is\text{-}hom\text{-}top\text{-}space}(L({v})_{w}, L({v})_{w}, (\lambda s \in L({v}, {v}).\; {s}{t})) \end{array}}}
\vnbprop{bop-left-mult-sot-continuous}{\vnbstmt{\begin{array}{@{}l@{}} \forall v, t. \\ \quad (\operatorname{is\text{-}complex\text{-}inner\text{-}product\text{-}space}(v) \wedge (t \in L({v}, {v}))) \Rightarrow  \\ \quad \quad \operatorname{is\text{-}hom\text{-}top\text{-}space}(L({v})_{s}, L({v})_{s}, (\lambda s \in L({v}, {v}).\; {t}{s})) \end{array}}}
\vnbprop{bop-left-mult-wot-continuous}{\vnbstmt{\begin{array}{@{}l@{}} \forall v, t. \\ \quad (\operatorname{is\text{-}complex\text{-}inner\text{-}product\text{-}space}(v) \wedge (\operatorname{is\text{-}hilbert\text{-}space}(v) \wedge (t \in L({v}, {v})))) \Rightarrow  \\ \quad \quad \operatorname{is\text{-}hom\text{-}top\text{-}space}(L({v})_{w}, L({v})_{w}, (\lambda s \in L({v}, {v}).\; {t}{s})) \end{array}}}
\vnbprop{vec-functional-type}{\vnbstmt{\begin{array}{@{}l@{}} \forall v, i, a. \\ \quad (\operatorname{is\text{-}complex\text{-}inner\text{-}product\text{-}space}(v) \wedge ((i \in \operatorname{vec}(v)) \wedge (a \in \operatorname{vec}(v)))) \Rightarrow  \\ \quad \quad \omega_{{i},{a}} \in (L({v}, {v}) \to \mathbb{C}) \end{array}}}
\vnbprop{vec-functional-apply}{\vnbstmt{\begin{array}{@{}l@{}} \forall v, i, a, t. \\ \quad (\operatorname{is\text{-}complex\text{-}inner\text{-}product\text{-}space}(v) \wedge ((i \in \operatorname{vec}(v)) \wedge ((a \in \operatorname{vec}(v)) \wedge (t \in L({v}, {v}))))) \Rightarrow  \\ \quad \quad \omega_{{i},{a}}(t) = \langle\cdot,\cdot\rangle_{v}(t(i), a) \end{array}}}
\vnbprop{vec-functional-linear}{\vnbstmt{\begin{array}{@{}l@{}} \forall v, i, a. \\ \quad (\operatorname{is\text{-}complex\text{-}inner\text{-}product\text{-}space}(v) \wedge ((i \in \operatorname{vec}(v)) \wedge (a \in \operatorname{vec}(v)))) \Rightarrow  \\ \quad \quad \operatorname{is\text{-}linear\text{-}map}(L({v}), \operatorname{cc\text{-}sn\text{-}module}, \omega_{{i},{a}}) \end{array}}}
\vnbprop{vec-functional-wot-continuous}{\vnbstmt{\begin{array}{@{}l@{}} \forall v, i, a. \\ \quad (\operatorname{is\text{-}complex\text{-}inner\text{-}product\text{-}space}(v) \wedge ((i \in \operatorname{vec}(v)) \wedge (a \in \operatorname{vec}(v)))) \Rightarrow  \\ \quad \quad \operatorname{is\text{-}hom\text{-}top\text{-}space}(L({v})_{w}, \operatorname{cc\text{-}sn\text{-}top}, \omega_{{i},{a}}) \end{array}}}
\vnbprop{vec-functional-sot-continuous}{\vnbstmt{\begin{array}{@{}l@{}} \forall v, i, a. \\ \quad (\operatorname{is\text{-}complex\text{-}inner\text{-}product\text{-}space}(v) \wedge ((i \in \operatorname{vec}(v)) \wedge (a \in \operatorname{vec}(v)))) \Rightarrow  \\ \quad \quad \operatorname{is\text{-}hom\text{-}top\text{-}space}(L({v})_{s}, \operatorname{cc\text{-}sn\text{-}top}, \omega_{{i},{a}}) \end{array}}}
\vnbprop{op-dual-unfold}{\vnbstmt{\begin{array}{@{}l@{}} \forall v, u. \\ \quad (\operatorname{op\text{-}dual}(v, u) \simeq \{f \in (L({v}, {v}) \to \mathbb{C}) : (\operatorname{is\text{-}linear\text{-}map}(L({v}), \operatorname{cc\text{-}sn\text{-}module}, f) \wedge \operatorname{is\text{-}hom\text{-}top\text{-}space}(u, \operatorname{cc\text{-}sn\text{-}top}, f))\}) \end{array}}}
\vnbprop{weak-dual-unfold}{\vnbstmt{\begin{array}{@{}l@{}} \forall v. \\ \quad (\operatorname{weak\text{-}dual}(v) \simeq \{f \in (L({v}, {v}) \to \mathbb{C}) : (\operatorname{is\text{-}linear\text{-}map}(L({v}), \operatorname{cc\text{-}sn\text{-}module}, f) \wedge \operatorname{is\text{-}hom\text{-}top\text{-}space}(L({v})_{w}, \operatorname{cc\text{-}sn\text{-}top}, f))\}) \end{array}}}
\vnbprop{strong-dual-unfold}{\vnbstmt{\forall v.\; (L_{\bullet}({v}) \simeq \operatorname{op\text{-}dual}(v, L({v})_{s}))}}
\vnbprop{vec-functional-in-weak-dual}{\vnbstmt{\begin{array}{@{}l@{}} \forall v, i, a. \\ \quad (\operatorname{is\text{-}complex\text{-}inner\text{-}product\text{-}space}(v) \wedge ((i \in \operatorname{vec}(v)) \wedge (a \in \operatorname{vec}(v)))) \Rightarrow  \\ \quad \quad \omega_{{i},{a}} \in \operatorname{weak\text{-}dual}(v) \end{array}}}
\vnbprop{vec-functional-in-strong-dual}{\vnbstmt{\begin{array}{@{}l@{}} \forall v, i, a. \\ \quad (\operatorname{is\text{-}complex\text{-}inner\text{-}product\text{-}space}(v) \wedge ((i \in \operatorname{vec}(v)) \wedge (a \in \operatorname{vec}(v)))) \Rightarrow  \\ \quad \quad \omega_{{i},{a}} \in L_{\bullet}({v}) \end{array}}}
\vnbprop{weak-dual-subset-strong-dual}{\vnbstmt{\begin{array}{@{}l@{}} \forall v, f. \\ \quad (\operatorname{is\text{-}complex\text{-}inner\text{-}product\text{-}space}(v) \wedge (f \in \operatorname{weak\text{-}dual}(v))) \Rightarrow  \\ \quad \quad f \in L_{\bullet}({v}) \end{array}}}

\vnbbookfile{theorem-library/hilbert-schmidt}{theorem-library/hilbert-schmidt}
\noindent{\small hilbert-schmidt.scm -- wave of 2026-11-09/10 (agent TH3-B), wired 2026-11-10 from scratchpad/th3-1109/b/th3b-hs.scm. Library: operator-algebras.}\par
\vnbprop{hs-sum-adjoint}{\vnbstmt{\begin{array}{@{}l@{}} \forall v, w.; \forall t \in L({v}, {w}), x, e, a, f. \\ \quad \operatorname{is\text{-}hilbert\text{-}space}(v) \wedge \\ \quad \operatorname{is\text{-}hilbert\text{-}space}(w) \wedge \\ \quad \operatorname{is\text{-}orthonormal\text{-}basis}(v, x, e) \wedge \\ \quad \operatorname{is\text{-}orthonormal\text{-}basis}(w, a, f) \Rightarrow \\ \quad \quad \operatorname{esum}((\lambda k \in x.\; (\lVert t(e(k)) \rVert \cdot \lVert t(e(k)) \rVert))) \\ \quad \quad =\; \operatorname{esum}((\lambda k \in a.\; (\lVert {t}^{*}(f(k)) \rVert \cdot \lVert {t}^{*}(f(k)) \rVert))) \end{array}}}
\vnbprop{hs-sum-basis-independent}{\vnbstmt{\begin{array}{@{}l@{}} \forall v, w.; \forall t \in L({v}, {w}), x, e, a, b. \\ \quad \operatorname{is\text{-}hilbert\text{-}space}(v) \wedge \\ \quad \operatorname{is\text{-}hilbert\text{-}space}(w) \wedge \\ \quad \operatorname{is\text{-}orthonormal\text{-}basis}(v, x, e) \wedge \\ \quad \operatorname{is\text{-}orthonormal\text{-}basis}(v, a, b) \Rightarrow \\ \quad \quad \sum_{i \in {x}} \lVert {t} {e}_i \rVert^2 = \sum_{i \in {a}} \lVert {t} {b}_i \rVert^2 \end{array}}}
\vnbprop{hs-sum-adjoint-eq}{\vnbstmt{\begin{array}{@{}l@{}} \forall v, w.; \forall t \in L({v}, {w}), x, e, a, f. \\ \quad \operatorname{is\text{-}hilbert\text{-}space}(v) \wedge \\ \quad \operatorname{is\text{-}hilbert\text{-}space}(w) \wedge \\ \quad \operatorname{is\text{-}orthonormal\text{-}basis}(v, x, e) \wedge \\ \quad \operatorname{is\text{-}orthonormal\text{-}basis}(w, a, f) \Rightarrow \\ \quad \quad \sum_{i \in {a}} \lVert {{t}^{*}} {f}_i \rVert^2 = \sum_{i \in {x}} \lVert {t} {e}_i \rVert^2 \end{array}}}
\vnbprop{hs-sum-real-any-basis}{\vnbstmt{\begin{array}{@{}l@{}} \forall v, w, t, x, e. \\ \quad \operatorname{is\text{-}hilbert\text{-}space}(v) \wedge \\ \quad \operatorname{is\text{-}hilbert\text{-}space}(w) \wedge \\ \quad \operatorname{is\text{-}hilbert\text{-}schmidt}(v, w, t) \wedge \\ \quad \operatorname{is\text{-}orthonormal\text{-}basis}(v, x, e) \Rightarrow \\ \quad \quad \sum_{i \in {x}} \lVert {t} {e}_i \rVert^2 \in \mathbb{R} \end{array}}}
\vnbprop{hs-adjoint}{\vnbstmt{\begin{array}{@{}l@{}} \forall v, w, t. \\ \quad (\operatorname{is\text{-}hilbert\text{-}space}(v) \wedge (\operatorname{is\text{-}hilbert\text{-}space}(w) \wedge \operatorname{is\text{-}hilbert\text{-}schmidt}(v, w, t))) \Rightarrow  \\ \quad \quad \operatorname{is\text{-}hilbert\text{-}schmidt}(w, v, {t}^{*}) \end{array}}}
\vnbprop{canon-onb-is-onb}{\vnbstmt{\begin{array}{@{}l@{}} \forall v. \\ \quad \operatorname{is\text{-}complex\text{-}inner\text{-}product\text{-}space}(v) \Rightarrow  \\ \quad \quad \operatorname{is\text{-}orthonormal\text{-}basis}(v, \operatorname{canon\text{-}onb}(v), (\lambda e \in \operatorname{canon\text{-}onb}(v).\; e)) \end{array}}}
\vnbprop{hs-norm-any-basis}{\vnbstmt{\begin{array}{@{}l@{}} \forall v, w.; \forall t \in L({v}, {w}), x, e. \\ \quad (\operatorname{is\text{-}hilbert\text{-}space}(v) \wedge (\operatorname{is\text{-}hilbert\text{-}space}(w) \wedge \operatorname{is\text{-}orthonormal\text{-}basis}(v, x, e))) \Rightarrow  \\ \quad \quad (\lVert {t} \rVert_{HS} \simeq \sqrt{\sum_{i \in {x}} \lVert {t} {e}_i \rVert^2}) \end{array}}}
\vnbprop{hs-sum-in-rps}{\vnbstmt{\begin{array}{@{}l@{}} \forall v, w.; \forall t \in L({v}, {w}), x, e. \\ \quad (\operatorname{is\text{-}hilbert\text{-}space}(v) \wedge (\operatorname{is\text{-}hilbert\text{-}space}(w) \wedge \operatorname{is\text{-}orthonormal\text{-}basis}(v, x, e))) \Rightarrow  \\ \quad \quad \sum_{i \in {x}} \lVert {t} {e}_i \rVert^2 \in \operatorname{rr\text{-}pos\text{-}star} \end{array}}}
\vnbprop{hs-sum-comp-left}{\vnbstmt{\begin{array}{@{}l@{}} \forall v, s, t, x, e. \\ \quad \operatorname{is\text{-}hilbert\text{-}space}(v) \wedge \\ \quad s \in L({v}, {v}) \wedge \\ \quad t \in L({v}, {v}) \wedge \\ \quad \operatorname{is\text{-}orthonormal\text{-}basis}(v, x, e) \Rightarrow \\ \quad \quad \sum_{i \in {x}} \lVert {{s}{t}} {e}_i \rVert^2 \\ \quad \quad \leq\; \operatorname{etimes}((\lVert {s} \rVert \cdot \lVert {s} \rVert), \sum_{i \in {x}} \lVert {t} {e}_i \rVert^2) \end{array}}}
\vnbprop{hs-comp-left}{\vnbstmt{\begin{array}{@{}l@{}} \forall v, s, t. \\ \quad (\operatorname{is\text{-}hilbert\text{-}space}(v) \wedge ((s \in L({v}, {v})) \wedge \operatorname{is\text{-}hilbert\text{-}schmidt}(v, v, t))) \Rightarrow  \\ \quad \quad \operatorname{is\text{-}hilbert\text{-}schmidt}(v, v, {s}{t}) \end{array}}}
\vnbprop{hs-comp-right}{\vnbstmt{\begin{array}{@{}l@{}} \forall v, s, t. \\ \quad (\operatorname{is\text{-}hilbert\text{-}space}(v) \wedge ((s \in L({v}, {v})) \wedge \operatorname{is\text{-}hilbert\text{-}schmidt}(v, v, t))) \Rightarrow  \\ \quad \quad \operatorname{is\text{-}hilbert\text{-}schmidt}(v, v, {t}{s}) \end{array}}}
\vnbprop{hs-vector-bound}{\vnbstmt{\begin{array}{@{}l@{}} \forall v, w.; \forall t \in L({v}, {w}), x, e, a, f.; \forall \mathit{tbx} \in \operatorname{vec}(v). \\ \quad \operatorname{is\text{-}hilbert\text{-}space}(v) \wedge \\ \quad \operatorname{is\text{-}hilbert\text{-}space}(w) \wedge \\ \quad \operatorname{is\text{-}orthonormal\text{-}basis}(v, x, e) \wedge \\ \quad \operatorname{is\text{-}orthonormal\text{-}basis}(w, a, f) \Rightarrow \\ \quad \quad \lVert t(\mathit{tbx}) \rVert \cdot \lVert t(\mathit{tbx}) \rVert \\ \quad \quad \leq\; \operatorname{etimes}((\lVert \mathit{tbx} \rVert \cdot \lVert \mathit{tbx} \rVert), \sum_{i \in {x}} \lVert {t} {e}_i \rVert^2) \end{array}}}
\vnbprop{hs-opnorm-le}{\vnbstmt{\begin{array}{@{}l@{}} \forall v, w, t. \\ \quad (\operatorname{is\text{-}hilbert\text{-}space}(v) \wedge (\operatorname{is\text{-}hilbert\text{-}space}(w) \wedge \operatorname{is\text{-}hilbert\text{-}schmidt}(v, w, t))) \Rightarrow  \\ \quad \quad \lVert {t} \rVert \leq \lVert {t} \rVert_{HS} \end{array}}}
\vnbprop{trace-abs-summable}{\vnbstmt{\begin{array}{@{}l@{}} \forall v, t, x, e. \\ \quad (\operatorname{is\text{-}hilbert\text{-}space}(v) \wedge (\operatorname{is\text{-}trace\text{-}class}(v, t) \wedge \operatorname{is\text{-}orthonormal\text{-}basis}(v, x, e))) \Rightarrow  \\ \quad \quad \operatorname{esum}((\lambda k \in x.\; \operatorname{magnitude}(\langle\cdot,\cdot\rangle_{v}(t(e(k)), e(k))))) \in \mathbb{R} \end{array}}}
\vnbprop{hs-sum-add-le}{\vnbstmt{\begin{array}{@{}l@{}} \forall v, w, s, t, x, e. \\ \quad \operatorname{is\text{-}hilbert\text{-}space}(v) \wedge \\ \quad \operatorname{is\text{-}hilbert\text{-}space}(w) \wedge \\ \quad s \in L({v}, {w}) \wedge \\ \quad t \in L({v}, {w}) \wedge \\ \quad \operatorname{is\text{-}orthonormal\text{-}basis}(v, x, e) \Rightarrow \\ \quad \quad \sum_{i \in {x}} \lVert {{s} + {t}} {e}_i \rVert^2 \\ \quad \quad \leq\; \operatorname{eplus}(\operatorname{etimes}(2, \sum_{i \in {x}} \lVert {s} {e}_i \rVert^2), \operatorname{etimes}(2, \sum_{i \in {x}} \lVert {t} {e}_i \rVert^2)) \end{array}}}
\vnbprop{hs-add}{\vnbstmt{\begin{array}{@{}l@{}} \forall v, w, s, t. \\ \quad \operatorname{is\text{-}hilbert\text{-}space}(v) \wedge \\ \quad \operatorname{is\text{-}hilbert\text{-}space}(w) \wedge \\ \quad \operatorname{is\text{-}hilbert\text{-}schmidt}(v, w, s) \wedge \\ \quad \operatorname{is\text{-}hilbert\text{-}schmidt}(v, w, t) \Rightarrow \\ \quad \quad \operatorname{is\text{-}hilbert\text{-}schmidt}(v, w, {s} + {t}) \end{array}}}
\vnbprop{hs-scale}{\vnbstmt{\begin{array}{@{}l@{}} \forall v, w, c, t. \\ \quad (\operatorname{is\text{-}hilbert\text{-}space}(v) \wedge (\operatorname{is\text{-}hilbert\text{-}space}(w) \wedge ((c \in \mathbb{C}) \wedge \operatorname{is\text{-}hilbert\text{-}schmidt}(v, w, t)))) \Rightarrow  \\ \quad \quad \operatorname{is\text{-}hilbert\text{-}schmidt}(v, w, {c}\,{t}) \end{array}}}

\vnbbookfile{theorem-library/th3a4-spectral}{theorem-library/th3a4-spectral}
\noindent{\small th3a4-spectral.scm -- wave of 2026-11-09/10 (agent TH3-A4), wired 2026-11-10 from scratchpad/th3-1109/a4/th3a4-spectral.scm. Library: operator-algebras.}\par
\vnbprop{sa-eigen-span-dense}{\vnbstmt{\begin{array}{@{}l@{}} \forall v, t, k. \\ \quad \operatorname{is\text{-}hilbert\text{-}space}(v) \wedge \\ \quad \operatorname{is\text{-}compact\text{-}op}(v, v, t) \wedge \\ \quad \operatorname{is\text{-}self\text{-}adjoint\text{-}op}(v, t) \wedge \\ \quad \operatorname{is\text{-}submodule}(v, k) \wedge \\ \quad \{z \in \operatorname{vec}(v) : \exists c \in \mathbb{C}.\; (z \in N_{c}({t}))\} \subseteq k \Rightarrow \\ \quad \quad \operatorname{closure}(\operatorname{ip\text{-}metric\text{-}space}(v), k) = \operatorname{vec}(v) \end{array}}}
\vnbprop{sa-eigen-perp-perp}{\vnbstmt{\begin{array}{@{}l@{}} \forall v, t, k. \\ \quad \operatorname{is\text{-}hilbert\text{-}space}(v) \wedge \\ \quad \operatorname{is\text{-}compact\text{-}op}(v, v, t) \wedge \\ \quad \operatorname{is\text{-}self\text{-}adjoint\text{-}op}(v, t) \wedge \\ \quad \operatorname{is\text{-}submodule}(v, k) \wedge \\ \quad \{z \in \operatorname{vec}(v) : \exists c \in \mathbb{C}.\; (z \in N_{c}({t}))\} \subseteq k \Rightarrow \\ \quad \quad {{k}^{\perp}}^{\perp} = \operatorname{vec}(v) \end{array}}}
\vnbprop{sa-eigen-closed-span}{\vnbstmt{\begin{array}{@{}l@{}} \forall v, t, k. \\ \quad \operatorname{is\text{-}hilbert\text{-}space}(v) \wedge \\ \quad \operatorname{is\text{-}compact\text{-}op}(v, v, t) \wedge \\ \quad \operatorname{is\text{-}self\text{-}adjoint\text{-}op}(v, t) \wedge \\ \quad \operatorname{is\text{-}submodule}(v, k) \wedge \\ \quad \{z \in \operatorname{vec}(v) : \exists c \in \mathbb{C}.\; (z \in N_{c}({t}))\} \subseteq k \wedge \\ \quad \operatorname{is\text{-}closed}(\operatorname{ip\text{-}metric\text{-}space}(v), k) \Rightarrow \\ \quad \quad k = \operatorname{vec}(v) \end{array}}}
\vnbprop{eigenspace-submodule}{\vnbstmt{\begin{array}{@{}l@{}} \forall v, t, c. \\ \quad (\operatorname{is\text{-}complex\text{-}inner\text{-}product\text{-}space}(v) \wedge ((t \in L({v}, {v})) \wedge (c \in \mathbb{C}))) \Rightarrow  \\ \quad \quad \operatorname{is\text{-}submodule}(v, N_{c}({t})) \end{array}}}
\vnbprop{eigen-canon-onb}{\vnbstmt{\begin{array}{@{}l@{}} \forall v, t, c. \\ \quad (\operatorname{is\text{-}complex\text{-}inner\text{-}product\text{-}space}(v) \wedge ((t \in L({v}, {v})) \wedge (c \in \mathbb{C}))) \Rightarrow  \\ \quad \quad \operatorname{canon\text{-}onb}(\operatorname{sub\text{-}cips}(v, N_{c}({t}))) \subseteq N_{c}({t}) \wedge \\ \quad \quad \operatorname{is\text{-}orthonormal\text{-}set}(v, \operatorname{canon\text{-}onb}(\operatorname{sub\text{-}cips}(v, N_{c}({t})))) \wedge \\ \quad \quad \forall y \in N_{c}({t}). \\ \quad \quad \quad \forall u \in \operatorname{canon\text{-}onb}(\operatorname{sub\text{-}cips}(v, N_{c}({t}))).\; (\langle\cdot,\cdot\rangle_{v}(y, u) = 0) \Rightarrow  \\ \quad \quad \quad \quad y = \operatorname{vzero}(v) \end{array}}}
\vnbprop{sa-eigen-onb}{\vnbstmt{\begin{array}{@{}l@{}} \forall v, t. \\ \quad (\operatorname{is\text{-}hilbert\text{-}space}(v) \wedge (\operatorname{is\text{-}compact\text{-}op}(v, v, t) \wedge \operatorname{is\text{-}self\text{-}adjoint\text{-}op}(v, t))) \Rightarrow  \\ \quad \quad \exists x, e. \\ \quad \quad \quad (\operatorname{is\text{-}orthonormal\text{-}basis}(v, x, e) \wedge \forall i \in x.\; (t(e(i)) = \operatorname{act}(v)(\langle\cdot,\cdot\rangle_{v}(t(e(i)), e(i)), e(i)))) \end{array}}}
\vnbprop{bop-hom-ag}{\vnbstmt{\begin{array}{@{}l@{}} \forall v, w, t. \\ \quad (\operatorname{is\text{-}complex\text{-}inner\text{-}product\text{-}space}(v) \wedge (\operatorname{is\text{-}complex\text{-}inner\text{-}product\text{-}space}(w) \wedge (t \in L({v}, {w})))) \Rightarrow  \\ \quad \quad \operatorname{is\text{-}hom\text{-}abelian\text{-}group}(\operatorname{normed\text{-}ag\text{-}as\text{-}abelian\text{-}group}(\operatorname{ip\text{-}normed\text{-}ag}(v)), \\ \quad \quad \quad \operatorname{normed\text{-}ag\text{-}as\text{-}abelian\text{-}group}(\operatorname{ip\text{-}normed\text{-}ag}(w)), \\ \quad \quad \quad t) \end{array}}}
\vnbprop{bop-sums-to}{\vnbstmt{\begin{array}{@{}l@{}} \forall v, w, t, x, m, r. \\ \quad \operatorname{is\text{-}complex\text{-}inner\text{-}product\text{-}space}(v) \wedge \\ \quad \operatorname{is\text{-}complex\text{-}inner\text{-}product\text{-}space}(w) \wedge \\ \quad t \in L({v}, {w}) \wedge \\ \quad x \in \mathbf{Set} \wedge \\ \quad m \in (x \to \operatorname{vec}(v)) \wedge \\ \quad r \in \operatorname{vec}(v) \wedge \\ \quad \operatorname{sums\text{-}to}(\operatorname{ip\text{-}normed\text{-}ag}(v), m, r) \Rightarrow \\ \quad \quad \operatorname{sums\text{-}to}(\operatorname{ip\text{-}normed\text{-}ag}(w), (\lambda k \in x.\; t(m(k))), t(r)) \end{array}}}
\vnbprop{op-eigen-expansion}{\vnbstmt{\begin{array}{@{}l@{}} \forall v, t, a, e, \ell.; \forall x \in \operatorname{vec}(v). \\ \quad \operatorname{is\text{-}hilbert\text{-}space}(v) \wedge \\ \quad t \in L({v}, {v}) \wedge \\ \quad \operatorname{is\text{-}orthonormal\text{-}basis}(v, a, e) \wedge \\ \quad \ell \in (a \to \mathbb{C}) \wedge \\ \quad \forall i \in a.\; (t(e(i)) = \operatorname{act}(v)(\ell(i), e(i))) \Rightarrow \\ \quad \quad \operatorname{sums\text{-}to}(\operatorname{ip\text{-}normed\text{-}ag}(v), (\lambda k \in a.\; \operatorname{act}(v)((\ell(k) \cdot \langle\cdot,\cdot\rangle_{v}(x, e(k))), e(k))), t(x)) \end{array}}}
\vnbprop{ord-nn-in-segment-of-infinite}{\vnbstmt{\forall l \in \mathrm{Ord}, n \in \mathbb{N}.\; (\neg (l \in \mathbb{N}) \Rightarrow (n \in \mathrm{OS}(l)))}}
\vnbprop{infinite-set-nn-injection}{\vnbstmt{\begin{array}{@{}l@{}} \forall s \in \mathbf{Set}. \\ \quad \neg (|s| \in \mathbb{N}) \Rightarrow  \\ \quad \quad \exists g \in (\mathbb{N} \to s). \forall a, b \in \mathbb{N}. \\ \quad \quad \quad (g(a) = g(b)) \Rightarrow (a = b) \end{array}}}
\vnbprop{ip-unit-norm-one}{\vnbstmt{\begin{array}{@{}l@{}} \forall v, x. \\ \quad (\operatorname{is\text{-}complex\text{-}inner\text{-}product\text{-}space}(v) \wedge ((x \in \operatorname{vec}(v)) \wedge (\langle\cdot,\cdot\rangle_{v}(x, x) = 1))) \Rightarrow  \\ \quad \quad \lVert x \rVert = 1 \end{array}}}
\vnbprop{compact-eigen-large-finite}{\vnbstmt{\begin{array}{@{}l@{}} \forall v, t, x, e, \ell, d. \\ \quad \operatorname{is\text{-}complex\text{-}inner\text{-}product\text{-}space}(v) \wedge \\ \quad \operatorname{is\text{-}compact\text{-}op}(v, v, t) \wedge \\ \quad \operatorname{is\text{-}orthonormal\text{-}family}(v, x, e) \wedge \\ \quad \ell \in (x \to \mathbb{C}) \wedge \\ \quad \forall i \in x.\; (t(e(i)) = \operatorname{act}(v)(\ell(i), e(i))) \wedge \\ \quad \operatorname{pos\text{-}rr}(d) \Rightarrow \\ \quad \quad |\{z \in x : (d \leq \operatorname{magnitude}(\ell(z)))\}| \in \mathbb{N} \end{array}}}
\vnbprop{compact-eigen-nonzero-countable}{\vnbstmt{\begin{array}{@{}l@{}} \forall v, t, x, e, \ell. \\ \quad \operatorname{is\text{-}complex\text{-}inner\text{-}product\text{-}space}(v) \wedge \\ \quad \operatorname{is\text{-}compact\text{-}op}(v, v, t) \wedge \\ \quad \operatorname{is\text{-}orthonormal\text{-}family}(v, x, e) \wedge \\ \quad \ell \in (x \to \mathbb{C}) \wedge \\ \quad \forall i \in x.\; (t(e(i)) = \operatorname{act}(v)(\ell(i), e(i))) \Rightarrow \\ \quad \quad \operatorname{is\text{-}countable}(\{z \in x : \neg (\ell(z) = 0)\}) \end{array}}}
\vnbprop{sa-eigen-expansion}{\vnbstmt{\begin{array}{@{}l@{}} \forall v, t. \\ \quad (\operatorname{is\text{-}hilbert\text{-}space}(v) \wedge (\operatorname{is\text{-}compact\text{-}op}(v, v, t) \wedge \operatorname{is\text{-}self\text{-}adjoint\text{-}op}(v, t))) \Rightarrow  \\ \quad \quad \exists \mathit{tfix}, e, m. \\ \quad \quad \quad \operatorname{is\text{-}orthonormal\text{-}basis}(v, \mathit{tfix}, e) \wedge \\ \quad \quad \quad m \in (\mathit{tfix} \to \mathbb{C}) \wedge \\ \quad \quad \quad \forall i \in \mathit{tfix}.\; ((m(i) = \operatorname{conjugate}(m(i))) \wedge (t(e(i)) = \operatorname{act}(v)(m(i), e(i)))) \wedge \\ \quad \quad \quad \forall x \in \operatorname{vec}(v). \\ \quad \quad \quad \quad \operatorname{sums\text{-}to}(\operatorname{ip\text{-}normed\text{-}ag}(v), (\lambda k \in \mathit{tfix}.\; \operatorname{act}(v)((m(k) \cdot \langle\cdot,\cdot\rangle_{v}(x, e(k))), e(k))), t(x)) \wedge \\ \quad \quad \quad \forall d.\; (\operatorname{pos\text{-}rr}(d) \Rightarrow (|\{z \in \mathit{tfix} : (d \leq \operatorname{magnitude}(m(z)))\}| \in \mathbb{N})) \wedge \\ \quad \quad \quad \operatorname{is\text{-}countable}(\{z \in \mathit{tfix} : \neg (m(z) = 0)\}) \end{array}}}
\vnbprop{eigenspace-onb-finite}{\vnbstmt{\begin{array}{@{}l@{}} \forall v, t, c. \\ \quad (\operatorname{is\text{-}complex\text{-}inner\text{-}product\text{-}space}(v) \wedge (\operatorname{is\text{-}compact\text{-}op}(v, v, t) \wedge ((c \in \mathbb{C}) \wedge \neg (c = 0)))) \Rightarrow  \\ \quad \quad |\operatorname{canon\text{-}onb}(\operatorname{sub\text{-}cips}(v, N_{c}({t})))| \in \mathbb{N} \end{array}}}

\vnbbookfile{theorem-library/hs-compact}{theorem-library/hs-compact}
\noindent{\small hs-compact.scm -- wave of 2026-11-09/10 (agent TH3-B3), wired 2026-11-10 from scratchpad/th3-1109/b3/tb3-trunc.scm. Library: operator-algebras.}\par
\vnbprop{cips-vadd-vzero-right}{\vnbstmt{\begin{array}{@{}l@{}} \forall v, x. \\ \quad (\operatorname{is\text{-}complex\text{-}inner\text{-}product\text{-}space}(v) \wedge (x \in \operatorname{vec}(v))) \Rightarrow  \\ \quad \quad \operatorname{vadd}(v)(x, \operatorname{vzero}(v)) = x \end{array}}}
\vnbprop{cips-vneg-vzero}{\vnbstmt{\begin{array}{@{}l@{}} \forall v. \\ \quad \operatorname{is\text{-}complex\text{-}inner\text{-}product\text{-}space}(v) \Rightarrow (\operatorname{vneg}(v)(\operatorname{vzero}(v)) = \operatorname{vzero}(v)) \end{array}}}
\vnbprop{cips-act-zero-scalar}{\vnbstmt{\begin{array}{@{}l@{}} \forall v, x. \\ \quad (\operatorname{is\text{-}complex\text{-}inner\text{-}product\text{-}space}(v) \wedge (x \in \operatorname{vec}(v))) \Rightarrow  \\ \quad \quad \operatorname{act}(v)(0, x) = \operatorname{vzero}(v) \end{array}}}
\vnbprop{cips-act-one}{\vnbstmt{\begin{array}{@{}l@{}} \forall v, x. \\ \quad (\operatorname{is\text{-}complex\text{-}inner\text{-}product\text{-}space}(v) \wedge (x \in \operatorname{vec}(v))) \Rightarrow (\operatorname{act}(v)(1, x) = x) \end{array}}}
\vnbprop{trunc-op-compact}{\vnbstmt{\begin{array}{@{}l@{}} \forall v, w, x, e, t, q. \\ \quad \operatorname{is\text{-}complex\text{-}inner\text{-}product\text{-}space}(v) \wedge \\ \quad \operatorname{is\text{-}complex\text{-}inner\text{-}product\text{-}space}(w) \wedge \\ \quad e \in (x \to \operatorname{vec}(v)) \wedge \\ \quad t \in L({v}, {w}) \wedge \\ \quad q \in \mathbf{Set} \wedge \\ \quad |q| \in \mathbb{N} \wedge \\ \quad q \subseteq x \Rightarrow \\ \quad \quad \operatorname{finsum}(\operatorname{normed\text{-}ag\text{-}as\text{-}abelian\text{-}group}(\operatorname{bop\text{-}normed\text{-}ag}(v, w)), \\ \quad \quad \quad (\lambda k \in q.\; \langle \cdot \mid {e(k)} \rangle {t(e(k))}), q) \\ \quad \quad \in\; K({v}, {w}) \end{array}}}
\vnbprop{trunc-op-bounded}{\vnbstmt{\begin{array}{@{}l@{}} \forall v, w, x, e, t, q. \\ \quad \operatorname{is\text{-}complex\text{-}inner\text{-}product\text{-}space}(v) \wedge \\ \quad \operatorname{is\text{-}complex\text{-}inner\text{-}product\text{-}space}(w) \wedge \\ \quad e \in (x \to \operatorname{vec}(v)) \wedge \\ \quad t \in L({v}, {w}) \wedge \\ \quad q \in \mathbf{Set} \wedge \\ \quad |q| \in \mathbb{N} \wedge \\ \quad q \subseteq x \Rightarrow \\ \quad \quad \operatorname{finsum}(\operatorname{normed\text{-}ag\text{-}as\text{-}abelian\text{-}group}(\operatorname{bop\text{-}normed\text{-}ag}(v, w)), \\ \quad \quad \quad (\lambda k \in q.\; \langle \cdot \mid {e(k)} \rangle {t(e(k))}), q) \\ \quad \quad \in\; L({v}, {w}) \end{array}}}
\vnbprop{trunc-op-basis-in}{\vnbstmt{\begin{array}{@{}l@{}} \forall v, w, x, e, t, q, z. \\ \quad \operatorname{is\text{-}complex\text{-}inner\text{-}product\text{-}space}(v) \wedge \\ \quad \operatorname{is\text{-}complex\text{-}inner\text{-}product\text{-}space}(w) \wedge \\ \quad \operatorname{is\text{-}orthonormal\text{-}family}(v, x, e) \wedge \\ \quad t \in L({v}, {w}) \wedge \\ \quad q \in \mathbf{Set} \wedge \\ \quad |q| \in \mathbb{N} \wedge \\ \quad q \subseteq x \wedge \\ \quad z \in q \Rightarrow \\ \quad \quad \operatorname{finsum}(\operatorname{normed\text{-}ag\text{-}as\text{-}abelian\text{-}group}(\operatorname{bop\text{-}normed\text{-}ag}(v, w)), (\lambda k \in q.\; \langle \cdot \mid {e(k)} \rangle {t(e(k))}), q)(e(z)) \\ \quad \quad =\; t(e(z)) \end{array}}}
\vnbprop{trunc-op-basis-out}{\vnbstmt{\begin{array}{@{}l@{}} \forall v, w, x, e, t, q, z. \\ \quad \operatorname{is\text{-}complex\text{-}inner\text{-}product\text{-}space}(v) \wedge \\ \quad \operatorname{is\text{-}complex\text{-}inner\text{-}product\text{-}space}(w) \wedge \\ \quad \operatorname{is\text{-}orthonormal\text{-}family}(v, x, e) \wedge \\ \quad t \in L({v}, {w}) \wedge \\ \quad q \in \mathbf{Set} \wedge \\ \quad |q| \in \mathbb{N} \wedge \\ \quad q \subseteq x \wedge \\ \quad z \in x \wedge \\ \quad \neg (z \in q) \Rightarrow \\ \quad \quad \operatorname{finsum}(\operatorname{normed\text{-}ag\text{-}as\text{-}abelian\text{-}group}(\operatorname{bop\text{-}normed\text{-}ag}(v, w)), (\lambda k \in q.\; \langle \cdot \mid {e(k)} \rangle {t(e(k))}), q)(e(z)) \\ \quad \quad =\; \operatorname{vzero}(w) \end{array}}}
\vnbprop{hs-sum-trunc-tail}{\vnbstmt{\begin{array}{@{}l@{}} \forall v, w, x, e, t, q. \\ \quad \operatorname{is\text{-}complex\text{-}inner\text{-}product\text{-}space}(v) \wedge \\ \quad \operatorname{is\text{-}complex\text{-}inner\text{-}product\text{-}space}(w) \wedge \\ \quad \operatorname{is\text{-}orthonormal\text{-}family}(v, x, e) \wedge \\ \quad t \in L({v}, {w}) \wedge \\ \quad q \in \mathbf{Set} \wedge \\ \quad |q| \in \mathbb{N} \wedge \\ \quad q \subseteq x \Rightarrow \\ \quad \quad \sum_{i \in {x}} \lVert {{t} + {-{\operatorname{finsum}(\operatorname{normed\text{-}ag\text{-}as\text{-}abelian\text{-}group}(\operatorname{bop\text{-}normed\text{-}ag}(v, w)), (\lambda k \in q.\; \langle \cdot \mid {e(k)} \rangle {t(e(k))}), q)}}} {e}_i \rVert^2 \\ \quad \quad =\; \operatorname{esum}((\lambda z \in x.\; \\ \quad \quad \quad \quad \left(\text{if }(z \in q)\text{ then }0\text{ else }(\lambda \mathit{kqk} \in x.\; (\lVert t(e(\mathit{kqk})) \rVert \cdot \lVert t(e(\mathit{kqk})) \rVert))(z)\right))) \end{array}}}
\vnbprop{hs-trunc-approx}{\vnbstmt{\begin{array}{@{}l@{}} \forall v, w, t, x, e, s. \\ \quad \operatorname{is\text{-}hilbert\text{-}space}(v) \wedge \\ \quad \operatorname{is\text{-}hilbert\text{-}space}(w) \wedge \\ \quad t \in L({v}, {w}) \wedge \\ \quad \operatorname{is\text{-}orthonormal\text{-}basis}(v, x, e) \wedge \\ \quad \sum_{i \in {x}} \lVert {t} {e}_i \rVert^2 \in \mathbb{R} \wedge \\ \quad \operatorname{pos\text{-}rr}(s) \Rightarrow \\ \quad \quad \exists q \in \mathbf{Set}. \\ \quad \quad \quad |q| \in \mathbb{N} \wedge \\ \quad \quad \quad q \subseteq x \wedge \\ \quad \quad \quad \forall \mathit{xq} \in \operatorname{vec}(v). \\ \quad \quad \quad \quad (\lVert \mathit{xq} \rVert \leq 1) \Rightarrow  \\ \quad \quad \quad \quad \quad \lVert \operatorname{vadd}(w)(t(\mathit{xq}), \operatorname{vneg}(w)(\operatorname{finsum}(\operatorname{normed\text{-}ag\text{-}as\text{-}abelian\text{-}group}(\operatorname{bop\text{-}normed\text{-}ag}(v, w)), (\lambda k \in q.\; \langle \cdot \mid {e(k)} \rangle {t(e(k))}), q)(\mathit{xq}))) \rVert \\ \quad \quad \quad \quad \quad \leq\; s \end{array}}}
\vnbprop{hs-op-compact}{\vnbstmt{\begin{array}{@{}l@{}} \forall v, w, t. \\ \quad (\operatorname{is\text{-}hilbert\text{-}space}(v) \wedge (\operatorname{is\text{-}hilbert\text{-}space}(w) \wedge \operatorname{is\text{-}hilbert\text{-}schmidt}(v, w, t))) \Rightarrow  \\ \quad \quad \operatorname{is\text{-}compact\text{-}op}(v, w, t) \end{array}}}
\vnbprop{hs-sum-product}{\vnbstmt{\begin{array}{@{}l@{}} \forall v, w, t, x, e, a, f. \\ \quad \operatorname{is\text{-}hilbert\text{-}space}(v) \wedge \\ \quad \operatorname{is\text{-}hilbert\text{-}space}(w) \wedge \\ \quad t \in L({v}, {w}) \wedge \\ \quad \operatorname{is\text{-}orthonormal\text{-}basis}(v, x, e) \wedge \\ \quad \operatorname{is\text{-}orthonormal\text{-}basis}(w, a, f) \Rightarrow \\ \quad \quad \sum_{i \in {x}} \lVert {t} {e}_i \rVert^2 \\ \quad \quad =\; \operatorname{esum}((\lambda \mathit{tbw} \in (x \times a).\; \\ \quad \quad \quad \quad \operatorname{magnitude}(\langle\cdot,\cdot\rangle_{w}(t(e(\operatorname{nth}(1, \mathit{tbw}))), f(\operatorname{nth}(2, \mathit{tbw})))) \cdot \operatorname{magnitude}(\langle\cdot,\cdot\rangle_{w}(t(e(\operatorname{nth}(1, \mathit{tbw}))), f(\operatorname{nth}(2, \mathit{tbw})))))) \end{array}}}

\vnbbookfile{theorem-library/th3-eigen-normalized}{theorem-library/th3-eigen-normalized}
\noindent{\small th3-eigen-normalized.scm -- wave of 2026-11-09/10 (agent TH3-SORT), wired 2026-11-10 from scratchpad/th3-1109/sort/th3-eigen-normalized.scm. Library: operator-algebras.}\par
\vnbprop{sa-eigen-normalized}{\vnbstmt{\begin{array}{@{}l@{}} \forall v, t. \\ \quad (\operatorname{is\text{-}hilbert\text{-}space}(v) \wedge (\operatorname{is\text{-}compact\text{-}op}(v, v, t) \wedge \operatorname{is\text{-}self\text{-}adjoint\text{-}op}(v, t))) \Rightarrow  \\ \quad \quad \exists \mathit{tnix}, e, m. \\ \quad \quad \quad \operatorname{is\text{-}orthonormal\text{-}basis}(v, \mathit{tnix}, e) \wedge \\ \quad \quad \quad m \in (\mathit{tnix} \to \mathbb{C}) \wedge \\ \quad \quad \quad \forall i \in \mathit{tnix}.\; ((m(i) = \operatorname{conjugate}(m(i))) \wedge (t(e(i)) = \operatorname{act}(v)(m(i), e(i)))) \wedge \\ \quad \quad \quad \forall x \in \operatorname{vec}(v). \\ \quad \quad \quad \quad \operatorname{sums\text{-}to}(\operatorname{ip\text{-}normed\text{-}ag}(v), (\lambda k \in \mathit{tnix}.\; \operatorname{act}(v)((m(k) \cdot \langle\cdot,\cdot\rangle_{v}(x, e(k))), e(k))), t(x)) \wedge \\ \quad \quad \quad |\{i \in \mathit{tnix} : (0 < \operatorname{magnitude}(m(i)))\}| \in \mathbb{N} \wedge \\ \quad \quad \quad \exists h \in \mathrm{Bij}(\mathrm{OS}(|\{i \in \mathit{tnix} : (0 < \operatorname{magnitude}(m(i)))\}|), \{i \in \mathit{tnix} : (0 < \operatorname{magnitude}(m(i)))\}). \forall n \in \mathbb{N}. \\ \quad \quad \quad \quad ((n + 1) < |\{i \in \mathit{tnix} : (0 < \operatorname{magnitude}(m(i)))\}|) \Rightarrow  \\ \quad \quad \quad \quad \quad \operatorname{magnitude}(m(h((n + 1)))) \leq \operatorname{magnitude}(m(h(n))) \vee \\ \quad \quad \quad \neg (|\{i \in \mathit{tnix} : (0 < \operatorname{magnitude}(m(i)))\}| \in \mathbb{N}) \wedge \\ \quad \quad \quad \exists h \in \mathrm{Bij}(\mathbb{N}, \{i \in \mathit{tnix} : (0 < \operatorname{magnitude}(m(i)))\}). \\ \quad \quad \quad \quad \forall n \in \mathbb{N}.\; (\operatorname{magnitude}(m(h((n + 1)))) \leq \operatorname{magnitude}(m(h(n)))) \wedge \\ \quad \quad \quad \quad \forall d. \\ \quad \quad \quad \quad \quad \operatorname{pos\text{-}rr}(d) \Rightarrow  \\ \quad \quad \quad \quad \quad \quad \exists j \in \mathbb{N}.\; \forall n \in \mathbb{N}.\; ((j \leq n) \Rightarrow (\operatorname{magnitude}(m(h(n))) < d)) \end{array}}}

\vnbbookfile{theorem-library/basw-top}{theorem-library/basw-top}
\noindent{\small basw-top.scm -- wave of 2026-11-09/10 (agent BA-SW), wired 2026-11-10 from scratchpad/basw-1109/basw-top.scm. Library: operator-algebras.}\par
\vnbprop{bsw-cont-local}{\vnbstmt{\begin{array}{@{}l@{}} \forall t, d, f.; \forall p \in \operatorname{pts}(t), s. \\ \quad (\operatorname{is\text{-}metric\text{-}space}(d) \wedge (\operatorname{is\text{-}hom\text{-}top\text{-}space}(t, \operatorname{metric\text{-}top}(d), f) \wedge \operatorname{pos\text{-}rr}(s))) \Rightarrow  \\ \quad \quad \exists u \in \operatorname{opens}(t).\; ((p \in u) \wedge \forall q \in u.\; (\operatorname{dist}(d)(f(p), f(q)) < s)) \end{array}}}
\vnbprop{bsw-local-cont}{\vnbstmt{\begin{array}{@{}l@{}} \forall t, d, f. \\ \quad \operatorname{is\text{-}top\text{-}space}(t) \wedge \\ \quad \operatorname{is\text{-}metric\text{-}space}(d) \wedge \\ \quad f \in (\operatorname{pts}(t) \to \operatorname{pts}(d)) \wedge \\ \quad \forall p \in \operatorname{pts}(t), s. \\ \quad \quad \operatorname{pos\text{-}rr}(s) \Rightarrow  \\ \quad \quad \quad \exists u \in \operatorname{opens}(t).\; ((p \in u) \wedge \forall q \in u.\; (\operatorname{dist}(d)(f(p), f(q)) < s)) \Rightarrow \\ \quad \quad \operatorname{is\text{-}hom\text{-}top\text{-}space}(t, \operatorname{metric\text{-}top}(d), f) \end{array}}}
\vnbprop{bsw-compact-induction}{\vnbstmt{\begin{array}{@{}l@{}} \forall t, a. \\ \quad \operatorname{is\text{-}compact\text{-}t}(t) \wedge \\ \quad \emptyset \in a \wedge \\ \quad \forall b, c \in a.\; ((b \cup c) \in a) \wedge \\ \quad \forall p \in \operatorname{pts}(t).\; \exists u \in \operatorname{opens}(t).\; ((p \in u) \wedge (u \in a)) \Rightarrow \\ \quad \quad \operatorname{pts}(t) \in a \end{array}}}
\vnbprop{bsw-rcont-unfold}{\vnbstmt{\begin{array}{@{}l@{}} \forall t. \\ \quad (C_{\mathbb{R}}({t}) \simeq \{f \in (\operatorname{pts}(t) \to \mathbb{R}) : \operatorname{is\text{-}hom\text{-}top\text{-}space}(t, \operatorname{metric\text{-}top}(\operatorname{rr\text{-}ms}), f)\}) \end{array}}}
\vnbprop{bsw-rcont-parts}{\vnbstmt{\begin{array}{@{}l@{}} \forall t.; \forall f \in C_{\mathbb{R}}({t}). \\ \quad ((f \in (\operatorname{pts}(t) \to \mathbb{R})) \wedge (\operatorname{is\text{-}top\text{-}space}(t) \wedge \operatorname{is\text{-}hom\text{-}top\text{-}space}(t, \operatorname{metric\text{-}top}(\operatorname{rr\text{-}ms}), f))) \end{array}}}
\vnbprop{bsw-rcont-elim}{\vnbstmt{\begin{array}{@{}l@{}} \forall t.; \forall f \in C_{\mathbb{R}}({t}), p \in \operatorname{pts}(t), s. \\ \quad \operatorname{pos\text{-}rr}(s) \Rightarrow  \\ \quad \quad \exists u \in \operatorname{opens}(t).\; ((p \in u) \wedge \forall q \in u.\; (\lvert (f(p) - f(q)) \rvert < s)) \end{array}}}
\vnbprop{bsw-rcont-intro}{\vnbstmt{\begin{array}{@{}l@{}} \forall t, f. \\ \quad \operatorname{is\text{-}top\text{-}space}(t) \wedge \\ \quad f \in (\operatorname{pts}(t) \to \mathbb{R}) \wedge \\ \quad \forall p \in \operatorname{pts}(t), s. \\ \quad \quad \operatorname{pos\text{-}rr}(s) \Rightarrow  \\ \quad \quad \quad \exists u \in \operatorname{opens}(t).\; ((p \in u) \wedge \forall q \in u.\; (\lvert (f(p) - f(q)) \rvert < s)) \Rightarrow \\ \quad \quad f \in C_{\mathbb{R}}({t}) \end{array}}}
\vnbprop{bsw-rcont-const}{\vnbstmt{\begin{array}{@{}l@{}} \forall t, c, h. \\ \quad (\operatorname{is\text{-}top\text{-}space}(t) \wedge ((c \in \mathbb{R}) \wedge ((h \in (\operatorname{pts}(t) \to \mathbb{R})) \wedge \forall z \in \operatorname{pts}(t).\; (h(z) = c)))) \Rightarrow  \\ \quad \quad h \in C_{\mathbb{R}}({t}) \end{array}}}
\vnbprop{bsw-rcont-add}{\vnbstmt{\begin{array}{@{}l@{}} \forall t, f, g, h. \\ \quad f \in C_{\mathbb{R}}({t}) \wedge \\ \quad g \in C_{\mathbb{R}}({t}) \wedge \\ \quad h \in (\operatorname{pts}(t) \to \mathbb{R}) \wedge \\ \quad \forall z \in \operatorname{pts}(t).\; (h(z) = (f(z) + g(z))) \Rightarrow \\ \quad \quad h \in C_{\mathbb{R}}({t}) \end{array}}}
\vnbprop{bsw-rcont-lin}{\vnbstmt{\begin{array}{@{}l@{}} \forall t, f, g, a, b, h. \\ \quad f \in C_{\mathbb{R}}({t}) \wedge \\ \quad g \in C_{\mathbb{R}}({t}) \wedge \\ \quad a \in \mathbb{R} \wedge \\ \quad b \in \mathbb{R} \wedge \\ \quad \lvert a \rvert \leq 1 \wedge \\ \quad \lvert b \rvert \leq 1 \wedge \\ \quad h \in (\operatorname{pts}(t) \to \mathbb{R}) \wedge \\ \quad \forall z \in \operatorname{pts}(t).\; (h(z) = ((f(z) \cdot a) + (g(z) \cdot b))) \Rightarrow \\ \quad \quad h \in C_{\mathbb{R}}({t}) \end{array}}}
\vnbprop{bsw-rcont-abs}{\vnbstmt{\begin{array}{@{}l@{}} \forall t, f, h. \\ \quad ((f \in C_{\mathbb{R}}({t})) \wedge ((h \in (\operatorname{pts}(t) \to \mathbb{R})) \wedge \forall z \in \operatorname{pts}(t).\; (h(z) = \lvert f(z) \rvert))) \Rightarrow  \\ \quad \quad h \in C_{\mathbb{R}}({t}) \end{array}}}
\vnbprop{bsw-rcont-bounded}{\vnbstmt{\begin{array}{@{}l@{}} \forall t, f. \\ \quad (\operatorname{is\text{-}compact\text{-}t}(t) \wedge (f \in C_{\mathbb{R}}({t}))) \Rightarrow  \\ \quad \quad \exists m \in \mathbb{R}.\; ((0 \leq m) \wedge \forall z \in \operatorname{pts}(t).\; (\lvert f(z) \rvert \leq m)) \end{array}}}
\vnbprop{bsw-rcont-closure-unfold}{\vnbstmt{\begin{array}{@{}l@{}} \forall t, a. \\ \quad (\operatorname{rcont\text{-}closure}(t, a) \simeq \{h \in C_{\mathbb{R}}({t}) : \forall e.\; (\operatorname{pos\text{-}rr}(e) \Rightarrow \exists g \in a.\; \forall y \in \operatorname{pts}(t).\; (\lvert (h(y) - g(y)) \rvert < e))\}) \end{array}}}
\vnbprop{bsw-rcont-set}{\vnbstmt{\forall t.\; (\operatorname{is\text{-}top\text{-}space}(t) \Rightarrow (C_{\mathbb{R}}({t}) \in \mathbf{Set}))}}
\vnbprop{bsw-rcont-closure-contains}{\vnbstmt{\begin{array}{@{}l@{}} \forall t, a, f. \\ \quad (\operatorname{is\text{-}top\text{-}space}(t) \wedge ((a \subseteq C_{\mathbb{R}}({t})) \wedge (f \in a))) \Rightarrow  \\ \quad \quad f \in \operatorname{rcont\text{-}closure}(t, a) \end{array}}}
\vnbprop{bsw-rcont-closure-closed}{\vnbstmt{\begin{array}{@{}l@{}} \forall t, a, h. \\ \quad a \subseteq C_{\mathbb{R}}({t}) \wedge \\ \quad h \in C_{\mathbb{R}}({t}) \wedge \\ \quad \forall e. \\ \quad \quad \operatorname{pos\text{-}rr}(e) \Rightarrow  \\ \quad \quad \quad \exists g \in \operatorname{rcont\text{-}closure}(t, a).\; \forall y \in \operatorname{pts}(t).\; (\lvert (h(y) - g(y)) \rvert < e) \Rightarrow \\ \quad \quad h \in \operatorname{rcont\text{-}closure}(t, a) \end{array}}}
\vnbprop{bsw-sa-const}{\vnbstmt{\begin{array}{@{}l@{}} \forall t, a, c. \\ \quad (\operatorname{is\text{-}rcont\text{-}subalgebra}(t, a) \wedge (c \in \mathbb{R})) \Rightarrow  \\ \quad \quad \exists k \in a.\; \forall z \in \operatorname{pts}(t).\; (k(z) = c) \end{array}}}
\vnbprop{bsw-sa-add}{\vnbstmt{\begin{array}{@{}l@{}} \forall t, a, f, g. \\ \quad (\operatorname{is\text{-}rcont\text{-}subalgebra}(t, a) \wedge ((f \in a) \wedge (g \in a))) \Rightarrow  \\ \quad \quad \exists k \in a.\; \forall z \in \operatorname{pts}(t).\; (k(z) = (f(z) + g(z))) \end{array}}}
\vnbprop{bsw-sa-mul}{\vnbstmt{\begin{array}{@{}l@{}} \forall t, a, f, g. \\ \quad (\operatorname{is\text{-}rcont\text{-}subalgebra}(t, a) \wedge ((f \in a) \wedge (g \in a))) \Rightarrow  \\ \quad \quad \exists k \in a.\; \forall z \in \operatorname{pts}(t).\; (k(z) = (f(z) \cdot g(z))) \end{array}}}
\vnbprop{bsw-sa-poly}{\vnbstmt{\begin{array}{@{}l@{}} \forall t, a, g.; \forall n \in \mathbb{N}, c \in (\mathbb{N} \to \mathbb{R}). \\ \quad (\operatorname{is\text{-}rcont\text{-}subalgebra}(t, a) \wedge ((g \in a) \wedge (g \in C_{\mathbb{R}}({t})))) \Rightarrow  \\ \quad \quad \exists k \in a. \forall z \in \operatorname{pts}(t). \\ \quad \quad \quad k(z) = \operatorname{cc\text{-}series\text{-}partial\text{-}sum}((\lambda j \in \mathbb{N}.\; (c(j) \cdot \operatorname{power}(g(z), j))), n) \end{array}}}
\vnbprop{bsw-abs-cont-on}{\vnbstmt{\begin{array}{@{}l@{}} \forall a, b. \\ \quad \operatorname{is\text{-}continuous\text{-}on}((\lambda w \in \operatorname{ccint}(a, b).\; \lvert w \rvert), \operatorname{ccint}(a, b)) \end{array}}}
\vnbprop{bsw-abs-poly-approx}{\vnbstmt{\begin{array}{@{}l@{}} \forall m \in \mathbb{R}, e. \\ \quad ((0 \leq m) \wedge \operatorname{pos\text{-}rr}(e)) \Rightarrow  \\ \quad \quad \exists n \in \mathbb{N}, c \in (\mathbb{N} \to \mathbb{R}). \forall x \in \mathbb{R}. \\ \quad \quad \quad (\lvert x \rvert \leq m) \Rightarrow  \\ \quad \quad \quad \quad \lvert (\lvert x \rvert - \operatorname{cc\text{-}series\text{-}partial\text{-}sum}((\lambda j \in \mathbb{N}.\; (c(j) \cdot \operatorname{power}(x, j))), n)) \rvert < e \end{array}}}
\vnbprop{bsw-ucl-abs}{\vnbstmt{\begin{array}{@{}l@{}} \forall t, a, f, h. \\ \quad \operatorname{is\text{-}compact\text{-}t}(t) \wedge \\ \quad \operatorname{is\text{-}rcont\text{-}subalgebra}(t, a) \wedge \\ \quad f \in \operatorname{rcont\text{-}closure}(t, a) \wedge \\ \quad h \in (\operatorname{pts}(t) \to \mathbb{R}) \wedge \\ \quad \forall z \in \operatorname{pts}(t).\; (h(z) = \lvert f(z) \rvert) \Rightarrow \\ \quad \quad h \in \operatorname{rcont\text{-}closure}(t, a) \end{array}}}
\vnbprop{bsw-ucl-lin}{\vnbstmt{\begin{array}{@{}l@{}} \forall t, a, f, g, b, c, h. \\ \quad \operatorname{is\text{-}rcont\text{-}subalgebra}(t, a) \wedge \\ \quad f \in \operatorname{rcont\text{-}closure}(t, a) \wedge \\ \quad g \in \operatorname{rcont\text{-}closure}(t, a) \wedge \\ \quad b \in \mathbb{R} \wedge \\ \quad c \in \mathbb{R} \wedge \\ \quad \lvert b \rvert \leq 1 \wedge \\ \quad \lvert c \rvert \leq 1 \wedge \\ \quad h \in (\operatorname{pts}(t) \to \mathbb{R}) \wedge \\ \quad \forall z \in \operatorname{pts}(t).\; (h(z) = ((f(z) \cdot b) + (g(z) \cdot c))) \Rightarrow \\ \quad \quad h \in \operatorname{rcont\text{-}closure}(t, a) \end{array}}}
\vnbprop{bsw-ucl-max}{\vnbstmt{\begin{array}{@{}l@{}} \forall t, a, f, g. \\ \quad \operatorname{is\text{-}compact\text{-}t}(t) \wedge \\ \quad \operatorname{is\text{-}rcont\text{-}subalgebra}(t, a) \wedge \\ \quad f \in \operatorname{rcont\text{-}closure}(t, a) \wedge \\ \quad g \in \operatorname{rcont\text{-}closure}(t, a) \Rightarrow \\ \quad \quad \exists m \in \operatorname{rcont\text{-}closure}(t, a). \forall z \in \operatorname{pts}(t). \\ \quad \quad \quad (f(z) \leq m(z)) \wedge ((g(z) \leq m(z)) \wedge ((m(z) = f(z)) \vee (m(z) = g(z)))) \end{array}}}
\vnbprop{bsw-ucl-min}{\vnbstmt{\begin{array}{@{}l@{}} \forall t, a, f, g. \\ \quad \operatorname{is\text{-}compact\text{-}t}(t) \wedge \\ \quad \operatorname{is\text{-}rcont\text{-}subalgebra}(t, a) \wedge \\ \quad f \in \operatorname{rcont\text{-}closure}(t, a) \wedge \\ \quad g \in \operatorname{rcont\text{-}closure}(t, a) \Rightarrow \\ \quad \quad \exists m \in \operatorname{rcont\text{-}closure}(t, a). \forall z \in \operatorname{pts}(t). \\ \quad \quad \quad (m(z) \leq f(z)) \wedge ((m(z) \leq g(z)) \wedge ((m(z) = f(z)) \vee (m(z) = g(z)))) \end{array}}}
\vnbprop{bsw-sa-interp}{\vnbstmt{\begin{array}{@{}l@{}} \forall t, c, p, q, a, b. \\ \quad \operatorname{is\text{-}rcont\text{-}subalgebra}(t, c) \wedge \\ \quad \operatorname{is\text{-}point\text{-}separating}(t, c) \wedge \\ \quad p \in \operatorname{pts}(t) \wedge \\ \quad q \in \operatorname{pts}(t) \wedge \\ \quad \neg (p = q) \wedge \\ \quad a \in \mathbb{R} \wedge \\ \quad b \in \mathbb{R} \Rightarrow \\ \quad \quad \exists k \in c.\; ((k(p) = a) \wedge (k(q) = b)) \end{array}}}
\vnbprop{bsw-sa-interp-f}{\vnbstmt{\begin{array}{@{}l@{}} \forall t, a, f, p, q. \\ \quad \operatorname{is\text{-}rcont\text{-}subalgebra}(t, a) \wedge \\ \quad \operatorname{is\text{-}point\text{-}separating}(t, a) \wedge \\ \quad f \in C_{\mathbb{R}}({t}) \wedge \\ \quad p \in \operatorname{pts}(t) \wedge \\ \quad q \in \operatorname{pts}(t) \Rightarrow \\ \quad \quad \exists k \in a.\; ((k(p) = f(p)) \wedge (k(q) = f(q))) \end{array}}}
\vnbprop{bsw-sw-step1}{\vnbstmt{\begin{array}{@{}l@{}} \forall t, a, f, x, e. \\ \quad \operatorname{is\text{-}compact\text{-}t}(t) \wedge \\ \quad \operatorname{is\text{-}rcont\text{-}subalgebra}(t, a) \wedge \\ \quad \operatorname{is\text{-}point\text{-}separating}(t, a) \wedge \\ \quad f \in C_{\mathbb{R}}({t}) \wedge \\ \quad x \in \operatorname{pts}(t) \wedge \\ \quad \operatorname{pos\text{-}rr}(e) \Rightarrow \\ \quad \quad \exists g \in \operatorname{rcont\text{-}closure}(t, a). \\ \quad \quad \quad (g(x) = f(x)) \wedge \forall z \in \operatorname{pts}(t).\; (g(z) < (f(z) + e)) \end{array}}}
\vnbprop{bsw-sw-step2}{\vnbstmt{\begin{array}{@{}l@{}} \forall t, a, f, e. \\ \quad \operatorname{is\text{-}compact\text{-}t}(t) \wedge \\ \quad \operatorname{is\text{-}rcont\text{-}subalgebra}(t, a) \wedge \\ \quad \operatorname{is\text{-}point\text{-}separating}(t, a) \wedge \\ \quad f \in C_{\mathbb{R}}({t}) \wedge \\ \quad \operatorname{pos\text{-}rr}(e) \Rightarrow \\ \quad \quad \exists g \in \operatorname{rcont\text{-}closure}(t, a). \forall z \in \operatorname{pts}(t). \\ \quad \quad \quad (g(z) < (f(z) + e)) \wedge ((f(z) - e) < g(z)) \end{array}}}
\vnbprop{bsw-stone-weierstrass-real}{\vnbstmt{\begin{array}{@{}l@{}} \forall t, a, f, e. \\ \quad \operatorname{is\text{-}compact\text{-}t}(t) \wedge \\ \quad \operatorname{is\text{-}rcont\text{-}subalgebra}(t, a) \wedge \\ \quad \operatorname{is\text{-}point\text{-}separating}(t, a) \wedge \\ \quad f \in C_{\mathbb{R}}({t}) \wedge \\ \quad \operatorname{pos\text{-}rr}(e) \Rightarrow \\ \quad \quad \exists g \in a.\; \forall z \in \operatorname{pts}(t).\; (\lvert (f(z) - g(z)) \rvert < e) \end{array}}}
\vnbprop{bsw-stone-weierstrass-real-dense}{\vnbstmt{\begin{array}{@{}l@{}} \forall t, a, f. \\ \quad (\operatorname{is\text{-}compact\text{-}t}(t) \wedge (\operatorname{is\text{-}rcont\text{-}subalgebra}(t, a) \wedge (\operatorname{is\text{-}point\text{-}separating}(t, a) \wedge (f \in C_{\mathbb{R}}({t}))))) \Rightarrow  \\ \quad \quad f \in \operatorname{rcont\text{-}closure}(t, a) \end{array}}}
\vnbprop{bsw-cc-elim}{\vnbstmt{\begin{array}{@{}l@{}} \forall t, f.; \forall p \in \operatorname{pts}(t), s. \\ \quad (\operatorname{is\text{-}top\text{-}space}(t) \wedge ((f \in (\operatorname{pts}(t) \to \mathbb{C})) \wedge (\operatorname{is\text{-}hom\text{-}top\text{-}space}(t, \operatorname{metric\text{-}top}(\operatorname{cc\text{-}ms}), f) \wedge \operatorname{pos\text{-}rr}(s)))) \Rightarrow  \\ \quad \quad \exists u \in \operatorname{opens}(t). \\ \quad \quad \quad (p \in u) \wedge \forall q \in u.\; (\operatorname{magnitude}((f(p) - f(q))) < s) \end{array}}}
\vnbprop{bsw-rcont-of-cc}{\vnbstmt{\begin{array}{@{}l@{}} \forall t, g, h. \\ \quad \operatorname{is\text{-}top\text{-}space}(t) \wedge \\ \quad g \in (\operatorname{pts}(t) \to \mathbb{C}) \wedge \\ \quad \operatorname{is\text{-}hom\text{-}top\text{-}space}(t, \operatorname{metric\text{-}top}(\operatorname{cc\text{-}ms}), g) \wedge \\ \quad h \in (\operatorname{pts}(t) \to \mathbb{R}) \wedge \\ \quad \forall z \in \operatorname{pts}(t).\; (h(z) = g(z)) \Rightarrow \\ \quad \quad h \in C_{\mathbb{R}}({t}) \end{array}}}
\vnbprop{bsw-rcont-re}{\vnbstmt{\begin{array}{@{}l@{}} \forall t, f, h. \\ \quad \operatorname{is\text{-}top\text{-}space}(t) \wedge \\ \quad f \in (\operatorname{pts}(t) \to \mathbb{C}) \wedge \\ \quad \operatorname{is\text{-}hom\text{-}top\text{-}space}(t, \operatorname{metric\text{-}top}(\operatorname{cc\text{-}ms}), f) \wedge \\ \quad h \in (\operatorname{pts}(t) \to \mathbb{R}) \wedge \\ \quad \forall z \in \operatorname{pts}(t).\; (h(z) = \operatorname{real\text{-}part}(f(z))) \Rightarrow \\ \quad \quad h \in C_{\mathbb{R}}({t}) \end{array}}}
\vnbprop{bsw-rcont-im}{\vnbstmt{\begin{array}{@{}l@{}} \forall t, f, h. \\ \quad \operatorname{is\text{-}top\text{-}space}(t) \wedge \\ \quad f \in (\operatorname{pts}(t) \to \mathbb{C}) \wedge \\ \quad \operatorname{is\text{-}hom\text{-}top\text{-}space}(t, \operatorname{metric\text{-}top}(\operatorname{cc\text{-}ms}), f) \wedge \\ \quad h \in (\operatorname{pts}(t) \to \mathbb{R}) \wedge \\ \quad \forall z \in \operatorname{pts}(t).\; (h(z) = \operatorname{imag\text{-}part}(f(z))) \Rightarrow \\ \quad \quad h \in C_{\mathbb{R}}({t}) \end{array}}}
\vnbprop{bsw-cs-const}{\vnbstmt{\begin{array}{@{}l@{}} \forall t, a, c. \\ \quad (\operatorname{is\text{-}ccont\text{-}star\text{-}subalgebra}(t, a) \wedge (c \in \mathbb{C})) \Rightarrow  \\ \quad \quad \exists k \in a.\; \forall z \in \operatorname{pts}(t).\; (k(z) = c) \end{array}}}
\vnbprop{bsw-cs-add}{\vnbstmt{\begin{array}{@{}l@{}} \forall t, a, f, g. \\ \quad (\operatorname{is\text{-}ccont\text{-}star\text{-}subalgebra}(t, a) \wedge ((f \in a) \wedge (g \in a))) \Rightarrow  \\ \quad \quad \exists k \in a.\; \forall z \in \operatorname{pts}(t).\; (k(z) = (f(z) + g(z))) \end{array}}}
\vnbprop{bsw-cs-mul}{\vnbstmt{\begin{array}{@{}l@{}} \forall t, a, f, g. \\ \quad (\operatorname{is\text{-}ccont\text{-}star\text{-}subalgebra}(t, a) \wedge ((f \in a) \wedge (g \in a))) \Rightarrow  \\ \quad \quad \exists k \in a.\; \forall z \in \operatorname{pts}(t).\; (k(z) = (f(z) \cdot g(z))) \end{array}}}
\vnbprop{bsw-cs-conj}{\vnbstmt{\begin{array}{@{}l@{}} \forall t, a, f. \\ \quad (\operatorname{is\text{-}ccont\text{-}star\text{-}subalgebra}(t, a) \wedge (f \in a)) \Rightarrow  \\ \quad \quad \exists k \in a.\; \forall z \in \operatorname{pts}(t).\; (k(z) = \operatorname{conjugate}(f(z))) \end{array}}}
\vnbprop{bsw-cs-real-sa}{\vnbstmt{\begin{array}{@{}l@{}} \forall t, a. \\ \quad (\operatorname{is\text{-}top\text{-}space}(t) \wedge \operatorname{is\text{-}ccont\text{-}star\text{-}subalgebra}(t, a)) \Rightarrow  \\ \quad \quad \operatorname{is\text{-}rcont\text{-}subalgebra}(t, \{h \in C_{\mathbb{R}}({t}) : \exists g \in a.\; \forall z \in \operatorname{pts}(t).\; (h(z) = g(z))\}) \end{array}}}
\vnbprop{bsw-cs-real-sep}{\vnbstmt{\begin{array}{@{}l@{}} \forall t, a. \\ \quad (\operatorname{is\text{-}top\text{-}space}(t) \wedge (\operatorname{is\text{-}ccont\text{-}star\text{-}subalgebra}(t, a) \wedge \operatorname{is\text{-}point\text{-}separating}(t, a))) \Rightarrow  \\ \quad \quad \operatorname{is\text{-}point\text{-}separating}(t, \{h \in C_{\mathbb{R}}({t}) : \exists g \in a.\; \forall z \in \operatorname{pts}(t).\; (h(z) = g(z))\}) \end{array}}}
\vnbprop{bsw-stone-weierstrass-complex}{\vnbstmt{\begin{array}{@{}l@{}} \forall t, a, f, e. \\ \quad \operatorname{is\text{-}compact\text{-}t}(t) \wedge \\ \quad \operatorname{is\text{-}ccont\text{-}star\text{-}subalgebra}(t, a) \wedge \\ \quad \operatorname{is\text{-}point\text{-}separating}(t, a) \wedge \\ \quad f \in (\operatorname{pts}(t) \to \mathbb{C}) \wedge \\ \quad \operatorname{is\text{-}hom\text{-}top\text{-}space}(t, \operatorname{metric\text{-}top}(\operatorname{cc\text{-}ms}), f) \wedge \\ \quad \operatorname{pos\text{-}rr}(e) \Rightarrow \\ \quad \quad \exists g \in a.\; \forall z \in \operatorname{pts}(t).\; (\operatorname{magnitude}((f(z) - g(z))) < e) \end{array}}}

\vnbbookfile{theorem-library/cna-normed-ag-typing}{theorem-library/cna-normed-ag-typing}
\noindent{\small theorem-library/cna-normed-ag-typing.scm -- THE NORMED GROUPS OF A COMPLEX NORMED ALGEBRA AND OF A UNITAL BANACH ALGEBRA (D8, 2026-11-15; the user's decision "Prove views").}\par
\vnbprop{cna-as-normed-ag-is-normed-ag}{\vnbstmt{\begin{array}{@{}l@{}} \forall r. \\ \quad \operatorname{is\text{-}complex\text{-}normed\text{-}algebra}(r) \Rightarrow \operatorname{is\text{-}normed\text{-}ag}(\operatorname{cna\text{-}as\text{-}normed\text{-}ag}(r)) \end{array}}}
\vnbprop{uba-as-normed-ag-is-normed-ag}{\vnbstmt{\begin{array}{@{}l@{}} \forall r. \\ \quad \operatorname{is\text{-}unital\text{-}banach\text{-}algebra}(r) \Rightarrow \operatorname{is\text{-}normed\text{-}ag}(\operatorname{uba\text{-}as\text{-}normed\text{-}ag}(r)) \end{array}}}

\vnbbookfile{theorem-library/ba-characters}{theorem-library/ba-characters}
\noindent{\small ba-characters.scm -- wave of 2026-11-09/10 (agent BA-1G), wired 2026-11-10 from scratchpad/ba1g-1109/ba-characters.scm. Library: operator-algebras.}\par
\vnbprop{bch-comm-uba}{\vnbstmt{\begin{array}{@{}l@{}} \forall a. \\ \quad \operatorname{is\text{-}commutative\text{-}banach\text{-}algebra}(a) \Rightarrow \operatorname{is\text{-}unital\text{-}banach\text{-}algebra}(a) \end{array}}}
\vnbprop{bch-comm}{\vnbstmt{\begin{array}{@{}l@{}} \forall a.; \forall x, y \in \operatorname{vec}(a). \\ \quad \operatorname{is\text{-}commutative\text{-}banach\text{-}algebra}(a) \Rightarrow (\cdot_{a}(x, y) = \cdot_{a}(y, x)) \end{array}}}
\vnbprop{bch-left-ideal-two-sided}{\vnbstmt{\begin{array}{@{}l@{}} \forall a, j. \\ \quad (\operatorname{is\text{-}commutative\text{-}banach\text{-}algebra}(a) \wedge \operatorname{is\text{-}left\text{-}ideal}(a, j)) \Rightarrow  \\ \quad \quad \operatorname{is\text{-}two\text{-}sided\text{-}ideal}(a, j) \end{array}}}
\vnbprop{commutative-maximal-ideal-left}{\vnbstmt{\begin{array}{@{}l@{}} \forall a, m. \\ \quad (\operatorname{is\text{-}commutative\text{-}banach\text{-}algebra}(a) \wedge \operatorname{is\text{-}maximal\text{-}ideal}(a, m)) \Rightarrow  \\ \quad \quad \operatorname{is\text{-}maximal\text{-}left\text{-}ideal}(a, m) \end{array}}}
\vnbprop{commutative-maximal-ideal-closed}{\vnbstmt{\begin{array}{@{}l@{}} \forall a, m. \\ \quad (\operatorname{is\text{-}commutative\text{-}banach\text{-}algebra}(a) \wedge \operatorname{is\text{-}maximal\text{-}ideal}(a, m)) \Rightarrow  \\ \quad \quad \operatorname{is\text{-}closed}(\operatorname{ba\text{-}metric\text{-}space}(a), m) \end{array}}}
\vnbprop{bch-class-zero-iff}{\vnbstmt{\begin{array}{@{}l@{}} \forall a, i.; \forall m \in \operatorname{vec}(a). \\ \quad (\operatorname{is\text{-}commutative\text{-}banach\text{-}algebra}(a) \wedge \operatorname{is\text{-}two\text{-}sided\text{-}ideal}(a, i)) \Rightarrow  \\ \quad \quad \operatorname{class}(\operatorname{qmod\text{-}setoid}(\operatorname{uba\text{-}as\text{-}module}(a), i), m) = \operatorname{class}(\operatorname{qmod\text{-}setoid}(\operatorname{uba\text{-}as\text{-}module}(a), i), \operatorname{vzero}(a)) \\ \quad \quad \Leftrightarrow\; m \in i \end{array}}}
\vnbprop{bch-gen-ideal}{\vnbstmt{\begin{array}{@{}l@{}} \forall a, i.; \forall x \in \operatorname{vec}(a). \\ \quad (\operatorname{is\text{-}commutative\text{-}banach\text{-}algebra}(a) \wedge \operatorname{is\text{-}two\text{-}sided\text{-}ideal}(a, i)) \Rightarrow  \\ \quad \quad \operatorname{is\text{-}two\text{-}sided\text{-}ideal}(a, \\ \quad \quad \quad \{w \in \operatorname{vec}(a) : \exists b \in \operatorname{vec}(a).\; (\operatorname{class}(\operatorname{qmod\text{-}setoid}(\operatorname{uba\text{-}as\text{-}module}(a), i), w) = \operatorname{class}(\operatorname{qmod\text{-}setoid}(\operatorname{uba\text{-}as\text{-}module}(a), i), \cdot_{a}(b, x)))\}) \end{array}}}
\vnbprop{bch-gen-contains}{\vnbstmt{\begin{array}{@{}l@{}} \forall a, i.; \forall x \in \operatorname{vec}(a). \\ \quad (\operatorname{is\text{-}commutative\text{-}banach\text{-}algebra}(a) \wedge \operatorname{is\text{-}two\text{-}sided\text{-}ideal}(a, i)) \Rightarrow  \\ \quad \quad i \\ \quad \quad \subseteq\; \{w \in \operatorname{vec}(a) : \exists b \in \operatorname{vec}(a).\; (\operatorname{class}(\operatorname{qmod\text{-}setoid}(\operatorname{uba\text{-}as\text{-}module}(a), i), w) = \operatorname{class}(\operatorname{qmod\text{-}setoid}(\operatorname{uba\text{-}as\text{-}module}(a), i), \cdot_{a}(b, x)))\} \end{array}}}
\vnbprop{bch-gen-mem}{\vnbstmt{\begin{array}{@{}l@{}} \forall a, i.; \forall x \in \operatorname{vec}(a). \\ \quad (\operatorname{is\text{-}commutative\text{-}banach\text{-}algebra}(a) \wedge \operatorname{is\text{-}two\text{-}sided\text{-}ideal}(a, i)) \Rightarrow  \\ \quad \quad x \\ \quad \quad \in\; \{w \in \operatorname{vec}(a) : \exists b \in \operatorname{vec}(a).\; (\operatorname{class}(\operatorname{qmod\text{-}setoid}(\operatorname{uba\text{-}as\text{-}module}(a), i), w) = \operatorname{class}(\operatorname{qmod\text{-}setoid}(\operatorname{uba\text{-}as\text{-}module}(a), i), \cdot_{a}(b, x)))\} \end{array}}}
\vnbprop{maximal-quotient-field}{\vnbstmt{\begin{array}{@{}l@{}} \forall a, i.; \forall c \in \operatorname{vec}({a}/{i}). \\ \quad (\operatorname{is\text{-}commutative\text{-}banach\text{-}algebra}(a) \wedge (\operatorname{is\text{-}maximal\text{-}ideal}(a, i) \wedge \neg (c = \operatorname{vzero}({a}/{i})))) \Rightarrow  \\ \quad \quad \operatorname{is\text{-}invertible\text{-}el}({a}/{i}, c) \end{array}}}
\vnbprop{maximal-quotient-scalar}{\vnbstmt{\begin{array}{@{}l@{}} \forall a, i.; \forall x \in \operatorname{vec}(a). \\ \quad (\operatorname{is\text{-}commutative\text{-}banach\text{-}algebra}(a) \wedge \operatorname{is\text{-}maximal\text{-}ideal}(a, i)) \Rightarrow  \\ \quad \quad \exists z \in \mathbb{C}. \\ \quad \quad \quad \operatorname{class}(\operatorname{qmod\text{-}setoid}(\operatorname{uba\text{-}as\text{-}module}(a), i), x) \\ \quad \quad \quad =\; \operatorname{class}(\operatorname{qmod\text{-}setoid}(\operatorname{uba\text{-}as\text{-}module}(a), i), \operatorname{act}(a)(z, 1_{a})) \end{array}}}
\vnbprop{maximal-quotient-scalar-unique}{\vnbstmt{\begin{array}{@{}l@{}} \forall a, i.; \forall z, w \in \mathbb{C}. \\ \quad \operatorname{is\text{-}commutative\text{-}banach\text{-}algebra}(a) \wedge \\ \quad \operatorname{is\text{-}maximal\text{-}ideal}(a, i) \wedge \\ \quad \operatorname{class}(\operatorname{qmod\text{-}setoid}(\operatorname{uba\text{-}as\text{-}module}(a), i), \operatorname{act}(a)(z, 1_{a})) \\ \quad =\; \operatorname{class}(\operatorname{qmod\text{-}setoid}(\operatorname{uba\text{-}as\text{-}module}(a), i), \operatorname{act}(a)(w, 1_{a})) \Rightarrow \\ \quad \quad z = w \end{array}}}
\vnbprop{bch-chi-value}{\vnbstmt{\begin{array}{@{}l@{}} \forall a, i.; \forall t \in \operatorname{vec}(a). \\ \quad (\operatorname{is\text{-}commutative\text{-}banach\text{-}algebra}(a) \wedge \operatorname{is\text{-}maximal\text{-}ideal}(a, i)) \Rightarrow  \\ \quad \quad (\iota\, v.\;((v \in \mathbb{C}) \wedge (\operatorname{class}(\operatorname{qmod\text{-}setoid}(\operatorname{uba\text{-}as\text{-}module}(a), i), t) = \operatorname{class}(\operatorname{qmod\text{-}setoid}(\operatorname{uba\text{-}as\text{-}module}(a), i), \operatorname{act}(a)(v, 1_{a}))))) \\ \quad \quad \in\; \mathbb{C} \wedge \\ \quad \quad \operatorname{class}(\operatorname{qmod\text{-}setoid}(\operatorname{uba\text{-}as\text{-}module}(a), i), t) \\ \quad \quad =\; \operatorname{class}(\operatorname{qmod\text{-}setoid}(\operatorname{uba\text{-}as\text{-}module}(a), i), \\ \quad \quad \quad \operatorname{act}(a)((\iota\, v.\;((v \in \mathbb{C}) \wedge (\operatorname{class}(\operatorname{qmod\text{-}setoid}(\operatorname{uba\text{-}as\text{-}module}(a), i), t) = \operatorname{class}(\operatorname{qmod\text{-}setoid}(\operatorname{uba\text{-}as\text{-}module}(a), i), \operatorname{act}(a)(v, 1_{a}))))), \\ \quad \quad \quad \quad 1_{a})) \end{array}}}
\vnbprop{bch-chi-unique}{\vnbstmt{\begin{array}{@{}l@{}} \forall a, i.; \forall t \in \operatorname{vec}(a), w \in \mathbb{C}. \\ \quad \operatorname{is\text{-}commutative\text{-}banach\text{-}algebra}(a) \wedge \\ \quad \operatorname{is\text{-}maximal\text{-}ideal}(a, i) \wedge \\ \quad \operatorname{class}(\operatorname{qmod\text{-}setoid}(\operatorname{uba\text{-}as\text{-}module}(a), i), t) \\ \quad =\; \operatorname{class}(\operatorname{qmod\text{-}setoid}(\operatorname{uba\text{-}as\text{-}module}(a), i), \operatorname{act}(a)(w, 1_{a})) \Rightarrow \\ \quad \quad (\iota\, v.\;((v \in \mathbb{C}) \wedge (\operatorname{class}(\operatorname{qmod\text{-}setoid}(\operatorname{uba\text{-}as\text{-}module}(a), i), t) = \operatorname{class}(\operatorname{qmod\text{-}setoid}(\operatorname{uba\text{-}as\text{-}module}(a), i), \operatorname{act}(a)(v, 1_{a}))))) \\ \quad \quad =\; w \end{array}}}
\vnbprop{maximal-ideal-character}{\vnbstmt{\begin{array}{@{}l@{}} \forall a, i. \\ \quad (\operatorname{is\text{-}commutative\text{-}banach\text{-}algebra}(a) \wedge \operatorname{is\text{-}maximal\text{-}ideal}(a, i)) \Rightarrow  \\ \quad \quad \exists b \in \operatorname{characters}(a).\; \forall x \in \operatorname{vec}(a).\; ((b(x) = 0) \Leftrightarrow (x \in i)) \end{array}}}
\vnbprop{maximal-quotient-norm-magnitude}{\vnbstmt{\begin{array}{@{}l@{}} \forall a, i.; \forall x \in \operatorname{vec}(a), z \in \mathbb{C}. \\ \quad \operatorname{is\text{-}commutative\text{-}banach\text{-}algebra}(a) \wedge \\ \quad \operatorname{is\text{-}maximal\text{-}ideal}(a, i) \wedge \\ \quad \operatorname{class}(\operatorname{qmod\text{-}setoid}(\operatorname{uba\text{-}as\text{-}module}(a), i), x) \\ \quad =\; \operatorname{class}(\operatorname{qmod\text{-}setoid}(\operatorname{uba\text{-}as\text{-}module}(a), i), \operatorname{act}(a)(z, 1_{a})) \Rightarrow \\ \quad \quad \operatorname{vnrm}({a}/{i})(\operatorname{class}(\operatorname{qmod\text{-}setoid}(\operatorname{uba\text{-}as\text{-}module}(a), i), x)) = \operatorname{magnitude}(z) \end{array}}}
\vnbprop{character-shift-value}{\vnbstmt{\begin{array}{@{}l@{}} \forall a.; \forall c \in \operatorname{characters}(a), z \in \mathbb{C}, x \in \operatorname{vec}(a). \\ \quad \operatorname{is\text{-}unital\text{-}banach\text{-}algebra}(a) \Rightarrow (c(\operatorname{shift\text{-}el}(a, z, x)) = (z + -c(x))) \end{array}}}
\vnbprop{character-kernel-maximal}{\vnbstmt{\begin{array}{@{}l@{}} \forall a.; \forall c \in \operatorname{characters}(a). \\ \quad \operatorname{is\text{-}unital\text{-}banach\text{-}algebra}(a) \Rightarrow \operatorname{is\text{-}maximal\text{-}ideal}(a, \{x \in \operatorname{vec}(a) : (c(x) = 0)\}) \end{array}}}
\vnbprop{character-kernel-injective}{\vnbstmt{\begin{array}{@{}l@{}} \forall a.; \forall c, d \in \operatorname{characters}(a). \\ \quad (\operatorname{is\text{-}unital\text{-}banach\text{-}algebra}(a) \wedge (\{x \in \operatorname{vec}(a) : (c(x) = 0)\} = \{x \in \operatorname{vec}(a) : (d(x) = 0)\})) \Rightarrow  \\ \quad \quad c = d \end{array}}}
\vnbprop{maximal-ideal-kernel}{\vnbstmt{\begin{array}{@{}l@{}} \forall a, i. \\ \quad (\operatorname{is\text{-}commutative\text{-}banach\text{-}algebra}(a) \wedge \operatorname{is\text{-}maximal\text{-}ideal}(a, i)) \Rightarrow  \\ \quad \quad \exists b \in \operatorname{characters}(a).\; (\{x \in \operatorname{vec}(a) : (b(x) = 0)\} = i) \end{array}}}
\vnbprop{principal-ideal-two-sided}{\vnbstmt{\begin{array}{@{}l@{}} \forall a.; \forall s \in \operatorname{vec}(a). \\ \quad \operatorname{is\text{-}commutative\text{-}banach\text{-}algebra}(a) \Rightarrow  \\ \quad \quad \operatorname{is\text{-}two\text{-}sided\text{-}ideal}(a, \{w \in \operatorname{vec}(a) : \exists b \in \operatorname{vec}(a).\; (w = \cdot_{a}(b, s))\}) \end{array}}}
\vnbprop{noninvertible-in-maximal-ideal}{\vnbstmt{\begin{array}{@{}l@{}} \forall a.; \forall s \in \operatorname{vec}(a). \\ \quad (\operatorname{is\text{-}commutative\text{-}banach\text{-}algebra}(a) \wedge \neg \operatorname{is\text{-}invertible\text{-}el}(a, s)) \Rightarrow  \\ \quad \quad \exists m.\; (\operatorname{is\text{-}maximal\text{-}ideal}(a, m) \wedge (s \in m)) \end{array}}}
\vnbprop{spectrum-character-value}{\vnbstmt{\begin{array}{@{}l@{}} \forall a.; \forall x \in \operatorname{vec}(a), z \in \sigma({x}). \\ \quad \operatorname{is\text{-}commutative\text{-}banach\text{-}algebra}(a) \Rightarrow \exists i \in \operatorname{characters}(a).\; (i(x) = z) \end{array}}}
\vnbprop{characters-in-set}{\vnbstmt{\forall a.\; (\operatorname{is\text{-}unital\text{-}banach\text{-}algebra}(a) \Rightarrow (\operatorname{characters}(a) \in \mathbf{Set}))}}
\vnbprop{gelfand-transform-in-fun}{\vnbstmt{\begin{array}{@{}l@{}} \forall a.; \forall x \in \operatorname{vec}(a). \\ \quad \operatorname{is\text{-}unital\text{-}banach\text{-}algebra}(a) \Rightarrow  \\ \quad \quad \operatorname{gelfand\text{-}transform}(a)(x) \in (\operatorname{characters}(a) \to \mathbb{C}) \end{array}}}
\vnbprop{gelfand-range-spectrum}{\vnbstmt{\begin{array}{@{}l@{}} \forall a.; \forall x \in \operatorname{vec}(a). \\ \quad \operatorname{is\text{-}commutative\text{-}banach\text{-}algebra}(a) \Rightarrow  \\ \quad \quad \sigma({x}) = \operatorname{image}(\operatorname{gelfand\text{-}transform}(a)(x), \operatorname{characters}(a)) \end{array}}}
\vnbprop{character-norm-one}{\vnbstmt{\begin{array}{@{}l@{}} \forall a.; \forall c \in \operatorname{characters}(a). \\ \quad \operatorname{is\text{-}unital\text{-}banach\text{-}algebra}(a) \Rightarrow (\operatorname{cdual\text{-}norm}(a, c) = 1) \end{array}}}

\vnbbookfile{theorem-library/ba2-cstar}{theorem-library/ba2-cstar}
\noindent{\small ba2-cstar.scm -- wave of 2026-11-09/10 (agent BA-2), wired 2026-11-10 from scratchpad/ba2-1109/ba2-cstar.scm. Library: operator-algebras.}\par
\vnbprop{bcp-ro-scal}{\vnbstmt{\begin{array}{@{}l@{}} \forall a. \\ \quad \operatorname{is\text{-}unital\text{-}involutive\text{-}banach\text{-}algebra}(a) \Rightarrow (\operatorname{scal}(\operatorname{uiba\text{-}as\text{-}uba}(a)) \simeq \operatorname{scal}(a)) \end{array}}}
\vnbprop{bcp-ro-vec}{\vnbstmt{\begin{array}{@{}l@{}} \forall a. \\ \quad \operatorname{is\text{-}unital\text{-}involutive\text{-}banach\text{-}algebra}(a) \Rightarrow (\operatorname{vec}(\operatorname{uiba\text{-}as\text{-}uba}(a)) \simeq \operatorname{vec}(a)) \end{array}}}
\vnbprop{bcp-ro-vadd}{\vnbstmt{\begin{array}{@{}l@{}} \forall a. \\ \quad \operatorname{is\text{-}unital\text{-}involutive\text{-}banach\text{-}algebra}(a) \Rightarrow (\operatorname{vadd}(\operatorname{uiba\text{-}as\text{-}uba}(a)) \simeq \operatorname{vadd}(a)) \end{array}}}
\vnbprop{bcp-ro-vzero}{\vnbstmt{\begin{array}{@{}l@{}} \forall a. \\ \quad \operatorname{is\text{-}unital\text{-}involutive\text{-}banach\text{-}algebra}(a) \Rightarrow (\operatorname{vzero}(\operatorname{uiba\text{-}as\text{-}uba}(a)) \simeq \operatorname{vzero}(a)) \end{array}}}
\vnbprop{bcp-ro-vneg}{\vnbstmt{\begin{array}{@{}l@{}} \forall a. \\ \quad \operatorname{is\text{-}unital\text{-}involutive\text{-}banach\text{-}algebra}(a) \Rightarrow (\operatorname{vneg}(\operatorname{uiba\text{-}as\text{-}uba}(a)) \simeq \operatorname{vneg}(a)) \end{array}}}
\vnbprop{bcp-ro-act}{\vnbstmt{\begin{array}{@{}l@{}} \forall a. \\ \quad \operatorname{is\text{-}unital\text{-}involutive\text{-}banach\text{-}algebra}(a) \Rightarrow (\operatorname{act}(\operatorname{uiba\text{-}as\text{-}uba}(a)) \simeq \operatorname{act}(a)) \end{array}}}
\vnbprop{bcp-ro-vnrm}{\vnbstmt{\begin{array}{@{}l@{}} \forall a. \\ \quad \operatorname{is\text{-}unital\text{-}involutive\text{-}banach\text{-}algebra}(a) \Rightarrow (\operatorname{vnrm}(\operatorname{uiba\text{-}as\text{-}uba}(a)) \simeq \operatorname{vnrm}(a)) \end{array}}}
\vnbprop{bcp-ro-vmul}{\vnbstmt{\begin{array}{@{}l@{}} \forall a. \\ \quad \operatorname{is\text{-}unital\text{-}involutive\text{-}banach\text{-}algebra}(a) \Rightarrow (\cdot_{\operatorname{uiba\text{-}as\text{-}uba}(a)} \simeq \cdot_{a}) \end{array}}}
\vnbprop{bcp-ro-vone}{\vnbstmt{\begin{array}{@{}l@{}} \forall a. \\ \quad \operatorname{is\text{-}unital\text{-}involutive\text{-}banach\text{-}algebra}(a) \Rightarrow (1_{\operatorname{uiba\text{-}as\text{-}uba}(a)} \simeq 1_{a}) \end{array}}}
\vnbprop{bcp-inverse-unfold}{\vnbstmt{\begin{array}{@{}l@{}} \forall a, x. \\ \quad ({x}^{-1} \simeq (\iota\, u.\;((u \in \operatorname{vec}(a)) \wedge ((\cdot_{a}(u, x) = 1_{a}) \wedge (\cdot_{a}(x, u) = 1_{a}))))) \end{array}}}
\vnbprop{bcp-u-bae-vnrm-type}{\vnbstmt{\begin{array}{@{}l@{}} \forall a.; \forall x \in \operatorname{vec}(a). \\ \quad \operatorname{is\text{-}unital\text{-}involutive\text{-}banach\text{-}algebra}(a) \Rightarrow (\operatorname{vnrm}(a)(x) \in \mathbb{R}) \end{array}}}
\vnbprop{bcp-u-bae-inverse-prop}{\vnbstmt{\begin{array}{@{}l@{}} \forall a, x. \\ \quad (\operatorname{is\text{-}unital\text{-}involutive\text{-}banach\text{-}algebra}(a) \wedge \operatorname{is\text{-}invertible\text{-}el}(a, x)) \Rightarrow  \\ \quad \quad {x}^{-1} \in \operatorname{vec}(a) \wedge \\ \quad \quad \cdot_{a}({x}^{-1}, x) = 1_{a} \wedge \\ \quad \quad \cdot_{a}(x, {x}^{-1}) = 1_{a} \end{array}}}
\vnbprop{bcp-u-bae-inverse-in}{\vnbstmt{\begin{array}{@{}l@{}} \forall a, x. \\ \quad (\operatorname{is\text{-}unital\text{-}involutive\text{-}banach\text{-}algebra}(a) \wedge \operatorname{is\text{-}invertible\text{-}el}(a, x)) \Rightarrow  \\ \quad \quad {x}^{-1} \in \operatorname{vec}(a) \end{array}}}
\vnbprop{bcp-u-bae-inverse-left}{\vnbstmt{\begin{array}{@{}l@{}} \forall a, x. \\ \quad (\operatorname{is\text{-}unital\text{-}involutive\text{-}banach\text{-}algebra}(a) \wedge \operatorname{is\text{-}invertible\text{-}el}(a, x)) \Rightarrow  \\ \quad \quad \cdot_{a}({x}^{-1}, x) = 1_{a} \end{array}}}
\vnbprop{bcp-u-bae-inverse-right}{\vnbstmt{\begin{array}{@{}l@{}} \forall a, x. \\ \quad (\operatorname{is\text{-}unital\text{-}involutive\text{-}banach\text{-}algebra}(a) \wedge \operatorname{is\text{-}invertible\text{-}el}(a, x)) \Rightarrow  \\ \quad \quad \cdot_{a}(x, {x}^{-1}) = 1_{a} \end{array}}}
\vnbprop{bcp-u-bae-invertible-intro}{\vnbstmt{\begin{array}{@{}l@{}} \forall a.; \forall x, u \in \operatorname{vec}(a). \\ \quad (\operatorname{is\text{-}unital\text{-}involutive\text{-}banach\text{-}algebra}(a) \wedge ((\cdot_{a}(u, x) = 1_{a}) \wedge (\cdot_{a}(x, u) = 1_{a}))) \Rightarrow  \\ \quad \quad \operatorname{is\text{-}invertible\text{-}el}(a, x) \end{array}}}
\vnbprop{bcp-u-bae-inverse-eq}{\vnbstmt{\begin{array}{@{}l@{}} \forall a.; \forall x, u \in \operatorname{vec}(a). \\ \quad (\operatorname{is\text{-}unital\text{-}involutive\text{-}banach\text{-}algebra}(a) \wedge ((\cdot_{a}(u, x) = 1_{a}) \wedge (\cdot_{a}(x, u) = 1_{a}))) \Rightarrow  \\ \quad \quad {x}^{-1} = u \end{array}}}
\vnbprop{bcp-u-bae-inv-unique}{\vnbstmt{\begin{array}{@{}l@{}} \forall a.; \forall x, u, w \in \operatorname{vec}(a). \\ \quad (\operatorname{is\text{-}unital\text{-}involutive\text{-}banach\text{-}algebra}(a) \wedge ((\cdot_{a}(u, x) = 1_{a}) \wedge (\cdot_{a}(x, w) = 1_{a}))) \Rightarrow  \\ \quad \quad u = w \end{array}}}
\vnbprop{bcp-u-bae-one-invertible}{\vnbstmt{\begin{array}{@{}l@{}} \forall a. \\ \quad \operatorname{is\text{-}unital\text{-}involutive\text{-}banach\text{-}algebra}(a) \Rightarrow \operatorname{is\text{-}invertible\text{-}el}(a, 1_{a}) \end{array}}}
\vnbprop{bcp-u-bae-act-type}{\vnbstmt{\begin{array}{@{}l@{}} \forall a.; \forall z \in \mathbb{C}, x \in \operatorname{vec}(a). \\ \quad \operatorname{is\text{-}unital\text{-}involutive\text{-}banach\text{-}algebra}(a) \Rightarrow (\operatorname{act}(a)(z, x) \in \operatorname{vec}(a)) \end{array}}}
\vnbprop{bcp-u-bal-uba-vnrm-nonneg}{\vnbstmt{\begin{array}{@{}l@{}} \forall a.; \forall x \in \operatorname{vec}(a). \\ \quad \operatorname{is\text{-}unital\text{-}involutive\text{-}banach\text{-}algebra}(a) \Rightarrow (0 \leq \operatorname{vnrm}(a)(x)) \end{array}}}
\vnbprop{bcp-u-bal-uba-vnrm-triangle}{\vnbstmt{\begin{array}{@{}l@{}} \forall a.; \forall x, y \in \operatorname{vec}(a). \\ \quad \operatorname{is\text{-}unital\text{-}involutive\text{-}banach\text{-}algebra}(a) \Rightarrow  \\ \quad \quad \operatorname{vnrm}(a)(\operatorname{vadd}(a)(x, y)) \leq (\operatorname{vnrm}(a)(x) + \operatorname{vnrm}(a)(y)) \end{array}}}
\vnbprop{bcp-u-bal-uba-vnrm-homog}{\vnbstmt{\begin{array}{@{}l@{}} \forall a.; \forall r \in \operatorname{carr}(\operatorname{scal}(a)), x \in \operatorname{vec}(a). \\ \quad \operatorname{is\text{-}unital\text{-}involutive\text{-}banach\text{-}algebra}(a) \Rightarrow  \\ \quad \quad \operatorname{vnrm}(a)(\operatorname{act}(a)(r, x)) = (\operatorname{magnitude}(r) \cdot \operatorname{vnrm}(a)(x)) \end{array}}}
\vnbprop{bcp-u-bal-uba-vnrm-submult}{\vnbstmt{\begin{array}{@{}l@{}} \forall a.; \forall x, y \in \operatorname{vec}(a). \\ \quad \operatorname{is\text{-}unital\text{-}involutive\text{-}banach\text{-}algebra}(a) \Rightarrow  \\ \quad \quad \operatorname{vnrm}(a)(\cdot_{a}(x, y)) \leq (\operatorname{vnrm}(a)(x) \cdot \operatorname{vnrm}(a)(y)) \end{array}}}
\vnbprop{bcp-u-bal-uba-vone-norm}{\vnbstmt{\begin{array}{@{}l@{}} \forall a. \\ \quad \operatorname{is\text{-}unital\text{-}involutive\text{-}banach\text{-}algebra}(a) \Rightarrow (\operatorname{vnrm}(a)(1_{a}) = 1) \end{array}}}
\vnbprop{bcp-u-bal-uba-vmul-act-left}{\vnbstmt{\begin{array}{@{}l@{}} \forall a.; \forall r \in \operatorname{carr}(\operatorname{scal}(a)), x, y \in \operatorname{vec}(a). \\ \quad \operatorname{is\text{-}unital\text{-}involutive\text{-}banach\text{-}algebra}(a) \Rightarrow  \\ \quad \quad \cdot_{a}(\operatorname{act}(a)(r, x), y) = \operatorname{act}(a)(r, \cdot_{a}(x, y)) \end{array}}}
\vnbprop{bcp-u-bal-uba-vmul-act-right}{\vnbstmt{\begin{array}{@{}l@{}} \forall a.; \forall r \in \operatorname{carr}(\operatorname{scal}(a)), x, y \in \operatorname{vec}(a). \\ \quad \operatorname{is\text{-}unital\text{-}involutive\text{-}banach\text{-}algebra}(a) \Rightarrow  \\ \quad \quad \cdot_{a}(x, \operatorname{act}(a)(r, y)) = \operatorname{act}(a)(r, \cdot_{a}(x, y)) \end{array}}}
\vnbprop{bcp-u-bal-uba-act-one}{\vnbstmt{\begin{array}{@{}l@{}} \forall a.; \forall x \in \operatorname{vec}(a). \\ \quad \operatorname{is\text{-}unital\text{-}involutive\text{-}banach\text{-}algebra}(a) \Rightarrow (\operatorname{act}(a)(\operatorname{one}(\operatorname{scal}(a)), x) = x) \end{array}}}
\vnbprop{bcp-u-bas-shift-type}{\vnbstmt{\begin{array}{@{}l@{}} \forall a.; \forall z \in \mathbb{C}, x \in \operatorname{vec}(a). \\ \quad \operatorname{is\text{-}unital\text{-}involutive\text{-}banach\text{-}algebra}(a) \Rightarrow (\operatorname{shift\text{-}el}(a, z, x) \in \operatorname{vec}(a)) \end{array}}}
\vnbprop{bcp-u-bas-shift-diff}{\vnbstmt{\begin{array}{@{}l@{}} \forall a.; \forall u, z \in \mathbb{C}, x \in \operatorname{vec}(a). \\ \quad \operatorname{is\text{-}unital\text{-}involutive\text{-}banach\text{-}algebra}(a) \Rightarrow  \\ \quad \quad \operatorname{vadd}(a)(\operatorname{shift\text{-}el}(a, u, x), \operatorname{vneg}(a)(\operatorname{shift\text{-}el}(a, z, x))) = \operatorname{act}(a)((u + -z), 1_{a}) \end{array}}}
\vnbprop{bcp-u-bas-act-zero}{\vnbstmt{\begin{array}{@{}l@{}} \forall a.; \forall v \in \operatorname{vec}(a). \\ \quad \operatorname{is\text{-}unital\text{-}involutive\text{-}banach\text{-}algebra}(a) \Rightarrow (\operatorname{act}(a)(0, v) = \operatorname{vzero}(a)) \end{array}}}
\vnbprop{bcp-u-bas-act-one}{\vnbstmt{\begin{array}{@{}l@{}} \forall a.; \forall v \in \operatorname{vec}(a). \\ \quad \operatorname{is\text{-}unital\text{-}involutive\text{-}banach\text{-}algebra}(a) \Rightarrow (\operatorname{act}(a)(1, v) = v) \end{array}}}
\vnbprop{bcp-u-bas-scal-neg-eq}{\vnbstmt{\begin{array}{@{}l@{}} \forall a.; \forall z \in \mathbb{C}. \\ \quad \operatorname{is\text{-}unital\text{-}involutive\text{-}banach\text{-}algebra}(a) \Rightarrow (\operatorname{neg}(\operatorname{scal}(a))(z) = -z) \end{array}}}
\vnbprop{bcp-u-inverse-near-bound}{\vnbstmt{\begin{array}{@{}l@{}} \forall a, b.; \forall v \in \operatorname{vec}(a). \\ \quad \operatorname{is\text{-}unital\text{-}involutive\text{-}banach\text{-}algebra}(a) \wedge \\ \quad \operatorname{is\text{-}invertible\text{-}el}(a, b) \wedge \\ \quad (\operatorname{vnrm}(a)(\operatorname{vadd}(a)(b, \operatorname{vneg}(a)(v))) \cdot \operatorname{vnrm}(a)({b}^{-1})) \leq 1/2 \Rightarrow \\ \quad \quad (\operatorname{is\text{-}invertible\text{-}el}(a, v) \wedge (\operatorname{vnrm}(a)({v}^{-1}) \leq (2 \cdot \operatorname{vnrm}(a)({b}^{-1})))) \end{array}}}
\vnbprop{bcp-u-invertible-near-invertible}{\vnbstmt{\begin{array}{@{}l@{}} \forall a, b.; \forall x \in \operatorname{vec}(a). \\ \quad \operatorname{is\text{-}unital\text{-}involutive\text{-}banach\text{-}algebra}(a) \wedge \\ \quad \operatorname{is\text{-}invertible\text{-}el}(a, b) \wedge \\ \quad \operatorname{vnrm}(a)(\operatorname{vadd}(a)(b, \operatorname{vneg}(a)(x))) < \operatorname{recip}(\operatorname{vnrm}(a)({b}^{-1})) \Rightarrow \\ \quad \quad \operatorname{is\text{-}invertible\text{-}el}(a, x) \end{array}}}
\vnbprop{bcp-u-ban-inverse-norm-pos}{\vnbstmt{\begin{array}{@{}l@{}} \forall a, b. \\ \quad (\operatorname{is\text{-}unital\text{-}involutive\text{-}banach\text{-}algebra}(a) \wedge \operatorname{is\text{-}invertible\text{-}el}(a, b)) \Rightarrow  \\ \quad \quad 0 < \operatorname{vnrm}(a)({b}^{-1}) \end{array}}}
\vnbprop{bcp-u-ban-neg-norm}{\vnbstmt{\begin{array}{@{}l@{}} \forall a.; \forall q \in \operatorname{vec}(a). \\ \quad \operatorname{is\text{-}unital\text{-}involutive\text{-}banach\text{-}algebra}(a) \Rightarrow  \\ \quad \quad \operatorname{vnrm}(a)(\operatorname{vneg}(a)(q)) = \operatorname{vnrm}(a)(q) \end{array}}}
\vnbprop{bcp-u-bal-uba-scal-pin}{\vnbstmt{\begin{array}{@{}l@{}} \forall a. \\ \quad \operatorname{is\text{-}unital\text{-}involutive\text{-}banach\text{-}algebra}(a) \Rightarrow  \\ \quad \quad \operatorname{scal}(a) = \operatorname{normed\text{-}field\text{-}as\text{-}commutative\text{-}ring}(\operatorname{cc\text{-}normed\text{-}field}) \end{array}}}
\vnbprop{bcp-invertible-view}{\vnbstmt{\begin{array}{@{}l@{}} \forall a, y. \\ \quad \operatorname{is\text{-}unital\text{-}involutive\text{-}banach\text{-}algebra}(a) \Rightarrow  \\ \quad \quad \operatorname{is\text{-}invertible\text{-}el}(\operatorname{uiba\text{-}as\text{-}uba}(a), y) \Leftrightarrow \operatorname{is\text{-}invertible\text{-}el}(a, y) \end{array}}}
\vnbprop{bcp-shift-view}{\vnbstmt{\begin{array}{@{}l@{}} \forall a, z, x. \\ \quad \operatorname{is\text{-}unital\text{-}involutive\text{-}banach\text{-}algebra}(a) \Rightarrow  \\ \quad \quad (\operatorname{shift\text{-}el}(\operatorname{uiba\text{-}as\text{-}uba}(a), z, x) \simeq \operatorname{shift\text{-}el}(a, z, x)) \end{array}}}
\vnbprop{bcp-inverse-view}{\vnbstmt{\begin{array}{@{}l@{}} \forall a, y. \\ \quad \operatorname{is\text{-}unital\text{-}involutive\text{-}banach\text{-}algebra}(a) \Rightarrow  \\ \quad \quad ({y}^{-1} \simeq {y}^{-1}) \end{array}}}
\vnbprop{bcp-u-shift-invertible-outside-ball}{\vnbstmt{\begin{array}{@{}l@{}} \forall a.; \forall x \in \operatorname{vec}(a), z \in \mathbb{C}. \\ \quad (\operatorname{is\text{-}unital\text{-}involutive\text{-}banach\text{-}algebra}(a) \wedge (\operatorname{vnrm}(a)(x) < \operatorname{magnitude}(z))) \Rightarrow  \\ \quad \quad \operatorname{is\text{-}invertible\text{-}el}(a, \operatorname{shift\text{-}el}(a, z, x)) \end{array}}}
\vnbprop{bcp-u-bae-invertible-mul}{\vnbstmt{\begin{array}{@{}l@{}} \forall a, x, y. \\ \quad (\operatorname{is\text{-}unital\text{-}involutive\text{-}banach\text{-}algebra}(a) \wedge (\operatorname{is\text{-}invertible\text{-}el}(a, x) \wedge \operatorname{is\text{-}invertible\text{-}el}(a, y))) \Rightarrow  \\ \quad \quad \operatorname{is\text{-}invertible\text{-}el}(a, \cdot_{a}(x, y)) \end{array}}}
\vnbprop{bcp-star-fun}{\vnbstmt{\begin{array}{@{}l@{}} \forall a. \\ \quad \operatorname{is\text{-}unital\text{-}involutive\text{-}banach\text{-}algebra}(a) \Rightarrow ({}^{*_{a}} \in (\operatorname{vec}(a) \to \operatorname{vec}(a))) \end{array}}}
\vnbprop{bcp-star-add}{\vnbstmt{\begin{array}{@{}l@{}} \forall a.; \forall x, y \in \operatorname{vec}(a). \\ \quad \operatorname{is\text{-}unital\text{-}involutive\text{-}banach\text{-}algebra}(a) \Rightarrow  \\ \quad \quad {}^{*_{a}}(\operatorname{vadd}(a)(x, y)) = \operatorname{vadd}(a)({}^{*_{a}}(x), {}^{*_{a}}(y)) \end{array}}}
\vnbprop{bcp-star-act}{\vnbstmt{\begin{array}{@{}l@{}} \forall a.; \forall r \in \mathbb{C}, x \in \operatorname{vec}(a). \\ \quad \operatorname{is\text{-}unital\text{-}involutive\text{-}banach\text{-}algebra}(a) \Rightarrow  \\ \quad \quad {}^{*_{a}}(\operatorname{act}(a)(r, x)) = \operatorname{act}(a)(\operatorname{conjugate}(r), {}^{*_{a}}(x)) \end{array}}}
\vnbprop{bcp-star-star}{\vnbstmt{\begin{array}{@{}l@{}} \forall a.; \forall x \in \operatorname{vec}(a). \\ \quad \operatorname{is\text{-}unital\text{-}involutive\text{-}banach\text{-}algebra}(a) \Rightarrow ({}^{*_{a}}({}^{*_{a}}(x)) = x) \end{array}}}
\vnbprop{bcp-star-mul}{\vnbstmt{\begin{array}{@{}l@{}} \forall a.; \forall x, y \in \operatorname{vec}(a). \\ \quad \operatorname{is\text{-}unital\text{-}involutive\text{-}banach\text{-}algebra}(a) \Rightarrow  \\ \quad \quad {}^{*_{a}}(\cdot_{a}(x, y)) = \cdot_{a}({}^{*_{a}}(y), {}^{*_{a}}(x)) \end{array}}}
\vnbprop{bcp-star-norm}{\vnbstmt{\begin{array}{@{}l@{}} \forall a.; \forall x \in \operatorname{vec}(a). \\ \quad \operatorname{is\text{-}unital\text{-}involutive\text{-}banach\text{-}algebra}(a) \Rightarrow  \\ \quad \quad \operatorname{vnrm}(a)({}^{*_{a}}(x)) = \operatorname{vnrm}(a)(x) \end{array}}}
\vnbprop{bcp-ucsa-uiba}{\vnbstmt{\begin{array}{@{}l@{}} \forall a. \\ \quad \operatorname{is\text{-}unital\text{-}c\text{-}star\text{-}algebra}(a) \Rightarrow \operatorname{is\text{-}unital\text{-}involutive\text{-}banach\text{-}algebra}(a) \end{array}}}
\vnbprop{bcp-ucsa-cstar}{\vnbstmt{\begin{array}{@{}l@{}} \forall a.; \forall x \in \operatorname{vec}(a). \\ \quad \operatorname{is\text{-}unital\text{-}c\text{-}star\text{-}algebra}(a) \Rightarrow  \\ \quad \quad \operatorname{vnrm}(a)(\cdot_{a}({}^{*_{a}}(x), x)) = (\operatorname{vnrm}(a)(x) \cdot \operatorname{vnrm}(a)(x)) \end{array}}}
\vnbprop{bcp-star-one}{\vnbstmt{\begin{array}{@{}l@{}} \forall a. \\ \quad \operatorname{is\text{-}unital\text{-}involutive\text{-}banach\text{-}algebra}(a) \Rightarrow ({}^{*_{a}}(1_{a}) = 1_{a}) \end{array}}}
\vnbprop{bcp-star-zero}{\vnbstmt{\begin{array}{@{}l@{}} \forall a. \\ \quad \operatorname{is\text{-}unital\text{-}involutive\text{-}banach\text{-}algebra}(a) \Rightarrow ({}^{*_{a}}(\operatorname{vzero}(a)) = \operatorname{vzero}(a)) \end{array}}}
\vnbprop{bcp-star-neg}{\vnbstmt{\begin{array}{@{}l@{}} \forall a.; \forall x \in \operatorname{vec}(a). \\ \quad \operatorname{is\text{-}unital\text{-}involutive\text{-}banach\text{-}algebra}(a) \Rightarrow  \\ \quad \quad {}^{*_{a}}(\operatorname{vneg}(a)(x)) = \operatorname{vneg}(a)({}^{*_{a}}(x)) \end{array}}}
\vnbprop{bcp-shift-translate}{\vnbstmt{\begin{array}{@{}l@{}} \forall a.; \forall z, w \in \mathbb{C}, x \in \operatorname{vec}(a). \\ \quad \operatorname{is\text{-}unital\text{-}involutive\text{-}banach\text{-}algebra}(a) \Rightarrow  \\ \quad \quad \operatorname{shift\text{-}el}(a, (z + w), \operatorname{vadd}(a)(x, \operatorname{act}(a)(w, 1_{a}))) = \operatorname{shift\text{-}el}(a, z, x) \end{array}}}
\vnbprop{bcp-shift-zero-neg}{\vnbstmt{\begin{array}{@{}l@{}} \forall a.; \forall y \in \operatorname{vec}(a). \\ \quad \operatorname{is\text{-}unital\text{-}involutive\text{-}banach\text{-}algebra}(a) \Rightarrow (\operatorname{shift\text{-}el}(a, 0, \operatorname{vneg}(a)(y)) = y) \end{array}}}
\vnbprop{bcp-quad-expand}{\vnbstmt{\begin{array}{@{}l@{}} \forall a.; \forall x \in \operatorname{vec}(a), t \in \mathbb{R}. \\ \quad \operatorname{is\text{-}unital\text{-}involutive\text{-}banach\text{-}algebra}(a) \Rightarrow  \\ \quad \quad \cdot_{a}(\operatorname{vadd}(a)(x, \operatorname{act}(a)((-+i \cdot t), 1_{a})), \operatorname{vadd}(a)(x, \operatorname{act}(a)((+i \cdot t), 1_{a}))) \\ \quad \quad =\; \operatorname{vadd}(a)(\cdot_{a}(x, x), \operatorname{act}(a)((t \cdot t), 1_{a})) \end{array}}}
\vnbprop{bcp-sa-noninvertible-real}{\vnbstmt{\begin{array}{@{}l@{}} \forall a, h.; \forall z \in \mathbb{C}. \\ \quad \operatorname{is\text{-}unital\text{-}c\text{-}star\text{-}algebra}(a) \wedge \\ \quad \operatorname{is\text{-}self\text{-}adjoint\text{-}el}(a, h) \wedge \\ \quad \neg \operatorname{is\text{-}invertible\text{-}el}(a, \operatorname{shift\text{-}el}(a, z, h)) \Rightarrow \\ \quad \quad z \in \mathbb{R} \end{array}}}
\vnbprop{bcp-sa-shift-invertible}{\vnbstmt{\begin{array}{@{}l@{}} \forall a, h.; \forall z \in \mathbb{C}. \\ \quad (\operatorname{is\text{-}unital\text{-}c\text{-}star\text{-}algebra}(a) \wedge (\operatorname{is\text{-}self\text{-}adjoint\text{-}el}(a, h) \wedge \neg (z \in \mathbb{R}))) \Rightarrow  \\ \quad \quad \operatorname{is\text{-}invertible\text{-}el}(a, \operatorname{shift\text{-}el}(a, z, h)) \end{array}}}
\vnbprop{bcp-sa-spectrum-real}{\vnbstmt{\begin{array}{@{}l@{}} \forall a, h.; \forall z \in \sigma({h}). \\ \quad (\operatorname{is\text{-}unital\text{-}c\text{-}star\text{-}algebra}(a) \wedge \operatorname{is\text{-}self\text{-}adjoint\text{-}el}(a, h)) \Rightarrow (z \in \mathbb{R}) \end{array}}}
\vnbprop{bcp-star-invertible}{\vnbstmt{\begin{array}{@{}l@{}} \forall a, y. \\ \quad (\operatorname{is\text{-}unital\text{-}involutive\text{-}banach\text{-}algebra}(a) \wedge \operatorname{is\text{-}invertible\text{-}el}(a, y)) \Rightarrow  \\ \quad \quad \operatorname{is\text{-}invertible\text{-}el}(a, {}^{*_{a}}(y)) \end{array}}}
\vnbprop{bcp-sub-ucsa-a}{\vnbstmt{\begin{array}{@{}l@{}} \forall a, b. \\ \quad \operatorname{is\text{-}unital\text{-}c\text{-}star\text{-}subalgebra}(a, b) \Rightarrow \operatorname{is\text{-}unital\text{-}c\text{-}star\text{-}algebra}(a) \end{array}}}
\vnbprop{bcp-sub-ucsa-b}{\vnbstmt{\begin{array}{@{}l@{}} \forall a, b. \\ \quad \operatorname{is\text{-}unital\text{-}c\text{-}star\text{-}subalgebra}(a, b) \Rightarrow \operatorname{is\text{-}unital\text{-}c\text{-}star\text{-}algebra}(b) \end{array}}}
\vnbprop{bcp-sub-vec}{\vnbstmt{\begin{array}{@{}l@{}} \forall a, b.; \forall x \in \operatorname{vec}(a). \\ \quad \operatorname{is\text{-}unital\text{-}c\text{-}star\text{-}subalgebra}(a, b) \Rightarrow (x \in \operatorname{vec}(b)) \end{array}}}
\vnbprop{bcp-sub-one}{\vnbstmt{\forall a, b.\; (\operatorname{is\text{-}unital\text{-}c\text{-}star\text{-}subalgebra}(a, b) \Rightarrow (1_{a} = 1_{b}))}}
\vnbprop{bcp-sub-vadd}{\vnbstmt{\begin{array}{@{}l@{}} \forall a, b.; \forall x, y \in \operatorname{vec}(a). \\ \quad \operatorname{is\text{-}unital\text{-}c\text{-}star\text{-}subalgebra}(a, b) \Rightarrow (\operatorname{vadd}(a)(x, y) = \operatorname{vadd}(b)(x, y)) \end{array}}}
\vnbprop{bcp-sub-vmul}{\vnbstmt{\begin{array}{@{}l@{}} \forall a, b.; \forall x, y \in \operatorname{vec}(a). \\ \quad \operatorname{is\text{-}unital\text{-}c\text{-}star\text{-}subalgebra}(a, b) \Rightarrow (\cdot_{a}(x, y) = \cdot_{b}(x, y)) \end{array}}}
\vnbprop{bcp-sub-act}{\vnbstmt{\begin{array}{@{}l@{}} \forall a, b.; \forall z \in \mathbb{C}, x \in \operatorname{vec}(a). \\ \quad \operatorname{is\text{-}unital\text{-}c\text{-}star\text{-}subalgebra}(a, b) \Rightarrow (\operatorname{act}(a)(z, x) = \operatorname{act}(b)(z, x)) \end{array}}}
\vnbprop{bcp-sub-vneg}{\vnbstmt{\begin{array}{@{}l@{}} \forall a, b.; \forall x \in \operatorname{vec}(a). \\ \quad \operatorname{is\text{-}unital\text{-}c\text{-}star\text{-}subalgebra}(a, b) \Rightarrow (\operatorname{vneg}(a)(x) = \operatorname{vneg}(b)(x)) \end{array}}}
\vnbprop{bcp-sub-vstar}{\vnbstmt{\begin{array}{@{}l@{}} \forall a, b.; \forall x \in \operatorname{vec}(a). \\ \quad \operatorname{is\text{-}unital\text{-}c\text{-}star\text{-}subalgebra}(a, b) \Rightarrow ({}^{*_{a}}(x) = {}^{*_{b}}(x)) \end{array}}}
\vnbprop{bcp-sub-vnrm}{\vnbstmt{\begin{array}{@{}l@{}} \forall a, b.; \forall x \in \operatorname{vec}(a). \\ \quad \operatorname{is\text{-}unital\text{-}c\text{-}star\text{-}subalgebra}(a, b) \Rightarrow (\operatorname{vnrm}(a)(x) = \operatorname{vnrm}(b)(x)) \end{array}}}
\vnbprop{bcp-sub-shift}{\vnbstmt{\begin{array}{@{}l@{}} \forall a, b.; \forall z \in \mathbb{C}, x \in \operatorname{vec}(a). \\ \quad \operatorname{is\text{-}unital\text{-}c\text{-}star\text{-}subalgebra}(a, b) \Rightarrow (\operatorname{shift\text{-}el}(a, z, x) = \operatorname{shift\text{-}el}(b, z, x)) \end{array}}}
\vnbprop{bcp-sub-invertible-up}{\vnbstmt{\begin{array}{@{}l@{}} \forall a, b, y. \\ \quad (\operatorname{is\text{-}unital\text{-}c\text{-}star\text{-}subalgebra}(a, b) \wedge \operatorname{is\text{-}invertible\text{-}el}(a, y)) \Rightarrow  \\ \quad \quad \operatorname{is\text{-}invertible\text{-}el}(b, y) \end{array}}}
\vnbprop{bcp-sub-inverse-up}{\vnbstmt{\begin{array}{@{}l@{}} \forall a, b, y. \\ \quad (\operatorname{is\text{-}unital\text{-}c\text{-}star\text{-}subalgebra}(a, b) \wedge \operatorname{is\text{-}invertible\text{-}el}(a, y)) \Rightarrow  \\ \quad \quad {y}^{-1} = {y}^{-1} \end{array}}}
\vnbprop{bcp-sub-sa-shift-down}{\vnbstmt{\begin{array}{@{}l@{}} \forall a, b, h.; \forall z \in \mathbb{R}. \\ \quad \operatorname{is\text{-}unital\text{-}c\text{-}star\text{-}subalgebra}(a, b) \wedge \\ \quad \operatorname{is\text{-}self\text{-}adjoint\text{-}el}(a, h) \wedge \\ \quad \operatorname{is\text{-}invertible\text{-}el}(b, \operatorname{shift\text{-}el}(b, z, h)) \Rightarrow \\ \quad \quad \operatorname{is\text{-}invertible\text{-}el}(a, \operatorname{shift\text{-}el}(a, z, h)) \end{array}}}
\vnbprop{bcp-sub-sa-invertible-down}{\vnbstmt{\begin{array}{@{}l@{}} \forall a, b, h. \\ \quad (\operatorname{is\text{-}unital\text{-}c\text{-}star\text{-}subalgebra}(a, b) \wedge (\operatorname{is\text{-}self\text{-}adjoint\text{-}el}(a, h) \wedge \operatorname{is\text{-}invertible\text{-}el}(b, h))) \Rightarrow  \\ \quad \quad \operatorname{is\text{-}invertible\text{-}el}(a, h) \end{array}}}
\vnbprop{bcp-sub-invertible-down}{\vnbstmt{\begin{array}{@{}l@{}} \forall a, b.; \forall y \in \operatorname{vec}(a). \\ \quad (\operatorname{is\text{-}unital\text{-}c\text{-}star\text{-}subalgebra}(a, b) \wedge \operatorname{is\text{-}invertible\text{-}el}(b, y)) \Rightarrow  \\ \quad \quad \operatorname{is\text{-}invertible\text{-}el}(a, y) \end{array}}}
\vnbprop{bcp-sub-invertible-iff}{\vnbstmt{\begin{array}{@{}l@{}} \forall a, b.; \forall y \in \operatorname{vec}(a). \\ \quad \operatorname{is\text{-}unital\text{-}c\text{-}star\text{-}subalgebra}(a, b) \Rightarrow  \\ \quad \quad \operatorname{is\text{-}invertible\text{-}el}(a, y) \Leftrightarrow \operatorname{is\text{-}invertible\text{-}el}(b, y) \end{array}}}
\vnbprop{spectral-permanence}{\vnbstmt{\begin{array}{@{}l@{}} \forall a, b.; \forall x \in \operatorname{vec}(a). \\ \quad \operatorname{is\text{-}unital\text{-}c\text{-}star\text{-}subalgebra}(a, b) \Rightarrow (\sigma({x}) = \sigma({x})) \end{array}}}

\vnbbookfile{theorem-library/ba2b-gen}{theorem-library/ba2b-gen}
\noindent{\small ba2b-gen.scm -- wave of 2026-11-09/10 (agent BA-2B), wired 2026-11-10 from scratchpad/ba2b-1109/ba2b-gen.scm. Library: operator-algebras.}\par
\vnbprop{b2n-vec-closed-star}{\vnbstmt{\begin{array}{@{}l@{}} \forall a. \\ \quad \operatorname{is\text{-}unital\text{-}involutive\text{-}banach\text{-}algebra}(a) \Rightarrow  \\ \quad \quad \operatorname{is\text{-}closed\text{-}star\text{-}subalgebra\text{-}set}(a, \operatorname{vec}(a)) \end{array}}}
\vnbprop{cstar-gen-contains}{\vnbstmt{\begin{array}{@{}l@{}} \forall a.; \forall x \in \operatorname{vec}(a). \\ \quad \operatorname{is\text{-}unital\text{-}involutive\text{-}banach\text{-}algebra}(a) \Rightarrow (x \in \operatorname{cstar\text{-}gen}(a, x)) \end{array}}}
\vnbprop{cstar-gen-least}{\vnbstmt{\begin{array}{@{}l@{}} \forall a, x, s. \\ \quad (\operatorname{is\text{-}unital\text{-}involutive\text{-}banach\text{-}algebra}(a) \wedge (\operatorname{is\text{-}closed\text{-}star\text{-}subalgebra\text{-}set}(a, s) \wedge (x \in s))) \Rightarrow  \\ \quad \quad \operatorname{cstar\text{-}gen}(a, x) \subseteq s \end{array}}}
\vnbprop{cstar-gen-subset}{\vnbstmt{\begin{array}{@{}l@{}} \forall a.; \forall x \in \operatorname{vec}(a). \\ \quad \operatorname{is\text{-}unital\text{-}involutive\text{-}banach\text{-}algebra}(a) \Rightarrow (\operatorname{cstar\text{-}gen}(a, x) \subseteq \operatorname{vec}(a)) \end{array}}}

\vnbbookfile{theorem-library/basw2-cx}{theorem-library/basw2-cx}
\noindent{\small basw2-cx.scm -- wave of 2026-11-09/10 (agent BA-SW2), wired 2026-11-10 from scratchpad/basw2-1109/basw2-cx.scm. Library: operator-algebras.}\par
\vnbprop{ccont-unfold}{\vnbstmt{\begin{array}{@{}l@{}} \forall t. \\ \quad (C({t}) \simeq \{f \in (\operatorname{pts}(t) \to \mathbb{C}) : \operatorname{is\text{-}hom\text{-}top\text{-}space}(t, \operatorname{metric\text{-}top}(\operatorname{cc\text{-}ms}), f)\}) \end{array}}}
\vnbprop{ccont-parts}{\vnbstmt{\begin{array}{@{}l@{}} \forall t.; \forall f \in C({t}). \\ \quad ((f \in (\operatorname{pts}(t) \to \mathbb{C})) \wedge (\operatorname{is\text{-}top\text{-}space}(t) \wedge \operatorname{is\text{-}hom\text{-}top\text{-}space}(t, \operatorname{metric\text{-}top}(\operatorname{cc\text{-}ms}), f))) \end{array}}}
\vnbprop{ccont-local}{\vnbstmt{\begin{array}{@{}l@{}} \forall t.; \forall f \in C({t}), p \in \operatorname{pts}(t), s. \\ \quad \operatorname{pos\text{-}rr}(s) \Rightarrow  \\ \quad \quad \exists u \in \operatorname{opens}(t). \\ \quad \quad \quad (p \in u) \wedge \forall q \in u.\; (\operatorname{magnitude}((f(p) - f(q))) < s) \end{array}}}
\vnbprop{ccont-intro}{\vnbstmt{\begin{array}{@{}l@{}} \forall t, f. \\ \quad \operatorname{is\text{-}top\text{-}space}(t) \wedge \\ \quad f \in (\operatorname{pts}(t) \to \mathbb{C}) \wedge \\ \quad \forall p \in \operatorname{pts}(t), s. \\ \quad \quad \operatorname{pos\text{-}rr}(s) \Rightarrow  \\ \quad \quad \quad \exists u \in \operatorname{opens}(t). \\ \quad \quad \quad \quad (p \in u) \wedge \forall q \in u.\; (\operatorname{magnitude}((f(p) - f(q))) < s) \Rightarrow \\ \quad \quad f \in C({t}) \end{array}}}
\vnbprop{ccont-compose}{\vnbstmt{\begin{array}{@{}l@{}} \forall t, f, g, h. \\ \quad f \in C({t}) \wedge \\ \quad g \in C({t}) \wedge \\ \quad h \in (\operatorname{pts}(t) \to \mathbb{C}) \wedge \\ \quad \forall p \in \operatorname{pts}(t), e. \\ \quad \quad \operatorname{pos\text{-}rr}(e) \Rightarrow  \\ \quad \quad \quad \exists d. \\ \quad \quad \quad \quad \operatorname{pos\text{-}rr}(d) \wedge \\ \quad \quad \quad \quad \forall q \in \operatorname{pts}(t). \\ \quad \quad \quad \quad \quad ((\operatorname{magnitude}((f(p) - f(q))) < d) \wedge (\operatorname{magnitude}((g(p) - g(q))) < d)) \Rightarrow  \\ \quad \quad \quad \quad \quad \quad \operatorname{magnitude}((h(p) - h(q))) < e \Rightarrow \\ \quad \quad h \in C({t}) \end{array}}}
\vnbprop{cc-add-near}{\vnbstmt{\begin{array}{@{}l@{}} \forall a, b \in \mathbb{C}, s. \\ \quad \operatorname{pos\text{-}rr}(s) \Rightarrow  \\ \quad \quad \exists d. \\ \quad \quad \quad \operatorname{pos\text{-}rr}(d) \wedge \\ \quad \quad \quad \forall x, y \in \mathbb{C}. \\ \quad \quad \quad \quad ((\operatorname{magnitude}((a - x)) < d) \wedge (\operatorname{magnitude}((b - y)) < d)) \Rightarrow  \\ \quad \quad \quad \quad \quad \operatorname{magnitude}(((a + b) - (x + y))) < s \end{array}}}
\vnbprop{cc-mul-near}{\vnbstmt{\begin{array}{@{}l@{}} \forall a, b \in \mathbb{C}, s. \\ \quad \operatorname{pos\text{-}rr}(s) \Rightarrow  \\ \quad \quad \exists d. \\ \quad \quad \quad \operatorname{pos\text{-}rr}(d) \wedge \\ \quad \quad \quad \forall x, y \in \mathbb{C}. \\ \quad \quad \quad \quad ((\operatorname{magnitude}((a - x)) < d) \wedge (\operatorname{magnitude}((b - y)) < d)) \Rightarrow  \\ \quad \quad \quad \quad \quad \operatorname{magnitude}(((a \cdot b) - (x \cdot y))) < s \end{array}}}
\vnbprop{ccont-compose-fun}{\vnbstmt{\begin{array}{@{}l@{}} \forall t, p, f, g, h. \\ \quad f \in C({t}) \wedge \\ \quad g \in C({t}) \wedge \\ \quad h \in (\operatorname{pts}(t) \to \mathbb{C}) \wedge \\ \quad \forall z \in \operatorname{pts}(t).\; (h(z) = p(f(z), g(z))) \wedge \\ \quad \forall a, b \in \mathbb{C}, s. \\ \quad \quad \operatorname{pos\text{-}rr}(s) \Rightarrow  \\ \quad \quad \quad \exists d. \\ \quad \quad \quad \quad \operatorname{pos\text{-}rr}(d) \wedge \\ \quad \quad \quad \quad \forall x, y \in \mathbb{C}. \\ \quad \quad \quad \quad \quad ((\operatorname{magnitude}((a - x)) < d) \wedge (\operatorname{magnitude}((b - y)) < d)) \Rightarrow  \\ \quad \quad \quad \quad \quad \quad \operatorname{magnitude}((p(a, b) - p(x, y))) < s \Rightarrow \\ \quad \quad h \in C({t}) \end{array}}}
\vnbprop{cc-neg-near}{\vnbstmt{\begin{array}{@{}l@{}} \forall a \in \mathbb{C}, s. \\ \quad \operatorname{pos\text{-}rr}(s) \Rightarrow  \\ \quad \quad \exists d. \\ \quad \quad \quad (\operatorname{pos\text{-}rr}(d) \wedge \forall x \in \mathbb{C}.\; ((\operatorname{magnitude}((a - x)) < d) \Rightarrow (\operatorname{magnitude}((-a - -x)) < s))) \end{array}}}
\vnbprop{cc-conj-near}{\vnbstmt{\begin{array}{@{}l@{}} \forall a \in \mathbb{C}, s. \\ \quad \operatorname{pos\text{-}rr}(s) \Rightarrow  \\ \quad \quad \exists d. \\ \quad \quad \quad \operatorname{pos\text{-}rr}(d) \wedge \\ \quad \quad \quad \forall x \in \mathbb{C}. \\ \quad \quad \quad \quad (\operatorname{magnitude}((a - x)) < d) \Rightarrow (\operatorname{magnitude}((\operatorname{conjugate}(a) - \operatorname{conjugate}(x))) < s) \end{array}}}
\vnbprop{ccont-add-ptw}{\vnbstmt{\begin{array}{@{}l@{}} \forall t, f, g, h. \\ \quad f \in C({t}) \wedge \\ \quad g \in C({t}) \wedge \\ \quad h \in (\operatorname{pts}(t) \to \mathbb{C}) \wedge \\ \quad \forall z \in \operatorname{pts}(t).\; (h(z) = (f(z) + g(z))) \Rightarrow \\ \quad \quad h \in C({t}) \end{array}}}
\vnbprop{ccont-mul-ptw}{\vnbstmt{\begin{array}{@{}l@{}} \forall t, f, g, h. \\ \quad f \in C({t}) \wedge \\ \quad g \in C({t}) \wedge \\ \quad h \in (\operatorname{pts}(t) \to \mathbb{C}) \wedge \\ \quad \forall z \in \operatorname{pts}(t).\; (h(z) = (f(z) \cdot g(z))) \Rightarrow \\ \quad \quad h \in C({t}) \end{array}}}
\vnbprop{ccont-neg-ptw}{\vnbstmt{\begin{array}{@{}l@{}} \forall t, f, h. \\ \quad ((f \in C({t})) \wedge ((h \in (\operatorname{pts}(t) \to \mathbb{C})) \wedge \forall z \in \operatorname{pts}(t).\; (h(z) = -f(z)))) \Rightarrow  \\ \quad \quad h \in C({t}) \end{array}}}
\vnbprop{ccont-conj-ptw}{\vnbstmt{\begin{array}{@{}l@{}} \forall t, f, h. \\ \quad ((f \in C({t})) \wedge ((h \in (\operatorname{pts}(t) \to \mathbb{C})) \wedge \forall z \in \operatorname{pts}(t).\; (h(z) = \operatorname{conjugate}(f(z))))) \Rightarrow  \\ \quad \quad h \in C({t}) \end{array}}}
\vnbprop{ccont-scal-ptw}{\vnbstmt{\begin{array}{@{}l@{}} \forall t, c, f, h. \\ \quad ((c \in \mathbb{C}) \wedge ((f \in C({t})) \wedge ((h \in (\operatorname{pts}(t) \to \mathbb{C})) \wedge \forall z \in \operatorname{pts}(t).\; (h(z) = (c \cdot f(z)))))) \Rightarrow  \\ \quad \quad h \in C({t}) \end{array}}}
\vnbprop{ccont-const-ptw}{\vnbstmt{\begin{array}{@{}l@{}} \forall t, c, h. \\ \quad (\operatorname{is\text{-}top\text{-}space}(t) \wedge ((c \in \mathbb{C}) \wedge ((h \in (\operatorname{pts}(t) \to \mathbb{C})) \wedge \forall z \in \operatorname{pts}(t).\; (h(z) = c)))) \Rightarrow  \\ \quad \quad h \in C({t}) \end{array}}}
\vnbprop{ccont-is-set}{\vnbstmt{\forall t.\; (\operatorname{is\text{-}top\text{-}space}(t) \Rightarrow (C({t}) \in \mathbf{Set}))}}
\vnbprop{ccont-vadd-type}{\vnbstmt{\begin{array}{@{}l@{}} \forall t. \\ \quad \operatorname{is\text{-}top\text{-}space}(t) \Rightarrow  \\ \quad \quad \operatorname{ccont\text{-}vadd}(t) \in ((C({t}) \times C({t})) \to C({t})) \end{array}}}
\vnbprop{ccont-vmul-type}{\vnbstmt{\begin{array}{@{}l@{}} \forall t. \\ \quad \operatorname{is\text{-}top\text{-}space}(t) \Rightarrow  \\ \quad \quad \operatorname{ccont\text{-}vmul}(t) \in ((C({t}) \times C({t})) \to C({t})) \end{array}}}
\vnbprop{ccont-vneg-type}{\vnbstmt{\forall t.\; (\operatorname{is\text{-}top\text{-}space}(t) \Rightarrow (\operatorname{ccont\text{-}vneg}(t) \in (C({t}) \to C({t}))))}}
\vnbprop{ccont-vstar-type}{\vnbstmt{\forall t.\; (\operatorname{is\text{-}top\text{-}space}(t) \Rightarrow (\operatorname{ccont\text{-}vstar}(t) \in (C({t}) \to C({t}))))}}
\vnbprop{ccont-act-type}{\vnbstmt{\begin{array}{@{}l@{}} \forall t. \\ \quad \operatorname{is\text{-}top\text{-}space}(t) \Rightarrow (\operatorname{ccont\text{-}act}(t) \in ((\mathbb{C} \times C({t})) \to C({t}))) \end{array}}}
\vnbprop{ccont-vzero-type}{\vnbstmt{\forall t.\; (\operatorname{is\text{-}top\text{-}space}(t) \Rightarrow (\operatorname{ccont\text{-}vzero}(t) \in C({t})))}}
\vnbprop{ccont-vone-type}{\vnbstmt{\forall t.\; (\operatorname{is\text{-}top\text{-}space}(t) \Rightarrow (\operatorname{ccont\text{-}vone}(t) \in C({t})))}}
\vnbprop{ccont-magnitude-rcont}{\vnbstmt{\begin{array}{@{}l@{}} \forall t, f, h. \\ \quad ((f \in C({t})) \wedge ((h \in (\operatorname{pts}(t) \to \mathbb{R})) \wedge \forall z \in \operatorname{pts}(t).\; (h(z) = \operatorname{magnitude}(f(z))))) \Rightarrow  \\ \quad \quad h \in C_{\mathbb{R}}({t}) \end{array}}}
\vnbprop{ccont-bounded}{\vnbstmt{\begin{array}{@{}l@{}} \forall t, f. \\ \quad (\operatorname{is\text{-}compact\text{-}t}(t) \wedge (f \in C({t}))) \Rightarrow  \\ \quad \quad \exists m \in \mathbb{R}.\; ((0 \leq m) \wedge \forall z \in \operatorname{pts}(t).\; (\operatorname{magnitude}(f(z)) \leq m)) \end{array}}}
\vnbprop{sup-norm-t-unfold}{\vnbstmt{\forall t, f.\; (\lVert f \rVert_{\infty} \simeq \operatorname{cc\text{-}sup\text{-}norm}(f, \operatorname{pts}(t)))}}
\vnbprop{sup-norm-t-upper}{\vnbstmt{\begin{array}{@{}l@{}} \forall t.; \forall f \in (\operatorname{pts}(t) \to \mathbb{C}), x \in \operatorname{pts}(t). \\ \quad \operatorname{magnitude}(f(x)) \leq \lVert f \rVert_{\infty} \end{array}}}
\vnbprop{sup-norm-t-least}{\vnbstmt{\begin{array}{@{}l@{}} \forall t.; \forall f \in (\operatorname{pts}(t) \to \mathbb{C}), m \in \mathbb{R}. \\ \quad ((0 \leq m) \wedge \forall x \in \operatorname{pts}(t).\; (\operatorname{magnitude}(f(x)) \leq m)) \Rightarrow  \\ \quad \quad (\lVert f \rVert_{\infty} \in \mathbb{R}) \wedge ((0 \leq \lVert f \rVert_{\infty}) \wedge (\lVert f \rVert_{\infty} \leq m)) \end{array}}}
\vnbprop{sup-norm-t-real}{\vnbstmt{\begin{array}{@{}l@{}} \forall t, f. \\ \quad (\operatorname{is\text{-}compact\text{-}t}(t) \wedge (f \in C({t}))) \Rightarrow  \\ \quad \quad (\lVert f \rVert_{\infty} \in \mathbb{R}) \wedge (0 \leq \lVert f \rVert_{\infty}) \end{array}}}
\vnbprop{ccont-vnrm-type}{\vnbstmt{\forall t.\; (\operatorname{is\text{-}compact\text{-}t}(t) \Rightarrow (\operatorname{ccont\text{-}vnrm}(t) \in (C({t}) \to \mathbb{R})))}}
\vnbprop{sup-norm-t-add-le}{\vnbstmt{\begin{array}{@{}l@{}} \forall t, f, g, h. \\ \quad \operatorname{is\text{-}compact\text{-}t}(t) \wedge \\ \quad f \in C({t}) \wedge \\ \quad g \in C({t}) \wedge \\ \quad h \in (\operatorname{pts}(t) \to \mathbb{C}) \wedge \\ \quad \forall z \in \operatorname{pts}(t).\; (h(z) = (f(z) + g(z))) \Rightarrow \\ \quad \quad \lVert h \rVert_{\infty} \leq (\lVert f \rVert_{\infty} + \lVert g \rVert_{\infty}) \end{array}}}
\vnbprop{sup-norm-t-mul-le}{\vnbstmt{\begin{array}{@{}l@{}} \forall t, f, g, h. \\ \quad \operatorname{is\text{-}compact\text{-}t}(t) \wedge \\ \quad f \in C({t}) \wedge \\ \quad g \in C({t}) \wedge \\ \quad h \in (\operatorname{pts}(t) \to \mathbb{C}) \wedge \\ \quad \forall z \in \operatorname{pts}(t).\; (h(z) = (f(z) \cdot g(z))) \Rightarrow \\ \quad \quad \lVert h \rVert_{\infty} \leq (\lVert f \rVert_{\infty} \cdot \lVert g \rVert_{\infty}) \end{array}}}
\vnbprop{sup-norm-t-scal-le}{\vnbstmt{\begin{array}{@{}l@{}} \forall t, c, f, h. \\ \quad \operatorname{is\text{-}compact\text{-}t}(t) \wedge \\ \quad c \in \mathbb{C} \wedge \\ \quad f \in C({t}) \wedge \\ \quad h \in (\operatorname{pts}(t) \to \mathbb{C}) \wedge \\ \quad \forall z \in \operatorname{pts}(t).\; (h(z) = (c \cdot f(z))) \Rightarrow \\ \quad \quad \lVert h \rVert_{\infty} \leq (\operatorname{magnitude}(c) \cdot \lVert f \rVert_{\infty}) \end{array}}}
\vnbprop{sup-norm-t-scal}{\vnbstmt{\begin{array}{@{}l@{}} \forall t, c, f, h. \\ \quad \operatorname{is\text{-}compact\text{-}t}(t) \wedge \\ \quad c \in \mathbb{C} \wedge \\ \quad f \in C({t}) \wedge \\ \quad h \in (\operatorname{pts}(t) \to \mathbb{C}) \wedge \\ \quad \forall z \in \operatorname{pts}(t).\; (h(z) = (c \cdot f(z))) \Rightarrow \\ \quad \quad \lVert h \rVert_{\infty} = (\operatorname{magnitude}(c) \cdot \lVert f \rVert_{\infty}) \end{array}}}
\vnbprop{sup-norm-t-conj}{\vnbstmt{\begin{array}{@{}l@{}} \forall t, f, h. \\ \quad \operatorname{is\text{-}compact\text{-}t}(t) \wedge \\ \quad f \in C({t}) \wedge \\ \quad h \in (\operatorname{pts}(t) \to \mathbb{C}) \wedge \\ \quad \forall z \in \operatorname{pts}(t).\; (h(z) = \operatorname{conjugate}(f(z))) \Rightarrow \\ \quad \quad \lVert h \rVert_{\infty} = \lVert f \rVert_{\infty} \end{array}}}
\vnbprop{sup-norm-t-cstar}{\vnbstmt{\begin{array}{@{}l@{}} \forall t, f, h. \\ \quad \operatorname{is\text{-}compact\text{-}t}(t) \wedge \\ \quad f \in C({t}) \wedge \\ \quad h \in (\operatorname{pts}(t) \to \mathbb{C}) \wedge \\ \quad \forall z \in \operatorname{pts}(t).\; (h(z) = (\operatorname{conjugate}(f(z)) \cdot f(z))) \Rightarrow \\ \quad \quad \lVert h \rVert_{\infty} = (\lVert f \rVert_{\infty} \cdot \lVert f \rVert_{\infty}) \end{array}}}
\vnbprop{sup-norm-t-zero-elim}{\vnbstmt{\begin{array}{@{}l@{}} \forall t, f.; \forall z \in \operatorname{pts}(t). \\ \quad (\operatorname{is\text{-}compact\text{-}t}(t) \wedge ((f \in C({t})) \wedge (\lVert f \rVert_{\infty} = 0))) \Rightarrow (f(z) = 0) \end{array}}}
\vnbprop{sup-norm-t-zero-intro}{\vnbstmt{\begin{array}{@{}l@{}} \forall t, h. \\ \quad (\operatorname{is\text{-}compact\text{-}t}(t) \wedge ((h \in (\operatorname{pts}(t) \to \mathbb{C})) \wedge \forall z \in \operatorname{pts}(t).\; (h(z) = 0))) \Rightarrow  \\ \quad \quad \lVert h \rVert_{\infty} = 0 \end{array}}}
\vnbprop{sup-norm-t-one}{\vnbstmt{\begin{array}{@{}l@{}} \forall t, h. \\ \quad \operatorname{is\text{-}compact\text{-}t}(t) \wedge \\ \quad \exists p.\; (p \in \operatorname{pts}(t)) \wedge \\ \quad h \in (\operatorname{pts}(t) \to \mathbb{C}) \wedge \\ \quad \forall z \in \operatorname{pts}(t).\; (h(z) = 1) \Rightarrow \\ \quad \quad \lVert h \rVert_{\infty} = 1 \end{array}}}
\vnbprop{ccont-pointwise-limit}{\vnbstmt{\begin{array}{@{}l@{}} \forall t, s.; \forall z \in \operatorname{pts}(t). \\ \quad \operatorname{is\text{-}top\text{-}space}(t) \wedge \\ \quad s \in (\mathbb{N} \to C({t})) \wedge \\ \quad \forall a. \\ \quad \quad \operatorname{pos\text{-}rr}(a) \Rightarrow  \\ \quad \quad \quad \exists p \in \mathbb{N}. \forall i, j \in \mathbb{N}, z \in \operatorname{pts}(t). \\ \quad \quad \quad \quad ((p \leq i) \wedge (p \leq j)) \Rightarrow (\operatorname{magnitude}((s(i)(z) - s(j)(z))) < a) \Rightarrow \\ \quad \quad \exists l \in \mathbb{C}.\; \operatorname{converges\text{-}to}(\operatorname{cc\text{-}ms}, (\lambda \mathit{cxi} \in \mathbb{N}.\; s(\mathit{cxi})(z)), l) \end{array}}}
\vnbprop{ccont-limit-function}{\vnbstmt{\begin{array}{@{}l@{}} \forall t, s. \\ \quad \operatorname{is\text{-}top\text{-}space}(t) \wedge \\ \quad s \in (\mathbb{N} \to C({t})) \wedge \\ \quad \forall a. \\ \quad \quad \operatorname{pos\text{-}rr}(a) \Rightarrow  \\ \quad \quad \quad \exists p \in \mathbb{N}. \forall i, j \in \mathbb{N}, z \in \operatorname{pts}(t). \\ \quad \quad \quad \quad ((p \leq i) \wedge (p \leq j)) \Rightarrow (\operatorname{magnitude}((s(i)(z) - s(j)(z))) < a) \Rightarrow \\ \quad \quad (\lambda \mathit{cxz} \in \operatorname{pts}(t).\; \\ \quad \quad \quad \varepsilon(\{l \in \mathbb{C} : \operatorname{converges\text{-}to}(\operatorname{cc\text{-}ms}, (\lambda \mathit{cxi} \in \mathbb{N}.\; s(\mathit{cxi})(\mathit{cxz})), l)\})) \\ \quad \quad \in\; (\operatorname{pts}(t) \to \mathbb{C}) \wedge \\ \quad \quad \forall z \in \operatorname{pts}(t). \\ \quad \quad \quad \operatorname{converges\text{-}to}(\operatorname{cc\text{-}ms}, (\lambda \mathit{cxi} \in \mathbb{N}.\; s(\mathit{cxi})(z)), \\ \quad \quad \quad \quad (\lambda \mathit{cxz} \in \operatorname{pts}(t).\; \varepsilon(\{l \in \mathbb{C} : \operatorname{converges\text{-}to}(\operatorname{cc\text{-}ms}, (\lambda \mathit{cxi} \in \mathbb{N}.\; s(\mathit{cxi})(\mathit{cxz})), l)\}))(z)) \end{array}}}
\vnbprop{ccont-limit-uniform}{\vnbstmt{\begin{array}{@{}l@{}} \forall t, s, a. \\ \quad \operatorname{is\text{-}top\text{-}space}(t) \wedge \\ \quad s \in (\mathbb{N} \to C({t})) \wedge \\ \quad \forall a. \\ \quad \quad \operatorname{pos\text{-}rr}(a) \Rightarrow  \\ \quad \quad \quad \exists p \in \mathbb{N}. \forall i, j \in \mathbb{N}, z \in \operatorname{pts}(t). \\ \quad \quad \quad \quad ((p \leq i) \wedge (p \leq j)) \Rightarrow (\operatorname{magnitude}((s(i)(z) - s(j)(z))) < a) \wedge \\ \quad \operatorname{pos\text{-}rr}(a) \Rightarrow \\ \quad \quad \exists p \in \mathbb{N}. \forall i \in \mathbb{N}, z \in \operatorname{pts}(t). \\ \quad \quad \quad (p \leq i) \Rightarrow  \\ \quad \quad \quad \quad \operatorname{magnitude}((s(i)(z) - (\lambda \mathit{cxz} \in \operatorname{pts}(t).\; \varepsilon(\{l \in \mathbb{C} : \operatorname{converges\text{-}to}(\operatorname{cc\text{-}ms}, (\lambda \mathit{cxi} \in \mathbb{N}.\; s(\mathit{cxi})(\mathit{cxz})), l)\}))(z))) \\ \quad \quad \quad \quad \leq\; a \end{array}}}
\vnbprop{cc-four-eps}{\vnbstmt{\begin{array}{@{}l@{}} \forall a, b, c, w \in \mathbb{C}, r, d, e \in \mathbb{R}. \\ \quad 0 < r \wedge \\ \quad (r + r) = d \wedge \\ \quad (d + d) = e \wedge \\ \quad \operatorname{magnitude}((b - c)) < r \wedge \\ \quad \operatorname{magnitude}((b - a)) \leq r \wedge \\ \quad \operatorname{magnitude}((c - w)) \leq r \Rightarrow \\ \quad \quad \operatorname{magnitude}((a - w)) < e \end{array}}}
\vnbprop{ccont-limit-continuous}{\vnbstmt{\begin{array}{@{}l@{}} \forall t, s. \\ \quad \operatorname{is\text{-}top\text{-}space}(t) \wedge \\ \quad s \in (\mathbb{N} \to C({t})) \wedge \\ \quad \forall a. \\ \quad \quad \operatorname{pos\text{-}rr}(a) \Rightarrow  \\ \quad \quad \quad \exists p \in \mathbb{N}. \forall i, j \in \mathbb{N}, z \in \operatorname{pts}(t). \\ \quad \quad \quad \quad ((p \leq i) \wedge (p \leq j)) \Rightarrow (\operatorname{magnitude}((s(i)(z) - s(j)(z))) < a) \Rightarrow \\ \quad \quad (\lambda \mathit{cxz} \in \operatorname{pts}(t).\; \\ \quad \quad \quad \varepsilon(\{l \in \mathbb{C} : \operatorname{converges\text{-}to}(\operatorname{cc\text{-}ms}, (\lambda \mathit{cxi} \in \mathbb{N}.\; s(\mathit{cxi})(\mathit{cxz})), l)\})) \\ \quad \quad \in\; C({t}) \end{array}}}
\vnbprop{ccont-uniformly-complete}{\vnbstmt{\begin{array}{@{}l@{}} \forall t, s. \\ \quad \operatorname{is\text{-}top\text{-}space}(t) \wedge \\ \quad s \in (\mathbb{N} \to C({t})) \wedge \\ \quad \forall a. \\ \quad \quad \operatorname{pos\text{-}rr}(a) \Rightarrow  \\ \quad \quad \quad \exists p \in \mathbb{N}. \forall i, j \in \mathbb{N}, z \in \operatorname{pts}(t). \\ \quad \quad \quad \quad ((p \leq i) \wedge (p \leq j)) \Rightarrow (\operatorname{magnitude}((s(i)(z) - s(j)(z))) < a) \Rightarrow \\ \quad \quad \exists g \in C({t}). \forall a. \\ \quad \quad \quad \operatorname{pos\text{-}rr}(a) \Rightarrow  \\ \quad \quad \quad \quad \exists p \in \mathbb{N}. \forall i \in \mathbb{N}, z \in \operatorname{pts}(t). \\ \quad \quad \quad \quad \quad (p \leq i) \Rightarrow (\operatorname{magnitude}((s(i)(z) - g(z))) \leq a) \end{array}}}
\vnbprop{ccont-vadd-in}{\vnbstmt{\begin{array}{@{}l@{}} \forall t.; \forall f, g \in C({t}). \\ \quad \operatorname{is\text{-}top\text{-}space}(t) \Rightarrow (\operatorname{ccont\text{-}vadd}(t)(f, g) \in C({t})) \end{array}}}
\vnbprop{ccont-vadd-val}{\vnbstmt{\begin{array}{@{}l@{}} \forall t.; \forall f, g \in C({t}), p \in \operatorname{pts}(t). \\ \quad \operatorname{is\text{-}top\text{-}space}(t) \Rightarrow (\operatorname{ccont\text{-}vadd}(t)(f, g)(p) = (f(p) + g(p))) \end{array}}}
\vnbprop{ccont-vmul-in}{\vnbstmt{\begin{array}{@{}l@{}} \forall t.; \forall f, g \in C({t}). \\ \quad \operatorname{is\text{-}top\text{-}space}(t) \Rightarrow (\operatorname{ccont\text{-}vmul}(t)(f, g) \in C({t})) \end{array}}}
\vnbprop{ccont-vmul-val}{\vnbstmt{\begin{array}{@{}l@{}} \forall t.; \forall f, g \in C({t}), p \in \operatorname{pts}(t). \\ \quad \operatorname{is\text{-}top\text{-}space}(t) \Rightarrow (\operatorname{ccont\text{-}vmul}(t)(f, g)(p) = (f(p) \cdot g(p))) \end{array}}}
\vnbprop{ccont-vneg-in}{\vnbstmt{\begin{array}{@{}l@{}} \forall t.; \forall f \in C({t}). \\ \quad \operatorname{is\text{-}top\text{-}space}(t) \Rightarrow (\operatorname{ccont\text{-}vneg}(t)(f) \in C({t})) \end{array}}}
\vnbprop{ccont-vneg-val}{\vnbstmt{\begin{array}{@{}l@{}} \forall t.; \forall f \in C({t}), p \in \operatorname{pts}(t). \\ \quad \operatorname{is\text{-}top\text{-}space}(t) \Rightarrow (\operatorname{ccont\text{-}vneg}(t)(f)(p) = -f(p)) \end{array}}}
\vnbprop{ccont-vstar-in}{\vnbstmt{\begin{array}{@{}l@{}} \forall t.; \forall f \in C({t}). \\ \quad \operatorname{is\text{-}top\text{-}space}(t) \Rightarrow (\operatorname{ccont\text{-}vstar}(t)(f) \in C({t})) \end{array}}}
\vnbprop{ccont-vstar-val}{\vnbstmt{\begin{array}{@{}l@{}} \forall t.; \forall f \in C({t}), p \in \operatorname{pts}(t). \\ \quad \operatorname{is\text{-}top\text{-}space}(t) \Rightarrow (\operatorname{ccont\text{-}vstar}(t)(f)(p) = \operatorname{conjugate}(f(p))) \end{array}}}
\vnbprop{ccont-act-in}{\vnbstmt{\begin{array}{@{}l@{}} \forall t.; \forall c \in \mathbb{C}, f \in C({t}). \\ \quad \operatorname{is\text{-}top\text{-}space}(t) \Rightarrow (\operatorname{ccont\text{-}act}(t)(c, f) \in C({t})) \end{array}}}
\vnbprop{ccont-act-val}{\vnbstmt{\begin{array}{@{}l@{}} \forall t.; \forall c \in \mathbb{C}, f \in C({t}), p \in \operatorname{pts}(t). \\ \quad \operatorname{is\text{-}top\text{-}space}(t) \Rightarrow (\operatorname{ccont\text{-}act}(t)(c, f)(p) = (c \cdot f(p))) \end{array}}}
\vnbprop{ccont-vzero-val}{\vnbstmt{\forall t.; \forall p \in \operatorname{pts}(t).\; (\operatorname{is\text{-}top\text{-}space}(t) \Rightarrow (\operatorname{ccont\text{-}vzero}(t)(p) = 0))}}
\vnbprop{ccont-vone-val}{\vnbstmt{\forall t.; \forall p \in \operatorname{pts}(t).\; (\operatorname{is\text{-}top\text{-}space}(t) \Rightarrow (\operatorname{ccont\text{-}vone}(t)(p) = 1))}}
\vnbprop{ccont-vnrm-val}{\vnbstmt{\begin{array}{@{}l@{}} \forall t.; \forall f \in C({t}). \\ \quad \operatorname{is\text{-}compact\text{-}t}(t) \Rightarrow (\operatorname{ccont\text{-}vnrm}(t)(f) = \lVert f \rVert_{\infty}) \end{array}}}
\vnbprop{ccont-vadd-assoc}{\vnbstmt{\forall t.\; (\operatorname{is\text{-}top\text{-}space}(t) \Rightarrow \operatorname{is\text{-}associative}(\operatorname{ccont\text{-}vadd}(t), C({t})))}}
\vnbprop{ccont-vadd-comm}{\vnbstmt{\forall t.\; (\operatorname{is\text{-}top\text{-}space}(t) \Rightarrow \operatorname{is\text{-}commutative}(\operatorname{ccont\text{-}vadd}(t), C({t})))}}
\vnbprop{ccont-vadd-ident}{\vnbstmt{\begin{array}{@{}l@{}} \forall t. \\ \quad \operatorname{is\text{-}top\text{-}space}(t) \Rightarrow \operatorname{is\text{-}identity}(\operatorname{ccont\text{-}vadd}(t), \operatorname{ccont\text{-}vzero}(t), C({t})) \end{array}}}
\vnbprop{ccont-vadd-inv}{\vnbstmt{\begin{array}{@{}l@{}} \forall t. \\ \quad \operatorname{is\text{-}top\text{-}space}(t) \Rightarrow  \\ \quad \quad \operatorname{has\text{-}inverses}(\operatorname{ccont\text{-}vadd}(t), \operatorname{ccont\text{-}vzero}(t), \operatorname{ccont\text{-}vneg}(t), C({t})) \end{array}}}
\vnbprop{ccont-vmul-assoc}{\vnbstmt{\forall t.\; (\operatorname{is\text{-}top\text{-}space}(t) \Rightarrow \operatorname{is\text{-}associative}(\operatorname{ccont\text{-}vmul}(t), C({t})))}}
\vnbprop{ccont-act-vadd}{\vnbstmt{\begin{array}{@{}l@{}} \forall t.; \forall r \in \operatorname{carr}(\operatorname{normed\text{-}field\text{-}as\text{-}commutative\text{-}ring}(\operatorname{cc\text{-}normed\text{-}field})), x, y \in C({t}). \\ \quad \operatorname{is\text{-}top\text{-}space}(t) \Rightarrow  \\ \quad \quad \operatorname{ccont\text{-}act}(t)(r, \operatorname{ccont\text{-}vadd}(t)(x, y)) \\ \quad \quad =\; \operatorname{ccont\text{-}vadd}(t)(\operatorname{ccont\text{-}act}(t)(r, x), \operatorname{ccont\text{-}act}(t)(r, y)) \end{array}}}
\vnbprop{ccont-act-add}{\vnbstmt{\begin{array}{@{}l@{}} \forall t.; \forall r, s \in \operatorname{carr}(\operatorname{normed\text{-}field\text{-}as\text{-}commutative\text{-}ring}(\operatorname{cc\text{-}normed\text{-}field})), x \in C({t}). \\ \quad \operatorname{is\text{-}top\text{-}space}(t) \Rightarrow  \\ \quad \quad \operatorname{ccont\text{-}act}(t)(\operatorname{add}(\operatorname{normed\text{-}field\text{-}as\text{-}commutative\text{-}ring}(\operatorname{cc\text{-}normed\text{-}field}))(r, s), x) \\ \quad \quad =\; \operatorname{ccont\text{-}vadd}(t)(\operatorname{ccont\text{-}act}(t)(r, x), \operatorname{ccont\text{-}act}(t)(s, x)) \end{array}}}
\vnbprop{ccont-act-mul}{\vnbstmt{\begin{array}{@{}l@{}} \forall t.; \forall r, s \in \operatorname{carr}(\operatorname{normed\text{-}field\text{-}as\text{-}commutative\text{-}ring}(\operatorname{cc\text{-}normed\text{-}field})), x \in C({t}). \\ \quad \operatorname{is\text{-}top\text{-}space}(t) \Rightarrow  \\ \quad \quad \operatorname{ccont\text{-}act}(t)(\operatorname{mul}(\operatorname{normed\text{-}field\text{-}as\text{-}commutative\text{-}ring}(\operatorname{cc\text{-}normed\text{-}field}))(r, s), x) \\ \quad \quad =\; \operatorname{ccont\text{-}act}(t)(r, \operatorname{ccont\text{-}act}(t)(s, x)) \end{array}}}
\vnbprop{ccont-act-one}{\vnbstmt{\begin{array}{@{}l@{}} \forall t.; \forall x \in C({t}). \\ \quad \operatorname{is\text{-}top\text{-}space}(t) \Rightarrow  \\ \quad \quad \operatorname{ccont\text{-}act}(t)(\operatorname{one}(\operatorname{normed\text{-}field\text{-}as\text{-}commutative\text{-}ring}(\operatorname{cc\text{-}normed\text{-}field})), x) = x \end{array}}}
\vnbprop{ccont-vmul-vadd-right}{\vnbstmt{\begin{array}{@{}l@{}} \forall t.; \forall x, y, z \in C({t}). \\ \quad \operatorname{is\text{-}top\text{-}space}(t) \Rightarrow  \\ \quad \quad \operatorname{ccont\text{-}vmul}(t)(x, \operatorname{ccont\text{-}vadd}(t)(y, z)) \\ \quad \quad =\; \operatorname{ccont\text{-}vadd}(t)(\operatorname{ccont\text{-}vmul}(t)(x, y), \operatorname{ccont\text{-}vmul}(t)(x, z)) \end{array}}}
\vnbprop{ccont-vmul-vadd-left}{\vnbstmt{\begin{array}{@{}l@{}} \forall t.; \forall x, y, z \in C({t}). \\ \quad \operatorname{is\text{-}top\text{-}space}(t) \Rightarrow  \\ \quad \quad \operatorname{ccont\text{-}vmul}(t)(\operatorname{ccont\text{-}vadd}(t)(x, y), z) \\ \quad \quad =\; \operatorname{ccont\text{-}vadd}(t)(\operatorname{ccont\text{-}vmul}(t)(x, z), \operatorname{ccont\text{-}vmul}(t)(y, z)) \end{array}}}
\vnbprop{ccont-vmul-act-left}{\vnbstmt{\begin{array}{@{}l@{}} \forall t.; \forall r \in \operatorname{carr}(\operatorname{normed\text{-}field\text{-}as\text{-}commutative\text{-}ring}(\operatorname{cc\text{-}normed\text{-}field})), x, y \in C({t}). \\ \quad \operatorname{is\text{-}top\text{-}space}(t) \Rightarrow  \\ \quad \quad \operatorname{ccont\text{-}vmul}(t)(\operatorname{ccont\text{-}act}(t)(r, x), y) = \operatorname{ccont\text{-}act}(t)(r, \operatorname{ccont\text{-}vmul}(t)(x, y)) \end{array}}}
\vnbprop{ccont-vmul-act-right}{\vnbstmt{\begin{array}{@{}l@{}} \forall t.; \forall r \in \operatorname{carr}(\operatorname{normed\text{-}field\text{-}as\text{-}commutative\text{-}ring}(\operatorname{cc\text{-}normed\text{-}field})), x, y \in C({t}). \\ \quad \operatorname{is\text{-}top\text{-}space}(t) \Rightarrow  \\ \quad \quad \operatorname{ccont\text{-}vmul}(t)(x, \operatorname{ccont\text{-}act}(t)(r, y)) = \operatorname{ccont\text{-}act}(t)(r, \operatorname{ccont\text{-}vmul}(t)(x, y)) \end{array}}}
\vnbprop{ccont-vnrm-nonneg}{\vnbstmt{\forall t.; \forall x \in C({t}).\; (\operatorname{is\text{-}compact\text{-}t}(t) \Rightarrow (0 \leq \operatorname{ccont\text{-}vnrm}(t)(x)))}}
\vnbprop{ccont-vnrm-zero-iff}{\vnbstmt{\begin{array}{@{}l@{}} \forall t.; \forall x \in C({t}). \\ \quad \operatorname{is\text{-}compact\text{-}t}(t) \Rightarrow ((\operatorname{ccont\text{-}vnrm}(t)(x) = 0) \Leftrightarrow (x = \operatorname{ccont\text{-}vzero}(t))) \end{array}}}
\vnbprop{ccont-vnrm-act}{\vnbstmt{\begin{array}{@{}l@{}} \forall t.; \forall r \in \operatorname{carr}(\operatorname{normed\text{-}field\text{-}as\text{-}commutative\text{-}ring}(\operatorname{cc\text{-}normed\text{-}field})), x \in C({t}). \\ \quad \operatorname{is\text{-}compact\text{-}t}(t) \Rightarrow  \\ \quad \quad \operatorname{ccont\text{-}vnrm}(t)(\operatorname{ccont\text{-}act}(t)(r, x)) = (\operatorname{magnitude}(r) \cdot \operatorname{ccont\text{-}vnrm}(t)(x)) \end{array}}}
\vnbprop{ccont-vnrm-triangle}{\vnbstmt{\begin{array}{@{}l@{}} \forall t.; \forall x, y \in C({t}). \\ \quad \operatorname{is\text{-}compact\text{-}t}(t) \Rightarrow  \\ \quad \quad \operatorname{ccont\text{-}vnrm}(t)(\operatorname{ccont\text{-}vadd}(t)(x, y)) \leq (\operatorname{ccont\text{-}vnrm}(t)(x) + \operatorname{ccont\text{-}vnrm}(t)(y)) \end{array}}}
\vnbprop{ccont-vnrm-submult}{\vnbstmt{\begin{array}{@{}l@{}} \forall t.; \forall x, y \in C({t}). \\ \quad \operatorname{is\text{-}compact\text{-}t}(t) \Rightarrow  \\ \quad \quad \operatorname{ccont\text{-}vnrm}(t)(\operatorname{ccont\text{-}vmul}(t)(x, y)) \leq (\operatorname{ccont\text{-}vnrm}(t)(x) \cdot \operatorname{ccont\text{-}vnrm}(t)(y)) \end{array}}}
\vnbprop{ccont-algebra-is-cna}{\vnbstmt{\begin{array}{@{}l@{}} \forall t. \\ \quad \operatorname{is\text{-}compact\text{-}t}(t) \Rightarrow  \\ \quad \quad \operatorname{is\text{-}complex\text{-}normed\text{-}algebra}({[}\operatorname{normed\text{-}field\text{-}as\text{-}commutative\text{-}ring}(\operatorname{cc\text{-}normed\text{-}field}), C({t}), \operatorname{ccont\text{-}vadd}(t), \operatorname{ccont\text{-}vzero}(t), \operatorname{ccont\text{-}vneg}(t), \operatorname{ccont\text{-}act}(t), \operatorname{ccont\text{-}vnrm}(t), \operatorname{ccont\text{-}vmul}(t)]) \end{array}}}
\vnbprop{ccont-tuple8-scal}{\vnbstmt{\begin{array}{@{}l@{}} \forall t. \\ \quad \operatorname{is\text{-}compact\text{-}t}(t) \Rightarrow  \\ \quad \quad (\operatorname{scal}({[}\operatorname{normed\text{-}field\text{-}as\text{-}commutative\text{-}ring}(\operatorname{cc\text{-}normed\text{-}field}), C({t}), \operatorname{ccont\text{-}vadd}(t), \operatorname{ccont\text{-}vzero}(t), \operatorname{ccont\text{-}vneg}(t), \operatorname{ccont\text{-}act}(t), \operatorname{ccont\text{-}vnrm}(t), \operatorname{ccont\text{-}vmul}(t)]) \simeq \operatorname{normed\text{-}field\text{-}as\text{-}commutative\text{-}ring}(\operatorname{cc\text{-}normed\text{-}field})) \end{array}}}
\vnbprop{ccont-tuple8-vec}{\vnbstmt{\begin{array}{@{}l@{}} \forall t. \\ \quad \operatorname{is\text{-}compact\text{-}t}(t) \Rightarrow  \\ \quad \quad (\operatorname{vec}({[}\operatorname{normed\text{-}field\text{-}as\text{-}commutative\text{-}ring}(\operatorname{cc\text{-}normed\text{-}field}), C({t}), \operatorname{ccont\text{-}vadd}(t), \operatorname{ccont\text{-}vzero}(t), \operatorname{ccont\text{-}vneg}(t), \operatorname{ccont\text{-}act}(t), \operatorname{ccont\text{-}vnrm}(t), \operatorname{ccont\text{-}vmul}(t)]) \simeq C({t})) \end{array}}}
\vnbprop{ccont-tuple8-vadd}{\vnbstmt{\begin{array}{@{}l@{}} \forall t. \\ \quad \operatorname{is\text{-}compact\text{-}t}(t) \Rightarrow  \\ \quad \quad (\operatorname{vadd}({[}\operatorname{normed\text{-}field\text{-}as\text{-}commutative\text{-}ring}(\operatorname{cc\text{-}normed\text{-}field}), C({t}), \operatorname{ccont\text{-}vadd}(t), \operatorname{ccont\text{-}vzero}(t), \operatorname{ccont\text{-}vneg}(t), \operatorname{ccont\text{-}act}(t), \operatorname{ccont\text{-}vnrm}(t), \operatorname{ccont\text{-}vmul}(t)]) \simeq \operatorname{ccont\text{-}vadd}(t)) \end{array}}}
\vnbprop{ccont-tuple8-vzero}{\vnbstmt{\begin{array}{@{}l@{}} \forall t. \\ \quad \operatorname{is\text{-}compact\text{-}t}(t) \Rightarrow  \\ \quad \quad (\operatorname{vzero}({[}\operatorname{normed\text{-}field\text{-}as\text{-}commutative\text{-}ring}(\operatorname{cc\text{-}normed\text{-}field}), C({t}), \operatorname{ccont\text{-}vadd}(t), \operatorname{ccont\text{-}vzero}(t), \operatorname{ccont\text{-}vneg}(t), \operatorname{ccont\text{-}act}(t), \operatorname{ccont\text{-}vnrm}(t), \operatorname{ccont\text{-}vmul}(t)]) \simeq \operatorname{ccont\text{-}vzero}(t)) \end{array}}}
\vnbprop{ccont-tuple8-vneg}{\vnbstmt{\begin{array}{@{}l@{}} \forall t. \\ \quad \operatorname{is\text{-}compact\text{-}t}(t) \Rightarrow  \\ \quad \quad (\operatorname{vneg}({[}\operatorname{normed\text{-}field\text{-}as\text{-}commutative\text{-}ring}(\operatorname{cc\text{-}normed\text{-}field}), C({t}), \operatorname{ccont\text{-}vadd}(t), \operatorname{ccont\text{-}vzero}(t), \operatorname{ccont\text{-}vneg}(t), \operatorname{ccont\text{-}act}(t), \operatorname{ccont\text{-}vnrm}(t), \operatorname{ccont\text{-}vmul}(t)]) \simeq \operatorname{ccont\text{-}vneg}(t)) \end{array}}}
\vnbprop{ccont-tuple8-act}{\vnbstmt{\begin{array}{@{}l@{}} \forall t. \\ \quad \operatorname{is\text{-}compact\text{-}t}(t) \Rightarrow  \\ \quad \quad (\operatorname{act}({[}\operatorname{normed\text{-}field\text{-}as\text{-}commutative\text{-}ring}(\operatorname{cc\text{-}normed\text{-}field}), C({t}), \operatorname{ccont\text{-}vadd}(t), \operatorname{ccont\text{-}vzero}(t), \operatorname{ccont\text{-}vneg}(t), \operatorname{ccont\text{-}act}(t), \operatorname{ccont\text{-}vnrm}(t), \operatorname{ccont\text{-}vmul}(t)]) \simeq \operatorname{ccont\text{-}act}(t)) \end{array}}}
\vnbprop{ccont-tuple8-vnrm}{\vnbstmt{\begin{array}{@{}l@{}} \forall t. \\ \quad \operatorname{is\text{-}compact\text{-}t}(t) \Rightarrow  \\ \quad \quad (\operatorname{vnrm}({[}\operatorname{normed\text{-}field\text{-}as\text{-}commutative\text{-}ring}(\operatorname{cc\text{-}normed\text{-}field}), C({t}), \operatorname{ccont\text{-}vadd}(t), \operatorname{ccont\text{-}vzero}(t), \operatorname{ccont\text{-}vneg}(t), \operatorname{ccont\text{-}act}(t), \operatorname{ccont\text{-}vnrm}(t), \operatorname{ccont\text{-}vmul}(t)]) \simeq \operatorname{ccont\text{-}vnrm}(t)) \end{array}}}
\vnbprop{ccont-tuple8-vmul}{\vnbstmt{\begin{array}{@{}l@{}} \forall t. \\ \quad \operatorname{is\text{-}compact\text{-}t}(t) \Rightarrow  \\ \quad \quad (\cdot_{{[}\operatorname{normed\text{-}field\text{-}as\text{-}commutative\text{-}ring}(\operatorname{cc\text{-}normed\text{-}field}), C({t}), \operatorname{ccont\text{-}vadd}(t), \operatorname{ccont\text{-}vzero}(t), \operatorname{ccont\text{-}vneg}(t), \operatorname{ccont\text{-}act}(t), \operatorname{ccont\text{-}vnrm}(t), \operatorname{ccont\text{-}vmul}(t)]} \simeq \operatorname{ccont\text{-}vmul}(t)) \end{array}}}
\vnbprop{ccont-algebra-dist}{\vnbstmt{\begin{array}{@{}l@{}} \forall t.; \forall u, v \in C({t}). \\ \quad \operatorname{is\text{-}compact\text{-}t}(t) \Rightarrow  \\ \quad \quad \operatorname{dist}(\operatorname{nag\text{-}metric\text{-}space}(\operatorname{cna\text{-}as\text{-}normed\text{-}ag}({[}\operatorname{normed\text{-}field\text{-}as\text{-}commutative\text{-}ring}(\operatorname{cc\text{-}normed\text{-}field}), C({t}), \operatorname{ccont\text{-}vadd}(t), \operatorname{ccont\text{-}vzero}(t), \operatorname{ccont\text{-}vneg}(t), \operatorname{ccont\text{-}act}(t), \operatorname{ccont\text{-}vnrm}(t), \operatorname{ccont\text{-}vmul}(t)])))(u, v) \\ \quad \quad =\; \lVert \operatorname{ccont\text{-}vadd}(t)(u, \operatorname{ccont\text{-}vneg}(t)(v)) \rVert_{\infty} \end{array}}}
\vnbprop{ccont-algebra-complete}{\vnbstmt{\begin{array}{@{}l@{}} \forall t. \\ \quad \operatorname{is\text{-}compact\text{-}t}(t) \Rightarrow  \\ \quad \quad \operatorname{is\text{-}complete}(\operatorname{nag\text{-}metric\text{-}space}(\operatorname{cna\text{-}as\text{-}normed\text{-}ag}({[}\operatorname{normed\text{-}field\text{-}as\text{-}commutative\text{-}ring}(\operatorname{cc\text{-}normed\text{-}field}), C({t}), \operatorname{ccont\text{-}vadd}(t), \operatorname{ccont\text{-}vzero}(t), \operatorname{ccont\text{-}vneg}(t), \operatorname{ccont\text{-}act}(t), \operatorname{ccont\text{-}vnrm}(t), \operatorname{ccont\text{-}vmul}(t)]))) \end{array}}}
\vnbprop{ccont-vmul-unit-left}{\vnbstmt{\begin{array}{@{}l@{}} \forall t.; \forall x \in C({t}). \\ \quad \operatorname{is\text{-}top\text{-}space}(t) \Rightarrow (\operatorname{ccont\text{-}vmul}(t)(\operatorname{ccont\text{-}vone}(t), x) = x) \end{array}}}
\vnbprop{ccont-vmul-unit-right}{\vnbstmt{\begin{array}{@{}l@{}} \forall t.; \forall x \in C({t}). \\ \quad \operatorname{is\text{-}top\text{-}space}(t) \Rightarrow (\operatorname{ccont\text{-}vmul}(t)(x, \operatorname{ccont\text{-}vone}(t)) = x) \end{array}}}
\vnbprop{ccont-vstar-vadd}{\vnbstmt{\begin{array}{@{}l@{}} \forall t.; \forall x, y \in C({t}). \\ \quad \operatorname{is\text{-}top\text{-}space}(t) \Rightarrow  \\ \quad \quad \operatorname{ccont\text{-}vstar}(t)(\operatorname{ccont\text{-}vadd}(t)(x, y)) = \operatorname{ccont\text{-}vadd}(t)(\operatorname{ccont\text{-}vstar}(t)(x), \operatorname{ccont\text{-}vstar}(t)(y)) \end{array}}}
\vnbprop{ccont-vstar-act}{\vnbstmt{\begin{array}{@{}l@{}} \forall t.; \forall r \in \mathbb{C}, x \in C({t}). \\ \quad \operatorname{is\text{-}top\text{-}space}(t) \Rightarrow  \\ \quad \quad \operatorname{ccont\text{-}vstar}(t)(\operatorname{ccont\text{-}act}(t)(r, x)) = \operatorname{ccont\text{-}act}(t)(\operatorname{conjugate}(r), \operatorname{ccont\text{-}vstar}(t)(x)) \end{array}}}
\vnbprop{ccont-vstar-vstar}{\vnbstmt{\begin{array}{@{}l@{}} \forall t.; \forall x \in C({t}). \\ \quad \operatorname{is\text{-}top\text{-}space}(t) \Rightarrow (\operatorname{ccont\text{-}vstar}(t)(\operatorname{ccont\text{-}vstar}(t)(x)) = x) \end{array}}}
\vnbprop{ccont-vstar-vmul}{\vnbstmt{\begin{array}{@{}l@{}} \forall t.; \forall x, y \in C({t}). \\ \quad \operatorname{is\text{-}top\text{-}space}(t) \Rightarrow  \\ \quad \quad \operatorname{ccont\text{-}vstar}(t)(\operatorname{ccont\text{-}vmul}(t)(x, y)) = \operatorname{ccont\text{-}vmul}(t)(\operatorname{ccont\text{-}vstar}(t)(y), \operatorname{ccont\text{-}vstar}(t)(x)) \end{array}}}
\vnbprop{ccont-vmul-comm}{\vnbstmt{\begin{array}{@{}l@{}} \forall t.; \forall x, y \in C({t}). \\ \quad \operatorname{is\text{-}top\text{-}space}(t) \Rightarrow (\operatorname{ccont\text{-}vmul}(t)(x, y) = \operatorname{ccont\text{-}vmul}(t)(y, x)) \end{array}}}
\vnbprop{ccont-vnrm-vstar}{\vnbstmt{\begin{array}{@{}l@{}} \forall t.; \forall x \in C({t}). \\ \quad \operatorname{is\text{-}compact\text{-}t}(t) \Rightarrow  \\ \quad \quad \operatorname{ccont\text{-}vnrm}(t)(\operatorname{ccont\text{-}vstar}(t)(x)) = \operatorname{ccont\text{-}vnrm}(t)(x) \end{array}}}
\vnbprop{ccont-vnrm-cstar}{\vnbstmt{\begin{array}{@{}l@{}} \forall t.; \forall x \in C({t}). \\ \quad \operatorname{is\text{-}compact\text{-}t}(t) \Rightarrow  \\ \quad \quad \operatorname{ccont\text{-}vnrm}(t)(\operatorname{ccont\text{-}vmul}(t)(\operatorname{ccont\text{-}vstar}(t)(x), x)) = (\operatorname{ccont\text{-}vnrm}(t)(x) \cdot \operatorname{ccont\text{-}vnrm}(t)(x)) \end{array}}}
\vnbprop{ccont-vnrm-vone}{\vnbstmt{\begin{array}{@{}l@{}} \forall t. \\ \quad (\operatorname{is\text{-}compact\text{-}t}(t) \wedge \exists p.\; (p \in \operatorname{pts}(t))) \Rightarrow  \\ \quad \quad \operatorname{ccont\text{-}vnrm}(t)(\operatorname{ccont\text{-}vone}(t)) = 1 \end{array}}}
\vnbprop{ccont-algebra-is-cba}{\vnbstmt{\begin{array}{@{}l@{}} \forall t. \\ \quad \operatorname{is\text{-}compact\text{-}t}(t) \Rightarrow  \\ \quad \quad \operatorname{is\text{-}complex\text{-}banach\text{-}algebra}({[}\operatorname{normed\text{-}field\text{-}as\text{-}commutative\text{-}ring}(\operatorname{cc\text{-}normed\text{-}field}), C({t}), \operatorname{ccont\text{-}vadd}(t), \operatorname{ccont\text{-}vzero}(t), \operatorname{ccont\text{-}vneg}(t), \operatorname{ccont\text{-}act}(t), \operatorname{ccont\text{-}vnrm}(t), \operatorname{ccont\text{-}vmul}(t)]) \end{array}}}
\vnbprop{ccont-algebra-is-uba}{\vnbstmt{\begin{array}{@{}l@{}} \forall t. \\ \quad (\operatorname{is\text{-}compact\text{-}t}(t) \wedge \exists p.\; (p \in \operatorname{pts}(t))) \Rightarrow  \\ \quad \quad \operatorname{is\text{-}unital\text{-}banach\text{-}algebra}({[}\operatorname{normed\text{-}field\text{-}as\text{-}commutative\text{-}ring}(\operatorname{cc\text{-}normed\text{-}field}), C({t}), \operatorname{ccont\text{-}vadd}(t), \operatorname{ccont\text{-}vzero}(t), \operatorname{ccont\text{-}vneg}(t), \operatorname{ccont\text{-}act}(t), \operatorname{ccont\text{-}vnrm}(t), \operatorname{ccont\text{-}vmul}(t), \operatorname{ccont\text{-}vone}(t)]) \end{array}}}
\vnbprop{ccont-algebra-is-uiba}{\vnbstmt{\begin{array}{@{}l@{}} \forall t. \\ \quad (\operatorname{is\text{-}compact\text{-}t}(t) \wedge \exists p.\; (p \in \operatorname{pts}(t))) \Rightarrow  \\ \quad \quad \operatorname{is\text{-}unital\text{-}involutive\text{-}banach\text{-}algebra}({[}\operatorname{normed\text{-}field\text{-}as\text{-}commutative\text{-}ring}(\operatorname{cc\text{-}normed\text{-}field}), C({t}), \operatorname{ccont\text{-}vadd}(t), \operatorname{ccont\text{-}vzero}(t), \operatorname{ccont\text{-}vneg}(t), \operatorname{ccont\text{-}act}(t), \operatorname{ccont\text{-}vnrm}(t), \operatorname{ccont\text{-}vmul}(t), \operatorname{ccont\text{-}vone}(t), \operatorname{ccont\text{-}vstar}(t)]) \end{array}}}
\vnbprop{ccont-algebra-is-ucsa-tuple}{\vnbstmt{\begin{array}{@{}l@{}} \forall t. \\ \quad (\operatorname{is\text{-}compact\text{-}t}(t) \wedge \exists p.\; (p \in \operatorname{pts}(t))) \Rightarrow  \\ \quad \quad \operatorname{is\text{-}unital\text{-}c\text{-}star\text{-}algebra}({[}\operatorname{normed\text{-}field\text{-}as\text{-}commutative\text{-}ring}(\operatorname{cc\text{-}normed\text{-}field}), C({t}), \operatorname{ccont\text{-}vadd}(t), \operatorname{ccont\text{-}vzero}(t), \operatorname{ccont\text{-}vneg}(t), \operatorname{ccont\text{-}act}(t), \operatorname{ccont\text{-}vnrm}(t), \operatorname{ccont\text{-}vmul}(t), \operatorname{ccont\text{-}vone}(t), \operatorname{ccont\text{-}vstar}(t)]) \end{array}}}
\vnbprop{ccont-algebra-is-ccsa-tuple}{\vnbstmt{\begin{array}{@{}l@{}} \forall t. \\ \quad (\operatorname{is\text{-}compact\text{-}t}(t) \wedge \exists p.\; (p \in \operatorname{pts}(t))) \Rightarrow  \\ \quad \quad \operatorname{is\text{-}commutative\text{-}c\text{-}star\text{-}algebra}({[}\operatorname{normed\text{-}field\text{-}as\text{-}commutative\text{-}ring}(\operatorname{cc\text{-}normed\text{-}field}), C({t}), \operatorname{ccont\text{-}vadd}(t), \operatorname{ccont\text{-}vzero}(t), \operatorname{ccont\text{-}vneg}(t), \operatorname{ccont\text{-}act}(t), \operatorname{ccont\text{-}vnrm}(t), \operatorname{ccont\text{-}vmul}(t), \operatorname{ccont\text{-}vone}(t), \operatorname{ccont\text{-}vstar}(t)]) \end{array}}}
\vnbprop{ccont-algebra-unital-c-star}{\vnbstmt{\begin{array}{@{}l@{}} \forall t. \\ \quad (\operatorname{is\text{-}compact\text{-}t}(t) \wedge \exists p.\; (p \in \operatorname{pts}(t))) \Rightarrow  \\ \quad \quad \operatorname{is\text{-}unital\text{-}c\text{-}star\text{-}algebra}(\operatorname{ccont\text{-}algebra}(t)) \end{array}}}
\vnbprop{ccont-algebra-commutative-c-star}{\vnbstmt{\begin{array}{@{}l@{}} \forall t. \\ \quad (\operatorname{is\text{-}compact\text{-}t}(t) \wedge \exists p.\; (p \in \operatorname{pts}(t))) \Rightarrow  \\ \quad \quad \operatorname{is\text{-}commutative\text{-}c\text{-}star\text{-}algebra}(\operatorname{ccont\text{-}algebra}(t)) \end{array}}}
\vnbprop{top-open-union2}{\vnbstmt{\begin{array}{@{}l@{}} \forall t.; \forall u, v \in \operatorname{opens}(t). \\ \quad \operatorname{is\text{-}top\text{-}space}(t) \Rightarrow ((u \cup v) \in \operatorname{opens}(t)) \end{array}}}
\vnbprop{top-hausdorff-point-closed-sep}{\vnbstmt{\begin{array}{@{}l@{}} \forall t, k.; \forall p \in \operatorname{pts}(t). \\ \quad (\operatorname{is\text{-}compact\text{-}t}(t) \wedge (\operatorname{is\text{-}hausdorff}(t) \wedge (\operatorname{is\text{-}top\text{-}closed}(t, k) \wedge \neg (p \in k)))) \Rightarrow  \\ \quad \quad \exists u, v \in \operatorname{opens}(t). \\ \quad \quad \quad (p \in u) \wedge ((k \subseteq v) \wedge \forall x \in u.\; \neg (x \in v)) \end{array}}}
\vnbprop{top-compact-hausdorff-normal}{\vnbstmt{\begin{array}{@{}l@{}} \forall t, a, b. \\ \quad \operatorname{is\text{-}compact\text{-}t}(t) \wedge \\ \quad \operatorname{is\text{-}hausdorff}(t) \wedge \\ \quad \operatorname{is\text{-}top\text{-}closed}(t, a) \wedge \\ \quad \operatorname{is\text{-}top\text{-}closed}(t, b) \wedge \\ \quad \forall x \in a.\; \neg (x \in b) \Rightarrow \\ \quad \quad \exists u, v \in \operatorname{opens}(t). \\ \quad \quad \quad (a \subseteq u) \wedge ((b \subseteq v) \wedge \forall x \in u.\; \neg (x \in v)) \end{array}}}
\vnbprop{top-compact-hausdorff-shrink}{\vnbstmt{\begin{array}{@{}l@{}} \forall t, a.; \forall u \in \operatorname{opens}(t). \\ \quad (\operatorname{is\text{-}compact\text{-}t}(t) \wedge (\operatorname{is\text{-}hausdorff}(t) \wedge (\operatorname{is\text{-}top\text{-}closed}(t, a) \wedge (a \subseteq u)))) \Rightarrow  \\ \quad \quad \exists v \in \operatorname{opens}(t), f. \\ \quad \quad \quad \operatorname{is\text{-}top\text{-}closed}(t, f) \wedge ((a \subseteq v) \wedge ((v \subseteq f) \wedge (f \subseteq u))) \end{array}}}
\vnbprop{ccont-algebra-scal}{\vnbstmt{\begin{array}{@{}l@{}} \forall t. \\ \quad \operatorname{is\text{-}compact\text{-}t}(t) \Rightarrow  \\ \quad \quad (\operatorname{scal}(\operatorname{ccont\text{-}algebra}(t)) \simeq \operatorname{normed\text{-}field\text{-}as\text{-}commutative\text{-}ring}(\operatorname{cc\text{-}normed\text{-}field})) \end{array}}}
\vnbprop{ccont-algebra-vec}{\vnbstmt{\forall t.\; (\operatorname{is\text{-}compact\text{-}t}(t) \Rightarrow (\operatorname{vec}(\operatorname{ccont\text{-}algebra}(t)) \simeq C({t})))}}
\vnbprop{ccont-algebra-vadd}{\vnbstmt{\forall t.\; (\operatorname{is\text{-}compact\text{-}t}(t) \Rightarrow (\operatorname{vadd}(\operatorname{ccont\text{-}algebra}(t)) \simeq \operatorname{ccont\text{-}vadd}(t)))}}
\vnbprop{ccont-algebra-vzero}{\vnbstmt{\forall t.\; (\operatorname{is\text{-}compact\text{-}t}(t) \Rightarrow (\operatorname{vzero}(\operatorname{ccont\text{-}algebra}(t)) \simeq \operatorname{ccont\text{-}vzero}(t)))}}
\vnbprop{ccont-algebra-vneg}{\vnbstmt{\forall t.\; (\operatorname{is\text{-}compact\text{-}t}(t) \Rightarrow (\operatorname{vneg}(\operatorname{ccont\text{-}algebra}(t)) \simeq \operatorname{ccont\text{-}vneg}(t)))}}
\vnbprop{ccont-algebra-act}{\vnbstmt{\forall t.\; (\operatorname{is\text{-}compact\text{-}t}(t) \Rightarrow (\operatorname{act}(\operatorname{ccont\text{-}algebra}(t)) \simeq \operatorname{ccont\text{-}act}(t)))}}
\vnbprop{ccont-algebra-vnrm}{\vnbstmt{\forall t.\; (\operatorname{is\text{-}compact\text{-}t}(t) \Rightarrow (\operatorname{vnrm}(\operatorname{ccont\text{-}algebra}(t)) \simeq \operatorname{ccont\text{-}vnrm}(t)))}}
\vnbprop{ccont-algebra-vmul}{\vnbstmt{\forall t.\; (\operatorname{is\text{-}compact\text{-}t}(t) \Rightarrow (\cdot_{\operatorname{ccont\text{-}algebra}(t)} \simeq \operatorname{ccont\text{-}vmul}(t)))}}
\vnbprop{ccont-algebra-vone}{\vnbstmt{\forall t.\; (\operatorname{is\text{-}compact\text{-}t}(t) \Rightarrow (1_{\operatorname{ccont\text{-}algebra}(t)} \simeq \operatorname{ccont\text{-}vone}(t)))}}
\vnbprop{ccont-algebra-vstar}{\vnbstmt{\forall t.\; (\operatorname{is\text{-}compact\text{-}t}(t) \Rightarrow ({}^{*_{\operatorname{ccont\text{-}algebra}(t)}} \simeq \operatorname{ccont\text{-}vstar}(t)))}}
\vnbprop{pos-cone-unfold}{\vnbstmt{\begin{array}{@{}l@{}} \forall a. \\ \quad ({a}_{+} \simeq \{x \in \operatorname{vec}(a) : \forall c \in \operatorname{power}(\operatorname{vec}(a)).\; ((\forall y \in \operatorname{vec}(a).\; (\cdot_{a}({}^{*_{a}}(y), y) \in c) \wedge \forall u, v \in c.\; (\operatorname{vadd}(a)(u, v) \in c)) \Rightarrow (x \in c))\}) \end{array}}}
\vnbprop{ccont-pos-cone-nonneg}{\vnbstmt{\begin{array}{@{}l@{}} \forall t.; \forall f \in {\operatorname{ccont\text{-}algebra}(t)}_{+}. \\ \quad \operatorname{is\text{-}compact\text{-}t}(t) \Rightarrow  \\ \quad \quad (f \in C({t})) \wedge \forall x \in \operatorname{pts}(t).\; ((f(x) \in \mathbb{R}) \wedge (0 \leq f(x))) \end{array}}}
\vnbprop{rr-sq-diff-lt}{\vnbstmt{\begin{array}{@{}l@{}} \forall x, y, e \in \mathbb{R}. \\ \quad ((0 \leq y) \wedge ((y \leq x) \wedge ((0 < e) \wedge (((x \cdot x) - (y \cdot y)) < (e \cdot e))))) \Rightarrow ((x - y) < e) \end{array}}}
\vnbprop{rr-sqrt-near}{\vnbstmt{\begin{array}{@{}l@{}} \forall p, q, e \in \mathbb{R}. \\ \quad ((0 \leq p) \wedge ((0 \leq q) \wedge ((0 < e) \wedge (\lvert (p - q) \rvert < (e \cdot e))))) \Rightarrow  \\ \quad \quad \lvert (\sqrt{p} - \sqrt{q}) \rvert < e \end{array}}}
\vnbprop{cc-sqrt-re-near}{\vnbstmt{\begin{array}{@{}l@{}} \forall a \in \mathbb{C}, s. \\ \quad \operatorname{pos\text{-}rr}(s) \Rightarrow  \\ \quad \quad \exists d. \\ \quad \quad \quad \operatorname{pos\text{-}rr}(d) \wedge \\ \quad \quad \quad \forall x \in \mathbb{C}. \\ \quad \quad \quad \quad (\operatorname{magnitude}((a - x)) < d) \Rightarrow  \\ \quad \quad \quad \quad \quad \operatorname{magnitude}((\sqrt{\lvert \operatorname{real\text{-}part}(a) \rvert} - \sqrt{\lvert \operatorname{real\text{-}part}(x) \rvert})) < s \end{array}}}
\vnbprop{ccont-sqrt-re-ptw}{\vnbstmt{\begin{array}{@{}l@{}} \forall t, f, h. \\ \quad ((f \in C({t})) \wedge ((h \in (\operatorname{pts}(t) \to \mathbb{C})) \wedge \forall z \in \operatorname{pts}(t).\; (h(z) = \sqrt{\lvert \operatorname{real\text{-}part}(f(z)) \rvert}))) \Rightarrow  \\ \quad \quad h \in C({t}) \end{array}}}
\vnbprop{ccont-nonneg-pos-cone}{\vnbstmt{\begin{array}{@{}l@{}} \forall t.; \forall f \in C({t}). \\ \quad (\operatorname{is\text{-}compact\text{-}t}(t) \wedge \forall x \in \operatorname{pts}(t).\; ((f(x) \in \mathbb{R}) \wedge (0 \leq f(x)))) \Rightarrow  \\ \quad \quad f \in {\operatorname{ccont\text{-}algebra}(t)}_{+} \end{array}}}
\vnbprop{ccont-positive-iff-nonneg}{\vnbstmt{\begin{array}{@{}l@{}} \forall t.; \forall f \in C({t}). \\ \quad \operatorname{is\text{-}compact\text{-}t}(t) \Rightarrow  \\ \quad \quad (f \in {\operatorname{ccont\text{-}algebra}(t)}_{+}) \Leftrightarrow \forall x \in \operatorname{pts}(t).\; ((f(x) \in \mathbb{R}) \wedge (0 \leq f(x))) \end{array}}}

\vnbbookfile{theorem-library/kal-cstar-122}{theorem-library/kal-cstar-122}
\noindent{\small kal-cstar-122.scm -- wave of 2026-11-09/10 (agent KIT-ALG), wired 2026-11-10 from scratchpad/kitalg-1109/kal-cstar-122.scm. Library: operator-algebras.}\par
\vnbprop{kcs-character-sa-real}{\vnbstmt{\begin{array}{@{}l@{}} \forall a.; \forall c \in \operatorname{characters}(a), h. \\ \quad (\operatorname{is\text{-}unital\text{-}c\text{-}star\text{-}algebra}(a) \wedge \operatorname{is\text{-}self\text{-}adjoint\text{-}el}(a, h)) \Rightarrow (c(h) \in \mathbb{R}) \end{array}}}
\vnbprop{character-preserves-involution}{\vnbstmt{\begin{array}{@{}l@{}} \forall a.; \forall c \in \operatorname{characters}(a), x \in \operatorname{vec}(a). \\ \quad \operatorname{is\text{-}unital\text{-}c\text{-}star\text{-}algebra}(a) \Rightarrow (c({}^{*_{a}}(x)) = \operatorname{conjugate}(c(x))) \end{array}}}
\vnbprop{ccsa-character-preserves-involution}{\vnbstmt{\begin{array}{@{}l@{}} \forall a.; \forall c \in \operatorname{characters}(a), x \in \operatorname{vec}(a). \\ \quad \operatorname{is\text{-}commutative\text{-}c\text{-}star\text{-}algebra}(a) \Rightarrow (c({}^{*_{a}}(x)) = \operatorname{conjugate}(c(x))) \end{array}}}

\vnbbookfile{theorem-library/basw3-urysohn}{theorem-library/basw3-urysohn}
\noindent{\small basw3-urysohn.scm -- wave of 2026-11-09/10 (agent BA-SW3), wired 2026-11-10 from scratchpad/basw3-1109/basw3-urysohn.scm. Library: operator-algebras.}\par
\vnbprop{ury-cu-unfold}{\vnbstmt{\begin{array}{@{}l@{}} \forall t. \\ \quad (\operatorname{cu\text{-}pairs}(t) \simeq \{w \in (\operatorname{power}(\operatorname{pts}(t)) \times \operatorname{opens}(t)) : (\operatorname{is\text{-}top\text{-}closed}(t, \operatorname{nth}(1, w)) \wedge (\operatorname{nth}(1, w) \subseteq \operatorname{nth}(2, w)))\}) \end{array}}}
\vnbprop{ury-cu-set}{\vnbstmt{\forall t.\; (\operatorname{is\text{-}top\text{-}space}(t) \Rightarrow (\operatorname{cu\text{-}pairs}(t) \in \mathbf{Set}))}}
\vnbprop{ury-cu-parts}{\vnbstmt{\begin{array}{@{}l@{}} \forall t.; \forall c \in \operatorname{cu\text{-}pairs}(t). \\ \quad \operatorname{is\text{-}top\text{-}closed}(t, \operatorname{nth}(1, c)) \wedge \\ \quad \operatorname{nth}(2, c) \in \operatorname{opens}(t) \wedge \\ \quad \operatorname{nth}(1, c) \subseteq \operatorname{nth}(2, c) \wedge \\ \quad c = {[}\operatorname{nth}(1, c), \operatorname{nth}(2, c)] \end{array}}}
\vnbprop{ury-cu-intro}{\vnbstmt{\begin{array}{@{}l@{}} \forall t, a, u. \\ \quad (\operatorname{is\text{-}top\text{-}space}(t) \wedge (\operatorname{is\text{-}top\text{-}closed}(t, a) \wedge ((u \in \operatorname{opens}(t)) \wedge (a \subseteq u)))) \Rightarrow  \\ \quad \quad {[}a, u] \in \operatorname{cu\text{-}pairs}(t) \end{array}}}
\vnbprop{ury-v-spec}{\vnbstmt{\begin{array}{@{}l@{}} \forall t, c. \\ \quad (\operatorname{is\text{-}compact\text{-}t}(t) \wedge (\operatorname{is\text{-}hausdorff}(t) \wedge (c \in \operatorname{cu\text{-}pairs}(t)))) \Rightarrow  \\ \quad \quad \operatorname{ury\text{-}v}(t, c) \in \operatorname{opens}(t) \wedge \\ \quad \quad \exists k \in \operatorname{power}(\operatorname{pts}(t)). \\ \quad \quad \quad (\operatorname{is\text{-}top\text{-}closed}(t, k) \wedge ((\operatorname{nth}(1, c) \subseteq \operatorname{ury\text{-}v}(t, c)) \wedge ((\operatorname{ury\text{-}v}(t, c) \subseteq k) \wedge (k \subseteq \operatorname{nth}(2, c))))) \end{array}}}
\vnbprop{ury-k-spec}{\vnbstmt{\begin{array}{@{}l@{}} \forall t, c. \\ \quad (\operatorname{is\text{-}compact\text{-}t}(t) \wedge (\operatorname{is\text{-}hausdorff}(t) \wedge (c \in \operatorname{cu\text{-}pairs}(t)))) \Rightarrow  \\ \quad \quad \operatorname{ury\text{-}k}(t, c) \in \operatorname{power}(\operatorname{pts}(t)) \wedge \\ \quad \quad \operatorname{is\text{-}top\text{-}closed}(t, \operatorname{ury\text{-}k}(t, c)) \wedge \\ \quad \quad \operatorname{ury\text{-}v}(t, c) \subseteq \operatorname{ury\text{-}k}(t, c) \wedge \\ \quad \quad \operatorname{ury\text{-}k}(t, c) \subseteq \operatorname{nth}(2, c) \end{array}}}
\vnbprop{ury-l-in}{\vnbstmt{\begin{array}{@{}l@{}} \forall t, c. \\ \quad (\operatorname{is\text{-}compact\text{-}t}(t) \wedge (\operatorname{is\text{-}hausdorff}(t) \wedge (c \in \operatorname{cu\text{-}pairs}(t)))) \Rightarrow  \\ \quad \quad \operatorname{ury\text{-}l}(t, c) \in \operatorname{cu\text{-}pairs}(t) \end{array}}}
\vnbprop{ury-r-in}{\vnbstmt{\begin{array}{@{}l@{}} \forall t, c. \\ \quad (\operatorname{is\text{-}compact\text{-}t}(t) \wedge (\operatorname{is\text{-}hausdorff}(t) \wedge (c \in \operatorname{cu\text{-}pairs}(t)))) \Rightarrow  \\ \quad \quad \operatorname{ury\text{-}r}(t, c) \in \operatorname{cu\text{-}pairs}(t) \end{array}}}
\vnbprop{ury-lr-nth}{\vnbstmt{\begin{array}{@{}l@{}} \forall t, c. \\ \quad (\operatorname{is\text{-}compact\text{-}t}(t) \wedge (\operatorname{is\text{-}hausdorff}(t) \wedge (c \in \operatorname{cu\text{-}pairs}(t)))) \Rightarrow  \\ \quad \quad \operatorname{nth}(1, \operatorname{ury\text{-}l}(t, c)) = \operatorname{nth}(1, c) \wedge \\ \quad \quad \operatorname{nth}(2, \operatorname{ury\text{-}l}(t, c)) = \operatorname{ury\text{-}v}(t, c) \wedge \\ \quad \quad \operatorname{nth}(1, \operatorname{ury\text{-}r}(t, c)) = \operatorname{ury\text{-}k}(t, c) \wedge \\ \quad \quad \operatorname{nth}(2, \operatorname{ury\text{-}r}(t, c)) = \operatorname{nth}(2, c) \end{array}}}
\vnbprop{ury-apx-type}{\vnbstmt{\begin{array}{@{}l@{}} \forall n \in \mathbb{N}, t, x. \\ \quad (\operatorname{is\text{-}compact\text{-}t}(t) \wedge \operatorname{is\text{-}hausdorff}(t)) \Rightarrow  \\ \quad \quad \operatorname{ury\text{-}apx}(t, x, n) \in (\operatorname{cu\text{-}pairs}(t) \to \mathbb{R}) \end{array}}}
\vnbprop{ury-apx-zero-in}{\vnbstmt{\begin{array}{@{}l@{}} \forall t, x, c. \\ \quad (\operatorname{is\text{-}compact\text{-}t}(t) \wedge (\operatorname{is\text{-}hausdorff}(t) \wedge ((c \in \operatorname{cu\text{-}pairs}(t)) \wedge (x \in \operatorname{nth}(2, c))))) \Rightarrow  \\ \quad \quad \operatorname{ury\text{-}apx}(t, x, 0)(c) = 0 \end{array}}}
\vnbprop{ury-apx-zero-out}{\vnbstmt{\begin{array}{@{}l@{}} \forall t, x, c. \\ \quad (\operatorname{is\text{-}compact\text{-}t}(t) \wedge (\operatorname{is\text{-}hausdorff}(t) \wedge ((c \in \operatorname{cu\text{-}pairs}(t)) \wedge \neg (x \in \operatorname{nth}(2, c))))) \Rightarrow  \\ \quad \quad \operatorname{ury\text{-}apx}(t, x, 0)(c) = 1 \end{array}}}
\vnbprop{ury-apx-succ-val}{\vnbstmt{\begin{array}{@{}l@{}} \forall n \in \mathbb{N}, t, x, c. \\ \quad (\operatorname{is\text{-}compact\text{-}t}(t) \wedge (\operatorname{is\text{-}hausdorff}(t) \wedge (c \in \operatorname{cu\text{-}pairs}(t)))) \Rightarrow  \\ \quad \quad \operatorname{ury\text{-}apx}(t, x, (n + 1))(c) \\ \quad \quad =\; 1/2 \cdot (\operatorname{ury\text{-}apx}(t, x, n)(\operatorname{ury\text{-}l}(t, c)) + \operatorname{ury\text{-}apx}(t, x, n)(\operatorname{ury\text{-}r}(t, c))) \end{array}}}
\vnbprop{ury-apx-bounds}{\vnbstmt{\begin{array}{@{}l@{}} \forall n \in \mathbb{N}, t, x.; \forall c \in \operatorname{cu\text{-}pairs}(t). \\ \quad (\operatorname{is\text{-}compact\text{-}t}(t) \wedge \operatorname{is\text{-}hausdorff}(t)) \Rightarrow  \\ \quad \quad (0 \leq \operatorname{ury\text{-}apx}(t, x, n)(c)) \wedge (\operatorname{ury\text{-}apx}(t, x, n)(c) \leq 1) \end{array}}}
\vnbprop{ury-apx-on-c}{\vnbstmt{\begin{array}{@{}l@{}} \forall n \in \mathbb{N}, t, x.; \forall c \in \operatorname{cu\text{-}pairs}(t). \\ \quad (\operatorname{is\text{-}compact\text{-}t}(t) \wedge (\operatorname{is\text{-}hausdorff}(t) \wedge (x \in \operatorname{nth}(1, c)))) \Rightarrow  \\ \quad \quad \operatorname{ury\text{-}apx}(t, x, n)(c) = 0 \end{array}}}
\vnbprop{ury-apx-off-u}{\vnbstmt{\begin{array}{@{}l@{}} \forall n \in \mathbb{N}, t, x.; \forall c \in \operatorname{cu\text{-}pairs}(t). \\ \quad (\operatorname{is\text{-}compact\text{-}t}(t) \wedge (\operatorname{is\text{-}hausdorff}(t) \wedge \neg (x \in \operatorname{nth}(2, c)))) \Rightarrow  \\ \quad \quad \operatorname{ury\text{-}apx}(t, x, n)(c) = 1 \end{array}}}
\vnbprop{ury-apx-le-of-sub}{\vnbstmt{\begin{array}{@{}l@{}} \forall t, x.; \forall a, b \in \operatorname{cu\text{-}pairs}(t), m, k \in \mathbb{N}. \\ \quad (\operatorname{is\text{-}compact\text{-}t}(t) \wedge (\operatorname{is\text{-}hausdorff}(t) \wedge (\operatorname{nth}(2, a) \subseteq \operatorname{nth}(1, b)))) \Rightarrow  \\ \quad \quad \operatorname{ury\text{-}apx}(t, x, k)(b) \leq \operatorname{ury\text{-}apx}(t, x, m)(a) \end{array}}}
\vnbprop{ury-apx-mid}{\vnbstmt{\begin{array}{@{}l@{}} \forall n \in \mathbb{N}, t, x.; \forall c \in \operatorname{cu\text{-}pairs}(t). \\ \quad (\operatorname{is\text{-}compact\text{-}t}(t) \wedge \operatorname{is\text{-}hausdorff}(t)) \Rightarrow  \\ \quad \quad \operatorname{ury\text{-}apx}(t, x, n)(\operatorname{ury\text{-}r}(t, c)) \leq \operatorname{ury\text{-}apx}(t, x, n)(c) \wedge \\ \quad \quad \operatorname{ury\text{-}apx}(t, x, n)(c) \leq \operatorname{ury\text{-}apx}(t, x, n)(\operatorname{ury\text{-}l}(t, c)) \end{array}}}
\vnbprop{ury-apx-mono}{\vnbstmt{\begin{array}{@{}l@{}} \forall n \in \mathbb{N}, t, x.; \forall c \in \operatorname{cu\text{-}pairs}(t). \\ \quad (\operatorname{is\text{-}compact\text{-}t}(t) \wedge \operatorname{is\text{-}hausdorff}(t)) \Rightarrow  \\ \quad \quad \operatorname{ury\text{-}apx}(t, x, n)(c) \leq \operatorname{ury\text{-}apx}(t, x, (n + 1))(c) \end{array}}}
\vnbprop{ury-apx-mono-add}{\vnbstmt{\begin{array}{@{}l@{}} \forall d \in \mathbb{N}, t, x.; \forall c \in \operatorname{cu\text{-}pairs}(t), k \in \mathbb{N}. \\ \quad (\operatorname{is\text{-}compact\text{-}t}(t) \wedge \operatorname{is\text{-}hausdorff}(t)) \Rightarrow  \\ \quad \quad \operatorname{ury\text{-}apx}(t, x, k)(c) \leq \operatorname{ury\text{-}apx}(t, x, (k + d))(c) \end{array}}}
\vnbprop{ury-lim-unfold}{\vnbstmt{\begin{array}{@{}l@{}} \forall t, c, x. \\ \quad (\operatorname{ury\text{-}lim}(t, c, x) \simeq \operatorname{esup}(\{r \in \mathbb{R} : \exists n \in \mathbb{N}.\; (r = \operatorname{ury\text{-}apx}(t, x, n)(c))\})) \end{array}}}
\vnbprop{ury-lim-sub}{\vnbstmt{\begin{array}{@{}l@{}} \forall t, x.; \forall c \in \operatorname{cu\text{-}pairs}(t). \\ \quad (\operatorname{is\text{-}compact\text{-}t}(t) \wedge \operatorname{is\text{-}hausdorff}(t)) \Rightarrow  \\ \quad \quad \{r \in \mathbb{R} : \exists n \in \mathbb{N}.\; (r = \operatorname{ury\text{-}apx}(t, x, n)(c))\} \subseteq \operatorname{rr\text{-}pos\text{-}star} \end{array}}}
\vnbprop{ury-lim-upper}{\vnbstmt{\begin{array}{@{}l@{}} \forall t, x.; \forall c \in \operatorname{cu\text{-}pairs}(t), n \in \mathbb{N}. \\ \quad (\operatorname{is\text{-}compact\text{-}t}(t) \wedge \operatorname{is\text{-}hausdorff}(t)) \Rightarrow  \\ \quad \quad \operatorname{ury\text{-}apx}(t, x, n)(c) \leq \operatorname{ury\text{-}lim}(t, c, x) \end{array}}}
\vnbprop{ury-lim-least}{\vnbstmt{\begin{array}{@{}l@{}} \forall t, x.; \forall c \in \operatorname{cu\text{-}pairs}(t), b \in \mathbb{R}. \\ \quad (\operatorname{is\text{-}compact\text{-}t}(t) \wedge (\operatorname{is\text{-}hausdorff}(t) \wedge \forall q \in \mathbb{N}.\; (\operatorname{ury\text{-}apx}(t, x, q)(c) \leq b))) \Rightarrow  \\ \quad \quad \operatorname{ury\text{-}lim}(t, c, x) \leq b \end{array}}}
\vnbprop{ury-lim-real}{\vnbstmt{\begin{array}{@{}l@{}} \forall t, x.; \forall c \in \operatorname{cu\text{-}pairs}(t). \\ \quad (\operatorname{is\text{-}compact\text{-}t}(t) \wedge \operatorname{is\text{-}hausdorff}(t)) \Rightarrow  \\ \quad \quad (\operatorname{ury\text{-}lim}(t, c, x) \in \mathbb{R}) \wedge ((0 \leq \operatorname{ury\text{-}lim}(t, c, x)) \wedge (\operatorname{ury\text{-}lim}(t, c, x) \leq 1)) \end{array}}}
\vnbprop{ury-lim-on-c}{\vnbstmt{\begin{array}{@{}l@{}} \forall t, x.; \forall c \in \operatorname{cu\text{-}pairs}(t). \\ \quad (\operatorname{is\text{-}compact\text{-}t}(t) \wedge (\operatorname{is\text{-}hausdorff}(t) \wedge (x \in \operatorname{nth}(1, c)))) \Rightarrow  \\ \quad \quad \operatorname{ury\text{-}lim}(t, c, x) = 0 \end{array}}}
\vnbprop{ury-lim-off-u}{\vnbstmt{\begin{array}{@{}l@{}} \forall t, x.; \forall c \in \operatorname{cu\text{-}pairs}(t). \\ \quad (\operatorname{is\text{-}compact\text{-}t}(t) \wedge (\operatorname{is\text{-}hausdorff}(t) \wedge \neg (x \in \operatorname{nth}(2, c)))) \Rightarrow  \\ \quad \quad \operatorname{ury\text{-}lim}(t, c, x) = 1 \end{array}}}
\vnbprop{ury-lim-mid}{\vnbstmt{\begin{array}{@{}l@{}} \forall t, x.; \forall c \in \operatorname{cu\text{-}pairs}(t). \\ \quad (\operatorname{is\text{-}compact\text{-}t}(t) \wedge \operatorname{is\text{-}hausdorff}(t)) \Rightarrow  \\ \quad \quad \operatorname{ury\text{-}lim}(t, c, x) = (1/2 \cdot (\operatorname{ury\text{-}lim}(t, \operatorname{ury\text{-}l}(t, c), x) + \operatorname{ury\text{-}lim}(t, \operatorname{ury\text{-}r}(t, c), x))) \end{array}}}
\vnbprop{top-open-inter2}{\vnbstmt{\begin{array}{@{}l@{}} \forall t.; \forall u, v \in \operatorname{opens}(t). \\ \quad \operatorname{is\text{-}top\text{-}space}(t) \Rightarrow ((u \cap v) \in \operatorname{opens}(t)) \end{array}}}
\vnbprop{ury-lim-osc}{\vnbstmt{\begin{array}{@{}l@{}} \forall n \in \mathbb{N}, t, x.; \forall c \in \operatorname{cu\text{-}pairs}(t). \\ \quad (\operatorname{is\text{-}compact\text{-}t}(t) \wedge (\operatorname{is\text{-}hausdorff}(t) \wedge (x \in \operatorname{pts}(t)))) \Rightarrow  \\ \quad \quad \exists w \in \operatorname{opens}(t). \\ \quad \quad \quad ((x \in w) \wedge \forall y \in w.\; (\lvert (\operatorname{ury\text{-}lim}(t, c, y) - \operatorname{ury\text{-}lim}(t, c, x)) \rvert \leq \operatorname{power}(3/4, n))) \end{array}}}
\vnbprop{urysohn-compact-hausdorff}{\vnbstmt{\begin{array}{@{}l@{}} \forall t, a, b. \\ \quad \operatorname{is\text{-}compact\text{-}t}(t) \wedge \\ \quad \operatorname{is\text{-}hausdorff}(t) \wedge \\ \quad \operatorname{is\text{-}top\text{-}closed}(t, a) \wedge \\ \quad \operatorname{is\text{-}top\text{-}closed}(t, b) \wedge \\ \quad \forall x \in a.\; \neg (x \in b) \Rightarrow \\ \quad \quad \exists f \in C_{\mathbb{R}}({t}). \\ \quad \quad \quad \forall z \in \operatorname{pts}(t).\; ((0 \leq f(z)) \wedge (f(z) \leq 1)) \wedge \\ \quad \quad \quad \forall z \in a.\; (f(z) = 0) \wedge \\ \quad \quad \quad \forall z \in b.\; (f(z) = 1) \end{array}}}
\vnbprop{ccx-characters-unfold}{\vnbstmt{\begin{array}{@{}l@{}} \forall t. \\ \quad (\operatorname{characters}(\operatorname{ccont\text{-}algebra}(t)) \simeq \{c \in (\operatorname{vec}(\operatorname{ccont\text{-}algebra}(t)) \to \mathbb{C}) : (\operatorname{is\text{-}cc\text{-}linear\text{-}on}(\operatorname{ccont\text{-}algebra}(t), c) \wedge (\forall x, y \in \operatorname{vec}(\operatorname{ccont\text{-}algebra}(t)).\; (c(\cdot_{\operatorname{ccont\text{-}algebra}(t)}(x, y)) = (c(x) \cdot c(y))) \wedge (c(1_{\operatorname{ccont\text{-}algebra}(t)}) = 1)))\}) \end{array}}}
\vnbprop{ccont-character-laws}{\vnbstmt{\begin{array}{@{}l@{}} \forall t.; \forall h \in \operatorname{characters}(\operatorname{ccont\text{-}algebra}(t)). \\ \quad \operatorname{is\text{-}compact\text{-}t}(t) \Rightarrow  \\ \quad \quad h \in (C({t}) \to \mathbb{C}) \wedge \\ \quad \quad \forall x, y \in C({t}).\; (h(\operatorname{ccont\text{-}vadd}(t)(x, y)) = (h(x) + h(y))) \wedge \\ \quad \quad \forall z \in \mathbb{C}, x \in C({t}).\; (h(\operatorname{ccont\text{-}act}(t)(z, x)) = (z \cdot h(x))) \wedge \\ \quad \quad \forall x, y \in C({t}).\; (h(\operatorname{ccont\text{-}vmul}(t)(x, y)) = (h(x) \cdot h(y))) \wedge \\ \quad \quad h(\operatorname{ccont\text{-}vone}(t)) = 1 \end{array}}}
\vnbprop{point-eval-val}{\vnbstmt{\forall t.; \forall x \in \operatorname{pts}(t), f \in C({t}).\; (\operatorname{point\text{-}eval}(t, x)(f) = f(x))}}
\vnbprop{point-eval-character}{\vnbstmt{\begin{array}{@{}l@{}} \forall t.; \forall x \in \operatorname{pts}(t). \\ \quad \operatorname{is\text{-}compact\text{-}t}(t) \Rightarrow (\operatorname{point\text{-}eval}(t, x) \in \operatorname{characters}(\operatorname{ccont\text{-}algebra}(t))) \end{array}}}
\vnbprop{ccont-character-eval-of-kernel}{\vnbstmt{\begin{array}{@{}l@{}} \forall t.; \forall h \in \operatorname{characters}(\operatorname{ccont\text{-}algebra}(t)), x \in \operatorname{pts}(t), f \in C({t}). \\ \quad (\operatorname{is\text{-}compact\text{-}t}(t) \wedge \forall g \in C({t}).\; ((h(g) = 0) \Rightarrow (g(x) = 0))) \Rightarrow  \\ \quad \quad h(f) = f(x) \end{array}}}
\vnbprop{rr-recip-lipschitz-ge-one}{\vnbstmt{\begin{array}{@{}l@{}} \forall a, b \in \mathbb{R}. \\ \quad ((1 \leq a) \wedge (1 \leq b)) \Rightarrow (\lvert (\operatorname{recip}(a) - \operatorname{recip}(b)) \rvert \leq \lvert (a - b) \rvert) \end{array}}}
\vnbprop{ccont-character-bump}{\vnbstmt{\begin{array}{@{}l@{}} \forall t.; \forall h \in \operatorname{characters}(\operatorname{ccont\text{-}algebra}(t)), g \in C({t}), p \in \operatorname{pts}(t). \\ \quad (\operatorname{is\text{-}compact\text{-}t}(t) \wedge ((h(g) = 0) \wedge \neg (g(p) = 0))) \Rightarrow  \\ \quad \quad \exists u \in \operatorname{opens}(t). \\ \quad \quad \quad p \in u \wedge \\ \quad \quad \quad u \\ \quad \quad \quad \in\; \{w \in \operatorname{power}(\operatorname{pts}(t)) : \exists \mathit{uhh} \in C({t}).\; ((h(\mathit{uhh}) = 0) \wedge (\forall z \in \operatorname{pts}(t).\; ((\operatorname{uhh}(z) \in \mathbb{R}) \wedge (0 \leq \operatorname{uhh}(z))) \wedge \forall z \in w.\; (1 \leq \operatorname{uhh}(z))))\} \end{array}}}
\vnbprop{ccont-character-no-unit}{\vnbstmt{\begin{array}{@{}l@{}} \forall t.; \forall \mathit{uch} \in \operatorname{characters}(\operatorname{ccont\text{-}algebra}(t)), h \in C({t}). \\ \quad (\operatorname{is\text{-}compact\text{-}t}(t) \wedge \forall z \in \operatorname{pts}(t).\; ((h(z) \in \mathbb{R}) \wedge (1 \leq h(z)))) \Rightarrow  \\ \quad \quad \neg (\operatorname{uch}(h) = 0) \end{array}}}
\vnbprop{ccont-character-cover}{\vnbstmt{\begin{array}{@{}l@{}} \forall t.; \forall h \in \operatorname{characters}(\operatorname{ccont\text{-}algebra}(t)). \\ \quad (\operatorname{is\text{-}compact\text{-}t}(t) \wedge \forall x \in \operatorname{pts}(t).\; \exists g \in C({t}).\; ((h(g) = 0) \wedge \neg (g(x) = 0))) \Rightarrow  \\ \quad \quad \operatorname{pts}(t) \\ \quad \quad \in\; \{w \in \operatorname{power}(\operatorname{pts}(t)) : \exists \mathit{uhh} \in C({t}).\; ((h(\mathit{uhh}) = 0) \wedge (\forall z \in \operatorname{pts}(t).\; ((\operatorname{uhh}(z) \in \mathbb{R}) \wedge (0 \leq \operatorname{uhh}(z))) \wedge \forall z \in w.\; (1 \leq \operatorname{uhh}(z))))\} \end{array}}}
\vnbprop{ccont-character-kernel-point}{\vnbstmt{\begin{array}{@{}l@{}} \forall t.; \forall h \in \operatorname{characters}(\operatorname{ccont\text{-}algebra}(t)). \\ \quad \operatorname{is\text{-}compact\text{-}t}(t) \Rightarrow  \\ \quad \quad \exists x \in \operatorname{pts}(t).\; \forall g \in C({t}).\; ((h(g) = 0) \Rightarrow (g(x) = 0)) \end{array}}}
\vnbprop{ccont-character-is-point-eval}{\vnbstmt{\begin{array}{@{}l@{}} \forall t.; \forall h \in \operatorname{characters}(\operatorname{ccont\text{-}algebra}(t)). \\ \quad \operatorname{is\text{-}compact\text{-}t}(t) \Rightarrow \exists x \in \operatorname{pts}(t).\; (h = \operatorname{point\text{-}eval}(t, x)) \end{array}}}
\vnbprop{top-hausdorff-singleton-closed}{\vnbstmt{\forall t.; \forall x \in \operatorname{pts}(t).\; (\operatorname{is\text{-}hausdorff}(t) \Rightarrow \operatorname{is\text{-}top\text{-}closed}(t, \operatorname{singleton}(x)))}}
\vnbprop{point-eval-injective}{\vnbstmt{\begin{array}{@{}l@{}} \forall t.; \forall x, y \in \operatorname{pts}(t). \\ \quad (\operatorname{is\text{-}compact\text{-}t}(t) \wedge (\operatorname{is\text{-}hausdorff}(t) \wedge \neg (x = y))) \Rightarrow  \\ \quad \quad \neg (\operatorname{point\text{-}eval}(t, x) = \operatorname{point\text{-}eval}(t, y)) \end{array}}}

\vnbbookfile{theorem-library/basw4-cx}{theorem-library/basw4-cx}
\noindent{\small basw4-cx.scm -- wave of 2026-11-09/10 (agent BA-SW4), wired 2026-11-10 from scratchpad/basw4-1109/basw4-cx.scm. Library: operator-algebras.}\par
\vnbprop{ccont-algebra-uba}{\vnbstmt{\begin{array}{@{}l@{}} \forall t. \\ \quad (\operatorname{is\text{-}compact\text{-}t}(t) \wedge \exists p.\; (p \in \operatorname{pts}(t))) \Rightarrow  \\ \quad \quad \operatorname{is\text{-}unital\text{-}banach\text{-}algebra}(\operatorname{uiba\text{-}as\text{-}uba}(\operatorname{ccont\text{-}algebra}(t))) \end{array}}}
\vnbprop{ccont-uba-vec}{\vnbstmt{\begin{array}{@{}l@{}} \forall t. \\ \quad (\operatorname{is\text{-}compact\text{-}t}(t) \wedge \exists p.\; (p \in \operatorname{pts}(t))) \Rightarrow  \\ \quad \quad (\operatorname{vec}(\operatorname{uiba\text{-}as\text{-}uba}(\operatorname{ccont\text{-}algebra}(t))) \simeq C({t})) \end{array}}}
\vnbprop{ccont-uba-vnrm}{\vnbstmt{\begin{array}{@{}l@{}} \forall t. \\ \quad (\operatorname{is\text{-}compact\text{-}t}(t) \wedge \exists p.\; (p \in \operatorname{pts}(t))) \Rightarrow  \\ \quad \quad (\operatorname{vnrm}(\operatorname{uiba\text{-}as\text{-}uba}(\operatorname{ccont\text{-}algebra}(t))) \simeq \operatorname{ccont\text{-}vnrm}(t)) \end{array}}}
\vnbprop{ccont-characters-view}{\vnbstmt{\begin{array}{@{}l@{}} \forall t. \\ \quad (\operatorname{is\text{-}compact\text{-}t}(t) \wedge \exists p.\; (p \in \operatorname{pts}(t))) \Rightarrow  \\ \quad \quad (\operatorname{characters}(\operatorname{uiba\text{-}as\text{-}uba}(\operatorname{ccont\text{-}algebra}(t))) \simeq \operatorname{characters}(\operatorname{ccont\text{-}algebra}(t))) \end{array}}}
\vnbprop{ccont-character-bounded}{\vnbstmt{\begin{array}{@{}l@{}} \forall t.; \forall h \in \operatorname{characters}(\operatorname{ccont\text{-}algebra}(t)), f \in C({t}). \\ \quad (\operatorname{is\text{-}compact\text{-}t}(t) \wedge \exists p.\; (p \in \operatorname{pts}(t))) \Rightarrow  \\ \quad \quad \operatorname{magnitude}(h(f)) \leq \lVert f \rVert_{\infty} \end{array}}}
\vnbprop{uiba-characters-view}{\vnbstmt{\begin{array}{@{}l@{}} \forall a. \\ \quad \operatorname{is\text{-}unital\text{-}involutive\text{-}banach\text{-}algebra}(a) \Rightarrow  \\ \quad \quad (\operatorname{characters}(\operatorname{uiba\text{-}as\text{-}uba}(a)) \simeq \operatorname{characters}(a)) \end{array}}}
\vnbprop{ccont-ideal-bump}{\vnbstmt{\begin{array}{@{}l@{}} \forall t, i.; \forall g \in i, p \in \operatorname{pts}(t). \\ \quad \operatorname{is\text{-}compact\text{-}t}(t) \wedge \\ \quad i \subseteq C({t}) \wedge \\ \quad \operatorname{ccont\text{-}vzero}(t) \in i \wedge \\ \quad \forall f, \mathit{bmg} \in i.\; (\operatorname{ccont\text{-}vadd}(t)(f, \mathit{bmg}) \in i) \wedge \\ \quad \forall z \in \mathbb{C}, f \in i.\; (\operatorname{ccont\text{-}act}(t)(z, f) \in i) \wedge \\ \quad \forall \mathit{bmg} \in C({t}), f \in i.\; (\operatorname{ccont\text{-}vmul}(t)(f, \mathit{bmg}) \in i) \wedge \\ \quad \neg (g(p) = 0) \Rightarrow \\ \quad \quad \exists u \in \operatorname{opens}(t). \\ \quad \quad \quad p \in u \wedge \\ \quad \quad \quad u \\ \quad \quad \quad \in\; \{w \in \operatorname{power}(\operatorname{pts}(t)) : \exists h \in C({t}).\; ((h \in i) \wedge (\forall \mathit{uhz} \in \operatorname{pts}(t).\; ((h(\mathit{uhz}) \in \mathbb{R}) \wedge (0 \leq h(\mathit{uhz}))) \wedge \forall \mathit{uhz} \in w.\; (1 \leq h(\mathit{uhz}))))\} \end{array}}}
\vnbprop{ccont-ideal-one}{\vnbstmt{\begin{array}{@{}l@{}} \forall t, i.; \forall h \in i. \\ \quad \operatorname{is\text{-}compact\text{-}t}(t) \wedge \\ \quad i \subseteq C({t}) \wedge \\ \quad \operatorname{ccont\text{-}vzero}(t) \in i \wedge \\ \quad \forall f, g \in i.\; (\operatorname{ccont\text{-}vadd}(t)(f, g) \in i) \wedge \\ \quad \forall z \in \mathbb{C}, f \in i.\; (\operatorname{ccont\text{-}act}(t)(z, f) \in i) \wedge \\ \quad \forall g \in C({t}), f \in i.\; (\operatorname{ccont\text{-}vmul}(t)(f, g) \in i) \wedge \\ \quad \forall \mathit{ulz} \in \operatorname{pts}(t).\; ((h(\mathit{ulz}) \in \mathbb{R}) \wedge (1 \leq h(\mathit{ulz}))) \Rightarrow \\ \quad \quad \operatorname{ccont\text{-}vone}(t) \in i \end{array}}}
\vnbprop{ccont-ideal-cover}{\vnbstmt{\begin{array}{@{}l@{}} \forall t, i. \\ \quad \operatorname{is\text{-}compact\text{-}t}(t) \wedge \\ \quad i \subseteq C({t}) \wedge \\ \quad \operatorname{ccont\text{-}vzero}(t) \in i \wedge \\ \quad \forall f, g \in i.\; (\operatorname{ccont\text{-}vadd}(t)(f, g) \in i) \wedge \\ \quad \forall z \in \mathbb{C}, f \in i.\; (\operatorname{ccont\text{-}act}(t)(z, f) \in i) \wedge \\ \quad \forall g \in C({t}), f \in i.\; (\operatorname{ccont\text{-}vmul}(t)(f, g) \in i) \wedge \\ \quad \forall x \in \operatorname{pts}(t).\; \exists \mathit{ucg} \in i.\; \neg (\operatorname{ucg}(x) = 0) \Rightarrow \\ \quad \quad \operatorname{pts}(t) \\ \quad \quad \in\; \{w \in \operatorname{power}(\operatorname{pts}(t)) : \exists h \in C({t}).\; ((h \in i) \wedge (\forall \mathit{uhz} \in \operatorname{pts}(t).\; ((h(\mathit{uhz}) \in \mathbb{R}) \wedge (0 \leq h(\mathit{uhz}))) \wedge \forall \mathit{uhz} \in w.\; (1 \leq h(\mathit{uhz}))))\} \end{array}}}
\vnbprop{ccont-ideal-common-zero}{\vnbstmt{\begin{array}{@{}l@{}} \forall t, i. \\ \quad \operatorname{is\text{-}compact\text{-}t}(t) \wedge \\ \quad i \subseteq C({t}) \wedge \\ \quad \operatorname{ccont\text{-}vzero}(t) \in i \wedge \\ \quad \forall f, g \in i.\; (\operatorname{ccont\text{-}vadd}(t)(f, g) \in i) \wedge \\ \quad \forall z \in \mathbb{C}, f \in i.\; (\operatorname{ccont\text{-}act}(t)(z, f) \in i) \wedge \\ \quad \forall g \in C({t}), f \in i.\; (\operatorname{ccont\text{-}vmul}(t)(f, g) \in i) \wedge \\ \quad \neg (\operatorname{ccont\text{-}vone}(t) \in i) \Rightarrow \\ \quad \quad \exists x \in \operatorname{pts}(t).\; \forall \mathit{ukg} \in i.\; (\operatorname{ukg}(x) = 0) \end{array}}}
\vnbprop{ccont-ideal-unit-all}{\vnbstmt{\begin{array}{@{}l@{}} \forall t, i. \\ \quad \operatorname{is\text{-}compact\text{-}t}(t) \wedge \\ \quad i \subseteq C({t}) \wedge \\ \quad \operatorname{ccont\text{-}vzero}(t) \in i \wedge \\ \quad \forall f, g \in i.\; (\operatorname{ccont\text{-}vadd}(t)(f, g) \in i) \wedge \\ \quad \forall z \in \mathbb{C}, f \in i.\; (\operatorname{ccont\text{-}act}(t)(z, f) \in i) \wedge \\ \quad \forall g \in C({t}), f \in i.\; (\operatorname{ccont\text{-}vmul}(t)(f, g) \in i) \wedge \\ \quad \operatorname{ccont\text{-}vone}(t) \in i \Rightarrow \\ \quad \quad i = C({t}) \end{array}}}
\vnbprop{ccont-ideal-laws}{\vnbstmt{\begin{array}{@{}l@{}} \forall t, i. \\ \quad (\operatorname{is\text{-}compact\text{-}t}(t) \wedge \operatorname{is\text{-}two\text{-}sided\text{-}ideal}(\operatorname{ccont\text{-}algebra}(t), i)) \Rightarrow  \\ \quad \quad i \subseteq C({t}) \wedge \\ \quad \quad \operatorname{ccont\text{-}vzero}(t) \in i \wedge \\ \quad \quad \forall f, g \in i.\; (\operatorname{ccont\text{-}vadd}(t)(f, g) \in i) \wedge \\ \quad \quad \forall z \in \mathbb{C}, f \in i.\; (\operatorname{ccont\text{-}act}(t)(z, f) \in i) \wedge \\ \quad \quad \forall g \in C({t}), f \in i.\; (\operatorname{ccont\text{-}vmul}(t)(f, g) \in i) \end{array}}}
\vnbprop{ccont-point-ideal-two-sided}{\vnbstmt{\begin{array}{@{}l@{}} \forall t.; \forall x \in \operatorname{pts}(t). \\ \quad \operatorname{is\text{-}compact\text{-}t}(t) \Rightarrow  \\ \quad \quad \operatorname{is\text{-}two\text{-}sided\text{-}ideal}(\operatorname{ccont\text{-}algebra}(t), \{h \in C({t}) : (h(x) = 0)\}) \end{array}}}
\vnbprop{ccont-point-ideal-proper}{\vnbstmt{\begin{array}{@{}l@{}} \forall t.; \forall x \in \operatorname{pts}(t). \\ \quad \operatorname{is\text{-}compact\text{-}t}(t) \Rightarrow \neg (\{h \in C({t}) : (h(x) = 0)\} = C({t})) \end{array}}}
\vnbprop{ccont-point-ideal-maximal}{\vnbstmt{\begin{array}{@{}l@{}} \forall t.; \forall x \in \operatorname{pts}(t). \\ \quad \operatorname{is\text{-}compact\text{-}t}(t) \Rightarrow  \\ \quad \quad \operatorname{is\text{-}maximal\text{-}ideal}(\operatorname{ccont\text{-}algebra}(t), \{h \in C({t}) : (h(x) = 0)\}) \end{array}}}
\vnbprop{ccont-maximal-ideal-is-point}{\vnbstmt{\begin{array}{@{}l@{}} \forall t, i. \\ \quad (\operatorname{is\text{-}compact\text{-}t}(t) \wedge \operatorname{is\text{-}maximal\text{-}ideal}(\operatorname{ccont\text{-}algebra}(t), i)) \Rightarrow  \\ \quad \quad \exists x \in \operatorname{pts}(t).\; (i = \{h \in C({t}) : (h(x) = 0)\}) \end{array}}}
\vnbprop{ccont-maximal-ideals}{\vnbstmt{\begin{array}{@{}l@{}} \forall t, i. \\ \quad \operatorname{is\text{-}compact\text{-}t}(t) \Rightarrow  \\ \quad \quad \operatorname{is\text{-}maximal\text{-}ideal}(\operatorname{ccont\text{-}algebra}(t), i) \Leftrightarrow \exists x \in \operatorname{pts}(t).\; (i = \{h \in C({t}) : (h(x) = 0)\}) \end{array}}}
\vnbprop{ccont-point-ideal-injective}{\vnbstmt{\begin{array}{@{}l@{}} \forall t.; \forall x, y \in \operatorname{pts}(t). \\ \quad (\operatorname{is\text{-}compact\text{-}t}(t) \wedge (\operatorname{is\text{-}hausdorff}(t) \wedge \neg (x = y))) \Rightarrow  \\ \quad \quad \neg (\{h \in C({t}) : (h(x) = 0)\} = \{h \in C({t}) : (h(y) = 0)\}) \end{array}}}

\vnbbookfile{theorem-library/ba-gelfand-formula}{theorem-library/ba-gelfand-formula}
\noindent{\small ba-gelfand-formula.scm -- wave of 2026-11-09/10 (agent BA-1H), wired 2026-11-10 from scratchpad/ba1h-1109/ba-gelfand-formula.scm. Library: operator-algebras.}\par
\vnbprop{bgf-psi-zero}{\vnbstmt{\begin{array}{@{}l@{}} \forall a, p. \\ \quad (\operatorname{is\text{-}unital\text{-}banach\text{-}algebra}(a) \wedge \operatorname{is\text{-}cc\text{-}linear\text{-}on}(a, p)) \Rightarrow (p(\operatorname{vzero}(a)) = 0) \end{array}}}
\vnbprop{bgf-psi-nsum}{\vnbstmt{\begin{array}{@{}l@{}} \forall n \in \mathbb{N}, a, p.; \forall y \in \operatorname{vec}(a), h \in (\mathbb{N} \to \mathbb{C}). \\ \quad \operatorname{is\text{-}unital\text{-}banach\text{-}algebra}(a) \wedge \\ \quad \operatorname{is\text{-}cc\text{-}linear\text{-}on}(a, p) \wedge \\ \quad \forall k \in \mathbb{N}.\; (h(k) = p({y}^{k})) \Rightarrow \\ \quad \quad \operatorname{cc\text{-}series\text{-}partial\text{-}sum}(h, n) = p(\sum_{k < n} {y}^{k}) \end{array}}}
\vnbprop{bgf-psi-conv}{\vnbstmt{\begin{array}{@{}l@{}} \forall a, p.; \forall s \in (\mathbb{N} \to \operatorname{vec}(a)), l \in \operatorname{vec}(a), h \in (\mathbb{N} \to \mathbb{C}). \\ \quad \operatorname{is\text{-}unital\text{-}banach\text{-}algebra}(a) \wedge \\ \quad \operatorname{is\text{-}cc\text{-}linear\text{-}on}(a, p) \wedge \\ \quad \forall v \in \operatorname{vec}(a).\; (\operatorname{magnitude}(p(v)) \leq \operatorname{vnrm}(a)(v)) \wedge \\ \quad \operatorname{converges\text{-}to}(\operatorname{ba\text{-}metric\text{-}space}(a), s, l) \wedge \\ \quad \forall n \in \mathbb{N}.\; (h(n) = p(s(n))) \Rightarrow \\ \quad \quad \operatorname{converges\text{-}to}(\operatorname{cc\text{-}ms}, h, p(l)) \end{array}}}
\vnbprop{bgf-coef-family}{\vnbstmt{\begin{array}{@{}l@{}} \forall a, p.; \forall x \in \operatorname{vec}(a). \\ \quad (\operatorname{is\text{-}unital\text{-}banach\text{-}algebra}(a) \wedge \operatorname{is\text{-}cc\text{-}linear\text{-}on}(a, p)) \Rightarrow  \\ \quad \quad \exists c \in (\mathbb{Z} \to \mathbb{C}). \\ \quad \quad \quad (\forall k \in \mathbb{N}.\; (c(k) = 0) \wedge \forall k \in \mathbb{N}.\; (c(-(k + 1)) = p({x}^{k}))) \end{array}}}
\vnbprop{bgf-laurent-resolvent}{\vnbstmt{\begin{array}{@{}l@{}} \forall a, p.; \forall x \in \operatorname{vec}(a), c \in (\mathbb{Z} \to \mathbb{C}), z \in \mathbb{C}. \\ \quad \operatorname{is\text{-}unital\text{-}banach\text{-}algebra}(a) \wedge \\ \quad \operatorname{is\text{-}cc\text{-}linear\text{-}on}(a, p) \wedge \\ \quad \forall v \in \operatorname{vec}(a).\; (\operatorname{magnitude}(p(v)) \leq \operatorname{vnrm}(a)(v)) \wedge \\ \quad \forall k \in \mathbb{N}.\; (c(k) = 0) \wedge \\ \quad \forall k \in \mathbb{N}.\; (c(-(k + 1)) = p({x}^{k})) \wedge \\ \quad \operatorname{vnrm}(a)(x) < \operatorname{magnitude}(z) \Rightarrow \\ \quad \quad \operatorname{laurent\text{-}sum}(c, 0, z) = p({\operatorname{shift\text{-}el}(a, z, x)}^{-1}) \end{array}}}
\vnbprop{bgf-functional-estimate}{\vnbstmt{\begin{array}{@{}l@{}} \forall a, p.; \forall x \in \operatorname{vec}(a), s, m \in \mathbb{R}, n \in \mathbb{N}. \\ \quad \operatorname{is\text{-}unital\text{-}banach\text{-}algebra}(a) \wedge \\ \quad \operatorname{is\text{-}cc\text{-}linear\text{-}on}(a, p) \wedge \\ \quad \forall v \in \operatorname{vec}(a).\; (\operatorname{magnitude}(p(v)) \leq \operatorname{vnrm}(a)(v)) \wedge \\ \quad 0 < s \wedge \\ \quad \forall z \in \mathbb{C}. \\ \quad \quad (s \leq \operatorname{magnitude}(z)) \Rightarrow \operatorname{is\text{-}invertible\text{-}el}(a, \operatorname{shift\text{-}el}(a, z, x)) \wedge \\ \quad 0 \leq m \wedge \\ \quad \forall z \in \mathbb{C}. \\ \quad \quad (\operatorname{magnitude}(z) = s) \Rightarrow (\operatorname{vnrm}(a)({\operatorname{shift\text{-}el}(a, z, x)}^{-1}) \leq m) \Rightarrow \\ \quad \quad \operatorname{magnitude}(p({x}^{n})) \leq ((s \cdot \operatorname{power}(s, n)) \cdot m) \end{array}}}
\vnbprop{bgf-vnrm-reverse}{\vnbstmt{\begin{array}{@{}l@{}} \forall a.; \forall u, v \in \operatorname{vec}(a). \\ \quad \operatorname{is\text{-}unital\text{-}banach\text{-}algebra}(a) \Rightarrow  \\ \quad \quad \lvert (\operatorname{vnrm}(a)(u) - \operatorname{vnrm}(a)(v)) \rvert \leq \operatorname{vnrm}(a)(\operatorname{vadd}(a)(u, \operatorname{vneg}(a)(v))) \end{array}}}
\vnbprop{bgf-resolvent-circle-bound}{\vnbstmt{\begin{array}{@{}l@{}} \forall a.; \forall x \in \operatorname{vec}(a), s \in \mathbb{R}. \\ \quad \operatorname{is\text{-}unital\text{-}banach\text{-}algebra}(a) \wedge \\ \quad 0 < s \wedge \\ \quad \forall z \in \mathbb{C}. \\ \quad \quad (s \leq \operatorname{magnitude}(z)) \Rightarrow \operatorname{is\text{-}invertible\text{-}el}(a, \operatorname{shift\text{-}el}(a, z, x)) \Rightarrow \\ \quad \quad \exists m \in \mathbb{R}. \\ \quad \quad \quad 0 \leq m \wedge \\ \quad \quad \quad \forall z \in \mathbb{C}. \\ \quad \quad \quad \quad (\operatorname{magnitude}(z) = s) \Rightarrow (\operatorname{vnrm}(a)({\operatorname{shift\text{-}el}(a, z, x)}^{-1}) \leq m) \end{array}}}
\vnbprop{bgf-power-estimate}{\vnbstmt{\begin{array}{@{}l@{}} \forall a.; \forall x \in \operatorname{vec}(a), s, m \in \mathbb{R}, n \in \mathbb{N}. \\ \quad \operatorname{is\text{-}unital\text{-}banach\text{-}algebra}(a) \wedge \\ \quad 0 < s \wedge \\ \quad \forall z \in \mathbb{C}. \\ \quad \quad (s \leq \operatorname{magnitude}(z)) \Rightarrow \operatorname{is\text{-}invertible\text{-}el}(a, \operatorname{shift\text{-}el}(a, z, x)) \wedge \\ \quad 0 \leq m \wedge \\ \quad \forall z \in \mathbb{C}. \\ \quad \quad (\operatorname{magnitude}(z) = s) \Rightarrow (\operatorname{vnrm}(a)({\operatorname{shift\text{-}el}(a, z, x)}^{-1}) \leq m) \Rightarrow \\ \quad \quad \operatorname{vnrm}(a)({x}^{n}) \leq ((s \cdot \operatorname{power}(s, n)) \cdot m) \end{array}}}
\vnbprop{bgf-radius-le-root}{\vnbstmt{\begin{array}{@{}l@{}} \forall a.; \forall x \in \operatorname{vec}(a), n \in \mathbb{N}. \\ \quad (\operatorname{is\text{-}unital\text{-}banach\text{-}algebra}(a) \wedge (1 \leq n)) \Rightarrow  \\ \quad \quad \rho({x}) \leq {\operatorname{vnrm}(a)({x}^{n})}^{{n}^{\ast\!-1}} \end{array}}}
\vnbprop{bgf-radius-from-estimate}{\vnbstmt{\begin{array}{@{}l@{}} \forall a.; \forall x \in \operatorname{vec}(a), s, m, t \in \mathbb{R}. \\ \quad \operatorname{is\text{-}unital\text{-}banach\text{-}algebra}(a) \wedge \\ \quad 0 < s \wedge \\ \quad 0 \leq m \wedge \\ \quad \forall n \in \mathbb{N}.\; (\operatorname{vnrm}(a)({x}^{n}) \leq ((s \cdot \operatorname{power}(s, n)) \cdot m)) \wedge \\ \quad s < t \Rightarrow \\ \quad \quad \rho({x}) \leq t \end{array}}}
\vnbprop{gelfand-spectral-radius-le}{\vnbstmt{\begin{array}{@{}l@{}} \forall a.; \forall x \in \operatorname{vec}(a), s \in \mathbb{R}. \\ \quad \operatorname{is\text{-}unital\text{-}banach\text{-}algebra}(a) \wedge \\ \quad 0 < s \wedge \\ \quad \forall z \in \mathbb{C}. \\ \quad \quad (s \leq \operatorname{magnitude}(z)) \Rightarrow \operatorname{is\text{-}invertible\text{-}el}(a, \operatorname{shift\text{-}el}(a, z, x)) \Rightarrow \\ \quad \quad \rho({x}) \leq s \end{array}}}
\vnbprop{gelfand-spectral-radius-max}{\vnbstmt{\begin{array}{@{}l@{}} \forall a.; \forall x \in \operatorname{vec}(a). \\ \quad \operatorname{is\text{-}unital\text{-}banach\text{-}algebra}(a) \Rightarrow  \\ \quad \quad \exists b \in \sigma({x}). \\ \quad \quad \quad ((\operatorname{magnitude}(b) = \rho({x})) \wedge \forall z \in \sigma({x}).\; (\operatorname{magnitude}(z) \leq \operatorname{magnitude}(b))) \end{array}}}
\vnbprop{gelfand-formula}{\vnbstmt{\begin{array}{@{}l@{}} \forall a.; \forall x \in \operatorname{vec}(a). \\ \quad \operatorname{is\text{-}unital\text{-}banach\text{-}algebra}(a) \Rightarrow  \\ \quad \quad \exists b \in \sigma({x}). \\ \quad \quad \quad \forall z \in \sigma({x}).\; (\operatorname{magnitude}(z) \leq \operatorname{magnitude}(b)) \wedge \\ \quad \quad \quad \operatorname{converges\text{-}to}(\operatorname{rr\text{-}ms}, \\ \quad \quad \quad \quad (\lambda n \in \mathbb{N}.\; {\operatorname{vnrm}(a)({x}^{(n + 1)})}^{{(n + 1)}^{\ast\!-1}}), \\ \quad \quad \quad \quad \operatorname{magnitude}(b)) \end{array}}}

\vnbbookfile{theorem-library/ba2b-radius}{theorem-library/ba2b-radius}
\noindent{\small ba2b-radius.scm -- wave of 2026-11-09/10 (agent BA-2B), wired 2026-11-10 from scratchpad/ba2b-1109/ba2b-radius.scm. Library: operator-algebras.}\par
\vnbprop{b2r-el-pow-view}{\vnbstmt{\begin{array}{@{}l@{}} \forall n \in \mathbb{N}, a.; \forall x \in \operatorname{vec}(a). \\ \quad \operatorname{is\text{-}unital\text{-}involutive\text{-}banach\text{-}algebra}(a) \Rightarrow  \\ \quad \quad ({x}^{n} \simeq {x}^{n}) \end{array}}}
\vnbprop{b2r-spectral-radius-view}{\vnbstmt{\begin{array}{@{}l@{}} \forall a.; \forall x \in \operatorname{vec}(a). \\ \quad \operatorname{is\text{-}unital\text{-}involutive\text{-}banach\text{-}algebra}(a) \Rightarrow  \\ \quad \quad (\rho({x}) \simeq \rho({x})) \end{array}}}
\vnbprop{b2r-el-pow-sa}{\vnbstmt{\begin{array}{@{}l@{}} \forall n \in \mathbb{N}, a, h. \\ \quad (\operatorname{is\text{-}unital\text{-}involutive\text{-}banach\text{-}algebra}(a) \wedge \operatorname{is\text{-}self\text{-}adjoint\text{-}el}(a, h)) \Rightarrow  \\ \quad \quad {}^{*_{a}}({h}^{n}) = {h}^{n} \end{array}}}
\vnbprop{b2r-el-pow-double}{\vnbstmt{\begin{array}{@{}l@{}} \forall n \in \mathbb{N}, a, h. \\ \quad (\operatorname{is\text{-}unital\text{-}c\text{-}star\text{-}algebra}(a) \wedge \operatorname{is\text{-}self\text{-}adjoint\text{-}el}(a, h)) \Rightarrow  \\ \quad \quad \operatorname{vnrm}(a)({h}^{(n + n)}) = (\operatorname{vnrm}(a)({h}^{n}) \cdot \operatorname{vnrm}(a)({h}^{n})) \end{array}}}
\vnbprop{b2r-sa-pow-frequent}{\vnbstmt{\begin{array}{@{}l@{}} \forall b \in \mathbb{N}, a, h. \\ \quad (\operatorname{is\text{-}unital\text{-}c\text{-}star\text{-}algebra}(a) \wedge (\operatorname{is\text{-}self\text{-}adjoint\text{-}el}(a, h) \wedge (0 < \operatorname{vnrm}(a)(h)))) \Rightarrow  \\ \quad \quad \exists k \in \mathbb{N}. \\ \quad \quad \quad (b \leq k) \wedge (\operatorname{vnrm}(a)({h}^{(k + 1)}) = {\operatorname{vnrm}(a)(h)}^{(k + 1)}) \end{array}}}
\vnbprop{cstar-sa-spectral-radius}{\vnbstmt{\begin{array}{@{}l@{}} \forall a, h. \\ \quad (\operatorname{is\text{-}unital\text{-}c\text{-}star\text{-}algebra}(a) \wedge \operatorname{is\text{-}self\text{-}adjoint\text{-}el}(a, h)) \Rightarrow  \\ \quad \quad \rho({h}) = \operatorname{vnrm}(a)(h) \end{array}}}

\vnbbookfile{theorem-library/ba2b-hom}{theorem-library/ba2b-hom}
\noindent{\small ba2b-hom.scm -- wave of 2026-11-09/10 (agent BA-2B), wired 2026-11-10 from scratchpad/ba2b-1109/ba2b-hom.scm. Library: operator-algebras.}\par
\vnbprop{b2h-hom-zero}{\vnbstmt{\begin{array}{@{}l@{}} \forall a, b, f. \\ \quad (\operatorname{is\text{-}unital\text{-}banach\text{-}algebra}(a) \wedge (\operatorname{is\text{-}unital\text{-}banach\text{-}algebra}(b) \wedge \operatorname{is\text{-}alg\text{-}hom}(a, b, f))) \Rightarrow  \\ \quad \quad f(\operatorname{vzero}(a)) = \operatorname{vzero}(b) \end{array}}}
\vnbprop{b2h-hom-neg}{\vnbstmt{\begin{array}{@{}l@{}} \forall a, b, f.; \forall x \in \operatorname{vec}(a). \\ \quad (\operatorname{is\text{-}unital\text{-}banach\text{-}algebra}(a) \wedge (\operatorname{is\text{-}unital\text{-}banach\text{-}algebra}(b) \wedge \operatorname{is\text{-}alg\text{-}hom}(a, b, f))) \Rightarrow  \\ \quad \quad f(\operatorname{vneg}(a)(x)) = \operatorname{vneg}(b)(f(x)) \end{array}}}
\vnbprop{b2h-hom-shift}{\vnbstmt{\begin{array}{@{}l@{}} \forall a, b, f.; \forall z \in \mathbb{C}, x \in \operatorname{vec}(a). \\ \quad (\operatorname{is\text{-}unital\text{-}banach\text{-}algebra}(a) \wedge (\operatorname{is\text{-}unital\text{-}banach\text{-}algebra}(b) \wedge \operatorname{is\text{-}alg\text{-}hom}(a, b, f))) \Rightarrow  \\ \quad \quad f(\operatorname{shift\text{-}el}(a, z, x)) = \operatorname{shift\text{-}el}(b, z, f(x)) \end{array}}}
\vnbprop{b2h-hom-invertible}{\vnbstmt{\begin{array}{@{}l@{}} \forall a, b, f, y. \\ \quad \operatorname{is\text{-}unital\text{-}banach\text{-}algebra}(a) \wedge \\ \quad \operatorname{is\text{-}unital\text{-}banach\text{-}algebra}(b) \wedge \\ \quad \operatorname{is\text{-}alg\text{-}hom}(a, b, f) \wedge \\ \quad \operatorname{is\text{-}invertible\text{-}el}(a, y) \Rightarrow \\ \quad \quad \operatorname{is\text{-}invertible\text{-}el}(b, f(y)) \end{array}}}
\vnbprop{hom-spectral-radius-le}{\vnbstmt{\begin{array}{@{}l@{}} \forall a, b, f.; \forall x \in \operatorname{vec}(a). \\ \quad (\operatorname{is\text{-}unital\text{-}banach\text{-}algebra}(a) \wedge (\operatorname{is\text{-}unital\text{-}banach\text{-}algebra}(b) \wedge \operatorname{is\text{-}alg\text{-}hom}(a, b, f))) \Rightarrow  \\ \quad \quad \rho({f(x)}) \leq \rho({x}) \end{array}}}
\vnbprop{b2h-alg-hom-view}{\vnbstmt{\begin{array}{@{}l@{}} \forall a, b, f. \\ \quad \operatorname{is\text{-}unital\text{-}involutive\text{-}banach\text{-}algebra}(a) \wedge \\ \quad \operatorname{is\text{-}unital\text{-}involutive\text{-}banach\text{-}algebra}(b) \wedge \\ \quad \operatorname{is\text{-}alg\text{-}hom}(a, b, f) \Rightarrow \\ \quad \quad \operatorname{is\text{-}alg\text{-}hom}(\operatorname{uiba\text{-}as\text{-}uba}(a), \operatorname{uiba\text{-}as\text{-}uba}(b), f) \end{array}}}
\vnbprop{b2h-hom-sa-contractive}{\vnbstmt{\begin{array}{@{}l@{}} \forall a, b, f, h. \\ \quad \operatorname{is\text{-}unital\text{-}c\text{-}star\text{-}algebra}(a) \wedge \\ \quad \operatorname{is\text{-}unital\text{-}c\text{-}star\text{-}algebra}(b) \wedge \\ \quad \operatorname{is\text{-}star\text{-}hom}(a, b, f) \wedge \\ \quad \operatorname{is\text{-}self\text{-}adjoint\text{-}el}(a, h) \Rightarrow \\ \quad \quad \operatorname{vnrm}(b)(f(h)) \leq \operatorname{vnrm}(a)(h) \end{array}}}
\vnbprop{cstar-hom-contractive}{\vnbstmt{\begin{array}{@{}l@{}} \forall a, b, f.; \forall x \in \operatorname{vec}(a). \\ \quad (\operatorname{is\text{-}unital\text{-}c\text{-}star\text{-}algebra}(a) \wedge (\operatorname{is\text{-}unital\text{-}c\text{-}star\text{-}algebra}(b) \wedge \operatorname{is\text{-}star\text{-}hom}(a, b, f))) \Rightarrow  \\ \quad \quad \operatorname{vnrm}(b)(f(x)) \leq \operatorname{vnrm}(a)(x) \end{array}}}
\vnbprop{cstar-hom-norm-one}{\vnbstmt{\begin{array}{@{}l@{}} \forall a, b, f. \\ \quad (\operatorname{is\text{-}unital\text{-}c\text{-}star\text{-}algebra}(a) \wedge (\operatorname{is\text{-}unital\text{-}c\text{-}star\text{-}algebra}(b) \wedge \operatorname{is\text{-}star\text{-}hom}(a, b, f))) \Rightarrow  \\ \quad \quad \operatorname{vnrm}(b)(f(1_{a})) = 1 \end{array}}}

\vnbbookfile{theorem-library/ba2b-gelfand}{theorem-library/ba2b-gelfand}
\noindent{\small ba2b-gelfand.scm -- wave of 2026-11-09/10 (agent BA-2B), wired 2026-11-10 from scratchpad/ba2b-1109/ba2b-gelfand.scm. Library: operator-algebras.}\par
\vnbprop{gelfand-space-pts}{\vnbstmt{\begin{array}{@{}l@{}} \forall a. \\ \quad \operatorname{is\text{-}unital\text{-}banach\text{-}algebra}(a) \Rightarrow (\operatorname{pts}(\widehat{a}) \simeq \operatorname{characters}(a)) \end{array}}}
\vnbprop{b2g-gelfand-transform-view}{\vnbstmt{\begin{array}{@{}l@{}} \forall a. \\ \quad \operatorname{is\text{-}unital\text{-}involutive\text{-}banach\text{-}algebra}(a) \Rightarrow  \\ \quad \quad (\operatorname{gelfand\text{-}transform}(\operatorname{uiba\text{-}as\text{-}uba}(a)) \simeq \operatorname{gelfand\text{-}transform}(a)) \end{array}}}
\vnbprop{b2g-gelfand-space-view}{\vnbstmt{\begin{array}{@{}l@{}} \forall a. \\ \quad \operatorname{is\text{-}unital\text{-}involutive\text{-}banach\text{-}algebra}(a) \Rightarrow  \\ \quad \quad (\widehat{\operatorname{uiba\text{-}as\text{-}uba}(a)} \simeq \widehat{a}) \end{array}}}
\vnbprop{b2g-ccsa-commutative-view}{\vnbstmt{\begin{array}{@{}l@{}} \forall a. \\ \quad \operatorname{is\text{-}commutative\text{-}c\text{-}star\text{-}algebra}(a) \Rightarrow  \\ \quad \quad \operatorname{is\text{-}commutative\text{-}banach\text{-}algebra}(\operatorname{uiba\text{-}as\text{-}uba}(a)) \end{array}}}
\vnbprop{gelfand-sup-le-norm}{\vnbstmt{\begin{array}{@{}l@{}} \forall a.; \forall x \in \operatorname{vec}(a). \\ \quad \operatorname{is\text{-}unital\text{-}banach\text{-}algebra}(a) \Rightarrow  \\ \quad \quad \lVert \operatorname{gelfand\text{-}transform}(a)(x) \rVert_{\infty} \in \mathbb{R} \wedge \\ \quad \quad 0 \leq \lVert \operatorname{gelfand\text{-}transform}(a)(x) \rVert_{\infty} \wedge \\ \quad \quad \lVert \operatorname{gelfand\text{-}transform}(a)(x) \rVert_{\infty} \leq \operatorname{vnrm}(a)(x) \end{array}}}
\vnbprop{gelfand-sup-spectral-radius}{\vnbstmt{\begin{array}{@{}l@{}} \forall a.; \forall x \in \operatorname{vec}(a). \\ \quad \operatorname{is\text{-}commutative\text{-}banach\text{-}algebra}(a) \Rightarrow  \\ \quad \quad \lVert \operatorname{gelfand\text{-}transform}(a)(x) \rVert_{\infty} = \rho({x}) \end{array}}}
\vnbprop{gelfand-transform-in-ccont}{\vnbstmt{\begin{array}{@{}l@{}} \forall a.; \forall x \in \operatorname{vec}(a). \\ \quad \operatorname{is\text{-}unital\text{-}banach\text{-}algebra}(a) \Rightarrow  \\ \quad \quad \operatorname{gelfand\text{-}transform}(a)(x) \in C({\widehat{a}}) \end{array}}}
\vnbprop{gelfand-isometric}{\vnbstmt{\begin{array}{@{}l@{}} \forall a.; \forall x \in \operatorname{vec}(a). \\ \quad \operatorname{is\text{-}commutative\text{-}c\text{-}star\text{-}algebra}(a) \Rightarrow  \\ \quad \quad \lVert \operatorname{gelfand\text{-}transform}(a)(x) \rVert_{\infty} = \operatorname{vnrm}(a)(x) \end{array}}}

\vnbbookfile{theorem-library/ba2b-unique}{theorem-library/ba2b-unique}
\noindent{\small ba2b-unique.scm -- wave of 2026-11-09/10 (agent BA-2B), wired 2026-11-10 from scratchpad/ba2b-1109/ba2b-unique.scm. Library: operator-algebras.}\par
\vnbprop{b2u-id-star-hom}{\vnbstmt{\begin{array}{@{}l@{}} \forall a, c. \\ \quad \operatorname{is\text{-}unital\text{-}involutive\text{-}banach\text{-}algebra}(a) \wedge \\ \quad \operatorname{is\text{-}unital\text{-}involutive\text{-}banach\text{-}algebra}(c) \wedge \\ \quad (\operatorname{vec}(c) \simeq \operatorname{vec}(a)) \wedge \\ \quad (\operatorname{vadd}(c) \simeq \operatorname{vadd}(a)) \wedge \\ \quad (\operatorname{act}(c) \simeq \operatorname{act}(a)) \wedge \\ \quad (\cdot_{c} \simeq \cdot_{a}) \wedge \\ \quad (1_{c} \simeq 1_{a}) \wedge \\ \quad ({}^{*_{c}} \simeq {}^{*_{a}}) \Rightarrow \\ \quad \quad \operatorname{is\text{-}star\text{-}hom}(a, c, (\lambda u \in \operatorname{vec}(a).\; u)) \end{array}}}
\vnbprop{cstar-norm-unique}{\vnbstmt{\begin{array}{@{}l@{}} \forall a, c.; \forall x \in \operatorname{vec}(a). \\ \quad \operatorname{is\text{-}unital\text{-}c\text{-}star\text{-}algebra}(a) \wedge \\ \quad \operatorname{is\text{-}unital\text{-}c\text{-}star\text{-}algebra}(c) \wedge \\ \quad (\operatorname{vec}(c) \simeq \operatorname{vec}(a)) \wedge \\ \quad (\operatorname{vadd}(c) \simeq \operatorname{vadd}(a)) \wedge \\ \quad (\operatorname{act}(c) \simeq \operatorname{act}(a)) \wedge \\ \quad (\cdot_{c} \simeq \cdot_{a}) \wedge \\ \quad (1_{c} \simeq 1_{a}) \wedge \\ \quad ({}^{*_{c}} \simeq {}^{*_{a}}) \Rightarrow \\ \quad \quad \operatorname{vnrm}(c)(x) = \operatorname{vnrm}(a)(x) \end{array}}}
\vnbprop{unitization-cstar-norm-unique}{\vnbstmt{\begin{array}{@{}l@{}} \forall a, b, c.; \forall u \in \operatorname{unit\text{-}vec}(a). \\ \quad \operatorname{is\text{-}unital\text{-}c\text{-}star\text{-}algebra}({[}\operatorname{normed\text{-}field\text{-}as\text{-}commutative\text{-}ring}(\operatorname{cc\text{-}normed\text{-}field}), \operatorname{unit\text{-}vec}(a), \operatorname{unit\text{-}vadd}(a), {[}0, \operatorname{vzero}(a)], \operatorname{unit\text{-}vneg}(a), \operatorname{unit\text{-}act}(a), b, \operatorname{unit\text{-}vmul}(a), {[}1, \operatorname{vzero}(a)], \operatorname{unit\text{-}istar}(a)]) \wedge \\ \quad \operatorname{is\text{-}unital\text{-}c\text{-}star\text{-}algebra}({[}\operatorname{normed\text{-}field\text{-}as\text{-}commutative\text{-}ring}(\operatorname{cc\text{-}normed\text{-}field}), \operatorname{unit\text{-}vec}(a), \operatorname{unit\text{-}vadd}(a), {[}0, \operatorname{vzero}(a)], \operatorname{unit\text{-}vneg}(a), \operatorname{unit\text{-}act}(a), c, \operatorname{unit\text{-}vmul}(a), {[}1, \operatorname{vzero}(a)], \operatorname{unit\text{-}istar}(a)]) \Rightarrow \\ \quad \quad c(u) = b(u) \end{array}}}

\vnbbookfile{structure-library/functoriality}{structure-library/functoriality}
\noindent{\small A file shared across books; the entries that belong to this one.}\par
\noindent{\small functoriality.scm -- every view-as is a FUNCTOR, and here is the proof.}\par
\vnbprop{cna-as-module-functorial}{\vnbstmt{\begin{array}{@{}l@{}} \forall a, b, f. \\ \quad \operatorname{is\text{-}hom\text{-}complex\text{-}normed\text{-}algebra}(a, b, f) \Rightarrow  \\ \quad \quad \operatorname{is\text{-}hom\text{-}module}(\operatorname{cna\text{-}as\text{-}module}(a), \operatorname{cna\text{-}as\text{-}module}(b), f) \end{array}}}
\vnbprop{cna-as-normed-ag-functorial}{\vnbstmt{\begin{array}{@{}l@{}} \forall a, b, f. \\ \quad \operatorname{is\text{-}hom\text{-}complex\text{-}normed\text{-}algebra}(a, b, f) \Rightarrow  \\ \quad \quad \operatorname{is\text{-}hom\text{-}normed\text{-}ag}(\operatorname{cna\text{-}as\text{-}normed\text{-}ag}(a), \operatorname{cna\text{-}as\text{-}normed\text{-}ag}(b), f) \end{array}}}
\vnbprop{nuiba-as-cba-functorial}{\vnbstmt{\begin{array}{@{}l@{}} \forall a, b, f. \\ \quad \operatorname{is\text{-}hom\text{-}nu\text{-}involutive\text{-}banach\text{-}algebra}(a, b, f) \Rightarrow  \\ \quad \quad \operatorname{is\text{-}hom\text{-}complex\text{-}banach\text{-}algebra}(\operatorname{nuiba\text{-}as\text{-}cba}(a), \operatorname{nuiba\text{-}as\text{-}cba}(b), f) \end{array}}}
\vnbprop{uba-as-cba-functorial}{\vnbstmt{\begin{array}{@{}l@{}} \forall a, b, f. \\ \quad \operatorname{is\text{-}hom\text{-}unital\text{-}banach\text{-}algebra}(a, b, f) \Rightarrow  \\ \quad \quad \operatorname{is\text{-}hom\text{-}complex\text{-}banach\text{-}algebra}(\operatorname{uba\text{-}as\text{-}cba}(a), \operatorname{uba\text{-}as\text{-}cba}(b), f) \end{array}}}
\vnbprop{uba-as-module-functorial}{\vnbstmt{\begin{array}{@{}l@{}} \forall a, b, f. \\ \quad \operatorname{is\text{-}hom\text{-}unital\text{-}banach\text{-}algebra}(a, b, f) \Rightarrow  \\ \quad \quad \operatorname{is\text{-}hom\text{-}module}(\operatorname{uba\text{-}as\text{-}module}(a), \operatorname{uba\text{-}as\text{-}module}(b), f) \end{array}}}
\vnbprop{uba-as-normed-ag-functorial}{\vnbstmt{\begin{array}{@{}l@{}} \forall a, b, f. \\ \quad \operatorname{is\text{-}hom\text{-}unital\text{-}banach\text{-}algebra}(a, b, f) \Rightarrow  \\ \quad \quad \operatorname{is\text{-}hom\text{-}normed\text{-}ag}(\operatorname{uba\text{-}as\text{-}normed\text{-}ag}(a), \operatorname{uba\text{-}as\text{-}normed\text{-}ag}(b), f) \end{array}}}
\vnbprop{uba-as-ring-functorial}{\vnbstmt{\begin{array}{@{}l@{}} \forall a, b, f. \\ \quad \operatorname{is\text{-}hom\text{-}unital\text{-}banach\text{-}algebra}(a, b, f) \Rightarrow  \\ \quad \quad \operatorname{is\text{-}hom\text{-}ring}(\operatorname{uba\text{-}as\text{-}ring}(a), \operatorname{uba\text{-}as\text{-}ring}(b), f) \end{array}}}
\vnbprop{uiba-as-nuiba-functorial}{\vnbstmt{\begin{array}{@{}l@{}} \forall a, b, f. \\ \quad \operatorname{is\text{-}hom\text{-}unital\text{-}involutive\text{-}banach\text{-}algebra}(a, b, f) \Rightarrow  \\ \quad \quad \operatorname{is\text{-}hom\text{-}nu\text{-}involutive\text{-}banach\text{-}algebra}(\operatorname{uiba\text{-}as\text{-}nuiba}(a), \operatorname{uiba\text{-}as\text{-}nuiba}(b), f) \end{array}}}
\vnbprop{uiba-as-uba-functorial}{\vnbstmt{\begin{array}{@{}l@{}} \forall a, b, f. \\ \quad \operatorname{is\text{-}hom\text{-}unital\text{-}involutive\text{-}banach\text{-}algebra}(a, b, f) \Rightarrow  \\ \quad \quad \operatorname{is\text{-}hom\text{-}unital\text{-}banach\text{-}algebra}(\operatorname{uiba\text{-}as\text{-}uba}(a), \operatorname{uiba\text{-}as\text{-}uba}(b), f) \end{array}}}

\vnbbookfile{theorem-library/family-nn-enum}{theorem-library/family-nn-enum}
\noindent{\small family-nn-enum.scm -- agent BA-1117-TRACE (round 2), 2026-11-17. The support of an absolutely summable family is enumerated by a subset of NN, and the family transported along the enumeration (zero off it) is an NN-indexed family with the same unordered sum and a convergent series with that sum (finding F2 of ...}\par
\vnbprop{esum-inj-le}{\vnbstmt{\begin{array}{@{}l@{}} \forall j, x, a, g, h. \\ \quad \forall \mathit{nej} \in j.\; (a(\mathit{nej}) \in x) \wedge \\ \quad \forall b, y \in j.\; ((a(b) = a(y)) \Rightarrow (b = y)) \wedge \\ \quad g \in (x \to \operatorname{rr\text{-}pos\text{-}star}) \wedge \\ \quad h \in (j \to \operatorname{rr\text{-}pos\text{-}star}) \wedge \\ \quad \forall b \in j.\; (h(b) = g(a(b))) \Rightarrow \\ \quad \quad \operatorname{esum}(h) \leq \operatorname{esum}(g) \end{array}}}
\vnbprop{esum-superlevel-finite}{\vnbstmt{\begin{array}{@{}l@{}} \forall x.; \forall g \in (x \to \operatorname{rr\text{-}pos\text{-}star}), d. \\ \quad ((\operatorname{esum}(g) \in \mathbb{R}) \wedge \operatorname{pos\text{-}rr}(d)) \Rightarrow (|\{\mathit{nex} \in x : (d \leq g(\mathit{nex}))\}| \in \mathbb{N}) \end{array}}}
\vnbprop{esum-support-enum}{\vnbstmt{\begin{array}{@{}l@{}} \forall x.; \forall g \in (x \to \mathbb{R}). \\ \quad ((g \in (x \to \operatorname{rr\text{-}pos\text{-}star})) \wedge (\operatorname{esum}(g) \in \mathbb{R})) \Rightarrow  \\ \quad \quad \exists m \in \mathbf{Set}, h. \\ \quad \quad \quad \forall n \in m.\; (n \in \mathbb{N}) \wedge (h \in \mathrm{Bij}(m, \{\mathit{nex} \in x : (0 < g(\mathit{nex}))\})) \end{array}}}
\vnbprop{finsum-ag-pad-ptwise}{\vnbstmt{\begin{array}{@{}l@{}} \forall g.; \forall t, r \in \mathbf{Set}, h. \\ \quad \operatorname{is\text{-}abelian\text{-}group}(g) \wedge \\ \quad |t| \in \mathbb{N} \wedge \\ \quad \forall z \in r.\; (z \in t) \wedge \\ \quad \forall z \in t.\; (h(z) \in \operatorname{carr}(g)) \wedge \\ \quad \forall z \in t.\; (\neg (z \in r) \Rightarrow (h(z) = \operatorname{iden}(g))) \Rightarrow \\ \quad \quad \operatorname{finsum}(g, h, t) = \operatorname{finsum}(g, h, r) \end{array}}}
\vnbprop{card-of-inj-into-finite}{\vnbstmt{\begin{array}{@{}l@{}} \forall x, y, h. \\ \quad x \in \mathbf{Set} \wedge \\ \quad y \in \mathbf{Set} \wedge \\ \quad |y| \in \mathbb{N} \wedge \\ \quad \forall j \in x.\; (h(j) \in y) \wedge \\ \quad \forall a, b \in x.\; ((h(a) = h(b)) \Rightarrow (a = b)) \Rightarrow \\ \quad \quad |x| \in \mathbb{N} \end{array}}}
\vnbprop{finsum-ag-image-inj}{\vnbstmt{\begin{array}{@{}l@{}} \forall g.; \forall t \in \mathbf{Set}, i, f, k. \\ \quad \operatorname{is\text{-}abelian\text{-}group}(g) \wedge \\ \quad |t| \in \mathbb{N} \wedge \\ \quad \forall x \in t.\; (i(x) \in \mathbf{Set}) \wedge \\ \quad \forall x, y \in t.\; ((i(x) = i(y)) \Rightarrow (x = y)) \wedge \\ \quad \forall z \in \operatorname{image}(i, t).\; (f(z) \in \operatorname{carr}(g)) \wedge \\ \quad \forall x \in t.\; (k(x) = f(i(x))) \Rightarrow \\ \quad \quad \operatorname{finsum}(g, k, t) = \operatorname{finsum}(g, f, \operatorname{image}(i, t)) \end{array}}}
\vnbprop{sums-to-nn-transfer}{\vnbstmt{\begin{array}{@{}l@{}} \forall a, x, m, p, h, f, g, r. \\ \quad \operatorname{is\text{-}normed\text{-}ag}(a) \wedge \\ \quad m \in \mathbf{Set} \wedge \\ \quad \forall n \in m.\; (n \in \mathbb{N}) \wedge \\ \quad h \in \mathrm{Bij}(m, p) \wedge \\ \quad \forall \mathit{nep} \in p.\; (\mathit{nep} \in x) \wedge \\ \quad f \in (x \to \operatorname{carr}(\operatorname{normed\text{-}ag\text{-}as\text{-}abelian\text{-}group}(a))) \wedge \\ \quad \forall \mathit{nep} \in x. \\ \quad \quad \neg (\mathit{nep} \in p) \Rightarrow (f(\mathit{nep}) = \operatorname{iden}(\operatorname{normed\text{-}ag\text{-}as\text{-}abelian\text{-}group}(a))) \wedge \\ \quad g \in (\mathbb{N} \to \operatorname{carr}(\operatorname{normed\text{-}ag\text{-}as\text{-}abelian\text{-}group}(a))) \wedge \\ \quad \forall n \in m.\; (g(n) = f(h(n))) \wedge \\ \quad \forall n \in \mathbb{N}. \\ \quad \quad \neg (n \in m) \Rightarrow (g(n) = \operatorname{iden}(\operatorname{normed\text{-}ag\text{-}as\text{-}abelian\text{-}group}(a))) \wedge \\ \quad \operatorname{sums\text{-}to}(a, f, r) \Rightarrow \\ \quad \quad \operatorname{sums\text{-}to}(a, g, r) \end{array}}}
\vnbprop{cc-nag-dist}{\vnbstmt{\begin{array}{@{}l@{}} \forall a, b \in \mathbb{C}. \\ \quad \operatorname{nrm}(\operatorname{cc\text{-}nag})(\operatorname{opr}(\operatorname{cc\text{-}nag})(a, \operatorname{inv}(\operatorname{cc\text{-}nag})(b))) = \operatorname{magnitude}((a - b)) \end{array}}}
\vnbprop{finsum-ord-segment-partial-sum}{\vnbstmt{\begin{array}{@{}l@{}} \forall k \in \mathbb{N}, g \in (\mathbb{N} \to \mathbb{C}). \\ \quad \operatorname{finsum}(\operatorname{normed\text{-}ag\text{-}as\text{-}abelian\text{-}group}(\operatorname{cc\text{-}nag}), g, \mathrm{OS}(k)) = \operatorname{cc\text{-}series\text{-}partial\text{-}sum}(g, k) \end{array}}}
\vnbprop{nn-sums-to-series}{\vnbstmt{\begin{array}{@{}l@{}} \forall g, r. \\ \quad ((g \in (\mathbb{N} \to \mathbb{C})) \wedge ((r \in \mathbb{C}) \wedge \operatorname{sums\text{-}to}(\operatorname{cc\text{-}nag}, g, r))) \Rightarrow  \\ \quad \quad \operatorname{cc\text{-}series\text{-}converges\text{-}to}(g, r) \end{array}}}
\vnbprop{esum-nn-transfer-le}{\vnbstmt{\begin{array}{@{}l@{}} \forall x, m, a, g, h. \\ \quad m \in \mathbf{Set} \wedge \\ \quad \forall n \in m.\; (n \in \mathbb{N}) \wedge \\ \quad \forall j \in m.\; (a(j) \in x) \wedge \\ \quad \forall b, y \in m.\; ((a(b) = a(y)) \Rightarrow (b = y)) \wedge \\ \quad g \in (x \to \operatorname{rr\text{-}pos\text{-}star}) \wedge \\ \quad h \in (\mathbb{N} \to \operatorname{rr\text{-}pos\text{-}star}) \wedge \\ \quad \forall n \in m.\; (h(n) = g(a(n))) \wedge \\ \quad \forall n \in \mathbb{N}.\; (\neg (n \in m) \Rightarrow (h(n) = 0)) \Rightarrow \\ \quad \quad \operatorname{esum}(h) \leq \operatorname{esum}(g) \end{array}}}
\vnbprop{cc-family-nn-series}{\vnbstmt{\begin{array}{@{}l@{}} \forall x, a, m, r. \\ \quad a \in (x \to \mathbb{C}) \wedge \\ \quad m \in (x \to \mathbb{R}) \wedge \\ \quad m \in (x \to \operatorname{rr\text{-}pos\text{-}star}) \wedge \\ \quad \forall \mathit{nex} \in x.\; (m(\mathit{nex}) = \operatorname{magnitude}(a(\mathit{nex}))) \wedge \\ \quad \operatorname{esum}(m) \in \mathbb{R} \wedge \\ \quad r \in \mathbb{C} \wedge \\ \quad \operatorname{sums\text{-}to}(\operatorname{cc\text{-}nag}, a, r) \Rightarrow \\ \quad \quad \exists \mathit{nem} \in \mathbf{Set}, h. \\ \quad \quad \quad \forall n \in \mathit{nem}.\; (n \in \mathbb{N}) \wedge \\ \quad \quad \quad h \in \mathrm{Bij}(\mathit{nem}, \{\mathit{nex} \in x : (0 < m(\mathit{nex}))\}) \wedge \\ \quad \quad \quad \operatorname{cc\text{-}series\text{-}converges\text{-}to}((\lambda n \in \mathbb{N}.\; \left(\text{if }(n \in \mathit{nem})\text{ then }a(h(n))\text{ else }0\right)), r) \wedge \\ \quad \quad \quad \operatorname{esum}((\lambda n \in \mathbb{N}.\; \left(\text{if }(n \in \mathit{nem})\text{ then }m(h(n))\text{ else }0\right))) \leq \operatorname{esum}(m) \end{array}}}

\vnbbookfile{theorem-library/sums-to-algebra}{theorem-library/sums-to-algebra}
\noindent{\small sums-to-algebra.scm -- wave of 2026-11-09/10 (agent BA-ONB3), wired 2026-11-10 from scratchpad/baonb3-1109/b3-sums.scm. Library: operator-algebras.}\par
\vnbprop{sums-to-add}{\vnbstmt{\begin{array}{@{}l@{}} \forall p.; \forall f, b \in (\operatorname{dom}(f) \to \operatorname{carr}(p)), r, s \in \operatorname{carr}(p). \\ \quad (\operatorname{is\text{-}normed\text{-}ag}(p) \wedge (\operatorname{sums\text{-}to}(p, f, r) \wedge \operatorname{sums\text{-}to}(p, b, s))) \Rightarrow  \\ \quad \quad \operatorname{sums\text{-}to}(p, (\lambda z \in \operatorname{dom}(f).\; \operatorname{opr}(p)(f(z), b(z))), \operatorname{opr}(p)(r, s)) \end{array}}}
\vnbprop{cc-sums-to-bounded-additive}{\vnbstmt{\begin{array}{@{}l@{}} \forall h \in (\mathbb{C} \to \mathbb{C}), c \in \mathbb{R}, f \in (\operatorname{dom}(f) \to \mathbb{C}), r \in \mathbb{C}. \\ \quad \forall a, b \in \mathbb{C}.\; (h((a + b)) = (h(a) + h(b))) \wedge \\ \quad \forall a \in \mathbb{C}.\; (h(-a) = -h(a)) \wedge \\ \quad h(0) = 0 \wedge \\ \quad 0 \leq c \wedge \\ \quad \forall a \in \mathbb{C}.\; (\operatorname{magnitude}(h(a)) \leq (c \cdot \operatorname{magnitude}(a))) \wedge \\ \quad \operatorname{sums\text{-}to}(\operatorname{cc\text{-}nag}, f, r) \Rightarrow \\ \quad \quad \operatorname{sums\text{-}to}(\operatorname{cc\text{-}nag}, (\lambda z \in \operatorname{dom}(f).\; h(f(z))), h(r)) \end{array}}}
\vnbprop{cc-sums-to-scale}{\vnbstmt{\begin{array}{@{}l@{}} \forall c \in \mathbb{C}, f \in (\operatorname{dom}(f) \to \mathbb{C}), r \in \mathbb{C}. \\ \quad \operatorname{sums\text{-}to}(\operatorname{cc\text{-}nag}, f, r) \Rightarrow  \\ \quad \quad \operatorname{sums\text{-}to}(\operatorname{cc\text{-}nag}, (\lambda z \in \operatorname{dom}(f).\; (c \cdot f(z))), (c \cdot r)) \end{array}}}
\vnbprop{cc-sums-to-conjugate}{\vnbstmt{\begin{array}{@{}l@{}} \forall f \in (\operatorname{dom}(f) \to \mathbb{C}), r \in \mathbb{C}. \\ \quad \operatorname{sums\text{-}to}(\operatorname{cc\text{-}nag}, f, r) \Rightarrow  \\ \quad \quad \operatorname{sums\text{-}to}(\operatorname{cc\text{-}nag}, (\lambda z \in \operatorname{dom}(f).\; \operatorname{conjugate}(f(z))), \operatorname{conjugate}(r)) \end{array}}}
\vnbprop{rps-finsum-eq-cc-finsum}{\vnbstmt{\begin{array}{@{}l@{}} \forall g \in (\operatorname{dom}(g) \to \mathbb{R}), w \in \mathbf{Set}. \\ \quad (\forall z \in \operatorname{dom}(g).\; (0 \leq g(z)) \wedge ((|w| \in \mathbb{N}) \wedge (w \subseteq \operatorname{dom}(g)))) \Rightarrow  \\ \quad \quad \operatorname{finsum}(\operatorname{rr\text{-}pos\text{-}star\text{-}add\text{-}monoid}, g, w) \in \mathbb{R} \wedge \\ \quad \quad 0 \leq \operatorname{finsum}(\operatorname{rr\text{-}pos\text{-}star\text{-}add\text{-}monoid}, g, w) \wedge \\ \quad \quad \operatorname{finsum}(\operatorname{rr\text{-}pos\text{-}star\text{-}add\text{-}monoid}, g, w) = \operatorname{finsum}(\operatorname{normed\text{-}ag\text{-}as\text{-}abelian\text{-}group}(\operatorname{cc\text{-}nag}), g, w) \end{array}}}
\vnbprop{esum-sums-to-cc}{\vnbstmt{\begin{array}{@{}l@{}} \forall g \in (\operatorname{dom}(g) \to \mathbb{R}). \\ \quad (\forall z \in \operatorname{dom}(g).\; (0 \leq g(z)) \wedge (\operatorname{esum}(g) \in \mathbb{R})) \Rightarrow  \\ \quad \quad \operatorname{sums\text{-}to}(\operatorname{cc\text{-}nag}, g, \operatorname{esum}(g)) \end{array}}}

\vnbbookfile{theorem-library/l2-families-2}{theorem-library/l2-families-2}
\noindent{\small l2-families-2.scm -- wave of 2026-11-09/10 (agent BA-ONB3), wired 2026-11-10 from scratchpad/baonb3-1109/b3-l2.scm. Library: operator-algebras.}\par
\vnbprop{esum-real-le-ptwise}{\vnbstmt{\begin{array}{@{}l@{}} \forall d, f, g. \\ \quad f \in (d \to \mathbb{R}) \wedge \\ \quad g \in (d \to \mathbb{R}) \wedge \\ \quad \forall z \in d.\; (0 \leq f(z)) \wedge \\ \quad \forall z \in d.\; (f(z) \leq g(z)) \wedge \\ \quad \operatorname{esum}(g) \in \mathbb{R} \Rightarrow \\ \quad \quad (\operatorname{esum}(f) \in \mathbb{R}) \wedge (\operatorname{esum}(f) \leq \operatorname{esum}(g)) \end{array}}}
\vnbprop{esum-real-add-le}{\vnbstmt{\begin{array}{@{}l@{}} \forall d, f, g, h. \\ \quad f \in (d \to \mathbb{R}) \wedge \\ \quad g \in (d \to \mathbb{R}) \wedge \\ \quad h \in (d \to \mathbb{R}) \wedge \\ \quad \forall z \in d.\; (0 \leq f(z)) \wedge \\ \quad \forall z \in d.\; (0 \leq g(z)) \wedge \\ \quad \forall z \in d.\; (0 \leq h(z)) \wedge \\ \quad \forall z \in d.\; (h(z) \leq (f(z) + g(z))) \wedge \\ \quad \operatorname{esum}(f) \in \mathbb{R} \wedge \\ \quad \operatorname{esum}(g) \in \mathbb{R} \Rightarrow \\ \quad \quad (\operatorname{esum}(h) \in \mathbb{R}) \wedge (\operatorname{esum}(h) \leq (\operatorname{esum}(f) + \operatorname{esum}(g))) \end{array}}}
\vnbprop{l2-families-unfold}{\vnbstmt{\begin{array}{@{}l@{}} \forall x. \\ \quad (\operatorname{l2\text{-}families}(x) \simeq \{c \in (x \to \mathbb{C}) : (\operatorname{esum}((\lambda k \in x.\; (\operatorname{magnitude}(c(k)) \cdot \operatorname{magnitude}(c(k))))) \in \mathbb{R})\}) \end{array}}}
\vnbprop{l2-families-elim}{\vnbstmt{\begin{array}{@{}l@{}} \forall x.; \forall c \in \operatorname{l2\text{-}families}(x). \\ \quad ((c \in (x \to \mathbb{C})) \wedge (\operatorname{esum}((\lambda k \in x.\; (\operatorname{magnitude}(c(k)) \cdot \operatorname{magnitude}(c(k))))) \in \mathbb{R})) \end{array}}}
\vnbprop{l2-families-intro}{\vnbstmt{\begin{array}{@{}l@{}} \forall x.; \forall c \in (x \to \mathbb{C}). \\ \quad (\operatorname{esum}((\lambda k \in x.\; (\operatorname{magnitude}(c(k)) \cdot \operatorname{magnitude}(c(k))))) \in \mathbb{R}) \Rightarrow  \\ \quad \quad c \in \operatorname{l2\text{-}families}(x) \end{array}}}
\vnbprop{l2-families-add-closed}{\vnbstmt{\begin{array}{@{}l@{}} \forall x, a, b. \\ \quad ((a \in \operatorname{l2\text{-}families}(x)) \wedge (b \in \operatorname{l2\text{-}families}(x))) \Rightarrow  \\ \quad \quad (\lambda \mathit{b3a} \in x.\; (a(\mathit{b3a}) + b(\mathit{b3a}))) \in \operatorname{l2\text{-}families}(x) \end{array}}}
\vnbprop{l2-families-act-closed}{\vnbstmt{\begin{array}{@{}l@{}} \forall x, c, a. \\ \quad ((c \in \mathbb{C}) \wedge (a \in \operatorname{l2\text{-}families}(x))) \Rightarrow  \\ \quad \quad (\lambda \mathit{b3a} \in x.\; (c \cdot a(\mathit{b3a}))) \in \operatorname{l2\text{-}families}(x) \end{array}}}
\vnbprop{l2-ip-family-sums}{\vnbstmt{\begin{array}{@{}l@{}} \forall x, a, b. \\ \quad ((a \in \operatorname{l2\text{-}families}(x)) \wedge (b \in \operatorname{l2\text{-}families}(x))) \Rightarrow  \\ \quad \quad \operatorname{cc\text{-}family\text{-}sum}((\lambda k \in x.\; (a(k) \cdot \operatorname{conjugate}(b(k))))) \in \mathbb{C} \wedge \\ \quad \quad \operatorname{sums\text{-}to}(\operatorname{cc\text{-}nag}, (\lambda k \in x.\; (a(k) \cdot \operatorname{conjugate}(b(k)))), \\ \quad \quad \quad \operatorname{cc\text{-}family\text{-}sum}((\lambda k \in x.\; (a(k) \cdot \operatorname{conjugate}(b(k)))))) \end{array}}}
\vnbprop{l2-ip-apply}{\vnbstmt{\begin{array}{@{}l@{}} \forall x, a, b. \\ \quad ((a \in \operatorname{l2\text{-}families}(x)) \wedge (b \in \operatorname{l2\text{-}families}(x))) \Rightarrow  \\ \quad \quad \operatorname{l2\text{-}ip}(x)(a, b) = \operatorname{cc\text{-}family\text{-}sum}((\lambda k \in x.\; (a(k) \cdot \operatorname{conjugate}(b(k))))) \end{array}}}
\vnbprop{l2-ip-add-left}{\vnbstmt{\begin{array}{@{}l@{}} \forall x, a, b, c. \\ \quad ((a \in \operatorname{l2\text{-}families}(x)) \wedge ((b \in \operatorname{l2\text{-}families}(x)) \wedge (c \in \operatorname{l2\text{-}families}(x)))) \Rightarrow  \\ \quad \quad \operatorname{l2\text{-}ip}(x)(\operatorname{l2f\text{-}vadd}(x)(a, b), c) = (\operatorname{l2\text{-}ip}(x)(a, c) + \operatorname{l2\text{-}ip}(x)(b, c)) \end{array}}}
\vnbprop{l2-ip-act-left}{\vnbstmt{\begin{array}{@{}l@{}} \forall x, c, a, b. \\ \quad ((c \in \mathbb{C}) \wedge ((a \in \operatorname{l2\text{-}families}(x)) \wedge (b \in \operatorname{l2\text{-}families}(x)))) \Rightarrow  \\ \quad \quad \operatorname{l2\text{-}ip}(x)(\operatorname{l2f\text{-}act}(x)(c, a), b) = (c \cdot \operatorname{l2\text{-}ip}(x)(a, b)) \end{array}}}
\vnbprop{l2-ip-conj-sym}{\vnbstmt{\begin{array}{@{}l@{}} \forall x, a, b. \\ \quad ((a \in \operatorname{l2\text{-}families}(x)) \wedge (b \in \operatorname{l2\text{-}families}(x))) \Rightarrow  \\ \quad \quad \operatorname{l2\text{-}ip}(x)(b, a) = \operatorname{conjugate}(\operatorname{l2\text{-}ip}(x)(a, b)) \end{array}}}
\vnbprop{l2-ip-self}{\vnbstmt{\begin{array}{@{}l@{}} \forall x.; \forall a \in \operatorname{l2\text{-}families}(x). \\ \quad \operatorname{l2\text{-}ip}(x)(a, a) = \operatorname{esum}((\lambda k \in x.\; (\operatorname{magnitude}(a(k)) \cdot \operatorname{magnitude}(a(k))))) \wedge \\ \quad \operatorname{l2\text{-}ip}(x)(a, a) \in \mathbb{R} \wedge \\ \quad 0 \leq \operatorname{l2\text{-}ip}(x)(a, a) \end{array}}}
\vnbprop{l2-ip-self-zero}{\vnbstmt{\forall x.; \forall a \in \operatorname{l2\text{-}families}(x).\; ((\operatorname{l2\text{-}ip}(x)(a, a) = 0) \Leftrightarrow (a = \operatorname{l2f\text{-}vzero}(x)))}}

\vnbbookfile{theorem-library/riesz-fischer}{theorem-library/riesz-fischer}
\noindent{\small riesz-fischer.scm -- wave of 2026-11-09/10 (agent BA-ONB4), wired 2026-11-10 from scratchpad/baonb4-1109/rf-1.scm, scratchpad/baonb4-1109/rf-2.scm, scratchpad/baonb4-1109/rf-3.scm. Library: operator-algebras.}\par
\vnbprop{rf-partial-coef}{\vnbstmt{\begin{array}{@{}l@{}} \forall v, x, e.; \forall c \in (x \to \mathbb{C}), w \in \mathbf{Set}. \\ \quad \operatorname{is\text{-}complex\text{-}inner\text{-}product\text{-}space}(v) \wedge \\ \quad \operatorname{is\text{-}orthonormal\text{-}family}(v, x, e) \wedge \\ \quad |w| \in \mathbb{N} \wedge \\ \quad w \subseteq x \Rightarrow \\ \quad \quad \operatorname{finsum}(\operatorname{normed\text{-}ag\text{-}as\text{-}abelian\text{-}group}(\operatorname{ip\text{-}normed\text{-}ag}(v)), (\lambda k \in x.\; \operatorname{act}(v)(c(k), e(k))), w) \\ \quad \quad \in\; \operatorname{vec}(v) \wedge \\ \quad \quad \forall q \in w. \\ \quad \quad \quad \langle\cdot,\cdot\rangle_{v}(\operatorname{finsum}(\operatorname{normed\text{-}ag\text{-}as\text{-}abelian\text{-}group}(\operatorname{ip\text{-}normed\text{-}ag}(v)), (\lambda k \in x.\; \operatorname{act}(v)(c(k), e(k))), w), \\ \quad \quad \quad \quad e(q)) \\ \quad \quad \quad =\; c(q) \end{array}}}
\vnbprop{rf-partial-diff}{\vnbstmt{\begin{array}{@{}l@{}} \forall v, x, e.; \forall c \in (x \to \mathbb{C}), a, b \in \mathbf{Set}. \\ \quad \operatorname{is\text{-}complex\text{-}inner\text{-}product\text{-}space}(v) \wedge \\ \quad \operatorname{is\text{-}orthonormal\text{-}family}(v, x, e) \wedge \\ \quad |a| \in \mathbb{N} \wedge \\ \quad a \subseteq x \wedge \\ \quad |b| \in \mathbb{N} \wedge \\ \quad b \subseteq a \Rightarrow \\ \quad \quad \operatorname{finsum}(\operatorname{rr\text{-}pos\text{-}star\text{-}add\text{-}monoid}, (\lambda k \in x.\; (\operatorname{magnitude}(c(k)) \cdot \operatorname{magnitude}(c(k)))), a) \in \mathbb{R} \wedge \\ \quad \quad \operatorname{finsum}(\operatorname{rr\text{-}pos\text{-}star\text{-}add\text{-}monoid}, (\lambda k \in x.\; (\operatorname{magnitude}(c(k)) \cdot \operatorname{magnitude}(c(k)))), b) \in \mathbb{R} \wedge \\ \quad \quad \langle\cdot,\cdot\rangle_{v}(\operatorname{vadd}(v)(\operatorname{finsum}(\operatorname{normed\text{-}ag\text{-}as\text{-}abelian\text{-}group}(\operatorname{ip\text{-}normed\text{-}ag}(v)), (\lambda \mathit{rfk} \in x.\; \operatorname{act}(v)(c(\mathit{rfk}), e(\mathit{rfk}))), a), \\ \quad \quad \quad \quad \operatorname{vneg}(v)(\operatorname{finsum}(\operatorname{normed\text{-}ag\text{-}as\text{-}abelian\text{-}group}(\operatorname{ip\text{-}normed\text{-}ag}(v)), (\lambda \mathit{rfk} \in x.\; \operatorname{act}(v)(c(\mathit{rfk}), e(\mathit{rfk}))), b))), \\ \quad \quad \quad \operatorname{vadd}(v)(\operatorname{finsum}(\operatorname{normed\text{-}ag\text{-}as\text{-}abelian\text{-}group}(\operatorname{ip\text{-}normed\text{-}ag}(v)), (\lambda \mathit{rfk} \in x.\; \operatorname{act}(v)(c(\mathit{rfk}), e(\mathit{rfk}))), a), \\ \quad \quad \quad \quad \operatorname{vneg}(v)(\operatorname{finsum}(\operatorname{normed\text{-}ag\text{-}as\text{-}abelian\text{-}group}(\operatorname{ip\text{-}normed\text{-}ag}(v)), (\lambda \mathit{rfk} \in x.\; \operatorname{act}(v)(c(\mathit{rfk}), e(\mathit{rfk}))), b)))) \\ \quad \quad =\; \operatorname{finsum}(\operatorname{rr\text{-}pos\text{-}star\text{-}add\text{-}monoid}, (\lambda k \in x.\; (\operatorname{magnitude}(c(k)) \cdot \operatorname{magnitude}(c(k)))), a) + -\operatorname{finsum}(\operatorname{rr\text{-}pos\text{-}star\text{-}add\text{-}monoid}, (\lambda k \in x.\; (\operatorname{magnitude}(c(k)) \cdot \operatorname{magnitude}(c(k)))), b) \end{array}}}
\vnbprop{rf-cauchy-net}{\vnbstmt{\begin{array}{@{}l@{}} \forall v, x, e.; \forall c \in (x \to \mathbb{C}), s. \\ \quad \operatorname{is\text{-}complex\text{-}inner\text{-}product\text{-}space}(v) \wedge \\ \quad \operatorname{is\text{-}orthonormal\text{-}family}(v, x, e) \wedge \\ \quad \operatorname{esum}((\lambda k \in x.\; (\operatorname{magnitude}(c(k)) \cdot \operatorname{magnitude}(c(k))))) \in \mathbb{R} \wedge \\ \quad \operatorname{pos\text{-}rr}(s) \Rightarrow \\ \quad \quad \exists n \in \mathbf{Set}. \\ \quad \quad \quad |n| \in \mathbb{N} \wedge \\ \quad \quad \quad n \subseteq \operatorname{dom}((\lambda \mathit{rfk} \in x.\; \operatorname{act}(v)(c(\mathit{rfk}), e(\mathit{rfk})))) \wedge \\ \quad \quad \quad \forall a, b \in \mathbf{Set}. \\ \quad \quad \quad \quad |a| \in \mathbb{N} \wedge \\ \quad \quad \quad \quad n \subseteq a \wedge \\ \quad \quad \quad \quad a \subseteq \operatorname{dom}((\lambda \mathit{rfk} \in x.\; \operatorname{act}(v)(c(\mathit{rfk}), e(\mathit{rfk})))) \wedge \\ \quad \quad \quad \quad |b| \in \mathbb{N} \wedge \\ \quad \quad \quad \quad n \subseteq b \wedge \\ \quad \quad \quad \quad b \subseteq \operatorname{dom}((\lambda \mathit{rfk} \in x.\; \operatorname{act}(v)(c(\mathit{rfk}), e(\mathit{rfk})))) \Rightarrow \\ \quad \quad \quad \quad \quad \operatorname{dist}(\operatorname{nag\text{-}metric\text{-}space}(\operatorname{ip\text{-}normed\text{-}ag}(v)))(\operatorname{finsum}(\operatorname{normed\text{-}ag\text{-}as\text{-}abelian\text{-}group}(\operatorname{ip\text{-}normed\text{-}ag}(v)), (\lambda \mathit{rfk} \in x.\; \operatorname{act}(v)(c(\mathit{rfk}), e(\mathit{rfk}))), a), \\ \quad \quad \quad \quad \quad \quad \operatorname{finsum}(\operatorname{normed\text{-}ag\text{-}as\text{-}abelian\text{-}group}(\operatorname{ip\text{-}normed\text{-}ag}(v)), (\lambda \mathit{rfk} \in x.\; \operatorname{act}(v)(c(\mathit{rfk}), e(\mathit{rfk}))), b)) \\ \quad \quad \quad \quad \quad <\; s \end{array}}}
\vnbprop{rf-coef-limit}{\vnbstmt{\begin{array}{@{}l@{}} \forall v, x, e.; \forall c \in (x \to \mathbb{C}), l \in \operatorname{vec}(v), p \in x. \\ \quad \operatorname{is\text{-}complex\text{-}inner\text{-}product\text{-}space}(v) \wedge \\ \quad \operatorname{is\text{-}orthonormal\text{-}family}(v, x, e) \wedge \\ \quad \operatorname{sums\text{-}to}(\operatorname{ip\text{-}normed\text{-}ag}(v), (\lambda k \in x.\; \operatorname{act}(v)(c(k), e(k))), l) \Rightarrow \\ \quad \quad \langle\cdot,\cdot\rangle_{v}(l, e(p)) = c(p) \end{array}}}
\vnbprop{riesz-fischer-sum}{\vnbstmt{\begin{array}{@{}l@{}} \forall v, a, e.; \forall c \in \operatorname{l2\text{-}families}(a). \\ \quad (\operatorname{is\text{-}hilbert\text{-}space}(v) \wedge \operatorname{is\text{-}orthonormal\text{-}family}(v, a, e)) \Rightarrow  \\ \quad \quad \exists x \in \operatorname{vec}(v). \\ \quad \quad \quad \operatorname{sums\text{-}to}(\operatorname{ip\text{-}normed\text{-}ag}(v), (\lambda k \in a.\; \operatorname{act}(v)(c(k), e(k))), x) \wedge \\ \quad \quad \quad \forall p \in a.\; (\langle\cdot,\cdot\rangle_{v}(x, e(p)) = c(p)) \end{array}}}
\vnbprop{riesz-fischer}{\vnbstmt{\begin{array}{@{}l@{}} \forall v, a, e.; \forall c \in \operatorname{l2\text{-}families}(a). \\ \quad (\operatorname{is\text{-}hilbert\text{-}space}(v) \wedge \operatorname{is\text{-}orthonormal\text{-}basis}(v, a, e)) \Rightarrow  \\ \quad \quad \exists x \in \operatorname{vec}(v).\; ((\lambda k \in a.\; \langle\cdot,\cdot\rangle_{v}(x, e(k))) = c) \end{array}}}

\vnbbookfile{theorem-library/l2-families-cips}{theorem-library/l2-families-cips}
\noindent{\small l2-families-cips.scm -- wave of 2026-11-09/10 (agent BA-ONB4), wired 2026-11-10 from scratchpad/baonb4-1109/l2p-1.scm, scratchpad/baonb4-1109/l2p-2.scm. Library: operator-algebras.}\par
\vnbprop{l2-vadd-apply}{\vnbstmt{\begin{array}{@{}l@{}} \forall x, a, b.; \forall k \in x. \\ \quad ((a \in \operatorname{l2\text{-}families}(x)) \wedge (b \in \operatorname{l2\text{-}families}(x))) \Rightarrow  \\ \quad \quad \operatorname{l2f\text{-}vadd}(x)(a, b)(k) = (a(k) + b(k)) \end{array}}}
\vnbprop{l2-act-apply}{\vnbstmt{\begin{array}{@{}l@{}} \forall x, z, a.; \forall k \in x. \\ \quad ((z \in \mathbb{C}) \wedge (a \in \operatorname{l2\text{-}families}(x))) \Rightarrow (\operatorname{l2f\text{-}act}(x)(z, a)(k) = (z \cdot a(k))) \end{array}}}
\vnbprop{l2-vneg-apply}{\vnbstmt{\forall x.; \forall a \in \operatorname{l2\text{-}families}(x), k \in x.\; (\operatorname{l2f\text{-}vneg}(x)(a)(k) = -a(k))}}
\vnbprop{l2-vzero-apply}{\vnbstmt{\forall x.; \forall k \in x.\; (\operatorname{l2f\text{-}vzero}(x)(k) = 0)}}
\vnbprop{l2-vadd-type}{\vnbstmt{\begin{array}{@{}l@{}} \forall x, a, b. \\ \quad ((a \in \operatorname{l2\text{-}families}(x)) \wedge (b \in \operatorname{l2\text{-}families}(x))) \Rightarrow  \\ \quad \quad \operatorname{l2f\text{-}vadd}(x)(a, b) \in \operatorname{l2\text{-}families}(x) \end{array}}}
\vnbprop{l2-act-type}{\vnbstmt{\begin{array}{@{}l@{}} \forall x, z, a. \\ \quad ((z \in \mathbb{C}) \wedge (a \in \operatorname{l2\text{-}families}(x))) \Rightarrow (\operatorname{l2f\text{-}act}(x)(z, a) \in \operatorname{l2\text{-}families}(x)) \end{array}}}
\vnbprop{l2-families-neg-closed}{\vnbstmt{\forall x.; \forall a \in \operatorname{l2\text{-}families}(x).\; ((\lambda k \in x.\; -a(k)) \in \operatorname{l2\text{-}families}(x))}}
\vnbprop{l2-vneg-type}{\vnbstmt{\forall x.; \forall a \in \operatorname{l2\text{-}families}(x).\; (\operatorname{l2f\text{-}vneg}(x)(a) \in \operatorname{l2\text{-}families}(x))}}
\vnbprop{l2-vzero-in}{\vnbstmt{\forall x \in \mathbf{Set}.\; (\operatorname{l2f\text{-}vzero}(x) \in \operatorname{l2\text{-}families}(x))}}
\vnbprop{l2-families-ext}{\vnbstmt{\begin{array}{@{}l@{}} \forall x, a, b. \\ \quad ((a \in \operatorname{l2\text{-}families}(x)) \wedge ((b \in \operatorname{l2\text{-}families}(x)) \wedge \forall e \in x.\; (a(e) = b(e)))) \Rightarrow  \\ \quad \quad a = b \end{array}}}
\vnbprop{l2-families-is-set}{\vnbstmt{\forall x \in \mathbf{Set}.\; (\operatorname{l2\text{-}families}(x) \in \mathbf{Set})}}
\vnbprop{l2-vadd-fun}{\vnbstmt{\begin{array}{@{}l@{}} \forall x \in \mathbf{Set}. \\ \quad \operatorname{l2f\text{-}vadd}(x) \in ((\operatorname{l2\text{-}families}(x) \times \operatorname{l2\text{-}families}(x)) \to \operatorname{l2\text{-}families}(x)) \end{array}}}
\vnbprop{l2-act-fun}{\vnbstmt{\begin{array}{@{}l@{}} \forall x \in \mathbf{Set}. \\ \quad \operatorname{l2f\text{-}act}(x) \in ((\mathbb{C} \times \operatorname{l2\text{-}families}(x)) \to \operatorname{l2\text{-}families}(x)) \end{array}}}
\vnbprop{l2-vneg-fun}{\vnbstmt{\forall x \in \mathbf{Set}.\; (\operatorname{l2f\text{-}vneg}(x) \in (\operatorname{l2\text{-}families}(x) \to \operatorname{l2\text{-}families}(x)))}}
\vnbprop{l2-ip-fun}{\vnbstmt{\begin{array}{@{}l@{}} \forall x \in \mathbf{Set}. \\ \quad \operatorname{l2\text{-}ip}(x) \in ((\operatorname{l2\text{-}families}(x) \times \operatorname{l2\text{-}families}(x)) \to \mathbb{C}) \end{array}}}
\vnbprop{l2-cips-scal}{\vnbstmt{\begin{array}{@{}l@{}} \forall x \in \mathbf{Set}. \\ \quad (\operatorname{scal}(\operatorname{l2\text{-}cips}(x)) \simeq \operatorname{normed\text{-}field\text{-}as\text{-}commutative\text{-}ring}(\operatorname{cc\text{-}normed\text{-}field})) \end{array}}}
\vnbprop{l2-cips-vec}{\vnbstmt{\forall x \in \mathbf{Set}.\; (\operatorname{vec}(\operatorname{l2\text{-}cips}(x)) \simeq \operatorname{l2\text{-}families}(x))}}
\vnbprop{l2-cips-vadd}{\vnbstmt{\forall x \in \mathbf{Set}.\; (\operatorname{vadd}(\operatorname{l2\text{-}cips}(x)) \simeq \operatorname{l2f\text{-}vadd}(x))}}
\vnbprop{l2-cips-vzero}{\vnbstmt{\forall x \in \mathbf{Set}.\; (\operatorname{vzero}(\operatorname{l2\text{-}cips}(x)) \simeq \operatorname{l2f\text{-}vzero}(x))}}
\vnbprop{l2-cips-vneg}{\vnbstmt{\forall x \in \mathbf{Set}.\; (\operatorname{vneg}(\operatorname{l2\text{-}cips}(x)) \simeq \operatorname{l2f\text{-}vneg}(x))}}
\vnbprop{l2-cips-act}{\vnbstmt{\forall x \in \mathbf{Set}.\; (\operatorname{act}(\operatorname{l2\text{-}cips}(x)) \simeq \operatorname{l2f\text{-}act}(x))}}
\vnbprop{l2-cips-ip}{\vnbstmt{\forall x \in \mathbf{Set}.\; (\langle\cdot,\cdot\rangle_{\operatorname{l2\text{-}cips}(x)} \simeq \operatorname{l2\text{-}ip}(x))}}
\vnbprop{l2-cips-is-cips}{\vnbstmt{\forall x \in \mathbf{Set}.\; \operatorname{is\text{-}complex\text{-}inner\text{-}product\text{-}space}(\operatorname{l2\text{-}cips}(x))}}

\vnbbookfile{theorem-library/hilbert-separable-l2}{theorem-library/hilbert-separable-l2}
\noindent{\small hilbert-separable-l2.scm -- wave of 2026-11-09/10 (agent BA-ONB5), wired 2026-11-10 from scratchpad/baonb5-1109/hilbert-separable-l2.scm. Library: operator-algebras.}\par
\vnbprop{b5-vsub-interchange}{\vnbstmt{\begin{array}{@{}l@{}} \forall v.; \forall a, b, c, d \in \operatorname{vec}(v). \\ \quad \operatorname{is\text{-}complex\text{-}inner\text{-}product\text{-}space}(v) \Rightarrow  \\ \quad \quad \operatorname{vadd}(v)(\operatorname{vadd}(v)(a, b), \operatorname{vneg}(v)(\operatorname{vadd}(v)(c, d))) \\ \quad \quad =\; \operatorname{vadd}(v)(\operatorname{vadd}(v)(a, \operatorname{vneg}(v)(c)), \operatorname{vadd}(v)(b, \operatorname{vneg}(v)(d))) \end{array}}}
\vnbprop{b5-act-sub-norm}{\vnbstmt{\begin{array}{@{}l@{}} \forall v.; \forall y \in \operatorname{vec}(v), a, b \in \mathbb{C}. \\ \quad (\operatorname{is\text{-}complex\text{-}inner\text{-}product\text{-}space}(v) \wedge (\langle\cdot,\cdot\rangle_{v}(y, y) = 1)) \Rightarrow  \\ \quad \quad \lVert \operatorname{vadd}(v)(\operatorname{act}(v)(a, y), \operatorname{vneg}(v)(\operatorname{act}(v)(b, y))) \rVert = \operatorname{magnitude}((a + -b)) \end{array}}}
\vnbprop{b5-finite-approx}{\vnbstmt{\begin{array}{@{}l@{}} \forall v, x.; \forall e \in (x \to \operatorname{vec}(v)), q \in (\mathbb{N} \to \mathbb{C}), u.; \forall f \in \mathbf{Set}, d. \\ \quad \operatorname{is\text{-}complex\text{-}inner\text{-}product\text{-}space}(v) \wedge \\ \quad \forall k \in x.\; (u(k) \in \operatorname{vec}(v)) \wedge \\ \quad \forall k \in x, \mathit{b5e}. \\ \quad \quad \operatorname{pos\text{-}rr}(\mathit{b5e}) \Rightarrow  \\ \quad \quad \quad \exists j \in \mathbb{N}. \\ \quad \quad \quad \quad \lVert \operatorname{vadd}(v)(u(k), \operatorname{vneg}(v)(\operatorname{act}(v)(q(j), e(k)))) \rVert < \mathit{b5e} \wedge \\ \quad |f| \in \mathbb{N} \wedge \\ \quad f \subseteq x \wedge \\ \quad \operatorname{pos\text{-}rr}(d) \Rightarrow \\ \quad \quad \exists s \in \operatorname{fin\text{-}subsets}((x \times \mathbb{N})). \\ \quad \quad \quad \forall p \in s.\; (\operatorname{nth}(1, p) \in f) \wedge \\ \quad \quad \quad \lVert \operatorname{vadd}(v)(\operatorname{finsum}(\operatorname{normed\text{-}ag\text{-}as\text{-}abelian\text{-}group}(\operatorname{ip\text{-}normed\text{-}ag}(v)), u, f), \operatorname{vneg}(v)(\operatorname{finsum}(\operatorname{normed\text{-}ag\text{-}as\text{-}abelian\text{-}group}(\operatorname{ip\text{-}normed\text{-}ag}(v)), (\lambda \mathit{b5q} \in (x \times \mathbb{N}).\; \operatorname{act}(v)(q(\operatorname{nth}(2, \mathit{b5q})), e(\operatorname{nth}(1, \mathit{b5q})))), s))) \rVert \\ \quad \quad \quad <\; d \end{array}}}
\vnbprop{onb-countable-separable}{\vnbstmt{\begin{array}{@{}l@{}} \forall v, x, e. \\ \quad (\operatorname{is\text{-}hilbert\text{-}space}(v) \wedge (\operatorname{is\text{-}orthonormal\text{-}basis}(v, x, e) \wedge \operatorname{is\text{-}countable}(x))) \Rightarrow  \\ \quad \quad \operatorname{is\text{-}separable}(\operatorname{ip\text{-}metric\text{-}space}(v)) \end{array}}}
\vnbprop{hilbert-separable-iff-countable-basis}{\vnbstmt{\begin{array}{@{}l@{}} \forall v. \\ \quad \operatorname{is\text{-}hilbert\text{-}space}(v) \Rightarrow  \\ \quad \quad \operatorname{is\text{-}separable}(\operatorname{ip\text{-}metric\text{-}space}(v)) \\ \quad \quad \Leftrightarrow\; \exists x, e.\; (\operatorname{is\text{-}orthonormal\text{-}basis}(v, x, e) \wedge \operatorname{is\text{-}countable}(x)) \end{array}}}
\vnbprop{cc-limit-finsum}{\vnbstmt{\begin{array}{@{}l@{}} \forall x.; \forall w \in \mathbf{Set}, u \in (\mathbb{N} \to (x \to \mathbb{C})), v \in (x \to \mathbb{C}). \\ \quad |w| \in \mathbb{N} \wedge \\ \quad w \subseteq x \wedge \\ \quad \forall k \in w.\; \operatorname{converges\text{-}to}(\operatorname{cc\text{-}ms}, (\lambda m \in \mathbb{N}.\; u(m)(k)), v(k)) \Rightarrow \\ \quad \quad \operatorname{converges\text{-}to}(\operatorname{cc\text{-}ms}, \\ \quad \quad \quad (\lambda m \in \mathbb{N}.\; \operatorname{finsum}(\operatorname{normed\text{-}ag\text{-}as\text{-}abelian\text{-}group}(\operatorname{cc\text{-}nag}), u(m), w)), \\ \quad \quad \quad \operatorname{finsum}(\operatorname{normed\text{-}ag\text{-}as\text{-}abelian\text{-}group}(\operatorname{cc\text{-}nag}), v, w)) \end{array}}}
\vnbprop{cc-limit-magnitude-sq}{\vnbstmt{\begin{array}{@{}l@{}} \forall s \in (\mathbb{N} \to \mathbb{C}), v \in \mathbb{C}, h \in (\mathbb{N} \to \mathbb{R}). \\ \quad (\forall j \in \mathbb{N}.\; (h(j) = (\operatorname{magnitude}(s(j)) \cdot \operatorname{magnitude}(s(j)))) \wedge \operatorname{converges\text{-}to}(\operatorname{cc\text{-}ms}, s, v)) \Rightarrow  \\ \quad \quad \operatorname{converges\text{-}to}(\operatorname{cc\text{-}ms}, h, (\operatorname{magnitude}(v) \cdot \operatorname{magnitude}(v))) \end{array}}}
\vnbprop{l2-dist-esum}{\vnbstmt{\begin{array}{@{}l@{}} \forall x \in \mathbf{Set}, a, c \in \operatorname{l2\text{-}families}(x). \\ \quad \operatorname{dist}(\operatorname{ip\text{-}metric\text{-}space}(\operatorname{l2\text{-}cips}(x)))(a, c) \\ \quad =\; \sqrt{\operatorname{esum}((\lambda z \in x.\; (\operatorname{magnitude}((a(z) - c(z))) \cdot \operatorname{magnitude}((a(z) - c(z))))))} \wedge \\ \quad \operatorname{esum}((\lambda z \in x.\; (\operatorname{magnitude}((a(z) - c(z))) \cdot \operatorname{magnitude}((a(z) - c(z)))))) \in \mathbb{R} \wedge \\ \quad 0 \leq \operatorname{esum}((\lambda z \in x.\; (\operatorname{magnitude}((a(z) - c(z))) \cdot \operatorname{magnitude}((a(z) - c(z)))))) \end{array}}}
\vnbprop{l2-partial-le}{\vnbstmt{\begin{array}{@{}l@{}} \forall x \in \mathbf{Set}, a, c \in \operatorname{l2\text{-}families}(x), w \in \mathbf{Set}. \\ \quad ((|w| \in \mathbb{N}) \wedge (w \subseteq x)) \Rightarrow  \\ \quad \quad \operatorname{finsum}(\operatorname{normed\text{-}ag\text{-}as\text{-}abelian\text{-}group}(\operatorname{cc\text{-}nag}), \\ \quad \quad \quad (\lambda z \in x.\; (\operatorname{magnitude}((a(z) - c(z))) \cdot \operatorname{magnitude}((a(z) - c(z))))), \\ \quad \quad \quad w) \\ \quad \quad \in\; \mathbb{R} \wedge \\ \quad \quad \operatorname{finsum}(\operatorname{normed\text{-}ag\text{-}as\text{-}abelian\text{-}group}(\operatorname{cc\text{-}nag}), \\ \quad \quad \quad (\lambda z \in x.\; (\operatorname{magnitude}((a(z) - c(z))) \cdot \operatorname{magnitude}((a(z) - c(z))))), \\ \quad \quad \quad w) \\ \quad \quad \leq\; \operatorname{esum}((\lambda z \in x.\; (\operatorname{magnitude}((a(z) - c(z))) \cdot \operatorname{magnitude}((a(z) - c(z)))))) \end{array}}}
\vnbprop{l2-coord-le}{\vnbstmt{\begin{array}{@{}l@{}} \forall x \in \mathbf{Set}, a, c \in \operatorname{l2\text{-}families}(x), k \in x. \\ \quad \operatorname{magnitude}((a(k) - c(k))) \leq \operatorname{dist}(\operatorname{ip\text{-}metric\text{-}space}(\operatorname{l2\text{-}cips}(x)))(a, c) \end{array}}}
\vnbprop{cc-const-converges}{\vnbstmt{\forall c \in \mathbb{C}.\; \operatorname{converges\text{-}to}(\operatorname{cc\text{-}ms}, (\lambda m \in \mathbb{N}.\; c), c)}}
\vnbprop{l2-cauchy-esum-tail}{\vnbstmt{\begin{array}{@{}l@{}} \forall x \in \mathbf{Set}, f \in (\mathbb{N} \to \operatorname{l2\text{-}families}(x)), b \in (x \to \mathbb{C}), e. \\ \quad \forall k \in x.\; \operatorname{converges\text{-}to}(\operatorname{cc\text{-}ms}, (\lambda \mathit{b5m} \in \mathbb{N}.\; f(\mathit{b5m})(k)), b(k)) \wedge \\ \quad \forall s. \\ \quad \quad \operatorname{pos\text{-}rr}(s) \Rightarrow  \\ \quad \quad \quad \exists n \in \mathbb{N}. \forall m, a \in \mathbb{N}. \\ \quad \quad \quad \quad ((n \leq m) \wedge (n \leq a)) \Rightarrow (\operatorname{dist}(\operatorname{ip\text{-}metric\text{-}space}(\operatorname{l2\text{-}cips}(x)))(f(m), f(a)) \leq s) \wedge \\ \quad \operatorname{pos\text{-}rr}(e) \Rightarrow \\ \quad \quad \exists \mathit{b5n} \in \mathbb{N}. \forall i \in \mathbb{N}. \\ \quad \quad \quad (\mathit{b5n} \leq i) \Rightarrow  \\ \quad \quad \quad \quad \operatorname{esum}((\lambda z \in x.\; (\operatorname{magnitude}((f(i)(z) - b(z))) \cdot \operatorname{magnitude}((f(i)(z) - b(z)))))) \in \mathbb{R} \wedge \\ \quad \quad \quad \quad \operatorname{esum}((\lambda z \in x.\; (\operatorname{magnitude}((f(i)(z) - b(z))) \cdot \operatorname{magnitude}((f(i)(z) - b(z)))))) \leq (e \cdot e) \end{array}}}
\vnbprop{l2-is-hilbert-space}{\vnbstmt{\forall x \in \mathbf{Set}.\; \operatorname{is\text{-}hilbert\text{-}space}(\operatorname{l2\text{-}cips}(x))}}

\vnbbookfile{theorem-library/cc-sums-ptwise}{theorem-library/cc-sums-ptwise}
\noindent{\small cc-sums-ptwise.scm -- wave of 2026-11-09/10 (agent TH3-B2), wired 2026-11-10 from scratchpad/th3-1109/b2/tb2-sums.scm. Library: operator-algebras.}\par
\vnbprop{cc-sums-to-add-ptwise}{\vnbstmt{\begin{array}{@{}l@{}} \forall d, f, g, h, r, s. \\ \quad f \in (d \to \mathbb{C}) \wedge \\ \quad g \in (d \to \mathbb{C}) \wedge \\ \quad h \in (d \to \mathbb{C}) \wedge \\ \quad r \in \mathbb{C} \wedge \\ \quad s \in \mathbb{C} \wedge \\ \quad \operatorname{sums\text{-}to}(\operatorname{cc\text{-}nag}, f, r) \wedge \\ \quad \operatorname{sums\text{-}to}(\operatorname{cc\text{-}nag}, g, s) \wedge \\ \quad \forall z \in d.\; (h(z) = (f(z) + g(z))) \Rightarrow \\ \quad \quad \operatorname{sums\text{-}to}(\operatorname{cc\text{-}nag}, h, (r + s)) \end{array}}}
\vnbprop{cc-sums-to-scale-ptwise}{\vnbstmt{\begin{array}{@{}l@{}} \forall d, c, f, h, r. \\ \quad c \in \mathbb{C} \wedge \\ \quad f \in (d \to \mathbb{C}) \wedge \\ \quad h \in (d \to \mathbb{C}) \wedge \\ \quad r \in \mathbb{C} \wedge \\ \quad \operatorname{sums\text{-}to}(\operatorname{cc\text{-}nag}, f, r) \wedge \\ \quad \forall z \in d.\; (h(z) = (c \cdot f(z))) \Rightarrow \\ \quad \quad \operatorname{sums\text{-}to}(\operatorname{cc\text{-}nag}, h, (c \cdot r)) \end{array}}}
\vnbprop{cc-sums-to-cscale-ptwise}{\vnbstmt{\begin{array}{@{}l@{}} \forall d, c, f, h, r. \\ \quad c \in \mathbb{C} \wedge \\ \quad f \in (d \to \mathbb{C}) \wedge \\ \quad h \in (d \to \mathbb{C}) \wedge \\ \quad r \in \mathbb{C} \wedge \\ \quad \operatorname{sums\text{-}to}(\operatorname{cc\text{-}nag}, f, r) \wedge \\ \quad \forall z \in d.\; (h(z) = (c \cdot \operatorname{conjugate}(f(z)))) \Rightarrow \\ \quad \quad \operatorname{sums\text{-}to}(\operatorname{cc\text{-}nag}, h, (c \cdot \operatorname{conjugate}(r))) \end{array}}}
\vnbprop{cc-sums-to-comb}{\vnbstmt{\begin{array}{@{}l@{}} \forall d, a, b, c, h, f, g, k, m, n, p, q. \\ \quad a \in (d \to \mathbb{C}) \wedge \\ \quad b \in (d \to \mathbb{C}) \wedge \\ \quad c \in (d \to \mathbb{C}) \wedge \\ \quad h \in (d \to \mathbb{C}) \wedge \\ \quad f \in \mathbb{C} \wedge \\ \quad g \in \mathbb{C} \wedge \\ \quad k \in \mathbb{C} \wedge \\ \quad m \in \mathbb{C} \wedge \\ \quad n \in \mathbb{C} \wedge \\ \quad p \in \mathbb{C} \wedge \\ \quad q \in \mathbb{C} \wedge \\ \quad \operatorname{sums\text{-}to}(\operatorname{cc\text{-}nag}, a, f) \wedge \\ \quad \operatorname{sums\text{-}to}(\operatorname{cc\text{-}nag}, b, g) \wedge \\ \quad \operatorname{sums\text{-}to}(\operatorname{cc\text{-}nag}, c, k) \wedge \\ \quad \forall z \in d. \\ \quad \quad h(z) = (((m \cdot a(z)) + (n \cdot b(z))) + ((p \cdot \operatorname{conjugate}(b(z))) + (q \cdot c(z)))) \Rightarrow \\ \quad \quad \operatorname{sums\text{-}to}(\operatorname{cc\text{-}nag}, h, (((m \cdot f) + (n \cdot g)) + ((p \cdot \operatorname{conjugate}(g)) + (q \cdot k)))) \end{array}}}

\vnbbookfile{theorem-library/hs-inner-product}{theorem-library/hs-inner-product}
\noindent{\small hs-inner-product.scm -- wave of 2026-11-09/10 (agent TH3-B2), wired 2026-11-10 from scratchpad/th3-1109/b2/tb2-hip.scm. Library: operator-algebras.}\par
\vnbprop{hs-ip-family-sums}{\vnbstmt{\begin{array}{@{}l@{}} \forall v, w, s, t, x, e. \\ \quad \operatorname{is\text{-}hilbert\text{-}space}(v) \wedge \\ \quad \operatorname{is\text{-}hilbert\text{-}space}(w) \wedge \\ \quad \operatorname{is\text{-}hilbert\text{-}schmidt}(v, w, s) \wedge \\ \quad \operatorname{is\text{-}hilbert\text{-}schmidt}(v, w, t) \wedge \\ \quad \operatorname{is\text{-}orthonormal\text{-}basis}(v, x, e) \Rightarrow \\ \quad \quad \operatorname{cc\text{-}family\text{-}sum}((\lambda k \in x.\; \langle\cdot,\cdot\rangle_{w}(s(e(k)), t(e(k))))) \in \mathbb{C} \wedge \\ \quad \quad \operatorname{sums\text{-}to}(\operatorname{cc\text{-}nag}, (\lambda k \in x.\; \langle\cdot,\cdot\rangle_{w}(s(e(k)), t(e(k)))), \\ \quad \quad \quad \operatorname{cc\text{-}family\text{-}sum}((\lambda k \in x.\; \langle\cdot,\cdot\rangle_{w}(s(e(k)), t(e(k)))))) \end{array}}}
\vnbprop{hs-sum-polar}{\vnbstmt{\begin{array}{@{}l@{}} \forall v, w, s, t.; \forall c \in \mathbb{C}, x, e. \\ \quad \operatorname{is\text{-}hilbert\text{-}space}(v) \wedge \\ \quad \operatorname{is\text{-}hilbert\text{-}space}(w) \wedge \\ \quad \operatorname{is\text{-}hilbert\text{-}schmidt}(v, w, s) \wedge \\ \quad \operatorname{is\text{-}hilbert\text{-}schmidt}(v, w, t) \wedge \\ \quad \operatorname{is\text{-}orthonormal\text{-}basis}(v, x, e) \Rightarrow \\ \quad \quad \sum_{i \in {x}} \lVert {{s} + {{c}\,{t}}} {e}_i \rVert^2 \\ \quad \quad =\; (\sum_{i \in {x}} \lVert {s} {e}_i \rVert^2 + (\operatorname{conjugate}(c) \cdot \operatorname{cc\text{-}family\text{-}sum}((\lambda k \in x.\; \langle\cdot,\cdot\rangle_{w}(s(e(k)), t(e(k))))))) + ((c \cdot \operatorname{conjugate}(\operatorname{cc\text{-}family\text{-}sum}((\lambda k \in x.\; \langle\cdot,\cdot\rangle_{w}(s(e(k)), t(e(k))))))) + ((c \cdot \operatorname{conjugate}(c)) \cdot \sum_{i \in {x}} \lVert {t} {e}_i \rVert^2)) \end{array}}}

\vnbbookfile{theorem-library/hs-inner-product-2}{theorem-library/hs-inner-product-2}
\noindent{\small hs-inner-product-2.scm -- wave of 2026-11-09/10 (agent TH3-B2), wired 2026-11-10 from scratchpad/th3-1109/b2/tb2-hip2.scm. Library: operator-algebras.}\par
\vnbprop{hs-ip-basis-independent}{\vnbstmt{\begin{array}{@{}l@{}} \forall v, w, s, t, x, e, a, f. \\ \quad \operatorname{is\text{-}hilbert\text{-}space}(v) \wedge \\ \quad \operatorname{is\text{-}hilbert\text{-}space}(w) \wedge \\ \quad \operatorname{is\text{-}hilbert\text{-}schmidt}(v, w, s) \wedge \\ \quad \operatorname{is\text{-}hilbert\text{-}schmidt}(v, w, t) \wedge \\ \quad \operatorname{is\text{-}orthonormal\text{-}basis}(v, x, e) \wedge \\ \quad \operatorname{is\text{-}orthonormal\text{-}basis}(v, a, f) \Rightarrow \\ \quad \quad \operatorname{cc\text{-}family\text{-}sum}((\lambda k \in x.\; \langle\cdot,\cdot\rangle_{w}(s(e(k)), t(e(k))))) \\ \quad \quad =\; \operatorname{cc\text{-}family\text{-}sum}((\lambda k \in a.\; \langle\cdot,\cdot\rangle_{w}(s(f(k)), t(f(k))))) \end{array}}}
\vnbprop{hs-ip-any-basis}{\vnbstmt{\begin{array}{@{}l@{}} \forall v, w, s, t, x, e. \\ \quad \operatorname{is\text{-}hilbert\text{-}space}(v) \wedge \\ \quad \operatorname{is\text{-}hilbert\text{-}space}(w) \wedge \\ \quad \operatorname{is\text{-}hilbert\text{-}schmidt}(v, w, s) \wedge \\ \quad \operatorname{is\text{-}hilbert\text{-}schmidt}(v, w, t) \wedge \\ \quad \operatorname{is\text{-}orthonormal\text{-}basis}(v, x, e) \Rightarrow \\ \quad \quad \langle {s} \mid {t} \rangle_{HS} = \operatorname{cc\text{-}family\text{-}sum}((\lambda k \in x.\; \langle\cdot,\cdot\rangle_{w}(s(e(k)), t(e(k))))) \end{array}}}
\vnbprop{hs-ip-family-conj}{\vnbstmt{\begin{array}{@{}l@{}} \forall v, w, s, t, x, e. \\ \quad \operatorname{is\text{-}hilbert\text{-}space}(v) \wedge \\ \quad \operatorname{is\text{-}hilbert\text{-}space}(w) \wedge \\ \quad \operatorname{is\text{-}hilbert\text{-}schmidt}(v, w, s) \wedge \\ \quad \operatorname{is\text{-}hilbert\text{-}schmidt}(v, w, t) \wedge \\ \quad \operatorname{is\text{-}orthonormal\text{-}basis}(v, x, e) \Rightarrow \\ \quad \quad \operatorname{cc\text{-}family\text{-}sum}((\lambda k \in x.\; \langle\cdot,\cdot\rangle_{w}(t(e(k)), s(e(k))))) \\ \quad \quad =\; \operatorname{conjugate}(\operatorname{cc\text{-}family\text{-}sum}((\lambda k \in x.\; \langle\cdot,\cdot\rangle_{w}(s(e(k)), t(e(k)))))) \end{array}}}
\vnbprop{hs-ip-adjoint-swap}{\vnbstmt{\begin{array}{@{}l@{}} \forall v, w, s, t, x, e, a, f. \\ \quad \operatorname{is\text{-}hilbert\text{-}space}(v) \wedge \\ \quad \operatorname{is\text{-}hilbert\text{-}space}(w) \wedge \\ \quad \operatorname{is\text{-}hilbert\text{-}schmidt}(v, w, s) \wedge \\ \quad \operatorname{is\text{-}hilbert\text{-}schmidt}(v, w, t) \wedge \\ \quad \operatorname{is\text{-}orthonormal\text{-}basis}(v, x, e) \wedge \\ \quad \operatorname{is\text{-}orthonormal\text{-}basis}(w, a, f) \Rightarrow \\ \quad \quad \operatorname{cc\text{-}family\text{-}sum}((\lambda k \in x.\; \langle\cdot,\cdot\rangle_{w}(s(e(k)), t(e(k))))) \\ \quad \quad =\; \operatorname{cc\text{-}family\text{-}sum}((\lambda k \in a.\; \langle\cdot,\cdot\rangle_{v}({t}^{*}(f(k)), {s}^{*}(f(k))))) \end{array}}}
\vnbprop{hs-ip-conj-sym}{\vnbstmt{\begin{array}{@{}l@{}} \forall v, w, s, t. \\ \quad \operatorname{is\text{-}hilbert\text{-}space}(v) \wedge \\ \quad \operatorname{is\text{-}hilbert\text{-}space}(w) \wedge \\ \quad \operatorname{is\text{-}hilbert\text{-}schmidt}(v, w, s) \wedge \\ \quad \operatorname{is\text{-}hilbert\text{-}schmidt}(v, w, t) \Rightarrow \\ \quad \quad \langle {t} \mid {s} \rangle_{HS} = \operatorname{conjugate}(\langle {s} \mid {t} \rangle_{HS}) \end{array}}}
\vnbprop{hs-ip-adjoint}{\vnbstmt{\begin{array}{@{}l@{}} \forall v, w, s, t. \\ \quad \operatorname{is\text{-}hilbert\text{-}space}(v) \wedge \\ \quad \operatorname{is\text{-}hilbert\text{-}space}(w) \wedge \\ \quad \operatorname{is\text{-}hilbert\text{-}schmidt}(v, w, s) \wedge \\ \quad \operatorname{is\text{-}hilbert\text{-}schmidt}(v, w, t) \Rightarrow \\ \quad \quad \langle {s} \mid {t} \rangle_{HS} = \langle {{t}^{*}} \mid {{s}^{*}} \rangle_{HS} \end{array}}}
\vnbprop{trace-family-comp}{\vnbstmt{\begin{array}{@{}l@{}} \forall v, s, t, x, e. \\ \quad \operatorname{is\text{-}hilbert\text{-}space}(v) \wedge \\ \quad \operatorname{is\text{-}hilbert\text{-}schmidt}(v, v, s) \wedge \\ \quad \operatorname{is\text{-}hilbert\text{-}schmidt}(v, v, t) \wedge \\ \quad \operatorname{is\text{-}orthonormal\text{-}basis}(v, x, e) \Rightarrow \\ \quad \quad \operatorname{sums\text{-}to}(\operatorname{cc\text{-}nag}, (\lambda k \in x.\; \langle\cdot,\cdot\rangle_{v}({s}{t}(e(k)), e(k))), \\ \quad \quad \quad \operatorname{cc\text{-}family\text{-}sum}((\lambda k \in x.\; \langle\cdot,\cdot\rangle_{v}(t(e(k)), {s}^{*}(e(k)))))) \wedge \\ \quad \quad \operatorname{cc\text{-}family\text{-}sum}((\lambda k \in x.\; \langle\cdot,\cdot\rangle_{v}({s}{t}(e(k)), e(k)))) \\ \quad \quad =\; \operatorname{cc\text{-}family\text{-}sum}((\lambda k \in x.\; \langle\cdot,\cdot\rangle_{v}(t(e(k)), {s}^{*}(e(k))))) \end{array}}}
\vnbprop{op-trace-canon}{\vnbstmt{\begin{array}{@{}l@{}} \forall v.; \forall t \in L({v}, {v}). \\ \quad \operatorname{is\text{-}hilbert\text{-}space}(v) \Rightarrow  \\ \quad \quad (\mathrm{tr}\,{t} \simeq \operatorname{cc\text{-}family\text{-}sum}((\lambda k \in \operatorname{canon\text{-}onb}(v).\; \langle\cdot,\cdot\rangle_{v}(t((\lambda e \in \operatorname{canon\text{-}onb}(v).\; e)(k)), (\lambda e \in \operatorname{canon\text{-}onb}(v).\; e)(k))))) \end{array}}}
\vnbprop{op-trace-sums}{\vnbstmt{\begin{array}{@{}l@{}} \forall v, t, x, e. \\ \quad (\operatorname{is\text{-}hilbert\text{-}space}(v) \wedge (\operatorname{is\text{-}trace\text{-}class}(v, t) \wedge \operatorname{is\text{-}orthonormal\text{-}basis}(v, x, e))) \Rightarrow  \\ \quad \quad ((\mathrm{tr}\,{t} \in \mathbb{C}) \wedge \operatorname{sums\text{-}to}(\operatorname{cc\text{-}nag}, (\lambda k \in x.\; \langle\cdot,\cdot\rangle_{v}(t(e(k)), e(k))), \mathrm{tr}\,{t})) \end{array}}}
\vnbprop{op-trace-add}{\vnbstmt{\begin{array}{@{}l@{}} \forall v, s, t. \\ \quad (\operatorname{is\text{-}hilbert\text{-}space}(v) \wedge (\operatorname{is\text{-}trace\text{-}class}(v, s) \wedge \operatorname{is\text{-}trace\text{-}class}(v, t))) \Rightarrow  \\ \quad \quad \mathrm{tr}\,{{s} + {t}} = (\mathrm{tr}\,{s} + \mathrm{tr}\,{t}) \end{array}}}
\vnbprop{op-trace-comm}{\vnbstmt{\begin{array}{@{}l@{}} \forall v, s, t. \\ \quad (\operatorname{is\text{-}hilbert\text{-}space}(v) \wedge (\operatorname{is\text{-}hilbert\text{-}schmidt}(v, v, s) \wedge \operatorname{is\text{-}hilbert\text{-}schmidt}(v, v, t))) \Rightarrow  \\ \quad \quad \mathrm{tr}\,{{s}{t}} = \mathrm{tr}\,{{t}{s}} \end{array}}}

\vnbbookfile{theorem-library/hs-inner-product-3}{theorem-library/hs-inner-product-3}
\noindent{\small hs-inner-product-3.scm -- wave of 2026-11-09/10 (agent TH3-B2), wired 2026-11-10 from scratchpad/th3-1109/b2/tb2-hip3.scm. Library: operator-algebras.}\par
\vnbprop{hs-ip-self}{\vnbstmt{\begin{array}{@{}l@{}} \forall v, w, t, x, e. \\ \quad \operatorname{is\text{-}hilbert\text{-}space}(v) \wedge \\ \quad \operatorname{is\text{-}hilbert\text{-}space}(w) \wedge \\ \quad \operatorname{is\text{-}hilbert\text{-}schmidt}(v, w, t) \wedge \\ \quad \operatorname{is\text{-}orthonormal\text{-}basis}(v, x, e) \Rightarrow \\ \quad \quad \langle {t} \mid {t} \rangle_{HS} = \sum_{i \in {x}} \lVert {t} {e}_i \rVert^2 \end{array}}}
\vnbprop{hs-ip-add-left}{\vnbstmt{\begin{array}{@{}l@{}} \forall v, w, s, t, u. \\ \quad \operatorname{is\text{-}hilbert\text{-}space}(v) \wedge \\ \quad \operatorname{is\text{-}hilbert\text{-}space}(w) \wedge \\ \quad \operatorname{is\text{-}hilbert\text{-}schmidt}(v, w, s) \wedge \\ \quad \operatorname{is\text{-}hilbert\text{-}schmidt}(v, w, t) \wedge \\ \quad \operatorname{is\text{-}hilbert\text{-}schmidt}(v, w, u) \Rightarrow \\ \quad \quad \langle {{s} + {t}} \mid {u} \rangle_{HS} = (\langle {s} \mid {u} \rangle_{HS} + \langle {t} \mid {u} \rangle_{HS}) \end{array}}}
\vnbprop{trace-class-comp}{\vnbstmt{\begin{array}{@{}l@{}} \forall v, s, t. \\ \quad (\operatorname{is\text{-}hilbert\text{-}space}(v) \wedge (\operatorname{is\text{-}hilbert\text{-}schmidt}(v, v, s) \wedge \operatorname{is\text{-}hilbert\text{-}schmidt}(v, v, t))) \Rightarrow  \\ \quad \quad \operatorname{is\text{-}trace\text{-}class}(v, {s}{t}) \end{array}}}
\vnbprop{op-trace-scale}{\vnbstmt{\begin{array}{@{}l@{}} \forall v, c, t. \\ \quad (\operatorname{is\text{-}hilbert\text{-}space}(v) \wedge ((c \in \mathbb{C}) \wedge \operatorname{is\text{-}trace\text{-}class}(v, t))) \Rightarrow  \\ \quad \quad \mathrm{tr}\,{{c}\,{t}} = (c \cdot \mathrm{tr}\,{t}) \end{array}}}

\vnbbookfile{structure-library/ipf-vocab}{structure-library/ipf-vocab}
\noindent{\small ipf-vocab.scm -- INNER PRODUCT FORMS ON A COMPLEX MODULE, the module carrying such a form, and the quotient of a module by the null space of a positive semi-definite form.  Agent BA-1111-A2, 2026-11-11.  Destination: structure-library/ipf-vocab.scm.  Library: analysis. DEFINITIONS ONLY: no proof, no support, no axiom.}\par
\par\medskip\noindent\textbf{Definition} (\texttt{is-hermitian-form}).\;For all $\operatorname{ipfm\_}$ and $\operatorname{ipff\_}$, 
\begin{equation*}
\begin{array}{@{}l@{}} \operatorname{is\text{-}hermitian\text{-}form}(\operatorname{ipfm\_}, \operatorname{ipff\_}) \\ \Leftrightarrow\; \operatorname{ipff\_} \in ((\operatorname{vec}(\operatorname{ipfm\_}) \times \operatorname{vec}(\operatorname{ipfm\_})) \to \mathbb{C}) \wedge \\ \forall \operatorname{ipfx\_}, \operatorname{ipfy\_}, \operatorname{ipfz\_} \in \operatorname{vec}(\operatorname{ipfm\_}). \\ \quad \operatorname{ipff\_}(\operatorname{vadd}(\operatorname{ipfm\_})(\operatorname{ipfx\_}, \operatorname{ipfy\_}), \operatorname{ipfz\_}) = +(\operatorname{ipff\_}(\operatorname{ipfx\_}, \operatorname{ipfz\_}), \operatorname{ipff\_}(\operatorname{ipfy\_}, \operatorname{ipfz\_})) \wedge \\ \forall \operatorname{ipfa\_} \in \mathbb{C}, \operatorname{ipfx\_}, \operatorname{ipfy\_} \in \operatorname{vec}(\operatorname{ipfm\_}). \\ \quad \operatorname{ipff\_}(\operatorname{act}(\operatorname{ipfm\_})(\operatorname{ipfa\_}, \operatorname{ipfx\_}), \operatorname{ipfy\_}) = \cdot(\operatorname{ipfa\_}, \operatorname{ipff\_}(\operatorname{ipfx\_}, \operatorname{ipfy\_})) \wedge \\ \forall \operatorname{ipfx\_}, \operatorname{ipfy\_} \in \operatorname{vec}(\operatorname{ipfm\_}). \\ \quad \operatorname{ipff\_}(\operatorname{ipfy\_}, \operatorname{ipfx\_}) = \operatorname{conjugate}(\operatorname{ipff\_}(\operatorname{ipfx\_}, \operatorname{ipfy\_})) \end{array}
\end{equation*}

\par\medskip\noindent\textbf{Definition} (\texttt{is-psd-form}).\;For all $\operatorname{ipfm\_}$ and $\operatorname{ipff\_}$, 
\begin{equation*}
\begin{array}{@{}l@{}} \operatorname{is\text{-}psd\text{-}form}(\operatorname{ipfm\_}, \operatorname{ipff\_}) \\ \Leftrightarrow\; \operatorname{is\text{-}hermitian\text{-}form}(\operatorname{ipfm\_}, \operatorname{ipff\_}) \wedge \\ \forall \operatorname{ipfx\_} \in \operatorname{vec}(\operatorname{ipfm\_}).\; (\operatorname{ipff\_}(\operatorname{ipfx\_}, \operatorname{ipfx\_}) \in \mathbb{R}) \wedge \\ \forall \operatorname{ipfx\_} \in \operatorname{vec}(\operatorname{ipfm\_}).\; (0 \leq \operatorname{ipff\_}(\operatorname{ipfx\_}, \operatorname{ipfx\_})) \end{array}
\end{equation*}

\par\medskip\noindent\textbf{Definition} (\texttt{is-inner-product-form}).\;For all $\operatorname{ipfm\_}$ and $\operatorname{ipff\_}$, 
\begin{equation*}
\begin{array}{@{}l@{}} \operatorname{is\text{-}inner\text{-}product\text{-}form}(\operatorname{ipfm\_}, \operatorname{ipff\_}) \\ \Leftrightarrow\; \operatorname{is\text{-}psd\text{-}form}(\operatorname{ipfm\_}, \operatorname{ipff\_}) \wedge \\ \forall \operatorname{ipfx\_} \in \operatorname{vec}(\operatorname{ipfm\_}). \\ \quad (\operatorname{ipff\_}(\operatorname{ipfx\_}, \operatorname{ipfx\_}) = 0) \Rightarrow (\operatorname{ipfx\_} = \operatorname{vzero}(\operatorname{ipfm\_})) \end{array}
\end{equation*}

\par\medskip\noindent\textbf{Definition} (\texttt{module-with-form}).\;$({\operatorname{ipfm\_}}, {\operatorname{ipff\_}}) :=$
\begin{equation*}
{[}\operatorname{scal}(\operatorname{ipfm\_}), \operatorname{vec}(\operatorname{ipfm\_}), \operatorname{vadd}(\operatorname{ipfm\_}), \operatorname{vzero}(\operatorname{ipfm\_}), \operatorname{vneg}(\operatorname{ipfm\_}), \operatorname{act}(\operatorname{ipfm\_}), \operatorname{ipff\_}]
\end{equation*}

\par\medskip\noindent\textbf{Definition} (\texttt{form-null}).\;$N_{\operatorname{ipff\_}} :=$ $\{\operatorname{ipfx\_} \in \operatorname{vec}(\operatorname{ipfm\_}) : (\operatorname{ipff\_}(\operatorname{ipfx\_}, \operatorname{ipfx\_}) = 0)\}$
\par\medskip\noindent\textbf{Definition} (\texttt{form-qip}).\;$\operatorname{form\text{-}qip}(\operatorname{ipfm\_}, \operatorname{ipff\_}) :=$ $\operatorname{descend2}(\operatorname{qmod\text{-}setoid}(\operatorname{ipfm\_}, N_{\operatorname{ipff\_}}), \operatorname{ipff\_})$
\par\medskip\noindent\textbf{Definition} (\texttt{form-quotient}).\;${\operatorname{ipfm\_}}/N_{\operatorname{ipff\_}} :=$
\begin{equation*}
({\operatorname{quotient\text{-}module}(\operatorname{ipfm\_}, N_{\operatorname{ipff\_}})}, {\operatorname{form\text{-}qip}(\operatorname{ipfm\_}, \operatorname{ipff\_})})
\end{equation*}


\vnbbookfile{structure-library/hilbert-completion-vocab}{structure-library/hilbert-completion-vocab}
\noindent{\small hc-vocab.scm -- THE HILBERT COMPLETION OF A PRE-HILBERT SPACE, CONCRETELY: the Cauchy sequences of the space as a module with pointwise operations, the limit form on them, and the quotient by its null space.  Agent BA-1111-A2, 2026-11-11.  Destination: structure-library/hc-vocab.scm. Library: analysis.  DEFINITIONS ...}\par
\par\medskip\noindent\textbf{Definition} (\texttt{form-ms}).\;$\operatorname{form\text{-}ms}(\mathit{md}, \mathit{fq}) :=$ $\operatorname{ip\text{-}metric\text{-}space}(({\mathit{md}}, {\mathit{fq}}))$
\par\medskip\noindent\textbf{Definition} (\texttt{form-cseq}).\;$\operatorname{form\text{-}cseq}(\mathit{md}, \mathit{fq}) :=$ $\operatorname{cseq}(\operatorname{form\text{-}ms}(\mathit{md}, \mathit{fq}))$
\par\medskip\noindent\textbf{Definition} (\texttt{seq-vadd}).\;$\operatorname{seq\text{-}vadd}(\mathit{md}, f, g) :=$ $(\lambda \operatorname{hcn\_} \in \mathbb{N}.\; \operatorname{vadd}(\mathit{md})(f(\operatorname{hcn\_}), g(\operatorname{hcn\_})))$
\par\medskip\noindent\textbf{Definition} (\texttt{seq-vneg}).\;$\operatorname{seq\text{-}vneg}(\mathit{md}, f) :=$ $(\lambda \operatorname{hcn\_} \in \mathbb{N}.\; \operatorname{vneg}(\mathit{md})(f(\operatorname{hcn\_})))$
\par\medskip\noindent\textbf{Definition} (\texttt{seq-act}).\;$\operatorname{seq\text{-}act}(\mathit{md}, a, f) :=$ $(\lambda \operatorname{hcn\_} \in \mathbb{N}.\; \operatorname{act}(\mathit{md})(a, f(\operatorname{hcn\_})))$
\par\medskip\noindent\textbf{Definition} (\texttt{seq-zero}).\;$\operatorname{seq\text{-}zero}(\mathit{md}) :=$ $(\lambda \operatorname{hcn\_} \in \mathbb{N}.\; \operatorname{vzero}(\mathit{md}))$
\par\medskip\noindent\textbf{Definition} (\texttt{seq-form}).\;$\operatorname{seq\text{-}form}(\mathit{fq}, f, g) :=$ $(\lambda \operatorname{hcn\_} \in \mathbb{N}.\; \operatorname{fq}(f(\operatorname{hcn\_}), g(\operatorname{hcn\_})))$
\par\medskip\noindent\textbf{Definition} (\texttt{cseq-module}).\;$\operatorname{cseq\text{-}module}(\mathit{md}, \mathit{fq}) :=$
\begin{equation*}
{[}\operatorname{scal}(\mathit{md}), \operatorname{form\text{-}cseq}(\mathit{md}, \mathit{fq}), (\lambda {[}\operatorname{hcf\_}, \operatorname{hcg\_}] \in (\operatorname{form\text{-}cseq}(\mathit{md}, \mathit{fq}) \times \operatorname{form\text{-}cseq}(\mathit{md}, \mathit{fq})).\; \operatorname{seq\text{-}vadd}(\mathit{md}, \operatorname{hcf\_}, \operatorname{hcg\_})), \operatorname{seq\text{-}zero}(\mathit{md}), (\lambda \operatorname{hcf\_} \in \operatorname{form\text{-}cseq}(\mathit{md}, \mathit{fq}).\; \operatorname{seq\text{-}vneg}(\mathit{md}, \operatorname{hcf\_})), (\lambda {[}\operatorname{hca\_}, \operatorname{hcf\_}] \in (\mathbb{C} \times \operatorname{form\text{-}cseq}(\mathit{md}, \mathit{fq})).\; \operatorname{seq\text{-}act}(\mathit{md}, \operatorname{hca\_}, \operatorname{hcf\_}))]
\end{equation*}

\par\medskip\noindent\textbf{Definition} (\texttt{limit-form}).\;$\operatorname{limit\text{-}form}(\mathit{md}, \mathit{fq}) :=$
\begin{equation*}
\begin{array}{@{}l@{}} (\lambda {[}\operatorname{hcf\_}, \operatorname{hcg\_}] \in (\operatorname{form\text{-}cseq}(\mathit{md}, \mathit{fq}) \times \operatorname{form\text{-}cseq}(\mathit{md}, \mathit{fq})).\; \\ \quad (\iota\, \operatorname{hcc\_}.\;((\operatorname{hcc\_} \in \mathbb{C}) \wedge \operatorname{converges\text{-}to}(\operatorname{cc\text{-}ms}, \operatorname{seq\text{-}form}(\mathit{fq}, \operatorname{hcf\_}, \operatorname{hcg\_}), \operatorname{hcc\_})))) \end{array}
\end{equation*}

\par\medskip\noindent\textbf{Definition} (\texttt{hilbert-completion}).\;$\operatorname{hilbert\text{-}completion}(\mathit{md}, \mathit{fq}) :=$ ${\operatorname{cseq\text{-}module}(\mathit{md}, \mathit{fq})}/N_{\operatorname{limit\text{-}form}(\mathit{md}, \mathit{fq})}$
\par\medskip\noindent\textbf{Definition} (\texttt{hc-embed}).\;$\operatorname{hc\text{-}embed}(\mathit{md}, \mathit{fq}) :=$
\begin{equation*}
\begin{array}{@{}l@{}} (\lambda \operatorname{hcx\_} \in \operatorname{vec}(\mathit{md}).\; \\ \quad \operatorname{class}(\operatorname{qmod\text{-}setoid}(\operatorname{cseq\text{-}module}(\mathit{md}, \mathit{fq}), N_{\operatorname{limit\text{-}form}(\mathit{md}, \mathit{fq})}), \\ \quad \quad \operatorname{embed\text{-}seq}(\operatorname{form\text{-}ms}(\mathit{md}, \mathit{fq}), \operatorname{hcx\_}))) \end{array}
\end{equation*}


\vnbbookfile{theorem-library/ipf-laws}{theorem-library/ipf-laws}
\noindent{\small ipf-laws.scm -- A MODULE OVER CC WITH AN INNER PRODUCT FORM IS A COMPLEX INNER PRODUCT SPACE. Agent BA-1111-A2, 2026-11-11.  Destination: theorem-library/ipf-laws.scm.  Library: analysis. Helper prefix ipl-.}\par
\vnbprop{module-with-form-scal}{\vnbstmt{\begin{array}{@{}l@{}} \forall d.; \forall q \in ((\operatorname{vec}(d) \times \operatorname{vec}(d)) \to \mathbb{C}). \\ \quad \operatorname{is\text{-}module}(d) \Rightarrow (\operatorname{scal}(({d}, {q})) \simeq \operatorname{scal}(d)) \end{array}}}
\vnbprop{module-with-form-vec}{\vnbstmt{\begin{array}{@{}l@{}} \forall d.; \forall q \in ((\operatorname{vec}(d) \times \operatorname{vec}(d)) \to \mathbb{C}). \\ \quad \operatorname{is\text{-}module}(d) \Rightarrow (\operatorname{vec}(({d}, {q})) \simeq \operatorname{vec}(d)) \end{array}}}
\vnbprop{module-with-form-vadd}{\vnbstmt{\begin{array}{@{}l@{}} \forall d.; \forall q \in ((\operatorname{vec}(d) \times \operatorname{vec}(d)) \to \mathbb{C}). \\ \quad \operatorname{is\text{-}module}(d) \Rightarrow (\operatorname{vadd}(({d}, {q})) \simeq \operatorname{vadd}(d)) \end{array}}}
\vnbprop{module-with-form-vzero}{\vnbstmt{\begin{array}{@{}l@{}} \forall d.; \forall q \in ((\operatorname{vec}(d) \times \operatorname{vec}(d)) \to \mathbb{C}). \\ \quad \operatorname{is\text{-}module}(d) \Rightarrow (\operatorname{vzero}(({d}, {q})) \simeq \operatorname{vzero}(d)) \end{array}}}
\vnbprop{module-with-form-vneg}{\vnbstmt{\begin{array}{@{}l@{}} \forall d.; \forall q \in ((\operatorname{vec}(d) \times \operatorname{vec}(d)) \to \mathbb{C}). \\ \quad \operatorname{is\text{-}module}(d) \Rightarrow (\operatorname{vneg}(({d}, {q})) \simeq \operatorname{vneg}(d)) \end{array}}}
\vnbprop{module-with-form-act}{\vnbstmt{\begin{array}{@{}l@{}} \forall d.; \forall q \in ((\operatorname{vec}(d) \times \operatorname{vec}(d)) \to \mathbb{C}). \\ \quad \operatorname{is\text{-}module}(d) \Rightarrow (\operatorname{act}(({d}, {q})) \simeq \operatorname{act}(d)) \end{array}}}
\vnbprop{module-with-form-ip}{\vnbstmt{\begin{array}{@{}l@{}} \forall d.; \forall q \in ((\operatorname{vec}(d) \times \operatorname{vec}(d)) \to \mathbb{C}). \\ \quad \operatorname{is\text{-}module}(d) \Rightarrow (\langle\cdot,\cdot\rangle_{({d}, {q})} \simeq q) \end{array}}}
\vnbprop{psd-form-zero}{\vnbstmt{\begin{array}{@{}l@{}} \forall d, q. \\ \quad (\operatorname{is\text{-}module}(d) \wedge \operatorname{is\text{-}psd\text{-}form}(d, q)) \Rightarrow (q(\operatorname{vzero}(d), \operatorname{vzero}(d)) = 0) \end{array}}}
\vnbprop{cips-of-inner-product-form}{\vnbstmt{\begin{array}{@{}l@{}} \forall d, q. \\ \quad \operatorname{is\text{-}module}(d) \wedge \\ \quad \operatorname{scal}(d) = \operatorname{normed\text{-}field\text{-}as\text{-}commutative\text{-}ring}(\operatorname{cc\text{-}normed\text{-}field}) \wedge \\ \quad \operatorname{is\text{-}inner\text{-}product\text{-}form}(d, q) \Rightarrow \\ \quad \quad \operatorname{is\text{-}complex\text{-}inner\text{-}product\text{-}space}(({d}, {q})) \end{array}}}

\vnbbookfile{theorem-library/ipf-quotient}{theorem-library/ipf-quotient}
\noindent{\small ipf-quotient.scm -- HERMITIAN FORMS ON A COMPLEX MODULE: the right-hand laws, the Cauchy-Schwarz consequence for a positive semi-definite form (a null vector is orthogonal to everything), the null space as a submodule, and THE QUOTIENT BY THE NULL SPACE AS A COMPLEX INNER PRODUCT SPACE (the pre-Hilbert step of the GNS ...}\par
\vnbprop{herm-form-type}{\vnbstmt{\begin{array}{@{}l@{}} \forall d, q. \\ \quad \operatorname{is\text{-}module}(d) \wedge \\ \quad \operatorname{scal}(d) = \operatorname{normed\text{-}field\text{-}as\text{-}commutative\text{-}ring}(\operatorname{cc\text{-}normed\text{-}field}) \wedge \\ \quad \operatorname{is\text{-}hermitian\text{-}form}(d, q) \Rightarrow \\ \quad \quad q \in ((\operatorname{vec}(d) \times \operatorname{vec}(d)) \to \mathbb{C}) \end{array}}}
\vnbprop{herm-form-add-left}{\vnbstmt{\begin{array}{@{}l@{}} \forall d, q, x, y, z. \\ \quad \operatorname{is\text{-}module}(d) \wedge \\ \quad \operatorname{scal}(d) = \operatorname{normed\text{-}field\text{-}as\text{-}commutative\text{-}ring}(\operatorname{cc\text{-}normed\text{-}field}) \wedge \\ \quad \operatorname{is\text{-}hermitian\text{-}form}(d, q) \wedge \\ \quad x \in \operatorname{vec}(d) \wedge \\ \quad y \in \operatorname{vec}(d) \wedge \\ \quad z \in \operatorname{vec}(d) \Rightarrow \\ \quad \quad q(\operatorname{vadd}(d)(x, y), z) = (q(x, z) + q(y, z)) \end{array}}}
\vnbprop{herm-form-homog-left}{\vnbstmt{\begin{array}{@{}l@{}} \forall d, q, a, x, y. \\ \quad \operatorname{is\text{-}module}(d) \wedge \\ \quad \operatorname{scal}(d) = \operatorname{normed\text{-}field\text{-}as\text{-}commutative\text{-}ring}(\operatorname{cc\text{-}normed\text{-}field}) \wedge \\ \quad \operatorname{is\text{-}hermitian\text{-}form}(d, q) \wedge \\ \quad a \in \mathbb{C} \wedge \\ \quad x \in \operatorname{vec}(d) \wedge \\ \quad y \in \operatorname{vec}(d) \Rightarrow \\ \quad \quad q(\operatorname{act}(d)(a, x), y) = (a \cdot q(x, y)) \end{array}}}
\vnbprop{herm-form-conj-sym}{\vnbstmt{\begin{array}{@{}l@{}} \forall d, q, x, y. \\ \quad \operatorname{is\text{-}module}(d) \wedge \\ \quad \operatorname{scal}(d) = \operatorname{normed\text{-}field\text{-}as\text{-}commutative\text{-}ring}(\operatorname{cc\text{-}normed\text{-}field}) \wedge \\ \quad \operatorname{is\text{-}hermitian\text{-}form}(d, q) \wedge \\ \quad x \in \operatorname{vec}(d) \wedge \\ \quad y \in \operatorname{vec}(d) \Rightarrow \\ \quad \quad q(y, x) = \operatorname{conjugate}(q(x, y)) \end{array}}}
\vnbprop{herm-form-value-type}{\vnbstmt{\begin{array}{@{}l@{}} \forall d, q, x, y. \\ \quad \operatorname{is\text{-}module}(d) \wedge \\ \quad \operatorname{scal}(d) = \operatorname{normed\text{-}field\text{-}as\text{-}commutative\text{-}ring}(\operatorname{cc\text{-}normed\text{-}field}) \wedge \\ \quad \operatorname{is\text{-}hermitian\text{-}form}(d, q) \wedge \\ \quad x \in \operatorname{vec}(d) \wedge \\ \quad y \in \operatorname{vec}(d) \Rightarrow \\ \quad \quad q(x, y) \in \mathbb{C} \end{array}}}
\vnbprop{psd-form-hermitian}{\vnbstmt{\begin{array}{@{}l@{}} \forall d, q. \\ \quad (\operatorname{is\text{-}module}(d) \wedge ((\operatorname{scal}(d) = \operatorname{normed\text{-}field\text{-}as\text{-}commutative\text{-}ring}(\operatorname{cc\text{-}normed\text{-}field})) \wedge \operatorname{is\text{-}psd\text{-}form}(d, q))) \Rightarrow  \\ \quad \quad \operatorname{is\text{-}hermitian\text{-}form}(d, q) \end{array}}}
\vnbprop{psd-form-real}{\vnbstmt{\begin{array}{@{}l@{}} \forall d, q, x. \\ \quad \operatorname{is\text{-}module}(d) \wedge \\ \quad \operatorname{scal}(d) = \operatorname{normed\text{-}field\text{-}as\text{-}commutative\text{-}ring}(\operatorname{cc\text{-}normed\text{-}field}) \wedge \\ \quad \operatorname{is\text{-}psd\text{-}form}(d, q) \wedge \\ \quad x \in \operatorname{vec}(d) \Rightarrow \\ \quad \quad q(x, x) \in \mathbb{R} \end{array}}}
\vnbprop{psd-form-nonneg}{\vnbstmt{\begin{array}{@{}l@{}} \forall d, q, x. \\ \quad \operatorname{is\text{-}module}(d) \wedge \\ \quad \operatorname{scal}(d) = \operatorname{normed\text{-}field\text{-}as\text{-}commutative\text{-}ring}(\operatorname{cc\text{-}normed\text{-}field}) \wedge \\ \quad \operatorname{is\text{-}psd\text{-}form}(d, q) \wedge \\ \quad x \in \operatorname{vec}(d) \Rightarrow \\ \quad \quad 0 \leq q(x, x) \end{array}}}
\vnbprop{herm-form-add-right}{\vnbstmt{\begin{array}{@{}l@{}} \forall d, q, x, y, z. \\ \quad \operatorname{is\text{-}module}(d) \wedge \\ \quad \operatorname{scal}(d) = \operatorname{normed\text{-}field\text{-}as\text{-}commutative\text{-}ring}(\operatorname{cc\text{-}normed\text{-}field}) \wedge \\ \quad \operatorname{is\text{-}hermitian\text{-}form}(d, q) \wedge \\ \quad x \in \operatorname{vec}(d) \wedge \\ \quad y \in \operatorname{vec}(d) \wedge \\ \quad z \in \operatorname{vec}(d) \Rightarrow \\ \quad \quad q(x, \operatorname{vadd}(d)(y, z)) = (q(x, y) + q(x, z)) \end{array}}}
\vnbprop{herm-form-homog-right}{\vnbstmt{\begin{array}{@{}l@{}} \forall d, q, a, x, y. \\ \quad \operatorname{is\text{-}module}(d) \wedge \\ \quad \operatorname{scal}(d) = \operatorname{normed\text{-}field\text{-}as\text{-}commutative\text{-}ring}(\operatorname{cc\text{-}normed\text{-}field}) \wedge \\ \quad \operatorname{is\text{-}hermitian\text{-}form}(d, q) \wedge \\ \quad a \in \mathbb{C} \wedge \\ \quad x \in \operatorname{vec}(d) \wedge \\ \quad y \in \operatorname{vec}(d) \Rightarrow \\ \quad \quad q(x, \operatorname{act}(d)(a, y)) = (\operatorname{conjugate}(a) \cdot q(x, y)) \end{array}}}
\vnbprop{psd-form-null-orth}{\vnbstmt{\begin{array}{@{}l@{}} \forall d, q, n, y. \\ \quad \operatorname{is\text{-}module}(d) \wedge \\ \quad \operatorname{scal}(d) = \operatorname{normed\text{-}field\text{-}as\text{-}commutative\text{-}ring}(\operatorname{cc\text{-}normed\text{-}field}) \wedge \\ \quad \operatorname{is\text{-}psd\text{-}form}(d, q) \wedge \\ \quad n \in \operatorname{vec}(d) \wedge \\ \quad q(n, n) = 0 \wedge \\ \quad y \in \operatorname{vec}(d) \Rightarrow \\ \quad \quad q(n, y) = 0 \end{array}}}
\vnbprop{psd-form-zero-left}{\vnbstmt{\begin{array}{@{}l@{}} \forall d, q, y. \\ \quad \operatorname{is\text{-}module}(d) \wedge \\ \quad \operatorname{scal}(d) = \operatorname{normed\text{-}field\text{-}as\text{-}commutative\text{-}ring}(\operatorname{cc\text{-}normed\text{-}field}) \wedge \\ \quad \operatorname{is\text{-}psd\text{-}form}(d, q) \wedge \\ \quad y \in \operatorname{vec}(d) \Rightarrow \\ \quad \quad q(\operatorname{vzero}(d), y) = 0 \end{array}}}
\vnbprop{psd-form-null-right}{\vnbstmt{\begin{array}{@{}l@{}} \forall d, q, n, y. \\ \quad \operatorname{is\text{-}module}(d) \wedge \\ \quad \operatorname{scal}(d) = \operatorname{normed\text{-}field\text{-}as\text{-}commutative\text{-}ring}(\operatorname{cc\text{-}normed\text{-}field}) \wedge \\ \quad \operatorname{is\text{-}psd\text{-}form}(d, q) \wedge \\ \quad n \in \operatorname{vec}(d) \wedge \\ \quad q(n, n) = 0 \wedge \\ \quad y \in \operatorname{vec}(d) \Rightarrow \\ \quad \quad q(y, n) = 0 \end{array}}}
\vnbprop{psd-form-neg-left}{\vnbstmt{\begin{array}{@{}l@{}} \forall d, q, x, y. \\ \quad \operatorname{is\text{-}module}(d) \wedge \\ \quad \operatorname{scal}(d) = \operatorname{normed\text{-}field\text{-}as\text{-}commutative\text{-}ring}(\operatorname{cc\text{-}normed\text{-}field}) \wedge \\ \quad \operatorname{is\text{-}psd\text{-}form}(d, q) \wedge \\ \quad x \in \operatorname{vec}(d) \wedge \\ \quad y \in \operatorname{vec}(d) \Rightarrow \\ \quad \quad q(\operatorname{vneg}(d)(x), y) = -q(x, y) \end{array}}}
\vnbprop{psd-form-neg-right}{\vnbstmt{\begin{array}{@{}l@{}} \forall d, q, x, y. \\ \quad \operatorname{is\text{-}module}(d) \wedge \\ \quad \operatorname{scal}(d) = \operatorname{normed\text{-}field\text{-}as\text{-}commutative\text{-}ring}(\operatorname{cc\text{-}normed\text{-}field}) \wedge \\ \quad \operatorname{is\text{-}psd\text{-}form}(d, q) \wedge \\ \quad x \in \operatorname{vec}(d) \wedge \\ \quad y \in \operatorname{vec}(d) \Rightarrow \\ \quad \quad q(x, \operatorname{vneg}(d)(y)) = -q(x, y) \end{array}}}
\vnbprop{form-null-unfold}{\vnbstmt{\forall d, q.\; (N_{q} \simeq \{x \in \operatorname{vec}(d) : (q(x, x) = 0)\})}}
\vnbprop{form-null-in}{\vnbstmt{\forall d, q.; \forall x \in \operatorname{vec}(d).\; ((q(x, x) = 0) \Rightarrow (x \in N_{q}))}}
\vnbprop{form-null-out}{\vnbstmt{\forall d, q.; \forall x \in N_{q}.\; ((x \in \operatorname{vec}(d)) \wedge (q(x, x) = 0))}}
\vnbprop{form-null-submodule}{\vnbstmt{\begin{array}{@{}l@{}} \forall d, q. \\ \quad (\operatorname{is\text{-}module}(d) \wedge ((\operatorname{scal}(d) = \operatorname{normed\text{-}field\text{-}as\text{-}commutative\text{-}ring}(\operatorname{cc\text{-}normed\text{-}field})) \wedge \operatorname{is\text{-}psd\text{-}form}(d, q))) \Rightarrow  \\ \quad \quad \operatorname{is\text{-}submodule}(d, N_{q}) \end{array}}}
\vnbprop{form-respects2}{\vnbstmt{\begin{array}{@{}l@{}} \forall d, q. \\ \quad (\operatorname{is\text{-}module}(d) \wedge ((\operatorname{scal}(d) = \operatorname{normed\text{-}field\text{-}as\text{-}commutative\text{-}ring}(\operatorname{cc\text{-}normed\text{-}field})) \wedge \operatorname{is\text{-}psd\text{-}form}(d, q))) \Rightarrow  \\ \quad \quad \operatorname{respects2}(\operatorname{qmod\text{-}setoid}(d, N_{q}), q) \end{array}}}
\vnbprop{form-qip-type}{\vnbstmt{\begin{array}{@{}l@{}} \forall d, q. \\ \quad (\operatorname{is\text{-}module}(d) \wedge ((\operatorname{scal}(d) = \operatorname{normed\text{-}field\text{-}as\text{-}commutative\text{-}ring}(\operatorname{cc\text{-}normed\text{-}field})) \wedge \operatorname{is\text{-}psd\text{-}form}(d, q))) \Rightarrow  \\ \quad \quad \operatorname{form\text{-}qip}(d, q) \\ \quad \quad \in\; ((\operatorname{quotient}(\operatorname{qmod\text{-}setoid}(d, N_{q})) \times \operatorname{quotient}(\operatorname{qmod\text{-}setoid}(d, N_{q}))) \to \mathbb{C}) \end{array}}}
\vnbprop{form-qip-computes}{\vnbstmt{\begin{array}{@{}l@{}} \forall d, q, x, y. \\ \quad \operatorname{is\text{-}module}(d) \wedge \\ \quad \operatorname{scal}(d) = \operatorname{normed\text{-}field\text{-}as\text{-}commutative\text{-}ring}(\operatorname{cc\text{-}normed\text{-}field}) \wedge \\ \quad \operatorname{is\text{-}psd\text{-}form}(d, q) \wedge \\ \quad x \in \operatorname{vec}(d) \wedge \\ \quad y \in \operatorname{vec}(d) \Rightarrow \\ \quad \quad \operatorname{form\text{-}qip}(d, q)(\operatorname{class}(\operatorname{qmod\text{-}setoid}(d, N_{q}), x), \\ \quad \quad \quad \operatorname{class}(\operatorname{qmod\text{-}setoid}(d, N_{q}), y)) \\ \quad \quad =\; q(x, y) \end{array}}}
\vnbprop{form-quotient-scal}{\vnbstmt{\begin{array}{@{}l@{}} \forall d, q. \\ \quad (\operatorname{is\text{-}module}(d) \wedge ((\operatorname{scal}(d) = \operatorname{normed\text{-}field\text{-}as\text{-}commutative\text{-}ring}(\operatorname{cc\text{-}normed\text{-}field})) \wedge \operatorname{is\text{-}psd\text{-}form}(d, q))) \Rightarrow  \\ \quad \quad \operatorname{scal}(\operatorname{quotient\text{-}module}(d, N_{q})) = \operatorname{normed\text{-}field\text{-}as\text{-}commutative\text{-}ring}(\operatorname{cc\text{-}normed\text{-}field}) \end{array}}}
\vnbprop{form-quotient-inner-product-form}{\vnbstmt{\begin{array}{@{}l@{}} \forall d, q. \\ \quad (\operatorname{is\text{-}module}(d) \wedge ((\operatorname{scal}(d) = \operatorname{normed\text{-}field\text{-}as\text{-}commutative\text{-}ring}(\operatorname{cc\text{-}normed\text{-}field})) \wedge \operatorname{is\text{-}psd\text{-}form}(d, q))) \Rightarrow  \\ \quad \quad \operatorname{is\text{-}inner\text{-}product\text{-}form}(\operatorname{quotient\text{-}module}(d, N_{q}), \operatorname{form\text{-}qip}(d, q)) \end{array}}}
\vnbprop{form-quotient-is-cips}{\vnbstmt{\begin{array}{@{}l@{}} \forall d, q. \\ \quad (\operatorname{is\text{-}module}(d) \wedge ((\operatorname{scal}(d) = \operatorname{normed\text{-}field\text{-}as\text{-}commutative\text{-}ring}(\operatorname{cc\text{-}normed\text{-}field})) \wedge \operatorname{is\text{-}psd\text{-}form}(d, q))) \Rightarrow  \\ \quad \quad \operatorname{is\text{-}complex\text{-}inner\text{-}product\text{-}space}({d}/N_{q}) \end{array}}}

\vnbbookfile{theorem-library/ipf-form-plus}{theorem-library/ipf-form-plus}
\noindent{\small ipf-form-plus.scm -- H\_Q+ IS A COMPLEX INNER PRODUCT SPACE (BanachAlgebras.pdf ch. 3, the space of Prop 3.16 / 3.17): the instance of cips-of-inner-product-form at FORM-PLUS-SPACE. Agent BA-1111-A2, 2026-11-11.  Destination: theorem-library/ipf-form-plus.scm.  Library: operator-algebras (FORM-PLUS-SPACE is ...}\par
\vnbprop{form-plus-module-scal}{\vnbstmt{\begin{array}{@{}l@{}} \forall v, d. \\ \quad (\operatorname{is\text{-}complex\text{-}inner\text{-}product\text{-}space}(v) \wedge \operatorname{is\text{-}submodule}(v, d)) \Rightarrow  \\ \quad \quad (\operatorname{scal}(\operatorname{complex\text{-}inner\text{-}product\text{-}space\text{-}as\text{-}module}(\operatorname{sub\text{-}cips}(v, d))) \simeq \operatorname{normed\text{-}field\text{-}as\text{-}commutative\text{-}ring}(\operatorname{cc\text{-}normed\text{-}field})) \end{array}}}
\vnbprop{form-plus-module-vec}{\vnbstmt{\begin{array}{@{}l@{}} \forall v, d. \\ \quad (\operatorname{is\text{-}complex\text{-}inner\text{-}product\text{-}space}(v) \wedge \operatorname{is\text{-}submodule}(v, d)) \Rightarrow  \\ \quad \quad (\operatorname{vec}(\operatorname{complex\text{-}inner\text{-}product\text{-}space\text{-}as\text{-}module}(\operatorname{sub\text{-}cips}(v, d))) \simeq d) \end{array}}}
\vnbprop{form-plus-module-vadd}{\vnbstmt{\begin{array}{@{}l@{}} \forall v, d. \\ \quad (\operatorname{is\text{-}complex\text{-}inner\text{-}product\text{-}space}(v) \wedge \operatorname{is\text{-}submodule}(v, d)) \Rightarrow  \\ \quad \quad (\operatorname{vadd}(\operatorname{complex\text{-}inner\text{-}product\text{-}space\text{-}as\text{-}module}(\operatorname{sub\text{-}cips}(v, d))) \simeq (\lambda {[}\operatorname{tsbx\_}, \operatorname{tsby\_}] \in (d \times d).\; \operatorname{vadd}(v)(\operatorname{tsbx\_}, \operatorname{tsby\_}))) \end{array}}}
\vnbprop{form-plus-module-vzero}{\vnbstmt{\begin{array}{@{}l@{}} \forall v, d. \\ \quad (\operatorname{is\text{-}complex\text{-}inner\text{-}product\text{-}space}(v) \wedge \operatorname{is\text{-}submodule}(v, d)) \Rightarrow  \\ \quad \quad (\operatorname{vzero}(\operatorname{complex\text{-}inner\text{-}product\text{-}space\text{-}as\text{-}module}(\operatorname{sub\text{-}cips}(v, d))) \simeq \operatorname{vzero}(v)) \end{array}}}
\vnbprop{form-plus-module-vneg}{\vnbstmt{\begin{array}{@{}l@{}} \forall v, d. \\ \quad (\operatorname{is\text{-}complex\text{-}inner\text{-}product\text{-}space}(v) \wedge \operatorname{is\text{-}submodule}(v, d)) \Rightarrow  \\ \quad \quad (\operatorname{vneg}(\operatorname{complex\text{-}inner\text{-}product\text{-}space\text{-}as\text{-}module}(\operatorname{sub\text{-}cips}(v, d))) \simeq (\lambda x \in d.\; \operatorname{vneg}(v)(x))) \end{array}}}
\vnbprop{form-plus-module-act}{\vnbstmt{\begin{array}{@{}l@{}} \forall v, d. \\ \quad (\operatorname{is\text{-}complex\text{-}inner\text{-}product\text{-}space}(v) \wedge \operatorname{is\text{-}submodule}(v, d)) \Rightarrow  \\ \quad \quad (\operatorname{act}(\operatorname{complex\text{-}inner\text{-}product\text{-}space\text{-}as\text{-}module}(\operatorname{sub\text{-}cips}(v, d))) \simeq (\lambda {[}\operatorname{tsbc\_}, \operatorname{tsbx\_}] \in (\mathbb{C} \times d).\; \operatorname{act}(v)(\operatorname{tsbc\_}, \operatorname{tsbx\_}))) \end{array}}}
\vnbprop{form-plus-space-as-module-with-form}{\vnbstmt{\begin{array}{@{}l@{}} \forall v, d, q. \\ \quad (\operatorname{is\text{-}complex\text{-}inner\text{-}product\text{-}space}(v) \wedge \operatorname{is\text{-}submodule}(v, d)) \Rightarrow  \\ \quad \quad (H_{{q}+} \simeq ({\operatorname{complex\text{-}inner\text{-}product\text{-}space\text{-}as\text{-}module}(\operatorname{sub\text{-}cips}(v, d))}, {{q}_{+}})) \end{array}}}
\vnbprop{form-plus-inner-product-form}{\vnbstmt{\begin{array}{@{}l@{}} \forall v, d, q. \\ \quad (\operatorname{is\text{-}complex\text{-}inner\text{-}product\text{-}space}(v) \wedge (\operatorname{is\text{-}submodule}(v, d) \wedge \operatorname{is\text{-}positive\text{-}form}(v, d, q))) \Rightarrow  \\ \quad \quad \operatorname{is\text{-}inner\text{-}product\text{-}form}(\operatorname{complex\text{-}inner\text{-}product\text{-}space\text{-}as\text{-}module}(\operatorname{sub\text{-}cips}(v, d)), {q}_{+}) \end{array}}}
\vnbprop{form-plus-space-is-cips}{\vnbstmt{\begin{array}{@{}l@{}} \forall v, d, q. \\ \quad (\operatorname{is\text{-}complex\text{-}inner\text{-}product\text{-}space}(v) \wedge (\operatorname{is\text{-}submodule}(v, d) \wedge \operatorname{is\text{-}positive\text{-}form}(v, d, q))) \Rightarrow  \\ \quad \quad \operatorname{is\text{-}complex\text{-}inner\text{-}product\text{-}space}(H_{{q}+}) \end{array}}}

\vnbbookfile{theorem-library/hilbert-completion}{theorem-library/hilbert-completion}
\noindent{\small hilbert-completion.scm -- THE HILBERT COMPLETION OF A PRE-HILBERT SPACE (md, fq).  Agent BA-1111-A2, 2026-11-11.  Vocabulary: structure-library/hc-vocab.scm.  Destination: theorem-library/hilbert-completion.scm. Library: analysis.  Helper prefix hcp-.  Nothing asserted.}\par
\vnbprop{ipf-form-psd}{\vnbstmt{\begin{array}{@{}l@{}} \forall d, q. \\ \quad \operatorname{is\text{-}module}(d) \wedge \\ \quad \operatorname{scal}(d) = \operatorname{normed\text{-}field\text{-}as\text{-}commutative\text{-}ring}(\operatorname{cc\text{-}normed\text{-}field}) \wedge \\ \quad \operatorname{is\text{-}inner\text{-}product\text{-}form}(d, q) \Rightarrow \\ \quad \quad \operatorname{is\text{-}psd\text{-}form}(d, q) \end{array}}}
\vnbprop{form-ms-is-metric-space}{\vnbstmt{\begin{array}{@{}l@{}} \forall d, q. \\ \quad \operatorname{is\text{-}module}(d) \wedge \\ \quad \operatorname{scal}(d) = \operatorname{normed\text{-}field\text{-}as\text{-}commutative\text{-}ring}(\operatorname{cc\text{-}normed\text{-}field}) \wedge \\ \quad \operatorname{is\text{-}inner\text{-}product\text{-}form}(d, q) \Rightarrow \\ \quad \quad \operatorname{is\text{-}metric\text{-}space}(\operatorname{form\text{-}ms}(d, q)) \end{array}}}
\vnbprop{form-ms-pts}{\vnbstmt{\begin{array}{@{}l@{}} \forall d, q. \\ \quad \operatorname{is\text{-}module}(d) \wedge \\ \quad \operatorname{scal}(d) = \operatorname{normed\text{-}field\text{-}as\text{-}commutative\text{-}ring}(\operatorname{cc\text{-}normed\text{-}field}) \wedge \\ \quad \operatorname{is\text{-}inner\text{-}product\text{-}form}(d, q) \Rightarrow \\ \quad \quad (\operatorname{pts}(\operatorname{form\text{-}ms}(d, q)) \simeq \operatorname{vec}(d)) \end{array}}}
\vnbprop{form-ms-dist}{\vnbstmt{\begin{array}{@{}l@{}} \forall d, q, x, y. \\ \quad \operatorname{is\text{-}module}(d) \wedge \\ \quad \operatorname{scal}(d) = \operatorname{normed\text{-}field\text{-}as\text{-}commutative\text{-}ring}(\operatorname{cc\text{-}normed\text{-}field}) \wedge \\ \quad \operatorname{is\text{-}inner\text{-}product\text{-}form}(d, q) \wedge \\ \quad x \in \operatorname{vec}(d) \wedge \\ \quad y \in \operatorname{vec}(d) \Rightarrow \\ \quad \quad \operatorname{dist}(\operatorname{form\text{-}ms}(d, q))(x, y) = \sqrt{q(\operatorname{vadd}(d)(x, \operatorname{vneg}(d)(y)), \operatorname{vadd}(d)(x, \operatorname{vneg}(d)(y)))} \end{array}}}
\vnbprop{form-norm-nonneg}{\vnbstmt{\begin{array}{@{}l@{}} \forall d, q, z. \\ \quad \operatorname{is\text{-}module}(d) \wedge \\ \quad \operatorname{scal}(d) = \operatorname{normed\text{-}field\text{-}as\text{-}commutative\text{-}ring}(\operatorname{cc\text{-}normed\text{-}field}) \wedge \\ \quad \operatorname{is\text{-}inner\text{-}product\text{-}form}(d, q) \wedge \\ \quad z \in \operatorname{vec}(d) \Rightarrow \\ \quad \quad (\sqrt{q(z, z)} \in \mathbb{R}) \wedge (0 \leq \sqrt{q(z, z)}) \end{array}}}
\vnbprop{form-minkowski}{\vnbstmt{\begin{array}{@{}l@{}} \forall d, q, x, y. \\ \quad \operatorname{is\text{-}module}(d) \wedge \\ \quad \operatorname{scal}(d) = \operatorname{normed\text{-}field\text{-}as\text{-}commutative\text{-}ring}(\operatorname{cc\text{-}normed\text{-}field}) \wedge \\ \quad \operatorname{is\text{-}inner\text{-}product\text{-}form}(d, q) \wedge \\ \quad x \in \operatorname{vec}(d) \wedge \\ \quad y \in \operatorname{vec}(d) \Rightarrow \\ \quad \quad \sqrt{q(\operatorname{vadd}(d)(x, y), \operatorname{vadd}(d)(x, y))} \leq (\sqrt{q(x, x)} + \sqrt{q(y, y)}) \end{array}}}
\vnbprop{form-schwarz}{\vnbstmt{\begin{array}{@{}l@{}} \forall d, q, x, y. \\ \quad \operatorname{is\text{-}module}(d) \wedge \\ \quad \operatorname{scal}(d) = \operatorname{normed\text{-}field\text{-}as\text{-}commutative\text{-}ring}(\operatorname{cc\text{-}normed\text{-}field}) \wedge \\ \quad \operatorname{is\text{-}inner\text{-}product\text{-}form}(d, q) \wedge \\ \quad x \in \operatorname{vec}(d) \wedge \\ \quad y \in \operatorname{vec}(d) \Rightarrow \\ \quad \quad \operatorname{magnitude}(q(x, y)) \leq (\sqrt{q(x, x)} \cdot \sqrt{q(y, y)}) \end{array}}}
\vnbprop{form-norm-sum-diff}{\vnbstmt{\begin{array}{@{}l@{}} \forall d, q, x, y, u, w. \\ \quad \operatorname{is\text{-}module}(d) \wedge \\ \quad \operatorname{scal}(d) = \operatorname{normed\text{-}field\text{-}as\text{-}commutative\text{-}ring}(\operatorname{cc\text{-}normed\text{-}field}) \wedge \\ \quad \operatorname{is\text{-}inner\text{-}product\text{-}form}(d, q) \wedge \\ \quad x \in \operatorname{vec}(d) \wedge \\ \quad y \in \operatorname{vec}(d) \wedge \\ \quad u \in \operatorname{vec}(d) \wedge \\ \quad w \in \operatorname{vec}(d) \Rightarrow \\ \quad \quad \sqrt{q(\operatorname{vadd}(d)(\operatorname{vadd}(d)(x, y), \operatorname{vneg}(d)(\operatorname{vadd}(d)(u, w))), \operatorname{vadd}(d)(\operatorname{vadd}(d)(x, y), \operatorname{vneg}(d)(\operatorname{vadd}(d)(u, w))))} \\ \quad \quad \leq\; \sqrt{q(\operatorname{vadd}(d)(x, \operatorname{vneg}(d)(u)), \operatorname{vadd}(d)(x, \operatorname{vneg}(d)(u)))} + \sqrt{q(\operatorname{vadd}(d)(y, \operatorname{vneg}(d)(w)), \operatorname{vadd}(d)(y, \operatorname{vneg}(d)(w)))} \end{array}}}
\vnbprop{form-norm-neg-diff}{\vnbstmt{\begin{array}{@{}l@{}} \forall d, q, x, y. \\ \quad \operatorname{is\text{-}module}(d) \wedge \\ \quad \operatorname{scal}(d) = \operatorname{normed\text{-}field\text{-}as\text{-}commutative\text{-}ring}(\operatorname{cc\text{-}normed\text{-}field}) \wedge \\ \quad \operatorname{is\text{-}inner\text{-}product\text{-}form}(d, q) \wedge \\ \quad x \in \operatorname{vec}(d) \wedge \\ \quad y \in \operatorname{vec}(d) \Rightarrow \\ \quad \quad \sqrt{q(\operatorname{vadd}(d)(\operatorname{vneg}(d)(x), \operatorname{vneg}(d)(\operatorname{vneg}(d)(y))), \operatorname{vadd}(d)(\operatorname{vneg}(d)(x), \operatorname{vneg}(d)(\operatorname{vneg}(d)(y))))} \\ \quad \quad =\; \sqrt{q(\operatorname{vadd}(d)(x, \operatorname{vneg}(d)(y)), \operatorname{vadd}(d)(x, \operatorname{vneg}(d)(y)))} \end{array}}}
\vnbprop{form-norm-scale-diff}{\vnbstmt{\begin{array}{@{}l@{}} \forall d, q, a, x, y. \\ \quad \operatorname{is\text{-}module}(d) \wedge \\ \quad \operatorname{scal}(d) = \operatorname{normed\text{-}field\text{-}as\text{-}commutative\text{-}ring}(\operatorname{cc\text{-}normed\text{-}field}) \wedge \\ \quad \operatorname{is\text{-}inner\text{-}product\text{-}form}(d, q) \wedge \\ \quad a \in \mathbb{C} \wedge \\ \quad x \in \operatorname{vec}(d) \wedge \\ \quad y \in \operatorname{vec}(d) \Rightarrow \\ \quad \quad \sqrt{q(\operatorname{vadd}(d)(\operatorname{act}(d)(a, x), \operatorname{vneg}(d)(\operatorname{act}(d)(a, y))), \operatorname{vadd}(d)(\operatorname{act}(d)(a, x), \operatorname{vneg}(d)(\operatorname{act}(d)(a, y))))} \\ \quad \quad =\; \operatorname{magnitude}(a) \cdot \sqrt{q(\operatorname{vadd}(d)(x, \operatorname{vneg}(d)(y)), \operatorname{vadd}(d)(x, \operatorname{vneg}(d)(y)))} \end{array}}}
\vnbprop{form-norm-bound}{\vnbstmt{\begin{array}{@{}l@{}} \forall d, q, x, y. \\ \quad \operatorname{is\text{-}module}(d) \wedge \\ \quad \operatorname{scal}(d) = \operatorname{normed\text{-}field\text{-}as\text{-}commutative\text{-}ring}(\operatorname{cc\text{-}normed\text{-}field}) \wedge \\ \quad \operatorname{is\text{-}inner\text{-}product\text{-}form}(d, q) \wedge \\ \quad x \in \operatorname{vec}(d) \wedge \\ \quad y \in \operatorname{vec}(d) \Rightarrow \\ \quad \quad \sqrt{q(x, x)} \leq (\sqrt{q(\operatorname{vadd}(d)(x, \operatorname{vneg}(d)(y)), \operatorname{vadd}(d)(x, \operatorname{vneg}(d)(y)))} + \sqrt{q(y, y)}) \end{array}}}
\vnbprop{form-diff-estimate}{\vnbstmt{\begin{array}{@{}l@{}} \forall d, q, x, y, u, w. \\ \quad \operatorname{is\text{-}module}(d) \wedge \\ \quad \operatorname{scal}(d) = \operatorname{normed\text{-}field\text{-}as\text{-}commutative\text{-}ring}(\operatorname{cc\text{-}normed\text{-}field}) \wedge \\ \quad \operatorname{is\text{-}inner\text{-}product\text{-}form}(d, q) \wedge \\ \quad x \in \operatorname{vec}(d) \wedge \\ \quad y \in \operatorname{vec}(d) \wedge \\ \quad u \in \operatorname{vec}(d) \wedge \\ \quad w \in \operatorname{vec}(d) \Rightarrow \\ \quad \quad \operatorname{magnitude}((q(x, y) + -q(u, w))) \\ \quad \quad \leq\; (\sqrt{q(\operatorname{vadd}(d)(x, \operatorname{vneg}(d)(u)), \operatorname{vadd}(d)(x, \operatorname{vneg}(d)(u)))} \cdot \sqrt{q(y, y)}) + (\sqrt{q(u, u)} \cdot \sqrt{q(\operatorname{vadd}(d)(y, \operatorname{vneg}(d)(w)), \operatorname{vadd}(d)(y, \operatorname{vneg}(d)(w)))}) \end{array}}}
\vnbprop{form-cauchy-intro}{\vnbstmt{\begin{array}{@{}l@{}} \forall d, q, f. \\ \quad \operatorname{is\text{-}module}(d) \wedge \\ \quad \operatorname{scal}(d) = \operatorname{normed\text{-}field\text{-}as\text{-}commutative\text{-}ring}(\operatorname{cc\text{-}normed\text{-}field}) \wedge \\ \quad \operatorname{is\text{-}inner\text{-}product\text{-}form}(d, q) \wedge \\ \quad f \in (\mathbb{N} \to \operatorname{vec}(d)) \wedge \\ \quad \forall s. \\ \quad \quad \operatorname{pos\text{-}rr}(s) \Rightarrow  \\ \quad \quad \quad \exists b \in \mathbb{N}. \forall a, c \in \mathbb{N}. \\ \quad \quad \quad \quad ((b \leq a) \wedge (b \leq c)) \Rightarrow  \\ \quad \quad \quad \quad \quad \sqrt{q(\operatorname{vadd}(d)(f(a), \operatorname{vneg}(d)(f(c))), \operatorname{vadd}(d)(f(a), \operatorname{vneg}(d)(f(c))))} \leq s \Rightarrow \\ \quad \quad \operatorname{is\text{-}cauchy\text{-}seq}(\operatorname{form\text{-}ms}(d, q), f) \end{array}}}
\vnbprop{form-cauchy-fun}{\vnbstmt{\begin{array}{@{}l@{}} \forall d, q, f. \\ \quad \operatorname{is\text{-}module}(d) \wedge \\ \quad \operatorname{scal}(d) = \operatorname{normed\text{-}field\text{-}as\text{-}commutative\text{-}ring}(\operatorname{cc\text{-}normed\text{-}field}) \wedge \\ \quad \operatorname{is\text{-}inner\text{-}product\text{-}form}(d, q) \wedge \\ \quad \operatorname{is\text{-}cauchy\text{-}seq}(\operatorname{form\text{-}ms}(d, q), f) \Rightarrow \\ \quad \quad f \in (\mathbb{N} \to \operatorname{vec}(d)) \end{array}}}
\vnbprop{form-cauchy-elim}{\vnbstmt{\begin{array}{@{}l@{}} \forall d, q, f, s. \\ \quad \operatorname{is\text{-}module}(d) \wedge \\ \quad \operatorname{scal}(d) = \operatorname{normed\text{-}field\text{-}as\text{-}commutative\text{-}ring}(\operatorname{cc\text{-}normed\text{-}field}) \wedge \\ \quad \operatorname{is\text{-}inner\text{-}product\text{-}form}(d, q) \wedge \\ \quad \operatorname{is\text{-}cauchy\text{-}seq}(\operatorname{form\text{-}ms}(d, q), f) \wedge \\ \quad \operatorname{pos\text{-}rr}(s) \Rightarrow \\ \quad \quad \exists b \in \mathbb{N}. \forall a, c \in \mathbb{N}. \\ \quad \quad \quad ((b \leq a) \wedge (b \leq c)) \Rightarrow  \\ \quad \quad \quad \quad \sqrt{q(\operatorname{vadd}(d)(f(a), \operatorname{vneg}(d)(f(c))), \operatorname{vadd}(d)(f(a), \operatorname{vneg}(d)(f(c))))} \leq s \end{array}}}
\vnbprop{form-cseq-unfold}{\vnbstmt{\forall d, q.\; (\operatorname{form\text{-}cseq}(d, q) \simeq \operatorname{cseq}(\operatorname{form\text{-}ms}(d, q)))}}
\vnbprop{form-cseq-in}{\vnbstmt{\begin{array}{@{}l@{}} \forall d, q, f. \\ \quad \operatorname{is\text{-}module}(d) \wedge \\ \quad \operatorname{scal}(d) = \operatorname{normed\text{-}field\text{-}as\text{-}commutative\text{-}ring}(\operatorname{cc\text{-}normed\text{-}field}) \wedge \\ \quad \operatorname{is\text{-}inner\text{-}product\text{-}form}(d, q) \wedge \\ \quad \operatorname{is\text{-}cauchy\text{-}seq}(\operatorname{form\text{-}ms}(d, q), f) \Rightarrow \\ \quad \quad f \in \operatorname{form\text{-}cseq}(d, q) \end{array}}}
\vnbprop{form-cseq-out}{\vnbstmt{\begin{array}{@{}l@{}} \forall d, q, f. \\ \quad \operatorname{is\text{-}module}(d) \wedge \\ \quad \operatorname{scal}(d) = \operatorname{normed\text{-}field\text{-}as\text{-}commutative\text{-}ring}(\operatorname{cc\text{-}normed\text{-}field}) \wedge \\ \quad \operatorname{is\text{-}inner\text{-}product\text{-}form}(d, q) \wedge \\ \quad f \in \operatorname{form\text{-}cseq}(d, q) \Rightarrow \\ \quad \quad \operatorname{is\text{-}cauchy\text{-}seq}(\operatorname{form\text{-}ms}(d, q), f) \end{array}}}
\vnbprop{seq-vadd-type}{\vnbstmt{\begin{array}{@{}l@{}} \forall d, q, f, g. \\ \quad \operatorname{is\text{-}module}(d) \wedge \\ \quad \operatorname{scal}(d) = \operatorname{normed\text{-}field\text{-}as\text{-}commutative\text{-}ring}(\operatorname{cc\text{-}normed\text{-}field}) \wedge \\ \quad \operatorname{is\text{-}inner\text{-}product\text{-}form}(d, q) \wedge \\ \quad f \in (\mathbb{N} \to \operatorname{vec}(d)) \wedge \\ \quad g \in (\mathbb{N} \to \operatorname{vec}(d)) \Rightarrow \\ \quad \quad \operatorname{seq\text{-}vadd}(d, f, g) \in (\mathbb{N} \to \operatorname{vec}(d)) \end{array}}}
\vnbprop{seq-vadd-at}{\vnbstmt{\begin{array}{@{}l@{}} \forall d, q, f, g, n. \\ \quad \operatorname{is\text{-}module}(d) \wedge \\ \quad \operatorname{scal}(d) = \operatorname{normed\text{-}field\text{-}as\text{-}commutative\text{-}ring}(\operatorname{cc\text{-}normed\text{-}field}) \wedge \\ \quad \operatorname{is\text{-}inner\text{-}product\text{-}form}(d, q) \wedge \\ \quad f \in (\mathbb{N} \to \operatorname{vec}(d)) \wedge \\ \quad g \in (\mathbb{N} \to \operatorname{vec}(d)) \wedge \\ \quad n \in \mathbb{N} \Rightarrow \\ \quad \quad \operatorname{seq\text{-}vadd}(d, f, g)(n) = \operatorname{vadd}(d)(f(n), g(n)) \end{array}}}
\vnbprop{seq-vneg-type}{\vnbstmt{\begin{array}{@{}l@{}} \forall d, q, f. \\ \quad \operatorname{is\text{-}module}(d) \wedge \\ \quad \operatorname{scal}(d) = \operatorname{normed\text{-}field\text{-}as\text{-}commutative\text{-}ring}(\operatorname{cc\text{-}normed\text{-}field}) \wedge \\ \quad \operatorname{is\text{-}inner\text{-}product\text{-}form}(d, q) \wedge \\ \quad f \in (\mathbb{N} \to \operatorname{vec}(d)) \Rightarrow \\ \quad \quad \operatorname{seq\text{-}vneg}(d, f) \in (\mathbb{N} \to \operatorname{vec}(d)) \end{array}}}
\vnbprop{seq-vneg-at}{\vnbstmt{\begin{array}{@{}l@{}} \forall d, q, f, n. \\ \quad \operatorname{is\text{-}module}(d) \wedge \\ \quad \operatorname{scal}(d) = \operatorname{normed\text{-}field\text{-}as\text{-}commutative\text{-}ring}(\operatorname{cc\text{-}normed\text{-}field}) \wedge \\ \quad \operatorname{is\text{-}inner\text{-}product\text{-}form}(d, q) \wedge \\ \quad f \in (\mathbb{N} \to \operatorname{vec}(d)) \wedge \\ \quad n \in \mathbb{N} \Rightarrow \\ \quad \quad \operatorname{seq\text{-}vneg}(d, f)(n) = \operatorname{vneg}(d)(f(n)) \end{array}}}
\vnbprop{seq-act-type}{\vnbstmt{\begin{array}{@{}l@{}} \forall d, q, a, f. \\ \quad \operatorname{is\text{-}module}(d) \wedge \\ \quad \operatorname{scal}(d) = \operatorname{normed\text{-}field\text{-}as\text{-}commutative\text{-}ring}(\operatorname{cc\text{-}normed\text{-}field}) \wedge \\ \quad \operatorname{is\text{-}inner\text{-}product\text{-}form}(d, q) \wedge \\ \quad a \in \mathbb{C} \wedge \\ \quad f \in (\mathbb{N} \to \operatorname{vec}(d)) \Rightarrow \\ \quad \quad \operatorname{seq\text{-}act}(d, a, f) \in (\mathbb{N} \to \operatorname{vec}(d)) \end{array}}}
\vnbprop{seq-act-at}{\vnbstmt{\begin{array}{@{}l@{}} \forall d, q, a, f, n. \\ \quad \operatorname{is\text{-}module}(d) \wedge \\ \quad \operatorname{scal}(d) = \operatorname{normed\text{-}field\text{-}as\text{-}commutative\text{-}ring}(\operatorname{cc\text{-}normed\text{-}field}) \wedge \\ \quad \operatorname{is\text{-}inner\text{-}product\text{-}form}(d, q) \wedge \\ \quad a \in \mathbb{C} \wedge \\ \quad f \in (\mathbb{N} \to \operatorname{vec}(d)) \wedge \\ \quad n \in \mathbb{N} \Rightarrow \\ \quad \quad \operatorname{seq\text{-}act}(d, a, f)(n) = \operatorname{act}(d)(a, f(n)) \end{array}}}
\vnbprop{seq-zero-type}{\vnbstmt{\begin{array}{@{}l@{}} \forall d, q. \\ \quad \operatorname{is\text{-}module}(d) \wedge \\ \quad \operatorname{scal}(d) = \operatorname{normed\text{-}field\text{-}as\text{-}commutative\text{-}ring}(\operatorname{cc\text{-}normed\text{-}field}) \wedge \\ \quad \operatorname{is\text{-}inner\text{-}product\text{-}form}(d, q) \Rightarrow \\ \quad \quad \operatorname{seq\text{-}zero}(d) \in (\mathbb{N} \to \operatorname{vec}(d)) \end{array}}}
\vnbprop{seq-zero-at}{\vnbstmt{\begin{array}{@{}l@{}} \forall d, q, n. \\ \quad \operatorname{is\text{-}module}(d) \wedge \\ \quad \operatorname{scal}(d) = \operatorname{normed\text{-}field\text{-}as\text{-}commutative\text{-}ring}(\operatorname{cc\text{-}normed\text{-}field}) \wedge \\ \quad \operatorname{is\text{-}inner\text{-}product\text{-}form}(d, q) \wedge \\ \quad n \in \mathbb{N} \Rightarrow \\ \quad \quad \operatorname{seq\text{-}zero}(d)(n) = \operatorname{vzero}(d) \end{array}}}
\vnbprop{seq-form-type}{\vnbstmt{\begin{array}{@{}l@{}} \forall d, q, f, g. \\ \quad \operatorname{is\text{-}module}(d) \wedge \\ \quad \operatorname{scal}(d) = \operatorname{normed\text{-}field\text{-}as\text{-}commutative\text{-}ring}(\operatorname{cc\text{-}normed\text{-}field}) \wedge \\ \quad \operatorname{is\text{-}inner\text{-}product\text{-}form}(d, q) \wedge \\ \quad f \in (\mathbb{N} \to \operatorname{vec}(d)) \wedge \\ \quad g \in (\mathbb{N} \to \operatorname{vec}(d)) \Rightarrow \\ \quad \quad \operatorname{seq\text{-}form}(q, f, g) \in (\mathbb{N} \to \mathbb{C}) \end{array}}}
\vnbprop{seq-form-at}{\vnbstmt{\begin{array}{@{}l@{}} \forall d, q, f, g, n. \\ \quad \operatorname{is\text{-}module}(d) \wedge \\ \quad \operatorname{scal}(d) = \operatorname{normed\text{-}field\text{-}as\text{-}commutative\text{-}ring}(\operatorname{cc\text{-}normed\text{-}field}) \wedge \\ \quad \operatorname{is\text{-}inner\text{-}product\text{-}form}(d, q) \wedge \\ \quad f \in (\mathbb{N} \to \operatorname{vec}(d)) \wedge \\ \quad g \in (\mathbb{N} \to \operatorname{vec}(d)) \wedge \\ \quad n \in \mathbb{N} \Rightarrow \\ \quad \quad \operatorname{seq\text{-}form}(q, f, g)(n) = q(f(n), g(n)) \end{array}}}
\vnbprop{seq-vadd-cauchy}{\vnbstmt{\begin{array}{@{}l@{}} \forall d, q, f, g. \\ \quad \operatorname{is\text{-}module}(d) \wedge \\ \quad \operatorname{scal}(d) = \operatorname{normed\text{-}field\text{-}as\text{-}commutative\text{-}ring}(\operatorname{cc\text{-}normed\text{-}field}) \wedge \\ \quad \operatorname{is\text{-}inner\text{-}product\text{-}form}(d, q) \wedge \\ \quad \operatorname{is\text{-}cauchy\text{-}seq}(\operatorname{form\text{-}ms}(d, q), f) \wedge \\ \quad \operatorname{is\text{-}cauchy\text{-}seq}(\operatorname{form\text{-}ms}(d, q), g) \Rightarrow \\ \quad \quad \operatorname{is\text{-}cauchy\text{-}seq}(\operatorname{form\text{-}ms}(d, q), \operatorname{seq\text{-}vadd}(d, f, g)) \end{array}}}
\vnbprop{seq-vneg-cauchy}{\vnbstmt{\begin{array}{@{}l@{}} \forall d, q, f. \\ \quad \operatorname{is\text{-}module}(d) \wedge \\ \quad \operatorname{scal}(d) = \operatorname{normed\text{-}field\text{-}as\text{-}commutative\text{-}ring}(\operatorname{cc\text{-}normed\text{-}field}) \wedge \\ \quad \operatorname{is\text{-}inner\text{-}product\text{-}form}(d, q) \wedge \\ \quad \operatorname{is\text{-}cauchy\text{-}seq}(\operatorname{form\text{-}ms}(d, q), f) \Rightarrow \\ \quad \quad \operatorname{is\text{-}cauchy\text{-}seq}(\operatorname{form\text{-}ms}(d, q), \operatorname{seq\text{-}vneg}(d, f)) \end{array}}}
\vnbprop{seq-act-cauchy}{\vnbstmt{\begin{array}{@{}l@{}} \forall d, q, a, f. \\ \quad \operatorname{is\text{-}module}(d) \wedge \\ \quad \operatorname{scal}(d) = \operatorname{normed\text{-}field\text{-}as\text{-}commutative\text{-}ring}(\operatorname{cc\text{-}normed\text{-}field}) \wedge \\ \quad \operatorname{is\text{-}inner\text{-}product\text{-}form}(d, q) \wedge \\ \quad a \in \mathbb{C} \wedge \\ \quad \operatorname{is\text{-}cauchy\text{-}seq}(\operatorname{form\text{-}ms}(d, q), f) \Rightarrow \\ \quad \quad \operatorname{is\text{-}cauchy\text{-}seq}(\operatorname{form\text{-}ms}(d, q), \operatorname{seq\text{-}act}(d, a, f)) \end{array}}}
\vnbprop{seq-zero-embed}{\vnbstmt{\begin{array}{@{}l@{}} \forall d, q. \\ \quad \operatorname{is\text{-}module}(d) \wedge \\ \quad \operatorname{scal}(d) = \operatorname{normed\text{-}field\text{-}as\text{-}commutative\text{-}ring}(\operatorname{cc\text{-}normed\text{-}field}) \wedge \\ \quad \operatorname{is\text{-}inner\text{-}product\text{-}form}(d, q) \Rightarrow \\ \quad \quad (\operatorname{seq\text{-}zero}(d) \simeq \operatorname{embed\text{-}seq}(\operatorname{form\text{-}ms}(d, q), \operatorname{vzero}(d))) \end{array}}}
\vnbprop{seq-zero-cauchy}{\vnbstmt{\begin{array}{@{}l@{}} \forall d, q. \\ \quad \operatorname{is\text{-}module}(d) \wedge \\ \quad \operatorname{scal}(d) = \operatorname{normed\text{-}field\text{-}as\text{-}commutative\text{-}ring}(\operatorname{cc\text{-}normed\text{-}field}) \wedge \\ \quad \operatorname{is\text{-}inner\text{-}product\text{-}form}(d, q) \Rightarrow \\ \quad \quad \operatorname{is\text{-}cauchy\text{-}seq}(\operatorname{form\text{-}ms}(d, q), \operatorname{seq\text{-}zero}(d)) \end{array}}}
\vnbprop{cseq-module-scal}{\vnbstmt{\begin{array}{@{}l@{}} \forall d, q. \\ \quad \operatorname{is\text{-}module}(d) \wedge \\ \quad \operatorname{scal}(d) = \operatorname{normed\text{-}field\text{-}as\text{-}commutative\text{-}ring}(\operatorname{cc\text{-}normed\text{-}field}) \wedge \\ \quad \operatorname{is\text{-}inner\text{-}product\text{-}form}(d, q) \Rightarrow \\ \quad \quad (\operatorname{scal}(\operatorname{cseq\text{-}module}(d, q)) \simeq \operatorname{scal}(d)) \end{array}}}
\vnbprop{cseq-module-vec}{\vnbstmt{\begin{array}{@{}l@{}} \forall d, q. \\ \quad \operatorname{is\text{-}module}(d) \wedge \\ \quad \operatorname{scal}(d) = \operatorname{normed\text{-}field\text{-}as\text{-}commutative\text{-}ring}(\operatorname{cc\text{-}normed\text{-}field}) \wedge \\ \quad \operatorname{is\text{-}inner\text{-}product\text{-}form}(d, q) \Rightarrow \\ \quad \quad (\operatorname{vec}(\operatorname{cseq\text{-}module}(d, q)) \simeq \operatorname{form\text{-}cseq}(d, q)) \end{array}}}
\vnbprop{cseq-module-vadd}{\vnbstmt{\begin{array}{@{}l@{}} \forall d, q. \\ \quad \operatorname{is\text{-}module}(d) \wedge \\ \quad \operatorname{scal}(d) = \operatorname{normed\text{-}field\text{-}as\text{-}commutative\text{-}ring}(\operatorname{cc\text{-}normed\text{-}field}) \wedge \\ \quad \operatorname{is\text{-}inner\text{-}product\text{-}form}(d, q) \Rightarrow \\ \quad \quad (\operatorname{vadd}(\operatorname{cseq\text{-}module}(d, q)) \simeq (\lambda {[}\operatorname{hcf\_}, \operatorname{hcg\_}] \in (\operatorname{form\text{-}cseq}(d, q) \times \operatorname{form\text{-}cseq}(d, q)).\; \operatorname{seq\text{-}vadd}(d, \operatorname{hcf\_}, \operatorname{hcg\_}))) \end{array}}}
\vnbprop{cseq-module-vzero}{\vnbstmt{\begin{array}{@{}l@{}} \forall d, q. \\ \quad \operatorname{is\text{-}module}(d) \wedge \\ \quad \operatorname{scal}(d) = \operatorname{normed\text{-}field\text{-}as\text{-}commutative\text{-}ring}(\operatorname{cc\text{-}normed\text{-}field}) \wedge \\ \quad \operatorname{is\text{-}inner\text{-}product\text{-}form}(d, q) \Rightarrow \\ \quad \quad (\operatorname{vzero}(\operatorname{cseq\text{-}module}(d, q)) \simeq \operatorname{seq\text{-}zero}(d)) \end{array}}}
\vnbprop{cseq-module-vneg}{\vnbstmt{\begin{array}{@{}l@{}} \forall d, q. \\ \quad \operatorname{is\text{-}module}(d) \wedge \\ \quad \operatorname{scal}(d) = \operatorname{normed\text{-}field\text{-}as\text{-}commutative\text{-}ring}(\operatorname{cc\text{-}normed\text{-}field}) \wedge \\ \quad \operatorname{is\text{-}inner\text{-}product\text{-}form}(d, q) \Rightarrow \\ \quad \quad (\operatorname{vneg}(\operatorname{cseq\text{-}module}(d, q)) \simeq (\lambda f \in \operatorname{form\text{-}cseq}(d, q).\; \operatorname{seq\text{-}vneg}(d, f))) \end{array}}}
\vnbprop{cseq-module-act}{\vnbstmt{\begin{array}{@{}l@{}} \forall d, q. \\ \quad \operatorname{is\text{-}module}(d) \wedge \\ \quad \operatorname{scal}(d) = \operatorname{normed\text{-}field\text{-}as\text{-}commutative\text{-}ring}(\operatorname{cc\text{-}normed\text{-}field}) \wedge \\ \quad \operatorname{is\text{-}inner\text{-}product\text{-}form}(d, q) \Rightarrow \\ \quad \quad (\operatorname{act}(\operatorname{cseq\text{-}module}(d, q)) \simeq (\lambda {[}\operatorname{hca\_}, \operatorname{hcf\_}] \in (\mathbb{C} \times \operatorname{form\text{-}cseq}(d, q)).\; \operatorname{seq\text{-}act}(d, \operatorname{hca\_}, \operatorname{hcf\_}))) \end{array}}}
\vnbprop{form-cseq-fun}{\vnbstmt{\begin{array}{@{}l@{}} \forall d, q, f. \\ \quad \operatorname{is\text{-}module}(d) \wedge \\ \quad \operatorname{scal}(d) = \operatorname{normed\text{-}field\text{-}as\text{-}commutative\text{-}ring}(\operatorname{cc\text{-}normed\text{-}field}) \wedge \\ \quad \operatorname{is\text{-}inner\text{-}product\text{-}form}(d, q) \wedge \\ \quad f \in \operatorname{form\text{-}cseq}(d, q) \Rightarrow \\ \quad \quad f \in (\mathbb{N} \to \operatorname{vec}(d)) \end{array}}}
\vnbprop{cseq-vadd-closed}{\vnbstmt{\begin{array}{@{}l@{}} \forall d, q, f, g. \\ \quad \operatorname{is\text{-}module}(d) \wedge \\ \quad \operatorname{scal}(d) = \operatorname{normed\text{-}field\text{-}as\text{-}commutative\text{-}ring}(\operatorname{cc\text{-}normed\text{-}field}) \wedge \\ \quad \operatorname{is\text{-}inner\text{-}product\text{-}form}(d, q) \wedge \\ \quad f \in \operatorname{form\text{-}cseq}(d, q) \wedge \\ \quad g \in \operatorname{form\text{-}cseq}(d, q) \Rightarrow \\ \quad \quad \operatorname{seq\text{-}vadd}(d, f, g) \in \operatorname{form\text{-}cseq}(d, q) \end{array}}}
\vnbprop{cseq-vneg-closed}{\vnbstmt{\begin{array}{@{}l@{}} \forall d, q, f. \\ \quad \operatorname{is\text{-}module}(d) \wedge \\ \quad \operatorname{scal}(d) = \operatorname{normed\text{-}field\text{-}as\text{-}commutative\text{-}ring}(\operatorname{cc\text{-}normed\text{-}field}) \wedge \\ \quad \operatorname{is\text{-}inner\text{-}product\text{-}form}(d, q) \wedge \\ \quad f \in \operatorname{form\text{-}cseq}(d, q) \Rightarrow \\ \quad \quad \operatorname{seq\text{-}vneg}(d, f) \in \operatorname{form\text{-}cseq}(d, q) \end{array}}}
\vnbprop{cseq-act-closed}{\vnbstmt{\begin{array}{@{}l@{}} \forall d, q, a, f. \\ \quad \operatorname{is\text{-}module}(d) \wedge \\ \quad \operatorname{scal}(d) = \operatorname{normed\text{-}field\text{-}as\text{-}commutative\text{-}ring}(\operatorname{cc\text{-}normed\text{-}field}) \wedge \\ \quad \operatorname{is\text{-}inner\text{-}product\text{-}form}(d, q) \wedge \\ \quad a \in \mathbb{C} \wedge \\ \quad f \in \operatorname{form\text{-}cseq}(d, q) \Rightarrow \\ \quad \quad \operatorname{seq\text{-}act}(d, a, f) \in \operatorname{form\text{-}cseq}(d, q) \end{array}}}
\vnbprop{cseq-zero-closed}{\vnbstmt{\begin{array}{@{}l@{}} \forall d, q. \\ \quad \operatorname{is\text{-}module}(d) \wedge \\ \quad \operatorname{scal}(d) = \operatorname{normed\text{-}field\text{-}as\text{-}commutative\text{-}ring}(\operatorname{cc\text{-}normed\text{-}field}) \wedge \\ \quad \operatorname{is\text{-}inner\text{-}product\text{-}form}(d, q) \Rightarrow \\ \quad \quad \operatorname{seq\text{-}zero}(d) \in \operatorname{form\text{-}cseq}(d, q) \end{array}}}
\vnbprop{cseq-module-vadd-computes}{\vnbstmt{\begin{array}{@{}l@{}} \forall d, q, f, g. \\ \quad \operatorname{is\text{-}module}(d) \wedge \\ \quad \operatorname{scal}(d) = \operatorname{normed\text{-}field\text{-}as\text{-}commutative\text{-}ring}(\operatorname{cc\text{-}normed\text{-}field}) \wedge \\ \quad \operatorname{is\text{-}inner\text{-}product\text{-}form}(d, q) \wedge \\ \quad f \in \operatorname{form\text{-}cseq}(d, q) \wedge \\ \quad g \in \operatorname{form\text{-}cseq}(d, q) \Rightarrow \\ \quad \quad \operatorname{vadd}(\operatorname{cseq\text{-}module}(d, q))(f, g) = \operatorname{seq\text{-}vadd}(d, f, g) \end{array}}}
\vnbprop{cseq-module-vneg-computes}{\vnbstmt{\begin{array}{@{}l@{}} \forall d, q, f. \\ \quad \operatorname{is\text{-}module}(d) \wedge \\ \quad \operatorname{scal}(d) = \operatorname{normed\text{-}field\text{-}as\text{-}commutative\text{-}ring}(\operatorname{cc\text{-}normed\text{-}field}) \wedge \\ \quad \operatorname{is\text{-}inner\text{-}product\text{-}form}(d, q) \wedge \\ \quad f \in \operatorname{form\text{-}cseq}(d, q) \Rightarrow \\ \quad \quad \operatorname{vneg}(\operatorname{cseq\text{-}module}(d, q))(f) = \operatorname{seq\text{-}vneg}(d, f) \end{array}}}
\vnbprop{cseq-module-act-computes}{\vnbstmt{\begin{array}{@{}l@{}} \forall d, q, a, f. \\ \quad \operatorname{is\text{-}module}(d) \wedge \\ \quad \operatorname{scal}(d) = \operatorname{normed\text{-}field\text{-}as\text{-}commutative\text{-}ring}(\operatorname{cc\text{-}normed\text{-}field}) \wedge \\ \quad \operatorname{is\text{-}inner\text{-}product\text{-}form}(d, q) \wedge \\ \quad a \in \mathbb{C} \wedge \\ \quad f \in \operatorname{form\text{-}cseq}(d, q) \Rightarrow \\ \quad \quad \operatorname{act}(\operatorname{cseq\text{-}module}(d, q))(a, f) = \operatorname{seq\text{-}act}(d, a, f) \end{array}}}
\vnbprop{cseq-module-is-module}{\vnbstmt{\begin{array}{@{}l@{}} \forall d, q. \\ \quad \operatorname{is\text{-}module}(d) \wedge \\ \quad \operatorname{scal}(d) = \operatorname{normed\text{-}field\text{-}as\text{-}commutative\text{-}ring}(\operatorname{cc\text{-}normed\text{-}field}) \wedge \\ \quad \operatorname{is\text{-}inner\text{-}product\text{-}form}(d, q) \Rightarrow \\ \quad \quad \operatorname{is\text{-}module}(\operatorname{cseq\text{-}module}(d, q)) \end{array}}}
\vnbprop{seq-form-cauchy}{\vnbstmt{\begin{array}{@{}l@{}} \forall d, q, f, g. \\ \quad \operatorname{is\text{-}module}(d) \wedge \\ \quad \operatorname{scal}(d) = \operatorname{normed\text{-}field\text{-}as\text{-}commutative\text{-}ring}(\operatorname{cc\text{-}normed\text{-}field}) \wedge \\ \quad \operatorname{is\text{-}inner\text{-}product\text{-}form}(d, q) \wedge \\ \quad f \in \operatorname{form\text{-}cseq}(d, q) \wedge \\ \quad g \in \operatorname{form\text{-}cseq}(d, q) \Rightarrow \\ \quad \quad \operatorname{is\text{-}cauchy\text{-}seq}(\operatorname{cc\text{-}ms}, \operatorname{seq\text{-}form}(q, f, g)) \end{array}}}
\vnbprop{limit-form-iota}{\vnbstmt{\begin{array}{@{}l@{}} \forall d, q, f, g. \\ \quad \operatorname{is\text{-}module}(d) \wedge \\ \quad \operatorname{scal}(d) = \operatorname{normed\text{-}field\text{-}as\text{-}commutative\text{-}ring}(\operatorname{cc\text{-}normed\text{-}field}) \wedge \\ \quad \operatorname{is\text{-}inner\text{-}product\text{-}form}(d, q) \wedge \\ \quad f \in \operatorname{form\text{-}cseq}(d, q) \wedge \\ \quad g \in \operatorname{form\text{-}cseq}(d, q) \Rightarrow \\ \quad \quad (\iota\, c.\;((c \in \mathbb{C}) \wedge \operatorname{converges\text{-}to}(\operatorname{cc\text{-}ms}, \operatorname{seq\text{-}form}(q, f, g), c))) \in \mathbb{C} \wedge \\ \quad \quad \operatorname{converges\text{-}to}(\operatorname{cc\text{-}ms}, \operatorname{seq\text{-}form}(q, f, g), \\ \quad \quad \quad (\iota\, c.\;((c \in \mathbb{C}) \wedge \operatorname{converges\text{-}to}(\operatorname{cc\text{-}ms}, \operatorname{seq\text{-}form}(q, f, g), c)))) \end{array}}}
\vnbprop{limit-form-prop}{\vnbstmt{\begin{array}{@{}l@{}} \forall d, q, f, g. \\ \quad \operatorname{is\text{-}module}(d) \wedge \\ \quad \operatorname{scal}(d) = \operatorname{normed\text{-}field\text{-}as\text{-}commutative\text{-}ring}(\operatorname{cc\text{-}normed\text{-}field}) \wedge \\ \quad \operatorname{is\text{-}inner\text{-}product\text{-}form}(d, q) \wedge \\ \quad f \in \operatorname{form\text{-}cseq}(d, q) \wedge \\ \quad g \in \operatorname{form\text{-}cseq}(d, q) \Rightarrow \\ \quad \quad ((\operatorname{limit\text{-}form}(d, q)(f, g) \in \mathbb{C}) \wedge \operatorname{converges\text{-}to}(\operatorname{cc\text{-}ms}, \operatorname{seq\text{-}form}(q, f, g), \operatorname{limit\text{-}form}(d, q)(f, g))) \end{array}}}
\vnbprop{limit-form-value}{\vnbstmt{\begin{array}{@{}l@{}} \forall d, q, f, g, c. \\ \quad \operatorname{is\text{-}module}(d) \wedge \\ \quad \operatorname{scal}(d) = \operatorname{normed\text{-}field\text{-}as\text{-}commutative\text{-}ring}(\operatorname{cc\text{-}normed\text{-}field}) \wedge \\ \quad \operatorname{is\text{-}inner\text{-}product\text{-}form}(d, q) \wedge \\ \quad f \in \operatorname{form\text{-}cseq}(d, q) \wedge \\ \quad g \in \operatorname{form\text{-}cseq}(d, q) \wedge \\ \quad c \in \mathbb{C} \wedge \\ \quad \operatorname{converges\text{-}to}(\operatorname{cc\text{-}ms}, \operatorname{seq\text{-}form}(q, f, g), c) \Rightarrow \\ \quad \quad \operatorname{limit\text{-}form}(d, q)(f, g) = c \end{array}}}
\vnbprop{limit-form-type}{\vnbstmt{\begin{array}{@{}l@{}} \forall d, q. \\ \quad \operatorname{is\text{-}module}(d) \wedge \\ \quad \operatorname{scal}(d) = \operatorname{normed\text{-}field\text{-}as\text{-}commutative\text{-}ring}(\operatorname{cc\text{-}normed\text{-}field}) \wedge \\ \quad \operatorname{is\text{-}inner\text{-}product\text{-}form}(d, q) \Rightarrow \\ \quad \quad \operatorname{limit\text{-}form}(d, q) \in ((\operatorname{form\text{-}cseq}(d, q) \times \operatorname{form\text{-}cseq}(d, q)) \to \mathbb{C}) \end{array}}}
\vnbprop{limit-form-add-left}{\vnbstmt{\begin{array}{@{}l@{}} \forall d, q, f, g, h. \\ \quad \operatorname{is\text{-}module}(d) \wedge \\ \quad \operatorname{scal}(d) = \operatorname{normed\text{-}field\text{-}as\text{-}commutative\text{-}ring}(\operatorname{cc\text{-}normed\text{-}field}) \wedge \\ \quad \operatorname{is\text{-}inner\text{-}product\text{-}form}(d, q) \wedge \\ \quad f \in \operatorname{form\text{-}cseq}(d, q) \wedge \\ \quad g \in \operatorname{form\text{-}cseq}(d, q) \wedge \\ \quad h \in \operatorname{form\text{-}cseq}(d, q) \Rightarrow \\ \quad \quad \operatorname{limit\text{-}form}(d, q)(\operatorname{seq\text{-}vadd}(d, f, g), h) = (\operatorname{limit\text{-}form}(d, q)(f, h) + \operatorname{limit\text{-}form}(d, q)(g, h)) \end{array}}}
\vnbprop{limit-form-homog-left}{\vnbstmt{\begin{array}{@{}l@{}} \forall d, q, a, f, h. \\ \quad \operatorname{is\text{-}module}(d) \wedge \\ \quad \operatorname{scal}(d) = \operatorname{normed\text{-}field\text{-}as\text{-}commutative\text{-}ring}(\operatorname{cc\text{-}normed\text{-}field}) \wedge \\ \quad \operatorname{is\text{-}inner\text{-}product\text{-}form}(d, q) \wedge \\ \quad a \in \mathbb{C} \wedge \\ \quad f \in \operatorname{form\text{-}cseq}(d, q) \wedge \\ \quad h \in \operatorname{form\text{-}cseq}(d, q) \Rightarrow \\ \quad \quad \operatorname{limit\text{-}form}(d, q)(\operatorname{seq\text{-}act}(d, a, f), h) = (a \cdot \operatorname{limit\text{-}form}(d, q)(f, h)) \end{array}}}
\vnbprop{cc-limit-conjugate}{\vnbstmt{\begin{array}{@{}l@{}} \forall s, h \in (\mathbb{N} \to \mathbb{C}), l \in \mathbb{C}. \\ \quad (\forall j \in \mathbb{N}.\; (h(j) = \operatorname{conjugate}(s(j))) \wedge \operatorname{converges\text{-}to}(\operatorname{cc\text{-}ms}, s, l)) \Rightarrow  \\ \quad \quad \operatorname{converges\text{-}to}(\operatorname{cc\text{-}ms}, h, \operatorname{conjugate}(l)) \end{array}}}
\vnbprop{limit-form-conj-sym}{\vnbstmt{\begin{array}{@{}l@{}} \forall d, q, f, g. \\ \quad \operatorname{is\text{-}module}(d) \wedge \\ \quad \operatorname{scal}(d) = \operatorname{normed\text{-}field\text{-}as\text{-}commutative\text{-}ring}(\operatorname{cc\text{-}normed\text{-}field}) \wedge \\ \quad \operatorname{is\text{-}inner\text{-}product\text{-}form}(d, q) \wedge \\ \quad f \in \operatorname{form\text{-}cseq}(d, q) \wedge \\ \quad g \in \operatorname{form\text{-}cseq}(d, q) \Rightarrow \\ \quad \quad \operatorname{limit\text{-}form}(d, q)(g, f) = \operatorname{conjugate}(\operatorname{limit\text{-}form}(d, q)(f, g)) \end{array}}}
\vnbprop{limit-form-diag}{\vnbstmt{\begin{array}{@{}l@{}} \forall d, q, f. \\ \quad \operatorname{is\text{-}module}(d) \wedge \\ \quad \operatorname{scal}(d) = \operatorname{normed\text{-}field\text{-}as\text{-}commutative\text{-}ring}(\operatorname{cc\text{-}normed\text{-}field}) \wedge \\ \quad \operatorname{is\text{-}inner\text{-}product\text{-}form}(d, q) \wedge \\ \quad f \in \operatorname{form\text{-}cseq}(d, q) \Rightarrow \\ \quad \quad (\operatorname{limit\text{-}form}(d, q)(f, f) \in \mathbb{R}) \wedge (0 \leq \operatorname{limit\text{-}form}(d, q)(f, f)) \end{array}}}
\vnbprop{limit-form-psd}{\vnbstmt{\begin{array}{@{}l@{}} \forall d, q. \\ \quad \operatorname{is\text{-}module}(d) \wedge \\ \quad \operatorname{scal}(d) = \operatorname{normed\text{-}field\text{-}as\text{-}commutative\text{-}ring}(\operatorname{cc\text{-}normed\text{-}field}) \wedge \\ \quad \operatorname{is\text{-}inner\text{-}product\text{-}form}(d, q) \Rightarrow \\ \quad \quad \operatorname{is\text{-}psd\text{-}form}(\operatorname{cseq\text{-}module}(d, q), \operatorname{limit\text{-}form}(d, q)) \end{array}}}
\vnbprop{hilbert-completion-is-cips}{\vnbstmt{\begin{array}{@{}l@{}} \forall d, q. \\ \quad \operatorname{is\text{-}module}(d) \wedge \\ \quad \operatorname{scal}(d) = \operatorname{normed\text{-}field\text{-}as\text{-}commutative\text{-}ring}(\operatorname{cc\text{-}normed\text{-}field}) \wedge \\ \quad \operatorname{is\text{-}inner\text{-}product\text{-}form}(d, q) \Rightarrow \\ \quad \quad \operatorname{is\text{-}complex\text{-}inner\text{-}product\text{-}space}(\operatorname{hilbert\text{-}completion}(d, q)) \end{array}}}
\vnbprop{form-quotient-ip}{\vnbstmt{\begin{array}{@{}l@{}} \forall d, q. \\ \quad (\operatorname{is\text{-}module}(d) \wedge ((\operatorname{scal}(d) = \operatorname{normed\text{-}field\text{-}as\text{-}commutative\text{-}ring}(\operatorname{cc\text{-}normed\text{-}field})) \wedge \operatorname{is\text{-}psd\text{-}form}(d, q))) \Rightarrow  \\ \quad \quad (\langle\cdot,\cdot\rangle_{{d}/N_{q}} \simeq \operatorname{form\text{-}qip}(d, q)) \end{array}}}
\vnbprop{limit-form-embed}{\vnbstmt{\begin{array}{@{}l@{}} \forall d, q, x, y. \\ \quad \operatorname{is\text{-}module}(d) \wedge \\ \quad \operatorname{scal}(d) = \operatorname{normed\text{-}field\text{-}as\text{-}commutative\text{-}ring}(\operatorname{cc\text{-}normed\text{-}field}) \wedge \\ \quad \operatorname{is\text{-}inner\text{-}product\text{-}form}(d, q) \wedge \\ \quad x \in \operatorname{vec}(d) \wedge \\ \quad y \in \operatorname{vec}(d) \Rightarrow \\ \quad \quad \operatorname{limit\text{-}form}(d, q)(\operatorname{embed\text{-}seq}(\operatorname{form\text{-}ms}(d, q), x), \operatorname{embed\text{-}seq}(\operatorname{form\text{-}ms}(d, q), y)) = q(x, y) \end{array}}}
\vnbprop{hc-embed-ip}{\vnbstmt{\begin{array}{@{}l@{}} \forall d, q, x, y. \\ \quad \operatorname{is\text{-}module}(d) \wedge \\ \quad \operatorname{scal}(d) = \operatorname{normed\text{-}field\text{-}as\text{-}commutative\text{-}ring}(\operatorname{cc\text{-}normed\text{-}field}) \wedge \\ \quad \operatorname{is\text{-}inner\text{-}product\text{-}form}(d, q) \wedge \\ \quad x \in \operatorname{vec}(d) \wedge \\ \quad y \in \operatorname{vec}(d) \Rightarrow \\ \quad \quad \langle\cdot,\cdot\rangle_{\operatorname{hilbert\text{-}completion}(d, q)}(\operatorname{hc\text{-}embed}(d, q)(x), \operatorname{hc\text{-}embed}(d, q)(y)) = q(x, y) \end{array}}}
\vnbprop{form-quotient-vadd}{\vnbstmt{\begin{array}{@{}l@{}} \forall d, q. \\ \quad (\operatorname{is\text{-}module}(d) \wedge ((\operatorname{scal}(d) = \operatorname{normed\text{-}field\text{-}as\text{-}commutative\text{-}ring}(\operatorname{cc\text{-}normed\text{-}field})) \wedge \operatorname{is\text{-}psd\text{-}form}(d, q))) \Rightarrow  \\ \quad \quad (\operatorname{vadd}({d}/N_{q}) \simeq \operatorname{vadd}(\operatorname{quotient\text{-}module}(d, N_{q}))) \end{array}}}
\vnbprop{form-quotient-act}{\vnbstmt{\begin{array}{@{}l@{}} \forall d, q. \\ \quad (\operatorname{is\text{-}module}(d) \wedge ((\operatorname{scal}(d) = \operatorname{normed\text{-}field\text{-}as\text{-}commutative\text{-}ring}(\operatorname{cc\text{-}normed\text{-}field})) \wedge \operatorname{is\text{-}psd\text{-}form}(d, q))) \Rightarrow  \\ \quad \quad (\operatorname{act}({d}/N_{q}) \simeq \operatorname{act}(\operatorname{quotient\text{-}module}(d, N_{q}))) \end{array}}}
\vnbprop{form-quotient-vec}{\vnbstmt{\begin{array}{@{}l@{}} \forall d, q. \\ \quad (\operatorname{is\text{-}module}(d) \wedge ((\operatorname{scal}(d) = \operatorname{normed\text{-}field\text{-}as\text{-}commutative\text{-}ring}(\operatorname{cc\text{-}normed\text{-}field})) \wedge \operatorname{is\text{-}psd\text{-}form}(d, q))) \Rightarrow  \\ \quad \quad (\operatorname{vec}({d}/N_{q}) \simeq \operatorname{vec}(\operatorname{quotient\text{-}module}(d, N_{q}))) \end{array}}}
\vnbprop{form-quotient-vneg}{\vnbstmt{\begin{array}{@{}l@{}} \forall d, q. \\ \quad (\operatorname{is\text{-}module}(d) \wedge ((\operatorname{scal}(d) = \operatorname{normed\text{-}field\text{-}as\text{-}commutative\text{-}ring}(\operatorname{cc\text{-}normed\text{-}field})) \wedge \operatorname{is\text{-}psd\text{-}form}(d, q))) \Rightarrow  \\ \quad \quad (\operatorname{vneg}({d}/N_{q}) \simeq \operatorname{vneg}(\operatorname{quotient\text{-}module}(d, N_{q}))) \end{array}}}
\vnbprop{hc-embed-vadd}{\vnbstmt{\begin{array}{@{}l@{}} \forall d, q, x, y. \\ \quad \operatorname{is\text{-}module}(d) \wedge \\ \quad \operatorname{scal}(d) = \operatorname{normed\text{-}field\text{-}as\text{-}commutative\text{-}ring}(\operatorname{cc\text{-}normed\text{-}field}) \wedge \\ \quad \operatorname{is\text{-}inner\text{-}product\text{-}form}(d, q) \wedge \\ \quad x \in \operatorname{vec}(d) \wedge \\ \quad y \in \operatorname{vec}(d) \Rightarrow \\ \quad \quad \operatorname{hc\text{-}embed}(d, q)(\operatorname{vadd}(d)(x, y)) \\ \quad \quad =\; \operatorname{vadd}(\operatorname{hilbert\text{-}completion}(d, q))(\operatorname{hc\text{-}embed}(d, q)(x), \operatorname{hc\text{-}embed}(d, q)(y)) \end{array}}}
\vnbprop{hc-embed-act}{\vnbstmt{\begin{array}{@{}l@{}} \forall d, q, a, x. \\ \quad \operatorname{is\text{-}module}(d) \wedge \\ \quad \operatorname{scal}(d) = \operatorname{normed\text{-}field\text{-}as\text{-}commutative\text{-}ring}(\operatorname{cc\text{-}normed\text{-}field}) \wedge \\ \quad \operatorname{is\text{-}inner\text{-}product\text{-}form}(d, q) \wedge \\ \quad a \in \mathbb{C} \wedge \\ \quad x \in \operatorname{vec}(d) \Rightarrow \\ \quad \quad \operatorname{hc\text{-}embed}(d, q)(\operatorname{act}(d)(a, x)) = \operatorname{act}(\operatorname{hilbert\text{-}completion}(d, q))(a, \operatorname{hc\text{-}embed}(d, q)(x)) \end{array}}}
\vnbprop{sqrt-diff-le}{\vnbstmt{\begin{array}{@{}l@{}} \forall x, y \in \mathbb{R}. \\ \quad ((0 \leq x) \wedge (0 \leq y)) \Rightarrow (\lvert (\sqrt{x} - \sqrt{y}) \rvert \leq \sqrt{\lvert (x - y) \rvert}) \end{array}}}
\vnbprop{rr-limit-sqrt}{\vnbstmt{\begin{array}{@{}l@{}} \forall a, b \in (\mathbb{N} \to \mathbb{R}), l \in \mathbb{R}. \\ \quad \forall j \in \mathbb{N}.\; (0 \leq a(j)) \wedge \\ \quad 0 \leq l \wedge \\ \quad \forall j \in \mathbb{N}.\; (b(j) = \sqrt{a(j)}) \wedge \\ \quad \operatorname{converges\text{-}to}(\operatorname{rr\text{-}ms}, a, l) \Rightarrow \\ \quad \quad \operatorname{converges\text{-}to}(\operatorname{rr\text{-}ms}, b, \sqrt{l}) \end{array}}}
\vnbprop{limit-form-rr}{\vnbstmt{\begin{array}{@{}l@{}} \forall d, q, f. \\ \quad \operatorname{is\text{-}module}(d) \wedge \\ \quad \operatorname{scal}(d) = \operatorname{normed\text{-}field\text{-}as\text{-}commutative\text{-}ring}(\operatorname{cc\text{-}normed\text{-}field}) \wedge \\ \quad \operatorname{is\text{-}inner\text{-}product\text{-}form}(d, q) \wedge \\ \quad f \in \operatorname{form\text{-}cseq}(d, q) \Rightarrow \\ \quad \quad \operatorname{converges\text{-}to}(\operatorname{rr\text{-}ms}, \operatorname{seq\text{-}form}(q, f, f), \operatorname{limit\text{-}form}(d, q)(f, f)) \end{array}}}
\vnbprop{dist-seq-limit}{\vnbstmt{\begin{array}{@{}l@{}} \forall d, q, f, g. \\ \quad \operatorname{is\text{-}module}(d) \wedge \\ \quad \operatorname{scal}(d) = \operatorname{normed\text{-}field\text{-}as\text{-}commutative\text{-}ring}(\operatorname{cc\text{-}normed\text{-}field}) \wedge \\ \quad \operatorname{is\text{-}inner\text{-}product\text{-}form}(d, q) \wedge \\ \quad f \in \operatorname{form\text{-}cseq}(d, q) \wedge \\ \quad g \in \operatorname{form\text{-}cseq}(d, q) \Rightarrow \\ \quad \quad \operatorname{converges\text{-}to}(\operatorname{rr\text{-}ms}, \operatorname{dist\text{-}seq}(\operatorname{form\text{-}ms}(d, q), f, g), \\ \quad \quad \quad \sqrt{\operatorname{limit\text{-}form}(d, q)(\operatorname{seq\text{-}vadd}(d, f, \operatorname{seq\text{-}vneg}(d, g)), \operatorname{seq\text{-}vadd}(d, f, \operatorname{seq\text{-}vneg}(d, g)))}) \end{array}}}
\vnbprop{hc-rel-fwd}{\vnbstmt{\begin{array}{@{}l@{}} \forall d, q, p, a. \\ \quad \operatorname{is\text{-}module}(d) \wedge \\ \quad \operatorname{scal}(d) = \operatorname{normed\text{-}field\text{-}as\text{-}commutative\text{-}ring}(\operatorname{cc\text{-}normed\text{-}field}) \wedge \\ \quad \operatorname{is\text{-}inner\text{-}product\text{-}form}(d, q) \wedge \\ \quad p \in \operatorname{form\text{-}cseq}(d, q) \wedge \\ \quad a \in \operatorname{form\text{-}cseq}(d, q) \wedge \\ \quad \operatorname{vadd}(\operatorname{cseq\text{-}module}(d, q))(p, \operatorname{vneg}(\operatorname{cseq\text{-}module}(d, q))(a)) \\ \quad \in\; N_{\operatorname{limit\text{-}form}(d, q)} \Rightarrow \\ \quad \quad \operatorname{cseq\text{-}equiv}(\operatorname{form\text{-}ms}(d, q), p, a) \end{array}}}
\vnbprop{hc-rel-bwd}{\vnbstmt{\begin{array}{@{}l@{}} \forall d, q, p, a. \\ \quad \operatorname{is\text{-}module}(d) \wedge \\ \quad \operatorname{scal}(d) = \operatorname{normed\text{-}field\text{-}as\text{-}commutative\text{-}ring}(\operatorname{cc\text{-}normed\text{-}field}) \wedge \\ \quad \operatorname{is\text{-}inner\text{-}product\text{-}form}(d, q) \wedge \\ \quad p \in \operatorname{form\text{-}cseq}(d, q) \wedge \\ \quad a \in \operatorname{form\text{-}cseq}(d, q) \wedge \\ \quad \operatorname{cseq\text{-}equiv}(\operatorname{form\text{-}ms}(d, q), p, a) \Rightarrow \\ \quad \quad \operatorname{vadd}(\operatorname{cseq\text{-}module}(d, q))(p, \operatorname{vneg}(\operatorname{cseq\text{-}module}(d, q))(a)) \\ \quad \quad \in\; N_{\operatorname{limit\text{-}form}(d, q)} \end{array}}}
\vnbprop{hc-rel-eq}{\vnbstmt{\begin{array}{@{}l@{}} \forall d, q. \\ \quad \operatorname{is\text{-}module}(d) \wedge \\ \quad \operatorname{scal}(d) = \operatorname{normed\text{-}field\text{-}as\text{-}commutative\text{-}ring}(\operatorname{cc\text{-}normed\text{-}field}) \wedge \\ \quad \operatorname{is\text{-}inner\text{-}product\text{-}form}(d, q) \Rightarrow \\ \quad \quad \operatorname{qmod\text{-}rel}(\operatorname{cseq\text{-}module}(d, q), N_{\operatorname{limit\text{-}form}(d, q)}) = \operatorname{crel}(\operatorname{form\text{-}ms}(d, q)) \end{array}}}
\vnbprop{hc-setoid-eq}{\vnbstmt{\begin{array}{@{}l@{}} \forall d, q. \\ \quad \operatorname{is\text{-}module}(d) \wedge \\ \quad \operatorname{scal}(d) = \operatorname{normed\text{-}field\text{-}as\text{-}commutative\text{-}ring}(\operatorname{cc\text{-}normed\text{-}field}) \wedge \\ \quad \operatorname{is\text{-}inner\text{-}product\text{-}form}(d, q) \Rightarrow \\ \quad \quad \operatorname{qmod\text{-}setoid}(\operatorname{cseq\text{-}module}(d, q), N_{\operatorname{limit\text{-}form}(d, q)}) \\ \quad \quad =\; \operatorname{cauchy\text{-}setoid}(\operatorname{form\text{-}ms}(d, q)) \end{array}}}
\vnbprop{hc-vec}{\vnbstmt{\begin{array}{@{}l@{}} \forall d, q. \\ \quad \operatorname{is\text{-}module}(d) \wedge \\ \quad \operatorname{scal}(d) = \operatorname{normed\text{-}field\text{-}as\text{-}commutative\text{-}ring}(\operatorname{cc\text{-}normed\text{-}field}) \wedge \\ \quad \operatorname{is\text{-}inner\text{-}product\text{-}form}(d, q) \Rightarrow \\ \quad \quad (\operatorname{vec}(\operatorname{hilbert\text{-}completion}(d, q)) \simeq \operatorname{quotient}(\operatorname{cauchy\text{-}setoid}(\operatorname{form\text{-}ms}(d, q)))) \end{array}}}
\vnbprop{hc-pts}{\vnbstmt{\begin{array}{@{}l@{}} \forall d, q. \\ \quad \operatorname{is\text{-}module}(d) \wedge \\ \quad \operatorname{scal}(d) = \operatorname{normed\text{-}field\text{-}as\text{-}commutative\text{-}ring}(\operatorname{cc\text{-}normed\text{-}field}) \wedge \\ \quad \operatorname{is\text{-}inner\text{-}product\text{-}form}(d, q) \Rightarrow \\ \quad \quad (\operatorname{pts}(\operatorname{ip\text{-}metric\text{-}space}(\operatorname{hilbert\text{-}completion}(d, q))) \simeq \operatorname{pts}(\operatorname{completion}(\operatorname{form\text{-}ms}(d, q)))) \end{array}}}
\vnbprop{hc-dist}{\vnbstmt{\begin{array}{@{}l@{}} \forall d, q, x, y. \\ \quad \operatorname{is\text{-}module}(d) \wedge \\ \quad \operatorname{scal}(d) = \operatorname{normed\text{-}field\text{-}as\text{-}commutative\text{-}ring}(\operatorname{cc\text{-}normed\text{-}field}) \wedge \\ \quad \operatorname{is\text{-}inner\text{-}product\text{-}form}(d, q) \wedge \\ \quad x \in \operatorname{pts}(\operatorname{completion}(\operatorname{form\text{-}ms}(d, q))) \wedge \\ \quad y \in \operatorname{pts}(\operatorname{completion}(\operatorname{form\text{-}ms}(d, q))) \Rightarrow \\ \quad \quad \operatorname{dist}(\operatorname{ip\text{-}metric\text{-}space}(\operatorname{hilbert\text{-}completion}(d, q)))(x, y) = \operatorname{dist}(\operatorname{completion}(\operatorname{form\text{-}ms}(d, q)))(x, y) \end{array}}}
\vnbprop{hilbert-completion-complete}{\vnbstmt{\begin{array}{@{}l@{}} \forall d, q. \\ \quad \operatorname{is\text{-}module}(d) \wedge \\ \quad \operatorname{scal}(d) = \operatorname{normed\text{-}field\text{-}as\text{-}commutative\text{-}ring}(\operatorname{cc\text{-}normed\text{-}field}) \wedge \\ \quad \operatorname{is\text{-}inner\text{-}product\text{-}form}(d, q) \Rightarrow \\ \quad \quad \operatorname{is\text{-}complete}(\operatorname{ip\text{-}metric\text{-}space}(\operatorname{hilbert\text{-}completion}(d, q))) \end{array}}}
\vnbprop{hilbert-completion-is-hilbert}{\vnbstmt{\begin{array}{@{}l@{}} \forall d, q. \\ \quad \operatorname{is\text{-}module}(d) \wedge \\ \quad \operatorname{scal}(d) = \operatorname{normed\text{-}field\text{-}as\text{-}commutative\text{-}ring}(\operatorname{cc\text{-}normed\text{-}field}) \wedge \\ \quad \operatorname{is\text{-}inner\text{-}product\text{-}form}(d, q) \Rightarrow \\ \quad \quad \operatorname{is\text{-}hilbert\text{-}space}(\operatorname{hilbert\text{-}completion}(d, q)) \end{array}}}
\vnbprop{hc-embed-type}{\vnbstmt{\begin{array}{@{}l@{}} \forall d, q, x. \\ \quad \operatorname{is\text{-}module}(d) \wedge \\ \quad \operatorname{scal}(d) = \operatorname{normed\text{-}field\text{-}as\text{-}commutative\text{-}ring}(\operatorname{cc\text{-}normed\text{-}field}) \wedge \\ \quad \operatorname{is\text{-}inner\text{-}product\text{-}form}(d, q) \wedge \\ \quad x \in \operatorname{vec}(d) \Rightarrow \\ \quad \quad \operatorname{hc\text{-}embed}(d, q)(x) \in \operatorname{vec}(\operatorname{hilbert\text{-}completion}(d, q)) \end{array}}}
\vnbprop{hc-embed-near}{\vnbstmt{\begin{array}{@{}l@{}} \forall d, q, x, s. \\ \quad \operatorname{is\text{-}module}(d) \wedge \\ \quad \operatorname{scal}(d) = \operatorname{normed\text{-}field\text{-}as\text{-}commutative\text{-}ring}(\operatorname{cc\text{-}normed\text{-}field}) \wedge \\ \quad \operatorname{is\text{-}inner\text{-}product\text{-}form}(d, q) \wedge \\ \quad x \in \operatorname{vec}(\operatorname{hilbert\text{-}completion}(d, q)) \wedge \\ \quad \operatorname{pos\text{-}rr}(s) \Rightarrow \\ \quad \quad \exists z \in \operatorname{vec}(d). \\ \quad \quad \quad \operatorname{dist}(\operatorname{ip\text{-}metric\text{-}space}(\operatorname{hilbert\text{-}completion}(d, q)))(x, \operatorname{hc\text{-}embed}(d, q)(z)) \leq s \end{array}}}
\vnbprop{hc-embed-dense}{\vnbstmt{\begin{array}{@{}l@{}} \forall d, q. \\ \quad \operatorname{is\text{-}module}(d) \wedge \\ \quad \operatorname{scal}(d) = \operatorname{normed\text{-}field\text{-}as\text{-}commutative\text{-}ring}(\operatorname{cc\text{-}normed\text{-}field}) \wedge \\ \quad \operatorname{is\text{-}inner\text{-}product\text{-}form}(d, q) \Rightarrow \\ \quad \quad \operatorname{is\text{-}dense}(\operatorname{ip\text{-}metric\text{-}space}(\operatorname{hilbert\text{-}completion}(d, q)), \operatorname{image}(\operatorname{hc\text{-}embed}(d, q), \operatorname{vec}(d))) \end{array}}}
\vnbprop{hc-embed-isometry}{\vnbstmt{\begin{array}{@{}l@{}} \forall d, q, x, y. \\ \quad \operatorname{is\text{-}module}(d) \wedge \\ \quad \operatorname{scal}(d) = \operatorname{normed\text{-}field\text{-}as\text{-}commutative\text{-}ring}(\operatorname{cc\text{-}normed\text{-}field}) \wedge \\ \quad \operatorname{is\text{-}inner\text{-}product\text{-}form}(d, q) \wedge \\ \quad x \in \operatorname{vec}(d) \wedge \\ \quad y \in \operatorname{vec}(d) \Rightarrow \\ \quad \quad \operatorname{dist}(\operatorname{ip\text{-}metric\text{-}space}(\operatorname{hilbert\text{-}completion}(d, q)))(\operatorname{hc\text{-}embed}(d, q)(x), \operatorname{hc\text{-}embed}(d, q)(y)) \\ \quad \quad =\; \operatorname{dist}(\operatorname{form\text{-}ms}(d, q))(x, y) \end{array}}}
\vnbprop{hc-embed-injective}{\vnbstmt{\begin{array}{@{}l@{}} \forall d, q, x, y. \\ \quad \operatorname{is\text{-}module}(d) \wedge \\ \quad \operatorname{scal}(d) = \operatorname{normed\text{-}field\text{-}as\text{-}commutative\text{-}ring}(\operatorname{cc\text{-}normed\text{-}field}) \wedge \\ \quad \operatorname{is\text{-}inner\text{-}product\text{-}form}(d, q) \wedge \\ \quad x \in \operatorname{vec}(d) \wedge \\ \quad y \in \operatorname{vec}(d) \wedge \\ \quad \operatorname{hc\text{-}embed}(d, q)(x) = \operatorname{hc\text{-}embed}(d, q)(y) \Rightarrow \\ \quad \quad x = y \end{array}}}
\vnbprop{hc-embed-inverse}{\vnbstmt{\begin{array}{@{}l@{}} \forall d, q, x. \\ \quad \operatorname{is\text{-}module}(d) \wedge \\ \quad \operatorname{scal}(d) = \operatorname{normed\text{-}field\text{-}as\text{-}commutative\text{-}ring}(\operatorname{cc\text{-}normed\text{-}field}) \wedge \\ \quad \operatorname{is\text{-}inner\text{-}product\text{-}form}(d, q) \wedge \\ \quad x \in \operatorname{vec}(d) \Rightarrow \\ \quad \quad (\iota\, u.\;((u \in \operatorname{vec}(d)) \wedge (\operatorname{hc\text{-}embed}(d, q)(u) = \operatorname{hc\text{-}embed}(d, q)(x)))) = x \end{array}}}
\vnbprop{hilbert-completion-universal}{\vnbstmt{\begin{array}{@{}l@{}} \forall d, q, k, t, b. \\ \quad \operatorname{is\text{-}module}(d) \wedge \\ \quad \operatorname{scal}(d) = \operatorname{normed\text{-}field\text{-}as\text{-}commutative\text{-}ring}(\operatorname{cc\text{-}normed\text{-}field}) \wedge \\ \quad \operatorname{is\text{-}inner\text{-}product\text{-}form}(d, q) \wedge \\ \quad \operatorname{is\text{-}hilbert\text{-}space}(k) \wedge \\ \quad t \in (\operatorname{vec}(d) \to \operatorname{vec}(k)) \wedge \\ \quad b \in \mathbb{R} \wedge \\ \quad 0 \leq b \wedge \\ \quad \forall x, y \in \operatorname{vec}(d). \\ \quad \quad \operatorname{dist}(\operatorname{ip\text{-}metric\text{-}space}(k))(t(x), t(y)) \leq (b \cdot \operatorname{dist}(\operatorname{form\text{-}ms}(d, q))(x, y)) \Rightarrow \\ \quad \quad \exists g. \\ \quad \quad \quad \operatorname{is\text{-}uniformly\text{-}continuous}(\operatorname{ip\text{-}metric\text{-}space}(\operatorname{hilbert\text{-}completion}(d, q)), \operatorname{ip\text{-}metric\text{-}space}(k), g) \wedge \\ \quad \quad \quad \forall x \in \operatorname{vec}(d).\; (g(\operatorname{hc\text{-}embed}(d, q)(x)) = t(x)) \wedge \\ \quad \quad \quad \forall h. \\ \quad \quad \quad \quad \operatorname{is\text{-}uniformly\text{-}continuous}(\operatorname{ip\text{-}metric\text{-}space}(\operatorname{hilbert\text{-}completion}(d, q)), \operatorname{ip\text{-}metric\text{-}space}(k), h) \wedge \\ \quad \quad \quad \quad \forall x \in \operatorname{vec}(d).\; (h(\operatorname{hc\text{-}embed}(d, q)(x)) = t(x)) \Rightarrow \\ \quad \quad \quad \quad \quad h = g \end{array}}}
\vnbprop{form-dist-add-le}{\vnbstmt{\begin{array}{@{}l@{}} \forall d, q, x, y, u, w. \\ \quad \operatorname{is\text{-}module}(d) \wedge \\ \quad \operatorname{scal}(d) = \operatorname{normed\text{-}field\text{-}as\text{-}commutative\text{-}ring}(\operatorname{cc\text{-}normed\text{-}field}) \wedge \\ \quad \operatorname{is\text{-}inner\text{-}product\text{-}form}(d, q) \wedge \\ \quad x \in \operatorname{vec}(d) \wedge \\ \quad y \in \operatorname{vec}(d) \wedge \\ \quad u \in \operatorname{vec}(d) \wedge \\ \quad w \in \operatorname{vec}(d) \Rightarrow \\ \quad \quad \operatorname{dist}(\operatorname{form\text{-}ms}(d, q))(\operatorname{vadd}(d)(x, y), \operatorname{vadd}(d)(u, w)) \\ \quad \quad \leq\; \operatorname{dist}(\operatorname{form\text{-}ms}(d, q))(x, u) + \operatorname{dist}(\operatorname{form\text{-}ms}(d, q))(y, w) \end{array}}}
\vnbprop{hc-ms-form}{\vnbstmt{\begin{array}{@{}l@{}} \forall d, q. \\ \quad \operatorname{is\text{-}module}(d) \wedge \\ \quad \operatorname{scal}(d) = \operatorname{normed\text{-}field\text{-}as\text{-}commutative\text{-}ring}(\operatorname{cc\text{-}normed\text{-}field}) \wedge \\ \quad \operatorname{is\text{-}inner\text{-}product\text{-}form}(d, q) \Rightarrow \\ \quad \quad (\operatorname{ip\text{-}metric\text{-}space}(\operatorname{hilbert\text{-}completion}(d, q)) \simeq \operatorname{form\text{-}ms}(\operatorname{quotient\text{-}module}(\operatorname{cseq\text{-}module}(d, q), N_{\operatorname{limit\text{-}form}(d, q)}), \operatorname{form\text{-}qip}(\operatorname{cseq\text{-}module}(d, q), \operatorname{limit\text{-}form}(d, q)))) \end{array}}}
\vnbprop{hc-vadd-form}{\vnbstmt{\begin{array}{@{}l@{}} \forall d, q. \\ \quad \operatorname{is\text{-}module}(d) \wedge \\ \quad \operatorname{scal}(d) = \operatorname{normed\text{-}field\text{-}as\text{-}commutative\text{-}ring}(\operatorname{cc\text{-}normed\text{-}field}) \wedge \\ \quad \operatorname{is\text{-}inner\text{-}product\text{-}form}(d, q) \Rightarrow \\ \quad \quad (\operatorname{vadd}(\operatorname{hilbert\text{-}completion}(d, q)) \simeq \operatorname{vadd}(\operatorname{quotient\text{-}module}(\operatorname{cseq\text{-}module}(d, q), N_{\operatorname{limit\text{-}form}(d, q)}))) \end{array}}}
\vnbprop{hc-vec-form}{\vnbstmt{\begin{array}{@{}l@{}} \forall d, q. \\ \quad \operatorname{is\text{-}module}(d) \wedge \\ \quad \operatorname{scal}(d) = \operatorname{normed\text{-}field\text{-}as\text{-}commutative\text{-}ring}(\operatorname{cc\text{-}normed\text{-}field}) \wedge \\ \quad \operatorname{is\text{-}inner\text{-}product\text{-}form}(d, q) \Rightarrow \\ \quad \quad (\operatorname{vec}(\operatorname{hilbert\text{-}completion}(d, q)) \simeq \operatorname{vec}(\operatorname{quotient\text{-}module}(\operatorname{cseq\text{-}module}(d, q), N_{\operatorname{limit\text{-}form}(d, q)}))) \end{array}}}
\vnbprop{hc-act-form}{\vnbstmt{\begin{array}{@{}l@{}} \forall d, q. \\ \quad \operatorname{is\text{-}module}(d) \wedge \\ \quad \operatorname{scal}(d) = \operatorname{normed\text{-}field\text{-}as\text{-}commutative\text{-}ring}(\operatorname{cc\text{-}normed\text{-}field}) \wedge \\ \quad \operatorname{is\text{-}inner\text{-}product\text{-}form}(d, q) \Rightarrow \\ \quad \quad (\operatorname{act}(\operatorname{hilbert\text{-}completion}(d, q)) \simeq \operatorname{act}(\operatorname{quotient\text{-}module}(\operatorname{cseq\text{-}module}(d, q), N_{\operatorname{limit\text{-}form}(d, q)}))) \end{array}}}
\vnbprop{hc-vadd-type}{\vnbstmt{\begin{array}{@{}l@{}} \forall d, q, x, y. \\ \quad \operatorname{is\text{-}module}(d) \wedge \\ \quad \operatorname{scal}(d) = \operatorname{normed\text{-}field\text{-}as\text{-}commutative\text{-}ring}(\operatorname{cc\text{-}normed\text{-}field}) \wedge \\ \quad \operatorname{is\text{-}inner\text{-}product\text{-}form}(d, q) \wedge \\ \quad x \in \operatorname{vec}(\operatorname{hilbert\text{-}completion}(d, q)) \wedge \\ \quad y \in \operatorname{vec}(\operatorname{hilbert\text{-}completion}(d, q)) \Rightarrow \\ \quad \quad \operatorname{vadd}(\operatorname{hilbert\text{-}completion}(d, q))(x, y) \in \operatorname{vec}(\operatorname{hilbert\text{-}completion}(d, q)) \end{array}}}
\vnbprop{hc-dist-add-le}{\vnbstmt{\begin{array}{@{}l@{}} \forall d, q, x, y, u, w. \\ \quad \operatorname{is\text{-}module}(d) \wedge \\ \quad \operatorname{scal}(d) = \operatorname{normed\text{-}field\text{-}as\text{-}commutative\text{-}ring}(\operatorname{cc\text{-}normed\text{-}field}) \wedge \\ \quad \operatorname{is\text{-}inner\text{-}product\text{-}form}(d, q) \wedge \\ \quad x \in \operatorname{vec}(\operatorname{hilbert\text{-}completion}(d, q)) \wedge \\ \quad y \in \operatorname{vec}(\operatorname{hilbert\text{-}completion}(d, q)) \wedge \\ \quad u \in \operatorname{vec}(\operatorname{hilbert\text{-}completion}(d, q)) \wedge \\ \quad w \in \operatorname{vec}(\operatorname{hilbert\text{-}completion}(d, q)) \Rightarrow \\ \quad \quad \operatorname{dist}(\operatorname{ip\text{-}metric\text{-}space}(\operatorname{hilbert\text{-}completion}(d, q)))(\operatorname{vadd}(\operatorname{hilbert\text{-}completion}(d, q))(x, y), \\ \quad \quad \quad \operatorname{vadd}(\operatorname{hilbert\text{-}completion}(d, q))(u, w)) \\ \quad \quad \leq\; \operatorname{dist}(\operatorname{ip\text{-}metric\text{-}space}(\operatorname{hilbert\text{-}completion}(d, q)))(x, u) + \operatorname{dist}(\operatorname{ip\text{-}metric\text{-}space}(\operatorname{hilbert\text{-}completion}(d, q)))(y, w) \end{array}}}
\vnbprop{hc-ext-additive}{\vnbstmt{\begin{array}{@{}l@{}} \forall d, q, m, \mathit{hcq}, t, g.; \forall x, y \in \operatorname{vec}(\operatorname{hilbert\text{-}completion}(d, q)). \\ \quad \operatorname{is\text{-}module}(d) \wedge \\ \quad \operatorname{scal}(d) = \operatorname{normed\text{-}field\text{-}as\text{-}commutative\text{-}ring}(\operatorname{cc\text{-}normed\text{-}field}) \wedge \\ \quad \operatorname{is\text{-}inner\text{-}product\text{-}form}(d, q) \wedge \\ \quad \operatorname{is\text{-}module}(m) \wedge \\ \quad \operatorname{scal}(m) = \operatorname{normed\text{-}field\text{-}as\text{-}commutative\text{-}ring}(\operatorname{cc\text{-}normed\text{-}field}) \wedge \\ \quad \operatorname{is\text{-}inner\text{-}product\text{-}form}(m, \mathit{hcq}) \wedge \\ \quad t \in (\operatorname{vec}(d) \to \operatorname{vec}(m)) \wedge \\ \quad \forall \mathit{hcx}, \mathit{hcy} \in \operatorname{vec}(d). \\ \quad \quad t(\operatorname{vadd}(d)(\mathit{hcx}, \mathit{hcy})) = \operatorname{vadd}(m)(t(\mathit{hcx}), t(\mathit{hcy})) \wedge \\ \quad \operatorname{is\text{-}uniformly\text{-}continuous}(\operatorname{ip\text{-}metric\text{-}space}(\operatorname{hilbert\text{-}completion}(d, q)), \operatorname{form\text{-}ms}(m, \mathit{hcq}), g) \wedge \\ \quad \forall \mathit{hcx} \in \operatorname{vec}(d).\; (g(\operatorname{hc\text{-}embed}(d, q)(\mathit{hcx})) = t(\mathit{hcx})) \Rightarrow \\ \quad \quad g(\operatorname{vadd}(\operatorname{hilbert\text{-}completion}(d, q))(x, y)) = \operatorname{vadd}(m)(g(x), g(y)) \end{array}}}
\vnbprop{form-dist-scale}{\vnbstmt{\begin{array}{@{}l@{}} \forall d, q, a, x, y. \\ \quad \operatorname{is\text{-}module}(d) \wedge \\ \quad \operatorname{scal}(d) = \operatorname{normed\text{-}field\text{-}as\text{-}commutative\text{-}ring}(\operatorname{cc\text{-}normed\text{-}field}) \wedge \\ \quad \operatorname{is\text{-}inner\text{-}product\text{-}form}(d, q) \wedge \\ \quad a \in \mathbb{C} \wedge \\ \quad x \in \operatorname{vec}(d) \wedge \\ \quad y \in \operatorname{vec}(d) \Rightarrow \\ \quad \quad \operatorname{dist}(\operatorname{form\text{-}ms}(d, q))(\operatorname{act}(d)(a, x), \operatorname{act}(d)(a, y)) = (\operatorname{magnitude}(a) \cdot \operatorname{dist}(\operatorname{form\text{-}ms}(d, q))(x, y)) \end{array}}}
\vnbprop{hc-act-type}{\vnbstmt{\begin{array}{@{}l@{}} \forall d, q, a, x. \\ \quad \operatorname{is\text{-}module}(d) \wedge \\ \quad \operatorname{scal}(d) = \operatorname{normed\text{-}field\text{-}as\text{-}commutative\text{-}ring}(\operatorname{cc\text{-}normed\text{-}field}) \wedge \\ \quad \operatorname{is\text{-}inner\text{-}product\text{-}form}(d, q) \wedge \\ \quad a \in \mathbb{C} \wedge \\ \quad x \in \operatorname{vec}(\operatorname{hilbert\text{-}completion}(d, q)) \Rightarrow \\ \quad \quad \operatorname{act}(\operatorname{hilbert\text{-}completion}(d, q))(a, x) \in \operatorname{vec}(\operatorname{hilbert\text{-}completion}(d, q)) \end{array}}}
\vnbprop{hc-dist-scale}{\vnbstmt{\begin{array}{@{}l@{}} \forall d, q, a, x, u. \\ \quad \operatorname{is\text{-}module}(d) \wedge \\ \quad \operatorname{scal}(d) = \operatorname{normed\text{-}field\text{-}as\text{-}commutative\text{-}ring}(\operatorname{cc\text{-}normed\text{-}field}) \wedge \\ \quad \operatorname{is\text{-}inner\text{-}product\text{-}form}(d, q) \wedge \\ \quad a \in \mathbb{C} \wedge \\ \quad x \in \operatorname{vec}(\operatorname{hilbert\text{-}completion}(d, q)) \wedge \\ \quad u \in \operatorname{vec}(\operatorname{hilbert\text{-}completion}(d, q)) \Rightarrow \\ \quad \quad \operatorname{dist}(\operatorname{ip\text{-}metric\text{-}space}(\operatorname{hilbert\text{-}completion}(d, q)))(\operatorname{act}(\operatorname{hilbert\text{-}completion}(d, q))(a, x), \\ \quad \quad \quad \operatorname{act}(\operatorname{hilbert\text{-}completion}(d, q))(a, u)) \\ \quad \quad =\; \operatorname{magnitude}(a) \cdot \operatorname{dist}(\operatorname{ip\text{-}metric\text{-}space}(\operatorname{hilbert\text{-}completion}(d, q)))(x, u) \end{array}}}
\vnbprop{hc-ext-homogeneous}{\vnbstmt{\begin{array}{@{}l@{}} \forall d, q, m, \mathit{hcq}, t, g.; \forall a \in \mathbb{C}, x \in \operatorname{vec}(\operatorname{hilbert\text{-}completion}(d, q)). \\ \quad \operatorname{is\text{-}module}(d) \wedge \\ \quad \operatorname{scal}(d) = \operatorname{normed\text{-}field\text{-}as\text{-}commutative\text{-}ring}(\operatorname{cc\text{-}normed\text{-}field}) \wedge \\ \quad \operatorname{is\text{-}inner\text{-}product\text{-}form}(d, q) \wedge \\ \quad \operatorname{is\text{-}module}(m) \wedge \\ \quad \operatorname{scal}(m) = \operatorname{normed\text{-}field\text{-}as\text{-}commutative\text{-}ring}(\operatorname{cc\text{-}normed\text{-}field}) \wedge \\ \quad \operatorname{is\text{-}inner\text{-}product\text{-}form}(m, \mathit{hcq}) \wedge \\ \quad t \in (\operatorname{vec}(d) \to \operatorname{vec}(m)) \wedge \\ \quad \forall \mathit{hca} \in \mathbb{C}, \mathit{hcx} \in \operatorname{vec}(d). \\ \quad \quad t(\operatorname{act}(d)(\mathit{hca}, \mathit{hcx})) = \operatorname{act}(m)(\mathit{hca}, t(\mathit{hcx})) \wedge \\ \quad \operatorname{is\text{-}uniformly\text{-}continuous}(\operatorname{ip\text{-}metric\text{-}space}(\operatorname{hilbert\text{-}completion}(d, q)), \operatorname{form\text{-}ms}(m, \mathit{hcq}), g) \wedge \\ \quad \forall \mathit{hcx} \in \operatorname{vec}(d).\; (g(\operatorname{hc\text{-}embed}(d, q)(\mathit{hcx})) = t(\mathit{hcx})) \Rightarrow \\ \quad \quad g(\operatorname{act}(\operatorname{hilbert\text{-}completion}(d, q))(a, x)) = \operatorname{act}(m)(a, g(x)) \end{array}}}

\vnbbookfile{theorem-library/ba-trace-class}{theorem-library/ba-trace-class}
\noindent{\small bta-trace.scm -- BA-1111-A, 2026-11-11.  BanachAlgebras.pdf ch. 2 s.2-3: the Hilbert-Schmidt inner product as an inner product, HS norm of an operator diagonal in a basis (2.11), and the trace-class laws 2.13-2.23 on top of the TH3 development (th3-vocab: IS-TRACE-CLASS is Def. 2.13, T = A B with A, B Hilbert-Schmidt; ...}\par
\vnbprop{trace-class-is-hs}{\vnbstmt{\begin{array}{@{}l@{}} \forall v, t. \\ \quad (\operatorname{is\text{-}hilbert\text{-}space}(v) \wedge \operatorname{is\text{-}trace\text{-}class}(v, t)) \Rightarrow \operatorname{is\text{-}hilbert\text{-}schmidt}(v, v, t) \end{array}}}
\vnbprop{trace-class-comp-left}{\vnbstmt{\begin{array}{@{}l@{}} \forall v, s, t. \\ \quad (\operatorname{is\text{-}hilbert\text{-}space}(v) \wedge ((s \in L({v}, {v})) \wedge \operatorname{is\text{-}trace\text{-}class}(v, t))) \Rightarrow  \\ \quad \quad \operatorname{is\text{-}trace\text{-}class}(v, {s}{t}) \end{array}}}
\vnbprop{trace-class-comp-right}{\vnbstmt{\begin{array}{@{}l@{}} \forall v, s, t. \\ \quad (\operatorname{is\text{-}hilbert\text{-}space}(v) \wedge ((s \in L({v}, {v})) \wedge \operatorname{is\text{-}trace\text{-}class}(v, t))) \Rightarrow  \\ \quad \quad \operatorname{is\text{-}trace\text{-}class}(v, {t}{s}) \end{array}}}
\vnbprop{op-trace-factor}{\vnbstmt{\begin{array}{@{}l@{}} \forall v, a, b. \\ \quad (\operatorname{is\text{-}hilbert\text{-}space}(v) \wedge (\operatorname{is\text{-}hilbert\text{-}schmidt}(v, v, a) \wedge \operatorname{is\text{-}hilbert\text{-}schmidt}(v, v, b))) \Rightarrow  \\ \quad \quad \mathrm{tr}\,{{a}{b}} = \langle {b} \mid {{a}^{*}} \rangle_{HS} \end{array}}}
\vnbprop{op-trace-comm-bounded}{\vnbstmt{\begin{array}{@{}l@{}} \forall v, t, s. \\ \quad (\operatorname{is\text{-}hilbert\text{-}space}(v) \wedge (\operatorname{is\text{-}trace\text{-}class}(v, t) \wedge (s \in L({v}, {v})))) \Rightarrow  \\ \quad \quad \mathrm{tr}\,{{t}{s}} = \mathrm{tr}\,{{s}{t}} \end{array}}}
\vnbprop{trace-class-adjoint}{\vnbstmt{\begin{array}{@{}l@{}} \forall v, t. \\ \quad (\operatorname{is\text{-}hilbert\text{-}space}(v) \wedge \operatorname{is\text{-}trace\text{-}class}(v, t)) \Rightarrow  \\ \quad \quad \operatorname{is\text{-}trace\text{-}class}(v, {t}^{*}) \end{array}}}
\vnbprop{op-trace-adjoint}{\vnbstmt{\begin{array}{@{}l@{}} \forall v, t. \\ \quad (\operatorname{is\text{-}hilbert\text{-}space}(v) \wedge \operatorname{is\text{-}trace\text{-}class}(v, t)) \Rightarrow  \\ \quad \quad \mathrm{tr}\,{{t}^{*}} = \operatorname{conjugate}(\mathrm{tr}\,{t}) \end{array}}}
\vnbprop{hs-ip-homog-left}{\vnbstmt{\begin{array}{@{}l@{}} \forall v, w, c, s, t. \\ \quad \operatorname{is\text{-}hilbert\text{-}space}(v) \wedge \\ \quad \operatorname{is\text{-}hilbert\text{-}space}(w) \wedge \\ \quad c \in \mathbb{C} \wedge \\ \quad \operatorname{is\text{-}hilbert\text{-}schmidt}(v, w, s) \wedge \\ \quad \operatorname{is\text{-}hilbert\text{-}schmidt}(v, w, t) \Rightarrow \\ \quad \quad \langle {{c}\,{s}} \mid {t} \rangle_{HS} = (c \cdot \langle {s} \mid {t} \rangle_{HS}) \end{array}}}
\vnbprop{hs-ip-self-nonneg}{\vnbstmt{\begin{array}{@{}l@{}} \forall v, w, t. \\ \quad (\operatorname{is\text{-}hilbert\text{-}space}(v) \wedge (\operatorname{is\text{-}hilbert\text{-}space}(w) \wedge \operatorname{is\text{-}hilbert\text{-}schmidt}(v, w, t))) \Rightarrow  \\ \quad \quad (\langle {t} \mid {t} \rangle_{HS} \in \mathbb{R}) \wedge (0 \leq \langle {t} \mid {t} \rangle_{HS}) \end{array}}}
\vnbprop{hs-sum-zero-vanishes}{\vnbstmt{\begin{array}{@{}l@{}} \forall v, w, t, x, e.; \forall a \in \operatorname{vec}(v). \\ \quad \operatorname{is\text{-}hilbert\text{-}space}(v) \wedge \\ \quad \operatorname{is\text{-}hilbert\text{-}space}(w) \wedge \\ \quad t \in L({v}, {w}) \wedge \\ \quad \operatorname{is\text{-}orthonormal\text{-}basis}(v, x, e) \wedge \\ \quad \sum_{i \in {x}} \lVert {t} {e}_i \rVert^2 = 0 \Rightarrow \\ \quad \quad t(a) = \operatorname{vzero}(w) \end{array}}}
\vnbprop{hs-ip-definite}{\vnbstmt{\begin{array}{@{}l@{}} \forall v, w, t.; \forall x \in \operatorname{vec}(v). \\ \quad \operatorname{is\text{-}hilbert\text{-}space}(v) \wedge \\ \quad \operatorname{is\text{-}hilbert\text{-}space}(w) \wedge \\ \quad \operatorname{is\text{-}hilbert\text{-}schmidt}(v, w, t) \wedge \\ \quad \langle {t} \mid {t} \rangle_{HS} = 0 \Rightarrow \\ \quad \quad t(x) = \operatorname{vzero}(w) \end{array}}}
\vnbprop{trace-pairing-nondeg-left}{\vnbstmt{\begin{array}{@{}l@{}} \forall v, t.; \forall x \in \operatorname{vec}(v). \\ \quad \operatorname{is\text{-}hilbert\text{-}space}(v) \wedge \\ \quad \operatorname{is\text{-}trace\text{-}class}(v, t) \wedge \\ \quad \forall s \in L({v}, {v}).\; (\mathrm{tr}\,{{s}{t}} = 0) \Rightarrow \\ \quad \quad t(x) = \operatorname{vzero}(v) \end{array}}}
\vnbprop{hs-sum-eigenbasis}{\vnbstmt{\begin{array}{@{}l@{}} \forall v, t, x, e, \ell. \\ \quad \operatorname{is\text{-}hilbert\text{-}space}(v) \wedge \\ \quad t \in L({v}, {v}) \wedge \\ \quad \operatorname{is\text{-}orthonormal\text{-}basis}(v, x, e) \wedge \\ \quad \ell \in (x \to \mathbb{C}) \wedge \\ \quad \forall k \in x.\; (t(e(k)) = \operatorname{act}(v)(\ell(k), e(k))) \Rightarrow \\ \quad \quad \sum_{i \in {x}} \lVert {t} {e}_i \rVert^2 = \operatorname{esum}((\lambda \mathit{kqk} \in x.\; (\operatorname{magnitude}(\ell(\mathit{kqk})) \cdot \operatorname{magnitude}(\ell(\mathit{kqk}))))) \end{array}}}
\vnbprop{hs-sum-compact-sa}{\vnbstmt{\begin{array}{@{}l@{}} \forall v, t. \\ \quad (\operatorname{is\text{-}hilbert\text{-}space}(v) \wedge (\operatorname{is\text{-}compact\text{-}op}(v, v, t) \wedge \operatorname{is\text{-}self\text{-}adjoint\text{-}op}(v, t))) \Rightarrow  \\ \quad \quad \exists i, e, l. \\ \quad \quad \quad \operatorname{is\text{-}orthonormal\text{-}basis}(v, i, e) \wedge \\ \quad \quad \quad l \in (i \to \mathbb{C}) \wedge \\ \quad \quad \quad \forall k \in i.\; (t(e(k)) = \operatorname{act}(v)(l(k), e(k))) \wedge \\ \quad \quad \quad \sum_{i \in {i}} \lVert {t} {e}_i \rVert^2 = \operatorname{esum}((\lambda \mathit{kqk} \in i.\; (\operatorname{magnitude}(l(\mathit{kqk})) \cdot \operatorname{magnitude}(l(\mathit{kqk}))))) \end{array}}}
\vnbprop{rank-one-op-hs}{\vnbstmt{\begin{array}{@{}l@{}} \forall v, a, b. \\ \quad (\operatorname{is\text{-}hilbert\text{-}space}(v) \wedge ((a \in \operatorname{vec}(v)) \wedge (b \in \operatorname{vec}(v)))) \Rightarrow  \\ \quad \quad \operatorname{is\text{-}hilbert\text{-}schmidt}(v, v, \langle \cdot \mid {a} \rangle {b}) \end{array}}}
\vnbprop{trace-pairing-nondeg-right}{\vnbstmt{\begin{array}{@{}l@{}} \forall v.; \forall s \in L({v}, {v}), x \in \operatorname{vec}(v). \\ \quad (\operatorname{is\text{-}hilbert\text{-}space}(v) \wedge \forall t.\; (\operatorname{is\text{-}trace\text{-}class}(v, t) \Rightarrow (\mathrm{tr}\,{{s}{t}} = 0))) \Rightarrow  \\ \quad \quad s(x) = \operatorname{vzero}(v) \end{array}}}

\vnbbookfile{structure-library/ba-trace-dual-vocab}{structure-library/ba-trace-dual-vocab}
\noindent{\small ba-trace-dual-vocab.scm -- Thayer, Banach Algebras and Applications, chapter 5 s.6-s.8 (spaces of ultraweakly continuous functionals, duals, transpose and module structure of L\_*(H)). DEFINITIONS ONLY: no proof, no support, no axiom. Agent BA-1117-TRACE, 2026-11-17.  DESTINATION: ...}\par
\par\medskip\noindent\textbf{Definition} (\texttt{tau-op}).\;$\tau_{\operatorname{tds\_}} :=$
\begin{equation*}
(\lambda \operatorname{tda\_} \in L({\operatorname{tdv\_}}, {\operatorname{tdv\_}}).\; \mathrm{tr}\,{{\operatorname{tda\_}}{\operatorname{tds\_}}})
\end{equation*}

\par\medskip\noindent\textbf{Definition} (\texttt{functional-star}).\;${\operatorname{tdf\_}}^{*} :=$
\begin{equation*}
(\lambda \operatorname{tda\_} \in L({\operatorname{tdv\_}}, {\operatorname{tdv\_}}).\; \operatorname{conjugate}(\operatorname{tdf\_}({\operatorname{tda\_}}^{*})))
\end{equation*}

\par\medskip\noindent\textbf{Definition} (\texttt{is-self-adjoint-functional}).\;For all $\operatorname{tdv\_}$ and $\operatorname{tdf\_}$, 
\begin{equation*}
\begin{array}{@{}l@{}} \operatorname{is\text{-}self\text{-}adjoint\text{-}functional}(\operatorname{tdv\_}, \operatorname{tdf\_}) \\ \Leftrightarrow\; \forall \operatorname{tda\_} \in L({\operatorname{tdv\_}}, {\operatorname{tdv\_}}). \\ \quad \operatorname{conjugate}(\operatorname{tdf\_}({\operatorname{tda\_}}^{*})) = \operatorname{tdf\_}(\operatorname{tda\_}) \end{array}
\end{equation*}

\par\medskip\noindent\textbf{Definition} (\texttt{is-positive-op-functional}).\;For all $\operatorname{tdv\_}$ and $\operatorname{tdf\_}$, 
\begin{equation*}
\begin{array}{@{}l@{}} \operatorname{is\text{-}positive\text{-}op\text{-}functional}(\operatorname{tdv\_}, \operatorname{tdf\_}) \\ \Leftrightarrow\; \forall \operatorname{tda\_}. \\ \quad \operatorname{is\text{-}positive\text{-}op}(\operatorname{tdv\_}, \operatorname{tda\_}) \Rightarrow ((\operatorname{tdf\_}(\operatorname{tda\_}) \in \mathbb{R}) \wedge (0 \leq \operatorname{tdf\_}(\operatorname{tda\_}))) \end{array}
\end{equation*}


\vnbbookfile{theorem-library/ba-trace-dual}{theorem-library/ba-trace-dual}
\noindent{\small ba-trace-dual.scm -- BA-1117-TRACE, 2026-11-17.  BanachAlgebras.pdf ch. 5 s.6-s.8: the trace functional tau\_S (Def 5.25), its adjoint (Prop 5.32), self-adjoint and positive trace functionals (Prop 5.33), the module structure of L\_*(H) (eqs. 145-146), and the trace of a rank-one operator. Vocabulary: ...}\par
\vnbprop{tau-op-unfold}{\vnbstmt{\begin{array}{@{}l@{}} \forall v, s. \\ \quad (\tau_{s} \simeq (\lambda a \in L({v}, {v}).\; \mathrm{tr}\,{{a}{s}})) \end{array}}}
\vnbprop{functional-star-unfold}{\vnbstmt{\begin{array}{@{}l@{}} \forall v, f. \\ \quad ({f}^{*} \simeq (\lambda a \in L({v}, {v}).\; \operatorname{conjugate}(f({a}^{*})))) \end{array}}}
\vnbprop{tau-op-apply}{\vnbstmt{\begin{array}{@{}l@{}} \forall v, s, a. \\ \quad (\operatorname{is\text{-}hilbert\text{-}space}(v) \wedge (\operatorname{is\text{-}trace\text{-}class}(v, s) \wedge (a \in L({v}, {v})))) \Rightarrow  \\ \quad \quad \tau_{s}(a) = \mathrm{tr}\,{{a}{s}} \end{array}}}
\vnbprop{op-trace-rank-one}{\vnbstmt{\begin{array}{@{}l@{}} \forall v, a, c. \\ \quad (\operatorname{is\text{-}hilbert\text{-}space}(v) \wedge ((a \in \operatorname{vec}(v)) \wedge (c \in \operatorname{vec}(v)))) \Rightarrow  \\ \quad \quad \mathrm{tr}\,{\langle \cdot \mid {a} \rangle {c}} = \langle\cdot,\cdot\rangle_{v}(c, a) \end{array}}}
\vnbprop{tau-op-type}{\vnbstmt{\begin{array}{@{}l@{}} \forall v, s. \\ \quad (\operatorname{is\text{-}hilbert\text{-}space}(v) \wedge \operatorname{is\text{-}trace\text{-}class}(v, s)) \Rightarrow  \\ \quad \quad \tau_{s} \in (L({v}, {v}) \to \mathbb{C}) \end{array}}}
\vnbprop{tau-op-star}{\vnbstmt{\begin{array}{@{}l@{}} \forall v, s, a. \\ \quad (\operatorname{is\text{-}hilbert\text{-}space}(v) \wedge (\operatorname{is\text{-}trace\text{-}class}(v, s) \wedge (a \in L({v}, {v})))) \Rightarrow  \\ \quad \quad \operatorname{conjugate}(\tau_{s}({a}^{*})) = \tau_{{s}^{*}}(a) \end{array}}}
\vnbprop{tau-op-star-functional}{\vnbstmt{\begin{array}{@{}l@{}} \forall v, s, a. \\ \quad (\operatorname{is\text{-}hilbert\text{-}space}(v) \wedge (\operatorname{is\text{-}trace\text{-}class}(v, s) \wedge (a \in L({v}, {v})))) \Rightarrow  \\ \quad \quad {\tau_{s}}^{*}(a) = \tau_{{s}^{*}}(a) \end{array}}}
\vnbprop{tau-op-module-right}{\vnbstmt{\begin{array}{@{}l@{}} \forall v, s, t, a. \\ \quad (\operatorname{is\text{-}hilbert\text{-}space}(v) \wedge ((s \in L({v}, {v})) \wedge (\operatorname{is\text{-}trace\text{-}class}(v, t) \wedge (a \in L({v}, {v}))))) \Rightarrow  \\ \quad \quad \tau_{t}({a}{s}) = \tau_{{s}{t}}(a) \end{array}}}
\vnbprop{tau-op-module-left}{\vnbstmt{\begin{array}{@{}l@{}} \forall v, s, t, a. \\ \quad (\operatorname{is\text{-}hilbert\text{-}space}(v) \wedge ((s \in L({v}, {v})) \wedge (\operatorname{is\text{-}trace\text{-}class}(v, t) \wedge (a \in L({v}, {v}))))) \Rightarrow  \\ \quad \quad \tau_{t}({s}{a}) = \tau_{{t}{s}}(a) \end{array}}}
\vnbprop{op-trace-comp-rank-one}{\vnbstmt{\begin{array}{@{}l@{}} \forall v, s, a, c. \\ \quad (\operatorname{is\text{-}hilbert\text{-}space}(v) \wedge ((s \in L({v}, {v})) \wedge ((a \in \operatorname{vec}(v)) \wedge (c \in \operatorname{vec}(v))))) \Rightarrow  \\ \quad \quad \mathrm{tr}\,{{s}{\langle \cdot \mid {a} \rangle {c}}} = \langle\cdot,\cdot\rangle_{v}(s(c), a) \end{array}}}
\vnbprop{op-trace-nonneg-basis}{\vnbstmt{\begin{array}{@{}l@{}} \forall v, t, x, e. \\ \quad \operatorname{is\text{-}hilbert\text{-}space}(v) \wedge \\ \quad \operatorname{is\text{-}trace\text{-}class}(v, t) \wedge \\ \quad \operatorname{is\text{-}orthonormal\text{-}basis}(v, x, e) \wedge \\ \quad \forall k \in x.\; ((\langle\cdot,\cdot\rangle_{v}(t(e(k)), e(k)) \in \mathbb{R}) \wedge (0 \leq \langle\cdot,\cdot\rangle_{v}(t(e(k)), e(k)))) \Rightarrow \\ \quad \quad (\mathrm{tr}\,{t} \in \mathbb{R}) \wedge (0 \leq \mathrm{tr}\,{t}) \end{array}}}
\vnbprop{op-trace-positive}{\vnbstmt{\begin{array}{@{}l@{}} \forall v, t. \\ \quad (\operatorname{is\text{-}hilbert\text{-}space}(v) \wedge (\operatorname{is\text{-}trace\text{-}class}(v, t) \wedge \operatorname{is\text{-}positive\text{-}op}(v, t))) \Rightarrow  \\ \quad \quad (\mathrm{tr}\,{t} \in \mathbb{R}) \wedge (0 \leq \mathrm{tr}\,{t}) \end{array}}}
\vnbprop{tau-op-positive}{\vnbstmt{\begin{array}{@{}l@{}} \forall v, s. \\ \quad (\operatorname{is\text{-}hilbert\text{-}space}(v) \wedge (\operatorname{is\text{-}trace\text{-}class}(v, s) \wedge \operatorname{is\text{-}positive\text{-}op}(v, s))) \Rightarrow  \\ \quad \quad \operatorname{is\text{-}positive\text{-}op\text{-}functional}(v, \tau_{s}) \end{array}}}
\vnbprop{self-adjoint-op-adjoint-eq}{\vnbstmt{\begin{array}{@{}l@{}} \forall v, t. \\ \quad (\operatorname{is\text{-}hilbert\text{-}space}(v) \wedge \operatorname{is\text{-}self\text{-}adjoint\text{-}op}(v, t)) \Rightarrow ({t}^{*} = t) \end{array}}}
\vnbprop{tau-op-self-adjoint-of-op}{\vnbstmt{\begin{array}{@{}l@{}} \forall v, s. \\ \quad (\operatorname{is\text{-}hilbert\text{-}space}(v) \wedge (\operatorname{is\text{-}trace\text{-}class}(v, s) \wedge \operatorname{is\text{-}self\text{-}adjoint\text{-}op}(v, s))) \Rightarrow  \\ \quad \quad \operatorname{is\text{-}self\text{-}adjoint\text{-}functional}(v, \tau_{s}) \end{array}}}
\vnbprop{self-adjoint-op-of-tau}{\vnbstmt{\begin{array}{@{}l@{}} \forall v, s. \\ \quad (\operatorname{is\text{-}hilbert\text{-}space}(v) \wedge (\operatorname{is\text{-}trace\text{-}class}(v, s) \wedge \operatorname{is\text{-}self\text{-}adjoint\text{-}functional}(v, \tau_{s}))) \Rightarrow  \\ \quad \quad \operatorname{is\text{-}self\text{-}adjoint\text{-}op}(v, s) \end{array}}}
\vnbprop{tau-op-self-adjoint-iff}{\vnbstmt{\begin{array}{@{}l@{}} \forall v, s. \\ \quad (\operatorname{is\text{-}hilbert\text{-}space}(v) \wedge \operatorname{is\text{-}trace\text{-}class}(v, s)) \Rightarrow  \\ \quad \quad \operatorname{is\text{-}self\text{-}adjoint\text{-}functional}(v, \tau_{s}) \Leftrightarrow \operatorname{is\text{-}self\text{-}adjoint\text{-}op}(v, s) \end{array}}}
\vnbprop{self-adjoint-of-real-form}{\vnbstmt{\begin{array}{@{}l@{}} \forall v.; \forall t \in L({v}, {v}). \\ \quad (\operatorname{is\text{-}hilbert\text{-}space}(v) \wedge \forall x \in \operatorname{vec}(v).\; (\langle\cdot,\cdot\rangle_{v}(t(x), x) \in \mathbb{R})) \Rightarrow  \\ \quad \quad \operatorname{is\text{-}self\text{-}adjoint\text{-}op}(v, t) \end{array}}}
\vnbprop{rank-one-op-positive}{\vnbstmt{\begin{array}{@{}l@{}} \forall v.; \forall a \in \operatorname{vec}(v). \\ \quad \operatorname{is\text{-}hilbert\text{-}space}(v) \Rightarrow \operatorname{is\text{-}positive\text{-}op}(v, \langle \cdot \mid {a} \rangle {a}) \end{array}}}
\vnbprop{positive-op-of-tau}{\vnbstmt{\begin{array}{@{}l@{}} \forall v, s. \\ \quad (\operatorname{is\text{-}hilbert\text{-}space}(v) \wedge (\operatorname{is\text{-}trace\text{-}class}(v, s) \wedge \operatorname{is\text{-}positive\text{-}op\text{-}functional}(v, \tau_{s}))) \Rightarrow  \\ \quad \quad \operatorname{is\text{-}positive\text{-}op}(v, s) \end{array}}}
\vnbprop{tau-op-positive-iff}{\vnbstmt{\begin{array}{@{}l@{}} \forall v, s. \\ \quad (\operatorname{is\text{-}hilbert\text{-}space}(v) \wedge \operatorname{is\text{-}trace\text{-}class}(v, s)) \Rightarrow  \\ \quad \quad \operatorname{is\text{-}positive\text{-}op\text{-}functional}(v, \tau_{s}) \Leftrightarrow \operatorname{is\text{-}positive\text{-}op}(v, s) \end{array}}}
\vnbprop{onb-coef-separates}{\vnbstmt{\begin{array}{@{}l@{}} \forall v, x, e, y, z. \\ \quad \operatorname{is\text{-}hilbert\text{-}space}(v) \wedge \\ \quad \operatorname{is\text{-}orthonormal\text{-}basis}(v, x, e) \wedge \\ \quad y \in \operatorname{vec}(v) \wedge \\ \quad z \in \operatorname{vec}(v) \wedge \\ \quad \forall k \in x.\; (\langle\cdot,\cdot\rangle_{v}(y, e(k)) = \langle\cdot,\cdot\rangle_{v}(z, e(k))) \Rightarrow \\ \quad \quad y = z \end{array}}}

\vnbbookfile{theorem-library/ba-graph-commutation}{theorem-library/ba-graph-commutation}
\noindent{\small btg-graphs.scm -- BA-1111-A, 2026-11-11.  BanachAlgebras.pdf ch. 3 s.5-6 for graphs (decision D4: an unbounded operator is its graph G, a subset of H x H; ba-unbdd-vocab).}\par
\vnbprop{graph-compose-pair-iff}{\vnbstmt{\begin{array}{@{}l@{}} \forall v, g, h, a, b. \\ \quad (\operatorname{is\text{-}complex\text{-}inner\text{-}product\text{-}space}(v) \wedge ((a \in \operatorname{vec}(v)) \wedge (b \in \operatorname{vec}(v)))) \Rightarrow  \\ \quad \quad ({[}a, b] \in {g}{h}) \Leftrightarrow \exists y \in \operatorname{vec}(v).\; (({[}a, y] \in h) \wedge ({[}y, b] \in g)) \end{array}}}
\vnbprop{graph-adjoint-compose-pair-iff}{\vnbstmt{\begin{array}{@{}l@{}} \forall v, g, a, b. \\ \quad \operatorname{is\text{-}complex\text{-}inner\text{-}product\text{-}space}(v) \wedge \\ \quad g \subseteq (\operatorname{vec}(v) \times \operatorname{vec}(v)) \wedge \\ \quad a \in \operatorname{vec}(v) \wedge \\ \quad b \in \operatorname{vec}(v) \Rightarrow \\ \quad \quad {[}a, b] \in {{g}^{*}}{g} \\ \quad \quad \Leftrightarrow\; \exists z \in \operatorname{vec}(v). \\ \quad \quad \quad {[}a, z] \in g \wedge \\ \quad \quad \quad \forall c, d \in \operatorname{vec}(v). \\ \quad \quad \quad \quad ({[}c, d] \in g) \Rightarrow (\langle\cdot,\cdot\rangle_{v}(c, b) = \langle\cdot,\cdot\rangle_{v}(d, z)) \end{array}}}
\vnbprop{commutes-adjoint-compose-iff}{\vnbstmt{\begin{array}{@{}l@{}} \forall v, g, s. \\ \quad (\operatorname{is\text{-}complex\text{-}inner\text{-}product\text{-}space}(v) \wedge ((g \subseteq (\operatorname{vec}(v) \times \operatorname{vec}(v))) \wedge (s \in L({v}, {v})))) \Rightarrow  \\ \quad \quad \forall a, b \in \operatorname{vec}(v). \\ \quad \quad \quad ({[}a, b] \in {{g}^{*}}{g}) \Rightarrow  \\ \quad \quad \quad \quad {[}s(a), s(b)] \in {{g}^{*}}{g} \\ \quad \quad \Leftrightarrow\; \forall a, b \in \operatorname{vec}(v). \\ \quad \quad \quad \exists z \in \operatorname{vec}(v). \\ \quad \quad \quad \quad {[}a, z] \in g \wedge \\ \quad \quad \quad \quad \forall c, d \in \operatorname{vec}(v). \\ \quad \quad \quad \quad \quad ({[}c, d] \in g) \Rightarrow (\langle\cdot,\cdot\rangle_{v}(c, b) = \langle\cdot,\cdot\rangle_{v}(d, z)) \Rightarrow \\ \quad \quad \quad \quad \exists z \in \operatorname{vec}(v). \\ \quad \quad \quad \quad \quad {[}s(a), z] \in g \wedge \\ \quad \quad \quad \quad \quad \forall c, d \in \operatorname{vec}(v). \\ \quad \quad \quad \quad \quad \quad ({[}c, d] \in g) \Rightarrow (\langle\cdot,\cdot\rangle_{v}(c, s(b)) = \langle\cdot,\cdot\rangle_{v}(d, z)) \end{array}}}
\vnbprop{commutes-adjoint-compose-of}{\vnbstmt{\begin{array}{@{}l@{}} \forall v, g, r, s.; \forall a, b \in \operatorname{vec}(v). \\ \quad \operatorname{is\text{-}hilbert\text{-}space}(v) \wedge \\ \quad g \subseteq (\operatorname{vec}(v) \times \operatorname{vec}(v)) \wedge \\ \quad r \in L({v}, {v}) \wedge \\ \quad s \in L({v}, {v}) \wedge \\ \quad \forall c, d \in \operatorname{vec}(v). \\ \quad \quad ({[}c, d] \in g) \Rightarrow  \\ \quad \quad \quad ({[}s(c), r(d)] \in g) \wedge ({[}{s}^{*}(c), {r}^{*}(d)] \in g) \wedge \\ \quad {[}a, b] \in {{g}^{*}}{g} \Rightarrow \\ \quad \quad {[}s(a), s(b)] \in {{g}^{*}}{g} \end{array}}}
\vnbprop{graph-of-map-pair-iff}{\vnbstmt{\begin{array}{@{}l@{}} \forall v, d, f, a, b. \\ \quad (\operatorname{is\text{-}complex\text{-}inner\text{-}product\text{-}space}(v) \wedge ((a \in \operatorname{vec}(v)) \wedge (b \in \operatorname{vec}(v)))) \Rightarrow  \\ \quad \quad {[}a, b] \in \{p \in (\operatorname{vec}(v) \times \operatorname{vec}(v)) : ((\operatorname{nth}(1, p) \in d) \wedge (\operatorname{nth}(2, p) = f(\operatorname{nth}(1, p))))\} \\ \quad \quad \Leftrightarrow\; (a \in d) \wedge (b = f(a)) \end{array}}}
\vnbprop{form-operator-pair-iff}{\vnbstmt{\begin{array}{@{}l@{}} \forall v, d, q, a, b. \\ \quad (\operatorname{is\text{-}complex\text{-}inner\text{-}product\text{-}space}(v) \wedge ((a \in \operatorname{vec}(v)) \wedge (b \in \operatorname{vec}(v)))) \Rightarrow  \\ \quad \quad ({[}a, b] \in T_{q}) \Leftrightarrow ((a \in d) \wedge \forall x \in d.\; (\langle\cdot,\cdot\rangle_{v}(x, b) = q(x, a))) \end{array}}}
\vnbprop{graph-compose-subset}{\vnbstmt{\begin{array}{@{}l@{}} \forall v, g, h. \\ \quad \operatorname{is\text{-}complex\text{-}inner\text{-}product\text{-}space}(v) \Rightarrow  \\ \quad \quad {g}{h} \subseteq (\operatorname{vec}(v) \times \operatorname{vec}(v)) \end{array}}}
\vnbprop{form-operator-subset}{\vnbstmt{\begin{array}{@{}l@{}} \forall v, d, q. \\ \quad \operatorname{is\text{-}complex\text{-}inner\text{-}product\text{-}space}(v) \Rightarrow  \\ \quad \quad T_{q} \subseteq (\operatorname{vec}(v) \times \operatorname{vec}(v)) \end{array}}}
\vnbprop{adjoint-compose-form-operator-pair-iff}{\vnbstmt{\begin{array}{@{}l@{}} \forall v, d, f, c, \delta. \\ \quad \operatorname{is\text{-}complex\text{-}inner\text{-}product\text{-}space}(v) \wedge \\ \quad d \subseteq \operatorname{vec}(v) \wedge \\ \quad f \in (d \to \operatorname{vec}(v)) \wedge \\ \quad c \in \operatorname{vec}(v) \wedge \\ \quad \delta \in \operatorname{vec}(v) \Rightarrow \\ \quad \quad {[}c, \delta] \\ \quad \quad \in\; {{\{p \in (\operatorname{vec}(v) \times \operatorname{vec}(v)) : ((\operatorname{nth}(1, p) \in d) \wedge (\operatorname{nth}(2, p) = f(\operatorname{nth}(1, p))))\}}^{*}}{\{p \in (\operatorname{vec}(v) \times \operatorname{vec}(v)) : ((\operatorname{nth}(1, p) \in d) \wedge (\operatorname{nth}(2, p) = f(\operatorname{nth}(1, p))))\}} \\ \quad \quad \Leftrightarrow\; {[}c, \delta] \\ \quad \quad \in\; T_{(\lambda {[}\operatorname{bgx\_}, \operatorname{bgy\_}] \in (d \times d).\; \langle\cdot,\cdot\rangle_{v}(f(\operatorname{bgx\_}), f(\operatorname{bgy\_})))} \end{array}}}
\vnbprop{adjoint-compose-form-operator-pointwise}{\vnbstmt{\begin{array}{@{}l@{}} \forall v, d, f.; \forall c, \mathit{ugd} \in \operatorname{vec}(v). \\ \quad (\operatorname{is\text{-}complex\text{-}inner\text{-}product\text{-}space}(v) \wedge ((d \subseteq \operatorname{vec}(v)) \wedge (f \in (d \to \operatorname{vec}(v))))) \Rightarrow  \\ \quad \quad {[}c, \mathit{ugd}] \\ \quad \quad \in\; {{\{p \in (\operatorname{vec}(v) \times \operatorname{vec}(v)) : ((\operatorname{nth}(1, p) \in d) \wedge (\operatorname{nth}(2, p) = f(\operatorname{nth}(1, p))))\}}^{*}}{\{p \in (\operatorname{vec}(v) \times \operatorname{vec}(v)) : ((\operatorname{nth}(1, p) \in d) \wedge (\operatorname{nth}(2, p) = f(\operatorname{nth}(1, p))))\}} \\ \quad \quad \Leftrightarrow\; {[}c, \mathit{ugd}] \\ \quad \quad \in\; T_{(\lambda {[}\operatorname{bgx\_}, \operatorname{bgy\_}] \in (d \times d).\; \langle\cdot,\cdot\rangle_{v}(f(\operatorname{bgx\_}), f(\operatorname{bgy\_})))} \end{array}}}
\vnbprop{adjoint-compose-form-operator}{\vnbstmt{\begin{array}{@{}l@{}} \forall v, d, f. \\ \quad (\operatorname{is\text{-}complex\text{-}inner\text{-}product\text{-}space}(v) \wedge ((d \subseteq \operatorname{vec}(v)) \wedge (f \in (d \to \operatorname{vec}(v))))) \Rightarrow  \\ \quad \quad {{\{p \in (\operatorname{vec}(v) \times \operatorname{vec}(v)) : ((\operatorname{nth}(1, p) \in d) \wedge (\operatorname{nth}(2, p) = f(\operatorname{nth}(1, p))))\}}^{*}}{\{p \in (\operatorname{vec}(v) \times \operatorname{vec}(v)) : ((\operatorname{nth}(1, p) \in d) \wedge (\operatorname{nth}(2, p) = f(\operatorname{nth}(1, p))))\}} \\ \quad \quad =\; T_{(\lambda {[}\operatorname{bgx\_}, \operatorname{bgy\_}] \in (d \times d).\; \langle\cdot,\cdot\rangle_{v}(f(\operatorname{bgx\_}), f(\operatorname{bgy\_})))} \end{array}}}

\vnbbookfile{structure-library/ba-rep-vocab}{structure-library/ba-rep-vocab}
\noindent{\small ba-rep-vocab.scm -- POSITIVE FUNCTIONALS, STATES AND REPRESENTATIONS of a unital involutive Banach algebra: F. J. Thayer, "Banach Algebras and Applications" (docs/BanachAlgebras.pdf), chapter 4 (Definitions 4.1, 4.17, 4.19, 4.22, formula (86)) and Definition 1.57. DEFINITIONS ONLY: no proof, no support, no axiom other ...}\par
\par\medskip\noindent\textbf{Definition} (\texttt{states}).\;$\Sigma({\operatorname{g4a\_}}) :=$
\begin{equation*}
\{\operatorname{g4f\_} \in (\operatorname{vec}(\operatorname{g4a\_}) \to \mathbb{C}) : (\operatorname{is\text{-}positive\text{-}lf}(\operatorname{g4a\_}, \operatorname{g4f\_}) \wedge (\operatorname{g4f\_}(1_{\operatorname{g4a\_}}) = 1))\}
\end{equation*}

\par\medskip\noindent\textbf{Definition} (\texttt{is-positive-lf}).\;For all $\operatorname{g4a\_}$ and $\operatorname{g4f\_}$, 
\begin{equation*}
\begin{array}{@{}l@{}} \operatorname{is\text{-}positive\text{-}lf}(\operatorname{g4a\_}, \operatorname{g4f\_}) \\ \Leftrightarrow\; \operatorname{is\text{-}cc\text{-}linear\text{-}on}(\operatorname{g4a\_}, \operatorname{g4f\_}) \wedge \\ \forall \operatorname{g4x\_} \in \operatorname{vec}(\operatorname{g4a\_}). \\ \quad \operatorname{g4f\_}(\cdot_{\operatorname{g4a\_}}({}^{*_{\operatorname{g4a\_}}}(\operatorname{g4x\_}), \operatorname{g4x\_})) \in \mathbb{R} \wedge \\ \quad 0 \leq \operatorname{g4f\_}(\cdot_{\operatorname{g4a\_}}({}^{*_{\operatorname{g4a\_}}}(\operatorname{g4x\_}), \operatorname{g4x\_})) \end{array}
\end{equation*}

\par\medskip\noindent\textbf{Definition} (\texttt{functional-le}).\;For all $\operatorname{g4a\_}$, $\operatorname{g4g\_}$ and $\operatorname{g4f\_}$, 
\begin{equation*}
\begin{array}{@{}l@{}} \operatorname{functional\text{-}le}(\operatorname{g4a\_}, \operatorname{g4g\_}, \operatorname{g4f\_}) \\ \Leftrightarrow\; \forall \operatorname{g4x\_} \in \operatorname{vec}(\operatorname{g4a\_}). \\ \quad \operatorname{g4g\_}(\cdot_{\operatorname{g4a\_}}({}^{*_{\operatorname{g4a\_}}}(\operatorname{g4x\_}), \operatorname{g4x\_})) \leq \operatorname{g4f\_}(\cdot_{\operatorname{g4a\_}}({}^{*_{\operatorname{g4a\_}}}(\operatorname{g4x\_}), \operatorname{g4x\_})) \end{array}
\end{equation*}

\par\medskip\noindent\textbf{Definition} (\texttt{is-pure-state}).\;For all $\operatorname{g4a\_}$ and $\operatorname{g4f\_}$, 
\begin{equation*}
\begin{array}{@{}l@{}} \operatorname{is\text{-}pure\text{-}state}(\operatorname{g4a\_}, \operatorname{g4f\_}) \\ \Leftrightarrow\; \operatorname{g4f\_} \in \Sigma({\operatorname{g4a\_}}) \wedge \\ \forall \operatorname{g4g\_}. \\ \quad \operatorname{is\text{-}positive\text{-}lf}(\operatorname{g4a\_}, \operatorname{g4g\_}) \Rightarrow  \\ \quad \operatorname{functional\text{-}le}(\operatorname{g4a\_}, \operatorname{g4g\_}, \operatorname{g4f\_}) \Rightarrow  \\ \quad \quad \exists \operatorname{g4l\_} \in \mathbb{R}. \\ \quad \quad \quad ((0 \leq \operatorname{g4l\_}) \wedge ((\operatorname{g4l\_} \leq 1) \wedge \forall \operatorname{g4x\_} \in \operatorname{vec}(\operatorname{g4a\_}).\; (\operatorname{g4g\_}(\operatorname{g4x\_}) = \cdot(\operatorname{g4l\_}, \operatorname{g4f\_}(\operatorname{g4x\_}))))) \end{array}
\end{equation*}

\par\medskip\noindent\textbf{Definition} (\texttt{is-representation}).\;For all $\operatorname{g4a\_}$, $\operatorname{g4v\_}$ and $\operatorname{g4p\_}$, 
\begin{equation*}
\begin{array}{@{}l@{}} \operatorname{is\text{-}representation}(\operatorname{g4a\_}, \operatorname{g4v\_}, \operatorname{g4p\_}) \\ \Leftrightarrow\; \operatorname{is\text{-}hilbert\text{-}space}(\operatorname{g4v\_}) \wedge \operatorname{is\text{-}star\text{-}hom}(\operatorname{g4a\_}, L({\operatorname{g4v\_}}), \operatorname{g4p\_}) \end{array}
\end{equation*}

\par\medskip\noindent\textbf{Definition} (\texttt{rep-functional}).\;$f_{{\operatorname{g4p\_}},{\operatorname{g4e\_}}} :=$
\begin{equation*}
(\lambda \operatorname{g4x\_} \in \operatorname{vec}(\operatorname{g4a\_}).\; \langle\cdot,\cdot\rangle_{\operatorname{g4v\_}}(\operatorname{g4p\_}(\operatorname{g4x\_})(\operatorname{g4e\_}), \operatorname{g4e\_}))
\end{equation*}

\par\medskip\noindent\textbf{Definition} (\texttt{rep-orbit}).\;$\operatorname{rep\text{-}orbit}(\operatorname{g4a\_}, \operatorname{g4v\_}, \operatorname{g4p\_}, \operatorname{g4e\_}) :=$
\begin{equation*}
\{\operatorname{g4y\_} \in \operatorname{vec}(\operatorname{g4v\_}) : \exists \operatorname{g4x\_} \in \operatorname{vec}(\operatorname{g4a\_}).\; (\operatorname{g4y\_} = \operatorname{g4p\_}(\operatorname{g4x\_})(\operatorname{g4e\_}))\}
\end{equation*}

\par\medskip\noindent\textbf{Definition} (\texttt{is-cyclic-vector-rep}).\;For all $\operatorname{g4a\_}$, $\operatorname{g4v\_}$, $\operatorname{g4p\_}$ and $\operatorname{g4e\_}$, 
\begin{equation*}
\begin{array}{@{}l@{}} \operatorname{is\text{-}cyclic\text{-}vector\text{-}rep}(\operatorname{g4a\_}, \operatorname{g4v\_}, \operatorname{g4p\_}, \operatorname{g4e\_}) \\ \Leftrightarrow\; ((\operatorname{g4e\_} \in \operatorname{vec}(\operatorname{g4v\_})) \wedge (\operatorname{closure}(\operatorname{ip\text{-}metric\text{-}space}(\operatorname{g4v\_}), \operatorname{rep\text{-}orbit}(\operatorname{g4a\_}, \operatorname{g4v\_}, \operatorname{g4p\_}, \operatorname{g4e\_})) = \operatorname{vec}(\operatorname{g4v\_}))) \end{array}
\end{equation*}

\par\medskip\noindent\textbf{Definition} (\texttt{is-invariant-subspace-rep}).\;For all $\operatorname{g4a\_}$, $\operatorname{g4v\_}$, $\operatorname{g4p\_}$ and $\operatorname{g4k\_}$, 
\begin{equation*}
\begin{array}{@{}l@{}} \operatorname{is\text{-}invariant\text{-}subspace\text{-}rep}(\operatorname{g4a\_}, \operatorname{g4v\_}, \operatorname{g4p\_}, \operatorname{g4k\_}) \\ \Leftrightarrow\; \operatorname{is\text{-}submodule}(\operatorname{g4v\_}, \operatorname{g4k\_}) \wedge \\ \operatorname{is\text{-}closed}(\operatorname{ip\text{-}metric\text{-}space}(\operatorname{g4v\_}), \operatorname{g4k\_}) \wedge \\ \forall \operatorname{g4x\_} \in \operatorname{vec}(\operatorname{g4a\_}), \operatorname{g4y\_} \in \operatorname{g4k\_}.\; (\operatorname{g4p\_}(\operatorname{g4x\_})(\operatorname{g4y\_}) \in \operatorname{g4k\_}) \end{array}
\end{equation*}

\par\medskip\noindent\textbf{Definition} (\texttt{is-irreducible-rep}).\;For all $\operatorname{g4a\_}$, $\operatorname{g4v\_}$ and $\operatorname{g4p\_}$, 
\begin{equation*}
\begin{array}{@{}l@{}} \operatorname{is\text{-}irreducible\text{-}rep}(\operatorname{g4a\_}, \operatorname{g4v\_}, \operatorname{g4p\_}) \\ \Leftrightarrow\; \operatorname{is\text{-}representation}(\operatorname{g4a\_}, \operatorname{g4v\_}, \operatorname{g4p\_}) \wedge \\ \forall \operatorname{g4k\_}. \\ \quad \operatorname{is\text{-}invariant\text{-}subspace\text{-}rep}(\operatorname{g4a\_}, \operatorname{g4v\_}, \operatorname{g4p\_}, \operatorname{g4k\_}) \Rightarrow  \\ \quad \quad (\operatorname{g4k\_} = \operatorname{vec}(\operatorname{g4v\_})) \vee (\operatorname{g4k\_} = \operatorname{make\text{-}set}({[}\operatorname{vzero}(\operatorname{g4v\_})])) \end{array}
\end{equation*}

\par\medskip\noindent\textbf{Definition} (\texttt{half-sq-iter}).\;For all $\operatorname{g4a\_}$ and $\operatorname{g4y\_}$,  $(\operatorname{half\text{-}sq\text{-}iter}(\operatorname{g4a\_}, \operatorname{g4y\_}, 0) \simeq \operatorname{vzero}(\operatorname{g4a\_}))$\par
Let $\operatorname{g4a\_}$ and $\operatorname{g4y\_}$ be arbitrary. Let $\operatorname{g4n\_} \in \mathbb{N}$. Then
\begin{equation*}
(\operatorname{half\text{-}sq\text{-}iter}(\operatorname{g4a\_}, \operatorname{g4y\_}, \operatorname{succ}(\operatorname{g4n\_})) \simeq \operatorname{act}(\operatorname{g4a\_})(1/2, \operatorname{vadd}(\operatorname{g4a\_})(\operatorname{g4y\_}, \cdot_{\operatorname{g4a\_}}(\operatorname{half\text{-}sq\text{-}iter}(\operatorname{g4a\_}, \operatorname{g4y\_}, \operatorname{g4n\_}), \operatorname{half\text{-}sq\text{-}iter}(\operatorname{g4a\_}, \operatorname{g4y\_}, \operatorname{g4n\_})))))
\end{equation*}


\vnbbookfile{structure-library/ba-gns-vocab}{structure-library/ba-gns-vocab}
\noindent{\small ba-gns-vocab.scm -- THE GNS CONSTRUCTION, the pre-Hilbert stage: BanachAlgebras.pdf Theorem 4.13 and its proof.  Agent BA-1111-B, 2026-11-11.  Destination: structure-library/ba-gns-vocab.scm. Library: operator-algebras.  DEFINITIONS ONLY: no proof, no support, no axiom. Load window: after structure-library/ba-rep-vocab ...}\par
\par\medskip\noindent\textbf{Definition} (\texttt{gns-module}).\;$\operatorname{gns\text{-}module}(\operatorname{g4ga\_}) :=$ $\operatorname{uba\text{-}as\text{-}module}(\operatorname{uiba\text{-}as\text{-}uba}(\operatorname{g4ga\_}))$
\par\medskip\noindent\textbf{Definition} (\texttt{gns-form}).\;$\operatorname{gns\text{-}form}(\operatorname{g4ga\_}, \operatorname{g4gf\_}) :=$
\begin{equation*}
\begin{array}{@{}l@{}} (\lambda {[}\operatorname{g4gx\_}, \operatorname{g4gy\_}] \in (\operatorname{vec}(\operatorname{g4ga\_}) \times \operatorname{vec}(\operatorname{g4ga\_})).\; \\ \quad \operatorname{g4gf\_}(\cdot_{\operatorname{g4ga\_}}({}^{*_{\operatorname{g4ga\_}}}(\operatorname{g4gy\_}), \operatorname{g4gx\_}))) \end{array}
\end{equation*}

\par\medskip\noindent\textbf{Definition} (\texttt{gns-setoid}).\;$\operatorname{gns\text{-}setoid}(\operatorname{g4ga\_}, \operatorname{g4gf\_}) :=$
\begin{equation*}
\operatorname{qmod\text{-}setoid}(\operatorname{gns\text{-}module}(\operatorname{g4ga\_}), N_{\operatorname{gns\text{-}form}(\operatorname{g4ga\_}, \operatorname{g4gf\_})})
\end{equation*}

\par\medskip\noindent\textbf{Definition} (\texttt{gns-pre}).\;$\operatorname{gns\text{-}pre}(\operatorname{g4ga\_}, \operatorname{g4gf\_}) :=$ ${\operatorname{gns\text{-}module}(\operatorname{g4ga\_})}/N_{\operatorname{gns\text{-}form}(\operatorname{g4ga\_}, \operatorname{g4gf\_})}$
\par\medskip\noindent\textbf{Definition} (\texttt{gns-vector}).\;$\operatorname{gns\text{-}vector}(\operatorname{g4ga\_}, \operatorname{g4gf\_}) :=$ $\operatorname{class}(\operatorname{gns\text{-}setoid}(\operatorname{g4ga\_}, \operatorname{g4gf\_}), 1_{\operatorname{g4ga\_}})$
\par\medskip\noindent\textbf{Definition} (\texttt{gns-op}).\;$\operatorname{gns\text{-}op}(\operatorname{g4ga\_}, \operatorname{g4gf\_}, \operatorname{g4gx\_}) :=$
\begin{equation*}
\begin{array}{@{}l@{}} \operatorname{descend}(\operatorname{gns\text{-}setoid}(\operatorname{g4ga\_}, \operatorname{g4gf\_}), \\ \quad (\lambda \operatorname{g4gb\_} \in \operatorname{vec}(\operatorname{g4ga\_}).\; \operatorname{class}(\operatorname{gns\text{-}setoid}(\operatorname{g4ga\_}, \operatorname{g4gf\_}), \cdot_{\operatorname{g4ga\_}}(\operatorname{g4gx\_}, \operatorname{g4gb\_})))) \end{array}
\end{equation*}


\vnbbookfile{structure-library/ba-gns2-vocab}{structure-library/ba-gns2-vocab}
\noindent{\small ba-gns2-vocab.scm -- THE GNS CONSTRUCTION ON THE HILBERT COMPLETION and THE RADICAL: F. J. Thayer, "Banach Algebras and Applications" (docs/BanachAlgebras.pdf), Theorem 4.13 and Definition 4.25.  Agent BA-1117-GNS2, wave 2026-11-17. Destination: structure-library/ba-gns2-vocab.scm.  Library: operator-algebras.}\par
\par\medskip\noindent\textbf{Definition} (\texttt{radical}).\;$\sqrt{\operatorname{g4ha\_}} :=$
\begin{equation*}
\{\operatorname{g4hx\_} \in \operatorname{vec}(\operatorname{g4ha\_}) : \forall \operatorname{g4hv\_}, \operatorname{g4hq\_}.\; (\operatorname{is\text{-}representation}(\operatorname{g4ha\_}, \operatorname{g4hv\_}, \operatorname{g4hq\_}) \Rightarrow (\operatorname{g4hq\_}(\operatorname{g4hx\_}) = 0))\}
\end{equation*}

\par\medskip\noindent\textbf{Definition} (\texttt{seq-map}).\;$\operatorname{seq\text{-}map}(\operatorname{g4ht\_}, \operatorname{g4hf\_}) :=$ $(\lambda \operatorname{g4hn\_} \in \mathbb{N}.\; \operatorname{g4ht\_}(\operatorname{g4hf\_}(\operatorname{g4hn\_})))$
\par\medskip\noindent\textbf{Definition} (\texttt{completion-lift}).\;$\operatorname{completion\text{-}lift}(\operatorname{g4hm\_}, \operatorname{g4ht\_}) :=$
\begin{equation*}
\begin{array}{@{}l@{}} \operatorname{descend}(\operatorname{cauchy\text{-}setoid}(\operatorname{g4hm\_}), \\ \quad (\lambda \operatorname{g4hf\_} \in \operatorname{cseq}(\operatorname{g4hm\_}).\; \operatorname{class}(\operatorname{cauchy\text{-}setoid}(\operatorname{g4hm\_}), \operatorname{seq\text{-}map}(\operatorname{g4ht\_}, \operatorname{g4hf\_})))) \end{array}
\end{equation*}

\par\medskip\noindent\textbf{Definition} (\texttt{gns-hmod}).\;$\operatorname{gns\text{-}hmod}(\operatorname{g4ha\_}, \operatorname{g4hp\_}) :=$
\begin{equation*}
\operatorname{quotient\text{-}module}(\operatorname{gns\text{-}module}(\operatorname{g4ha\_}), N_{\operatorname{gns\text{-}form}(\operatorname{g4ha\_}, \operatorname{g4hp\_})})
\end{equation*}

\par\medskip\noindent\textbf{Definition} (\texttt{gns-hform}).\;$\operatorname{gns\text{-}hform}(\operatorname{g4ha\_}, \operatorname{g4hp\_}) :=$ $\operatorname{form\text{-}qip}(\operatorname{gns\text{-}module}(\operatorname{g4ha\_}), \operatorname{gns\text{-}form}(\operatorname{g4ha\_}, \operatorname{g4hp\_}))$
\par\medskip\noindent\textbf{Definition} (\texttt{gns-space}).\;$H_{\operatorname{g4hp\_}} :=$
\begin{equation*}
\operatorname{hilbert\text{-}completion}(\operatorname{gns\text{-}hmod}(\operatorname{g4ha\_}, \operatorname{g4hp\_}), \operatorname{gns\text{-}hform}(\operatorname{g4ha\_}, \operatorname{g4hp\_}))
\end{equation*}

\par\medskip\noindent\textbf{Definition} (\texttt{gns-rep}).\;$\pi_{\operatorname{g4hp\_}} :=$
\begin{equation*}
\begin{array}{@{}l@{}} (\lambda \operatorname{g4hx\_} \in \operatorname{vec}(\operatorname{g4ha\_}).\; \\ \quad \operatorname{completion\text{-}lift}(\operatorname{form\text{-}ms}(\operatorname{gns\text{-}hmod}(\operatorname{g4ha\_}, \operatorname{g4hp\_}), \operatorname{gns\text{-}hform}(\operatorname{g4ha\_}, \operatorname{g4hp\_})), \\ \quad \quad \operatorname{gns\text{-}op}(\operatorname{g4ha\_}, \operatorname{g4hp\_}, \operatorname{g4hx\_}))) \end{array}
\end{equation*}

\par\medskip\noindent\textbf{Definition} (\texttt{gns-xi}).\;$\xi_{\operatorname{g4hp\_}} :=$
\begin{equation*}
\operatorname{hc\text{-}embed}(\operatorname{gns\text{-}hmod}(\operatorname{g4ha\_}, \operatorname{g4hp\_}), \operatorname{gns\text{-}hform}(\operatorname{g4ha\_}, \operatorname{g4hp\_}))(\operatorname{gns\text{-}vector}(\operatorname{g4ha\_}, \operatorname{g4hp\_}))
\end{equation*}

\par\medskip\noindent\textbf{Definition} (\texttt{is-extreme-state}).\;For all $\operatorname{g4ha\_}$ and $\operatorname{g4hp\_}$, 
\begin{equation*}
\begin{array}{@{}l@{}} \operatorname{is\text{-}extreme\text{-}state}(\operatorname{g4ha\_}, \operatorname{g4hp\_}) \\ \Leftrightarrow\; \operatorname{g4hp\_} \in \Sigma({\operatorname{g4ha\_}}) \wedge \\ \forall \operatorname{g4hg\_}, \operatorname{g4hh\_} \in \Sigma({\operatorname{g4ha\_}}), \operatorname{g4hr\_} \in \mathbb{R}. \\ \quad (0 < \operatorname{g4hr\_}) \Rightarrow  \\ \quad (\operatorname{g4hr\_} < 1) \Rightarrow  \\ \quad \forall \operatorname{g4hy\_} \in \operatorname{vec}(\operatorname{g4ha\_}).\; (\operatorname{g4hp\_}(\operatorname{g4hy\_}) = +(\cdot(\operatorname{g4hr\_}, \operatorname{g4hg\_}(\operatorname{g4hy\_})), \cdot(-(1, \operatorname{g4hr\_}), \operatorname{g4hh\_}(\operatorname{g4hy\_})))) \Rightarrow  \\ \quad \quad \forall \operatorname{g4hz\_} \in \operatorname{vec}(\operatorname{g4ha\_}). \\ \quad \quad \quad (\operatorname{g4hg\_}(\operatorname{g4hz\_}) = \operatorname{g4hp\_}(\operatorname{g4hz\_})) \wedge (\operatorname{g4hh\_}(\operatorname{g4hz\_}) = \operatorname{g4hp\_}(\operatorname{g4hz\_})) \end{array}
\end{equation*}

\par\medskip\noindent\textbf{Definition} (\texttt{cdual-weak-star-top}).\;$\operatorname{cdual\text{-}weak\text{-}star\text{-}top}(\operatorname{g4ha\_}) :=$
\begin{equation*}
\begin{array}{@{}l@{}} \operatorname{subspace\text{-}top}(\operatorname{tproduct\text{-}top}(\operatorname{vec}(\operatorname{g4ha\_}), (\lambda \operatorname{bgx\_} \in \operatorname{vec}(\operatorname{g4ha\_}).\; \operatorname{metric\text{-}top}(\operatorname{cc\text{-}ms}))), \\ \quad \operatorname{cdual\text{-}vec}(\operatorname{g4ha\_})) \end{array}
\end{equation*}

\par\medskip\noindent\textbf{Definition} (\texttt{rfn-module}).\;$\operatorname{rfn\text{-}module}(\operatorname{g4hu\_}) :=$
\begin{equation*}
{[}\operatorname{normed\text{-}field\text{-}as\text{-}commutative\text{-}ring}(\operatorname{rr\text{-}normed\text{-}field}), (\operatorname{g4hu\_} \to \mathbb{C}), (\lambda {[}\operatorname{g4hf\_}, \operatorname{g4hg\_}] \in ((\operatorname{g4hu\_} \to \mathbb{C}) \times (\operatorname{g4hu\_} \to \mathbb{C})).\; (\lambda \operatorname{g4hy\_} \in \operatorname{g4hu\_}.\; (\operatorname{g4hf\_}(\operatorname{g4hy\_}) + \operatorname{g4hg\_}(\operatorname{g4hy\_})))), (\lambda \operatorname{g4hy\_} \in \operatorname{g4hu\_}.\; 0), (\lambda \operatorname{g4hf\_} \in (\operatorname{g4hu\_} \to \mathbb{C}).\; (\lambda \operatorname{g4hy\_} \in \operatorname{g4hu\_}.\; -\operatorname{g4hf\_}(\operatorname{g4hy\_}))), (\lambda {[}\operatorname{g4hr\_}, \operatorname{g4hf\_}] \in (\mathbb{R} \times (\operatorname{g4hu\_} \to \mathbb{C})).\; (\lambda \operatorname{g4hy\_} \in \operatorname{g4hu\_}.\; (\operatorname{g4hr\_} \cdot \operatorname{g4hf\_}(\operatorname{g4hy\_}))))]
\end{equation*}

\par\medskip\noindent\textbf{Definition} (\texttt{eval-seminorms}).\;$\operatorname{eval\text{-}seminorms}(\operatorname{g4hu\_}) :=$
\begin{equation*}
(\lambda \operatorname{g4hy\_} \in \operatorname{g4hu\_}.\; (\lambda \operatorname{g4hf\_} \in (\operatorname{g4hu\_} \to \mathbb{C}).\; \operatorname{magnitude}(\operatorname{g4hf\_}(\operatorname{g4hy\_}))))
\end{equation*}

\par\medskip\noindent\textbf{Definition} (\texttt{rfn-top}).\;$\operatorname{rfn\text{-}top}(\operatorname{g4hu\_}) :=$
\begin{equation*}
\operatorname{seminorm\text{-}top}(\operatorname{rfn\text{-}module}(\operatorname{g4hu\_}), \operatorname{g4hu\_}, \operatorname{eval\text{-}seminorms}(\operatorname{g4hu\_}))
\end{equation*}


\vnbbookfile{theorem-library/completion-lift}{theorem-library/completion-lift}
\noindent{\small completion-lift.scm -- THE COMPLETION OF A LIPSCHITZ MAP, and its Hilbert-space laws. Agent BA-1117-GNS2, 2026-11-17.  Destination: theorem-library/completion-lift.scm. Library: operator-algebras.  Load window: after hilbert-completion (hc-setoid-eq, hc-vec, form-ms-*), rake-setoid2 (descend-computes, class-eq-iff, the ...}\par
\vnbprop{seq-map-type}{\vnbstmt{\begin{array}{@{}l@{}} \forall d, e, q, f. \\ \quad ((q \in (d \to e)) \wedge (f \in (\mathbb{N} \to d))) \Rightarrow (\operatorname{seq\text{-}map}(q, f) \in (\mathbb{N} \to e)) \end{array}}}
\vnbprop{seq-map-at}{\vnbstmt{\begin{array}{@{}l@{}} \forall d, e, q, f, n. \\ \quad ((q \in (d \to e)) \wedge ((f \in (\mathbb{N} \to d)) \wedge (n \in \mathbb{N}))) \Rightarrow  \\ \quad \quad \operatorname{seq\text{-}map}(q, f)(n) = q(f(n)) \end{array}}}
\vnbprop{seq-map-cauchy}{\vnbstmt{\begin{array}{@{}l@{}} \forall q, a, b, f. \\ \quad \operatorname{is\text{-}metric\text{-}space}(q) \wedge \\ \quad a \in (\operatorname{pts}(q) \to \operatorname{pts}(q)) \wedge \\ \quad b \in \mathbb{R} \wedge \\ \quad 0 \leq b \wedge \\ \quad \forall u, w \in \operatorname{pts}(q).\; (\operatorname{dist}(q)(a(u), a(w)) \leq (b \cdot \operatorname{dist}(q)(u, w))) \wedge \\ \quad \operatorname{is\text{-}cauchy\text{-}seq}(q, f) \Rightarrow \\ \quad \quad \operatorname{is\text{-}cauchy\text{-}seq}(q, \operatorname{seq\text{-}map}(a, f)) \end{array}}}
\vnbprop{seq-map-equiv}{\vnbstmt{\begin{array}{@{}l@{}} \forall q, a, b, f, g. \\ \quad \operatorname{is\text{-}metric\text{-}space}(q) \wedge \\ \quad a \in (\operatorname{pts}(q) \to \operatorname{pts}(q)) \wedge \\ \quad b \in \mathbb{R} \wedge \\ \quad 0 \leq b \wedge \\ \quad \forall u, w \in \operatorname{pts}(q).\; (\operatorname{dist}(q)(a(u), a(w)) \leq (b \cdot \operatorname{dist}(q)(u, w))) \wedge \\ \quad f \in (\mathbb{N} \to \operatorname{pts}(q)) \wedge \\ \quad g \in (\mathbb{N} \to \operatorname{pts}(q)) \wedge \\ \quad \operatorname{cseq\text{-}equiv}(q, f, g) \Rightarrow \\ \quad \quad \operatorname{cseq\text{-}equiv}(q, \operatorname{seq\text{-}map}(a, f), \operatorname{seq\text{-}map}(a, g)) \end{array}}}
\vnbprop{completion-lift-ok}{\vnbstmt{\begin{array}{@{}l@{}} \forall q, a, b. \\ \quad \operatorname{is\text{-}metric\text{-}space}(q) \wedge \\ \quad a \in (\operatorname{pts}(q) \to \operatorname{pts}(q)) \wedge \\ \quad b \in \mathbb{R} \wedge \\ \quad 0 \leq b \wedge \\ \quad \forall u, w \in \operatorname{pts}(q).\; (\operatorname{dist}(q)(a(u), a(w)) \leq (b \cdot \operatorname{dist}(q)(u, w))) \Rightarrow \\ \quad \quad (\lambda f \in \operatorname{cseq}(q).\; \operatorname{class}(\operatorname{cauchy\text{-}setoid}(q), \operatorname{seq\text{-}map}(a, f))) \\ \quad \quad \in\; (\operatorname{pts}(\operatorname{cauchy\text{-}setoid}(q)) \to \operatorname{quotient}(\operatorname{cauchy\text{-}setoid}(q))) \wedge \\ \quad \quad \operatorname{respects}(\operatorname{cauchy\text{-}setoid}(q), (\lambda f \in \operatorname{cseq}(q).\; \operatorname{class}(\operatorname{cauchy\text{-}setoid}(q), \operatorname{seq\text{-}map}(a, f)))) \end{array}}}
\vnbprop{completion-lift-computes}{\vnbstmt{\begin{array}{@{}l@{}} \forall q, a, b.; \forall f \in \operatorname{cseq}(q). \\ \quad \operatorname{is\text{-}metric\text{-}space}(q) \wedge \\ \quad a \in (\operatorname{pts}(q) \to \operatorname{pts}(q)) \wedge \\ \quad b \in \mathbb{R} \wedge \\ \quad 0 \leq b \wedge \\ \quad \forall u, w \in \operatorname{pts}(q).\; (\operatorname{dist}(q)(a(u), a(w)) \leq (b \cdot \operatorname{dist}(q)(u, w))) \Rightarrow \\ \quad \quad \operatorname{completion\text{-}lift}(q, a)(\operatorname{class}(\operatorname{cauchy\text{-}setoid}(q), f)) = \operatorname{class}(\operatorname{cauchy\text{-}setoid}(q), \operatorname{seq\text{-}map}(a, f)) \end{array}}}
\vnbprop{completion-lift-type}{\vnbstmt{\begin{array}{@{}l@{}} \forall q, a, b. \\ \quad \operatorname{is\text{-}metric\text{-}space}(q) \wedge \\ \quad a \in (\operatorname{pts}(q) \to \operatorname{pts}(q)) \wedge \\ \quad b \in \mathbb{R} \wedge \\ \quad 0 \leq b \wedge \\ \quad \forall u, w \in \operatorname{pts}(q).\; (\operatorname{dist}(q)(a(u), a(w)) \leq (b \cdot \operatorname{dist}(q)(u, w))) \Rightarrow \\ \quad \quad \operatorname{completion\text{-}lift}(q, a) \in (\operatorname{quotient}(\operatorname{cauchy\text{-}setoid}(q)) \to \operatorname{quotient}(\operatorname{cauchy\text{-}setoid}(q))) \end{array}}}
\vnbprop{hc-vadd-classes}{\vnbstmt{\begin{array}{@{}l@{}} \forall d, q.; \forall f, g \in \operatorname{cseq}(\operatorname{form\text{-}ms}(d, q)). \\ \quad \operatorname{is\text{-}module}(d) \wedge \\ \quad \operatorname{scal}(d) = \operatorname{normed\text{-}field\text{-}as\text{-}commutative\text{-}ring}(\operatorname{cc\text{-}normed\text{-}field}) \wedge \\ \quad \operatorname{is\text{-}inner\text{-}product\text{-}form}(d, q) \Rightarrow \\ \quad \quad \operatorname{vadd}(\operatorname{hilbert\text{-}completion}(d, q))(\operatorname{class}(\operatorname{cauchy\text{-}setoid}(\operatorname{form\text{-}ms}(d, q)), f), \\ \quad \quad \quad \operatorname{class}(\operatorname{cauchy\text{-}setoid}(\operatorname{form\text{-}ms}(d, q)), g)) \\ \quad \quad =\; \operatorname{class}(\operatorname{cauchy\text{-}setoid}(\operatorname{form\text{-}ms}(d, q)), \operatorname{seq\text{-}vadd}(d, f, g)) \end{array}}}
\vnbprop{hc-act-classes}{\vnbstmt{\begin{array}{@{}l@{}} \forall d, q.; \forall c \in \mathbb{C}, f \in \operatorname{cseq}(\operatorname{form\text{-}ms}(d, q)). \\ \quad \operatorname{is\text{-}module}(d) \wedge \\ \quad \operatorname{scal}(d) = \operatorname{normed\text{-}field\text{-}as\text{-}commutative\text{-}ring}(\operatorname{cc\text{-}normed\text{-}field}) \wedge \\ \quad \operatorname{is\text{-}inner\text{-}product\text{-}form}(d, q) \Rightarrow \\ \quad \quad \operatorname{act}(\operatorname{hilbert\text{-}completion}(d, q))(c, \operatorname{class}(\operatorname{cauchy\text{-}setoid}(\operatorname{form\text{-}ms}(d, q)), f)) \\ \quad \quad =\; \operatorname{class}(\operatorname{cauchy\text{-}setoid}(\operatorname{form\text{-}ms}(d, q)), \operatorname{seq\text{-}act}(d, c, f)) \end{array}}}
\vnbprop{hc-ip-classes}{\vnbstmt{\begin{array}{@{}l@{}} \forall d, q.; \forall f, g \in \operatorname{cseq}(\operatorname{form\text{-}ms}(d, q)). \\ \quad \operatorname{is\text{-}module}(d) \wedge \\ \quad \operatorname{scal}(d) = \operatorname{normed\text{-}field\text{-}as\text{-}commutative\text{-}ring}(\operatorname{cc\text{-}normed\text{-}field}) \wedge \\ \quad \operatorname{is\text{-}inner\text{-}product\text{-}form}(d, q) \Rightarrow \\ \quad \quad \langle\cdot,\cdot\rangle_{\operatorname{hilbert\text{-}completion}(d, q)}(\operatorname{class}(\operatorname{cauchy\text{-}setoid}(\operatorname{form\text{-}ms}(d, q)), f), \\ \quad \quad \quad \operatorname{class}(\operatorname{cauchy\text{-}setoid}(\operatorname{form\text{-}ms}(d, q)), g)) \\ \quad \quad =\; \operatorname{limit\text{-}form}(d, q)(f, g) \end{array}}}
\vnbprop{hc-class-rep}{\vnbstmt{\begin{array}{@{}l@{}} \forall d, q.; \forall x \in \operatorname{vec}(\operatorname{hilbert\text{-}completion}(d, q)). \\ \quad \operatorname{is\text{-}module}(d) \wedge \\ \quad \operatorname{scal}(d) = \operatorname{normed\text{-}field\text{-}as\text{-}commutative\text{-}ring}(\operatorname{cc\text{-}normed\text{-}field}) \wedge \\ \quad \operatorname{is\text{-}inner\text{-}product\text{-}form}(d, q) \Rightarrow \\ \quad \quad \exists f \in \operatorname{cseq}(\operatorname{form\text{-}ms}(d, q)).\; (x = \operatorname{class}(\operatorname{cauchy\text{-}setoid}(\operatorname{form\text{-}ms}(d, q)), f)) \end{array}}}
\vnbprop{completion-lift-embed}{\vnbstmt{\begin{array}{@{}l@{}} \forall d, q, a, b.; \forall u \in \operatorname{vec}(d). \\ \quad \operatorname{is\text{-}module}(d) \wedge \\ \quad \operatorname{scal}(d) = \operatorname{normed\text{-}field\text{-}as\text{-}commutative\text{-}ring}(\operatorname{cc\text{-}normed\text{-}field}) \wedge \\ \quad \operatorname{is\text{-}inner\text{-}product\text{-}form}(d, q) \wedge \\ \quad a \in (\operatorname{pts}(\operatorname{form\text{-}ms}(d, q)) \to \operatorname{pts}(\operatorname{form\text{-}ms}(d, q))) \wedge \\ \quad b \in \mathbb{R} \wedge \\ \quad 0 \leq b \wedge \\ \quad \forall \mathit{g4u}, w \in \operatorname{pts}(\operatorname{form\text{-}ms}(d, q)). \\ \quad \quad \operatorname{dist}(\operatorname{form\text{-}ms}(d, q))(a(\mathit{g4u}), a(w)) \leq (b \cdot \operatorname{dist}(\operatorname{form\text{-}ms}(d, q))(\mathit{g4u}, w)) \Rightarrow \\ \quad \quad \operatorname{completion\text{-}lift}(\operatorname{form\text{-}ms}(d, q), a)(\operatorname{hc\text{-}embed}(d, q)(u)) = \operatorname{hc\text{-}embed}(d, q)(a(u)) \end{array}}}
\vnbprop{completion-lift-compose}{\vnbstmt{\begin{array}{@{}l@{}} \forall d, q, a, b, c, f, g, h.; \forall x \in \operatorname{vec}(\operatorname{hilbert\text{-}completion}(d, q)). \\ \quad \operatorname{is\text{-}module}(d) \wedge \\ \quad \operatorname{scal}(d) = \operatorname{normed\text{-}field\text{-}as\text{-}commutative\text{-}ring}(\operatorname{cc\text{-}normed\text{-}field}) \wedge \\ \quad \operatorname{is\text{-}inner\text{-}product\text{-}form}(d, q) \wedge \\ \quad a \in (\operatorname{pts}(\operatorname{form\text{-}ms}(d, q)) \to \operatorname{pts}(\operatorname{form\text{-}ms}(d, q))) \wedge \\ \quad b \in \mathbb{R} \wedge \\ \quad 0 \leq b \wedge \\ \quad \forall u, w \in \operatorname{pts}(\operatorname{form\text{-}ms}(d, q)). \\ \quad \quad \operatorname{dist}(\operatorname{form\text{-}ms}(d, q))(a(u), a(w)) \leq (b \cdot \operatorname{dist}(\operatorname{form\text{-}ms}(d, q))(u, w)) \wedge \\ \quad c \in (\operatorname{pts}(\operatorname{form\text{-}ms}(d, q)) \to \operatorname{pts}(\operatorname{form\text{-}ms}(d, q))) \wedge \\ \quad f \in \mathbb{R} \wedge \\ \quad 0 \leq f \wedge \\ \quad \forall u, w \in \operatorname{pts}(\operatorname{form\text{-}ms}(d, q)). \\ \quad \quad \operatorname{dist}(\operatorname{form\text{-}ms}(d, q))(c(u), c(w)) \leq (f \cdot \operatorname{dist}(\operatorname{form\text{-}ms}(d, q))(u, w)) \wedge \\ \quad g \in (\operatorname{pts}(\operatorname{form\text{-}ms}(d, q)) \to \operatorname{pts}(\operatorname{form\text{-}ms}(d, q))) \wedge \\ \quad h \in \mathbb{R} \wedge \\ \quad 0 \leq h \wedge \\ \quad \forall u, w \in \operatorname{pts}(\operatorname{form\text{-}ms}(d, q)). \\ \quad \quad \operatorname{dist}(\operatorname{form\text{-}ms}(d, q))(g(u), g(w)) \leq (h \cdot \operatorname{dist}(\operatorname{form\text{-}ms}(d, q))(u, w)) \wedge \\ \quad \forall v \in \operatorname{pts}(\operatorname{form\text{-}ms}(d, q)).\; (g(v) = a(c(v))) \Rightarrow \\ \quad \quad \operatorname{completion\text{-}lift}(\operatorname{form\text{-}ms}(d, q), g)(x) \\ \quad \quad =\; \operatorname{completion\text{-}lift}(\operatorname{form\text{-}ms}(d, q), a)(\operatorname{completion\text{-}lift}(\operatorname{form\text{-}ms}(d, q), c)(x)) \end{array}}}
\vnbprop{completion-lift-id}{\vnbstmt{\begin{array}{@{}l@{}} \forall d, q, a, b.; \forall x \in \operatorname{vec}(\operatorname{hilbert\text{-}completion}(d, q)). \\ \quad \operatorname{is\text{-}module}(d) \wedge \\ \quad \operatorname{scal}(d) = \operatorname{normed\text{-}field\text{-}as\text{-}commutative\text{-}ring}(\operatorname{cc\text{-}normed\text{-}field}) \wedge \\ \quad \operatorname{is\text{-}inner\text{-}product\text{-}form}(d, q) \wedge \\ \quad a \in (\operatorname{pts}(\operatorname{form\text{-}ms}(d, q)) \to \operatorname{pts}(\operatorname{form\text{-}ms}(d, q))) \wedge \\ \quad b \in \mathbb{R} \wedge \\ \quad 0 \leq b \wedge \\ \quad \forall u, w \in \operatorname{pts}(\operatorname{form\text{-}ms}(d, q)). \\ \quad \quad \operatorname{dist}(\operatorname{form\text{-}ms}(d, q))(a(u), a(w)) \leq (b \cdot \operatorname{dist}(\operatorname{form\text{-}ms}(d, q))(u, w)) \wedge \\ \quad \forall v \in \operatorname{pts}(\operatorname{form\text{-}ms}(d, q)).\; (a(v) = v) \Rightarrow \\ \quad \quad \operatorname{completion\text{-}lift}(\operatorname{form\text{-}ms}(d, q), a)(x) = x \end{array}}}
\vnbprop{completion-lift-sum}{\vnbstmt{\begin{array}{@{}l@{}} \forall d, q, a, b, c, f, g, h.; \forall x \in \operatorname{vec}(\operatorname{hilbert\text{-}completion}(d, q)). \\ \quad \operatorname{is\text{-}module}(d) \wedge \\ \quad \operatorname{scal}(d) = \operatorname{normed\text{-}field\text{-}as\text{-}commutative\text{-}ring}(\operatorname{cc\text{-}normed\text{-}field}) \wedge \\ \quad \operatorname{is\text{-}inner\text{-}product\text{-}form}(d, q) \wedge \\ \quad a \in (\operatorname{pts}(\operatorname{form\text{-}ms}(d, q)) \to \operatorname{pts}(\operatorname{form\text{-}ms}(d, q))) \wedge \\ \quad b \in \mathbb{R} \wedge \\ \quad 0 \leq b \wedge \\ \quad \forall u, w \in \operatorname{pts}(\operatorname{form\text{-}ms}(d, q)). \\ \quad \quad \operatorname{dist}(\operatorname{form\text{-}ms}(d, q))(a(u), a(w)) \leq (b \cdot \operatorname{dist}(\operatorname{form\text{-}ms}(d, q))(u, w)) \wedge \\ \quad c \in (\operatorname{pts}(\operatorname{form\text{-}ms}(d, q)) \to \operatorname{pts}(\operatorname{form\text{-}ms}(d, q))) \wedge \\ \quad f \in \mathbb{R} \wedge \\ \quad 0 \leq f \wedge \\ \quad \forall u, w \in \operatorname{pts}(\operatorname{form\text{-}ms}(d, q)). \\ \quad \quad \operatorname{dist}(\operatorname{form\text{-}ms}(d, q))(c(u), c(w)) \leq (f \cdot \operatorname{dist}(\operatorname{form\text{-}ms}(d, q))(u, w)) \wedge \\ \quad g \in (\operatorname{pts}(\operatorname{form\text{-}ms}(d, q)) \to \operatorname{pts}(\operatorname{form\text{-}ms}(d, q))) \wedge \\ \quad h \in \mathbb{R} \wedge \\ \quad 0 \leq h \wedge \\ \quad \forall u, w \in \operatorname{pts}(\operatorname{form\text{-}ms}(d, q)). \\ \quad \quad \operatorname{dist}(\operatorname{form\text{-}ms}(d, q))(g(u), g(w)) \leq (h \cdot \operatorname{dist}(\operatorname{form\text{-}ms}(d, q))(u, w)) \wedge \\ \quad \forall v \in \operatorname{pts}(\operatorname{form\text{-}ms}(d, q)).\; (g(v) = \operatorname{vadd}(d)(a(v), c(v))) \Rightarrow \\ \quad \quad \operatorname{completion\text{-}lift}(\operatorname{form\text{-}ms}(d, q), g)(x) \\ \quad \quad =\; \operatorname{vadd}(\operatorname{hilbert\text{-}completion}(d, q))(\operatorname{completion\text{-}lift}(\operatorname{form\text{-}ms}(d, q), a)(x), \\ \quad \quad \quad \operatorname{completion\text{-}lift}(\operatorname{form\text{-}ms}(d, q), c)(x)) \end{array}}}
\vnbprop{completion-lift-scale}{\vnbstmt{\begin{array}{@{}l@{}} \forall d, q, a, b, f, g.; \forall c \in \mathbb{C}, x \in \operatorname{vec}(\operatorname{hilbert\text{-}completion}(d, q)). \\ \quad \operatorname{is\text{-}module}(d) \wedge \\ \quad \operatorname{scal}(d) = \operatorname{normed\text{-}field\text{-}as\text{-}commutative\text{-}ring}(\operatorname{cc\text{-}normed\text{-}field}) \wedge \\ \quad \operatorname{is\text{-}inner\text{-}product\text{-}form}(d, q) \wedge \\ \quad a \in (\operatorname{pts}(\operatorname{form\text{-}ms}(d, q)) \to \operatorname{pts}(\operatorname{form\text{-}ms}(d, q))) \wedge \\ \quad b \in \mathbb{R} \wedge \\ \quad 0 \leq b \wedge \\ \quad \forall u, w \in \operatorname{pts}(\operatorname{form\text{-}ms}(d, q)). \\ \quad \quad \operatorname{dist}(\operatorname{form\text{-}ms}(d, q))(a(u), a(w)) \leq (b \cdot \operatorname{dist}(\operatorname{form\text{-}ms}(d, q))(u, w)) \wedge \\ \quad f \in (\operatorname{pts}(\operatorname{form\text{-}ms}(d, q)) \to \operatorname{pts}(\operatorname{form\text{-}ms}(d, q))) \wedge \\ \quad g \in \mathbb{R} \wedge \\ \quad 0 \leq g \wedge \\ \quad \forall u, w \in \operatorname{pts}(\operatorname{form\text{-}ms}(d, q)). \\ \quad \quad \operatorname{dist}(\operatorname{form\text{-}ms}(d, q))(f(u), f(w)) \leq (g \cdot \operatorname{dist}(\operatorname{form\text{-}ms}(d, q))(u, w)) \wedge \\ \quad \forall v \in \operatorname{pts}(\operatorname{form\text{-}ms}(d, q)).\; (f(v) = \operatorname{act}(d)(c, a(v))) \Rightarrow \\ \quad \quad \operatorname{completion\text{-}lift}(\operatorname{form\text{-}ms}(d, q), f)(x) \\ \quad \quad =\; \operatorname{act}(\operatorname{hilbert\text{-}completion}(d, q))(c, \operatorname{completion\text{-}lift}(\operatorname{form\text{-}ms}(d, q), a)(x)) \end{array}}}
\vnbprop{completion-lift-additive}{\vnbstmt{\begin{array}{@{}l@{}} \forall d, q, a, b.; \forall x, y \in \operatorname{vec}(\operatorname{hilbert\text{-}completion}(d, q)). \\ \quad \operatorname{is\text{-}module}(d) \wedge \\ \quad \operatorname{scal}(d) = \operatorname{normed\text{-}field\text{-}as\text{-}commutative\text{-}ring}(\operatorname{cc\text{-}normed\text{-}field}) \wedge \\ \quad \operatorname{is\text{-}inner\text{-}product\text{-}form}(d, q) \wedge \\ \quad a \in (\operatorname{pts}(\operatorname{form\text{-}ms}(d, q)) \to \operatorname{pts}(\operatorname{form\text{-}ms}(d, q))) \wedge \\ \quad b \in \mathbb{R} \wedge \\ \quad 0 \leq b \wedge \\ \quad \forall u, w \in \operatorname{pts}(\operatorname{form\text{-}ms}(d, q)). \\ \quad \quad \operatorname{dist}(\operatorname{form\text{-}ms}(d, q))(a(u), a(w)) \leq (b \cdot \operatorname{dist}(\operatorname{form\text{-}ms}(d, q))(u, w)) \wedge \\ \quad \forall v, z \in \operatorname{pts}(\operatorname{form\text{-}ms}(d, q)). \\ \quad \quad a(\operatorname{vadd}(d)(v, z)) = \operatorname{vadd}(d)(a(v), a(z)) \Rightarrow \\ \quad \quad \operatorname{completion\text{-}lift}(\operatorname{form\text{-}ms}(d, q), a)(\operatorname{vadd}(\operatorname{hilbert\text{-}completion}(d, q))(x, y)) \\ \quad \quad =\; \operatorname{vadd}(\operatorname{hilbert\text{-}completion}(d, q))(\operatorname{completion\text{-}lift}(\operatorname{form\text{-}ms}(d, q), a)(x), \\ \quad \quad \quad \operatorname{completion\text{-}lift}(\operatorname{form\text{-}ms}(d, q), a)(y)) \end{array}}}
\vnbprop{completion-lift-homogeneous}{\vnbstmt{\begin{array}{@{}l@{}} \forall d, q, a, b.; \forall c \in \mathbb{C}, x \in \operatorname{vec}(\operatorname{hilbert\text{-}completion}(d, q)). \\ \quad \operatorname{is\text{-}module}(d) \wedge \\ \quad \operatorname{scal}(d) = \operatorname{normed\text{-}field\text{-}as\text{-}commutative\text{-}ring}(\operatorname{cc\text{-}normed\text{-}field}) \wedge \\ \quad \operatorname{is\text{-}inner\text{-}product\text{-}form}(d, q) \wedge \\ \quad a \in (\operatorname{pts}(\operatorname{form\text{-}ms}(d, q)) \to \operatorname{pts}(\operatorname{form\text{-}ms}(d, q))) \wedge \\ \quad b \in \mathbb{R} \wedge \\ \quad 0 \leq b \wedge \\ \quad \forall u, w \in \operatorname{pts}(\operatorname{form\text{-}ms}(d, q)). \\ \quad \quad \operatorname{dist}(\operatorname{form\text{-}ms}(d, q))(a(u), a(w)) \leq (b \cdot \operatorname{dist}(\operatorname{form\text{-}ms}(d, q))(u, w)) \wedge \\ \quad \forall \mathit{g4c} \in \mathbb{C}, v \in \operatorname{pts}(\operatorname{form\text{-}ms}(d, q)). \\ \quad \quad a(\operatorname{act}(d)(\mathit{g4c}, v)) = \operatorname{act}(d)(\mathit{g4c}, a(v)) \Rightarrow \\ \quad \quad \operatorname{completion\text{-}lift}(\operatorname{form\text{-}ms}(d, q), a)(\operatorname{act}(\operatorname{hilbert\text{-}completion}(d, q))(c, x)) \\ \quad \quad =\; \operatorname{act}(\operatorname{hilbert\text{-}completion}(d, q))(c, \operatorname{completion\text{-}lift}(\operatorname{form\text{-}ms}(d, q), a)(x)) \end{array}}}
\vnbprop{completion-lift-adjoint}{\vnbstmt{\begin{array}{@{}l@{}} \forall d, q, a, b, c, f.; \forall x, y \in \operatorname{vec}(\operatorname{hilbert\text{-}completion}(d, q)). \\ \quad \operatorname{is\text{-}module}(d) \wedge \\ \quad \operatorname{scal}(d) = \operatorname{normed\text{-}field\text{-}as\text{-}commutative\text{-}ring}(\operatorname{cc\text{-}normed\text{-}field}) \wedge \\ \quad \operatorname{is\text{-}inner\text{-}product\text{-}form}(d, q) \wedge \\ \quad a \in (\operatorname{pts}(\operatorname{form\text{-}ms}(d, q)) \to \operatorname{pts}(\operatorname{form\text{-}ms}(d, q))) \wedge \\ \quad b \in \mathbb{R} \wedge \\ \quad 0 \leq b \wedge \\ \quad \forall u, w \in \operatorname{pts}(\operatorname{form\text{-}ms}(d, q)). \\ \quad \quad \operatorname{dist}(\operatorname{form\text{-}ms}(d, q))(a(u), a(w)) \leq (b \cdot \operatorname{dist}(\operatorname{form\text{-}ms}(d, q))(u, w)) \wedge \\ \quad c \in (\operatorname{pts}(\operatorname{form\text{-}ms}(d, q)) \to \operatorname{pts}(\operatorname{form\text{-}ms}(d, q))) \wedge \\ \quad f \in \mathbb{R} \wedge \\ \quad 0 \leq f \wedge \\ \quad \forall u, w \in \operatorname{pts}(\operatorname{form\text{-}ms}(d, q)). \\ \quad \quad \operatorname{dist}(\operatorname{form\text{-}ms}(d, q))(c(u), c(w)) \leq (f \cdot \operatorname{dist}(\operatorname{form\text{-}ms}(d, q))(u, w)) \wedge \\ \quad \forall v, z \in \operatorname{pts}(\operatorname{form\text{-}ms}(d, q)). \\ \quad \quad q(a(v), z) = q(v, c(z)) \Rightarrow \\ \quad \quad \langle\cdot,\cdot\rangle_{\operatorname{hilbert\text{-}completion}(d, q)}(\operatorname{completion\text{-}lift}(\operatorname{form\text{-}ms}(d, q), a)(x), y) \\ \quad \quad =\; \langle\cdot,\cdot\rangle_{\operatorname{hilbert\text{-}completion}(d, q)}(x, \operatorname{completion\text{-}lift}(\operatorname{form\text{-}ms}(d, q), c)(y)) \end{array}}}
\vnbprop{completion-lift-bound}{\vnbstmt{\begin{array}{@{}l@{}} \forall d, q, a, b.; \forall c \in \mathbb{R}, x \in \operatorname{vec}(\operatorname{hilbert\text{-}completion}(d, q)). \\ \quad \operatorname{is\text{-}module}(d) \wedge \\ \quad \operatorname{scal}(d) = \operatorname{normed\text{-}field\text{-}as\text{-}commutative\text{-}ring}(\operatorname{cc\text{-}normed\text{-}field}) \wedge \\ \quad \operatorname{is\text{-}inner\text{-}product\text{-}form}(d, q) \wedge \\ \quad a \in (\operatorname{pts}(\operatorname{form\text{-}ms}(d, q)) \to \operatorname{pts}(\operatorname{form\text{-}ms}(d, q))) \wedge \\ \quad b \in \mathbb{R} \wedge \\ \quad 0 \leq b \wedge \\ \quad \forall u, w \in \operatorname{pts}(\operatorname{form\text{-}ms}(d, q)). \\ \quad \quad \operatorname{dist}(\operatorname{form\text{-}ms}(d, q))(a(u), a(w)) \leq (b \cdot \operatorname{dist}(\operatorname{form\text{-}ms}(d, q))(u, w)) \wedge \\ \quad \forall v \in \operatorname{pts}(\operatorname{form\text{-}ms}(d, q)).\; (q(a(v), a(v)) \leq (c \cdot q(v, v))) \Rightarrow \\ \quad \quad \langle\cdot,\cdot\rangle_{\operatorname{hilbert\text{-}completion}(d, q)}(\operatorname{completion\text{-}lift}(\operatorname{form\text{-}ms}(d, q), a)(x), \\ \quad \quad \quad \operatorname{completion\text{-}lift}(\operatorname{form\text{-}ms}(d, q), a)(x)) \\ \quad \quad \leq\; c \cdot \langle\cdot,\cdot\rangle_{\operatorname{hilbert\text{-}completion}(d, q)}(x, x) \end{array}}}
\vnbprop{completion-lift-bounded}{\vnbstmt{\begin{array}{@{}l@{}} \forall d, q, a, b. \\ \quad \operatorname{is\text{-}module}(d) \wedge \\ \quad \operatorname{scal}(d) = \operatorname{normed\text{-}field\text{-}as\text{-}commutative\text{-}ring}(\operatorname{cc\text{-}normed\text{-}field}) \wedge \\ \quad \operatorname{is\text{-}inner\text{-}product\text{-}form}(d, q) \wedge \\ \quad a \in (\operatorname{pts}(\operatorname{form\text{-}ms}(d, q)) \to \operatorname{pts}(\operatorname{form\text{-}ms}(d, q))) \wedge \\ \quad b \in \mathbb{R} \wedge \\ \quad 0 \leq b \wedge \\ \quad \forall u, w \in \operatorname{pts}(\operatorname{form\text{-}ms}(d, q)). \\ \quad \quad \operatorname{dist}(\operatorname{form\text{-}ms}(d, q))(a(u), a(w)) \leq (b \cdot \operatorname{dist}(\operatorname{form\text{-}ms}(d, q))(u, w)) \wedge \\ \quad \forall v, z \in \operatorname{pts}(\operatorname{form\text{-}ms}(d, q)). \\ \quad \quad a(\operatorname{vadd}(d)(v, z)) = \operatorname{vadd}(d)(a(v), a(z)) \wedge \\ \quad \forall c \in \mathbb{C}, v \in \operatorname{pts}(\operatorname{form\text{-}ms}(d, q)). \\ \quad \quad a(\operatorname{act}(d)(c, v)) = \operatorname{act}(d)(c, a(v)) \wedge \\ \quad \forall v \in \operatorname{pts}(\operatorname{form\text{-}ms}(d, q)).\; (q(a(v), a(v)) \leq ((b \cdot b) \cdot q(v, v))) \Rightarrow \\ \quad \quad \operatorname{completion\text{-}lift}(\operatorname{form\text{-}ms}(d, q), a) \in L({\operatorname{hilbert\text{-}completion}(d, q)}, {\operatorname{hilbert\text{-}completion}(d, q)}) \end{array}}}

\vnbbookfile{theorem-library/ba4-sqrt}{theorem-library/ba4-sqrt}
\noindent{\small ba4-sqrt.scm -- 1 - y IS A SQUARE for ||y|| $<$ 1, by iteration (the square root behind Prop 4.7). Agent BA-1111-B, 2026-11-11.  Destination: theorem-library/ba4-sqrt.scm.  Library: operator-algebras. Load window: after structure-library/ba-rep-vocab (HALF-SQ-ITER), theorem-library/kal-laws, ba-laws, ba-elem ...}\par
\vnbprop{g4s-iter-type}{\vnbstmt{\begin{array}{@{}l@{}} \forall n \in \mathbb{N}, a.; \forall y \in \operatorname{vec}(a). \\ \quad \operatorname{is\text{-}unital\text{-}banach\text{-}algebra}(a) \Rightarrow (\operatorname{half\text{-}sq\text{-}iter}(a, y, n) \in \operatorname{vec}(a)) \end{array}}}
\vnbprop{g4s-contract}{\vnbstmt{\begin{array}{@{}l@{}} \forall a.; \forall y, u, w \in \operatorname{vec}(a). \\ \quad \operatorname{is\text{-}unital\text{-}banach\text{-}algebra}(a) \Rightarrow  \\ \quad \quad \operatorname{vnrm}(a)(\operatorname{vadd}(a)(\operatorname{act}(a)(1/2, \operatorname{vadd}(a)(y, \cdot_{a}(u, u))), \\ \quad \quad \quad \quad \operatorname{vneg}(a)(\operatorname{act}(a)(1/2, \operatorname{vadd}(a)(y, \cdot_{a}(w, w)))))) \\ \quad \quad \leq\; (1/2 \cdot (\operatorname{vnrm}(a)(u) + \operatorname{vnrm}(a)(w))) \cdot \operatorname{vnrm}(a)(\operatorname{vadd}(a)(u, \operatorname{vneg}(a)(w))) \end{array}}}
\vnbprop{g4s-iter-bound}{\vnbstmt{\begin{array}{@{}l@{}} \forall n \in \mathbb{N}, a.; \forall y \in \operatorname{vec}(a), s \in \mathbb{R}. \\ \quad (\operatorname{is\text{-}unital\text{-}banach\text{-}algebra}(a) \wedge ((\operatorname{vnrm}(a)(y) < 1) \wedge (s = (1/2 \cdot (1 + \operatorname{vnrm}(a)(y)))))) \Rightarrow  \\ \quad \quad \operatorname{vnrm}(a)(\operatorname{half\text{-}sq\text{-}iter}(a, y, n)) \leq s \end{array}}}
\vnbprop{g4s-iter-incr}{\vnbstmt{\begin{array}{@{}l@{}} \forall n \in \mathbb{N}, a.; \forall y \in \operatorname{vec}(a), s \in \mathbb{R}. \\ \quad (\operatorname{is\text{-}unital\text{-}banach\text{-}algebra}(a) \wedge ((\operatorname{vnrm}(a)(y) < 1) \wedge (s = (1/2 \cdot (1 + \operatorname{vnrm}(a)(y)))))) \Rightarrow  \\ \quad \quad \operatorname{vnrm}(a)(\operatorname{vadd}(a)(\operatorname{half\text{-}sq\text{-}iter}(a, y, n), \operatorname{vneg}(a)(\operatorname{half\text{-}sq\text{-}iter}(a, y, (n + 1))))) \leq \operatorname{power}(s, n) \end{array}}}
\vnbprop{half-sq-fixed-point}{\vnbstmt{\begin{array}{@{}l@{}} \forall a.; \forall y \in \operatorname{vec}(a). \\ \quad (\operatorname{is\text{-}unital\text{-}banach\text{-}algebra}(a) \wedge (\operatorname{vnrm}(a)(y) < 1)) \Rightarrow  \\ \quad \quad \exists z \in \operatorname{vec}(a). \\ \quad \quad \quad (\operatorname{vnrm}(a)(z) < 1) \wedge (z = \operatorname{act}(a)(1/2, \operatorname{vadd}(a)(y, \cdot_{a}(z, z)))) \end{array}}}
\vnbprop{half-sq-fixed-unique}{\vnbstmt{\begin{array}{@{}l@{}} \forall a.; \forall y, u, w \in \operatorname{vec}(a). \\ \quad \operatorname{is\text{-}unital\text{-}banach\text{-}algebra}(a) \wedge \\ \quad \operatorname{vnrm}(a)(u) < 1 \wedge \\ \quad u = \operatorname{act}(a)(1/2, \operatorname{vadd}(a)(y, \cdot_{a}(u, u))) \wedge \\ \quad \operatorname{vnrm}(a)(w) < 1 \wedge \\ \quad w = \operatorname{act}(a)(1/2, \operatorname{vadd}(a)(y, \cdot_{a}(w, w))) \Rightarrow \\ \quad \quad u = w \end{array}}}
\vnbprop{one-minus-square}{\vnbstmt{\begin{array}{@{}l@{}} \forall a.; \forall y \in \operatorname{vec}(a). \\ \quad (\operatorname{is\text{-}unital\text{-}banach\text{-}algebra}(a) \wedge (\operatorname{vnrm}(a)(y) < 1)) \Rightarrow  \\ \quad \quad \exists b \in \operatorname{vec}(a).\; (\cdot_{a}(b, b) = \operatorname{vadd}(a)(1_{a}, \operatorname{vneg}(a)(y))) \end{array}}}
\vnbprop{uiba-one-minus-sa-square}{\vnbstmt{\begin{array}{@{}l@{}} \forall a.; \forall y \in \operatorname{vec}(a). \\ \quad (\operatorname{is\text{-}unital\text{-}involutive\text{-}banach\text{-}algebra}(a) \wedge (({}^{*_{a}}(y) = y) \wedge (\operatorname{vnrm}(a)(y) < 1))) \Rightarrow  \\ \quad \quad \exists b \in \operatorname{vec}(a). \\ \quad \quad \quad ({}^{*_{a}}(b) = b) \wedge (\cdot_{a}(b, b) = \operatorname{vadd}(a)(1_{a}, \operatorname{vneg}(a)(y))) \end{array}}}

\vnbbookfile{structure-library/gn-vocab}{structure-library/gn-vocab}
\noindent{\small gn-vocab.scm -- BANACH ALGEBRAS: the natural order of an involutive algebra (Def. 1.34, (15)) and the continuous functional calculus in a commutative C*-algebra (Prop. 1.31). F. J. Thayer, "Banach Algebras and Applications", ch. 1 ss. 3-5. DEFINITIONS ONLY: no proof, no support, no axiom.  Agent BA-1117-GN, wave ...}\par
\par\medskip\noindent\textbf{Definition} (\texttt{ale}).\;For all $\operatorname{gnva\_}$, $\operatorname{gnvx\_}$ and $\operatorname{gnvy\_}$, 
\begin{equation*}
{\operatorname{gnvx\_}} \le_{\operatorname{gnva\_}} {\operatorname{gnvy\_}} \Leftrightarrow (\operatorname{vadd}(\operatorname{gnva\_})(\operatorname{gnvy\_}, \operatorname{vneg}(\operatorname{gnva\_})(\operatorname{gnvx\_})) \in {\operatorname{gnva\_}}_{+})
\end{equation*}

\par\medskip\noindent\textbf{Definition} (\texttt{cfc}).\;${\operatorname{gnvf\_}}({\operatorname{gnvx\_}}) :=$
\begin{equation*}
\varepsilon(\{\operatorname{gnvy\_} \in \operatorname{vec}(\operatorname{gnva\_}) : (\operatorname{gelfand\text{-}transform}(\operatorname{gnva\_})(\operatorname{gnvy\_}) = \operatorname{compose}(\operatorname{gnvf\_}, \operatorname{gelfand\text{-}transform}(\operatorname{gnva\_})(\operatorname{gnvx\_})))\})
\end{equation*}

\par\medskip\noindent\textbf{Definition} (\texttt{spec-top}).\;$\operatorname{spec\text{-}top}(\operatorname{gnva\_}, \operatorname{gnvx\_}) :=$ $\operatorname{subspace\text{-}top}(\operatorname{metric\text{-}top}(\operatorname{cc\text{-}ms}), \sigma({\operatorname{gnvx\_}}))$
\par\medskip\noindent\textbf{Definition} (\texttt{dual-map}).\;$\widehat{{\operatorname{gnvf\_}}} :=$ $(\lambda \operatorname{gnvw\_} \in \operatorname{characters}(\operatorname{gnvb\_}).\; \operatorname{compose}(\operatorname{gnvw\_}, \operatorname{gnvf\_}))$
\par\medskip\noindent\textbf{Definition} (\texttt{dual-pull}).\;$\operatorname{dual\text{-}pull}(\operatorname{gnva\_}, \operatorname{gnvb\_}, \operatorname{gnvg\_}) :=$
\begin{equation*}
\begin{array}{@{}l@{}} (\lambda \operatorname{gnvx\_} \in \operatorname{vec}(\operatorname{gnva\_}).\; \\ \quad \varepsilon(\{\operatorname{gnvy\_} \in \operatorname{vec}(\operatorname{gnvb\_}) : (\operatorname{gelfand\text{-}transform}(\operatorname{gnvb\_})(\operatorname{gnvy\_}) = \operatorname{compose}(\operatorname{gelfand\text{-}transform}(\operatorname{gnva\_})(\operatorname{gnvx\_}), \operatorname{gnvg\_}))\})) \end{array}
\end{equation*}

\par\medskip\noindent\textbf{Definition} (\texttt{sub-ucsa}).\;$\operatorname{sub\text{-}ucsa}(\operatorname{gnvb\_}, \operatorname{gnvs\_}) :=$
\begin{equation*}
{[}\operatorname{normed\text{-}field\text{-}as\text{-}commutative\text{-}ring}(\operatorname{cc\text{-}normed\text{-}field}), \operatorname{gnvs\_}, (\lambda {[}\operatorname{gnvu\_}, \operatorname{gnvv\_}] \in (\operatorname{gnvs\_} \times \operatorname{gnvs\_}).\; \operatorname{vadd}(\operatorname{gnvb\_})(\operatorname{gnvu\_}, \operatorname{gnvv\_})), \operatorname{vzero}(\operatorname{gnvb\_}), (\lambda \operatorname{gnvu\_} \in \operatorname{gnvs\_}.\; \operatorname{vneg}(\operatorname{gnvb\_})(\operatorname{gnvu\_})), (\lambda {[}\operatorname{gnvc\_}, \operatorname{gnvu\_}] \in (\mathbb{C} \times \operatorname{gnvs\_}).\; \operatorname{act}(\operatorname{gnvb\_})(\operatorname{gnvc\_}, \operatorname{gnvu\_})), (\lambda \operatorname{gnvu\_} \in \operatorname{gnvs\_}.\; \operatorname{vnrm}(\operatorname{gnvb\_})(\operatorname{gnvu\_})), (\lambda {[}\operatorname{gnvu\_}, \operatorname{gnvv\_}] \in (\operatorname{gnvs\_} \times \operatorname{gnvs\_}).\; \cdot_{\operatorname{gnvb\_}}(\operatorname{gnvu\_}, \operatorname{gnvv\_})), 1_{\operatorname{gnvb\_}}, (\lambda \operatorname{gnvu\_} \in \operatorname{gnvs\_}.\; {}^{*_{\operatorname{gnvb\_}}}(\operatorname{gnvu\_}))]
\end{equation*}

\par\medskip\noindent\textbf{Definition} (\texttt{commutant-el}).\;$\operatorname{commutant\text{-}el}(\operatorname{gnvb\_}, \operatorname{gnvt\_}) :=$
\begin{equation*}
\{\operatorname{gnvy\_} \in \operatorname{vec}(\operatorname{gnvb\_}) : \forall \operatorname{gnvz\_} \in \operatorname{gnvt\_}.\; (\cdot_{\operatorname{gnvb\_}}(\operatorname{gnvy\_}, \operatorname{gnvz\_}) = \cdot_{\operatorname{gnvb\_}}(\operatorname{gnvz\_}, \operatorname{gnvy\_}))\}
\end{equation*}

\par\medskip\noindent\textbf{Definition} (\texttt{cfc-n}).\;$\operatorname{cfc\text{-}n}(\operatorname{gnvb\_}, \operatorname{gnvx\_}, \operatorname{gnvf\_}) :=$ ${\operatorname{gnvf\_}}({\operatorname{gnvx\_}})$

\vnbbookfile{theorem-library/gn-order}{theorem-library/gn-order}
\noindent{\small gn-order.scm -- the positive cone and the order unit (Def. 1.34, Prop. 1.36) of F. J. Thayer, "Banach Algebras and Applications", ch. 1 s. 4.  Agent BA-1117-GN, wave 2026-11-17.  All proofs modulo 0. Destination: theorem-library/gn-order.scm, after structure-library/gn-vocab.  Library: operator-algebras. Helper prefix ...}\par
\vnbprop{pos-cone-in-vec}{\vnbstmt{\forall a.; \forall x \in {a}_{+}.\; (x \in \operatorname{vec}(a))}}
\vnbprop{pos-cone-star-square}{\vnbstmt{\begin{array}{@{}l@{}} \forall a.; \forall y \in \operatorname{vec}(a). \\ \quad \operatorname{is\text{-}unital\text{-}involutive\text{-}banach\text{-}algebra}(a) \Rightarrow  \\ \quad \quad \cdot_{a}({}^{*_{a}}(y), y) \in {a}_{+} \end{array}}}
\vnbprop{pos-cone-add-closed}{\vnbstmt{\begin{array}{@{}l@{}} \forall a.; \forall x, y \in {a}_{+}. \\ \quad \operatorname{is\text{-}unital\text{-}involutive\text{-}banach\text{-}algebra}(a) \Rightarrow (\operatorname{vadd}(a)(x, y) \in {a}_{+}) \end{array}}}
\vnbprop{pos-cone-induct}{\vnbstmt{\begin{array}{@{}l@{}} \forall a, s.; \forall x \in {a}_{+}. \\ \quad \operatorname{is\text{-}unital\text{-}involutive\text{-}banach\text{-}algebra}(a) \wedge \\ \quad s \subseteq \operatorname{vec}(a) \wedge \\ \quad \forall z \in \operatorname{vec}(a).\; (\cdot_{a}({}^{*_{a}}(z), z) \in s) \wedge \\ \quad \forall u, v \in s.\; (\operatorname{vadd}(a)(u, v) \in s) \Rightarrow \\ \quad \quad x \in s \end{array}}}
\vnbprop{sa-part-unfold}{\vnbstmt{\forall a.\; (\operatorname{sa\text{-}part}(a) \simeq \{v \in \operatorname{vec}(a) : ({}^{*_{a}}(v) = v)\})}}
\vnbprop{sa-squares-unfold}{\vnbstmt{\begin{array}{@{}l@{}} \forall a. \\ \quad (\operatorname{sa\text{-}squares}(a) \simeq \{v \in \operatorname{vec}(a) : \exists w \in \operatorname{sa\text{-}part}(a).\; (v = \cdot_{a}(w, w))\}) \end{array}}}
\vnbprop{pos-cone-self-adjoint}{\vnbstmt{\begin{array}{@{}l@{}} \forall a.; \forall x \in {a}_{+}. \\ \quad \operatorname{is\text{-}unital\text{-}involutive\text{-}banach\text{-}algebra}(a) \Rightarrow (x \in \operatorname{sa\text{-}part}(a)) \end{array}}}
\vnbprop{pos-cone-one}{\vnbstmt{\begin{array}{@{}l@{}} \forall a. \\ \quad \operatorname{is\text{-}unital\text{-}involutive\text{-}banach\text{-}algebra}(a) \Rightarrow (1_{a} \in {a}_{+}) \end{array}}}
\vnbprop{order-unit}{\vnbstmt{\begin{array}{@{}l@{}} \forall a.; \forall x \in \operatorname{sa\text{-}part}(a). \\ \quad \operatorname{is\text{-}unital\text{-}involutive\text{-}banach\text{-}algebra}(a) \Rightarrow  \\ \quad \quad {x} \le_{a} {\operatorname{act}(a)(((\operatorname{vnrm}(a)(x) + 1) \cdot (\operatorname{vnrm}(a)(x) + 1)), 1_{a})} \end{array}}}
\vnbprop{order-unit-exists}{\vnbstmt{\begin{array}{@{}l@{}} \forall a.; \forall x \in \operatorname{sa\text{-}part}(a). \\ \quad \operatorname{is\text{-}unital\text{-}involutive\text{-}banach\text{-}algebra}(a) \Rightarrow  \\ \quad \quad \exists l \in \mathbb{R}.\; {x} \le_{a} {\operatorname{act}(a)(l, 1_{a})} \end{array}}}
\vnbprop{order-unit-strict}{\vnbstmt{\begin{array}{@{}l@{}} \forall a.; \forall x \in \operatorname{sa\text{-}part}(a), l \in \mathbb{R}. \\ \quad (\operatorname{is\text{-}unital\text{-}involutive\text{-}banach\text{-}algebra}(a) \wedge (\operatorname{vnrm}(a)(x) < l)) \Rightarrow  \\ \quad \quad {x} \le_{a} {\operatorname{act}(a)(l, 1_{a})} \end{array}}}
\vnbprop{pos-cone-conjugate}{\vnbstmt{\begin{array}{@{}l@{}} \forall a.; \forall z \in \operatorname{vec}(a), x \in {a}_{+}. \\ \quad \operatorname{is\text{-}unital\text{-}involutive\text{-}banach\text{-}algebra}(a) \Rightarrow  \\ \quad \quad \cdot_{a}(\cdot_{a}({}^{*_{a}}(z), x), z) \in {a}_{+} \end{array}}}

\vnbbookfile{theorem-library/gn-onto}{theorem-library/gn-onto}
\noindent{\small gn-onto.scm -- Theorem 1.23 of F. J. Thayer, "Banach Algebras and Applications", ch. 1: the ONTO half of commutative Gelfand-Naimark.  Agent BA-1117-GN, wave 2026-11-17.  All proofs modulo 0. Destination: theorem-library/gn-onto.scm.  Library: operator-algebras.  Helper prefix gno-. Binders gn*\_ (gnc\_ gnk\_ gnu\_ gny\_ ...}\par
\vnbprop{gno-char-linear}{\vnbstmt{\begin{array}{@{}l@{}} \forall a.; \forall c \in \operatorname{characters}(a). \\ \quad \operatorname{is\text{-}unital\text{-}banach\text{-}algebra}(a) \Rightarrow \operatorname{is\text{-}cc\text{-}linear\text{-}on}(a, c) \end{array}}}
\vnbprop{gno-char-sub}{\vnbstmt{\begin{array}{@{}l@{}} \forall a.; \forall c \in \operatorname{characters}(a), x, y \in \operatorname{vec}(a). \\ \quad \operatorname{is\text{-}unital\text{-}banach\text{-}algebra}(a) \Rightarrow (c(\operatorname{vadd}(a)(x, \operatorname{vneg}(a)(y))) = (c(x) - c(y))) \end{array}}}
\vnbprop{gno-closed-range}{\vnbstmt{\begin{array}{@{}l@{}} \forall a.; \forall f \in C({\widehat{a}}). \\ \quad \operatorname{is\text{-}unital\text{-}banach\text{-}algebra}(a) \wedge \\ \quad \forall x \in \operatorname{vec}(a). \\ \quad \quad \lVert \operatorname{gelfand\text{-}transform}(a)(x) \rVert_{\infty} = \operatorname{vnrm}(a)(x) \wedge \\ \quad \forall e. \\ \quad \quad \operatorname{pos\text{-}rr}(e) \Rightarrow  \\ \quad \quad \quad \exists y \in \operatorname{vec}(a).\; \forall c \in \operatorname{characters}(a).\; (\operatorname{magnitude}((c(y) - f(c))) < e) \Rightarrow \\ \quad \quad \exists x \in \operatorname{vec}(a).\; (\operatorname{gelfand\text{-}transform}(a)(x) = f) \end{array}}}
\vnbprop{gno-u-pts}{\vnbstmt{\begin{array}{@{}l@{}} \forall a. \\ \quad \operatorname{is\text{-}unital\text{-}involutive\text{-}banach\text{-}algebra}(a) \Rightarrow  \\ \quad \quad (\operatorname{pts}(\widehat{a}) \simeq \operatorname{characters}(a)) \end{array}}}
\vnbprop{gno-u-pts-set}{\vnbstmt{\begin{array}{@{}l@{}} \forall a. \\ \quad \operatorname{is\text{-}unital\text{-}involutive\text{-}banach\text{-}algebra}(a) \Rightarrow (\operatorname{pts}(\widehat{a}) \in \mathbf{Set}) \end{array}}}
\vnbprop{gno-u-compact}{\vnbstmt{\begin{array}{@{}l@{}} \forall a. \\ \quad \operatorname{is\text{-}unital\text{-}involutive\text{-}banach\text{-}algebra}(a) \Rightarrow \operatorname{is\text{-}compact\text{-}t}(\widehat{a}) \end{array}}}
\vnbprop{gno-u-in-ccont}{\vnbstmt{\begin{array}{@{}l@{}} \forall a.; \forall x \in \operatorname{vec}(a). \\ \quad \operatorname{is\text{-}unital\text{-}involutive\text{-}banach\text{-}algebra}(a) \Rightarrow  \\ \quad \quad \operatorname{gelfand\text{-}transform}(a)(x) \in C({\widehat{a}}) \end{array}}}
\vnbprop{gno-u-value}{\vnbstmt{\begin{array}{@{}l@{}} \forall a.; \forall x \in \operatorname{vec}(a), c \in \operatorname{characters}(a). \\ \quad \operatorname{is\text{-}unital\text{-}involutive\text{-}banach\text{-}algebra}(a) \Rightarrow  \\ \quad \quad \operatorname{gelfand\text{-}transform}(a)(x)(c) = c(x) \end{array}}}
\vnbprop{gno-range-star-subalgebra}{\vnbstmt{\begin{array}{@{}l@{}} \forall a. \\ \quad \operatorname{is\text{-}unital\text{-}c\text{-}star\text{-}algebra}(a) \Rightarrow  \\ \quad \quad \operatorname{is\text{-}ccont\text{-}star\text{-}subalgebra}(\widehat{a}, \\ \quad \quad \quad \{g \in C({\widehat{a}}) : \exists y \in \operatorname{vec}(a).\; (g = \operatorname{gelfand\text{-}transform}(a)(y))\}) \end{array}}}
\vnbprop{gno-range-separating}{\vnbstmt{\begin{array}{@{}l@{}} \forall a. \\ \quad \operatorname{is\text{-}unital\text{-}involutive\text{-}banach\text{-}algebra}(a) \Rightarrow  \\ \quad \quad \operatorname{is\text{-}point\text{-}separating}(\widehat{a}, \\ \quad \quad \quad \{g \in C({\widehat{a}}) : \exists y \in \operatorname{vec}(a).\; (g = \operatorname{gelfand\text{-}transform}(a)(y))\}) \end{array}}}
\vnbprop{gelfand-dense}{\vnbstmt{\begin{array}{@{}l@{}} \forall a.; \forall f \in C({\widehat{a}}), e. \\ \quad (\operatorname{is\text{-}unital\text{-}c\text{-}star\text{-}algebra}(a) \wedge \operatorname{pos\text{-}rr}(e)) \Rightarrow  \\ \quad \quad \exists y \in \operatorname{vec}(a).\; \forall c \in \operatorname{characters}(a).\; (\operatorname{magnitude}((c(y) - f(c))) < e) \end{array}}}
\vnbprop{gelfand-onto}{\vnbstmt{\begin{array}{@{}l@{}} \forall a.; \forall f \in C({\widehat{a}}). \\ \quad \operatorname{is\text{-}commutative\text{-}c\text{-}star\text{-}algebra}(a) \Rightarrow  \\ \quad \quad \exists x \in \operatorname{vec}(a).\; (\operatorname{gelfand\text{-}transform}(a)(x) = f) \end{array}}}
\vnbprop{gelfand-star-iso}{\vnbstmt{\begin{array}{@{}l@{}} \forall a. \\ \quad \operatorname{is\text{-}commutative\text{-}c\text{-}star\text{-}algebra}(a) \Rightarrow  \\ \quad \quad \operatorname{is\text{-}hom\text{-}commutative\text{-}c\text{-}star\text{-}algebra}(a, \operatorname{ccont\text{-}algebra}(\widehat{a}), \operatorname{gelfand\text{-}transform}(a)) \end{array}}}
\vnbprop{spectrum-unfold}{\vnbstmt{\begin{array}{@{}l@{}} \forall a, x. \\ \quad (\sigma({x}) \simeq \{z \in \mathbb{C} : \neg \operatorname{is\text{-}invertible\text{-}el}(a, \operatorname{shift\text{-}el}(a, z, x))\}) \end{array}}}
\vnbprop{gelfand-value-in-spectrum}{\vnbstmt{\begin{array}{@{}l@{}} \forall a.; \forall x \in \operatorname{vec}(a), c \in \operatorname{characters}(a). \\ \quad \operatorname{is\text{-}unital\text{-}involutive\text{-}banach\text{-}algebra}(a) \Rightarrow (c(x) \in \sigma({x})) \end{array}}}
\vnbprop{spec-top-unfold}{\vnbstmt{\begin{array}{@{}l@{}} \forall a, x. \\ \quad (\operatorname{spec\text{-}top}(a, x) \simeq \operatorname{subspace\text{-}top}(\operatorname{metric\text{-}top}(\operatorname{cc\text{-}ms}), \sigma({x}))) \end{array}}}
\vnbprop{gelfand-into-spectrum-continuous}{\vnbstmt{\begin{array}{@{}l@{}} \forall a.; \forall x \in \operatorname{vec}(a). \\ \quad \operatorname{is\text{-}unital\text{-}involutive\text{-}banach\text{-}algebra}(a) \Rightarrow  \\ \quad \quad \operatorname{is\text{-}hom\text{-}top\text{-}space}(\widehat{a}, \operatorname{spec\text{-}top}(a, x), \operatorname{gelfand\text{-}transform}(a)(x)) \end{array}}}
\vnbprop{compose-gelfand-in-ccont}{\vnbstmt{\begin{array}{@{}l@{}} \forall a.; \forall x \in \operatorname{vec}(a), f \in C({\operatorname{spec\text{-}top}(a, x)}). \\ \quad \operatorname{is\text{-}unital\text{-}involutive\text{-}banach\text{-}algebra}(a) \Rightarrow  \\ \quad \quad \operatorname{compose}(f, \operatorname{gelfand\text{-}transform}(a)(x)) \in C({\widehat{a}}) \end{array}}}
\vnbprop{uiba-sub-zero-eq}{\vnbstmt{\begin{array}{@{}l@{}} \forall a.; \forall y, z \in \operatorname{vec}(a). \\ \quad (\operatorname{is\text{-}unital\text{-}involutive\text{-}banach\text{-}algebra}(a) \wedge (\operatorname{vadd}(a)(y, \operatorname{vneg}(a)(z)) = \operatorname{vzero}(a))) \Rightarrow  \\ \quad \quad y = z \end{array}}}
\vnbprop{gelfand-injective}{\vnbstmt{\begin{array}{@{}l@{}} \forall a.; \forall y, z \in \operatorname{vec}(a). \\ \quad (\operatorname{is\text{-}commutative\text{-}c\text{-}star\text{-}algebra}(a) \wedge (\operatorname{gelfand\text{-}transform}(a)(y) = \operatorname{gelfand\text{-}transform}(a)(z))) \Rightarrow  \\ \quad \quad y = z \end{array}}}
\vnbprop{cfc-unfold}{\vnbstmt{\begin{array}{@{}l@{}} \forall a, x, f. \\ \quad ({f}({x}) \simeq \varepsilon(\{y \in \operatorname{vec}(a) : (\operatorname{gelfand\text{-}transform}(a)(y) = \operatorname{compose}(f, \operatorname{gelfand\text{-}transform}(a)(x)))\})) \end{array}}}
\vnbprop{cfc-spec}{\vnbstmt{\begin{array}{@{}l@{}} \forall a.; \forall x \in \operatorname{vec}(a), f \in C({\operatorname{spec\text{-}top}(a, x)}). \\ \quad \operatorname{is\text{-}commutative\text{-}c\text{-}star\text{-}algebra}(a) \Rightarrow  \\ \quad \quad {f}({x}) \in \operatorname{vec}(a) \wedge \\ \quad \quad \operatorname{gelfand\text{-}transform}(a)({f}({x})) = \operatorname{compose}(f, \operatorname{gelfand\text{-}transform}(a)(x)) \end{array}}}
\vnbprop{cfc-unique}{\vnbstmt{\begin{array}{@{}l@{}} \forall a.; \forall x \in \operatorname{vec}(a), f \in C({\operatorname{spec\text{-}top}(a, x)}), y \in \operatorname{vec}(a). \\ \quad \operatorname{is\text{-}commutative\text{-}c\text{-}star\text{-}algebra}(a) \wedge \\ \quad \operatorname{gelfand\text{-}transform}(a)(y) = \operatorname{compose}(f, \operatorname{gelfand\text{-}transform}(a)(x)) \Rightarrow \\ \quad \quad y = {f}({x}) \end{array}}}
\vnbprop{cfc-value}{\vnbstmt{\begin{array}{@{}l@{}} \forall a.; \forall x \in \operatorname{vec}(a), f \in C({\operatorname{spec\text{-}top}(a, x)}), c \in \operatorname{characters}(a). \\ \quad \operatorname{is\text{-}commutative\text{-}c\text{-}star\text{-}algebra}(a) \Rightarrow (c({f}({x})) = f(c(x))) \end{array}}}
\vnbprop{cfc-add}{\vnbstmt{\begin{array}{@{}l@{}} \forall a.; \forall x \in \operatorname{vec}(a), f, g \in C({\operatorname{spec\text{-}top}(a, x)}). \\ \quad \operatorname{is\text{-}commutative\text{-}c\text{-}star\text{-}algebra}(a) \Rightarrow  \\ \quad \quad {\operatorname{ccont\text{-}vadd}(\operatorname{spec\text{-}top}(a, x))(f, g)}({x}) = \operatorname{vadd}(a)({f}({x}), {g}({x})) \end{array}}}
\vnbprop{cfc-mul}{\vnbstmt{\begin{array}{@{}l@{}} \forall a.; \forall x \in \operatorname{vec}(a), f, g \in C({\operatorname{spec\text{-}top}(a, x)}). \\ \quad \operatorname{is\text{-}commutative\text{-}c\text{-}star\text{-}algebra}(a) \Rightarrow  \\ \quad \quad {\operatorname{ccont\text{-}vmul}(\operatorname{spec\text{-}top}(a, x))(f, g)}({x}) = \cdot_{a}({f}({x}), {g}({x})) \end{array}}}
\vnbprop{cfc-act}{\vnbstmt{\begin{array}{@{}l@{}} \forall a.; \forall x \in \operatorname{vec}(a), f \in C({\operatorname{spec\text{-}top}(a, x)}), z \in \mathbb{C}. \\ \quad \operatorname{is\text{-}commutative\text{-}c\text{-}star\text{-}algebra}(a) \Rightarrow  \\ \quad \quad {\operatorname{ccont\text{-}act}(\operatorname{spec\text{-}top}(a, x))(z, f)}({x}) = \operatorname{act}(a)(z, {f}({x})) \end{array}}}
\vnbprop{cfc-star}{\vnbstmt{\begin{array}{@{}l@{}} \forall a.; \forall x \in \operatorname{vec}(a), f \in C({\operatorname{spec\text{-}top}(a, x)}). \\ \quad \operatorname{is\text{-}commutative\text{-}c\text{-}star\text{-}algebra}(a) \Rightarrow  \\ \quad \quad {\operatorname{ccont\text{-}vstar}(\operatorname{spec\text{-}top}(a, x))(f)}({x}) = {}^{*_{a}}({f}({x})) \end{array}}}
\vnbprop{cfc-one}{\vnbstmt{\begin{array}{@{}l@{}} \forall a.; \forall x \in \operatorname{vec}(a). \\ \quad \operatorname{is\text{-}commutative\text{-}c\text{-}star\text{-}algebra}(a) \Rightarrow  \\ \quad \quad {\operatorname{ccont\text{-}vone}(\operatorname{spec\text{-}top}(a, x))}({x}) = 1_{a} \end{array}}}
\vnbprop{cfc-id}{\vnbstmt{\begin{array}{@{}l@{}} \forall a.; \forall x \in \operatorname{vec}(a). \\ \quad \operatorname{is\text{-}commutative\text{-}c\text{-}star\text{-}algebra}(a) \Rightarrow  \\ \quad \quad (\lambda z \in \sigma({x}).\; z) \in C({\operatorname{spec\text{-}top}(a, x)}) \wedge \\ \quad \quad {(\lambda z \in \sigma({x}).\; z)}({x}) = x \end{array}}}
\vnbprop{cfc-normal-in}{\vnbstmt{\begin{array}{@{}l@{}} \forall a, b.; \forall x \in \operatorname{vec}(a), f \in C({\operatorname{spec\text{-}top}(b, x)}). \\ \quad (\operatorname{is\text{-}commutative\text{-}c\text{-}star\text{-}algebra}(a) \wedge \operatorname{is\text{-}unital\text{-}c\text{-}star\text{-}subalgebra}(a, b)) \Rightarrow  \\ \quad \quad {f}({x}) \in \operatorname{vec}(b) \end{array}}}
\vnbprop{cfc-normal-add}{\vnbstmt{\begin{array}{@{}l@{}} \forall a, b.; \forall x \in \operatorname{vec}(a), f, g \in C({\operatorname{spec\text{-}top}(b, x)}). \\ \quad (\operatorname{is\text{-}commutative\text{-}c\text{-}star\text{-}algebra}(a) \wedge \operatorname{is\text{-}unital\text{-}c\text{-}star\text{-}subalgebra}(a, b)) \Rightarrow  \\ \quad \quad {\operatorname{ccont\text{-}vadd}(\operatorname{spec\text{-}top}(b, x))(f, g)}({x}) = \operatorname{vadd}(b)({f}({x}), {g}({x})) \end{array}}}
\vnbprop{cfc-normal-mul}{\vnbstmt{\begin{array}{@{}l@{}} \forall a, b.; \forall x \in \operatorname{vec}(a), f, g \in C({\operatorname{spec\text{-}top}(b, x)}). \\ \quad (\operatorname{is\text{-}commutative\text{-}c\text{-}star\text{-}algebra}(a) \wedge \operatorname{is\text{-}unital\text{-}c\text{-}star\text{-}subalgebra}(a, b)) \Rightarrow  \\ \quad \quad {\operatorname{ccont\text{-}vmul}(\operatorname{spec\text{-}top}(b, x))(f, g)}({x}) = \cdot_{b}({f}({x}), {g}({x})) \end{array}}}
\vnbprop{cfc-normal-act}{\vnbstmt{\begin{array}{@{}l@{}} \forall a, b.; \forall x \in \operatorname{vec}(a), f \in C({\operatorname{spec\text{-}top}(b, x)}), z \in \mathbb{C}. \\ \quad (\operatorname{is\text{-}commutative\text{-}c\text{-}star\text{-}algebra}(a) \wedge \operatorname{is\text{-}unital\text{-}c\text{-}star\text{-}subalgebra}(a, b)) \Rightarrow  \\ \quad \quad {\operatorname{ccont\text{-}act}(\operatorname{spec\text{-}top}(b, x))(z, f)}({x}) = \operatorname{act}(b)(z, {f}({x})) \end{array}}}
\vnbprop{cfc-normal-star}{\vnbstmt{\begin{array}{@{}l@{}} \forall a, b.; \forall x \in \operatorname{vec}(a), f \in C({\operatorname{spec\text{-}top}(b, x)}). \\ \quad (\operatorname{is\text{-}commutative\text{-}c\text{-}star\text{-}algebra}(a) \wedge \operatorname{is\text{-}unital\text{-}c\text{-}star\text{-}subalgebra}(a, b)) \Rightarrow  \\ \quad \quad {\operatorname{ccont\text{-}vstar}(\operatorname{spec\text{-}top}(b, x))(f)}({x}) = {}^{*_{b}}({f}({x})) \end{array}}}
\vnbprop{cfc-normal-one-id}{\vnbstmt{\begin{array}{@{}l@{}} \forall a, b.; \forall x \in \operatorname{vec}(a). \\ \quad (\operatorname{is\text{-}commutative\text{-}c\text{-}star\text{-}algebra}(a) \wedge \operatorname{is\text{-}unital\text{-}c\text{-}star\text{-}subalgebra}(a, b)) \Rightarrow  \\ \quad \quad (({\operatorname{ccont\text{-}vone}(\operatorname{spec\text{-}top}(b, x))}({x}) = 1_{b}) \wedge ({(\lambda z \in \sigma({x}).\; z)}({x}) = x)) \end{array}}}
\vnbprop{char-vadd}{\vnbstmt{\begin{array}{@{}l@{}} \forall a.; \forall c \in \operatorname{characters}(a), x, y \in \operatorname{vec}(a). \\ \quad \operatorname{is\text{-}unital\text{-}involutive\text{-}banach\text{-}algebra}(a) \Rightarrow (c(\operatorname{vadd}(a)(x, y)) = (c(x) + c(y))) \end{array}}}
\vnbprop{char-act}{\vnbstmt{\begin{array}{@{}l@{}} \forall a.; \forall c \in \operatorname{characters}(a), z \in \mathbb{C}, x \in \operatorname{vec}(a). \\ \quad \operatorname{is\text{-}unital\text{-}involutive\text{-}banach\text{-}algebra}(a) \Rightarrow (c(\operatorname{act}(a)(z, x)) = (z \cdot c(x))) \end{array}}}
\vnbprop{char-vmul}{\vnbstmt{\begin{array}{@{}l@{}} \forall a.; \forall c \in \operatorname{characters}(a), x, y \in \operatorname{vec}(a). \\ \quad \operatorname{is\text{-}unital\text{-}involutive\text{-}banach\text{-}algebra}(a) \Rightarrow (c(\cdot_{a}(x, y)) = (c(x) \cdot c(y))) \end{array}}}
\vnbprop{char-vone}{\vnbstmt{\begin{array}{@{}l@{}} \forall a.; \forall c \in \operatorname{characters}(a). \\ \quad \operatorname{is\text{-}unital\text{-}involutive\text{-}banach\text{-}algebra}(a) \Rightarrow (c(1_{a}) = 1) \end{array}}}
\vnbprop{char-vneg}{\vnbstmt{\begin{array}{@{}l@{}} \forall a.; \forall c \in \operatorname{characters}(a), x \in \operatorname{vec}(a). \\ \quad \operatorname{is\text{-}unital\text{-}involutive\text{-}banach\text{-}algebra}(a) \Rightarrow (c(\operatorname{vneg}(a)(x)) = -c(x)) \end{array}}}
\vnbprop{gelfand-space-nonempty}{\vnbstmt{\begin{array}{@{}l@{}} \forall a. \\ \quad \operatorname{is\text{-}commutative\text{-}c\text{-}star\text{-}algebra}(a) \Rightarrow  \\ \quad \quad \exists p.\; (p \in \operatorname{pts}(\widehat{a})) \wedge \\ \quad \quad \operatorname{is\text{-}commutative\text{-}c\text{-}star\text{-}algebra}(\operatorname{ccont\text{-}algebra}(\widehat{a})) \end{array}}}

\vnbbookfile{theorem-library/gn-jac}{theorem-library/gn-jac}
\noindent{\small gn-jac.scm -- Proposition 1.41 of F. J. Thayer, "Banach Algebras and Applications", ch. 1: in a unital involutive Banach algebra, sigma(xy) and sigma(yx) agree off 0.  Agent BA-1117-GN, wave 2026-11-17. All proofs modulo 0.  Destination: theorem-library/gn-jac.scm, after gn-onto (spectrum-unfold). Library: ...}\par
\vnbprop{jacobson-shift-core}{\vnbstmt{\begin{array}{@{}l@{}} \forall a.; \forall x, y \in \operatorname{vec}(a), z \in \mathbb{C}, u \in \operatorname{vec}(a). \\ \quad \operatorname{is\text{-}unital\text{-}involutive\text{-}banach\text{-}algebra}(a) \wedge \\ \quad \neg (z = 0) \wedge \\ \quad \cdot_{a}(u, \operatorname{shift\text{-}el}(a, z, \cdot_{a}(y, x))) = 1_{a} \wedge \\ \quad \cdot_{a}(\operatorname{shift\text{-}el}(a, z, \cdot_{a}(y, x)), u) = 1_{a} \Rightarrow \\ \quad \quad \operatorname{is\text{-}invertible\text{-}el}(a, \operatorname{shift\text{-}el}(a, z, \cdot_{a}(x, y))) \end{array}}}
\vnbprop{jacobson-shift}{\vnbstmt{\begin{array}{@{}l@{}} \forall a.; \forall x, y \in \operatorname{vec}(a), z \in \mathbb{C}. \\ \quad \operatorname{is\text{-}unital\text{-}involutive\text{-}banach\text{-}algebra}(a) \wedge \\ \quad \neg (z = 0) \wedge \\ \quad \operatorname{is\text{-}invertible\text{-}el}(a, \operatorname{shift\text{-}el}(a, z, \cdot_{a}(y, x))) \Rightarrow \\ \quad \quad \operatorname{is\text{-}invertible\text{-}el}(a, \operatorname{shift\text{-}el}(a, z, \cdot_{a}(x, y))) \end{array}}}
\vnbprop{spectrum-mul-commute-off-zero}{\vnbstmt{\begin{array}{@{}l@{}} \forall a.; \forall x, y \in \operatorname{vec}(a), z \in \mathbb{C}. \\ \quad (\operatorname{is\text{-}unital\text{-}involutive\text{-}banach\text{-}algebra}(a) \wedge \neg (z = 0)) \Rightarrow  \\ \quad \quad (z \in \sigma({\cdot_{a}(x, y)})) \Leftrightarrow (z \in \sigma({\cdot_{a}(y, x)})) \end{array}}}

\vnbbookfile{theorem-library/gn-pos}{theorem-library/gn-pos}
\noindent{\small gn-pos.scm -- positivity through *-homomorphisms, and positivity in a commutative C*-algebra read off the characters (the commutative core of 1.38, 1.39, 1.43 of F. J. Thayer, "Banach Algebras and Applications", ch. 1 s. 5).  Agent BA-1117-GN, wave 2026-11-17.  All proofs modulo 0. Destination: ...}\par
\vnbprop{hom-pos-cone}{\vnbstmt{\begin{array}{@{}l@{}} \forall a, b, f.; \forall x \in {a}_{+}. \\ \quad \operatorname{is\text{-}hom\text{-}unital\text{-}involutive\text{-}banach\text{-}algebra}(a, b, f) \Rightarrow (f(x) \in {b}_{+}) \end{array}}}
\vnbprop{hom-pos-cone-reflect}{\vnbstmt{\begin{array}{@{}l@{}} \forall a, b, f.; \forall x \in \operatorname{vec}(a). \\ \quad \operatorname{is\text{-}hom\text{-}unital\text{-}involutive\text{-}banach\text{-}algebra}(a, b, f) \wedge \\ \quad \forall y \in \operatorname{vec}(b).\; \exists \mathit{gnpx} \in \operatorname{vec}(a).\; (f(\mathit{gnpx}) = y) \wedge \\ \quad \forall \mathit{gnpx}, y \in \operatorname{vec}(a).\; ((f(\mathit{gnpx}) = f(y)) \Rightarrow (\mathit{gnpx} = y)) \wedge \\ \quad f(x) \in {b}_{+} \Rightarrow \\ \quad \quad x \in {a}_{+} \end{array}}}
\vnbprop{gelfand-hom-uiba}{\vnbstmt{\begin{array}{@{}l@{}} \forall a. \\ \quad \operatorname{is\text{-}commutative\text{-}c\text{-}star\text{-}algebra}(a) \Rightarrow  \\ \quad \quad \operatorname{is\text{-}hom\text{-}unital\text{-}involutive\text{-}banach\text{-}algebra}(a, \operatorname{ccont\text{-}algebra}(\widehat{a}), \operatorname{gelfand\text{-}transform}(a)) \end{array}}}
\vnbprop{gelfand-pos-cone-iff}{\vnbstmt{\begin{array}{@{}l@{}} \forall a.; \forall x \in \operatorname{vec}(a). \\ \quad \operatorname{is\text{-}commutative\text{-}c\text{-}star\text{-}algebra}(a) \Rightarrow  \\ \quad \quad (x \in {a}_{+}) \Leftrightarrow (\operatorname{gelfand\text{-}transform}(a)(x) \in {\operatorname{ccont\text{-}algebra}(\widehat{a})}_{+}) \end{array}}}
\vnbprop{gelfand-positive-iff}{\vnbstmt{\begin{array}{@{}l@{}} \forall a.; \forall x \in \operatorname{vec}(a). \\ \quad \operatorname{is\text{-}commutative\text{-}c\text{-}star\text{-}algebra}(a) \Rightarrow  \\ \quad \quad (x \in {a}_{+}) \Leftrightarrow \forall c \in \operatorname{characters}(a).\; ((c(x) \in \mathbb{R}) \wedge (0 \leq c(x))) \end{array}}}
\vnbprop{ale-unfold}{\vnbstmt{\forall a, x, y.\; ({x} \le_{a} {y} \Leftrightarrow (\operatorname{vadd}(a)(y, \operatorname{vneg}(a)(x)) \in {a}_{+}))}}
\vnbprop{ccsa-order-unit-le}{\vnbstmt{\begin{array}{@{}l@{}} \forall a.; \forall x \in \operatorname{sa\text{-}part}(a), l \in \mathbb{R}. \\ \quad (\operatorname{is\text{-}commutative\text{-}c\text{-}star\text{-}algebra}(a) \wedge (\operatorname{vnrm}(a)(x) \leq l)) \Rightarrow  \\ \quad \quad {x} \le_{a} {\operatorname{act}(a)(l, 1_{a})} \end{array}}}
\vnbprop{ccsa-norm-monotone}{\vnbstmt{\begin{array}{@{}l@{}} \forall a.; \forall x, y \in {a}_{+}. \\ \quad (\operatorname{is\text{-}commutative\text{-}c\text{-}star\text{-}algebra}(a) \wedge {x} \le_{a} {y}) \Rightarrow  \\ \quad \quad \operatorname{vnrm}(a)(x) \leq \operatorname{vnrm}(a)(y) \end{array}}}
\vnbprop{ccsa-one-minus-pos-norm}{\vnbstmt{\begin{array}{@{}l@{}} \forall a.; \forall y \in {a}_{+}. \\ \quad (\operatorname{is\text{-}commutative\text{-}c\text{-}star\text{-}algebra}(a) \wedge (\operatorname{vnrm}(a)(y) \leq 1)) \Rightarrow  \\ \quad \quad \operatorname{vnrm}(a)(\operatorname{vadd}(a)(1_{a}, \operatorname{vneg}(a)(y))) \leq 1 \end{array}}}
\vnbprop{ccsa-pos-norm-test}{\vnbstmt{\begin{array}{@{}l@{}} \forall a.; \forall x \in \operatorname{sa\text{-}part}(a). \\ \quad \operatorname{is\text{-}commutative\text{-}c\text{-}star\text{-}algebra}(a) \Rightarrow  \\ \quad \quad x \in {a}_{+} \\ \quad \quad \Leftrightarrow\; \operatorname{vnrm}(a)(\operatorname{vadd}(a)(\operatorname{act}(a)(\operatorname{vnrm}(a)(x), 1_{a}), \operatorname{vneg}(a)(x))) \leq \operatorname{vnrm}(a)(x) \end{array}}}
\vnbprop{ccsa-pos-cone-antisym}{\vnbstmt{\begin{array}{@{}l@{}} \forall a.; \forall x \in {a}_{+}. \\ \quad (\operatorname{is\text{-}commutative\text{-}c\text{-}star\text{-}algebra}(a) \wedge (\operatorname{vneg}(a)(x) \in {a}_{+})) \Rightarrow  \\ \quad \quad x = \operatorname{vzero}(a) \end{array}}}
\vnbprop{ccsa-sa-le-one}{\vnbstmt{\begin{array}{@{}l@{}} \forall a.; \forall x \in \operatorname{sa\text{-}part}(a). \\ \quad (\operatorname{is\text{-}commutative\text{-}c\text{-}star\text{-}algebra}(a) \wedge (\operatorname{vnrm}(a)(x) \leq 1)) \Rightarrow  \\ \quad \quad {x} \le_{a} {1_{a}} \end{array}}}
\vnbprop{sub-pos-cone}{\vnbstmt{\begin{array}{@{}l@{}} \forall a, b.; \forall x \in {a}_{+}. \\ \quad \operatorname{is\text{-}unital\text{-}c\text{-}star\text{-}subalgebra}(a, b) \Rightarrow (x \in {b}_{+}) \end{array}}}
\vnbprop{sub-order-unit-le}{\vnbstmt{\begin{array}{@{}l@{}} \forall a, b.; \forall x \in \operatorname{vec}(a), l \in \mathbb{R}. \\ \quad \operatorname{is\text{-}commutative\text{-}c\text{-}star\text{-}algebra}(a) \wedge \\ \quad \operatorname{is\text{-}unital\text{-}c\text{-}star\text{-}subalgebra}(a, b) \wedge \\ \quad {}^{*_{b}}(x) = x \wedge \\ \quad \operatorname{vnrm}(b)(x) \leq l \Rightarrow \\ \quad \quad {x} \le_{b} {\operatorname{act}(b)(l, 1_{b})} \end{array}}}
\vnbprop{ccsa-pos-sqrt}{\vnbstmt{\begin{array}{@{}l@{}} \forall a.; \forall x \in {a}_{+}. \\ \quad \operatorname{is\text{-}commutative\text{-}c\text{-}star\text{-}algebra}(a) \Rightarrow  \\ \quad \quad \exists h \in {a}_{+}.\; ((h \in \operatorname{sa\text{-}part}(a)) \wedge (\cdot_{a}(h, h) = x)) \end{array}}}
\vnbprop{sa-square-pos}{\vnbstmt{\begin{array}{@{}l@{}} \forall a.; \forall x \in \operatorname{sa\text{-}squares}(a). \\ \quad \operatorname{is\text{-}unital\text{-}involutive\text{-}banach\text{-}algebra}(a) \Rightarrow (x \in {a}_{+}) \end{array}}}
\vnbprop{ccsa-pos-cone-sa-squares}{\vnbstmt{\begin{array}{@{}l@{}} \forall a.; \forall x \in \operatorname{vec}(a). \\ \quad \operatorname{is\text{-}commutative\text{-}c\text{-}star\text{-}algebra}(a) \Rightarrow ((x \in {a}_{+}) \Leftrightarrow (x \in \operatorname{sa\text{-}squares}(a))) \end{array}}}
\vnbprop{ccsa-one-minus-sa-square}{\vnbstmt{\begin{array}{@{}l@{}} \forall a.; \forall x \in \operatorname{sa\text{-}part}(a). \\ \quad (\operatorname{is\text{-}commutative\text{-}c\text{-}star\text{-}algebra}(a) \wedge (\operatorname{vnrm}(a)(x) \leq 1)) \Rightarrow  \\ \quad \quad \operatorname{vadd}(a)(1_{a}, \operatorname{vneg}(a)(x)) \in \operatorname{sa\text{-}squares}(a) \end{array}}}
\vnbprop{sub-one-minus-sa-square}{\vnbstmt{\begin{array}{@{}l@{}} \forall a, b.; \forall x \in \operatorname{vec}(a). \\ \quad \operatorname{is\text{-}commutative\text{-}c\text{-}star\text{-}algebra}(a) \wedge \\ \quad \operatorname{is\text{-}unital\text{-}c\text{-}star\text{-}subalgebra}(a, b) \wedge \\ \quad {}^{*_{b}}(x) = x \wedge \\ \quad \operatorname{vnrm}(b)(x) \leq 1 \Rightarrow \\ \quad \quad \operatorname{vadd}(b)(1_{b}, \operatorname{vneg}(b)(x)) \in \operatorname{sa\text{-}squares}(b) \end{array}}}
\vnbprop{sub-pos-norm-test-suff}{\vnbstmt{\begin{array}{@{}l@{}} \forall a, b.; \forall x \in \operatorname{vec}(a). \\ \quad \operatorname{is\text{-}commutative\text{-}c\text{-}star\text{-}algebra}(a) \wedge \\ \quad \operatorname{is\text{-}unital\text{-}c\text{-}star\text{-}subalgebra}(a, b) \wedge \\ \quad {}^{*_{b}}(x) = x \wedge \\ \quad \operatorname{vnrm}(b)(\operatorname{vadd}(b)(\operatorname{act}(b)(\operatorname{vnrm}(b)(x), 1_{b}), \operatorname{vneg}(b)(x))) \leq \operatorname{vnrm}(b)(x) \Rightarrow \\ \quad \quad x \in {b}_{+} \end{array}}}
\vnbprop{sub-pos-norm-test-nec}{\vnbstmt{\begin{array}{@{}l@{}} \forall a, b.; \forall x \in {a}_{+}. \\ \quad (\operatorname{is\text{-}commutative\text{-}c\text{-}star\text{-}algebra}(a) \wedge \operatorname{is\text{-}unital\text{-}c\text{-}star\text{-}subalgebra}(a, b)) \Rightarrow  \\ \quad \quad \operatorname{vnrm}(b)(\operatorname{vadd}(b)(\operatorname{act}(b)(\operatorname{vnrm}(b)(x), 1_{b}), \operatorname{vneg}(b)(x))) \leq \operatorname{vnrm}(b)(x) \end{array}}}
\vnbprop{sub-sa-squares}{\vnbstmt{\begin{array}{@{}l@{}} \forall a, b.; \forall x \in \operatorname{sa\text{-}squares}(a). \\ \quad \operatorname{is\text{-}unital\text{-}c\text{-}star\text{-}subalgebra}(a, b) \Rightarrow (x \in \operatorname{sa\text{-}squares}(b)) \end{array}}}
\vnbprop{sub-pos-norm-test-sq}{\vnbstmt{\begin{array}{@{}l@{}} \forall a, b.; \forall x \in \operatorname{vec}(a). \\ \quad \operatorname{is\text{-}commutative\text{-}c\text{-}star\text{-}algebra}(a) \wedge \\ \quad \operatorname{is\text{-}unital\text{-}c\text{-}star\text{-}subalgebra}(a, b) \wedge \\ \quad {}^{*_{b}}(x) = x \wedge \\ \quad \operatorname{vnrm}(b)(\operatorname{vadd}(b)(\operatorname{act}(b)(\operatorname{vnrm}(b)(x), 1_{b}), \operatorname{vneg}(b)(x))) \leq \operatorname{vnrm}(b)(x) \Rightarrow \\ \quad \quad x \in \operatorname{sa\text{-}squares}(b) \end{array}}}
\vnbprop{sub-one-minus-pos-norm}{\vnbstmt{\begin{array}{@{}l@{}} \forall a, b.; \forall y \in {a}_{+}. \\ \quad (\operatorname{is\text{-}commutative\text{-}c\text{-}star\text{-}algebra}(a) \wedge (\operatorname{is\text{-}unital\text{-}c\text{-}star\text{-}subalgebra}(a, b) \wedge (\operatorname{vnrm}(b)(y) \leq 1))) \Rightarrow  \\ \quad \quad \operatorname{vnrm}(b)(\operatorname{vadd}(b)(1_{b}, \operatorname{vneg}(b)(y))) \leq 1 \end{array}}}

\vnbbookfile{theorem-library/gn-dual}{theorem-library/gn-dual}
\noindent{\small gn-dual.scm -- Proposition 1.26 of F. J. Thayer, "Banach Algebras and Applications", ch. 1: a unital algebra homomorphism f : A -$>$ B induces the continuous map phi-hat : B-hat -$>$ A-hat, omega |-$>$ omega o f. Agent BA-1117-GN, wave 2026-11-17.  All proofs modulo 0.  Destination: theorem-library/gn-dual.scm, after ...}\par
\vnbprop{dual-map-character}{\vnbstmt{\begin{array}{@{}l@{}} \forall a, b, f.; \forall w \in \operatorname{characters}(b). \\ \quad \operatorname{is\text{-}unital\text{-}banach\text{-}algebra}(a) \wedge \\ \quad \operatorname{is\text{-}unital\text{-}banach\text{-}algebra}(b) \wedge \\ \quad f \in (\operatorname{vec}(a) \to \operatorname{vec}(b)) \wedge \\ \quad \forall x, y \in \operatorname{vec}(a).\; (f(\operatorname{vadd}(a)(x, y)) = \operatorname{vadd}(b)(f(x), f(y))) \wedge \\ \quad \forall z \in \mathbb{C}, x \in \operatorname{vec}(a).\; (f(\operatorname{act}(a)(z, x)) = \operatorname{act}(b)(z, f(x))) \wedge \\ \quad \forall x, y \in \operatorname{vec}(a).\; (f(\cdot_{a}(x, y)) = \cdot_{b}(f(x), f(y))) \wedge \\ \quad f(1_{a}) = 1_{b} \Rightarrow \\ \quad \quad \operatorname{compose}(w, f) \in \operatorname{characters}(a) \end{array}}}
\vnbprop{gelfand-space-unfold}{\vnbstmt{\begin{array}{@{}l@{}} \forall a. \\ \quad (\widehat{a} \simeq \operatorname{subspace\text{-}top}(\operatorname{tproduct\text{-}top}(\operatorname{vec}(a), (\lambda x \in \operatorname{vec}(a).\; \operatorname{metric\text{-}top}(\operatorname{cc\text{-}ms}))), \operatorname{characters}(a))) \end{array}}}
\vnbprop{dual-map-unfold}{\vnbstmt{\forall b, f.\; (\widehat{{f}} \simeq (\lambda w \in \operatorname{characters}(b).\; \operatorname{compose}(w, f)))}}
\vnbprop{dual-map-continuous}{\vnbstmt{\begin{array}{@{}l@{}} \forall a, b, f. \\ \quad \operatorname{is\text{-}unital\text{-}banach\text{-}algebra}(a) \wedge \\ \quad \operatorname{is\text{-}unital\text{-}banach\text{-}algebra}(b) \wedge \\ \quad f \in (\operatorname{vec}(a) \to \operatorname{vec}(b)) \wedge \\ \quad \forall x, y \in \operatorname{vec}(a).\; (f(\operatorname{vadd}(a)(x, y)) = \operatorname{vadd}(b)(f(x), f(y))) \wedge \\ \quad \forall z \in \mathbb{C}, x \in \operatorname{vec}(a).\; (f(\operatorname{act}(a)(z, x)) = \operatorname{act}(b)(z, f(x))) \wedge \\ \quad \forall x, y \in \operatorname{vec}(a).\; (f(\cdot_{a}(x, y)) = \cdot_{b}(f(x), f(y))) \wedge \\ \quad f(1_{a}) = 1_{b} \Rightarrow \\ \quad \quad \forall w \in \operatorname{characters}(b).\; (\widehat{{f}}(w) = \operatorname{compose}(w, f)) \wedge \\ \quad \quad \operatorname{is\text{-}hom\text{-}top\text{-}space}(\widehat{b}, \widehat{a}, \widehat{{f}}) \end{array}}}
\vnbprop{dual-pull-unfold}{\vnbstmt{\begin{array}{@{}l@{}} \forall a, b, g. \\ \quad (\operatorname{dual\text{-}pull}(a, b, g) \simeq (\lambda x \in \operatorname{vec}(a).\; \varepsilon(\{y \in \operatorname{vec}(b) : (\operatorname{gelfand\text{-}transform}(b)(y) = \operatorname{compose}(\operatorname{gelfand\text{-}transform}(a)(x), g))\}))) \end{array}}}
\vnbprop{dual-pull-spec}{\vnbstmt{\begin{array}{@{}l@{}} \forall a, b, g.; \forall x \in \operatorname{vec}(a). \\ \quad \operatorname{is\text{-}commutative\text{-}c\text{-}star\text{-}algebra}(a) \wedge \\ \quad \operatorname{is\text{-}commutative\text{-}c\text{-}star\text{-}algebra}(b) \wedge \\ \quad \operatorname{is\text{-}hom\text{-}top\text{-}space}(\widehat{b}, \widehat{a}, g) \Rightarrow \\ \quad \quad \operatorname{dual\text{-}pull}(a, b, g)(x) \in \operatorname{vec}(b) \wedge \\ \quad \quad \operatorname{gelfand\text{-}transform}(b)(\operatorname{dual\text{-}pull}(a, b, g)(x)) = \operatorname{compose}(\operatorname{gelfand\text{-}transform}(a)(x), g) \end{array}}}
\vnbprop{dual-pull-value}{\vnbstmt{\begin{array}{@{}l@{}} \forall a, b, g.; \forall x \in \operatorname{vec}(a), w \in \operatorname{characters}(b). \\ \quad \operatorname{is\text{-}commutative\text{-}c\text{-}star\text{-}algebra}(a) \wedge \\ \quad \operatorname{is\text{-}commutative\text{-}c\text{-}star\text{-}algebra}(b) \wedge \\ \quad \operatorname{is\text{-}hom\text{-}top\text{-}space}(\widehat{b}, \widehat{a}, g) \Rightarrow \\ \quad \quad w(\operatorname{dual\text{-}pull}(a, b, g)(x)) = g(w)(x) \end{array}}}
\vnbprop{dual-pull-fun}{\vnbstmt{\begin{array}{@{}l@{}} \forall a, b, g. \\ \quad \operatorname{is\text{-}commutative\text{-}c\text{-}star\text{-}algebra}(a) \wedge \\ \quad \operatorname{is\text{-}commutative\text{-}c\text{-}star\text{-}algebra}(b) \wedge \\ \quad \operatorname{is\text{-}hom\text{-}top\text{-}space}(\widehat{b}, \widehat{a}, g) \Rightarrow \\ \quad \quad \operatorname{dual\text{-}pull}(a, b, g) \in (\operatorname{vec}(a) \to \operatorname{vec}(b)) \end{array}}}
\vnbprop{dual-pull-vadd}{\vnbstmt{\begin{array}{@{}l@{}} \forall a, b, g.; \forall x, y \in \operatorname{vec}(a). \\ \quad \operatorname{is\text{-}commutative\text{-}c\text{-}star\text{-}algebra}(a) \wedge \\ \quad \operatorname{is\text{-}commutative\text{-}c\text{-}star\text{-}algebra}(b) \wedge \\ \quad \operatorname{is\text{-}hom\text{-}top\text{-}space}(\widehat{b}, \widehat{a}, g) \Rightarrow \\ \quad \quad \operatorname{dual\text{-}pull}(a, b, g)(\operatorname{vadd}(a)(x, y)) = \operatorname{vadd}(b)(\operatorname{dual\text{-}pull}(a, b, g)(x), \operatorname{dual\text{-}pull}(a, b, g)(y)) \end{array}}}
\vnbprop{dual-pull-vmul}{\vnbstmt{\begin{array}{@{}l@{}} \forall a, b, g.; \forall x, y \in \operatorname{vec}(a). \\ \quad \operatorname{is\text{-}commutative\text{-}c\text{-}star\text{-}algebra}(a) \wedge \\ \quad \operatorname{is\text{-}commutative\text{-}c\text{-}star\text{-}algebra}(b) \wedge \\ \quad \operatorname{is\text{-}hom\text{-}top\text{-}space}(\widehat{b}, \widehat{a}, g) \Rightarrow \\ \quad \quad \operatorname{dual\text{-}pull}(a, b, g)(\cdot_{a}(x, y)) = \cdot_{b}(\operatorname{dual\text{-}pull}(a, b, g)(x), \operatorname{dual\text{-}pull}(a, b, g)(y)) \end{array}}}
\vnbprop{dual-pull-act}{\vnbstmt{\begin{array}{@{}l@{}} \forall a, b, g.; \forall z \in \mathbb{C}, x \in \operatorname{vec}(a). \\ \quad \operatorname{is\text{-}commutative\text{-}c\text{-}star\text{-}algebra}(a) \wedge \\ \quad \operatorname{is\text{-}commutative\text{-}c\text{-}star\text{-}algebra}(b) \wedge \\ \quad \operatorname{is\text{-}hom\text{-}top\text{-}space}(\widehat{b}, \widehat{a}, g) \Rightarrow \\ \quad \quad \operatorname{dual\text{-}pull}(a, b, g)(\operatorname{act}(a)(z, x)) = \operatorname{act}(b)(z, \operatorname{dual\text{-}pull}(a, b, g)(x)) \end{array}}}
\vnbprop{dual-pull-star}{\vnbstmt{\begin{array}{@{}l@{}} \forall a, b, g.; \forall x \in \operatorname{vec}(a). \\ \quad \operatorname{is\text{-}commutative\text{-}c\text{-}star\text{-}algebra}(a) \wedge \\ \quad \operatorname{is\text{-}commutative\text{-}c\text{-}star\text{-}algebra}(b) \wedge \\ \quad \operatorname{is\text{-}hom\text{-}top\text{-}space}(\widehat{b}, \widehat{a}, g) \Rightarrow \\ \quad \quad \operatorname{dual\text{-}pull}(a, b, g)({}^{*_{a}}(x)) = {}^{*_{b}}(\operatorname{dual\text{-}pull}(a, b, g)(x)) \end{array}}}
\vnbprop{dual-pull-one}{\vnbstmt{\begin{array}{@{}l@{}} \forall a, b, g. \\ \quad \operatorname{is\text{-}commutative\text{-}c\text{-}star\text{-}algebra}(a) \wedge \\ \quad \operatorname{is\text{-}commutative\text{-}c\text{-}star\text{-}algebra}(b) \wedge \\ \quad \operatorname{is\text{-}hom\text{-}top\text{-}space}(\widehat{b}, \widehat{a}, g) \Rightarrow \\ \quad \quad \operatorname{dual\text{-}pull}(a, b, g)(1_{a}) = 1_{b} \end{array}}}
\vnbprop{dual-pull-dual}{\vnbstmt{\begin{array}{@{}l@{}} \forall a, b, g.; \forall w \in \operatorname{characters}(b). \\ \quad \operatorname{is\text{-}commutative\text{-}c\text{-}star\text{-}algebra}(a) \wedge \\ \quad \operatorname{is\text{-}commutative\text{-}c\text{-}star\text{-}algebra}(b) \wedge \\ \quad \operatorname{is\text{-}hom\text{-}top\text{-}space}(\widehat{b}, \widehat{a}, g) \Rightarrow \\ \quad \quad \widehat{{\operatorname{dual\text{-}pull}(a, b, g)}}(w) = g(w) \end{array}}}
\vnbprop{dual-pull-unique}{\vnbstmt{\begin{array}{@{}l@{}} \forall a, b, g.; \forall f \in (\operatorname{vec}(a) \to \operatorname{vec}(b)), x \in \operatorname{vec}(a). \\ \quad \operatorname{is\text{-}commutative\text{-}c\text{-}star\text{-}algebra}(a) \wedge \\ \quad \operatorname{is\text{-}commutative\text{-}c\text{-}star\text{-}algebra}(b) \wedge \\ \quad \operatorname{is\text{-}hom\text{-}top\text{-}space}(\widehat{b}, \widehat{a}, g) \wedge \\ \quad \forall w \in \operatorname{characters}(b), x \in \operatorname{vec}(a).\; (w(f(x)) = g(w)(x)) \Rightarrow \\ \quad \quad f(x) = \operatorname{dual\text{-}pull}(a, b, g)(x) \end{array}}}
\vnbprop{dual-pull-injective-of-onto}{\vnbstmt{\begin{array}{@{}l@{}} \forall a, b, g.; \forall x, y \in \operatorname{vec}(a). \\ \quad \operatorname{is\text{-}commutative\text{-}c\text{-}star\text{-}algebra}(a) \wedge \\ \quad \operatorname{is\text{-}commutative\text{-}c\text{-}star\text{-}algebra}(b) \wedge \\ \quad \operatorname{is\text{-}hom\text{-}top\text{-}space}(\widehat{b}, \widehat{a}, g) \wedge \\ \quad \forall c \in \operatorname{characters}(a).\; \exists w \in \operatorname{characters}(b).\; (g(w) = c) \wedge \\ \quad \operatorname{dual\text{-}pull}(a, b, g)(x) = \operatorname{dual\text{-}pull}(a, b, g)(y) \Rightarrow \\ \quad \quad x = y \end{array}}}

\vnbbookfile{theorem-library/gn-cgen}{theorem-library/gn-cgen}
\noindent{\small gn-cgen.scm -- the C*-subalgebra C*(x, 1) generated by a normal element x of a unital C*-algebra b, as a STRUCTURE: CSTAR-GEN(b, x) (ba2b-gen-vocab) is a norm-closed unital *-subalgebra set (cstar-gen-closed-star-subalgebra); for any such set S, SUB-UCSA(b, S) (gn-vocab) is a unital C*-algebra and a unital ...}\par
\vnbprop{cstar-gen-mem-elim}{\vnbstmt{\begin{array}{@{}l@{}} \forall b, x.; \forall y \in \operatorname{cstar\text{-}gen}(b, x), s. \\ \quad (\operatorname{is\text{-}unital\text{-}c\text{-}star\text{-}algebra}(b) \wedge ((x \in \operatorname{vec}(b)) \wedge (\operatorname{is\text{-}closed\text{-}star\text{-}subalgebra\text{-}set}(b, s) \wedge (x \in s)))) \Rightarrow  \\ \quad \quad y \in s \end{array}}}
\vnbprop{cstar-gen-mem-intro}{\vnbstmt{\begin{array}{@{}l@{}} \forall b, x, y. \\ \quad \operatorname{is\text{-}unital\text{-}c\text{-}star\text{-}algebra}(b) \wedge \\ \quad x \in \operatorname{vec}(b) \wedge \\ \quad \forall s.\; ((\operatorname{is\text{-}closed\text{-}star\text{-}subalgebra\text{-}set}(b, s) \wedge (x \in s)) \Rightarrow (y \in s)) \Rightarrow \\ \quad \quad y \in \operatorname{cstar\text{-}gen}(b, x) \end{array}}}
\vnbprop{cstar-gen-closed-star-subalgebra}{\vnbstmt{\begin{array}{@{}l@{}} \forall b, x. \\ \quad (\operatorname{is\text{-}unital\text{-}c\text{-}star\text{-}algebra}(b) \wedge (x \in \operatorname{vec}(b))) \Rightarrow  \\ \quad \quad \operatorname{is\text{-}closed\text{-}star\text{-}subalgebra\text{-}set}(b, \operatorname{cstar\text{-}gen}(b, x)) \end{array}}}
\vnbprop{commutant-el-unfold}{\vnbstmt{\begin{array}{@{}l@{}} \forall b, t. \\ \quad (\operatorname{commutant\text{-}el}(b, t) \simeq \{y \in \operatorname{vec}(b) : \forall z \in t.\; (\cdot_{b}(y, z) = \cdot_{b}(z, y))\}) \end{array}}}
\vnbprop{commutant-el-closed-star}{\vnbstmt{\begin{array}{@{}l@{}} \forall b, t. \\ \quad (\operatorname{is\text{-}unital\text{-}involutive\text{-}banach\text{-}algebra}(b) \wedge ((t \subseteq \operatorname{vec}(b)) \wedge \forall z \in t.\; ({}^{*_{b}}(z) \in t))) \Rightarrow  \\ \quad \quad \operatorname{is\text{-}closed\text{-}star\text{-}subalgebra\text{-}set}(b, \operatorname{commutant\text{-}el}(b, t)) \end{array}}}
\vnbprop{cstar-gen-commutative}{\vnbstmt{\begin{array}{@{}l@{}} \forall b, x.; \forall y, z \in \operatorname{cstar\text{-}gen}(b, x). \\ \quad (\operatorname{is\text{-}unital\text{-}c\text{-}star\text{-}algebra}(b) \wedge \operatorname{is\text{-}normal\text{-}el}(b, x)) \Rightarrow  \\ \quad \quad \cdot_{b}(y, z) = \cdot_{b}(z, y) \end{array}}}
\vnbprop{sub-ucsa-vadd-apply}{\vnbstmt{\begin{array}{@{}l@{}} \forall b, k.; \forall u, v \in k. \\ \quad (\operatorname{is\text{-}unital\text{-}c\text{-}star\text{-}algebra}(b) \wedge \operatorname{is\text{-}closed\text{-}star\text{-}subalgebra\text{-}set}(b, k)) \Rightarrow  \\ \quad \quad (\lambda {[}\operatorname{gnvu\_}, \operatorname{gnvv\_}] \in (k \times k).\; \operatorname{vadd}(b)(\operatorname{gnvu\_}, \operatorname{gnvv\_}))(u, v) = \operatorname{vadd}(b)(u, v) \end{array}}}
\vnbprop{sub-ucsa-vmul-apply}{\vnbstmt{\begin{array}{@{}l@{}} \forall b, k.; \forall u, v \in k. \\ \quad (\operatorname{is\text{-}unital\text{-}c\text{-}star\text{-}algebra}(b) \wedge \operatorname{is\text{-}closed\text{-}star\text{-}subalgebra\text{-}set}(b, k)) \Rightarrow  \\ \quad \quad (\lambda {[}\operatorname{gnvu\_}, \operatorname{gnvv\_}] \in (k \times k).\; \cdot_{b}(\operatorname{gnvu\_}, \operatorname{gnvv\_}))(u, v) = \cdot_{b}(u, v) \end{array}}}
\vnbprop{sub-ucsa-act-apply}{\vnbstmt{\begin{array}{@{}l@{}} \forall b, k.; \forall c \in \mathbb{C}, u \in k. \\ \quad (\operatorname{is\text{-}unital\text{-}c\text{-}star\text{-}algebra}(b) \wedge \operatorname{is\text{-}closed\text{-}star\text{-}subalgebra\text{-}set}(b, k)) \Rightarrow  \\ \quad \quad (\lambda {[}\operatorname{gnvc\_}, \operatorname{gnvu\_}] \in (\mathbb{C} \times k).\; \operatorname{act}(b)(\operatorname{gnvc\_}, \operatorname{gnvu\_}))(c, u) = \operatorname{act}(b)(c, u) \end{array}}}
\vnbprop{sub-ucsa-vneg-apply}{\vnbstmt{\begin{array}{@{}l@{}} \forall b, k.; \forall u \in k. \\ \quad (\operatorname{is\text{-}unital\text{-}c\text{-}star\text{-}algebra}(b) \wedge \operatorname{is\text{-}closed\text{-}star\text{-}subalgebra\text{-}set}(b, k)) \Rightarrow  \\ \quad \quad (\lambda \mathit{gnvu} \in k.\; \operatorname{vneg}(b)(\mathit{gnvu}))(u) = \operatorname{vneg}(b)(u) \end{array}}}
\vnbprop{sub-ucsa-vnrm-apply}{\vnbstmt{\begin{array}{@{}l@{}} \forall b, k.; \forall u \in k. \\ \quad (\operatorname{is\text{-}unital\text{-}c\text{-}star\text{-}algebra}(b) \wedge \operatorname{is\text{-}closed\text{-}star\text{-}subalgebra\text{-}set}(b, k)) \Rightarrow  \\ \quad \quad (\lambda \mathit{gnvu} \in k.\; \operatorname{vnrm}(b)(\mathit{gnvu}))(u) = \operatorname{vnrm}(b)(u) \end{array}}}
\vnbprop{sub-ucsa-vstar-apply}{\vnbstmt{\begin{array}{@{}l@{}} \forall b, k.; \forall u \in k. \\ \quad (\operatorname{is\text{-}unital\text{-}c\text{-}star\text{-}algebra}(b) \wedge \operatorname{is\text{-}closed\text{-}star\text{-}subalgebra\text{-}set}(b, k)) \Rightarrow  \\ \quad \quad (\lambda \mathit{gnvu} \in k.\; {}^{*_{b}}(\mathit{gnvu}))(u) = {}^{*_{b}}(u) \end{array}}}
\vnbprop{sub-ucsa-tuple8-scal}{\vnbstmt{\begin{array}{@{}l@{}} \forall b, k. \\ \quad (\operatorname{is\text{-}unital\text{-}c\text{-}star\text{-}algebra}(b) \wedge \operatorname{is\text{-}closed\text{-}star\text{-}subalgebra\text{-}set}(b, k)) \Rightarrow  \\ \quad \quad (\operatorname{scal}({[}\operatorname{normed\text{-}field\text{-}as\text{-}commutative\text{-}ring}(\operatorname{cc\text{-}normed\text{-}field}), k, (\lambda {[}u, \operatorname{gnvv\_}] \in (k \times k).\; \operatorname{vadd}(b)(u, \operatorname{gnvv\_})), \operatorname{vzero}(b), (\lambda u \in k.\; \operatorname{vneg}(b)(u)), (\lambda {[}\operatorname{gnvc\_}, u] \in (\mathbb{C} \times k).\; \operatorname{act}(b)(\operatorname{gnvc\_}, u)), (\lambda u \in k.\; \operatorname{vnrm}(b)(u)), (\lambda {[}u, \operatorname{gnvv\_}] \in (k \times k).\; \cdot_{b}(u, \operatorname{gnvv\_}))]) \simeq \operatorname{normed\text{-}field\text{-}as\text{-}commutative\text{-}ring}(\operatorname{cc\text{-}normed\text{-}field})) \end{array}}}
\vnbprop{sub-ucsa-tuple8-vec}{\vnbstmt{\begin{array}{@{}l@{}} \forall b, k. \\ \quad (\operatorname{is\text{-}unital\text{-}c\text{-}star\text{-}algebra}(b) \wedge \operatorname{is\text{-}closed\text{-}star\text{-}subalgebra\text{-}set}(b, k)) \Rightarrow  \\ \quad \quad (\operatorname{vec}({[}\operatorname{normed\text{-}field\text{-}as\text{-}commutative\text{-}ring}(\operatorname{cc\text{-}normed\text{-}field}), k, (\lambda {[}u, \operatorname{gnvv\_}] \in (k \times k).\; \operatorname{vadd}(b)(u, \operatorname{gnvv\_})), \operatorname{vzero}(b), (\lambda u \in k.\; \operatorname{vneg}(b)(u)), (\lambda {[}\operatorname{gnvc\_}, u] \in (\mathbb{C} \times k).\; \operatorname{act}(b)(\operatorname{gnvc\_}, u)), (\lambda u \in k.\; \operatorname{vnrm}(b)(u)), (\lambda {[}u, \operatorname{gnvv\_}] \in (k \times k).\; \cdot_{b}(u, \operatorname{gnvv\_}))]) \simeq k) \end{array}}}
\vnbprop{sub-ucsa-tuple8-vadd}{\vnbstmt{\begin{array}{@{}l@{}} \forall b, k. \\ \quad (\operatorname{is\text{-}unital\text{-}c\text{-}star\text{-}algebra}(b) \wedge \operatorname{is\text{-}closed\text{-}star\text{-}subalgebra\text{-}set}(b, k)) \Rightarrow  \\ \quad \quad (\operatorname{vadd}({[}\operatorname{normed\text{-}field\text{-}as\text{-}commutative\text{-}ring}(\operatorname{cc\text{-}normed\text{-}field}), k, (\lambda {[}u, \operatorname{gnvv\_}] \in (k \times k).\; \operatorname{vadd}(b)(u, \operatorname{gnvv\_})), \operatorname{vzero}(b), (\lambda u \in k.\; \operatorname{vneg}(b)(u)), (\lambda {[}\operatorname{gnvc\_}, u] \in (\mathbb{C} \times k).\; \operatorname{act}(b)(\operatorname{gnvc\_}, u)), (\lambda u \in k.\; \operatorname{vnrm}(b)(u)), (\lambda {[}u, \operatorname{gnvv\_}] \in (k \times k).\; \cdot_{b}(u, \operatorname{gnvv\_}))]) \simeq (\lambda {[}u, \operatorname{gnvv\_}] \in (k \times k).\; \operatorname{vadd}(b)(u, \operatorname{gnvv\_}))) \end{array}}}
\vnbprop{sub-ucsa-tuple8-vzero}{\vnbstmt{\begin{array}{@{}l@{}} \forall b, k. \\ \quad (\operatorname{is\text{-}unital\text{-}c\text{-}star\text{-}algebra}(b) \wedge \operatorname{is\text{-}closed\text{-}star\text{-}subalgebra\text{-}set}(b, k)) \Rightarrow  \\ \quad \quad (\operatorname{vzero}({[}\operatorname{normed\text{-}field\text{-}as\text{-}commutative\text{-}ring}(\operatorname{cc\text{-}normed\text{-}field}), k, (\lambda {[}u, \operatorname{gnvv\_}] \in (k \times k).\; \operatorname{vadd}(b)(u, \operatorname{gnvv\_})), \operatorname{vzero}(b), (\lambda u \in k.\; \operatorname{vneg}(b)(u)), (\lambda {[}\operatorname{gnvc\_}, u] \in (\mathbb{C} \times k).\; \operatorname{act}(b)(\operatorname{gnvc\_}, u)), (\lambda u \in k.\; \operatorname{vnrm}(b)(u)), (\lambda {[}u, \operatorname{gnvv\_}] \in (k \times k).\; \cdot_{b}(u, \operatorname{gnvv\_}))]) \simeq \operatorname{vzero}(b)) \end{array}}}
\vnbprop{sub-ucsa-tuple8-vneg}{\vnbstmt{\begin{array}{@{}l@{}} \forall b, k. \\ \quad (\operatorname{is\text{-}unital\text{-}c\text{-}star\text{-}algebra}(b) \wedge \operatorname{is\text{-}closed\text{-}star\text{-}subalgebra\text{-}set}(b, k)) \Rightarrow  \\ \quad \quad (\operatorname{vneg}({[}\operatorname{normed\text{-}field\text{-}as\text{-}commutative\text{-}ring}(\operatorname{cc\text{-}normed\text{-}field}), k, (\lambda {[}u, \operatorname{gnvv\_}] \in (k \times k).\; \operatorname{vadd}(b)(u, \operatorname{gnvv\_})), \operatorname{vzero}(b), (\lambda u \in k.\; \operatorname{vneg}(b)(u)), (\lambda {[}\operatorname{gnvc\_}, u] \in (\mathbb{C} \times k).\; \operatorname{act}(b)(\operatorname{gnvc\_}, u)), (\lambda u \in k.\; \operatorname{vnrm}(b)(u)), (\lambda {[}u, \operatorname{gnvv\_}] \in (k \times k).\; \cdot_{b}(u, \operatorname{gnvv\_}))]) \simeq (\lambda u \in k.\; \operatorname{vneg}(b)(u))) \end{array}}}
\vnbprop{sub-ucsa-tuple8-act}{\vnbstmt{\begin{array}{@{}l@{}} \forall b, k. \\ \quad (\operatorname{is\text{-}unital\text{-}c\text{-}star\text{-}algebra}(b) \wedge \operatorname{is\text{-}closed\text{-}star\text{-}subalgebra\text{-}set}(b, k)) \Rightarrow  \\ \quad \quad (\operatorname{act}({[}\operatorname{normed\text{-}field\text{-}as\text{-}commutative\text{-}ring}(\operatorname{cc\text{-}normed\text{-}field}), k, (\lambda {[}u, \operatorname{gnvv\_}] \in (k \times k).\; \operatorname{vadd}(b)(u, \operatorname{gnvv\_})), \operatorname{vzero}(b), (\lambda u \in k.\; \operatorname{vneg}(b)(u)), (\lambda {[}\operatorname{gnvc\_}, u] \in (\mathbb{C} \times k).\; \operatorname{act}(b)(\operatorname{gnvc\_}, u)), (\lambda u \in k.\; \operatorname{vnrm}(b)(u)), (\lambda {[}u, \operatorname{gnvv\_}] \in (k \times k).\; \cdot_{b}(u, \operatorname{gnvv\_}))]) \simeq (\lambda {[}\operatorname{gnvc\_}, u] \in (\mathbb{C} \times k).\; \operatorname{act}(b)(\operatorname{gnvc\_}, u))) \end{array}}}
\vnbprop{sub-ucsa-tuple8-vnrm}{\vnbstmt{\begin{array}{@{}l@{}} \forall b, k. \\ \quad (\operatorname{is\text{-}unital\text{-}c\text{-}star\text{-}algebra}(b) \wedge \operatorname{is\text{-}closed\text{-}star\text{-}subalgebra\text{-}set}(b, k)) \Rightarrow  \\ \quad \quad (\operatorname{vnrm}({[}\operatorname{normed\text{-}field\text{-}as\text{-}commutative\text{-}ring}(\operatorname{cc\text{-}normed\text{-}field}), k, (\lambda {[}u, \operatorname{gnvv\_}] \in (k \times k).\; \operatorname{vadd}(b)(u, \operatorname{gnvv\_})), \operatorname{vzero}(b), (\lambda u \in k.\; \operatorname{vneg}(b)(u)), (\lambda {[}\operatorname{gnvc\_}, u] \in (\mathbb{C} \times k).\; \operatorname{act}(b)(\operatorname{gnvc\_}, u)), (\lambda u \in k.\; \operatorname{vnrm}(b)(u)), (\lambda {[}u, \operatorname{gnvv\_}] \in (k \times k).\; \cdot_{b}(u, \operatorname{gnvv\_}))]) \simeq (\lambda u \in k.\; \operatorname{vnrm}(b)(u))) \end{array}}}
\vnbprop{sub-ucsa-tuple8-vmul}{\vnbstmt{\begin{array}{@{}l@{}} \forall b, k. \\ \quad (\operatorname{is\text{-}unital\text{-}c\text{-}star\text{-}algebra}(b) \wedge \operatorname{is\text{-}closed\text{-}star\text{-}subalgebra\text{-}set}(b, k)) \Rightarrow  \\ \quad \quad (\cdot_{{[}\operatorname{normed\text{-}field\text{-}as\text{-}commutative\text{-}ring}(\operatorname{cc\text{-}normed\text{-}field}), k, (\lambda {[}u, \operatorname{gnvv\_}] \in (k \times k).\; \operatorname{vadd}(b)(u, \operatorname{gnvv\_})), \operatorname{vzero}(b), (\lambda u \in k.\; \operatorname{vneg}(b)(u)), (\lambda {[}\operatorname{gnvc\_}, u] \in (\mathbb{C} \times k).\; \operatorname{act}(b)(\operatorname{gnvc\_}, u)), (\lambda u \in k.\; \operatorname{vnrm}(b)(u)), (\lambda {[}u, \operatorname{gnvv\_}] \in (k \times k).\; \cdot_{b}(u, \operatorname{gnvv\_}))]} \simeq (\lambda {[}u, \operatorname{gnvv\_}] \in (k \times k).\; \cdot_{b}(u, \operatorname{gnvv\_}))) \end{array}}}
\vnbprop{sub-ucsa-tuple9-scal}{\vnbstmt{\begin{array}{@{}l@{}} \forall b, k. \\ \quad (\operatorname{is\text{-}unital\text{-}c\text{-}star\text{-}algebra}(b) \wedge \operatorname{is\text{-}closed\text{-}star\text{-}subalgebra\text{-}set}(b, k)) \Rightarrow  \\ \quad \quad (\operatorname{scal}({[}\operatorname{normed\text{-}field\text{-}as\text{-}commutative\text{-}ring}(\operatorname{cc\text{-}normed\text{-}field}), k, (\lambda {[}u, \operatorname{gnvv\_}] \in (k \times k).\; \operatorname{vadd}(b)(u, \operatorname{gnvv\_})), \operatorname{vzero}(b), (\lambda u \in k.\; \operatorname{vneg}(b)(u)), (\lambda {[}\operatorname{gnvc\_}, u] \in (\mathbb{C} \times k).\; \operatorname{act}(b)(\operatorname{gnvc\_}, u)), (\lambda u \in k.\; \operatorname{vnrm}(b)(u)), (\lambda {[}u, \operatorname{gnvv\_}] \in (k \times k).\; \cdot_{b}(u, \operatorname{gnvv\_})), 1_{b}]) \simeq \operatorname{normed\text{-}field\text{-}as\text{-}commutative\text{-}ring}(\operatorname{cc\text{-}normed\text{-}field})) \end{array}}}
\vnbprop{sub-ucsa-tuple9-vec}{\vnbstmt{\begin{array}{@{}l@{}} \forall b, k. \\ \quad (\operatorname{is\text{-}unital\text{-}c\text{-}star\text{-}algebra}(b) \wedge \operatorname{is\text{-}closed\text{-}star\text{-}subalgebra\text{-}set}(b, k)) \Rightarrow  \\ \quad \quad (\operatorname{vec}({[}\operatorname{normed\text{-}field\text{-}as\text{-}commutative\text{-}ring}(\operatorname{cc\text{-}normed\text{-}field}), k, (\lambda {[}u, \operatorname{gnvv\_}] \in (k \times k).\; \operatorname{vadd}(b)(u, \operatorname{gnvv\_})), \operatorname{vzero}(b), (\lambda u \in k.\; \operatorname{vneg}(b)(u)), (\lambda {[}\operatorname{gnvc\_}, u] \in (\mathbb{C} \times k).\; \operatorname{act}(b)(\operatorname{gnvc\_}, u)), (\lambda u \in k.\; \operatorname{vnrm}(b)(u)), (\lambda {[}u, \operatorname{gnvv\_}] \in (k \times k).\; \cdot_{b}(u, \operatorname{gnvv\_})), 1_{b}]) \simeq k) \end{array}}}
\vnbprop{sub-ucsa-tuple9-vadd}{\vnbstmt{\begin{array}{@{}l@{}} \forall b, k. \\ \quad (\operatorname{is\text{-}unital\text{-}c\text{-}star\text{-}algebra}(b) \wedge \operatorname{is\text{-}closed\text{-}star\text{-}subalgebra\text{-}set}(b, k)) \Rightarrow  \\ \quad \quad (\operatorname{vadd}({[}\operatorname{normed\text{-}field\text{-}as\text{-}commutative\text{-}ring}(\operatorname{cc\text{-}normed\text{-}field}), k, (\lambda {[}u, \operatorname{gnvv\_}] \in (k \times k).\; \operatorname{vadd}(b)(u, \operatorname{gnvv\_})), \operatorname{vzero}(b), (\lambda u \in k.\; \operatorname{vneg}(b)(u)), (\lambda {[}\operatorname{gnvc\_}, u] \in (\mathbb{C} \times k).\; \operatorname{act}(b)(\operatorname{gnvc\_}, u)), (\lambda u \in k.\; \operatorname{vnrm}(b)(u)), (\lambda {[}u, \operatorname{gnvv\_}] \in (k \times k).\; \cdot_{b}(u, \operatorname{gnvv\_})), 1_{b}]) \simeq (\lambda {[}u, \operatorname{gnvv\_}] \in (k \times k).\; \operatorname{vadd}(b)(u, \operatorname{gnvv\_}))) \end{array}}}
\vnbprop{sub-ucsa-tuple9-vzero}{\vnbstmt{\begin{array}{@{}l@{}} \forall b, k. \\ \quad (\operatorname{is\text{-}unital\text{-}c\text{-}star\text{-}algebra}(b) \wedge \operatorname{is\text{-}closed\text{-}star\text{-}subalgebra\text{-}set}(b, k)) \Rightarrow  \\ \quad \quad (\operatorname{vzero}({[}\operatorname{normed\text{-}field\text{-}as\text{-}commutative\text{-}ring}(\operatorname{cc\text{-}normed\text{-}field}), k, (\lambda {[}u, \operatorname{gnvv\_}] \in (k \times k).\; \operatorname{vadd}(b)(u, \operatorname{gnvv\_})), \operatorname{vzero}(b), (\lambda u \in k.\; \operatorname{vneg}(b)(u)), (\lambda {[}\operatorname{gnvc\_}, u] \in (\mathbb{C} \times k).\; \operatorname{act}(b)(\operatorname{gnvc\_}, u)), (\lambda u \in k.\; \operatorname{vnrm}(b)(u)), (\lambda {[}u, \operatorname{gnvv\_}] \in (k \times k).\; \cdot_{b}(u, \operatorname{gnvv\_})), 1_{b}]) \simeq \operatorname{vzero}(b)) \end{array}}}
\vnbprop{sub-ucsa-tuple9-vneg}{\vnbstmt{\begin{array}{@{}l@{}} \forall b, k. \\ \quad (\operatorname{is\text{-}unital\text{-}c\text{-}star\text{-}algebra}(b) \wedge \operatorname{is\text{-}closed\text{-}star\text{-}subalgebra\text{-}set}(b, k)) \Rightarrow  \\ \quad \quad (\operatorname{vneg}({[}\operatorname{normed\text{-}field\text{-}as\text{-}commutative\text{-}ring}(\operatorname{cc\text{-}normed\text{-}field}), k, (\lambda {[}u, \operatorname{gnvv\_}] \in (k \times k).\; \operatorname{vadd}(b)(u, \operatorname{gnvv\_})), \operatorname{vzero}(b), (\lambda u \in k.\; \operatorname{vneg}(b)(u)), (\lambda {[}\operatorname{gnvc\_}, u] \in (\mathbb{C} \times k).\; \operatorname{act}(b)(\operatorname{gnvc\_}, u)), (\lambda u \in k.\; \operatorname{vnrm}(b)(u)), (\lambda {[}u, \operatorname{gnvv\_}] \in (k \times k).\; \cdot_{b}(u, \operatorname{gnvv\_})), 1_{b}]) \simeq (\lambda u \in k.\; \operatorname{vneg}(b)(u))) \end{array}}}
\vnbprop{sub-ucsa-tuple9-act}{\vnbstmt{\begin{array}{@{}l@{}} \forall b, k. \\ \quad (\operatorname{is\text{-}unital\text{-}c\text{-}star\text{-}algebra}(b) \wedge \operatorname{is\text{-}closed\text{-}star\text{-}subalgebra\text{-}set}(b, k)) \Rightarrow  \\ \quad \quad (\operatorname{act}({[}\operatorname{normed\text{-}field\text{-}as\text{-}commutative\text{-}ring}(\operatorname{cc\text{-}normed\text{-}field}), k, (\lambda {[}u, \operatorname{gnvv\_}] \in (k \times k).\; \operatorname{vadd}(b)(u, \operatorname{gnvv\_})), \operatorname{vzero}(b), (\lambda u \in k.\; \operatorname{vneg}(b)(u)), (\lambda {[}\operatorname{gnvc\_}, u] \in (\mathbb{C} \times k).\; \operatorname{act}(b)(\operatorname{gnvc\_}, u)), (\lambda u \in k.\; \operatorname{vnrm}(b)(u)), (\lambda {[}u, \operatorname{gnvv\_}] \in (k \times k).\; \cdot_{b}(u, \operatorname{gnvv\_})), 1_{b}]) \simeq (\lambda {[}\operatorname{gnvc\_}, u] \in (\mathbb{C} \times k).\; \operatorname{act}(b)(\operatorname{gnvc\_}, u))) \end{array}}}
\vnbprop{sub-ucsa-tuple9-vnrm}{\vnbstmt{\begin{array}{@{}l@{}} \forall b, k. \\ \quad (\operatorname{is\text{-}unital\text{-}c\text{-}star\text{-}algebra}(b) \wedge \operatorname{is\text{-}closed\text{-}star\text{-}subalgebra\text{-}set}(b, k)) \Rightarrow  \\ \quad \quad (\operatorname{vnrm}({[}\operatorname{normed\text{-}field\text{-}as\text{-}commutative\text{-}ring}(\operatorname{cc\text{-}normed\text{-}field}), k, (\lambda {[}u, \operatorname{gnvv\_}] \in (k \times k).\; \operatorname{vadd}(b)(u, \operatorname{gnvv\_})), \operatorname{vzero}(b), (\lambda u \in k.\; \operatorname{vneg}(b)(u)), (\lambda {[}\operatorname{gnvc\_}, u] \in (\mathbb{C} \times k).\; \operatorname{act}(b)(\operatorname{gnvc\_}, u)), (\lambda u \in k.\; \operatorname{vnrm}(b)(u)), (\lambda {[}u, \operatorname{gnvv\_}] \in (k \times k).\; \cdot_{b}(u, \operatorname{gnvv\_})), 1_{b}]) \simeq (\lambda u \in k.\; \operatorname{vnrm}(b)(u))) \end{array}}}
\vnbprop{sub-ucsa-tuple9-vmul}{\vnbstmt{\begin{array}{@{}l@{}} \forall b, k. \\ \quad (\operatorname{is\text{-}unital\text{-}c\text{-}star\text{-}algebra}(b) \wedge \operatorname{is\text{-}closed\text{-}star\text{-}subalgebra\text{-}set}(b, k)) \Rightarrow  \\ \quad \quad (\cdot_{{[}\operatorname{normed\text{-}field\text{-}as\text{-}commutative\text{-}ring}(\operatorname{cc\text{-}normed\text{-}field}), k, (\lambda {[}u, \operatorname{gnvv\_}] \in (k \times k).\; \operatorname{vadd}(b)(u, \operatorname{gnvv\_})), \operatorname{vzero}(b), (\lambda u \in k.\; \operatorname{vneg}(b)(u)), (\lambda {[}\operatorname{gnvc\_}, u] \in (\mathbb{C} \times k).\; \operatorname{act}(b)(\operatorname{gnvc\_}, u)), (\lambda u \in k.\; \operatorname{vnrm}(b)(u)), (\lambda {[}u, \operatorname{gnvv\_}] \in (k \times k).\; \cdot_{b}(u, \operatorname{gnvv\_})), 1_{b}]} \simeq (\lambda {[}u, \operatorname{gnvv\_}] \in (k \times k).\; \cdot_{b}(u, \operatorname{gnvv\_}))) \end{array}}}
\vnbprop{sub-ucsa-tuple9-vone}{\vnbstmt{\begin{array}{@{}l@{}} \forall b, k. \\ \quad (\operatorname{is\text{-}unital\text{-}c\text{-}star\text{-}algebra}(b) \wedge \operatorname{is\text{-}closed\text{-}star\text{-}subalgebra\text{-}set}(b, k)) \Rightarrow  \\ \quad \quad (1_{{[}\operatorname{normed\text{-}field\text{-}as\text{-}commutative\text{-}ring}(\operatorname{cc\text{-}normed\text{-}field}), k, (\lambda {[}u, \operatorname{gnvv\_}] \in (k \times k).\; \operatorname{vadd}(b)(u, \operatorname{gnvv\_})), \operatorname{vzero}(b), (\lambda u \in k.\; \operatorname{vneg}(b)(u)), (\lambda {[}\operatorname{gnvc\_}, u] \in (\mathbb{C} \times k).\; \operatorname{act}(b)(\operatorname{gnvc\_}, u)), (\lambda u \in k.\; \operatorname{vnrm}(b)(u)), (\lambda {[}u, \operatorname{gnvv\_}] \in (k \times k).\; \cdot_{b}(u, \operatorname{gnvv\_})), 1_{b}]} \simeq 1_{b}) \end{array}}}
\vnbprop{sub-ucsa-tuple10-scal}{\vnbstmt{\begin{array}{@{}l@{}} \forall b, k. \\ \quad (\operatorname{is\text{-}unital\text{-}c\text{-}star\text{-}algebra}(b) \wedge \operatorname{is\text{-}closed\text{-}star\text{-}subalgebra\text{-}set}(b, k)) \Rightarrow  \\ \quad \quad (\operatorname{scal}({[}\operatorname{normed\text{-}field\text{-}as\text{-}commutative\text{-}ring}(\operatorname{cc\text{-}normed\text{-}field}), k, (\lambda {[}u, \operatorname{gnvv\_}] \in (k \times k).\; \operatorname{vadd}(b)(u, \operatorname{gnvv\_})), \operatorname{vzero}(b), (\lambda u \in k.\; \operatorname{vneg}(b)(u)), (\lambda {[}\operatorname{gnvc\_}, u] \in (\mathbb{C} \times k).\; \operatorname{act}(b)(\operatorname{gnvc\_}, u)), (\lambda u \in k.\; \operatorname{vnrm}(b)(u)), (\lambda {[}u, \operatorname{gnvv\_}] \in (k \times k).\; \cdot_{b}(u, \operatorname{gnvv\_})), 1_{b}, (\lambda u \in k.\; {}^{*_{b}}(u))]) \simeq \operatorname{normed\text{-}field\text{-}as\text{-}commutative\text{-}ring}(\operatorname{cc\text{-}normed\text{-}field})) \end{array}}}
\vnbprop{sub-ucsa-tuple10-vec}{\vnbstmt{\begin{array}{@{}l@{}} \forall b, k. \\ \quad (\operatorname{is\text{-}unital\text{-}c\text{-}star\text{-}algebra}(b) \wedge \operatorname{is\text{-}closed\text{-}star\text{-}subalgebra\text{-}set}(b, k)) \Rightarrow  \\ \quad \quad (\operatorname{vec}({[}\operatorname{normed\text{-}field\text{-}as\text{-}commutative\text{-}ring}(\operatorname{cc\text{-}normed\text{-}field}), k, (\lambda {[}u, \operatorname{gnvv\_}] \in (k \times k).\; \operatorname{vadd}(b)(u, \operatorname{gnvv\_})), \operatorname{vzero}(b), (\lambda u \in k.\; \operatorname{vneg}(b)(u)), (\lambda {[}\operatorname{gnvc\_}, u] \in (\mathbb{C} \times k).\; \operatorname{act}(b)(\operatorname{gnvc\_}, u)), (\lambda u \in k.\; \operatorname{vnrm}(b)(u)), (\lambda {[}u, \operatorname{gnvv\_}] \in (k \times k).\; \cdot_{b}(u, \operatorname{gnvv\_})), 1_{b}, (\lambda u \in k.\; {}^{*_{b}}(u))]) \simeq k) \end{array}}}
\vnbprop{sub-ucsa-tuple10-vadd}{\vnbstmt{\begin{array}{@{}l@{}} \forall b, k. \\ \quad (\operatorname{is\text{-}unital\text{-}c\text{-}star\text{-}algebra}(b) \wedge \operatorname{is\text{-}closed\text{-}star\text{-}subalgebra\text{-}set}(b, k)) \Rightarrow  \\ \quad \quad (\operatorname{vadd}({[}\operatorname{normed\text{-}field\text{-}as\text{-}commutative\text{-}ring}(\operatorname{cc\text{-}normed\text{-}field}), k, (\lambda {[}u, \operatorname{gnvv\_}] \in (k \times k).\; \operatorname{vadd}(b)(u, \operatorname{gnvv\_})), \operatorname{vzero}(b), (\lambda u \in k.\; \operatorname{vneg}(b)(u)), (\lambda {[}\operatorname{gnvc\_}, u] \in (\mathbb{C} \times k).\; \operatorname{act}(b)(\operatorname{gnvc\_}, u)), (\lambda u \in k.\; \operatorname{vnrm}(b)(u)), (\lambda {[}u, \operatorname{gnvv\_}] \in (k \times k).\; \cdot_{b}(u, \operatorname{gnvv\_})), 1_{b}, (\lambda u \in k.\; {}^{*_{b}}(u))]) \simeq (\lambda {[}u, \operatorname{gnvv\_}] \in (k \times k).\; \operatorname{vadd}(b)(u, \operatorname{gnvv\_}))) \end{array}}}
\vnbprop{sub-ucsa-tuple10-vzero}{\vnbstmt{\begin{array}{@{}l@{}} \forall b, k. \\ \quad (\operatorname{is\text{-}unital\text{-}c\text{-}star\text{-}algebra}(b) \wedge \operatorname{is\text{-}closed\text{-}star\text{-}subalgebra\text{-}set}(b, k)) \Rightarrow  \\ \quad \quad (\operatorname{vzero}({[}\operatorname{normed\text{-}field\text{-}as\text{-}commutative\text{-}ring}(\operatorname{cc\text{-}normed\text{-}field}), k, (\lambda {[}u, \operatorname{gnvv\_}] \in (k \times k).\; \operatorname{vadd}(b)(u, \operatorname{gnvv\_})), \operatorname{vzero}(b), (\lambda u \in k.\; \operatorname{vneg}(b)(u)), (\lambda {[}\operatorname{gnvc\_}, u] \in (\mathbb{C} \times k).\; \operatorname{act}(b)(\operatorname{gnvc\_}, u)), (\lambda u \in k.\; \operatorname{vnrm}(b)(u)), (\lambda {[}u, \operatorname{gnvv\_}] \in (k \times k).\; \cdot_{b}(u, \operatorname{gnvv\_})), 1_{b}, (\lambda u \in k.\; {}^{*_{b}}(u))]) \simeq \operatorname{vzero}(b)) \end{array}}}
\vnbprop{sub-ucsa-tuple10-vneg}{\vnbstmt{\begin{array}{@{}l@{}} \forall b, k. \\ \quad (\operatorname{is\text{-}unital\text{-}c\text{-}star\text{-}algebra}(b) \wedge \operatorname{is\text{-}closed\text{-}star\text{-}subalgebra\text{-}set}(b, k)) \Rightarrow  \\ \quad \quad (\operatorname{vneg}({[}\operatorname{normed\text{-}field\text{-}as\text{-}commutative\text{-}ring}(\operatorname{cc\text{-}normed\text{-}field}), k, (\lambda {[}u, \operatorname{gnvv\_}] \in (k \times k).\; \operatorname{vadd}(b)(u, \operatorname{gnvv\_})), \operatorname{vzero}(b), (\lambda u \in k.\; \operatorname{vneg}(b)(u)), (\lambda {[}\operatorname{gnvc\_}, u] \in (\mathbb{C} \times k).\; \operatorname{act}(b)(\operatorname{gnvc\_}, u)), (\lambda u \in k.\; \operatorname{vnrm}(b)(u)), (\lambda {[}u, \operatorname{gnvv\_}] \in (k \times k).\; \cdot_{b}(u, \operatorname{gnvv\_})), 1_{b}, (\lambda u \in k.\; {}^{*_{b}}(u))]) \simeq (\lambda u \in k.\; \operatorname{vneg}(b)(u))) \end{array}}}
\vnbprop{sub-ucsa-tuple10-act}{\vnbstmt{\begin{array}{@{}l@{}} \forall b, k. \\ \quad (\operatorname{is\text{-}unital\text{-}c\text{-}star\text{-}algebra}(b) \wedge \operatorname{is\text{-}closed\text{-}star\text{-}subalgebra\text{-}set}(b, k)) \Rightarrow  \\ \quad \quad (\operatorname{act}({[}\operatorname{normed\text{-}field\text{-}as\text{-}commutative\text{-}ring}(\operatorname{cc\text{-}normed\text{-}field}), k, (\lambda {[}u, \operatorname{gnvv\_}] \in (k \times k).\; \operatorname{vadd}(b)(u, \operatorname{gnvv\_})), \operatorname{vzero}(b), (\lambda u \in k.\; \operatorname{vneg}(b)(u)), (\lambda {[}\operatorname{gnvc\_}, u] \in (\mathbb{C} \times k).\; \operatorname{act}(b)(\operatorname{gnvc\_}, u)), (\lambda u \in k.\; \operatorname{vnrm}(b)(u)), (\lambda {[}u, \operatorname{gnvv\_}] \in (k \times k).\; \cdot_{b}(u, \operatorname{gnvv\_})), 1_{b}, (\lambda u \in k.\; {}^{*_{b}}(u))]) \simeq (\lambda {[}\operatorname{gnvc\_}, u] \in (\mathbb{C} \times k).\; \operatorname{act}(b)(\operatorname{gnvc\_}, u))) \end{array}}}
\vnbprop{sub-ucsa-tuple10-vnrm}{\vnbstmt{\begin{array}{@{}l@{}} \forall b, k. \\ \quad (\operatorname{is\text{-}unital\text{-}c\text{-}star\text{-}algebra}(b) \wedge \operatorname{is\text{-}closed\text{-}star\text{-}subalgebra\text{-}set}(b, k)) \Rightarrow  \\ \quad \quad (\operatorname{vnrm}({[}\operatorname{normed\text{-}field\text{-}as\text{-}commutative\text{-}ring}(\operatorname{cc\text{-}normed\text{-}field}), k, (\lambda {[}u, \operatorname{gnvv\_}] \in (k \times k).\; \operatorname{vadd}(b)(u, \operatorname{gnvv\_})), \operatorname{vzero}(b), (\lambda u \in k.\; \operatorname{vneg}(b)(u)), (\lambda {[}\operatorname{gnvc\_}, u] \in (\mathbb{C} \times k).\; \operatorname{act}(b)(\operatorname{gnvc\_}, u)), (\lambda u \in k.\; \operatorname{vnrm}(b)(u)), (\lambda {[}u, \operatorname{gnvv\_}] \in (k \times k).\; \cdot_{b}(u, \operatorname{gnvv\_})), 1_{b}, (\lambda u \in k.\; {}^{*_{b}}(u))]) \simeq (\lambda u \in k.\; \operatorname{vnrm}(b)(u))) \end{array}}}
\vnbprop{sub-ucsa-tuple10-vmul}{\vnbstmt{\begin{array}{@{}l@{}} \forall b, k. \\ \quad (\operatorname{is\text{-}unital\text{-}c\text{-}star\text{-}algebra}(b) \wedge \operatorname{is\text{-}closed\text{-}star\text{-}subalgebra\text{-}set}(b, k)) \Rightarrow  \\ \quad \quad (\cdot_{{[}\operatorname{normed\text{-}field\text{-}as\text{-}commutative\text{-}ring}(\operatorname{cc\text{-}normed\text{-}field}), k, (\lambda {[}u, \operatorname{gnvv\_}] \in (k \times k).\; \operatorname{vadd}(b)(u, \operatorname{gnvv\_})), \operatorname{vzero}(b), (\lambda u \in k.\; \operatorname{vneg}(b)(u)), (\lambda {[}\operatorname{gnvc\_}, u] \in (\mathbb{C} \times k).\; \operatorname{act}(b)(\operatorname{gnvc\_}, u)), (\lambda u \in k.\; \operatorname{vnrm}(b)(u)), (\lambda {[}u, \operatorname{gnvv\_}] \in (k \times k).\; \cdot_{b}(u, \operatorname{gnvv\_})), 1_{b}, (\lambda u \in k.\; {}^{*_{b}}(u))]} \simeq (\lambda {[}u, \operatorname{gnvv\_}] \in (k \times k).\; \cdot_{b}(u, \operatorname{gnvv\_}))) \end{array}}}
\vnbprop{sub-ucsa-tuple10-vone}{\vnbstmt{\begin{array}{@{}l@{}} \forall b, k. \\ \quad (\operatorname{is\text{-}unital\text{-}c\text{-}star\text{-}algebra}(b) \wedge \operatorname{is\text{-}closed\text{-}star\text{-}subalgebra\text{-}set}(b, k)) \Rightarrow  \\ \quad \quad (1_{{[}\operatorname{normed\text{-}field\text{-}as\text{-}commutative\text{-}ring}(\operatorname{cc\text{-}normed\text{-}field}), k, (\lambda {[}u, \operatorname{gnvv\_}] \in (k \times k).\; \operatorname{vadd}(b)(u, \operatorname{gnvv\_})), \operatorname{vzero}(b), (\lambda u \in k.\; \operatorname{vneg}(b)(u)), (\lambda {[}\operatorname{gnvc\_}, u] \in (\mathbb{C} \times k).\; \operatorname{act}(b)(\operatorname{gnvc\_}, u)), (\lambda u \in k.\; \operatorname{vnrm}(b)(u)), (\lambda {[}u, \operatorname{gnvv\_}] \in (k \times k).\; \cdot_{b}(u, \operatorname{gnvv\_})), 1_{b}, (\lambda u \in k.\; {}^{*_{b}}(u))]} \simeq 1_{b}) \end{array}}}
\vnbprop{sub-ucsa-tuple10-vstar}{\vnbstmt{\begin{array}{@{}l@{}} \forall b, k. \\ \quad (\operatorname{is\text{-}unital\text{-}c\text{-}star\text{-}algebra}(b) \wedge \operatorname{is\text{-}closed\text{-}star\text{-}subalgebra\text{-}set}(b, k)) \Rightarrow  \\ \quad \quad ({}^{*_{{[}\operatorname{normed\text{-}field\text{-}as\text{-}commutative\text{-}ring}(\operatorname{cc\text{-}normed\text{-}field}), k, (\lambda {[}u, \operatorname{gnvv\_}] \in (k \times k).\; \operatorname{vadd}(b)(u, \operatorname{gnvv\_})), \operatorname{vzero}(b), (\lambda u \in k.\; \operatorname{vneg}(b)(u)), (\lambda {[}\operatorname{gnvc\_}, u] \in (\mathbb{C} \times k).\; \operatorname{act}(b)(\operatorname{gnvc\_}, u)), (\lambda u \in k.\; \operatorname{vnrm}(b)(u)), (\lambda {[}u, \operatorname{gnvv\_}] \in (k \times k).\; \cdot_{b}(u, \operatorname{gnvv\_})), 1_{b}, (\lambda u \in k.\; {}^{*_{b}}(u))]}} \simeq (\lambda u \in k.\; {}^{*_{b}}(u))) \end{array}}}
\vnbprop{sub-ucsa-tuple-cna}{\vnbstmt{\begin{array}{@{}l@{}} \forall b, k. \\ \quad (\operatorname{is\text{-}unital\text{-}c\text{-}star\text{-}algebra}(b) \wedge \operatorname{is\text{-}closed\text{-}star\text{-}subalgebra\text{-}set}(b, k)) \Rightarrow  \\ \quad \quad \operatorname{is\text{-}complex\text{-}normed\text{-}algebra}({[}\operatorname{normed\text{-}field\text{-}as\text{-}commutative\text{-}ring}(\operatorname{cc\text{-}normed\text{-}field}), k, (\lambda {[}u, \operatorname{gnvv\_}] \in (k \times k).\; \operatorname{vadd}(b)(u, \operatorname{gnvv\_})), \operatorname{vzero}(b), (\lambda u \in k.\; \operatorname{vneg}(b)(u)), (\lambda {[}\operatorname{gnvc\_}, u] \in (\mathbb{C} \times k).\; \operatorname{act}(b)(\operatorname{gnvc\_}, u)), (\lambda u \in k.\; \operatorname{vnrm}(b)(u)), (\lambda {[}u, \operatorname{gnvv\_}] \in (k \times k).\; \cdot_{b}(u, \operatorname{gnvv\_}))]) \end{array}}}
\vnbprop{sub-ucsa-tuple-complete}{\vnbstmt{\begin{array}{@{}l@{}} \forall b, k. \\ \quad (\operatorname{is\text{-}unital\text{-}c\text{-}star\text{-}algebra}(b) \wedge \operatorname{is\text{-}closed\text{-}star\text{-}subalgebra\text{-}set}(b, k)) \Rightarrow  \\ \quad \quad \operatorname{is\text{-}complete}(\operatorname{nag\text{-}metric\text{-}space}(\operatorname{cna\text{-}as\text{-}normed\text{-}ag}({[}\operatorname{normed\text{-}field\text{-}as\text{-}commutative\text{-}ring}(\operatorname{cc\text{-}normed\text{-}field}), k, (\lambda {[}u, \operatorname{gnvv\_}] \in (k \times k).\; \operatorname{vadd}(b)(u, \operatorname{gnvv\_})), \operatorname{vzero}(b), (\lambda u \in k.\; \operatorname{vneg}(b)(u)), (\lambda {[}\operatorname{gnvc\_}, u] \in (\mathbb{C} \times k).\; \operatorname{act}(b)(\operatorname{gnvc\_}, u)), (\lambda u \in k.\; \operatorname{vnrm}(b)(u)), (\lambda {[}u, \operatorname{gnvv\_}] \in (k \times k).\; \cdot_{b}(u, \operatorname{gnvv\_}))]))) \end{array}}}
\vnbprop{sub-ucsa-tuple-cba}{\vnbstmt{\begin{array}{@{}l@{}} \forall b, k. \\ \quad (\operatorname{is\text{-}unital\text{-}c\text{-}star\text{-}algebra}(b) \wedge \operatorname{is\text{-}closed\text{-}star\text{-}subalgebra\text{-}set}(b, k)) \Rightarrow  \\ \quad \quad \operatorname{is\text{-}complex\text{-}banach\text{-}algebra}({[}\operatorname{normed\text{-}field\text{-}as\text{-}commutative\text{-}ring}(\operatorname{cc\text{-}normed\text{-}field}), k, (\lambda {[}u, \operatorname{gnvv\_}] \in (k \times k).\; \operatorname{vadd}(b)(u, \operatorname{gnvv\_})), \operatorname{vzero}(b), (\lambda u \in k.\; \operatorname{vneg}(b)(u)), (\lambda {[}\operatorname{gnvc\_}, u] \in (\mathbb{C} \times k).\; \operatorname{act}(b)(\operatorname{gnvc\_}, u)), (\lambda u \in k.\; \operatorname{vnrm}(b)(u)), (\lambda {[}u, \operatorname{gnvv\_}] \in (k \times k).\; \cdot_{b}(u, \operatorname{gnvv\_}))]) \end{array}}}
\vnbprop{sub-ucsa-tuple-uba}{\vnbstmt{\begin{array}{@{}l@{}} \forall b, k. \\ \quad (\operatorname{is\text{-}unital\text{-}c\text{-}star\text{-}algebra}(b) \wedge \operatorname{is\text{-}closed\text{-}star\text{-}subalgebra\text{-}set}(b, k)) \Rightarrow  \\ \quad \quad \operatorname{is\text{-}unital\text{-}banach\text{-}algebra}({[}\operatorname{normed\text{-}field\text{-}as\text{-}commutative\text{-}ring}(\operatorname{cc\text{-}normed\text{-}field}), k, (\lambda {[}u, \operatorname{gnvv\_}] \in (k \times k).\; \operatorname{vadd}(b)(u, \operatorname{gnvv\_})), \operatorname{vzero}(b), (\lambda u \in k.\; \operatorname{vneg}(b)(u)), (\lambda {[}\operatorname{gnvc\_}, u] \in (\mathbb{C} \times k).\; \operatorname{act}(b)(\operatorname{gnvc\_}, u)), (\lambda u \in k.\; \operatorname{vnrm}(b)(u)), (\lambda {[}u, \operatorname{gnvv\_}] \in (k \times k).\; \cdot_{b}(u, \operatorname{gnvv\_})), 1_{b}]) \end{array}}}
\vnbprop{sub-ucsa-tuple-uiba}{\vnbstmt{\begin{array}{@{}l@{}} \forall b, k. \\ \quad (\operatorname{is\text{-}unital\text{-}c\text{-}star\text{-}algebra}(b) \wedge \operatorname{is\text{-}closed\text{-}star\text{-}subalgebra\text{-}set}(b, k)) \Rightarrow  \\ \quad \quad \operatorname{is\text{-}unital\text{-}involutive\text{-}banach\text{-}algebra}({[}\operatorname{normed\text{-}field\text{-}as\text{-}commutative\text{-}ring}(\operatorname{cc\text{-}normed\text{-}field}), k, (\lambda {[}u, \operatorname{gnvv\_}] \in (k \times k).\; \operatorname{vadd}(b)(u, \operatorname{gnvv\_})), \operatorname{vzero}(b), (\lambda u \in k.\; \operatorname{vneg}(b)(u)), (\lambda {[}\operatorname{gnvc\_}, u] \in (\mathbb{C} \times k).\; \operatorname{act}(b)(\operatorname{gnvc\_}, u)), (\lambda u \in k.\; \operatorname{vnrm}(b)(u)), (\lambda {[}u, \operatorname{gnvv\_}] \in (k \times k).\; \cdot_{b}(u, \operatorname{gnvv\_})), 1_{b}, (\lambda u \in k.\; {}^{*_{b}}(u))]) \end{array}}}
\vnbprop{sub-ucsa-tuple-ucsa}{\vnbstmt{\begin{array}{@{}l@{}} \forall b, k. \\ \quad (\operatorname{is\text{-}unital\text{-}c\text{-}star\text{-}algebra}(b) \wedge \operatorname{is\text{-}closed\text{-}star\text{-}subalgebra\text{-}set}(b, k)) \Rightarrow  \\ \quad \quad \operatorname{is\text{-}unital\text{-}c\text{-}star\text{-}algebra}({[}\operatorname{normed\text{-}field\text{-}as\text{-}commutative\text{-}ring}(\operatorname{cc\text{-}normed\text{-}field}), k, (\lambda {[}u, \operatorname{gnvv\_}] \in (k \times k).\; \operatorname{vadd}(b)(u, \operatorname{gnvv\_})), \operatorname{vzero}(b), (\lambda u \in k.\; \operatorname{vneg}(b)(u)), (\lambda {[}\operatorname{gnvc\_}, u] \in (\mathbb{C} \times k).\; \operatorname{act}(b)(\operatorname{gnvc\_}, u)), (\lambda u \in k.\; \operatorname{vnrm}(b)(u)), (\lambda {[}u, \operatorname{gnvv\_}] \in (k \times k).\; \cdot_{b}(u, \operatorname{gnvv\_})), 1_{b}, (\lambda u \in k.\; {}^{*_{b}}(u))]) \end{array}}}
\vnbprop{sub-ucsa-tuple-ccsa}{\vnbstmt{\begin{array}{@{}l@{}} \forall b, k. \\ \quad \operatorname{is\text{-}unital\text{-}c\text{-}star\text{-}algebra}(b) \wedge \\ \quad \operatorname{is\text{-}closed\text{-}star\text{-}subalgebra\text{-}set}(b, k) \wedge \\ \quad \forall m, n \in k.\; (\cdot_{b}(m, n) = \cdot_{b}(n, m)) \Rightarrow \\ \quad \quad \operatorname{is\text{-}commutative\text{-}c\text{-}star\text{-}algebra}({[}\operatorname{normed\text{-}field\text{-}as\text{-}commutative\text{-}ring}(\operatorname{cc\text{-}normed\text{-}field}), k, (\lambda {[}u, \operatorname{gnvv\_}] \in (k \times k).\; \operatorname{vadd}(b)(u, \operatorname{gnvv\_})), \operatorname{vzero}(b), (\lambda u \in k.\; \operatorname{vneg}(b)(u)), (\lambda {[}\operatorname{gnvc\_}, u] \in (\mathbb{C} \times k).\; \operatorname{act}(b)(\operatorname{gnvc\_}, u)), (\lambda u \in k.\; \operatorname{vnrm}(b)(u)), (\lambda {[}u, \operatorname{gnvv\_}] \in (k \times k).\; \cdot_{b}(u, \operatorname{gnvv\_})), 1_{b}, (\lambda u \in k.\; {}^{*_{b}}(u))]) \end{array}}}
\vnbprop{sub-ucsa-vec}{\vnbstmt{\begin{array}{@{}l@{}} \forall b, k. \\ \quad (\operatorname{is\text{-}unital\text{-}c\text{-}star\text{-}algebra}(b) \wedge \operatorname{is\text{-}closed\text{-}star\text{-}subalgebra\text{-}set}(b, k)) \Rightarrow  \\ \quad \quad (\operatorname{vec}(\operatorname{sub\text{-}ucsa}(b, k)) \simeq k) \end{array}}}
\vnbprop{sub-ucsa-unital-c-star}{\vnbstmt{\begin{array}{@{}l@{}} \forall b, k. \\ \quad (\operatorname{is\text{-}unital\text{-}c\text{-}star\text{-}algebra}(b) \wedge \operatorname{is\text{-}closed\text{-}star\text{-}subalgebra\text{-}set}(b, k)) \Rightarrow  \\ \quad \quad \operatorname{is\text{-}unital\text{-}c\text{-}star\text{-}algebra}(\operatorname{sub\text{-}ucsa}(b, k)) \end{array}}}
\vnbprop{sub-ucsa-commutative-c-star}{\vnbstmt{\begin{array}{@{}l@{}} \forall b, k. \\ \quad \operatorname{is\text{-}unital\text{-}c\text{-}star\text{-}algebra}(b) \wedge \\ \quad \operatorname{is\text{-}closed\text{-}star\text{-}subalgebra\text{-}set}(b, k) \wedge \\ \quad \forall m, n \in k.\; (\cdot_{b}(m, n) = \cdot_{b}(n, m)) \Rightarrow \\ \quad \quad \operatorname{is\text{-}commutative\text{-}c\text{-}star\text{-}algebra}(\operatorname{sub\text{-}ucsa}(b, k)) \end{array}}}
\vnbprop{sub-ucsa-subalgebra}{\vnbstmt{\begin{array}{@{}l@{}} \forall b, k. \\ \quad (\operatorname{is\text{-}unital\text{-}c\text{-}star\text{-}algebra}(b) \wedge \operatorname{is\text{-}closed\text{-}star\text{-}subalgebra\text{-}set}(b, k)) \Rightarrow  \\ \quad \quad \operatorname{is\text{-}unital\text{-}c\text{-}star\text{-}subalgebra}(\operatorname{sub\text{-}ucsa}(b, k), b) \end{array}}}
\vnbprop{cstar-gen-ucsa}{\vnbstmt{\begin{array}{@{}l@{}} \forall b, x. \\ \quad (\operatorname{is\text{-}unital\text{-}c\text{-}star\text{-}algebra}(b) \wedge (x \in \operatorname{vec}(b))) \Rightarrow  \\ \quad \quad \operatorname{is\text{-}unital\text{-}c\text{-}star\text{-}subalgebra}(\operatorname{sub\text{-}ucsa}(b, \operatorname{cstar\text{-}gen}(b, x)), b) \wedge \\ \quad \quad x \in \operatorname{vec}(\operatorname{sub\text{-}ucsa}(b, \operatorname{cstar\text{-}gen}(b, x))) \end{array}}}
\vnbprop{cstar-gen-commutative-c-star}{\vnbstmt{\begin{array}{@{}l@{}} \forall b, x. \\ \quad (\operatorname{is\text{-}unital\text{-}c\text{-}star\text{-}algebra}(b) \wedge \operatorname{is\text{-}normal\text{-}el}(b, x)) \Rightarrow  \\ \quad \quad \operatorname{is\text{-}commutative\text{-}c\text{-}star\text{-}algebra}(\operatorname{sub\text{-}ucsa}(b, \operatorname{cstar\text{-}gen}(b, x))) \end{array}}}

\vnbbookfile{theorem-library/gn-gen}{theorem-library/gn-gen}
\noindent{\small gn-gen.scm -- the notes' general (non-commutative) forms of Propositions 1.31, 1.36 and 1.39(a) of F. J. Thayer, "Banach Algebras and Applications", ch. 1, by instantiating the relative forms (gn-onto cfc-normal-*, gn-pos sub-order-unit-le, sub-one-minus-sa-square) at a := C*(x, 1) = SUB-UCSA(b, CSTAR-GEN(b, x)), a ...}\par
\vnbprop{cfc-n-in}{\vnbstmt{\begin{array}{@{}l@{}} \forall b, x.; \forall f \in C({\operatorname{spec\text{-}top}(b, x)}). \\ \quad (\operatorname{is\text{-}unital\text{-}c\text{-}star\text{-}algebra}(b) \wedge \operatorname{is\text{-}normal\text{-}el}(b, x)) \Rightarrow  \\ \quad \quad \operatorname{cfc\text{-}n}(b, x, f) \in \operatorname{vec}(b) \end{array}}}
\vnbprop{cfc-n-add}{\vnbstmt{\begin{array}{@{}l@{}} \forall b, x.; \forall f, g \in C({\operatorname{spec\text{-}top}(b, x)}). \\ \quad (\operatorname{is\text{-}unital\text{-}c\text{-}star\text{-}algebra}(b) \wedge \operatorname{is\text{-}normal\text{-}el}(b, x)) \Rightarrow  \\ \quad \quad \operatorname{cfc\text{-}n}(b, x, \operatorname{ccont\text{-}vadd}(\operatorname{spec\text{-}top}(b, x))(f, g)) = \operatorname{vadd}(b)(\operatorname{cfc\text{-}n}(b, x, f), \operatorname{cfc\text{-}n}(b, x, g)) \end{array}}}
\vnbprop{cfc-n-mul}{\vnbstmt{\begin{array}{@{}l@{}} \forall b, x.; \forall f, g \in C({\operatorname{spec\text{-}top}(b, x)}). \\ \quad (\operatorname{is\text{-}unital\text{-}c\text{-}star\text{-}algebra}(b) \wedge \operatorname{is\text{-}normal\text{-}el}(b, x)) \Rightarrow  \\ \quad \quad \operatorname{cfc\text{-}n}(b, x, \operatorname{ccont\text{-}vmul}(\operatorname{spec\text{-}top}(b, x))(f, g)) = \cdot_{b}(\operatorname{cfc\text{-}n}(b, x, f), \operatorname{cfc\text{-}n}(b, x, g)) \end{array}}}
\vnbprop{cfc-n-act}{\vnbstmt{\begin{array}{@{}l@{}} \forall b, x.; \forall f \in C({\operatorname{spec\text{-}top}(b, x)}), z \in \mathbb{C}. \\ \quad (\operatorname{is\text{-}unital\text{-}c\text{-}star\text{-}algebra}(b) \wedge \operatorname{is\text{-}normal\text{-}el}(b, x)) \Rightarrow  \\ \quad \quad \operatorname{cfc\text{-}n}(b, x, \operatorname{ccont\text{-}act}(\operatorname{spec\text{-}top}(b, x))(z, f)) = \operatorname{act}(b)(z, \operatorname{cfc\text{-}n}(b, x, f)) \end{array}}}
\vnbprop{cfc-n-star}{\vnbstmt{\begin{array}{@{}l@{}} \forall b, x.; \forall f \in C({\operatorname{spec\text{-}top}(b, x)}). \\ \quad (\operatorname{is\text{-}unital\text{-}c\text{-}star\text{-}algebra}(b) \wedge \operatorname{is\text{-}normal\text{-}el}(b, x)) \Rightarrow  \\ \quad \quad \operatorname{cfc\text{-}n}(b, x, \operatorname{ccont\text{-}vstar}(\operatorname{spec\text{-}top}(b, x))(f)) = {}^{*_{b}}(\operatorname{cfc\text{-}n}(b, x, f)) \end{array}}}
\vnbprop{cfc-n-one-id}{\vnbstmt{\begin{array}{@{}l@{}} \forall b, x. \\ \quad (\operatorname{is\text{-}unital\text{-}c\text{-}star\text{-}algebra}(b) \wedge \operatorname{is\text{-}normal\text{-}el}(b, x)) \Rightarrow  \\ \quad \quad \operatorname{cfc\text{-}n}(b, x, \operatorname{ccont\text{-}vone}(\operatorname{spec\text{-}top}(b, x))) = 1_{b} \wedge \\ \quad \quad \operatorname{cfc\text{-}n}(b, x, (\lambda z \in \sigma({x}).\; z)) = x \end{array}}}
\vnbprop{self-adjoint-normal}{\vnbstmt{\begin{array}{@{}l@{}} \forall b.; \forall x \in \operatorname{vec}(b). \\ \quad (\operatorname{is\text{-}unital\text{-}involutive\text{-}banach\text{-}algebra}(b) \wedge ({}^{*_{b}}(x) = x)) \Rightarrow  \\ \quad \quad \operatorname{is\text{-}normal\text{-}el}(b, x) \end{array}}}
\vnbprop{order-unit-le}{\vnbstmt{\begin{array}{@{}l@{}} \forall b.; \forall x \in \operatorname{vec}(b), l \in \mathbb{R}. \\ \quad (\operatorname{is\text{-}unital\text{-}c\text{-}star\text{-}algebra}(b) \wedge (({}^{*_{b}}(x) = x) \wedge (\operatorname{vnrm}(b)(x) \leq l))) \Rightarrow  \\ \quad \quad {x} \le_{b} {\operatorname{act}(b)(l, 1_{b})} \end{array}}}
\vnbprop{one-minus-sa-square}{\vnbstmt{\begin{array}{@{}l@{}} \forall b.; \forall x \in \operatorname{vec}(b). \\ \quad (\operatorname{is\text{-}unital\text{-}c\text{-}star\text{-}algebra}(b) \wedge (({}^{*_{b}}(x) = x) \wedge (\operatorname{vnrm}(b)(x) \leq 1))) \Rightarrow  \\ \quad \quad \operatorname{vadd}(b)(1_{b}, \operatorname{vneg}(b)(x)) \in \operatorname{sa\text{-}squares}(b) \end{array}}}
\vnbprop{pos-norm-test-suff}{\vnbstmt{\begin{array}{@{}l@{}} \forall b.; \forall x \in \operatorname{vec}(b). \\ \quad \operatorname{is\text{-}unital\text{-}c\text{-}star\text{-}algebra}(b) \wedge \\ \quad {}^{*_{b}}(x) = x \wedge \\ \quad \operatorname{vnrm}(b)(\operatorname{vadd}(b)(\operatorname{act}(b)(\operatorname{vnrm}(b)(x), 1_{b}), \operatorname{vneg}(b)(x))) \leq \operatorname{vnrm}(b)(x) \Rightarrow \\ \quad \quad x \in {b}_{+} \end{array}}}
\vnbprop{sa-square-norm-test}{\vnbstmt{\begin{array}{@{}l@{}} \forall b.; \forall y \in \operatorname{sa\text{-}squares}(b). \\ \quad \operatorname{is\text{-}unital\text{-}c\text{-}star\text{-}algebra}(b) \Rightarrow  \\ \quad \quad \operatorname{vnrm}(b)(\operatorname{vadd}(b)(\operatorname{act}(b)(\operatorname{vnrm}(b)(y), 1_{b}), \operatorname{vneg}(b)(y))) \leq \operatorname{vnrm}(b)(y) \end{array}}}
\vnbprop{pos-norm-test-sq}{\vnbstmt{\begin{array}{@{}l@{}} \forall b.; \forall x \in \operatorname{vec}(b). \\ \quad \operatorname{is\text{-}unital\text{-}c\text{-}star\text{-}algebra}(b) \wedge \\ \quad {}^{*_{b}}(x) = x \wedge \\ \quad \operatorname{vnrm}(b)(\operatorname{vadd}(b)(\operatorname{act}(b)(\operatorname{vnrm}(b)(x), 1_{b}), \operatorname{vneg}(b)(x))) \leq \operatorname{vnrm}(b)(x) \Rightarrow \\ \quad \quad x \in \operatorname{sa\text{-}squares}(b) \end{array}}}
\vnbprop{one-minus-sq-norm}{\vnbstmt{\begin{array}{@{}l@{}} \forall b.; \forall y \in \operatorname{sa\text{-}squares}(b). \\ \quad (\operatorname{is\text{-}unital\text{-}c\text{-}star\text{-}algebra}(b) \wedge (\operatorname{vnrm}(b)(y) \leq 1)) \Rightarrow  \\ \quad \quad \operatorname{vnrm}(b)(\operatorname{vadd}(b)(1_{b}, \operatorname{vneg}(b)(y))) \leq 1 \end{array}}}
\vnbprop{pos-cone-act-nonneg}{\vnbstmt{\begin{array}{@{}l@{}} \forall a.; \forall t \in \mathbb{R}, x \in {a}_{+}. \\ \quad (\operatorname{is\text{-}unital\text{-}involutive\text{-}banach\text{-}algebra}(a) \wedge (0 \leq t)) \Rightarrow  \\ \quad \quad \operatorname{act}(a)(t, x) \in {a}_{+} \end{array}}}

\vnbbookfile{theorem-library/ba4-pos}{theorem-library/ba4-pos}
\noindent{\small ba4-pos.scm -- POSITIVE FUNCTIONALS: BanachAlgebras.pdf Lemma 4.4, Corollary 4.6. Agent BA-1111-B, 2026-11-11.  Destination: theorem-library/ba4-pos.scm.  Library: operator-algebras. Load window: after structure-library/ba-rep-vocab (this wave), theorem-library/kal-laws (the law set kal-uiba of dk-alg-norm!), ...}\par
\vnbprop{pos-lf-conj-sym}{\vnbstmt{\begin{array}{@{}l@{}} \forall a, f.; \forall x, y \in \operatorname{vec}(a). \\ \quad (\operatorname{is\text{-}unital\text{-}involutive\text{-}banach\text{-}algebra}(a) \wedge \operatorname{is\text{-}positive\text{-}lf}(a, f)) \Rightarrow  \\ \quad \quad f(\cdot_{a}({}^{*_{a}}(x), y)) = \operatorname{conjugate}(f(\cdot_{a}({}^{*_{a}}(y), x))) \end{array}}}
\vnbprop{pos-lf-sesquilinear}{\vnbstmt{\begin{array}{@{}l@{}} \forall a, f. \\ \quad (\operatorname{is\text{-}unital\text{-}involutive\text{-}banach\text{-}algebra}(a) \wedge \operatorname{is\text{-}positive\text{-}lf}(a, f)) \Rightarrow  \\ \quad \quad \forall x, y, z \in \operatorname{vec}(a). \\ \quad \quad \quad f(\cdot_{a}({}^{*_{a}}(z), \operatorname{vadd}(a)(x, y))) \\ \quad \quad \quad =\; f(\cdot_{a}({}^{*_{a}}(z), x)) + f(\cdot_{a}({}^{*_{a}}(z), y)) \wedge \\ \quad \quad \forall c \in \mathbb{C}, x, y \in \operatorname{vec}(a). \\ \quad \quad \quad f(\cdot_{a}({}^{*_{a}}(y), \operatorname{act}(a)(c, x))) = (c \cdot f(\cdot_{a}({}^{*_{a}}(y), x))) \wedge \\ \quad \quad \forall x, y \in \operatorname{vec}(a). \\ \quad \quad \quad f(\cdot_{a}({}^{*_{a}}(x), y)) = \operatorname{conjugate}(f(\cdot_{a}({}^{*_{a}}(y), x))) \wedge \\ \quad \quad \forall x \in \operatorname{vec}(a). \\ \quad \quad \quad (f(\cdot_{a}({}^{*_{a}}(x), x)) \in \mathbb{R}) \wedge (0 \leq f(\cdot_{a}({}^{*_{a}}(x), x))) \end{array}}}
\vnbprop{pos-lf-hermitian}{\vnbstmt{\begin{array}{@{}l@{}} \forall a, f.; \forall x \in \operatorname{vec}(a). \\ \quad (\operatorname{is\text{-}unital\text{-}involutive\text{-}banach\text{-}algebra}(a) \wedge \operatorname{is\text{-}positive\text{-}lf}(a, f)) \Rightarrow  \\ \quad \quad f({}^{*_{a}}(x)) = \operatorname{conjugate}(f(x)) \end{array}}}
\vnbprop{pos-lf-expand-two}{\vnbstmt{\begin{array}{@{}l@{}} \forall a, f.; \forall x, y \in \operatorname{vec}(a), r, s \in \mathbb{C}. \\ \quad (\operatorname{is\text{-}unital\text{-}involutive\text{-}banach\text{-}algebra}(a) \wedge \operatorname{is\text{-}positive\text{-}lf}(a, f)) \Rightarrow  \\ \quad \quad f(\cdot_{a}({}^{*_{a}}(\operatorname{vadd}(a)(\operatorname{act}(a)(r, x), \operatorname{act}(a)(s, y))), \\ \quad \quad \quad \quad \operatorname{vadd}(a)(\operatorname{act}(a)(r, x), \operatorname{act}(a)(s, y)))) \\ \quad \quad =\; ((((\operatorname{conjugate}(r) \cdot r) \cdot f(\cdot_{a}({}^{*_{a}}(x), x))) + ((\operatorname{conjugate}(r) \cdot s) \cdot \operatorname{conjugate}(f(\cdot_{a}({}^{*_{a}}(y), x))))) + ((\operatorname{conjugate}(s) \cdot r) \cdot f(\cdot_{a}({}^{*_{a}}(y), x)))) + ((\operatorname{conjugate}(s) \cdot s) \cdot f(\cdot_{a}({}^{*_{a}}(y), y))) \end{array}}}
\vnbprop{pos-lf-schwarz}{\vnbstmt{\begin{array}{@{}l@{}} \forall a, f.; \forall x, y \in \operatorname{vec}(a). \\ \quad (\operatorname{is\text{-}unital\text{-}involutive\text{-}banach\text{-}algebra}(a) \wedge \operatorname{is\text{-}positive\text{-}lf}(a, f)) \Rightarrow  \\ \quad \quad f(\cdot_{a}({}^{*_{a}}(y), x)) \cdot \operatorname{conjugate}(f(\cdot_{a}({}^{*_{a}}(y), x))) \\ \quad \quad \leq\; f(\cdot_{a}({}^{*_{a}}(y), y)) \cdot f(\cdot_{a}({}^{*_{a}}(x), x)) \end{array}}}
\vnbprop{pos-lf-schwarz-one}{\vnbstmt{\begin{array}{@{}l@{}} \forall a, f.; \forall x \in \operatorname{vec}(a). \\ \quad (\operatorname{is\text{-}unital\text{-}involutive\text{-}banach\text{-}algebra}(a) \wedge \operatorname{is\text{-}positive\text{-}lf}(a, f)) \Rightarrow  \\ \quad \quad (f(x) \cdot \operatorname{conjugate}(f(x))) \leq (f(1_{a}) \cdot f(\cdot_{a}({}^{*_{a}}(x), x))) \end{array}}}
\vnbprop{pos-lf-unit}{\vnbstmt{\begin{array}{@{}l@{}} \forall a, f. \\ \quad (\operatorname{is\text{-}unital\text{-}involutive\text{-}banach\text{-}algebra}(a) \wedge \operatorname{is\text{-}positive\text{-}lf}(a, f)) \Rightarrow  \\ \quad \quad (f(1_{a}) \in \mathbb{R}) \wedge (0 \leq f(1_{a})) \end{array}}}
\vnbprop{pos-lf-ball}{\vnbstmt{\begin{array}{@{}l@{}} \forall a, f.; \forall x \in \operatorname{vec}(a). \\ \quad (\operatorname{is\text{-}unital\text{-}involutive\text{-}banach\text{-}algebra}(a) \wedge (\operatorname{is\text{-}positive\text{-}lf}(a, f) \wedge (\operatorname{vnrm}(a)(x) < 1))) \Rightarrow  \\ \quad \quad f(\cdot_{a}({}^{*_{a}}(x), x)) \leq f(1_{a}) \end{array}}}
\vnbprop{pos-lf-ball-bound}{\vnbstmt{\begin{array}{@{}l@{}} \forall a, f.; \forall x \in \operatorname{vec}(a). \\ \quad (\operatorname{is\text{-}unital\text{-}involutive\text{-}banach\text{-}algebra}(a) \wedge (\operatorname{is\text{-}positive\text{-}lf}(a, f) \wedge (\operatorname{vnrm}(a)(x) < 1))) \Rightarrow  \\ \quad \quad \operatorname{magnitude}(f(x)) \leq f(1_{a}) \end{array}}}
\vnbprop{pos-lf-bounded}{\vnbstmt{\begin{array}{@{}l@{}} \forall a, f.; \forall x \in \operatorname{vec}(a). \\ \quad (\operatorname{is\text{-}unital\text{-}involutive\text{-}banach\text{-}algebra}(a) \wedge \operatorname{is\text{-}positive\text{-}lf}(a, f)) \Rightarrow  \\ \quad \quad \operatorname{magnitude}(f(x)) \leq (f(1_{a}) \cdot \operatorname{vnrm}(a)(x)) \end{array}}}
\vnbprop{pos-lf-norm}{\vnbstmt{\begin{array}{@{}l@{}} \forall a, f. \\ \quad (\operatorname{is\text{-}unital\text{-}involutive\text{-}banach\text{-}algebra}(a) \wedge \operatorname{is\text{-}positive\text{-}lf}(a, f)) \Rightarrow  \\ \quad \quad (f \in \operatorname{cdual\text{-}vec}(a)) \wedge (\operatorname{cdual\text{-}norm}(a, f) = f(1_{a})) \end{array}}}
\vnbprop{pos-lf-sandwich}{\vnbstmt{\begin{array}{@{}l@{}} \forall a, f.; \forall y \in \operatorname{vec}(a). \\ \quad (\operatorname{is\text{-}unital\text{-}involutive\text{-}banach\text{-}algebra}(a) \wedge \operatorname{is\text{-}positive\text{-}lf}(a, f)) \Rightarrow  \\ \quad \quad \operatorname{is\text{-}positive\text{-}lf}(a, (\lambda u \in \operatorname{vec}(a).\; f(\cdot_{a}(\cdot_{a}({}^{*_{a}}(y), u), y)))) \end{array}}}
\vnbprop{pos-lf-sandwich-bound}{\vnbstmt{\begin{array}{@{}l@{}} \forall a, f.; \forall x, y \in \operatorname{vec}(a). \\ \quad (\operatorname{is\text{-}unital\text{-}involutive\text{-}banach\text{-}algebra}(a) \wedge \operatorname{is\text{-}positive\text{-}lf}(a, f)) \Rightarrow  \\ \quad \quad \operatorname{magnitude}(f(\cdot_{a}(\cdot_{a}({}^{*_{a}}(y), x), y))) \\ \quad \quad \leq\; \operatorname{vnrm}(a)(x) \cdot f(\cdot_{a}({}^{*_{a}}(y), y)) \end{array}}}
\vnbprop{uiba-vnrm-real}{\vnbstmt{\begin{array}{@{}l@{}} \forall a.; \forall x \in \operatorname{vec}(a). \\ \quad \operatorname{is\text{-}unital\text{-}involutive\text{-}banach\text{-}algebra}(a) \Rightarrow  \\ \quad \quad (\operatorname{vnrm}(a)(x) \in \mathbb{R}) \wedge (0 \leq \operatorname{vnrm}(a)(x)) \end{array}}}
\vnbprop{uiba-star-mul-le}{\vnbstmt{\begin{array}{@{}l@{}} \forall a.; \forall x \in \operatorname{vec}(a). \\ \quad \operatorname{is\text{-}unital\text{-}involutive\text{-}banach\text{-}algebra}(a) \Rightarrow  \\ \quad \quad \operatorname{vnrm}(a)(\cdot_{a}({}^{*_{a}}(x), x)) \leq (\operatorname{vnrm}(a)(x) \cdot \operatorname{vnrm}(a)(x)) \end{array}}}

\vnbbookfile{theorem-library/ba4-rep}{theorem-library/ba4-rep}
\noindent{\small ba4-rep.scm -- THE FUNCTIONAL OF A REPRESENTATION: BanachAlgebras.pdf Proposition 4.11. Agent BA-1111-B, 2026-11-11.  Destination: theorem-library/ba4-rep.scm.  Library: operator-algebras. Load window: after ba4-pos (pos-lf-norm), ba-bh2 (B(H) and its value lemmas), hilbert-adjoint (op-adjoint-ip, -bounded, ...}\par
\vnbprop{bh-ro-vec}{\vnbstmt{\forall v.\; (\operatorname{is\text{-}hilbert\text{-}space}(v) \Rightarrow (\operatorname{vec}(L({v})) \simeq L({v}, {v})))}}
\vnbprop{bh-ro-vadd}{\vnbstmt{\forall v.\; (\operatorname{is\text{-}hilbert\text{-}space}(v) \Rightarrow (\operatorname{vadd}(L({v})) \simeq \operatorname{bh\text{-}vadd}(v)))}}
\vnbprop{bh-ro-act}{\vnbstmt{\forall v.\; (\operatorname{is\text{-}hilbert\text{-}space}(v) \Rightarrow (\operatorname{act}(L({v})) \simeq \operatorname{bh\text{-}act}(v)))}}
\vnbprop{bh-ro-vmul}{\vnbstmt{\forall v.\; (\operatorname{is\text{-}hilbert\text{-}space}(v) \Rightarrow (\cdot_{L({v})} \simeq \operatorname{bh\text{-}vmul}(v)))}}
\vnbprop{bh-ro-vstar}{\vnbstmt{\forall v.\; (\operatorname{is\text{-}hilbert\text{-}space}(v) \Rightarrow ({}^{*_{L({v})}} \simeq \operatorname{bh\text{-}vstar}(v)))}}
\vnbprop{bh-ro-vone}{\vnbstmt{\forall v.\; (\operatorname{is\text{-}hilbert\text{-}space}(v) \Rightarrow (1_{L({v})} \simeq 1))}}
\vnbprop{bh-ro-vnrm}{\vnbstmt{\forall v.\; (\operatorname{is\text{-}hilbert\text{-}space}(v) \Rightarrow (\operatorname{vnrm}(L({v})) \simeq \operatorname{bh\text{-}vnrm}(v)))}}
\vnbprop{rep-functional-value}{\vnbstmt{\begin{array}{@{}l@{}} \forall a, v, p.; \forall e \in \operatorname{vec}(v), x \in \operatorname{vec}(a). \\ \quad (\operatorname{is\text{-}unital\text{-}involutive\text{-}banach\text{-}algebra}(a) \wedge \operatorname{is\text{-}representation}(a, v, p)) \Rightarrow  \\ \quad \quad f_{{p},{e}}(x) = \langle\cdot,\cdot\rangle_{v}(p(x)(e), e) \end{array}}}
\vnbprop{rep-functional-positive}{\vnbstmt{\begin{array}{@{}l@{}} \forall a, v, p.; \forall e \in \operatorname{vec}(v). \\ \quad (\operatorname{is\text{-}unital\text{-}involutive\text{-}banach\text{-}algebra}(a) \wedge \operatorname{is\text{-}representation}(a, v, p)) \Rightarrow  \\ \quad \quad \operatorname{is\text{-}positive\text{-}lf}(a, f_{{p},{e}}) \end{array}}}
\vnbprop{rep-functional-norm}{\vnbstmt{\begin{array}{@{}l@{}} \forall a, v, p.; \forall e \in \operatorname{vec}(v). \\ \quad (\operatorname{is\text{-}unital\text{-}involutive\text{-}banach\text{-}algebra}(a) \wedge \operatorname{is\text{-}representation}(a, v, p)) \Rightarrow  \\ \quad \quad f_{{p},{e}} \in \operatorname{cdual\text{-}vec}(a) \wedge \\ \quad \quad \operatorname{cdual\text{-}norm}(a, f_{{p},{e}}) = (\lVert e \rVert \cdot \lVert e \rVert) \end{array}}}
\vnbprop{rep-functional-square}{\vnbstmt{\begin{array}{@{}l@{}} \forall a, v, p.; \forall e \in \operatorname{vec}(v), x \in \operatorname{vec}(a). \\ \quad (\operatorname{is\text{-}unital\text{-}involutive\text{-}banach\text{-}algebra}(a) \wedge \operatorname{is\text{-}representation}(a, v, p)) \Rightarrow  \\ \quad \quad f_{{p},{e}}(\cdot_{a}({}^{*_{a}}(x), x)) = \langle\cdot,\cdot\rangle_{v}(p(x)(e), p(x)(e)) \end{array}}}
\vnbprop{rep-functional-unit}{\vnbstmt{\begin{array}{@{}l@{}} \forall a, v, p.; \forall e \in \operatorname{vec}(v). \\ \quad (\operatorname{is\text{-}unital\text{-}involutive\text{-}banach\text{-}algebra}(a) \wedge \operatorname{is\text{-}representation}(a, v, p)) \Rightarrow  \\ \quad \quad f_{{p},{e}}(1_{a}) = \langle\cdot,\cdot\rangle_{v}(e, e) \end{array}}}
\vnbprop{rep-contractive}{\vnbstmt{\begin{array}{@{}l@{}} \forall a, v, p.; \forall x \in \operatorname{vec}(a). \\ \quad (\operatorname{is\text{-}unital\text{-}involutive\text{-}banach\text{-}algebra}(a) \wedge \operatorname{is\text{-}representation}(a, v, p)) \Rightarrow  \\ \quad \quad \operatorname{vnrm}(L({v}))(p(x)) \leq \operatorname{vnrm}(a)(x) \end{array}}}
\vnbprop{rep-norm-one}{\vnbstmt{\begin{array}{@{}l@{}} \forall a, v, p. \\ \quad \operatorname{is\text{-}unital\text{-}involutive\text{-}banach\text{-}algebra}(a) \wedge \\ \quad \operatorname{is\text{-}representation}(a, v, p) \wedge \\ \quad \exists x \in \operatorname{vec}(v).\; \neg (x = \operatorname{vzero}(v)) \Rightarrow \\ \quad \quad \operatorname{vnrm}(L({v}))(p(1_{a})) = 1 \end{array}}}

\vnbbookfile{theorem-library/gns-unique}{theorem-library/gns-unique}
\noindent{\small gns-unique.scm -- THE UNIQUENESS HALF OF THE GNS THEOREM, first step: F. J. Thayer, "Banach Algebras and Applications" (docs/BanachAlgebras.pdf), Theorem 4.13 (the unitary W, [b] |-$>$ rho(b) eta). Agent BA-1117-GNS2, wave 2026-11-17 (round 2). Destination: theorem-library/gns-unique.scm.  Library: operator-algebras.  ...}\par
\vnbprop{rep-functional-polar}{\vnbstmt{\begin{array}{@{}l@{}} \forall a, v, p.; \forall e \in \operatorname{vec}(v), b, c \in \operatorname{vec}(a). \\ \quad (\operatorname{is\text{-}unital\text{-}involutive\text{-}banach\text{-}algebra}(a) \wedge \operatorname{is\text{-}representation}(a, v, p)) \Rightarrow  \\ \quad \quad f_{{p},{e}}(\cdot_{a}({}^{*_{a}}(c), b)) = \langle\cdot,\cdot\rangle_{v}(p(b)(e), p(c)(e)) \end{array}}}

\vnbbookfile{theorem-library/ba4-gns}{theorem-library/ba4-gns}
\noindent{\small ba4-gns.scm -- THE GNS CONSTRUCTION, pre-Hilbert stage: BanachAlgebras.pdf Theorem 4.13 (87) and the first half of its proof.  Agent BA-1111-B, 2026-11-11.  Destination: theorem-library/ba4-gns.scm. Library: operator-algebras.  Load window: after ba4-pos (Lemma 4.4, Corollary 4.10), ba4-rep, ...}\par
\vnbprop{gns-md-vec}{\vnbstmt{\begin{array}{@{}l@{}} \forall a. \\ \quad \operatorname{is\text{-}unital\text{-}involutive\text{-}banach\text{-}algebra}(a) \Rightarrow (\operatorname{vec}(\operatorname{gns\text{-}module}(a)) \simeq \operatorname{vec}(a)) \end{array}}}
\vnbprop{gns-md-vadd}{\vnbstmt{\begin{array}{@{}l@{}} \forall a. \\ \quad \operatorname{is\text{-}unital\text{-}involutive\text{-}banach\text{-}algebra}(a) \Rightarrow (\operatorname{vadd}(\operatorname{gns\text{-}module}(a)) \simeq \operatorname{vadd}(a)) \end{array}}}
\vnbprop{gns-md-act}{\vnbstmt{\begin{array}{@{}l@{}} \forall a. \\ \quad \operatorname{is\text{-}unital\text{-}involutive\text{-}banach\text{-}algebra}(a) \Rightarrow (\operatorname{act}(\operatorname{gns\text{-}module}(a)) \simeq \operatorname{act}(a)) \end{array}}}
\vnbprop{gns-md-vzero}{\vnbstmt{\begin{array}{@{}l@{}} \forall a. \\ \quad \operatorname{is\text{-}unital\text{-}involutive\text{-}banach\text{-}algebra}(a) \Rightarrow (\operatorname{vzero}(\operatorname{gns\text{-}module}(a)) \simeq \operatorname{vzero}(a)) \end{array}}}
\vnbprop{gns-md-vneg}{\vnbstmt{\begin{array}{@{}l@{}} \forall a. \\ \quad \operatorname{is\text{-}unital\text{-}involutive\text{-}banach\text{-}algebra}(a) \Rightarrow (\operatorname{vneg}(\operatorname{gns\text{-}module}(a)) \simeq \operatorname{vneg}(a)) \end{array}}}
\vnbprop{gns-md-scal}{\vnbstmt{\begin{array}{@{}l@{}} \forall a. \\ \quad \operatorname{is\text{-}unital\text{-}involutive\text{-}banach\text{-}algebra}(a) \Rightarrow  \\ \quad \quad \operatorname{scal}(\operatorname{gns\text{-}module}(a)) = \operatorname{normed\text{-}field\text{-}as\text{-}commutative\text{-}ring}(\operatorname{cc\text{-}normed\text{-}field}) \end{array}}}
\vnbprop{gns-form-value}{\vnbstmt{\begin{array}{@{}l@{}} \forall a, f.; \forall x, y \in \operatorname{vec}(a). \\ \quad (\operatorname{is\text{-}unital\text{-}involutive\text{-}banach\text{-}algebra}(a) \wedge \operatorname{is\text{-}positive\text{-}lf}(a, f)) \Rightarrow  \\ \quad \quad \operatorname{gns\text{-}form}(a, f)(x, y) = f(\cdot_{a}({}^{*_{a}}(y), x)) \end{array}}}
\vnbprop{gns-form-psd}{\vnbstmt{\begin{array}{@{}l@{}} \forall a, f. \\ \quad (\operatorname{is\text{-}unital\text{-}involutive\text{-}banach\text{-}algebra}(a) \wedge \operatorname{is\text{-}positive\text{-}lf}(a, f)) \Rightarrow  \\ \quad \quad \operatorname{is\text{-}psd\text{-}form}(\operatorname{gns\text{-}module}(a), \operatorname{gns\text{-}form}(a, f)) \end{array}}}
\vnbprop{gns-module-is-module}{\vnbstmt{\begin{array}{@{}l@{}} \forall a. \\ \quad \operatorname{is\text{-}unital\text{-}involutive\text{-}banach\text{-}algebra}(a) \Rightarrow \operatorname{is\text{-}module}(\operatorname{gns\text{-}module}(a)) \end{array}}}
\vnbprop{gns-pre-is-cips}{\vnbstmt{\begin{array}{@{}l@{}} \forall a, f. \\ \quad (\operatorname{is\text{-}unital\text{-}involutive\text{-}banach\text{-}algebra}(a) \wedge \operatorname{is\text{-}positive\text{-}lf}(a, f)) \Rightarrow  \\ \quad \quad \operatorname{is\text{-}complex\text{-}inner\text{-}product\text{-}space}(\operatorname{gns\text{-}pre}(a, f)) \end{array}}}
\vnbprop{gns-setoid-is-setoid}{\vnbstmt{\begin{array}{@{}l@{}} \forall a, f. \\ \quad (\operatorname{is\text{-}unital\text{-}involutive\text{-}banach\text{-}algebra}(a) \wedge \operatorname{is\text{-}positive\text{-}lf}(a, f)) \Rightarrow  \\ \quad \quad \operatorname{is\text{-}setoid}(\operatorname{gns\text{-}setoid}(a, f)) \end{array}}}
\vnbprop{gns-setoid-pts}{\vnbstmt{\begin{array}{@{}l@{}} \forall a, f. \\ \quad (\operatorname{is\text{-}unital\text{-}involutive\text{-}banach\text{-}algebra}(a) \wedge \operatorname{is\text{-}positive\text{-}lf}(a, f)) \Rightarrow  \\ \quad \quad (\operatorname{pts}(\operatorname{gns\text{-}setoid}(a, f)) \simeq \operatorname{vec}(a)) \end{array}}}
\vnbprop{gns-null-invariant}{\vnbstmt{\begin{array}{@{}l@{}} \forall a, f.; \forall x, b \in \operatorname{vec}(a). \\ \quad \operatorname{is\text{-}unital\text{-}involutive\text{-}banach\text{-}algebra}(a) \wedge \\ \quad \operatorname{is\text{-}positive\text{-}lf}(a, f) \wedge \\ \quad f(\cdot_{a}({}^{*_{a}}(b), b)) = 0 \Rightarrow \\ \quad \quad f(\cdot_{a}({}^{*_{a}}(\cdot_{a}(x, b)), \cdot_{a}(x, b))) = 0 \end{array}}}
\vnbprop{gns-related-iff}{\vnbstmt{\begin{array}{@{}l@{}} \forall a, f.; \forall b, c \in \operatorname{vec}(a). \\ \quad (\operatorname{is\text{-}unital\text{-}involutive\text{-}banach\text{-}algebra}(a) \wedge \operatorname{is\text{-}positive\text{-}lf}(a, f)) \Rightarrow  \\ \quad \quad \operatorname{related}(\operatorname{gns\text{-}setoid}(a, f), b, c) \\ \quad \quad \Leftrightarrow\; f(\cdot_{a}({}^{*_{a}}(\operatorname{vadd}(a)(b, \operatorname{vneg}(a)(c))), \operatorname{vadd}(a)(b, \operatorname{vneg}(a)(c)))) = 0 \end{array}}}
\vnbprop{gns-pre-vec}{\vnbstmt{\begin{array}{@{}l@{}} \forall a, f. \\ \quad (\operatorname{is\text{-}unital\text{-}involutive\text{-}banach\text{-}algebra}(a) \wedge \operatorname{is\text{-}positive\text{-}lf}(a, f)) \Rightarrow  \\ \quad \quad (\operatorname{vec}(\operatorname{gns\text{-}pre}(a, f)) \simeq \operatorname{quotient}(\operatorname{gns\text{-}setoid}(a, f))) \end{array}}}
\vnbprop{gns-pre-ip}{\vnbstmt{\begin{array}{@{}l@{}} \forall a, f. \\ \quad (\operatorname{is\text{-}unital\text{-}involutive\text{-}banach\text{-}algebra}(a) \wedge \operatorname{is\text{-}positive\text{-}lf}(a, f)) \Rightarrow  \\ \quad \quad (\langle\cdot,\cdot\rangle_{\operatorname{gns\text{-}pre}(a, f)} \simeq \operatorname{form\text{-}qip}(\operatorname{gns\text{-}module}(a), \operatorname{gns\text{-}form}(a, f))) \end{array}}}
\vnbprop{gns-ip-computes}{\vnbstmt{\begin{array}{@{}l@{}} \forall a, f.; \forall b, c \in \operatorname{vec}(a). \\ \quad (\operatorname{is\text{-}unital\text{-}involutive\text{-}banach\text{-}algebra}(a) \wedge \operatorname{is\text{-}positive\text{-}lf}(a, f)) \Rightarrow  \\ \quad \quad \langle\cdot,\cdot\rangle_{\operatorname{gns\text{-}pre}(a, f)}(\operatorname{class}(\operatorname{gns\text{-}setoid}(a, f), b), \operatorname{class}(\operatorname{gns\text{-}setoid}(a, f), c)) \\ \quad \quad =\; f(\cdot_{a}({}^{*_{a}}(c), b)) \end{array}}}
\vnbprop{gns-op-ok}{\vnbstmt{\begin{array}{@{}l@{}} \forall a, f.; \forall x \in \operatorname{vec}(a). \\ \quad (\operatorname{is\text{-}unital\text{-}involutive\text{-}banach\text{-}algebra}(a) \wedge \operatorname{is\text{-}positive\text{-}lf}(a, f)) \Rightarrow  \\ \quad \quad (\lambda b \in \operatorname{vec}(a).\; \operatorname{class}(\operatorname{gns\text{-}setoid}(a, f), \cdot_{a}(x, b))) \\ \quad \quad \in\; (\operatorname{pts}(\operatorname{gns\text{-}setoid}(a, f)) \to \operatorname{quotient}(\operatorname{gns\text{-}setoid}(a, f))) \wedge \\ \quad \quad \operatorname{respects}(\operatorname{gns\text{-}setoid}(a, f), (\lambda b \in \operatorname{vec}(a).\; \operatorname{class}(\operatorname{gns\text{-}setoid}(a, f), \cdot_{a}(x, b)))) \end{array}}}
\vnbprop{gns-op-computes}{\vnbstmt{\begin{array}{@{}l@{}} \forall a, f.; \forall x, b \in \operatorname{vec}(a). \\ \quad (\operatorname{is\text{-}unital\text{-}involutive\text{-}banach\text{-}algebra}(a) \wedge \operatorname{is\text{-}positive\text{-}lf}(a, f)) \Rightarrow  \\ \quad \quad \operatorname{gns\text{-}op}(a, f, x)(\operatorname{class}(\operatorname{gns\text{-}setoid}(a, f), b)) = \operatorname{class}(\operatorname{gns\text{-}setoid}(a, f), \cdot_{a}(x, b)) \end{array}}}
\vnbprop{gns-op-type}{\vnbstmt{\begin{array}{@{}l@{}} \forall a, f.; \forall x \in \operatorname{vec}(a). \\ \quad (\operatorname{is\text{-}unital\text{-}involutive\text{-}banach\text{-}algebra}(a) \wedge \operatorname{is\text{-}positive\text{-}lf}(a, f)) \Rightarrow  \\ \quad \quad \operatorname{gns\text{-}op}(a, f, x) \in (\operatorname{vec}(\operatorname{gns\text{-}pre}(a, f)) \to \operatorname{vec}(\operatorname{gns\text{-}pre}(a, f))) \end{array}}}
\vnbprop{gns-vector-in}{\vnbstmt{\begin{array}{@{}l@{}} \forall a, f. \\ \quad (\operatorname{is\text{-}unital\text{-}involutive\text{-}banach\text{-}algebra}(a) \wedge \operatorname{is\text{-}positive\text{-}lf}(a, f)) \Rightarrow  \\ \quad \quad \operatorname{gns\text{-}vector}(a, f) \in \operatorname{vec}(\operatorname{gns\text{-}pre}(a, f)) \end{array}}}
\vnbprop{gns-represents}{\vnbstmt{\begin{array}{@{}l@{}} \forall a, f.; \forall x \in \operatorname{vec}(a). \\ \quad (\operatorname{is\text{-}unital\text{-}involutive\text{-}banach\text{-}algebra}(a) \wedge \operatorname{is\text{-}positive\text{-}lf}(a, f)) \Rightarrow  \\ \quad \quad f(x) = \langle\cdot,\cdot\rangle_{\operatorname{gns\text{-}pre}(a, f)}(\operatorname{gns\text{-}op}(a, f, x)(\operatorname{gns\text{-}vector}(a, f)), \operatorname{gns\text{-}vector}(a, f)) \end{array}}}
\vnbprop{gns-cyclic}{\vnbstmt{\begin{array}{@{}l@{}} \forall a, f.; \forall u \in \operatorname{vec}(\operatorname{gns\text{-}pre}(a, f)). \\ \quad (\operatorname{is\text{-}unital\text{-}involutive\text{-}banach\text{-}algebra}(a) \wedge \operatorname{is\text{-}positive\text{-}lf}(a, f)) \Rightarrow  \\ \quad \quad \exists b \in \operatorname{vec}(a).\; (u = \operatorname{gns\text{-}op}(a, f, b)(\operatorname{gns\text{-}vector}(a, f))) \end{array}}}
\vnbprop{gns-op-adjoint}{\vnbstmt{\begin{array}{@{}l@{}} \forall a, f.; \forall x \in \operatorname{vec}(a), u, w \in \operatorname{vec}(\operatorname{gns\text{-}pre}(a, f)). \\ \quad (\operatorname{is\text{-}unital\text{-}involutive\text{-}banach\text{-}algebra}(a) \wedge \operatorname{is\text{-}positive\text{-}lf}(a, f)) \Rightarrow  \\ \quad \quad \langle\cdot,\cdot\rangle_{\operatorname{gns\text{-}pre}(a, f)}(\operatorname{gns\text{-}op}(a, f, x)(u), w) \\ \quad \quad =\; \langle\cdot,\cdot\rangle_{\operatorname{gns\text{-}pre}(a, f)}(u, \operatorname{gns\text{-}op}(a, f, {}^{*_{a}}(x))(w)) \end{array}}}
\vnbprop{gns-op-mul}{\vnbstmt{\begin{array}{@{}l@{}} \forall a, f.; \forall x, y \in \operatorname{vec}(a), u \in \operatorname{vec}(\operatorname{gns\text{-}pre}(a, f)). \\ \quad (\operatorname{is\text{-}unital\text{-}involutive\text{-}banach\text{-}algebra}(a) \wedge \operatorname{is\text{-}positive\text{-}lf}(a, f)) \Rightarrow  \\ \quad \quad \operatorname{gns\text{-}op}(a, f, \cdot_{a}(x, y))(u) = \operatorname{gns\text{-}op}(a, f, x)(\operatorname{gns\text{-}op}(a, f, y)(u)) \end{array}}}
\vnbprop{gns-op-unit}{\vnbstmt{\begin{array}{@{}l@{}} \forall a, f.; \forall u \in \operatorname{vec}(\operatorname{gns\text{-}pre}(a, f)). \\ \quad (\operatorname{is\text{-}unital\text{-}involutive\text{-}banach\text{-}algebra}(a) \wedge \operatorname{is\text{-}positive\text{-}lf}(a, f)) \Rightarrow  \\ \quad \quad \operatorname{gns\text{-}op}(a, f, 1_{a})(u) = u \end{array}}}
\vnbprop{gns-op-bounded}{\vnbstmt{\begin{array}{@{}l@{}} \forall a, f.; \forall x \in \operatorname{vec}(a), u \in \operatorname{vec}(\operatorname{gns\text{-}pre}(a, f)). \\ \quad (\operatorname{is\text{-}unital\text{-}involutive\text{-}banach\text{-}algebra}(a) \wedge \operatorname{is\text{-}positive\text{-}lf}(a, f)) \Rightarrow  \\ \quad \quad \langle\cdot,\cdot\rangle_{\operatorname{gns\text{-}pre}(a, f)}(\operatorname{gns\text{-}op}(a, f, x)(u), \operatorname{gns\text{-}op}(a, f, x)(u)) \\ \quad \quad \leq\; (\operatorname{vnrm}(a)(x) \cdot \operatorname{vnrm}(a)(x)) \cdot \langle\cdot,\cdot\rangle_{\operatorname{gns\text{-}pre}(a, f)}(u, u) \end{array}}}
\vnbprop{gns-pre-vadd}{\vnbstmt{\begin{array}{@{}l@{}} \forall a, f. \\ \quad (\operatorname{is\text{-}unital\text{-}involutive\text{-}banach\text{-}algebra}(a) \wedge \operatorname{is\text{-}positive\text{-}lf}(a, f)) \Rightarrow  \\ \quad \quad (\operatorname{vadd}(\operatorname{gns\text{-}pre}(a, f)) \simeq \operatorname{vadd}(\operatorname{quotient\text{-}module}(\operatorname{gns\text{-}module}(a), N_{\operatorname{gns\text{-}form}(a, f)}))) \end{array}}}
\vnbprop{gns-pre-act}{\vnbstmt{\begin{array}{@{}l@{}} \forall a, f. \\ \quad (\operatorname{is\text{-}unital\text{-}involutive\text{-}banach\text{-}algebra}(a) \wedge \operatorname{is\text{-}positive\text{-}lf}(a, f)) \Rightarrow  \\ \quad \quad (\operatorname{act}(\operatorname{gns\text{-}pre}(a, f)) \simeq \operatorname{act}(\operatorname{quotient\text{-}module}(\operatorname{gns\text{-}module}(a), N_{\operatorname{gns\text{-}form}(a, f)}))) \end{array}}}
\vnbprop{gns-pre-vadd-computes}{\vnbstmt{\begin{array}{@{}l@{}} \forall a, f.; \forall b, c \in \operatorname{vec}(a). \\ \quad (\operatorname{is\text{-}unital\text{-}involutive\text{-}banach\text{-}algebra}(a) \wedge \operatorname{is\text{-}positive\text{-}lf}(a, f)) \Rightarrow  \\ \quad \quad \operatorname{vadd}(\operatorname{gns\text{-}pre}(a, f))(\operatorname{class}(\operatorname{gns\text{-}setoid}(a, f), b), \operatorname{class}(\operatorname{gns\text{-}setoid}(a, f), c)) \\ \quad \quad =\; \operatorname{class}(\operatorname{gns\text{-}setoid}(a, f), \operatorname{vadd}(a)(b, c)) \end{array}}}
\vnbprop{gns-pre-act-computes}{\vnbstmt{\begin{array}{@{}l@{}} \forall a, f.; \forall c \in \mathbb{C}, b \in \operatorname{vec}(a). \\ \quad (\operatorname{is\text{-}unital\text{-}involutive\text{-}banach\text{-}algebra}(a) \wedge \operatorname{is\text{-}positive\text{-}lf}(a, f)) \Rightarrow  \\ \quad \quad \operatorname{act}(\operatorname{gns\text{-}pre}(a, f))(c, \operatorname{class}(\operatorname{gns\text{-}setoid}(a, f), b)) = \operatorname{class}(\operatorname{gns\text{-}setoid}(a, f), \operatorname{act}(a)(c, b)) \end{array}}}
\vnbprop{gns-op-vadd}{\vnbstmt{\begin{array}{@{}l@{}} \forall a, f.; \forall x \in \operatorname{vec}(a), u, w \in \operatorname{vec}(\operatorname{gns\text{-}pre}(a, f)). \\ \quad (\operatorname{is\text{-}unital\text{-}involutive\text{-}banach\text{-}algebra}(a) \wedge \operatorname{is\text{-}positive\text{-}lf}(a, f)) \Rightarrow  \\ \quad \quad \operatorname{gns\text{-}op}(a, f, x)(\operatorname{vadd}(\operatorname{gns\text{-}pre}(a, f))(u, w)) \\ \quad \quad =\; \operatorname{vadd}(\operatorname{gns\text{-}pre}(a, f))(\operatorname{gns\text{-}op}(a, f, x)(u), \operatorname{gns\text{-}op}(a, f, x)(w)) \end{array}}}
\vnbprop{gns-op-act}{\vnbstmt{\begin{array}{@{}l@{}} \forall a, f.; \forall x \in \operatorname{vec}(a), c \in \mathbb{C}, u \in \operatorname{vec}(\operatorname{gns\text{-}pre}(a, f)). \\ \quad (\operatorname{is\text{-}unital\text{-}involutive\text{-}banach\text{-}algebra}(a) \wedge \operatorname{is\text{-}positive\text{-}lf}(a, f)) \Rightarrow  \\ \quad \quad \operatorname{gns\text{-}op}(a, f, x)(\operatorname{act}(\operatorname{gns\text{-}pre}(a, f))(c, u)) = \operatorname{act}(\operatorname{gns\text{-}pre}(a, f))(c, \operatorname{gns\text{-}op}(a, f, x)(u)) \end{array}}}
\vnbprop{gns-op-sum}{\vnbstmt{\begin{array}{@{}l@{}} \forall a, f.; \forall x, y \in \operatorname{vec}(a), u \in \operatorname{vec}(\operatorname{gns\text{-}pre}(a, f)). \\ \quad (\operatorname{is\text{-}unital\text{-}involutive\text{-}banach\text{-}algebra}(a) \wedge \operatorname{is\text{-}positive\text{-}lf}(a, f)) \Rightarrow  \\ \quad \quad \operatorname{gns\text{-}op}(a, f, \operatorname{vadd}(a)(x, y))(u) \\ \quad \quad =\; \operatorname{vadd}(\operatorname{gns\text{-}pre}(a, f))(\operatorname{gns\text{-}op}(a, f, x)(u), \operatorname{gns\text{-}op}(a, f, y)(u)) \end{array}}}
\vnbprop{gns-op-scale}{\vnbstmt{\begin{array}{@{}l@{}} \forall a, f.; \forall c \in \mathbb{C}, x \in \operatorname{vec}(a), u \in \operatorname{vec}(\operatorname{gns\text{-}pre}(a, f)). \\ \quad (\operatorname{is\text{-}unital\text{-}involutive\text{-}banach\text{-}algebra}(a) \wedge \operatorname{is\text{-}positive\text{-}lf}(a, f)) \Rightarrow  \\ \quad \quad \operatorname{gns\text{-}op}(a, f, \operatorname{act}(a)(c, x))(u) = \operatorname{act}(\operatorname{gns\text{-}pre}(a, f))(c, \operatorname{gns\text{-}op}(a, f, x)(u)) \end{array}}}
\vnbprop{gns-matrix-coeff}{\vnbstmt{\begin{array}{@{}l@{}} \forall a, f.; \forall x, y, z \in \operatorname{vec}(a). \\ \quad (\operatorname{is\text{-}unital\text{-}involutive\text{-}banach\text{-}algebra}(a) \wedge \operatorname{is\text{-}positive\text{-}lf}(a, f)) \Rightarrow  \\ \quad \quad \langle\cdot,\cdot\rangle_{\operatorname{gns\text{-}pre}(a, f)}(\operatorname{gns\text{-}op}(a, f, x)(\operatorname{gns\text{-}op}(a, f, y)(\operatorname{gns\text{-}vector}(a, f))), \\ \quad \quad \quad \operatorname{gns\text{-}op}(a, f, z)(\operatorname{gns\text{-}vector}(a, f))) \\ \quad \quad =\; f(\cdot_{a}({}^{*_{a}}(z), \cdot_{a}(x, y))) \end{array}}}

\vnbbookfile{theorem-library/ba4-states}{theorem-library/ba4-states}
\noindent{\small ba4-states.scm -- STATES: BanachAlgebras.pdf Corollary 4.9 (states have norm 1), the convexity of Sigma(A) and "a pure state is an extreme point" (Proposition 4.24, first half of its proof). Agent BA-1111-B, 2026-11-11.  Destination: theorem-library/ba4-states.scm.  Library: operator-algebras. Load window: after ...}\par
\vnbprop{states-unfold}{\vnbstmt{\begin{array}{@{}l@{}} \forall a. \\ \quad (\Sigma({a}) \simeq \{f \in (\operatorname{vec}(a) \to \mathbb{C}) : (\operatorname{is\text{-}positive\text{-}lf}(a, f) \wedge (f(1_{a}) = 1))\}) \end{array}}}
\vnbprop{pos-lf-combo}{\vnbstmt{\begin{array}{@{}l@{}} \forall a, g, h.; \forall b, c \in \mathbb{R}. \\ \quad \operatorname{is\text{-}unital\text{-}involutive\text{-}banach\text{-}algebra}(a) \wedge \\ \quad \operatorname{is\text{-}positive\text{-}lf}(a, g) \wedge \\ \quad \operatorname{is\text{-}positive\text{-}lf}(a, h) \wedge \\ \quad 0 \leq b \wedge \\ \quad 0 \leq c \Rightarrow \\ \quad \quad \operatorname{is\text{-}positive\text{-}lf}(a, (\lambda w \in \operatorname{vec}(a).\; ((b \cdot g(w)) + (c \cdot h(w))))) \end{array}}}
\vnbprop{states-norm-one}{\vnbstmt{\begin{array}{@{}l@{}} \forall a.; \forall g \in \Sigma({a}). \\ \quad \operatorname{is\text{-}unital\text{-}involutive\text{-}banach\text{-}algebra}(a) \Rightarrow  \\ \quad \quad (g \in \operatorname{cdual\text{-}vec}(a)) \wedge (\operatorname{cdual\text{-}norm}(a, g) = 1) \end{array}}}
\vnbprop{states-convex}{\vnbstmt{\begin{array}{@{}l@{}} \forall a.; \forall g, h \in \Sigma({a}), r \in \mathbb{R}. \\ \quad (\operatorname{is\text{-}unital\text{-}involutive\text{-}banach\text{-}algebra}(a) \wedge ((0 \leq r) \wedge (r \leq 1))) \Rightarrow  \\ \quad \quad (\lambda w \in \operatorname{vec}(a).\; ((r \cdot g(w)) + ((1 - r) \cdot h(w)))) \in \Sigma({a}) \end{array}}}
\vnbprop{pure-state-extreme}{\vnbstmt{\begin{array}{@{}l@{}} \forall a, f.; \forall g, h \in \Sigma({a}), r \in \mathbb{R}, x \in \operatorname{vec}(a). \\ \quad \operatorname{is\text{-}unital\text{-}involutive\text{-}banach\text{-}algebra}(a) \wedge \\ \quad \operatorname{is\text{-}pure\text{-}state}(a, f) \wedge \\ \quad 0 < r \wedge \\ \quad r < 1 \wedge \\ \quad \forall y \in \operatorname{vec}(a).\; (f(y) = ((r \cdot g(y)) + ((1 - r) \cdot h(y)))) \Rightarrow \\ \quad \quad (g(x) = f(x)) \wedge (h(x) = f(x)) \end{array}}}
\vnbprop{lf-diff-positive}{\vnbstmt{\begin{array}{@{}l@{}} \forall a, f, g. \\ \quad \operatorname{is\text{-}unital\text{-}involutive\text{-}banach\text{-}algebra}(a) \wedge \\ \quad \operatorname{is\text{-}positive\text{-}lf}(a, f) \wedge \\ \quad \operatorname{is\text{-}positive\text{-}lf}(a, g) \wedge \\ \quad \operatorname{functional\text{-}le}(a, g, f) \Rightarrow \\ \quad \quad \operatorname{is\text{-}positive\text{-}lf}(a, (\lambda w \in \operatorname{vec}(a).\; (f(w) + (-1 \cdot g(w))))) \end{array}}}
\vnbprop{pos-lf-vanish}{\vnbstmt{\begin{array}{@{}l@{}} \forall a, f.; \forall x \in \operatorname{vec}(a). \\ \quad (\operatorname{is\text{-}unital\text{-}involutive\text{-}banach\text{-}algebra}(a) \wedge (\operatorname{is\text{-}positive\text{-}lf}(a, f) \wedge (f(1_{a}) = 0))) \Rightarrow  \\ \quad \quad f(x) = 0 \end{array}}}
\vnbprop{extreme-state-pure}{\vnbstmt{\begin{array}{@{}l@{}} \forall a.; \forall f \in \Sigma({a}). \\ \quad \operatorname{is\text{-}unital\text{-}involutive\text{-}banach\text{-}algebra}(a) \wedge \\ \quad \forall g, h \in \Sigma({a}), r \in \mathbb{R}, z \in \operatorname{vec}(a). \\ \quad \quad ((0 < r) \wedge ((r < 1) \wedge \forall y \in \operatorname{vec}(a).\; (f(y) = ((r \cdot g(y)) + ((1 - r) \cdot h(y)))))) \Rightarrow  \\ \quad \quad \quad g(z) = f(z) \Rightarrow \\ \quad \quad \operatorname{is\text{-}pure\text{-}state}(a, f) \end{array}}}

\vnbbookfile{theorem-library/state-extreme}{theorem-library/state-extreme}
\noindent{\small state-extreme.scm -- PROPOSITION 4.24, second sentence: the extreme points of Sigma(A) are the pure states.  F. J. Thayer, "Banach Algebras and Applications", Proposition 4.24. Agent BA-1117-GNS2, 2026-11-17.  Destination: theorem-library/state-extreme.scm. Library: operator-algebras.  Load window: after ba4-states ...}\par
\vnbprop{pure-state-is-extreme}{\vnbstmt{\begin{array}{@{}l@{}} \forall a, f. \\ \quad (\operatorname{is\text{-}unital\text{-}involutive\text{-}banach\text{-}algebra}(a) \wedge \operatorname{is\text{-}pure\text{-}state}(a, f)) \Rightarrow  \\ \quad \quad \operatorname{is\text{-}extreme\text{-}state}(a, f) \end{array}}}
\vnbprop{extreme-state-is-pure}{\vnbstmt{\begin{array}{@{}l@{}} \forall a, f. \\ \quad (\operatorname{is\text{-}unital\text{-}involutive\text{-}banach\text{-}algebra}(a) \wedge \operatorname{is\text{-}extreme\text{-}state}(a, f)) \Rightarrow  \\ \quad \quad \operatorname{is\text{-}pure\text{-}state}(a, f) \end{array}}}
\vnbprop{pure-state-iff-extreme}{\vnbstmt{\begin{array}{@{}l@{}} \forall a, f. \\ \quad \operatorname{is\text{-}unital\text{-}involutive\text{-}banach\text{-}algebra}(a) \Rightarrow  \\ \quad \quad \operatorname{is\text{-}pure\text{-}state}(a, f) \Leftrightarrow \operatorname{is\text{-}extreme\text{-}state}(a, f) \end{array}}}

\vnbbookfile{theorem-library/states-compact}{theorem-library/states-compact}
\noindent{\small states-compact.scm -- THE STATE SPACE IS WEAK* COMPACT: F. J. Thayer, "Banach Algebras and Applications" (docs/BanachAlgebras.pdf), Corollary 4.9 and the first sentence of Proposition 4.24. Agent BA-1117-GNS2, wave 2026-11-17 (round 2). Destination: theorem-library/states-compact.scm.  Library: operator-algebras.}\par
\vnbprop{cc-im-of-real}{\vnbstmt{\forall z \in \mathbb{R}.\; (\operatorname{imag\text{-}part}(z) = 0)}}
\vnbprop{states-closed-in-product}{\vnbstmt{\begin{array}{@{}l@{}} \forall a. \\ \quad \operatorname{is\text{-}unital\text{-}involutive\text{-}banach\text{-}algebra}(a) \Rightarrow  \\ \quad \quad \operatorname{is\text{-}top\text{-}closed}(\operatorname{tproduct\text{-}top}(\operatorname{vec}(a), (\lambda x \in \operatorname{vec}(a).\; \operatorname{metric\text{-}top}(\operatorname{cc\text{-}ms}))), \Sigma({a})) \end{array}}}
\vnbprop{cdual-vec-unfold}{\vnbstmt{\begin{array}{@{}l@{}} \forall a. \\ \quad (\operatorname{cdual\text{-}vec}(a) \simeq \{f \in (\operatorname{vec}(a) \to \mathbb{C}) : (\operatorname{is\text{-}cc\text{-}linear\text{-}on}(a, f) \wedge \exists m \in \mathbb{R}.\; \forall x \in \operatorname{vec}(a).\; (\operatorname{magnitude}(f(x)) \leq (m \cdot \operatorname{vnrm}(a)(x))))\}) \end{array}}}
\vnbprop{cdual-vec-in-fun}{\vnbstmt{\forall a.; \forall f \in \operatorname{cdual\text{-}vec}(a).\; (f \in (\operatorname{vec}(a) \to \mathbb{C}))}}
\vnbprop{states-product-compact}{\vnbstmt{\begin{array}{@{}l@{}} \forall a. \\ \quad \operatorname{is\text{-}unital\text{-}involutive\text{-}banach\text{-}algebra}(a) \Rightarrow  \\ \quad \quad \operatorname{is\text{-}compact\text{-}t}(\operatorname{subspace\text{-}top}(\operatorname{tproduct\text{-}top}(\operatorname{vec}(a), (\lambda x \in \operatorname{vec}(a).\; \operatorname{metric\text{-}top}(\operatorname{cc\text{-}ms}))), \Sigma({a}))) \end{array}}}
\vnbprop{states-weak-star-compact}{\vnbstmt{\begin{array}{@{}l@{}} \forall a. \\ \quad \operatorname{is\text{-}unital\text{-}involutive\text{-}banach\text{-}algebra}(a) \Rightarrow  \\ \quad \quad \operatorname{is\text{-}compact\text{-}t}(\operatorname{subspace\text{-}top}(\operatorname{cdual\text{-}weak\text{-}star\text{-}top}(a), \Sigma({a}))) \end{array}}}

\vnbbookfile{theorem-library/pure-states}{theorem-library/pure-states}
\noindent{\small pure-states.scm -- PURE STATES EXIST, AND Sigma(A) IS THE CLOSED CONVEX HULL OF THEM: F. J. Thayer, "Banach Algebras and Applications" (docs/BanachAlgebras.pdf), Proposition 4.24 (Krein-Milman applied to the state space).  Agent BA-1117-GNS2, wave 2026-11-17 (round 2). Destination: theorem-library/pure-states.scm.  ...}\par
\vnbprop{rfn-module-scal}{\vnbstmt{\begin{array}{@{}l@{}} \forall u \in \mathbf{Set}. \\ \quad (\operatorname{scal}(\operatorname{rfn\text{-}module}(u)) \simeq \operatorname{normed\text{-}field\text{-}as\text{-}commutative\text{-}ring}(\operatorname{rr\text{-}normed\text{-}field})) \end{array}}}
\vnbprop{rfn-module-vec}{\vnbstmt{\forall u \in \mathbf{Set}.\; (\operatorname{vec}(\operatorname{rfn\text{-}module}(u)) \simeq (u \to \mathbb{C}))}}
\vnbprop{rfn-module-vadd}{\vnbstmt{\begin{array}{@{}l@{}} \forall u \in \mathbf{Set}. \\ \quad (\operatorname{vadd}(\operatorname{rfn\text{-}module}(u)) \simeq (\lambda {[}\operatorname{g4hf\_}, \operatorname{g4hg\_}] \in ((u \to \mathbb{C}) \times (u \to \mathbb{C})).\; (\lambda y \in u.\; (\operatorname{g4hf\_}(y) + \operatorname{g4hg\_}(y))))) \end{array}}}
\vnbprop{rfn-module-vzero}{\vnbstmt{\forall u \in \mathbf{Set}.\; (\operatorname{vzero}(\operatorname{rfn\text{-}module}(u)) \simeq (\lambda y \in u.\; 0))}}
\vnbprop{rfn-module-vneg}{\vnbstmt{\begin{array}{@{}l@{}} \forall u \in \mathbf{Set}. \\ \quad (\operatorname{vneg}(\operatorname{rfn\text{-}module}(u)) \simeq (\lambda f \in (u \to \mathbb{C}).\; (\lambda y \in u.\; -f(y)))) \end{array}}}
\vnbprop{rfn-module-act}{\vnbstmt{\begin{array}{@{}l@{}} \forall u \in \mathbf{Set}. \\ \quad (\operatorname{act}(\operatorname{rfn\text{-}module}(u)) \simeq (\lambda {[}\operatorname{g4hr\_}, \operatorname{g4hf\_}] \in (\mathbb{R} \times (u \to \mathbb{C})).\; (\lambda y \in u.\; (\operatorname{g4hr\_} \cdot \operatorname{g4hf\_}(y))))) \end{array}}}
\vnbprop{rfn-module-is-module}{\vnbstmt{\forall u \in \mathbf{Set}.\; \operatorname{is\text{-}module}(\operatorname{rfn\text{-}module}(u))}}
\vnbprop{rfn-vadd-value}{\vnbstmt{\begin{array}{@{}l@{}} \forall u \in \mathbf{Set}, f, g \in (u \to \mathbb{C}), y \in u. \\ \quad \operatorname{vadd}(\operatorname{rfn\text{-}module}(u))(f, g)(y) = (f(y) + g(y)) \end{array}}}
\vnbprop{rfn-act-value}{\vnbstmt{\begin{array}{@{}l@{}} \forall u \in \mathbf{Set}, r \in \mathbb{R}, f \in (u \to \mathbb{C}), y \in u. \\ \quad \operatorname{act}(\operatorname{rfn\text{-}module}(u))(r, f)(y) = (r \cdot f(y)) \end{array}}}
\vnbprop{rfn-vzero-value}{\vnbstmt{\forall u \in \mathbf{Set}, y \in u.\; (\operatorname{vzero}(\operatorname{rfn\text{-}module}(u))(y) = 0)}}
\vnbprop{eval-seminorms-value}{\vnbstmt{\begin{array}{@{}l@{}} \forall u \in \mathbf{Set}, y \in u, f \in (u \to \mathbb{C}). \\ \quad \operatorname{eval\text{-}seminorms}(u)(y)(f) = \operatorname{magnitude}(f(y)) \end{array}}}
\vnbprop{eval-seminorms-family}{\vnbstmt{\begin{array}{@{}l@{}} \forall u \in \mathbf{Set}. \\ \quad \operatorname{is\text{-}seminorm\text{-}family\text{-}on}(\operatorname{rr\text{-}normed\text{-}field}, \operatorname{rfn\text{-}module}(u), u, \operatorname{eval\text{-}seminorms}(u)) \end{array}}}
\vnbprop{rfn-top-unfold}{\vnbstmt{\forall u.\; (\operatorname{rfn\text{-}top}(u) \simeq \operatorname{seminorm\text{-}top}(\operatorname{rfn\text{-}module}(u), u, \operatorname{eval\text{-}seminorms}(u)))}}
\vnbprop{rfn-top-locally-convex}{\vnbstmt{\forall u \in \mathbf{Set}.\; \operatorname{is\text{-}locally\text{-}convex\text{-}tvs}(\operatorname{rfn\text{-}module}(u), \operatorname{rfn\text{-}top}(u))}}
\vnbprop{rfn-vneg-value}{\vnbstmt{\begin{array}{@{}l@{}} \forall u \in \mathbf{Set}, f \in (u \to \mathbb{C}), y \in u. \\ \quad \operatorname{vneg}(\operatorname{rfn\text{-}module}(u))(f)(y) = -f(y) \end{array}}}
\vnbprop{cc-product-ball-constraint}{\vnbstmt{\begin{array}{@{}l@{}} \forall u \in \mathbf{Set}, y \in u, c \in \mathbb{C}, r. \\ \quad \operatorname{pos\text{-}rr}(r) \Rightarrow  \\ \quad \quad {[}y, \operatorname{ball}(\operatorname{cc\text{-}ms}, c, r)] \in \operatorname{tprod\text{-}constr}(u, (\lambda x \in u.\; \operatorname{metric\text{-}top}(\operatorname{cc\text{-}ms}))) \end{array}}}
\vnbprop{rfn-open-product-open}{\vnbstmt{\begin{array}{@{}l@{}} \forall u \in \mathbf{Set}, o \in \operatorname{opens}(\operatorname{rfn\text{-}top}(u)). \\ \quad o \in \operatorname{opens}(\operatorname{tproduct\text{-}top}(u, (\lambda x \in u.\; \operatorname{metric\text{-}top}(\operatorname{cc\text{-}ms})))) \end{array}}}
\vnbprop{compact-t-coarser-subspace}{\vnbstmt{\begin{array}{@{}l@{}} \forall a, b, s. \\ \quad \operatorname{is\text{-}compact\text{-}t}(\operatorname{subspace\text{-}top}(a, s)) \wedge \\ \quad \operatorname{is\text{-}top\text{-}space}(b) \wedge \\ \quad s \subseteq \operatorname{pts}(b) \wedge \\ \quad s \in \mathbf{Set} \wedge \\ \quad \forall w \in \operatorname{opens}(b).\; (w \in \operatorname{opens}(a)) \Rightarrow \\ \quad \quad \operatorname{is\text{-}compact\text{-}t}(\operatorname{subspace\text{-}top}(b, s)) \end{array}}}
\vnbprop{states-rfn-compact}{\vnbstmt{\begin{array}{@{}l@{}} \forall a. \\ \quad \operatorname{is\text{-}unital\text{-}involutive\text{-}banach\text{-}algebra}(a) \Rightarrow  \\ \quad \quad \operatorname{is\text{-}compact\text{-}t}(\operatorname{subspace\text{-}top}(\operatorname{rfn\text{-}top}(\operatorname{vec}(a)), \Sigma({a}))) \end{array}}}
\vnbprop{states-rfn-convex}{\vnbstmt{\begin{array}{@{}l@{}} \forall a. \\ \quad \operatorname{is\text{-}unital\text{-}involutive\text{-}banach\text{-}algebra}(a) \Rightarrow  \\ \quad \quad \operatorname{is\text{-}convex\text{-}subset}(\operatorname{rfn\text{-}module}(\operatorname{vec}(a)), \Sigma({a})) \end{array}}}
\vnbprop{extreme-point-is-extreme-state}{\vnbstmt{\begin{array}{@{}l@{}} \forall a.; \forall x \in \operatorname{extreme\text{-}points}(\operatorname{rfn\text{-}module}(\operatorname{vec}(a)), \Sigma({a})). \\ \quad \operatorname{is\text{-}unital\text{-}involutive\text{-}banach\text{-}algebra}(a) \Rightarrow \operatorname{is\text{-}extreme\text{-}state}(a, x) \end{array}}}
\vnbprop{pure-state-exists}{\vnbstmt{\begin{array}{@{}l@{}} \forall a. \\ \quad (\operatorname{is\text{-}unital\text{-}involutive\text{-}banach\text{-}algebra}(a) \wedge \exists z.\; (z \in \Sigma({a}))) \Rightarrow  \\ \quad \quad \exists p.\; \operatorname{is\text{-}pure\text{-}state}(a, p) \end{array}}}
\vnbprop{states-krein-milman}{\vnbstmt{\begin{array}{@{}l@{}} \forall a. \\ \quad \operatorname{is\text{-}unital\text{-}involutive\text{-}banach\text{-}algebra}(a) \Rightarrow  \\ \quad \quad \Sigma({a}) \\ \quad \quad =\; \operatorname{closed\text{-}conv\text{-}hull}(\operatorname{rfn\text{-}module}(\operatorname{vec}(a)), \operatorname{rfn\text{-}top}(\operatorname{vec}(a)), \operatorname{extreme\text{-}points}(\operatorname{rfn\text{-}module}(\operatorname{vec}(a)), \Sigma({a}))) \end{array}}}
\vnbprop{extreme-state-is-extreme-point}{\vnbstmt{\begin{array}{@{}l@{}} \forall a, x. \\ \quad (\operatorname{is\text{-}unital\text{-}involutive\text{-}banach\text{-}algebra}(a) \wedge \operatorname{is\text{-}extreme\text{-}state}(a, x)) \Rightarrow  \\ \quad \quad x \in \operatorname{extreme\text{-}points}(\operatorname{rfn\text{-}module}(\operatorname{vec}(a)), \Sigma({a})) \end{array}}}
\vnbprop{extreme-points-are-pure-states}{\vnbstmt{\begin{array}{@{}l@{}} \forall a, x. \\ \quad \operatorname{is\text{-}unital\text{-}involutive\text{-}banach\text{-}algebra}(a) \Rightarrow  \\ \quad \quad (x \in \operatorname{extreme\text{-}points}(\operatorname{rfn\text{-}module}(\operatorname{vec}(a)), \Sigma({a}))) \Leftrightarrow \operatorname{is\text{-}pure\text{-}state}(a, x) \end{array}}}
\vnbprop{extreme-points-pure-set}{\vnbstmt{\begin{array}{@{}l@{}} \forall a. \\ \quad \operatorname{is\text{-}unital\text{-}involutive\text{-}banach\text{-}algebra}(a) \Rightarrow  \\ \quad \quad \operatorname{extreme\text{-}points}(\operatorname{rfn\text{-}module}(\operatorname{vec}(a)), \Sigma({a})) = \{p \in (\operatorname{vec}(a) \to \mathbb{C}) : \operatorname{is\text{-}pure\text{-}state}(a, p)\} \end{array}}}
\vnbprop{states-hull-of-pure-states}{\vnbstmt{\begin{array}{@{}l@{}} \forall a. \\ \quad \operatorname{is\text{-}unital\text{-}involutive\text{-}banach\text{-}algebra}(a) \Rightarrow  \\ \quad \quad \Sigma({a}) \\ \quad \quad =\; \operatorname{closed\text{-}conv\text{-}hull}(\operatorname{rfn\text{-}module}(\operatorname{vec}(a)), \operatorname{rfn\text{-}top}(\operatorname{vec}(a)), \{p \in (\operatorname{vec}(a) \to \mathbb{C}) : \operatorname{is\text{-}pure\text{-}state}(a, p)\}) \end{array}}}

\vnbbookfile{theorem-library/ba-optop-adjoint}{theorem-library/ba-optop-adjoint}
\noindent{\small ba-optop-adjoint.scm -- Thayer, Banach Algebras and Applications, Prop 5.11, third clause: S |-$>$ S* is continuous for the weak operator topology.  The proof is the notes': p\_(xi,eta)(S*) = |$<$S* xi | eta$>$| = |$<$xi | S eta$>$| = |$<$S eta | xi$>$| = p\_(eta,xi)(S), with the seminorm continuity condition for a map that is only ...}\par
\vnbprop{op-adjoint-neg}{\vnbstmt{\begin{array}{@{}l@{}} \forall v, t. \\ \quad (\operatorname{is\text{-}hilbert\text{-}space}(v) \wedge (t \in L({v}, {v}))) \Rightarrow  \\ \quad \quad {-{t}}^{*} = -{{t}^{*}} \end{array}}}
\vnbprop{bh-adjoint-map-type}{\vnbstmt{\begin{array}{@{}l@{}} \forall v. \\ \quad \operatorname{is\text{-}hilbert\text{-}space}(v) \Rightarrow  \\ \quad \quad (\lambda s \in L({v}, {v}).\; {s}^{*}) \in (\operatorname{vec}(L({v})) \to \operatorname{vec}(L({v}))) \end{array}}}
\vnbprop{bh-adjoint-map-sub}{\vnbstmt{\begin{array}{@{}l@{}} \forall v.; \forall a, b \in \operatorname{vec}(L({v})). \\ \quad \operatorname{is\text{-}hilbert\text{-}space}(v) \Rightarrow  \\ \quad \quad (\lambda s \in L({v}, {v}).\; {s}^{*})(\operatorname{vadd}(L({v}))(a, \operatorname{vneg}(L({v}))(b))) \\ \quad \quad =\; \operatorname{vadd}(L({v}))((\lambda s \in L({v}, {v}).\; {s}^{*})(a), \\ \quad \quad \quad \operatorname{vneg}(L({v}))((\lambda s \in L({v}, {v}).\; {s}^{*})(b))) \end{array}}}
\vnbprop{bop-adjoint-wot-continuous}{\vnbstmt{\begin{array}{@{}l@{}} \forall v. \\ \quad \operatorname{is\text{-}hilbert\text{-}space}(v) \Rightarrow  \\ \quad \quad \operatorname{is\text{-}hom\text{-}top\text{-}space}(L({v})_{w}, L({v})_{w}, (\lambda s \in L({v}, {v}).\; {s}^{*})) \end{array}}}
\vnbprop{strong-dual-unfold-sot}{\vnbstmt{\begin{array}{@{}l@{}} \forall v. \\ \quad (L_{\bullet}({v}) \simeq \{f \in (L({v}, {v}) \to \mathbb{C}) : (\operatorname{is\text{-}linear\text{-}map}(L({v}), \operatorname{cc\text{-}sn\text{-}module}, f) \wedge \operatorname{is\text{-}hom\text{-}top\text{-}space}(L({v})_{s}, \operatorname{cc\text{-}sn\text{-}top}, f))\}) \end{array}}}
\vnbprop{sot-unfold}{\vnbstmt{\forall v.\; (L({v})_{s} \simeq \operatorname{seminorm\text{-}top}(L({v}), \operatorname{vec}(v), \operatorname{sot\text{-}pf}(v)))}}
\vnbprop{strong-dual-bound}{\vnbstmt{\begin{array}{@{}l@{}} \forall v, f. \\ \quad (\operatorname{is\text{-}complex\text{-}inner\text{-}product\text{-}space}(v) \wedge (f \in L_{\bullet}({v}))) \Rightarrow  \\ \quad \quad \exists \mathit{sbf} \in \mathbf{Set}. \\ \quad \quad \quad |\mathit{sbf}| \in \mathbb{N} \wedge \\ \quad \quad \quad \mathit{sbf} \subseteq \operatorname{vec}(v) \wedge \\ \quad \quad \quad \exists c \in \mathbb{R}. \\ \quad \quad \quad \quad \operatorname{pos\text{-}rr}(c) \wedge \\ \quad \quad \quad \quad \forall x \in L({v}, {v}), t \in \mathbb{R}. \\ \quad \quad \quad \quad \quad ((0 \leq t) \wedge \forall i \in \mathit{sbf}.\; (\lVert x(i) \rVert \leq t)) \Rightarrow  \\ \quad \quad \quad \quad \quad \quad \operatorname{magnitude}(f(x)) \leq (c \cdot t) \end{array}}}
\vnbprop{wot-unfold}{\vnbstmt{\begin{array}{@{}l@{}} \forall v. \\ \quad (L({v})_{w} \simeq \operatorname{seminorm\text{-}top}(L({v}), (\operatorname{vec}(v) \times \operatorname{vec}(v)), \operatorname{wot\text{-}pf}(v))) \end{array}}}
\vnbprop{weak-dual-bound}{\vnbstmt{\begin{array}{@{}l@{}} \forall v, f. \\ \quad (\operatorname{is\text{-}complex\text{-}inner\text{-}product\text{-}space}(v) \wedge (f \in \operatorname{weak\text{-}dual}(v))) \Rightarrow  \\ \quad \quad \exists \mathit{sbf} \in \mathbf{Set}. \\ \quad \quad \quad |\mathit{sbf}| \in \mathbb{N} \wedge \\ \quad \quad \quad \mathit{sbf} \subseteq (\operatorname{vec}(v) \times \operatorname{vec}(v)) \wedge \\ \quad \quad \quad \exists c \in \mathbb{R}. \\ \quad \quad \quad \quad \operatorname{pos\text{-}rr}(c) \wedge \\ \quad \quad \quad \quad \forall x \in L({v}, {v}), t \in \mathbb{R}. \\ \quad \quad \quad \quad \quad ((0 \leq t) \wedge \forall i \in \mathit{sbf}.\; (\operatorname{magnitude}(\langle\cdot,\cdot\rangle_{v}(x(\operatorname{nth}(1, i)), \operatorname{nth}(2, i))) \leq t)) \Rightarrow  \\ \quad \quad \quad \quad \quad \quad \operatorname{magnitude}(f(x)) \leq (c \cdot t) \end{array}}}

\vnbbookfile{structure-library/ba-vn-vocab}{structure-library/ba-vn-vocab}
\noindent{\small ba-vn-vocab.scm -- Thayer, Banach Algebras and Applications, chapter 6 (density theorems): the commutant and bicommutant of a set of operators, involutive operator algebras, von Neumann algebras, cyclic and separating vectors.  DEFINITIONS ONLY: no proof, no support, no axiom. Agent BA-1111-C, 2026-11-11.  DESTINATION: ...}\par
\par\medskip\noindent\textbf{Definition} (\texttt{commutant}).\;${\operatorname{vns\_}}' :=$
\begin{equation*}
\{\operatorname{vnt\_} \in L({\operatorname{vnv\_}}, {\operatorname{vnv\_}}) : \forall \operatorname{vna\_} \in \operatorname{vns\_}.\; ({\operatorname{vnt\_}}{\operatorname{vna\_}} = {\operatorname{vna\_}}{\operatorname{vnt\_}})\}
\end{equation*}

\par\medskip\noindent\textbf{Definition} (\texttt{bicommutant}).\;${\operatorname{vns\_}}'' :=$ ${{\operatorname{vns\_}}'}'$
\par\medskip\noindent\textbf{Definition} (\texttt{is-star-op-algebra}).\;For all $\operatorname{vnv\_}$ and $\operatorname{vns\_}$, 
\begin{equation*}
\begin{array}{@{}l@{}} \operatorname{is\text{-}star\text{-}op\text{-}algebra}(\operatorname{vnv\_}, \operatorname{vns\_}) \\ \Leftrightarrow\; \operatorname{vns\_} \subseteq L({\operatorname{vnv\_}}, {\operatorname{vnv\_}}) \wedge \\ 1 \in \operatorname{vns\_} \wedge \\ \forall \operatorname{vna\_}, \operatorname{vnb\_} \in \operatorname{vns\_}.\; ({\operatorname{vna\_}} + {\operatorname{vnb\_}} \in \operatorname{vns\_}) \wedge \\ \forall \operatorname{vnc\_} \in \mathbb{C}, \operatorname{vna\_} \in \operatorname{vns\_}.\; ({\operatorname{vnc\_}}\,{\operatorname{vna\_}} \in \operatorname{vns\_}) \wedge \\ \forall \operatorname{vna\_}, \operatorname{vnb\_} \in \operatorname{vns\_}.\; ({\operatorname{vna\_}}{\operatorname{vnb\_}} \in \operatorname{vns\_}) \wedge \\ \forall \operatorname{vna\_} \in \operatorname{vns\_}.\; ({\operatorname{vna\_}}^{*} \in \operatorname{vns\_}) \end{array}
\end{equation*}

\par\medskip\noindent\textbf{Definition} (\texttt{is-von-neumann-algebra}).\;For all $\operatorname{vnv\_}$ and $\operatorname{vns\_}$, 
\begin{equation*}
\begin{array}{@{}l@{}} \operatorname{is\text{-}von\text{-}neumann\text{-}algebra}(\operatorname{vnv\_}, \operatorname{vns\_}) \\ \Leftrightarrow\; \operatorname{is\text{-}star\text{-}op\text{-}algebra}(\operatorname{vnv\_}, \operatorname{vns\_}) \wedge \operatorname{is\text{-}top\text{-}closed}(L({\operatorname{vnv\_}})_{s}, \operatorname{vns\_}) \end{array}
\end{equation*}

\par\medskip\noindent\textbf{Definition} (\texttt{is-cyclic-for}).\;For all $\operatorname{vnv\_}$, $\operatorname{vns\_}$ and $\operatorname{vnx\_}$, 
\begin{equation*}
\begin{array}{@{}l@{}} \operatorname{is\text{-}cyclic\text{-}for}(\operatorname{vnv\_}, \operatorname{vns\_}, \operatorname{vnx\_}) \\ \Leftrightarrow\; \operatorname{vnx\_} \in \operatorname{vec}(\operatorname{vnv\_}) \wedge \\ \forall \operatorname{vny\_} \in \operatorname{vec}(\operatorname{vnv\_}), \operatorname{vne\_}. \\ \quad \operatorname{pos\text{-}rr}(\operatorname{vne\_}) \Rightarrow  \\ \quad \quad \exists \operatorname{vna\_} \in \operatorname{vns\_}. \\ \quad \quad \quad \lVert \operatorname{vadd}(\operatorname{vnv\_})(\operatorname{vny\_}, \operatorname{vneg}(\operatorname{vnv\_})(\operatorname{vna\_}(\operatorname{vnx\_}))) \rVert < \operatorname{vne\_} \end{array}
\end{equation*}

\par\medskip\noindent\textbf{Definition} (\texttt{is-separating-for}).\;For all $\operatorname{vnv\_}$, $\operatorname{vns\_}$ and $\operatorname{vnx\_}$, 
\begin{equation*}
\begin{array}{@{}l@{}} \operatorname{is\text{-}separating\text{-}for}(\operatorname{vnv\_}, \operatorname{vns\_}, \operatorname{vnx\_}) \\ \Leftrightarrow\; \operatorname{vnx\_} \in \operatorname{vec}(\operatorname{vnv\_}) \wedge \\ \forall \operatorname{vna\_} \in \operatorname{vns\_}. \\ \quad (\operatorname{vna\_}(\operatorname{vnx\_}) = \operatorname{vzero}(\operatorname{vnv\_})) \Rightarrow (\operatorname{vna\_} = 0) \end{array}
\end{equation*}


\vnbbookfile{theorem-library/ba-vn}{theorem-library/ba-vn}
\noindent{\small ba-vn.scm -- Thayer, Banach Algebras and Applications, chapter 6: commutants and bicommutants, the strong closedness of a commutant (Thm 6.1, first half), the one-vector core of the bicommutant theorem (T xi in the closure of A xi for T in A''), the commutant of a self-adjoint set as a von Neumann algebra, Prop 6.8 ...}\par
\vnbprop{commutant-unfold}{\vnbstmt{\begin{array}{@{}l@{}} \forall v, s. \\ \quad ({s}' \simeq \{t \in L({v}, {v}) : \forall a \in s.\; ({t}{a} = {a}{t})\}) \end{array}}}
\vnbprop{bicommutant-unfold}{\vnbstmt{\forall v, s.\; ({s}'' \simeq {{s}'}')}}
\vnbprop{commutant-membership}{\vnbstmt{\begin{array}{@{}l@{}} \forall v, s, t. \\ \quad \operatorname{is\text{-}complex\text{-}inner\text{-}product\text{-}space}(v) \Rightarrow  \\ \quad \quad t \in {s}' \\ \quad \quad \Leftrightarrow\; ((t \in L({v}, {v})) \wedge \forall a \in s.\; ({t}{a} = {a}{t})) \end{array}}}
\vnbprop{subset-bicommutant}{\vnbstmt{\begin{array}{@{}l@{}} \forall v, s. \\ \quad (\operatorname{is\text{-}complex\text{-}inner\text{-}product\text{-}space}(v) \wedge (s \subseteq L({v}, {v}))) \Rightarrow  \\ \quad \quad s \subseteq {s}'' \end{array}}}
\vnbprop{cyclic-commutant-separating}{\vnbstmt{\begin{array}{@{}l@{}} \forall v, s, i. \\ \quad \operatorname{is\text{-}complex\text{-}inner\text{-}product\text{-}space}(v) \wedge \\ \quad s \subseteq L({v}, {v}) \wedge \\ \quad \operatorname{is\text{-}cyclic\text{-}for}(v, {s}', i) \Rightarrow \\ \quad \quad \operatorname{is\text{-}separating\text{-}for}(v, s, i) \end{array}}}
\vnbprop{cc-conjugate-zero}{\vnbstmt{\operatorname{conjugate}(0) = 0}}
\vnbprop{separating-commutant-cyclic}{\vnbstmt{\begin{array}{@{}l@{}} \forall v, s, i. \\ \quad (\operatorname{is\text{-}hilbert\text{-}space}(v) \wedge (\operatorname{is\text{-}star\text{-}op\text{-}algebra}(v, s) \wedge \operatorname{is\text{-}separating\text{-}for}(v, {s}', i))) \Rightarrow  \\ \quad \quad \operatorname{is\text{-}cyclic\text{-}for}(v, s, i) \end{array}}}
\vnbprop{commutant-subset-bo}{\vnbstmt{\begin{array}{@{}l@{}} \forall v, s. \\ \quad \operatorname{is\text{-}complex\text{-}inner\text{-}product\text{-}space}(v) \Rightarrow  \\ \quad \quad {s}' \subseteq L({v}, {v}) \end{array}}}
\vnbprop{cyclic-superset}{\vnbstmt{\begin{array}{@{}l@{}} \forall v, s, t, i. \\ \quad ((s \subseteq t) \wedge \operatorname{is\text{-}cyclic\text{-}for}(v, s, i)) \Rightarrow \operatorname{is\text{-}cyclic\text{-}for}(v, t, i) \end{array}}}
\vnbprop{separating-subset}{\vnbstmt{\begin{array}{@{}l@{}} \forall v, s, t, i. \\ \quad ((s \subseteq t) \wedge \operatorname{is\text{-}separating\text{-}for}(v, t, i)) \Rightarrow \operatorname{is\text{-}separating\text{-}for}(v, s, i) \end{array}}}
\vnbprop{cyclic-iff-commutant-separating}{\vnbstmt{\begin{array}{@{}l@{}} \forall v, s, i. \\ \quad (\operatorname{is\text{-}hilbert\text{-}space}(v) \wedge \operatorname{is\text{-}star\text{-}op\text{-}algebra}(v, s)) \Rightarrow  \\ \quad \quad \operatorname{is\text{-}cyclic\text{-}for}(v, s, i) \Leftrightarrow \operatorname{is\text{-}separating\text{-}for}(v, {s}', i) \end{array}}}
\vnbprop{commutant-star-algebra}{\vnbstmt{\begin{array}{@{}l@{}} \forall v, s. \\ \quad (\operatorname{is\text{-}hilbert\text{-}space}(v) \wedge ((s \subseteq L({v}, {v})) \wedge \forall a \in s.\; ({a}^{*} \in s))) \Rightarrow  \\ \quad \quad \operatorname{is\text{-}star\text{-}op\text{-}algebra}(v, {s}') \end{array}}}
\vnbprop{separating-iff-commutant-cyclic}{\vnbstmt{\begin{array}{@{}l@{}} \forall v, s, i. \\ \quad (\operatorname{is\text{-}hilbert\text{-}space}(v) \wedge (\operatorname{is\text{-}star\text{-}op\text{-}algebra}(v, s) \wedge ({s}'' = s))) \Rightarrow  \\ \quad \quad \operatorname{is\text{-}separating\text{-}for}(v, s, i) \Leftrightarrow \operatorname{is\text{-}cyclic\text{-}for}(v, {s}', i) \end{array}}}
\vnbprop{commutant-sot-closed}{\vnbstmt{\begin{array}{@{}l@{}} \forall v, s. \\ \quad (\operatorname{is\text{-}complex\text{-}inner\text{-}product\text{-}space}(v) \wedge (s \subseteq L({v}, {v}))) \Rightarrow  \\ \quad \quad \operatorname{is\text{-}top\text{-}closed}(L({v})_{s}, {s}') \end{array}}}
\vnbprop{commutant-von-neumann}{\vnbstmt{\begin{array}{@{}l@{}} \forall v, s. \\ \quad (\operatorname{is\text{-}hilbert\text{-}space}(v) \wedge ((s \subseteq L({v}, {v})) \wedge \forall a \in s.\; ({a}^{*} \in s))) \Rightarrow  \\ \quad \quad \operatorname{is\text{-}von\text{-}neumann\text{-}algebra}(v, {s}') \end{array}}}
\vnbprop{bicommutant-von-neumann}{\vnbstmt{\begin{array}{@{}l@{}} \forall v, s. \\ \quad (\operatorname{is\text{-}hilbert\text{-}space}(v) \wedge ((s \subseteq L({v}, {v})) \wedge \forall a \in s.\; ({a}^{*} \in s))) \Rightarrow  \\ \quad \quad \operatorname{is\text{-}von\text{-}neumann\text{-}algebra}(v, {s}'') \end{array}}}
\vnbprop{von-neumann-of-bicommutant-eq}{\vnbstmt{\begin{array}{@{}l@{}} \forall v, s. \\ \quad (\operatorname{is\text{-}hilbert\text{-}space}(v) \wedge (\operatorname{is\text{-}star\text{-}op\text{-}algebra}(v, s) \wedge ({s}'' = s))) \Rightarrow  \\ \quad \quad \operatorname{is\text{-}von\text{-}neumann\text{-}algebra}(v, s) \end{array}}}
\vnbprop{orbit-projection-commutant}{\vnbstmt{\begin{array}{@{}l@{}} \forall v, s, i. \\ \quad (\operatorname{is\text{-}hilbert\text{-}space}(v) \wedge (\operatorname{is\text{-}star\text{-}op\text{-}algebra}(v, s) \wedge (i \in \operatorname{vec}(v)))) \Rightarrow  \\ \quad \quad \operatorname{is\text{-}submodule}(v, \{m \in \operatorname{vec}(v) : \exists b \in s.\; (m = b(i))\}) \wedge \\ \quad \quad E_{{\{m \in \operatorname{vec}(v) : \exists b \in s.\; (m = b(i))\}}^{\perp}} \in {s}' \wedge \\ \quad \quad E_{{\{m \in \operatorname{vec}(v) : \exists b \in s.\; (m = b(i))\}}^{\perp}}(i) = \operatorname{vzero}(v) \end{array}}}
\vnbprop{bicommutant-vector-closure}{\vnbstmt{\begin{array}{@{}l@{}} \forall v, s, i, t. \\ \quad (\operatorname{is\text{-}hilbert\text{-}space}(v) \wedge (\operatorname{is\text{-}star\text{-}op\text{-}algebra}(v, s) \wedge ((i \in \operatorname{vec}(v)) \wedge (t \in {s}'')))) \Rightarrow  \\ \quad \quad t(i) \in \operatorname{closure}(\operatorname{ip\text{-}metric\text{-}space}(v), \{m \in \operatorname{vec}(v) : \exists b \in s.\; (m = b(i))\}) \end{array}}}

\vnbbookfile{theorem-library/ba-forms}{theorem-library/ba-forms}
\noindent{\small ba-forms.scm -- BanachAlgebras.pdf ch. 3, sec. 4-5: closed quadratic forms, T*T, Friedrichs. 3.16 form-closed-iff-hilbert: a positive form Q is closed iff H\_Q+ is a Hilbert space. 3.17 form-operator-self-adjoint / -positive / -functional / -dom-dense: the operator T\_Q of a closed positive form is a densely defined ...}\par
\vnbprop{form-plus-space-vec}{\vnbstmt{\begin{array}{@{}l@{}} \forall v, d, q. \\ \quad (\operatorname{is\text{-}complex\text{-}inner\text{-}product\text{-}space}(v) \wedge (\operatorname{is\text{-}submodule}(v, d) \wedge \operatorname{is\text{-}positive\text{-}form}(v, d, q))) \Rightarrow  \\ \quad \quad (\operatorname{vec}(H_{{q}+}) \simeq d) \end{array}}}
\vnbprop{form-plus-space-ip}{\vnbstmt{\begin{array}{@{}l@{}} \forall v, d, q. \\ \quad (\operatorname{is\text{-}complex\text{-}inner\text{-}product\text{-}space}(v) \wedge (\operatorname{is\text{-}submodule}(v, d) \wedge \operatorname{is\text{-}positive\text{-}form}(v, d, q))) \Rightarrow  \\ \quad \quad (\langle\cdot,\cdot\rangle_{H_{{q}+}} \simeq {q}_{+}) \end{array}}}
\vnbprop{form-plus-space-vadd}{\vnbstmt{\begin{array}{@{}l@{}} \forall v, d, q. \\ \quad (\operatorname{is\text{-}complex\text{-}inner\text{-}product\text{-}space}(v) \wedge (\operatorname{is\text{-}submodule}(v, d) \wedge \operatorname{is\text{-}positive\text{-}form}(v, d, q))) \Rightarrow  \\ \quad \quad (\operatorname{vadd}(H_{{q}+}) \simeq (\lambda {[}\operatorname{tsbx\_}, \operatorname{tsby\_}] \in (d \times d).\; \operatorname{vadd}(v)(\operatorname{tsbx\_}, \operatorname{tsby\_}))) \end{array}}}
\vnbprop{form-plus-space-vneg}{\vnbstmt{\begin{array}{@{}l@{}} \forall v, d, q. \\ \quad (\operatorname{is\text{-}complex\text{-}inner\text{-}product\text{-}space}(v) \wedge (\operatorname{is\text{-}submodule}(v, d) \wedge \operatorname{is\text{-}positive\text{-}form}(v, d, q))) \Rightarrow  \\ \quad \quad (\operatorname{vneg}(H_{{q}+}) \simeq (\lambda x \in d.\; \operatorname{vneg}(v)(x))) \end{array}}}
\vnbprop{form-plus-space-ip-apply}{\vnbstmt{\begin{array}{@{}l@{}} \forall v, d, q, x, y. \\ \quad \operatorname{is\text{-}complex\text{-}inner\text{-}product\text{-}space}(v) \wedge \\ \quad \operatorname{is\text{-}submodule}(v, d) \wedge \\ \quad \operatorname{is\text{-}positive\text{-}form}(v, d, q) \wedge \\ \quad x \in d \wedge \\ \quad y \in d \Rightarrow \\ \quad \quad \langle\cdot,\cdot\rangle_{H_{{q}+}}(x, y) = (q(x, y) + \langle\cdot,\cdot\rangle_{v}(x, y)) \end{array}}}
\vnbprop{form-plus-space-vadd-apply}{\vnbstmt{\begin{array}{@{}l@{}} \forall v, d, q, x, y. \\ \quad \operatorname{is\text{-}complex\text{-}inner\text{-}product\text{-}space}(v) \wedge \\ \quad \operatorname{is\text{-}submodule}(v, d) \wedge \\ \quad \operatorname{is\text{-}positive\text{-}form}(v, d, q) \wedge \\ \quad x \in d \wedge \\ \quad y \in d \Rightarrow \\ \quad \quad \operatorname{vadd}(H_{{q}+})(x, y) = \operatorname{vadd}(v)(x, y) \end{array}}}
\vnbprop{form-plus-space-vneg-apply}{\vnbstmt{\begin{array}{@{}l@{}} \forall v, d, q, x. \\ \quad (\operatorname{is\text{-}complex\text{-}inner\text{-}product\text{-}space}(v) \wedge (\operatorname{is\text{-}submodule}(v, d) \wedge (\operatorname{is\text{-}positive\text{-}form}(v, d, q) \wedge (x \in d)))) \Rightarrow  \\ \quad \quad \operatorname{vneg}(H_{{q}+})(x) = \operatorname{vneg}(v)(x) \end{array}}}
\vnbprop{form-plus-space-dist}{\vnbstmt{\begin{array}{@{}l@{}} \forall v, d, q, x, y. \\ \quad \operatorname{is\text{-}complex\text{-}inner\text{-}product\text{-}space}(v) \wedge \\ \quad \operatorname{is\text{-}submodule}(v, d) \wedge \\ \quad \operatorname{is\text{-}positive\text{-}form}(v, d, q) \wedge \\ \quad x \in d \wedge \\ \quad y \in d \Rightarrow \\ \quad \quad \operatorname{dist}(\operatorname{ip\text{-}metric\text{-}space}(H_{{q}+}))(x, y) \\ \quad \quad =\; \sqrt{(q(\operatorname{vadd}(v)(x, \operatorname{vneg}(v)(y)), \operatorname{vadd}(v)(x, \operatorname{vneg}(v)(y))) + \langle\cdot,\cdot\rangle_{v}(\operatorname{vadd}(v)(x, \operatorname{vneg}(v)(y)), \operatorname{vadd}(v)(x, \operatorname{vneg}(v)(y))))} \end{array}}}
\vnbprop{form-plus-space-dist-ge}{\vnbstmt{\begin{array}{@{}l@{}} \forall v, d, q, x, y. \\ \quad \operatorname{is\text{-}complex\text{-}inner\text{-}product\text{-}space}(v) \wedge \\ \quad \operatorname{is\text{-}submodule}(v, d) \wedge \\ \quad \operatorname{is\text{-}positive\text{-}form}(v, d, q) \wedge \\ \quad x \in d \wedge \\ \quad y \in d \Rightarrow \\ \quad \quad \operatorname{dist}(\operatorname{ip\text{-}metric\text{-}space}(v))(x, y) \leq \operatorname{dist}(\operatorname{ip\text{-}metric\text{-}space}(H_{{q}+}))(x, y) \end{array}}}
\vnbprop{form-plus-space-form-le}{\vnbstmt{\begin{array}{@{}l@{}} \forall v, d, q, x, y. \\ \quad \operatorname{is\text{-}complex\text{-}inner\text{-}product\text{-}space}(v) \wedge \\ \quad \operatorname{is\text{-}submodule}(v, d) \wedge \\ \quad \operatorname{is\text{-}positive\text{-}form}(v, d, q) \wedge \\ \quad x \in d \wedge \\ \quad y \in d \Rightarrow \\ \quad \quad \sqrt{q(\operatorname{vadd}(v)(x, \operatorname{vneg}(v)(y)), \operatorname{vadd}(v)(x, \operatorname{vneg}(v)(y)))} \\ \quad \quad \leq\; \operatorname{dist}(\operatorname{ip\text{-}metric\text{-}space}(H_{{q}+}))(x, y) \end{array}}}
\vnbprop{form-plus-space-dist-le}{\vnbstmt{\begin{array}{@{}l@{}} \forall v, d, q, x, y. \\ \quad \operatorname{is\text{-}complex\text{-}inner\text{-}product\text{-}space}(v) \wedge \\ \quad \operatorname{is\text{-}submodule}(v, d) \wedge \\ \quad \operatorname{is\text{-}positive\text{-}form}(v, d, q) \wedge \\ \quad x \in d \wedge \\ \quad y \in d \Rightarrow \\ \quad \quad \operatorname{dist}(\operatorname{ip\text{-}metric\text{-}space}(H_{{q}+}))(x, y) \\ \quad \quad \leq\; \sqrt{q(\operatorname{vadd}(v)(x, \operatorname{vneg}(v)(y)), \operatorname{vadd}(v)(x, \operatorname{vneg}(v)(y)))} + \operatorname{dist}(\operatorname{ip\text{-}metric\text{-}space}(v))(x, y) \end{array}}}
\vnbprop{rr-sqrt-pos-rr}{\vnbstmt{\forall e.\; (\operatorname{pos\text{-}rr}(e) \Rightarrow \operatorname{pos\text{-}rr}(\sqrt{e}))}}
\vnbprop{form-plus-space-small-form}{\vnbstmt{\begin{array}{@{}l@{}} \forall v, d, q, x, y, h. \\ \quad \operatorname{is\text{-}complex\text{-}inner\text{-}product\text{-}space}(v) \wedge \\ \quad \operatorname{is\text{-}submodule}(v, d) \wedge \\ \quad \operatorname{is\text{-}positive\text{-}form}(v, d, q) \wedge \\ \quad x \in d \wedge \\ \quad y \in d \wedge \\ \quad h \in \mathbb{R} \wedge \\ \quad 0 \leq h \wedge \\ \quad \operatorname{dist}(\operatorname{ip\text{-}metric\text{-}space}(H_{{q}+}))(x, y) \leq h \Rightarrow \\ \quad \quad \operatorname{magnitude}(q(\operatorname{vadd}(v)(x, \operatorname{vneg}(v)(y)), \operatorname{vadd}(v)(x, \operatorname{vneg}(v)(y)))) \leq (h \cdot h) \end{array}}}
\vnbprop{form-plus-space-small-dist}{\vnbstmt{\begin{array}{@{}l@{}} \forall v, d, q, x, y, h. \\ \quad \operatorname{is\text{-}complex\text{-}inner\text{-}product\text{-}space}(v) \wedge \\ \quad \operatorname{is\text{-}submodule}(v, d) \wedge \\ \quad \operatorname{is\text{-}positive\text{-}form}(v, d, q) \wedge \\ \quad x \in d \wedge \\ \quad y \in d \wedge \\ \quad h \in \mathbb{R} \wedge \\ \quad 0 \leq h \wedge \\ \quad \operatorname{magnitude}(q(\operatorname{vadd}(v)(x, \operatorname{vneg}(v)(y)), \operatorname{vadd}(v)(x, \operatorname{vneg}(v)(y)))) \leq (h \cdot h) \wedge \\ \quad \operatorname{dist}(\operatorname{ip\text{-}metric\text{-}space}(v))(x, y) \leq h \Rightarrow \\ \quad \quad \operatorname{dist}(\operatorname{ip\text{-}metric\text{-}space}(H_{{q}+}))(x, y) \leq (h + h) \end{array}}}
\vnbprop{form-closed-plus-cauchy-converges}{\vnbstmt{\begin{array}{@{}l@{}} \forall v, d, q, s. \\ \quad \operatorname{is\text{-}hilbert\text{-}space}(v) \wedge \\ \quad \operatorname{is\text{-}closed\text{-}form}(v, d, q) \wedge \\ \quad \operatorname{is\text{-}cauchy\text{-}seq}(\operatorname{ip\text{-}metric\text{-}space}(H_{{q}+}), s) \Rightarrow \\ \quad \quad \operatorname{converges}(\operatorname{ip\text{-}metric\text{-}space}(H_{{q}+}), s) \end{array}}}
\vnbprop{form-plus-hilbert-closed}{\vnbstmt{\begin{array}{@{}l@{}} \forall v, d, q. \\ \quad (\operatorname{is\text{-}hilbert\text{-}space}(v) \wedge (\operatorname{is\text{-}positive\text{-}form}(v, d, q) \wedge \operatorname{is\text{-}hilbert\text{-}space}(H_{{q}+}))) \Rightarrow  \\ \quad \quad \operatorname{is\text{-}closed\text{-}form}(v, d, q) \end{array}}}
\vnbprop{form-closed-iff-hilbert}{\vnbstmt{\begin{array}{@{}l@{}} \forall v, d, q. \\ \quad (\operatorname{is\text{-}hilbert\text{-}space}(v) \wedge \operatorname{is\text{-}positive\text{-}form}(v, d, q)) \Rightarrow  \\ \quad \quad \operatorname{is\text{-}closed\text{-}form}(v, d, q) \Leftrightarrow \operatorname{is\text{-}hilbert\text{-}space}(H_{{q}+}) \end{array}}}
\vnbprop{form-plus-space-act}{\vnbstmt{\begin{array}{@{}l@{}} \forall v, d, q. \\ \quad (\operatorname{is\text{-}complex\text{-}inner\text{-}product\text{-}space}(v) \wedge (\operatorname{is\text{-}submodule}(v, d) \wedge \operatorname{is\text{-}positive\text{-}form}(v, d, q))) \Rightarrow  \\ \quad \quad (\operatorname{act}(H_{{q}+}) \simeq (\lambda {[}\operatorname{tsbc\_}, \operatorname{tsbx\_}] \in (\mathbb{C} \times d).\; \operatorname{act}(v)(\operatorname{tsbc\_}, \operatorname{tsbx\_}))) \end{array}}}
\vnbprop{form-plus-space-act-apply}{\vnbstmt{\begin{array}{@{}l@{}} \forall v, d, q, c, x. \\ \quad \operatorname{is\text{-}complex\text{-}inner\text{-}product\text{-}space}(v) \wedge \\ \quad \operatorname{is\text{-}submodule}(v, d) \wedge \\ \quad \operatorname{is\text{-}positive\text{-}form}(v, d, q) \wedge \\ \quad c \in \mathbb{C} \wedge \\ \quad x \in d \Rightarrow \\ \quad \quad \operatorname{act}(H_{{q}+})(c, x) = \operatorname{act}(v)(c, x) \end{array}}}
\vnbprop{form-plus-space-norm-ge}{\vnbstmt{\begin{array}{@{}l@{}} \forall v, d, q, x. \\ \quad (\operatorname{is\text{-}complex\text{-}inner\text{-}product\text{-}space}(v) \wedge (\operatorname{is\text{-}submodule}(v, d) \wedge (\operatorname{is\text{-}positive\text{-}form}(v, d, q) \wedge (x \in d)))) \Rightarrow  \\ \quad \quad \lVert x \rVert \leq \lVert x \rVert \end{array}}}
\vnbprop{form-operator-symmetric}{\vnbstmt{\begin{array}{@{}l@{}} \forall v, d, q. \\ \quad (\operatorname{is\text{-}complex\text{-}inner\text{-}product\text{-}space}(v) \wedge (\operatorname{is\text{-}submodule}(v, d) \wedge \operatorname{is\text{-}positive\text{-}form}(v, d, q))) \Rightarrow  \\ \quad \quad \operatorname{is\text{-}symmetric\text{-}graph}(v, T_{q}) \end{array}}}
\vnbprop{form-operator-positive}{\vnbstmt{\begin{array}{@{}l@{}} \forall v, d, q, c, \mathit{bfd}. \\ \quad \operatorname{is\text{-}complex\text{-}inner\text{-}product\text{-}space}(v) \wedge \\ \quad \operatorname{is\text{-}submodule}(v, d) \wedge \\ \quad \operatorname{is\text{-}positive\text{-}form}(v, d, q) \wedge \\ \quad c \in \operatorname{vec}(v) \wedge \\ \quad \mathit{bfd} \in \operatorname{vec}(v) \wedge \\ \quad {[}c, \mathit{bfd}] \in T_{q} \Rightarrow \\ \quad \quad (\langle\cdot,\cdot\rangle_{v}(c, \mathit{bfd}) \in \mathbb{R}) \wedge (0 \leq \langle\cdot,\cdot\rangle_{v}(c, \mathit{bfd})) \end{array}}}
\vnbprop{cc-sub-cancel-eq}{\vnbstmt{\begin{array}{@{}l@{}} \forall a, p, c, r, e \in \mathbb{C}. \\ \quad ((a = (p + -c)) \wedge (((a + c) + -e) = (r + -e))) \Rightarrow (p = r) \end{array}}}
\vnbprop{form-operator-onto-plus}{\vnbstmt{\begin{array}{@{}l@{}} \forall v, d, q, p. \\ \quad (\operatorname{is\text{-}hilbert\text{-}space}(v) \wedge (\operatorname{is\text{-}closed\text{-}form}(v, d, q) \wedge (p \in \operatorname{vec}(v)))) \Rightarrow  \\ \quad \quad \exists e \in d.\; ({[}e, \operatorname{vadd}(v)(p, \operatorname{vneg}(v)(e))] \in T_{q}) \end{array}}}
\vnbprop{form-operator-self-adjoint}{\vnbstmt{\begin{array}{@{}l@{}} \forall v, d, q. \\ \quad (\operatorname{is\text{-}hilbert\text{-}space}(v) \wedge \operatorname{is\text{-}closed\text{-}form}(v, d, q)) \Rightarrow  \\ \quad \quad \operatorname{is\text{-}self\text{-}adjoint\text{-}graph}(v, T_{q}) \end{array}}}
\vnbprop{graph-form-value}{\vnbstmt{\begin{array}{@{}l@{}} \forall v, d, f, x, y. \\ \quad \operatorname{is\text{-}complex\text{-}inner\text{-}product\text{-}space}(v) \wedge \\ \quad \operatorname{is\text{-}submodule}(v, d) \wedge \\ \quad f \in (d \to \operatorname{vec}(v)) \wedge \\ \quad \forall \mathit{bax}, \mathit{bay} \in d.\; (f(\operatorname{vadd}(v)(\mathit{bax}, \mathit{bay})) = \operatorname{vadd}(v)(f(\mathit{bax}), f(\mathit{bay}))) \wedge \\ \quad \forall c \in \mathbb{C}, \mathit{bax} \in d.\; (f(\operatorname{act}(v)(c, \mathit{bax})) = \operatorname{act}(v)(c, f(\mathit{bax}))) \wedge \\ \quad x \in d \wedge \\ \quad y \in d \Rightarrow \\ \quad \quad (\lambda {[}\operatorname{bgx\_}, \operatorname{bgy\_}] \in (d \times d).\; \langle\cdot,\cdot\rangle_{v}(f(\operatorname{bgx\_}), f(\operatorname{bgy\_})))(x, y) = \langle\cdot,\cdot\rangle_{v}(f(x), f(y)) \end{array}}}
\vnbprop{graph-map-vneg}{\vnbstmt{\begin{array}{@{}l@{}} \forall v, d, f, y. \\ \quad \operatorname{is\text{-}complex\text{-}inner\text{-}product\text{-}space}(v) \wedge \\ \quad \operatorname{is\text{-}submodule}(v, d) \wedge \\ \quad f \in (d \to \operatorname{vec}(v)) \wedge \\ \quad \forall x, \mathit{bay} \in d.\; (f(\operatorname{vadd}(v)(x, \mathit{bay})) = \operatorname{vadd}(v)(f(x), f(\mathit{bay}))) \wedge \\ \quad \forall c \in \mathbb{C}, x \in d.\; (f(\operatorname{act}(v)(c, x)) = \operatorname{act}(v)(c, f(x))) \wedge \\ \quad y \in d \Rightarrow \\ \quad \quad f(\operatorname{vneg}(v)(y)) = \operatorname{vneg}(v)(f(y)) \end{array}}}
\vnbprop{graph-map-sub}{\vnbstmt{\begin{array}{@{}l@{}} \forall v, d, f, x, y. \\ \quad \operatorname{is\text{-}complex\text{-}inner\text{-}product\text{-}space}(v) \wedge \\ \quad \operatorname{is\text{-}submodule}(v, d) \wedge \\ \quad f \in (d \to \operatorname{vec}(v)) \wedge \\ \quad \forall \mathit{bax}, \mathit{bay} \in d.\; (f(\operatorname{vadd}(v)(\mathit{bax}, \mathit{bay})) = \operatorname{vadd}(v)(f(\mathit{bax}), f(\mathit{bay}))) \wedge \\ \quad \forall c \in \mathbb{C}, \mathit{bax} \in d.\; (f(\operatorname{act}(v)(c, \mathit{bax})) = \operatorname{act}(v)(c, f(\mathit{bax}))) \wedge \\ \quad x \in d \wedge \\ \quad y \in d \Rightarrow \\ \quad \quad f(\operatorname{vadd}(v)(x, \operatorname{vneg}(v)(y))) = \operatorname{vadd}(v)(f(x), \operatorname{vneg}(v)(f(y))) \end{array}}}
\vnbprop{graph-form-sub-value}{\vnbstmt{\begin{array}{@{}l@{}} \forall v, d, f, x, y. \\ \quad \operatorname{is\text{-}complex\text{-}inner\text{-}product\text{-}space}(v) \wedge \\ \quad \operatorname{is\text{-}submodule}(v, d) \wedge \\ \quad f \in (d \to \operatorname{vec}(v)) \wedge \\ \quad \forall \mathit{bax}, \mathit{bay} \in d.\; (f(\operatorname{vadd}(v)(\mathit{bax}, \mathit{bay})) = \operatorname{vadd}(v)(f(\mathit{bax}), f(\mathit{bay}))) \wedge \\ \quad \forall c \in \mathbb{C}, \mathit{bax} \in d.\; (f(\operatorname{act}(v)(c, \mathit{bax})) = \operatorname{act}(v)(c, f(\mathit{bax}))) \wedge \\ \quad x \in d \wedge \\ \quad y \in d \Rightarrow \\ \quad \quad (\lambda {[}\operatorname{bgx\_}, \operatorname{bgy\_}] \in (d \times d).\; \langle\cdot,\cdot\rangle_{v}(f(\operatorname{bgx\_}), f(\operatorname{bgy\_})))(\operatorname{vadd}(v)(x, \operatorname{vneg}(v)(y)), \operatorname{vadd}(v)(x, \operatorname{vneg}(v)(y))) \\ \quad \quad =\; \langle\cdot,\cdot\rangle_{v}(\operatorname{vadd}(v)(f(x), \operatorname{vneg}(v)(f(y))), \operatorname{vadd}(v)(f(x), \operatorname{vneg}(v)(f(y)))) \end{array}}}
\vnbprop{cips-dist-le-of-sq-le}{\vnbstmt{\begin{array}{@{}l@{}} \forall v, x, y, h. \\ \quad \operatorname{is\text{-}complex\text{-}inner\text{-}product\text{-}space}(v) \wedge \\ \quad x \in \operatorname{vec}(v) \wedge \\ \quad y \in \operatorname{vec}(v) \wedge \\ \quad h \in \mathbb{R} \wedge \\ \quad 0 \leq h \wedge \\ \quad \operatorname{magnitude}(\langle\cdot,\cdot\rangle_{v}(\operatorname{vadd}(v)(x, \operatorname{vneg}(v)(y)), \operatorname{vadd}(v)(x, \operatorname{vneg}(v)(y)))) \leq (h \cdot h) \Rightarrow \\ \quad \quad \operatorname{dist}(\operatorname{ip\text{-}metric\text{-}space}(v))(x, y) \leq h \end{array}}}
\vnbprop{cips-sq-le-of-dist-le}{\vnbstmt{\begin{array}{@{}l@{}} \forall v, x, y, h. \\ \quad \operatorname{is\text{-}complex\text{-}inner\text{-}product\text{-}space}(v) \wedge \\ \quad x \in \operatorname{vec}(v) \wedge \\ \quad y \in \operatorname{vec}(v) \wedge \\ \quad h \in \mathbb{R} \wedge \\ \quad \operatorname{dist}(\operatorname{ip\text{-}metric\text{-}space}(v))(x, y) \leq h \Rightarrow \\ \quad \quad \operatorname{magnitude}(\langle\cdot,\cdot\rangle_{v}(\operatorname{vadd}(v)(x, \operatorname{vneg}(v)(y)), \operatorname{vadd}(v)(x, \operatorname{vneg}(v)(y)))) \leq (h \cdot h) \end{array}}}
\vnbprop{graph-form-positive}{\vnbstmt{\begin{array}{@{}l@{}} \forall v, d, f. \\ \quad \operatorname{is\text{-}complex\text{-}inner\text{-}product\text{-}space}(v) \wedge \\ \quad \operatorname{is\text{-}submodule}(v, d) \wedge \\ \quad f \in (d \to \operatorname{vec}(v)) \wedge \\ \quad \forall x, y \in d.\; (f(\operatorname{vadd}(v)(x, y)) = \operatorname{vadd}(v)(f(x), f(y))) \wedge \\ \quad \forall c \in \mathbb{C}, x \in d.\; (f(\operatorname{act}(v)(c, x)) = \operatorname{act}(v)(c, f(x))) \wedge \\ \quad \operatorname{is\text{-}dense}(\operatorname{ip\text{-}metric\text{-}space}(v), d) \Rightarrow \\ \quad \quad \operatorname{is\text{-}positive\text{-}form}(v, d, (\lambda {[}\operatorname{bgx\_}, \operatorname{bgy\_}] \in (d \times d).\; \langle\cdot,\cdot\rangle_{v}(f(\operatorname{bgx\_}), f(\operatorname{bgy\_})))) \end{array}}}
\vnbprop{graph-form-closed}{\vnbstmt{\begin{array}{@{}l@{}} \forall v, d, f. \\ \quad \operatorname{is\text{-}hilbert\text{-}space}(v) \wedge \\ \quad \operatorname{is\text{-}submodule}(v, d) \wedge \\ \quad \operatorname{is\text{-}dense}(\operatorname{ip\text{-}metric\text{-}space}(v), d) \wedge \\ \quad f \in (d \to \operatorname{vec}(v)) \wedge \\ \quad \forall x, y \in d.\; (f(\operatorname{vadd}(v)(x, y)) = \operatorname{vadd}(v)(f(x), f(y))) \wedge \\ \quad \forall c \in \mathbb{C}, x \in d.\; (f(\operatorname{act}(v)(c, x)) = \operatorname{act}(v)(c, f(x))) \wedge \\ \quad \operatorname{is\text{-}closed\text{-}graph}(v, \{p \in (\operatorname{vec}(v) \times \operatorname{vec}(v)) : ((\operatorname{nth}(1, p) \in d) \wedge (\operatorname{nth}(2, p) = f(\operatorname{nth}(1, p))))\}) \Rightarrow \\ \quad \quad \operatorname{is\text{-}closed\text{-}form}(v, d, (\lambda {[}\operatorname{bgx\_}, \operatorname{bgy\_}] \in (d \times d).\; \langle\cdot,\cdot\rangle_{v}(f(\operatorname{bgx\_}), f(\operatorname{bgy\_})))) \end{array}}}
\vnbprop{closed-op-adjoint-compose-self-adjoint}{\vnbstmt{\begin{array}{@{}l@{}} \forall v, d, f. \\ \quad \operatorname{is\text{-}hilbert\text{-}space}(v) \wedge \\ \quad \operatorname{is\text{-}submodule}(v, d) \wedge \\ \quad \operatorname{is\text{-}dense}(\operatorname{ip\text{-}metric\text{-}space}(v), d) \wedge \\ \quad f \in (d \to \operatorname{vec}(v)) \wedge \\ \quad \forall x, y \in d.\; (f(\operatorname{vadd}(v)(x, y)) = \operatorname{vadd}(v)(f(x), f(y))) \wedge \\ \quad \forall c \in \mathbb{C}, x \in d.\; (f(\operatorname{act}(v)(c, x)) = \operatorname{act}(v)(c, f(x))) \wedge \\ \quad \operatorname{is\text{-}closed\text{-}graph}(v, \{p \in (\operatorname{vec}(v) \times \operatorname{vec}(v)) : ((\operatorname{nth}(1, p) \in d) \wedge (\operatorname{nth}(2, p) = f(\operatorname{nth}(1, p))))\}) \Rightarrow \\ \quad \quad \operatorname{is\text{-}self\text{-}adjoint\text{-}graph}(v, \\ \quad \quad \quad {{\{p \in (\operatorname{vec}(v) \times \operatorname{vec}(v)) : ((\operatorname{nth}(1, p) \in d) \wedge (\operatorname{nth}(2, p) = f(\operatorname{nth}(1, p))))\}}^{*}}{\{p \in (\operatorname{vec}(v) \times \operatorname{vec}(v)) : ((\operatorname{nth}(1, p) \in d) \wedge (\operatorname{nth}(2, p) = f(\operatorname{nth}(1, p))))\}}) \end{array}}}
\vnbprop{closed-op-adjoint-compose-positive}{\vnbstmt{\begin{array}{@{}l@{}} \forall v, d, f, c, \mathit{bfd}. \\ \quad \operatorname{is\text{-}hilbert\text{-}space}(v) \wedge \\ \quad \operatorname{is\text{-}submodule}(v, d) \wedge \\ \quad \operatorname{is\text{-}dense}(\operatorname{ip\text{-}metric\text{-}space}(v), d) \wedge \\ \quad f \in (d \to \operatorname{vec}(v)) \wedge \\ \quad \forall x, y \in d.\; (f(\operatorname{vadd}(v)(x, y)) = \operatorname{vadd}(v)(f(x), f(y))) \wedge \\ \quad \forall \mathit{bac} \in \mathbb{C}, x \in d.\; (f(\operatorname{act}(v)(\mathit{bac}, x)) = \operatorname{act}(v)(\mathit{bac}, f(x))) \wedge \\ \quad \operatorname{is\text{-}closed\text{-}graph}(v, \{p \in (\operatorname{vec}(v) \times \operatorname{vec}(v)) : ((\operatorname{nth}(1, p) \in d) \wedge (\operatorname{nth}(2, p) = f(\operatorname{nth}(1, p))))\}) \wedge \\ \quad c \in \operatorname{vec}(v) \wedge \\ \quad \mathit{bfd} \in \operatorname{vec}(v) \wedge \\ \quad {[}c, \mathit{bfd}] \\ \quad \in\; {{\{p \in (\operatorname{vec}(v) \times \operatorname{vec}(v)) : ((\operatorname{nth}(1, p) \in d) \wedge (\operatorname{nth}(2, p) = f(\operatorname{nth}(1, p))))\}}^{*}}{\{p \in (\operatorname{vec}(v) \times \operatorname{vec}(v)) : ((\operatorname{nth}(1, p) \in d) \wedge (\operatorname{nth}(2, p) = f(\operatorname{nth}(1, p))))\}} \Rightarrow \\ \quad \quad (\langle\cdot,\cdot\rangle_{v}(c, \mathit{bfd}) \in \mathbb{R}) \wedge (0 \leq \langle\cdot,\cdot\rangle_{v}(c, \mathit{bfd})) \end{array}}}
\vnbprop{cc-add-neg-zero-eq}{\vnbstmt{\forall a, b \in \mathbb{C}.\; (((a + -b) = 0) \Rightarrow (a = b))}}
\vnbprop{form-operator-functional}{\vnbstmt{\begin{array}{@{}l@{}} \forall v, d, q. \\ \quad (\operatorname{is\text{-}hilbert\text{-}space}(v) \wedge \operatorname{is\text{-}positive\text{-}form}(v, d, q)) \Rightarrow  \\ \quad \quad \operatorname{is\text{-}functional\text{-}graph}(v, T_{q}) \end{array}}}
\vnbprop{form-operator-dom-dense}{\vnbstmt{\begin{array}{@{}l@{}} \forall v, d, q. \\ \quad (\operatorname{is\text{-}hilbert\text{-}space}(v) \wedge \operatorname{is\text{-}closed\text{-}form}(v, d, q)) \Rightarrow  \\ \quad \quad \operatorname{is\text{-}dense}(\operatorname{ip\text{-}metric\text{-}space}(v), \mathrm{dom}\,{T_{q}}) \end{array}}}
\vnbprop{closed-op-adjoint-compose-dom-dense}{\vnbstmt{\begin{array}{@{}l@{}} \forall v, d, f. \\ \quad \operatorname{is\text{-}hilbert\text{-}space}(v) \wedge \\ \quad \operatorname{is\text{-}submodule}(v, d) \wedge \\ \quad \operatorname{is\text{-}dense}(\operatorname{ip\text{-}metric\text{-}space}(v), d) \wedge \\ \quad f \in (d \to \operatorname{vec}(v)) \wedge \\ \quad \forall x, y \in d.\; (f(\operatorname{vadd}(v)(x, y)) = \operatorname{vadd}(v)(f(x), f(y))) \wedge \\ \quad \forall c \in \mathbb{C}, x \in d.\; (f(\operatorname{act}(v)(c, x)) = \operatorname{act}(v)(c, f(x))) \wedge \\ \quad \operatorname{is\text{-}closed\text{-}graph}(v, \{p \in (\operatorname{vec}(v) \times \operatorname{vec}(v)) : ((\operatorname{nth}(1, p) \in d) \wedge (\operatorname{nth}(2, p) = f(\operatorname{nth}(1, p))))\}) \Rightarrow \\ \quad \quad \operatorname{is\text{-}dense}(\operatorname{ip\text{-}metric\text{-}space}(v), \\ \quad \quad \quad \mathrm{dom}\,{{{\{p \in (\operatorname{vec}(v) \times \operatorname{vec}(v)) : ((\operatorname{nth}(1, p) \in d) \wedge (\operatorname{nth}(2, p) = f(\operatorname{nth}(1, p))))\}}^{*}}{\{p \in (\operatorname{vec}(v) \times \operatorname{vec}(v)) : ((\operatorname{nth}(1, p) \in d) \wedge (\operatorname{nth}(2, p) = f(\operatorname{nth}(1, p))))\}}}) \end{array}}}
\vnbprop{symmetric-op-form-positive}{\vnbstmt{\begin{array}{@{}l@{}} \forall v, d, f. \\ \quad \operatorname{is\text{-}complex\text{-}inner\text{-}product\text{-}space}(v) \wedge \\ \quad \operatorname{is\text{-}submodule}(v, d) \wedge \\ \quad f \in (d \to \operatorname{vec}(v)) \wedge \\ \quad \forall x, y \in d.\; (f(\operatorname{vadd}(v)(x, y)) = \operatorname{vadd}(v)(f(x), f(y))) \wedge \\ \quad \forall c \in \mathbb{C}, x \in d.\; (f(\operatorname{act}(v)(c, x)) = \operatorname{act}(v)(c, f(x))) \wedge \\ \quad \operatorname{is\text{-}dense}(\operatorname{ip\text{-}metric\text{-}space}(v), d) \wedge \\ \quad \forall \mathit{bsx}, \mathit{bsy} \in d.\; (\langle\cdot,\cdot\rangle_{v}(f(\mathit{bsx}), \mathit{bsy}) = \langle\cdot,\cdot\rangle_{v}(\mathit{bsx}, f(\mathit{bsy}))) \wedge \\ \quad \forall \mathit{bsx} \in d.\; ((\langle\cdot,\cdot\rangle_{v}(\mathit{bsx}, f(\mathit{bsx})) \in \mathbb{R}) \wedge (0 \leq \langle\cdot,\cdot\rangle_{v}(\mathit{bsx}, f(\mathit{bsx})))) \Rightarrow \\ \quad \quad \operatorname{is\text{-}positive\text{-}form}(v, d, (\lambda {[}\operatorname{bgx\_}, \operatorname{bgy\_}] \in (d \times d).\; \langle\cdot,\cdot\rangle_{v}(\operatorname{bgx\_}, f(\operatorname{bgy\_})))) \end{array}}}
\vnbprop{cips-module-scal}{\vnbstmt{\begin{array}{@{}l@{}} \forall v. \\ \quad \operatorname{is\text{-}complex\text{-}inner\text{-}product\text{-}space}(v) \Rightarrow  \\ \quad \quad (\operatorname{scal}(\operatorname{complex\text{-}inner\text{-}product\text{-}space\text{-}as\text{-}module}(v)) \simeq \operatorname{scal}(v)) \end{array}}}
\vnbprop{cips-module-vec}{\vnbstmt{\begin{array}{@{}l@{}} \forall v. \\ \quad \operatorname{is\text{-}complex\text{-}inner\text{-}product\text{-}space}(v) \Rightarrow  \\ \quad \quad (\operatorname{vec}(\operatorname{complex\text{-}inner\text{-}product\text{-}space\text{-}as\text{-}module}(v)) \simeq \operatorname{vec}(v)) \end{array}}}
\vnbprop{cips-module-vadd}{\vnbstmt{\begin{array}{@{}l@{}} \forall v. \\ \quad \operatorname{is\text{-}complex\text{-}inner\text{-}product\text{-}space}(v) \Rightarrow  \\ \quad \quad (\operatorname{vadd}(\operatorname{complex\text{-}inner\text{-}product\text{-}space\text{-}as\text{-}module}(v)) \simeq \operatorname{vadd}(v)) \end{array}}}
\vnbprop{cips-module-vzero}{\vnbstmt{\begin{array}{@{}l@{}} \forall v. \\ \quad \operatorname{is\text{-}complex\text{-}inner\text{-}product\text{-}space}(v) \Rightarrow  \\ \quad \quad (\operatorname{vzero}(\operatorname{complex\text{-}inner\text{-}product\text{-}space\text{-}as\text{-}module}(v)) \simeq \operatorname{vzero}(v)) \end{array}}}
\vnbprop{cips-module-vneg}{\vnbstmt{\begin{array}{@{}l@{}} \forall v. \\ \quad \operatorname{is\text{-}complex\text{-}inner\text{-}product\text{-}space}(v) \Rightarrow  \\ \quad \quad (\operatorname{vneg}(\operatorname{complex\text{-}inner\text{-}product\text{-}space\text{-}as\text{-}module}(v)) \simeq \operatorname{vneg}(v)) \end{array}}}
\vnbprop{cips-module-act}{\vnbstmt{\begin{array}{@{}l@{}} \forall v. \\ \quad \operatorname{is\text{-}complex\text{-}inner\text{-}product\text{-}space}(v) \Rightarrow  \\ \quad \quad (\operatorname{act}(\operatorname{complex\text{-}inner\text{-}product\text{-}space\text{-}as\text{-}module}(v)) \simeq \operatorname{act}(v)) \end{array}}}
\vnbprop{cips-module-scal-pin}{\vnbstmt{\begin{array}{@{}l@{}} \forall v. \\ \quad \operatorname{is\text{-}complex\text{-}inner\text{-}product\text{-}space}(v) \Rightarrow  \\ \quad \quad \operatorname{scal}(\operatorname{complex\text{-}inner\text{-}product\text{-}space\text{-}as\text{-}module}(v)) = \operatorname{normed\text{-}field\text{-}as\text{-}commutative\text{-}ring}(\operatorname{cc\text{-}normed\text{-}field}) \end{array}}}
\vnbprop{cips-module-inner-product-form}{\vnbstmt{\begin{array}{@{}l@{}} \forall v. \\ \quad \operatorname{is\text{-}complex\text{-}inner\text{-}product\text{-}space}(v) \Rightarrow  \\ \quad \quad \operatorname{is\text{-}inner\text{-}product\text{-}form}(\operatorname{complex\text{-}inner\text{-}product\text{-}space\text{-}as\text{-}module}(v), \langle\cdot,\cdot\rangle_{v}) \end{array}}}
\vnbprop{form-ms-cips-module}{\vnbstmt{\begin{array}{@{}l@{}} \forall v. \\ \quad \operatorname{is\text{-}complex\text{-}inner\text{-}product\text{-}space}(v) \Rightarrow  \\ \quad \quad (\operatorname{form\text{-}ms}(\operatorname{complex\text{-}inner\text{-}product\text{-}space\text{-}as\text{-}module}(v), \langle\cdot,\cdot\rangle_{v}) \simeq \operatorname{ip\text{-}metric\text{-}space}(v)) \end{array}}}
\vnbprop{form-ms-form-plus}{\vnbstmt{\begin{array}{@{}l@{}} \forall v, d, q. \\ \quad (\operatorname{is\text{-}complex\text{-}inner\text{-}product\text{-}space}(v) \wedge (\operatorname{is\text{-}submodule}(v, d) \wedge \operatorname{is\text{-}positive\text{-}form}(v, d, q))) \Rightarrow  \\ \quad \quad (\operatorname{form\text{-}ms}(\operatorname{complex\text{-}inner\text{-}product\text{-}space\text{-}as\text{-}module}(\operatorname{sub\text{-}cips}(v, d)), {q}_{+}) \simeq \operatorname{ip\text{-}metric\text{-}space}(H_{{q}+})) \end{array}}}
\vnbprop{fr-g-type}{\vnbstmt{\begin{array}{@{}l@{}} \forall v, d, q, g. \\ \quad \operatorname{is\text{-}hilbert\text{-}space}(v) \wedge \\ \quad \operatorname{is\text{-}positive\text{-}form}(v, d, q) \wedge \\ \quad \operatorname{is\text{-}uniformly\text{-}continuous}(\operatorname{ip\text{-}metric\text{-}space}(\operatorname{hilbert\text{-}completion}(\operatorname{complex\text{-}inner\text{-}product\text{-}space\text{-}as\text{-}module}(\operatorname{sub\text{-}cips}(v, d)), {q}_{+})), \\ \quad \quad \operatorname{ip\text{-}metric\text{-}space}(v), g) \wedge \\ \quad \forall x \in d. \\ \quad \quad g(\operatorname{hc\text{-}embed}(\operatorname{complex\text{-}inner\text{-}product\text{-}space\text{-}as\text{-}module}(\operatorname{sub\text{-}cips}(v, d)), {q}_{+})(x)) = x \Rightarrow \\ \quad \quad g \\ \quad \quad \in\; (\operatorname{vec}(\operatorname{hilbert\text{-}completion}(\operatorname{complex\text{-}inner\text{-}product\text{-}space\text{-}as\text{-}module}(\operatorname{sub\text{-}cips}(v, d)), {q}_{+})) \to \operatorname{vec}(v)) \end{array}}}
\vnbprop{fr-g-additive}{\vnbstmt{\begin{array}{@{}l@{}} \forall v, d, q, g, u, w. \\ \quad \operatorname{is\text{-}hilbert\text{-}space}(v) \wedge \\ \quad \operatorname{is\text{-}positive\text{-}form}(v, d, q) \wedge \\ \quad \operatorname{is\text{-}uniformly\text{-}continuous}(\operatorname{ip\text{-}metric\text{-}space}(\operatorname{hilbert\text{-}completion}(\operatorname{complex\text{-}inner\text{-}product\text{-}space\text{-}as\text{-}module}(\operatorname{sub\text{-}cips}(v, d)), {q}_{+})), \\ \quad \quad \operatorname{ip\text{-}metric\text{-}space}(v), g) \wedge \\ \quad \forall x \in d. \\ \quad \quad g(\operatorname{hc\text{-}embed}(\operatorname{complex\text{-}inner\text{-}product\text{-}space\text{-}as\text{-}module}(\operatorname{sub\text{-}cips}(v, d)), {q}_{+})(x)) = x \wedge \\ \quad u \\ \quad \in\; \operatorname{vec}(\operatorname{hilbert\text{-}completion}(\operatorname{complex\text{-}inner\text{-}product\text{-}space\text{-}as\text{-}module}(\operatorname{sub\text{-}cips}(v, d)), {q}_{+})) \wedge \\ \quad w \\ \quad \in\; \operatorname{vec}(\operatorname{hilbert\text{-}completion}(\operatorname{complex\text{-}inner\text{-}product\text{-}space\text{-}as\text{-}module}(\operatorname{sub\text{-}cips}(v, d)), {q}_{+})) \Rightarrow \\ \quad \quad g(\operatorname{vadd}(\operatorname{hilbert\text{-}completion}(\operatorname{complex\text{-}inner\text{-}product\text{-}space\text{-}as\text{-}module}(\operatorname{sub\text{-}cips}(v, d)), {q}_{+}))(u, w)) \\ \quad \quad =\; \operatorname{vadd}(v)(g(u), g(w)) \end{array}}}
\vnbprop{fr-g-homogeneous}{\vnbstmt{\begin{array}{@{}l@{}} \forall v, d, q, g, c, u. \\ \quad \operatorname{is\text{-}hilbert\text{-}space}(v) \wedge \\ \quad \operatorname{is\text{-}positive\text{-}form}(v, d, q) \wedge \\ \quad \operatorname{is\text{-}uniformly\text{-}continuous}(\operatorname{ip\text{-}metric\text{-}space}(\operatorname{hilbert\text{-}completion}(\operatorname{complex\text{-}inner\text{-}product\text{-}space\text{-}as\text{-}module}(\operatorname{sub\text{-}cips}(v, d)), {q}_{+})), \\ \quad \quad \operatorname{ip\text{-}metric\text{-}space}(v), g) \wedge \\ \quad \forall x \in d. \\ \quad \quad g(\operatorname{hc\text{-}embed}(\operatorname{complex\text{-}inner\text{-}product\text{-}space\text{-}as\text{-}module}(\operatorname{sub\text{-}cips}(v, d)), {q}_{+})(x)) = x \wedge \\ \quad c \in \mathbb{C} \wedge \\ \quad u \\ \quad \in\; \operatorname{vec}(\operatorname{hilbert\text{-}completion}(\operatorname{complex\text{-}inner\text{-}product\text{-}space\text{-}as\text{-}module}(\operatorname{sub\text{-}cips}(v, d)), {q}_{+})) \Rightarrow \\ \quad \quad g(\operatorname{act}(\operatorname{hilbert\text{-}completion}(\operatorname{complex\text{-}inner\text{-}product\text{-}space\text{-}as\text{-}module}(\operatorname{sub\text{-}cips}(v, d)), {q}_{+}))(c, u)) \\ \quad \quad =\; \operatorname{act}(v)(c, g(u)) \end{array}}}
\vnbprop{cips-ip-sub-left}{\vnbstmt{\begin{array}{@{}l@{}} \forall v, x, y, z. \\ \quad (\operatorname{is\text{-}complex\text{-}inner\text{-}product\text{-}space}(v) \wedge ((x \in \operatorname{vec}(v)) \wedge ((y \in \operatorname{vec}(v)) \wedge (z \in \operatorname{vec}(v))))) \Rightarrow  \\ \quad \quad (\langle\cdot,\cdot\rangle_{v}(x, z) - \langle\cdot,\cdot\rangle_{v}(y, z)) = \langle\cdot,\cdot\rangle_{v}(\operatorname{vadd}(v)(x, \operatorname{vneg}(v)(y)), z) \end{array}}}
\vnbprop{cips-ip-sub-left-bound}{\vnbstmt{\begin{array}{@{}l@{}} \forall v, x, y, z. \\ \quad (\operatorname{is\text{-}complex\text{-}inner\text{-}product\text{-}space}(v) \wedge ((x \in \operatorname{vec}(v)) \wedge ((y \in \operatorname{vec}(v)) \wedge (z \in \operatorname{vec}(v))))) \Rightarrow  \\ \quad \quad \operatorname{magnitude}((\langle\cdot,\cdot\rangle_{v}(x, z) - \langle\cdot,\cdot\rangle_{v}(y, z))) \leq (\operatorname{dist}(\operatorname{ip\text{-}metric\text{-}space}(v))(x, y) \cdot \lVert z \rVert) \end{array}}}
\vnbprop{cips-norm-dist-zero}{\vnbstmt{\begin{array}{@{}l@{}} \forall v, x. \\ \quad (\operatorname{is\text{-}complex\text{-}inner\text{-}product\text{-}space}(v) \wedge (x \in \operatorname{vec}(v))) \Rightarrow  \\ \quad \quad \lVert x \rVert = \operatorname{dist}(\operatorname{ip\text{-}metric\text{-}space}(v))(x, \operatorname{vzero}(v)) \end{array}}}
\vnbprop{cc-eq-of-small}{\vnbstmt{\begin{array}{@{}l@{}} \forall a, b \in \mathbb{C}, h \in \mathbb{R}. \\ \quad ((0 \leq h) \wedge \forall e.\; (\operatorname{pos\text{-}rr}(e) \Rightarrow (\operatorname{magnitude}((a - b)) \leq (e \cdot h)))) \Rightarrow  \\ \quad \quad a = b \end{array}}}
\vnbprop{fr-approx}{\vnbstmt{\begin{array}{@{}l@{}} \forall v, d, q, g, u, e. \\ \quad \operatorname{is\text{-}hilbert\text{-}space}(v) \wedge \\ \quad \operatorname{is\text{-}positive\text{-}form}(v, d, q) \wedge \\ \quad \operatorname{is\text{-}uniformly\text{-}continuous}(\operatorname{ip\text{-}metric\text{-}space}(\operatorname{hilbert\text{-}completion}(\operatorname{complex\text{-}inner\text{-}product\text{-}space\text{-}as\text{-}module}(\operatorname{sub\text{-}cips}(v, d)), {q}_{+})), \\ \quad \quad \operatorname{ip\text{-}metric\text{-}space}(v), g) \wedge \\ \quad \forall x \in d. \\ \quad \quad g(\operatorname{hc\text{-}embed}(\operatorname{complex\text{-}inner\text{-}product\text{-}space\text{-}as\text{-}module}(\operatorname{sub\text{-}cips}(v, d)), {q}_{+})(x)) = x \wedge \\ \quad u \\ \quad \in\; \operatorname{vec}(\operatorname{hilbert\text{-}completion}(\operatorname{complex\text{-}inner\text{-}product\text{-}space\text{-}as\text{-}module}(\operatorname{sub\text{-}cips}(v, d)), {q}_{+})) \wedge \\ \quad \operatorname{pos\text{-}rr}(e) \Rightarrow \\ \quad \quad \exists z \in d. \\ \quad \quad \quad \operatorname{dist}(\operatorname{ip\text{-}metric\text{-}space}(\operatorname{hilbert\text{-}completion}(\operatorname{complex\text{-}inner\text{-}product\text{-}space\text{-}as\text{-}module}(\operatorname{sub\text{-}cips}(v, d)), {q}_{+})))(u, \\ \quad \quad \quad \quad \operatorname{hc\text{-}embed}(\operatorname{complex\text{-}inner\text{-}product\text{-}space\text{-}as\text{-}module}(\operatorname{sub\text{-}cips}(v, d)), {q}_{+})(z)) \\ \quad \quad \quad \leq\; e \wedge \\ \quad \quad \quad \operatorname{dist}(\operatorname{ip\text{-}metric\text{-}space}(v))(g(u), z) \leq e \end{array}}}
\vnbprop{fr-embed-ip}{\vnbstmt{\begin{array}{@{}l@{}} \forall v, d, q, x, y. \\ \quad \operatorname{is\text{-}complex\text{-}inner\text{-}product\text{-}space}(v) \wedge \\ \quad \operatorname{is\text{-}submodule}(v, d) \wedge \\ \quad \operatorname{is\text{-}positive\text{-}form}(v, d, q) \wedge \\ \quad x \in d \wedge \\ \quad y \in d \Rightarrow \\ \quad \quad \langle\cdot,\cdot\rangle_{\operatorname{hilbert\text{-}completion}(\operatorname{complex\text{-}inner\text{-}product\text{-}space\text{-}as\text{-}module}(\operatorname{sub\text{-}cips}(v, d)), {q}_{+})}(\operatorname{hc\text{-}embed}(\operatorname{complex\text{-}inner\text{-}product\text{-}space\text{-}as\text{-}module}(\operatorname{sub\text{-}cips}(v, d)), {q}_{+})(x), \\ \quad \quad \quad \operatorname{hc\text{-}embed}(\operatorname{complex\text{-}inner\text{-}product\text{-}space\text{-}as\text{-}module}(\operatorname{sub\text{-}cips}(v, d)), {q}_{+})(y)) \\ \quad \quad =\; q(x, y) + \langle\cdot,\cdot\rangle_{v}(x, y) \end{array}}}
\vnbprop{fr-key}{\vnbstmt{\begin{array}{@{}l@{}} \forall v, d, q, g, f, u, y. \\ \quad \operatorname{is\text{-}hilbert\text{-}space}(v) \wedge \\ \quad \operatorname{is\text{-}positive\text{-}form}(v, d, q) \wedge \\ \quad \operatorname{is\text{-}uniformly\text{-}continuous}(\operatorname{ip\text{-}metric\text{-}space}(\operatorname{hilbert\text{-}completion}(\operatorname{complex\text{-}inner\text{-}product\text{-}space\text{-}as\text{-}module}(\operatorname{sub\text{-}cips}(v, d)), {q}_{+})), \\ \quad \quad \operatorname{ip\text{-}metric\text{-}space}(v), g) \wedge \\ \quad \forall x \in d. \\ \quad \quad g(\operatorname{hc\text{-}embed}(\operatorname{complex\text{-}inner\text{-}product\text{-}space\text{-}as\text{-}module}(\operatorname{sub\text{-}cips}(v, d)), {q}_{+})(x)) = x \wedge \\ \quad f \in (d \to \operatorname{vec}(v)) \wedge \\ \quad \forall \mathit{bty}, \mathit{btx} \in d.\; (q(\mathit{btx}, \mathit{bty}) = \langle\cdot,\cdot\rangle_{v}(\mathit{btx}, f(\mathit{bty}))) \wedge \\ \quad u \\ \quad \in\; \operatorname{vec}(\operatorname{hilbert\text{-}completion}(\operatorname{complex\text{-}inner\text{-}product\text{-}space\text{-}as\text{-}module}(\operatorname{sub\text{-}cips}(v, d)), {q}_{+})) \wedge \\ \quad y \in d \Rightarrow \\ \quad \quad \langle\cdot,\cdot\rangle_{\operatorname{hilbert\text{-}completion}(\operatorname{complex\text{-}inner\text{-}product\text{-}space\text{-}as\text{-}module}(\operatorname{sub\text{-}cips}(v, d)), {q}_{+})}(u, \\ \quad \quad \quad \operatorname{hc\text{-}embed}(\operatorname{complex\text{-}inner\text{-}product\text{-}space\text{-}as\text{-}module}(\operatorname{sub\text{-}cips}(v, d)), {q}_{+})(y)) \\ \quad \quad =\; \langle\cdot,\cdot\rangle_{v}(g(u), \operatorname{vadd}(v)(f(y), y)) \end{array}}}
\vnbprop{fr-separates}{\vnbstmt{\begin{array}{@{}l@{}} \forall v, d, q, u, w. \\ \quad \operatorname{is\text{-}complex\text{-}inner\text{-}product\text{-}space}(v) \wedge \\ \quad \operatorname{is\text{-}submodule}(v, d) \wedge \\ \quad \operatorname{is\text{-}positive\text{-}form}(v, d, q) \wedge \\ \quad u \\ \quad \in\; \operatorname{vec}(\operatorname{hilbert\text{-}completion}(\operatorname{complex\text{-}inner\text{-}product\text{-}space\text{-}as\text{-}module}(\operatorname{sub\text{-}cips}(v, d)), {q}_{+})) \wedge \\ \quad w \\ \quad \in\; \operatorname{vec}(\operatorname{hilbert\text{-}completion}(\operatorname{complex\text{-}inner\text{-}product\text{-}space\text{-}as\text{-}module}(\operatorname{sub\text{-}cips}(v, d)), {q}_{+})) \wedge \\ \quad \forall y \in d. \\ \quad \quad \langle\cdot,\cdot\rangle_{\operatorname{hilbert\text{-}completion}(\operatorname{complex\text{-}inner\text{-}product\text{-}space\text{-}as\text{-}module}(\operatorname{sub\text{-}cips}(v, d)), {q}_{+})}(u, \\ \quad \quad \quad \operatorname{hc\text{-}embed}(\operatorname{complex\text{-}inner\text{-}product\text{-}space\text{-}as\text{-}module}(\operatorname{sub\text{-}cips}(v, d)), {q}_{+})(y)) \\ \quad \quad =\; \langle\cdot,\cdot\rangle_{\operatorname{hilbert\text{-}completion}(\operatorname{complex\text{-}inner\text{-}product\text{-}space\text{-}as\text{-}module}(\operatorname{sub\text{-}cips}(v, d)), {q}_{+})}(w, \\ \quad \quad \quad \operatorname{hc\text{-}embed}(\operatorname{complex\text{-}inner\text{-}product\text{-}space\text{-}as\text{-}module}(\operatorname{sub\text{-}cips}(v, d)), {q}_{+})(y)) \Rightarrow \\ \quad \quad u = w \end{array}}}
\vnbprop{fr-g-injective}{\vnbstmt{\begin{array}{@{}l@{}} \forall v, d, q, g, f, u, w. \\ \quad \operatorname{is\text{-}hilbert\text{-}space}(v) \wedge \\ \quad \operatorname{is\text{-}positive\text{-}form}(v, d, q) \wedge \\ \quad \operatorname{is\text{-}uniformly\text{-}continuous}(\operatorname{ip\text{-}metric\text{-}space}(\operatorname{hilbert\text{-}completion}(\operatorname{complex\text{-}inner\text{-}product\text{-}space\text{-}as\text{-}module}(\operatorname{sub\text{-}cips}(v, d)), {q}_{+})), \\ \quad \quad \operatorname{ip\text{-}metric\text{-}space}(v), g) \wedge \\ \quad \forall x \in d. \\ \quad \quad g(\operatorname{hc\text{-}embed}(\operatorname{complex\text{-}inner\text{-}product\text{-}space\text{-}as\text{-}module}(\operatorname{sub\text{-}cips}(v, d)), {q}_{+})(x)) = x \wedge \\ \quad f \in (d \to \operatorname{vec}(v)) \wedge \\ \quad \forall y, \mathit{btx} \in d.\; (q(\mathit{btx}, y) = \langle\cdot,\cdot\rangle_{v}(\mathit{btx}, f(y))) \wedge \\ \quad u \\ \quad \in\; \operatorname{vec}(\operatorname{hilbert\text{-}completion}(\operatorname{complex\text{-}inner\text{-}product\text{-}space\text{-}as\text{-}module}(\operatorname{sub\text{-}cips}(v, d)), {q}_{+})) \wedge \\ \quad w \\ \quad \in\; \operatorname{vec}(\operatorname{hilbert\text{-}completion}(\operatorname{complex\text{-}inner\text{-}product\text{-}space\text{-}as\text{-}module}(\operatorname{sub\text{-}cips}(v, d)), {q}_{+})) \wedge \\ \quad g(u) = g(w) \Rightarrow \\ \quad \quad u = w \end{array}}}
\vnbprop{complete-of-isometric-cover}{\vnbstmt{\begin{array}{@{}l@{}} \forall s, t, f, h. \\ \quad \operatorname{is\text{-}complete}(s) \wedge \\ \quad \operatorname{is\text{-}metric\text{-}space}(t) \wedge \\ \quad f \in (\operatorname{pts}(s) \to \operatorname{pts}(t)) \wedge \\ \quad h \in (\operatorname{pts}(t) \to \operatorname{pts}(s)) \wedge \\ \quad \forall y \in \operatorname{pts}(t).\; (f(h(y)) = y) \wedge \\ \quad \forall x, z \in \operatorname{pts}(s).\; (\operatorname{dist}(t)(f(x), f(z)) = \operatorname{dist}(s)(x, z)) \Rightarrow \\ \quad \quad \operatorname{is\text{-}complete}(t) \end{array}}}
\vnbprop{fr-gi-left}{\vnbstmt{\begin{array}{@{}l@{}} \forall v, d, q, g, f, u. \\ \quad \operatorname{is\text{-}hilbert\text{-}space}(v) \wedge \\ \quad \operatorname{is\text{-}positive\text{-}form}(v, d, q) \wedge \\ \quad \operatorname{is\text{-}uniformly\text{-}continuous}(\operatorname{ip\text{-}metric\text{-}space}(\operatorname{hilbert\text{-}completion}(\operatorname{complex\text{-}inner\text{-}product\text{-}space\text{-}as\text{-}module}(\operatorname{sub\text{-}cips}(v, d)), {q}_{+})), \\ \quad \quad \operatorname{ip\text{-}metric\text{-}space}(v), g) \wedge \\ \quad \forall x \in d. \\ \quad \quad g(\operatorname{hc\text{-}embed}(\operatorname{complex\text{-}inner\text{-}product\text{-}space\text{-}as\text{-}module}(\operatorname{sub\text{-}cips}(v, d)), {q}_{+})(x)) = x \wedge \\ \quad f \in (d \to \operatorname{vec}(v)) \wedge \\ \quad \forall y, \mathit{btx} \in d.\; (q(\mathit{btx}, y) = \langle\cdot,\cdot\rangle_{v}(\mathit{btx}, f(y))) \wedge \\ \quad u \\ \quad \in\; \operatorname{vec}(\operatorname{hilbert\text{-}completion}(\operatorname{complex\text{-}inner\text{-}product\text{-}space\text{-}as\text{-}module}(\operatorname{sub\text{-}cips}(v, d)), {q}_{+})) \Rightarrow \\ \quad \quad (\iota\, \mathit{bgu}.\;((\mathit{bgu} \in \operatorname{vec}(\operatorname{hilbert\text{-}completion}(\operatorname{complex\text{-}inner\text{-}product\text{-}space\text{-}as\text{-}module}(\operatorname{sub\text{-}cips}(v, d)), {q}_{+}))) \wedge (g(\mathit{bgu}) = g(u)))) \\ \quad \quad =\; u \end{array}}}
\vnbprop{fr-d1-of}{\vnbstmt{\begin{array}{@{}l@{}} \forall v, d, q, g, f, u. \\ \quad \operatorname{is\text{-}hilbert\text{-}space}(v) \wedge \\ \quad \operatorname{is\text{-}positive\text{-}form}(v, d, q) \wedge \\ \quad \operatorname{is\text{-}uniformly\text{-}continuous}(\operatorname{ip\text{-}metric\text{-}space}(\operatorname{hilbert\text{-}completion}(\operatorname{complex\text{-}inner\text{-}product\text{-}space\text{-}as\text{-}module}(\operatorname{sub\text{-}cips}(v, d)), {q}_{+})), \\ \quad \quad \operatorname{ip\text{-}metric\text{-}space}(v), g) \wedge \\ \quad \forall x \in d. \\ \quad \quad g(\operatorname{hc\text{-}embed}(\operatorname{complex\text{-}inner\text{-}product\text{-}space\text{-}as\text{-}module}(\operatorname{sub\text{-}cips}(v, d)), {q}_{+})(x)) = x \wedge \\ \quad f \in (d \to \operatorname{vec}(v)) \wedge \\ \quad \forall y, \mathit{btx} \in d.\; (q(\mathit{btx}, y) = \langle\cdot,\cdot\rangle_{v}(\mathit{btx}, f(y))) \wedge \\ \quad u \\ \quad \in\; \operatorname{vec}(\operatorname{hilbert\text{-}completion}(\operatorname{complex\text{-}inner\text{-}product\text{-}space\text{-}as\text{-}module}(\operatorname{sub\text{-}cips}(v, d)), {q}_{+})) \Rightarrow \\ \quad \quad g(u) \\ \quad \quad \in\; \operatorname{image}(g, \\ \quad \quad \quad \operatorname{vec}(\operatorname{hilbert\text{-}completion}(\operatorname{complex\text{-}inner\text{-}product\text{-}space\text{-}as\text{-}module}(\operatorname{sub\text{-}cips}(v, d)), {q}_{+}))) \end{array}}}
\vnbprop{fr-gi-right}{\vnbstmt{\begin{array}{@{}l@{}} \forall v, d, q, g, f, a. \\ \quad \operatorname{is\text{-}hilbert\text{-}space}(v) \wedge \\ \quad \operatorname{is\text{-}positive\text{-}form}(v, d, q) \wedge \\ \quad \operatorname{is\text{-}uniformly\text{-}continuous}(\operatorname{ip\text{-}metric\text{-}space}(\operatorname{hilbert\text{-}completion}(\operatorname{complex\text{-}inner\text{-}product\text{-}space\text{-}as\text{-}module}(\operatorname{sub\text{-}cips}(v, d)), {q}_{+})), \\ \quad \quad \operatorname{ip\text{-}metric\text{-}space}(v), g) \wedge \\ \quad \forall x \in d. \\ \quad \quad g(\operatorname{hc\text{-}embed}(\operatorname{complex\text{-}inner\text{-}product\text{-}space\text{-}as\text{-}module}(\operatorname{sub\text{-}cips}(v, d)), {q}_{+})(x)) = x \wedge \\ \quad f \in (d \to \operatorname{vec}(v)) \wedge \\ \quad \forall y, \mathit{btx} \in d.\; (q(\mathit{btx}, y) = \langle\cdot,\cdot\rangle_{v}(\mathit{btx}, f(y))) \wedge \\ \quad a \\ \quad \in\; \operatorname{image}(g, \\ \quad \quad \operatorname{vec}(\operatorname{hilbert\text{-}completion}(\operatorname{complex\text{-}inner\text{-}product\text{-}space\text{-}as\text{-}module}(\operatorname{sub\text{-}cips}(v, d)), {q}_{+}))) \Rightarrow \\ \quad \quad (\iota\, u.\;((u \in \operatorname{vec}(\operatorname{hilbert\text{-}completion}(\operatorname{complex\text{-}inner\text{-}product\text{-}space\text{-}as\text{-}module}(\operatorname{sub\text{-}cips}(v, d)), {q}_{+}))) \wedge (g(u) = a))) \\ \quad \quad \in\; \operatorname{vec}(\operatorname{hilbert\text{-}completion}(\operatorname{complex\text{-}inner\text{-}product\text{-}space\text{-}as\text{-}module}(\operatorname{sub\text{-}cips}(v, d)), {q}_{+})) \wedge \\ \quad \quad g((\iota\, u.\;((u \in \operatorname{vec}(\operatorname{hilbert\text{-}completion}(\operatorname{complex\text{-}inner\text{-}product\text{-}space\text{-}as\text{-}module}(\operatorname{sub\text{-}cips}(v, d)), {q}_{+}))) \wedge (g(u) = a)))) \\ \quad \quad =\; a \end{array}}}
\vnbprop{fr-d1-subset}{\vnbstmt{\begin{array}{@{}l@{}} \forall v, d, q, g, f. \\ \quad \operatorname{is\text{-}hilbert\text{-}space}(v) \wedge \\ \quad \operatorname{is\text{-}positive\text{-}form}(v, d, q) \wedge \\ \quad \operatorname{is\text{-}uniformly\text{-}continuous}(\operatorname{ip\text{-}metric\text{-}space}(\operatorname{hilbert\text{-}completion}(\operatorname{complex\text{-}inner\text{-}product\text{-}space\text{-}as\text{-}module}(\operatorname{sub\text{-}cips}(v, d)), {q}_{+})), \\ \quad \quad \operatorname{ip\text{-}metric\text{-}space}(v), g) \wedge \\ \quad \forall x \in d. \\ \quad \quad g(\operatorname{hc\text{-}embed}(\operatorname{complex\text{-}inner\text{-}product\text{-}space\text{-}as\text{-}module}(\operatorname{sub\text{-}cips}(v, d)), {q}_{+})(x)) = x \wedge \\ \quad f \in (d \to \operatorname{vec}(v)) \wedge \\ \quad \forall y, \mathit{btx} \in d.\; (q(\mathit{btx}, y) = \langle\cdot,\cdot\rangle_{v}(\mathit{btx}, f(y))) \Rightarrow \\ \quad \quad \operatorname{image}(g, \\ \quad \quad \quad \operatorname{vec}(\operatorname{hilbert\text{-}completion}(\operatorname{complex\text{-}inner\text{-}product\text{-}space\text{-}as\text{-}module}(\operatorname{sub\text{-}cips}(v, d)), {q}_{+}))) \\ \quad \quad \subseteq\; \operatorname{vec}(v) \end{array}}}
\vnbprop{fr-dd-in-d1}{\vnbstmt{\begin{array}{@{}l@{}} \forall v, d, q, g, f, y. \\ \quad \operatorname{is\text{-}hilbert\text{-}space}(v) \wedge \\ \quad \operatorname{is\text{-}positive\text{-}form}(v, d, q) \wedge \\ \quad \operatorname{is\text{-}uniformly\text{-}continuous}(\operatorname{ip\text{-}metric\text{-}space}(\operatorname{hilbert\text{-}completion}(\operatorname{complex\text{-}inner\text{-}product\text{-}space\text{-}as\text{-}module}(\operatorname{sub\text{-}cips}(v, d)), {q}_{+})), \\ \quad \quad \operatorname{ip\text{-}metric\text{-}space}(v), g) \wedge \\ \quad \forall x \in d. \\ \quad \quad g(\operatorname{hc\text{-}embed}(\operatorname{complex\text{-}inner\text{-}product\text{-}space\text{-}as\text{-}module}(\operatorname{sub\text{-}cips}(v, d)), {q}_{+})(x)) = x \wedge \\ \quad f \in (d \to \operatorname{vec}(v)) \wedge \\ \quad \forall \mathit{bty}, \mathit{btx} \in d.\; (q(\mathit{btx}, \mathit{bty}) = \langle\cdot,\cdot\rangle_{v}(\mathit{btx}, f(\mathit{bty}))) \wedge \\ \quad y \in d \Rightarrow \\ \quad \quad y \\ \quad \quad \in\; \operatorname{image}(g, \\ \quad \quad \quad \operatorname{vec}(\operatorname{hilbert\text{-}completion}(\operatorname{complex\text{-}inner\text{-}product\text{-}space\text{-}as\text{-}module}(\operatorname{sub\text{-}cips}(v, d)), {q}_{+}))) \wedge \\ \quad \quad (\iota\, u.\;((u \in \operatorname{vec}(\operatorname{hilbert\text{-}completion}(\operatorname{complex\text{-}inner\text{-}product\text{-}space\text{-}as\text{-}module}(\operatorname{sub\text{-}cips}(v, d)), {q}_{+}))) \wedge (g(u) = y))) \\ \quad \quad =\; \operatorname{hc\text{-}embed}(\operatorname{complex\text{-}inner\text{-}product\text{-}space\text{-}as\text{-}module}(\operatorname{sub\text{-}cips}(v, d)), {q}_{+})(y) \end{array}}}
\vnbprop{fr-g-vadd-gi}{\vnbstmt{\begin{array}{@{}l@{}} \forall v, d, q, g, f, a, b. \\ \quad \operatorname{is\text{-}hilbert\text{-}space}(v) \wedge \\ \quad \operatorname{is\text{-}positive\text{-}form}(v, d, q) \wedge \\ \quad \operatorname{is\text{-}uniformly\text{-}continuous}(\operatorname{ip\text{-}metric\text{-}space}(\operatorname{hilbert\text{-}completion}(\operatorname{complex\text{-}inner\text{-}product\text{-}space\text{-}as\text{-}module}(\operatorname{sub\text{-}cips}(v, d)), {q}_{+})), \\ \quad \quad \operatorname{ip\text{-}metric\text{-}space}(v), g) \wedge \\ \quad \forall x \in d. \\ \quad \quad g(\operatorname{hc\text{-}embed}(\operatorname{complex\text{-}inner\text{-}product\text{-}space\text{-}as\text{-}module}(\operatorname{sub\text{-}cips}(v, d)), {q}_{+})(x)) = x \wedge \\ \quad f \in (d \to \operatorname{vec}(v)) \wedge \\ \quad \forall y, \mathit{btx} \in d.\; (q(\mathit{btx}, y) = \langle\cdot,\cdot\rangle_{v}(\mathit{btx}, f(y))) \wedge \\ \quad a \\ \quad \in\; \operatorname{image}(g, \\ \quad \quad \operatorname{vec}(\operatorname{hilbert\text{-}completion}(\operatorname{complex\text{-}inner\text{-}product\text{-}space\text{-}as\text{-}module}(\operatorname{sub\text{-}cips}(v, d)), {q}_{+}))) \wedge \\ \quad b \\ \quad \in\; \operatorname{image}(g, \\ \quad \quad \operatorname{vec}(\operatorname{hilbert\text{-}completion}(\operatorname{complex\text{-}inner\text{-}product\text{-}space\text{-}as\text{-}module}(\operatorname{sub\text{-}cips}(v, d)), {q}_{+}))) \Rightarrow \\ \quad \quad g(\operatorname{vadd}(\operatorname{hilbert\text{-}completion}(\operatorname{complex\text{-}inner\text{-}product\text{-}space\text{-}as\text{-}module}(\operatorname{sub\text{-}cips}(v, d)), {q}_{+}))((\iota\, u.\;((u \in \operatorname{vec}(\operatorname{hilbert\text{-}completion}(\operatorname{complex\text{-}inner\text{-}product\text{-}space\text{-}as\text{-}module}(\operatorname{sub\text{-}cips}(v, d)), {q}_{+}))) \wedge (g(u) = a))), \\ \quad \quad \quad \quad (\iota\, u.\;((u \in \operatorname{vec}(\operatorname{hilbert\text{-}completion}(\operatorname{complex\text{-}inner\text{-}product\text{-}space\text{-}as\text{-}module}(\operatorname{sub\text{-}cips}(v, d)), {q}_{+}))) \wedge (g(u) = b))))) \\ \quad \quad =\; \operatorname{vadd}(v)(a, b) \end{array}}}
\vnbprop{fr-g-act-gi}{\vnbstmt{\begin{array}{@{}l@{}} \forall v, d, q, g, f, c, a. \\ \quad \operatorname{is\text{-}hilbert\text{-}space}(v) \wedge \\ \quad \operatorname{is\text{-}positive\text{-}form}(v, d, q) \wedge \\ \quad \operatorname{is\text{-}uniformly\text{-}continuous}(\operatorname{ip\text{-}metric\text{-}space}(\operatorname{hilbert\text{-}completion}(\operatorname{complex\text{-}inner\text{-}product\text{-}space\text{-}as\text{-}module}(\operatorname{sub\text{-}cips}(v, d)), {q}_{+})), \\ \quad \quad \operatorname{ip\text{-}metric\text{-}space}(v), g) \wedge \\ \quad \forall x \in d. \\ \quad \quad g(\operatorname{hc\text{-}embed}(\operatorname{complex\text{-}inner\text{-}product\text{-}space\text{-}as\text{-}module}(\operatorname{sub\text{-}cips}(v, d)), {q}_{+})(x)) = x \wedge \\ \quad f \in (d \to \operatorname{vec}(v)) \wedge \\ \quad \forall y, \mathit{btx} \in d.\; (q(\mathit{btx}, y) = \langle\cdot,\cdot\rangle_{v}(\mathit{btx}, f(y))) \wedge \\ \quad c \in \mathbb{C} \wedge \\ \quad a \\ \quad \in\; \operatorname{image}(g, \\ \quad \quad \operatorname{vec}(\operatorname{hilbert\text{-}completion}(\operatorname{complex\text{-}inner\text{-}product\text{-}space\text{-}as\text{-}module}(\operatorname{sub\text{-}cips}(v, d)), {q}_{+}))) \Rightarrow \\ \quad \quad g(\operatorname{act}(\operatorname{hilbert\text{-}completion}(\operatorname{complex\text{-}inner\text{-}product\text{-}space\text{-}as\text{-}module}(\operatorname{sub\text{-}cips}(v, d)), {q}_{+}))(c, \\ \quad \quad \quad \quad (\iota\, u.\;((u \in \operatorname{vec}(\operatorname{hilbert\text{-}completion}(\operatorname{complex\text{-}inner\text{-}product\text{-}space\text{-}as\text{-}module}(\operatorname{sub\text{-}cips}(v, d)), {q}_{+}))) \wedge (g(u) = a))))) \\ \quad \quad =\; \operatorname{act}(v)(c, a) \end{array}}}
\vnbprop{fr-gi-add}{\vnbstmt{\begin{array}{@{}l@{}} \forall v, d, q, g, f, a, b. \\ \quad \operatorname{is\text{-}hilbert\text{-}space}(v) \wedge \\ \quad \operatorname{is\text{-}positive\text{-}form}(v, d, q) \wedge \\ \quad \operatorname{is\text{-}uniformly\text{-}continuous}(\operatorname{ip\text{-}metric\text{-}space}(\operatorname{hilbert\text{-}completion}(\operatorname{complex\text{-}inner\text{-}product\text{-}space\text{-}as\text{-}module}(\operatorname{sub\text{-}cips}(v, d)), {q}_{+})), \\ \quad \quad \operatorname{ip\text{-}metric\text{-}space}(v), g) \wedge \\ \quad \forall x \in d. \\ \quad \quad g(\operatorname{hc\text{-}embed}(\operatorname{complex\text{-}inner\text{-}product\text{-}space\text{-}as\text{-}module}(\operatorname{sub\text{-}cips}(v, d)), {q}_{+})(x)) = x \wedge \\ \quad f \in (d \to \operatorname{vec}(v)) \wedge \\ \quad \forall y, \mathit{btx} \in d.\; (q(\mathit{btx}, y) = \langle\cdot,\cdot\rangle_{v}(\mathit{btx}, f(y))) \wedge \\ \quad a \\ \quad \in\; \operatorname{image}(g, \\ \quad \quad \operatorname{vec}(\operatorname{hilbert\text{-}completion}(\operatorname{complex\text{-}inner\text{-}product\text{-}space\text{-}as\text{-}module}(\operatorname{sub\text{-}cips}(v, d)), {q}_{+}))) \wedge \\ \quad b \\ \quad \in\; \operatorname{image}(g, \\ \quad \quad \operatorname{vec}(\operatorname{hilbert\text{-}completion}(\operatorname{complex\text{-}inner\text{-}product\text{-}space\text{-}as\text{-}module}(\operatorname{sub\text{-}cips}(v, d)), {q}_{+}))) \Rightarrow \\ \quad \quad (\iota\, u.\;((u \in \operatorname{vec}(\operatorname{hilbert\text{-}completion}(\operatorname{complex\text{-}inner\text{-}product\text{-}space\text{-}as\text{-}module}(\operatorname{sub\text{-}cips}(v, d)), {q}_{+}))) \wedge (g(u) = \operatorname{vadd}(v)(a, b)))) \\ \quad \quad =\; \operatorname{vadd}(\operatorname{hilbert\text{-}completion}(\operatorname{complex\text{-}inner\text{-}product\text{-}space\text{-}as\text{-}module}(\operatorname{sub\text{-}cips}(v, d)), {q}_{+}))((\iota\, u.\;((u \in \operatorname{vec}(\operatorname{hilbert\text{-}completion}(\operatorname{complex\text{-}inner\text{-}product\text{-}space\text{-}as\text{-}module}(\operatorname{sub\text{-}cips}(v, d)), {q}_{+}))) \wedge (g(u) = a))), \\ \quad \quad \quad (\iota\, u.\;((u \in \operatorname{vec}(\operatorname{hilbert\text{-}completion}(\operatorname{complex\text{-}inner\text{-}product\text{-}space\text{-}as\text{-}module}(\operatorname{sub\text{-}cips}(v, d)), {q}_{+}))) \wedge (g(u) = b)))) \end{array}}}
\vnbprop{fr-gi-act}{\vnbstmt{\begin{array}{@{}l@{}} \forall v, d, q, g, f, c, a. \\ \quad \operatorname{is\text{-}hilbert\text{-}space}(v) \wedge \\ \quad \operatorname{is\text{-}positive\text{-}form}(v, d, q) \wedge \\ \quad \operatorname{is\text{-}uniformly\text{-}continuous}(\operatorname{ip\text{-}metric\text{-}space}(\operatorname{hilbert\text{-}completion}(\operatorname{complex\text{-}inner\text{-}product\text{-}space\text{-}as\text{-}module}(\operatorname{sub\text{-}cips}(v, d)), {q}_{+})), \\ \quad \quad \operatorname{ip\text{-}metric\text{-}space}(v), g) \wedge \\ \quad \forall x \in d. \\ \quad \quad g(\operatorname{hc\text{-}embed}(\operatorname{complex\text{-}inner\text{-}product\text{-}space\text{-}as\text{-}module}(\operatorname{sub\text{-}cips}(v, d)), {q}_{+})(x)) = x \wedge \\ \quad f \in (d \to \operatorname{vec}(v)) \wedge \\ \quad \forall y, \mathit{btx} \in d.\; (q(\mathit{btx}, y) = \langle\cdot,\cdot\rangle_{v}(\mathit{btx}, f(y))) \wedge \\ \quad c \in \mathbb{C} \wedge \\ \quad a \\ \quad \in\; \operatorname{image}(g, \\ \quad \quad \operatorname{vec}(\operatorname{hilbert\text{-}completion}(\operatorname{complex\text{-}inner\text{-}product\text{-}space\text{-}as\text{-}module}(\operatorname{sub\text{-}cips}(v, d)), {q}_{+}))) \Rightarrow \\ \quad \quad (\iota\, u.\;((u \in \operatorname{vec}(\operatorname{hilbert\text{-}completion}(\operatorname{complex\text{-}inner\text{-}product\text{-}space\text{-}as\text{-}module}(\operatorname{sub\text{-}cips}(v, d)), {q}_{+}))) \wedge (g(u) = \operatorname{act}(v)(c, a)))) \\ \quad \quad =\; \operatorname{act}(\operatorname{hilbert\text{-}completion}(\operatorname{complex\text{-}inner\text{-}product\text{-}space\text{-}as\text{-}module}(\operatorname{sub\text{-}cips}(v, d)), {q}_{+}))(c, \\ \quad \quad \quad (\iota\, u.\;((u \in \operatorname{vec}(\operatorname{hilbert\text{-}completion}(\operatorname{complex\text{-}inner\text{-}product\text{-}space\text{-}as\text{-}module}(\operatorname{sub\text{-}cips}(v, d)), {q}_{+}))) \wedge (g(u) = a)))) \end{array}}}
\vnbprop{fr-d1-submodule}{\vnbstmt{\begin{array}{@{}l@{}} \forall v, d, q, g, f. \\ \quad \operatorname{is\text{-}hilbert\text{-}space}(v) \wedge \\ \quad \operatorname{is\text{-}positive\text{-}form}(v, d, q) \wedge \\ \quad \operatorname{is\text{-}uniformly\text{-}continuous}(\operatorname{ip\text{-}metric\text{-}space}(\operatorname{hilbert\text{-}completion}(\operatorname{complex\text{-}inner\text{-}product\text{-}space\text{-}as\text{-}module}(\operatorname{sub\text{-}cips}(v, d)), {q}_{+})), \\ \quad \quad \operatorname{ip\text{-}metric\text{-}space}(v), g) \wedge \\ \quad \forall x \in d. \\ \quad \quad g(\operatorname{hc\text{-}embed}(\operatorname{complex\text{-}inner\text{-}product\text{-}space\text{-}as\text{-}module}(\operatorname{sub\text{-}cips}(v, d)), {q}_{+})(x)) = x \wedge \\ \quad f \in (d \to \operatorname{vec}(v)) \wedge \\ \quad \forall y, \mathit{btx} \in d.\; (q(\mathit{btx}, y) = \langle\cdot,\cdot\rangle_{v}(\mathit{btx}, f(y))) \Rightarrow \\ \quad \quad \operatorname{is\text{-}submodule}(v, \\ \quad \quad \quad \operatorname{image}(g, \\ \quad \quad \quad \quad \operatorname{vec}(\operatorname{hilbert\text{-}completion}(\operatorname{complex\text{-}inner\text{-}product\text{-}space\text{-}as\text{-}module}(\operatorname{sub\text{-}cips}(v, d)), {q}_{+})))) \end{array}}}
\vnbprop{fr-d1-dense}{\vnbstmt{\begin{array}{@{}l@{}} \forall v, d, q, g, f. \\ \quad \operatorname{is\text{-}hilbert\text{-}space}(v) \wedge \\ \quad \operatorname{is\text{-}positive\text{-}form}(v, d, q) \wedge \\ \quad \operatorname{is\text{-}uniformly\text{-}continuous}(\operatorname{ip\text{-}metric\text{-}space}(\operatorname{hilbert\text{-}completion}(\operatorname{complex\text{-}inner\text{-}product\text{-}space\text{-}as\text{-}module}(\operatorname{sub\text{-}cips}(v, d)), {q}_{+})), \\ \quad \quad \operatorname{ip\text{-}metric\text{-}space}(v), g) \wedge \\ \quad \forall x \in d. \\ \quad \quad g(\operatorname{hc\text{-}embed}(\operatorname{complex\text{-}inner\text{-}product\text{-}space\text{-}as\text{-}module}(\operatorname{sub\text{-}cips}(v, d)), {q}_{+})(x)) = x \wedge \\ \quad f \in (d \to \operatorname{vec}(v)) \wedge \\ \quad \forall y, \mathit{btx} \in d.\; (q(\mathit{btx}, y) = \langle\cdot,\cdot\rangle_{v}(\mathit{btx}, f(y))) \Rightarrow \\ \quad \quad \operatorname{is\text{-}dense}(\operatorname{ip\text{-}metric\text{-}space}(v), \\ \quad \quad \quad \operatorname{image}(g, \\ \quad \quad \quad \quad \operatorname{vec}(\operatorname{hilbert\text{-}completion}(\operatorname{complex\text{-}inner\text{-}product\text{-}space\text{-}as\text{-}module}(\operatorname{sub\text{-}cips}(v, d)), {q}_{+})))) \end{array}}}
\vnbprop{fr-norm-le-embed}{\vnbstmt{\begin{array}{@{}l@{}} \forall v, d, q, x. \\ \quad (\operatorname{is\text{-}complex\text{-}inner\text{-}product\text{-}space}(v) \wedge (\operatorname{is\text{-}submodule}(v, d) \wedge (\operatorname{is\text{-}positive\text{-}form}(v, d, q) \wedge (x \in d)))) \Rightarrow  \\ \quad \quad \lVert x \rVert \\ \quad \quad \leq\; \lVert \operatorname{hc\text{-}embed}(\operatorname{complex\text{-}inner\text{-}product\text{-}space\text{-}as\text{-}module}(\operatorname{sub\text{-}cips}(v, d)), {q}_{+})(x) \rVert \end{array}}}
\vnbprop{fr-g-norm-le}{\vnbstmt{\begin{array}{@{}l@{}} \forall v, d, q, g, u. \\ \quad \operatorname{is\text{-}hilbert\text{-}space}(v) \wedge \\ \quad \operatorname{is\text{-}positive\text{-}form}(v, d, q) \wedge \\ \quad \operatorname{is\text{-}uniformly\text{-}continuous}(\operatorname{ip\text{-}metric\text{-}space}(\operatorname{hilbert\text{-}completion}(\operatorname{complex\text{-}inner\text{-}product\text{-}space\text{-}as\text{-}module}(\operatorname{sub\text{-}cips}(v, d)), {q}_{+})), \\ \quad \quad \operatorname{ip\text{-}metric\text{-}space}(v), g) \wedge \\ \quad \forall x \in d. \\ \quad \quad g(\operatorname{hc\text{-}embed}(\operatorname{complex\text{-}inner\text{-}product\text{-}space\text{-}as\text{-}module}(\operatorname{sub\text{-}cips}(v, d)), {q}_{+})(x)) = x \wedge \\ \quad u \\ \quad \in\; \operatorname{vec}(\operatorname{hilbert\text{-}completion}(\operatorname{complex\text{-}inner\text{-}product\text{-}space\text{-}as\text{-}module}(\operatorname{sub\text{-}cips}(v, d)), {q}_{+})) \Rightarrow \\ \quad \quad \lVert g(u) \rVert \\ \quad \quad \leq\; \lVert u \rVert \end{array}}}
\vnbprop{fr-g-ip-le}{\vnbstmt{\begin{array}{@{}l@{}} \forall v, d, q, g, u. \\ \quad \operatorname{is\text{-}hilbert\text{-}space}(v) \wedge \\ \quad \operatorname{is\text{-}positive\text{-}form}(v, d, q) \wedge \\ \quad \operatorname{is\text{-}uniformly\text{-}continuous}(\operatorname{ip\text{-}metric\text{-}space}(\operatorname{hilbert\text{-}completion}(\operatorname{complex\text{-}inner\text{-}product\text{-}space\text{-}as\text{-}module}(\operatorname{sub\text{-}cips}(v, d)), {q}_{+})), \\ \quad \quad \operatorname{ip\text{-}metric\text{-}space}(v), g) \wedge \\ \quad \forall x \in d. \\ \quad \quad g(\operatorname{hc\text{-}embed}(\operatorname{complex\text{-}inner\text{-}product\text{-}space\text{-}as\text{-}module}(\operatorname{sub\text{-}cips}(v, d)), {q}_{+})(x)) = x \wedge \\ \quad u \\ \quad \in\; \operatorname{vec}(\operatorname{hilbert\text{-}completion}(\operatorname{complex\text{-}inner\text{-}product\text{-}space\text{-}as\text{-}module}(\operatorname{sub\text{-}cips}(v, d)), {q}_{+})) \Rightarrow \\ \quad \quad \langle\cdot,\cdot\rangle_{v}(g(u), g(u)) \\ \quad \quad \leq\; \langle\cdot,\cdot\rangle_{\operatorname{hilbert\text{-}completion}(\operatorname{complex\text{-}inner\text{-}product\text{-}space\text{-}as\text{-}module}(\operatorname{sub\text{-}cips}(v, d)), {q}_{+})}(u, u) \end{array}}}
\vnbprop{pb-value}{\vnbstmt{\begin{array}{@{}l@{}} \forall v, d, w, i, x, y. \\ \quad \operatorname{is\text{-}hilbert\text{-}space}(v) \wedge \\ \quad \operatorname{is\text{-}hilbert\text{-}space}(w) \wedge \\ \quad \operatorname{is\text{-}submodule}(v, d) \wedge \\ \quad \operatorname{is\text{-}dense}(\operatorname{ip\text{-}metric\text{-}space}(v), d) \wedge \\ \quad i \in (d \to \operatorname{vec}(w)) \wedge \\ \quad \forall \mathit{bpx}, \mathit{bpy} \in d.\; (i(\operatorname{vadd}(v)(\mathit{bpx}, \mathit{bpy})) = \operatorname{vadd}(w)(i(\mathit{bpx}), i(\mathit{bpy}))) \wedge \\ \quad \forall c \in \mathbb{C}, \mathit{bpx} \in d.\; (i(\operatorname{act}(v)(c, \mathit{bpx})) = \operatorname{act}(w)(c, i(\mathit{bpx}))) \wedge \\ \quad \forall \mathit{bpx} \in d.\; (\langle\cdot,\cdot\rangle_{v}(\mathit{bpx}, \mathit{bpx}) \leq \langle\cdot,\cdot\rangle_{w}(i(\mathit{bpx}), i(\mathit{bpx}))) \wedge \\ \quad x \in d \wedge \\ \quad y \in d \Rightarrow \\ \quad \quad (\lambda {[}\operatorname{bga\_}, \operatorname{bgb\_}] \in (d \times d).\; (\langle\cdot,\cdot\rangle_{w}(i(\operatorname{bga\_}), i(\operatorname{bgb\_})) - \langle\cdot,\cdot\rangle_{v}(\operatorname{bga\_}, \operatorname{bgb\_})))(x, y) \\ \quad \quad =\; (\langle\cdot,\cdot\rangle_{w}(i(x), i(y)) - \langle\cdot,\cdot\rangle_{v}(x, y)) \end{array}}}
\vnbprop{pb-positive}{\vnbstmt{\begin{array}{@{}l@{}} \forall v, d, w, i. \\ \quad \operatorname{is\text{-}hilbert\text{-}space}(v) \wedge \\ \quad \operatorname{is\text{-}hilbert\text{-}space}(w) \wedge \\ \quad \operatorname{is\text{-}submodule}(v, d) \wedge \\ \quad \operatorname{is\text{-}dense}(\operatorname{ip\text{-}metric\text{-}space}(v), d) \wedge \\ \quad i \in (d \to \operatorname{vec}(w)) \wedge \\ \quad \forall x, y \in d.\; (i(\operatorname{vadd}(v)(x, y)) = \operatorname{vadd}(w)(i(x), i(y))) \wedge \\ \quad \forall c \in \mathbb{C}, x \in d.\; (i(\operatorname{act}(v)(c, x)) = \operatorname{act}(w)(c, i(x))) \wedge \\ \quad \forall x \in d.\; (\langle\cdot,\cdot\rangle_{v}(x, x) \leq \langle\cdot,\cdot\rangle_{w}(i(x), i(x))) \Rightarrow \\ \quad \quad \operatorname{is\text{-}positive\text{-}form}(v, d, \\ \quad \quad \quad (\lambda {[}\operatorname{bga\_}, \operatorname{bgb\_}] \in (d \times d).\; (\langle\cdot,\cdot\rangle_{w}(i(\operatorname{bga\_}), i(\operatorname{bgb\_})) - \langle\cdot,\cdot\rangle_{v}(\operatorname{bga\_}, \operatorname{bgb\_})))) \end{array}}}
\vnbprop{pb-gi-vneg}{\vnbstmt{\begin{array}{@{}l@{}} \forall v, d, w, i, y. \\ \quad \operatorname{is\text{-}hilbert\text{-}space}(v) \wedge \\ \quad \operatorname{is\text{-}hilbert\text{-}space}(w) \wedge \\ \quad \operatorname{is\text{-}submodule}(v, d) \wedge \\ \quad \operatorname{is\text{-}dense}(\operatorname{ip\text{-}metric\text{-}space}(v), d) \wedge \\ \quad i \in (d \to \operatorname{vec}(w)) \wedge \\ \quad \forall x, \mathit{bpy} \in d.\; (i(\operatorname{vadd}(v)(x, \mathit{bpy})) = \operatorname{vadd}(w)(i(x), i(\mathit{bpy}))) \wedge \\ \quad \forall c \in \mathbb{C}, x \in d.\; (i(\operatorname{act}(v)(c, x)) = \operatorname{act}(w)(c, i(x))) \wedge \\ \quad \forall x \in d.\; (\langle\cdot,\cdot\rangle_{v}(x, x) \leq \langle\cdot,\cdot\rangle_{w}(i(x), i(x))) \wedge \\ \quad y \in d \Rightarrow \\ \quad \quad i(\operatorname{vneg}(v)(y)) = \operatorname{vneg}(w)(i(y)) \end{array}}}
\vnbprop{pb-isometry}{\vnbstmt{\begin{array}{@{}l@{}} \forall v, d, w, i, h, x, z. \\ \quad \operatorname{is\text{-}hilbert\text{-}space}(v) \wedge \\ \quad \operatorname{is\text{-}hilbert\text{-}space}(w) \wedge \\ \quad \operatorname{is\text{-}submodule}(v, d) \wedge \\ \quad \operatorname{is\text{-}dense}(\operatorname{ip\text{-}metric\text{-}space}(v), d) \wedge \\ \quad i \in (d \to \operatorname{vec}(w)) \wedge \\ \quad \forall \mathit{bpx}, y \in d.\; (i(\operatorname{vadd}(v)(\mathit{bpx}, y)) = \operatorname{vadd}(w)(i(\mathit{bpx}), i(y))) \wedge \\ \quad \forall c \in \mathbb{C}, \mathit{bpx} \in d.\; (i(\operatorname{act}(v)(c, \mathit{bpx})) = \operatorname{act}(w)(c, i(\mathit{bpx}))) \wedge \\ \quad \forall \mathit{bpx} \in d.\; (\langle\cdot,\cdot\rangle_{v}(\mathit{bpx}, \mathit{bpx}) \leq \langle\cdot,\cdot\rangle_{w}(i(\mathit{bpx}), i(\mathit{bpx}))) \wedge \\ \quad h \in (\operatorname{vec}(w) \to d) \wedge \\ \quad \forall y \in \operatorname{vec}(w).\; (i(h(y)) = y) \wedge \\ \quad \forall \mathit{bpx} \in d.\; (h(i(\mathit{bpx})) = \mathit{bpx}) \wedge \\ \quad x \in \operatorname{vec}(w) \wedge \\ \quad z \in \operatorname{vec}(w) \Rightarrow \\ \quad \quad \operatorname{dist}(\operatorname{ip\text{-}metric\text{-}space}(H_{{(\lambda {[}\operatorname{bga\_}, \operatorname{bgb\_}] \in (d \times d).\; (\langle\cdot,\cdot\rangle_{w}(i(\operatorname{bga\_}), i(\operatorname{bgb\_})) - \langle\cdot,\cdot\rangle_{v}(\operatorname{bga\_}, \operatorname{bgb\_})))}+}))(h(x), h(z)) \\ \quad \quad =\; \operatorname{dist}(\operatorname{ip\text{-}metric\text{-}space}(w))(x, z) \end{array}}}
\vnbprop{pb-hilbert}{\vnbstmt{\begin{array}{@{}l@{}} \forall v, d, w, i, h. \\ \quad \operatorname{is\text{-}hilbert\text{-}space}(v) \wedge \\ \quad \operatorname{is\text{-}hilbert\text{-}space}(w) \wedge \\ \quad \operatorname{is\text{-}submodule}(v, d) \wedge \\ \quad \operatorname{is\text{-}dense}(\operatorname{ip\text{-}metric\text{-}space}(v), d) \wedge \\ \quad i \in (d \to \operatorname{vec}(w)) \wedge \\ \quad \forall x, y \in d.\; (i(\operatorname{vadd}(v)(x, y)) = \operatorname{vadd}(w)(i(x), i(y))) \wedge \\ \quad \forall c \in \mathbb{C}, x \in d.\; (i(\operatorname{act}(v)(c, x)) = \operatorname{act}(w)(c, i(x))) \wedge \\ \quad \forall x \in d.\; (\langle\cdot,\cdot\rangle_{v}(x, x) \leq \langle\cdot,\cdot\rangle_{w}(i(x), i(x))) \wedge \\ \quad h \in (\operatorname{vec}(w) \to d) \wedge \\ \quad \forall y \in \operatorname{vec}(w).\; (i(h(y)) = y) \wedge \\ \quad \forall x \in d.\; (h(i(x)) = x) \Rightarrow \\ \quad \quad \operatorname{is\text{-}hilbert\text{-}space}(H_{{(\lambda {[}\operatorname{bga\_}, \operatorname{bgb\_}] \in (d \times d).\; (\langle\cdot,\cdot\rangle_{w}(i(\operatorname{bga\_}), i(\operatorname{bgb\_})) - \langle\cdot,\cdot\rangle_{v}(\operatorname{bga\_}, \operatorname{bgb\_})))}+}) \end{array}}}
\vnbprop{pb-closed}{\vnbstmt{\begin{array}{@{}l@{}} \forall v, d, w, i, h. \\ \quad \operatorname{is\text{-}hilbert\text{-}space}(v) \wedge \\ \quad \operatorname{is\text{-}hilbert\text{-}space}(w) \wedge \\ \quad \operatorname{is\text{-}submodule}(v, d) \wedge \\ \quad \operatorname{is\text{-}dense}(\operatorname{ip\text{-}metric\text{-}space}(v), d) \wedge \\ \quad i \in (d \to \operatorname{vec}(w)) \wedge \\ \quad \forall x, y \in d.\; (i(\operatorname{vadd}(v)(x, y)) = \operatorname{vadd}(w)(i(x), i(y))) \wedge \\ \quad \forall c \in \mathbb{C}, x \in d.\; (i(\operatorname{act}(v)(c, x)) = \operatorname{act}(w)(c, i(x))) \wedge \\ \quad \forall x \in d.\; (\langle\cdot,\cdot\rangle_{v}(x, x) \leq \langle\cdot,\cdot\rangle_{w}(i(x), i(x))) \wedge \\ \quad h \in (\operatorname{vec}(w) \to d) \wedge \\ \quad \forall y \in \operatorname{vec}(w).\; (i(h(y)) = y) \wedge \\ \quad \forall x \in d.\; (h(i(x)) = x) \Rightarrow \\ \quad \quad \operatorname{is\text{-}closed\text{-}form}(v, d, \\ \quad \quad \quad (\lambda {[}\operatorname{bga\_}, \operatorname{bgb\_}] \in (d \times d).\; (\langle\cdot,\cdot\rangle_{w}(i(\operatorname{bga\_}), i(\operatorname{bgb\_})) - \langle\cdot,\cdot\rangle_{v}(\operatorname{bga\_}, \operatorname{bgb\_})))) \end{array}}}
\vnbprop{fr-gil-type}{\vnbstmt{\begin{array}{@{}l@{}} \forall v, d, q, g, f. \\ \quad \operatorname{is\text{-}hilbert\text{-}space}(v) \wedge \\ \quad \operatorname{is\text{-}positive\text{-}form}(v, d, q) \wedge \\ \quad \operatorname{is\text{-}uniformly\text{-}continuous}(\operatorname{ip\text{-}metric\text{-}space}(\operatorname{hilbert\text{-}completion}(\operatorname{complex\text{-}inner\text{-}product\text{-}space\text{-}as\text{-}module}(\operatorname{sub\text{-}cips}(v, d)), {q}_{+})), \\ \quad \quad \operatorname{ip\text{-}metric\text{-}space}(v), g) \wedge \\ \quad \forall x \in d. \\ \quad \quad g(\operatorname{hc\text{-}embed}(\operatorname{complex\text{-}inner\text{-}product\text{-}space\text{-}as\text{-}module}(\operatorname{sub\text{-}cips}(v, d)), {q}_{+})(x)) = x \wedge \\ \quad f \in (d \to \operatorname{vec}(v)) \wedge \\ \quad \forall y, \mathit{btx} \in d.\; (q(\mathit{btx}, y) = \langle\cdot,\cdot\rangle_{v}(\mathit{btx}, f(y))) \Rightarrow \\ \quad \quad (\lambda i \in \operatorname{image}(g, \operatorname{vec}(\operatorname{hilbert\text{-}completion}(\operatorname{complex\text{-}inner\text{-}product\text{-}space\text{-}as\text{-}module}(\operatorname{sub\text{-}cips}(v, d)), {q}_{+}))).\; \\ \quad \quad \quad (\iota\, u.\;((u \in \operatorname{vec}(\operatorname{hilbert\text{-}completion}(\operatorname{complex\text{-}inner\text{-}product\text{-}space\text{-}as\text{-}module}(\operatorname{sub\text{-}cips}(v, d)), {q}_{+}))) \wedge (g(u) = i)))) \\ \quad \quad \in\; (\operatorname{image}(g, \operatorname{vec}(\operatorname{hilbert\text{-}completion}(\operatorname{complex\text{-}inner\text{-}product\text{-}space\text{-}as\text{-}module}(\operatorname{sub\text{-}cips}(v, d)), {q}_{+}))) \to \operatorname{vec}(\operatorname{hilbert\text{-}completion}(\operatorname{complex\text{-}inner\text{-}product\text{-}space\text{-}as\text{-}module}(\operatorname{sub\text{-}cips}(v, d)), {q}_{+}))) \end{array}}}
\vnbprop{fr-gil-value}{\vnbstmt{\begin{array}{@{}l@{}} \forall v, d, q, g, f, a. \\ \quad \operatorname{is\text{-}hilbert\text{-}space}(v) \wedge \\ \quad \operatorname{is\text{-}positive\text{-}form}(v, d, q) \wedge \\ \quad \operatorname{is\text{-}uniformly\text{-}continuous}(\operatorname{ip\text{-}metric\text{-}space}(\operatorname{hilbert\text{-}completion}(\operatorname{complex\text{-}inner\text{-}product\text{-}space\text{-}as\text{-}module}(\operatorname{sub\text{-}cips}(v, d)), {q}_{+})), \\ \quad \quad \operatorname{ip\text{-}metric\text{-}space}(v), g) \wedge \\ \quad \forall x \in d. \\ \quad \quad g(\operatorname{hc\text{-}embed}(\operatorname{complex\text{-}inner\text{-}product\text{-}space\text{-}as\text{-}module}(\operatorname{sub\text{-}cips}(v, d)), {q}_{+})(x)) = x \wedge \\ \quad f \in (d \to \operatorname{vec}(v)) \wedge \\ \quad \forall y, \mathit{btx} \in d.\; (q(\mathit{btx}, y) = \langle\cdot,\cdot\rangle_{v}(\mathit{btx}, f(y))) \wedge \\ \quad a \\ \quad \in\; \operatorname{image}(g, \\ \quad \quad \operatorname{vec}(\operatorname{hilbert\text{-}completion}(\operatorname{complex\text{-}inner\text{-}product\text{-}space\text{-}as\text{-}module}(\operatorname{sub\text{-}cips}(v, d)), {q}_{+}))) \Rightarrow \\ \quad \quad (\lambda i \in \operatorname{image}(g, \operatorname{vec}(\operatorname{hilbert\text{-}completion}(\operatorname{complex\text{-}inner\text{-}product\text{-}space\text{-}as\text{-}module}(\operatorname{sub\text{-}cips}(v, d)), {q}_{+}))).\; (\iota\, u.\;((u \in \operatorname{vec}(\operatorname{hilbert\text{-}completion}(\operatorname{complex\text{-}inner\text{-}product\text{-}space\text{-}as\text{-}module}(\operatorname{sub\text{-}cips}(v, d)), {q}_{+}))) \wedge (g(u) = i))))(a) \\ \quad \quad =\; (\iota\, u.\;((u \in \operatorname{vec}(\operatorname{hilbert\text{-}completion}(\operatorname{complex\text{-}inner\text{-}product\text{-}space\text{-}as\text{-}module}(\operatorname{sub\text{-}cips}(v, d)), {q}_{+}))) \wedge (g(u) = a))) \end{array}}}
\vnbprop{fr-gl-type}{\vnbstmt{\begin{array}{@{}l@{}} \forall v, d, q, g, f. \\ \quad \operatorname{is\text{-}hilbert\text{-}space}(v) \wedge \\ \quad \operatorname{is\text{-}positive\text{-}form}(v, d, q) \wedge \\ \quad \operatorname{is\text{-}uniformly\text{-}continuous}(\operatorname{ip\text{-}metric\text{-}space}(\operatorname{hilbert\text{-}completion}(\operatorname{complex\text{-}inner\text{-}product\text{-}space\text{-}as\text{-}module}(\operatorname{sub\text{-}cips}(v, d)), {q}_{+})), \\ \quad \quad \operatorname{ip\text{-}metric\text{-}space}(v), g) \wedge \\ \quad \forall x \in d. \\ \quad \quad g(\operatorname{hc\text{-}embed}(\operatorname{complex\text{-}inner\text{-}product\text{-}space\text{-}as\text{-}module}(\operatorname{sub\text{-}cips}(v, d)), {q}_{+})(x)) = x \wedge \\ \quad f \in (d \to \operatorname{vec}(v)) \wedge \\ \quad \forall y, \mathit{btx} \in d.\; (q(\mathit{btx}, y) = \langle\cdot,\cdot\rangle_{v}(\mathit{btx}, f(y))) \Rightarrow \\ \quad \quad (\lambda j \in \operatorname{vec}(\operatorname{hilbert\text{-}completion}(\operatorname{complex\text{-}inner\text{-}product\text{-}space\text{-}as\text{-}module}(\operatorname{sub\text{-}cips}(v, d)), {q}_{+})).\; \\ \quad \quad \quad g(j)) \\ \quad \quad \in\; (\operatorname{vec}(\operatorname{hilbert\text{-}completion}(\operatorname{complex\text{-}inner\text{-}product\text{-}space\text{-}as\text{-}module}(\operatorname{sub\text{-}cips}(v, d)), {q}_{+})) \to \operatorname{image}(g, \operatorname{vec}(\operatorname{hilbert\text{-}completion}(\operatorname{complex\text{-}inner\text{-}product\text{-}space\text{-}as\text{-}module}(\operatorname{sub\text{-}cips}(v, d)), {q}_{+})))) \end{array}}}
\vnbprop{fr-gl-value}{\vnbstmt{\begin{array}{@{}l@{}} \forall v, d, q, g, f, u. \\ \quad \operatorname{is\text{-}hilbert\text{-}space}(v) \wedge \\ \quad \operatorname{is\text{-}positive\text{-}form}(v, d, q) \wedge \\ \quad \operatorname{is\text{-}uniformly\text{-}continuous}(\operatorname{ip\text{-}metric\text{-}space}(\operatorname{hilbert\text{-}completion}(\operatorname{complex\text{-}inner\text{-}product\text{-}space\text{-}as\text{-}module}(\operatorname{sub\text{-}cips}(v, d)), {q}_{+})), \\ \quad \quad \operatorname{ip\text{-}metric\text{-}space}(v), g) \wedge \\ \quad \forall x \in d. \\ \quad \quad g(\operatorname{hc\text{-}embed}(\operatorname{complex\text{-}inner\text{-}product\text{-}space\text{-}as\text{-}module}(\operatorname{sub\text{-}cips}(v, d)), {q}_{+})(x)) = x \wedge \\ \quad f \in (d \to \operatorname{vec}(v)) \wedge \\ \quad \forall y, \mathit{btx} \in d.\; (q(\mathit{btx}, y) = \langle\cdot,\cdot\rangle_{v}(\mathit{btx}, f(y))) \wedge \\ \quad u \\ \quad \in\; \operatorname{vec}(\operatorname{hilbert\text{-}completion}(\operatorname{complex\text{-}inner\text{-}product\text{-}space\text{-}as\text{-}module}(\operatorname{sub\text{-}cips}(v, d)), {q}_{+})) \Rightarrow \\ \quad \quad (\lambda j \in \operatorname{vec}(\operatorname{hilbert\text{-}completion}(\operatorname{complex\text{-}inner\text{-}product\text{-}space\text{-}as\text{-}module}(\operatorname{sub\text{-}cips}(v, d)), {q}_{+})).\; g(j))(u) \\ \quad \quad =\; g(u) \end{array}}}
\vnbprop{fr-gil-add}{\vnbstmt{\begin{array}{@{}l@{}} \forall v, d, q, g, f.; \forall x, y \in \operatorname{image}(g, \operatorname{vec}(\operatorname{hilbert\text{-}completion}(\operatorname{complex\text{-}inner\text{-}product\text{-}space\text{-}as\text{-}module}(\operatorname{sub\text{-}cips}(v, d)), {q}_{+}))). \\ \quad \operatorname{is\text{-}hilbert\text{-}space}(v) \wedge \\ \quad \operatorname{is\text{-}positive\text{-}form}(v, d, q) \wedge \\ \quad \operatorname{is\text{-}uniformly\text{-}continuous}(\operatorname{ip\text{-}metric\text{-}space}(\operatorname{hilbert\text{-}completion}(\operatorname{complex\text{-}inner\text{-}product\text{-}space\text{-}as\text{-}module}(\operatorname{sub\text{-}cips}(v, d)), {q}_{+})), \\ \quad \quad \operatorname{ip\text{-}metric\text{-}space}(v), g) \wedge \\ \quad \forall \mathit{bex} \in d. \\ \quad \quad g(\operatorname{hc\text{-}embed}(\operatorname{complex\text{-}inner\text{-}product\text{-}space\text{-}as\text{-}module}(\operatorname{sub\text{-}cips}(v, d)), {q}_{+})(\mathit{bex})) = \mathit{bex} \wedge \\ \quad f \in (d \to \operatorname{vec}(v)) \wedge \\ \quad \forall \mathit{bty}, \mathit{btx} \in d.\; (q(\mathit{btx}, \mathit{bty}) = \langle\cdot,\cdot\rangle_{v}(\mathit{btx}, f(\mathit{bty}))) \Rightarrow \\ \quad \quad (\lambda i \in \operatorname{image}(g, \operatorname{vec}(\operatorname{hilbert\text{-}completion}(\operatorname{complex\text{-}inner\text{-}product\text{-}space\text{-}as\text{-}module}(\operatorname{sub\text{-}cips}(v, d)), {q}_{+}))).\; (\iota\, u.\;((u \in \operatorname{vec}(\operatorname{hilbert\text{-}completion}(\operatorname{complex\text{-}inner\text{-}product\text{-}space\text{-}as\text{-}module}(\operatorname{sub\text{-}cips}(v, d)), {q}_{+}))) \wedge (g(u) = i))))(\operatorname{vadd}(v)(x, y)) \\ \quad \quad =\; \operatorname{vadd}(\operatorname{hilbert\text{-}completion}(\operatorname{complex\text{-}inner\text{-}product\text{-}space\text{-}as\text{-}module}(\operatorname{sub\text{-}cips}(v, d)), {q}_{+}))((\lambda i \in \operatorname{image}(g, \operatorname{vec}(\operatorname{hilbert\text{-}completion}(\operatorname{complex\text{-}inner\text{-}product\text{-}space\text{-}as\text{-}module}(\operatorname{sub\text{-}cips}(v, d)), {q}_{+}))).\; (\iota\, u.\;((u \in \operatorname{vec}(\operatorname{hilbert\text{-}completion}(\operatorname{complex\text{-}inner\text{-}product\text{-}space\text{-}as\text{-}module}(\operatorname{sub\text{-}cips}(v, d)), {q}_{+}))) \wedge (g(u) = i))))(x), \\ \quad \quad \quad (\lambda i \in \operatorname{image}(g, \operatorname{vec}(\operatorname{hilbert\text{-}completion}(\operatorname{complex\text{-}inner\text{-}product\text{-}space\text{-}as\text{-}module}(\operatorname{sub\text{-}cips}(v, d)), {q}_{+}))).\; (\iota\, u.\;((u \in \operatorname{vec}(\operatorname{hilbert\text{-}completion}(\operatorname{complex\text{-}inner\text{-}product\text{-}space\text{-}as\text{-}module}(\operatorname{sub\text{-}cips}(v, d)), {q}_{+}))) \wedge (g(u) = i))))(y)) \end{array}}}
\vnbprop{fr-gil-hom}{\vnbstmt{\begin{array}{@{}l@{}} \forall v, d, q, g, f.; \forall c \in \mathbb{C}, x \in \operatorname{image}(g, \operatorname{vec}(\operatorname{hilbert\text{-}completion}(\operatorname{complex\text{-}inner\text{-}product\text{-}space\text{-}as\text{-}module}(\operatorname{sub\text{-}cips}(v, d)), {q}_{+}))). \\ \quad \operatorname{is\text{-}hilbert\text{-}space}(v) \wedge \\ \quad \operatorname{is\text{-}positive\text{-}form}(v, d, q) \wedge \\ \quad \operatorname{is\text{-}uniformly\text{-}continuous}(\operatorname{ip\text{-}metric\text{-}space}(\operatorname{hilbert\text{-}completion}(\operatorname{complex\text{-}inner\text{-}product\text{-}space\text{-}as\text{-}module}(\operatorname{sub\text{-}cips}(v, d)), {q}_{+})), \\ \quad \quad \operatorname{ip\text{-}metric\text{-}space}(v), g) \wedge \\ \quad \forall \mathit{bex} \in d. \\ \quad \quad g(\operatorname{hc\text{-}embed}(\operatorname{complex\text{-}inner\text{-}product\text{-}space\text{-}as\text{-}module}(\operatorname{sub\text{-}cips}(v, d)), {q}_{+})(\mathit{bex})) = \mathit{bex} \wedge \\ \quad f \in (d \to \operatorname{vec}(v)) \wedge \\ \quad \forall y, \mathit{btx} \in d.\; (q(\mathit{btx}, y) = \langle\cdot,\cdot\rangle_{v}(\mathit{btx}, f(y))) \Rightarrow \\ \quad \quad (\lambda i \in \operatorname{image}(g, \operatorname{vec}(\operatorname{hilbert\text{-}completion}(\operatorname{complex\text{-}inner\text{-}product\text{-}space\text{-}as\text{-}module}(\operatorname{sub\text{-}cips}(v, d)), {q}_{+}))).\; (\iota\, u.\;((u \in \operatorname{vec}(\operatorname{hilbert\text{-}completion}(\operatorname{complex\text{-}inner\text{-}product\text{-}space\text{-}as\text{-}module}(\operatorname{sub\text{-}cips}(v, d)), {q}_{+}))) \wedge (g(u) = i))))(\operatorname{act}(v)(c, x)) \\ \quad \quad =\; \operatorname{act}(\operatorname{hilbert\text{-}completion}(\operatorname{complex\text{-}inner\text{-}product\text{-}space\text{-}as\text{-}module}(\operatorname{sub\text{-}cips}(v, d)), {q}_{+}))(c, \\ \quad \quad \quad (\lambda i \in \operatorname{image}(g, \operatorname{vec}(\operatorname{hilbert\text{-}completion}(\operatorname{complex\text{-}inner\text{-}product\text{-}space\text{-}as\text{-}module}(\operatorname{sub\text{-}cips}(v, d)), {q}_{+}))).\; (\iota\, u.\;((u \in \operatorname{vec}(\operatorname{hilbert\text{-}completion}(\operatorname{complex\text{-}inner\text{-}product\text{-}space\text{-}as\text{-}module}(\operatorname{sub\text{-}cips}(v, d)), {q}_{+}))) \wedge (g(u) = i))))(x)) \end{array}}}
\vnbprop{fr-gil-bnd}{\vnbstmt{\begin{array}{@{}l@{}} \forall v, d, q, g, f.; \forall x \in \operatorname{image}(g, \operatorname{vec}(\operatorname{hilbert\text{-}completion}(\operatorname{complex\text{-}inner\text{-}product\text{-}space\text{-}as\text{-}module}(\operatorname{sub\text{-}cips}(v, d)), {q}_{+}))). \\ \quad \operatorname{is\text{-}hilbert\text{-}space}(v) \wedge \\ \quad \operatorname{is\text{-}positive\text{-}form}(v, d, q) \wedge \\ \quad \operatorname{is\text{-}uniformly\text{-}continuous}(\operatorname{ip\text{-}metric\text{-}space}(\operatorname{hilbert\text{-}completion}(\operatorname{complex\text{-}inner\text{-}product\text{-}space\text{-}as\text{-}module}(\operatorname{sub\text{-}cips}(v, d)), {q}_{+})), \\ \quad \quad \operatorname{ip\text{-}metric\text{-}space}(v), g) \wedge \\ \quad \forall \mathit{bex} \in d. \\ \quad \quad g(\operatorname{hc\text{-}embed}(\operatorname{complex\text{-}inner\text{-}product\text{-}space\text{-}as\text{-}module}(\operatorname{sub\text{-}cips}(v, d)), {q}_{+})(\mathit{bex})) = \mathit{bex} \wedge \\ \quad f \in (d \to \operatorname{vec}(v)) \wedge \\ \quad \forall y, \mathit{btx} \in d.\; (q(\mathit{btx}, y) = \langle\cdot,\cdot\rangle_{v}(\mathit{btx}, f(y))) \Rightarrow \\ \quad \quad \langle\cdot,\cdot\rangle_{v}(x, x) \\ \quad \quad \leq\; \langle\cdot,\cdot\rangle_{\operatorname{hilbert\text{-}completion}(\operatorname{complex\text{-}inner\text{-}product\text{-}space\text{-}as\text{-}module}(\operatorname{sub\text{-}cips}(v, d)), {q}_{+})}((\lambda i \in \operatorname{image}(g, \operatorname{vec}(\operatorname{hilbert\text{-}completion}(\operatorname{complex\text{-}inner\text{-}product\text{-}space\text{-}as\text{-}module}(\operatorname{sub\text{-}cips}(v, d)), {q}_{+}))).\; (\iota\, u.\;((u \in \operatorname{vec}(\operatorname{hilbert\text{-}completion}(\operatorname{complex\text{-}inner\text{-}product\text{-}space\text{-}as\text{-}module}(\operatorname{sub\text{-}cips}(v, d)), {q}_{+}))) \wedge (g(u) = i))))(x), \\ \quad \quad \quad (\lambda i \in \operatorname{image}(g, \operatorname{vec}(\operatorname{hilbert\text{-}completion}(\operatorname{complex\text{-}inner\text{-}product\text{-}space\text{-}as\text{-}module}(\operatorname{sub\text{-}cips}(v, d)), {q}_{+}))).\; (\iota\, u.\;((u \in \operatorname{vec}(\operatorname{hilbert\text{-}completion}(\operatorname{complex\text{-}inner\text{-}product\text{-}space\text{-}as\text{-}module}(\operatorname{sub\text{-}cips}(v, d)), {q}_{+}))) \wedge (g(u) = i))))(x)) \end{array}}}
\vnbprop{fr-hl}{\vnbstmt{\begin{array}{@{}l@{}} \forall v, d, q, g, f.; \forall y \in \operatorname{vec}(\operatorname{hilbert\text{-}completion}(\operatorname{complex\text{-}inner\text{-}product\text{-}space\text{-}as\text{-}module}(\operatorname{sub\text{-}cips}(v, d)), {q}_{+})). \\ \quad \operatorname{is\text{-}hilbert\text{-}space}(v) \wedge \\ \quad \operatorname{is\text{-}positive\text{-}form}(v, d, q) \wedge \\ \quad \operatorname{is\text{-}uniformly\text{-}continuous}(\operatorname{ip\text{-}metric\text{-}space}(\operatorname{hilbert\text{-}completion}(\operatorname{complex\text{-}inner\text{-}product\text{-}space\text{-}as\text{-}module}(\operatorname{sub\text{-}cips}(v, d)), {q}_{+})), \\ \quad \quad \operatorname{ip\text{-}metric\text{-}space}(v), g) \wedge \\ \quad \forall x \in d. \\ \quad \quad g(\operatorname{hc\text{-}embed}(\operatorname{complex\text{-}inner\text{-}product\text{-}space\text{-}as\text{-}module}(\operatorname{sub\text{-}cips}(v, d)), {q}_{+})(x)) = x \wedge \\ \quad f \in (d \to \operatorname{vec}(v)) \wedge \\ \quad \forall \mathit{bty}, \mathit{btx} \in d.\; (q(\mathit{btx}, \mathit{bty}) = \langle\cdot,\cdot\rangle_{v}(\mathit{btx}, f(\mathit{bty}))) \Rightarrow \\ \quad \quad (\lambda i \in \operatorname{image}(g, \operatorname{vec}(\operatorname{hilbert\text{-}completion}(\operatorname{complex\text{-}inner\text{-}product\text{-}space\text{-}as\text{-}module}(\operatorname{sub\text{-}cips}(v, d)), {q}_{+}))).\; (\iota\, u.\;((u \in \operatorname{vec}(\operatorname{hilbert\text{-}completion}(\operatorname{complex\text{-}inner\text{-}product\text{-}space\text{-}as\text{-}module}(\operatorname{sub\text{-}cips}(v, d)), {q}_{+}))) \wedge (g(u) = i))))((\lambda j \in \operatorname{vec}(\operatorname{hilbert\text{-}completion}(\operatorname{complex\text{-}inner\text{-}product\text{-}space\text{-}as\text{-}module}(\operatorname{sub\text{-}cips}(v, d)), {q}_{+})).\; g(j))(y)) \\ \quad \quad =\; y \end{array}}}
\vnbprop{fr-hr}{\vnbstmt{\begin{array}{@{}l@{}} \forall v, d, q, g, f.; \forall x \in \operatorname{image}(g, \operatorname{vec}(\operatorname{hilbert\text{-}completion}(\operatorname{complex\text{-}inner\text{-}product\text{-}space\text{-}as\text{-}module}(\operatorname{sub\text{-}cips}(v, d)), {q}_{+}))). \\ \quad \operatorname{is\text{-}hilbert\text{-}space}(v) \wedge \\ \quad \operatorname{is\text{-}positive\text{-}form}(v, d, q) \wedge \\ \quad \operatorname{is\text{-}uniformly\text{-}continuous}(\operatorname{ip\text{-}metric\text{-}space}(\operatorname{hilbert\text{-}completion}(\operatorname{complex\text{-}inner\text{-}product\text{-}space\text{-}as\text{-}module}(\operatorname{sub\text{-}cips}(v, d)), {q}_{+})), \\ \quad \quad \operatorname{ip\text{-}metric\text{-}space}(v), g) \wedge \\ \quad \forall \mathit{bex} \in d. \\ \quad \quad g(\operatorname{hc\text{-}embed}(\operatorname{complex\text{-}inner\text{-}product\text{-}space\text{-}as\text{-}module}(\operatorname{sub\text{-}cips}(v, d)), {q}_{+})(\mathit{bex})) = \mathit{bex} \wedge \\ \quad f \in (d \to \operatorname{vec}(v)) \wedge \\ \quad \forall y, \mathit{btx} \in d.\; (q(\mathit{btx}, y) = \langle\cdot,\cdot\rangle_{v}(\mathit{btx}, f(y))) \Rightarrow \\ \quad \quad (\lambda j \in \operatorname{vec}(\operatorname{hilbert\text{-}completion}(\operatorname{complex\text{-}inner\text{-}product\text{-}space\text{-}as\text{-}module}(\operatorname{sub\text{-}cips}(v, d)), {q}_{+})).\; g(j))((\lambda i \in \operatorname{image}(g, \operatorname{vec}(\operatorname{hilbert\text{-}completion}(\operatorname{complex\text{-}inner\text{-}product\text{-}space\text{-}as\text{-}module}(\operatorname{sub\text{-}cips}(v, d)), {q}_{+}))).\; (\iota\, u.\;((u \in \operatorname{vec}(\operatorname{hilbert\text{-}completion}(\operatorname{complex\text{-}inner\text{-}product\text{-}space\text{-}as\text{-}module}(\operatorname{sub\text{-}cips}(v, d)), {q}_{+}))) \wedge (g(u) = i))))(x)) \\ \quad \quad =\; x \end{array}}}
\vnbprop{fr-q1-closed}{\vnbstmt{\begin{array}{@{}l@{}} \forall v, d, q, g, f. \\ \quad \operatorname{is\text{-}hilbert\text{-}space}(v) \wedge \\ \quad \operatorname{is\text{-}positive\text{-}form}(v, d, q) \wedge \\ \quad \operatorname{is\text{-}uniformly\text{-}continuous}(\operatorname{ip\text{-}metric\text{-}space}(\operatorname{hilbert\text{-}completion}(\operatorname{complex\text{-}inner\text{-}product\text{-}space\text{-}as\text{-}module}(\operatorname{sub\text{-}cips}(v, d)), {q}_{+})), \\ \quad \quad \operatorname{ip\text{-}metric\text{-}space}(v), g) \wedge \\ \quad \forall x \in d. \\ \quad \quad g(\operatorname{hc\text{-}embed}(\operatorname{complex\text{-}inner\text{-}product\text{-}space\text{-}as\text{-}module}(\operatorname{sub\text{-}cips}(v, d)), {q}_{+})(x)) = x \wedge \\ \quad f \in (d \to \operatorname{vec}(v)) \wedge \\ \quad \forall y, \mathit{btx} \in d.\; (q(\mathit{btx}, y) = \langle\cdot,\cdot\rangle_{v}(\mathit{btx}, f(y))) \Rightarrow \\ \quad \quad \operatorname{is\text{-}closed\text{-}form}(v, \\ \quad \quad \quad \operatorname{image}(g, \\ \quad \quad \quad \quad \operatorname{vec}(\operatorname{hilbert\text{-}completion}(\operatorname{complex\text{-}inner\text{-}product\text{-}space\text{-}as\text{-}module}(\operatorname{sub\text{-}cips}(v, d)), {q}_{+}))), \\ \quad \quad \quad (\lambda {[}\operatorname{bga\_}, \operatorname{bgb\_}] \in (\operatorname{image}(g, \operatorname{vec}(\operatorname{hilbert\text{-}completion}(\operatorname{complex\text{-}inner\text{-}product\text{-}space\text{-}as\text{-}module}(\operatorname{sub\text{-}cips}(v, d)), {q}_{+}))) \times \operatorname{image}(g, \operatorname{vec}(\operatorname{hilbert\text{-}completion}(\operatorname{complex\text{-}inner\text{-}product\text{-}space\text{-}as\text{-}module}(\operatorname{sub\text{-}cips}(v, d)), {q}_{+})))).\; \\ \quad \quad \quad \quad (\langle\cdot,\cdot\rangle_{\operatorname{hilbert\text{-}completion}(\operatorname{complex\text{-}inner\text{-}product\text{-}space\text{-}as\text{-}module}(\operatorname{sub\text{-}cips}(v, d)), {q}_{+})}((\lambda i \in \operatorname{image}(g, \operatorname{vec}(\operatorname{hilbert\text{-}completion}(\operatorname{complex\text{-}inner\text{-}product\text{-}space\text{-}as\text{-}module}(\operatorname{sub\text{-}cips}(v, d)), {q}_{+}))).\; (\iota\, u.\;((u \in \operatorname{vec}(\operatorname{hilbert\text{-}completion}(\operatorname{complex\text{-}inner\text{-}product\text{-}space\text{-}as\text{-}module}(\operatorname{sub\text{-}cips}(v, d)), {q}_{+}))) \wedge (g(u) = i))))(\operatorname{bga\_}), (\lambda i \in \operatorname{image}(g, \operatorname{vec}(\operatorname{hilbert\text{-}completion}(\operatorname{complex\text{-}inner\text{-}product\text{-}space\text{-}as\text{-}module}(\operatorname{sub\text{-}cips}(v, d)), {q}_{+}))).\; (\iota\, u.\;((u \in \operatorname{vec}(\operatorname{hilbert\text{-}completion}(\operatorname{complex\text{-}inner\text{-}product\text{-}space\text{-}as\text{-}module}(\operatorname{sub\text{-}cips}(v, d)), {q}_{+}))) \wedge (g(u) = i))))(\operatorname{bgb\_})) - \langle\cdot,\cdot\rangle_{v}(\operatorname{bga\_}, \operatorname{bgb\_})))) \end{array}}}
\vnbprop{fr-q1-extends}{\vnbstmt{\begin{array}{@{}l@{}} \forall v, d, q, g, f, x, y. \\ \quad \operatorname{is\text{-}hilbert\text{-}space}(v) \wedge \\ \quad \operatorname{is\text{-}positive\text{-}form}(v, d, q) \wedge \\ \quad \operatorname{is\text{-}uniformly\text{-}continuous}(\operatorname{ip\text{-}metric\text{-}space}(\operatorname{hilbert\text{-}completion}(\operatorname{complex\text{-}inner\text{-}product\text{-}space\text{-}as\text{-}module}(\operatorname{sub\text{-}cips}(v, d)), {q}_{+})), \\ \quad \quad \operatorname{ip\text{-}metric\text{-}space}(v), g) \wedge \\ \quad \forall \mathit{bex} \in d. \\ \quad \quad g(\operatorname{hc\text{-}embed}(\operatorname{complex\text{-}inner\text{-}product\text{-}space\text{-}as\text{-}module}(\operatorname{sub\text{-}cips}(v, d)), {q}_{+})(\mathit{bex})) = \mathit{bex} \wedge \\ \quad f \in (d \to \operatorname{vec}(v)) \wedge \\ \quad \forall \mathit{bty}, \mathit{btx} \in d.\; (q(\mathit{btx}, \mathit{bty}) = \langle\cdot,\cdot\rangle_{v}(\mathit{btx}, f(\mathit{bty}))) \wedge \\ \quad x \in d \wedge \\ \quad y \in d \Rightarrow \\ \quad \quad (\lambda {[}\operatorname{bga\_}, \operatorname{bgb\_}] \in (\operatorname{image}(g, \operatorname{vec}(\operatorname{hilbert\text{-}completion}(\operatorname{complex\text{-}inner\text{-}product\text{-}space\text{-}as\text{-}module}(\operatorname{sub\text{-}cips}(v, d)), {q}_{+}))) \times \operatorname{image}(g, \operatorname{vec}(\operatorname{hilbert\text{-}completion}(\operatorname{complex\text{-}inner\text{-}product\text{-}space\text{-}as\text{-}module}(\operatorname{sub\text{-}cips}(v, d)), {q}_{+})))).\; (\langle\cdot,\cdot\rangle_{\operatorname{hilbert\text{-}completion}(\operatorname{complex\text{-}inner\text{-}product\text{-}space\text{-}as\text{-}module}(\operatorname{sub\text{-}cips}(v, d)), {q}_{+})}((\lambda i \in \operatorname{image}(g, \operatorname{vec}(\operatorname{hilbert\text{-}completion}(\operatorname{complex\text{-}inner\text{-}product\text{-}space\text{-}as\text{-}module}(\operatorname{sub\text{-}cips}(v, d)), {q}_{+}))).\; (\iota\, u.\;((u \in \operatorname{vec}(\operatorname{hilbert\text{-}completion}(\operatorname{complex\text{-}inner\text{-}product\text{-}space\text{-}as\text{-}module}(\operatorname{sub\text{-}cips}(v, d)), {q}_{+}))) \wedge (g(u) = i))))(\operatorname{bga\_}), (\lambda i \in \operatorname{image}(g, \operatorname{vec}(\operatorname{hilbert\text{-}completion}(\operatorname{complex\text{-}inner\text{-}product\text{-}space\text{-}as\text{-}module}(\operatorname{sub\text{-}cips}(v, d)), {q}_{+}))).\; (\iota\, u.\;((u \in \operatorname{vec}(\operatorname{hilbert\text{-}completion}(\operatorname{complex\text{-}inner\text{-}product\text{-}space\text{-}as\text{-}module}(\operatorname{sub\text{-}cips}(v, d)), {q}_{+}))) \wedge (g(u) = i))))(\operatorname{bgb\_})) - \langle\cdot,\cdot\rangle_{v}(\operatorname{bga\_}, \operatorname{bgb\_})))(x, y) \\ \quad \quad =\; q(x, y) \end{array}}}
\vnbprop{fr-dd-subset-d1}{\vnbstmt{\begin{array}{@{}l@{}} \forall v, d, q, g, f. \\ \quad \operatorname{is\text{-}hilbert\text{-}space}(v) \wedge \\ \quad \operatorname{is\text{-}positive\text{-}form}(v, d, q) \wedge \\ \quad \operatorname{is\text{-}uniformly\text{-}continuous}(\operatorname{ip\text{-}metric\text{-}space}(\operatorname{hilbert\text{-}completion}(\operatorname{complex\text{-}inner\text{-}product\text{-}space\text{-}as\text{-}module}(\operatorname{sub\text{-}cips}(v, d)), {q}_{+})), \\ \quad \quad \operatorname{ip\text{-}metric\text{-}space}(v), g) \wedge \\ \quad \forall x \in d. \\ \quad \quad g(\operatorname{hc\text{-}embed}(\operatorname{complex\text{-}inner\text{-}product\text{-}space\text{-}as\text{-}module}(\operatorname{sub\text{-}cips}(v, d)), {q}_{+})(x)) = x \wedge \\ \quad f \in (d \to \operatorname{vec}(v)) \wedge \\ \quad \forall y, \mathit{btx} \in d.\; (q(\mathit{btx}, y) = \langle\cdot,\cdot\rangle_{v}(\mathit{btx}, f(y))) \Rightarrow \\ \quad \quad d \\ \quad \quad \subseteq\; \operatorname{image}(g, \\ \quad \quad \quad \operatorname{vec}(\operatorname{hilbert\text{-}completion}(\operatorname{complex\text{-}inner\text{-}product\text{-}space\text{-}as\text{-}module}(\operatorname{sub\text{-}cips}(v, d)), {q}_{+}))) \end{array}}}
\vnbprop{fr-graph-in-t1}{\vnbstmt{\begin{array}{@{}l@{}} \forall v, d, q, g, f, y. \\ \quad \operatorname{is\text{-}hilbert\text{-}space}(v) \wedge \\ \quad \operatorname{is\text{-}positive\text{-}form}(v, d, q) \wedge \\ \quad \operatorname{is\text{-}uniformly\text{-}continuous}(\operatorname{ip\text{-}metric\text{-}space}(\operatorname{hilbert\text{-}completion}(\operatorname{complex\text{-}inner\text{-}product\text{-}space\text{-}as\text{-}module}(\operatorname{sub\text{-}cips}(v, d)), {q}_{+})), \\ \quad \quad \operatorname{ip\text{-}metric\text{-}space}(v), g) \wedge \\ \quad \forall x \in d. \\ \quad \quad g(\operatorname{hc\text{-}embed}(\operatorname{complex\text{-}inner\text{-}product\text{-}space\text{-}as\text{-}module}(\operatorname{sub\text{-}cips}(v, d)), {q}_{+})(x)) = x \wedge \\ \quad f \in (d \to \operatorname{vec}(v)) \wedge \\ \quad \forall \mathit{bty}, \mathit{btx} \in d.\; (q(\mathit{btx}, \mathit{bty}) = \langle\cdot,\cdot\rangle_{v}(\mathit{btx}, f(\mathit{bty}))) \wedge \\ \quad y \in d \Rightarrow \\ \quad \quad {[}y, f(y)] \\ \quad \quad \in\; T_{(\lambda {[}\operatorname{bga\_}, \operatorname{bgb\_}] \in (\operatorname{image}(g, \operatorname{vec}(\operatorname{hilbert\text{-}completion}(\operatorname{complex\text{-}inner\text{-}product\text{-}space\text{-}as\text{-}module}(\operatorname{sub\text{-}cips}(v, d)), {q}_{+}))) \times \operatorname{image}(g, \operatorname{vec}(\operatorname{hilbert\text{-}completion}(\operatorname{complex\text{-}inner\text{-}product\text{-}space\text{-}as\text{-}module}(\operatorname{sub\text{-}cips}(v, d)), {q}_{+})))).\; (\langle\cdot,\cdot\rangle_{\operatorname{hilbert\text{-}completion}(\operatorname{complex\text{-}inner\text{-}product\text{-}space\text{-}as\text{-}module}(\operatorname{sub\text{-}cips}(v, d)), {q}_{+})}((\lambda i \in \operatorname{image}(g, \operatorname{vec}(\operatorname{hilbert\text{-}completion}(\operatorname{complex\text{-}inner\text{-}product\text{-}space\text{-}as\text{-}module}(\operatorname{sub\text{-}cips}(v, d)), {q}_{+}))).\; (\iota\, u.\;((u \in \operatorname{vec}(\operatorname{hilbert\text{-}completion}(\operatorname{complex\text{-}inner\text{-}product\text{-}space\text{-}as\text{-}module}(\operatorname{sub\text{-}cips}(v, d)), {q}_{+}))) \wedge (g(u) = i))))(\operatorname{bga\_}), (\lambda i \in \operatorname{image}(g, \operatorname{vec}(\operatorname{hilbert\text{-}completion}(\operatorname{complex\text{-}inner\text{-}product\text{-}space\text{-}as\text{-}module}(\operatorname{sub\text{-}cips}(v, d)), {q}_{+}))).\; (\iota\, u.\;((u \in \operatorname{vec}(\operatorname{hilbert\text{-}completion}(\operatorname{complex\text{-}inner\text{-}product\text{-}space\text{-}as\text{-}module}(\operatorname{sub\text{-}cips}(v, d)), {q}_{+}))) \wedge (g(u) = i))))(\operatorname{bgb\_})) - \langle\cdot,\cdot\rangle_{v}(\operatorname{bga\_}, \operatorname{bgb\_})))} \end{array}}}
\vnbprop{fr-t1-self-adjoint}{\vnbstmt{\begin{array}{@{}l@{}} \forall v, d, q, g, f. \\ \quad \operatorname{is\text{-}hilbert\text{-}space}(v) \wedge \\ \quad \operatorname{is\text{-}positive\text{-}form}(v, d, q) \wedge \\ \quad \operatorname{is\text{-}uniformly\text{-}continuous}(\operatorname{ip\text{-}metric\text{-}space}(\operatorname{hilbert\text{-}completion}(\operatorname{complex\text{-}inner\text{-}product\text{-}space\text{-}as\text{-}module}(\operatorname{sub\text{-}cips}(v, d)), {q}_{+})), \\ \quad \quad \operatorname{ip\text{-}metric\text{-}space}(v), g) \wedge \\ \quad \forall x \in d. \\ \quad \quad g(\operatorname{hc\text{-}embed}(\operatorname{complex\text{-}inner\text{-}product\text{-}space\text{-}as\text{-}module}(\operatorname{sub\text{-}cips}(v, d)), {q}_{+})(x)) = x \wedge \\ \quad f \in (d \to \operatorname{vec}(v)) \wedge \\ \quad \forall y, \mathit{btx} \in d.\; (q(\mathit{btx}, y) = \langle\cdot,\cdot\rangle_{v}(\mathit{btx}, f(y))) \Rightarrow \\ \quad \quad \operatorname{is\text{-}self\text{-}adjoint\text{-}graph}(v, \\ \quad \quad \quad T_{(\lambda {[}\operatorname{bga\_}, \operatorname{bgb\_}] \in (\operatorname{image}(g, \operatorname{vec}(\operatorname{hilbert\text{-}completion}(\operatorname{complex\text{-}inner\text{-}product\text{-}space\text{-}as\text{-}module}(\operatorname{sub\text{-}cips}(v, d)), {q}_{+}))) \times \operatorname{image}(g, \operatorname{vec}(\operatorname{hilbert\text{-}completion}(\operatorname{complex\text{-}inner\text{-}product\text{-}space\text{-}as\text{-}module}(\operatorname{sub\text{-}cips}(v, d)), {q}_{+})))).\; (\langle\cdot,\cdot\rangle_{\operatorname{hilbert\text{-}completion}(\operatorname{complex\text{-}inner\text{-}product\text{-}space\text{-}as\text{-}module}(\operatorname{sub\text{-}cips}(v, d)), {q}_{+})}((\lambda i \in \operatorname{image}(g, \operatorname{vec}(\operatorname{hilbert\text{-}completion}(\operatorname{complex\text{-}inner\text{-}product\text{-}space\text{-}as\text{-}module}(\operatorname{sub\text{-}cips}(v, d)), {q}_{+}))).\; (\iota\, u.\;((u \in \operatorname{vec}(\operatorname{hilbert\text{-}completion}(\operatorname{complex\text{-}inner\text{-}product\text{-}space\text{-}as\text{-}module}(\operatorname{sub\text{-}cips}(v, d)), {q}_{+}))) \wedge (g(u) = i))))(\operatorname{bga\_}), (\lambda i \in \operatorname{image}(g, \operatorname{vec}(\operatorname{hilbert\text{-}completion}(\operatorname{complex\text{-}inner\text{-}product\text{-}space\text{-}as\text{-}module}(\operatorname{sub\text{-}cips}(v, d)), {q}_{+}))).\; (\iota\, u.\;((u \in \operatorname{vec}(\operatorname{hilbert\text{-}completion}(\operatorname{complex\text{-}inner\text{-}product\text{-}space\text{-}as\text{-}module}(\operatorname{sub\text{-}cips}(v, d)), {q}_{+}))) \wedge (g(u) = i))))(\operatorname{bgb\_})) - \langle\cdot,\cdot\rangle_{v}(\operatorname{bga\_}, \operatorname{bgb\_})))}) \end{array}}}
\vnbprop{fr-near-plus}{\vnbstmt{\begin{array}{@{}l@{}} \forall v, d, q, g, f, x, e. \\ \quad \operatorname{is\text{-}hilbert\text{-}space}(v) \wedge \\ \quad \operatorname{is\text{-}positive\text{-}form}(v, d, q) \wedge \\ \quad \operatorname{is\text{-}uniformly\text{-}continuous}(\operatorname{ip\text{-}metric\text{-}space}(\operatorname{hilbert\text{-}completion}(\operatorname{complex\text{-}inner\text{-}product\text{-}space\text{-}as\text{-}module}(\operatorname{sub\text{-}cips}(v, d)), {q}_{+})), \\ \quad \quad \operatorname{ip\text{-}metric\text{-}space}(v), g) \wedge \\ \quad \forall \mathit{bex} \in d. \\ \quad \quad g(\operatorname{hc\text{-}embed}(\operatorname{complex\text{-}inner\text{-}product\text{-}space\text{-}as\text{-}module}(\operatorname{sub\text{-}cips}(v, d)), {q}_{+})(\mathit{bex})) = \mathit{bex} \wedge \\ \quad f \in (d \to \operatorname{vec}(v)) \wedge \\ \quad \forall y, \mathit{btx} \in d.\; (q(\mathit{btx}, y) = \langle\cdot,\cdot\rangle_{v}(\mathit{btx}, f(y))) \wedge \\ \quad x \\ \quad \in\; \operatorname{pts}(\operatorname{ip\text{-}metric\text{-}space}(H_{{(\lambda {[}\operatorname{bga\_}, \operatorname{bgb\_}] \in (\operatorname{image}(g, \operatorname{vec}(\operatorname{hilbert\text{-}completion}(\operatorname{complex\text{-}inner\text{-}product\text{-}space\text{-}as\text{-}module}(\operatorname{sub\text{-}cips}(v, d)), {q}_{+}))) \times \operatorname{image}(g, \operatorname{vec}(\operatorname{hilbert\text{-}completion}(\operatorname{complex\text{-}inner\text{-}product\text{-}space\text{-}as\text{-}module}(\operatorname{sub\text{-}cips}(v, d)), {q}_{+})))).\; (\langle\cdot,\cdot\rangle_{\operatorname{hilbert\text{-}completion}(\operatorname{complex\text{-}inner\text{-}product\text{-}space\text{-}as\text{-}module}(\operatorname{sub\text{-}cips}(v, d)), {q}_{+})}((\lambda i \in \operatorname{image}(g, \operatorname{vec}(\operatorname{hilbert\text{-}completion}(\operatorname{complex\text{-}inner\text{-}product\text{-}space\text{-}as\text{-}module}(\operatorname{sub\text{-}cips}(v, d)), {q}_{+}))).\; (\iota\, u.\;((u \in \operatorname{vec}(\operatorname{hilbert\text{-}completion}(\operatorname{complex\text{-}inner\text{-}product\text{-}space\text{-}as\text{-}module}(\operatorname{sub\text{-}cips}(v, d)), {q}_{+}))) \wedge (g(u) = i))))(\operatorname{bga\_}), (\lambda i \in \operatorname{image}(g, \operatorname{vec}(\operatorname{hilbert\text{-}completion}(\operatorname{complex\text{-}inner\text{-}product\text{-}space\text{-}as\text{-}module}(\operatorname{sub\text{-}cips}(v, d)), {q}_{+}))).\; (\iota\, u.\;((u \in \operatorname{vec}(\operatorname{hilbert\text{-}completion}(\operatorname{complex\text{-}inner\text{-}product\text{-}space\text{-}as\text{-}module}(\operatorname{sub\text{-}cips}(v, d)), {q}_{+}))) \wedge (g(u) = i))))(\operatorname{bgb\_})) - \langle\cdot,\cdot\rangle_{v}(\operatorname{bga\_}, \operatorname{bgb\_})))}+})) \wedge \\ \quad \operatorname{pos\text{-}rr}(e) \Rightarrow \\ \quad \quad \exists a \in d. \\ \quad \quad \quad \operatorname{dist}(\operatorname{ip\text{-}metric\text{-}space}(H_{{(\lambda {[}\operatorname{bga\_}, \operatorname{bgb\_}] \in (\operatorname{image}(g, \operatorname{vec}(\operatorname{hilbert\text{-}completion}(\operatorname{complex\text{-}inner\text{-}product\text{-}space\text{-}as\text{-}module}(\operatorname{sub\text{-}cips}(v, d)), {q}_{+}))) \times \operatorname{image}(g, \operatorname{vec}(\operatorname{hilbert\text{-}completion}(\operatorname{complex\text{-}inner\text{-}product\text{-}space\text{-}as\text{-}module}(\operatorname{sub\text{-}cips}(v, d)), {q}_{+})))).\; (\langle\cdot,\cdot\rangle_{\operatorname{hilbert\text{-}completion}(\operatorname{complex\text{-}inner\text{-}product\text{-}space\text{-}as\text{-}module}(\operatorname{sub\text{-}cips}(v, d)), {q}_{+})}((\lambda i \in \operatorname{image}(g, \operatorname{vec}(\operatorname{hilbert\text{-}completion}(\operatorname{complex\text{-}inner\text{-}product\text{-}space\text{-}as\text{-}module}(\operatorname{sub\text{-}cips}(v, d)), {q}_{+}))).\; (\iota\, u.\;((u \in \operatorname{vec}(\operatorname{hilbert\text{-}completion}(\operatorname{complex\text{-}inner\text{-}product\text{-}space\text{-}as\text{-}module}(\operatorname{sub\text{-}cips}(v, d)), {q}_{+}))) \wedge (g(u) = i))))(\operatorname{bga\_}), (\lambda i \in \operatorname{image}(g, \operatorname{vec}(\operatorname{hilbert\text{-}completion}(\operatorname{complex\text{-}inner\text{-}product\text{-}space\text{-}as\text{-}module}(\operatorname{sub\text{-}cips}(v, d)), {q}_{+}))).\; (\iota\, u.\;((u \in \operatorname{vec}(\operatorname{hilbert\text{-}completion}(\operatorname{complex\text{-}inner\text{-}product\text{-}space\text{-}as\text{-}module}(\operatorname{sub\text{-}cips}(v, d)), {q}_{+}))) \wedge (g(u) = i))))(\operatorname{bgb\_})) - \langle\cdot,\cdot\rangle_{v}(\operatorname{bga\_}, \operatorname{bgb\_})))}+}))(x, a) \\ \quad \quad \quad <\; e \end{array}}}
\vnbprop{fr-dd-dense-plus}{\vnbstmt{\begin{array}{@{}l@{}} \forall v, d, q, g, f. \\ \quad \operatorname{is\text{-}hilbert\text{-}space}(v) \wedge \\ \quad \operatorname{is\text{-}positive\text{-}form}(v, d, q) \wedge \\ \quad \operatorname{is\text{-}uniformly\text{-}continuous}(\operatorname{ip\text{-}metric\text{-}space}(\operatorname{hilbert\text{-}completion}(\operatorname{complex\text{-}inner\text{-}product\text{-}space\text{-}as\text{-}module}(\operatorname{sub\text{-}cips}(v, d)), {q}_{+})), \\ \quad \quad \operatorname{ip\text{-}metric\text{-}space}(v), g) \wedge \\ \quad \forall x \in d. \\ \quad \quad g(\operatorname{hc\text{-}embed}(\operatorname{complex\text{-}inner\text{-}product\text{-}space\text{-}as\text{-}module}(\operatorname{sub\text{-}cips}(v, d)), {q}_{+})(x)) = x \wedge \\ \quad f \in (d \to \operatorname{vec}(v)) \wedge \\ \quad \forall y, \mathit{btx} \in d.\; (q(\mathit{btx}, y) = \langle\cdot,\cdot\rangle_{v}(\mathit{btx}, f(y))) \Rightarrow \\ \quad \quad \operatorname{is\text{-}dense}(\operatorname{ip\text{-}metric\text{-}space}(H_{{(\lambda {[}\operatorname{bga\_}, \operatorname{bgb\_}] \in (\operatorname{image}(g, \operatorname{vec}(\operatorname{hilbert\text{-}completion}(\operatorname{complex\text{-}inner\text{-}product\text{-}space\text{-}as\text{-}module}(\operatorname{sub\text{-}cips}(v, d)), {q}_{+}))) \times \operatorname{image}(g, \operatorname{vec}(\operatorname{hilbert\text{-}completion}(\operatorname{complex\text{-}inner\text{-}product\text{-}space\text{-}as\text{-}module}(\operatorname{sub\text{-}cips}(v, d)), {q}_{+})))).\; (\langle\cdot,\cdot\rangle_{\operatorname{hilbert\text{-}completion}(\operatorname{complex\text{-}inner\text{-}product\text{-}space\text{-}as\text{-}module}(\operatorname{sub\text{-}cips}(v, d)), {q}_{+})}((\lambda i \in \operatorname{image}(g, \operatorname{vec}(\operatorname{hilbert\text{-}completion}(\operatorname{complex\text{-}inner\text{-}product\text{-}space\text{-}as\text{-}module}(\operatorname{sub\text{-}cips}(v, d)), {q}_{+}))).\; (\iota\, u.\;((u \in \operatorname{vec}(\operatorname{hilbert\text{-}completion}(\operatorname{complex\text{-}inner\text{-}product\text{-}space\text{-}as\text{-}module}(\operatorname{sub\text{-}cips}(v, d)), {q}_{+}))) \wedge (g(u) = i))))(\operatorname{bga\_}), (\lambda i \in \operatorname{image}(g, \operatorname{vec}(\operatorname{hilbert\text{-}completion}(\operatorname{complex\text{-}inner\text{-}product\text{-}space\text{-}as\text{-}module}(\operatorname{sub\text{-}cips}(v, d)), {q}_{+}))).\; (\iota\, u.\;((u \in \operatorname{vec}(\operatorname{hilbert\text{-}completion}(\operatorname{complex\text{-}inner\text{-}product\text{-}space\text{-}as\text{-}module}(\operatorname{sub\text{-}cips}(v, d)), {q}_{+}))) \wedge (g(u) = i))))(\operatorname{bgb\_})) - \langle\cdot,\cdot\rangle_{v}(\operatorname{bga\_}, \operatorname{bgb\_})))}+}), \\ \quad \quad \quad d) \end{array}}}
\vnbprop{fr-g-exists}{\vnbstmt{\begin{array}{@{}l@{}} \forall v, d, q. \\ \quad (\operatorname{is\text{-}hilbert\text{-}space}(v) \wedge \operatorname{is\text{-}positive\text{-}form}(v, d, q)) \Rightarrow  \\ \quad \quad \exists g. \\ \quad \quad \quad \operatorname{is\text{-}uniformly\text{-}continuous}(\operatorname{ip\text{-}metric\text{-}space}(\operatorname{hilbert\text{-}completion}(\operatorname{complex\text{-}inner\text{-}product\text{-}space\text{-}as\text{-}module}(\operatorname{sub\text{-}cips}(v, d)), {q}_{+})), \\ \quad \quad \quad \quad \operatorname{ip\text{-}metric\text{-}space}(v), g) \wedge \\ \quad \quad \quad \forall x \in d. \\ \quad \quad \quad \quad g(\operatorname{hc\text{-}embed}(\operatorname{complex\text{-}inner\text{-}product\text{-}space\text{-}as\text{-}module}(\operatorname{sub\text{-}cips}(v, d)), {q}_{+})(x)) = x \end{array}}}
\vnbprop{friedrichs-extension-of-form}{\vnbstmt{\begin{array}{@{}l@{}} \forall v, d, q, f. \\ \quad \operatorname{is\text{-}hilbert\text{-}space}(v) \wedge \\ \quad \operatorname{is\text{-}positive\text{-}form}(v, d, q) \wedge \\ \quad f \in (d \to \operatorname{vec}(v)) \wedge \\ \quad \forall y, x \in d.\; (q(x, y) = \langle\cdot,\cdot\rangle_{v}(x, f(y))) \Rightarrow \\ \quad \quad \exists d_{1}, q_{1}. \\ \quad \quad \quad \operatorname{is\text{-}closed\text{-}form}(v, d_{1}, q_{1}) \wedge \\ \quad \quad \quad \operatorname{is\text{-}dense}(\operatorname{ip\text{-}metric\text{-}space}(H_{{q_{1}}+}), d) \wedge \\ \quad \quad \quad d \subseteq d_{1} \wedge \\ \quad \quad \quad \forall \mathit{brx}, \mathit{bry} \in d.\; (\operatorname{q1}(\mathit{brx}, \mathit{bry}) = q(\mathit{brx}, \mathit{bry})) \wedge \\ \quad \quad \quad \forall \mathit{bry} \in d.\; ({[}\mathit{bry}, f(\mathit{bry})] \in T_{q_{1}}) \wedge \\ \quad \quad \quad \operatorname{is\text{-}self\text{-}adjoint\text{-}graph}(v, T_{q_{1}}) \wedge \\ \quad \quad \quad \forall c, \mathit{brd} \in \operatorname{vec}(v). \\ \quad \quad \quad \quad ({[}c, \mathit{brd}] \in T_{q_{1}}) \Rightarrow  \\ \quad \quad \quad \quad \quad (\langle\cdot,\cdot\rangle_{v}(c, \mathit{brd}) \in \mathbb{R}) \wedge (0 \leq \langle\cdot,\cdot\rangle_{v}(c, \mathit{brd})) \end{array}}}
\vnbprop{friedrichs-extension}{\vnbstmt{\begin{array}{@{}l@{}} \forall v, d, f. \\ \quad \operatorname{is\text{-}hilbert\text{-}space}(v) \wedge \\ \quad \operatorname{is\text{-}submodule}(v, d) \wedge \\ \quad \operatorname{is\text{-}dense}(\operatorname{ip\text{-}metric\text{-}space}(v), d) \wedge \\ \quad f \in (d \to \operatorname{vec}(v)) \wedge \\ \quad \forall x, y \in d.\; (f(\operatorname{vadd}(v)(x, y)) = \operatorname{vadd}(v)(f(x), f(y))) \wedge \\ \quad \forall c \in \mathbb{C}, x \in d.\; (f(\operatorname{act}(v)(c, x)) = \operatorname{act}(v)(c, f(x))) \wedge \\ \quad \forall \mathit{bsx}, \mathit{bsy} \in d.\; (\langle\cdot,\cdot\rangle_{v}(f(\mathit{bsx}), \mathit{bsy}) = \langle\cdot,\cdot\rangle_{v}(\mathit{bsx}, f(\mathit{bsy}))) \wedge \\ \quad \forall \mathit{bsx} \in d.\; ((\langle\cdot,\cdot\rangle_{v}(\mathit{bsx}, f(\mathit{bsx})) \in \mathbb{R}) \wedge (0 \leq \langle\cdot,\cdot\rangle_{v}(\mathit{bsx}, f(\mathit{bsx})))) \Rightarrow \\ \quad \quad \exists d_{1}, q_{1}. \\ \quad \quad \quad \operatorname{is\text{-}closed\text{-}form}(v, d_{1}, q_{1}) \wedge \\ \quad \quad \quad \operatorname{is\text{-}dense}(\operatorname{ip\text{-}metric\text{-}space}(H_{{q_{1}}+}), d) \wedge \\ \quad \quad \quad d \subseteq d_{1} \wedge \\ \quad \quad \quad \forall \mathit{brx}, \mathit{bry} \in d.\; (\operatorname{q1}(\mathit{brx}, \mathit{bry}) = \langle\cdot,\cdot\rangle_{v}(\mathit{brx}, f(\mathit{bry}))) \wedge \\ \quad \quad \quad \forall \mathit{bry} \in d.\; ({[}\mathit{bry}, f(\mathit{bry})] \in T_{q_{1}}) \wedge \\ \quad \quad \quad \operatorname{is\text{-}self\text{-}adjoint\text{-}graph}(v, T_{q_{1}}) \wedge \\ \quad \quad \quad \forall \mathit{brc}, \mathit{brd} \in \operatorname{vec}(v). \\ \quad \quad \quad \quad ({[}\mathit{brc}, \mathit{brd}] \in T_{q_{1}}) \Rightarrow  \\ \quad \quad \quad \quad \quad (\langle\cdot,\cdot\rangle_{v}(\mathit{brc}, \mathit{brd}) \in \mathbb{R}) \wedge (0 \leq \langle\cdot,\cdot\rangle_{v}(\mathit{brc}, \mathit{brd})) \end{array}}}

\vnbbookfile{structure-library/ba-amp-vocab}{structure-library/ba-amp-vocab}
\noindent{\small ba-amp-vocab.scm -- Thayer, Banach Algebras and Applications, ch. 6 (the amplification in the proof of the bicommutant theorem 6.1): the canonical injections, projections and diagonal operators of the Hilbert direct sum HSUM(w, v) (ba-ops-vocab.scm), the one-step amplification phi |-$>$ (T |-$>$ phi(T) (+) T), the identity ...}\par
\par\medskip\noindent\textbf{Definition} (\texttt{hsum-in1}).\;$\iota_1 :=$ $(\lambda \operatorname{amx\_} \in \operatorname{vec}(\operatorname{amw\_}).\; {[}\operatorname{amx\_}, \operatorname{vzero}(\operatorname{amv\_})])$
\par\medskip\noindent\textbf{Definition} (\texttt{hsum-in2}).\;$\iota_2 :=$ $(\lambda \operatorname{amy\_} \in \operatorname{vec}(\operatorname{amv\_}).\; {[}\operatorname{vzero}(\operatorname{amw\_}), \operatorname{amy\_}])$
\par\medskip\noindent\textbf{Definition} (\texttt{hsum-pr1}).\;$\pi_1 :=$
\begin{equation*}
(\lambda \operatorname{amp\_} \in (\operatorname{vec}(\operatorname{amw\_}) \times \operatorname{vec}(\operatorname{amv\_})).\; \operatorname{nth}(1, \operatorname{amp\_}))
\end{equation*}

\par\medskip\noindent\textbf{Definition} (\texttt{hsum-pr2}).\;$\pi_2 :=$
\begin{equation*}
(\lambda \operatorname{amp\_} \in (\operatorname{vec}(\operatorname{amw\_}) \times \operatorname{vec}(\operatorname{amv\_})).\; \operatorname{nth}(2, \operatorname{amp\_}))
\end{equation*}

\par\medskip\noindent\textbf{Definition} (\texttt{hsum-diag}).\;${\operatorname{ama\_}} \oplus {\operatorname{amb\_}} :=$
\begin{equation*}
(\lambda \operatorname{amp\_} \in (\operatorname{vec}(\operatorname{amw\_}) \times \operatorname{vec}(\operatorname{amv\_})).\; {[}\operatorname{ama\_}(\operatorname{nth}(1, \operatorname{amp\_})), \operatorname{amb\_}(\operatorname{nth}(2, \operatorname{amp\_}))])
\end{equation*}

\par\medskip\noindent\textbf{Definition} (\texttt{amp-step}).\;${\operatorname{amf\_}} \oplus \mathrm{id} :=$
\begin{equation*}
(\lambda \operatorname{amt\_} \in L({\operatorname{amv\_}}, {\operatorname{amv\_}}).\; {\operatorname{amf\_}(\operatorname{amt\_})} \oplus {\operatorname{amt\_}})
\end{equation*}

\par\medskip\noindent\textbf{Definition} (\texttt{id-rep}).\;$\mathrm{id} :=$ $(\lambda \operatorname{amt\_} \in L({\operatorname{amv\_}}, {\operatorname{amv\_}}).\; \operatorname{amt\_})$
\par\medskip\noindent\textbf{Definition} (\texttt{rep-image}).\;${\operatorname{amf\_}}({\operatorname{ams\_}}) :=$
\begin{equation*}
\{\operatorname{amb\_} \in L({\operatorname{amw\_}}, {\operatorname{amw\_}}) : \exists \operatorname{amt\_} \in \operatorname{ams\_}.\; (\operatorname{amb\_} = \operatorname{amf\_}(\operatorname{amt\_}))\}
\end{equation*}

\par\medskip\noindent\textbf{Definition} (\texttt{is-star-rep}).\;For all $\operatorname{amv\_}$, $\operatorname{amw\_}$ and $\operatorname{amf\_}$, 
\begin{equation*}
\begin{array}{@{}l@{}} \operatorname{is\text{-}star\text{-}rep}(\operatorname{amv\_}, \operatorname{amw\_}, \operatorname{amf\_}) \\ \Leftrightarrow\; \operatorname{amf\_} \in (L({\operatorname{amv\_}}, {\operatorname{amv\_}}) \to L({\operatorname{amw\_}}, {\operatorname{amw\_}})) \wedge \\ \forall \operatorname{ama\_}, \operatorname{amb\_} \in L({\operatorname{amv\_}}, {\operatorname{amv\_}}). \\ \quad \operatorname{amf\_}({\operatorname{ama\_}} + {\operatorname{amb\_}}) = {\operatorname{amf\_}(\operatorname{ama\_})} + {\operatorname{amf\_}(\operatorname{amb\_})} \wedge \\ \forall \operatorname{amc\_} \in \mathbb{C}, \operatorname{ama\_} \in L({\operatorname{amv\_}}, {\operatorname{amv\_}}). \\ \quad \operatorname{amf\_}({\operatorname{amc\_}}\,{\operatorname{ama\_}}) = {\operatorname{amc\_}}\,{\operatorname{amf\_}(\operatorname{ama\_})} \wedge \\ \forall \operatorname{ama\_}, \operatorname{amb\_} \in L({\operatorname{amv\_}}, {\operatorname{amv\_}}). \\ \quad \operatorname{amf\_}({\operatorname{ama\_}}{\operatorname{amb\_}}) = {\operatorname{amf\_}(\operatorname{ama\_})}{\operatorname{amf\_}(\operatorname{amb\_})} \wedge \\ \operatorname{amf\_}(1) = 1 \wedge \\ \forall \operatorname{ama\_} \in L({\operatorname{amv\_}}, {\operatorname{amv\_}}). \\ \quad \operatorname{amf\_}({\operatorname{ama\_}}^{*}) = {\operatorname{amf\_}(\operatorname{ama\_})}^{*} \end{array}
\end{equation*}

\par\medskip\noindent\textbf{Definition} (\texttt{is-amplification}).\;For all $\operatorname{amv\_}$, $\operatorname{ams\_}$, $\operatorname{amw\_}$ and $\operatorname{amf\_}$, 
\begin{equation*}
\begin{array}{@{}l@{}} \operatorname{is\text{-}amplification}(\operatorname{amv\_}, \operatorname{ams\_}, \operatorname{amw\_}, \operatorname{amf\_}) \\ \Leftrightarrow\; \operatorname{is\text{-}hilbert\text{-}space}(\operatorname{amw\_}) \wedge \\ \operatorname{is\text{-}star\text{-}rep}(\operatorname{amv\_}, \operatorname{amw\_}, \operatorname{amf\_}) \wedge \\ \operatorname{is\text{-}intertwined}(\operatorname{amv\_}, \operatorname{amw\_}, \mathrm{id}, \operatorname{amf\_}, \operatorname{amv\_}, \operatorname{ams\_}) \wedge \\ \operatorname{is\text{-}intertwined}(\operatorname{amw\_}, \operatorname{amv\_}, \operatorname{amf\_}, \mathrm{id}, \operatorname{amv\_}, \operatorname{ams\_}) \wedge \\ \operatorname{is\text{-}intertwined}(\operatorname{amw\_}, \operatorname{amw\_}, \operatorname{amf\_}, \operatorname{amf\_}, \operatorname{amv\_}, \operatorname{ams\_}) \end{array}
\end{equation*}

\par\medskip\noindent\textbf{Definition} (\texttt{is-intertwined}).\;For all $\operatorname{amu\_}$, $\operatorname{amz\_}$, $\operatorname{amr\_}$, $\operatorname{amq\_}$, $\operatorname{amv\_}$ and $\operatorname{ams\_}$, 
\begin{equation*}
\begin{array}{@{}l@{}} \operatorname{is\text{-}intertwined}(\operatorname{amu\_}, \operatorname{amz\_}, \operatorname{amr\_}, \operatorname{amq\_}, \operatorname{amv\_}, \operatorname{ams\_}) \\ \Leftrightarrow\; \forall \operatorname{amk\_} \in L({\operatorname{amu\_}}, {\operatorname{amz\_}}). \\ \quad \forall \operatorname{ama\_} \in \operatorname{ams\_}, \operatorname{amx\_} \in \operatorname{vec}(\operatorname{amu\_}).\; (\operatorname{amk\_}(\operatorname{amr\_}(\operatorname{ama\_})(\operatorname{amx\_})) = \operatorname{amq\_}(\operatorname{ama\_})(\operatorname{amk\_}(\operatorname{amx\_}))) \Rightarrow  \\ \quad \quad \forall \operatorname{amt\_} \in {\operatorname{ams\_}}'', \operatorname{amx\_} \in \operatorname{vec}(\operatorname{amu\_}). \\ \quad \quad \quad \operatorname{amk\_}(\operatorname{amr\_}(\operatorname{amt\_})(\operatorname{amx\_})) = \operatorname{amq\_}(\operatorname{amt\_})(\operatorname{amk\_}(\operatorname{amx\_})) \end{array}
\end{equation*}


\vnbbookfile{theorem-library/ba-amp}{theorem-library/ba-amp}
\noindent{\small ba-amp.scm -- Thayer, Banach Algebras and Applications, ch. 6: von Neumann's bicommutant theorem (Theorem 6.1) by amplification, and Corollary 6.9 for a von Neumann algebra. Part A  the injections, projections and diagonal operators of the Hilbert direct sum HSUM(w, v), their values and bounds; composites of bounded ...}\par
\vnbprop{hsum-vec-eq}{\vnbstmt{\begin{array}{@{}l@{}} \forall w, v. \\ \quad (\operatorname{is\text{-}complex\text{-}inner\text{-}product\text{-}space}(w) \wedge \operatorname{is\text{-}complex\text{-}inner\text{-}product\text{-}space}(v)) \Rightarrow  \\ \quad \quad \operatorname{vec}({w} \oplus {v}) = (\operatorname{vec}(w) \times \operatorname{vec}(v)) \end{array}}}
\vnbprop{cips-act-vzero}{\vnbstmt{\begin{array}{@{}l@{}} \forall v.; \forall c \in \mathbb{C}. \\ \quad \operatorname{is\text{-}complex\text{-}inner\text{-}product\text{-}space}(v) \Rightarrow (\operatorname{act}(v)(c, \operatorname{vzero}(v)) = \operatorname{vzero}(v)) \end{array}}}
\vnbprop{hsum-in1-apply}{\vnbstmt{\begin{array}{@{}l@{}} \forall w, v, x. \\ \quad (\operatorname{is\text{-}complex\text{-}inner\text{-}product\text{-}space}(w) \wedge (\operatorname{is\text{-}complex\text{-}inner\text{-}product\text{-}space}(v) \wedge (x \in \operatorname{vec}(w)))) \Rightarrow  \\ \quad \quad \iota_1(x) = {[}x, \operatorname{vzero}(v)] \end{array}}}
\vnbprop{hsum-in2-apply}{\vnbstmt{\begin{array}{@{}l@{}} \forall w, v, y. \\ \quad (\operatorname{is\text{-}complex\text{-}inner\text{-}product\text{-}space}(w) \wedge (\operatorname{is\text{-}complex\text{-}inner\text{-}product\text{-}space}(v) \wedge (y \in \operatorname{vec}(v)))) \Rightarrow  \\ \quad \quad \iota_2(y) = {[}\operatorname{vzero}(w), y] \end{array}}}
\vnbprop{hsum-pr1-apply}{\vnbstmt{\begin{array}{@{}l@{}} \forall w, v, x, y. \\ \quad \operatorname{is\text{-}complex\text{-}inner\text{-}product\text{-}space}(w) \wedge \\ \quad \operatorname{is\text{-}complex\text{-}inner\text{-}product\text{-}space}(v) \wedge \\ \quad x \in \operatorname{vec}(w) \wedge \\ \quad y \in \operatorname{vec}(v) \Rightarrow \\ \quad \quad \pi_1({[}x, y]) = x \end{array}}}
\vnbprop{hsum-pr2-apply}{\vnbstmt{\begin{array}{@{}l@{}} \forall w, v, x, y. \\ \quad \operatorname{is\text{-}complex\text{-}inner\text{-}product\text{-}space}(w) \wedge \\ \quad \operatorname{is\text{-}complex\text{-}inner\text{-}product\text{-}space}(v) \wedge \\ \quad x \in \operatorname{vec}(w) \wedge \\ \quad y \in \operatorname{vec}(v) \Rightarrow \\ \quad \quad \pi_2({[}x, y]) = y \end{array}}}
\vnbprop{hsum-diag-apply}{\vnbstmt{\begin{array}{@{}l@{}} \forall w, v, a, b, x, y. \\ \quad \operatorname{is\text{-}complex\text{-}inner\text{-}product\text{-}space}(w) \wedge \\ \quad \operatorname{is\text{-}complex\text{-}inner\text{-}product\text{-}space}(v) \wedge \\ \quad a \in (\operatorname{vec}(w) \to \operatorname{vec}(w)) \wedge \\ \quad b \in (\operatorname{vec}(v) \to \operatorname{vec}(v)) \wedge \\ \quad x \in \operatorname{vec}(w) \wedge \\ \quad y \in \operatorname{vec}(v) \Rightarrow \\ \quad \quad {a} \oplus {b}({[}x, y]) = {[}a(x), b(y)] \end{array}}}
\vnbprop{hsum-pair-split}{\vnbstmt{\begin{array}{@{}l@{}} \forall w, v, x, y. \\ \quad \operatorname{is\text{-}complex\text{-}inner\text{-}product\text{-}space}(w) \wedge \\ \quad \operatorname{is\text{-}complex\text{-}inner\text{-}product\text{-}space}(v) \wedge \\ \quad x \in \operatorname{vec}(w) \wedge \\ \quad y \in \operatorname{vec}(v) \Rightarrow \\ \quad \quad {[}x, y] = \operatorname{vadd}({w} \oplus {v})({[}x, \operatorname{vzero}(v)], {[}\operatorname{vzero}(w), y]) \end{array}}}
\vnbprop{hsum-norm-pair-sq}{\vnbstmt{\begin{array}{@{}l@{}} \forall w, v, a, b. \\ \quad \operatorname{is\text{-}complex\text{-}inner\text{-}product\text{-}space}(w) \wedge \\ \quad \operatorname{is\text{-}complex\text{-}inner\text{-}product\text{-}space}(v) \wedge \\ \quad a \in \operatorname{vec}(w) \wedge \\ \quad b \in \operatorname{vec}(v) \Rightarrow \\ \quad \quad \lVert {[}a, b] \rVert \cdot \lVert {[}a, b] \rVert \\ \quad \quad =\; (\lVert a \rVert \cdot \lVert a \rVert) + (\lVert b \rVert \cdot \lVert b \rVert) \end{array}}}
\vnbprop{hsum-norm-ge-1}{\vnbstmt{\begin{array}{@{}l@{}} \forall w, v, a, b. \\ \quad \operatorname{is\text{-}complex\text{-}inner\text{-}product\text{-}space}(w) \wedge \\ \quad \operatorname{is\text{-}complex\text{-}inner\text{-}product\text{-}space}(v) \wedge \\ \quad a \in \operatorname{vec}(w) \wedge \\ \quad b \in \operatorname{vec}(v) \Rightarrow \\ \quad \quad \lVert a \rVert \leq \lVert {[}a, b] \rVert \end{array}}}
\vnbprop{hsum-norm-ge-2}{\vnbstmt{\begin{array}{@{}l@{}} \forall w, v, a, b. \\ \quad \operatorname{is\text{-}complex\text{-}inner\text{-}product\text{-}space}(w) \wedge \\ \quad \operatorname{is\text{-}complex\text{-}inner\text{-}product\text{-}space}(v) \wedge \\ \quad a \in \operatorname{vec}(w) \wedge \\ \quad b \in \operatorname{vec}(v) \Rightarrow \\ \quad \quad \lVert b \rVert \leq \lVert {[}a, b] \rVert \end{array}}}
\vnbprop{hsum-norm-le-sum}{\vnbstmt{\begin{array}{@{}l@{}} \forall w, v, a, b. \\ \quad \operatorname{is\text{-}complex\text{-}inner\text{-}product\text{-}space}(w) \wedge \\ \quad \operatorname{is\text{-}complex\text{-}inner\text{-}product\text{-}space}(v) \wedge \\ \quad a \in \operatorname{vec}(w) \wedge \\ \quad b \in \operatorname{vec}(v) \Rightarrow \\ \quad \quad \lVert {[}a, b] \rVert \leq (\lVert a \rVert + \lVert b \rVert) \end{array}}}
\vnbprop{bop-comp-apply-gen}{\vnbstmt{\begin{array}{@{}l@{}} \forall a, b, c, s, t, x. \\ \quad ((t \in (\operatorname{vec}(a) \to \operatorname{vec}(b))) \wedge ((s \in (\operatorname{vec}(b) \to \operatorname{vec}(c))) \wedge (x \in \operatorname{vec}(a)))) \Rightarrow  \\ \quad \quad {s}{t}(x) = s(t(x)) \end{array}}}
\vnbprop{hsum-in1-bounded}{\vnbstmt{\begin{array}{@{}l@{}} \forall w, v. \\ \quad (\operatorname{is\text{-}complex\text{-}inner\text{-}product\text{-}space}(w) \wedge \operatorname{is\text{-}complex\text{-}inner\text{-}product\text{-}space}(v)) \Rightarrow  \\ \quad \quad \iota_1 \in L({w}, {{w} \oplus {v}}) \end{array}}}
\vnbprop{hsum-in2-bounded}{\vnbstmt{\begin{array}{@{}l@{}} \forall w, v. \\ \quad (\operatorname{is\text{-}complex\text{-}inner\text{-}product\text{-}space}(w) \wedge \operatorname{is\text{-}complex\text{-}inner\text{-}product\text{-}space}(v)) \Rightarrow  \\ \quad \quad \iota_2 \in L({v}, {{w} \oplus {v}}) \end{array}}}
\vnbprop{hsum-pr1-bounded}{\vnbstmt{\begin{array}{@{}l@{}} \forall w, v. \\ \quad (\operatorname{is\text{-}complex\text{-}inner\text{-}product\text{-}space}(w) \wedge \operatorname{is\text{-}complex\text{-}inner\text{-}product\text{-}space}(v)) \Rightarrow  \\ \quad \quad \pi_1 \in L({{w} \oplus {v}}, {w}) \end{array}}}
\vnbprop{hsum-pr2-bounded}{\vnbstmt{\begin{array}{@{}l@{}} \forall w, v. \\ \quad (\operatorname{is\text{-}complex\text{-}inner\text{-}product\text{-}space}(w) \wedge \operatorname{is\text{-}complex\text{-}inner\text{-}product\text{-}space}(v)) \Rightarrow  \\ \quad \quad \pi_2 \in L({{w} \oplus {v}}, {v}) \end{array}}}
\vnbprop{bop-comp-bounded-gen}{\vnbstmt{\begin{array}{@{}l@{}} \forall a, b, c, s, t. \\ \quad \operatorname{is\text{-}complex\text{-}inner\text{-}product\text{-}space}(a) \wedge \\ \quad \operatorname{is\text{-}complex\text{-}inner\text{-}product\text{-}space}(b) \wedge \\ \quad \operatorname{is\text{-}complex\text{-}inner\text{-}product\text{-}space}(c) \wedge \\ \quad t \in L({a}, {b}) \wedge \\ \quad s \in L({b}, {c}) \Rightarrow \\ \quad \quad {s}{t} \in L({a}, {c}) \end{array}}}
\vnbprop{hsum-diag-bounded}{\vnbstmt{\begin{array}{@{}l@{}} \forall w, v, a, b. \\ \quad \operatorname{is\text{-}complex\text{-}inner\text{-}product\text{-}space}(w) \wedge \\ \quad \operatorname{is\text{-}complex\text{-}inner\text{-}product\text{-}space}(v) \wedge \\ \quad a \in L({w}, {w}) \wedge \\ \quad b \in L({v}, {v}) \Rightarrow \\ \quad \quad {a} \oplus {b} \in L({{w} \oplus {v}}, {{w} \oplus {v}}) \end{array}}}
\vnbprop{hsum-pr1-value}{\vnbstmt{\begin{array}{@{}l@{}} \forall w, v, p. \\ \quad \operatorname{is\text{-}complex\text{-}inner\text{-}product\text{-}space}(w) \wedge \\ \quad \operatorname{is\text{-}complex\text{-}inner\text{-}product\text{-}space}(v) \wedge \\ \quad p \in (\operatorname{vec}(w) \times \operatorname{vec}(v)) \Rightarrow \\ \quad \quad \pi_1(p) = \operatorname{nth}(1, p) \end{array}}}
\vnbprop{hsum-pr2-value}{\vnbstmt{\begin{array}{@{}l@{}} \forall w, v, p. \\ \quad \operatorname{is\text{-}complex\text{-}inner\text{-}product\text{-}space}(w) \wedge \\ \quad \operatorname{is\text{-}complex\text{-}inner\text{-}product\text{-}space}(v) \wedge \\ \quad p \in (\operatorname{vec}(w) \times \operatorname{vec}(v)) \Rightarrow \\ \quad \quad \pi_2(p) = \operatorname{nth}(2, p) \end{array}}}
\vnbprop{hsum-diag-value}{\vnbstmt{\begin{array}{@{}l@{}} \forall w, v, a, b, p. \\ \quad \operatorname{is\text{-}complex\text{-}inner\text{-}product\text{-}space}(w) \wedge \\ \quad \operatorname{is\text{-}complex\text{-}inner\text{-}product\text{-}space}(v) \wedge \\ \quad a \in (\operatorname{vec}(w) \to \operatorname{vec}(w)) \wedge \\ \quad b \in (\operatorname{vec}(v) \to \operatorname{vec}(v)) \wedge \\ \quad p \in (\operatorname{vec}(w) \times \operatorname{vec}(v)) \Rightarrow \\ \quad \quad {a} \oplus {b}(p) = {[}a(\operatorname{nth}(1, p)), b(\operatorname{nth}(2, p))] \end{array}}}
\vnbprop{id-rep-apply}{\vnbstmt{\forall v.; \forall t \in L({v}, {v}).\; (\mathrm{id}(t) = t)}}
\vnbprop{amp-step-apply}{\vnbstmt{\begin{array}{@{}l@{}} \forall w, v, f.; \forall t \in L({v}, {v}). \\ \quad ({f} \oplus \mathrm{id}(t) \simeq {f(t)} \oplus {t}) \end{array}}}
\vnbprop{id-rep-type}{\vnbstmt{\begin{array}{@{}l@{}} \forall v. \\ \quad \operatorname{is\text{-}complex\text{-}inner\text{-}product\text{-}space}(v) \Rightarrow  \\ \quad \quad \mathrm{id} \in (L({v}, {v}) \to L({v}, {v})) \end{array}}}
\vnbprop{amp-step-type}{\vnbstmt{\begin{array}{@{}l@{}} \forall w, v, f. \\ \quad \operatorname{is\text{-}complex\text{-}inner\text{-}product\text{-}space}(w) \wedge \\ \quad \operatorname{is\text{-}complex\text{-}inner\text{-}product\text{-}space}(v) \wedge \\ \quad f \in (L({v}, {v}) \to L({w}, {w})) \Rightarrow \\ \quad \quad {f} \oplus \mathrm{id} \in (L({v}, {v}) \to L({{w} \oplus {v}}, {{w} \oplus {v}})) \end{array}}}
\vnbprop{intertwined-base}{\vnbstmt{\begin{array}{@{}l@{}} \forall v, s. \\ \quad (\operatorname{is\text{-}complex\text{-}inner\text{-}product\text{-}space}(v) \wedge (s \subseteq L({v}, {v}))) \Rightarrow  \\ \quad \quad \operatorname{is\text{-}intertwined}(v, v, \mathrm{id}, \mathrm{id}, v, s) \end{array}}}
\vnbprop{intertwined-split-codomain}{\vnbstmt{\begin{array}{@{}l@{}} \forall u, w, v, f, r, s. \\ \quad \operatorname{is\text{-}complex\text{-}inner\text{-}product\text{-}space}(u) \wedge \\ \quad \operatorname{is\text{-}complex\text{-}inner\text{-}product\text{-}space}(w) \wedge \\ \quad \operatorname{is\text{-}complex\text{-}inner\text{-}product\text{-}space}(v) \wedge \\ \quad f \in (L({v}, {v}) \to L({w}, {w})) \wedge \\ \quad r \in (L({v}, {v}) \to L({u}, {u})) \wedge \\ \quad s \subseteq L({v}, {v}) \wedge \\ \quad \operatorname{is\text{-}intertwined}(u, w, r, f, v, s) \wedge \\ \quad \operatorname{is\text{-}intertwined}(u, v, r, \mathrm{id}, v, s) \Rightarrow \\ \quad \quad \operatorname{is\text{-}intertwined}(u, {w} \oplus {v}, r, {f} \oplus \mathrm{id}, v, s) \end{array}}}
\vnbprop{intertwined-split-domain}{\vnbstmt{\begin{array}{@{}l@{}} \forall w, v, u, f, r, s. \\ \quad \operatorname{is\text{-}complex\text{-}inner\text{-}product\text{-}space}(w) \wedge \\ \quad \operatorname{is\text{-}complex\text{-}inner\text{-}product\text{-}space}(v) \wedge \\ \quad \operatorname{is\text{-}complex\text{-}inner\text{-}product\text{-}space}(u) \wedge \\ \quad f \in (L({v}, {v}) \to L({w}, {w})) \wedge \\ \quad r \in (L({v}, {v}) \to L({u}, {u})) \wedge \\ \quad s \subseteq L({v}, {v}) \wedge \\ \quad \operatorname{is\text{-}intertwined}(w, u, f, r, v, s) \wedge \\ \quad \operatorname{is\text{-}intertwined}(v, u, \mathrm{id}, r, v, s) \Rightarrow \\ \quad \quad \operatorname{is\text{-}intertwined}({w} \oplus {v}, u, {f} \oplus \mathrm{id}, r, v, s) \end{array}}}
\vnbprop{hsum-diag-sum}{\vnbstmt{\begin{array}{@{}l@{}} \forall w, v, a, b, c, d. \\ \quad \operatorname{is\text{-}complex\text{-}inner\text{-}product\text{-}space}(w) \wedge \\ \quad \operatorname{is\text{-}complex\text{-}inner\text{-}product\text{-}space}(v) \wedge \\ \quad a \in L({w}, {w}) \wedge \\ \quad b \in L({v}, {v}) \wedge \\ \quad c \in L({w}, {w}) \wedge \\ \quad d \in L({v}, {v}) \Rightarrow \\ \quad \quad {{a} \oplus {b}} + {{c} \oplus {d}} \\ \quad \quad =\; {{a} + {c}} \oplus {{b} + {d}} \end{array}}}
\vnbprop{hsum-diag-scale}{\vnbstmt{\begin{array}{@{}l@{}} \forall w, v, a, b, k. \\ \quad \operatorname{is\text{-}complex\text{-}inner\text{-}product\text{-}space}(w) \wedge \\ \quad \operatorname{is\text{-}complex\text{-}inner\text{-}product\text{-}space}(v) \wedge \\ \quad a \in L({w}, {w}) \wedge \\ \quad b \in L({v}, {v}) \wedge \\ \quad k \in \mathbb{C} \Rightarrow \\ \quad \quad {k}\,{{a} \oplus {b}} \\ \quad \quad =\; {{k}\,{a}} \oplus {{k}\,{b}} \end{array}}}
\vnbprop{hsum-diag-comp}{\vnbstmt{\begin{array}{@{}l@{}} \forall w, v, a, b, c, d. \\ \quad \operatorname{is\text{-}complex\text{-}inner\text{-}product\text{-}space}(w) \wedge \\ \quad \operatorname{is\text{-}complex\text{-}inner\text{-}product\text{-}space}(v) \wedge \\ \quad a \in L({w}, {w}) \wedge \\ \quad b \in L({v}, {v}) \wedge \\ \quad c \in L({w}, {w}) \wedge \\ \quad d \in L({v}, {v}) \Rightarrow \\ \quad \quad {{a} \oplus {b}}{{c} \oplus {d}} \\ \quad \quad =\; {{a}{c}} \oplus {{b}{d}} \end{array}}}
\vnbprop{hsum-diag-id}{\vnbstmt{\begin{array}{@{}l@{}} \forall w, v. \\ \quad (\operatorname{is\text{-}complex\text{-}inner\text{-}product\text{-}space}(w) \wedge \operatorname{is\text{-}complex\text{-}inner\text{-}product\text{-}space}(v)) \Rightarrow  \\ \quad \quad {1} \oplus {1} = 1 \end{array}}}
\vnbprop{hsum-diag-adjoint}{\vnbstmt{\begin{array}{@{}l@{}} \forall w, v, a, b. \\ \quad (\operatorname{is\text{-}hilbert\text{-}space}(w) \wedge (\operatorname{is\text{-}hilbert\text{-}space}(v) \wedge ((a \in L({w}, {w})) \wedge (b \in L({v}, {v}))))) \Rightarrow  \\ \quad \quad {{a} \oplus {b}}^{*} \\ \quad \quad =\; {{a}^{*}} \oplus {{b}^{*}} \end{array}}}
\vnbprop{id-rep-star-rep}{\vnbstmt{\forall v.\; (\operatorname{is\text{-}hilbert\text{-}space}(v) \Rightarrow \operatorname{is\text{-}star\text{-}rep}(v, v, \mathrm{id}))}}
\vnbprop{amp-step-star-rep}{\vnbstmt{\begin{array}{@{}l@{}} \forall w, v, f. \\ \quad (\operatorname{is\text{-}hilbert\text{-}space}(w) \wedge (\operatorname{is\text{-}hilbert\text{-}space}(v) \wedge \operatorname{is\text{-}star\text{-}rep}(v, w, f))) \Rightarrow  \\ \quad \quad \operatorname{is\text{-}star\text{-}rep}(v, {w} \oplus {v}, {f} \oplus \mathrm{id}) \end{array}}}
\vnbprop{amplification-base}{\vnbstmt{\begin{array}{@{}l@{}} \forall v, s. \\ \quad (\operatorname{is\text{-}hilbert\text{-}space}(v) \wedge (s \subseteq L({v}, {v}))) \Rightarrow  \\ \quad \quad \operatorname{is\text{-}amplification}(v, s, v, \mathrm{id}) \end{array}}}
\vnbprop{amplification-step}{\vnbstmt{\begin{array}{@{}l@{}} \forall v, s, w, f. \\ \quad (\operatorname{is\text{-}hilbert\text{-}space}(v) \wedge ((s \subseteq L({v}, {v})) \wedge \operatorname{is\text{-}amplification}(v, s, w, f))) \Rightarrow  \\ \quad \quad \operatorname{is\text{-}amplification}(v, s, {w} \oplus {v}, {f} \oplus \mathrm{id}) \end{array}}}
\vnbprop{rep-image-unfold}{\vnbstmt{\begin{array}{@{}l@{}} \forall v, w, f, s. \\ \quad ({f}({s}) \simeq \{b \in L({w}, {w}) : \exists t \in s.\; (b = f(t))\}) \end{array}}}
\vnbprop{rep-image-star-algebra}{\vnbstmt{\begin{array}{@{}l@{}} \forall v, w, f, s. \\ \quad (\operatorname{is\text{-}hilbert\text{-}space}(v) \wedge (\operatorname{is\text{-}hilbert\text{-}space}(w) \wedge (\operatorname{is\text{-}star\text{-}rep}(v, w, f) \wedge \operatorname{is\text{-}star\text{-}op\text{-}algebra}(v, s)))) \Rightarrow  \\ \quad \quad \operatorname{is\text{-}star\text{-}op\text{-}algebra}(w, {f}({s})) \end{array}}}
\vnbprop{rep-image-bicommutant}{\vnbstmt{\begin{array}{@{}l@{}} \forall v, s, w, f, t. \\ \quad \operatorname{is\text{-}hilbert\text{-}space}(v) \wedge \\ \quad \operatorname{is\text{-}star\text{-}op\text{-}algebra}(v, s) \wedge \\ \quad \operatorname{is\text{-}amplification}(v, s, w, f) \wedge \\ \quad t \in {s}'' \Rightarrow \\ \quad \quad f(t) \in {{f}({s})}'' \end{array}}}
\vnbprop{amp-step-norm-ge-1}{\vnbstmt{\begin{array}{@{}l@{}} \forall w, v, f, a, b, c, d. \\ \quad \operatorname{is\text{-}complex\text{-}inner\text{-}product\text{-}space}(w) \wedge \\ \quad \operatorname{is\text{-}complex\text{-}inner\text{-}product\text{-}space}(v) \wedge \\ \quad f \in (L({v}, {v}) \to L({w}, {w})) \wedge \\ \quad a \in L({v}, {v}) \wedge \\ \quad b \in L({v}, {v}) \wedge \\ \quad c \in \operatorname{vec}(w) \wedge \\ \quad d \in \operatorname{vec}(v) \Rightarrow \\ \quad \quad \lVert \operatorname{vadd}(w)(f(a)(c), \operatorname{vneg}(w)(f(b)(c))) \rVert \\ \quad \quad \leq\; \lVert \operatorname{vadd}({w} \oplus {v})({f} \oplus \mathrm{id}(a)({[}c, d]), \operatorname{vneg}({w} \oplus {v})({f} \oplus \mathrm{id}(b)({[}c, d]))) \rVert \end{array}}}
\vnbprop{amp-step-norm-ge-2}{\vnbstmt{\begin{array}{@{}l@{}} \forall w, v, f, a, b, c, d. \\ \quad \operatorname{is\text{-}complex\text{-}inner\text{-}product\text{-}space}(w) \wedge \\ \quad \operatorname{is\text{-}complex\text{-}inner\text{-}product\text{-}space}(v) \wedge \\ \quad f \in (L({v}, {v}) \to L({w}, {w})) \wedge \\ \quad a \in L({v}, {v}) \wedge \\ \quad b \in L({v}, {v}) \wedge \\ \quad c \in \operatorname{vec}(w) \wedge \\ \quad d \in \operatorname{vec}(v) \Rightarrow \\ \quad \quad \lVert \operatorname{vadd}(v)(a(d), \operatorname{vneg}(v)(b(d))) \rVert \\ \quad \quad \leq\; \lVert \operatorname{vadd}({w} \oplus {v})({f} \oplus \mathrm{id}(a)({[}c, d]), \operatorname{vneg}({w} \oplus {v})({f} \oplus \mathrm{id}(b)({[}c, d]))) \rVert \end{array}}}
\vnbprop{amp-vector-step}{\vnbstmt{\begin{array}{@{}l@{}} \forall v, s, w, f, x, u, z. \\ \quad \operatorname{is\text{-}hilbert\text{-}space}(v) \wedge \\ \quad s \subseteq L({v}, {v}) \wedge \\ \quad \operatorname{is\text{-}amplification}(v, s, w, f) \wedge \\ \quad x \in \operatorname{vec}(w) \wedge \\ \quad z \in \operatorname{vec}(v) \wedge \\ \quad u \subseteq \operatorname{vec}(v) \wedge \\ \quad \forall a, b \in L({v}, {v}), y \in u. \\ \quad \quad \lVert \operatorname{vadd}(v)(a(y), \operatorname{vneg}(v)(b(y))) \rVert \\ \quad \quad \leq\; \lVert \operatorname{vadd}(w)(f(a)(x), \operatorname{vneg}(w)(f(b)(x))) \rVert \Rightarrow \\ \quad \quad \exists \mathit{cqw}, \mathit{cqf} \in \mathbf{Set}. \\ \quad \quad \quad \operatorname{is\text{-}amplification}(v, s, \mathit{cqw}, \mathit{cqf}) \wedge \\ \quad \quad \quad \exists \mathit{cqx} \in \operatorname{vec}(\mathit{cqw}). \forall a, b \in L({v}, {v}), y \in (u \cup \{z\}). \\ \quad \quad \quad \quad \lVert \operatorname{vadd}(v)(a(y), \operatorname{vneg}(v)(b(y))) \rVert \\ \quad \quad \quad \quad \leq\; \lVert \operatorname{vadd}(\mathit{cqw})(\operatorname{cqf}(a)(\mathit{cqx}), \operatorname{vneg}(\mathit{cqw})(\operatorname{cqf}(b)(\mathit{cqx}))) \rVert \end{array}}}
\vnbprop{amp-vector-exists}{\vnbstmt{\begin{array}{@{}l@{}} \forall v, s, g. \\ \quad \operatorname{is\text{-}hilbert\text{-}space}(v) \wedge \\ \quad s \subseteq L({v}, {v}) \wedge \\ \quad g \in \mathbf{Set} \wedge \\ \quad |g| \in \mathbb{N} \wedge \\ \quad g \subseteq \operatorname{vec}(v) \Rightarrow \\ \quad \quad \exists w, f \in \mathbf{Set}. \\ \quad \quad \quad \operatorname{is\text{-}amplification}(v, s, w, f) \wedge \\ \quad \quad \quad \exists x \in \operatorname{vec}(w). \forall a, b \in L({v}, {v}), y \in g. \\ \quad \quad \quad \quad \lVert \operatorname{vadd}(v)(a(y), \operatorname{vneg}(v)(b(y))) \rVert \\ \quad \quad \quad \quad \leq\; \lVert \operatorname{vadd}(w)(f(a)(x), \operatorname{vneg}(w)(f(b)(x))) \rVert \end{array}}}
\vnbprop{bicommutant-approx-at}{\vnbstmt{\begin{array}{@{}l@{}} \forall v, s, w, f, x, g, t, a. \\ \quad \operatorname{is\text{-}hilbert\text{-}space}(v) \wedge \\ \quad \operatorname{is\text{-}star\text{-}op\text{-}algebra}(v, s) \wedge \\ \quad \operatorname{is\text{-}amplification}(v, s, w, f) \wedge \\ \quad x \in \operatorname{vec}(w) \wedge \\ \quad g \subseteq \operatorname{vec}(v) \wedge \\ \quad \forall \mathit{cqa}, b \in L({v}, {v}), y \in g. \\ \quad \quad \lVert \operatorname{vadd}(v)(\operatorname{cqa}(y), \operatorname{vneg}(v)(b(y))) \rVert \\ \quad \quad \leq\; \lVert \operatorname{vadd}(w)(f(\mathit{cqa})(x), \operatorname{vneg}(w)(f(b)(x))) \rVert \wedge \\ \quad t \in {s}'' \wedge \\ \quad \operatorname{pos\text{-}rr}(a) \Rightarrow \\ \quad \quad \exists \mathit{cqa} \in s. \forall y \in g. \\ \quad \quad \quad \lVert \operatorname{vadd}(v)(t(y), \operatorname{vneg}(v)(\operatorname{cqa}(y))) \rVert < a \end{array}}}
\vnbprop{bicommutant-finite-approx}{\vnbstmt{\begin{array}{@{}l@{}} \forall v, s, g, t, a. \\ \quad \operatorname{is\text{-}hilbert\text{-}space}(v) \wedge \\ \quad \operatorname{is\text{-}star\text{-}op\text{-}algebra}(v, s) \wedge \\ \quad g \in \mathbf{Set} \wedge \\ \quad |g| \in \mathbb{N} \wedge \\ \quad g \subseteq \operatorname{vec}(v) \wedge \\ \quad t \in {s}'' \wedge \\ \quad \operatorname{pos\text{-}rr}(a) \Rightarrow \\ \quad \quad \exists \mathit{cqa} \in s. \forall y \in g. \\ \quad \quad \quad \lVert \operatorname{vadd}(v)(t(y), \operatorname{vneg}(v)(\operatorname{cqa}(y))) \rVert < a \end{array}}}
\vnbprop{ip-norm-diff-sym}{\vnbstmt{\begin{array}{@{}l@{}} \forall v, x, y. \\ \quad (\operatorname{is\text{-}complex\text{-}inner\text{-}product\text{-}space}(v) \wedge ((x \in \operatorname{vec}(v)) \wedge (y \in \operatorname{vec}(v)))) \Rightarrow  \\ \quad \quad \lVert \operatorname{vadd}(v)(x, \operatorname{vneg}(v)(y)) \rVert = \lVert \operatorname{vadd}(v)(y, \operatorname{vneg}(v)(x)) \rVert \end{array}}}
\vnbprop{top-closure-unfold}{\vnbstmt{\begin{array}{@{}l@{}} \forall q, a. \\ \quad (\overline{a} \simeq \{x \in \operatorname{pts}(q) : \forall u \in \operatorname{opens}(q).\; ((x \in u) \Rightarrow \exists y \in a.\; (y \in u))\}) \end{array}}}
\vnbprop{sot-closure-in-bicommutant}{\vnbstmt{\begin{array}{@{}l@{}} \forall v, s. \\ \quad (\operatorname{is\text{-}hilbert\text{-}space}(v) \wedge (s \subseteq L({v}, {v}))) \Rightarrow  \\ \quad \quad \overline{s} \subseteq {s}'' \end{array}}}
\vnbprop{bicommutant-in-sot-closure}{\vnbstmt{\begin{array}{@{}l@{}} \forall v, s. \\ \quad (\operatorname{is\text{-}hilbert\text{-}space}(v) \wedge \operatorname{is\text{-}star\text{-}op\text{-}algebra}(v, s)) \Rightarrow  \\ \quad \quad {s}'' \subseteq \overline{s} \end{array}}}
\vnbprop{bicommutant-is-sot-closure}{\vnbstmt{\begin{array}{@{}l@{}} \forall v, s. \\ \quad (\operatorname{is\text{-}hilbert\text{-}space}(v) \wedge \operatorname{is\text{-}star\text{-}op\text{-}algebra}(v, s)) \Rightarrow  \\ \quad \quad {s}'' = \overline{s} \end{array}}}
\vnbprop{top-closure-closed-mem}{\vnbstmt{\forall q, a, x.\; ((\operatorname{is\text{-}top\text{-}closed}(q, a) \wedge (x \in \overline{a})) \Rightarrow (x \in a))}}
\vnbprop{von-neumann-bicommutant-eq}{\vnbstmt{\begin{array}{@{}l@{}} \forall v, s. \\ \quad (\operatorname{is\text{-}hilbert\text{-}space}(v) \wedge \operatorname{is\text{-}von\text{-}neumann\text{-}algebra}(v, s)) \Rightarrow  \\ \quad \quad {s}'' = s \end{array}}}
\vnbprop{von-neumann-cyclic-separating}{\vnbstmt{\begin{array}{@{}l@{}} \forall v, s, i. \\ \quad (\operatorname{is\text{-}hilbert\text{-}space}(v) \wedge \operatorname{is\text{-}von\text{-}neumann\text{-}algebra}(v, s)) \Rightarrow  \\ \quad \quad \operatorname{is\text{-}cyclic\text{-}for}(v, s, i) \Leftrightarrow \operatorname{is\text{-}separating\text{-}for}(v, {s}', i) \wedge \\ \quad \quad \operatorname{is\text{-}separating\text{-}for}(v, s, i) \Leftrightarrow \operatorname{is\text{-}cyclic\text{-}for}(v, {s}', i) \end{array}}}

\vnbbookfile{theorem-library/gns-completion}{theorem-library/gns-completion}
\noindent{\small gns-completion.scm -- THEOREM 4.13 ON THE HILBERT COMPLETION: the GNS representation of a unital involutive Banach algebra on H\_phi, the operators of the pre-Hilbert stage (theorem-library/ba4-gns.scm) "extended by continuity" through COMPLETION-LIFT (theorem-library/completion-lift.scm).  F. J. Thayer, "Banach ...}\par
\vnbprop{gns-hmod-is-module}{\vnbstmt{\begin{array}{@{}l@{}} \forall a, f. \\ \quad (\operatorname{is\text{-}unital\text{-}involutive\text{-}banach\text{-}algebra}(a) \wedge \operatorname{is\text{-}positive\text{-}lf}(a, f)) \Rightarrow  \\ \quad \quad \operatorname{is\text{-}module}(\operatorname{gns\text{-}hmod}(a, f)) \end{array}}}
\vnbprop{gns-hmod-scal}{\vnbstmt{\begin{array}{@{}l@{}} \forall a, f. \\ \quad (\operatorname{is\text{-}unital\text{-}involutive\text{-}banach\text{-}algebra}(a) \wedge \operatorname{is\text{-}positive\text{-}lf}(a, f)) \Rightarrow  \\ \quad \quad \operatorname{scal}(\operatorname{gns\text{-}hmod}(a, f)) = \operatorname{normed\text{-}field\text{-}as\text{-}commutative\text{-}ring}(\operatorname{cc\text{-}normed\text{-}field}) \end{array}}}
\vnbprop{gns-hform-ipf}{\vnbstmt{\begin{array}{@{}l@{}} \forall a, f. \\ \quad (\operatorname{is\text{-}unital\text{-}involutive\text{-}banach\text{-}algebra}(a) \wedge \operatorname{is\text{-}positive\text{-}lf}(a, f)) \Rightarrow  \\ \quad \quad \operatorname{is\text{-}inner\text{-}product\text{-}form}(\operatorname{gns\text{-}hmod}(a, f), \operatorname{gns\text{-}hform}(a, f)) \end{array}}}
\vnbprop{gns-ms-eq}{\vnbstmt{\begin{array}{@{}l@{}} \forall a, f. \\ \quad (\operatorname{is\text{-}unital\text{-}involutive\text{-}banach\text{-}algebra}(a) \wedge \operatorname{is\text{-}positive\text{-}lf}(a, f)) \Rightarrow  \\ \quad \quad (\operatorname{form\text{-}ms}(\operatorname{gns\text{-}hmod}(a, f), \operatorname{gns\text{-}hform}(a, f)) \simeq \operatorname{ip\text{-}metric\text{-}space}(\operatorname{gns\text{-}pre}(a, f))) \end{array}}}
\vnbprop{gns-hmod-vec}{\vnbstmt{\begin{array}{@{}l@{}} \forall a, f. \\ \quad (\operatorname{is\text{-}unital\text{-}involutive\text{-}banach\text{-}algebra}(a) \wedge \operatorname{is\text{-}positive\text{-}lf}(a, f)) \Rightarrow  \\ \quad \quad (\operatorname{vec}(\operatorname{gns\text{-}hmod}(a, f)) \simeq \operatorname{vec}(\operatorname{gns\text{-}pre}(a, f))) \end{array}}}
\vnbprop{gns-hmod-vadd}{\vnbstmt{\begin{array}{@{}l@{}} \forall a, f. \\ \quad (\operatorname{is\text{-}unital\text{-}involutive\text{-}banach\text{-}algebra}(a) \wedge \operatorname{is\text{-}positive\text{-}lf}(a, f)) \Rightarrow  \\ \quad \quad (\operatorname{vadd}(\operatorname{gns\text{-}hmod}(a, f)) \simeq \operatorname{vadd}(\operatorname{gns\text{-}pre}(a, f))) \end{array}}}
\vnbprop{gns-hmod-act}{\vnbstmt{\begin{array}{@{}l@{}} \forall a, f. \\ \quad (\operatorname{is\text{-}unital\text{-}involutive\text{-}banach\text{-}algebra}(a) \wedge \operatorname{is\text{-}positive\text{-}lf}(a, f)) \Rightarrow  \\ \quad \quad (\operatorname{act}(\operatorname{gns\text{-}hmod}(a, f)) \simeq \operatorname{act}(\operatorname{gns\text{-}pre}(a, f))) \end{array}}}
\vnbprop{gns-hmod-vneg}{\vnbstmt{\begin{array}{@{}l@{}} \forall a, f. \\ \quad (\operatorname{is\text{-}unital\text{-}involutive\text{-}banach\text{-}algebra}(a) \wedge \operatorname{is\text{-}positive\text{-}lf}(a, f)) \Rightarrow  \\ \quad \quad (\operatorname{vneg}(\operatorname{gns\text{-}hmod}(a, f)) \simeq \operatorname{vneg}(\operatorname{gns\text{-}pre}(a, f))) \end{array}}}
\vnbprop{gns-hform-ip}{\vnbstmt{\begin{array}{@{}l@{}} \forall a, f. \\ \quad (\operatorname{is\text{-}unital\text{-}involutive\text{-}banach\text{-}algebra}(a) \wedge \operatorname{is\text{-}positive\text{-}lf}(a, f)) \Rightarrow  \\ \quad \quad (\operatorname{gns\text{-}hform}(a, f) \simeq \langle\cdot,\cdot\rangle_{\operatorname{gns\text{-}pre}(a, f)}) \end{array}}}
\vnbprop{gns-ms-pts}{\vnbstmt{\begin{array}{@{}l@{}} \forall a, f. \\ \quad (\operatorname{is\text{-}unital\text{-}involutive\text{-}banach\text{-}algebra}(a) \wedge \operatorname{is\text{-}positive\text{-}lf}(a, f)) \Rightarrow  \\ \quad \quad (\operatorname{pts}(\operatorname{form\text{-}ms}(\operatorname{gns\text{-}hmod}(a, f), \operatorname{gns\text{-}hform}(a, f))) \simeq \operatorname{vec}(\operatorname{gns\text{-}pre}(a, f))) \end{array}}}
\vnbprop{gns-op-vneg}{\vnbstmt{\begin{array}{@{}l@{}} \forall a, f.; \forall x \in \operatorname{vec}(a), u \in \operatorname{vec}(\operatorname{gns\text{-}pre}(a, f)). \\ \quad (\operatorname{is\text{-}unital\text{-}involutive\text{-}banach\text{-}algebra}(a) \wedge \operatorname{is\text{-}positive\text{-}lf}(a, f)) \Rightarrow  \\ \quad \quad \operatorname{gns\text{-}op}(a, f, x)(\operatorname{vneg}(\operatorname{gns\text{-}pre}(a, f))(u)) = \operatorname{vneg}(\operatorname{gns\text{-}pre}(a, f))(\operatorname{gns\text{-}op}(a, f, x)(u)) \end{array}}}
\vnbprop{gns-op-norm-bound}{\vnbstmt{\begin{array}{@{}l@{}} \forall a, f.; \forall x \in \operatorname{vec}(a), u \in \operatorname{vec}(\operatorname{gns\text{-}pre}(a, f)). \\ \quad (\operatorname{is\text{-}unital\text{-}involutive\text{-}banach\text{-}algebra}(a) \wedge \operatorname{is\text{-}positive\text{-}lf}(a, f)) \Rightarrow  \\ \quad \quad \lVert \operatorname{gns\text{-}op}(a, f, x)(u) \rVert \leq (\operatorname{vnrm}(a)(x) \cdot \lVert u \rVert) \end{array}}}
\vnbprop{gns-op-pm-type}{\vnbstmt{\begin{array}{@{}l@{}} \forall a, f.; \forall x \in \operatorname{vec}(a). \\ \quad (\operatorname{is\text{-}unital\text{-}involutive\text{-}banach\text{-}algebra}(a) \wedge \operatorname{is\text{-}positive\text{-}lf}(a, f)) \Rightarrow  \\ \quad \quad \operatorname{gns\text{-}op}(a, f, x) \\ \quad \quad \in\; (\operatorname{pts}(\operatorname{form\text{-}ms}(\operatorname{gns\text{-}hmod}(a, f), \operatorname{gns\text{-}hform}(a, f))) \to \operatorname{pts}(\operatorname{form\text{-}ms}(\operatorname{gns\text{-}hmod}(a, f), \operatorname{gns\text{-}hform}(a, f)))) \end{array}}}
\vnbprop{gns-op-lipschitz}{\vnbstmt{\begin{array}{@{}l@{}} \forall a, f.; \forall x \in \operatorname{vec}(a), u, w \in \operatorname{pts}(\operatorname{form\text{-}ms}(\operatorname{gns\text{-}hmod}(a, f), \operatorname{gns\text{-}hform}(a, f))). \\ \quad (\operatorname{is\text{-}unital\text{-}involutive\text{-}banach\text{-}algebra}(a) \wedge \operatorname{is\text{-}positive\text{-}lf}(a, f)) \Rightarrow  \\ \quad \quad \operatorname{dist}(\operatorname{form\text{-}ms}(\operatorname{gns\text{-}hmod}(a, f), \operatorname{gns\text{-}hform}(a, f)))(\operatorname{gns\text{-}op}(a, f, x)(u), \operatorname{gns\text{-}op}(a, f, x)(w)) \\ \quad \quad \leq\; \operatorname{vnrm}(a)(x) \cdot \operatorname{dist}(\operatorname{form\text{-}ms}(\operatorname{gns\text{-}hmod}(a, f), \operatorname{gns\text{-}hform}(a, f)))(u, w) \end{array}}}
\vnbprop{gns-op-add-pm}{\vnbstmt{\begin{array}{@{}l@{}} \forall a, f.; \forall x \in \operatorname{vec}(a), v, z \in \operatorname{pts}(\operatorname{form\text{-}ms}(\operatorname{gns\text{-}hmod}(a, f), \operatorname{gns\text{-}hform}(a, f))). \\ \quad (\operatorname{is\text{-}unital\text{-}involutive\text{-}banach\text{-}algebra}(a) \wedge \operatorname{is\text{-}positive\text{-}lf}(a, f)) \Rightarrow  \\ \quad \quad \operatorname{gns\text{-}op}(a, f, x)(\operatorname{vadd}(\operatorname{gns\text{-}hmod}(a, f))(v, z)) \\ \quad \quad =\; \operatorname{vadd}(\operatorname{gns\text{-}hmod}(a, f))(\operatorname{gns\text{-}op}(a, f, x)(v), \operatorname{gns\text{-}op}(a, f, x)(z)) \end{array}}}
\vnbprop{gns-op-hom-pm}{\vnbstmt{\begin{array}{@{}l@{}} \forall a, f.; \forall x \in \operatorname{vec}(a), c \in \mathbb{C}, v \in \operatorname{pts}(\operatorname{form\text{-}ms}(\operatorname{gns\text{-}hmod}(a, f), \operatorname{gns\text{-}hform}(a, f))). \\ \quad (\operatorname{is\text{-}unital\text{-}involutive\text{-}banach\text{-}algebra}(a) \wedge \operatorname{is\text{-}positive\text{-}lf}(a, f)) \Rightarrow  \\ \quad \quad \operatorname{gns\text{-}op}(a, f, x)(\operatorname{act}(\operatorname{gns\text{-}hmod}(a, f))(c, v)) = \operatorname{act}(\operatorname{gns\text{-}hmod}(a, f))(c, \operatorname{gns\text{-}op}(a, f, x)(v)) \end{array}}}
\vnbprop{gns-op-bound-pm}{\vnbstmt{\begin{array}{@{}l@{}} \forall a, f.; \forall x \in \operatorname{vec}(a), v \in \operatorname{pts}(\operatorname{form\text{-}ms}(\operatorname{gns\text{-}hmod}(a, f), \operatorname{gns\text{-}hform}(a, f))). \\ \quad (\operatorname{is\text{-}unital\text{-}involutive\text{-}banach\text{-}algebra}(a) \wedge \operatorname{is\text{-}positive\text{-}lf}(a, f)) \Rightarrow  \\ \quad \quad \operatorname{gns\text{-}hform}(a, f)(\operatorname{gns\text{-}op}(a, f, x)(v), \operatorname{gns\text{-}op}(a, f, x)(v)) \\ \quad \quad \leq\; (\operatorname{vnrm}(a)(x) \cdot \operatorname{vnrm}(a)(x)) \cdot \operatorname{gns\text{-}hform}(a, f)(v, v) \end{array}}}
\vnbprop{gns-op-sum-pm}{\vnbstmt{\begin{array}{@{}l@{}} \forall a, f.; \forall x, y \in \operatorname{vec}(a), v \in \operatorname{pts}(\operatorname{form\text{-}ms}(\operatorname{gns\text{-}hmod}(a, f), \operatorname{gns\text{-}hform}(a, f))). \\ \quad (\operatorname{is\text{-}unital\text{-}involutive\text{-}banach\text{-}algebra}(a) \wedge \operatorname{is\text{-}positive\text{-}lf}(a, f)) \Rightarrow  \\ \quad \quad \operatorname{gns\text{-}op}(a, f, \operatorname{vadd}(a)(x, y))(v) \\ \quad \quad =\; \operatorname{vadd}(\operatorname{gns\text{-}hmod}(a, f))(\operatorname{gns\text{-}op}(a, f, x)(v), \operatorname{gns\text{-}op}(a, f, y)(v)) \end{array}}}
\vnbprop{gns-op-mul-pm}{\vnbstmt{\begin{array}{@{}l@{}} \forall a, f.; \forall x, y \in \operatorname{vec}(a), v \in \operatorname{pts}(\operatorname{form\text{-}ms}(\operatorname{gns\text{-}hmod}(a, f), \operatorname{gns\text{-}hform}(a, f))). \\ \quad (\operatorname{is\text{-}unital\text{-}involutive\text{-}banach\text{-}algebra}(a) \wedge \operatorname{is\text{-}positive\text{-}lf}(a, f)) \Rightarrow  \\ \quad \quad \operatorname{gns\text{-}op}(a, f, \cdot_{a}(x, y))(v) = \operatorname{gns\text{-}op}(a, f, x)(\operatorname{gns\text{-}op}(a, f, y)(v)) \end{array}}}
\vnbprop{gns-op-unit-pm}{\vnbstmt{\begin{array}{@{}l@{}} \forall a, f.; \forall v \in \operatorname{pts}(\operatorname{form\text{-}ms}(\operatorname{gns\text{-}hmod}(a, f), \operatorname{gns\text{-}hform}(a, f))). \\ \quad (\operatorname{is\text{-}unital\text{-}involutive\text{-}banach\text{-}algebra}(a) \wedge \operatorname{is\text{-}positive\text{-}lf}(a, f)) \Rightarrow  \\ \quad \quad \operatorname{gns\text{-}op}(a, f, 1_{a})(v) = v \end{array}}}
\vnbprop{gns-op-scale-pm}{\vnbstmt{\begin{array}{@{}l@{}} \forall a, f.; \forall c \in \mathbb{C}, x \in \operatorname{vec}(a), v \in \operatorname{pts}(\operatorname{form\text{-}ms}(\operatorname{gns\text{-}hmod}(a, f), \operatorname{gns\text{-}hform}(a, f))). \\ \quad (\operatorname{is\text{-}unital\text{-}involutive\text{-}banach\text{-}algebra}(a) \wedge \operatorname{is\text{-}positive\text{-}lf}(a, f)) \Rightarrow  \\ \quad \quad \operatorname{gns\text{-}op}(a, f, \operatorname{act}(a)(c, x))(v) = \operatorname{act}(\operatorname{gns\text{-}hmod}(a, f))(c, \operatorname{gns\text{-}op}(a, f, x)(v)) \end{array}}}
\vnbprop{gns-op-adj-pm}{\vnbstmt{\begin{array}{@{}l@{}} \forall a, f.; \forall x \in \operatorname{vec}(a), v, z \in \operatorname{pts}(\operatorname{form\text{-}ms}(\operatorname{gns\text{-}hmod}(a, f), \operatorname{gns\text{-}hform}(a, f))). \\ \quad (\operatorname{is\text{-}unital\text{-}involutive\text{-}banach\text{-}algebra}(a) \wedge \operatorname{is\text{-}positive\text{-}lf}(a, f)) \Rightarrow  \\ \quad \quad \operatorname{gns\text{-}hform}(a, f)(\operatorname{gns\text{-}op}(a, f, x)(v), z) = \operatorname{gns\text{-}hform}(a, f)(v, \operatorname{gns\text{-}op}(a, f, {}^{*_{a}}(x))(z)) \end{array}}}
\vnbprop{gns-lift-bounded}{\vnbstmt{\begin{array}{@{}l@{}} \forall a, f.; \forall x \in \operatorname{vec}(a). \\ \quad (\operatorname{is\text{-}unital\text{-}involutive\text{-}banach\text{-}algebra}(a) \wedge \operatorname{is\text{-}positive\text{-}lf}(a, f)) \Rightarrow  \\ \quad \quad \operatorname{completion\text{-}lift}(\operatorname{form\text{-}ms}(\operatorname{gns\text{-}hmod}(a, f), \operatorname{gns\text{-}hform}(a, f)), \operatorname{gns\text{-}op}(a, f, x)) \\ \quad \quad \in\; L({\operatorname{hilbert\text{-}completion}(\operatorname{gns\text{-}hmod}(a, f), \operatorname{gns\text{-}hform}(a, f))}, {\operatorname{hilbert\text{-}completion}(\operatorname{gns\text{-}hmod}(a, f), \operatorname{gns\text{-}hform}(a, f))}) \end{array}}}
\vnbprop{gns-space-eq}{\vnbstmt{\begin{array}{@{}l@{}} \forall a, f. \\ \quad (\operatorname{is\text{-}unital\text{-}involutive\text{-}banach\text{-}algebra}(a) \wedge \operatorname{is\text{-}positive\text{-}lf}(a, f)) \Rightarrow  \\ \quad \quad (H_{f} \simeq \operatorname{hilbert\text{-}completion}(\operatorname{gns\text{-}hmod}(a, f), \operatorname{gns\text{-}hform}(a, f))) \end{array}}}
\vnbprop{gns-space-is-hilbert}{\vnbstmt{\begin{array}{@{}l@{}} \forall a, f. \\ \quad (\operatorname{is\text{-}unital\text{-}involutive\text{-}banach\text{-}algebra}(a) \wedge \operatorname{is\text{-}positive\text{-}lf}(a, f)) \Rightarrow  \\ \quad \quad \operatorname{is\text{-}hilbert\text{-}space}(H_{f}) \end{array}}}
\vnbprop{gns-space-is-cips}{\vnbstmt{\begin{array}{@{}l@{}} \forall a, f. \\ \quad (\operatorname{is\text{-}unital\text{-}involutive\text{-}banach\text{-}algebra}(a) \wedge \operatorname{is\text{-}positive\text{-}lf}(a, f)) \Rightarrow  \\ \quad \quad \operatorname{is\text{-}complex\text{-}inner\text{-}product\text{-}space}(H_{f}) \end{array}}}
\vnbprop{gns-rep-value}{\vnbstmt{\begin{array}{@{}l@{}} \forall a, f.; \forall x \in \operatorname{vec}(a). \\ \quad (\operatorname{is\text{-}unital\text{-}involutive\text{-}banach\text{-}algebra}(a) \wedge \operatorname{is\text{-}positive\text{-}lf}(a, f)) \Rightarrow  \\ \quad \quad \pi_{f}(x) = \operatorname{completion\text{-}lift}(\operatorname{form\text{-}ms}(\operatorname{gns\text{-}hmod}(a, f), \operatorname{gns\text{-}hform}(a, f)), \operatorname{gns\text{-}op}(a, f, x)) \end{array}}}
\vnbprop{gns-rep-bounded}{\vnbstmt{\begin{array}{@{}l@{}} \forall a, f.; \forall x \in \operatorname{vec}(a). \\ \quad (\operatorname{is\text{-}unital\text{-}involutive\text{-}banach\text{-}algebra}(a) \wedge \operatorname{is\text{-}positive\text{-}lf}(a, f)) \Rightarrow  \\ \quad \quad \pi_{f}(x) \in L({H_{f}}, {H_{f}}) \end{array}}}
\vnbprop{gns-rep-type}{\vnbstmt{\begin{array}{@{}l@{}} \forall a, f. \\ \quad (\operatorname{is\text{-}unital\text{-}involutive\text{-}banach\text{-}algebra}(a) \wedge \operatorname{is\text{-}positive\text{-}lf}(a, f)) \Rightarrow  \\ \quad \quad \pi_{f} \in (\operatorname{vec}(a) \to L({H_{f}}, {H_{f}})) \end{array}}}
\vnbprop{gns-rep-sum}{\vnbstmt{\begin{array}{@{}l@{}} \forall a, f.; \forall x, y \in \operatorname{vec}(a). \\ \quad (\operatorname{is\text{-}unital\text{-}involutive\text{-}banach\text{-}algebra}(a) \wedge \operatorname{is\text{-}positive\text{-}lf}(a, f)) \Rightarrow  \\ \quad \quad \pi_{f}(\operatorname{vadd}(a)(x, y)) \\ \quad \quad =\; {\pi_{f}(x)} + {\pi_{f}(y)} \end{array}}}
\vnbprop{gns-rep-scale}{\vnbstmt{\begin{array}{@{}l@{}} \forall a, f.; \forall c \in \mathbb{C}, x \in \operatorname{vec}(a). \\ \quad (\operatorname{is\text{-}unital\text{-}involutive\text{-}banach\text{-}algebra}(a) \wedge \operatorname{is\text{-}positive\text{-}lf}(a, f)) \Rightarrow  \\ \quad \quad \pi_{f}(\operatorname{act}(a)(c, x)) = {c}\,{\pi_{f}(x)} \end{array}}}
\vnbprop{gns-rep-mul}{\vnbstmt{\begin{array}{@{}l@{}} \forall a, f.; \forall x, y \in \operatorname{vec}(a). \\ \quad (\operatorname{is\text{-}unital\text{-}involutive\text{-}banach\text{-}algebra}(a) \wedge \operatorname{is\text{-}positive\text{-}lf}(a, f)) \Rightarrow  \\ \quad \quad \pi_{f}(\cdot_{a}(x, y)) = {\pi_{f}(x)}{\pi_{f}(y)} \end{array}}}
\vnbprop{gns-rep-unit}{\vnbstmt{\begin{array}{@{}l@{}} \forall a, f. \\ \quad (\operatorname{is\text{-}unital\text{-}involutive\text{-}banach\text{-}algebra}(a) \wedge \operatorname{is\text{-}positive\text{-}lf}(a, f)) \Rightarrow  \\ \quad \quad \pi_{f}(1_{a}) = 1 \end{array}}}
\vnbprop{gns-rep-star}{\vnbstmt{\begin{array}{@{}l@{}} \forall a, f.; \forall x \in \operatorname{vec}(a). \\ \quad (\operatorname{is\text{-}unital\text{-}involutive\text{-}banach\text{-}algebra}(a) \wedge \operatorname{is\text{-}positive\text{-}lf}(a, f)) \Rightarrow  \\ \quad \quad \pi_{f}({}^{*_{a}}(x)) = {\pi_{f}(x)}^{*} \end{array}}}
\vnbprop{gns-rep-is-representation}{\vnbstmt{\begin{array}{@{}l@{}} \forall a, f. \\ \quad (\operatorname{is\text{-}unital\text{-}involutive\text{-}banach\text{-}algebra}(a) \wedge \operatorname{is\text{-}positive\text{-}lf}(a, f)) \Rightarrow  \\ \quad \quad \operatorname{is\text{-}representation}(a, H_{f}, \pi_{f}) \end{array}}}
\vnbprop{gns-rep-represents}{\vnbstmt{\begin{array}{@{}l@{}} \forall a, f.; \forall x \in \operatorname{vec}(a). \\ \quad (\operatorname{is\text{-}unital\text{-}involutive\text{-}banach\text{-}algebra}(a) \wedge \operatorname{is\text{-}positive\text{-}lf}(a, f)) \Rightarrow  \\ \quad \quad f(x) = \langle\cdot,\cdot\rangle_{H_{f}}(\pi_{f}(x)(\xi_{f}), \xi_{f}) \end{array}}}
\vnbprop{gns-rep-cyclic}{\vnbstmt{\begin{array}{@{}l@{}} \forall a, f. \\ \quad (\operatorname{is\text{-}unital\text{-}involutive\text{-}banach\text{-}algebra}(a) \wedge \operatorname{is\text{-}positive\text{-}lf}(a, f)) \Rightarrow  \\ \quad \quad \operatorname{is\text{-}cyclic\text{-}vector\text{-}rep}(a, H_{f}, \pi_{f}, \xi_{f}) \end{array}}}
\vnbprop{gns-theorem}{\vnbstmt{\begin{array}{@{}l@{}} \forall a, f. \\ \quad (\operatorname{is\text{-}unital\text{-}involutive\text{-}banach\text{-}algebra}(a) \wedge \operatorname{is\text{-}positive\text{-}lf}(a, f)) \Rightarrow  \\ \quad \quad \exists v, p, e. \\ \quad \quad \quad \operatorname{is\text{-}representation}(a, v, p) \wedge \\ \quad \quad \quad \operatorname{is\text{-}cyclic\text{-}vector\text{-}rep}(a, v, p, e) \wedge \\ \quad \quad \quad \forall x \in \operatorname{vec}(a).\; (f(x) = \langle\cdot,\cdot\rangle_{v}(p(x)(e), e)) \end{array}}}
\vnbprop{radical-unfold}{\vnbstmt{\begin{array}{@{}l@{}} \forall a. \\ \quad (\sqrt{a} \simeq \{x \in \operatorname{vec}(a) : \forall v, q.\; (\operatorname{is\text{-}representation}(a, v, q) \Rightarrow (q(x) = 0))\}) \end{array}}}
\vnbprop{radical-positive-vanishes}{\vnbstmt{\begin{array}{@{}l@{}} \forall a, f.; \forall x \in \sqrt{a}. \\ \quad (\operatorname{is\text{-}unital\text{-}involutive\text{-}banach\text{-}algebra}(a) \wedge \operatorname{is\text{-}positive\text{-}lf}(a, f)) \Rightarrow (f(x) = 0) \end{array}}}

\vnbbookfile{theorem-library/amp-sum}{theorem-library/amp-sum}
\noindent{\small amp-sum.scm -- agent BA-1117-TRACE (round 3), 2026-11-17. A companion of amp-vector-exists (theorem-library/ba-amp.scm, unchanged) that exposes what the iterated Hilbert sum H (+) ... (+) H with f(T) = T (+) ... (+) T gives: for all X, Y in w, $<$f(T) X | Y$>$\_w = sum\_(k $<$ n) $<$T a\_k | b\_k$>$ for some n in NN and sequences a, ...}\par
\vnbprop{amp-vector-step-explicit}{\vnbstmt{\begin{array}{@{}l@{}} \forall v, s, w, f, x, u, z. \\ \quad \operatorname{is\text{-}hilbert\text{-}space}(v) \wedge \\ \quad s \subseteq L({v}, {v}) \wedge \\ \quad \operatorname{is\text{-}amplification}(v, s, w, f) \wedge \\ \quad x \in \operatorname{vec}(w) \wedge \\ \quad z \in \operatorname{vec}(v) \wedge \\ \quad u \subseteq \operatorname{vec}(v) \wedge \\ \quad \forall a, b \in L({v}, {v}), y \in u. \\ \quad \quad \lVert \operatorname{vadd}(v)(a(y), \operatorname{vneg}(v)(b(y))) \rVert \\ \quad \quad \leq\; \lVert \operatorname{vadd}(w)(f(a)(x), \operatorname{vneg}(w)(f(b)(x))) \rVert \Rightarrow \\ \quad \quad \operatorname{is\text{-}amplification}(v, s, {w} \oplus {v}, {f} \oplus \mathrm{id}) \wedge \\ \quad \quad {[}x, z] \in \operatorname{vec}({w} \oplus {v}) \wedge \\ \quad \quad \forall a, b \in L({v}, {v}), y \in (u \cup \{z\}). \\ \quad \quad \quad \lVert \operatorname{vadd}(v)(a(y), \operatorname{vneg}(v)(b(y))) \rVert \\ \quad \quad \quad \leq\; \lVert \operatorname{vadd}({w} \oplus {v})({f} \oplus \mathrm{id}(a)({[}x, z]), \operatorname{vneg}({w} \oplus {v})({f} \oplus \mathrm{id}(b)({[}x, z]))) \rVert \end{array}}}
\vnbprop{amp-rep-base}{\vnbstmt{\begin{array}{@{}l@{}} \forall v.; \forall p, q \in \operatorname{vec}(v). \\ \quad \operatorname{is\text{-}complex\text{-}inner\text{-}product\text{-}space}(v) \Rightarrow  \\ \quad \quad \exists n \in \mathbb{N}, a, b \in (\mathbb{N} \to \operatorname{vec}(v)). \forall t \in L({v}, {v}). \\ \quad \quad \quad \langle\cdot,\cdot\rangle_{v}(\mathrm{id}(t)(p), q) \\ \quad \quad \quad =\; \operatorname{finsum}(\operatorname{normed\text{-}ag\text{-}as\text{-}abelian\text{-}group}(\operatorname{cc\text{-}nag}), \\ \quad \quad \quad \quad (\lambda i \in \mathrm{OS}(n).\; \langle\cdot,\cdot\rangle_{v}(t(a(i)), b(i))), \\ \quad \quad \quad \quad \mathrm{OS}(n)) \end{array}}}
\vnbprop{amp-rep-step}{\vnbstmt{\begin{array}{@{}l@{}} \forall w, v, f.; \forall p, q \in \operatorname{vec}({w} \oplus {v}). \\ \quad \operatorname{is\text{-}complex\text{-}inner\text{-}product\text{-}space}(w) \wedge \\ \quad \operatorname{is\text{-}complex\text{-}inner\text{-}product\text{-}space}(v) \wedge \\ \quad f \in (L({v}, {v}) \to L({w}, {w})) \wedge \\ \quad \forall p, q \in \operatorname{vec}(w). \exists n \in \mathbb{N}, a, b \in (\mathbb{N} \to \operatorname{vec}(v)). \forall t \in L({v}, {v}). \\ \quad \quad \langle\cdot,\cdot\rangle_{w}(f(t)(p), q) \\ \quad \quad =\; \operatorname{finsum}(\operatorname{normed\text{-}ag\text{-}as\text{-}abelian\text{-}group}(\operatorname{cc\text{-}nag}), \\ \quad \quad \quad (\lambda i \in \mathrm{OS}(n).\; \langle\cdot,\cdot\rangle_{v}(t(a(i)), b(i))), \\ \quad \quad \quad \mathrm{OS}(n)) \Rightarrow \\ \quad \quad \exists n \in \mathbb{N}, a, b \in (\mathbb{N} \to \operatorname{vec}(v)). \forall t \in L({v}, {v}). \\ \quad \quad \quad \langle\cdot,\cdot\rangle_{{w} \oplus {v}}({f} \oplus \mathrm{id}(t)(p), q) \\ \quad \quad \quad =\; \operatorname{finsum}(\operatorname{normed\text{-}ag\text{-}as\text{-}abelian\text{-}group}(\operatorname{cc\text{-}nag}), \\ \quad \quad \quad \quad (\lambda i \in \mathrm{OS}(n).\; \langle\cdot,\cdot\rangle_{v}(t(a(i)), b(i))), \\ \quad \quad \quad \quad \mathrm{OS}(n)) \end{array}}}
\vnbprop{amp-vector-exists-sum}{\vnbstmt{\begin{array}{@{}l@{}} \forall v, s, g. \\ \quad \operatorname{is\text{-}hilbert\text{-}space}(v) \wedge \\ \quad s \subseteq L({v}, {v}) \wedge \\ \quad g \in \mathbf{Set} \wedge \\ \quad |g| \in \mathbb{N} \wedge \\ \quad g \subseteq \operatorname{vec}(v) \Rightarrow \\ \quad \quad \exists w, f \in \mathbf{Set}. \\ \quad \quad \quad \operatorname{is\text{-}amplification}(v, s, w, f) \wedge \\ \quad \quad \quad \forall p, q \in \operatorname{vec}(w). \exists n \in \mathbb{N}, a, b \in (\mathbb{N} \to \operatorname{vec}(v)). \forall t \in L({v}, {v}). \\ \quad \quad \quad \quad \langle\cdot,\cdot\rangle_{w}(f(t)(p), q) \\ \quad \quad \quad \quad =\; \operatorname{finsum}(\operatorname{normed\text{-}ag\text{-}as\text{-}abelian\text{-}group}(\operatorname{cc\text{-}nag}), \\ \quad \quad \quad \quad \quad (\lambda i \in \mathrm{OS}(n).\; \langle\cdot,\cdot\rangle_{v}(t(a(i)), b(i))), \\ \quad \quad \quad \quad \quad \mathrm{OS}(n)) \wedge \\ \quad \quad \quad \exists x \in \operatorname{vec}(w). \forall \mathit{cqa}, \mathit{cqb} \in L({v}, {v}), y \in g. \\ \quad \quad \quad \quad \lVert \operatorname{vadd}(v)(\operatorname{cqa}(y), \operatorname{vneg}(v)(\operatorname{cqb}(y))) \rVert \\ \quad \quad \quad \quad \leq\; \lVert \operatorname{vadd}(w)(f(\mathit{cqa})(x), \operatorname{vneg}(w)(f(\mathit{cqb})(x))) \rVert \end{array}}}
\vnbprop{strong-dual-linear}{\vnbstmt{\begin{array}{@{}l@{}} \forall v, i. \\ \quad (\operatorname{is\text{-}complex\text{-}inner\text{-}product\text{-}space}(v) \wedge (i \in L_{\bullet}({v}))) \Rightarrow  \\ \quad \quad i \in (L({v}, {v}) \to \mathbb{C}) \wedge \\ \quad \quad \forall s, u \in L({v}, {v}). \\ \quad \quad \quad i({s} + {u}) = (i(s) + i(u)) \wedge \\ \quad \quad \forall c \in \mathbb{C}, s \in L({v}, {v}).\; (i({c}\,{s}) = (c \cdot i(s))) \end{array}}}
\vnbprop{strong-dual-finite-sum}{\vnbstmt{\begin{array}{@{}l@{}} \forall v, i. \\ \quad (\operatorname{is\text{-}hilbert\text{-}space}(v) \wedge (i \in L_{\bullet}({v}))) \Rightarrow  \\ \quad \quad \exists n \in \mathbb{N}, a, b \in (\mathbb{N} \to \operatorname{vec}(v)). \forall t \in L({v}, {v}). \\ \quad \quad \quad i(t) \\ \quad \quad \quad =\; \operatorname{finsum}(\operatorname{normed\text{-}ag\text{-}as\text{-}abelian\text{-}group}(\operatorname{cc\text{-}nag}), \\ \quad \quad \quad \quad (\lambda \mathit{avi} \in \mathrm{OS}(n).\; \langle\cdot,\cdot\rangle_{v}(t(a(\mathit{avi})), b(\mathit{avi}))), \\ \quad \quad \quad \quad \mathrm{OS}(n)) \end{array}}}
\vnbprop{finsum-vf-bound}{\vnbstmt{\begin{array}{@{}l@{}} \forall v, u, a, b, t.; \forall n \in \mathbb{N}. \\ \quad \operatorname{is\text{-}complex\text{-}inner\text{-}product\text{-}space}(v) \wedge \\ \quad u \in L({v}, {v}) \wedge \\ \quad a \in (\mathbb{N} \to \operatorname{vec}(v)) \wedge \\ \quad b \in (\mathbb{N} \to \operatorname{vec}(v)) \wedge \\ \quad t \in \mathbb{R} \wedge \\ \quad 0 \leq t \wedge \\ \quad \forall k \in \mathrm{OS}(n).\; (\operatorname{magnitude}(\langle\cdot,\cdot\rangle_{v}(u(a(k)), b(k))) \leq t) \Rightarrow \\ \quad \quad \operatorname{magnitude}(\operatorname{finsum}(\operatorname{normed\text{-}ag\text{-}as\text{-}abelian\text{-}group}(\operatorname{cc\text{-}nag}), \\ \quad \quad \quad \quad (\lambda i \in \mathrm{OS}(n).\; \langle\cdot,\cdot\rangle_{v}(u(a(i)), b(i))), \\ \quad \quad \quad \quad \mathrm{OS}(n))) \\ \quad \quad \leq\; n \cdot t \end{array}}}
\vnbprop{finite-sum-in-weak-dual}{\vnbstmt{\begin{array}{@{}l@{}} \forall v, i, n, a, b. \\ \quad \operatorname{is\text{-}complex\text{-}inner\text{-}product\text{-}space}(v) \wedge \\ \quad i \in (L({v}, {v}) \to \mathbb{C}) \wedge \\ \quad \operatorname{is\text{-}linear\text{-}map}(L({v}), \operatorname{cc\text{-}sn\text{-}module}, i) \wedge \\ \quad n \in \mathbb{N} \wedge \\ \quad a \in (\mathbb{N} \to \operatorname{vec}(v)) \wedge \\ \quad b \in (\mathbb{N} \to \operatorname{vec}(v)) \wedge \\ \quad \forall t \in L({v}, {v}). \\ \quad \quad i(t) \\ \quad \quad =\; \operatorname{finsum}(\operatorname{normed\text{-}ag\text{-}as\text{-}abelian\text{-}group}(\operatorname{cc\text{-}nag}), \\ \quad \quad \quad (\lambda \mathit{avi} \in \mathrm{OS}(n).\; \langle\cdot,\cdot\rangle_{v}(t(a(\mathit{avi})), b(\mathit{avi}))), \\ \quad \quad \quad \mathrm{OS}(n)) \Rightarrow \\ \quad \quad i \in \operatorname{weak\text{-}dual}(v) \end{array}}}
\vnbprop{finite-sum-in-strong-dual}{\vnbstmt{\begin{array}{@{}l@{}} \forall v, i, n, a, b. \\ \quad \operatorname{is\text{-}complex\text{-}inner\text{-}product\text{-}space}(v) \wedge \\ \quad i \in (L({v}, {v}) \to \mathbb{C}) \wedge \\ \quad \operatorname{is\text{-}linear\text{-}map}(L({v}), \operatorname{cc\text{-}sn\text{-}module}, i) \wedge \\ \quad n \in \mathbb{N} \wedge \\ \quad a \in (\mathbb{N} \to \operatorname{vec}(v)) \wedge \\ \quad b \in (\mathbb{N} \to \operatorname{vec}(v)) \wedge \\ \quad \forall t \in L({v}, {v}). \\ \quad \quad i(t) \\ \quad \quad =\; \operatorname{finsum}(\operatorname{normed\text{-}ag\text{-}as\text{-}abelian\text{-}group}(\operatorname{cc\text{-}nag}), \\ \quad \quad \quad (\lambda \mathit{avi} \in \mathrm{OS}(n).\; \langle\cdot,\cdot\rangle_{v}(t(a(\mathit{avi})), b(\mathit{avi}))), \\ \quad \quad \quad \mathrm{OS}(n)) \Rightarrow \\ \quad \quad i \in L_{\bullet}({v}) \end{array}}}
\vnbprop{strong-dual-subset-weak-dual}{\vnbstmt{\begin{array}{@{}l@{}} \forall v, i. \\ \quad (\operatorname{is\text{-}hilbert\text{-}space}(v) \wedge (i \in L_{\bullet}({v}))) \Rightarrow (i \in \operatorname{weak\text{-}dual}(v)) \end{array}}}
\vnbprop{strong-dual-iff-weak-dual}{\vnbstmt{\begin{array}{@{}l@{}} \forall v, i. \\ \quad \operatorname{is\text{-}hilbert\text{-}space}(v) \Rightarrow ((i \in L_{\bullet}({v})) \Leftrightarrow (i \in \operatorname{weak\text{-}dual}(v))) \end{array}}}
\vnbprop{weak-pairing-pf-apply}{\vnbstmt{\begin{array}{@{}l@{}} \forall v, f, t. \\ \quad (\operatorname{is\text{-}complex\text{-}inner\text{-}product\text{-}space}(v) \wedge ((f \in L_{\bullet}({v})) \wedge (t \in L({v}, {v})))) \Rightarrow  \\ \quad \quad \operatorname{weak\text{-}pairing\text{-}pf}(L({v}), L_{\bullet}({v}))(f)(t) = \operatorname{magnitude}(f(t)) \end{array}}}
\vnbprop{weak-pairing-seminorm-family}{\vnbstmt{\begin{array}{@{}l@{}} \forall v. \\ \quad \operatorname{is\text{-}complex\text{-}inner\text{-}product\text{-}space}(v) \Rightarrow  \\ \quad \quad \operatorname{is\text{-}seminorm\text{-}family\text{-}on}(\operatorname{cc\text{-}normed\text{-}field}, L({v}), L_{\bullet}({v}), \\ \quad \quad \quad \operatorname{weak\text{-}pairing\text{-}pf}(L({v}), L_{\bullet}({v}))) \end{array}}}
\vnbprop{weak-pairing-to-wot-continuous}{\vnbstmt{\begin{array}{@{}l@{}} \forall v. \\ \quad \operatorname{is\text{-}complex\text{-}inner\text{-}product\text{-}space}(v) \Rightarrow  \\ \quad \quad \operatorname{is\text{-}hom\text{-}top\text{-}space}(w\langle {L({v})}, {L_{\bullet}({v})}\rangle, L({v})_{w}, \operatorname{id\text{-}fun}(L({v}, {v}))) \end{array}}}
\vnbprop{wot-to-weak-pairing-continuous}{\vnbstmt{\begin{array}{@{}l@{}} \forall v. \\ \quad \operatorname{is\text{-}hilbert\text{-}space}(v) \Rightarrow  \\ \quad \quad \operatorname{is\text{-}hom\text{-}top\text{-}space}(L({v})_{w}, w\langle {L({v})}, {L_{\bullet}({v})}\rangle, \operatorname{id\text{-}fun}(L({v}, {v}))) \end{array}}}
\vnbprop{wot-is-strong-dual-pairing}{\vnbstmt{\begin{array}{@{}l@{}} \forall v. \\ \quad \operatorname{is\text{-}hilbert\text{-}space}(v) \Rightarrow  \\ \quad \quad \operatorname{is\text{-}hom\text{-}top\text{-}space}(L({v})_{w}, w\langle {L({v})}, {L_{\bullet}({v})}\rangle, \operatorname{id\text{-}fun}(L({v}, {v}))) \wedge \\ \quad \quad \operatorname{is\text{-}hom\text{-}top\text{-}space}(w\langle {L({v})}, {L_{\bullet}({v})}\rangle, L({v})_{w}, \operatorname{id\text{-}fun}(L({v}, {v}))) \end{array}}}

\vnbbookfile{theorem-library/cc-separation}{theorem-library/cc-separation}
\noindent{\small cc-separation.scm -- agent BA-1117-TRACE (round 3), 2026-11-17. The complex geometric Hahn-Banach theorem for a topology defined by a separating family of seminorms over CC: a point of an open set u is separated from a nonempty convex set missing u by the real part of a continuous CC-linear functional.  The real-part ...}\par
\vnbprop{realification-seminorm-top}{\vnbstmt{\begin{array}{@{}l@{}} \forall m, r, x, f. \\ \quad \operatorname{is\text{-}module}(m) \wedge \\ \quad \operatorname{scal}(m) = \operatorname{normed\text{-}field\text{-}as\text{-}commutative\text{-}ring}(\operatorname{cc\text{-}normed\text{-}field}) \wedge \\ \quad \operatorname{is\text{-}module}(r) \wedge \\ \quad \operatorname{scal}(r) = \operatorname{normed\text{-}field\text{-}as\text{-}commutative\text{-}ring}(\operatorname{rr\text{-}normed\text{-}field}) \wedge \\ \quad \operatorname{vec}(r) = \operatorname{vec}(m) \wedge \\ \quad \operatorname{vadd}(r) = \operatorname{vadd}(m) \wedge \\ \quad \operatorname{vzero}(r) = \operatorname{vzero}(m) \wedge \\ \quad \operatorname{vneg}(r) = \operatorname{vneg}(m) \wedge \\ \quad \forall a \in \mathbb{R}, b \in \operatorname{vec}(m).\; (\operatorname{act}(r)(a, b) = \operatorname{act}(m)(a, b)) \Rightarrow \\ \quad \quad (\operatorname{seminorm\text{-}top}(r, x, f) \simeq \operatorname{seminorm\text{-}top}(m, x, f)) \end{array}}}
\vnbprop{realification-seminorm-family}{\vnbstmt{\begin{array}{@{}l@{}} \forall m, r, x, f. \\ \quad \operatorname{is\text{-}module}(m) \wedge \\ \quad \operatorname{scal}(m) = \operatorname{normed\text{-}field\text{-}as\text{-}commutative\text{-}ring}(\operatorname{cc\text{-}normed\text{-}field}) \wedge \\ \quad \operatorname{is\text{-}module}(r) \wedge \\ \quad \operatorname{scal}(r) = \operatorname{normed\text{-}field\text{-}as\text{-}commutative\text{-}ring}(\operatorname{rr\text{-}normed\text{-}field}) \wedge \\ \quad \operatorname{vec}(r) = \operatorname{vec}(m) \wedge \\ \quad \operatorname{vadd}(r) = \operatorname{vadd}(m) \wedge \\ \quad \operatorname{vzero}(r) = \operatorname{vzero}(m) \wedge \\ \quad \operatorname{vneg}(r) = \operatorname{vneg}(m) \wedge \\ \quad \forall a \in \mathbb{R}, b \in \operatorname{vec}(m).\; (\operatorname{act}(r)(a, b) = \operatorname{act}(m)(a, b)) \wedge \\ \quad \operatorname{is\text{-}seminorm\text{-}family\text{-}on}(\operatorname{cc\text{-}normed\text{-}field}, m, x, f) \Rightarrow \\ \quad \quad \operatorname{is\text{-}seminorm\text{-}family\text{-}on}(\operatorname{rr\text{-}normed\text{-}field}, r, x, f) \end{array}}}
\vnbprop{realification-convex}{\vnbstmt{\begin{array}{@{}l@{}} \forall m, r, c. \\ \quad \operatorname{is\text{-}module}(m) \wedge \\ \quad \operatorname{scal}(m) = \operatorname{normed\text{-}field\text{-}as\text{-}commutative\text{-}ring}(\operatorname{cc\text{-}normed\text{-}field}) \wedge \\ \quad \operatorname{is\text{-}module}(r) \wedge \\ \quad \operatorname{scal}(r) = \operatorname{normed\text{-}field\text{-}as\text{-}commutative\text{-}ring}(\operatorname{rr\text{-}normed\text{-}field}) \wedge \\ \quad \operatorname{vec}(r) = \operatorname{vec}(m) \wedge \\ \quad \operatorname{vadd}(r) = \operatorname{vadd}(m) \wedge \\ \quad \operatorname{vzero}(r) = \operatorname{vzero}(m) \wedge \\ \quad \operatorname{vneg}(r) = \operatorname{vneg}(m) \wedge \\ \quad \forall a \in \mathbb{R}, x \in \operatorname{vec}(m).\; (\operatorname{act}(r)(a, x) = \operatorname{act}(m)(a, x)) \wedge \\ \quad \operatorname{is\text{-}convex\text{-}subset}(m, c) \Rightarrow \\ \quad \quad \operatorname{is\text{-}convex\text{-}subset}(r, c) \end{array}}}
\vnbprop{rr-nvs-act-apply}{\vnbstmt{\forall v, w \in \mathbb{R}.\; (\operatorname{act}(\operatorname{rr\text{-}nvs})(v, w) = (v \cdot w))}}
\vnbprop{rr-nvs-scal}{\vnbstmt{(\operatorname{scal}(\operatorname{rr\text{-}nvs}) \simeq \operatorname{normed\text{-}field\text{-}as\text{-}commutative\text{-}ring}(\operatorname{rr\text{-}normed\text{-}field}))}}
\vnbprop{linear-functional-linear-map}{\vnbstmt{\begin{array}{@{}l@{}} \forall r, f. \\ \quad \operatorname{is\text{-}module}(r) \wedge \\ \quad \operatorname{scal}(r) = \operatorname{normed\text{-}field\text{-}as\text{-}commutative\text{-}ring}(\operatorname{rr\text{-}normed\text{-}field}) \wedge \\ \quad \operatorname{is\text{-}linear\text{-}functional}(r, f) \Rightarrow \\ \quad \quad \operatorname{is\text{-}linear\text{-}map}(r, \operatorname{normed\text{-}vector\text{-}space\text{-}as\text{-}module}(\operatorname{rr\text{-}nvs}), f) \end{array}}}
\vnbprop{k-linear-functional-linear-map}{\vnbstmt{\begin{array}{@{}l@{}} \forall m, f. \\ \quad (\operatorname{is\text{-}module}(m) \wedge ((\operatorname{scal}(m) = \operatorname{normed\text{-}field\text{-}as\text{-}commutative\text{-}ring}(\operatorname{cc\text{-}normed\text{-}field})) \wedge \operatorname{is\text{-}k\text{-}linear}(m, f))) \Rightarrow  \\ \quad \quad \operatorname{is\text{-}linear\text{-}map}(m, \operatorname{cc\text{-}sn\text{-}module}, f) \end{array}}}
\vnbprop{complexified-functional-continuous}{\vnbstmt{\begin{array}{@{}l@{}} \forall m, x, f.; \forall n \in (\operatorname{vec}(m) \to \mathbb{R}), a \in (\operatorname{vec}(m) \to \mathbb{C}). \\ \quad \operatorname{is\text{-}seminorm\text{-}family\text{-}on}(\operatorname{cc\text{-}normed\text{-}field}, m, x, f) \wedge \\ \quad \forall \mathit{rax}, y \in \operatorname{vec}(m).\; (n(\operatorname{vadd}(m)(\mathit{rax}, y)) = (n(\mathit{rax}) + n(y))) \wedge \\ \quad \forall r \in \mathbb{R}, \mathit{rhx} \in \operatorname{vec}(m).\; (n(\operatorname{act}(m)(r, \mathit{rhx})) = (r \cdot n(\mathit{rhx}))) \wedge \\ \quad \forall \mathit{cvy} \in \operatorname{vec}(m).\; (a(\mathit{cvy}) \simeq (n(\mathit{cvy}) + (-n(\operatorname{act}(m)(+i, \mathit{cvy})) \cdot +i))) \wedge \\ \quad \exists \mathit{sbf} \in \mathbf{Set}. \\ \quad \quad |\mathit{sbf}| \in \mathbb{N} \wedge \\ \quad \quad \mathit{sbf} \subseteq x \wedge \\ \quad \quad \exists c \in \mathbb{R}. \\ \quad \quad \quad \operatorname{pos\text{-}rr}(c) \wedge \\ \quad \quad \quad \forall \mathit{sbx} \in \operatorname{vec}(m), t \in \mathbb{R}. \\ \quad \quad \quad \quad ((0 \leq t) \wedge \forall i \in \mathit{sbf}.\; (f(i)(\mathit{sbx}) \leq t)) \Rightarrow (\lvert n(\mathit{sbx}) \rvert \leq (c \cdot t)) \Rightarrow \\ \quad \quad \operatorname{is\text{-}hom\text{-}top\text{-}space}(\operatorname{seminorm\text{-}top}(m, x, f), \operatorname{cc\text{-}sn\text{-}top}, a) \end{array}}}
\vnbprop{cc-locally-convex-separation}{\vnbstmt{\begin{array}{@{}l@{}} \forall m, x, f, u, a, c. \\ \quad \operatorname{is\text{-}seminorm\text{-}family\text{-}on}(\operatorname{cc\text{-}normed\text{-}field}, m, x, f) \wedge \\ \quad u \in \operatorname{opens}(\operatorname{seminorm\text{-}top}(m, x, f)) \wedge \\ \quad a \in u \wedge \\ \quad \operatorname{is\text{-}convex\text{-}subset}(m, c) \wedge \\ \quad \exists c_{0}.\; (c_{0} \in c) \wedge \\ \quad \forall z \in c.\; \neg (z \in u) \Rightarrow \\ \quad \quad \exists \mathit{ccf}. \\ \quad \quad \quad \operatorname{is\text{-}k\text{-}linear}(m, \mathit{ccf}) \wedge \\ \quad \quad \quad \operatorname{is\text{-}hom\text{-}top\text{-}space}(\operatorname{seminorm\text{-}top}(m, x, f), \operatorname{cc\text{-}sn\text{-}top}, \mathit{ccf}) \wedge \\ \quad \quad \quad \forall \mathit{ccz} \in c.\; (\operatorname{real\text{-}part}(\operatorname{ccf}(\mathit{ccz})) < \operatorname{real\text{-}part}(\operatorname{ccf}(a))) \end{array}}}

\vnbbookfile{theorem-library/ba-sesq}{theorem-library/ba-sesq}
\noindent{\small ba-sesq.scm -- Thayer, Banach Algebras and Applications, ch. 5 (used for 5.10): a bounded sesquilinear form on a Hilbert space is $<$S x | y$>$ for a unique bounded operator S.  The form is CURRIED, bf in FUN(VEC v, FUN(VEC v, CC)), linear in its first and conjugate-linear in its second argument. Route: for each y, x |-$>$ ...}\par
\vnbprop{sesq-riesz-point}{\vnbstmt{\begin{array}{@{}l@{}} \forall v, f, c, y. \\ \quad \operatorname{is\text{-}hilbert\text{-}space}(v) \wedge \\ \quad f \in (\operatorname{vec}(v) \to (\operatorname{vec}(v) \to \mathbb{C})) \wedge \\ \quad \forall a, b, \mathit{bsy} \in \operatorname{vec}(v). \\ \quad \quad f(\operatorname{vadd}(v)(a, b))(\mathit{bsy}) = (f(a)(\mathit{bsy}) + f(b)(\mathit{bsy})) \wedge \\ \quad \forall k \in \mathbb{C}, a, \mathit{bsy} \in \operatorname{vec}(v). \\ \quad \quad f(\operatorname{act}(v)(k, a))(\mathit{bsy}) = (k \cdot f(a)(\mathit{bsy})) \wedge \\ \quad c \in \mathbb{R} \wedge \\ \quad 0 \leq c \wedge \\ \quad \forall a, \mathit{bsy} \in \operatorname{vec}(v). \\ \quad \quad \operatorname{magnitude}(f(a)(\mathit{bsy})) \leq (c \cdot (\lVert a \rVert \cdot \lVert \mathit{bsy} \rVert)) \wedge \\ \quad y \in \operatorname{vec}(v) \Rightarrow \\ \quad \quad \exists e \in \operatorname{vec}(v).\; \forall x \in \operatorname{vec}(v).\; (f(x)(y) = \langle\cdot,\cdot\rangle_{v}(x, e)) \end{array}}}
\vnbprop{sesq-riesz-map}{\vnbstmt{\begin{array}{@{}l@{}} \forall v, f, c. \\ \quad \operatorname{is\text{-}hilbert\text{-}space}(v) \wedge \\ \quad f \in (\operatorname{vec}(v) \to (\operatorname{vec}(v) \to \mathbb{C})) \wedge \\ \quad \forall a, b, y \in \operatorname{vec}(v). \\ \quad \quad f(\operatorname{vadd}(v)(a, b))(y) = (f(a)(y) + f(b)(y)) \wedge \\ \quad \forall k \in \mathbb{C}, a, y \in \operatorname{vec}(v). \\ \quad \quad f(\operatorname{act}(v)(k, a))(y) = (k \cdot f(a)(y)) \wedge \\ \quad c \in \mathbb{R} \wedge \\ \quad 0 \leq c \wedge \\ \quad \forall a, y \in \operatorname{vec}(v). \\ \quad \quad \operatorname{magnitude}(f(a)(y)) \leq (c \cdot (\lVert a \rVert \cdot \lVert y \rVert)) \Rightarrow \\ \quad \quad \exists r \in (\operatorname{vec}(v) \to \operatorname{vec}(v)). \forall y, x \in \operatorname{vec}(v). \\ \quad \quad \quad f(x)(y) = \langle\cdot,\cdot\rangle_{v}(x, r(y)) \end{array}}}
\vnbprop{sesq-riesz-bounded}{\vnbstmt{\begin{array}{@{}l@{}} \forall v, f, c, r. \\ \quad \operatorname{is\text{-}hilbert\text{-}space}(v) \wedge \\ \quad f \in (\operatorname{vec}(v) \to (\operatorname{vec}(v) \to \mathbb{C})) \wedge \\ \quad \forall a, b, y \in \operatorname{vec}(v). \\ \quad \quad f(\operatorname{vadd}(v)(a, b))(y) = (f(a)(y) + f(b)(y)) \wedge \\ \quad \forall k \in \mathbb{C}, a, y \in \operatorname{vec}(v). \\ \quad \quad f(\operatorname{act}(v)(k, a))(y) = (k \cdot f(a)(y)) \wedge \\ \quad c \in \mathbb{R} \wedge \\ \quad 0 \leq c \wedge \\ \quad \forall a, y \in \operatorname{vec}(v). \\ \quad \quad \operatorname{magnitude}(f(a)(y)) \leq (c \cdot (\lVert a \rVert \cdot \lVert y \rVert)) \wedge \\ \quad \forall a, b, y \in \operatorname{vec}(v). \\ \quad \quad f(a)(\operatorname{vadd}(v)(b, y)) = (f(a)(b) + f(a)(y)) \wedge \\ \quad \forall k \in \mathbb{C}, a, y \in \operatorname{vec}(v). \\ \quad \quad f(a)(\operatorname{act}(v)(k, y)) = (\operatorname{conjugate}(k) \cdot f(a)(y)) \wedge \\ \quad r \in (\operatorname{vec}(v) \to \operatorname{vec}(v)) \wedge \\ \quad \forall y, x \in \operatorname{vec}(v).\; (f(x)(y) = \langle\cdot,\cdot\rangle_{v}(x, r(y))) \Rightarrow \\ \quad \quad r \in L({v}, {v}) \end{array}}}
\vnbprop{bounded-sesquilinear-form-operator}{\vnbstmt{\begin{array}{@{}l@{}} \forall v, f, c. \\ \quad \operatorname{is\text{-}hilbert\text{-}space}(v) \wedge \\ \quad f \in (\operatorname{vec}(v) \to (\operatorname{vec}(v) \to \mathbb{C})) \wedge \\ \quad \forall a, b, y \in \operatorname{vec}(v). \\ \quad \quad f(\operatorname{vadd}(v)(a, b))(y) = (f(a)(y) + f(b)(y)) \wedge \\ \quad \forall k \in \mathbb{C}, a, y \in \operatorname{vec}(v). \\ \quad \quad f(\operatorname{act}(v)(k, a))(y) = (k \cdot f(a)(y)) \wedge \\ \quad c \in \mathbb{R} \wedge \\ \quad 0 \leq c \wedge \\ \quad \forall a, y \in \operatorname{vec}(v). \\ \quad \quad \operatorname{magnitude}(f(a)(y)) \leq (c \cdot (\lVert a \rVert \cdot \lVert y \rVert)) \wedge \\ \quad \forall a, b, y \in \operatorname{vec}(v). \\ \quad \quad f(a)(\operatorname{vadd}(v)(b, y)) = (f(a)(b) + f(a)(y)) \wedge \\ \quad \forall k \in \mathbb{C}, a, y \in \operatorname{vec}(v). \\ \quad \quad f(a)(\operatorname{act}(v)(k, y)) = (\operatorname{conjugate}(k) \cdot f(a)(y)) \Rightarrow \\ \quad \quad \exists s \in L({v}, {v}). \forall x, y \in \operatorname{vec}(v). \\ \quad \quad \quad f(x)(y) = \langle\cdot,\cdot\rangle_{v}(s(x), y) \end{array}}}
\vnbprop{bounded-sesquilinear-form-unique}{\vnbstmt{\begin{array}{@{}l@{}} \forall v, f, s_{1}, s_{2}. \\ \quad \operatorname{is\text{-}hilbert\text{-}space}(v) \wedge \\ \quad f \in (\operatorname{vec}(v) \to (\operatorname{vec}(v) \to \mathbb{C})) \wedge \\ \quad s_{1} \in L({v}, {v}) \wedge \\ \quad s_{2} \in L({v}, {v}) \wedge \\ \quad \forall x, y \in \operatorname{vec}(v).\; (f(x)(y) = \langle\cdot,\cdot\rangle_{v}(\operatorname{s1}(x), y)) \wedge \\ \quad \forall x, y \in \operatorname{vec}(v).\; (f(x)(y) = \langle\cdot,\cdot\rangle_{v}(\operatorname{s2}(x), y)) \Rightarrow \\ \quad \quad s_{1} = s_{2} \end{array}}}

\vnbbookfile{structure-library/wot-compact-vocab}{structure-library/wot-compact-vocab}
\noindent{\small wot-compact-vocab.scm -- agent BA-1117-TRACE (round 2), 2026-11-17. THE UNIT BALL OF L(H): the bounded operators of operator norm at most 1 (BanachAlgebras.pdf Proposition 5.10).  DEFINITIONS ONLY: no proof, no support, no axiom. Destination: structure-library/wot-compact-vocab.scm.  Library: operator-algebras.}\par
\par\medskip\noindent\textbf{Definition} (\texttt{op-unit-ball}).\;$\operatorname{op\text{-}unit\text{-}ball}(\operatorname{wuv\_}) :=$
\begin{equation*}
\{\operatorname{wub\_} \in L({\operatorname{wuv\_}}, {\operatorname{wuv\_}}) : (\lVert {\operatorname{wub\_}} \rVert \leq 1)\}
\end{equation*}


\vnbbookfile{theorem-library/wot-compact}{theorem-library/wot-compact}
\noindent{\small wot-compact.scm -- agent BA-1117-TRACE (round 2), 2026-11-17. BanachAlgebras.pdf Proposition 5.10: the unit ball of L(H) is compact in the weak operator topology. (1) sntop-hom-of-local, the initial-topology lemma (finding F-5.10 of BA-1111-C2-REPORT.md): a map g from a topological space z into a space with the ...}\par
\vnbprop{sntop-hom-of-local}{\vnbstmt{\begin{array}{@{}l@{}} \forall d, m, x, f, z, g. \\ \quad \operatorname{is\text{-}seminorm\text{-}family\text{-}on}(d, m, x, f) \wedge \\ \quad \operatorname{is\text{-}top\text{-}space}(z) \wedge \\ \quad g \in (\operatorname{pts}(z) \to \operatorname{vec}(m)) \wedge \\ \quad \forall q \in \operatorname{pts}(z), \mathit{slf} \in \mathbf{Set}, e. \\ \quad \quad ((|\mathit{slf}| \in \mathbb{N}) \wedge ((\mathit{slf} \subseteq x) \wedge \operatorname{pos\text{-}rr}(e))) \Rightarrow  \\ \quad \quad \quad \exists w \in \operatorname{opens}(z). \\ \quad \quad \quad \quad ((q \in w) \wedge \forall p \in w, i \in \mathit{slf}.\; (f(i)(\operatorname{vadd}(m)(g(p), \operatorname{vneg}(m)(g(q)))) \leq e)) \Rightarrow \\ \quad \quad \operatorname{is\text{-}hom\text{-}top\text{-}space}(z, \operatorname{seminorm\text{-}top}(m, x, f), g) \end{array}}}
\vnbprop{op-unit-ball-unfold}{\vnbstmt{\forall v.\; (\operatorname{op\text{-}unit\text{-}ball}(v) \simeq \{b \in L({v}, {v}) : (\lVert {b} \rVert \leq 1)\})}}
\vnbprop{op-unit-ball-form-bound}{\vnbstmt{\begin{array}{@{}l@{}} \forall v, t, x, y. \\ \quad (\operatorname{is\text{-}complex\text{-}inner\text{-}product\text{-}space}(v) \wedge ((t \in \operatorname{op\text{-}unit\text{-}ball}(v)) \wedge ((x \in \operatorname{vec}(v)) \wedge (y \in \operatorname{vec}(v))))) \Rightarrow  \\ \quad \quad \operatorname{magnitude}(\langle\cdot,\cdot\rangle_{v}(t(x), y)) \leq (\lVert x \rVert \cdot \lVert y \rVert) \end{array}}}
\vnbprop{op-unit-ball-of-form-bound}{\vnbstmt{\begin{array}{@{}l@{}} \forall v, t. \\ \quad \operatorname{is\text{-}complex\text{-}inner\text{-}product\text{-}space}(v) \wedge \\ \quad t \in L({v}, {v}) \wedge \\ \quad \forall x, y \in \operatorname{vec}(v). \\ \quad \quad \operatorname{magnitude}(\langle\cdot,\cdot\rangle_{v}(t(x), y)) \leq (\lVert x \rVert \cdot \lVert y \rVert) \Rightarrow \\ \quad \quad t \in \operatorname{op\text{-}unit\text{-}ball}(v) \end{array}}}
\vnbprop{bounded-forms-closed-in-product}{\vnbstmt{\begin{array}{@{}l@{}} \forall v. \\ \quad \operatorname{is\text{-}complex\text{-}inner\text{-}product\text{-}space}(v) \Rightarrow  \\ \quad \quad \operatorname{is\text{-}top\text{-}closed}(\operatorname{tproduct\text{-}top}((\operatorname{vec}(v) \times \operatorname{vec}(v)), (\lambda x \in (\operatorname{vec}(v) \times \operatorname{vec}(v)).\; \operatorname{metric\text{-}top}(\operatorname{cc\text{-}ms}))), \\ \quad \quad \quad \{h \in \operatorname{tprod\text{-}pts}((\operatorname{vec}(v) \times \operatorname{vec}(v)), (\lambda x \in (\operatorname{vec}(v) \times \operatorname{vec}(v)).\; \operatorname{metric\text{-}top}(\operatorname{cc\text{-}ms}))) : (\forall a, y \in \operatorname{vec}(v).\; (\operatorname{magnitude}(h({[}a, y])) \leq (\lVert a \rVert \cdot \lVert y \rVert)) \wedge (\forall a, b, y \in \operatorname{vec}(v).\; (h({[}\operatorname{vadd}(v)(a, b), y]) = (h({[}a, y]) + h({[}b, y]))) \wedge (\forall k \in \mathbb{C}, a, y \in \operatorname{vec}(v).\; (h({[}\operatorname{act}(v)(k, a), y]) = (k \cdot h({[}a, y]))) \wedge (\forall a, b, y \in \operatorname{vec}(v).\; (h({[}a, \operatorname{vadd}(v)(b, y)]) = (h({[}a, b]) + h({[}a, y]))) \wedge \forall k \in \mathbb{C}, a, y \in \operatorname{vec}(v).\; (h({[}a, \operatorname{act}(v)(k, y)]) = (\operatorname{conjugate}(k) \cdot h({[}a, y])))))))\}) \end{array}}}
\vnbprop{bounded-forms-compact}{\vnbstmt{\begin{array}{@{}l@{}} \forall v. \\ \quad \operatorname{is\text{-}complex\text{-}inner\text{-}product\text{-}space}(v) \Rightarrow  \\ \quad \quad \operatorname{is\text{-}compact\text{-}t}(\operatorname{subspace\text{-}top}(\operatorname{tproduct\text{-}top}((\operatorname{vec}(v) \times \operatorname{vec}(v)), (\lambda x \in (\operatorname{vec}(v) \times \operatorname{vec}(v)).\; \operatorname{metric\text{-}top}(\operatorname{cc\text{-}ms}))), \\ \quad \quad \quad \quad \{h \in \operatorname{tprod\text{-}pts}((\operatorname{vec}(v) \times \operatorname{vec}(v)), (\lambda x \in (\operatorname{vec}(v) \times \operatorname{vec}(v)).\; \operatorname{metric\text{-}top}(\operatorname{cc\text{-}ms}))) : (\forall a, y \in \operatorname{vec}(v).\; (\operatorname{magnitude}(h({[}a, y])) \leq (\lVert a \rVert \cdot \lVert y \rVert)) \wedge (\forall a, b, y \in \operatorname{vec}(v).\; (h({[}\operatorname{vadd}(v)(a, b), y]) = (h({[}a, y]) + h({[}b, y]))) \wedge (\forall k \in \mathbb{C}, a, y \in \operatorname{vec}(v).\; (h({[}\operatorname{act}(v)(k, a), y]) = (k \cdot h({[}a, y]))) \wedge (\forall a, b, y \in \operatorname{vec}(v).\; (h({[}a, \operatorname{vadd}(v)(b, y)]) = (h({[}a, b]) + h({[}a, y]))) \wedge \forall k \in \mathbb{C}, a, y \in \operatorname{vec}(v).\; (h({[}a, \operatorname{act}(v)(k, y)]) = (\operatorname{conjugate}(k) \cdot h({[}a, y])))))))\})) \end{array}}}
\vnbprop{bounded-form-operator-in-unit-ball}{\vnbstmt{\begin{array}{@{}l@{}} \forall v, h. \\ \quad \operatorname{is\text{-}hilbert\text{-}space}(v) \wedge \\ \quad h \\ \quad \in\; \{\mathit{wkh} \in \operatorname{tprod\text{-}pts}((\operatorname{vec}(v) \times \operatorname{vec}(v)), (\lambda x \in (\operatorname{vec}(v) \times \operatorname{vec}(v)).\; \operatorname{metric\text{-}top}(\operatorname{cc\text{-}ms}))) : (\forall a, y \in \operatorname{vec}(v).\; (\operatorname{magnitude}(\operatorname{wkh}({[}a, y])) \leq (\lVert a \rVert \cdot \lVert y \rVert)) \wedge (\forall a, b, y \in \operatorname{vec}(v).\; (\operatorname{wkh}({[}\operatorname{vadd}(v)(a, b), y]) = (\operatorname{wkh}({[}a, y]) + \operatorname{wkh}({[}b, y]))) \wedge (\forall k \in \mathbb{C}, a, y \in \operatorname{vec}(v).\; (\operatorname{wkh}({[}\operatorname{act}(v)(k, a), y]) = (k \cdot \operatorname{wkh}({[}a, y]))) \wedge (\forall a, b, y \in \operatorname{vec}(v).\; (\operatorname{wkh}({[}a, \operatorname{vadd}(v)(b, y)]) = (\operatorname{wkh}({[}a, b]) + \operatorname{wkh}({[}a, y]))) \wedge \forall k \in \mathbb{C}, a, y \in \operatorname{vec}(v).\; (\operatorname{wkh}({[}a, \operatorname{act}(v)(k, y)]) = (\operatorname{conjugate}(k) \cdot \operatorname{wkh}({[}a, y])))))))\} \Rightarrow \\ \quad \quad \exists s \in \operatorname{op\text{-}unit\text{-}ball}(v). \forall a, y \in \operatorname{vec}(v). \\ \quad \quad \quad h({[}a, y]) = \langle\cdot,\cdot\rangle_{v}(s(a), y) \end{array}}}
\vnbprop{operator-of-form-unique}{\vnbstmt{\begin{array}{@{}l@{}} \forall v, a, b. \\ \quad \operatorname{is\text{-}hilbert\text{-}space}(v) \wedge \\ \quad a \in L({v}, {v}) \wedge \\ \quad b \in L({v}, {v}) \wedge \\ \quad \forall \mathit{wka}, y \in \operatorname{vec}(v).\; (\langle\cdot,\cdot\rangle_{v}(a(\mathit{wka}), y) = \langle\cdot,\cdot\rangle_{v}(b(\mathit{wka}), y)) \Rightarrow \\ \quad \quad a = b \end{array}}}
\vnbprop{operator-form-in-bounded-forms}{\vnbstmt{\begin{array}{@{}l@{}} \forall v, t. \\ \quad (\operatorname{is\text{-}complex\text{-}inner\text{-}product\text{-}space}(v) \wedge (t \in \operatorname{op\text{-}unit\text{-}ball}(v))) \Rightarrow  \\ \quad \quad (\lambda h \in (\operatorname{vec}(v) \times \operatorname{vec}(v)).\; \langle\cdot,\cdot\rangle_{v}(t(\operatorname{nth}(1, h)), \operatorname{nth}(2, h))) \\ \quad \quad \in\; \{\mathit{wkh} \in \operatorname{tprod\text{-}pts}((\operatorname{vec}(v) \times \operatorname{vec}(v)), (\lambda x \in (\operatorname{vec}(v) \times \operatorname{vec}(v)).\; \operatorname{metric\text{-}top}(\operatorname{cc\text{-}ms}))) : (\forall a, y \in \operatorname{vec}(v).\; (\operatorname{magnitude}(\operatorname{wkh}({[}a, y])) \leq (\lVert a \rVert \cdot \lVert y \rVert)) \wedge (\forall a, b, y \in \operatorname{vec}(v).\; (\operatorname{wkh}({[}\operatorname{vadd}(v)(a, b), y]) = (\operatorname{wkh}({[}a, y]) + \operatorname{wkh}({[}b, y]))) \wedge (\forall k \in \mathbb{C}, a, y \in \operatorname{vec}(v).\; (\operatorname{wkh}({[}\operatorname{act}(v)(k, a), y]) = (k \cdot \operatorname{wkh}({[}a, y]))) \wedge (\forall a, b, y \in \operatorname{vec}(v).\; (\operatorname{wkh}({[}a, \operatorname{vadd}(v)(b, y)]) = (\operatorname{wkh}({[}a, b]) + \operatorname{wkh}({[}a, y]))) \wedge \forall k \in \mathbb{C}, a, y \in \operatorname{vec}(v).\; (\operatorname{wkh}({[}a, \operatorname{act}(v)(k, y)]) = (\operatorname{conjugate}(k) \cdot \operatorname{wkh}({[}a, y])))))))\} \end{array}}}
\vnbprop{bounded-forms-inverse-wot-continuous}{\vnbstmt{\begin{array}{@{}l@{}} \forall v, f. \\ \quad \operatorname{is\text{-}complex\text{-}inner\text{-}product\text{-}space}(v) \wedge \\ \quad f \\ \quad \in\; (\{h \in \operatorname{tprod\text{-}pts}((\operatorname{vec}(v) \times \operatorname{vec}(v)), (\lambda x \in (\operatorname{vec}(v) \times \operatorname{vec}(v)).\; \operatorname{metric\text{-}top}(\operatorname{cc\text{-}ms}))) : (\forall a, y \in \operatorname{vec}(v).\; (\operatorname{magnitude}(h({[}a, y])) \leq (\lVert a \rVert \cdot \lVert y \rVert)) \wedge (\forall a, b, y \in \operatorname{vec}(v).\; (h({[}\operatorname{vadd}(v)(a, b), y]) = (h({[}a, y]) + h({[}b, y]))) \wedge (\forall k \in \mathbb{C}, a, y \in \operatorname{vec}(v).\; (h({[}\operatorname{act}(v)(k, a), y]) = (k \cdot h({[}a, y]))) \wedge (\forall a, b, y \in \operatorname{vec}(v).\; (h({[}a, \operatorname{vadd}(v)(b, y)]) = (h({[}a, b]) + h({[}a, y]))) \wedge \forall k \in \mathbb{C}, a, y \in \operatorname{vec}(v).\; (h({[}a, \operatorname{act}(v)(k, y)]) = (\operatorname{conjugate}(k) \cdot h({[}a, y])))))))\} \to L({v}, {v})) \wedge \\ \quad \forall \mathit{wqh} \in \{h \in \operatorname{tprod\text{-}pts}((\operatorname{vec}(v) \times \operatorname{vec}(v)), (\lambda x \in (\operatorname{vec}(v) \times \operatorname{vec}(v)).\; \operatorname{metric\text{-}top}(\operatorname{cc\text{-}ms}))) : (\forall a, y \in \operatorname{vec}(v).\; (\operatorname{magnitude}(h({[}a, y])) \leq (\lVert a \rVert \cdot \lVert y \rVert)) \wedge (\forall a, b, y \in \operatorname{vec}(v).\; (h({[}\operatorname{vadd}(v)(a, b), y]) = (h({[}a, y]) + h({[}b, y]))) \wedge (\forall k \in \mathbb{C}, a, y \in \operatorname{vec}(v).\; (h({[}\operatorname{act}(v)(k, a), y]) = (k \cdot h({[}a, y]))) \wedge (\forall a, b, y \in \operatorname{vec}(v).\; (h({[}a, \operatorname{vadd}(v)(b, y)]) = (h({[}a, b]) + h({[}a, y]))) \wedge \forall k \in \mathbb{C}, a, y \in \operatorname{vec}(v).\; (h({[}a, \operatorname{act}(v)(k, y)]) = (\operatorname{conjugate}(k) \cdot h({[}a, y])))))))\}, a, y \in \operatorname{vec}(v). \\ \quad \quad \operatorname{wqh}({[}a, y]) = \langle\cdot,\cdot\rangle_{v}(f(\mathit{wqh})(a), y) \Rightarrow \\ \quad \quad \operatorname{is\text{-}hom\text{-}top\text{-}space}(\operatorname{subspace\text{-}top}(\operatorname{tproduct\text{-}top}((\operatorname{vec}(v) \times \operatorname{vec}(v)), (\lambda x \in (\operatorname{vec}(v) \times \operatorname{vec}(v)).\; \operatorname{metric\text{-}top}(\operatorname{cc\text{-}ms}))), \\ \quad \quad \quad \quad \{h \in \operatorname{tprod\text{-}pts}((\operatorname{vec}(v) \times \operatorname{vec}(v)), (\lambda x \in (\operatorname{vec}(v) \times \operatorname{vec}(v)).\; \operatorname{metric\text{-}top}(\operatorname{cc\text{-}ms}))) : (\forall a, y \in \operatorname{vec}(v).\; (\operatorname{magnitude}(h({[}a, y])) \leq (\lVert a \rVert \cdot \lVert y \rVert)) \wedge (\forall a, b, y \in \operatorname{vec}(v).\; (h({[}\operatorname{vadd}(v)(a, b), y]) = (h({[}a, y]) + h({[}b, y]))) \wedge (\forall k \in \mathbb{C}, a, y \in \operatorname{vec}(v).\; (h({[}\operatorname{act}(v)(k, a), y]) = (k \cdot h({[}a, y]))) \wedge (\forall a, b, y \in \operatorname{vec}(v).\; (h({[}a, \operatorname{vadd}(v)(b, y)]) = (h({[}a, b]) + h({[}a, y]))) \wedge \forall k \in \mathbb{C}, a, y \in \operatorname{vec}(v).\; (h({[}a, \operatorname{act}(v)(k, y)]) = (\operatorname{conjugate}(k) \cdot h({[}a, y])))))))\}), \\ \quad \quad \quad L({v})_{w}, f) \end{array}}}
\vnbprop{op-unit-ball-wot-compact}{\vnbstmt{\begin{array}{@{}l@{}} \forall v. \\ \quad \operatorname{is\text{-}hilbert\text{-}space}(v) \Rightarrow \operatorname{is\text{-}compact\text{-}t}(\operatorname{subspace\text{-}top}(L({v})_{w}, \operatorname{op\text{-}unit\text{-}ball}(v))) \end{array}}}

\part{Strand Spaces}
\noindent Messages as tokens, lists and encryptions; strands, strand spaces, nodes and the two node relations; bundles and causal precedence; origination and entry points; the adversary and the bound on what it can originate.
\par\noindent\textit{Builds on:} Foundations.
\par\noindent{\small 21 files.}

\vnbbookfile{structure-library/strand-messages}{structure-library/strand-messages}
\noindent{\small strand-messages.scm -- THE MESSAGES OF STRAND SPACES, ON THE USER'S MODEL, WITH ENCRYPTION (\textasciitilde{}/notes/notes-60.txt, item (a): "Let T be a set of tokens.  A message is a nested list of tokens, i.e. a token is a message and any list of messages is a message.  Subterms are defined recursively in the obvious way."  The user, ...}\par
\par\medskip\noindent\textbf{Definition} (\texttt{encr}).\;$\operatorname{encr}(\operatorname{smg\_}, \operatorname{smk\_}) :=$ ${[}2, \operatorname{smg\_}, \operatorname{smk\_}]$
\par\medskip\noindent\textbf{Definition} (\texttt{lst}).\;$\operatorname{lst}(\operatorname{sml\_}) :=$ ${[}1, \operatorname{sml\_}]$
\par\medskip\noindent\textbf{Definition} (\texttt{is-token-set}).\;For every $\operatorname{smt\_}$, 
\begin{equation*}
\operatorname{is\text{-}token\text{-}set}(\operatorname{smt\_}) \Leftrightarrow ((\operatorname{smt\_} \in \mathbf{Set}) \wedge \forall \operatorname{smx\_} \in \operatorname{smt\_}.\; \neg (\operatorname{smx\_} \in \operatorname{tuples}(\mathbf{Set})))
\end{equation*}

\par\medskip\noindent\textbf{Definition} (\texttt{msg-set}).\;$\operatorname{msg\text{-}set}(\operatorname{smt\_}, \operatorname{smk\_}) :=$ $\operatorname{big\text{-}union}(\operatorname{smn\_}, \mathbb{N}, \operatorname{msg\text{-}stage}(\operatorname{smt\_}, \operatorname{smk\_}, \operatorname{smn\_}))$
\par\medskip\noindent\textbf{Definition} (\texttt{msg-stage}).\;For all $\operatorname{smt\_}$ and $\operatorname{smk\_}$,  $(\operatorname{msg\text{-}stage}(\operatorname{smt\_}, \operatorname{smk\_}, 0) \simeq \operatorname{smt\_})$\par
Let $\operatorname{smt\_}$ and $\operatorname{smk\_}$ be arbitrary. Let $\operatorname{smn\_} \in \mathbb{N}$. Then
\begin{equation*}
(\operatorname{msg\text{-}stage}(\operatorname{smt\_}, \operatorname{smk\_}, \operatorname{succ}(\operatorname{smn\_})) \simeq (\operatorname{msg\text{-}stage}(\operatorname{smt\_}, \operatorname{smk\_}, \operatorname{smn\_}) \cup ((\{1\} \times \operatorname{tuples}(\operatorname{msg\text{-}stage}(\operatorname{smt\_}, \operatorname{smk\_}, \operatorname{smn\_}))) \cup (\{2\} \times \operatorname{msg\text{-}stage}(\operatorname{smt\_}, \operatorname{smk\_}, \operatorname{smn\_})))))
\end{equation*}

\par\medskip\noindent\textbf{Definition} (\texttt{subterm-rel}).\;$\operatorname{subterm\text{-}rel}(\operatorname{smt\_}, \operatorname{smk\_}) :=$
\begin{equation*}
\{\operatorname{smp\_} \in (\operatorname{msg\text{-}set}(\operatorname{smt\_}, \operatorname{smk\_}) \times \operatorname{msg\text{-}set}(\operatorname{smt\_}, \operatorname{smk\_})) : \forall \operatorname{smq\_} \in \operatorname{power}((\operatorname{msg\text{-}set}(\operatorname{smt\_}, \operatorname{smk\_}) \times \operatorname{msg\text{-}set}(\operatorname{smt\_}, \operatorname{smk\_}))).\; (\operatorname{is\text{-}subterm\text{-}closed}(\operatorname{smt\_}, \operatorname{smk\_}, \operatorname{smq\_}) \Rightarrow (\operatorname{smp\_} \in \operatorname{smq\_}))\}
\end{equation*}

\par\medskip\noindent\textbf{Definition} (\texttt{is-subterm}).\;For all $\operatorname{smt\_}$, $\operatorname{smk\_}$, $\operatorname{sma\_}$ and $\operatorname{smb\_}$, 
\begin{equation*}
\operatorname{is\text{-}subterm}(\operatorname{smt\_}, \operatorname{smk\_}, \operatorname{sma\_}, \operatorname{smb\_}) \Leftrightarrow ({[}\operatorname{sma\_}, \operatorname{smb\_}] \in \operatorname{subterm\text{-}rel}(\operatorname{smt\_}, \operatorname{smk\_}))
\end{equation*}

\par\medskip\noindent\textbf{Definition} (\texttt{msg-rank}).\;$\operatorname{msg\text{-}rank}(\operatorname{smt\_}, \operatorname{smk\_}, \operatorname{smm\_}) :=$
\begin{equation*}
(\iota\, \operatorname{smr\_}.\;((\operatorname{smr\_} \in \mathbb{N}) \wedge ((\operatorname{smm\_} \in \operatorname{msg\text{-}stage}(\operatorname{smt\_}, \operatorname{smk\_}, \operatorname{smr\_})) \wedge \forall \operatorname{smj\_} \in \mathbb{N}.\; ((\operatorname{smm\_} \in \operatorname{msg\text{-}stage}(\operatorname{smt\_}, \operatorname{smk\_}, \operatorname{smj\_})) \Rightarrow (\operatorname{smr\_} \leq \operatorname{smj\_})))))
\end{equation*}

\par\medskip\noindent\textbf{Definition} (\texttt{is-subterm-closed}).\;For all $\operatorname{smt\_}$, $\operatorname{smk\_}$ and $\operatorname{smr\_}$, 
\begin{equation*}
\begin{array}{@{}l@{}} \operatorname{is\text{-}subterm\text{-}closed}(\operatorname{smt\_}, \operatorname{smk\_}, \operatorname{smr\_}) \\ \Leftrightarrow\; \forall \operatorname{sma\_} \in \operatorname{msg\text{-}set}(\operatorname{smt\_}, \operatorname{smk\_}).\; ({[}\operatorname{sma\_}, \operatorname{sma\_}] \in \operatorname{smr\_}) \wedge \\ \forall \operatorname{sma\_}.; \forall \operatorname{sml\_} \in \operatorname{tuples}(\operatorname{msg\text{-}set}(\operatorname{smt\_}, \operatorname{smk\_})), \operatorname{smi\_} \in \mathbb{N}. \\ \quad (1 \leq \operatorname{smi\_}) \Rightarrow  \\ \quad (\operatorname{smi\_} \leq \operatorname{length}(\operatorname{sml\_})) \Rightarrow  \\ \quad ({[}\operatorname{sma\_}, \operatorname{nth}(\operatorname{smi\_}, \operatorname{sml\_})] \in \operatorname{smr\_}) \Rightarrow ({[}\operatorname{sma\_}, \operatorname{lst}(\operatorname{sml\_})] \in \operatorname{smr\_}) \wedge \\ \forall \operatorname{sma\_}.; \forall \operatorname{smg\_} \in \operatorname{msg\text{-}set}(\operatorname{smt\_}, \operatorname{smk\_}), \operatorname{smy\_} \in \operatorname{smk\_}. \\ \quad ({[}\operatorname{sma\_}, \operatorname{smg\_}] \in \operatorname{smr\_}) \Rightarrow ({[}\operatorname{sma\_}, \operatorname{encr}(\operatorname{smg\_}, \operatorname{smy\_})] \in \operatorname{smr\_}) \end{array}
\end{equation*}


\vnbbookfile{theorem-library/strand-messages-laws}{theorem-library/strand-messages-laws}
\noindent{\small strand-messages-laws.scm -- THE MESSAGES OF STRAND SPACES (the user's model, with ENCRYPTION): their laws. (\textasciitilde{}/notes/notes-60.txt item (a), the user's decision of 2026-10-06 16:40Z adding \{m\}\_k; the definitions are structure-library/strand-messages.scm.)}\par
\vnbprop{lst-unfold}{\vnbstmt{\forall l.\; (\operatorname{lst}(l) \simeq {[}1, l])}}
\vnbprop{encr-unfold}{\vnbstmt{\forall g, y.\; (\operatorname{encr}(g, y) \simeq {[}2, g, y])}}
\vnbprop{encr-ne-list}{\vnbstmt{\forall g, y, l.\; \neg (\operatorname{encr}(g, y) = \operatorname{lst}(l))}}
\vnbprop{lst-ne-encr}{\vnbstmt{\forall l, g, y.\; \neg (\operatorname{lst}(l) = \operatorname{encr}(g, y))}}
\vnbprop{lst-injective}{\vnbstmt{\forall l, a.\; ((\operatorname{lst}(l) = \operatorname{lst}(a)) \Rightarrow (l = a))}}
\vnbprop{encr-injective}{\vnbstmt{\forall g, y, u, z.\; ((\operatorname{encr}(g, y) = \operatorname{encr}(u, z)) \Rightarrow ((g = u) \wedge (y = z)))}}
\vnbprop{msg-stage-token}{\vnbstmt{\forall t, k.; \forall a \in t.\; (a \in \operatorname{msg\text{-}stage}(t, k, 0))}}
\vnbprop{msg-stage-keep}{\vnbstmt{\begin{array}{@{}l@{}} \forall t, k, n, m. \\ \quad ((n \in \mathbb{N}) \wedge (m \in \operatorname{msg\text{-}stage}(t, k, n))) \Rightarrow (m \in \operatorname{msg\text{-}stage}(t, k, (n + 1))) \end{array}}}
\vnbprop{msg-stage-list}{\vnbstmt{\begin{array}{@{}l@{}} \forall t, k, n, l. \\ \quad ((n \in \mathbb{N}) \wedge (l \in \operatorname{tuples}(\operatorname{msg\text{-}stage}(t, k, n)))) \Rightarrow  \\ \quad \quad \operatorname{lst}(l) \in \operatorname{msg\text{-}stage}(t, k, (n + 1)) \end{array}}}
\vnbprop{msg-stage-encr}{\vnbstmt{\begin{array}{@{}l@{}} \forall t, k, n, g, y. \\ \quad ((n \in \mathbb{N}) \wedge ((g \in \operatorname{msg\text{-}stage}(t, k, n)) \wedge (y \in k))) \Rightarrow  \\ \quad \quad \operatorname{encr}(g, y) \in \operatorname{msg\text{-}stage}(t, k, (n + 1)) \end{array}}}
\vnbprop{msg-stage-zero-elim}{\vnbstmt{\forall t, k.; \forall m \in \operatorname{msg\text{-}stage}(t, k, 0).\; (m \in t)}}
\vnbprop{msg-stage-succ-cases}{\vnbstmt{\begin{array}{@{}l@{}} \forall t, k, n, m. \\ \quad ((n \in \mathbb{N}) \wedge (m \in \operatorname{msg\text{-}stage}(t, k, (n + 1)))) \Rightarrow  \\ \quad \quad m \in \operatorname{msg\text{-}stage}(t, k, n) \vee \\ \quad \quad \exists l \in \operatorname{tuples}(\operatorname{msg\text{-}stage}(t, k, n)).\; (m = \operatorname{lst}(l)) \vee \\ \quad \quad \exists g \in \operatorname{msg\text{-}stage}(t, k, n), y \in k.\; (m = \operatorname{encr}(g, y)) \end{array}}}
\vnbprop{msg-stage-add}{\vnbstmt{\begin{array}{@{}l@{}} \forall d \in \mathbb{N}, t, k, j, m. \\ \quad ((j \in \mathbb{N}) \wedge (m \in \operatorname{msg\text{-}stage}(t, k, j))) \Rightarrow (m \in \operatorname{msg\text{-}stage}(t, k, (j + d))) \end{array}}}
\vnbprop{msg-set-unfold}{\vnbstmt{\forall t, k.\; (\operatorname{msg\text{-}set}(t, k) \simeq \operatorname{big\text{-}union}(n, \mathbb{N}, \operatorname{msg\text{-}stage}(t, k, n)))}}
\vnbprop{msg-set-intro}{\vnbstmt{\begin{array}{@{}l@{}} \forall t, k, n, m. \\ \quad ((n \in \mathbb{N}) \wedge (m \in \operatorname{msg\text{-}stage}(t, k, n))) \Rightarrow (m \in \operatorname{msg\text{-}set}(t, k)) \end{array}}}
\vnbprop{msg-set-elim}{\vnbstmt{\forall t, k.; \forall m \in \operatorname{msg\text{-}set}(t, k).\; \exists n \in \mathbb{N}.\; (m \in \operatorname{msg\text{-}stage}(t, k, n))}}
\vnbprop{token-in-msgs}{\vnbstmt{\forall t, k.; \forall a \in t.\; (a \in \operatorname{msg\text{-}set}(t, k))}}
\vnbprop{msg-stage-set}{\vnbstmt{\forall n \in \mathbb{N}, t, k.\; (((t \in \mathbf{Set}) \wedge (k \in \mathbf{Set})) \Rightarrow (\operatorname{msg\text{-}stage}(t, k, n) \in \mathbf{Set}))}}
\vnbprop{msg-set-is-set}{\vnbstmt{\forall t, k.\; (((t \in \mathbf{Set}) \wedge (k \in \mathbf{Set})) \Rightarrow (\operatorname{msg\text{-}set}(t, k) \in \mathbf{Set}))}}
\vnbprop{msg-lists-are-lists}{\vnbstmt{\forall t, k.; \forall l \in \operatorname{tuples}(\operatorname{msg\text{-}set}(t, k)).\; (l \in \operatorname{tuples}(\mathbf{Set}))}}
\vnbprop{list-in-some-stage}{\vnbstmt{\begin{array}{@{}l@{}} \forall t, k.; \forall l \in \operatorname{tuples}(\operatorname{msg\text{-}set}(t, k)). \exists n \in \mathbb{N}. \\ \quad l \in \operatorname{tuples}(\operatorname{msg\text{-}stage}(t, k, n)) \end{array}}}
\vnbprop{list-in-msgs}{\vnbstmt{\forall t, k.; \forall l \in \operatorname{tuples}(\operatorname{msg\text{-}set}(t, k)).\; (\operatorname{lst}(l) \in \operatorname{msg\text{-}set}(t, k))}}
\vnbprop{encr-in-msgs}{\vnbstmt{\begin{array}{@{}l@{}} \forall t, k, g, y. \\ \quad ((g \in \operatorname{msg\text{-}set}(t, k)) \wedge (y \in k)) \Rightarrow (\operatorname{encr}(g, y) \in \operatorname{msg\text{-}set}(t, k)) \end{array}}}
\vnbprop{lst-in-tuples-set}{\vnbstmt{\forall l \in \mathbf{Set}.\; (\operatorname{lst}(l) \in \operatorname{tuples}(\mathbf{Set}))}}
\vnbprop{encr-in-tuples-set}{\vnbstmt{\forall g, y.\; (((g \in \mathbf{Set}) \wedge (y \in \mathbf{Set})) \Rightarrow (\operatorname{encr}(g, y) \in \operatorname{tuples}(\mathbf{Set})))}}
\vnbprop{token-set-is-set}{\vnbstmt{\forall t.\; (\operatorname{is\text{-}token\text{-}set}(t) \Rightarrow (t \in \mathbf{Set}))}}
\vnbprop{token-not-list}{\vnbstmt{\forall t, k.; \forall a \in t.\; (\operatorname{is\text{-}token\text{-}set}(t) \Rightarrow \neg (a \in \operatorname{tuples}(\operatorname{msg\text{-}set}(t, k))))}}
\vnbprop{token-ne-list}{\vnbstmt{\begin{array}{@{}l@{}} \forall t, k, l. \\ \quad (\operatorname{is\text{-}token\text{-}set}(t) \wedge ((k \in t) \wedge (l \in \mathbf{Set}))) \Rightarrow \neg (k = \operatorname{lst}(l)) \end{array}}}
\vnbprop{encr-ne-token}{\vnbstmt{\begin{array}{@{}l@{}} \forall t, k, g, y. \\ \quad (\operatorname{is\text{-}token\text{-}set}(t) \wedge ((k \in t) \wedge ((g \in \mathbf{Set}) \wedge (y \in \mathbf{Set})))) \Rightarrow \neg (k = \operatorname{encr}(g, y)) \end{array}}}
\vnbprop{msg-induction}{\vnbstmt{\begin{array}{@{}l@{}} \forall t, k, c.; \forall m \in \operatorname{msg\text{-}set}(t, k). \\ \quad \forall a \in t.\; (a \in c) \wedge \\ \quad \forall l \in \operatorname{tuples}(\operatorname{msg\text{-}set}(t, k)). \\ \quad \quad \forall i \in \mathbb{N}.\; (((1 \leq i) \wedge (i \leq \operatorname{length}(l))) \Rightarrow (\operatorname{nth}(i, l) \in c)) \Rightarrow  \\ \quad \quad \quad \operatorname{lst}(l) \in c \wedge \\ \quad \forall g, y.\; (((g \in \operatorname{msg\text{-}set}(t, k)) \wedge ((y \in k) \wedge (g \in c))) \Rightarrow (\operatorname{encr}(g, y) \in c)) \Rightarrow \\ \quad \quad m \in c \end{array}}}
\vnbprop{msg-cases}{\vnbstmt{\begin{array}{@{}l@{}} \forall t, k.; \forall m \in \operatorname{msg\text{-}set}(t, k). \\ \quad m \in t \vee \\ \quad \exists l \in \operatorname{tuples}(\operatorname{msg\text{-}set}(t, k)).\; (m = \operatorname{lst}(l)) \vee \\ \quad \exists g \in \operatorname{msg\text{-}set}(t, k), y \in k.\; (m = \operatorname{encr}(g, y)) \end{array}}}
\vnbprop{subterm-rel-unfold}{\vnbstmt{\begin{array}{@{}l@{}} \forall t, k. \\ \quad (\operatorname{subterm\text{-}rel}(t, k) \simeq \{p \in (\operatorname{msg\text{-}set}(t, k) \times \operatorname{msg\text{-}set}(t, k)) : \forall q \in \operatorname{power}((\operatorname{msg\text{-}set}(t, k) \times \operatorname{msg\text{-}set}(t, k))).\; (\operatorname{is\text{-}subterm\text{-}closed}(t, k, q) \Rightarrow (p \in q))\}) \end{array}}}
\vnbprop{subterm-in-msgs}{\vnbstmt{\begin{array}{@{}l@{}} \forall t, k, a, b. \\ \quad \operatorname{is\text{-}subterm}(t, k, a, b) \Rightarrow ((a \in \operatorname{msg\text{-}set}(t, k)) \wedge (b \in \operatorname{msg\text{-}set}(t, k))) \end{array}}}
\vnbprop{subterm-least}{\vnbstmt{\begin{array}{@{}l@{}} \forall t, k.; \forall r \in \operatorname{power}((\operatorname{msg\text{-}set}(t, k) \times \operatorname{msg\text{-}set}(t, k))), a, b. \\ \quad (\operatorname{is\text{-}subterm\text{-}closed}(t, k, r) \wedge \operatorname{is\text{-}subterm}(t, k, a, b)) \Rightarrow ({[}a, b] \in r) \end{array}}}
\vnbprop{subterm-refl}{\vnbstmt{\forall t, k.; \forall a \in \operatorname{msg\text{-}set}(t, k).\; \operatorname{is\text{-}subterm}(t, k, a, a)}}
\vnbprop{subterm-component}{\vnbstmt{\begin{array}{@{}l@{}} \forall t, k, a, l, i. \\ \quad l \in \operatorname{tuples}(\operatorname{msg\text{-}set}(t, k)) \wedge \\ \quad i \in \mathbb{N} \wedge \\ \quad 1 \leq i \wedge \\ \quad i \leq \operatorname{length}(l) \wedge \\ \quad \operatorname{is\text{-}subterm}(t, k, a, \operatorname{nth}(i, l)) \Rightarrow \\ \quad \quad \operatorname{is\text{-}subterm}(t, k, a, \operatorname{lst}(l)) \end{array}}}
\vnbprop{subterm-encr}{\vnbstmt{\begin{array}{@{}l@{}} \forall t, k, a, g, y. \\ \quad ((g \in \operatorname{msg\text{-}set}(t, k)) \wedge ((y \in k) \wedge \operatorname{is\text{-}subterm}(t, k, a, g))) \Rightarrow  \\ \quad \quad \operatorname{is\text{-}subterm}(t, k, a, \operatorname{encr}(g, y)) \end{array}}}
\vnbprop{subterm-cases}{\vnbstmt{\begin{array}{@{}l@{}} \forall t, k, a, b. \\ \quad ((t \in \mathbf{Set}) \wedge ((k \in \mathbf{Set}) \wedge \operatorname{is\text{-}subterm}(t, k, a, b))) \Rightarrow  \\ \quad \quad a = b \vee \\ \quad \quad \exists l \in \operatorname{tuples}(\operatorname{msg\text{-}set}(t, k)). \\ \quad \quad \quad ((b = \operatorname{lst}(l)) \wedge \exists i \in \mathbb{N}.\; ((1 \leq i) \wedge ((i \leq \operatorname{length}(l)) \wedge \operatorname{is\text{-}subterm}(t, k, a, \operatorname{nth}(i, l))))) \vee \\ \quad \quad \exists g \in \operatorname{msg\text{-}set}(t, k), y \in k.\; ((b = \operatorname{encr}(g, y)) \wedge \operatorname{is\text{-}subterm}(t, k, a, g)) \end{array}}}
\vnbprop{subterm-trans}{\vnbstmt{\begin{array}{@{}l@{}} \forall t, k, a, b, c. \\ \quad ((t \in \mathbf{Set}) \wedge ((k \in \mathbf{Set}) \wedge (\operatorname{is\text{-}subterm}(t, k, a, b) \wedge \operatorname{is\text{-}subterm}(t, k, b, c)))) \Rightarrow  \\ \quad \quad \operatorname{is\text{-}subterm}(t, k, a, c) \end{array}}}
\vnbprop{subterm-of-token}{\vnbstmt{\begin{array}{@{}l@{}} \forall t, k, a.; \forall b \in t. \\ \quad (\operatorname{is\text{-}token\text{-}set}(t) \wedge ((k \in \mathbf{Set}) \wedge \operatorname{is\text{-}subterm}(t, k, a, b))) \Rightarrow (a = b) \end{array}}}
\vnbprop{subterm-of-list}{\vnbstmt{\begin{array}{@{}l@{}} \forall t, k, a.; \forall l \in \operatorname{tuples}(\operatorname{msg\text{-}set}(t, k)). \\ \quad ((t \in \mathbf{Set}) \wedge ((k \in \mathbf{Set}) \wedge \operatorname{is\text{-}subterm}(t, k, a, \operatorname{lst}(l)))) \Rightarrow  \\ \quad \quad ((a = \operatorname{lst}(l)) \vee \exists i \in \mathbb{N}.\; ((1 \leq i) \wedge ((i \leq \operatorname{length}(l)) \wedge \operatorname{is\text{-}subterm}(t, k, a, \operatorname{nth}(i, l))))) \end{array}}}
\vnbprop{subterm-of-encr}{\vnbstmt{\begin{array}{@{}l@{}} \forall t, k, a, g, y. \\ \quad ((t \in \mathbf{Set}) \wedge ((k \in \mathbf{Set}) \wedge ((g \in \operatorname{msg\text{-}set}(t, k)) \wedge ((y \in k) \wedge \operatorname{is\text{-}subterm}(t, k, a, \operatorname{encr}(g, y)))))) \Rightarrow  \\ \quad \quad (a = \operatorname{encr}(g, y)) \vee \operatorname{is\text{-}subterm}(t, k, a, g) \end{array}}}
\vnbprop{token-subterm-of-list}{\vnbstmt{\begin{array}{@{}l@{}} \forall t, k, a, l. \\ \quad \operatorname{is\text{-}token\text{-}set}(t) \wedge \\ \quad k \in \mathbf{Set} \wedge \\ \quad a \in t \wedge \\ \quad l \in \operatorname{tuples}(\operatorname{msg\text{-}set}(t, k)) \wedge \\ \quad \operatorname{is\text{-}subterm}(t, k, a, \operatorname{lst}(l)) \Rightarrow \\ \quad \quad \exists i \in \mathbb{N}.\; ((1 \leq i) \wedge ((i \leq \operatorname{length}(l)) \wedge \operatorname{is\text{-}subterm}(t, k, a, \operatorname{nth}(i, l)))) \end{array}}}
\vnbprop{token-subterm-of-encr}{\vnbstmt{\begin{array}{@{}l@{}} \forall t, k, a, g, y. \\ \quad \operatorname{is\text{-}token\text{-}set}(t) \wedge \\ \quad k \in \mathbf{Set} \wedge \\ \quad a \in t \wedge \\ \quad g \in \operatorname{msg\text{-}set}(t, k) \wedge \\ \quad y \in k \wedge \\ \quad \operatorname{is\text{-}subterm}(t, k, a, \operatorname{encr}(g, y)) \Rightarrow \\ \quad \quad \operatorname{is\text{-}subterm}(t, k, a, g) \end{array}}}
\vnbprop{subterm-encr-other-key}{\vnbstmt{\begin{array}{@{}l@{}} \forall t, k, h, a, y, b. \\ \quad t \in \mathbf{Set} \wedge \\ \quad k \in \mathbf{Set} \wedge \\ \quad h \in \operatorname{msg\text{-}set}(t, k) \wedge \\ \quad a \in \operatorname{msg\text{-}set}(t, k) \wedge \\ \quad y \in k \wedge \\ \quad b \in k \wedge \\ \quad \neg (y = b) \wedge \\ \quad \operatorname{is\text{-}subterm}(t, k, \operatorname{encr}(a, b), \operatorname{encr}(h, y)) \Rightarrow \\ \quad \quad \operatorname{is\text{-}subterm}(t, k, \operatorname{encr}(a, b), h) \end{array}}}
\vnbprop{msg-rank-prop}{\vnbstmt{\begin{array}{@{}l@{}} \forall t, k.; \forall m \in \operatorname{msg\text{-}set}(t, k). \\ \quad \operatorname{msg\text{-}rank}(t, k, m) \in \mathbb{N} \wedge \\ \quad m \in \operatorname{msg\text{-}stage}(t, k, \operatorname{msg\text{-}rank}(t, k, m)) \wedge \\ \quad \forall j \in \mathbb{N}.\; ((m \in \operatorname{msg\text{-}stage}(t, k, j)) \Rightarrow (\operatorname{msg\text{-}rank}(t, k, m) \leq j)) \end{array}}}
\vnbprop{msg-stage-list-inv}{\vnbstmt{\begin{array}{@{}l@{}} \forall n \in \mathbb{N}, t, k.; \forall l \in \mathbf{Set}. \\ \quad (\operatorname{is\text{-}token\text{-}set}(t) \wedge (\operatorname{lst}(l) \in \operatorname{msg\text{-}stage}(t, k, n))) \Rightarrow  \\ \quad \quad \exists j \in \mathbb{N}.\; (((j + 1) \leq n) \wedge (l \in \operatorname{tuples}(\operatorname{msg\text{-}stage}(t, k, j)))) \end{array}}}
\vnbprop{msg-stage-encr-inv}{\vnbstmt{\begin{array}{@{}l@{}} \forall n \in \mathbb{N}, t, k, g, y. \\ \quad (\operatorname{is\text{-}token\text{-}set}(t) \wedge ((g \in \mathbf{Set}) \wedge ((y \in \mathbf{Set}) \wedge (\operatorname{encr}(g, y) \in \operatorname{msg\text{-}stage}(t, k, n))))) \Rightarrow  \\ \quad \quad \exists j \in \mathbb{N}.\; (((j + 1) \leq n) \wedge (g \in \operatorname{msg\text{-}stage}(t, k, j))) \end{array}}}
\vnbprop{msg-rank-token}{\vnbstmt{\forall t, k.; \forall a \in t.\; (\operatorname{msg\text{-}rank}(t, k, a) = 0)}}
\vnbprop{msg-rank-component}{\vnbstmt{\begin{array}{@{}l@{}} \forall t, k, l, i. \\ \quad (\operatorname{is\text{-}token\text{-}set}(t) \wedge ((l \in \operatorname{tuples}(\operatorname{msg\text{-}set}(t, k))) \wedge ((i \in \mathbb{N}) \wedge ((1 \leq i) \wedge (i \leq \operatorname{length}(l)))))) \Rightarrow  \\ \quad \quad (\operatorname{msg\text{-}rank}(t, k, \operatorname{nth}(i, l)) + 1) \leq \operatorname{msg\text{-}rank}(t, k, \operatorname{lst}(l)) \end{array}}}
\vnbprop{msg-rank-encr}{\vnbstmt{\begin{array}{@{}l@{}} \forall t, k, g, y. \\ \quad (\operatorname{is\text{-}token\text{-}set}(t) \wedge ((g \in \operatorname{msg\text{-}set}(t, k)) \wedge (y \in k))) \Rightarrow  \\ \quad \quad (\operatorname{msg\text{-}rank}(t, k, g) + 1) \leq \operatorname{msg\text{-}rank}(t, k, \operatorname{encr}(g, y)) \end{array}}}
\vnbprop{subterm-rank}{\vnbstmt{\begin{array}{@{}l@{}} \forall t, k, a, b. \\ \quad (\operatorname{is\text{-}token\text{-}set}(t) \wedge ((k \in \mathbf{Set}) \wedge \operatorname{is\text{-}subterm}(t, k, a, b))) \Rightarrow  \\ \quad \quad (a = b) \vee ((\operatorname{msg\text{-}rank}(t, k, a) + 1) \leq \operatorname{msg\text{-}rank}(t, k, b)) \end{array}}}
\vnbprop{subterm-antisym}{\vnbstmt{\begin{array}{@{}l@{}} \forall t, k, a, b. \\ \quad (\operatorname{is\text{-}token\text{-}set}(t) \wedge ((k \in \mathbf{Set}) \wedge (\operatorname{is\text{-}subterm}(t, k, a, b) \wedge \operatorname{is\text{-}subterm}(t, k, b, a)))) \Rightarrow  \\ \quad \quad a = b \end{array}}}

\vnbbookfile{structure-library/strand-spaces}{structure-library/strand-spaces}
\noindent{\small strand-spaces.scm -- STRAND SPACES, NODES, BUNDLES ON THE USER'S MODEL (\textasciitilde{}/notes/notes-60.txt items (b)-(g); the user's decisions of 2026-10-06 recorded in scratchpad/triage/WEEK3.md, section S-2; the paper's sections 2.1-2.4 supply the THEOREMS, not the definitions: Thayer, Herzog, Guttman, ...}\par
\vnbstructdef{bundle}{A bundle $s$ is a 3-tuple whose components are listed below.\par
\begin{enumerate}[nosep]
\item $\operatorname{bspace}(s) = \operatorname{nth}(1, s)$, a substructure (\texttt{substructure}): $\operatorname{is\text{-}strand\text{-}space}(\operatorname{bspace}(s))$.
\item $\operatorname{bnodes}(s) = \operatorname{nth}(2, s)$, a carrier (\texttt{carriers}): $\operatorname{bnodes}(s) \in \mathbf{Set}$.
\item $\operatorname{bsend}(s) = \operatorname{nth}(3, s)$, a relation (\texttt{relation}): $\operatorname{bsend}(s) \in \operatorname{power}((\operatorname{bnodes}(s) \times \operatorname{bnodes}(s)))$.
\end{enumerate}
The laws (the \texttt{property} and \texttt{law} clauses):
\begin{itemize}[nosep]
\item $\forall \operatorname{ssx\_} \in \operatorname{bnodes}(s).\; (\operatorname{ssx\_} \in \operatorname{nodes}(\operatorname{bspace}(s)))$
\item $|\operatorname{bnodes}(s)| \in \mathbb{N}$
\item $\forall \operatorname{ssx\_}, \operatorname{ssy\_}.\; (({[}\operatorname{ssx\_}, \operatorname{ssy\_}] \in \operatorname{bsend}(s)) \Rightarrow \operatorname{sends\text{-}to}(\operatorname{bspace}(s), \operatorname{ssx\_}, \operatorname{ssy\_}))$
\item
\begin{equation*}
\begin{array}{@{}l@{}} \forall \operatorname{ssy\_} \in \operatorname{bnodes}(s). \\ \quad \operatorname{is\text{-}received}(\operatorname{bspace}(s), \operatorname{ssy\_}) \Rightarrow  \\ \quad \quad \exists \operatorname{ssx\_}. \\ \quad \quad \quad (({[}\operatorname{ssx\_}, \operatorname{ssy\_}] \in \operatorname{bsend}(s)) \wedge \forall \operatorname{ssz\_}.\; (({[}\operatorname{ssz\_}, \operatorname{ssy\_}] \in \operatorname{bsend}(s)) \Rightarrow (\operatorname{ssz\_} = \operatorname{ssx\_}))) \end{array}
\end{equation*}

\item
\begin{equation*}
\begin{array}{@{}l@{}} \forall \operatorname{ssy\_} \in \operatorname{bnodes}(s), \operatorname{ssx\_}. \\ \quad \operatorname{same\text{-}strand\text{-}before}(\operatorname{bspace}(s), \operatorname{ssx\_}, \operatorname{ssy\_}) \Rightarrow (\operatorname{ssx\_} \in \operatorname{bnodes}(s)) \end{array}
\end{equation*}

\item $\forall \operatorname{ssx\_}.\; \neg ({[}\operatorname{ssx\_}, \operatorname{ssx\_}] \in \operatorname{prec\text{-}strict}(\operatorname{bspace}(s), \operatorname{bnodes}(s), \operatorname{bsend}(s)))$
\end{itemize}
The characteristic law (the axiom \texttt{is-bundle}): $\operatorname{is\text{-}bundle}(s)$ holds if and only if $\operatorname{length}(s) = 3$, every component has the typing listed with it, every law holds, and $s \in \mathbf{Set}$ (the structure is a set).
}
\vnbstructdef{strand-space}{A strand space $s$ is a 5-tuple whose components are listed below.\par
\begin{enumerate}[nosep]
\item $\operatorname{tokens}(s) = \operatorname{nth}(1, s)$, a carrier (\texttt{carriers}): $\operatorname{tokens}(s) \in \mathbf{Set}$.
\item $\operatorname{ids}(s) = \operatorname{nth}(2, s)$, a carrier (\texttt{carriers}): $\operatorname{ids}(s) \in \mathbf{Set}$.
\item $\operatorname{keys}(s) = \operatorname{nth}(3, s)$, a carrier (\texttt{carriers}): $\operatorname{keys}(s) \in \mathbf{Set}$.
\item $\operatorname{strand\text{-}map}(s) = \operatorname{nth}(4, s)$, an operation (\texttt{op}): $\operatorname{strand\text{-}map}(s) \in (\operatorname{ids}(s) \to \operatorname{strands}(\operatorname{tokens}(s), \operatorname{keys}(s)))$.
\item $\operatorname{kinv}(s) = \operatorname{nth}(5, s)$, an operation (\texttt{op}): $\operatorname{kinv}(s) \in (\operatorname{keys}(s) \to \operatorname{keys}(s))$.
\end{enumerate}
The laws (the \texttt{property} and \texttt{law} clauses):
\begin{itemize}[nosep]
\item $\operatorname{is\text{-}token\text{-}set}(\operatorname{tokens}(s))$
\item $\forall \operatorname{ssx\_} \in \operatorname{keys}(s).\; (\operatorname{ssx\_} \in \operatorname{tokens}(s))$
\item $\forall \operatorname{ssx\_} \in \operatorname{keys}(s).\; (\operatorname{kinv}(s)(\operatorname{kinv}(s)(\operatorname{ssx\_})) = \operatorname{ssx\_})$
\end{itemize}
The characteristic law (the axiom \texttt{is-strand-space}): $\operatorname{is\text{-}strand\text{-}space}(s)$ holds if and only if $\operatorname{length}(s) = 5$, every component has the typing listed with it, every law holds, and $s \in \mathbf{Set}$ (the structure is a set).
}
\par\medskip\noindent\textbf{Definition} (\texttt{nodes}).\;$\operatorname{nodes}(\operatorname{ssp\_}) :=$
\begin{equation*}
\{\operatorname{ssn\_} \in (\mathbb{N} \times \operatorname{ids}(\operatorname{ssp\_})) : ((1 \leq \operatorname{nth}(1, \operatorname{ssn\_})) \wedge (\operatorname{nth}(1, \operatorname{ssn\_}) \leq \operatorname{length}(\operatorname{strand\text{-}map}(\operatorname{ssp\_})(\operatorname{nth}(2, \operatorname{ssn\_})))))\}
\end{equation*}

\par\medskip\noindent\textbf{Definition} (\texttt{directed}).\;$\operatorname{directed}(\operatorname{sst\_}, \operatorname{ssk\_}) :=$ $(\{\mathit{sent}, \mathit{received}\} \times \operatorname{msg\text{-}set}(\operatorname{sst\_}, \operatorname{ssk\_}))$
\par\medskip\noindent\textbf{Definition} (\texttt{sent}).\; $\mathit{sent} = 0$
\par\medskip\noindent\textbf{Definition} (\texttt{received}).\; $\mathit{received} = 1$
\par\medskip\noindent\textbf{Definition} (\texttt{dir-of}).\;$\operatorname{dir\text{-}of}(\operatorname{ssd\_}) :=$ $\operatorname{nth}(1, \operatorname{ssd\_})$
\par\medskip\noindent\textbf{Definition} (\texttt{content-of}).\;$\operatorname{content\text{-}of}(\operatorname{ssd\_}) :=$ $\operatorname{nth}(2, \operatorname{ssd\_})$
\par\medskip\noindent\textbf{Definition} (\texttt{strands}).\;$\operatorname{strands}(\operatorname{sst\_}, \operatorname{ssk\_}) :=$ $\operatorname{tuples}(\operatorname{directed}(\operatorname{sst\_}, \operatorname{ssk\_}))$
\par\medskip\noindent\textbf{Definition} (\texttt{index-of}).\;$\operatorname{index\text{-}of}(\operatorname{ssn\_}) :=$ $\operatorname{nth}(1, \operatorname{ssn\_})$
\par\medskip\noindent\textbf{Definition} (\texttt{id-of}).\;$\operatorname{id\text{-}of}(\operatorname{ssn\_}) :=$ $\operatorname{nth}(2, \operatorname{ssn\_})$
\par\medskip\noindent\textbf{Definition} (\texttt{node-content}).\;$\operatorname{node\text{-}content}(\operatorname{ssp\_}, \operatorname{ssn\_}) :=$ $\operatorname{nth}(2, \operatorname{nth}(\operatorname{nth}(1, \operatorname{ssn\_}), \operatorname{strand\text{-}map}(\operatorname{ssp\_})(\operatorname{nth}(2, \operatorname{ssn\_}))))$
\par\medskip\noindent\textbf{Definition} (\texttt{node-dir}).\;$\operatorname{node\text{-}dir}(\operatorname{ssp\_}, \operatorname{ssn\_}) :=$ $\operatorname{nth}(1, \operatorname{nth}(\operatorname{nth}(1, \operatorname{ssn\_}), \operatorname{strand\text{-}map}(\operatorname{ssp\_})(\operatorname{nth}(2, \operatorname{ssn\_}))))$
\par\medskip\noindent\textbf{Definition} (\texttt{is-sent}).\;For all $\operatorname{ssp\_}$ and $\operatorname{ssn\_}$,  $\operatorname{is\text{-}sent}(\operatorname{ssp\_}, \operatorname{ssn\_}) \Leftrightarrow (\operatorname{node\text{-}dir}(\operatorname{ssp\_}, \operatorname{ssn\_}) = \mathit{sent})$
\par\medskip\noindent\textbf{Definition} (\texttt{is-received}).\;For all $\operatorname{ssp\_}$ and $\operatorname{ssn\_}$,  $\operatorname{is\text{-}received}(\operatorname{ssp\_}, \operatorname{ssn\_}) \Leftrightarrow (\operatorname{node\text{-}dir}(\operatorname{ssp\_}, \operatorname{ssn\_}) = \mathit{received})$
\par\medskip\noindent\textbf{Definition} (\texttt{same-strand-before}).\;For all $\operatorname{ssp\_}$, $\operatorname{ssn\_}$ and $\operatorname{ssm\_}$, 
\begin{equation*}
\begin{array}{@{}l@{}} \operatorname{same\text{-}strand\text{-}before}(\operatorname{ssp\_}, \operatorname{ssn\_}, \operatorname{ssm\_}) \\ \Leftrightarrow\; \operatorname{ssn\_} \in \operatorname{nodes}(\operatorname{ssp\_}) \wedge \\ \operatorname{ssm\_} \in \operatorname{nodes}(\operatorname{ssp\_}) \wedge \\ \operatorname{nth}(2, \operatorname{ssn\_}) = \operatorname{nth}(2, \operatorname{ssm\_}) \wedge \\ \operatorname{nth}(1, \operatorname{ssn\_}) < \operatorname{nth}(1, \operatorname{ssm\_}) \end{array}
\end{equation*}

\par\medskip\noindent\textbf{Definition} (\texttt{same-strand-upto}).\;For all $\operatorname{ssp\_}$, $\operatorname{ssn\_}$ and $\operatorname{ssm\_}$, 
\begin{equation*}
\begin{array}{@{}l@{}} \operatorname{same\text{-}strand\text{-}upto}(\operatorname{ssp\_}, \operatorname{ssn\_}, \operatorname{ssm\_}) \\ \Leftrightarrow\; \operatorname{ssn\_} \in \operatorname{nodes}(\operatorname{ssp\_}) \wedge \\ \operatorname{ssm\_} \in \operatorname{nodes}(\operatorname{ssp\_}) \wedge \\ \operatorname{nth}(2, \operatorname{ssn\_}) = \operatorname{nth}(2, \operatorname{ssm\_}) \wedge \\ \operatorname{nth}(1, \operatorname{ssn\_}) \leq \operatorname{nth}(1, \operatorname{ssm\_}) \end{array}
\end{equation*}

\par\medskip\noindent\textbf{Definition} (\texttt{sends-to}).\;For all $\operatorname{ssp\_}$, $\operatorname{ssn\_}$ and $\operatorname{ssm\_}$, 
\begin{equation*}
\begin{array}{@{}l@{}} \operatorname{sends\text{-}to}(\operatorname{ssp\_}, \operatorname{ssn\_}, \operatorname{ssm\_}) \\ \Leftrightarrow\; \operatorname{ssn\_} \in \operatorname{nodes}(\operatorname{ssp\_}) \wedge \\ \operatorname{ssm\_} \in \operatorname{nodes}(\operatorname{ssp\_}) \wedge \\ \operatorname{is\text{-}sent}(\operatorname{ssp\_}, \operatorname{ssn\_}) \wedge \\ \operatorname{is\text{-}received}(\operatorname{ssp\_}, \operatorname{ssm\_}) \wedge \\ \operatorname{node\text{-}content}(\operatorname{ssp\_}, \operatorname{ssn\_}) = \operatorname{node\text{-}content}(\operatorname{ssp\_}, \operatorname{ssm\_}) \end{array}
\end{equation*}

\par\medskip\noindent\textbf{Definition} (\texttt{occurs-in-space}).\;For all $\operatorname{ssp\_}$ and $\operatorname{ssu\_}$, 
\begin{equation*}
\begin{array}{@{}l@{}} \operatorname{occurs\text{-}in\text{-}space}(\operatorname{ssp\_}, \operatorname{ssu\_}) \\ \Leftrightarrow\; \exists \operatorname{ssn\_} \in \operatorname{nodes}(\operatorname{ssp\_}). \\ \quad \operatorname{is\text{-}subterm}(\operatorname{tokens}(\operatorname{ssp\_}), \operatorname{keys}(\operatorname{ssp\_}), \operatorname{ssu\_}, \operatorname{node\text{-}content}(\operatorname{ssp\_}, \operatorname{ssn\_})) \end{array}
\end{equation*}

\par\medskip\noindent\textbf{Definition} (\texttt{is-entry-point}).\;For all $\operatorname{ssp\_}$, $\operatorname{ssn\_}$ and $\operatorname{ssi\_}$, 
\begin{equation*}
\begin{array}{@{}l@{}} \operatorname{is\text{-}entry\text{-}point}(\operatorname{ssp\_}, \operatorname{ssn\_}, \operatorname{ssi\_}) \\ \Leftrightarrow\; \operatorname{ssn\_} \in \operatorname{nodes}(\operatorname{ssp\_}) \wedge \\ \operatorname{is\text{-}sent}(\operatorname{ssp\_}, \operatorname{ssn\_}) \wedge \\ \operatorname{node\text{-}content}(\operatorname{ssp\_}, \operatorname{ssn\_}) \in \operatorname{ssi\_} \wedge \\ \forall \operatorname{ssm\_}. \\ \quad \operatorname{same\text{-}strand\text{-}before}(\operatorname{ssp\_}, \operatorname{ssm\_}, \operatorname{ssn\_}) \Rightarrow  \\ \quad \quad \neg (\operatorname{node\text{-}content}(\operatorname{ssp\_}, \operatorname{ssm\_}) \in \operatorname{ssi\_}) \end{array}
\end{equation*}

\par\medskip\noindent\textbf{Definition} (\texttt{originates}).\;For all $\operatorname{ssp\_}$, $\operatorname{ssu\_}$ and $\operatorname{ssn\_}$, 
\begin{equation*}
\begin{array}{@{}l@{}} \operatorname{originates}(\operatorname{ssp\_}, \operatorname{ssu\_}, \operatorname{ssn\_}) \\ \Leftrightarrow\; \operatorname{is\text{-}entry\text{-}point}(\operatorname{ssp\_}, \operatorname{ssn\_}, \\ \quad \{\operatorname{ssw\_} \in \operatorname{msg\text{-}set}(\operatorname{tokens}(\operatorname{ssp\_}), \operatorname{keys}(\operatorname{ssp\_})) : \operatorname{is\text{-}subterm}(\operatorname{tokens}(\operatorname{ssp\_}), \operatorname{keys}(\operatorname{ssp\_}), \operatorname{ssu\_}, \operatorname{ssw\_})\}) \end{array}
\end{equation*}

\par\medskip\noindent\textbf{Definition} (\texttt{uniquely-originates}).\;For all $\operatorname{ssp\_}$ and $\operatorname{ssu\_}$, 
\begin{equation*}
\begin{array}{@{}l@{}} \operatorname{uniquely\text{-}originates}(\operatorname{ssp\_}, \operatorname{ssu\_}) \\ \Leftrightarrow\; \exists \operatorname{ssn\_}. \\ \quad (\operatorname{originates}(\operatorname{ssp\_}, \operatorname{ssu\_}, \operatorname{ssn\_}) \wedge \forall \operatorname{ssm\_}.\; (\operatorname{originates}(\operatorname{ssp\_}, \operatorname{ssu\_}, \operatorname{ssm\_}) \Rightarrow (\operatorname{ssm\_} = \operatorname{ssn\_}))) \end{array}
\end{equation*}

\par\medskip\noindent\textbf{Definition} (\texttt{is-prec-closed}).\;For all $\operatorname{spc\_}$, $\operatorname{snd\_}$, $\operatorname{ssd\_}$ and $\operatorname{srl\_}$, 
\begin{equation*}
\begin{array}{@{}l@{}} \operatorname{is\text{-}prec\text{-}closed}(\operatorname{spc\_}, \operatorname{snd\_}, \operatorname{ssd\_}, \operatorname{srl\_}) \\ \Leftrightarrow\; \forall \operatorname{ssx\_}, \operatorname{ssy\_}.\; (({[}\operatorname{ssx\_}, \operatorname{ssy\_}] \in \operatorname{ssd\_}) \Rightarrow ({[}\operatorname{ssx\_}, \operatorname{ssy\_}] \in \operatorname{srl\_})) \wedge \\ \forall \operatorname{ssx\_}, \operatorname{ssy\_} \in \operatorname{snd\_}. \\ \quad \operatorname{same\text{-}strand\text{-}before}(\operatorname{spc\_}, \operatorname{ssx\_}, \operatorname{ssy\_}) \Rightarrow ({[}\operatorname{ssx\_}, \operatorname{ssy\_}] \in \operatorname{srl\_}) \wedge \\ \forall \operatorname{ssx\_}, \operatorname{ssy\_}, \operatorname{ssz\_}. \\ \quad ({[}\operatorname{ssx\_}, \operatorname{ssy\_}] \in \operatorname{srl\_}) \Rightarrow (({[}\operatorname{ssy\_}, \operatorname{ssz\_}] \in \operatorname{srl\_}) \Rightarrow ({[}\operatorname{ssx\_}, \operatorname{ssz\_}] \in \operatorname{srl\_})) \end{array}
\end{equation*}

\par\medskip\noindent\textbf{Definition} (\texttt{prec-strict}).\;$\operatorname{prec\text{-}strict}(\operatorname{spc\_}, \operatorname{snd\_}, \operatorname{ssd\_}) :=$
\begin{equation*}
\{\operatorname{ssp\_} \in (\operatorname{snd\_} \times \operatorname{snd\_}) : \forall \operatorname{ssq\_} \in \operatorname{power}((\operatorname{snd\_} \times \operatorname{snd\_})).\; (\operatorname{is\text{-}prec\text{-}closed}(\operatorname{spc\_}, \operatorname{snd\_}, \operatorname{ssd\_}, \operatorname{ssq\_}) \Rightarrow (\operatorname{ssp\_} \in \operatorname{ssq\_}))\}
\end{equation*}

\par\medskip\noindent\textbf{Definition} (\texttt{precedes-strict}).\;$\operatorname{precedes\text{-}strict}(\operatorname{ssc\_}) :=$ $\operatorname{prec\text{-}strict}(\operatorname{bspace}(\operatorname{ssc\_}), \operatorname{bnodes}(\operatorname{ssc\_}), \operatorname{bsend}(\operatorname{ssc\_}))$
\par\medskip\noindent\textbf{Definition} (\texttt{precedes}).\;$\operatorname{precedes}(\operatorname{ssc\_}) :=$
\begin{equation*}
(\operatorname{precedes\text{-}strict}(\operatorname{ssc\_}) \cup \{\operatorname{ssp\_} \in (\operatorname{bnodes}(\operatorname{ssc\_}) \times \operatorname{bnodes}(\operatorname{ssc\_})) : (\operatorname{nth}(1, \operatorname{ssp\_}) = \operatorname{nth}(2, \operatorname{ssp\_}))\})
\end{equation*}

\par\medskip\noindent\textbf{Definition} (\texttt{hom-bundle}).\;$\operatorname{hom\text{-}bundle}(a, b) :=$
\begin{equation*}
\{\operatorname{homf\_} \in (\operatorname{bnodes}(a) \to \operatorname{bnodes}(b)) : \operatorname{is\text{-}hom\text{-}bundle}(a, b, \operatorname{homf\_})\}
\end{equation*}

\vnbdef{bundle-class}{\vnbstmt{\forall s.\; ((s \in \mathit{bundle}) \Leftrightarrow \operatorname{is\text{-}bundle}(s))}}
\vnbdef{is-hom-bundle-def}{\vnbstmt{\begin{array}{@{}l@{}} \forall a, b, f. \\ \quad \operatorname{is\text{-}hom\text{-}bundle}(a, b, f) \\ \quad \Leftrightarrow\; \operatorname{is\text{-}bundle}(a) \wedge \\ \quad \operatorname{is\text{-}bundle}(b) \wedge \\ \quad f \in (\operatorname{bnodes}(a) \to \operatorname{bnodes}(b)) \wedge \\ \quad \operatorname{bspace}(a) = \operatorname{bspace}(b) \wedge \\ \quad \forall x_{1}, x_{2} \in \operatorname{bnodes}(a). \\ \quad \quad ({[}x_{1}, x_{2}] \in \operatorname{bsend}(a)) \Rightarrow ({[}f(x_{1}), f(x_{2})] \in \operatorname{bsend}(b)) \end{array}}}
\vnbdef{is-hom-strand-space-def}{\vnbstmt{\begin{array}{@{}l@{}} \forall a, b, c, d, f. \\ \quad \operatorname{is\text{-}hom\text{-}strand\text{-}space}(a, b, c, d, f) \\ \quad \Leftrightarrow\; \operatorname{is\text{-}strand\text{-}space}(a) \wedge \\ \quad \operatorname{is\text{-}strand\text{-}space}(b) \wedge \\ \quad c \in (\operatorname{tokens}(a) \to \operatorname{tokens}(b)) \wedge \\ \quad d \in (\operatorname{ids}(a) \to \operatorname{ids}(b)) \wedge \\ \quad f \in (\operatorname{keys}(a) \to \operatorname{keys}(b)) \wedge \\ \quad \forall x_{1} \in \operatorname{ids}(a).\; (\operatorname{strand\text{-}map}(a)(x_{1}) = \operatorname{strand\text{-}map}(b)(d(x_{1}))) \wedge \\ \quad \forall x_{1} \in \operatorname{keys}(a).\; (f(\operatorname{kinv}(a)(x_{1})) = \operatorname{kinv}(b)(f(x_{1}))) \end{array}}}
\vnbdef{strand-space-class}{\vnbstmt{\forall s.\; ((s \in \operatorname{strand\text{-}space}) \Leftrightarrow \operatorname{is\text{-}strand\text{-}space}(s))}}
\vnbprop{is-strand-space-is-set}{\vnbstmt{\forall s.\; (\operatorname{is\text{-}strand\text{-}space}(s) \Rightarrow (s \in \mathbf{Set}))}}
\vnbprop{is-bundle-is-set}{\vnbstmt{\forall s.\; (\operatorname{is\text{-}bundle}(s) \Rightarrow (s \in \mathbf{Set}))}}

\vnbbookfile{theorem-library/strand-space-laws}{theorem-library/strand-space-laws}
\noindent{\small strand-space-laws.scm -- STRAND SPACES, NODES, BUNDLES, CAUSAL PRECEDENCE: the laws, and the paper's Lemmas 2.7, 2.8, 2.9 on the user's model. (\textasciitilde{}/notes/notes-60.txt; the definitions are structure-library/strand-spaces.scm; the paper is Thayer, Herzog, Guttman, \textasciitilde{}/docs/jcs\_strand\_spaces.pdf, sections 2.2-2.4.)}\par
\vnbprop{nodes-unfold}{\vnbstmt{\begin{array}{@{}l@{}} \forall a. \\ \quad (\operatorname{nodes}(a) \simeq \{n \in (\mathbb{N} \times \operatorname{ids}(a)) : ((1 \leq \operatorname{nth}(1, n)) \wedge (\operatorname{nth}(1, n) \leq \operatorname{length}(\operatorname{strand\text{-}map}(a)(\operatorname{nth}(2, n)))))\}) \end{array}}}
\vnbprop{node-content-unfold}{\vnbstmt{\begin{array}{@{}l@{}} \forall a, n. \\ \quad (\operatorname{node\text{-}content}(a, n) \simeq \operatorname{nth}(2, \operatorname{nth}(\operatorname{nth}(1, n), \operatorname{strand\text{-}map}(a)(\operatorname{nth}(2, n))))) \end{array}}}
\vnbprop{node-dir-unfold}{\vnbstmt{\begin{array}{@{}l@{}} \forall a, n. \\ \quad (\operatorname{node\text{-}dir}(a, n) \simeq \operatorname{nth}(1, \operatorname{nth}(\operatorname{nth}(1, n), \operatorname{strand\text{-}map}(a)(\operatorname{nth}(2, n))))) \end{array}}}
\vnbprop{directed-unfold}{\vnbstmt{\forall t, k.\; (\operatorname{directed}(t, k) \simeq (\{\mathit{sent}, \mathit{received}\} \times \operatorname{msg\text{-}set}(t, k)))}}
\vnbprop{strands-unfold}{\vnbstmt{\forall t, k.\; (\operatorname{strands}(t, k) \simeq \operatorname{tuples}(\operatorname{directed}(t, k)))}}
\vnbprop{index-id-unfold}{\vnbstmt{\forall n.\; ((\operatorname{index\text{-}of}(n) \simeq \operatorname{nth}(1, n)) \wedge (\operatorname{id\text{-}of}(n) \simeq \operatorname{nth}(2, n)))}}
\vnbprop{prec-strict-unfold}{\vnbstmt{\begin{array}{@{}l@{}} \forall a, d, b. \\ \quad (\operatorname{prec\text{-}strict}(a, d, b) \simeq \{p \in (d \times d) : \forall q \in \operatorname{power}((d \times d)).\; (\operatorname{is\text{-}prec\text{-}closed}(a, d, b, q) \Rightarrow (p \in q))\}) \end{array}}}
\vnbprop{precedes-strict-unfold}{\vnbstmt{\forall c.\; (\operatorname{precedes\text{-}strict}(c) \simeq \operatorname{prec\text{-}strict}(\operatorname{bspace}(c), \operatorname{bnodes}(c), \operatorname{bsend}(c)))}}
\vnbprop{precedes-unfold}{\vnbstmt{\begin{array}{@{}l@{}} \forall c. \\ \quad (\operatorname{precedes}(c) \simeq (\operatorname{precedes\text{-}strict}(c) \cup \{p \in (\operatorname{bnodes}(c) \times \operatorname{bnodes}(c)) : (\operatorname{nth}(1, p) = \operatorname{nth}(2, p))\})) \end{array}}}
\vnbprop{strand-space-parts}{\vnbstmt{\begin{array}{@{}l@{}} \forall a. \\ \quad \operatorname{is\text{-}strand\text{-}space}(a) \Rightarrow  \\ \quad \quad \operatorname{is\text{-}token\text{-}set}(\operatorname{tokens}(a)) \wedge \\ \quad \quad \operatorname{tokens}(a) \in \mathbf{Set} \wedge \\ \quad \quad \operatorname{ids}(a) \in \mathbf{Set} \wedge \\ \quad \quad \operatorname{keys}(a) \in \mathbf{Set} \wedge \\ \quad \quad \operatorname{strand\text{-}map}(a) \in (\operatorname{ids}(a) \to \operatorname{strands}(\operatorname{tokens}(a), \operatorname{keys}(a))) \end{array}}}
\vnbprop{keys-in-tokens}{\vnbstmt{\forall a.; \forall x \in \operatorname{keys}(a).\; (\operatorname{is\text{-}strand\text{-}space}(a) \Rightarrow (x \in \operatorname{tokens}(a)))}}
\vnbprop{kinv-in-fun}{\vnbstmt{\forall a.\; (\operatorname{is\text{-}strand\text{-}space}(a) \Rightarrow (\operatorname{kinv}(a) \in (\operatorname{keys}(a) \to \operatorname{keys}(a))))}}
\vnbprop{kinv-involution}{\vnbstmt{\begin{array}{@{}l@{}} \forall a.; \forall y \in \operatorname{keys}(a). \\ \quad \operatorname{is\text{-}strand\text{-}space}(a) \Rightarrow (\operatorname{kinv}(a)(\operatorname{kinv}(a)(y)) = y) \end{array}}}
\vnbprop{kinv-injective}{\vnbstmt{\begin{array}{@{}l@{}} \forall a, y, z. \\ \quad (\operatorname{is\text{-}strand\text{-}space}(a) \wedge ((y \in \operatorname{keys}(a)) \wedge ((z \in \operatorname{keys}(a)) \wedge (\operatorname{kinv}(a)(y) = \operatorname{kinv}(a)(z))))) \Rightarrow  \\ \quad \quad y = z \end{array}}}
\vnbprop{strand-at-id}{\vnbstmt{\begin{array}{@{}l@{}} \forall b, a. \\ \quad (\operatorname{is\text{-}strand\text{-}space}(b) \wedge (a \in \operatorname{ids}(b))) \Rightarrow  \\ \quad \quad \operatorname{strand\text{-}map}(b)(a) \in \operatorname{tuples}(\operatorname{directed}(\operatorname{tokens}(b), \operatorname{keys}(b))) \end{array}}}
\vnbprop{node-parts}{\vnbstmt{\begin{array}{@{}l@{}} \forall a.; \forall n \in \operatorname{nodes}(a). \\ \quad \operatorname{nth}(1, n) \in \mathbb{N} \wedge \\ \quad \operatorname{nth}(2, n) \in \operatorname{ids}(a) \wedge \\ \quad 1 \leq \operatorname{nth}(1, n) \wedge \\ \quad \operatorname{nth}(1, n) \leq \operatorname{length}(\operatorname{strand\text{-}map}(a)(\operatorname{nth}(2, n))) \wedge \\ \quad n = {[}\operatorname{nth}(1, n), \operatorname{nth}(2, n)] \end{array}}}
\vnbprop{node-intro}{\vnbstmt{\begin{array}{@{}l@{}} \forall b, k, a. \\ \quad ((k \in \mathbb{N}) \wedge ((a \in \operatorname{ids}(b)) \wedge ((1 \leq k) \wedge (k \leq \operatorname{length}(\operatorname{strand\text{-}map}(b)(a)))))) \Rightarrow  \\ \quad \quad {[}k, a] \in \operatorname{nodes}(b) \end{array}}}
\vnbprop{node-ext}{\vnbstmt{\begin{array}{@{}l@{}} \forall a, x, y. \\ \quad ((x \in \operatorname{nodes}(a)) \wedge ((y \in \operatorname{nodes}(a)) \wedge ((\operatorname{nth}(1, x) = \operatorname{nth}(1, y)) \wedge (\operatorname{nth}(2, x) = \operatorname{nth}(2, y))))) \Rightarrow  \\ \quad \quad x = y \end{array}}}
\vnbprop{directed-parts}{\vnbstmt{\begin{array}{@{}l@{}} \forall t, k.; \forall d \in \operatorname{directed}(t, k). \\ \quad (\operatorname{nth}(1, d) = \mathit{sent}) \vee (\operatorname{nth}(1, d) = \mathit{received}) \wedge \\ \quad \operatorname{nth}(2, d) \in \operatorname{msg\text{-}set}(t, k) \wedge \\ \quad d = {[}\operatorname{nth}(1, d), \operatorname{nth}(2, d)] \end{array}}}
\vnbprop{node-step-in-directed}{\vnbstmt{\begin{array}{@{}l@{}} \forall a, n. \\ \quad (\operatorname{is\text{-}strand\text{-}space}(a) \wedge (n \in \operatorname{nodes}(a))) \Rightarrow  \\ \quad \quad \operatorname{nth}(\operatorname{nth}(1, n), \operatorname{strand\text{-}map}(a)(\operatorname{nth}(2, n))) \in \operatorname{directed}(\operatorname{tokens}(a), \operatorname{keys}(a)) \end{array}}}
\vnbprop{node-content-in-msgs}{\vnbstmt{\begin{array}{@{}l@{}} \forall a, n. \\ \quad (\operatorname{is\text{-}strand\text{-}space}(a) \wedge (n \in \operatorname{nodes}(a))) \Rightarrow  \\ \quad \quad \operatorname{node\text{-}content}(a, n) \in \operatorname{msg\text{-}set}(\operatorname{tokens}(a), \operatorname{keys}(a)) \end{array}}}
\vnbprop{node-dir-cases}{\vnbstmt{\begin{array}{@{}l@{}} \forall a, n. \\ \quad (\operatorname{is\text{-}strand\text{-}space}(a) \wedge (n \in \operatorname{nodes}(a))) \Rightarrow  \\ \quad \quad (\operatorname{node\text{-}dir}(a, n) = \mathit{sent}) \vee (\operatorname{node\text{-}dir}(a, n) = \mathit{received}) \end{array}}}
\vnbprop{sent-ne-received}{\vnbstmt{\neg (\mathit{sent} = \mathit{received})}}
\vnbprop{node-sent-or-received}{\vnbstmt{\begin{array}{@{}l@{}} \forall a, n. \\ \quad (\operatorname{is\text{-}strand\text{-}space}(a) \wedge (n \in \operatorname{nodes}(a))) \Rightarrow (\operatorname{is\text{-}sent}(a, n) \vee \operatorname{is\text{-}received}(a, n)) \end{array}}}
\vnbprop{node-not-sent-and-received}{\vnbstmt{\forall a, n.\; (\operatorname{is\text{-}sent}(a, n) \Rightarrow \neg \operatorname{is\text{-}received}(a, n))}}
\vnbprop{node-content-eq}{\vnbstmt{\begin{array}{@{}l@{}} \forall a, n, k, d. \\ \quad ((\operatorname{nth}(1, n) = k) \wedge (\operatorname{nth}(k, \operatorname{strand\text{-}map}(a)(\operatorname{nth}(2, n))) = d)) \Rightarrow  \\ \quad \quad (\operatorname{node\text{-}content}(a, n) \simeq \operatorname{nth}(2, d)) \wedge (\operatorname{node\text{-}dir}(a, n) \simeq \operatorname{nth}(1, d)) \end{array}}}
\vnbprop{ssb-irreflexive}{\vnbstmt{\forall a, n.\; \neg \operatorname{same\text{-}strand\text{-}before}(a, n, n)}}
\vnbprop{ssb-transitive}{\vnbstmt{\begin{array}{@{}l@{}} \forall a, x, y, z. \\ \quad (\operatorname{same\text{-}strand\text{-}before}(a, x, y) \wedge \operatorname{same\text{-}strand\text{-}before}(a, y, z)) \Rightarrow  \\ \quad \quad \operatorname{same\text{-}strand\text{-}before}(a, x, z) \end{array}}}
\vnbprop{ssu-of-ssb}{\vnbstmt{\forall a, x, y.\; (\operatorname{same\text{-}strand\text{-}before}(a, x, y) \Rightarrow \operatorname{same\text{-}strand\text{-}upto}(a, x, y))}}
\vnbprop{ssu-refl}{\vnbstmt{\forall a.; \forall x \in \operatorname{nodes}(a).\; \operatorname{same\text{-}strand\text{-}upto}(a, x, x)}}
\vnbprop{ssu-cases}{\vnbstmt{\begin{array}{@{}l@{}} \forall a, x, y. \\ \quad \operatorname{same\text{-}strand\text{-}upto}(a, x, y) \Rightarrow (\operatorname{same\text{-}strand\text{-}before}(a, x, y) \vee (x = y)) \end{array}}}
\vnbprop{sends-to-dirs}{\vnbstmt{\begin{array}{@{}l@{}} \forall a, x, y. \\ \quad \operatorname{sends\text{-}to}(a, x, y) \Rightarrow  \\ \quad \quad \operatorname{is\text{-}sent}(a, x) \wedge (\operatorname{is\text{-}received}(a, y) \wedge (\operatorname{node\text{-}content}(a, x) = \operatorname{node\text{-}content}(a, y))) \end{array}}}
\vnbprop{bundle-space}{\vnbstmt{\forall c.\; (\operatorname{is\text{-}bundle}(c) \Rightarrow \operatorname{is\text{-}strand\text{-}space}(\operatorname{bspace}(c)))}}
\vnbprop{bundle-nodes-set}{\vnbstmt{\forall c.\; (\operatorname{is\text{-}bundle}(c) \Rightarrow (\operatorname{bnodes}(c) \in \mathbf{Set}))}}
\vnbprop{bundle-send-typing}{\vnbstmt{\begin{array}{@{}l@{}} \forall c. \\ \quad \operatorname{is\text{-}bundle}(c) \Rightarrow (\operatorname{bsend}(c) \in \operatorname{power}((\operatorname{bnodes}(c) \times \operatorname{bnodes}(c)))) \end{array}}}
\vnbprop{bundle-nodes-in-space}{\vnbstmt{\forall c.; \forall x \in \operatorname{bnodes}(c).\; (\operatorname{is\text{-}bundle}(c) \Rightarrow (x \in \operatorname{nodes}(\operatorname{bspace}(c))))}}
\vnbprop{bundle-finite}{\vnbstmt{\forall c.\; (\operatorname{is\text{-}bundle}(c) \Rightarrow (|\operatorname{bnodes}(c)| \in \mathbb{N}))}}
\vnbprop{bundle-send-sends}{\vnbstmt{\begin{array}{@{}l@{}} \forall c, x, y. \\ \quad (\operatorname{is\text{-}bundle}(c) \wedge ({[}x, y] \in \operatorname{bsend}(c))) \Rightarrow \operatorname{sends\text{-}to}(\operatorname{bspace}(c), x, y) \end{array}}}
\vnbprop{bundle-received-sender}{\vnbstmt{\begin{array}{@{}l@{}} \forall c.; \forall y \in \operatorname{bnodes}(c). \\ \quad (\operatorname{is\text{-}bundle}(c) \wedge \operatorname{is\text{-}received}(\operatorname{bspace}(c), y)) \Rightarrow  \\ \quad \quad \exists x. \\ \quad \quad \quad ({[}x, y] \in \operatorname{bsend}(c)) \wedge \forall z.\; (({[}z, y] \in \operatorname{bsend}(c)) \Rightarrow (z = x)) \end{array}}}
\vnbprop{bundle-ssb-closed}{\vnbstmt{\begin{array}{@{}l@{}} \forall c.; \forall y \in \operatorname{bnodes}(c), x. \\ \quad (\operatorname{is\text{-}bundle}(c) \wedge \operatorname{same\text{-}strand\text{-}before}(\operatorname{bspace}(c), x, y)) \Rightarrow (x \in \operatorname{bnodes}(c)) \end{array}}}
\vnbprop{bundle-acyclic}{\vnbstmt{\begin{array}{@{}l@{}} \forall c, x. \\ \quad \operatorname{is\text{-}bundle}(c) \Rightarrow \neg ({[}x, x] \in \operatorname{prec\text{-}strict}(\operatorname{bspace}(c), \operatorname{bnodes}(c), \operatorname{bsend}(c))) \end{array}}}
\vnbprop{pair-in-power-box}{\vnbstmt{\begin{array}{@{}l@{}} \forall r, d, x, y. \\ \quad ((r \in \operatorname{power}((d \times d))) \wedge ({[}x, y] \in r)) \Rightarrow ((x \in d) \wedge (y \in d)) \end{array}}}
\vnbprop{bundle-send-nodes}{\vnbstmt{\begin{array}{@{}l@{}} \forall c, x, y. \\ \quad (\operatorname{is\text{-}bundle}(c) \wedge ({[}x, y] \in \operatorname{bsend}(c))) \Rightarrow ((x \in \operatorname{bnodes}(c)) \wedge (y \in \operatorname{bnodes}(c))) \end{array}}}
\vnbprop{ps-in-box}{\vnbstmt{\begin{array}{@{}l@{}} \forall a, d, b, x, y. \\ \quad ({[}x, y] \in \operatorname{prec\text{-}strict}(a, d, b)) \Rightarrow ((x \in d) \wedge (y \in d)) \end{array}}}
\vnbprop{ps-least}{\vnbstmt{\begin{array}{@{}l@{}} \forall a, d, b.; \forall r \in \operatorname{power}((d \times d)), x, y. \\ \quad (\operatorname{is\text{-}prec\text{-}closed}(a, d, b, r) \wedge ({[}x, y] \in \operatorname{prec\text{-}strict}(a, d, b))) \Rightarrow  \\ \quad \quad {[}x, y] \in r \end{array}}}
\vnbprop{ps-send}{\vnbstmt{\begin{array}{@{}l@{}} \forall a, d, b, x, y. \\ \quad ((b \in \operatorname{power}((d \times d))) \wedge ({[}x, y] \in b)) \Rightarrow  \\ \quad \quad {[}x, y] \in \operatorname{prec\text{-}strict}(a, d, b) \end{array}}}
\vnbprop{ps-ssb}{\vnbstmt{\begin{array}{@{}l@{}} \forall a, d, b, x, y. \\ \quad ((x \in d) \wedge ((y \in d) \wedge \operatorname{same\text{-}strand\text{-}before}(a, x, y))) \Rightarrow  \\ \quad \quad {[}x, y] \in \operatorname{prec\text{-}strict}(a, d, b) \end{array}}}
\vnbprop{ps-trans}{\vnbstmt{\begin{array}{@{}l@{}} \forall a, d, b, x, y, z. \\ \quad (({[}x, y] \in \operatorname{prec\text{-}strict}(a, d, b)) \wedge ({[}y, z] \in \operatorname{prec\text{-}strict}(a, d, b))) \Rightarrow  \\ \quad \quad {[}x, z] \in \operatorname{prec\text{-}strict}(a, d, b) \end{array}}}
\vnbprop{precedes-strict-trans}{\vnbstmt{\begin{array}{@{}l@{}} \forall c, x, y, z. \\ \quad (\operatorname{is\text{-}bundle}(c) \wedge (({[}x, y] \in \operatorname{precedes\text{-}strict}(c)) \wedge ({[}y, z] \in \operatorname{precedes\text{-}strict}(c)))) \Rightarrow  \\ \quad \quad {[}x, z] \in \operatorname{precedes\text{-}strict}(c) \end{array}}}
\vnbprop{precedes-strict-irreflexive}{\vnbstmt{\forall c, x.\; (\operatorname{is\text{-}bundle}(c) \Rightarrow \neg ({[}x, x] \in \operatorname{precedes\text{-}strict}(c)))}}
\vnbprop{bundle-send-precedes}{\vnbstmt{\begin{array}{@{}l@{}} \forall c, x, y. \\ \quad (\operatorname{is\text{-}bundle}(c) \wedge ({[}x, y] \in \operatorname{bsend}(c))) \Rightarrow ({[}x, y] \in \operatorname{precedes\text{-}strict}(c)) \end{array}}}
\vnbprop{bundle-ssb-precedes}{\vnbstmt{\begin{array}{@{}l@{}} \forall c, x, y. \\ \quad (\operatorname{is\text{-}bundle}(c) \wedge ((y \in \operatorname{bnodes}(c)) \wedge \operatorname{same\text{-}strand\text{-}before}(\operatorname{bspace}(c), x, y))) \Rightarrow  \\ \quad \quad {[}x, y] \in \operatorname{precedes\text{-}strict}(c) \end{array}}}
\vnbprop{precedes-intro-strict}{\vnbstmt{\begin{array}{@{}l@{}} \forall c, x, y. \\ \quad (\operatorname{is\text{-}bundle}(c) \wedge ({[}x, y] \in \operatorname{precedes\text{-}strict}(c))) \Rightarrow ({[}x, y] \in \operatorname{precedes}(c)) \end{array}}}
\vnbprop{precedes-refl}{\vnbstmt{\forall c.; \forall x \in \operatorname{bnodes}(c).\; ({[}x, x] \in \operatorname{precedes}(c))}}
\vnbprop{precedes-elim}{\vnbstmt{\begin{array}{@{}l@{}} \forall c, x, y. \\ \quad ({[}x, y] \in \operatorname{precedes}(c)) \Rightarrow  \\ \quad \quad ({[}x, y] \in \operatorname{precedes\text{-}strict}(c)) \vee ((x = y) \wedge (x \in \operatorname{bnodes}(c))) \end{array}}}
\vnbprop{ps-last-step}{\vnbstmt{\begin{array}{@{}l@{}} \forall a, d, b, x, y. \\ \quad ((d \in \mathbf{Set}) \wedge ((b \in \operatorname{power}((d \times d))) \wedge ({[}x, y] \in \operatorname{prec\text{-}strict}(a, d, b)))) \Rightarrow  \\ \quad \quad \exists z. \\ \quad \quad \quad ({[}z, y] \in b) \vee ((z \in d) \wedge \operatorname{same\text{-}strand\text{-}before}(a, z, y)) \wedge \\ \quad \quad \quad (z = x) \vee ({[}x, z] \in \operatorname{prec\text{-}strict}(a, d, b)) \end{array}}}
\vnbprop{precedes-trans}{\vnbstmt{\begin{array}{@{}l@{}} \forall c, x, y, z. \\ \quad (\operatorname{is\text{-}bundle}(c) \wedge (({[}x, y] \in \operatorname{precedes}(c)) \wedge ({[}y, z] \in \operatorname{precedes}(c)))) \Rightarrow  \\ \quad \quad {[}x, z] \in \operatorname{precedes}(c) \end{array}}}
\vnbprop{precedes-antisym}{\vnbstmt{\begin{array}{@{}l@{}} \forall c, x, y. \\ \quad (\operatorname{is\text{-}bundle}(c) \wedge (({[}x, y] \in \operatorname{precedes}(c)) \wedge ({[}y, x] \in \operatorname{precedes}(c)))) \Rightarrow (x = y) \end{array}}}
\vnbprop{precedes-set}{\vnbstmt{\forall c.\; (\operatorname{is\text{-}bundle}(c) \Rightarrow (\operatorname{precedes}(c) \in \mathbf{Set}))}}
\vnbprop{precedes-partial-order}{\vnbstmt{\forall c.\; (\operatorname{is\text{-}bundle}(c) \Rightarrow \operatorname{is\text{-}partial\text{-}order}(\operatorname{bnodes}(c), \operatorname{precedes}(c)))}}
\vnbprop{finite-strict-minimal}{\vnbstmt{\begin{array}{@{}l@{}} \forall r.; \forall s \in \mathbf{Set}. \\ \quad \forall x, y, z.\; ((({[}x, y] \in r) \wedge ({[}y, z] \in r)) \Rightarrow ({[}x, z] \in r)) \wedge \\ \quad \forall x.\; \neg ({[}x, x] \in r) \wedge \\ \quad |s| \in \mathbb{N} \wedge \\ \quad \exists w.\; (w \in s) \Rightarrow \\ \quad \quad \exists m \in s.\; \forall z \in s.\; \neg ({[}z, m] \in r) \end{array}}}
\vnbprop{bundle-minimal}{\vnbstmt{\begin{array}{@{}l@{}} \forall c, s. \\ \quad (\operatorname{is\text{-}bundle}(c) \wedge (\forall z \in s.\; (z \in \operatorname{bnodes}(c)) \wedge \exists w.\; (w \in s))) \Rightarrow  \\ \quad \quad \exists m \in s.\; \forall z \in s.\; \neg ({[}z, m] \in \operatorname{precedes\text{-}strict}(c)) \end{array}}}
\vnbprop{bundle-minimal-precedes}{\vnbstmt{\begin{array}{@{}l@{}} \forall c, s. \\ \quad (\operatorname{is\text{-}bundle}(c) \wedge (\forall z \in s.\; (z \in \operatorname{bnodes}(c)) \wedge \exists w.\; (w \in s))) \Rightarrow  \\ \quad \quad \exists m \in s. \forall z \in s. \\ \quad \quad \quad \neg (z = m) \Rightarrow \neg ({[}z, m] \in \operatorname{precedes}(c)) \end{array}}}
\vnbprop{bundle-minimal-sent}{\vnbstmt{\begin{array}{@{}l@{}} \forall c, s, n. \\ \quad \operatorname{is\text{-}bundle}(c) \wedge \\ \quad \forall z \in s.\; (z \in \operatorname{bnodes}(c)) \wedge \\ \quad \forall x, y. \\ \quad \quad x \in \operatorname{bnodes}(c) \wedge \\ \quad \quad y \in \operatorname{bnodes}(c) \wedge \\ \quad \quad \operatorname{node\text{-}content}(\operatorname{bspace}(c), x) = \operatorname{node\text{-}content}(\operatorname{bspace}(c), y) \wedge \\ \quad \quad y \in s \Rightarrow \\ \quad \quad \quad x \in s \wedge \\ \quad n \in s \wedge \\ \quad \forall z \in s.\; \neg ({[}z, n] \in \operatorname{precedes\text{-}strict}(c)) \Rightarrow \\ \quad \quad \operatorname{is\text{-}sent}(\operatorname{bspace}(c), n) \end{array}}}
\vnbprop{originates-elim}{\vnbstmt{\begin{array}{@{}l@{}} \forall a, u, n. \\ \quad (\operatorname{is\text{-}strand\text{-}space}(a) \wedge \operatorname{originates}(a, u, n)) \Rightarrow  \\ \quad \quad n \in \operatorname{nodes}(a) \wedge \\ \quad \quad \operatorname{is\text{-}sent}(a, n) \wedge \\ \quad \quad \operatorname{is\text{-}subterm}(\operatorname{tokens}(a), \operatorname{keys}(a), u, \operatorname{node\text{-}content}(a, n)) \wedge \\ \quad \quad \forall m. \\ \quad \quad \quad \operatorname{same\text{-}strand\text{-}before}(a, m, n) \Rightarrow  \\ \quad \quad \quad \quad \neg \operatorname{is\text{-}subterm}(\operatorname{tokens}(a), \operatorname{keys}(a), u, \operatorname{node\text{-}content}(a, m)) \end{array}}}
\vnbprop{originates-intro}{\vnbstmt{\begin{array}{@{}l@{}} \forall a, u, n. \\ \quad \operatorname{is\text{-}strand\text{-}space}(a) \wedge \\ \quad n \in \operatorname{nodes}(a) \wedge \\ \quad \operatorname{is\text{-}sent}(a, n) \wedge \\ \quad \operatorname{is\text{-}subterm}(\operatorname{tokens}(a), \operatorname{keys}(a), u, \operatorname{node\text{-}content}(a, n)) \wedge \\ \quad \forall m. \\ \quad \quad \operatorname{same\text{-}strand\text{-}before}(a, m, n) \Rightarrow  \\ \quad \quad \quad \neg \operatorname{is\text{-}subterm}(\operatorname{tokens}(a), \operatorname{keys}(a), u, \operatorname{node\text{-}content}(a, m)) \Rightarrow \\ \quad \quad \operatorname{originates}(a, u, n) \end{array}}}
\vnbprop{originates-sent}{\vnbstmt{\begin{array}{@{}l@{}} \forall a, u, n. \\ \quad (\operatorname{is\text{-}strand\text{-}space}(a) \wedge \operatorname{originates}(a, u, n)) \Rightarrow  \\ \quad \quad \operatorname{is\text{-}sent}(a, n) \wedge \operatorname{is\text{-}subterm}(\operatorname{tokens}(a), \operatorname{keys}(a), u, \operatorname{node\text{-}content}(a, n)) \end{array}}}
\vnbprop{originates-no-earlier}{\vnbstmt{\begin{array}{@{}l@{}} \forall a, u, n, m. \\ \quad (\operatorname{is\text{-}strand\text{-}space}(a) \wedge (\operatorname{originates}(a, u, n) \wedge \operatorname{same\text{-}strand\text{-}before}(a, m, n))) \Rightarrow  \\ \quad \quad \neg \operatorname{is\text{-}subterm}(\operatorname{tokens}(a), \operatorname{keys}(a), u, \operatorname{node\text{-}content}(a, m)) \end{array}}}
\vnbprop{originates-unique}{\vnbstmt{\begin{array}{@{}l@{}} \forall a, u, n, m. \\ \quad (\operatorname{uniquely\text{-}originates}(a, u) \wedge (\operatorname{originates}(a, u, n) \wedge \operatorname{originates}(a, u, m))) \Rightarrow  \\ \quad \quad m = n \end{array}}}
\vnbprop{bundle-minimal-originates}{\vnbstmt{\begin{array}{@{}l@{}} \forall c, u, n. \\ \quad \operatorname{is\text{-}bundle}(c) \wedge \\ \quad n \in \operatorname{bnodes}(c) \wedge \\ \quad \operatorname{is\text{-}subterm}(\operatorname{tokens}(\operatorname{bspace}(c)), \operatorname{keys}(\operatorname{bspace}(c)), u, \operatorname{node\text{-}content}(\operatorname{bspace}(c), n)) \wedge \\ \quad \forall z \in \operatorname{bnodes}(c). \\ \quad \quad \operatorname{is\text{-}subterm}(\operatorname{tokens}(\operatorname{bspace}(c)), \operatorname{keys}(\operatorname{bspace}(c)), u, \operatorname{node\text{-}content}(\operatorname{bspace}(c), z)) \Rightarrow  \\ \quad \quad \quad \neg ({[}z, n] \in \operatorname{precedes\text{-}strict}(c)) \Rightarrow \\ \quad \quad \operatorname{originates}(\operatorname{bspace}(c), u, n) \end{array}}}

\vnbbookfile{structure-library/strand-ideals}{structure-library/strand-ideals}
\noindent{\small strand-ideals.scm -- IDEALS OF THE MESSAGE ALGEBRA (Thayer, Herzog, Guttman, "Strand spaces: proving security protocols correct", section 6.1, Definitions 6.1, 6.2, 6.5 and the stages of Lemma 6.4), ON THE USER'S MODEL (notes-60 with encryption: tokens, tagged n-ary lists LST(l), encryptions ENCR(m, y); ...}\par
\par\medskip\noindent\textbf{Definition} (\texttt{is-key-ideal}).\;For all $\operatorname{sbit\_}$, $\operatorname{sbik\_}$, $\operatorname{sbiks\_}$ and $\operatorname{sbii\_}$, 
\begin{equation*}
\begin{array}{@{}l@{}} \operatorname{is\text{-}key\text{-}ideal}(\operatorname{sbit\_}, \operatorname{sbik\_}, \operatorname{sbiks\_}, \operatorname{sbii\_}) \\ \Leftrightarrow\; \forall \operatorname{sbix\_} \in \operatorname{sbii\_}.\; (\operatorname{sbix\_} \in \operatorname{msg\text{-}set}(\operatorname{sbit\_}, \operatorname{sbik\_})) \wedge \\ \forall \operatorname{sbil\_} \in \operatorname{tuples}(\operatorname{msg\text{-}set}(\operatorname{sbit\_}, \operatorname{sbik\_})), \operatorname{sbij\_} \in \mathbb{N}. \\ \quad (1 \leq \operatorname{sbij\_}) \Rightarrow  \\ \quad (\operatorname{sbij\_} \leq \operatorname{length}(\operatorname{sbil\_})) \Rightarrow  \\ \quad (\operatorname{nth}(\operatorname{sbij\_}, \operatorname{sbil\_}) \in \operatorname{sbii\_}) \Rightarrow (\operatorname{lst}(\operatorname{sbil\_}) \in \operatorname{sbii\_}) \wedge \\ \forall \operatorname{sbih\_} \in \operatorname{sbii\_}, \operatorname{sbiy\_} \in \operatorname{sbiks\_}. \\ \quad \operatorname{encr}(\operatorname{sbih\_}, \operatorname{sbiy\_}) \in \operatorname{sbii\_} \end{array}
\end{equation*}

\par\medskip\noindent\textbf{Definition} (\texttt{key-ideal-gen}).\;$\operatorname{key\text{-}ideal\text{-}gen}(\operatorname{sbgt\_}, \operatorname{sbgk\_}, \operatorname{sbgks\_}, \operatorname{sbgs\_}) :=$
\begin{equation*}
\{\operatorname{sbgm\_} \in \operatorname{msg\text{-}set}(\operatorname{sbgt\_}, \operatorname{sbgk\_}) : \forall \operatorname{sbgq\_} \in \operatorname{power}(\operatorname{msg\text{-}set}(\operatorname{sbgt\_}, \operatorname{sbgk\_})).\; (\operatorname{is\text{-}key\text{-}ideal}(\operatorname{sbgt\_}, \operatorname{sbgk\_}, \operatorname{sbgks\_}, \operatorname{sbgq\_}) \Rightarrow (\forall \operatorname{sbgz\_} \in \operatorname{sbgs\_}.\; (\operatorname{sbgz\_} \in \operatorname{sbgq\_}) \Rightarrow (\operatorname{sbgm\_} \in \operatorname{sbgq\_})))\}
\end{equation*}

\par\medskip\noindent\textbf{Definition} (\texttt{is-simple-msg}).\;For all $\operatorname{sbst\_}$, $\operatorname{sbsk\_}$ and $\operatorname{sbsm\_}$, 
\begin{equation*}
\begin{array}{@{}l@{}} \operatorname{is\text{-}simple\text{-}msg}(\operatorname{sbst\_}, \operatorname{sbsk\_}, \operatorname{sbsm\_}) \\ \Leftrightarrow\; \neg \exists \operatorname{sbsl\_} \in \operatorname{tuples}(\operatorname{msg\text{-}set}(\operatorname{sbst\_}, \operatorname{sbsk\_})).\; (\operatorname{sbsm\_} = \operatorname{lst}(\operatorname{sbsl\_})) \end{array}
\end{equation*}

\par\medskip\noindent\textbf{Definition} (\texttt{ideal-stage}).\;For all $\operatorname{sbrt\_}$, $\operatorname{sbrk\_}$, $\operatorname{sbrks\_}$ and $\operatorname{sbrs\_}$,  $(\operatorname{ideal\text{-}stage}(\operatorname{sbrt\_}, \operatorname{sbrk\_}, \operatorname{sbrks\_}, \operatorname{sbrs\_}, 0) \simeq \operatorname{sbrs\_})$\par
Let $\operatorname{sbrt\_}$, $\operatorname{sbrk\_}$, $\operatorname{sbrks\_}$ and $\operatorname{sbrs\_}$ be arbitrary. Let $\operatorname{sbrn\_} \in \mathbb{N}$. Then
\begin{equation*}
(\operatorname{ideal\text{-}stage}(\operatorname{sbrt\_}, \operatorname{sbrk\_}, \operatorname{sbrks\_}, \operatorname{sbrs\_}, \operatorname{succ}(\operatorname{sbrn\_})) \simeq \{\operatorname{sbrm\_} \in \operatorname{msg\text{-}set}(\operatorname{sbrt\_}, \operatorname{sbrk\_}) : \exists \operatorname{sbrg\_} \in \operatorname{key\text{-}ideal\text{-}gen}(\operatorname{sbrt\_}, \operatorname{sbrk\_}, \emptyset, \operatorname{ideal\text{-}stage}(\operatorname{sbrt\_}, \operatorname{sbrk\_}, \operatorname{sbrks\_}, \operatorname{sbrs\_}, \operatorname{sbrn\_})), \operatorname{sbry\_} \in \operatorname{sbrks\_}.\; (\operatorname{sbrm\_} = \operatorname{encr}(\operatorname{sbrg\_}, \operatorname{sbry\_}))\})
\end{equation*}


\vnbbookfile{theorem-library/strand-ideals-laws}{theorem-library/strand-ideals-laws}
\noindent{\small strand-ideals-laws.scm -- IDEALS OF THE MESSAGE ALGEBRA: Propositions 6.3, 6.6, 6.7, 6.8, Lemma 6.4 and Proposition 6.9 of Thayer, Herzog, Guttman, "Strand spaces: proving security protocols correct" (\textasciitilde{}/docs/jcs\_strand\_spaces.pdf, section 6), on the user's model (notes-60 with encryption).  Vocabulary: ...}\par
\vnbprop{key-ideal-gen-unfold}{\vnbstmt{\begin{array}{@{}l@{}} \forall t, k, a, s. \\ \quad (\operatorname{key\text{-}ideal\text{-}gen}(t, k, a, s) \simeq \{m \in \operatorname{msg\text{-}set}(t, k) : \forall q \in \operatorname{power}(\operatorname{msg\text{-}set}(t, k)).\; ((\operatorname{is\text{-}key\text{-}ideal}(t, k, a, q) \wedge \forall z \in s.\; (z \in q)) \Rightarrow (m \in q))\}) \end{array}}}
\vnbprop{key-ideal-gen-elim}{\vnbstmt{\begin{array}{@{}l@{}} \forall t, k, a, s, m, i. \\ \quad m \in \operatorname{key\text{-}ideal\text{-}gen}(t, k, a, s) \wedge \\ \quad i \in \operatorname{power}(\operatorname{msg\text{-}set}(t, k)) \wedge \\ \quad \operatorname{is\text{-}key\text{-}ideal}(t, k, a, i) \wedge \\ \quad \forall z \in s.\; (z \in i) \Rightarrow \\ \quad \quad m \in i \end{array}}}
\vnbprop{key-ideal-gen-in-msgs}{\vnbstmt{\forall t, k, a, s.; \forall m \in \operatorname{key\text{-}ideal\text{-}gen}(t, k, a, s).\; (m \in \operatorname{msg\text{-}set}(t, k))}}
\vnbprop{key-ideal-gen-least}{\vnbstmt{\begin{array}{@{}l@{}} \forall t, k, a, s, i.; \forall m \in \operatorname{key\text{-}ideal\text{-}gen}(t, k, a, s). \\ \quad ((t \in \mathbf{Set}) \wedge ((k \in \mathbf{Set}) \wedge (\operatorname{is\text{-}key\text{-}ideal}(t, k, a, i) \wedge \forall z \in s.\; (z \in i)))) \Rightarrow  \\ \quad \quad m \in i \end{array}}}
\vnbprop{key-ideal-gen-contains}{\vnbstmt{\begin{array}{@{}l@{}} \forall t, k, a, s.; \forall x \in s. \\ \quad (x \in \operatorname{msg\text{-}set}(t, k)) \Rightarrow (x \in \operatorname{key\text{-}ideal\text{-}gen}(t, k, a, s)) \end{array}}}
\vnbprop{key-ideal-gen-is-ideal}{\vnbstmt{\begin{array}{@{}l@{}} \forall t, k, a, s. \\ \quad ((t \in \mathbf{Set}) \wedge ((k \in \mathbf{Set}) \wedge \forall y \in a.\; (y \in k))) \Rightarrow  \\ \quad \quad \operatorname{is\text{-}key\text{-}ideal}(t, k, a, \operatorname{key\text{-}ideal\text{-}gen}(t, k, a, s)) \end{array}}}
\vnbprop{key-ideal-gen-self}{\vnbstmt{\forall t, k, s.; \forall x \in \operatorname{msg\text{-}set}(t, k).\; (x \in \operatorname{key\text{-}ideal\text{-}gen}(t, k, s, \operatorname{singleton}(x)))}}
\vnbprop{key-ideal-gen-singleton-least}{\vnbstmt{\begin{array}{@{}l@{}} \forall t, k, s, i, x, m. \\ \quad t \in \mathbf{Set} \wedge \\ \quad k \in \mathbf{Set} \wedge \\ \quad \operatorname{is\text{-}key\text{-}ideal}(t, k, s, i) \wedge \\ \quad x \in i \wedge \\ \quad m \in \operatorname{key\text{-}ideal\text{-}gen}(t, k, s, \operatorname{singleton}(x)) \Rightarrow \\ \quad \quad m \in i \end{array}}}
\vnbprop{key-ideal-gen-union}{\vnbstmt{\begin{array}{@{}l@{}} \forall t, k, a, s. \\ \quad ((t \in \mathbf{Set}) \wedge ((k \in \mathbf{Set}) \wedge (\forall y \in a.\; (y \in k) \wedge \forall z \in s.\; (z \in \operatorname{msg\text{-}set}(t, k))))) \Rightarrow  \\ \quad \quad \operatorname{key\text{-}ideal\text{-}gen}(t, k, a, s) = \operatorname{big\text{-}union}(x, s, \operatorname{key\text{-}ideal\text{-}gen}(t, k, a, \operatorname{singleton}(x))) \end{array}}}
\vnbprop{key-ideal-empty-keys}{\vnbstmt{\begin{array}{@{}l@{}} \forall t, k, s, i. \\ \quad \operatorname{is\text{-}key\text{-}ideal}(t, k, s, i) \Rightarrow \operatorname{is\text{-}key\text{-}ideal}(t, k, \emptyset, i) \end{array}}}
\vnbprop{ideal-stage-in-power}{\vnbstmt{\begin{array}{@{}l@{}} \forall i \in \mathbb{N}, t, k, a, s. \\ \quad ((t \in \mathbf{Set}) \wedge ((k \in \mathbf{Set}) \wedge \forall z \in s.\; (z \in \operatorname{msg\text{-}set}(t, k)))) \Rightarrow  \\ \quad \quad \operatorname{ideal\text{-}stage}(t, k, a, s, i) \in \operatorname{power}(\operatorname{msg\text{-}set}(t, k)) \end{array}}}
\vnbprop{ideal-stage-in-gen}{\vnbstmt{\begin{array}{@{}l@{}} \forall i \in \mathbb{N}, t, k, a, s.; \forall z \in \operatorname{ideal\text{-}stage}(t, k, a, s, i). \\ \quad ((t \in \mathbf{Set}) \wedge ((k \in \mathbf{Set}) \wedge (\forall y \in a.\; (y \in k) \wedge \forall z \in s.\; (z \in \operatorname{msg\text{-}set}(t, k))))) \Rightarrow  \\ \quad \quad z \in \operatorname{key\text{-}ideal\text{-}gen}(t, k, a, s) \end{array}}}
\vnbprop{key-ideal-gen-stages}{\vnbstmt{\begin{array}{@{}l@{}} \forall t, k, a, s. \\ \quad ((t \in \mathbf{Set}) \wedge ((k \in \mathbf{Set}) \wedge (\forall y \in a.\; (y \in k) \wedge \forall z \in s.\; (z \in \operatorname{msg\text{-}set}(t, k))))) \Rightarrow  \\ \quad \quad \operatorname{key\text{-}ideal\text{-}gen}(t, k, a, s) \\ \quad \quad =\; \operatorname{big\text{-}union}(i, \mathbb{N}, \operatorname{key\text{-}ideal\text{-}gen}(t, k, \emptyset, \operatorname{ideal\text{-}stage}(t, k, a, s, i))) \end{array}}}
\vnbprop{key-ideal-gen-cases}{\vnbstmt{\begin{array}{@{}l@{}} \forall t, k, a, s, x. \\ \quad t \in \mathbf{Set} \wedge \\ \quad k \in \mathbf{Set} \wedge \\ \quad \forall y \in a.\; (y \in k) \wedge \\ \quad \forall z \in s.\; (z \in \operatorname{msg\text{-}set}(t, k)) \wedge \\ \quad x \in \operatorname{key\text{-}ideal\text{-}gen}(t, k, a, s) \Rightarrow \\ \quad \quad x \in s \vee \\ \quad \quad \exists l \in \operatorname{tuples}(\operatorname{msg\text{-}set}(t, k)). \\ \quad \quad \quad x = \operatorname{lst}(l) \wedge \\ \quad \quad \quad \exists j \in \mathbb{N}. \\ \quad \quad \quad \quad (1 \leq j) \wedge ((j \leq \operatorname{length}(l)) \wedge (\operatorname{nth}(j, l) \in \operatorname{key\text{-}ideal\text{-}gen}(t, k, a, s))) \vee \\ \quad \quad \exists h \in \operatorname{key\text{-}ideal\text{-}gen}(t, k, a, s), y \in a.\; (x = \operatorname{encr}(h, y)) \end{array}}}
\vnbprop{key-ideal-gen-decrypt}{\vnbstmt{\begin{array}{@{}l@{}} \forall t, k, a, s, h, y. \\ \quad t \in \mathbf{Set} \wedge \\ \quad k \in \mathbf{Set} \wedge \\ \quad \forall w \in a.\; (w \in k) \wedge \\ \quad y \in k \wedge \\ \quad \forall z \in s.\; (z \in \operatorname{msg\text{-}set}(t, k)) \wedge \\ \quad \forall u \in s.\; \operatorname{is\text{-}simple\text{-}msg}(t, k, u) \wedge \\ \quad \forall u \in s, g.\; \neg (u = \operatorname{encr}(g, y)) \wedge \\ \quad \operatorname{encr}(h, y) \in \operatorname{key\text{-}ideal\text{-}gen}(t, k, a, s) \Rightarrow \\ \quad \quad h \in \operatorname{key\text{-}ideal\text{-}gen}(t, k, a, s) \end{array}}}
\vnbprop{key-ideal-gen-encr-key}{\vnbstmt{\begin{array}{@{}l@{}} \forall t, k, a, s, h, y. \\ \quad t \in \mathbf{Set} \wedge \\ \quad k \in \mathbf{Set} \wedge \\ \quad \forall w \in a.\; (w \in k) \wedge \\ \quad y \in k \wedge \\ \quad \forall z \in s.\; (z \in \operatorname{msg\text{-}set}(t, k)) \wedge \\ \quad \forall u \in s.\; \operatorname{is\text{-}simple\text{-}msg}(t, k, u) \wedge \\ \quad \forall u \in s, g.\; \neg (u = \operatorname{encr}(g, y)) \wedge \\ \quad \operatorname{encr}(h, y) \in \operatorname{key\text{-}ideal\text{-}gen}(t, k, a, s) \Rightarrow \\ \quad \quad y \in a \end{array}}}
\vnbprop{key-ideal-gen-list-component}{\vnbstmt{\begin{array}{@{}l@{}} \forall t, k, a, s, l. \\ \quad t \in \mathbf{Set} \wedge \\ \quad k \in \mathbf{Set} \wedge \\ \quad \forall w \in a.\; (w \in k) \wedge \\ \quad \forall z \in s.\; (z \in \operatorname{msg\text{-}set}(t, k)) \wedge \\ \quad \forall u \in s.\; \operatorname{is\text{-}simple\text{-}msg}(t, k, u) \wedge \\ \quad l \in \operatorname{tuples}(\operatorname{msg\text{-}set}(t, k)) \wedge \\ \quad \operatorname{lst}(l) \in \operatorname{key\text{-}ideal\text{-}gen}(t, k, a, s) \Rightarrow \\ \quad \quad \exists j \in \mathbb{N}. \\ \quad \quad \quad (1 \leq j) \wedge ((j \leq \operatorname{length}(l)) \wedge (\operatorname{nth}(j, l) \in \operatorname{key\text{-}ideal\text{-}gen}(t, k, a, s))) \end{array}}}
\vnbprop{bundle-minimal-entry-point}{\vnbstmt{\begin{array}{@{}l@{}} \forall c, i, n. \\ \quad \operatorname{is\text{-}bundle}(c) \wedge \\ \quad n \in \operatorname{bnodes}(c) \wedge \\ \quad \operatorname{node\text{-}content}(\operatorname{bspace}(c), n) \in i \wedge \\ \quad \forall z \in \operatorname{bnodes}(c). \\ \quad \quad (\operatorname{node\text{-}content}(\operatorname{bspace}(c), z) \in i) \Rightarrow \neg ({[}z, n] \in \operatorname{precedes\text{-}strict}(c)) \Rightarrow \\ \quad \quad \operatorname{is\text{-}entry\text{-}point}(\operatorname{bspace}(c), n, i) \end{array}}}

\vnbbookfile{theorem-library/strand-tokens-nn}{theorem-library/strand-tokens-nn}
\noindent{\small sb5-tokens.scm -- every subset of NN is a token set; a nonempty token set exists. Agent SB-5, 2026-10-26.  Destination: theorem-library/strand-tokens-nn.scm (library strand). Load window: after structure-library/strand-messages (IS-TOKEN-SET) and list-recursion (nn-not-a-list).  Helper prefix sbt-.}\par
\vnbprop{nn-subset-token-set}{\vnbstmt{\forall t.\; ((t \subseteq \mathbb{N}) \Rightarrow \operatorname{is\text{-}token\text{-}set}(t))}}
\vnbprop{nn-is-token-set}{\vnbstmt{\operatorname{is\text{-}token\text{-}set}(\mathbb{N})}}
\vnbprop{token-set-nonempty-exists}{\vnbstmt{\exists t.\; (\operatorname{is\text{-}token\text{-}set}(t) \wedge \exists x.\; (x \in t))}}

\vnbbookfile{structure-library/strand-adversary}{structure-library/strand-adversary}
\noindent{\small strand-adversary.scm -- THE ADVERSARY: its traces, the infiltrated strand space (The paper's Definitions 3.1 and 3.2, Thayer, Herzog, Guttman, \textasciitilde{}/docs/jcs\_strand\_spaces.pdf, section 3.1, on the user's model of notes-60 with ENCRYPTION (the user's decision of 2026-10-06 16:40Z): the binary C and S traces made n-ary, the ...}\par
\vnbstructdef{infiltrated-space}{An infiltrated space $s$ is a 3-tuple whose components are listed below.\par
\begin{enumerate}[nosep]
\item $\operatorname{ispace}(s) = \operatorname{nth}(1, s)$, a substructure (\texttt{substructure}): $\operatorname{is\text{-}strand\text{-}space}(\operatorname{ispace}(s))$.
\item $\operatorname{adv\text{-}ids}(s) = \operatorname{nth}(2, s)$, a carrier (\texttt{carriers}): $\operatorname{adv\text{-}ids}(s) \in \mathbf{Set}$.
\item $\operatorname{adv\text{-}tokens}(s) = \operatorname{nth}(3, s)$, a carrier (\texttt{carriers}): $\operatorname{adv\text{-}tokens}(s) \in \mathbf{Set}$.
\end{enumerate}
The laws (the \texttt{property} and \texttt{law} clauses):
\begin{itemize}[nosep]
\item $\forall \operatorname{sap\_} \in \operatorname{adv\text{-}ids}(s).\; (\operatorname{sap\_} \in \operatorname{ids}(\operatorname{ispace}(s)))$
\item $\forall \operatorname{sat\_} \in \operatorname{adv\text{-}tokens}(s).\; (\operatorname{sat\_} \in \operatorname{tokens}(\operatorname{ispace}(s)))$
\item
\begin{equation*}
\begin{array}{@{}l@{}} \forall \operatorname{sap\_} \in \operatorname{adv\text{-}ids}(s). \\ \quad \operatorname{is\text{-}adversary\text{-}trace}(\operatorname{ispace}(s), \operatorname{adv\text{-}tokens}(s), \operatorname{strand\text{-}map}(\operatorname{ispace}(s))(\operatorname{sap\_})) \end{array}
\end{equation*}

\end{itemize}
The characteristic law (the axiom \texttt{is-infiltrated-space}): $\operatorname{is\text{-}infiltrated\text{-}space}(s)$ holds if and only if $\operatorname{length}(s) = 3$, every component has the typing listed with it, every law holds, and $s \in \mathbf{Set}$ (the structure is a set).
}
\par\medskip\noindent\textbf{Definition} (\texttt{is-text-trace}).\;For all $\operatorname{sasp\_}$, $\operatorname{sap\_}$ and $\operatorname{sas\_}$, 
\begin{equation*}
\operatorname{is\text{-}text\text{-}trace}(\operatorname{sasp\_}, \operatorname{sap\_}, \operatorname{sas\_}) \Leftrightarrow \exists \operatorname{saa\_} \in \operatorname{sap\_}.\; (\operatorname{sas\_} = {[}{[}\mathit{sent}, \operatorname{saa\_}]])
\end{equation*}

\par\medskip\noindent\textbf{Definition} (\texttt{is-flush-trace}).\;For all $\operatorname{sasp\_}$ and $\operatorname{sas\_}$, 
\begin{equation*}
\begin{array}{@{}l@{}} \operatorname{is\text{-}flush\text{-}trace}(\operatorname{sasp\_}, \operatorname{sas\_}) \\ \Leftrightarrow\; \exists \operatorname{sag\_} \in \operatorname{msg\text{-}set}(\operatorname{tokens}(\operatorname{sasp\_}), \operatorname{keys}(\operatorname{sasp\_})). \\ \quad \operatorname{sas\_} = {[}{[}\mathit{received}, \operatorname{sag\_}]] \end{array}
\end{equation*}

\par\medskip\noindent\textbf{Definition} (\texttt{is-tee-trace}).\;For all $\operatorname{sasp\_}$ and $\operatorname{sas\_}$, 
\begin{equation*}
\begin{array}{@{}l@{}} \operatorname{is\text{-}tee\text{-}trace}(\operatorname{sasp\_}, \operatorname{sas\_}) \\ \Leftrightarrow\; \exists \operatorname{sag\_} \in \operatorname{msg\text{-}set}(\operatorname{tokens}(\operatorname{sasp\_}), \operatorname{keys}(\operatorname{sasp\_})). \\ \quad \operatorname{sas\_} = {[}{[}\mathit{received}, \operatorname{sag\_}], {[}\mathit{sent}, \operatorname{sag\_}], {[}\mathit{sent}, \operatorname{sag\_}]] \end{array}
\end{equation*}

\par\medskip\noindent\textbf{Definition} (\texttt{is-concat-trace}).\;For all $\operatorname{sasp\_}$ and $\operatorname{sas\_}$, 
\begin{equation*}
\begin{array}{@{}l@{}} \operatorname{is\text{-}concat\text{-}trace}(\operatorname{sasp\_}, \operatorname{sas\_}) \\ \Leftrightarrow\; \exists \operatorname{sal\_} \in \operatorname{tuples}(\operatorname{msg\text{-}set}(\operatorname{tokens}(\operatorname{sasp\_}), \operatorname{keys}(\operatorname{sasp\_}))). \\ \quad \operatorname{length}(\operatorname{sas\_}) = \operatorname{succ}(\operatorname{length}(\operatorname{sal\_})) \wedge \\ \quad \forall \operatorname{sai\_} \in \mathbb{N}. \\ \quad \quad (1 \leq \operatorname{sai\_}) \Rightarrow  \\ \quad \quad (\operatorname{sai\_} \leq \operatorname{length}(\operatorname{sal\_})) \Rightarrow (\operatorname{nth}(\operatorname{sai\_}, \operatorname{sas\_}) = {[}\mathit{received}, \operatorname{nth}(\operatorname{sai\_}, \operatorname{sal\_})]) \wedge \\ \quad \operatorname{nth}(\operatorname{succ}(\operatorname{length}(\operatorname{sal\_})), \operatorname{sas\_}) = {[}\mathit{sent}, \operatorname{lst}(\operatorname{sal\_})] \end{array}
\end{equation*}

\par\medskip\noindent\textbf{Definition} (\texttt{is-separate-trace}).\;For all $\operatorname{sasp\_}$ and $\operatorname{sas\_}$, 
\begin{equation*}
\begin{array}{@{}l@{}} \operatorname{is\text{-}separate\text{-}trace}(\operatorname{sasp\_}, \operatorname{sas\_}) \\ \Leftrightarrow\; \exists \operatorname{sal\_} \in \operatorname{tuples}(\operatorname{msg\text{-}set}(\operatorname{tokens}(\operatorname{sasp\_}), \operatorname{keys}(\operatorname{sasp\_}))). \\ \quad \operatorname{length}(\operatorname{sas\_}) = \operatorname{succ}(\operatorname{length}(\operatorname{sal\_})) \wedge \\ \quad \operatorname{nth}(1, \operatorname{sas\_}) = {[}\mathit{received}, \operatorname{lst}(\operatorname{sal\_})] \wedge \\ \quad \forall \operatorname{sai\_} \in \mathbb{N}. \\ \quad \quad (1 \leq \operatorname{sai\_}) \Rightarrow  \\ \quad \quad (\operatorname{sai\_} \leq \operatorname{length}(\operatorname{sal\_})) \Rightarrow (\operatorname{nth}(\operatorname{succ}(\operatorname{sai\_}), \operatorname{sas\_}) = {[}\mathit{sent}, \operatorname{nth}(\operatorname{sai\_}, \operatorname{sal\_})]) \end{array}
\end{equation*}

\par\medskip\noindent\textbf{Definition} (\texttt{is-encrypt-trace}).\;For all $\operatorname{sasp\_}$ and $\operatorname{sas\_}$, 
\begin{equation*}
\begin{array}{@{}l@{}} \operatorname{is\text{-}encrypt\text{-}trace}(\operatorname{sasp\_}, \operatorname{sas\_}) \\ \Leftrightarrow\; \exists \operatorname{sak\_} \in \operatorname{keys}(\operatorname{sasp\_}), \operatorname{sag\_} \in \operatorname{msg\text{-}set}(\operatorname{tokens}(\operatorname{sasp\_}), \operatorname{keys}(\operatorname{sasp\_})). \\ \quad \operatorname{sas\_} = {[}{[}\mathit{received}, \operatorname{sak\_}], {[}\mathit{received}, \operatorname{sag\_}], {[}\mathit{sent}, \operatorname{encr}(\operatorname{sag\_}, \operatorname{sak\_})]] \end{array}
\end{equation*}

\par\medskip\noindent\textbf{Definition} (\texttt{is-decrypt-trace}).\;For all $\operatorname{sasp\_}$ and $\operatorname{sas\_}$, 
\begin{equation*}
\begin{array}{@{}l@{}} \operatorname{is\text{-}decrypt\text{-}trace}(\operatorname{sasp\_}, \operatorname{sas\_}) \\ \Leftrightarrow\; \exists \operatorname{sak\_} \in \operatorname{keys}(\operatorname{sasp\_}), \operatorname{sag\_} \in \operatorname{msg\text{-}set}(\operatorname{tokens}(\operatorname{sasp\_}), \operatorname{keys}(\operatorname{sasp\_})). \\ \quad \operatorname{sas\_} = {[}{[}\mathit{received}, \operatorname{kinv}(\operatorname{sasp\_})(\operatorname{sak\_})], {[}\mathit{received}, \operatorname{encr}(\operatorname{sag\_}, \operatorname{sak\_})], {[}\mathit{sent}, \operatorname{sag\_}]] \end{array}
\end{equation*}

\par\medskip\noindent\textbf{Definition} (\texttt{is-adversary-trace}).\;For all $\operatorname{sasp\_}$, $\operatorname{sap\_}$ and $\operatorname{sas\_}$, 
\begin{equation*}
\begin{array}{@{}l@{}} \operatorname{is\text{-}adversary\text{-}trace}(\operatorname{sasp\_}, \operatorname{sap\_}, \operatorname{sas\_}) \\ \Leftrightarrow\; \operatorname{is\text{-}text\text{-}trace}(\operatorname{sasp\_}, \operatorname{sap\_}, \operatorname{sas\_}) \vee \\ \operatorname{is\text{-}flush\text{-}trace}(\operatorname{sasp\_}, \operatorname{sas\_}) \vee \\ \operatorname{is\text{-}tee\text{-}trace}(\operatorname{sasp\_}, \operatorname{sas\_}) \vee \\ \operatorname{is\text{-}concat\text{-}trace}(\operatorname{sasp\_}, \operatorname{sas\_}) \vee \\ \operatorname{is\text{-}separate\text{-}trace}(\operatorname{sasp\_}, \operatorname{sas\_}) \vee \\ \operatorname{is\text{-}encrypt\text{-}trace}(\operatorname{sasp\_}, \operatorname{sas\_}) \vee \\ \operatorname{is\text{-}decrypt\text{-}trace}(\operatorname{sasp\_}, \operatorname{sas\_}) \end{array}
\end{equation*}

\par\medskip\noindent\textbf{Definition} (\texttt{is-adversary-node}).\;For all $\operatorname{sai\_}$ and $\operatorname{san\_}$, 
\begin{equation*}
\operatorname{is\text{-}adversary\text{-}node}(\operatorname{sai\_}, \operatorname{san\_}) \Leftrightarrow (\operatorname{nth}(2, \operatorname{san\_}) \in \operatorname{adv\text{-}ids}(\operatorname{sai\_}))
\end{equation*}

\par\medskip\noindent\textbf{Definition} (\texttt{is-regular-node}).\;For all $\operatorname{sai\_}$ and $\operatorname{san\_}$, 
\begin{equation*}
\operatorname{is\text{-}regular\text{-}node}(\operatorname{sai\_}, \operatorname{san\_}) \Leftrightarrow \neg (\operatorname{nth}(2, \operatorname{san\_}) \in \operatorname{adv\text{-}ids}(\operatorname{sai\_}))
\end{equation*}

\par\medskip\noindent\textbf{Definition} (\texttt{hom-infiltrated-space}).\;$\operatorname{hom\text{-}infiltrated\text{-}space}(a, b) :=$
\begin{equation*}
\{\operatorname{homp\_} \in ((\operatorname{adv\text{-}ids}(a) \to \operatorname{adv\text{-}ids}(b)) \times (\operatorname{adv\text{-}tokens}(a) \to \operatorname{adv\text{-}tokens}(b))) : \operatorname{is\text{-}hom\text{-}infiltrated\text{-}space}(a, b, \operatorname{nth}(1, \operatorname{homp\_}), \operatorname{nth}(2, \operatorname{homp\_}))\}
\end{equation*}

\vnbdef{infiltrated-space-class}{\vnbstmt{\forall s.\; ((s \in \operatorname{infiltrated\text{-}space}) \Leftrightarrow \operatorname{is\text{-}infiltrated\text{-}space}(s))}}
\vnbdef{is-hom-infiltrated-space-def}{\vnbstmt{\begin{array}{@{}l@{}} \forall a, b, c, d. \\ \quad \operatorname{is\text{-}hom\text{-}infiltrated\text{-}space}(a, b, c, d) \\ \quad \Leftrightarrow\; \operatorname{is\text{-}infiltrated\text{-}space}(a) \wedge \\ \quad \operatorname{is\text{-}infiltrated\text{-}space}(b) \wedge \\ \quad c \in (\operatorname{adv\text{-}ids}(a) \to \operatorname{adv\text{-}ids}(b)) \wedge \\ \quad d \in (\operatorname{adv\text{-}tokens}(a) \to \operatorname{adv\text{-}tokens}(b)) \wedge \\ \quad \operatorname{ispace}(a) = \operatorname{ispace}(b) \end{array}}}
\vnbprop{is-infiltrated-space-is-set}{\vnbstmt{\forall s.\; (\operatorname{is\text{-}infiltrated\text{-}space}(s) \Rightarrow (s \in \mathbf{Set}))}}

\vnbbookfile{theorem-library/strand-adversary-bound}{theorem-library/strand-adversary-bound}
\noindent{\small strand-adversary-bound.scm -- THE ADVERSARY BOUND (the paper's Proposition 3.3 on the user's model with encryption): a token the adversary does not know, and which originates on no regular node of a bundle, is a subterm of the content of NO node of the bundle. (Thayer, Herzog, Guttman, \textasciitilde{}/docs/jcs\_strand\_spaces.pdf, ...}\par
\vnbprop{text-trace-origin}{\vnbstmt{\begin{array}{@{}l@{}} \forall a, p, u, n. \\ \quad \operatorname{is\text{-}strand\text{-}space}(a) \wedge \\ \quad \forall x \in p.\; (x \in \operatorname{tokens}(a)) \wedge \\ \quad n \in \operatorname{nodes}(a) \wedge \\ \quad \operatorname{is\text{-}text\text{-}trace}(a, p, \operatorname{strand\text{-}map}(a)(\operatorname{nth}(2, n))) \wedge \\ \quad \operatorname{originates}(a, u, n) \Rightarrow \\ \quad \quad u \in p \end{array}}}
\vnbprop{flush-trace-no-sent}{\vnbstmt{\begin{array}{@{}l@{}} \forall a, n. \\ \quad (\operatorname{is\text{-}strand\text{-}space}(a) \wedge ((n \in \operatorname{nodes}(a)) \wedge \operatorname{is\text{-}flush\text{-}trace}(a, \operatorname{strand\text{-}map}(a)(\operatorname{nth}(2, n))))) \Rightarrow  \\ \quad \quad \neg \operatorname{is\text{-}sent}(a, n) \end{array}}}
\vnbprop{tee-trace-no-origin}{\vnbstmt{\begin{array}{@{}l@{}} \forall a, u, n. \\ \quad (\operatorname{is\text{-}strand\text{-}space}(a) \wedge ((n \in \operatorname{nodes}(a)) \wedge \operatorname{is\text{-}tee\text{-}trace}(a, \operatorname{strand\text{-}map}(a)(\operatorname{nth}(2, n))))) \Rightarrow  \\ \quad \quad \neg \operatorname{originates}(a, u, n) \end{array}}}
\vnbprop{concat-trace-no-origin}{\vnbstmt{\begin{array}{@{}l@{}} \forall a, u, n. \\ \quad \operatorname{is\text{-}strand\text{-}space}(a) \wedge \\ \quad u \in \operatorname{tokens}(a) \wedge \\ \quad n \in \operatorname{nodes}(a) \wedge \\ \quad \operatorname{is\text{-}concat\text{-}trace}(a, \operatorname{strand\text{-}map}(a)(\operatorname{nth}(2, n))) \Rightarrow \\ \quad \quad \neg \operatorname{originates}(a, u, n) \end{array}}}
\vnbprop{separate-trace-no-origin}{\vnbstmt{\begin{array}{@{}l@{}} \forall a, u, n. \\ \quad (\operatorname{is\text{-}strand\text{-}space}(a) \wedge ((n \in \operatorname{nodes}(a)) \wedge \operatorname{is\text{-}separate\text{-}trace}(a, \operatorname{strand\text{-}map}(a)(\operatorname{nth}(2, n))))) \Rightarrow  \\ \quad \quad \neg \operatorname{originates}(a, u, n) \end{array}}}
\vnbprop{encrypt-trace-no-origin}{\vnbstmt{\begin{array}{@{}l@{}} \forall a, u, n. \\ \quad \operatorname{is\text{-}strand\text{-}space}(a) \wedge \\ \quad u \in \operatorname{tokens}(a) \wedge \\ \quad n \in \operatorname{nodes}(a) \wedge \\ \quad \operatorname{is\text{-}encrypt\text{-}trace}(a, \operatorname{strand\text{-}map}(a)(\operatorname{nth}(2, n))) \Rightarrow \\ \quad \quad \neg \operatorname{originates}(a, u, n) \end{array}}}
\vnbprop{decrypt-trace-no-origin}{\vnbstmt{\begin{array}{@{}l@{}} \forall a, u, n. \\ \quad (\operatorname{is\text{-}strand\text{-}space}(a) \wedge ((n \in \operatorname{nodes}(a)) \wedge \operatorname{is\text{-}decrypt\text{-}trace}(a, \operatorname{strand\text{-}map}(a)(\operatorname{nth}(2, n))))) \Rightarrow  \\ \quad \quad \neg \operatorname{originates}(a, u, n) \end{array}}}
\vnbprop{adversary-node-origin}{\vnbstmt{\begin{array}{@{}l@{}} \forall p, u, n. \\ \quad \operatorname{is\text{-}infiltrated\text{-}space}(p) \wedge \\ \quad u \in \operatorname{tokens}(\operatorname{ispace}(p)) \wedge \\ \quad n \in \operatorname{nodes}(\operatorname{ispace}(p)) \wedge \\ \quad \operatorname{is\text{-}adversary\text{-}node}(p, n) \wedge \\ \quad \operatorname{originates}(\operatorname{ispace}(p), u, n) \Rightarrow \\ \quad \quad u \in \operatorname{adv\text{-}tokens}(p) \end{array}}}
\vnbprop{adversary-bound}{\vnbstmt{\begin{array}{@{}l@{}} \forall p, c, u.; \forall n \in \operatorname{bnodes}(c). \\ \quad \operatorname{is\text{-}infiltrated\text{-}space}(p) \wedge \\ \quad \operatorname{is\text{-}bundle}(c) \wedge \\ \quad \operatorname{bspace}(c) = \operatorname{ispace}(p) \wedge \\ \quad u \in \operatorname{tokens}(\operatorname{ispace}(p)) \wedge \\ \quad \neg (u \in \operatorname{adv\text{-}tokens}(p)) \wedge \\ \quad \forall n \in \operatorname{bnodes}(c). \\ \quad \quad \operatorname{is\text{-}regular\text{-}node}(p, n) \Rightarrow \neg \operatorname{originates}(\operatorname{ispace}(p), u, n) \Rightarrow \\ \quad \quad \neg \operatorname{is\text{-}subterm}(\operatorname{tokens}(\operatorname{ispace}(p)), \operatorname{keys}(\operatorname{ispace}(p)), u, \operatorname{node\text{-}content}(\operatorname{ispace}(p), n)) \end{array}}}
\vnbprop{adversary-bound-adversary-nodes}{\vnbstmt{\begin{array}{@{}l@{}} \forall p, c, u.; \forall n \in \operatorname{bnodes}(c). \\ \quad \operatorname{is\text{-}infiltrated\text{-}space}(p) \wedge \\ \quad \operatorname{is\text{-}bundle}(c) \wedge \\ \quad \operatorname{bspace}(c) = \operatorname{ispace}(p) \wedge \\ \quad u \in \operatorname{tokens}(\operatorname{ispace}(p)) \wedge \\ \quad \neg (u \in \operatorname{adv\text{-}tokens}(p)) \wedge \\ \quad \forall n \in \operatorname{bnodes}(c). \\ \quad \quad \operatorname{is\text{-}regular\text{-}node}(p, n) \Rightarrow \neg \operatorname{originates}(\operatorname{ispace}(p), u, n) \wedge \\ \quad \operatorname{is\text{-}adversary\text{-}node}(p, n) \Rightarrow \\ \quad \quad \neg \operatorname{is\text{-}subterm}(\operatorname{tokens}(\operatorname{ispace}(p)), \operatorname{keys}(\operatorname{ispace}(p)), u, \operatorname{node\text{-}content}(\operatorname{ispace}(p), n)) \end{array}}}
\vnbprop{sent-received-in-pair}{\vnbstmt{(\mathit{sent} \in \{\mathit{sent}, \mathit{received}\}) \wedge (\mathit{received} \in \{\mathit{sent}, \mathit{received}\})}}
\vnbprop{directed-intro}{\vnbstmt{\begin{array}{@{}l@{}} \forall t, k.; \forall g \in \operatorname{msg\text{-}set}(t, k). \\ \quad ({[}\mathit{sent}, g] \in \operatorname{directed}(t, k)) \wedge ({[}\mathit{received}, g] \in \operatorname{directed}(t, k)) \end{array}}}
\vnbprop{separate-trace-exists}{\vnbstmt{\begin{array}{@{}l@{}} \forall a, l. \\ \quad (\operatorname{is\text{-}strand\text{-}space}(a) \wedge (l \in \operatorname{tuples}(\operatorname{msg\text{-}set}(\operatorname{tokens}(a), \operatorname{keys}(a))))) \Rightarrow  \\ \quad \quad \exists s \in \operatorname{strands}(\operatorname{tokens}(a), \operatorname{keys}(a)).\; \operatorname{is\text{-}separate\text{-}trace}(a, s) \end{array}}}
\vnbprop{concat-trace-exists}{\vnbstmt{\begin{array}{@{}l@{}} \forall a, l. \\ \quad (\operatorname{is\text{-}strand\text{-}space}(a) \wedge (l \in \operatorname{tuples}(\operatorname{msg\text{-}set}(\operatorname{tokens}(a), \operatorname{keys}(a))))) \Rightarrow  \\ \quad \quad \exists s \in \operatorname{strands}(\operatorname{tokens}(a), \operatorname{keys}(a)).\; \operatorname{is\text{-}concat\text{-}trace}(a, s) \end{array}}}
\vnbprop{encrypt-trace-strand}{\vnbstmt{\begin{array}{@{}l@{}} \forall a, k, m. \\ \quad (\operatorname{is\text{-}strand\text{-}space}(a) \wedge ((k \in \operatorname{keys}(a)) \wedge (m \in \operatorname{msg\text{-}set}(\operatorname{tokens}(a), \operatorname{keys}(a))))) \Rightarrow  \\ \quad \quad {[}{[}\mathit{received}, k], {[}\mathit{received}, m], {[}\mathit{sent}, \operatorname{encr}(m, k)]] \in \operatorname{strands}(\operatorname{tokens}(a), \operatorname{keys}(a)) \wedge \\ \quad \quad \operatorname{is\text{-}encrypt\text{-}trace}(a, {[}{[}\mathit{received}, k], {[}\mathit{received}, m], {[}\mathit{sent}, \operatorname{encr}(m, k)]]) \end{array}}}
\vnbprop{decrypt-trace-strand}{\vnbstmt{\begin{array}{@{}l@{}} \forall a, k, m. \\ \quad (\operatorname{is\text{-}strand\text{-}space}(a) \wedge ((k \in \operatorname{keys}(a)) \wedge (m \in \operatorname{msg\text{-}set}(\operatorname{tokens}(a), \operatorname{keys}(a))))) \Rightarrow  \\ \quad \quad {[}{[}\mathit{received}, \operatorname{kinv}(a)(k)], {[}\mathit{received}, \operatorname{encr}(m, k)], {[}\mathit{sent}, m]] \in \operatorname{strands}(\operatorname{tokens}(a), \operatorname{keys}(a)) \wedge \\ \quad \quad \operatorname{is\text{-}decrypt\text{-}trace}(a, {[}{[}\mathit{received}, \operatorname{kinv}(a)(k)], {[}\mathit{received}, \operatorname{encr}(m, k)], {[}\mathit{sent}, m]]) \end{array}}}

\vnbbookfile{structure-library/nsl}{structure-library/nsl}
\noindent{\small nsl.scm -- THE NEEDHAM-SCHROEDER-LOWE PROTOCOL ON THE USER'S STRAND MODEL: names, public keys, the initiator's and the responder's strands, NSL spaces. (Thayer, Herzog, Guttman, \textasciitilde{}/docs/jcs\_strand\_spaces.pdf, section 5: the paper's Definition 5.1 and the term-algebra specialisation that precedes it; the model is ...}\par
\vnbstructdef{pk-space}{A pk space $s$ is a 3-tuple whose components are listed below.\par
\begin{enumerate}[nosep]
\item $\operatorname{pk\text{-}inf}(s) = \operatorname{nth}(1, s)$, a substructure (\texttt{substructure}): $\operatorname{is\text{-}infiltrated\text{-}space}(\operatorname{pk\text{-}inf}(s))$.
\item $\operatorname{names}(s) = \operatorname{nth}(2, s)$, a carrier (\texttt{carriers}): $\operatorname{names}(s) \in \mathbf{Set}$.
\item $\operatorname{pubkey}(s) = \operatorname{nth}(3, s)$, an operation (\texttt{op}): $\operatorname{pubkey}(s) \in (\operatorname{names}(s) \to \operatorname{keys}(\operatorname{ispace}(\operatorname{pk\text{-}inf}(s))))$.
\end{enumerate}
The laws (the \texttt{property} and \texttt{law} clauses):
\begin{itemize}[nosep]
\item $\forall \operatorname{nlx\_} \in \operatorname{names}(s).\; (\operatorname{nlx\_} \in \operatorname{tokens}(\operatorname{ispace}(\operatorname{pk\text{-}inf}(s))))$
\item $\forall \operatorname{nlx\_} \in \operatorname{names}(s).\; \neg (\operatorname{nlx\_} \in \operatorname{keys}(\operatorname{ispace}(\operatorname{pk\text{-}inf}(s))))$
\item
\begin{equation*}
\begin{array}{@{}l@{}} \forall \operatorname{nlx\_}, \operatorname{nly\_} \in \operatorname{names}(s). \\ \quad (\operatorname{pubkey}(s)(\operatorname{nlx\_}) = \operatorname{pubkey}(s)(\operatorname{nly\_})) \Rightarrow (\operatorname{nlx\_} = \operatorname{nly\_}) \end{array}
\end{equation*}

\end{itemize}
The characteristic law (the axiom \texttt{is-pk-space}): $\operatorname{is\text{-}pk\text{-}space}(s)$ holds if and only if $\operatorname{length}(s) = 3$, every component has the typing listed with it, every law holds, and $s \in \mathbf{Set}$ (the structure is a set).
}
\par\medskip\noindent\textbf{Definition} (\texttt{is-init-strand}).\;For all $\operatorname{nlp\_}$, $\operatorname{nls\_}$, $\operatorname{nla\_}$, $\operatorname{nlb\_}$, $\operatorname{nlm\_}$ and $\operatorname{nln\_}$, 
\begin{equation*}
\begin{array}{@{}l@{}} \operatorname{is\text{-}init\text{-}strand}(\operatorname{nlp\_}, \operatorname{nls\_}, \operatorname{nla\_}, \operatorname{nlb\_}, \operatorname{nlm\_}, \operatorname{nln\_}) \\ \Leftrightarrow\; \operatorname{nla\_} \in \operatorname{names}(\operatorname{nlp\_}) \wedge \\ \operatorname{nlb\_} \in \operatorname{names}(\operatorname{nlp\_}) \wedge \\ \operatorname{nlm\_} \in \operatorname{tokens}(\operatorname{ispace}(\operatorname{pk\text{-}inf}(\operatorname{nlp\_}))) \wedge \\ \operatorname{nln\_} \in \operatorname{tokens}(\operatorname{ispace}(\operatorname{pk\text{-}inf}(\operatorname{nlp\_}))) \wedge \\ \neg (\operatorname{nlm\_} \in \operatorname{names}(\operatorname{nlp\_})) \wedge \\ \neg (\operatorname{nlm\_} \in \operatorname{keys}(\operatorname{ispace}(\operatorname{pk\text{-}inf}(\operatorname{nlp\_})))) \wedge \\ \neg (\operatorname{nln\_} \in \operatorname{keys}(\operatorname{ispace}(\operatorname{pk\text{-}inf}(\operatorname{nlp\_})))) \wedge \\ \operatorname{nls\_} \\ =\; {[}{[}\mathit{sent}, \operatorname{encr}(\operatorname{lst}({[}\operatorname{nlm\_}, \operatorname{nla\_}]), \operatorname{pubkey}(\operatorname{nlp\_})(\operatorname{nlb\_}))], {[}\mathit{received}, \operatorname{encr}(\operatorname{lst}({[}\operatorname{nlm\_}, \operatorname{nln\_}, \operatorname{nlb\_}]), \operatorname{pubkey}(\operatorname{nlp\_})(\operatorname{nla\_}))], {[}\mathit{sent}, \operatorname{encr}(\operatorname{nln\_}, \operatorname{pubkey}(\operatorname{nlp\_})(\operatorname{nlb\_}))]] \end{array}
\end{equation*}

\par\medskip\noindent\textbf{Definition} (\texttt{is-resp-strand}).\;For all $\operatorname{nlp\_}$, $\operatorname{nls\_}$, $\operatorname{nla\_}$, $\operatorname{nlb\_}$, $\operatorname{nlm\_}$ and $\operatorname{nln\_}$, 
\begin{equation*}
\begin{array}{@{}l@{}} \operatorname{is\text{-}resp\text{-}strand}(\operatorname{nlp\_}, \operatorname{nls\_}, \operatorname{nla\_}, \operatorname{nlb\_}, \operatorname{nlm\_}, \operatorname{nln\_}) \\ \Leftrightarrow\; \operatorname{nla\_} \in \operatorname{names}(\operatorname{nlp\_}) \wedge \\ \operatorname{nlb\_} \in \operatorname{names}(\operatorname{nlp\_}) \wedge \\ \operatorname{nlm\_} \in \operatorname{tokens}(\operatorname{ispace}(\operatorname{pk\text{-}inf}(\operatorname{nlp\_}))) \wedge \\ \operatorname{nln\_} \in \operatorname{tokens}(\operatorname{ispace}(\operatorname{pk\text{-}inf}(\operatorname{nlp\_}))) \wedge \\ \neg (\operatorname{nln\_} \in \operatorname{names}(\operatorname{nlp\_})) \wedge \\ \neg (\operatorname{nlm\_} \in \operatorname{keys}(\operatorname{ispace}(\operatorname{pk\text{-}inf}(\operatorname{nlp\_})))) \wedge \\ \neg (\operatorname{nln\_} \in \operatorname{keys}(\operatorname{ispace}(\operatorname{pk\text{-}inf}(\operatorname{nlp\_})))) \wedge \\ \operatorname{nls\_} \\ =\; {[}{[}\mathit{received}, \operatorname{encr}(\operatorname{lst}({[}\operatorname{nlm\_}, \operatorname{nla\_}]), \operatorname{pubkey}(\operatorname{nlp\_})(\operatorname{nlb\_}))], {[}\mathit{sent}, \operatorname{encr}(\operatorname{lst}({[}\operatorname{nlm\_}, \operatorname{nln\_}, \operatorname{nlb\_}]), \operatorname{pubkey}(\operatorname{nlp\_})(\operatorname{nla\_}))], {[}\mathit{received}, \operatorname{encr}(\operatorname{nln\_}, \operatorname{pubkey}(\operatorname{nlp\_})(\operatorname{nlb\_}))]] \end{array}
\end{equation*}

\par\medskip\noindent\textbf{Definition} (\texttt{is-nsl-space}).\;For every $\operatorname{nlp\_}$, 
\begin{equation*}
\begin{array}{@{}l@{}} \operatorname{is\text{-}nsl\text{-}space}(\operatorname{nlp\_}) \\ \Leftrightarrow\; \operatorname{is\text{-}pk\text{-}space}(\operatorname{nlp\_}) \wedge \\ \forall \operatorname{nli\_} \in \operatorname{ids}(\operatorname{ispace}(\operatorname{pk\text{-}inf}(\operatorname{nlp\_}))). \\ \quad \operatorname{nli\_} \in \operatorname{adv\text{-}ids}(\operatorname{pk\text{-}inf}(\operatorname{nlp\_})) \vee \\ \quad \exists \operatorname{nla\_}, \operatorname{nlb\_}, \operatorname{nlm\_}, \operatorname{nln\_}. \\ \quad \quad \operatorname{is\text{-}init\text{-}strand}(\operatorname{nlp\_}, \operatorname{strand\text{-}map}(\operatorname{ispace}(\operatorname{pk\text{-}inf}(\operatorname{nlp\_})))(\operatorname{nli\_}), \operatorname{nla\_}, \operatorname{nlb\_}, \operatorname{nlm\_}, \operatorname{nln\_}) \vee \\ \quad \exists \operatorname{nla\_}, \operatorname{nlb\_}, \operatorname{nlm\_}, \operatorname{nln\_}. \\ \quad \quad \operatorname{is\text{-}resp\text{-}strand}(\operatorname{nlp\_}, \operatorname{strand\text{-}map}(\operatorname{ispace}(\operatorname{pk\text{-}inf}(\operatorname{nlp\_})))(\operatorname{nli\_}), \operatorname{nla\_}, \operatorname{nlb\_}, \operatorname{nlm\_}, \operatorname{nln\_}) \end{array}
\end{equation*}

\par\medskip\noindent\textbf{Definition} (\texttt{hom-pk-space}).\;$\operatorname{hom\text{-}pk\text{-}space}(a, b) :=$
\begin{equation*}
\{\operatorname{homf\_} \in (\operatorname{names}(a) \to \operatorname{names}(b)) : \operatorname{is\text{-}hom\text{-}pk\text{-}space}(a, b, \operatorname{homf\_})\}
\end{equation*}

\vnbdef{is-hom-pk-space-def}{\vnbstmt{\begin{array}{@{}l@{}} \forall a, b, f. \\ \quad \operatorname{is\text{-}hom\text{-}pk\text{-}space}(a, b, f) \\ \quad \Leftrightarrow\; \operatorname{is\text{-}pk\text{-}space}(a) \wedge \\ \quad \operatorname{is\text{-}pk\text{-}space}(b) \wedge \\ \quad f \in (\operatorname{names}(a) \to \operatorname{names}(b)) \wedge \\ \quad \operatorname{pk\text{-}inf}(a) = \operatorname{pk\text{-}inf}(b) \wedge \\ \quad \forall x_{1} \in \operatorname{names}(a).\; (\operatorname{pubkey}(a)(x_{1}) = \operatorname{pubkey}(b)(f(x_{1}))) \end{array}}}
\vnbdef{pk-space-class}{\vnbstmt{\forall s.\; ((s \in \operatorname{pk\text{-}space}) \Leftrightarrow \operatorname{is\text{-}pk\text{-}space}(s))}}
\vnbprop{is-pk-space-is-set}{\vnbstmt{\forall s.\; (\operatorname{is\text{-}pk\text{-}space}(s) \Rightarrow (s \in \mathbf{Set}))}}

\vnbbookfile{theorem-library/nsl}{theorem-library/nsl}
\noindent{\small nsl.scm -- THE NEEDHAM-SCHROEDER-LOWE PROTOCOL ON THE USER'S STRAND MODEL: the responder's guarantee (agreement, Proposition 5.2) and the secrecy of both nonces, the injectivity of the initiator's strand, the initiator's agreement (Propositions 5.8, 5.10, 5.13, 5.14). (Thayer, Herzog, Guttman, ...}\par
\vnbprop{pk-space-parts}{\vnbstmt{\begin{array}{@{}l@{}} \forall p. \\ \quad \operatorname{is\text{-}pk\text{-}space}(p) \Rightarrow  \\ \quad \quad \operatorname{is\text{-}infiltrated\text{-}space}(\operatorname{pk\text{-}inf}(p)) \wedge \\ \quad \quad \operatorname{is\text{-}strand\text{-}space}(\operatorname{ispace}(\operatorname{pk\text{-}inf}(p))) \wedge \\ \quad \quad \operatorname{names}(p) \in \mathbf{Set} \wedge \\ \quad \quad \operatorname{pubkey}(p) \in (\operatorname{names}(p) \to \operatorname{keys}(\operatorname{ispace}(\operatorname{pk\text{-}inf}(p)))) \wedge \\ \quad \quad \forall x \in \operatorname{names}(p).\; (x \in \operatorname{tokens}(\operatorname{ispace}(\operatorname{pk\text{-}inf}(p)))) \wedge \\ \quad \quad \forall x \in \operatorname{names}(p).\; \neg (x \in \operatorname{keys}(\operatorname{ispace}(\operatorname{pk\text{-}inf}(p)))) \wedge \\ \quad \quad \forall x, y. \\ \quad \quad \quad ((x \in \operatorname{names}(p)) \wedge ((y \in \operatorname{names}(p)) \wedge (\operatorname{pubkey}(p)(x) = \operatorname{pubkey}(p)(y)))) \Rightarrow  \\ \quad \quad \quad \quad x = y \end{array}}}
\vnbprop{three-strand-node}{\vnbstmt{\begin{array}{@{}l@{}} \forall a, n, b, c, d. \\ \quad ((n \in \operatorname{nodes}(a)) \wedge (\operatorname{strand\text{-}map}(a)(\operatorname{nth}(2, n)) = {[}b, c, d])) \Rightarrow  \\ \quad \quad ((\operatorname{nth}(1, n) = 1) \wedge ((\operatorname{node\text{-}content}(a, n) \simeq \operatorname{nth}(2, b)) \wedge (\operatorname{node\text{-}dir}(a, n) \simeq \operatorname{nth}(1, b)))) \vee \\ \quad \quad ((\operatorname{nth}(1, n) = 2) \wedge ((\operatorname{node\text{-}content}(a, n) \simeq \operatorname{nth}(2, c)) \wedge (\operatorname{node\text{-}dir}(a, n) \simeq \operatorname{nth}(1, c)))) \vee \\ \quad \quad ((\operatorname{nth}(1, n) = 3) \wedge ((\operatorname{node\text{-}content}(a, n) \simeq \operatorname{nth}(2, d)) \wedge (\operatorname{node\text{-}dir}(a, n) \simeq \operatorname{nth}(1, d)))) \end{array}}}
\vnbprop{three-strand-nodes}{\vnbstmt{\begin{array}{@{}l@{}} \forall b, a, c, d, f. \\ \quad ((a \in \operatorname{ids}(b)) \wedge (\operatorname{strand\text{-}map}(b)(a) = {[}c, d, f])) \Rightarrow  \\ \quad \quad {[}1, a] \in \operatorname{nodes}(b) \wedge \\ \quad \quad (\operatorname{node\text{-}content}(b, {[}1, a]) \simeq \operatorname{nth}(2, c)) \wedge \\ \quad \quad (\operatorname{node\text{-}dir}(b, {[}1, a]) \simeq \operatorname{nth}(1, c)) \wedge \\ \quad \quad {[}2, a] \in \operatorname{nodes}(b) \wedge \\ \quad \quad (\operatorname{node\text{-}content}(b, {[}2, a]) \simeq \operatorname{nth}(2, d)) \wedge \\ \quad \quad (\operatorname{node\text{-}dir}(b, {[}2, a]) \simeq \operatorname{nth}(1, d)) \wedge \\ \quad \quad {[}3, a] \in \operatorname{nodes}(b) \wedge \\ \quad \quad (\operatorname{node\text{-}content}(b, {[}3, a]) \simeq \operatorname{nth}(2, f)) \wedge \\ \quad \quad (\operatorname{node\text{-}dir}(b, {[}3, a]) \simeq \operatorname{nth}(1, f)) \end{array}}}
\vnbprop{init-strand-is-strand}{\vnbstmt{\begin{array}{@{}l@{}} \forall p, s, a, b, c, d. \\ \quad (\operatorname{is\text{-}pk\text{-}space}(p) \wedge \operatorname{is\text{-}init\text{-}strand}(p, s, a, b, c, d)) \Rightarrow  \\ \quad \quad s \in \operatorname{strands}(\operatorname{tokens}(\operatorname{ispace}(\operatorname{pk\text{-}inf}(p))), \operatorname{keys}(\operatorname{ispace}(\operatorname{pk\text{-}inf}(p)))) \end{array}}}
\vnbprop{resp-strand-is-strand}{\vnbstmt{\begin{array}{@{}l@{}} \forall p, s, a, b, c, d. \\ \quad (\operatorname{is\text{-}pk\text{-}space}(p) \wedge \operatorname{is\text{-}resp\text{-}strand}(p, s, a, b, c, d)) \Rightarrow  \\ \quad \quad s \in \operatorname{strands}(\operatorname{tokens}(\operatorname{ispace}(\operatorname{pk\text{-}inf}(p))), \operatorname{keys}(\operatorname{ispace}(\operatorname{pk\text{-}inf}(p)))) \end{array}}}
\vnbprop{init-strand-nodes}{\vnbstmt{\begin{array}{@{}l@{}} \forall p, a, b, c, d, f. \\ \quad \operatorname{is\text{-}pk\text{-}space}(p) \wedge \\ \quad a \in \operatorname{ids}(\operatorname{ispace}(\operatorname{pk\text{-}inf}(p))) \wedge \\ \quad \operatorname{is\text{-}init\text{-}strand}(p, \operatorname{strand\text{-}map}(\operatorname{ispace}(\operatorname{pk\text{-}inf}(p)))(a), b, c, d, f) \Rightarrow \\ \quad \quad {[}1, a] \in \operatorname{nodes}(\operatorname{ispace}(\operatorname{pk\text{-}inf}(p))) \wedge \\ \quad \quad \operatorname{is\text{-}sent}(\operatorname{ispace}(\operatorname{pk\text{-}inf}(p)), {[}1, a]) \wedge \\ \quad \quad \operatorname{node\text{-}content}(\operatorname{ispace}(\operatorname{pk\text{-}inf}(p)), {[}1, a]) = \operatorname{encr}(\operatorname{lst}({[}d, b]), \operatorname{pubkey}(p)(c)) \wedge \\ \quad \quad {[}2, a] \in \operatorname{nodes}(\operatorname{ispace}(\operatorname{pk\text{-}inf}(p))) \wedge \\ \quad \quad \operatorname{is\text{-}received}(\operatorname{ispace}(\operatorname{pk\text{-}inf}(p)), {[}2, a]) \wedge \\ \quad \quad \operatorname{node\text{-}content}(\operatorname{ispace}(\operatorname{pk\text{-}inf}(p)), {[}2, a]) = \operatorname{encr}(\operatorname{lst}({[}d, f, c]), \operatorname{pubkey}(p)(b)) \wedge \\ \quad \quad {[}3, a] \in \operatorname{nodes}(\operatorname{ispace}(\operatorname{pk\text{-}inf}(p))) \wedge \\ \quad \quad \operatorname{is\text{-}sent}(\operatorname{ispace}(\operatorname{pk\text{-}inf}(p)), {[}3, a]) \wedge \\ \quad \quad \operatorname{node\text{-}content}(\operatorname{ispace}(\operatorname{pk\text{-}inf}(p)), {[}3, a]) = \operatorname{encr}(f, \operatorname{pubkey}(p)(c)) \end{array}}}
\vnbprop{resp-strand-nodes}{\vnbstmt{\begin{array}{@{}l@{}} \forall p, a, b, c, d, f. \\ \quad \operatorname{is\text{-}pk\text{-}space}(p) \wedge \\ \quad a \in \operatorname{ids}(\operatorname{ispace}(\operatorname{pk\text{-}inf}(p))) \wedge \\ \quad \operatorname{is\text{-}resp\text{-}strand}(p, \operatorname{strand\text{-}map}(\operatorname{ispace}(\operatorname{pk\text{-}inf}(p)))(a), b, c, d, f) \Rightarrow \\ \quad \quad {[}1, a] \in \operatorname{nodes}(\operatorname{ispace}(\operatorname{pk\text{-}inf}(p))) \wedge \\ \quad \quad \operatorname{is\text{-}received}(\operatorname{ispace}(\operatorname{pk\text{-}inf}(p)), {[}1, a]) \wedge \\ \quad \quad \operatorname{node\text{-}content}(\operatorname{ispace}(\operatorname{pk\text{-}inf}(p)), {[}1, a]) = \operatorname{encr}(\operatorname{lst}({[}d, b]), \operatorname{pubkey}(p)(c)) \wedge \\ \quad \quad {[}2, a] \in \operatorname{nodes}(\operatorname{ispace}(\operatorname{pk\text{-}inf}(p))) \wedge \\ \quad \quad \operatorname{is\text{-}sent}(\operatorname{ispace}(\operatorname{pk\text{-}inf}(p)), {[}2, a]) \wedge \\ \quad \quad \operatorname{node\text{-}content}(\operatorname{ispace}(\operatorname{pk\text{-}inf}(p)), {[}2, a]) = \operatorname{encr}(\operatorname{lst}({[}d, f, c]), \operatorname{pubkey}(p)(b)) \wedge \\ \quad \quad {[}3, a] \in \operatorname{nodes}(\operatorname{ispace}(\operatorname{pk\text{-}inf}(p))) \wedge \\ \quad \quad \operatorname{is\text{-}received}(\operatorname{ispace}(\operatorname{pk\text{-}inf}(p)), {[}3, a]) \wedge \\ \quad \quad \operatorname{node\text{-}content}(\operatorname{ispace}(\operatorname{pk\text{-}inf}(p)), {[}3, a]) = \operatorname{encr}(f, \operatorname{pubkey}(p)(c)) \end{array}}}
\vnbprop{token-not-encr}{\vnbstmt{\begin{array}{@{}l@{}} \forall t, g, y. \\ \quad (\operatorname{is\text{-}token\text{-}set}(t) \wedge ((g \in \mathbf{Set}) \wedge (y \in \mathbf{Set}))) \Rightarrow \neg (\operatorname{encr}(g, y) \in t) \end{array}}}
\vnbprop{token-not-lst}{\vnbstmt{\forall t, l.\; ((\operatorname{is\text{-}token\text{-}set}(t) \wedge (l \in \mathbf{Set})) \Rightarrow \neg (\operatorname{lst}(l) \in t))}}
\vnbprop{subterm-of-flat-list}{\vnbstmt{\begin{array}{@{}l@{}} \forall t, k, d, w, l. \\ \quad \operatorname{is\text{-}token\text{-}set}(t) \wedge \\ \quad k \in \mathbf{Set} \wedge \\ \quad \forall x \in d.\; (x \in t) \wedge \\ \quad l \in \operatorname{tuples}(d) \wedge \\ \quad \operatorname{is\text{-}subterm}(t, k, w, \operatorname{lst}(l)) \Rightarrow \\ \quad \quad (w = \operatorname{lst}(l)) \vee (w \in d) \end{array}}}
\vnbprop{subterm-of-flat-encr-list}{\vnbstmt{\begin{array}{@{}l@{}} \forall t, k, d, w, l, y. \\ \quad \operatorname{is\text{-}token\text{-}set}(t) \wedge \\ \quad k \in \mathbf{Set} \wedge \\ \quad \forall x \in d.\; (x \in t) \wedge \\ \quad l \in \operatorname{tuples}(d) \wedge \\ \quad y \in k \wedge \\ \quad \operatorname{is\text{-}subterm}(t, k, w, \operatorname{encr}(\operatorname{lst}(l), y)) \Rightarrow \\ \quad \quad (w = \operatorname{encr}(\operatorname{lst}(l), y)) \vee ((w = \operatorname{lst}(l)) \vee (w \in d)) \end{array}}}
\vnbprop{subterm-of-flat-encr-token}{\vnbstmt{\begin{array}{@{}l@{}} \forall t, k, d, w, x, y. \\ \quad \operatorname{is\text{-}token\text{-}set}(t) \wedge \\ \quad k \in \mathbf{Set} \wedge \\ \quad \forall z \in d.\; (z \in t) \wedge \\ \quad x \in d \wedge \\ \quad y \in k \wedge \\ \quad \operatorname{is\text{-}subterm}(t, k, w, \operatorname{encr}(x, y)) \Rightarrow \\ \quad \quad (w = \operatorname{encr}(x, y)) \vee (w \in d) \end{array}}}
\vnbprop{nsl-regular-content}{\vnbstmt{\begin{array}{@{}l@{}} \forall p, n. \\ \quad (\operatorname{is\text{-}nsl\text{-}space}(p) \wedge ((n \in \operatorname{nodes}(\operatorname{ispace}(\operatorname{pk\text{-}inf}(p)))) \wedge \operatorname{is\text{-}regular\text{-}node}(\operatorname{pk\text{-}inf}(p), n))) \Rightarrow  \\ \quad \quad \exists k \in \operatorname{keys}(\operatorname{ispace}(\operatorname{pk\text{-}inf}(p))). \\ \quad \quad \quad \exists x \in \{d \in \operatorname{tokens}(\operatorname{ispace}(\operatorname{pk\text{-}inf}(p))) : \neg (d \in \operatorname{keys}(\operatorname{ispace}(\operatorname{pk\text{-}inf}(p))))\}. \\ \quad \quad \quad \quad \operatorname{node\text{-}content}(\operatorname{ispace}(\operatorname{pk\text{-}inf}(p)), n) = \operatorname{encr}(x, k) \vee \\ \quad \quad \quad \exists l \in \operatorname{tuples}(\{d \in \operatorname{tokens}(\operatorname{ispace}(\operatorname{pk\text{-}inf}(p))) : \neg (d \in \operatorname{keys}(\operatorname{ispace}(\operatorname{pk\text{-}inf}(p))))\}). \\ \quad \quad \quad \quad \operatorname{node\text{-}content}(\operatorname{ispace}(\operatorname{pk\text{-}inf}(p)), n) = \operatorname{encr}(\operatorname{lst}(l), k) \end{array}}}
\vnbprop{nsl-regular-encr-subterm}{\vnbstmt{\begin{array}{@{}l@{}} \forall p, n, g, y. \\ \quad \operatorname{is\text{-}nsl\text{-}space}(p) \wedge \\ \quad n \in \operatorname{nodes}(\operatorname{ispace}(\operatorname{pk\text{-}inf}(p))) \wedge \\ \quad \operatorname{is\text{-}regular\text{-}node}(\operatorname{pk\text{-}inf}(p), n) \wedge \\ \quad g \in \operatorname{msg\text{-}set}(\operatorname{tokens}(\operatorname{ispace}(\operatorname{pk\text{-}inf}(p))), \operatorname{keys}(\operatorname{ispace}(\operatorname{pk\text{-}inf}(p)))) \wedge \\ \quad y \in \operatorname{keys}(\operatorname{ispace}(\operatorname{pk\text{-}inf}(p))) \wedge \\ \quad \operatorname{is\text{-}subterm}(\operatorname{tokens}(\operatorname{ispace}(\operatorname{pk\text{-}inf}(p))), \operatorname{keys}(\operatorname{ispace}(\operatorname{pk\text{-}inf}(p))), \operatorname{encr}(g, y), \\ \quad \quad \operatorname{node\text{-}content}(\operatorname{ispace}(\operatorname{pk\text{-}inf}(p)), n)) \Rightarrow \\ \quad \quad \operatorname{encr}(g, y) = \operatorname{node\text{-}content}(\operatorname{ispace}(\operatorname{pk\text{-}inf}(p)), n) \end{array}}}
\vnbprop{nsl-regular-list-subterm}{\vnbstmt{\begin{array}{@{}l@{}} \forall p, n, m. \\ \quad \operatorname{is\text{-}nsl\text{-}space}(p) \wedge \\ \quad n \in \operatorname{nodes}(\operatorname{ispace}(\operatorname{pk\text{-}inf}(p))) \wedge \\ \quad \operatorname{is\text{-}regular\text{-}node}(\operatorname{pk\text{-}inf}(p), n) \wedge \\ \quad m \in \operatorname{tuples}(\operatorname{msg\text{-}set}(\operatorname{tokens}(\operatorname{ispace}(\operatorname{pk\text{-}inf}(p))), \operatorname{keys}(\operatorname{ispace}(\operatorname{pk\text{-}inf}(p))))) \wedge \\ \quad \operatorname{is\text{-}subterm}(\operatorname{tokens}(\operatorname{ispace}(\operatorname{pk\text{-}inf}(p))), \operatorname{keys}(\operatorname{ispace}(\operatorname{pk\text{-}inf}(p))), \operatorname{lst}(m), \\ \quad \quad \operatorname{node\text{-}content}(\operatorname{ispace}(\operatorname{pk\text{-}inf}(p)), n)) \Rightarrow \\ \quad \quad m \in \operatorname{tuples}(\operatorname{tokens}(\operatorname{ispace}(\operatorname{pk\text{-}inf}(p)))) \end{array}}}
\vnbprop{nsl-regular-no-key}{\vnbstmt{\begin{array}{@{}l@{}} \forall p, n, y. \\ \quad \operatorname{is\text{-}nsl\text{-}space}(p) \wedge \\ \quad n \in \operatorname{nodes}(\operatorname{ispace}(\operatorname{pk\text{-}inf}(p))) \wedge \\ \quad \operatorname{is\text{-}regular\text{-}node}(\operatorname{pk\text{-}inf}(p), n) \wedge \\ \quad y \in \operatorname{keys}(\operatorname{ispace}(\operatorname{pk\text{-}inf}(p))) \Rightarrow \\ \quad \quad \neg \operatorname{is\text{-}subterm}(\operatorname{tokens}(\operatorname{ispace}(\operatorname{pk\text{-}inf}(p))), \operatorname{keys}(\operatorname{ispace}(\operatorname{pk\text{-}inf}(p))), y, \operatorname{node\text{-}content}(\operatorname{ispace}(\operatorname{pk\text{-}inf}(p)), n)) \end{array}}}
\vnbprop{adversary-encr-loss}{\vnbstmt{\begin{array}{@{}l@{}} \forall p, m, x, g, y. \\ \quad \operatorname{is\text{-}infiltrated\text{-}space}(p) \wedge \\ \quad m \in \operatorname{nodes}(\operatorname{ispace}(p)) \wedge \\ \quad \operatorname{is\text{-}adversary\text{-}node}(p, m) \wedge \\ \quad \operatorname{is\text{-}sent}(\operatorname{ispace}(p), m) \wedge \\ \quad \operatorname{same\text{-}strand\text{-}before}(\operatorname{ispace}(p), x, m) \wedge \\ \quad g \in \operatorname{msg\text{-}set}(\operatorname{tokens}(\operatorname{ispace}(p)), \operatorname{keys}(\operatorname{ispace}(p))) \wedge \\ \quad y \in \operatorname{keys}(\operatorname{ispace}(p)) \wedge \\ \quad \operatorname{is\text{-}subterm}(\operatorname{tokens}(\operatorname{ispace}(p)), \operatorname{keys}(\operatorname{ispace}(p)), \operatorname{encr}(g, y), \operatorname{node\text{-}content}(\operatorname{ispace}(p), x)) \wedge \\ \quad \neg \operatorname{is\text{-}subterm}(\operatorname{tokens}(\operatorname{ispace}(p)), \operatorname{keys}(\operatorname{ispace}(p)), \operatorname{encr}(g, y), \operatorname{node\text{-}content}(\operatorname{ispace}(p), m)) \Rightarrow \\ \quad \quad \operatorname{same\text{-}strand\text{-}before}(\operatorname{ispace}(p), {[}1, \operatorname{nth}(2, m)], m) \wedge \\ \quad \quad \exists l \in \operatorname{tuples}(\operatorname{msg\text{-}set}(\operatorname{tokens}(\operatorname{ispace}(p)), \operatorname{keys}(\operatorname{ispace}(p)))). \\ \quad \quad \quad \operatorname{node\text{-}content}(\operatorname{ispace}(p), {[}1, \operatorname{nth}(2, m)]) = \operatorname{lst}(l) \wedge \\ \quad \quad \quad \operatorname{is\text{-}subterm}(\operatorname{tokens}(\operatorname{ispace}(p)), \operatorname{keys}(\operatorname{ispace}(p)), \operatorname{encr}(g, y), \operatorname{lst}(l)) \wedge \\ \quad \quad \quad \exists q \in \mathbb{N}. \\ \quad \quad \quad \quad (1 \leq q) \wedge ((q \leq \operatorname{length}(l)) \wedge (\operatorname{node\text{-}content}(\operatorname{ispace}(p), m) = \operatorname{nth}(q, l))) \vee \\ \quad \quad \operatorname{node\text{-}content}(\operatorname{ispace}(p), {[}1, \operatorname{nth}(2, m)]) = \operatorname{kinv}(\operatorname{ispace}(p))(y) \end{array}}}
\vnbprop{adversary-sent-source}{\vnbstmt{\begin{array}{@{}l@{}} \forall p, m, w. \\ \quad \operatorname{is\text{-}infiltrated\text{-}space}(p) \wedge \\ \quad m \in \operatorname{nodes}(\operatorname{ispace}(p)) \wedge \\ \quad \operatorname{is\text{-}adversary\text{-}node}(p, m) \wedge \\ \quad \operatorname{is\text{-}sent}(\operatorname{ispace}(p), m) \wedge \\ \quad \operatorname{is\text{-}subterm}(\operatorname{tokens}(\operatorname{ispace}(p)), \operatorname{keys}(\operatorname{ispace}(p)), w, \operatorname{node\text{-}content}(\operatorname{ispace}(p), m)) \Rightarrow \\ \quad \quad \exists x. \\ \quad \quad \quad \operatorname{same\text{-}strand\text{-}before}(\operatorname{ispace}(p), x, m) \wedge \\ \quad \quad \quad \operatorname{is\text{-}subterm}(\operatorname{tokens}(\operatorname{ispace}(p)), \operatorname{keys}(\operatorname{ispace}(p)), w, \operatorname{node\text{-}content}(\operatorname{ispace}(p), x)) \vee \\ \quad \quad w = \operatorname{node\text{-}content}(\operatorname{ispace}(p), m) \wedge \\ \quad \quad w \in \operatorname{adv\text{-}tokens}(p) \vee \\ \quad \quad \exists l \in \operatorname{tuples}(\operatorname{msg\text{-}set}(\operatorname{tokens}(\operatorname{ispace}(p)), \operatorname{keys}(\operatorname{ispace}(p)))). \\ \quad \quad \quad \operatorname{node\text{-}content}(\operatorname{ispace}(p), m) = \operatorname{lst}(l) \wedge \\ \quad \quad \quad \forall i \in \mathbb{N}. \\ \quad \quad \quad \quad ((1 \leq i) \wedge (i \leq \operatorname{length}(l))) \Rightarrow  \\ \quad \quad \quad \quad \quad \exists x. \\ \quad \quad \quad \quad \quad \quad \operatorname{same\text{-}strand\text{-}before}(\operatorname{ispace}(p), x, m) \wedge (\operatorname{node\text{-}content}(\operatorname{ispace}(p), x) = \operatorname{nth}(i, l)) \vee \\ \quad \quad \exists h.; \exists k \in \operatorname{keys}(\operatorname{ispace}(p)). \\ \quad \quad \quad \operatorname{node\text{-}content}(\operatorname{ispace}(p), m) = \operatorname{encr}(h, k) \wedge \\ \quad \quad \quad \exists x. \\ \quad \quad \quad \quad \operatorname{same\text{-}strand\text{-}before}(\operatorname{ispace}(p), x, m) \wedge (\operatorname{node\text{-}content}(\operatorname{ispace}(p), x) = h) \end{array}}}
\vnbprop{bundle-send-content}{\vnbstmt{\begin{array}{@{}l@{}} \forall c, d, a, b. \\ \quad (\operatorname{is\text{-}bundle}(c) \wedge ((\operatorname{bspace}(c) = d) \wedge ({[}a, b] \in \operatorname{bsend}(c)))) \Rightarrow  \\ \quad \quad (a \in \operatorname{bnodes}(c)) \wedge ((b \in \operatorname{bnodes}(c)) \wedge (\operatorname{node\text{-}content}(d, a) = \operatorname{node\text{-}content}(d, b))) \end{array}}}
\vnbprop{bundle-minimal-sent-back}{\vnbstmt{\begin{array}{@{}l@{}} \forall c, s, n. \\ \quad \operatorname{is\text{-}bundle}(c) \wedge \\ \quad \forall z \in s.\; (z \in \operatorname{bnodes}(c)) \wedge \\ \quad \forall x, y.\; ((({[}x, y] \in \operatorname{bsend}(c)) \wedge (y \in s)) \Rightarrow (x \in s)) \wedge \\ \quad n \in s \wedge \\ \quad \forall z \in s.\; \neg ({[}z, n] \in \operatorname{precedes\text{-}strict}(c)) \Rightarrow \\ \quad \quad \operatorname{is\text{-}sent}(\operatorname{bspace}(c), n) \end{array}}}
\vnbprop{bundle-minimal-earlier}{\vnbstmt{\begin{array}{@{}l@{}} \forall c, a, s, u, b, n. \\ \quad \operatorname{is\text{-}bundle}(c) \wedge \\ \quad \operatorname{bspace}(c) = a \wedge \\ \quad \forall z \in s.\; (z \in \operatorname{bnodes}(c)) \wedge \\ \quad \operatorname{uniquely\text{-}originates}(a, u) \wedge \\ \quad \operatorname{originates}(a, u, b) \wedge \\ \quad \neg (b \in s) \wedge \\ \quad n \in s \wedge \\ \quad \operatorname{is\text{-}sent}(a, n) \wedge \\ \quad \operatorname{is\text{-}subterm}(\operatorname{tokens}(a), \operatorname{keys}(a), u, \operatorname{node\text{-}content}(a, n)) \Rightarrow \\ \quad \quad \exists x \in \operatorname{bnodes}(c). \\ \quad \quad \quad {[}x, n] \in \operatorname{precedes\text{-}strict}(c) \wedge \\ \quad \quad \quad \operatorname{same\text{-}strand\text{-}before}(a, x, n) \wedge \\ \quad \quad \quad \operatorname{is\text{-}subterm}(\operatorname{tokens}(a), \operatorname{keys}(a), u, \operatorname{node\text{-}content}(a, x)) \end{array}}}
\vnbprop{nsl-key-secret}{\vnbstmt{\begin{array}{@{}l@{}} \forall p, c, y.; \forall n \in \operatorname{bnodes}(c). \\ \quad \operatorname{is\text{-}nsl\text{-}space}(p) \wedge \\ \quad \operatorname{is\text{-}bundle}(c) \wedge \\ \quad \operatorname{bspace}(c) = \operatorname{ispace}(\operatorname{pk\text{-}inf}(p)) \wedge \\ \quad y \in \operatorname{keys}(\operatorname{ispace}(\operatorname{pk\text{-}inf}(p))) \wedge \\ \quad \neg (y \in \operatorname{adv\text{-}tokens}(\operatorname{pk\text{-}inf}(p))) \Rightarrow \\ \quad \quad \neg \operatorname{is\text{-}subterm}(\operatorname{tokens}(\operatorname{ispace}(\operatorname{pk\text{-}inf}(p))), \operatorname{keys}(\operatorname{ispace}(\operatorname{pk\text{-}inf}(p))), y, \operatorname{node\text{-}content}(\operatorname{ispace}(\operatorname{pk\text{-}inf}(p)), n)) \end{array}}}
\vnbprop{nsl-list-source}{\vnbstmt{\begin{array}{@{}l@{}} \forall p, c, a, l, q. \\ \quad \operatorname{is\text{-}nsl\text{-}space}(p) \wedge \\ \quad \operatorname{is\text{-}bundle}(c) \wedge \\ \quad \operatorname{bspace}(c) = \operatorname{ispace}(\operatorname{pk\text{-}inf}(p)) \wedge \\ \quad a \in \operatorname{bnodes}(c) \wedge \\ \quad l \in \operatorname{tuples}(\operatorname{msg\text{-}set}(\operatorname{tokens}(\operatorname{ispace}(\operatorname{pk\text{-}inf}(p))), \operatorname{keys}(\operatorname{ispace}(\operatorname{pk\text{-}inf}(p))))) \wedge \\ \quad q \in \mathbb{N} \wedge \\ \quad 1 \leq q \wedge \\ \quad q \leq \operatorname{length}(l) \wedge \\ \quad \operatorname{node\text{-}content}(\operatorname{ispace}(\operatorname{pk\text{-}inf}(p)), a) = \operatorname{nth}(q, l) \wedge \\ \quad \neg (l \in \operatorname{tuples}(\operatorname{tokens}(\operatorname{ispace}(\operatorname{pk\text{-}inf}(p))))) \wedge \\ \quad \exists x \in \operatorname{bnodes}(c). \\ \quad \quad {[}x, a] \in \operatorname{precedes\text{-}strict}(c) \wedge \\ \quad \quad \operatorname{is\text{-}subterm}(\operatorname{tokens}(\operatorname{ispace}(\operatorname{pk\text{-}inf}(p))), \operatorname{keys}(\operatorname{ispace}(\operatorname{pk\text{-}inf}(p))), \operatorname{lst}(l), \\ \quad \quad \quad \operatorname{node\text{-}content}(\operatorname{ispace}(\operatorname{pk\text{-}inf}(p)), x)) \Rightarrow \\ \quad \quad \exists z \in \operatorname{bnodes}(c). \\ \quad \quad \quad {[}z, a] \in \operatorname{precedes\text{-}strict}(c) \wedge \\ \quad \quad \quad \operatorname{node\text{-}content}(\operatorname{ispace}(\operatorname{pk\text{-}inf}(p)), z) = \operatorname{node\text{-}content}(\operatorname{ispace}(\operatorname{pk\text{-}inf}(p)), a) \end{array}}}
\vnbprop{nsl-adversary-loss}{\vnbstmt{\begin{array}{@{}l@{}} \forall p, c, m, x, g, y. \\ \quad \operatorname{is\text{-}nsl\text{-}space}(p) \wedge \\ \quad \operatorname{is\text{-}bundle}(c) \wedge \\ \quad \operatorname{bspace}(c) = \operatorname{ispace}(\operatorname{pk\text{-}inf}(p)) \wedge \\ \quad m \in \operatorname{bnodes}(c) \wedge \\ \quad \operatorname{is\text{-}adversary\text{-}node}(\operatorname{pk\text{-}inf}(p), m) \wedge \\ \quad \operatorname{is\text{-}sent}(\operatorname{ispace}(\operatorname{pk\text{-}inf}(p)), m) \wedge \\ \quad \operatorname{same\text{-}strand\text{-}before}(\operatorname{ispace}(\operatorname{pk\text{-}inf}(p)), x, m) \wedge \\ \quad g \in \operatorname{msg\text{-}set}(\operatorname{tokens}(\operatorname{ispace}(\operatorname{pk\text{-}inf}(p))), \operatorname{keys}(\operatorname{ispace}(\operatorname{pk\text{-}inf}(p)))) \wedge \\ \quad y \in \operatorname{keys}(\operatorname{ispace}(\operatorname{pk\text{-}inf}(p))) \wedge \\ \quad \neg (\operatorname{kinv}(\operatorname{ispace}(\operatorname{pk\text{-}inf}(p)))(y) \in \operatorname{adv\text{-}tokens}(\operatorname{pk\text{-}inf}(p))) \wedge \\ \quad \operatorname{is\text{-}subterm}(\operatorname{tokens}(\operatorname{ispace}(\operatorname{pk\text{-}inf}(p))), \operatorname{keys}(\operatorname{ispace}(\operatorname{pk\text{-}inf}(p))), \operatorname{encr}(g, y), \\ \quad \quad \operatorname{node\text{-}content}(\operatorname{ispace}(\operatorname{pk\text{-}inf}(p)), x)) \wedge \\ \quad \neg \operatorname{is\text{-}subterm}(\operatorname{tokens}(\operatorname{ispace}(\operatorname{pk\text{-}inf}(p))), \operatorname{keys}(\operatorname{ispace}(\operatorname{pk\text{-}inf}(p))), \operatorname{encr}(g, y), \operatorname{node\text{-}content}(\operatorname{ispace}(\operatorname{pk\text{-}inf}(p)), m)) \Rightarrow \\ \quad \quad \exists z \in \operatorname{bnodes}(c). \\ \quad \quad \quad {[}z, m] \in \operatorname{precedes\text{-}strict}(c) \wedge \\ \quad \quad \quad \operatorname{node\text{-}content}(\operatorname{ispace}(\operatorname{pk\text{-}inf}(p)), z) = \operatorname{node\text{-}content}(\operatorname{ispace}(\operatorname{pk\text{-}inf}(p)), m) \end{array}}}
\vnbprop{nsl-resp-nonce-originates}{\vnbstmt{\begin{array}{@{}l@{}} \forall p, s, a, b, c, d. \\ \quad \operatorname{is\text{-}pk\text{-}space}(p) \wedge \\ \quad s \in \operatorname{ids}(\operatorname{ispace}(\operatorname{pk\text{-}inf}(p))) \wedge \\ \quad \operatorname{is\text{-}resp\text{-}strand}(p, \operatorname{strand\text{-}map}(\operatorname{ispace}(\operatorname{pk\text{-}inf}(p)))(s), a, b, c, d) \wedge \\ \quad \neg (c = d) \Rightarrow \\ \quad \quad \operatorname{originates}(\operatorname{ispace}(\operatorname{pk\text{-}inf}(p)), d, {[}2, s]) \end{array}}}
\vnbprop{nsl-init-nonce-originates}{\vnbstmt{\begin{array}{@{}l@{}} \forall p, t, a, b, c, d. \\ \quad \operatorname{is\text{-}pk\text{-}space}(p) \wedge \\ \quad t \in \operatorname{ids}(\operatorname{ispace}(\operatorname{pk\text{-}inf}(p))) \wedge \\ \quad \operatorname{is\text{-}init\text{-}strand}(p, \operatorname{strand\text{-}map}(\operatorname{ispace}(\operatorname{pk\text{-}inf}(p)))(t), a, b, c, d) \Rightarrow \\ \quad \quad \operatorname{originates}(\operatorname{ispace}(\operatorname{pk\text{-}inf}(p)), c, {[}1, t]) \end{array}}}
\vnbprop{nsl-init-unique}{\vnbstmt{\begin{array}{@{}l@{}} \forall p, a, b, c, d, f, g, a_{2}, b_{2}, h. \\ \quad \operatorname{is\text{-}pk\text{-}space}(p) \wedge \\ \quad \operatorname{uniquely\text{-}originates}(\operatorname{ispace}(\operatorname{pk\text{-}inf}(p)), f) \wedge \\ \quad a \in \operatorname{ids}(\operatorname{ispace}(\operatorname{pk\text{-}inf}(p))) \wedge \\ \quad b \in \operatorname{ids}(\operatorname{ispace}(\operatorname{pk\text{-}inf}(p))) \wedge \\ \quad \operatorname{is\text{-}init\text{-}strand}(p, \operatorname{strand\text{-}map}(\operatorname{ispace}(\operatorname{pk\text{-}inf}(p)))(a), c, d, f, g) \wedge \\ \quad \operatorname{is\text{-}init\text{-}strand}(p, \operatorname{strand\text{-}map}(\operatorname{ispace}(\operatorname{pk\text{-}inf}(p)))(b), a_{2}, b_{2}, f, h) \Rightarrow \\ \quad \quad a = b \end{array}}}
\vnbprop{nsl-lemma-5-4}{\vnbstmt{\begin{array}{@{}l@{}} \forall p, c, s, a, b, d, f. \\ \quad \operatorname{is\text{-}nsl\text{-}space}(p) \wedge \\ \quad \operatorname{is\text{-}bundle}(c) \wedge \\ \quad \operatorname{bspace}(c) = \operatorname{ispace}(\operatorname{pk\text{-}inf}(p)) \wedge \\ \quad s \in \operatorname{ids}(\operatorname{ispace}(\operatorname{pk\text{-}inf}(p))) \wedge \\ \quad \operatorname{is\text{-}resp\text{-}strand}(p, \operatorname{strand\text{-}map}(\operatorname{ispace}(\operatorname{pk\text{-}inf}(p)))(s), a, b, d, f) \wedge \\ \quad {[}3, s] \in \operatorname{bnodes}(c) \wedge \\ \quad \neg (\operatorname{kinv}(\operatorname{ispace}(\operatorname{pk\text{-}inf}(p)))(\operatorname{pubkey}(p)(a)) \in \operatorname{adv\text{-}tokens}(\operatorname{pk\text{-}inf}(p))) \wedge \\ \quad \neg (d = f) \wedge \\ \quad \operatorname{uniquely\text{-}originates}(\operatorname{ispace}(\operatorname{pk\text{-}inf}(p)), f) \Rightarrow \\ \quad \quad \exists g \in \operatorname{bnodes}(c). \\ \quad \quad \quad \operatorname{is\text{-}subterm}(\operatorname{tokens}(\operatorname{ispace}(\operatorname{pk\text{-}inf}(p))), \operatorname{keys}(\operatorname{ispace}(\operatorname{pk\text{-}inf}(p))), f, \\ \quad \quad \quad \quad \operatorname{node\text{-}content}(\operatorname{ispace}(\operatorname{pk\text{-}inf}(p)), g)) \wedge \\ \quad \quad \quad \neg \operatorname{is\text{-}subterm}(\operatorname{tokens}(\operatorname{ispace}(\operatorname{pk\text{-}inf}(p))), \operatorname{keys}(\operatorname{ispace}(\operatorname{pk\text{-}inf}(p))), \operatorname{encr}(\operatorname{lst}({[}d, f, b]), \operatorname{pubkey}(p)(a)), \operatorname{node\text{-}content}(\operatorname{ispace}(\operatorname{pk\text{-}inf}(p)), g)) \wedge \\ \quad \quad \quad \forall z \in \operatorname{bnodes}(c). \\ \quad \quad \quad \quad \operatorname{is\text{-}subterm}(\operatorname{tokens}(\operatorname{ispace}(\operatorname{pk\text{-}inf}(p))), \operatorname{keys}(\operatorname{ispace}(\operatorname{pk\text{-}inf}(p))), f, \\ \quad \quad \quad \quad \quad \operatorname{node\text{-}content}(\operatorname{ispace}(\operatorname{pk\text{-}inf}(p)), z)) \wedge \\ \quad \quad \quad \quad \neg \operatorname{is\text{-}subterm}(\operatorname{tokens}(\operatorname{ispace}(\operatorname{pk\text{-}inf}(p))), \operatorname{keys}(\operatorname{ispace}(\operatorname{pk\text{-}inf}(p))), \operatorname{encr}(\operatorname{lst}({[}d, f, b]), \operatorname{pubkey}(p)(a)), \operatorname{node\text{-}content}(\operatorname{ispace}(\operatorname{pk\text{-}inf}(p)), z)) \Rightarrow \\ \quad \quad \quad \quad \quad \neg ({[}z, g] \in \operatorname{precedes\text{-}strict}(c)) \wedge \\ \quad \quad \quad \operatorname{is\text{-}regular\text{-}node}(\operatorname{pk\text{-}inf}(p), g) \wedge \\ \quad \quad \quad \operatorname{is\text{-}sent}(\operatorname{ispace}(\operatorname{pk\text{-}inf}(p)), g) \end{array}}}
\vnbprop{nsl-lemma-5-6}{\vnbstmt{\begin{array}{@{}l@{}} \forall p, c, s, a, b, d, f.; \forall g \in \operatorname{bnodes}(c). \\ \quad \operatorname{is\text{-}nsl\text{-}space}(p) \wedge \\ \quad \operatorname{is\text{-}bundle}(c) \wedge \\ \quad \operatorname{bspace}(c) = \operatorname{ispace}(\operatorname{pk\text{-}inf}(p)) \wedge \\ \quad s \in \operatorname{ids}(\operatorname{ispace}(\operatorname{pk\text{-}inf}(p))) \wedge \\ \quad \operatorname{is\text{-}resp\text{-}strand}(p, \operatorname{strand\text{-}map}(\operatorname{ispace}(\operatorname{pk\text{-}inf}(p)))(s), a, b, d, f) \wedge \\ \quad {[}3, s] \in \operatorname{bnodes}(c) \wedge \\ \quad \neg (\operatorname{kinv}(\operatorname{ispace}(\operatorname{pk\text{-}inf}(p)))(\operatorname{pubkey}(p)(a)) \in \operatorname{adv\text{-}tokens}(\operatorname{pk\text{-}inf}(p))) \wedge \\ \quad \neg (d = f) \wedge \\ \quad \operatorname{uniquely\text{-}originates}(\operatorname{ispace}(\operatorname{pk\text{-}inf}(p)), f) \wedge \\ \quad \operatorname{is\text{-}subterm}(\operatorname{tokens}(\operatorname{ispace}(\operatorname{pk\text{-}inf}(p))), \operatorname{keys}(\operatorname{ispace}(\operatorname{pk\text{-}inf}(p))), f, \\ \quad \quad \operatorname{node\text{-}content}(\operatorname{ispace}(\operatorname{pk\text{-}inf}(p)), g)) \wedge \\ \quad \neg \operatorname{is\text{-}subterm}(\operatorname{tokens}(\operatorname{ispace}(\operatorname{pk\text{-}inf}(p))), \operatorname{keys}(\operatorname{ispace}(\operatorname{pk\text{-}inf}(p))), \operatorname{encr}(\operatorname{lst}({[}d, f, b]), \operatorname{pubkey}(p)(a)), \operatorname{node\text{-}content}(\operatorname{ispace}(\operatorname{pk\text{-}inf}(p)), g)) \wedge \\ \quad \forall z \in \operatorname{bnodes}(c). \\ \quad \quad \operatorname{is\text{-}subterm}(\operatorname{tokens}(\operatorname{ispace}(\operatorname{pk\text{-}inf}(p))), \operatorname{keys}(\operatorname{ispace}(\operatorname{pk\text{-}inf}(p))), f, \\ \quad \quad \quad \operatorname{node\text{-}content}(\operatorname{ispace}(\operatorname{pk\text{-}inf}(p)), z)) \wedge \\ \quad \quad \neg \operatorname{is\text{-}subterm}(\operatorname{tokens}(\operatorname{ispace}(\operatorname{pk\text{-}inf}(p))), \operatorname{keys}(\operatorname{ispace}(\operatorname{pk\text{-}inf}(p))), \operatorname{encr}(\operatorname{lst}({[}d, f, b]), \operatorname{pubkey}(p)(a)), \operatorname{node\text{-}content}(\operatorname{ispace}(\operatorname{pk\text{-}inf}(p)), z)) \Rightarrow \\ \quad \quad \quad \neg ({[}z, g] \in \operatorname{precedes\text{-}strict}(c)) \wedge \\ \quad \operatorname{is\text{-}regular\text{-}node}(\operatorname{pk\text{-}inf}(p), g) \wedge \\ \quad \operatorname{is\text{-}sent}(\operatorname{ispace}(\operatorname{pk\text{-}inf}(p)), g) \Rightarrow \\ \quad \quad \exists h. \\ \quad \quad \quad \operatorname{same\text{-}strand\text{-}before}(\operatorname{ispace}(\operatorname{pk\text{-}inf}(p)), h, g) \wedge \\ \quad \quad \quad \operatorname{node\text{-}content}(\operatorname{ispace}(\operatorname{pk\text{-}inf}(p)), h) = \operatorname{encr}(\operatorname{lst}({[}d, f, b]), \operatorname{pubkey}(p)(a)) \end{array}}}
\vnbprop{nsl-lemma-5-7}{\vnbstmt{\begin{array}{@{}l@{}} \forall p, a, b, c, d, f, g. \\ \quad \operatorname{is\text{-}nsl\text{-}space}(p) \wedge \\ \quad c \in \operatorname{names}(p) \wedge \\ \quad b \in \operatorname{nodes}(\operatorname{ispace}(\operatorname{pk\text{-}inf}(p))) \wedge \\ \quad \operatorname{is\text{-}regular\text{-}node}(\operatorname{pk\text{-}inf}(p), b) \wedge \\ \quad \operatorname{is\text{-}sent}(\operatorname{ispace}(\operatorname{pk\text{-}inf}(p)), b) \wedge \\ \quad \operatorname{same\text{-}strand\text{-}before}(\operatorname{ispace}(\operatorname{pk\text{-}inf}(p)), a, b) \wedge \\ \quad \operatorname{node\text{-}content}(\operatorname{ispace}(\operatorname{pk\text{-}inf}(p)), a) = \operatorname{encr}(\operatorname{lst}({[}f, g, d]), \operatorname{pubkey}(p)(c)) \Rightarrow \\ \quad \quad (\operatorname{is\text{-}init\text{-}strand}(p, \operatorname{strand\text{-}map}(\operatorname{ispace}(\operatorname{pk\text{-}inf}(p)))(\operatorname{nth}(2, b)), c, d, f, g) \wedge (\operatorname{nth}(1, b) = 3)) \end{array}}}
\vnbprop{nsl-responder-agreement}{\vnbstmt{\begin{array}{@{}l@{}} \forall p, c, s, a, b, d, f. \\ \quad \operatorname{is\text{-}nsl\text{-}space}(p) \wedge \\ \quad \operatorname{is\text{-}bundle}(c) \wedge \\ \quad \operatorname{bspace}(c) = \operatorname{ispace}(\operatorname{pk\text{-}inf}(p)) \wedge \\ \quad s \in \operatorname{ids}(\operatorname{ispace}(\operatorname{pk\text{-}inf}(p))) \wedge \\ \quad \operatorname{is\text{-}resp\text{-}strand}(p, \operatorname{strand\text{-}map}(\operatorname{ispace}(\operatorname{pk\text{-}inf}(p)))(s), a, b, d, f) \wedge \\ \quad {[}3, s] \in \operatorname{bnodes}(c) \wedge \\ \quad \neg (\operatorname{kinv}(\operatorname{ispace}(\operatorname{pk\text{-}inf}(p)))(\operatorname{pubkey}(p)(a)) \in \operatorname{adv\text{-}tokens}(\operatorname{pk\text{-}inf}(p))) \wedge \\ \quad \neg (d = f) \wedge \\ \quad \operatorname{uniquely\text{-}originates}(\operatorname{ispace}(\operatorname{pk\text{-}inf}(p)), f) \Rightarrow \\ \quad \quad \exists t \in \operatorname{ids}(\operatorname{ispace}(\operatorname{pk\text{-}inf}(p))). \\ \quad \quad \quad (\operatorname{is\text{-}init\text{-}strand}(p, \operatorname{strand\text{-}map}(\operatorname{ispace}(\operatorname{pk\text{-}inf}(p)))(t), a, b, d, f) \wedge ({[}3, t] \in \operatorname{bnodes}(c))) \end{array}}}
\vnbprop{nsl-lemma-5-12}{\vnbstmt{\begin{array}{@{}l@{}} \forall p, c, s, a, b, d, f.; \forall g \in \operatorname{bnodes}(c). \\ \quad \operatorname{is\text{-}nsl\text{-}space}(p) \wedge \\ \quad \operatorname{is\text{-}bundle}(c) \wedge \\ \quad \operatorname{bspace}(c) = \operatorname{ispace}(\operatorname{pk\text{-}inf}(p)) \wedge \\ \quad s \in \operatorname{ids}(\operatorname{ispace}(\operatorname{pk\text{-}inf}(p))) \wedge \\ \quad \operatorname{is\text{-}resp\text{-}strand}(p, \operatorname{strand\text{-}map}(\operatorname{ispace}(\operatorname{pk\text{-}inf}(p)))(s), a, b, d, f) \wedge \\ \quad \neg (\operatorname{kinv}(\operatorname{ispace}(\operatorname{pk\text{-}inf}(p)))(\operatorname{pubkey}(p)(a)) \in \operatorname{adv\text{-}tokens}(\operatorname{pk\text{-}inf}(p))) \wedge \\ \quad \neg (\operatorname{kinv}(\operatorname{ispace}(\operatorname{pk\text{-}inf}(p)))(\operatorname{pubkey}(p)(b)) \in \operatorname{adv\text{-}tokens}(\operatorname{pk\text{-}inf}(p))) \wedge \\ \quad \neg (d = f) \wedge \\ \quad \operatorname{uniquely\text{-}originates}(\operatorname{ispace}(\operatorname{pk\text{-}inf}(p)), f) \wedge \\ \quad \operatorname{is\text{-}subterm}(\operatorname{tokens}(\operatorname{ispace}(\operatorname{pk\text{-}inf}(p))), \operatorname{keys}(\operatorname{ispace}(\operatorname{pk\text{-}inf}(p))), f, \\ \quad \quad \operatorname{node\text{-}content}(\operatorname{ispace}(\operatorname{pk\text{-}inf}(p)), g)) \wedge \\ \quad \neg \operatorname{is\text{-}subterm}(\operatorname{tokens}(\operatorname{ispace}(\operatorname{pk\text{-}inf}(p))), \operatorname{keys}(\operatorname{ispace}(\operatorname{pk\text{-}inf}(p))), \operatorname{encr}(\operatorname{lst}({[}d, f, b]), \operatorname{pubkey}(p)(a)), \operatorname{node\text{-}content}(\operatorname{ispace}(\operatorname{pk\text{-}inf}(p)), g)) \wedge \\ \quad \neg \operatorname{is\text{-}subterm}(\operatorname{tokens}(\operatorname{ispace}(\operatorname{pk\text{-}inf}(p))), \operatorname{keys}(\operatorname{ispace}(\operatorname{pk\text{-}inf}(p))), \operatorname{encr}(f, \operatorname{pubkey}(p)(b)), \operatorname{node\text{-}content}(\operatorname{ispace}(\operatorname{pk\text{-}inf}(p)), g)) \wedge \\ \quad \forall z \in \operatorname{bnodes}(c). \\ \quad \quad \operatorname{is\text{-}subterm}(\operatorname{tokens}(\operatorname{ispace}(\operatorname{pk\text{-}inf}(p))), \operatorname{keys}(\operatorname{ispace}(\operatorname{pk\text{-}inf}(p))), f, \\ \quad \quad \quad \operatorname{node\text{-}content}(\operatorname{ispace}(\operatorname{pk\text{-}inf}(p)), z)) \wedge \\ \quad \quad \neg \operatorname{is\text{-}subterm}(\operatorname{tokens}(\operatorname{ispace}(\operatorname{pk\text{-}inf}(p))), \operatorname{keys}(\operatorname{ispace}(\operatorname{pk\text{-}inf}(p))), \operatorname{encr}(\operatorname{lst}({[}d, f, b]), \operatorname{pubkey}(p)(a)), \operatorname{node\text{-}content}(\operatorname{ispace}(\operatorname{pk\text{-}inf}(p)), z)) \wedge \\ \quad \quad \neg \operatorname{is\text{-}subterm}(\operatorname{tokens}(\operatorname{ispace}(\operatorname{pk\text{-}inf}(p))), \operatorname{keys}(\operatorname{ispace}(\operatorname{pk\text{-}inf}(p))), \operatorname{encr}(f, \operatorname{pubkey}(p)(b)), \operatorname{node\text{-}content}(\operatorname{ispace}(\operatorname{pk\text{-}inf}(p)), z)) \Rightarrow \\ \quad \quad \quad \neg ({[}z, g] \in \operatorname{precedes\text{-}strict}(c)) \Rightarrow \\ \quad \quad \neg \operatorname{is\text{-}adversary\text{-}node}(\operatorname{pk\text{-}inf}(p), g) \end{array}}}
\vnbprop{nsl-lemma-5-11}{\vnbstmt{\begin{array}{@{}l@{}} \forall p, c, s, a, b, d, f.; \forall g \in \operatorname{bnodes}(c). \\ \quad \operatorname{is\text{-}nsl\text{-}space}(p) \wedge \\ \quad \operatorname{is\text{-}bundle}(c) \wedge \\ \quad \operatorname{bspace}(c) = \operatorname{ispace}(\operatorname{pk\text{-}inf}(p)) \wedge \\ \quad s \in \operatorname{ids}(\operatorname{ispace}(\operatorname{pk\text{-}inf}(p))) \wedge \\ \quad \operatorname{is\text{-}resp\text{-}strand}(p, \operatorname{strand\text{-}map}(\operatorname{ispace}(\operatorname{pk\text{-}inf}(p)))(s), a, b, d, f) \wedge \\ \quad \neg (\operatorname{kinv}(\operatorname{ispace}(\operatorname{pk\text{-}inf}(p)))(\operatorname{pubkey}(p)(a)) \in \operatorname{adv\text{-}tokens}(\operatorname{pk\text{-}inf}(p))) \wedge \\ \quad \neg (\operatorname{kinv}(\operatorname{ispace}(\operatorname{pk\text{-}inf}(p)))(\operatorname{pubkey}(p)(b)) \in \operatorname{adv\text{-}tokens}(\operatorname{pk\text{-}inf}(p))) \wedge \\ \quad \neg (d = f) \wedge \\ \quad \operatorname{uniquely\text{-}originates}(\operatorname{ispace}(\operatorname{pk\text{-}inf}(p)), f) \wedge \\ \quad \operatorname{is\text{-}subterm}(\operatorname{tokens}(\operatorname{ispace}(\operatorname{pk\text{-}inf}(p))), \operatorname{keys}(\operatorname{ispace}(\operatorname{pk\text{-}inf}(p))), f, \\ \quad \quad \operatorname{node\text{-}content}(\operatorname{ispace}(\operatorname{pk\text{-}inf}(p)), g)) \wedge \\ \quad \neg \operatorname{is\text{-}subterm}(\operatorname{tokens}(\operatorname{ispace}(\operatorname{pk\text{-}inf}(p))), \operatorname{keys}(\operatorname{ispace}(\operatorname{pk\text{-}inf}(p))), \operatorname{encr}(\operatorname{lst}({[}d, f, b]), \operatorname{pubkey}(p)(a)), \operatorname{node\text{-}content}(\operatorname{ispace}(\operatorname{pk\text{-}inf}(p)), g)) \wedge \\ \quad \neg \operatorname{is\text{-}subterm}(\operatorname{tokens}(\operatorname{ispace}(\operatorname{pk\text{-}inf}(p))), \operatorname{keys}(\operatorname{ispace}(\operatorname{pk\text{-}inf}(p))), \operatorname{encr}(f, \operatorname{pubkey}(p)(b)), \operatorname{node\text{-}content}(\operatorname{ispace}(\operatorname{pk\text{-}inf}(p)), g)) \wedge \\ \quad \forall z \in \operatorname{bnodes}(c). \\ \quad \quad \operatorname{is\text{-}subterm}(\operatorname{tokens}(\operatorname{ispace}(\operatorname{pk\text{-}inf}(p))), \operatorname{keys}(\operatorname{ispace}(\operatorname{pk\text{-}inf}(p))), f, \\ \quad \quad \quad \operatorname{node\text{-}content}(\operatorname{ispace}(\operatorname{pk\text{-}inf}(p)), z)) \wedge \\ \quad \quad \neg \operatorname{is\text{-}subterm}(\operatorname{tokens}(\operatorname{ispace}(\operatorname{pk\text{-}inf}(p))), \operatorname{keys}(\operatorname{ispace}(\operatorname{pk\text{-}inf}(p))), \operatorname{encr}(\operatorname{lst}({[}d, f, b]), \operatorname{pubkey}(p)(a)), \operatorname{node\text{-}content}(\operatorname{ispace}(\operatorname{pk\text{-}inf}(p)), z)) \wedge \\ \quad \quad \neg \operatorname{is\text{-}subterm}(\operatorname{tokens}(\operatorname{ispace}(\operatorname{pk\text{-}inf}(p))), \operatorname{keys}(\operatorname{ispace}(\operatorname{pk\text{-}inf}(p))), \operatorname{encr}(f, \operatorname{pubkey}(p)(b)), \operatorname{node\text{-}content}(\operatorname{ispace}(\operatorname{pk\text{-}inf}(p)), z)) \Rightarrow \\ \quad \quad \quad \neg ({[}z, g] \in \operatorname{precedes\text{-}strict}(c)) \Rightarrow \\ \quad \quad \neg \operatorname{is\text{-}regular\text{-}node}(\operatorname{pk\text{-}inf}(p), g) \end{array}}}
\vnbprop{nsl-responder-secrecy}{\vnbstmt{\begin{array}{@{}l@{}} \forall p, c, s, a, b, d, f.; \forall m \in \operatorname{bnodes}(c). \\ \quad \operatorname{is\text{-}nsl\text{-}space}(p) \wedge \\ \quad \operatorname{is\text{-}bundle}(c) \wedge \\ \quad \operatorname{bspace}(c) = \operatorname{ispace}(\operatorname{pk\text{-}inf}(p)) \wedge \\ \quad s \in \operatorname{ids}(\operatorname{ispace}(\operatorname{pk\text{-}inf}(p))) \wedge \\ \quad \operatorname{is\text{-}resp\text{-}strand}(p, \operatorname{strand\text{-}map}(\operatorname{ispace}(\operatorname{pk\text{-}inf}(p)))(s), a, b, d, f) \wedge \\ \quad \neg (\operatorname{kinv}(\operatorname{ispace}(\operatorname{pk\text{-}inf}(p)))(\operatorname{pubkey}(p)(a)) \in \operatorname{adv\text{-}tokens}(\operatorname{pk\text{-}inf}(p))) \wedge \\ \quad \neg (\operatorname{kinv}(\operatorname{ispace}(\operatorname{pk\text{-}inf}(p)))(\operatorname{pubkey}(p)(b)) \in \operatorname{adv\text{-}tokens}(\operatorname{pk\text{-}inf}(p))) \wedge \\ \quad \neg (d = f) \wedge \\ \quad \operatorname{uniquely\text{-}originates}(\operatorname{ispace}(\operatorname{pk\text{-}inf}(p)), f) \wedge \\ \quad \operatorname{is\text{-}subterm}(\operatorname{tokens}(\operatorname{ispace}(\operatorname{pk\text{-}inf}(p))), \operatorname{keys}(\operatorname{ispace}(\operatorname{pk\text{-}inf}(p))), f, \\ \quad \quad \operatorname{node\text{-}content}(\operatorname{ispace}(\operatorname{pk\text{-}inf}(p)), m)) \Rightarrow \\ \quad \quad \operatorname{is\text{-}subterm}(\operatorname{tokens}(\operatorname{ispace}(\operatorname{pk\text{-}inf}(p))), \operatorname{keys}(\operatorname{ispace}(\operatorname{pk\text{-}inf}(p))), \\ \quad \quad \quad \operatorname{encr}(\operatorname{lst}({[}d, f, b]), \operatorname{pubkey}(p)(a)), \\ \quad \quad \quad \operatorname{node\text{-}content}(\operatorname{ispace}(\operatorname{pk\text{-}inf}(p)), m)) \vee \\ \quad \quad \operatorname{is\text{-}subterm}(\operatorname{tokens}(\operatorname{ispace}(\operatorname{pk\text{-}inf}(p))), \operatorname{keys}(\operatorname{ispace}(\operatorname{pk\text{-}inf}(p))), \\ \quad \quad \quad \operatorname{encr}(f, \operatorname{pubkey}(p)(b)), \operatorname{node\text{-}content}(\operatorname{ispace}(\operatorname{pk\text{-}inf}(p)), m)) \end{array}}}
\vnbprop{nsl-responder-nonce-not-sent}{\vnbstmt{\begin{array}{@{}l@{}} \forall p, c, s, a, b, d, f.; \forall m \in \operatorname{bnodes}(c). \\ \quad \operatorname{is\text{-}nsl\text{-}space}(p) \wedge \\ \quad \operatorname{is\text{-}bundle}(c) \wedge \\ \quad \operatorname{bspace}(c) = \operatorname{ispace}(\operatorname{pk\text{-}inf}(p)) \wedge \\ \quad s \in \operatorname{ids}(\operatorname{ispace}(\operatorname{pk\text{-}inf}(p))) \wedge \\ \quad \operatorname{is\text{-}resp\text{-}strand}(p, \operatorname{strand\text{-}map}(\operatorname{ispace}(\operatorname{pk\text{-}inf}(p)))(s), a, b, d, f) \wedge \\ \quad \neg (\operatorname{kinv}(\operatorname{ispace}(\operatorname{pk\text{-}inf}(p)))(\operatorname{pubkey}(p)(a)) \in \operatorname{adv\text{-}tokens}(\operatorname{pk\text{-}inf}(p))) \wedge \\ \quad \neg (\operatorname{kinv}(\operatorname{ispace}(\operatorname{pk\text{-}inf}(p)))(\operatorname{pubkey}(p)(b)) \in \operatorname{adv\text{-}tokens}(\operatorname{pk\text{-}inf}(p))) \wedge \\ \quad \neg (d = f) \wedge \\ \quad \operatorname{uniquely\text{-}originates}(\operatorname{ispace}(\operatorname{pk\text{-}inf}(p)), f) \Rightarrow \\ \quad \quad \neg (\operatorname{node\text{-}content}(\operatorname{ispace}(\operatorname{pk\text{-}inf}(p)), m) = f) \end{array}}}
\vnbprop{nsl-initiator-secrecy-adversary}{\vnbstmt{\begin{array}{@{}l@{}} \forall p, c, s, a, b, d, f.; \forall g \in \operatorname{bnodes}(c). \\ \quad \operatorname{is\text{-}nsl\text{-}space}(p) \wedge \\ \quad \operatorname{is\text{-}bundle}(c) \wedge \\ \quad \operatorname{bspace}(c) = \operatorname{ispace}(\operatorname{pk\text{-}inf}(p)) \wedge \\ \quad s \in \operatorname{ids}(\operatorname{ispace}(\operatorname{pk\text{-}inf}(p))) \wedge \\ \quad \operatorname{is\text{-}init\text{-}strand}(p, \operatorname{strand\text{-}map}(\operatorname{ispace}(\operatorname{pk\text{-}inf}(p)))(s), a, b, d, f) \wedge \\ \quad {[}3, s] \in \operatorname{bnodes}(c) \wedge \\ \quad \neg (\operatorname{kinv}(\operatorname{ispace}(\operatorname{pk\text{-}inf}(p)))(\operatorname{pubkey}(p)(a)) \in \operatorname{adv\text{-}tokens}(\operatorname{pk\text{-}inf}(p))) \wedge \\ \quad \neg (\operatorname{kinv}(\operatorname{ispace}(\operatorname{pk\text{-}inf}(p)))(\operatorname{pubkey}(p)(b)) \in \operatorname{adv\text{-}tokens}(\operatorname{pk\text{-}inf}(p))) \wedge \\ \quad \neg (d = f) \wedge \\ \quad \operatorname{uniquely\text{-}originates}(\operatorname{ispace}(\operatorname{pk\text{-}inf}(p)), d) \wedge \\ \quad \operatorname{is\text{-}subterm}(\operatorname{tokens}(\operatorname{ispace}(\operatorname{pk\text{-}inf}(p))), \operatorname{keys}(\operatorname{ispace}(\operatorname{pk\text{-}inf}(p))), d, \\ \quad \quad \operatorname{node\text{-}content}(\operatorname{ispace}(\operatorname{pk\text{-}inf}(p)), g)) \wedge \\ \quad \neg \operatorname{is\text{-}subterm}(\operatorname{tokens}(\operatorname{ispace}(\operatorname{pk\text{-}inf}(p))), \operatorname{keys}(\operatorname{ispace}(\operatorname{pk\text{-}inf}(p))), \operatorname{encr}(\operatorname{lst}({[}d, a]), \operatorname{pubkey}(p)(b)), \operatorname{node\text{-}content}(\operatorname{ispace}(\operatorname{pk\text{-}inf}(p)), g)) \wedge \\ \quad \forall w \in \operatorname{tokens}(\operatorname{ispace}(\operatorname{pk\text{-}inf}(p))). \\ \quad \quad \neg \operatorname{is\text{-}subterm}(\operatorname{tokens}(\operatorname{ispace}(\operatorname{pk\text{-}inf}(p))), \operatorname{keys}(\operatorname{ispace}(\operatorname{pk\text{-}inf}(p))), \operatorname{encr}(\operatorname{lst}({[}d, w, b]), \operatorname{pubkey}(p)(a)), \operatorname{node\text{-}content}(\operatorname{ispace}(\operatorname{pk\text{-}inf}(p)), g)) \wedge \\ \quad \forall z \in \operatorname{bnodes}(c). \\ \quad \quad \operatorname{is\text{-}subterm}(\operatorname{tokens}(\operatorname{ispace}(\operatorname{pk\text{-}inf}(p))), \operatorname{keys}(\operatorname{ispace}(\operatorname{pk\text{-}inf}(p))), d, \\ \quad \quad \quad \operatorname{node\text{-}content}(\operatorname{ispace}(\operatorname{pk\text{-}inf}(p)), z)) \wedge \\ \quad \quad \neg \operatorname{is\text{-}subterm}(\operatorname{tokens}(\operatorname{ispace}(\operatorname{pk\text{-}inf}(p))), \operatorname{keys}(\operatorname{ispace}(\operatorname{pk\text{-}inf}(p))), \operatorname{encr}(\operatorname{lst}({[}d, a]), \operatorname{pubkey}(p)(b)), \operatorname{node\text{-}content}(\operatorname{ispace}(\operatorname{pk\text{-}inf}(p)), z)) \wedge \\ \quad \quad \forall w \in \operatorname{tokens}(\operatorname{ispace}(\operatorname{pk\text{-}inf}(p))). \\ \quad \quad \quad \neg \operatorname{is\text{-}subterm}(\operatorname{tokens}(\operatorname{ispace}(\operatorname{pk\text{-}inf}(p))), \operatorname{keys}(\operatorname{ispace}(\operatorname{pk\text{-}inf}(p))), \operatorname{encr}(\operatorname{lst}({[}d, w, b]), \operatorname{pubkey}(p)(a)), \operatorname{node\text{-}content}(\operatorname{ispace}(\operatorname{pk\text{-}inf}(p)), z)) \Rightarrow \\ \quad \quad \quad \neg ({[}z, g] \in \operatorname{precedes\text{-}strict}(c)) \Rightarrow \\ \quad \quad \neg \operatorname{is\text{-}adversary\text{-}node}(\operatorname{pk\text{-}inf}(p), g) \end{array}}}
\vnbprop{nsl-initiator-secrecy-regular}{\vnbstmt{\begin{array}{@{}l@{}} \forall p, c, s, a, b, d, f.; \forall g \in \operatorname{bnodes}(c). \\ \quad \operatorname{is\text{-}nsl\text{-}space}(p) \wedge \\ \quad \operatorname{is\text{-}bundle}(c) \wedge \\ \quad \operatorname{bspace}(c) = \operatorname{ispace}(\operatorname{pk\text{-}inf}(p)) \wedge \\ \quad s \in \operatorname{ids}(\operatorname{ispace}(\operatorname{pk\text{-}inf}(p))) \wedge \\ \quad \operatorname{is\text{-}init\text{-}strand}(p, \operatorname{strand\text{-}map}(\operatorname{ispace}(\operatorname{pk\text{-}inf}(p)))(s), a, b, d, f) \wedge \\ \quad {[}3, s] \in \operatorname{bnodes}(c) \wedge \\ \quad \neg (\operatorname{kinv}(\operatorname{ispace}(\operatorname{pk\text{-}inf}(p)))(\operatorname{pubkey}(p)(a)) \in \operatorname{adv\text{-}tokens}(\operatorname{pk\text{-}inf}(p))) \wedge \\ \quad \neg (\operatorname{kinv}(\operatorname{ispace}(\operatorname{pk\text{-}inf}(p)))(\operatorname{pubkey}(p)(b)) \in \operatorname{adv\text{-}tokens}(\operatorname{pk\text{-}inf}(p))) \wedge \\ \quad \neg (d = f) \wedge \\ \quad \operatorname{uniquely\text{-}originates}(\operatorname{ispace}(\operatorname{pk\text{-}inf}(p)), d) \wedge \\ \quad \operatorname{is\text{-}subterm}(\operatorname{tokens}(\operatorname{ispace}(\operatorname{pk\text{-}inf}(p))), \operatorname{keys}(\operatorname{ispace}(\operatorname{pk\text{-}inf}(p))), d, \\ \quad \quad \operatorname{node\text{-}content}(\operatorname{ispace}(\operatorname{pk\text{-}inf}(p)), g)) \wedge \\ \quad \neg \operatorname{is\text{-}subterm}(\operatorname{tokens}(\operatorname{ispace}(\operatorname{pk\text{-}inf}(p))), \operatorname{keys}(\operatorname{ispace}(\operatorname{pk\text{-}inf}(p))), \operatorname{encr}(\operatorname{lst}({[}d, a]), \operatorname{pubkey}(p)(b)), \operatorname{node\text{-}content}(\operatorname{ispace}(\operatorname{pk\text{-}inf}(p)), g)) \wedge \\ \quad \forall w \in \operatorname{tokens}(\operatorname{ispace}(\operatorname{pk\text{-}inf}(p))). \\ \quad \quad \neg \operatorname{is\text{-}subterm}(\operatorname{tokens}(\operatorname{ispace}(\operatorname{pk\text{-}inf}(p))), \operatorname{keys}(\operatorname{ispace}(\operatorname{pk\text{-}inf}(p))), \operatorname{encr}(\operatorname{lst}({[}d, w, b]), \operatorname{pubkey}(p)(a)), \operatorname{node\text{-}content}(\operatorname{ispace}(\operatorname{pk\text{-}inf}(p)), g)) \wedge \\ \quad \forall z \in \operatorname{bnodes}(c). \\ \quad \quad \operatorname{is\text{-}subterm}(\operatorname{tokens}(\operatorname{ispace}(\operatorname{pk\text{-}inf}(p))), \operatorname{keys}(\operatorname{ispace}(\operatorname{pk\text{-}inf}(p))), d, \\ \quad \quad \quad \operatorname{node\text{-}content}(\operatorname{ispace}(\operatorname{pk\text{-}inf}(p)), z)) \wedge \\ \quad \quad \neg \operatorname{is\text{-}subterm}(\operatorname{tokens}(\operatorname{ispace}(\operatorname{pk\text{-}inf}(p))), \operatorname{keys}(\operatorname{ispace}(\operatorname{pk\text{-}inf}(p))), \operatorname{encr}(\operatorname{lst}({[}d, a]), \operatorname{pubkey}(p)(b)), \operatorname{node\text{-}content}(\operatorname{ispace}(\operatorname{pk\text{-}inf}(p)), z)) \wedge \\ \quad \quad \forall w \in \operatorname{tokens}(\operatorname{ispace}(\operatorname{pk\text{-}inf}(p))). \\ \quad \quad \quad \neg \operatorname{is\text{-}subterm}(\operatorname{tokens}(\operatorname{ispace}(\operatorname{pk\text{-}inf}(p))), \operatorname{keys}(\operatorname{ispace}(\operatorname{pk\text{-}inf}(p))), \operatorname{encr}(\operatorname{lst}({[}d, w, b]), \operatorname{pubkey}(p)(a)), \operatorname{node\text{-}content}(\operatorname{ispace}(\operatorname{pk\text{-}inf}(p)), z)) \Rightarrow \\ \quad \quad \quad \neg ({[}z, g] \in \operatorname{precedes\text{-}strict}(c)) \Rightarrow \\ \quad \quad \neg \operatorname{is\text{-}regular\text{-}node}(\operatorname{pk\text{-}inf}(p), g) \end{array}}}
\vnbprop{nsl-initiator-secrecy}{\vnbstmt{\begin{array}{@{}l@{}} \forall p, c, s, a, b, d, f.; \forall m \in \operatorname{bnodes}(c). \\ \quad \operatorname{is\text{-}nsl\text{-}space}(p) \wedge \\ \quad \operatorname{is\text{-}bundle}(c) \wedge \\ \quad \operatorname{bspace}(c) = \operatorname{ispace}(\operatorname{pk\text{-}inf}(p)) \wedge \\ \quad s \in \operatorname{ids}(\operatorname{ispace}(\operatorname{pk\text{-}inf}(p))) \wedge \\ \quad \operatorname{is\text{-}init\text{-}strand}(p, \operatorname{strand\text{-}map}(\operatorname{ispace}(\operatorname{pk\text{-}inf}(p)))(s), a, b, d, f) \wedge \\ \quad {[}3, s] \in \operatorname{bnodes}(c) \wedge \\ \quad \neg (\operatorname{kinv}(\operatorname{ispace}(\operatorname{pk\text{-}inf}(p)))(\operatorname{pubkey}(p)(a)) \in \operatorname{adv\text{-}tokens}(\operatorname{pk\text{-}inf}(p))) \wedge \\ \quad \neg (\operatorname{kinv}(\operatorname{ispace}(\operatorname{pk\text{-}inf}(p)))(\operatorname{pubkey}(p)(b)) \in \operatorname{adv\text{-}tokens}(\operatorname{pk\text{-}inf}(p))) \wedge \\ \quad \neg (d = f) \wedge \\ \quad \operatorname{uniquely\text{-}originates}(\operatorname{ispace}(\operatorname{pk\text{-}inf}(p)), d) \wedge \\ \quad \operatorname{is\text{-}subterm}(\operatorname{tokens}(\operatorname{ispace}(\operatorname{pk\text{-}inf}(p))), \operatorname{keys}(\operatorname{ispace}(\operatorname{pk\text{-}inf}(p))), d, \\ \quad \quad \operatorname{node\text{-}content}(\operatorname{ispace}(\operatorname{pk\text{-}inf}(p)), m)) \Rightarrow \\ \quad \quad \operatorname{is\text{-}subterm}(\operatorname{tokens}(\operatorname{ispace}(\operatorname{pk\text{-}inf}(p))), \operatorname{keys}(\operatorname{ispace}(\operatorname{pk\text{-}inf}(p))), \\ \quad \quad \quad \operatorname{encr}(\operatorname{lst}({[}d, a]), \operatorname{pubkey}(p)(b)), \\ \quad \quad \quad \operatorname{node\text{-}content}(\operatorname{ispace}(\operatorname{pk\text{-}inf}(p)), m)) \vee \\ \quad \quad \exists w \in \operatorname{tokens}(\operatorname{ispace}(\operatorname{pk\text{-}inf}(p))). \\ \quad \quad \quad \operatorname{is\text{-}subterm}(\operatorname{tokens}(\operatorname{ispace}(\operatorname{pk\text{-}inf}(p))), \operatorname{keys}(\operatorname{ispace}(\operatorname{pk\text{-}inf}(p))), \\ \quad \quad \quad \quad \operatorname{encr}(\operatorname{lst}({[}d, w, b]), \operatorname{pubkey}(p)(a)), \\ \quad \quad \quad \quad \operatorname{node\text{-}content}(\operatorname{ispace}(\operatorname{pk\text{-}inf}(p)), m)) \end{array}}}
\vnbprop{nsl-initiator-nonce-not-sent}{\vnbstmt{\begin{array}{@{}l@{}} \forall p, c, s, a, b, d, f.; \forall m \in \operatorname{bnodes}(c). \\ \quad \operatorname{is\text{-}nsl\text{-}space}(p) \wedge \\ \quad \operatorname{is\text{-}bundle}(c) \wedge \\ \quad \operatorname{bspace}(c) = \operatorname{ispace}(\operatorname{pk\text{-}inf}(p)) \wedge \\ \quad s \in \operatorname{ids}(\operatorname{ispace}(\operatorname{pk\text{-}inf}(p))) \wedge \\ \quad \operatorname{is\text{-}init\text{-}strand}(p, \operatorname{strand\text{-}map}(\operatorname{ispace}(\operatorname{pk\text{-}inf}(p)))(s), a, b, d, f) \wedge \\ \quad {[}3, s] \in \operatorname{bnodes}(c) \wedge \\ \quad \neg (\operatorname{kinv}(\operatorname{ispace}(\operatorname{pk\text{-}inf}(p)))(\operatorname{pubkey}(p)(a)) \in \operatorname{adv\text{-}tokens}(\operatorname{pk\text{-}inf}(p))) \wedge \\ \quad \neg (\operatorname{kinv}(\operatorname{ispace}(\operatorname{pk\text{-}inf}(p)))(\operatorname{pubkey}(p)(b)) \in \operatorname{adv\text{-}tokens}(\operatorname{pk\text{-}inf}(p))) \wedge \\ \quad \neg (d = f) \wedge \\ \quad \operatorname{uniquely\text{-}originates}(\operatorname{ispace}(\operatorname{pk\text{-}inf}(p)), d) \Rightarrow \\ \quad \quad \neg (\operatorname{node\text{-}content}(\operatorname{ispace}(\operatorname{pk\text{-}inf}(p)), m) = d) \end{array}}}
\vnbprop{nsl-initiator-agreement}{\vnbstmt{\begin{array}{@{}l@{}} \forall p, c, s, a, b, d, f. \\ \quad \operatorname{is\text{-}nsl\text{-}space}(p) \wedge \\ \quad \operatorname{is\text{-}bundle}(c) \wedge \\ \quad \operatorname{bspace}(c) = \operatorname{ispace}(\operatorname{pk\text{-}inf}(p)) \wedge \\ \quad s \in \operatorname{ids}(\operatorname{ispace}(\operatorname{pk\text{-}inf}(p))) \wedge \\ \quad \operatorname{is\text{-}init\text{-}strand}(p, \operatorname{strand\text{-}map}(\operatorname{ispace}(\operatorname{pk\text{-}inf}(p)))(s), a, b, d, f) \wedge \\ \quad {[}3, s] \in \operatorname{bnodes}(c) \wedge \\ \quad \neg (\operatorname{kinv}(\operatorname{ispace}(\operatorname{pk\text{-}inf}(p)))(\operatorname{pubkey}(p)(a)) \in \operatorname{adv\text{-}tokens}(\operatorname{pk\text{-}inf}(p))) \wedge \\ \quad \neg (\operatorname{kinv}(\operatorname{ispace}(\operatorname{pk\text{-}inf}(p)))(\operatorname{pubkey}(p)(b)) \in \operatorname{adv\text{-}tokens}(\operatorname{pk\text{-}inf}(p))) \wedge \\ \quad \neg (d = f) \wedge \\ \quad \operatorname{uniquely\text{-}originates}(\operatorname{ispace}(\operatorname{pk\text{-}inf}(p)), d) \Rightarrow \\ \quad \quad \exists t \in \operatorname{ids}(\operatorname{ispace}(\operatorname{pk\text{-}inf}(p))). \\ \quad \quad \quad (\operatorname{is\text{-}resp\text{-}strand}(p, \operatorname{strand\text{-}map}(\operatorname{ispace}(\operatorname{pk\text{-}inf}(p)))(t), a, b, d, f) \wedge ({[}2, t] \in \operatorname{bnodes}(c))) \end{array}}}

\vnbbookfile{structure-library/flat-regular}{structure-library/flat-regular}
\noindent{\small flat-regular.scm -- PUBLIC-KEY SPACES WHOSE REGULAR CONTENTS ARE FLAT: the one property of the Needham-Schroeder-Lowe space (structure-library/nsl.scm) that sections 4-11 of theorem-library/nsl.scm use, abstracted so that NSL, the original Needham-Schroeder protocol (structure-library/nso.scm) and later protocols are ...}\par
\par\medskip\noindent\textbf{Definition} (\texttt{is-flat-regular}).\;For every $\operatorname{frp\_}$, 
\begin{equation*}
\begin{array}{@{}l@{}} \operatorname{is\text{-}flat\text{-}regular}(\operatorname{frp\_}) \\ \Leftrightarrow\; \operatorname{is\text{-}pk\text{-}space}(\operatorname{frp\_}) \wedge \\ \forall \operatorname{frn\_} \in \operatorname{nodes}(\operatorname{ispace}(\operatorname{pk\text{-}inf}(\operatorname{frp\_}))). \\ \quad \operatorname{is\text{-}regular\text{-}node}(\operatorname{pk\text{-}inf}(\operatorname{frp\_}), \operatorname{frn\_}) \Rightarrow  \\ \quad \quad \exists \operatorname{frk\_} \in \operatorname{keys}(\operatorname{ispace}(\operatorname{pk\text{-}inf}(\operatorname{frp\_}))). \\ \quad \quad \quad \exists \operatorname{frx\_} \in \{\operatorname{frd\_} \in \operatorname{tokens}(\operatorname{ispace}(\operatorname{pk\text{-}inf}(\operatorname{frp\_}))) : \neg (\operatorname{frd\_} \in \operatorname{keys}(\operatorname{ispace}(\operatorname{pk\text{-}inf}(\operatorname{frp\_}))))\}. \\ \quad \quad \quad \quad \operatorname{node\text{-}content}(\operatorname{ispace}(\operatorname{pk\text{-}inf}(\operatorname{frp\_})), \operatorname{frn\_}) = \operatorname{encr}(\operatorname{frx\_}, \operatorname{frk\_}) \vee \\ \quad \quad \quad \exists \operatorname{frl\_} \in \operatorname{tuples}(\{\operatorname{frd\_} \in \operatorname{tokens}(\operatorname{ispace}(\operatorname{pk\text{-}inf}(\operatorname{frp\_}))) : \neg (\operatorname{frd\_} \in \operatorname{keys}(\operatorname{ispace}(\operatorname{pk\text{-}inf}(\operatorname{frp\_}))))\}). \\ \quad \quad \quad \quad \operatorname{node\text{-}content}(\operatorname{ispace}(\operatorname{pk\text{-}inf}(\operatorname{frp\_})), \operatorname{frn\_}) = \operatorname{encr}(\operatorname{lst}(\operatorname{frl\_}), \operatorname{frk\_}) \end{array}
\end{equation*}


\vnbbookfile{structure-library/nso}{structure-library/nso}
\noindent{\small nso.scm -- THE ORIGINAL NEEDHAM-SCHROEDER PROTOCOL (Thayer, Herzog, Guttman, section 5.3), ON THE USER'S MODEL.  The NSL vocabulary of structure-library/nsl.scm (the public-key space PK-SPACE is reused as it is) with the second message \{Na Nb\}\_KA: the responder's name is NOT in it (the paper: "in the original protocol, ...}\par
\par\medskip\noindent\textbf{Definition} (\texttt{is-nso-init-strand}).\;For all $\operatorname{nlp\_}$, $\operatorname{nls\_}$, $\operatorname{nla\_}$, $\operatorname{nlb\_}$, $\operatorname{nlm\_}$ and $\operatorname{nln\_}$, 
\begin{equation*}
\begin{array}{@{}l@{}} \operatorname{is\text{-}nso\text{-}init\text{-}strand}(\operatorname{nlp\_}, \operatorname{nls\_}, \operatorname{nla\_}, \operatorname{nlb\_}, \operatorname{nlm\_}, \operatorname{nln\_}) \\ \Leftrightarrow\; \operatorname{nla\_} \in \operatorname{names}(\operatorname{nlp\_}) \wedge \\ \operatorname{nlb\_} \in \operatorname{names}(\operatorname{nlp\_}) \wedge \\ \operatorname{nlm\_} \in \operatorname{tokens}(\operatorname{ispace}(\operatorname{pk\text{-}inf}(\operatorname{nlp\_}))) \wedge \\ \operatorname{nln\_} \in \operatorname{tokens}(\operatorname{ispace}(\operatorname{pk\text{-}inf}(\operatorname{nlp\_}))) \wedge \\ \neg (\operatorname{nlm\_} \in \operatorname{names}(\operatorname{nlp\_})) \wedge \\ \neg (\operatorname{nlm\_} \in \operatorname{keys}(\operatorname{ispace}(\operatorname{pk\text{-}inf}(\operatorname{nlp\_})))) \wedge \\ \neg (\operatorname{nln\_} \in \operatorname{keys}(\operatorname{ispace}(\operatorname{pk\text{-}inf}(\operatorname{nlp\_})))) \wedge \\ \operatorname{nls\_} \\ =\; {[}{[}\mathit{sent}, \operatorname{encr}(\operatorname{lst}({[}\operatorname{nlm\_}, \operatorname{nla\_}]), \operatorname{pubkey}(\operatorname{nlp\_})(\operatorname{nlb\_}))], {[}\mathit{received}, \operatorname{encr}(\operatorname{lst}({[}\operatorname{nlm\_}, \operatorname{nln\_}]), \operatorname{pubkey}(\operatorname{nlp\_})(\operatorname{nla\_}))], {[}\mathit{sent}, \operatorname{encr}(\operatorname{nln\_}, \operatorname{pubkey}(\operatorname{nlp\_})(\operatorname{nlb\_}))]] \end{array}
\end{equation*}

\par\medskip\noindent\textbf{Definition} (\texttt{is-nso-resp-strand}).\;For all $\operatorname{nlp\_}$, $\operatorname{nls\_}$, $\operatorname{nla\_}$, $\operatorname{nlb\_}$, $\operatorname{nlm\_}$ and $\operatorname{nln\_}$, 
\begin{equation*}
\begin{array}{@{}l@{}} \operatorname{is\text{-}nso\text{-}resp\text{-}strand}(\operatorname{nlp\_}, \operatorname{nls\_}, \operatorname{nla\_}, \operatorname{nlb\_}, \operatorname{nlm\_}, \operatorname{nln\_}) \\ \Leftrightarrow\; \operatorname{nla\_} \in \operatorname{names}(\operatorname{nlp\_}) \wedge \\ \operatorname{nlb\_} \in \operatorname{names}(\operatorname{nlp\_}) \wedge \\ \operatorname{nlm\_} \in \operatorname{tokens}(\operatorname{ispace}(\operatorname{pk\text{-}inf}(\operatorname{nlp\_}))) \wedge \\ \operatorname{nln\_} \in \operatorname{tokens}(\operatorname{ispace}(\operatorname{pk\text{-}inf}(\operatorname{nlp\_}))) \wedge \\ \neg (\operatorname{nln\_} \in \operatorname{names}(\operatorname{nlp\_})) \wedge \\ \neg (\operatorname{nlm\_} \in \operatorname{keys}(\operatorname{ispace}(\operatorname{pk\text{-}inf}(\operatorname{nlp\_})))) \wedge \\ \neg (\operatorname{nln\_} \in \operatorname{keys}(\operatorname{ispace}(\operatorname{pk\text{-}inf}(\operatorname{nlp\_})))) \wedge \\ \operatorname{nls\_} \\ =\; {[}{[}\mathit{received}, \operatorname{encr}(\operatorname{lst}({[}\operatorname{nlm\_}, \operatorname{nla\_}]), \operatorname{pubkey}(\operatorname{nlp\_})(\operatorname{nlb\_}))], {[}\mathit{sent}, \operatorname{encr}(\operatorname{lst}({[}\operatorname{nlm\_}, \operatorname{nln\_}]), \operatorname{pubkey}(\operatorname{nlp\_})(\operatorname{nla\_}))], {[}\mathit{received}, \operatorname{encr}(\operatorname{nln\_}, \operatorname{pubkey}(\operatorname{nlp\_})(\operatorname{nlb\_}))]] \end{array}
\end{equation*}

\par\medskip\noindent\textbf{Definition} (\texttt{is-nso-space}).\;For every $\operatorname{nlp\_}$, 
\begin{equation*}
\begin{array}{@{}l@{}} \operatorname{is\text{-}nso\text{-}space}(\operatorname{nlp\_}) \\ \Leftrightarrow\; \operatorname{is\text{-}pk\text{-}space}(\operatorname{nlp\_}) \wedge \\ \forall \operatorname{nli\_} \in \operatorname{ids}(\operatorname{ispace}(\operatorname{pk\text{-}inf}(\operatorname{nlp\_}))). \\ \quad \operatorname{nli\_} \in \operatorname{adv\text{-}ids}(\operatorname{pk\text{-}inf}(\operatorname{nlp\_})) \vee \\ \quad \exists \operatorname{nla\_}, \operatorname{nlb\_}, \operatorname{nlm\_}, \operatorname{nln\_}. \\ \quad \quad \operatorname{is\text{-}nso\text{-}init\text{-}strand}(\operatorname{nlp\_}, \operatorname{strand\text{-}map}(\operatorname{ispace}(\operatorname{pk\text{-}inf}(\operatorname{nlp\_})))(\operatorname{nli\_}), \operatorname{nla\_}, \operatorname{nlb\_}, \operatorname{nlm\_}, \operatorname{nln\_}) \vee \\ \quad \exists \operatorname{nla\_}, \operatorname{nlb\_}, \operatorname{nlm\_}, \operatorname{nln\_}. \\ \quad \quad \operatorname{is\text{-}nso\text{-}resp\text{-}strand}(\operatorname{nlp\_}, \operatorname{strand\text{-}map}(\operatorname{ispace}(\operatorname{pk\text{-}inf}(\operatorname{nlp\_})))(\operatorname{nli\_}), \operatorname{nla\_}, \operatorname{nlb\_}, \operatorname{nlm\_}, \operatorname{nln\_}) \end{array}
\end{equation*}


\vnbbookfile{theorem-library/flat-regular-laws}{theorem-library/flat-regular-laws}
\noindent{\small flat-regular-laws.scm -- THE NSL ARGUMENT, ONCE, OVER ANY PUBLIC-KEY SPACE WHOSE REGULAR CONTENTS ARE FLAT (IS-FLAT-REGULAR, structure-library/flat-regular.scm).  Thayer, Herzog, Guttman, \textasciitilde{}/docs/jcs\_strand\_spaces.pdf, section 5: the lemmas behind Lemmas 5.3, 5.4, 5.6 (and, for the original protocol, 5.9), stated for a ...}\par
\vnbprop{flat-regular-content}{\vnbstmt{\begin{array}{@{}l@{}} \forall p, n. \\ \quad (\operatorname{is\text{-}flat\text{-}regular}(p) \wedge ((n \in \operatorname{nodes}(\operatorname{ispace}(\operatorname{pk\text{-}inf}(p)))) \wedge \operatorname{is\text{-}regular\text{-}node}(\operatorname{pk\text{-}inf}(p), n))) \Rightarrow  \\ \quad \quad \exists k \in \operatorname{keys}(\operatorname{ispace}(\operatorname{pk\text{-}inf}(p))). \\ \quad \quad \quad \exists x \in \{d \in \operatorname{tokens}(\operatorname{ispace}(\operatorname{pk\text{-}inf}(p))) : \neg (d \in \operatorname{keys}(\operatorname{ispace}(\operatorname{pk\text{-}inf}(p))))\}. \\ \quad \quad \quad \quad \operatorname{node\text{-}content}(\operatorname{ispace}(\operatorname{pk\text{-}inf}(p)), n) = \operatorname{encr}(x, k) \vee \\ \quad \quad \quad \exists l \in \operatorname{tuples}(\{d \in \operatorname{tokens}(\operatorname{ispace}(\operatorname{pk\text{-}inf}(p))) : \neg (d \in \operatorname{keys}(\operatorname{ispace}(\operatorname{pk\text{-}inf}(p))))\}). \\ \quad \quad \quad \quad \operatorname{node\text{-}content}(\operatorname{ispace}(\operatorname{pk\text{-}inf}(p)), n) = \operatorname{encr}(\operatorname{lst}(l), k) \end{array}}}
\vnbprop{flat-regular-encr-subterm}{\vnbstmt{\begin{array}{@{}l@{}} \forall p, n, g, y. \\ \quad \operatorname{is\text{-}flat\text{-}regular}(p) \wedge \\ \quad n \in \operatorname{nodes}(\operatorname{ispace}(\operatorname{pk\text{-}inf}(p))) \wedge \\ \quad \operatorname{is\text{-}regular\text{-}node}(\operatorname{pk\text{-}inf}(p), n) \wedge \\ \quad g \in \operatorname{msg\text{-}set}(\operatorname{tokens}(\operatorname{ispace}(\operatorname{pk\text{-}inf}(p))), \operatorname{keys}(\operatorname{ispace}(\operatorname{pk\text{-}inf}(p)))) \wedge \\ \quad y \in \operatorname{keys}(\operatorname{ispace}(\operatorname{pk\text{-}inf}(p))) \wedge \\ \quad \operatorname{is\text{-}subterm}(\operatorname{tokens}(\operatorname{ispace}(\operatorname{pk\text{-}inf}(p))), \operatorname{keys}(\operatorname{ispace}(\operatorname{pk\text{-}inf}(p))), \operatorname{encr}(g, y), \\ \quad \quad \operatorname{node\text{-}content}(\operatorname{ispace}(\operatorname{pk\text{-}inf}(p)), n)) \Rightarrow \\ \quad \quad \operatorname{encr}(g, y) = \operatorname{node\text{-}content}(\operatorname{ispace}(\operatorname{pk\text{-}inf}(p)), n) \end{array}}}
\vnbprop{flat-regular-list-subterm}{\vnbstmt{\begin{array}{@{}l@{}} \forall p, n, m. \\ \quad \operatorname{is\text{-}flat\text{-}regular}(p) \wedge \\ \quad n \in \operatorname{nodes}(\operatorname{ispace}(\operatorname{pk\text{-}inf}(p))) \wedge \\ \quad \operatorname{is\text{-}regular\text{-}node}(\operatorname{pk\text{-}inf}(p), n) \wedge \\ \quad m \in \operatorname{tuples}(\operatorname{msg\text{-}set}(\operatorname{tokens}(\operatorname{ispace}(\operatorname{pk\text{-}inf}(p))), \operatorname{keys}(\operatorname{ispace}(\operatorname{pk\text{-}inf}(p))))) \wedge \\ \quad \operatorname{is\text{-}subterm}(\operatorname{tokens}(\operatorname{ispace}(\operatorname{pk\text{-}inf}(p))), \operatorname{keys}(\operatorname{ispace}(\operatorname{pk\text{-}inf}(p))), \operatorname{lst}(m), \\ \quad \quad \operatorname{node\text{-}content}(\operatorname{ispace}(\operatorname{pk\text{-}inf}(p)), n)) \Rightarrow \\ \quad \quad m \in \operatorname{tuples}(\operatorname{tokens}(\operatorname{ispace}(\operatorname{pk\text{-}inf}(p)))) \end{array}}}
\vnbprop{flat-regular-no-key}{\vnbstmt{\begin{array}{@{}l@{}} \forall p, n, y. \\ \quad \operatorname{is\text{-}flat\text{-}regular}(p) \wedge \\ \quad n \in \operatorname{nodes}(\operatorname{ispace}(\operatorname{pk\text{-}inf}(p))) \wedge \\ \quad \operatorname{is\text{-}regular\text{-}node}(\operatorname{pk\text{-}inf}(p), n) \wedge \\ \quad y \in \operatorname{keys}(\operatorname{ispace}(\operatorname{pk\text{-}inf}(p))) \Rightarrow \\ \quad \quad \neg \operatorname{is\text{-}subterm}(\operatorname{tokens}(\operatorname{ispace}(\operatorname{pk\text{-}inf}(p))), \operatorname{keys}(\operatorname{ispace}(\operatorname{pk\text{-}inf}(p))), y, \operatorname{node\text{-}content}(\operatorname{ispace}(\operatorname{pk\text{-}inf}(p)), n)) \end{array}}}
\vnbprop{flat-key-secret}{\vnbstmt{\begin{array}{@{}l@{}} \forall p, c, y.; \forall n \in \operatorname{bnodes}(c). \\ \quad \operatorname{is\text{-}flat\text{-}regular}(p) \wedge \\ \quad \operatorname{is\text{-}bundle}(c) \wedge \\ \quad \operatorname{bspace}(c) = \operatorname{ispace}(\operatorname{pk\text{-}inf}(p)) \wedge \\ \quad y \in \operatorname{keys}(\operatorname{ispace}(\operatorname{pk\text{-}inf}(p))) \wedge \\ \quad \neg (y \in \operatorname{adv\text{-}tokens}(\operatorname{pk\text{-}inf}(p))) \Rightarrow \\ \quad \quad \neg \operatorname{is\text{-}subterm}(\operatorname{tokens}(\operatorname{ispace}(\operatorname{pk\text{-}inf}(p))), \operatorname{keys}(\operatorname{ispace}(\operatorname{pk\text{-}inf}(p))), y, \operatorname{node\text{-}content}(\operatorname{ispace}(\operatorname{pk\text{-}inf}(p)), n)) \end{array}}}
\vnbprop{flat-list-source}{\vnbstmt{\begin{array}{@{}l@{}} \forall p, c, a, l, q. \\ \quad \operatorname{is\text{-}flat\text{-}regular}(p) \wedge \\ \quad \operatorname{is\text{-}bundle}(c) \wedge \\ \quad \operatorname{bspace}(c) = \operatorname{ispace}(\operatorname{pk\text{-}inf}(p)) \wedge \\ \quad a \in \operatorname{bnodes}(c) \wedge \\ \quad l \in \operatorname{tuples}(\operatorname{msg\text{-}set}(\operatorname{tokens}(\operatorname{ispace}(\operatorname{pk\text{-}inf}(p))), \operatorname{keys}(\operatorname{ispace}(\operatorname{pk\text{-}inf}(p))))) \wedge \\ \quad q \in \mathbb{N} \wedge \\ \quad 1 \leq q \wedge \\ \quad q \leq \operatorname{length}(l) \wedge \\ \quad \operatorname{node\text{-}content}(\operatorname{ispace}(\operatorname{pk\text{-}inf}(p)), a) = \operatorname{nth}(q, l) \wedge \\ \quad \neg (l \in \operatorname{tuples}(\operatorname{tokens}(\operatorname{ispace}(\operatorname{pk\text{-}inf}(p))))) \wedge \\ \quad \exists x \in \operatorname{bnodes}(c). \\ \quad \quad {[}x, a] \in \operatorname{precedes\text{-}strict}(c) \wedge \\ \quad \quad \operatorname{is\text{-}subterm}(\operatorname{tokens}(\operatorname{ispace}(\operatorname{pk\text{-}inf}(p))), \operatorname{keys}(\operatorname{ispace}(\operatorname{pk\text{-}inf}(p))), \operatorname{lst}(l), \\ \quad \quad \quad \operatorname{node\text{-}content}(\operatorname{ispace}(\operatorname{pk\text{-}inf}(p)), x)) \Rightarrow \\ \quad \quad \exists z \in \operatorname{bnodes}(c). \\ \quad \quad \quad {[}z, a] \in \operatorname{precedes\text{-}strict}(c) \wedge \\ \quad \quad \quad \operatorname{node\text{-}content}(\operatorname{ispace}(\operatorname{pk\text{-}inf}(p)), z) = \operatorname{node\text{-}content}(\operatorname{ispace}(\operatorname{pk\text{-}inf}(p)), a) \end{array}}}
\vnbprop{flat-adversary-loss}{\vnbstmt{\begin{array}{@{}l@{}} \forall p, c, m, x, g, y. \\ \quad \operatorname{is\text{-}flat\text{-}regular}(p) \wedge \\ \quad \operatorname{is\text{-}bundle}(c) \wedge \\ \quad \operatorname{bspace}(c) = \operatorname{ispace}(\operatorname{pk\text{-}inf}(p)) \wedge \\ \quad m \in \operatorname{bnodes}(c) \wedge \\ \quad \operatorname{is\text{-}adversary\text{-}node}(\operatorname{pk\text{-}inf}(p), m) \wedge \\ \quad \operatorname{is\text{-}sent}(\operatorname{ispace}(\operatorname{pk\text{-}inf}(p)), m) \wedge \\ \quad \operatorname{same\text{-}strand\text{-}before}(\operatorname{ispace}(\operatorname{pk\text{-}inf}(p)), x, m) \wedge \\ \quad g \in \operatorname{msg\text{-}set}(\operatorname{tokens}(\operatorname{ispace}(\operatorname{pk\text{-}inf}(p))), \operatorname{keys}(\operatorname{ispace}(\operatorname{pk\text{-}inf}(p)))) \wedge \\ \quad y \in \operatorname{keys}(\operatorname{ispace}(\operatorname{pk\text{-}inf}(p))) \wedge \\ \quad \neg (\operatorname{kinv}(\operatorname{ispace}(\operatorname{pk\text{-}inf}(p)))(y) \in \operatorname{adv\text{-}tokens}(\operatorname{pk\text{-}inf}(p))) \wedge \\ \quad \operatorname{is\text{-}subterm}(\operatorname{tokens}(\operatorname{ispace}(\operatorname{pk\text{-}inf}(p))), \operatorname{keys}(\operatorname{ispace}(\operatorname{pk\text{-}inf}(p))), \operatorname{encr}(g, y), \\ \quad \quad \operatorname{node\text{-}content}(\operatorname{ispace}(\operatorname{pk\text{-}inf}(p)), x)) \wedge \\ \quad \neg \operatorname{is\text{-}subterm}(\operatorname{tokens}(\operatorname{ispace}(\operatorname{pk\text{-}inf}(p))), \operatorname{keys}(\operatorname{ispace}(\operatorname{pk\text{-}inf}(p))), \operatorname{encr}(g, y), \operatorname{node\text{-}content}(\operatorname{ispace}(\operatorname{pk\text{-}inf}(p)), m)) \Rightarrow \\ \quad \quad \exists z \in \operatorname{bnodes}(c). \\ \quad \quad \quad {[}z, m] \in \operatorname{precedes\text{-}strict}(c) \wedge \\ \quad \quad \quad \operatorname{node\text{-}content}(\operatorname{ispace}(\operatorname{pk\text{-}inf}(p)), z) = \operatorname{node\text{-}content}(\operatorname{ispace}(\operatorname{pk\text{-}inf}(p)), m) \end{array}}}
\vnbprop{resp-trace-nodes}{\vnbstmt{\begin{array}{@{}l@{}} \forall p, a, b, c, d. \\ \quad \operatorname{is\text{-}pk\text{-}space}(p) \wedge \\ \quad a \in \operatorname{ids}(\operatorname{ispace}(\operatorname{pk\text{-}inf}(p))) \wedge \\ \quad \operatorname{strand\text{-}map}(\operatorname{ispace}(\operatorname{pk\text{-}inf}(p)))(a) = {[}{[}\mathit{received}, b], {[}\mathit{sent}, c], {[}\mathit{received}, d]] \Rightarrow \\ \quad \quad {[}1, a] \in \operatorname{nodes}(\operatorname{ispace}(\operatorname{pk\text{-}inf}(p))) \wedge \\ \quad \quad \operatorname{is\text{-}received}(\operatorname{ispace}(\operatorname{pk\text{-}inf}(p)), {[}1, a]) \wedge \\ \quad \quad \operatorname{node\text{-}content}(\operatorname{ispace}(\operatorname{pk\text{-}inf}(p)), {[}1, a]) = b \wedge \\ \quad \quad {[}2, a] \in \operatorname{nodes}(\operatorname{ispace}(\operatorname{pk\text{-}inf}(p))) \wedge \\ \quad \quad \operatorname{is\text{-}sent}(\operatorname{ispace}(\operatorname{pk\text{-}inf}(p)), {[}2, a]) \wedge \\ \quad \quad \operatorname{node\text{-}content}(\operatorname{ispace}(\operatorname{pk\text{-}inf}(p)), {[}2, a]) = c \wedge \\ \quad \quad {[}3, a] \in \operatorname{nodes}(\operatorname{ispace}(\operatorname{pk\text{-}inf}(p))) \wedge \\ \quad \quad \operatorname{is\text{-}received}(\operatorname{ispace}(\operatorname{pk\text{-}inf}(p)), {[}3, a]) \wedge \\ \quad \quad \operatorname{node\text{-}content}(\operatorname{ispace}(\operatorname{pk\text{-}inf}(p)), {[}3, a]) = d \end{array}}}
\vnbprop{flat-resp-nonce-originates}{\vnbstmt{\begin{array}{@{}l@{}} \forall p, s, a, b, c, d, f. \\ \quad \operatorname{is\text{-}pk\text{-}space}(p) \wedge \\ \quad s \in \operatorname{ids}(\operatorname{ispace}(\operatorname{pk\text{-}inf}(p))) \wedge \\ \quad \operatorname{strand\text{-}map}(\operatorname{ispace}(\operatorname{pk\text{-}inf}(p)))(s) \\ \quad =\; {[}{[}\mathit{received}, a], {[}\mathit{sent}, \operatorname{encr}(\operatorname{lst}(b), c)], {[}\mathit{received}, \operatorname{encr}(f, d)]] \wedge \\ \quad b \in \operatorname{tuples}(\operatorname{msg\text{-}set}(\operatorname{tokens}(\operatorname{ispace}(\operatorname{pk\text{-}inf}(p))), \operatorname{keys}(\operatorname{ispace}(\operatorname{pk\text{-}inf}(p))))) \wedge \\ \quad c \in \operatorname{keys}(\operatorname{ispace}(\operatorname{pk\text{-}inf}(p))) \wedge \\ \quad d \in \operatorname{keys}(\operatorname{ispace}(\operatorname{pk\text{-}inf}(p))) \wedge \\ \quad f \in \operatorname{tokens}(\operatorname{ispace}(\operatorname{pk\text{-}inf}(p))) \wedge \\ \quad \operatorname{is\text{-}subterm}(\operatorname{tokens}(\operatorname{ispace}(\operatorname{pk\text{-}inf}(p))), \operatorname{keys}(\operatorname{ispace}(\operatorname{pk\text{-}inf}(p))), f, \operatorname{lst}(b)) \wedge \\ \quad \neg \operatorname{is\text{-}subterm}(\operatorname{tokens}(\operatorname{ispace}(\operatorname{pk\text{-}inf}(p))), \operatorname{keys}(\operatorname{ispace}(\operatorname{pk\text{-}inf}(p))), f, a) \Rightarrow \\ \quad \quad \operatorname{originates}(\operatorname{ispace}(\operatorname{pk\text{-}inf}(p)), f, {[}2, s]) \end{array}}}
\vnbprop{flat-lemma-5-4}{\vnbstmt{\begin{array}{@{}l@{}} \forall p, c, s, a, b, d, f, g. \\ \quad \operatorname{is\text{-}flat\text{-}regular}(p) \wedge \\ \quad \operatorname{is\text{-}bundle}(c) \wedge \\ \quad \operatorname{bspace}(c) = \operatorname{ispace}(\operatorname{pk\text{-}inf}(p)) \wedge \\ \quad s \in \operatorname{ids}(\operatorname{ispace}(\operatorname{pk\text{-}inf}(p))) \wedge \\ \quad \operatorname{strand\text{-}map}(\operatorname{ispace}(\operatorname{pk\text{-}inf}(p)))(s) \\ \quad =\; {[}{[}\mathit{received}, a], {[}\mathit{sent}, \operatorname{encr}(\operatorname{lst}(b), d)], {[}\mathit{received}, \operatorname{encr}(g, f)]] \wedge \\ \quad {[}3, s] \in \operatorname{bnodes}(c) \wedge \\ \quad \neg (\operatorname{kinv}(\operatorname{ispace}(\operatorname{pk\text{-}inf}(p)))(d) \in \operatorname{adv\text{-}tokens}(\operatorname{pk\text{-}inf}(p))) \wedge \\ \quad b \in \operatorname{tuples}(\operatorname{msg\text{-}set}(\operatorname{tokens}(\operatorname{ispace}(\operatorname{pk\text{-}inf}(p))), \operatorname{keys}(\operatorname{ispace}(\operatorname{pk\text{-}inf}(p))))) \wedge \\ \quad d \in \operatorname{keys}(\operatorname{ispace}(\operatorname{pk\text{-}inf}(p))) \wedge \\ \quad f \in \operatorname{keys}(\operatorname{ispace}(\operatorname{pk\text{-}inf}(p))) \wedge \\ \quad g \in \operatorname{tokens}(\operatorname{ispace}(\operatorname{pk\text{-}inf}(p))) \wedge \\ \quad \neg (g \in \operatorname{keys}(\operatorname{ispace}(\operatorname{pk\text{-}inf}(p)))) \wedge \\ \quad \operatorname{is\text{-}subterm}(\operatorname{tokens}(\operatorname{ispace}(\operatorname{pk\text{-}inf}(p))), \operatorname{keys}(\operatorname{ispace}(\operatorname{pk\text{-}inf}(p))), g, \operatorname{lst}(b)) \wedge \\ \quad \neg \operatorname{is\text{-}subterm}(\operatorname{tokens}(\operatorname{ispace}(\operatorname{pk\text{-}inf}(p))), \operatorname{keys}(\operatorname{ispace}(\operatorname{pk\text{-}inf}(p))), g, a) \wedge \\ \quad \operatorname{uniquely\text{-}originates}(\operatorname{ispace}(\operatorname{pk\text{-}inf}(p)), g) \Rightarrow \\ \quad \quad \exists h \in \operatorname{bnodes}(c). \\ \quad \quad \quad \operatorname{is\text{-}subterm}(\operatorname{tokens}(\operatorname{ispace}(\operatorname{pk\text{-}inf}(p))), \operatorname{keys}(\operatorname{ispace}(\operatorname{pk\text{-}inf}(p))), g, \\ \quad \quad \quad \quad \operatorname{node\text{-}content}(\operatorname{ispace}(\operatorname{pk\text{-}inf}(p)), h)) \wedge \\ \quad \quad \quad \neg \operatorname{is\text{-}subterm}(\operatorname{tokens}(\operatorname{ispace}(\operatorname{pk\text{-}inf}(p))), \operatorname{keys}(\operatorname{ispace}(\operatorname{pk\text{-}inf}(p))), \operatorname{encr}(\operatorname{lst}(b), d), \operatorname{node\text{-}content}(\operatorname{ispace}(\operatorname{pk\text{-}inf}(p)), h)) \wedge \\ \quad \quad \quad \forall z \in \operatorname{bnodes}(c). \\ \quad \quad \quad \quad \operatorname{is\text{-}subterm}(\operatorname{tokens}(\operatorname{ispace}(\operatorname{pk\text{-}inf}(p))), \operatorname{keys}(\operatorname{ispace}(\operatorname{pk\text{-}inf}(p))), g, \\ \quad \quad \quad \quad \quad \operatorname{node\text{-}content}(\operatorname{ispace}(\operatorname{pk\text{-}inf}(p)), z)) \wedge \\ \quad \quad \quad \quad \neg \operatorname{is\text{-}subterm}(\operatorname{tokens}(\operatorname{ispace}(\operatorname{pk\text{-}inf}(p))), \operatorname{keys}(\operatorname{ispace}(\operatorname{pk\text{-}inf}(p))), \operatorname{encr}(\operatorname{lst}(b), d), \operatorname{node\text{-}content}(\operatorname{ispace}(\operatorname{pk\text{-}inf}(p)), z)) \Rightarrow \\ \quad \quad \quad \quad \quad \neg ({[}z, h] \in \operatorname{precedes\text{-}strict}(c)) \wedge \\ \quad \quad \quad \operatorname{is\text{-}regular\text{-}node}(\operatorname{pk\text{-}inf}(p), h) \wedge \\ \quad \quad \quad \operatorname{is\text{-}sent}(\operatorname{ispace}(\operatorname{pk\text{-}inf}(p)), h) \end{array}}}
\vnbprop{flat-lemma-5-6}{\vnbstmt{\begin{array}{@{}l@{}} \forall p, c, s, a, b, d, f, g.; \forall h \in \operatorname{bnodes}(c). \\ \quad \operatorname{is\text{-}flat\text{-}regular}(p) \wedge \\ \quad \operatorname{is\text{-}bundle}(c) \wedge \\ \quad \operatorname{bspace}(c) = \operatorname{ispace}(\operatorname{pk\text{-}inf}(p)) \wedge \\ \quad s \in \operatorname{ids}(\operatorname{ispace}(\operatorname{pk\text{-}inf}(p))) \wedge \\ \quad \operatorname{strand\text{-}map}(\operatorname{ispace}(\operatorname{pk\text{-}inf}(p)))(s) \\ \quad =\; {[}{[}\mathit{received}, a], {[}\mathit{sent}, \operatorname{encr}(\operatorname{lst}(b), d)], {[}\mathit{received}, \operatorname{encr}(g, f)]] \wedge \\ \quad {[}3, s] \in \operatorname{bnodes}(c) \wedge \\ \quad \neg (\operatorname{kinv}(\operatorname{ispace}(\operatorname{pk\text{-}inf}(p)))(d) \in \operatorname{adv\text{-}tokens}(\operatorname{pk\text{-}inf}(p))) \wedge \\ \quad b \in \operatorname{tuples}(\operatorname{msg\text{-}set}(\operatorname{tokens}(\operatorname{ispace}(\operatorname{pk\text{-}inf}(p))), \operatorname{keys}(\operatorname{ispace}(\operatorname{pk\text{-}inf}(p))))) \wedge \\ \quad d \in \operatorname{keys}(\operatorname{ispace}(\operatorname{pk\text{-}inf}(p))) \wedge \\ \quad f \in \operatorname{keys}(\operatorname{ispace}(\operatorname{pk\text{-}inf}(p))) \wedge \\ \quad g \in \operatorname{tokens}(\operatorname{ispace}(\operatorname{pk\text{-}inf}(p))) \wedge \\ \quad \neg (g \in \operatorname{keys}(\operatorname{ispace}(\operatorname{pk\text{-}inf}(p)))) \wedge \\ \quad \operatorname{is\text{-}subterm}(\operatorname{tokens}(\operatorname{ispace}(\operatorname{pk\text{-}inf}(p))), \operatorname{keys}(\operatorname{ispace}(\operatorname{pk\text{-}inf}(p))), g, \operatorname{lst}(b)) \wedge \\ \quad \neg \operatorname{is\text{-}subterm}(\operatorname{tokens}(\operatorname{ispace}(\operatorname{pk\text{-}inf}(p))), \operatorname{keys}(\operatorname{ispace}(\operatorname{pk\text{-}inf}(p))), g, a) \wedge \\ \quad \operatorname{uniquely\text{-}originates}(\operatorname{ispace}(\operatorname{pk\text{-}inf}(p)), g) \wedge \\ \quad \operatorname{is\text{-}subterm}(\operatorname{tokens}(\operatorname{ispace}(\operatorname{pk\text{-}inf}(p))), \operatorname{keys}(\operatorname{ispace}(\operatorname{pk\text{-}inf}(p))), g, \\ \quad \quad \operatorname{node\text{-}content}(\operatorname{ispace}(\operatorname{pk\text{-}inf}(p)), h)) \wedge \\ \quad \neg \operatorname{is\text{-}subterm}(\operatorname{tokens}(\operatorname{ispace}(\operatorname{pk\text{-}inf}(p))), \operatorname{keys}(\operatorname{ispace}(\operatorname{pk\text{-}inf}(p))), \operatorname{encr}(\operatorname{lst}(b), d), \operatorname{node\text{-}content}(\operatorname{ispace}(\operatorname{pk\text{-}inf}(p)), h)) \wedge \\ \quad \forall z \in \operatorname{bnodes}(c). \\ \quad \quad \operatorname{is\text{-}subterm}(\operatorname{tokens}(\operatorname{ispace}(\operatorname{pk\text{-}inf}(p))), \operatorname{keys}(\operatorname{ispace}(\operatorname{pk\text{-}inf}(p))), g, \\ \quad \quad \quad \operatorname{node\text{-}content}(\operatorname{ispace}(\operatorname{pk\text{-}inf}(p)), z)) \wedge \\ \quad \quad \neg \operatorname{is\text{-}subterm}(\operatorname{tokens}(\operatorname{ispace}(\operatorname{pk\text{-}inf}(p))), \operatorname{keys}(\operatorname{ispace}(\operatorname{pk\text{-}inf}(p))), \operatorname{encr}(\operatorname{lst}(b), d), \operatorname{node\text{-}content}(\operatorname{ispace}(\operatorname{pk\text{-}inf}(p)), z)) \Rightarrow \\ \quad \quad \quad \neg ({[}z, h] \in \operatorname{precedes\text{-}strict}(c)) \wedge \\ \quad \operatorname{is\text{-}regular\text{-}node}(\operatorname{pk\text{-}inf}(p), h) \wedge \\ \quad \operatorname{is\text{-}sent}(\operatorname{ispace}(\operatorname{pk\text{-}inf}(p)), h) \Rightarrow \\ \quad \quad \exists k. \\ \quad \quad \quad (\operatorname{same\text{-}strand\text{-}before}(\operatorname{ispace}(\operatorname{pk\text{-}inf}(p)), k, h) \wedge (\operatorname{node\text{-}content}(\operatorname{ispace}(\operatorname{pk\text{-}inf}(p)), k) = \operatorname{encr}(\operatorname{lst}(b), d))) \end{array}}}

\vnbbookfile{theorem-library/nso-laws}{theorem-library/nso-laws}
\noindent{\small nso-laws.scm -- THE ORIGINAL NEEDHAM-SCHROEDER PROTOCOL: Lemma 5.9 of Thayer, Herzog, Guttman (\textasciitilde{}/docs/jcs\_strand\_spaces.pdf, section 5.3): "In the original Needham-Schroeder protocol, a node n1 precedes n2 on the same regular strand t, and term(n1) = \{Na Nb\}\_KA."  Version 2 (agent SB-2, 2026-10-26): the theorems are ...}\par
\vnbprop{nso-regular-content}{\vnbstmt{\begin{array}{@{}l@{}} \forall p, n. \\ \quad (\operatorname{is\text{-}nso\text{-}space}(p) \wedge ((n \in \operatorname{nodes}(\operatorname{ispace}(\operatorname{pk\text{-}inf}(p)))) \wedge \operatorname{is\text{-}regular\text{-}node}(\operatorname{pk\text{-}inf}(p), n))) \Rightarrow  \\ \quad \quad \exists k \in \operatorname{keys}(\operatorname{ispace}(\operatorname{pk\text{-}inf}(p))). \\ \quad \quad \quad \exists x \in \{d \in \operatorname{tokens}(\operatorname{ispace}(\operatorname{pk\text{-}inf}(p))) : \neg (d \in \operatorname{keys}(\operatorname{ispace}(\operatorname{pk\text{-}inf}(p))))\}. \\ \quad \quad \quad \quad \operatorname{node\text{-}content}(\operatorname{ispace}(\operatorname{pk\text{-}inf}(p)), n) = \operatorname{encr}(x, k) \vee \\ \quad \quad \quad \exists l \in \operatorname{tuples}(\{d \in \operatorname{tokens}(\operatorname{ispace}(\operatorname{pk\text{-}inf}(p))) : \neg (d \in \operatorname{keys}(\operatorname{ispace}(\operatorname{pk\text{-}inf}(p))))\}). \\ \quad \quad \quad \quad \operatorname{node\text{-}content}(\operatorname{ispace}(\operatorname{pk\text{-}inf}(p)), n) = \operatorname{encr}(\operatorname{lst}(l), k) \end{array}}}
\vnbprop{nso-flat}{\vnbstmt{\forall p.\; (\operatorname{is\text{-}nso\text{-}space}(p) \Rightarrow \operatorname{is\text{-}flat\text{-}regular}(p))}}
\vnbprop{nso-regular-encr-subterm}{\vnbstmt{\begin{array}{@{}l@{}} \forall p, n, g, y. \\ \quad \operatorname{is\text{-}nso\text{-}space}(p) \wedge \\ \quad n \in \operatorname{nodes}(\operatorname{ispace}(\operatorname{pk\text{-}inf}(p))) \wedge \\ \quad \operatorname{is\text{-}regular\text{-}node}(\operatorname{pk\text{-}inf}(p), n) \wedge \\ \quad g \in \operatorname{msg\text{-}set}(\operatorname{tokens}(\operatorname{ispace}(\operatorname{pk\text{-}inf}(p))), \operatorname{keys}(\operatorname{ispace}(\operatorname{pk\text{-}inf}(p)))) \wedge \\ \quad y \in \operatorname{keys}(\operatorname{ispace}(\operatorname{pk\text{-}inf}(p))) \wedge \\ \quad \operatorname{is\text{-}subterm}(\operatorname{tokens}(\operatorname{ispace}(\operatorname{pk\text{-}inf}(p))), \operatorname{keys}(\operatorname{ispace}(\operatorname{pk\text{-}inf}(p))), \operatorname{encr}(g, y), \\ \quad \quad \operatorname{node\text{-}content}(\operatorname{ispace}(\operatorname{pk\text{-}inf}(p)), n)) \Rightarrow \\ \quad \quad \operatorname{encr}(g, y) = \operatorname{node\text{-}content}(\operatorname{ispace}(\operatorname{pk\text{-}inf}(p)), n) \end{array}}}
\vnbprop{nso-regular-list-subterm}{\vnbstmt{\begin{array}{@{}l@{}} \forall p, n, m. \\ \quad \operatorname{is\text{-}nso\text{-}space}(p) \wedge \\ \quad n \in \operatorname{nodes}(\operatorname{ispace}(\operatorname{pk\text{-}inf}(p))) \wedge \\ \quad \operatorname{is\text{-}regular\text{-}node}(\operatorname{pk\text{-}inf}(p), n) \wedge \\ \quad m \in \operatorname{tuples}(\operatorname{msg\text{-}set}(\operatorname{tokens}(\operatorname{ispace}(\operatorname{pk\text{-}inf}(p))), \operatorname{keys}(\operatorname{ispace}(\operatorname{pk\text{-}inf}(p))))) \wedge \\ \quad \operatorname{is\text{-}subterm}(\operatorname{tokens}(\operatorname{ispace}(\operatorname{pk\text{-}inf}(p))), \operatorname{keys}(\operatorname{ispace}(\operatorname{pk\text{-}inf}(p))), \operatorname{lst}(m), \\ \quad \quad \operatorname{node\text{-}content}(\operatorname{ispace}(\operatorname{pk\text{-}inf}(p)), n)) \Rightarrow \\ \quad \quad m \in \operatorname{tuples}(\operatorname{tokens}(\operatorname{ispace}(\operatorname{pk\text{-}inf}(p)))) \end{array}}}
\vnbprop{nso-regular-no-key}{\vnbstmt{\begin{array}{@{}l@{}} \forall p, n, y. \\ \quad \operatorname{is\text{-}nso\text{-}space}(p) \wedge \\ \quad n \in \operatorname{nodes}(\operatorname{ispace}(\operatorname{pk\text{-}inf}(p))) \wedge \\ \quad \operatorname{is\text{-}regular\text{-}node}(\operatorname{pk\text{-}inf}(p), n) \wedge \\ \quad y \in \operatorname{keys}(\operatorname{ispace}(\operatorname{pk\text{-}inf}(p))) \Rightarrow \\ \quad \quad \neg \operatorname{is\text{-}subterm}(\operatorname{tokens}(\operatorname{ispace}(\operatorname{pk\text{-}inf}(p))), \operatorname{keys}(\operatorname{ispace}(\operatorname{pk\text{-}inf}(p))), y, \operatorname{node\text{-}content}(\operatorname{ispace}(\operatorname{pk\text{-}inf}(p)), n)) \end{array}}}
\vnbprop{nso-key-secret}{\vnbstmt{\begin{array}{@{}l@{}} \forall p, c, y.; \forall n \in \operatorname{bnodes}(c). \\ \quad \operatorname{is\text{-}nso\text{-}space}(p) \wedge \\ \quad \operatorname{is\text{-}bundle}(c) \wedge \\ \quad \operatorname{bspace}(c) = \operatorname{ispace}(\operatorname{pk\text{-}inf}(p)) \wedge \\ \quad y \in \operatorname{keys}(\operatorname{ispace}(\operatorname{pk\text{-}inf}(p))) \wedge \\ \quad \neg (y \in \operatorname{adv\text{-}tokens}(\operatorname{pk\text{-}inf}(p))) \Rightarrow \\ \quad \quad \neg \operatorname{is\text{-}subterm}(\operatorname{tokens}(\operatorname{ispace}(\operatorname{pk\text{-}inf}(p))), \operatorname{keys}(\operatorname{ispace}(\operatorname{pk\text{-}inf}(p))), y, \operatorname{node\text{-}content}(\operatorname{ispace}(\operatorname{pk\text{-}inf}(p)), n)) \end{array}}}
\vnbprop{nso-list-source}{\vnbstmt{\begin{array}{@{}l@{}} \forall p, c, a, l, q. \\ \quad \operatorname{is\text{-}nso\text{-}space}(p) \wedge \\ \quad \operatorname{is\text{-}bundle}(c) \wedge \\ \quad \operatorname{bspace}(c) = \operatorname{ispace}(\operatorname{pk\text{-}inf}(p)) \wedge \\ \quad a \in \operatorname{bnodes}(c) \wedge \\ \quad l \in \operatorname{tuples}(\operatorname{msg\text{-}set}(\operatorname{tokens}(\operatorname{ispace}(\operatorname{pk\text{-}inf}(p))), \operatorname{keys}(\operatorname{ispace}(\operatorname{pk\text{-}inf}(p))))) \wedge \\ \quad q \in \mathbb{N} \wedge \\ \quad 1 \leq q \wedge \\ \quad q \leq \operatorname{length}(l) \wedge \\ \quad \operatorname{node\text{-}content}(\operatorname{ispace}(\operatorname{pk\text{-}inf}(p)), a) = \operatorname{nth}(q, l) \wedge \\ \quad \neg (l \in \operatorname{tuples}(\operatorname{tokens}(\operatorname{ispace}(\operatorname{pk\text{-}inf}(p))))) \wedge \\ \quad \exists x \in \operatorname{bnodes}(c). \\ \quad \quad {[}x, a] \in \operatorname{precedes\text{-}strict}(c) \wedge \\ \quad \quad \operatorname{is\text{-}subterm}(\operatorname{tokens}(\operatorname{ispace}(\operatorname{pk\text{-}inf}(p))), \operatorname{keys}(\operatorname{ispace}(\operatorname{pk\text{-}inf}(p))), \operatorname{lst}(l), \\ \quad \quad \quad \operatorname{node\text{-}content}(\operatorname{ispace}(\operatorname{pk\text{-}inf}(p)), x)) \Rightarrow \\ \quad \quad \exists z \in \operatorname{bnodes}(c). \\ \quad \quad \quad {[}z, a] \in \operatorname{precedes\text{-}strict}(c) \wedge \\ \quad \quad \quad \operatorname{node\text{-}content}(\operatorname{ispace}(\operatorname{pk\text{-}inf}(p)), z) = \operatorname{node\text{-}content}(\operatorname{ispace}(\operatorname{pk\text{-}inf}(p)), a) \end{array}}}
\vnbprop{nso-adversary-loss}{\vnbstmt{\begin{array}{@{}l@{}} \forall p, c, m, x, g, y. \\ \quad \operatorname{is\text{-}nso\text{-}space}(p) \wedge \\ \quad \operatorname{is\text{-}bundle}(c) \wedge \\ \quad \operatorname{bspace}(c) = \operatorname{ispace}(\operatorname{pk\text{-}inf}(p)) \wedge \\ \quad m \in \operatorname{bnodes}(c) \wedge \\ \quad \operatorname{is\text{-}adversary\text{-}node}(\operatorname{pk\text{-}inf}(p), m) \wedge \\ \quad \operatorname{is\text{-}sent}(\operatorname{ispace}(\operatorname{pk\text{-}inf}(p)), m) \wedge \\ \quad \operatorname{same\text{-}strand\text{-}before}(\operatorname{ispace}(\operatorname{pk\text{-}inf}(p)), x, m) \wedge \\ \quad g \in \operatorname{msg\text{-}set}(\operatorname{tokens}(\operatorname{ispace}(\operatorname{pk\text{-}inf}(p))), \operatorname{keys}(\operatorname{ispace}(\operatorname{pk\text{-}inf}(p)))) \wedge \\ \quad y \in \operatorname{keys}(\operatorname{ispace}(\operatorname{pk\text{-}inf}(p))) \wedge \\ \quad \neg (\operatorname{kinv}(\operatorname{ispace}(\operatorname{pk\text{-}inf}(p)))(y) \in \operatorname{adv\text{-}tokens}(\operatorname{pk\text{-}inf}(p))) \wedge \\ \quad \operatorname{is\text{-}subterm}(\operatorname{tokens}(\operatorname{ispace}(\operatorname{pk\text{-}inf}(p))), \operatorname{keys}(\operatorname{ispace}(\operatorname{pk\text{-}inf}(p))), \operatorname{encr}(g, y), \\ \quad \quad \operatorname{node\text{-}content}(\operatorname{ispace}(\operatorname{pk\text{-}inf}(p)), x)) \wedge \\ \quad \neg \operatorname{is\text{-}subterm}(\operatorname{tokens}(\operatorname{ispace}(\operatorname{pk\text{-}inf}(p))), \operatorname{keys}(\operatorname{ispace}(\operatorname{pk\text{-}inf}(p))), \operatorname{encr}(g, y), \operatorname{node\text{-}content}(\operatorname{ispace}(\operatorname{pk\text{-}inf}(p)), m)) \Rightarrow \\ \quad \quad \exists z \in \operatorname{bnodes}(c). \\ \quad \quad \quad {[}z, m] \in \operatorname{precedes\text{-}strict}(c) \wedge \\ \quad \quad \quad \operatorname{node\text{-}content}(\operatorname{ispace}(\operatorname{pk\text{-}inf}(p)), z) = \operatorname{node\text{-}content}(\operatorname{ispace}(\operatorname{pk\text{-}inf}(p)), m) \end{array}}}
\vnbprop{nso-resp-nonce-originates}{\vnbstmt{\begin{array}{@{}l@{}} \forall p, s, a, b, c, d. \\ \quad \operatorname{is\text{-}pk\text{-}space}(p) \wedge \\ \quad s \in \operatorname{ids}(\operatorname{ispace}(\operatorname{pk\text{-}inf}(p))) \wedge \\ \quad \operatorname{is\text{-}nso\text{-}resp\text{-}strand}(p, \operatorname{strand\text{-}map}(\operatorname{ispace}(\operatorname{pk\text{-}inf}(p)))(s), a, b, c, d) \wedge \\ \quad \neg (c = d) \Rightarrow \\ \quad \quad \operatorname{originates}(\operatorname{ispace}(\operatorname{pk\text{-}inf}(p)), d, {[}2, s]) \end{array}}}
\vnbprop{nso-lemma-5-4}{\vnbstmt{\begin{array}{@{}l@{}} \forall p, c, s, a, b, d, f. \\ \quad \operatorname{is\text{-}nso\text{-}space}(p) \wedge \\ \quad \operatorname{is\text{-}bundle}(c) \wedge \\ \quad \operatorname{bspace}(c) = \operatorname{ispace}(\operatorname{pk\text{-}inf}(p)) \wedge \\ \quad s \in \operatorname{ids}(\operatorname{ispace}(\operatorname{pk\text{-}inf}(p))) \wedge \\ \quad \operatorname{is\text{-}nso\text{-}resp\text{-}strand}(p, \operatorname{strand\text{-}map}(\operatorname{ispace}(\operatorname{pk\text{-}inf}(p)))(s), a, b, d, f) \wedge \\ \quad {[}3, s] \in \operatorname{bnodes}(c) \wedge \\ \quad \neg (\operatorname{kinv}(\operatorname{ispace}(\operatorname{pk\text{-}inf}(p)))(\operatorname{pubkey}(p)(a)) \in \operatorname{adv\text{-}tokens}(\operatorname{pk\text{-}inf}(p))) \wedge \\ \quad \neg (d = f) \wedge \\ \quad \operatorname{uniquely\text{-}originates}(\operatorname{ispace}(\operatorname{pk\text{-}inf}(p)), f) \Rightarrow \\ \quad \quad \exists g \in \operatorname{bnodes}(c). \\ \quad \quad \quad \operatorname{is\text{-}subterm}(\operatorname{tokens}(\operatorname{ispace}(\operatorname{pk\text{-}inf}(p))), \operatorname{keys}(\operatorname{ispace}(\operatorname{pk\text{-}inf}(p))), f, \\ \quad \quad \quad \quad \operatorname{node\text{-}content}(\operatorname{ispace}(\operatorname{pk\text{-}inf}(p)), g)) \wedge \\ \quad \quad \quad \neg \operatorname{is\text{-}subterm}(\operatorname{tokens}(\operatorname{ispace}(\operatorname{pk\text{-}inf}(p))), \operatorname{keys}(\operatorname{ispace}(\operatorname{pk\text{-}inf}(p))), \operatorname{encr}(\operatorname{lst}({[}d, f]), \operatorname{pubkey}(p)(a)), \operatorname{node\text{-}content}(\operatorname{ispace}(\operatorname{pk\text{-}inf}(p)), g)) \wedge \\ \quad \quad \quad \forall z \in \operatorname{bnodes}(c). \\ \quad \quad \quad \quad \operatorname{is\text{-}subterm}(\operatorname{tokens}(\operatorname{ispace}(\operatorname{pk\text{-}inf}(p))), \operatorname{keys}(\operatorname{ispace}(\operatorname{pk\text{-}inf}(p))), f, \\ \quad \quad \quad \quad \quad \operatorname{node\text{-}content}(\operatorname{ispace}(\operatorname{pk\text{-}inf}(p)), z)) \wedge \\ \quad \quad \quad \quad \neg \operatorname{is\text{-}subterm}(\operatorname{tokens}(\operatorname{ispace}(\operatorname{pk\text{-}inf}(p))), \operatorname{keys}(\operatorname{ispace}(\operatorname{pk\text{-}inf}(p))), \operatorname{encr}(\operatorname{lst}({[}d, f]), \operatorname{pubkey}(p)(a)), \operatorname{node\text{-}content}(\operatorname{ispace}(\operatorname{pk\text{-}inf}(p)), z)) \Rightarrow \\ \quad \quad \quad \quad \quad \neg ({[}z, g] \in \operatorname{precedes\text{-}strict}(c)) \wedge \\ \quad \quad \quad \operatorname{is\text{-}regular\text{-}node}(\operatorname{pk\text{-}inf}(p), g) \wedge \\ \quad \quad \quad \operatorname{is\text{-}sent}(\operatorname{ispace}(\operatorname{pk\text{-}inf}(p)), g) \end{array}}}
\vnbprop{nso-lemma-5-9}{\vnbstmt{\begin{array}{@{}l@{}} \forall p, c, s, a, b, d, f.; \forall g \in \operatorname{bnodes}(c). \\ \quad \operatorname{is\text{-}nso\text{-}space}(p) \wedge \\ \quad \operatorname{is\text{-}bundle}(c) \wedge \\ \quad \operatorname{bspace}(c) = \operatorname{ispace}(\operatorname{pk\text{-}inf}(p)) \wedge \\ \quad s \in \operatorname{ids}(\operatorname{ispace}(\operatorname{pk\text{-}inf}(p))) \wedge \\ \quad \operatorname{is\text{-}nso\text{-}resp\text{-}strand}(p, \operatorname{strand\text{-}map}(\operatorname{ispace}(\operatorname{pk\text{-}inf}(p)))(s), a, b, d, f) \wedge \\ \quad {[}3, s] \in \operatorname{bnodes}(c) \wedge \\ \quad \neg (\operatorname{kinv}(\operatorname{ispace}(\operatorname{pk\text{-}inf}(p)))(\operatorname{pubkey}(p)(a)) \in \operatorname{adv\text{-}tokens}(\operatorname{pk\text{-}inf}(p))) \wedge \\ \quad \neg (d = f) \wedge \\ \quad \operatorname{uniquely\text{-}originates}(\operatorname{ispace}(\operatorname{pk\text{-}inf}(p)), f) \wedge \\ \quad \operatorname{is\text{-}subterm}(\operatorname{tokens}(\operatorname{ispace}(\operatorname{pk\text{-}inf}(p))), \operatorname{keys}(\operatorname{ispace}(\operatorname{pk\text{-}inf}(p))), f, \\ \quad \quad \operatorname{node\text{-}content}(\operatorname{ispace}(\operatorname{pk\text{-}inf}(p)), g)) \wedge \\ \quad \neg \operatorname{is\text{-}subterm}(\operatorname{tokens}(\operatorname{ispace}(\operatorname{pk\text{-}inf}(p))), \operatorname{keys}(\operatorname{ispace}(\operatorname{pk\text{-}inf}(p))), \operatorname{encr}(\operatorname{lst}({[}d, f]), \operatorname{pubkey}(p)(a)), \operatorname{node\text{-}content}(\operatorname{ispace}(\operatorname{pk\text{-}inf}(p)), g)) \wedge \\ \quad \forall z \in \operatorname{bnodes}(c). \\ \quad \quad \operatorname{is\text{-}subterm}(\operatorname{tokens}(\operatorname{ispace}(\operatorname{pk\text{-}inf}(p))), \operatorname{keys}(\operatorname{ispace}(\operatorname{pk\text{-}inf}(p))), f, \\ \quad \quad \quad \operatorname{node\text{-}content}(\operatorname{ispace}(\operatorname{pk\text{-}inf}(p)), z)) \wedge \\ \quad \quad \neg \operatorname{is\text{-}subterm}(\operatorname{tokens}(\operatorname{ispace}(\operatorname{pk\text{-}inf}(p))), \operatorname{keys}(\operatorname{ispace}(\operatorname{pk\text{-}inf}(p))), \operatorname{encr}(\operatorname{lst}({[}d, f]), \operatorname{pubkey}(p)(a)), \operatorname{node\text{-}content}(\operatorname{ispace}(\operatorname{pk\text{-}inf}(p)), z)) \Rightarrow \\ \quad \quad \quad \neg ({[}z, g] \in \operatorname{precedes\text{-}strict}(c)) \wedge \\ \quad \operatorname{is\text{-}regular\text{-}node}(\operatorname{pk\text{-}inf}(p), g) \wedge \\ \quad \operatorname{is\text{-}sent}(\operatorname{ispace}(\operatorname{pk\text{-}inf}(p)), g) \Rightarrow \\ \quad \quad \exists h. \\ \quad \quad \quad \operatorname{same\text{-}strand\text{-}before}(\operatorname{ispace}(\operatorname{pk\text{-}inf}(p)), h, g) \wedge \\ \quad \quad \quad \operatorname{node\text{-}content}(\operatorname{ispace}(\operatorname{pk\text{-}inf}(p)), h) = \operatorname{encr}(\operatorname{lst}({[}d, f]), \operatorname{pubkey}(p)(a)) \end{array}}}

\vnbbookfile{structure-library/strand-honesty}{structure-library/strand-honesty}
\noindent{\small strand-honesty.scm -- HONEST IDEALS (Thayer, Herzog, Guttman, \textasciitilde{}/docs/jcs\_strand\_spaces.pdf, Definition 6.10), on the user's model: an infiltrated space isp, a bundle c, a set I of messages.  The paper: "I is honest relative to a bundle C iff whenever a penetrator node p is an entry point for I, p is an M node or a K ...}\par
\par\medskip\noindent\textbf{Definition} (\texttt{is-honest}).\;For all $\operatorname{shp\_}$, $\operatorname{shc\_}$ and $\operatorname{shi\_}$, 
\begin{equation*}
\begin{array}{@{}l@{}} \operatorname{is\text{-}honest}(\operatorname{shp\_}, \operatorname{shc\_}, \operatorname{shi\_}) \\ \Leftrightarrow\; \forall \operatorname{shn\_} \in \operatorname{bnodes}(\operatorname{shc\_}). \\ \quad \operatorname{is\text{-}adversary\text{-}node}(\operatorname{shp\_}, \operatorname{shn\_}) \Rightarrow  \\ \quad \operatorname{is\text{-}entry\text{-}point}(\operatorname{ispace}(\operatorname{shp\_}), \operatorname{shn\_}, \operatorname{shi\_}) \Rightarrow  \\ \quad \quad \operatorname{is\text{-}text\text{-}trace}(\operatorname{ispace}(\operatorname{shp\_}), \operatorname{adv\text{-}tokens}(\operatorname{shp\_}), \operatorname{strand\text{-}map}(\operatorname{ispace}(\operatorname{shp\_}))(\operatorname{nth}(2, \operatorname{shn\_}))) \end{array}
\end{equation*}


\vnbbookfile{theorem-library/strand-honesty-laws}{theorem-library/strand-honesty-laws}
\noindent{\small strand-honesty-laws.scm -- HONESTY (Thayer, Herzog, Guttman, \textasciitilde{}/docs/jcs\_strand\_spaces.pdf, section 6.2-6.3): Theorem 6.11 (ideals generated by atoms are honest), Corollaries 6.12 and 6.13, and Proposition 3.3 re-derived from 6.12 as the paper says ("a special case of this, in which \{K\} is chosen as S and K is chosen as ...}\par
\vnbprop{adversary-sent-shapes}{\vnbstmt{\begin{array}{@{}l@{}} \forall p, m. \\ \quad \operatorname{is\text{-}infiltrated\text{-}space}(p) \wedge \\ \quad m \in \operatorname{nodes}(\operatorname{ispace}(p)) \wedge \\ \quad \operatorname{is\text{-}adversary\text{-}node}(p, m) \wedge \\ \quad \operatorname{is\text{-}sent}(\operatorname{ispace}(p), m) \Rightarrow \\ \quad \quad \operatorname{is\text{-}text\text{-}trace}(\operatorname{ispace}(p), \operatorname{adv\text{-}tokens}(p), \operatorname{strand\text{-}map}(\operatorname{ispace}(p))(\operatorname{nth}(2, m))) \wedge \\ \quad \quad \operatorname{node\text{-}content}(\operatorname{ispace}(p), m) \in \operatorname{adv\text{-}tokens}(p) \vee \\ \quad \quad \exists x. \\ \quad \quad \quad (\operatorname{same\text{-}strand\text{-}before}(\operatorname{ispace}(p), x, m) \wedge (\operatorname{node\text{-}content}(\operatorname{ispace}(p), x) = \operatorname{node\text{-}content}(\operatorname{ispace}(p), m))) \vee \\ \quad \quad \exists l \in \operatorname{tuples}(\operatorname{msg\text{-}set}(\operatorname{tokens}(\operatorname{ispace}(p)), \operatorname{keys}(\operatorname{ispace}(p)))). \\ \quad \quad \quad \operatorname{node\text{-}content}(\operatorname{ispace}(p), m) = \operatorname{lst}(l) \wedge \\ \quad \quad \quad \forall i \in \mathbb{N}. \\ \quad \quad \quad \quad ((1 \leq i) \wedge (i \leq \operatorname{length}(l))) \Rightarrow  \\ \quad \quad \quad \quad \quad \exists x. \\ \quad \quad \quad \quad \quad \quad \operatorname{same\text{-}strand\text{-}before}(\operatorname{ispace}(p), x, m) \wedge (\operatorname{node\text{-}content}(\operatorname{ispace}(p), x) = \operatorname{nth}(i, l)) \vee \\ \quad \quad \exists l \in \operatorname{tuples}(\operatorname{msg\text{-}set}(\operatorname{tokens}(\operatorname{ispace}(p)), \operatorname{keys}(\operatorname{ispace}(p)))), q \in \mathbb{N}. \\ \quad \quad \quad 1 \leq q \wedge \\ \quad \quad \quad q \leq \operatorname{length}(l) \wedge \\ \quad \quad \quad \operatorname{node\text{-}content}(\operatorname{ispace}(p), m) = \operatorname{nth}(q, l) \wedge \\ \quad \quad \quad \exists x. \\ \quad \quad \quad \quad \operatorname{same\text{-}strand\text{-}before}(\operatorname{ispace}(p), x, m) \wedge (\operatorname{node\text{-}content}(\operatorname{ispace}(p), x) = \operatorname{lst}(l)) \vee \\ \quad \quad \exists h \in \operatorname{msg\text{-}set}(\operatorname{tokens}(\operatorname{ispace}(p)), \operatorname{keys}(\operatorname{ispace}(p))), k \in \operatorname{keys}(\operatorname{ispace}(p)). \\ \quad \quad \quad \operatorname{node\text{-}content}(\operatorname{ispace}(p), m) = \operatorname{encr}(h, k) \wedge \\ \quad \quad \quad \exists x. \\ \quad \quad \quad \quad \operatorname{same\text{-}strand\text{-}before}(\operatorname{ispace}(p), x, m) \wedge (\operatorname{node\text{-}content}(\operatorname{ispace}(p), x) = k) \wedge \\ \quad \quad \quad \exists x. \\ \quad \quad \quad \quad \operatorname{same\text{-}strand\text{-}before}(\operatorname{ispace}(p), x, m) \wedge (\operatorname{node\text{-}content}(\operatorname{ispace}(p), x) = h) \vee \\ \quad \quad \exists h \in \operatorname{msg\text{-}set}(\operatorname{tokens}(\operatorname{ispace}(p)), \operatorname{keys}(\operatorname{ispace}(p))), k \in \operatorname{keys}(\operatorname{ispace}(p)). \\ \quad \quad \quad \operatorname{node\text{-}content}(\operatorname{ispace}(p), m) = h \wedge \\ \quad \quad \quad \exists x. \\ \quad \quad \quad \quad (\operatorname{same\text{-}strand\text{-}before}(\operatorname{ispace}(p), x, m) \wedge (\operatorname{node\text{-}content}(\operatorname{ispace}(p), x) = \operatorname{kinv}(\operatorname{ispace}(p))(k))) \wedge \\ \quad \quad \quad \exists x. \\ \quad \quad \quad \quad \operatorname{same\text{-}strand\text{-}before}(\operatorname{ispace}(p), x, m) \wedge (\operatorname{node\text{-}content}(\operatorname{ispace}(p), x) = \operatorname{encr}(h, k)) \end{array}}}
\vnbprop{tokens-are-simple}{\vnbstmt{\begin{array}{@{}l@{}} \forall t, k, s.; \forall u \in s. \\ \quad (\operatorname{is\text{-}token\text{-}set}(t) \wedge ((k \in \mathbf{Set}) \wedge \forall z \in s.\; (z \in t))) \Rightarrow  \\ \quad \quad \operatorname{is\text{-}simple\text{-}msg}(t, k, u) \end{array}}}
\vnbprop{key-ideal-gen-decrypt-atoms}{\vnbstmt{\begin{array}{@{}l@{}} \forall t, k, a, s, h, y. \\ \quad \operatorname{is\text{-}token\text{-}set}(t) \wedge \\ \quad k \in \mathbf{Set} \wedge \\ \quad \forall w \in a.\; (w \in k) \wedge \\ \quad \forall z \in s.\; (z \in t) \wedge \\ \quad h \in \operatorname{msg\text{-}set}(t, k) \wedge \\ \quad y \in k \wedge \\ \quad \operatorname{encr}(h, y) \in \operatorname{key\text{-}ideal\text{-}gen}(t, k, a, s) \Rightarrow \\ \quad \quad h \in \operatorname{key\text{-}ideal\text{-}gen}(t, k, a, s) \end{array}}}
\vnbprop{key-ideal-gen-honest}{\vnbstmt{\begin{array}{@{}l@{}} \forall p, c, a, s. \\ \quad \operatorname{is\text{-}infiltrated\text{-}space}(p) \wedge \\ \quad \operatorname{is\text{-}bundle}(c) \wedge \\ \quad \operatorname{bspace}(c) = \operatorname{ispace}(p) \wedge \\ \quad \forall y \in a.\; (y \in \operatorname{keys}(\operatorname{ispace}(p))) \wedge \\ \quad \forall y \in \operatorname{keys}(\operatorname{ispace}(p)).\; ((y \in s) \vee (\operatorname{kinv}(\operatorname{ispace}(p))(y) \in a)) \wedge \\ \quad \forall z \in s.\; (z \in \operatorname{tokens}(\operatorname{ispace}(p))) \Rightarrow \\ \quad \quad \operatorname{is\text{-}honest}(p, c, \operatorname{key\text{-}ideal\text{-}gen}(\operatorname{tokens}(\operatorname{ispace}(p)), \operatorname{keys}(\operatorname{ispace}(p)), a, s)) \end{array}}}
\vnbprop{text-node-content}{\vnbstmt{\begin{array}{@{}l@{}} \forall p, m. \\ \quad \operatorname{is\text{-}infiltrated\text{-}space}(p) \wedge \\ \quad m \in \operatorname{nodes}(\operatorname{ispace}(p)) \wedge \\ \quad \operatorname{is\text{-}text\text{-}trace}(\operatorname{ispace}(p), \operatorname{adv\text{-}tokens}(p), \operatorname{strand\text{-}map}(\operatorname{ispace}(p))(\operatorname{nth}(2, m))) \Rightarrow \\ \quad \quad \operatorname{node\text{-}content}(\operatorname{ispace}(p), m) \in \operatorname{adv\text{-}tokens}(p) \end{array}}}
\vnbprop{key-ideal-gen-token}{\vnbstmt{\begin{array}{@{}l@{}} \forall t, k, a, s, x. \\ \quad \operatorname{is\text{-}token\text{-}set}(t) \wedge \\ \quad k \in \mathbf{Set} \wedge \\ \quad \forall w \in a.\; (w \in k) \wedge \\ \quad \forall z \in s.\; (z \in \operatorname{msg\text{-}set}(t, k)) \wedge \\ \quad x \in t \wedge \\ \quad x \in \operatorname{key\text{-}ideal\text{-}gen}(t, k, a, s) \Rightarrow \\ \quad \quad x \in s \end{array}}}
\vnbprop{key-ideal-gen-regular-entry}{\vnbstmt{\begin{array}{@{}l@{}} \forall p, c, a, s. \\ \quad \operatorname{is\text{-}infiltrated\text{-}space}(p) \wedge \\ \quad \operatorname{is\text{-}bundle}(c) \wedge \\ \quad \operatorname{bspace}(c) = \operatorname{ispace}(p) \wedge \\ \quad \forall y \in a.\; (y \in \operatorname{keys}(\operatorname{ispace}(p))) \wedge \\ \quad \forall y \in \operatorname{keys}(\operatorname{ispace}(p)).\; ((y \in s) \vee (\operatorname{kinv}(\operatorname{ispace}(p))(y) \in a)) \wedge \\ \quad \forall z \in s.\; (z \in \operatorname{keys}(\operatorname{ispace}(p))) \wedge \\ \quad \forall z \in s.\; \neg (z \in \operatorname{adv\text{-}tokens}(p)) \wedge \\ \quad \exists m \in \operatorname{bnodes}(c). \\ \quad \quad \operatorname{node\text{-}content}(\operatorname{ispace}(p), m) \in \operatorname{key\text{-}ideal\text{-}gen}(\operatorname{tokens}(\operatorname{ispace}(p)), \operatorname{keys}(\operatorname{ispace}(p)), a, s) \Rightarrow \\ \quad \quad \exists n \in \operatorname{bnodes}(c). \\ \quad \quad \quad \operatorname{is\text{-}regular\text{-}node}(p, n) \wedge \\ \quad \quad \quad \operatorname{is\text{-}entry\text{-}point}(\operatorname{ispace}(p), n, \operatorname{key\text{-}ideal\text{-}gen}(\operatorname{tokens}(\operatorname{ispace}(p)), \operatorname{keys}(\operatorname{ispace}(p)), a, s)) \end{array}}}
\vnbprop{key-ideal-gen-no-adversary-origin}{\vnbstmt{\begin{array}{@{}l@{}} \forall p, c, a, s.; \forall n \in \operatorname{bnodes}(c), g, y. \\ \quad \operatorname{is\text{-}infiltrated\text{-}space}(p) \wedge \\ \quad \operatorname{is\text{-}bundle}(c) \wedge \\ \quad \operatorname{bspace}(c) = \operatorname{ispace}(p) \wedge \\ \quad \forall y \in a.\; (y \in \operatorname{keys}(\operatorname{ispace}(p))) \wedge \\ \quad \forall y \in \operatorname{keys}(\operatorname{ispace}(p)).\; ((y \in s) \vee (\operatorname{kinv}(\operatorname{ispace}(p))(y) \in a)) \wedge \\ \quad \forall z \in s.\; (z \in \operatorname{keys}(\operatorname{ispace}(p))) \wedge \\ \quad \forall z \in s.\; \neg (z \in \operatorname{adv\text{-}tokens}(p)) \wedge \\ \quad \forall n \in \operatorname{bnodes}(c). \\ \quad \quad \operatorname{is\text{-}regular\text{-}node}(p, n) \Rightarrow  \\ \quad \quad \quad \neg \operatorname{is\text{-}entry\text{-}point}(\operatorname{ispace}(p), n, \operatorname{key\text{-}ideal\text{-}gen}(\operatorname{tokens}(\operatorname{ispace}(p)), \operatorname{keys}(\operatorname{ispace}(p)), a, s)) \wedge \\ \quad \operatorname{is\text{-}adversary\text{-}node}(p, n) \wedge \\ \quad g \in \operatorname{msg\text{-}set}(\operatorname{tokens}(\operatorname{ispace}(p)), \operatorname{keys}(\operatorname{ispace}(p))) \wedge \\ \quad y \in s \Rightarrow \\ \quad \quad \neg \operatorname{originates}(\operatorname{ispace}(p), \operatorname{encr}(g, y), n) \end{array}}}
\vnbprop{key-ideal-gen-subterms}{\vnbstmt{\begin{array}{@{}l@{}} \forall t, k, t_{0}. \\ \quad (\operatorname{is\text{-}token\text{-}set}(t) \wedge ((k \in \mathbf{Set}) \wedge (t_{0} \in \operatorname{msg\text{-}set}(t, k)))) \Rightarrow  \\ \quad \quad \operatorname{key\text{-}ideal\text{-}gen}(t, k, k, \operatorname{singleton}(t_{0})) = \{w \in \operatorname{msg\text{-}set}(t, k) : \operatorname{is\text{-}subterm}(t, k, t_{0}, w)\} \end{array}}}
\vnbprop{adversary-bound-via-ideals}{\vnbstmt{\begin{array}{@{}l@{}} \forall p, c, u.; \forall n \in \operatorname{bnodes}(c). \\ \quad \operatorname{is\text{-}infiltrated\text{-}space}(p) \wedge \\ \quad \operatorname{is\text{-}bundle}(c) \wedge \\ \quad \operatorname{bspace}(c) = \operatorname{ispace}(p) \wedge \\ \quad u \in \operatorname{keys}(\operatorname{ispace}(p)) \wedge \\ \quad \neg (u \in \operatorname{adv\text{-}tokens}(p)) \wedge \\ \quad \forall n \in \operatorname{bnodes}(c). \\ \quad \quad \operatorname{is\text{-}regular\text{-}node}(p, n) \Rightarrow \neg \operatorname{originates}(\operatorname{ispace}(p), u, n) \Rightarrow \\ \quad \quad \neg \operatorname{is\text{-}subterm}(\operatorname{tokens}(\operatorname{ispace}(p)), \operatorname{keys}(\operatorname{ispace}(p)), u, \operatorname{node\text{-}content}(\operatorname{ispace}(p), n)) \end{array}}}

\vnbbookfile{structure-library/otway-rees}{structure-library/otway-rees}
\noindent{\small otway-rees.scm -- THE OTWAY-REES PROTOCOL ON THE USER'S STRAND MODEL (Thayer, Herzog, Guttman, \textasciitilde{}/docs/jcs\_strand\_spaces.pdf, section 7.1, Definitions 7.1 and 7.3). DEFINITIONS ONLY; the laws are in theorem-library/otway-rees-laws.scm.  Agent SB-3, wave MAP-1026, 2026-10-26.  Destination: ...}\par
\vnbstructdef{sk-space}{A sk space $s$ is a 3-tuple whose components are listed below.\par
\begin{enumerate}[nosep]
\item $\operatorname{sk\text{-}inf}(s) = \operatorname{nth}(1, s)$, a substructure (\texttt{substructure}): $\operatorname{is\text{-}infiltrated\text{-}space}(\operatorname{sk\text{-}inf}(s))$.
\item $\operatorname{sk\text{-}names}(s) = \operatorname{nth}(2, s)$, a carrier (\texttt{carriers}): $\operatorname{sk\text{-}names}(s) \in \mathbf{Set}$.
\item $\operatorname{ltkey}(s) = \operatorname{nth}(3, s)$, an operation (\texttt{op}): $\operatorname{ltkey}(s) \in (\operatorname{sk\text{-}names}(s) \to \operatorname{keys}(\operatorname{ispace}(\operatorname{sk\text{-}inf}(s))))$.
\end{enumerate}
The laws (the \texttt{property} and \texttt{law} clauses):
\begin{itemize}[nosep]
\item $\forall \operatorname{orx\_} \in \operatorname{sk\text{-}names}(s).\; (\operatorname{orx\_} \in \operatorname{tokens}(\operatorname{ispace}(\operatorname{sk\text{-}inf}(s))))$
\item $\forall \operatorname{orx\_} \in \operatorname{sk\text{-}names}(s).\; \neg (\operatorname{orx\_} \in \operatorname{keys}(\operatorname{ispace}(\operatorname{sk\text{-}inf}(s))))$
\item
\begin{equation*}
\begin{array}{@{}l@{}} \forall \operatorname{orx\_}, \operatorname{ory\_} \in \operatorname{sk\text{-}names}(s). \\ \quad (\operatorname{ltkey}(s)(\operatorname{orx\_}) = \operatorname{ltkey}(s)(\operatorname{ory\_})) \Rightarrow (\operatorname{orx\_} = \operatorname{ory\_}) \end{array}
\end{equation*}

\item
\begin{equation*}
\begin{array}{@{}l@{}} \forall \operatorname{orx\_} \in \operatorname{sk\text{-}names}(s). \\ \quad \operatorname{kinv}(\operatorname{ispace}(\operatorname{sk\text{-}inf}(s)))(\operatorname{ltkey}(s)(\operatorname{orx\_})) = \operatorname{ltkey}(s)(\operatorname{orx\_}) \end{array}
\end{equation*}

\end{itemize}
The characteristic law (the axiom \texttt{is-sk-space}): $\operatorname{is\text{-}sk\text{-}space}(s)$ holds if and only if $\operatorname{length}(s) = 3$, every component has the typing listed with it, every law holds, and $s \in \mathbf{Set}$ (the structure is a set).
}
\par\medskip\noindent\textbf{Definition} (\texttt{is-or-text}).\;For all $\operatorname{ortp\_}$ and $\operatorname{ortx\_}$, 
\begin{equation*}
\begin{array}{@{}l@{}} \operatorname{is\text{-}or\text{-}text}(\operatorname{ortp\_}, \operatorname{ortx\_}) \\ \Leftrightarrow\; \operatorname{ortx\_} \in \operatorname{tokens}(\operatorname{ispace}(\operatorname{sk\text{-}inf}(\operatorname{ortp\_}))) \wedge \\ \neg (\operatorname{ortx\_} \in \operatorname{sk\text{-}names}(\operatorname{ortp\_})) \wedge \\ \neg (\operatorname{ortx\_} \in \operatorname{keys}(\operatorname{ispace}(\operatorname{sk\text{-}inf}(\operatorname{ortp\_})))) \end{array}
\end{equation*}

\par\medskip\noindent\textbf{Definition} (\texttt{is-or-init}).\;For all $\operatorname{orip\_}$, $\operatorname{oris\_}$, $\operatorname{oria\_}$, $\operatorname{orib\_}$, $\operatorname{orin\_}$, $\operatorname{orim\_}$ and $\operatorname{orik\_}$, 
\begin{equation*}
\begin{array}{@{}l@{}} \operatorname{is\text{-}or\text{-}init}(\operatorname{orip\_}, \operatorname{oris\_}, \operatorname{oria\_}, \operatorname{orib\_}, \operatorname{orin\_}, \operatorname{orim\_}, \operatorname{orik\_}) \\ \Leftrightarrow\; \operatorname{oria\_} \in \operatorname{sk\text{-}names}(\operatorname{orip\_}) \wedge \\ \operatorname{orib\_} \in \operatorname{sk\text{-}names}(\operatorname{orip\_}) \wedge \\ \operatorname{is\text{-}or\text{-}text}(\operatorname{orip\_}, \operatorname{orin\_}) \wedge \\ \operatorname{is\text{-}or\text{-}text}(\operatorname{orip\_}, \operatorname{orim\_}) \wedge \\ \operatorname{orik\_} \in \operatorname{keys}(\operatorname{ispace}(\operatorname{sk\text{-}inf}(\operatorname{orip\_}))) \wedge \\ \operatorname{oris\_} \\ =\; {[}{[}\mathit{sent}, \operatorname{lst}({[}\operatorname{orim\_}, \operatorname{oria\_}, \operatorname{orib\_}, \operatorname{encr}(\operatorname{lst}({[}\operatorname{orin\_}, \operatorname{orim\_}, \operatorname{oria\_}, \operatorname{orib\_}]), \operatorname{ltkey}(\operatorname{orip\_})(\operatorname{oria\_}))])], {[}\mathit{received}, \operatorname{lst}({[}\operatorname{orim\_}, \operatorname{encr}(\operatorname{lst}({[}\operatorname{orin\_}, \operatorname{orik\_}]), \operatorname{ltkey}(\operatorname{orip\_})(\operatorname{oria\_}))])]] \end{array}
\end{equation*}

\par\medskip\noindent\textbf{Definition} (\texttt{is-or-resp}).\;For all $\operatorname{orrp\_}$, $\operatorname{orrs\_}$, $\operatorname{orra\_}$, $\operatorname{orrb\_}$, $\operatorname{orrn\_}$, $\operatorname{orrm\_}$, $\operatorname{orrk\_}$, $\operatorname{orrh\_}$ and $\operatorname{orrg\_}$, 
\begin{equation*}
\begin{array}{@{}l@{}} \operatorname{is\text{-}or\text{-}resp}(\operatorname{orrp\_}, \operatorname{orrs\_}, \operatorname{orra\_}, \operatorname{orrb\_}, \operatorname{orrn\_}, \operatorname{orrm\_}, \operatorname{orrk\_}, \operatorname{orrh\_}, \operatorname{orrg\_}) \\ \Leftrightarrow\; \operatorname{orra\_} \in \operatorname{sk\text{-}names}(\operatorname{orrp\_}) \wedge \\ \operatorname{orrb\_} \in \operatorname{sk\text{-}names}(\operatorname{orrp\_}) \wedge \\ \operatorname{is\text{-}or\text{-}text}(\operatorname{orrp\_}, \operatorname{orrn\_}) \wedge \\ \operatorname{is\text{-}or\text{-}text}(\operatorname{orrp\_}, \operatorname{orrm\_}) \wedge \\ \operatorname{orrk\_} \in \operatorname{keys}(\operatorname{ispace}(\operatorname{sk\text{-}inf}(\operatorname{orrp\_}))) \wedge \\ \operatorname{orrh\_} \in \operatorname{msg\text{-}set}(\operatorname{tokens}(\operatorname{ispace}(\operatorname{sk\text{-}inf}(\operatorname{orrp\_}))), \operatorname{keys}(\operatorname{ispace}(\operatorname{sk\text{-}inf}(\operatorname{orrp\_})))) \wedge \\ \operatorname{orrg\_} \in \operatorname{msg\text{-}set}(\operatorname{tokens}(\operatorname{ispace}(\operatorname{sk\text{-}inf}(\operatorname{orrp\_}))), \operatorname{keys}(\operatorname{ispace}(\operatorname{sk\text{-}inf}(\operatorname{orrp\_})))) \wedge \\ \neg (\operatorname{orrn\_} = \operatorname{orrm\_}) \wedge \\ \neg \operatorname{is\text{-}subterm}(\operatorname{tokens}(\operatorname{ispace}(\operatorname{sk\text{-}inf}(\operatorname{orrp\_}))), \operatorname{keys}(\operatorname{ispace}(\operatorname{sk\text{-}inf}(\operatorname{orrp\_}))), \operatorname{orrn\_}, \operatorname{orrh\_}) \wedge \\ \operatorname{orrs\_} \\ =\; {[}{[}\mathit{received}, \operatorname{lst}({[}\operatorname{orrm\_}, \operatorname{orra\_}, \operatorname{orrb\_}, \operatorname{orrh\_}])], {[}\mathit{sent}, \operatorname{lst}({[}\operatorname{orrm\_}, \operatorname{orra\_}, \operatorname{orrb\_}, \operatorname{orrh\_}, \operatorname{encr}(\operatorname{lst}({[}\operatorname{orrn\_}, \operatorname{orrm\_}, \operatorname{orra\_}, \operatorname{orrb\_}]), \operatorname{ltkey}(\operatorname{orrp\_})(\operatorname{orrb\_}))])], {[}\mathit{received}, \operatorname{lst}({[}\operatorname{orrm\_}, \operatorname{orrg\_}, \operatorname{encr}(\operatorname{lst}({[}\operatorname{orrn\_}, \operatorname{orrk\_}]), \operatorname{ltkey}(\operatorname{orrp\_})(\operatorname{orrb\_}))])], {[}\mathit{sent}, \operatorname{lst}({[}\operatorname{orrm\_}, \operatorname{orrg\_}])]] \end{array}
\end{equation*}

\par\medskip\noindent\textbf{Definition} (\texttt{is-or-serv}).\;For all $\operatorname{orsp\_}$, $\operatorname{orss\_}$, $\operatorname{orsa\_}$, $\operatorname{orsb\_}$, $\operatorname{orsn\_}$, $\operatorname{orso\_}$, $\operatorname{orsm\_}$ and $\operatorname{orsk\_}$, 
\begin{equation*}
\begin{array}{@{}l@{}} \operatorname{is\text{-}or\text{-}serv}(\operatorname{orsp\_}, \operatorname{orss\_}, \operatorname{orsa\_}, \operatorname{orsb\_}, \operatorname{orsn\_}, \operatorname{orso\_}, \operatorname{orsm\_}, \operatorname{orsk\_}) \\ \Leftrightarrow\; \operatorname{orsa\_} \in \operatorname{sk\text{-}names}(\operatorname{orsp\_}) \wedge \\ \operatorname{orsb\_} \in \operatorname{sk\text{-}names}(\operatorname{orsp\_}) \wedge \\ \operatorname{is\text{-}or\text{-}text}(\operatorname{orsp\_}, \operatorname{orsn\_}) \wedge \\ \operatorname{is\text{-}or\text{-}text}(\operatorname{orsp\_}, \operatorname{orso\_}) \wedge \\ \operatorname{is\text{-}or\text{-}text}(\operatorname{orsp\_}, \operatorname{orsm\_}) \wedge \\ \operatorname{orsk\_} \in \operatorname{keys}(\operatorname{ispace}(\operatorname{sk\text{-}inf}(\operatorname{orsp\_}))) \wedge \\ \neg (\operatorname{orsk\_} \in \operatorname{adv\text{-}tokens}(\operatorname{sk\text{-}inf}(\operatorname{orsp\_}))) \wedge \\ \forall \operatorname{orsx\_} \in \operatorname{sk\text{-}names}(\operatorname{orsp\_}).\; \neg (\operatorname{orsk\_} = \operatorname{ltkey}(\operatorname{orsp\_})(\operatorname{orsx\_})) \wedge \\ \operatorname{kinv}(\operatorname{ispace}(\operatorname{sk\text{-}inf}(\operatorname{orsp\_})))(\operatorname{orsk\_}) = \operatorname{orsk\_} \wedge \\ \operatorname{orss\_} \\ =\; {[}{[}\mathit{received}, \operatorname{lst}({[}\operatorname{orsm\_}, \operatorname{orsa\_}, \operatorname{orsb\_}, \operatorname{encr}(\operatorname{lst}({[}\operatorname{orsn\_}, \operatorname{orsm\_}, \operatorname{orsa\_}, \operatorname{orsb\_}]), \operatorname{ltkey}(\operatorname{orsp\_})(\operatorname{orsa\_})), \operatorname{encr}(\operatorname{lst}({[}\operatorname{orso\_}, \operatorname{orsm\_}, \operatorname{orsa\_}, \operatorname{orsb\_}]), \operatorname{ltkey}(\operatorname{orsp\_})(\operatorname{orsb\_}))])], {[}\mathit{sent}, \operatorname{lst}({[}\operatorname{orsm\_}, \operatorname{encr}(\operatorname{lst}({[}\operatorname{orsn\_}, \operatorname{orsk\_}]), \operatorname{ltkey}(\operatorname{orsp\_})(\operatorname{orsa\_})), \operatorname{encr}(\operatorname{lst}({[}\operatorname{orso\_}, \operatorname{orsk\_}]), \operatorname{ltkey}(\operatorname{orsp\_})(\operatorname{orsb\_}))])]] \end{array}
\end{equation*}

\par\medskip\noindent\textbf{Definition} (\texttt{is-or-space}).\;For every $\operatorname{orqp\_}$, 
\begin{equation*}
\begin{array}{@{}l@{}} \operatorname{is\text{-}or\text{-}space}(\operatorname{orqp\_}) \\ \Leftrightarrow\; \operatorname{is\text{-}sk\text{-}space}(\operatorname{orqp\_}) \wedge \\ \forall \operatorname{orqi\_} \in \operatorname{ids}(\operatorname{ispace}(\operatorname{sk\text{-}inf}(\operatorname{orqp\_}))). \\ \quad \operatorname{orqi\_} \in \operatorname{adv\text{-}ids}(\operatorname{sk\text{-}inf}(\operatorname{orqp\_})) \vee \\ \quad \exists \operatorname{orqa\_}, \operatorname{orqb\_}, \operatorname{orqn\_}, \operatorname{orqm\_}, \operatorname{orqk\_}. \\ \quad \quad \operatorname{is\text{-}or\text{-}init}(\operatorname{orqp\_}, \operatorname{strand\text{-}map}(\operatorname{ispace}(\operatorname{sk\text{-}inf}(\operatorname{orqp\_})))(\operatorname{orqi\_}), \operatorname{orqa\_}, \operatorname{orqb\_}, \operatorname{orqn\_}, \operatorname{orqm\_}, \operatorname{orqk\_}) \vee \\ \quad \exists \operatorname{orqa\_}, \operatorname{orqb\_}, \operatorname{orqn\_}, \operatorname{orqm\_}, \operatorname{orqk\_}, \operatorname{orqh\_}, \operatorname{orqg\_}. \\ \quad \quad \operatorname{is\text{-}or\text{-}resp}(\operatorname{orqp\_}, \\ \quad \quad \quad \operatorname{strand\text{-}map}(\operatorname{ispace}(\operatorname{sk\text{-}inf}(\operatorname{orqp\_})))(\operatorname{orqi\_}), \\ \quad \quad \quad \operatorname{orqa\_}, \operatorname{orqb\_}, \operatorname{orqn\_}, \operatorname{orqm\_}, \operatorname{orqk\_}, \operatorname{orqh\_}, \operatorname{orqg\_}) \vee \\ \quad \exists \operatorname{orqa\_}, \operatorname{orqb\_}, \operatorname{orqn\_}, \operatorname{orqo\_}, \operatorname{orqm\_}, \operatorname{orqk\_}. \\ \quad \quad \operatorname{is\text{-}or\text{-}serv}(\operatorname{orqp\_}, \\ \quad \quad \quad \operatorname{strand\text{-}map}(\operatorname{ispace}(\operatorname{sk\text{-}inf}(\operatorname{orqp\_})))(\operatorname{orqi\_}), \\ \quad \quad \quad \operatorname{orqa\_}, \operatorname{orqb\_}, \operatorname{orqn\_}, \operatorname{orqo\_}, \operatorname{orqm\_}, \operatorname{orqk\_}) \end{array}
\end{equation*}

\par\medskip\noindent\textbf{Definition} (\texttt{hom-sk-space}).\;$\operatorname{hom\text{-}sk\text{-}space}(a, b) :=$
\begin{equation*}
\{\operatorname{homf\_} \in (\operatorname{sk\text{-}names}(a) \to \operatorname{sk\text{-}names}(b)) : \operatorname{is\text{-}hom\text{-}sk\text{-}space}(a, b, \operatorname{homf\_})\}
\end{equation*}

\vnbdef{sk-space-class}{\vnbstmt{\forall s.\; ((s \in \operatorname{sk\text{-}space}) \Leftrightarrow \operatorname{is\text{-}sk\text{-}space}(s))}}
\vnbdef{is-hom-sk-space-def}{\vnbstmt{\begin{array}{@{}l@{}} \forall a, b, f. \\ \quad \operatorname{is\text{-}hom\text{-}sk\text{-}space}(a, b, f) \\ \quad \Leftrightarrow\; \operatorname{is\text{-}sk\text{-}space}(a) \wedge \\ \quad \operatorname{is\text{-}sk\text{-}space}(b) \wedge \\ \quad f \in (\operatorname{sk\text{-}names}(a) \to \operatorname{sk\text{-}names}(b)) \wedge \\ \quad \operatorname{sk\text{-}inf}(a) = \operatorname{sk\text{-}inf}(b) \wedge \\ \quad \forall x_{1} \in \operatorname{sk\text{-}names}(a).\; (\operatorname{ltkey}(a)(x_{1}) = \operatorname{ltkey}(b)(f(x_{1}))) \end{array}}}
\vnbprop{is-sk-space-is-set}{\vnbstmt{\forall s.\; (\operatorname{is\text{-}sk\text{-}space}(s) \Rightarrow (s \in \mathbf{Set}))}}

\vnbbookfile{theorem-library/otway-rees-laws}{theorem-library/otway-rees-laws}
\noindent{\small otway-rees-laws.scm -- THE OTWAY-REES PROTOCOL: Lemma 7.2, Theorem 7.4, Propositions 7.5, 7.6, Corollary 7.7 (Thayer, Herzog, Guttman, \textasciitilde{}/docs/jcs\_strand\_spaces.pdf, section 7), on the user's model. Vocabulary: structure-library/otway-rees.scm.  Agent SB-3, wave MAP-1026, 2026-10-26.  Destination: ...}\par
\vnbprop{sk-space-parts}{\vnbstmt{\begin{array}{@{}l@{}} \forall p. \\ \quad \operatorname{is\text{-}sk\text{-}space}(p) \Rightarrow  \\ \quad \quad \operatorname{is\text{-}infiltrated\text{-}space}(\operatorname{sk\text{-}inf}(p)) \wedge \\ \quad \quad \operatorname{is\text{-}strand\text{-}space}(\operatorname{ispace}(\operatorname{sk\text{-}inf}(p))) \wedge \\ \quad \quad \operatorname{sk\text{-}names}(p) \in \mathbf{Set} \wedge \\ \quad \quad \operatorname{ltkey}(p) \in (\operatorname{sk\text{-}names}(p) \to \operatorname{keys}(\operatorname{ispace}(\operatorname{sk\text{-}inf}(p)))) \wedge \\ \quad \quad \forall x \in \operatorname{sk\text{-}names}(p).\; (x \in \operatorname{tokens}(\operatorname{ispace}(\operatorname{sk\text{-}inf}(p)))) \wedge \\ \quad \quad \forall x \in \operatorname{sk\text{-}names}(p).\; \neg (x \in \operatorname{keys}(\operatorname{ispace}(\operatorname{sk\text{-}inf}(p)))) \wedge \\ \quad \quad \forall x, y. \\ \quad \quad \quad ((x \in \operatorname{sk\text{-}names}(p)) \wedge ((y \in \operatorname{sk\text{-}names}(p)) \wedge (\operatorname{ltkey}(p)(x) = \operatorname{ltkey}(p)(y)))) \Rightarrow  \\ \quad \quad \quad \quad x = y \wedge \\ \quad \quad \forall x \in \operatorname{sk\text{-}names}(p). \\ \quad \quad \quad \operatorname{kinv}(\operatorname{ispace}(\operatorname{sk\text{-}inf}(p)))(\operatorname{ltkey}(p)(x)) = \operatorname{ltkey}(p)(x) \end{array}}}
\vnbprop{or-init-resp-disjoint}{\vnbstmt{\begin{array}{@{}l@{}} \forall p, s, a, b, n, m, k, c, d, f, q, r, h, g. \\ \quad (\operatorname{is\text{-}sk\text{-}space}(p) \wedge \operatorname{is\text{-}or\text{-}init}(p, s, a, b, n, m, k)) \Rightarrow  \\ \quad \quad \neg \operatorname{is\text{-}or\text{-}resp}(p, s, c, d, f, q, r, h, g) \end{array}}}
\vnbprop{or-init-serv-disjoint}{\vnbstmt{\begin{array}{@{}l@{}} \forall p, s, a, b, n, m, k, c, d, f, g, h, q. \\ \quad (\operatorname{is\text{-}sk\text{-}space}(p) \wedge \operatorname{is\text{-}or\text{-}init}(p, s, a, b, n, m, k)) \Rightarrow  \\ \quad \quad \neg \operatorname{is\text{-}or\text{-}serv}(p, s, c, d, f, g, h, q) \end{array}}}
\vnbprop{or-resp-serv-disjoint}{\vnbstmt{\begin{array}{@{}l@{}} \forall p, s, a, b, n, m, k, h, g, c, d, f, q, r, t. \\ \quad (\operatorname{is\text{-}sk\text{-}space}(p) \wedge \operatorname{is\text{-}or\text{-}resp}(p, s, a, b, n, m, k, h, g)) \Rightarrow  \\ \quad \quad \neg \operatorname{is\text{-}or\text{-}serv}(p, s, c, d, f, q, r, t) \end{array}}}
\vnbprop{key-ideal-gen-encr-atoms}{\vnbstmt{\begin{array}{@{}l@{}} \forall t, k, a, s, h, y. \\ \quad \operatorname{is\text{-}token\text{-}set}(t) \wedge \\ \quad k \in \mathbf{Set} \wedge \\ \quad \forall w \in a.\; (w \in k) \wedge \\ \quad \forall z \in s.\; (z \in t) \wedge \\ \quad h \in \operatorname{msg\text{-}set}(t, k) \wedge \\ \quad y \in k \wedge \\ \quad \operatorname{encr}(h, y) \in \operatorname{key\text{-}ideal\text{-}gen}(t, k, a, s) \Rightarrow \\ \quad \quad (h \in \operatorname{key\text{-}ideal\text{-}gen}(t, k, a, s)) \wedge (y \in a) \end{array}}}
\vnbprop{or-serv-key-originates}{\vnbstmt{\begin{array}{@{}l@{}} \forall p, s, a, b, c, d, m, k. \\ \quad \operatorname{is\text{-}sk\text{-}space}(p) \wedge \\ \quad s \in \operatorname{ids}(\operatorname{ispace}(\operatorname{sk\text{-}inf}(p))) \wedge \\ \quad \operatorname{is\text{-}or\text{-}serv}(p, \operatorname{strand\text{-}map}(\operatorname{ispace}(\operatorname{sk\text{-}inf}(p)))(s), a, b, c, d, m, k) \Rightarrow \\ \quad \quad \operatorname{originates}(\operatorname{ispace}(\operatorname{sk\text{-}inf}(p)), k, {[}2, s]) \end{array}}}
\vnbprop{or-session-key-secrecy}{\vnbstmt{\begin{array}{@{}l@{}} \forall p, c, v, a, b, d, f, m, k.; \forall z \in \operatorname{bnodes}(c). \\ \quad \operatorname{is\text{-}or\text{-}space}(p) \wedge \\ \quad \operatorname{is\text{-}bundle}(c) \wedge \\ \quad \operatorname{bspace}(c) = \operatorname{ispace}(\operatorname{sk\text{-}inf}(p)) \wedge \\ \quad v \in \operatorname{ids}(\operatorname{ispace}(\operatorname{sk\text{-}inf}(p))) \wedge \\ \quad \operatorname{is\text{-}or\text{-}serv}(p, \operatorname{strand\text{-}map}(\operatorname{ispace}(\operatorname{sk\text{-}inf}(p)))(v), a, b, d, f, m, k) \wedge \\ \quad \operatorname{uniquely\text{-}originates}(\operatorname{ispace}(\operatorname{sk\text{-}inf}(p)), k) \wedge \\ \quad \neg (\operatorname{ltkey}(p)(a) \in \operatorname{adv\text{-}tokens}(\operatorname{sk\text{-}inf}(p))) \wedge \\ \quad \neg (\operatorname{ltkey}(p)(b) \in \operatorname{adv\text{-}tokens}(\operatorname{sk\text{-}inf}(p))) \Rightarrow \\ \quad \quad \neg (\operatorname{node\text{-}content}(\operatorname{ispace}(\operatorname{sk\text{-}inf}(p)), z) \in \operatorname{key\text{-}ideal\text{-}gen}(\operatorname{tokens}(\operatorname{ispace}(\operatorname{sk\text{-}inf}(p))), \operatorname{keys}(\operatorname{ispace}(\operatorname{sk\text{-}inf}(p))), \{\mathit{orv} \in \operatorname{keys}(\operatorname{ispace}(\operatorname{sk\text{-}inf}(p))) : \neg (\mathit{orv} \in \{w \in \operatorname{keys}(\operatorname{ispace}(\operatorname{sk\text{-}inf}(p))) : ((w = \operatorname{ltkey}(p)(a)) \vee ((w = \operatorname{ltkey}(p)(b)) \vee (w = k)))\})\}, \operatorname{singleton}(k))) \end{array}}}
\vnbprop{or-ltkey-not-regular-origin}{\vnbstmt{\begin{array}{@{}l@{}} \forall p, x, n. \\ \quad (\operatorname{is\text{-}or\text{-}space}(p) \wedge ((x \in \operatorname{sk\text{-}names}(p)) \wedge \operatorname{is\text{-}regular\text{-}node}(\operatorname{sk\text{-}inf}(p), n))) \Rightarrow  \\ \quad \quad \neg \operatorname{originates}(\operatorname{ispace}(\operatorname{sk\text{-}inf}(p)), \operatorname{ltkey}(p)(x), n) \end{array}}}
\vnbprop{or-ltkey-no-adversary-origin}{\vnbstmt{\begin{array}{@{}l@{}} \forall p, c, x.; \forall n \in \operatorname{bnodes}(c), g \in \operatorname{msg\text{-}set}(\operatorname{tokens}(\operatorname{ispace}(\operatorname{sk\text{-}inf}(p))), \operatorname{keys}(\operatorname{ispace}(\operatorname{sk\text{-}inf}(p)))). \\ \quad \operatorname{is\text{-}or\text{-}space}(p) \wedge \\ \quad \operatorname{is\text{-}bundle}(c) \wedge \\ \quad \operatorname{bspace}(c) = \operatorname{ispace}(\operatorname{sk\text{-}inf}(p)) \wedge \\ \quad x \in \operatorname{sk\text{-}names}(p) \wedge \\ \quad \neg (\operatorname{ltkey}(p)(x) \in \operatorname{adv\text{-}tokens}(\operatorname{sk\text{-}inf}(p))) \wedge \\ \quad \operatorname{is\text{-}adversary\text{-}node}(\operatorname{sk\text{-}inf}(p), n) \Rightarrow \\ \quad \quad \neg \operatorname{originates}(\operatorname{ispace}(\operatorname{sk\text{-}inf}(p)), \operatorname{encr}(g, \operatorname{ltkey}(p)(x)), n) \end{array}}}
\vnbprop{or-ltkey-injective-at}{\vnbstmt{\begin{array}{@{}l@{}} \forall p, x, z. \\ \quad (\operatorname{is\text{-}sk\text{-}space}(p) \wedge ((x \in \operatorname{sk\text{-}names}(p)) \wedge ((z \in \operatorname{sk\text{-}names}(p)) \wedge (\operatorname{ltkey}(p)(x) = \operatorname{ltkey}(p)(z))))) \Rightarrow  \\ \quad \quad x = z \end{array}}}
\vnbprop{or-regular-encr-shapes}{\vnbstmt{\begin{array}{@{}l@{}} \forall p, x, n, h. \\ \quad \operatorname{is\text{-}or\text{-}space}(p) \wedge \\ \quad x \in \operatorname{sk\text{-}names}(p) \wedge \\ \quad \operatorname{is\text{-}regular\text{-}node}(\operatorname{sk\text{-}inf}(p), n) \wedge \\ \quad h \in \operatorname{msg\text{-}set}(\operatorname{tokens}(\operatorname{ispace}(\operatorname{sk\text{-}inf}(p))), \operatorname{keys}(\operatorname{ispace}(\operatorname{sk\text{-}inf}(p)))) \wedge \\ \quad \operatorname{originates}(\operatorname{ispace}(\operatorname{sk\text{-}inf}(p)), \operatorname{encr}(h, \operatorname{ltkey}(p)(x)), n) \Rightarrow \\ \quad \quad \exists a, b. \\ \quad \quad \quad \operatorname{is\text{-}or\text{-}text}(p, a) \wedge ((b \in \operatorname{keys}(\operatorname{ispace}(\operatorname{sk\text{-}inf}(p)))) \wedge (h = \operatorname{lst}({[}a, b]))) \vee \\ \quad \quad \exists a, c, d. \\ \quad \quad \quad (\operatorname{is\text{-}or\text{-}text}(p, a) \wedge (\operatorname{is\text{-}or\text{-}text}(p, c) \wedge ((d \in \operatorname{sk\text{-}names}(p)) \wedge (h = \operatorname{lst}({[}a, c, x, d]))))) \vee \\ \quad \quad \exists a, c, d. \\ \quad \quad \quad (\operatorname{is\text{-}or\text{-}text}(p, a) \wedge (\operatorname{is\text{-}or\text{-}text}(p, c) \wedge ((d \in \operatorname{sk\text{-}names}(p)) \wedge (h = \operatorname{lst}({[}a, c, d, x]))))) \end{array}}}
\vnbprop{or-init-nonce-originates}{\vnbstmt{\begin{array}{@{}l@{}} \forall p, t, a, b, n, m, k. \\ \quad \operatorname{is\text{-}sk\text{-}space}(p) \wedge \\ \quad t \in \operatorname{ids}(\operatorname{ispace}(\operatorname{sk\text{-}inf}(p))) \wedge \\ \quad \operatorname{is\text{-}or\text{-}init}(p, \operatorname{strand\text{-}map}(\operatorname{ispace}(\operatorname{sk\text{-}inf}(p)))(t), a, b, n, m, k) \Rightarrow \\ \quad \quad \operatorname{originates}(\operatorname{ispace}(\operatorname{sk\text{-}inf}(p)), n, {[}1, t]) \end{array}}}
\vnbprop{or-resp-nonce-originates}{\vnbstmt{\begin{array}{@{}l@{}} \forall p, t, a, b, n, m, k, h, g. \\ \quad \operatorname{is\text{-}sk\text{-}space}(p) \wedge \\ \quad t \in \operatorname{ids}(\operatorname{ispace}(\operatorname{sk\text{-}inf}(p))) \wedge \\ \quad \operatorname{is\text{-}or\text{-}resp}(p, \operatorname{strand\text{-}map}(\operatorname{ispace}(\operatorname{sk\text{-}inf}(p)))(t), a, b, n, m, k, h, g) \Rightarrow \\ \quad \quad \operatorname{originates}(\operatorname{ispace}(\operatorname{sk\text{-}inf}(p)), n, {[}2, t]) \end{array}}}
\vnbprop{or-77-server}{\vnbstmt{\begin{array}{@{}l@{}} \forall p, x, n, c, k. \\ \quad \operatorname{is\text{-}or\text{-}space}(p) \wedge \\ \quad x \in \operatorname{sk\text{-}names}(p) \wedge \\ \quad \operatorname{is\text{-}regular\text{-}node}(\operatorname{sk\text{-}inf}(p), n) \wedge \\ \quad c \in \operatorname{tokens}(\operatorname{ispace}(\operatorname{sk\text{-}inf}(p))) \wedge \\ \quad k \in \operatorname{keys}(\operatorname{ispace}(\operatorname{sk\text{-}inf}(p))) \wedge \\ \quad \operatorname{originates}(\operatorname{ispace}(\operatorname{sk\text{-}inf}(p)), \operatorname{encr}(\operatorname{lst}({[}c, k]), \operatorname{ltkey}(p)(x)), n) \Rightarrow \\ \quad \quad \operatorname{nth}(1, n) = 2 \wedge \\ \quad \quad \exists a, b, d. \\ \quad \quad \quad \operatorname{is\text{-}or\text{-}serv}(p, \operatorname{strand\text{-}map}(\operatorname{ispace}(\operatorname{sk\text{-}inf}(p)))(\operatorname{nth}(2, n)), a, x, b, c, d, k) \vee \\ \quad \quad \exists f, b, d. \\ \quad \quad \quad \operatorname{is\text{-}or\text{-}serv}(p, \operatorname{strand\text{-}map}(\operatorname{ispace}(\operatorname{sk\text{-}inf}(p)))(\operatorname{nth}(2, n)), x, f, c, b, d, k) \wedge \\ \quad \quad \operatorname{originates}(\operatorname{ispace}(\operatorname{sk\text{-}inf}(p)), k, {[}2, \operatorname{nth}(2, n)]) \end{array}}}
\vnbprop{or-77-initiator}{\vnbstmt{\begin{array}{@{}l@{}} \forall p, n, c, d, a, b. \\ \quad \operatorname{is\text{-}or\text{-}space}(p) \wedge \\ \quad a \in \operatorname{sk\text{-}names}(p) \wedge \\ \quad b \in \operatorname{sk\text{-}names}(p) \wedge \\ \quad \neg (a = b) \wedge \\ \quad \operatorname{is\text{-}regular\text{-}node}(\operatorname{sk\text{-}inf}(p), n) \wedge \\ \quad c \in \operatorname{tokens}(\operatorname{ispace}(\operatorname{sk\text{-}inf}(p))) \wedge \\ \quad d \in \operatorname{tokens}(\operatorname{ispace}(\operatorname{sk\text{-}inf}(p))) \wedge \\ \quad \operatorname{originates}(\operatorname{ispace}(\operatorname{sk\text{-}inf}(p)), \operatorname{encr}(\operatorname{lst}({[}c, d, a, b]), \operatorname{ltkey}(p)(a)), n) \Rightarrow \\ \quad \quad \operatorname{nth}(1, n) = 1 \wedge \\ \quad \quad \exists f. \\ \quad \quad \quad \operatorname{is\text{-}or\text{-}init}(p, \operatorname{strand\text{-}map}(\operatorname{ispace}(\operatorname{sk\text{-}inf}(p)))(\operatorname{nth}(2, n)), a, b, c, d, f) \wedge \\ \quad \quad \operatorname{originates}(\operatorname{ispace}(\operatorname{sk\text{-}inf}(p)), c, {[}1, \operatorname{nth}(2, n)]) \end{array}}}
\vnbprop{or-77-responder}{\vnbstmt{\begin{array}{@{}l@{}} \forall p, n, c, d, a, b. \\ \quad \operatorname{is\text{-}or\text{-}space}(p) \wedge \\ \quad a \in \operatorname{sk\text{-}names}(p) \wedge \\ \quad b \in \operatorname{sk\text{-}names}(p) \wedge \\ \quad \neg (a = b) \wedge \\ \quad \operatorname{is\text{-}regular\text{-}node}(\operatorname{sk\text{-}inf}(p), n) \wedge \\ \quad c \in \operatorname{tokens}(\operatorname{ispace}(\operatorname{sk\text{-}inf}(p))) \wedge \\ \quad d \in \operatorname{tokens}(\operatorname{ispace}(\operatorname{sk\text{-}inf}(p))) \wedge \\ \quad \operatorname{originates}(\operatorname{ispace}(\operatorname{sk\text{-}inf}(p)), \operatorname{encr}(\operatorname{lst}({[}c, d, a, b]), \operatorname{ltkey}(p)(b)), n) \Rightarrow \\ \quad \quad \operatorname{nth}(1, n) = 2 \wedge \\ \quad \quad \exists f, g, h. \\ \quad \quad \quad \operatorname{is\text{-}or\text{-}resp}(p, \operatorname{strand\text{-}map}(\operatorname{ispace}(\operatorname{sk\text{-}inf}(p)))(\operatorname{nth}(2, n)), a, b, c, d, f, g, h) \wedge \\ \quad \quad \operatorname{originates}(\operatorname{ispace}(\operatorname{sk\text{-}inf}(p)), c, {[}2, \operatorname{nth}(2, n)]) \end{array}}}
\vnbprop{bundle-subterm-originates}{\vnbstmt{\begin{array}{@{}l@{}} \forall c, a, u, z. \\ \quad \operatorname{is\text{-}bundle}(c) \wedge \\ \quad \operatorname{bspace}(c) = a \wedge \\ \quad z \in \operatorname{bnodes}(c) \wedge \\ \quad \operatorname{is\text{-}subterm}(\operatorname{tokens}(a), \operatorname{keys}(a), u, \operatorname{node\text{-}content}(a, z)) \Rightarrow \\ \quad \quad \exists r \in \operatorname{bnodes}(c).\; \operatorname{originates}(a, u, r) \end{array}}}
\vnbprop{or-initiator-nonce-source}{\vnbstmt{\begin{array}{@{}l@{}} \forall p, s, a, b, c, m, k, h, g, d, f, u, v. \\ \quad \operatorname{is\text{-}or\text{-}space}(p) \wedge \\ \quad s \in \operatorname{ids}(\operatorname{ispace}(\operatorname{sk\text{-}inf}(p))) \wedge \\ \quad \operatorname{is\text{-}or\text{-}init}(p, \operatorname{strand\text{-}map}(\operatorname{ispace}(\operatorname{sk\text{-}inf}(p)))(s), a, b, c, m, k) \wedge \\ \quad \operatorname{uniquely\text{-}originates}(\operatorname{ispace}(\operatorname{sk\text{-}inf}(p)), c) \wedge \\ \quad \operatorname{is\text{-}regular\text{-}node}(\operatorname{sk\text{-}inf}(p), d) \wedge \\ \quad f \in \operatorname{tokens}(\operatorname{ispace}(\operatorname{sk\text{-}inf}(p))) \wedge \\ \quad u \in \operatorname{sk\text{-}names}(p) \wedge \\ \quad v \in \operatorname{sk\text{-}names}(p) \wedge \\ \quad \operatorname{originates}(\operatorname{ispace}(\operatorname{sk\text{-}inf}(p)), \operatorname{encr}(\operatorname{lst}({[}c, f, u, v]), \operatorname{ltkey}(p)(a)), d) \Rightarrow \\ \quad \quad (f = m) \wedge ((u = a) \wedge (v = b)) \end{array}}}
\vnbprop{or-responder-nonce-source}{\vnbstmt{\begin{array}{@{}l@{}} \forall p, s, a, b, c, m, k, h, g, d, f, u, v. \\ \quad \operatorname{is\text{-}or\text{-}space}(p) \wedge \\ \quad s \in \operatorname{ids}(\operatorname{ispace}(\operatorname{sk\text{-}inf}(p))) \wedge \\ \quad \operatorname{is\text{-}or\text{-}resp}(p, \operatorname{strand\text{-}map}(\operatorname{ispace}(\operatorname{sk\text{-}inf}(p)))(s), a, b, c, m, k, h, g) \wedge \\ \quad \operatorname{uniquely\text{-}originates}(\operatorname{ispace}(\operatorname{sk\text{-}inf}(p)), c) \wedge \\ \quad \operatorname{is\text{-}regular\text{-}node}(\operatorname{sk\text{-}inf}(p), d) \wedge \\ \quad f \in \operatorname{tokens}(\operatorname{ispace}(\operatorname{sk\text{-}inf}(p))) \wedge \\ \quad u \in \operatorname{sk\text{-}names}(p) \wedge \\ \quad v \in \operatorname{sk\text{-}names}(p) \wedge \\ \quad \operatorname{originates}(\operatorname{ispace}(\operatorname{sk\text{-}inf}(p)), \operatorname{encr}(\operatorname{lst}({[}c, f, u, v]), \operatorname{ltkey}(p)(b)), d) \Rightarrow \\ \quad \quad (f = m) \wedge ((u = a) \wedge (v = b)) \end{array}}}
\vnbprop{or-initiator-guarantee}{\vnbstmt{\begin{array}{@{}l@{}} \forall p, c, s, a, b, d, m, k. \\ \quad \operatorname{is\text{-}or\text{-}space}(p) \wedge \\ \quad \operatorname{is\text{-}bundle}(c) \wedge \\ \quad \operatorname{bspace}(c) = \operatorname{ispace}(\operatorname{sk\text{-}inf}(p)) \wedge \\ \quad s \in \operatorname{ids}(\operatorname{ispace}(\operatorname{sk\text{-}inf}(p))) \wedge \\ \quad \neg (a = b) \wedge \\ \quad \operatorname{uniquely\text{-}originates}(\operatorname{ispace}(\operatorname{sk\text{-}inf}(p)), d) \wedge \\ \quad \neg (\operatorname{ltkey}(p)(a) \in \operatorname{adv\text{-}tokens}(\operatorname{sk\text{-}inf}(p))) \wedge \\ \quad \neg (\operatorname{ltkey}(p)(b) \in \operatorname{adv\text{-}tokens}(\operatorname{sk\text{-}inf}(p))) \wedge \\ \quad \operatorname{is\text{-}or\text{-}init}(p, \operatorname{strand\text{-}map}(\operatorname{ispace}(\operatorname{sk\text{-}inf}(p)))(s), a, b, d, m, k) \wedge \\ \quad {[}2, s] \in \operatorname{bnodes}(c) \Rightarrow \\ \quad \quad \exists f \in \operatorname{tokens}(\operatorname{ispace}(\operatorname{sk\text{-}inf}(p))). \\ \quad \quad \quad \exists r \in \operatorname{ids}(\operatorname{ispace}(\operatorname{sk\text{-}inf}(p))). \\ \quad \quad \quad \quad \neg (r \in \operatorname{adv\text{-}ids}(\operatorname{sk\text{-}inf}(p))) \wedge \\ \quad \quad \quad \quad {[}2, r] \in \operatorname{bnodes}(c) \wedge \\ \quad \quad \quad \quad \exists g, h, n. \\ \quad \quad \quad \quad \quad \operatorname{is\text{-}or\text{-}resp}(p, \operatorname{strand\text{-}map}(\operatorname{ispace}(\operatorname{sk\text{-}inf}(p)))(r), a, b, f, m, g, h, n) \wedge \\ \quad \quad \quad \exists v \in \operatorname{ids}(\operatorname{ispace}(\operatorname{sk\text{-}inf}(p))). \\ \quad \quad \quad \quad \neg (v \in \operatorname{adv\text{-}ids}(\operatorname{sk\text{-}inf}(p))) \wedge \\ \quad \quad \quad \quad {[}2, v] \in \operatorname{bnodes}(c) \wedge \\ \quad \quad \quad \quad \operatorname{is\text{-}or\text{-}serv}(p, \operatorname{strand\text{-}map}(\operatorname{ispace}(\operatorname{sk\text{-}inf}(p)))(v), a, b, d, f, m, k) \end{array}}}
\vnbprop{or-responder-guarantee}{\vnbstmt{\begin{array}{@{}l@{}} \forall p, c, s, a, b, d, m, k, h, g. \\ \quad \operatorname{is\text{-}or\text{-}space}(p) \wedge \\ \quad \operatorname{is\text{-}bundle}(c) \wedge \\ \quad \operatorname{bspace}(c) = \operatorname{ispace}(\operatorname{sk\text{-}inf}(p)) \wedge \\ \quad s \in \operatorname{ids}(\operatorname{ispace}(\operatorname{sk\text{-}inf}(p))) \wedge \\ \quad \neg (a = b) \wedge \\ \quad \operatorname{uniquely\text{-}originates}(\operatorname{ispace}(\operatorname{sk\text{-}inf}(p)), d) \wedge \\ \quad \neg (\operatorname{ltkey}(p)(a) \in \operatorname{adv\text{-}tokens}(\operatorname{sk\text{-}inf}(p))) \wedge \\ \quad \neg (\operatorname{ltkey}(p)(b) \in \operatorname{adv\text{-}tokens}(\operatorname{sk\text{-}inf}(p))) \wedge \\ \quad \operatorname{is\text{-}or\text{-}resp}(p, \operatorname{strand\text{-}map}(\operatorname{ispace}(\operatorname{sk\text{-}inf}(p)))(s), a, b, d, m, k, h, g) \wedge \\ \quad {[}3, s] \in \operatorname{bnodes}(c) \Rightarrow \\ \quad \quad \exists i \in \operatorname{ids}(\operatorname{ispace}(\operatorname{sk\text{-}inf}(p))). \\ \quad \quad \quad \neg (i \in \operatorname{adv\text{-}ids}(\operatorname{sk\text{-}inf}(p))) \wedge \\ \quad \quad \quad {[}1, i] \in \operatorname{bnodes}(c) \wedge \\ \quad \quad \quad \exists f, n. \\ \quad \quad \quad \quad \operatorname{is\text{-}or\text{-}init}(p, \operatorname{strand\text{-}map}(\operatorname{ispace}(\operatorname{sk\text{-}inf}(p)))(i), a, b, f, m, n) \wedge \\ \quad \quad \exists v \in \operatorname{ids}(\operatorname{ispace}(\operatorname{sk\text{-}inf}(p))). \\ \quad \quad \quad \neg (v \in \operatorname{adv\text{-}ids}(\operatorname{sk\text{-}inf}(p))) \wedge \\ \quad \quad \quad {[}2, v] \in \operatorname{bnodes}(c) \wedge \\ \quad \quad \quad \exists f. \\ \quad \quad \quad \quad \operatorname{is\text{-}or\text{-}serv}(p, \operatorname{strand\text{-}map}(\operatorname{ispace}(\operatorname{sk\text{-}inf}(p)))(v), a, b, f, d, m, k) \end{array}}}

\vnbbookfile{structure-library/or-missing-guarantee-instance}{structure-library/or-missing-guarantee-instance}
\noindent{\small cx-vocab.scm -- THE BUNDLE OF SECTION 7.3.3 (Figure 9) OF \textasciitilde{}/docs/jcs\_strand\_spaces.pdf: a concrete Otway-Rees strand space and bundle in which the initiator and the responder complete their strands and agree on A, B and M but not on the session key.  DEFINITIONS ONLY (four declare-instance! constants); the laws are in ...}\par
\par\medskip\noindent\textbf{Definition} (\texttt{cx-sp}).\;
\begin{equation*}
\begin{array}{@{}l@{}} \operatorname{cx\text{-}sp} \\ =\; {[}\mathbb{N}, \operatorname{interval}(1, 8), \operatorname{interval}(30, 33), (\lambda \operatorname{cxi\_} \in \operatorname{interval}(1, 8).\; \operatorname{nth}(\operatorname{cxi\_}, {[}{[}{[}\mathit{sent}, \operatorname{lst}({[}22, 10, 11, \operatorname{encr}(\operatorname{lst}({[}20, 22, 10, 11]), 30)])], {[}\mathit{received}, \operatorname{lst}({[}22, \operatorname{encr}(\operatorname{lst}({[}20, 33]), 30)])]], {[}{[}\mathit{received}, \operatorname{lst}({[}22, 10, 11, \operatorname{encr}(\operatorname{lst}({[}20, 22, 10, 11]), 30)])], {[}\mathit{sent}, \operatorname{lst}({[}22, 10, 11, \operatorname{encr}(\operatorname{lst}({[}20, 22, 10, 11]), 30), \operatorname{encr}(\operatorname{lst}({[}21, 22, 10, 11]), 31)])], {[}\mathit{received}, \operatorname{lst}({[}22, \operatorname{encr}(\operatorname{lst}({[}20, 32]), 30), \operatorname{encr}(\operatorname{lst}({[}21, 32]), 31)])], {[}\mathit{sent}, \operatorname{lst}({[}22, \operatorname{encr}(\operatorname{lst}({[}20, 32]), 30)])]], {[}{[}\mathit{received}, \operatorname{lst}({[}22, 10, 11, \operatorname{encr}(\operatorname{lst}({[}20, 22, 10, 11]), 30), \operatorname{encr}(\operatorname{lst}({[}21, 22, 10, 11]), 31)])], {[}\mathit{sent}, \operatorname{lst}({[}22, \operatorname{encr}(\operatorname{lst}({[}20, 32]), 30), \operatorname{encr}(\operatorname{lst}({[}21, 32]), 31)])]], {[}{[}\mathit{received}, \operatorname{lst}({[}22, 10, 11, \operatorname{encr}(\operatorname{lst}({[}20, 22, 10, 11]), 30), \operatorname{encr}(\operatorname{lst}({[}21, 22, 10, 11]), 31)])], {[}\mathit{sent}, \operatorname{lst}({[}22, \operatorname{encr}(\operatorname{lst}({[}20, 33]), 30), \operatorname{encr}(\operatorname{lst}({[}21, 33]), 31)])]], {[}{[}\mathit{received}, \operatorname{lst}({[}22, 10, 11, \operatorname{encr}(\operatorname{lst}({[}20, 22, 10, 11]), 30), \operatorname{encr}(\operatorname{lst}({[}21, 22, 10, 11]), 31)])], {[}\mathit{sent}, \operatorname{lst}({[}22, 10, 11, \operatorname{encr}(\operatorname{lst}({[}20, 22, 10, 11]), 30), \operatorname{encr}(\operatorname{lst}({[}21, 22, 10, 11]), 31)])], {[}\mathit{sent}, \operatorname{lst}({[}22, 10, 11, \operatorname{encr}(\operatorname{lst}({[}20, 22, 10, 11]), 30), \operatorname{encr}(\operatorname{lst}({[}21, 22, 10, 11]), 31)])]], {[}{[}\mathit{received}, \operatorname{lst}({[}22, \operatorname{encr}(\operatorname{lst}({[}20, 32]), 30)])]], {[}{[}\mathit{received}, \operatorname{lst}({[}22, \operatorname{encr}(\operatorname{lst}({[}20, 33]), 30), \operatorname{encr}(\operatorname{lst}({[}21, 33]), 31)])], {[}\mathit{sent}, 22], {[}\mathit{sent}, \operatorname{encr}(\operatorname{lst}({[}20, 33]), 30)], {[}\mathit{sent}, \operatorname{encr}(\operatorname{lst}({[}21, 33]), 31)]], {[}{[}\mathit{received}, 22], {[}\mathit{received}, \operatorname{encr}(\operatorname{lst}({[}20, 33]), 30)], {[}\mathit{sent}, \operatorname{lst}({[}22, \operatorname{encr}(\operatorname{lst}({[}20, 33]), 30)])]]])), (\lambda \operatorname{cxk\_} \in \operatorname{interval}(30, 33).\; \operatorname{cxk\_})] \end{array}
\end{equation*}

\par\medskip\noindent\textbf{Definition} (\texttt{cx-isp}).\; $\operatorname{cx\text{-}isp} = {[}\operatorname{cx\text{-}sp}, \operatorname{interval}(5, 8), \emptyset]$
\par\medskip\noindent\textbf{Definition} (\texttt{cx-orp}).\;
\begin{equation*}
\operatorname{cx\text{-}orp} = {[}\operatorname{cx\text{-}isp}, \operatorname{interval}(10, 11), (\lambda \operatorname{cxn\_} \in \operatorname{interval}(10, 11).\; +(\operatorname{cxn\_}, 20))]
\end{equation*}

\par\medskip\noindent\textbf{Definition} (\texttt{cx-c}).\;
\begin{equation*}
\begin{array}{@{}l@{}} \operatorname{cx\text{-}c} \\ =\; {[}\operatorname{cx\text{-}sp}, \operatorname{nodes}(\operatorname{cx\text{-}sp}), \{\operatorname{cxe\_} \in (\operatorname{nodes}(\operatorname{cx\text{-}sp}) \times \operatorname{nodes}(\operatorname{cx\text{-}sp})) : ((\operatorname{cxe\_} = {[}{[}1, 1], {[}1, 2]]) \vee ((\operatorname{cxe\_} = {[}{[}2, 2], {[}1, 5]]) \vee ((\operatorname{cxe\_} = {[}{[}2, 5], {[}1, 3]]) \vee ((\operatorname{cxe\_} = {[}{[}3, 5], {[}1, 4]]) \vee ((\operatorname{cxe\_} = {[}{[}2, 3], {[}3, 2]]) \vee ((\operatorname{cxe\_} = {[}{[}4, 2], {[}1, 6]]) \vee ((\operatorname{cxe\_} = {[}{[}2, 4], {[}1, 7]]) \vee ((\operatorname{cxe\_} = {[}{[}2, 7], {[}1, 8]]) \vee ((\operatorname{cxe\_} = {[}{[}3, 7], {[}2, 8]]) \vee (\operatorname{cxe\_} = {[}{[}3, 8], {[}2, 1]]))))))))))\}] \end{array}
\end{equation*}

\vnbdef{cx-isp@adv-tokens}{\vnbstmt{(\operatorname{adv\text{-}tokens}(\operatorname{cx\text{-}isp}) \simeq \emptyset)}}
\vnbdef{cx-orp@ltkey}{\vnbstmt{(\operatorname{ltkey}(\operatorname{cx\text{-}orp}) \simeq (\lambda n \in \operatorname{interval}(10, 11).\; (n + 20)))}}
\vnbdef{cx-isp@adv-ids}{\vnbstmt{(\operatorname{adv\text{-}ids}(\operatorname{cx\text{-}isp}) \simeq \operatorname{interval}(5, 8))}}
\vnbdef{cx-orp@sk-inf}{\vnbstmt{(\operatorname{sk\text{-}inf}(\operatorname{cx\text{-}orp}) \simeq \operatorname{cx\text{-}isp})}}
\vnbdef{cx-isp@ispace}{\vnbstmt{(\operatorname{ispace}(\operatorname{cx\text{-}isp}) \simeq \operatorname{cx\text{-}sp})}}
\vnbdef{cx-c@bsend}{\vnbstmt{(\operatorname{bsend}(\operatorname{cx\text{-}c}) \simeq \{e \in (\operatorname{nodes}(\operatorname{cx\text{-}sp}) \times \operatorname{nodes}(\operatorname{cx\text{-}sp})) : ((e = {[}{[}1, 1], {[}1, 2]]) \vee ((e = {[}{[}2, 2], {[}1, 5]]) \vee ((e = {[}{[}2, 5], {[}1, 3]]) \vee ((e = {[}{[}3, 5], {[}1, 4]]) \vee ((e = {[}{[}2, 3], {[}3, 2]]) \vee ((e = {[}{[}4, 2], {[}1, 6]]) \vee ((e = {[}{[}2, 4], {[}1, 7]]) \vee ((e = {[}{[}2, 7], {[}1, 8]]) \vee ((e = {[}{[}3, 7], {[}2, 8]]) \vee (e = {[}{[}3, 8], {[}2, 1]]))))))))))\})}}
\vnbdef{cx-c@bspace}{\vnbstmt{(\operatorname{bspace}(\operatorname{cx\text{-}c}) \simeq \operatorname{cx\text{-}sp})}}
\vnbdef{cx-sp@keys}{\vnbstmt{(\operatorname{keys}(\operatorname{cx\text{-}sp}) \simeq \operatorname{interval}(30, 33))}}
\vnbdef{cx-sp@tokens}{\vnbstmt{(\operatorname{tokens}(\operatorname{cx\text{-}sp}) \simeq \mathbb{N})}}
\vnbdef{cx-sp@kinv}{\vnbstmt{(\operatorname{kinv}(\operatorname{cx\text{-}sp}) \simeq (\lambda k \in \operatorname{interval}(30, 33).\; k))}}
\vnbdef{cx-sp@ids}{\vnbstmt{(\operatorname{ids}(\operatorname{cx\text{-}sp}) \simeq \operatorname{interval}(1, 8))}}
\vnbdef{cx-orp@sk-names}{\vnbstmt{(\operatorname{sk\text{-}names}(\operatorname{cx\text{-}orp}) \simeq \operatorname{interval}(10, 11))}}
\vnbdef{cx-sp@strand-map}{\vnbstmt{(\operatorname{strand\text{-}map}(\operatorname{cx\text{-}sp}) \simeq (\lambda i \in \operatorname{interval}(1, 8).\; \operatorname{nth}(i, {[}{[}{[}\mathit{sent}, \operatorname{lst}({[}22, 10, 11, \operatorname{encr}(\operatorname{lst}({[}20, 22, 10, 11]), 30)])], {[}\mathit{received}, \operatorname{lst}({[}22, \operatorname{encr}(\operatorname{lst}({[}20, 33]), 30)])]], {[}{[}\mathit{received}, \operatorname{lst}({[}22, 10, 11, \operatorname{encr}(\operatorname{lst}({[}20, 22, 10, 11]), 30)])], {[}\mathit{sent}, \operatorname{lst}({[}22, 10, 11, \operatorname{encr}(\operatorname{lst}({[}20, 22, 10, 11]), 30), \operatorname{encr}(\operatorname{lst}({[}21, 22, 10, 11]), 31)])], {[}\mathit{received}, \operatorname{lst}({[}22, \operatorname{encr}(\operatorname{lst}({[}20, 32]), 30), \operatorname{encr}(\operatorname{lst}({[}21, 32]), 31)])], {[}\mathit{sent}, \operatorname{lst}({[}22, \operatorname{encr}(\operatorname{lst}({[}20, 32]), 30)])]], {[}{[}\mathit{received}, \operatorname{lst}({[}22, 10, 11, \operatorname{encr}(\operatorname{lst}({[}20, 22, 10, 11]), 30), \operatorname{encr}(\operatorname{lst}({[}21, 22, 10, 11]), 31)])], {[}\mathit{sent}, \operatorname{lst}({[}22, \operatorname{encr}(\operatorname{lst}({[}20, 32]), 30), \operatorname{encr}(\operatorname{lst}({[}21, 32]), 31)])]], {[}{[}\mathit{received}, \operatorname{lst}({[}22, 10, 11, \operatorname{encr}(\operatorname{lst}({[}20, 22, 10, 11]), 30), \operatorname{encr}(\operatorname{lst}({[}21, 22, 10, 11]), 31)])], {[}\mathit{sent}, \operatorname{lst}({[}22, \operatorname{encr}(\operatorname{lst}({[}20, 33]), 30), \operatorname{encr}(\operatorname{lst}({[}21, 33]), 31)])]], {[}{[}\mathit{received}, \operatorname{lst}({[}22, 10, 11, \operatorname{encr}(\operatorname{lst}({[}20, 22, 10, 11]), 30), \operatorname{encr}(\operatorname{lst}({[}21, 22, 10, 11]), 31)])], {[}\mathit{sent}, \operatorname{lst}({[}22, 10, 11, \operatorname{encr}(\operatorname{lst}({[}20, 22, 10, 11]), 30), \operatorname{encr}(\operatorname{lst}({[}21, 22, 10, 11]), 31)])], {[}\mathit{sent}, \operatorname{lst}({[}22, 10, 11, \operatorname{encr}(\operatorname{lst}({[}20, 22, 10, 11]), 30), \operatorname{encr}(\operatorname{lst}({[}21, 22, 10, 11]), 31)])]], {[}{[}\mathit{received}, \operatorname{lst}({[}22, \operatorname{encr}(\operatorname{lst}({[}20, 32]), 30)])]], {[}{[}\mathit{received}, \operatorname{lst}({[}22, \operatorname{encr}(\operatorname{lst}({[}20, 33]), 30), \operatorname{encr}(\operatorname{lst}({[}21, 33]), 31)])], {[}\mathit{sent}, 22], {[}\mathit{sent}, \operatorname{encr}(\operatorname{lst}({[}20, 33]), 30)], {[}\mathit{sent}, \operatorname{encr}(\operatorname{lst}({[}21, 33]), 31)]], {[}{[}\mathit{received}, 22], {[}\mathit{received}, \operatorname{encr}(\operatorname{lst}({[}20, 33]), 30)], {[}\mathit{sent}, \operatorname{lst}({[}22, \operatorname{encr}(\operatorname{lst}({[}20, 33]), 30)])]]])))}}
\vnbdef{cx-c@bnodes}{\vnbstmt{(\operatorname{bnodes}(\operatorname{cx\text{-}c}) \simeq \operatorname{nodes}(\operatorname{cx\text{-}sp}))}}

\vnbbookfile{theorem-library/or-missing-guarantee}{theorem-library/or-missing-guarantee}
\noindent{\small cx-laws.scm -- THE MISSING GUARANTEE OF OTWAY-REES (strand-spaces section 7.3.3, Figure 9): the concrete bundle CX-C of structure-library/or-missing-guarantee-instance.scm satisfies every hypothesis of Theorem 7.8 (or-initiator-guarantee) for the initiator with session key K' = 33, yet no regular responder strand of ...}\par
\vnbprop{cx-strand-list-typed}{\vnbstmt{\begin{array}{@{}l@{}} {[}{[}{[}\mathit{sent}, \operatorname{lst}({[}22, 10, 11, \operatorname{encr}(\operatorname{lst}({[}20, 22, 10, 11]), 30)])], {[}\mathit{received}, \operatorname{lst}({[}22, \operatorname{encr}(\operatorname{lst}({[}20, 33]), 30)])]], {[}{[}\mathit{received}, \operatorname{lst}({[}22, 10, 11, \operatorname{encr}(\operatorname{lst}({[}20, 22, 10, 11]), 30)])], {[}\mathit{sent}, \operatorname{lst}({[}22, 10, 11, \operatorname{encr}(\operatorname{lst}({[}20, 22, 10, 11]), 30), \operatorname{encr}(\operatorname{lst}({[}21, 22, 10, 11]), 31)])], {[}\mathit{received}, \operatorname{lst}({[}22, \operatorname{encr}(\operatorname{lst}({[}20, 32]), 30), \operatorname{encr}(\operatorname{lst}({[}21, 32]), 31)])], {[}\mathit{sent}, \operatorname{lst}({[}22, \operatorname{encr}(\operatorname{lst}({[}20, 32]), 30)])]], {[}{[}\mathit{received}, \operatorname{lst}({[}22, 10, 11, \operatorname{encr}(\operatorname{lst}({[}20, 22, 10, 11]), 30), \operatorname{encr}(\operatorname{lst}({[}21, 22, 10, 11]), 31)])], {[}\mathit{sent}, \operatorname{lst}({[}22, \operatorname{encr}(\operatorname{lst}({[}20, 32]), 30), \operatorname{encr}(\operatorname{lst}({[}21, 32]), 31)])]], {[}{[}\mathit{received}, \operatorname{lst}({[}22, 10, 11, \operatorname{encr}(\operatorname{lst}({[}20, 22, 10, 11]), 30), \operatorname{encr}(\operatorname{lst}({[}21, 22, 10, 11]), 31)])], {[}\mathit{sent}, \operatorname{lst}({[}22, \operatorname{encr}(\operatorname{lst}({[}20, 33]), 30), \operatorname{encr}(\operatorname{lst}({[}21, 33]), 31)])]], {[}{[}\mathit{received}, \operatorname{lst}({[}22, 10, 11, \operatorname{encr}(\operatorname{lst}({[}20, 22, 10, 11]), 30), \operatorname{encr}(\operatorname{lst}({[}21, 22, 10, 11]), 31)])], {[}\mathit{sent}, \operatorname{lst}({[}22, 10, 11, \operatorname{encr}(\operatorname{lst}({[}20, 22, 10, 11]), 30), \operatorname{encr}(\operatorname{lst}({[}21, 22, 10, 11]), 31)])], {[}\mathit{sent}, \operatorname{lst}({[}22, 10, 11, \operatorname{encr}(\operatorname{lst}({[}20, 22, 10, 11]), 30), \operatorname{encr}(\operatorname{lst}({[}21, 22, 10, 11]), 31)])]], {[}{[}\mathit{received}, \operatorname{lst}({[}22, \operatorname{encr}(\operatorname{lst}({[}20, 32]), 30)])]], {[}{[}\mathit{received}, \operatorname{lst}({[}22, \operatorname{encr}(\operatorname{lst}({[}20, 33]), 30), \operatorname{encr}(\operatorname{lst}({[}21, 33]), 31)])], {[}\mathit{sent}, 22], {[}\mathit{sent}, \operatorname{encr}(\operatorname{lst}({[}20, 33]), 30)], {[}\mathit{sent}, \operatorname{encr}(\operatorname{lst}({[}21, 33]), 31)]], {[}{[}\mathit{received}, 22], {[}\mathit{received}, \operatorname{encr}(\operatorname{lst}({[}20, 33]), 30)], {[}\mathit{sent}, \operatorname{lst}({[}22, \operatorname{encr}(\operatorname{lst}({[}20, 33]), 30)])]]] \\ \in\; \operatorname{tuples}(\operatorname{strands}(\mathbb{N}, \operatorname{interval}(30, 33))) \end{array}}}
\vnbprop{cx-map-1}{\vnbstmt{(\operatorname{strand\text{-}map}(\operatorname{cx\text{-}sp})(1) \simeq {[}{[}\mathit{sent}, \operatorname{lst}({[}22, 10, 11, \operatorname{encr}(\operatorname{lst}({[}20, 22, 10, 11]), 30)])], {[}\mathit{received}, \operatorname{lst}({[}22, \operatorname{encr}(\operatorname{lst}({[}20, 33]), 30)])]])}}
\vnbprop{cx-map-2}{\vnbstmt{(\operatorname{strand\text{-}map}(\operatorname{cx\text{-}sp})(2) \simeq {[}{[}\mathit{received}, \operatorname{lst}({[}22, 10, 11, \operatorname{encr}(\operatorname{lst}({[}20, 22, 10, 11]), 30)])], {[}\mathit{sent}, \operatorname{lst}({[}22, 10, 11, \operatorname{encr}(\operatorname{lst}({[}20, 22, 10, 11]), 30), \operatorname{encr}(\operatorname{lst}({[}21, 22, 10, 11]), 31)])], {[}\mathit{received}, \operatorname{lst}({[}22, \operatorname{encr}(\operatorname{lst}({[}20, 32]), 30), \operatorname{encr}(\operatorname{lst}({[}21, 32]), 31)])], {[}\mathit{sent}, \operatorname{lst}({[}22, \operatorname{encr}(\operatorname{lst}({[}20, 32]), 30)])]])}}
\vnbprop{cx-map-3}{\vnbstmt{(\operatorname{strand\text{-}map}(\operatorname{cx\text{-}sp})(3) \simeq {[}{[}\mathit{received}, \operatorname{lst}({[}22, 10, 11, \operatorname{encr}(\operatorname{lst}({[}20, 22, 10, 11]), 30), \operatorname{encr}(\operatorname{lst}({[}21, 22, 10, 11]), 31)])], {[}\mathit{sent}, \operatorname{lst}({[}22, \operatorname{encr}(\operatorname{lst}({[}20, 32]), 30), \operatorname{encr}(\operatorname{lst}({[}21, 32]), 31)])]])}}
\vnbprop{cx-map-4}{\vnbstmt{(\operatorname{strand\text{-}map}(\operatorname{cx\text{-}sp})(4) \simeq {[}{[}\mathit{received}, \operatorname{lst}({[}22, 10, 11, \operatorname{encr}(\operatorname{lst}({[}20, 22, 10, 11]), 30), \operatorname{encr}(\operatorname{lst}({[}21, 22, 10, 11]), 31)])], {[}\mathit{sent}, \operatorname{lst}({[}22, \operatorname{encr}(\operatorname{lst}({[}20, 33]), 30), \operatorname{encr}(\operatorname{lst}({[}21, 33]), 31)])]])}}
\vnbprop{cx-map-5}{\vnbstmt{(\operatorname{strand\text{-}map}(\operatorname{cx\text{-}sp})(5) \simeq {[}{[}\mathit{received}, \operatorname{lst}({[}22, 10, 11, \operatorname{encr}(\operatorname{lst}({[}20, 22, 10, 11]), 30), \operatorname{encr}(\operatorname{lst}({[}21, 22, 10, 11]), 31)])], {[}\mathit{sent}, \operatorname{lst}({[}22, 10, 11, \operatorname{encr}(\operatorname{lst}({[}20, 22, 10, 11]), 30), \operatorname{encr}(\operatorname{lst}({[}21, 22, 10, 11]), 31)])], {[}\mathit{sent}, \operatorname{lst}({[}22, 10, 11, \operatorname{encr}(\operatorname{lst}({[}20, 22, 10, 11]), 30), \operatorname{encr}(\operatorname{lst}({[}21, 22, 10, 11]), 31)])]])}}
\vnbprop{cx-map-6}{\vnbstmt{(\operatorname{strand\text{-}map}(\operatorname{cx\text{-}sp})(6) \simeq {[}{[}\mathit{received}, \operatorname{lst}({[}22, \operatorname{encr}(\operatorname{lst}({[}20, 32]), 30)])]])}}
\vnbprop{cx-map-7}{\vnbstmt{(\operatorname{strand\text{-}map}(\operatorname{cx\text{-}sp})(7) \simeq {[}{[}\mathit{received}, \operatorname{lst}({[}22, \operatorname{encr}(\operatorname{lst}({[}20, 33]), 30), \operatorname{encr}(\operatorname{lst}({[}21, 33]), 31)])], {[}\mathit{sent}, 22], {[}\mathit{sent}, \operatorname{encr}(\operatorname{lst}({[}20, 33]), 30)], {[}\mathit{sent}, \operatorname{encr}(\operatorname{lst}({[}21, 33]), 31)]])}}
\vnbprop{cx-map-8}{\vnbstmt{(\operatorname{strand\text{-}map}(\operatorname{cx\text{-}sp})(8) \simeq {[}{[}\mathit{received}, 22], {[}\mathit{received}, \operatorname{encr}(\operatorname{lst}({[}20, 33]), 30)], {[}\mathit{sent}, \operatorname{lst}({[}22, \operatorname{encr}(\operatorname{lst}({[}20, 33]), 30)])]])}}
\vnbprop{cx-sp-is-strand-space}{\vnbstmt{\operatorname{is\text{-}strand\text{-}space}(\operatorname{cx\text{-}sp})}}
\vnbprop{cx-isp-is-infiltrated-space}{\vnbstmt{\operatorname{is\text{-}infiltrated\text{-}space}(\operatorname{cx\text{-}isp})}}
\vnbprop{cx-orp-is-sk-space}{\vnbstmt{\operatorname{is\text{-}sk\text{-}space}(\operatorname{cx\text{-}orp})}}
\vnbprop{cx-init-1}{\vnbstmt{\operatorname{is\text{-}or\text{-}init}(\operatorname{cx\text{-}orp}, \operatorname{strand\text{-}map}(\operatorname{ispace}(\operatorname{sk\text{-}inf}(\operatorname{cx\text{-}orp})))(1), 10, 11, 20, 22, 33)}}
\vnbprop{cx-resp-2}{\vnbstmt{\begin{array}{@{}l@{}} \operatorname{is\text{-}or\text{-}resp}(\operatorname{cx\text{-}orp}, \\ \quad \operatorname{strand\text{-}map}(\operatorname{ispace}(\operatorname{sk\text{-}inf}(\operatorname{cx\text{-}orp})))(2), \\ \quad 10, 11, 21, 22, 32, \operatorname{encr}(\operatorname{lst}({[}20, 22, 10, 11]), 30), \\ \quad \operatorname{encr}(\operatorname{lst}({[}20, 32]), 30)) \end{array}}}
\vnbprop{cx-serv-3}{\vnbstmt{\operatorname{is\text{-}or\text{-}serv}(\operatorname{cx\text{-}orp}, \operatorname{strand\text{-}map}(\operatorname{ispace}(\operatorname{sk\text{-}inf}(\operatorname{cx\text{-}orp})))(3), 10, 11, 20, 21, 22, 32)}}
\vnbprop{cx-serv-4}{\vnbstmt{\operatorname{is\text{-}or\text{-}serv}(\operatorname{cx\text{-}orp}, \operatorname{strand\text{-}map}(\operatorname{ispace}(\operatorname{sk\text{-}inf}(\operatorname{cx\text{-}orp})))(4), 10, 11, 20, 21, 22, 33)}}
\vnbprop{cx-orp-is-or-space}{\vnbstmt{\operatorname{is\text{-}or\text{-}space}(\operatorname{cx\text{-}orp})}}
\vnbprop{cx-nodes-set}{\vnbstmt{\operatorname{nodes}(\operatorname{cx\text{-}sp}) \in \mathbf{Set}}}
\vnbprop{cx-edge-1}{\vnbstmt{\operatorname{sends\text{-}to}(\operatorname{cx\text{-}sp}, {[}1, 1], {[}1, 2])}}
\vnbprop{cx-edge-2}{\vnbstmt{\operatorname{sends\text{-}to}(\operatorname{cx\text{-}sp}, {[}2, 2], {[}1, 5])}}
\vnbprop{cx-edge-3}{\vnbstmt{\operatorname{sends\text{-}to}(\operatorname{cx\text{-}sp}, {[}2, 5], {[}1, 3])}}
\vnbprop{cx-edge-4}{\vnbstmt{\operatorname{sends\text{-}to}(\operatorname{cx\text{-}sp}, {[}3, 5], {[}1, 4])}}
\vnbprop{cx-edge-5}{\vnbstmt{\operatorname{sends\text{-}to}(\operatorname{cx\text{-}sp}, {[}2, 3], {[}3, 2])}}
\vnbprop{cx-edge-6}{\vnbstmt{\operatorname{sends\text{-}to}(\operatorname{cx\text{-}sp}, {[}4, 2], {[}1, 6])}}
\vnbprop{cx-edge-7}{\vnbstmt{\operatorname{sends\text{-}to}(\operatorname{cx\text{-}sp}, {[}2, 4], {[}1, 7])}}
\vnbprop{cx-edge-8}{\vnbstmt{\operatorname{sends\text{-}to}(\operatorname{cx\text{-}sp}, {[}2, 7], {[}1, 8])}}
\vnbprop{cx-edge-9}{\vnbstmt{\operatorname{sends\text{-}to}(\operatorname{cx\text{-}sp}, {[}3, 7], {[}2, 8])}}
\vnbprop{cx-edge-10}{\vnbstmt{\operatorname{sends\text{-}to}(\operatorname{cx\text{-}sp}, {[}3, 8], {[}2, 1])}}
\vnbprop{cx-nodes-in-box}{\vnbstmt{\forall y \in \operatorname{nodes}(\operatorname{cx\text{-}sp}).\; (y \in (\operatorname{interval}(1, 4) \times \operatorname{interval}(1, 8)))}}
\vnbprop{cx-nodes-finite}{\vnbstmt{|\operatorname{nodes}(\operatorname{cx\text{-}sp})| \in \mathbb{N}}}
\vnbprop{cx-unique-sender}{\vnbstmt{\begin{array}{@{}l@{}} \forall y \in \operatorname{nodes}(\operatorname{cx\text{-}sp}). \\ \quad \operatorname{is\text{-}received}(\operatorname{cx\text{-}sp}, y) \Rightarrow  \\ \quad \quad \exists x. \\ \quad \quad \quad {[}x, y] \\ \quad \quad \quad \in\; \{e \in (\operatorname{nodes}(\operatorname{cx\text{-}sp}) \times \operatorname{nodes}(\operatorname{cx\text{-}sp})) : ((e = {[}{[}1, 1], {[}1, 2]]) \vee ((e = {[}{[}2, 2], {[}1, 5]]) \vee ((e = {[}{[}2, 5], {[}1, 3]]) \vee ((e = {[}{[}3, 5], {[}1, 4]]) \vee ((e = {[}{[}2, 3], {[}3, 2]]) \vee ((e = {[}{[}4, 2], {[}1, 6]]) \vee ((e = {[}{[}2, 4], {[}1, 7]]) \vee ((e = {[}{[}2, 7], {[}1, 8]]) \vee ((e = {[}{[}3, 7], {[}2, 8]]) \vee (e = {[}{[}3, 8], {[}2, 1]]))))))))))\} \wedge \\ \quad \quad \quad \forall z. \\ \quad \quad \quad \quad {[}z, y] \\ \quad \quad \quad \quad \in\; \{e \in (\operatorname{nodes}(\operatorname{cx\text{-}sp}) \times \operatorname{nodes}(\operatorname{cx\text{-}sp})) : ((e = {[}{[}1, 1], {[}1, 2]]) \vee ((e = {[}{[}2, 2], {[}1, 5]]) \vee ((e = {[}{[}2, 5], {[}1, 3]]) \vee ((e = {[}{[}3, 5], {[}1, 4]]) \vee ((e = {[}{[}2, 3], {[}3, 2]]) \vee ((e = {[}{[}4, 2], {[}1, 6]]) \vee ((e = {[}{[}2, 4], {[}1, 7]]) \vee ((e = {[}{[}2, 7], {[}1, 8]]) \vee ((e = {[}{[}3, 7], {[}2, 8]]) \vee (e = {[}{[}3, 8], {[}2, 1]]))))))))))\} \Rightarrow \\ \quad \quad \quad \quad \quad z = x \end{array}}}
\vnbprop{cx-rank-type}{\vnbstmt{\begin{array}{@{}l@{}} (\lambda r \in \operatorname{nodes}(\operatorname{cx\text{-}sp}).\; \\ \quad \operatorname{nth}(\operatorname{nth}(2, r), {[}6, 20, 43, 43, 40, 60, 45, 48]) + (\operatorname{nth}(\operatorname{nth}(2, r), {[}23, 10, 1, 1, 1, 1, 1, 1]) \cdot \operatorname{nth}(1, r))) \\ \in\; (\operatorname{nodes}(\operatorname{cx\text{-}sp}) \to \mathbb{R}) \end{array}}}
\vnbprop{cx-rank-rel-power}{\vnbstmt{\begin{array}{@{}l@{}} \{q \in (\operatorname{nodes}(\operatorname{cx\text{-}sp}) \times \operatorname{nodes}(\operatorname{cx\text{-}sp})) : ((\lambda r \in \operatorname{nodes}(\operatorname{cx\text{-}sp}).\; (\operatorname{nth}(\operatorname{nth}(2, r), {[}6, 20, 43, 43, 40, 60, 45, 48]) + (\operatorname{nth}(\operatorname{nth}(2, r), {[}23, 10, 1, 1, 1, 1, 1, 1]) \cdot \operatorname{nth}(1, r))))(\operatorname{nth}(1, q)) < (\lambda r \in \operatorname{nodes}(\operatorname{cx\text{-}sp}).\; (\operatorname{nth}(\operatorname{nth}(2, r), {[}6, 20, 43, 43, 40, 60, 45, 48]) + (\operatorname{nth}(\operatorname{nth}(2, r), {[}23, 10, 1, 1, 1, 1, 1, 1]) \cdot \operatorname{nth}(1, r))))(\operatorname{nth}(2, q)))\} \\ \in\; \operatorname{power}((\operatorname{nodes}(\operatorname{cx\text{-}sp}) \times \operatorname{nodes}(\operatorname{cx\text{-}sp}))) \end{array}}}
\vnbprop{cx-rank-closed}{\vnbstmt{\begin{array}{@{}l@{}} \operatorname{is\text{-}prec\text{-}closed}(\operatorname{cx\text{-}sp}, \operatorname{nodes}(\operatorname{cx\text{-}sp}), \\ \quad \{e \in (\operatorname{nodes}(\operatorname{cx\text{-}sp}) \times \operatorname{nodes}(\operatorname{cx\text{-}sp})) : ((e = {[}{[}1, 1], {[}1, 2]]) \vee ((e = {[}{[}2, 2], {[}1, 5]]) \vee ((e = {[}{[}2, 5], {[}1, 3]]) \vee ((e = {[}{[}3, 5], {[}1, 4]]) \vee ((e = {[}{[}2, 3], {[}3, 2]]) \vee ((e = {[}{[}4, 2], {[}1, 6]]) \vee ((e = {[}{[}2, 4], {[}1, 7]]) \vee ((e = {[}{[}2, 7], {[}1, 8]]) \vee ((e = {[}{[}3, 7], {[}2, 8]]) \vee (e = {[}{[}3, 8], {[}2, 1]]))))))))))\}, \\ \quad \{q \in (\operatorname{nodes}(\operatorname{cx\text{-}sp}) \times \operatorname{nodes}(\operatorname{cx\text{-}sp})) : ((\lambda r \in \operatorname{nodes}(\operatorname{cx\text{-}sp}).\; (\operatorname{nth}(\operatorname{nth}(2, r), {[}6, 20, 43, 43, 40, 60, 45, 48]) + (\operatorname{nth}(\operatorname{nth}(2, r), {[}23, 10, 1, 1, 1, 1, 1, 1]) \cdot \operatorname{nth}(1, r))))(\operatorname{nth}(1, q)) < (\lambda r \in \operatorname{nodes}(\operatorname{cx\text{-}sp}).\; (\operatorname{nth}(\operatorname{nth}(2, r), {[}6, 20, 43, 43, 40, 60, 45, 48]) + (\operatorname{nth}(\operatorname{nth}(2, r), {[}23, 10, 1, 1, 1, 1, 1, 1]) \cdot \operatorname{nth}(1, r))))(\operatorname{nth}(2, q)))\}) \end{array}}}
\vnbprop{cx-acyclic}{\vnbstmt{\begin{array}{@{}l@{}} \forall x. \\ \quad \neg ({[}x, x] \in \operatorname{prec\text{-}strict}(\operatorname{cx\text{-}sp}, \operatorname{nodes}(\operatorname{cx\text{-}sp}), \{e \in (\operatorname{nodes}(\operatorname{cx\text{-}sp}) \times \operatorname{nodes}(\operatorname{cx\text{-}sp})) : ((e = {[}{[}1, 1], {[}1, 2]]) \vee ((e = {[}{[}2, 2], {[}1, 5]]) \vee ((e = {[}{[}2, 5], {[}1, 3]]) \vee ((e = {[}{[}3, 5], {[}1, 4]]) \vee ((e = {[}{[}2, 3], {[}3, 2]]) \vee ((e = {[}{[}4, 2], {[}1, 6]]) \vee ((e = {[}{[}2, 4], {[}1, 7]]) \vee ((e = {[}{[}2, 7], {[}1, 8]]) \vee ((e = {[}{[}3, 7], {[}2, 8]]) \vee (e = {[}{[}3, 8], {[}2, 1]]))))))))))\})) \end{array}}}
\vnbprop{cx-c-is-bundle}{\vnbstmt{\operatorname{is\text{-}bundle}(\operatorname{cx\text{-}c})}}
\vnbprop{cx-na-uniquely-originates}{\vnbstmt{\operatorname{uniquely\text{-}originates}(\operatorname{cx\text{-}sp}, 20)}}
\vnbprop{or-missing-guarantee-bundle}{\vnbstmt{\begin{array}{@{}l@{}} \operatorname{is\text{-}or\text{-}init}(\operatorname{cx\text{-}orp}, \operatorname{strand\text{-}map}(\operatorname{ispace}(\operatorname{sk\text{-}inf}(\operatorname{cx\text{-}orp})))(1), 10, 11, 20, 22, 33) \wedge \\ {[}2, 1] \in \operatorname{bnodes}(\operatorname{cx\text{-}c}) \wedge \\ \operatorname{is\text{-}or\text{-}resp}(\operatorname{cx\text{-}orp}, \\ \quad \operatorname{strand\text{-}map}(\operatorname{ispace}(\operatorname{sk\text{-}inf}(\operatorname{cx\text{-}orp})))(2), \\ \quad 10, 11, 21, 22, 32, \operatorname{encr}(\operatorname{lst}({[}20, 22, 10, 11]), 30), \\ \quad \operatorname{encr}(\operatorname{lst}({[}20, 32]), 30)) \wedge \\ {[}4, 2] \in \operatorname{bnodes}(\operatorname{cx\text{-}c}) \wedge \\ \neg (32 = 33) \end{array}}}
\vnbprop{or-initiator-guarantee-key-not-determined}{\vnbstmt{\begin{array}{@{}l@{}} \exists p, c, s, a, b, d, m, k. \\ \quad \operatorname{is\text{-}or\text{-}space}(p) \wedge \\ \quad \operatorname{is\text{-}bundle}(c) \wedge \\ \quad \operatorname{bspace}(c) = \operatorname{ispace}(\operatorname{sk\text{-}inf}(p)) \wedge \\ \quad s \in \operatorname{ids}(\operatorname{ispace}(\operatorname{sk\text{-}inf}(p))) \wedge \\ \quad \neg (a = b) \wedge \\ \quad \operatorname{uniquely\text{-}originates}(\operatorname{ispace}(\operatorname{sk\text{-}inf}(p)), d) \wedge \\ \quad \neg (\operatorname{ltkey}(p)(a) \in \operatorname{adv\text{-}tokens}(\operatorname{sk\text{-}inf}(p))) \wedge \\ \quad \neg (\operatorname{ltkey}(p)(b) \in \operatorname{adv\text{-}tokens}(\operatorname{sk\text{-}inf}(p))) \wedge \\ \quad \operatorname{is\text{-}or\text{-}init}(p, \operatorname{strand\text{-}map}(\operatorname{ispace}(\operatorname{sk\text{-}inf}(p)))(s), a, b, d, m, k) \wedge \\ \quad {[}2, s] \in \operatorname{bnodes}(c) \wedge \\ \quad \neg \exists r \in \operatorname{ids}(\operatorname{ispace}(\operatorname{sk\text{-}inf}(p))).\; (\neg (r \in \operatorname{adv\text{-}ids}(\operatorname{sk\text{-}inf}(p))) \wedge (({[}2, r] \in \operatorname{bnodes}(c)) \wedge \exists f, g, h.\; \operatorname{is\text{-}or\text{-}resp}(p, \operatorname{strand\text{-}map}(\operatorname{ispace}(\operatorname{sk\text{-}inf}(p)))(r), a, b, f, m, k, g, h))) \end{array}}}
\vnbprop{cx-nb-uniquely-originates}{\vnbstmt{\operatorname{uniquely\text{-}originates}(\operatorname{cx\text{-}sp}, 21)}}
\vnbprop{or-responder-guarantee-key-not-determined}{\vnbstmt{\begin{array}{@{}l@{}} \exists p, c, s, a, b, d, m, k, h, g. \\ \quad \operatorname{is\text{-}or\text{-}space}(p) \wedge \\ \quad \operatorname{is\text{-}bundle}(c) \wedge \\ \quad \operatorname{bspace}(c) = \operatorname{ispace}(\operatorname{sk\text{-}inf}(p)) \wedge \\ \quad s \in \operatorname{ids}(\operatorname{ispace}(\operatorname{sk\text{-}inf}(p))) \wedge \\ \quad \neg (a = b) \wedge \\ \quad \operatorname{uniquely\text{-}originates}(\operatorname{ispace}(\operatorname{sk\text{-}inf}(p)), d) \wedge \\ \quad \neg (\operatorname{ltkey}(p)(a) \in \operatorname{adv\text{-}tokens}(\operatorname{sk\text{-}inf}(p))) \wedge \\ \quad \neg (\operatorname{ltkey}(p)(b) \in \operatorname{adv\text{-}tokens}(\operatorname{sk\text{-}inf}(p))) \wedge \\ \quad \operatorname{is\text{-}or\text{-}resp}(p, \operatorname{strand\text{-}map}(\operatorname{ispace}(\operatorname{sk\text{-}inf}(p)))(s), a, b, d, m, k, h, g) \wedge \\ \quad {[}3, s] \in \operatorname{bnodes}(c) \wedge \\ \quad \neg \exists i \in \operatorname{ids}(\operatorname{ispace}(\operatorname{sk\text{-}inf}(p))).\; (\neg (i \in \operatorname{adv\text{-}ids}(\operatorname{sk\text{-}inf}(p))) \wedge (({[}1, i] \in \operatorname{bnodes}(c)) \wedge \exists f.\; \operatorname{is\text{-}or\text{-}init}(p, \operatorname{strand\text{-}map}(\operatorname{ispace}(\operatorname{sk\text{-}inf}(p)))(i), a, b, f, m, k))) \end{array}}}

\part{The Kernel Encoded}
\noindent Expressions of the VNB language encoded as sets (syntax, free variables, substitution, rank); the core calculus as an inductively defined set DERIVABLE(T) of sequents; contextual rewriting and the rewriting rules as admissible rules of the core; the induction rules and the definedness clauses of the kernel against the core (K2, docs/meta-k2-admissibility-2026-10-08.md).
\par\noindent\textit{Builds on:} Foundations.
\par\noindent{\small 12 files.}

\vnbbookfile{structure-library/meta-syntax}{structure-library/meta-syntax}
\noindent{\small meta-syntax.scm -- THE SYNTAX OF THE KERNEL AS A VNB OBJECT: locally nameless expressions, their stages, and the four structural operations FV, SUBST, OPEN, LC. (docs/meta-k2-admissibility-2026-10-08.md, sections 5 and 6, phase 1: the cost probe.) DEFINITIONS ONLY; the laws are in theorem-library/meta-syntax-laws.scm.}\par
\par\medskip\noindent\textbf{Definition} (\texttt{efvar}).\;$\operatorname{efvar}(\operatorname{mxn\_}) :=$ ${[}0, \operatorname{mxn\_}]$
\par\medskip\noindent\textbf{Definition} (\texttt{ebvar}).\;$\operatorname{ebvar}(\operatorname{mxk\_}) :=$ ${[}1, \operatorname{mxk\_}]$
\par\medskip\noindent\textbf{Definition} (\texttt{econst}).\;$\operatorname{econst}(\operatorname{mxc\_}) :=$ ${[}2, \operatorname{mxc\_}]$
\par\medskip\noindent\textbf{Definition} (\texttt{eapp}).\;$\operatorname{eapp}(\operatorname{mxh\_}, \operatorname{mxl\_}) :=$ ${[}3, \operatorname{mxh\_}, \operatorname{mxl\_}]$
\par\medskip\noindent\textbf{Definition} (\texttt{ebind}).\;$\operatorname{ebind}(\operatorname{mxh\_}, \operatorname{mxn\_}, \operatorname{mxa\_}, \operatorname{mxb\_}) :=$ ${[}4, \operatorname{mxh\_}, \operatorname{mxn\_}, \operatorname{mxa\_}, \operatorname{mxb\_}]$
\par\medskip\noindent\textbf{Definition} (\texttt{expr-fv}).\;$\operatorname{expr\text{-}fv}(\operatorname{mxe\_}) :=$ $\operatorname{fv\text{-}stage}(\operatorname{expr\text{-}rank}(\operatorname{mxe\_}))(\operatorname{mxe\_})$
\par\medskip\noindent\textbf{Definition} (\texttt{expr-subst}).\;$\operatorname{expr\text{-}subst}(\operatorname{mxe\_}, \operatorname{mxm\_}, \operatorname{mxu\_}) :=$ $\operatorname{subst\text{-}stage}(\operatorname{mxm\_}, \operatorname{mxu\_}, \operatorname{expr\text{-}rank}(\operatorname{mxe\_}))(\operatorname{mxe\_})$
\par\medskip\noindent\textbf{Definition} (\texttt{expr-open}).\;$\operatorname{expr\text{-}open}(\operatorname{mxe\_}, \operatorname{mxk\_}, \operatorname{mxu\_}) :=$ $\operatorname{open\text{-}stage}(\operatorname{mxu\_}, \operatorname{expr\text{-}rank}(\operatorname{mxe\_}))(\operatorname{mxe\_}, \operatorname{mxk\_})$
\par\medskip\noindent\textbf{Definition} (\texttt{expr-lc}).\;For every $\operatorname{mxe\_}$,  $\operatorname{expr\text{-}lc}(\operatorname{mxe\_}) \Leftrightarrow (0 \in \operatorname{expr\text{-}lcd}(\operatorname{mxe\_}))$
\par\medskip\noindent\textbf{Definition} (\texttt{expr-set}).\; $(\operatorname{expr\text{-}set} \simeq \operatorname{big\text{-}union}(\operatorname{mxn\_}, \mathbb{N}, \operatorname{expr\text{-}stage}(\operatorname{mxn\_})))$
\par\medskip\noindent\textbf{Definition} (\texttt{expr-stage}).\;
\begin{equation*}
(\operatorname{expr\text{-}stage}(0) \simeq ((\{0\} \times \mathbb{N}) \cup ((\{1\} \times \mathbb{N}) \cup (\{2\} \times \mathbb{N}))))
\end{equation*}
\par
Suppose $\operatorname{mxn\_} \in \mathbb{N}$. Then
\begin{equation*}
(\operatorname{expr\text{-}stage}(\operatorname{succ}(\operatorname{mxn\_})) \simeq (\operatorname{expr\text{-}stage}(\operatorname{mxn\_}) \cup ((\{3\} \times \mathbb{N}) \cup (\{4\} \times \mathbb{N}))))
\end{equation*}

\par\medskip\noindent\textbf{Definition} (\texttt{expr-rank}).\;$\operatorname{expr\text{-}rank}(\operatorname{mxe\_}) :=$
\begin{equation*}
(\iota\, \operatorname{mxr\_}.\;((\operatorname{mxr\_} \in \mathbb{N}) \wedge ((\operatorname{mxe\_} \in \operatorname{expr\text{-}stage}(\operatorname{mxr\_})) \wedge \forall \operatorname{mxj\_} \in \mathbb{N}.\; ((\operatorname{mxe\_} \in \operatorname{expr\text{-}stage}(\operatorname{mxj\_})) \Rightarrow (\operatorname{mxr\_} \leq \operatorname{mxj\_})))))
\end{equation*}

\par\medskip\noindent\textbf{Definition} (\texttt{mapl}).\;$\operatorname{mapl}(\operatorname{mxf\_}, \operatorname{mxl\_}) :=$
\begin{equation*}
(\iota\, \operatorname{mxr\_}.\;((\operatorname{mxr\_} \in \operatorname{tuples}(\mathbf{Set})) \wedge ((\operatorname{length}(\operatorname{mxr\_}) = \operatorname{length}(\operatorname{mxl\_})) \wedge \forall \operatorname{mxi\_} \in \mathbb{N}.\; ((1 \leq \operatorname{mxi\_}) \Rightarrow ((\operatorname{mxi\_} \leq \operatorname{length}(\operatorname{mxl\_})) \Rightarrow (\operatorname{nth}(\operatorname{mxi\_}, \operatorname{mxr\_}) = \operatorname{mxf\_}(\operatorname{nth}(\operatorname{mxi\_}, \operatorname{mxl\_}))))))))
\end{equation*}

\par\medskip\noindent\textbf{Definition} (\texttt{maplk}).\;$\operatorname{maplk}(\operatorname{mxf\_}, \operatorname{mxl\_}, \operatorname{mxk\_}) :=$
\begin{equation*}
(\iota\, \operatorname{mxr\_}.\;((\operatorname{mxr\_} \in \operatorname{tuples}(\mathbf{Set})) \wedge ((\operatorname{length}(\operatorname{mxr\_}) = \operatorname{length}(\operatorname{mxl\_})) \wedge \forall \operatorname{mxi\_} \in \mathbb{N}.\; ((1 \leq \operatorname{mxi\_}) \Rightarrow ((\operatorname{mxi\_} \leq \operatorname{length}(\operatorname{mxl\_})) \Rightarrow (\operatorname{nth}(\operatorname{mxi\_}, \operatorname{mxr\_}) = \operatorname{mxf\_}(\operatorname{nth}(\operatorname{mxi\_}, \operatorname{mxl\_}), \operatorname{mxk\_})))))))
\end{equation*}

\par\medskip\noindent\textbf{Definition} (\texttt{fv-stage}).\;
\begin{equation*}
(\operatorname{fv\text{-}stage}(0) \simeq (\lambda \operatorname{mxe\_} \in \operatorname{expr\text{-}stage}(0).\; \left(\text{if }(\operatorname{nth}(1, \operatorname{mxe\_}) = 0)\text{ then }\{\operatorname{nth}(2, \operatorname{mxe\_})\}\text{ else }\emptyset\right)))
\end{equation*}
\par
Suppose $\operatorname{mxn\_} \in \mathbb{N}$. Then
\begin{equation*}
(\operatorname{fv\text{-}stage}(\operatorname{succ}(\operatorname{mxn\_})) \simeq (\lambda \operatorname{mxe\_} \in \operatorname{expr\text{-}stage}(\operatorname{succ}(\operatorname{mxn\_})).\; \left(\text{if }(\operatorname{mxe\_} \in \operatorname{expr\text{-}stage}(\operatorname{mxn\_}))\text{ then }\operatorname{fv\text{-}stage}(\operatorname{mxn\_})(\operatorname{mxe\_})\text{ else }\left(\text{if }(\operatorname{nth}(1, \operatorname{mxe\_}) = 3)\text{ then }\operatorname{big\text{-}union}(\operatorname{mxw\_}, \operatorname{make\text{-}set}(\operatorname{nth}(3, \operatorname{mxe\_})), \operatorname{fv\text{-}stage}(\operatorname{mxn\_})(\operatorname{mxw\_}))\text{ else }(\operatorname{big\text{-}union}(\operatorname{mxw\_}, \operatorname{make\text{-}set}(\operatorname{nth}(4, \operatorname{mxe\_})), \operatorname{fv\text{-}stage}(\operatorname{mxn\_})(\operatorname{mxw\_})) \cup \operatorname{fv\text{-}stage}(\operatorname{mxn\_})(\operatorname{nth}(5, \operatorname{mxe\_})))\right)\right)))
\end{equation*}

\par\medskip\noindent\textbf{Definition} (\texttt{subst-stage}).\;For all $\operatorname{mxm\_}$ and $\operatorname{mxu\_}$, 
\begin{equation*}
(\operatorname{subst\text{-}stage}(\operatorname{mxm\_}, \operatorname{mxu\_}, 0) \simeq (\lambda \operatorname{mxe\_} \in \operatorname{expr\text{-}stage}(0).\; \left(\text{if }(\operatorname{mxe\_} = {[}0, \operatorname{mxm\_}])\text{ then }\operatorname{mxu\_}\text{ else }\operatorname{mxe\_}\right)))
\end{equation*}
\par
Let $\operatorname{mxm\_}$ and $\operatorname{mxu\_}$ be arbitrary. Let $\operatorname{mxn\_} \in \mathbb{N}$. Then
\begin{equation*}
(\operatorname{subst\text{-}stage}(\operatorname{mxm\_}, \operatorname{mxu\_}, \operatorname{succ}(\operatorname{mxn\_})) \simeq (\lambda \operatorname{mxe\_} \in \operatorname{expr\text{-}stage}(\operatorname{succ}(\operatorname{mxn\_})).\; \left(\text{if }(\operatorname{mxe\_} \in \operatorname{expr\text{-}stage}(\operatorname{mxn\_}))\text{ then }\operatorname{subst\text{-}stage}(\operatorname{mxm\_}, \operatorname{mxu\_}, \operatorname{mxn\_})(\operatorname{mxe\_})\text{ else }\left(\text{if }(\operatorname{nth}(1, \operatorname{mxe\_}) = 3)\text{ then }{[}3, \operatorname{nth}(2, \operatorname{mxe\_}), \operatorname{mapl}(\operatorname{subst\text{-}stage}(\operatorname{mxm\_}, \operatorname{mxu\_}, \operatorname{mxn\_}), \operatorname{nth}(3, \operatorname{mxe\_}))]\text{ else }{[}4, \operatorname{nth}(2, \operatorname{mxe\_}), \operatorname{nth}(3, \operatorname{mxe\_}), \operatorname{mapl}(\operatorname{subst\text{-}stage}(\operatorname{mxm\_}, \operatorname{mxu\_}, \operatorname{mxn\_}), \operatorname{nth}(4, \operatorname{mxe\_})), \operatorname{subst\text{-}stage}(\operatorname{mxm\_}, \operatorname{mxu\_}, \operatorname{mxn\_})(\operatorname{nth}(5, \operatorname{mxe\_}))]\right)\right)))
\end{equation*}

\par\medskip\noindent\textbf{Definition} (\texttt{open-stage}).\;For every $\operatorname{mxu\_}$, 
\begin{equation*}
(\operatorname{open\text{-}stage}(\operatorname{mxu\_}, 0) \simeq (\lambda {[}\operatorname{mxe\_}, \operatorname{mxk\_}] \in (\operatorname{expr\text{-}stage}(0) \times \mathbb{N}).\; \left(\text{if }(\operatorname{mxe\_} = {[}1, \operatorname{mxk\_}])\text{ then }\operatorname{mxu\_}\text{ else }\operatorname{mxe\_}\right)))
\end{equation*}
\par
Let $\operatorname{mxu\_}$ be arbitrary. Let $\operatorname{mxn\_} \in \mathbb{N}$. Then
\begin{equation*}
(\operatorname{open\text{-}stage}(\operatorname{mxu\_}, \operatorname{succ}(\operatorname{mxn\_})) \simeq (\lambda {[}\operatorname{mxe\_}, \operatorname{mxk\_}] \in (\operatorname{expr\text{-}stage}(\operatorname{succ}(\operatorname{mxn\_})) \times \mathbb{N}).\; \left(\text{if }(\operatorname{mxe\_} \in \operatorname{expr\text{-}stage}(\operatorname{mxn\_}))\text{ then }\operatorname{open\text{-}stage}(\operatorname{mxu\_}, \operatorname{mxn\_})(\operatorname{mxe\_}, \operatorname{mxk\_})\text{ else }\left(\text{if }(\operatorname{nth}(1, \operatorname{mxe\_}) = 3)\text{ then }{[}3, \operatorname{nth}(2, \operatorname{mxe\_}), \operatorname{maplk}(\operatorname{open\text{-}stage}(\operatorname{mxu\_}, \operatorname{mxn\_}), \operatorname{nth}(3, \operatorname{mxe\_}), \operatorname{mxk\_})]\text{ else }{[}4, \operatorname{nth}(2, \operatorname{mxe\_}), \operatorname{nth}(3, \operatorname{mxe\_}), \operatorname{maplk}(\operatorname{open\text{-}stage}(\operatorname{mxu\_}, \operatorname{mxn\_}), \operatorname{nth}(4, \operatorname{mxe\_}), \operatorname{mxk\_}), \operatorname{open\text{-}stage}(\operatorname{mxu\_}, \operatorname{mxn\_})(\operatorname{nth}(5, \operatorname{mxe\_}), +(\operatorname{mxk\_}, \operatorname{nth}(3, \operatorname{mxe\_})))]\right)\right)))
\end{equation*}

\par\medskip\noindent\textbf{Definition} (\texttt{expr-lcd}).\;$\operatorname{expr\text{-}lcd}(\operatorname{mxe\_}) :=$ $\operatorname{lcd\text{-}stage}(\operatorname{expr\text{-}rank}(\operatorname{mxe\_}))(\operatorname{mxe\_})$
\par\medskip\noindent\textbf{Definition} (\texttt{lcd-stage}).\;
\begin{equation*}
(\operatorname{lcd\text{-}stage}(0) \simeq (\lambda \operatorname{mxe\_} \in \operatorname{expr\text{-}stage}(0).\; \left(\text{if }(\operatorname{nth}(1, \operatorname{mxe\_}) = 1)\text{ then }\{\operatorname{mxd\_} \in \mathbb{N} : (\operatorname{nth}(2, \operatorname{mxe\_}) < \operatorname{mxd\_})\}\text{ else }\mathbb{N}\right)))
\end{equation*}
\par
Suppose $\operatorname{mxn\_} \in \mathbb{N}$. Then
\begin{equation*}
(\operatorname{lcd\text{-}stage}(\operatorname{succ}(\operatorname{mxn\_})) \simeq (\lambda \operatorname{mxe\_} \in \operatorname{expr\text{-}stage}(\operatorname{succ}(\operatorname{mxn\_})).\; \left(\text{if }(\operatorname{mxe\_} \in \operatorname{expr\text{-}stage}(\operatorname{mxn\_}))\text{ then }\operatorname{lcd\text{-}stage}(\operatorname{mxn\_})(\operatorname{mxe\_})\text{ else }\left(\text{if }(\operatorname{nth}(1, \operatorname{mxe\_}) = 3)\text{ then }\{\operatorname{mxd\_} \in \mathbb{N} : \forall \operatorname{mxw\_} \in \operatorname{make\text{-}set}(\operatorname{nth}(3, \operatorname{mxe\_})).\; (\operatorname{mxd\_} \in \operatorname{lcd\text{-}stage}(\operatorname{mxn\_})(\operatorname{mxw\_}))\}\text{ else }\{\operatorname{mxd\_} \in \mathbb{N} : (\forall \operatorname{mxw\_} \in \operatorname{make\text{-}set}(\operatorname{nth}(4, \operatorname{mxe\_})).\; (\operatorname{mxd\_} \in \operatorname{lcd\text{-}stage}(\operatorname{mxn\_})(\operatorname{mxw\_})) \wedge (+(\operatorname{mxd\_}, \operatorname{nth}(3, \operatorname{mxe\_})) \in \operatorname{lcd\text{-}stage}(\operatorname{mxn\_})(\operatorname{nth}(5, \operatorname{mxe\_}))))\}\right)\right)))
\end{equation*}

\par\medskip\noindent\textbf{Definition} (\texttt{subst-fun}).\;$\operatorname{subst\text{-}fun}(\operatorname{mxm\_}, \operatorname{mxu\_}) :=$ $(\lambda \operatorname{mxw\_} \in \operatorname{expr\text{-}set}.\; \operatorname{expr\text{-}subst}(\operatorname{mxw\_}, \operatorname{mxm\_}, \operatorname{mxu\_}))$
\par\medskip\noindent\textbf{Definition} (\texttt{open-fun}).\;$\operatorname{open\text{-}fun}(\operatorname{mxu\_}) :=$
\begin{equation*}
(\lambda {[}\operatorname{mxw\_}, \operatorname{mxj\_}] \in (\operatorname{expr\text{-}set} \times \mathbb{N}).\; \operatorname{expr\text{-}open}(\operatorname{mxw\_}, \operatorname{mxj\_}, \operatorname{mxu\_}))
\end{equation*}


\vnbbookfile{theorem-library/meta-syntax-laws}{theorem-library/meta-syntax-laws}
\noindent{\small meta-syntax-laws.scm -- THE SYNTAX OF THE KERNEL AS A VNB OBJECT: the laws of the locally nameless expressions (structure-library/meta-syntax.scm). (docs/meta-k2-admissibility-2026-10-08.md, section 6, phase 1: the cost probe.) Agent META-1, 2026-10-08.  Template: theorem-library/strand-messages-laws.scm, whose local}\par
\vnbprop{efvar-unfold}{\vnbstmt{\forall n.\; (\operatorname{efvar}(n) \simeq {[}0, n])}}
\vnbprop{ebvar-unfold}{\vnbstmt{\forall k.\; (\operatorname{ebvar}(k) \simeq {[}1, k])}}
\vnbprop{econst-unfold}{\vnbstmt{\forall c.\; (\operatorname{econst}(c) \simeq {[}2, c])}}
\vnbprop{eapp-unfold}{\vnbstmt{\forall h, l.\; (\operatorname{eapp}(h, l) \simeq {[}3, h, l])}}
\vnbprop{ebind-unfold}{\vnbstmt{\forall h, c, a, b.\; (\operatorname{ebind}(h, c, a, b) \simeq {[}4, h, c, a, b])}}
\vnbprop{eapp-injective}{\vnbstmt{\forall h, l, a, b.\; ((\operatorname{eapp}(h, l) = \operatorname{eapp}(a, b)) \Rightarrow ((h = a) \wedge (l = b)))}}
\vnbprop{ebind-injective}{\vnbstmt{\begin{array}{@{}l@{}} \forall h, c, a, b, d, f, g, k. \\ \quad (\operatorname{ebind}(h, c, a, b) = \operatorname{ebind}(d, f, g, k)) \Rightarrow  \\ \quad \quad (h = d) \wedge ((c = f) \wedge ((a = g) \wedge (b = k))) \end{array}}}
\vnbprop{expr-stage-fvar}{\vnbstmt{\forall n \in \mathbb{N}.\; (\operatorname{efvar}(n) \in \operatorname{expr\text{-}stage}(0))}}
\vnbprop{expr-stage-bvar}{\vnbstmt{\forall k \in \mathbb{N}.\; (\operatorname{ebvar}(k) \in \operatorname{expr\text{-}stage}(0))}}
\vnbprop{expr-stage-const}{\vnbstmt{\forall c \in \mathbb{N}.\; (\operatorname{econst}(c) \in \operatorname{expr\text{-}stage}(0))}}
\vnbprop{expr-stage-keep}{\vnbstmt{\forall n, m.\; (((n \in \mathbb{N}) \wedge (m \in \operatorname{expr\text{-}stage}(n))) \Rightarrow (m \in \operatorname{expr\text{-}stage}((n + 1))))}}
\vnbprop{expr-stage-app}{\vnbstmt{\begin{array}{@{}l@{}} \forall n, h, l. \\ \quad ((n \in \mathbb{N}) \wedge ((h \in \mathbb{N}) \wedge (l \in \operatorname{tuples}(\operatorname{expr\text{-}stage}(n))))) \Rightarrow  \\ \quad \quad \operatorname{eapp}(h, l) \in \operatorname{expr\text{-}stage}((n + 1)) \end{array}}}
\vnbprop{expr-stage-bind}{\vnbstmt{\begin{array}{@{}l@{}} \forall n, h, c, a, b. \\ \quad ((n \in \mathbb{N}) \wedge ((h \in \mathbb{N}) \wedge ((c \in \mathbb{N}) \wedge ((a \in \operatorname{tuples}(\operatorname{expr\text{-}stage}(n))) \wedge (b \in \operatorname{expr\text{-}stage}(n)))))) \Rightarrow  \\ \quad \quad \operatorname{ebind}(h, c, a, b) \in \operatorname{expr\text{-}stage}((n + 1)) \end{array}}}
\vnbprop{expr-stage-zero-cases}{\vnbstmt{\begin{array}{@{}l@{}} \forall m \in \operatorname{expr\text{-}stage}(0). \\ \quad \exists c \in \mathbb{N}.\; (m = \operatorname{efvar}(c)) \vee \\ \quad \exists c \in \mathbb{N}.\; (m = \operatorname{ebvar}(c)) \vee \\ \quad \exists c \in \mathbb{N}.\; (m = \operatorname{econst}(c)) \end{array}}}
\vnbprop{expr-stage-succ-cases}{\vnbstmt{\begin{array}{@{}l@{}} \forall n, m. \\ \quad ((n \in \mathbb{N}) \wedge (m \in \operatorname{expr\text{-}stage}((n + 1)))) \Rightarrow  \\ \quad \quad m \in \operatorname{expr\text{-}stage}(n) \vee \\ \quad \quad \exists h \in \mathbb{N}, l \in \operatorname{tuples}(\operatorname{expr\text{-}stage}(n)).\; (m = \operatorname{eapp}(h, l)) \vee \\ \quad \quad \exists h, c \in \mathbb{N}, a \in \operatorname{tuples}(\operatorname{expr\text{-}stage}(n)), b \in \operatorname{expr\text{-}stage}(n). \\ \quad \quad \quad m = \operatorname{ebind}(h, c, a, b) \end{array}}}
\vnbprop{expr-stage-add}{\vnbstmt{\begin{array}{@{}l@{}} \forall d \in \mathbb{N}, j, m. \\ \quad ((j \in \mathbb{N}) \wedge (m \in \operatorname{expr\text{-}stage}(j))) \Rightarrow (m \in \operatorname{expr\text{-}stage}((j + d))) \end{array}}}
\vnbprop{expr-set-intro}{\vnbstmt{\forall n, m.\; (((n \in \mathbb{N}) \wedge (m \in \operatorname{expr\text{-}stage}(n))) \Rightarrow (m \in \operatorname{expr\text{-}set}))}}
\vnbprop{expr-set-elim}{\vnbstmt{\forall m \in \operatorname{expr\text{-}set}.\; \exists n \in \mathbb{N}.\; (m \in \operatorname{expr\text{-}stage}(n))}}
\vnbprop{cartesian5-set}{\vnbstmt{\begin{array}{@{}l@{}} \forall a_{1}, a_{2}, a_{3}, a_{4}, a_{5}. \\ \quad ((a_{1} \in \mathbf{Set}) \wedge ((a_{2} \in \mathbf{Set}) \wedge ((a_{3} \in \mathbf{Set}) \wedge ((a_{4} \in \mathbf{Set}) \wedge (a_{5} \in \mathbf{Set}))))) \Rightarrow  \\ \quad \quad (a_{1} \times a_{2}) \in \mathbf{Set} \end{array}}}
\vnbprop{expr-stage-set}{\vnbstmt{\forall n \in \mathbb{N}.\; (\operatorname{expr\text{-}stage}(n) \in \mathbf{Set})}}
\vnbprop{expr-set-is-set}{\vnbstmt{\operatorname{expr\text{-}set} \in \mathbf{Set}}}
\vnbprop{expr-list-in-some-stage}{\vnbstmt{\forall l \in \operatorname{tuples}(\operatorname{expr\text{-}set}).\; \exists n \in \mathbb{N}.\; (l \in \operatorname{tuples}(\operatorname{expr\text{-}stage}(n)))}}
\vnbprop{efvar-in-expr}{\vnbstmt{\forall n \in \mathbb{N}.\; (\operatorname{efvar}(n) \in \operatorname{expr\text{-}set})}}
\vnbprop{ebvar-in-expr}{\vnbstmt{\forall k \in \mathbb{N}.\; (\operatorname{ebvar}(k) \in \operatorname{expr\text{-}set})}}
\vnbprop{econst-in-expr}{\vnbstmt{\forall c \in \mathbb{N}.\; (\operatorname{econst}(c) \in \operatorname{expr\text{-}set})}}
\vnbprop{eapp-in-expr}{\vnbstmt{\forall h, l.\; (((h \in \mathbb{N}) \wedge (l \in \operatorname{tuples}(\operatorname{expr\text{-}set}))) \Rightarrow (\operatorname{eapp}(h, l) \in \operatorname{expr\text{-}set}))}}
\vnbprop{ebind-in-expr}{\vnbstmt{\begin{array}{@{}l@{}} \forall h, c, a, b. \\ \quad ((h \in \mathbb{N}) \wedge ((c \in \mathbb{N}) \wedge ((a \in \operatorname{tuples}(\operatorname{expr\text{-}set})) \wedge (b \in \operatorname{expr\text{-}set})))) \Rightarrow  \\ \quad \quad \operatorname{ebind}(h, c, a, b) \in \operatorname{expr\text{-}set} \end{array}}}
\vnbprop{expr-induction}{\vnbstmt{\begin{array}{@{}l@{}} \forall s.; \forall e \in \operatorname{expr\text{-}set}. \\ \quad \forall n \in \mathbb{N}.\; (\operatorname{efvar}(n) \in s) \wedge \\ \quad \forall k \in \mathbb{N}.\; (\operatorname{ebvar}(k) \in s) \wedge \\ \quad \forall c \in \mathbb{N}.\; (\operatorname{econst}(c) \in s) \wedge \\ \quad \forall h, l. \\ \quad \quad h \in \mathbb{N} \wedge \\ \quad \quad l \in \operatorname{tuples}(\operatorname{expr\text{-}set}) \wedge \\ \quad \quad \forall i \in \mathbb{N}.\; (((1 \leq i) \wedge (i \leq \operatorname{length}(l))) \Rightarrow (\operatorname{nth}(i, l) \in s)) \Rightarrow \\ \quad \quad \quad \operatorname{eapp}(h, l) \in s \wedge \\ \quad \forall h, b, a, \mathit{mxjb}. \\ \quad \quad h \in \mathbb{N} \wedge \\ \quad \quad b \in \mathbb{N} \wedge \\ \quad \quad a \in \operatorname{tuples}(\operatorname{expr\text{-}set}) \wedge \\ \quad \quad \mathit{mxjb} \in \operatorname{expr\text{-}set} \wedge \\ \quad \quad \forall i \in \mathbb{N}.\; (((1 \leq i) \wedge (i \leq \operatorname{length}(a))) \Rightarrow (\operatorname{nth}(i, a) \in s)) \wedge \\ \quad \quad \mathit{mxjb} \in s \Rightarrow \\ \quad \quad \quad \operatorname{ebind}(h, b, a, \mathit{mxjb}) \in s \Rightarrow \\ \quad \quad e \in s \end{array}}}
\vnbprop{expr-rank-prop}{\vnbstmt{\begin{array}{@{}l@{}} \forall e \in \operatorname{expr\text{-}set}. \\ \quad \operatorname{expr\text{-}rank}(e) \in \mathbb{N} \wedge \\ \quad e \in \operatorname{expr\text{-}stage}(\operatorname{expr\text{-}rank}(e)) \wedge \\ \quad \forall j \in \mathbb{N}.\; ((e \in \operatorname{expr\text{-}stage}(j)) \Rightarrow (\operatorname{expr\text{-}rank}(e) \leq j)) \end{array}}}
\vnbprop{expr-rank-fvar}{\vnbstmt{\forall n \in \mathbb{N}.\; (\operatorname{expr\text{-}rank}(\operatorname{efvar}(n)) = 0)}}
\vnbprop{expr-rank-bvar}{\vnbstmt{\forall n \in \mathbb{N}.\; (\operatorname{expr\text{-}rank}(\operatorname{ebvar}(n)) = 0)}}
\vnbprop{expr-rank-const}{\vnbstmt{\forall n \in \mathbb{N}.\; (\operatorname{expr\text{-}rank}(\operatorname{econst}(n)) = 0)}}
\vnbprop{expr-rank-app}{\vnbstmt{\begin{array}{@{}l@{}} \forall h, l. \\ \quad ((h \in \mathbb{N}) \wedge (l \in \operatorname{tuples}(\operatorname{expr\text{-}set}))) \Rightarrow  \\ \quad \quad \exists k \in \mathbb{N}. \\ \quad \quad \quad ((\operatorname{expr\text{-}rank}(\operatorname{eapp}(h, l)) = (k + 1)) \wedge ((l \in \operatorname{tuples}(\operatorname{expr\text{-}stage}(k))) \wedge \neg (\operatorname{eapp}(h, l) \in \operatorname{expr\text{-}stage}(k)))) \end{array}}}
\vnbprop{expr-rank-bind}{\vnbstmt{\begin{array}{@{}l@{}} \forall h, c, a, b. \\ \quad ((h \in \mathbb{N}) \wedge ((c \in \mathbb{N}) \wedge ((a \in \operatorname{tuples}(\operatorname{expr\text{-}set})) \wedge (b \in \operatorname{expr\text{-}set})))) \Rightarrow  \\ \quad \quad \exists k \in \mathbb{N}. \\ \quad \quad \quad \operatorname{expr\text{-}rank}(\operatorname{ebind}(h, c, a, b)) = (k + 1) \wedge \\ \quad \quad \quad a \in \operatorname{tuples}(\operatorname{expr\text{-}stage}(k)) \wedge \\ \quad \quad \quad b \in \operatorname{expr\text{-}stage}(k) \wedge \\ \quad \quad \quad \neg (\operatorname{ebind}(h, c, a, b) \in \operatorname{expr\text{-}stage}(k)) \end{array}}}
\vnbprop{makeset-member-in}{\vnbstmt{\forall s, l, w.\; (((l \in \operatorname{tuples}(s)) \wedge (w \in \operatorname{make\text{-}set}(l))) \Rightarrow (w \in s))}}
\vnbprop{fv-stage-step}{\vnbstmt{\begin{array}{@{}l@{}} \forall k, e. \\ \quad ((k \in \mathbb{N}) \wedge (e \in \operatorname{expr\text{-}stage}(k))) \Rightarrow  \\ \quad \quad (\operatorname{fv\text{-}stage}((k + 1))(e) \simeq \operatorname{fv\text{-}stage}(k)(e)) \end{array}}}
\vnbprop{fv-stage-coh}{\vnbstmt{\begin{array}{@{}l@{}} \forall d \in \mathbb{N}, j, e. \\ \quad ((j \in \mathbb{N}) \wedge (e \in \operatorname{expr\text{-}stage}(j))) \Rightarrow (\operatorname{fv\text{-}stage}((j + d))(e) \simeq \operatorname{fv\text{-}stage}(j)(e)) \end{array}}}
\vnbprop{expr-fv-at}{\vnbstmt{\begin{array}{@{}l@{}} \forall j, e. \\ \quad ((j \in \mathbb{N}) \wedge (e \in \operatorname{expr\text{-}stage}(j))) \Rightarrow (\operatorname{expr\text{-}fv}(e) \simeq \operatorname{fv\text{-}stage}(j)(e)) \end{array}}}
\vnbprop{fv-efvar}{\vnbstmt{\forall n \in \mathbb{N}.\; (\operatorname{expr\text{-}fv}(\operatorname{efvar}(n)) \simeq \{n\})}}
\vnbprop{fv-ebvar}{\vnbstmt{\forall k \in \mathbb{N}.\; (\operatorname{expr\text{-}fv}(\operatorname{ebvar}(k)) \simeq \emptyset)}}
\vnbprop{fv-econst}{\vnbstmt{\forall c \in \mathbb{N}.\; (\operatorname{expr\text{-}fv}(\operatorname{econst}(c)) \simeq \emptyset)}}
\vnbprop{fv-eapp}{\vnbstmt{\begin{array}{@{}l@{}} \forall h, l. \\ \quad ((h \in \mathbb{N}) \wedge (l \in \operatorname{tuples}(\operatorname{expr\text{-}set}))) \Rightarrow  \\ \quad \quad (\operatorname{expr\text{-}fv}(\operatorname{eapp}(h, l)) \simeq \operatorname{big\text{-}union}(w, \operatorname{make\text{-}set}(l), \operatorname{expr\text{-}fv}(w))) \end{array}}}
\vnbprop{fv-ebind}{\vnbstmt{\begin{array}{@{}l@{}} \forall h, c, a, b. \\ \quad ((h \in \mathbb{N}) \wedge ((c \in \mathbb{N}) \wedge ((a \in \operatorname{tuples}(\operatorname{expr\text{-}set})) \wedge (b \in \operatorname{expr\text{-}set})))) \Rightarrow  \\ \quad \quad (\operatorname{expr\text{-}fv}(\operatorname{ebind}(h, c, a, b)) \simeq (\operatorname{big\text{-}union}(w, \operatorname{make\text{-}set}(a), \operatorname{expr\text{-}fv}(w)) \cup \operatorname{expr\text{-}fv}(b))) \end{array}}}
\vnbprop{bu-finite-nn}{\vnbstmt{\begin{array}{@{}l@{}} \forall s \in \mathbf{Set}, m. \\ \quad ((|s| \in \mathbb{N}) \wedge \forall v_{5} \in s.\; ((m(v_{5}) \in \mathbf{Set}) \wedge ((m(v_{5}) \subseteq \mathbb{N}) \wedge (|m(v_{5})| \in \mathbb{N})))) \Rightarrow  \\ \quad \quad \operatorname{big\text{-}union}(w, s, m(w)) \in \mathbf{Set} \wedge \\ \quad \quad \operatorname{big\text{-}union}(w, s, m(w)) \subseteq \mathbb{N} \wedge \\ \quad \quad |\operatorname{big\text{-}union}(w, s, m(w))| \in \mathbb{N} \end{array}}}
\vnbprop{fv-finite}{\vnbstmt{\begin{array}{@{}l@{}} \forall e \in \operatorname{expr\text{-}set}. \\ \quad (\operatorname{expr\text{-}fv}(e) \in \mathbf{Set}) \wedge ((\operatorname{expr\text{-}fv}(e) \subseteq \mathbb{N}) \wedge (|\operatorname{expr\text{-}fv}(e)| \in \mathbb{N})) \end{array}}}
\vnbprop{mapl-prop}{\vnbstmt{\begin{array}{@{}l@{}} \forall s, g, l. \\ \quad ((\operatorname{length}(l) \in \mathbb{N}) \wedge \forall i \in \mathbb{N}.\; (((1 \leq i) \wedge (i \leq \operatorname{length}(l))) \Rightarrow (g(\operatorname{nth}(i, l)) \in s))) \Rightarrow  \\ \quad \quad \operatorname{mapl}(g, l) \in \operatorname{tuples}(s) \wedge \\ \quad \quad \operatorname{length}(\operatorname{mapl}(g, l)) = \operatorname{length}(l) \wedge \\ \quad \quad \forall i \in \mathbb{N}. \\ \quad \quad \quad ((1 \leq i) \wedge (i \leq \operatorname{length}(l))) \Rightarrow (\operatorname{nth}(i, \operatorname{mapl}(g, l)) = g(\operatorname{nth}(i, l))) \end{array}}}
\vnbprop{maplk-prop}{\vnbstmt{\begin{array}{@{}l@{}} \forall s, g, l, x. \\ \quad ((\operatorname{length}(l) \in \mathbb{N}) \wedge \forall i \in \mathbb{N}.\; (((1 \leq i) \wedge (i \leq \operatorname{length}(l))) \Rightarrow (g(\operatorname{nth}(i, l), x) \in s))) \Rightarrow  \\ \quad \quad \operatorname{maplk}(g, l, x) \in \operatorname{tuples}(s) \wedge \\ \quad \quad \operatorname{length}(\operatorname{maplk}(g, l, x)) = \operatorname{length}(l) \wedge \\ \quad \quad \forall i \in \mathbb{N}. \\ \quad \quad \quad ((1 \leq i) \wedge (i \leq \operatorname{length}(l))) \Rightarrow (\operatorname{nth}(i, \operatorname{maplk}(g, l, x)) = g(\operatorname{nth}(i, l), x)) \end{array}}}
\vnbprop{subst-stage-fun}{\vnbstmt{\begin{array}{@{}l@{}} \forall n \in \mathbb{N}, m.; \forall u \in \operatorname{expr\text{-}set}. \\ \quad \operatorname{subst\text{-}stage}(m, u, n) \in (\operatorname{expr\text{-}stage}(n) \to \operatorname{expr\text{-}set}) \end{array}}}
\vnbprop{subst-stage-step}{\vnbstmt{\begin{array}{@{}l@{}} \forall m, u, k, e. \\ \quad ((k \in \mathbb{N}) \wedge (e \in \operatorname{expr\text{-}stage}(k))) \Rightarrow  \\ \quad \quad (\operatorname{subst\text{-}stage}(m, u, (k + 1))(e) \simeq \operatorname{subst\text{-}stage}(m, u, k)(e)) \end{array}}}
\vnbprop{subst-stage-coh}{\vnbstmt{\begin{array}{@{}l@{}} \forall d \in \mathbb{N}, m, u, j, e. \\ \quad ((j \in \mathbb{N}) \wedge (e \in \operatorname{expr\text{-}stage}(j))) \Rightarrow  \\ \quad \quad (\operatorname{subst\text{-}stage}(m, u, (j + d))(e) \simeq \operatorname{subst\text{-}stage}(m, u, j)(e)) \end{array}}}
\vnbprop{expr-subst-at}{\vnbstmt{\begin{array}{@{}l@{}} \forall m, u, j, e. \\ \quad ((j \in \mathbb{N}) \wedge (e \in \operatorname{expr\text{-}stage}(j))) \Rightarrow  \\ \quad \quad (\operatorname{expr\text{-}subst}(e, m, u) \simeq \operatorname{subst\text{-}stage}(m, u, j)(e)) \end{array}}}
\vnbprop{subst-in-expr}{\vnbstmt{\forall m.; \forall u, e \in \operatorname{expr\text{-}set}.\; (\operatorname{expr\text{-}subst}(e, m, u) \in \operatorname{expr\text{-}set})}}
\vnbprop{subst-efvar-hit}{\vnbstmt{\forall m \in \mathbb{N}, u \in \operatorname{expr\text{-}set}.\; (\operatorname{expr\text{-}subst}(\operatorname{efvar}(m), m, u) \simeq u)}}
\vnbprop{subst-efvar-miss}{\vnbstmt{\begin{array}{@{}l@{}} \forall m, u, n. \\ \quad ((u \in \operatorname{expr\text{-}set}) \wedge ((n \in \mathbb{N}) \wedge \neg (n = m))) \Rightarrow  \\ \quad \quad (\operatorname{expr\text{-}subst}(\operatorname{efvar}(n), m, u) \simeq \operatorname{efvar}(n)) \end{array}}}
\vnbprop{subst-ebvar}{\vnbstmt{\begin{array}{@{}l@{}} \forall m, u, k. \\ \quad ((u \in \operatorname{expr\text{-}set}) \wedge (k \in \mathbb{N})) \Rightarrow (\operatorname{expr\text{-}subst}(\operatorname{ebvar}(k), m, u) \simeq \operatorname{ebvar}(k)) \end{array}}}
\vnbprop{subst-econst}{\vnbstmt{\begin{array}{@{}l@{}} \forall m, u, c. \\ \quad ((u \in \operatorname{expr\text{-}set}) \wedge (c \in \mathbb{N})) \Rightarrow (\operatorname{expr\text{-}subst}(\operatorname{econst}(c), m, u) \simeq \operatorname{econst}(c)) \end{array}}}
\vnbprop{subst-eapp}{\vnbstmt{\begin{array}{@{}l@{}} \forall m, u, h, l. \\ \quad ((u \in \operatorname{expr\text{-}set}) \wedge ((h \in \mathbb{N}) \wedge (l \in \operatorname{tuples}(\operatorname{expr\text{-}set})))) \Rightarrow  \\ \quad \quad (\operatorname{expr\text{-}subst}(\operatorname{eapp}(h, l), m, u) \simeq \operatorname{eapp}(h, \operatorname{mapl}(\operatorname{subst\text{-}fun}(m, u), l))) \end{array}}}
\vnbprop{subst-ebind}{\vnbstmt{\begin{array}{@{}l@{}} \forall m, u, h, c, a, b. \\ \quad ((u \in \operatorname{expr\text{-}set}) \wedge ((h \in \mathbb{N}) \wedge ((c \in \mathbb{N}) \wedge ((a \in \operatorname{tuples}(\operatorname{expr\text{-}set})) \wedge (b \in \operatorname{expr\text{-}set}))))) \Rightarrow  \\ \quad \quad (\operatorname{expr\text{-}subst}(\operatorname{ebind}(h, c, a, b), m, u) \simeq \operatorname{ebind}(h, c, \operatorname{mapl}(\operatorname{subst\text{-}fun}(m, u), a), \operatorname{expr\text{-}subst}(b, m, u))) \end{array}}}
\vnbprop{subst-fresh}{\vnbstmt{\begin{array}{@{}l@{}} \forall m.; \forall u, e \in \operatorname{expr\text{-}set}. \\ \quad \neg (m \in \operatorname{expr\text{-}fv}(e)) \Rightarrow (\operatorname{expr\text{-}subst}(e, m, u) = e) \end{array}}}
\vnbprop{subst-fun-unfold}{\vnbstmt{\forall m, u.\; (\operatorname{subst\text{-}fun}(m, u) \simeq (\lambda w \in \operatorname{expr\text{-}set}.\; \operatorname{expr\text{-}subst}(w, m, u)))}}
\vnbprop{fv-subst}{\vnbstmt{\begin{array}{@{}l@{}} \forall m.; \forall u, e \in \operatorname{expr\text{-}set}, x \in \operatorname{expr\text{-}fv}(\operatorname{expr\text{-}subst}(e, m, u)). \\ \quad ((x \in \operatorname{expr\text{-}fv}(e)) \wedge \neg (x = m)) \vee (x \in \operatorname{expr\text{-}fv}(u)) \end{array}}}
\vnbprop{open-stage-fun}{\vnbstmt{\begin{array}{@{}l@{}} \forall n \in \mathbb{N}, u \in \operatorname{expr\text{-}set}. \\ \quad \operatorname{open\text{-}stage}(u, n) \in ((\operatorname{expr\text{-}stage}(n) \times \mathbb{N}) \to \operatorname{expr\text{-}set}) \end{array}}}
\vnbprop{open-stage-step}{\vnbstmt{\begin{array}{@{}l@{}} \forall u, k, e, x. \\ \quad ((k \in \mathbb{N}) \wedge ((e \in \operatorname{expr\text{-}stage}(k)) \wedge (x \in \mathbb{N}))) \Rightarrow  \\ \quad \quad (\operatorname{open\text{-}stage}(u, (k + 1))(e, x) \simeq \operatorname{open\text{-}stage}(u, k)(e, x)) \end{array}}}
\vnbprop{open-stage-coh}{\vnbstmt{\begin{array}{@{}l@{}} \forall d \in \mathbb{N}, u, j, e, x. \\ \quad ((j \in \mathbb{N}) \wedge ((e \in \operatorname{expr\text{-}stage}(j)) \wedge (x \in \mathbb{N}))) \Rightarrow  \\ \quad \quad (\operatorname{open\text{-}stage}(u, (j + d))(e, x) \simeq \operatorname{open\text{-}stage}(u, j)(e, x)) \end{array}}}
\vnbprop{expr-open-at}{\vnbstmt{\begin{array}{@{}l@{}} \forall u, j, e, x. \\ \quad ((j \in \mathbb{N}) \wedge ((e \in \operatorname{expr\text{-}stage}(j)) \wedge (x \in \mathbb{N}))) \Rightarrow  \\ \quad \quad (\operatorname{expr\text{-}open}(e, x, u) \simeq \operatorname{open\text{-}stage}(u, j)(e, x)) \end{array}}}
\vnbprop{open-in-expr}{\vnbstmt{\forall u \in \operatorname{expr\text{-}set}, x \in \mathbb{N}, e \in \operatorname{expr\text{-}set}.\; (\operatorname{expr\text{-}open}(e, x, u) \in \operatorname{expr\text{-}set})}}
\vnbprop{open-ebvar-hit}{\vnbstmt{\forall u, x.\; (((u \in \operatorname{expr\text{-}set}) \wedge (x \in \mathbb{N})) \Rightarrow (\operatorname{expr\text{-}open}(\operatorname{ebvar}(x), x, u) \simeq u))}}
\vnbprop{open-ebvar-miss}{\vnbstmt{\begin{array}{@{}l@{}} \forall u, x, j. \\ \quad ((u \in \operatorname{expr\text{-}set}) \wedge ((x \in \mathbb{N}) \wedge ((j \in \mathbb{N}) \wedge \neg (j = x)))) \Rightarrow  \\ \quad \quad (\operatorname{expr\text{-}open}(\operatorname{ebvar}(j), x, u) \simeq \operatorname{ebvar}(j)) \end{array}}}
\vnbprop{open-efvar}{\vnbstmt{\begin{array}{@{}l@{}} \forall u, x, n. \\ \quad ((u \in \operatorname{expr\text{-}set}) \wedge ((x \in \mathbb{N}) \wedge (n \in \mathbb{N}))) \Rightarrow  \\ \quad \quad (\operatorname{expr\text{-}open}(\operatorname{efvar}(n), x, u) \simeq \operatorname{efvar}(n)) \end{array}}}
\vnbprop{open-econst}{\vnbstmt{\begin{array}{@{}l@{}} \forall u, x, c. \\ \quad ((u \in \operatorname{expr\text{-}set}) \wedge ((x \in \mathbb{N}) \wedge (c \in \mathbb{N}))) \Rightarrow  \\ \quad \quad (\operatorname{expr\text{-}open}(\operatorname{econst}(c), x, u) \simeq \operatorname{econst}(c)) \end{array}}}
\vnbprop{open-eapp}{\vnbstmt{\begin{array}{@{}l@{}} \forall u, x, h, l. \\ \quad ((u \in \operatorname{expr\text{-}set}) \wedge ((x \in \mathbb{N}) \wedge ((h \in \mathbb{N}) \wedge (l \in \operatorname{tuples}(\operatorname{expr\text{-}set}))))) \Rightarrow  \\ \quad \quad (\operatorname{expr\text{-}open}(\operatorname{eapp}(h, l), x, u) \simeq \operatorname{eapp}(h, \operatorname{maplk}(\operatorname{open\text{-}fun}(u), l, x))) \end{array}}}
\vnbprop{open-ebind}{\vnbstmt{\begin{array}{@{}l@{}} \forall u, x, h, c, a, b. \\ \quad ((u \in \operatorname{expr\text{-}set}) \wedge ((x \in \mathbb{N}) \wedge ((h \in \mathbb{N}) \wedge ((c \in \mathbb{N}) \wedge ((a \in \operatorname{tuples}(\operatorname{expr\text{-}set})) \wedge (b \in \operatorname{expr\text{-}set})))))) \Rightarrow  \\ \quad \quad (\operatorname{expr\text{-}open}(\operatorname{ebind}(h, c, a, b), x, u) \simeq \operatorname{ebind}(h, c, \operatorname{maplk}(\operatorname{open\text{-}fun}(u), a, x), \operatorname{expr\text{-}open}(b, (x + c), u))) \end{array}}}
\vnbprop{lcd-stage-step}{\vnbstmt{\begin{array}{@{}l@{}} \forall k, e. \\ \quad ((k \in \mathbb{N}) \wedge (e \in \operatorname{expr\text{-}stage}(k))) \Rightarrow  \\ \quad \quad (\operatorname{lcd\text{-}stage}((k + 1))(e) \simeq \operatorname{lcd\text{-}stage}(k)(e)) \end{array}}}
\vnbprop{lcd-stage-coh}{\vnbstmt{\begin{array}{@{}l@{}} \forall d \in \mathbb{N}, j, e. \\ \quad ((j \in \mathbb{N}) \wedge (e \in \operatorname{expr\text{-}stage}(j))) \Rightarrow  \\ \quad \quad (\operatorname{lcd\text{-}stage}((j + d))(e) \simeq \operatorname{lcd\text{-}stage}(j)(e)) \end{array}}}
\vnbprop{expr-lcd-at}{\vnbstmt{\begin{array}{@{}l@{}} \forall j, e. \\ \quad ((j \in \mathbb{N}) \wedge (e \in \operatorname{expr\text{-}stage}(j))) \Rightarrow (\operatorname{expr\text{-}lcd}(e) \simeq \operatorname{lcd\text{-}stage}(j)(e)) \end{array}}}
\vnbprop{lcd-efvar}{\vnbstmt{\forall n \in \mathbb{N}.\; (\operatorname{expr\text{-}lcd}(\operatorname{efvar}(n)) \simeq \mathbb{N})}}
\vnbprop{lcd-econst}{\vnbstmt{\forall c \in \mathbb{N}.\; (\operatorname{expr\text{-}lcd}(\operatorname{econst}(c)) \simeq \mathbb{N})}}
\vnbprop{lcd-ebvar}{\vnbstmt{\forall k \in \mathbb{N}.\; (\operatorname{expr\text{-}lcd}(\operatorname{ebvar}(k)) \simeq \{d \in \mathbb{N} : (k < d)\})}}
\vnbprop{lcd-eapp}{\vnbstmt{\begin{array}{@{}l@{}} \forall h, l. \\ \quad ((h \in \mathbb{N}) \wedge (l \in \operatorname{tuples}(\operatorname{expr\text{-}set}))) \Rightarrow  \\ \quad \quad (\operatorname{expr\text{-}lcd}(\operatorname{eapp}(h, l)) \simeq \{d \in \mathbb{N} : \forall w \in \operatorname{make\text{-}set}(l).\; (d \in \operatorname{expr\text{-}lcd}(w))\}) \end{array}}}
\vnbprop{lcd-ebind}{\vnbstmt{\begin{array}{@{}l@{}} \forall h, c, a, b. \\ \quad ((h \in \mathbb{N}) \wedge ((c \in \mathbb{N}) \wedge ((a \in \operatorname{tuples}(\operatorname{expr\text{-}set})) \wedge (b \in \operatorname{expr\text{-}set})))) \Rightarrow  \\ \quad \quad (\operatorname{expr\text{-}lcd}(\operatorname{ebind}(h, c, a, b)) \simeq \{d \in \mathbb{N} : (\forall w \in \operatorname{make\text{-}set}(a).\; (d \in \operatorname{expr\text{-}lcd}(w)) \wedge ((d + c) \in \operatorname{expr\text{-}lcd}(b)))\}) \end{array}}}
\vnbprop{lcd-open}{\vnbstmt{\begin{array}{@{}l@{}} \forall u, e \in \operatorname{expr\text{-}set}, x, d. \\ \quad ((x \in \mathbb{N}) \wedge ((d \in \mathbb{N}) \wedge ((d \leq x) \wedge (d \in \operatorname{expr\text{-}lcd}(e))))) \Rightarrow  \\ \quad \quad \operatorname{expr\text{-}open}(e, x, u) = e \end{array}}}
\vnbprop{lc-open}{\vnbstmt{\begin{array}{@{}l@{}} \forall u, e \in \operatorname{expr\text{-}set}, x \in \mathbb{N}. \\ \quad \operatorname{expr\text{-}lc}(e) \Rightarrow (\operatorname{expr\text{-}open}(e, x, u) = e) \end{array}}}
\vnbprop{subst-open-commute}{\vnbstmt{\begin{array}{@{}l@{}} \forall n, u, v.; \forall e \in \operatorname{expr\text{-}set}, x \in \mathbb{N}. \\ \quad ((u \in \operatorname{expr\text{-}set}) \wedge ((v \in \operatorname{expr\text{-}set}) \wedge \operatorname{expr\text{-}lc}(v))) \Rightarrow  \\ \quad \quad \operatorname{expr\text{-}subst}(\operatorname{expr\text{-}open}(e, x, u), n, v) = \operatorname{expr\text{-}open}(\operatorname{expr\text{-}subst}(e, n, v), x, \operatorname{expr\text{-}subst}(u, n, v)) \end{array}}}

\vnbbookfile{structure-library/meta-core}{structure-library/meta-core}
\noindent{\small meta-core.scm -- THE CORE OF K2: well-formedness, sequents, the core rule relations, DERIVABLE, and the eight kernel relations transcribed (three in phase 2, five by META-KIT). (docs/meta-k2-admissibility-2026-10-08.md, sections 3, 4 and 5, phase 2.) DEFINITIONS ONLY; the laws are in theorem-library/meta-core-laws.scm.}\par
\par\medskip\noindent\textbf{Definition} (\texttt{derivable}).\;$\operatorname{derivable}(\operatorname{mct\_}) :=$
\begin{equation*}
\{\operatorname{mcq\_} \in \operatorname{seq\text{-}set} : \forall \operatorname{mcs\_} \in \operatorname{power}(\operatorname{seq\text{-}set}).\; (\operatorname{core\text{-}closed}(\operatorname{mct\_}, \operatorname{mcs\_}) \Rightarrow (\operatorname{mcq\_} \in \operatorname{mcs\_}))\}
\end{equation*}

\par\medskip\noindent\textbf{Definition} (\texttt{enc-num}).\;$\operatorname{enc\text{-}num}(\operatorname{mck\_}) :=$ $\operatorname{econst}((2 \cdot \operatorname{mck\_}))$
\par\medskip\noindent\textbf{Definition} (\texttt{enc-setc}).\; $(\operatorname{enc\text{-}setc} \simeq \operatorname{econst}(1))$
\par\medskip\noindent\textbf{Definition} (\texttt{pred-code}).\;For every $\operatorname{mcpk\_}$, 
\begin{equation*}
\operatorname{pred\text{-}code}(\operatorname{mcpk\_}) \Leftrightarrow ((\operatorname{mcpk\_} = 7) \vee ((\operatorname{mcpk\_} = 9) \vee ((\operatorname{mcpk\_} = 11) \vee (\operatorname{mcpk\_} = 13))))
\end{equation*}

\par\medskip\noindent\textbf{Definition} (\texttt{hd-res}).\;$\operatorname{hd\text{-}res}(\operatorname{mch\_}) :=$ $\left(\text{if }(\operatorname{mch\_} \leq 10)\text{ then }2\text{ else }1\right)$
\par\medskip\noindent\textbf{Definition} (\texttt{app-ar-ok}).\;For all $\operatorname{mch\_}$ and $\operatorname{mcn\_}$, 
\begin{equation*}
\begin{array}{@{}l@{}} \operatorname{app\text{-}ar\text{-}ok}(\operatorname{mch\_}, \operatorname{mcn\_}) \\ \Leftrightarrow\; ((10 \leq \operatorname{mch\_}) \vee (((\operatorname{mch\_} \leq 1) \wedge (\operatorname{mcn\_} = 0)) \vee (((\operatorname{mch\_} = 2) \wedge (\operatorname{mcn\_} = 1)) \vee ((3 \leq \operatorname{mch\_}) \wedge (\operatorname{mcn\_} = 2))))) \end{array}
\end{equation*}

\par\medskip\noindent\textbf{Definition} (\texttt{hd-argk}).\;$\operatorname{hd\text{-}argk}(\operatorname{mch\_}) :=$ $\left(\text{if }(\operatorname{mch\_} \leq 6)\text{ then }2\text{ else }1\right)$
\par\medskip\noindent\textbf{Definition} (\texttt{hd-bres}).\;$\operatorname{hd\text{-}bres}(\operatorname{mch\_}) :=$ $\left(\text{if }(\operatorname{mch\_} \leq 21)\text{ then }2\text{ else }1\right)$
\par\medskip\noindent\textbf{Definition} (\texttt{bind-sig}).\;For all $\operatorname{mch\_}$, $\operatorname{mcnb\_}$ and $\operatorname{mcn\_}$, 
\begin{equation*}
\begin{array}{@{}l@{}} \operatorname{bind\text{-}sig}(\operatorname{mch\_}, \operatorname{mcnb\_}, \operatorname{mcn\_}) \\ \Leftrightarrow\; (20 \leq \operatorname{mch\_}) \wedge ((\operatorname{mch\_} \leq 21) \wedge ((\operatorname{mcnb\_} = 1) \wedge (\operatorname{mcn\_} = 0))) \vee \\ ((\operatorname{mch\_} = 22) \vee ((\operatorname{mch\_} = 24) \vee (\operatorname{mch\_} = 26))) \wedge ((\operatorname{mcnb\_} = 1) \wedge (\operatorname{mcn\_} = 1)) \vee \\ ((\operatorname{mch\_} = 23) \vee (\operatorname{mch\_} = 25)) \wedge ((\operatorname{mcnb\_} = 1) \wedge (\operatorname{mcn\_} = 0)) \vee \\ (\operatorname{mch\_} = 27) \wedge ((\operatorname{mcnb\_} = 0) \wedge (\operatorname{mcn\_} = 2)) \end{array}
\end{equation*}

\par\medskip\noindent\textbf{Definition} (\texttt{hd-body}).\;$\operatorname{hd\text{-}body}(\operatorname{mch\_}) :=$ $\left(\text{if }((\operatorname{mch\_} = 24) \vee (\operatorname{mch\_} = 26))\text{ then }1\text{ else }2\right)$
\par\medskip\noindent\textbf{Definition} (\texttt{wf-term}).\;For every $\operatorname{mce\_}$,  $\operatorname{wf\text{-}term}(\operatorname{mce\_}) \Leftrightarrow (1 \in \operatorname{expr\text{-}wks}(\operatorname{mce\_}))$
\par\medskip\noindent\textbf{Definition} (\texttt{wf-fml}).\;For every $\operatorname{mce\_}$,  $\operatorname{wf\text{-}fml}(\operatorname{mce\_}) \Leftrightarrow (2 \in \operatorname{expr\text{-}wks}(\operatorname{mce\_}))$
\par\medskip\noindent\textbf{Definition} (\texttt{clf}).\; $(\mathit{clf} \simeq \{\operatorname{mce\_} \in \operatorname{expr\text{-}set} : (\operatorname{wf\text{-}fml}(\operatorname{mce\_}) \wedge \operatorname{expr\text{-}lc}(\operatorname{mce\_}))\})$
\par\medskip\noindent\textbf{Definition} (\texttt{clt}).\; $(\mathit{clt} \simeq \{\operatorname{mce\_} \in \operatorname{expr\text{-}set} : (\operatorname{wf\text{-}term}(\operatorname{mce\_}) \wedge \operatorname{expr\text{-}lc}(\operatorname{mce\_}))\})$
\par\medskip\noindent\textbf{Definition} (\texttt{ctx-set}).\; $(\operatorname{ctx\text{-}set} \simeq \{\operatorname{mcg\_} \in \operatorname{power}(\mathit{clf}) : (|\operatorname{mcg\_}| \in \mathbb{N})\})$
\par\medskip\noindent\textbf{Definition} (\texttt{seq-set}).\; $(\operatorname{seq\text{-}set} \simeq (\operatorname{ctx\text{-}set} \times \mathit{clf}))$
\par\medskip\noindent\textbf{Definition} (\texttt{ctx-add}).\;$\operatorname{ctx\text{-}add}(\operatorname{mcg\_}, \operatorname{mca\_}) :=$ $(\operatorname{mcg\_} \cup \{\operatorname{mca\_}\})$
\par\medskip\noindent\textbf{Definition} (\texttt{ctx-fv}).\;$\operatorname{ctx\text{-}fv}(\operatorname{mcg\_}) :=$ $\operatorname{big\text{-}union}(\operatorname{mcw\_}, \operatorname{mcg\_}, \operatorname{expr\text{-}fv}(\operatorname{mcw\_}))$
\par\medskip\noindent\textbf{Definition} (\texttt{core-closed}).\;For all $\operatorname{mct\_}$ and $\operatorname{mcs\_}$, 
\begin{equation*}
\begin{array}{@{}l@{}} \operatorname{core\text{-}closed}(\operatorname{mct\_}, \operatorname{mcs\_}) \\ \Leftrightarrow\; \forall \operatorname{mchs\_}, \operatorname{mcc\_}. \\ \quad \operatorname{core\text{-}rule}(\operatorname{mct\_}, \operatorname{mchs\_}, \operatorname{mcc\_}) \Rightarrow  \\ \quad \operatorname{hyps\text{-}in}(\operatorname{mchs\_}, \operatorname{mcs\_}) \Rightarrow ((\operatorname{mcc\_} \in \operatorname{seq\text{-}set}) \Rightarrow (\operatorname{mcc\_} \in \operatorname{mcs\_})) \end{array}
\end{equation*}

\par\medskip\noindent\textbf{Definition} (\texttt{core-rule}).\;For all $\operatorname{mct\_}$, $\operatorname{mchs\_}$ and $\operatorname{mcc\_}$, 
\begin{equation*}
\begin{array}{@{}l@{}} \operatorname{core\text{-}rule}(\operatorname{mct\_}, \operatorname{mchs\_}, \operatorname{mcc\_}) \\ \Leftrightarrow\; \operatorname{cr\text{-}and\text{-}intro}(\operatorname{mct\_}, \operatorname{mchs\_}, \operatorname{mcc\_}) \vee \\ \operatorname{cr\text{-}and\text{-}elim}(\operatorname{mct\_}, \operatorname{mchs\_}, \operatorname{mcc\_}) \vee \\ \operatorname{cr\text{-}or\text{-}intro\text{-}left}(\operatorname{mct\_}, \operatorname{mchs\_}, \operatorname{mcc\_}) \vee \\ \operatorname{cr\text{-}or\text{-}intro\text{-}right}(\operatorname{mct\_}, \operatorname{mchs\_}, \operatorname{mcc\_}) \vee \\ \operatorname{cr\text{-}or\text{-}elim}(\operatorname{mct\_}, \operatorname{mchs\_}, \operatorname{mcc\_}) \vee \\ \operatorname{cr\text{-}implies\text{-}intro}(\operatorname{mct\_}, \operatorname{mchs\_}, \operatorname{mcc\_}) \vee \\ \operatorname{cr\text{-}backchain}(\operatorname{mct\_}, \operatorname{mchs\_}, \operatorname{mcc\_}) \vee \\ \operatorname{cr\text{-}not\text{-}intro}(\operatorname{mct\_}, \operatorname{mchs\_}, \operatorname{mcc\_}) \vee \\ \operatorname{cr\text{-}not\text{-}elim}(\operatorname{mct\_}, \operatorname{mchs\_}, \operatorname{mcc\_}) \vee \\ \operatorname{cr\text{-}iff\text{-}intro}(\operatorname{mct\_}, \operatorname{mchs\_}, \operatorname{mcc\_}) \vee \\ \operatorname{cr\text{-}iff\text{-}elim}(\operatorname{mct\_}, \operatorname{mchs\_}, \operatorname{mcc\_}) \vee \\ \operatorname{cr\text{-}forall\text{-}intro}(\operatorname{mct\_}, \operatorname{mchs\_}, \operatorname{mcc\_}) \vee \\ \operatorname{cr\text{-}forall\text{-}elim}(\operatorname{mct\_}, \operatorname{mchs\_}, \operatorname{mcc\_}) \vee \\ \operatorname{cr\text{-}forsome\text{-}intro}(\operatorname{mct\_}, \operatorname{mchs\_}, \operatorname{mcc\_}) \vee \\ \operatorname{cr\text{-}forsome\text{-}elim}(\operatorname{mct\_}, \operatorname{mchs\_}, \operatorname{mcc\_}) \vee \\ \operatorname{cr\text{-}assumption}(\operatorname{mct\_}, \operatorname{mchs\_}, \operatorname{mcc\_}) \vee \\ \operatorname{cr\text{-}theorem\text{-}assumption}(\operatorname{mct\_}, \operatorname{mchs\_}, \operatorname{mcc\_}) \vee \\ \operatorname{cr\text{-}cut}(\operatorname{mct\_}, \operatorname{mchs\_}, \operatorname{mcc\_}) \vee \\ \operatorname{cr\text{-}weakening}(\operatorname{mct\_}, \operatorname{mchs\_}, \operatorname{mcc\_}) \vee \\ \operatorname{cr\text{-}contradiction}(\operatorname{mct\_}, \operatorname{mchs\_}, \operatorname{mcc\_}) \vee \\ \operatorname{cr\text{-}truth\text{-}intro}(\operatorname{mct\_}, \operatorname{mchs\_}, \operatorname{mcc\_}) \vee \\ \operatorname{cr\text{-}eq\text{-}subst}(\operatorname{mct\_}, \operatorname{mchs\_}, \operatorname{mcc\_}) \vee \\ \operatorname{cr\text{-}eq\text{-}subst\text{-}sym}(\operatorname{mct\_}, \operatorname{mchs\_}, \operatorname{mcc\_}) \vee \\ \operatorname{cr\text{-}reflexivity}(\operatorname{mct\_}, \operatorname{mchs\_}, \operatorname{mcc\_}) \vee \\ \operatorname{cr\text{-}quasi\text{-}reflexivity}(\operatorname{mct\_}, \operatorname{mchs\_}, \operatorname{mcc\_}) \vee \\ \operatorname{cr\text{-}union\text{-}mem}(\operatorname{mct\_}, \operatorname{mchs\_}, \operatorname{mcc\_}) \vee \\ \operatorname{cr\text{-}intersection\text{-}mem}(\operatorname{mct\_}, \operatorname{mchs\_}, \operatorname{mcc\_}) \vee \\ \operatorname{cr\text{-}cartesian\text{-}mem}(\operatorname{mct\_}, \operatorname{mchs\_}, \operatorname{mcc\_}) \vee \\ \operatorname{cr\text{-}tuples\text{-}mem}(\operatorname{mct\_}, \operatorname{mchs\_}, \operatorname{mcc\_}) \vee \\ \operatorname{cr\text{-}nth\text{-}eq}(\operatorname{mct\_}, \operatorname{mchs\_}, \operatorname{mcc\_}) \vee \\ \operatorname{cr\text{-}length\text{-}eq}(\operatorname{mct\_}, \operatorname{mchs\_}, \operatorname{mcc\_}) \vee \\ \operatorname{cr\text{-}tuple\text{-}eq}(\operatorname{mct\_}, \operatorname{mchs\_}, \operatorname{mcc\_}) \vee \\ \operatorname{cr\text{-}if\text{-}true}(\operatorname{mct\_}, \operatorname{mchs\_}, \operatorname{mcc\_}) \vee \\ \operatorname{cr\text{-}if\text{-}false}(\operatorname{mct\_}, \operatorname{mchs\_}, \operatorname{mcc\_}) \vee \\ \operatorname{cr\text{-}sep\text{-}sethood}(\operatorname{mct\_}, \operatorname{mchs\_}, \operatorname{mcc\_}) \vee \\ \operatorname{cr\text{-}sep\text{-}mem}(\operatorname{mct\_}, \operatorname{mchs\_}, \operatorname{mcc\_}) \vee \\ \operatorname{cr\text{-}comp\text{-}mem}(\operatorname{mct\_}, \operatorname{mchs\_}, \operatorname{mcc\_}) \vee \\ \operatorname{cr\text{-}big\text{-}union\text{-}sethood}(\operatorname{mct\_}, \operatorname{mchs\_}, \operatorname{mcc\_}) \vee \\ \operatorname{cr\text{-}big\text{-}union\text{-}mem}(\operatorname{mct\_}, \operatorname{mchs\_}, \operatorname{mcc\_}) \vee \\ \operatorname{cr\text{-}iota\text{-}def}(\operatorname{mct\_}, \operatorname{mchs\_}, \operatorname{mcc\_}) \vee \\ \operatorname{cr\text{-}iota\text{-}in\text{-}elim}(\operatorname{mct\_}, \operatorname{mchs\_}, \operatorname{mcc\_}) \vee \\ \operatorname{cr\text{-}lambda\text{-}type}(\operatorname{mct\_}, \operatorname{mchs\_}, \operatorname{mcc\_}) \vee \\ \operatorname{cr\text{-}lambda\text{-}beta\text{-}eq}(\operatorname{mct\_}, \operatorname{mchs\_}, \operatorname{mcc\_}) \vee \\ \operatorname{cr\text{-}def\text{-}strict\text{-}in\text{-}l}(\operatorname{mct\_}, \operatorname{mchs\_}, \operatorname{mcc\_}) \vee \\ \operatorname{cr\text{-}def\text{-}strict\text{-}in\text{-}r}(\operatorname{mct\_}, \operatorname{mchs\_}, \operatorname{mcc\_}) \vee \\ \operatorname{cr\text{-}def\text{-}total}(\operatorname{mct\_}, \operatorname{mchs\_}, \operatorname{mcc\_}) \vee \\ \operatorname{cr\text{-}def\text{-}binder}(\operatorname{mct\_}, \operatorname{mchs\_}, \operatorname{mcc\_}) \vee \\ \operatorname{cr\text{-}def\text{-}comp}(\operatorname{mct\_}, \operatorname{mchs\_}, \operatorname{mcc\_}) \vee \\ \operatorname{cr\text{-}def\text{-}list}(\operatorname{mct\_}, \operatorname{mchs\_}, \operatorname{mcc\_}) \vee \\ \operatorname{cr\text{-}def\text{-}strict\text{-}arg}(\operatorname{mct\_}, \operatorname{mchs\_}, \operatorname{mcc\_}) \vee \\ \operatorname{cr\text{-}def\text{-}strict\text{-}dom}(\operatorname{mct\_}, \operatorname{mchs\_}, \operatorname{mcc\_}) \vee \\ \operatorname{cr\text{-}def\text{-}strict\text{-}pred}(\operatorname{mct\_}, \operatorname{mchs\_}, \operatorname{mcc\_}) \vee \\ \operatorname{cr\text{-}iota\text{-}cr}(\operatorname{mct\_}, \operatorname{mchs\_}, \operatorname{mcc\_}) \vee \\ \operatorname{cr\text{-}tbind\text{-}cr}(\operatorname{mct\_}, \operatorname{mchs\_}, \operatorname{mcc\_}) \vee \\ \operatorname{cr\text{-}functoid\text{-}beta\text{-}eq}(\operatorname{mct\_}, \operatorname{mchs\_}, \operatorname{mcc\_}) \vee \\ \operatorname{cr\text{-}tuple\text{-}eta}(\operatorname{mct\_}, \operatorname{mchs\_}, \operatorname{mcc\_}) \end{array}
\end{equation*}

\par\medskip\noindent\textbf{Definition} (\texttt{wks-stage}).\; $(\operatorname{wks\text{-}stage}(0) \simeq (\lambda \operatorname{mce\_} \in \operatorname{expr\text{-}stage}(0).\; \{1\}))$\par
Suppose $\operatorname{mcn\_} \in \mathbb{N}$. Then
\begin{equation*}
(\operatorname{wks\text{-}stage}(\operatorname{succ}(\operatorname{mcn\_})) \simeq (\lambda \operatorname{mce\_} \in \operatorname{expr\text{-}stage}(\operatorname{succ}(\operatorname{mcn\_})).\; \left(\text{if }(\operatorname{mce\_} \in \operatorname{expr\text{-}stage}(\operatorname{mcn\_}))\text{ then }\operatorname{wks\text{-}stage}(\operatorname{mcn\_})(\operatorname{mce\_})\text{ else }\left(\text{if }(\operatorname{nth}(1, \operatorname{mce\_}) = 3)\text{ then }\{\operatorname{mckk\_} \in \mathbb{N} : ((\operatorname{mckk\_} = \operatorname{hd\text{-}res}(\operatorname{nth}(2, \operatorname{mce\_}))) \wedge (\operatorname{app\text{-}ar\text{-}ok}(\operatorname{nth}(2, \operatorname{mce\_}), \operatorname{length}(\operatorname{nth}(3, \operatorname{mce\_}))) \wedge \forall \operatorname{mcw\_} \in \operatorname{make\text{-}set}(\operatorname{nth}(3, \operatorname{mce\_})).\; (\operatorname{hd\text{-}argk}(\operatorname{nth}(2, \operatorname{mce\_})) \in \operatorname{wks\text{-}stage}(\operatorname{mcn\_})(\operatorname{mcw\_}))))\}\text{ else }\{\operatorname{mckk\_} \in \mathbb{N} : ((\operatorname{mckk\_} = \operatorname{hd\text{-}bres}(\operatorname{nth}(2, \operatorname{mce\_}))) \wedge (\operatorname{bind\text{-}sig}(\operatorname{nth}(2, \operatorname{mce\_}), \operatorname{nth}(3, \operatorname{mce\_}), \operatorname{length}(\operatorname{nth}(4, \operatorname{mce\_}))) \wedge (\forall \operatorname{mcw\_} \in \operatorname{make\text{-}set}(\operatorname{nth}(4, \operatorname{mce\_})).\; (1 \in \operatorname{wks\text{-}stage}(\operatorname{mcn\_})(\operatorname{mcw\_})) \wedge (\operatorname{hd\text{-}body}(\operatorname{nth}(2, \operatorname{mce\_})) \in \operatorname{wks\text{-}stage}(\operatorname{mcn\_})(\operatorname{nth}(5, \operatorname{mce\_}))))))\}\right)\right)))
\end{equation*}

\par\medskip\noindent\textbf{Definition} (\texttt{expr-wks}).\;$\operatorname{expr\text{-}wks}(\operatorname{mce\_}) :=$ $\operatorname{wks\text{-}stage}(\operatorname{expr\text{-}rank}(\operatorname{mce\_}))(\operatorname{mce\_})$
\par\medskip\noindent\textbf{Definition} (\texttt{enc-truth}).\; $(\operatorname{enc\text{-}truth} \simeq \operatorname{eapp}(0, {[}]))$
\par\medskip\noindent\textbf{Definition} (\texttt{enc-falsity}).\; $(\operatorname{enc\text{-}falsity} \simeq \operatorname{eapp}(1, {[}]))$
\par\medskip\noindent\textbf{Definition} (\texttt{enc-not}).\;$\operatorname{enc\text{-}not}(\operatorname{mcp\_}) :=$ $\operatorname{eapp}(2, {[}\operatorname{mcp\_}])$
\par\medskip\noindent\textbf{Definition} (\texttt{enc-and}).\;$\operatorname{enc\text{-}and}(\operatorname{mcp\_}, \operatorname{mcq\_}) :=$ $\operatorname{eapp}(3, {[}\operatorname{mcp\_}, \operatorname{mcq\_}])$
\par\medskip\noindent\textbf{Definition} (\texttt{enc-or}).\;$\operatorname{enc\text{-}or}(\operatorname{mcp\_}, \operatorname{mcq\_}) :=$ $\operatorname{eapp}(4, {[}\operatorname{mcp\_}, \operatorname{mcq\_}])$
\par\medskip\noindent\textbf{Definition} (\texttt{enc-imp}).\;$\operatorname{enc\text{-}imp}(\operatorname{mcp\_}, \operatorname{mcq\_}) :=$ $\operatorname{eapp}(5, {[}\operatorname{mcp\_}, \operatorname{mcq\_}])$
\par\medskip\noindent\textbf{Definition} (\texttt{enc-iff}).\;$\operatorname{enc\text{-}iff}(\operatorname{mcp\_}, \operatorname{mcq\_}) :=$ $\operatorname{eapp}(6, {[}\operatorname{mcp\_}, \operatorname{mcq\_}])$
\par\medskip\noindent\textbf{Definition} (\texttt{enc-in}).\;$\operatorname{enc\text{-}in}(\operatorname{mcp\_}, \operatorname{mcq\_}) :=$ $\operatorname{eapp}(7, {[}\operatorname{mcp\_}, \operatorname{mcq\_}])$
\par\medskip\noindent\textbf{Definition} (\texttt{enc-eq}).\;$\operatorname{enc\text{-}eq}(\operatorname{mcp\_}, \operatorname{mcq\_}) :=$ $\operatorname{eapp}(8, {[}\operatorname{mcp\_}, \operatorname{mcq\_}])$
\par\medskip\noindent\textbf{Definition} (\texttt{enc-qeq}).\;$\operatorname{enc\text{-}qeq}(\operatorname{mcp\_}, \operatorname{mcq\_}) :=$ $\operatorname{eapp}(9, {[}\operatorname{mcp\_}, \operatorname{mcq\_}])$
\par\medskip\noindent\textbf{Definition} (\texttt{enc-union}).\;$\operatorname{enc\text{-}union}(\operatorname{mcl\_}) :=$ $\operatorname{eapp}(11, \operatorname{mcl\_})$
\par\medskip\noindent\textbf{Definition} (\texttt{enc-inter}).\;$\operatorname{enc\text{-}inter}(\operatorname{mcl\_}) :=$ $\operatorname{eapp}(12, \operatorname{mcl\_})$
\par\medskip\noindent\textbf{Definition} (\texttt{enc-cart}).\;$\operatorname{enc\text{-}cart}(\operatorname{mcl\_}) :=$ $\operatorname{eapp}(13, \operatorname{mcl\_})$
\par\medskip\noindent\textbf{Definition} (\texttt{enc-tuples}).\;$\operatorname{enc\text{-}tuples}(\operatorname{mcp\_}) :=$ $\operatorname{eapp}(14, {[}\operatorname{mcp\_}])$
\par\medskip\noindent\textbf{Definition} (\texttt{enc-nth}).\;$\operatorname{enc\text{-}nth}(\operatorname{mcp\_}, \operatorname{mcq\_}) :=$ $\operatorname{eapp}(15, {[}\operatorname{mcp\_}, \operatorname{mcq\_}])$
\par\medskip\noindent\textbf{Definition} (\texttt{enc-length}).\;$\operatorname{enc\text{-}length}(\operatorname{mcp\_}) :=$ $\operatorname{eapp}(16, {[}\operatorname{mcp\_}])$
\par\medskip\noindent\textbf{Definition} (\texttt{enc-list}).\;$\operatorname{enc\text{-}list}(\operatorname{mcl\_}) :=$ $\operatorname{eapp}(17, \operatorname{mcl\_})$
\par\medskip\noindent\textbf{Definition} (\texttt{enc-apply}).\;$\operatorname{enc\text{-}apply}(\operatorname{mcp\_}, \operatorname{mcq\_}) :=$ $\operatorname{eapp}(18, {[}\operatorname{mcp\_}, \operatorname{mcq\_}])$
\par\medskip\noindent\textbf{Definition} (\texttt{enc-fun}).\;$\operatorname{enc\text{-}fun}(\operatorname{mcp\_}, \operatorname{mcq\_}) :=$ $\operatorname{eapp}(19, {[}\operatorname{mcp\_}, \operatorname{mcq\_}])$
\par\medskip\noindent\textbf{Definition} (\texttt{enc-all}).\;$\operatorname{enc\text{-}all}(\operatorname{mcb\_}) :=$ $\operatorname{ebind}(20, 1, {[}], \operatorname{mcb\_})$
\par\medskip\noindent\textbf{Definition} (\texttt{enc-ex}).\;$\operatorname{enc\text{-}ex}(\operatorname{mcb\_}) :=$ $\operatorname{ebind}(21, 1, {[}], \operatorname{mcb\_})$
\par\medskip\noindent\textbf{Definition} (\texttt{enc-sep}).\;$\operatorname{enc\text{-}sep}(\operatorname{mca\_}, \operatorname{mcb\_}) :=$ $\operatorname{ebind}(22, 1, {[}\operatorname{mca\_}], \operatorname{mcb\_})$
\par\medskip\noindent\textbf{Definition} (\texttt{enc-comp}).\;$\operatorname{enc\text{-}comp}(\operatorname{mcb\_}) :=$ $\operatorname{ebind}(23, 1, {[}], \operatorname{mcb\_})$
\par\medskip\noindent\textbf{Definition} (\texttt{enc-bu}).\;$\operatorname{enc\text{-}bu}(\operatorname{mca\_}, \operatorname{mcb\_}) :=$ $\operatorname{ebind}(24, 1, {[}\operatorname{mca\_}], \operatorname{mcb\_})$
\par\medskip\noindent\textbf{Definition} (\texttt{enc-iota}).\;$\operatorname{enc\text{-}iota}(\operatorname{mcb\_}) :=$ $\operatorname{ebind}(25, 1, {[}], \operatorname{mcb\_})$
\par\medskip\noindent\textbf{Definition} (\texttt{enc-lam}).\;$\operatorname{enc\text{-}lam}(\operatorname{mca\_}, \operatorname{mcb\_}) :=$ $\operatorname{ebind}(26, 1, {[}\operatorname{mca\_}], \operatorname{mcb\_})$
\par\medskip\noindent\textbf{Definition} (\texttt{enc-if}).\;$\operatorname{enc\text{-}if}(\operatorname{mcp\_}, \operatorname{mca\_}, \operatorname{mcb\_}) :=$ $\operatorname{ebind}(27, 0, {[}\operatorname{mca\_}, \operatorname{mcb\_}], \operatorname{mcp\_})$
\par\medskip\noindent\textbf{Definition} (\texttt{enc-orl}).\;For all $\operatorname{mcx\_}$ and $\operatorname{mcl\_}$,  $(\operatorname{enc\text{-}orl}(\operatorname{mcx\_}, \operatorname{mcl\_}, 0) \simeq \operatorname{enc\text{-}falsity})$\par
Let $\operatorname{mcx\_}$ and $\operatorname{mcl\_}$ be arbitrary. Let $\operatorname{mcn\_} \in \mathbb{N}$. Then
\begin{equation*}
(\operatorname{enc\text{-}orl}(\operatorname{mcx\_}, \operatorname{mcl\_}, \operatorname{succ}(\operatorname{mcn\_})) \simeq \operatorname{enc\text{-}or}(\operatorname{enc\text{-}orl}(\operatorname{mcx\_}, \operatorname{mcl\_}, \operatorname{mcn\_}), \operatorname{enc\text{-}in}(\operatorname{mcx\_}, \operatorname{nth}(\operatorname{succ}(\operatorname{mcn\_}), \operatorname{mcl\_}))))
\end{equation*}

\par\medskip\noindent\textbf{Definition} (\texttt{enc-andl}).\;For all $\operatorname{mcx\_}$ and $\operatorname{mcl\_}$,  $(\operatorname{enc\text{-}andl}(\operatorname{mcx\_}, \operatorname{mcl\_}, 0) \simeq \operatorname{enc\text{-}truth})$\par
Let $\operatorname{mcx\_}$ and $\operatorname{mcl\_}$ be arbitrary. Let $\operatorname{mcn\_} \in \mathbb{N}$. Then
\begin{equation*}
(\operatorname{enc\text{-}andl}(\operatorname{mcx\_}, \operatorname{mcl\_}, \operatorname{succ}(\operatorname{mcn\_})) \simeq \operatorname{enc\text{-}and}(\operatorname{enc\text{-}andl}(\operatorname{mcx\_}, \operatorname{mcl\_}, \operatorname{mcn\_}), \operatorname{enc\text{-}in}(\operatorname{mcx\_}, \operatorname{nth}(\operatorname{succ}(\operatorname{mcn\_}), \operatorname{mcl\_}))))
\end{equation*}

\par\medskip\noindent\textbf{Definition} (\texttt{enc-ande}).\;For all $\operatorname{mcl\_}$ and $\operatorname{mca\_}$,  $(\operatorname{enc\text{-}ande}(\operatorname{mcl\_}, \operatorname{mca\_}, 0) \simeq \operatorname{enc\text{-}truth})$\par
Let $\operatorname{mcl\_}$ and $\operatorname{mca\_}$ be arbitrary. Let $\operatorname{mcn\_} \in \mathbb{N}$. Then
\begin{equation*}
(\operatorname{enc\text{-}ande}(\operatorname{mcl\_}, \operatorname{mca\_}, \operatorname{succ}(\operatorname{mcn\_})) \simeq \operatorname{enc\text{-}and}(\operatorname{enc\text{-}ande}(\operatorname{mcl\_}, \operatorname{mca\_}, \operatorname{mcn\_}), \operatorname{enc\text{-}in}(\operatorname{nth}(\operatorname{succ}(\operatorname{mcn\_}), \operatorname{mcl\_}), \operatorname{mca\_})))
\end{equation*}

\par\medskip\noindent\textbf{Definition} (\texttt{enc-cartl}).\;For all $\operatorname{mcx\_}$ and $\operatorname{mcl\_}$,  $(\operatorname{enc\text{-}cartl}(\operatorname{mcx\_}, \operatorname{mcl\_}, 0) \simeq \operatorname{enc\text{-}truth})$\par
Let $\operatorname{mcx\_}$ and $\operatorname{mcl\_}$ be arbitrary. Let $\operatorname{mcn\_} \in \mathbb{N}$. Then
\begin{equation*}
(\operatorname{enc\text{-}cartl}(\operatorname{mcx\_}, \operatorname{mcl\_}, \operatorname{succ}(\operatorname{mcn\_})) \simeq \operatorname{enc\text{-}and}(\operatorname{enc\text{-}cartl}(\operatorname{mcx\_}, \operatorname{mcl\_}, \operatorname{mcn\_}), \operatorname{enc\text{-}in}(\operatorname{enc\text{-}nth}(\operatorname{enc\text{-}num}(\operatorname{succ}(\operatorname{mcn\_})), \operatorname{mcx\_}), \operatorname{nth}(\operatorname{succ}(\operatorname{mcn\_}), \operatorname{mcl\_}))))
\end{equation*}

\par\medskip\noindent\textbf{Definition} (\texttt{enc-eql}).\;For all $\operatorname{mcl\_}$ and $\operatorname{mca\_}$,  $(\operatorname{enc\text{-}eql}(\operatorname{mcl\_}, \operatorname{mca\_}, 0) \simeq \operatorname{enc\text{-}truth})$\par
Let $\operatorname{mcl\_}$ and $\operatorname{mca\_}$ be arbitrary. Let $\operatorname{mcn\_} \in \mathbb{N}$. Then
\begin{equation*}
(\operatorname{enc\text{-}eql}(\operatorname{mcl\_}, \operatorname{mca\_}, \operatorname{succ}(\operatorname{mcn\_})) \simeq \operatorname{enc\text{-}and}(\operatorname{enc\text{-}eql}(\operatorname{mcl\_}, \operatorname{mca\_}, \operatorname{mcn\_}), \operatorname{enc\text{-}eq}(\operatorname{nth}(\operatorname{succ}(\operatorname{mcn\_}), \operatorname{mcl\_}), \operatorname{nth}(\operatorname{succ}(\operatorname{mcn\_}), \operatorname{mca\_}))))
\end{equation*}

\par\medskip\noindent\textbf{Definition} (\texttt{cr-and-intro}).\;For all $\operatorname{mct\_}$, $\operatorname{mchs\_}$ and $\operatorname{mcc\_}$, 
\begin{equation*}
\begin{array}{@{}l@{}} \operatorname{cr\text{-}and\text{-}intro}(\operatorname{mct\_}, \operatorname{mchs\_}, \operatorname{mcc\_}) \\ \Leftrightarrow\; \exists \operatorname{mcg\_}, \operatorname{mca\_}, \operatorname{mcb\_}. \\ \quad (\operatorname{mcc\_} = {[}\operatorname{mcg\_}, \operatorname{enc\text{-}and}(\operatorname{mca\_}, \operatorname{mcb\_})]) \wedge (\operatorname{mchs\_} = {[}{[}\operatorname{mcg\_}, \operatorname{mca\_}], {[}\operatorname{mcg\_}, \operatorname{mcb\_}]]) \end{array}
\end{equation*}

\par\medskip\noindent\textbf{Definition} (\texttt{cr-and-elim}).\;For all $\operatorname{mct\_}$, $\operatorname{mchs\_}$ and $\operatorname{mcc\_}$, 
\begin{equation*}
\begin{array}{@{}l@{}} \operatorname{cr\text{-}and\text{-}elim}(\operatorname{mct\_}, \operatorname{mchs\_}, \operatorname{mcc\_}) \\ \Leftrightarrow\; \exists \operatorname{mcg\_}, \operatorname{mcf\_}, \operatorname{mca\_}, \operatorname{mcb\_}. \\ \quad \operatorname{mcc\_} = {[}\operatorname{mcg\_}, \operatorname{mcf\_}] \wedge \\ \quad \operatorname{enc\text{-}and}(\operatorname{mca\_}, \operatorname{mcb\_}) \in \operatorname{mcg\_} \wedge \\ \quad \operatorname{mchs\_} = {[}{[}\operatorname{ctx\text{-}add}(\operatorname{ctx\text{-}add}(\operatorname{mcg\_}, \operatorname{mca\_}), \operatorname{mcb\_}), \operatorname{mcf\_}]] \end{array}
\end{equation*}

\par\medskip\noindent\textbf{Definition} (\texttt{cr-or-intro-left}).\;For all $\operatorname{mct\_}$, $\operatorname{mchs\_}$ and $\operatorname{mcc\_}$, 
\begin{equation*}
\begin{array}{@{}l@{}} \operatorname{cr\text{-}or\text{-}intro\text{-}left}(\operatorname{mct\_}, \operatorname{mchs\_}, \operatorname{mcc\_}) \\ \Leftrightarrow\; \exists \operatorname{mcg\_}, \operatorname{mca\_}, \operatorname{mcb\_}. \\ \quad (\operatorname{mcc\_} = {[}\operatorname{mcg\_}, \operatorname{enc\text{-}or}(\operatorname{mca\_}, \operatorname{mcb\_})]) \wedge (\operatorname{mchs\_} = {[}{[}\operatorname{mcg\_}, \operatorname{mca\_}]]) \end{array}
\end{equation*}

\par\medskip\noindent\textbf{Definition} (\texttt{cr-or-intro-right}).\;For all $\operatorname{mct\_}$, $\operatorname{mchs\_}$ and $\operatorname{mcc\_}$, 
\begin{equation*}
\begin{array}{@{}l@{}} \operatorname{cr\text{-}or\text{-}intro\text{-}right}(\operatorname{mct\_}, \operatorname{mchs\_}, \operatorname{mcc\_}) \\ \Leftrightarrow\; \exists \operatorname{mcg\_}, \operatorname{mca\_}, \operatorname{mcb\_}. \\ \quad (\operatorname{mcc\_} = {[}\operatorname{mcg\_}, \operatorname{enc\text{-}or}(\operatorname{mca\_}, \operatorname{mcb\_})]) \wedge (\operatorname{mchs\_} = {[}{[}\operatorname{mcg\_}, \operatorname{mcb\_}]]) \end{array}
\end{equation*}

\par\medskip\noindent\textbf{Definition} (\texttt{cr-or-elim}).\;For all $\operatorname{mct\_}$, $\operatorname{mchs\_}$ and $\operatorname{mcc\_}$, 
\begin{equation*}
\begin{array}{@{}l@{}} \operatorname{cr\text{-}or\text{-}elim}(\operatorname{mct\_}, \operatorname{mchs\_}, \operatorname{mcc\_}) \\ \Leftrightarrow\; \exists \operatorname{mcg\_}, \operatorname{mcf\_}, \operatorname{mca\_}, \operatorname{mcb\_}. \\ \quad \operatorname{mcc\_} = {[}\operatorname{mcg\_}, \operatorname{mcf\_}] \wedge \\ \quad \operatorname{enc\text{-}or}(\operatorname{mca\_}, \operatorname{mcb\_}) \in \operatorname{mcg\_} \wedge \\ \quad \operatorname{mchs\_} = {[}{[}\operatorname{ctx\text{-}add}(\operatorname{mcg\_}, \operatorname{mca\_}), \operatorname{mcf\_}], {[}\operatorname{ctx\text{-}add}(\operatorname{mcg\_}, \operatorname{mcb\_}), \operatorname{mcf\_}]] \end{array}
\end{equation*}

\par\medskip\noindent\textbf{Definition} (\texttt{cr-implies-intro}).\;For all $\operatorname{mct\_}$, $\operatorname{mchs\_}$ and $\operatorname{mcc\_}$, 
\begin{equation*}
\begin{array}{@{}l@{}} \operatorname{cr\text{-}implies\text{-}intro}(\operatorname{mct\_}, \operatorname{mchs\_}, \operatorname{mcc\_}) \\ \Leftrightarrow\; \exists \operatorname{mcg\_}, \operatorname{mca\_}, \operatorname{mcb\_}. \\ \quad (\operatorname{mcc\_} = {[}\operatorname{mcg\_}, \operatorname{enc\text{-}imp}(\operatorname{mca\_}, \operatorname{mcb\_})]) \wedge (\operatorname{mchs\_} = {[}{[}\operatorname{ctx\text{-}add}(\operatorname{mcg\_}, \operatorname{mca\_}), \operatorname{mcb\_}]]) \end{array}
\end{equation*}

\par\medskip\noindent\textbf{Definition} (\texttt{cr-backchain}).\;For all $\operatorname{mct\_}$, $\operatorname{mchs\_}$ and $\operatorname{mcc\_}$, 
\begin{equation*}
\begin{array}{@{}l@{}} \operatorname{cr\text{-}backchain}(\operatorname{mct\_}, \operatorname{mchs\_}, \operatorname{mcc\_}) \\ \Leftrightarrow\; \exists \operatorname{mcg\_}, \operatorname{mcf\_}, \operatorname{mca\_}. \\ \quad (\operatorname{mcc\_} = {[}\operatorname{mcg\_}, \operatorname{mcf\_}]) \wedge ((\operatorname{enc\text{-}imp}(\operatorname{mca\_}, \operatorname{mcf\_}) \in \operatorname{mcg\_}) \wedge (\operatorname{mchs\_} = {[}{[}\operatorname{mcg\_}, \operatorname{mca\_}]])) \end{array}
\end{equation*}

\par\medskip\noindent\textbf{Definition} (\texttt{cr-not-intro}).\;For all $\operatorname{mct\_}$, $\operatorname{mchs\_}$ and $\operatorname{mcc\_}$, 
\begin{equation*}
\begin{array}{@{}l@{}} \operatorname{cr\text{-}not\text{-}intro}(\operatorname{mct\_}, \operatorname{mchs\_}, \operatorname{mcc\_}) \\ \Leftrightarrow\; \exists \operatorname{mcg\_}, \operatorname{mca\_}. \\ \quad (\operatorname{mcc\_} = {[}\operatorname{mcg\_}, \operatorname{enc\text{-}not}(\operatorname{mca\_})]) \wedge (\operatorname{mchs\_} = {[}{[}\operatorname{ctx\text{-}add}(\operatorname{mcg\_}, \operatorname{mca\_}), \operatorname{enc\text{-}falsity}]]) \end{array}
\end{equation*}

\par\medskip\noindent\textbf{Definition} (\texttt{cr-not-elim}).\;For all $\operatorname{mct\_}$, $\operatorname{mchs\_}$ and $\operatorname{mcc\_}$, 
\begin{equation*}
\begin{array}{@{}l@{}} \operatorname{cr\text{-}not\text{-}elim}(\operatorname{mct\_}, \operatorname{mchs\_}, \operatorname{mcc\_}) \\ \Leftrightarrow\; \exists \operatorname{mcg\_}, \operatorname{mcf\_}, \operatorname{mca\_}. \\ \quad (\operatorname{mcc\_} = {[}\operatorname{mcg\_}, \operatorname{mcf\_}]) \wedge ((\operatorname{enc\text{-}not}(\operatorname{mca\_}) \in \operatorname{mcg\_}) \wedge ((\operatorname{mca\_} \in \operatorname{mcg\_}) \wedge (\operatorname{mchs\_} = {[}]))) \end{array}
\end{equation*}

\par\medskip\noindent\textbf{Definition} (\texttt{cr-iff-intro}).\;For all $\operatorname{mct\_}$, $\operatorname{mchs\_}$ and $\operatorname{mcc\_}$, 
\begin{equation*}
\begin{array}{@{}l@{}} \operatorname{cr\text{-}iff\text{-}intro}(\operatorname{mct\_}, \operatorname{mchs\_}, \operatorname{mcc\_}) \\ \Leftrightarrow\; \exists \operatorname{mcg\_}, \operatorname{mca\_}, \operatorname{mcb\_}. \\ \quad \operatorname{mcc\_} = {[}\operatorname{mcg\_}, \operatorname{enc\text{-}iff}(\operatorname{mca\_}, \operatorname{mcb\_})] \wedge \\ \quad \operatorname{mchs\_} = {[}{[}\operatorname{ctx\text{-}add}(\operatorname{mcg\_}, \operatorname{mca\_}), \operatorname{mcb\_}], {[}\operatorname{ctx\text{-}add}(\operatorname{mcg\_}, \operatorname{mcb\_}), \operatorname{mca\_}]] \end{array}
\end{equation*}

\par\medskip\noindent\textbf{Definition} (\texttt{cr-iff-elim}).\;For all $\operatorname{mct\_}$, $\operatorname{mchs\_}$ and $\operatorname{mcc\_}$, 
\begin{equation*}
\begin{array}{@{}l@{}} \operatorname{cr\text{-}iff\text{-}elim}(\operatorname{mct\_}, \operatorname{mchs\_}, \operatorname{mcc\_}) \\ \Leftrightarrow\; \exists \operatorname{mcg\_}, \operatorname{mcf\_}, \operatorname{mca\_}, \operatorname{mcb\_}. \\ \quad \operatorname{mcc\_} = {[}\operatorname{mcg\_}, \operatorname{mcf\_}] \wedge \\ \quad \operatorname{enc\text{-}iff}(\operatorname{mca\_}, \operatorname{mcb\_}) \in \operatorname{mcg\_} \wedge \\ \quad \operatorname{mchs\_} = {[}{[}\operatorname{ctx\text{-}add}(\operatorname{ctx\text{-}add}(\operatorname{mcg\_}, \operatorname{enc\text{-}imp}(\operatorname{mca\_}, \operatorname{mcb\_})), \operatorname{enc\text{-}imp}(\operatorname{mcb\_}, \operatorname{mca\_})), \operatorname{mcf\_}]] \end{array}
\end{equation*}

\par\medskip\noindent\textbf{Definition} (\texttt{cr-forall-intro}).\;For all $\operatorname{mct\_}$, $\operatorname{mchs\_}$ and $\operatorname{mcc\_}$, 
\begin{equation*}
\begin{array}{@{}l@{}} \operatorname{cr\text{-}forall\text{-}intro}(\operatorname{mct\_}, \operatorname{mchs\_}, \operatorname{mcc\_}) \\ \Leftrightarrow\; \exists \operatorname{mcg\_}, \operatorname{mcb\_}, \operatorname{mcy\_}. \\ \quad \operatorname{mcc\_} = {[}\operatorname{mcg\_}, \operatorname{enc\text{-}all}(\operatorname{mcb\_})] \wedge \\ \quad \operatorname{mcy\_} \in \mathbb{N} \wedge \\ \quad \neg (\operatorname{mcy\_} \in (\operatorname{ctx\text{-}fv}(\operatorname{mcg\_}) \cup \operatorname{expr\text{-}fv}(\operatorname{mcb\_}))) \wedge \\ \quad \operatorname{mchs\_} = {[}{[}\operatorname{mcg\_}, \operatorname{expr\text{-}open}(\operatorname{mcb\_}, 0, \operatorname{efvar}(\operatorname{mcy\_}))]] \end{array}
\end{equation*}

\par\medskip\noindent\textbf{Definition} (\texttt{cr-forall-elim}).\;For all $\operatorname{mct\_}$, $\operatorname{mchs\_}$ and $\operatorname{mcc\_}$, 
\begin{equation*}
\begin{array}{@{}l@{}} \operatorname{cr\text{-}forall\text{-}elim}(\operatorname{mct\_}, \operatorname{mchs\_}, \operatorname{mcc\_}) \\ \Leftrightarrow\; \exists \operatorname{mcg\_}, \operatorname{mcf\_}, \operatorname{mcb\_}, \operatorname{mctm\_}. \\ \quad \operatorname{mcc\_} = {[}\operatorname{mcg\_}, \operatorname{mcf\_}] \wedge \\ \quad \operatorname{enc\text{-}all}(\operatorname{mcb\_}) \in \operatorname{mcg\_} \wedge \\ \quad \operatorname{mctm\_} \in \mathit{clt} \wedge \\ \quad \operatorname{mchs\_} = {[}{[}\operatorname{ctx\text{-}add}(\operatorname{mcg\_}, \operatorname{expr\text{-}open}(\operatorname{mcb\_}, 0, \operatorname{mctm\_})), \operatorname{mcf\_}], {[}\operatorname{mcg\_}, \operatorname{enc\text{-}eq}(\operatorname{mctm\_}, \operatorname{mctm\_})]] \end{array}
\end{equation*}

\par\medskip\noindent\textbf{Definition} (\texttt{cr-forsome-intro}).\;For all $\operatorname{mct\_}$, $\operatorname{mchs\_}$ and $\operatorname{mcc\_}$, 
\begin{equation*}
\begin{array}{@{}l@{}} \operatorname{cr\text{-}forsome\text{-}intro}(\operatorname{mct\_}, \operatorname{mchs\_}, \operatorname{mcc\_}) \\ \Leftrightarrow\; \exists \operatorname{mcg\_}, \operatorname{mcb\_}, \operatorname{mctm\_}. \\ \quad \operatorname{mcc\_} = {[}\operatorname{mcg\_}, \operatorname{enc\text{-}ex}(\operatorname{mcb\_})] \wedge \\ \quad \operatorname{mctm\_} \in \mathit{clt} \wedge \\ \quad \operatorname{mchs\_} = {[}{[}\operatorname{mcg\_}, \operatorname{expr\text{-}open}(\operatorname{mcb\_}, 0, \operatorname{mctm\_})], {[}\operatorname{mcg\_}, \operatorname{enc\text{-}eq}(\operatorname{mctm\_}, \operatorname{mctm\_})]] \end{array}
\end{equation*}

\par\medskip\noindent\textbf{Definition} (\texttt{cr-forsome-elim}).\;For all $\operatorname{mct\_}$, $\operatorname{mchs\_}$ and $\operatorname{mcc\_}$, 
\begin{equation*}
\begin{array}{@{}l@{}} \operatorname{cr\text{-}forsome\text{-}elim}(\operatorname{mct\_}, \operatorname{mchs\_}, \operatorname{mcc\_}) \\ \Leftrightarrow\; \exists \operatorname{mcg\_}, \operatorname{mcf\_}, \operatorname{mcb\_}, \operatorname{mcy\_}. \\ \quad \operatorname{mcc\_} = {[}\operatorname{mcg\_}, \operatorname{mcf\_}] \wedge \\ \quad \operatorname{enc\text{-}ex}(\operatorname{mcb\_}) \in \operatorname{mcg\_} \wedge \\ \quad \operatorname{mcy\_} \in \mathbb{N} \wedge \\ \quad \neg (\operatorname{mcy\_} \in (\operatorname{ctx\text{-}fv}(\operatorname{mcg\_}) \cup \operatorname{expr\text{-}fv}(\operatorname{mcf\_}))) \wedge \\ \quad \operatorname{mchs\_} = {[}{[}\operatorname{ctx\text{-}add}(\operatorname{mcg\_}, \operatorname{expr\text{-}open}(\operatorname{mcb\_}, 0, \operatorname{efvar}(\operatorname{mcy\_}))), \operatorname{mcf\_}]] \end{array}
\end{equation*}

\par\medskip\noindent\textbf{Definition} (\texttt{cr-assumption}).\;For all $\operatorname{mct\_}$, $\operatorname{mchs\_}$ and $\operatorname{mcc\_}$, 
\begin{equation*}
\begin{array}{@{}l@{}} \operatorname{cr\text{-}assumption}(\operatorname{mct\_}, \operatorname{mchs\_}, \operatorname{mcc\_}) \\ \Leftrightarrow\; \exists \operatorname{mcg\_}, \operatorname{mca\_}.\; ((\operatorname{mcc\_} = {[}\operatorname{mcg\_}, \operatorname{mca\_}]) \wedge ((\operatorname{mca\_} \in \operatorname{mcg\_}) \wedge (\operatorname{mchs\_} = {[}]))) \end{array}
\end{equation*}

\par\medskip\noindent\textbf{Definition} (\texttt{cr-theorem-assumption}).\;For all $\operatorname{mct\_}$, $\operatorname{mchs\_}$ and $\operatorname{mcc\_}$, 
\begin{equation*}
\begin{array}{@{}l@{}} \operatorname{cr\text{-}theorem\text{-}assumption}(\operatorname{mct\_}, \operatorname{mchs\_}, \operatorname{mcc\_}) \\ \Leftrightarrow\; \exists \operatorname{mcg\_}, \operatorname{mca\_}.\; ((\operatorname{mcc\_} = {[}\operatorname{mcg\_}, \operatorname{mca\_}]) \wedge ((\operatorname{mca\_} \in \operatorname{mct\_}) \wedge (\operatorname{mchs\_} = {[}]))) \end{array}
\end{equation*}

\par\medskip\noindent\textbf{Definition} (\texttt{cr-cut}).\;For all $\operatorname{mct\_}$, $\operatorname{mchs\_}$ and $\operatorname{mcc\_}$, 
\begin{equation*}
\begin{array}{@{}l@{}} \operatorname{cr\text{-}cut}(\operatorname{mct\_}, \operatorname{mchs\_}, \operatorname{mcc\_}) \\ \Leftrightarrow\; \exists \operatorname{mcg\_}, \operatorname{mcf\_}, \operatorname{mca\_}. \\ \quad (\operatorname{mcc\_} = {[}\operatorname{mcg\_}, \operatorname{mcf\_}]) \wedge (\operatorname{mchs\_} = {[}{[}\operatorname{mcg\_}, \operatorname{mca\_}], {[}\operatorname{ctx\text{-}add}(\operatorname{mcg\_}, \operatorname{mca\_}), \operatorname{mcf\_}]]) \end{array}
\end{equation*}

\par\medskip\noindent\textbf{Definition} (\texttt{cr-weakening}).\;For all $\operatorname{mct\_}$, $\operatorname{mchs\_}$ and $\operatorname{mcc\_}$, 
\begin{equation*}
\begin{array}{@{}l@{}} \operatorname{cr\text{-}weakening}(\operatorname{mct\_}, \operatorname{mchs\_}, \operatorname{mcc\_}) \\ \Leftrightarrow\; \exists \operatorname{mcg\_}, \operatorname{mcf\_}, \operatorname{mcd\_}. \\ \quad (\operatorname{mcc\_} = {[}\operatorname{mcg\_}, \operatorname{mcf\_}]) \wedge ((\operatorname{mcd\_} \subseteq \operatorname{mcg\_}) \wedge (\operatorname{mchs\_} = {[}{[}\operatorname{mcd\_}, \operatorname{mcf\_}]])) \end{array}
\end{equation*}

\par\medskip\noindent\textbf{Definition} (\texttt{cr-contradiction}).\;For all $\operatorname{mct\_}$, $\operatorname{mchs\_}$ and $\operatorname{mcc\_}$, 
\begin{equation*}
\begin{array}{@{}l@{}} \operatorname{cr\text{-}contradiction}(\operatorname{mct\_}, \operatorname{mchs\_}, \operatorname{mcc\_}) \\ \Leftrightarrow\; \exists \operatorname{mcg\_}, \operatorname{mcf\_}. \\ \quad (\operatorname{mcc\_} = {[}\operatorname{mcg\_}, \operatorname{mcf\_}]) \wedge (\operatorname{mchs\_} = {[}{[}\operatorname{ctx\text{-}add}(\operatorname{mcg\_}, \operatorname{enc\text{-}not}(\operatorname{mcf\_})), \operatorname{enc\text{-}falsity}]]) \end{array}
\end{equation*}

\par\medskip\noindent\textbf{Definition} (\texttt{cr-truth-intro}).\;For all $\operatorname{mct\_}$, $\operatorname{mchs\_}$ and $\operatorname{mcc\_}$, 
\begin{equation*}
\operatorname{cr\text{-}truth\text{-}intro}(\operatorname{mct\_}, \operatorname{mchs\_}, \operatorname{mcc\_}) \Leftrightarrow \exists \operatorname{mcg\_}.\; ((\operatorname{mcc\_} = {[}\operatorname{mcg\_}, \operatorname{enc\text{-}truth}]) \wedge (\operatorname{mchs\_} = {[}]))
\end{equation*}

\par\medskip\noindent\textbf{Definition} (\texttt{cr-eq-subst}).\;For all $\operatorname{mct\_}$, $\operatorname{mchs\_}$ and $\operatorname{mcc\_}$, 
\begin{equation*}
\begin{array}{@{}l@{}} \operatorname{cr\text{-}eq\text{-}subst}(\operatorname{mct\_}, \operatorname{mchs\_}, \operatorname{mcc\_}) \\ \Leftrightarrow\; \exists \operatorname{mcg\_}, \operatorname{mcp\_}, \operatorname{mcz\_}, \operatorname{mcs\_}, \operatorname{mctm\_}. \\ \quad \operatorname{mcc\_} = {[}\operatorname{mcg\_}, \operatorname{expr\text{-}subst}(\operatorname{mcp\_}, \operatorname{mcz\_}, \operatorname{mcs\_})] \wedge \\ \quad \operatorname{mcz\_} \in \mathbb{N} \wedge \\ \quad \operatorname{mcp\_} \in \operatorname{expr\text{-}set} \wedge \\ \quad (\operatorname{enc\text{-}eq}(\operatorname{mcs\_}, \operatorname{mctm\_}) \in \operatorname{mcg\_}) \vee (\operatorname{enc\text{-}qeq}(\operatorname{mcs\_}, \operatorname{mctm\_}) \in \operatorname{mcg\_}) \wedge \\ \quad \operatorname{mchs\_} = {[}{[}\operatorname{mcg\_}, \operatorname{expr\text{-}subst}(\operatorname{mcp\_}, \operatorname{mcz\_}, \operatorname{mctm\_})]] \end{array}
\end{equation*}

\par\medskip\noindent\textbf{Definition} (\texttt{cr-eq-subst-sym}).\;For all $\operatorname{mct\_}$, $\operatorname{mchs\_}$ and $\operatorname{mcc\_}$, 
\begin{equation*}
\begin{array}{@{}l@{}} \operatorname{cr\text{-}eq\text{-}subst\text{-}sym}(\operatorname{mct\_}, \operatorname{mchs\_}, \operatorname{mcc\_}) \\ \Leftrightarrow\; \exists \operatorname{mcg\_}, \operatorname{mcp\_}, \operatorname{mcz\_}, \operatorname{mcs\_}, \operatorname{mctm\_}. \\ \quad \operatorname{mcc\_} = {[}\operatorname{mcg\_}, \operatorname{expr\text{-}subst}(\operatorname{mcp\_}, \operatorname{mcz\_}, \operatorname{mcs\_})] \wedge \\ \quad \operatorname{mcz\_} \in \mathbb{N} \wedge \\ \quad \operatorname{mcp\_} \in \operatorname{expr\text{-}set} \wedge \\ \quad (\operatorname{enc\text{-}eq}(\operatorname{mctm\_}, \operatorname{mcs\_}) \in \operatorname{mcg\_}) \vee (\operatorname{enc\text{-}qeq}(\operatorname{mctm\_}, \operatorname{mcs\_}) \in \operatorname{mcg\_}) \wedge \\ \quad \operatorname{mchs\_} = {[}{[}\operatorname{mcg\_}, \operatorname{expr\text{-}subst}(\operatorname{mcp\_}, \operatorname{mcz\_}, \operatorname{mctm\_})]] \end{array}
\end{equation*}

\par\medskip\noindent\textbf{Definition} (\texttt{cr-reflexivity}).\;For all $\operatorname{mct\_}$, $\operatorname{mchs\_}$ and $\operatorname{mcc\_}$, 
\begin{equation*}
\begin{array}{@{}l@{}} \operatorname{cr\text{-}reflexivity}(\operatorname{mct\_}, \operatorname{mchs\_}, \operatorname{mcc\_}) \\ \Leftrightarrow\; \exists \operatorname{mcg\_}, \operatorname{mctm\_}, \operatorname{mcn\_}. \\ \quad \operatorname{mcc\_} = {[}\operatorname{mcg\_}, \operatorname{enc\text{-}eq}(\operatorname{mctm\_}, \operatorname{mctm\_})] \wedge \\ \quad \operatorname{mcn\_} \in \mathbb{N} \wedge \\ \quad (\operatorname{mctm\_} = \operatorname{efvar}(\operatorname{mcn\_})) \vee (\operatorname{mctm\_} = \operatorname{econst}(\operatorname{mcn\_})) \wedge \\ \quad \operatorname{mchs\_} = {[}] \end{array}
\end{equation*}

\par\medskip\noindent\textbf{Definition} (\texttt{cr-quasi-reflexivity}).\;For all $\operatorname{mct\_}$, $\operatorname{mchs\_}$ and $\operatorname{mcc\_}$, 
\begin{equation*}
\begin{array}{@{}l@{}} \operatorname{cr\text{-}quasi\text{-}reflexivity}(\operatorname{mct\_}, \operatorname{mchs\_}, \operatorname{mcc\_}) \\ \Leftrightarrow\; \exists \operatorname{mcg\_}, \operatorname{mctm\_}.\; ((\operatorname{mcc\_} = {[}\operatorname{mcg\_}, \operatorname{enc\text{-}qeq}(\operatorname{mctm\_}, \operatorname{mctm\_})]) \wedge (\operatorname{mchs\_} = {[}])) \end{array}
\end{equation*}

\par\medskip\noindent\textbf{Definition} (\texttt{strict-head}).\;For every $\operatorname{mchd\_}$, 
\begin{equation*}
\begin{array}{@{}l@{}} \operatorname{strict\text{-}head}(\operatorname{mchd\_}) \\ \Leftrightarrow\; (((7 \leq \operatorname{mchd\_}) \wedge (\operatorname{mchd\_} \leq 8)) \vee (((11 \leq \operatorname{mchd\_}) \wedge (\operatorname{mchd\_} \leq 19)) \vee ((30 \leq \operatorname{mchd\_}) \wedge (\operatorname{mchd\_} \leq 31)))) \end{array}
\end{equation*}

\par\medskip\noindent\textbf{Definition} (\texttt{cr-union-mem}).\;For all $\operatorname{mct\_}$, $\operatorname{mchs\_}$ and $\operatorname{mcc\_}$, 
\begin{equation*}
\begin{array}{@{}l@{}} \operatorname{cr\text{-}union\text{-}mem}(\operatorname{mct\_}, \operatorname{mchs\_}, \operatorname{mcc\_}) \\ \Leftrightarrow\; \exists \operatorname{mcg\_}, \operatorname{mcx\_}, \operatorname{mcl\_}. \\ \quad \operatorname{mcc\_} = {[}\operatorname{mcg\_}, \operatorname{enc\text{-}iff}(\operatorname{enc\text{-}in}(\operatorname{mcx\_}, \operatorname{enc\text{-}union}(\operatorname{mcl\_})), \operatorname{enc\text{-}orl}(\operatorname{mcx\_}, \operatorname{mcl\_}, \operatorname{length}(\operatorname{mcl\_})))] \wedge \\ \quad \operatorname{length}(\operatorname{mchs\_}) = \operatorname{length}(\operatorname{mcl\_}) \wedge \\ \quad \forall \operatorname{mci\_} \in \mathbb{N}. \\ \quad \quad (1 \leq \operatorname{mci\_}) \Rightarrow  \\ \quad \quad (\operatorname{mci\_} \leq \operatorname{length}(\operatorname{mcl\_})) \Rightarrow  \\ \quad \quad \quad \operatorname{nth}(\operatorname{mci\_}, \operatorname{mchs\_}) = {[}\operatorname{mcg\_}, \operatorname{enc\text{-}eq}(\operatorname{nth}(\operatorname{mci\_}, \operatorname{mcl\_}), \operatorname{nth}(\operatorname{mci\_}, \operatorname{mcl\_}))] \end{array}
\end{equation*}

\par\medskip\noindent\textbf{Definition} (\texttt{cr-intersection-mem}).\;For all $\operatorname{mct\_}$, $\operatorname{mchs\_}$ and $\operatorname{mcc\_}$, 
\begin{equation*}
\begin{array}{@{}l@{}} \operatorname{cr\text{-}intersection\text{-}mem}(\operatorname{mct\_}, \operatorname{mchs\_}, \operatorname{mcc\_}) \\ \Leftrightarrow\; \exists \operatorname{mcg\_}, \operatorname{mcx\_}, \operatorname{mcl\_}. \\ \quad \operatorname{mcc\_} = {[}\operatorname{mcg\_}, \operatorname{enc\text{-}iff}(\operatorname{enc\text{-}in}(\operatorname{mcx\_}, \operatorname{enc\text{-}inter}(\operatorname{mcl\_})), \operatorname{enc\text{-}andl}(\operatorname{mcx\_}, \operatorname{mcl\_}, \operatorname{length}(\operatorname{mcl\_})))] \wedge \\ \quad 1 \leq \operatorname{length}(\operatorname{mcl\_}) \wedge \\ \quad \operatorname{mchs\_} = {[}] \end{array}
\end{equation*}

\par\medskip\noindent\textbf{Definition} (\texttt{cr-cartesian-mem}).\;For all $\operatorname{mct\_}$, $\operatorname{mchs\_}$ and $\operatorname{mcc\_}$, 
\begin{equation*}
\begin{array}{@{}l@{}} \operatorname{cr\text{-}cartesian\text{-}mem}(\operatorname{mct\_}, \operatorname{mchs\_}, \operatorname{mcc\_}) \\ \Leftrightarrow\; \exists \operatorname{mcg\_}, \operatorname{mcx\_}, \operatorname{mcl\_}. \\ \quad \operatorname{mcc\_} \\ \quad =\; {[}\operatorname{mcg\_}, \operatorname{enc\text{-}iff}(\operatorname{enc\text{-}in}(\operatorname{mcx\_}, \operatorname{enc\text{-}cart}(\operatorname{mcl\_})), \operatorname{enc\text{-}and}(\operatorname{enc\text{-}in}(\operatorname{mcx\_}, \operatorname{enc\text{-}tuples}(\operatorname{enc\text{-}setc})), \operatorname{enc\text{-}and}(\operatorname{enc\text{-}eq}(\operatorname{enc\text{-}length}(\operatorname{mcx\_}), \operatorname{enc\text{-}num}(\operatorname{length}(\operatorname{mcl\_}))), \operatorname{enc\text{-}cartl}(\operatorname{mcx\_}, \operatorname{mcl\_}, \operatorname{length}(\operatorname{mcl\_})))))] \wedge \\ \quad 1 \leq \operatorname{length}(\operatorname{mcl\_}) \wedge \\ \quad \operatorname{mchs\_} = {[}] \end{array}
\end{equation*}

\par\medskip\noindent\textbf{Definition} (\texttt{cr-tuples-mem}).\;For all $\operatorname{mct\_}$, $\operatorname{mchs\_}$ and $\operatorname{mcc\_}$, 
\begin{equation*}
\begin{array}{@{}l@{}} \operatorname{cr\text{-}tuples\text{-}mem}(\operatorname{mct\_}, \operatorname{mchs\_}, \operatorname{mcc\_}) \\ \Leftrightarrow\; \exists \operatorname{mcg\_}, \operatorname{mcl\_}, \operatorname{mca\_}. \\ \quad \operatorname{mcc\_} = {[}\operatorname{mcg\_}, \operatorname{enc\text{-}iff}(\operatorname{enc\text{-}in}(\operatorname{enc\text{-}list}(\operatorname{mcl\_}), \operatorname{enc\text{-}tuples}(\operatorname{mca\_})), \operatorname{enc\text{-}ande}(\operatorname{mcl\_}, \operatorname{mca\_}, \operatorname{length}(\operatorname{mcl\_})))] \wedge \\ \quad \operatorname{mchs\_} = {[}{[}\operatorname{mcg\_}, \operatorname{enc\text{-}eq}(\operatorname{mca\_}, \operatorname{mca\_})]] \end{array}
\end{equation*}

\par\medskip\noindent\textbf{Definition} (\texttt{cr-nth-eq}).\;For all $\operatorname{mct\_}$, $\operatorname{mchs\_}$ and $\operatorname{mcc\_}$, 
\begin{equation*}
\begin{array}{@{}l@{}} \operatorname{cr\text{-}nth\text{-}eq}(\operatorname{mct\_}, \operatorname{mchs\_}, \operatorname{mcc\_}) \\ \Leftrightarrow\; \exists \operatorname{mcg\_}, \operatorname{mcl\_}, \operatorname{mck\_}. \\ \quad \operatorname{mcc\_} = {[}\operatorname{mcg\_}, \operatorname{enc\text{-}qeq}(\operatorname{enc\text{-}nth}(\operatorname{enc\text{-}num}(\operatorname{mck\_}), \operatorname{enc\text{-}list}(\operatorname{mcl\_})), \operatorname{nth}(\operatorname{mck\_}, \operatorname{mcl\_}))] \wedge \\ \quad \operatorname{mck\_} \in \mathbb{N} \wedge \\ \quad 1 \leq \operatorname{mck\_} \wedge \\ \quad \operatorname{mck\_} \leq \operatorname{length}(\operatorname{mcl\_}) \wedge \\ \quad \operatorname{mchs\_} = {[}{[}\operatorname{mcg\_}, \operatorname{enc\text{-}eq}(\operatorname{enc\text{-}list}(\operatorname{mcl\_}), \operatorname{enc\text{-}list}(\operatorname{mcl\_}))]] \end{array}
\end{equation*}

\par\medskip\noindent\textbf{Definition} (\texttt{cr-length-eq}).\;For all $\operatorname{mct\_}$, $\operatorname{mchs\_}$ and $\operatorname{mcc\_}$, 
\begin{equation*}
\begin{array}{@{}l@{}} \operatorname{cr\text{-}length\text{-}eq}(\operatorname{mct\_}, \operatorname{mchs\_}, \operatorname{mcc\_}) \\ \Leftrightarrow\; \exists \operatorname{mcg\_}, \operatorname{mcl\_}. \\ \quad \operatorname{mcc\_} = {[}\operatorname{mcg\_}, \operatorname{enc\text{-}qeq}(\operatorname{enc\text{-}length}(\operatorname{enc\text{-}list}(\operatorname{mcl\_})), \operatorname{enc\text{-}num}(\operatorname{length}(\operatorname{mcl\_})))] \wedge \\ \quad \operatorname{mchs\_} = {[}{[}\operatorname{mcg\_}, \operatorname{enc\text{-}eq}(\operatorname{enc\text{-}list}(\operatorname{mcl\_}), \operatorname{enc\text{-}list}(\operatorname{mcl\_}))]] \end{array}
\end{equation*}

\par\medskip\noindent\textbf{Definition} (\texttt{cr-tuple-eq}).\;For all $\operatorname{mct\_}$, $\operatorname{mchs\_}$ and $\operatorname{mcc\_}$, 
\begin{equation*}
\begin{array}{@{}l@{}} \operatorname{cr\text{-}tuple\text{-}eq}(\operatorname{mct\_}, \operatorname{mchs\_}, \operatorname{mcc\_}) \\ \Leftrightarrow\; \exists \operatorname{mcg\_}, \operatorname{mcl\_}, \operatorname{mca\_}. \\ \quad \operatorname{mcc\_} = {[}\operatorname{mcg\_}, \operatorname{enc\text{-}iff}(\operatorname{enc\text{-}eq}(\operatorname{enc\text{-}list}(\operatorname{mcl\_}), \operatorname{enc\text{-}list}(\operatorname{mca\_})), \operatorname{enc\text{-}eql}(\operatorname{mcl\_}, \operatorname{mca\_}, \operatorname{length}(\operatorname{mcl\_})))] \wedge \\ \quad \operatorname{length}(\operatorname{mcl\_}) = \operatorname{length}(\operatorname{mca\_}) \wedge \\ \quad \operatorname{mchs\_} \\ \quad =\; {[}{[}\operatorname{mcg\_}, \operatorname{enc\text{-}eq}(\operatorname{enc\text{-}list}(\operatorname{mcl\_}), \operatorname{enc\text{-}list}(\operatorname{mcl\_}))], {[}\operatorname{mcg\_}, \operatorname{enc\text{-}eq}(\operatorname{enc\text{-}list}(\operatorname{mca\_}), \operatorname{enc\text{-}list}(\operatorname{mca\_}))]] \end{array}
\end{equation*}

\par\medskip\noindent\textbf{Definition} (\texttt{cr-if-true}).\;For all $\operatorname{mct\_}$, $\operatorname{mchs\_}$ and $\operatorname{mcc\_}$, 
\begin{equation*}
\begin{array}{@{}l@{}} \operatorname{cr\text{-}if\text{-}true}(\operatorname{mct\_}, \operatorname{mchs\_}, \operatorname{mcc\_}) \\ \Leftrightarrow\; \exists \operatorname{mcg\_}, \operatorname{mcf\_}, \operatorname{mcp\_}, \operatorname{mca\_}, \operatorname{mcb\_}. \\ \quad \operatorname{mcc\_} = {[}\operatorname{mcg\_}, \operatorname{mcf\_}] \wedge \\ \quad \operatorname{mchs\_} = {[}{[}\operatorname{mcg\_}, \operatorname{mcp\_}], {[}\operatorname{ctx\text{-}add}(\operatorname{mcg\_}, \operatorname{enc\text{-}qeq}(\operatorname{enc\text{-}if}(\operatorname{mcp\_}, \operatorname{mca\_}, \operatorname{mcb\_}), \operatorname{mca\_})), \operatorname{mcf\_}]] \wedge \\ \quad \operatorname{mca\_} \in \mathit{clt} \wedge \\ \quad \operatorname{mcb\_} \in \mathit{clt} \end{array}
\end{equation*}

\par\medskip\noindent\textbf{Definition} (\texttt{cr-if-false}).\;For all $\operatorname{mct\_}$, $\operatorname{mchs\_}$ and $\operatorname{mcc\_}$, 
\begin{equation*}
\begin{array}{@{}l@{}} \operatorname{cr\text{-}if\text{-}false}(\operatorname{mct\_}, \operatorname{mchs\_}, \operatorname{mcc\_}) \\ \Leftrightarrow\; \exists \operatorname{mcg\_}, \operatorname{mcf\_}, \operatorname{mcp\_}, \operatorname{mca\_}, \operatorname{mcb\_}. \\ \quad \operatorname{mcc\_} = {[}\operatorname{mcg\_}, \operatorname{mcf\_}] \wedge \\ \quad \operatorname{mchs\_} = {[}{[}\operatorname{mcg\_}, \operatorname{enc\text{-}not}(\operatorname{mcp\_})], {[}\operatorname{ctx\text{-}add}(\operatorname{mcg\_}, \operatorname{enc\text{-}qeq}(\operatorname{enc\text{-}if}(\operatorname{mcp\_}, \operatorname{mca\_}, \operatorname{mcb\_}), \operatorname{mcb\_})), \operatorname{mcf\_}]] \wedge \\ \quad \operatorname{mca\_} \in \mathit{clt} \wedge \\ \quad \operatorname{mcb\_} \in \mathit{clt} \end{array}
\end{equation*}

\par\medskip\noindent\textbf{Definition} (\texttt{fun-def}).\;For all $\operatorname{mcfh\_}$, $\operatorname{mcfn\_}$ and $\operatorname{mcfb\_}$, 
\begin{equation*}
\begin{array}{@{}l@{}} \operatorname{fun\text{-}def}(\operatorname{mcfh\_}, \operatorname{mcfn\_}, \operatorname{mcfb\_}) \\ \Leftrightarrow\; ((\operatorname{mcfh\_} = 32) \wedge ((\operatorname{mcfn\_} = 2) \wedge (\operatorname{mcfb\_} = \operatorname{enc\text{-}sep}(\operatorname{ebvar}(0), \operatorname{enc\text{-}in}(\operatorname{ebvar}(0), \operatorname{ebvar}(2)))))) \end{array}
\end{equation*}

\par\medskip\noindent\textbf{Definition} (\texttt{cr-sep-sethood}).\;For all $\operatorname{mct\_}$, $\operatorname{mchs\_}$ and $\operatorname{mcc\_}$, 
\begin{equation*}
\begin{array}{@{}l@{}} \operatorname{cr\text{-}sep\text{-}sethood}(\operatorname{mct\_}, \operatorname{mchs\_}, \operatorname{mcc\_}) \\ \Leftrightarrow\; \exists \operatorname{mcg\_}, \operatorname{mca\_}, \operatorname{mcp\_}. \\ \quad \operatorname{mcc\_} = {[}\operatorname{mcg\_}, \operatorname{enc\text{-}in}(\operatorname{enc\text{-}sep}(\operatorname{mca\_}, \operatorname{mcp\_}), \operatorname{enc\text{-}setc})] \wedge \\ \quad \operatorname{nbg\text{-}ok}(\operatorname{mcp\_}) \wedge \\ \quad \operatorname{mchs\_} = {[}{[}\operatorname{mcg\_}, \operatorname{enc\text{-}in}(\operatorname{mca\_}, \operatorname{enc\text{-}setc})]] \end{array}
\end{equation*}

\par\medskip\noindent\textbf{Definition} (\texttt{cr-sep-mem}).\;For all $\operatorname{mct\_}$, $\operatorname{mchs\_}$ and $\operatorname{mcc\_}$, 
\begin{equation*}
\begin{array}{@{}l@{}} \operatorname{cr\text{-}sep\text{-}mem}(\operatorname{mct\_}, \operatorname{mchs\_}, \operatorname{mcc\_}) \\ \Leftrightarrow\; \exists \operatorname{mcg\_}, \operatorname{mca\_}, \operatorname{mcp\_}, \operatorname{mctm\_}. \\ \quad \operatorname{mcc\_} \\ \quad =\; {[}\operatorname{mcg\_}, \operatorname{enc\text{-}iff}(\operatorname{enc\text{-}in}(\operatorname{mctm\_}, \operatorname{enc\text{-}sep}(\operatorname{mca\_}, \operatorname{mcp\_})), \operatorname{enc\text{-}and}(\operatorname{enc\text{-}in}(\operatorname{mctm\_}, \operatorname{mca\_}), \operatorname{expr\text{-}open}(\operatorname{mcp\_}, 0, \operatorname{mctm\_})))] \wedge \\ \quad \operatorname{nbg\text{-}ok}(\operatorname{mcp\_}) \wedge \\ \quad \operatorname{mchs\_} = {[}] \end{array}
\end{equation*}

\par\medskip\noindent\textbf{Definition} (\texttt{cr-comp-mem}).\;For all $\operatorname{mct\_}$, $\operatorname{mchs\_}$ and $\operatorname{mcc\_}$, 
\begin{equation*}
\begin{array}{@{}l@{}} \operatorname{cr\text{-}comp\text{-}mem}(\operatorname{mct\_}, \operatorname{mchs\_}, \operatorname{mcc\_}) \\ \Leftrightarrow\; \exists \operatorname{mcg\_}, \operatorname{mcp\_}, \operatorname{mctm\_}. \\ \quad \operatorname{mcc\_} \\ \quad =\; {[}\operatorname{mcg\_}, \operatorname{enc\text{-}iff}(\operatorname{enc\text{-}in}(\operatorname{mctm\_}, \operatorname{enc\text{-}comp}(\operatorname{mcp\_})), \operatorname{enc\text{-}and}(\operatorname{enc\text{-}in}(\operatorname{mctm\_}, \operatorname{enc\text{-}setc}), \operatorname{expr\text{-}open}(\operatorname{mcp\_}, 0, \operatorname{mctm\_})))] \wedge \\ \quad \operatorname{nbg\text{-}ok}(\operatorname{mcp\_}) \wedge \\ \quad \operatorname{mchs\_} = {[}] \end{array}
\end{equation*}

\par\medskip\noindent\textbf{Definition} (\texttt{cr-big-union-sethood}).\;For all $\operatorname{mct\_}$, $\operatorname{mchs\_}$ and $\operatorname{mcc\_}$, 
\begin{equation*}
\begin{array}{@{}l@{}} \operatorname{cr\text{-}big\text{-}union\text{-}sethood}(\operatorname{mct\_}, \operatorname{mchs\_}, \operatorname{mcc\_}) \\ \Leftrightarrow\; \exists \operatorname{mcg\_}, \operatorname{mca\_}, \operatorname{mcb\_}. \\ \quad \operatorname{mcc\_} = {[}\operatorname{mcg\_}, \operatorname{enc\text{-}in}(\operatorname{enc\text{-}bu}(\operatorname{mca\_}, \operatorname{mcb\_}), \operatorname{enc\text{-}setc})] \wedge \\ \quad \operatorname{nbg\text{-}ok}(\operatorname{mcb\_}) \wedge \\ \quad \operatorname{mchs\_} \\ \quad =\; {[}{[}\operatorname{mcg\_}, \operatorname{enc\text{-}in}(\operatorname{mca\_}, \operatorname{enc\text{-}setc})], {[}\operatorname{mcg\_}, \operatorname{enc\text{-}all}(\operatorname{enc\text{-}imp}(\operatorname{enc\text{-}in}(\operatorname{ebvar}(0), \operatorname{mca\_}), \operatorname{enc\text{-}in}(\operatorname{mcb\_}, \operatorname{enc\text{-}setc})))]] \end{array}
\end{equation*}

\par\medskip\noindent\textbf{Definition} (\texttt{cr-big-union-mem}).\;For all $\operatorname{mct\_}$, $\operatorname{mchs\_}$ and $\operatorname{mcc\_}$, 
\begin{equation*}
\begin{array}{@{}l@{}} \operatorname{cr\text{-}big\text{-}union\text{-}mem}(\operatorname{mct\_}, \operatorname{mchs\_}, \operatorname{mcc\_}) \\ \Leftrightarrow\; \exists \operatorname{mcg\_}, \operatorname{mca\_}, \operatorname{mcb\_}, \operatorname{mcx\_}. \\ \quad \operatorname{mcc\_} \\ \quad =\; {[}\operatorname{mcg\_}, \operatorname{enc\text{-}iff}(\operatorname{enc\text{-}in}(\operatorname{mcx\_}, \operatorname{enc\text{-}bu}(\operatorname{mca\_}, \operatorname{mcb\_})), \operatorname{enc\text{-}ex}(\operatorname{enc\text{-}and}(\operatorname{enc\text{-}in}(\operatorname{ebvar}(0), \operatorname{mca\_}), \operatorname{enc\text{-}in}(\operatorname{mcx\_}, \operatorname{mcb\_}))))] \wedge \\ \quad \operatorname{nbg\text{-}ok}(\operatorname{mcb\_}) \wedge \\ \quad \operatorname{mchs\_} = {[}] \end{array}
\end{equation*}

\par\medskip\noindent\textbf{Definition} (\texttt{cr-lambda-type}).\;For all $\operatorname{mct\_}$, $\operatorname{mchs\_}$ and $\operatorname{mcc\_}$, 
\begin{equation*}
\begin{array}{@{}l@{}} \operatorname{cr\text{-}lambda\text{-}type}(\operatorname{mct\_}, \operatorname{mchs\_}, \operatorname{mcc\_}) \\ \Leftrightarrow\; \exists \operatorname{mcg\_}, \operatorname{mca\_}, \operatorname{mcb\_}, \operatorname{mcd\_}. \\ \quad \operatorname{mcc\_} = {[}\operatorname{mcg\_}, \operatorname{enc\text{-}in}(\operatorname{enc\text{-}lam}(\operatorname{mca\_}, \operatorname{mcb\_}), \operatorname{enc\text{-}fun}(\operatorname{mca\_}, \operatorname{mcd\_}))] \wedge \\ \quad \operatorname{nbg\text{-}ok}(\operatorname{mcb\_}) \wedge \\ \quad \operatorname{mchs\_} \\ \quad =\; {[}{[}\operatorname{mcg\_}, \operatorname{enc\text{-}all}(\operatorname{enc\text{-}imp}(\operatorname{enc\text{-}in}(\operatorname{ebvar}(0), \operatorname{mca\_}), \operatorname{enc\text{-}in}(\operatorname{mcb\_}, \operatorname{mcd\_})))], {[}\operatorname{mcg\_}, \operatorname{enc\text{-}in}(\operatorname{mca\_}, \operatorname{enc\text{-}setc})], {[}\operatorname{mcg\_}, \operatorname{enc\text{-}eq}(\operatorname{mcd\_}, \operatorname{mcd\_})]] \end{array}
\end{equation*}

\par\medskip\noindent\textbf{Definition} (\texttt{cr-lambda-beta-eq}).\;For all $\operatorname{mct\_}$, $\operatorname{mchs\_}$ and $\operatorname{mcc\_}$, 
\begin{equation*}
\begin{array}{@{}l@{}} \operatorname{cr\text{-}lambda\text{-}beta\text{-}eq}(\operatorname{mct\_}, \operatorname{mchs\_}, \operatorname{mcc\_}) \\ \Leftrightarrow\; \exists \operatorname{mcg\_}, \operatorname{mca\_}, \operatorname{mcb\_}, \operatorname{mctm\_}. \\ \quad \operatorname{mcc\_} = {[}\operatorname{mcg\_}, \operatorname{enc\text{-}qeq}(\operatorname{enc\text{-}apply}(\operatorname{enc\text{-}lam}(\operatorname{mca\_}, \operatorname{mcb\_}), \operatorname{mctm\_}), \operatorname{expr\text{-}open}(\operatorname{mcb\_}, 0, \operatorname{mctm\_}))] \wedge \\ \quad \operatorname{nbg\text{-}ok}(\operatorname{mcb\_}) \wedge \\ \quad \operatorname{mchs\_} = {[}{[}\operatorname{mcg\_}, \operatorname{enc\text{-}in}(\operatorname{mctm\_}, \operatorname{mca\_})], {[}\operatorname{mcg\_}, \operatorname{enc\text{-}in}(\operatorname{expr\text{-}open}(\operatorname{mcb\_}, 0, \operatorname{mctm\_}), \operatorname{enc\text{-}setc})]] \end{array}
\end{equation*}

\par\medskip\noindent\textbf{Definition} (\texttt{cr-def-binder}).\;For all $\operatorname{mct\_}$, $\operatorname{mchs\_}$ and $\operatorname{mcc\_}$, 
\begin{equation*}
\begin{array}{@{}l@{}} \operatorname{cr\text{-}def\text{-}binder}(\operatorname{mct\_}, \operatorname{mchs\_}, \operatorname{mcc\_}) \\ \Leftrightarrow\; \exists \operatorname{mcg\_}, \operatorname{mchd\_}, \operatorname{mca\_}, \operatorname{mcb\_}. \\ \quad \operatorname{mcc\_} = {[}\operatorname{mcg\_}, \operatorname{enc\text{-}eq}(\operatorname{ebind}(\operatorname{mchd\_}, 1, {[}\operatorname{mca\_}], \operatorname{mcb\_}), \operatorname{ebind}(\operatorname{mchd\_}, 1, {[}\operatorname{mca\_}], \operatorname{mcb\_}))] \wedge \\ \quad (\operatorname{mchd\_} = 22) \vee ((\operatorname{mchd\_} = 24) \vee (\operatorname{mchd\_} = 26)) \wedge \\ \quad \operatorname{nbg\text{-}ok}(\operatorname{mcb\_}) \wedge \\ \quad \operatorname{mchs\_} = {[}{[}\operatorname{mcg\_}, \operatorname{enc\text{-}eq}(\operatorname{mca\_}, \operatorname{mca\_})]] \end{array}
\end{equation*}

\par\medskip\noindent\textbf{Definition} (\texttt{cr-def-comp}).\;For all $\operatorname{mct\_}$, $\operatorname{mchs\_}$ and $\operatorname{mcc\_}$, 
\begin{equation*}
\begin{array}{@{}l@{}} \operatorname{cr\text{-}def\text{-}comp}(\operatorname{mct\_}, \operatorname{mchs\_}, \operatorname{mcc\_}) \\ \Leftrightarrow\; \exists \operatorname{mcg\_}, \operatorname{mcp\_}. \\ \quad ((\operatorname{mcc\_} = {[}\operatorname{mcg\_}, \operatorname{enc\text{-}eq}(\operatorname{enc\text{-}comp}(\operatorname{mcp\_}), \operatorname{enc\text{-}comp}(\operatorname{mcp\_}))]) \wedge (\operatorname{nbg\text{-}ok}(\operatorname{mcp\_}) \wedge (\operatorname{mchs\_} = {[}]))) \end{array}
\end{equation*}

\par\medskip\noindent\textbf{Definition} (\texttt{nbg-ok}).\;For every $\operatorname{mnbp\_}$, 
\begin{equation*}
\operatorname{nbg\text{-}ok}(\operatorname{mnbp\_}) \Leftrightarrow \forall \operatorname{mnboe\_} \in \operatorname{nbg\text{-}walk}(\operatorname{mnbp\_}).\; \operatorname{qnode\text{-}ok}(\operatorname{mnboe\_})
\end{equation*}

\par\medskip\noindent\textbf{Definition} (\texttt{all-spine}).\;$\operatorname{all\text{-}spine}(\operatorname{mnbb\_}) :=$
\begin{equation*}
\{\operatorname{mnbz\_} \in (\operatorname{expr\text{-}set} \times \mathbb{N}) : \forall \operatorname{mnbv\_} \in \operatorname{power}((\operatorname{expr\text{-}set} \times \mathbb{N})).\; (({[}\operatorname{mnbb\_}, 0] \in \operatorname{mnbv\_}) \Rightarrow (\operatorname{all\text{-}spine\text{-}closed}(\operatorname{mnbv\_}) \Rightarrow (\operatorname{mnbz\_} \in \operatorname{mnbv\_})))\}
\end{equation*}

\par\medskip\noindent\textbf{Definition} (\texttt{ex-spine}).\;$\operatorname{ex\text{-}spine}(\operatorname{mnbb\_}) :=$
\begin{equation*}
\{\operatorname{mnbz\_} \in (\operatorname{expr\text{-}set} \times \mathbb{N}) : \forall \operatorname{mnbv\_} \in \operatorname{power}((\operatorname{expr\text{-}set} \times \mathbb{N})).\; (({[}\operatorname{mnbb\_}, 0] \in \operatorname{mnbv\_}) \Rightarrow (\operatorname{ex\text{-}spine\text{-}closed}(\operatorname{mnbv\_}) \Rightarrow (\operatorname{mnbz\_} \in \operatorname{mnbv\_})))\}
\end{equation*}

\par\medskip\noindent\textbf{Definition} (\texttt{nbg-walk}).\;$\operatorname{nbg\text{-}walk}(\operatorname{mnbp\_}) :=$
\begin{equation*}
\{\operatorname{mnbz\_} \in \operatorname{expr\text{-}set} : \forall \operatorname{mnbv\_} \in \operatorname{power}(\operatorname{expr\text{-}set}).\; ((\operatorname{mnbp\_} \in \operatorname{mnbv\_}) \Rightarrow (\operatorname{nbg\text{-}closed}(\operatorname{mnbv\_}) \Rightarrow (\operatorname{mnbz\_} \in \operatorname{mnbv\_})))\}
\end{equation*}

\par\medskip\noindent\textbf{Definition} (\texttt{nbg-step}).\;For all $\operatorname{mnbe\_}$ and $\operatorname{mnbf\_}$, 
\begin{equation*}
\begin{array}{@{}l@{}} \operatorname{nbg\text{-}step}(\operatorname{mnbe\_}, \operatorname{mnbf\_}) \\ \Leftrightarrow\; \exists \operatorname{mnb1h\_} \in \mathbb{N}, \operatorname{mnb1l\_} \in \mathbf{Set}, \operatorname{mnb1i\_} \in \mathbb{N}. \\ \quad \operatorname{mnbe\_} = \operatorname{eapp}(\operatorname{mnb1h\_}, \operatorname{mnb1l\_}) \wedge \\ \quad 1 \leq \operatorname{mnb1i\_} \wedge \\ \quad \operatorname{mnb1i\_} \leq \operatorname{length}(\operatorname{mnb1l\_}) \wedge \\ \quad \operatorname{mnbf\_} = \operatorname{nth}(\operatorname{mnb1i\_}, \operatorname{mnb1l\_}) \vee \\ \exists \operatorname{mnb2h\_}, \operatorname{mnb2n\_} \in \mathbb{N}, \operatorname{mnb2a\_} \in \mathbf{Set}, \operatorname{mnb2b\_} \in \operatorname{expr\text{-}set}, \operatorname{mnb2i\_} \in \mathbb{N}. \\ \quad \operatorname{mnbe\_} = \operatorname{ebind}(\operatorname{mnb2h\_}, \operatorname{mnb2n\_}, \operatorname{mnb2a\_}, \operatorname{mnb2b\_}) \wedge \\ \quad 1 \leq \operatorname{mnb2i\_} \wedge \\ \quad \operatorname{mnb2i\_} \leq \operatorname{length}(\operatorname{mnb2a\_}) \wedge \\ \quad \operatorname{mnbf\_} = \operatorname{nth}(\operatorname{mnb2i\_}, \operatorname{mnb2a\_}) \vee \\ \exists \operatorname{mnb3h\_}, \operatorname{mnb3n\_} \in \mathbb{N}, \operatorname{mnb3a\_} \in \mathbf{Set}. \\ \quad \operatorname{mnbe\_} = \operatorname{ebind}(\operatorname{mnb3h\_}, \operatorname{mnb3n\_}, \operatorname{mnb3a\_}, \operatorname{mnbf\_}) \vee \\ \exists \operatorname{mnb4h\_} \in \mathbb{N}, \operatorname{mnb4l\_} \in \mathbf{Set}. \\ \quad (\operatorname{mnbe\_} = \operatorname{eapp}(\operatorname{mnb4h\_}, \operatorname{mnb4l\_})) \wedge \operatorname{fun\text{-}def}(\operatorname{mnb4h\_}, \operatorname{length}(\operatorname{mnb4l\_}), \operatorname{mnbf\_}) \vee \\ \exists \operatorname{mnb5l\_} \in \mathbf{Set}, \operatorname{mnb5k\_} \in \mathbb{N}. \\ \quad \operatorname{mnbe\_} = \operatorname{eapp}(10, \operatorname{mnb5l\_}) \wedge \\ \quad \operatorname{nth}(1, \operatorname{mnb5l\_}) = \operatorname{econst}(\operatorname{mnb5k\_}) \wedge \\ \quad \operatorname{pred\text{-}def}(\operatorname{mnb5k\_}, \operatorname{length}(\operatorname{mnb5l\_}), \operatorname{mnbf\_}) \end{array}
\end{equation*}

\par\medskip\noindent\textbf{Definition} (\texttt{pred-def}).\;For all $\operatorname{mcdk\_}$, $\operatorname{mcdn\_}$ and $\operatorname{mcdb\_}$, 
\begin{equation*}
\operatorname{pred\text{-}def}(\operatorname{mcdk\_}, \operatorname{mcdn\_}, \operatorname{mcdb\_}) \Leftrightarrow ((\operatorname{mcdk\_} \in \mathbb{N}) \wedge ((\operatorname{mcdn\_} \in \mathbb{N}) \wedge ((\operatorname{mcdb\_} \in \operatorname{expr\text{-}set}) \wedge (0 = 1))))
\end{equation*}

\par\medskip\noindent\textbf{Definition} (\texttt{cr-functoid-beta-eq}).\;For all $\operatorname{mct\_}$, $\operatorname{mchs\_}$ and $\operatorname{mcc\_}$, 
\begin{equation*}
\begin{array}{@{}l@{}} \operatorname{cr\text{-}functoid\text{-}beta\text{-}eq}(\operatorname{mct\_}, \operatorname{mchs\_}, \operatorname{mcc\_}) \\ \Leftrightarrow\; \exists \operatorname{mcg\_}, \operatorname{mchd\_}, \operatorname{mcl\_}, \operatorname{mcb\_}. \\ \quad \operatorname{mcc\_} = {[}\operatorname{mcg\_}, \operatorname{enc\text{-}qeq}(\operatorname{eapp}(\operatorname{mchd\_}, \operatorname{mcl\_}), \operatorname{openl}(\operatorname{mcb\_}, \operatorname{mcl\_}, \operatorname{length}(\operatorname{mcl\_})))] \wedge \\ \quad \operatorname{fun\text{-}def}(\operatorname{mchd\_}, \operatorname{length}(\operatorname{mcl\_}), \operatorname{mcb\_}) \wedge \\ \quad \operatorname{mchs\_} = {[}] \end{array}
\end{equation*}

\par\medskip\noindent\textbf{Definition} (\texttt{nbg-closed}).\;For every $\operatorname{mnbw\_}$, 
\begin{equation*}
\begin{array}{@{}l@{}} \operatorname{nbg\text{-}closed}(\operatorname{mnbw\_}) \\ \Leftrightarrow\; \forall \operatorname{mnbce\_} \in \operatorname{mnbw\_}, \operatorname{mnbcf\_} \in \operatorname{expr\text{-}set}. \\ \quad \operatorname{nbg\text{-}step}(\operatorname{mnbce\_}, \operatorname{mnbcf\_}) \Rightarrow (\operatorname{mnbcf\_} \in \operatorname{mnbw\_}) \end{array}
\end{equation*}

\par\medskip\noindent\textbf{Definition} (\texttt{conj-step}).\;For all $\operatorname{mnbx\_}$ and $\operatorname{mnby\_}$, 
\begin{equation*}
\begin{array}{@{}l@{}} \operatorname{conj\text{-}step}(\operatorname{mnbx\_}, \operatorname{mnby\_}) \\ \Leftrightarrow\; \exists \operatorname{mnbcp\_}, \operatorname{mnbcq\_} \in \operatorname{expr\text{-}set}. \\ \quad (\operatorname{mnbx\_} = \operatorname{enc\text{-}and}(\operatorname{mnbcp\_}, \operatorname{mnbcq\_})) \wedge ((\operatorname{mnby\_} = \operatorname{mnbcp\_}) \vee (\operatorname{mnby\_} = \operatorname{mnbcq\_})) \end{array}
\end{equation*}

\par\medskip\noindent\textbf{Definition} (\texttt{conj-closed}).\;For every $\operatorname{mnbw\_}$, 
\begin{equation*}
\begin{array}{@{}l@{}} \operatorname{conj\text{-}closed}(\operatorname{mnbw\_}) \\ \Leftrightarrow\; \forall \operatorname{mnbce\_} \in \operatorname{mnbw\_}, \operatorname{mnbcf\_} \in \operatorname{expr\text{-}set}. \\ \quad \operatorname{conj\text{-}step}(\operatorname{mnbce\_}, \operatorname{mnbcf\_}) \Rightarrow (\operatorname{mnbcf\_} \in \operatorname{mnbw\_}) \end{array}
\end{equation*}

\par\medskip\noindent\textbf{Definition} (\texttt{conj-reach}).\;$\operatorname{conj\text{-}reach}(\operatorname{mnbx\_}) :=$
\begin{equation*}
\{\operatorname{mnbz\_} \in \operatorname{expr\text{-}set} : \forall \operatorname{mnbv\_} \in \operatorname{power}(\operatorname{expr\text{-}set}).\; ((\operatorname{mnbx\_} \in \operatorname{mnbv\_}) \Rightarrow (\operatorname{conj\text{-}closed}(\operatorname{mnbv\_}) \Rightarrow (\operatorname{mnbz\_} \in \operatorname{mnbv\_})))\}
\end{equation*}

\par\medskip\noindent\textbf{Definition} (\texttt{all-spine-step}).\;For all $\operatorname{mnbu\_}$ and $\operatorname{mnbt\_}$, 
\begin{equation*}
\begin{array}{@{}l@{}} \operatorname{all\text{-}spine\text{-}step}(\operatorname{mnbu\_}, \operatorname{mnbt\_}) \\ \Leftrightarrow\; \exists \operatorname{mnbsc\_} \in \operatorname{expr\text{-}set}, \operatorname{mnbsd\_} \in \mathbb{N}. \\ \quad (\operatorname{mnbu\_} = {[}\operatorname{enc\text{-}all}(\operatorname{mnbsc\_}), \operatorname{mnbsd\_}]) \wedge (\operatorname{mnbt\_} = {[}\operatorname{mnbsc\_}, +(\operatorname{mnbsd\_}, 1)]) \vee \\ \exists \operatorname{mnbsx\_}, \operatorname{mnbsy\_} \in \operatorname{expr\text{-}set}, \operatorname{mnbse\_} \in \mathbb{N}. \\ \quad (\operatorname{mnbu\_} = {[}\operatorname{enc\text{-}imp}(\operatorname{mnbsx\_}, \operatorname{mnbsy\_}), \operatorname{mnbse\_}]) \wedge (\operatorname{mnbt\_} = {[}\operatorname{mnbsy\_}, \operatorname{mnbse\_}]) \end{array}
\end{equation*}

\par\medskip\noindent\textbf{Definition} (\texttt{all-spine-closed}).\;For every $\operatorname{mnbw\_}$, 
\begin{equation*}
\begin{array}{@{}l@{}} \operatorname{all\text{-}spine\text{-}closed}(\operatorname{mnbw\_}) \\ \Leftrightarrow\; \forall \operatorname{mnbce\_} \in \operatorname{mnbw\_}, \operatorname{mnbcf\_} \in (\operatorname{expr\text{-}set} \times \mathbb{N}). \\ \quad \operatorname{all\text{-}spine\text{-}step}(\operatorname{mnbce\_}, \operatorname{mnbcf\_}) \Rightarrow (\operatorname{mnbcf\_} \in \operatorname{mnbw\_}) \end{array}
\end{equation*}

\par\medskip\noindent\textbf{Definition} (\texttt{ex-spine-step}).\;For all $\operatorname{mnbu\_}$ and $\operatorname{mnbt\_}$, 
\begin{equation*}
\begin{array}{@{}l@{}} \operatorname{ex\text{-}spine\text{-}step}(\operatorname{mnbu\_}, \operatorname{mnbt\_}) \\ \Leftrightarrow\; \exists \operatorname{mnbsc\_} \in \operatorname{expr\text{-}set}, \operatorname{mnbsd\_} \in \mathbb{N}. \\ \quad (\operatorname{mnbu\_} = {[}\operatorname{enc\text{-}ex}(\operatorname{mnbsc\_}), \operatorname{mnbsd\_}]) \wedge (\operatorname{mnbt\_} = {[}\operatorname{mnbsc\_}, +(\operatorname{mnbsd\_}, 1)]) \vee \\ \exists \operatorname{mnbsx\_}, \operatorname{mnbsy\_} \in \operatorname{expr\text{-}set}, \operatorname{mnbse\_} \in \mathbb{N}. \\ \quad ((\operatorname{mnbu\_} = {[}\operatorname{enc\text{-}and}(\operatorname{mnbsx\_}, \operatorname{mnbsy\_}), \operatorname{mnbse\_}]) \wedge ((\operatorname{mnbt\_} = {[}\operatorname{mnbsx\_}, \operatorname{mnbse\_}]) \vee (\operatorname{mnbt\_} = {[}\operatorname{mnbsy\_}, \operatorname{mnbse\_}]))) \end{array}
\end{equation*}

\par\medskip\noindent\textbf{Definition} (\texttt{ex-spine-closed}).\;For every $\operatorname{mnbw\_}$, 
\begin{equation*}
\begin{array}{@{}l@{}} \operatorname{ex\text{-}spine\text{-}closed}(\operatorname{mnbw\_}) \\ \Leftrightarrow\; \forall \operatorname{mnbce\_} \in \operatorname{mnbw\_}, \operatorname{mnbcf\_} \in (\operatorname{expr\text{-}set} \times \mathbb{N}). \\ \quad \operatorname{ex\text{-}spine\text{-}step}(\operatorname{mnbce\_}, \operatorname{mnbcf\_}) \Rightarrow (\operatorname{mnbcf\_} \in \operatorname{mnbw\_}) \end{array}
\end{equation*}

\par\medskip\noindent\textbf{Definition} (\texttt{all-confined}).\;For every $\operatorname{mnbb\_}$, 
\begin{equation*}
\begin{array}{@{}l@{}} \operatorname{all\text{-}confined}(\operatorname{mnbb\_}) \\ \Leftrightarrow\; \exists \operatorname{mnbqx\_}, \operatorname{mnbqy\_} \in \operatorname{expr\text{-}set}, \operatorname{mnbqd\_} \in \mathbb{N}, \operatorname{mnbqc\_} \in \operatorname{expr\text{-}set}. \\ \quad {[}\operatorname{enc\text{-}imp}(\operatorname{mnbqx\_}, \operatorname{mnbqy\_}), \operatorname{mnbqd\_}] \in \operatorname{all\text{-}spine}(\operatorname{mnbb\_}) \wedge \\ \quad \operatorname{enc\text{-}in}(\operatorname{ebvar}(\operatorname{mnbqd\_}), \operatorname{mnbqc\_}) \in \operatorname{conj\text{-}reach}(\operatorname{mnbqx\_}) \end{array}
\end{equation*}

\par\medskip\noindent\textbf{Definition} (\texttt{ex-confined}).\;For every $\operatorname{mnbb\_}$, 
\begin{equation*}
\begin{array}{@{}l@{}} \operatorname{ex\text{-}confined}(\operatorname{mnbb\_}) \\ \Leftrightarrow\; \exists \operatorname{mnbqd\_} \in \mathbb{N}, \operatorname{mnbqc\_} \in \operatorname{expr\text{-}set}. \\ \quad {[}\operatorname{enc\text{-}in}(\operatorname{ebvar}(\operatorname{mnbqd\_}), \operatorname{mnbqc\_}), \operatorname{mnbqd\_}] \in \operatorname{ex\text{-}spine}(\operatorname{mnbb\_}) \end{array}
\end{equation*}

\par\medskip\noindent\textbf{Definition} (\texttt{iota-confined}).\;For every $\operatorname{mnbb\_}$, 
\begin{equation*}
\begin{array}{@{}l@{}} \operatorname{iota\text{-}confined}(\operatorname{mnbb\_}) \\ \Leftrightarrow\; \exists \operatorname{mnbqc\_} \in \operatorname{expr\text{-}set}.\; (\operatorname{enc\text{-}in}(\operatorname{ebvar}(0), \operatorname{mnbqc\_}) \in \operatorname{conj\text{-}reach}(\operatorname{mnbb\_})) \end{array}
\end{equation*}

\par\medskip\noindent\textbf{Definition} (\texttt{qnode-ok}).\;For every $\operatorname{mnbe\_}$, 
\begin{equation*}
\begin{array}{@{}l@{}} \operatorname{qnode\text{-}ok}(\operatorname{mnbe\_}) \\ \Leftrightarrow\; \forall \operatorname{mnbqb\_} \in \operatorname{expr\text{-}set}. \\ \quad (\operatorname{mnbe\_} = \operatorname{enc\text{-}all}(\operatorname{mnbqb\_})) \Rightarrow \operatorname{all\text{-}confined}(\operatorname{mnbqb\_}) \wedge \\ \quad (\operatorname{mnbe\_} = \operatorname{enc\text{-}ex}(\operatorname{mnbqb\_})) \Rightarrow \operatorname{ex\text{-}confined}(\operatorname{mnbqb\_}) \wedge \\ \quad (\operatorname{mnbe\_} = \operatorname{enc\text{-}iota}(\operatorname{mnbqb\_})) \Rightarrow \operatorname{iota\text{-}confined}(\operatorname{mnbqb\_}) \end{array}
\end{equation*}

\par\medskip\noindent\textbf{Definition} (\texttt{cr-iota-def}).\;For all $\operatorname{mct\_}$, $\operatorname{mchs\_}$ and $\operatorname{mcc\_}$, 
\begin{equation*}
\begin{array}{@{}l@{}} \operatorname{cr\text{-}iota\text{-}def}(\operatorname{mct\_}, \operatorname{mchs\_}, \operatorname{mcc\_}) \\ \Leftrightarrow\; \exists \operatorname{mcg\_}, \operatorname{mcf\_}, \operatorname{mcp\_}. \\ \quad \operatorname{mcc\_} = {[}\operatorname{mcg\_}, \operatorname{mcf\_}] \wedge \\ \quad \operatorname{mchs\_} \\ \quad =\; {[}{[}\operatorname{mcg\_}, \operatorname{enc\text{-}ex}(\operatorname{enc\text{-}and}(\operatorname{mcp\_}, \operatorname{enc\text{-}all}(\operatorname{enc\text{-}imp}(\operatorname{mcp\_}, \operatorname{enc\text{-}eq}(\operatorname{ebvar}(1), \operatorname{ebvar}(0))))))], {[}\operatorname{ctx\text{-}add}(\operatorname{mcg\_}, \operatorname{expr\text{-}open}(\operatorname{mcp\_}, 0, \operatorname{enc\text{-}iota}(\operatorname{mcp\_}))), \operatorname{mcf\_}]] \end{array}
\end{equation*}

\par\medskip\noindent\textbf{Definition} (\texttt{cr-iota-in-elim}).\;For all $\operatorname{mct\_}$, $\operatorname{mchs\_}$ and $\operatorname{mcc\_}$, 
\begin{equation*}
\begin{array}{@{}l@{}} \operatorname{cr\text{-}iota\text{-}in\text{-}elim}(\operatorname{mct\_}, \operatorname{mchs\_}, \operatorname{mcc\_}) \\ \Leftrightarrow\; \exists \operatorname{mcg\_}, \operatorname{mcf\_}, \operatorname{mcp\_}. \\ \quad \operatorname{mcc\_} = {[}\operatorname{mcg\_}, \operatorname{mcf\_}] \wedge \\ \quad \operatorname{mchs\_} \\ \quad =\; {[}{[}\operatorname{mcg\_}, \operatorname{enc\text{-}eq}(\operatorname{enc\text{-}iota}(\operatorname{mcp\_}), \operatorname{enc\text{-}iota}(\operatorname{mcp\_}))], {[}\operatorname{ctx\text{-}add}(\operatorname{mcg\_}, \operatorname{expr\text{-}open}(\operatorname{mcp\_}, 0, \operatorname{enc\text{-}iota}(\operatorname{mcp\_}))), \operatorname{mcf\_}]] \end{array}
\end{equation*}

\par\medskip\noindent\textbf{Definition} (\texttt{cr-def-total}).\;For all $\operatorname{mct\_}$, $\operatorname{mchs\_}$ and $\operatorname{mcc\_}$, 
\begin{equation*}
\begin{array}{@{}l@{}} \operatorname{cr\text{-}def\text{-}total}(\operatorname{mct\_}, \operatorname{mchs\_}, \operatorname{mcc\_}) \\ \Leftrightarrow\; \exists \operatorname{mcg\_}, \operatorname{mchd\_}, \operatorname{mcl\_}. \\ \quad \operatorname{mcc\_} = {[}\operatorname{mcg\_}, \operatorname{enc\text{-}eq}(\operatorname{eapp}(\operatorname{mchd\_}, \operatorname{mcl\_}), \operatorname{eapp}(\operatorname{mchd\_}, \operatorname{mcl\_}))] \wedge \\ \quad \operatorname{total\text{-}head}(\operatorname{mchd\_}) \wedge \\ \quad \operatorname{length}(\operatorname{mchs\_}) = \operatorname{length}(\operatorname{mcl\_}) \wedge \\ \quad \forall \operatorname{mci\_} \in \mathbb{N}. \\ \quad \quad (1 \leq \operatorname{mci\_}) \Rightarrow  \\ \quad \quad (\operatorname{mci\_} \leq \operatorname{length}(\operatorname{mcl\_})) \Rightarrow  \\ \quad \quad \quad \operatorname{nth}(\operatorname{mci\_}, \operatorname{mchs\_}) = {[}\operatorname{mcg\_}, \operatorname{enc\text{-}eq}(\operatorname{nth}(\operatorname{mci\_}, \operatorname{mcl\_}), \operatorname{nth}(\operatorname{mci\_}, \operatorname{mcl\_}))] \end{array}
\end{equation*}

\par\medskip\noindent\textbf{Definition} (\texttt{total-head}).\;For every $\operatorname{mchd\_}$, 
\begin{equation*}
\operatorname{total\text{-}head}(\operatorname{mchd\_}) \Leftrightarrow (((11 \leq \operatorname{mchd\_}) \wedge (\operatorname{mchd\_} \leq 14)) \vee ((\operatorname{mchd\_} = 19) \vee ((30 \leq \operatorname{mchd\_}) \wedge (\operatorname{mchd\_} \leq 31))))
\end{equation*}

\par\medskip\noindent\textbf{Definition} (\texttt{cr-def-strict-in-l}).\;For all $\operatorname{mct\_}$, $\operatorname{mchs\_}$ and $\operatorname{mcc\_}$, 
\begin{equation*}
\begin{array}{@{}l@{}} \operatorname{cr\text{-}def\text{-}strict\text{-}in\text{-}l}(\operatorname{mct\_}, \operatorname{mchs\_}, \operatorname{mcc\_}) \\ \Leftrightarrow\; \exists \operatorname{mcg\_}, \operatorname{mcw\_}, \operatorname{mca\_}. \\ \quad (\operatorname{mcc\_} = {[}\operatorname{mcg\_}, \operatorname{enc\text{-}eq}(\operatorname{mcw\_}, \operatorname{mcw\_})]) \wedge (\operatorname{mchs\_} = {[}{[}\operatorname{mcg\_}, \operatorname{enc\text{-}in}(\operatorname{mcw\_}, \operatorname{mca\_})]]) \end{array}
\end{equation*}

\par\medskip\noindent\textbf{Definition} (\texttt{cr-def-strict-in-r}).\;For all $\operatorname{mct\_}$, $\operatorname{mchs\_}$ and $\operatorname{mcc\_}$, 
\begin{equation*}
\begin{array}{@{}l@{}} \operatorname{cr\text{-}def\text{-}strict\text{-}in\text{-}r}(\operatorname{mct\_}, \operatorname{mchs\_}, \operatorname{mcc\_}) \\ \Leftrightarrow\; \exists \operatorname{mcg\_}, \operatorname{mcw\_}, \operatorname{mca\_}. \\ \quad (\operatorname{mcc\_} = {[}\operatorname{mcg\_}, \operatorname{enc\text{-}eq}(\operatorname{mca\_}, \operatorname{mca\_})]) \wedge (\operatorname{mchs\_} = {[}{[}\operatorname{mcg\_}, \operatorname{enc\text{-}in}(\operatorname{mcw\_}, \operatorname{mca\_})]]) \end{array}
\end{equation*}

\par\medskip\noindent\textbf{Definition} (\texttt{cr-def-list}).\;For all $\operatorname{mct\_}$, $\operatorname{mchs\_}$ and $\operatorname{mcc\_}$, 
\begin{equation*}
\begin{array}{@{}l@{}} \operatorname{cr\text{-}def\text{-}list}(\operatorname{mct\_}, \operatorname{mchs\_}, \operatorname{mcc\_}) \\ \Leftrightarrow\; \exists \operatorname{mcg\_}, \operatorname{mcl\_}. \\ \quad \operatorname{mcc\_} = {[}\operatorname{mcg\_}, \operatorname{enc\text{-}eq}(\operatorname{enc\text{-}list}(\operatorname{mcl\_}), \operatorname{enc\text{-}list}(\operatorname{mcl\_}))] \wedge \\ \quad \operatorname{length}(\operatorname{mchs\_}) = \operatorname{length}(\operatorname{mcl\_}) \wedge \\ \quad \forall \operatorname{mci\_} \in \mathbb{N}. \\ \quad \quad (1 \leq \operatorname{mci\_}) \Rightarrow  \\ \quad \quad (\operatorname{mci\_} \leq \operatorname{length}(\operatorname{mcl\_})) \Rightarrow  \\ \quad \quad \quad \operatorname{nth}(\operatorname{mci\_}, \operatorname{mchs\_}) = {[}\operatorname{mcg\_}, \operatorname{enc\text{-}in}(\operatorname{nth}(\operatorname{mci\_}, \operatorname{mcl\_}), \operatorname{enc\text{-}setc})] \end{array}
\end{equation*}

\par\medskip\noindent\textbf{Definition} (\texttt{cr-def-strict-arg}).\;For all $\operatorname{mct\_}$, $\operatorname{mchs\_}$ and $\operatorname{mcc\_}$, 
\begin{equation*}
\begin{array}{@{}l@{}} \operatorname{cr\text{-}def\text{-}strict\text{-}arg}(\operatorname{mct\_}, \operatorname{mchs\_}, \operatorname{mcc\_}) \\ \Leftrightarrow\; \exists \operatorname{mcg\_}, \operatorname{mchd\_}, \operatorname{mcl\_}, \operatorname{mci\_}. \\ \quad \operatorname{mcc\_} = {[}\operatorname{mcg\_}, \operatorname{enc\text{-}eq}(\operatorname{nth}(\operatorname{mci\_}, \operatorname{mcl\_}), \operatorname{nth}(\operatorname{mci\_}, \operatorname{mcl\_}))] \wedge \\ \quad \operatorname{strict\text{-}head}(\operatorname{mchd\_}) \wedge \\ \quad \operatorname{mci\_} \in \mathbb{N} \wedge \\ \quad 1 \leq \operatorname{mci\_} \wedge \\ \quad \operatorname{mci\_} \leq \operatorname{length}(\operatorname{mcl\_}) \wedge \\ \quad \operatorname{mchs\_} = {[}{[}\operatorname{mcg\_}, \operatorname{enc\text{-}eq}(\operatorname{eapp}(\operatorname{mchd\_}, \operatorname{mcl\_}), \operatorname{eapp}(\operatorname{mchd\_}, \operatorname{mcl\_}))]] \end{array}
\end{equation*}

\par\medskip\noindent\textbf{Definition} (\texttt{cr-def-strict-dom}).\;For all $\operatorname{mct\_}$, $\operatorname{mchs\_}$ and $\operatorname{mcc\_}$, 
\begin{equation*}
\begin{array}{@{}l@{}} \operatorname{cr\text{-}def\text{-}strict\text{-}dom}(\operatorname{mct\_}, \operatorname{mchs\_}, \operatorname{mcc\_}) \\ \Leftrightarrow\; \exists \operatorname{mcg\_}, \operatorname{mchd\_}, \operatorname{mca\_}, \operatorname{mcb\_}. \\ \quad \operatorname{mcc\_} = {[}\operatorname{mcg\_}, \operatorname{enc\text{-}eq}(\operatorname{mca\_}, \operatorname{mca\_})] \wedge \\ \quad (\operatorname{mchd\_} = 22) \vee ((\operatorname{mchd\_} = 24) \vee (\operatorname{mchd\_} = 26)) \wedge \\ \quad \operatorname{mchs\_} = {[}{[}\operatorname{mcg\_}, \operatorname{enc\text{-}eq}(\operatorname{ebind}(\operatorname{mchd\_}, 1, {[}\operatorname{mca\_}], \operatorname{mcb\_}), \operatorname{ebind}(\operatorname{mchd\_}, 1, {[}\operatorname{mca\_}], \operatorname{mcb\_}))]] \end{array}
\end{equation*}

\par\medskip\noindent\textbf{Definition} (\texttt{cr-def-strict-pred}).\;For all $\operatorname{mct\_}$, $\operatorname{mchs\_}$ and $\operatorname{mcc\_}$, 
\begin{equation*}
\begin{array}{@{}l@{}} \operatorname{cr\text{-}def\text{-}strict\text{-}pred}(\operatorname{mct\_}, \operatorname{mchs\_}, \operatorname{mcc\_}) \\ \Leftrightarrow\; \exists \operatorname{mcg\_}, \operatorname{mcl\_}, \operatorname{mck\_}, \operatorname{mci\_}. \\ \quad \operatorname{mcc\_} = {[}\operatorname{mcg\_}, \operatorname{enc\text{-}eq}(\operatorname{nth}(\operatorname{mci\_}, \operatorname{mcl\_}), \operatorname{nth}(\operatorname{mci\_}, \operatorname{mcl\_}))] \wedge \\ \quad \operatorname{nth}(1, \operatorname{mcl\_}) = \operatorname{econst}(\operatorname{mck\_}) \wedge \\ \quad \operatorname{pred\text{-}code}(\operatorname{mck\_}) \wedge \\ \quad \operatorname{mci\_} \in \mathbb{N} \wedge \\ \quad 2 \leq \operatorname{mci\_} \wedge \\ \quad \operatorname{mci\_} \leq \operatorname{length}(\operatorname{mcl\_}) \wedge \\ \quad \operatorname{mchs\_} = {[}{[}\operatorname{mcg\_}, \operatorname{eapp}(10, \operatorname{mcl\_})]] \end{array}
\end{equation*}

\par\medskip\noindent\textbf{Definition} (\texttt{cr-iota-cr}).\;For all $\operatorname{mct\_}$, $\operatorname{mchs\_}$ and $\operatorname{mcc\_}$, 
\begin{equation*}
\begin{array}{@{}l@{}} \operatorname{cr\text{-}iota\text{-}cr}(\operatorname{mct\_}, \operatorname{mchs\_}, \operatorname{mcc\_}) \\ \Leftrightarrow\; \exists \operatorname{mcg\_}, \operatorname{mcp\_}, \operatorname{mcq\_}, \operatorname{mcy\_}. \\ \quad \operatorname{mcc\_} = {[}\operatorname{mcg\_}, \operatorname{enc\text{-}qeq}(\operatorname{enc\text{-}iota}(\operatorname{mcp\_}), \operatorname{enc\text{-}iota}(\operatorname{mcq\_}))] \wedge \\ \quad \operatorname{mcy\_} \in \mathbb{N} \wedge \\ \quad \neg (\operatorname{mcy\_} \in (\operatorname{ctx\text{-}fv}(\operatorname{mcg\_}) \cup \operatorname{expr\text{-}fv}(\operatorname{enc\text{-}qeq}(\operatorname{enc\text{-}iota}(\operatorname{mcp\_}), \operatorname{enc\text{-}iota}(\operatorname{mcq\_}))))) \wedge \\ \quad \operatorname{mchs\_} = {[}{[}\operatorname{mcg\_}, \operatorname{enc\text{-}iff}(\operatorname{expr\text{-}open}(\operatorname{mcp\_}, 0, \operatorname{efvar}(\operatorname{mcy\_})), \operatorname{expr\text{-}open}(\operatorname{mcq\_}, 0, \operatorname{efvar}(\operatorname{mcy\_})))]] \end{array}
\end{equation*}

\par\medskip\noindent\textbf{Definition} (\texttt{cr-tbind-cr}).\;For all $\operatorname{mct\_}$, $\operatorname{mchs\_}$ and $\operatorname{mcc\_}$, 
\begin{equation*}
\begin{array}{@{}l@{}} \operatorname{cr\text{-}tbind\text{-}cr}(\operatorname{mct\_}, \operatorname{mchs\_}, \operatorname{mcc\_}) \\ \Leftrightarrow\; \exists \operatorname{mcg\_}, \operatorname{mchd\_}, \operatorname{mcol\_}, \operatorname{mcp\_}, \operatorname{mcq\_}, \operatorname{mcy\_}, \operatorname{mcf\_}. \\ \quad \operatorname{mcc\_} = {[}\operatorname{mcg\_}, \operatorname{enc\text{-}qeq}(\operatorname{ebind}(\operatorname{mchd\_}, 1, \operatorname{mcol\_}, \operatorname{mcp\_}), \operatorname{ebind}(\operatorname{mchd\_}, 1, \operatorname{mcol\_}, \operatorname{mcq\_}))] \wedge \\ \quad \operatorname{mcy\_} \in \mathbb{N} \wedge \\ \quad \neg (\operatorname{mcy\_} \in (\operatorname{ctx\text{-}fv}(\operatorname{mcg\_}) \cup \operatorname{expr\text{-}fv}(\operatorname{enc\text{-}qeq}(\operatorname{ebind}(\operatorname{mchd\_}, 1, \operatorname{mcol\_}, \operatorname{mcp\_}), \operatorname{ebind}(\operatorname{mchd\_}, 1, \operatorname{mcol\_}, \operatorname{mcq\_}))))) \wedge \\ \quad (\operatorname{mchd\_} = 22) \vee (\operatorname{mchd\_} = 23) \wedge \\ \quad \operatorname{mcf\_} = \operatorname{enc\text{-}iff}(\operatorname{expr\text{-}open}(\operatorname{mcp\_}, 0, \operatorname{efvar}(\operatorname{mcy\_})), \operatorname{expr\text{-}open}(\operatorname{mcq\_}, 0, \operatorname{efvar}(\operatorname{mcy\_}))) \vee \\ \quad (\operatorname{mchd\_} = 24) \vee (\operatorname{mchd\_} = 26) \wedge \\ \quad \operatorname{mcf\_} = \operatorname{enc\text{-}qeq}(\operatorname{expr\text{-}open}(\operatorname{mcp\_}, 0, \operatorname{efvar}(\operatorname{mcy\_})), \operatorname{expr\text{-}open}(\operatorname{mcq\_}, 0, \operatorname{efvar}(\operatorname{mcy\_}))) \wedge \\ \quad \operatorname{mchs\_} = {[}{[}\operatorname{mcg\_}, \operatorname{mcf\_}]] \end{array}
\end{equation*}

\par\medskip\noindent\textbf{Definition} (\texttt{openl}).\;For all $\operatorname{mcob\_}$ and $\operatorname{mcol\_}$,  $(\operatorname{openl}(\operatorname{mcob\_}, \operatorname{mcol\_}, 0) \simeq \operatorname{mcob\_})$\par
Let $\operatorname{mcob\_}$ and $\operatorname{mcol\_}$ be arbitrary. Let $\operatorname{mcon\_} \in \mathbb{N}$. Then
\begin{equation*}
(\operatorname{openl}(\operatorname{mcob\_}, \operatorname{mcol\_}, \operatorname{succ}(\operatorname{mcon\_})) \simeq \operatorname{expr\text{-}open}(\operatorname{openl}(\operatorname{mcob\_}, \operatorname{mcol\_}, \operatorname{mcon\_}), \operatorname{mcon\_}, \operatorname{nth}(\operatorname{succ}(\operatorname{mcon\_}), \operatorname{mcol\_})))
\end{equation*}

\par\medskip\noindent\textbf{Definition} (\texttt{cr-tuple-eta}).\;For all $\operatorname{mct\_}$, $\operatorname{mchs\_}$ and $\operatorname{mcc\_}$, 
\begin{equation*}
\begin{array}{@{}l@{}} \operatorname{cr\text{-}tuple\text{-}eta}(\operatorname{mct\_}, \operatorname{mchs\_}, \operatorname{mcc\_}) \\ \Leftrightarrow\; \exists \operatorname{mcg\_}, \operatorname{mcx\_}, \operatorname{mca\_}, \operatorname{mcl\_}. \\ \quad \operatorname{mcc\_} = {[}\operatorname{mcg\_}, \operatorname{enc\text{-}eq}(\operatorname{mcx\_}, \operatorname{enc\text{-}list}(\operatorname{mcl\_}))] \wedge \\ \quad 1 \leq \operatorname{length}(\operatorname{mca\_}) \wedge \\ \quad \operatorname{length}(\operatorname{mcl\_}) = \operatorname{length}(\operatorname{mca\_}) \wedge \\ \quad \forall \operatorname{mci\_} \in \mathbb{N}. \\ \quad \quad (1 \leq \operatorname{mci\_}) \Rightarrow  \\ \quad \quad (\operatorname{mci\_} \leq \operatorname{length}(\operatorname{mca\_})) \Rightarrow (\operatorname{nth}(\operatorname{mci\_}, \operatorname{mcl\_}) = \operatorname{enc\text{-}nth}(\operatorname{enc\text{-}num}(\operatorname{mci\_}), \operatorname{mcx\_})) \wedge \\ \quad \operatorname{mchs\_} = {[}{[}\operatorname{mcg\_}, \operatorname{enc\text{-}in}(\operatorname{mcx\_}, \operatorname{enc\text{-}cart}(\operatorname{mca\_}))]] \end{array}
\end{equation*}

\par\medskip\noindent\textbf{Definition} (\texttt{hyps-in}).\;For all $\operatorname{mchs\_}$ and $\operatorname{mcs\_}$, 
\begin{equation*}
\begin{array}{@{}l@{}} \operatorname{hyps\text{-}in}(\operatorname{mchs\_}, \operatorname{mcs\_}) \\ \Leftrightarrow\; \operatorname{mchs\_} \in \mathbf{Set} \wedge \\ \forall \operatorname{mci\_} \in \mathbb{N}. \\ \quad (1 \leq \operatorname{mci\_}) \Rightarrow ((\operatorname{mci\_} \leq \operatorname{length}(\operatorname{mchs\_})) \Rightarrow (\operatorname{nth}(\operatorname{mci\_}, \operatorname{mchs\_}) \in \operatorname{mcs\_})) \end{array}
\end{equation*}

\par\medskip\noindent\textbf{Definition} (\texttt{rel-detach}).\;For all $\operatorname{mct\_}$, $\operatorname{mchs\_}$ and $\operatorname{mcc\_}$, 
\begin{equation*}
\begin{array}{@{}l@{}} \operatorname{rel\text{-}detach}(\operatorname{mct\_}, \operatorname{mchs\_}, \operatorname{mcc\_}) \\ \Leftrightarrow\; \operatorname{mcc\_} \in \operatorname{seq\text{-}set} \wedge \\ \exists \operatorname{mcg\_}, \operatorname{mcf\_}, \operatorname{mcp\_}, \operatorname{mcq\_}. \\ \quad \operatorname{mcc\_} = {[}\operatorname{mcg\_}, \operatorname{mcf\_}] \wedge \\ \quad \operatorname{enc\text{-}imp}(\operatorname{mcp\_}, \operatorname{mcq\_}) \in \operatorname{mcg\_} \wedge \\ \quad \operatorname{mcp\_} \in \operatorname{mcg\_} \wedge \\ \quad \operatorname{mchs\_} = {[}{[}\operatorname{ctx\text{-}add}(\operatorname{mcg\_}, \operatorname{mcq\_}), \operatorname{mcf\_}]] \end{array}
\end{equation*}

\par\medskip\noindent\textbf{Definition} (\texttt{rel-union-intro}).\;For all $\operatorname{mct\_}$, $\operatorname{mchs\_}$ and $\operatorname{mcc\_}$, 
\begin{equation*}
\begin{array}{@{}l@{}} \operatorname{rel\text{-}union\text{-}intro}(\operatorname{mct\_}, \operatorname{mchs\_}, \operatorname{mcc\_}) \\ \Leftrightarrow\; \operatorname{mcc\_} \in \operatorname{seq\text{-}set} \wedge \\ \exists \operatorname{mcg\_}, \operatorname{mcx\_}, \operatorname{mcl\_}, \operatorname{mck\_}. \\ \quad \operatorname{mcc\_} = {[}\operatorname{mcg\_}, \operatorname{enc\text{-}in}(\operatorname{mcx\_}, \operatorname{enc\text{-}union}(\operatorname{mcl\_}))] \wedge \\ \quad \operatorname{mck\_} \in \mathbb{N} \wedge \\ \quad 1 \leq \operatorname{mck\_} \wedge \\ \quad \operatorname{mck\_} \leq \operatorname{length}(\operatorname{mcl\_}) \wedge \\ \quad \operatorname{length}(\operatorname{mchs\_}) = \operatorname{succ}(\operatorname{length}(\operatorname{mcl\_})) \wedge \\ \quad \operatorname{nth}(1, \operatorname{mchs\_}) = {[}\operatorname{mcg\_}, \operatorname{enc\text{-}in}(\operatorname{mcx\_}, \operatorname{nth}(\operatorname{mck\_}, \operatorname{mcl\_}))] \wedge \\ \quad \forall \operatorname{mci\_} \in \mathbb{N}. \\ \quad \quad (1 \leq \operatorname{mci\_}) \Rightarrow  \\ \quad \quad (\operatorname{mci\_} \leq \operatorname{length}(\operatorname{mcl\_})) \Rightarrow  \\ \quad \quad \quad \operatorname{nth}(\operatorname{succ}(\operatorname{mci\_}), \operatorname{mchs\_}) = {[}\operatorname{mcg\_}, \operatorname{enc\text{-}eq}(\operatorname{nth}(\operatorname{mci\_}, \operatorname{mcl\_}), \operatorname{nth}(\operatorname{mci\_}, \operatorname{mcl\_}))] \end{array}
\end{equation*}

\par\medskip\noindent\textbf{Definition} (\texttt{rel-union-elim}).\;For all $\operatorname{mct\_}$, $\operatorname{mchs\_}$ and $\operatorname{mcc\_}$, 
\begin{equation*}
\begin{array}{@{}l@{}} \operatorname{rel\text{-}union\text{-}elim}(\operatorname{mct\_}, \operatorname{mchs\_}, \operatorname{mcc\_}) \\ \Leftrightarrow\; \operatorname{mcc\_} \in \operatorname{seq\text{-}set} \wedge \\ \exists \operatorname{mcg\_}, \operatorname{mcf\_}, \operatorname{mcx\_}, \operatorname{mcl\_}. \\ \quad \operatorname{mcc\_} = {[}\operatorname{mcg\_}, \operatorname{mcf\_}] \wedge \\ \quad \operatorname{enc\text{-}in}(\operatorname{mcx\_}, \operatorname{enc\text{-}union}(\operatorname{mcl\_})) \in \operatorname{mcg\_} \wedge \\ \quad \operatorname{length}(\operatorname{mchs\_}) = \operatorname{length}(\operatorname{mcl\_}) \wedge \\ \quad \forall \operatorname{mci\_} \in \mathbb{N}. \\ \quad \quad (1 \leq \operatorname{mci\_}) \Rightarrow  \\ \quad \quad (\operatorname{mci\_} \leq \operatorname{length}(\operatorname{mcl\_})) \Rightarrow  \\ \quad \quad \quad \operatorname{nth}(\operatorname{mci\_}, \operatorname{mchs\_}) \\ \quad \quad \quad =\; {[}\operatorname{ctx\text{-}add}(\operatorname{difference}(\operatorname{mcg\_}, \{\operatorname{enc\text{-}in}(\operatorname{mcx\_}, \operatorname{enc\text{-}union}(\operatorname{mcl\_}))\}), \operatorname{enc\text{-}in}(\operatorname{mcx\_}, \operatorname{nth}(\operatorname{mci\_}, \operatorname{mcl\_}))), \operatorname{mcf\_}] \end{array}
\end{equation*}

\par\medskip\noindent\textbf{Definition} (\texttt{rel-intersection-intro}).\;For all $\operatorname{mct\_}$, $\operatorname{mchs\_}$ and $\operatorname{mcc\_}$, 
\begin{equation*}
\begin{array}{@{}l@{}} \operatorname{rel\text{-}intersection\text{-}intro}(\operatorname{mct\_}, \operatorname{mchs\_}, \operatorname{mcc\_}) \\ \Leftrightarrow\; \operatorname{mcc\_} \in \operatorname{seq\text{-}set} \wedge \\ \exists \operatorname{mcg\_}, \operatorname{mcx\_}, \operatorname{mcl\_}. \\ \quad \operatorname{mcc\_} = {[}\operatorname{mcg\_}, \operatorname{enc\text{-}in}(\operatorname{mcx\_}, \operatorname{enc\text{-}inter}(\operatorname{mcl\_}))] \wedge \\ \quad 1 \leq \operatorname{length}(\operatorname{mcl\_}) \wedge \\ \quad \operatorname{length}(\operatorname{mchs\_}) = \operatorname{length}(\operatorname{mcl\_}) \wedge \\ \quad \forall \operatorname{mci\_} \in \mathbb{N}. \\ \quad \quad (1 \leq \operatorname{mci\_}) \Rightarrow  \\ \quad \quad (\operatorname{mci\_} \leq \operatorname{length}(\operatorname{mcl\_})) \Rightarrow  \\ \quad \quad \quad \operatorname{nth}(\operatorname{mci\_}, \operatorname{mchs\_}) = {[}\operatorname{mcg\_}, \operatorname{enc\text{-}in}(\operatorname{mcx\_}, \operatorname{nth}(\operatorname{mci\_}, \operatorname{mcl\_}))] \end{array}
\end{equation*}

\par\medskip\noindent\textbf{Definition} (\texttt{rel-intersection-elim}).\;For all $\operatorname{mct\_}$, $\operatorname{mchs\_}$ and $\operatorname{mcc\_}$, 
\begin{equation*}
\begin{array}{@{}l@{}} \operatorname{rel\text{-}intersection\text{-}elim}(\operatorname{mct\_}, \operatorname{mchs\_}, \operatorname{mcc\_}) \\ \Leftrightarrow\; \operatorname{mcc\_} \in \operatorname{seq\text{-}set} \wedge \\ \exists \operatorname{mcg\_}, \operatorname{mcf\_}, \operatorname{mcx\_}, \operatorname{mcl\_}, \operatorname{mck\_}. \\ \quad \operatorname{mcc\_} = {[}\operatorname{mcg\_}, \operatorname{mcf\_}] \wedge \\ \quad \operatorname{enc\text{-}in}(\operatorname{mcx\_}, \operatorname{enc\text{-}inter}(\operatorname{mcl\_})) \in \operatorname{mcg\_} \wedge \\ \quad \operatorname{mck\_} \in \mathbb{N} \wedge \\ \quad 1 \leq \operatorname{mck\_} \wedge \\ \quad \operatorname{mck\_} \leq \operatorname{length}(\operatorname{mcl\_}) \wedge \\ \quad \operatorname{mchs\_} = {[}{[}\operatorname{ctx\text{-}add}(\operatorname{mcg\_}, \operatorname{enc\text{-}in}(\operatorname{mcx\_}, \operatorname{nth}(\operatorname{mck\_}, \operatorname{mcl\_}))), \operatorname{mcf\_}]] \end{array}
\end{equation*}

\par\medskip\noindent\textbf{Definition} (\texttt{rel-tuples-intro-0}).\;For all $\operatorname{mct\_}$, $\operatorname{mchs\_}$ and $\operatorname{mcc\_}$, 
\begin{equation*}
\begin{array}{@{}l@{}} \operatorname{rel\text{-}tuples\text{-}intro\text{-}0}(\operatorname{mct\_}, \operatorname{mchs\_}, \operatorname{mcc\_}) \\ \Leftrightarrow\; \operatorname{mcc\_} \in \operatorname{seq\text{-}set} \wedge \\ \exists \operatorname{mcg\_}, \operatorname{mcl\_}, \operatorname{mca\_}. \\ \quad \operatorname{mcc\_} = {[}\operatorname{mcg\_}, \operatorname{enc\text{-}in}(\operatorname{enc\text{-}list}(\operatorname{mcl\_}), \operatorname{enc\text{-}tuples}(\operatorname{mca\_}))] \wedge \\ \quad \operatorname{length}(\operatorname{mcl\_}) = 0 \wedge \\ \quad \operatorname{mchs\_} = {[}{[}\operatorname{mcg\_}, \operatorname{enc\text{-}eq}(\operatorname{mca\_}, \operatorname{mca\_})]] \end{array}
\end{equation*}

\par\medskip\noindent\textbf{Definition} (\texttt{rel-tuples-intro}).\;For all $\operatorname{mct\_}$, $\operatorname{mchs\_}$ and $\operatorname{mcc\_}$, 
\begin{equation*}
\begin{array}{@{}l@{}} \operatorname{rel\text{-}tuples\text{-}intro}(\operatorname{mct\_}, \operatorname{mchs\_}, \operatorname{mcc\_}) \\ \Leftrightarrow\; \operatorname{mcc\_} \in \operatorname{seq\text{-}set} \wedge \\ \exists \operatorname{mcg\_}, \operatorname{mcl\_}, \operatorname{mca\_}. \\ \quad \operatorname{mcc\_} = {[}\operatorname{mcg\_}, \operatorname{enc\text{-}in}(\operatorname{enc\text{-}list}(\operatorname{mcl\_}), \operatorname{enc\text{-}tuples}(\operatorname{mca\_}))] \wedge \\ \quad 1 \leq \operatorname{length}(\operatorname{mcl\_}) \wedge \\ \quad \operatorname{length}(\operatorname{mchs\_}) = \operatorname{length}(\operatorname{mcl\_}) \wedge \\ \quad \forall \operatorname{mci\_} \in \mathbb{N}. \\ \quad \quad (1 \leq \operatorname{mci\_}) \Rightarrow  \\ \quad \quad (\operatorname{mci\_} \leq \operatorname{length}(\operatorname{mcl\_})) \Rightarrow  \\ \quad \quad \quad \operatorname{nth}(\operatorname{mci\_}, \operatorname{mchs\_}) = {[}\operatorname{mcg\_}, \operatorname{enc\text{-}in}(\operatorname{nth}(\operatorname{mci\_}, \operatorname{mcl\_}), \operatorname{mca\_})] \end{array}
\end{equation*}

\par\medskip\noindent\textbf{Definition} (\texttt{rel-tuples-elim}).\;For all $\operatorname{mct\_}$, $\operatorname{mchs\_}$ and $\operatorname{mcc\_}$, 
\begin{equation*}
\begin{array}{@{}l@{}} \operatorname{rel\text{-}tuples\text{-}elim}(\operatorname{mct\_}, \operatorname{mchs\_}, \operatorname{mcc\_}) \\ \Leftrightarrow\; \operatorname{mcc\_} \in \operatorname{seq\text{-}set} \wedge \\ \exists \operatorname{mcg\_}, \operatorname{mcf\_}, \operatorname{mcl\_}, \operatorname{mca\_}, \operatorname{mck\_}. \\ \quad \operatorname{mcc\_} = {[}\operatorname{mcg\_}, \operatorname{mcf\_}] \wedge \\ \quad \operatorname{enc\text{-}in}(\operatorname{enc\text{-}list}(\operatorname{mcl\_}), \operatorname{enc\text{-}tuples}(\operatorname{mca\_})) \in \operatorname{mcg\_} \wedge \\ \quad \operatorname{mck\_} \in \mathbb{N} \wedge \\ \quad 1 \leq \operatorname{mck\_} \wedge \\ \quad \operatorname{mck\_} \leq \operatorname{length}(\operatorname{mcl\_}) \wedge \\ \quad \operatorname{mchs\_} = {[}{[}\operatorname{ctx\text{-}add}(\operatorname{mcg\_}, \operatorname{enc\text{-}in}(\operatorname{nth}(\operatorname{mck\_}, \operatorname{mcl\_}), \operatorname{mca\_})), \operatorname{mcf\_}]] \end{array}
\end{equation*}


\vnbbookfile{theorem-library/meta-core-laws}{theorem-library/meta-core-laws}
\noindent{\small meta-core-laws.scm -- THE CORE OF K2: the laws of well-formedness, sequents, DERIVABLE, and the admissibility of eight kernel rules (structure-library/meta-core.scm): detach, contraposition, union-intro (section H, phase 2; their proofs rewritten as dk-derive! trees by META-KIT), and union-elim, intersection-intro, ...}\par
\vnbprop{enc-not-unfold}{\vnbstmt{\forall p.\; (\operatorname{enc\text{-}not}(p) \simeq \operatorname{eapp}(2, {[}p]))}}
\vnbprop{enc-and-unfold}{\vnbstmt{\forall p, q.\; (\operatorname{enc\text{-}and}(p, q) \simeq \operatorname{eapp}(3, {[}p, q]))}}
\vnbprop{enc-or-unfold}{\vnbstmt{\forall p, q.\; (\operatorname{enc\text{-}or}(p, q) \simeq \operatorname{eapp}(4, {[}p, q]))}}
\vnbprop{enc-imp-unfold}{\vnbstmt{\forall p, q.\; (\operatorname{enc\text{-}imp}(p, q) \simeq \operatorname{eapp}(5, {[}p, q]))}}
\vnbprop{enc-iff-unfold}{\vnbstmt{\forall p, q.\; (\operatorname{enc\text{-}iff}(p, q) \simeq \operatorname{eapp}(6, {[}p, q]))}}
\vnbprop{enc-in-unfold}{\vnbstmt{\forall p, q.\; (\operatorname{enc\text{-}in}(p, q) \simeq \operatorname{eapp}(7, {[}p, q]))}}
\vnbprop{enc-eq-unfold}{\vnbstmt{\forall p, q.\; (\operatorname{enc\text{-}eq}(p, q) \simeq \operatorname{eapp}(8, {[}p, q]))}}
\vnbprop{enc-qeq-unfold}{\vnbstmt{\forall p, q.\; (\operatorname{enc\text{-}qeq}(p, q) \simeq \operatorname{eapp}(9, {[}p, q]))}}
\vnbprop{enc-union-unfold}{\vnbstmt{\forall l.\; (\operatorname{enc\text{-}union}(l) \simeq \operatorname{eapp}(11, l))}}
\vnbprop{ctx-add-unfold}{\vnbstmt{\forall g, a.\; (\operatorname{ctx\text{-}add}(g, a) \simeq (g \cup \{a\}))}}
\vnbprop{ctx-fv-unfold}{\vnbstmt{\forall g.\; (\operatorname{ctx\text{-}fv}(g) \simeq \operatorname{big\text{-}union}(w, g, \operatorname{expr\text{-}fv}(w)))}}
\vnbprop{derivable-unfold}{\vnbstmt{\begin{array}{@{}l@{}} \forall t. \\ \quad (\operatorname{derivable}(t) \simeq \{q \in \operatorname{seq\text{-}set} : \forall s \in \operatorname{power}(\operatorname{seq\text{-}set}).\; (\operatorname{core\text{-}closed}(t, s) \Rightarrow (q \in s))\}) \end{array}}}
\vnbprop{wks-stage-step}{\vnbstmt{\begin{array}{@{}l@{}} \forall k, e. \\ \quad ((k \in \mathbb{N}) \wedge (e \in \operatorname{expr\text{-}stage}(k))) \Rightarrow  \\ \quad \quad (\operatorname{wks\text{-}stage}((k + 1))(e) \simeq \operatorname{wks\text{-}stage}(k)(e)) \end{array}}}
\vnbprop{wks-stage-coh}{\vnbstmt{\begin{array}{@{}l@{}} \forall d \in \mathbb{N}, j, e. \\ \quad ((j \in \mathbb{N}) \wedge (e \in \operatorname{expr\text{-}stage}(j))) \Rightarrow  \\ \quad \quad (\operatorname{wks\text{-}stage}((j + d))(e) \simeq \operatorname{wks\text{-}stage}(j)(e)) \end{array}}}
\vnbprop{expr-wks-at}{\vnbstmt{\begin{array}{@{}l@{}} \forall j, e. \\ \quad ((j \in \mathbb{N}) \wedge (e \in \operatorname{expr\text{-}stage}(j))) \Rightarrow (\operatorname{expr\text{-}wks}(e) \simeq \operatorname{wks\text{-}stage}(j)(e)) \end{array}}}
\vnbprop{wks-efvar}{\vnbstmt{\forall n \in \mathbb{N}.\; (\operatorname{expr\text{-}wks}(\operatorname{efvar}(n)) \simeq \{1\})}}
\vnbprop{wks-ebvar}{\vnbstmt{\forall n \in \mathbb{N}.\; (\operatorname{expr\text{-}wks}(\operatorname{ebvar}(n)) \simeq \{1\})}}
\vnbprop{wks-econst}{\vnbstmt{\forall n \in \mathbb{N}.\; (\operatorname{expr\text{-}wks}(\operatorname{econst}(n)) \simeq \{1\})}}
\vnbprop{wks-eapp}{\vnbstmt{\begin{array}{@{}l@{}} \forall h, l. \\ \quad ((h \in \mathbb{N}) \wedge (l \in \operatorname{tuples}(\operatorname{expr\text{-}set}))) \Rightarrow  \\ \quad \quad (\operatorname{expr\text{-}wks}(\operatorname{eapp}(h, l)) \simeq \{k \in \mathbb{N} : ((k = \operatorname{hd\text{-}res}(h)) \wedge (\operatorname{app\text{-}ar\text{-}ok}(h, \operatorname{length}(l)) \wedge \forall w \in \operatorname{make\text{-}set}(l).\; (\operatorname{hd\text{-}argk}(h) \in \operatorname{expr\text{-}wks}(w))))\}) \end{array}}}
\vnbprop{wks-ebind}{\vnbstmt{\begin{array}{@{}l@{}} \forall h, c, a, b. \\ \quad ((h \in \mathbb{N}) \wedge ((c \in \mathbb{N}) \wedge ((a \in \operatorname{tuples}(\operatorname{expr\text{-}set})) \wedge (b \in \operatorname{expr\text{-}set})))) \Rightarrow  \\ \quad \quad (\operatorname{expr\text{-}wks}(\operatorname{ebind}(h, c, a, b)) \simeq \{k \in \mathbb{N} : ((k = \operatorname{hd\text{-}bres}(h)) \wedge (\operatorname{bind\text{-}sig}(h, c, \operatorname{length}(a)) \wedge (\forall w \in \operatorname{make\text{-}set}(a).\; (1 \in \operatorname{expr\text{-}wks}(w)) \wedge (\operatorname{hd\text{-}body}(h) \in \operatorname{expr\text{-}wks}(b)))))\}) \end{array}}}
\vnbprop{eapp-stage-inv}{\vnbstmt{\begin{array}{@{}l@{}} \forall n \in \mathbb{N}, h, l. \\ \quad (\operatorname{eapp}(h, l) \in \operatorname{expr\text{-}stage}(n)) \Rightarrow ((h \in \mathbb{N}) \wedge (l \in \operatorname{tuples}(\operatorname{expr\text{-}set}))) \end{array}}}
\vnbprop{eapp-expr-inv}{\vnbstmt{\forall h, l.\; ((\operatorname{eapp}(h, l) \in \operatorname{expr\text{-}set}) \Rightarrow ((h \in \mathbb{N}) \wedge (l \in \operatorname{tuples}(\operatorname{expr\text{-}set}))))}}
\vnbprop{ms2-intro}{\vnbstmt{\forall a, b.\; (((a \in \mathbf{Set}) \wedge (b \in \mathbf{Set})) \Rightarrow ((a \in \operatorname{make\text{-}set}({[}a, b])) \wedge (b \in \operatorname{make\text{-}set}({[}a, b]))))}}
\vnbprop{ms1-intro}{\vnbstmt{\forall a \in \mathbf{Set}.\; (a \in \operatorname{make\text{-}set}({[}a]))}}
\vnbprop{ms2-cases}{\vnbstmt{\forall a, b.; \forall w \in \operatorname{make\text{-}set}({[}a, b]).\; ((w = a) \vee (w = b))}}
\vnbprop{ms1-cases}{\vnbstmt{\forall a.; \forall w \in \operatorname{make\text{-}set}({[}a]).\; (w = a)}}
\vnbprop{clf-iff}{\vnbstmt{\forall e.\; ((e \in \mathit{clf}) \Leftrightarrow ((e \in \operatorname{expr\text{-}set}) \wedge (\operatorname{wf\text{-}fml}(e) \wedge \operatorname{expr\text{-}lc}(e))))}}
\vnbprop{clt-iff}{\vnbstmt{\forall e.\; ((e \in \mathit{clt}) \Leftrightarrow ((e \in \operatorname{expr\text{-}set}) \wedge (\operatorname{wf\text{-}term}(e) \wedge \operatorname{expr\text{-}lc}(e))))}}
\vnbprop{wks-eapp-iff}{\vnbstmt{\begin{array}{@{}l@{}} \forall h, l, k. \\ \quad ((h \in \mathbb{N}) \wedge ((l \in \operatorname{tuples}(\operatorname{expr\text{-}set})) \wedge (k \in \mathbb{N}))) \Rightarrow  \\ \quad \quad k \in \operatorname{expr\text{-}wks}(\operatorname{eapp}(h, l)) \\ \quad \quad \Leftrightarrow\; ((k = \operatorname{hd\text{-}res}(h)) \wedge (\operatorname{app\text{-}ar\text{-}ok}(h, \operatorname{length}(l)) \wedge \forall w \in \operatorname{make\text{-}set}(l).\; (\operatorname{hd\text{-}argk}(h) \in \operatorname{expr\text{-}wks}(w)))) \end{array}}}
\vnbprop{lc-eapp-iff}{\vnbstmt{\begin{array}{@{}l@{}} \forall h, l. \\ \quad ((h \in \mathbb{N}) \wedge (l \in \operatorname{tuples}(\operatorname{expr\text{-}set}))) \Rightarrow  \\ \quad \quad \operatorname{expr\text{-}lc}(\operatorname{eapp}(h, l)) \Leftrightarrow \forall w \in \operatorname{make\text{-}set}(l).\; \operatorname{expr\text{-}lc}(w) \end{array}}}
\vnbprop{clf-conn2}{\vnbstmt{\begin{array}{@{}l@{}} \forall h, a, b. \\ \quad ((h \in \mathbb{N}) \wedge ((3 \leq h) \wedge ((h \leq 6) \wedge ((a \in \operatorname{expr\text{-}set}) \wedge (b \in \operatorname{expr\text{-}set}))))) \Rightarrow  \\ \quad \quad (\operatorname{eapp}(h, {[}a, b]) \in \mathit{clf}) \Leftrightarrow ((a \in \mathit{clf}) \wedge (b \in \mathit{clf})) \end{array}}}
\vnbprop{clf-rel2}{\vnbstmt{\begin{array}{@{}l@{}} \forall h, a, b. \\ \quad ((h \in \mathbb{N}) \wedge ((7 \leq h) \wedge ((h \leq 9) \wedge ((a \in \operatorname{expr\text{-}set}) \wedge (b \in \operatorname{expr\text{-}set}))))) \Rightarrow  \\ \quad \quad (\operatorname{eapp}(h, {[}a, b]) \in \mathit{clf}) \Leftrightarrow ((a \in \mathit{clt}) \wedge (b \in \mathit{clt})) \end{array}}}
\vnbprop{clf-not1}{\vnbstmt{\forall a \in \operatorname{expr\text{-}set}.\; ((\operatorname{eapp}(2, {[}a]) \in \mathit{clf}) \Leftrightarrow (a \in \mathit{clf}))}}
\vnbprop{clt-union}{\vnbstmt{\begin{array}{@{}l@{}} \forall l \in \operatorname{tuples}(\operatorname{expr\text{-}set}). \\ \quad (\operatorname{eapp}(11, l) \in \mathit{clt}) \Leftrightarrow \forall u \in \operatorname{make\text{-}set}(l).\; (u \in \mathit{clt}) \end{array}}}
\vnbprop{clf-falsity0}{\vnbstmt{\operatorname{eapp}(1, {[}]) \in \mathit{clf}}}
\vnbprop{clf-truth0}{\vnbstmt{\operatorname{eapp}(0, {[}]) \in \mathit{clf}}}
\vnbprop{clf-imp-elim}{\vnbstmt{\forall p, q.\; ((\operatorname{enc\text{-}imp}(p, q) \in \mathit{clf}) \Rightarrow ((p \in \mathit{clf}) \wedge (q \in \mathit{clf})))}}
\vnbprop{clf-imp-intro}{\vnbstmt{\forall p, q.\; (((p \in \mathit{clf}) \wedge (q \in \mathit{clf})) \Rightarrow (\operatorname{enc\text{-}imp}(p, q) \in \mathit{clf}))}}
\vnbprop{clf-or-intro}{\vnbstmt{\forall p, q.\; (((p \in \mathit{clf}) \wedge (q \in \mathit{clf})) \Rightarrow (\operatorname{enc\text{-}or}(p, q) \in \mathit{clf}))}}
\vnbprop{clf-iff-intro}{\vnbstmt{\forall p, q.\; (((p \in \mathit{clf}) \wedge (q \in \mathit{clf})) \Rightarrow (\operatorname{enc\text{-}iff}(p, q) \in \mathit{clf}))}}
\vnbprop{clf-not-intro}{\vnbstmt{\forall p \in \mathit{clf}.\; (\operatorname{enc\text{-}not}(p) \in \mathit{clf})}}
\vnbprop{clf-in-elim}{\vnbstmt{\forall x, y.\; ((\operatorname{enc\text{-}in}(x, y) \in \mathit{clf}) \Rightarrow ((x \in \mathit{clt}) \wedge (y \in \mathit{clt})))}}
\vnbprop{clf-in-intro}{\vnbstmt{\forall x, y.\; (((x \in \mathit{clt}) \wedge (y \in \mathit{clt})) \Rightarrow (\operatorname{enc\text{-}in}(x, y) \in \mathit{clf}))}}
\vnbprop{clf-eq-c}{\vnbstmt{\forall x, y.\; (((x \in \mathit{clt}) \wedge (y \in \mathit{clt})) \Rightarrow (\operatorname{enc\text{-}eq}(x, y) \in \mathit{clf}))}}
\vnbprop{clt-union-elim}{\vnbstmt{\begin{array}{@{}l@{}} \forall l. \\ \quad (\operatorname{enc\text{-}union}(l) \in \mathit{clt}) \Rightarrow  \\ \quad \quad (l \in \operatorname{tuples}(\operatorname{expr\text{-}set})) \wedge \forall u \in \operatorname{make\text{-}set}(l).\; (u \in \mathit{clt}) \end{array}}}
\vnbprop{clf-falsity}{\vnbstmt{\operatorname{enc\text{-}falsity} \in \mathit{clf}}}
\vnbprop{clf-and-intro}{\vnbstmt{\forall p, q.\; (((p \in \mathit{clf}) \wedge (q \in \mathit{clf})) \Rightarrow (\operatorname{enc\text{-}and}(p, q) \in \mathit{clf}))}}
\vnbprop{clf-and-elim}{\vnbstmt{\forall p, q.\; ((\operatorname{enc\text{-}and}(p, q) \in \mathit{clf}) \Rightarrow ((p \in \mathit{clf}) \wedge (q \in \mathit{clf})))}}
\vnbprop{clf-or-elim}{\vnbstmt{\forall p, q.\; ((\operatorname{enc\text{-}or}(p, q) \in \mathit{clf}) \Rightarrow ((p \in \mathit{clf}) \wedge (q \in \mathit{clf})))}}
\vnbprop{clf-iff-elim}{\vnbstmt{\forall p, q.\; ((\operatorname{enc\text{-}iff}(p, q) \in \mathit{clf}) \Rightarrow ((p \in \mathit{clf}) \wedge (q \in \mathit{clf})))}}
\vnbprop{clf-not-elim}{\vnbstmt{\forall p.\; ((\operatorname{enc\text{-}not}(p) \in \mathit{clf}) \Rightarrow (p \in \mathit{clf}))}}
\vnbprop{clf-truth}{\vnbstmt{\operatorname{enc\text{-}truth} \in \mathit{clf}}}
\vnbprop{seq-iff}{\vnbstmt{\forall g, f.\; (({[}g, f] \in \operatorname{seq\text{-}set}) \Leftrightarrow ((g \in \operatorname{ctx\text{-}set}) \wedge (f \in \mathit{clf})))}}
\vnbprop{ctx-iff}{\vnbstmt{\forall g.\; ((g \in \operatorname{ctx\text{-}set}) \Leftrightarrow ((g \in \operatorname{power}(\mathit{clf})) \wedge (|g| \in \mathbb{N})))}}
\vnbprop{ctx-mem}{\vnbstmt{\forall g, a.\; (((g \in \operatorname{ctx\text{-}set}) \wedge (a \in g)) \Rightarrow (a \in \mathit{clf}))}}
\vnbprop{ctx-add-has}{\vnbstmt{\forall g.; \forall a \in \mathbf{Set}.\; (a \in \operatorname{ctx\text{-}add}(g, a))}}
\vnbprop{ctx-add-keeps}{\vnbstmt{\forall g, a, b.\; (((a \in \mathbf{Set}) \wedge (b \in g)) \Rightarrow (b \in \operatorname{ctx\text{-}add}(g, a)))}}
\vnbprop{ctx-add-ctx}{\vnbstmt{\forall g, a.\; (((g \in \operatorname{ctx\text{-}set}) \wedge (a \in \mathit{clf})) \Rightarrow (\operatorname{ctx\text{-}add}(g, a) \in \operatorname{ctx\text{-}set}))}}
\vnbprop{ctx-add-sub}{\vnbstmt{\forall g.; \forall a \in \mathbf{Set}.\; (g \subseteq \operatorname{ctx\text{-}add}(g, a))}}
\vnbprop{ctx-sub-cons}{\vnbstmt{\forall d, g.; \forall a \in \mathbf{Set}.\; (((d \subseteq g) \wedge (a \in g)) \Rightarrow (\operatorname{ctx\text{-}add}(d, a) \subseteq g))}}
\vnbprop{ctx-sub-snoc}{\vnbstmt{\forall d, g.; \forall b \in \mathbf{Set}.\; ((d \subseteq g) \Rightarrow (d \subseteq \operatorname{ctx\text{-}add}(g, b)))}}
\vnbprop{der-in-seq}{\vnbstmt{\forall t.; \forall s \in \operatorname{derivable}(t).\; (s \in \operatorname{seq\text{-}set})}}
\vnbprop{derivable-induction}{\vnbstmt{\begin{array}{@{}l@{}} \forall t, \mathit{mczs}, s. \\ \quad ((\mathit{mczs} \in \operatorname{power}(\operatorname{seq\text{-}set})) \wedge (\operatorname{core\text{-}closed}(t, \mathit{mczs}) \wedge (s \in \operatorname{derivable}(t)))) \Rightarrow  \\ \quad \quad s \in \mathit{mczs} \end{array}}}
\vnbprop{core-closed-apply}{\vnbstmt{\begin{array}{@{}l@{}} \forall t, s, a, c. \\ \quad (\operatorname{core\text{-}closed}(t, s) \wedge (\operatorname{core\text{-}rule}(t, a, c) \wedge (\operatorname{hyps\text{-}in}(a, s) \wedge (c \in \operatorname{seq\text{-}set})))) \Rightarrow  \\ \quad \quad c \in s \end{array}}}
\vnbprop{hyps-in-set}{\vnbstmt{\forall s, \mathit{mczs}.\; (\operatorname{hyps\text{-}in}(s, \mathit{mczs}) \Rightarrow (s \in \mathbf{Set}))}}
\vnbprop{hyps-in-all}{\vnbstmt{\begin{array}{@{}l@{}} \forall s, \mathit{mczs}.; \forall i \in \mathbb{N}. \\ \quad (\operatorname{hyps\text{-}in}(s, \mathit{mczs}) \wedge ((1 \leq i) \wedge (i \leq \operatorname{length}(s)))) \Rightarrow (\operatorname{nth}(i, s) \in \mathit{mczs}) \end{array}}}
\vnbprop{derivable-closed}{\vnbstmt{\begin{array}{@{}l@{}} \forall t, s, c. \\ \quad (\operatorname{core\text{-}rule}(t, s, c) \wedge (\operatorname{hyps\text{-}in}(s, \operatorname{derivable}(t)) \wedge (c \in \operatorname{seq\text{-}set}))) \Rightarrow  \\ \quad \quad c \in \operatorname{derivable}(t) \end{array}}}
\vnbprop{hyps-in-0}{\vnbstmt{\forall s.\; \operatorname{hyps\text{-}in}({[}], s)}}
\vnbprop{hyps-in-1}{\vnbstmt{\forall a, s.\; ((a \in s) \Rightarrow \operatorname{hyps\text{-}in}({[}a], s))}}
\vnbprop{hyps-in-2}{\vnbstmt{\forall a, b, s.\; (((a \in s) \wedge (b \in s)) \Rightarrow \operatorname{hyps\text{-}in}({[}a, b], s))}}
\vnbprop{hyps-in-3}{\vnbstmt{\forall a, b, c, s.\; (((a \in s) \wedge ((b \in s) \wedge (c \in s))) \Rightarrow \operatorname{hyps\text{-}in}({[}a, b, c], s))}}
\vnbprop{hyps-in-1-elim}{\vnbstmt{\forall a, s.\; ((({[}a] = {[}a]) \wedge \operatorname{hyps\text{-}in}({[}a], s)) \Rightarrow (a \in s))}}
\vnbprop{derivable-and-intro}{\vnbstmt{\begin{array}{@{}l@{}} \forall t, s, c. \\ \quad (\operatorname{cr\text{-}and\text{-}intro}(t, s, c) \wedge (\operatorname{hyps\text{-}in}(s, \operatorname{derivable}(t)) \wedge (c \in \operatorname{seq\text{-}set}))) \Rightarrow  \\ \quad \quad c \in \operatorname{derivable}(t) \end{array}}}
\vnbprop{derivable-and-elim}{\vnbstmt{\begin{array}{@{}l@{}} \forall t, s, c. \\ \quad (\operatorname{cr\text{-}and\text{-}elim}(t, s, c) \wedge (\operatorname{hyps\text{-}in}(s, \operatorname{derivable}(t)) \wedge (c \in \operatorname{seq\text{-}set}))) \Rightarrow  \\ \quad \quad c \in \operatorname{derivable}(t) \end{array}}}
\vnbprop{derivable-or-intro-left}{\vnbstmt{\begin{array}{@{}l@{}} \forall t, s, c. \\ \quad (\operatorname{cr\text{-}or\text{-}intro\text{-}left}(t, s, c) \wedge (\operatorname{hyps\text{-}in}(s, \operatorname{derivable}(t)) \wedge (c \in \operatorname{seq\text{-}set}))) \Rightarrow  \\ \quad \quad c \in \operatorname{derivable}(t) \end{array}}}
\vnbprop{derivable-or-intro-right}{\vnbstmt{\begin{array}{@{}l@{}} \forall t, s, c. \\ \quad (\operatorname{cr\text{-}or\text{-}intro\text{-}right}(t, s, c) \wedge (\operatorname{hyps\text{-}in}(s, \operatorname{derivable}(t)) \wedge (c \in \operatorname{seq\text{-}set}))) \Rightarrow  \\ \quad \quad c \in \operatorname{derivable}(t) \end{array}}}
\vnbprop{derivable-or-elim}{\vnbstmt{\begin{array}{@{}l@{}} \forall t, s, c. \\ \quad (\operatorname{cr\text{-}or\text{-}elim}(t, s, c) \wedge (\operatorname{hyps\text{-}in}(s, \operatorname{derivable}(t)) \wedge (c \in \operatorname{seq\text{-}set}))) \Rightarrow  \\ \quad \quad c \in \operatorname{derivable}(t) \end{array}}}
\vnbprop{derivable-implies-intro}{\vnbstmt{\begin{array}{@{}l@{}} \forall t, s, c. \\ \quad (\operatorname{cr\text{-}implies\text{-}intro}(t, s, c) \wedge (\operatorname{hyps\text{-}in}(s, \operatorname{derivable}(t)) \wedge (c \in \operatorname{seq\text{-}set}))) \Rightarrow  \\ \quad \quad c \in \operatorname{derivable}(t) \end{array}}}
\vnbprop{derivable-backchain}{\vnbstmt{\begin{array}{@{}l@{}} \forall t, s, c. \\ \quad (\operatorname{cr\text{-}backchain}(t, s, c) \wedge (\operatorname{hyps\text{-}in}(s, \operatorname{derivable}(t)) \wedge (c \in \operatorname{seq\text{-}set}))) \Rightarrow  \\ \quad \quad c \in \operatorname{derivable}(t) \end{array}}}
\vnbprop{derivable-not-intro}{\vnbstmt{\begin{array}{@{}l@{}} \forall t, s, c. \\ \quad (\operatorname{cr\text{-}not\text{-}intro}(t, s, c) \wedge (\operatorname{hyps\text{-}in}(s, \operatorname{derivable}(t)) \wedge (c \in \operatorname{seq\text{-}set}))) \Rightarrow  \\ \quad \quad c \in \operatorname{derivable}(t) \end{array}}}
\vnbprop{derivable-not-elim}{\vnbstmt{\begin{array}{@{}l@{}} \forall t, s, c. \\ \quad (\operatorname{cr\text{-}not\text{-}elim}(t, s, c) \wedge (\operatorname{hyps\text{-}in}(s, \operatorname{derivable}(t)) \wedge (c \in \operatorname{seq\text{-}set}))) \Rightarrow  \\ \quad \quad c \in \operatorname{derivable}(t) \end{array}}}
\vnbprop{derivable-iff-intro}{\vnbstmt{\begin{array}{@{}l@{}} \forall t, s, c. \\ \quad (\operatorname{cr\text{-}iff\text{-}intro}(t, s, c) \wedge (\operatorname{hyps\text{-}in}(s, \operatorname{derivable}(t)) \wedge (c \in \operatorname{seq\text{-}set}))) \Rightarrow  \\ \quad \quad c \in \operatorname{derivable}(t) \end{array}}}
\vnbprop{derivable-iff-elim}{\vnbstmt{\begin{array}{@{}l@{}} \forall t, s, c. \\ \quad (\operatorname{cr\text{-}iff\text{-}elim}(t, s, c) \wedge (\operatorname{hyps\text{-}in}(s, \operatorname{derivable}(t)) \wedge (c \in \operatorname{seq\text{-}set}))) \Rightarrow  \\ \quad \quad c \in \operatorname{derivable}(t) \end{array}}}
\vnbprop{derivable-forall-intro}{\vnbstmt{\begin{array}{@{}l@{}} \forall t, s, c. \\ \quad (\operatorname{cr\text{-}forall\text{-}intro}(t, s, c) \wedge (\operatorname{hyps\text{-}in}(s, \operatorname{derivable}(t)) \wedge (c \in \operatorname{seq\text{-}set}))) \Rightarrow  \\ \quad \quad c \in \operatorname{derivable}(t) \end{array}}}
\vnbprop{derivable-forall-elim}{\vnbstmt{\begin{array}{@{}l@{}} \forall t, s, c. \\ \quad (\operatorname{cr\text{-}forall\text{-}elim}(t, s, c) \wedge (\operatorname{hyps\text{-}in}(s, \operatorname{derivable}(t)) \wedge (c \in \operatorname{seq\text{-}set}))) \Rightarrow  \\ \quad \quad c \in \operatorname{derivable}(t) \end{array}}}
\vnbprop{derivable-forsome-intro}{\vnbstmt{\begin{array}{@{}l@{}} \forall t, s, c. \\ \quad (\operatorname{cr\text{-}forsome\text{-}intro}(t, s, c) \wedge (\operatorname{hyps\text{-}in}(s, \operatorname{derivable}(t)) \wedge (c \in \operatorname{seq\text{-}set}))) \Rightarrow  \\ \quad \quad c \in \operatorname{derivable}(t) \end{array}}}
\vnbprop{derivable-forsome-elim}{\vnbstmt{\begin{array}{@{}l@{}} \forall t, s, c. \\ \quad (\operatorname{cr\text{-}forsome\text{-}elim}(t, s, c) \wedge (\operatorname{hyps\text{-}in}(s, \operatorname{derivable}(t)) \wedge (c \in \operatorname{seq\text{-}set}))) \Rightarrow  \\ \quad \quad c \in \operatorname{derivable}(t) \end{array}}}
\vnbprop{derivable-assumption}{\vnbstmt{\begin{array}{@{}l@{}} \forall t, s, c. \\ \quad (\operatorname{cr\text{-}assumption}(t, s, c) \wedge (\operatorname{hyps\text{-}in}(s, \operatorname{derivable}(t)) \wedge (c \in \operatorname{seq\text{-}set}))) \Rightarrow  \\ \quad \quad c \in \operatorname{derivable}(t) \end{array}}}
\vnbprop{derivable-theorem-assumption}{\vnbstmt{\begin{array}{@{}l@{}} \forall t, s, c. \\ \quad (\operatorname{cr\text{-}theorem\text{-}assumption}(t, s, c) \wedge (\operatorname{hyps\text{-}in}(s, \operatorname{derivable}(t)) \wedge (c \in \operatorname{seq\text{-}set}))) \Rightarrow  \\ \quad \quad c \in \operatorname{derivable}(t) \end{array}}}
\vnbprop{derivable-cut}{\vnbstmt{\begin{array}{@{}l@{}} \forall t, s, c. \\ \quad (\operatorname{cr\text{-}cut}(t, s, c) \wedge (\operatorname{hyps\text{-}in}(s, \operatorname{derivable}(t)) \wedge (c \in \operatorname{seq\text{-}set}))) \Rightarrow  \\ \quad \quad c \in \operatorname{derivable}(t) \end{array}}}
\vnbprop{derivable-weakening}{\vnbstmt{\begin{array}{@{}l@{}} \forall t, s, c. \\ \quad (\operatorname{cr\text{-}weakening}(t, s, c) \wedge (\operatorname{hyps\text{-}in}(s, \operatorname{derivable}(t)) \wedge (c \in \operatorname{seq\text{-}set}))) \Rightarrow  \\ \quad \quad c \in \operatorname{derivable}(t) \end{array}}}
\vnbprop{derivable-contradiction}{\vnbstmt{\begin{array}{@{}l@{}} \forall t, s, c. \\ \quad (\operatorname{cr\text{-}contradiction}(t, s, c) \wedge (\operatorname{hyps\text{-}in}(s, \operatorname{derivable}(t)) \wedge (c \in \operatorname{seq\text{-}set}))) \Rightarrow  \\ \quad \quad c \in \operatorname{derivable}(t) \end{array}}}
\vnbprop{derivable-truth-intro}{\vnbstmt{\begin{array}{@{}l@{}} \forall t, s, c. \\ \quad (\operatorname{cr\text{-}truth\text{-}intro}(t, s, c) \wedge (\operatorname{hyps\text{-}in}(s, \operatorname{derivable}(t)) \wedge (c \in \operatorname{seq\text{-}set}))) \Rightarrow  \\ \quad \quad c \in \operatorname{derivable}(t) \end{array}}}
\vnbprop{derivable-eq-subst}{\vnbstmt{\begin{array}{@{}l@{}} \forall t, s, c. \\ \quad (\operatorname{cr\text{-}eq\text{-}subst}(t, s, c) \wedge (\operatorname{hyps\text{-}in}(s, \operatorname{derivable}(t)) \wedge (c \in \operatorname{seq\text{-}set}))) \Rightarrow  \\ \quad \quad c \in \operatorname{derivable}(t) \end{array}}}
\vnbprop{derivable-reflexivity}{\vnbstmt{\begin{array}{@{}l@{}} \forall t, s, c. \\ \quad (\operatorname{cr\text{-}reflexivity}(t, s, c) \wedge (\operatorname{hyps\text{-}in}(s, \operatorname{derivable}(t)) \wedge (c \in \operatorname{seq\text{-}set}))) \Rightarrow  \\ \quad \quad c \in \operatorname{derivable}(t) \end{array}}}
\vnbprop{derivable-quasi-reflexivity}{\vnbstmt{\begin{array}{@{}l@{}} \forall t, s, c. \\ \quad (\operatorname{cr\text{-}quasi\text{-}reflexivity}(t, s, c) \wedge (\operatorname{hyps\text{-}in}(s, \operatorname{derivable}(t)) \wedge (c \in \operatorname{seq\text{-}set}))) \Rightarrow  \\ \quad \quad c \in \operatorname{derivable}(t) \end{array}}}
\vnbprop{derivable-union-mem}{\vnbstmt{\begin{array}{@{}l@{}} \forall t, s, c. \\ \quad (\operatorname{cr\text{-}union\text{-}mem}(t, s, c) \wedge (\operatorname{hyps\text{-}in}(s, \operatorname{derivable}(t)) \wedge (c \in \operatorname{seq\text{-}set}))) \Rightarrow  \\ \quad \quad c \in \operatorname{derivable}(t) \end{array}}}
\vnbprop{derivable-intersection-mem}{\vnbstmt{\begin{array}{@{}l@{}} \forall t, s, c. \\ \quad (\operatorname{cr\text{-}intersection\text{-}mem}(t, s, c) \wedge (\operatorname{hyps\text{-}in}(s, \operatorname{derivable}(t)) \wedge (c \in \operatorname{seq\text{-}set}))) \Rightarrow  \\ \quad \quad c \in \operatorname{derivable}(t) \end{array}}}
\vnbprop{derivable-cartesian-mem}{\vnbstmt{\begin{array}{@{}l@{}} \forall t, s, c. \\ \quad (\operatorname{cr\text{-}cartesian\text{-}mem}(t, s, c) \wedge (\operatorname{hyps\text{-}in}(s, \operatorname{derivable}(t)) \wedge (c \in \operatorname{seq\text{-}set}))) \Rightarrow  \\ \quad \quad c \in \operatorname{derivable}(t) \end{array}}}
\vnbprop{derivable-tuples-mem}{\vnbstmt{\begin{array}{@{}l@{}} \forall t, s, c. \\ \quad (\operatorname{cr\text{-}tuples\text{-}mem}(t, s, c) \wedge (\operatorname{hyps\text{-}in}(s, \operatorname{derivable}(t)) \wedge (c \in \operatorname{seq\text{-}set}))) \Rightarrow  \\ \quad \quad c \in \operatorname{derivable}(t) \end{array}}}
\vnbprop{derivable-nth-eq}{\vnbstmt{\begin{array}{@{}l@{}} \forall t, s, c. \\ \quad (\operatorname{cr\text{-}nth\text{-}eq}(t, s, c) \wedge (\operatorname{hyps\text{-}in}(s, \operatorname{derivable}(t)) \wedge (c \in \operatorname{seq\text{-}set}))) \Rightarrow  \\ \quad \quad c \in \operatorname{derivable}(t) \end{array}}}
\vnbprop{derivable-length-eq}{\vnbstmt{\begin{array}{@{}l@{}} \forall t, s, c. \\ \quad (\operatorname{cr\text{-}length\text{-}eq}(t, s, c) \wedge (\operatorname{hyps\text{-}in}(s, \operatorname{derivable}(t)) \wedge (c \in \operatorname{seq\text{-}set}))) \Rightarrow  \\ \quad \quad c \in \operatorname{derivable}(t) \end{array}}}
\vnbprop{derivable-tuple-eq}{\vnbstmt{\begin{array}{@{}l@{}} \forall t, s, c. \\ \quad (\operatorname{cr\text{-}tuple\text{-}eq}(t, s, c) \wedge (\operatorname{hyps\text{-}in}(s, \operatorname{derivable}(t)) \wedge (c \in \operatorname{seq\text{-}set}))) \Rightarrow  \\ \quad \quad c \in \operatorname{derivable}(t) \end{array}}}
\vnbprop{derivable-if-true}{\vnbstmt{\begin{array}{@{}l@{}} \forall t, s, c. \\ \quad (\operatorname{cr\text{-}if\text{-}true}(t, s, c) \wedge (\operatorname{hyps\text{-}in}(s, \operatorname{derivable}(t)) \wedge (c \in \operatorname{seq\text{-}set}))) \Rightarrow  \\ \quad \quad c \in \operatorname{derivable}(t) \end{array}}}
\vnbprop{derivable-if-false}{\vnbstmt{\begin{array}{@{}l@{}} \forall t, s, c. \\ \quad (\operatorname{cr\text{-}if\text{-}false}(t, s, c) \wedge (\operatorname{hyps\text{-}in}(s, \operatorname{derivable}(t)) \wedge (c \in \operatorname{seq\text{-}set}))) \Rightarrow  \\ \quad \quad c \in \operatorname{derivable}(t) \end{array}}}
\vnbprop{derivable-sep-sethood}{\vnbstmt{\begin{array}{@{}l@{}} \forall t, s, c. \\ \quad (\operatorname{cr\text{-}sep\text{-}sethood}(t, s, c) \wedge (\operatorname{hyps\text{-}in}(s, \operatorname{derivable}(t)) \wedge (c \in \operatorname{seq\text{-}set}))) \Rightarrow  \\ \quad \quad c \in \operatorname{derivable}(t) \end{array}}}
\vnbprop{derivable-sep-mem}{\vnbstmt{\begin{array}{@{}l@{}} \forall t, s, c. \\ \quad (\operatorname{cr\text{-}sep\text{-}mem}(t, s, c) \wedge (\operatorname{hyps\text{-}in}(s, \operatorname{derivable}(t)) \wedge (c \in \operatorname{seq\text{-}set}))) \Rightarrow  \\ \quad \quad c \in \operatorname{derivable}(t) \end{array}}}
\vnbprop{derivable-comp-mem}{\vnbstmt{\begin{array}{@{}l@{}} \forall t, s, c. \\ \quad (\operatorname{cr\text{-}comp\text{-}mem}(t, s, c) \wedge (\operatorname{hyps\text{-}in}(s, \operatorname{derivable}(t)) \wedge (c \in \operatorname{seq\text{-}set}))) \Rightarrow  \\ \quad \quad c \in \operatorname{derivable}(t) \end{array}}}
\vnbprop{derivable-big-union-sethood}{\vnbstmt{\begin{array}{@{}l@{}} \forall t, s, c. \\ \quad (\operatorname{cr\text{-}big\text{-}union\text{-}sethood}(t, s, c) \wedge (\operatorname{hyps\text{-}in}(s, \operatorname{derivable}(t)) \wedge (c \in \operatorname{seq\text{-}set}))) \Rightarrow  \\ \quad \quad c \in \operatorname{derivable}(t) \end{array}}}
\vnbprop{derivable-big-union-mem}{\vnbstmt{\begin{array}{@{}l@{}} \forall t, s, c. \\ \quad (\operatorname{cr\text{-}big\text{-}union\text{-}mem}(t, s, c) \wedge (\operatorname{hyps\text{-}in}(s, \operatorname{derivable}(t)) \wedge (c \in \operatorname{seq\text{-}set}))) \Rightarrow  \\ \quad \quad c \in \operatorname{derivable}(t) \end{array}}}
\vnbprop{derivable-iota-def}{\vnbstmt{\begin{array}{@{}l@{}} \forall t, s, c. \\ \quad (\operatorname{cr\text{-}iota\text{-}def}(t, s, c) \wedge (\operatorname{hyps\text{-}in}(s, \operatorname{derivable}(t)) \wedge (c \in \operatorname{seq\text{-}set}))) \Rightarrow  \\ \quad \quad c \in \operatorname{derivable}(t) \end{array}}}
\vnbprop{derivable-iota-in-elim}{\vnbstmt{\begin{array}{@{}l@{}} \forall t, s, c. \\ \quad (\operatorname{cr\text{-}iota\text{-}in\text{-}elim}(t, s, c) \wedge (\operatorname{hyps\text{-}in}(s, \operatorname{derivable}(t)) \wedge (c \in \operatorname{seq\text{-}set}))) \Rightarrow  \\ \quad \quad c \in \operatorname{derivable}(t) \end{array}}}
\vnbprop{derivable-lambda-type}{\vnbstmt{\begin{array}{@{}l@{}} \forall t, s, c. \\ \quad (\operatorname{cr\text{-}lambda\text{-}type}(t, s, c) \wedge (\operatorname{hyps\text{-}in}(s, \operatorname{derivable}(t)) \wedge (c \in \operatorname{seq\text{-}set}))) \Rightarrow  \\ \quad \quad c \in \operatorname{derivable}(t) \end{array}}}
\vnbprop{derivable-lambda-beta-eq}{\vnbstmt{\begin{array}{@{}l@{}} \forall t, s, c. \\ \quad (\operatorname{cr\text{-}lambda\text{-}beta\text{-}eq}(t, s, c) \wedge (\operatorname{hyps\text{-}in}(s, \operatorname{derivable}(t)) \wedge (c \in \operatorname{seq\text{-}set}))) \Rightarrow  \\ \quad \quad c \in \operatorname{derivable}(t) \end{array}}}
\vnbprop{derivable-eq-subst-sym}{\vnbstmt{\begin{array}{@{}l@{}} \forall t, s, c. \\ \quad (\operatorname{cr\text{-}eq\text{-}subst\text{-}sym}(t, s, c) \wedge (\operatorname{hyps\text{-}in}(s, \operatorname{derivable}(t)) \wedge (c \in \operatorname{seq\text{-}set}))) \Rightarrow  \\ \quad \quad c \in \operatorname{derivable}(t) \end{array}}}
\vnbprop{derivable-def-strict-in-l}{\vnbstmt{\begin{array}{@{}l@{}} \forall t, s, c. \\ \quad (\operatorname{cr\text{-}def\text{-}strict\text{-}in\text{-}l}(t, s, c) \wedge (\operatorname{hyps\text{-}in}(s, \operatorname{derivable}(t)) \wedge (c \in \operatorname{seq\text{-}set}))) \Rightarrow  \\ \quad \quad c \in \operatorname{derivable}(t) \end{array}}}
\vnbprop{derivable-def-strict-in-r}{\vnbstmt{\begin{array}{@{}l@{}} \forall t, s, c. \\ \quad (\operatorname{cr\text{-}def\text{-}strict\text{-}in\text{-}r}(t, s, c) \wedge (\operatorname{hyps\text{-}in}(s, \operatorname{derivable}(t)) \wedge (c \in \operatorname{seq\text{-}set}))) \Rightarrow  \\ \quad \quad c \in \operatorname{derivable}(t) \end{array}}}
\vnbprop{derivable-def-total}{\vnbstmt{\begin{array}{@{}l@{}} \forall t, s, c. \\ \quad (\operatorname{cr\text{-}def\text{-}total}(t, s, c) \wedge (\operatorname{hyps\text{-}in}(s, \operatorname{derivable}(t)) \wedge (c \in \operatorname{seq\text{-}set}))) \Rightarrow  \\ \quad \quad c \in \operatorname{derivable}(t) \end{array}}}
\vnbprop{derivable-def-binder}{\vnbstmt{\begin{array}{@{}l@{}} \forall t, s, c. \\ \quad (\operatorname{cr\text{-}def\text{-}binder}(t, s, c) \wedge (\operatorname{hyps\text{-}in}(s, \operatorname{derivable}(t)) \wedge (c \in \operatorname{seq\text{-}set}))) \Rightarrow  \\ \quad \quad c \in \operatorname{derivable}(t) \end{array}}}
\vnbprop{derivable-def-comp}{\vnbstmt{\begin{array}{@{}l@{}} \forall t, s, c. \\ \quad (\operatorname{cr\text{-}def\text{-}comp}(t, s, c) \wedge (\operatorname{hyps\text{-}in}(s, \operatorname{derivable}(t)) \wedge (c \in \operatorname{seq\text{-}set}))) \Rightarrow  \\ \quad \quad c \in \operatorname{derivable}(t) \end{array}}}
\vnbprop{derivable-def-list}{\vnbstmt{\begin{array}{@{}l@{}} \forall t, s, c. \\ \quad (\operatorname{cr\text{-}def\text{-}list}(t, s, c) \wedge (\operatorname{hyps\text{-}in}(s, \operatorname{derivable}(t)) \wedge (c \in \operatorname{seq\text{-}set}))) \Rightarrow  \\ \quad \quad c \in \operatorname{derivable}(t) \end{array}}}
\vnbprop{derivable-def-strict-arg}{\vnbstmt{\begin{array}{@{}l@{}} \forall t, s, c. \\ \quad (\operatorname{cr\text{-}def\text{-}strict\text{-}arg}(t, s, c) \wedge (\operatorname{hyps\text{-}in}(s, \operatorname{derivable}(t)) \wedge (c \in \operatorname{seq\text{-}set}))) \Rightarrow  \\ \quad \quad c \in \operatorname{derivable}(t) \end{array}}}
\vnbprop{derivable-def-strict-dom}{\vnbstmt{\begin{array}{@{}l@{}} \forall t, s, c. \\ \quad (\operatorname{cr\text{-}def\text{-}strict\text{-}dom}(t, s, c) \wedge (\operatorname{hyps\text{-}in}(s, \operatorname{derivable}(t)) \wedge (c \in \operatorname{seq\text{-}set}))) \Rightarrow  \\ \quad \quad c \in \operatorname{derivable}(t) \end{array}}}
\vnbprop{derivable-def-strict-pred}{\vnbstmt{\begin{array}{@{}l@{}} \forall t, s, c. \\ \quad (\operatorname{cr\text{-}def\text{-}strict\text{-}pred}(t, s, c) \wedge (\operatorname{hyps\text{-}in}(s, \operatorname{derivable}(t)) \wedge (c \in \operatorname{seq\text{-}set}))) \Rightarrow  \\ \quad \quad c \in \operatorname{derivable}(t) \end{array}}}
\vnbprop{derivable-iota-cr}{\vnbstmt{\begin{array}{@{}l@{}} \forall t, s, c. \\ \quad (\operatorname{cr\text{-}iota\text{-}cr}(t, s, c) \wedge (\operatorname{hyps\text{-}in}(s, \operatorname{derivable}(t)) \wedge (c \in \operatorname{seq\text{-}set}))) \Rightarrow  \\ \quad \quad c \in \operatorname{derivable}(t) \end{array}}}
\vnbprop{derivable-tbind-cr}{\vnbstmt{\begin{array}{@{}l@{}} \forall t, s, c. \\ \quad (\operatorname{cr\text{-}tbind\text{-}cr}(t, s, c) \wedge (\operatorname{hyps\text{-}in}(s, \operatorname{derivable}(t)) \wedge (c \in \operatorname{seq\text{-}set}))) \Rightarrow  \\ \quad \quad c \in \operatorname{derivable}(t) \end{array}}}
\vnbprop{derivable-functoid-beta-eq}{\vnbstmt{\begin{array}{@{}l@{}} \forall t, s, c. \\ \quad (\operatorname{cr\text{-}functoid\text{-}beta\text{-}eq}(t, s, c) \wedge (\operatorname{hyps\text{-}in}(s, \operatorname{derivable}(t)) \wedge (c \in \operatorname{seq\text{-}set}))) \Rightarrow  \\ \quad \quad c \in \operatorname{derivable}(t) \end{array}}}
\vnbprop{derivable-tuple-eta}{\vnbstmt{\begin{array}{@{}l@{}} \forall t, s, c. \\ \quad (\operatorname{cr\text{-}tuple\text{-}eta}(t, s, c) \wedge (\operatorname{hyps\text{-}in}(s, \operatorname{derivable}(t)) \wedge (c \in \operatorname{seq\text{-}set}))) \Rightarrow  \\ \quad \quad c \in \operatorname{derivable}(t) \end{array}}}
\vnbprop{iota-cr-refuses-eq}{\vnbstmt{\begin{array}{@{}l@{}} \forall t, s, g, a, b. \\ \quad ((g \in \operatorname{ctx\text{-}set}) \wedge ((\operatorname{enc\text{-}iota}(a) \in \mathit{clt}) \wedge (\operatorname{enc\text{-}iota}(b) \in \mathit{clt}))) \Rightarrow  \\ \quad \quad \neg \operatorname{cr\text{-}iota\text{-}cr}(t, s, {[}g, \operatorname{enc\text{-}eq}(\operatorname{enc\text{-}iota}(a), \operatorname{enc\text{-}iota}(b))]) \end{array}}}
\vnbprop{tbind-cr-refuses-eq}{\vnbstmt{\begin{array}{@{}l@{}} \forall t, s, g, h, d, a, b. \\ \quad ((g \in \operatorname{ctx\text{-}set}) \wedge ((\operatorname{ebind}(h, 1, {[}d], a) \in \mathit{clt}) \wedge (\operatorname{ebind}(h, 1, {[}d], b) \in \mathit{clt}))) \Rightarrow  \\ \quad \quad \neg \operatorname{cr\text{-}tbind\text{-}cr}(t, s, {[}g, \operatorname{enc\text{-}eq}(\operatorname{ebind}(h, 1, {[}d], a), \operatorname{ebind}(h, 1, {[}d], b))]) \end{array}}}
\vnbprop{tbind-cr-comp-refuses-eq}{\vnbstmt{\begin{array}{@{}l@{}} \forall t, s, g, a, b. \\ \quad ((g \in \operatorname{ctx\text{-}set}) \wedge ((\operatorname{ebind}(23, 1, {[}], a) \in \mathit{clt}) \wedge (\operatorname{ebind}(23, 1, {[}], b) \in \mathit{clt}))) \Rightarrow  \\ \quad \quad \neg \operatorname{cr\text{-}tbind\text{-}cr}(t, s, {[}g, \operatorname{enc\text{-}eq}(\operatorname{ebind}(23, 1, {[}], a), \operatorname{ebind}(23, 1, {[}], b))]) \end{array}}}
\vnbprop{functoid-beta-eq-refuses-eq}{\vnbstmt{\begin{array}{@{}l@{}} \forall t, s, g, h, l, b. \\ \quad ((g \in \operatorname{ctx\text{-}set}) \wedge ((\operatorname{eapp}(h, l) \in \mathit{clt}) \wedge (\operatorname{openl}(b, l, \operatorname{length}(l)) \in \mathit{clt}))) \Rightarrow  \\ \quad \quad \neg \operatorname{cr\text{-}functoid\text{-}beta\text{-}eq}(t, s, {[}g, \operatorname{enc\text{-}eq}(\operatorname{eapp}(h, l), \operatorname{openl}(b, l, \operatorname{length}(l)))]) \end{array}}}
\vnbprop{fun-def-heads}{\vnbstmt{\begin{array}{@{}l@{}} \forall h, n, b. \\ \quad \operatorname{fun\text{-}def}(h, n, b) \Rightarrow ((32 \leq h) \wedge (\neg \operatorname{strict\text{-}head}(h) \wedge \neg \operatorname{total\text{-}head}(h))) \end{array}}}
\vnbprop{derivable-assumption-c}{\vnbstmt{\forall t, g, a.\; ((({[}g, a] \in \operatorname{seq\text{-}set}) \wedge (a \in g)) \Rightarrow ({[}g, a] \in \operatorname{derivable}(t)))}}
\vnbprop{derivable-backchain-c}{\vnbstmt{\begin{array}{@{}l@{}} \forall t, g, f, a. \\ \quad (({[}g, f] \in \operatorname{seq\text{-}set}) \wedge ((\operatorname{enc\text{-}imp}(a, f) \in g) \wedge ({[}g, a] \in \operatorname{derivable}(t)))) \Rightarrow  \\ \quad \quad {[}g, f] \in \operatorname{derivable}(t) \end{array}}}
\vnbprop{derivable-cut-c}{\vnbstmt{\begin{array}{@{}l@{}} \forall t, g, f, a. \\ \quad (({[}g, f] \in \operatorname{seq\text{-}set}) \wedge (({[}g, a] \in \operatorname{derivable}(t)) \wedge ({[}\operatorname{ctx\text{-}add}(g, a), f] \in \operatorname{derivable}(t)))) \Rightarrow  \\ \quad \quad {[}g, f] \in \operatorname{derivable}(t) \end{array}}}
\vnbprop{derivable-weakening-c}{\vnbstmt{\begin{array}{@{}l@{}} \forall t, g, f, d. \\ \quad (({[}g, f] \in \operatorname{seq\text{-}set}) \wedge ((d \subseteq g) \wedge ({[}d, f] \in \operatorname{derivable}(t)))) \Rightarrow  \\ \quad \quad {[}g, f] \in \operatorname{derivable}(t) \end{array}}}
\vnbprop{derivable-not-intro-c}{\vnbstmt{\begin{array}{@{}l@{}} \forall t, g, a. \\ \quad (({[}g, \operatorname{enc\text{-}not}(a)] \in \operatorname{seq\text{-}set}) \wedge ({[}\operatorname{ctx\text{-}add}(g, a), \operatorname{enc\text{-}falsity}] \in \operatorname{derivable}(t))) \Rightarrow  \\ \quad \quad {[}g, \operatorname{enc\text{-}not}(a)] \in \operatorname{derivable}(t) \end{array}}}
\vnbprop{derivable-not-elim-c}{\vnbstmt{\begin{array}{@{}l@{}} \forall t, g, f, a. \\ \quad (({[}g, f] \in \operatorname{seq\text{-}set}) \wedge ((\operatorname{enc\text{-}not}(a) \in g) \wedge (a \in g))) \Rightarrow ({[}g, f] \in \operatorname{derivable}(t)) \end{array}}}
\vnbprop{derivable-or-intro-left-c}{\vnbstmt{\begin{array}{@{}l@{}} \forall t, g, a, b. \\ \quad (({[}g, \operatorname{enc\text{-}or}(a, b)] \in \operatorname{seq\text{-}set}) \wedge ({[}g, a] \in \operatorname{derivable}(t))) \Rightarrow  \\ \quad \quad {[}g, \operatorname{enc\text{-}or}(a, b)] \in \operatorname{derivable}(t) \end{array}}}
\vnbprop{derivable-or-intro-right-c}{\vnbstmt{\begin{array}{@{}l@{}} \forall t, g, a, b. \\ \quad (({[}g, \operatorname{enc\text{-}or}(a, b)] \in \operatorname{seq\text{-}set}) \wedge ({[}g, b] \in \operatorname{derivable}(t))) \Rightarrow  \\ \quad \quad {[}g, \operatorname{enc\text{-}or}(a, b)] \in \operatorname{derivable}(t) \end{array}}}
\vnbprop{derivable-iff-elim-c}{\vnbstmt{\begin{array}{@{}l@{}} \forall t, g, f, a, b. \\ \quad {[}g, f] \in \operatorname{seq\text{-}set} \wedge \\ \quad \operatorname{enc\text{-}iff}(a, b) \in g \wedge \\ \quad {[}\operatorname{ctx\text{-}add}(\operatorname{ctx\text{-}add}(g, \operatorname{enc\text{-}imp}(a, b)), \operatorname{enc\text{-}imp}(b, a)), f] \in \operatorname{derivable}(t) \Rightarrow \\ \quad \quad {[}g, f] \in \operatorname{derivable}(t) \end{array}}}
\vnbprop{derivable-union-mem-c}{\vnbstmt{\begin{array}{@{}l@{}} \forall t, g, x, l. \\ \quad {[}g, \operatorname{enc\text{-}iff}(\operatorname{enc\text{-}in}(x, \operatorname{enc\text{-}union}(l)), \operatorname{enc\text{-}orl}(x, l, \operatorname{length}(l)))] \in \operatorname{seq\text{-}set} \wedge \\ \quad \forall d \in \mathbb{N}. \\ \quad \quad ((1 \leq d) \wedge (d \leq \operatorname{length}(l))) \Rightarrow  \\ \quad \quad \quad {[}g, \operatorname{enc\text{-}eq}(\operatorname{nth}(d, l), \operatorname{nth}(d, l))] \in \operatorname{derivable}(t) \Rightarrow \\ \quad \quad {[}g, \operatorname{enc\text{-}iff}(\operatorname{enc\text{-}in}(x, \operatorname{enc\text{-}union}(l)), \operatorname{enc\text{-}orl}(x, l, \operatorname{length}(l)))] \in \operatorname{derivable}(t) \end{array}}}
\vnbprop{derivable-union-mem-u}{\vnbstmt{\begin{array}{@{}l@{}} \forall t, g, x, l. \\ \quad {[}g, \operatorname{enc\text{-}iff}(\operatorname{enc\text{-}in}(x, \operatorname{enc\text{-}union}(l)), \operatorname{enc\text{-}orl}(x, l, \operatorname{length}(l)))] \in \operatorname{seq\text{-}set} \wedge \\ \quad {[}g, \operatorname{enc\text{-}eq}(\operatorname{enc\text{-}union}(l), \operatorname{enc\text{-}union}(l))] \in \operatorname{derivable}(t) \Rightarrow \\ \quad \quad {[}g, \operatorname{enc\text{-}iff}(\operatorname{enc\text{-}in}(x, \operatorname{enc\text{-}union}(l)), \operatorname{enc\text{-}orl}(x, l, \operatorname{length}(l)))] \in \operatorname{derivable}(t) \end{array}}}
\vnbprop{adm-detach}{\vnbstmt{\begin{array}{@{}l@{}} \forall t, s, c. \\ \quad (\operatorname{rel\text{-}detach}(t, s, c) \wedge \operatorname{hyps\text{-}in}(s, \operatorname{derivable}(t))) \Rightarrow (c \in \operatorname{derivable}(t)) \end{array}}}
\vnbprop{orl-clf}{\vnbstmt{\begin{array}{@{}l@{}} \forall n \in \mathbb{N}, x, l. \\ \quad ((x \in \mathit{clt}) \wedge ((l \in \operatorname{tuples}(\operatorname{expr\text{-}set})) \wedge (\forall u \in \operatorname{make\text{-}set}(l).\; (u \in \mathit{clt}) \wedge (n \leq \operatorname{length}(l))))) \Rightarrow  \\ \quad \quad \operatorname{enc\text{-}orl}(x, l, n) \in \mathit{clf} \end{array}}}
\vnbprop{orl-der}{\vnbstmt{\begin{array}{@{}l@{}} \forall n \in \mathbb{N}, t, g, x, l, k. \\ \quad {[}g, \operatorname{enc\text{-}in}(x, \operatorname{nth}(k, l))] \in \operatorname{derivable}(t) \wedge \\ \quad g \in \operatorname{ctx\text{-}set} \wedge \\ \quad x \in \mathit{clt} \wedge \\ \quad l \in \operatorname{tuples}(\operatorname{expr\text{-}set}) \wedge \\ \quad \forall u \in \operatorname{make\text{-}set}(l).\; (u \in \mathit{clt}) \wedge \\ \quad k \in \mathbb{N} \wedge \\ \quad 1 \leq k \wedge \\ \quad k \leq n \wedge \\ \quad n \leq \operatorname{length}(l) \Rightarrow \\ \quad \quad {[}g, \operatorname{enc\text{-}orl}(x, l, n)] \in \operatorname{derivable}(t) \end{array}}}
\vnbprop{adm-union-intro}{\vnbstmt{\begin{array}{@{}l@{}} \forall t, s, c. \\ \quad (\operatorname{rel\text{-}union\text{-}intro}(t, s, c) \wedge \operatorname{hyps\text{-}in}(s, \operatorname{derivable}(t))) \Rightarrow (c \in \operatorname{derivable}(t)) \end{array}}}
\vnbprop{clf-is-set}{\vnbstmt{\mathit{clf} \in \mathbf{Set}}}
\vnbprop{seq-set-is-set}{\vnbstmt{\operatorname{seq\text{-}set} \in \mathbf{Set}}}
\vnbprop{ctx-fv-finite}{\vnbstmt{\begin{array}{@{}l@{}} \forall g \in \operatorname{ctx\text{-}set}. \\ \quad (\operatorname{ctx\text{-}fv}(g) \in \mathbf{Set}) \wedge ((\operatorname{ctx\text{-}fv}(g) \subseteq \mathbb{N}) \wedge (|\operatorname{ctx\text{-}fv}(g)| \in \mathbb{N})) \end{array}}}
\vnbprop{fresh-name-exists}{\vnbstmt{\begin{array}{@{}l@{}} \forall g, f. \\ \quad ((g \in \operatorname{ctx\text{-}set}) \wedge (f \in \operatorname{expr\text{-}set})) \Rightarrow  \\ \quad \quad \exists n \in \mathbb{N}.\; \neg (n \in (\operatorname{ctx\text{-}fv}(g) \cup \operatorname{expr\text{-}fv}(f))) \end{array}}}
\vnbprop{fv-app2}{\vnbstmt{\begin{array}{@{}l@{}} \forall h, a, b. \\ \quad ((h \in \mathbb{N}) \wedge ((a \in \operatorname{expr\text{-}set}) \wedge (b \in \operatorname{expr\text{-}set}))) \Rightarrow  \\ \quad \quad \operatorname{expr\text{-}fv}(\operatorname{eapp}(h, {[}a, b])) = (\operatorname{expr\text{-}fv}(a) \cup \operatorname{expr\text{-}fv}(b)) \end{array}}}
\vnbprop{fv-app1}{\vnbstmt{\begin{array}{@{}l@{}} \forall h, a. \\ \quad ((h \in \mathbb{N}) \wedge (a \in \operatorname{expr\text{-}set})) \Rightarrow (\operatorname{expr\text{-}fv}(\operatorname{eapp}(h, {[}a])) = \operatorname{expr\text{-}fv}(a)) \end{array}}}
\vnbprop{fv-app0}{\vnbstmt{\forall h \in \mathbb{N}.\; (\operatorname{expr\text{-}fv}(\operatorname{eapp}(h, {[}])) = \emptyset)}}
\vnbprop{fv-bind0}{\vnbstmt{\begin{array}{@{}l@{}} \forall h, a, b. \\ \quad ((h \in \mathbb{N}) \wedge ((a \in \mathbb{N}) \wedge (b \in \operatorname{expr\text{-}set}))) \Rightarrow  \\ \quad \quad \operatorname{expr\text{-}fv}(\operatorname{ebind}(h, a, {[}], b)) = \operatorname{expr\text{-}fv}(b) \end{array}}}
\vnbprop{wks-eapp-transfer}{\vnbstmt{\begin{array}{@{}l@{}} \forall h, l, a. \\ \quad h \in \mathbb{N} \wedge \\ \quad l \in \operatorname{tuples}(\operatorname{expr\text{-}set}) \wedge \\ \quad a \in \operatorname{tuples}(\operatorname{expr\text{-}set}) \wedge \\ \quad \operatorname{length}(a) = \operatorname{length}(l) \wedge \\ \quad \forall i \in \mathbb{N}. \\ \quad \quad ((1 \leq i) \wedge (i \leq \operatorname{length}(l))) \Rightarrow (\operatorname{expr\text{-}wks}(\operatorname{nth}(i, a)) = \operatorname{expr\text{-}wks}(\operatorname{nth}(i, l))) \Rightarrow \\ \quad \quad \operatorname{expr\text{-}wks}(\operatorname{eapp}(h, a)) = \operatorname{expr\text{-}wks}(\operatorname{eapp}(h, l)) \end{array}}}
\vnbprop{wks-ebind-transfer}{\vnbstmt{\begin{array}{@{}l@{}} \forall h, c, a, d, b, f. \\ \quad h \in \mathbb{N} \wedge \\ \quad c \in \mathbb{N} \wedge \\ \quad a \in \operatorname{tuples}(\operatorname{expr\text{-}set}) \wedge \\ \quad d \in \operatorname{tuples}(\operatorname{expr\text{-}set}) \wedge \\ \quad b \in \operatorname{expr\text{-}set} \wedge \\ \quad f \in \operatorname{expr\text{-}set} \wedge \\ \quad \operatorname{length}(d) = \operatorname{length}(a) \wedge \\ \quad \forall i \in \mathbb{N}. \\ \quad \quad ((1 \leq i) \wedge (i \leq \operatorname{length}(a))) \Rightarrow (\operatorname{expr\text{-}wks}(\operatorname{nth}(i, d)) = \operatorname{expr\text{-}wks}(\operatorname{nth}(i, a))) \wedge \\ \quad \operatorname{expr\text{-}wks}(f) = \operatorname{expr\text{-}wks}(b) \Rightarrow \\ \quad \quad \operatorname{expr\text{-}wks}(\operatorname{ebind}(h, c, d, f)) = \operatorname{expr\text{-}wks}(\operatorname{ebind}(h, c, a, b)) \end{array}}}
\vnbprop{wks-subst}{\vnbstmt{\begin{array}{@{}l@{}} \forall m.; \forall u, e \in \operatorname{expr\text{-}set}. \\ \quad (\operatorname{expr\text{-}wks}(u) = \{1\}) \Rightarrow (\operatorname{expr\text{-}wks}(\operatorname{expr\text{-}subst}(e, m, u)) = \operatorname{expr\text{-}wks}(e)) \end{array}}}
\vnbprop{wks-open}{\vnbstmt{\begin{array}{@{}l@{}} \forall u, e \in \operatorname{expr\text{-}set}, x \in \mathbb{N}. \\ \quad (\operatorname{expr\text{-}wks}(u) = \{1\}) \Rightarrow (\operatorname{expr\text{-}wks}(\operatorname{expr\text{-}open}(e, x, u)) = \operatorname{expr\text{-}wks}(e)) \end{array}}}
\vnbprop{enc-inter-unfold}{\vnbstmt{\forall l.\; (\operatorname{enc\text{-}inter}(l) \simeq \operatorname{eapp}(12, l))}}
\vnbprop{enc-tuples-unfold}{\vnbstmt{\forall p.\; (\operatorname{enc\text{-}tuples}(p) \simeq \operatorname{eapp}(14, {[}p]))}}
\vnbprop{enc-list-unfold}{\vnbstmt{\forall l.\; (\operatorname{enc\text{-}list}(l) \simeq \operatorname{eapp}(17, l))}}
\vnbprop{clt-inter}{\vnbstmt{\begin{array}{@{}l@{}} \forall l \in \operatorname{tuples}(\operatorname{expr\text{-}set}). \\ \quad (\operatorname{eapp}(12, l) \in \mathit{clt}) \Leftrightarrow \forall u \in \operatorname{make\text{-}set}(l).\; (u \in \mathit{clt}) \end{array}}}
\vnbprop{clt-list}{\vnbstmt{\begin{array}{@{}l@{}} \forall l \in \operatorname{tuples}(\operatorname{expr\text{-}set}). \\ \quad (\operatorname{eapp}(17, l) \in \mathit{clt}) \Leftrightarrow \forall u \in \operatorname{make\text{-}set}(l).\; (u \in \mathit{clt}) \end{array}}}
\vnbprop{clt-tuples1}{\vnbstmt{\forall a \in \operatorname{expr\text{-}set}.\; ((\operatorname{eapp}(14, {[}a]) \in \mathit{clt}) \Leftrightarrow (a \in \mathit{clt}))}}
\vnbprop{clt-inter-elim}{\vnbstmt{\begin{array}{@{}l@{}} \forall l. \\ \quad (\operatorname{enc\text{-}inter}(l) \in \mathit{clt}) \Rightarrow  \\ \quad \quad (l \in \operatorname{tuples}(\operatorname{expr\text{-}set})) \wedge \forall u \in \operatorname{make\text{-}set}(l).\; (u \in \mathit{clt}) \end{array}}}
\vnbprop{clt-list-elim}{\vnbstmt{\begin{array}{@{}l@{}} \forall l. \\ \quad (\operatorname{enc\text{-}list}(l) \in \mathit{clt}) \Rightarrow  \\ \quad \quad (l \in \operatorname{tuples}(\operatorname{expr\text{-}set})) \wedge \forall u \in \operatorname{make\text{-}set}(l).\; (u \in \mathit{clt}) \end{array}}}
\vnbprop{enc-cart-unfold}{\vnbstmt{\forall l.\; (\operatorname{enc\text{-}cart}(l) \simeq \operatorname{eapp}(13, l))}}
\vnbprop{clt-cart}{\vnbstmt{\begin{array}{@{}l@{}} \forall l \in \operatorname{tuples}(\operatorname{expr\text{-}set}). \\ \quad (\operatorname{eapp}(13, l) \in \mathit{clt}) \Leftrightarrow \forall u \in \operatorname{make\text{-}set}(l).\; (u \in \mathit{clt}) \end{array}}}
\vnbprop{clt-cart-elim}{\vnbstmt{\begin{array}{@{}l@{}} \forall l. \\ \quad (\operatorname{enc\text{-}cart}(l) \in \mathit{clt}) \Rightarrow  \\ \quad \quad (l \in \operatorname{tuples}(\operatorname{expr\text{-}set})) \wedge \forall u \in \operatorname{make\text{-}set}(l).\; (u \in \mathit{clt}) \end{array}}}
\vnbprop{clt-tuples-elim}{\vnbstmt{\forall p.\; ((\operatorname{enc\text{-}tuples}(p) \in \mathit{clt}) \Rightarrow (p \in \mathit{clt}))}}
\vnbprop{ctx-diff-sub}{\vnbstmt{\forall g, b.\; (\operatorname{difference}(g, b) \subseteq g)}}
\vnbprop{enc-andl-clf}{\vnbstmt{\begin{array}{@{}l@{}} \forall n \in \mathbb{N}, x, l. \\ \quad ((x \in \mathit{clt}) \wedge ((l \in \operatorname{tuples}(\operatorname{expr\text{-}set})) \wedge (\forall u \in \operatorname{make\text{-}set}(l).\; (u \in \mathit{clt}) \wedge (n \leq \operatorname{length}(l))))) \Rightarrow  \\ \quad \quad \operatorname{enc\text{-}andl}(x, l, n) \in \mathit{clf} \end{array}}}
\vnbprop{enc-ande-clf}{\vnbstmt{\begin{array}{@{}l@{}} \forall n \in \mathbb{N}, l.; \forall a \in \mathit{clt}. \\ \quad ((l \in \operatorname{tuples}(\operatorname{expr\text{-}set})) \wedge (\forall u \in \operatorname{make\text{-}set}(l).\; (u \in \mathit{clt}) \wedge (n \leq \operatorname{length}(l)))) \Rightarrow  \\ \quad \quad \operatorname{enc\text{-}ande}(l, a, n) \in \mathit{clf} \end{array}}}
\vnbprop{enc-andl-der}{\vnbstmt{\begin{array}{@{}l@{}} \forall n \in \mathbb{N}, t, g, x, l. \\ \quad g \in \operatorname{ctx\text{-}set} \wedge \\ \quad x \in \mathit{clt} \wedge \\ \quad l \in \operatorname{tuples}(\operatorname{expr\text{-}set}) \wedge \\ \quad \forall u \in \operatorname{make\text{-}set}(l).\; (u \in \mathit{clt}) \wedge \\ \quad \forall k \in \mathbb{N}. \\ \quad \quad ((1 \leq k) \wedge (k \leq n)) \Rightarrow ({[}g, \operatorname{enc\text{-}in}(x, \operatorname{nth}(k, l))] \in \operatorname{derivable}(t)) \wedge \\ \quad n \leq \operatorname{length}(l) \Rightarrow \\ \quad \quad {[}g, \operatorname{enc\text{-}andl}(x, l, n)] \in \operatorname{derivable}(t) \end{array}}}
\vnbprop{enc-ande-der}{\vnbstmt{\begin{array}{@{}l@{}} \forall n \in \mathbb{N}, t, g, l, a. \\ \quad g \in \operatorname{ctx\text{-}set} \wedge \\ \quad a \in \mathit{clt} \wedge \\ \quad l \in \operatorname{tuples}(\operatorname{expr\text{-}set}) \wedge \\ \quad \forall u \in \operatorname{make\text{-}set}(l).\; (u \in \mathit{clt}) \wedge \\ \quad \forall k \in \mathbb{N}. \\ \quad \quad ((1 \leq k) \wedge (k \leq n)) \Rightarrow ({[}g, \operatorname{enc\text{-}in}(\operatorname{nth}(k, l), a)] \in \operatorname{derivable}(t)) \wedge \\ \quad n \leq \operatorname{length}(l) \Rightarrow \\ \quad \quad {[}g, \operatorname{enc\text{-}ande}(l, a, n)] \in \operatorname{derivable}(t) \end{array}}}
\vnbprop{enc-andl-elim-der}{\vnbstmt{\begin{array}{@{}l@{}} \forall n \in \mathbb{N}, t, g, x, l, k. \\ \quad {[}g, \operatorname{enc\text{-}andl}(x, l, n)] \in \operatorname{derivable}(t) \wedge \\ \quad g \in \operatorname{ctx\text{-}set} \wedge \\ \quad x \in \mathit{clt} \wedge \\ \quad l \in \operatorname{tuples}(\operatorname{expr\text{-}set}) \wedge \\ \quad \forall u \in \operatorname{make\text{-}set}(l).\; (u \in \mathit{clt}) \wedge \\ \quad k \in \mathbb{N} \wedge \\ \quad 1 \leq k \wedge \\ \quad k \leq n \wedge \\ \quad n \leq \operatorname{length}(l) \Rightarrow \\ \quad \quad {[}g, \operatorname{enc\text{-}in}(x, \operatorname{nth}(k, l))] \in \operatorname{derivable}(t) \end{array}}}
\vnbprop{enc-ande-elim-der}{\vnbstmt{\begin{array}{@{}l@{}} \forall n \in \mathbb{N}, t, g, l, a, k. \\ \quad {[}g, \operatorname{enc\text{-}ande}(l, a, n)] \in \operatorname{derivable}(t) \wedge \\ \quad g \in \operatorname{ctx\text{-}set} \wedge \\ \quad a \in \mathit{clt} \wedge \\ \quad l \in \operatorname{tuples}(\operatorname{expr\text{-}set}) \wedge \\ \quad \forall u \in \operatorname{make\text{-}set}(l).\; (u \in \mathit{clt}) \wedge \\ \quad k \in \mathbb{N} \wedge \\ \quad 1 \leq k \wedge \\ \quad k \leq n \wedge \\ \quad n \leq \operatorname{length}(l) \Rightarrow \\ \quad \quad {[}g, \operatorname{enc\text{-}in}(\operatorname{nth}(k, l), a)] \in \operatorname{derivable}(t) \end{array}}}
\vnbprop{orl-elim-der}{\vnbstmt{\begin{array}{@{}l@{}} \forall n \in \mathbb{N}, t, g, f, x, l. \\ \quad g \in \operatorname{ctx\text{-}set} \wedge \\ \quad f \in \mathit{clf} \wedge \\ \quad x \in \mathit{clt} \wedge \\ \quad l \in \operatorname{tuples}(\operatorname{expr\text{-}set}) \wedge \\ \quad \forall u \in \operatorname{make\text{-}set}(l).\; (u \in \mathit{clt}) \wedge \\ \quad \forall k \in \mathbb{N}. \\ \quad \quad ((1 \leq k) \wedge (k \leq n)) \Rightarrow ({[}\operatorname{ctx\text{-}add}(g, \operatorname{enc\text{-}in}(x, \operatorname{nth}(k, l))), f] \in \operatorname{derivable}(t)) \wedge \\ \quad n \leq \operatorname{length}(l) \Rightarrow \\ \quad \quad {[}\operatorname{ctx\text{-}add}(g, \operatorname{enc\text{-}orl}(x, l, n)), f] \in \operatorname{derivable}(t) \end{array}}}
\vnbprop{adm-union-elim}{\vnbstmt{\begin{array}{@{}l@{}} \forall t, s, c. \\ \quad (\operatorname{rel\text{-}union\text{-}elim}(t, s, c) \wedge \operatorname{hyps\text{-}in}(s, \operatorname{derivable}(t))) \Rightarrow (c \in \operatorname{derivable}(t)) \end{array}}}
\vnbprop{adm-intersection-intro}{\vnbstmt{\begin{array}{@{}l@{}} \forall t, s, c. \\ \quad (\operatorname{rel\text{-}intersection\text{-}intro}(t, s, c) \wedge \operatorname{hyps\text{-}in}(s, \operatorname{derivable}(t))) \Rightarrow  \\ \quad \quad c \in \operatorname{derivable}(t) \end{array}}}
\vnbprop{adm-intersection-elim}{\vnbstmt{\begin{array}{@{}l@{}} \forall t, s, c. \\ \quad (\operatorname{rel\text{-}intersection\text{-}elim}(t, s, c) \wedge \operatorname{hyps\text{-}in}(s, \operatorname{derivable}(t))) \Rightarrow  \\ \quad \quad c \in \operatorname{derivable}(t) \end{array}}}
\vnbprop{adm-tuples-intro}{\vnbstmt{\begin{array}{@{}l@{}} \forall t, s, c. \\ \quad (\operatorname{rel\text{-}tuples\text{-}intro}(t, s, c) \wedge \operatorname{hyps\text{-}in}(s, \operatorname{derivable}(t))) \Rightarrow (c \in \operatorname{derivable}(t)) \end{array}}}
\vnbprop{adm-tuples-intro-0}{\vnbstmt{\begin{array}{@{}l@{}} \forall t, s, c. \\ \quad (\operatorname{rel\text{-}tuples\text{-}intro\text{-}0}(t, s, c) \wedge \operatorname{hyps\text{-}in}(s, \operatorname{derivable}(t))) \Rightarrow  \\ \quad \quad c \in \operatorname{derivable}(t) \end{array}}}
\vnbprop{adm-tuples-elim}{\vnbstmt{\begin{array}{@{}l@{}} \forall t, s, c. \\ \quad (\operatorname{rel\text{-}tuples\text{-}elim}(t, s, c) \wedge \operatorname{hyps\text{-}in}(s, \operatorname{derivable}(t))) \Rightarrow (c \in \operatorname{derivable}(t)) \end{array}}}

\vnbbookfile{structure-library/meta-cr}{structure-library/meta-cr}
\noindent{\small meta-cr.scm -- CONTEXTUAL REWRITING (CR) over the encoded core: the one-step rewrite relation with its local contexts, the equivalence it is to be proved inside, and the kernel relations of the seven rewriting rules, transcribed. (docs/meta-k2-admissibility-2026-10-08.md, sections 4.5 and 6, phase 3.)}\par
\par\medskip\noindent\textbf{Definition} (\texttt{rewr}).\;$\operatorname{rewr}(\operatorname{mrt\_}, \operatorname{mrk\_}) :=$
\begin{equation*}
\{\operatorname{mrx\_} \in \operatorname{rw\text{-}univ} : \forall \operatorname{mrs\_} \in \operatorname{power}(\operatorname{rw\text{-}univ}).\; (\operatorname{rw\text{-}closed}(\operatorname{mrt\_}, \operatorname{mrk\_}, \operatorname{mrs\_}) \Rightarrow (\operatorname{mrx\_} \in \operatorname{mrs\_}))\}
\end{equation*}

\par\medskip\noindent\textbf{Definition} (\texttt{eqv}).\;$\operatorname{eqv}(\operatorname{mrt\_}) :=$ $\{\operatorname{mrx\_} \in \operatorname{rw\text{-}univ} : \operatorname{eqv\text{-}ok}(\operatorname{mrt\_}, \operatorname{mrx\_})\}$
\par\medskip\noindent\textbf{Definition} (\texttt{rw-univ}).\; $(\operatorname{rw\text{-}univ} \simeq (\mathbb{N} \times \operatorname{ctx\text{-}set}))$
\par\medskip\noindent\textbf{Definition} (\texttt{eqv-ok}).\;For all $\operatorname{mrt\_}$ and $\operatorname{mrx\_}$, 
\begin{equation*}
\begin{array}{@{}l@{}} \operatorname{eqv\text{-}ok}(\operatorname{mrt\_}, \operatorname{mrx\_}) \\ \Leftrightarrow\; \exists \operatorname{mrg\_}, \operatorname{mre\_}, \operatorname{mrf\_}. \\ \quad \operatorname{mrx\_} = {[}2, \operatorname{mrg\_}, \operatorname{mre\_}, \operatorname{mrf\_}] \wedge \\ \quad \operatorname{mre\_} \in \mathit{clf} \wedge \\ \quad \operatorname{mrf\_} \in \mathit{clf} \wedge \\ \quad {[}\operatorname{mrg\_}, \operatorname{enc\text{-}iff}(\operatorname{mre\_}, \operatorname{mrf\_})] \in \operatorname{derivable}(\operatorname{mrt\_}) \vee \\ \quad \operatorname{mrx\_} = {[}1, \operatorname{mrg\_}, \operatorname{mre\_}, \operatorname{mrf\_}] \wedge \\ \quad \operatorname{mre\_} \in \mathit{clt} \wedge \\ \quad \operatorname{mrf\_} \in \mathit{clt} \wedge \\ \quad {[}\operatorname{mrg\_}, \operatorname{enc\text{-}qeq}(\operatorname{mre\_}, \operatorname{mrf\_})] \in \operatorname{derivable}(\operatorname{mrt\_}) \end{array}
\end{equation*}

\par\medskip\noindent\textbf{Definition} (\texttt{rw-rule}).\;For all $\operatorname{mrt\_}$, $\operatorname{mrk\_}$, $\operatorname{mrs\_}$ and $\operatorname{mrx\_}$, 
\begin{equation*}
\begin{array}{@{}l@{}} \operatorname{rw\text{-}rule}(\operatorname{mrt\_}, \operatorname{mrk\_}, \operatorname{mrs\_}, \operatorname{mrx\_}) \\ \Leftrightarrow\; (1 \leq \operatorname{mrk\_}) \wedge \operatorname{rw\text{-}top}(\operatorname{mrt\_}, \operatorname{mrs\_}, \operatorname{mrx\_}) \vee \\ (1 \leq \operatorname{mrk\_}) \wedge \operatorname{rw\text{-}not}(\operatorname{mrt\_}, \operatorname{mrs\_}, \operatorname{mrx\_}) \vee \\ (1 \leq \operatorname{mrk\_}) \wedge \operatorname{rw\text{-}and\text{-}l}(\operatorname{mrt\_}, \operatorname{mrs\_}, \operatorname{mrx\_}) \vee \\ (1 \leq \operatorname{mrk\_}) \wedge \operatorname{rw\text{-}and\text{-}r}(\operatorname{mrt\_}, \operatorname{mrs\_}, \operatorname{mrx\_}) \vee \\ (1 \leq \operatorname{mrk\_}) \wedge \operatorname{rw\text{-}or\text{-}l}(\operatorname{mrt\_}, \operatorname{mrs\_}, \operatorname{mrx\_}) \vee \\ (1 \leq \operatorname{mrk\_}) \wedge \operatorname{rw\text{-}or\text{-}r}(\operatorname{mrt\_}, \operatorname{mrs\_}, \operatorname{mrx\_}) \vee \\ (1 \leq \operatorname{mrk\_}) \wedge \operatorname{rw\text{-}imp\text{-}l}(\operatorname{mrt\_}, \operatorname{mrs\_}, \operatorname{mrx\_}) \vee \\ (1 \leq \operatorname{mrk\_}) \wedge \operatorname{rw\text{-}imp\text{-}r}(\operatorname{mrt\_}, \operatorname{mrs\_}, \operatorname{mrx\_}) \vee \\ (1 \leq \operatorname{mrk\_}) \wedge \operatorname{rw\text{-}iff\text{-}l}(\operatorname{mrt\_}, \operatorname{mrs\_}, \operatorname{mrx\_}) \vee \\ (1 \leq \operatorname{mrk\_}) \wedge \operatorname{rw\text{-}iff\text{-}r}(\operatorname{mrt\_}, \operatorname{mrs\_}, \operatorname{mrx\_}) \vee \\ (1 \leq \operatorname{mrk\_}) \wedge \operatorname{rw\text{-}rel\text{-}l}(\operatorname{mrt\_}, \operatorname{mrs\_}, \operatorname{mrx\_}) \vee \\ (1 \leq \operatorname{mrk\_}) \wedge \operatorname{rw\text{-}rel\text{-}r}(\operatorname{mrt\_}, \operatorname{mrs\_}, \operatorname{mrx\_}) \vee \\ (2 \leq \operatorname{mrk\_}) \wedge \operatorname{rw\text{-}all}(\operatorname{mrt\_}, \operatorname{mrs\_}, \operatorname{mrx\_}) \vee \\ (2 \leq \operatorname{mrk\_}) \wedge \operatorname{rw\text{-}ex}(\operatorname{mrt\_}, \operatorname{mrs\_}, \operatorname{mrx\_}) \vee \\ (4 \leq \operatorname{mrk\_}) \wedge \operatorname{rw\text{-}pred}(\operatorname{mrt\_}, \operatorname{mrs\_}, \operatorname{mrx\_}) \vee \\ (4 \leq \operatorname{mrk\_}) \wedge \operatorname{rw\text{-}tapp}(\operatorname{mrt\_}, \operatorname{mrs\_}, \operatorname{mrx\_}) \vee \\ (5 \leq \operatorname{mrk\_}) \wedge \operatorname{rw\text{-}bind\text{-}dom}(\operatorname{mrt\_}, \operatorname{mrs\_}, \operatorname{mrx\_}) \vee \\ (6 \leq \operatorname{mrk\_}) \wedge \operatorname{rw\text{-}sep\text{-}body}(\operatorname{mrt\_}, \operatorname{mrs\_}, \operatorname{mrx\_}) \vee \\ (6 \leq \operatorname{mrk\_}) \wedge \operatorname{rw\text{-}tbind\text{-}body}(\operatorname{mrt\_}, \operatorname{mrs\_}, \operatorname{mrx\_}) \vee \\ (7 \leq \operatorname{mrk\_}) \wedge \operatorname{rw\text{-}desc\text{-}body}(\operatorname{mrt\_}, \operatorname{mrs\_}, \operatorname{mrx\_}) \vee \\ (3 \leq \operatorname{mrk\_}) \wedge \operatorname{rw\text{-}if\text{-}cond}(\operatorname{mrt\_}, \operatorname{mrs\_}, \operatorname{mrx\_}) \vee \\ (3 \leq \operatorname{mrk\_}) \wedge \operatorname{rw\text{-}if\text{-}then}(\operatorname{mrt\_}, \operatorname{mrs\_}, \operatorname{mrx\_}) \vee \\ (3 \leq \operatorname{mrk\_}) \wedge \operatorname{rw\text{-}if\text{-}else}(\operatorname{mrt\_}, \operatorname{mrs\_}, \operatorname{mrx\_}) \end{array}
\end{equation*}

\par\medskip\noindent\textbf{Definition} (\texttt{rw-closed}).\;For all $\operatorname{mrt\_}$, $\operatorname{mrk\_}$ and $\operatorname{mrs\_}$, 
\begin{equation*}
\begin{array}{@{}l@{}} \operatorname{rw\text{-}closed}(\operatorname{mrt\_}, \operatorname{mrk\_}, \operatorname{mrs\_}) \\ \Leftrightarrow\; \forall \operatorname{mrx\_}. \\ \quad \operatorname{rw\text{-}rule}(\operatorname{mrt\_}, \operatorname{mrk\_}, \operatorname{mrs\_}, \operatorname{mrx\_}) \Rightarrow ((\operatorname{mrx\_} \in \operatorname{rw\text{-}univ}) \Rightarrow (\operatorname{mrx\_} \in \operatorname{mrs\_})) \end{array}
\end{equation*}

\par\medskip\noindent\textbf{Definition} (\texttt{rw-top}).\;For all $\operatorname{mrt\_}$, $\operatorname{mrs\_}$ and $\operatorname{mrx\_}$,  $\operatorname{rw\text{-}top}(\operatorname{mrt\_}, \operatorname{mrs\_}, \operatorname{mrx\_}) \Leftrightarrow (\operatorname{mrx\_} \in \operatorname{mrt\_})$
\par\medskip\noindent\textbf{Definition} (\texttt{rw-not}).\;For all $\operatorname{mrt\_}$, $\operatorname{mrs\_}$ and $\operatorname{mrx\_}$, 
\begin{equation*}
\begin{array}{@{}l@{}} \operatorname{rw\text{-}not}(\operatorname{mrt\_}, \operatorname{mrs\_}, \operatorname{mrx\_}) \\ \Leftrightarrow\; \exists \operatorname{mrg\_}, \operatorname{mra\_}, \operatorname{mrb\_}. \\ \quad (\operatorname{mrx\_} = {[}2, \operatorname{mrg\_}, \operatorname{enc\text{-}not}(\operatorname{mra\_}), \operatorname{enc\text{-}not}(\operatorname{mrb\_})]) \wedge ({[}2, \operatorname{mrg\_}, \operatorname{mra\_}, \operatorname{mrb\_}] \in \operatorname{mrs\_}) \end{array}
\end{equation*}

\par\medskip\noindent\textbf{Definition} (\texttt{rw-and-l}).\;For all $\operatorname{mrt\_}$, $\operatorname{mrs\_}$ and $\operatorname{mrx\_}$, 
\begin{equation*}
\begin{array}{@{}l@{}} \operatorname{rw\text{-}and\text{-}l}(\operatorname{mrt\_}, \operatorname{mrs\_}, \operatorname{mrx\_}) \\ \Leftrightarrow\; \exists \operatorname{mrg\_}, \operatorname{mra\_}, \operatorname{mrb\_}, \operatorname{mrc\_}. \\ \quad \operatorname{mrx\_} = {[}2, \operatorname{mrg\_}, \operatorname{enc\text{-}and}(\operatorname{mra\_}, \operatorname{mrc\_}), \operatorname{enc\text{-}and}(\operatorname{mrb\_}, \operatorname{mrc\_})] \wedge \\ \quad \operatorname{mrc\_} \in \mathit{clf} \wedge \\ \quad {[}2, \operatorname{ctx\text{-}add}(\operatorname{mrg\_}, \operatorname{mrc\_}), \operatorname{mra\_}, \operatorname{mrb\_}] \in \operatorname{mrs\_} \end{array}
\end{equation*}

\par\medskip\noindent\textbf{Definition} (\texttt{rw-and-r}).\;For all $\operatorname{mrt\_}$, $\operatorname{mrs\_}$ and $\operatorname{mrx\_}$, 
\begin{equation*}
\begin{array}{@{}l@{}} \operatorname{rw\text{-}and\text{-}r}(\operatorname{mrt\_}, \operatorname{mrs\_}, \operatorname{mrx\_}) \\ \Leftrightarrow\; \exists \operatorname{mrg\_}, \operatorname{mra\_}, \operatorname{mrb\_}, \operatorname{mrc\_}. \\ \quad \operatorname{mrx\_} = {[}2, \operatorname{mrg\_}, \operatorname{enc\text{-}and}(\operatorname{mrc\_}, \operatorname{mra\_}), \operatorname{enc\text{-}and}(\operatorname{mrc\_}, \operatorname{mrb\_})] \wedge \\ \quad \operatorname{mrc\_} \in \mathit{clf} \wedge \\ \quad {[}2, \operatorname{ctx\text{-}add}(\operatorname{mrg\_}, \operatorname{mrc\_}), \operatorname{mra\_}, \operatorname{mrb\_}] \in \operatorname{mrs\_} \end{array}
\end{equation*}

\par\medskip\noindent\textbf{Definition} (\texttt{rw-or-l}).\;For all $\operatorname{mrt\_}$, $\operatorname{mrs\_}$ and $\operatorname{mrx\_}$, 
\begin{equation*}
\begin{array}{@{}l@{}} \operatorname{rw\text{-}or\text{-}l}(\operatorname{mrt\_}, \operatorname{mrs\_}, \operatorname{mrx\_}) \\ \Leftrightarrow\; \exists \operatorname{mrg\_}, \operatorname{mra\_}, \operatorname{mrb\_}, \operatorname{mrc\_}. \\ \quad \operatorname{mrx\_} = {[}2, \operatorname{mrg\_}, \operatorname{enc\text{-}or}(\operatorname{mra\_}, \operatorname{mrc\_}), \operatorname{enc\text{-}or}(\operatorname{mrb\_}, \operatorname{mrc\_})] \wedge \\ \quad \operatorname{mrc\_} \in \mathit{clf} \wedge \\ \quad {[}2, \operatorname{ctx\text{-}add}(\operatorname{mrg\_}, \operatorname{enc\text{-}not}(\operatorname{mrc\_})), \operatorname{mra\_}, \operatorname{mrb\_}] \in \operatorname{mrs\_} \end{array}
\end{equation*}

\par\medskip\noindent\textbf{Definition} (\texttt{rw-or-r}).\;For all $\operatorname{mrt\_}$, $\operatorname{mrs\_}$ and $\operatorname{mrx\_}$, 
\begin{equation*}
\begin{array}{@{}l@{}} \operatorname{rw\text{-}or\text{-}r}(\operatorname{mrt\_}, \operatorname{mrs\_}, \operatorname{mrx\_}) \\ \Leftrightarrow\; \exists \operatorname{mrg\_}, \operatorname{mra\_}, \operatorname{mrb\_}, \operatorname{mrc\_}. \\ \quad \operatorname{mrx\_} = {[}2, \operatorname{mrg\_}, \operatorname{enc\text{-}or}(\operatorname{mrc\_}, \operatorname{mra\_}), \operatorname{enc\text{-}or}(\operatorname{mrc\_}, \operatorname{mrb\_})] \wedge \\ \quad \operatorname{mrc\_} \in \mathit{clf} \wedge \\ \quad {[}2, \operatorname{ctx\text{-}add}(\operatorname{mrg\_}, \operatorname{enc\text{-}not}(\operatorname{mrc\_})), \operatorname{mra\_}, \operatorname{mrb\_}] \in \operatorname{mrs\_} \end{array}
\end{equation*}

\par\medskip\noindent\textbf{Definition} (\texttt{rw-imp-l}).\;For all $\operatorname{mrt\_}$, $\operatorname{mrs\_}$ and $\operatorname{mrx\_}$, 
\begin{equation*}
\begin{array}{@{}l@{}} \operatorname{rw\text{-}imp\text{-}l}(\operatorname{mrt\_}, \operatorname{mrs\_}, \operatorname{mrx\_}) \\ \Leftrightarrow\; \exists \operatorname{mrg\_}, \operatorname{mra\_}, \operatorname{mrb\_}, \operatorname{mrc\_}. \\ \quad \operatorname{mrx\_} = {[}2, \operatorname{mrg\_}, \operatorname{enc\text{-}imp}(\operatorname{mra\_}, \operatorname{mrc\_}), \operatorname{enc\text{-}imp}(\operatorname{mrb\_}, \operatorname{mrc\_})] \wedge \\ \quad \operatorname{mrc\_} \in \mathit{clf} \wedge \\ \quad {[}2, \operatorname{mrg\_}, \operatorname{mra\_}, \operatorname{mrb\_}] \in \operatorname{mrs\_} \end{array}
\end{equation*}

\par\medskip\noindent\textbf{Definition} (\texttt{rw-imp-r}).\;For all $\operatorname{mrt\_}$, $\operatorname{mrs\_}$ and $\operatorname{mrx\_}$, 
\begin{equation*}
\begin{array}{@{}l@{}} \operatorname{rw\text{-}imp\text{-}r}(\operatorname{mrt\_}, \operatorname{mrs\_}, \operatorname{mrx\_}) \\ \Leftrightarrow\; \exists \operatorname{mrg\_}, \operatorname{mra\_}, \operatorname{mrb\_}, \operatorname{mrc\_}. \\ \quad \operatorname{mrx\_} = {[}2, \operatorname{mrg\_}, \operatorname{enc\text{-}imp}(\operatorname{mrc\_}, \operatorname{mra\_}), \operatorname{enc\text{-}imp}(\operatorname{mrc\_}, \operatorname{mrb\_})] \wedge \\ \quad \operatorname{mrc\_} \in \mathit{clf} \wedge \\ \quad {[}2, \operatorname{ctx\text{-}add}(\operatorname{mrg\_}, \operatorname{mrc\_}), \operatorname{mra\_}, \operatorname{mrb\_}] \in \operatorname{mrs\_} \end{array}
\end{equation*}

\par\medskip\noindent\textbf{Definition} (\texttt{rw-iff-l}).\;For all $\operatorname{mrt\_}$, $\operatorname{mrs\_}$ and $\operatorname{mrx\_}$, 
\begin{equation*}
\begin{array}{@{}l@{}} \operatorname{rw\text{-}iff\text{-}l}(\operatorname{mrt\_}, \operatorname{mrs\_}, \operatorname{mrx\_}) \\ \Leftrightarrow\; \exists \operatorname{mrg\_}, \operatorname{mra\_}, \operatorname{mrb\_}, \operatorname{mrc\_}. \\ \quad \operatorname{mrx\_} = {[}2, \operatorname{mrg\_}, \operatorname{enc\text{-}iff}(\operatorname{mra\_}, \operatorname{mrc\_}), \operatorname{enc\text{-}iff}(\operatorname{mrb\_}, \operatorname{mrc\_})] \wedge \\ \quad \operatorname{mrc\_} \in \mathit{clf} \wedge \\ \quad {[}2, \operatorname{mrg\_}, \operatorname{mra\_}, \operatorname{mrb\_}] \in \operatorname{mrs\_} \end{array}
\end{equation*}

\par\medskip\noindent\textbf{Definition} (\texttt{rw-iff-r}).\;For all $\operatorname{mrt\_}$, $\operatorname{mrs\_}$ and $\operatorname{mrx\_}$, 
\begin{equation*}
\begin{array}{@{}l@{}} \operatorname{rw\text{-}iff\text{-}r}(\operatorname{mrt\_}, \operatorname{mrs\_}, \operatorname{mrx\_}) \\ \Leftrightarrow\; \exists \operatorname{mrg\_}, \operatorname{mra\_}, \operatorname{mrb\_}, \operatorname{mrc\_}. \\ \quad \operatorname{mrx\_} = {[}2, \operatorname{mrg\_}, \operatorname{enc\text{-}iff}(\operatorname{mrc\_}, \operatorname{mra\_}), \operatorname{enc\text{-}iff}(\operatorname{mrc\_}, \operatorname{mrb\_})] \wedge \\ \quad \operatorname{mrc\_} \in \mathit{clf} \wedge \\ \quad {[}2, \operatorname{mrg\_}, \operatorname{mra\_}, \operatorname{mrb\_}] \in \operatorname{mrs\_} \end{array}
\end{equation*}

\par\medskip\noindent\textbf{Definition} (\texttt{rw-rel-l}).\;For all $\operatorname{mrt\_}$, $\operatorname{mrs\_}$ and $\operatorname{mrx\_}$, 
\begin{equation*}
\begin{array}{@{}l@{}} \operatorname{rw\text{-}rel\text{-}l}(\operatorname{mrt\_}, \operatorname{mrs\_}, \operatorname{mrx\_}) \\ \Leftrightarrow\; \exists \operatorname{mrg\_}, \operatorname{mrh\_}, \operatorname{mra\_}, \operatorname{mrb\_}, \operatorname{mrc\_}. \\ \quad \operatorname{mrx\_} = {[}2, \operatorname{mrg\_}, \operatorname{eapp}(\operatorname{mrh\_}, {[}\operatorname{mra\_}, \operatorname{mrc\_}]), \operatorname{eapp}(\operatorname{mrh\_}, {[}\operatorname{mrb\_}, \operatorname{mrc\_}])] \wedge \\ \quad \operatorname{mrh\_} \in \mathbb{N} \wedge \\ \quad 7 \leq \operatorname{mrh\_} \wedge \\ \quad \operatorname{mrh\_} \leq 9 \wedge \\ \quad \operatorname{mrc\_} \in \mathit{clt} \wedge \\ \quad {[}1, \operatorname{mrg\_}, \operatorname{mra\_}, \operatorname{mrb\_}] \in \operatorname{mrs\_} \end{array}
\end{equation*}

\par\medskip\noindent\textbf{Definition} (\texttt{rw-rel-r}).\;For all $\operatorname{mrt\_}$, $\operatorname{mrs\_}$ and $\operatorname{mrx\_}$, 
\begin{equation*}
\begin{array}{@{}l@{}} \operatorname{rw\text{-}rel\text{-}r}(\operatorname{mrt\_}, \operatorname{mrs\_}, \operatorname{mrx\_}) \\ \Leftrightarrow\; \exists \operatorname{mrg\_}, \operatorname{mrh\_}, \operatorname{mra\_}, \operatorname{mrb\_}, \operatorname{mrc\_}. \\ \quad \operatorname{mrx\_} = {[}2, \operatorname{mrg\_}, \operatorname{eapp}(\operatorname{mrh\_}, {[}\operatorname{mrc\_}, \operatorname{mra\_}]), \operatorname{eapp}(\operatorname{mrh\_}, {[}\operatorname{mrc\_}, \operatorname{mrb\_}])] \wedge \\ \quad \operatorname{mrh\_} \in \mathbb{N} \wedge \\ \quad 7 \leq \operatorname{mrh\_} \wedge \\ \quad \operatorname{mrh\_} \leq 9 \wedge \\ \quad \operatorname{mrc\_} \in \mathit{clt} \wedge \\ \quad {[}1, \operatorname{mrg\_}, \operatorname{mra\_}, \operatorname{mrb\_}] \in \operatorname{mrs\_} \end{array}
\end{equation*}

\par\medskip\noindent\textbf{Definition} (\texttt{rw-all}).\;For all $\operatorname{mrt\_}$, $\operatorname{mrs\_}$ and $\operatorname{mrx\_}$, 
\begin{equation*}
\begin{array}{@{}l@{}} \operatorname{rw\text{-}all}(\operatorname{mrt\_}, \operatorname{mrs\_}, \operatorname{mrx\_}) \\ \Leftrightarrow\; \exists \operatorname{mrg\_}, \operatorname{mra\_}, \operatorname{mrb\_}. \\ \quad \operatorname{mrx\_} = {[}2, \operatorname{mrg\_}, \operatorname{enc\text{-}all}(\operatorname{mra\_}), \operatorname{enc\text{-}all}(\operatorname{mrb\_})] \wedge \\ \quad \operatorname{enc\text{-}all}(\operatorname{mra\_}) \in \mathit{clf} \wedge \\ \quad \operatorname{enc\text{-}all}(\operatorname{mrb\_}) \in \mathit{clf} \wedge \\ \quad \exists \operatorname{mrfs\_} \in \mathbf{Set}. \\ \quad \quad \operatorname{mrfs\_} \subseteq \mathbb{N} \wedge \\ \quad \quad |\operatorname{mrfs\_}| \in \mathbb{N} \wedge \\ \quad \quad \forall \operatorname{mry\_} \in \mathbb{N}. \\ \quad \quad \quad \neg (\operatorname{mry\_} \in \operatorname{mrfs\_}) \Rightarrow  \\ \quad \quad \quad \quad {[}2, \operatorname{mrg\_}, \operatorname{expr\text{-}open}(\operatorname{mra\_}, 0, \operatorname{efvar}(\operatorname{mry\_})), \operatorname{expr\text{-}open}(\operatorname{mrb\_}, 0, \operatorname{efvar}(\operatorname{mry\_}))] \in \operatorname{mrs\_} \end{array}
\end{equation*}

\par\medskip\noindent\textbf{Definition} (\texttt{rw-ex}).\;For all $\operatorname{mrt\_}$, $\operatorname{mrs\_}$ and $\operatorname{mrx\_}$, 
\begin{equation*}
\begin{array}{@{}l@{}} \operatorname{rw\text{-}ex}(\operatorname{mrt\_}, \operatorname{mrs\_}, \operatorname{mrx\_}) \\ \Leftrightarrow\; \exists \operatorname{mrg\_}, \operatorname{mra\_}, \operatorname{mrb\_}. \\ \quad \operatorname{mrx\_} = {[}2, \operatorname{mrg\_}, \operatorname{enc\text{-}ex}(\operatorname{mra\_}), \operatorname{enc\text{-}ex}(\operatorname{mrb\_})] \wedge \\ \quad \operatorname{enc\text{-}ex}(\operatorname{mra\_}) \in \mathit{clf} \wedge \\ \quad \operatorname{enc\text{-}ex}(\operatorname{mrb\_}) \in \mathit{clf} \wedge \\ \quad \exists \operatorname{mrfs\_} \in \mathbf{Set}. \\ \quad \quad \operatorname{mrfs\_} \subseteq \mathbb{N} \wedge \\ \quad \quad |\operatorname{mrfs\_}| \in \mathbb{N} \wedge \\ \quad \quad \forall \operatorname{mry\_} \in \mathbb{N}. \\ \quad \quad \quad \neg (\operatorname{mry\_} \in \operatorname{mrfs\_}) \Rightarrow  \\ \quad \quad \quad \quad {[}2, \operatorname{mrg\_}, \operatorname{expr\text{-}open}(\operatorname{mra\_}, 0, \operatorname{efvar}(\operatorname{mry\_})), \operatorname{expr\text{-}open}(\operatorname{mrb\_}, 0, \operatorname{efvar}(\operatorname{mry\_}))] \in \operatorname{mrs\_} \end{array}
\end{equation*}

\par\medskip\noindent\textbf{Definition} (\texttt{rw-pred}).\;For all $\operatorname{mrt\_}$, $\operatorname{mrs\_}$ and $\operatorname{mrx\_}$, 
\begin{equation*}
\begin{array}{@{}l@{}} \operatorname{rw\text{-}pred}(\operatorname{mrt\_}, \operatorname{mrs\_}, \operatorname{mrx\_}) \\ \Leftrightarrow\; \exists \operatorname{mrg\_}, \operatorname{mrh\_}, \operatorname{mra\_}, \operatorname{mrb\_}, \operatorname{mri\_}. \\ \quad \operatorname{mrx\_} = {[}2, \operatorname{mrg\_}, \operatorname{eapp}(\operatorname{mrh\_}, \operatorname{mra\_}), \operatorname{eapp}(\operatorname{mrh\_}, \operatorname{mrb\_})] \wedge \\ \quad \operatorname{mrh\_} \in \mathbb{N} \wedge \\ \quad \operatorname{mrh\_} = 10 \wedge \\ \quad \operatorname{eapp}(\operatorname{mrh\_}, \operatorname{mra\_}) \in \mathit{clf} \wedge \\ \quad \operatorname{eapp}(\operatorname{mrh\_}, \operatorname{mrb\_}) \in \mathit{clf} \wedge \\ \quad \operatorname{mra\_} \in \operatorname{tuples}(\operatorname{expr\text{-}set}) \wedge \\ \quad \operatorname{mrb\_} \in \operatorname{tuples}(\operatorname{expr\text{-}set}) \wedge \\ \quad \operatorname{length}(\operatorname{mrb\_}) = \operatorname{length}(\operatorname{mra\_}) \wedge \\ \quad \operatorname{mri\_} \in \mathbb{N} \wedge \\ \quad 1 \leq \operatorname{mri\_} \wedge \\ \quad \operatorname{mri\_} \leq \operatorname{length}(\operatorname{mra\_}) \wedge \\ \quad \forall \operatorname{mrj\_} \in \mathbb{N}. \\ \quad \quad (1 \leq \operatorname{mrj\_}) \Rightarrow  \\ \quad \quad (\operatorname{mrj\_} \leq \operatorname{length}(\operatorname{mra\_})) \Rightarrow  \\ \quad \quad \neg (\operatorname{mrj\_} = \operatorname{mri\_}) \Rightarrow (\operatorname{nth}(\operatorname{mrj\_}, \operatorname{mrb\_}) = \operatorname{nth}(\operatorname{mrj\_}, \operatorname{mra\_})) \wedge \\ \quad {[}1, \operatorname{mrg\_}, \operatorname{nth}(\operatorname{mri\_}, \operatorname{mra\_}), \operatorname{nth}(\operatorname{mri\_}, \operatorname{mrb\_})] \in \operatorname{mrs\_} \end{array}
\end{equation*}

\par\medskip\noindent\textbf{Definition} (\texttt{rw-tapp}).\;For all $\operatorname{mrt\_}$, $\operatorname{mrs\_}$ and $\operatorname{mrx\_}$, 
\begin{equation*}
\begin{array}{@{}l@{}} \operatorname{rw\text{-}tapp}(\operatorname{mrt\_}, \operatorname{mrs\_}, \operatorname{mrx\_}) \\ \Leftrightarrow\; \exists \operatorname{mrg\_}, \operatorname{mrh\_}, \operatorname{mra\_}, \operatorname{mrb\_}, \operatorname{mri\_}. \\ \quad \operatorname{mrx\_} = {[}1, \operatorname{mrg\_}, \operatorname{eapp}(\operatorname{mrh\_}, \operatorname{mra\_}), \operatorname{eapp}(\operatorname{mrh\_}, \operatorname{mrb\_})] \wedge \\ \quad \operatorname{mrh\_} \in \mathbb{N} \wedge \\ \quad 11 \leq \operatorname{mrh\_} \wedge \\ \quad \operatorname{eapp}(\operatorname{mrh\_}, \operatorname{mra\_}) \in \mathit{clt} \wedge \\ \quad \operatorname{eapp}(\operatorname{mrh\_}, \operatorname{mrb\_}) \in \mathit{clt} \wedge \\ \quad \operatorname{mra\_} \in \operatorname{tuples}(\operatorname{expr\text{-}set}) \wedge \\ \quad \operatorname{mrb\_} \in \operatorname{tuples}(\operatorname{expr\text{-}set}) \wedge \\ \quad \operatorname{length}(\operatorname{mrb\_}) = \operatorname{length}(\operatorname{mra\_}) \wedge \\ \quad \operatorname{mri\_} \in \mathbb{N} \wedge \\ \quad 1 \leq \operatorname{mri\_} \wedge \\ \quad \operatorname{mri\_} \leq \operatorname{length}(\operatorname{mra\_}) \wedge \\ \quad \forall \operatorname{mrj\_} \in \mathbb{N}. \\ \quad \quad (1 \leq \operatorname{mrj\_}) \Rightarrow  \\ \quad \quad (\operatorname{mrj\_} \leq \operatorname{length}(\operatorname{mra\_})) \Rightarrow  \\ \quad \quad \neg (\operatorname{mrj\_} = \operatorname{mri\_}) \Rightarrow (\operatorname{nth}(\operatorname{mrj\_}, \operatorname{mrb\_}) = \operatorname{nth}(\operatorname{mrj\_}, \operatorname{mra\_})) \wedge \\ \quad {[}1, \operatorname{mrg\_}, \operatorname{nth}(\operatorname{mri\_}, \operatorname{mra\_}), \operatorname{nth}(\operatorname{mri\_}, \operatorname{mrb\_})] \in \operatorname{mrs\_} \end{array}
\end{equation*}

\par\medskip\noindent\textbf{Definition} (\texttt{rw-bind-dom}).\;For all $\operatorname{mrt\_}$, $\operatorname{mrs\_}$ and $\operatorname{mrx\_}$, 
\begin{equation*}
\begin{array}{@{}l@{}} \operatorname{rw\text{-}bind\text{-}dom}(\operatorname{mrt\_}, \operatorname{mrs\_}, \operatorname{mrx\_}) \\ \Leftrightarrow\; \exists \operatorname{mrg\_}, \operatorname{mrh\_}, \operatorname{mra\_}, \operatorname{mrb\_}, \operatorname{mrc\_}. \\ \quad \operatorname{mrx\_} = {[}1, \operatorname{mrg\_}, \operatorname{ebind}(\operatorname{mrh\_}, 1, {[}\operatorname{mra\_}], \operatorname{mrc\_}), \operatorname{ebind}(\operatorname{mrh\_}, 1, {[}\operatorname{mrb\_}], \operatorname{mrc\_})] \wedge \\ \quad (\operatorname{mrh\_} = 22) \vee ((\operatorname{mrh\_} = 24) \vee (\operatorname{mrh\_} = 26)) \wedge \\ \quad \operatorname{ebind}(\operatorname{mrh\_}, 1, {[}\operatorname{mra\_}], \operatorname{mrc\_}) \in \mathit{clt} \wedge \\ \quad \operatorname{ebind}(\operatorname{mrh\_}, 1, {[}\operatorname{mrb\_}], \operatorname{mrc\_}) \in \mathit{clt} \wedge \\ \quad {[}1, \operatorname{mrg\_}, \operatorname{mra\_}, \operatorname{mrb\_}] \in \operatorname{mrs\_} \end{array}
\end{equation*}

\par\medskip\noindent\textbf{Definition} (\texttt{rw-sep-body}).\;For all $\operatorname{mrt\_}$, $\operatorname{mrs\_}$ and $\operatorname{mrx\_}$, 
\begin{equation*}
\begin{array}{@{}l@{}} \operatorname{rw\text{-}sep\text{-}body}(\operatorname{mrt\_}, \operatorname{mrs\_}, \operatorname{mrx\_}) \\ \Leftrightarrow\; \exists \operatorname{mrg\_}, \operatorname{mra\_}, \operatorname{mrb\_}, \operatorname{mrc\_}. \\ \quad \operatorname{mrx\_} = {[}1, \operatorname{mrg\_}, \operatorname{enc\text{-}sep}(\operatorname{mrc\_}, \operatorname{mra\_}), \operatorname{enc\text{-}sep}(\operatorname{mrc\_}, \operatorname{mrb\_})] \wedge \\ \quad \operatorname{enc\text{-}sep}(\operatorname{mrc\_}, \operatorname{mra\_}) \in \mathit{clt} \wedge \\ \quad \operatorname{enc\text{-}sep}(\operatorname{mrc\_}, \operatorname{mrb\_}) \in \mathit{clt} \wedge \\ \quad \exists \operatorname{mrfs\_} \in \mathbf{Set}. \\ \quad \quad \operatorname{mrfs\_} \subseteq \mathbb{N} \wedge \\ \quad \quad |\operatorname{mrfs\_}| \in \mathbb{N} \wedge \\ \quad \quad \forall \operatorname{mry\_} \in \mathbb{N}. \\ \quad \quad \quad \neg (\operatorname{mry\_} \in \operatorname{mrfs\_}) \Rightarrow  \\ \quad \quad \quad \quad {[}2, \operatorname{mrg\_}, \operatorname{expr\text{-}open}(\operatorname{mra\_}, 0, \operatorname{efvar}(\operatorname{mry\_})), \operatorname{expr\text{-}open}(\operatorname{mrb\_}, 0, \operatorname{efvar}(\operatorname{mry\_}))] \in \operatorname{mrs\_} \end{array}
\end{equation*}

\par\medskip\noindent\textbf{Definition} (\texttt{rw-tbind-body}).\;For all $\operatorname{mrt\_}$, $\operatorname{mrs\_}$ and $\operatorname{mrx\_}$, 
\begin{equation*}
\begin{array}{@{}l@{}} \operatorname{rw\text{-}tbind\text{-}body}(\operatorname{mrt\_}, \operatorname{mrs\_}, \operatorname{mrx\_}) \\ \Leftrightarrow\; \exists \operatorname{mrg\_}, \operatorname{mrh\_}, \operatorname{mra\_}, \operatorname{mrb\_}, \operatorname{mrc\_}. \\ \quad \operatorname{mrx\_} = {[}1, \operatorname{mrg\_}, \operatorname{ebind}(\operatorname{mrh\_}, 1, {[}\operatorname{mrc\_}], \operatorname{mra\_}), \operatorname{ebind}(\operatorname{mrh\_}, 1, {[}\operatorname{mrc\_}], \operatorname{mrb\_})] \wedge \\ \quad (\operatorname{mrh\_} = 24) \vee (\operatorname{mrh\_} = 26) \wedge \\ \quad \operatorname{ebind}(\operatorname{mrh\_}, 1, {[}\operatorname{mrc\_}], \operatorname{mra\_}) \in \mathit{clt} \wedge \\ \quad \operatorname{ebind}(\operatorname{mrh\_}, 1, {[}\operatorname{mrc\_}], \operatorname{mrb\_}) \in \mathit{clt} \wedge \\ \quad \exists \operatorname{mrfs\_} \in \mathbf{Set}. \\ \quad \quad \operatorname{mrfs\_} \subseteq \mathbb{N} \wedge \\ \quad \quad |\operatorname{mrfs\_}| \in \mathbb{N} \wedge \\ \quad \quad \forall \operatorname{mry\_} \in \mathbb{N}. \\ \quad \quad \quad \neg (\operatorname{mry\_} \in \operatorname{mrfs\_}) \Rightarrow  \\ \quad \quad \quad \quad {[}1, \operatorname{mrg\_}, \operatorname{expr\text{-}open}(\operatorname{mra\_}, 0, \operatorname{efvar}(\operatorname{mry\_})), \operatorname{expr\text{-}open}(\operatorname{mrb\_}, 0, \operatorname{efvar}(\operatorname{mry\_}))] \in \operatorname{mrs\_} \end{array}
\end{equation*}

\par\medskip\noindent\textbf{Definition} (\texttt{rw-desc-body}).\;For all $\operatorname{mrt\_}$, $\operatorname{mrs\_}$ and $\operatorname{mrx\_}$, 
\begin{equation*}
\begin{array}{@{}l@{}} \operatorname{rw\text{-}desc\text{-}body}(\operatorname{mrt\_}, \operatorname{mrs\_}, \operatorname{mrx\_}) \\ \Leftrightarrow\; \exists \operatorname{mrg\_}, \operatorname{mrh\_}, \operatorname{mra\_}, \operatorname{mrb\_}. \\ \quad \operatorname{mrx\_} = {[}1, \operatorname{mrg\_}, \operatorname{ebind}(\operatorname{mrh\_}, 1, {[}], \operatorname{mra\_}), \operatorname{ebind}(\operatorname{mrh\_}, 1, {[}], \operatorname{mrb\_})] \wedge \\ \quad (\operatorname{mrh\_} = 23) \vee (\operatorname{mrh\_} = 25) \wedge \\ \quad \operatorname{ebind}(\operatorname{mrh\_}, 1, {[}], \operatorname{mra\_}) \in \mathit{clt} \wedge \\ \quad \operatorname{ebind}(\operatorname{mrh\_}, 1, {[}], \operatorname{mrb\_}) \in \mathit{clt} \wedge \\ \quad \exists \operatorname{mrfs\_} \in \mathbf{Set}. \\ \quad \quad \operatorname{mrfs\_} \subseteq \mathbb{N} \wedge \\ \quad \quad |\operatorname{mrfs\_}| \in \mathbb{N} \wedge \\ \quad \quad \forall \operatorname{mry\_} \in \mathbb{N}. \\ \quad \quad \quad \neg (\operatorname{mry\_} \in \operatorname{mrfs\_}) \Rightarrow  \\ \quad \quad \quad \quad {[}2, \operatorname{mrg\_}, \operatorname{expr\text{-}open}(\operatorname{mra\_}, 0, \operatorname{efvar}(\operatorname{mry\_})), \operatorname{expr\text{-}open}(\operatorname{mrb\_}, 0, \operatorname{efvar}(\operatorname{mry\_}))] \in \operatorname{mrs\_} \end{array}
\end{equation*}

\par\medskip\noindent\textbf{Definition} (\texttt{rw-if-cond}).\;For all $\operatorname{mrt\_}$, $\operatorname{mrs\_}$ and $\operatorname{mrx\_}$, 
\begin{equation*}
\begin{array}{@{}l@{}} \operatorname{rw\text{-}if\text{-}cond}(\operatorname{mrt\_}, \operatorname{mrs\_}, \operatorname{mrx\_}) \\ \Leftrightarrow\; \exists \operatorname{mrg\_}, \operatorname{mra\_}, \operatorname{mrb\_}, \operatorname{mrc\_}, \operatorname{mrd\_}. \\ \quad \operatorname{mrx\_} = {[}1, \operatorname{mrg\_}, \operatorname{enc\text{-}if}(\operatorname{mra\_}, \operatorname{mrc\_}, \operatorname{mrd\_}), \operatorname{enc\text{-}if}(\operatorname{mrb\_}, \operatorname{mrc\_}, \operatorname{mrd\_})] \wedge \\ \quad \operatorname{enc\text{-}if}(\operatorname{mra\_}, \operatorname{mrc\_}, \operatorname{mrd\_}) \in \mathit{clt} \wedge \\ \quad \operatorname{enc\text{-}if}(\operatorname{mrb\_}, \operatorname{mrc\_}, \operatorname{mrd\_}) \in \mathit{clt} \wedge \\ \quad \operatorname{mrc\_} \in \mathit{clt} \wedge \\ \quad \operatorname{mrd\_} \in \mathit{clt} \wedge \\ \quad {[}2, \operatorname{mrg\_}, \operatorname{mra\_}, \operatorname{mrb\_}] \in \operatorname{mrs\_} \end{array}
\end{equation*}

\par\medskip\noindent\textbf{Definition} (\texttt{rw-if-then}).\;For all $\operatorname{mrt\_}$, $\operatorname{mrs\_}$ and $\operatorname{mrx\_}$, 
\begin{equation*}
\begin{array}{@{}l@{}} \operatorname{rw\text{-}if\text{-}then}(\operatorname{mrt\_}, \operatorname{mrs\_}, \operatorname{mrx\_}) \\ \Leftrightarrow\; \exists \operatorname{mrg\_}, \operatorname{mra\_}, \operatorname{mrb\_}, \operatorname{mrc\_}, \operatorname{mrd\_}. \\ \quad \operatorname{mrx\_} = {[}1, \operatorname{mrg\_}, \operatorname{enc\text{-}if}(\operatorname{mrc\_}, \operatorname{mra\_}, \operatorname{mrd\_}), \operatorname{enc\text{-}if}(\operatorname{mrc\_}, \operatorname{mrb\_}, \operatorname{mrd\_})] \wedge \\ \quad \operatorname{enc\text{-}if}(\operatorname{mrc\_}, \operatorname{mra\_}, \operatorname{mrd\_}) \in \mathit{clt} \wedge \\ \quad \operatorname{enc\text{-}if}(\operatorname{mrc\_}, \operatorname{mrb\_}, \operatorname{mrd\_}) \in \mathit{clt} \wedge \\ \quad \operatorname{mrc\_} \in \mathit{clf} \wedge \\ \quad \operatorname{mrd\_} \in \mathit{clt} \wedge \\ \quad {[}1, \operatorname{mrg\_}, \operatorname{mra\_}, \operatorname{mrb\_}] \in \operatorname{mrs\_} \end{array}
\end{equation*}

\par\medskip\noindent\textbf{Definition} (\texttt{rw-if-else}).\;For all $\operatorname{mrt\_}$, $\operatorname{mrs\_}$ and $\operatorname{mrx\_}$, 
\begin{equation*}
\begin{array}{@{}l@{}} \operatorname{rw\text{-}if\text{-}else}(\operatorname{mrt\_}, \operatorname{mrs\_}, \operatorname{mrx\_}) \\ \Leftrightarrow\; \exists \operatorname{mrg\_}, \operatorname{mra\_}, \operatorname{mrb\_}, \operatorname{mrc\_}, \operatorname{mrd\_}. \\ \quad \operatorname{mrx\_} = {[}1, \operatorname{mrg\_}, \operatorname{enc\text{-}if}(\operatorname{mrc\_}, \operatorname{mrd\_}, \operatorname{mra\_}), \operatorname{enc\text{-}if}(\operatorname{mrc\_}, \operatorname{mrd\_}, \operatorname{mrb\_})] \wedge \\ \quad \operatorname{enc\text{-}if}(\operatorname{mrc\_}, \operatorname{mrd\_}, \operatorname{mra\_}) \in \mathit{clt} \wedge \\ \quad \operatorname{enc\text{-}if}(\operatorname{mrc\_}, \operatorname{mrd\_}, \operatorname{mrb\_}) \in \mathit{clt} \wedge \\ \quad \operatorname{mrc\_} \in \mathit{clf} \wedge \\ \quad \operatorname{mrd\_} \in \mathit{clt} \wedge \\ \quad {[}1, \operatorname{mrg\_}, \operatorname{mra\_}, \operatorname{mrb\_}] \in \operatorname{mrs\_} \end{array}
\end{equation*}

\par\medskip\noindent\textbf{Definition} (\texttt{rel-rewrite-1}).\;For all $\operatorname{mrt\_}$, $\operatorname{mrhs\_}$ and $\operatorname{mrc\_}$, 
\begin{equation*}
\begin{array}{@{}l@{}} \operatorname{rel\text{-}rewrite\text{-}1}(\operatorname{mrt\_}, \operatorname{mrhs\_}, \operatorname{mrc\_}) \\ \Leftrightarrow\; \operatorname{mrc\_} \in \operatorname{seq\text{-}set} \wedge \\ \exists \operatorname{mrg\_}, \operatorname{mra\_}, \operatorname{mrb\_}. \\ \quad ((\operatorname{mrc\_} = {[}\operatorname{mrg\_}, \operatorname{mra\_}]) \wedge ((\operatorname{mrhs\_} = {[}{[}\operatorname{mrg\_}, \operatorname{mrb\_}]]) \wedge ({[}2, \operatorname{mrg\_}, \operatorname{mra\_}, \operatorname{mrb\_}] \in \operatorname{rewr}(\operatorname{mrt\_}, 5)))) \end{array}
\end{equation*}

\par\medskip\noindent\textbf{Definition} (\texttt{rel-rewrite-hyp-1}).\;For all $\operatorname{mrt\_}$, $\operatorname{mrhs\_}$ and $\operatorname{mrc\_}$, 
\begin{equation*}
\begin{array}{@{}l@{}} \operatorname{rel\text{-}rewrite\text{-}hyp\text{-}1}(\operatorname{mrt\_}, \operatorname{mrhs\_}, \operatorname{mrc\_}) \\ \Leftrightarrow\; \operatorname{mrc\_} \in \operatorname{seq\text{-}set} \wedge \\ \exists \operatorname{mrg\_}, \operatorname{mrf\_}, \operatorname{mra\_}, \operatorname{mrb\_}. \\ \quad \operatorname{mrc\_} = {[}\operatorname{mrg\_}, \operatorname{mrf\_}] \wedge \\ \quad \operatorname{mra\_} \in \operatorname{mrg\_} \wedge \\ \quad \operatorname{mrhs\_} = {[}{[}\operatorname{ctx\text{-}add}(\operatorname{difference}(\operatorname{mrg\_}, \{\operatorname{mra\_}\}), \operatorname{mrb\_}), \operatorname{mrf\_}]] \wedge \\ \quad {[}2, \operatorname{mrg\_}, \operatorname{mra\_}, \operatorname{mrb\_}] \in \operatorname{rewr}(\operatorname{mrt\_}, 5) \end{array}
\end{equation*}

\par\medskip\noindent\textbf{Definition} (\texttt{rel-rewrite-7}).\;For all $\operatorname{mrt\_}$, $\operatorname{mrhs\_}$ and $\operatorname{mrc\_}$, 
\begin{equation*}
\begin{array}{@{}l@{}} \operatorname{rel\text{-}rewrite\text{-}7}(\operatorname{mrt\_}, \operatorname{mrhs\_}, \operatorname{mrc\_}) \\ \Leftrightarrow\; \operatorname{mrc\_} \in \operatorname{seq\text{-}set} \wedge \\ \exists \operatorname{mrg\_}, \operatorname{mra\_}, \operatorname{mrb\_}. \\ \quad ((\operatorname{mrc\_} = {[}\operatorname{mrg\_}, \operatorname{mra\_}]) \wedge ((\operatorname{mrhs\_} = {[}{[}\operatorname{mrg\_}, \operatorname{mrb\_}]]) \wedge ({[}2, \operatorname{mrg\_}, \operatorname{mra\_}, \operatorname{mrb\_}] \in \operatorname{rewr}(\operatorname{mrt\_}, 7)))) \end{array}
\end{equation*}

\par\medskip\noindent\textbf{Definition} (\texttt{j-nth}).\;$\operatorname{j\text{-}nth}(\operatorname{mrt\_}) :=$
\begin{equation*}
\{\operatorname{mrx\_} \in \operatorname{rw\text{-}univ} : \exists \operatorname{mrg\_}, \operatorname{mrl\_} \in \mathbf{Set}, \operatorname{mrk\_}.\; ((\operatorname{mrx\_} = {[}1, \operatorname{mrg\_}, \operatorname{enc\text{-}nth}(\operatorname{enc\text{-}num}(\operatorname{mrk\_}), \operatorname{enc\text{-}list}(\operatorname{mrl\_})), \operatorname{nth}(\operatorname{mrk\_}, \operatorname{mrl\_})]) \wedge ((\operatorname{mrk\_} \in \mathbb{N}) \wedge ((1 \leq \operatorname{mrk\_}) \wedge ((\operatorname{mrk\_} \leq \operatorname{length}(\operatorname{mrl\_})) \wedge ((\operatorname{enc\text{-}nth}(\operatorname{enc\text{-}num}(\operatorname{mrk\_}), \operatorname{enc\text{-}list}(\operatorname{mrl\_})) \in \mathit{clt}) \wedge ((\operatorname{nth}(\operatorname{mrk\_}, \operatorname{mrl\_}) \in \mathit{clt}) \wedge ({[}\operatorname{mrg\_}, \operatorname{enc\text{-}eq}(\operatorname{enc\text{-}list}(\operatorname{mrl\_}), \operatorname{enc\text{-}list}(\operatorname{mrl\_}))] \in \operatorname{derivable}(\operatorname{mrt\_}))))))))\}
\end{equation*}

\par\medskip\noindent\textbf{Definition} (\texttt{j-len}).\;$\operatorname{j\text{-}len}(\operatorname{mrt\_}) :=$
\begin{equation*}
\{\operatorname{mrx\_} \in \operatorname{rw\text{-}univ} : \exists \operatorname{mrg\_}, \operatorname{mrl\_} \in \mathbf{Set}.\; ((\operatorname{mrx\_} = {[}1, \operatorname{mrg\_}, \operatorname{enc\text{-}length}(\operatorname{enc\text{-}list}(\operatorname{mrl\_})), \operatorname{enc\text{-}num}(\operatorname{length}(\operatorname{mrl\_}))]) \wedge ((\operatorname{enc\text{-}length}(\operatorname{enc\text{-}list}(\operatorname{mrl\_})) \in \mathit{clt}) \wedge ((\operatorname{enc\text{-}num}(\operatorname{length}(\operatorname{mrl\_})) \in \mathit{clt}) \wedge ({[}\operatorname{mrg\_}, \operatorname{enc\text{-}eq}(\operatorname{enc\text{-}list}(\operatorname{mrl\_}), \operatorname{enc\text{-}list}(\operatorname{mrl\_}))] \in \operatorname{derivable}(\operatorname{mrt\_})))))\}
\end{equation*}

\par\medskip\noindent\textbf{Definition} (\texttt{j-lb}).\;$\operatorname{j\text{-}lb}(\operatorname{mrt\_}) :=$
\begin{equation*}
\{\operatorname{mrx\_} \in \operatorname{rw\text{-}univ} : \exists \operatorname{mrg\_}, \operatorname{mra\_}, \operatorname{mrb\_}, \operatorname{mre\_} \in \mathbf{Set}.\; ((\operatorname{mrx\_} = {[}1, \operatorname{mrg\_}, \operatorname{enc\text{-}apply}(\operatorname{enc\text{-}lam}(\operatorname{mra\_}, \operatorname{mrb\_}), \operatorname{mre\_}), \operatorname{expr\text{-}open}(\operatorname{mrb\_}, 0, \operatorname{mre\_})]) \wedge (\operatorname{nbg\text{-}ok}(\operatorname{mrb\_}) \wedge ((\operatorname{enc\text{-}apply}(\operatorname{enc\text{-}lam}(\operatorname{mra\_}, \operatorname{mrb\_}), \operatorname{mre\_}) \in \mathit{clt}) \wedge ((\operatorname{expr\text{-}open}(\operatorname{mrb\_}, 0, \operatorname{mre\_}) \in \mathit{clt}) \wedge (({[}\operatorname{mrg\_}, \operatorname{enc\text{-}in}(\operatorname{mre\_}, \operatorname{mra\_})] \in \operatorname{derivable}(\operatorname{mrt\_})) \wedge ({[}\operatorname{mrg\_}, \operatorname{enc\text{-}in}(\operatorname{expr\text{-}open}(\operatorname{mrb\_}, 0, \operatorname{mre\_}), \operatorname{enc\text{-}setc})] \in \operatorname{derivable}(\operatorname{mrt\_})))))))\}
\end{equation*}

\par\medskip\noindent\textbf{Definition} (\texttt{j-mac}).\;$\operatorname{j\text{-}mac}(\operatorname{mrt\_}) :=$
\begin{equation*}
\{\operatorname{mrx\_} \in \operatorname{rw\text{-}univ} : \exists \operatorname{mrg\_}, \operatorname{mrc\_}, \operatorname{mra\_}, \operatorname{mrb\_} \in \mathbf{Set}, \operatorname{mre\_} \in \mathit{clt}.\; (({[}\operatorname{mrg\_}, \operatorname{expr\text{-}open}(\operatorname{mrc\_}, 0, \operatorname{mre\_})] \in \operatorname{derivable}(\operatorname{mrt\_})) \wedge (({[}\operatorname{mrg\_}, \operatorname{enc\text{-}eq}(\operatorname{mre\_}, \operatorname{mre\_})] \in \operatorname{derivable}(\operatorname{mrt\_})) \wedge (((\operatorname{mrx\_} = {[}1, \operatorname{mrg\_}, \operatorname{expr\text{-}open}(\operatorname{mra\_}, 0, \operatorname{mre\_}), \operatorname{expr\text{-}open}(\operatorname{mrb\_}, 0, \operatorname{mre\_})]) \wedge ((\operatorname{enc\text{-}all}(\operatorname{enc\text{-}imp}(\operatorname{mrc\_}, \operatorname{enc\text{-}qeq}(\operatorname{mra\_}, \operatorname{mrb\_}))) \in \operatorname{mrt\_}) \wedge ((\operatorname{enc\text{-}all}(\operatorname{enc\text{-}imp}(\operatorname{mrc\_}, \operatorname{enc\text{-}qeq}(\operatorname{mra\_}, \operatorname{mrb\_}))) \in \mathit{clf}) \wedge ((\operatorname{expr\text{-}open}(\operatorname{mra\_}, 0, \operatorname{mre\_}) \in \mathit{clt}) \wedge (\operatorname{expr\text{-}open}(\operatorname{mrb\_}, 0, \operatorname{mre\_}) \in \mathit{clt}))))) \vee ((\operatorname{mrx\_} = {[}2, \operatorname{mrg\_}, \operatorname{expr\text{-}open}(\operatorname{mra\_}, 0, \operatorname{mre\_}), \operatorname{expr\text{-}open}(\operatorname{mrb\_}, 0, \operatorname{mre\_})]) \wedge ((\operatorname{enc\text{-}all}(\operatorname{enc\text{-}imp}(\operatorname{mrc\_}, \operatorname{enc\text{-}iff}(\operatorname{mra\_}, \operatorname{mrb\_}))) \in \operatorname{mrt\_}) \wedge ((\operatorname{enc\text{-}all}(\operatorname{enc\text{-}imp}(\operatorname{mrc\_}, \operatorname{enc\text{-}iff}(\operatorname{mra\_}, \operatorname{mrb\_}))) \in \mathit{clf}) \wedge ((\operatorname{expr\text{-}open}(\operatorname{mra\_}, 0, \operatorname{mre\_}) \in \mathit{clf}) \wedge (\operatorname{expr\text{-}open}(\operatorname{mrb\_}, 0, \operatorname{mre\_}) \in \mathit{clf}))))))))\}
\end{equation*}

\par\medskip\noindent\textbf{Definition} (\texttt{j-fb}).\;$\operatorname{j\text{-}fb}(\operatorname{mrt\_}) :=$
\begin{equation*}
\{\operatorname{mrx\_} \in \operatorname{rw\text{-}univ} : \exists \operatorname{mrg\_}, \operatorname{mrh\_}, \operatorname{mrl\_}, \operatorname{mrb\_} \in \mathbf{Set}.\; ((\operatorname{mrx\_} = {[}1, \operatorname{mrg\_}, \operatorname{eapp}(\operatorname{mrh\_}, \operatorname{mrl\_}), \operatorname{openl}(\operatorname{mrb\_}, \operatorname{mrl\_}, \operatorname{length}(\operatorname{mrl\_}))]) \wedge (\operatorname{fun\text{-}def}(\operatorname{mrh\_}, \operatorname{length}(\operatorname{mrl\_}), \operatorname{mrb\_}) \wedge ((\operatorname{eapp}(\operatorname{mrh\_}, \operatorname{mrl\_}) \in \mathit{clt}) \wedge (\operatorname{openl}(\operatorname{mrb\_}, \operatorname{mrl\_}, \operatorname{length}(\operatorname{mrl\_})) \in \mathit{clt}))))\}
\end{equation*}

\par\medskip\noindent\textbf{Definition} (\texttt{rel-macete-1}).\;For all $\operatorname{mrt\_}$, $\operatorname{mrhs\_}$ and $\operatorname{mrc\_}$, 
\begin{equation*}
\operatorname{rel\text{-}macete\text{-}1}(\operatorname{mrt\_}, \operatorname{mrhs\_}, \operatorname{mrc\_}) \Leftrightarrow \operatorname{rel\text{-}rewrite\text{-}1}(\operatorname{j\text{-}mac}(\operatorname{mrt\_}), \operatorname{mrhs\_}, \operatorname{mrc\_})
\end{equation*}

\par\medskip\noindent\textbf{Definition} (\texttt{rel-functoid-beta-1}).\;For all $\operatorname{mrt\_}$, $\operatorname{mrhs\_}$ and $\operatorname{mrc\_}$, 
\begin{equation*}
\operatorname{rel\text{-}functoid\text{-}beta\text{-}1}(\operatorname{mrt\_}, \operatorname{mrhs\_}, \operatorname{mrc\_}) \Leftrightarrow \operatorname{rel\text{-}rewrite\text{-}1}(\operatorname{j\text{-}fb}(\operatorname{mrt\_}), \operatorname{mrhs\_}, \operatorname{mrc\_})
\end{equation*}

\par\medskip\noindent\textbf{Definition} (\texttt{rel-nth-reduce-1}).\;For all $\operatorname{mrt\_}$, $\operatorname{mrhs\_}$ and $\operatorname{mrc\_}$, 
\begin{equation*}
\operatorname{rel\text{-}nth\text{-}reduce\text{-}1}(\operatorname{mrt\_}, \operatorname{mrhs\_}, \operatorname{mrc\_}) \Leftrightarrow \operatorname{rel\text{-}rewrite\text{-}1}(\operatorname{j\text{-}nth}(\operatorname{mrt\_}), \operatorname{mrhs\_}, \operatorname{mrc\_})
\end{equation*}

\par\medskip\noindent\textbf{Definition} (\texttt{rel-length-reduce-1}).\;For all $\operatorname{mrt\_}$, $\operatorname{mrhs\_}$ and $\operatorname{mrc\_}$, 
\begin{equation*}
\operatorname{rel\text{-}length\text{-}reduce\text{-}1}(\operatorname{mrt\_}, \operatorname{mrhs\_}, \operatorname{mrc\_}) \Leftrightarrow \operatorname{rel\text{-}rewrite\text{-}1}(\operatorname{j\text{-}len}(\operatorname{mrt\_}), \operatorname{mrhs\_}, \operatorname{mrc\_})
\end{equation*}

\par\medskip\noindent\textbf{Definition} (\texttt{rel-lambda-beta-1}).\;For all $\operatorname{mrt\_}$, $\operatorname{mrhs\_}$ and $\operatorname{mrc\_}$, 
\begin{equation*}
\operatorname{rel\text{-}lambda\text{-}beta\text{-}1}(\operatorname{mrt\_}, \operatorname{mrhs\_}, \operatorname{mrc\_}) \Leftrightarrow \operatorname{rel\text{-}rewrite\text{-}1}(\operatorname{j\text{-}lb}(\operatorname{mrt\_}), \operatorname{mrhs\_}, \operatorname{mrc\_})
\end{equation*}

\par\medskip\noindent\textbf{Definition} (\texttt{rel-lambda-beta-hyp-1}).\;For all $\operatorname{mrt\_}$, $\operatorname{mrhs\_}$ and $\operatorname{mrc\_}$, 
\begin{equation*}
\operatorname{rel\text{-}lambda\text{-}beta\text{-}hyp\text{-}1}(\operatorname{mrt\_}, \operatorname{mrhs\_}, \operatorname{mrc\_}) \Leftrightarrow \operatorname{rel\text{-}rewrite\text{-}hyp\text{-}1}(\operatorname{j\text{-}lb}(\operatorname{mrt\_}), \operatorname{mrhs\_}, \operatorname{mrc\_})
\end{equation*}

\par\medskip\noindent\textbf{Definition} (\texttt{rel-macete-hyp-1}).\;For all $\operatorname{mrt\_}$, $\operatorname{mrhs\_}$ and $\operatorname{mrc\_}$, 
\begin{equation*}
\operatorname{rel\text{-}macete\text{-}hyp\text{-}1}(\operatorname{mrt\_}, \operatorname{mrhs\_}, \operatorname{mrc\_}) \Leftrightarrow \operatorname{rel\text{-}rewrite\text{-}hyp\text{-}1}(\operatorname{j\text{-}mac}(\operatorname{mrt\_}), \operatorname{mrhs\_}, \operatorname{mrc\_})
\end{equation*}

\par\medskip\noindent\textbf{Definition} (\texttt{rel-rewrite-hyp-7}).\;For all $\operatorname{mrt\_}$, $\operatorname{mrhs\_}$ and $\operatorname{mrc\_}$, 
\begin{equation*}
\begin{array}{@{}l@{}} \operatorname{rel\text{-}rewrite\text{-}hyp\text{-}7}(\operatorname{mrt\_}, \operatorname{mrhs\_}, \operatorname{mrc\_}) \\ \Leftrightarrow\; \operatorname{mrc\_} \in \operatorname{seq\text{-}set} \wedge \\ \exists \operatorname{mrg\_}, \operatorname{mrf\_}, \operatorname{mra\_}, \operatorname{mrb\_}. \\ \quad \operatorname{mrc\_} = {[}\operatorname{mrg\_}, \operatorname{mrf\_}] \wedge \\ \quad \operatorname{mra\_} \in \operatorname{mrg\_} \wedge \\ \quad \operatorname{mrhs\_} = {[}{[}\operatorname{ctx\text{-}add}(\operatorname{difference}(\operatorname{mrg\_}, \{\operatorname{mra\_}\}), \operatorname{mrb\_}), \operatorname{mrf\_}]] \wedge \\ \quad {[}2, \operatorname{mrg\_}, \operatorname{mra\_}, \operatorname{mrb\_}] \in \operatorname{rewr}(\operatorname{mrt\_}, 7) \end{array}
\end{equation*}

\par\medskip\noindent\textbf{Definition} (\texttt{rel-nth-reduce-7}).\;For all $\operatorname{mrt\_}$, $\operatorname{mrhs\_}$ and $\operatorname{mrc\_}$, 
\begin{equation*}
\operatorname{rel\text{-}nth\text{-}reduce\text{-}7}(\operatorname{mrt\_}, \operatorname{mrhs\_}, \operatorname{mrc\_}) \Leftrightarrow \operatorname{rel\text{-}rewrite\text{-}7}(\operatorname{j\text{-}nth}(\operatorname{mrt\_}), \operatorname{mrhs\_}, \operatorname{mrc\_})
\end{equation*}

\par\medskip\noindent\textbf{Definition} (\texttt{rel-length-reduce-7}).\;For all $\operatorname{mrt\_}$, $\operatorname{mrhs\_}$ and $\operatorname{mrc\_}$, 
\begin{equation*}
\operatorname{rel\text{-}length\text{-}reduce\text{-}7}(\operatorname{mrt\_}, \operatorname{mrhs\_}, \operatorname{mrc\_}) \Leftrightarrow \operatorname{rel\text{-}rewrite\text{-}7}(\operatorname{j\text{-}len}(\operatorname{mrt\_}), \operatorname{mrhs\_}, \operatorname{mrc\_})
\end{equation*}

\par\medskip\noindent\textbf{Definition} (\texttt{rel-lambda-beta-7}).\;For all $\operatorname{mrt\_}$, $\operatorname{mrhs\_}$ and $\operatorname{mrc\_}$, 
\begin{equation*}
\operatorname{rel\text{-}lambda\text{-}beta\text{-}7}(\operatorname{mrt\_}, \operatorname{mrhs\_}, \operatorname{mrc\_}) \Leftrightarrow \operatorname{rel\text{-}rewrite\text{-}7}(\operatorname{j\text{-}lb}(\operatorname{mrt\_}), \operatorname{mrhs\_}, \operatorname{mrc\_})
\end{equation*}

\par\medskip\noindent\textbf{Definition} (\texttt{rel-lambda-beta-hyp-7}).\;For all $\operatorname{mrt\_}$, $\operatorname{mrhs\_}$ and $\operatorname{mrc\_}$, 
\begin{equation*}
\operatorname{rel\text{-}lambda\text{-}beta\text{-}hyp\text{-}7}(\operatorname{mrt\_}, \operatorname{mrhs\_}, \operatorname{mrc\_}) \Leftrightarrow \operatorname{rel\text{-}rewrite\text{-}hyp\text{-}7}(\operatorname{j\text{-}lb}(\operatorname{mrt\_}), \operatorname{mrhs\_}, \operatorname{mrc\_})
\end{equation*}

\par\medskip\noindent\textbf{Definition} (\texttt{rel-macete-7}).\;For all $\operatorname{mrt\_}$, $\operatorname{mrhs\_}$ and $\operatorname{mrc\_}$, 
\begin{equation*}
\operatorname{rel\text{-}macete\text{-}7}(\operatorname{mrt\_}, \operatorname{mrhs\_}, \operatorname{mrc\_}) \Leftrightarrow \operatorname{rel\text{-}rewrite\text{-}7}(\operatorname{j\text{-}mac}(\operatorname{mrt\_}), \operatorname{mrhs\_}, \operatorname{mrc\_})
\end{equation*}

\par\medskip\noindent\textbf{Definition} (\texttt{rel-macete-hyp-7}).\;For all $\operatorname{mrt\_}$, $\operatorname{mrhs\_}$ and $\operatorname{mrc\_}$, 
\begin{equation*}
\operatorname{rel\text{-}macete\text{-}hyp\text{-}7}(\operatorname{mrt\_}, \operatorname{mrhs\_}, \operatorname{mrc\_}) \Leftrightarrow \operatorname{rel\text{-}rewrite\text{-}hyp\text{-}7}(\operatorname{j\text{-}mac}(\operatorname{mrt\_}), \operatorname{mrhs\_}, \operatorname{mrc\_})
\end{equation*}

\par\medskip\noindent\textbf{Definition} (\texttt{rel-functoid-beta-7}).\;For all $\operatorname{mrt\_}$, $\operatorname{mrhs\_}$ and $\operatorname{mrc\_}$, 
\begin{equation*}
\operatorname{rel\text{-}functoid\text{-}beta\text{-}7}(\operatorname{mrt\_}, \operatorname{mrhs\_}, \operatorname{mrc\_}) \Leftrightarrow \operatorname{rel\text{-}rewrite\text{-}7}(\operatorname{j\text{-}fb}(\operatorname{mrt\_}), \operatorname{mrhs\_}, \operatorname{mrc\_})
\end{equation*}


\vnbbookfile{theorem-library/meta-cr-laws}{theorem-library/meta-cr-laws}
\noindent{\small meta-cr-laws.scm -- CONTEXTUAL REWRITING (CR) as an admissibility theorem over the encoded core, and the admissibility of the rewriting rules at one position (structure-library/meta-cr.scm). (docs/meta-k2-admissibility-2026-10-08.md, section 6, phase 3.) Agent META-3, 2026-10-09.  Builds on ...}\par
\vnbprop{cartesian4-set}{\vnbstmt{\begin{array}{@{}l@{}} \forall a_{1}, a_{2}, a_{3}, a_{4}. \\ \quad ((a_{1} \in \mathbf{Set}) \wedge ((a_{2} \in \mathbf{Set}) \wedge ((a_{3} \in \mathbf{Set}) \wedge (a_{4} \in \mathbf{Set})))) \Rightarrow  \\ \quad \quad (a_{1} \times a_{2}) \in \mathbf{Set} \end{array}}}
\vnbprop{rw-univ-is-set}{\vnbstmt{\operatorname{rw\text{-}univ} \in \mathbf{Set}}}
\vnbprop{eqv-unfold}{\vnbstmt{\forall t.\; (\operatorname{eqv}(t) \simeq \{x \in \operatorname{rw\text{-}univ} : \operatorname{eqv\text{-}ok}(t, x)\})}}
\vnbprop{rewr-unfold}{\vnbstmt{\begin{array}{@{}l@{}} \forall t, k. \\ \quad (\operatorname{rewr}(t, k) \simeq \{x \in \operatorname{rw\text{-}univ} : \forall s \in \operatorname{power}(\operatorname{rw\text{-}univ}).\; (\operatorname{rw\text{-}closed}(t, k, s) \Rightarrow (x \in s))\}) \end{array}}}
\vnbprop{eqv-in-power}{\vnbstmt{\forall t.\; (\operatorname{eqv}(t) \in \operatorname{power}(\operatorname{rw\text{-}univ}))}}
\vnbprop{rewr-least}{\vnbstmt{\begin{array}{@{}l@{}} \forall t, k.; \forall s \in \operatorname{power}(\operatorname{rw\text{-}univ}), x \in \operatorname{rewr}(t, k). \\ \quad \operatorname{rw\text{-}closed}(t, k, s) \Rightarrow (x \in s) \end{array}}}
\vnbprop{eqv-f-open}{\vnbstmt{\begin{array}{@{}l@{}} \forall t, g, a, b. \\ \quad ({[}2, g, a, b] \in \operatorname{eqv}(t)) \Rightarrow  \\ \quad \quad (a \in \mathit{clf}) \wedge ((b \in \mathit{clf}) \wedge ({[}g, \operatorname{enc\text{-}iff}(a, b)] \in \operatorname{derivable}(t))) \end{array}}}
\vnbprop{eqv-t-open}{\vnbstmt{\begin{array}{@{}l@{}} \forall t, g, a, b. \\ \quad ({[}1, g, a, b] \in \operatorname{eqv}(t)) \Rightarrow  \\ \quad \quad (a \in \mathit{clt}) \wedge ((b \in \mathit{clt}) \wedge ({[}g, \operatorname{enc\text{-}qeq}(a, b)] \in \operatorname{derivable}(t))) \end{array}}}
\vnbprop{eqv-f-intro}{\vnbstmt{\begin{array}{@{}l@{}} \forall t, g, a, b. \\ \quad ((a \in \mathit{clf}) \wedge ((b \in \mathit{clf}) \wedge ({[}g, \operatorname{enc\text{-}iff}(a, b)] \in \operatorname{derivable}(t)))) \Rightarrow  \\ \quad \quad {[}2, g, a, b] \in \operatorname{eqv}(t) \end{array}}}
\vnbprop{eqv-t-intro}{\vnbstmt{\begin{array}{@{}l@{}} \forall t, g, a, b. \\ \quad ((a \in \mathit{clt}) \wedge ((b \in \mathit{clt}) \wedge ({[}g, \operatorname{enc\text{-}qeq}(a, b)] \in \operatorname{derivable}(t)))) \Rightarrow  \\ \quad \quad {[}1, g, a, b] \in \operatorname{eqv}(t) \end{array}}}
\vnbprop{rw-univ-ctx}{\vnbstmt{\forall k, g, a, b.\; (({[}k, g, a, b] \in \operatorname{rw\text{-}univ}) \Rightarrow (g \in \operatorname{ctx\text{-}set}))}}
\vnbprop{der-iff-r}{\vnbstmt{\begin{array}{@{}l@{}} \forall t, g, a, b. \\ \quad (({[}g, \operatorname{enc\text{-}iff}(a, b)] \in \operatorname{derivable}(t)) \wedge ({[}g, a] \in \operatorname{derivable}(t))) \Rightarrow  \\ \quad \quad {[}g, b] \in \operatorname{derivable}(t) \end{array}}}
\vnbprop{der-iff-l}{\vnbstmt{\begin{array}{@{}l@{}} \forall t, g, a, b. \\ \quad (({[}g, \operatorname{enc\text{-}iff}(a, b)] \in \operatorname{derivable}(t)) \wedge ({[}g, b] \in \operatorname{derivable}(t))) \Rightarrow  \\ \quad \quad {[}g, a] \in \operatorname{derivable}(t) \end{array}}}
\vnbprop{eqv-top}{\vnbstmt{\begin{array}{@{}l@{}} \forall j, t, x. \\ \quad ((j \subseteq \operatorname{eqv}(t)) \wedge (\operatorname{rw\text{-}top}(j, \operatorname{eqv}(t), x) \wedge (x \in \operatorname{rw\text{-}univ}))) \Rightarrow (x \in \operatorname{eqv}(t)) \end{array}}}
\vnbprop{eqv-not}{\vnbstmt{\forall j, t, x.\; ((\operatorname{rw\text{-}not}(j, \operatorname{eqv}(t), x) \wedge (x \in \operatorname{rw\text{-}univ})) \Rightarrow (x \in \operatorname{eqv}(t)))}}
\vnbprop{eqv-and-l}{\vnbstmt{\forall j, t, x.\; ((\operatorname{rw\text{-}and\text{-}l}(j, \operatorname{eqv}(t), x) \wedge (x \in \operatorname{rw\text{-}univ})) \Rightarrow (x \in \operatorname{eqv}(t)))}}
\vnbprop{eqv-and-r}{\vnbstmt{\forall j, t, x.\; ((\operatorname{rw\text{-}and\text{-}r}(j, \operatorname{eqv}(t), x) \wedge (x \in \operatorname{rw\text{-}univ})) \Rightarrow (x \in \operatorname{eqv}(t)))}}
\vnbprop{eqv-or-l}{\vnbstmt{\forall j, t, x.\; ((\operatorname{rw\text{-}or\text{-}l}(j, \operatorname{eqv}(t), x) \wedge (x \in \operatorname{rw\text{-}univ})) \Rightarrow (x \in \operatorname{eqv}(t)))}}
\vnbprop{eqv-or-r}{\vnbstmt{\forall j, t, x.\; ((\operatorname{rw\text{-}or\text{-}r}(j, \operatorname{eqv}(t), x) \wedge (x \in \operatorname{rw\text{-}univ})) \Rightarrow (x \in \operatorname{eqv}(t)))}}
\vnbprop{eqv-imp-l}{\vnbstmt{\forall j, t, x.\; ((\operatorname{rw\text{-}imp\text{-}l}(j, \operatorname{eqv}(t), x) \wedge (x \in \operatorname{rw\text{-}univ})) \Rightarrow (x \in \operatorname{eqv}(t)))}}
\vnbprop{eqv-imp-r}{\vnbstmt{\forall j, t, x.\; ((\operatorname{rw\text{-}imp\text{-}r}(j, \operatorname{eqv}(t), x) \wedge (x \in \operatorname{rw\text{-}univ})) \Rightarrow (x \in \operatorname{eqv}(t)))}}
\vnbprop{eqv-iff-l}{\vnbstmt{\forall j, t, x.\; ((\operatorname{rw\text{-}iff\text{-}l}(j, \operatorname{eqv}(t), x) \wedge (x \in \operatorname{rw\text{-}univ})) \Rightarrow (x \in \operatorname{eqv}(t)))}}
\vnbprop{eqv-iff-r}{\vnbstmt{\forall j, t, x.\; ((\operatorname{rw\text{-}iff\text{-}r}(j, \operatorname{eqv}(t), x) \wedge (x \in \operatorname{rw\text{-}univ})) \Rightarrow (x \in \operatorname{eqv}(t)))}}
\vnbprop{subst-app2}{\vnbstmt{\begin{array}{@{}l@{}} \forall m, u, h, a, b. \\ \quad ((u \in \operatorname{expr\text{-}set}) \wedge ((h \in \mathbb{N}) \wedge ((a \in \operatorname{expr\text{-}set}) \wedge (b \in \operatorname{expr\text{-}set})))) \Rightarrow  \\ \quad \quad \operatorname{expr\text{-}subst}(\operatorname{eapp}(h, {[}a, b]), m, u) = \operatorname{eapp}(h, {[}\operatorname{expr\text{-}subst}(a, m, u), \operatorname{expr\text{-}subst}(b, m, u)]) \end{array}}}
\vnbprop{subst-hole-l}{\vnbstmt{\begin{array}{@{}l@{}} \forall z, s, c, h. \\ \quad ((z \in \mathbb{N}) \wedge ((s \in \operatorname{expr\text{-}set}) \wedge ((c \in \operatorname{expr\text{-}set}) \wedge ((h \in \mathbb{N}) \wedge \neg (z \in \operatorname{expr\text{-}fv}(c)))))) \Rightarrow  \\ \quad \quad \operatorname{expr\text{-}subst}(\operatorname{eapp}(h, {[}\operatorname{efvar}(z), c]), z, s) = \operatorname{eapp}(h, {[}s, c]) \end{array}}}
\vnbprop{subst-hole-r}{\vnbstmt{\begin{array}{@{}l@{}} \forall z, s, c, h. \\ \quad ((z \in \mathbb{N}) \wedge ((s \in \operatorname{expr\text{-}set}) \wedge ((c \in \operatorname{expr\text{-}set}) \wedge ((h \in \mathbb{N}) \wedge \neg (z \in \operatorname{expr\text{-}fv}(c)))))) \Rightarrow  \\ \quad \quad \operatorname{expr\text{-}subst}(\operatorname{eapp}(h, {[}c, \operatorname{efvar}(z)]), z, s) = \operatorname{eapp}(h, {[}c, s]) \end{array}}}
\vnbprop{der-atom-l}{\vnbstmt{\begin{array}{@{}l@{}} \forall a, g, h, s, t, c. \\ \quad h \in \mathbb{N} \wedge \\ \quad s \in \mathit{clt} \wedge \\ \quad t \in \mathit{clt} \wedge \\ \quad c \in \mathit{clt} \wedge \\ \quad \operatorname{eapp}(h, {[}s, c]) \in \mathit{clf} \wedge \\ \quad {[}g, \operatorname{enc\text{-}qeq}(s, t)] \in \operatorname{derivable}(a) \wedge \\ \quad {[}g, \operatorname{eapp}(h, {[}t, c])] \in \operatorname{derivable}(a) \Rightarrow \\ \quad \quad {[}g, \operatorname{eapp}(h, {[}s, c])] \in \operatorname{derivable}(a) \end{array}}}
\vnbprop{der-atom-r}{\vnbstmt{\begin{array}{@{}l@{}} \forall a, g, h, s, t, c. \\ \quad h \in \mathbb{N} \wedge \\ \quad s \in \mathit{clt} \wedge \\ \quad t \in \mathit{clt} \wedge \\ \quad c \in \mathit{clt} \wedge \\ \quad \operatorname{eapp}(h, {[}c, s]) \in \mathit{clf} \wedge \\ \quad {[}g, \operatorname{enc\text{-}qeq}(s, t)] \in \operatorname{derivable}(a) \wedge \\ \quad {[}g, \operatorname{eapp}(h, {[}c, t])] \in \operatorname{derivable}(a) \Rightarrow \\ \quad \quad {[}g, \operatorname{eapp}(h, {[}c, s])] \in \operatorname{derivable}(a) \end{array}}}
\vnbprop{der-qeq-sym}{\vnbstmt{\begin{array}{@{}l@{}} \forall a, g, s, t. \\ \quad ((s \in \mathit{clt}) \wedge ((t \in \mathit{clt}) \wedge ({[}g, \operatorname{enc\text{-}qeq}(s, t)] \in \operatorname{derivable}(a)))) \Rightarrow  \\ \quad \quad {[}g, \operatorname{enc\text{-}qeq}(t, s)] \in \operatorname{derivable}(a) \end{array}}}
\vnbprop{eqv-rel-l}{\vnbstmt{\forall j, t, x.\; ((\operatorname{rw\text{-}rel\text{-}l}(j, \operatorname{eqv}(t), x) \wedge (x \in \operatorname{rw\text{-}univ})) \Rightarrow (x \in \operatorname{eqv}(t)))}}
\vnbprop{eqv-rel-r}{\vnbstmt{\forall j, t, x.\; ((\operatorname{rw\text{-}rel\text{-}r}(j, \operatorname{eqv}(t), x) \wedge (x \in \operatorname{rw\text{-}univ})) \Rightarrow (x \in \operatorname{eqv}(t)))}}
\vnbprop{enc-all-unfold}{\vnbstmt{\forall p.\; (\operatorname{enc\text{-}all}(p) \simeq \operatorname{ebind}(20, 1, {[}], p))}}
\vnbprop{enc-ex-unfold}{\vnbstmt{\forall p.\; (\operatorname{enc\text{-}ex}(p) \simeq \operatorname{ebind}(21, 1, {[}], p))}}
\vnbprop{clt-efvar}{\vnbstmt{\forall n \in \mathbb{N}.\; (\operatorname{efvar}(n) \in \mathit{clt})}}
\vnbprop{clf-eq-intro}{\vnbstmt{\forall x, y.\; (((x \in \mathit{clt}) \wedge (y \in \mathit{clt})) \Rightarrow (\operatorname{enc\text{-}eq}(x, y) \in \mathit{clf}))}}
\vnbprop{ebind-stage-inv}{\vnbstmt{\begin{array}{@{}l@{}} \forall n \in \mathbb{N}, h, c, a, b. \\ \quad (\operatorname{ebind}(h, c, a, b) \in \operatorname{expr\text{-}stage}(n)) \Rightarrow (b \in \operatorname{expr\text{-}set}) \end{array}}}
\vnbprop{ebind-body-expr}{\vnbstmt{\forall h, c, a, b.\; ((\operatorname{ebind}(h, c, a, b) \in \operatorname{expr\text{-}set}) \Rightarrow (b \in \operatorname{expr\text{-}set}))}}
\vnbprop{mr-fresh}{\vnbstmt{\begin{array}{@{}l@{}} \forall s, g, f. \\ \quad ((s \in \mathbf{Set}) \wedge ((s \subseteq \mathbb{N}) \wedge ((|s| \in \mathbb{N}) \wedge ((g \in \operatorname{ctx\text{-}set}) \wedge (f \in \operatorname{expr\text{-}set}))))) \Rightarrow  \\ \quad \quad \exists n \in \mathbb{N}.\; (\neg (n \in s) \wedge \neg (n \in (\operatorname{ctx\text{-}fv}(g) \cup \operatorname{expr\text{-}fv}(f)))) \end{array}}}
\vnbprop{eqv-all}{\vnbstmt{\forall j, t, x.\; ((\operatorname{rw\text{-}all}(j, \operatorname{eqv}(t), x) \wedge (x \in \operatorname{rw\text{-}univ})) \Rightarrow (x \in \operatorname{eqv}(t)))}}
\vnbprop{eqv-ex}{\vnbstmt{\forall j, t, x.\; ((\operatorname{rw\text{-}ex}(j, \operatorname{eqv}(t), x) \wedge (x \in \operatorname{rw\text{-}univ})) \Rightarrow (x \in \operatorname{eqv}(t)))}}
\vnbprop{clf-qeq-intro}{\vnbstmt{\forall x, y.\; (((x \in \mathit{clt}) \wedge (y \in \mathit{clt})) \Rightarrow (\operatorname{enc\text{-}qeq}(x, y) \in \mathit{clf}))}}
\vnbprop{der-qeq-join}{\vnbstmt{\begin{array}{@{}l@{}} \forall t, g, x, y, c. \\ \quad x \in \mathit{clt} \wedge \\ \quad y \in \mathit{clt} \wedge \\ \quad c \in \mathit{clt} \wedge \\ \quad {[}g, \operatorname{enc\text{-}qeq}(x, c)] \in \operatorname{derivable}(t) \wedge \\ \quad {[}g, \operatorname{enc\text{-}qeq}(y, c)] \in \operatorname{derivable}(t) \Rightarrow \\ \quad \quad {[}g, \operatorname{enc\text{-}qeq}(x, y)] \in \operatorname{derivable}(t) \end{array}}}
\vnbprop{eqv-if-cond}{\vnbstmt{\forall j, t, x.\; ((\operatorname{rw\text{-}if\text{-}cond}(j, \operatorname{eqv}(t), x) \wedge (x \in \operatorname{rw\text{-}univ})) \Rightarrow (x \in \operatorname{eqv}(t)))}}
\vnbprop{eqv-if-then}{\vnbstmt{\forall j, t, x.\; ((\operatorname{rw\text{-}if\text{-}then}(j, \operatorname{eqv}(t), x) \wedge (x \in \operatorname{rw\text{-}univ})) \Rightarrow (x \in \operatorname{eqv}(t)))}}
\vnbprop{eqv-if-else}{\vnbstmt{\forall j, t, x.\; ((\operatorname{rw\text{-}if\text{-}else}(j, \operatorname{eqv}(t), x) \wedge (x \in \operatorname{rw\text{-}univ})) \Rightarrow (x \in \operatorname{eqv}(t)))}}
\vnbprop{cr-closed-1}{\vnbstmt{\forall t, j.\; ((j \subseteq \operatorname{eqv}(t)) \Rightarrow \operatorname{rw\text{-}closed}(j, 1, \operatorname{eqv}(t)))}}
\vnbprop{cr-closed-2}{\vnbstmt{\forall t, j.\; ((j \subseteq \operatorname{eqv}(t)) \Rightarrow \operatorname{rw\text{-}closed}(j, 2, \operatorname{eqv}(t)))}}
\vnbprop{cr-closed-3}{\vnbstmt{\forall t, j.\; ((j \subseteq \operatorname{eqv}(t)) \Rightarrow \operatorname{rw\text{-}closed}(j, 3, \operatorname{eqv}(t)))}}
\vnbprop{cr-1}{\vnbstmt{\forall t, j, x.\; (((j \subseteq \operatorname{eqv}(t)) \wedge (x \in \operatorname{rewr}(j, 1))) \Rightarrow (x \in \operatorname{eqv}(t)))}}
\vnbprop{cr-2}{\vnbstmt{\forall t, j, x.\; (((j \subseteq \operatorname{eqv}(t)) \wedge (x \in \operatorname{rewr}(j, 2))) \Rightarrow (x \in \operatorname{eqv}(t)))}}
\vnbprop{cr-3}{\vnbstmt{\forall t, j, x.\; (((j \subseteq \operatorname{eqv}(t)) \wedge (x \in \operatorname{rewr}(j, 3))) \Rightarrow (x \in \operatorname{eqv}(t)))}}
\vnbprop{cr-goal-1}{\vnbstmt{\begin{array}{@{}l@{}} \forall t, j, g, o, o_{2}. \\ \quad ((j \subseteq \operatorname{eqv}(t)) \wedge (({[}2, g, o, o_{2}] \in \operatorname{rewr}(j, 1)) \wedge ({[}g, o_{2}] \in \operatorname{derivable}(t)))) \Rightarrow  \\ \quad \quad {[}g, o] \in \operatorname{derivable}(t) \end{array}}}
\vnbprop{cr-goal-2}{\vnbstmt{\begin{array}{@{}l@{}} \forall t, j, g, o, o_{2}. \\ \quad ((j \subseteq \operatorname{eqv}(t)) \wedge (({[}2, g, o, o_{2}] \in \operatorname{rewr}(j, 2)) \wedge ({[}g, o_{2}] \in \operatorname{derivable}(t)))) \Rightarrow  \\ \quad \quad {[}g, o] \in \operatorname{derivable}(t) \end{array}}}
\vnbprop{cr-goal-3}{\vnbstmt{\begin{array}{@{}l@{}} \forall t, j, g, o, o_{2}. \\ \quad ((j \subseteq \operatorname{eqv}(t)) \wedge (({[}2, g, o, o_{2}] \in \operatorname{rewr}(j, 3)) \wedge ({[}g, o_{2}] \in \operatorname{derivable}(t)))) \Rightarrow  \\ \quad \quad {[}g, o] \in \operatorname{derivable}(t) \end{array}}}
\vnbprop{fv-entry}{\vnbstmt{\begin{array}{@{}l@{}} \forall h, l, k, z. \\ \quad h \in \mathbb{N} \wedge \\ \quad l \in \operatorname{tuples}(\operatorname{expr\text{-}set}) \wedge \\ \quad k \in \mathbb{N} \wedge \\ \quad 1 \leq k \wedge \\ \quad k \leq \operatorname{length}(l) \wedge \\ \quad z \in \operatorname{expr\text{-}fv}(\operatorname{nth}(k, l)) \Rightarrow \\ \quad \quad z \in \operatorname{expr\text{-}fv}(\operatorname{eapp}(h, l)) \end{array}}}
\vnbprop{hole-list}{\vnbstmt{\begin{array}{@{}l@{}} \forall z, l, i. \\ \quad ((z \in \mathbb{N}) \wedge ((l \in \operatorname{tuples}(\operatorname{expr\text{-}set})) \wedge ((i \in \mathbb{N}) \wedge ((1 \leq i) \wedge (i \leq \operatorname{length}(l)))))) \Rightarrow  \\ \quad \quad \exists \mathit{mrlz} \in \operatorname{tuples}(\operatorname{expr\text{-}set}). \\ \quad \quad \quad \operatorname{length}(\mathit{mrlz}) = \operatorname{length}(l) \wedge \\ \quad \quad \quad \operatorname{nth}(i, \mathit{mrlz}) = \operatorname{efvar}(z) \wedge \\ \quad \quad \quad \forall j \in \mathbb{N}. \\ \quad \quad \quad \quad ((1 \leq j) \wedge ((j \leq \operatorname{length}(l)) \wedge \neg (j = i))) \Rightarrow (\operatorname{nth}(j, \mathit{mrlz}) = \operatorname{nth}(j, l)) \end{array}}}
\vnbprop{subst-hole-list}{\vnbstmt{\begin{array}{@{}l@{}} \forall z, h, l, a, b, i. \\ \quad z \in \mathbb{N} \wedge \\ \quad h \in \mathbb{N} \wedge \\ \quad l \in \operatorname{tuples}(\operatorname{expr\text{-}set}) \wedge \\ \quad a \in \operatorname{tuples}(\operatorname{expr\text{-}set}) \wedge \\ \quad b \in \operatorname{tuples}(\operatorname{expr\text{-}set}) \wedge \\ \quad \operatorname{length}(a) = \operatorname{length}(l) \wedge \\ \quad \operatorname{length}(b) = \operatorname{length}(l) \wedge \\ \quad i \in \mathbb{N} \wedge \\ \quad 1 \leq i \wedge \\ \quad i \leq \operatorname{length}(l) \wedge \\ \quad \operatorname{nth}(i, a) = \operatorname{efvar}(z) \wedge \\ \quad \forall j \in \mathbb{N}. \\ \quad \quad ((1 \leq j) \wedge ((j \leq \operatorname{length}(l)) \wedge \neg (j = i))) \Rightarrow (\operatorname{nth}(j, a) = \operatorname{nth}(j, l)) \wedge \\ \quad \forall j \in \mathbb{N}. \\ \quad \quad ((1 \leq j) \wedge ((j \leq \operatorname{length}(l)) \wedge \neg (j = i))) \Rightarrow (\operatorname{nth}(j, b) = \operatorname{nth}(j, l)) \wedge \\ \quad \neg (z \in \operatorname{expr\text{-}fv}(\operatorname{eapp}(h, l))) \Rightarrow \\ \quad \quad \operatorname{expr\text{-}subst}(\operatorname{eapp}(h, a), z, \operatorname{nth}(i, b)) = \operatorname{eapp}(h, b) \end{array}}}
\vnbprop{agree-sym}{\vnbstmt{\begin{array}{@{}l@{}} \forall l, a, i.; \forall j \in \mathbb{N}. \\ \quad \operatorname{length}(a) = \operatorname{length}(l) \wedge \\ \quad \forall j \in \mathbb{N}. \\ \quad \quad ((1 \leq j) \wedge ((j \leq \operatorname{length}(l)) \wedge \neg (j = i))) \Rightarrow (\operatorname{nth}(j, a) = \operatorname{nth}(j, l)) \wedge \\ \quad 1 \leq j \wedge \\ \quad j \leq \operatorname{length}(a) \wedge \\ \quad \neg (j = i) \Rightarrow \\ \quad \quad \operatorname{nth}(j, l) = \operatorname{nth}(j, a) \end{array}}}
\vnbprop{der-list-hole}{\vnbstmt{\begin{array}{@{}l@{}} \forall t, g, h, l, a, i. \\ \quad h \in \mathbb{N} \wedge \\ \quad l \in \operatorname{tuples}(\operatorname{expr\text{-}set}) \wedge \\ \quad a \in \operatorname{tuples}(\operatorname{expr\text{-}set}) \wedge \\ \quad \operatorname{length}(a) = \operatorname{length}(l) \wedge \\ \quad i \in \mathbb{N} \wedge \\ \quad 1 \leq i \wedge \\ \quad i \leq \operatorname{length}(l) \wedge \\ \quad \forall j \in \mathbb{N}. \\ \quad \quad ((1 \leq j) \wedge ((j \leq \operatorname{length}(l)) \wedge \neg (j = i))) \Rightarrow (\operatorname{nth}(j, a) = \operatorname{nth}(j, l)) \wedge \\ \quad \operatorname{eapp}(h, l) \in \mathit{clf} \wedge \\ \quad \operatorname{eapp}(h, a) \in \mathit{clf} \wedge \\ \quad {[}g, \operatorname{enc\text{-}qeq}(\operatorname{nth}(i, a), \operatorname{nth}(i, l))] \in \operatorname{derivable}(t) \wedge \\ \quad {[}g, \operatorname{eapp}(h, l)] \in \operatorname{derivable}(t) \Rightarrow \\ \quad \quad {[}g, \operatorname{eapp}(h, a)] \in \operatorname{derivable}(t) \end{array}}}
\vnbprop{der-tapp}{\vnbstmt{\begin{array}{@{}l@{}} \forall t, g, h, l, a, i. \\ \quad h \in \mathbb{N} \wedge \\ \quad l \in \operatorname{tuples}(\operatorname{expr\text{-}set}) \wedge \\ \quad a \in \operatorname{tuples}(\operatorname{expr\text{-}set}) \wedge \\ \quad \operatorname{length}(a) = \operatorname{length}(l) \wedge \\ \quad i \in \mathbb{N} \wedge \\ \quad 1 \leq i \wedge \\ \quad i \leq \operatorname{length}(l) \wedge \\ \quad \forall j \in \mathbb{N}. \\ \quad \quad ((1 \leq j) \wedge ((j \leq \operatorname{length}(l)) \wedge \neg (j = i))) \Rightarrow (\operatorname{nth}(j, a) = \operatorname{nth}(j, l)) \wedge \\ \quad \operatorname{eapp}(h, l) \in \mathit{clt} \wedge \\ \quad \operatorname{eapp}(h, a) \in \mathit{clt} \wedge \\ \quad {[}g, \operatorname{enc\text{-}qeq}(\operatorname{nth}(i, l), \operatorname{nth}(i, a))] \in \operatorname{derivable}(t) \Rightarrow \\ \quad \quad {[}g, \operatorname{enc\text{-}qeq}(\operatorname{eapp}(h, l), \operatorname{eapp}(h, a))] \in \operatorname{derivable}(t) \end{array}}}
\vnbprop{eqv-tapp}{\vnbstmt{\forall j, t, x.\; ((\operatorname{rw\text{-}tapp}(j, \operatorname{eqv}(t), x) \wedge (x \in \operatorname{rw\text{-}univ})) \Rightarrow (x \in \operatorname{eqv}(t)))}}
\vnbprop{eqv-pred}{\vnbstmt{\forall j, t, x.\; ((\operatorname{rw\text{-}pred}(j, \operatorname{eqv}(t), x) \wedge (x \in \operatorname{rw\text{-}univ})) \Rightarrow (x \in \operatorname{eqv}(t)))}}
\vnbprop{cr-closed-4}{\vnbstmt{\forall t, j.\; ((j \subseteq \operatorname{eqv}(t)) \Rightarrow \operatorname{rw\text{-}closed}(j, 4, \operatorname{eqv}(t)))}}
\vnbprop{cr-4}{\vnbstmt{\forall t, j, x.\; (((j \subseteq \operatorname{eqv}(t)) \wedge (x \in \operatorname{rewr}(j, 4))) \Rightarrow (x \in \operatorname{eqv}(t)))}}
\vnbprop{cr-goal-4}{\vnbstmt{\begin{array}{@{}l@{}} \forall t, j, g, o, o_{2}. \\ \quad ((j \subseteq \operatorname{eqv}(t)) \wedge (({[}2, g, o, o_{2}] \in \operatorname{rewr}(j, 4)) \wedge ({[}g, o_{2}] \in \operatorname{derivable}(t)))) \Rightarrow  \\ \quad \quad {[}g, o] \in \operatorname{derivable}(t) \end{array}}}
\vnbprop{subst-bind1}{\vnbstmt{\begin{array}{@{}l@{}} \forall m, u, h, c, a, b. \\ \quad ((u \in \operatorname{expr\text{-}set}) \wedge ((h \in \mathbb{N}) \wedge ((c \in \mathbb{N}) \wedge ((a \in \operatorname{expr\text{-}set}) \wedge (b \in \operatorname{expr\text{-}set}))))) \Rightarrow  \\ \quad \quad \operatorname{expr\text{-}subst}(\operatorname{ebind}(h, c, {[}a], b), m, u) = \operatorname{ebind}(h, c, {[}\operatorname{expr\text{-}subst}(a, m, u)], \operatorname{expr\text{-}subst}(b, m, u)) \end{array}}}
\vnbprop{fv-bind-body}{\vnbstmt{\begin{array}{@{}l@{}} \forall h, c, a, b, z. \\ \quad ((h \in \mathbb{N}) \wedge ((c \in \mathbb{N}) \wedge ((a \in \operatorname{tuples}(\operatorname{expr\text{-}set})) \wedge ((b \in \operatorname{expr\text{-}set}) \wedge (z \in \operatorname{expr\text{-}fv}(b)))))) \Rightarrow  \\ \quad \quad z \in \operatorname{expr\text{-}fv}(\operatorname{ebind}(h, c, a, b)) \end{array}}}
\vnbprop{der-bind-dom}{\vnbstmt{\begin{array}{@{}l@{}} \forall t, g, h, a, b, c. \\ \quad h \in \mathbb{N} \wedge \\ \quad a \in \mathit{clt} \wedge \\ \quad b \in \mathit{clt} \wedge \\ \quad \operatorname{ebind}(h, 1, {[}a], c) \in \mathit{clt} \wedge \\ \quad \operatorname{ebind}(h, 1, {[}b], c) \in \mathit{clt} \wedge \\ \quad {[}g, \operatorname{enc\text{-}qeq}(a, b)] \in \operatorname{derivable}(t) \Rightarrow \\ \quad \quad {[}g, \operatorname{enc\text{-}qeq}(\operatorname{ebind}(h, 1, {[}a], c), \operatorname{ebind}(h, 1, {[}b], c))] \in \operatorname{derivable}(t) \end{array}}}
\vnbprop{eqv-bind-dom}{\vnbstmt{\begin{array}{@{}l@{}} \forall j, t, x. \\ \quad (\operatorname{rw\text{-}bind\text{-}dom}(j, \operatorname{eqv}(t), x) \wedge (x \in \operatorname{rw\text{-}univ})) \Rightarrow (x \in \operatorname{eqv}(t)) \end{array}}}
\vnbprop{cr-closed-5}{\vnbstmt{\forall t, j.\; ((j \subseteq \operatorname{eqv}(t)) \Rightarrow \operatorname{rw\text{-}closed}(j, 5, \operatorname{eqv}(t)))}}
\vnbprop{cr-5}{\vnbstmt{\forall t, j, x.\; (((j \subseteq \operatorname{eqv}(t)) \wedge (x \in \operatorname{rewr}(j, 5))) \Rightarrow (x \in \operatorname{eqv}(t)))}}
\vnbprop{cr-goal-5}{\vnbstmt{\begin{array}{@{}l@{}} \forall t, j, g, o, o_{2}. \\ \quad ((j \subseteq \operatorname{eqv}(t)) \wedge (({[}2, g, o, o_{2}] \in \operatorname{rewr}(j, 5)) \wedge ({[}g, o_{2}] \in \operatorname{derivable}(t)))) \Rightarrow  \\ \quad \quad {[}g, o] \in \operatorname{derivable}(t) \end{array}}}
\vnbprop{adm-rewrite-1}{\vnbstmt{\begin{array}{@{}l@{}} \forall t, j, s, c. \\ \quad ((j \subseteq \operatorname{eqv}(t)) \wedge (\operatorname{rel\text{-}rewrite\text{-}1}(j, s, c) \wedge \operatorname{hyps\text{-}in}(s, \operatorname{derivable}(t)))) \Rightarrow  \\ \quad \quad c \in \operatorname{derivable}(t) \end{array}}}
\vnbprop{adm-rewrite-hyp-1}{\vnbstmt{\begin{array}{@{}l@{}} \forall t, j, s, c. \\ \quad ((j \subseteq \operatorname{eqv}(t)) \wedge (\operatorname{rel\text{-}rewrite\text{-}hyp\text{-}1}(j, s, c) \wedge \operatorname{hyps\text{-}in}(s, \operatorname{derivable}(t)))) \Rightarrow  \\ \quad \quad c \in \operatorname{derivable}(t) \end{array}}}
\vnbprop{j-nth-sub}{\vnbstmt{\forall t.\; (\operatorname{j\text{-}nth}(t) \subseteq \operatorname{eqv}(t))}}
\vnbprop{j-len-sub}{\vnbstmt{\forall t.\; (\operatorname{j\text{-}len}(t) \subseteq \operatorname{eqv}(t))}}
\vnbprop{j-lb-sub}{\vnbstmt{\forall t.\; (\operatorname{j\text{-}lb}(t) \subseteq \operatorname{eqv}(t))}}
\vnbprop{open-fun-unfold}{\vnbstmt{\begin{array}{@{}l@{}} \forall u. \\ \quad (\operatorname{open\text{-}fun}(u) \simeq (\lambda {[}\operatorname{mxw\_}, \operatorname{mxj\_}] \in (\operatorname{expr\text{-}set} \times \mathbb{N}).\; \operatorname{expr\text{-}open}(\operatorname{mxw\_}, \operatorname{mxj\_}, u))) \end{array}}}
\vnbprop{open-app2}{\vnbstmt{\begin{array}{@{}l@{}} \forall u, k, h, a, b. \\ \quad ((u \in \operatorname{expr\text{-}set}) \wedge ((k \in \mathbb{N}) \wedge ((h \in \mathbb{N}) \wedge ((a \in \operatorname{expr\text{-}set}) \wedge (b \in \operatorname{expr\text{-}set}))))) \Rightarrow  \\ \quad \quad \operatorname{expr\text{-}open}(\operatorname{eapp}(h, {[}a, b]), k, u) = \operatorname{eapp}(h, {[}\operatorname{expr\text{-}open}(a, k, u), \operatorname{expr\text{-}open}(b, k, u)]) \end{array}}}
\vnbprop{j-mac-sub}{\vnbstmt{\forall t.\; (\operatorname{j\text{-}mac}(t) \subseteq \operatorname{eqv}(t))}}
\vnbprop{adm-nth-reduce-1}{\vnbstmt{\begin{array}{@{}l@{}} \forall t, s, c. \\ \quad (\operatorname{rel\text{-}nth\text{-}reduce\text{-}1}(t, s, c) \wedge \operatorname{hyps\text{-}in}(s, \operatorname{derivable}(t))) \Rightarrow (c \in \operatorname{derivable}(t)) \end{array}}}
\vnbprop{adm-length-reduce-1}{\vnbstmt{\begin{array}{@{}l@{}} \forall t, s, c. \\ \quad (\operatorname{rel\text{-}length\text{-}reduce\text{-}1}(t, s, c) \wedge \operatorname{hyps\text{-}in}(s, \operatorname{derivable}(t))) \Rightarrow  \\ \quad \quad c \in \operatorname{derivable}(t) \end{array}}}
\vnbprop{adm-lambda-beta-1}{\vnbstmt{\begin{array}{@{}l@{}} \forall t, s, c. \\ \quad (\operatorname{rel\text{-}lambda\text{-}beta\text{-}1}(t, s, c) \wedge \operatorname{hyps\text{-}in}(s, \operatorname{derivable}(t))) \Rightarrow  \\ \quad \quad c \in \operatorname{derivable}(t) \end{array}}}
\vnbprop{adm-lambda-beta-hyp-1}{\vnbstmt{\begin{array}{@{}l@{}} \forall t, s, c. \\ \quad (\operatorname{rel\text{-}lambda\text{-}beta\text{-}hyp\text{-}1}(t, s, c) \wedge \operatorname{hyps\text{-}in}(s, \operatorname{derivable}(t))) \Rightarrow  \\ \quad \quad c \in \operatorname{derivable}(t) \end{array}}}
\vnbprop{adm-macete-1}{\vnbstmt{\begin{array}{@{}l@{}} \forall t, s, c. \\ \quad (\operatorname{rel\text{-}macete\text{-}1}(t, s, c) \wedge \operatorname{hyps\text{-}in}(s, \operatorname{derivable}(t))) \Rightarrow (c \in \operatorname{derivable}(t)) \end{array}}}
\vnbprop{adm-macete-hyp-1}{\vnbstmt{\begin{array}{@{}l@{}} \forall t, s, c. \\ \quad (\operatorname{rel\text{-}macete\text{-}hyp\text{-}1}(t, s, c) \wedge \operatorname{hyps\text{-}in}(s, \operatorname{derivable}(t))) \Rightarrow (c \in \operatorname{derivable}(t)) \end{array}}}
\vnbprop{enc-iota-unfold}{\vnbstmt{\forall p.\; (\operatorname{enc\text{-}iota}(p) \simeq \operatorname{ebind}(25, 1, {[}], p))}}
\vnbprop{eqv-iota-body}{\vnbstmt{\begin{array}{@{}l@{}} \forall j, t, g, a, b. \\ \quad \operatorname{rw\text{-}desc\text{-}body}(j, \operatorname{eqv}(t), {[}1, g, \operatorname{enc\text{-}iota}(a), \operatorname{enc\text{-}iota}(b)]) \wedge \\ \quad {[}1, g, \operatorname{enc\text{-}iota}(a), \operatorname{enc\text{-}iota}(b)] \in \operatorname{rw\text{-}univ} \Rightarrow \\ \quad \quad {[}1, g, \operatorname{enc\text{-}iota}(a), \operatorname{enc\text{-}iota}(b)] \in \operatorname{eqv}(t) \end{array}}}
\vnbprop{der-eq-qeq}{\vnbstmt{\begin{array}{@{}l@{}} \forall t, g, x, y. \\ \quad ((x \in \mathit{clt}) \wedge ((y \in \mathit{clt}) \wedge ({[}g, \operatorname{enc\text{-}eq}(x, y)] \in \operatorname{derivable}(t)))) \Rightarrow  \\ \quad \quad {[}g, \operatorname{enc\text{-}qeq}(x, y)] \in \operatorname{derivable}(t) \end{array}}}
\vnbprop{eqv-sep-body}{\vnbstmt{\begin{array}{@{}l@{}} \forall j, t, x. \\ \quad (\operatorname{rw\text{-}sep\text{-}body}(j, \operatorname{eqv}(t), x) \wedge (x \in \operatorname{rw\text{-}univ})) \Rightarrow (x \in \operatorname{eqv}(t)) \end{array}}}
\vnbprop{eqv-tbind-body}{\vnbstmt{\begin{array}{@{}l@{}} \forall j, t, x. \\ \quad (\operatorname{rw\text{-}tbind\text{-}body}(j, \operatorname{eqv}(t), x) \wedge (x \in \operatorname{rw\text{-}univ})) \Rightarrow (x \in \operatorname{eqv}(t)) \end{array}}}
\vnbprop{cr-closed-6}{\vnbstmt{\forall t, j.\; ((j \subseteq \operatorname{eqv}(t)) \Rightarrow \operatorname{rw\text{-}closed}(j, 6, \operatorname{eqv}(t)))}}
\vnbprop{cr-6}{\vnbstmt{\forall t, j, x.\; (((j \subseteq \operatorname{eqv}(t)) \wedge (x \in \operatorname{rewr}(j, 6))) \Rightarrow (x \in \operatorname{eqv}(t)))}}
\vnbprop{cr-goal-6}{\vnbstmt{\begin{array}{@{}l@{}} \forall t, j, g, o, o_{2}. \\ \quad ((j \subseteq \operatorname{eqv}(t)) \wedge (({[}2, g, o, o_{2}] \in \operatorname{rewr}(j, 6)) \wedge ({[}g, o_{2}] \in \operatorname{derivable}(t)))) \Rightarrow  \\ \quad \quad {[}g, o] \in \operatorname{derivable}(t) \end{array}}}
\vnbprop{j-fb-sub}{\vnbstmt{\forall t.\; (\operatorname{j\text{-}fb}(t) \subseteq \operatorname{eqv}(t))}}
\vnbprop{adm-functoid-beta-1}{\vnbstmt{\begin{array}{@{}l@{}} \forall t, s, c. \\ \quad (\operatorname{rel\text{-}functoid\text{-}beta\text{-}1}(t, s, c) \wedge \operatorname{hyps\text{-}in}(s, \operatorname{derivable}(t))) \Rightarrow  \\ \quad \quad c \in \operatorname{derivable}(t) \end{array}}}

\vnbbookfile{structure-library/meta-ind}{structure-library/meta-ind}
\noindent{\small meta-ind.scm -- K2 PHASE 4: the kernel relations of the induction rules, the set-formation rules and the structural rules, TRANSCRIBED from their checkers (rule-checkers-logic.scm, rule-checkers-schema.scm, rule-checkers-rewrite.scm of tree 95bf8fe, 2026-10-12), and the encoded catalog axioms the induction rules are ...}\par
\par\medskip\noindent\textbf{Definition} (\texttt{rel-sep-sethood}).\;For all $\operatorname{mit\_}$, $\operatorname{mihs\_}$ and $\operatorname{mic\_}$, 
\begin{equation*}
\begin{array}{@{}l@{}} \operatorname{rel\text{-}sep\text{-}sethood}(\operatorname{mit\_}, \operatorname{mihs\_}, \operatorname{mic\_}) \\ \Leftrightarrow\; \operatorname{mic\_} \in \operatorname{seq\text{-}set} \wedge \\ \exists \operatorname{mig\_}, \operatorname{mia\_}, \operatorname{mip\_}. \\ \quad \operatorname{mic\_} = {[}\operatorname{mig\_}, \operatorname{enc\text{-}in}(\operatorname{enc\text{-}sep}(\operatorname{mia\_}, \operatorname{mip\_}), \operatorname{enc\text{-}setc})] \wedge \\ \quad \operatorname{nbg\text{-}ok}(\operatorname{mip\_}) \wedge \\ \quad \operatorname{mihs\_} = {[}{[}\operatorname{mig\_}, \operatorname{enc\text{-}in}(\operatorname{mia\_}, \operatorname{enc\text{-}setc})]] \end{array}
\end{equation*}

\par\medskip\noindent\textbf{Definition} (\texttt{rel-sep-mem-intro}).\;For all $\operatorname{mit\_}$, $\operatorname{mihs\_}$ and $\operatorname{mic\_}$, 
\begin{equation*}
\begin{array}{@{}l@{}} \operatorname{rel\text{-}sep\text{-}mem\text{-}intro}(\operatorname{mit\_}, \operatorname{mihs\_}, \operatorname{mic\_}) \\ \Leftrightarrow\; \operatorname{mic\_} \in \operatorname{seq\text{-}set} \wedge \\ \exists \operatorname{mig\_}, \operatorname{mia\_}, \operatorname{mip\_}, \operatorname{miy\_}. \\ \quad \operatorname{mic\_} = {[}\operatorname{mig\_}, \operatorname{enc\text{-}in}(\operatorname{miy\_}, \operatorname{enc\text{-}sep}(\operatorname{mia\_}, \operatorname{mip\_}))] \wedge \\ \quad \operatorname{nbg\text{-}ok}(\operatorname{mip\_}) \wedge \\ \quad \operatorname{mihs\_} = {[}{[}\operatorname{mig\_}, \operatorname{enc\text{-}in}(\operatorname{miy\_}, \operatorname{mia\_})], {[}\operatorname{mig\_}, \operatorname{expr\text{-}open}(\operatorname{mip\_}, 0, \operatorname{miy\_})]] \end{array}
\end{equation*}

\par\medskip\noindent\textbf{Definition} (\texttt{rel-sep-mem-elim}).\;For all $\operatorname{mit\_}$, $\operatorname{mihs\_}$ and $\operatorname{mic\_}$, 
\begin{equation*}
\begin{array}{@{}l@{}} \operatorname{rel\text{-}sep\text{-}mem\text{-}elim}(\operatorname{mit\_}, \operatorname{mihs\_}, \operatorname{mic\_}) \\ \Leftrightarrow\; \operatorname{mic\_} \in \operatorname{seq\text{-}set} \wedge \\ \exists \operatorname{mig\_}, \operatorname{mif\_}, \operatorname{mia\_}, \operatorname{mip\_}, \operatorname{miy\_}. \\ \quad \operatorname{mic\_} = {[}\operatorname{mig\_}, \operatorname{mif\_}] \wedge \\ \quad \operatorname{enc\text{-}in}(\operatorname{miy\_}, \operatorname{enc\text{-}sep}(\operatorname{mia\_}, \operatorname{mip\_})) \in \operatorname{mig\_} \wedge \\ \quad \operatorname{nbg\text{-}ok}(\operatorname{mip\_}) \wedge \\ \quad \operatorname{mihs\_} \\ \quad =\; {[}{[}\operatorname{ctx\text{-}add}(\operatorname{ctx\text{-}add}(\operatorname{difference}(\operatorname{mig\_}, \{\operatorname{enc\text{-}in}(\operatorname{miy\_}, \operatorname{enc\text{-}sep}(\operatorname{mia\_}, \operatorname{mip\_}))\}), \operatorname{enc\text{-}in}(\operatorname{miy\_}, \operatorname{mia\_})), \operatorname{expr\text{-}open}(\operatorname{mip\_}, 0, \operatorname{miy\_})), \operatorname{mif\_}]] \end{array}
\end{equation*}

\par\medskip\noindent\textbf{Definition} (\texttt{rel-comp-mem-intro}).\;For all $\operatorname{mit\_}$, $\operatorname{mihs\_}$ and $\operatorname{mic\_}$, 
\begin{equation*}
\begin{array}{@{}l@{}} \operatorname{rel\text{-}comp\text{-}mem\text{-}intro}(\operatorname{mit\_}, \operatorname{mihs\_}, \operatorname{mic\_}) \\ \Leftrightarrow\; \operatorname{mic\_} \in \operatorname{seq\text{-}set} \wedge \\ \exists \operatorname{mig\_}, \operatorname{mip\_}, \operatorname{miy\_}. \\ \quad \operatorname{mic\_} = {[}\operatorname{mig\_}, \operatorname{enc\text{-}in}(\operatorname{miy\_}, \operatorname{enc\text{-}comp}(\operatorname{mip\_}))] \wedge \\ \quad \operatorname{nbg\text{-}ok}(\operatorname{mip\_}) \wedge \\ \quad \operatorname{mihs\_} = {[}{[}\operatorname{mig\_}, \operatorname{enc\text{-}in}(\operatorname{miy\_}, \operatorname{enc\text{-}setc})], {[}\operatorname{mig\_}, \operatorname{expr\text{-}open}(\operatorname{mip\_}, 0, \operatorname{miy\_})]] \end{array}
\end{equation*}

\par\medskip\noindent\textbf{Definition} (\texttt{rel-comp-mem-elim}).\;For all $\operatorname{mit\_}$, $\operatorname{mihs\_}$ and $\operatorname{mic\_}$, 
\begin{equation*}
\begin{array}{@{}l@{}} \operatorname{rel\text{-}comp\text{-}mem\text{-}elim}(\operatorname{mit\_}, \operatorname{mihs\_}, \operatorname{mic\_}) \\ \Leftrightarrow\; \operatorname{mic\_} \in \operatorname{seq\text{-}set} \wedge \\ \exists \operatorname{mig\_}, \operatorname{mif\_}, \operatorname{mip\_}, \operatorname{miy\_}. \\ \quad \operatorname{mic\_} = {[}\operatorname{mig\_}, \operatorname{mif\_}] \wedge \\ \quad \operatorname{enc\text{-}in}(\operatorname{miy\_}, \operatorname{enc\text{-}comp}(\operatorname{mip\_})) \in \operatorname{mig\_} \wedge \\ \quad \operatorname{nbg\text{-}ok}(\operatorname{mip\_}) \wedge \\ \quad \operatorname{mihs\_} \\ \quad =\; {[}{[}\operatorname{ctx\text{-}add}(\operatorname{ctx\text{-}add}(\operatorname{difference}(\operatorname{mig\_}, \{\operatorname{enc\text{-}in}(\operatorname{miy\_}, \operatorname{enc\text{-}comp}(\operatorname{mip\_}))\}), \operatorname{enc\text{-}in}(\operatorname{miy\_}, \operatorname{enc\text{-}setc})), \operatorname{expr\text{-}open}(\operatorname{mip\_}, 0, \operatorname{miy\_})), \operatorname{mif\_}]] \end{array}
\end{equation*}

\par\medskip\noindent\textbf{Definition} (\texttt{rel-big-union-sethood}).\;For all $\operatorname{mit\_}$, $\operatorname{mihs\_}$ and $\operatorname{mic\_}$, 
\begin{equation*}
\begin{array}{@{}l@{}} \operatorname{rel\text{-}big\text{-}union\text{-}sethood}(\operatorname{mit\_}, \operatorname{mihs\_}, \operatorname{mic\_}) \\ \Leftrightarrow\; \operatorname{mic\_} \in \operatorname{seq\text{-}set} \wedge \\ \exists \operatorname{mig\_}, \operatorname{mia\_}, \operatorname{mib\_}. \\ \quad \operatorname{mic\_} = {[}\operatorname{mig\_}, \operatorname{enc\text{-}in}(\operatorname{enc\text{-}bu}(\operatorname{mia\_}, \operatorname{mib\_}), \operatorname{enc\text{-}setc})] \wedge \\ \quad \operatorname{nbg\text{-}ok}(\operatorname{mib\_}) \wedge \\ \quad \operatorname{mihs\_} \\ \quad =\; {[}{[}\operatorname{mig\_}, \operatorname{enc\text{-}in}(\operatorname{mia\_}, \operatorname{enc\text{-}setc})], {[}\operatorname{mig\_}, \operatorname{enc\text{-}all}(\operatorname{enc\text{-}imp}(\operatorname{enc\text{-}in}(\operatorname{ebvar}(0), \operatorname{mia\_}), \operatorname{enc\text{-}in}(\operatorname{mib\_}, \operatorname{enc\text{-}setc})))]] \end{array}
\end{equation*}

\par\medskip\noindent\textbf{Definition} (\texttt{rel-big-union-mem-intro}).\;For all $\operatorname{mit\_}$, $\operatorname{mihs\_}$ and $\operatorname{mic\_}$, 
\begin{equation*}
\begin{array}{@{}l@{}} \operatorname{rel\text{-}big\text{-}union\text{-}mem\text{-}intro}(\operatorname{mit\_}, \operatorname{mihs\_}, \operatorname{mic\_}) \\ \Leftrightarrow\; \operatorname{mic\_} \in \operatorname{seq\text{-}set} \wedge \\ \exists \operatorname{mig\_}, \operatorname{mia\_}, \operatorname{mib\_}, \operatorname{mix\_}, \operatorname{miw\_}. \\ \quad \operatorname{mic\_} = {[}\operatorname{mig\_}, \operatorname{enc\text{-}in}(\operatorname{mix\_}, \operatorname{enc\text{-}bu}(\operatorname{mia\_}, \operatorname{mib\_}))] \wedge \\ \quad \operatorname{nbg\text{-}ok}(\operatorname{mib\_}) \wedge \\ \quad \operatorname{mihs\_} = {[}{[}\operatorname{mig\_}, \operatorname{enc\text{-}in}(\operatorname{miw\_}, \operatorname{mia\_})], {[}\operatorname{mig\_}, \operatorname{enc\text{-}in}(\operatorname{mix\_}, \operatorname{expr\text{-}open}(\operatorname{mib\_}, 0, \operatorname{miw\_}))]] \end{array}
\end{equation*}

\par\medskip\noindent\textbf{Definition} (\texttt{rel-big-union-mem-elim}).\;For all $\operatorname{mit\_}$, $\operatorname{mihs\_}$ and $\operatorname{mic\_}$, 
\begin{equation*}
\begin{array}{@{}l@{}} \operatorname{rel\text{-}big\text{-}union\text{-}mem\text{-}elim}(\operatorname{mit\_}, \operatorname{mihs\_}, \operatorname{mic\_}) \\ \Leftrightarrow\; \operatorname{mic\_} \in \operatorname{seq\text{-}set} \wedge \\ \exists \operatorname{mig\_}, \operatorname{mif\_}, \operatorname{mia\_}, \operatorname{mib\_}, \operatorname{mix\_}, \operatorname{mie\_}. \\ \quad \operatorname{mic\_} = {[}\operatorname{mig\_}, \operatorname{mif\_}] \wedge \\ \quad \operatorname{enc\text{-}in}(\operatorname{mix\_}, \operatorname{enc\text{-}bu}(\operatorname{mia\_}, \operatorname{mib\_})) \in \operatorname{mig\_} \wedge \\ \quad \operatorname{nbg\text{-}ok}(\operatorname{mib\_}) \wedge \\ \quad \operatorname{mie\_} \in \mathbb{N} \wedge \\ \quad \neg (\operatorname{mie\_} \in (\operatorname{ctx\text{-}fv}(\operatorname{mig\_}) \cup \operatorname{expr\text{-}fv}(\operatorname{mif\_}))) \wedge \\ \quad \operatorname{mihs\_} \\ \quad =\; {[}{[}\operatorname{ctx\text{-}add}(\operatorname{ctx\text{-}add}(\operatorname{difference}(\operatorname{mig\_}, \{\operatorname{enc\text{-}in}(\operatorname{mix\_}, \operatorname{enc\text{-}bu}(\operatorname{mia\_}, \operatorname{mib\_}))\}), \operatorname{enc\text{-}in}(\operatorname{efvar}(\operatorname{mie\_}), \operatorname{mia\_})), \operatorname{enc\text{-}in}(\operatorname{mix\_}, \operatorname{expr\text{-}open}(\operatorname{mib\_}, 0, \operatorname{efvar}(\operatorname{mie\_})))), \operatorname{mif\_}]] \end{array}
\end{equation*}

\par\medskip\noindent\textbf{Definition} (\texttt{enc-nnc}).\; $(\operatorname{enc\text{-}nnc} \simeq \operatorname{econst}(3))$
\par\medskip\noindent\textbf{Definition} (\texttt{enc-ordc}).\; $(\operatorname{enc\text{-}ordc} \simeq \operatorname{econst}(5))$
\par\medskip\noindent\textbf{Definition} (\texttt{enc-succ}).\;$\operatorname{enc\text{-}succ}(\operatorname{mip\_}) :=$ $\operatorname{eapp}(30, {[}\operatorname{mip\_}])$
\par\medskip\noindent\textbf{Definition} (\texttt{enc-succo}).\;$\operatorname{enc\text{-}succo}(\operatorname{mip\_}) :=$ $\operatorname{eapp}(31, {[}\operatorname{mip\_}])$
\par\medskip\noindent\textbf{Definition} (\texttt{enc-ordlt}).\;$\operatorname{enc\text{-}ordlt}(\operatorname{mip\_}, \operatorname{miq\_}) :=$ $\operatorname{eapp}(10, {[}\operatorname{econst}(7), \operatorname{mip\_}, \operatorname{miq\_}])$
\par\medskip\noindent\textbf{Definition} (\texttt{enc-limit}).\;$\operatorname{enc\text{-}limit}(\operatorname{mip\_}) :=$ $\operatorname{eapp}(10, {[}\operatorname{econst}(9), \operatorname{mip\_}])$
\par\medskip\noindent\textbf{Definition} (\texttt{ax-nn-ind}).\;
\begin{equation*}
(\operatorname{ax\text{-}nn\text{-}ind} \simeq \operatorname{enc\text{-}all}(\operatorname{enc\text{-}imp}(\operatorname{enc\text{-}and}(\operatorname{enc\text{-}in}(\operatorname{enc\text{-}num}(0), \operatorname{ebvar}(0)), \operatorname{enc\text{-}all}(\operatorname{enc\text{-}imp}(\operatorname{enc\text{-}and}(\operatorname{enc\text{-}in}(\operatorname{ebvar}(0), \operatorname{enc\text{-}nnc}), \operatorname{enc\text{-}in}(\operatorname{ebvar}(0), \operatorname{ebvar}(1))), \operatorname{enc\text{-}in}(\operatorname{enc\text{-}succ}(\operatorname{ebvar}(0)), \operatorname{ebvar}(1))))), \operatorname{enc\text{-}all}(\operatorname{enc\text{-}imp}(\operatorname{enc\text{-}in}(\operatorname{ebvar}(0), \operatorname{enc\text{-}nnc}), \operatorname{enc\text{-}in}(\operatorname{ebvar}(0), \operatorname{ebvar}(1)))))))
\end{equation*}

\par\medskip\noindent\textbf{Definition} (\texttt{ax-tfi}).\;
\begin{equation*}
(\operatorname{ax\text{-}tfi} \simeq \operatorname{enc\text{-}all}(\operatorname{enc\text{-}imp}(\operatorname{enc\text{-}all}(\operatorname{enc\text{-}imp}(\operatorname{enc\text{-}and}(\operatorname{enc\text{-}in}(\operatorname{ebvar}(0), \operatorname{enc\text{-}ordc}), \operatorname{enc\text{-}all}(\operatorname{enc\text{-}imp}(\operatorname{enc\text{-}ordlt}(\operatorname{ebvar}(0), \operatorname{ebvar}(1)), \operatorname{enc\text{-}in}(\operatorname{ebvar}(0), \operatorname{ebvar}(2))))), \operatorname{enc\text{-}in}(\operatorname{ebvar}(0), \operatorname{ebvar}(1)))), \operatorname{enc\text{-}all}(\operatorname{enc\text{-}imp}(\operatorname{enc\text{-}in}(\operatorname{ebvar}(0), \operatorname{enc\text{-}ordc}), \operatorname{enc\text{-}in}(\operatorname{ebvar}(0), \operatorname{ebvar}(1)))))))
\end{equation*}

\par\medskip\noindent\textbf{Definition} (\texttt{ax-mem-set}).\;
\begin{equation*}
(\operatorname{ax\text{-}mem\text{-}set} \simeq \operatorname{enc\text{-}all}(\operatorname{enc\text{-}all}(\operatorname{enc\text{-}imp}(\operatorname{enc\text{-}in}(\operatorname{ebvar}(1), \operatorname{ebvar}(0)), \operatorname{enc\text{-}in}(\operatorname{ebvar}(1), \operatorname{enc\text{-}setc})))))
\end{equation*}

\par\medskip\noindent\textbf{Definition} (\texttt{ax-class-ext}).\;
\begin{equation*}
(\operatorname{ax\text{-}class\text{-}ext} \simeq \operatorname{enc\text{-}all}(\operatorname{enc\text{-}all}(\operatorname{enc\text{-}imp}(\operatorname{enc\text{-}all}(\operatorname{enc\text{-}iff}(\operatorname{enc\text{-}in}(\operatorname{ebvar}(0), \operatorname{ebvar}(2)), \operatorname{enc\text{-}in}(\operatorname{ebvar}(0), \operatorname{ebvar}(1)))), \operatorname{enc\text{-}eq}(\operatorname{ebvar}(1), \operatorname{ebvar}(0))))))
\end{equation*}

\par\medskip\noindent\textbf{Definition} (\texttt{rel-nn-induction}).\;For all $\operatorname{mit\_}$, $\operatorname{mihs\_}$ and $\operatorname{mic\_}$, 
\begin{equation*}
\begin{array}{@{}l@{}} \operatorname{rel\text{-}nn\text{-}induction}(\operatorname{mit\_}, \operatorname{mihs\_}, \operatorname{mic\_}) \\ \Leftrightarrow\; \operatorname{mic\_} \in \operatorname{seq\text{-}set} \wedge \\ \exists \operatorname{mig\_}, \operatorname{mip\_}, \operatorname{mis\_}, \operatorname{mifs\_}. \\ \quad \operatorname{mic\_} = {[}\operatorname{mig\_}, \operatorname{enc\text{-}all}(\operatorname{enc\text{-}imp}(\operatorname{enc\text{-}in}(\operatorname{ebvar}(0), \operatorname{enc\text{-}nnc}), \operatorname{mip\_}))] \wedge \\ \quad \operatorname{mihs\_} = {[}{[}\operatorname{mig\_}, \operatorname{expr\text{-}open}(\operatorname{mip\_}, 0, \operatorname{enc\text{-}num}(0))], {[}\operatorname{mig\_}, \operatorname{enc\text{-}all}(\operatorname{mis\_})]] \wedge \\ \quad \operatorname{mifs\_} \in \mathbf{Set} \wedge \\ \quad \operatorname{mifs\_} \subseteq \mathbb{N} \wedge \\ \quad |\operatorname{mifs\_}| \in \mathbb{N} \wedge \\ \quad \forall \operatorname{miy\_} \in \mathbb{N}. \\ \quad \quad \neg (\operatorname{miy\_} \in \operatorname{mifs\_}) \Rightarrow  \\ \quad \quad \quad \operatorname{expr\text{-}open}(\operatorname{mis\_}, 0, \operatorname{efvar}(\operatorname{miy\_})) \\ \quad \quad \quad =\; \operatorname{enc\text{-}imp}(\operatorname{enc\text{-}in}(\operatorname{efvar}(\operatorname{miy\_}), \operatorname{enc\text{-}nnc}), \\ \quad \quad \quad \quad \operatorname{enc\text{-}imp}(\operatorname{expr\text{-}open}(\operatorname{mip\_}, 0, \operatorname{efvar}(\operatorname{miy\_})), \operatorname{expr\text{-}open}(\operatorname{mip\_}, 0, \operatorname{enc\text{-}succ}(\operatorname{efvar}(\operatorname{miy\_}))))) \end{array}
\end{equation*}

\par\medskip\noindent\textbf{Definition} (\texttt{rel-transfinite-induction}).\;For all $\operatorname{mit\_}$, $\operatorname{mihs\_}$ and $\operatorname{mic\_}$, 
\begin{equation*}
\begin{array}{@{}l@{}} \operatorname{rel\text{-}transfinite\text{-}induction}(\operatorname{mit\_}, \operatorname{mihs\_}, \operatorname{mic\_}) \\ \Leftrightarrow\; \operatorname{mic\_} \in \operatorname{seq\text{-}set} \wedge \\ \exists \operatorname{mig\_}, \operatorname{mip\_}. \\ \quad \operatorname{mic\_} = {[}\operatorname{mig\_}, \operatorname{enc\text{-}all}(\operatorname{enc\text{-}imp}(\operatorname{enc\text{-}in}(\operatorname{ebvar}(0), \operatorname{enc\text{-}ordc}), \operatorname{mip\_}))] \wedge \\ \quad \operatorname{mihs\_} \\ \quad =\; {[}{[}\operatorname{mig\_}, \operatorname{enc\text{-}all}(\operatorname{enc\text{-}imp}(\operatorname{enc\text{-}and}(\operatorname{enc\text{-}in}(\operatorname{ebvar}(0), \operatorname{enc\text{-}ordc}), \operatorname{enc\text{-}all}(\operatorname{enc\text{-}imp}(\operatorname{enc\text{-}ordlt}(\operatorname{ebvar}(0), \operatorname{ebvar}(1)), \operatorname{mip\_}))), \operatorname{mip\_}))]] \end{array}
\end{equation*}

\par\medskip\noindent\textbf{Definition} (\texttt{rel-tfi-3cases}).\;For all $\operatorname{mit\_}$, $\operatorname{mihs\_}$ and $\operatorname{mic\_}$, 
\begin{equation*}
\begin{array}{@{}l@{}} \operatorname{rel\text{-}tfi\text{-}3cases}(\operatorname{mit\_}, \operatorname{mihs\_}, \operatorname{mic\_}) \\ \Leftrightarrow\; \operatorname{mic\_} \in \operatorname{seq\text{-}set} \wedge \\ \exists \operatorname{mig\_}, \operatorname{mip\_}, \operatorname{mis\_}, \operatorname{mifs\_}. \\ \quad \operatorname{mic\_} = {[}\operatorname{mig\_}, \operatorname{enc\text{-}all}(\operatorname{enc\text{-}imp}(\operatorname{enc\text{-}in}(\operatorname{ebvar}(0), \operatorname{enc\text{-}ordc}), \operatorname{mip\_}))] \wedge \\ \quad \operatorname{mihs\_} \\ \quad =\; {[}{[}\operatorname{mig\_}, \operatorname{expr\text{-}open}(\operatorname{mip\_}, 0, \operatorname{enc\text{-}num}(0))], {[}\operatorname{mig\_}, \operatorname{enc\text{-}all}(\operatorname{mis\_})], {[}\operatorname{mig\_}, \operatorname{enc\text{-}all}(\operatorname{enc\text{-}imp}(\operatorname{enc\text{-}and}(\operatorname{enc\text{-}limit}(\operatorname{ebvar}(0)), \operatorname{enc\text{-}all}(\operatorname{enc\text{-}imp}(\operatorname{enc\text{-}ordlt}(\operatorname{ebvar}(0), \operatorname{ebvar}(1)), \operatorname{mip\_}))), \operatorname{mip\_}))]] \wedge \\ \quad \operatorname{mifs\_} \in \mathbf{Set} \wedge \\ \quad \operatorname{mifs\_} \subseteq \mathbb{N} \wedge \\ \quad |\operatorname{mifs\_}| \in \mathbb{N} \wedge \\ \quad \forall \operatorname{miy\_} \in \mathbb{N}. \\ \quad \quad \neg (\operatorname{miy\_} \in \operatorname{mifs\_}) \Rightarrow  \\ \quad \quad \quad \operatorname{expr\text{-}open}(\operatorname{mis\_}, 0, \operatorname{efvar}(\operatorname{miy\_})) \\ \quad \quad \quad =\; \operatorname{enc\text{-}imp}(\operatorname{enc\text{-}and}(\operatorname{enc\text{-}in}(\operatorname{efvar}(\operatorname{miy\_}), \operatorname{enc\text{-}ordc}), \operatorname{expr\text{-}open}(\operatorname{mip\_}, 0, \operatorname{efvar}(\operatorname{miy\_}))), \\ \quad \quad \quad \quad \operatorname{expr\text{-}open}(\operatorname{mip\_}, 0, \operatorname{enc\text{-}succo}(\operatorname{efvar}(\operatorname{miy\_})))) \end{array}
\end{equation*}

\par\medskip\noindent\textbf{Definition} (\texttt{rel-nn-induction-nbg}).\;For all $\operatorname{mit\_}$, $\operatorname{mihs\_}$ and $\operatorname{mic\_}$, 
\begin{equation*}
\begin{array}{@{}l@{}} \operatorname{rel\text{-}nn\text{-}induction\text{-}nbg}(\operatorname{mit\_}, \operatorname{mihs\_}, \operatorname{mic\_}) \\ \Leftrightarrow\; \operatorname{mic\_} \in \operatorname{seq\text{-}set} \wedge \\ \exists \operatorname{mig\_}, \operatorname{mip\_}, \operatorname{mis\_}, \operatorname{mifs\_}. \\ \quad \operatorname{mic\_} = {[}\operatorname{mig\_}, \operatorname{enc\text{-}all}(\operatorname{enc\text{-}imp}(\operatorname{enc\text{-}in}(\operatorname{ebvar}(0), \operatorname{enc\text{-}nnc}), \operatorname{mip\_}))] \wedge \\ \quad \operatorname{nbg\text{-}ok}(\operatorname{mip\_}) \wedge \\ \quad \operatorname{mihs\_} = {[}{[}\operatorname{mig\_}, \operatorname{expr\text{-}open}(\operatorname{mip\_}, 0, \operatorname{enc\text{-}num}(0))], {[}\operatorname{mig\_}, \operatorname{enc\text{-}all}(\operatorname{mis\_})]] \wedge \\ \quad \operatorname{mifs\_} \in \mathbf{Set} \wedge \\ \quad \operatorname{mifs\_} \subseteq \mathbb{N} \wedge \\ \quad |\operatorname{mifs\_}| \in \mathbb{N} \wedge \\ \quad \forall \operatorname{miy\_} \in \mathbb{N}. \\ \quad \quad \neg (\operatorname{miy\_} \in \operatorname{mifs\_}) \Rightarrow  \\ \quad \quad \quad \operatorname{expr\text{-}open}(\operatorname{mis\_}, 0, \operatorname{efvar}(\operatorname{miy\_})) \\ \quad \quad \quad =\; \operatorname{enc\text{-}imp}(\operatorname{enc\text{-}in}(\operatorname{efvar}(\operatorname{miy\_}), \operatorname{enc\text{-}nnc}), \\ \quad \quad \quad \quad \operatorname{enc\text{-}imp}(\operatorname{expr\text{-}open}(\operatorname{mip\_}, 0, \operatorname{efvar}(\operatorname{miy\_})), \operatorname{expr\text{-}open}(\operatorname{mip\_}, 0, \operatorname{enc\text{-}succ}(\operatorname{efvar}(\operatorname{miy\_}))))) \end{array}
\end{equation*}

\par\medskip\noindent\textbf{Definition} (\texttt{rel-transfinite-induction-nbg}).\;For all $\operatorname{mit\_}$, $\operatorname{mihs\_}$ and $\operatorname{mic\_}$, 
\begin{equation*}
\begin{array}{@{}l@{}} \operatorname{rel\text{-}transfinite\text{-}induction\text{-}nbg}(\operatorname{mit\_}, \operatorname{mihs\_}, \operatorname{mic\_}) \\ \Leftrightarrow\; \operatorname{mic\_} \in \operatorname{seq\text{-}set} \wedge \\ \exists \operatorname{mig\_}, \operatorname{mip\_}. \\ \quad \operatorname{mic\_} = {[}\operatorname{mig\_}, \operatorname{enc\text{-}all}(\operatorname{enc\text{-}imp}(\operatorname{enc\text{-}in}(\operatorname{ebvar}(0), \operatorname{enc\text{-}ordc}), \operatorname{mip\_}))] \wedge \\ \quad \operatorname{nbg\text{-}ok}(\operatorname{mip\_}) \wedge \\ \quad \operatorname{mihs\_} \\ \quad =\; {[}{[}\operatorname{mig\_}, \operatorname{enc\text{-}all}(\operatorname{enc\text{-}imp}(\operatorname{enc\text{-}and}(\operatorname{enc\text{-}in}(\operatorname{ebvar}(0), \operatorname{enc\text{-}ordc}), \operatorname{enc\text{-}all}(\operatorname{enc\text{-}imp}(\operatorname{enc\text{-}ordlt}(\operatorname{ebvar}(0), \operatorname{ebvar}(1)), \operatorname{mip\_}))), \operatorname{mip\_}))]] \end{array}
\end{equation*}

\par\medskip\noindent\textbf{Definition} (\texttt{rel-tfi-3cases-nbg}).\;For all $\operatorname{mit\_}$, $\operatorname{mihs\_}$ and $\operatorname{mic\_}$, 
\begin{equation*}
\begin{array}{@{}l@{}} \operatorname{rel\text{-}tfi\text{-}3cases\text{-}nbg}(\operatorname{mit\_}, \operatorname{mihs\_}, \operatorname{mic\_}) \\ \Leftrightarrow\; \operatorname{mic\_} \in \operatorname{seq\text{-}set} \wedge \\ \exists \operatorname{mig\_}, \operatorname{mip\_}, \operatorname{mis\_}, \operatorname{mifs\_}. \\ \quad \operatorname{mic\_} = {[}\operatorname{mig\_}, \operatorname{enc\text{-}all}(\operatorname{enc\text{-}imp}(\operatorname{enc\text{-}in}(\operatorname{ebvar}(0), \operatorname{enc\text{-}ordc}), \operatorname{mip\_}))] \wedge \\ \quad \operatorname{nbg\text{-}ok}(\operatorname{mip\_}) \wedge \\ \quad \operatorname{mihs\_} \\ \quad =\; {[}{[}\operatorname{mig\_}, \operatorname{expr\text{-}open}(\operatorname{mip\_}, 0, \operatorname{enc\text{-}num}(0))], {[}\operatorname{mig\_}, \operatorname{enc\text{-}all}(\operatorname{mis\_})], {[}\operatorname{mig\_}, \operatorname{enc\text{-}all}(\operatorname{enc\text{-}imp}(\operatorname{enc\text{-}and}(\operatorname{enc\text{-}limit}(\operatorname{ebvar}(0)), \operatorname{enc\text{-}all}(\operatorname{enc\text{-}imp}(\operatorname{enc\text{-}ordlt}(\operatorname{ebvar}(0), \operatorname{ebvar}(1)), \operatorname{mip\_}))), \operatorname{mip\_}))]] \wedge \\ \quad \operatorname{mifs\_} \in \mathbf{Set} \wedge \\ \quad \operatorname{mifs\_} \subseteq \mathbb{N} \wedge \\ \quad |\operatorname{mifs\_}| \in \mathbb{N} \wedge \\ \quad \forall \operatorname{miy\_} \in \mathbb{N}. \\ \quad \quad \neg (\operatorname{miy\_} \in \operatorname{mifs\_}) \Rightarrow  \\ \quad \quad \quad \operatorname{expr\text{-}open}(\operatorname{mis\_}, 0, \operatorname{efvar}(\operatorname{miy\_})) \\ \quad \quad \quad =\; \operatorname{enc\text{-}imp}(\operatorname{enc\text{-}and}(\operatorname{enc\text{-}in}(\operatorname{efvar}(\operatorname{miy\_}), \operatorname{enc\text{-}ordc}), \operatorname{expr\text{-}open}(\operatorname{mip\_}, 0, \operatorname{efvar}(\operatorname{miy\_}))), \\ \quad \quad \quad \quad \operatorname{expr\text{-}open}(\operatorname{mip\_}, 0, \operatorname{enc\text{-}succo}(\operatorname{efvar}(\operatorname{miy\_})))) \end{array}
\end{equation*}

\par\medskip\noindent\textbf{Definition} (\texttt{rel-cartesian-intro}).\;For all $\operatorname{mit\_}$, $\operatorname{mihs\_}$ and $\operatorname{mic\_}$, 
\begin{equation*}
\begin{array}{@{}l@{}} \operatorname{rel\text{-}cartesian\text{-}intro}(\operatorname{mit\_}, \operatorname{mihs\_}, \operatorname{mic\_}) \\ \Leftrightarrow\; \operatorname{mic\_} \in \operatorname{seq\text{-}set} \wedge \\ \exists \operatorname{mig\_}, \operatorname{mie\_}, \operatorname{mia\_}. \\ \quad \operatorname{mic\_} = {[}\operatorname{mig\_}, \operatorname{enc\text{-}in}(\operatorname{enc\text{-}list}(\operatorname{mie\_}), \operatorname{enc\text{-}cart}(\operatorname{mia\_}))] \wedge \\ \quad \operatorname{length}(\operatorname{mie\_}) = \operatorname{length}(\operatorname{mia\_}) \wedge \\ \quad 1 \leq \operatorname{length}(\operatorname{mia\_}) \wedge \\ \quad \operatorname{length}(\operatorname{mihs\_}) = \operatorname{length}(\operatorname{mie\_}) \wedge \\ \quad \forall \operatorname{mii\_} \in \mathbb{N}. \\ \quad \quad (1 \leq \operatorname{mii\_}) \Rightarrow  \\ \quad \quad (\operatorname{mii\_} \leq \operatorname{length}(\operatorname{mie\_})) \Rightarrow  \\ \quad \quad \quad \operatorname{nth}(\operatorname{mii\_}, \operatorname{mihs\_}) = {[}\operatorname{mig\_}, \operatorname{enc\text{-}in}(\operatorname{nth}(\operatorname{mii\_}, \operatorname{mie\_}), \operatorname{nth}(\operatorname{mii\_}, \operatorname{mia\_}))] \end{array}
\end{equation*}

\par\medskip\noindent\textbf{Definition} (\texttt{rel-cartesian-elim}).\;For all $\operatorname{mit\_}$, $\operatorname{mihs\_}$ and $\operatorname{mic\_}$, 
\begin{equation*}
\begin{array}{@{}l@{}} \operatorname{rel\text{-}cartesian\text{-}elim}(\operatorname{mit\_}, \operatorname{mihs\_}, \operatorname{mic\_}) \\ \Leftrightarrow\; \operatorname{mic\_} \in \operatorname{seq\text{-}set} \wedge \\ \exists \operatorname{mig\_}, \operatorname{mif\_}, \operatorname{mie\_}, \operatorname{mia\_}, \operatorname{mik\_}. \\ \quad \operatorname{mic\_} = {[}\operatorname{mig\_}, \operatorname{mif\_}] \wedge \\ \quad \operatorname{enc\text{-}in}(\operatorname{enc\text{-}list}(\operatorname{mie\_}), \operatorname{enc\text{-}cart}(\operatorname{mia\_})) \in \operatorname{mig\_} \wedge \\ \quad \operatorname{length}(\operatorname{mie\_}) = \operatorname{length}(\operatorname{mia\_}) \wedge \\ \quad \operatorname{mik\_} \in \mathbb{N} \wedge \\ \quad 1 \leq \operatorname{mik\_} \wedge \\ \quad \operatorname{mik\_} \leq \operatorname{length}(\operatorname{mie\_}) \wedge \\ \quad \operatorname{mihs\_} = {[}{[}\operatorname{ctx\text{-}add}(\operatorname{mig\_}, \operatorname{enc\text{-}in}(\operatorname{nth}(\operatorname{mik\_}, \operatorname{mie\_}), \operatorname{nth}(\operatorname{mik\_}, \operatorname{mia\_}))), \operatorname{mif\_}]] \end{array}
\end{equation*}

\par\medskip\noindent\textbf{Definition} (\texttt{idx-list}).\;$\operatorname{idx\text{-}list}(\operatorname{min\_}) :=$
\begin{equation*}
(\iota\, \operatorname{mir\_}.\;((\operatorname{mir\_} \in \operatorname{tuples}(\mathbf{Set})) \wedge ((\operatorname{length}(\operatorname{mir\_}) = \operatorname{min\_}) \wedge \forall \operatorname{mii\_} \in \mathbb{N}.\; ((1 \leq \operatorname{mii\_}) \Rightarrow ((\operatorname{mii\_} \leq \operatorname{min\_}) \Rightarrow (\operatorname{nth}(\operatorname{mii\_}, \operatorname{mir\_}) = \operatorname{ebvar}((\operatorname{min\_} - \operatorname{mii\_}))))))))
\end{equation*}

\par\medskip\noindent\textbf{Definition} (\texttt{enc-cch}).\;For all $\operatorname{mix\_}$, $\operatorname{mia\_}$ and $\operatorname{min\_}$, 
\begin{equation*}
(\operatorname{enc\text{-}cch}(\operatorname{mix\_}, \operatorname{mia\_}, \operatorname{min\_}, 0) \simeq \operatorname{enc\text{-}eq}(\operatorname{mix\_}, \operatorname{enc\text{-}list}(\operatorname{idx\text{-}list}(\operatorname{min\_}))))
\end{equation*}
\par
Let $\operatorname{mix\_}$, $\operatorname{mia\_}$ and $\operatorname{min\_}$ be arbitrary. Let $\operatorname{mik\_} \in \mathbb{N}$. Then
\begin{equation*}
(\operatorname{enc\text{-}cch}(\operatorname{mix\_}, \operatorname{mia\_}, \operatorname{min\_}, \operatorname{succ}(\operatorname{mik\_})) \simeq \operatorname{enc\text{-}ex}(\operatorname{enc\text{-}and}(\operatorname{enc\text{-}in}(\operatorname{ebvar}(0), \operatorname{nth}(-(\operatorname{min\_}, \operatorname{mik\_}), \operatorname{mia\_})), \operatorname{enc\text{-}cch}(\operatorname{mix\_}, \operatorname{mia\_}, \operatorname{min\_}, \operatorname{mik\_}))))
\end{equation*}

\par\medskip\noindent\textbf{Definition} (\texttt{j-teq}).\;$\operatorname{j\text{-}teq}(\operatorname{mit\_}) :=$
\begin{equation*}
\{\operatorname{mix\_} \in \operatorname{rw\text{-}univ} : \exists \operatorname{mig\_}.; \exists \operatorname{mia\_}, \operatorname{mib\_}, \operatorname{mif\_} \in \mathbf{Set}.\; ((\operatorname{mix\_} = {[}2, \operatorname{mig\_}, \operatorname{enc\text{-}eq}(\operatorname{enc\text{-}list}(\operatorname{mia\_}), \operatorname{enc\text{-}list}(\operatorname{mib\_})), \operatorname{mif\_}]) \wedge ((\operatorname{mig\_} \in \operatorname{ctx\text{-}set}) \wedge ((\operatorname{length}(\operatorname{mia\_}) = \operatorname{length}(\operatorname{mib\_})) \wedge ((\operatorname{enc\text{-}eq}(\operatorname{enc\text{-}list}(\operatorname{mia\_}), \operatorname{enc\text{-}list}(\operatorname{mib\_})) \in \mathit{clf}) \wedge ((\operatorname{mif\_} \in \mathit{clf}) \wedge ((((\operatorname{length}(\operatorname{mia\_}) = 0) \wedge (\operatorname{mif\_} = \operatorname{enc\text{-}truth})) \vee ({[}0, \operatorname{mif\_}] \in \operatorname{tf\text{-}set}(\operatorname{mia\_}, \operatorname{mib\_}))) \wedge (({[}\operatorname{mig\_}, \operatorname{enc\text{-}eq}(\operatorname{enc\text{-}list}(\operatorname{mia\_}), \operatorname{enc\text{-}list}(\operatorname{mia\_}))] \in \operatorname{derivable}(\operatorname{mit\_})) \wedge ({[}\operatorname{mig\_}, \operatorname{enc\text{-}eq}(\operatorname{enc\text{-}list}(\operatorname{mib\_}), \operatorname{enc\text{-}list}(\operatorname{mib\_}))] \in \operatorname{derivable}(\operatorname{mit\_})))))))))\}
\end{equation*}

\par\medskip\noindent\textbf{Definition} (\texttt{j-cart}).\;$\operatorname{j\text{-}cart}(\operatorname{mit\_}) :=$
\begin{equation*}
\{\operatorname{mix\_} \in \operatorname{rw\text{-}univ} : \exists \operatorname{mig\_}.; \exists \operatorname{miz\_}, \operatorname{mia\_} \in \mathbf{Set}.\; ((\operatorname{mix\_} = {[}2, \operatorname{mig\_}, \operatorname{enc\text{-}in}(\operatorname{miz\_}, \operatorname{enc\text{-}cart}(\operatorname{mia\_})), \operatorname{enc\text{-}cch}(\operatorname{miz\_}, \operatorname{mia\_}, \operatorname{length}(\operatorname{mia\_}), \operatorname{length}(\operatorname{mia\_}))]) \wedge ((\operatorname{mig\_} \in \operatorname{ctx\text{-}set}) \wedge ((1 \leq \operatorname{length}(\operatorname{mia\_})) \wedge ((\operatorname{enc\text{-}in}(\operatorname{miz\_}, \operatorname{enc\text{-}cart}(\operatorname{mia\_})) \in \mathit{clf}) \wedge (\operatorname{enc\text{-}cch}(\operatorname{miz\_}, \operatorname{mia\_}, \operatorname{length}(\operatorname{mia\_}), \operatorname{length}(\operatorname{mia\_})) \in \mathit{clf})))))\}
\end{equation*}

\par\medskip\noindent\textbf{Definition} (\texttt{cchs}).\;For all $\operatorname{m57x\_}$, $\operatorname{m57a\_}$, $\operatorname{m57n\_}$ and $\operatorname{m57l\_}$, 
\begin{equation*}
(\operatorname{cchs}(\operatorname{m57x\_}, \operatorname{m57a\_}, \operatorname{m57n\_}, \operatorname{m57l\_}, 0) \simeq \operatorname{enc\text{-}eq}(\operatorname{m57x\_}, \operatorname{enc\text{-}list}(\operatorname{m57l\_})))
\end{equation*}
\par
Let $\operatorname{m57x\_}$, $\operatorname{m57a\_}$, $\operatorname{m57n\_}$ and $\operatorname{m57l\_}$ be arbitrary. Let $\operatorname{m57k\_} \in \mathbb{N}$. Then
\begin{equation*}
(\operatorname{cchs}(\operatorname{m57x\_}, \operatorname{m57a\_}, \operatorname{m57n\_}, \operatorname{m57l\_}, \operatorname{succ}(\operatorname{m57k\_})) \simeq \operatorname{enc\text{-}ex}(\operatorname{enc\text{-}and}(\operatorname{enc\text{-}in}(\operatorname{ebvar}(0), \operatorname{nth}(-(\operatorname{m57n\_}, \operatorname{m57k\_}), \operatorname{m57a\_})), \operatorname{cchs}(\operatorname{m57x\_}, \operatorname{m57a\_}, \operatorname{m57n\_}, \operatorname{m57l\_}, \operatorname{m57k\_}))))
\end{equation*}

\par\medskip\noindent\textbf{Definition} (\texttt{rel-cartesian-decompose-1}).\;For all $\operatorname{mit\_}$, $\operatorname{mihs\_}$ and $\operatorname{mic\_}$, 
\begin{equation*}
\operatorname{rel\text{-}cartesian\text{-}decompose\text{-}1}(\operatorname{mit\_}, \operatorname{mihs\_}, \operatorname{mic\_}) \Leftrightarrow \operatorname{rel\text{-}rewrite\text{-}1}(\operatorname{j\text{-}cart}(\operatorname{mit\_}), \operatorname{mihs\_}, \operatorname{mic\_})
\end{equation*}

\par\medskip\noindent\textbf{Definition} (\texttt{rel-cartesian-decompose-7}).\;For all $\operatorname{mit\_}$, $\operatorname{mihs\_}$ and $\operatorname{mic\_}$, 
\begin{equation*}
\operatorname{rel\text{-}cartesian\text{-}decompose\text{-}7}(\operatorname{mit\_}, \operatorname{mihs\_}, \operatorname{mic\_}) \Leftrightarrow \operatorname{rel\text{-}rewrite\text{-}7}(\operatorname{j\text{-}cart}(\operatorname{mit\_}), \operatorname{mihs\_}, \operatorname{mic\_})
\end{equation*}

\par\medskip\noindent\textbf{Definition} (\texttt{tf-set}).\;$\operatorname{tf\text{-}set}(\operatorname{mitfa\_}, \operatorname{mitfb\_}) :=$
\begin{equation*}
\{\operatorname{mitfp\_} \in (\mathbb{N} \times \operatorname{expr\text{-}set}) : \forall \operatorname{mitfw\_} \in \operatorname{power}((\mathbb{N} \times \operatorname{expr\text{-}set})).\; (\operatorname{tf\text{-}closed}(\operatorname{mitfa\_}, \operatorname{mitfb\_}, \operatorname{mitfw\_}) \Rightarrow (\operatorname{mitfp\_} \in \operatorname{mitfw\_}))\}
\end{equation*}

\par\medskip\noindent\textbf{Definition} (\texttt{tf-closed}).\;For all $\operatorname{mitfa\_}$, $\operatorname{mitfb\_}$ and $\operatorname{mitfs\_}$, 
\begin{equation*}
\begin{array}{@{}l@{}} \operatorname{tf\text{-}closed}(\operatorname{mitfa\_}, \operatorname{mitfb\_}, \operatorname{mitfs\_}) \\ \Leftrightarrow\; \forall \operatorname{mitfj\_} \in \mathbb{N}. \\ \quad (\operatorname{succ}(\operatorname{mitfj\_}) = \operatorname{length}(\operatorname{mitfa\_})) \Rightarrow  \\ \quad \quad {[}\operatorname{mitfj\_}, \operatorname{enc\text{-}eq}(\operatorname{nth}(\operatorname{succ}(\operatorname{mitfj\_}), \operatorname{mitfa\_}), \operatorname{nth}(\operatorname{succ}(\operatorname{mitfj\_}), \operatorname{mitfb\_}))] \in \operatorname{mitfs\_} \wedge \\ \forall \operatorname{mitfk\_} \in \mathbb{N}, \operatorname{mitff\_}. \\ \quad ({[}\operatorname{succ}(\operatorname{mitfk\_}), \operatorname{mitff\_}] \in \operatorname{mitfs\_}) \Rightarrow  \\ \quad \quad {[}\operatorname{mitfk\_}, \operatorname{enc\text{-}and}(\operatorname{enc\text{-}eq}(\operatorname{nth}(\operatorname{succ}(\operatorname{mitfk\_}), \operatorname{mitfa\_}), \operatorname{nth}(\operatorname{succ}(\operatorname{mitfk\_}), \operatorname{mitfb\_})), \operatorname{mitff\_})] \in \operatorname{mitfs\_} \end{array}
\end{equation*}

\par\medskip\noindent\textbf{Definition} (\texttt{rel-tuple-eq-decompose-1}).\;For all $\operatorname{mit\_}$, $\operatorname{mihs\_}$ and $\operatorname{mic\_}$, 
\begin{equation*}
\operatorname{rel\text{-}tuple\text{-}eq\text{-}decompose\text{-}1}(\operatorname{mit\_}, \operatorname{mihs\_}, \operatorname{mic\_}) \Leftrightarrow \operatorname{rel\text{-}rewrite\text{-}1}(\operatorname{j\text{-}teq}(\operatorname{mit\_}), \operatorname{mihs\_}, \operatorname{mic\_})
\end{equation*}

\par\medskip\noindent\textbf{Definition} (\texttt{rel-tuple-eq-decompose-7}).\;For all $\operatorname{mit\_}$, $\operatorname{mihs\_}$ and $\operatorname{mic\_}$, 
\begin{equation*}
\operatorname{rel\text{-}tuple\text{-}eq\text{-}decompose\text{-}7}(\operatorname{mit\_}, \operatorname{mihs\_}, \operatorname{mic\_}) \Leftrightarrow \operatorname{rel\text{-}rewrite\text{-}7}(\operatorname{j\text{-}teq}(\operatorname{mit\_}), \operatorname{mihs\_}, \operatorname{mic\_})
\end{equation*}

\par\medskip\noindent\textbf{Definition} (\texttt{def-comp-closed}).\;For every $\operatorname{mit\_}$, 
\begin{equation*}
\begin{array}{@{}l@{}} \operatorname{def\text{-}comp\text{-}closed}(\operatorname{mit\_}) \\ \Leftrightarrow\; \forall \operatorname{mig\_} \in \operatorname{ctx\text{-}set}, \operatorname{miq\_}. \\ \quad (\operatorname{enc\text{-}comp}(\operatorname{miq\_}) \in \mathit{clt}) \Rightarrow  \\ \quad \operatorname{nbg\text{-}ok}(\operatorname{miq\_}) \Rightarrow  \\ \quad \quad {[}\operatorname{mig\_}, \operatorname{enc\text{-}eq}(\operatorname{enc\text{-}comp}(\operatorname{miq\_}), \operatorname{enc\text{-}comp}(\operatorname{miq\_}))] \in \operatorname{derivable}(\operatorname{mit\_}) \end{array}
\end{equation*}

\par\medskip\noindent\textbf{Definition} (\texttt{def-in-closed}).\;For every $\operatorname{mit\_}$, 
\begin{equation*}
\begin{array}{@{}l@{}} \operatorname{def\text{-}in\text{-}closed}(\operatorname{mit\_}) \\ \Leftrightarrow\; \forall \operatorname{mig\_}, \operatorname{miw\_}, \operatorname{mia\_}. \\ \quad ({[}\operatorname{mig\_}, \operatorname{enc\text{-}in}(\operatorname{miw\_}, \operatorname{mia\_})] \in \operatorname{derivable}(\operatorname{mit\_})) \Rightarrow  \\ \quad \quad {[}\operatorname{mig\_}, \operatorname{enc\text{-}eq}(\operatorname{miw\_}, \operatorname{miw\_})] \in \operatorname{derivable}(\operatorname{mit\_}) \end{array}
\end{equation*}

\par\medskip\noindent\textbf{Definition} (\texttt{ax-zero-nn}).\; $(\operatorname{ax\text{-}zero\text{-}nn} \simeq \operatorname{enc\text{-}in}(\operatorname{enc\text{-}num}(0), \operatorname{enc\text{-}nnc}))$
\par\medskip\noindent\textbf{Definition} (\texttt{ax-succ-nn}).\;
\begin{equation*}
(\operatorname{ax\text{-}succ\text{-}nn} \simeq \operatorname{enc\text{-}all}(\operatorname{enc\text{-}imp}(\operatorname{enc\text{-}in}(\operatorname{ebvar}(0), \operatorname{enc\text{-}nnc}), \operatorname{enc\text{-}in}(\operatorname{enc\text{-}succ}(\operatorname{ebvar}(0)), \operatorname{enc\text{-}nnc}))))
\end{equation*}

\par\medskip\noindent\textbf{Definition} (\texttt{ax-limit-iff}).\;
\begin{equation*}
(\operatorname{ax\text{-}limit\text{-}iff} \simeq \operatorname{enc\text{-}all}(\operatorname{enc\text{-}iff}(\operatorname{enc\text{-}limit}(\operatorname{ebvar}(0)), \operatorname{enc\text{-}and}(\operatorname{enc\text{-}in}(\operatorname{ebvar}(0), \operatorname{enc\text{-}ordc}), \operatorname{enc\text{-}and}(\operatorname{enc\text{-}not}(\operatorname{enc\text{-}eq}(\operatorname{ebvar}(0), \operatorname{enc\text{-}num}(0))), \operatorname{enc\text{-}not}(\operatorname{enc\text{-}ex}(\operatorname{enc\text{-}and}(\operatorname{enc\text{-}in}(\operatorname{ebvar}(0), \operatorname{enc\text{-}ordc}), \operatorname{enc\text{-}eq}(\operatorname{ebvar}(1), \operatorname{enc\text{-}succo}(\operatorname{ebvar}(0)))))))))))
\end{equation*}

\par\medskip\noindent\textbf{Definition} (\texttt{ax-ord-succ-above}).\;
\begin{equation*}
(\operatorname{ax\text{-}ord\text{-}succ\text{-}above} \simeq \operatorname{enc\text{-}all}(\operatorname{enc\text{-}imp}(\operatorname{enc\text{-}in}(\operatorname{ebvar}(0), \operatorname{enc\text{-}ordc}), \operatorname{enc\text{-}ordlt}(\operatorname{ebvar}(0), \operatorname{enc\text{-}succo}(\operatorname{ebvar}(0))))))
\end{equation*}

\par\medskip\noindent\textbf{Definition} (\texttt{rel-lambda-type-n}).\;For all $\operatorname{mit\_}$, $\operatorname{mihs\_}$ and $\operatorname{mic\_}$, 
\begin{equation*}
\begin{array}{@{}l@{}} \operatorname{rel\text{-}lambda\text{-}type\text{-}n}(\operatorname{mit\_}, \operatorname{mihs\_}, \operatorname{mic\_}) \\ \Leftrightarrow\; \operatorname{mic\_} \in \operatorname{seq\text{-}set} \wedge \\ \exists \operatorname{mig\_}, \operatorname{mia\_}, \operatorname{mib\_}, \operatorname{mid\_}, \operatorname{mie\_}. \\ \quad \operatorname{mic\_} = {[}\operatorname{mig\_}, \operatorname{enc\text{-}in}(\operatorname{enc\text{-}lam}(\operatorname{enc\text{-}cart}(\operatorname{mia\_}), \operatorname{mie\_}), \operatorname{enc\text{-}fun}(\operatorname{enc\text{-}cart}(\operatorname{mia\_}), \operatorname{mid\_}))] \wedge \\ \quad 1 \leq \operatorname{length}(\operatorname{mia\_}) \wedge \\ \quad \operatorname{mib\_} \in \operatorname{expr\text{-}set} \wedge \\ \quad \operatorname{nbg\text{-}ok}(\operatorname{mie\_}) \wedge \\ \quad \forall \operatorname{miy\_} \in \mathbb{N}. \\ \quad \quad \operatorname{expr\text{-}open}(\operatorname{mie\_}, 0, \operatorname{efvar}(\operatorname{miy\_})) = \operatorname{ltop}(\operatorname{mib\_}, \operatorname{miy\_}, \operatorname{length}(\operatorname{mia\_}), \operatorname{length}(\operatorname{mia\_})) \wedge \\ \quad \operatorname{mihs\_} \\ \quad =\; {[}{[}\operatorname{mig\_}, \operatorname{ltch}(\operatorname{mia\_}, \operatorname{mib\_}, \operatorname{mid\_}, \operatorname{length}(\operatorname{mia\_}), \operatorname{length}(\operatorname{mia\_}))], {[}\operatorname{mig\_}, \operatorname{enc\text{-}in}(\operatorname{enc\text{-}cart}(\operatorname{mia\_}), \operatorname{enc\text{-}setc})], {[}\operatorname{mig\_}, \operatorname{enc\text{-}eq}(\operatorname{mid\_}, \operatorname{mid\_})]] \end{array}
\end{equation*}

\par\medskip\noindent\textbf{Definition} (\texttt{ltch}).\;For all $\operatorname{m52a\_}$, $\operatorname{m52b\_}$, $\operatorname{m52d\_}$ and $\operatorname{m52n\_}$,  $(\operatorname{ltch}(\operatorname{m52a\_}, \operatorname{m52b\_}, \operatorname{m52d\_}, \operatorname{m52n\_}, 0) \simeq \operatorname{enc\text{-}in}(\operatorname{m52b\_}, \operatorname{m52d\_}))$\par
Let $\operatorname{m52a\_}$, $\operatorname{m52b\_}$, $\operatorname{m52d\_}$ and $\operatorname{m52n\_}$ be arbitrary. Let $\operatorname{m52k\_} \in \mathbb{N}$. Then
\begin{equation*}
(\operatorname{ltch}(\operatorname{m52a\_}, \operatorname{m52b\_}, \operatorname{m52d\_}, \operatorname{m52n\_}, \operatorname{succ}(\operatorname{m52k\_})) \simeq \operatorname{enc\text{-}all}(\operatorname{enc\text{-}imp}(\operatorname{enc\text{-}in}(\operatorname{ebvar}(0), \operatorname{nth}(-(\operatorname{m52n\_}, \operatorname{m52k\_}), \operatorname{m52a\_})), \operatorname{ltch}(\operatorname{m52a\_}, \operatorname{m52b\_}, \operatorname{m52d\_}, \operatorname{m52n\_}, \operatorname{m52k\_}))))
\end{equation*}

\par\medskip\noindent\textbf{Definition} (\texttt{ltop}).\;For all $\operatorname{m52b\_}$, $\operatorname{m52y\_}$ and $\operatorname{m52n\_}$,  $(\operatorname{ltop}(\operatorname{m52b\_}, \operatorname{m52y\_}, \operatorname{m52n\_}, 0) \simeq \operatorname{m52b\_})$\par
Let $\operatorname{m52b\_}$, $\operatorname{m52y\_}$ and $\operatorname{m52n\_}$ be arbitrary. Let $\operatorname{m52k\_} \in \mathbb{N}$. Then
\begin{equation*}
(\operatorname{ltop}(\operatorname{m52b\_}, \operatorname{m52y\_}, \operatorname{m52n\_}, \operatorname{succ}(\operatorname{m52k\_})) \simeq \operatorname{expr\text{-}open}(\operatorname{ltop}(\operatorname{m52b\_}, \operatorname{m52y\_}, \operatorname{m52n\_}, \operatorname{m52k\_}), -(\operatorname{m52n\_}, \operatorname{succ}(\operatorname{m52k\_})), \operatorname{enc\text{-}nth}(\operatorname{enc\text{-}num}(\operatorname{succ}(\operatorname{m52k\_})), \operatorname{efvar}(\operatorname{m52y\_}))))
\end{equation*}

\par\medskip\noindent\textbf{Definition} (\texttt{rel-lambda-type-4}).\;For all $\operatorname{mit\_}$, $\operatorname{mihs\_}$ and $\operatorname{mic\_}$, 
\begin{equation*}
\begin{array}{@{}l@{}} \operatorname{rel\text{-}lambda\text{-}type\text{-}4}(\operatorname{mit\_}, \operatorname{mihs\_}, \operatorname{mic\_}) \\ \Leftrightarrow\; \operatorname{mic\_} \in \operatorname{seq\text{-}set} \wedge \\ \exists \operatorname{mig\_}, \operatorname{mia1\_}, \operatorname{mia2\_}, \operatorname{mia3\_}, \operatorname{mia4\_}, \operatorname{mib\_}, \operatorname{mid\_}, \operatorname{mie\_}. \\ \quad \operatorname{mic\_} \\ \quad =\; {[}\operatorname{mig\_}, \operatorname{enc\text{-}in}(\operatorname{enc\text{-}lam}(\operatorname{enc\text{-}cart}({[}\operatorname{mia1\_}, \operatorname{mia2\_}, \operatorname{mia3\_}, \operatorname{mia4\_}]), \operatorname{mie\_}), \operatorname{enc\text{-}fun}(\operatorname{enc\text{-}cart}({[}\operatorname{mia1\_}, \operatorname{mia2\_}, \operatorname{mia3\_}, \operatorname{mia4\_}]), \operatorname{mid\_}))] \wedge \\ \quad \operatorname{nbg\text{-}ok}(\operatorname{mie\_}) \wedge \\ \quad \forall \operatorname{miy\_} \in \mathbb{N}. \\ \quad \quad \operatorname{expr\text{-}open}(\operatorname{mie\_}, 0, \operatorname{efvar}(\operatorname{miy\_})) \\ \quad \quad =\; \operatorname{expr\text{-}open}(\operatorname{expr\text{-}open}(\operatorname{expr\text{-}open}(\operatorname{expr\text{-}open}(\operatorname{mib\_}, 3, \operatorname{enc\text{-}nth}(\operatorname{enc\text{-}num}(1), \operatorname{efvar}(\operatorname{miy\_}))), 2, \operatorname{enc\text{-}nth}(\operatorname{enc\text{-}num}(2), \operatorname{efvar}(\operatorname{miy\_}))), \\ \quad \quad \quad \quad 1, \\ \quad \quad \quad \quad \operatorname{enc\text{-}nth}(\operatorname{enc\text{-}num}(3), \operatorname{efvar}(\operatorname{miy\_}))), \\ \quad \quad \quad 0, \\ \quad \quad \quad \operatorname{enc\text{-}nth}(\operatorname{enc\text{-}num}(4), \operatorname{efvar}(\operatorname{miy\_}))) \wedge \\ \quad \operatorname{mihs\_} \\ \quad =\; {[}{[}\operatorname{mig\_}, \operatorname{enc\text{-}all}(\operatorname{enc\text{-}imp}(\operatorname{enc\text{-}in}(\operatorname{ebvar}(0), \operatorname{mia1\_}), \operatorname{enc\text{-}all}(\operatorname{enc\text{-}imp}(\operatorname{enc\text{-}in}(\operatorname{ebvar}(0), \operatorname{mia2\_}), \operatorname{enc\text{-}all}(\operatorname{enc\text{-}imp}(\operatorname{enc\text{-}in}(\operatorname{ebvar}(0), \operatorname{mia3\_}), \operatorname{enc\text{-}all}(\operatorname{enc\text{-}imp}(\operatorname{enc\text{-}in}(\operatorname{ebvar}(0), \operatorname{mia4\_}), \operatorname{enc\text{-}in}(\operatorname{mib\_}, \operatorname{mid\_})))))))))], {[}\operatorname{mig\_}, \operatorname{enc\text{-}in}(\operatorname{enc\text{-}cart}({[}\operatorname{mia1\_}, \operatorname{mia2\_}, \operatorname{mia3\_}, \operatorname{mia4\_}]), \operatorname{enc\text{-}setc})], {[}\operatorname{mig\_}, \operatorname{enc\text{-}eq}(\operatorname{mid\_}, \operatorname{mid\_})]] \end{array}
\end{equation*}

\par\medskip\noindent\textbf{Definition} (\texttt{rel-lambda-type-3}).\;For all $\operatorname{mit\_}$, $\operatorname{mihs\_}$ and $\operatorname{mic\_}$, 
\begin{equation*}
\begin{array}{@{}l@{}} \operatorname{rel\text{-}lambda\text{-}type\text{-}3}(\operatorname{mit\_}, \operatorname{mihs\_}, \operatorname{mic\_}) \\ \Leftrightarrow\; \operatorname{mic\_} \in \operatorname{seq\text{-}set} \wedge \\ \exists \operatorname{mig\_}, \operatorname{mia1\_}, \operatorname{mia2\_}, \operatorname{mia3\_}, \operatorname{mib\_}, \operatorname{mid\_}, \operatorname{mie\_}. \\ \quad \operatorname{mic\_} \\ \quad =\; {[}\operatorname{mig\_}, \operatorname{enc\text{-}in}(\operatorname{enc\text{-}lam}(\operatorname{enc\text{-}cart}({[}\operatorname{mia1\_}, \operatorname{mia2\_}, \operatorname{mia3\_}]), \operatorname{mie\_}), \operatorname{enc\text{-}fun}(\operatorname{enc\text{-}cart}({[}\operatorname{mia1\_}, \operatorname{mia2\_}, \operatorname{mia3\_}]), \operatorname{mid\_}))] \wedge \\ \quad \operatorname{nbg\text{-}ok}(\operatorname{mie\_}) \wedge \\ \quad \forall \operatorname{miy\_} \in \mathbb{N}. \\ \quad \quad \operatorname{expr\text{-}open}(\operatorname{mie\_}, 0, \operatorname{efvar}(\operatorname{miy\_})) \\ \quad \quad =\; \operatorname{expr\text{-}open}(\operatorname{expr\text{-}open}(\operatorname{expr\text{-}open}(\operatorname{mib\_}, 2, \operatorname{enc\text{-}nth}(\operatorname{enc\text{-}num}(1), \operatorname{efvar}(\operatorname{miy\_}))), 1, \operatorname{enc\text{-}nth}(\operatorname{enc\text{-}num}(2), \operatorname{efvar}(\operatorname{miy\_}))), \\ \quad \quad \quad 0, \\ \quad \quad \quad \operatorname{enc\text{-}nth}(\operatorname{enc\text{-}num}(3), \operatorname{efvar}(\operatorname{miy\_}))) \wedge \\ \quad \operatorname{mihs\_} \\ \quad =\; {[}{[}\operatorname{mig\_}, \operatorname{enc\text{-}all}(\operatorname{enc\text{-}imp}(\operatorname{enc\text{-}in}(\operatorname{ebvar}(0), \operatorname{mia1\_}), \operatorname{enc\text{-}all}(\operatorname{enc\text{-}imp}(\operatorname{enc\text{-}in}(\operatorname{ebvar}(0), \operatorname{mia2\_}), \operatorname{enc\text{-}all}(\operatorname{enc\text{-}imp}(\operatorname{enc\text{-}in}(\operatorname{ebvar}(0), \operatorname{mia3\_}), \operatorname{enc\text{-}in}(\operatorname{mib\_}, \operatorname{mid\_})))))))], {[}\operatorname{mig\_}, \operatorname{enc\text{-}in}(\operatorname{enc\text{-}cart}({[}\operatorname{mia1\_}, \operatorname{mia2\_}, \operatorname{mia3\_}]), \operatorname{enc\text{-}setc})], {[}\operatorname{mig\_}, \operatorname{enc\text{-}eq}(\operatorname{mid\_}, \operatorname{mid\_})]] \end{array}
\end{equation*}

\par\medskip\noindent\textbf{Definition} (\texttt{rel-lambda-type-2}).\;For all $\operatorname{mit\_}$, $\operatorname{mihs\_}$ and $\operatorname{mic\_}$, 
\begin{equation*}
\begin{array}{@{}l@{}} \operatorname{rel\text{-}lambda\text{-}type\text{-}2}(\operatorname{mit\_}, \operatorname{mihs\_}, \operatorname{mic\_}) \\ \Leftrightarrow\; \operatorname{mic\_} \in \operatorname{seq\text{-}set} \wedge \\ \exists \operatorname{mig\_}, \operatorname{mia1\_}, \operatorname{mia2\_}, \operatorname{mib\_}, \operatorname{mid\_}, \operatorname{mie\_}. \\ \quad \operatorname{mic\_} \\ \quad =\; {[}\operatorname{mig\_}, \operatorname{enc\text{-}in}(\operatorname{enc\text{-}lam}(\operatorname{enc\text{-}cart}({[}\operatorname{mia1\_}, \operatorname{mia2\_}]), \operatorname{mie\_}), \operatorname{enc\text{-}fun}(\operatorname{enc\text{-}cart}({[}\operatorname{mia1\_}, \operatorname{mia2\_}]), \operatorname{mid\_}))] \wedge \\ \quad \operatorname{nbg\text{-}ok}(\operatorname{mie\_}) \wedge \\ \quad \forall \operatorname{miy\_} \in \mathbb{N}. \\ \quad \quad \operatorname{expr\text{-}open}(\operatorname{mie\_}, 0, \operatorname{efvar}(\operatorname{miy\_})) \\ \quad \quad =\; \operatorname{expr\text{-}open}(\operatorname{expr\text{-}open}(\operatorname{mib\_}, 1, \operatorname{enc\text{-}nth}(\operatorname{enc\text{-}num}(1), \operatorname{efvar}(\operatorname{miy\_}))), 0, \operatorname{enc\text{-}nth}(\operatorname{enc\text{-}num}(2), \operatorname{efvar}(\operatorname{miy\_}))) \wedge \\ \quad \operatorname{mihs\_} \\ \quad =\; {[}{[}\operatorname{mig\_}, \operatorname{enc\text{-}all}(\operatorname{enc\text{-}imp}(\operatorname{enc\text{-}in}(\operatorname{ebvar}(0), \operatorname{mia1\_}), \operatorname{enc\text{-}all}(\operatorname{enc\text{-}imp}(\operatorname{enc\text{-}in}(\operatorname{ebvar}(0), \operatorname{mia2\_}), \operatorname{enc\text{-}in}(\operatorname{mib\_}, \operatorname{mid\_})))))], {[}\operatorname{mig\_}, \operatorname{enc\text{-}in}(\operatorname{enc\text{-}cart}({[}\operatorname{mia1\_}, \operatorname{mia2\_}]), \operatorname{enc\text{-}setc})], {[}\operatorname{mig\_}, \operatorname{enc\text{-}eq}(\operatorname{mid\_}, \operatorname{mid\_})]] \end{array}
\end{equation*}

\par\medskip\noindent\textbf{Definition} (\texttt{rel-lambda-type-1}).\;For all $\operatorname{mit\_}$, $\operatorname{mihs\_}$ and $\operatorname{mic\_}$, 
\begin{equation*}
\begin{array}{@{}l@{}} \operatorname{rel\text{-}lambda\text{-}type\text{-}1}(\operatorname{mit\_}, \operatorname{mihs\_}, \operatorname{mic\_}) \\ \Leftrightarrow\; \operatorname{mic\_} \in \operatorname{seq\text{-}set} \wedge \\ \exists \operatorname{mig\_}, \operatorname{mia1\_}, \operatorname{mib\_}, \operatorname{mid\_}, \operatorname{mie\_}. \\ \quad \operatorname{mic\_} = {[}\operatorname{mig\_}, \operatorname{enc\text{-}in}(\operatorname{enc\text{-}lam}(\operatorname{enc\text{-}cart}({[}\operatorname{mia1\_}]), \operatorname{mie\_}), \operatorname{enc\text{-}fun}(\operatorname{enc\text{-}cart}({[}\operatorname{mia1\_}]), \operatorname{mid\_}))] \wedge \\ \quad \operatorname{nbg\text{-}ok}(\operatorname{mie\_}) \wedge \\ \quad \forall \operatorname{miy\_} \in \mathbb{N}. \\ \quad \quad \operatorname{expr\text{-}open}(\operatorname{mie\_}, 0, \operatorname{efvar}(\operatorname{miy\_})) = \operatorname{expr\text{-}open}(\operatorname{mib\_}, 0, \operatorname{enc\text{-}nth}(\operatorname{enc\text{-}num}(1), \operatorname{efvar}(\operatorname{miy\_}))) \wedge \\ \quad \operatorname{mihs\_} \\ \quad =\; {[}{[}\operatorname{mig\_}, \operatorname{enc\text{-}all}(\operatorname{enc\text{-}imp}(\operatorname{enc\text{-}in}(\operatorname{ebvar}(0), \operatorname{mia1\_}), \operatorname{enc\text{-}in}(\operatorname{mib\_}, \operatorname{mid\_})))], {[}\operatorname{mig\_}, \operatorname{enc\text{-}in}(\operatorname{enc\text{-}cart}({[}\operatorname{mia1\_}]), \operatorname{enc\text{-}setc})], {[}\operatorname{mig\_}, \operatorname{enc\text{-}eq}(\operatorname{mid\_}, \operatorname{mid\_})]] \end{array}
\end{equation*}


\vnbbookfile{structure-library/meta-def}{structure-library/meta-def}
\noindent{\small meta-def.scm -- K2 PHASE 5: the clauses of the kernel's definedness certificate pi--defined? (primitive-inferences.scm 1164-1235 and its helpers, tree f7e2246, 2026-10-12), each TRANSCRIBED as a rule DEF-k over the encoding: hypotheses the context facts the clause reads, conclusion [g, t = t]. ...}\par
\par\medskip\noindent\textbf{Definition} (\texttt{def-strict-in-l}).\;For all $\operatorname{mdt\_}$, $\operatorname{mdhs\_}$ and $\operatorname{mdc\_}$, 
\begin{equation*}
\begin{array}{@{}l@{}} \operatorname{def\text{-}strict\text{-}in\text{-}l}(\operatorname{mdt\_}, \operatorname{mdhs\_}, \operatorname{mdc\_}) \\ \Leftrightarrow\; \operatorname{mdc\_} \in \operatorname{seq\text{-}set} \wedge \\ \exists \operatorname{mdg\_}, \operatorname{mdw\_}, \operatorname{mda\_}. \\ \quad (\operatorname{mdc\_} = {[}\operatorname{mdg\_}, \operatorname{enc\text{-}eq}(\operatorname{mdw\_}, \operatorname{mdw\_})]) \wedge (\operatorname{mdhs\_} = {[}{[}\operatorname{mdg\_}, \operatorname{enc\text{-}in}(\operatorname{mdw\_}, \operatorname{mda\_})]]) \end{array}
\end{equation*}

\par\medskip\noindent\textbf{Definition} (\texttt{def-strict-eq-l}).\;For all $\operatorname{mdt\_}$, $\operatorname{mdhs\_}$ and $\operatorname{mdc\_}$, 
\begin{equation*}
\begin{array}{@{}l@{}} \operatorname{def\text{-}strict\text{-}eq\text{-}l}(\operatorname{mdt\_}, \operatorname{mdhs\_}, \operatorname{mdc\_}) \\ \Leftrightarrow\; \operatorname{mdc\_} \in \operatorname{seq\text{-}set} \wedge \\ \exists \operatorname{mdg\_}, \operatorname{mds\_}, \operatorname{mdu\_}. \\ \quad (\operatorname{mdc\_} = {[}\operatorname{mdg\_}, \operatorname{enc\text{-}eq}(\operatorname{mds\_}, \operatorname{mds\_})]) \wedge (\operatorname{mdhs\_} = {[}{[}\operatorname{mdg\_}, \operatorname{enc\text{-}eq}(\operatorname{mds\_}, \operatorname{mdu\_})]]) \end{array}
\end{equation*}

\par\medskip\noindent\textbf{Definition} (\texttt{def-strict-eq-r}).\;For all $\operatorname{mdt\_}$, $\operatorname{mdhs\_}$ and $\operatorname{mdc\_}$, 
\begin{equation*}
\begin{array}{@{}l@{}} \operatorname{def\text{-}strict\text{-}eq\text{-}r}(\operatorname{mdt\_}, \operatorname{mdhs\_}, \operatorname{mdc\_}) \\ \Leftrightarrow\; \operatorname{mdc\_} \in \operatorname{seq\text{-}set} \wedge \\ \exists \operatorname{mdg\_}, \operatorname{mds\_}, \operatorname{mdu\_}. \\ \quad (\operatorname{mdc\_} = {[}\operatorname{mdg\_}, \operatorname{enc\text{-}eq}(\operatorname{mdu\_}, \operatorname{mdu\_})]) \wedge (\operatorname{mdhs\_} = {[}{[}\operatorname{mdg\_}, \operatorname{enc\text{-}eq}(\operatorname{mds\_}, \operatorname{mdu\_})]]) \end{array}
\end{equation*}

\par\medskip\noindent\textbf{Definition} (\texttt{def-strict-arg}).\;For all $\operatorname{mdt\_}$, $\operatorname{mdhs\_}$ and $\operatorname{mdc\_}$, 
\begin{equation*}
\begin{array}{@{}l@{}} \operatorname{def\text{-}strict\text{-}arg}(\operatorname{mdt\_}, \operatorname{mdhs\_}, \operatorname{mdc\_}) \\ \Leftrightarrow\; \operatorname{mdc\_} \in \operatorname{seq\text{-}set} \wedge \\ \exists \operatorname{mdg\_}, \operatorname{mdh\_}, \operatorname{mdl\_}, \operatorname{mdi\_}. \\ \quad \operatorname{mdc\_} = {[}\operatorname{mdg\_}, \operatorname{enc\text{-}eq}(\operatorname{nth}(\operatorname{mdi\_}, \operatorname{mdl\_}), \operatorname{nth}(\operatorname{mdi\_}, \operatorname{mdl\_}))] \wedge \\ \quad \operatorname{strict\text{-}head}(\operatorname{mdh\_}) \wedge \\ \quad \operatorname{mdi\_} \in \mathbb{N} \wedge \\ \quad 1 \leq \operatorname{mdi\_} \wedge \\ \quad \operatorname{mdi\_} \leq \operatorname{length}(\operatorname{mdl\_}) \wedge \\ \quad \operatorname{mdhs\_} = {[}{[}\operatorname{mdg\_}, \operatorname{enc\text{-}eq}(\operatorname{eapp}(\operatorname{mdh\_}, \operatorname{mdl\_}), \operatorname{eapp}(\operatorname{mdh\_}, \operatorname{mdl\_}))]] \end{array}
\end{equation*}

\par\medskip\noindent\textbf{Definition} (\texttt{def-strict-hyp}).\;For all $\operatorname{mdt\_}$, $\operatorname{mdhs\_}$ and $\operatorname{mdc\_}$, 
\begin{equation*}
\begin{array}{@{}l@{}} \operatorname{def\text{-}strict\text{-}hyp}(\operatorname{mdt\_}, \operatorname{mdhs\_}, \operatorname{mdc\_}) \\ \Leftrightarrow\; \operatorname{mdc\_} \in \operatorname{seq\text{-}set} \wedge \\ \exists \operatorname{mdg\_}, \operatorname{mdf\_}, \operatorname{md7h2\_}, \operatorname{mdw\_}. \\ \quad \operatorname{mdc\_} = {[}\operatorname{mdg\_}, \operatorname{enc\text{-}eq}(\operatorname{mdw\_}, \operatorname{mdw\_})] \wedge \\ \quad \operatorname{mdhs\_} = {[}{[}\operatorname{mdg\_}, \operatorname{mdf\_}]] \wedge \\ \quad \operatorname{md7h2\_} \in \operatorname{true\text{-}reach}(\operatorname{mdf\_}) \wedge \\ \quad \operatorname{strict\text{-}atom}(\operatorname{md7h2\_}, \operatorname{mdw\_}) \end{array}
\end{equation*}

\par\medskip\noindent\textbf{Definition} (\texttt{def-strict-pred}).\;For all $\operatorname{mdt\_}$, $\operatorname{mdhs\_}$ and $\operatorname{mdc\_}$, 
\begin{equation*}
\begin{array}{@{}l@{}} \operatorname{def\text{-}strict\text{-}pred}(\operatorname{mdt\_}, \operatorname{mdhs\_}, \operatorname{mdc\_}) \\ \Leftrightarrow\; \operatorname{mdc\_} \in \operatorname{seq\text{-}set} \wedge \\ \exists \operatorname{mdg\_}, \operatorname{mdl\_}, \operatorname{mdk\_}, \operatorname{mdi\_}. \\ \quad \operatorname{mdc\_} = {[}\operatorname{mdg\_}, \operatorname{enc\text{-}eq}(\operatorname{nth}(\operatorname{mdi\_}, \operatorname{mdl\_}), \operatorname{nth}(\operatorname{mdi\_}, \operatorname{mdl\_}))] \wedge \\ \quad \operatorname{nth}(1, \operatorname{mdl\_}) = \operatorname{econst}(\operatorname{mdk\_}) \wedge \\ \quad \operatorname{pred\text{-}code}(\operatorname{mdk\_}) \wedge \\ \quad \operatorname{mdi\_} \in \mathbb{N} \wedge \\ \quad 2 \leq \operatorname{mdi\_} \wedge \\ \quad \operatorname{mdi\_} \leq \operatorname{length}(\operatorname{mdl\_}) \wedge \\ \quad \operatorname{mdhs\_} = {[}{[}\operatorname{mdg\_}, \operatorname{eapp}(10, \operatorname{mdl\_})]] \end{array}
\end{equation*}

\par\medskip\noindent\textbf{Definition} (\texttt{def-struct-hyp}).\;For all $\operatorname{mdt\_}$, $\operatorname{mdhs\_}$ and $\operatorname{mdc\_}$, 
\begin{equation*}
\begin{array}{@{}l@{}} \operatorname{def\text{-}struct\text{-}hyp}(\operatorname{mdt\_}, \operatorname{mdhs\_}, \operatorname{mdc\_}) \\ \Leftrightarrow\; \operatorname{mdc\_} \in \operatorname{seq\text{-}set} \wedge \\ \exists \operatorname{mdg\_}, \operatorname{mdk\_}, \operatorname{mds\_}. \\ \quad \operatorname{mdc\_} = {[}\operatorname{mdg\_}, \operatorname{enc\text{-}eq}(\operatorname{mds\_}, \operatorname{mds\_})] \wedge \\ \quad \operatorname{pred\text{-}code}(\operatorname{mdk\_}) \wedge \\ \quad \operatorname{mdhs\_} = {[}{[}\operatorname{mdg\_}, \operatorname{enc\text{-}pred1}(\operatorname{econst}(\operatorname{mdk\_}), \operatorname{mds\_})]] \end{array}
\end{equation*}

\par\medskip\noindent\textbf{Definition} (\texttt{def-list}).\;For all $\operatorname{mdt\_}$, $\operatorname{mdhs\_}$ and $\operatorname{mdc\_}$, 
\begin{equation*}
\begin{array}{@{}l@{}} \operatorname{def\text{-}list}(\operatorname{mdt\_}, \operatorname{mdhs\_}, \operatorname{mdc\_}) \\ \Leftrightarrow\; \operatorname{mdc\_} \in \operatorname{seq\text{-}set} \wedge \\ \exists \operatorname{mdg\_}, \operatorname{mdl\_}. \\ \quad \operatorname{mdc\_} = {[}\operatorname{mdg\_}, \operatorname{enc\text{-}eq}(\operatorname{enc\text{-}list}(\operatorname{mdl\_}), \operatorname{enc\text{-}list}(\operatorname{mdl\_}))] \wedge \\ \quad \operatorname{length}(\operatorname{mdhs\_}) = \operatorname{length}(\operatorname{mdl\_}) \wedge \\ \quad \forall \operatorname{mdi\_} \in \mathbb{N}. \\ \quad \quad (1 \leq \operatorname{mdi\_}) \Rightarrow  \\ \quad \quad (\operatorname{mdi\_} \leq \operatorname{length}(\operatorname{mdl\_})) \Rightarrow  \\ \quad \quad \quad \operatorname{nth}(\operatorname{mdi\_}, \operatorname{mdhs\_}) = {[}\operatorname{mdg\_}, \operatorname{enc\text{-}in}(\operatorname{nth}(\operatorname{mdi\_}, \operatorname{mdl\_}), \operatorname{enc\text{-}setc})] \end{array}
\end{equation*}

\par\medskip\noindent\textbf{Definition} (\texttt{def-total}).\;For all $\operatorname{mdt\_}$, $\operatorname{mdhs\_}$ and $\operatorname{mdc\_}$, 
\begin{equation*}
\begin{array}{@{}l@{}} \operatorname{def\text{-}total}(\operatorname{mdt\_}, \operatorname{mdhs\_}, \operatorname{mdc\_}) \\ \Leftrightarrow\; \operatorname{mdc\_} \in \operatorname{seq\text{-}set} \wedge \\ \exists \operatorname{mdg\_}, \operatorname{mdh\_}, \operatorname{mdl\_}. \\ \quad \operatorname{mdc\_} = {[}\operatorname{mdg\_}, \operatorname{enc\text{-}eq}(\operatorname{eapp}(\operatorname{mdh\_}, \operatorname{mdl\_}), \operatorname{eapp}(\operatorname{mdh\_}, \operatorname{mdl\_}))] \wedge \\ \quad \operatorname{total\text{-}head}(\operatorname{mdh\_}) \wedge \\ \quad \operatorname{length}(\operatorname{mdhs\_}) = \operatorname{length}(\operatorname{mdl\_}) \wedge \\ \quad \forall \operatorname{mdi\_} \in \mathbb{N}. \\ \quad \quad (1 \leq \operatorname{mdi\_}) \Rightarrow  \\ \quad \quad (\operatorname{mdi\_} \leq \operatorname{length}(\operatorname{mdl\_})) \Rightarrow  \\ \quad \quad \quad \operatorname{nth}(\operatorname{mdi\_}, \operatorname{mdhs\_}) = {[}\operatorname{mdg\_}, \operatorname{enc\text{-}eq}(\operatorname{nth}(\operatorname{mdi\_}, \operatorname{mdl\_}), \operatorname{nth}(\operatorname{mdi\_}, \operatorname{mdl\_}))] \end{array}
\end{equation*}

\par\medskip\noindent\textbf{Definition} (\texttt{def-binder}).\;For all $\operatorname{mdt\_}$, $\operatorname{mdhs\_}$ and $\operatorname{mdc\_}$, 
\begin{equation*}
\begin{array}{@{}l@{}} \operatorname{def\text{-}binder}(\operatorname{mdt\_}, \operatorname{mdhs\_}, \operatorname{mdc\_}) \\ \Leftrightarrow\; \operatorname{mdc\_} \in \operatorname{seq\text{-}set} \wedge \\ \exists \operatorname{mdg\_}, \operatorname{mdh\_}, \operatorname{mda\_}, \operatorname{mdb\_}. \\ \quad \operatorname{mdc\_} = {[}\operatorname{mdg\_}, \operatorname{enc\text{-}eq}(\operatorname{ebind}(\operatorname{mdh\_}, 1, {[}\operatorname{mda\_}], \operatorname{mdb\_}), \operatorname{ebind}(\operatorname{mdh\_}, 1, {[}\operatorname{mda\_}], \operatorname{mdb\_}))] \wedge \\ \quad (\operatorname{mdh\_} = 22) \vee ((\operatorname{mdh\_} = 24) \vee (\operatorname{mdh\_} = 26)) \wedge \\ \quad \operatorname{nbg\text{-}ok}(\operatorname{mdb\_}) \wedge \\ \quad \operatorname{mdhs\_} = {[}{[}\operatorname{mdg\_}, \operatorname{enc\text{-}eq}(\operatorname{mda\_}, \operatorname{mda\_})]] \end{array}
\end{equation*}

\par\medskip\noindent\textbf{Definition} (\texttt{def-comp}).\;For all $\operatorname{mdt\_}$, $\operatorname{mdhs\_}$ and $\operatorname{mdc\_}$, 
\begin{equation*}
\begin{array}{@{}l@{}} \operatorname{def\text{-}comp}(\operatorname{mdt\_}, \operatorname{mdhs\_}, \operatorname{mdc\_}) \\ \Leftrightarrow\; \operatorname{mdc\_} \in \operatorname{seq\text{-}set} \wedge \\ \exists \operatorname{mdg\_}, \operatorname{mdp\_}. \\ \quad ((\operatorname{mdc\_} = {[}\operatorname{mdg\_}, \operatorname{enc\text{-}eq}(\operatorname{enc\text{-}comp}(\operatorname{mdp\_}), \operatorname{enc\text{-}comp}(\operatorname{mdp\_}))]) \wedge (\operatorname{nbg\text{-}ok}(\operatorname{mdp\_}) \wedge (\operatorname{mdhs\_} = {[}]))) \end{array}
\end{equation*}

\par\medskip\noindent\textbf{Definition} (\texttt{def-nth-lit}).\;For all $\operatorname{mdt\_}$, $\operatorname{mdhs\_}$ and $\operatorname{mdc\_}$, 
\begin{equation*}
\begin{array}{@{}l@{}} \operatorname{def\text{-}nth\text{-}lit}(\operatorname{mdt\_}, \operatorname{mdhs\_}, \operatorname{mdc\_}) \\ \Leftrightarrow\; \operatorname{mdc\_} \in \operatorname{seq\text{-}set} \wedge \\ \exists \operatorname{mdg\_}, \operatorname{mdl\_}, \operatorname{mdk\_}. \\ \quad \operatorname{mdc\_} \\ \quad =\; {[}\operatorname{mdg\_}, \operatorname{enc\text{-}eq}(\operatorname{enc\text{-}nth}(\operatorname{enc\text{-}num}(\operatorname{mdk\_}), \operatorname{enc\text{-}list}(\operatorname{mdl\_})), \operatorname{enc\text{-}nth}(\operatorname{enc\text{-}num}(\operatorname{mdk\_}), \operatorname{enc\text{-}list}(\operatorname{mdl\_})))] \wedge \\ \quad \operatorname{mdk\_} \in \mathbb{N} \wedge \\ \quad 1 \leq \operatorname{mdk\_} \wedge \\ \quad \operatorname{mdk\_} \leq \operatorname{length}(\operatorname{mdl\_}) \wedge \\ \quad \operatorname{mdhs\_} = {[}{[}\operatorname{mdg\_}, \operatorname{enc\text{-}eq}(\operatorname{nth}(\operatorname{mdk\_}, \operatorname{mdl\_}), \operatorname{nth}(\operatorname{mdk\_}, \operatorname{mdl\_}))]] \end{array}
\end{equation*}

\par\medskip\noindent\textbf{Definition} (\texttt{def-nth-list}).\;For all $\operatorname{mdt\_}$, $\operatorname{mdhs\_}$ and $\operatorname{mdc\_}$, 
\begin{equation*}
\begin{array}{@{}l@{}} \operatorname{def\text{-}nth\text{-}list}(\operatorname{mdt\_}, \operatorname{mdhs\_}, \operatorname{mdc\_}) \\ \Leftrightarrow\; \operatorname{mdc\_} \in \operatorname{seq\text{-}set} \wedge \\ \exists \operatorname{mdg\_}, \operatorname{mdl\_}, \operatorname{mdk\_}. \\ \quad \operatorname{mdc\_} \\ \quad =\; {[}\operatorname{mdg\_}, \operatorname{enc\text{-}eq}(\operatorname{enc\text{-}nth}(\operatorname{enc\text{-}num}(\operatorname{mdk\_}), \operatorname{enc\text{-}list}(\operatorname{mdl\_})), \operatorname{enc\text{-}nth}(\operatorname{enc\text{-}num}(\operatorname{mdk\_}), \operatorname{enc\text{-}list}(\operatorname{mdl\_})))] \wedge \\ \quad \operatorname{mdk\_} \in \mathbb{N} \wedge \\ \quad 1 \leq \operatorname{mdk\_} \wedge \\ \quad \operatorname{mdk\_} \leq \operatorname{length}(\operatorname{mdl\_}) \wedge \\ \quad \operatorname{length}(\operatorname{mdhs\_}) = \operatorname{length}(\operatorname{mdl\_}) \wedge \\ \quad \forall \operatorname{mdi\_} \in \mathbb{N}. \\ \quad \quad (1 \leq \operatorname{mdi\_}) \Rightarrow  \\ \quad \quad (\operatorname{mdi\_} \leq \operatorname{length}(\operatorname{mdl\_})) \Rightarrow  \\ \quad \quad \quad \operatorname{nth}(\operatorname{mdi\_}, \operatorname{mdhs\_}) = {[}\operatorname{mdg\_}, \operatorname{enc\text{-}in}(\operatorname{nth}(\operatorname{mdi\_}, \operatorname{mdl\_}), \operatorname{enc\text{-}setc})] \end{array}
\end{equation*}

\par\medskip\noindent\textbf{Definition} (\texttt{def-app}).\;For all $\operatorname{mdt\_}$, $\operatorname{mdhs\_}$ and $\operatorname{mdc\_}$, 
\begin{equation*}
\begin{array}{@{}l@{}} \operatorname{def\text{-}app}(\operatorname{mdt\_}, \operatorname{mdhs\_}, \operatorname{mdc\_}) \\ \Leftrightarrow\; \operatorname{mdc\_} \in \operatorname{seq\text{-}set} \wedge \\ \exists \operatorname{mdg\_}, \operatorname{mdf\_}, \operatorname{mdd\_}, \operatorname{mda\_}. \\ \quad \operatorname{mdc\_} = {[}\operatorname{mdg\_}, \operatorname{enc\text{-}eq}(\operatorname{enc\text{-}apply}(\operatorname{mdf\_}, \operatorname{mda\_}), \operatorname{enc\text{-}apply}(\operatorname{mdf\_}, \operatorname{mda\_}))] \wedge \\ \quad \operatorname{mdhs\_} = {[}{[}\operatorname{mdg\_}, \operatorname{enc\text{-}in}(\operatorname{mdf\_}, \operatorname{enc\text{-}fun1}(\operatorname{mdd\_}))], {[}\operatorname{mdg\_}, \operatorname{enc\text{-}in}(\operatorname{mda\_}, \operatorname{mdd\_})]] \end{array}
\end{equation*}

\par\medskip\noindent\textbf{Definition} (\texttt{def-pred-unfold}).\;For all $\operatorname{mdt\_}$, $\operatorname{mdhs\_}$ and $\operatorname{mdc\_}$, 
\begin{equation*}
\begin{array}{@{}l@{}} \operatorname{def\text{-}pred\text{-}unfold}(\operatorname{mdt\_}, \operatorname{mdhs\_}, \operatorname{mdc\_}) \\ \Leftrightarrow\; \operatorname{mdc\_} \in \operatorname{seq\text{-}set} \wedge \\ \exists \operatorname{mdg\_}, \operatorname{mdk\_}, \operatorname{mds\_}, \operatorname{mdb\_}. \\ \quad \operatorname{mdc\_} = {[}\operatorname{mdg\_}, \operatorname{expr\text{-}open}(\operatorname{mdb\_}, 0, \operatorname{mds\_})] \wedge \\ \quad \operatorname{pred\text{-}code}(\operatorname{mdk\_}) \wedge \\ \quad \operatorname{mdhs\_} \\ \quad =\; {[}{[}\operatorname{mdg\_}, \operatorname{enc\text{-}pred1}(\operatorname{econst}(\operatorname{mdk\_}), \operatorname{mds\_})], {[}\operatorname{mdg\_}, \operatorname{enc\text{-}all}(\operatorname{enc\text{-}iff}(\operatorname{enc\text{-}pred1}(\operatorname{econst}(\operatorname{mdk\_}), \operatorname{ebvar}(0)), \operatorname{mdb\_}))]] \end{array}
\end{equation*}

\par\medskip\noindent\textbf{Definition} (\texttt{def-pred-hyp}).\;For all $\operatorname{mdt\_}$, $\operatorname{mdhs\_}$ and $\operatorname{mdc\_}$, 
\begin{equation*}
\begin{array}{@{}l@{}} \operatorname{def\text{-}pred\text{-}hyp}(\operatorname{mdt\_}, \operatorname{mdhs\_}, \operatorname{mdc\_}) \\ \Leftrightarrow\; \operatorname{mdc\_} \in \operatorname{seq\text{-}set} \wedge \\ \exists \operatorname{mdg\_}, \operatorname{mdk\_}, \operatorname{mds\_}, \operatorname{mdb\_}, \operatorname{md10f\_}, \operatorname{mdw\_}. \\ \quad \operatorname{mdc\_} = {[}\operatorname{mdg\_}, \operatorname{enc\text{-}eq}(\operatorname{mdw\_}, \operatorname{mdw\_})] \wedge \\ \quad \operatorname{pred\text{-}code}(\operatorname{mdk\_}) \wedge \\ \quad \operatorname{mdhs\_} \\ \quad =\; {[}{[}\operatorname{mdg\_}, \operatorname{enc\text{-}pred1}(\operatorname{econst}(\operatorname{mdk\_}), \operatorname{mds\_})], {[}\operatorname{mdg\_}, \operatorname{enc\text{-}all}(\operatorname{enc\text{-}iff}(\operatorname{enc\text{-}pred1}(\operatorname{econst}(\operatorname{mdk\_}), \operatorname{ebvar}(0)), \operatorname{mdb\_}))]] \wedge \\ \quad \operatorname{md10f\_} \in \operatorname{true\text{-}reach}(\operatorname{expr\text{-}open}(\operatorname{mdb\_}, 0, \operatorname{mds\_})) \wedge \\ \quad \operatorname{strict\text{-}atom}(\operatorname{md10f\_}, \operatorname{mdw\_}) \end{array}
\end{equation*}

\par\medskip\noindent\textbf{Definition} (\texttt{def-pred-unfold-n}).\;For all $\operatorname{mdt\_}$, $\operatorname{mdhs\_}$ and $\operatorname{mdc\_}$, 
\begin{equation*}
\begin{array}{@{}l@{}} \operatorname{def\text{-}pred\text{-}unfold\text{-}n}(\operatorname{mdt\_}, \operatorname{mdhs\_}, \operatorname{mdc\_}) \\ \Leftrightarrow\; \operatorname{mdc\_} \in \operatorname{seq\text{-}set} \wedge \\ \exists \operatorname{mdg\_}, \operatorname{mdl\_}, \operatorname{mdk\_}, \operatorname{mdd\_}, \operatorname{mdb\_}. \\ \quad \operatorname{mdc\_} = {[}\operatorname{mdg\_}, \operatorname{mdb\_}] \wedge \\ \quad \operatorname{nth}(1, \operatorname{mdl\_}) = \operatorname{econst}(\operatorname{mdk\_}) \wedge \\ \quad \operatorname{pred\text{-}code}(\operatorname{mdk\_}) \wedge \\ \quad \operatorname{enc\text{-}iff}(\operatorname{eapp}(10, \operatorname{mdl\_}), \operatorname{mdb\_}) \in \operatorname{inst\text{-}reach}(\operatorname{mdd\_}, \operatorname{mdl\_}) \wedge \\ \quad \operatorname{mdhs\_} = {[}{[}\operatorname{mdg\_}, \operatorname{eapp}(10, \operatorname{mdl\_})], {[}\operatorname{mdg\_}, \operatorname{mdd\_}]] \end{array}
\end{equation*}

\par\medskip\noindent\textbf{Definition} (\texttt{def-pred-hyp-n}).\;For all $\operatorname{mdt\_}$, $\operatorname{mdhs\_}$ and $\operatorname{mdc\_}$, 
\begin{equation*}
\begin{array}{@{}l@{}} \operatorname{def\text{-}pred\text{-}hyp\text{-}n}(\operatorname{mdt\_}, \operatorname{mdhs\_}, \operatorname{mdc\_}) \\ \Leftrightarrow\; \operatorname{mdc\_} \in \operatorname{seq\text{-}set} \wedge \\ \exists \operatorname{mdg\_}, \operatorname{mdl\_}, \operatorname{mdk\_}, \operatorname{mdd\_}, \operatorname{mdb\_}, \operatorname{mdw\_}. \\ \quad \operatorname{mdc\_} = {[}\operatorname{mdg\_}, \operatorname{enc\text{-}eq}(\operatorname{mdw\_}, \operatorname{mdw\_})] \wedge \\ \quad \operatorname{nth}(1, \operatorname{mdl\_}) = \operatorname{econst}(\operatorname{mdk\_}) \wedge \\ \quad \operatorname{pred\text{-}code}(\operatorname{mdk\_}) \wedge \\ \quad \operatorname{enc\text{-}iff}(\operatorname{eapp}(10, \operatorname{mdl\_}), \operatorname{mdb\_}) \in \operatorname{inst\text{-}reach}(\operatorname{mdd\_}, \operatorname{mdl\_}) \wedge \\ \quad \operatorname{mdb\_} \in \operatorname{cert\text{-}set}(\operatorname{mdt\_}, \operatorname{mdg\_}, \operatorname{mdw\_}) \wedge \\ \quad \operatorname{mdhs\_} = {[}{[}\operatorname{mdg\_}, \operatorname{eapp}(10, \operatorname{mdl\_})], {[}\operatorname{mdg\_}, \operatorname{mdd\_}]] \end{array}
\end{equation*}

\par\medskip\noindent\textbf{Definition} (\texttt{def-cert-hyp}).\;For all $\operatorname{mdt\_}$, $\operatorname{mdhs\_}$ and $\operatorname{mdc\_}$, 
\begin{equation*}
\begin{array}{@{}l@{}} \operatorname{def\text{-}cert\text{-}hyp}(\operatorname{mdt\_}, \operatorname{mdhs\_}, \operatorname{mdc\_}) \\ \Leftrightarrow\; \operatorname{mdc\_} \in \operatorname{seq\text{-}set} \wedge \\ \exists \operatorname{mdg\_}, \operatorname{mdf\_}, \operatorname{mdw\_}. \\ \quad ((\operatorname{mdc\_} = {[}\operatorname{mdg\_}, \operatorname{enc\text{-}eq}(\operatorname{mdw\_}, \operatorname{mdw\_})]) \wedge ((\operatorname{mdf\_} \in \operatorname{cert\text{-}set}(\operatorname{mdt\_}, \operatorname{mdg\_}, \operatorname{mdw\_})) \wedge (\operatorname{mdhs\_} = {[}{[}\operatorname{mdg\_}, \operatorname{mdf\_}]]))) \end{array}
\end{equation*}

\par\medskip\noindent\textbf{Definition} (\texttt{cert-set}).\;$\operatorname{cert\text{-}set}(\operatorname{md12t\_}, \operatorname{md12cg\_}, \operatorname{md12cw\_}) :=$
\begin{equation*}
\{\operatorname{md12zz\_} \in \operatorname{expr\text{-}set} : \forall \operatorname{md12v\_} \in \operatorname{power}(\operatorname{expr\text{-}set}).\; (\operatorname{cert\text{-}closed}(\operatorname{md12t\_}, \operatorname{md12cg\_}, \operatorname{md12cw\_}, \operatorname{md12v\_}) \Rightarrow (\operatorname{md12zz\_} \in \operatorname{md12v\_}))\}
\end{equation*}

\par\medskip\noindent\textbf{Definition} (\texttt{def-strict-in-r}).\;For all $\operatorname{mdt\_}$, $\operatorname{mdhs\_}$ and $\operatorname{mdc\_}$, 
\begin{equation*}
\begin{array}{@{}l@{}} \operatorname{def\text{-}strict\text{-}in\text{-}r}(\operatorname{mdt\_}, \operatorname{mdhs\_}, \operatorname{mdc\_}) \\ \Leftrightarrow\; \operatorname{mdc\_} \in \operatorname{seq\text{-}set} \wedge \\ \exists \operatorname{mdg\_}, \operatorname{mdw\_}, \operatorname{mda\_}. \\ \quad (\operatorname{mdc\_} = {[}\operatorname{mdg\_}, \operatorname{enc\text{-}eq}(\operatorname{mda\_}, \operatorname{mda\_})]) \wedge (\operatorname{mdhs\_} = {[}{[}\operatorname{mdg\_}, \operatorname{enc\text{-}in}(\operatorname{mdw\_}, \operatorname{mda\_})]]) \end{array}
\end{equation*}

\par\medskip\noindent\textbf{Definition} (\texttt{def-strict-dom}).\;For all $\operatorname{mdt\_}$, $\operatorname{mdhs\_}$ and $\operatorname{mdc\_}$, 
\begin{equation*}
\begin{array}{@{}l@{}} \operatorname{def\text{-}strict\text{-}dom}(\operatorname{mdt\_}, \operatorname{mdhs\_}, \operatorname{mdc\_}) \\ \Leftrightarrow\; \operatorname{mdc\_} \in \operatorname{seq\text{-}set} \wedge \\ \exists \operatorname{mdg\_}, \operatorname{mdh\_}, \operatorname{mda\_}, \operatorname{mdb\_}. \\ \quad \operatorname{mdc\_} = {[}\operatorname{mdg\_}, \operatorname{enc\text{-}eq}(\operatorname{mda\_}, \operatorname{mda\_})] \wedge \\ \quad (\operatorname{mdh\_} = 22) \vee ((\operatorname{mdh\_} = 24) \vee (\operatorname{mdh\_} = 26)) \wedge \\ \quad \operatorname{mdhs\_} = {[}{[}\operatorname{mdg\_}, \operatorname{enc\text{-}eq}(\operatorname{ebind}(\operatorname{mdh\_}, 1, {[}\operatorname{mda\_}], \operatorname{mdb\_}), \operatorname{ebind}(\operatorname{mdh\_}, 1, {[}\operatorname{mda\_}], \operatorname{mdb\_}))]] \end{array}
\end{equation*}

\par\medskip\noindent\textbf{Definition} (\texttt{strict-atom}).\;For all $\operatorname{md7af\_}$ and $\operatorname{md7at\_}$, 
\begin{equation*}
\begin{array}{@{}l@{}} \operatorname{strict\text{-}atom}(\operatorname{md7af\_}, \operatorname{md7at\_}) \\ \Leftrightarrow\; \exists \operatorname{md7ax\_}, \operatorname{md7ay\_}. \\ \quad (\operatorname{md7af\_} = \operatorname{enc\text{-}in}(\operatorname{md7ax\_}, \operatorname{md7ay\_})) \vee (\operatorname{md7af\_} = \operatorname{enc\text{-}eq}(\operatorname{md7ax\_}, \operatorname{md7ay\_})) \wedge \\ \quad (\operatorname{md7at\_} \in \operatorname{strict\text{-}reach}(\operatorname{md7ax\_})) \vee (\operatorname{md7at\_} \in \operatorname{strict\text{-}reach}(\operatorname{md7ay\_})) \vee \\ \exists \operatorname{md7pl\_}, \operatorname{md7pk\_}, \operatorname{md7pi\_}. \\ \quad \operatorname{md7af\_} = \operatorname{eapp}(10, \operatorname{md7pl\_}) \wedge \\ \quad \operatorname{nth}(1, \operatorname{md7pl\_}) = \operatorname{econst}(\operatorname{md7pk\_}) \wedge \\ \quad \operatorname{pred\text{-}code}(\operatorname{md7pk\_}) \wedge \\ \quad \operatorname{md7pi\_} \in \mathbb{N} \wedge \\ \quad 2 \leq \operatorname{md7pi\_} \wedge \\ \quad \operatorname{md7pi\_} \leq \operatorname{length}(\operatorname{md7pl\_}) \wedge \\ \quad \operatorname{md7at\_} \in \operatorname{strict\text{-}reach}(\operatorname{nth}(\operatorname{md7pi\_}, \operatorname{md7pl\_})) \end{array}
\end{equation*}

\par\medskip\noindent\textbf{Definition} (\texttt{enc-le}).\;$\operatorname{enc\text{-}le}(\operatorname{md10x\_}, \operatorname{md10y\_}) :=$ $\operatorname{eapp}(10, {[}\operatorname{econst}(11), \operatorname{md10x\_}, \operatorname{md10y\_}])$
\par\medskip\noindent\textbf{Definition} (\texttt{enc-lt}).\;$\operatorname{enc\text{-}lt}(\operatorname{md10x\_}, \operatorname{md10y\_}) :=$ $\operatorname{eapp}(10, {[}\operatorname{econst}(13), \operatorname{md10x\_}, \operatorname{md10y\_}])$
\par\medskip\noindent\textbf{Definition} (\texttt{true-reach}).\;$\operatorname{true\text{-}reach}(\operatorname{md7j\_}) :=$
\begin{equation*}
\{\operatorname{md7zz\_} \in \operatorname{expr\text{-}set} : \forall \operatorname{md7ww\_} \in \operatorname{power}(\operatorname{expr\text{-}set}).\; ((\operatorname{md7j\_} \in \operatorname{md7ww\_}) \Rightarrow (\operatorname{true\text{-}closed}(\operatorname{md7ww\_}) \Rightarrow (\operatorname{md7zz\_} \in \operatorname{md7ww\_})))\}
\end{equation*}

\par\medskip\noindent\textbf{Definition} (\texttt{strict-step}).\;For all $\operatorname{md7s\_}$ and $\operatorname{md7t\_}$, 
\begin{equation*}
\begin{array}{@{}l@{}} \operatorname{strict\text{-}step}(\operatorname{md7s\_}, \operatorname{md7t\_}) \\ \Leftrightarrow\; \exists \operatorname{md7h\_}, \operatorname{md7l\_}, \operatorname{md7i\_}. \\ \quad \operatorname{md7s\_} = \operatorname{eapp}(\operatorname{md7h\_}, \operatorname{md7l\_}) \wedge \\ \quad \operatorname{strict\text{-}head}(\operatorname{md7h\_}) \wedge \\ \quad \operatorname{md7i\_} \in \mathbb{N} \wedge \\ \quad 1 \leq \operatorname{md7i\_} \wedge \\ \quad \operatorname{md7i\_} \leq \operatorname{length}(\operatorname{md7l\_}) \wedge \\ \quad \operatorname{md7t\_} = \operatorname{nth}(\operatorname{md7i\_}, \operatorname{md7l\_}) \vee \\ \exists \operatorname{md7k\_}, \operatorname{md7a\_}, \operatorname{md7b\_}. \\ \quad (\operatorname{md7k\_} = 22) \vee ((\operatorname{md7k\_} = 24) \vee (\operatorname{md7k\_} = 26)) \wedge \\ \quad \operatorname{md7s\_} = \operatorname{ebind}(\operatorname{md7k\_}, 1, {[}\operatorname{md7a\_}], \operatorname{md7b\_}) \wedge \\ \quad \operatorname{md7t\_} = \operatorname{md7a\_} \end{array}
\end{equation*}

\par\medskip\noindent\textbf{Definition} (\texttt{strict-reach}).\;$\operatorname{strict\text{-}reach}(\operatorname{md7e\_}) :=$
\begin{equation*}
\{\operatorname{md7z\_} \in \operatorname{expr\text{-}set} : \forall \operatorname{md7w\_} \in \operatorname{power}(\operatorname{expr\text{-}set}).\; ((\operatorname{md7e\_} \in \operatorname{md7w\_}) \Rightarrow (\operatorname{strict\text{-}closed}(\operatorname{md7w\_}) \Rightarrow (\operatorname{md7z\_} \in \operatorname{md7w\_})))\}
\end{equation*}

\par\medskip\noindent\textbf{Definition} (\texttt{enc-fun1}).\;$\operatorname{enc\text{-}fun1}(\operatorname{mdp\_}) :=$ $\operatorname{eapp}(19, {[}\operatorname{mdp\_}])$
\par\medskip\noindent\textbf{Definition} (\texttt{adm-def-strict-in-l}).\;For every $\operatorname{mdt\_}$, 
\begin{equation*}
\begin{array}{@{}l@{}} \operatorname{adm\text{-}def\text{-}strict\text{-}in\text{-}l}(\operatorname{mdt\_}) \\ \Leftrightarrow\; \forall \operatorname{mdhs\_}, \operatorname{mdc\_}. \\ \quad \operatorname{def\text{-}strict\text{-}in\text{-}l}(\operatorname{mdt\_}, \operatorname{mdhs\_}, \operatorname{mdc\_}) \Rightarrow  \\ \quad \operatorname{hyps\text{-}in}(\operatorname{mdhs\_}, \operatorname{derivable}(\operatorname{mdt\_})) \Rightarrow (\operatorname{mdc\_} \in \operatorname{derivable}(\operatorname{mdt\_})) \end{array}
\end{equation*}

\par\medskip\noindent\textbf{Definition} (\texttt{adm-def-comp}).\;For every $\operatorname{mdt\_}$, 
\begin{equation*}
\begin{array}{@{}l@{}} \operatorname{adm\text{-}def\text{-}comp}(\operatorname{mdt\_}) \\ \Leftrightarrow\; \forall \operatorname{mdhs\_}, \operatorname{mdc\_}. \\ \quad \operatorname{def\text{-}comp}(\operatorname{mdt\_}, \operatorname{mdhs\_}, \operatorname{mdc\_}) \Rightarrow  \\ \quad \operatorname{hyps\text{-}in}(\operatorname{mdhs\_}, \operatorname{derivable}(\operatorname{mdt\_})) \Rightarrow (\operatorname{mdc\_} \in \operatorname{derivable}(\operatorname{mdt\_})) \end{array}
\end{equation*}

\par\medskip\noindent\textbf{Definition} (\texttt{adm-def-strict-in-r}).\;For every $\operatorname{mdt\_}$, 
\begin{equation*}
\begin{array}{@{}l@{}} \operatorname{adm\text{-}def\text{-}strict\text{-}in\text{-}r}(\operatorname{mdt\_}) \\ \Leftrightarrow\; \forall \operatorname{mdhs\_}, \operatorname{mdc\_}. \\ \quad \operatorname{def\text{-}strict\text{-}in\text{-}r}(\operatorname{mdt\_}, \operatorname{mdhs\_}, \operatorname{mdc\_}) \Rightarrow  \\ \quad \operatorname{hyps\text{-}in}(\operatorname{mdhs\_}, \operatorname{derivable}(\operatorname{mdt\_})) \Rightarrow (\operatorname{mdc\_} \in \operatorname{derivable}(\operatorname{mdt\_})) \end{array}
\end{equation*}

\par\medskip\noindent\textbf{Definition} (\texttt{adm-def-strict-arg}).\;For every $\operatorname{mdt\_}$, 
\begin{equation*}
\begin{array}{@{}l@{}} \operatorname{adm\text{-}def\text{-}strict\text{-}arg}(\operatorname{mdt\_}) \\ \Leftrightarrow\; \forall \operatorname{mdhs\_}, \operatorname{mdc\_}. \\ \quad \operatorname{def\text{-}strict\text{-}arg}(\operatorname{mdt\_}, \operatorname{mdhs\_}, \operatorname{mdc\_}) \Rightarrow  \\ \quad \operatorname{hyps\text{-}in}(\operatorname{mdhs\_}, \operatorname{derivable}(\operatorname{mdt\_})) \Rightarrow (\operatorname{mdc\_} \in \operatorname{derivable}(\operatorname{mdt\_})) \end{array}
\end{equation*}

\par\medskip\noindent\textbf{Definition} (\texttt{adm-def-strict-dom}).\;For every $\operatorname{mdt\_}$, 
\begin{equation*}
\begin{array}{@{}l@{}} \operatorname{adm\text{-}def\text{-}strict\text{-}dom}(\operatorname{mdt\_}) \\ \Leftrightarrow\; \forall \operatorname{mdhs\_}, \operatorname{mdc\_}. \\ \quad \operatorname{def\text{-}strict\text{-}dom}(\operatorname{mdt\_}, \operatorname{mdhs\_}, \operatorname{mdc\_}) \Rightarrow  \\ \quad \operatorname{hyps\text{-}in}(\operatorname{mdhs\_}, \operatorname{derivable}(\operatorname{mdt\_})) \Rightarrow (\operatorname{mdc\_} \in \operatorname{derivable}(\operatorname{mdt\_})) \end{array}
\end{equation*}

\par\medskip\noindent\textbf{Definition} (\texttt{ax-fun-app}).\;
\begin{equation*}
(\operatorname{ax\text{-}fun\text{-}app} \simeq \operatorname{enc\text{-}all}(\operatorname{enc\text{-}all}(\operatorname{enc\text{-}all}(\operatorname{enc\text{-}imp}(\operatorname{enc\text{-}in}(\operatorname{ebvar}(1), \operatorname{enc\text{-}fun1}(\operatorname{ebvar}(2))), \operatorname{enc\text{-}iff}(\operatorname{enc\text{-}eq}(\operatorname{enc\text{-}apply}(\operatorname{ebvar}(1), \operatorname{ebvar}(0)), \operatorname{enc\text{-}apply}(\operatorname{ebvar}(1), \operatorname{ebvar}(0))), \operatorname{enc\text{-}in}(\operatorname{ebvar}(0), \operatorname{ebvar}(2))))))))
\end{equation*}

\par\medskip\noindent\textbf{Definition} (\texttt{def-len-tuples}).\;For all $\operatorname{mdt\_}$, $\operatorname{mdhs\_}$ and $\operatorname{mdc\_}$, 
\begin{equation*}
\begin{array}{@{}l@{}} \operatorname{def\text{-}len\text{-}tuples}(\operatorname{mdt\_}, \operatorname{mdhs\_}, \operatorname{mdc\_}) \\ \Leftrightarrow\; \operatorname{mdc\_} \in \operatorname{seq\text{-}set} \wedge \\ \exists \operatorname{mdg\_}, \operatorname{mdw\_}, \operatorname{mda\_}. \\ \quad \operatorname{mdc\_} = {[}\operatorname{mdg\_}, \operatorname{enc\text{-}eq}(\operatorname{enc\text{-}length}(\operatorname{mdw\_}), \operatorname{enc\text{-}length}(\operatorname{mdw\_}))] \wedge \\ \quad \operatorname{mdhs\_} = {[}{[}\operatorname{mdg\_}, \operatorname{enc\text{-}in}(\operatorname{mdw\_}, \operatorname{enc\text{-}tuples}(\operatorname{mda\_}))]] \end{array}
\end{equation*}

\par\medskip\noindent\textbf{Definition} (\texttt{ax-len-nn}).\;
\begin{equation*}
(\operatorname{ax\text{-}len\text{-}nn} \simeq \operatorname{enc\text{-}all}(\operatorname{enc\text{-}all}(\operatorname{enc\text{-}imp}(\operatorname{enc\text{-}in}(\operatorname{ebvar}(0), \operatorname{enc\text{-}tuples}(\operatorname{ebvar}(1))), \operatorname{enc\text{-}in}(\operatorname{enc\text{-}length}(\operatorname{ebvar}(0)), \operatorname{enc\text{-}nnc})))))
\end{equation*}

\par\medskip\noindent\textbf{Definition} (\texttt{strict-closed}).\;For every $\operatorname{md7x\_}$, 
\begin{equation*}
\begin{array}{@{}l@{}} \operatorname{strict\text{-}closed}(\operatorname{md7x\_}) \\ \Leftrightarrow\; \forall \operatorname{md7p\_} \in \operatorname{md7x\_}, \operatorname{md7q\_} \in \mathbf{Set}. \\ \quad \operatorname{strict\text{-}step}(\operatorname{md7p\_}, \operatorname{md7q\_}) \Rightarrow (\operatorname{md7q\_} \in \operatorname{md7x\_}) \end{array}
\end{equation*}

\par\medskip\noindent\textbf{Definition} (\texttt{true-step}).\;For all $\operatorname{md7f\_}$ and $\operatorname{md7g\_}$, 
\begin{equation*}
\begin{array}{@{}l@{}} \operatorname{true\text{-}step}(\operatorname{md7f\_}, \operatorname{md7g\_}) \\ \Leftrightarrow\; \exists \operatorname{md7c\_}, \operatorname{md7d\_}. \\ \quad (\operatorname{md7f\_} = \operatorname{enc\text{-}and}(\operatorname{md7c\_}, \operatorname{md7d\_})) \wedge ((\operatorname{md7g\_} = \operatorname{md7c\_}) \vee (\operatorname{md7g\_} = \operatorname{md7d\_})) \vee \\ \exists \operatorname{md7y\_}, \operatorname{md7o\_}, \operatorname{md7r\_}. \\ \quad \operatorname{md7f\_} = \operatorname{enc\text{-}in}(\operatorname{md7y\_}, \operatorname{enc\text{-}sep}(\operatorname{md7o\_}, \operatorname{md7r\_})) \wedge \\ \quad \operatorname{nbg\text{-}ok}(\operatorname{md7r\_}) \wedge \\ \quad \operatorname{md7g\_} = \operatorname{expr\text{-}open}(\operatorname{md7r\_}, 0, \operatorname{md7y\_}) \end{array}
\end{equation*}

\par\medskip\noindent\textbf{Definition} (\texttt{true-closed}).\;For every $\operatorname{md7v\_}$, 
\begin{equation*}
\begin{array}{@{}l@{}} \operatorname{true\text{-}closed}(\operatorname{md7v\_}) \\ \Leftrightarrow\; \forall \operatorname{md7m\_} \in \operatorname{md7v\_}, \operatorname{md7n\_} \in \mathbf{Set}. \\ \quad \operatorname{true\text{-}step}(\operatorname{md7m\_}, \operatorname{md7n\_}) \Rightarrow (\operatorname{md7n\_} \in \operatorname{md7v\_}) \end{array}
\end{equation*}

\par\medskip\noindent\textbf{Definition} (\texttt{adm-def-strict-hyp}).\;For every $\operatorname{mdt\_}$, 
\begin{equation*}
\begin{array}{@{}l@{}} \operatorname{adm\text{-}def\text{-}strict\text{-}hyp}(\operatorname{mdt\_}) \\ \Leftrightarrow\; \forall \operatorname{mdhs\_}, \operatorname{mdc\_}. \\ \quad \operatorname{def\text{-}strict\text{-}hyp}(\operatorname{mdt\_}, \operatorname{mdhs\_}, \operatorname{mdc\_}) \Rightarrow  \\ \quad \operatorname{hyps\text{-}in}(\operatorname{mdhs\_}, \operatorname{derivable}(\operatorname{mdt\_})) \Rightarrow (\operatorname{mdc\_} \in \operatorname{derivable}(\operatorname{mdt\_})) \end{array}
\end{equation*}

\par\medskip\noindent\textbf{Definition} (\texttt{adm-def-nth-list}).\;For every $\operatorname{mdt\_}$, 
\begin{equation*}
\begin{array}{@{}l@{}} \operatorname{adm\text{-}def\text{-}nth\text{-}list}(\operatorname{mdt\_}) \\ \Leftrightarrow\; \forall \operatorname{mdhs\_}, \operatorname{mdc\_}. \\ \quad \operatorname{def\text{-}nth\text{-}list}(\operatorname{mdt\_}, \operatorname{mdhs\_}, \operatorname{mdc\_}) \Rightarrow  \\ \quad \operatorname{hyps\text{-}in}(\operatorname{mdhs\_}, \operatorname{derivable}(\operatorname{mdt\_})) \Rightarrow (\operatorname{mdc\_} \in \operatorname{derivable}(\operatorname{mdt\_})) \end{array}
\end{equation*}

\par\medskip\noindent\textbf{Definition} (\texttt{adm-def-strict-pred}).\;For every $\operatorname{mdt\_}$, 
\begin{equation*}
\begin{array}{@{}l@{}} \operatorname{adm\text{-}def\text{-}strict\text{-}pred}(\operatorname{mdt\_}) \\ \Leftrightarrow\; \forall \operatorname{mdhs\_}, \operatorname{mdc\_}. \\ \quad \operatorname{def\text{-}strict\text{-}pred}(\operatorname{mdt\_}, \operatorname{mdhs\_}, \operatorname{mdc\_}) \Rightarrow  \\ \quad \operatorname{hyps\text{-}in}(\operatorname{mdhs\_}, \operatorname{derivable}(\operatorname{mdt\_})) \Rightarrow (\operatorname{mdc\_} \in \operatorname{derivable}(\operatorname{mdt\_})) \end{array}
\end{equation*}

\par\medskip\noindent\textbf{Definition} (\texttt{enc-pred1}).\;$\operatorname{enc\text{-}pred1}(\operatorname{md10n\_}, \operatorname{md10p\_}) :=$ $\operatorname{eapp}(10, {[}\operatorname{md10n\_}, \operatorname{md10p\_}])$
\par\medskip\noindent\textbf{Definition} (\texttt{adm-def-struct-hyp}).\;For every $\operatorname{mdt\_}$, 
\begin{equation*}
\begin{array}{@{}l@{}} \operatorname{adm\text{-}def\text{-}struct\text{-}hyp}(\operatorname{mdt\_}) \\ \Leftrightarrow\; \forall \operatorname{mdhs\_}, \operatorname{mdc\_}. \\ \quad \operatorname{def\text{-}struct\text{-}hyp}(\operatorname{mdt\_}, \operatorname{mdhs\_}, \operatorname{mdc\_}) \Rightarrow  \\ \quad \operatorname{hyps\text{-}in}(\operatorname{mdhs\_}, \operatorname{derivable}(\operatorname{mdt\_})) \Rightarrow (\operatorname{mdc\_} \in \operatorname{derivable}(\operatorname{mdt\_})) \end{array}
\end{equation*}

\par\medskip\noindent\textbf{Definition} (\texttt{adm-def-pred-unfold}).\;For every $\operatorname{mdt\_}$, 
\begin{equation*}
\begin{array}{@{}l@{}} \operatorname{adm\text{-}def\text{-}pred\text{-}unfold}(\operatorname{mdt\_}) \\ \Leftrightarrow\; \forall \operatorname{mdhs\_}, \operatorname{mdc\_}. \\ \quad \operatorname{def\text{-}pred\text{-}unfold}(\operatorname{mdt\_}, \operatorname{mdhs\_}, \operatorname{mdc\_}) \Rightarrow  \\ \quad \operatorname{hyps\text{-}in}(\operatorname{mdhs\_}, \operatorname{derivable}(\operatorname{mdt\_})) \Rightarrow (\operatorname{mdc\_} \in \operatorname{derivable}(\operatorname{mdt\_})) \end{array}
\end{equation*}

\par\medskip\noindent\textbf{Definition} (\texttt{adm-def-pred-hyp}).\;For every $\operatorname{mdt\_}$, 
\begin{equation*}
\begin{array}{@{}l@{}} \operatorname{adm\text{-}def\text{-}pred\text{-}hyp}(\operatorname{mdt\_}) \\ \Leftrightarrow\; \forall \operatorname{mdhs\_}, \operatorname{mdc\_}. \\ \quad \operatorname{def\text{-}pred\text{-}hyp}(\operatorname{mdt\_}, \operatorname{mdhs\_}, \operatorname{mdc\_}) \Rightarrow  \\ \quad \operatorname{hyps\text{-}in}(\operatorname{mdhs\_}, \operatorname{derivable}(\operatorname{mdt\_})) \Rightarrow (\operatorname{mdc\_} \in \operatorname{derivable}(\operatorname{mdt\_})) \end{array}
\end{equation*}

\par\medskip\noindent\textbf{Definition} (\texttt{def-fun-inst}).\;For all $\operatorname{mdt\_}$, $\operatorname{mdhs\_}$ and $\operatorname{mdc\_}$, 
\begin{equation*}
\begin{array}{@{}l@{}} \operatorname{def\text{-}fun\text{-}inst}(\operatorname{mdt\_}, \operatorname{mdhs\_}, \operatorname{mdc\_}) \\ \Leftrightarrow\; \operatorname{mdc\_} \in \operatorname{seq\text{-}set} \wedge \\ \exists \operatorname{mdg\_}, \operatorname{mdh\_}, \operatorname{mdl\_}, \operatorname{mdb\_}. \\ \quad \operatorname{mdc\_} = {[}\operatorname{mdg\_}, \operatorname{enc\text{-}eq}(\operatorname{eapp}(\operatorname{mdh\_}, \operatorname{mdl\_}), \operatorname{eapp}(\operatorname{mdh\_}, \operatorname{mdl\_}))] \wedge \\ \quad \operatorname{fun\text{-}def}(\operatorname{mdh\_}, \operatorname{length}(\operatorname{mdl\_}), \operatorname{mdb\_}) \wedge \\ \quad \operatorname{mdhs\_} = {[}{[}\operatorname{mdg\_}, \operatorname{enc\text{-}eq}(\operatorname{openl}(\operatorname{mdb\_}, \operatorname{mdl\_}, \operatorname{length}(\operatorname{mdl\_})), \operatorname{openl}(\operatorname{mdb\_}, \operatorname{mdl\_}, \operatorname{length}(\operatorname{mdl\_})))]] \end{array}
\end{equation*}

\par\medskip\noindent\textbf{Definition} (\texttt{def-fun-nth}).\;For all $\operatorname{mdt\_}$, $\operatorname{mdhs\_}$ and $\operatorname{mdc\_}$, 
\begin{equation*}
\begin{array}{@{}l@{}} \operatorname{def\text{-}fun\text{-}nth}(\operatorname{mdt\_}, \operatorname{mdhs\_}, \operatorname{mdc\_}) \\ \Leftrightarrow\; \operatorname{mdc\_} \in \operatorname{seq\text{-}set} \wedge \\ \exists \operatorname{mdg\_}, \operatorname{mdh\_}, \operatorname{mdl\_}, \operatorname{mdb\_}, \operatorname{mdk\_}. \\ \quad \operatorname{mdc\_} \\ \quad =\; {[}\operatorname{mdg\_}, \operatorname{enc\text{-}eq}(\operatorname{enc\text{-}nth}(\operatorname{enc\text{-}num}(\operatorname{mdk\_}), \operatorname{eapp}(\operatorname{mdh\_}, \operatorname{mdl\_})), \operatorname{enc\text{-}nth}(\operatorname{enc\text{-}num}(\operatorname{mdk\_}), \operatorname{eapp}(\operatorname{mdh\_}, \operatorname{mdl\_})))] \wedge \\ \quad \operatorname{fun\text{-}def}(\operatorname{mdh\_}, \operatorname{length}(\operatorname{mdl\_}), \operatorname{mdb\_}) \wedge \\ \quad \operatorname{mdk\_} \in \mathbb{N} \wedge \\ \quad \operatorname{mdhs\_} \\ \quad =\; {[}{[}\operatorname{mdg\_}, \operatorname{enc\text{-}eq}(\operatorname{enc\text{-}nth}(\operatorname{enc\text{-}num}(\operatorname{mdk\_}), \operatorname{openl}(\operatorname{mdb\_}, \operatorname{mdl\_}, \operatorname{length}(\operatorname{mdl\_}))), \operatorname{enc\text{-}nth}(\operatorname{enc\text{-}num}(\operatorname{mdk\_}), \operatorname{openl}(\operatorname{mdb\_}, \operatorname{mdl\_}, \operatorname{length}(\operatorname{mdl\_}))))]] \end{array}
\end{equation*}

\par\medskip\noindent\textbf{Definition} (\texttt{def-fun-strict}).\;For all $\operatorname{mdt\_}$, $\operatorname{mdhs\_}$ and $\operatorname{mdc\_}$, 
\begin{equation*}
\begin{array}{@{}l@{}} \operatorname{def\text{-}fun\text{-}strict}(\operatorname{mdt\_}, \operatorname{mdhs\_}, \operatorname{mdc\_}) \\ \Leftrightarrow\; \operatorname{mdc\_} \in \operatorname{seq\text{-}set} \wedge \\ \exists \operatorname{mdg\_}, \operatorname{mdh\_}, \operatorname{mdl\_}, \operatorname{mdb\_}. \\ \quad \operatorname{mdc\_} = {[}\operatorname{mdg\_}, \operatorname{enc\text{-}eq}(\operatorname{openl}(\operatorname{mdb\_}, \operatorname{mdl\_}, \operatorname{length}(\operatorname{mdl\_})), \operatorname{openl}(\operatorname{mdb\_}, \operatorname{mdl\_}, \operatorname{length}(\operatorname{mdl\_})))] \wedge \\ \quad \operatorname{fun\text{-}def}(\operatorname{mdh\_}, \operatorname{length}(\operatorname{mdl\_}), \operatorname{mdb\_}) \wedge \\ \quad \operatorname{mdhs\_} = {[}{[}\operatorname{mdg\_}, \operatorname{enc\text{-}eq}(\operatorname{eapp}(\operatorname{mdh\_}, \operatorname{mdl\_}), \operatorname{eapp}(\operatorname{mdh\_}, \operatorname{mdl\_}))]] \end{array}
\end{equation*}

\par\medskip\noindent\textbf{Definition} (\texttt{def-fun-sep}).\;For all $\operatorname{mdt\_}$, $\operatorname{mdhs\_}$ and $\operatorname{mdc\_}$, 
\begin{equation*}
\begin{array}{@{}l@{}} \operatorname{def\text{-}fun\text{-}sep}(\operatorname{mdt\_}, \operatorname{mdhs\_}, \operatorname{mdc\_}) \\ \Leftrightarrow\; \operatorname{mdc\_} \in \operatorname{seq\text{-}set} \wedge \\ \exists \operatorname{mdg\_}, \operatorname{mdh\_}, \operatorname{mdl\_}, \operatorname{mdb\_}, \operatorname{md11y\_}. \\ \quad \operatorname{mdc\_} = {[}\operatorname{mdg\_}, \operatorname{enc\text{-}in}(\operatorname{md11y\_}, \operatorname{openl}(\operatorname{mdb\_}, \operatorname{mdl\_}, \operatorname{length}(\operatorname{mdl\_})))] \wedge \\ \quad \operatorname{fun\text{-}def}(\operatorname{mdh\_}, \operatorname{length}(\operatorname{mdl\_}), \operatorname{mdb\_}) \wedge \\ \quad \operatorname{mdhs\_} = {[}{[}\operatorname{mdg\_}, \operatorname{enc\text{-}in}(\operatorname{md11y\_}, \operatorname{eapp}(\operatorname{mdh\_}, \operatorname{mdl\_}))]] \end{array}
\end{equation*}

\par\medskip\noindent\textbf{Definition} (\texttt{adm-def-fun-inst}).\;For every $\operatorname{mdt\_}$, 
\begin{equation*}
\begin{array}{@{}l@{}} \operatorname{adm\text{-}def\text{-}fun\text{-}inst}(\operatorname{mdt\_}) \\ \Leftrightarrow\; \forall \operatorname{mdhs\_}, \operatorname{mdc\_}. \\ \quad \operatorname{def\text{-}fun\text{-}inst}(\operatorname{mdt\_}, \operatorname{mdhs\_}, \operatorname{mdc\_}) \Rightarrow  \\ \quad \operatorname{hyps\text{-}in}(\operatorname{mdhs\_}, \operatorname{derivable}(\operatorname{mdt\_})) \Rightarrow (\operatorname{mdc\_} \in \operatorname{derivable}(\operatorname{mdt\_})) \end{array}
\end{equation*}

\par\medskip\noindent\textbf{Definition} (\texttt{adm-def-fun-strict}).\;For every $\operatorname{mdt\_}$, 
\begin{equation*}
\begin{array}{@{}l@{}} \operatorname{adm\text{-}def\text{-}fun\text{-}strict}(\operatorname{mdt\_}) \\ \Leftrightarrow\; \forall \operatorname{mdhs\_}, \operatorname{mdc\_}. \\ \quad \operatorname{def\text{-}fun\text{-}strict}(\operatorname{mdt\_}, \operatorname{mdhs\_}, \operatorname{mdc\_}) \Rightarrow  \\ \quad \operatorname{hyps\text{-}in}(\operatorname{mdhs\_}, \operatorname{derivable}(\operatorname{mdt\_})) \Rightarrow (\operatorname{mdc\_} \in \operatorname{derivable}(\operatorname{mdt\_})) \end{array}
\end{equation*}

\par\medskip\noindent\textbf{Definition} (\texttt{adm-def-fun-nth}).\;For every $\operatorname{mdt\_}$, 
\begin{equation*}
\begin{array}{@{}l@{}} \operatorname{adm\text{-}def\text{-}fun\text{-}nth}(\operatorname{mdt\_}) \\ \Leftrightarrow\; \forall \operatorname{mdhs\_}, \operatorname{mdc\_}. \\ \quad \operatorname{def\text{-}fun\text{-}nth}(\operatorname{mdt\_}, \operatorname{mdhs\_}, \operatorname{mdc\_}) \Rightarrow  \\ \quad \operatorname{hyps\text{-}in}(\operatorname{mdhs\_}, \operatorname{derivable}(\operatorname{mdt\_})) \Rightarrow (\operatorname{mdc\_} \in \operatorname{derivable}(\operatorname{mdt\_})) \end{array}
\end{equation*}

\par\medskip\noindent\textbf{Definition} (\texttt{adm-def-fun-sep}).\;For every $\operatorname{mdt\_}$, 
\begin{equation*}
\begin{array}{@{}l@{}} \operatorname{adm\text{-}def\text{-}fun\text{-}sep}(\operatorname{mdt\_}) \\ \Leftrightarrow\; \forall \operatorname{mdhs\_}, \operatorname{mdc\_}. \\ \quad \operatorname{def\text{-}fun\text{-}sep}(\operatorname{mdt\_}, \operatorname{mdhs\_}, \operatorname{mdc\_}) \Rightarrow  \\ \quad \operatorname{hyps\text{-}in}(\operatorname{mdhs\_}, \operatorname{derivable}(\operatorname{mdt\_})) \Rightarrow (\operatorname{mdc\_} \in \operatorname{derivable}(\operatorname{mdt\_})) \end{array}
\end{equation*}

\par\medskip\noindent\textbf{Definition} (\texttt{inst-reach}).\;$\operatorname{inst\text{-}reach}(\operatorname{md12d\_}, \operatorname{md12l\_}) :=$
\begin{equation*}
\{\operatorname{md12z\_} \in \operatorname{expr\text{-}set} : \forall \operatorname{md12w\_} \in \operatorname{power}(\operatorname{expr\text{-}set}).\; ((\operatorname{md12d\_} \in \operatorname{md12w\_}) \Rightarrow (\operatorname{inst\text{-}closed}(\operatorname{md12l\_}, \operatorname{md12w\_}) \Rightarrow (\operatorname{md12z\_} \in \operatorname{md12w\_})))\}
\end{equation*}

\par\medskip\noindent\textbf{Definition} (\texttt{cert-step}).\;For all $\operatorname{md12t\_}$, $\operatorname{md12cg\_}$, $\operatorname{md12cw\_}$, $\operatorname{md12x\_}$ and $\operatorname{md12f\_}$, 
\begin{equation*}
\begin{array}{@{}l@{}} \operatorname{cert\text{-}step}(\operatorname{md12t\_}, \operatorname{md12cg\_}, \operatorname{md12cw\_}, \operatorname{md12x\_}, \operatorname{md12f\_}) \\ \Leftrightarrow\; \operatorname{strict\text{-}atom}(\operatorname{md12f\_}, \operatorname{md12cw\_}) \vee \\ \exists \operatorname{md12f2\_} \in \operatorname{true\text{-}reach}(\operatorname{md12f\_}).\; (\operatorname{md12f2\_} \in \operatorname{md12x\_}) \vee \\ \exists \operatorname{md12b\_}, \operatorname{md12y\_}. \\ \quad \operatorname{md12f\_} = \operatorname{enc\text{-}ex}(\operatorname{md12b\_}) \wedge \\ \quad \operatorname{md12y\_} \in \mathbb{N} \wedge \\ \quad \neg (\operatorname{md12y\_} \in (\operatorname{ctx\text{-}fv}(\operatorname{ctx\text{-}add}(\operatorname{md12cg\_}, \operatorname{md12f\_})) \cup \operatorname{expr\text{-}fv}(\operatorname{enc\text{-}eq}(\operatorname{md12cw\_}, \operatorname{md12cw\_})))) \wedge \\ \quad \operatorname{expr\text{-}open}(\operatorname{md12b\_}, 0, \operatorname{efvar}(\operatorname{md12y\_})) \in \operatorname{md12x\_} \vee \\ \exists \operatorname{md12l\_}, \operatorname{md12k\_}, \operatorname{md12d\_}, \operatorname{md12b2\_}. \\ \quad \operatorname{md12f\_} = \operatorname{eapp}(10, \operatorname{md12l\_}) \wedge \\ \quad \operatorname{nth}(1, \operatorname{md12l\_}) = \operatorname{econst}(\operatorname{md12k\_}) \wedge \\ \quad \operatorname{pred\text{-}code}(\operatorname{md12k\_}) \wedge \\ \quad {[}\operatorname{md12cg\_}, \operatorname{md12d\_}] \in \operatorname{derivable}(\operatorname{md12t\_}) \wedge \\ \quad \operatorname{enc\text{-}iff}(\operatorname{md12f\_}, \operatorname{md12b2\_}) \in \operatorname{inst\text{-}reach}(\operatorname{md12d\_}, \operatorname{md12l\_}) \wedge \\ \quad \operatorname{md12b2\_} \in \operatorname{md12x\_} \vee \\ \exists \operatorname{md12a\_}, \operatorname{md12s\_}, \operatorname{md12d3\_}, \operatorname{md12b3\_}. \\ \quad \operatorname{md12f\_} = \operatorname{enc\text{-}in}(\operatorname{md12s\_}, \operatorname{md12a\_}) \wedge \\ \quad {[}\operatorname{md12cg\_}, \operatorname{md12d3\_}] \in \operatorname{derivable}(\operatorname{md12t\_}) \wedge \\ \quad \operatorname{enc\text{-}iff}(\operatorname{md12f\_}, \operatorname{md12b3\_}) \in \operatorname{inst\text{-}reach}(\operatorname{md12d3\_}, {[}\operatorname{md12a\_}, \operatorname{md12s\_}]) \wedge \\ \quad \operatorname{md12b3\_} \in \operatorname{md12x\_} \end{array}
\end{equation*}

\par\medskip\noindent\textbf{Definition} (\texttt{inst-step}).\;For all $\operatorname{md12l\_}$, $\operatorname{md12f\_}$ and $\operatorname{md12g\_}$, 
\begin{equation*}
\begin{array}{@{}l@{}} \operatorname{inst\text{-}step}(\operatorname{md12l\_}, \operatorname{md12f\_}, \operatorname{md12g\_}) \\ \Leftrightarrow\; \exists \operatorname{md12b\_}, \operatorname{md12i\_}. \\ \quad \operatorname{md12f\_} = \operatorname{enc\text{-}all}(\operatorname{md12b\_}) \wedge \\ \quad \operatorname{md12i\_} \in \mathbb{N} \wedge \\ \quad 2 \leq \operatorname{md12i\_} \wedge \\ \quad \operatorname{md12i\_} \leq \operatorname{length}(\operatorname{md12l\_}) \wedge \\ \quad \operatorname{md12g\_} = \operatorname{expr\text{-}open}(\operatorname{md12b\_}, 0, \operatorname{nth}(\operatorname{md12i\_}, \operatorname{md12l\_})) \end{array}
\end{equation*}

\par\medskip\noindent\textbf{Definition} (\texttt{inst-closed}).\;For all $\operatorname{md12l\_}$ and $\operatorname{md12x\_}$, 
\begin{equation*}
\begin{array}{@{}l@{}} \operatorname{inst\text{-}closed}(\operatorname{md12l\_}, \operatorname{md12x\_}) \\ \Leftrightarrow\; \forall \operatorname{md12p\_} \in \operatorname{md12x\_}, \operatorname{md12q\_} \in \mathbf{Set}. \\ \quad \operatorname{inst\text{-}step}(\operatorname{md12l\_}, \operatorname{md12p\_}, \operatorname{md12q\_}) \Rightarrow (\operatorname{md12q\_} \in \operatorname{md12x\_}) \end{array}
\end{equation*}

\par\medskip\noindent\textbf{Definition} (\texttt{adm-def-pred-unfold-n}).\;For every $\operatorname{mdt\_}$, 
\begin{equation*}
\begin{array}{@{}l@{}} \operatorname{adm\text{-}def\text{-}pred\text{-}unfold\text{-}n}(\operatorname{mdt\_}) \\ \Leftrightarrow\; \forall \operatorname{mdhs\_}, \operatorname{mdc\_}. \\ \quad \operatorname{def\text{-}pred\text{-}unfold\text{-}n}(\operatorname{mdt\_}, \operatorname{mdhs\_}, \operatorname{mdc\_}) \Rightarrow  \\ \quad \operatorname{hyps\text{-}in}(\operatorname{mdhs\_}, \operatorname{derivable}(\operatorname{mdt\_})) \Rightarrow (\operatorname{mdc\_} \in \operatorname{derivable}(\operatorname{mdt\_})) \end{array}
\end{equation*}

\par\medskip\noindent\textbf{Definition} (\texttt{cert-closed}).\;For all $\operatorname{md12t\_}$, $\operatorname{md12cg\_}$, $\operatorname{md12cw\_}$ and $\operatorname{md12x\_}$, 
\begin{equation*}
\begin{array}{@{}l@{}} \operatorname{cert\text{-}closed}(\operatorname{md12t\_}, \operatorname{md12cg\_}, \operatorname{md12cw\_}, \operatorname{md12x\_}) \\ \Leftrightarrow\; \forall \operatorname{md12e\_} \in \operatorname{expr\text{-}set}. \\ \quad \operatorname{cert\text{-}step}(\operatorname{md12t\_}, \operatorname{md12cg\_}, \operatorname{md12cw\_}, \operatorname{md12x\_}, \operatorname{md12e\_}) \Rightarrow (\operatorname{md12e\_} \in \operatorname{md12x\_}) \end{array}
\end{equation*}

\par\medskip\noindent\textbf{Definition} (\texttt{adm-def-pred-hyp-n}).\;For every $\operatorname{mdt\_}$, 
\begin{equation*}
\begin{array}{@{}l@{}} \operatorname{adm\text{-}def\text{-}pred\text{-}hyp\text{-}n}(\operatorname{mdt\_}) \\ \Leftrightarrow\; \forall \operatorname{mdhs\_}, \operatorname{mdc\_}. \\ \quad \operatorname{def\text{-}pred\text{-}hyp\text{-}n}(\operatorname{mdt\_}, \operatorname{mdhs\_}, \operatorname{mdc\_}) \Rightarrow  \\ \quad \operatorname{hyps\text{-}in}(\operatorname{mdhs\_}, \operatorname{derivable}(\operatorname{mdt\_})) \Rightarrow (\operatorname{mdc\_} \in \operatorname{derivable}(\operatorname{mdt\_})) \end{array}
\end{equation*}

\par\medskip\noindent\textbf{Definition} (\texttt{adm-def-cert-hyp}).\;For every $\operatorname{mdt\_}$, 
\begin{equation*}
\begin{array}{@{}l@{}} \operatorname{adm\text{-}def\text{-}cert\text{-}hyp}(\operatorname{mdt\_}) \\ \Leftrightarrow\; \forall \operatorname{mdhs\_}, \operatorname{mdc\_}. \\ \quad \operatorname{def\text{-}cert\text{-}hyp}(\operatorname{mdt\_}, \operatorname{mdhs\_}, \operatorname{mdc\_}) \Rightarrow  \\ \quad \operatorname{hyps\text{-}in}(\operatorname{mdhs\_}, \operatorname{derivable}(\operatorname{mdt\_})) \Rightarrow (\operatorname{mdc\_} \in \operatorname{derivable}(\operatorname{mdt\_})) \end{array}
\end{equation*}

\par\medskip\noindent\textbf{Definition} (\texttt{j-macn}).\;$\operatorname{j\text{-}macn}(\operatorname{md13t\_}) :=$
\begin{equation*}
\{\operatorname{md13x\_} \in \operatorname{rw\text{-}univ} : \exists \operatorname{md13g\_} \in \mathbf{Set}, \operatorname{md13m\_} \in \operatorname{md13t\_}, \operatorname{md13l\_} \in \mathbf{Set}, \operatorname{md13i\_}.; \exists \operatorname{md13r\_}, \operatorname{md13a\_}, \operatorname{md13b\_} \in \mathbf{Set}.\; ((\operatorname{md13m\_} \in \mathit{clf}) \wedge ((\operatorname{md13i\_} \in \operatorname{inst\text{-}reach}(\operatorname{md13m\_}, \operatorname{md13l\_})) \wedge (\forall \operatorname{md13j\_} \in \mathbb{N}.\; ((2 \leq \operatorname{md13j\_}) \Rightarrow ((\operatorname{md13j\_} \leq \operatorname{length}(\operatorname{md13l\_})) \Rightarrow ({[}\operatorname{md13g\_}, \operatorname{enc\text{-}eq}(\operatorname{nth}(\operatorname{md13j\_}, \operatorname{md13l\_}), \operatorname{nth}(\operatorname{md13j\_}, \operatorname{md13l\_}))] \in \operatorname{derivable}(\operatorname{md13t\_})))) \wedge (((\operatorname{md13i\_} = \operatorname{md13r\_}) \vee \exists \operatorname{md13c\_} \in \mathbf{Set}.\; ((\operatorname{md13i\_} = \operatorname{enc\text{-}imp}(\operatorname{md13c\_}, \operatorname{md13r\_})) \wedge ({[}\operatorname{md13g\_}, \operatorname{md13c\_}] \in \operatorname{derivable}(\operatorname{md13t\_})))) \wedge (((\operatorname{md13r\_} = \operatorname{enc\text{-}qeq}(\operatorname{md13a\_}, \operatorname{md13b\_})) \wedge ((\operatorname{md13x\_} = {[}1, \operatorname{md13g\_}, \operatorname{md13a\_}, \operatorname{md13b\_}]) \wedge ((\operatorname{md13a\_} \in \mathit{clt}) \wedge (\operatorname{md13b\_} \in \mathit{clt})))) \vee ((\operatorname{md13r\_} = \operatorname{enc\text{-}iff}(\operatorname{md13a\_}, \operatorname{md13b\_})) \wedge ((\operatorname{md13x\_} = {[}2, \operatorname{md13g\_}, \operatorname{md13a\_}, \operatorname{md13b\_}]) \wedge ((\operatorname{md13a\_} \in \mathit{clf}) \wedge (\operatorname{md13b\_} \in \mathit{clf})))))))))\}
\end{equation*}

\par\medskip\noindent\textbf{Definition} (\texttt{rel-macete-n-1}).\;For all $\operatorname{md13t\_}$, $\operatorname{md13hs\_}$ and $\operatorname{md13c\_}$, 
\begin{equation*}
\operatorname{rel\text{-}macete\text{-}n\text{-}1}(\operatorname{md13t\_}, \operatorname{md13hs\_}, \operatorname{md13c\_}) \Leftrightarrow \operatorname{rel\text{-}rewrite\text{-}1}(\operatorname{j\text{-}macn}(\operatorname{md13t\_}), \operatorname{md13hs\_}, \operatorname{md13c\_})
\end{equation*}

\par\medskip\noindent\textbf{Definition} (\texttt{rel-macete-n-hyp-1}).\;For all $\operatorname{md13t\_}$, $\operatorname{md13hs\_}$ and $\operatorname{md13c\_}$, 
\begin{equation*}
\operatorname{rel\text{-}macete\text{-}n\text{-}hyp\text{-}1}(\operatorname{md13t\_}, \operatorname{md13hs\_}, \operatorname{md13c\_}) \Leftrightarrow \operatorname{rel\text{-}rewrite\text{-}hyp\text{-}1}(\operatorname{j\text{-}macn}(\operatorname{md13t\_}), \operatorname{md13hs\_}, \operatorname{md13c\_})
\end{equation*}

\par\medskip\noindent\textbf{Definition} (\texttt{rel-macete-n-7}).\;For all $\operatorname{md13t\_}$, $\operatorname{md13hs\_}$ and $\operatorname{md13c\_}$, 
\begin{equation*}
\operatorname{rel\text{-}macete\text{-}n\text{-}7}(\operatorname{md13t\_}, \operatorname{md13hs\_}, \operatorname{md13c\_}) \Leftrightarrow \operatorname{rel\text{-}rewrite\text{-}7}(\operatorname{j\text{-}macn}(\operatorname{md13t\_}), \operatorname{md13hs\_}, \operatorname{md13c\_})
\end{equation*}

\par\medskip\noindent\textbf{Definition} (\texttt{rel-macete-n-hyp-7}).\;For all $\operatorname{md13t\_}$, $\operatorname{md13hs\_}$ and $\operatorname{md13c\_}$, 
\begin{equation*}
\operatorname{rel\text{-}macete\text{-}n\text{-}hyp\text{-}7}(\operatorname{md13t\_}, \operatorname{md13hs\_}, \operatorname{md13c\_}) \Leftrightarrow \operatorname{rel\text{-}rewrite\text{-}hyp\text{-}7}(\operatorname{j\text{-}macn}(\operatorname{md13t\_}), \operatorname{md13hs\_}, \operatorname{md13c\_})
\end{equation*}

\par\medskip\noindent\textbf{Definition} (\texttt{rf-set}).\;$\operatorname{rf\text{-}set}(\operatorname{md14h\_}, \operatorname{md14x\_}, \operatorname{md14l\_}) :=$
\begin{equation*}
\{\operatorname{md14p\_} \in (\mathbb{N} \times \operatorname{expr\text{-}set}) : \forall \operatorname{md14w\_} \in \operatorname{power}((\mathbb{N} \times \operatorname{expr\text{-}set})).\; (\operatorname{rf\text{-}closed}(\operatorname{md14h\_}, \operatorname{md14x\_}, \operatorname{md14l\_}, \operatorname{md14w\_}) \Rightarrow (\operatorname{md14p\_} \in \operatorname{md14w\_}))\}
\end{equation*}

\par\medskip\noindent\textbf{Definition} (\texttt{j-udec}).\;$\operatorname{j\text{-}udec}(\operatorname{md14t\_}) :=$
\begin{equation*}
\{\operatorname{md14y\_} \in \operatorname{rw\text{-}univ} : \exists \operatorname{md14g\_}, \operatorname{md14x2\_}, \operatorname{md14l2\_} \in \mathbf{Set}, \operatorname{md14f2\_}.\; ((\operatorname{md14y\_} = {[}2, \operatorname{md14g\_}, \operatorname{enc\text{-}in}(\operatorname{md14x2\_}, \operatorname{enc\text{-}union}(\operatorname{md14l2\_})), \operatorname{md14f2\_}]) \wedge (({[}0, \operatorname{md14f2\_}] \in \operatorname{rf\text{-}set}(4, \operatorname{md14x2\_}, \operatorname{md14l2\_})) \wedge ((\operatorname{enc\text{-}in}(\operatorname{md14x2\_}, \operatorname{enc\text{-}union}(\operatorname{md14l2\_})) \in \mathit{clf}) \wedge ((\operatorname{md14f2\_} \in \mathit{clf}) \wedge \forall \operatorname{md14i\_} \in \mathbb{N}.\; ((1 \leq \operatorname{md14i\_}) \Rightarrow ((\operatorname{md14i\_} \leq \operatorname{length}(\operatorname{md14l2\_})) \Rightarrow ({[}\operatorname{md14g\_}, \operatorname{enc\text{-}eq}(\operatorname{nth}(\operatorname{md14i\_}, \operatorname{md14l2\_}), \operatorname{nth}(\operatorname{md14i\_}, \operatorname{md14l2\_}))] \in \operatorname{derivable}(\operatorname{md14t\_}))))))))\}
\end{equation*}

\par\medskip\noindent\textbf{Definition} (\texttt{j-idec}).\;$\operatorname{j\text{-}idec}(\operatorname{md14t\_}) :=$
\begin{equation*}
\{\operatorname{md14y\_} \in \operatorname{rw\text{-}univ} : \exists \operatorname{md14g\_}, \operatorname{md14x2\_}, \operatorname{md14l2\_} \in \mathbf{Set}, \operatorname{md14f2\_}.\; ((\operatorname{md14y\_} = {[}2, \operatorname{md14g\_}, \operatorname{enc\text{-}in}(\operatorname{md14x2\_}, \operatorname{enc\text{-}inter}(\operatorname{md14l2\_})), \operatorname{md14f2\_}]) \wedge (({[}0, \operatorname{md14f2\_}] \in \operatorname{rf\text{-}set}(3, \operatorname{md14x2\_}, \operatorname{md14l2\_})) \wedge ((\operatorname{enc\text{-}in}(\operatorname{md14x2\_}, \operatorname{enc\text{-}inter}(\operatorname{md14l2\_})) \in \mathit{clf}) \wedge (\operatorname{md14f2\_} \in \mathit{clf}))))\}
\end{equation*}

\par\medskip\noindent\textbf{Definition} (\texttt{rf-closed}).\;For all $\operatorname{md14h\_}$, $\operatorname{md14x\_}$, $\operatorname{md14l\_}$ and $\operatorname{md14s\_}$, 
\begin{equation*}
\begin{array}{@{}l@{}} \operatorname{rf\text{-}closed}(\operatorname{md14h\_}, \operatorname{md14x\_}, \operatorname{md14l\_}, \operatorname{md14s\_}) \\ \Leftrightarrow\; \forall \operatorname{md14j\_} \in \mathbb{N}. \\ \quad (\operatorname{succ}(\operatorname{md14j\_}) = \operatorname{length}(\operatorname{md14l\_})) \Rightarrow  \\ \quad \quad {[}\operatorname{md14j\_}, \operatorname{enc\text{-}in}(\operatorname{md14x\_}, \operatorname{nth}(\operatorname{succ}(\operatorname{md14j\_}), \operatorname{md14l\_}))] \in \operatorname{md14s\_} \wedge \\ \forall \operatorname{md14k\_} \in \mathbb{N}, \operatorname{md14f\_}. \\ \quad ({[}\operatorname{succ}(\operatorname{md14k\_}), \operatorname{md14f\_}] \in \operatorname{md14s\_}) \Rightarrow  \\ \quad \quad {[}\operatorname{md14k\_}, \operatorname{eapp}(\operatorname{md14h\_}, {[}\operatorname{enc\text{-}in}(\operatorname{md14x\_}, \operatorname{nth}(\operatorname{succ}(\operatorname{md14k\_}), \operatorname{md14l\_})), \operatorname{md14f\_}])] \in \operatorname{md14s\_} \end{array}
\end{equation*}

\par\medskip\noindent\textbf{Definition} (\texttt{rel-union-decompose-1}).\;For all $\operatorname{md14t\_}$, $\operatorname{md14hs\_}$ and $\operatorname{md14c\_}$, 
\begin{equation*}
\operatorname{rel\text{-}union\text{-}decompose\text{-}1}(\operatorname{md14t\_}, \operatorname{md14hs\_}, \operatorname{md14c\_}) \Leftrightarrow \operatorname{rel\text{-}rewrite\text{-}1}(\operatorname{j\text{-}udec}(\operatorname{md14t\_}), \operatorname{md14hs\_}, \operatorname{md14c\_})
\end{equation*}

\par\medskip\noindent\textbf{Definition} (\texttt{rel-union-decompose-hyp-1}).\;For all $\operatorname{md14t\_}$, $\operatorname{md14hs\_}$ and $\operatorname{md14c\_}$, 
\begin{equation*}
\begin{array}{@{}l@{}} \operatorname{rel\text{-}union\text{-}decompose\text{-}hyp\text{-}1}(\operatorname{md14t\_}, \operatorname{md14hs\_}, \operatorname{md14c\_}) \\ \Leftrightarrow\; \operatorname{rel\text{-}rewrite\text{-}hyp\text{-}1}(\operatorname{j\text{-}udec}(\operatorname{md14t\_}), \operatorname{md14hs\_}, \operatorname{md14c\_}) \end{array}
\end{equation*}

\par\medskip\noindent\textbf{Definition} (\texttt{rel-union-decompose-7}).\;For all $\operatorname{md14t\_}$, $\operatorname{md14hs\_}$ and $\operatorname{md14c\_}$, 
\begin{equation*}
\operatorname{rel\text{-}union\text{-}decompose\text{-}7}(\operatorname{md14t\_}, \operatorname{md14hs\_}, \operatorname{md14c\_}) \Leftrightarrow \operatorname{rel\text{-}rewrite\text{-}7}(\operatorname{j\text{-}udec}(\operatorname{md14t\_}), \operatorname{md14hs\_}, \operatorname{md14c\_})
\end{equation*}

\par\medskip\noindent\textbf{Definition} (\texttt{rel-union-decompose-hyp-7}).\;For all $\operatorname{md14t\_}$, $\operatorname{md14hs\_}$ and $\operatorname{md14c\_}$, 
\begin{equation*}
\begin{array}{@{}l@{}} \operatorname{rel\text{-}union\text{-}decompose\text{-}hyp\text{-}7}(\operatorname{md14t\_}, \operatorname{md14hs\_}, \operatorname{md14c\_}) \\ \Leftrightarrow\; \operatorname{rel\text{-}rewrite\text{-}hyp\text{-}7}(\operatorname{j\text{-}udec}(\operatorname{md14t\_}), \operatorname{md14hs\_}, \operatorname{md14c\_}) \end{array}
\end{equation*}

\par\medskip\noindent\textbf{Definition} (\texttt{rel-intersection-decompose-1}).\;For all $\operatorname{md14t\_}$, $\operatorname{md14hs\_}$ and $\operatorname{md14c\_}$, 
\begin{equation*}
\begin{array}{@{}l@{}} \operatorname{rel\text{-}intersection\text{-}decompose\text{-}1}(\operatorname{md14t\_}, \operatorname{md14hs\_}, \operatorname{md14c\_}) \\ \Leftrightarrow\; \operatorname{rel\text{-}rewrite\text{-}1}(\operatorname{j\text{-}idec}(\operatorname{md14t\_}), \operatorname{md14hs\_}, \operatorname{md14c\_}) \end{array}
\end{equation*}

\par\medskip\noindent\textbf{Definition} (\texttt{rel-intersection-decompose-hyp-1}).\;For all $\operatorname{md14t\_}$, $\operatorname{md14hs\_}$ and $\operatorname{md14c\_}$, 
\begin{equation*}
\begin{array}{@{}l@{}} \operatorname{rel\text{-}intersection\text{-}decompose\text{-}hyp\text{-}1}(\operatorname{md14t\_}, \operatorname{md14hs\_}, \operatorname{md14c\_}) \\ \Leftrightarrow\; \operatorname{rel\text{-}rewrite\text{-}hyp\text{-}1}(\operatorname{j\text{-}idec}(\operatorname{md14t\_}), \operatorname{md14hs\_}, \operatorname{md14c\_}) \end{array}
\end{equation*}

\par\medskip\noindent\textbf{Definition} (\texttt{rel-intersection-decompose-7}).\;For all $\operatorname{md14t\_}$, $\operatorname{md14hs\_}$ and $\operatorname{md14c\_}$, 
\begin{equation*}
\begin{array}{@{}l@{}} \operatorname{rel\text{-}intersection\text{-}decompose\text{-}7}(\operatorname{md14t\_}, \operatorname{md14hs\_}, \operatorname{md14c\_}) \\ \Leftrightarrow\; \operatorname{rel\text{-}rewrite\text{-}7}(\operatorname{j\text{-}idec}(\operatorname{md14t\_}), \operatorname{md14hs\_}, \operatorname{md14c\_}) \end{array}
\end{equation*}

\par\medskip\noindent\textbf{Definition} (\texttt{rel-intersection-decompose-hyp-7}).\;For all $\operatorname{md14t\_}$, $\operatorname{md14hs\_}$ and $\operatorname{md14c\_}$, 
\begin{equation*}
\begin{array}{@{}l@{}} \operatorname{rel\text{-}intersection\text{-}decompose\text{-}hyp\text{-}7}(\operatorname{md14t\_}, \operatorname{md14hs\_}, \operatorname{md14c\_}) \\ \Leftrightarrow\; \operatorname{rel\text{-}rewrite\text{-}hyp\text{-}7}(\operatorname{j\text{-}idec}(\operatorname{md14t\_}), \operatorname{md14hs\_}, \operatorname{md14c\_}) \end{array}
\end{equation*}


\vnbbookfile{theorem-library/meta-ind-laws}{theorem-library/meta-ind-laws}
\noindent{\small meta-ind-laws.scm -- K2 PHASE 4: the admissibility of the kernel's set-formation rules and of the three induction rules over the core, the typing machinery an instance at a compound term needs, and (section J) the one phase-5 derivation that needs the OPEN calculator (docs/meta-k2-admissibility-2026-10-08.md, section ...}\par
\vnbprop{hyps-in-2-elim}{\vnbstmt{\begin{array}{@{}l@{}} \forall a, b, s. \\ \quad (({[}a, b] = {[}a, b]) \wedge \operatorname{hyps\text{-}in}({[}a, b], s)) \Rightarrow ((a \in s) \wedge (b \in s)) \end{array}}}
\vnbprop{adm-sep-sethood}{\vnbstmt{\begin{array}{@{}l@{}} \forall t, s, c. \\ \quad (\operatorname{rel\text{-}sep\text{-}sethood}(t, s, c) \wedge \operatorname{hyps\text{-}in}(s, \operatorname{derivable}(t))) \Rightarrow (c \in \operatorname{derivable}(t)) \end{array}}}
\vnbprop{adm-big-union-sethood}{\vnbstmt{\begin{array}{@{}l@{}} \forall t, s, c. \\ \quad (\operatorname{rel\text{-}big\text{-}union\text{-}sethood}(t, s, c) \wedge \operatorname{hyps\text{-}in}(s, \operatorname{derivable}(t))) \Rightarrow  \\ \quad \quad c \in \operatorname{derivable}(t) \end{array}}}
\vnbprop{adm-sep-mem-intro}{\vnbstmt{\begin{array}{@{}l@{}} \forall t, s, c. \\ \quad (\operatorname{rel\text{-}sep\text{-}mem\text{-}intro}(t, s, c) \wedge \operatorname{hyps\text{-}in}(s, \operatorname{derivable}(t))) \Rightarrow  \\ \quad \quad c \in \operatorname{derivable}(t) \end{array}}}
\vnbprop{adm-comp-mem-intro}{\vnbstmt{\begin{array}{@{}l@{}} \forall t, s, c. \\ \quad (\operatorname{rel\text{-}comp\text{-}mem\text{-}intro}(t, s, c) \wedge \operatorname{hyps\text{-}in}(s, \operatorname{derivable}(t))) \Rightarrow  \\ \quad \quad c \in \operatorname{derivable}(t) \end{array}}}
\vnbprop{enc-sep-unfold}{\vnbstmt{\forall a, p.\; (\operatorname{enc\text{-}sep}(a, p) \simeq \operatorname{ebind}(22, 1, {[}a], p))}}
\vnbprop{enc-comp-unfold}{\vnbstmt{\forall p.\; (\operatorname{enc\text{-}comp}(p) \simeq \operatorname{ebind}(23, 1, {[}], p))}}
\vnbprop{enc-bu-unfold}{\vnbstmt{\forall a, p.\; (\operatorname{enc\text{-}bu}(a, p) \simeq \operatorname{ebind}(24, 1, {[}a], p))}}
\vnbprop{adm-sep-mem-elim}{\vnbstmt{\begin{array}{@{}l@{}} \forall t, s, c. \\ \quad (\operatorname{rel\text{-}sep\text{-}mem\text{-}elim}(t, s, c) \wedge \operatorname{hyps\text{-}in}(s, \operatorname{derivable}(t))) \Rightarrow (c \in \operatorname{derivable}(t)) \end{array}}}
\vnbprop{adm-comp-mem-elim}{\vnbstmt{\begin{array}{@{}l@{}} \forall t, s, c. \\ \quad (\operatorname{rel\text{-}comp\text{-}mem\text{-}elim}(t, s, c) \wedge \operatorname{hyps\text{-}in}(s, \operatorname{derivable}(t))) \Rightarrow  \\ \quad \quad c \in \operatorname{derivable}(t) \end{array}}}
\vnbprop{wks-single}{\vnbstmt{\forall e \in \operatorname{expr\text{-}set}.\; ((1 \in \operatorname{expr\text{-}wks}(e)) \Rightarrow (\operatorname{expr\text{-}wks}(e) = \{1\}))}}
\vnbprop{lcd-mono}{\vnbstmt{\begin{array}{@{}l@{}} \forall e \in \operatorname{expr\text{-}set}, d, a. \\ \quad ((d \in \mathbb{N}) \wedge ((a \in \mathbb{N}) \wedge ((d \leq a) \wedge (d \in \operatorname{expr\text{-}lcd}(e))))) \Rightarrow (a \in \operatorname{expr\text{-}lcd}(e)) \end{array}}}
\vnbprop{lcd-open-closes}{\vnbstmt{\begin{array}{@{}l@{}} \forall u, e \in \operatorname{expr\text{-}set}, x \in \mathbb{N}. \\ \quad (((x + 1) \in \operatorname{expr\text{-}lcd}(e)) \wedge (x \in \operatorname{expr\text{-}lcd}(u))) \Rightarrow  \\ \quad \quad x \in \operatorname{expr\text{-}lcd}(\operatorname{expr\text{-}open}(e, x, u)) \end{array}}}
\vnbprop{enc-succ-unfold}{\vnbstmt{\forall p.\; (\operatorname{enc\text{-}succ}(p) \simeq \operatorname{eapp}(30, {[}p]))}}
\vnbprop{enc-succo-unfold}{\vnbstmt{\forall p.\; (\operatorname{enc\text{-}succo}(p) \simeq \operatorname{eapp}(31, {[}p]))}}
\vnbprop{enc-ordlt-unfold}{\vnbstmt{\forall p, q.\; (\operatorname{enc\text{-}ordlt}(p, q) \simeq \operatorname{eapp}(10, {[}\operatorname{econst}(7), p, q]))}}
\vnbprop{enc-limit-unfold}{\vnbstmt{\forall p.\; (\operatorname{enc\text{-}limit}(p) \simeq \operatorname{eapp}(10, {[}\operatorname{econst}(9), p]))}}
\vnbprop{enc-num-unfold}{\vnbstmt{\forall k \in \mathbb{N}.\; (\operatorname{enc\text{-}num}(k) = \operatorname{econst}((2 \cdot k)))}}
\vnbprop{ebind-args-stage}{\vnbstmt{\begin{array}{@{}l@{}} \forall n \in \mathbb{N}, h, c, a, b. \\ \quad (\operatorname{ebind}(h, c, a, b) \in \operatorname{expr\text{-}stage}(n)) \Rightarrow (a \in \operatorname{tuples}(\operatorname{expr\text{-}set})) \end{array}}}
\vnbprop{ebind-args-expr}{\vnbstmt{\begin{array}{@{}l@{}} \forall h, c, a, b. \\ \quad (\operatorname{ebind}(h, c, a, b) \in \operatorname{expr\text{-}set}) \Rightarrow (a \in \operatorname{tuples}(\operatorname{expr\text{-}set})) \end{array}}}
\vnbprop{comp-of-body-clt}{\vnbstmt{\forall q, p.\; ((\operatorname{enc\text{-}all}(\operatorname{enc\text{-}imp}(q, p)) \in \mathit{clf}) \Rightarrow (\operatorname{enc\text{-}comp}(p) \in \mathit{clt}))}}
\vnbprop{open-app1}{\vnbstmt{\begin{array}{@{}l@{}} \forall u, k, h, a_{1}. \\ \quad ((u \in \operatorname{expr\text{-}set}) \wedge ((k \in \mathbb{N}) \wedge ((h \in \mathbb{N}) \wedge (a_{1} \in \operatorname{expr\text{-}set})))) \Rightarrow  \\ \quad \quad \operatorname{expr\text{-}open}(\operatorname{eapp}(h, {[}a_{1}]), k, u) = \operatorname{eapp}(h, {[}\operatorname{expr\text{-}open}(a_{1}, k, u)]) \end{array}}}
\vnbprop{open-app3}{\vnbstmt{\begin{array}{@{}l@{}} \forall u, k, h, a_{1}, a_{2}, a_{3}. \\ \quad ((u \in \operatorname{expr\text{-}set}) \wedge ((k \in \mathbb{N}) \wedge ((h \in \mathbb{N}) \wedge ((a_{1} \in \operatorname{expr\text{-}set}) \wedge ((a_{2} \in \operatorname{expr\text{-}set}) \wedge (a_{3} \in \operatorname{expr\text{-}set})))))) \Rightarrow  \\ \quad \quad \operatorname{expr\text{-}open}(\operatorname{eapp}(h, {[}a_{1}, a_{2}, a_{3}]), k, u) \\ \quad \quad =\; \operatorname{eapp}(h, {[}\operatorname{expr\text{-}open}(a_{1}, k, u), \operatorname{expr\text{-}open}(a_{2}, k, u), \operatorname{expr\text{-}open}(a_{3}, k, u)]) \end{array}}}
\vnbprop{maplk-nil}{\vnbstmt{\forall u, k.\; (((u \in \operatorname{expr\text{-}set}) \wedge (k \in \mathbb{N})) \Rightarrow (\operatorname{maplk}(\operatorname{open\text{-}fun}(u), {[}], k) = {[}]))}}
\vnbprop{der-comp-out}{\vnbstmt{\begin{array}{@{}l@{}} \forall a, h, q, t. \\ \quad (\operatorname{nbg\text{-}ok}(q) \wedge ({[}h, \operatorname{enc\text{-}in}(t, \operatorname{enc\text{-}comp}(q))] \in \operatorname{derivable}(a))) \Rightarrow  \\ \quad \quad {[}h, \operatorname{expr\text{-}open}(q, 0, t)] \in \operatorname{derivable}(a) \end{array}}}
\vnbprop{der-comp-in}{\vnbstmt{\begin{array}{@{}l@{}} \forall a, h, q, t. \\ \quad \operatorname{enc\text{-}comp}(q) \in \mathit{clt} \wedge \\ \quad \operatorname{nbg\text{-}ok}(q) \wedge \\ \quad {[}h, \operatorname{enc\text{-}in}(t, \operatorname{enc\text{-}setc})] \in \operatorname{derivable}(a) \wedge \\ \quad {[}h, \operatorname{expr\text{-}open}(q, 0, t)] \in \operatorname{derivable}(a) \Rightarrow \\ \quad \quad {[}h, \operatorname{enc\text{-}in}(t, \operatorname{enc\text{-}comp}(q))] \in \operatorname{derivable}(a) \end{array}}}
\vnbprop{der-name-set}{\vnbstmt{\begin{array}{@{}l@{}} \forall t, h, n, m. \\ \quad ((\operatorname{ax\text{-}mem\text{-}set} \in t) \wedge ((n \in \mathbb{N}) \wedge ((m \in \mathbb{N}) \wedge ({[}h, \operatorname{enc\text{-}in}(\operatorname{efvar}(n), \operatorname{econst}(m))] \in \operatorname{derivable}(t))))) \Rightarrow  \\ \quad \quad {[}h, \operatorname{enc\text{-}in}(\operatorname{efvar}(n), \operatorname{enc\text{-}setc})] \in \operatorname{derivable}(t) \end{array}}}
\vnbprop{adm-big-union-mem-intro}{\vnbstmt{\begin{array}{@{}l@{}} \forall t, s, c. \\ \quad (\operatorname{def\text{-}in\text{-}closed}(t) \wedge (\operatorname{rel\text{-}big\text{-}union\text{-}mem\text{-}intro}(t, s, c) \wedge \operatorname{hyps\text{-}in}(s, \operatorname{derivable}(t)))) \Rightarrow  \\ \quad \quad c \in \operatorname{derivable}(t) \end{array}}}
\vnbprop{ctx-fv-member}{\vnbstmt{\forall g, h, y.\; (((h \in g) \wedge (y \in \operatorname{expr\text{-}fv}(h))) \Rightarrow (y \in \operatorname{ctx\text{-}fv}(g)))}}
\vnbprop{ctx-fv-add}{\vnbstmt{\begin{array}{@{}l@{}} \forall g, h, y. \\ \quad ((h \in \mathbf{Set}) \wedge (y \in \operatorname{ctx\text{-}fv}(\operatorname{ctx\text{-}add}(g, h)))) \Rightarrow  \\ \quad \quad (y \in \operatorname{ctx\text{-}fv}(g)) \vee (y \in \operatorname{expr\text{-}fv}(h)) \end{array}}}
\vnbprop{fv-bind1-dom}{\vnbstmt{\begin{array}{@{}l@{}} \forall h, c, a, b, z. \\ \quad ((h \in \mathbb{N}) \wedge ((c \in \mathbb{N}) \wedge ((a \in \operatorname{expr\text{-}set}) \wedge ((b \in \operatorname{expr\text{-}set}) \wedge (z \in \operatorname{expr\text{-}fv}(a)))))) \Rightarrow  \\ \quad \quad z \in \operatorname{expr\text{-}fv}(\operatorname{ebind}(h, c, {[}a], b)) \end{array}}}
\vnbprop{fv-ex-bu}{\vnbstmt{\begin{array}{@{}l@{}} \forall a, b, x, z. \\ \quad a \in \operatorname{expr\text{-}set} \wedge \\ \quad b \in \operatorname{expr\text{-}set} \wedge \\ \quad x \in \operatorname{expr\text{-}set} \wedge \\ \quad z \in \operatorname{expr\text{-}fv}(\operatorname{enc\text{-}ex}(\operatorname{enc\text{-}and}(\operatorname{enc\text{-}in}(\operatorname{ebvar}(0), a), \operatorname{enc\text{-}in}(x, b)))) \Rightarrow \\ \quad \quad z \in \operatorname{expr\text{-}fv}(\operatorname{enc\text{-}in}(x, \operatorname{enc\text{-}bu}(a, b))) \end{array}}}
\vnbprop{adm-big-union-mem-elim}{\vnbstmt{\begin{array}{@{}l@{}} \forall t, s, c. \\ \quad (\operatorname{rel\text{-}big\text{-}union\text{-}mem\text{-}elim}(t, s, c) \wedge \operatorname{hyps\text{-}in}(s, \operatorname{derivable}(t))) \Rightarrow  \\ \quad \quad c \in \operatorname{derivable}(t) \end{array}}}
\vnbprop{clf-open-inst}{\vnbstmt{\forall b, t.\; (((\operatorname{enc\text{-}all}(b) \in \mathit{clf}) \wedge (t \in \mathit{clt})) \Rightarrow (\operatorname{expr\text{-}open}(b, 0, t) \in \mathit{clf}))}}
\vnbprop{clf-open-inst-ex}{\vnbstmt{\forall b, t.\; (((\operatorname{enc\text{-}ex}(b) \in \mathit{clf}) \wedge (t \in \mathit{clt})) \Rightarrow (\operatorname{expr\text{-}open}(b, 0, t) \in \mathit{clf}))}}
\vnbprop{adm-transfinite-induction}{\vnbstmt{\begin{array}{@{}l@{}} \forall t, s, c. \\ \quad t \subseteq \mathit{clf} \wedge \\ \quad \operatorname{ax\text{-}tfi} \in t \wedge \\ \quad \operatorname{ax\text{-}mem\text{-}set} \in t \wedge \\ \quad \operatorname{def\text{-}comp\text{-}closed}(t) \wedge \\ \quad \operatorname{rel\text{-}transfinite\text{-}induction\text{-}nbg}(t, s, c) \wedge \\ \quad \operatorname{hyps\text{-}in}(s, \operatorname{derivable}(t)) \Rightarrow \\ \quad \quad c \in \operatorname{derivable}(t) \end{array}}}
\vnbprop{der-mem-set}{\vnbstmt{\begin{array}{@{}l@{}} \forall t, h, x, m. \\ \quad t \subseteq \mathit{clf} \wedge \\ \quad \operatorname{ax\text{-}mem\text{-}set} \in t \wedge \\ \quad m \in \mathbb{N} \wedge \\ \quad {[}h, \operatorname{enc\text{-}in}(x, \operatorname{econst}(m))] \in \operatorname{derivable}(t) \wedge \\ \quad {[}h, \operatorname{enc\text{-}eq}(x, x)] \in \operatorname{derivable}(t) \Rightarrow \\ \quad \quad {[}h, \operatorname{enc\text{-}in}(x, \operatorname{enc\text{-}setc})] \in \operatorname{derivable}(t) \end{array}}}
\vnbprop{adm-nn-induction}{\vnbstmt{\begin{array}{@{}l@{}} \forall t, s, c. \\ \quad t \subseteq \mathit{clf} \wedge \\ \quad \operatorname{ax\text{-}nn\text{-}ind} \in t \wedge \\ \quad \operatorname{ax\text{-}mem\text{-}set} \in t \wedge \\ \quad \operatorname{ax\text{-}zero\text{-}nn} \in t \wedge \\ \quad \operatorname{ax\text{-}succ\text{-}nn} \in t \wedge \\ \quad \operatorname{def\text{-}comp\text{-}closed}(t) \wedge \\ \quad \operatorname{def\text{-}in\text{-}closed}(t) \wedge \\ \quad \operatorname{rel\text{-}nn\text{-}induction\text{-}nbg}(t, s, c) \wedge \\ \quad \operatorname{hyps\text{-}in}(s, \operatorname{derivable}(t)) \Rightarrow \\ \quad \quad c \in \operatorname{derivable}(t) \end{array}}}
\vnbprop{der-em}{\vnbstmt{\begin{array}{@{}l@{}} \forall t, h, x. \\ \quad ((h \in \operatorname{ctx\text{-}set}) \wedge (x \in \mathit{clf})) \Rightarrow ({[}h, \operatorname{enc\text{-}or}(x, \operatorname{enc\text{-}not}(x))] \in \operatorname{derivable}(t)) \end{array}}}
\vnbprop{subst-open-name}{\vnbstmt{\begin{array}{@{}l@{}} \forall y, e, v. \\ \quad ((y \in \mathbb{N}) \wedge ((e \in \operatorname{expr\text{-}set}) \wedge ((v \in \operatorname{expr\text{-}set}) \wedge (\operatorname{expr\text{-}lc}(v) \wedge \neg (y \in \operatorname{expr\text{-}fv}(e)))))) \Rightarrow  \\ \quad \quad \operatorname{expr\text{-}subst}(\operatorname{expr\text{-}open}(e, 0, \operatorname{efvar}(y)), y, v) = \operatorname{expr\text{-}open}(e, 0, v) \end{array}}}
\vnbprop{hyps-in-3-elim}{\vnbstmt{\begin{array}{@{}l@{}} \forall a, b, c, s. \\ \quad (({[}a, b, c] = {[}a, b, c]) \wedge \operatorname{hyps\text{-}in}({[}a, b, c], s)) \Rightarrow  \\ \quad \quad (a \in s) \wedge ((b \in s) \wedge (c \in s)) \end{array}}}
\vnbprop{adm-tfi-3cases}{\vnbstmt{\begin{array}{@{}l@{}} \forall t, s, c. \\ \quad t \subseteq \mathit{clf} \wedge \\ \quad \operatorname{ax\text{-}tfi} \in t \wedge \\ \quad \operatorname{ax\text{-}mem\text{-}set} \in t \wedge \\ \quad \operatorname{ax\text{-}limit\text{-}iff} \in t \wedge \\ \quad \operatorname{ax\text{-}ord\text{-}succ\text{-}above} \in t \wedge \\ \quad \operatorname{def\text{-}comp\text{-}closed}(t) \wedge \\ \quad \operatorname{rel\text{-}tfi\text{-}3cases\text{-}nbg}(t, s, c) \wedge \\ \quad \operatorname{hyps\text{-}in}(s, \operatorname{derivable}(t)) \Rightarrow \\ \quad \quad c \in \operatorname{derivable}(t) \end{array}}}
\vnbprop{enc-apply-unfold}{\vnbstmt{\forall p, q.\; (\operatorname{enc\text{-}apply}(p, q) \simeq \operatorname{eapp}(18, {[}p, q]))}}
\vnbprop{enc-fun1-unfold}{\vnbstmt{\forall p.\; (\operatorname{enc\text{-}fun1}(p) \simeq \operatorname{eapp}(19, {[}p]))}}
\vnbprop{adm-def-app}{\vnbstmt{\begin{array}{@{}l@{}} \forall t, s, c. \\ \quad t \subseteq \mathit{clf} \wedge \\ \quad \operatorname{ax\text{-}fun\text{-}app} \in t \wedge \\ \quad \operatorname{adm\text{-}def\text{-}strict\text{-}in\text{-}l}(t) \wedge \\ \quad \operatorname{adm\text{-}def\text{-}strict\text{-}in\text{-}r}(t) \wedge \\ \quad \operatorname{def\text{-}app}(t, s, c) \wedge \\ \quad \operatorname{hyps\text{-}in}(s, \operatorname{derivable}(t)) \Rightarrow \\ \quad \quad c \in \operatorname{derivable}(t) \end{array}}}
\vnbprop{enc-pred1-unfold}{\vnbstmt{\forall p, q.\; (\operatorname{enc\text{-}pred1}(p, q) \simeq \operatorname{eapp}(10, {[}p, q]))}}
\vnbprop{enc-le-unfold}{\vnbstmt{\forall p, q.\; (\operatorname{enc\text{-}le}(p, q) \simeq \operatorname{eapp}(10, {[}\operatorname{econst}(11), p, q]))}}
\vnbprop{enc-lt-unfold}{\vnbstmt{\forall p, q.\; (\operatorname{enc\text{-}lt}(p, q) \simeq \operatorname{eapp}(10, {[}\operatorname{econst}(13), p, q]))}}
\vnbprop{pred-code-nn}{\vnbstmt{\forall k.\; (\operatorname{pred\text{-}code}(k) \Rightarrow (k \in \mathbb{N}))}}
\vnbprop{clf-pred-arg}{\vnbstmt{\begin{array}{@{}l@{}} \forall l, i. \\ \quad ((\operatorname{eapp}(10, l) \in \mathit{clf}) \wedge ((i \in \mathbb{N}) \wedge ((1 \leq i) \wedge (i \leq \operatorname{length}(l))))) \Rightarrow  \\ \quad \quad \operatorname{nth}(i, l) \in \mathit{clt} \end{array}}}
\vnbprop{adm-pred1-arg}{\vnbstmt{\begin{array}{@{}l@{}} \forall t, g, k, s. \\ \quad (\operatorname{pred\text{-}code}(k) \wedge ({[}g, \operatorname{enc\text{-}pred1}(\operatorname{econst}(k), s)] \in \operatorname{derivable}(t))) \Rightarrow  \\ \quad \quad {[}g, \operatorname{enc\text{-}eq}(s, s)] \in \operatorname{derivable}(t) \end{array}}}
\vnbprop{adm-pred1-unfold}{\vnbstmt{\begin{array}{@{}l@{}} \forall t, g, k, s, b. \\ \quad \operatorname{pred\text{-}code}(k) \wedge \\ \quad {[}g, \operatorname{enc\text{-}pred1}(\operatorname{econst}(k), s)] \in \operatorname{derivable}(t) \wedge \\ \quad {[}g, \operatorname{enc\text{-}all}(\operatorname{enc\text{-}iff}(\operatorname{enc\text{-}pred1}(\operatorname{econst}(k), \operatorname{ebvar}(0)), b))] \in \operatorname{derivable}(t) \Rightarrow \\ \quad \quad {[}g, \operatorname{expr\text{-}open}(b, 0, s)] \in \operatorname{derivable}(t) \end{array}}}
\vnbprop{def-pred-unfold-derived}{\vnbstmt{\begin{array}{@{}l@{}} \forall t, s, c. \\ \quad (\operatorname{def\text{-}pred\text{-}unfold}(t, s, c) \wedge \operatorname{hyps\text{-}in}(s, \operatorname{derivable}(t))) \Rightarrow (c \in \operatorname{derivable}(t)) \end{array}}}
\vnbprop{der-comp-cr}{\vnbstmt{\begin{array}{@{}l@{}} \forall t, g, a, b, y. \\ \quad t \subseteq \mathit{clf} \wedge \\ \quad \operatorname{ax\text{-}class\text{-}ext} \in t \wedge \\ \quad g \in \operatorname{ctx\text{-}set} \wedge \\ \quad \operatorname{enc\text{-}comp}(a) \in \mathit{clt} \wedge \\ \quad \operatorname{enc\text{-}comp}(b) \in \mathit{clt} \wedge \\ \quad \operatorname{nbg\text{-}ok}(a) \wedge \\ \quad \operatorname{nbg\text{-}ok}(b) \wedge \\ \quad y \in \mathbb{N} \wedge \\ \quad \neg (y \in (\operatorname{ctx\text{-}fv}(g) \cup \operatorname{expr\text{-}fv}(\operatorname{enc\text{-}iff}(\operatorname{enc\text{-}in}(\operatorname{ebvar}(0), \operatorname{enc\text{-}comp}(a)), \operatorname{enc\text{-}in}(\operatorname{ebvar}(0), \operatorname{enc\text{-}comp}(b)))))) \wedge \\ \quad {[}g, \operatorname{enc\text{-}iff}(\operatorname{expr\text{-}open}(a, 0, \operatorname{efvar}(y)), \operatorname{expr\text{-}open}(b, 0, \operatorname{efvar}(y)))] \in \operatorname{derivable}(t) \Rightarrow \\ \quad \quad {[}g, \operatorname{enc\text{-}qeq}(\operatorname{enc\text{-}comp}(a), \operatorname{enc\text{-}comp}(b))] \in \operatorname{derivable}(t) \end{array}}}
\vnbprop{eqv-desc-body}{\vnbstmt{\begin{array}{@{}l@{}} \forall j, t, x. \\ \quad (\operatorname{rw\text{-}desc\text{-}body}(j, \operatorname{eqv}(t), x) \wedge (x \in \operatorname{rw\text{-}univ})) \Rightarrow (x \in \operatorname{eqv}(t)) \end{array}}}
\vnbprop{cr-closed-7}{\vnbstmt{\forall t, j.\; ((j \subseteq \operatorname{eqv}(t)) \Rightarrow \operatorname{rw\text{-}closed}(j, 7, \operatorname{eqv}(t)))}}
\vnbprop{cr-7}{\vnbstmt{\forall t, j, x.\; (((j \subseteq \operatorname{eqv}(t)) \wedge (x \in \operatorname{rewr}(j, 7))) \Rightarrow (x \in \operatorname{eqv}(t)))}}
\vnbprop{cr-goal-7}{\vnbstmt{\begin{array}{@{}l@{}} \forall t, j, g, o, o_{2}. \\ \quad ((j \subseteq \operatorname{eqv}(t)) \wedge (({[}2, g, o, o_{2}] \in \operatorname{rewr}(j, 7)) \wedge ({[}g, o_{2}] \in \operatorname{derivable}(t)))) \Rightarrow  \\ \quad \quad {[}g, o] \in \operatorname{derivable}(t) \end{array}}}
\vnbprop{adm-rewrite-7}{\vnbstmt{\begin{array}{@{}l@{}} \forall t, j, s, c. \\ \quad ((j \subseteq \operatorname{eqv}(t)) \wedge (\operatorname{rel\text{-}rewrite\text{-}7}(j, s, c) \wedge \operatorname{hyps\text{-}in}(s, \operatorname{derivable}(t)))) \Rightarrow  \\ \quad \quad c \in \operatorname{derivable}(t) \end{array}}}
\vnbprop{adm-rewrite-hyp-7}{\vnbstmt{\begin{array}{@{}l@{}} \forall t, j, s, c. \\ \quad ((j \subseteq \operatorname{eqv}(t)) \wedge (\operatorname{rel\text{-}rewrite\text{-}hyp\text{-}7}(j, s, c) \wedge \operatorname{hyps\text{-}in}(s, \operatorname{derivable}(t)))) \Rightarrow  \\ \quad \quad c \in \operatorname{derivable}(t) \end{array}}}
\vnbprop{adm-nth-reduce-7}{\vnbstmt{\begin{array}{@{}l@{}} \forall t, s, c. \\ \quad (\operatorname{rel\text{-}nth\text{-}reduce\text{-}7}(t, s, c) \wedge \operatorname{hyps\text{-}in}(s, \operatorname{derivable}(t))) \Rightarrow (c \in \operatorname{derivable}(t)) \end{array}}}
\vnbprop{adm-length-reduce-7}{\vnbstmt{\begin{array}{@{}l@{}} \forall t, s, c. \\ \quad (\operatorname{rel\text{-}length\text{-}reduce\text{-}7}(t, s, c) \wedge \operatorname{hyps\text{-}in}(s, \operatorname{derivable}(t))) \Rightarrow  \\ \quad \quad c \in \operatorname{derivable}(t) \end{array}}}
\vnbprop{adm-lambda-beta-7}{\vnbstmt{\begin{array}{@{}l@{}} \forall t, s, c. \\ \quad (\operatorname{rel\text{-}lambda\text{-}beta\text{-}7}(t, s, c) \wedge \operatorname{hyps\text{-}in}(s, \operatorname{derivable}(t))) \Rightarrow  \\ \quad \quad c \in \operatorname{derivable}(t) \end{array}}}
\vnbprop{adm-lambda-beta-hyp-7}{\vnbstmt{\begin{array}{@{}l@{}} \forall t, s, c. \\ \quad (\operatorname{rel\text{-}lambda\text{-}beta\text{-}hyp\text{-}7}(t, s, c) \wedge \operatorname{hyps\text{-}in}(s, \operatorname{derivable}(t))) \Rightarrow  \\ \quad \quad c \in \operatorname{derivable}(t) \end{array}}}
\vnbprop{adm-macete-7}{\vnbstmt{\begin{array}{@{}l@{}} \forall t, s, c. \\ \quad (\operatorname{rel\text{-}macete\text{-}7}(t, s, c) \wedge \operatorname{hyps\text{-}in}(s, \operatorname{derivable}(t))) \Rightarrow (c \in \operatorname{derivable}(t)) \end{array}}}
\vnbprop{adm-macete-hyp-7}{\vnbstmt{\begin{array}{@{}l@{}} \forall t, s, c. \\ \quad (\operatorname{rel\text{-}macete\text{-}hyp\text{-}7}(t, s, c) \wedge \operatorname{hyps\text{-}in}(s, \operatorname{derivable}(t))) \Rightarrow (c \in \operatorname{derivable}(t)) \end{array}}}
\vnbprop{adm-functoid-beta-7}{\vnbstmt{\begin{array}{@{}l@{}} \forall t, s, c. \\ \quad (\operatorname{rel\text{-}functoid\text{-}beta\text{-}7}(t, s, c) \wedge \operatorname{hyps\text{-}in}(s, \operatorname{derivable}(t))) \Rightarrow  \\ \quad \quad c \in \operatorname{derivable}(t) \end{array}}}
\vnbprop{enc-nth-unfold}{\vnbstmt{\forall p_{7}, q_{7}.\; (\operatorname{enc\text{-}nth}(p_{7}, q_{7}) \simeq \operatorname{eapp}(15, {[}p_{7}, q_{7}]))}}
\vnbprop{enc-length-unfold}{\vnbstmt{\forall p.\; (\operatorname{enc\text{-}length}(p) \simeq \operatorname{eapp}(16, {[}p]))}}
\vnbprop{enc-cartl-clf}{\vnbstmt{\begin{array}{@{}l@{}} \forall n \in \mathbb{N}, x, l. \\ \quad ((x \in \mathit{clt}) \wedge ((l \in \operatorname{tuples}(\operatorname{expr\text{-}set})) \wedge (\forall u \in \operatorname{make\text{-}set}(l).\; (u \in \mathit{clt}) \wedge (n \leq \operatorname{length}(l))))) \Rightarrow  \\ \quad \quad \operatorname{enc\text{-}cartl}(x, l, n) \in \mathit{clf} \end{array}}}
\vnbprop{enc-cartl-elim-der}{\vnbstmt{\begin{array}{@{}l@{}} \forall n \in \mathbb{N}, t, g, x, l, k. \\ \quad {[}g, \operatorname{enc\text{-}cartl}(x, l, n)] \in \operatorname{derivable}(t) \wedge \\ \quad g \in \operatorname{ctx\text{-}set} \wedge \\ \quad x \in \mathit{clt} \wedge \\ \quad l \in \operatorname{tuples}(\operatorname{expr\text{-}set}) \wedge \\ \quad \forall u \in \operatorname{make\text{-}set}(l).\; (u \in \mathit{clt}) \wedge \\ \quad k \in \mathbb{N} \wedge \\ \quad 1 \leq k \wedge \\ \quad k \leq n \wedge \\ \quad n \leq \operatorname{length}(l) \Rightarrow \\ \quad \quad {[}g, \operatorname{enc\text{-}in}(\operatorname{enc\text{-}nth}(\operatorname{enc\text{-}num}(k), x), \operatorname{nth}(k, l))] \in \operatorname{derivable}(t) \end{array}}}
\vnbprop{enc-cartl-der}{\vnbstmt{\begin{array}{@{}l@{}} \forall n \in \mathbb{N}, t, g, x, l. \\ \quad g \in \operatorname{ctx\text{-}set} \wedge \\ \quad x \in \mathit{clt} \wedge \\ \quad l \in \operatorname{tuples}(\operatorname{expr\text{-}set}) \wedge \\ \quad \forall u \in \operatorname{make\text{-}set}(l).\; (u \in \mathit{clt}) \wedge \\ \quad \forall k \in \mathbb{N}. \\ \quad \quad ((1 \leq k) \wedge (k \leq n)) \Rightarrow  \\ \quad \quad \quad {[}g, \operatorname{enc\text{-}in}(\operatorname{enc\text{-}nth}(\operatorname{enc\text{-}num}(k), x), \operatorname{nth}(k, l))] \in \operatorname{derivable}(t) \wedge \\ \quad n \leq \operatorname{length}(l) \Rightarrow \\ \quad \quad {[}g, \operatorname{enc\text{-}cartl}(x, l, n)] \in \operatorname{derivable}(t) \end{array}}}
\vnbprop{der-and-parts}{\vnbstmt{\begin{array}{@{}l@{}} \forall t, g, a, b. \\ \quad ({[}g, \operatorname{enc\text{-}and}(a, b)] \in \operatorname{derivable}(t)) \Rightarrow  \\ \quad \quad ({[}g, a] \in \operatorname{derivable}(t)) \wedge ({[}g, b] \in \operatorname{derivable}(t)) \end{array}}}
\vnbprop{der-mem-set-gen}{\vnbstmt{\begin{array}{@{}l@{}} \forall t, h, x, y. \\ \quad ((t \subseteq \mathit{clf}) \wedge ((\operatorname{ax\text{-}mem\text{-}set} \in t) \wedge ({[}h, \operatorname{enc\text{-}in}(x, y)] \in \operatorname{derivable}(t)))) \Rightarrow  \\ \quad \quad {[}h, \operatorname{enc\text{-}in}(x, \operatorname{enc\text{-}setc})] \in \operatorname{derivable}(t) \end{array}}}
\vnbprop{adm-cartesian-elim}{\vnbstmt{\begin{array}{@{}l@{}} \forall t, s, c. \\ \quad (\operatorname{rel\text{-}cartesian\text{-}elim}(t, s, c) \wedge \operatorname{hyps\text{-}in}(s, \operatorname{derivable}(t))) \Rightarrow  \\ \quad \quad c \in \operatorname{derivable}(t) \end{array}}}
\vnbprop{adm-cartesian-intro}{\vnbstmt{\begin{array}{@{}l@{}} \forall t, s, c. \\ \quad ((t \subseteq \mathit{clf}) \wedge ((\operatorname{ax\text{-}mem\text{-}set} \in t) \wedge (\operatorname{rel\text{-}cartesian\text{-}intro}(t, s, c) \wedge \operatorname{hyps\text{-}in}(s, \operatorname{derivable}(t))))) \Rightarrow  \\ \quad \quad c \in \operatorname{derivable}(t) \end{array}}}
\vnbprop{idx-minus-in-nn}{\vnbstmt{\forall n, k \in \mathbb{N}.\; ((k \leq n) \Rightarrow ((n - k) \in \mathbb{N}))}}
\vnbprop{idx-list-prop}{\vnbstmt{\begin{array}{@{}l@{}} \forall n \in \mathbb{N}. \\ \quad \operatorname{idx\text{-}list}(n) \in \operatorname{tuples}(\operatorname{expr\text{-}set}) \wedge \\ \quad \operatorname{length}(\operatorname{idx\text{-}list}(n)) = n \wedge \\ \quad \forall i \in \mathbb{N}.\; (((1 \leq i) \wedge (i \leq n)) \Rightarrow (\operatorname{nth}(i, \operatorname{idx\text{-}list}(n)) = \operatorname{ebvar}((n - i)))) \end{array}}}
\vnbprop{cchs-idx}{\vnbstmt{\begin{array}{@{}l@{}} \forall k \in \mathbb{N}, x, a.; \forall n \in \mathbb{N}. \\ \quad (\operatorname{enc\text{-}cch}(x, a, n, k) \simeq \operatorname{cchs}(x, a, n, \operatorname{idx\text{-}list}(n), k)) \end{array}}}
\vnbprop{cchs-expr}{\vnbstmt{\begin{array}{@{}l@{}} \forall k \in \mathbb{N}, x, a, n, l. \\ \quad x \in \mathit{clt} \wedge \\ \quad a \in \operatorname{tuples}(\operatorname{expr\text{-}set}) \wedge \\ \quad \forall u \in \operatorname{make\text{-}set}(a).\; (u \in \mathit{clt}) \wedge \\ \quad n \in \mathbb{N} \wedge \\ \quad n \leq \operatorname{length}(a) \wedge \\ \quad l \in \operatorname{tuples}(\operatorname{expr\text{-}set}) \wedge \\ \quad k \leq n \Rightarrow \\ \quad \quad \operatorname{cchs}(x, a, n, l, k) \in \operatorname{expr\text{-}set} \end{array}}}
\vnbprop{maplk-open}{\vnbstmt{\begin{array}{@{}l@{}} \forall u, l, x. \\ \quad ((u \in \operatorname{expr\text{-}set}) \wedge ((l \in \operatorname{tuples}(\operatorname{expr\text{-}set})) \wedge (x \in \mathbb{N}))) \Rightarrow  \\ \quad \quad \operatorname{maplk}(\operatorname{open\text{-}fun}(u), l, x) \in \operatorname{tuples}(\operatorname{expr\text{-}set}) \wedge \\ \quad \quad \operatorname{length}(\operatorname{maplk}(\operatorname{open\text{-}fun}(u), l, x)) = \operatorname{length}(l) \wedge \\ \quad \quad \forall i \in \mathbb{N}. \\ \quad \quad \quad ((1 \leq i) \wedge (i \leq \operatorname{length}(l))) \Rightarrow  \\ \quad \quad \quad \quad \operatorname{nth}(i, \operatorname{maplk}(\operatorname{open\text{-}fun}(u), l, x)) = \operatorname{expr\text{-}open}(\operatorname{nth}(i, l), x, u) \end{array}}}
\vnbprop{cchs-open-body}{\vnbstmt{\begin{array}{@{}l@{}} \forall t, a, c. \\ \quad ((t \in \operatorname{expr\text{-}set}) \wedge ((a \in \mathit{clt}) \wedge (c \in \operatorname{expr\text{-}set}))) \Rightarrow  \\ \quad \quad (\operatorname{expr\text{-}open}(\operatorname{enc\text{-}and}(\operatorname{enc\text{-}in}(\operatorname{ebvar}(0), a), c), 0, t) \simeq \operatorname{enc\text{-}and}(\operatorname{enc\text{-}in}(t, a), \operatorname{expr\text{-}open}(c, 0, t))) \end{array}}}
\vnbprop{cchs-open}{\vnbstmt{\begin{array}{@{}l@{}} \forall k \in \mathbb{N}, x, a, n, l, u.; \forall d \in \mathbb{N}. \\ \quad x \in \mathit{clt} \wedge \\ \quad a \in \operatorname{tuples}(\operatorname{expr\text{-}set}) \wedge \\ \quad \forall \mathit{mczu} \in \operatorname{make\text{-}set}(a).\; (\mathit{mczu} \in \mathit{clt}) \wedge \\ \quad n \in \mathbb{N} \wedge \\ \quad n \leq \operatorname{length}(a) \wedge \\ \quad l \in \operatorname{tuples}(\operatorname{expr\text{-}set}) \wedge \\ \quad u \in \operatorname{expr\text{-}set} \wedge \\ \quad k \leq n \Rightarrow \\ \quad \quad (\operatorname{expr\text{-}open}(\operatorname{cchs}(x, a, n, l, k), d, u) \simeq \operatorname{cchs}(x, a, n, \operatorname{maplk}(\operatorname{open\text{-}fun}(u), l, (d + k)), k)) \end{array}}}
\vnbprop{cchs-fwd}{\vnbstmt{\begin{array}{@{}l@{}} \forall k \in \mathbb{N}, t, g, x, a, n.; \forall l \in \operatorname{tuples}(\operatorname{expr\text{-}set}). \\ \quad g \in \operatorname{ctx\text{-}set} \wedge \\ \quad x \in \mathit{clt} \wedge \\ \quad a \in \operatorname{tuples}(\operatorname{expr\text{-}set}) \wedge \\ \quad \forall u \in \operatorname{make\text{-}set}(a).\; (u \in \mathit{clt}) \wedge \\ \quad n \in \mathbb{N} \wedge \\ \quad \operatorname{length}(a) = n \wedge \\ \quad 1 \leq n \wedge \\ \quad {[}g, \operatorname{enc\text{-}in}(x, \operatorname{enc\text{-}cart}(a))] \in \operatorname{derivable}(t) \wedge \\ \quad \forall j \in \mathbb{N}. \\ \quad \quad ((1 \leq j) \wedge (j \leq n)) \Rightarrow  \\ \quad \quad \quad {[}g, \operatorname{enc\text{-}in}(\operatorname{enc\text{-}nth}(\operatorname{enc\text{-}num}(j), x), \operatorname{nth}(j, a))] \in \operatorname{derivable}(t) \wedge \\ \quad \operatorname{length}(l) = n \wedge \\ \quad k \leq n \wedge \\ \quad \forall j \in \mathbb{N}. \\ \quad \quad ((1 \leq j) \wedge ((j \leq n) \wedge ((j + k) \leq n))) \Rightarrow (\operatorname{nth}(j, l) = \operatorname{enc\text{-}nth}(\operatorname{enc\text{-}num}(j), x)) \wedge \\ \quad \forall j \in \mathbb{N}. \\ \quad \quad ((1 \leq j) \wedge ((j \leq n) \wedge (n < (j + k)))) \Rightarrow (\operatorname{nth}(j, l) = \operatorname{ebvar}((n - j))) \wedge \\ \quad \operatorname{cchs}(x, a, n, l, k) \in \mathit{clf} \Rightarrow \\ \quad \quad {[}g, \operatorname{cchs}(x, a, n, l, k)] \in \operatorname{derivable}(t) \end{array}}}
\vnbprop{cchs-bwd}{\vnbstmt{\begin{array}{@{}l@{}} \forall k \in \mathbb{N}, t, x, a, n.; \forall h \in \operatorname{ctx\text{-}set}, l \in \operatorname{tuples}(\operatorname{expr\text{-}set}). \\ \quad t \subseteq \mathit{clf} \wedge \\ \quad \operatorname{ax\text{-}mem\text{-}set} \in t \wedge \\ \quad x \in \mathit{clt} \wedge \\ \quad a \in \operatorname{tuples}(\operatorname{expr\text{-}set}) \wedge \\ \quad \forall u \in \operatorname{make\text{-}set}(a).\; (u \in \mathit{clt}) \wedge \\ \quad n \in \mathbb{N} \wedge \\ \quad \operatorname{length}(a) = n \wedge \\ \quad 1 \leq n \wedge \\ \quad \operatorname{enc\text{-}in}(x, \operatorname{enc\text{-}cart}(a)) \in \mathit{clf} \wedge \\ \quad \operatorname{length}(l) = n \wedge \\ \quad k \leq n \wedge \\ \quad \forall j \in \mathbb{N}. \\ \quad \quad ((1 \leq j) \wedge ((j \leq n) \wedge (n < (j + k)))) \Rightarrow (\operatorname{nth}(j, l) = \operatorname{ebvar}((n - j))) \wedge \\ \quad \forall j \in \mathbb{N}. \\ \quad \quad ((1 \leq j) \wedge ((j \leq n) \wedge ((j + k) \leq n))) \Rightarrow  \\ \quad \quad \quad {[}h, \operatorname{enc\text{-}in}(\operatorname{nth}(j, l), \operatorname{nth}(j, a))] \in \operatorname{derivable}(t) \wedge \\ \quad {[}h, \operatorname{cchs}(x, a, n, l, k)] \in \operatorname{derivable}(t) \Rightarrow \\ \quad \quad {[}h, \operatorname{enc\text{-}in}(x, \operatorname{enc\text{-}cart}(a))] \in \operatorname{derivable}(t) \end{array}}}
\vnbprop{der-cartesian-decompose}{\vnbstmt{\begin{array}{@{}l@{}} \forall t, g, x, a. \\ \quad t \subseteq \mathit{clf} \wedge \\ \quad \operatorname{ax\text{-}mem\text{-}set} \in t \wedge \\ \quad g \in \operatorname{ctx\text{-}set} \wedge \\ \quad 1 \leq \operatorname{length}(a) \wedge \\ \quad \operatorname{enc\text{-}in}(x, \operatorname{enc\text{-}cart}(a)) \in \mathit{clf} \wedge \\ \quad \operatorname{enc\text{-}cch}(x, a, \operatorname{length}(a), \operatorname{length}(a)) \in \mathit{clf} \Rightarrow \\ \quad \quad {[}g, \operatorname{enc\text{-}iff}(\operatorname{enc\text{-}in}(x, \operatorname{enc\text{-}cart}(a)), \operatorname{enc\text{-}cch}(x, a, \operatorname{length}(a), \operatorname{length}(a)))] \in \operatorname{derivable}(t) \end{array}}}
\vnbprop{j-cart-sub}{\vnbstmt{\forall t.\; (((t \subseteq \mathit{clf}) \wedge (\operatorname{ax\text{-}mem\text{-}set} \in t)) \Rightarrow (\operatorname{j\text{-}cart}(t) \subseteq \operatorname{eqv}(t)))}}
\vnbprop{adm-cartesian-decompose-1}{\vnbstmt{\begin{array}{@{}l@{}} \forall t, s, c. \\ \quad ((t \subseteq \mathit{clf}) \wedge ((\operatorname{ax\text{-}mem\text{-}set} \in t) \wedge (\operatorname{rel\text{-}cartesian\text{-}decompose\text{-}1}(t, s, c) \wedge \operatorname{hyps\text{-}in}(s, \operatorname{derivable}(t))))) \Rightarrow  \\ \quad \quad c \in \operatorname{derivable}(t) \end{array}}}
\vnbprop{adm-cartesian-decompose-7}{\vnbstmt{\begin{array}{@{}l@{}} \forall t, s, c. \\ \quad ((t \subseteq \mathit{clf}) \wedge ((\operatorname{ax\text{-}mem\text{-}set} \in t) \wedge (\operatorname{rel\text{-}cartesian\text{-}decompose\text{-}7}(t, s, c) \wedge \operatorname{hyps\text{-}in}(s, \operatorname{derivable}(t))))) \Rightarrow  \\ \quad \quad c \in \operatorname{derivable}(t) \end{array}}}
\vnbprop{adm-def-len-tuples}{\vnbstmt{\begin{array}{@{}l@{}} \forall t, s, c. \\ \quad t \subseteq \mathit{clf} \wedge \\ \quad \operatorname{ax\text{-}len\text{-}nn} \in t \wedge \\ \quad \operatorname{adm\text{-}def\text{-}strict\text{-}in\text{-}l}(t) \wedge \\ \quad \operatorname{adm\text{-}def\text{-}strict\text{-}in\text{-}r}(t) \wedge \\ \quad \operatorname{def\text{-}len\text{-}tuples}(t, s, c) \wedge \\ \quad \operatorname{hyps\text{-}in}(s, \operatorname{derivable}(t)) \Rightarrow \\ \quad \quad c \in \operatorname{derivable}(t) \end{array}}}
\vnbprop{clt-eapp-arg}{\vnbstmt{\begin{array}{@{}l@{}} \forall h, l, i. \\ \quad ((\operatorname{eapp}(h, l) \in \mathit{clt}) \wedge ((i \in \mathbb{N}) \wedge ((1 \leq i) \wedge (i \leq \operatorname{length}(l))))) \Rightarrow  \\ \quad \quad \operatorname{nth}(i, l) \in \mathit{clt} \end{array}}}
\vnbprop{enc-fun-unfold}{\vnbstmt{\forall p, q.\; (\operatorname{enc\text{-}fun}(p, q) \simeq \operatorname{eapp}(19, {[}p, q]))}}
\vnbprop{enc-lam-unfold}{\vnbstmt{\forall p, q.\; (\operatorname{enc\text{-}lam}(p, q) \simeq \operatorname{ebind}(26, 1, {[}p], q))}}
\vnbprop{adm-lambda-type-1}{\vnbstmt{\begin{array}{@{}l@{}} \forall t, s, c. \\ \quad (\operatorname{rel\text{-}lambda\text{-}type\text{-}1}(t, s, c) \wedge \operatorname{hyps\text{-}in}(s, \operatorname{derivable}(t))) \Rightarrow  \\ \quad \quad c \in \operatorname{derivable}(t) \end{array}}}
\vnbprop{adm-lambda-type-2}{\vnbstmt{\begin{array}{@{}l@{}} \forall t, s, c. \\ \quad (\operatorname{rel\text{-}lambda\text{-}type\text{-}2}(t, s, c) \wedge \operatorname{hyps\text{-}in}(s, \operatorname{derivable}(t))) \Rightarrow  \\ \quad \quad c \in \operatorname{derivable}(t) \end{array}}}
\vnbprop{adm-lambda-type-3}{\vnbstmt{\begin{array}{@{}l@{}} \forall t, s, c. \\ \quad (\operatorname{rel\text{-}lambda\text{-}type\text{-}3}(t, s, c) \wedge \operatorname{hyps\text{-}in}(s, \operatorname{derivable}(t))) \Rightarrow  \\ \quad \quad c \in \operatorname{derivable}(t) \end{array}}}
\vnbprop{adm-lambda-type-4}{\vnbstmt{\begin{array}{@{}l@{}} \forall t, s, c. \\ \quad (\operatorname{rel\text{-}lambda\text{-}type\text{-}4}(t, s, c) \wedge \operatorname{hyps\text{-}in}(s, \operatorname{derivable}(t))) \Rightarrow  \\ \quad \quad c \in \operatorname{derivable}(t) \end{array}}}
\vnbprop{ltch-expr}{\vnbstmt{\begin{array}{@{}l@{}} \forall k \in \mathbb{N}, a, b, d, n. \\ \quad a \in \operatorname{tuples}(\operatorname{expr\text{-}set}) \wedge \\ \quad \forall u \in \operatorname{make\text{-}set}(a).\; (u \in \mathit{clt}) \wedge \\ \quad n \in \mathbb{N} \wedge \\ \quad n \leq \operatorname{length}(a) \wedge \\ \quad b \in \operatorname{expr\text{-}set} \wedge \\ \quad d \in \mathit{clt} \wedge \\ \quad k \leq n \Rightarrow \\ \quad \quad \operatorname{ltch}(a, b, d, n, k) \in \operatorname{expr\text{-}set} \end{array}}}
\vnbprop{ltop-expr}{\vnbstmt{\begin{array}{@{}l@{}} \forall k \in \mathbb{N}, b, z, n. \\ \quad ((b \in \operatorname{expr\text{-}set}) \wedge ((z \in \mathbb{N}) \wedge ((n \in \mathbb{N}) \wedge (k \leq n)))) \Rightarrow  \\ \quad \quad \operatorname{ltop}(b, z, n, k) \in \operatorname{expr\text{-}set} \end{array}}}
\vnbprop{ltch-open}{\vnbstmt{\begin{array}{@{}l@{}} \forall k \in \mathbb{N}, a, b, d, n, u.; \forall j \in \mathbb{N}. \\ \quad a \in \operatorname{tuples}(\operatorname{expr\text{-}set}) \wedge \\ \quad \forall \mathit{mczu} \in \operatorname{make\text{-}set}(a).\; (\mathit{mczu} \in \mathit{clt}) \wedge \\ \quad n \in \mathbb{N} \wedge \\ \quad n \leq \operatorname{length}(a) \wedge \\ \quad b \in \operatorname{expr\text{-}set} \wedge \\ \quad d \in \mathit{clt} \wedge \\ \quad u \in \operatorname{expr\text{-}set} \wedge \\ \quad k \leq n \Rightarrow \\ \quad \quad (\operatorname{expr\text{-}open}(\operatorname{ltch}(a, b, d, n, k), j, u) \simeq \operatorname{ltch}(a, \operatorname{expr\text{-}open}(b, (j + k), u), d, n, k)) \end{array}}}
\vnbprop{der-imp-mp}{\vnbstmt{\begin{array}{@{}l@{}} \forall t, g, a, c. \\ \quad (({[}g, \operatorname{enc\text{-}imp}(a, c)] \in \operatorname{derivable}(t)) \wedge ({[}g, a] \in \operatorname{derivable}(t))) \Rightarrow  \\ \quad \quad {[}g, c] \in \operatorname{derivable}(t) \end{array}}}
\vnbprop{ltch-inst}{\vnbstmt{\begin{array}{@{}l@{}} \forall k \in \mathbb{N}, t, g, a, b, d, n, z.; \forall c \in \operatorname{expr\text{-}set}. \\ \quad g \in \operatorname{ctx\text{-}set} \wedge \\ \quad a \in \operatorname{tuples}(\operatorname{expr\text{-}set}) \wedge \\ \quad \forall u \in \operatorname{make\text{-}set}(a).\; (u \in \mathit{clt}) \wedge \\ \quad n \in \mathbb{N} \wedge \\ \quad n \leq \operatorname{length}(a) \wedge \\ \quad b \in \operatorname{expr\text{-}set} \wedge \\ \quad d \in \mathit{clt} \wedge \\ \quad z \in \mathbb{N} \wedge \\ \quad \forall i \in \mathbb{N}. \\ \quad \quad ((1 \leq i) \wedge (i \leq n)) \Rightarrow  \\ \quad \quad \quad {[}g, \operatorname{enc\text{-}in}(\operatorname{enc\text{-}nth}(\operatorname{enc\text{-}num}(i), \operatorname{efvar}(z)), \operatorname{nth}(i, a))] \in \operatorname{derivable}(t) \wedge \\ \quad k \leq n \wedge \\ \quad c = \operatorname{ltop}(b, z, n, (n - k)) \wedge \\ \quad {[}g, \operatorname{ltch}(a, c, d, n, k)] \in \operatorname{derivable}(t) \Rightarrow \\ \quad \quad {[}g, \operatorname{enc\text{-}in}(\operatorname{ltop}(b, z, n, n), d)] \in \operatorname{derivable}(t) \end{array}}}
\vnbprop{adm-lambda-type-n}{\vnbstmt{\begin{array}{@{}l@{}} \forall t, s, c. \\ \quad (\operatorname{rel\text{-}lambda\text{-}type\text{-}n}(t, s, c) \wedge \operatorname{hyps\text{-}in}(s, \operatorname{derivable}(t))) \Rightarrow  \\ \quad \quad c \in \operatorname{derivable}(t) \end{array}}}

\vnbbookfile{theorem-library/meta-def-laws}{theorem-library/meta-def-laws}
\noindent{\small meta-def-laws.scm -- K2 PHASE 5: the definedness rules of structure-library/meta-def.scm against the core (docs/meta-k2-admissibility-2026-10-08.md, section 4.7; report scratchpad/triage/META-5-REPORT.md). Agent META-5, 2026-10-12.  Every theorem here is (qed ...); helper procedures are local (prefix md-). (1) the two ...}\par
\vnbprop{def-in-closed-from-strict}{\vnbstmt{\forall t.\; (\operatorname{adm\text{-}def\text{-}strict\text{-}in\text{-}l}(t) \Rightarrow \operatorname{def\text{-}in\text{-}closed}(t))}}
\vnbprop{def-comp-closed-from-schema}{\vnbstmt{\forall t.\; (\operatorname{adm\text{-}def\text{-}comp}(t) \Rightarrow \operatorname{def\text{-}comp\text{-}closed}(t))}}
\vnbprop{clf-eq-elim}{\vnbstmt{\forall x, y.\; ((\operatorname{enc\text{-}eq}(x, y) \in \mathit{clf}) \Rightarrow ((x \in \mathit{clt}) \wedge (y \in \mathit{clt})))}}
\vnbprop{adm-def-strict-eq-l}{\vnbstmt{\begin{array}{@{}l@{}} \forall t, s, c. \\ \quad (\operatorname{def\text{-}strict\text{-}eq\text{-}l}(t, s, c) \wedge \operatorname{hyps\text{-}in}(s, \operatorname{derivable}(t))) \Rightarrow (c \in \operatorname{derivable}(t)) \end{array}}}
\vnbprop{adm-def-strict-eq-r}{\vnbstmt{\begin{array}{@{}l@{}} \forall t, s, c. \\ \quad (\operatorname{def\text{-}strict\text{-}eq\text{-}r}(t, s, c) \wedge \operatorname{hyps\text{-}in}(s, \operatorname{derivable}(t))) \Rightarrow (c \in \operatorname{derivable}(t)) \end{array}}}
\vnbprop{def-strict-in-l-admitted}{\vnbstmt{\begin{array}{@{}l@{}} \forall t, s, c. \\ \quad (\operatorname{def\text{-}strict\text{-}in\text{-}l}(t, s, c) \wedge \operatorname{hyps\text{-}in}(s, \operatorname{derivable}(t))) \Rightarrow (c \in \operatorname{derivable}(t)) \end{array}}}
\vnbprop{def-strict-in-r-admitted}{\vnbstmt{\begin{array}{@{}l@{}} \forall t, s, c. \\ \quad (\operatorname{def\text{-}strict\text{-}in\text{-}r}(t, s, c) \wedge \operatorname{hyps\text{-}in}(s, \operatorname{derivable}(t))) \Rightarrow (c \in \operatorname{derivable}(t)) \end{array}}}
\vnbprop{def-total-admitted}{\vnbstmt{\begin{array}{@{}l@{}} \forall t, s, c. \\ \quad (\operatorname{def\text{-}total}(t, s, c) \wedge \operatorname{hyps\text{-}in}(s, \operatorname{derivable}(t))) \Rightarrow (c \in \operatorname{derivable}(t)) \end{array}}}
\vnbprop{def-binder-admitted}{\vnbstmt{\begin{array}{@{}l@{}} \forall t, s, c. \\ \quad (\operatorname{def\text{-}binder}(t, s, c) \wedge \operatorname{hyps\text{-}in}(s, \operatorname{derivable}(t))) \Rightarrow (c \in \operatorname{derivable}(t)) \end{array}}}
\vnbprop{def-comp-admitted}{\vnbstmt{\begin{array}{@{}l@{}} \forall t, s, c. \\ \quad (\operatorname{def\text{-}comp}(t, s, c) \wedge \operatorname{hyps\text{-}in}(s, \operatorname{derivable}(t))) \Rightarrow (c \in \operatorname{derivable}(t)) \end{array}}}
\vnbprop{def-list-admitted}{\vnbstmt{\begin{array}{@{}l@{}} \forall t, s, c. \\ \quad (\operatorname{def\text{-}list}(t, s, c) \wedge \operatorname{hyps\text{-}in}(s, \operatorname{derivable}(t))) \Rightarrow (c \in \operatorname{derivable}(t)) \end{array}}}
\vnbprop{adm-def-strict-in-l-holds}{\vnbstmt{\forall t.\; \operatorname{adm\text{-}def\text{-}strict\text{-}in\text{-}l}(t)}}
\vnbprop{adm-def-strict-in-r-holds}{\vnbstmt{\forall t.\; \operatorname{adm\text{-}def\text{-}strict\text{-}in\text{-}r}(t)}}
\vnbprop{adm-def-comp-holds}{\vnbstmt{\forall t.\; \operatorname{adm\text{-}def\text{-}comp}(t)}}
\vnbprop{def-in-closed-holds}{\vnbstmt{\forall t.\; \operatorname{def\text{-}in\text{-}closed}(t)}}
\vnbprop{def-comp-closed-holds}{\vnbstmt{\forall t.\; \operatorname{def\text{-}comp\text{-}closed}(t)}}
\vnbprop{adm-big-union-mem-intro-core}{\vnbstmt{\begin{array}{@{}l@{}} \forall t, s, c. \\ \quad (\operatorname{rel\text{-}big\text{-}union\text{-}mem\text{-}intro}(t, s, c) \wedge \operatorname{hyps\text{-}in}(s, \operatorname{derivable}(t))) \Rightarrow  \\ \quad \quad c \in \operatorname{derivable}(t) \end{array}}}
\vnbprop{adm-transfinite-induction-core}{\vnbstmt{\begin{array}{@{}l@{}} \forall t, s, c. \\ \quad t \subseteq \mathit{clf} \wedge \\ \quad \operatorname{ax\text{-}tfi} \in t \wedge \\ \quad \operatorname{ax\text{-}mem\text{-}set} \in t \wedge \\ \quad \operatorname{rel\text{-}transfinite\text{-}induction\text{-}nbg}(t, s, c) \wedge \\ \quad \operatorname{hyps\text{-}in}(s, \operatorname{derivable}(t)) \Rightarrow \\ \quad \quad c \in \operatorname{derivable}(t) \end{array}}}
\vnbprop{adm-nn-induction-core}{\vnbstmt{\begin{array}{@{}l@{}} \forall t, s, c. \\ \quad t \subseteq \mathit{clf} \wedge \\ \quad \operatorname{ax\text{-}nn\text{-}ind} \in t \wedge \\ \quad \operatorname{ax\text{-}mem\text{-}set} \in t \wedge \\ \quad \operatorname{ax\text{-}zero\text{-}nn} \in t \wedge \\ \quad \operatorname{ax\text{-}succ\text{-}nn} \in t \wedge \\ \quad \operatorname{rel\text{-}nn\text{-}induction\text{-}nbg}(t, s, c) \wedge \\ \quad \operatorname{hyps\text{-}in}(s, \operatorname{derivable}(t)) \Rightarrow \\ \quad \quad c \in \operatorname{derivable}(t) \end{array}}}
\vnbprop{adm-tfi-3cases-core}{\vnbstmt{\begin{array}{@{}l@{}} \forall t, s, c. \\ \quad t \subseteq \mathit{clf} \wedge \\ \quad \operatorname{ax\text{-}tfi} \in t \wedge \\ \quad \operatorname{ax\text{-}mem\text{-}set} \in t \wedge \\ \quad \operatorname{ax\text{-}limit\text{-}iff} \in t \wedge \\ \quad \operatorname{ax\text{-}ord\text{-}succ\text{-}above} \in t \wedge \\ \quad \operatorname{rel\text{-}tfi\text{-}3cases\text{-}nbg}(t, s, c) \wedge \\ \quad \operatorname{hyps\text{-}in}(s, \operatorname{derivable}(t)) \Rightarrow \\ \quad \quad c \in \operatorname{derivable}(t) \end{array}}}
\vnbprop{def-strict-arg-admitted}{\vnbstmt{\begin{array}{@{}l@{}} \forall t, s, c. \\ \quad (\operatorname{def\text{-}strict\text{-}arg}(t, s, c) \wedge \operatorname{hyps\text{-}in}(s, \operatorname{derivable}(t))) \Rightarrow (c \in \operatorname{derivable}(t)) \end{array}}}
\vnbprop{def-strict-dom-admitted}{\vnbstmt{\begin{array}{@{}l@{}} \forall t, s, c. \\ \quad (\operatorname{def\text{-}strict\text{-}dom}(t, s, c) \wedge \operatorname{hyps\text{-}in}(s, \operatorname{derivable}(t))) \Rightarrow (c \in \operatorname{derivable}(t)) \end{array}}}
\vnbprop{adm-def-strict-arg-holds}{\vnbstmt{\forall t.\; \operatorname{adm\text{-}def\text{-}strict\text{-}arg}(t)}}
\vnbprop{adm-def-strict-dom-holds}{\vnbstmt{\forall t.\; \operatorname{adm\text{-}def\text{-}strict\text{-}dom}(t)}}
\vnbprop{clt-ebind-dom}{\vnbstmt{\forall h, a, b.\; (((h \in \mathbb{N}) \wedge (\operatorname{ebind}(h, 1, {[}a], b) \in \mathit{clt})) \Rightarrow (a \in \mathit{clt}))}}
\vnbprop{clf-sep-open}{\vnbstmt{\begin{array}{@{}l@{}} \forall y, d, p. \\ \quad (\operatorname{enc\text{-}in}(y, \operatorname{enc\text{-}sep}(d, p)) \in \mathit{clf}) \Rightarrow (\operatorname{expr\text{-}open}(p, 0, y) \in \mathit{clf}) \end{array}}}
\vnbprop{strict-reach-unfold}{\vnbstmt{\begin{array}{@{}l@{}} \forall e. \\ \quad (\operatorname{strict\text{-}reach}(e) \simeq \{z \in \operatorname{expr\text{-}set} : \forall w \in \operatorname{power}(\operatorname{expr\text{-}set}).\; (((e \in w) \wedge \operatorname{strict\text{-}closed}(w)) \Rightarrow (z \in w))\}) \end{array}}}
\vnbprop{true-reach-unfold}{\vnbstmt{\begin{array}{@{}l@{}} \forall j. \\ \quad (\operatorname{true\text{-}reach}(j) \simeq \{z \in \operatorname{expr\text{-}set} : \forall w \in \operatorname{power}(\operatorname{expr\text{-}set}).\; (((j \in w) \wedge \operatorname{true\text{-}closed}(w)) \Rightarrow (z \in w))\}) \end{array}}}
\vnbprop{strict-reach-least}{\vnbstmt{\begin{array}{@{}l@{}} \forall e.; \forall s \in \operatorname{power}(\operatorname{expr\text{-}set}), x \in \operatorname{strict\text{-}reach}(e). \\ \quad ((e \in s) \wedge \operatorname{strict\text{-}closed}(s)) \Rightarrow (x \in s) \end{array}}}
\vnbprop{true-reach-least}{\vnbstmt{\begin{array}{@{}l@{}} \forall e.; \forall s \in \operatorname{power}(\operatorname{expr\text{-}set}), x \in \operatorname{true\text{-}reach}(e). \\ \quad ((e \in s) \wedge \operatorname{true\text{-}closed}(s)) \Rightarrow (x \in s) \end{array}}}
\vnbprop{adm-strict-walk}{\vnbstmt{\begin{array}{@{}l@{}} \forall t, g, s, x. \\ \quad (({[}g, \operatorname{enc\text{-}eq}(s, s)] \in \operatorname{derivable}(t)) \wedge (x \in \operatorname{strict\text{-}reach}(s))) \Rightarrow  \\ \quad \quad {[}g, \operatorname{enc\text{-}eq}(x, x)] \in \operatorname{derivable}(t) \end{array}}}
\vnbprop{adm-true-walk}{\vnbstmt{\begin{array}{@{}l@{}} \forall t, g, f, x. \\ \quad (({[}g, f] \in \operatorname{derivable}(t)) \wedge (x \in \operatorname{true\text{-}reach}(f))) \Rightarrow ({[}g, x] \in \operatorname{derivable}(t)) \end{array}}}
\vnbprop{adm-strict-atom}{\vnbstmt{\begin{array}{@{}l@{}} \forall t, g, f, w. \\ \quad (({[}g, f] \in \operatorname{derivable}(t)) \wedge \operatorname{strict\text{-}atom}(f, w)) \Rightarrow  \\ \quad \quad {[}g, \operatorname{enc\text{-}eq}(w, w)] \in \operatorname{derivable}(t) \end{array}}}
\vnbprop{def-strict-hyp-derived}{\vnbstmt{\begin{array}{@{}l@{}} \forall t, s, c. \\ \quad (\operatorname{def\text{-}strict\text{-}hyp}(t, s, c) \wedge \operatorname{hyps\text{-}in}(s, \operatorname{derivable}(t))) \Rightarrow (c \in \operatorname{derivable}(t)) \end{array}}}
\vnbprop{adm-def-strict-hyp-holds}{\vnbstmt{\forall t.\; \operatorname{adm\text{-}def\text{-}strict\text{-}hyp}(t)}}
\vnbprop{def-nth-list-derived}{\vnbstmt{\begin{array}{@{}l@{}} \forall t, s, c. \\ \quad (\operatorname{def\text{-}nth\text{-}list}(t, s, c) \wedge \operatorname{hyps\text{-}in}(s, \operatorname{derivable}(t))) \Rightarrow (c \in \operatorname{derivable}(t)) \end{array}}}
\vnbprop{adm-def-nth-list-holds}{\vnbstmt{\forall t.\; \operatorname{adm\text{-}def\text{-}nth\text{-}list}(t)}}
\vnbprop{def-strict-pred-admitted}{\vnbstmt{\begin{array}{@{}l@{}} \forall t, s, c. \\ \quad (\operatorname{def\text{-}strict\text{-}pred}(t, s, c) \wedge \operatorname{hyps\text{-}in}(s, \operatorname{derivable}(t))) \Rightarrow (c \in \operatorname{derivable}(t)) \end{array}}}
\vnbprop{adm-def-strict-pred-holds}{\vnbstmt{\forall t.\; \operatorname{adm\text{-}def\text{-}strict\text{-}pred}(t)}}
\vnbprop{def-strict-pred-refuses-qeq}{\vnbstmt{\begin{array}{@{}l@{}} \forall t, g, a, b, c. \\ \quad ((g \in \operatorname{ctx\text{-}set}) \wedge ((a \in \mathit{clt}) \wedge (b \in \mathit{clt}))) \Rightarrow  \\ \quad \quad \neg \operatorname{cr\text{-}def\text{-}strict\text{-}pred}(t, {[}{[}g, \operatorname{enc\text{-}qeq}(a, b)]], c) \end{array}}}
\vnbprop{def-strict-pred-refuses-not}{\vnbstmt{\begin{array}{@{}l@{}} \forall t, g, a, c. \\ \quad ((g \in \operatorname{ctx\text{-}set}) \wedge (a \in \mathit{clt})) \Rightarrow  \\ \quad \quad \neg \operatorname{cr\text{-}def\text{-}strict\text{-}pred}(t, {[}{[}g, \operatorname{enc\text{-}not}(\operatorname{enc\text{-}eq}(a, a))]], c) \end{array}}}
\vnbprop{strict-atom-pred2}{\vnbstmt{\begin{array}{@{}l@{}} \forall k, x, y, w. \\ \quad (\operatorname{pred\text{-}code}(k) \wedge ((x \in \operatorname{expr\text{-}set}) \wedge ((y \in \operatorname{expr\text{-}set}) \wedge ((w \in \operatorname{strict\text{-}reach}(x)) \vee (w \in \operatorname{strict\text{-}reach}(y)))))) \Rightarrow  \\ \quad \quad \operatorname{strict\text{-}atom}(\operatorname{eapp}(10, {[}\operatorname{econst}(k), x, y]), w) \end{array}}}
\vnbprop{strict-atom-le}{\vnbstmt{\begin{array}{@{}l@{}} \forall x, y, w. \\ \quad ((x \in \operatorname{expr\text{-}set}) \wedge ((y \in \operatorname{expr\text{-}set}) \wedge ((w \in \operatorname{strict\text{-}reach}(x)) \vee (w \in \operatorname{strict\text{-}reach}(y))))) \Rightarrow  \\ \quad \quad \operatorname{strict\text{-}atom}(\operatorname{enc\text{-}le}(x, y), w) \end{array}}}
\vnbprop{strict-atom-lt}{\vnbstmt{\begin{array}{@{}l@{}} \forall x, y, w. \\ \quad ((x \in \operatorname{expr\text{-}set}) \wedge ((y \in \operatorname{expr\text{-}set}) \wedge ((w \in \operatorname{strict\text{-}reach}(x)) \vee (w \in \operatorname{strict\text{-}reach}(y))))) \Rightarrow  \\ \quad \quad \operatorname{strict\text{-}atom}(\operatorname{enc\text{-}lt}(x, y), w) \end{array}}}
\vnbprop{def-struct-hyp-derived}{\vnbstmt{\begin{array}{@{}l@{}} \forall t, s, c. \\ \quad (\operatorname{def\text{-}struct\text{-}hyp}(t, s, c) \wedge \operatorname{hyps\text{-}in}(s, \operatorname{derivable}(t))) \Rightarrow (c \in \operatorname{derivable}(t)) \end{array}}}
\vnbprop{adm-def-struct-hyp-holds}{\vnbstmt{\forall t.\; \operatorname{adm\text{-}def\text{-}struct\text{-}hyp}(t)}}
\vnbprop{adm-def-pred-unfold-holds}{\vnbstmt{\forall t.\; \operatorname{adm\text{-}def\text{-}pred\text{-}unfold}(t)}}
\vnbprop{def-pred-hyp-derived}{\vnbstmt{\begin{array}{@{}l@{}} \forall t, s, c. \\ \quad (\operatorname{def\text{-}pred\text{-}hyp}(t, s, c) \wedge \operatorname{hyps\text{-}in}(s, \operatorname{derivable}(t))) \Rightarrow (c \in \operatorname{derivable}(t)) \end{array}}}
\vnbprop{adm-def-pred-hyp-holds}{\vnbstmt{\forall t.\; \operatorname{adm\text{-}def\text{-}pred\text{-}hyp}(t)}}
\vnbprop{der-qeq-hole}{\vnbstmt{\begin{array}{@{}l@{}} \forall a, g, s, t, p.; \forall z \in \mathbb{N}. \\ \quad p \in \operatorname{expr\text{-}set} \wedge \\ \quad \operatorname{expr\text{-}subst}(p, z, s) \in \mathit{clf} \wedge \\ \quad {[}g, \operatorname{enc\text{-}qeq}(s, t)] \in \operatorname{derivable}(a) \wedge \\ \quad {[}g, \operatorname{expr\text{-}subst}(p, z, t)] \in \operatorname{derivable}(a) \Rightarrow \\ \quad \quad {[}g, \operatorname{expr\text{-}subst}(p, z, s)] \in \operatorname{derivable}(a) \end{array}}}
\vnbprop{def-fun-inst-derived}{\vnbstmt{\begin{array}{@{}l@{}} \forall t, s, c. \\ \quad (\operatorname{def\text{-}fun\text{-}inst}(t, s, c) \wedge \operatorname{hyps\text{-}in}(s, \operatorname{derivable}(t))) \Rightarrow (c \in \operatorname{derivable}(t)) \end{array}}}
\vnbprop{def-fun-strict-derived}{\vnbstmt{\begin{array}{@{}l@{}} \forall t, s, c. \\ \quad (\operatorname{def\text{-}fun\text{-}strict}(t, s, c) \wedge \operatorname{hyps\text{-}in}(s, \operatorname{derivable}(t))) \Rightarrow (c \in \operatorname{derivable}(t)) \end{array}}}
\vnbprop{def-fun-nth-derived}{\vnbstmt{\begin{array}{@{}l@{}} \forall t, s, c. \\ \quad (\operatorname{def\text{-}fun\text{-}nth}(t, s, c) \wedge \operatorname{hyps\text{-}in}(s, \operatorname{derivable}(t))) \Rightarrow (c \in \operatorname{derivable}(t)) \end{array}}}
\vnbprop{def-fun-sep-derived}{\vnbstmt{\begin{array}{@{}l@{}} \forall t, s, c. \\ \quad (\operatorname{def\text{-}fun\text{-}sep}(t, s, c) \wedge \operatorname{hyps\text{-}in}(s, \operatorname{derivable}(t))) \Rightarrow (c \in \operatorname{derivable}(t)) \end{array}}}
\vnbprop{adm-pred-arg}{\vnbstmt{\begin{array}{@{}l@{}} \forall t, g, l, k, i. \\ \quad \operatorname{nth}(1, l) = \operatorname{econst}(k) \wedge \\ \quad \operatorname{pred\text{-}code}(k) \wedge \\ \quad {[}g, \operatorname{eapp}(10, l)] \in \operatorname{derivable}(t) \wedge \\ \quad i \in \mathbb{N} \wedge \\ \quad 2 \leq i \wedge \\ \quad i \leq \operatorname{length}(l) \Rightarrow \\ \quad \quad {[}g, \operatorname{enc\text{-}eq}(\operatorname{nth}(i, l), \operatorname{nth}(i, l))] \in \operatorname{derivable}(t) \end{array}}}
\vnbprop{der-forall-inst}{\vnbstmt{\begin{array}{@{}l@{}} \forall a, g, b, t. \\ \quad (({[}g, \operatorname{enc\text{-}all}(b)] \in \operatorname{derivable}(a)) \wedge ({[}g, \operatorname{enc\text{-}eq}(t, t)] \in \operatorname{derivable}(a))) \Rightarrow  \\ \quad \quad {[}g, \operatorname{expr\text{-}open}(b, 0, t)] \in \operatorname{derivable}(a) \end{array}}}
\vnbprop{inst-reach-unfold}{\vnbstmt{\begin{array}{@{}l@{}} \forall d, l. \\ \quad (\operatorname{inst\text{-}reach}(d, l) \simeq \{z \in \operatorname{expr\text{-}set} : \forall w \in \operatorname{power}(\operatorname{expr\text{-}set}).\; (((d \in w) \wedge \operatorname{inst\text{-}closed}(l, w)) \Rightarrow (z \in w))\}) \end{array}}}
\vnbprop{inst-reach-least}{\vnbstmt{\begin{array}{@{}l@{}} \forall d, l.; \forall s \in \operatorname{power}(\operatorname{expr\text{-}set}), x \in \operatorname{inst\text{-}reach}(d, l). \\ \quad ((d \in s) \wedge \operatorname{inst\text{-}closed}(l, s)) \Rightarrow (x \in s) \end{array}}}
\vnbprop{adm-inst-reach}{\vnbstmt{\begin{array}{@{}l@{}} \forall t, g, l, d, x. \\ \quad {[}g, d] \in \operatorname{derivable}(t) \wedge \\ \quad \forall i \in \mathbb{N}. \\ \quad \quad ((2 \leq i) \wedge (i \leq \operatorname{length}(l))) \Rightarrow  \\ \quad \quad \quad {[}g, \operatorname{enc\text{-}eq}(\operatorname{nth}(i, l), \operatorname{nth}(i, l))] \in \operatorname{derivable}(t) \wedge \\ \quad x \in \operatorname{inst\text{-}reach}(d, l) \Rightarrow \\ \quad \quad {[}g, x] \in \operatorname{derivable}(t) \end{array}}}
\vnbprop{adm-pred-unfold-n}{\vnbstmt{\begin{array}{@{}l@{}} \forall t, g, l, k, d, b. \\ \quad \operatorname{nth}(1, l) = \operatorname{econst}(k) \wedge \\ \quad \operatorname{pred\text{-}code}(k) \wedge \\ \quad {[}g, \operatorname{eapp}(10, l)] \in \operatorname{derivable}(t) \wedge \\ \quad {[}g, d] \in \operatorname{derivable}(t) \wedge \\ \quad \operatorname{enc\text{-}iff}(\operatorname{eapp}(10, l), b) \in \operatorname{inst\text{-}reach}(d, l) \Rightarrow \\ \quad \quad {[}g, b] \in \operatorname{derivable}(t) \end{array}}}
\vnbprop{def-pred-unfold-n-derived}{\vnbstmt{\begin{array}{@{}l@{}} \forall t, s, c. \\ \quad (\operatorname{def\text{-}pred\text{-}unfold\text{-}n}(t, s, c) \wedge \operatorname{hyps\text{-}in}(s, \operatorname{derivable}(t))) \Rightarrow  \\ \quad \quad c \in \operatorname{derivable}(t) \end{array}}}
\vnbprop{adm-def-pred-unfold-n-holds}{\vnbstmt{\forall t.\; \operatorname{adm\text{-}def\text{-}pred\text{-}unfold\text{-}n}(t)}}
\vnbprop{cert-set-unfold}{\vnbstmt{\begin{array}{@{}l@{}} \forall t, g, w. \\ \quad (\operatorname{cert\text{-}set}(t, g, w) \simeq \{z \in \operatorname{expr\text{-}set} : \forall v \in \operatorname{power}(\operatorname{expr\text{-}set}).\; (\operatorname{cert\text{-}closed}(t, g, w, v) \Rightarrow (z \in v))\}) \end{array}}}
\vnbprop{cert-set-least}{\vnbstmt{\begin{array}{@{}l@{}} \forall t, g, w.; \forall s \in \operatorname{power}(\operatorname{expr\text{-}set}), x \in \operatorname{cert\text{-}set}(t, g, w). \\ \quad \operatorname{cert\text{-}closed}(t, g, w, s) \Rightarrow (x \in s) \end{array}}}
\vnbprop{adm-cert}{\vnbstmt{\begin{array}{@{}l@{}} \forall t, g, w, f. \\ \quad (({[}g, f] \in \operatorname{derivable}(t)) \wedge (f \in \operatorname{cert\text{-}set}(t, g, w))) \Rightarrow  \\ \quad \quad {[}g, \operatorname{enc\text{-}eq}(w, w)] \in \operatorname{derivable}(t) \end{array}}}
\vnbprop{def-pred-hyp-n-derived}{\vnbstmt{\begin{array}{@{}l@{}} \forall t, s, c. \\ \quad (\operatorname{def\text{-}pred\text{-}hyp\text{-}n}(t, s, c) \wedge \operatorname{hyps\text{-}in}(s, \operatorname{derivable}(t))) \Rightarrow (c \in \operatorname{derivable}(t)) \end{array}}}
\vnbprop{adm-def-pred-hyp-n-holds}{\vnbstmt{\forall t.\; \operatorname{adm\text{-}def\text{-}pred\text{-}hyp\text{-}n}(t)}}
\vnbprop{def-cert-hyp-derived}{\vnbstmt{\begin{array}{@{}l@{}} \forall t, s, c. \\ \quad (\operatorname{def\text{-}cert\text{-}hyp}(t, s, c) \wedge \operatorname{hyps\text{-}in}(s, \operatorname{derivable}(t))) \Rightarrow (c \in \operatorname{derivable}(t)) \end{array}}}
\vnbprop{adm-def-cert-hyp-holds}{\vnbstmt{\forall t.\; \operatorname{adm\text{-}def\text{-}cert\text{-}hyp}(t)}}
\vnbprop{j-macn-sub}{\vnbstmt{\forall t.\; (\operatorname{j\text{-}macn}(t) \subseteq \operatorname{eqv}(t))}}
\vnbprop{adm-macete-n-1}{\vnbstmt{\begin{array}{@{}l@{}} \forall t, s, c. \\ \quad (\operatorname{rel\text{-}macete\text{-}n\text{-}1}(t, s, c) \wedge \operatorname{hyps\text{-}in}(s, \operatorname{derivable}(t))) \Rightarrow (c \in \operatorname{derivable}(t)) \end{array}}}
\vnbprop{adm-macete-n-hyp-1}{\vnbstmt{\begin{array}{@{}l@{}} \forall t, s, c. \\ \quad (\operatorname{rel\text{-}macete\text{-}n\text{-}hyp\text{-}1}(t, s, c) \wedge \operatorname{hyps\text{-}in}(s, \operatorname{derivable}(t))) \Rightarrow  \\ \quad \quad c \in \operatorname{derivable}(t) \end{array}}}
\vnbprop{adm-macete-n-7}{\vnbstmt{\begin{array}{@{}l@{}} \forall t, s, c. \\ \quad (\operatorname{rel\text{-}macete\text{-}n\text{-}7}(t, s, c) \wedge \operatorname{hyps\text{-}in}(s, \operatorname{derivable}(t))) \Rightarrow (c \in \operatorname{derivable}(t)) \end{array}}}
\vnbprop{adm-macete-n-hyp-7}{\vnbstmt{\begin{array}{@{}l@{}} \forall t, s, c. \\ \quad (\operatorname{rel\text{-}macete\text{-}n\text{-}hyp\text{-}7}(t, s, c) \wedge \operatorname{hyps\text{-}in}(s, \operatorname{derivable}(t))) \Rightarrow  \\ \quad \quad c \in \operatorname{derivable}(t) \end{array}}}
\vnbprop{der-iff-refl}{\vnbstmt{\begin{array}{@{}l@{}} \forall t, g, a. \\ \quad ((g \in \operatorname{ctx\text{-}set}) \wedge (a \in \mathit{clf})) \Rightarrow ({[}g, \operatorname{enc\text{-}iff}(a, a)] \in \operatorname{derivable}(t)) \end{array}}}
\vnbprop{der-iff-trans}{\vnbstmt{\begin{array}{@{}l@{}} \forall t, g, a, b, c. \\ \quad (({[}g, \operatorname{enc\text{-}iff}(a, b)] \in \operatorname{derivable}(t)) \wedge ({[}g, \operatorname{enc\text{-}iff}(b, c)] \in \operatorname{derivable}(t))) \Rightarrow  \\ \quad \quad {[}g, \operatorname{enc\text{-}iff}(a, c)] \in \operatorname{derivable}(t) \end{array}}}
\vnbprop{der-or-assoc}{\vnbstmt{\begin{array}{@{}l@{}} \forall t, g, a, b, c. \\ \quad ((g \in \operatorname{ctx\text{-}set}) \wedge ((a \in \mathit{clf}) \wedge ((b \in \mathit{clf}) \wedge (c \in \mathit{clf})))) \Rightarrow  \\ \quad \quad {[}g, \operatorname{enc\text{-}iff}(\operatorname{enc\text{-}or}(\operatorname{enc\text{-}or}(a, b), c), \operatorname{enc\text{-}or}(a, \operatorname{enc\text{-}or}(b, c)))] \in \operatorname{derivable}(t) \end{array}}}
\vnbprop{der-and-assoc}{\vnbstmt{\begin{array}{@{}l@{}} \forall t, g, a, b, c. \\ \quad ((g \in \operatorname{ctx\text{-}set}) \wedge ((a \in \mathit{clf}) \wedge ((b \in \mathit{clf}) \wedge (c \in \mathit{clf})))) \Rightarrow  \\ \quad \quad {[}g, \operatorname{enc\text{-}iff}(\operatorname{enc\text{-}and}(\operatorname{enc\text{-}and}(a, b), c), \operatorname{enc\text{-}and}(a, \operatorname{enc\text{-}and}(b, c)))] \in \operatorname{derivable}(t) \end{array}}}
\vnbprop{der-or-falsity}{\vnbstmt{\begin{array}{@{}l@{}} \forall t, g, f. \\ \quad ((g \in \operatorname{ctx\text{-}set}) \wedge (f \in \mathit{clf})) \Rightarrow  \\ \quad \quad {[}g, \operatorname{enc\text{-}iff}(\operatorname{enc\text{-}or}(\operatorname{enc\text{-}falsity}, f), f)] \in \operatorname{derivable}(t) \end{array}}}
\vnbprop{der-and-truth}{\vnbstmt{\begin{array}{@{}l@{}} \forall t, g, f. \\ \quad ((g \in \operatorname{ctx\text{-}set}) \wedge (f \in \mathit{clf})) \Rightarrow  \\ \quad \quad {[}g, \operatorname{enc\text{-}iff}(\operatorname{enc\text{-}and}(\operatorname{enc\text{-}truth}, f), f)] \in \operatorname{derivable}(t) \end{array}}}
\vnbprop{rf-set-unfold}{\vnbstmt{\begin{array}{@{}l@{}} \forall h, x, l. \\ \quad (\operatorname{rf\text{-}set}(h, x, l) \simeq \{p \in (\mathbb{N} \times \operatorname{expr\text{-}set}) : \forall w \in \operatorname{power}((\mathbb{N} \times \operatorname{expr\text{-}set})).\; (\operatorname{rf\text{-}closed}(h, x, l, w) \Rightarrow (p \in w))\}) \end{array}}}
\vnbprop{rf-set-least}{\vnbstmt{\begin{array}{@{}l@{}} \forall h, x, l.; \forall s \in \operatorname{power}((\mathbb{N} \times \operatorname{expr\text{-}set})), z \in \operatorname{rf\text{-}set}(h, x, l). \\ \quad \operatorname{rf\text{-}closed}(h, x, l, s) \Rightarrow (z \in s) \end{array}}}
\vnbprop{rf-orl}{\vnbstmt{\begin{array}{@{}l@{}} \forall t, g, x, l, j, f. \\ \quad g \in \operatorname{ctx\text{-}set} \wedge \\ \quad x \in \mathit{clt} \wedge \\ \quad l \in \operatorname{tuples}(\operatorname{expr\text{-}set}) \wedge \\ \quad \forall u \in \operatorname{make\text{-}set}(l).\; (u \in \mathit{clt}) \wedge \\ \quad {[}j, f] \in \operatorname{rf\text{-}set}(4, x, l) \Rightarrow \\ \quad \quad (j + 1) \leq \operatorname{length}(l) \wedge \\ \quad \quad {[}g, \operatorname{enc\text{-}iff}(\operatorname{enc\text{-}orl}(x, l, \operatorname{length}(l)), \operatorname{enc\text{-}or}(\operatorname{enc\text{-}orl}(x, l, j), f))] \in \operatorname{derivable}(t) \end{array}}}
\vnbprop{rf-andl}{\vnbstmt{\begin{array}{@{}l@{}} \forall t, g, x, l, j, f. \\ \quad g \in \operatorname{ctx\text{-}set} \wedge \\ \quad x \in \mathit{clt} \wedge \\ \quad l \in \operatorname{tuples}(\operatorname{expr\text{-}set}) \wedge \\ \quad \forall u \in \operatorname{make\text{-}set}(l).\; (u \in \mathit{clt}) \wedge \\ \quad {[}j, f] \in \operatorname{rf\text{-}set}(3, x, l) \Rightarrow \\ \quad \quad (j + 1) \leq \operatorname{length}(l) \wedge \\ \quad \quad {[}g, \operatorname{enc\text{-}iff}(\operatorname{enc\text{-}andl}(x, l, \operatorname{length}(l)), \operatorname{enc\text{-}and}(\operatorname{enc\text{-}andl}(x, l, j), f))] \in \operatorname{derivable}(t) \end{array}}}
\vnbprop{der-union-decompose}{\vnbstmt{\begin{array}{@{}l@{}} \forall t, g, x, l, f. \\ \quad g \in \operatorname{ctx\text{-}set} \wedge \\ \quad \operatorname{enc\text{-}in}(x, \operatorname{enc\text{-}union}(l)) \in \mathit{clf} \wedge \\ \quad {[}0, f] \in \operatorname{rf\text{-}set}(4, x, l) \wedge \\ \quad \forall i \in \mathbb{N}. \\ \quad \quad ((1 \leq i) \wedge (i \leq \operatorname{length}(l))) \Rightarrow  \\ \quad \quad \quad {[}g, \operatorname{enc\text{-}eq}(\operatorname{nth}(i, l), \operatorname{nth}(i, l))] \in \operatorname{derivable}(t) \Rightarrow \\ \quad \quad {[}g, \operatorname{enc\text{-}iff}(\operatorname{enc\text{-}in}(x, \operatorname{enc\text{-}union}(l)), f)] \in \operatorname{derivable}(t) \end{array}}}
\vnbprop{der-intersection-decompose}{\vnbstmt{\begin{array}{@{}l@{}} \forall t, g, x, l, f. \\ \quad ((g \in \operatorname{ctx\text{-}set}) \wedge ((\operatorname{enc\text{-}in}(x, \operatorname{enc\text{-}inter}(l)) \in \mathit{clf}) \wedge ({[}0, f] \in \operatorname{rf\text{-}set}(3, x, l)))) \Rightarrow  \\ \quad \quad {[}g, \operatorname{enc\text{-}iff}(\operatorname{enc\text{-}in}(x, \operatorname{enc\text{-}inter}(l)), f)] \in \operatorname{derivable}(t) \end{array}}}
\vnbprop{j-udec-sub}{\vnbstmt{\forall t.\; (\operatorname{j\text{-}udec}(t) \subseteq \operatorname{eqv}(t))}}
\vnbprop{j-idec-sub}{\vnbstmt{\forall t.\; (\operatorname{j\text{-}idec}(t) \subseteq \operatorname{eqv}(t))}}
\vnbprop{adm-union-decompose-1}{\vnbstmt{\begin{array}{@{}l@{}} \forall t, s, c. \\ \quad (\operatorname{rel\text{-}union\text{-}decompose\text{-}1}(t, s, c) \wedge \operatorname{hyps\text{-}in}(s, \operatorname{derivable}(t))) \Rightarrow  \\ \quad \quad c \in \operatorname{derivable}(t) \end{array}}}
\vnbprop{adm-union-decompose-hyp-1}{\vnbstmt{\begin{array}{@{}l@{}} \forall t, s, c. \\ \quad (\operatorname{rel\text{-}union\text{-}decompose\text{-}hyp\text{-}1}(t, s, c) \wedge \operatorname{hyps\text{-}in}(s, \operatorname{derivable}(t))) \Rightarrow  \\ \quad \quad c \in \operatorname{derivable}(t) \end{array}}}
\vnbprop{adm-union-decompose-7}{\vnbstmt{\begin{array}{@{}l@{}} \forall t, s, c. \\ \quad (\operatorname{rel\text{-}union\text{-}decompose\text{-}7}(t, s, c) \wedge \operatorname{hyps\text{-}in}(s, \operatorname{derivable}(t))) \Rightarrow  \\ \quad \quad c \in \operatorname{derivable}(t) \end{array}}}
\vnbprop{adm-union-decompose-hyp-7}{\vnbstmt{\begin{array}{@{}l@{}} \forall t, s, c. \\ \quad (\operatorname{rel\text{-}union\text{-}decompose\text{-}hyp\text{-}7}(t, s, c) \wedge \operatorname{hyps\text{-}in}(s, \operatorname{derivable}(t))) \Rightarrow  \\ \quad \quad c \in \operatorname{derivable}(t) \end{array}}}
\vnbprop{adm-intersection-decompose-1}{\vnbstmt{\begin{array}{@{}l@{}} \forall t, s, c. \\ \quad (\operatorname{rel\text{-}intersection\text{-}decompose\text{-}1}(t, s, c) \wedge \operatorname{hyps\text{-}in}(s, \operatorname{derivable}(t))) \Rightarrow  \\ \quad \quad c \in \operatorname{derivable}(t) \end{array}}}
\vnbprop{adm-intersection-decompose-hyp-1}{\vnbstmt{\begin{array}{@{}l@{}} \forall t, s, c. \\ \quad (\operatorname{rel\text{-}intersection\text{-}decompose\text{-}hyp\text{-}1}(t, s, c) \wedge \operatorname{hyps\text{-}in}(s, \operatorname{derivable}(t))) \Rightarrow  \\ \quad \quad c \in \operatorname{derivable}(t) \end{array}}}
\vnbprop{adm-intersection-decompose-7}{\vnbstmt{\begin{array}{@{}l@{}} \forall t, s, c. \\ \quad (\operatorname{rel\text{-}intersection\text{-}decompose\text{-}7}(t, s, c) \wedge \operatorname{hyps\text{-}in}(s, \operatorname{derivable}(t))) \Rightarrow  \\ \quad \quad c \in \operatorname{derivable}(t) \end{array}}}
\vnbprop{adm-intersection-decompose-hyp-7}{\vnbstmt{\begin{array}{@{}l@{}} \forall t, s, c. \\ \quad (\operatorname{rel\text{-}intersection\text{-}decompose\text{-}hyp\text{-}7}(t, s, c) \wedge \operatorname{hyps\text{-}in}(s, \operatorname{derivable}(t))) \Rightarrow  \\ \quad \quad c \in \operatorname{derivable}(t) \end{array}}}
\vnbprop{der-union-binary}{\vnbstmt{\begin{array}{@{}l@{}} \forall t, g, x, a, b. \\ \quad g \in \operatorname{ctx\text{-}set} \wedge \\ \quad x \in \mathit{clt} \wedge \\ \quad a \in \mathit{clt} \wedge \\ \quad b \in \mathit{clt} \wedge \\ \quad {[}g, \operatorname{enc\text{-}eq}(a, a)] \in \operatorname{derivable}(t) \wedge \\ \quad {[}g, \operatorname{enc\text{-}eq}(b, b)] \in \operatorname{derivable}(t) \Rightarrow \\ \quad \quad {[}g, \operatorname{enc\text{-}iff}(\operatorname{enc\text{-}in}(x, \operatorname{enc\text{-}union}({[}a, b])), \operatorname{enc\text{-}or}(\operatorname{enc\text{-}in}(x, a), \operatorname{enc\text{-}in}(x, b)))] \in \operatorname{derivable}(t) \end{array}}}
\vnbprop{der-intersection-binary}{\vnbstmt{\begin{array}{@{}l@{}} \forall t, g, x, a, b. \\ \quad ((g \in \operatorname{ctx\text{-}set}) \wedge ((x \in \mathit{clt}) \wedge ((a \in \mathit{clt}) \wedge (b \in \mathit{clt})))) \Rightarrow  \\ \quad \quad {[}g, \operatorname{enc\text{-}iff}(\operatorname{enc\text{-}in}(x, \operatorname{enc\text{-}inter}({[}a, b])), \operatorname{enc\text{-}and}(\operatorname{enc\text{-}in}(x, a), \operatorname{enc\text{-}in}(x, b)))] \in \operatorname{derivable}(t) \end{array}}}

\vnbbookfile{structure-library/meta-nbg}{structure-library/meta-nbg}
\noindent{\small meta-nbg.scm -- K2, agent METACORE (2026-10-29): the kernel forms of the ten C+ rows of the K2 ledger (scratchpad/K2-LEDGER.md of \textasciitilde{}/prover-mf, section 1, rows 8, 10, 11, 13, 14, 15, 24, 26, 27, 51), TRANSCRIBED from their checkers of a320cb6f.  DEFINITIONS ONLY; the laws (the BRIDGES: each kernel form admissible over ...}\par
\par\medskip\noindent\textbf{Definition} (\texttt{rel-and-elim-k}).\;For all $\operatorname{mnkt\_}$, $\operatorname{mnkhs\_}$ and $\operatorname{mnkc\_}$, 
\begin{equation*}
\begin{array}{@{}l@{}} \operatorname{rel\text{-}and\text{-}elim\text{-}k}(\operatorname{mnkt\_}, \operatorname{mnkhs\_}, \operatorname{mnkc\_}) \\ \Leftrightarrow\; \operatorname{mnkc\_} \in \operatorname{seq\text{-}set} \wedge \\ \exists \operatorname{mnkg\_}, \operatorname{mnkf\_}, \operatorname{mnka\_}, \operatorname{mnkb\_}. \\ \quad \operatorname{mnkc\_} = {[}\operatorname{mnkg\_}, \operatorname{mnkf\_}] \wedge \\ \quad \operatorname{enc\text{-}and}(\operatorname{mnka\_}, \operatorname{mnkb\_}) \in \operatorname{mnkg\_} \wedge \\ \quad \operatorname{mnkhs\_} \\ \quad =\; {[}{[}\operatorname{ctx\text{-}add}(\operatorname{ctx\text{-}add}(\operatorname{difference}(\operatorname{mnkg\_}, \{\operatorname{enc\text{-}and}(\operatorname{mnka\_}, \operatorname{mnkb\_})\}), \operatorname{mnka\_}), \operatorname{mnkb\_}), \operatorname{mnkf\_}]] \end{array}
\end{equation*}

\par\medskip\noindent\textbf{Definition} (\texttt{rel-or-elim-k}).\;For all $\operatorname{mnkt\_}$, $\operatorname{mnkhs\_}$ and $\operatorname{mnkc\_}$, 
\begin{equation*}
\begin{array}{@{}l@{}} \operatorname{rel\text{-}or\text{-}elim\text{-}k}(\operatorname{mnkt\_}, \operatorname{mnkhs\_}, \operatorname{mnkc\_}) \\ \Leftrightarrow\; \operatorname{mnkc\_} \in \operatorname{seq\text{-}set} \wedge \\ \exists \operatorname{mnkg\_}, \operatorname{mnkf\_}, \operatorname{mnka\_}, \operatorname{mnkb\_}. \\ \quad \operatorname{mnkc\_} = {[}\operatorname{mnkg\_}, \operatorname{mnkf\_}] \wedge \\ \quad \operatorname{enc\text{-}or}(\operatorname{mnka\_}, \operatorname{mnkb\_}) \in \operatorname{mnkg\_} \wedge \\ \quad \operatorname{mnkhs\_} \\ \quad =\; {[}{[}\operatorname{ctx\text{-}add}(\operatorname{difference}(\operatorname{mnkg\_}, \{\operatorname{enc\text{-}or}(\operatorname{mnka\_}, \operatorname{mnkb\_})\}), \operatorname{mnka\_}), \operatorname{mnkf\_}], {[}\operatorname{ctx\text{-}add}(\operatorname{difference}(\operatorname{mnkg\_}, \{\operatorname{enc\text{-}or}(\operatorname{mnka\_}, \operatorname{mnkb\_})\}), \operatorname{mnkb\_}), \operatorname{mnkf\_}]] \end{array}
\end{equation*}

\par\medskip\noindent\textbf{Definition} (\texttt{rel-iff-elim-k}).\;For all $\operatorname{mnkt\_}$, $\operatorname{mnkhs\_}$ and $\operatorname{mnkc\_}$, 
\begin{equation*}
\begin{array}{@{}l@{}} \operatorname{rel\text{-}iff\text{-}elim\text{-}k}(\operatorname{mnkt\_}, \operatorname{mnkhs\_}, \operatorname{mnkc\_}) \\ \Leftrightarrow\; \operatorname{mnkc\_} \in \operatorname{seq\text{-}set} \wedge \\ \exists \operatorname{mnkg\_}, \operatorname{mnkf\_}, \operatorname{mnka\_}, \operatorname{mnkb\_}. \\ \quad \operatorname{mnkc\_} = {[}\operatorname{mnkg\_}, \operatorname{mnkf\_}] \wedge \\ \quad \operatorname{enc\text{-}iff}(\operatorname{mnka\_}, \operatorname{mnkb\_}) \in \operatorname{mnkg\_} \wedge \\ \quad \operatorname{mnkhs\_} \\ \quad =\; {[}{[}\operatorname{ctx\text{-}add}(\operatorname{ctx\text{-}add}(\operatorname{difference}(\operatorname{mnkg\_}, \{\operatorname{enc\text{-}iff}(\operatorname{mnka\_}, \operatorname{mnkb\_})\}), \operatorname{enc\text{-}imp}(\operatorname{mnka\_}, \operatorname{mnkb\_})), \operatorname{enc\text{-}imp}(\operatorname{mnkb\_}, \operatorname{mnka\_})), \operatorname{mnkf\_}]] \end{array}
\end{equation*}

\par\medskip\noindent\textbf{Definition} (\texttt{rel-forsome-elim-k}).\;For all $\operatorname{mnkt\_}$, $\operatorname{mnkhs\_}$ and $\operatorname{mnkc\_}$, 
\begin{equation*}
\begin{array}{@{}l@{}} \operatorname{rel\text{-}forsome\text{-}elim\text{-}k}(\operatorname{mnkt\_}, \operatorname{mnkhs\_}, \operatorname{mnkc\_}) \\ \Leftrightarrow\; \operatorname{mnkc\_} \in \operatorname{seq\text{-}set} \wedge \\ \exists \operatorname{mnkg\_}, \operatorname{mnkf\_}, \operatorname{mnkb\_}, \operatorname{mnky\_}. \\ \quad \operatorname{mnkc\_} = {[}\operatorname{mnkg\_}, \operatorname{mnkf\_}] \wedge \\ \quad \operatorname{enc\text{-}ex}(\operatorname{mnkb\_}) \in \operatorname{mnkg\_} \wedge \\ \quad \operatorname{mnky\_} \in \mathbb{N} \wedge \\ \quad \neg (\operatorname{mnky\_} \in (\operatorname{ctx\text{-}fv}(\operatorname{difference}(\operatorname{mnkg\_}, \{\operatorname{enc\text{-}ex}(\operatorname{mnkb\_})\})) \cup (\operatorname{expr\text{-}fv}(\operatorname{mnkf\_}) \cup \operatorname{expr\text{-}fv}(\operatorname{enc\text{-}ex}(\operatorname{mnkb\_}))))) \wedge \\ \quad \operatorname{mnkhs\_} \\ \quad =\; {[}{[}\operatorname{ctx\text{-}add}(\operatorname{difference}(\operatorname{mnkg\_}, \{\operatorname{enc\text{-}ex}(\operatorname{mnkb\_})\}), \operatorname{expr\text{-}open}(\operatorname{mnkb\_}, 0, \operatorname{efvar}(\operatorname{mnky\_}))), \operatorname{mnkf\_}]] \end{array}
\end{equation*}

\par\medskip\noindent\textbf{Definition} (\texttt{rel-forsome-intro-k}).\;For all $\operatorname{mnkt\_}$, $\operatorname{mnkhs\_}$ and $\operatorname{mnkc\_}$, 
\begin{equation*}
\begin{array}{@{}l@{}} \operatorname{rel\text{-}forsome\text{-}intro\text{-}k}(\operatorname{mnkt\_}, \operatorname{mnkhs\_}, \operatorname{mnkc\_}) \\ \Leftrightarrow\; \operatorname{mnkc\_} \in \operatorname{seq\text{-}set} \wedge \\ \exists \operatorname{mnkg\_}, \operatorname{mnkb\_}, \operatorname{mnktm\_}. \\ \quad \operatorname{mnkc\_} = {[}\operatorname{mnkg\_}, \operatorname{enc\text{-}ex}(\operatorname{mnkb\_})] \wedge \\ \quad \operatorname{mnktm\_} \in \mathit{clt} \wedge \\ \quad \operatorname{mnkhs\_} = {[}{[}\operatorname{mnkg\_}, \operatorname{expr\text{-}open}(\operatorname{mnkb\_}, 0, \operatorname{mnktm\_})]] \wedge \\ \quad \exists \operatorname{mnkf2\_} \in \operatorname{true\text{-}reach}(\operatorname{expr\text{-}open}(\operatorname{mnkb\_}, 0, \operatorname{mnktm\_})). \\ \quad \quad \operatorname{strict\text{-}atom}(\operatorname{mnkf2\_}, \operatorname{mnktm\_}) \vee \\ \quad \exists \operatorname{mnkh\_} \in \operatorname{mnkg\_}.\; (\operatorname{mnkh\_} \in \operatorname{cert\text{-}set}(\operatorname{mnkt\_}, \operatorname{mnkg\_}, \operatorname{mnktm\_})) \end{array}
\end{equation*}

\par\medskip\noindent\textbf{Definition} (\texttt{rel-forall-elim-k}).\;For all $\operatorname{mnkt\_}$, $\operatorname{mnkhs\_}$ and $\operatorname{mnkc\_}$, 
\begin{equation*}
\begin{array}{@{}l@{}} \operatorname{rel\text{-}forall\text{-}elim\text{-}k}(\operatorname{mnkt\_}, \operatorname{mnkhs\_}, \operatorname{mnkc\_}) \\ \Leftrightarrow\; \operatorname{mnkc\_} \in \operatorname{seq\text{-}set} \wedge \\ \exists \operatorname{mnkg\_}, \operatorname{mnkf\_}, \operatorname{mnkb\_}, \operatorname{mnktm\_}, \operatorname{mnkh\_}. \\ \quad \operatorname{mnkc\_} = {[}\operatorname{mnkg\_}, \operatorname{mnkf\_}] \wedge \\ \quad \operatorname{enc\text{-}all}(\operatorname{mnkb\_}) \in \operatorname{mnkg\_} \wedge \\ \quad \operatorname{mnktm\_} \in \mathit{clt} \wedge \\ \quad \operatorname{mnkh\_} \in \operatorname{mnkg\_} \wedge \\ \quad \operatorname{mnkh\_} \in \operatorname{cert\text{-}set}(\operatorname{mnkt\_}, \operatorname{mnkg\_}, \operatorname{mnktm\_}) \wedge \\ \quad \operatorname{mnkhs\_} = {[}{[}\operatorname{ctx\text{-}add}(\operatorname{mnkg\_}, \operatorname{expr\text{-}open}(\operatorname{mnkb\_}, 0, \operatorname{mnktm\_})), \operatorname{mnkf\_}]] \end{array}
\end{equation*}

\par\medskip\noindent\textbf{Definition} (\texttt{rel-reflexivity-k}).\;For all $\operatorname{mnkt\_}$, $\operatorname{mnkhs\_}$ and $\operatorname{mnkc\_}$, 
\begin{equation*}
\begin{array}{@{}l@{}} \operatorname{rel\text{-}reflexivity\text{-}k}(\operatorname{mnkt\_}, \operatorname{mnkhs\_}, \operatorname{mnkc\_}) \\ \Leftrightarrow\; \operatorname{mnkc\_} \in \operatorname{seq\text{-}set} \wedge \\ \exists \operatorname{mnkg\_}, \operatorname{mnktm\_}, \operatorname{mnkh\_}. \\ \quad \operatorname{mnkc\_} = {[}\operatorname{mnkg\_}, \operatorname{enc\text{-}eq}(\operatorname{mnktm\_}, \operatorname{mnktm\_})] \wedge \\ \quad \operatorname{mnktm\_} \in \mathit{clt} \wedge \\ \quad \operatorname{mnkh\_} \in \operatorname{mnkg\_} \wedge \\ \quad \operatorname{mnkh\_} \in \operatorname{cert\text{-}set}(\operatorname{mnkt\_}, \operatorname{mnkg\_}, \operatorname{mnktm\_}) \wedge \\ \quad \operatorname{mnkhs\_} = {[}] \end{array}
\end{equation*}

\par\medskip\noindent\textbf{Definition} (\texttt{rel-if-true-k}).\;For all $\operatorname{mnkt\_}$, $\operatorname{mnkhs\_}$ and $\operatorname{mnkc\_}$, 
\begin{equation*}
\begin{array}{@{}l@{}} \operatorname{rel\text{-}if\text{-}true\text{-}k}(\operatorname{mnkt\_}, \operatorname{mnkhs\_}, \operatorname{mnkc\_}) \\ \Leftrightarrow\; \operatorname{mnkc\_} \in \operatorname{seq\text{-}set} \wedge \\ \exists \operatorname{mnkg\_}, \operatorname{mnkf\_}, \operatorname{mnkp\_}, \operatorname{mnka\_}, \operatorname{mnkb\_}, \operatorname{mnkh\_}. \\ \quad \operatorname{mnkc\_} = {[}\operatorname{mnkg\_}, \operatorname{mnkf\_}] \wedge \\ \quad \operatorname{mnkh\_} \in \operatorname{mnkg\_} \wedge \\ \quad \operatorname{mnkh\_} \in \operatorname{cert\text{-}set}(\operatorname{mnkt\_}, \operatorname{mnkg\_}, \operatorname{mnka\_}) \wedge \\ \quad \operatorname{enc\text{-}if}(\operatorname{mnkp\_}, \operatorname{mnka\_}, \operatorname{mnkb\_}) \in \mathit{clt} \wedge \\ \quad \operatorname{mnka\_} \in \mathit{clt} \wedge \\ \quad \operatorname{mnkb\_} \in \mathit{clt} \wedge \\ \quad \operatorname{mnkhs\_} = {[}{[}\operatorname{mnkg\_}, \operatorname{mnkp\_}], {[}\operatorname{ctx\text{-}add}(\operatorname{mnkg\_}, \operatorname{enc\text{-}eq}(\operatorname{enc\text{-}if}(\operatorname{mnkp\_}, \operatorname{mnka\_}, \operatorname{mnkb\_}), \operatorname{mnka\_})), \operatorname{mnkf\_}]] \end{array}
\end{equation*}

\par\medskip\noindent\textbf{Definition} (\texttt{rel-if-false-k}).\;For all $\operatorname{mnkt\_}$, $\operatorname{mnkhs\_}$ and $\operatorname{mnkc\_}$, 
\begin{equation*}
\begin{array}{@{}l@{}} \operatorname{rel\text{-}if\text{-}false\text{-}k}(\operatorname{mnkt\_}, \operatorname{mnkhs\_}, \operatorname{mnkc\_}) \\ \Leftrightarrow\; \operatorname{mnkc\_} \in \operatorname{seq\text{-}set} \wedge \\ \exists \operatorname{mnkg\_}, \operatorname{mnkf\_}, \operatorname{mnkp\_}, \operatorname{mnka\_}, \operatorname{mnkb\_}, \operatorname{mnkh\_}. \\ \quad \operatorname{mnkc\_} = {[}\operatorname{mnkg\_}, \operatorname{mnkf\_}] \wedge \\ \quad \operatorname{mnkh\_} \in \operatorname{mnkg\_} \wedge \\ \quad \operatorname{mnkh\_} \in \operatorname{cert\text{-}set}(\operatorname{mnkt\_}, \operatorname{mnkg\_}, \operatorname{mnkb\_}) \wedge \\ \quad \operatorname{enc\text{-}if}(\operatorname{mnkp\_}, \operatorname{mnka\_}, \operatorname{mnkb\_}) \in \mathit{clt} \wedge \\ \quad \operatorname{mnka\_} \in \mathit{clt} \wedge \\ \quad \operatorname{mnkb\_} \in \mathit{clt} \wedge \\ \quad \operatorname{mnkhs\_} \\ \quad =\; {[}{[}\operatorname{mnkg\_}, \operatorname{enc\text{-}not}(\operatorname{mnkp\_})], {[}\operatorname{ctx\text{-}add}(\operatorname{mnkg\_}, \operatorname{enc\text{-}eq}(\operatorname{enc\text{-}if}(\operatorname{mnkp\_}, \operatorname{mnka\_}, \operatorname{mnkb\_}), \operatorname{mnkb\_})), \operatorname{mnkf\_}]] \end{array}
\end{equation*}

\par\medskip\noindent\textbf{Definition} (\texttt{rel-iota-in-elim-k}).\;For all $\operatorname{mnkt\_}$, $\operatorname{mnkhs\_}$ and $\operatorname{mnkc\_}$, 
\begin{equation*}
\begin{array}{@{}l@{}} \operatorname{rel\text{-}iota\text{-}in\text{-}elim\text{-}k}(\operatorname{mnkt\_}, \operatorname{mnkhs\_}, \operatorname{mnkc\_}) \\ \Leftrightarrow\; \operatorname{mnkc\_} \in \operatorname{seq\text{-}set} \wedge \\ \exists \operatorname{mnkg\_}, \operatorname{mnkf\_}, \operatorname{mnkp\_}. \\ \quad \operatorname{mnkc\_} = {[}\operatorname{mnkg\_}, \operatorname{mnkf\_}] \wedge \\ \quad \exists \operatorname{mnka\_}.\; (\operatorname{enc\text{-}in}(\operatorname{enc\text{-}iota}(\operatorname{mnkp\_}), \operatorname{mnka\_}) \in \operatorname{mnkg\_}) \vee \\ \quad \exists \operatorname{mnku\_}. \\ \quad \quad ((\operatorname{enc\text{-}eq}(\operatorname{enc\text{-}iota}(\operatorname{mnkp\_}), \operatorname{mnku\_}) \in \operatorname{mnkg\_}) \vee (\operatorname{enc\text{-}eq}(\operatorname{mnku\_}, \operatorname{enc\text{-}iota}(\operatorname{mnkp\_})) \in \operatorname{mnkg\_})) \wedge \\ \quad \operatorname{mnkhs\_} = {[}{[}\operatorname{ctx\text{-}add}(\operatorname{mnkg\_}, \operatorname{expr\text{-}open}(\operatorname{mnkp\_}, 0, \operatorname{enc\text{-}iota}(\operatorname{mnkp\_}))), \operatorname{mnkf\_}]] \end{array}
\end{equation*}

\par\medskip\noindent\textbf{Definition} (\texttt{def-rule}).\;For all $\operatorname{mnkt\_}$, $\operatorname{mnkhs\_}$ and $\operatorname{mnkc\_}$, 
\begin{equation*}
\begin{array}{@{}l@{}} \operatorname{def\text{-}rule}(\operatorname{mnkt\_}, \operatorname{mnkhs\_}, \operatorname{mnkc\_}) \\ \Leftrightarrow\; \operatorname{cr\text{-}assumption}(\operatorname{mnkt\_}, \operatorname{mnkhs\_}, \operatorname{mnkc\_}) \vee \\ \operatorname{cr\text{-}theorem\text{-}assumption}(\operatorname{mnkt\_}, \operatorname{mnkhs\_}, \operatorname{mnkc\_}) \vee \\ \operatorname{cr\text{-}reflexivity}(\operatorname{mnkt\_}, \operatorname{mnkhs\_}, \operatorname{mnkc\_}) \vee \\ \operatorname{def\text{-}strict\text{-}in\text{-}l}(\operatorname{mnkt\_}, \operatorname{mnkhs\_}, \operatorname{mnkc\_}) \vee \\ \operatorname{def\text{-}strict\text{-}in\text{-}r}(\operatorname{mnkt\_}, \operatorname{mnkhs\_}, \operatorname{mnkc\_}) \vee \\ \operatorname{def\text{-}strict\text{-}eq\text{-}l}(\operatorname{mnkt\_}, \operatorname{mnkhs\_}, \operatorname{mnkc\_}) \vee \\ \operatorname{def\text{-}strict\text{-}eq\text{-}r}(\operatorname{mnkt\_}, \operatorname{mnkhs\_}, \operatorname{mnkc\_}) \vee \\ \operatorname{def\text{-}strict\text{-}arg}(\operatorname{mnkt\_}, \operatorname{mnkhs\_}, \operatorname{mnkc\_}) \vee \\ \operatorname{def\text{-}strict\text{-}dom}(\operatorname{mnkt\_}, \operatorname{mnkhs\_}, \operatorname{mnkc\_}) \vee \\ \operatorname{def\text{-}total}(\operatorname{mnkt\_}, \operatorname{mnkhs\_}, \operatorname{mnkc\_}) \vee \\ \operatorname{def\text{-}binder}(\operatorname{mnkt\_}, \operatorname{mnkhs\_}, \operatorname{mnkc\_}) \vee \\ \operatorname{def\text{-}comp}(\operatorname{mnkt\_}, \operatorname{mnkhs\_}, \operatorname{mnkc\_}) \vee \\ \operatorname{def\text{-}list}(\operatorname{mnkt\_}, \operatorname{mnkhs\_}, \operatorname{mnkc\_}) \vee \\ \operatorname{def\text{-}strict\text{-}hyp}(\operatorname{mnkt\_}, \operatorname{mnkhs\_}, \operatorname{mnkc\_}) \vee \\ \operatorname{def\text{-}nth\text{-}list}(\operatorname{mnkt\_}, \operatorname{mnkhs\_}, \operatorname{mnkc\_}) \vee \\ \operatorname{def\text{-}strict\text{-}pred}(\operatorname{mnkt\_}, \operatorname{mnkhs\_}, \operatorname{mnkc\_}) \vee \\ \operatorname{def\text{-}struct\text{-}hyp}(\operatorname{mnkt\_}, \operatorname{mnkhs\_}, \operatorname{mnkc\_}) \vee \\ \operatorname{def\text{-}pred\text{-}unfold}(\operatorname{mnkt\_}, \operatorname{mnkhs\_}, \operatorname{mnkc\_}) \vee \\ \operatorname{def\text{-}pred\text{-}hyp}(\operatorname{mnkt\_}, \operatorname{mnkhs\_}, \operatorname{mnkc\_}) \vee \\ \operatorname{def\text{-}fun\text{-}inst}(\operatorname{mnkt\_}, \operatorname{mnkhs\_}, \operatorname{mnkc\_}) \vee \\ \operatorname{def\text{-}fun\text{-}strict}(\operatorname{mnkt\_}, \operatorname{mnkhs\_}, \operatorname{mnkc\_}) \vee \\ \operatorname{def\text{-}fun\text{-}nth}(\operatorname{mnkt\_}, \operatorname{mnkhs\_}, \operatorname{mnkc\_}) \vee \\ \operatorname{def\text{-}fun\text{-}sep}(\operatorname{mnkt\_}, \operatorname{mnkhs\_}, \operatorname{mnkc\_}) \vee \\ \operatorname{def\text{-}pred\text{-}unfold\text{-}n}(\operatorname{mnkt\_}, \operatorname{mnkhs\_}, \operatorname{mnkc\_}) \vee \\ \operatorname{def\text{-}pred\text{-}hyp\text{-}n}(\operatorname{mnkt\_}, \operatorname{mnkhs\_}, \operatorname{mnkc\_}) \vee \\ \operatorname{def\text{-}cert\text{-}hyp}(\operatorname{mnkt\_}, \operatorname{mnkhs\_}, \operatorname{mnkc\_}) \vee \\ (\operatorname{mnkt\_} \subseteq \mathit{clf}) \wedge ((\operatorname{ax\text{-}fun\text{-}app} \in \operatorname{mnkt\_}) \wedge \operatorname{def\text{-}app}(\operatorname{mnkt\_}, \operatorname{mnkhs\_}, \operatorname{mnkc\_})) \vee \\ (\operatorname{mnkt\_} \subseteq \mathit{clf}) \wedge ((\operatorname{ax\text{-}len\text{-}nn} \in \operatorname{mnkt\_}) \wedge \operatorname{def\text{-}len\text{-}tuples}(\operatorname{mnkt\_}, \operatorname{mnkhs\_}, \operatorname{mnkc\_})) \end{array}
\end{equation*}

\par\medskip\noindent\textbf{Definition} (\texttt{def-derivable}).\;$\operatorname{def\text{-}derivable}(\operatorname{mnkt\_}) :=$
\begin{equation*}
\{\operatorname{mnkq\_} \in \operatorname{seq\text{-}set} : \forall \operatorname{mnks\_} \in \operatorname{power}(\operatorname{seq\text{-}set}).\; (\operatorname{def\text{-}closed}(\operatorname{mnkt\_}, \operatorname{mnks\_}) \Rightarrow (\operatorname{mnkq\_} \in \operatorname{mnks\_}))\}
\end{equation*}

\par\medskip\noindent\textbf{Definition} (\texttt{def-closed}).\;For all $\operatorname{mnkt\_}$ and $\operatorname{mnks\_}$, 
\begin{equation*}
\begin{array}{@{}l@{}} \operatorname{def\text{-}closed}(\operatorname{mnkt\_}, \operatorname{mnks\_}) \\ \Leftrightarrow\; \forall \operatorname{mnkhs\_}, \operatorname{mnkc\_}. \\ \quad \operatorname{def\text{-}rule}(\operatorname{mnkt\_}, \operatorname{mnkhs\_}, \operatorname{mnkc\_}) \Rightarrow  \\ \quad \operatorname{hyps\text{-}in}(\operatorname{mnkhs\_}, \operatorname{mnks\_}) \Rightarrow ((\operatorname{mnkc\_} \in \operatorname{seq\text{-}set}) \Rightarrow (\operatorname{mnkc\_} \in \operatorname{mnks\_})) \end{array}
\end{equation*}

\par\medskip\noindent\textbf{Definition} (\texttt{rel-forall-elim-kd}).\;For all $\operatorname{mnkt\_}$, $\operatorname{mnkhs\_}$ and $\operatorname{mnkc\_}$, 
\begin{equation*}
\begin{array}{@{}l@{}} \operatorname{rel\text{-}forall\text{-}elim\text{-}kd}(\operatorname{mnkt\_}, \operatorname{mnkhs\_}, \operatorname{mnkc\_}) \\ \Leftrightarrow\; \operatorname{mnkc\_} \in \operatorname{seq\text{-}set} \wedge \\ \exists \operatorname{mnkg\_}, \operatorname{mnkf\_}, \operatorname{mnkb\_}, \operatorname{mnktm\_}. \\ \quad \operatorname{mnkc\_} = {[}\operatorname{mnkg\_}, \operatorname{mnkf\_}] \wedge \\ \quad \operatorname{enc\text{-}all}(\operatorname{mnkb\_}) \in \operatorname{mnkg\_} \wedge \\ \quad \operatorname{mnktm\_} \in \mathit{clt} \wedge \\ \quad {[}\operatorname{mnkg\_}, \operatorname{enc\text{-}eq}(\operatorname{mnktm\_}, \operatorname{mnktm\_})] \in \operatorname{def\text{-}derivable}(\operatorname{mnkt\_}) \wedge \\ \quad \operatorname{mnkhs\_} = {[}{[}\operatorname{ctx\text{-}add}(\operatorname{mnkg\_}, \operatorname{expr\text{-}open}(\operatorname{mnkb\_}, 0, \operatorname{mnktm\_})), \operatorname{mnkf\_}]] \end{array}
\end{equation*}

\par\medskip\noindent\textbf{Definition} (\texttt{rel-forsome-intro-kd}).\;For all $\operatorname{mnkt\_}$, $\operatorname{mnkhs\_}$ and $\operatorname{mnkc\_}$, 
\begin{equation*}
\begin{array}{@{}l@{}} \operatorname{rel\text{-}forsome\text{-}intro\text{-}kd}(\operatorname{mnkt\_}, \operatorname{mnkhs\_}, \operatorname{mnkc\_}) \\ \Leftrightarrow\; \operatorname{mnkc\_} \in \operatorname{seq\text{-}set} \wedge \\ \exists \operatorname{mnkg\_}, \operatorname{mnkb\_}, \operatorname{mnktm\_}. \\ \quad \operatorname{mnkc\_} = {[}\operatorname{mnkg\_}, \operatorname{enc\text{-}ex}(\operatorname{mnkb\_})] \wedge \\ \quad \operatorname{mnktm\_} \in \mathit{clt} \wedge \\ \quad {[}\operatorname{mnkg\_}, \operatorname{enc\text{-}eq}(\operatorname{mnktm\_}, \operatorname{mnktm\_})] \in \operatorname{def\text{-}derivable}(\operatorname{mnkt\_}) \wedge \\ \quad \operatorname{mnkhs\_} = {[}{[}\operatorname{mnkg\_}, \operatorname{expr\text{-}open}(\operatorname{mnkb\_}, 0, \operatorname{mnktm\_})]] \end{array}
\end{equation*}

\par\medskip\noindent\textbf{Definition} (\texttt{rel-reflexivity-kd}).\;For all $\operatorname{mnkt\_}$, $\operatorname{mnkhs\_}$ and $\operatorname{mnkc\_}$, 
\begin{equation*}
\begin{array}{@{}l@{}} \operatorname{rel\text{-}reflexivity\text{-}kd}(\operatorname{mnkt\_}, \operatorname{mnkhs\_}, \operatorname{mnkc\_}) \\ \Leftrightarrow\; \operatorname{mnkc\_} \in \operatorname{seq\text{-}set} \wedge \\ \exists \operatorname{mnkg\_}, \operatorname{mnktm\_}. \\ \quad \operatorname{mnkc\_} = {[}\operatorname{mnkg\_}, \operatorname{enc\text{-}eq}(\operatorname{mnktm\_}, \operatorname{mnktm\_})] \wedge \\ \quad {[}\operatorname{mnkg\_}, \operatorname{enc\text{-}eq}(\operatorname{mnktm\_}, \operatorname{mnktm\_})] \in \operatorname{def\text{-}derivable}(\operatorname{mnkt\_}) \wedge \\ \quad \operatorname{mnkhs\_} = {[}] \end{array}
\end{equation*}

\par\medskip\noindent\textbf{Definition} (\texttt{rel-reflexivity-kt}).\;For all $\operatorname{mnkt\_}$, $\operatorname{mnkhs\_}$ and $\operatorname{mnkc\_}$, 
\begin{equation*}
\begin{array}{@{}l@{}} \operatorname{rel\text{-}reflexivity\text{-}kt}(\operatorname{mnkt\_}, \operatorname{mnkhs\_}, \operatorname{mnkc\_}) \\ \Leftrightarrow\; ((\operatorname{mnkc\_} \in \operatorname{seq\text{-}set}) \wedge \exists \operatorname{mnkg\_}.\; ((\operatorname{mnkc\_} = {[}\operatorname{mnkg\_}, \operatorname{enc\text{-}truth}]) \wedge (\operatorname{mnkhs\_} = {[}]))) \end{array}
\end{equation*}

\par\medskip\noindent\textbf{Definition} (\texttt{rel-if-true-kd}).\;For all $\operatorname{mnkt\_}$, $\operatorname{mnkhs\_}$ and $\operatorname{mnkc\_}$, 
\begin{equation*}
\begin{array}{@{}l@{}} \operatorname{rel\text{-}if\text{-}true\text{-}kd}(\operatorname{mnkt\_}, \operatorname{mnkhs\_}, \operatorname{mnkc\_}) \\ \Leftrightarrow\; \operatorname{mnkc\_} \in \operatorname{seq\text{-}set} \wedge \\ \exists \operatorname{mnkg\_}, \operatorname{mnkf\_}, \operatorname{mnkp\_}, \operatorname{mnka\_}, \operatorname{mnkb\_}. \\ \quad \operatorname{mnkc\_} = {[}\operatorname{mnkg\_}, \operatorname{mnkf\_}] \wedge \\ \quad {[}\operatorname{mnkg\_}, \operatorname{enc\text{-}eq}(\operatorname{mnka\_}, \operatorname{mnka\_})] \in \operatorname{def\text{-}derivable}(\operatorname{mnkt\_}) \wedge \\ \quad \operatorname{enc\text{-}if}(\operatorname{mnkp\_}, \operatorname{mnka\_}, \operatorname{mnkb\_}) \in \mathit{clt} \wedge \\ \quad \operatorname{mnka\_} \in \mathit{clt} \wedge \\ \quad \operatorname{mnkb\_} \in \mathit{clt} \wedge \\ \quad \operatorname{mnkhs\_} = {[}{[}\operatorname{mnkg\_}, \operatorname{mnkp\_}], {[}\operatorname{ctx\text{-}add}(\operatorname{mnkg\_}, \operatorname{enc\text{-}eq}(\operatorname{enc\text{-}if}(\operatorname{mnkp\_}, \operatorname{mnka\_}, \operatorname{mnkb\_}), \operatorname{mnka\_})), \operatorname{mnkf\_}]] \end{array}
\end{equation*}

\par\medskip\noindent\textbf{Definition} (\texttt{rel-if-false-kd}).\;For all $\operatorname{mnkt\_}$, $\operatorname{mnkhs\_}$ and $\operatorname{mnkc\_}$, 
\begin{equation*}
\begin{array}{@{}l@{}} \operatorname{rel\text{-}if\text{-}false\text{-}kd}(\operatorname{mnkt\_}, \operatorname{mnkhs\_}, \operatorname{mnkc\_}) \\ \Leftrightarrow\; \operatorname{mnkc\_} \in \operatorname{seq\text{-}set} \wedge \\ \exists \operatorname{mnkg\_}, \operatorname{mnkf\_}, \operatorname{mnkp\_}, \operatorname{mnka\_}, \operatorname{mnkb\_}. \\ \quad \operatorname{mnkc\_} = {[}\operatorname{mnkg\_}, \operatorname{mnkf\_}] \wedge \\ \quad {[}\operatorname{mnkg\_}, \operatorname{enc\text{-}eq}(\operatorname{mnkb\_}, \operatorname{mnkb\_})] \in \operatorname{def\text{-}derivable}(\operatorname{mnkt\_}) \wedge \\ \quad \operatorname{enc\text{-}if}(\operatorname{mnkp\_}, \operatorname{mnka\_}, \operatorname{mnkb\_}) \in \mathit{clt} \wedge \\ \quad \operatorname{mnka\_} \in \mathit{clt} \wedge \\ \quad \operatorname{mnkb\_} \in \mathit{clt} \wedge \\ \quad \operatorname{mnkhs\_} \\ \quad =\; {[}{[}\operatorname{mnkg\_}, \operatorname{enc\text{-}not}(\operatorname{mnkp\_})], {[}\operatorname{ctx\text{-}add}(\operatorname{mnkg\_}, \operatorname{enc\text{-}eq}(\operatorname{enc\text{-}if}(\operatorname{mnkp\_}, \operatorname{mnka\_}, \operatorname{mnkb\_}), \operatorname{mnkb\_})), \operatorname{mnkf\_}]] \end{array}
\end{equation*}

\par\medskip\noindent\textbf{Definition} (\texttt{rel-adds-nothing-k}).\;For all $\operatorname{mnkt\_}$, $\operatorname{mnkhs\_}$ and $\operatorname{mnkc\_}$, 
\begin{equation*}
\begin{array}{@{}l@{}} \operatorname{rel\text{-}adds\text{-}nothing\text{-}k}(\operatorname{mnkt\_}, \operatorname{mnkhs\_}, \operatorname{mnkc\_}) \\ \Leftrightarrow\; \operatorname{mnkc\_} \in \operatorname{seq\text{-}set} \wedge \\ \exists \operatorname{mnkg\_}, \operatorname{mnkf\_}, \operatorname{mnkd\_}, \operatorname{mnki\_}. \\ \quad \operatorname{mnkc\_} = {[}\operatorname{mnkg\_}, \operatorname{mnkf\_}] \wedge \\ \quad \operatorname{mnkd\_} \subseteq \operatorname{mnkg\_} \wedge \\ \quad \operatorname{mnki\_} \in \mathbb{N} \wedge \\ \quad 1 \leq \operatorname{mnki\_} \wedge \\ \quad \operatorname{mnki\_} \leq \operatorname{length}(\operatorname{mnkhs\_}) \wedge \\ \quad \operatorname{nth}(\operatorname{mnki\_}, \operatorname{mnkhs\_}) = {[}\operatorname{mnkd\_}, \operatorname{mnkf\_}] \end{array}
\end{equation*}

\par\medskip\noindent\textbf{Definition} (\texttt{rewrs}).\;$\operatorname{rewrs}(\operatorname{mpoj\_}) :=$
\begin{equation*}
\{\operatorname{mpox\_} \in \operatorname{rw\text{-}univ} : \forall \operatorname{mpos\_} \in \operatorname{power}(\operatorname{rw\text{-}univ}).\; (\operatorname{rws\text{-}closed}(\operatorname{mpoj\_}, \operatorname{mpos\_}) \Rightarrow (\operatorname{mpox\_} \in \operatorname{mpos\_}))\}
\end{equation*}

\par\medskip\noindent\textbf{Definition} (\texttt{rws-rule}).\;For all $\operatorname{mpoj\_}$, $\operatorname{mpos\_}$ and $\operatorname{mpox\_}$, 
\begin{equation*}
\begin{array}{@{}l@{}} \operatorname{rws\text{-}rule}(\operatorname{mpoj\_}, \operatorname{mpos\_}, \operatorname{mpox\_}) \\ \Leftrightarrow\; \operatorname{mpox\_} \in \operatorname{rewr}(\operatorname{mpoj\_}, 7) \vee \\ \exists \operatorname{mpog\_} \in \operatorname{ctx\text{-}set}, \operatorname{mpoe\_} \in \mathit{clf}. \\ \quad \operatorname{mpox\_} = {[}2, \operatorname{mpog\_}, \operatorname{mpoe\_}, \operatorname{mpoe\_}] \vee \\ \exists \operatorname{mpog\_} \in \operatorname{ctx\text{-}set}, \operatorname{mpoa\_}, \operatorname{mpob\_}, \operatorname{mpoc\_} \in \operatorname{expr\text{-}set}. \\ \quad \operatorname{mpox\_} = {[}2, \operatorname{mpog\_}, \operatorname{mpoa\_}, \operatorname{mpoc\_}] \wedge \\ \quad {[}2, \operatorname{mpog\_}, \operatorname{mpoa\_}, \operatorname{mpob\_}] \in \operatorname{mpos\_} \wedge \\ \quad {[}2, \operatorname{mpog\_}, \operatorname{mpob\_}, \operatorname{mpoc\_}] \in \operatorname{mpos\_} \end{array}
\end{equation*}

\par\medskip\noindent\textbf{Definition} (\texttt{rws-closed}).\;For all $\operatorname{mpoj\_}$ and $\operatorname{mpos\_}$, 
\begin{equation*}
\begin{array}{@{}l@{}} \operatorname{rws\text{-}closed}(\operatorname{mpoj\_}, \operatorname{mpos\_}) \\ \Leftrightarrow\; \forall \operatorname{mpox\_} \in \operatorname{rw\text{-}univ}.\; (\operatorname{rws\text{-}rule}(\operatorname{mpoj\_}, \operatorname{mpos\_}, \operatorname{mpox\_}) \Rightarrow (\operatorname{mpox\_} \in \operatorname{mpos\_})) \end{array}
\end{equation*}

\par\medskip\noindent\textbf{Definition} (\texttt{rel-rewrite-s}).\;For all $\operatorname{mpoj\_}$, $\operatorname{mpohs\_}$ and $\operatorname{mpocc\_}$, 
\begin{equation*}
\begin{array}{@{}l@{}} \operatorname{rel\text{-}rewrite\text{-}s}(\operatorname{mpoj\_}, \operatorname{mpohs\_}, \operatorname{mpocc\_}) \\ \Leftrightarrow\; \operatorname{mpocc\_} \in \operatorname{seq\text{-}set} \wedge \\ \exists \operatorname{mpog\_}, \operatorname{mpoa\_}, \operatorname{mpob\_}. \\ \quad ((\operatorname{mpocc\_} = {[}\operatorname{mpog\_}, \operatorname{mpoa\_}]) \wedge ((\operatorname{mpohs\_} = {[}{[}\operatorname{mpog\_}, \operatorname{mpob\_}]]) \wedge ({[}2, \operatorname{mpog\_}, \operatorname{mpoa\_}, \operatorname{mpob\_}] \in \operatorname{rewrs}(\operatorname{mpoj\_})))) \end{array}
\end{equation*}

\par\medskip\noindent\textbf{Definition} (\texttt{rel-rewrite-hyp-s}).\;For all $\operatorname{mpoj\_}$, $\operatorname{mpohs\_}$ and $\operatorname{mpocc\_}$, 
\begin{equation*}
\begin{array}{@{}l@{}} \operatorname{rel\text{-}rewrite\text{-}hyp\text{-}s}(\operatorname{mpoj\_}, \operatorname{mpohs\_}, \operatorname{mpocc\_}) \\ \Leftrightarrow\; \operatorname{mpocc\_} \in \operatorname{seq\text{-}set} \wedge \\ \exists \operatorname{mpog\_}, \operatorname{mpof\_}, \operatorname{mpoa\_}, \operatorname{mpob\_}. \\ \quad \operatorname{mpocc\_} = {[}\operatorname{mpog\_}, \operatorname{mpof\_}] \wedge \\ \quad \operatorname{mpoa\_} \in \operatorname{mpog\_} \wedge \\ \quad \operatorname{mpohs\_} = {[}{[}\operatorname{ctx\text{-}add}(\operatorname{difference}(\operatorname{mpog\_}, \{\operatorname{mpoa\_}\}), \operatorname{mpob\_}), \operatorname{mpof\_}]] \wedge \\ \quad {[}2, \operatorname{mpog\_}, \operatorname{mpoa\_}, \operatorname{mpob\_}] \in \operatorname{rewrs}(\operatorname{mpoj\_}) \end{array}
\end{equation*}

\par\medskip\noindent\textbf{Definition} (\texttt{rel-nth-reduce-s}).\;For all $\operatorname{mpot\_}$, $\operatorname{mpohs\_}$ and $\operatorname{mpocc\_}$, 
\begin{equation*}
\operatorname{rel\text{-}nth\text{-}reduce\text{-}s}(\operatorname{mpot\_}, \operatorname{mpohs\_}, \operatorname{mpocc\_}) \Leftrightarrow \operatorname{rel\text{-}rewrite\text{-}s}(\operatorname{j\text{-}nth}(\operatorname{mpot\_}), \operatorname{mpohs\_}, \operatorname{mpocc\_})
\end{equation*}

\par\medskip\noindent\textbf{Definition} (\texttt{rel-length-reduce-s}).\;For all $\operatorname{mpot\_}$, $\operatorname{mpohs\_}$ and $\operatorname{mpocc\_}$, 
\begin{equation*}
\operatorname{rel\text{-}length\text{-}reduce\text{-}s}(\operatorname{mpot\_}, \operatorname{mpohs\_}, \operatorname{mpocc\_}) \Leftrightarrow \operatorname{rel\text{-}rewrite\text{-}s}(\operatorname{j\text{-}len}(\operatorname{mpot\_}), \operatorname{mpohs\_}, \operatorname{mpocc\_})
\end{equation*}

\par\medskip\noindent\textbf{Definition} (\texttt{rel-functoid-beta-s}).\;For all $\operatorname{mpot\_}$, $\operatorname{mpohs\_}$ and $\operatorname{mpocc\_}$, 
\begin{equation*}
\operatorname{rel\text{-}functoid\text{-}beta\text{-}s}(\operatorname{mpot\_}, \operatorname{mpohs\_}, \operatorname{mpocc\_}) \Leftrightarrow \operatorname{rel\text{-}rewrite\text{-}s}(\operatorname{j\text{-}fb}(\operatorname{mpot\_}), \operatorname{mpohs\_}, \operatorname{mpocc\_})
\end{equation*}

\par\medskip\noindent\textbf{Definition} (\texttt{rel-lambda-beta-s}).\;For all $\operatorname{mpot\_}$, $\operatorname{mpohs\_}$ and $\operatorname{mpocc\_}$, 
\begin{equation*}
\operatorname{rel\text{-}lambda\text{-}beta\text{-}s}(\operatorname{mpot\_}, \operatorname{mpohs\_}, \operatorname{mpocc\_}) \Leftrightarrow \operatorname{rel\text{-}rewrite\text{-}s}(\operatorname{j\text{-}lb}(\operatorname{mpot\_}), \operatorname{mpohs\_}, \operatorname{mpocc\_})
\end{equation*}

\par\medskip\noindent\textbf{Definition} (\texttt{rel-lambda-beta-hyp-s}).\;For all $\operatorname{mpot\_}$, $\operatorname{mpohs\_}$ and $\operatorname{mpocc\_}$, 
\begin{equation*}
\operatorname{rel\text{-}lambda\text{-}beta\text{-}hyp\text{-}s}(\operatorname{mpot\_}, \operatorname{mpohs\_}, \operatorname{mpocc\_}) \Leftrightarrow \operatorname{rel\text{-}rewrite\text{-}hyp\text{-}s}(\operatorname{j\text{-}lb}(\operatorname{mpot\_}), \operatorname{mpohs\_}, \operatorname{mpocc\_})
\end{equation*}

\par\medskip\noindent\textbf{Definition} (\texttt{rel-macete-s}).\;For all $\operatorname{mpot\_}$, $\operatorname{mpohs\_}$ and $\operatorname{mpocc\_}$, 
\begin{equation*}
\operatorname{rel\text{-}macete\text{-}s}(\operatorname{mpot\_}, \operatorname{mpohs\_}, \operatorname{mpocc\_}) \Leftrightarrow \operatorname{rel\text{-}rewrite\text{-}s}(\operatorname{j\text{-}mac}(\operatorname{mpot\_}), \operatorname{mpohs\_}, \operatorname{mpocc\_})
\end{equation*}

\par\medskip\noindent\textbf{Definition} (\texttt{rel-macete-hyp-s}).\;For all $\operatorname{mpot\_}$, $\operatorname{mpohs\_}$ and $\operatorname{mpocc\_}$, 
\begin{equation*}
\operatorname{rel\text{-}macete\text{-}hyp\text{-}s}(\operatorname{mpot\_}, \operatorname{mpohs\_}, \operatorname{mpocc\_}) \Leftrightarrow \operatorname{rel\text{-}rewrite\text{-}hyp\text{-}s}(\operatorname{j\text{-}mac}(\operatorname{mpot\_}), \operatorname{mpohs\_}, \operatorname{mpocc\_})
\end{equation*}

\par\medskip\noindent\textbf{Definition} (\texttt{rel-macete-n-s}).\;For all $\operatorname{mpot\_}$, $\operatorname{mpohs\_}$ and $\operatorname{mpocc\_}$, 
\begin{equation*}
\operatorname{rel\text{-}macete\text{-}n\text{-}s}(\operatorname{mpot\_}, \operatorname{mpohs\_}, \operatorname{mpocc\_}) \Leftrightarrow \operatorname{rel\text{-}rewrite\text{-}s}(\operatorname{j\text{-}macn}(\operatorname{mpot\_}), \operatorname{mpohs\_}, \operatorname{mpocc\_})
\end{equation*}

\par\medskip\noindent\textbf{Definition} (\texttt{rel-macete-n-hyp-s}).\;For all $\operatorname{mpot\_}$, $\operatorname{mpohs\_}$ and $\operatorname{mpocc\_}$, 
\begin{equation*}
\operatorname{rel\text{-}macete\text{-}n\text{-}hyp\text{-}s}(\operatorname{mpot\_}, \operatorname{mpohs\_}, \operatorname{mpocc\_}) \Leftrightarrow \operatorname{rel\text{-}rewrite\text{-}hyp\text{-}s}(\operatorname{j\text{-}macn}(\operatorname{mpot\_}), \operatorname{mpohs\_}, \operatorname{mpocc\_})
\end{equation*}

\par\medskip\noindent\textbf{Definition} (\texttt{rel-union-decompose-s}).\;For all $\operatorname{mpot\_}$, $\operatorname{mpohs\_}$ and $\operatorname{mpocc\_}$, 
\begin{equation*}
\operatorname{rel\text{-}union\text{-}decompose\text{-}s}(\operatorname{mpot\_}, \operatorname{mpohs\_}, \operatorname{mpocc\_}) \Leftrightarrow \operatorname{rel\text{-}rewrite\text{-}s}(\operatorname{j\text{-}udec}(\operatorname{mpot\_}), \operatorname{mpohs\_}, \operatorname{mpocc\_})
\end{equation*}

\par\medskip\noindent\textbf{Definition} (\texttt{rel-union-decompose-hyp-s}).\;For all $\operatorname{mpot\_}$, $\operatorname{mpohs\_}$ and $\operatorname{mpocc\_}$, 
\begin{equation*}
\begin{array}{@{}l@{}} \operatorname{rel\text{-}union\text{-}decompose\text{-}hyp\text{-}s}(\operatorname{mpot\_}, \operatorname{mpohs\_}, \operatorname{mpocc\_}) \\ \Leftrightarrow\; \operatorname{rel\text{-}rewrite\text{-}hyp\text{-}s}(\operatorname{j\text{-}udec}(\operatorname{mpot\_}), \operatorname{mpohs\_}, \operatorname{mpocc\_}) \end{array}
\end{equation*}

\par\medskip\noindent\textbf{Definition} (\texttt{rel-intersection-decompose-s}).\;For all $\operatorname{mpot\_}$, $\operatorname{mpohs\_}$ and $\operatorname{mpocc\_}$, 
\begin{equation*}
\operatorname{rel\text{-}intersection\text{-}decompose\text{-}s}(\operatorname{mpot\_}, \operatorname{mpohs\_}, \operatorname{mpocc\_}) \Leftrightarrow \operatorname{rel\text{-}rewrite\text{-}s}(\operatorname{j\text{-}idec}(\operatorname{mpot\_}), \operatorname{mpohs\_}, \operatorname{mpocc\_})
\end{equation*}

\par\medskip\noindent\textbf{Definition} (\texttt{rel-intersection-decompose-hyp-s}).\;For all $\operatorname{mpot\_}$, $\operatorname{mpohs\_}$ and $\operatorname{mpocc\_}$, 
\begin{equation*}
\begin{array}{@{}l@{}} \operatorname{rel\text{-}intersection\text{-}decompose\text{-}hyp\text{-}s}(\operatorname{mpot\_}, \operatorname{mpohs\_}, \operatorname{mpocc\_}) \\ \Leftrightarrow\; \operatorname{rel\text{-}rewrite\text{-}hyp\text{-}s}(\operatorname{j\text{-}idec}(\operatorname{mpot\_}), \operatorname{mpohs\_}, \operatorname{mpocc\_}) \end{array}
\end{equation*}

\par\medskip\noindent\textbf{Definition} (\texttt{rel-cartesian-decompose-s}).\;For all $\operatorname{mpot\_}$, $\operatorname{mpohs\_}$ and $\operatorname{mpocc\_}$, 
\begin{equation*}
\operatorname{rel\text{-}cartesian\text{-}decompose\text{-}s}(\operatorname{mpot\_}, \operatorname{mpohs\_}, \operatorname{mpocc\_}) \Leftrightarrow \operatorname{rel\text{-}rewrite\text{-}s}(\operatorname{j\text{-}cart}(\operatorname{mpot\_}), \operatorname{mpohs\_}, \operatorname{mpocc\_})
\end{equation*}

\par\medskip\noindent\textbf{Definition} (\texttt{rel-tuple-eq-decompose-s}).\;For all $\operatorname{mpot\_}$, $\operatorname{mpohs\_}$ and $\operatorname{mpocc\_}$, 
\begin{equation*}
\operatorname{rel\text{-}tuple\text{-}eq\text{-}decompose\text{-}s}(\operatorname{mpot\_}, \operatorname{mpohs\_}, \operatorname{mpocc\_}) \Leftrightarrow \operatorname{rel\text{-}rewrite\text{-}s}(\operatorname{j\text{-}teq}(\operatorname{mpot\_}), \operatorname{mpohs\_}, \operatorname{mpocc\_})
\end{equation*}

\par\medskip\noindent\textbf{Definition} (\texttt{rel-tuples-intro-0k}).\;For all $\operatorname{mnkt\_}$, $\operatorname{mnkhs\_}$ and $\operatorname{mnkc\_}$, 
\begin{equation*}
\begin{array}{@{}l@{}} \operatorname{rel\text{-}tuples\text{-}intro\text{-}0k}(\operatorname{mnkt\_}, \operatorname{mnkhs\_}, \operatorname{mnkc\_}) \\ \Leftrightarrow\; \operatorname{mnkc\_} \in \operatorname{seq\text{-}set} \wedge \\ \exists \operatorname{mnkg\_}, \operatorname{mnkl\_}, \operatorname{mnka\_}, \operatorname{mnkh\_}. \\ \quad \operatorname{mnkc\_} = {[}\operatorname{mnkg\_}, \operatorname{enc\text{-}in}(\operatorname{enc\text{-}list}(\operatorname{mnkl\_}), \operatorname{enc\text{-}tuples}(\operatorname{mnka\_}))] \wedge \\ \quad \operatorname{length}(\operatorname{mnkl\_}) = 0 \wedge \\ \quad \operatorname{mnka\_} \in \mathit{clt} \wedge \\ \quad \operatorname{mnkh\_} \in \operatorname{mnkg\_} \wedge \\ \quad \operatorname{mnkh\_} \in \operatorname{cert\text{-}set}(\operatorname{mnkt\_}, \operatorname{mnkg\_}, \operatorname{mnka\_}) \wedge \\ \quad \operatorname{mnkhs\_} = {[}] \end{array}
\end{equation*}

\par\medskip\noindent\textbf{Definition} (\texttt{rel-lambda-type-k}).\;For all $\operatorname{mnkt\_}$, $\operatorname{mnkhs\_}$ and $\operatorname{mnkc\_}$, 
\begin{equation*}
\begin{array}{@{}l@{}} \operatorname{rel\text{-}lambda\text{-}type\text{-}k}(\operatorname{mnkt\_}, \operatorname{mnkhs\_}, \operatorname{mnkc\_}) \\ \Leftrightarrow\; \operatorname{mnkc\_} \in \operatorname{seq\text{-}set} \wedge \\ \exists \operatorname{mnkg\_}, \operatorname{mnka\_}, \operatorname{mnkb\_}, \operatorname{mnkd\_}, \operatorname{mnkh\_}. \\ \quad \operatorname{mnkc\_} = {[}\operatorname{mnkg\_}, \operatorname{enc\text{-}in}(\operatorname{enc\text{-}lam}(\operatorname{mnka\_}, \operatorname{mnkb\_}), \operatorname{enc\text{-}fun}(\operatorname{mnka\_}, \operatorname{mnkd\_}))] \wedge \\ \quad \operatorname{nbg\text{-}ok}(\operatorname{mnkb\_}) \wedge \\ \quad \operatorname{mnkd\_} \in \mathit{clt} \wedge \\ \quad \operatorname{mnkh\_} \in \operatorname{mnkg\_} \wedge \\ \quad \operatorname{mnkh\_} \in \operatorname{cert\text{-}set}(\operatorname{mnkt\_}, \operatorname{mnkg\_}, \operatorname{mnkd\_}) \wedge \\ \quad \operatorname{mnkhs\_} \\ \quad =\; {[}{[}\operatorname{mnkg\_}, \operatorname{enc\text{-}all}(\operatorname{enc\text{-}imp}(\operatorname{enc\text{-}in}(\operatorname{ebvar}(0), \operatorname{mnka\_}), \operatorname{enc\text{-}in}(\operatorname{mnkb\_}, \operatorname{mnkd\_})))], {[}\operatorname{mnkg\_}, \operatorname{enc\text{-}in}(\operatorname{mnka\_}, \operatorname{enc\text{-}setc})]] \end{array}
\end{equation*}


\vnbbookfile{theorem-library/meta-nbg-laws}{theorem-library/meta-nbg-laws}
\noindent{\small meta-nbg-laws.scm -- K2, agent METACORE (2026-10-29).  Three parts: A. THE NBG CONTROLS of the predicativity relation NBG-OK (structure-library/meta-core.scm 4.3a): the walk's unfold and least-set laws; a predicative condition (TRUTH, the condition of the universal class) passes; an unconfined universal (forall X. X = ...}\par
\vnbprop{nbg-walk-unfold}{\vnbstmt{\begin{array}{@{}l@{}} \forall p. \\ \quad (\operatorname{nbg\text{-}walk}(p) \simeq \{z \in \operatorname{expr\text{-}set} : \forall v \in \operatorname{power}(\operatorname{expr\text{-}set}).\; (((p \in v) \wedge \operatorname{nbg\text{-}closed}(v)) \Rightarrow (z \in v))\}) \end{array}}}
\vnbprop{nbg-walk-least}{\vnbstmt{\begin{array}{@{}l@{}} \forall p.; \forall s \in \operatorname{power}(\operatorname{expr\text{-}set}), x \in \operatorname{nbg\text{-}walk}(p). \\ \quad ((p \in s) \wedge \operatorname{nbg\text{-}closed}(s)) \Rightarrow (x \in s) \end{array}}}
\vnbprop{nbg-walk-base}{\vnbstmt{\forall p \in \operatorname{expr\text{-}set}.\; (p \in \operatorname{nbg\text{-}walk}(p))}}
\vnbprop{all-spine-unfold}{\vnbstmt{\begin{array}{@{}l@{}} \forall b. \\ \quad (\operatorname{all\text{-}spine}(b) \simeq \{z \in (\operatorname{expr\text{-}set} \times \mathbb{N}) : \forall v \in \operatorname{power}((\operatorname{expr\text{-}set} \times \mathbb{N})).\; ((({[}b, 0] \in v) \wedge \operatorname{all\text{-}spine\text{-}closed}(v)) \Rightarrow (z \in v))\}) \end{array}}}
\vnbprop{all-spine-least}{\vnbstmt{\begin{array}{@{}l@{}} \forall b.; \forall s \in \operatorname{power}((\operatorname{expr\text{-}set} \times \mathbb{N})), x \in \operatorname{all\text{-}spine}(b). \\ \quad (({[}b, 0] \in s) \wedge \operatorname{all\text{-}spine\text{-}closed}(s)) \Rightarrow (x \in s) \end{array}}}
\vnbprop{nbg-ok-truth}{\vnbstmt{\operatorname{nbg\text{-}ok}(\operatorname{enc\text{-}truth})}}
\vnbprop{nbg-ok-refuses-unguarded}{\vnbstmt{\neg \operatorname{nbg\text{-}ok}(\operatorname{enc\text{-}all}(\operatorname{enc\text{-}eq}(\operatorname{ebvar}(0), \operatorname{ebvar}(0))))}}
\vnbprop{conj-reach-unfold}{\vnbstmt{\begin{array}{@{}l@{}} \forall x. \\ \quad (\operatorname{conj\text{-}reach}(x) \simeq \{z \in \operatorname{expr\text{-}set} : \forall v \in \operatorname{power}(\operatorname{expr\text{-}set}).\; (((x \in v) \wedge \operatorname{conj\text{-}closed}(v)) \Rightarrow (z \in v))\}) \end{array}}}
\vnbprop{ex-spine-unfold}{\vnbstmt{\begin{array}{@{}l@{}} \forall b. \\ \quad (\operatorname{ex\text{-}spine}(b) \simeq \{z \in (\operatorname{expr\text{-}set} \times \mathbb{N}) : \forall v \in \operatorname{power}((\operatorname{expr\text{-}set} \times \mathbb{N})).\; ((({[}b, 0] \in v) \wedge \operatorname{ex\text{-}spine\text{-}closed}(v)) \Rightarrow (z \in v))\}) \end{array}}}
\vnbprop{all-confined-guarded}{\vnbstmt{\begin{array}{@{}l@{}} \forall c, b. \\ \quad ((c \in \operatorname{expr\text{-}set}) \wedge (b \in \operatorname{expr\text{-}set})) \Rightarrow  \\ \quad \quad \operatorname{all\text{-}confined}(\operatorname{enc\text{-}imp}(\operatorname{enc\text{-}in}(\operatorname{ebvar}(0), c), b)) \end{array}}}
\vnbprop{ex-confined-guarded}{\vnbstmt{\begin{array}{@{}l@{}} \forall c, b. \\ \quad ((c \in \operatorname{expr\text{-}set}) \wedge (b \in \operatorname{expr\text{-}set})) \Rightarrow  \\ \quad \quad \operatorname{ex\text{-}confined}(\operatorname{enc\text{-}and}(\operatorname{enc\text{-}in}(\operatorname{ebvar}(0), c), b)) \end{array}}}
\vnbprop{adm-and-elim-k}{\vnbstmt{\begin{array}{@{}l@{}} \forall t, s, c. \\ \quad (\operatorname{rel\text{-}and\text{-}elim\text{-}k}(t, s, c) \wedge \operatorname{hyps\text{-}in}(s, \operatorname{derivable}(t))) \Rightarrow (c \in \operatorname{derivable}(t)) \end{array}}}
\vnbprop{adm-or-elim-k}{\vnbstmt{\begin{array}{@{}l@{}} \forall t, s, c. \\ \quad (\operatorname{rel\text{-}or\text{-}elim\text{-}k}(t, s, c) \wedge \operatorname{hyps\text{-}in}(s, \operatorname{derivable}(t))) \Rightarrow (c \in \operatorname{derivable}(t)) \end{array}}}
\vnbprop{adm-iff-elim-k}{\vnbstmt{\begin{array}{@{}l@{}} \forall t, s, c. \\ \quad (\operatorname{rel\text{-}iff\text{-}elim\text{-}k}(t, s, c) \wedge \operatorname{hyps\text{-}in}(s, \operatorname{derivable}(t))) \Rightarrow (c \in \operatorname{derivable}(t)) \end{array}}}
\vnbprop{ctx-fv-minus}{\vnbstmt{\begin{array}{@{}l@{}} \forall g, e, y. \\ \quad ((e \in \mathbf{Set}) \wedge (y \in \operatorname{ctx\text{-}fv}(g))) \Rightarrow  \\ \quad \quad (y \in \operatorname{ctx\text{-}fv}(\operatorname{difference}(g, \{e\}))) \vee (y \in \operatorname{expr\text{-}fv}(e)) \end{array}}}
\vnbprop{adm-forsome-elim-k}{\vnbstmt{\begin{array}{@{}l@{}} \forall t, s, c. \\ \quad (\operatorname{rel\text{-}forsome\text{-}elim\text{-}k}(t, s, c) \wedge \operatorname{hyps\text{-}in}(s, \operatorname{derivable}(t))) \Rightarrow  \\ \quad \quad c \in \operatorname{derivable}(t) \end{array}}}
\vnbprop{adm-forsome-intro-k}{\vnbstmt{\begin{array}{@{}l@{}} \forall t, s, c. \\ \quad (\operatorname{rel\text{-}forsome\text{-}intro\text{-}k}(t, s, c) \wedge \operatorname{hyps\text{-}in}(s, \operatorname{derivable}(t))) \Rightarrow  \\ \quad \quad c \in \operatorname{derivable}(t) \end{array}}}
\vnbprop{adm-forall-elim-k}{\vnbstmt{\begin{array}{@{}l@{}} \forall t, s, c. \\ \quad (\operatorname{rel\text{-}forall\text{-}elim\text{-}k}(t, s, c) \wedge \operatorname{hyps\text{-}in}(s, \operatorname{derivable}(t))) \Rightarrow  \\ \quad \quad c \in \operatorname{derivable}(t) \end{array}}}
\vnbprop{adm-reflexivity-k}{\vnbstmt{\begin{array}{@{}l@{}} \forall t, s, c. \\ \quad (\operatorname{rel\text{-}reflexivity\text{-}k}(t, s, c) \wedge \operatorname{hyps\text{-}in}(s, \operatorname{derivable}(t))) \Rightarrow  \\ \quad \quad c \in \operatorname{derivable}(t) \end{array}}}
\vnbprop{adm-if-true-k}{\vnbstmt{\begin{array}{@{}l@{}} \forall t, s, c. \\ \quad (\operatorname{rel\text{-}if\text{-}true\text{-}k}(t, s, c) \wedge \operatorname{hyps\text{-}in}(s, \operatorname{derivable}(t))) \Rightarrow (c \in \operatorname{derivable}(t)) \end{array}}}
\vnbprop{adm-if-false-k}{\vnbstmt{\begin{array}{@{}l@{}} \forall t, s, c. \\ \quad (\operatorname{rel\text{-}if\text{-}false\text{-}k}(t, s, c) \wedge \operatorname{hyps\text{-}in}(s, \operatorname{derivable}(t))) \Rightarrow (c \in \operatorname{derivable}(t)) \end{array}}}
\vnbprop{adm-iota-in-elim-k}{\vnbstmt{\begin{array}{@{}l@{}} \forall t, s, c. \\ \quad (\operatorname{rel\text{-}iota\text{-}in\text{-}elim\text{-}k}(t, s, c) \wedge \operatorname{hyps\text{-}in}(s, \operatorname{derivable}(t))) \Rightarrow  \\ \quad \quad c \in \operatorname{derivable}(t) \end{array}}}
\vnbprop{tf-set-unfold}{\vnbstmt{\begin{array}{@{}l@{}} \forall a, b. \\ \quad (\operatorname{tf\text{-}set}(a, b) \simeq \{p \in (\mathbb{N} \times \operatorname{expr\text{-}set}) : \forall w \in \operatorname{power}((\mathbb{N} \times \operatorname{expr\text{-}set})).\; (\operatorname{tf\text{-}closed}(a, b, w) \Rightarrow (p \in w))\}) \end{array}}}
\vnbprop{tf-set-least}{\vnbstmt{\begin{array}{@{}l@{}} \forall a, b.; \forall s \in \operatorname{power}((\mathbb{N} \times \operatorname{expr\text{-}set})), z \in \operatorname{tf\text{-}set}(a, b). \\ \quad \operatorname{tf\text{-}closed}(a, b, s) \Rightarrow (z \in s) \end{array}}}
\vnbprop{enc-eql-clf}{\vnbstmt{\begin{array}{@{}l@{}} \forall n \in \mathbb{N}, a, b. \\ \quad a \in \operatorname{tuples}(\operatorname{expr\text{-}set}) \wedge \\ \quad b \in \operatorname{tuples}(\operatorname{expr\text{-}set}) \wedge \\ \quad \forall u \in \operatorname{make\text{-}set}(a).\; (u \in \mathit{clt}) \wedge \\ \quad \forall v \in \operatorname{make\text{-}set}(b).\; (v \in \mathit{clt}) \wedge \\ \quad n \leq \operatorname{length}(a) \wedge \\ \quad n \leq \operatorname{length}(b) \Rightarrow \\ \quad \quad \operatorname{enc\text{-}eql}(a, b, n) \in \mathit{clf} \end{array}}}
\vnbprop{tf-eql}{\vnbstmt{\begin{array}{@{}l@{}} \forall t, g, a, b, j, f. \\ \quad g \in \operatorname{ctx\text{-}set} \wedge \\ \quad a \in \operatorname{tuples}(\operatorname{expr\text{-}set}) \wedge \\ \quad b \in \operatorname{tuples}(\operatorname{expr\text{-}set}) \wedge \\ \quad \operatorname{length}(a) = \operatorname{length}(b) \wedge \\ \quad \forall u \in \operatorname{make\text{-}set}(a).\; (u \in \mathit{clt}) \wedge \\ \quad \forall v \in \operatorname{make\text{-}set}(b).\; (v \in \mathit{clt}) \wedge \\ \quad {[}j, f] \in \operatorname{tf\text{-}set}(a, b) \Rightarrow \\ \quad \quad (j + 1) \leq \operatorname{length}(a) \wedge \\ \quad \quad {[}g, \operatorname{enc\text{-}iff}(\operatorname{enc\text{-}eql}(a, b, \operatorname{length}(a)), \operatorname{enc\text{-}and}(\operatorname{enc\text{-}eql}(a, b, j), f))] \in \operatorname{derivable}(t) \end{array}}}
\vnbprop{der-tuple-eq-decompose}{\vnbstmt{\begin{array}{@{}l@{}} \forall t, g, a, b, f. \\ \quad g \in \operatorname{ctx\text{-}set} \wedge \\ \quad \operatorname{enc\text{-}eq}(\operatorname{enc\text{-}list}(a), \operatorname{enc\text{-}list}(b)) \in \mathit{clf} \wedge \\ \quad f \in \mathit{clf} \wedge \\ \quad \operatorname{length}(a) = \operatorname{length}(b) \wedge \\ \quad {[}g, \operatorname{enc\text{-}eq}(\operatorname{enc\text{-}list}(a), \operatorname{enc\text{-}list}(a))] \in \operatorname{derivable}(t) \wedge \\ \quad {[}g, \operatorname{enc\text{-}eq}(\operatorname{enc\text{-}list}(b), \operatorname{enc\text{-}list}(b))] \in \operatorname{derivable}(t) \wedge \\ \quad ((\operatorname{length}(a) = 0) \wedge (f = \operatorname{enc\text{-}truth})) \vee ({[}0, f] \in \operatorname{tf\text{-}set}(a, b)) \Rightarrow \\ \quad \quad {[}g, \operatorname{enc\text{-}iff}(\operatorname{enc\text{-}eq}(\operatorname{enc\text{-}list}(a), \operatorname{enc\text{-}list}(b)), f)] \in \operatorname{derivable}(t) \end{array}}}
\vnbprop{j-teq-sub}{\vnbstmt{\forall t.\; (\operatorname{j\text{-}teq}(t) \subseteq \operatorname{eqv}(t))}}
\vnbprop{adm-tuple-eq-decompose-1}{\vnbstmt{\begin{array}{@{}l@{}} \forall t, s, c. \\ \quad (\operatorname{rel\text{-}tuple\text{-}eq\text{-}decompose\text{-}1}(t, s, c) \wedge \operatorname{hyps\text{-}in}(s, \operatorname{derivable}(t))) \Rightarrow  \\ \quad \quad c \in \operatorname{derivable}(t) \end{array}}}
\vnbprop{adm-tuple-eq-decompose-7}{\vnbstmt{\begin{array}{@{}l@{}} \forall t, s, c. \\ \quad (\operatorname{rel\text{-}tuple\text{-}eq\text{-}decompose\text{-}7}(t, s, c) \wedge \operatorname{hyps\text{-}in}(s, \operatorname{derivable}(t))) \Rightarrow  \\ \quad \quad c \in \operatorname{derivable}(t) \end{array}}}
\vnbprop{def-derivable-unfold}{\vnbstmt{\begin{array}{@{}l@{}} \forall t. \\ \quad (\operatorname{def\text{-}derivable}(t) \simeq \{q \in \operatorname{seq\text{-}set} : \forall s \in \operatorname{power}(\operatorname{seq\text{-}set}).\; (\operatorname{def\text{-}closed}(t, s) \Rightarrow (q \in s))\}) \end{array}}}
\vnbprop{derivable-in-power}{\vnbstmt{\forall t.\; (\operatorname{derivable}(t) \in \operatorname{power}(\operatorname{seq\text{-}set}))}}
\vnbprop{def-closed-derivable}{\vnbstmt{\forall t.\; \operatorname{def\text{-}closed}(t, \operatorname{derivable}(t))}}
\vnbprop{def-derivable-sub}{\vnbstmt{\forall t.; \forall s \in \operatorname{def\text{-}derivable}(t).\; (s \in \operatorname{derivable}(t))}}
\vnbprop{adm-forall-elim-kd}{\vnbstmt{\begin{array}{@{}l@{}} \forall t, s, c. \\ \quad (\operatorname{rel\text{-}forall\text{-}elim\text{-}kd}(t, s, c) \wedge \operatorname{hyps\text{-}in}(s, \operatorname{derivable}(t))) \Rightarrow  \\ \quad \quad c \in \operatorname{derivable}(t) \end{array}}}
\vnbprop{adm-forsome-intro-kd}{\vnbstmt{\begin{array}{@{}l@{}} \forall t, s, c. \\ \quad (\operatorname{rel\text{-}forsome\text{-}intro\text{-}kd}(t, s, c) \wedge \operatorname{hyps\text{-}in}(s, \operatorname{derivable}(t))) \Rightarrow  \\ \quad \quad c \in \operatorname{derivable}(t) \end{array}}}
\vnbprop{adm-reflexivity-kd}{\vnbstmt{\begin{array}{@{}l@{}} \forall t, s, c. \\ \quad (\operatorname{rel\text{-}reflexivity\text{-}kd}(t, s, c) \wedge \operatorname{hyps\text{-}in}(s, \operatorname{derivable}(t))) \Rightarrow  \\ \quad \quad c \in \operatorname{derivable}(t) \end{array}}}
\vnbprop{adm-reflexivity-kt}{\vnbstmt{\begin{array}{@{}l@{}} \forall t, s, c. \\ \quad (\operatorname{rel\text{-}reflexivity\text{-}kt}(t, s, c) \wedge \operatorname{hyps\text{-}in}(s, \operatorname{derivable}(t))) \Rightarrow  \\ \quad \quad c \in \operatorname{derivable}(t) \end{array}}}
\vnbprop{adm-if-true-kd}{\vnbstmt{\begin{array}{@{}l@{}} \forall t, s, c. \\ \quad (\operatorname{rel\text{-}if\text{-}true\text{-}kd}(t, s, c) \wedge \operatorname{hyps\text{-}in}(s, \operatorname{derivable}(t))) \Rightarrow (c \in \operatorname{derivable}(t)) \end{array}}}
\vnbprop{adm-if-false-kd}{\vnbstmt{\begin{array}{@{}l@{}} \forall t, s, c. \\ \quad (\operatorname{rel\text{-}if\text{-}false\text{-}kd}(t, s, c) \wedge \operatorname{hyps\text{-}in}(s, \operatorname{derivable}(t))) \Rightarrow (c \in \operatorname{derivable}(t)) \end{array}}}
\vnbprop{adm-adds-nothing-k}{\vnbstmt{\begin{array}{@{}l@{}} \forall t, s, c. \\ \quad (\operatorname{rel\text{-}adds\text{-}nothing\text{-}k}(t, s, c) \wedge \operatorname{hyps\text{-}in}(s, \operatorname{derivable}(t))) \Rightarrow  \\ \quad \quad c \in \operatorname{derivable}(t) \end{array}}}
\vnbprop{rewrs-unfold}{\vnbstmt{\begin{array}{@{}l@{}} \forall j. \\ \quad (\operatorname{rewrs}(j) \simeq \{x \in \operatorname{rw\text{-}univ} : \forall s \in \operatorname{power}(\operatorname{rw\text{-}univ}).\; (\operatorname{rws\text{-}closed}(j, s) \Rightarrow (x \in s))\}) \end{array}}}
\vnbprop{rewrs-least}{\vnbstmt{\forall j.; \forall s \in \operatorname{power}(\operatorname{rw\text{-}univ}), x \in \operatorname{rewrs}(j).\; (\operatorname{rws\text{-}closed}(j, s) \Rightarrow (x \in s))}}
\vnbprop{eqv-rws-closed}{\vnbstmt{\forall t, j.\; ((j \subseteq \operatorname{eqv}(t)) \Rightarrow \operatorname{rws\text{-}closed}(j, \operatorname{eqv}(t)))}}
\vnbprop{crs-s}{\vnbstmt{\forall t, j, x.\; (((j \subseteq \operatorname{eqv}(t)) \wedge (x \in \operatorname{rewrs}(j))) \Rightarrow (x \in \operatorname{eqv}(t)))}}
\vnbprop{rewrs-base}{\vnbstmt{\forall j.; \forall x \in \operatorname{rewr}(j, 7).\; (x \in \operatorname{rewrs}(j))}}
\vnbprop{rewrs-refl}{\vnbstmt{\forall j, g, e.\; (((g \in \operatorname{ctx\text{-}set}) \wedge (e \in \mathit{clf})) \Rightarrow ({[}2, g, e, e] \in \operatorname{rewrs}(j)))}}
\vnbprop{rewrs-trans}{\vnbstmt{\begin{array}{@{}l@{}} \forall j, g, a, b, c. \\ \quad g \in \operatorname{ctx\text{-}set} \wedge \\ \quad a \in \operatorname{expr\text{-}set} \wedge \\ \quad b \in \operatorname{expr\text{-}set} \wedge \\ \quad c \in \operatorname{expr\text{-}set} \wedge \\ \quad {[}2, g, a, b] \in \operatorname{rewrs}(j) \wedge \\ \quad {[}2, g, b, c] \in \operatorname{rewrs}(j) \Rightarrow \\ \quad \quad {[}2, g, a, c] \in \operatorname{rewrs}(j) \end{array}}}
\vnbprop{adm-rewrite-s}{\vnbstmt{\begin{array}{@{}l@{}} \forall t, j, s, c. \\ \quad ((j \subseteq \operatorname{eqv}(t)) \wedge (\operatorname{rel\text{-}rewrite\text{-}s}(j, s, c) \wedge \operatorname{hyps\text{-}in}(s, \operatorname{derivable}(t)))) \Rightarrow  \\ \quad \quad c \in \operatorname{derivable}(t) \end{array}}}
\vnbprop{adm-rewrite-hyp-s}{\vnbstmt{\begin{array}{@{}l@{}} \forall t, j, s, c. \\ \quad ((j \subseteq \operatorname{eqv}(t)) \wedge (\operatorname{rel\text{-}rewrite\text{-}hyp\text{-}s}(j, s, c) \wedge \operatorname{hyps\text{-}in}(s, \operatorname{derivable}(t)))) \Rightarrow  \\ \quad \quad c \in \operatorname{derivable}(t) \end{array}}}
\vnbprop{adm-nth-reduce-s}{\vnbstmt{\begin{array}{@{}l@{}} \forall t, s, c. \\ \quad (\operatorname{rel\text{-}nth\text{-}reduce\text{-}s}(t, s, c) \wedge \operatorname{hyps\text{-}in}(s, \operatorname{derivable}(t))) \Rightarrow (c \in \operatorname{derivable}(t)) \end{array}}}
\vnbprop{adm-length-reduce-s}{\vnbstmt{\begin{array}{@{}l@{}} \forall t, s, c. \\ \quad (\operatorname{rel\text{-}length\text{-}reduce\text{-}s}(t, s, c) \wedge \operatorname{hyps\text{-}in}(s, \operatorname{derivable}(t))) \Rightarrow  \\ \quad \quad c \in \operatorname{derivable}(t) \end{array}}}
\vnbprop{adm-functoid-beta-s}{\vnbstmt{\begin{array}{@{}l@{}} \forall t, s, c. \\ \quad (\operatorname{rel\text{-}functoid\text{-}beta\text{-}s}(t, s, c) \wedge \operatorname{hyps\text{-}in}(s, \operatorname{derivable}(t))) \Rightarrow  \\ \quad \quad c \in \operatorname{derivable}(t) \end{array}}}
\vnbprop{adm-lambda-beta-s}{\vnbstmt{\begin{array}{@{}l@{}} \forall t, s, c. \\ \quad (\operatorname{rel\text{-}lambda\text{-}beta\text{-}s}(t, s, c) \wedge \operatorname{hyps\text{-}in}(s, \operatorname{derivable}(t))) \Rightarrow  \\ \quad \quad c \in \operatorname{derivable}(t) \end{array}}}
\vnbprop{adm-lambda-beta-hyp-s}{\vnbstmt{\begin{array}{@{}l@{}} \forall t, s, c. \\ \quad (\operatorname{rel\text{-}lambda\text{-}beta\text{-}hyp\text{-}s}(t, s, c) \wedge \operatorname{hyps\text{-}in}(s, \operatorname{derivable}(t))) \Rightarrow  \\ \quad \quad c \in \operatorname{derivable}(t) \end{array}}}
\vnbprop{adm-macete-s}{\vnbstmt{\begin{array}{@{}l@{}} \forall t, s, c. \\ \quad (\operatorname{rel\text{-}macete\text{-}s}(t, s, c) \wedge \operatorname{hyps\text{-}in}(s, \operatorname{derivable}(t))) \Rightarrow (c \in \operatorname{derivable}(t)) \end{array}}}
\vnbprop{adm-macete-hyp-s}{\vnbstmt{\begin{array}{@{}l@{}} \forall t, s, c. \\ \quad (\operatorname{rel\text{-}macete\text{-}hyp\text{-}s}(t, s, c) \wedge \operatorname{hyps\text{-}in}(s, \operatorname{derivable}(t))) \Rightarrow (c \in \operatorname{derivable}(t)) \end{array}}}
\vnbprop{adm-macete-n-s}{\vnbstmt{\begin{array}{@{}l@{}} \forall t, s, c. \\ \quad (\operatorname{rel\text{-}macete\text{-}n\text{-}s}(t, s, c) \wedge \operatorname{hyps\text{-}in}(s, \operatorname{derivable}(t))) \Rightarrow (c \in \operatorname{derivable}(t)) \end{array}}}
\vnbprop{adm-macete-n-hyp-s}{\vnbstmt{\begin{array}{@{}l@{}} \forall t, s, c. \\ \quad (\operatorname{rel\text{-}macete\text{-}n\text{-}hyp\text{-}s}(t, s, c) \wedge \operatorname{hyps\text{-}in}(s, \operatorname{derivable}(t))) \Rightarrow  \\ \quad \quad c \in \operatorname{derivable}(t) \end{array}}}
\vnbprop{adm-union-decompose-s}{\vnbstmt{\begin{array}{@{}l@{}} \forall t, s, c. \\ \quad (\operatorname{rel\text{-}union\text{-}decompose\text{-}s}(t, s, c) \wedge \operatorname{hyps\text{-}in}(s, \operatorname{derivable}(t))) \Rightarrow  \\ \quad \quad c \in \operatorname{derivable}(t) \end{array}}}
\vnbprop{adm-union-decompose-hyp-s}{\vnbstmt{\begin{array}{@{}l@{}} \forall t, s, c. \\ \quad (\operatorname{rel\text{-}union\text{-}decompose\text{-}hyp\text{-}s}(t, s, c) \wedge \operatorname{hyps\text{-}in}(s, \operatorname{derivable}(t))) \Rightarrow  \\ \quad \quad c \in \operatorname{derivable}(t) \end{array}}}
\vnbprop{adm-intersection-decompose-s}{\vnbstmt{\begin{array}{@{}l@{}} \forall t, s, c. \\ \quad (\operatorname{rel\text{-}intersection\text{-}decompose\text{-}s}(t, s, c) \wedge \operatorname{hyps\text{-}in}(s, \operatorname{derivable}(t))) \Rightarrow  \\ \quad \quad c \in \operatorname{derivable}(t) \end{array}}}
\vnbprop{adm-intersection-decompose-hyp-s}{\vnbstmt{\begin{array}{@{}l@{}} \forall t, s, c. \\ \quad (\operatorname{rel\text{-}intersection\text{-}decompose\text{-}hyp\text{-}s}(t, s, c) \wedge \operatorname{hyps\text{-}in}(s, \operatorname{derivable}(t))) \Rightarrow  \\ \quad \quad c \in \operatorname{derivable}(t) \end{array}}}
\vnbprop{adm-tuple-eq-decompose-s}{\vnbstmt{\begin{array}{@{}l@{}} \forall t, s, c. \\ \quad (\operatorname{rel\text{-}tuple\text{-}eq\text{-}decompose\text{-}s}(t, s, c) \wedge \operatorname{hyps\text{-}in}(s, \operatorname{derivable}(t))) \Rightarrow  \\ \quad \quad c \in \operatorname{derivable}(t) \end{array}}}
\vnbprop{adm-cartesian-decompose-s}{\vnbstmt{\begin{array}{@{}l@{}} \forall t, s, c. \\ \quad ((t \subseteq \mathit{clf}) \wedge ((\operatorname{ax\text{-}mem\text{-}set} \in t) \wedge (\operatorname{rel\text{-}cartesian\text{-}decompose\text{-}s}(t, s, c) \wedge \operatorname{hyps\text{-}in}(s, \operatorname{derivable}(t))))) \Rightarrow  \\ \quad \quad c \in \operatorname{derivable}(t) \end{array}}}
\vnbprop{adm-tuples-intro-0k}{\vnbstmt{\begin{array}{@{}l@{}} \forall t, s, c. \\ \quad (\operatorname{rel\text{-}tuples\text{-}intro\text{-}0k}(t, s, c) \wedge \operatorname{hyps\text{-}in}(s, \operatorname{derivable}(t))) \Rightarrow  \\ \quad \quad c \in \operatorname{derivable}(t) \end{array}}}
\vnbprop{adm-lambda-type-k}{\vnbstmt{\begin{array}{@{}l@{}} \forall t, s, c. \\ \quad (\operatorname{rel\text{-}lambda\text{-}type\text{-}k}(t, s, c) \wedge \operatorname{hyps\text{-}in}(s, \operatorname{derivable}(t))) \Rightarrow  \\ \quad \quad c \in \operatorname{derivable}(t) \end{array}}}

